\input{preamble.tex}
% Override color setting
\definecolor{bookheadingcolor}{HTML}{9B3A3F} % 9B3A3F - Deep red

\title{Η ΑΓΙΑ ΓΡΑΦΗ}
\author{}
\date{}

\begin{document}
\begin{spacing}{1.1}
\maketitle

\cleardoublepage
\begin{titlepage}
  \begin{center}
    \textcolor{bookheadingcolor}{\Huge Preface}\par
  \end{center}
  \vspace{2em}
  
  This project was undertaken in love and respect for the Bible, with a desire to have an accessible and
  beautiful Greek Bible available in print to anyone who would like one. While there are many great Greek
  New Testaments available in print, the Septuagint has been less accessible, particularly in a format that
  is both compact and minimalist. Most of the Septuagints available in print are quite large. Additionally,
  there are almost no complete Greek Bibles available to purchase for a reasonable price. I have undertaken
  this project for those out there who, like me, want to take a physical Greek Bible along 
  with them where ever they want to go.

  When I started this project, I went hunting for open and public domain editions of the Septuagint and the NT
  that I could use as the texts for this Bible. While there are several great options out there, I settled on
  the Brenton Septuagint and the OpenGNT new testament. The reason for choosing Brenton's Septuagint was simple:
  I found a great open source project that had already digitized the text and prepared it for print: 
  https://github.com/mrgreekgeek/Brenton-LXX-Latex-print-project/. Starting with this baseline, I was able to
  style the text in a way that I liked. Then I had to find and prepare a NT Text.

  For the NT I chose the Open GNT (https://opengnt.com/) which was prepared by Eliran Wong and released under
  the Creative Commons Attribution 4.0 International License (CC BY 4.0). This project was created "to offer a 
  FREE NA-equivalent text of Greek New Testament, compiled from open-resources" and provided access to the text
  in a format that I could adapt to my needs.

  As for formatting, I was inspired by some of the beautiful minimalist reader Bibles available in English. As
  much as possible, I wanted to keep the text front and center, eliminating distractions and unnecessary elements.
  I have tried to mitigate the distraction from things like section headings, spacing between chapters, and even
  chapter numbers to some degree. I ultimately decided to leave verse numbers in place, because I think navigating
  the Old Testament may have been more difficult without them; however, I tried to minimize their visual impact.
  My goal is to facilitate a novel-like reading experience, free of distractions.

  The source code that I have used to extract, process, and format the texts used for this Bible is available 
  free of charge at https://github.com/jjorloff1/lxx-nt-greek-bible-builder.

  I hope that this Greek Bible will serve you well as you study and meditate on the Scriptures.
  Glory to God!

  \vfill
  \begin{flushright}
    {\large\textit{Jesse Orloff}\par}
    {\large www.jesseorloff.com\par}
    {\large August 2025\par}
  \end{flushright}  
\end{titlepage}

\cleardoublepage
\pagestyle{empty}
\begingroup
\centering
{\huge \textcolor{bookheadingcolor}{Table of Contents} \par}
\endgroup

\begin{multicols}{2}
\makeatletter
\renewcommand{\tableofcontents}{\@starttoc{toc}}
\makeatother
\tableofcontents
\end{multicols}
\pagestyle{fancy}

\cleardoublepage
\thispagestyle{empty}
\vspace*{3cm}
\phantomsection
\addcontentsline{toc}{part}{Η ΠΑΛΑΙΑ ΔΙΑΘΗΚΗ}
\begin{center}
  {\Huge Η ΠΑΛΑΙΑ ΔΙΑΘΗΚΗ}\\[2em]
  {\large Ἡ μετάφρασις τῶν Ἑβδομήκοντα}
\end{center}
\newpage
\thispagestyle{empty}
\null

\def\book{ΓΕΝΕΣΙΣ}
\biblebook{ΓΕΝΕΣΙΣ}


\lettrine[lines=2, loversize=0.2, nindent=0em, findent=.25em]{\textcolor{bookheadingcolor}{Ἐ}}{Ν} ἀρχῇ ἐποίησεν ὁ Θεὸς τὸν οὐρανὸν καὶ τὴν γῆν.
\vs{2}Ἡ δὲ γῆ ἦν ἀόρατος καὶ ἀκατασκεύαστος, καὶ σκότος ἐπάνω τῆς ἀβύσσου· καὶ πνεῦμα Θεοῦ ἐπεφέρετο ἐπάνω τοῦ ὕδατος.
\vs{3}Καὶ εἶπεν ὁ Θεὸς, γενηθήτω φῶς· καὶ ἐγένετο φῶς.
\vs{4}Καὶ εἶδεν ὁ Θεὸς τὸ φῶς, ὅτι καλόν· καὶ διεχώρισεν ὁ Θεὸς ἀνὰ μέσον τοῦ φωτὸς, καὶ ἀνὰ μέσον τοῦ σκότους.
\vs{5}Καὶ ἐκάλεσεν ὁ Θεὸς τὸ φῶς ἡμέραν, καὶ τὸ σκότος ἐκάλεσε νύκτα. Καὶ ἐγένετο ἑσπέρα, καὶ ἐγένετο πρωῒ, ἡμέρα μία.

\vs{6}Καὶ εἶπεν ὁ Θεὸς, γενηθήτω στερέωμα ἐν μέσῳ τοῦ ὕδατος· καὶ ἔστω διαχωρίζον ἀνὰ μέσον ὕδατος καὶ ὕδατος· καὶ ἐγένετο οὕτως.
\vs{7}Καὶ ἐποίησεν ὁ Θεὸς τὸ στερέωμα· καὶ διεχώρισεν ὁ Θεὸς ἀνὰ μέσον τοῦ ὕδατος, ὃ ἦν ὑποκάτω τοῦ στερεώματος, καὶ ἀνὰ μέσον τοῦ ὕδατος, τοῦ ἐπάνω τοῦ στερεώματος.
\vs{8}Καὶ ἐκάλεσεν ὁ Θεὸς τὸ στερέωμα οὐρανόν· καὶ εἶδεν ὁ Θεὸς ὅτι καλόν· καὶ ἐγένετο ἑσπέρα, καὶ ἐγένετο πρωῒ, ἡμέρα δευτέρα.

\vs{9}Καὶ εἶπεν ὁ Θεὸς, συναχθήτω τὸ ὕδωρ τὸ ὑποκάτω τοῦ οὐρανοῦ εἰς συναγωγὴν μίαν, καὶ ὀφθήτω ἡ ξηρά· καὶ ἐγένετο οὕτως· καὶ συνήχθη τὸ ὕδωρ τὸ ὑποκάτω τοῦ οὐρανοῦ εἰς τὰς συναγωγὰς αὐτῶν, καὶ ὤφθη ἡ ξηρά.
\vs{10}Καὶ ἐκάλεσεν ὁ Θεὸς τὴν ξηρὰν, γῆν· καὶ τὰ συστήματα τῶν ὑδάτων ἐκάλεσε θαλάσσας· καὶ εἶδεν ὁ Θεὸς ὅτι καλόν.
\vs{11}Καὶ εἶπεν ὁ Θεὸς, βλαστησάτω ἡ γῆ βοτάνην χόρτου, σπεῖρον σπέρμα κατὰ γένος καὶ καθʼ ὁμοιότητα, καὶ ξύλον κάρπιμον ποιοῦν καρπὸν, οὗ τὸ σπέρμα αὐτοῦ ἐν αὐτῷ κατὰ γένος ἐπὶ τῆς γῆς· καὶ ἐγένετο οὕτως.
\vs{12}Καὶ ἐξήνεγκεν ἡ γῆ βοτάνην χόρτου, σπεῖρον σπέρμα κατὰ γένος καὶ καθʼ ὁμοιότητα, καὶ ξύλον κάρπιμον ποιοῦν καρπὸν, οὗ τὸ σπέρμα αὐτοῦ ἐν αὐτῷ κατὰ γένος ἐπὶ τῆς γῆς· καὶ εἶδεν ὁ Θεὸς ὅτι καλόν.
\vs{13}Καὶ ἐγένετο ἑσπέρα, καὶ ἐγένετο πρωῒ, ἡμέρα τρίτη.

\vs{14}Καὶ εἶπεν ὁ Θεὸς, γενηθήτωσαν φωστῆρες ἐν τῷ στερεώματι τοῦ οὐρανοῦ εἰς φαῦσιν ἐπὶ τῆς γῆς, τοῦ διαχωρίζειν ἀνὰ μέσον τῆς ἡμέρας καὶ ἀνὰ μέσον τῆς νυκτός· καὶ ἔστωσαν εἰς σημεῖα, καὶ εἰς καιροὺς, καὶ εἰς ἡμέρας, καὶ εἰς ἐνιαυτούς.
\vs{15}Καὶ ἔστωσαν εἰς φαῦσιν ἐν τῷ στερεώματι τοῦ οὐρανοῦ, ὥστε φαίνειν ἐπὶ τῆς γῆς· καὶ ἐγένετο οὕτως.
\vs{16}Καὶ ἐποίησεν ὁ Θεὸς τοὺς δύο φωστῆρας τοὺς μεγάλους· τὸν φωστῆρα τὸν μέγαν εἰς ἀρχὰς τῆς ἡμέρας, καὶ τὸν φωστῆρα τὸν ἐλάσσω εἰς ἀρχὰς τῆς νυκτὸς, καὶ τοὺς ἀστέρας.
\vs{17}Καὶ ἔθετο αὐτοὺς ὁ Θεὸς ἐν τῷ στερεώματι τοῦ οὐρανοῦ, ὥστε φαίνειν ἐπὶ τῆς γῆς,
\vs{18}καὶ ἄρχειν τῆς ἡμέρας καὶ τῆς νυκτὸς, καὶ διαχωρίζειν ἀνὰ μέσον τοῦ φωτὸς, καὶ ἀνὰ μέσον τοῦ σκότους· καὶ εἶδεν ὁ Θεὸς ὅτι καλόν.
\vs{19}Καὶ ἐγένετο ἑσπέρα καὶ ἐγένετο πρωῒ, ἡμέρα τετάρτη.

\vs{20}Καὶ εἶπεν ὁ Θεὸς, ἐξαγαγέτω τὰ ὕδατα ἑρπετὰ ψυχῶν ζωσῶν, καὶ πετεινὰ πετόμενα ἐπὶ τῆς γῆς κατὰ τὸ στερέωμα τοῦ οὐρανοῦ· καὶ ἐγένετο οὕτως.
\vs{21}Καὶ ἐποίησεν ὁ Θεὸς τὰ κήτη τὰ μεγάλα, καὶ πᾶσαν ψυχὴν ζώων ἑρπετῶν, ἃ ἐξήγαγε τὰ ὕδατα κατὰ γένη αὐτῶν, καὶ πᾶν πετεινὸν πτερωτὸν κατὰ γένος· καὶ εἶδεν ὁ Θεὸς ὅτι καλά.
\vs{22}Καὶ εὐλόγησεν αὐτὰ ὁ Θεὸς, λέγων, αὐξάνεσθε καὶ πληθύνεσθε, καὶ πληρώσατε τὰ ὕδατα ἐν ταῖς θαλάσσαις, καὶ τὰ πετεινὰ πληθυνέσθωσαν ἐπὶ τῆς γῆς.
\vs{23}Καὶ ἐγένετο ἑσπέρα, καὶ ἐγένετο πρωῒ, ἡμέρα πέμπτη.

\vs{24}Καὶ εἶπεν ὁ Θεὸς, ἐξαγαγέτω ἡ γῆ ψυχὴν ζῶσαν κατὰ γένος, τετράποδα, καὶ ἑρπετὰ, καὶ θηρία τῆς γῆς κατὰ γένος· καὶ ἐγένετο οὕτως.
\vs{25}Καὶ ἐποίησεν ὁ Θεὸς τὰ θηρία τῆς γῆς κατὰ γένος, καὶ τὰ κτήνη κατὰ γένος αὐτῶν, καὶ πάντα τὰ ἑρπετὰ τῆς γῆς κατὰ γένος· καὶ εἶδεν ὁ Θεὸς ὅτι καλά.

\vs{26}Καὶ εἶπεν ὁ Θεός, Ποιήσωμεν ἄνθρωπον κατʼ εἰκόνα ἡμετέραν καὶ καθʼ ὁμοίωσιν· καὶ ἀρχέτωσαν τῶν ἰχθύων τῆς θαλάσσης, καὶ τῶν πετεινῶν τοῦ οὐρανοῦ, καὶ τῶν κτηνῶν, καὶ πάσης τῆς γῆς, καὶ πάντων τῶν ἑρπετῶν τῶν ἑρπόντων ἐπὶ τῆς γῆς.
\vs{27}Καὶ ἐποιήσεν ὁ Θεὸς τὸν ἄνθρωπον· κατʼ εἰκόνα Θεοῦ ἐποίησεν αὐτόν· ἄρσεν καὶ θῆλυ ἐποίησεν αὐτούς.
\vs{28}Καὶ εὐλόγησεν αὐτοὺς ὁ Θεὸς, λέγων, αὐξάνεσθε καὶ πληθύνεσθε, καὶ πληρώσατε τὴν γῆν, καὶ κατακυριεύσατε αὐτῆς· καὶ ἄρχετε τῶν ἰχθύων τῆς θαλάσσης, καὶ τῶν πετεινῶν τοῦ οὐρανοῦ, καὶ πάντων τῶν κτηνῶν, καὶ πάσης τῆς γῆς, καὶ πάντων τῶν ἑρπετῶν τῶν ἑρπόντων ἐπὶ τῆς γῆς.
\vs{29}Καὶ εἶπεν ὁ Θεός, Ἰδοὺ δέδωκα ὑμῖν πάντα χόρτον σπόριμον σπεῖρον σπέρμα, ὅ ἐστιν ἐπάνω πάσης τῆς γῆς· καὶ πᾶν ξύλον, ὃ ἔχει ἐν ἑαυτῷ καρπὸν σπέρματος σπορίμου, ὑμῖν ἔσται εἰς βρῶσιν,
\vs{30}καὶ πᾶσι τοῖς θηρίοις τῆς γῆς, καὶ πᾶσι τοῖς πετεινοῖς τοῦ οὐρανοῦ, καὶ παντὶ ἑρπετῷ ἕρποντι ἐπὶ τῆς γῆς, ὃ ἔχει ἐν ἑαυτῷ ψυχὴν ζωῆς, καὶ πάντα χόρτον χλωρὸν εἰς βρῶσιν· καὶ ἐγένετο οὕτως.
\vs{31}Καὶ εἶδεν ὁ Θεὸς τὰ πάντα, ὅσα ἐποίησε, καὶ ἰδοὺ καλὰ λίαν· καὶ ἐγένετο ἑσπέρα, καὶ ἐγένετο πρωῒ, ἡμέρα ἕκτη.

\ch{2}Καὶ συνετελέσθησαν ὁ οὐρανὸς καὶ ἡ γῆ, καὶ πᾶς ὁ κόσμος αὐτῶν.

\vs{2}Καὶ συνετέλεσεν ὁ Θεὸς ἐν τῇ ἡμέρᾳ τῇ ἕκτῃ τὰ ἔργα αὐτοῦ, ἃ ἐποίησε· καὶ κατέπαυσε τῇ ἡμέρᾳ τῇ ἑβδόμῃ ἀπὸ πάντων τῶν ἔργων αὐτοῦ, ὧν ἐποίησε.
\vs{3}Καὶ εὐλόγησεν ὁ Θεὸς τὴν ἡμέραν τὴν ἑβδόμην, καὶ ἡγίασεν αὐτήν, ὅτι ἐν αὐτῇ κατέπαυσεν ἀπὸ πάντων τῶν ἔργων αὐτοῦ, ὧν ἤρξατο ὁ Θεὸς ποιῆσαι.

\vs{4}Αὕτη ἡ βίβλος γενέσεως οὐρανοῦ καὶ γῆς, ὅτε ἐγένετο, ᾗ ἡμέρᾳ ἐποίησε Κύριος ὁ Θεὸς τὸν οὐρανὸν καὶ τὴν γῆν,
\vs{5}καὶ πᾶν χλωρὸν ἀγροῦ πρὸ τοῦ γενέσθαι ἐπὶ τῆς γῆς, καὶ πάντα χόρτον ἀγροῦ πρὸ τοῦ ἀνατεῖλαι· οὐ γὰρ ἔβρεξεν ὁ Θεὸς ἐπὶ τὴν γῆν, καὶ ἄνθρωπος οὐκ ἦν ἐργάζεσθαι αὐτήν.
\vs{6}Πηγὴ δὲ ἀνέβαινεν ἐκ τῆς γῆς, καὶ ἐπότιζε πᾶν τὸ πρόσωπον τῆς γῆς.
\vs{7}Καὶ ἔπλασεν ὁ Θεὸς τὸν ἄνθρωπον, χοῦν ἀπὸ τῆς γῆς· καὶ ἐνεφύσησεν εἰς τὸ πρόσωπον αὐτοῦ πνοὴν ζωῆς, καὶ ἐγένετο ὁ ἄνθρωπος εἰς ψυχὴν ζῶσαν.

\vs{8}Καὶ ἐφύτευσεν ὁ Θεὸς παράδεισον ἐν Ἐδὲμ κατὰ ἀνατολάς· καὶ ἔθετο ἐκεῖ τὸν ἄνθρωπον, ὃν ἔπλασε.
\vs{9}Καὶ ἐξανέτειλεν ὁ Θεὸς ἔτι ἐκ τῆς γῆς πᾶν ξύλον ὡραῖον εἰς ὅρασιν, καὶ καλὸν εἰς βρῶσιν, καὶ τὸ ξύλον τῆς ζωῆς ἐν μέσῳ τοῦ παραδείσου, καὶ τὸ ξύλον τοῦ εἰδέναι γνωστὸν καλοῦ καὶ πονηροῦ.
\vs{10}Ποταμὸς δὲ ἐκπορεύεται ἐξ Ἐδὲμ ποτίζειν τὸν παράδεισον· ἐκεῖθεν ἀφορίζεται εἰς τέσσαρας ἀρχάς.
\vs{11}Ὄνομα τῷ ἑνὶ, Φισῶν· οὗτος ὁ κυκλῶν πᾶσαν τὴν γῆν Εὐιλάτ· ἐκεῖ οὗ ἐστι τὸ χρυσίον.
\vs{12}Τὸ δὲ χρυσίον τῆς γῆς ἐκείνης καλόν· καὶ ἐκεῖ ἐστιν ὁ ἄνθραξ, καὶ ὁ λίθος ὁ πράσινος.
\vs{13}Καὶ ὄνομα τῷ ποταμῷ τῷ δευτέρῳ, Γεῶν· οὗτος ὁ κυκλῶν πᾶσαν τὴν γὴν Αἰθιοπίας.
\vs{14}Καὶ ὁ ποταμὸς ὁ τρίτος, Τίγρις· οὗτος ὁ προπορευόμενος κατέναντι Ἀσσυρίων· ὁ δὲ ποταμὸς ὁ τέταρτος, Εὐφράτης.
\vs{15}Καὶ ἔλαβε Κύριος ὁ Θεὸς τὸν ἄνθρωπον ὃν ἔπλασε, καὶ ἔθετο αὐτὸν ἐν τῷ παραδείσῳ τῆς τρυφῆς, ἐργάζεσθαι αὐτὸν καὶ φυλάσσειν.
\vs{16}Καὶ ἐνετείλατο Κύριος ὁ Θεὸς τῷ Ἀδὰμ, λέγων, ἀπὸ παντὸς ξύλου τοῦ ἐν τῷ παραδείσῳ βρώσει φαγῇ.
\vs{17}Ἀπὸ δὲ τοῦ ξύλου τοῦ γινώσκειν καλὸν καὶ πονηρὸν, οὐ φάγεσθε ἀπʼ αὐτοῦ· ᾗ δʼ ἂν ἡμέρᾳ φάγητε ἀπʼ αὐτοῦ, θανάτῳ ἀποθανεῖσθε.

\vs{18}Καὶ εἶπε Κύριος ὁ Θεὸς, οὐ καλὸν εἶναι τὸν ἄνθρωπον μόνον· ποιήσωμεν αὐτῷ βοηθὸν κατʼ αὐτόν.
\vs{19}Καὶ ἔπλασεν ὁ Θεὸς ἔτι ἐκ τῆς γῆς πάντα τὰ θηρία τοῦ ἀγροῦ, καὶ πάντα τὰ πετεινὰ τοῦ οὐρανοῦ· καὶ ἤγαγεν αὐτὰ πρὸς τὸν Ἀδὰμ, ἰδεῖν τί καλέσει αὐτά· καὶ πᾶν ὃ ἐὰν ἐκάλεσεν αὐτὸ Ἀδὰμ ψυχὴν ζῶσαν, τοῦτο ὄνομα αὐτῷ.
\vs{20}Καὶ ἐκάλεσεν Ἀδὰμ ὀνόματα πᾶσι τοῖς κτήνεσι, καὶ πᾶσι τοῖς πετεινοῖς τοῦ οὐρανοῦ, καὶ πᾶσι τοῖς θηρίοις τοῦ ἀγροῦ· τῷ δὲ Ἀδὰμ οὐχ εὑρέθη βοηθὸς ὅμοιος αὐτῷ.
\vs{21}Καὶ ἐπέβαλεν ὁ Θεὸς ἔκστασιν ἐπὶ τὸν Ἀδὰμ, καὶ ὕπνωσε· καὶ ἔλαβε μίαν τῶν πλευρῶν αὐτοῦ, καὶ ἀνεπλήρωσε σάρκα ἀντʼ αὐτῆς.
\vs{22}Καὶ ᾠκοδόμησεν ὁ Θεὸς τὴν πλευρὰν, ἣν ἔλαβεν ἀπὸ τοῦ Ἀδὰμ εἰς γυναῖκα· καὶ ἤγαγεν αὐτὴν πρὸς τὸν Ἀδάμ.
\vs{23}Καὶ εἶπεν Ἀδάμ· τοῦτο νῦν ὀστοῦν ἐκ τῶν ὀστέων μου, καὶ σὰρξ ἐκ τῆς σαρκός μου· αὕτη κληθήσεται γυνὴ, ὅτι ἐκ τοῦ ἀνδρὸς αὐτῆς ἐλήφθη.
\vs{24}Ἕνεκεν τούτου καταλείψει ἄνθρωπος τὸν πατέρα αὐτοῦ καὶ τὴν μητέρα, καὶ προσκολληθήσεται πρὸς τὴν γυναῖκα αὐτοῦ· καὶ ἔσονται οἱ δύο εἰς σάρκα μίαν.
\vs{25}Καὶ ἦσαν οἱ δύο γυμνοὶ, ὅ, τε Ἀδὰμ καὶ ἡ γυνὴ αὐτοῦ, καὶ οὐκ ᾐσχύνοντο.

\ch{3}
Ὁ δὲ ὄφις ἦν φρονιμώτατος πάντων τῶν θηρίων τῶν ἐπὶ τῆς γῆς, ὧν ἐποίησε Κύριος ὁ Θεός· καὶ εἶπεν ὁ ὄφις τῇ γυναικὶ, τί ὅτι εἶπεν ὁ Θεός, οὐ μὴ φάγητε ἀπὸ παντὸς ξύλου τοῦ παραδείσου;
\vs{2}Καὶ εἶπεν ἡ γυνὴ τῷ ὄφει, ἀπὸ καρποῦ τοῦ ξύλου τοῦ παραδείσου φαγούμεθα·
\vs{3}Ἀπὸ δὲ τοῦ καρποῦ τοῦ ξύλου, ὅ ἐστιν ἐν μέσῳ τοῦ παραδείσου, εἶπεν ὁ Θεός, οὐ φάγεσθε ἀπʼ αὐτοῦ, οὐδὲ μὴ ἅψησθε αὐτοῦ, ἵνα μὴ ἀποθάνητε.
\vs{4}Καὶ εἶπεν ὁ ὄφις τῇ γυναικί· οὐ θανάτῳ ἀποθανεῖσθε·
\vs{5}Ἤδει γὰρ ὁ Θεὸς, ὅτι ᾗ ἂν ἡμέρᾳ φάγητε ἀπʼ αὐτοῦ, διανοιχθήσονται ὑμῶν οἱ ὀφθαλμοί, καὶ ἔσεσθε ὡς θεοί, γινώσκοντες καλὸν καὶ πονηρόν.
\vs{6}Καὶ εἶδεν ἡ γυνὴ, ὅτι καλὸν τὸ ξύλον εἰς βρῶσιν, καὶ ὅτι ἀρεστὸν τοῖς ὀφθαλμοῖς ἰδεῖν, καὶ ὡραῖόν ἐστι τοῦ κατανοῆσαι· καὶ λαβοῦσα ἀπὸ τοῦ καρποῦ αὐτοῦ, ἔφαγε· καὶ ἔδωκε καὶ τῷ ἀνδρὶ αὐτῆς μετʼ αὐτῆς, καὶ ἔφαγον.
\vs{7}Καὶ διηνοίχθησαν οἱ ὀφθαλμοὶ τῶν δύο, καὶ ἔγνωσαν ὅτι γυμνοὶ ἦσαν· καὶ ἔῤῥαψαν φύλλα συκῆς, καὶ ἐποίησαν ἑαυτοῖς περιζώματα.
\vs{8}Καὶ ἤκουσαν τὴς φωνὴς Κυρίου τοῦ Θεοῦ περιπατοῦντος ἐν τῷ παραδείσῳ τὸ δειλινόν· καὶ ἐκρύβησαν ὅ, τε Ἀδὰμ καὶ ἡ γυνὴ αὐτοῦ ἀπὸ προσώπου Κυρίου τοῦ Θεοῦ ἐν μέσῳ τοῦ ξύλου τοῦ παραδείσου.
\vs{9}Καὶ ἐκάλεσεν Κύριος ὁ Θεὸς τὸν Ἀδὰμ, καὶ εἶπεν αὐτῷ· Ἀδὰμ ποῦ εἶ;
\vs{10}Καὶ εἶπεν αὐτῷ· τὴς φωνῆς σου ἤκουσα περιπατοῦντος ἐν τῷ παραδείσῳ, καὶ ἐφοβήθην ὅτι γυμνός εἰμι, καὶ ἐκρύβην.
\vs{11}Καὶ εἶπεν αὐτῷ ὁ Θεὸς, τὶς ἀνήγγειλέ σοι ὅτι γυμνὸς εἶ, εἰ μὴ ἀπὸ τοῦ ξύλου, οὗ ἐνετειλάμην σοι τούτου μόνου μὴ φαγεῖν, ἀπʼ αὐτοῦ ἔφαγες;
\vs{12}Καὶ εἶπεν ὁ Ἀδάμ· ἡ γυνή, ἣν ἔδωκας μετʼ ἐμοῦ, αὕτη μοι ἔδωκεν ἀπὸ τοῦ ξύλου, καὶ ἔφαγον.
\vs{13}Καὶ εἶπε Κύριος ὁ Θεὸς τῇ γυναικί· τί τοῦτο ἐποιήσας; καὶ εἶπεν ἡ γυνὴ, ὁ ὄφις ἠπάτησέ με, καὶ ἔφαγον.

\vs{14}Καὶ εἶπε Κύριος ὁ Θεὸς τῷ ὄφει· ὅτι ἐποίησας τοῦτο, ἐπικατάρατος σὺ ἀπὸ πάντων τῶν κτηνῶν, καὶ ἀπὸ πάντων τῶν θηρίων τῶν ἐπὶ τῆς γῆς· ἐπὶ τῷ στήθει σου καὶ τῇ κοιλίᾳ πορεύσῃ, καὶ γῆν φαγῃ πάσας τὰς ἡμέρας τῆς ζωῆς σου.
\vs{15}Καὶ ἔχθραν θήσω ἀνὰ μέσον σοῦ καὶ ἀνὰ μέσον τῆς γυναικὸς, καὶ ἀνὰ μέσον τοῦ σπέρματός σου, καὶ ἀνὰ μέσον τοῦ σπέρματος αὐτῆς· αὐτός σοῦ τηρήσει κεφαλὴν, καὶ σὺ τηρήσεις αὐτοῦ πτέρναν.
\vs{16}Καὶ τῇ γυναικὶ εἶπε· πληθύνων πληθυνῶ τὰς λύπας σου, καὶ τὸν στεναγμόν σου· ἐν λύπαις τέξῃ τέκνα, καὶ πρὸς τὸν ἄνδρα σου ἡ ἀποστροφή σου· καὶ αὐτός σου κυριεύσει.
\vs{17}Τῷ δὲ Ἀδὰμ εἶπεν· ὅτι ἤκουσας τῆς φωνῆς τῆς γυναικός σου, καὶ ἔφαγες ἀπὸ τοῦ ξύλου, οὗ ἐνετειλάμην σοι τούτου μόνου μὴ φαγεῖν, ἀπʼ αὐτοῦ ἔφαγες, ἐπικατάρατος ἡ γῆ ἐν τοῖς ἔργοις σου· ἐν λύπαις φάγῃ αὐτὴν πάσας τὰς ἡμέρας τῆς ζωῆς σου.
\vs{18}Ἀκάνθας καὶ τριβόλους ἀνατελεῖ σοι, καὶ φαγῇ τὸν χόρτον τοῦ ἀγροῦ.
\vs{19}Ἐν ἱδρῶτι τοῦ προσώπου σου φαγῃ τὸν ἄρτον σου, ἕως τοῦ ἀποστρέψαι σε εἰς τὴν γῆν ἐξ ἧς ἐλήμφθης· ὅτι γῆ εἶ, καὶ εἰς γῆν ἀπελεύσῃ.
\vs{20}Καὶ ἐκάλεσεν Ἀδὰμ τὸ ὄνομα τῆς γυναικὸς αὐτοῦ Ζωή, ὅτι μήτηρ πάντων τῶν ζώντων.
\vs{21}Καὶ ἐποίησε Κύριος ὁ Θεὸς τῷ Ἀδὰμ, καὶ τῇ γυναικὶ αὐτοῦ χιτῶνας δερματίνους, καὶ ἐνέδυσεν αὐτούς.

\vs{22}Καὶ εἶπεν ὁ Θεός, ἰδοὺ Ἀδὰμ γέγονεν ὡς εἷς ἐξ ἡμῶν, τοῦ γινώσκειν καλὸν καὶ πονηρόν· καὶ νῦν μή ποτε ἐκτείνῃ τὴν χεῖρα αὐτοῦ, καὶ λάβῃ τοῦ ξύλου τῆς ζωῆς καὶ φάγῃ, καὶ ζήσεται εἰς τὸν αἰῶνα.
\vs{23}Καὶ ἐξαπέστειλεν αὐτὸν Κύριος ὁ Θεὸς ἐκ τοῦ παραδείσου τῆς τρυφῆς, ἐργάζεσθαι τὴν γῆν ἐξ ἧς ἐλήμφθη.
\vs{24}Καὶ ἐξέβαλεν τὸν Ἀδὰμ, καὶ κατῴκισεν αὐτὸν ἀπέναντι τοῦ παραδείσου τῆς τρυφῆς· καὶ ἔταξε τὰ χερουβὶμ· καὶ τὴν φλογίνην ῥομφαίαν τὴν στρεφομένην, φυλάσσειν τὴν ὁδὸν τοῦ ξύλου τῆς ζωῆς.

\ch{4}
Ἀδὰμ δὲ ἔγνω Εὔαν τὴν γυναῖκα αὐτοῦ, καὶ συλλαβοῦσα ἔτεκε τὸν Κάϊν· καὶ εἶπεν, ἐκτησάμην ἄνθρωπον διὰ τοῦ Θεοῦ.
\vs{2}Καὶ προσέθηκε τεκεῖν τὸν ἀδελφὸν αὐτοῦ τὸν Ἄβελ· καὶ ἐγένετο Ἄβελ ποιμὴν προβάτων, Κάϊν δὲ ἦν ἐργαζόμενος τὴν γῆν.
\vs{3}Καὶ ἐγένετο μεθʼ ἡμέρας ἤνεγκε Κάϊν ἀπὸ τῶν καρπῶν τῆς γῆς θυσίαν τῷ Κυρίῳ·
\vs{4}Καὶ Ἄβελ ἤνεγκε καὶ αὐτὸς ἀπὸ τῶν πρωτοτόκων τῶν προβάτων αὐτοῦ, καὶ ἀπὸ τῶν στεάτων αὐτῶν· καὶ ἐπεῖδεν ὁ Θεὸς ἐπὶ Ἄβελ, καὶ ἐπὶ τοῖς δώροις αὐτοῦ.
\vs{5}Ἐπὶ δὲ Κάϊν, καὶ ἐπὶ ταῖς θυσίαις αὐτοῦ, οὐ προσέσχε· καὶ ἐλυπήθη Κάϊν λίαν, καὶ συνέπεσε τῷ προσώπῳ αὐτοῦ.
\vs{6}Καὶ εἶπε Κύριος ὁ Θεὸς τῷ Κάϊν, ἵνα τί περίλυπος ἐγένου, καὶ ἵνα τί συνέπεσε τὸ πρόσωπόν σου;
\vs{7}Οὐκ ἐὰν ὀρθῶς προσενέγκῃς, ὀρθῶς δὲ μὴ διέλῃς, ἥμαρτες; ἡσυχασον· πρός σὲ ἡ ἀποστροφὴ αὐτοῦ, καὶ σὺ ἄρξεις αὐτοῦ.

\vs{8}Καὶ εἶπεν Κάϊν πρὸς Ἄβελ τὸν ἀδελφὸν αὐτοῦ, διέλθωμεν εἰς τὸ πεδίον· καὶ ἐγένετο ἐν τῷ εἶναι αὐτοὺς ἐν τῷ πεδίῳ, ἀνέστη Κάϊν ἐπὶ Ἄβελ τὸν ἀδελφὸν αὐτοῦ, καὶ ἀπέκτεινεν αὐτόν.
\vs{9}Καὶ εἶπε Κύπιος ὁ Θεὸς πρὸς Κάϊν· ποῦ ἔστιν Ἄβελ ὁ ἀδελφός σου; καὶ εἶπεν, οὐ γινώσκω· μὴ φύλαξ τοῦ ἀδελφοῦ μου εἰμὶ ἐγώ;
\vs{10}Καὶ εἶπε Κύριος, τί πεποίηκας; φωνὴ αἵματος τοῦ ἀδελφοῦ σου βοᾷ πρός με ἐκ τῆς γῆς.
\vs{11}Καὶ νῦν ἐπικατάρατος σὺ ἀπὸ τῆς γῆς, ἣ ἔχανε τὸ στόμα αὐτῆς δέξασθαι τὸ αἷμα τοῦ ἀδελφοῦ σου ἐκ τῆς χειρός σου.
\vs{12}Ὅτε ἐργᾷ τὴν γῆν, καὶ οὐ προσθήσει τὴν ἰσχὺν αὐτῆς δοῦναί σοι· στένων καὶ τρέμων ἐσῃ ἐπὶ τῆς γῆς.
\vs{13}Καὶ εἶπε Κάϊν πρὸς Κύριον τὸν Θεὸν, μείζων ἡ αἰτία μου τοῦ ἀφεθῆναί με.
\vs{14}Εἰ ἐκβάλλεις με σήμερον ἀπὸ προσώπου τῆς γῆς, καὶ ἀπὸ τοῦ προσώπου σου κρυβήσομαι, καὶ ἔσομαι στένων καὶ τρέμων ἐπὶ τῆς γῆς, καὶ ἔσται πᾶς ὁ εὑρίσκων με, ἀποκτενεῖ με.
\vs{15}Καὶ εἴπεν αὐτῷ Κύριος ὁ Θεὸς, οὐχ οὕτω· πᾶς ὁ ἀποκτείνας Κάϊν, ἑπτὰ ἐκδικούμενα παραλύσει. Καὶ ἔθετο Κύριος ὁ Θεὸς σημεῖον τῷ Κάϊν, τοῦ μὴ ἀνελεῖν αὐτὸν πάντα τὸν εὑρίσκοντα αὐτόν.
\vs{16}Ἐξῆλθεν δὲ Κάϊν ἀπὸ προσώπου τοῦ Θεοῦ, καὶ ᾤκησεν ἐν γῇ Ναὶδ κατέναντι Ἐδέμ.

\vs{17}Καὶ ἔγνω Κάϊν τὴν γυναῖκα αὐτοῦ· καὶ συλλαβοῦσα ἔτεκε τὸν Ἐνώχ. Καὶ ἦν οἰκοδομῶν πόλιν· καὶ ἐπῳνόμασε τὴν πόλιν ἐπὶ τῷ ὀνόματι τοῦ υἱοῦ αὐτοῦ, Ἐνώχ.
\vs{18}Ἐγενήθη δὲ τῷ Ἐνὼχ Γαϊδάδ· καὶ Γαϊδὰδ ἐγέννησε τὸν Μαλελεὴλ· καὶ Μαλελεὴλ ἐγέννησε τὸν Μαθουσάλα· καὶ Μαθουσάλα ἐγέννησε τὸν Λάμεχ.

\vs{19}Καὶ ἔλαβεν ἑαυτῷ Λάμεχ δύο γυναῖκας· ὄνομα τῇ μιᾷ, Ἀδά· καὶ ὄνομα τῇ δευτέρᾳ, Σελλά.
\vs{20}Καὶ ἔτεκεν Ἀδὰ τὸν Ἰωβήλ· οὗτος ἦν πατὴρ οἰκούντων ἐν σκηναῖς κτηνοτρόφων.
\vs{21}Καὶ ὄνομα τῷ ἀδελφῷ αὐτοῦ, Ἰουβάλ· οὗτος ἦν ὁ καταδείξας ψαλτήριον καὶ κιθάραν.
\vs{22}Σελλὰ δὲ καὶ αὐτὴ ἔτεκε τὸν Θόβελ· καὶ ἦν σφυροκόπος χαλκεὺς χαλκοῦ καὶ σιδήρου. ἀδελφὴ δὲ Θόβελ, Νοεμά.
\vs{23}Εἶπε δὲ Λάμεχ ταῖς ἑαυτοῦ γυναιξίν, Ἀδὰ καὶ Σελλὰ, ἀκούσατέ μου τῆς φωνῆς, γυναῖκες Λάμεχ, ἐνωτίσασθέ μου τοὺς λόγους· ὅτι ἄνδρα ἀπέκτεινα εἰς τραῦμα ἐμοὶ, καὶ νεανίσκον εἰς μώλωπα ἐμοί.
\vs{24}Ὅτι ἑπτάκις ἐκδεδίκηται ἐκ Κάϊν· ἐκ δὲ Λάμεχ, ἑβδομηκοντάκις ἑπτά.

\vs{25}Ἔγνω δὲ Ἀδὰμ Εὔαν τὴν γυναῖκα αὐτοῦ· καὶ συλλαβοῦσα ἔτεκεν υἱόν· καὶ ἐπωνόμασε τὸ ὄνομα αὐτοῦ Σὴθ, λέγουσα, ἐξανέστησε γάρ μοι ὁ Θεὸς σπέρμα ἕτερον ἀντὶ Ἄβελ, ὃν ἀπέκτεινε Κάϊν.
\vs{26}Καὶ τῷ Σὴθ ἐγένετο υἱός· ἐπωνόμασε δὲ τὸ ὄνομα αὐτοῦ, Ἑνώς· οὗτος ἤλπισεν ἐπικαλεῖσθαι τὸ ὄνομα Κυρίου τοῦ Θεοῦ.

\ch{5}
Αὕτη ἡ βίβλος γενέσεως ἀνθρώπων· ᾗ ἡμέρᾳ ἐποίησεν ὁ Θεὸς τὸν Ἀδὰμ, κατʼ εἰκόνα Θεοῦ ἐποίησεν αὐτόν·
\vs{2}Ἄρσεν καὶ θῆλυ ἐποίησεν αὐτούς· καὶ εὐλόγησεν αὐτούς· καὶ ἐπωνόμασε τὸ ὄνομα αὐτοῦ Ἀδὰμ, ᾗ ἡμέρᾳ ἐποίησεν αὐτούς.
\vs{3}Ἔζησεν δὲ Ἀδὰμ τριάκοντα καὶ διακόσια ἔτη, καὶ ἐγέννησε κατὰ τὴν ἰδέαν αὐτοῦ, καὶ κατὰ τὴν εἰκόνα αὐτοῦ, καὶ ἐπωνόμασε τὸ ὄνομα αὐτοῦ, Σήθ.
\vs{4}Ἐγένοντο δὲ αἱ ἡμέραι Ἀδὰμ, ἃς ἔζησε μετὰ τὸ γεννῆσαι αὐτὸν τὸν Σὴθ, ἔτη ἑπτακόσια· καὶ ἐγέννησεν υἱοὺς καὶ θυγατέρας.
\vs{5}Καὶ ἐγένοντο πᾶσαι αἱ ἡμέραι Ἀδὰμ, ἃς ἔζησε, τριάκοντα καὶ ἐννακόσια ἔτη· καὶ ἀπέθανεν.
\vs{6}Ἔζησε δὲ Σὴθ πέντε καὶ διακόσια ἔτη· καὶ ἐγέννησε τὸν Ἐνώς.
\vs{7}Καὶ ἔζησε Σὴθ μετὰ τὸ γεννῆσαι αὐτὸν τὸν Ἐνὼς, ἑπτὰ ἔτη καὶ ἑπτακόσια· καὶ ἐγέννησεν υἱοὺς καὶ θυγατέρας.
\vs{8}Καὶ ἐγένοντο πᾶσαι αἱ ἡμέραι Σὴθ δώδεκα καὶ ἐννακόσια ἔτη· καὶ ἀπέθανε.
\vs{9}Καὶ ἔζησεν Ἐνὼς ἔτη ἑκατὸν ἐννεήκοντα· καὶ ἐγέννησε τὸν Καϊνᾶν.
\vs{10}Καὶ ἔζησεν Ἐνὼς μετὰ τὸ γεννῆσαι αὐτὸν τὸν Καϊνᾶν, πεντεκαίδεκα ἔτη καὶ ἑπτκόσια· καὶ ἐγέννησεν υἱοὺς καὶ θυγατέρας.
\vs{11}Καὶ ἐγένοντο πᾶσαι αἱ ἡμέραι Ἐνὼς πέντε ἔτη καὶ ἐννακόσια· καὶ ἀπέθανε.
\vs{12}Καὶ ἔζησεν Καϊνᾶν ἑβδομήκοντα καὶ ἑκατὸν ἔτη· καὶ ἐγέννησε τὸν Μαλελεήλ.
\vs{13}Καὶ ἔζησε Καϊνᾶν μετὰ τὸ γεννῆσαι αὐτὸν τὸν Μαλελεὴλ, τεσσεράκοντα καὶ ἑπτακόσια ἔτη· καὶ ἐγέννησεν υἱοὺς καὶ θυγατέρας.
\vs{14}Καὶ ἐγένοντο πᾶσαι αἱ ἡμέραι Καϊνᾶν δέκα ἔτη καὶ ἐννακόσια· καὶ ἀπέθανε.

\vs{15}Καὶ ἔζησε Μαλελεὴλ πέντε καὶ ἑξήκοντα καὶ ἑκατὸν ἔτη· καὶ ἐγέννησε τὸν Ἰάρεδ.
\vs{16}Καὶ ἔζησε Μαλελεὴλ μετὰ τὸ γεννῆσαι αὐτὸν τὸν Ἰάρεδ, ἔτη τριάκοντα καὶ ἑπτακόσια· καὶ ἐγέννησεν υἱοὺς καὶ θυγατέρας.
\vs{17}Καὶ ἐγένοντο πᾶσαι αἱ ἡμέραι Μαλελεὴλ, ἔτη πέντε καὶ ἐννενήκοντα καὶ ὀκτακόσια· καὶ ἀπέθανε.
\vs{18}Καὶ ἔζησεν Ἰάρεδ δύο καὶ ἑξήκοντα ἔτη καὶ ἑκατὸν· καὶ ἐγέννησε τὸν Ἐνώχ.
\vs{19}Καὶ ἔζησεν Ἰάρεδ μετὰ τὸ γεννῆσαι αὐτὸν τὸν Ἐνὼχ, ὀκτακόσια ἔτη· καὶ ἐγέννησεν υἱοὺς καὶ θυγατέρας.
\vs{20}Καὶ ἐγένοντο πᾶσαι αἱ ἡμέραι Ἰάρεδ, δύο καὶ ἑξήκοντα καὶ ἐννακόσια ἔτη· καὶ ἀπέθανε.
\vs{21}Καὶ ἔζησεν Ἐνὼχ πέντε καὶ ἑξήκοντα καὶ ἑκατὸν ἔτη· καὶ ἐγέννησε τὸν Μαθουσάλα.
\vs{22}Εὐηρέστησε δὲ Ἐνὼχ τῷ Θεῷ μετὰ τὸ γεννῆσαι αὐτὸν τὸν Μαθουσάλα, διακόσια ἔτη· καὶ ἐγέννησεν υἱοὺς καὶ θυγατέρας.
\vs{23}Καὶ ἐγένοντο πᾶσαι αἱ ἡμέραι Ἐνὼχ, πέντε καὶ ἑξήκοντα καὶ τριακόσια ἔτη.
\vs{24}Καὶ εὐηρέστησεν Ἐνὼχ τῷ Θεῷ· καὶ οὐχ εὑρίσκετο, ὅτι μετέθηκεν αὐτὸν ὁ Θεός.
\vs{25}Καὶ ἔζησε Μαθουσάλα ἑπτὰ ἔτη καὶ ἑξήκοντα καὶ ἑκατόν· καὶ ἐγέννησε τὸν Λάμεχ.
\vs{26}Καὶ ἔζησε Μαθουσάλα μετὰ τὸ γεννῆσαι αὐτὸν τὸν Λάμεχ, δύο καὶ ὀκτακόσια ἔτη· καὶ ἐγέννησεν υἱοὺς καὶ θυγατέρας.
\vs{27}Καὶ ἐγένοντο πᾶσαι αἱ ἡμέραι Μαθουσάλα ἃς ἔζησεν, ἐννέα καὶ ἑξήκοντα καὶ ἐννακόσια ἔτη· καὶ ἀπέθανε.
\vs{28}Καὶ ἔζησε Λάμεχ ὀκτὼ καὶ ὀγδοήκοντα καὶ ἑκατὸν ἔτη· καὶ ἐγέννησεν υἱόν.
\vs{29}Καὶ ἐπωνόμασε τὸ ὄνομα αὐτοῦ Νῶε, λέγων, οὗτος διαναπαύσει ἡμᾶς ἀπὸ τῶν ἔργων ἡμῶν, καὶ ἀπὸ τῶν λυπῶν τῶν χειρῶν ἡμῶν, καὶ ἀπὸ τῆς γῆς, ἧς κατηράσατο Κύριος ὁ Θεός.
\vs{30}Καὶ ἔζησε Λάμεχ μετὰ τὸ γεννῆσαι αὐτὸν τὸν Νῶε, πεντακόσια καὶ ἑξήκοντα καὶ πέντε ἔτη· καὶ ἐγέννησεν υἱοὺς καὶ θυγατέρας.
\vs{31}Καὶ ἐγένοντο πᾶσαι αἱ ἡμέραι Λάμεχ, ἑπτακόσια καὶ πεντήκοντα τρία ἔτη· καὶ ἀπέθανε.
\vs{32}Καὶ ἦν Νῶε ἐτῶν πεντακοσίων· καὶ ἐγέννησε τρεῖς υἱοὺς, τὸν Σὴμ, τὸν Χὰμ, τὸν Ἰάφεθ.

\ch{6}
Καὶ ἐγένετο ἡνίκα ἤρξαντο οἱ ἄνθρωποι πολλοὶ γίνεσθαι ἐπὶ τῆς γῆς, καὶ θυγατέρες ἐγεννήθησαν αὐτοῖς.
\vs{2}Ἰδόντες δὲ υἱοὶ τοῦ Θεοῦ τὰς θυγατέρας τῶν ἀνθρώπων, ὅτι καλαί εἰσιν, ἔλαβον ἑαυτοῖς γυναῖκας ἀπὸ πασῶν, ὧν ἐξελέξαντο.
\vs{3}Καὶ εἶπε Κύριος ὁ Θεὸς, οὐ μὴ καταμείνῃ τὸ πνεῦμά μου ἐν τοῖς ἀνθρώποις τούτοις εἰς τὸν αἰῶνα, διὰ τὸ εἶναι αὐτοὺς σάρκας· ἔσονται δὲ αἱ ἡμέραι αὐτῶν, ἑκατὸν εἴκοσιν ἔτη.
\vs{4}Οἱ δὲ γίγαντες ἦσαν ἐπὶ τῆς γῆς ἐν ταῖς ἡμέραις ἐκείναις, καὶ μετʼ ἐκεῖνο, ὡς ἂν εἰσεπορεύοντο οἱ υἱοὶ τοῦ Θεοῦ πρὸς τὰς θυγατέρας τῶν ἀνθρώπων, καὶ ἐγεννῶσαν αὐτοῖς· ἐκεῖνοι ἦσαν οἱ γίγαντες οἱ ἀπʼ αἰῶνος, οἱ ἄνθρωποι οἱ ὀνομαστοί.

\vs{5}Ἰδὼν δὲ Κύριος ὁ Θεὸς, ὅτι ἐπληθύνθησαν αἱ κακίαι τῶν ἀνθρώπων ἐπὶ τῆς γῆς, καὶ πᾶς τις διανοεῖται ἐν τῇ καρδίᾳ αὐτοῦ ἐπιμελῶς ἐπὶ τὰ πονηρὰ πάσας τὰς ἡμέρας·
\vs{6}Καὶ ἐνεθυμήθη ὁ Θεὸς, ὅτι ἐποίησε τὸν ἄνθρωπον ἐπὶ τῆς γῆς, καὶ διενοήθη.
\vs{7}Καὶ εἶπεν ὁ Θεὸς, ἀπαλείψω τὸν ἄνθρωπον, ὃν ἐποίησα, ἀπὸ προσώπου τῆς γῆς, ἀπὸ ἀνθρώπου ἕως κτήνους, καὶ ἀπὸ ἑρπετῶν ἕως πετεινῶν τοῦ οὐρανοῦ· ὅτι ἐνεθυμήθην, ὅτι ἐποίησα αὐτούς.

\vs{8}Νῶε δὲ εὗρε χάριν ἐναντίον Κυρίου τοῦ Θεοῦ.
\vs{9}Αὗται δὲ αἱ γενέσεις Νῶε. Νῶε ἄνθρωπος δίκαιος, τέλειος ὢν ἐν τῇ γενεᾷ αὐτοῦ, τῷ Θεῷ εὐηρέστησε Νῶε.
\vs{10}Ἐγέννησε δὲ Νῶε τρεῖς υἱοὺς, τὸν Σὴμ, τὸν Χὰμ, τὸν Ἰάφεθ.
\vs{11}Ἐφθάρη δὲ ἡ γῆ ἐναντίον τοῦ Θεοῦ, καὶ ἐπλήσθη ἡ γῆ ἀδικίας.
\vs{12}Καὶ εἶδε Κύριος ὁ Θεὸς τὴν γῆν, καὶ ἦν κατεφθαρμένη· ὅτι κατέφθειρε πᾶσα σὰρξ τὴν ὁδὸν αὐτοῦ ἐπὶ τῆς γῆς.
\vs{13}Καὶ εἶπε Κύριος ὁ Θεὸς τῷ Νῶε, καιρὸς παντὸς ἀνθρώπου ἥκει ἐναντίον μου, ὅτι ἐπλήσθη ἡ γῆ ἀδικίας ἀπʼ αὐτῶν· καὶ ἰδοὺ ἐγὼ καταφθείρω αὐτοὺς καὶ τὴν γῆν.

\vs{14}Ποίησον οὖν σεαυτῷ κιβωτὸν ἐκ ξύλων τετραγώνων· νοσσιὰς ποιήσεις τὴν κιβωτόν· καὶ ἀσφαλτώσεις αὐτὴν ἔσωθεν καὶ ἔξωθεν τῇ ἀσφάλτῳ.
\vs{15}Καὶ οὕτω ποιήσεις τὴν κιβωτόν· τριακοσίων πήχεων τὸ μῆκος τῆς κιβωτοῦ, καὶ πεντήκοντα πήχεων τὸ πλάτος, καὶ τριάκοντα πήχεων τὸ ὕψος αὐτῆς.
\vs{16}Ἐπισυνάγων ποιήσεις τὴν κιβωτὸν, καὶ εἰς πῆχυν συντελέσεις αὐτὴν ἄνωθεν· τὴν δὲ θύραν τῆς κιβωτοῦ ποιήσεις ἐκ πλαγίων· κατάγαια διώροφα καὶ τριώροφα ποιήσεις αὐτήν.
\vs{17}Ἐγὼ δὲ ἰδοὺ ἐπάγω τὸν κατακλυσμὸν, ὕδωρ ἐπὶ τὴν γῆν, καταφθεῖραι πᾶσαν σάρκα, ἐν ᾗ ἐστι πνεῦμα ζωῆς ὑποκάτω τοῦ οὐρανοῦ· καὶ ὅσα ἂν ᾖ ἐπὶ τῆς γῆς, τελευτήσει.

\vs{18}Καὶ στήσω τὴν διαθήκην μου μετά σου· εἰσελεύσῃ δὲ εἰς τὴν κιβωτὸν σὺ, καὶ οἱ υἱοί σου, καὶ ἡ γυνή σου, καὶ αἱ γυναῖκες τῶν υἱῶν σου μετά σου.
\vs{19}Καὶ ἀπὸ πάντων τῶν κτηνῶν, καὶ ἀπὸ πάντων τῶν ἑρπετῶν, καὶ ἀπὸ πάντων τῶν θηρίων, καὶ ἀπὸ πάσης σαρκὸς δύο δύο ἀπὸ πάντων εἰσάξεις εἰς τὴν κιβωτὸν, ἵνα τρέφῃς μετὰ σεαυτοῦ· ἄρσεν καὶ θῆλυ ἔσονται.
\vs{20}Ἀπὸ πάντων τῶν ὀρνέων τῶν πετεινῶν κατὰ γένος, καὶ ἀπὸ πάντων τῶν κτηνῶν κατὰ γένος, καὶ ἀπὸ πάντων τῶν ἑρπετῶν τῶν ἑρπόντων ἐπὶ τῆς γῆς κατὰ γένος αὐτῶν, δύο δύο ἀπὸ πάντων εἰσελεύσονται πρὸς σὲ τρέφεσθαι μετά σου, ἄρσεν καὶ θῆλυ.
\vs{21}Σὺ δὲ λήψῃ σεαυτῷ ἀπὸ πάντων τῶν βρωμάτων ἃ ἔδεσθε, καὶ συνάξεις πρὸς σεαυτὸν, καὶ ἔσται σοι καὶ ἐκείνοις φαγεῖν.
\vs{22}Καὶ ἐποίησε Νῶε πάντα ὅσα ἐνετείλατο αὐτῷ Κύριος ὁ Θεὸς, οὕτως ἐποίησε.

\ch{7}
Καὶ εἶπε Κύριος ὁ Θεὸς πρὸς Νῶε, εἴσελθε σὺ καὶ πᾶς ὁ οἶκός σου εἰς τὴν κιβωτὸν, ὅτι σὲ εἶδον δίκαιον ἐναντίον μου ἐν τῇ γενεᾷ ταύτῃ.
\vs{2}Ἀπὸ δὲ τῶν κτηνῶν τῶν καθαρῶν εἰσάγαγε πρὸς σὲ ἑπτὰ ἑπτὰ ἄρσεν καὶ θῆλυ, ἀπὸ δὲ τῶν κτηνῶν τῶν μὴ καθαρῶν δύο δύο ἄρσεν καὶ θῆλυ.
\vs{3}Καὶ ἀπὸ τῶν πετεινῶν τοῦ οὐρανοῦ τῶν καθαρῶν ἑπτὰ ἑπτὰ ἄρσεν καὶ θῆλυ, καὶ ἀπὸ πάντων τῶν πετεινῶν τῶν μὴ καθαρῶν δύο δύο ἄρσεν καὶ θῆλυ, διαθρέψαι σπέρμα ἐπὶ πᾶσαν τὴν γῆν.
\vs{4}Ἔτι γὰρ ἡμερῶν ἑπτὰ ἐγὼ ἐπάγω ὑετὸν ἐπὶ τὴν γῆν, τεσσαράκοντα ἡμέρας καὶ τεσσαράκοντα νύκτας· καὶ ἐξαλείψω πᾶν τὸ ἀνάστημα, ὃ ἐποίησα ἀπὸ προσώπου πάσης τῆς γῆς.
\vs{5}Καὶ ἐποίησε Νῶε πάντα, ὅσα ἐνετείλατο αὐτῷ Κύριος ὁ Θεός.
\vs{6}Νῶε δὲ ἦν ἐτῶν ἑξακοσίων, καὶ ὁ κατακλυσμὸς τοῦ ὕδατος ἐγένετο ἐπὶ τῆς γῆς.
\vs{7}Εἰσῆλθε δὲ Νῶε καὶ οἱ υἱοὶ αὐτοῦ, καὶ ἡ γυνὴ αὐτοῦ, καὶ αἱ γυναῖκες τῶν υἱῶν αὐτοῦ μετʼ αὐτοῦ εἰς τὴν κιβωτὸν, διὰ τὸ ὕδωρ τοῦ κατακλυσμοῦ.
\vs{8}Καὶ ἀπὸ τῶν πετεινῶν τῶν καθαρῶν, καὶ ἀπὸ τῶν πετεινῶν τῶν μὴ καθαρῶν, καὶ ἀπὸ τῶν κτηνῶν τῶν καθαρῶν, καὶ ἀπὸ τῶν κτηνῶν τῶν μὴ καθαρῶν, καὶ ἀπὸ πάντων τῶν ἑρπόντων ἐπὶ τῆς γῆς,
\vs{9}δύο δύο εἰσῆλθον πρὸς Νῶε εἰς τὴν κιβωτὸν ἄρσεν καὶ θῆλυ, καθὰ ἐνετείλατο ὁ Θεὸς τῷ Νῶε.
\vs{10}Καὶ ἐγένετο μετὰ τὰς ἑπτὰ ἡμέρας, καὶ τὸ ὕδωρ τοῦ κατακλυσμοῦ ἐγένετο ἐπὶ τῆς γῆς.
\vs{11}Ἐν τῷ ἑξακοσιοστῷ ἔτει ἐν τῇ ζωῇ τοῦ Νῶε, τοῦ δευτέρου μηνὸς, ἑβδόμῃ καὶ εἰκάδι τοῦ μηνὸς, τῇ ἡμέρᾳ ταύτῃ ἐῤῥάγησαν πᾶσαι αἱ πηγαὶ τῆς ἀβύσσου, καὶ οἱ καταῤῥάκται τοῦ οὐρανοῦ ἠνεῴχθησαν.
\vs{12}Καὶ ἐγένετο ὁ ὑετὸς ἐπὶ τῆς γῆς τεσσαράκοντα ἡμέρας καὶ τεσσαράκοντα νύκτας.
\vs{13}Ἐν τῇ ἡμέρᾳ ταύτῃ εἰσῆλθε Νῶε, Σὴμ, Χὰμ, Ἰάφεθ, οἱ υἱοὶ Νῶε, καὶ ἡ γυνὴ Νῶε, καὶ αἱ τρεῖς γυναῖκες τῶν υἱῶν αὐτοῦ μετʼ αὐτοῦ, εἰς τὴν κιβωτόν.
\vs{14}Καὶ πάντα τὰ θηρία κατὰ γένος, καὶ πάντα τὰ κτήνη κατὰ γένος, καὶ πᾶν ἑρπετὸν κινούμενον ἐπὶ τῆς γῆς κατὰ γένος, καὶ πᾶν ὄρνεον πετεινὸν κατὰ γένος αὐτοῦ,
\vs{15}εἰσῆλθον πρὸς Νῶε εἰς τὴν κιβωτὸν, δύο δύο ἄρσεν καὶ θῆλυ ἀπὸ πάσης σαρκὸς, ἐν ᾧ ἐστι πνεῦμα ζωῆς.
\vs{16}Καὶ τὰ εἰσπορευόμενα ἄρσεν καὶ θῆλυ ἀπὸ πάσης σαρκὸς εἰσῆλθε, καθὰ ἐνετείλατο ὁ Θεὸς τῷ Νῶε· καὶ ἔκλεισε Κύριος ὁ Θεὸς τὴν κιβωτὸν ἔξωθεν αὐτοῦ.

\vs{17}Καὶ ἐγένετο ὁ κατακλυσμὸς τεσσαράκοντα ἡμέρας καὶ τεσσαράκοντα νύκτας ἐπὶ τῆς γῆς· καὶ ἐπεπληθύνθη τὸ ὕδωρ· καὶ ἐπῇρε τὴν κιβωτὸν, καὶ ὑψώθη ἀπὸ τῆς γῆς.
\vs{18}Καὶ ἐπεκράτει τὸ ὕδωρ, καὶ ἐπληθύνετο σφόδρα ἐπὶ τῆς γῆς· καὶ ἐπεφέρετο ἡ κιβωτὸς ἐπάνω τοῦ ὕδατος.
\vs{19}Τὸ δὲ ὕδωρ ἐπεκράτει σφόδρα σφόδρα ἐπὶ τῆς γῆς· καὶ ἐκάλυψε πάντα τὰ ὄρη τὰ ὑψηλὰ, ἃ ἦν ὑποκάτω τοῦ οὐρανοῦ.
\vs{20}Πεντεκαίδεκα πήχεις ὑπεράνω ὑψώθη τὸ ὕδωρ· καὶ ἐπεκάλυψε πάντα τὰ ὄρη τὰ ὑψηλά.
\vs{21}Καὶ ἀπέθανε πᾶσα σὰρξ κινουμένη ἐπὶ τῆς γῆς τῶν πετεινῶν, καὶ τῶν κτηνῶν, καὶ τῶν θηρίων· καὶ πᾶν ἑρπετὸν κινούμενον ἐπὶ τῆς γῆς, καὶ πᾶς ἄνθρωπος.
\vs{22}Καὶ πάντα ὅσα ἔχει πνοὴν ζωῆς, καὶ πᾶν ὃ ἦν ἐπὶ τῆς ξηρᾶς, ἀπέθανε.
\vs{23}Καὶ ἐξήλειψε πᾶν τὸ ἀνάστημα, ὃ ἦν ἐπὶ προσώπου τῆς γῆς, ἀπὸ ἀνθρώπου ἕως κτήνους, καὶ ἑρπετῶν, καὶ τῶν πετεινῶν τοῦ οὐρανοῦ· καὶ ἐξηλείφησαν ἀπὸ τῆς γῆς· καὶ κατελείφθη μόνος Νῶε, καὶ οἱ μετʼ αὐτοῦ ἐν τῇ κιβωτῷ.
\vs{24}Καὶ ὑψώθη τὸ ὕδωρ ἐπὶ τῆς γῆς ἡμέρας ἑκατὸν πεντήκοντα.

\ch{8}
Καὶ ἀνεμνήσθη ὁ Θεὸς τοῦ Νῶε, καὶ πάντων τῶν θηρίων, καὶ πάντων τῶν κτηνῶν, καὶ πάντων τῶν πετεινῶν, καὶ πάντων τῶν ἑρπετῶν τῶν ἑρπόντων, ὅσα ἦν μετʼ αὐτοῦ ἐν τῇ κιβωτῷ· καὶ ἐπήγαγεν ὁ Θεὸς πνεῦμα ἐπὶ τὴν γῆν, καὶ ἐκόπασε τὸ ὕδωρ.
\vs{2}Καὶ ἐπεκαλύφθησαν αἱ πηγαὶ τῆς ἀβύσσου, καὶ οἱ καταῤῥάκται τοῦ οὐρανοῦ, καὶ συνεσχέθη ὁ ὑετὸς ἀπὸ τοῦ οὐρανοῦ.
\vs{3}Καὶ ἐνεδίδου τὸ ὕδωρ πορευόμενον ἀπὸ τῆς γῆς· καὶ ἠλαττονοῦτο τὸ ὕδωρ μετὰ πεντήκοντα καὶ ἑκατὸν ἡμέρας.
\vs{4}Καὶ ἐκάθισεν ἡ κιβωτὸς ἐν μηνὶ τῷ ἑβδόμῳ, ἑβδόμῃ καὶ εἰκάδι τοῦ μηνὸς, ἐπὶ τὰ ὄρη τὰ Ἀραράτ.
\vs{5}Τὸ δὲ ὕδωρ ἠλαττονοῦτο ἕως τοῦ δεκάτου μηνός. Καὶ ἐν τῷ δεκάτῳ μηνὶ, τῇ πρώτῃ τοῦ μηνὸς, ὤφθησαν αἱ κεφαλαὶ τῶν ὀρέων.
\vs{6}Καὶ ἐγένετο μετὰ τεσσαράκοντα ἡμέρας ἠνέῳξε Νῶε τὴν θυρίδα τῆς κιβωτοῦ, ἣν ἐποίησε.
\vs{7}Καὶ ἀπέστειλε τὸν κόρακα· καὶ ἐξελθὼν, οὐκ ἀνέστρεψεν ἕως τοῦ ξηρανθῆναι τὸ ὕδωρ ἀπὸ τῆς γῆς.
\vs{8}Καὶ ἀπέστειλε τὴν περιστερὰν ὀπίσω αὐτοῦ, ἰδεῖν εἰ κεκόπακε τὸ ὕδωρ ἀπὸ τῆς γῆς.
\vs{9}Καὶ οὐχ εὑροῦσα ἡ περιστερὰ ἀνάπαυσιν τοῖς ποσὶν αὐτῆς, ἀνέστρεψε πρὸς αὐτὸν εἰς τὴν κιβωτὸν, ὅτι ὕδωρ ἦν ἐπὶ πᾶν τὸ πρόσωπον τῆς γῆς· καὶ ἐκτείνας τὴν χεῖρα ἔλαβεν αὐτὴν, καὶ εἰσήγαγεν αὐτὴν πρὸς ἑαυτὸν εἰς τὴν κιβωτόν.
\vs{10}Καὶ ἐπισχὼν ἔτι ἡμέρας ἑπτὰ ἑτέρας, πάλιν ἐξαπέστειλε τὴν περιστερὰν ἐκ τῆς κιβωτοῦ.
\vs{11}Καὶ ἀνέστρεψε πρὸς αὐτὸν ἡ περιστερὰ τὸ πρὸς ἑσπέραν· καὶ εἶχε φύλλον ἐλαίας κάρφος ἐν τῷ στόματι αὐτῆς· καὶ ἔγνω Νῶε, ὅτι κεκόπακε τὸ ὕδωρ ἀπὸ τῆς γῆς.
\vs{12}Καὶ ἐπισχὼν ἔτι ἡμέρας ἑπτὰ ἑτέρας, πάλιν ἐξαπέστειλε τὴν περιστερὰν, καὶ οὐ προσέθετο τοῦ ἐπιστρέψαι πρὸς αὐτὸν ἔτι.
\vs{13}Καὶ ἐγένετο ἐν τῷ ἑνὶ καὶ ἑξακοσιοστῷ ἔτει ἐν τῇ ζωῇ τοῦ Νῶε, τοῦ πρώτου μηνὸς, μιᾷ τοῦ μηνὸς, ἐξέλιπε τὸ ὕδωρ ἀπὸ τῆς γῆς. Καὶ ἀπεκάλυψε Νῶε τὴν στέγην τῆς κιβωτοῦ, ἣν ἐποίησε· καὶ εἶδεν ὅτι ἐξέλιπε τὸ ὕδωρ ἀπὸ προσώπου τῆς γῆς.
\vs{14}Ἐν δὲ τῷ δευτέρῳ μηνὶ ἐξηράνθη ἡ γῆ, ἑβδόμῃ καὶ εἰκάδι τοῦ μηνός.

\vs{15}Καὶ εἶπε Κύριος ὁ Θεὸς πρὸς Νῶε, λέγων,
\vs{16}Ἔξελθε ἐκ τῆς κιβωτοῦ σὺ, καὶ ἡ γυνή σου, καὶ οἱ υἱοί σου, καὶ αἱ γυναῖκες τῶν υἱῶν σου μετὰ σοῦ,
\vs{17}Καὶ πάντα τὰ θηρία ὅσα ἐστὶ μετὰ σοῦ, καὶ πᾶσα σὰρξ ἀπὸ πετεινῶν ἕως κτηνῶν, καὶ πᾶν ἑρπετὸν κινούμενον ἐπὶ τῆς γῆς, ἐξάγαγε μετὰ σεαυτοῦ. καὶ αὐξάνεσθε καὶ πληθύνεσθε ἐπὶ τῆς γῆς.
\vs{18}Καὶ ἐξῆλθε Νῶε, καὶ ἡ γυνὴ αὐτοῦ, καὶ οἱ υἱοὶ αὐτοῦ, καὶ αἱ γυναῖκες τῶν υἱῶν αὐτοῦ μετʼ αὐτοῦ·
\vs{19}Καὶ πάντα τὰ θηρία, καὶ πάντα τὰ κτήνη, καὶ πᾶν πετεινὸν, καὶ πᾶν ἑρπετὸν κινούμενον ἐπὶ τῆς γῆς κατὰ γένος αὐτῶν, ἐξήλθοσαν ἐκ τῆς κιβωτοῦ.

\vs{20}Καὶ ᾠκοδόμησε Νῶε θυσιαστήριον τῷ Κυρίῳ· καὶ ἔλαβεν ἀπὸ πάντων τῶν κτηνῶν τῶν καθαρῶν, καὶ ἀπὸ πάντων τῶν πετεινῶν τῶν καθαρῶν, καὶ ἀνήνεγκεν εἰς ὁλοκάρπωσιν ἐπὶ τὸ θυσιαστήριον.
\vs{21}Καὶ ὠσφράνθη Κύριος ὁ Θεὸς ὀσμὴν εὐωδίας. Καὶ εἶπε Κύριος ὁ Θεὸς διανοηθείς, οὐ προσθήσω ἔτι καταράσασθαι τὴν γῆν διὰ τὰ ἔργα τῶν ἀνθρώπων· ὅτι ἔγκειται ἡ διάνοια τοῦ ἀνθρώπου ἐπιμελῶς ἐπὶ τὰ πονηρὰ ἐκ νεότητος αὐτοῦ· οὐ προσθήσω οὖν ἔτι πατάξαι πᾶσαν σάρκα ζῶσαν, καθὼς ἐποίησα.
\vs{22}Πάσας τὰς ἡμέρας τῆς γῆς, σπέρμα καὶ θερισμὸς, ψύχος καὶ καῦμα, θέρος καὶ ἔαρ, ἡμέραν καὶ νύκτα, οὐ καταπαύσουσι.

\ch{9}
Καὶ εὐλόγησεν ὁ Θεὸς τὸν Νῶε, καὶ τοὺς υἱοὺς αὐτοῦ· καὶ εἶπεν αὐτοῖς· αὐξάνεσθε καὶ πληθύνεσθε, καὶ πληρώσατε τὴν γῆν, καὶ κατακυριεύσατε αὐτῆς.
\vs{2}Καὶ ὁ τρόμος, καὶ ὁ φόβος ὑμῶν, ἔσται ἐπὶ πᾶσι τοῖς θηρίοις τῆς γῆς, ἐπὶ πάντα τὰ πετεινὰ τοῦ οὐρανοῦ, καὶ ἐπὶ πάντα τὰ κινούμενα ἐπὶ τῆς γῆς, καὶ ἐπὶ πάντας τοὺς ἰχθύας τῆς θαλάσσης· ὑπὸ χεῖρας ὑμῖν δέδωκα.
\vs{3}Καὶ πᾶν ἑρπετὸν, ὅ ἐστι ζῶν, ὑμῖν ἔσται εἰς βρῶσιν· ὡς λάχανα χόρτου δέδωκα ὑμῖν τὰ πάντα.
\vs{4}Πλὴν κρέας ἐν αἵματι ψυχῆς οὐ φάγεσθε.
\vs{5}Καὶ γὰρ τὸ ὑμέτερον αἷμα τῶν ψυχῶν ὑμῶν ἐκ χειρὸς πάντων τῶν θηρίων ἐκζητήσω αὐτό· καὶ ἐκ χειρὸς ἀνθρώπου ἀδελφοῦ ἐκζητήσω τὴν ψυχὴν τοῦ ἀνθρώπου.
\vs{6}Ὁ ἐκχέων αἷμα ἀνθρώπου, ἀντὶ τοῦ αἵματος αὐτοῦ ἐκχυθήσεται, ὅτι ἐν εἰκόνι Θεοῦ ἐποίησα τὸν ἄνθρωπον.
\vs{7}Ὑμεῖς δὲ αὐξάνεσθε, καὶ πληθύνεσθε, καὶ πληρώσατε τὴν γῆν, καὶ κατακυριεύσατε αὐτῆς.

\vs{8}Καὶ εἶπεν ὁ Θεὸς τῷ Νῶε καὶ τοῖς υἱοῖς αὐτοῦ, μετʼ αὐτοῦ λέγων,
\vs{9}καὶ ἰδοὺ ἐγὼ ἀνίστημι τὴν διαθήκην μου ὑμῖν, καὶ τῷ σπέρματι ὑμῶν μεθʼ ὑμᾶς,
\vs{10}καὶ πάσῃ ψυχῇ ζώσῃ μεθʼ ὑμῶν, ἀπὸ ὀρνέων, καὶ ἀπὸ κτηνῶν· καὶ πᾶσι τοῖς θηρίοις τῆς γῆς, ὅσα ἐστὶ μεθʼ ὑμῶν ἀπὸ πάντων τῶν ἐξελθόντων ἐκ τῆς κιβωτοῦ.
\vs{11}Καὶ στήσω τὴν διαθήκην μου πρὸς ὑμᾶς· καὶ οὐκ ἀποθανεῖται πᾶσα σὰρξ ἔτι ἀπὸ τοῦ ὕδατος τοῦ κατακλυσμοῦ· καὶ οὐκ ἔτι ἔσται κατακλυσμὸς ὕδατος, καταφθεῖραι πᾶσαν τὴν γῆν.
\vs{12}Καὶ εἶπε Κύριος ὁ Θεὸς πρὸς Νῶε· τοῦτο τὸ σημεῖον τῆς διαθήκης, ὃ ἐγὼ δίδωμι ἀνὰ μέσον ἐμοῦ καὶ ὑμῶν, καὶ ἀνὰ μέσον πάσης ψυχῆς ζώσης, ἥ ἐστι μεθʼ ὑμῶν εἰς γενεὰς αἰωνίους.
\vs{13}Τὸ τόξον μου τίθημι ἐν τῇ νεφέλῃ, καὶ ἔσται εἰς σημεῖον διαθήκης ἀνὰ μέσον ἐμοῦ καὶ τῆς γῆς.
\vs{14}Καὶ ἔσται ἐν τῷ συννεφεῖν με νεφέλας ἐπὶ τὴν γῆν, ὀφθήσεται τὸ τόξον ἐν τῇ νεφέλῃ.
\vs{15}Καὶ μνησθήσομαι τῆς διαθήκης μου, ἥ ἐστιν ἀνὰ μέσον ἐμοῦ καὶ ὑμῶν, καὶ ἀνὰ μέσον πάσης ψυχῆς ζώσης ἐν πάσῃ σαρκί· καὶ οὐκ ἔσται ἔτι τὸ ὕδωρ εἰς κατακλυσμὸν, ὥστε ἐξαλεῖψαι πᾶσαν σάρκα.
\vs{16}Καὶ ἔσται τὸ τόξον μου ἐν τῇ νεφέλῃ· καὶ ὄψομαι τοῦ μνησθῆναι διαθήκην αἰώνιον ἀνὰ μέσον ἐμοῦ καὶ τῆς γῆς, καὶ ἀνὰ μέσον ψυχῆς ζώσης ἐν πάσῃ σαρκὶ, ἥ ἐστιν ἐπὶ τῆς γῆς.
\vs{17}Καὶ εἶπεν ὁ Θεὸς τῷ Νῶε, τοῦτο τὸ σημεῖον τῆς διαθήκης, ἧς διεθέμην ἀνὰ μέσον ἐμοῦ, καὶ ἀνὰ μέσον πάσης σαρκὸς, ἥ ἐστιν ἐπὶ τῆς γῆς.

\vs{18}Ἦσαν δὲ οἱ υἱοὶ Νῶε, οἱ ἐξελθόντες ἐκ τῆς κιβωτοῦ, Σὴμ, Χὰμ, Ἰάφεθ. Χὰμ δὲ ἦν πατὴρ Χαναάν.
\vs{19}Τρεῖς οὗτοί εἰσιν υἱοὶ Νῶε· ἀπὸ τούτων διεσπάρησαν ἐπὶ πᾶσαν τὴν γῆν.
\vs{20}Καὶ ἤρξατο Νῶε ἄνθρωπος γεωργὸς γῆς, καὶ ἐφύτευσεν ἀμπελῶνα.
\vs{21}Καὶ ἔπιεν ἐκ τοῦ οἴνου, καὶ ἐμεθύσθη, καὶ ἐγυμνώθη ἐν τῷ οἴκῳ αὐτοῦ.
\vs{22}Καὶ εἶδε Χὰμ ὁ πατὴρ Χαναὰν τὴν γύμνωσιν τοῦ πατρὸς αὐτοῦ, καὶ ἐξελθὼν ἀνήγγειλε τοῖς δυσὶν ἀδελφοῖς αὐτοῦ ἔξω.
\vs{23}Καὶ λαβόντες Σὴμ καὶ Ἰάφεθ τὸ ἱμάτιον, ἐπέθεντο ἐπὶ τὰ δύο νῶτα αὐτῶν, καὶ ἐπορεύθησαν ὀπισθοφανῶς, καὶ συνεκάλυψαν τὴν γύμνωσιν τοῦ πατρὸς αὐτῶν· καὶ τὸ πρόσωπον αὐτῶν ὀπισθοφανῶς, καὶ τὴν γύμνωσιν τοῦ πατρὸς αὐτῶν οὐκ εἶδον.
\vs{24}Ἐξένηψε δὲ Νῶε ἀπὸ τοῦ οἴνου, καὶ ἔγνω ὅσα ἐποίησεν αὐτῷ ὁ υἱὸς αὐτοῦ ὁ νεώτερος.
\vs{25}Καὶ εἶπεν, ἐπικατάρατος Χαναὰν παῖς· οἰκέτης ἔσται τοῖς ἀδελφοῖς αὐτοῦ.
\vs{26}Καὶ εἶπεν, εὐλογητὸς Κύριος ὁ Θεὸς τοῦ Σήμ· καὶ ἔσται Χαναὰν παῖς οἰκέτης αὐτοῦ.
\vs{27}Πλατύναι ὁ Θεὸς τῷ Ἰάφεθ, καὶ κατοικησάτω ἐν τοῖς οἴκοις τοῦ Σήμ· καὶ γενηθήτω Χαναὰν παῖς αὐτοῦ.

\vs{28}Ἔζησε δὲ Νῶε μετὰ τὸν κατακλυσμὸν ἔτη τριακόσια πεντήκοντα.
\vs{29}Καὶ ἐγένοντο πᾶσαι αἱ ἡμέραι Νῶε ἐννακόσια πεντήκοντα ἔτη· καὶ ἀπέθανεν.

\ch{10}
Αὗται δὲ αἱ γενέσεις τῶν υἱῶν Νῶε, Σὴμ, Χὰμ, Ἰάφεθ· καὶ ἐγεννήθησαν αὐτοῖς υἱοὶ μετὰ τὸν κατακλυσμόν.

\vs{2}Υἱοὶ Ἰάφεθ, Γαμὲρ, καὶ Μαγὼγ, καὶ Μαδοὶ, καὶ Ἰωύαν, καὶ Ἐλισὰ, καὶ Θοβὲλ, καὶ Μοσὸχ, καὶ Θείρας.
\vs{3}Καὶ υἱοὶ Γαμὲρ, Ἀσχανὰζ, καὶ Ῥιφὰθ, καὶ Θοργαμά.
\vs{4}Καὶ υἱοὶ Ἰωύαν, Ἐλισὰ, καὶ Θάρσεις, Κήτιοι, Ῥόδὶοι.
\vs{5}Ἐκ τούτων ἀφωρίσθησαν νῆσοι τῶν ἐθνῶν ἐν τῇ γῇ αὐτῶν· ἕκαστος κατὰ γλῶσσαν ἐν ταῖς φυλαῖς αὐτῶν, καὶ ἐν τοῖς ἔθνεσιν αὐτῶν.

\vs{6}Υἱοὶ δὲ Χὰμ, Χοὺς, καὶ Μεσραῒν, Φοὺδ, καὶ Χαναάν.
\vs{7}Υἱοὶ δὲ Χοὺς, Σαβὰ, καὶ Εὐϊλὰ, καὶ Σαβαθὰ, καὶ Ῥεγμὰ, καὶ Σαβαθακά· υἱοὶ δὲ Ῥεγμὰ, Σαβὰ, καὶ Δαδάν.
\vs{8}Χοὺς δὲ ἐγέννησε τὸν Νεβρώδ· οὗτος ἤρξατο εἶναι γίγας ἐπὶ τῆς γῆς.
\vs{9}Οὗτος ἦν γίγας κυνηγὸς ἐναντίον Κυρίου τοῦ Θεοῦ· διὰ τοῦτο ἐροῦσιν, ὡς Νεβρὼδ γίγας κυνηγὸς ἐναντίον Κυρίου.
\vs{10}Καὶ ἐγένετο ἀρχὴ τῆς βασιλείας αὐτοῦ Βαβυλὼν, καὶ Ὀρὲχ, καὶ Ἀρχὰδ, καὶ Χαλάννη, ἐν τῇ γῇ Σεναάρ.
\vs{11}Ἐκ τῆς γῆς ἐκείνης ἐξῆλθεν Ἀσσούρ· καὶ ᾠκοδόμησε τὴν Νινευῒ, καὶ τὴν Ῥοωβὼθ πόλιν, καὶ τὴν Χαλὰχ,
\vs{12}καὶ τὴν Δασὴ ἀνὰ μέσον Νινευῒ, καὶ ἀνὰ μέσον Χαλάχ· αὕτη ἡ πόλις μεγάλη.
\vs{13}Καὶ Μεσραῒν ἐγέννησε τοὺς Λουδιεὶμ, καὶ τοὺς Νεφθαλεὶμ, καὶ τοὺς Ἐνεμετιεὶμ, καὶ τοὺς Λαβιεὶμ,
\vs{14}καὶ τοὺς Πατροσωνιεὶμ, καὶ τοὺς Χασμωνιεὶμ, ὅθεν ἐξῆλθε Φυλιστιεὶμ, καὶ τοὺς Γαφθοριείμ.
\vs{15}Χαναὰν δὲ ἐγέννησε τὸν Σιδῶνα πρωτότοκον αὐτοῦ, καὶ τὸν Χετταῖον,
\vs{16}καὶ τὸν Ἰεβουσαῖον, καὶ τὸν Ἀμοῤῥαῖον, καὶ τὸν Γεργεσαῖον,
\vs{17}καὶ τὸν Εὐαῖον, καὶ τὸν Ἀρουκαῖον, καὶ τὸν Ἀσενναῖον,
\vs{18}καὶ τον Ἀράδιον, καὶ τὸν Σαμαραῖον, καὶ τὸν Ἀμαθί. Καὶ μετὰ τοῦτο διεσπάρησαν αἱ φυλαὶ τῶν Χαναναίων.
\vs{19}Καὶ ἐγένετο τὰ ὅρια τῶν Χαναναίων ἀπὸ Σιδῶνος ἕως ἐλθεῖν εἰς Γεραρὰ καὶ Γαζὰν, ἕως ἐλθεῖν ἕως Σοδόμων καὶ Γομόῤῥας, Ἀδαμὰ καὶ Σεβωῒμ ἕως Δασά.
\vs{20}Οὗτοι υἱοὶ Χὰμ, ἐν ταῖς φυλαῖς αὐτῶν, κατὰ γλώσσας αὐτῶν, ἐν ταῖς χώραις αὐτῶν, καὶ ἐν τοῖς ἔθνεσιν αὐτῶν.

\vs{21}Καὶ τῷ Σὴμ ἐγεννήθη καὶ αὐτῷ πατρὶ πάντων τῶν υἱῶν Ἕβερ, ἀδελφῷ Ἰάφεθ τοῦ μείζονος.
\vs{22}Υἱοὶ Σὴμ, Ἐλὰμ, καὶ Ἀσσοὺρ, καὶ Ἀρφαξὰδ, καὶ Λοὺδ, καὶ Ἀρὰμ, καὶ Καϊνᾶν.
\vs{23}Καὶ υἱοὶ Ἀρὰμ, Οὒζ, καὶ Οὒλ, καὶ Γατὲρ, καὶ Μοσόχ.
\vs{24}Καὶ Ἀρφαξὰδ ἐγέννησε τὸν Καϊνᾶν, καὶ Καϊνᾶν ἐγέννησε τὸν Σαλά· Σαλὰ δὲ ἐγέννησε τὸν Ἕβερ.
\vs{25}Καὶ τῷ Ἕβερ ἐγεννήθησαν δύο υἱοί· ὄνομα τῷ ἑνὶ, Φαλὲγ, ὅτι ἐν ταῖς ἡμέραις αὐτοῦ διεμερίσθη ἡ γῆ· καὶ ὄνομα τῷ ἀδελφῷ αὐτοῦ Ἰεκτάν.
\vs{26}Ἰεκτὰν δὲ ἐγέννησε τὸν Ἐλμωδὰδ, καὶ Σαλὲθ, καὶ τὸν Σαρμὼθ, καὶ Ἰαρὰχ,
\vs{27}καὶ Ὁδοῤῥὰ, καὶ Αἰβὴλ, καὶ Δεκλὰ, καὶ Εὐὰλ,
\vs{28}καὶ Ἀβιμαὲλ, καὶ Σαβὰ,
\vs{29}καὶ Οὐφεὶρ, καὶ Εὑεϊλὰ, καὶ Ἰωβάβ· πάντες οὗτοι υἱοὶ Ἰεκτάν.
\vs{30}Καὶ ἐγένετο ἡ κατοίκησις αὐτῶν, ἀπὸ Μασσῆ ἕως ἐλθεῖν εἰς Σαφηρὰ ὄρος ἀνατολῶν.
\vs{31}Οὗτοι υἱοὶ Σὴμ, ἐν ταῖς φυλαῖς αὐτῶν, κατὰ γλώσσας αὐτῶν, ἐν ταῖς χώραις αὐτῶν, καὶ ἐν τοῖς ἔθνεσιν αὐτῶν.
\vs{32}Αὗται αἱ φυλαὶ υἱῶν Νῶε κατὰ γενέσεις αὐτῶν, κατὰ ἔθνη αὐτῶν· ἀπὸ τούτων διεσπάρησαν νῆσοι τῶν ἐθνῶν ἐπὶ τῆς γῆς μετὰ τὸν κατακλυσμόν.

\ch{11}
Καὶ ἦν πᾶσα ἡ γῆ χεῖλος ἓν, καὶ φωνὴ μία πᾶσι.
\vs{2}Καὶ ἐγένετο ἐν τῷ κινῆσαι αὐτοὺς ἀπὸ ἀνατολῶν, εὗρον πεδίον ἐν γῇ Σεναὰρ, καὶ κατῴκησαν ἐκεῖ.
\vs{3}Καὶ εἶπεν ἄνθρωπος τῷ πλησίον αὐτοῦ, δεῦτε πλινθεύσωμεν πλίνθους, καὶ ὀπτήσωμεν αὐτὰς πυρί· καὶ ἐγένετο αὐτοῖς ἡ πλίνθος εἰς λίθον, καὶ ἄσφαλτος ἦν αὐτοῖς ὁ πηλός.
\vs{4}Καὶ εἶπαν, δεῦτε οἰκοδομήσωμεν ἑαυτοῖς πόλιν καὶ πύργον, οὗ ἔσται ἡ κεφαλὴ ἕως τοῦ οὐρανοῦ, καὶ ποιήσωμεν ἑαυτοῖς ὄνομα, πρὸ τοῦ διασπαρῆναι ἡμᾶς ἐπὶ προσώπου πάσης τῆς γῆς.
\vs{5}Καὶ κατέβη Κύριος ἰδεῖν τὴν πόλιν καὶ τὸν πύργον, ὃν ᾠκοδόμησαν οἱ υἱοὶ τῶν ἀνθρώπων.
\vs{6}Καὶ εἶπε Κύριος, ἰδοὺ γένος ἓν, καὶ χεῖλος ἓν πάντων, καὶ τοῦτο ἤρξαντο ποιῆσαι, καὶ νῦν οὐκ ἐκλείψει ἀπʼ αὐτῶν πάντα ὅσα ἂν ἐπιθῶνται ποιεῖν.
\vs{7}Δεῦτε, καὶ καταβάντες συγχέωμεν αὐτῶν ἐκεῖ τὴν γλῶσσαν, ἵνα μὴ ἀκούσωσιν ἕκαστος τὴν φωνὴν τοῦ πλησίον.
\vs{8}Καὶ διέσπειρεν αὐτοὺς Κύριος ἐκεῖθεν ἐπὶ πρόσωπον πάσης τῆς γῆς· καὶ ἐπαύσαντο οἰκοδομοῦντες τῆν πόλιν καὶ τὸν πύργον.
\vs{9}Διὰ τοῦτο ἐκλήθη τὸ ὄνομα αὐτῆς, Σύγχυσις, ὅτι ἐκεῖ συνέχεε Κύριος τὰ χείλη πάσης τῆς γῆς, καὶ ἐκεῖθεν διέσπειρεν αὐτοὺς Κύριος ἐπὶ πρόσωπον πάσης τῆς γῆς.

\vs{10}Καὶ αὗται αἱ γενέσεις Σήμ· καὶ ἦν Σὴμ υἱὸς ἑκατὸν ἐτῶν, ὅτε ἐγέννησε τὸν Ἀρφαξὰδ, δευτέρου ἔτους μετὰ τὸν κατακλυσμόν.
\vs{11}Καὶ ἔζησε Σὴμ, μετὰ τὸ γεννῆσαι αὐτὸν τὸν Ἀρφαξὰδ, ἔτη πεντακόσια, καὶ ἐγέννησεν υἱοὺς καὶ θυγατέρας, καὶ ἀπέθανε.
\vs{12}Καὶ ἔζησεν Ἀρφαξὰδ ἑκατὸν τριακονταπέντε ἔτη, καὶ ἐγέννησε τὸν Καϊνᾶν.
\vs{13}Καὶ ἔζησεν Ἀρφαξὰδ, μετὰ τὸ γεννῆσαι αὐτὸν τὸν Καϊνᾶν, ἔτη τετρακόσια, καὶ ἐγέννησεν υἱοὺς καὶ θυγατέρας, καὶ ἀπέθανε. Καὶ ἔζησε Καϊνᾶν ἑκατὸν καὶ τριάκοντα ἔτη, καὶ ἐγέννησε τὸν Σαλά· καὶ ἔξησε Καϊνᾶν, μετὰ τὸ γεννῆσαι αὐτὸν τὸν Σαλὰ, ἔτη τριακόσια τριάκοντα, καὶ ἐγέννησεν υἱοὺς καὶ θυγατέρας, καὶ ἀπέθανε.
\vs{14}Καὶ ἔζησε Σαλὰ ἑκατὸν τριάκοντα ἔτη, καὶ ἐγέννησε τὸν Ἕβερ.
\vs{15}Καὶ ἔζησε Σαλὰ μετὰ τὸ γεννῆσαι αὐτὸν τὸν Ἕβερ, τριακόσια τριάκοντα ἔτη, καὶ ἐγέννησεν υἱοὺς καὶ θυγατέρας· καὶ ἀπέθανε.
\vs{16}Καὶ ἔζησεν Ἕβερ ἑκατὸν τριάκοντα τέσσαρα ἔτη, καὶ ἐγέννησε τὸν Φαλέγ.
\vs{17}Καὶ ἔξησεν Ἕβερ, μετὰ τὸ γεννῆσαι αὐτὸν τὸν Φαλὲγ, ἔτη διακόσια ἑβδομήκοντα, καὶ ἐγέννησεν υἱοὺς καὶ θυγατέρας, καὶ ἀπέθανε.
\vs{18}Καὶ ἔζησε Φαλὲγ τριάκοντα καὶ ἑκατὸν ἔτη, καὶ ἐγέννησε τὸν Ῥαγαῦ.
\vs{19}Καὶ ἔζησε Φαλὲγ, μετὰ τὸ γεννῆσαι αὐτὸν τὸν Ῥαγαῦ, ἐννέα καὶ διακόσια ἔτη, καὶ ἐγέννησεν υἱοὺς και θυγατέρας, καὶ ἀπέθανε.
\vs{20}Καὶ ἔζησε Ῥαγαὺ ἑκατὸν τριάκοντα καὶ δύο ἔτη, καὶ ἐγέννησε τὸν Σερούχ.
\vs{21}Καὶ ἔζησε Ῥαγαῦ, μετὰ τὸ γεννῆσαι αὐτὸν τὸν Σεροὺχ, διακόσια ἑπτὰ ἔτη, καὶ ἐγέννησεν υἱοὺς καὶ θυγατέρας, καὶ ἀπέθανε.
\vs{22}Καὶ ἔζησε Σεροὺχ ἑκατὸν τριάκοντα ἔτη, καὶ ἐγέννησε τὸν Ναχώρ.
\vs{23}Καὶ ἔζησε Σεροὺχ, μετὰ τὸ γεννῆσαι αὐτὸν τὸν Ναχὼρ, ἔτη διακόσια, καὶ ἐγέννησεν υἱοὺς καὶ θυγατέρας, καὶ ἀπέθανε.
\vs{24}Καὶ ἔζησε Ναχὼρ ἔτη ἑκατὸν ἑβδομηκονταεννέα, καὶ ἐγέννησε τὸν Θάῤῥα.
\vs{25}Καὶ ἔζησε Ναχὼρ, μετὰ τὸ γεννῆσαι αὐτὸν τὸν Θάῤῥα, ἔτη ἑκατὸν εἰκοσιπὲντε, καὶ ἐγεννησεν υἱοὺς καὶ θυγατέρας, καὶ ἀπέθανε.
\vs{26}Καὶ ἔζησε Θάῤῥα ἑβδομήκοντα ἔτη, καὶ ἐγέννησε τὸν Ἄβραμ, καὶ τὸν Ναχὼρ, καὶ τὸν Ἀῤῥάν.

\vs{27}Αὗται δὲ αἱ γενέσεις Θάῤῥα· Θάῤῥα ἐγέννησε τὸν Ἅβραμ, καὶ τὸν Ναχὼρ, καὶ τὸν Ἀῤῥάν· καὶ Ἀῤῥὰν ἐγέννησε τὸν Λώτ.
\vs{28}Καὶ ἀπέθανεν Ἀῤῥὰν ἐνώπιον Θάῤῥα τοῦ πατρὸς αὐτοῦ ἐν τῇ γῇ ᾗ ἐγενήθη, ἐν τῇ χώρᾳ τῶν Χαλδαίων.
\vs{29}Καὶ ἔλαβον Ἅβραμ καὶ Ναχὼρ ἑαυτοῖς γυναῖκας· ὄνομα τῇ γυναικὶ Ἅβραμ, Σάρα, καὶ ὄνομα τῇ γυναικὶ Ναχὼρ, Μελχά, θυγάτηρ Ἀῤῥάν· καὶ πατὴρ Μελχὰ, καὶ πατὴρ Ἰεσχά.
\vs{30}Καὶ ἦν Σάρα στεῖρα, καὶ οὐκ ἐτεκνοποίει.
\vs{31}Καὶ ἔλαβε Θάῤῥα τὸν Ἅβραμ υἱὸν αὐτοῦ, καὶ τὸν Λὼτ υἱὸν Ἀῤῥάν, υἱὸν τοῦ υἱοῦ αὐτοῦ, καὶ τὴν Σάραν τὴν νύμφην αὐτοῦ, γυναῖκα Ἅβραμ τοῦ υἱοῦ αὐτοῦ, καὶ ἐξήγαγεν αὐτοὺς ἐκ τῆς χώρας τῶν Χαλδαίων, πορευθῆναι εἰς γῆν Χαναάν· καὶ ἦλθον ἕως Χαῤῥὰν, καὶ κατῴκησεν ἐκεῖ.
\vs{32}Καὶ ἐγένοντο πᾶσαι αἱ ἡμέραι Θάῤῥα ἐν γῇ Χαῤῥὰν, διακόσια πέντε ἔτη· καὶ ἀπέθανε Θάῤῥα ἐν Χαῤῥάν.

\ch{12}
Καὶ εἶπε Κύριος τῷ Ἅβραμ, ἔξελθε ἐκ τῆς γῆς σου, καὶ ἐκ τῆς συγγενείας σου, καὶ ἐκ τοῦ οἴκου τοῦ πατρός σου, καὶ δεῦρο εἰς τὴν γῆν, ἣν ἄν σοι δείξω.
\vs{2}Καὶ ποιήσω σε εἰς ἔθνος μέγα, καὶ εὐλογήσω σε, καὶ μεγαλυνῶ τὸ ὄνομά σου, καὶ ἔσῃ εὐλογημένος.
\vs{3}Καὶ εὐλογήσω τοὺς εὐλογοῦντάς σε, καὶ τοὺς καταρωμένους σε καταράσομαι, καὶ ἐνευλογηθήσονται ἐν σοὶ πᾶσαι αἱ φυλαὶ τῆς γῆς.
\vs{4}Καὶ ἐπορεύθη Ἅβραμ, καθάπερ ἐλάλησεν αὐτῷ Κύριος, καὶ ᾤχετο μετʼ αὐτοῦ Λώτ· Ἅβραμ δὲ ἦν ἐτῶν ἑβδομηκονταπέντε, ὅτε ἐξῆλθεν ἐκ Χαῤῥάν.
\vs{5}Καὶ ἔλαβεν Ἅβραμ Σάραν τὴν γυναῖκα αὐτοῦ, καὶ τὸν Λὼτ υἱὸν τοῦ ἀδελφοῦ αὐτοῦ, καὶ πάντα τὰ ὑπάρχοντα αὐτῶν ὅσα ἐκτήσαντο, καὶ πᾶσαν ψυχὴν ἣν ἐκτήσαντο, ἐκ Χαῤῥάν, καὶ ἐξήλθοσαν πορευθῆναι εἰς γῆν Χανάαν.
\vs{6}Καὶ διώδευσεν Ἅβραμ τὴν γῆν εἰς τὸ μῆκος αὐτῆς ἕως τοῦ τόπου Συχέμ, ἐπὶ τὴν δρῦν τὴν ὑψηλήν· οἱ δὲ Χαναναῖοι τότε κατῴκουν τὴν γῆν.
\vs{7}Καὶ ὤφθη Κύριος τῷ Ἅβραμ, καὶ εἶπεν αὐτῷ, τῷ σπέρματί σου δώσω τὴν γῆν ταύτην· καὶ ᾠκοδόμησεν ἐκεῖ Ἅβραμ θυσιαστήριον Κυρίῳ τῷ ὀφθέντι αὐτῷ.
\vs{8}Καὶ ἀπέστη ἐκεῖθεν εἰς τὸ ὄρος κατὰ ἀνατολὰς Βαιθήλ· καὶ ἔστησεν ἐκεῖ τὴν σκηνὴν αὐτοῦ ἐν Βαιθὴλ κατὰ θάλασσαν, καὶ Ἀγγαὶ κατὰ ἀνατολάς· καὶ ᾠκοδόμησεν ἐκεῖ θυσιαστήριον τῷ Κυρίῳ, καὶ ἐπεκαλέσατο ἐπὶ τῷ ὀνόματι Κυρίου.
\vs{9}Καὶ ἀπῇρεν Ἅβραμ, καὶ πορευθεὶς ἐστρατοπέδευσεν ἐν τῇ ἐρήμῳ.

\vs{10}Καὶ ἐγένετο λιμὸς ἐπὶ τῆς γῆς· καὶ κατέβη Ἅβραμ εἰς Αἴγυπτον παροικῆσαι ἐκεῖ, ὅτι ἐνίσχυσεν ὁ λιμὸς ἐπὶ τῆς γῆς.
\vs{11}Ἐγένετο δὲ ἡνίκα ἤγγισεν Ἅβραμ εἰσελθεῖν εἰς Αἴγυπτον, εἶπεν Ἅβραμ Σάρα τῇ γυναικὶ, γινώσκω ἐγὼ, ὅτι γυνὴ εὐπρόσωπος εἶ.
\vs{12}Ἔσται οὖν ὡς ἂν ἴδωσί σε οἱ Αἰγύπτιοι, ἐροῦσιν ὅτι γυνὴ αὐτοῦ ἐστιν αὐτὴ, καὶ ἀποκτενοῦσί με, σὲ δὲ περιποιήσονται.
\vs{13}Εἶπον οὖν, ὅτι ἀδελφὴ αὐτοῦ εἰμι, ὅπως ἄν εὖ μοι γένηται διὰ σὲ, καὶ ζήσεται ἡ ψυχή μου ἕνεκέν σου.
\vs{14}Ἐγένετο δὲ, ἡνίκα εἰσῆλθεν Ἅβραμ εἰς Αἴγυπτον, ἰδόντες οἱ Αἰγύπτιοι τὴν γυναῖκα αὐτοῦ, ὅτι καλὴ ἦν σφόδρα.
\vs{15}Καὶ ἴδον αὐτὴν οἱ ἄρχοντες Φαραὼ, καὶ ἐπῄνεσαν αὐτὴν πρὸς Φαραὼ, καὶ εἰσήγαγον αὐτὴν εἰς τὸν οἶκον Φαραώ.
\vs{16}Καὶ τῷ Ἅβραμ εὖ ἐχρήσαντο διʼ αὐτήν· καὶ ἐγένοντο αὐτῷ πρόβατα, καὶ μόσχοι, καὶ ὄνοι, καὶ παῖδες, καὶ παιδίσκαι, καὶ ἡμίονοι, καὶ κάμηλοι.
\vs{17}Καὶ ἤτασεν ὁ Θεὸς τὸν Φαραὼ ἐτασμοῖς μεγάλοις καὶ πονηροῖς, καὶ τὸν οἶκον αὐτοῦ, περὶ Σάρας τῆς γυναικὸς Ἅβραμ.
\vs{18}Καλέσας δὲ Φαραὼ τὸν Ἅβραμ, εἶπεν, τί τοῦτο ἐποίησάς μοι, ὅτι οὐκ ἀπήγγειλάς μοι, ὅτι γυνή σου ἐστίν;
\vs{19}Ἱνατί εἶπας ὅτι ἀδελφή μου ἐστίν; καὶ ἔλαβον αὐτὴν ἐμαυτῷ γυναῖκα· καὶ νῦν ἰδοὺ ἡ γυνή σου ἔναντί σου, λαβὼν ἀπότρεχε.
\vs{20}Καὶ ἐνετείλατο Φαραὼ ἀνδράσι περὶ Ἅβραμ συμπροπέμψαι αὐτὸν, καὶ τὴν γυναῖκα αὐτοῦ, καὶ πάντα ὅσα ἦν αὐτῷ.

\ch{13}
Ἀνέβη δὲ Ἅβραμ ἐξ Αἰγύπτου αὐτὸς, καὶ ἡ γυνὴ αὐτοῦ, καὶ πάντα τὰ αὐτοῦ, καὶ Λὼτ μετʼ αὐτοῦ, εἰς τὴν ἔρημον.
\vs{2}Ἅβραμ δὲ ἦν πλούσιος σφόδρα κτήνεσι, καὶ ἀργυρίῳ, καὶ χρυσίῳ.
\vs{3}Καὶ ἐπορεύθη ὅθεν ἦλθεν εἰς τὴν ἔρημον ἕως Βαιθὴλ, ἕως τοῦ τόπου οὗ ἦν ἡ σκηνὴ αὐτοῦ τὸ πρότερον, ἀνὰ μέσον Βαιθὴλ καὶ ἀνὰ μέσον Ἀγγαί,
\vs{4}εἰς τὸν τόπον τοῦ θυσιαστηρίου, οὗ ἐποίησεν ἐκεῖ τὴν ἀρχὴν, καὶ ἐπεκαλέσατο ἐκεῖ Ἅβραμ τὸ ὄνομα τοῦ Κυρίου.
\vs{5}Καὶ Λὼτ τῷ συμπορευομένῳ μετὰ Ἅβραμ ἦν πρόβατα, καὶ βόες, καὶ σκηναί.
\vs{6}Καὶ οὐκ ἐχώρει αὐτοὺς ἡ γῆ κατοικεῖν ἅμα, ὅτι ἦν τὰ ὑπάρχοντα αὐτῶν πολλά· καὶ οὐκ ἐχώρει αὐτοὺδ ἡ γῆ κατοικεῖν ἅμα.
\vs{7}Καὶ ἐγενετο μάχη ἀνὰ μέσον τῶν ποιμένων τῶν κτηνῶν τοῦ Ἅβραμ, καὶ ἀνὰ μέσον τῶν ποιμένων τῶν κτηνῶν τοῦ Λώτ· οἱ δὲ Χαναναῖοι καὶ οἱ Φερεζαῖοι τότε κατῴκουν τὴν γῆν.
\vs{8}Εἶπε δὲ Ἅβραμ τῷ Λὼτ, μὴ ἔστω μάχη ἀνὰ μέσον ἐμοῦ καὶ σοῦ, καὶ ἀνὰ μέσον τῶν ποιμένων μου καὶ ἀνὰ μέσον τῶν ποιμένων σοῦ, ὅτι ἄνθρωποι ἀδελφοὶ ἐσμὲν ἡμεῖς.
\vs{9}Οὐκ ἰδοὺ πᾶσα ἡ γῆ ἐναντίον σου ἐστί; διαχωρίσθητι ἀπʼ ἐμοῦ· εἰ σὺ εἰς ἀριστερὰ, ἐγὼ εἰς δεξιά· εἰ δὲ σὺ εἰς δεξιὰ, ἐγὼ εἰς ἀριστερά.
\vs{10}Καὶ ἐπάρας Λὼτ τοὺς ὀφθαλμοὺς αὐτοῦ, ἐπεῖδε πᾶσαν τὴν περίχωρον τοῦ Ἰορδάνου, ὅτι πᾶσα ἦν ποτιζομένη, πρὸ τοῦ καταστρέψαι τὸν Θεὸν Σόδομα καὶ Γόμοῤῥα, ὡς ὁ παράδεισος τοῦ Θεοῦ, καὶ ὡς ἡ γῆ Αἰγύπτου, ἕως ἐλθεῖν εἰς Ζόγορα.
\vs{11}Καὶ ἐξελέξατο ἑαυτῷ Λὼτ πᾶσαν τὴν περίχωρον τοῦ Ἰορδάνου· καὶ ἀπῇρε Λὼτ ἀπὸ ἀνατολῶν· καὶ διεχωρίσθησαν ἕκαστος ἀπὸ τοῦ ἀδελφοῦ αὐτοῦ.
\vs{12}Ἅβραμ δὲ κατῴκησεν ἐν γῇ Χαναάν· Λὼτ δὲ κατῴκησεν ἐν πόλει τῶν περιχώρων, καὶ ἐσκήνωσεν ἐν Σοδόμοις.
\vs{13}Οἱ δὲ ἄνθρωποι οἱ ἐν Σοδόμοις πονηροὶ καὶ ἁμαρτωλοὶ ἐναντίον τοῦ Θεοῦ σφόδρα.
\vs{14}Ὁ δὲ Θεὸς εἶπε τῷ Ἅβραμ μετὰ τὸ διαχωρισθῆναι τὸν Λὼτ ἀπʼ αὐτοῦ, ἀνάβλεψον τοῖς ὀφθαλμοῖς σου, καὶ ἴδε ἀπὸ τοῦ τόπου οὗ νῦν σὺ εἶ πρὸς βοῤῥὰν καὶ λίβα καὶ ἀνατολὰς καὶ θάλασσαν·
\vs{15}ὅτι πᾶσαν τὴν γῆν, ἣν σὺ ὁρᾷς, σοὶ δώσω αὐτὴν καὶ τῷ σπέρματί σου ἕως αἰῶνος.
\vs{16}Καὶ ποιήσω τὸ σπέρμα σου, ὡς τὴν ἄμμον τῆς γῆς· εἰ δύναταί τις ἐξαριθμῆσαι τὴν ἄμμον τῆς γῆς, καὶ τὸ σπέρμα σου ἐξαριθμηθήσεται.
\vs{17}Ἀναστὰς διόδευσον τὴν γῆν εἴς τε τὸ μῆκος αὐτῆς καὶ εἰς τὸ πλάτος· ὅτι σοι δώσω αὐτὴν καὶ τῷ σπέρματί σου εἰς τὸν αἰῶνα.
\vs{18}Καὶ ἀποσκηνώσας Ἅβραμ, ἐλθὼν κατῴκησε παρὰ τὴν δρῦν τὴν Μαμβρῆ, ἣ ἦν ἐν Χεβρὼμ, καὶ ᾠκοδόμησεν ἐκεῖ θυσιαστήριον τῷ Κυρίῳ.

\ch{14}
Ἐγένετο δὲ ἐν τῇ βασιλείᾳ τῇ Ἀμαρφὰλ βασιλέως Σενναὰρ, καὶ Ἀριὼχ βασιλέως Ἑλλασὰρ, Χοδολλογομὸρ βασιλεὺς Ἐλὰμ, καὶ Θαργὰλ βασιλεὺς ἐθνῶν,
\vs{2}ἐποίησαν πόλεμον μετὰ Βαλλὰ βασιλέως Σοδόμων, καὶ μετὰ Βαρσὰ βασιλέως Γομόῤῥας, καὶ μετὰ Σενναὰρ βασιλέως Ἀδαμὰ, καὶ μετὰ Συμοβὸρ βασιλέως Σεβωεὶμ, καὶ βασιλέως Βαλάκ· αὕτη ἐστὶ Σηγώρ.
\vs{3}Πάντες οὗτοι συνεφώνησαν ἐπὶ τὴν φάραγγα τὴν ἁλυκήν· αὕτη ἡ θάλασσα τῶν ἁλῶν.
\vs{4}Δώδεκα ἔτη αὐτοὶ ἐδούλευσαν τῷ Χοδολλογομόρ· τῷ δὲ τρισκαιδεκάτῳ ἔτει ἀπέστησαν.
\vs{5}Ἐν δὲ τῷ τεσσαρεσκαιδεκάτῳ ἔτει ἦλθε Χοδολλογομὸρ καὶ οἱ βασιλεῖς μετʼ αὐτοῦ, καὶ κατέκοψαν τοὺς γίγαντας τοὺς ἐν Ἀσταρὼθ, καὶ Καρναῒν, καὶ ἔθνη ἰσχυρὰ ἅμα αὐτοῖς, καὶ τοὺς Ὀμμαίους τοὺς ἐν Σαυῇ τῇ πόλει.
\vs{6}Καὶ τοὺς Χοῤῥαίους τοὺς ἐν τοῖς ὄρεσι Σηεὶρ, ἕως τῆς τερεβίνθου τῆς Φαρὰν, ἥ ἐστιν ἐν τῇ ἐρήμῳ.
\vs{7}Καὶ ἀναστρέψαντες ἦλθον ἐπὶ τὴν πηγὴν τῆς κρίσεως· αὕτη ἐστὶ Κάδης· καὶ κατέκοψαν πάντας τοὺς ἄρχοντας Ἀμαλὴκ, καὶ τοὺς Ἀμοῤῥαίους τοὺς κατοικοῦντας ἐν ʼΑσασονθαμὰρ
\vs{8}Ἐξῆλθε δὲ βασιλεὺς Σοδόμων, καὶ βασιλεὺς Γομόῤῥας, καὶ βασιλεὺς Ἀδαμὰ, καὶ βασιλεὺς Σεβωεὶμ, καὶ βασιλεὺς Βαλάκ· αὕτη ἐστὶ Σηγώρ· καὶ παρετάξαντο αὐτοῖς εἰς πόλεμον ἐν τῇ κοιλάδι, τῇ ἁλυκῇ,
\vs{9}πρὸς Χοδολλογομὸρ βασιλέα Ἐλὰμ, καὶ Θαπγὰλ βασιλέα ἐθνῶν, καὶ Ἀμαρφὰλ βασιλέα Σενναὰρ, καὶ Ἀριὼχ βασιλέα Ἑλλασὰρ, οἱ τέσσαρες βασιλεῖς πρὸς τοὺς πέντε.
\vs{10}Ἡ δὲ κοιλὰς ἡ ἁλυκὴ, φρέατα ἀσφάλτου· ἔφυγε δὲ βασιλεὺς Σοδόμων καὶ βασιλεὺς Γομόῤῥας, καὶ ἐνέπεσαν ἐκεῖ· οἱ δὲ καταλειφθέντες εἰς τὴν ὀρεινὴν ἔφυγον.
\vs{11}Ἔλαβον δὲ τὴν ἵππον πᾶσαν τὴν Σοδόμων καὶ Γομόῤῥας, καὶ πάντα τὰ βρώματα αὐτῶν, καὶ ἀπῆλθον.
\vs{12}Ἔλαβον δὲ καὶ τὸν Λὼτ τὸν υἱὸν τοῦ ἀδελφοῦ Ἅβραμ, καὶ τὴν ἀποσκευὴν αὐτοῦ, καὶ ἀπῴχοντο· ἦν γὰρ κατοικῶν ἐν Σοδόμοις.

\vs{13}Παραγενόμενος δὲ τῶν ἀνασωθέντων τις ἀπήγγειλεν Ἅβραμ τῷ περάτῃ· αὐτὸς δὲ κατῴκει παρὰ τῇ δρυῒ τῇ Μαμβρῇ Ἀμοῤῥαίου τοῦ ἀδελφοῦ Ἐσχὼλ, καὶ τοῦ ἀδελφοῦ Αὐνὰν, οἳ ἦσαν συνωμόται τοῦ Ἅβραμ.
\vs{14}Ἀκούσας δὲ Ἅβραμ ὅτι ᾐχμαλώτευται Λὼτ ὁ ἀδελφοῦς αὐτοῦ, ἠρίθμησε τοὺς ἰδίους οἰκογενεῖς αὐτοῦ τριακοσίους δέκα καὶ ὀκτώ· καὶ κατεδίωξεν ὀπίσω αὐτῶν ἕως Δάν.
\vs{15}Καὶ ἐπέπεσεν ἐπʼ αὐτοὺς τὴν νύκτα αὐτὸς, καὶ οἱ παῖδες αὐτοῦ, καὶ ἐπάταξεν αὐτοὺς, καὶ κατεδίωξεν αὐτοὺς ἕως Χοβὰ, ἥ ἐστιν ἐν ἀριστερᾷ Δαμασκοῦ.
\vs{16}Καὶ ἀπέστρεψε πᾶσαν τὴν ἵππον Σοδόμων· καὶ Λὼτ τὸν ἀδελφιδοῦν αὐτοῦ ἀπέστρεψε, καὶ πάντα τὰ ὑπάρχοντα αὐτοῦ, καὶ τὰς γυναῖκας, καὶ τὸν λαόν.
\vs{17}Ἐξῆλθε δὲ βασιλεὺς Σοδόμων εἰς συνάντησιν αὐτῷ, μετὰ τὸ ὑποστρέψαι αὐτὸν ἀπὸ τῆς κοπῆς τοῦ Χοδολλογομὸρ, καὶ τῶν βασιλέων τῶν μετʼ αὐτοῦ εἰς τὴν κοιλάδα τοῦ Σαβύ· τοῦτο ἦν τὸ πεδίον τῶν βασιλέων.

\vs{18}Καὶ Μελχισεδὲκ βασιλεὺς Σαλὴμ ἐξήνεγκεν ἄρτους καὶ οἶνον· ἦν δὲ ἱερεὺς τοῦ Θεοῦ τοῦ ὑψίστου.
\vs{19}Καὶ εὐλόγησε τὸν Ἅβραμ, καὶ εἶπεν, εὐλογημένος Ἅβραμ τῷ Θεῷ τῷ ὑψίστῳ, ὃς ἔκτισε τὸν οὐρανὸν καὶ τὴν γῆν.
\vs{20}Καὶ εὐλογητὸς ὁ Θεὸς ὁ ὕψιστος, ὃς παρέδωκε τοὺς ἐχθρούς σου ὑποχειρίους σοι· καὶ ἔδωκεν αὐτῷ Ἅβραμ δεκάτην ἀπὸ πάντων.
\vs{21}Εἶπε δὲ βασιλεὺς Σοδόμων πρὸς Ἅβραμ, δός μοι τοὺς ἄνδρας, τὴν δὲ ἵππον λάβε σεαυτῷ.
\vs{22}Εἶπε δὲ Ἅβραμ πρὸς τὸν βασιλέα Σοδόμων, ἐκτενῶ τὴν χεῖρά μου πρὸς Κύπιον τὸν Θεὸν τὸν ὕψιστον, ὃς ἔκτισε τὸν οὐρανὸν καὶ τὴν γῆν,
\vs{23}εἰ ἀπὸ σπαρτίου ἕως σφυρωτῆρος ὑποδήματος λήψομαι ἀπὸ πάντων τῶν σῶν, ἵνα μὴ εἴπῃς, ὅτι ἐγὼ ἐπλούτισα τὸν Ἅβραμ.
\vs{24}Πλὴν ὧν ἔφαγον οἱ νεανίσκοι, καὶ τῆς μερίδος τῶν ἀνδρῶν τῶν συμπορευθέντων μετʼ ἐμοῦ Ἐσχὼλ, Αὐνᾶν, Μαμβρῆ· οὗτοι λήψονται μερίδα.

\ch{15}
Μετὰ δὲ τὰ ῥήματα ταῦτα ἐγενήθη ῥῆμα Κυρίου πρὸς Ἅβραμ ἐν ὁράματι, λέγων, μὴ φοβοῦ Ἅβραμ· ἐγὼ ὑπερασπίζω σου· ὁ μισθός σου πολὺς ἔσται σφόδρα.
\vs{2}Δέγει δὲ Ἅβραμ, Δέσποτα Κύριε, τί μοι δώσεις; ἐγὼ δὲ ἀπολύομαι ἄτεκνος· ὁ δὲ υἱὸς Μασὲκ τῆς οἰκογενοῦς μου, οὗτος Δαμασκὸς Ἐλιέζερ.
\vs{3}Καὶ εἶπεν Ἅβραμ, ἐπειδὴ ἐμοὶ οὐκ ἔδωκας σπέρμα, ὁ δὲ οἰκογενής μου κληρονομήσει με.
\vs{4}Καὶ εὐθὺς φωνὴ Κυρίου ἐγένετο πρὸς αὐτὸν, λέγουσα, οὐ κληρονομήσει σε οὗτος· ἀλλʼ ὃς ἐξελεύσεται ἐκ σοῦ, οὗτος κληρονομήσει σε.
\vs{5}Ἐξήγαγε δὲ αὐτὸν ἔξω, καὶ εἶπεν αὐτῷ, ἀνάβλεψον δὴ εἰς τὸν οὐρανὸν, καὶ ἀρίθμησον τοὺς ἀστέρας, εἰ δυνήσῃ ἐξαριθμῆσαι αὐτούς· καὶ εἶπεν, οὕτως ἔσται τὸ σπέρμα σου.
\vs{6}Καὶ ἐπίστευσεν Ἅβραμ τῷ Θεῷ, καὶ ἐλογίσθη αὐτῷ εἰς δικαιοσύνην.
\vs{7}Εἶπε δὲ πρὸς αὐτὸν, ἐγὼ ὁ Θεὸς ὁ ἐξαγαγών σε ἐκ χώρας Χαλδαίων, ὥστε δοῦναί σοι τὴν γῆν ταύτην κληρονομῆσαι.
\vs{8}Εἶπε δέ, Δέσποτα Κύριε, κατὰ τί γνώσομαι, ὅτι κληρονομήσω αὐτήν;
\vs{9}Εἶπε δὲ αὐτῷ, λάβε μοι δάμαλιν τριετίζουσαν, καὶ αἶγα τριετίζουσαν, καὶ κριὸν τριετίζοντα, καὶ τρυγόνα, καὶ περιστεράν.
\vs{10}Ἔλαβε δὲ αὐτῷ πάντα ταῦτα, καὶ διεῖλεν αὐτὰ μέσα, καὶ ἔθηκεν αὐτὰ ἀντιπρόσωπα ἀλλήλοις· τὰ δὲ ὄρνεα οὐ διεῖλε.
\vs{11}Κατέβη δὲ ὄρνεα ἐπὶ τὰ σώματα, ἐπὶ τὰ διχοτομήματα αὐτῶν· καὶ συνεκάθισεν αὐτοῖς Ἅβραμ.
\vs{12}Περὶ δὲ ἡλίου δυσμὰς ἔκστασις ἐπέπεσε τῷ Ἅβραμ, καὶ ἰδοὺ φόβος σκοτεινὸς μέγας ἐπιπίπτει αὐτῷ.
\vs{13}Καὶ ἐῤῥέθη πρὸς Ἅβραμ· γινώσκων γνώσῃ ὅτι πάροικον ἔσται τὸ σπέρμα σου ἐν γῇ οὐκ ἰδίᾳ, καὶ δουλώσουσιν αὐτοὺς, καὶ κακώσουσιν αὐτοὺς, καὶ ταπεινώσουσιν αὐτοὺς, τετρακόσια ἔτη.
\vs{14}Τὸ δὲ ἔθνος, ᾧ ἐὰν δουλεύσωσι, κρινῶ ἐγώ· μετὰ δὲ ταῦτα, ἐξελεύσονται ὧδε μετὰ ἀποσκευῆς πολλῆς.
\vs{15}Σὺ δὲ ἀπελεύσῃ πρὸς τοὺς πατέρας σου ἐν εἰρήνῃ, τραφεὶς ἐν γήρᾳ καλῷ.
\vs{16}Τετάρτῃ δὲ γενεᾷ ἀποστραφήσονται ὧδε· οὔπω γὰρ ἀναπεπλήρωνται αἱ ἁμαρτίαι τῶν Ἀμοῤῥαίων ἕως τοῦ νῦν.
\vs{17}Ἐπεὶ δὲ ὁ ἥλιος ἐγένετο πρὸς δυσμὰς, φλὸξ ἐγένετο· καὶ ἰδοὺ κλίβανος καπνιζόμενος καὶ λαμπάδες πυρός, αἳ διῆλθον ἀνὰ μέσον τῶν διχοτομημάτων τούτων.
\vs{18}Ἐν τῇ ἡμέρᾳ ἐκείνῃ διέθετο Κύριος τῷ Ἅβραμ διαθήκην, λέγων, τῷ σπέρματί σου δώσω τὴν γῆν ταύτην, ἀπὸ τοῦ ποταμοῦ Αἰγύπτου ἕως τοῦ ποταμοῦ τοῦ μεγάλου Εὐφράτου·
\vs{19}Τοὺς Κεναίους, καὶ τοὺς Κενεζαίους, καὶ τοὺς Κεδμωναίους,
\vs{20}καὶ τοὺς Χετταίους, καὶ τοὺς Φερεζαίους, καὶ τοὺς ʼΡαφαεὶν,
\vs{21}καὶ τοὺς Ἀμοῤῥαίους, καὶ τοὺς Χαναναίους, καὶ τοὺς Εὐαίους, καὶ τοὺς Γεργεσαίους, καὶ τοὺς Ἰεβουσαίους.

\ch{16}
Σάρα δὲ ἡ γυνὴ Ἅβραμ οὐκ ἔτικτεν αὐτῷ· ἦν δὲ αὐτῇ παιδίσκη Αἰγυπτία, ᾗ ὄνομα Ἄγαρ.
\vs{2}Εἶπε δὲ Σάρα πρὸς Ἅβραμ, ἰδοὺ συνέκλεισέ με Κύριος τοῦ μὴ τίκτειν· εἴσελθε οὖν πρὸς τὴν παιδίσκην μου, ἵνα τεκνοποιήσωμαι ἐξ αὐτῆς· ὑπήκουσελ δὲ Ἅβραμ τῆς φωνῆς Σάρας.
\vs{3}Καὶ λαβοῦσα Σάρα ἡ γυνὴ Ἅβραμ Ἄγαρ τὴν Αἰγυπτίαν τὴν ἑαυτῆς παιδίσκην, μετὰ δέκα ἔτη τοῦ οἰκῆσαι Ἅβραμ ἐν γῇ Χαναὰν, ἔδωκεν αὐτὴν τῷ Ἅβραμ ἀνδρὶ αὐτῆς αὐτῷ γυναῖκα.
\vs{4}Καὶ εἰσῆλθε πρὸς Ἄγαρ, καὶ συνέλαβε· καὶ εἶδεν ὅτι ἐν γαστρὶ ἔχει, καὶ ἠτιμάσθη ἡ κυρία ἐναντίον αὐτῆς.
\vs{5}Εἶπε δὲ Σάρα πρὸς Ἅβραμ, ἀδικοῦμαι ἐκ σοῦ· ἐγὼ δέδωκα τὴν παιδίσκην μου εἰς τὸν κόλπον σου, ἰδοῦσα δὲ ὅτι ἐν γαστρὶ ἔχει, ἠτιμάσθην ἐναντίον αὐτῆς. κρίναι ὁ Θεὸς ἀνὰ μέσον ἐμοῦ καὶ σου.
\vs{6}Εἶπε δὲ Ἅβραμ πρὸς Σάραν, ἰδοὺ ἡ παιδίσκη σου ἐν ταῖς χερσί σου, χρῶ αὐτῇ ὡς ἄν σοι ἀρεστὸν ᾖ. καὶ ἐκάκωσεν αὐτὴν Σάρα, καὶ ἀπέδρα ἀπὸ προσώπου αὐτῆς.

\vs{7}Εὗρε δὲ αὐτὴν ἄγγελος Κυρίου ἐπὶ τῆς πηγῆς τοῦ ὕδατος ἐν τῇ ἐρήμῳ, ἐπὶ τῆς πηγῆς ἐν τῇ ὁδῷ Σούρ.
\vs{8}Καὶ εἶπεν αὐτῇ ὁ ἄγγελος Κυρίου, Ἄγαρ παιδίσκη Σάρας, πόθεν ἔρχῃ; καὶ ποῦ πορεύῃ; καὶ εἶπεν· ἀπὸ προσώπου Σάρας τῆς κυρίας μου ἐγὼ ἀποδιδράσκω.
\vs{9}Εἶπε δὲ αὐτῇ ὁ ἄγγελος Κυρίου, ἀποστράφηθι πρὸς τὴν κυρίαν σου, καὶ ταπεινώθητι ὑπὸ τὰς χεῖρας αὐτῆς.
\vs{10}Καὶ εἶπεν αὐτῇ ὁ ἄγγελος Κυρίου, πληθύνων πληθυνῶ τὸ σπέρμα σου, καὶ οὐκ ἀριθμηθήσεται ὑπὸ τοῦ πλήθους.
\vs{11}Καὶ εἶπεν αὐτῇ ὁ ἄγγελος Κυρίου, ἰδοὺ σὺ ἐν γαστρὶ ἔχεις, καὶ τέξῃ υἱὸν, καὶ καλέσεις τὸ ὄνομα αὐτοῦ Ἰσμαὴλ, ὅτι ἐπήκουσε Κύριος τῇ ταπεινώσει σου.
\vs{12}Οὗτος ἔσται ἄγροικος ἄνθρωπος· αἱ χεῖρες αὐτοῦ ἐπὶ πάντας, καὶ αἱ χεῖρες πάντων ἐπʼ αὐτόν· καὶ κατὰ πρόσωπον πάντων τῶν ἀδελφῶν αὐτοῦ κατοικήσει.
\vs{13}Καὶ ἐκάλεσε τὸ ὄνομα Κυρίου τοῦ λαλοῦντος πρὸς αὐτὴν, σὺ ὁ Θεὸς ὁ ἐτιδών με· ὅτι εἶπε, καὶ γὰρ ἐνώπιον εἶδον ὀφθέντα μοι.
\vs{14}Ἕνεκεν τούτου ἐκάλεσε τὸ φρέαρ, φρέαρ οὗ ἐνώπιον εἶδον· ἰδοὺ ἀνὰ μέσον Κάδης καὶ ἀνὰ μέσον Βαράδ.
\vs{15}Καὶ ἔτεκεν Ἄγαρ τῷ Ἅβραμ υἱὸν, καὶ ἐκάλεσεν Ἅβραμ τὸ ὄνομα τοῦ υἱοῦ αὐτοῦ, ὃν ἔτεκεν αὐτῷ Ἄγαρ, Ἰσμαήλ.
\vs{16}Ἅβραμ δὲ ἦν ἐτῶν ὀγδοηκονταὲξ, ἡνίκα ἔτεκεν Ἄγαρ τῷ Ἅβραμ τὸν Ἰσμαήλ.

\ch{17}
Ἐγένετο δὲ Ἅβραμ ἐτῶν ἐννενηκονταεννέα. Καὶ ὤφθη Κύριος τῷ Ἅβραμ, καὶ εἶπεν αὐτῷ, ἐγώ εἰμι ὁ Θεός σου· εὐαρέστει ἐνώπιον ἐμοῦ, καὶ γίνου ἄμεμπτος.
\vs{2}Καὶ θήσομαι τὴν διαθήκην μου ἀνὰ μέσον ἐμοῦ, καὶ ἀνὰ μέσον σου, καὶ πληθυνῶ σε σφόδρα.
\vs{3}Καὶ ἔπεσεν Ἅβραμ ἐπὶ πρόσωπον αὐτοῦ.
\vs{4}Καὶ ἐλάλησεν αὐτῷ ὁ Θεὸς, λέγων, Καὶ ἐγὼ ἰδοὺ ἡ διαθήκη μου μετὰ σοῦ· καὶ ἔσῃ πατὴρ πλήθους ἐθνῶν.
\vs{5}Καὶ οὐ κληθήσεται ἔτι τὸ ὄνομά σου Ἅβραμ, ἀλλʼ ἔσται τὸ ὄνομά σου Ἁβραὰμ, ὅτι πατέρα πολλῶν ἐθνῶν τέθεικά σε.
\vs{6}Καὶ αὐξανῶ σε σφόδρα σφόδρα, καὶ θήσω σε εἰς ἔθνη· καὶ βασιλεῖς ἐκ σοῦ ἐξελεύσονται.
\vs{7}Καὶ στήσω τὴν διαθήκην μου ἀνὰ μέσον σου, καὶ ἀνὰ μέσον τοῦ σπέρματός σου μετὰ σὲ εἰς τὰς γενεὰς αὐτῶν, εἰς διαθήκην αἰώνιον εἶναί σου Θεὸς, καὶ τοῦ σπέρματός σου μετὰ σέ.
\vs{8}Καὶ δώσω σοι καὶ τῷ σπέρματί σου μετὰ σὲ τὴν γῆν, ἣν παροικεῖς, πᾶσαν τὴν γῆν Χαναὰν, εἰς κατάσχεσιν αἰώνιον· καὶ ἔσομαι αὐτοῖς εἰς Θεόν.
\vs{9}Καὶ εἶπεν ὁ Θεὸς πρὸς Ἁβραὰμ, σὺ δὲ τὴν. διαθήκην μου διατηρήσεις, σὺ καὶ τὸ σπέρμα σου μετὰ σὲ εἰς τὰς γενεὰς αὐτῶν.
\vs{10}Καὶ αὕτη ἡ διαθήκη, ἣν διατηρήσεις, ἀνὰ μέσον ἐμοῦ καὶ ὑμῶν, καὶ ἀνὰ μέσον τοῦ σπέρματός σου μετὰ σὲ εἰς τὰς γενεὰς αὐτῶν· περιτμηθήσεται ὑμῶν πᾶν ἀρσενικόν.
\vs{11}Καὶ περιτμηθήσεσθε τὴν σάρκα τῆς ἀκροβυστίας ὑμῶν, καὶ ἔσται εἰς σημεῖον διαθήκης ἀνὰ μέσον ἐμοῦ καὶ ὑμῶν.
\vs{12}Καὶ παιδίον ὀκτὼ ἡμερῶν περιτμηθήσεται ὑμῖν, πᾶν ἀρσενικὸν εἰς τὰς γενεὰς ὑμῶν· καὶ οἰκογενὴς καὶ ὁ ἀργυρώνητος ἀπὸ παντὸς υἱοῦ ἀλλοτρίου, ὃς οὐκ ἔστιν ἐκ τοῦ σπέρματός σου,
\vs{13}Περιτομῇ περιτμηθήσεται ὁ οἰκογενὴς τῆς οἰκίας σου, καὶ ὁ ἀργυρώνητος· καὶ ἔσται ἡ διαθήκη μου ἐπὶ τῆς σαρκὸς ὑμῶν εἰς διαθήκην αἰώνιον.
\vs{14}Καὶ ἀπερίτμητος ἄρσην, ὃς οὐ περιτμηθήσεται τὴν σάρκα τῆς ἀκροβυστίας αὐτοῦ τῇ ἡμέρᾳ τῇ ὀγδόῃ, ἐξολοθρευθήσεται ἡ ψυχὴ ἐκείνη ἐκ τοῦ γένους αὐτῆς, ὅτι τὴν διαθήκην μου διεσκέδασε.
\vs{15}Καὶ εἶπεν ὁ Θεὸς τῷ Ἁβραὰμ, Σάρα ἡ γυνή σου, οὐ κληθήσεται τὸ ὄνομα αὐτῆς Σάρα, Σάῤῥα ἔσται τὸ ὄνομα αὐτῆς.
\vs{16}Εὐλογήσω δὲ αὐτὴν, καὶ δώσω σοι ἐξ αὐτῆς τέκνον, καὶ εὐλογήσω αὐτὸ, καὶ ἔσται εἰς ἔθνη, καὶ βασιλεῖς ἐθνῶν ἐξ αὐτοῦ ἔσονται.
\vs{17}Καὶ ἔπεσεν Ἁβραὰμ ἐπὶ πρόσωπον αὐτοῦ, καὶ ἐγέλασε· καὶ εἶπεν ἐν τῇ διανοίᾳ αὐτοῦ, λέγων, εἰ τῷ ἑκατονταετεῖ γενήσεται υἱός; καὶ εἰ ἡ Σάῤῥα ἐννενήκοντα ἐτῶν τέξεται;
\vs{18}Εἶπε δὲ Ἁβραὰμ πρὸς τὸν Θεόν· Ἰσμαὴλ οὗτος ζήτω ἐναντίον σου.
\vs{19}Εἶπε δὲ ὁ Θεὸς πρὸς Ἁβραὰμ, ναί· ἰδοὺ Σάῤῥα ἡ γυνή σου τέξεταί σοι υἱὸν, καὶ καλέσεις τὸ ὄνομα αὐτοῦ Ἰσαάκ· καὶ στήσω τὴν διαθήκην μου πρὸς αὐτὸν, εἰς διαθήκην αἰώνιον, εἶναι αὐτῷ Θεὸς καὶ τῷ σπέρματι αὐτοῦ μετʼ αὐτόν.
\vs{20}Περὶ δὲ Ἰσμαὴλ ἰδοὺ ἐπήκουσά σου· καὶ ἰδοὺ εὐλόγηκα αὐτὸν, καὶ αὐξανῶ αὐτὸν, καὶ πληθυνῶ αὐτὸν σφόδρα δώδεκα ἔθνη γεννήσει, καὶ δώσω αὐτὸν εἰς ἔθνος μέγα.
\vs{21}Τὴν δὲ διαθήκην μου στήσω πρὸς Ἰσαὰκ, ὃν τέξεταί σοι Σάῤῥα εἰς τὸν καιρὸν τοῦτον, ἐν τῷ ἐνιαυτῷ τῷ ἑτέρῳ.
\vs{22}Συνετέλεσε δὲ λαλῶν πρὸς αὐτὸν, καὶ ἀνέβη ὁ Θεὸς ἀπὸ Ἁβραάμ.

\vs{23}Καὶ ἔλαβεν Ἁβραὰμ Ἰσμαὴλ τὸν υἱὸν ἑαυτοῦ, καὶ πάντας τοὺς οἰκογενεῖς αὐτοῦ, καὶ πάντας τοὺς ἀργυρωνήτους, καὶ πᾶν ἄρσεν τῶν ἀνδρῶν τῶν ἐν τῷ οἴκῳ Ἁβραὰμ, καὶ περιέτεμε τὰς ἀκροβυστίας αὐτῶν, ἐν τῷ καιρῷ τῆς ἡμέρας ἐκείνης, καθὰ ἐλάλησεν αὐτῷ ὁ Θεός.
\vs{24}Ἁβραὰμ δὲ ἐννενηκονταεννέα ἦν ἐτῶν, ἡνίκα περιετέμετο τὴν σάρκα τῆς ἀκροβυστίας αὐτοῦ.
\vs{25}Ἰσμαὴλ δὲ ὁ υἱὸς αὐτοῦ ἦν ἐτῶν δεκατριῶν, ἡνίκα περιετέμετο τὴν σάρκα τῆς ἀκροβυστίας αὐτοῦ.
\vs{26}Ἐν δὲ τῷ καιρῷ τῆς ἡμέρας ἐκείνης, περιετμήθη Ἁβραὰμ, καὶ Ἰσμαὴλ ὁ υἱὸς αὐτοῦ,
\vs{27}καὶ πάντες οἱ ἄνδρες τοῦ οἴκου αὐτοῦ, καὶ οἱ οἰκογενεῖς αὐτοῦ, καὶ οἱ ἀργυρώνητοι ἐξ ἀλλογενῶν ἐθνῶν.

\ch{18}
Ὤφθη δὲ αὐτῷ ὁ Θεὸς πρὸς τῇ δρυῒ τῇ Μαμβρῇ, καθημένου αὐτοῦ ἐπὶ τῆς θύρας τῆς σκηνῆς αὐτοῦ μεσημβρίας.
\vs{2}Ἀναβλέψας δὲ τοῖς ὀφθαλμοῖς αὐτοῦ εἶδε, καὶ ἰδοὺ τρεῖς ἄνδρες εἱστήκεισαν ἐπάνω αὐτοῦ· καὶ ἰδὼν, προσέδραμεν εἰς συνάντησιν αὐτοῖς ἀπὸ τῆς θύρας τῆς σκηνῆς αὐτοῦ, καὶ προσεκύνησεν ἐπὶ τὴν γῆν.
\vs{3}Καὶ εἶπε, Κύριε, εἰ ἄρα εὗρον χάριν ἐναντίον σου, μὴ παρέλθῃς τὸν παῖδά σου.
\vs{4}Ληφθήτω δὴ ὕδωρ, καὶ νιψάτωσαν τοὺς πόδας ὑμῶν, καὶ καταψύξατε ὑπὸ τὸ δένδρον.
\vs{5}Καὶ λήψομαι ἄρτον, καὶ φάγεσθε. Καὶ μετὰ τοῦτο παρελεύσεσθε εἰς τὴν ὁδὸν ὑμῶν, οὗ ἕνεκεν ἐξεκλίνατε πρὸς τὸν παῖδα ὑμῶν. Καὶ εἶπεν, οὕτω ποίησον, καθὼς εἴρηκας.
\vs{6}Καὶ ἔσπευσεν Ἁβραὰμ ἐπὶ τὴν σκηνὴν πρὸς Σάῤῥαν, καὶ εἶπεν αὐτῇ, σπεῦσον, καὶ φύρασον τρία μέτρα σεμιδάλεως, καὶ ποίησον ἐγκρυφίας.
\vs{7}Καὶ εἰς τὰς βόας ἔδραμεν Ἁβραὰμ, καὶ ἔλαβεν ἁπαλὸν μοσχάριον καὶ καλὸν, καὶ ἔδωκε τῷ παιδὶ, καὶ ἐτάχυνε τοῦ ποιῆσαι αὐτό.
\vs{8}Ἔλαβε δὲ βούτυρον, καὶ γάλα, καὶ τὸ μοσχάριον ὃ ἐποίησε, καὶ παρέθηκεν αὐτοῖς, καὶ ἔφαγον· αὐτὸς δὲ παρειστήκει αὐτοῖς ὑπὸ τὸ δένδρον.

\vs{9}Εἶπε δὲ πρὸς αὐτὸν, ποῦ Σάῤῥα ἡ γυνή σου; ὁ δὲ ἀποκριθεὶς εἶπεν, ἰδοὺ ἐν τῇ σκηνῇ.
\vs{10}Εἶπε δὲ, ἐπαναστρέφων ἥξω πρὸς σὲ κατὰ τὸν καιρὸν τοῦτον εἰς ὥρας, καὶ ἕξει υἱὸν Σάῤῥα ἡ γυνή σου. Σάῤῥα δὲ ἤκουσε πρὸς τῇ θύρᾳ τῆς σκηνῆς οὖσα ὄπισθεν αὐτοῦ.
\vs{11}Ἁβραὰμ δὲ καὶ Σάῤῥα πρεσβύτεροι προβεβηκότες ἡμερῶν· ἐξέλιπε δὲ τῇ Σάῤῥᾳ γίνεσθαι τὰ γυναικεια.
\vs{12}Ἐγέλασε δὲ Σάῤῥα ἐν ἑαυτῇ λέγουσα, οὔπω μέν μοι γέγονεν ἕως τοῦ νῦν· ὁ δὲ κύριός μου πρεσβύτερος.
\vs{13}Καὶ εἶπε Κύριος πρὸς Ἁβραὰμ, τί ὅτι ἐγέλασε Σάῤῥα ἐν ἑαυτῇ, λέγουσα, ἆρά γε ἀληθῶς τέξομαι; ἐγὼ δὲ γεγήρακα.
\vs{14}Μὴ ἀδυνατήσει παρὰ τῷ Θεῷ ῥῆμα; εἰς τὸν καιρὸν τοῦτον ἀναστρέψω πρὸς σὲ εἰς ὥρας, καὶ ἔσται τῇ Σάῤῥᾳ υἱός.
\vs{15}Ἠρνήσατο δὲ Σάῤῥα, λέγουσα, οὐκ ἐγέλασα· ἐφοβήθη γάρ. Καὶ εἶπεν αὐτῇ, οὐχὶ, ἀλλὰ ἐγέλασας.

\vs{16}Ἐξαναστάντες δὲ ἐκεῖθεν οἱ ἄνδρες κατέβλεψαν ἐπὶ πρόσωπον Σοδόμων καὶ Γομόῤῥας. Ἁβραὰμ δὲ συνεπορεύετο μετʼ αὐτῶν, συμπροπέμπων αὐτούς.
\vs{17}Ὁ δὲ Κύριος εἶπε, οὐ μὴ κρύψω ἐγὼ ἀπὸ Ἁβραὰμ τοῦ παιδός μου ἃ ἐγὼ ποιῶ.
\vs{18}Ἁβραὰμ δὲ γινόμενος ἔσται εἰς ἔθνος μέγα καὶ πολὺ, καὶ ἐνευλογηθήσονται ἐν αὐτῷ πάντα τὰ ἔθνη τῆς γῆς.
\vs{19}Ἤδειν γὰρ ὅτι συντάξει τοῖς υἱοῖς αὐτοῦ, καὶ τῷ οἴκῳ αὐτοῦ μετʼ αὐτὸν, καὶ φυλάξουσι τὰς ὁδοὺς Κυρίου, ποιεῖν δικαιοσύνην καὶ κρίσιν, ὅπως ἂν ἐπαγάγῃ Κύριος ἐπὶ Ἁβραὰμ πάντα ὅσα ἐλάλησε πρὸς αὐτόν.
\vs{20}Εἶπε δὲ Κύριος, κραυγὴ Σοδόμων καὶ Γομόῤῥας πεπλήθυνται πρὸς μὲ, καὶ αἱ ἁμαρτίαι αὐτῶν μεγάλαι σφόδρα.
\vs{21}Καταβὰς οὖν ὄψομαι, εἰ κατὰ τὴν κραυγὴν αὐτῶν τὴν ἐρχομενεην πρὸς μὲ, συντελοῦνται· εἰ δὲ μὴ, ἵνα γνῶ.
\vs{22}Καὶ ἀποστρέψαντες ἐκεῖθεν οἱ ἄνδρες, ἦλθον εἰς Σόδομα· Ἁβραὰμ δὲ ἔτι ἦν ἑστηκὼς ἐναντίον Κυρίου.
\vs{23}Καὶ ἐγγίσας Ἁβραὰμ, εἶπε, μὴ συναπολέσῃς δίκαιον μετὰ ἀσεβοῦς· καὶ ἔσται ὁ δίκαιος ὡς ὁ ἀσεβής.
\vs{24}Ἐὰν ὦσι πεντήκοντα δίκαιοι ἐν τῇ πόλει, ἀπολεῖς αὐτούς; οὐκ ἀνήσεις πάντα τὸν τόπον ἕνεκεν τῶν πεντήκοντα δικαίων, ἐὰν ὦσιν ἐν αὐτῇ;
\vs{25}Μηδαμῶς σὺ ποιήσεις ὡς τὸ ῥῆμα τοῦτο, τοῦ ἀποκτεῖναι δίκαιον μετὰ ἀσεβοῦς, καὶ ἔσται ὁ δίκαιος ὡς ὁ ἀσεβής· μηδαμῶς· ὁ κρίνων πᾶσαν τὴν γῆν, οὐ ποιήσεις κρίσιν;
\vs{26}Εἶπε δὲ Κύριος, ἐὰν ὦσιν ἐν Σοδόμοις πεντήκοντα δίκαιοι ἐν τῇ πόλει, ἀφήσω ὅλην τὴν πόλιν, καὶ πάντα τὸν τόπον διʼ αὐτούς.
\vs{27}Καὶ ἀποκριθεὶς Ἁβραὰμ εἶπε, νῦν ἠρξάμην λαλῆσαι πρὸς τὸν Κύριόν μου· ἐγὼ δὲ εἰμὶ γῆ καὶ σποδός.
\vs{28}Ἐὰν δὲ ἐλαττονωθῶσιν οἱ πεντήκοντα δίκαιοι εἰς τεσσαρακονταπέντε, ἀπολεῖς ἕνεκεν τῶν πέντε πᾶσαν τὴν πόλιν; καὶ εἶπεν, οὐ μὴ ἀπολέσω, ἐὰν εὕρω ἐκεῖ τεσσαρακονταπέντε.
\vs{29}Καὶ προσέθηκεν ἔτι λαλῆσαι πρὸς αὐτὸν, καὶ εἶπεν, ἐὰν δὲ εὑρεθῶσιν ἐκεῖ τεσσαράκοντα· καὶ εἶπεν, οὐ μὴ ἀπολέσω ἕνεκεν τῶν τεσσαράκοντα.
\vs{30}Καὶ εἶπε, μή τι Κύριε ἐὰν λαλήσω; ἐὰν δὲ εὑρεθῶσιν ἐκεῖ τριάκοντα; καὶ εἶπεν, οὐ μὴ ἀπολέσω ἕνεκεν τῶν τριάκοντα.
\vs{31}Καὶ εἶπεν, ἐπειδὴ ἔχω λαλῆσαι πρὸς τὸν Κύριον, ἐὰν δὲ εὑρεθῶσιν ἐκεῖ εἴκοσι; καὶ εἶπεν, οὐ μὴ ἀπολέσω, ἐὰν εὕρω ἐκεῖ εἴκοσι.
\vs{32}Καὶ εἶπε, μή τι Κύριε ἐὰν λαλήσω ἔτι ἅπαξ; ἐὰν δὲ εὑρεθῶσιν ἐκεῖ δέκα; καὶ εἶπεν, οὐ μὴ ἀπολέσω ἕνεκεν τῶν δέκα.
\vs{33}Ἀπῆλθε δὲ ὁ Κύριος, ὡς ἐπαύσατο λαλῶν τῷ Ἁβραάμ· καὶ Ἁβραὰμ ἀπέστρεψεν εἰς τὸν τόπον αὐτοῦ.

\ch{19}
Ἦλθον δε οἱ δύο ἄγγελοι εἰς Σόδομα ἑσπέρας. Λὼτ δὲ ἐκάθητο παρὰ τὴν πύλην Σοδόμων· ἰδὸν δὲ Λὼτ, ἐξανέστη εἰς συνάντησιν αὐτοῖς, καὶ προσεκύνησε τῷ προσώπῳ ἐπὶ τὴν γῆν.
\vs{2}Καὶ εἶπεν, ἰδοὺ, Κύριοι, ἐκκλίνατε εἰς τὸν οἶκον τοῦ παιδὸς ὑμῶν, καὶ καταλύσατε, καὶ νίψασθε τοὺς πόδας ὑμῶν, καὶ ὀρθρίσαντες ἀπελεύσεσθε εἰς τὴν ὁδὸν ὑμῶν. Καὶ εἶπαν, οὐχὶ, ἀλλʼ ἐν τῇ πλατείᾳ καταλύσομεν.
\vs{3}Καὶ κατεβιάσατο αὐτοὺς, καὶ ἐξέκλιναν πρὸς αὐτὸν, καὶ εἰσῆλθον εἰς τὸν οἶκον αὐτοῦ· καὶ ἐποίησεν αὐτοῖς πότον, καὶ ἀζύμους ἔπεψεν αὐτοῖς, καὶ ἔφαγον.
\vs{4}Πρὸ τοῦ κοιμηθῆναι δὲ, οἱ ἄνδρες τῆς πόλεως, οἱ Σοδομῖται περιεκύκλωσαν τὴν οἰκίαν, ἀπὸ νεανίσκου ἕως πρεσβυτέρου, ἅπας ὁ λαὸς ἅμα.
\vs{5}Καὶ ἐξεκαλοῦντο τὸν Λὼτ, καὶ ἔλεγον πρὸς αὐτὸν, ποῦ εἰσιν οἱ ἄνδρες οἱ εἰσελθόντες πρὸς σὲ τὴν νύκτα; ἐξάγαγε αὐτοὺς πρὸς ἡμᾶς, ἵνα συγγενώμεθα αὐτοῖς.
\vs{6}Ἐξῆλθε δὲ Λὼτ πρὸς αὐτοὺς πρὸς τὸ πρόθυρον, τὴν δὲ θύραν προσέῳξεν ὀπίσω αὐτοῦ.
\vs{7}Εἶπε δὲ πρὸς αὐτοὺς, μηδαμῶς ἀδελφοὶ μὴ πονηρεύσησθε.
\vs{8}Εἰσὶ δέ μοι δύο θυγατέρες, αἳ οὐκ ἔγνωσαν ἄνδρα· ἐξάξω αὐτὰς πρὸς ὑμᾶς, καὶ χρᾶσθε αὐταῖς καθὰ ἂν ἀρέσκοι ὑμῖν· μόνον εἰς τοὺς ἄνδρας τούτους μὴ ποιήσητε ἄδικον, οὗ εἵνεκεν εἰσῆλθον ὑπὸ τὴν σκέπην τῶν δοκῶν μου.
\vs{9}Εἶπαν δὲ αὐτῷ, ἀπόστα ἐκεῖ· εἰσῆλθες παροικεῖν, μὴ καὶ κρίσιν κρίνειν; νῦν οὖν σε κακώσωμεν μᾶλλον ἢ ἐκείνους. Καὶ παρεβιάζοντο τὸν ἄνδρα τὸν Λὼτ σφόδρα, καὶ ἤγγισαν συντρίψαι τὴν θύραν.
\vs{10}Ἐκτείναντες δὲ οἱ ἄνδρες τὰς χεῖρας εἰσεσπάσαντο τὸν Λὼτ πρὸς ἑαυτοὺς εἰς τὸν οἶκον, καὶ τὴν θύραν τοῦ οἴκου ἀπέκλεισαν.
\vs{11}Τοὺς δὲ ἄνδρας τοὺς ὄντας ἐπὶ τῆς θύρας τοῦ οἴκου ἐπάταξαν ἐν ἀορασίᾳ ἀπὸ μικροῦ ἕως μεγάλου· καὶ παρελύθησαν ζητοῦντες τὴν θύραν.
\vs{12}Εἶπαν δὲ οἱ ἄνδρες πρὸς τὸν Λὼτ, εἰσί σοι ὧδε γαμβροὶ, ἢ υἱοὶ, ἢ θυγατέρες; ἢ εἴτις σοι ἄλλος ἐστὶν ἐν τῇ πόλει, ἐξάγαγε ἐκ τοῦ τόπου τούτου,
\vs{13}Ὅτι ἡμεῖς ἀπόλλυμεν τὸν τόπον τοῦτον· ὅτι ὑψώθη ἡ κραυγὴ αὐτῶν ἔναντι Κυρίου, καὶ ἀπέστειλεν ἡμᾶς Κύριος ἐκτρίψαι αὐτήν.
\vs{14}Ἐξῆλθε δὲ Λῶτ, καὶ ἐλάλησε πρὸς τοὺς γαμβροὺς αὐτοῦ τοὺς εἰληφότας τὰς θυγατέρας αὐτοῦ, καὶ εἶπεν, ἀνάστητε, καὶ ἐξέλθετε ἐκ τοῦ τόπου τούτου, ὅτι ἐκτρίβει Κύριος τὴν πόλιν· ἔδοξε δὲ γελοιάζειν ἐναντίον τῶν γαμβρῶν αὐτοῦ.
\vs{15}Ἡνίκα δὲ ὄρθρος ἐγένετο, ἐσπούδαζον οἱ ἄγγελοι τὸν Λὼτ, λέγοντες, ἀναστὰς λάβε τὴν γυναῖκά σου, καὶ τὰς δύο θυγατέρας σου, ἃς ἔχεις, καὶ ἔξελθε, ἵνα μὴ καὶ σὺ συναπόλῃ ταῖς ἀνομίαις τῆς πόλεως.
\vs{16}Καὶ ἐταράχθησαν, καὶ ἐκράτησαν οἱ ἄγγελοι τῆς χειρὸς αὐτοῦ, καὶ τῆς χειρὸς τῆς γυναικὸς αὐτοῦ, καὶ τῶν χειρῶν τῶν δύο θυγατέρων αὐτοῦ, ἐν τῷ φείσασθαι Κύριον αὐτοῦ.

\vs{17}Καὶ ἐγένετο ἡνίκα ἐξήγαγον αὐτοὺς ἔξω, καὶ εἶπαν, σώζων σῶζε τὴν σεαυτοῦ ψυχήν· μὴ περιβλέψῃ εἰς τὰ ὀπίσω, μηδὲ στῇς ἐν πάσῃ τῇ περιχώρῳ· εἰς τὸ ὄρος σώζου, μή ποτε συμπαραληφθῇς.
\vs{18}Εἶπε δὲ Λὼτ πρὸς αὐτοὺς, δέομαι
\vs{19}Κύριε, ἐπειδὴ εὗρεν ὁ παῖς σου ἔλεος ἐναντίον σου, καὶ ἐμεγάλυνας τὴν δικαιοσύνην σου, ὃ ποιεῖς ἐπʼ ἐμὲ, τοῦ ζῆν τὴν ψυχήν μου· ἐγὼ δὲ οὐ δυνήσομαι διασωθῆναι εἰς τὸ ὄρος, μή ποτε καταλάβῃ με τὰ κακὰ, καὶ ἀποθάνω.
\vs{20}Ἰδοὺ πόλις αὕτη ἐγγὺς τοῦ καταφυγεῖν με ἐκεῖ, ἥ ἐστι μικρά· καὶ ἐκεῖ διασωθήσομαι· οὐ μικρά ἐστι; καὶ ζήσεται ἡ ψυχή μου ἕνεκέν σου.
\vs{21}Καὶ εἶπεν αὐτῷ, ἰδοὺ ἐθαύμασά σου τὸ πρόσωπον καὶ ἐπὶ τῷ ῥήματι τούτῳ, τοῦ μὴ καταστρέψαι τὴν πόλιν περὶ ἧς ἐλάλησας.
\vs{22}Σπεῦσον οὖν τοῦ σωθῆναι ἐκεῖ, οὐ γὰρ δυνήσομαι ποιῆσαι πρᾶγμα, ἕως τοῦ ἐλθεῖν σε ἐκεῖ. διὰ τοῦτο ἐκάλεσε τὸ ὄνομα τῆς πόλεως ἐκείνης, Σηγώρ.
\vs{23}Ὁ ἥλιος ἐξῆλθεν ἐπὶ τὴν γῆν, καὶ Λὼτ εἰσῆλθεν εἰς Σηγώρ.
\vs{24}Καὶ Κύριος ἔβρεξεν ἐπὶ Σόδομα καὶ Γόμοῤῥα θεῖον καὶ πῦρ παρὰ Κυρίου ἐξ οὐρανοῦ.
\vs{25}Καὶ κατέστρεψε τὰς πόλεις ταύτας, καὶ πᾶσαν τὴν περίχωρον, καὶ πάντας τοὺς κατοικοῦντας ἐν ταῖς πόλεσι, καὶ τὰ ἀνατέλλοντα ἐκ τῆς γῆς.
\vs{26}Καὶ ἐπέβλεψεν ἡ γυνὴ αὐτοῦ εἰς τὰ ὀπίσω, καὶ ἐγένετο στήλη ἁλός.
\vs{27}Ὤρθρισε δὲ Ἁβραὰμ τῷ πρωῒ εἰς τὸν τόπον, οὗ εἱστήκει ἐναντίον Κυρίου.
\vs{28}Καὶ ἐπέβλεψεν ἐπὶ πρόσωπον Σοδόμων καὶ Γομόῤῥας, καὶ ἐπὶ πρόσωπον τῆς περιχώρου, καὶ εἶδε, καὶ ἰδοὺ ἀνέβαινεν φλὸξ ἐκ τῆς γῆς, ὡσεὶ ἀτμὶς καμίνου.
\vs{29}Καὶ ἐγένετο ἐν τῷ ἐκτρίψαι τὸν Θεὸν πάσας τὰς πόλεις τῆς περιοίκου, ἐμνήσθη ὁ Θεὸς τοῦ Ἁβραάμ· καὶ ἐξαπέστειλε τὸν Λὼτ ἐκ μέσου τῆς καταστροφῆς, ἐν τῷ καταστρέψαι Κύριον τὰς πόλεις, ἐν αἷς κατῴκει ἐν αὐταῖς Λώτ.

\vs{30}Ἀνέβη δὲ Λὼτ ἐκ Σηγὼρ, καὶ ἐκάθητο ἐν τῷ ὄρει αὐτὸς, καὶ αἱ δύο θυγατέρες αὐτοῦ μετʼ αὐτοῦ· ἐφοβήθη γὰρ κατοικῆσαι ἐν Σηγώρ· καὶ κατῴκησεν ἐν τῷ σπηλαίῳ αὐτὸς, καὶ αἱ δύο θυγατέρες αὐτοῦ μετʼ αὐτοῦ.
\vs{31}Εἶπε δὲ ἡ πρεσβυτέρα πρὸς τὴν νεωτέραν, ὁ πατὴρ ἡμῶν πρεσβύτερος, καὶ οὐδείς ἐστιν ἐπὶ τῆς γῆς, ὃς εἰσελεύσεται πρὸς ἡμᾶς, ὡς καθήκει πάσῃ τῇ γῇ.
\vs{32}Δεῦρο καὶ ποτίσωμεν τὸν πατέρα ἡμῶν οἶνον, καὶ κοιμηθῶμεν μετʼ αὐτοῦ, καὶ ἐξαναστήσωμεν ἐκ τοῦ πατρὸς ἡμῶν σπέρμα.
\vs{33}Ἐπότισαν δὲ τὸν πατέρα αὐτῶν οἶνον ἐν τῇ νυκτὶ ἐκείνῃ, καὶ εἰσελθοῦσα ἡ πρεσβυτέρα ἐκοιμήθη μετὰ τοῦ πατρὸς αὐτῆς ἐν τῇ νυκτὶ ἐκείνῃ· καὶ οὐκ ᾔδει ἐν τῷ κοιμηθῆναι αὐτὸν, καὶ ἐν τῷ ἀναστῆναι.
\vs{34}Ἐγένετο δὲ ἐν τῇ ἐπαύριον, καὶ εἶπεν ἡ πρεσβυτέρα πρὸς τὴν νεωτέραν, ἰδοὺ ἐκοιμήθην χθὲς μετὰ τοῦ πατρὸς ἡμῶν· ποτίσωμεν αὐτὸν οἶνον καὶ ἐν τῇ νυκτὶ ταύτῃ, καὶ εἰσελθοῦσα κοιμήθητι μετʼ αὐτοῦ, καὶ ἐξαναστήσωμεν ἐκ τοῦ πατρὸς ἡμῶν σπέρμα.
\vs{35}Ἐπότισαν δὲ καὶ ἐν τῇ νυκτὶ ἐκείνῃ τὸν πατέρα αὐτῶν οἶνον, καὶ εἰσελθοῦσα ἡ νεωτέρα ἐκοιμήθη μετὰ τοῦ πατρὸς αὐτῆς· καὶ οὐκ ᾔδει ἐν τῷ κοιμηθῆναι αὐτὸν, καὶ ἀναστῆναι.
\vs{36}Καὶ συνέλαβον αἱ δύο θυγατέρες Λὼτ ἐκ τοῦ πατρὸς αὐτῶν.
\vs{37}Καὶ ἔτεκεν ἡ πρεσβυτέρα υἱὸν, καὶ ἐκάλεσε τὸ ὄνομα αὐτοῦ Μωὰβ, λέγουσα, ἐκ τοῦ πατρός μου· οὗτος πατὴρ Μωαβιτῶν ἕως τῆς σήμερον ἡμέρας.
\vs{38}Ἔτεκε δὲ καὶ ἡ νεωτέρα υἱὸν, καὶ ἐκάλεσε τὸ ὄνομα αὐτοῦ Ἀμμὰν, λέγουσα, υἱὸς γένους μου· οὗτος πατὴρ Ἀμμανιτῶν ἕως τῆς σήμερον ἡμέρας.

\ch{20}
Καὶ ἐκίνησεν ἐκεῖθεν Ἁβραὰμ εἰς γῆν πρὸς Λίβα· καὶ ᾤκησεν ἀνὰ μέσον Κάδης, καὶ ἀνὰ μέσον Σούρ· καὶ παρῴκησεν ἐν Γεράροις.
\vs{2}Εἶπε δὲ Ἁβραὰμ περὶ Σάῤῥας τῆς γυναικὸς αὐτοῦ, ὅτι ἀδελφή μου ἐστίν· ἐφοβήθη γὰρ εἰπεῖν ὅτι γυνή μου ἐστὶ, μή ποτε ἀποκτείνωσιν αὐτὸν οἱ ἄνδρες τῆς πόλεως διʼ αὐτήν· ἀπέστειλε δὲ Ἀβιμέλεχ βασιλεὺς Γεράρων, καὶ ἔλαβε τὴν Σάῤῥαν.
\vs{3}Καὶ εἰσῆλθεν ὁ Θεὸς πρὸς Ἀβιμέλεχ ἐν ὕπνῳ τὴν νύκτα, καὶ εἶπεν, ἰδοὺ σὺ ἀποθνήσκεις περὶ τῆς γυναικὸς, ἧς ἔλαβες· αὕτη δέ ἐστι συνῳκηκυῖα ἀνδρί.
\vs{4}Ἀβιμέλεχ δὲ οὐχ ἥψατο αὐτῆς· καὶ εἶπε, Κύριε, ἔθνος ἀγνοοῦν καὶ δίκαιον ἀπολεῖς;
\vs{5}Οὐκ αὐτός μοι εἶπεν, ἀδφή μου ἐστί; καὶ αὕτη μοι εἶπεν, ἀδελφός μου ἐστίν; ἐν καθαρᾷ καρδίᾳ καὶ ἐν δικαιοσύνῃ χειρῶν ἐποίησα τοῦτο.
\vs{6}Εἶπε δὲ αὐτῷ ὁ Θεὸς καθʼ ὕπνον, κᾀγὼ ἔγνων ὅτι ἐν καθαρᾷ καρδίᾳ ἐποίησας τοῦτο, καὶ ἐφεισάμην σου τοῦ μὴ ἁμαρτεῖν σε εἰς ἐμέ· ἕνεκα τούτου οὐκ ἀφῆκά σε ἅψασθαι αὐτῆς.
\vs{7}Νῦν δὲ ἀπόδος τὴν γυναῖκα τῷ ἀνθρώπῳ, ὅτι προφήτης ἐστὶ, καὶ προσεύξεται περὶ σοῦ, καὶ ζήσῃ· εἰ δὲ μὴ ἀποδίδως, γνώσῃ ὅτι ἀποθανῇ σὺ καὶ πάντα τὰ σὰ.
\vs{8}Καὶ ὤρθρισεν Ἀβιμέλεχ τῷ πρωῒ, καὶ ἐκάλεσε πάντας τοὺς παῖδας αὐτοῦ, καὶ ἐλάλησε πάντα τὰ ῥήματα ταῦτα εἰς τὰ ὦτα αὐτῶν· ἐφοβήθησαν δὲ πάντες οἱ ἄνθρωποι σφόδρα.
\vs{9}Καὶ ἐκάλεσεν Ἀβιμέλεχ τὸν Ἁβραὰμ καὶ εἶπεν αὐτῷ, τί τοῦτο ἐποίησας ἡμῖν; μήτι ἡμάρτομεν εἰς σὲ, ὅτι ἐπήγαγες ἐπʼ ἐμὲ καὶ ἐπὶ τὴν βασιλείαν μου ἁμαρτίαν μεγάλην; ἔργον ὃ οὐδεὶς ποιήσει, πεποίηκάς μοι.
\vs{10}Εἶπε δὲ Ἀβιμέλεχ τῷ Ἁβραὰμ, τί ἐνιδὼν ἐποίησας τοῦτο;
\vs{11}Εἶπε δὲ Ἁβραὰμ, εἶπα γὰρ, ἄρα οὐκ ἔστι θεοσέβεια ἐν τῷ τόπῳ τούτῳ, ἐμέ τε ἀποκτενοῦσιν ἕνεκεν τῆς γυναικός μου.
\vs{12}Καὶ γὰρ ἀληθῶς, ἀδελφή μου ἐστὶν ἐκ πατρὸς, ἀλλʼ οὐκ ἐκ μητρός· ἐγενήθη δέ μοι εἰς γυναῖκα.
\vs{13}Ἐγένετο δὲ ἡνίκα ἐξήγαγέ με ὁ Θεὸς ἐκ τοῦ οἴκου τοῦ πατρός μου, καὶ εἶπα αὐτῇ, ταύτην τὴν δικαιοσύνην ποιήσεις εἰς ἐμὲ, εἰς πάντα τόπον οὗ ἐὰν εἰσέλθωμεν ἐκεῖ, εἶπον ἐμὲ, ὅτι ἀδελφός μου ἐστίν.
\vs{14}Ἔλαβε δὲ Ἀβιμέλεχ χίλια δίδραγμα, καὶ πρόβατα, καὶ μόσχους, καὶ παῖδας, καὶ παιδίσκας, καὶ ἔδωκε τῷ Ἁβραάμ· καὶ ἀπέδωκεν αὐτῷ Σάῤῥαν τὴν γυναῖκα αὐτοῦ.
\vs{15}Καὶ εἶπεν Ἀβιμέλεχ τῷ Ἁβραὰμ, ἰδοὺ ἡ γῆ μου ἐναντίον σου· οὗ ἄν σοι ἀρέσκῃ, κατοίκει.
\vs{16}Τῇ δὲ Σάῤῥᾳ εἶπεν, ἰδοὺ δέδωκα χίλια δίδραγμα τῷ ἀδελφῷ σου· ταῦτα ἔσται σοι εἰς τιμὴν τοῦ προσώπου σου, καὶ πάσαις ταῖς μετὰ σοῦ· καὶ πάντα ἀλήθευσον.
\vs{17}Προσηύξατο δὲ Ἁβραὰμ πρὸς τὸν Θεὸν, καὶ ἰάσατο ὁ Θεὸς τὸν Ἀβιμέλεχ, καὶ τὴν γυναῖκα αὐτοῦ, καὶ τὰς παιδίσκας αὐτοῦ· καὶ ἔτεκον.
\vs{18}Ὅτι συγκλείων συνέκλεισε Κύριος ἔξωθεν πᾶσαν μήτραν ἐν τῷ οἴκῳ Ἀβιμέλεχ, ἕνεκεν Σάῤῥας τῆς γυναικὸς Ἁβραάμ.

\ch{21}
Καὶ Κύριος ἐπεσκέψατο τὴν Σάῤῥαν, καθὰ εἶπε· καὶ ἐποίησε Κύριος τῇ Σάῥῥᾳ, καθὰ ἐλάλησε.
\vs{2}Καὶ συλλαβοῦσα ἔτεκε τῷ Ἁβραὰμ υἱὸν εἰς τὸ γῆρας, εἰς τὸν καιρὸν καθὰ ἐλάλησεν αὐτῷ Κύριος.
\vs{3}Καὶ ἐκάλεσεν Ἁβραὰμ τὸ ὄνομα τοῦ υἱοῦ αὐτοῦ τοῦ γενομένου αὐτῷ, ὃν ἔτεκεν αὐτῷ Σάῤῥα, Ἰσαάκ·
\vs{4}Περιέτεμε δὲ Ἁβραὰμ τὸν Ἰσαὰκ τῇ ἡμέρᾳ τῇ ὀγδόῃ, καθὰ ἐνετείλατο αὐτῷ ὁ Θεός.
\vs{5}Καὶ Ἁβραὰμ ἦν ἑκατὸν ἐτῶν, ηνίκα ἐγένετο αὐτῷ Ἰσαὰκ ὁ υἱὸς αὐτοῦ.
\vs{6}Εἶπε δὲ Σάῤῥα, γέλωτά μοι ἐποίησε Κύριος· ὃς γὰρ ἂν ἀκούσῃ συγχαρεῖταί μοι.
\vs{7}Καὶ εἶπε τίς ἀναγγελεῖ τῷ Ἁβραὰμ ὅτι θηλάζει παιδίον Σάῤῥα; ὅτι ἔτεκον υἱὸν ἐν τῷ γήρᾳ μου.
\vs{8}Καὶ ηὐξήθη τὸ παιδίον, καὶ ἀπεγαλακτίσθη· καὶ ἐποίησεν Ἁβραὰμ δοχὴν μεγάλην, ᾗ ἡμέρᾳ ἀπεκγαλακτίσθη Ἰσαὰκ ὁ υἱὸς αὐτοῦ.
\vs{9}Ἰδοῦσα δὲ Σάῥῥα τὸν υἱὸν Ἄγαρ τῆς Αἰγυπτίας, ὃς ἐγένετο τῷ Ἁβραὰμ, παίζοντα μετὰ Ἰσαὰκ τοῦ υἱοῦ αὐτῆς,
\vs{10}καὶ εἶπε τῷ Ἁβραὰμ, ἔκβαλε τὴν παιδίσκην ταύτην, καὶ τὸν υἱὸν αὐτῆς· οὐ γὰρ μὴ κληρονομήσει ὁ υἱὸς τῆς παιδίσκης ταύτης μετὰ τοῦ υἱοῦ μου Ἰσαάκ.
\vs{11}Σκληρὸν δὲ ἐφάνη τὸ ῥῆμα σφόδρα ἐνατίον Ἁβραὰμ περὶ τοῦ υἱοῦ αὐτοῦ.
\vs{12}Εἶπε δὲ ὁ Θεὸς τῷ Ἁβραὰμ, μὴ σκληρὸν ἔστω ἐναντίον σου περὶ τοῦ παιδίου, καὶ περὶ τῆς παιδίσκης· πάντα ὅσα ἂν εἴπῃ σοι Σάῤῥα, ἄκουε τῆς φωνῆς αὐτῆς· ὅτι ἐν Ἰσαὰκ κληθήσεταί σοι σπέρμα.
\vs{13}Καὶ τὸν υἱὸν δὲ τῆς παιδίσκης ταύτης εἰς ἔθνος μέγα ποιήσω αὐτὸν, ὅτι σπέρμα σόν ἐστιν.
\vs{14}Ἀνέστη δὲ Ἁβραὰμ τὸ πρωῒ, καὶ ἔλαβεν ἄρτους καὶ ἀσκὸν ὕδατος, καὶ ἔδωκεν τῇ Ἄγαρ· καὶ ἐπέθηκεν ἐπὶ τὸν ὦμον αὐτῆς τὸ παιδίον, καὶ ἀπέστειλεν αὐτήν· Ἀπελθοῦσα δὲ ἐπλανᾶτο κατὰ τὴν ἔρημον, κατὰ τὸ φρέαρ τοῦ ὅρκου.
\vs{15}Ἐξέλιπε δὲ τὸ ὕδωρ ἐκ τοῦ ἀσκου· καὶ ἔῤῥιψε τὸ παιδίον ὑποκάτω μιᾶς ἐλάτης·
\vs{16}Ἀπελθοῦσα δὲ ἐκάθητο ἀπέναντι αὐτοῦ μακρόθεν, ὡσεὶ τόξου βολήν· εἶπε γὰρ, οὐ μὴ ἴδω τὸν θάνατον τοῦ παιδίου μου. καὶ ἐκάθισεν ἀπέναντι αὐτοῦ· ἀναβοῆσαν δὲ τὸ παιδίον ἔκλαυσεν.
\vs{17}Εἰσήκουσε δὲ ὁ Θεὸς τῆς φωνῆς τοῦ παιδίου ἐκ τοῦ τόπου οὗ ἦν· καὶ ἐκάλεσεν ἄγγελος Θεοῦ τὴν Ἄγαρ ἐκ τοῦ οὐρανοῦ, καὶ εἶπεν αὐτῇ, τί ἐστιν Ἄγαρ; μὴ φοβοῦ· ἐπακήκοε γὰρ ὁ Θεὸς τῆς φωνῆς τοῦ παιδίου ἐκ τοῦ τόπου οὗ ἐστιν.
\vs{18}Ἀνάστηθι καὶ λάβε τὸ παιδίον, καὶ κράτησον τῇ χειρί σου αὐτό· εἰς γὰρ ἔθνος μέγα ποιήσω αὐτό.
\vs{19}Καὶ ἀνέῳξεν ὁ Θεὸς τοὺς ὀφθαλμοὺς αὐτῆς· καὶ εἶδε φρέαρ ὕδατος ζῶντος, καὶ ἐπορεύθη, καὶ ἔπλησε τὸν ἀσκὸν ὕδατος, καὶ ἐπότισε τὸ παιδίον.
\vs{20}Καὶ ἦν ὁ Θεὸς μετὰ τοῦ παιδίου· καὶ ηὐξήθη, καὶ κατῴκησεν ἐν τῇ ἐρήμῳ· ἐγένετο δὲ τοξότης.
\vs{21}Καὶ κατῴκησεν ἐν τῇ ἐρήμῳ· καὶ ἔλαβεν αὐτῷ ἡ μήτηρ γυναῖκα ἐκ Φαρὰν Αἰγύπτου.

\vs{22}Ἐγένετο δὲ ἐν τῷ καιρῷ ἐκείνῳ, καὶ εἶπεν Ἀβιμέλεχ, καὶ Ὁχοζὰθ ὁ νυμφαγωγὸς αὐτοῦ, καὶ Φιχὸλ ὁ ἀρχιστράτηγος τῆς δυνάμεως αὐτοῦ, πρὸς Ἁβραὰμ, λέγων, ὁ Θεὸς μετὰ σοῦ ἐν πᾶσιν, οἷς ἐὰν ποιῇς.
\vs{23}Νῦν οὖν ὄμοσόν μοι τὸν Θεὸν μὴ ἀδικήσειν με, μηδὲ τὸ σπέρμα μου, μηδὲ τὸ ὄνομά μου· ἀλλὰ κατὰ τὴν δικαιοσύνην ἣν ἐποίησα μετὰ σοῦ, ποιήσεις μετʼ ἐμοῦ, καὶ τῇ γᾗ, ᾗ σὺ παρῴκησας ἐν αὐτῇ.
\vs{24}Καὶ εἶπεν Ἁβραὰμ, ἐγὼ ὀμοῦμαι.
\vs{25}Καὶ ἤλεγξεν Ἁβραὰμ τὸν Ἀβιμέλεχ περὶ τῶν φρεάτων τοῦ ὕδατος, ὧν ἀφείλοντο οἱ παῖδες τοῦ Ἀβιμέλεχ.
\vs{26}Καὶ εἶπεν αὐτῷ Ἀβιμέλεχ, οὐκ ἔγνων τίς ἐποίησέ σοι τὸ ῥῆμα τοῦτο· οὐδὲ σύ μοι ἀπήγγειλας, οὐδὲ ἐγὼ ἤκουσα, ἀλλʼ ἢ σήμερον.
\vs{27}Καὶ ἔλαβεν Ἁβραὰμ πρόβατα καὶ μόσχους, καὶ ἔδωκε τῷ Ἀβιμέλεχ· καὶ διέθεντο ἀμφότεροι διαθήκην.
\vs{28}Καὶ ἔστησεν Ἁβραὰμ, ἑπτὰ ἀμνάδας προβάτων μόνας.
\vs{29}Καὶ εἶπεν Ἀβιμέλεχ τῷ Ἁβραὰμ, τί εἰσιν αἱ ἑπτὰ ἀμνάδες τῶν προβάτων τούτων, ἃς ἔστησας μόνας;
\vs{30}Καὶ εἶπεν Ἁβραὰμ, ὅτι τὰς ἑπτὰ ἀμνάδας λήψῃ παρʼ ἐμοῦ, ἵνα ὦσι μοι εἰς μαρτύριον, ὅτι ἐγὼ ὤρυξα τό φρέαρ τοῦτο.
\vs{31}Διὰ τοῦτο ἐπωνόμασε τὸ ὄνομα τοῦ τόπου ἐκείνου, Φρέαρ ὁρκισμοῦ· ὅτι ἐκεῖ ὤμοσαν ἀμφότεροι.
\vs{32}Καὶ διέθεντο διαθήκην ἐν τῷ φρέατι τοῦ ὁρκισμου· ἀνέστη δὲ Ἀβιμέλεχ, Ὁχοζὰθ ὁ νυμφαγωγὸς αὐτοῦ, καὶ Φίχολ ὁ ἀρχιστράτηγος τῆς δυνάμεως αὐτοῦ, καὶ ἐπέστρεψαν εἰς τὴν γῆν τῶν Φυλιστιείμ.
\vs{33}Καὶ ἐφύτευσεν Ἁβραὰμ ἄρουραν ἐπὶ τῷ φρέατι τοῦ ὅρκου· καὶ ἐπεκαλέσατο ἐκεῖ τὸ ὄνομα Κυρίου, Θεὸς αἰώνιος.
\vs{34}Παρῴκησε δὲ Ἁβραὰμ ἐν τῇ γῇ τῶν Φυλιστιεὶμ ἡμέρας πολλάς.

\ch{22}
Καὶ ἐγένετο μετὰ τὰ ῥήματα ταῦτα ὁ Θεὸς ἐπείρασε τὸν Ἁβραὰμ, καὶ εἶπεν αὐτῷ, Ἁβραὰμ, Ἁβραάμ· καὶ εἶπεν, ἰδοὺ ἐγώ.
\vs{2}Καὶ εἶπε, λάβε τὸν υἱόν σου τὸν ἀγαπητὸν, ὃν ἠγάπησας, τὸν Ἰσαὰκ, καὶ πορεύθητι εἰς τὴν γῆν τὴν ὑψηλὴν, καὶ ἀνένεγκε αὐτὸν ἐκεῖ εἰς ὁλοκάρπωσιν ἐφʼ ἓν τῶν ὀρέων ὧν ἄν σοι εἴπω.
\vs{3}Ἀναστὰς δὲ Ἁβραὰμ τὸ πρωῒ, ἐπέσαξε τὴν ὄνον αὐτοῦ· παρέλαβε δὲ μεθʼ ἑαυτοῦ δύο παῖδας, καὶ Ἰσαὰκ τὸν υἱὸν αὐτοῦ· καὶ σχίσας ξύλα εἰς ὁλοκάρπωσιν, ἀναστὰς ἐπορεύθη, καὶ ἦλθεν ἐπὶ τὸν τόπον, ὃν εἶπεν αὐτῷ ὁ Θεὸς, τῇ ἡμέρᾳ τῇ τρίτῃ.
\vs{4}Καὶ ἀναβλέψας Ἁβραὰμ τοῖς ὀφθαλμοῖς αὐτοῦ, εἶδε τὸν τόπον μακρόθεν.
\vs{5}Καὶ εἶπεν Ἁβραὰμ τοῖς παισὶν αὐτοῦ, καθίσατε αὐτοῦ μετὰ τῆς ὄνου· ἐγὼ δὲ καὶ τὸ παιδάριον διελευσόμεθα ἕως ὧδε· καὶ προσκυνήσαντες ἀναστρέψομεν πρὸς ὑμᾶς.
\vs{6}Ἔλαβε δὲ Ἁβραὰμ τὰ ξύλα τῆς ὁλοκαρπώσεως, καὶ ἐπέθηκεν Ἰσαὰκ τῷ υἱῷ αὐτοῦ· ἔλαβε δὲ μετὰ χεῖρας καὶ τὸ πῦρ καὶ τὴν μάχαιραν, καὶ ἐπορεύθησαν οἱ δύο ἅμα.
\vs{7}Εἶπε δὲ Ἰσαὰκ πρὸς Ἁβραὰμ τὸν πατέρα αὐτοῦ, πάτερ· ὁ δὲ εἶπε, τί ἐστι, τέκνον; εἶπε, δὲ, ἰδοὺ τὸ πῦρ καὶ τὰ ξύλα, ποῦ ἐστὶ τὸ πρόβατον τὸ εἰς ὁλοκάρπωσιν;
\vs{8}Εἶπε δὲ Ἁβραὰμ, ὁ Θεὸς ὄψεται ἑαυτῷ πρόβατον εἰς ὁλοκάρπωσιν, τέκνον. πορευθέντες δὲ ἀμφότεροι ἅμα,
\vs{9}ἦλθον ἐπὶ τὸν τόπον, ὃν εἶπεν αὐτῷ ὁ Θεός· καὶ ᾠκοδόμησεν ἐκεῖ Ἁβραὰμ τὸ θυσιαστήριον, καὶ ἐπέθηκε τὰ ξύλα· καὶ συμποδίσας Ἰσαὰκ τὸν υἱὸν αὐτοῦ, ἐπέθηκεν αὐτὸν ἐπὶ τὸ θυσιαστήριον ἐπάνω τῶν ξύλων.
\vs{10}Καὶ ἐξέτεινεν Ἁβραὰμ τὴν χεῖρα αὐτοῦ λαβεῖν τὴν μάχαιραν, σφάξαι τὸν υἱὸν αὐτοῦ.
\vs{11}Καὶ ἐκάλεσεν αὐτὸν Ἄγγελος Κυρίου ἐκ τοῦ οὐρανοῦ, καὶ εἶπεν, Ἁβραὰμ, Ἁβραάμ· ὁ δὲ εἶπεν, ἰδοὺ ἐγώ.
\vs{12}Καὶ εἶπε, μὴ ἐπιβάλῃς τὴν χεῖρά σου ἐπὶ τὸ παιδάριον, μηδὲ ποιήσῃς αὐτῷ μηδέν· νῦν γὰρ ἔγνων, ὅτι φοβῇ σὺ τὸν Θεόν· καὶ οὐκ ἐφείσω τοῦ υἱοῦ σου τοῦ ἀγαπητοῦ διʼ ἐμέ.
\vs{13}Καὶ ἀναβλέψας Ἁβραὰμ τοῖς ὀφθαλμοῖς αὐτοῦ εἶδε, καὶ ἰδοὺ κριὸς εἷς κατεχόμενος ἐν φυτῷ Σαβὲκ τῶν κεράτων. Καὶ ἐπορεύθη Ἁβραὰμ, καὶ ἔλαβε τὸν κριὸν, καὶ ἁνήνεγκεν αὐτὸν εἰς ὁλοκάρπωσιν ἀντὶ Ἰσαὰκ τοῦ υἱοῦ αὐτοῦ.

\vs{14}Καὶ ἐκάλεσεν Ἁβραὰμ τὸ ὄνομα τοῦ τόπου ἐκείνου, Κύριος εἶδεν· ἵνα εἴπωσιν σήμερον, ἐν τῷ ὄρει Κύριος ὤφθη.
\vs{15}Καὶ ἐκάλεσεν Ἄγγελος Κυρίου τὸν Ἁβραὰμ δεύτερον ἐκ τοῦ οὐρανοῦ,
\vs{16}λέγων, Κατʼ ἐμαυτοῦ ὤμοσα, λέγει Κύριος, οὗ εἵνεκεν ἐποίησας τὸ ῥῆμα τοῦτο, καὶ οὐκ ἐφείσω τοῦ υἱοῦ σου τοῦ ἀγαπτοῦ διʼ ἐμὲ,
\vs{17}Ἦ μὴν εὐλογῶν εὐλογήσω σε, καὶ πληθύνων πληθυνῶ τὸ σπέρμα σου, ὡς τοὺς ἀστέρας τοῦ οὐρανοῦ, καὶ ὡς τὴν ἄμμον τὴν παρὰ τὸ χεῖλος τῆς θαλάσσης· καὶ κληρονομήσει τὸ σπέρμα σου τὰς πόλεις τῶν ὑπεναντίων.
\vs{18}Καὶ ἐνευλογηθήσονται ἐν τῷ σπέρματί σου πάντα τὰ ἔθνη τῆς γῆς, ἀνθʼ ὧν ὑπήκουσας τῆς ἐμῆς φωνῆς.
\vs{19}Ἀπεστράφη δὲ Ἁβραὰμ πρὸς τοὺς παῖδας αὐτοῦ· καὶ ἀναστάντες ἐπορεύθησαν ἅμα ἐπὶ τὸ φρέαρ τοῦ ὅρκου. Καὶ κατῴκησεν Ἁβραὰμ ἐπὶ τὸ φρέαρ τοῦ ὅρκου.

\vs{20}Ἐγένετο δὲ μετὰ τὰ ῥήματα ταῦτα, καὶ ἀνηγγέλη τῷ Ἁβραὰμ, λέγοντες, ἰδοὺ τέτοκε Μελχὰ καὶ αὐτὴ υἱοὺς τῷ Ναχὼρ τῷ ἀδελφῷ σου,
\vs{21}τὸν Οὒζ πρωτότοκον, καὶ τὸν Βαὺξ ἀδελφὸν αὐτοῦ, καὶ τὸν Καμουὴλ πατέρα Σύρων,
\vs{22}καὶ τὸν Χαζὰδ, καὶ Ἀζαῦ, καὶ τὸν Φαλδὲς, καὶ τὸν Ἰελδὰφ, καὶ τὸν Βαθουήλ.
\vs{23}Βαθουὴλ δὲ ἐγέννησε τὴν Ῥεβέκκαν. ὀκτὼ οὗτοι υἱοὶ, οὓς ἔτεκε Μελχὰ τῷ Ναχὼρ τῷ ἀδελφῷ Ἁβραάμ.
\vs{24}Καὶ ἡ παλλακὴ αὐτοῦ, ᾗ ὄνομα Ῥεύμα, ἔτεκε καὶ αὐτὴ τὸν Ταβὲκ, καὶ τὸν Ταὰμ, καὶ τὸν Τοχός, καὶ τὸν Μοχά.

\ch{23}
Ἐγένετο δὲ ἡ ζωὴ Σάῤῥας, ἔτη ἑκατὸν εἰκοσιεπτά.
\vs{2}Καὶ ἀπέθανε Σάῤῥα ἐν πόλει Ἀρβὸκ, ἥ ἐστιν ἐν τῷ κοιλώματι· αὕτη ἔστι Χεβρὼν ἐν τῇ γῇ Χαναάν ἦλθε δὲ Ἁβραὰμ κόψασθαι Σάῤῥαν, καὶ πενθῆσαι.
\vs{3}Καὶ ἀνέστη Ἁβραὰμ ἀπὸ τοῦ νεκροῦ αὐτοῦ· καὶ εἶπεν Ἁβραὰμ τοῖς υἱοῖς τοῦ Χὲτ, λέγων,
\vs{4}Πάροικος καὶ παρεπίδημος ἐγώ εἰμι μεθʼ ὑμῶν· δότε μοι οὖν κτῆσιν τάφου μεθʼ ὑμῶν, καὶ θάψω τὸν νεκρόν μου ἀπʼ ἐμοῦ.
\vs{5}Ἀπεκρίθησαν δὲ οἱ υἱοὶ Χὲτ πρὸς Ἁβραὰμ, λέγοντες, μὴ, κύριε.
\vs{6}Ἄκουσον δὲ ἡμῶν· βασιλεὺς παρὰ Θεοῦ σὺ εἶ ἐν ἡμῖν· ἐν τοῖς ἐκλεκτοῖς μνημείοις ἡμῶν θάψον τὸν νεκρόν σου· οὐδεὶς γὰρ ἡμῶν οὐ μὴ κωλύσει τὸ μνημεῖον αὐτοῦ ἀπὸ σοῦ, τοῦ θάψαι τὸν νεκρόν σου ἐκεῖ.
\vs{7}Ἀναστὰς δὲ Ἁβραὰμ προσεκύνησε τῷ λαῷ τῆς γῆς, τοῖς υἱοῖς τοῦ Χέτ.
\vs{8}Καὶ ἐλάλησε πρὸς αὐτοὺς Ἁβραὰμ, λέγων, εἰ ἔχετε τῇ ψυχῇ ὑμῶν, ὥστε θάψαι τὸν νεκρόν μου ἀπὸ προσώπου μου, ἀκούσατέ μου, καὶ λαλήσατε περὶ ἐμοῦ Ἐφρὼν τῷ τοῦ Σαάρ.
\vs{9}Καὶ δότω μοι τὸ σπήλαιον τὸ διπλοῦν, ὅ ἐστιν αὐτῷ, τὸ ὂν ἐν μέρει τοῦ ἀγροῦ αὐτοῦ· ἀργυρίου τοῦ ἀξίου δότε μοι αὐτὸ ἐν ὑμῖν εἰς κτῆσιν μνημείου.
\vs{10}Ἐφρὼν δὲ ἐκάθητο ἐν μέσῳ τῶν υἱῶν Χέτ· ἀποκριθεὶς δὲ Ἐφρὼν ὁ Χετταῖος πρὸς Ἁβραὰμ εἶπεν, ἀκουόντων τῶν υἱῶν Χὲτ, καὶ τῶν εἰσπορευομένων εἰς τὴν πόλιν πάντων, λέγων,
\vs{11}Παρʼ ἐμοὶ γενοῦ, κύριε, καὶ ἄκουσόν μου· τὸν ἀγρὸν, καὶ τὸ σπήλαιον τὸ ἐν αὐτῷ, σοὶ δίδωμι· ἐναντίον πάντων τῶν πολιτῶν μου δέδωκά σοι· θάψον τὸν νεκρόν σου.
\vs{12}Καὶ προσεκύνησεν Ἁβραὰμ ἐναντίον τοῦ λαοῦ τῆς γῆς.
\vs{13}Καὶ εἶπε τῷ Ἐφρὼν εἰς τὰ ὦτα ἐναντίον τοῦ λαοῦ τῆς γῆς, ἐπειδὴ πρὸς ἐμοῦ εἶ, ἄκουσόν μου· τὸ ἀργύριον τοῦ ἀγροῦ λάβε παρʼ ἐμοῦ, καὶ θάψω τὸν νεκρόν μου ἐκεῖ.
\vs{14}Ἀπεκρίθη δὲ Ἐφρὼν τῷ Ἁβραὰμ, λέγων,
\vs{15}Οὐχὶ, κύριε· ἀκήκοα γὰρ, γῆ τετρακοσίων διδράχμων ἀργύριου· ἀλλὰ τί ἂν εἴη τοῦτο ἀνὰ μέσον ἐμοῦ καὶ σοῦ; σὺ δὲ τὸν νεκρόν σου θάψον.
\vs{16}καὶ ἤκουσεν Ἁβραὰμ τοῦ Ἐφρών· καὶ ἀπεκατέστησεν Ἁβραὰμ τῷ Ἐφρὼν τὸ ἀργύριον, ὃ ἐλάλησεν εἰς τὰ ὦτα τῶν υἱῶν Χὲτ, τετρακόσια δίδραχμα ἀργυρίου δοκίμου ἐμπόροις.
\vs{17}Καὶ ἔστη ὁ ἀγρὸς Ἐφρών, ὃς ἦν ἐν τῷ διπλῷ σπηλαίῳ, ὅς ἐστι κατὰ πρόσωπον Μαμβρῆ, ὁ ἀγρὸς καὶ τὸ σπήλαιον, ὃ ἦν ἐν αὐτῷ, καὶ πᾶν δένδρον, ὃ ἦν ἐν τῷ ἀγρῷ, καὶ πᾶν ὅ ἐστιν ἐν τοῖς ὁρίοις αὐτοῦ κύκλῳ,
\vs{18}τῷ Ἁβραὰμ, εἰς κτῆσιν ἐναντίον τῶν υἱῶν Χὲτ, καὶ πάντων τῶν εἰσπορευομένων εἰς τὴν πόλιν.
\vs{19}Μετὰ ταῦτα ἔθαψεν Ἁβραὰμ Σάῤῥαν τὴν γυναῖκα αὐτοῦ ἐν τῷ σπηλαίῳ τοῦ ἀγροῦ τῷ διπλῷ, ὅ ἐστιν ἀπέναντι Μαμβρῆ· αὕτη ἐστὶ Χεβρὼν ἐν τῇ γῇ Χαναάν.
\vs{20}Καὶ ἐκυρώθη ὁ ἀγρὸς καὶ τὸ σπήλαιον ὃ ἦν ἐν αὐτῷ τῷ Ἁβραὰμ εἰς κτῆσιν τάφου, παρὰ τῶν υἱῶν Χέτ.

\ch{24}Καὶ Ἁβραὰμ ἦν πρεσβύτερος προβεβηκὼς ἡμερῶν· καὶ Κύριος ηὐλόγησε τὸν Ἁβραὰμ κατὰ πάντα.

\vs{2}Καὶ εἶπεν Ἁβραὰμ τῷ παιδὶ αὐτοῦ τῷ πρεσβυτέρῳ τῆς οἰκίας αὐτοῦ, τῷ ἄρχοντι πάντων τῶν αὐτοῦ, θὲς τὴν χεῖρά σου ὑπὸ τὸν μηρόν μου.
\vs{3}Καὶ ἐξορκιῶ σε Κύριον τὸν Θεὸν τοῦ οὐρανοῦ καὶ τὸν Θεὸν τῆς γῆς, ἵνα μὴ λάβῃς γυναῖκα τῷ υἱῷ μου Ἰσαὰκ ἀπὸ τῶν θυγατέρων τῶν Χαναναίων, μεθʼ ὧν ἐγὼ οἰκῶ ἐν αὐτοις.
\vs{4}Ἀλλʼ ἢ εἰς τὴν γῆν μου, οὗ ἐγεννήθην, πορεύσῃ, καὶ εἰς τὴν φυλήν μου, καὶ λήψῃ γυναῖκα τῷ υἱῷ μου Ἰσαὰκ ἐκεῖθεν.
\vs{5}Εἶπε δὲ πρὸς αὐτὸν ὁ παῖς, μή ποτε οὐ βούληται ἡ γυνὴ πορευθῆναι μετʼ ἐμοῦ ὀπίσω εἰς τὴν γῆν ταύτην, ἀποστρέψω τὸν υἱόν σου εἰς τὴν γῆν, ὅθεν ἐξῆλθες ἐκεῖθεν;
\vs{6}Εἶπε δὲ πρὸς αὐτὸν Ἁβραάμ, πρόσεχε σεαυτῷ μὴ ἀποστρέψῃς τὸν υἱόν μου ἐκεῖ.
\vs{7}Κύριος ὁ Θεὸς τοῦ οὐρανοῦ καὶ ὁ Θεὸς τῆς γῆς, ὃς ἔλαβέ με ἐκ τοῦ οἴκου τοῦ πατρός μου, καἰ ἐκ τῆς γῆς ἧς ἐγεννήθην, ὃς ἐλάλησέ μοι, καὶ ὃς ὤμοσέ μοι, λέγων, σοὶ δώσω τὴν γῆν ταύτην καὶ τῷ σπέρματί σου, αὐτὸς ἀποστελεῖ τὸν Ἄγγελον αὐτοῦ ἔμπροσθέν σου, καὶ λήψῃ γυναῖκα τῷ υἱῷ μου ἐκεῖθεν.
\vs{8}Ἐὰν δὲ μὴ θέλῃ ἡ γυνὴ πορευθῆναι μετὰ σοῦ εἰς τὴν γῆν ταύτην, καθαρὸς ἔσῃ ἀπὸ τοῦ ὅρκου μου· μόνον τὸν υἱόν μου μὴ ἀποστρέψῃς ἐκεῖ.
\vs{9}Καὶ ἔθηκεν ὁ παῖς τὴν χεῖρα αὐτοῦ ὑπὸ τὸν μηρὸν Ἁβραὰμ τοῦ κυρίου αὐτοῦ, καὶ ὤμοσεν αὐτῷ περὶ τοῦ ῥήματος τούτου.
\vs{10}Καὶ ἔλαβεν ὁ παῖς δέκα καμήλους ἀπὸ τῶν καμήλων τοῦ κυρίου αὐτοῦ, καὶ ἀπὸ πάντων τῶν ἀγαθῶν τοῦ κυρίου αὐτοῦ μεθʼ ἑαυτοῦ· καὶ ἀναστὰς ἐπορεύθη εἰς τὴν Μεσοποταμίαν εἰς τὴν πόλιν Ναχώρ.
\vs{11}Καὶ ἐκοίμησε τὰς καμήλους ἔξω τῆς πόλεως παρὰ τὸ φρέαρ τοῦ ὕδατος τὸ πρὸς ὀψέ, ἡνίκα ἐκπορεύονται αἱ ὑδρευόμεναι.

\vs{12}Καὶ εἶπε, Κύριε ὁ Θεὸς τοῦ κυρίου μου Ἁβραάμ, εὐόδωσον ἐναντίον ἐμοῦ σήμερον, καὶ ποίησον ἔλεος μετὰ τοῦ κυρίου μου Ἁβραάμ.
\vs{13}Ἰδοὺ ἐγὼ ἕστηκα ἐπὶ τῆς πηγῆς τοῦ ὕδατος· αἱ δὲ θυγατέρες τῶν οἰκούντων τὴν πόλιν ἐκπορεύονται ἀντλῆσαι ὕδωρ.
\vs{14}Καὶ ἔσται ἡ παρθένος ᾗ ἂν ἐγὼ εἴπω, ἐπίκλινον τὴν ὑδρίαν σου, ἵνα πίω, καὶ εἴπῃ μοι, πίε σύ, καὶ τὰς καμήλους σου ποτιῶ, ἕως ἂν παύσωνται πίνουσαι, ταύτην ἡτοίμασας τῷ παιδί σου τῷ Ἰσαάκ· καὶ ἐν τούτῳ γνώσομαι, ὅτι ἐποίησας ἔλεος μετὰ τοῦ κυρίου μου Ἁβραάμ.

\vs{15}Καὶ ἐγένετο πρὸ τοῦ συντελέσαι αὐτὸν λαλοῦντα ἐν τῇ διανοίᾳ αὐτοῦ, καὶ ἰδοὺ Ῥεβέκκα ἐξεπορεύετο ἡ τεχθεῖσα Βαθουήλ, υἱῷ Μελχὰς τῆς γυναικὸς Ναχώρ, ἀδελφοῦ δὲ Ἁβραάμ, ἔχουσα τὴν ὑδρίαν ἐπὶ τῶν ὤμων αὐτῆς.
\vs{16}Ἡ δὲ παρθένος ἦν καλὴ τῇ ὄψει σφόδρα· παρθένος ἦν, ἀνὴρ οὐκ ἔγνω αὐτήν· καταβᾶσα δὲ ἐπὶ τὴν πηγὴν, ἔπλησε τὴν ὑδρίαν αὐτῆς, καὶ ἀνέβη.
\vs{17}Ἐπέδραμε δὲ ὁ παῖς εἰς συνάντησιν αὐτῆς, καὶ εἶπε, Πότισόν με μικρὸν ὕδωρ ἐκ τῆς ὑδρίας σου.
\vs{18}Ἡ δὲ εἶπε, πίε, κύριε· καὶ ἔσπευσε καὶ καθεῖλε τὴν ὑδρίαν ἐπὶ τὸν βραχίονα αὐτῆς, καὶ ἐπότισεν αὐτὸν, ἕων ἐπαύσατο πίνων.
\vs{19}Καὶ εἶπε, καὶ ταῖς καμήλοις σου ὑδρεύσομαι, ἕως ἂν πᾶσαι πίωσι.
\vs{20}Καὶ ἔσπευσε καὶ ἐξεκένωσε τὴν ὑδρίαν εἰς τὸ ποτιστήριον· καὶ ἔδραμεν ἐπὶ τὸ φρέαρ ἀντλῆσαι πάλιν· καὶ ὑδρεύσατο πάσαις ταῖς καμήλοις.
\vs{21}Ὁ δὲ ἄνθρωπος κατεμάνθανεν αὐτήν· καὶ παρεσιώπα τοῦ γνῶναι εἰ εὐώδωκε Κύριος τὴν ὁδὸν αὐτοῦ, ἢ οὔ.
\vs{22}Ἐγένετο δὲ ἡνίκα ἐπαύσαντο πᾶσαι αἱ κάμηλοι πίνουσαι, ἔλαβεν ὁ ἄνθρωπος ἐνώτια χρυσᾶ ἀνὰ δραχμὴν ὁλκῆς, καὶ δύο ψέλλια ἐπὶ τὰς χεῖρας αὐτῆς, δέκα χρυσῶν ὁλκὴ αὐτῶν.
\vs{23}Καὶ ἐπηρώτησεν αὐτὴν, καὶ εἶπε, θυγάτηρ τίνος εἶ; ἀνάγγειλόν μοι, εἰ ἔστι παρὰ τῷ πατρί σου τόπος ἡμῖν του καταλῦσαι.
\vs{24}Ἡ δὲ εἶπεν αὐτῷ, θυγάτηρ Βαθουήλ εἰμι τοῦ Μελχάς, ὃν ἔτεκε τῷ Ναχώρ.
\vs{25}Καὶ εἶπεν αὐτῷ, Καὶ ἄχυρα καὶ χορτάσματα πολλὰ παρʼ ἡμῖν, καὶ τόπος τοῦ καταλῦσαι.
\vs{26}Καὶ εὐδοκήσας ὁ ἄνθρωπος προσεκύνησε τῷ Κυρίῳ
\vs{27}Καὶ εἶπεν, εὐλογητὸς Κύριος ὁ Θεὸς τοῦ κυρίου μου Ἁβραάμ, ὃς οὐκ ἐγκατέλειπε τὴν δικαιοσύνην αὐτοῦ, καὶ τὴν ἀλήθειαν, ἀπὸ τοῦ κυρίου μου· ἐμὲ τʼ εὐώδωκε Κύριος εἰς οἶκον τοῦ ἀδελφοῦ τοῦ κυρίου μου.
\vs{28}Καὶ δραμοῦσα ἡ παῖς ἀνήγγειλεν εἰς τὸν οἶκον τῆς μητρὸς αὐτῆς, κατὰ τὰ ῥήματα ταῦτα.
\vs{29}Τῇ δὲ Ῥεβέκκᾷ ἀδελφὸς ἦν, ᾧ ὄνομα Λάβαν· καὶ ἔδραμε Λάβαν πρὸς τὸν ἄνθρωπον ἔξω ἐπὶ τὴν πηγήν.
\vs{30}Καὶ ἐγένετο ἡνίκα εἶδε τὰ ἐνώτια, καὶ τὰ ψέλλια ἐν ταῖς χερσὶ τῆς ἀδελφῆς αὐτοῦ, καὶ ὅτε ἤκουσε τὰ ῥήματα Ῥεβέκκας τῆς ἀδελφῆς αὐτοῦ, λεγούσης, οὕτω λελάληκέ μοι ὁ ἄνθρωπος, καὶ ἦλθε πρὸς τὸν ἄνθρωπον, ἑστηκότος αὐτοῦ ἐπὶ τῶν καμήλων ἐπὶ τῆς πηγῆς.
\vs{31}Καὶ εἶπεν αὐτῷ, δεῦρο εἴσελθε, εὐλογητὸς Κυροίυ· ἱνατί ἕστηκας ἔξω; ἐγὼ δὲ ἡτοίμασα τὴν οἰκίαν, καὶ τόπον ταῖς καμήλοις.
\vs{32}Εἰσῆλθε δὲ ὁ ἄνθρωπος εἰς τὴν οἰκίαν, καὶ ἀπέσαξε τὰς καμήλους· καὶ ἔδωκεν ἄχυρα καὶ χορτάσματα ταῖς καμήλοις, καὶ ὕδωρ νίψασθαι τοῖς ποσὶν αὐτοῦ, καὶ τοῖς ποσὶ τῶν ἀνδρῶν τῶν μετʼ αὐτοῦ.
\vs{33}Καὶ παρέθηκεν αὐτοῖς ἄρτους φαγεῖν· καὶ εἶπεν, οὐ μὴ φάγω, ἕως τοῦ λαλῆσαί με τὰ ῥήματά μου· καὶ εἶπεν, λάλησον.

\vs{34}Καὶ εἶπε, παῖς Ἁβραὰμ ἐγώ εἰμι.
\vs{35}Κύριος δὲ ηὐλόγησε τὸν κύριόν μου σφόδρα, καὶ ὑψώθη· καὶ ἔδωκεν αὐτῷ πρόβατα, καὶ μόσχους, καὶ ἀργύριον, καὶ χρυσίον, παῖδας, καὶ παιδίσκας, καμήλους, καὶ ὄνους.
\vs{36}Καὶ ἔτεκε Σάῤῥα ἡ γυνὴ τοῦ κυρίου μου υἱὸν ἕνα τῷ κυρίῳ μου μετὰ τὸ γηράσαι αὐτόν· καὶ ἔδωκεν αὐτῷ ὅσα ἦν αὐτῷ.
\vs{37}Καὶ ὥρκισέ με ὁ κύριός μου, λέγων, οὐ λήμψῃ γυναῖκα τῷ υἱῷ μου ἀπὸ τῶν θυγατέρων τῶν Χαναναίων, ἐν οἷς ἐγὼ παροικῶ ἐν τῇ γῇ αὐτῶν.
\vs{38}Ἀλλʼ εἰς τὸν οἶκον τοῦ πατρός μου πορεύσῃ, καὶ εἰς τὴν φυλήν μου, καὶ λήψῃ γυναῖκα τῷ υἱῷ μου ἐκεῖθεν.
\vs{39}Εἶπα δὲ τῷ κυρίῳ μου, μήποτε οὐ πορεύσεται ἡ γυνὴ μετʼ ἐμοῦ.
\vs{40}Καὶ εἶπέ μοι, Κύριος ὁ Θεὸς ᾧ εὐηρέστησα ἐναντίον αὐτοῦ, αὐτὸς ἐξαποστελεῖ τὸν Ἀγγελον αὐτοῦ μετὰ σοῦ, καὶ εὐοδώσει τὴν ὁδόν σου· καὶ λήψῃ γυναῖκα τῷ υἱῷ μου ἐκ τῆς φυλῆς μου, καὶ ἐκ τοῦ οἴκου τοῦ πατρός μου.
\vs{41}Τότε ἀθῷος ἔσῃ ἀπὸ τῆς ἀρᾶς μου· ἡνίκα γὰρ ἐὰν ἔλθῃς εἰς τὴν φυλήν μου, καὶ μή σοι δῶσι, καὶ ἔσῃ ἀθῷος ἀπὸ τοῦ ὁρκισμοῦ μου.
\vs{42}Καὶ ἐλθὼν σήμερον ἐπὶ τὴν πηγὴν εἶπα, Κύριε ὁ Θεὸς τοῦ κυρίου μου Ἁβραὰμ, εἰ σὺ εὐοδοῖς τὴν ὁδόν μου, ἐν ᾗ νῦν ἐγὼ πορεύομαι ἐν αὐτῇ,
\vs{43}ἰδοὺ ἐγὼ ἐφέστηκα ἐπὶ τῆς πηγῆς τοῦ ὕδατος, καὶ αἱ θυγατέρες τῶν ἀνθρώπων τῆς πόλεως ἐκπορεύονται ἀντλῆσαι ὕδωρ· καὶ ἔσται ἡ παρθένος, ᾗ ἂν ἐγὼ εἴπω, πότισόν με ἐκ τῆς ὑδρίας σου μικρὸν ὕδωρ,
\vs{44}καὶ εἴπῃ μοι, καὶ σὺ πίε, καὶ ταῖς καμήλοις σου ὑδρεύσομαι, αὕτη ἡ γυνὴ ἣν ἡτοίμασε Κύριος τῷ ἑαυτοῦ θεράποντι Ἰσαάκ· καὶ ἐν τούτῳ γνώσομαι, ὅτι πεποίηκας ἔλεος τῷ κυρίῳ μου Ἁβραάμ.
\vs{45}Καὶ ἐγένετο πρὸ τοῦ συντελέσαι με λαλοῦντα ἐν τῇ διανοίᾳ μου, εὐθὺς Ῥεβέκκα ἐξεπορεύετο, ἔχουσα τὴν ὑδρίαν ἐπὶ τῶν ὤμων· καὶ κατέβη ἐπὶ τὴν πηγὴν, καὶ ὑδρεύσατο· εἶπα δὲ αὐτῇ, πότισόν με.
\vs{46}Καὶ σπεύσασα καθεῖλε τὴν ὑδρίαν ἐπὶ τὸν βραχίονα αὐτῆς ἀφʼ ἑαυτῆς, καὶ εἶπε, πίε σὺ, καὶ τὰς καμήλους σου ποτιῶ· καὶ ἔπιον, καὶ τὰς καμήλους ἐπότισε.
\vs{47}Καὶ ἠρώτησα αὐτὴν, καὶ εἶπα, θυγάτηρ τίνος εἶ, ἀναγγειλόν μοι· ἡ δὲ ἔφη, θυγάτηρ Βαθουὴλ εἰμὶ υἱοῦ τοῦ Ναχὼρ, ὃν ἔτεκεν αὐτῷ Μελχά· καὶ περιέθηκα αὐτῇ τὰ ἐνώτια, καὶ τὰ ψέλλια περὶ τὰς χεῖρας αὐτῆς.
\vs{48}Καὶ εὐδοκήσας προσεκύνησα τῷ Κυρίῳ, καὶ εὐλόγησα Κύριον τὸν Θεὸν τοῦ κυρίου μου Ἁβραὰμ, ὃς εὐώδωσέ με ἐν ὁδῷ ἀληθείας λαβεῖν τὴν θυγατέρα τοῦ ἀδελφοῦ τοῦ κυρίου μου τῷ υἱῷ αὐτοῦ.
\vs{49}Εἰ οὖν ποιεῖτε ὑμεῖς ἔλεος καὶ δικαιοσύνην πρὸς τὸν κύριόν μου· εἰ δὲ μὴ, ἀπαγγείλατέ μοι, ἵνα ἐπιστρέψω εἰς δεξιὰν ἤ ἀριστεράν.

\vs{50}Ἀποκριθεὶς δὲ Λάβαν καὶ Βαθουὴλ εἶπαν, παρὰ κυρίου ἐξῆλθε τὸ πρᾶγμα τοῦτο· οὐ δυνησόμεθά σοι ἀντειπεῖν κακὸν ἢ καλόν.
\vs{51}Ἰδοὺ Ῥεβέκκα ἐνώπιόν σου· λαβὼν ἀπότρεχε· καὶ ἔστω γυνὴ τῷ υἱῷ τοῦ κυρίου σου, καθὰ ἐλάλησε Κύριος.
\vs{52}Ἐγένετο δὲ ἐν τῷ ἀκοῦσαι τὸν παῖδα τοῦ Ἁβραὰμ τῶν ῥημάτων αὐτῶν, προσεκύνησεν ἐπὶ τὴν γῆν τῷ κυρίῳ.
\vs{53}καὶ ἐξενέγκας ὁ παῖς σκεύη ἀργυρᾶ καὶ χρυσᾶ καὶ ἱματισμὸν, ἔδωκε τῇ Ῥεβέκκᾳ· καὶ δῶρα ἔδωκε τῷ ἀδελφῷ αὐτῆς, καὶ τῇ μητρὶ αὐτῆς.
\vs{54}Καὶ ἔφαγον καὶ ἔπιον καὶ αὐτὸς καὶ οἱ ἄνδρες οἱ μετʼ αὐτοῦ ὄντες, καὶ ἐκοιμήθησαν· καὶ ἀναστὰς τὸ πρωῒ εἶπεν, ἐκπέμψατέ με, ἵνα ἀπέλθω πρὸς τὸν κύριόν μου.
\vs{55}Εἶπαν δὲ οἱ ἀδελφοὶ αὐτῆς, καὶ ἡ μήτηρ, μεινάτω ἡ παρθένος μεθʼ ἡμῶν ἡμέρας ὡσεὶ δέκα, καὶ μετὰ ταῦτα ἀπελεύσεται.
\vs{56}Ὁ δὲ εἶπε πρὸς αὐτοὺς, μὴ κατέχετέ με· καὶ Κύριος εὐώδωσε τὴν ὁδόν μου ἐν ἐμοί· ἐκπέμψατέ με, ἵνα ἀπέλθω πρὸς τὸν κύριόν μου.
\vs{57}Οἱ δὲ εἶπαν, Καλέσωμεν τὴν παῖδα, καὶ ἐρωτήσωμεν τὸ στόμα αὐτῆς.
\vs{58}Καὶ ἐκάλεσαν τὴν Ῥεβέκκαν, καὶ εἶπαν αὐτῇ, πορεύσῃ μετὰ τοῦ ἀνθρώπου τούτου; ἡ δὲ εἶπε, πορεύσομαι.
\vs{59}Καὶ ἐξέπεμψαν Ῥεβέκκαν τὴν ἀδελφὴν αὐτῶν, καὶ τὰ ὑπάρχοντα αὐτῆς, καὶ τὸν παῖδα τοῦ Ἁβραὰμ, καὶ τοὺς μετʼ αὐτοῦ.
\vs{60}Καὶ εὐλόγησαν Ῥεβέκκαν, καὶ εἶπαν αὐτῇ, ἀδελφὴ ἡμῶν εἶ, γίνου εἰς χιλιάδας μυριάδων, καὶ κληρονομησάτω τὸ σπέρμα σου τὰς πόλεις τῶν ὑπεναντίων.
\vs{61}Ἀναστᾶσα δὲ Ῥεβέκκα καὶ αἱ ἅβραι αὐτῆς, ἐπέβησαν ἐπὶ τὰς καμήλους, καὶ ἐπορεύθησαν μετὰ τοῦ ἀνθρώπου· καὶ ἀναλαβὼν ὁ παῖς τὴν Ῥεβέκκαν ἀπῆλθεν.

\vs{62}Ἰσαὰκ δὲ διεπορεύετο διὰ τῆς ἐρήμου κατὰ τὸ φρέαρ τῆς ὁράσεως· αὐτὸς δὲ κατῴκει ἐν τῇ γῇ τῇ πρὸς Λίβα.
\vs{63}Καὶ ἐξῆλθεν Ἰσαὰκ ἀδολεσχῆσαι εἰς τὸ πεδίον τὸ πρὸς δείλης, καὶ ἀναβλέψας τοῖς ὀφθαλμοῖς αὐτοῦ εἶδε καμήλους ἐρχομένας.
\vs{64}Καὶ ἀναβλέψασα Ῥεβέκκα τοῖς ὀφθαλμοῖς εἶδε τὸν Ἰσαάκ· καὶ κατεπήδησεν ἀπὸ τῆς καμήλου.
\vs{65}Καὶ εἶπε τῷ παιδὶ, τίς ἐστιν ὁ ἄνθρωπος ἐκεῖνος ὁ πορευόμενος ἐν τῷ πεδίῳ εἰς συνάντησιν ἡμῖν; εἶπε δὲ ὁ παῖς, οὗτός ἐστιν ὁ κύριός μου· ἡ δὲ λαβοῦσα τὸ θέριστρον, περιεβάλετο.
\vs{66}Καὶ διηγήσατο ὁ παῖς τῷ Ἰσαὰκ πάντα τὰ ῥήματα, ἃ ἐποίησεν.
\vs{67}Εἰσῆλθε δὲ Ἰσαὰκ εἰς τὸν οἶκον τῆς μητρὸς αὐτοῦ, καὶ ἔλαβε τὴν Ῥεβέκκαν, καὶ ἐγένετο αὐτοῦ γυνὴ, καὶ ἠγάπησεν αὐτήν· καὶ παρεκλήθη Ἰσαὰκ περὶ Σάῤῥας τῆς μητρὸς αὐτοῦ.

\ch{25}
Προσθέμενος δὲ Ἁβραὰμ ἔλαβε γυναῖκα, ᾗ ὄνομα Χεττούρα.
\vs{2}Ἔτεκε δὲ αὐτῷ τὸν Ζομβρᾶν, καὶ τὸν Ἰεζὰν, καὶ τὸν Μαδὰλ, καὶ τὸν Μαδιὰμ, καὶ τὸν Ἰεσβὼκ, καὶ τὸν Σωίε.
\vs{3}Ἰεζὰν δὲ ἐγέννησε τὸν Σαβὰ, καὶ τὸν Δεδάν· υἱοὶ δὲ Δεδὰν Ἀσσουριεὶμ, καὶ Λατουσιεὶμ, καὶ Λαωμείμ.
\vs{4}Υἱοὶ δὲ Μαδιὰμ Γεφὰρ, καὶ Ἀφεὶρ, καὶ Ἐνὼχ, καὶ Ἀβειδὰ, καὶ Ἐλδαγά· πάντες οὗτοι ἦσαν υἱοὶ Χεττούρας.
\vs{5}Ἔδωκε δὲ Ἁβραὰμ πάντα τὰ ὑπάρχοντα αὐτοῦ Ἰσαὰκ τῷ υἱῷ αὐτοῦ.
\vs{6}Καὶ τοῖς υἱοῖς τῶν παλλακῶν αὐτοῦ ἔδωκεν Ἁβραὰμ δόματα, καὶ ἐξαπέστειλεν αὐτοὺς ἀπὸ Ἰσαὰκ τοῦ υἱοῦ αὐτοῦ, ἔτι ζῶντος αὐτοῦ, πρὸς ἀνατολὰς εἰς γῆν ἀνατολῶν.
\vs{7}Ταῦτα δὲ τὰ ἔτη ἡμερῶν τῆς ζωῆς Ἁβραὰμ ὅσα ἔζησεν, ἑκατὸν ἑβδομηκονταπεντε ἔτη.
\vs{8}Καὶ ἐκλείπων ἀπέθανεν Ἁβραὰμ ἐν γήρᾳ καλῷ πρεσβύτης, καὶ πλήρης ἡμερῶν, καὶ προσετέθη πρὸς τὸν λαὸν αὐτοῦ.
\vs{9}Καὶ ἔθαψαν αὐτὸν Ἰσαὰκ καὶ Ἰσμαὴλ οἱ υἱοὶ αὐτοῦ εἰς τὸ σπήλαιον τὸ διπλοῦν, εἰς τὸν ἀγρὸν Ἐφρων τοῦ Σαὰρ τοῦ Χετταίου, ὅς ἐστιν ἀπέναντι Μαμβρῆ,
\vs{10}τὸν ἀγρὸν καὶ τὸ σπήλαιον, ὃ ἐκτήσατο Ἁβραὰμ παρὰ τῶν υἱῶν τοῦ Χέτ· ἐκεῖ ἔθαψαν Ἁβραὰμ, καὶ Σάῤῥαν τὴν γυναῖκα αὐτοῦ.
\vs{11}Ἐγένετο δὲ μετὰ τὸ ἀποθανεῖν Ἁβραὰμ, εὐλόγησεν ὁ Θεὸς τὸν Ἰσαὰκ υἱὸν αὐτοῦ· καὶ κατῴκησεν Ἰσαὰκ παρὰ τὸ φρέαρ τῆς ὁράσεως.
\vs{12}Αὗται δὲ αἱ γενέσεις Ἰσμαὴλ τοῦ υἱοῦ Ἁβραὰμ, ὃν ἔτεκεν Ἄγαρ ἡ Αἰγυπτία, ἡ παιδίσκη Σάῤῥας, τῷ Ἁβραάμ.
\vs{13}Καὶ ταῦτα τὰ ὀνόματα τῶν υἱῶν Ἰσμαὴλ, κατʼ ὀνόματα τῶν γενεῶν αὐτοῦ· πρωτότοκος Ἰσμαὴλ, καὶ Ναβαϊὼθ, καὶ Κηδὰρ, καὶ Ναβδεὴλ, καὶ Μασσὰμ,
\vs{14}καὶ Μασμὰ, καὶ Δουμὰ, καὶ Μασσῆ,
\vs{15}καὶ Χοδδὰν, καὶ Θαιμὰν, καὶ Ἰετοὺρ, καὶ Ναφὲς, καὶ Κεδμά.
\vs{16}οὗτοί εἰσιν οἱ υἱοὶ Ἰσμαὴλ, καὶ ταῦτα τὰ ὀνόματα αὐτῶν ἐν ταῖς σκηναῖς αὐτῶν, καὶ ἐν ταῖς ἐπαύλεσιν αὐτῶν· δώδεκα ἄρχοντες κατὰ ἔθνη αὐτῶν.
\vs{17}Καὶ ταῦτα τὰ ἔτη τῆς ζωῆς Ἰσμαὴλ, ἑκατὸν τριακονταεπτὰ ἔτη· καὶ ἐκλείπων ἀπέθανε, καὶ προσετέθη πρὸς τὸ γένος αὐτοῦ.
\vs{18}Κατῴκησε δὲ ἀπὸ Εὐϊλὰτ ἕως Σοὺρ, ἥ ἐστι κατὰ πρόσωπον Αἰγύπτου ἕως ἐλθεῖν πρὸς Ἀσσυρίους· κατὰ πρόσωπον πάντων τῶν ἀδελφῶν αὐτοῦ κατῴκησε.

\vs{19}Καὶ αὗται αἱ γενέσεις Ἰσαὰκ τοῦ υἱοῦ Ἁβραάμ· Ἁβραάμ ἐγέννησε τὸν Ἰσαάκ.
\vs{20}Ἦν δὲ Ἰσαὰκ ἐτῶν τεσσαράκοντα ὅτε ἔλαβε τὴν Ῥεβέκκαν θυγατέρα Βαθουὴλ τοῦ Σύρου ἐκ τῆς Μεσοποταμίας Συρίας, ἀδελφὴν Λάβαν τοῦ Σύρου, ἑαυτῷ εἰς γυναῖκα.
\vs{21}Ἐδέετο δὲ Ἰσαὰκ Κυρίου περὶ Ῥεβέκκας τῆς γυναικὸς αὐτοῦ, ὅτι στεῖρα ἦν· ἐπήκουσε δὲ αὐτοῦ ὁ Θεὸς, καὶ συνέλαβεν ἐν γαστρὶ Ῥεβέκκα ἡ γυνὴ αὐτοῦ.
\vs{22}Ἐσκίρτων δὲ τὰ παιδία ἐν αὐτῇ· εἶπε δὲ, εἰ οὕτω μοι μέλλει γίνεσθαι, ἵνα τί μοι τοῦτο; ἐπορεύθη δὲ πυθέσθαι παρὰ Κυρίου.
\vs{23}Καὶ εἶπε Κύριος αὐτῇ, δύο ἔθνη ἐν γαστρί σου εἰσὶ, καὶ δύο λαοὶ ἐκ τῆς κοιλίας σου διασταλήσονται· καὶ λαὸς λαοῦ ὑπερέξει, καὶ ὁ μείζων δουλεύσει τῷ ἐλάσσονι.
\vs{24}Καὶ ἐπληρώθησαν αἱ ἡμέραι τοῦ τεκεῖν αὐτήν· καὶ τῇδε ἦν δίδυμα ἐν τῇ κοιλίᾳ αὐτῆς.
\vs{25}Ἐξῆλθε δὲ ὁ πρωτότοκος πυῤῥάκης· ὅλος, ὡσεὶ δορὰ, δασύς· ἐπωνόμασε δὲ τὸ ὄνομα αὐτοῦ, Ἡσαῦ.
\vs{26}Καὶ μετὰ τοῦτο ἐξῆλθεν ὁ ἀδελφὸς αὐτοῦ, καὶ ἡ χεὶρ αὐτοῦ ἐπειλημμένη τῆς πτέρνης Ἡσαῦ· καὶ ἐκάλεσε τὸ ὄνομα αὐτοῦ, Ἰακώβ. Ἰσαὰκ δὲ ἦν ἐτῶν ἑξήκοντα, ὅτε ἔτεκεν αὐτοὺς Ῥεβέκκα.
\vs{27}Ηὐξήθησαν δὲ οἱ νεανίσκοι· καὶ ἦν Ἡσαῦ ἄνθρωπος εἰδὼς κυνηγεῖν, ἄγροικος· Ἰακὼβ δὲ ἄνθρωπος ἄπλαστος, οἰκῶν οἰκίαν.
\vs{28}Ἠγάπησε δὲ Ἰσαὰκ τὸν Ἡσαῦ, ὅτι ἡ θήρα αὐτοῦ βρῶσις αὐτῷ· Ῥεβέκκα δὲ ἠγάπα τὸν Ἰακώβ.

\vs{29}Ἥψησε δὲ Ἰακὼβ ἕψημα· ἦλθε δὲ Ἡσαῦ ἐκ τοῦ πεδίου ἐκλείπων.
\vs{30}Καὶ εἶπεν Ἡσαῦ τῷ Ἰακὼβ, γεῦσόν με ἀπὸ τοῦ ἑψήματος πυῤῥου τούτου, ὅτι ἐκλείπω· διὰ τοῦτο ἐκλήθη τὸ ὄνομα αὐτοῦ, Ἐδώμ.
\vs{31}Εἶπε δὲ Ἰακὼβ τῷ Ἡσαῦ, ἀπόδου μοι σήμερον τὰ πρωτοτόκιά σου ἐμοί.
\vs{32}Καὶ εἶπεν Ἡσαῦ, ἰδοὺ ἐγὼ πορεύομαι τελευτᾷν· καὶ ἵνα τί μοι ταῦτα τὰ πρωτοτόκια;
\vs{33}Καὶ εἶπεν αὐτῷ Ἰακὼβ, ὄμοσόν μοι σήμερον· καὶ ὤμοσεν αὐτῷ· ἀπέδοτο δὲ Ἡσαῦ τὰ πρωτοτόκια τῷ Ἰακώβ.
\vs{34}Ἰακὼβ δὲ ἔδωκε τῷ Ἠσαῦ ἄρτον, καὶ ἕψημα φακοῦ· καὶ ἔφαγε καὶ ἔπιε, καὶ ἀναστὰς ᾤχετο· καὶ ἐφαύλισεν Ἡσαῦ τὰ πρωτοτόκια.

\ch{26}
Ἐγένετο δὲ λιμὸς ἐπὶ τῆς γῆς, χωρὶς τοῦ λιμοῦ τοῦ πρότερον, ὃς ἐγένετο ἐν τῷ καιρῷ τοῦ Ἁβραάμ· ἐπορεύθη δὲ Ἰσαὰκ πρὸς Ἀβιμέλεχ βασιλέα Φυλιστιεὶμ εἰς Γέραρα.
\vs{2}Ὤφθη δὲ αὐτῷ Κύριος, καὶ εἶπε, μὴ καταβῇς εἰς Αἴγυπτον· κατοίκησον δὲ ἐν τῇ γῇ, ᾗ ἄν σοι εἴπω.
\vs{3}Καὶ παροίκει ἐν τῇ γῇ ταύτῃ, καὶ ἔσομαι μετὰ σοῦ, καὶ εὐλογήσω σε· σοὶ γὰρ καὶ τῷ σπέρματί σου δώσω πᾶσαν τὴν γῆν ταύτην· καὶ στήσω τὸν ὅρκον μου, ὅν ὤμοσα τῷ Ἁβραὰμ τῷ πατρί σου.
\vs{4}Καὶ πληθυνῶ τὸ σπέρμα σου, ὡς τοὺς ἀστέρας τοῦ οὐρανοῦ· καὶ δώσω τῷ σπέρματί σου πᾶσαν τὴν γῆν ταύτην· καὶ εὐλογηθήσονται ἐν τῷ σπέρματί σου πάντα τὰ ἔθνη τῆς γῆς.
\vs{5}Ἀνθʼ ὧν ὑπήκουσεν Ἁβραὰμ ὁ πατήρ σου τῆς ἐμῆς φωνῆς, καὶ ἐφύλαξε τὰ προστάγματά μου, καὶ τὰς ἐντολάς μου, καὶ τὰ δικαιώματά μου, καὶ τὰ νόμιμά μου.
\vs{6}Κατῴκησε δὲ Ἰσαὰκ ἐν Γεράροις.
\vs{7}Ἐπηρώτησαν δὲ οἱ ἄνδρες τοῦ τόπου περὶ Ῥεβέκκας τῆς γυναικὸς αὐτοῦ, καὶ εἶπεν, ἀδελφή μου ἐστίν· ἐφοβήθη γὰρ εἰπεῖν, ὅτι γυνή μου ἐστὶ, μή ποτε ἀποκτείνωσιν αὐτὸν οἱ ἄνδρες τοῦ τόπου περὶ Ῥεβέκκας, ὅτι ὡραία τῇ ὄψει ἦν.
\vs{8}Ἐγένετο δὲ πολυχρόνιος ἐκεῖ· καὶ παρακύψας Ἀβιμέλεχ ὁ βασιλεὺς Γεράρων διὰ τῆς θυρίδος, εἶδε τὸν Ἰσαὰκ παίζοντα μετὰ Ῥεβέκκας τῆς γυναικὸς αὐτοῦ.
\vs{9}Ἐκάλεσε δὲ Ἀβιμέλεχ τὸν Ἰσαὰκ, καὶ εἶπεν αὐτῷ, ἆρά γε γυνή σου ἐστί; τί ὅτι εἶπας, ἀδελφή μου ἐστίν; εἶπε δὲ αὐτῷ Ἰσαὰκ, εἶπα γὰρ, μή ποτε ἀποθάνω διʼ αὐτήν.
\vs{10}Εἶπε δὲ αὐτῷ Ἀβιμέλεχ, τί τοῦτο ἐποίησας ἡμῖν; μικροῦ ἐκοιμήθη τις ἐκ τοῦ γένους μου μετὰ τῆς γυναικός σου, καὶ ἐπήγαγες ἂν ἐφʼ ἡμᾶς ἄγνοιαν.
\vs{11}Συνέταξε δὲ Ἀβιμέλεχ παντὶ τῷ λαῷ αὐτοῦ, λέγων, πᾶς ὁ ἁψάμενος τοῦ ἀνθρώπου τούτου καὶ τῆς γυναικὸς αὐτοῦ, θανάτῳ ἔνοχος ἔσται.
\vs{12}Ἔσπειρε δὲ Ἰσαὰκ ἐν τῇ γῇ ἐκείνῃ, καὶ εὗρεν ἐν τῷ ἐνιαυτῷ ἐκείνῳ ἑκατοστεύουσαν κριθήν· εὐλόγησε δὲ αὐτὸν Κύριος.
\vs{13}Καὶ ὑψώθη ὁ ἄνθρωπος, καὶ προβαίνων μείζων ἐγένετο, ἕως οὗ μέγας ἐγένετο σφόδρα.
\vs{14}Ἐγένετο δὲ αὐτῷ κτήνη προβάτων, καὶ κτήνη βοῶν, καὶ γεώργια πολλά. ἐζήλωσαν δὲ αὐτὸν οἱ Φυλιστιείμ.
\vs{15}Καὶ πάντα τὰ φρέατα, ἃ ὤρυξαν οἱ παῖδες τοῦ πατρὸς αὐτοῦ ἐν τῷ χρόνῳ τοῦ πατρὸς αὐτοῦ, ἐνέφραξαν αὐτὰ οἱ Φυλιστιεὶμ, καὶ ἔπλησαν αὐτὰ γῆς.
\vs{16}Εἶπε δὲ Ἀβιμέλεχ πρὸς Ἰσαὰκ, ἄπελθε ἀφʼ ἡμῶν, ὅτι δυνατώτερος ἡμῶν ἐγένου σφόδρα.
\vs{17}Καὶ ἀπῆλθεν ἐκεῖθεν Ἰσαάκ· καὶ κατέλυσεν ἐν τῇ φάραγγι Γεράρων, καὶ κατῴκησεν ἐκεῖ.

\vs{18}Καὶ πάλιν Ἰσαὰκ ὤρυξε τὰ φρέατα τοῦ ὕδατος, ἃ ὤρυξαν οἱ παῖδες Ἁβραὰμ τοῦ πατρὸς αὐτοῦ, καὶ ἐνέφραξαν αὐτὰ οἱ Φυλιστιεὶμ μετὰ τὸ ἀποθανεῖν Ἁβραὰμ τὸν πατέρα αὐτοῦ· καὶ ἐπωνόμασεν αὐτοῖς ὀνόματα κατὰ τὰ ὀνόματα, ἃ ὠνόμασεν ὁ πατὴρ αὐτοῦ.
\vs{19}Καὶ ὤρυξαν οἱ παῖδες Ἰσαὰκ ἐν τῇ φάραγγι Γεράρων· καὶ εὗρον ἐκεῖ φρέαρ ὕδατος ζῶντος.
\vs{20}Καὶ ἐμαχέσαντο οἱ ποιμένες Γεράρων μετὰ τῶν ποιμένων Ἰσαὰκ, φάσκοντες αὐτῶν εἶναι τὸ ὕδωρ· καὶ ἐκάλεσαν τὸ ὄνομα τοῦ φρέατος, Ἀδικία· ἠδίκησαν γὰρ αὐτόν.
\vs{21}Ἀπᾴρας δὲ ἐκεῖθεν ὤρυξε φρέαρ ἕτερον· ἐκρίνοντο δὲ καὶ περὶ ἐκείνου· καὶ ἐπωνόμασε τὸ ὄνομα αὐτοῦ, Ἐχθρία.
\vs{22}Ἀπᾴρας δὲ ἐκεῖθεν ὤρυξε φρέαρ ἕτερον· καὶ οὐκ ἐμαχέσαντο περὶ αὐτοῦ· καὶ ἐπωνόμασε τὸ ὄνομα αὐτοῦ, Εὐρυχωρία, λέγων, διότι νῦν ἐπλάτυνε Κύριος ἡμῖν, καὶ ηὔξησεν ἡμᾶς ἐπὶ τῆς γῆς.

\vs{23}Ἀνέβη δὲ ἐκεῖθεν ἐπὶ τὸ φρέαρ τοῦ ὅρκου.
\vs{24}Καὶ ὤφθη αὐτῷ Κύριος ἐν τῇ νυκτὶ ἐκείνῃ, καὶ εἶπεν, ἐγώ εἰμι ὁ Θεὸς Ἁβραὰμ τοῦ πατρός σου· μὴ φοβοῦ, μετὰ σοῦ γάρ εἰμι, καὶ εὐλογήσω σε, καὶ πληθυνῶ τὸ σπέρμα σου διʼ Ἁβραὰμ τὸν πατέρα σου.
\vs{25}Καὶ ᾠκοδόμησεν ἐκεῖ θυσιαστήριον, καὶ ἐπεκαλέσατο τὸ ὄνομα Κυρίου, καὶ ἔπηξεν ἐκεῖ τὴν σκηνὴν αὐτοῦ· ὤρυξαν δὲ ἐκεῖ οἱ παῖδες Ἰσαὰκ φρέαρ ἐν τῇ φάραγγι Γεράρων.
\vs{26}Καὶ Ἀβιμέλεχ ἐπορεύθη πρὸς αὐτὸν ἀπὸ Γεράρων, καὶ Ὁχοζὰθ ὁ νυμφαγωγὸς αὐτοῦ, καὶ Φιχὼλ ὁ ἀρχιστράτηγος τῆς δυνάμεως αὐτοῦ.
\vs{27}Καὶ εἶπεν αὐτοῖς Ἰσαὰκ, ἵνα τί ἤλθετε πρός με; ὑμεῖς δὲ ἐμισήσατέ με, καὶ ἐξαπεστείλατέ με ἀφʼ ὑμῶν.
\vs{28}Οἱ δὲ εἶπαν, ἰδόντες ἑωράκαμεν ὅτι ἦν Κύριος μετὰ σοῦ· καὶ εἴπαμεν, γενέσθω ἀρὰ ἀνὰ μέσον ἡμῶν καὶ ἀνὰ μέσον σου, καὶ διαθησόμεθα μετὰ σοῦ διαθήκην,
\vs{29}Μὴ ποιήσαι μεθʼ ἡμῶν κακὸν, καθότι οὐκ ἐβδελυξάμεθά σε ἡμεῖς, καὶ ὃν τρόπον ἐχρησάμεθά σοι καλῶς, καὶ ἐξαπεστείλαμέν σε μετʼ εἰρήνης· καὶ νῦν εὐλογημένος σὺ ὑπὸ Κυρίου.
\vs{30}Καὶ ἐποίησεν αὐτοῖς δοχὴν, καὶ ἔφαγον καὶ ἔπιον.
\vs{31}Καὶ ἀναστάντες τὸ πρωῒ, ὤμοσεν ἕκαστος τῷ πλησίον· καὶ ἐξαπέστειλεν αὐτοὺς Ἰσαάκ· καὶ ἀπῴχοντο ἀπʼ αὐτοῦ μετὰ σωτηρίας.
\vs{32}Ἐγένετο δὲ ἐν τῇ ἡμέρᾳ ἐκείνῃ, καὶ παραγενόμενοι οἱ παῖδες Ἰσαὰκ ἀπήγγειλαν αὐτῷ περὶ τοῦ φρέατος οὗ ὤρυξαν, καὶ εἶπαν, οὐχ εὕρομεν ὕδωρ.
\vs{33}Καὶ ἐκάλεσεν αὐτὸ, Ὅρκος· διὰ τοῦτο ἐκάλεσεν ὄνομα τῇ πόλει ἐκείνῃ, Φρέαρ Ὅρκου, ἕως τῆς σήμερον ἡμέρας.

\vs{34}Ἦν δὲ Ἡσαῦ ἐτῶν τεσσαράκοντα, καὶ ἔλαβε γυναῖκα Ἰουδὶθ, θυγατέρα Βεὼχ τοῦ Χετταίου, καὶ τὴν Βασεμὰθ, θυγατέρα Ἑλὼν Χετταίου.
\vs{35}Καὶ ἦσαν ἐρίζουσαι τῷ Ἰσαὰκ καὶ τῇ Ῥεβέκκᾳ.

\ch{27}
Ἐγένετο δὲ μετὰ τὸ γηράσαι τὸν Ἰσαὰκ, καὶ ἠμβλύνθησαν οἱ ὀφθαλμοὶ αὐτοῦ τοῦ ὁρᾷν, καὶ ἐκάλεσεν Ἡσαῦ τὸν υἱὸν αὐτοῦ τὸν πρεσβύτερον, καὶ εἶπεν αὐτῷ, υἱέ μου· καὶ εἶπεν, ἰδοὺ ἐγώ.
\vs{2}Καὶ εἶπεν, ἰδοὺ γεγήρακα, καὶ οὐ γινώσκω τὴν ἡμέραν τῆς τελευτῆς μου.
\vs{3}Νῦν οὖν λάβε τὸ σκεῦός σου, τήν τε φαρέτραν, καὶ τὸ τόξον, καὶ ἔξελθε εἰς τὸ πεδίον, καὶ θήρευσόν μοι θήραν.
\vs{4}Καὶ ποίησόν μοι ἐδέσματα, ὡς φιλῶ ἐγὼ, καὶ ἔνεγκέ μοι, ἵνα φάγω, ὅπως εὐλογήσῃ σε ἡ ψυχή μου πρὶν ἀποθανεῖν με.
\vs{5}Ῥεβέκκα δὲ ἤκουσε λαλοῦντος Ἰσαὰκ πρὸς Ἡσαῦ τὸν υἱὸν αὐτοῦ· ἐπορεύθη δὲ Ἡσαῦ εἰς τὸ πεδίον θηρεῦσαι θήραν τῷ πατρὶ αὐτοῦ.
\vs{6}Ῥεβέκκα δὲ εἶπε πρὸς τὸν Ἰακὼβ τὸν υἱὸν αὐτῆς τὸν ἐλάσσω, ἴδε, ἤκουσα τοῦ πατρός σου λαλοῦντος πρὸς Ἡσαῦ τὸν ἀδελφόν σου, λέγοντος,
\vs{7}Ἔνεγκόν μοι θήραν, καὶ ποίησόν μοι ἐδέσματα, ἵνα φαγὼν εὐλογήσω σε ἐναντίον Κυρίου πρὸ τοῦ ἀποθανεῖν με.
\vs{8}Νῦν οὖν, υἱέ μου, ἄκουσόν μου, καθὰ ἐγώ σοι ἐντέλλομαι.
\vs{9}Καὶ πορευθεὶς εἰς τὰ πρόβατα, λάβε μοι ἐκεῖθεν δύο ἐρίφους ἁπαλοὺς καὶ καλοὺς, καὶ ποιήσω αὐτοὺς ἐδέσματα τῷ πατρί σου, ὡς φιλεῖ.
\vs{10}Καὶ εἰσοίσεις τῷ πατρί σου, καὶ φάγεται, ὅπως εὐλογήσῃ σε ὁ πατήρ σου πρὸ τοῦ ἀποθανεῖν αὐτόν.
\vs{11}Εἶπε δὲ Ἰακὼβ πρὸς Ῥεβέκκαν τὴν μητέρα αὐτοῦ, ἔστιν Ἡσαῦ ὁ ἀδελφός μου ἀνὴρ δασὺς, ἐγὼ δὲ ἀνὴρ λεῖος.
\vs{12}Μή ποτε ψηλαφήσῃ με ὁ πατὴρ, καὶ ἔσομαι ἐναντίον αὐτοῦ ὡς καταφρονῶν, καὶ ἐπάξω ἐπʼ ἐμαυτὸν κατάραν, καὶ οὐκ εὐλογίαν.
\vs{13}Εἶπε δὲ αὐτῷ ἡ μήτηρ, ἐπʼ ἐμὲ ἡ κατάρα σου, τέκνον· μόνον ἐπάκουσόν μου τῆς φωνῆς, καὶ πορευθεὶς ἔνεγκέ μοι.
\vs{14}Πορευθεὶς δὲ ἔλαβε, καὶ ἤνεγκε τῇ μητρί· καὶ ἐποίησεν ἡ μήτηρ αὐτοῦ ἐδέσματα, καθὰ ἐφίλει ὁ πατὴρ αὐτοῦ.

\vs{15}Καὶ λαβοῦσα Ῥεβέκκα τὴν στολὴν Ἡσαῦ τοῦ υἱοῦ αὐτῆς τοῦ πρεσβυτέρου τὴν καλὴν, ἣ ἦν παρʼ αὐτῇ ἐν τῷ οἴκῳ, ἐνέδυσεν αὐτὴν Ἰακὼβ τὸν υἱὸν αὐτῆς τὸν νεώτερον.
\vs{16}Καὶ τὰ δέρματα τῶν ἐρίφων περιέθηκεν ἐπὶ τοὺς βραχίονας αὐτοῦ, καὶ ἐπὶ τὰ γυμνὰ τοῦ τραχήλου αὐτοῦ.
\vs{17}Καὶ ἔδωκε τὰ ἐδέσματα, καὶ τοὺς ἄρτους οὓς ἐποίησεν, εἰς τὰς χεῖρας Ἰακὼβ τοῦ υἱοῦ αὐτῆς.
\vs{18}Καὶ εἰσήνεγκε τῷ πατρὶ αὐτοῦ· εἶπε δὲ, πάτερ· ὁ δὲ εἶπεν, ἰδοὺ ἐγώ· τίς εἶ σὺ, τέκνον;
\vs{19}Καὶ εἶπεν Ἰακὼβ τῷ πατρὶ, ἐγὼ Ἡσαῦ ὁ πρωτότοκός σου πεποίηκα καθὰ ἐλάλησάς μοι· ἀναστὰς κάθισον, καὶ φάγε ἀπὸ τῆς θήρας μου, ὅπως εὐλογήσῃ με ἡ ψυχή σου.
\vs{20}Εἶπε δὲ Ἰσαὰκ τῷ υἱῷ αὐτοῦ, τί τοῦτο, ὃ ταχὺ εὗρες, ὦ τέκνον; ὁ δὲ εἶπεν, ὃ παρέδωκε Κύριος ὁ Θεός σου ἐναντίον μου.
\vs{21}Εἶπε δὲ Ἰσαὰκ τῷ Ἰακὼβ, ἔγγισόν μοι, καὶ ψηλαφήσω σε, τέκνον, εἰ σὺ εἶ ὁ υἱός μου Ἡσαῦ, ἢ οὔ.
\vs{22}Ἤγγισε δὲ Ἰακὼβ πρὸς Ἰσαὰκ τὸν πατέρα αὐτοῦ· καὶ ἐψηλάφησεν αὐτὸν, καὶ εἶπεν, ἡ μὲν φωνὴ, φωνὴ Ἰακὼβ, αἱ δὲ χεῖρες, χεῖρες Ἡσαῦ.
\vs{23}Καὶ οὐκ ἐπέγνω αὐτὸν, ἦσαν γὰρ αἱ χεῖρες αὐτοῦ, ὡς αἱ χεῖρες Ἡσαῦ τοῦ ἀδελφοῦ αὐτοῦ, δασεῖαι· καὶ εὐλόγησεν αὐτὸν,
\vs{24}καὶ εἶπε, σὺ εἶ ὁ υἱός μου Ἡσαῦ; ὁ δὲ εἶπεν, ἐγώ.
\vs{25}Καὶ εἶπε, προσάγαγέ μοι, καὶ φάγομαι ἀπὸ τῆς θήρας σου, τέκνον, ἵνα εὐλογήσῃ σε ἡ ψυχή μου· καὶ προσήνεγκεν αὐτῷ, καὶ ἔφαγε· καὶ εἰσήνεγκεν αὐτῷ οἶνον, καὶ ἔπιε.
\vs{26}Καὶ εἴπεν αὐτῷ Ἰσαὰκ ὁ πατὴρ αὐτοῦ, ἔγγισόν μοι, καὶ φίλησόν με, τέκνον.
\vs{27}Καὶ ἐγγίσας ἐφίλησεν αὐτόν· καὶ ὠσφράνθη τὴν ὀσμὴν τῶν ἱματίων αὐτοῦ, καὶ εὐλόγησεν αὐτὸν, καὶ εἶπεν, ἰδοὺ ὀσμὴ τοῦ υἱοῦ μου, ὡς ὀσμὴ ἀγροῦ πλήρους, ὃν εὐλόγησε Κύριος.
\vs{28}Καὶ δῴη σοι ὁ Θεὸς ἀπὸ τῆς δρόσου τοῦ οὐρανοῦ, καὶ ἀπὸ τῆς πιότητος τῆς γῆς, καὶ πλῆθος σίτου καὶ οἴνου.
\vs{29}Καὶ δουλευσάτωσάν σοι ἔθνη, καὶ προσκυνησάτωσάν σοι ἄρχοντες· καὶ γίνου κύριος τοῦ ἀδελφοῦ σου, καὶ προσκυνήσουσί σοι οἱ υἱοὶ τοῦ πατρός σου· ὁ καταρώμενός σε, ἐπικατάρατος· ὁ δὲ εὐλογῶν σε, εὐλογημένος.

\vs{30}Καὶ ἐγένετο μετὰ τὸ παύσασθαι Ἰσαὰκ εὐλογοῦντα Ἰακὼβ τὸν υἱὸν αὐτοῦ, καὶ ἐγένετο, ὡς ἂν ἐξῆλθεν Ἰακὼβ ἀπὸ προσώπου Ἰσαὰκ τοῦ πατρὸς αὐτοῦ, καὶ Ἡσαῦ ὁ ἀδελφὸς αὐτοῦ ἦλθεν ἀπὸ τῆς θήρας.
\vs{31}Καὶ ἐποίησε καὶ αὐτὸς ἐδέσματα, καὶ προσήνεγκε τῷ πατρὶ αὐτοῦ· καὶ εἶπε τῷ πατρὶ, ἀναστήτω ὁ πατήρ μου, καὶ φαγέτω ἀπὸ τῆς θήρας τοῦ υἱοῦ αὐτοῦ, ὅπως εὐλογήσῃ με ἡ ψυχή σου.
\vs{32}Καὶ εἶπεν αὐτῷ Ἰσαὰκ ὁ πατὴρ αὐτοῦ, τίς εἶ σύ; ὁ δὲ εἶπεν, ἐγώ εἰμι ὁ υἱός σου ὁ πρωτότοκος Ἡσαῦ.
\vs{33}Ἐξέστη δὲ Ἰσαὰκ ἔκστασιν μεγάλην σφόδρα, καὶ εἶπε, τίς οὖν ὁ θηρεύσας μοι θήραν καὶ εἰσενέγκας μοι, καὶ ἔφαγον ἀπὸ πάντων πρὸ τοῦ ἐλθεῖν σε; καὶ εὐλόγησα αὐτὸν, καὶ εὐλογημένος ἔσται.
\vs{34}Ἐγένετο δὲ ἡνίκα ἤκουσεν Ἡσαῦ τὰ ῥήματα τοῦ πατρὸς αὐτοῦ Ἰσαὰκ, ἀνεβόησε φωνὴν μεγάλην καὶ πικρὰν σφόδρα· καὶ εἶπεν, εὐλόγησον δὴ κᾀμὲ, πάτερ.
\vs{35}Εἶπε δὲ αὐτῷ, ἐλθὼν ὁ ἀδελφός σου μετὰ δόλου ἔλαβε τὴν εὐλογίαν σου.
\vs{36}Καὶ εἶπε, δικαίως ἐκλήθη τὸ ὄνομα αὐτοῦ Ἰακὼβ, ἐπτέρνικε γάρ με ἰδοὺ δεύτερον τοῦτο· τά τε πρωτοτόκιά μου εἴληφε, καὶ νῦν ἔλαβε τὴν εὐλογίαν μου· καὶ εἶπεν Ἡσαῦ τῷ πατρὶ αὐτοῦ, οὐχ ὑπελίπου μοι εὐλογίαν, πάτερ;
\vs{37}Ἀποκριθεὶς δὲ Ἰσαὰκ εἶπε τῷ Ἡσαῦ, εἰ κύριον αὐτὸν πεποίηκά σου, καὶ πάντας τοὺς ἀδελφοὺς αὐτοῦ πεποίηκα αὐτοῦ οἰκέτας· σίτῳ καὶ οἴνῳ ἐστήριξα αὐτόν· σοὶ δὲ τί ποιήσω, τέκνον;
\vs{38}Εἶπε δὲ Ἡσαῦ πρὸς τὸν πατέρα αὐτοῦ, μὴ εὐλογία μία σοι ἔστι, πάτερ; εὐλόγησον δὴ κᾀμὲ, πάτερ· κατανυχθέντος δὲ Ἰσαὰκ, ἀνεβόησε φωνῇ Ἡσαῦ, καὶ ἔκλαυσεν.
\vs{39}Ἀποκοιθεὶς δὲ Ἰσαὰκ ὁ πατὴρ αὐτοῦ εἶπεν αὐτῷ, ἰδοὺ ἀπὸ τῆς πιότητος τῆς γῆς ἔσται ἡ κατοίκησίς σου, καὶ ἀπὸ τῆς δρόσου τοῦ οὐρανοῦ ἄνωθεν.
\vs{40}Καὶ ἐπὶ τῇ μαχαίρᾳ σου ζήσῃ, καὶ τῷ ἀδελφῷ σου δουλεύσεις· ἔσται δὲ ἡνίκα ἐὰν καθέλῃς καὶ ἐκλύσῃς τὸν ζυγὸν αὐτοῦ ἀπὸ τοῦ τραχήλου σου.

\vs{41}Καὶ ἐνεκότει Ἡσαῦ τῷ Ἰακὼβ περὶ τῆς εὐλογίας, ἧς εὐλόγησεν αὐτὸν ὁ πατὴρ αὐτοῦ· εἶπε δὲ Ἡσαῦ ἐν τῇ διανοίᾳ αὐτοῦ, ἐγγισάτωσαν αἱ ἡμέραι τοῦ πένθους τοῦ πατρός μου, ἵνα ἀποκτείνω Ἰακὼβ τὸν ἀδελφόν μου.
\vs{42}Ἀπηγγέλη δὲ Ῥεβέκκᾳ τὰ ῥήματα Ἡσαῦ τοῦ υἱοῦ αὐτῆς τοῦ πρεσβυτέρου· καὶ πέμψασα ἐκάλεσεν Ἰακὼβ τὸν υἱὸν αὐτῆς τὸν νεώτερον, καὶ εἶπεν αὐτῷ, ἰδοὺ Ἡσαῦ ὁ ἀδελφός σου ἀπειλεῖ σοι τοῦ ἀποκτεῖναί σε.
\vs{43}Νῦν οὖν, τέκνον, ἄκουσόν μου τῆς φωνῆς, καὶ ἀναστὰς ἀπόδραθι εἰς τὴν Μεσοποταμίαν πρὸς Λάβαν τὸν ἀδελφόν μου εἰς Χαῤῥάν.
\vs{44}Καὶ οἴκησον μετʼ αὐτοῦ ἡμέρας τινὰς, ἕως τοῦ ἀποστρέψαι τὸν θυμὸν,
\vs{45}καὶ τὴν ὀργὴν τοῦ ἀδελφοῦ σου ἀπὸ σοῦ, καὶ ἐπιλάθηται ἃ πεποίηκας αὐτῷ· καὶ ἀποστείλασα μεταπέμψομαί σε ἐκεῖθεν, μή ποτε ἀποτεκνωθῶ ἀπὸ τῶν δύο ὑμῶν ἐν ἡμέρᾳ μιᾷ.
\vs{46}Εἶπε δὲ Ῥεβέκκα πρὸς Ἰσαὰκ, προσώχθικα τῇ ζωῇ μου διὰ τὰς θυγατέρας τῶν υἱῶν Χέτ· εἰ λήψεται Ἰακὼβ γυναῖκα ἀπὸ τῶν θυγατέρων τῆς γῆς ταύτης, ἵνα τί μοι τὸ ζῇν;

\ch{28}
Προσκαλεσάμενος δὲ Ἰσαὰκ τὸν Ἰακὼβ, εὐλόγησεν αὐτὸν, καὶ ἐνετείλατο αὐτῷ, λέγων, οὐ λήψῃ γυναῖκα ἐκ τῶν θυγατέρων τῶν Χαναναίων.
\vs{2}Ἀναστὰς ἀπόδραθι εἰς τὴν Μεσοποταμίαν, εἰς τὸν οἶκον Βαθουὴλ τοῦ πατρὸς τῆς μητρός σου, καὶ λάβε σεαυτῷ ἐκεῖθεν γυναῖκα ἐκ τῶν θυγατέρων Λάβαν τοῦ ἀδελφοῦ τῆς μητρός σου.
\vs{3}Ὁ δὲ Θεός μου εὐλογήσαι σε, καὶ αὐξήσαι σε, καὶ πληθύναι σε· καὶ ἔσῃ εἰς συναγωγὰς ἐθνῶν.
\vs{4}Καὶ δῴη σοι τὴν εὐλόγιαν Ἁβραὰμ τοῦ πατρός μου, σοὶ καὶ τῷ σπέρματί σου μετὰ σὲ, κληρονομῆσαι τὴν γῆν τῆς παροικήσεώς σου, ἣν ἔδωκεν ὁ Θεὸς τῷ Ἁβραάμ.
\vs{5}Καὶ ἀπέστειλεν Ἰσαὰκ τὸν Ἰακώβ· καὶ ἐπορεύθη εἰς τὴν Μεσοποταμίαν πρὸς Λάβαν τὸν υἱὸν Βαθουὴλ τοῦ Σύρου, ἀδελφὸν Ῥεβέκκας τῆς μητρὸς Ἰακὼβ καὶ Ἡσαῦ.

\vs{6}Ἴδε δὲ Ἡσαῦ ὅτι εὐλόγησεν Ἰσαὰκ τὸν Ἰακὼβ, καὶ ἀπέστειλεν εἰς τὴν Μεσοποταμίαν Συρίας, λαβεῖν ἑαυτῷ γυναῖκα ἐκεῖθεν, ἐν τῷ εὐλογεῖν αὐτόν· καὶ ἐνετείλατο αὐτῷ, λέγων, οὐ λήψῃ γυναῖκα ἐκ τῶν θυγατέρων τῶν Χαναναίων.
\vs{7}Καὶ ἤκουσεν Ἰακὼβ τοῦ πατρὸς καὶ τῆς μητρὸς αὐτοῦ· καὶ ἐπορεύθη εἰς τὴν Μεσοποταμίαν Συρίας.
\vs{8}Ἰδὼν δὲ καὶ Ἡσαῦ ὅτι πονηραί εἰσιν αἱ θυγατέρες Χαναὰν ἐναντίον Ἰσαὰκ τοῦ πατρὸς αὐτοῦ,
\vs{9}ἐπορεύθη Ἡσαῦ πρὸς Ἰσμαήλ· καὶ ἔλαβε τὴν Μαελὲθ, θυγατέρα Ἰσμαὴλ τοῦ υἱοῦ Ἁβραὰμ, ἀδελφὴν Ναβεὼθ, πρὸς ταῖς γυναιξὶν αὐτοῦ γυναῖκα.

\vs{10}Καὶ ἐξῆλθεν Ἰακὼβ ἀπὸ τοῦ φρέατος τοῦ ὅρκου, καὶ ἐπορεύθη εἰς Χαῤῥάν.
\vs{11}Καὶ ἀπήντησε τόπῳ, καὶ ἐκοιμήθη ἐκεῖ, ἔδυ γὰρ ὁ ἥλιος· καὶ ἔλαβεν ἀπὸ τῶν λίθων τοῦ τόπου, καὶ ἔθηκε πρὸς κεφαλῆς αὐτοῦ· καὶ ἐκοιμήθη ἐν τῷ τόπῳ ἐκείνῳ.
\vs{12}Καὶ ἐνυπνιάσθη· καὶ ἰδοὺ κλίμαξ ἐστηριγμένη ἐν τῇ γῇ, ἧς ἡ κεφαλὴ ἀφικνεῖτο εἰς τὸν οὐρανόν· καὶ οἱ ἄγγελοι τοῦ θεοῦ ἀνέβαινον καὶ κατέβαινον ἐπʼ αὐτῇ.
\vs{13}Ὁ δὲ Κύριος ἐπεστήρικτο ἐπʼ αὐτῆς· καὶ εἶπεν, ἐγώ εἰμι ὁ Θεὸς Ἁβραὰμ τοῦ πατρός σου, καὶ ὁ Θεὸς Ἰσαάκ· μὴ φοβοῦ· ἡ γῆ ἐφʼ ἧς σὺ καθεύδεις ἐπʼ αὐτῆς, σοὶ δώσω αὐτὴν, καὶ τῷ σπέρματί σου.
\vs{14}Καὶ ἔσται τὸ σπέρμα σου ὡς ἡ ἄμμος τῆς γῆς, καὶ πλατυνθήσεται ἐπὶ θάλασσαν, καὶ Λίβα, καὶ Βοῤῥὰν, καὶ ἐπὶ ἀνατολάς· καὶ ἐνευλογηθήσονται ἐν σοὶ πᾶσαι αἱ φυλαὶ τῆς γῆς, καὶ ἐν τῷ σπέρματί σου.
\vs{15}Καὶ ἰδοὺ ἐγώ εἰμι μετὰ σοῦ, διαφυλάσσων σε ἐν τῇ ὁδῷ πάσῃ, οὗ ἂν πορευθῇς· καὶ ἀποστρέψω σε εἰς τὴν γῆν ταύτην· ὅτι οὐ μή σε ἐγκαταλίπω, ἕως τοῦ ποιῆσαί με πάντα ὅσα ἐλάλησά σοι.
\vs{16}Καὶ ἐξηγέρθη Ἰακὼβ ἐκ τοῦ ὕπνου αὐτοῦ, καὶ εἶπεν, ὅτι ἔστι Κύριος ἐν τῷ τόπῳ τούτῳ, ἐγὼ δὲ οὐκ ᾔδειν.
\vs{17}Καὶ ἐφοβήθη, καὶ εἶπεν, ὡς φοβερὸς ὁ τόπος οὗτος· οὐκ ἔστι τοῦτο ἀλλʼ ἢ οἶκος Θεοῦ, καὶ αὕτη ἡ πύλη τοῦ οὐρανοῦ.
\vs{18}Καὶ ἀνέστη Ἰακὼβ τὸ πρωῒ, καὶ ἔλαβε τὸν λίθον, ὃν ὑπέθηκεν ἐκεῖ πρὸς κεφαλῆς αὐτοῦ, καὶ ἔστησεν αὐτὸν στήλην, καὶ ἐπέχεεν ἔλαιον ἐπὶ τὸ ἄκρον αὐτῆς.
\vs{19}Καὶ ἐκάλεσε τὸ ὄνομα τοῦ τόπου ἐκείνου, οἶκος Θεοῦ· καὶ Οὐλαμλοὺζ ἦν ὄνομα τῇ πόλει τὸ πρότερον.
\vs{20}Καὶ ηὔξατο Ἰακὼβ εὐχὴν, λέγων, ἐὰν ᾖ Κύριος ὁ Θεὸς μετʼ ἐμοῦ, καὶ διαφυλάξῃ με ἐν τῇ ὁδῷ ταύτῃ, ᾗ ἐγὼ πορεύομαι, καὶ δῷ μοι ἄρτον φαγεῖν, καὶ ἱμάτιον περιβαλέσθαι,
\vs{21}καὶ ἀποστρέψῃ με μετὰ σωτηρίας εἰς τὸν οἶκον τοῦ πατρός μου, καὶ ἔσται Κύριός μοι εἰς Θεόν.
\vs{22}Καὶ ὁ λίθος οὗτος, ὃν ἔστησα στήλην, ἔσται μοι οἶκος Θεοῦ· καὶ πάντων ὧν ἐάν μοι δῷς, δεκάτην ἀποδεκατώσω αὐτά σοι.

\ch{29}
Καὶ ἐξᾴρας Ἰακὼβ τοὺς πόδας ἐπορεύθη εἰς γῆν ἀνατολῶν, πρὸς Λάβαν τὸν υἱὸν Βαθουὴλ τοῦ Σύρου, ἀδελφὸν δὲ Ῥεβέκκας, μητρὸς Ἰακὼβ καὶ Ἡσαῦ.
\vs{2}Καὶ ὁρᾷ, καὶ ἰδοὺ φρέαρ ἐν τῷ πεδίῳ· ἦσαν δὲ ἐκεῖ τρία ποίμνια προβάτων ἀναπαυόμενα ἐπʼ αὐτοῦ· ἐκ γὰρ τοῦ φρέατος ἐκείνου ἐπότιζον τὰ ποίμνια· λίθος δὲ ἦν μέγας ἐπὶ τῷ στόματι τοῦ φρέατος.
\vs{3}Καὶ συνήγοντο ἐκεῖ πάντα τὰ ποίμνια· καὶ ἀπεκύλιον τὸν λίθον ἀπὸ τοῦ στόματος τοῦ φρέατος, καὶ ἐπότιζον τὰ πρόβατα, καὶ ἀπεκαθίστων τὸν λίθον ἐπὶ τὸ στόμα τοῦ φρέατος εἰς τὸν τόπον αὐτοῦ.
\vs{4}Εἶπε δὲ αὐτοῖς Ἰακὼβ, ἀδελφοὶ, πόθεν ἐστὲ ὑμεῖς; οἱ δὲ εἶπαν, ἐκ Χαῤῥὰν ἐσμέν.
\vs{5}Εἶπε δὲ αὐτοῖς, γινώσκετε Λάβαν τὸν υἱὸν Ναχώρ; οἱ δὲ εἶπαν, γινώσκομεν·
\vs{6}Εἶπε δὲ αὐτοῖς, ὑγιαίνει; οἱ δὲ εἶπαν, ὑγιαίνει· καὶ ἰδοὺ Ῥαχὴλ ἡ θυγάτηρ αὐτοῦ ἤρχετο μετὰ τῶν προβάτων.
\vs{7}Καὶ εἶπεν Ἰακὼβ, ἔτι ἐστὶν ἡμέρα πολλὴ· οὔπω ὥρα συναχθῆναι τὰ κτήνη· ποτίσαντες τὰ πρόβατα, ἀπελθόντες βόσκετε.
\vs{8}Οἱ δὲ εἶπαν, οὐ δυνησόμεθα, ἕως τοῦ συναχθῆναι πάντας τοὺς ποιμένας, καὶ ἀποκυλίσουσι τὸν λίθον ἀπὸ τοῦ στόματος τοῦ φρέατος, καὶ ποτιοῦμεν τὰ πρόβατα.
\vs{9}Ἔτι αὐτοῦ λαλοῦντος αὐτοῖς, καὶ ἰδοὺ Ῥαχὴλ ἡ θυγάτηρ Λάβαν ἤρχετο μετὰ τῶν προβάτων τοῦ πατρὸς αὐτῆς· αὐτὴ γὰρ ἔβοσκε τὰ πρόβατα τοῦ πατρὸς αὐτῆς.
\vs{10}Ἐγένετο δὲ ὡς εἶδεν Ἰακὼβ τὴν Ῥαχὴλ τὴν θυγατέρα Λάβαν, τοῦ ἀδελφοῦ τῆς μητρὸς αὐτοῦ, καὶ τὰ πρόβατα Λάβαν τοῦ ἀδελφοῦ τῆς μητρὸς αὐτοῦ, καὶ προσελθὼν Ἰακὼβ ἀπεκύλισε τὸν λίθον ἀπὸ τοῦ στόματος τοῦ φρέατος, καὶ ἐπότιζε τὰ πρόβατα Λάβαν τοῦ ἀδελφοῦ τῆς μητρὸς αὐτοῦ.
\vs{11}Καὶ ἐφίλησεν Ἰακὼβ τὴν Ῥαχὴλ, καὶ βοήσας τῇ φωνῇ αὐτοῦ ἔκλαυσε.
\vs{12}Καὶ ἀπήγγειλε τῇ Ῥαχὴλ, ὅτι ἀδελφὸς τοῦ πατρὸς αὐτῆς ἐστι, καὶ ὅτι υἱὸς Ῥεβέκκας ἐστί· καὶ δραμοῦσα ἀπήγγειλε τῷ πατρὶ αὐτῆς κατὰ τὰ ῥήματα ταῦτα.
\vs{13}Ἐγένετο δὲ ὡς ἤκουσε Λάβαν τὸ ὄνομα Ἰακὼβ τοῦ υἱοῦ τῆς ἀδελφῆς αὐτοῦ, ἔδραμεν εἰς συνάντησιν αὐτῷ, καὶ περιλαβὼν αὐτὸν ἐφίλησε, καὶ εἰσήγαγεν αὐτὸν εἰς τὸν οἶκον αὐτοῦ· καὶ διηγήσατο τῷ Λάβαν πάντας τοὺς λόγους τούτους.
\vs{14}Καὶ εἶπεν αὐτῷ Λάβαν, ἐκ τῶν ὀστῶν μου καὶ ἐκ τῆς σαρκός μου εἶ σύ· καὶ ἦν μετʼ αὐτοῦ μῆνα ἡμερῶν.

\vs{15}Εἶπε δὲ Λάβαν τῷ Ἰακὼβ, ὅτι γὰρ ἀδελφός μου εἶ, οὐ δουλεύσεις μοι δωρεάν· ἀπάγγειλόν μοι τίς ὁ μισθός σου ἐστί;
\vs{16}Τῷ δὲ Λάβαν ἦσαν δύο θυγατέρες· ὄνομα τῇ μείζονι, Λεία, καὶ ὄνομα τῇ νεωτέρᾳ, Ῥαχήλ.
\vs{17}Οἱ δὲ ὀφθάλμοὶ Λείας, ἀσθενεῖς· Ῥαχῆλ δὲ ἦν καλὴ τῷ εἴδει, καὶ ὡραία τῇ ὄψει σφάδρα.
\vs{18}Ἠγάπησε δὲ Ἰακὼβ τὴν Ῥαχήλ· καὶ εἶπε, δουλεύσω σοι ἑπτὰ ἔτη περὶ τῆς Ῥαχὴλ τῆς θυγατρός σου τῆς νεωτέρας.
\vs{19}Εἶπε δὲ αὐτῷ Λάβαν, βέλτιον δοῦναί με αὐτήν σοι, ἢ δοῦναί με αὐτὴν ἀνδρὶ ἑτέρῳ· οἴκησον μετʼ ἐμοῦ.
\vs{20}Καὶ ἐδούλευσεν Ἰακὼβ περὶ Ῥαχὴλ ἑπτὰ ἔτη· καὶ ἤσαν ἐναντίον αὐτοῦ ὡς ἡμέραι ὀλίγαι, παρὰ τὸ ἀγαπᾷν αὐτὸν αὐτήν.
\vs{21}Εἶπε δὲ Ἰακὼβ τῷ Λάβαν, δός μοι τὴν γυναῖκά μου, πεπλήρωνται γὰρ αἱ ἡμέραι ὅπως εἰσέλθω πρὸς αὐτήν.
\vs{22}Συνήγαγε δὲ Λάβαν πάντας τοὺς ἄνδρας τοῦ τόπου, καὶ ἐποίησε γάμον.
\vs{23}Καὶ ἐγένετο ἑσπέρα, καὶ λαβὼν Λείαν τὴν θυγατέρα αὐτοῦ, εἰσήγαγεν πρὸς Ἰακὼβ, καὶ εἰσῆλθε πρὸς αὐτὴν Ἰακώβ.
\vs{24}Ἔδωκε δὲ Λάβαν Λείᾳ τῇ θυγατρὶ αὐτοῦ Ζελφὰν τὴν παιδίσκην αὐτοῦ, αὐτῇ παιδίσκην.
\vs{25}Ἐγένετο δὲ πρωῒ, καὶ ἰδοὺ ἦν Λεία· εἶπε δὲ Ἰακὼβ τῷ Λάβαν, τί τοῦτο ἐποίησάς μοι; οὐ περὶ Ῥαχὴλ ἐδούλευσα παρὰ σοι; καὶ ἱνατί παρελογίσω με;
\vs{26}Ἀπεκρίθη δὲ Λάβαν, οὐκ ἔστιν οὕτως ἐν τῷ τόπῳ ἡμῶν, δοῦναι τὴν νεωτέραν πρινὴ τὴν πρεσβυτέραν.
\vs{27}Συντέλεσον οὖν τὰ ἕβδομα ταύτης, καὶ δώσω σοι καὶ ταύτην ἀντὶ τῆς ἐργασίας, ἧς ἐργᾷ παρʼ ἐμοὶ ἔτι ἑπτὰ ἔτη ἕτερα.
\vs{28}Ἐποίησε δὲ Ἰακὼβ οὕτως, καὶ ἀνεπλήρωσε τὰ ἕβδομα ταύτης· καὶ ἔδωκεν αὐτῷ Λάβαν Ῥαχὴλ τὴν θυγατέρα αὐτοῦ αὐτῷ γυναῖκα.
\vs{29}Ἔδωκε δὲ Λάβαν τῇ θυγατρὶ αὐτοῦ Βαλλὰν τὴν παιδίσκην αὐτοῦ, αὐτῇ παιδίσκην.
\vs{30}Καὶ εἰσῆλθε πρὸς Ῥαχήλ· ἠγάπησε δὲ Ῥαχὴλ μᾶλλον ἢ Λείαν· καὶ ἐδούλευσεν αὐτῷ ἑπτὰ ἔτη ἕτερα.

\vs{31}Ἰδὼν δὲ Κύριος ὁ Θεὸς ὅτι ἐμισεῖτο Λεία, ἤνοιξε τὴν μήτραν αὐτῆς· Ῥαχὴλ δὲ ἦν στεῖρα.
\vs{32}Καὶ συνέλαβε Λεία, καὶ ἔτεκεν υἱὸν τῷ Ἰακώβ· ἐκάλεσε δὲ τὸ ὄνομα αὐτοῦ Ῥουβὴν, λέγουσα, διότι εἶδέ μου Κύριος τὴν ταπείνωσιν, καὶ ἔδωκέ μοι υἱόν· νῦν οὖν ἀγαπήσει με ὁ ἀνήρ μου.
\vs{33}Καὶ συνέλαβε πάλιν, καὶ ἔτεκεν υἱὸν δεύτερον τῷ Ἰακὼβ, καὶ εἶπεν, ὅτι ἤκουσε Κύριος ὅτι μισοῦμαι, καὶ προσέδωκέ μοι καὶ τοῦτον· καὶ ἐκάλεσε τὸ ὄνομα αὐτοῦ, Συμεών.
\vs{34}Καὶ συνέλαβεν ἔτι, καὶ ἔτεκεν υἱὸν, καὶ εἶπεν, ἐν τῷ νῦν καιρῷ πρὸς ἐμοῦ ἔσται ὁ ἀνήρ μου, τέτοκα γὰρ αὐτῷ τρεῖς υἱούς· διὰ τοῦτο ἐκάλεσε τὸ ὄνομα αὐτοῦ, Λευεί.
\vs{35}Καὶ συλλαβοῦσα ἔτι ἔτεκεν υἱὸν, καὶ εἶπε, νῦν ἔτι τοῦτο ἐξομολογήσομαι Κυρίῳ· διὰ τοῦτο ἐκάλεσε τὸ ὄνομα αὐτοῦ, Ἰούδαν· καὶ ἔστη τοῦ τίκτειν.

\ch{30}
Ἰδοῦσα δὲ Ῥαχὴλ, ὅτι οὐ τέτοκε τῷ Ἱακώβ· καὶ ἐζήλωσε Ῥαχὴλ τὴν ἀδελφὴν αὐτῆς· καὶ εἶπε τῷ Ἰακὼβ, δός μοι τέκνα· εἰ δὲ μὴ, τελευτήσω ἐγώ.
\vs{2}Θυμωθεὶς δὲ Ἰακὼβ τῇ Ῥαχὴλ εἶπεν αὐτῇ, μὴ ἀντὶ Θεοῦ ἐγώ εἰμι, ὃς ἐστέρησέ σε καρπὸν κοιλίας;
\vs{3}Εἶπε δὲ Ῥαχὴλ τῷ Ἰακὼβ, ἰδοὺ ἡ παιδίσκη μου Βαλλά· εἴσελθε πρὸς αὐτήν· καὶ τέξεται ἐπὶ τῶν γονάτων μου, καὶ τεκνοποιήσομαι κᾀγὼ ἐξ αὐτῆς.
\vs{4}Καὶ ἔδωκεν αὐτῷ Βαλλὰν τὴν παιδίσκην αὐτῆς, αὐτῷ γυναῖκα· καὶ εἰσῆλθε πρὸς αὐτὴν Ἰακώβ.
\vs{5}Καὶ συνέλαβε Βαλλὰ ἡ παιδίσκη Ῥαχὴλ, καὶ ἔτεκε τῷ Ἰακὼβ υἱόν.
\vs{6}Καὶ εἶπε Ῥαχὴλ, ἔκρινέ μοι ὁ Θεὸς, καὶ ἐπήκουσε τῆς φωνῆς μου, καὶ ἔδωκε μοι υἱόν· διὰ τοῦτο ἐκάλεσε τὸ ὄνομα αὐτοῦ, Δάν.
\vs{7}Καὶ συνέλαβεν ἔτι Βαλλὰ ἡ παιδίσκη Ῥαχὴλ, καὶ ἔτεκεν υἱὸν δεύτερον τῷ Ἰακώβ.
\vs{8}Καὶ εἶπε Ῥαχὴλ, συναντελάβετό μου ὁ Θεὸς, καὶ συνανεστράφην τῇ ἀδελφῇ μου, καὶ ἠδυνάσθην· καὶ ἐκάλεσε τὸ ὄνομα αὐτοῦ, Νεφθαλεί.
\vs{9}Εἶδε δὲ Λεία ὅτι ἔστη τοῦ τίκτειν· καὶ ἔλαβε Ζελφὰν τὴν παιδίσκην αὐτῆς, καὶ ἔδωκεν αὐτὴν τῷ Ἰακὼβ γυναῖκα· καὶ εἰσῆλθε πρὸς αὐτήν.
\vs{10}Καὶ συνέλαβε Ζελφὰ ἡ παιδίσκη Λείας, καὶ ἔτεκε τῷ Ἰακὼβ υἱόν.
\vs{11}Καὶ εἶπε Λεία, ἐν τύχῃ· καὶ ἐπωνόμασε τὸ ὄνομα αὐτοῦ, Γάδ.
\vs{12}Καὶ συνέλαβεν ἔτι Ζελφὰ ἡ παιδίσκη Λείας, καὶ ἔτεκε τῷ Ἰακὼβ υἱὸν δεύτερον.
\vs{13}Καὶ εἶπε Λεία, μακαρία ἐγὼ, ὅτι μακαριοῦσί με αἱ γυναῖκες· καὶ ἐκάλεσε τὸ ὄνομα αὐτοῦ, Ἀσήρ.
\vs{14}Ἐπορεύθη δὲ Ῥουβὴν ἐν ἡμέρᾳ θερισμοῦ πυρῶν, καὶ εὗρε μῆλα μανδραγορῶν ἐν τῷ ἀγρῷ, καὶ ἤνεγκεν αὐτὰ πρὸς Λείαν τὴν μητέρα αὐτοῦ· εἶπε δὲ Ῥαχὴλ τῇ Λείᾳ τῇ ἀδελφῇ αὐτῆς, δός μοι τῶν μανδραγορῶν τοῦ υἱοῦ σου.
\vs{15}Εἶπε δὲ Λεία, οὐχ ἱκανόν σοι ὅτι ἔλαβες τὸν ἄνδρα μου; μὴ καὶ τοὺς μανδραγόρας τοῦ υἱοῦ μου λήψῃ; εἶπε δὲ Ῥαχὴλ, οὐχ οὕτως· κοιμηθήτω μετὰ σοῦ τὴν νύκτα ταύτην ἀντὶ τῶν μανδραγορῶν τοῦ υἱοῦ σου.
\vs{16}Εἰσῆλθεν δὲ Ἰακὼβ ἐξ ἀγροῦ ἑσπέρας· καὶ ἐξῆλθε Λεία εἰς συνάντησιν αὐτῷ, καὶ εἶπε, πρὸς ἐμὲ εἰσελεύσῃ σήμερον· μεμίσθωμαι γάρ σε ἀντὶ τῶν μανδραγορῶν τοῦ υἱοῦ μου· καὶ ἐκοιμήθη μετʼ αὐτῆς τὴν νύκτα ἐκείνην.
\vs{17}Καὶ ἐπήκουσεν ὁ Θεὸς Λείας· καὶ συλλαβοῦσα ἔτεκε τῷ Ἰακὼβ υἱὸν πέμπτον.
\vs{18}Καὶ εἶπε Λεία, δέδωκέ μοι ὁ Θεὸς τὸν μισθόν μου, ἀνθʼ οὗ ἔδωκα τὴν παιδίσκην μου τῷ ἀνδρί μου· καὶ ἐκάλεσε τὸ ὄνομα αὐτοῦ, Ἰσσάχαρ, ὅ ἐστι μισθός.
\vs{19}Καὶ συνέλαβεν ἔτι Λεία, καὶ ἔτεκεν υἱὸν ἕκτον τῷ Ἰακώβ.
\vs{20}Καὶ εἶπε Λεία, δεδώρηται ὁ Θεός μοι δῶρον καλὸν ἐν τῷ νῦν καιρῷ· αἱρετιεῖ με ὁ ἀνήρ μου, τέτοκα γὰρ αὐτῷ υἱοὺς ἕξ· καὶ ἐκάλεσε τὸ ὄνομα αὐτοῦ, Ζαβουλών.
\vs{21}Καὶ μετὰ τοῦτο ἔτεκε θυγατέρα, καὶ ἐκάλεσε τὸ ὄνομα αὐτῆς, Δεῖνα.
\vs{22}Ἐμνήσθη δὲ ὁ Θεὸς τῆς Ῥαχὴλ, καὶ ἔπήκουσεν αὐτῆς ὁ Θεός· καὶ ἀνέῳξεν αὐτῆς τὴν μήτραν.
\vs{23}Καὶ συλλαβοῦσα ἔτεκε τῷ Ἰακὼβ υἱόν· εἶπε δὲ Ῥαχὴλ, ἀφεῖλεν ὁ Θεός μου τὸ ὄνειδος.
\vs{24}Καὶ ἐκάλεσε τὸ ὄνομα αὐτοῦ Ἰωσὴφ, λέγουσα, προσθέτω ὁ Θεός μοι υἱὸν ἕτερον.

\vs{25}Ἐγένετο δὲ ὡς ἔτεκε Ῥαχὴλ τὸν Ἰωσὴφ, εἶπεν Ἰακὼβ τῷ Λάβαν, ἀπόστειλόν με, ἵνα ἀπέλθω εἰς τὸν τόπον μου, καὶ εἰς τὴν γῆν μου.
\vs{26}Ἀπόδος τὰς γυναῖκας μου, καὶ τὰ παιδία μου, περὶ ὧν δεδούλευκά σοι, ἵνα ἀπέλθω· σὺ γὰρ γινώσκεις τὴν δουλείαν, ἣν δεδούλευκά σοι.
\vs{27}Εἶπε δὲ αὐτῷ Λάβαν, εἰ εὗρον χάριν ἐναντίον σου, οἰωνισάμην ἄν· εὐλόγησε γάρ με ὁ Θεὸς ἐπὶ τῇ σῇ εἰσόδῳ.
\vs{28}Διάστειλον τὸν μισθόν σου πρός με, καὶ δώσω.
\vs{29}Εἶπε δὲ Ἰακὼβ, σὺ γινώσκεις ἃ δεδούλευκά σοι, καὶ ὅσα ἦν κτήνη σου μετʼ ἐμοῦ.
\vs{30}Μικρὰ γὰρ ἦν ὅσα σοι ἐναντίον ἐμοῦ, καὶ ηὐξήθη εἰς πλῆθος· καὶ εὐλόγησέ σε Κύριος ὁ Θεὸς ἐπὶ τῷ ποδί μου· νῦν οὖν πότε ποιήσω κᾀγὼ ἐμαυτῷ οἶκον;
\vs{31}Καὶ εἶπεν αὐτῷ Λάβαν, τί σοι δώσω; Εἶπε δὲ αὐτῷ Ἰακὼβ, οὐ δώσεις μοι οὐθὲν, ἐὰν ποιήσῃς μοι τὸ ῥῆμα τοῦτο, πάλιν ποιμανῶ τὰ πρόβατά σου, καὶ φυλάξω.
\vs{32}Παρελθέτω πάντα τὰ πρόβατά σου σήμερον, καὶ διαχώρισον ἐκεῖθεν πᾶν πρόβατον φαιὸν ἐν τοῖς ἄρνασι, καὶ πᾶν διάλευκον καὶ ῥαντὸν ἐν ταῖς αἰξὶν, ἔσται μοι μισθός.
\vs{33}Καὶ ἐπακούσεταί μοι ἡ δικαιοσύνη μου ἐν τῇ ἡμέρᾳ τῇ ἐπαύριον, ὅτι ἐστὶν ὁ μισθός μου ἐνώπιόν σου· πᾶν ὃ ἐὰν μὴ ᾖ ῥαντὸν καὶ διάλευκον ἐν ταῖς αἰξὶ, καὶ φαιὸν ἐν τοῖς ἄρνασι, κεκλεμμένον ἔσται παρʼ ἐμοί.
\vs{34}Εἶπε δὲ αὐτῷ Λάβαν, ἔστω κατὰ τὸ ῥῆμά σου.
\vs{35}Καὶ διέστειλεν ἐν τῇ ἡμέρᾳ ἐκείνῃ τοὺς τράγους τοὺς ῥαντοὺς καὶ τοὺς διαλεύκους, καὶ πάσας τὰς αἶγας τὰς ῥαντὰς καὶ τὰς διαλεύκους, καὶ πᾶν ὃ ἦν φαιὸν ἐν τοῖς ἄρνασι, καὶ πᾶν ὃ ἦν λευκὸν ἐν αὐτοῖς, καὶ ἔδωκε διὰ χειρὸς τῶν υἱῶν αὐτοῦ.
\vs{36}Καὶ ἀπέστησεν ὁδὸν τριῶν ἡμερῶν, καὶ ἀνὰ μέσον αὐτῶν καὶ ἀνὰ μέσον Ἰακώβ· Ἰακὼβ δὲ ἐποίμαινε τὰ πρόβατα Λάβαν τὰ ὑπολειφθέντα.
\vs{37}Ἔλαβε δὲ ἑαυτῷ Ἰακὼβ ῥάβδον στυρακίνην χλωρὰν καὶ καρυΐνην καὶ πλατάνου· καὶ ἐλέπισεν αὐτὰς Ἰακὼβ λεπίσματα λευκά· καὶ περισύρων τὸ χλωρὸν, ἐφαίνετο ἐπὶ ταῖς ῥάβδοις τὸ λευκὸν, ὃ ἐλέπισε, ποικίλον.
\vs{38}Καὶ παρέθηκε τὰς ῥάβδους, ἃς ἐλέπισεν, ἐν τοῖς ληνοῖς τῶν ποτιστηρίων τοῦ ὕδατος, ἵνα ὡς ἂν ἔλθωσι τὰ πρόβατα πιεῖν, ἐνώπιον τῶν ῥάβδων ἐλθόντων αὐτῶν εἰς τὸ πιεῖν, ἐγκισσήσωσι τὰ πρόβατα εἰς τὰς ῥάβδους.
\vs{39}Καὶ ἐνεκίσσων τὰ πρόβατα εἰς τὰς ῥάβδους· καὶ ἔτικτον τὰ πρόβατα διάλευκα καὶ ποικίλα καὶ σποδοειδῆ ῥαντά.
\vs{40}Τοὺς δὲ ἀμνοὺς διέστειλεν Ἰακὼβ, καὶ ἔστησεν ἐναντίον τῶν προβάτων κριὸν διάλευκον, καὶ πᾶν ποικίλον ἐν τοῖς ἀμνοῖς· καὶ διεχώρισεν ἑαυτῷ ποίμνια καθʼ ἑαυτὸν, καὶ οὐκ ἔμιξεν αὐτὰ εἰς τὰ πρόβατα Λάβαν.
\vs{41}Ἐγένετο δὲ ἐν τῷ καιρῷ ᾧ ἐνεκίσσων τὰ πρόβατα ἐν γαστρὶ λαμβάνοντα, ἔθηκεν Ἰακὼβ τὰς ῥάβδους ἐναντίον τῶν προβάτων ἐν τοῖς ληνοῖς, τοῦ ἐγκισσῆσαι αὐτὰ κατὰ τὰς ῥάβδους.
\vs{42}Ἡνίκα δʼ ἂν ἔτεκε τὰ πρόβατα, οὐκ ἐτίθει· ἐγένετο δὲ τὰ μὲν ἄσημα τοῦ Λάβαν, τὰ δὲ ἐπίσημα τοῦ Ἰακώβ.
\vs{43}Καὶ ἐπλούτησεν ὁ ἄνθρωπος σφόδρα σφόδρα· καὶ ἐγένετο αὐτῷ κτήνη πολλὰ, καὶ βόες, καὶ παῖδες, καὶ παιδίσκαι, καὶ κάμηλοι, καὶ ὄνοι.

\ch{31}
Ἤκουσε δὲ Ἰακὼβ τὰ ῥήματα τῶν υἱῶν Λάβαν, λεγόντων, εἴληφεν Ἰακὼβ πάντα τὰ τοῦ πατρὸς ἡμῶν, καὶ ἐκ τῶν τοῦ πατρὸς ἡμῶν πεποίηκε πᾶσαν τὴν δόξαν ταύτην.
\vs{2}Καὶ εἶδεν Ἰακὼβ τὸ πρόσωπον τοῦ Λάβαν, καὶ ἰδοὺ οὐκ ἦν πρὸς αὐτὸν ὡσεὶ χθὲς καὶ τρίτην ἡμέραν.
\vs{3}Εἶπε δὲ Κύριος πρὸς Ἰακὼβ, ἀποστρέφου εἰς τὴν γῆν τοῦ πατρός σου, καὶ εἰς τὴν γενεάν σου, καὶ ἔσομαι μετὰ σοῦ.
\vs{4}Ἀποστείλας δὲ Ἰακὼβ ἐκάλεσε Λείαν καὶ Ῥαχὴλ εἰς τὸ πεδίον, οὗ ἦν τὰ ποίμνια.
\vs{5}Καὶ εἶπεν αὐταῖς, ὁρῶ ἐγὼ τὸ πρόσωπον τοῦ πατρὸς ὑμῶν, ὅτι οὐκ ἔστι πρὸς ἐμοῦ, ὡς ἐχθὲς καὶ τρίτην ἡμέραν· ὁ δὲ Θεὸς τοῦ πατρός μου ἦν μετʼ ἐμοῦ.
\vs{6}Καὶ αὐταὶ δὲ οἴδατε, ὅτι ἐν πάσῃ τῇ ἰσχύϊ μου δεδούλευκα τῷ πατρὶ ὑμῶν.
\vs{7}Ὁ δὲ πατὴρ ὑμῶν παρεκρούσατό με, καὶ ἤλλαξε τὸν μισθόν μου τῶν δέκα ἀμνῶν· καὶ οὐκ ἔδωκεν αὐτῷ ὁ Θεὸς κακοποιῆσαί με.
\vs{8}Ἐὰν οὕτως εἴπῃ, τὰ ποικίλα ἔσται σου μισθὸς, καὶ τέξεται πάντα τὰ πρόβατα ποικίλα· ἐὰν δὲ εἴπῃ, τὰ λευκὰ ἔσται σου μισθὸς, καὶ τέξεται πάντα τὰ πρόβατα λευκά.
\vs{9}Καὶ ἀφείλετο ὁ Θεὸς πάντα τὰ κτήνη τοῦ πατρὸς ὑμῶν, καὶ ἔδωκέ μοι αὐτά.
\vs{10}Καὶ ἐγένετο ἡνίκα ἐνεκίσσων τὰ πρόβατα ἐν γαστρὶ λαμβάνοντα, καὶ εἶδον τοῖς ὀφθαλμοῖς μου ἐν τῷ ὕπνῳ· καὶ ἰδοὺ οἱ τράγοι καὶ οἱ κριοὶ ἀναβαίνοντες ἐπὶ τὰ πρόβατα καὶ τὰς αἶγας, διάλευκοι καὶ ποικίλοι καὶ σποδοειδεῖς ῥαντοί.
\vs{11}Καὶ εἶπέ μοι ὁ Ἄγγελος τοῦ Θεοῦ καθʼ ὕπνον, Ἰακώβ· ἐγὼ δὲ εἶπα, τί ἐστι;
\vs{12}Καὶ εἶπεν, ἀνάβλεψον τοῖς ὀφθαλμοῖς σου, καὶ ἴδε τοὺς τράγους καὶ τοὺς κριοὺς ἀναβαίνοντας ἐπὶ τὰ πρόβατα καὶ τὰς αἶγας διαλεύκους καὶ ποικίλους καὶ σποδοειδεῖς ῥαντούς· ἑώρακα γὰρ ὅσα σοι Λάβαν ποιεῖ.
\vs{13}Ἐγώ εἰμι ὁ Θεὸς ὁ ὀφθείς σοι ἐν τόπῳ Θεοῦ, οὗ ἤλειψάς μοι ἐκεῖ στήλην, καὶ ηὔξω μοι ἐκεῖ εὐχήν· νῦν οὖν ἀνάστηθι, καὶ ἔξελθε ἐκ τῆς γῆς ταύτης, καὶ ἄπελθε εἰς τὴν γῆν τῆς γενέσεώς σου, καὶ ἔσομαι μετὰ σοῦ.
\vs{14}Καὶ ἀποκριθεῖσαι Ῥαχὴλ καὶ Λεία εἶπαν αὐτῷ, μὴ ἔστιν ἡμῖν ἔτι μερὶς ἢ κληρονομία ἐν τῷ οἴκῳ τοῦ πατρὸς ἡμῶν;
\vs{15}Οὐχ ὡς αἱ ἀλλότριαι λελογίσμεθα αὐτῷ; πέπρακε γὰρ ἡμᾶς, καὶ καταβρώσει κατέφαγε τὸ ἀργύριον ἡμῶν.
\vs{16}Πάντα τὸν πλοῦτον καὶ τὴν δόξαν, ἣν ἀφείλετο ὁ Θεὸς τοῦ πατρὸς ἡμῶν, ἡμῖν ἔσται καὶ τοῖς τέκνοις ἡμῶν· νῦν οὖν ὅσα σοι εἴρηκεν ὁ Θεὸς, ποίει.
\vs{17}Ἀναστὰς δὲ Ἰακὼβ ἔλαβε τὰς γυναῖκας αὐτοῦ καὶ τὰ παιδία αὐτοῦ ἐπὶ τὰς καμήλους·
\vs{18}Καὶ ἀπήγαγε πάντα τὰ ὑπάρχοντα αὐτῷ, καὶ πᾶσαν τὴν ἀποσκευὴν αὐτοῦ, ἣν περιεποιήσατο ἐν τῇ Μεσοποταμίᾳ, καὶ πάντα τὰ αὐτοῦ, ἀπελθεῖν πρὸς Ἰσαὰκ τὸν πατέρα αὐτοῦ εἰς γῆν Χαναάν.
\vs{19}Λάβαν δὲ ᾤχετο κεῖραι τὰ πρόβατα αὐτοῦ· ἔκλεψε δὲ Ῥαχὴλ τὰ εἴδωλα τοῦ πατρὸς αὐτῆς.
\vs{20}Ἔκρυψε δὲ Ἰακὼβ Λάβαν τὸν Σύρον, τοῦ μὴ ἀναγγεῖλαι αὐτῷ, ὅτι ἀποδιδράσκει.
\vs{21}Καὶ ἀπέδρα αὐτὸς, καὶ τὰ αὐτοῦ πάντα, καὶ διέβη τὸν ποταμὸν, καὶ ὥρμησεν εἰς τὸ ὄρος Γαλαάδ.
\vs{22}Ἀνηγγέλη δὲ Λάβαν τῷ Σύρῳ τῇ ἡμέρᾳ τῇ τρίτῃ, ὅτι ἀπέδρα Ἰακώβ.
\vs{23}Καὶ παραλαβὼν τοὺς ἀδελφοὺς αὐτοῦ μεθʼ ἑαντοῦ, ἐδίωξεν ὀπίσω αὐτοῦ ὁδὸν ἡμερῶν ἑπτά· καὶ κατέλαβεν αὐτὸν ἐν τῷ ὄρει Γαλαάδ.
\vs{24}Ἦλθε δὲ ὁ Θεὸς πρὸς Λάβαν τὸν Σύρον καθʼ ὕπνον τὴν νύκτα, καὶ εἶπεν αὐτῷ, Φύλαξαι σεαυτὸν μή ποτε λαλήσῃς μετὰ Ἰακὼβ πονηρά.
\vs{25}Καὶ κατέλαβε Λάβαν τὸν Ἰακώβ· Ἰακὼβ δὲ ἔπηξεν τὴν σκηνὴν αὐτοῦ ἐν τῷ ὄρει· Λάβαν δὲ ἔστησε τοὺς ἀδελφοὺς αὐτοῦ ἐν τῷ ὄρει Γαλαάδ.
\vs{26}Εἶπε δὲ Λάβαν τῷ Ἰακὼβ, τί ἐποίησας; ἱνατί κρυφῇ ἀπέδρας, καὶ ἐκλοποφόρησάς με, καὶ ἀπήγαγες τὰς θυγατέρας μου, ὡς αἰχμαλώτιδας μαχαίρᾳ;
\vs{27}Καὶ εἰ ἀνήγγειλάς μοι, ἐξαπέστειλα ἄν σε μετʼ εὐφροσύνης, καὶ μετὰ μουσικῶν, καὶ τυμπάνων, καὶ κιθάρας.
\vs{28}Καὶ οὐκ ἠξιώθην καταφιλῆσαι τὰ παιδία μου, καὶ τὰς θυγατέρας μου· νῦν δὲ ἀφρόνως ἔπραξας.
\vs{29}Καὶ νῦν ἰσχύει ἡ χείρ μου κακοποιῆσαί σε· ὁ δὲ Θεὸς τοῦ πατρός σου χθὲς εἶπε πρός με, λέγων, Φύλαξαι σεαυτὸν μή ποτε λαλήσῃς μετὰ Ἰακὼβ πονηρά.
\vs{30}Νῦν οὖν πεπόρευσαι· ἐπιθυμίᾳ γὰρ ἐπεθύμησας ἀπελθεῖν εἰς τὸν οἶκον τοῦ πατρός σου· ἱνατί ἔκλεψας τοὺς θεούς μου;
\vs{31}Ἀποκριθεὶς δὲ Ἰακὼβ εἶπε τῷ Λάβαν, ὅτι ἐφοβήθην· εἶπα γὰρ, μή ποτε ἀφέλῃ τὰς θυγατέρας σου ἀπʼ ἐμοῦ, καὶ πάντα τὰ ἐμά.
\vs{32}Καὶ εἶπεν Ἰακὼβ, παρʼ ᾧ ἂν εὕρῃς τοὺς θεούς σου, οὐ ζήσεται ἐναντίον τῶν ἀδελφῶν ἡμῶν· ἐπίγνωθι τί ἐστι παρʼ ἐμοὶ τῶν σῶν, καὶ λάβε· καὶ οὐκ ἐπέγνω παρʼ αὐτῷ οὐθέν· οὐκ ᾔδει δὲ Ἰακὼβ, ὅτι Ῥαχὴλ ἡ γυνὴ αὐτοῦ ἔκλεψεν αὐτούς.
\vs{33}Εἰσελθὼν δὲ Λάβαν ἠρεύνησεν εἰς τὸν οἶκον Λείας, καὶ οὐχ εὗρεν· καὶ ἐξῆλθεν ἐκ τοῦ οἴκου Λείας, καὶ ἠρεύνησε τὸν οἶκον Ἰακὼβ, καὶ ἐν τῷ οἴκῳ τῶν δύο παιδισκῶν, καὶ οὐχ εὗρεν· εἰσῆλθε δὲ καὶ εἰς τὸν οἶκον Ῥαχήλ.
\vs{34}Ῥαχὴλ δὲ ἔλαβε τὰ εἴδωλα, καὶ ἐνέβαλεν αὐτὰ εἰς τὰ σάγματα τῆς καμήλου, καὶ ἐπεκάθισεν αὐτοῖς.
\vs{35}Καὶ εἶπε τῷ πατρὶ αὐτῆς, μὴ βαρέως φέρε, κύριε· οὐ δυνάμαι ἀναστῆναι ἐνώπιόν σου, ὅτι τὰ κατʼ ἐθισμὸν τῶν γυναικῶν μοι ἐστίν· ἠρεύνησε Λάβαν ἐν ὅλῳ τῷ οἴκῳ, καὶ οὐχ εὗρε τὰ εἴδωλα.
\vs{36}Ὠργίσθη δὲ Ἰακὼβ, καὶ ἐμαχέσατο τῷ Λάβαν· ἀποκριθεὶς δὲ Ἰακὼβ εἶπε τῷ Λάβαν, τί τὸ ἀδίκημά μου; καὶ τί τὸ ἁμάρτημά μου, ὅτι κατεδίωξας ὀπίσω μου,
\vs{37}καὶ ὅτι ἠρεύνησας πάντα τὰ σκεύη τοῦ οἴκου μου; τί εὗρες ἀπὸ πάντων τῶν σκευῶν τοῦ οἴκου σου; θὲς ὧδε ἐνώπιον τῶν ἀδελφῶν σου καὶ τῶν ἀδελφῶν μου, καὶ ἐλεγξάτωσαν ἀνὰ μέσον τῶν δύο ἡμῶν.
\vs{38}Ταῦτά μοι εἴκοσι ἔτη ἐγώ εἰμι μετὰ σοῦ· τὰ πρόβατά σου καὶ αἱ αἶγές σου οὐκ ἠτεκνώθησαν· κριοὺς τῶν προβάτων σου οὐ κατέφαγον.
\vs{39}Θηριάλωτον οὐκ ἐνήνοχά σοι· ἐγὼ ἀπετίννυον παρʼ ἐμαυτοῦ κλέμματα ἡμέρας, καὶ κλέμματα νυκτός.
\vs{40}Ἐγενόμην τῆς ἡμέρας συγκαιόμενος τῷ καύματι, καὶ τῷ παγετῷ τῆς νυκτός· καὶ ἀφίστατο ὁ ὕπνος μου ἀπὸ τῶν ὀφθαλμῶν μου.
\vs{41}Ταῦτά μοι εἴκοσι ἔτη ἐγώ εἰμι ἐν τῇ οἰκίᾳ σου· ἐδούλευσά σοι δεκατέσσαρα ἔτη ἀντὶ τῶν δύο θυγατέρων σου, καὶ ἓξ ἔτη ἐν τοῖς προβάτοις σου, καὶ παρελογίσω τὸν μισθόν μου δέκα ἀμνάσιν.
\vs{42}Εἰ μὴ ὁ Θεὸς τοῦ πατρός μου Ἁβραὰμ, καὶ ὁ φόβος Ἰσαὰκ, ἦν μοι, νῦν ἂν κενόν με ἐξαπέστειλας· τὴν ταπείνωσίν μου, καὶ τὸν κόπον τῶν χειρῶν μου, εἶδεν ὁ Θεός· καὶ ἤλεγξέ σε χθές.

\vs{43}Ἀποκριθεὶς δὲ Λάβαν εἶπε τῷ Ἰακὼβ, αἱ θυγατέρες, θυγατέρες μου, καὶ υἱοὶ, υἱοί μου, καὶ τὰ κτήνη, κτήνη μου· καὶ πάντα ὅσα σὺ ὁρᾷς, ἐμά ἐστι, καὶ τῶν θυγατέρων μου· τί ποιήσω ταύταις σήμερον ἢ τοῖς τέκνοις αὐτῶν, οἷς ἔτεκον;
\vs{44}Νῦν οὖν δεῦρο διαθῶμαι διαθήκην ἐγώ τε καὶ σύ· καὶ ἔσται εἰς μαρτύριον ἀνὰ μέσον ἐμοῦ καὶ σοῦ· εἶπε δὲ αὐτῷ, ἰδοὺ οὐθεὶς μεθʼ ἡμῶν ἐστιν· ἴδε ὁ Θεὸς μάρτυς ἀνὰ μέσον ἐμοῦ καὶ σοῦ.
\vs{45}Λαβὼν δὲ Ἰακὼβ λίθον, ἔστησεν αὐτὸν στήλην.
\vs{46}Εἶπε δὲ Ἰακὼβ τοῖς ἀδελφοῖς αὐτοῦ, συλλέγετε λίθους· καὶ συνέλεξαν λίθους, καὶ ἐποίησαν βουνόν· καὶ ἔφαγον ἐκεῖ ἐπὶ τοῦ βουνοῦ· καὶ εἶπεν αὐτῷ Λάβαν, ὁ βουνὸς οὗτος μαρτυρεῖ ἀνὰ μέσον ἐμοῦ καὶ σοῦ σήμερον.
\vs{47}Καὶ ἐκάλεσεν αὐτὸν Λάβαν, βουνὸς τῆς μαρτυρίας· Ἰακὼβ δὲ ἐκάλεσεν αὐτὸν, βουνὸς μάρτυς.
\vs{48}Εἶπε δὲ Λάβαν τῷ Ἰακὼβ, ἰδοὺ ὁ βουνὸς οὗτος καὶ ἡ στήλη, ἣν ἔστησα ἀνὰ μέσον ἐμοῦ καὶ σοῦ· μαρτυρεῖ ὁ βουνὸς οὗτος, καὶ μαρτυρεῖ ἡ στήλη αὕτη· διὰ τοῦτο ἐκλήθη τὸ ὄνομα, βουνὸς μαρτυρεῖ.
\vs{49}Καὶ ἡ ὅρασις, ἣν εἶπεν, ἐπίδοι ὁ Θεὸς ἀνὰ μέσον ἐμοῦ καὶ σοῦ· ὅτι ἀποστησόμεθα ἕτερος ἀφʼ ἑτέρου.
\vs{50}Εἰ ταπεινώσεις τὰς θυγατέρας μου, εἰ λάβῃς γυναῖκας πρὸς ταῖς θυγατράσι μου, ὅρα, οὐθεὶς μεθʼ ἡμῶν ἐστιν ὁρῶν· Θεὸς μάρτυς μεταξὺ ἐμοῦ καὶ μεταξὺ σοῦ.
\vs{50a}Καὶ εἶπε Λάβαν τῷ Ἰακὼβ, ἰδοὺ ὁ βουνὸς οὗτος καὶ μάρτυς ἡ στήλη αὕτη.
\vs{52}Ἐάν τε γὰρ ἐγὼ μὴ διαβῶ πρός σε, μήτε σὺ διαβῇς πρός με τὸν βουνὸν τοῦτον καὶ τὴν στήλην ταύτην ἐπὶ κακίᾳ.
\vs{53}Ὁ Θεὸς Ἁβραὰμ καὶ ὁ Θεὸς Ναχὼρ κρίναι ἀνὰ μέσον ἡμῶν· καὶ ὤμοσεν Ἰακὼβ κατὰ τοῦ φόβου τοῦ πατρὸς αὐτοῦ Ἰσαάκ.
\vs{54}Καὶ ἔθυσεν θυσίαν ἐν τῷ ὄρει· καὶ ἐκάλεσε τοὺς ἀδελφοὺς αὐτοῦ, καὶ ἔφαγον καὶ ἔπιον, καὶ ἐκοιμήθησαν ἐν τῷ ὄρει.

\ch{32}Ἀναστὰς δὲ Λάβαν τὸ πρωῒ, κατεφίλησε τοὺς υἱοὺς καὶ τὰς θυγατέρας αὐτοῦ, καὶ εὐλόγησεν αὐτούς· καὶ ἀποστραφεὶς Λάβαν ἀπῆλθεν εἰς τὸν τόπον αὐτοῦ.

\vs{2}Καὶ Ἰακὼβ ἀπῆλθεν εἰς τὴν ὁδὸν ἑαυτοῦ· καὶ ἀναβλέψας εἶδε παρεμβολὴν Θεοῦ παρεμβεβληκυῖαν· καὶ συνήντησαν αὐτῷ οἱ Ἄγγελοι τοῦ Θεοῦ.
\vs{3}Εἶπε δὲ Ἰακὼβ, ἡνίκα εἶδεν αὐτοὺς, παρεμβολὴ Θεοῦ αὕτη· καὶ ἐκάλεσε τὸ ὄνομα τοῦ τόπου ἐκείνου, Παρεμβολαί.

\vs{4}Ἀπέστειλε δὲ Ἰακὼβ ἀγγέλους ἔμπροσθεν αὐτοῦ πρὸς Ἡσαῦ τὸν ἀδελφὸν αὐτοῦ εἰς γῆν Σηεὶρ, εἰς χώραν Ἐδώμ.
\vs{5}Καὶ ἐνετείλατο αὐτοῖς, λέγων, οὕτως ἐρεῖτε τῷ κυρίῳ μου Ἡσαῦ· οὕτως λέγει ὁ παῖς σου Ἰακώβ· μετὰ Λάβαν παρῴκησα, καὶ ἐχρόνισα ἕως τοῦ νῦν.
\vs{6}Καὶ ἐγένοντό μοι βόες, καὶ ὄνοι, καὶ πρόβατα, καὶ παῖδες, καὶ παιδίσκαι· καὶ ἀπέστειλα ἀναγγεῖλαι τῷ κυρίῳ μου Ἡσαῦ, ἵνα εὕρῃ ὁ παῖς σου χάριν ἐναντίον σου.
\vs{7}Καὶ ἀνέστρεψαν οἱ ἄγγελοι πρὸς Ἰακὼβ, λέγοντες, ἤλθομεν πρὸς τὸν ἀδελφόν σου Ἡσαυ· καὶ ἰδοὺ αὐτὸς ἔρχεται εἰς συνάντησίν σου, καὶ τετρακόσιοι ἄνδρες μεθʼ αὐτοῦ.
\vs{8}Ἐφοβήθη δὲ Ἰακὼβ σφόδρα, καὶ ἠπορεῖτο· καὶ διεῖλε τὸν λαὸν τὸν μεθʼ ἑαυτοῦ, καὶ τοὺς βόας, καὶ τὰς καμήλους, καὶ τὰ πρόβατα, εἰς δύο παρεμβολάς.
\vs{9}Καὶ εἶπεν Ἰακὼβ, ἐὰν ἔλθῃ Ἡσαῦ εἰς παρεμβολὴν μίαν, καὶ κόψῃ αὐτὴν, ἔσται ἡ παρεμβολὴ ἡ δευτέρα εἰς τὸ σώζεσθαι.
\vs{10}Εἶπε δὲ Ἰακὼβ, ὁ Θεὸς τοῦ πατρός μου Ἁβραὰμ, καὶ ὁ Θεὸς τοῦ πατρός μου Ἰσαὰκ, Κύριε σὺ ὁ εἰπών μοι, ἀπότρεχε εἰς τὴν γῆν τῆς γενέσεώς σου, καὶ εὖ σε ποιήσω·
\vs{11}Ἱκανούσθω μοι ἀπὸ πάσης δικαιοσύνης, καὶ ἀπὸ πάσης ἀληθείας, ἧς ἐποίησας τῷ παιδί σου· ἐν γὰρ τῇ ῥάβδῳ μου ταύτῃ διέβην τὸν Ἰορδάνην τοῦτον· νυνὶ δὲ γέγονα εἰς δύο παρεμβολάς.
\vs{12}Ἐξελοῦ με ἐκ χειρὸς τοῦ ἀδελφοῦ μου, ἐκ χειρὸς Ἡσαῦ· ὅτι φοβοῦμαι ἐγὼ αὐτὸν, μή ποτε ἐλθὼν πατάξῃ με, καὶ μητέρα ἐπὶ τέκνοις.
\vs{13}Σὺ δὲ εἶπας, εὐ σε ποιήσω, καὶ θήσω τὸ σπέρμα σου ὡς τὴν ἄμμον τῆς θαλάσσης, ἣ οὐκ ἀριθμηθήσεται ὑπὸ τοῦ πλήθους.
\vs{14}Καὶ ἐκοιμήθη ἐκεῖ τὴν νύκτα ἐκείνην· καὶ ἔλαβεν ὧν ἔφερεν δῶρα· καὶ ἐξαπέστειλεν Ἡσαῦ τῷ ἀδελφῷ αὐτοῦ,
\vs{15}αἶγας διακοσίας, τράγους εἴκοσι, πρόβατα διακόσια, κριοὺς εἴκοσι,
\vs{16}καμήλους θηλαζούσας καὶ τὰ παιδία αὐτῶν τριάκοντα, βόας τεσσαράκοντα, ταύρους δέκα, ὄνους εἴκοσι, καὶ πώλους δέκα.
\vs{17}Καὶ ἔδωκεν αὐτὰ τοῖς παισὶν αὐτοῦ ποίμνιον κατὰ μόνας· εἶπε δὲ τοῖς παισὶν αὐτοῦ, προπορεύεσθε ἔμπροσθέν μου, καὶ διάστημα ποιεῖτε ἀνὰ μέσον ποίμνης καὶ ποίμνης.
\vs{18}Καὶ ἐνετείλατο τῷ πρώτῳ, λέγων, ἐάν σοι συναντήσῃ Ἡσαῦ ὁ ἀδελφός μου, καὶ ἐρωτᾷ σε, λέγων, τίνος εἶ; καὶ ποῦ πορεύῃ; καὶ τίνος ταῦτα τὰ προπορευόμενά σου;
\vs{19}Ἐρεῖς, τοῦ παιδός σου Ἰακώβ· δῶρα ἀπέσταλκε τῷ κυρίῳ μου Ἡσαῦ· καὶ ἰδοὺ αὐτὸς ὀπίσω ἡμῶν.
\vs{20}Καὶ ἐνετείλατο τῷ πρώτῳ, καὶ τῷ δευτέρῳ, καὶ τῷ τρίτῳ, καὶ πᾶσι τοῖς προπορευομένοις ὀπίσω τῶν ποιμνίων τούτων, λέγων, κατὰ τὸ ῥῆμα τοῦτο λαλήσατε Ἡσαῦ ἐν τῷ εὑρεῖν ὑμᾶς αὐτόν·
\vs{21}Καὶ ἐρεῖτε, ἰδοὺ ὁ παῖς σου Ἰακὼβ παραγίνεται ὀπίσω ἡμῶν· εἶπε γὰρ, ἐξιλάσομαι τὸ πρόσωπον αὐτοῦ ἐν τοῖς δώροις τοῖς προπορευομένοις αὐτοῦ, καὶ μετὰ τοῦτο ὄψομαι τὸ πρόσωπον αὐτοῦ· ἴσως γὰρ προσδέξεται τὸ πρόσωπόν μου.
\vs{22}Καὶ προεπορεύετο τὰ δῶρα κατὰ πρόσωπον αὐτοῦ· αὐτὸς δὲ ἐκοιμήθη τὴν νύκτα ἐκείνην ἐν τῇ παρεμβολῇ.
\vs{23}Ἀναστὰς δὲ τὴν νύκτα ἐκείνην, ἔλαβε τὰς δύο γυναῖκας, καὶ τὰς δύο παιδίσκας, καὶ τὰ ἕνδεκα παιδία αὐτοῦ, καὶ διέβη τὴν διάβασιν τοῦ Ἰαβώχ.
\vs{24}Καὶ ἔλαβεν αὐτοὺς, καὶ διέβη τὸν χειμάῤῥουν, καὶ διεβίβασε πάντα τὰ αὐτοῦ.

\vs{25}Ὑπελείφθη δὲ Ἰακὼβ μόνος· καὶ ἐπάλαιεν ἄνθρωπος μετʼ αὐτοῦ ἕως πρωΐ.
\vs{26}Εἶδε δὲ ὅτι οὐ δύναται πρὸς αὐτόν· καὶ ἥψατο τοῦ πλάτους τοῦ μηροῦ αὐτοῦ, καὶ ἐνάρκησε τὸ πλάτος τοῦ μηροῦ Ἰακὼβ ἐν τῷ παλαίειν αὐτὸν μετʼ αὐτοῦ.
\vs{27}Καὶ εἶπεν αὐτῷ, ἀπόστειλόν με, ἀνέβη γὰρ ὁ ὄρθρος. ὁ δὲ εἶπεν, οὐ μή σε ἀποστείλω, ἐὰν μή με εὐλογήσῃς.
\vs{28}Εἶπε δὲ αὐτῷ, τί τὸ ὄνομά σου ἐστίν; ὁ δὲ εἶπεν, Ἰακώβ.
\vs{29}Καὶ εἶπεν αὐτῷ, οὐ κληθήσεται ἔτι τὸ ὄνομά σου Ἰακὼβ, ἀλλʼ Ἰσραὴλ ἔσται τὸ ὄνομά σου· ὅτι ἐνίσχυσας μετὰ Θεοῦ, καὶ μετὰ ἀνθρώπων δυνατὸς ἔσῃ.
\vs{30}Ἠρώτησε δὲ Ἰακὼβ, καὶ εἶπεν, ἀνάγγειλόν μοι τὸ ὄνομά σου· καὶ εἶπεν, ἱνατί τοῦτο ἐρωτᾷς σὺ τὸ ὄνομά μου; καὶ εὐλόγησεν αὐτὸν ἐκεῖ.
\vs{31}Καὶ ἐκάλεσεν Ἰακὼβ τὸ ὄνομα τοῦ τόπου ἐκείνου, εἶδος Θεοῦ· εἶδον γὰρ Θεὸν πρόσωπον πρὸς πρὸσωπον, καὶ ἐσώθη μου ἡ ψυχή.
\vs{32}Ἀνέτειλεν δὲ αὐτῷ ὁ ἥλιος, ἡνίκα παρῆλθε τὸ εἶδος τοῦ Θεοῦ· αὐτὸς δὲ ἐπέσκαζε τῷ μηρῷ αὐτοῦ.
\vs{33}Ἕνεκεν τούτου οὐ μὴ φάγωσιν υἱοὶ Ἰσραὴλ τὸ νεῦρον, ὃ ἐνάρκησεν, ὅ ἐστιν ἐπὶ τοῦ πλάτους τοῦ μηροῦ, ἕως τῆς ἡμέρας ταύτης, ὅτι ἥψατο τοῦ πλάτους τοῦ μηροῦ Ἰακὼβ τοῦ νεύρου, ὃ ἐνάρκησεν.

\ch{33}
Ἀναβλέψας δὲ Ἰακὼβ τοῖς ὀφθαλμοῖς αὐτοῦ εἶδε· καὶ ἰδοὺ Ἡσαῦ ὁ ἀδελφὸς αὐτοῦ ἐρχόμενος, καὶ τετρακόσιοι ἄνδρες μετʼ αὐτοῦ· καὶ διεῖλεν Ἰακὼβ τὰ παιδία ἐπὶ Λείαν, καὶ ἐπὶ Ῥαχὴλ, καὶ τὰς δύο παιδίσκας.
\vs{2}Καὶ ἔθετο τὰς δύο παιδίσκας καὶ τοὺς υἱοὺς αὐτῶν ἐν πρώτοις, καὶ Λείαν καὶ τὰ παιδία αὐτῆς ὀπίσω, καὶ Ῥαχὴλ καὶ Ἰωσὴφ ἐσχάτους.
\vs{3}Αὐτὸς δὲ προῆλθεν ἔμπροσθεν αὐτῶν· καὶ προσεκύνησεν ἐπὶ τὴν γῆν ἑπτάκις, ἕως τοῦ ἐγγίσαι τῷ ἀδελφῷ αὐτοῦ.
\vs{4}Καὶ προσέδραμεν Ἡσαῦ εἰς συνάντησιν αὐτῷ· καὶ περιλαβὼν αὐτὸν προσέπεσεν ἐπὶ τὸν τράχηλον αὐτοῦ, καὶ κατεφίλησεν αὐτόν· καὶ ἔκλαυσαν ἀμφότεροι.
\vs{5}Καὶ ἀναβλέψας Ἡσαῦ εἶδε τὰς γυναῖκας καὶ τὰ παιδία· καὶ εἶπε, τί ταῦτά σοι ἐστίν; ὁ δὲ εἶπε, τὰ παιδία, οἷς ἠλέησεν ὁ Θεὸς τὸν παῖδά σου.
\vs{6}Καὶ προσήγγισαν αἱ παιδίσκαι καὶ τὰ τέκνα αὐτῶν, καὶ προσεκύνησαν.
\vs{7}Καὶ προσήγγισε Λεία καὶ τὰ τέκνα αὐτῆς, καὶ προσεκύνησαν· καὶ μετὰ ταῦτα προσήγγισε Ῥαχὴλ καὶ Ἰωσὴφ, καὶ προσεκύνησαν.
\vs{8}Καὶ εἶπε, τί ταῦτά σοι ἐστὶν, πᾶσαι αἱ παρεμβολαὶ αὗται, αἷς ἀπήντηκα; ὁ δὲ εἶπεν, ἵνα εὕρῃ ὁ παῖς σου χάριν ἐναντίον σου, κύριε.
\vs{9}Εἶπε δὲ Ἡσαῦ, ἔστι μοι πολλὰ, ἀδελφέ· ἔστω σοι τὰ σά.
\vs{10}Εἶπε δὲ Ἰακὼβ, εἰ εὓρον χάριν ἐναντίον σου, δέξαι τὰ δῶρα διὰ τῶν ἐμῶν χειρῶν· ἕνεκεν τούτου εἶδον τὸ πρόσωπόν σου, ὡς ἄν τις ἴδοι πρόσωπον Θεοῦ, καὶ εὐδοκήσεις με.
\vs{11}Λάβε τὰς εὐλογίας μου, ἃς ἤνεγκά σοι, ὅτι ἠλέησέ με ὁ Θεὸς, καὶ ἔστι μοι πάντα· καὶ ἐβιάσατο αὐτὸν, καὶ ἔλαβε.
\vs{12}Καὶ εἶπεν, ἀπάραντες πορευσώμεθα ἐπʼ εὐθεῖαν.
\vs{13}Εἶπε δὲ αὐτῷ, ὁ κύριός μου γινώσκει, ὅτι τὰ παιδία ἁπαλώτερα, καὶ τὰ πρόβατα καὶ αἱ βόες λοχεύονται ἐπʼ ἐμέ· ἐὰν οὖν καταδιώξω αὐτὰ ἡμέραν μίαν, ἀποθανοῦνται πάντα τὰ κτήνη.
\vs{14}Προελθέτω ὁ κύριός μου ἔμπροσθεν τοῦ παιδὸς αὐτοῦ· ἐγὼ δὲ ἐνισχύσω ἐν τῇ ὁδῷ κατὰ σχολὴν τῆς πορεύσεως τῆς ἐναντίον μου, καὶ κατὰ πόδα τῶν παιδαρίων, ἕως τοῦ ἐλθεῖν με πρὸς τὸν κύριόν μου εἰς Σηείρ.
\vs{15}Εἶπε δὲ Ἡσαῦ, καταλείψω μετὰ σοῦ ἀπὸ τοῦ λαοῦ τοῦ μετʼ ἐμοῦ· ὁ δὲ εἶπεν, ἱνατί τοῦτο; ἱκανὸν ὅτι εὗρον χάριν ἐναντίον σου, κύριε.
\vs{16}Ἀπέστρεψε δὲ Ἡσαῦ ἐν τῇ ἡμέρᾳ ἐκείνῃ εἰς τὴν ὁδὸν αὐτοῦ εἰς Σηείρ.
\vs{17}Καὶ Ἰακὼβ ἀπαίρει εἰς σκηνὰς, καὶ ἐποίησεν ἑαυτῷ ἐκεῖ οἰκίας, καὶ τοῖς κτήνεσιν αὐτοῦ ἐποίησε σκηνάς· διὰ τοῦτο ἐκάλεσε τὸ ὄνομα τοῦ τόπου ἐκείνου, Σκηναί.

\vs{18}Καὶ ἦλθεν Ἰακὼβ εἰς Σαλὴμ, πόλιν Σηκίμων, ἥ ἐστιν ἐν γῇ Χαναὰν, ὅτε ἐπανῆλθεν ἐκ τῆς Μεσοποταμίας Συρίας· καὶ παρενέλαβε κατὰ πρόσωπον τῆς πόλεως.
\vs{19}Καὶ ἐκτήσατο τὴν μερίδα τοῦ ἀγροῦ, οὗ ἔστησεν ἐκεῖ τὴν σκηνὴν αὐτοῦ, παρὰ Ἐμμὼρ πατρὸς Συχὲμ, ἑκατὸν ἀμνῶν.
\vs{20}Καὶ ἔστησεν ἐκεῖ θυσιαστήριον, καὶ ἐπεκαλέσατο τὸν Θεὸν Ἰσραήλ.

\ch{34}
Ἐξῆλθε δὲ Δείνα, ἡ θυγάτηρ Λείας, ἣν ἔτεκε τῷ Ἰακώβ, καταμαθεῖν τὰς θυγατέρας τῶν ἐγχωρίων.
\vs{2}Καὶ εἶδεν αὐτὴν Συχὲμ ὁ υἱὸς Ἐμμὼρ ὁ Εὐαῖος, ὁ ἄρχων τῆς γῆς· καὶ λαβὼν αὐτὴν, ἐκοιμήθη μετʼ αὐτῆς, καὶ ἐταπείνωσεν αὐτήν.
\vs{3}Καὶ προσέσχε τῇ ψυχῇ Δείνας τῆς θυγατρὸς Ἰακώβ· καὶ ἠγάπησε τὴν παρθένον· καὶ ἐλάλησε κατὰ τὴν διάνοιαν τῆς παρθένου αυτῇ.
\vs{4}Εἶπε Συχὲμ πρὸς Ἐμμὼρ τὸν πατέρα αὐτοῦ, λέγων, λάβε μοι τὴν παῖδα ταύτην εἰς γυναῖκα.
\vs{5}Ἰακὼβ δὲ ἤκουσεν, ὅτι ἐμίανεν ὁ υἱὸς Ἐμμὼρ Δείναν τὴν θυγατέρα αὐτοῦ· οἱ δὲ υἱοὶ αὐτοῦ ἦσαν μετὰ τῶν κτηνῶν αὐτοῦ ἐν τῷ πεδίῳ· παρεσιώπησε δὲ Ἰακὼβ, ἕως τοῦ ἐλθεῖν αὐτούς.
\vs{6}Ἐξῆλθε δὲ Ἐμμὼρ ὁ πατὴρ Συχὲμ πρὸς Ἰακὼβ, λαλῆσαι αὐτῷ.
\vs{7}Οἱ δὲ υἱοὶ Ἰακὼβ ἦλθον ἐκ τοῦ πεδίου· ὡς δὲ ἤκουσαν, κατενύγησαν οἱ ἄνδρες, καὶ λυπηρὸν ἦν αὐτοῖς σφόδρα· ὅτι ἄσχημον ἐποίησεν ἐν Ἰσραὴλ, κοιμηθεὶς μετὰ τῆς θυγατρὸς Ἰακώβ· καὶ οὐχ οὕτως ἔσται.
\vs{8}Καὶ ἐλάλησεν Ἐμμὼρ αὐτοῖς, λέγων, Συχὲμ ὁ υἱός μου προείλετο τῇ ψυχῇ τὴν θυγατέρα ὑμῶν· δότε οὖν αὐτὴν αὐτῷ γυναῖκα,
\vs{9}καὶ ἐπιγαμβρεύσασθε ἡμῖν· τὰς θυγατέρας ὑμῶν δότε ἡμῖν, καὶ τὰς θυγατέρας ἡμῶν λάβετε τοῖς υἱοῖς ὑμῶν.
\vs{10}Καὶ ἐν ἡμῖν κατοικεῖτε· καὶ ἡ γῆ ἰδοὺ πλατεῖα ἐναντίον ὑμῶν· κατοικεῖτε, καὶ ἐμπορεύεσθε ἐπʼ αὐτῆς, καὶ ἐγκτᾶσθε ἐν αὐτῇ.
\vs{11}Εἶπε δὲ Συχὲμ πρὸς τὸν πατέρα αὐτῆς, καὶ πρὸς τοὺς ἀδελφοὺς αὐτῆς, εὕροιμι χάριν ἐναντίον ὑμῶν· καὶ ὃ ἐὰν εἴπητε, δώσομεν.
\vs{12}Πληθύνατε τὴν φερνὴν σφόδρα, καὶ δώσω καθότι ἂν εἴπητέ μοι, καὶ δώσετέ μοι τὴν παῖδα ταύτην εἰς γυναῖκα.

\vs{13}Ἀπεκρίθησαν δὲ οἱ υἱοὶ Ἰακὼβ τῷ Συχὲμ, καὶ Ἐμμὼρ τῷ πατρὶ αὐτοῦ, μετὰ δόλου· καὶ ἐλάλησαν αὐτοῖς, ὅτι ἐμίαναν Δείναν τὴν ἀδελφὴν αὐτῶν.
\vs{14}Καὶ εἶπαν αὐτοῖς Συμεὼν καὶ Λευὶ οἱ ἀδελφοὶ Δείνας, οὐ δυνησόμεθα ποιῆσαι τὸ ῥῆμα τοῦτο, δοῦναι τὴν ἀδελφὴν ἡμῶν ἀνθρώπῳ, ὃς ἔχει ἀκροβυστίαν· ἔστι γὰρ ὄνειδος ἡμῖν.
\vs{15}Μόνον ἐν τούτῳ ὁμοιωθησόμεθα ὑμῖν, καὶ κατοικήσομεν ἐν ὑμῖν, ἐὰν γένησθε ὡς ἡμεῖς καὶ ὑμεῖς, ἐν τῷ περιτμηθῆναι ὑμῶν πᾶν ἀρσενικόν.
\vs{16}Καὶ δώσομεν τὰς θυγατέρας ἡμῶν ὑμῖν, καὶ ἀπὸ τῶν θυγατέρων ὑμῶν ληψόμεθα ἡμῖν γυναῖκας, καὶ οἰκήσομεν παρʼ ὑμῖν, καὶ ἐσόμεθα ὡς γένος ἕν.
\vs{17}Ἐὰν δὲ μὴ εἰσακούσητε ἡμῶν τοῦ περιτεμέσθαι, λαβόντες τὴν θυγατέρα ἡμῶν ἀπελευσόμεθα.
\vs{18}Καὶ ἤρεσαν οἱ λόγοι ἐναντίον Ἐμμὼρ, καὶ ἐναντίον Συχὲμ τοῦ υἱοῦ Ἐμμώρ.
\vs{19}Καὶ οὐκ ἐχρόνισεν ὁ νεανίσκος τοῦ ποιῆσαι τὸ ῥῆμα τοῦτο· ἐνέκειτο γὰρ τῇ θυγατρὶ Ἰακώβ· αὐτὸς δὲ ἦν ἐνδοξότατος πάντων τῶν ἐν τῷ οἴκῳ τοῦ πατρὸς αὐτοῦ.
\vs{20}Ἦλθε δὲ Ἐμμὼρ καὶ Συχὲμ ὁ υἱὸς αὐτοῦ πρὸς τὴν πύλην τῆς πόλεως αὐτῶν, καὶ ἐλάλησαν πρὸς τοὺς ἄνδρας τῆς πόλεως αὐτῶν, λέγοντες,
\vs{21}Οἱ ἄνθρωποι οὗτοι εἰρήνικοί εἰσι, μεθʼ ἡμῶν οἰκείτωσαν επὶ τῆς γῆς, καὶ ἐμπορευέσθωσαν αὐτήν· ἡ δὲ γῆ ἰδοὺ πλατεῖα ἐναντίον αὐτῶν· τὰς θυγατέρας αὐτῶν ληψόμεθα ἡμῖν γυναῖκας, καὶ τὰς θυγατέρας ἡμῶν δώσομεν αὐτοῖς.
\vs{22}Ἐν τούτῳ μόνον ὁμοιωθήσονται ἡμῖν οἱ ἄνθρωποι τοῦ κατοικεῖν μεθʼ ἡμῶν, ὥστε εἶναι λαὸν ἕνα, ἐν τῷ περιτεμέσθαι ἡμῶν πᾶν ἀρσενικὸν, καθὰ καὶ αὐτοὶ περιτέτμηνται.
\vs{23}Καὶ τὰ κτήνη αὐτῶν, καὶ τὰ τετράποδα, καὶ τὰ ὑπάρχοντα αὐτῶν, οὐχ ἡμῶν ἔσται; μόνον ἐν τούτῳ ὁμοιωθῶμεν αὐτοῖς, καὶ οἰκήσουσι μεθʼ ἡμῶν.
\vs{24}Καὶ εἰσήκουσαν Ἐμμὼρ καὶ Συχὲμ τοῦ υἱοῦ αὐτοῦ πάντες οἱ ἐμπορευόμενοι τὴν πύλην τῆς πόλεως αὐτῶν· καὶ περιετέμοντο τὴν σάρκα τῆς ἀκροβυστίας αὐτῶν πᾶς ἄρσην.

\vs{25}Ἐγένετο δὲ ἐν τῇ ἡμέρᾳ τῇ τρίτῃ, ὅτε ἦσαν ἐν τῷ πόνῳ, ἔλαβον οἱ δύο υἱοὶ Ἰακὼβ Συμεὼν καὶ Λευὶ, ἀδελφοὶ Δείνας, ἕκαστος τὴν μάχαιραν αὐτοῦ, καὶ εἰσῆλθον εἰς τὴν πόλιν ἀσφαλὼς, καὶ ἀπέκτειναν πᾶν ἀρσενικόν.
\vs{26}Τόν τε Ἐμμὼρ καὶ Συχὲμ τὸν υἱὸν αὐτοῦ ἀπέκτειναν ἐν στόματι μαχαίρας· καὶ ἔλαβον τὴν Δείναν ἐκ τοῦ οἴκου τοῦ Συχὲμ, καὶ ἐξῆλθον.
\vs{27}Οἱ δὲ υἱοὶ Ἰακὼβ εἰσῆλθον ἐπὶ τοὺς τραυματίας, καὶ διήρπασαν τὴν πόλιν, ἐν ᾗ ἐμίαναν Δείναν τὴν ἀδελφὴν αὐτῶν.
\vs{28}Καὶ τὰ πρόβατα αὐτῶν, καὶ τοὺς βόας αὐτῶν, καὶ τοὺς ὄνους αὐτῶν, ὅσα τε ἦν ἐν τῇ πόλει, καὶ ὅσα ἦν ἐν τῷ πεδίῳ, ἔλαβον.
\vs{29}Καὶ πάντα τὰ σώματα αὐτῶν, καὶ πᾶσαν τὴν ἀποσκευὴν αὐτῶν, καὶ τὰς γυναῖκας αὐτῶν ἠχμαλώτευσαν· καὶ διήρπασαν ὅσα τε ἦν ἐν τῇ πόλει, καὶ ὅσα ἦν ἐν ταῖς οἰκίαις.
\vs{30}Εἶπε δὲ Ἰακὼβ πρὸς Συμεὼν καὶ Λευὶ, μισητόν με πεποιήκατε, ὥστε πονηρόν με εἶναι πᾶσι τοῖς κατοικοῦσι τὴν γῆν, ἔν τε τοῖς Χαναναίοις, καὶ ἐν τοῖς Φερεζαίοις· ἐγὼ δὲ ὀλιγοστός εἰμι ἐν ἀριθμῷ· καὶ συναχθέντες ἐπʼ ἐμὲ συγκόψουσί με, καὶ ἐκτριβήσομαι ἐγὼ, καὶ ὁ οἶκός μου.
\vs{31}Οἱ δὲ εἶπαν, ἀλλʼ ὡσεὶ πόρνῃ χρήσονται τῇ ἀδελφῇ ἡμῶν;

\ch{35}
Εἶπε δὲ ὁ Θεὸς πρὸς Ἰακὼβ, ἀναστὰς ἀνάβηθι εἰς τὸν τόπον Βαιθὴλ, καὶ οἴκει ἐκεῖ· καὶ ποίησον ἐκεῖ θυσιαστήριον τῷ Θεῷ τῷ ὀφθέντι σοι, ἐν τῷ ἀποδιδράσκειν σε ἀπὸ προσώπου Ἡσαῦ τοῦ ἀδελφοῦ σου.
\vs{2}Εἶπε δὲ Ἰακὼβ τῷ οἴκῳ αὐτοῦ, καὶ πᾶσι τοῖς μετʼ αὐτοῦ, ἄρατε τοὺς θεοὺς τοὺς ἀλλοτρίους τοὺς μεθʼ ὑμῶν ἐκ μέσου ὑμῶν, καὶ καθαρίσθητε, καὶ ἀλλάξατε τὰς στολὰς ὑμῶν.
\vs{3}Καὶ ἀναστάντες ἀναβῶμεν εἰς Βαιθὴλ, καὶ ποιήσωμεν ἐκεῖ θυσιαστήριον τῷ Θεῷ τῷ ἐπακούσαντί μου ἐν ἡμέρᾳ θλίψεως, ὃς ἦν μετʼ ἐμοῦ, καὶ διέσωσέ με ἐν τῇ ὁδῷ, ᾗ ἐπορεύθην.
\vs{4}Καὶ ἔδωκαν τῷ Ἰακὼβ τοὺς θεοὺς τοὺς ἀλλοτρίους, οἳ ἦσαν ἐν ταῖς χερσὶν αὐτῶν, καὶ τὰ ἐνώτια τὰ ἐν τοῖς ὠσὶν αὐτῶν· καὶ κατέκρυψεν αὐτὰ Ἰακὼβ ὑπὸ τὴν τερέβινθον τὴν ἐν Σηκίμοις· καὶ ἀπώλεσεν αὐτὰ, ἕως τῆς σήμερον ἡμέρας.
\vs{5}Καὶ ἐξῇρεν Ἰσραὴλ ἐκ Σηκίμων· καὶ ἐγένετο φόβος Θεοῦ ἐπὶ τὰς πόλεις τὰς κύκλῳ αὐτῶν, καὶ οὐ κατεδίωξαν ὀπίσω τῶν υἱῶν Ἰσραήλ.
\vs{6}Ἦλθε δὲ Ἰακὼβ εἰς Λουζὰ ἥ ἐστιν ἐν γῇ Χαναὰν, ἥ ἐστι Βαιθὴλ, αὐτὸς, καὶ πᾶς ὁ λαὸς, ὃς ἦν μετʼ αὐτοῦ.
\vs{7}Καὶ ᾠκοδόμησεν ἐκεῖ θυσιαστήριον, καὶ ἐκάλεσε τὸ ὄνομα τοῦ τόπου, Βαιθήλ· ἐκεῖ γὰρ ἐφάνη αὐτῷ ὁ Θεὸς, ἐν τῷ ἀποδιδράσκειν αὐτὸν ἀπὸ προσώπου Ἡσαῦ τοῦ ἀδελφοῦ αὐτοῦ.

\vs{8}Ἀπέθανε δὲ Δεβόῤῥα, ἡ τρόφος Ῥεβέκκας, καὶ ἐτάφη κατώτερον Βαιθὴλ ὑπὸ τὴν βάλανον· καὶ ἐκάλεσεν Ἰακὼβ τὸ ὄνομα αὐτῆς, βάλανος πένθους.
\vs{9}Ὤφθη δὲ ὁ Θεὸς τῷ Ἰακὼβ ἔτι ἐν Λουζᾷ, ὅτε παρεγένετο ἐκ Μεσοποταμίας τῆς Συρίας· καὶ εὐλόγησεν αὐτὸν ὁ Θεὸς.
\vs{10}Καὶ εἶπεν αὐτῷ ὁ Θεὸς, τὸ ὄνομά σου οὐ κληθήσεται ἔτι Ἰακὼβ, ἀλλʼ Ἰσραὴλ ἔσται τὸ ὄνομά σου· καὶ ἐκάλεσε τὸ ὄνομα αὐτοῦ Ἰσραήλ.
\vs{11}Εἶπε δὲ αὐτῷ ὁ Θεὸς, ἐγὼ ὁ Θεός σου· αὐξάνου, καὶ πληθύνου· ἔθνη καὶ συναγωγαὶ ἐθνῶν ἔσονται ἐκ σοῦ, καὶ βασιλεῖς ἐκ τῆς ὀσφύος σου ἐξελεύσονται.
\vs{12}Καὶ τὴν γῆν, ἣν ἔδωκα Ἁβραὰμ καὶ Ἰσαὰκ, σοὶ δέδωκα αὐτήν· σοὶ ἔσται· καὶ τῷ σπέρματί σου μετὰ σὲ δώσω τῆν γῆν ταύτην.
\vs{13}Ἀνέβη δὲ ὁ Θεὸς ἀπʼ αὐτοῦ ἐκ τοῦ τόπου, οὗ ἐλάλησε μετʼ αὐτοῦ.
\vs{14}Καὶ ἔστησεν Ἰακὼβ στήλην ἐν τῷ τόπῳ, ᾧ ἐλάλησε μετʼ αὐτοῦ ὁ Θεὸς, στήλην λιθίνην· καὶ ἔσπεισεν ἐπʼ αὐτὴν σπονδὴν, καὶ ἐπέχεεν ἐπʼ αὐτὴν ἔλαιον.
\vs{15}Καὶ ἐκάλεσεν Ἰακὼβ τὸ ὄνομα τοῦ τόπου, ἐν ᾧ ἐλάλησε μετʼ αὐτοῦ ἐκεῖ ὁ Θεὸς, Βαιθήλ.
\vs{16}Ἀπάρας δὲ Ἰακὼβ ἐκ Βαιθὴλ, ἔπηξε τὴν σκηνὴν αὐτοῦ ἐπέκεινα τοῦ πύργου Γαδέρ· ἐγένετο δὲ ἡνίκα ἤγγισεν εἰς Χαβραθὰ τοῦ ἐλθεῖν εἰς τὴν Ἐφραθᾶ, ἔτεκε Ῥαχήλ· καὶ ἐδυστόκησεν ἐν τῷ τοκετῷ.
\vs{17}Ἐγένετο δὲ ἐν τῷ σκληρὼς αὐτὴν τίκτειν, εἶπεν αὐτῇ ἡ μαῖα, θάρσει, καὶ γὰρ οὗτός σοι ἐστὶν υἱός.
\vs{18}Ἐγένετο δὲ ἐν τῷ ἀφιέναι αὐτὴν τὴν ψυχὴν, ἀπέθνησκε γὰρ, ἐκάλεσε τὸ ὄνομα αὐτοῦ, υἱὸς ὀδύνης μου· ὁ δὲ πατὴρ ἐκάλεσεν τὸ ὄνομα αὐτοῦ, Βενιαμίν.
\vs{19}Ἀπέθανε δὲ Ῥαχὴλ, καὶ ἐτάφη ἐν τῇ ὁδῷ τοῦ ἱπποδρόμου Ἐφραθᾶ· αὕτη ἐστὶ Βηθλεέμ.
\vs{20}Καὶ ἔστησεν Ἰακὼβ στήλην ἐπὶ τοῦ μνημείου αὐτῆς· αὕτη ἐστὶν ἡ στήλη ἐπὶ τοῦ μνημείου Ῥαχὴλ ἕως τῆς ἡμέρας ταύτης.
\vs{22}Ἐγένετο δὲ ἡνίκα κατῴκησεν Ἰσραὴλ ἐν τῇ γῇ ἐκείνῃ, ἐπορεύθη Ῥουβὴν, καὶ ἐκοιμήθη μετὰ Βαλλὰς, τῆς παλλακῆς τοῦ πατρὸς αὐτοῦ Ἰακώβ· καὶ ἤκουσεν Ἰσραὴλ, καὶ πονηρὸν ἐφάνη ἐναντίον αὐτοῦ.

Ἦσαν δὲ οἱ υἱοὶ Ἰακὼβ, δώδεκα.
\vs{23}Υἱοὶ Λείας, πρωτότοκος Ἰακὼβ, Ῥουβὴν, Συμεὼν, Λευὶ, Ἰούδας, Ἰσσάχαρ, Ζαβουλών.
\vs{24}Υἱοὶ δὲ Ῥαχὴλ, Ἰωσὴφ, καὶ Βενιαμίν.
\vs{25}Υἱοὶ δὲ Βαλλᾶς παιδίσκης Ῥαχὴλ, Δαν, καὶ Νεφθαλείμ.
\vs{26}Υἱοὶ δὲ Ζελφᾶς παιδίσκης Λείας, Γὰδ, καὶ Ἀσήρ· οὗτοι υἱοὶ Ἰακὼβ, οἳ ἐγένοντο αὐτῷ ἐν Μεσοποταμίᾳ τῆς Συρίας.
\vs{27}Ἦλθε δὲ Ἰακὼβ πρὸς Ἰσαὰκ τὸν πατέρα αὐτοῦ εἰς Μαμβρῆ, εἰς πόλιν τοῦ πεδίου· αὕτη ἐστὶ Χεβρὼν ἐν γῇ Χαναὰν, οὗ παρῴκησεν Ἁβραὰμ καὶ Ἰσαάκ.
\vs{28}Ἐγένοντο δὲ αἱ ἡμέραι Ἰσαὰκ, ἃς ἔζησεν, ἔτη ἑκατὸν ὀγδοήκοντα.
\vs{29}Καὶ ἐκλείπων Ἰσαὰκ ἀπέθανε, καὶ προσετέθη πρὸς τὸ γένος αὐτοῦ πρεσβύτερος καὶ πλήρης ἡμερῶν· καὶ ἔθαψαν αὐτὸν Ἡσαῦ καὶ Ἰακὼβ οἱ υἱοὶ αὐτοῦ.

\ch{36}
Αὗται δὲ αἱ γενέσεις Ἡσαῦ· αὐτός ἐστιν Ἐδώμ.
\vs{2}Ἡσαῦ δὲ ἔλαβε τὰς γυναῖκας ἑαυτῷ ἀπὸ τῶν θυγατέρων τῶν Χαναναίων· τὴν Ἀδὰ, θυγατέρα Αἰλὼμ τοῦ Χετταίου· καὶ τὴν Ὀλιβεμὰ, θυγατέρα Ἀνὰ τοῦ υἱοῦ Σεβεγὼν τοῦ Εὐαίου.
\vs{3}Καὶ τὴν Βασεμὰθ, θυγατέρα Ἰσμαὴλ, ἀδελφὴν Ναβαιώθ.
\vs{4}Ἔτεκε δὲ αὐτῷ Ἀδὰ τὸν Ἑλιφάς· καὶ Βασεμὰθ ἔτεκε τὸν Ῥαγουήλ.
\vs{5}Καὶ Ὀλιβεμὰ ἔτεκε τὸν Ἰεοὺς, καὶ τὸν Ἰεγλὸμ, καὶ τὸν Κορέ· οὗτοι υἱοὶ Ἡσαῦ, οἳ ἐγένοντο αὐτῷ ἐν γῇ Χαναάν.
\vs{6}Ἔλαβε δὲ Ἡσαῦ τὰς γυναῖκας αὐτοῦ, καὶ τοὺς υἱοὺς αὐτοῦ, καὶ τὰς θυγατέρας αὐτοῦ, καὶ πάντα τὰ σώματα τοῦ οἴκου αὐτοῦ, καὶ πάντα τὰ ὑπάρχοντα αὐτοῦ, καὶ πάντα τὰ κτήνη, καὶ πάντα ὅσα ἐκτήσατο, καὶ πάντα ὅσα περιεποιήσατο ἐν γῇ Χαναάν· καὶ ἐπορεύθη Ἡσαῦ ἐκ τῆς γῆς Χαναὰν ἀπὸ προσώπου Ἰακὼβ τοῦ ἀδελφοῦ αὐτοῦ.
\vs{7}Ἦν γὰρ αὐτῶν τὰ ὑπάρχοντα πολλὰ, τοῦ οἰκεῖν ἅμα· καὶ οὐκ ἠδύνατο ἡ γῆ τῆς παροικήσεως αὐτῶν φέρειν αὐτοὺς, ἀπὸ τοῦ πλήθους τῶν ὑπαρχόντων αὐτῶν.
\vs{8}Κατῴκησε δὲ Ἡσαῦ ἐν τῷ ὄρει Σηείρ· Ἡσαῦ αὐτός ἐστιν Ἐδώμ.
\vs{9}Αὗται δὲ αἱ γενέσεις Ἡσαῦ πατρὸς Ἐδὼμ ἐν τῷ ὄρει Σηείρ.
\vs{10}Καὶ ταῦτα τὰ ὀνόματα τῶν υἱῶν Ἡσαῦ· Ἑλιφὰς υἱὸς Ἀδὰς γυναικὸς Ἡσαῦ· καὶ Ῥαγουὴλ υἱὸς Βασεμὰθ γυναικὸς Ἡσαῦ.
\vs{11}Ἐγένοντο δὲ Ἑλιφὰς υἱοὶ, Θαιμὰν, Ὠμὰρ, Σωφὰρ, Γοθὼμ, καὶ Κενέζ.
\vs{12}Θαμνὰ δὲ ἦν παλλακὴ Ἑλιφὰς τοῦ υἱοῦ Ἡσαῦ· καὶ ἔτεκε τῷ Ἑλιφὰς τὸν Ἀμαλήκ· οὗτοι υἱοὶ Ἀδὰς γυναικὸς Ἡσαῦ.
\vs{13}Οὗτοι δὲ υἱοὶ Ῥαγουὴλ, Ναχὼθ, Ζαρὲ, Σομὲ, καὶ Μοζέ· οὗτοι ἦσαν υἱοὶ Βασεμὰθ γυναικὸς Ἡσαῦ.
\vs{14}Οὗτοι δὲ υἱοὶ Ὀλιβεμὰς θυγατρὸς Ἀνὰ τοῦ υἱοῦ Σεβεγὼν, γυναικὸς Ἡσαῦ· ἔτεκε δὲ τῷ Ἡσαῦ τὸν Ἰεοὺς, καὶ τὸν Ἰεγλὸμ, καὶ τὸν Κορέ.
\vs{15}Οὗτοι ἡγεμόνες υἱοὶ Ἡσαῦ· υἱοὶ Ἑλιφὰς πρωτοτόκου Ἡσαῦ· ἡγεμὼν Θαιμὰν, ἡγεμὼν Ὠμὰρ, ἡγεμὼν Σωφὰρ, ἡγεμὼν Κενὲζ,
\vs{16}ἡγεμὼν Κορὲ, ἡγεμὼν Γοθὼμ, ἡγεμὼν Ἀμαλήκ· οὗτοι ἡγεμόνες Ἑλιφὰς ἐν γῇ Ἰδουμαίᾳ· οὗτοι υἱοὶ Ἀδάς.
\vs{17}Καὶ οὗτοι υἱοὶ Ῥαγουὴλ υἱοῦ Ἡσαῦ· ἡγεμὼν Ναχὼθ, ἡγεμὼν Ζαρὲ, ἡγεμὼν Σομὲ, ἡγεμὼν Μοζέ· οὗτοι ἡγεμόνες Ῥαγουὴλ ἐν γῇ Ἐδώμ· οὗτοι υἱοὶ Βασεμὰθ γυναικὸς Ἡσαῦ.
\vs{18}Οὗτοι δὲ υἱοὶ Ὀλιβεμὰς γυναικὸς Ἡσαῦ· ἡγεμὼν Ἰεοὺς, ἡγεμὼν Ἰεγλὸμ, ἡγεμὼν Κορέ· οὗτοι ἡγεμόνες Ὀλιβεμὰς θυγατρὸς Ἀνὰ γυναικὸς Ἡσαῦ.
\vs{19}Οὗτοι υἱοὶ Ἡσαῦ, καὶ οὗτοι ἡγεμόνες αὐτῶν· οὗτοί εἰσιν υἱοὶ Ἐδώμ.
\vs{20}Οὗτοι δὲ υἱοὶ Σηεὶρ τοῦ Χοῤῥαίου, τοῦ κατοικοῦντος τὴν γῆν· Λωτὰν, Σωβὰλ, Σεβεγὼν, Ἀνὰ,
\vs{21}καὶ Δησὼν, καὶ Ἀσὰρ, καὶ Ῥισών· οὗτοι ἡγεμόνες τοῦ Χοῤῥαίου, τοῦ υἱοῦ Σηεὶρ ἐν τῇ γῇ Ἐδώμ.
\vs{22}Ἐγένοντο δὲ υἱοὶ Λωτάν· Χοῤῥὶ, καὶ Αἱμάν· ἀδελφὴ δὲ Λωτὰν, Θαμνά.
\vs{23}Οὗτοι δὲ υἱοὶ Σωβάλ· Γωλὰμ, καὶ Μαναχὰθ, καὶ Γαιβὴλ, καὶ Σωφὰρ, καὶ Ὠμάρ.
\vs{24}Καὶ οὗτοι υἱοὶ Σεβεγὼν, Ἀϊὲ, καὶ Ἀνά· οὗτός ἐστιν Ἀνὰ, ὃς εὗρε τὸν Ἰαμεὶν ἐν τῇ ἐρήμῳ, ὅτε ἔνεμε τὰ ὑποζύγια Σεβεγὼν τοῦ πατρὸς αὐτοῦ·
\vs{25}Οὗτοι δὲ υἱοὶ Ἀνά· Δησὼν, καὶ Ὀλιβεμὰ θυγάτηρ Ἀνά.
\vs{26}Οὗτοι δὲ υἱοὶ Δησών· Ἀμαδὰ, καὶ Ἀσβὰν, καὶ Ἰθρὰν, καὶ Χαῤῥάν.
\vs{27}Οὗτοι δὲ υἱοὶ Ἀσάρ· Βαλαὰμ, καὶ Ζουκὰμ, καὶ Ἰουκάμ.
\vs{28}Οὗτοι δὲ υἱοὶ Ῥισὼν, Ὧς, καὶ Ἀράν.
\vs{29}Οὗτοι δὲ ἡγεμόνες Χοῤῥί· ἡγεμὼν Λωτὰν, ἡγεμὼν Σωβὰλ, ἡγεμὼν Σεβεγὼν, ἡγεμὼν Ἀνὰ,
\vs{30}ἡγεμὼν Δησὼν, ἡγεμὼν Ἀσὰρ, ἡγεμὼν Ῥισών· οὗτοι ἡγεμόνες Χοῤῥὶ ἐν ταῖς ἡγεμονίαις αὐτῶν ἐν γῇ Ἐδώμ.

\vs{31}Καὶ οὗτοι οἱ βασιλεῖς οἱ βασιλεύσαντες ἐν Ἐδὼμ, πρὸ τοῦ βασιλεῦσαι βασιλέα ἐν Ἰσραήλ.
\vs{32}Καὶ ἐβασίλευσεν ἐν Ἐδὼμ Βαλὰκ υἱὸς Βεώρ· καὶ ὄνομα τῇ πόλει αὐτοῦ, Δενναβά.
\vs{33}Ἀπέθανε δὲ Βαλὰκ, καὶ ἐβασίλευσεν ἀντʼ αὐτοῦ Ἰωβὰβ υἱὸς Ζαρὰ ἐκ Βοσόῤῥας.
\vs{34}Ἀπέθανε δὲ Ἰωβὰβ, καὶ ἐβασίλευσεν ἀντʼ αὐτοῦ Ἀσὼμ ἐκ τῆς γῆς Θαιμανών.
\vs{35}Ἀπέθανε δὲ Ἀσὼμ, καὶ ἐβασίλευσεν ἀντʼ αὐτοῦ Ἀδὰδ υἱὸς Βαρὰδ ὁ ἐκκόψας Μαδιὰμ ἐν τῷ πεδίῳ Μωάβ· καὶ ὄνομα τῇ πόλει αὐτοῦ Γετθαίμ.
\vs{36}Ἀπέθανε δὲ Ἀδὰδ, καὶ ἐβασίλευσεν ἀντʼ αὐτοῦ Σαμαδὰ ἐκ Μασσεκκάς.
\vs{37}Ἀπέθανε δὲ Σαμαδὰ, καὶ ἐβασίλευσεν ἀντʼ αὐτοῦ Σαοὺλ ἐκ Ῥοωβὼθ τῆς παρὰ ποταμόν.
\vs{38}Ἀπέθανε δὲ Σαοὺλ, καὶ ἐβασίλευσεν ἀντʼ αὐτοῦ Βαλλενὼν υἱὸς Ἀχοβώρ.
\vs{39}Ἀπέθανε δὲ Βαλλενὼν υἱὸς Ἀχοβὼρ, καὶ ἐβασίλευσεν ἀντʼ αὐτοῦ Ἀρὰδ υἱὸς Βαράδ· καὶ ὄνομα τῇ πόλει αὐτοῦ Φογώρ· ὄνομα δὲ τῇ γυναικὶ αὐτοῦ Μετεβεὴλ, θυγάτηρ Ματραῒθ, υἱοῦ Μαιζοώβ.
\vs{40}Ταῦτα τὰ ὀνόματα τῶν ἡγεμόνων Ἡσαῦ, ἐν ταῖς φυλαῖς αὐτῶν, κατὰ τόπον αὐτῶν, ἐν ταῖς χώραις αὐτῶν, καὶ ἐν τοῖς ἔθνεσιν αὐτῶν· ἡγεμὼν Θαμνὰ, ἡγεμὼν Γωλὰ, ἡγεμὼν Ἰεθὲρ,
\vs{41}ἡγεμὼν Ὁλιβεμὰς, ἡγεμὼν Ἡλὰς, ἡγεμὼν Φινὼν,
\vs{42}ἡγεμὼν Κενὲζ, ἡγεμὼν Θαιμὰν, ἡγεμὼν Μαζὰρ,
\vs{43}ἡγεμὼν Μαγεδιὴλ, ἡγεμὼν Ζαφωίν· οὗτοι ἡγεμόνες Ἐδὼμ, ἐν ταῖς κατῳκοδομημέναις ἐν τῇ γῇ τῆς κτήσεως αὐτῶν· οὗτος Ἡσαῦ πατὴρ Ἐδώμ.

\ch{37}
Κατῴκει δὲ Ἰακὼβ ἐν τῇ γῇ, οὗ παρῴκησεν ὁ πατὴρ αὐτοῦ ἐν γῇ Χαναάν· αὗται δὲ αἱ γενέσεις Ἰακώβ.
\vs{2}Ἰωσὴφ δὲ δέκα καὶ ἑπτὰ ἐτῶν ἦν, ποιμαίνων τὰ πρόβατα τοῦ πατρὸς αὐτοῦ μετὰ τῶν ἀδελφῶν αὐτοῦ, ὢν νέος, μετὰ τῶν υἱῶν Βαλλᾶς, καὶ μετὰ τῶν υἱῶν Ζελφᾶς, τῶν γυναικῶν τοῦ πατρὸς αὐτοῦ· κατήνεγκαν δὲ Ἰωσὴφ ψόγον πονηρὸν πρὸς Ἰσραὴλ τὸν πατέρα αὐτῶν.
\vs{3}Ἰακὼβ δὲ ἠγάπα τὸν Ἰωσὴφ παρὰ πάντας τοὺς υἱοὺς αὐτοῦ, ὅτι υἱὸς γήρως ἦν αὐτῷ· ἐποίησε δὲ αὐτῷ χιτῶνα ποικίλον.
\vs{4}Ἰδόντες δὲ οἱ ἀδελφοὶ αὐτοῦ, ὅτι αὐτὸν ὁ πατὴρ φιλεῖ ἐκ πάντων τῶν υἱῶν αὐτοῦ, ἐμίσησαν αὐτὸν, καὶ οὐκ ἠδύναντο λαλεῖν αὐτῷ οὐδὲν εἰρηνικόν.
\vs{5}Ἐνυπνιασθεὶς δὲ Ἰωσὴφ ἐνύπνιον, ἀπήγγειλεν αὐτὸ τοῖς ἀδελφοῖς αὐτοῦ.
\vs{6}Καὶ εἶπεν αὐτοῖς, ἀκούσατε τοῦ ἐνυπνίου τούτου, οὗ ἐνυπνιάσθην.
\vs{7}Ὤμην ὑμᾶς δεσμεύειν δράγματα ἐν μέσῳ τῷ πεδίῳ· καὶ ἀνέστη τὸ ἐμὸν δράγμα, καὶ ὠρθώθη· περιστραφέντα δὲ τὰ δράγματα ὑμῶν, προσεκύνησαν τὸ ἐμὸν δράγμα.
\vs{8}Εἶπαν δὲ αὐτῷ οἱ ἀδελφοὶ αὐτοῦ, μὴ βασιλεύων βασιλεύσεις ἐφʼ ἡμᾶς, ἢ κυριεύων κυριεύσεις ἡμῶν, καὶ προσέθεντο ἔτι μισεῖν αὐτὸν ἕνεκεν τῶν ἐνυπνίων αὐτοῦ, καὶ ἕνεκεν τῶν ῥημάτων αὐτοῦ.
\vs{9}Εἶδε δὲ ἐνύπνιον ἕτερον, καὶ διηγήσατο αὐτὸ τῷ πατρὶ αὐτοῦ, καὶ τοῖς ἀδελφοῖς αὐτοῦ· καὶ εἶπεν, ἰδοὺ ἐνυπνιασάμην ἐνύπνιον ἕτερον· ὥσπερ ὁ ἥλιος, καὶ ἡ σελήνη, καὶ ἕνδεκα ἀστέρες προσεκύνουν με.
\vs{10}Καὶ ἐπετίμησεν αὐτῷ ὁ πατὴρ αὐτοῦ, καὶ εἶπεν αὐτῷ, τί τὸ ἐνύπνιον τοῦτο, ὃ ἐνυπνιάσθης; ἆρά γε ἐλθόντες ἐλευσόμεθα ἐγώ τε καὶ ἡ μήτηρ σου καὶ οἱ ἀδελφοί σου προσκυνῆσαί σοι ἐπὶ τὴν γῆν;
\vs{11}Ἐζήλωσαν δὲ αὐτὸν οἱ ἀδελφοὶ αὐτοῦ· ὁ δὲ πατὴρ αὐτοῦ διετήρησε τὸ ῥῆμα.
\vs{12}Ἐπορεύθησαν δὲ οἱ ἀδελφοὶ αὐτοῦ βόσκειν τὰ πρόβατα τοῦ πατρὸς αὐτῶν εἰς Συχέμ.
\vs{13}Καὶ εἶπεν Ἰσραὴλ πρὸς Ἰωσὴφ, οὐχὶ οἱ ἀδελφοί σου ποιμαίνουσιν εἰς Συχέμ; δεῦρο ἀποστείλω σε πρὸς αὐτούς· εἶπε δὲ αὐτῷ, ἰδοὺ ἐγώ.
\vs{14}Εἶπε δὲ αὐτῷ Ἰσραὴλ, πορευθεὶς ἴδε, εἰ ὑγιαίνουσιν οἱ ἀδελφοί σου, καὶ τὰ πρόβατα, καὶ ἀνάγγειλόν μοι· καὶ ἀπέστειλεν αὐτὸν ἐκ τῆς κοιλάδος τῆς Χεβρών· καὶ ἦλθεν εἰς Συχέμ.
\vs{15}Καὶ εὗρεν αὐτὸν ἄνθρωπος πλανώμενον ἐν τῷ πεδίῳ· ἠρώτησε δὲ αὐτὸν ὁ ἄνθρωπος, λέγων, τί ζητεῖς;
\vs{16}Ὁ δὲ εἶπε, τοὺς ἀδελφούς μου ζητῶ· ἀπάγγειλόν μοι ποῦ βόσκουσιν.
\vs{17}Εἶπε δὲ αὐτῷ ὁ ἄνθρωπος, ἀπῄρκασιν ἐντεῦθεν· ἤκουσα γὰρ αὐτῶν λεγόντων, πορευθῶμεν εἰς Δωθαείμ· καὶ ἐπορεύθη Ἰωσὴφ κατόπισθε τῶν ἀδελφῶν αὐτοῦ, καὶ εὗρεν αὐτοὺς ἐν Δωθαείμ.

\vs{18}Προεῖδον δὲ αὐτὸν μακρόθεν πρὸ τοῦ ἐγγίσαι αὐτὸν πρὸς αὐτούς· καὶ ἐπονηρεύοντο τοῦ ἀποκτεῖναι αὐτόν.
\vs{19}Εἶπε δὲ ἕκαστος πρὸς τὸν ἀδελφὸν αὐτοῦ, ἰδοὺ ὁ ἐνυπνιαστὴς ἐκεῖνος ἔρχεται.
\vs{20}Νῦν οὖν δεῦτε ἀποκτείνωμεν αὐτὸν, καὶ ῥίψωμεν αὐτὸν εἰς ἕνα τῶν λάκκων· καὶ ἐροῦμεν, θηρίον πονηρὸν κατέφαγεν αὐτόν· καὶ ὀψόμεθα, τί ἔσται τὰ ἐνύπνια αὐτοῦ.
\vs{21}Ἀκούσας δὲ Ῥουβὴν, ἐξείλετο αὐτὸν ἐκ τῶν χειρῶν αὐτῶν· καὶ εἶπεν, οὐ πατάξωμεν αὐτὸν εἰς ψυχήν.
\vs{22}Εἶπε δὲ αὐτοῖς Ῥουβὴν, μὴ ἐκχέητε αἷμα· ἐμβάλλετε αὐτὸν εἰς ἕνα τῶν λάκκων τούτων τῶν ἐν τῇ ἐρήμῳ, χεῖρα δὲ μὴ ἐπενέγκητε αὐτῷ· ὅπως ἐξέληται αὐτὸν ἐκ τῶν χειρῶν αὐτῶν, καὶ ἀποδῷ αὐτὸν τῷ πατρὶ αὐτοῦ.
\vs{23}Ἐγένετο δὲ ἡνίκα ἦλθεν Ἰωσὴφ πρὸς τοὺς ἀδελφοὺς αὐτοῦ, ἐξέδυσαν Ἰωσὴφ τὸν χιτῶνα τὸν ποικίλον τὸν περὶ αὐτόν.
\vs{24}Καὶ λαβόντες αὐτὸν, ἔῤῥιψαν εἰς τὸν λάκκον· ὁ δὲ λάκκος κενὸς, ὕδωρ οὐκ εἶχε.
\vs{25}Ἐκάθισαν δὲ φαγεῖν ἄρτον· καὶ ἀναβλέψαντες τοῖς ὀφθαλμοῖς εἶδον, καὶ ἰδοὺ ὁδοιπόροι Ἰσμαηλῖται ἤρχοντο ἐκ Γαλαάδ· καὶ αἱ κάμηλοι αὐτῶν ἔγεμον θυμιαμάτων καὶ ῥητίνης καὶ στακτῆς. ἐπορεύοντο δὲ καταγαγεῖν εἰς Αἴγυπτον.

\vs{26}Εἶπε δὲ Ἰούδας πρὸς τοὺς ἀδελφοὺς αὐτοῦ, τί χρήσιμον, ἐὰν ἀποκτείνωμεν τὸν ἀδελφὸν ἡμῶν, καὶ κρύψωμεν τὸ αἷμα αὐτοῦ;
\vs{27}Δεῦτε ἀποδώμεθα αὐτὸν τοῖς Ἰσμαηλίταις τούτοις· αἱ δὲ χεῖρες ἡμῶν μὴ ἔστωσαν ἐπʼ αὐτὸν, ὅτι ἀδελφὸς ἡμῶν καὶ σὰρξ ἡμῶν ἐστίν. Ἤκουσαν δὲ οἱ ἀδελφοὶ αὐτοῦ.
\vs{28}Καὶ παρεπορεύοντο οἱ ἄνθρωποι οἱ Μαδιηναῖοι ἔμποροι, καὶ ἐξείλκυσαν καὶ ἀνεβίβασαν τὸν Ἰωσὴφ ἐκ τοῦ λάκκου· καὶ ἀπέδοντο τὸν Ἰωσὴφ τοῖς Ἰσμαηλίταις εἴκοσι χρυσῶν. Καὶ κατήγαγον τὸν Ἰωσὴφ εἰς Αἴγυπτον.
\vs{29}Ἀνέστρεψε δὲ Ῥουβὴν ἐπὶ τὸν λάκκον, καὶ οὐχ ὁρᾷ τὸν Ἰωσὴφ ἐν τῷ λάκκῳ· καὶ διέῤῥηξε τὰ ἱμάτια αὐτοῦ.
\vs{30}Καὶ ἐπέστρεψε πρὸς τοὺς ἀδελφοὺς αὐτοῦ, καὶ εἶπε, τὸ παιδάριον οὐκ ἔστιν· ἐγὼ δὲ ποῦ πορεύομαι ἔτι;
\vs{31}Λαβόντες δὲ τὸν χιτῶνα τοῦ Ἰωσὴφ, ἔσφαξαν ἔριφον αἰγῶν, καὶ ἐμόλυναν τὸν χιτῶνα τῷ αἵματι.
\vs{32}Καὶ ἀπέστειλαν τὸν χιτῶνα τὸν ποικίλον, καὶ εἰσήνεγκαν τῷ πατρὶ αὐτῶν· καὶ εἶπαν, τοῦτον εὕρομεν, ἐπίγνωθι εἰ χιτὼν τοῦ υἱοῦ σου ἐστὶν, ἢ οὔ.
\vs{33}Καὶ ἐπέγνω αὐτὸν, καὶ εἶπε, χιτὼν τοῦ υἱοῦ μου ἐστί· θηρίον πονηρὸν κατέφαγεν αὐτόν· θηρίον ἥρπασε τὸν Ἰωσήφ.
\vs{34}Διέῤῥηξε δὲ Ἰακὼβ τὰ ἱμάτια αὐτοῦ, καὶ ἐπέθετο σάκκον ἐπὶ τὴν ὀσφῦν αὐτοῦ, καὶ ἐπένθει τὸν υἱὸν αὐτοῦ ἡμέρας πολλάς.
\vs{35}Συνήχθησαν δὲ πάντες οἱ υἱοὶ αὐτοῦ καὶ αἱ θυγατέρες, καὶ ἦλθον παρακαλέσαι αὐτόν· καὶ οὐκ ἤθελε παρακαλεῖσθαι, λέγων, ὅτι καταβήσομαι πρὸς τὸν υἱόν μου πενθῶν εἰς ᾅδου· καὶ ἔκλαυσεν αὐτὸν ὁ πατὴρ αὐτοῦ.
\vs{36}Οἱ δὲ Μαδιηναῖοι ἀπέδοντο τὸν Ἰωσὴφ εἰς Αἴγυπτον τῷ Πετεφρῇ τῷ σπάδοντι Φαραὼ ἀρχιμαγείρῳ.

\ch{38}
Ἐγένετο δὲ ἐν τῷ καιρῷ ἐκείνῳ, κατέβη Ἰούδας ἀπὸ τῶν ἀδελφῶν αὐτοῦ, καὶ ἀφίκετο ἕως πρὸς ἄνθρωπον τινὰ Ὀδολλαμίτην, ᾧ ὄνομα Εἰράς.
\vs{2}Καὶ εἶδεν ἐκεῖ Ἰούδας θυγατέρα ἀνθρώπου Χαναναίου, ᾗ ὄνομα Σαυά· καὶ ἔλαβεν αὐτὴν, καὶ εἰσῆλθε πρὸς αὐτήν.
\vs{3}Καὶ συλλαβοῦσα ἔτεκεν υἱὸν, καὶ ἐκάλεσε τὸ ὄνομα αὐτοῦ, Ἤρ.
\vs{4}Καὶ συλλαβοῦσα ἔτεκεν υἱὸν ἔτι, καὶ ἐκάλεσε τὸ ὄνομα αὐτοῦ, Αὐνάν.
\vs{5}Καὶ προσθεῖσα ἔτεκεν υἱὸν, καὶ ἐκάλεσε τὸ ὄνομα αὐτοῦ, Σηλώμ· αὕτη δὲ ἦν ἐν Χασβὶ, ἡνίκα ἔτεκεν αὐτούς.
\vs{6}Καὶ ἔλαβεν Ἰούδας γυναῖκα Ἢρ τῷ πρωτοτόκῳ αὐτοῦ, ᾗ ὄνομα Θάμαρ.
\vs{7}Ἐγένετο δὲ Ἢρ πρωτότοκος Ἰούδα πονηρὸς ἔναντι Κυρίου· καὶ ἀπέκτεινεν αὐτὸν ὁ Θεός.
\vs{8}Εἶπε δὲ Ἰούδας τῷ Αὐνάν· εἴσελθε πρὸς τὴν γυναῖκα τοῦ ἀδελφοῦ σου, καὶ ἐπιγάμβρευσαι αὐτὴν, καὶ ἀνάστησον σπέρμα τῷ ἀδελφῷ σου.
\vs{9}Γνοὺς δὲ Αὐνὰν, ὅτι οὐκ αὐτῷ ἔσται τὸ σπέρμα, ἐγένετο ὅταν εἰσήρχετο πρὸς τὴν γυναῖκα τοῦ ἀδελφοῦ αὐτου, ἐξέχεεν ἐπὶ τὴν γῆν, τοῦ μὴ δοῦναι σπέρμα τῷ ἀδελφῷ αὐτοῦ.
\vs{10}Πονηρὸν δὲ ἐφάνη ἐναντίον τοῦ Θεοῦ, ὅτι ἐποίησε τοῦτο· καὶ ἐθανάτωσε καὶ τοῦτον.

\vs{11}Εἶπε δὲ Ἰούδας Θάμαρ τῇ νύμφῃ αὐτοῦ, κάθου χήρα ἐν τῷ οἴκῳ τοῦ πατρός σου, ἕως μέγας γένηται Σηλὼμ ὁ υἱός μου· εἶπε γάρ, μή ποτε ἀποθάνῃ καὶ οὗτος, ὥσπερ καὶ οἱ ἀδελφοὶ αὐτοῦ. Ἀπελθοῦσα δὲ Θάμαρ ἐκάθητο ἐν τῷ οἴκῳ τοῦ πατρὸς αὐτῆς.
\vs{12}Ἐπληθύνθησαν δὲ αἱ ἡμέραι, καὶ ἀπέθανε Σαυὰ ἡ γυνὴ Ἰούδα· καὶ παρακληθεὶς Ἰούδας ἀνέβη ἐπὶ τοὺς κείροντας τὰ πρόβατα αὐτοῦ, αὐτὸς καὶ Εἰρὰς ὁ ποιμὴν αὐτοῦ ὁ Ὀδολλαμίτης εἰς Θαμνά.
\vs{13}Καὶ ἀπηγγέλε Θάμαρ τῇ νύμφῃ αὐτοῦ, λέγοντες, ἰδοὺ ὁ πενθερός σου ἀναβαίνει εἰς Θαμνὰ, κεῖραι τὰ πρόβατα αὐτοῦ.
\vs{14}Καὶ περιελομένη τὰ ἱμάτια τῆς χηρεύσεως ἀφʼ ἑαυτῆς, περιέβαλε τὸ θέριστρον, καὶ ἐκαλλωπίσατο, καὶ ἐκάθισε πρὸς ταῖς πύλαις Αἰνὰν, ἥ ἐστιν ἐν παρόδῳ Θαμνά· ἴδε γὰρ ὅτι μέγας γέγονε Σηλὼμ, αὐτὸς δὲ οὐκ ἔδωκεν αὐτὴν αὐτῷ γυναῖκα.
\vs{15}Καὶ ἰδὼν αὐτὴν Ἰούδας ἔδοξεν αὐτὴν πόρνην εἶναι· κατεκαλύψατο γὰρ τὸ πρόσωπον αὐτῆς καὶ οὐκ ἐπέγνω αὐτήν.
\vs{16}Ἐξέκλινε δὲ πρὸς αὐτὴν τὴν ὁδόν· καὶ εἶπεν αὐτῇ, ἔασόν με εἰσελθεῖν πρός σε· οὐ γὰρ ἔγνω, ὅτι νύμφη αὐτοῦ ἐστίν· ἡ δὲ εἶπε, τί μοι δώσεις, ἐὰν εἰσέλθῃς πρός με;
\vs{17}Ὁ δὲ εἶπεν, ἐγώ σοι ἀποστελλῶ ἔριφον αἰγῶν ἐκ τῶν προβάτων μον· ἡ δὲ εἶπεν, ἐὰν δῷς μοι ἀῤῥαβῶνα, ἕως τοῦ ἀποστεῖλαί σε.
\vs{18}Ὁ δὲ εἶπε, τίνα τὸν ἀῤῥαβῶνά σοι δώσω; ἡ δὲ εἶπε, τὸν δακτύλιόν σου, καὶ τὸν ὁρμίσκον, καὶ τὴν ῥάβδον τὴν ἐν τῇ χειρίσου. Καὶ ἔδωκεν αὐτῇ, καὶ εἰσῆλθε πρὸς αὐτήν· καὶ ἐν γαστρὶ ἔλαβεν ἐξ αὐτοῦ.
\vs{19}Καὶ ἀναστᾶσα ἀπῆλθε, καὶ περιείλετο τὸ θέριστρον αὐτῆς ἀφʼ ἑαυτῆς, καὶ ἐνεδύσατο τὰ ἱμάτια τῆς χηρεύσεως αὐτῆς.
\vs{20}Ἀπέστειλε δὲ Ἰούδας τὸν ἔριφον ἐξ αἰγῶν ἐν χειρὶ τοῦ ποιμένος αὐτοῦ τοῦ Ὀδολλαμείτου, κομίσασθαι παρὰ τῆς γυναικὸς τὸν ἀῤῥαβῶνα· καὶ οὐχ εὗρεν αὐτήν.
\vs{21}Ἐπηρώτησε δὲ τοὺς ἄνδρας τοὺς ἐκ τοῦ τόπου, ποῦ ἐστιν ἡ πόρνη ἡ γενομένη ἐν Αἰνὰν ἐπὶ τῆς ὁδοῦ; καὶ εἶπαν, οὐκ ἦν ἐνταῦθα πόρνη.
\vs{22}Καὶ ἀπεστράφη πρὸς Ἰούδαν, καὶ εἶπεν, οὐχ εὗρον· καὶ οἱ ἄνθρωποι οἱ ἐκ τοῦ τόπου λέγουσι, μὴ εἶναι ὧδε πόρνην.
\vs{23}Εἶπε δὲ Ἰούδας, ἐχέτω αὐτά· ἀλλὰ μή ποτε καταγελασθῶμεν· ἐγὼ μὲν ἀπέσταλκα τὸν ἔριφον τοῦτον, σὺ δὲ οὐχ εὕρηκας.
\vs{24}Ἐγένετο δὲ μετὰ τρίμηνον ἀνηγγέλη τῷ Ἰούδα, λέγοντες, ἐκπεπόρνευκε Θάμαρ ἡ νύμφη σου, καὶ ἰδοὺ ἐν γαστρὶ ἔχει ἐκ πορνείας· Εἶπε δὲ Ἰούδας, ἐξαγάγετε αὐτὴν, καὶ κατακαυθήτω.
\vs{25}Αὐτὴ δὲ ἀγομένη ἀπέστειλε πρὸς τὸν πενθερὸν αὐτὴς, λέγουσα, ἐκ τοῦ ἀνθρώπου οὕτινος ταῦτά ἐστιν, ἐγὼ ἐν γαστρὶ ἔχω· καὶ εἶπεν, ἐπίγνωθι τίνος ὁ δακτύλιος, καὶ ὁ ὁρμίσκος καὶ ἡ ῥάβδος αὕτη.
\vs{26}Ἐπέγνω δὲ Ἰούδας, καὶ εἶπε, δεδικαίωται Θάμαρ ἢ ἐγώ· οὗ ἕνεκεν οὐκ ἔδωκα αὐτὴν Σηλὼμ τῷ υἱῷ μου· Καὶ οὐ προσέθετο ἔτι τοῦ γνῶναι αὐτήν.
\vs{27}Ἐγένετο δὲ ἡνίκα ἔτικτε, καὶ τῇδε ἦν δίδυμα ἐν τῇ γαστρὶ αὐτῆς.
\vs{28}Ἐγένετο δὲ ἐν τῷ τίκτειν αὐτὴν, ὁ εἷς προεξήνεγκεν τὴν χεῖρα· λαβοῦσα δὲ ἡ μαῖα, ἔδησεν ἐπὶ τὴν χεῖρα αὐτοῦ κόκκινον, λέγουσα, οὗτος ἐξελεύσεται πρότερος.
\vs{29}Ὡς δὲ ἐπισυνήγαγε τὴν χεῖρα, καὶ εὐθὺς ἐξῆλθεν ὁ ἀδελφὸς αὐτοῦ· ἡ δὲ εἶπε, τί διεκόπη διὰ σὲ φραγμός; καὶ ἐκάλεσε τὸ ὄνομα αὐτοῦ, Φαρές.
\vs{30}Καὶ μετὰ τοῦτο ἐξῆλθεν ὁ ἀδελφὸς αὐτοῦ, ἐφʼ ᾧ ἦν ἐπὶ τῇ χειρὶ αὐτοῦ τὸ κόκκινον· καὶ ἐκάλεσε τὸ ὄνομα αὐτοῦ, Ζαρά.

\ch{39}
Ἰωσὴφ δὲ κατήχθη εἰς Αἴγυπτον· καὶ ἐκτήσατο αὐτὸν Πετεφρὴς ὁ εὐνοῦχος Φαραὼ, ὁ ἀρχιμάγειρος, ἀνὴρ Αἰγύπτιος, ἐκ χειρῶν τῶν Ἰσμαηλιτῶν, οἳ κατήγαγον αὐτὸν ἐκεῖ.
\vs{2}Καὶ ἦν Κύριος μετὰ Ἰωσήφ· καὶ ἦν ἀνὴρ ἐπιτυγχάνων· καὶ ἐγένετο ἐν τῷ οἴκῳ παρὰ τῷ κυρίῳ αὐτοῦ τῷ Αἰγυπτίῳ.
\vs{3}Ἤδει δὲ ὁ κύριος αὐτοῦ, ὅτι ὁ Κύριος ἦν μετʼ αὐτοῦ, καὶ ὅσα ἐὰν ποιῇ, Κύριος εὐοδοῖ ἐν ταῖς χερσὶν αὐτοῦ.
\vs{4}Καὶ εὗρεν Ἰωσὴφ χάριν ἐναντίον τοῦ κυρίου αὐτοῦ, καὶ εὐηρέστησεν αὐτῷ. Καὶ κατέστησε αὐτὸν ἐπὶ τοῦ οἴκου αὐτοῦ· καὶ πάντα ὅσα ἦν αὐτῷ, ἔδωκε διὰ χειρὸς Ἰωσήφ.
\vs{5}Ἐγένετο δὲ μετὰ τὸ καταστῆναι αὐτὸν ἐπὶ τοῦ οἴκου αὐτοῦ, καὶ ἐπὶ πάντα ὅσα ἦν αὐτῷ, καὶ ηὐλόγησε Κύριος τὸν οἶκον τοῦ Αἰγυπτίου διὰ Ἰωσήφ· καὶ ἐγενήθη εὐλογία Κυρίου ἐν πᾶσι τοῖς ὑπάρχουσιν αὐτῷ ἐν τῷ οἴκῳ, καὶ ἐν τῷ ἀγρῷ αὐτοῦ.
\vs{6}Καὶ ἐπέτρεψε πάντα ὅσα ἦν αὐτῷ, εἰς χεῖρας Ἰωσήφ· καὶ οὐκ ᾔδει τῶν καθʼ αὑτὸν οὐδὲν, πλὴν τοῦ ἄρτου, οὗ ἤσθιεν αὐτός. Καὶ ἦν Ἰωσὴφ καλὸς τῷ εἴδει, καὶ ὡραῖος τῇ ὄψει σφόδρα.
\vs{7}Καὶ ἐγένετο μετὰ τὰ ῥήματα ταῦτα, καὶ ἐπέβαλεν ἡ γυνὴ τοῦ κυρίου αὐτοῦ τοὺς ὀφθαλμοὺς αὐτῆς ἐπὶ Ἰωσήφ· καὶ εἶπεν, κοιμήθητι μετʼ ἐμοῦ.
\vs{8}Ὁ δὲ οὐκ ἤθελεν· εἶπε δὲ τῇ γυναικὶ τοῦ κυρίου αὐτοῦ, εἰ ὁ κύριός μου οὐ γινώσκει διʼ ἐμὲ οὐδὲν ἐν τῷ οἴκῳ αὐτοῦ, καὶ πάντα ὅσα ἐστὶν αὐτῷ ἔδωκεν εἰς τὰς χεῖράς μου,
\vs{9}καὶ οὐχ ὑπερέχει ἐν τῇ οἰκίᾳ ταύτῆ οὐθὲν ἐμοῦ, οὐδὲ ὑπεξῄρηται ἀπʼ ἐμοῦ οὐδὲν, πλὴν σοῦ, διὰ τὸ σὲ γυναῖκα αὐτοῦ εἶναι, καὶ πῶς ποιήσω τὸ ῥῆμα τὸ πονηρὸν τοῦτο, καὶ ἁμαρτήσομαι ἐναντίον τοῦ Θεοῦ;
\vs{10}Ἡνίκα δὲ ἐλάλει τῷ Ἰωσὴφ ἡμέραν ἐξ ἡμέρας, καὶ οὐχ ὑπήκουεν αὐτῇ καθεύδειν μετʼ αὐτῆς, τοῦ συγγενέσθαι αὐτῇ.
\vs{11}Ἐγένετο δὲ τοιαύτη τις ἡμέρα, καὶ εἰσῆλθεν Ἰωσὴφ εἰς τὴν οἰκίαν ποιεῖν τὰ ἔργα αὐτοῦ, καὶ οὐθεὶς ἦν τῶν ἐν τῇ οἰκίᾳ ἔσω.
\vs{12}Καὶ ἐπεσπάσατο αὐτὸν τῶν ἱματίων, λέγουσα, κοιμήθητι μετʼ ἐμοῦ· καὶ καταλιπὼν τὰ ἱμάτια αὐτοῦ ἐν ταῖς χερσὶν αὐτῆς ἔφυγε, καὶ ἐξῆλθεν ἔξω.
\vs{13}Καὶ ἐγένετο ὡς εἶδεν ὅτι καταλιπὼν τὰ ἱμάτια αὐτοῦ ἐν ταῖς χερσὶν αὐτῆς ἔφυγε, καὶ ἐξῆλθεν ἔξω,
\vs{14}καὶ ἐκάλεσε τοὺς ὄντας ἐν τῇ οἰκίᾳ, καὶ εἶπεν αὐτοῖς, λέγουσα, ἴδετε, εἰσήγαγε ἡμῖν παῖδα Ἐβραῖον, ἐμπαίζειν ἡμῖν· εἰσῆλθε πρός με, λέγων, κοιμήθητι μετʼ ἐμοῦ· καὶ ἐβόησα φωνῇ μεγάλῃ.
\vs{15}Ἐν δὲ τῷ ἀκοῦσαι αὐτὸν, ὅτι ὕψωσα τὴν φωνήν μου καὶ ἐβόησα, καταλιπὼν τὰ ἱμάτια αὐτοῦ παρʼ ἐμοὶ ἔφυγε, καὶ ἐξῆλθεν ἔξω.
\vs{16}Καὶ καταλιμπάνει τὰ ἱμάτια παρʼ ἑαυτῇ, ἕως ἦλθεν ὁ κύριος εἰς τὸν οἶκον αὐτοῦ.
\vs{17}Καὶ ἐλάλησεν αὐτῷ κατὰ τὰ ῥήματα ταῦτα, λέγουσα, εἰσῆλθε πρός με ὁ παῖς ὁ Ἑβραῖος, ὃν εἰσήγαγες πρὸς ἡμᾶς, ἐμπαῖξαί μοι· καὶ εἶπέ μοι, κοιμηθήσομαι μετὰ σοῦ.
\vs{18}Ὡς δὲ ἤκοῦσεν, ὅτι ὕψωσα τὴν φωνήν μου καὶ ἐβόησα, καταλιπὼν τὰ ἱμάτια αὐτοῦ παρʼ ἐμοὶ ἔφυγε, καὶ ἐξῆλθεν ἔξω.
\vs{19}Ἐγένετο δὲ, ὡς ἤκουσεν ὁ κύριος τὰ ῥήματα τῆς γυναικὸς αὐτοῦ, ὅσα ἐλάλησε πρὸς αὐτὸν, λέγουσα, οὕτως ἐποίησέ μοι ὁ παῖς σου, καὶ ἐθυμώθη ὀργῇ.

\vs{20}Καὶ λαβὼν ὁ κύριος Ἰωσὴφ, ἐνέβαλε αὐτὸν εἰς τὸ ὀχύρωμα, εἰς τὸν τόπον ἐν ᾧ οἱ δεσμῶται τοῦ βασιλέως κατέχονται ἐκεῖ ἐν τῷ ὀχυρώματι.
\vs{21}Καὶ ἦν Κύριος μετὰ Ἰωσὴφ, καὶ κατέχεεν αὐτοῦ ἔλεος· καὶ ἔδωκεν αὐτῷ χάριν ἐναντίον τοῦ ἀρχιδεσμοφύλακος.
\vs{22}Καὶ ἔδωκεν ὁ ἀρχιδεσμοφύλαξ τὸ δεσμωτήριον διὰ χειρὸς Ἰωσὴφ, καὶ πάντας τοὺς ἀπηγμένους ὅσοι ἐν τῷ δεσμωτηρίῳ, καὶ πάντα ὅσα ποιοῦσιν ἐκεῖ, αὐτὸς ἦν ποιῶν.
\vs{23}Οὐκ ἦν ὁ ἀρχιδεσμοφύλαξ τοῦ δεσμωτηρίου γινώσκον διʼ αὐτὸν οὐθέν· πάντα γὰρ ἦν διὰ χειρὸς Ἰωσὴφ, διὰ τὸ τὸν Κύριον μετʼ αὐτοῦ εἶναι· καὶ ὅσα αὐτὸς ἐποίει, ὁ Κύριος εὐώδο ἐν ταῖς χερσὶν αὐτοῦ.

\ch{40}
Ἐγένετο δὲ μετὰ τὰ ῥήματα ταῦτα, ἥμαρτεν ὁ ἀρχιοινοχόος τοῦ βασιλέως Αἰγύπτου, καὶ ὁ ἀρχισιτοποιὸς, τῷ κυρίῳ αὐτῶν βασιλεῖ Αἰγύπτου.
\vs{2}Καὶ ὠργίσθη Φαραὼ ἐπὶ τοῖς δυσὶν εὐνούχοις αὐτοῦ, ἐπὶ τῷ ἀρχιοινοχόῳ, καὶ ἐπὶ τῷ ἀρχισιτοποιῷ·
\vs{3}Καὶ ἔθετο αὐτοὺς ἐν φυλακῇ εἰς τὸ δεσμωτήριον, εἰς τὸν τόπον, οὗ Ἰωσὴφ ἀπῆκτο ἐκεῖ.
\vs{4}Καὶ συνέστησεν ὁ ἀρχιδεσμώτης τῷ Ἰωσὴφ αὐτούς· καὶ παρέστη αὐτοῖς· ἦσαν δὲ ἡμέρας ἐν τῇ φυλακῇ.
\vs{5}Καὶ εἶδον ἀμφότεροι ἐνύπνιον ἐν μιᾷ νυκτί· ἡ δὲ ὅρασις τοῦ ἐνυπνίου τοῦ ἀρχιοινοχόου καὶ ἀρχισιτοποιοῦ, οἳ ἦσαν τῷ βασιλεῖ Αἰγύπτου, οἱ ὄντες ἐν τῷ δεσμωτηρίῳ, ἦν αὕτη.
\vs{6}Εἰσῆλθε πρὸς αὐτοὺς Ἰωσὴφ τὸ πρωῒ, καὶ εἶδεν αὐτοὺς, καὶ ἦσαν τεταραγμένοι.
\vs{7}Καὶ ἠρώτα τοὺς εὐνούχους Φαραὼ, οἳ ἦσαν μετʼ αὐτοῦ ἐν τῇ φυλακῇ παρὰ τῷ κυρίῳ αὐτοῦ, λέγων, τί ὅτι τὰ πρόσωπα ὑμῶν σκυθρωπὰ σήμερον;
\vs{8}Οἱ δὲ εἶπαν αὐτῷ, ἐνύπνιον εἴδομεν, καὶ ὁ συγκρίνων οὐκ ἔστιν αὐτό· εἶπε δὲ αὐτοῖς Ἰωσὴφ, οὐχὶ διὰ τοῦ Θεοῦ ἡ διασάφησις αὐτῶν ἐστι; διηγήσασθε οὖν μοὶ.
\vs{9}Καὶ διηγήσατο ὁ ἀρχιοινοχόος τὸ ἐνύπνιον αὐτοῦ τῷ Ἰωσήφ· καὶ εἶπεν, ἐν τῷ ὕπνῳ μου ἦν ἄμπελος ἐναντίον μου.
\vs{10}Ἐν δὲ τῇ ἀμπέλῳ τρεῖς πυθμένες, καὶ αὐτὴ θάλλουσα, ἀνενηνοχυῖα βλαστούς· πέπειροι οἱ βότρυες σταφυλῆς.
\vs{11}Καὶ τὸ ποτήριον Φαραὼ ἐν τῇ χειρί μου· καὶ ἔλαβον τὴν σταφυλὴν, καὶ ἐξέθλιψα αὐτὴν εἰς τὸ ποτήριον, καὶ ἔδωκα τὸ ποτήριον εἰς τὴν χεῖρα Φαραώ.
\vs{12}Καὶ εἶπεν αὐτῷ Ἰωσὴφ, τοῦτο ἡ σύγκρίσις αὐτοῦ· οἱ τρεῖς πυθμένες, τρεῖς ἡμέραι εἰσίν.
\vs{13}Ετι τρεῖς ἡμέραι, καὶ μνησθήσεται Φαραὼ τῆς ἀρχῆς σου, καὶ ἀποκαταστήσει σε ἐπὶ τὴν ἀρχιοινοχοΐαν σου, καὶ δώσεις τὸ ποτήριον Φαραὼ εἰς τὴν χεῖρα αὐτοῦ κατὰ τὴν ἀρχήν σου τὴν προτέραν, ὡς ἦσθα οἰνοχοῶν.
\vs{14}Ἀλλὰ μνήσθητί μου διὰ σεαυτοῦ, ὅταν εὖ γενηταί σοι· καὶ ποιήσεις ἐν ἐμοὶ ἔλεος· καὶ μνησθήσῃ περὶ ἐμοῦ πρὸς Φαραὼ, καὶ ἐξάξεις με ἐκ τοῦ ὀχυρώματος τούτου.
\vs{15}Ὅτι κλοπῇ ἐκλάπην ἐκ γῆς Ἑβραίων, καὶ ὧδε οὐκ ἐποίησα οὐδὲν, ἀλλʼ ἐνέβαλόν με εἰς τὸν λάκκον τοῦτον.
\vs{16}Καὶ εἶδεν ὁ ἀρχισιτοποιὸς ὅτι ὀρθῶς συνέκριεν· καὶ εἶπε τῷ Ἰωσὴφ, κᾀγὼ εἶδον ἐνύπνιον· καὶ ᾤμην τρία κανᾶ χονδριτῶν αἴρειν ἐπὶ τῆς κεφαλῆς μου·
\vs{17}Ἐν δὲ κανῷ τῷ ἐπάνω ἀπὸ πάντων τῶν γενῶν, ὧν Φαραὼ ἐσθίει, ἔργον σιτοποιοῦ, καὶ τὰ πετεινὰ τοῦ οὐρανου κατήσθιεν αὐτὰ ἀπὸ τοῦ κανοῦ τοῦ ἐπάνω τῆς κεφαλῆς μου.
\vs{18}Ἀποκριθεὶς δὲ Ἰωσὴφ εἶπεν αὐτῷ, αὕτη ἡ σύγκρισις αὐτοῦ· τὰ τρία κανᾶ, τρεῖς ἡμέραι εἰσίν·
\vs{19}Ἔτι τριῶν ἡμερῶν, καὶ ἀφελεῖ Φαραὼ τὴν κεφαλήν σου ἀπὸ σου· καὶ κρεμάσει σε ἐπὶ ξύλου, καὶ φάγεται τὰ ὄρνεα τοῦ οὐρανοῦ τὰς σάρκας σου ἀπὸ σοῦ.
\vs{20}Ἐγένετο δὲ ἐν τῇ ἡμέρᾳ τῇ τρίτῃ, ἡμέρα γενέσεως ἦν Φαραὼ, καὶ ἐποίει πότον πᾶσι τοῖς παισὶν αὐτοῦ· καὶ ἐμνήσθη τῆς ἀρχῆς τοῦ οἰνοχόου καὶ τῆς ἀρχῆς τοῦ σιτοποιοῦ ἐν μέσῳ τῶν παίδων αὐτοῦ.
\vs{21}Καὶ ἀποκατέστησε τὸν ἀρχιοινοχόον ἐπὶ τὴν ἀρχὴν αὐτοῦ· καὶ ἔδωκε τὸ ποτήριον εἰς τὴν χεῖρα Φαραώ.
\vs{22}Τὸν δὲ ἀρχισιτοποιὸν ἐκρέμασεν, καθὰ συνέκρινεν αὐτοῖς Ἰωσήφ.
\vs{23}Καὶ οὐκ ἐμνήσθη ὁ ἀρχιοινοχόος τοῦ Ἰωσὴφ, ἀλλαʼ ἐπελάθετο αὐτοῦ.

\ch{41}
Ἐγένετο δὲ μετὰ δύο ἔτη ἡμερῶν, Φαραὼ εἶδεν ἐνύπνιον· ᾤετο ἑστάναι ἐπὶ τοῦ ποταμοῦ.
\vs{2}Καὶ ἰδοὺ ὥσπερ ἐκ τοῦ ποταμοῦ ἀνέβαινον ἐπτὰ βόες, καλαὶ τῷ εἴδει, καὶ ἐκλεκταὶ ταῖς σαρξὶ, καὶ ἐβόσκοντο ἐν τῷ Ἄχει.
\vs{3}Ἄλλαι δὲ ἑπτὰ βόες ἀνέβαινον μετὰ ταύτας ἐκ τοῦ ποταμοῦ, αἰσχραὶ τῷ εἴδει, καὶ λεπταὶ ταῖς σαρξὶ, καὶ ἐνέμοντο παρὰ τὰς βόας ἐπὶ τὸ χεῖλος τοῦ ποταμοῦ.
\vs{4}Καὶ κατέφαγον αἱ ἑπτὰ βόες αἱ αἰσχραὶ καὶ λεπταὶ ταῖς σαρξὶ τὰς ἑπτὰ βόας τὰς καλὰς τῷ εἴδει καὶ τὰς ἐκλεκτὰς ταῖς σαρξί· ἠγέρθη δὲ Φαραώ.
\vs{5}Καὶ ἐνυπνιάσθη τὸ δεύτερον· καὶ ἰδοὺ ἑπτὰ στάχυες ἀνέβαινον ἐν τῷ πυθμένι ἑνὶ ἐκλεκτοὶ καὶ καλοί.
\vs{6}Καὶ ἰδοὺ ἑπτὰ στάχυες λεπτοὶ καὶ ἀνεμόφθοροι ἀνεφύοντο μετʼ αὐτούς.
\vs{7}Καὶ κατέπιον οἱ ἑπτὰ στάχυες οἱ λεπτοὶ καὶ ἀνεμόφθοροι τοὺς ἑπτὰ στάχυας τοὺς ἐκλεκτοὺς καὶ τοὺς πλήρεις· ἠγέρθη δὲ Φαραὼ, καὶ ἦν ἐνύπνιον.
\vs{8}Ἐγένετο δὲ πρωῒ, καὶ ἐταράχθη ἡ ψυχὴ αὐτοῦ, καὶ ἀποστείλας ἐκάλεσε πάντας τοὺς ἐξηγητὰς Αἰγύπτου, καὶ πάντας τοὺς σοφοὺς αὐτῆς· καὶ διηγήσατο αὐτοῖς Φαραὼ τὸ ἐνύπνιον αὐτοῦ, καὶ οὐκ ἦν ὁ ἀπαγγέλλων αὐτὸ τῷ Φαραώ.
\vs{9}Καὶ ἐλάλησεν ὁ ἀρχιοινοχόος πρὸς Φαραὼ, λέγων, τὴν ἁμαρτίαν μου ἀναμιμνήσκω σήμερον.
\vs{10}Φαραὼ ὠργίσθη τοῖς παισὶν αὐτοῦ, καὶ ἔθετο ἡμᾶς ἐν φυλακῇ, ἐν τῷ οἴκῳ τοῦ ἀρχιμαγείρου, ἐμέ τε καὶ τὸν ἀρχισιτοποιόν.
\vs{11}Καὶ εἴδομεν ἐνύπνιον ἀμφότεροι ἐν νυκτὶ μιᾷ ἐγὼ καὶ αὐτὸς, ἕκαστος κατὰ τὸ αὐτοῦ ἐνύπνιον εἴδομεν.
\vs{12}Ἦν δὲ ἐκεῖ μεθʼ ἡμῶν νεανίσκος παῖς Ἑβραῖος τοῦ ἀρχιμαγείρου, καὶ διηγησάμεθα αὐτῷ, καὶ συνέκρινεν ἡμῖν.
\vs{13}Ἐγενήθη δὲ, καθὼς συνέκρινεν ἡμῖν οὕτω καὶ συνέβη, ἐμέ τε ἀποκατασταθῆναι ἐπὶ τὴν ἀρχήν μου, ἐκεῖνον δὲ κρεμασθῆναι.
\vs{14}Ἀποστείλας δὲ Φαραὼ ἐκάλεσε τὸν Ἰωσήφ· καὶ ἐξήγαγον αὐτὸν ἀπὸ τοῦ ὀχυρώματος, καὶ ἐξύρησαν αὐτὸν, καὶ ἤλλαξαν τὴν στολὴν αὐτοῦ· καὶ ἦλθε πρὸς Φαραώ.
\vs{15}Εἶπε δὲ Φαραὼ πρὸς Ἰωσὴφ, ἐνύπνιον ἑώρακα, καὶ ὁ συγκρίνων οὐκ ἔστιν αὐτό· ἐγὼ δὲ ἀκήκοα περὶ σοῦ λεγόντων, ἀκούσαντά σε ἐνύπνια, συγκρῖναι αὐτά.
\vs{16}Ἀποκριθεὶς δὲ Ἰωσὴφ τῷ Φαραὼ εἶπεν, ἄνευ τοῦ Θεοῦ οὐκ ἀποκριθήσεται τὸ σωτήριον Φαραώ.
\vs{17}Ἐλάλησε δὲ Φαραὼ τῷ Ἰωσὴφ, λέγων, ἐν τῷ ὕπνῳ μου ᾤμην ἑστάναι παρὰ τὸ χεῖλος τοῦ ποταμοῦ.
\vs{18}Καὶ ὥσπερ ἐκ τοῦ ποταμοῦ ἀνέβαινον ἑπτὰ βόες καλαὶ τῷ εἴδει καὶ ἐκλεκταὶ ταῖς σαρξὶ, καὶ ἐνέμοντο ἐν τῷ Ἄχει.
\vs{19}Καὶ ἰδοὺ ἑπτὰ βόες ἕτεραι ἀνέβαινον ὀπίσω αὐτῶν ἐκ τοῦ ποταμοῦ, πονηραὶ καὶ αἰσχραὶ τῷ εἴδει, καὶ λεπταὶ ταῖς σαρξὶν, οἵας οὐκ εἶδον τοιαύτας ἐν ὅλῃ γῇ Αἰγύπτου αἰσχροτέρας.
\vs{20}Καὶ κατέφαγον αἱ ἑπτὰ βόες αἱ αἰσχραὶ καὶ λεπταὶ τὰς ἑπτὰ βόας τὰς πρώτας τὰς καλὰς καὶ τὰς ἐκλεκτάς.
\vs{21}Καὶ εἰσῆλθον εἰς τὰς κοιλίας αὐτῶν· καὶ οὑ διάδηλοι ἐγένοντο, ὅτι εἰσῆλθον εἰς τὰς κοιλίας αὐτῶν· καὶ αἱ ὄψεις αὐτῶν αἰσχραὶ, καθὰ καὶ τὴν ἀρχήν· ἐξεγερθεὶς δὲ ἐκοιμήθην.
\vs{22}Καὶ εἶδον πάλιν ἐν τῷ ὕπνῳ μου, καὶ ὥσπερ ἑπτὰ στάχυες ἀνέβαινον ἐν πυθμένι ἑνὶ πλήρεις καὶ καλοί·
\vs{23}Ἄλλοι δὲ ἑπτὰ στάχυες λεπτοὶ καὶ ἀνεμόφθοροι ἀνεφύοντο ἐχόμενοι αὐτῶν.
\vs{24}Καὶ κατέπιον οἱ ἑπτὰ στάχυες οἱ λεπτοὶ καὶ ἀνεμόφθοροι τοὺς ἑπτὰ στάχυας τοὺς καλοὺς καὶ τοὺς πλήρεις· εἶπα οὖν τοῖς ἐξῆγηταῖς, καὶ οὐκ ἦν ὁ ἀπαγγέλλων μοι αὐτό.

\vs{25}Καὶ εἶπεν Ἰωσὴφ τῷ Φαραὼ, τὸ ἐνύπνιον Φαραὼ ἕν ἐστιν· ὅσα ὁ Θεὸς ποιεῖ, ἔδειξε τῷ Φαραώ.
\vs{26}Αἱ ἑπτὰ βόες αἱ καλαὶ, ἑπτὰ ἔτη ἐστί· καὶ οἱ ἑπτὰ στάχυες οἱ καλοὶ, ἑπτὰ ἔτη ἐστί· τὸ ἐνύπνιον Φαραὼ ἕν ἐστι.
\vs{27}Καὶ αἱ ἑπτὰ βόες αἱ λεπταὶ, αἱ ἀναβαίνουσαι ὀπίσω αὐτῶν, ἑπτὰ ἔτη ἐστί· καὶ οἱ ἑπτὰ στάχυες οἱ λεπτοὶ καὶ ἀνεμόφθοροι, ἑπτὰ ἔτη ἐστί· ἔσονται ἑπτὰ ἔτη λιμοῦ.
\vs{28}Τὸ δὲ ῥῆμα ὃ εἴρηκα Φαραὼ, ὅσα ὁ Θεὸς ποιεῖ, ἔδειξε τῷ Φαραώ.
\vs{29}Ἰδοὺ ἑπτὰ ἔτη ἔρχεται εὐθηνία πολλὴ ἐν πάσῃ γῇ Αἰγύπτου.
\vs{30}Ἥξει δὲ ἑπτὰ ἔτη λιμοῦ μετὰ ταῦτα· καὶ ἐπιλήσονται τῆς πλησμονῆς τῆς ἐσομένης ἐν ὅλῃ Αἰγύπτῳ· καὶ ἀναλώσει ὁ λιμὸς τῆν γῆν.
\vs{31}Καὶ οὐκ ἐπιγνωσθήσεται ἡ εὐθηνία ἐπὶ τῆς γῆς ἀπὸ τοῦ λιμοῦ τοῦ ἐσομένου μετὰ ταῦτα· ἰσχυρὸς γὰρ ἔσται σφόδρα.
\vs{32}Περὶ δὲ τοῦ δευτερῶσαι τὸ ἐνύπνιον Φαραὼ δὶς, ὅτι ἀληθὲς ἔσται τὸ ῥῆμα τὸ παρὰ τοῦ Θεοῦ· καὶ ταχυνεῖ ὁ Θεὸς τοῦ ποιῆσαι αὐτό.
\vs{33}Νῦν οὖν σκέψαι ἄνθρωπον φρόνιμον καὶ συνετὸν, καὶ κατάστησον αὐτὸν ἐπὶ γῆς Αἰγύπτου.
\vs{34}Καὶ ποιησάτω Φαραὼ καὶ καταστησάτω τοπάρχας ἐπὶ τῆς γῆς· καὶ ἀποπεμπτωσάτωσαν πάντα τὰ γεννήματα τῆς γῆς Αἰγύπτου τῶν ἑπτὰ ἐτῶν τῆς εὐθηνίας,
\vs{35}καὶ συναγαγέτωσαν πάντα τὰ βρώματα τῶν ἑπτὰ ἐτῶν τῶν ἐρχομένων τῶν καλῶν τούτων· καὶ συναχθήτω ὁ σῖτος ὑπὸ χεῖρα Φαραώ· βρώματα ἐν ταῖς πόλεσι φυλαχθήτω.
\vs{36}Καὶ ἔσται τὰ βρώματα τὰ πεφυλαγμένα τῇ γῇ εἰς τὰ ἑπτὰ ἔτη τοῦ λιμοῦ, ἃ ἔσονται ἐν γῇ Αἰγύπτου, καὶ οὐκ ἐκτριβήσεται ἡ γῆ ἐν τῷ λιμῷ.
\vs{37}Ἤρεσε δὲ τὸ ῥῆμα ἐναντίον Φαραὼ, καὶ ἐναντίον πάντων τῶν παίδων αὐτοῦ.

\vs{38}Καὶ εἶπε Φαραὼ πᾶσι τοῖς παισὶν αὐτοῦ, μῆ εὑρήσομεν ἄνθρωπον τοιοῦτον, ὃς ἔχει πνεῦμα Θεοῦ ἐν αὐτῷ;
\vs{39}Εἶπε δὲ Φαραὼ τῷ Ἰωσὴφ, ἐπειδὴ ἔδειξεν ὁ Θεός σοι πάντα ταῦτα, οὐκ ἔστιν ἄνθρωπος φρονιμώτερος καὶ συνετώτερός σου.
\vs{40}Σὺ ἔσῃ ἐπὶ τῷ οἴκῳ μου, καὶ ἐπὶ τῷ στόματί σου ὑπακούσεται πᾶς ὁ λαός μου· πλὴν τὸν θρόνον ὑπερέξω σου ἐγώ.
\vs{41}Εἶπε δὲ Φαραὼ τῷ Ἰωσὴφ, ἰδοὺ καθίστημί σε σήμερον ἐπὶ πάσῃ γῇ Αἰγύπτου.
\vs{42}Καὶ περιελόμενος Φαραὼ τὸν δακτύλιον ἀπὸ τῆς χειρὸς αὐτοῦ, περίεθηκεν αὐτὸν ἐπὶ τὴν χεῖρα Ἰωσὴφ, καὶ ἐνέδυσεν αὐτὸν στολὴν βυσσίνην, καὶ περιέθηκε κλοιὸν χρυσοῦν περὶ τὸν τράχηλον αὐτοῦ.
\vs{43}Καὶ ἀνεβίβασεν αὐτὸν ἐπὶ τὸ ἅρμα τὸ δεύτερον τῶν αὐτοῦ· καὶ ἐκήρυξεν ἔμπροσθεν αὐτοῦ κήρυξ· καὶ κατέστησεν αὐτὸν ἐφʼ ὅλης γῆς Αἰγύπτου.
\vs{44}Εἶπε δὲ Φαραὼ τῷ Ἰωσὴφ, ἐγὼ Φαραώ· ἄνευ σοῦ οὐκ ἐξαρεῖ οὐθεὶς τὴν χεῖρα αὐτοῦ ἐπὶ πάσης γῆς Αἰγύπτου.
\vs{45}Καὶ ἐκάλεσε Φαραὼ τὸ ὄνομα Ἰωσὴφ, Ψονθομφανήχ· καὶ ἔδωκεν αὐτῷ τὴν Ἀσενὲθ θυγατέρα Πετεφρῆ ἱερέως Ἡλιουπόλεως αὐτῷ εἰς γυναῖκα.
\vs{46}Ἰωσὴφ δὲ ἦν ἐτῶν τριάκοντα, ὅτε ἔστη ἐναντίον Φαραὼ βασιλέως Αἰγύπτου· ἐξῆλθε δὲ Ἰωσὴφ ἀπὸ προσώπου Φαραὼ, καὶ διῆλθε πᾶσαν γῆν Αἰγύπτου.
\vs{47}Καὶ ἐποίησεν ἡ γῆ ἐν τοῖς ἑπτὰ ἔτεσι τῆς εὐθηνίας δράγματα.
\vs{48}Καὶ συνήγαγε πάντα τὰ βρώματα τῶν ἑπτὰ ἐτῶν, ἐν οἷς ἦν ἡ εὐθηνία ἐν τῇ γῇ Αἰγύπτου· καὶ ἔθηκε τὰ βρώματα ἐν ταῖς πόλεσι· βρώματα τῶν πεδίων τῆς πόλεως τῶν κύκλῳ αὐτῆς ἔθηκεν ἐν αὐτῇ.
\vs{49}Καὶ συνήγαγεν Ἰωσὴφ σῖτον ὡσεὶ τὴν ἄμμον τῆς θαλάσσης πολὺν σφόδρα, ἕως οὐκ ἠδύνατο ἀριθμηθῆναι, οὐ γὰρ ἦν ἀριθμός.

\vs{50}Τῷ δὲ Ἰωσὴφ ἐγένοντο υἱοὶ δύο πρὸ τοῦ ἐλθεῖν τὰ ἑπτὰ ἔτη τοῦ λιμοῦ, οὓς ἔτεκεν αὐτῷ Ἀσενὲθ ἡ θυγάτηρ Πετεφρῆ ἱερέως Ἡλιουπόλεως.
\vs{51}Ἐκάλεσε δὲ Ἰωσὴφ τὸ ὄνομα τοῦ πρωτοτόκου, Μανασσῆ· ὅτι ἐπιλαθέσθαι με ἐποίησεν ὁ Θεὸς πάντων τῶν πόνων μου, καὶ πάντων τῶν τοῦ πατρός μου·
\vs{52}Τὸ δὲ ὄνομα τοῦ δευτέρου ἐκάλεσεν, Ἐφραίμ· ὅτι ηὔξησέ με ὁ Θεὸς ἐν γῇ ταπεινώσεώς μου.
\vs{53}Παρῆλθον δὲ τὰ ἑπτὰ ἔτη τῆς εὐθηνίας, ἃ ἐγένοντο ἐν τῇ γῇ Αἰγύπτου.
\vs{54}Καὶ ἤρξατο τὰ ἑπτὰ ἔτη τοῦ λιμοῦ ἔρχεσθαι, καθὰ εἶπεν Ἰωσήφ· καὶ ἐγένετο λιμὸς ἐν πάσῃ τῇ γῇ· ἐν δὲ πάσῃ τῇ γῇ Αἰγύπτου ἦσαν ἄρτοι.
\vs{55}Καὶ ἐπείνασε πᾶσα ἡ γῆ Αἰγύπτου· ἔκραξε δὲ ὁ λαὸς πρὸς Φαραὼ περὶ ἄρτων· εἶπε δὲ Φαραὼ πᾶσι τοῖς Αἰγυπτίοις, πορεύεσθε πρὸς Ἰωσὴφ, καὶ ὃ ἐὰν εἴπῃ ὑμῖν, ποιήσατε.
\vs{56}Καὶ ὁ λιμὸς ἦν ἐπὶ προσώπου πάσης τῆς γῆς· ἀνέῳξε δὲ Ἰωσὴφ πάντας τοὺς σιτοβολῶνας, καὶ ἐπώλει πᾶσι τοῖς Αἰγυπτίοις.
\vs{57}Καὶ πᾶσαι αἱ χῶραι ἦλθον εἰς Αἴγυπτον, ἀγοράζειν πρὸς Ἰωσήφ· ἐπεκράτησε γὰρ ὁ λιμὸς ἐν πάσῃ τῇ γῇ·

\ch{42}
Ἰδὼν δὲ Ἰακὼβ ὅτι ἐστὶ πράσις ἐν Αἰγύπτῳ, εἶπε τοῖς υἱοῖς αὐτοῦ, ἱνατί ῥαθυμεῖτε;
\vs{2}Ἰδοὺ ἀκήκοα, ὅτι ἐστὶ σῖτος ἐν Αἰγύπτῳ· κατάβητε ἐκεὶ, καὶ πρίασθε ἡμῖν μικρὰ βρώματα, ἵνα ζήσωμεν καὶ μὴ ἀποθάνωμεν.

\vs{3}Κατέβησαν δὲ οἱ ἀδελφοὶ Ἰωσὴφ οἱ δέκα, πρίασθαι σῖτον ἐξ Αἰγύπτου·
\vs{4}Τὸν δὲ Βενιαμὶν, τὸν ἀδελφὸν Ἰωσὴφ, οὐκ ἀπέστειλε μετὰ τῶν ἀδελφῶν αὐτοῦ· εἶπε γὰρ, μή ποτε συμβῇ αὐτῷ μαλακία.
\vs{5}Ἦλθον δὲ οἱ υἱοὶ Ἰσραὴλ ἀγοράζειν μετὰ τῶν ἐρχομένων· ἦν γὰρ ὁ λιμὸς ἐν γῇ Χαναάν.
\vs{6}Ἰωσὴφ δὲ ἦν ὁ ἄρχων τῆς γῆς· οὗτος ἐπώλει παντὶ τῷ λαῷ τῆς γῆς· ἐλθόντες δὲ οἱ ἀδελφοὶ Ἰωσὴφ προσεκύνησαν αὐτῷ ἐπὶ πρόσωπον ἐπὶ τὴν γῆν.
\vs{7}Ἰδὼν δὲ Ἰωσὴφ τοὺς ἀδελφοὺς αὐτοῦ, ἐπέγνω· καὶ ἠλλοτριοῦτο ἀπʼ αὐτῶν, καὶ ἐλάλησεν αὐτοῖς σκληρά· καὶ εἶπεν αὐτοῖς, πόθεν ἥκατε; οἱ δὲ εἶπον, ἐκ γῆς Χαναὰν, ἀγοράσαι βρώματα.
\vs{8}Ἐπέγνω δὲ Ἰωσὴφ τοὺς ἀδελφοὺς αὐτοῦ· αὐτοὶ δὲ οὐκ ἐπέγνωσαν αὐτόν·
\vs{9}Καὶ ἐμνήσθη Ἰωσὴφ τῶν ἐνυπνίων αὐτοῦ, ὧν εἶδεν αὐτός· καὶ εἶπεν αὐτοῖς, κατάσκοποί ἐστε, κατανοῆσαι τὰ ἴχνη τῆς χώρας ἥκατε.
\vs{10}Οἱ δὲ εἶπαν, οὐχὶ, κύριε· οἱ παῖδές σου ἤλθομεν πρίασθαι βρώματα.
\vs{11}Πάντες ἐσμὲν υἱοὶ ἑνὸς ἀνθρώπου· εἰρηνικοί ἐσμεν, οὐκ εἰσιν οἱ παῖδές σου κατάσκοποι.
\vs{12}Εἶπε δὲ αὐτοῖς, οὐχί· ἀλλὰ τὰ ἴχνη τῆς γῆς ἤλθετε ἰδεῖν.
\vs{13}Οἱ δὲ εἶπαν, δώδεκά ἐσμεν οἱ παῖδές σου ἀδελφοὶ ἐν γῇ Χαναάν· καὶ ἰδοὺ ὁ νεώτερος μετὰ τοῦ πατρὸς ἡμῶν σήμερον· ὁ δὲ ἕτερος οὐχ ὑπάρχει.
\vs{14}Εἶπε δὲ αὐτοῖς Ἰωσὴφ, τοῦτό ἐστιν ὃ εἴρηκα ὑμῖν, λέγων, ὅτι κατάσκοποί ἐστε.
\vs{15}Ἐν τούτῳ φανεῖσθε· νὴ τὴν ὑγίειαν Φαραὼ, οὐ μὴ ἐξέλθητε ἐντεῦθεν, ἐὰν μὴ ὁ ἀδελφὸς ὑμῶν ὁ νεώτερος ἔλθῃ ὧδε.
\vs{16}Ἀποστείλατε ἐξ ὑμῶν ἕνα, καὶ λάβετε τὸν ἀδελφὸν ὑμῶν· ὑμεῖς δὲ ἀπάχθητε ἕως τοῦ φανερὰ γενέσθαι τὰ ῥήματα ὑμῶν, εἰ ἀληθεύετε ἢ οὔ· εἰ δὲ μὴ, νὴ τὴν ὑγίειαν Φαραὼ, ἦ μὴν κατάσκοποί ἐστε.
\vs{17}Καὶ ἔθετο αὐτοὺς ἐν φυλακῇ ἡμέρας τρεῖς.
\vs{18}Εἶπε δὲ αὐτοῖς τῇ ἡμέρᾳ τῇ τρίτῃ, τοῦτο ποιήσατε, καὶ ζήσεσθε· τὸν Θεὸν γὰρ ἐγὼ φοβοῦμαι.
\vs{19}Εἰ εἰρηνικοί ἐστε, ἀδελφὸς ὑμῶν κατασχεθήτω εἷς ἐν τῇ φυλακῇ· αὐτοὶ δὲ βαδίσατε, καὶ ἀπαγάγετε τὸν ἀγορασμὸν τῆς σιτοδοσίας ὑμῶν.
\vs{20}Καὶ τὸν ἀδελφὸν ὑμῶν τὸν νεώτερον ἀγάγετε πρός με, καὶ πιστευθήσονται τὰ ῥήματα ὑμῶν· εἰ δὲ μὴ, ἀποθανεῖσθε. Ἐποίησαν δὲ οὕτως.
\vs{21}Καὶ εἶπεν ἕκαστος πρὸς τὸν ἀδελφὸν αὐτοῦ, ναὶ, ἐν ἁμαρτίαις γάρ ἐσμεν περὶ τοῦ ἀδελφοῦ ἡμῶν, ὅτι ὑπερίδομεν τὴν θλίψιν τῆς ψυχῆς αὐτοῦ, ὅτε κατεδέετο ἡμῶν, καὶ οὐκ εἰσηκούσαμεν αὐτοῦ· καὶ ἕνεκεν τούτου ἐπῆλθεν ἐφʼ ἡμᾶς ἡ θλίψις αὕτη.
\vs{22}Ἀποκριθεὶς δὲ Ῥουβὴν εἶπεν αὐτοῖς, οὐκ ἐλάλησα ὑμῖν, λέγων, μὴ ἀδικήσητε τὸ παιδάριον, καὶ οὐκ εἰσηκούσατέ μου; καὶ ἰδοὺ τὸ αἷμα αὐτοῦ ἐκζητεῖται.
\vs{23}Αὐτοὶ δὲ οὐκ ᾔδεισαν, ὅτι ἀκούει Ἰωσήφ· ὁ γὰρ ἑρμηνευτὴς ἀνὰ μέσον αὐτῶν ἦν·
\vs{24}Ἀποστραφεὶς δὲ ἀπʼ αὐτῶν ἔκλαυσεν Ἰωσήφ· καὶ πάλιν προσῆλθε πρὸς αὐτοὺς, καὶ εἶπεν αὐτοῖς· καὶ ἔλαβε τὸν Συμεὼν ἀπʼ αὐτῶν, καὶ ἔδησεν αὐτὸν ἐναντίον αὐτῶν.

\vs{25}Ἐνετείλατο δὲ Ἰωσὴφ ἐμπλῆσαι τὰ ἀγγεῖα αὐτῶν σίτου, καὶ ἀποδοῦναι τὸ ἀργύριον αὐτῶν ἑκάστῳ εἰς τὸν σάκκον αὐτοῦ, καὶ δοῦναι αὐτοῖς ἐπισιτισμὸν εἰς τὴν ὁδόν· καὶ ἐγενήθη αὐτοῖς οὕτως.
\vs{26}Καὶ ἐπιθέντες τὸν σῖτον ἐπὶ τοῦς ὄνους αὐτῶν, ἀπῆλθον ἐκεῖθεν.
\vs{27}Λύσας δὲ εἷς τὸν μάρσιππον αὐτοῦ, δοῦναι χορτάσματα τοῖς ὄνοις αὐτοῦ, οὗ κατέλυσαν, καὶ εἶδε τὸν δεσμὸν τοῦ ἀργυρίου αὐτοῦ, καὶ ἦν ἐπάνω τοῦ στόματος τοῦ μαρσίππου.
\vs{28}Καὶ εἶπε τοῖς ἀδελφοῖς αὐτοῦ, ἀπεδόθη μοι τὸ ἀργύριον, καὶ ἰδοὺ τοῦτο ἐν τῷ μαρσίππῳ μου· καὶ ἐξέστη ἡ καρδία αὐτῶν, καὶ ἐταράχθησαν πρὸς ἀλλήλους, λέγοντες, τί τοῦτο ἐποίησεν ὁ Θεὸς ἡμῖν;
\vs{29}Ἦλθον δὲ πρὸς Ἰακὼβ τὸν πατέρα αὐτῶν εἰς γὴν Χαναὰν, καὶ ἀπήγγειλαν αὐτῷ πάντα τὰ συμβάντα αὐτοῖς, λέγοντες,
\vs{30}Λελάληκεν ὁ ἄνθρωπος ὁ κύριος τῆς γῆς πρὸς ἡμᾶς σκληρὰ, καὶ ἔθετο ἡμᾶς ἐν φυλακῇ, ὡς κατασκοπεύοντας τὴν γῆν.
\vs{31}Εἴπαμεν δὲ αὐτῷ, εἰρήνικοί ἐσμεν, οὐκ ἐσμὲν κατάσκοποι.
\vs{32}Δώδεκα ἀδελφοί ἐσμεν, υἱοὶ τοῦ πατρὸς ἡμῶν· ὁ εἷς οὐχ ὑπάρχει· ὁ δὲ μικρὸς μετὰ τοῦ πατρὸς ἡμῶν σήμερον ἐν γῇ Χαναάν.
\vs{33}Εἶπε δὲ ἡμῖν ὁ ἄνθρωπος ὁ κύριος τῆς γῆς, ἐν τούτῳ γνώσομαι, ὅτι εἰρηνικοί ἐστε· ἀδελφὸν ἕνα ἄφετε ὧδε μετʼ ἐμοῦ· τὸν δὲ ἀγορασμὸν τῆς σιτοδοσίας τοῦ οἴκου ὑμῶν λαβόντες ἀπέλθατε.
\vs{34}Καὶ ἀγάγετε πρός με τὸν ἀδελφὸν ὑμῶν τὸν νεώτερον· καὶ γνώσομαι ὅτι οὐ κατάσκοποί ἐστε, ἀλλʼ ὅτι εἰρηνικοί ἐστε· καὶ τὸν ἀδελφὸν ὑμῶν ἀποδώσω ὑμῖν, καὶ τῇ γῇ ἐμπορεύσεσθε.
\vs{35}Ἐγένετο δὲ ἐν τῷ κατακενοῦν αὐτοὺς τοὺς σάκκους αὐτῶν, καὶ ἦν ἑκάστου ὁ δεσμὸς τοῦ ἀργυρίου ἐν τῷ σάκκῳ αὐτῶν· καὶ εἶδον τοὺς δεσμοὺς τοῦ ἀργυρίου αὐτῶν αὐτοὶ, καὶ ὁ πατὴρ αὐτῶν, καὶ ἐφοβήθησαν.
\vs{36}Εἶπε δὲ αὐτοῖς Ἰακὼβ ὁ πατὴρ αὐτῶν, ἐμὲ ἠτεκνώσατε· Ἰωσὴφ οὐκ ἔστι, Συμεὼν οὐκ ἔστι, καὶ τὸν Βενιαμὶν λήψεσθε; ἐπʼ ἐμὲ ἐγένετο ταῦτα πάντα.
\vs{37}Εἶπε δὲ Ῥουβὴν τῷ πατρὶ αὐτῶν, λέγων, τοὺς δύο υἱούς μου ἀπόκτεινον, ἐὰν μὴ ἀγάγω αὐτὸν πρὸς σέ· δὸς αὐτὸν εἰς τὴν χεῖρά μου, κᾀγὼ ἀνάξω αὐτὸν πρὸς σέ.
\vs{38}Ὁ δὲ εἶπεν, οὐ καταβήσεται ὁ υἱός μου μεθʼ ὑμῶν, ὅτι ὁ ἀδελφὸς αὐτοῦ ἀπέθανε, καὶ αὐτὸς μόνος καταλέλειπται· καὶ συμβήσεται αὐτὸν μαλακισθῆναι ἐν τῇ ὁδῷ, ᾗ ἐὰν πορεύησθε, καὶ κατάξετέ μου τὸ γῆρας μετὰ λύπης εἰς ᾅδοῦ.

\ch{43}
Ὁ δὲ λιμὸς ἐνίσχυσεν ἐπὶ τῆς γῆς.
\vs{2}Ἐγένετο δὲ ἡνίκα συνετέλεσαν καταφαγεῖν τὸν σῖτον, ὃν ἤνεγκαν ἐξ Αἰγύπτου, καὶ εἶπεν αὐτοῖς ὁ πατὴρ αὐτῶν, πάλιν πορευθέντες πρίασθε ἡμῖν μικρὰ βρώματα.
\vs{3}Εἶπε δὲ αὐτῷ Ἰούδας, λέγων, διαμαρτυρίᾳ μεμαρτύρηται ἡμῖν ὁ ἄνθρωπος ὁ κύριος τῆς γῆς, λέγων, οὐκ ὄψεσθε τὸ πρόσωπόν μου, ἐὰν μὴ ὁ ἀδελφὸς ὑμῶν ὁ νεώτερος μεθʼ ὑμῶν ᾖ.
\vs{4}Εἰ μὲν οὖν ἀποστέλλῃς τὸν ἀδελφὸν ἡμῶν μεθʼ ἡμῶν, καταβησόμεθα, καὶ ἀγοράσομέν σοι βρώματα.
\vs{5}Εἰ δὲ μὴ ἀποστέλλῃς τὸν ἀδελφὸν ἡμῶν μεθʼ ἡμῶν, οὐ πορευσόμεθα· ὁ γὰρ ἄνθρωπος εἶπεν ἡμῖν, λέγων, οὐκ ὄψεσθέ μου τὸ πρόσωπον, ἐὰν μὴ ὁ ἀδελφὸς ὑμῶν ὁ νεώτερος μεθʼ ὑμῶν ᾖ.
\vs{6}Εἶπε δὲ Ἰσραὴλ, τί ἐκακοποιήσατέ με, ἀναγγείλαντες τῷ ἀνθρώπῳ ὅτι ἐστὶν ὑμῖν ἀδελφός;
\vs{7}Οἱ δὲ εἶπαν, ἐρωτῶν ἐπηρώτησεν ἡμᾶς ὁ ἄνθρωπος καὶ τὴν γενεὰν ἡμῶν, λέγων, εἰ ἔτι ὁ πατὴρ ὑμῶν ζῇ, καὶ εἰ ἔστιν ὑμῖν ἀδελφός· καὶ ἀπηγγείλαμεν αὐτῷ κατὰ τὴν ἐπερώτησιν ταύτην· μὴ ᾔδειμεν ὅτι ἐρεῖ ἡμῖν, ἀγάγετε τὸν ἀδελφὸν ὑμῶν;
\vs{8}Εἶπε δὲ Ἰούδας πρὸς Ἰσραὴλ τὸν πατέρα αὐτοῦ, ἀπόστειλον τὸ παιδάριον μετʼ ἐμοῦ· καὶ ἀναστάντες πορευσόμεθα, ἵνα ζῶμεν καὶ μὴ ἀποθάνωμεν καὶ ἡμεῖς, καὶ σὺ, καὶ ἡ ἀποσκευὴ ἡμῶν.
\vs{9}Ἐγὼ δὲ ἐκδέχομαι αὐτόν· ἐκ χειρός μου ζήτησον αὐτόν· ἐὰν μὴ ἀγάγω αὐτὸν πρός σε, καὶ στήσω αὐτὸν ἐναντίον σου, ἡμαρτηκὼς ἔσομαι εἰς σὲ πάσας τὰς ἡμέραν.
\vs{10}Εἰ μὴ γὰρ ἐβραδύναμεν, ἤδη ἂν ὑπεστρέψαμεν δίς.
\vs{11}Εἶπε δὲ αὐτοῖς Ἰσραὴλ ὁ πατὴρ αὐτῶν, εἰ οὕτως ἐστὶ, τοῦτο ποιήσατε· λάβετε ἀπὸ τῶν καρπῶν τῆς γῆς ἐν τοῖς ἀγγείοις ὑμῶν, καὶ καταγάγετε τῷ ἀνθρώπῳ δῶρα τῆς ῥητίνης, καὶ τοῦ μέλιτος, θυμίαμά τε καὶ στακτὴν, καὶ τερέινθον, καὶ κάρυα.
\vs{12}Καὶ τὸ ἀργύριον δισσὸν λάβετε ἐν ταῖς χερσὶν ὑμῶν· τὸ ἀργύριον τὸ ἀποστραφὲν ἐν τοῖς μαρσίπποις ὑμῶν ἀποστρέψατε μεθʼ ὑμῶν· μή ποτε ἀγνόημά ἐστι.
\vs{13}Καὶ τὸν ἀδελφὸν ὑμῶν λάβετε· καὶ ἀναστάντες κατάβητε πρὸς τὸν ἄνθρωπον.
\vs{14}Ὁ δὲ Θεός μου δώῃ ὑμῖν χάριν ἐναντίον τοῦ ἀνθρώπου καὶ ἀποστείλαι τὸν ἀδελφὸν ὑμῶν τὸν ἕνα, καὶ τὸν Βενιαμίν· ἐγὼ μὲν γὰρ καθάπερ ἠτέκνωμαι, ἠτέκνωμαι.

\vs{15}Λαβόντες δὲ οἱ ἄνδρες τὰ δῶρα ταῦτα καὶ τὸ ἀργύριον διπλοῦν, ἔλαβον ἐν ταῖς χερσὶν αὐτῶν καὶ τὸν Βενιαμείν· καὶ ἀναστάντες κατέβησαν εἰς Αἴγυπτον· καὶ ἔστησαν ἐναντίον Ἰωσήφ.
\vs{16}Εἶδε δὲ Ἰωσὴφ αὐτοὺς, καὶ τὸν Βενιαμὶν τὸν ἀδελφὸν αὐτοῦ τὸν ὁμομήτριον· καὶ εἶπε τῷ ἐπὶ τῆς οἰκίας αὐτοῦ, εἰσάγαγε τοὺς ἀνθρώπους εἰς τὴν οἰκίαν, καὶ σφάξον θύματα, καὶ ἑτοίμασον· μετʼ ἐμοῦ γὰρ φάγονται οἱ ἄνθρωποι ἄρτους τὴν μεσημβρίαν.
\vs{17}Ἐποίησε δὲ ὁ ἄνθρωπος καθὰ εἶπεν Ἰωσήφ· καὶ εἰσήγαγε τοὺς ἀνθρώπους εἰς τὸν οἶκον Ἰωσήφ.
\vs{18}Ἰδόντες δὲ οἱ ἄνδρες ὅτι εἰσήχθησαν εἰς τὸν οἶκον τοῦ Ἰωσήφ, εἶπαν, διὰ τὸ ἀργύριον τὸ ἀποστραφὲν ἐν τοῖς μαρσίπποις ἡμῶν τὴν ἀρχὴν, ἡμεῖς εἰσαγόμεθα, τοῦ συκοφαντῆσαι ἡμᾶς καὶ ἐπιθέσθαι ἡμῖν, τοῦ λαβεῖν ἡμᾶς εἰς παῖδας, καὶ τοὺς ὄνους ἡμῶν.
\vs{19}Προσελθόντες δὲ πρὸς τὸν ἄνθρωπον τὸν ἐπὶ τοῦ οἴκου τοῦ Ἰωσὴφ, ἐλάλησαν αὐτῷ ἐν τῷ πυλῶνι τοῦ οἴκου,
\vs{20}λέγοντες, δεόμεθα, κύριε· κατέβηεν τὴν ἀρχὴν πρίασθαι βρώματα.
\vs{21}Ἐγένετο δὲ ἡνίκα ἤλθομεν εἰς τὸ καταλῦσαι, καὶ ἠνοίξαμεν τοὺς μαρσίππους ἡμῶν, καὶ τόδε τὸ ἀργύριον ἑκάστου ἐν τῷ μαρσίππῳ αὐτοῦ· τὸ ἀργύριον ἡμῶν ἐν σταθμῷ ἀπεστρέψαμεν νῦν ἐν ταῖς χερσὶν ἡμῶν.
\vs{22}Καὶ ἀργύριον ἕτερον ἠνέγκαμεν μεθʼ ἑαυτῶν, ἀγοράσαι βρώματα· οὐκ οἴδαμεν τίς ἐνέβαλεν τὸ ἀργύριον εἰς τοὺς μαρσίππους ἡμῶν.
\vs{23}Εἶπε δὲ αὐτοῖς, ἵλεως ὑμῖν, μὴ φοβεῖσθε· ὁ Θεὸς ὑμῶν, καὶ ὁ Θεὸς τῶν πατέρων ὑμῶν, ἔδωκεν ὑμῖν θησαυροὺς ἐν τοῖς μαρσίπποις ὑμῶν· καὶ τὸ ἀργύριον ὑμῶν εὐδοκιμοῦν ἀπέχω· καὶ ἐξήγαγε πρὸς αὐτοὺς τὸν Συμεών.
\vs{24}Καὶ ἤνεγκεν ὕδωρ νίψαι τοὺς πόδας αὐτῶν· καὶ ἔδωκε χορτάσματα τοῖς ὄνοις αὐτῶν.
\vs{25}Ἡτοίμασαν δὲ τὰ δῶρα, ἕως τοῦ ἐλθεῖν τὸν Ἰωσὴφ μεσημβρίας· ἤκουσαν γὰρ ὅτι ἐκεῖ μέλλει ἀριστᾷν.
\vs{26}Εἰσῆλθε δὲ Ἰωσὴφ εἰς τὴν οἰκίαν, καὶ προσήνεγκαν αὐτῷ τὰ δῶρα, ἃ εἶχον ἐν ταῖς χερσὶν αὐτῶν, εἰς τὸν οἶκον· καὶ προσεκύνησαν αὐτῷ ἐπὶ πρόσωπον ἐπὶ τὴν γῆν.
\vs{27}Ἠρώτησε δὲ αὐτοὺς, πῶς ἔχετε; καὶ εἶπεν αὐτοῖς, εἰ ὑγιαίνει ὁ πατὴρ ὑμῶν ὁ πρεσβύτης, ὃν εἴπατε; ἔτι ζῇ;
\vs{28}Οἱ δὲ εἶπαν, ὑγιαίνει ὁ παῖς σου ὁ πατὴρ ἡμῶν, ἔτι ζῇ. Καὶ εἶπεν, εὐλογημένος ὁ ἄνθρωπος ἐκεῖνος τῷ Θεῷ· καὶ κύψαντες προσεκύνησαν αὐτῷ.
\vs{29}Ἀναβλέψας δὲ τοῖς ὀφθαλμοῖς αὐτοῦ Ἰωσὴφ, εἶδε Βενιαμὶν τὸν ἀδελφὸν αὐτοῦ τὸν ὁμομήτριον· καὶ εἶπεν, οὗτος ὁ ἀδελφὸς ὑμῶν ὁ νεώτερος, ὃν εἴπατε πρός με ἀγαγεῖν; καὶ εἶπεν, ὁ Θεὸς ἐλεήσαι σε, τέκνον.
\vs{30}Ἐταράχθη δὲ Ἰωσήφ· συνεστρέφετο γὰρ τὰ ἔγκατα αὐτοῦ ἐπὶ τῷ ἀδελφῷ αὐτοῦ, καὶ ἐζήτει κλαῦσαι· εἰσελθὼν δὲ εἰς τὸ ταμεῖον, ἔκλαυσεν ἐκεῖ.

\vs{31}Καὶ νιψάμενος τὸ πρόσωπον, ἐξελθὼν ἐνεκρατεύσατο· καὶ εἶπε, παράθετε ἄρτους.
\vs{32}Καὶ παρέθηκαν αὐτῷ μόνῳ, καὶ αὐτοῖς καθʼ ἑαυτούς, καὶ τοῖς Αἰγυπτίοις τοῖς συνδειπνοῦσι μετʼ αὐτοῦ καθʼ ἑαυτούς· οὐ γὰρ ἐδύναντο οἱ Αἰγύπτιοι συνεσθίειν μετὰ τῶν Ἐβραίων ἄρτους· βδέλυγμα γάρ ἐστι τοῖς Αἰγυπτίοις.
\vs{33}Ἐκάθισαν δὲ ἐναντίον αὐτοῦ, ὁ πρωτότοκος κατὰ τὰ πρεσβεῖα αὐτοῦ, καὶ ὁ νεώτερος κατὰ τὴν νεότητα αὐτοῦ· ἐξίσταντο δὲ οἱ ἄνθρωποι ἕκαστος πρὸς τὸν ἀδελφὸν αὐτοῦ.
\vs{34}Ἦραν δὲ μερίδα παρʼ αὐτοῦ πρὸς ἑαυτούς· ἐμεγαλύνθη δὲ ἡ μερὶς Βενιαμεὶν παρὰ τὰς μερίδας πάντων πενταπλασίως πρὸς τὰς ἐκείνων· ἔπιον δὲ καὶ ἐμεθύσθησαν μετʼ αὐτοῦ.

\ch{44}Καὶ ἐνετείλατο ὁ Ἰωσὴφ τῷ ὄντι ἐπὶ τῆς οἰκίας αὐτοῦ, λέγων, πλήσατε τοὺς μαρσίππους τῶν ἀνθρώπων βρωμάτων, ὅσα ἐὰν δύνωνται ἆραι· καὶ ἐμβάλετε ἑκάστου τὸ ἀργύριον ἐπὶ τοῦ στόματος τοῦ μαρσίππου.
\vs{2}Καὶ τὸ κόνδυ μου τὸ ἀργυροῦν ἐμβάλετε εἰς τὸν μάρσιππον τοῦ νεωτέρου, καὶ τὴν τιμὴν τοῦ σίτου αὐτοῦ· ἐγενήθη δὲ κατὰ τὸ ῥῆμα Ἰωσὴφ, καθὼς εἶπε.

\vs{3}Τὸ πρωῒ διέφαυσε· καὶ οἱ ἄνθρωποι ἀπεστάλησαν, αὐτοὶ καὶ οἱ ὄνοι αὐτῶν.
\vs{4}Ἐξελθόντων δὲ αὐτῶν τὴν πόλιν, οὐκ ἀπέσχον μακράν· καὶ Ἰωσὴφ εἶπε τῷ ἐπὶ τῆς οἰκίας αὐτοῦ, ἀναστὰς ἐπιδίωξον ὀπίσω τῶν ἀνθρώπων, καὶ καταλήμψῃ αὐτοὺς, καὶ ἐρεῖς αὐτοῖς τί ὅτι ἀνταπεδώκατε πονηρὰ ἀντὶ καλῶν;
\vs{5}Ἱνατί ἐκλέψατέ μου τὸ κόνδυ τὸ ἀργυροῦν; οὐ τοῦτό ἐστιν, ἐν ᾧ πίνει ὁ κύριός μου; αὐτὸς δὲ οἰωνισμῷ οἰωνίζεται ἐν αὐτῷ. πονηρὰ συντετελέκατε ἃ πεποιήκατε.
\vs{6}Εὑρὼν δὲ αὐτοὺς, εἶπεν αὐτοῖς κατὰ τὰ ῥήματα ταῦτα.
\vs{7}Οἱ δὲ εἶπαν αὐτῷ, ἱνατί λαλεῖ ὁ κύριος κατὰ τὰ ῥήματα ταῦτα; μὴ γένοιτο τοῖς παισίν σου ποιῆσαι κατὰ τὸ ῥῆμα τοῦτο.
\vs{8}Εἰ τὸ μὲν ἀργύριον, ὃ εὕρομεν ἐν τοῖς μαρσίπποις ἡμῶν, ἀπεστρέψαμεν πρὸς σὲ ἐκ γῆς Χαναὰν, πῶς ἂν κλέψαιμεν ἐκ τοῦ οἴκου τοῦ κυρίου σου ἀργύριον ἢ χρυσίον;
\vs{9}Παρʼ ᾧ ἂν εὕρῃς τὸ κόνδυ τῶν παιδων σου, ἀποθνησκέτω· καὶ ἡμεῖς δὲ ἐσόμεθα παῖδες τῷ κυρίῳ ἡμῶν.
\vs{10}Ὁ δὲ εἶπε, καὶ νῦν, ὡς λέγετε, οὕτως ἔσται· παρʼ ᾧ ἂν εὑρεθῇ τὸ κόνδυ, ἔσται μου παῖς ὑμεῖς δὲ ἔσεσθε καθαροί.
\vs{11}Καὶ ἔσπευσαν, καὶ καθεῖλαν ἕκαστος τὸν μάρσιππον αὐτοῦ ἐπὶ τὴν γῆν· καὶ ἤνοιξαν ἕκαστος τὸν μάρσιππον αὐτοῦ.
\vs{12}Ἠρεύνησε δὲ ἀπὸ τοῦ πρεσβυτέρου ἀρξάμενος, ἕως ἦλθεν ἐπὶ τὸν νεώτερον. καὶ εὗρε τὸ κόνδυ ἐν τῷ μαρσίππῳ τοῦ Βενιαμίν.
\vs{13}Καὶ διέῤῥηξαν τὰ ἱμάτια αὐτῶν, καὶ ἐπέθηκαν ἕκαστος τὸν μαρσίππον αὐτοῦ ἐπὶ τὸν ὄνον αὐτοῦ, καὶ ἐπέστρεψαν εἰς τὴν πόλιν.

\vs{14}Εἰσῆλθε δὲ Ἰούδας καὶ οἱ ἀδελφοὶ αὐτοῦ πρὸς Ἰωσὴφ ἔτι αὐτοῦ ὄντος ἐκεῖ, καὶ ἔπεσον ἐναντίον αὐτοῦ ἐπὶ τὴν γῆν.
\vs{15}Εἶπε δὲ αὐτοῖς Ἰωσὴφ, τί τὸ πρᾶγμα τοῦτο ἐποιήσατε; οὐκ οἴδατε ὅτι οἰωνισμῷ οἰωνιεῖται ὁ ἄνθρωπος, οἷος ἐγώ;
\vs{16}Εἶπε δὲ Ἰούδας, τί ἀντεροῦμεν τῷ κυρίῳ, ἢ τί λαλήσομεν, ἢ τί δικαιωθῶμεν; ὁ Θεὸς δὲ εὗρε τὴν ἀδικίαν τῶν παίδων σου· ἰδού ἐσμεν οἰκέται τῷ κυρίῳ ἡμῶν, καὶ ἡμεῖς, καὶ παρʼ ᾧ εὑρέθη τὸ κόνδυ.
\vs{17}Εἶπε δὲ Ἰωσὴφ, μή μοι γένοιτο ποιῆσαι τὸ ῥῆμα τοῦτο· ὁ ἄνθρωπος παρʼ ᾧ εὑρέθη τὸ κόνδυ, αὐτὸς ἔσται μου παῖς· ὑμεῖς δὲ ἀνάβητε μετὰ σωτηρίας πρὸς τὸν πατέρα ὑμῶν.
\vs{18}Ἐγγίσας δὲ αὐτῷ Ἰούδας εἶπε, δέομαι, κύριε· λαλησάτω ὁ παῖς σου ῥῆμα ἐναντίον σου, καὶ μὴ θυμωθῇς τῷ παιδί σου, ὅτι σὺ εἶ μετὰ Φαραώ.
\vs{19}Κύριε, σὺ ἠρώτησας τοὺς παῖδάς σου, λέγων, εἰ ἔχετε πατέρα ἢ ἀδελφόν.
\vs{20}Καὶ εἴπαμεν τῷ κυρίῳ, ἔστιν ἡμῖν πατὴρ πρεσβύτερος, καὶ παιδίον γήρως νεώτερον αὐτῷ, καὶ ὁ ἀδελφὸς αὐτοῦ ἀπέθανεν, αὐτὸς δὲ μόνος ὑπελείφθη τῇ μητρὶ αὐτοῦ, ὁ δὲ πατὴρ αὐτὸν ἠγάπησεν·
\vs{21}Εἶπας δὲ τοῖς παισί σου, καταγάγετε αὐτὸν πρὸς μέ, καὶ ἐπιμελοῦμαι αὐτοῦ.
\vs{22}Καὶ εἴπαμεν τῷ κυρίῳ, οὐ δυνήσεται τὸ παιδίον καταλιπεῖν τὸν πατέρα αὐτοῦ· ἐὰν δὲ καταλείπῃ τὸν πατέρα, ἀποθανεῖται.
\vs{23}Σὺ δὲ εἶπας τοῖς παισί σου, ἐὰν μὴ καταβῇ ὁ ἀδελφὸς ὑμῶν ὁ νεώτερος μεθʼ ὑμῶν, οὐ προσθήσεσθε ἰδεῖν τὸ πρόσωπόν μου.
\vs{24}Ἐγένετο δὲ ἡνίκα ἀνέβημεν πρὸς τὸν παῖδά σου πατέρα ἡμῶν, ἀπηγγείλαμεν αὐτῷ τὰ ῥήματα τοῦ κυρίου ἡμῶν.
\vs{25}Εἶπε δὲ ὁ πατὴρ ἡμῶν, βαδίσατε πάλιν καὶ ἀγοράσατε ἡμῖν μικρὰ βρώματα.
\vs{26}Ἡμεῖς δὲ εἴπομεν, οὐ δυνησόμεθα καταβῆναι· ἀλλʼ εἰ μὲν ὁ ἀδελφὸς ἡμῶν ὁ νεώτερος καταβαίνει μεθʼ ἡμῶν, καταβησόμεθα· οὐ γὰρ δυνησόμεθα ἰδεῖν τὸ πρόσωπον τοῦ ἀνθρώπου, τοῦ ἀδελφοῦ ἡμῶν τοῦ νεωτέρου μὴ ὄντος μεθʼ ἡμῶν.
\vs{27}Εἶπε δὲ ὁ παῖς σου πατὴρ ἡμῶν πρὸς ἡμᾶς, ὑμεῖς γινώσκετε ὅτι δύο ἔτεκέ μοι ἡ γυνὴ,
\vs{28}καὶ ἐξῆλθεν ὁ εἷς ἀπʼ ἐμοῦ· καὶ εἴπατε ὅτι θηριόβρωτος γέγονεν, καὶ οὐκ ἴδον αὐτὸν ἄχρι νῦν.
\vs{29}Ἐὰν οὖν λάβητε καὶ τοῦτον ἐκ τοῦ προσώπου μου, καὶ συμβῇ αὐτῷ μαλακία ἐν τῇ ὁδῷ, καὶ κατάξετέ μου τὸ γῆρας μετὰ λύπης εἰς ᾅδου.
\vs{30}Νῦν οὖν ἐὰν εἰσπορεύωμαι πρὸς τὸν παῖδά σου, πατέρα δὲ ἡμῶν, καὶ τὸ παιδίον μὴ ᾖ μεθʼ ἡμῶν, ἡ δὲ ψυχὴ αὐτοῦ ἐκκρέμαται ἐκ τῆς τούτου ψυχῆς,
\vs{31}καὶ ἔσται ἐν τῷ ἰδεῖν αὐτὸν μὴ ὂν τὸ παιδίον μεθʼ ἡμῶν, τελευτήσει, καὶ κατάξουσιν οἱ παῖδές σου τὸ γῆρας τοῦ παιδός σου, πατρὸς δὲ ἡμῶν, μετὰ λύπης εἰς ᾅδου.
\vs{32}Ὁ γὰρ παῖς σου παρὰ τοῦ πατρὸς ἐκδέδεκται τὸ παιδίον, λέγων, ἐὰν μὴ ἀγάγω αὐτὸν πρὸς σὲ, καὶ στήσω αὐτὸν ἐνώπιόν σου, ἡμαρτηκὼς ἔσομαι εἰς τὸν πατέρα πάσας τὰς ἡμέρας.
\vs{33}Νῦν οὖν παραμενῶ σοι παῖς ἀντὶ τοῦ παιδίου, οἰκέτης τοῦ κυρίου· τὸ δὲ παιδίον ἀναβήτω μετὰ τῶν ἀδελφῶν αὐτοῦ.
\vs{34}Πῶς γὰρ ἀναβήσομαι πρὸς τὸν πατέρα, τοῦ παιδίου μὴ ὄντος μεθʼ ἡμῶν; ἵνα μὴ ἴδω τὰ κακὰ, ἃ εὑρήσει τὸν πατέρα μου.

\ch{45}
Καὶ οὐκ ἠδύνατο Ἰωσὴφ ἀνέχεσθαι πάντων τῶν παρεστηκότων αὐτῷ, ἀλλʼ εἶπεν, ἐξαποστείλατε πάντας ἀπʼ ἐμοῦ· καὶ οὐ παρειστήκει οὐδεὶς τῷ Ἰωσὴφ, ἡνίκα ἀνεγνωρίζετο τοῖς ἀδελφοῖς αὐτοῦ.
\vs{2}Καὶ ἀφῆκε φωνὴν μετὰ κλαυθμοῦ· ἤκουσαν δὲ πάντες οἱ Αἰγύπτιοι, καὶ ἀκουστὸν ἐγένετο εἰς τὸν οἶκον Φαραώ.
\vs{3}Εἶπε δὲ Ἰωσὴφ πρὸς τοὺς ἀδελφοὺς αὐτοῦ, ἐγώ εἰμι Ἰωσήφ· ἔτι ὁ πατήρ μου ζῇ; καὶ οὐκ ἠδύναντο οἱ ἀδελφοὶ ἀποκριθῆναι αὐτῷ· ἐταράχθησαν γάρ.
\vs{4}Εἶπε δὲ Ἰωσὴφ πρὸς τοὺς ἀδελφοὺς αὐτοῦ, ἐγγίσατε πρὸς μέ· καὶ ἤγγισαν· καὶ εἶπεν, ἐγώ εἰμι Ἰωσὴφ ὁ ἀδελφὸς ὑμῶν, ὃν ἀπέδοσθε εἰς Αἴγυπτον.
\vs{5}Νῦν οὖν μὴ λυπεῖσθε, μηδὲ σκληρὸν ὑμῖν φανήτω, ὅτι ἀπέδοσθέ με ὧδε· εἰς γὰρ ζωὴν ἀπέστειλέ με ὁ Θεὸς ἔμπροσθεν ὑμῶν.
\vs{6}Τοῦτο γὰρ δεύτερον ἔτος λιμὸς ἐπὶ τῆς γῆς, καὶ ἔτι λοιπὰ πέντε ἔτη, ἐν οἷς οὐκ ἔσται ἀροτρίασις, οὐδὲ ἀμητός·
\vs{7}Ἀπέστειλε γάρ με ὁ Θεὸς ἔμπροσθεν ὑμῶν, ὑπολείπεσθαι ὑμῶν κατάλειμμα ἐπὶ τῆς γῆς, καὶ ἐκθρέψαι ὑμῶν κατάλειψιν μεγάλην.
\vs{8}Νῦν οὖν οὐχ ὑμεῖς με ἀπεστάλκατε ὧδε, ἀλλὰ ὁ Θεός· καὶ ἐποίησέ με ὡς πατέρα Φαραὼ, καὶ κύριον παντὸς τοῦ οἴκου αὐτοῦ, καὶ ἄρχοντα πάσης γῆς Αἰγύπτου.
\vs{9}Σπεύσαντες οὖν ἀνάβητε πρὸς τὸν πατέρα μου, καὶ εἴπατε αὐτῷ, τάδε λέγει ὁ υἱός σου Ἰωσήφ· ἐποίησέ με ὁ Θεὸς κύριον πάσης γῆς Αἰγύπτου· κατάβηθι οὖν πρός με, καὶ μὴ μείνῃς·
\vs{10}Καὶ κατοικήσεις ἐν γῇ Γεσὲμ Ἀραβίας· καὶ ἔσῃ ἐγγύς μου σὺ, καὶ οἱ υἱοί σου, καὶ οἱ υἱοὶ τῶν υἱῶν σου, τὰ πρόβατά σου, καὶ αἱ βόες σου, καὶ ὅσα σοι ἐστί.
\vs{11}Καὶ ἐκθρέψω σε ἐκεῖ· ἔτι γὰρ πέντε ἔτη λιμός· ἵνα μὴ ἐκτριβῇς σὺ, καὶ οἱ υἱοί σου, καὶ πάντα τὰ ὑπάρχοντά σου.
\vs{12}Ἰδοὺ οἱ ὀφθαλμοὶ ὑμῶν βλέπουσι, καὶ οἱ ὀφθαλμοὶ Βενιαμεὶν τοῦ ἀδελφοῦ μου, ὅτι τὸ στόμα μου τὸ λαλοῦν πρὸς ὑμᾶς.
\vs{13}Ἀπαγγείλατε οὖν τῷ πατρί μου πᾶσαν τὴν δόξαν μου τὴν ἐν Αἰγύπτῳ, καὶ ὅσα ἴδετε· καὶ ταχύναντες καταγάγετε τὸν πατέρα μου ὧδε.
\vs{14}Καὶ ἐπιπεσὼν ἐπὶ τὸν τράχηλον Βενιαμὶν τοῦ ἀδελφοῦ αὐτοῦ, ἔκλαυσεν ἐπʼ αὐτῷ· καὶ Βενιαμὶν ἔκλαυσεν ἐπὶ τῷ τραχήλῳ αὐτοῦ.
\vs{15}Καὶ καταφιλήσας πάντας τοὺς ἀδελφοὺς αὐτοῦ, ἔκλαυσεν ἐπʼ αὐτοῖς· καὶ μετὰ ταῦτα ἐλάλησαν οἱ ἀδελφοὶ αὐτοῦ πρὸς αὐτόν.

\vs{16}Καὶ διεβοήθη ἡ φωνὴ εἰς τὸν οἶκον Φαραὼ, λέγοντες, ἥκασιν οἱ ἀδελφοὶ Ἰωσήφ· ἐχάρη δὲ Φαραὼ καὶ ἡ θεραπεία αὐτοῦ.
\vs{17}Εἶπε δὲ Φαραὼ πρὸς Ἰωσὴφ, εἰπον τοῖς ἀδελφοῖς σου, τοῦτο ποιήσατε, γεμίσατε τὰ φορεῖα ὑμῶν, καὶ ἀπέλθετε εἰς γῆν Χαναάν.
\vs{18}Καὶ ἀναλαβόντες τὸν πατέρα ὑμῶν, καὶ τὰ ὑπάρχοντα ὑμῶν, ἥκετε πρός με· καὶ δώσω ὑμῖν πάντων τῶν ἀγαθῶν Αἰγύπτου, καὶ φάγεσθε τὸν μυελὸν τῆς γῆς.
\vs{19}Σὺ δὲ ἔντειλαι ταῦτα· λαβεῖν αὐτοῖς ἁμάξας ἐκ γῆς Αἰγύπτου τοῖς παιδίοις ὑμῶν, καὶ ταῖς γυναιξὶν ὑμῶν· καὶ ἀναλαβόντες τὸν πατέρα ὑμῶν παραγίνεσθε.
\vs{20}Καὶ μὴ φείσησθε τοῖς ὀφθαλμοῖς τῶν σκευῶν ὑμῶν· τὰ γὰρ πάντα ἀγαθὰ Αἰγύπτου ὑμῖν ἔσται.
\vs{21}Ἐποίησαν δὲ οὕτως οἱ υἱοὶ Ἰσραήλ· ἔδωκε δὲ Ἰωσὴφ αὐτοῖς ἁμάξας κατὰ τὰ εἰρημένα ὑπὸ Φαραὼ τοῦ βασιλέως· καὶ ἔδωκεν αὐτοῖς ἐπισιτισμὸν εἰς τὴν ὁδόν·
\vs{22}Καὶ πᾶσιν ἔδωκε δισσὰς στολάς· τῷ δὲ Βενιαμὶν ἔδωκε τριακοσίους χρυσοὺς, καὶ πέντε ἐξαλλασσούσας στολάς.
\vs{23}Καὶ τῷ πατρὶ αὐτοῦ ἀπέστειλε κατὰ τὰ αὐτά· καὶ δέκα ὄνους, αἴροντας ἀπὸ πάντων τῶν ἀγαθῶν Αἰγύπτου, καὶ δέκα ἡμιόνους, αἰρούσας ἄρτους τῷ πατρὶ αὐτοῦ εἰς ὁδόν.
\vs{24}Ἐξαπέστειλε δὲ τοὺς ἀδελφοὺς αὐτοῦ, καὶ ἐπορεύθησαν· καὶ εἶπεν αὐτοῖς, μὴ ὀργίζεσθε ἐν τῇ ὁδῷ.
\vs{25}Καὶ ἀνέβησαν ἐξ Αἰγυπτου, καὶ ἦλθον εἰς γῆν Χαναὰν πρὸς Ἰακὼβ τὸν πατέρα αὐτῶν.
\vs{26}Καὶ ἀνήγγειλαν αὐτῷ λέγοντες, ὅτι ὁ υἱός σου Ἰωσὴφ ζῇ, καὶ αὐτὸς ἄρχει πάσης γῆς Αἰγύπτου· καὶ ἐξέστη τῇ διανοίᾳ Ἰακὼβ, οὐ γὰρ ἐπίστευσεν αὐτοῖς.
\vs{27}Ἐλάλησαν δὲ αὐτῷ πάντα τὰ ῥηθέντα ὑπὸ Ἰωσὴφ, ὅσα εἶπεν αὐτοῖς. Ἰδὼν δὲ τὰς ἁμάξας, ἃς ἀπέστειλεν Ἰωσὴφ ὥστε ἀναλαβεῖν αὐτὸν, ἀνεζωπύρησε τὸ πνεῦμα Ἰακὼβ τοῦ πατρὸς αὐτῶν.
\vs{28}Εἶπε δὲ Ἰσραὴλ, μέγα μοι ἐστὶν, εἰ ἔτι Ἰωσὴφ ὁ υἱός μου ζῇ· πορευθεὶς ὄψομαι αὐτὸν πρὸ τοῦ ἀποθανεῖν με.

\ch{46}
Ἀπᾴρας δὲ Ἰσραὴλ, αὐτὸς καὶ πάντα τὰ αὐτοῦ, ἦλθεν ἐπὶ τὸ φρέαρ τοῦ ὅρκου· καὶ ἔθυσε θυσίαν τῷ Θεῷ τοῦ πατρὸς αὐτοῦ Ἰσαάκ.
\vs{2}Εἶπε δὲ ὁ Θεὸς τῷ Ἰσραὴλ ἐν ὁράματι τῆς νυκτὸς, εἰπὼν, Ἰακὼβ, Ἰακώβ· ὁ δὲ εἶπε, τί ἐστιν;
\vs{3}Ὁ δὲ λέγει αὐτῷ, ἐγώ εἰμι ὁ Θεὸς τῶν πατέρων σου· μὴ φοβοῦ καταβῆναι εἰς Αἴγυπτον· εἰς γὰρ ἔθνος μέγα ποιήσω σε ἐκεῖ.
\vs{4}Καὶ ἐγὼ καταβήσομαι μετὰ σοῦ εἰς Αἴγυπτον, καὶ ἐγὼ ἀναβιβάσω σε εἰς τέλος· καὶ Ἰωσὴφ ἐπιβαλεῖ τὰς χεῖρας ἐπὶ τοὺς ὀφθαλμούς σου.
\vs{5}Ἀνέστη δὲ Ἰακὼβ ἀπὸ τοῦ φρέατος τοῦ ὅρκου· καὶ ἀνέλαβον οἱ υἱοὶ Ἰσραὴλ τὸν πατέρα αὐτῶν, καὶ τὴν ἀποσκευὴν, καὶ τὰς γυναῖκας αὐτῶν, ἐπὶ τὰς ἁμάξας, ἃς ἀπέστειλεν Ἰωσὴφ ἆραι αὐτόν.
\vs{6}Καὶ ἀναλαβόντες τὰ ὑπάρχοντα αὐτῶν, καὶ πᾶσαν τὴν κτῆσιν, ἣν ἐκτήσαντο ἐκ γῇ Χαναὰν, εἰσῆλθον εἰς Αἴγυπτον, Ἰακὼβ, καὶ πᾶν τὸ σπέρμα αὐτοῦ μετʼ αὐτοῦ.
\vs{7}Υἱοὶ, καὶ υἱοὶ τῶν υἱῶν αὐτοῦ μετʼ αὐτοῦ· θυγατέρες, καὶ θυγατέρες τῶν θυγατέρων αὐτοῦ· καὶ πᾶν τὸ σπέρμα αὐτοῦ ἤγαγεν εἰς Αἴγυπτον·
\vs{8}Ταῦτα δὲ τὰ ὀνόματα τῶν υἱῶν Ἰσραὴλ τῶν εἰσελθόντων εἰς Αἴγυπτον ἅμα Ἰακὼβ τῷ πατρὶ αὐτῶν. Ἰακὼβ καὶ οἱ υἱοὶ αὐτοῦ· πρωτότοκος Ἰακὼβ, Ῥουβήν.
\vs{9}γἱοὶ δὲ Ῥουβὴν, Ἑνὼχ, καὶ Φαλλὸς, Ἀσρὼν, καὶ Χαρμί.
\vs{10}Υἱοὶ δὲ Συμεὼν, Ἰεμουὴλ, καὶ Ἰαμεὶν, καὶ Ἀὼδ, καὶ Ἀχεὶν, καὶ Σαὰρ, καὶ Σαοὺλ υἱὸς τῆς Χανανίτιδος.
\vs{11}Υἱοὶ δὲ Λευὶ, Γηρσὼν, Κὰθ, καὶ Μεραρί.
\vs{12}Υἱοὶ δὲ Ἰούδα, Ἢρ, καὶ Αὐνὰν, καὶ Σηλὼμ, καὶ Φαρὲς, καὶ Ζαρά· ἀπέθανε δὲ Ἢρ καὶ Αὐνὰν ἐν γῇ Χαναάν· ἐγένοντο δὲ υἱοὶ Φαρὲς, Ἑσρὼν, καὶ Ἰεμουήλ.
\vs{13}Υἱοὶ δὲ Ἰσσάχαρ, Θωλὰ, καὶ Φουὰ, καὶ Ἀσοὺμ, καὶ Σαμβράν.
\vs{14}Υἱοὶ δὲ Ζαβουλὼν, Σερὲδ, καὶ Ἀλλὼν, καὶ Ἀχοήλ.
\vs{15}Οὗτοι υἱοὶ Λείας, οὓς ἔτεκε τῷ Ἰακὼβ ἐν Μεσοποταμίᾳ τῆς Συρίας, καὶ Δείναν τὴν θυγατέρα αὐτοῦ· πᾶσαι αἱ ψυχαί, υἱοὶ καὶ θυγατέρες, τριάκοντα τρεῖς.
\vs{16}Υἱοὶ δὲ Γάδ, Σαφὼν, καὶ Ἀγγὶς, καὶ Σαννὶς, καὶ Θασοβὰν, καὶ Ἀηδεὶς, καὶ Ἀροηδεὶς, καὶ Ἀρεηλείς.
\vs{17}Υἱοὶ δὲ Ἀσὴρ, Ἰεμνα, Ἰεσσουὰ, καὶ Ἰεοὺλ, καὶ βαριὰ, καὶ Σάρα ἀδελφὴ αὐτῶν. Υἱοὶ δὲ βαριὰ, Χοβὸρ, καὶ Μελχιΐλ.
\vs{18}Οὗτοι υἱοὶ Ζελφᾶς, ἣν ἔδωκε Λάβαν Λείᾳ τῇ θυγατρὶ αὐτοῦ, ἣ ἔτεκε τούτους τῷ Ἰακὼβ, δεκαὲξ ψυχάς.
\vs{19}Υἱοὶ δὲ Ῥαχὴλ γυναικὸς Ἰακὼβ, Ἰωσὴφ, καὶ Βενιαμείν.
\vs{20}Ἐγένοντο δὲ υἱοὶ Ἰωσὴφ ἐν γῇ Αἰγύπτου, οὓς ἔτεκεν αὐτῷ Ἀσενὲθ θυγάτηρ Πετεφρῆ ἱερέως Ἡλιουπόλεως, τὸν Μανασσῆ, καὶ τὸν Ἐφραίμ· ἐγένοντο δὲ υἱοὶ Μανασσῆ, οὓς ἔτεκεν αὐτῷ ἡ παλλακὴ ἡ Σύρα, τὸν Μαχίρ· Μαχὶρ δὲ ἐγέννησε τὸν Γαλαάδ· υἱοὶ δὲ Ἐφραὶμ ἀδελφοῦ Μανασσῆ, Σουταλαὰμ, καὶ Ταάμ· υἱοὶ δὲ Σουταλαὰμ, Ἐδώμ.
\vs{21}Υἱοὶ δὲ Βενιαμεὶν, Βαλὰ καὶ Βοχὸρ, καὶ Ἀσβήλ. Ἐγένοντο δὲ υἱοὶ Βαλὰ, Γηρὰ, καὶ Νοεμὰν, καὶ Ἀγχὶς, καὶ Ῥὼς, καὶ Μαμφίμ· Γηρὰ δὲ ἐγέννησε τὸν Ἀράδ.
\vs{22}Οὗτοι υἱοὶ Ῥαχὴλ, οὓς ἔτεκε τῷ Ἰακώβ· πᾶσαι αἱ ψυχαὶ δεκαοκτώ.
\vs{23}Υἱοὶ δὲ Δὰν, Ἀσόμ.
\vs{24}Καὶ υἱοὶ Νεφθαλὶ, Ἀσιὴλ, καὶ Γωνὶ, καὶ Ἰσσάαρ, καὶ Σολλήμ.
\vs{25}Οὗτοι υἱοὶ Βαλλὰς, ἣν ἔδωκε Λάβαν Ῥαχὴλ τῇ θυγατρὶ αὐτοῦ, ἣ ἔτεκε τούτους τῷ Ἰακὼβ, πᾶσαι αἱ ψυχαὶ ἑπτά.
\vs{26}Πᾶσαι δὲ ψυχαὶ αἱ εἰσελθοῦσαι μετὰ Ἰακὼβ εἰς Αἴγυπτον, οἱ ἐξελθόντες ἐκ τῶν μηρῶν αὐτοῦ, χωρὶς τῶν γυναικῶν υἱῶν Ἰακὼβ, πᾶσαι αἱ ψυχαὶ, ἑξηκονταέξ·
\vs{27}Υἱοὶ δὲ Ἰωσὴφ οἱ γενόμενοι αὐτῷ ἐν γῇ Αἰγύπτῳ, ψυχαὶ ἐννέα. Πᾶσαι ψυχαὶ οἴκου Ἰακὼβ, αἱ εἰσελθοῦσαι μετὰ Ἰακὼβ εἰς Αἴγυπτον, ψυχαὶ ἑβδομηκονταπέντε.

\vs{28}Τὸν δὲ Ἰούδαν ἀπέστειλεν ἔμπροσθεν αὐτοῦ πρὸς Ἰωσὴφ, συναντῆσαι αὐτῷ καθʼ Ἡρώων πόλιν, εἰς γῆν Ῥαμεσσῆ.
\vs{29}Ζεύξας δὲ Ἰωσὴφ τὰ ἅρματα αὐτοῦ, ἀνέβη εἰς συνάντησιν Ἰσραὴλ τῷ πατρὶ αὐτοῦ, καθʼ Ἡρώων πόλιν· καὶ ὀφθεὶς αὐτῷ ἐπέπεσεν ἐπὶ τὸν τράχηλον αὐτοῦ, καὶ ἔκλαυσε κλαυθμῷ πίονι.
\vs{30}Καὶ εἶπεν Ἰσραήλ πρὸς Ἰωσὴφ, ἀποθανοῦμαι ἀπὸ τοῦ νῦν, ἐπεὶ ἑώρακα τὸ πρόσωπόν σου· ἔτι γὰρ σὺ ζῇς.
\vs{31}Εἶπε δὲ Ἰωσὴφ πρὸς τοὺς ἀδελφοὺς αὐτοῦ, ἀναβὰς ἀπαγγελῶ τῷ Φαραῷ, καὶ ἐρῶ αὐτῷ, οἱ ἀδελφοί μου, καὶ ὁ οἶκος τοῦ πατρός μου, οἳ ἦσαν ἐν γῇ χαναὰν, ἥκασι πρός με.
\vs{32}Οἱ δὲ ἄνδρες εἰσὶ ποιμένες· ἄνδρες γὰρ κτηνοτρόφοι ἦσαν· καὶ τὰ κτήνη, καὶ τοὺς βόας, καὶ πάντα τὰ αὐτῶν ἀγηόχασιν.
\vs{33}Ἐὰν οὖν καλέσῃ ὑμᾶς Φαραὼ, καὶ εἴπῃ ὑμῖν, τί τὸ ἔργον ὑμῶν ἐστι;
\vs{34}Ἐρεῖτε, ἄνδρες κτηνοτρόφοι ἐσμὲν οἱ παῖδές σου ἐκ παιδὸς ἕως τοῦ νῦν, καὶ ἡμεῖς, καὶ οἱ πατέρες ἡμῶν· ἵνα κατοικήσητε ἐν γῇ Γεσὲμ Ἀραβίας· βδέλυγμα γάρ ἐστιν Αἰγυπτίοις πᾶς ποιμὴν προβάτων.

\ch{47}
Ἐλθῶν δὲ Ἰωσὴφ ἀπήγγειλε τῷ Φαραῶ, λέγων, ὁ πατὴρ μου, καὶ οἱ ἀδελφοί μου, καὶ τὰ κτήνη, καὶ οἱ βόες αὐτῶν, καὶ πάντα τὰ αὐτῶν, ἦλθον ἐκ γῆς Χαναάν· καὶ ἰδού εἰσιν ἐν γῇ Γεσέμ.
\vs{2}Ἀπὸ δὲ τῶν ἀδελφῶν αὐτοῦ παρέλαβε πέντε ἄνδρας, καὶ ἔστησεν αὐτοὺς ἐναντίον Φαραώ.
\vs{3}Καὶ εἶπε Φαραὼ τοῖς ἀδελφοῖς Ἰωσὴφ, Τί τὸ ἔργον ὑμῶν; οἱ δὲ εἶπαν τῷ Φαραῷ, ποιμένες προβάτων οἱ παῖδές σου, καὶ ἡμεῖς καὶ οἱ πατέρες ἡμῶν.
\vs{4}Εἶπαν δὲ τῷ Φαραῷ, παροικεῖν ἐν τῇ γῇ ἥκαμεν, οὐ γάρ ἐστι νομὴ τοῖς κτήνεσι τῶν παιδων σου, ἐνίσχυσε γὰρ ὁ λιμὸς ἐν γῇ Χανάαν· νῦν οὖν κατοικήσομεν ἐν γῇ Γεσέμ. Εἶπε δὲ Φαραὼ τῷ Ἰωσὴφ, Κατοικείτωσαν ἐν γῇ Γεσέμ· εἰ δὲ ἐπίστῃ, ὅτι εἰσὶν ἐν αὐτοῖς ἄνδρες δυνατοὶ, κατάστησον αὐτοὺς ἄρχοντας τῶν ἐμῶν κτηνῶν. Ἦλθον δὲ εἰς Αἴγυπτον πρὸς Ἰωσὴφ Ἰακὼβ, καὶ οἱ υἱοὶ αὐτοῦ· καὶ ἤκουσε Φαραὼ βασιλεὺς Αἰγύπτου.
\vs{5}Καὶ εἶπε Φαραὼ πρὸς Ἰωσὴφ, λέγων, ὁ πατήρ σου, καὶ οἱ ἀδελφοί σου, ἥκασι πρὸς σέ.
\vs{6}Ἰδοὺ ἡ γῆ Αἰγύπτου ἐναντίον σου ἐστίν· ἐν τῇ βελτίστῃ γῇ κατοίκισον τὸν πατέρα σου, καὶ τοὺς ἀδελφούς σου.
\vs{7}Εἰσήγαγε δὲ Ἰωσὴφ Ἰακὼβ τὸν πατέρα αὐτοῦ, καὶ ἔστησεν αὐτὸν ἐναντίον Φαραώ· καὶ ηὐλόγησεν Ἰακὼβ τὸν Φαραώ.
\vs{8}Εἶπε δὲ Φαραὼ τῷ Ἰακὼβ, πόσα ἔτη ἡμερῶν τῆς ζωῆς σου;
\vs{9}Καὶ εἶπεν Ἰακὼβ τῷ Φαραῷ, αἱ ἡμέραι τῶν ἐτῶν τῆς ζωῆς μου, ἃς παροικῶ, ἑκατὸν τριάκοντα ἔτη· μικραὶ καὶ πονηραὶ γεγόνασιν αἱ ἡμέραι τῶν ἐτῶν τῆς ζωῆς μου· οὐκ ἀφίκοντο εἰς τὰς ἡμέρας τῶν ἐτῶν τῆς ζῶης τῶν πατέρων μου, ἃς ἡμέρας παρῴκησαν.
\vs{10}Καὶ εὐλογήσας Ἰακὼβ τὸν Φαραὼ, ἐξῆλθεν ἀπʼ αὐτοῦ.

\vs{11}Καὶ κατῴκισεν Ἰωσὴφ τὸν πατέρα αὐτοῦ, καὶ τοὺς ἀδελφοὺς αὐτοῦ, καὶ ἔδωκεν αὐτοῖς κατάσχεσιν ἐν γῇ Αἰγύπτῳ, ἐν τῇ βελτίστῃ γῇ, ἐν γῇ Ῥαμεσσῆ, καθὰ προσέταξε Φαραώ.
\vs{12}Καὶ ἐσιτομέτρει Ἰωσὴφ τῷ πατρὶ αὐτοῦ, καὶ τοῖς ἀδελφοῖς, καὶ παντὶ τῷ οἴκῳ τοῦ πατρὸς αὐτοῦ, σῖτον κατὰ σῶμα.

\vs{13}Σῖτος δὲ οὐκ ἦν ἐν πάσῃ τῇ γῇ, ἐνίσχυσε γὰρ ὁ λιμὸς σφόδρα· ἐξέλιπε δὲ ἡ γῆ Αἰγύπτου καὶ ἡ γῆ Χαναὰν ἀπὸ τοῦ λιμοῦ.
\vs{14}Συνήγαγε δὲ Ἰωσὴφ πᾶν τὸ ἀργύριον τὸ εὑρεθὲν ἐν γῇ Αἰγύπτου καὶ ἐν γῇ Χαναὰν, τοῦ σίτου, οὗ ἠγόραζον, καὶ ἐσιτομέτρει αὐτοῖς, καὶ εἰσήνεγκεν Ἰωσὴφ πᾶν τὸ ἀργύριον εἰς τὸν οἶκον Φαραώ.
\vs{15}Καὶ ἐξέλιπε πᾶν τὸ ἀργύριον ἐκ γῆς Αἰγύπτου καὶ ἐκ γῆς Χαναάν· ἦλθον δὲ πάντες οἱ Αἰγύπτιοι πρὸς Ἰωσὴφ, λέγοντες, δὸς ἡμῖν ἄρτους, καὶ ἱνατί ἀποθνήσκομεν ἐναντίον σου; ἐκλέλοιπε γὰρ τὸ ἀργύριον ἡμῶν.
\vs{16}Εἶπε δὲ αὐτοῖς Ἰωσὴφ, φέρετε τὰ κτήνη ὑμῶν, καὶ δώσω ὑμῖν ἄρτους, ἀντὶ τῶν κτηνῶν ὑμῶν, εἰ ἐκλέλοιπε τὸ ἀργύριον ὑμῶν.
\vs{17}Ἤγαγον δὲ τὰ κτήνη αὐτῶν πρὸς Ἰωσήφ· καὶ ἔδωκεν αὐτοῖς Ἰωσὴφ ἄρτους ἀντὶ τῶν ἵππων, καὶ ἀντὶ τῶν προβάτων, καὶ ἀντὶ τῶν βοῶν, καὶ ἀντὶ τῶν ὄνων· καὶ ἐξέθρεψεν αὐτοὺς ἐν ἄρτοις ἀντὶ πάντων τῶν κτηνῶν αὐτῶν ἐν τῷ ἐνιαυτῷ ἐκείνῳ.
\vs{18}Ἐξῆλθε δὲ τὸ ἔτος ἐκεῖνο, καὶ ἦλθον πρὸς αὐτὸν ἐν τῷ ἔτει τῷ δευτέρῳ, καὶ εἶπαν αὐτῷ, μή ποτε ἐκτριβῶμεν ἀπὸ τοῦ κυρίου ἡμῶν; εἰ γὰρ ἐκλέλοιπε τὸ ἀργύριον ἡμῶν, καὶ τὰ ὑπάρχοντα καὶ τὰ κτήνη πρός σε τὸν κύριον, καὶ οὐχ ὑπολέλειπται ἡμῖν ἐναντίον τοῦ κυρίου ἡμῶν, ἀλλʼ ἢ τὸ ἴδιον σῶμα καὶ ἡ γῆ ἡμῶν,
\vs{19}ἵνα οὖν μὴ ἀποθάνωμεν ἐναντίον σου, καὶ ἡ γῆ ἐρημωθῇ, κτῆσαι ἡμᾶς καὶ τὴν γῆν ἡμῶν ἀντὶ ἄρτων, καὶ ἐσόμεθα ἡμεῖς καὶ ἡ γῆ ἡμῶν παῖδες τῷ Φαραώ· δὸς σπέρμα, ἵνα σπείρωμεν, καὶ ζῶμεν καὶ μὴ ἀποθάνωμεν, καὶ ἡ γῆ οὐκ ἐρημωθήσεται.
\vs{20}Καὶ ἐκτήσατο Ἰωσὴφ πᾶσαν τὴν γῆν τῶν Αἰγυπτίων τῷ Φαραώ· ἀπέδοντο γὰρ οἱ Αἰγύπτιοι τὴν γῆν αὐτῶν τῷ Φαραώ· ἐπεκράτησε γὰρ αὐτῶν ὁ λιμός· καὶ ἐγένετο ἡ γῆ τῷ Φαραώ.
\vs{21}Καὶ τὸν λαὸν κατεδουλώσατο αὐτῷ εἰς παῖδας, ἀπʼ ἄκρων ὁρίων Αἰγύπτου ἕως τῶν ἄκρων,
\vs{22}χωρὶς τῆς γῆς τῶν ἱερέων μόνον· οὐκ ἐκτήσατο ταύτην Ἰωσήφ· ἐν δόσει γὰρ ἔδωκε δόμα τοῖς ἱερεῦσι Φαραὼ, καὶ ἤσθιον τὴν δόσιν, ἣν ἔδωκεν αὐτοῖς Φαραώ· διὰ τοῦτο οὐκ ἀπέδοντο τὴν γῆν αὐτῶν.
\vs{23}Εἶπε δὲ Ἰωσὴφ πᾶσι τοῖς Αἰγυπτίοις, ἰδοὺ κέκτημαι ὑμᾶς καὶ τὴν γῆν ὑμῶν σήμερον τῷ Φαραῷ· λάβετε ἑαυτοῖς σπέρμα, καὶ σπείρατε τὴν γῆν.
\vs{24}Καὶ ἔσται τὰ γεννήματα αὐτῆς· καὶ δώσετε τὸ πεμπτὸν μέρος τῷ Φαραώ· τὰ δὲ τέσσαρα μέρη ἔσται ὑμῖν αὐτοῖς εἰς σπέρμα τῇ γῇ, καὶ εἰς βρῶσιν ὑμῖν, καὶ πᾶσι τοῖς ἐν τοῖς οἴκοις ὑμῶν.
\vs{25}Καὶ εἶπαν, σέσωκας ἡμᾶς· εὕρομεν χάριν ἐναντίον τοῦ κυρίου ἡμῶν, καὶ ἐσόμεθα παῖδες τῷ Φαραώ.
\vs{26}Καὶ ἔθετο αὐτοῖς Ἰωσὴφ εἰς πρόσταγμα ἕως τῆς ἡμέρας ταύτης, ἐπὶ γῆς Αἰγύπτου τῷ Φαραὼ ἀποπεμπτοῦν, χωρὶς τῆς γῆς τῶν ἱερέων μόνον· οὐκ ἧν τῷ Φαραώ.

\vs{27}Κατῶκησε δὲ Ἰσραὴλ ἐν γῇ Αἰγύπτῳ ἐπὶ γῆς Γεσὲμ, καὶ ἐκληρονόμησαν ἐπʼ αὐτῆς· καὶ ηὐξήθησαν καὶ ἐπληθύνθησαν σφόδρα.
\vs{28}Ἐπέζησε δὲ Ἰακὼβ ἐν γῇ Αἰγύπτῳ δεκαεπτὰ ἔτη· καὶ ἐγένοντο αἱ ἡμέραι Ἰακὼβ ἐνιαυτῶν τῆς ζωῆς αὐτοῦ ἑκατὸν τεσσαρακονταεπτὰ ἔτη.
\vs{29}Ἤγγισαν δὲ αἱ ἡμέραι Ἰσραὴλ τοῦ ἀποθανεῖν· καὶ ἐκάλεσε τὸν υἱὸν αὐτοῦ Ἰωσὴφ, καὶ εἶπεν αὐτῷ, εἰ εὕρηκα χάριν ἐναντίον σου, ὑπόθες τὴν χεῖρά σου ὑπὸ τὸν μηρόν μου, καὶ ποιήσεις ἐπʼ ἐμὲ ἐλεημοσύνην, καὶ ἀλήθειαν, τοῦ μή με θάψαι ἐν Αἰγύπτῳ·
\vs{30}Ἀλλὰ κοιμηθήσομαι μετὰ τῶν πατέρων μου· καὶ ἀρεῖς με ἐξ Αἰγύπτου, καὶ θάψεις με ἐν τῷ τάφῳ αὐτῶν· ὁ δὲ εἶπεν, ἐγὼ ποιήσω κατὰ τὸ ῥῆμά σου.
\vs{31}Εἶπε δὲ, ὄμοσόν μοι· καὶ ὤμοσεν αὐτῷ· καὶ προσεκύνησεν Ἰσραὴλ ἐπὶ τὸ ἄκρον τῆς ῥάβδου αὐτοῦ.

\ch{48}
Ἐγένετο δὲ μετὰ τὰ ῥήματα ταῦτα, καὶ ἀπηγγέλη τῷ Ἰωσὴφ, ὅτι ὁ πατήρ σου ἐνοχλεῖται· καὶ ἀναλαβὼν τοὺς δύο υἱοὺς αὐτοῦ τὸν Μανασσῆ καὶ τὸν Ἐφραὶμ, ἦλθε πρὸς Ἰακώβ.
\vs{2}Ἀπηγγέλη δὲ τῷ Ἰακὼβ, λέγοντες, ἰδοὺ ὁ υἱός σου Ἰωσὴφ ἔρχεται πρὸς σέ· καὶ ἐνισχύσας Ἰσραὴλ ἐκάθισεν ἐπὶ τὴν κλίνην.
\vs{3}Καὶ εἶπεν Ἰακὼβ τῷ Ἰωσὴφ, ὁ Θεός μου ὤφθη μοι ἐν Λουζᾷ ἐν γῇ Χαναὰν, καὶ εὐλόγησέ με,
\vs{4}καὶ εἶπέ μοι, ἰδοὺ ἐγὼ αὐξανῶ σε, καὶ πληθυνῶ σε, καὶ ποιήσω σε εἰς συναγωγὰς ἐθνῶν· καὶ δώσω σοι τὴν γῆν ταύτην, καὶ τῷ σπέρματί σου μετὰ σὲ, εἰς κατάσχεσιν αἰώνιον.
\vs{5}Νῦν οὖν οἱ δύο υἱοί σου, οἱ γενόμενοί σοι ἐν γῇ Αἰγύπτῳ πρὸ τοῦ με ἐλθεῖν πρός σε εἰς Αἴγυπτον, ἐμοί εἰσιν, Ἐφραὶμ καὶ Μανασσῆ· ὡς Ῥουβὴν καὶ Συμεὼν ἔσονταί μοι.
\vs{6}Τὰ δὲ ἔκγονα, ἃ ἐὰν γεννήσῃς μετὰ ταῦτα, ἔσονται ἐπὶ τῷ ὀνόματι τῶν ἀδελφῶν αὐτῶν· κληθήσονται ἐπὶ τοῖς ἐκείνων κλήροις.
\vs{7}Ἐγὼ δὲ ἡνίκα ἠρχόμην ἐκ Μεσοποταμίας τῆς Συρίας, ἀπέθανε Ῥαχὴλ ἡ μήτηρ σου ἐν γῇ Χαναὰν, ἐγγίζοντός μου κατὰ τὸν ἱππόδρομον Χαβραθὰ τῆς γῆς, τοῦ ἐλθεῖν Ἐφραθά· καὶ κατώρυξα αὐτὴν ἐν τῇ ὁδῷ τοῦ ἱπποδρόμου· αὕτη ἐστὶ Βηθλεέμ.

\vs{8}Ἰδὼν δὲ Ἰσραὴλ τοὺς υἱοὺς Ἰωσὴφ, εἶπε, τίνες σοι οὗτοι;
\vs{9}Εἶπε δὲ Ἰωσὴφ τῷ πατρὶ αὐτοῦ, υἱοί μου εἰσὶν, οὓς ἔδωκε μοι ὁ Θεὸς ἐνταῦθα. Καὶ εἶπεν Ἰακὼβ, προσάγαγέ μοι αὐτοὺς, ἵνα εὐλογήσω αὐτούς.
\vs{10}Οἱ ὀφθαλμοὶ δὲ Ἰσραὴλ ἐβαρυώπησαν ἀπὸ τοῦ γήρως, καὶ οὐκ ἠδύνατο βλέπειν· καὶ ἤγγισεν αὐτοὺς πρὸς αὐτὸν, καὶ ἐφίλησεν αὐτοὺς, καὶ περιέλαβεν αὐτους.
\vs{11}Καὶ εἶπεν Ἰσραὴλ πρὸς Ἰωσὴφ, ἰδοὺ τοῦ προσώπου σου οὐκ ἐστερήθην, καὶ ἰδοὺ ἔδειξέ μοι ὁ Θεὸς καὶ τὸ σπέρμα σου.
\vs{12}Καὶ ἐξήγαγεν αὐτοὺς Ἰωσὴφ ἀπὸ τῶν γονάτων αὐτοῦ· καὶ προσεκύνησαν αὐτῷ ἐπὶ πρόσωπον ἐπὶ τῆς γῆς.
\vs{13}Λαβὼν δὲ Ἰωσὴφ τοὺς δύο υἱοὺς αὐτοῦ, τόν τε Ἐφραὶμ ἐν τῇ δεξιᾷ, ἐξ ἀριστερῶν δὲ Ἰσραὴλ, τὸν δὲ Μανασσῆ ἐξ ἀριστερῶν, ἐκ δεξιῶν δὲ Ἰσραὴλ, ἤγγισεν αὐτοὺς αὐτῷ.
\vs{14}Ἐκτείνας δὲ Ἰσραὴλ τὴν χεῖρα τὴν δεξιὰν, ἐπέβαλεν ἐπὶ τὴν κεφαλὴν Ἐφραὶμ, οὗτος δὲ ἦν ὁ νεώτερος, καὶ τὴν ἀριστερὰν ἐπὶ τὴν κεφαλὴν Μανασσῆ, ἐναλλὰξ τὰς χεῖρας.

\vs{15}Καὶ εὐλόγησεν αὐτοὺς, καὶ εἶπεν, ὁ Θεὸς, ᾧ εὐηρέστησαν οἱ πατέρες μου ἐνώπιον αὐτοῦ, Ἁβραὰμ καὶ Ἰσαὰκ, ὁ Θεὸς ὁ τρέφων με ἐκ νεότητος ἕως τῆς ἡμέρας ταύτης,
\vs{16}ὁ Ἄγγελος ὁ ῥυόμενός με ἐκ πάντων τῶν κακῶν, εὐλογήσαι τὰ παιδία ταῦτα· καὶ ἐπικληθήσεται ἐν αὐτοις τὸ ὄνομά μου, καὶ τὸ ὄνομα τῶν πατέρων μου Ἁβραὰμ καὶ Ἰσαάκ· καὶ πληθυνθείησαν εἰς πλῆθος πολὺ ἐπὶ τῆς γῆς.
\vs{17}Ἰδὼν δὲ Ἰωσὴφ ὅτι ἐπέβαλεν ὁ πατὴρ αὐτοῦ τὴν χεῖρα τὴν δεξιὰν αὐτοῦ ἐπὶ τὴν κεφαλὴν Ἐφραὶμ, βαρὺ αὐτῷ κατεφάνη· καὶ ἀντελάβετο Ἰωσὴφ τῆς χειρὸς τοῦ πατρὸς αὐτοῦ, ἀφελεῖν αὐτὴν ἀπὸ τῆς κεφαλῆς Ἐφραὶμ ἐπὶ τὴν κεφαλὴν Μανασσῆ.
\vs{18}Εἶπε δὲ Ἰωσὴφ τῷ πατρὶ αὐτοῦ, οὐχ οὕτως, πατὴρ, οὗτος γὰρ ὁ πρωτότοκος· ἐπίθες τὴν δεξιάν σου ἐπὶ τὴν κεφαλὴν αὐτοῦ.
\vs{19}Καὶ οὐκ ἠθέλησεν, ἀλλʼ εἶπεν, οἶδα, τέκνον, οἶδα· καὶ οὗτος ἔσται εἰς λαὸν, καὶ οὗτος ὑψωθήσεται· ἀλλʼ ὁ ἀδελφὸς αὐτοῦ ὁ νεώτερος μείζον αὐτοῦ ἔσται, καὶ τὸ σπέρμα αὐτοῦ ἔσται εἰς πλῆθος ἐθνῶν.
\vs{20}Καὶ εὐλόγησεν αὐτοὺς ἐν τῇ ἡμέρᾳ ἐκείνῃ, λέγων, ἐν ὑμῖν εὐλογηθήσεται Ἰσραὴλ, λέγοντες, ποιήσαι σε ὁ Θεὸς ὡς Ἐφραὶμ καὶ ὡς Μανασσῆ· καὶ ἔθηκε τὸν Ἐφραὶμ ἔμπροσθεν τοῦ Μανασσῆ.
\vs{21}Εἶπε δὲ Ἰσραὴλ τῷ Ἰωσὴφ, ἰδοὺ ἐγὼ ἀποθνήσκω· καὶ ἔσται ὁ Θεὸς μεθʼ ὑμῶν, καὶ ἀποστρέψει ὑμᾶς εἰς τὴν γῆν τῶν πατέρων ὑμῶν.
\vs{22}Ἐγὼ δὲ δίδωμί σοι Σίκιμα ἐξαίρετον ὑπὲρ τοὺς ἀδελφούς σου, ἣν ἔλαβον ἐκ χειρὸς Ἀμοῤῥαίων ἐν μαχαίρᾳ μου καὶ τόξῳ.

\ch{49}
Ἐκάλεσε δὲ Ἰακὼβ τοὺς υἱοὺς αὐτοῦ, καὶ εἶπεν αὐτοῖς, συνάχθητε, ἵνα ἀναγγείλω ὑμῖν, τί ἀπαντήσει ὑμῖν ἐπʼ ἐσχάτων τῶν ἡμέρων.
\vs{2}Συνάχθητε, καὶ ἀκούσατέ μου, υἱοὶ Ἰακώβ· ἀκούσατε Ἰσραὴλ, ἀκούσατε τοῦ πατρὸς ὑμῶν.
\vs{3}Ῥουβὴν πρωτότοκός μου, σὺ ἰσχύς μου, καὶ ἀρχὴ τέκνων μου, σκληρὸς φέρεσθαι, καὶ σκληρὸς αὐθάδης.
\vs{4}Ἐξύβρισας ὡς ὕδωρ, μὴ ἐκζέσῃς, ἀνέβης γὰρ ἐπὶ τὴν κοίτην τοῦ πατρός σου· τότε ἐμίανας τὴν στρωμνὴν, οὗ ἀνέβης.
\vs{5}Συμεὼν καὶ Λευὶ ἀδελφοὶ συνετέλεσαν ἀδικίαν ἐξαιρέσεως αὐτῶν·
\vs{6}Εἰς βουλὴν αὐτῶν μὴ ἔλθοι ἡ ψυχή μου, καὶ ἐπὶ τῇ συστάσει αὐτῶν μὴ ἐρίσαι τὰ ἥπατά μου· ὅτι ἐν τῷ θυμῷ αὐτῶν ἀπέκτειναν ἀνθρώπους, καὶ ἐν τῇ ἐπιθυμίᾳ αὐτῶν ἐνευροκόπησαν ταῦρον.
\vs{7}Ἐπικατάρατος ὁ θυμὸς αὐτὼν, ὅτι αὐθάδης· καὶ ἡ μῆνις αὐτῶν, ὅτι ἐσκληρύνθη· διαμεριῷ αὐτοὺς ἐν Ἰακὼβ, καὶ διασπερῷ αὐτοὺς ἐν Ἰσραήλ.
\vs{8}Ἰούδα, σὲ αἰνέσαισαν οἱ ἀδελφοί σου· αἱ χεῖρές σου ἐπὶ νώτου τῶν ἐχθρῶν σου· προσκυνήσουσί σοι οἱ υἱοὶ τοῦ πατρός σου.
\vs{9}Σκύμνος λέοντος Ἰούδα· ἐκ βλαστοῦ, υἱέ μου, ἀνέβης· ἀναπεσὼν ἐκοιμήθης ὡς λέων καὶ ὡς σκύμνος· τίς ἐγερεῖ αὐτόν;
\vs{10}Οὐκ ἐκλείψει ἄρχων ἐξ Ἰούδα, καὶ ἡγούμενος ἐκ τῶν μηρῶν αὐτοῦ, ἕως ἐὰν ἔλθῃ τὰ ἀποκείμενα αὐτῷ· καὶ αὐτὸς προσδοκία ἐθνῶν.
\vs{11}Δεσμεύων πρὸς ἄμπελον τὸν πῶλον αὐτοῦ, καὶ τῇ ἕλικι τὸν πῶλον τῆς ὄνου αὐτοῦ, πλυνεῖ ἐν οἴνῳ τὴν στολὴν αὐτοῦ, καὶ ἐν αἵματι σταφυλῆς τὴν περιβολὴν αὐτοῦ.
\vs{12}Χαροποιοὶ οἱ ὀφθαλμοὶ αὐτοῦ ὑπὲρ οἶνον· καὶ λευκοὶ οἱ ὀδόντες αὐτοῦ ἢ γάλα.
\vs{13}Ζαβουλὼν παράλιος κατοικήσει καὶ αὐτὸς παρʼ ὅρμον πλοίων, καὶ παρατενεῖ ἕως Σιδῶνος.
\vs{14}Ἰσσάχαρ τὸ καλὸν ἐπεθύμησεν, ἀναπαυόμενος ἀνὰ μέσον τῶν κλήρων.
\vs{15}Καὶ ἰδὼν τὴν ἀνάπαυσιν ὅτι καλὴ, καὶ τὴν γῆν ὅτι πίων, ὑπέθηκε τὸν ὦμον αὐτοῦ εἰς τὸ πονεῖν, καὶ ἐγενήθη ἀνὴρ γεωργός.
\vs{16}Δὰν κρινεῖ τὸν λαὸν αὐτοῦ, ὡσεὶ καὶ μία φυλὴ ἐν Ἰσραήλ.
\vs{17}Καὶ γενηθητω Δὰν ὄφις ἐφʼ ὁδοῦ, ἐγκαθήμενος ἐπὶ τρίβου, δάκνων πτέρναν ἵππου· καὶ πεσεῖται ὁ ἱππεὺς εἰς τὰ ὀπίσω,
\vs{18}τὴν σωτηρίαν περιμένων Κυρίου.
\vs{19}Γὰδ, πειρατήριον πειρατεύσει αὐτόν· αὐτὸς δὲ πειράτεύσει αὐτὸν κατὰ πόδας.
\vs{20}Ἀσὴρ, πίων αὐτοῦ ὁ ἄρτος· καὶ αὐτὸς δώσει τρυφὴν ἄρχουσι.
\vs{21}Νεφθαλὶ στέλεχος ἀνειμένον, ἐπιδιδοὺς ἐν τῷ γεννήματι κάλλος.
\vs{22}Υἱὸς ηὐξημένος Ἰωσὴφ, υἱὸς ηὐξημένος μου ζηλωτὸς, υἱός μου νεώτατος· πρός με ἀνάστρεψον.
\vs{23}Εἰς ὃν διαβουλευόμενοι ἐλοιδόρουν, καὶ ἐνεῖχον αὐτῷ κύριοι τοξευμάτων.
\vs{24}Καὶ συνετρίβη μετὰ κράτους τὰ τόξα αὐτῶν· καὶ ἐξελύθη τὰ νεῦρα βραχιόνων χειρὸς αὐτῶν, διὰ χεῖρα δυνάστου Ἰακώβ· ἐκεῖθεν ὁ κατισχύσας Ἰσραὴλ παρὰ Θεοῦ τοῦ πατρός σου.
\vs{25}Καὶ ἐβοήθησέ σοι ὁ Θεὸς ὁ ἐμὸς, καὶ εὐλόγησέ σε εὐλογίαν οὐρανοῦ ἄνωθεν, καὶ εὐλογίαν γῆς ἐχούσης πάντα, εἵνεκεν εὐλογίας μαστῶν καὶ μήτρας,
\vs{26}εὐλογίας πατρός σου καὶ μητρός σου· ὑπερίσχυσεν ὑπὲρ εὐλογίας ὀρέων μονίμων, καὶ ἐπʼ εὐλογίαις θινῶν ἀενάων· ἔσονται ἐπὶ κεφαλὴν Ἰωσὴφ, καὶ ἐπὶ κορυφῆς ὧν ἡγήσατο ἀδελφῶν.
\vs{27}Βενιαμὶν λύκος ἅρπαξ, τὸ πρωϊνὸν ἔδεται ἔτι, καὶ εἰς τὸ ἑσπέρας δίδωσι τροφήν.
\vs{28}Πάντες οὕτοι υἱοὶ Ἰακὼβ δώδεκα· καὶ ταῦτα ἐλάλησεν αὐτοῖς ὁ πατὴρ αὐτῶν· καὶ εὐλόγησεν αὐτούς· ἕκαστον κατὰ τὴν εὐλογίαν αὐτοῦ εὐλόγησεν αὐτούς.
\vs{29}Καὶ εἶπεν αὐτοῖς, ἐγὼ προστίθεμαι πρὸς τὸν ἐμὸν λαόν· θάψτέ με μετὰ τῶν πατέρων μου ἐν τῷ σπηλαίῳ, ὅ ἐστιν ἐν τῷ ἀγρῷ Ἐφρὼν τοῦ Χετταίου,
\vs{30}ἐν τῷ σπηλαίῳ τῷ διπλῷ, τῷ ἀπέναντι Μαμβρῆ, ἐν γῇ Χαναὰν, ὃ ἐκτήσατο Ἁβραὰμ τὸ σπήλαιον παρὰ Ἐφρὼν τοῦ Χετταίου ἐν κτήσει μνημείου.
\vs{31}Ἐκεῖ ἔθαψαν Ἁβραὰμ καὶ Σάῤῥαν τὴν γυναῖκα αὐτοῦ· ἐκεῖ ἔθαψαν Ἰσαὰκ καὶ Ῥεβέκκαν τὴν γυναῖκα αὐτοῦ· ἐκεῖ ἔθαψαν Λείαν·
\vs{32}Ἐν κτήσει τοῦ ἀγροῦ καὶ τοῦ σπηλαίου τοῦ ὄντος ἐν αὐτῷ, παρὰ τῶν υἱῶν Χέτ.
\vs{33}Καὶ κατέπαυσεν Ἰακὼβ ἐπιτάσσων τοῖς υἱοῖς αὐτοῦ· καὶ ἐξᾴρας τοὺς πόδας αὐτοῦ ἐπὶ τὴν κλίνην, ἐξέλιπε· καὶ προσετέθη πρὸς τὸν λαὸν αὐτοῦ.

\ch{50}
Καὶ ἐπιπεσὼν Ἰωσὴφ ἐπὶ πρόσωπον τοῦ πατρὸς αὐτοῦ ἔκλαυσεν αὐτὸν, καὶ ἐφίλησεν αὐτόν.
\vs{2}Καὶ προσέταξεν Ἰωσὴφ τοῖς παισὶν αὐτοῦ τοῖς ἐνταφιασταῖς, ἐνταφιάσαι τὸν πατέρα αὐτοῦ· καὶ ἐνεταφίασαν οἱ ἐνταφιασταὶ τὸν Ἰσραήλ.
\vs{3}Καὶ ἐπλήρωσαν αὐτοῦ τεσσαράκοντα ἡμέρας· οὕτω γὰρ καταριθμοῦνται αἱ ἡμέραι τῆς ταφῆς· καὶ ἐπένθησεν αὐτὸν Αἴγυπτος ἑβδομήκοντα ἡμέρας.
\vs{4}Ἐπεὶ δὲ παρῆλθον αἱ ἡμέραι τοῦ πένθους, ἐλάλησεν Ἰωσὴφ πρὸς τοὺς δυνάστας Φαραὼ, λέγων, εἰ εὗρον χάριν ἐναντίον ὑμῶν, λαλήσατε περὶ ἐμοῦ εἰς τὰ ὦτα Φαραὼ, λέγοντες,
\vs{5}ὁ πατήρ μου ὥρκισέ με, λέγων, ἐν τῷ μνημείῳ, ᾧ ὤρυξα ἐμαυτῷ ἐν γῇ Χαναὰν, ἐκεῖ με θάψεις· νῦν οὖν ἀναβὰς· θάψω τὸν πατέρα μου, καὶ ἐπανελεύσομαι·
\vs{6}Καὶ εἶπε Φαραὼ τῷ Ἰωσὴφ, ἀνάβηθι, θάψον τὸν πατέρα σου, καθάπερ ὥρκισέ σε.
\vs{7}Καὶ ἀνέβη Ἰωσὴφ θάψαι τὸν πατέρα αὐτοῦ· καὶ συνανέβησαν μετʼ αὐτοῦ πάντες οἱ παῖδες Φαραὼ, καὶ οἱ πρεσβύτεροι τοῦ οἴκου αὐτοῦ, καὶ πάντες οἱ πρεσβύτεροι τῆς γῆς Αἰγύπτου,
\vs{8}καὶ πᾶσα ἡ πανοικία Ἰωσὴφ, καὶ οἱ ἀδελφοὶ αὐτοῦ, καὶ πᾶσα ἡ οἰκία ἡ πατρικὴ αὐτοῦ, καὶ ἡ συγγένεια αὐτοῦ· καὶ τὰ πρόβατα, καὶ τοὺς βόας ὑπελίποντο ἐν γῇ Γεσέμ.
\vs{9}Καὶ συνανέβησαν μετʼ αὐτοῦ καὶ ἅρματα καὶ ἱππεῖς, καὶ ἐγένετο ἡ παρεμβολὴ μεγάλη σφόδρα.
\vs{10}Καὶ παρεγένοντο εἰς ἅλωνα Ἀτὰδ, ὅ ἐστι πέραν τοῦ Ἰορδάνου· καὶ ἐκόψαντο αὐτὸν κοπετὸν μέγαν καὶ ἰσχυρὸν σφόδρα· καὶ ἐποίησε τὸ πένθος τῷ πατρὶ αὐτοῦ ἑπτὰ ἡμέρας.
\vs{11}Καὶ εἶδον οἱ κάτοικοι τῆς γῆς Χαναὰν τὸ πένθος ἐπὶ ἅλωνι Ἀτὰδ, καὶ εἶπαν, πένθος μέγα τοῦτό ἐστι τοῖς Αἰγυπτίοις· διὰ τοῦτο ἐκάλεσε τὸ ὄνομα αὐτοῦ, Πένθος Αἰγύπτου, ὅ ἐστι πέραν τοῦ Ἰορδάνου.
\vs{12}Καὶ ἐποίησαν αὐτῷ οὕτως οἱ υἱοὶ αὐτοῦ.
\vs{13}Καὶ ἀνέλαβον αὐτὸν οἱ υἱοὶ αὐτοῦ εἰς γῆν Χαναάν· καὶ ἔθαψαν αὐτὸν εἰς τὸ σπήλαιον τὸ διπλοῦν, ὃ ἐκτήσατο Ἁβραὰμ τὸ σπήλαιον ἐν κτήσει μνημείου παρὰ Ἐφρὼν τοῦ Χετταίου, κατέναντι Μαμβρή.
\vs{14}Καὶ ὑπέστρεψεν Ἰωσὴφ εἰς Αἴγυπτον, αὐτὸς καὶ οἱ ἀδελφοὶ αὐτοῦ, καὶ οἱ συναναβάντες θάψαι τὸν πατέρα αὐτοῦ.

\vs{15}Ἰδόντες δὲ οἱ ἀδελφοὶ Ἰωσὴφ, ὅτι τέθνηκεν ὁ πατὴρ αὐτῶν, εἶπαν, μή ποτε μνησικακήσῃ ἡμῖν Ἰωσὴφ, καὶ ἀνταπόδομα ἀνταποδῷ ἡμῖν πάντα τὰ κακὰ, ἃ ἐνεδειξάμεθα εἰς αὐτὸν.
\vs{16}Καὶ παραγενόμενοι πρὸς Ἰωσὴφ εἶπαν, ὁ πατήρ σου ὥρκισε πρὸ τοῦ τελευτῆσαι αὐτὸν, λέγων,
\vs{17}οὕτως εἴπατε Ἰωσήφ· ἄφες αὐτοῖς τὴν ἀδικίαν καὶ τὴν ἁμαρτίαν αὐτῶν, ὅτι πονηρά σοι ἐνεδείξαντο· καὶ νῦν δέξαι τὴν ἀδικίαν τῶν θεραπόντων τοῦ Θεοῦ τοῦ πατρός σου· καὶ ἔκλαυσεν Ἰωσὴφ λαλούντων αὐτῶν πρὸς αὐτόν.
\vs{18}Καὶ ἐλθόντες πρὸς αὐτὸν εἶπαν, οἵδε ἡμεῖς σοὶ οἰκέται.
\vs{19}Καὶ εἶπεν αὐτοῖς Ἰωσὴφ, μὴ φοβεῖσθε, τοῦ γὰρ Θεοῦ εἰμι ἐγώ.
\vs{20}Ὑμεῖς ἐβουλεύσασθε κατʼ ἐμοῦ εἰς πονηρὰ, ὁ δὲ Θεὸς ἐβουλεύσατο περὶ ἐμοῦ εἰς ἀγαθὰ, ὅπως ἂν γενηθῇ ὡς σήμερον, καὶ τραφῇ λαὸς πολύς.
\vs{21}Καὶ εἶπεν αὐτοῖς, μὴ φοβεῖσθε· ἐγὼ διαθρέψω ὑμᾶς, καὶ τὰς οἰκίας ὑμῶν· καὶ παρεκάλεσεν αὐτοὺς, καὶ ἐλάλησεν αὐτῶν εἰς τὴν καρδίαν.
\vs{22}Καὶ κατῴκησεν Ἰωσὴφ ἐν Αἰγύπτῳ, αὐτὸς καὶ οἱ ἀδελφοὶ αὐτοῦ, καὶ πᾶσα ἡ πανοικία τοῦ πατρὸς αὐτοῦ· καὶ ἔζησεν Ἰωσὴφ ἔτη ἑκατὸν δέκα.
\vs{23}Καὶ εἶδεν Ἰωσὴφ Ἐφραὶμ παιδία, ἕως τρίτης γενεᾶς· καὶ οἱ υἱοὶ Μαχεὶρ τοῦ υἱοῦ Μανασσῆ ἐτέχθησαν ἐπὶ μηρῶν Ἰωσήφ.
\vs{24}Καὶ εἶπεν Ἰωσὴφ τοῖς ἀδελφοῖς αὐτοῦ, λέγων, ἐγὼ ἀποθνήσκω· ἐπισκοπῇ δὲ ἐπισκέψεται ὁ Θεὸς ὑμᾶς, καὶ ἀνάξει ὑμᾶς ἐκ τῆς γῆς ταύτης εἰς τὴν γῆν, ἣν ὤμοσεν ὁ Θεὸς τοῖς πατράσιν ἡμῶν, Ἁβραὰμ, Ἰσαὰκ, καὶ Ἰακώβ.
\vs{25}Καὶ ὥρκισεν Ἰωσὴφ τοὺς υἱοὺς Ἰσραὴλ, λέγων, ἐν τῇ ἐπισκοπῇ ᾗ ἐπισκέψηται ὁ Θεὸς ὑμᾶς, καὶ συνανοίσετε τὰ ὀστᾶ μου ἐντεῦθεν μεθʼ ὑμῶν.
\vs{26}Καὶ ἐτελεύτησεν Ἰωσὴφ ἐτῶν ἑκατὸν δέκα· καὶ ἔθαψαν αὐτὸν, καὶ ἔθηκαν ἐν τῇ σορῷ ἐν Αἰγύπτῳ.


\def\book{ΕΞΟΔΟΣ}
\biblebook{ΕΞΟΔΟΣ}


\lettrine[lines=2, loversize=0.2, nindent=0em, findent=.25em]{\textcolor{bookheadingcolor}{Τ}}{ΑΥΤΑ} τὰ ὀνόματα τῶν υἱῶν Ἰσραὴλ τῶν εἰσπεπορευμένων εἰς Αἴγυπτον ἅμα Ἰακὼβ τῷ πατρὶ αὐτῶν, ἕκαστος πανοικὶ αὐτῶν εἰσήλθοσαν.
\vs{2}Ῥουβὴν, Συμεών, Λευὶ, Ἰούδας,
\vs{3}Ἰσσάχαρ, Ζαβουλὼν, Βενιαμὶν,
\vs{4}Δὰν, καὶ Νεφθαλὶ, Γὰδ, καὶ Ἀσήρ.
\vs{5}Ἰωσὴφ δὲ ἦν ἐν Αἰγύπτῳ· ἦσαν δὲ πᾶσαι ψυχαὶ ἐξ Ἰακὼβ, πέντε καὶ ἑβδομήκοντα.
\vs{6}Ἐτελεύτησε δὲ Ἰωσὴφ, καὶ πάντες οἱ ἀδελφοὶ αὐτοῦ, καὶ πᾶσα ἡ γενεὰ ἐκείνη.
\vs{7}Οἱ δὲ υἱοὶ Ἰσραὴλ ηὐξήθησαν, καὶ ἐπληθύνθησαν, καὶ χυδαῖοι ἐγένοντο, καὶ κατίσχυον σφόδρα σφόδρα· ἐπλήθυνε δὲ ἡ γῆ αὐτούς.
\vs{8}Ἀνέστη δὲ βασιλεὺς ἕτερος ἐπʼ Αἴγυπτον, ὃς οὐκ ᾔδει τὸν Ἰωσήφ.
\vs{9}Εἶπε δὲ τῷ ἔθνει αὐτοῦ, ἰδοὺ τὸ γένος τῶν υἱῶν Ἰσραὴλ μέγα πλῆθος, καὶ ἰσχύει ὑπὲρ ἡμᾶς.
\vs{10}Δεῦτε οὖν κατασοφισώμεθα αὐτοὺς, μήποτε πληθυνθῇ, καὶ ἡνίκα ἂν συμβῇ ἡμῖν πόλεμος, προστεθήσονται καὶ οὗτοι πρὸς τοὺς ὑπεναντίους, καὶ ἐκπολεμήσαντες ἡμᾶς, ἐξελεύσονται ἐκ τῆς γῆς.
\vs{11}Καὶ ἐπέστησεν αὐτοῖς ἐπιστάτας τῶν ἔργων, ἵνα κακώσωσιν αὐτοὺς ἐν τοῖς ἔργοις. Καὶ ᾠκοδόμησαν πόλεις ὀχυρὰς τῷ Φαραῷ, τήν τε Πειθὼ, καὶ Ῥαμεσσῆ, καὶ Ὢν, ἥ ἐστιν Ἡλιού πολις.
\vs{12}Καθότι δὲ αὐτοὺς ἐταπείνουν, τοσούτῳ πλείους ἐγίνοντο, καὶ ἴσχυον σφόδρα σφόδρα· καὶ ἐβδελύσσοντο οἱ Αἰγύπτιοι ἀπὸ τῶν υἱῶν Ἰσραήλ.
\vs{13}Καὶ κατεδυνάστευον οἱ Αἰγύπτιοι τοὺς υἱοὺς Ἰσραὴλ βίᾳ.
\vs{14}Καὶ κατωδύνων αὐτῶν τὴν ζωὴν ἐν τοῖς ἔργοις τοῖς σκληροῖς, τῷ πηλῷ καὶ τῇ πλινθείᾳ, καὶ πᾶσι τοῖς ἔργοις τοῖς ἐν τοῖς πεδίοις, κατὰ πάντα τὰ ἔργα, ὧν κατεδουλοῦντο αὐτοὺς μετὰ βίας.

\vs{15}Καὶ εἶπεν ὁ βασιλεὺς τῶν Αἰγυπτίων ταῖς μαίαις τῶν Ἐβραίων, τῇ μιᾷ αὐτῶν ὄνομα Σεπφώρα, καὶ τὸ ὄνομα τῆς δευτέρας Φουά·
\vs{16}Καὶ εἶπεν, ὅταν μαιοῦσθε τὰς Ἐβραίας, καὶ ὦσι πρὸς τῷ τίκτειν, ἐὰν μὲν ἄρσεν ᾖ, ἀποκτείνατε αὐτό· ἐὰν δὲ θῆλυ, περιποιεῖσθε αὐτό.
\vs{17}Ἐφοβήθησαν δὲ αἱ μαῖαι τὸν Θεὸν, καὶ οὐκ ἐποίησαν καθότι συνέταξεν αὐταῖς ὁ βασιλεὺς Αἰγύπτου, καὶ ἐζωογόνουν τὰ ἄρσενα.
\vs{18}Ἐκάλεσε δὲ ὁ βασιλεὺς Αἰγύπτου τὰς μαίας, καὶ εἶπεν αὐταῖς, τί ὅτι ἐποιήσατε τὸ πρᾶγμα τοῦτο, καὶ ἐζωογονεῖτε τὰ ἄρσενα;
\vs{19}Εἶπαν δὲ αἱ μαῖαι τῷ Φαραῷ, οὐχ ὡς γυναῖκες Αἰγύπτου αἱ Ἐβραῖαι· τίκτουσι γὰρ πρὶν ἢ εἰσελθεῖν πρὸς αὐτὰς τὰς μαίας· καὶ ἔτικτον.
\vs{20}Εὖ δὲ ἐποίει ὁ Θεὸς ταῖς μαίαις· καὶ ἐπλήθυνεν ὁ λαὸς, καὶ ἴσχυε σφόδπα.
\vs{21}Ἐπεὶ δὲ ἐφοβοῦντο αἱ μαῖαι τὸν Θεὸν, ἐποίησαν ἑαυταῖς οἰκίας.
\vs{22}Συνέταξε δὲ Φαραὼ παντὶ τῷ λαῷ αὐτοῦ, λέγων, πᾶν ἄρσεν, ὃ ἐὰν τεχθῇ τοῖς Ἑβραίοις, εἰς τὸν ποταμὸν ῥίψατε, καὶ πᾶν θῆλυ, ζωογονεῖτε αὐτό.

\ch{2}
Ἦν δέ τις ἐκ τῆς φυλῆς Λευὶ, ὃς ἔλαβεν τῶν θυγατέρων Λευί.
\vs{2}Καὶ ἐν γαστρὶ ἔλαβε, καὶ ἔτεκεν ἄρσεν· ἰδόντες δὲ αὐτὸ ἀστεῖον, ἐσπέπασαν αὐτὸ μῆνας τρεῖς.
\vs{3}Ἐπεὶ δὲ οὐκ ἐδύναντο αὐτὸ ἔτι κρύπτειν, ἔλαβεν αὐτῷ ἡ μήτηρ αὐτοῦ θῖβιν, καὶ κατέχρισεν αὐτὴν ἀσφαλτοπίσσῃ, καὶ ἐνέβαλε τὸ παιδίον εἰς αὐτήν, καὶ ἔθηκεν αὐτὴν εἰς τὸ ἕλος παρὰ τὸν ποταμόν.
\vs{4}Καὶ κατεσκόπευεν ἡ ἀδελφὴ αὐτοῦ μακρόθεν, μαθεῖν τί τὸ ἀποβησόμενον αὐτῷ.

\vs{5}Κατέβη δὲ ἡ θυγάτηρ Φαραὼ λούσασθαι ἐπὶ τὸν ποταμὸν, καὶ αἱ ἅβραι αὐτῆς παρεπορεύοντο παρὰ τὸν ποταμόν· καὶ ἰδοῦσα τὴν θίβιν ἐν τῷ ἕλει, ἀποστείλασα τὴν ἅβραν, ἀνείλατο αὐτήν.
\vs{6}Ἀνοίξασα δὲ ὁρᾷ παιδίον κλαῖον ἐν τῇ θίβει· καὶ ἐφείσατο αὐτοῦ ἡ θυγάτηρ Φαραὼ, καὶ ἔφη, ἀπὸ τῶν παιδίων τῶν Ἐβραίων τοῦτο.
\vs{7}Καὶ εἶπεν ἡ ἀδελφὴ αὐτοῦ τῇ θυγατρὶ Φαραὼ, θέλεις καλέσω σοι γυναῖκα τροφεύουσαν ἐκ τῶν Ἐβραίων, καὶ θηλάσει σαι τὸ παιδὶον σοι τὸ παιδίον;
\vs{8}Ἡ δὲ εἶπεν ἡ θυγάτηρ Φαραὼ, πορεύου· ἐλθοῦσα δὲ νεᾶνις ἐκάλεσε τὴν μητέρα τοῦ παιδίου.
\vs{9}Εἶπεν δὲ πρὸς αὐτὴν ἡ θυγάτηρ Φαραὼ, διατήρησόν μοι τὸ παιδίον τοῦτο, καὶ θήλασόν μοι αὐτὸ, ἐγὼ δὲ δώσω σοι τὸν μισθόν· ἔλαβε δὲ ἡ γυνὴ τὸ παιδίον, καὶ ἐθήλαζεν αὐτό.
\vs{10}Ἁδρυνθέντος δὲ τοῦ παιδίου, εἰσήγαγεν αὐτὸ πρὸς τὴν θυγατέρα Φαραὼ, καὶ ἐγενήθη αὐτῇ εἰς υἱόν· ἐπωνόμασε δὲ τὸ ὄνομα αὐτοῦ Μωυσῆν, λέγουσα, ἐκ τοῦ ὕδατος αὐτὸν ἀνειλόμην.

\vs{11}Ἐγένετο δὲ ἐν ταῖς ἡμέραις ταῖς πολλαῖς ἐκείναις μέγας γενόμενος Μωυσῆς, ἐξῆλθε πρὸς τοὺς ἀδελφοὺς αὐτοῦ τοὺς υἱοὺς Ἰσραήλ· κατανοήσας δὲ τὸν πόνον αὐτῶν, ὁρᾷ ἄνθρωπον Αἰγύπτιον τύπτοντα τινὰ Ἐβραῖον, τῶν ἑαυτοῦ ἀδελφῶν τῶν υἱῶν Ἰσραήλ.
\vs{12}Περιβλεψάμενος δὲ ὧδε καὶ ὧδε οὐχ ὁρᾷ οὐδένα, καὶ πατάξας τὸν Αἰγύπτιον, ἔκρυψεν αὐτὸν ἐν τῇ ἄμμῳ.
\vs{13}Ἐξελθὼν δὲ τῇ ἡμέρᾳ τῇ δευτέρᾳ, ὁρᾷ δύο ἄνδρας Ἐβραίους διαπληκτιζομένους· καὶ λέγει τῷ ἀδικοῦντι, διὰ τί σὺ τύπτεις τὸν πλησίον;
\vs{14}Ὁ δὲ εἶπε, τίς σε κατέστησεν ἄρχοντα καὶ δικαστὴν ἐφʼ ἡμῶν; μὴ ἀνελεῖν με σὺ θέλεις, ὃν τρόπον ἀνεῖλες χθὲς τὸν Αἰγύπτιον; ἐφοβήθη δὲ Μωυσῆς, καὶ εἶπεν, εἰ οὕτως ἐμφανὲς γέγονε τὸ ῥῆμα τοῦτο.
\vs{15}Ἤκουσε δὲ Φαραὼ τὸ ῥῆμα τοῦτο, καὶ ἐζήτει ἀνελεῖν Μωυσῆν. Ἀνεχώρησε δὲ Μωυσῆς ἀπὸ προσώπου Φαραὼ, καὶ ᾤκησεν ἐν γῇ Μαδιάμ· ἐλθὼν δὲ εἰς γῆν Μαδιὰμ, ἐκάθισεν ἐπὶ τοῦ φρέατος.
\vs{16}Τῷ δὲ ἱερεῖ Μαδιὰμ ἦσαν ἑπτὰ θυγατέρες, ποιμαίνουσαι τὰ πρόβατα τοῦ πατρὸς αὐτῶν Ἰοθόρ· παραγενόμεναι δὲ ἤντλουν, ἕως ἔπλησαν τὰς δεξαμενάς, ποτίσαι τὰ πρόβατα τοῦ πατρὸς αὐτῶν Ἰοθόρ.
\vs{17}Παραγενόμενοι δὲ οἱ ποιμένες ἐξέβαλλον αὐτάς· ἀναστὰς δὲ Μωυσῆς ἐῤῥύσατο αὐτὰς, καὶ ἤντλησεν αὐταῖς, καὶ ἐπότισε τὰ πρόβατα αὐτῶν.
\vs{18}Παρεγένοντο δὲ πρὸς Ῥαγουὴλ τὸν πατέρα αὐτῶν· ὁ δὲ εἶπεν αὐταῖς, διατί ἐταχύνατε τοῦ παραγενέσθαι σήμερον;
\vs{19}Αἱ δὲ εἶπαν, ἄνθρωπος Αἰγύπτιος ἐῤῥύσατο ἡμᾶς ἀπὸ τῶν ποιμένων, καὶ ἤντλησεν ἡμῖν, καὶ ἐπότισε τὰ πρόβατα ἡμῶν.
\vs{20}Ὁ δὲ εἶπε ταῖς θυγατράσιν αὐτοῦ, καὶ ποῦ ἐστιν; καὶ ἱνατί καταλελοίπατε τὸν ἄνθρωπον; καλέσατε οὖν αὐτὸν, ὅπως φάγῃ ἄρτον.
\vs{21}Κατῳκίσθη δὲ Μωυσῆς παρὰ τῷ ἀνθρώπῳ· καὶ ἐξέδοτο Σεπφώραν τὴν θυγατέρα αὐτοῦ Μωυσῇ γυναῖκα.
\vs{22}Ἐν γαστρὶ δὲ λαβοῦσα ἡ γυνὴ ἔτεκεν υἱόν· καὶ ἐπωνόμασε Μωυσῆς τὸ ὄνομα αὐτοῦ Γηρσάμ, λέγων, ὅτι παροικός εἰμι ἐν γῇ ἀλλοτρίᾳ.
\vs{23}Μετὰ δὲ τὰς ἡμέρας τὰς πολλὰς ἐκείνας, ἐτελεύτησεν ὁ βασιλεὺς Αἰγύπτου, καὶ κατεστέναξαν οἱ υἱοὶ Ἰσραὴλ ἀπὸ τῶν ἔργων, καὶ ἀνεβόησαν· καὶ ἀνέβη ἡ βοὴ αὐτῶν πρὸς τὸν Θεὸν ἀπὸ τῶν ἔργων.
\vs{24}Καὶ εἰσήκουσεν ὁ Θεὸς τὸν στεναγμὸν αὐτῶν· καὶ ἐμνήσθη ὁ Θεὸς τῆς διαθήκης αὐτοῦ τῆς πρὸς Ἀβραὰμ, καὶ Ἰσαὰκ, καὶ Ἰακώβ.
\vs{25}Καὶ ἐπεῖδεν ὁ Θεὸς τοὺς υἱοὺς Ἰσραὴλ, καὶ ἐγνώσθη αὐτοῖς.

\ch{3}
Καὶ Μωυσῆς ἦν ποιμαίνων τὰ πρόβατα Ἰοθὸρ τοῦ γαμβροῦ αὐτοῦ, τοῦ ἱερέως Μαδιὰμ, καὶ ἤγαγεν τὰ πρόβατα ὑπὸ τὴν ἔρημον, καὶ ἦλθεν εἰς τὸ ὄρος Χωρήβ.
\vs{2}Ὤφθη δὲ αὐτῷ Ἄγγελος Κυρίου ἐν πυρὶ φλογὸς ἐκ τοῦ βάτου· καὶ ὁρᾷ ὅτι ὁ βάτος καίεται πυρί, ὁ δὲ βάτος οὐ κατεκαίετο.
\vs{3}Εἶπε δὲ Μωυσῆς, παρελθὼν ὄψομαι τὸ ὅραμα τὸ μέγα τοῦτο, ὅτι οὐ κατακαίεται ὁ βάτος.
\vs{4}Ὡς δὲ εἶδεν Κύριος ὅτι προσάγει ἰδεῖν, ἐκάλεσεν αὐτὸν Κύριος ἐκ τοῦ βάτου, λέγων, Μωυσῆ, Μωυσῆ· ὁ δὲ εἶπε, τί ἐστιν;
\vs{5}Ὁ δὲ εἶπε, μὴ ἐγγίσῃς ὧδε· λύσαι τὸ ὑπόδημα ἐκ τῶν ποδῶν σου, ὁ γὰρ τόπος. ἐν ᾧ σὺ ἕστηκας, γῆ ἁγία ἐστί.
\vs{6}Καὶ εἶπεν, ἐγώ εἰμι ὁ Θεὸς τοῦ πατρός σου, Θεὸς Ἁβραὰμ, καὶ Θεὸς Ἰσαὰκ, καὶ Θεὸς Ἰακώβ· ἀπέστρεψε δὲ Μωυσῆς τὸ πρόσωπον αὐτοῦ, εὐλαβεῖτο γὰρ κατεμβλέψαι ἐνώπιον τοῦ Θεοῦ.
\vs{7}Εἶπε δὲ Κύριος πρὸς Μωυσῆν, ἰδὼν εἶδον τὴν κάκωσιν τοῦ λαοῦ μου τοῦ ἐν Αἰγύπτῳ, καὶ τῆς κραυγῆς αὐτῶν ἀκήκοα ἀπὸ τῶν ἐργοδιωκτῶν· οἶδα γὰρ τὴν ὀδύνην αὐτων,
\vs{8}καὶ κατέβην ἐξελέσθαι αὐτοὺς ἐκ χειρὸς τῶν Αἰγυπτίων, καὶ ἐξαγαγεῖν αὐτοὺς ἐκ τῆς γῆς ἐκείνης, καὶ εἰσαγαγεῖν αὐτοὺς εἰς γῆν ἀγαθὴν καὶ πολλήν, εἰς γῆν ῥέουσαν γάλα καὶ μέλι, εἰς τὸν τόπον τῶν Χαναναίων, καὶ Χετταίων, καὶ Ἀμοῤῥαίων, καὶ Φερεζαίων, καὶ Γεργεσαὶων, καὶ Εὐαίων, καὶ Ἰεβουσαίων.
\vs{9}Καὶ νῦν ἰδοὺ κραυγὴ τῶν υἱῶν Ἰσραὴλ ἥκει πρὸς με· κᾀγὼ ἑώρακα τὸν θλιμμὸν, ὃν οἱ Αἰγύπτιοι θλίβουσιν αὐτούς.
\vs{10}Καὶ νῦν δεῦρο, ἀποστείλω σε πρὸς Φαραὼ βασιλέα Αἰγύπτου, καὶ ἐξάξεις τὸν λαόν μου τοὺς υἱοὺς Ἰσραὴλ ἐκ γῆς Αἰγύπτου.

\vs{11}Καὶ εἶπε Μωυσῆς πρὸς τὸν Θεὸν, τίς εἰμι ἐγὼ, ὅτι πορεύσομαι πρὸς Φαραὼ βασιλέα Αἰγύπτου, καὶ ὅτι ἐξάξω τοὺς υἱοὺς Ἰσραὴλ ἐκ γῆς Αἰγύπτου;
\vs{12}Εἶπε δὲ ὁ Θεὸς Μωυσῇ, λέγων, ὅτι ἔσομαι μετὰ σοῦ· καὶ τοῦτό σοι τὸ σημεῖον ὅτι ἐγώ σε ἐξαποστελῶ, ἐν τῷ ἐξαγαγεῖν σε τὸν λαόν μου ἐξ Αἰγύπτου, καὶ λατρεύσετε τῷ Θεῷ ἐν τῷ ὄρει τοῦτῳ.
\vs{13}Καὶ εἶπε Μωυσῆς πρὸς τὸν Θεὸν, ἰδοὺ ἐγὼ ἐξελεύσομαι πρὸς τοὺς υἱοὺς Ἰσραὴλ, καὶ ἐρῶ πρὸς αὐτοὺς, ὁ Θεὸς τῶν πατέρων ἡμῶν ἀπέσταλκέ με πρὸς ὑμᾶς· ἐρωτήσουσί με, τί ὄνομα αὐτῷ; τί ἐρῶ πρὸς αὐτούς;
\vs{14}Καὶ εἶπεν ὁ Θεὸς πρὸς Μωυσῆν, λέγων, ἐγώ εἰμι ὁ Ὤν· καὶ εἶπεν, οὕτως ἐρεῖς τοῖς υἱοῖς Ἰσραὴλ, ὁ Ὢν ἀπέσταλκέ με πρὸς ὑμᾶς.
\vs{15}Καὶ εἶπεν ὁ Θεὸς πάλιν πρὸς Μωυσῆν, οὕτως ἐρεῖς τοῖς υἱοῖς Ἰσραήλ, Κύριος ὁ Θεὸς τῶν πατέρων ἡμῶν, Θεὸς Ἀβραὰμ, καὶ Θεὸς Ἰσαὰκ, καὶ Θεὸς Ἰακὼβ, ἀπέσταλκέ με πρὸς ὑμᾶς· τοῦτό μου ἐστὶν ὄνομα αἰώνιον, καὶ μνημόσυνον γενεῶν γενεαῖς.
\vs{16}Ἐλθὼν οὐν συνάγαγε τὴν γερουσίαν τῶν υἱῶν Ἰσραὴλ, καὶ ἐρεῖς πρὸς αὐτοὺς, Κύριος ὁ Θεὸς τῶν πατέρων ἡμων ὦπταί μοι, Θεὸς Ἀβραὰμ, καὶ Θεὸς Ἰσαὰκ, καὶ Θεὸς Ἰακὼβ, λέγων, ἐπισκοπῇ ἐπέσκεμμαι ὑμᾶς, καὶ ὅσα συμβέβηκεν ὑμῖν ἐν Αἰγύπτῳ.
\vs{17}Καὶ εἶπεν, ἀναβιβάσω ὑμᾶς ἐκ τῆς κακώσεως τῶν Αἰγυπτίων, εἰς τὴν γῆν τῶν Χαναναίων, καὶ Χετταίων, καὶ Ἀμοῤῥαίων, καὶ Φερεζαίων, καὶ Γεργεσαίων, καὶ Εὑαίων, καὶ Ἰεβουσαίων, εἰς γῆν ῥέουσαν γάλα καὶ μέλι.
\vs{18}Καὶ εἰσακούσονταί σου τῆς φωνῆς· καὶ εἰσελεύσῃ σὺ, καὶ ἡ γερουσία Ἰσραὴλ, πρὸς Φαραὼ βασιλέα Αἰγύπτου, καὶ ἐρεῖς πρὸς αὐτὸν ὁ Θεὸς τῶν Ἑβραίων προσκέκληται ἡμᾶς· πορευσόμεθα οὖν ὁδὸν τριῶν ἡμερῶν εἰς τὴν ἔρημον, ἵνα θύσωμεν τῷ Θεῷ ἡμῶν.
\vs{19}Ἐγὼ δὲ οἶδα ὅτι οὐ προήσεται ὑμᾶς Φαραὼ βασιλεὺς Αἰγύπτου πορευθῆναι, ἐὰν μὴ μετὰ χειρὸς κραταιᾶς.
\vs{20}Καὶ ἐκτείνας τὴν χεῖρα, πατάξω τοὺς Αἰγυπτίους ἐν πᾶσι τοῖς θαυμασίοις μου, οἷς ποιήσω ἐν αὐτοῖς· καὶ μετὰ ταῦτα ἐξαποστελεῖ ὑμᾶς.
\vs{21}Καὶ δώσω χάριν τῷ λαῷ τούτῳ ἐναντίον τῶν Αἰγυπτίων· ὅταν δὲ ἀποτρέχητε, οὐκ ἀπελεύσεσθε κενοί·
\vs{22}Ἀλλὰ αἰτήσει γυνὴ παρὰ γείτονος καὶ συσκήνου αὐτῆς σκεύη ἀργυρᾶ, καὶ χρυσᾶ, καὶ ἱματισμόν· καὶ ἐπιθήσετε ἐπὶ τοὺς υἱοὺς ὑμῶν, καὶ ἐπὶ τὰς θυγατέρας ὑμῶν, καὶ σκυλεύσατε τοὺς Αἰγυπτίους.

\ch{4}
Ἀπεκρίθη δὲ Μωυσῆς, καὶ εἶπεν, ἐὰν μὴ πιστεύσωσί μοι, μηδὲ εἰσακούσωσι τῆς φωνῆς μου, ἐροῦσι γὰρ, ὅτι οὐκ ὦπταί σοι ὁ Θεὸς, τί ἐρῶ πρὸς αὐτούς;
\vs{2}Εἶπε δὲ αὐτῳ Κύριος, τί τοῦτό ἐστι τὸ ἐν τῇ χειρί σου; ὁ δὲ εἶπε, ῥάβδος.
\vs{3}Καὶ εἶπε, ῥίψον αὐτὴν ἐπὶ τὴν γῆν· καὶ ἔῤῥιψεν αὐτὴν ἐπὶ τὴν γῆν, καὶ ἐγένετο ὄφις· καὶ ἔφυγε Μωυσῆς ἀπʼ αὐτοῦ.
\vs{4}Καὶ εἶπε Κύριος πρὸς Μωυσῆν, ἔκτεινον τῆν χεῖρα, καὶ ἐπιλαβοῦ τῆς κέρκου· ἐκτείνας οὖν τὴν χεῖρα ἐπελάβετο τῆς κέρκου· καὶ ἐγένετο ῥάβδος ἐν τῇ χειρὶ αὐτοῦ.
\vs{5}Ἵνα πιστεύσωσί σοι, ὅτι ὦπταί σοι ὁ Θεὸς τῶν πατέρων αὐτῶν, Θεὸς Ἀβραὰμ, καὶ Θεὸς Ἰσαὰκ, καὶ Θεὸς Ἰακώβ.
\vs{6}Εἶπε δὲ αὐτῷ Κύριος πάλιν, εἰσένεγκον τὴν χεῖρά σου εἰς τὸν κόλπον σου· καὶ εἰσήνεγκε τὴν χεῖρα αὐτοῦ εἰς τὸν κόλπον αὐτοῦ· καὶ ἐξήνεγκεν τὴν χεῖρα αὐτοῦ ἐκ τοῦ κόλπου αὐτοῦ, καὶ ἐγενήθη ἡ χεὶρ αὐτοῦ ὡσεὶ χιών.
\vs{7}Καὶ εἶπεν πάλιν, εἰσένεγκον τὴν χεῖρά σου εἰς τὸν κόλπον σου· καὶ εἰσήνεγκε τὴν χεῖρα εἰς τὸν κόλπον αὐτοῦ· καὶ ἐξήνεγκεν αὐτὴν ἐκ τοῦ κόλπου αὐτοῦ, καὶ πάλιν ἀπεκατέστη εἰς τὴν χρόαν τῆς σαρκὸς αὐτῆς.
\vs{8}Ἐὰν δὲ μὴ πιστεύσωσί σοι, μηδὲ εἰσακούσωσι τῆς φωνῆς τοῦ σημείου τοῦ πρώτου, πιστεύσουσί σοι τῆς φωνῆς τοῦ σημείου τοῦ δετέρου.
\vs{9}Καὶ ἔσται ἐὰν μὴ πιστεύσωσί σοι τοῖς δυσὶ σημείοις τούτοις, μηδὲ εἰσακούσωσι τῆς φωνῆς σου, λήψῃ ἀπὸ τοῦ ὕδατος τοῦ ποταμοῦ, καὶ ἐκχεεῖς ἐπὶ τὸ ξηρόν· καὶ ἔσται τὸ ὕδωρ, ὃ ἐὰν λάβῃς ἀπὸ τοῦ ποταμοῦ, αἷμα ἐπὶ τοῦ ξηροῦ.
\vs{10}Εἶπε δὲ Μωυσῆς πρὸς Κύριον, δέομαι, Κύριε· οὐχ ἱκανός εἰμι πρὸ τῆς χθὲς οὐδὲ πρὸ τῆς τρίτης ἡμέρας, οὐδὲ ἀφʼ οὗ ἤρξω λαλεῖν τῷ θεράποντί σου· ἰσχνόφωνος καὶ βραδύγλωσσος ἐγώ εἰμι.
\vs{11}Εἶπε δὲ Κύριος πρὸς Μωυσῆν, τίς ἔδωκε στόμα ἀνθρώπῳ; καὶ τίς ἐποιήσε δύσκωφον καὶ κωφὸν, βλέποντα καὶ τυφλόν; οὐκ ἐγὼ ὁ Θεός;
\vs{12}Καὶ νῦν πορεύου, καὶ ἐγὼ ἀνοίξω τὸ στόμα σου, καὶ συμβιβάσω σε ὃ μέλλεις λαλῆσαι.
\vs{13}Καὶ εἶπε Μωυσῆς, δέομαι, Κύριε· προχείρισαι δυνάμενον ἄλλον, ὃν ἀποστελεῖς.
\vs{14}Καὶ θυμωθεὶς ὀργῇ Κύριος ἐπὶ Μωυσῆν, εἶπεν, οὐκ ἰδοὺ Ἀαρὼν ὁ ἄδελφός σου ὁ Λευίτης; ἐπίσταμαι ὅτι λαλῶν λαλήσει αὐτός σοι· καὶ ἰδοὺ αὐτὸς ἐξελεύσεται εἰς συνάντησίν σοι, καὶ ἰδών σε χαρήσεται ἐν ἑαυτῷ.
\vs{15}Καὶ ἐρεῖς πρὸς αὐτὸν, καὶ δώσεις τὰ ῥήματά μου εἰς τὸ στόμα αὐτοῦ, καὶ ἐγὼ ἀνοίξω τὸ στόμα σου καὶ τὸ στόμα αὐτοῦ, καί συμβιβάσω ὑμᾶς ἃ ποιήσετε.
\vs{16}Καὶ αὐτός σοι λαλήσει πρὸς τὸν λαὸν, καὶ αὐτὸς ἔσται σου στόμα· σὺ δὲ αὐτῷ ἔσῃ τὰ πρὸς τὸν Θεόν.
\vs{17}Καὶ τὴν ῥάβδον ταύτην, τὴν στραφεῖσαν εἰς ὄφιν, λήψῃ ἐν τῇ χειρί σου, ἐν ᾗ ποιήσεις ἐν αὐτῇ τὰ σημεῖα.

\vs{18}Ἐπορεύθη δὲ Μωυσῆς, καὶ ἀπέστρεψε πρὸς Ἰοθὸρ τὸν γαμβρὸν αὐτοῦ, καὶ λέγει, πορεύσομαι καὶ ἀποστρέψω πρὸς τοὺς ἀδελφούς μου τοὺς ἐν Αἰγύπτῳ, καὶ ὄψομαι εἰ ἔτι ζῶσι· καὶ εἶπεν Ἰοθὸρ Μωυσῇ, βάδιζε ὑγιαίνων· μετὰ δὲ τὰς ἡμέρας τὰς πολλὰς ἐκείνας ἐτελεύτησεν ὁ βασιλεὺς Αἰγύπτου.
\vs{19}Εἶπε δὲ Κύριος πρὸς Μωυσῆν ἐν Μαδιὰμ, βάδιζε, ἄπελθε εἰς Αἴγυπτον, τεθνήκασι γὰρ πάντες οἱ ζητοῦντες σου τὴν ψυχήν.
\vs{20}Ἀναλαβὼν δὲ Μωυσῆς τὴν γυναῖκα καὶ τὰ παιδία, ἀνεβίβασεν αὐτὰ ἐπὶ τὰ ὑποζύγια, καὶ ἐπέστρεψεν εἰς Αἴγυπτον· ἔλαβε δὲ Μωυσῆς τὴν ῥάβδον τὴν παρὰ τοῦ Θεοῦ ἐν τῇ χειρὶ αὐτοῦ.
\vs{21}Εἶπε δὲ Κύριος πρὸς Μωυσῆν, πορευομένου σου καὶ ἀποστρέφοντος εἰς Αἴγυπτον, ὅρα πάντα τὰ τέρατα ἃ δέδωκα ἐν ταῖς χερσί σου, ποιήσεις αὐτὰ ἐναντίον Φαραώ· ἐγὼ δὲ σκληρυνῶ τὴν καρδίαν αὐτοῦ, καὶ οὐ μὴ ἐξαποστείλῃ τὸν λαόν.
\vs{22}Σὺ δὲ ἐρεῖς τῷ Φαραῴ, τάδε λέγει Κύριος, υἱὸς πρωτότοκός μου Ἰσραήλ.
\vs{23}Εἶπα δέ σοι, ἐξαπόστειλον τὸν λαόν μου, ἵνα μοι λατρεύσῃ· εἰ μὲν οὖν μὴ βούλει ἐξαποστεῖλαι αὐτούς, ὅρα οὖν, ἐγὼ ἀποκτένῶ τὸν υἱόν σου τὸν πρωτότοκον.
\vs{24}Ἐγένετο δὲ ἐν τῇ ὁδῷ ἐν τῷ καταλύματι συνήντησεν αὐτῷ Ἄγγελος Κυρίου, καὶ ἐζήτει αὐτὸν ἀποκτεῖναι.
\vs{25}Καὶ λαβοῦσα Σεπφώρα ψῆφον, περιέτεμε τὴν ἀκροβυστίαν τοῦ υἱοῦ αὐτῆς· καὶ προσέπεσε πρὸς τοὺς πόδας αὐτοῦ, καὶ εἶπεν, ἔστη τὸ αἷμα τῆς περιτομῆς τοῦ παιδίου μου.
\vs{26}Καὶ ἀπῆλθεν ἀπʼ αὐτοῦ, διότι εἶπεν, ἔστη τὸ αἷμα τῆς περιτομῆς τοῦ παιδίου μου.
\vs{27}Εἶπε δὲ Κύριος πρὸς Ἀαρὼν, πορεύθητι εἰς συνάντησιν Μωυσῇ εἰς τὴν ἔρημον· καὶ ἐπορεύθη, καὶ συνήντησεν αὐτῷ ἐν τῷ ὄρει τοῦ Θεοῦ, καὶ κατεφίλησαν ἀλλήλους.
\vs{28}Καὶ ἀνήγγειλε Μωυσῆς τῷ Ἀαρὼν πάντας τοὺς λόγους Κυρίου, οὓς ἀπέστειλε, καὶ πάντα τὰ ῥήματα, ἃ ἐνετείλατο αὐτῷ.
\vs{29}Ἐπορεύθη δὲ Μωυσῆς καὶ Ἀαρὼν, καὶ συνήγαγον τὴν γερουσίαν τῶν υἱῶν Ἰσραήλ.
\vs{30}Καὶ ἐλάλησεν Ἀαρὼν πάντα τὰ ῥήματα ταῦτα, ἃ ἐλάλησεν ὁ Θεὸς πρὸς Μωυσῆν, καὶ ἐποίησε τὰ σημεῖα ἐναντίον τοῦ λαοῦ.
\vs{31}Καὶ ἐπίστευσεν ὁ λαὸς καὶ ἐχάρη, ὅτι ἐπεσκέψατο ὁ Θεὸς τοὺς υἱοὺς Ἰσραὴλ, καὶ ὅτι εἶδεν αὐτῶν τὴν θλίψιν· κύψας δὲ ὁ λαὸς προσεκύνησε.

\ch{5}
Καὶ μετὰ ταῦτα εἰσῆλθε Μωυσῆς καὶ Ἀαρὼν πρὸς Φαραὼ, καὶ εἶπαν αὐτῷ, τάδε λέγει Κύριος ὁ Θεὸς Ἰσραὴλ, ἐξαπόστειλον τὸν λαόν μου, ἵνα μοι ἑορτάσωσιν ἐν τῇ ἐρήμῳ.
\vs{2}Καὶ εἶπε Φαραὼ, τίς ἐστιν οὗ εἰσακούσομαι τῆς φωνῆς αὐτοῦ, ὥστε ἐξαποστεῖλαι τοὺς υἱοὺς Ἰσραήλ; οὐκ οἶδα τὸν Κύριον, καὶ τὸν Ἰσραὴλ οὐκ ἐξαποστέλλω.
\vs{3}Καὶ λέγουσιν αὐτῷ, ὁ Θεὸς τῶν Ἐβραίων προσκέκληται ἡμᾶς· πορευσόμεθα οὖν ὁδὸν τριῶν ἡμερῶν εἰς τὴν ἔρημον, ὅπως θύσωμεν Κυρίῳ τῷ Θεῷ ἡμῶν, μή ποτε συναντήσῃ ἡμῖν θάνατος ἢ φόνος.
\vs{4}Καὶ εἶπεν αὐτοῖς ὁ βασιλεὺς Αἰγύπτου, ἱνατί Μωυσῆς καὶ Ἀαρών διαστρέφετε τὸν λαὸν ἀπὸ τῶν ἔργων; ἀπέλθατε ἕκαστος ὑμῶν πρὸς τὰ ἔργα αὐτοῦ.
\vs{5}Καὶ εἶπεν Φαραὼ, ἰδοὺ νῦν πολυπληθεῖ ὁ λαὸς, μὴ οὖν καταπαύσωμεν αὐτοὺς ἀπὸ τῶν ἔργων.
\vs{6}Συνέταξε δὲ Φαραὼ τοῖς ἐργοδιώκταις τοῦ λαοῦ, καὶ τοῖς γραμματεῦσι, λέγων,
\vs{7}οὐκέτι προστεθήσεσθε διδόναι ἄχυρον τῷ λαῷ εἰς τὴν πλινθουργίαν, καθάπερ χθὲς καὶ τρίτην ἡμέραν· ἀλλ αὐτοὶ πορευέσθωσαν καὶ συναγαγέτωσαν ἑαυτοῖς ἄχυρα.
\vs{8}Καὶ τὴν σύνταξιν τῆς πλινθείας, ἧς αὐτοὶ ποιοῦσι, καθʼ ἑκάστην ἡμέραν ἐπιβαλεῖς αὐτοῖς· οὐκ ἀφελεῖς οὐδέν· σχολάζουσι γάρ· διὰ τοῦτο κεκράγασι, λέγοντες, ἐγερθῶμεν, καὶ θύσωμεν τῷ Θεῷ ἡμῶν.
\vs{9}Βαρυνέσθω τὰ ἔργα τῶν ἀνθρώπων τούτων, καὶ μεριμνάτωσαν ταῦτα, καὶ μὴ μεριμνάτωσαν ἐν λόγοις κενοῖς.

\vs{10}Κατέσπευδον δὲ αὐτοὺς οἱ ἐργοδιῶκται καὶ οἱ γραμματεῖς, καὶ ἔλεγον πρὸς τὸν λαὸν, λέγοντες, τάδε λέγει Φαραὼ, οὐκέτι δίδωμι ὑμῖν ἄχυρα.
\vs{11}Αὐτοὶ ὑμεῖς πορευόμενοι συλλέγετε ἑαυτοῖς ἄχυρα, ὅθεν ἐὰν εὕρητε· οὐ γὰρ ἀφαιρεῖται ἀπὸ τῆς συντάξεως ὑμῶν οὐθέν.
\vs{12}Καὶ διεσπάρη ὁ λαὸς ἐν ὅλῃ γῇ Αἰγύπτῳ συναγαγεῖν καλάμην εἰς ἄχυρα.
\vs{13}Οἱ δὲ ἐργοδιῶκται κατέσπευδον αὐτοὺς, λέγοντες, συντελεῖτε τὰ ἔργα τὰ καθήκοντα καθʼ ἡμέραν, καθάπερ καὶ ὅτε τὸ ἄχυρον ἐδίδοτο ὑμῖν.
\vs{14}Καὶ ἐμαστιγώθησαν οἱ γραμματεῖς τοῦ γένους τῶν υἱῶν Ἰσραὴλ, οἱ κατασταθέντες ἐπʼ αὐτοὺς, ὑπὸ τῶν ἐπιστατῶν τοῦ Φαραὼ, λέγοντες, διατί οὐ συνετελέσατε τὰς συντάξεις ὑμῶν τῆς πλινθείας καθάπερ χθὲς καὶ τρίτην ἡμέραν, καὶ τὸ τῆς σήμερον;
\vs{15}Εἰσελθόντες δὲ οἱ γραμματεῖς τῶν υἱῶν Ἰσραὴλ κατεβόησαν πρὸς Φαραὼ, λέγοντες, ἱνατί σὺ οὕτως ποιεῖς τοῖς σοῖς οἰκέταις;
\vs{16}Ἄχυρον οὐ δίδοται τοῖς οἰκέταις σου, καὶ τὴν πλίνθον ἡμῖν λέγουσι ποιεῖν· καὶ ἰδοὺ οἱ παῖδές σου μεμαστίγωνται, ἀδικήσεις οὖν τὸν λαόν σου.
\vs{17}Καὶ εἶπεν αὐτοῖς, σχολάζετε, σχολασταί ἐστε· διὰ τοῦτο λέγετε, πορευθῶμεν, θύσωμεν τῷ Θεῷ ἡμῶν.
\vs{18}Νῦν οὖν πορευθέντες, ἐργάζεσθε· τὸ γὰρ ἄχυρον οὐ δοθήσεται ὑμῖν, καὶ τὴν σύνταξιν τῆς πλινθείας ἀποδώσετε.
\vs{19}Ἑώρων δὲ οἱ γραμματεῖς τῶν υἱῶν Ἰσραὴλ ἑαυτοὺς ἐν κακοῖς, λέγοντες, οὐκ ἀπολείψετε τῆς πλινθείας τὸ καθῆκον τῇ ἡμέρᾳ.
\vs{20}Συνήντησαν δὲ Μωυσῇ καὶ Ἀαρὼν ἐρχομένοις εἰς συνάντησιν αὐτοῖς, ἐκπορευομένων αὐτῶν ἀπὸ Φαραώ,
\vs{21}Καὶ εἶπαν αὐτοῖς, ἴδοι ὁ Θεὸς ὑμᾶς καὶ κρίναι, ὅτι ἐβδελύξατε τὴν ὀσμὴν ἡμῶν ἐναντίον Φαραὼ, καὶ ἐναντίον τῶν θεραπόντων αὐτοῦ, δοῦναι ῥομφαίαν εἰς τὰς χεῖρας αὐτοῦ, ἀποκτεῖναι ἡμᾶς.
\vs{22}Ἐπέστρεψε δὲ Μωυσῆς πρὸς Κύριον, καὶ εἶπε, δέομαι, Κύριε· τί ἐκάκωσας τὸν λαὸν τοῦτον; καὶ ἱνατί ἀπέσταλκάς με;
\vs{23}Καὶ ἀφʼ οὗ πεπόρευμαι πρὸς Φαραὼ, λαλῆσαι ἐπὶ τῷ σῷ ὀνόματι, ἐκάκωσε τὸν λαὸν τοῦτον· καὶ οὐκ ἐῤῥύσω τὸν λαόν σου.

\ch{6}
Καὶ εἶπε Κύριος πρὸς Μωυσῆν, ἤδη ὄψει ἃ ποιήσω τῷ Φαραῷ· ἐν γὰρ χειρὶ κραταίᾳ ἐξαποστελεῖ αὐτούς, καὶ ἐν βραχίονι ὑψηλῷ ἐκβαλεῖ αὐτοὺς ἐκ τῆς γῆς αὐτοῦ.
\vs{2}Ἐλάλησε δὲ ὁ Θεὸς πρὸς Μωυσῆν, καὶ εἶπε πρὸς αὐτὸν, ἐγὼ Κύριος.
\vs{3}Καὶ ὤφθην πρὸς Ἀβραὰμ καὶ Ἰσαὰκ καὶ Ἰακὼβ, Θεὸς ὢν αὐτῶν· καὶ τὸ ὄνομά μου Κύριος οὐκ ἐδήλωσα αὐτοῖς.
\vs{4}Καὶ ἔστησα τὴν διαθήκην μου πρὸς αὐτοὺς, ὥστε δοῦναι αὐτοῖς τὴν γῆν τῶν Χαναναίων, τὴν γῆν ἣν παρῳκήκασιν, ἐν ᾗ καὶ παρῴκησαν ἐπʼ αὐτῆς.
\vs{5}Καὶ ἐγὼ εἰσήκουσα τὸν στεναγμὸν τῶν υἱῶν Ἰσραήλ, ὃν οἱ Αἰγύπτιοι καταδουλοῦνται αὐτούς, καὶ ἐμνήσθην τῆς διαθήκης ὑμῶν.
\vs{6}Βάδιζε, εἶπον τοῖς υἱοῖς Ἰσραὴλ, λέγων, ἐγὼ Κύριος· καὶ ἐξάξω ὑμᾶς ἀπὸ τῆς δυναστείας τῶν Αἰγυπτίων, καὶ ῥύσομαι ὑμᾶς ἐκ τῆς δουλείας, καὶ λυτρώσομαι ὑμᾶς ἐν βραχίονι ὑψηλῷ καὶ κρίσει μεγάλῃ·
\vs{7}Καὶ λὴψομαι ἐμαυτῷ ὑμᾶς λαὸν ἐμοὶ, καὶ ἔσομαι ὑμῶν Θεός· καὶ γνώσεσθε ὅτι ἐγὼ Κύριος ὁ Θεὸς ὑμῶν, ὁ ἐξαγαγὼν ὑμᾶς ἐκ τῆς καταδυναστείας τῶν Αἰγυπτίων.
\vs{8}Καὶ εἰσάξω ὑμᾶς εἰς τὴν γῆν, εἰς ἣν ἐξέτεινα τὴν χεῖρά μου, δοῦναι αὐτὴν τῷ Ἀβραὰμ, καὶ Ἰσαὰκ, καὶ Ἰακὼβ, καὶ δώσω ὑμῖν αὐτὴν ἐν κληρῷ· ἐγὼ Κύριος.
\vs{9}Ἐλάλησε δὲ Μωυσῆς οὕτω τοῖς υἱοῖς Ἰσραήλ· καὶ οὐκ εἰσήκουσαν Μωυσῇ ἀπὸ τῆς ὀλιγοψυχίας, καὶ ἀπὸ τῶν ἔργων τῶν σκληρῶν.
\vs{10}Εἶπε δὲ Κύριος πρὸς Μωυσῆν λέγων,
\vs{11}εἴσελθε, λάλησον Φαραῷ βασιλεῖ Αἰγύπτου, ἵνα ἐξαποστείλῃ τοὺς υἱοὺς Ἰσραὴλ ἐκ τῆς γῆς αὐτοῦ.
\vs{12}Ἐλάλησε δὲ Μωυσῆς ἔναντι Κυρίου, λέγων, ἰδοὺ οἱ υἱοὶ Ἰσραὴλ οὐκ εἰσήκουσάν μου, καὶ πῶς εἰσακούσεταί μου Φαραώ; ἐγὼ δὲ ἄλογός εἰμι.
\vs{13}Εἶπε δὲ Κύριος πρὸς Μωυσῆν καὶ Ἀαρὼν, καὶ συνέταξεν αὐτοῖς πρὸς Φαραὼ βασιλέα Αἰγύπτου, ὥστε ἐξαποστεῖλαι τοὺς υἱοὺς Ἰσραὴλ ἐκ γῆς Αἰγύπτου.

\vs{14}Καὶ οὗτοι ἀρχηγοὶ οἴκων πατριῶν αὐτῶν· υἱοὶ Ῥουβὴν, πρωτοτόκου Ἰσραήλ· Ἑνὼχ, καὶ Φαλλοὺς, Ἀσρὼν, καὶ Χαρμεί· αὕτη ἡ συγγένεια Ῥουβήν.
\vs{15}Καὶ υἱοὶ Συμεών· Ἰεμουὴλ, καὶ Ἰαμεὶμ, καὶ Ἀὼδ, καὶ Ἰαχεὶν, καὶ Σαὰρ, καὶ Σαοὺλ ὁ ἐκ τῆς Φοινίσσης· αὗται αἱ πατριαὶ τῶν υἱῶν Συμεών.
\vs{16}Καὶ ταῦτα τὰ ὀνόματα τῶν υἱῶν Λευὶ κατὰ συγγενείας αὐτῶν· Γεδσὼν, Καὰθ, καὶ Μεραρεί· καὶ τὰ ἔτη τῆς ζωῆς Λευὶ ἑκατὸν τριάκοντα ἑπτά.
\vs{17}Καὶ οὗτοι υἱοὶ Γεδσών· Λοβενεὶ, καὶ Σεμεεί· οἶκοι πατριᾶς αὐτῶν.
\vs{18}Καὶ υἱοὶ Καάθ· Ἀμβρὰμ, καὶ Ἰσσαάρ, Χεβρὼν, καὶ Ὀζειήλ· καὶ τὰ ἔτη τῆς ζωῆς Καὰθ ἑκατὸν τριάκοντα τρία ἔτη.
\vs{19}Καὶ υἱοὶ Μεραρεί· Μοολεὶ, καὶ Ὀμουσεί. οὗτοι οἱ οἶκοι πατριῶν Λευὶ κατὰ συγγενείας αὐτῶν.
\vs{20}Καὶ ἔλαβεν Ἀμβρὰν τὴν Ἰωχαβὲδ, θυγατέρα τοῦ ἀδελφοῦ τοῦ πατρὸς αὐτοῦ, ἑαυτῷ εἰς γυναῖκα· καὶ ἐγέννησεν αὐτῷ τόν τε Ἀαρὼν καὶ τὸν Μωυσῆν, καὶ Μαριὰμ τὴν ἀδελφὴν αὐτῶν· τὰ δὲ ἔτη τῆς ζωῆς Ἀμβρὰμ, ἑκατὸν τριάκοντα δύο ἔτη.
\vs{21}Καὶ υἱοὶ Ἰσσαάρ· Κορὲ, καὶ Ναφὲκ, καὶ Ζεχρεί.
\vs{22}Καὶ υἱοὶ Ὀζειήλ· Μισαὴλ, καὶ Ἐλισαφὰν, καὶ Σεγρεί.
\vs{23}Ἔλαβε δὲ Ἀαρὼν τὴν Ἐλισαβὲθ θυγατέρα Ἀμιναδὰβ, ἀδελφὴν Ναασσὼν, αὐτῷ γυναῖκα· καὶ ἔτεκεν αὐτῷ τόν τε Ναδὰβ, καὶ Ἀβιοὺδ, καὶ τὸν Ἐλεάζαρ, καὶ Ἰθάμαρ.
\vs{24}Υἱοὶ δὲ Κορέ· Ἀσεὶρ, καὶ Ἑλκανὰ, καὶ Ἀβιασάρ· αὗται αἱ γενέσεις Κορέ.
\vs{25}Καὶ Ἐλεάζαρ ὁ τοῦ Ἀαρὼν ἔλαβε τῶν θυγατέρων Φουτιὴλ αὐτῷ γυναῖκα· καὶ ἔτεκεν αὐτῷ τὸν Φινεές· αὗται αἱ ἀρχαὶ πατριᾶς Λευιτῶν, κατὰ γενέσεις αὐτῶν.
\vs{26}Οὗτος Ἀαρὼν καὶ Μωυσῆς, οἷς εἶπεν αὐτοῖς ὁ Θεὸς ἐξαγαγεῖν τοὺς υἱοὺς Ἰσραὴλ ἐκ γῆς Αἰγύπτου σὺν δυνάμει αὐτῶν.
\vs{27}Οὗτοί εἰσιν οἱ διαλεγόμενοι πρὸς Φαραὼ βασιλέα Αἰγύπτου· καὶ ἐξήγαγον τοὺς υἱοὺς Ἰσραὴλ ἐκ γῆς Αἰγύπτου αὐτὸς Ἀαρὼν καὶ Μωυσὴς,
\vs{28}ᾗ ἡμέρᾳ ἐλάλησε Κύριος Μωυσῇ ἐν γῇ Αἰγύπτῳ.
\vs{29}Καὶ ἐλάλησε Κύριος πρὸς Μωυσῆν, λέγων, ἐγὼ Κύριος· λάλησον πρὸς Φαραὼ βασιλέα Αἰγύπτου ὅσα ἐγὼ λέγω πρὸς σέ.
\vs{30}Καὶ εἶπε Μωυσῆς ἐναντίον Κυρίου, ἰδοὺ ἐγὼ ἰσχνόφωνός εἰμι, καὶ πῶς εἰσακούσεταί μου Φαραώ,

\ch{7}
Καὶ εἶπε Κύριος πρὸς Μωυσῆν, λέγων, ἰδοὺ δέδωκά σε θεὸν Φαραὼ, καὶ Ἀαρὼν ὁ ἀδελφός σου ἔσται σου προφήτης.
\vs{2}Σὺ δὲ λαλήσεις αὐτῷ πάντα ὅσα σοι ἐντέλλομαι· ὁ δὲ Ἀαρὼν ὁ ἀδελφός σου λαλήσει πρὸς Φαραὼ, ὥστε ἐξαποστεῖλαι τοὺς υἱοὺς Ἰσραὴλ ἐκ τῆς γῆς αὐτοῦ.
\vs{3}Ἐγὼ δὲ σκληρυνῶ τὴν καρδίαν Φαραὼ, καὶ πληθυνῶ τὰ σημεῖά μου καὶ τὰ τέρατα ἐν γῇ Αἰγύπτῳ.
\vs{4}Καὶ οὐκ εἰσακούσεται ὑμῶν Φαραώ· καὶ ἐπιβαλῶ τὴν χεῖρά μου ἐπʼ Αἴγυπτον, καὶ ἐξάξω σὺν δυνάμει μου τὸν λαόν μου τοὺς υἱοὺς Ἰσραὴλ ἐκ γῆς Αἰγύπτου σὺν ἐκδικήσει μεγάλῃ.
\vs{5}Καὶ γνώσονται πάντες οἱ Αἰγύπτιοι ὅτι ἐγώ εἰμι Κύριος, ἐκτείνων τὴν χεῖρά μου ἐπʼ Αἴγυπτον, καὶ ἐξάξω τοὺς υἱοὺς Ἰσραὴλ ἐκ μέσον αὐτῶν.
\vs{6}Ἐποίησε δὲ Μωυσῆς καὶ Ἀαρὼν καθάπερ ἐνετείλατο αὐτοῖς Κύριος, οὕτως ἐποίησαν.
\vs{7}Μωυσῆς δὲ ἦν ἐτῶν ὀγδοήκοντα, Ἀαρὼν δὲ ὁ ἀδελφὸς αὐτοῦ ἐτῶν ὀγδοήκοντατριῶν, ἡνίκα ἐλάλησεν πρὸς Φαραώ.
\vs{8}Καὶ εἶπε Κύριος πρὸς Μωυσῆν καὶ Ἀαρὼν, λέγων,
\vs{9}καὶ ἐὰν λαλήσῃ πρὸς ὑμᾶς Φαραὼ, λέγων, δότε ἡμῖν σημεῖον ἢ τέρας, καὶ ἐρεῖς Ἀαρὼν τῷ ἀδελφῷ σου, λάβε τὴν ῥάβδον, καὶ ῥίψον ἐπὶ τὴν γῆν ἐναντίον Φαραὼ, καὶ ἐναντίον τῶν θεραπόντων αὐτοῦ, καὶ ἔσται δράκων.
\vs{10}Εἰσῆλθε δὲ Μωυσῆς καὶ Ἀαρὼν ἐναντίον Φαραὼ, καὶ τῶν θεραπόντων αὐτοῦ· καὶ ἐποίησαν οὕτως, καθάπερ ἐνετείλατο αὐτοῖς Κύριος· καὶ ἔῤῥιψεν Ἀαρὼν τὴν ῥάβδον ἐναντίον Φαραὼ, καὶ ἐναντίον τῶν θεραποντων αὐτοῦ, καὶ ἐγένετο δράκων.
\vs{11}Συνεκάλεσε δὲ Φαραὼ τοὺς σοφιστὰς Αἰγύπτου, καὶ τοὺς φαρμακούς· καὶ ἐποίησαν καὶ οἱ ἐπαοιδοὶ τῶν Αἰγυπτίων ταῖς φαρμακίαις αὐτῶν ὡσαύτως.
\vs{12}Καὶ ἔῤῥιψαν ἔκαστος τὴν ῥάβδον αὐτῶν, καὶ ἐγένοντο δράκοντες· καὶ κατέπιεν ἡ ῥάβδος ἡ Ἀαρὼν τὰς ἐκείνων ῥάβδους.
\vs{13}Καὶ κατίσχυσεν ἡ καρδία Φαραὼ, καὶ οὐκ εἰσήκουσεν αὐτῶν, καθάπερ ἐνετείλατο αὐτοῖς Κύριος.

\vs{14}Εἶπε δὲ Κύριος πρὸς Μωυσῆν, βεβάρηται ἡ καρδία Φαραὼ, τοῦ μὴ ἐξαποστεῖλαι τὸν λαόν.
\vs{15}Βάδισον πρὸς Φαραὼ τὸ πρωΐ· ἰδοὺ αὐτὸς ἐκπορεύεται ἐπὶ τὸ ὕδωρ, καὶ ἔσῃ συναντῶν αὐτῷ ἐπὶ τὸ χεῖλος τοῦ ποταμοῦ· καὶ τὴν ῥάβδον τὴν στραφεῖσαν εἰς ὄφιν λήψῃ ἐν τῇ χειρί σου.
\vs{16}Καὶ ἐρεῖς πρὸς αὐτὸν, Κύριος ὁ Θεὸς τῶν Ἐβραίων ἀπέσταλκέ με πρὸς σὲ, λέγων, ἐξαπόστειλον τὸν λαόν μου, ἵνα μοι λατρεύσῃ ἐν τῇ ἐρήμῳ· καὶ ἰδοὺ οὐκ εἰσήκουσας ἕως τούτου.
\vs{17}Τάδε λέγει Κύριος, ἐν τούτῳ γνώσῃ ὅτι ἐγὼ Κύριος· ἰδοὺ ἑγὼ τύπτω τῇ ῥάβδῳ τῇ ἐν τῇ χειρί μου ἐπὶ τὸ ὕδωρ τὸ ἐν τῷ ποταμῷ, καὶ μεταβαλεῖ εἰς αἷμα.
\vs{18}Καὶ οἱ ἰχθύες οἱ ἐν τῷ ποταμῷ τελευτήσουσι· καὶ ἐποζέσει ὁ ποταμὸς, καὶ οὐ δυνήσονται οἱ Αἰγύπτιοι πιεῖν ὕδωρ ἀπὸ τοῦ ποταμοῦ.
\vs{19}Εἶπε δὲ Κύριος πρὸς Μωυσῆν, εἶπὸν Ἀαρὼν τῷ ἀδελφῷ σου, λάβε τὴν ῥάβδον σου ἐν τῇ χειρί σου, καὶ ἔκτεινον τὴν χεῖρά σου ἐπὶ τὰ ὕδατα Αἰγύπτου, καὶ ἐπὶ τοὺς ποταμοὺς αὐτῶν, καὶ ἐπὶ τὰς διώρυγας αὐτῶν, καὶ ἐπὶ τὰ ἕλη αὐτῶν, καὶ ἐπὶ πᾶν συνεστηκὸς ὕδωρ αὐτῶν, καὶ ἔσται αἷμα· καὶ ἐγένετο αἷμα ἐν πάσῃ γῇ Αἰγύπτου, ἔν τε τοῖς ξύλοις καὶ ἐν τοῖς λίθοις.
\vs{20}Καὶ ἐποίησαν οὕτως Μωυσῆς καὶ Ἀαρὼν, καθάπερ ἐνετείλατο αὐτοῖς Κύριος· καὶ ἐπάρας τῇ ῥάβδῳ αὐτοῦ ἐπάταξε τὸ ὕδωρ τὸ ἐν τῷ ποταμῷ ἐναντίον Φαραὼ, καὶ ἐναντίον τῶν θεραπόντων αὐτοῦ· καὶ μετέβαλε πᾶν τὸ ὕδωρ τὸ ἐν τῷ ποταμῷ εἰς αἷμα.
\vs{21}Καὶ οἱ ἰχθύες οἱ ἐν τῷ ποταμῷ ἐτελεύτησαν· καὶ ἐπώζεσεν ὁ ποταμὸς, καὶ οὐκ ἠδύναντο οἱ Αἰγύπτιοι πιεῖν ὕδωρ ἐκ τοῦ ποταμοῦ· καὶ ἦν τὸ αἷμα ἐν πάσῃ γῇ Αἰγύπτου.
\vs{22}Ἐποίησαν δὲ ὡσαύτως καὶ οἱ ἐπαοιδοὶ τῶν Αἰγυπτίων ταῖς φαρμακίαις αὐτῶν· καὶ ἐσκληρύνθη ἡ καρδία Φαραὼ, καὶ οὐκ εἰσήκουσεν αὐτῶν, καθάπερ εἶπε Κύριος.
\vs{23}Ἐπιστραφεὶς δὲ Φαραὼ εἰσῆλθεν εἰς τὸν οἶκον αὐτοῦ· καὶ οὐκ ἐπέστησε τὸν νοῦν αὐτοῦ οὐδὲ ἐπὶ τούτῳ.
\vs{24}Ὤρυξαν δὲ πάντες οἱ Αἰγύπτιοι κύκλῳ τοῦ ποταμοῦ, ὥστε πιεῖν ὕδωρ· καὶ οὐκ ἠδύναντο πιεῖν ὕδωρ ἀπὸ τοῦ ποταμοῦ.
\vs{25}Καὶ ἀνεπληρώθησαν ἑπτὰ ἡμέραι, μετὰ τὸ πατάξαι Κύριον τὸν ποταμόν.

\vs{26}Εἶπε δὲ Κύριος πρὸς Μωυσὴν, εἴσελθε πρὸς Φαραὼ, καὶ ἐρεῖς πρὸς αὐτὸν, τάδε λεγέι Κύριος, ἐξαπόστειλον τὸν λαόν μου, ἵνα μοι λατρεύσωσιν.
\vs{27}Εἰ δὲ μὴ βούλει σὺ ἐξαποστεῖλαι, ἰδοὺ ἐγὼ τύπτω πάντα τὰ ὅριά σου τοῖς βατράχοις.
\vs{28}Καὶ ἐξερεύξεται ὁ ποταμὸς βατράχους· καὶ ἀναβάντες εἰσελεύσονται εἰς τοὺς οἴκους σου, καὶ εἰς τὰ ταμιεῖα τῶν κοιτώνων σου, καὶ ἐπὶ τῶν κλινῶν σου, καὶ ἐπὶ τοὺς οἴκους τῶν θεραπόντων σου, καὶ τοῦ λαοῦ σου, καὶ ἐν τοῖς φυράμασί σου, καὶ ἐν τοῖς κλιβάνοις σου.
\vs{29}Καὶ ἐπὶ σὲ, καὶ ἐπὶ τοὺς θεράποντάς σου, καὶ ἐπὶ τὸν λαόν σου, ἀναβήσονται οἱ βάτραχοι.

\ch{8}Εἶπε δὲ Κύριος πρὸς Μωυσῆν, εἶπον Ἀαρὼν τῷ ἀδελφῷ σου, ἔκτεινον τῇ χειρὶ τὴν ῥάβδον σου ἐπὶ τοὺς ποταμοὺς, καὶ ἐπὶ τὰς διώρυγας, καὶ ἐπὶ τὰ ἕλη, καὶ ἀνάγαγε τοὺς βατράχους.
\vs{2}Καὶ ἐξέτεινεν Ἀαρὼν τὴν χεῖρα ἐπὶ τὰ ὕδατα Αἰγύπτου, καὶ ἀνήγαγε τοὺς βατράχους· καὶ ἀνεβιβάσθη ὁ βάτραχος, καὶ ἐκάλυψε τὴν γῆν Αἰγύπτου.
\vs{3}Ἐποίησαν δὲ ὡσαύτως καὶ οἱ ἐπαοιδοὶ τῶν Αἰγυπτίων ταῖς φαρμακίαις αὐτῶν, καὶ ἀνήγαγον τοὺς βατράχους ἐπὶ γῆν Αἰγύπτου.
\vs{4}Καὶ ἐκάλεσε Φαραὼ Μωυσῆν καὶ Ἀαρὼν, καὶ εἶπεν, εὔξασθε περὶ ἐμοῦ πρὸς Κύριον, καὶ περιελέτω τοὺς βατράχους ἀπʼ ἐμοῦ, καὶ ἀπὸ τοῦ ἐμοῦ λαοῦ· καὶ ἐξαποστελῶ αὐτοὺς, καὶ θύσωσι τῷ Κυρίῳ.
\vs{5}Εἶπε δὲ Μωυσῆς πρὸς Φαραὼ, τάξαι πρὸς με πότε εὔξομαι περὶ σοῦ, καὶ περὶ τῶν θεραπόντων σου, καὶ τοῦ λαοῦ σου, ἀφανίσαι τοὺς βατράχους ἀπὸ σοῦ, καὶ ἀπὸ τοῦ λαοῦ σου, καὶ ἐκ τῶν οἰκιῶν ὑμῶν, πλὴν ἐν τῷ ποταμῷ ὑπολειφθήσονται.
\vs{6}Ὁ δὲ εἶπεν, εἰς αὔριον· εἶπεν οὖν, ὡς εἴρηκας· ἵνα εἰδῇς ὅτι οὐκ ἔστιν ἄλλος πλὴν Κυρίου·
\vs{7}Καὶ περιαιρεθήσονται οἱ βάτραχοι ἀπὸ σοῦ, καὶ ἀπὸ τῶν οἰκιῶν ὑμῶν, καὶ ἀπὸ τῶν ἐπαύλεων, καὶ ἀπὸ τῶν θεραπόντων σου, καὶ ἀπὸ τοῦ λαοῦ σου, πλὴν ἐν τῷ ποταμῷ ὑπολειφθήσονται.
\vs{8}Ἐξῆλθε δὲ Μωυσῆς καὶ Ἀαρὼν ἀπὸ Φαραώ· καὶ ἐβόησε Μωυσῆς πρὸς Κύριον περὶ τοῦ ὁρισμοῦ τῶν βατράχων, ὡς ἐτάξατο Φαραώ.
\vs{9}Ἐποιήσε δὲ Κύριος καθάπερ εἶπε Μωυσῆς· καὶ ἐτελεύτησαν οἱ βάτραχοι ἐκ τῶν οἰκιῶν, καὶ ἐκ τῶν ἐπαύλεων, καὶ ἐκ τῶν ἀγρῶν.
\vs{10}Καὶ συνήγαγον αὐτοὺς, θημωνίας θημωνίας· καὶ ὤζεσεν ἡ γῆ.
\vs{11}Ἰδὼν δὲ Φαραὼ ὅτι γέγονεν ἀνάψυξις, ἐβαρύνθη ἡ καρδία αὐτοῦ, καὶ οὐκ εἰσήκουσεν αὐτῶν, καθάπερ ἐλάλησε Κύριος.
\vs{12}Εἶπε δὲ Κύριος πρὸς Μωυσῆν, εἶπον Ἀαρὼν, ἔκτεινον τῇ χειρὶ τὴν ῥάβδον σου, καὶ πάταξον τὸ χῶμα τῆς γῆς· καὶ ἔσονται σκνίφες ἔν τε τοῖς ἀνθρώποις, καὶ ἐν τοῖς τετράποσι, καὶ ἐν πάσῃ γῇ Αἰγύπτου.
\vs{13}Ἐξέτεινεν οὖν Ἀαρὼν τῇ χειρὶ τὴν ῥάβδον, καὶ ἐπάταξε τὸ χῶμα τῆς γῆς· καὶ ἐγένοντο οἱ σκνίφες ἐν τοῖς ἀνθρώποις, ἔν τε τοῖς τετράποσι, καὶ ἐν παντὶ χώματι τῆς γῆς ἐγένοντο οἱ σκνίφες.
\vs{14}Ἐποίησαν δὲ ὡσαύτως καὶ οἱ ἐπαοιδοὶ ταῖς φαρμακίαις αὐτῶν, ἐξαγαγεῖν τὸν σκνῖφα, καὶ οὐκ ἠδύναντο· καὶ ἐγένοντο οἱ σκνίφες ἔν τε τοῖς ἀνθρώποις, καὶ ἐν τοῖς τετράποσιν.
\vs{15}Εἶπαν οὖν οἱ ἐπαοιδοὶ τῷ Φαραῷ, δάκτυλος Θεοῦ ἐστι τοῦτο· καὶ ἐσκληρύνθη ἡ καρδία Φαραὼ, καὶ οὐκ εἰσήκουσεν αὐτῶν, καθάπερ ἐλάλησε Κύριος.
\vs{16}Εἶπε δὲ Κύριος πρὸς Μωυσῆν, ὄρθρισον τὸ πρωΐ, καὶ στῆθι ἐναντίον Φαραώ· καὶ ἰδοὺ αὐτὸς ἐξελεύσεται ἐπὶ τὸ ὕδωρ· καὶ ἐρεῖς πρὸς αὐτὸν, τάδε λέγει Κύριος, ἐξαπόστειλον τὸν λαόν μου, ἵνα μοι λατρεύσωσιν ἐν τῇ ἐρήμῳ.
\vs{17}Ἐὰν δὲ μὴ βούλει ἐξαποστεῖλαι τὸν λαόν μου, ἰδοὺ ἐγὼ ἐξαποστέλλω ἐπὶ σὲ, καὶ ἐπὶ τοὺς θεράποντάς σου, καὶ ἐπὶ τὸν λαόν σου, καὶ ἐπὶ τοὺς οἴκους ὑμῶν, κυνόμυιαν· καὶ πλησθήσονται αἱ οἰκίαι τῶν Αἰγυπτίων τῆς κυνομυίης, καὶ εἰς τὴν γῆν ἐφʼ ἧς εἰσιν ἐπʼ αὐτῆς.
\vs{18}Καὶ παραδοξάσω ἐν τῇ ἡμέρᾳ ἐκείνῃ τὴν γῆν Γεσὲμ, ἐφʼ ἧς ὁ λαός μου ἔπεστιν ἐπʼ αὐτῆς, ἐφʼ ἧς οὐκ ἔσται ἐκεῖ ἡ κυνόμυια· ἵνα εἰδῇς ὅτι ἐγώ εἰμι Κύριος ὁ Θεὸς πάσης τῆς γῆς.
\vs{19}Καὶ δώσω διαστολὴν ἀνὰ μέσον τοῦ ἐμοῦ λαοῦ, καὶ ἀνὰ μέσον τοῦ σου λαοῦ· ἐν δὲ τῇ αὔριον ἔσται τοῦτο ἐπὶ τῆς γῆς.
\vs{20}Ἐποίησε δὲ Κύριος οὕτως· καὶ παρεγένετο ἡ κυνόμυια πλῆθος εἰς τοὺς οἴκους Φαραὼ, καὶ εἰς τοὺς οἴκους τῶν θεραπόντων αὐτοῦ, καὶ εἰς πᾶσαν τὴν γῆν Αἰγύπτου· καὶ ἐξωλοθρεύθη ἡ γῆ ἀπὸ τῆς κυνομυίης.

\vs{21}Ἐκάλεσε δὲ Φαραὼ Μωυσῆν καὶ Ἀαρὼν, λέγων, ἐλθόντες θύσατε Κυρίῳ τῷ Θεῷ ὑμῶν ἐν τῇ γῇ.
\vs{22}Καὶ εἶπε Μωυσῆς, οὐ δυνατὸν γενέσθαι οὕτως· τὰ γὰρ βδελύγματα τῶν Αἰγυπτίων θύσομεν Κυρίῳ τῷ Θεῷ ἡμῶν· ἐὰν γὰρ θύσωμεν τὰ βδελύγματα τῶν Αἰγυπτίων ἐναντίον αὐτῶν, λιθοβοληθησόμεθα.
\vs{23}Ὁδὸν τριῶν ἡμερῶν πορευσόμεθα εἰς τὴν ἔρημον· καὶ θύσομεν τῷ Θεῷ ἡμῶν, καθάπερ εἶπεν Κύριος ἡμῖν.
\vs{24}Καὶ εἶπε Φαραὼ, ἐγὼ ἀποστέλλω ὑμᾶς, καὶ θύσατε τῷ Θεῷ ὑμῶν ἐν τῇ ἐρήμῳ· ἀλλʼ οὐ μακρὰν ἀποτενεῖτε πορευθῆναι· εὔξασθε οὖν περὶ ἐμοῦ πρὸς Κύριον.
\vs{25}Εἶπε δὲ Μωυσῆς, ὁ δὲ ἐγὼ ἐξελεύσομαι ἀπὸ σοῦ, καὶ εὔξομαι πρὸς τὸν Θεὸν, καὶ ἀπελεύσεται ἡ κυνόμυια καὶ ἀπὸ τῶν θεραπόντων σου, καὶ ἀπὸ τοῦ λαοῦ σου αὔριον· μὴ προσθῇς ἔτι Φαραὼ ἐξαπατῆσαι, τοῦ μὴ ἐξαποστεῖλαι τὸν λαὸν θῦσαι Κυρίῳ.
\vs{26}Ἐξῆλθε δὲ Μωυσῆς ἀπὸ Φαραὼ, καὶ ηὔξατο πρὸς τὸν Θεόν.
\vs{27}Ἐποίησε δὲ Κύριος καθάπερ εἶπε Μωυσῆς· καὶ περιεῖλε τὴν κυνόμυιαν ἀπὸ Φαραὼ, καὶ τῶν θεραπόντων αὐτοῦ, καὶ τοῦ λαοῦ αὐτοῦ, καὶ οὐ κατελείφθη οὐδεμία.
\vs{28}Καὶ ἐβάρυνε Φαραὼ τὴν καρδίαν αὐτοῦ καὶ ἐπὶ τοῦ καιροῦ τούτου, καὶ οὐκ ἠθέλησεν ἐξαποστεῖλαι τὸν λαόν.

\ch{9}
Εἶπε δὲ Κύριος πρὸς Μωυσῆν, εἴσελθε πρὸς Φαραὼ, καὶ ἐρεῖς αὐτῷ, τάδε λέγει Κύριος ὁ Θεὸς τῶν Ἑβραίων, ἐξαπόστειλον τὸν λαόν μου, ἵνα μοι λατρεύσωσι.
\vs{2}Εἰ μὲν οὖν μὴ βούλει ἐξαποστεῖλαι τὸν λαόν μου, ἀλλὰ ἔτι ἐγκρατεῖς αὐτοῦ,
\vs{3}Ἰδοὺ, χεὶρ Κυρίου ἐπέσται ἐν τοῖς κτήνεσί σου τοῖς ἐν τοῖς πεδίοις, ἔν τε τοῖς ἵπποις, καὶ ἐν τοῖς ὑποζυγίοις, καὶ ταῖς καμήλοις, καὶ βουσὶ, καὶ προβάτοις, θάνατος μέγας σφόδρα.
\vs{4}Καὶ παραδοξάσω ἐγὼ ἐν τῷ καιρῷ ἐκείνῳ ἀνὰ μέσον τῶν κτηνῶν τῶν Αἰγυπτίων, καὶ ἀνὰ μέσον τῶν κτηνῶν τῶν υἱῶν Ἰσραήλ· οὐ τελευτήσει ἀπὸ πάντων τῶν τοῦ Ἰσραὴλ υἱῶν ῥητόν.
\vs{5}Καὶ ἔδωκεν ὁ Θεὸς ὅρον, λέγων, ἐν τῇ αὔριον ποιήσει Κύριος τὸ ῥῆμα τοῦτο ἐπὶ τῆς γῆς.
\vs{6}Καὶ ἐποίησε Κύριος τὸ ῥῆμα τοῦτο τῇ ἐπαύριον· καὶ ἐτελεύτησε πάντα τὰ κτήνη τῶν Αἰγυπτίων· ἀπὸ δὲ τῶν κτηνῶν τῶν υἱῶν Ἰσραὴλ οὐκ ἐτελεύτησεν οὐδέν.
\vs{7}Ἰδὼν δὲ Φαραὼ ὅτι οὐκ ἐτελεύτησεν ἀπὸ πάντων τῶν κτηνῶν τῶν υἱῶν Ἰσραὴλ οὐδὲν, ἐβαρύνθη ἡ καρδία Φαραὼ, καὶ οὐκ ἐξαπέστειλε τὸν λαόν.
\vs{8}Εἶπε δὲ Κύριος πρὸς Μωυσῆν καὶ Ἀαρὼν, λέγων, λάβετε ὑμεῖς πληρεῖς τὰς χεῖρας αἰθάλης καμιναίας, καὶ πασάτω Μωυσῆς εἰς τὸν οὐρανὸν ἐναντίον Φαραὼ, καὶ ἐναντίον τῶν θεραπόντων αὐτοῦ.
\vs{9}Καὶ γενηθήτω κονιορτὸς ἐπὶ πᾶσαν τὴν γῆν Αἰγύπτου· καὶ ἔσται ἐπὶ τοὺς ἀνθρώπους, καὶ ἐπὶ τὰ τετράποδα, ἕλκη, φλυκτίδες ἀναζέουσαι ἔν τε τοῖς ἀνθρώποις, καὶ ἐν τοῖς τετράποσιν, ἐν πάσῃ γῇ Αἰγύπτου.
\vs{10}Καὶ ἔλαβεν τὴν αἰθάλην τῆς καμιναίας ἐναντίον Φαραὼ, καὶ ἔπασεν αὐτὴν Μωυσῆς εἰς τὸν οὐρανόν· καὶ ἐγένετο ἕλκη, φλυκτίδες ἀναζέουσαι, ἔν τε τοῖς ἀνθρώποις, καὶ ἐν τοῖς τετράποσι.
\vs{11}Καὶ οὐκ ἠδύναντο οἱ φαρμακοὶ στῆναι ἐναντίον Μωυσῆ διὰ τὰ ἕλκη· ἐγένετο γὰρ τὰ ἕλκη ἐν τοῖς φαρμακοῖς, καὶ ἐν πάσῃ γῇ Αἰγύπτου.
\vs{12}Ἐσκλήρυνε δὲ Κύριος τὴν καρδίαν Φαραὼ, καὶ οὐκ εἰσήκουσεν αὐτῶν, καθὰ συνέταξε Κύριος.
\vs{13}Εἶπε δὲ Κύριος πρὸς Μωυσῆν, ὄρθρισον τὸ πρωῒ, καὶ στῆθι ἐναντίον Φαραὼ, καὶ ἐρεῖς πρὸς αὐτὸν, τάδε λέγει Κύριος ὁ Θεὸς τῶν Ἑβραίων, ἐξαπόστειλον τὸν λαόν μου, ἵνα λατρεύσωσί μοι.
\vs{14}Ἐν τῷ γὰρ νῦν καιρῷ ἐγὼ ἐξαποστέλλω πάντα τὰ συναντήματά μου εἰς τὴν καρδίαν σου, καὶ τῶν θεραπόντων σου, καὶ τοῦ λαοῦ σου, ἵνα εἴδῇς ὅτι οὐκ ἔστιν, ὡς ἐγὼ, ἄλλος ἐν πάσῃ τῇ γῇ.
\vs{15}Νῦν γὰρ ἀποστείλας τὴν χεῖρα πατάξω σε, καὶ τὸν λαόν σου θανατώσω, καὶ ἐκτριβήσῃ ἀπὸ τῆς γῆς.
\vs{16}Καὶ ἕνεκεν τούτου διετηρήθης, ἵνα ἐνδείξωμαι ἐν σοὶ τὴν ἰσχύν μου, καὶ ὅπως διαγγελῇ τὸ ὄνομά μου ἐν πάσῃ τῇ γῇ.
\vs{17}Ἔτι οὖν σὺ ἐνποιῇ τοῦ λαοῦ μου, τοῦ μὴ ἐξαποστεῖλαι αὐτούς;
\vs{18}Ἰδοὺ ἐγὼ ὕω ταύτην τὴν ὥραν αὔριον χάλαζαν πολλὴν σφόδρα, ἥτις τοιαύτη οὐ γέγονεν ἐν Αἰγύπτῳ, ἀφʼ ἧς ἡμέρας ἔκτισται, ἕως τῆς ἡμέρας ταύτης.
\vs{19}Νῦν οὖν κατάσπευσον συναγαγεῖν τὰ κτήνη σου, καὶ ὅσα σοι ἐστὶν ἐν τῷ πεδίῳ· πάντες γὰρ οἱ ἄνθρωποι, καὶ τὰ κτήνη, ὅσα σοί ἐστιν ἐν τῷ πεδίῳ· πὰντες γὰρ οἱ ἄνθρωποι, καὶ τὰ κτήνη, ὅσα ἐὰν εὑρεθῇ ἐν τοῖς πεδίοις, καὶ μὴ εἰσέλθῃ εἰς οἰκίαν, πεσῇ δὲ ἐπʼ αὐτὰ ἡ χάλαζα, τελευτήσει.
\vs{20}Ὁ φοβούμενος τὸ ῥῆμα Κυρίου τῶν θεραπόντων Φαραὼ, συνήγαγε τὰ κτήνη αὐτοῦ εἰς τοὺς οἴκους.
\vs{21}Ὃς δὲ μὴ πρόσεσχεν τῇ διανοίᾳ εἰς τὸ ῥῆμα Κυρίου, ἀφῆκε τὰ κτήνη ἐν τοῖς πεδίοις.

\vs{22}Εἶπε δὲ Κύριος πρὸς Μωυσῆν, ἔκτεινον τὴν χεῖρά σου εἰς τὸν οὐρανὸν, καὶ ἔσται χάλαζα ἐπὶ πᾶσαν γῆν Αἰγύπτου, ἐπί τε τοὺς ἀνθρώπους, καὶ τὰ κτήνη, καὶ ἐπὶ πᾶσαν βοτάνην τὴν ἐπὶ τῆς γῆς.
\vs{23}Ἐξέτεινε δὲ Μωυσῆς τὴν χεῖρα εἰς τὸν οὐρανὸν, καὶ Κύριος ἔδωκε φωνὰς καὶ χάλαζαν· καὶ διέτρεχε τὸ πῦρ ἐπὶ τῆς γῆς· καὶ ἔβρεξε Κύριος χάλαζαν ἐπὶ πᾶσαν γῆν Αἰγύπτου.
\vs{24}Ἦν δὲ ἡ χάλαζα καὶ τὸ πῦρ φλογίζον ἐν τῇ χαλάζῃ· ἡ δὲ χάλαζα πολλὴ σφόδρα, ἥτις τοιαύτη οὐ γέγονεν ἐν Αἰγύπτῳ, ἀφʼ ἧς ἡμέρας γεγένηται ἐπʼ αὐτῆς ἔθνος.
\vs{25}Ἐπάταξε δὲ ἡ χάλαζα ἐν πάσῃ γῇ Αἰγύπτου, ἀπὸ ἀνθρώπου ἕως κτήνους· καὶ πᾶσαν βοτάνην τὴν ἐν τῷ πεδίῳ ἐπάταξεν ἡ χάλαζα· καὶ πάντα τὰ ξύλα τὰ ἐν τοῖς πεδίοις συνέτριψεν ἡ χάλαζα.
\vs{26}Πλὴν ἐν γῇ Γεσὲμ, οὗ ἦσαν οἱ υἱοὶ Ἰσραὴλ, οὐκ ἐγένετο ἡ χάλαζα.
\vs{27}Ἀποστείλας δὲ Φαραὼ ἐκάλεσε Μωυσῆν καὶ Ἀαρὼν, καὶ εἶπεν αὐτοῖς, ἡμάρτηκα τὸ νῦν· ὁ Κύριος δίκαιος, ἐγὼ δὲ καὶ ὁ λαός μου ἀσεβεῖς.
\vs{28}Εὔξασθε οὖν περὶ ἐμοῦ πρὸς Κύριον, καὶ παυσάσθω τοῦ γενηθῆναι φωνὰς Θεοῦ, καὶ χάλαζαν, καὶ πῦρ· καὶ ἐξαποστελῶ ὑμᾶς, καὶ οὐκέτι προστεθήσεσθε μένειν.
\vs{29}Εἶπε δὲ αὐτῷ Μωυσῆς, ὡς ἂν ἐξέλθω τὴν πόλιν, ἐκπετάσω τὰς χεῖράς μου πρὸς τὸν Κύριον, καὶ αἱ φωναὶ παύσονται, καὶ ἡ χὰλαζα καὶ ὁ ὑετὸς οὐκ ἔσται ἔτι, ἵνα γνῷς ὅτι τοῦ Κυρίου ἡ γῆ.
\vs{30}Καὶ σὺ καὶ οἱ θεράποντές σου, ἐπίσταμαι ὅτι οὐδέπω πεφόβησθε τὸν Κύριον.
\vs{31}Τὸ δὲ λίνον καὶ ἡ κριθὴ ἐπλήγη· ἡ γὰρ κριθὴ παρεστηκυῖα, τὸ δὲ λίνον σπερματίζον.
\vs{32}Ὁ δὲ πυρὸς καὶ ἡ ὀλύρα οὐκ ἐπληγησαν, ὄψιμα γὰρ ἦν.
\vs{33}Ἐξῆλθε δὲ Μωυσῆς ἀπὸ Φαραὼ ἐκτὸς τῆς πόλεως, καὶ ἐξέτεινε τὰς χεῖρας πρὸς Κύριον· καὶ αἱ φωναὶ ἐπαύσαντο, καὶ ἡ χάλαζα καὶ ὁ ὑετὸς οὐκ ἔσταξεν ἔτι ἐπὶ τὴν γῆν.
\vs{34}Ἰδὼν δὲ Φαραὼ ὅτι πέπαυται ὁ ὑετὸς καὶ ἡ χάλαζα καὶ αἱ φωναὶ, προσέθετο τοῦ ἁμαρτάνειν· καὶ ἐβάρυνεν αὐτοῦ τὴν καρδίαν, καὶ τῶν θεραπόντων αὐτοῦ.
\vs{35}Καὶ ἐσκληρύνθη ἡ καρδία Φαραὼ, καὶ οὐκ ἐξαπέστειλε τοὺς υἱοὺς Ἰσραὴλ, καθάπερ ἐλάλησε Κύριος τῷ Μωυσῇ.

\ch{10}
Εἶπε δὲ Κύριος πρὸς Μωυσῆν, λέγων, εἴσελθε πρὸς Φαραὼ, ἐγὼ γὰρ ἐσκλήρυνα αὐτοῦ τὴν καρδίαν καὶ τῶν θεραπόντων αὐτοῦ, ἵνα ἑξῆς ἐπέλθῃ τὰ σημεῖα ταῦτα ἐπʼ αὐτούς·
\vs{2}ὅπως διηγήσησθε εἰς τὰ ὦτα τῶν τέκνων ὑμῶν, καὶ τοῖς τέκνοις τῶν τέκνων ὑμῶν, ὅσα ἐμπέπαιχα τοῖς Αἰγυπτίοις, καὶ τὰ σημεῖά μου, ἃ ἐποίησα ἐν αὐτοῖς· καὶ γνώσεσθε ὅτι ἐγὼ Κύριος.
\vs{3}Εἰσῆλθε δὲ Μωυσῆς καὶ Ἀαρὼν ἐναντίον Φαραὼ, καὶ εἶπαν αὐτῷ, τάδε λέγει Κύριος ὁ Θεὸς τῶν Ἐβραίων, ἕως τίνος οὐ βούλει ἐντραπῆναί με; ἔξαπόστειλον τὸν λαόν μου, ἵνα λατρεύσωσί μοι.
\vs{4}Ἐὰν δὲ μὴ θέλῃς σὺ ἐξαποστεῖλαι τὸν λαόν μου, ἰδοὺ ἐγὼ ἐπάγω ταύτην τὴν ὥραν αὔριον ἀκρίδα πολλὴν ἐπὶ πάντα τὰ ὅριά σου.
\vs{5}Καὶ καλύψει τὴν ὄψιν τῆς γῆς, καὶ οὐ δυνήσῃ κατιδεῖν τὴν γῆν· καὶ κατέδεται πᾶν τὸ περισσὸν τῆς γῆς τὸ καταλειφθὲν, ὃ κατέλιπεν ὑμῖν ἡ χάλαζα, καὶ κατέδεται πᾶν ξύλον τὸ φυόμενον ὑμῖν ἐπὶ τῆς γῆς.
\vs{6}Καὶ πλησθήσονταί σου αἱ οἰκίαι, καὶ αἱ οἰκίαι τῶν θεραπόντων σου, καὶ πᾶσαι αἱ οἰκίαι ἐν πάσῃ γῇ τῶν Αἰγυπτίων· ἃ οὐδέποτε ἑωράκασιν οἱ πατέρες σου, οὐδʼ οἱ πρόπαπποι αὐτῶν, ἀφʼ ἧς ἡμέρας γεγόνασιν ἐπὶ τῆς γῆς, ἔως τῆς ἡμέρας ταύτης· καὶ ἐκκλίνας Μωυσῆς ἐξῆλθεν ἀπὸ Φαραώ.
\vs{7}Καὶ λέγουσιν οἱ θεράποντες Φαραὼ πρὸς αὐτὸν, ἕως τίνος ἔσται τοῦτο ἡμῖν σκῶλον; ἐξαπόστειλον τοὺς ἀνθρώπους, ὅπως λατρεύσωσι τῷ Θεῷ αὐτῶν· ἢ εἰδέναι βούλει ὅτι ἀπόλωλεν Αἴγυπτος;
\vs{8}Καὶ ἀπέστρεψαν τόν τε Μωυσῆν καὶ Ἀαρὼν πρὸς Φαραὼ, καὶ εἶπεν αὐτοῖς, πορεύεσθε καὶ λατρεύσατε Κυρίῳ τῷ Θεῷ ὑμῶν· τίνες δὲ καὶ τίνες εἰσιν οἱ πορευόμενοι;
\vs{9}Καὶ λέγει Μωυσῆς, σὺν τοῖς νεανίσκοις καὶ πρεσβυτέροις πορευσόμεθα, σὺν τοῖς υἱοῖς καὶ θυγατράσι, καὶ προβάτοις, καὶ βουσὶν ἡμῶν· ἔστι γὰρ ἑορτὴ Κυρίου.
\vs{10}Καὶ εἶπε πρὸς αὐτοὺς, ἔστω οὕτω Κύριος μεθʼ ὑμῶν· καθότι ἀποστέλλω ὑμᾶς, μὴ καὶ τὴν ἀποσκευὴν ὑμῶν; ἴδετε ὅτι πονηρία πρόσκειται ὑμῖν.
\vs{11}Μὴ οὕτως· πορευέσθωσαν δὲ οἱ ἄνδρες, καὶ λατρευσάτωσαν τῷ Θεῷ· τοῦτο γὰρ αὐτοὶ ἐκζητεῖτε· ἐξέβαλον δὲ αὐτοὺς ἀπὸ προσώπου Φαραώ.
\vs{12}Εἶπε δὲ Κύριος πρὸς Μωυσῆν, ἔκτεινον τὴν χεῖρα ἐπὶ γῆν Αἰγύπτου· καὶ ἀναβήτω ἀκρὶς ἐπὶ τὴν γῆν, καὶ κατέδεται πᾶσαν βοτάνην τῆς γῆς, καὶ πάντα τὸν καρπὸν τῶν ξύλων, ὃν ὑπελίπετο ἡ χάλαζα.
\vs{13}Καὶ ἐπῇρε Μωυσῆς τὴν ῥάβδον εἰς τὸν οὐρανὸν, καὶ Κύριος ἐπήγαγεν ἄνεμον νότον ἐπὶ τὴν γῆν, ὅλην τὴν ἡμέραν ἐκείνην, καὶ ὅλην τὴν νύκτα· τὸ πρωῒ ἐγενήθη, καὶ ὁ ἄνεμος ὁ νότος ἀνέλαβεν τὴν ἀκρίδα,
\vs{14}καὶ ἀνήγαγεν αὐτὴν ἐπὶ πᾶσαν γῆν Αἰγύπτου· καὶ κατέπαυσεν ἐπὶ πάντα τὰ ὅρια Αἰγύπτου πολλὴ σφόδρα· προτέρα αὐτῆς οὐ γέγονε τοιαύτη ἀκρὶς, καὶ μετὰ ταῦτα οὐκ ἔσται οὕτως.
\vs{15}Καὶ ἐκάλυψε τὴν ὄψιν τῆς γῆς, καὶ ἐφθάρη ἡ γῆ· καὶ κατέφαγε πᾶσαν βοτάνην τῆς γῆς, καὶ πάντα τόν καρπὸν τῶν ξύλων, ὃς ὑπελείφθη ἀπὸ τῆς χαλάζης· οὐχ ὑπελείφθη χλωρὸν οὐδὲν ἐν τοῖς ξύλοις, καὶ ἐν πάσῃ βοτάνῃ τοῦ πεδίου, ἐν πάσηῃ γῇ Αἰγύπτου.

\vs{16}Κατέσπευδε δὲ Φαραὼ καλέσαι Μωυσῆν καὶ Ἀαρὼν, λέγων, ἡμάρτηκα ἐναντίον Κυρίου τοῦ Θεοῦ ὑμῶν, καὶ εἰς ὑμᾶς.
\vs{17}Προσδέξασθε οὖν μου τὴν ἁμαρτίαν ἔτι νῦν, καὶ προσεύξασθε πρὸς Κύριον τὸν Θεὸν ὑμῶν, καὶ περιελέτω ἀπʼ ἐμοῦ τὸν θάνατον τοῦτον.
\vs{18}Ἐξῆλθε δὲ Μωυσῆς ἀπὸ Φαραὼ, καὶ ηὔξατο πρὸς τὸν Θεόν.
\vs{19}Καὶ μετέβαλε Κύριος ἄνεμον ἀπὸ θαλάσσης σφοδρὸν, καὶ ἀνέλαβε τὴν ἀκρίδα, καὶ ἔβαλεν αὐτὴν εἰς τὴν ἐρυθρὰν θαλάσσαν· καὶ οὐχ ὑπελείφη ἀκρὶς μία ἐν πάσῃ γῇ Αἰγύπτου.
\vs{20}Καὶ ἐσκλήρυνε Κύριος τὴν καρδίαν Φαραὼ, καὶ οὐκ ἐξαπέστειλε τοὺς υἱοὺς Ἰσραήλ.
\vs{21}Εἶπε δὲ Κύριος πρὸς Μωυσῆν, ἔκτεινον τὴν χεῖρά σου εἰς τὸν οὐρανὸν, καὶ γενηθήτω σκότος ἐπὶ γῆς Αἰγύπτου, ψηλαφητὸν σκότος.
\vs{22}Ἐξέτεινε δὲ Μωυσῆς τὴν χεῖρα εἰς τὸν οὐρανόν· καὶ ἐγένετο σκότος γνόφος, θύελλα ἐπὶ πᾶσαν γῆν Αἰγύπτου τρεῖς ἡμέρας·
\vs{23}Καὶ οὐκ εἶδεν οὐδεὶς τὸν ἀδελφὸν αὐτοῦ τρεῖς ἡμέρας· καὶ οὐκ ἐξανέστη οὐδεὶς ἐκ τῆς κοίτης αὐτοῦ τρεῖς ἡμέρας· πᾶσι δὲ τοῖς υἱοῖς Ἰσραὴλ φῶς ἦν ἐν πᾶσιν οἷς κατεγίνοντο.
\vs{24}Καὶ ἐκάλεσε Φαραὼ Μωυσῆν καὶ Ἀαρὼν, λέγων, Βαδίζετε, λατρεύσατε Κυρίῳ τῷ Θεῷ ὑμῶν, πλὴν τῶν προβάτων καὶ τῶν βοῶν ὑπολείπεσθε· καὶ ἡ ἀποσκευὴ ὑμῶν ἀποτρεχέτω μεθʼ ὑμῶν.
\vs{25}Καὶ εἶπε Μωυσῆς, ἀλλὰ καὶ σὺ δώσεις ἡμῖν ὁλοκαυτώματα καὶ θυσίας, ἂ ποιήσομεν Κυρίῳ τῷ Θεῷ ἡμῶν.
\vs{26}Καὶ τὰ κτήνη ἡμῶν πορεύσεται μεθʼ ἡμῶν, καὶ οὐχ ὑπολειψόμεθα ὁπλήν· ἀπʼ αὐτῶν γὰρ ληψόμεθα λατρεῦσαι Κυρίῳ τῷ Θεῷ ἡμῶν· ἡμεῖς δὲ οὐκ οἴδαμεν τί λατρεύσομεν Κυρίῳ τῷ Θεῷ ἡμῶν, ἕως τοῦ ἐλθεῖν ἡμᾶς ἐκεῖ.
\vs{27}Ἐσκλήρυνε δὲ Κύριος τὴν καρδίαν Φαραὼ, καὶ οὐκ ἐβουλήθη ἐξαποστεῖλαι αὐτούς.
\vs{28}Καὶ λέγει Φαραὼ, ἄπελθε ἀπʼ ἐμοῦ· πρόσεχε σεαυτῷ ἔτι προσθεῖναι ἰδεῖν μου τὸ πρόσωπον· ᾗ δʼ ἂν ἡμέρᾳ ὀφθῇς μοι, ἀποθανῇ.
\vs{29}Λέγει δὲ Μωυσῆς, εἴρηκας· οὐκ ἔτι ὀφθήσομαί σοι εἰς πρόσωπον.

\ch{11}
Εἶπε δὲ Κύριος πρὸς Μωυσῆν, ἔτι μίαν πληγὴν ἐγὼ ἐπάξω ἐπὶ Φαραὼ, καὶ ἐπʼ Αἴγυπτον, καὶ μετὰ ταῦτα ἐξαποστελεῖ ὑμᾶς ἐντεῦθεν· ὅταν δὲ ἐξαποστέλλῃ ὑμᾶς σὺν παντὶ, ἐκβαλεῖ ὑμᾶς ἐκβολῇ.
\vs{2}Λάλησον οὖν κρυφῇ εἰς τὰ ὦτα τοῦ λαοῦ, καὶ αἰτησάτω ἕκαστος παρὰ τοῦ πλησίον σκεύη ἀργυρᾶ καὶ χρυσὰ καὶ ἱματισμόν.
\vs{3}Κύριος δὲ ἔδωκε τὴν χάριν τῷ λαῷ αὐτοῦ ἐναντίον τῶν Αἰγυπτίων, καὶ ἔχρησαν αὐτοῖς· καὶ ὁ ἄνθρωπος Μωυσῆς μέγας ἐγενήθη σφόδρα ἐναντίον τῶν Αἰγυπτίων, καὶ ἐναντίον Φαραὼ, καὶ ἐναντίον τῶν θεραπόντων ἀτοῦ.
\vs{4}Καὶ εἶπε Μωυσῆς, τάδε λέγει Κύριος, περὶ μέσας νύκτας ἐγὼ εἰσπορεύομαι εἰς μέσον Αἰγύπτου·
\vs{5}Καὶ τελευτήσει πᾶν πρωτότοκον ἐν γῇ Αἰγύπτῳ, ἀπὸ πρωτοτόκου Φαραὼ, ὃς κάθηται ἐπὶ τοῦ θρόνου, καὶ ἕως πρωτοτόκου τῆς θεραπαίνης τῆς παρὰ τὸν μύλον, καὶ ἕως πρωτοτοκου παντος κτήνους.
\vs{6}Καὶ ἔσται κραυγὴ μεγάλη κατὰ πᾶσαν γῆν Αἰγύπτου, ἥτις τοιαύτη οὐ γέγονε, καὶ τοιαύτη οὐκ ἔτι προστεθήσεται.
\vs{7}Καὶ ἐν πᾶσι τοῖς υἱοῖς Ἰσραὴλ οὐ γρύξει κύων τῇ γλώσσῃ αὐτοῦ, ἀπὸ ἀνθρώπου ἕως κτήνους· ὅπως εἰδῇς ὅσα παραδοξάσει Κύριος ἀνὰ μέσον τῶν Αἰγυπτίων καὶ τοῦ Ἰσραήλ.
\vs{8}Καὶ καταβήσονται πάντες οἱ παῖδές σου οὗτοι πρός με, καὶ προσκυνήσουσί με, λέγοντες, ἔξελθε σὺ, καὶ πᾶς ὁ λαός σου, οὗ σὺ ἀφηγῇ· καὶ μετὰ ταῦτα ἐξελεύσομαι· ἐξῆλθε δὲ Μωυσῆς ἀπὸ Φαραὼ μετὰ θυμοῦ.
\vs{9}Εἶπε δὲ Κύριος πρὸς Μωυσῆν, οὐκ εἰσακούσεται ὑμῶν Φαραὼ, ἵνα πληθύνων πληθυνῶ μου τὰ σημεῖα, καὶ τὰ τέρατα ἐν γῇ Αἰγύπτῳ.
\vs{10}Μωσῆς δὲ καὶ Ἀαρὼν ἐποίησαν πάντα τὰ σημεῖα καὶ τὰ τέρατα ταῦτα ἐν γῇ Αἰγύπτῳ ἐναντίον Φαραώ· ἐσκλήρυνε δὲ Κύριος τὴν καρδίαν Φαραὼ, καὶ οὐκ εἰσήκουσεν ἐξαποστεῖλαι τοὺς υἱοὺς Ἰσραὴλ ἐκ γῆς Αἰγύπτου.

\ch{12}
Εἶπε δὲ Κύριος πρὸς Μευσῆν καὶ Ἀαρὼν ἐν γῇ Αἰγύπτου, λέγων,
\vs{2}ὁ μὴν οὗτος ὑμῖν ἀρχὴ μηνῶν· πρῶτός ἐστιν ὑμῖν ἐν τοῖς μησὶ τοῦ ἐνιαυτοῦ.
\vs{3}Λάλησον πρὸς πᾶσαν συναγωγὴν υἱῶν Ἰσραὴλ, λέγων, τῇ δεκάτῃ τοῦ μηνὸς τούτου λαβέτωσαν ἕκαστος πρόβατον κατʼ οἴκους πατριῶν, ἕκαστος πρόβατον κατʼ οἰκίαν.
\vs{4}Ἐὰν δὲ ὀλιγοστοὶ ὦσιν ἐν τῇ οἰκίᾳ, ὥστε μὴ εἶναι ἱκανοὺς εἰς πρόβατον, συλλήψεται μεθʼ ἑαυτοῦ τὸν γείτονα τὸν πλησίον αὐτοῦ· κατὰ ἀριθμὸν ψυχῶν, ἕκαστος τὸ ἀρκοῦν αὐτῷ συναριθμήσεται εἰς πρόβατον.
\vs{5}Πρόβατον τέλειον, ἄρσεν, ἐνιαύσιον ἔσται ὑμῖν· ἀπὸ τῶν ἀρνῶν καὶ τῶν ἐρίφων λὴψεσθε.
\vs{6}Καὶ ἔσται ὑμῖν διατετηρημένον ἕως τῆς τεσσαρεσκαιδεκάτης τοῦ μηνὸς τούτου· καὶ σφάξουσιν αὐτὸ πᾶν τὸ πλῆθος συναγωγῆς υἱῶν Ἰσραὴλ πρὸς ἑσπέραν.
\vs{7}Καὶ λήψονται ἀπὸ τοῦ αἵματος, καὶ θήσουσιν ἐπὶ τῶν δύο σταθμῶν καὶ ἐπὶ τὴν φλιὰν, ἐν τοῖς οἴκοις ἐν οἷς ἐὰν φάγωσιν αὐτὰ ἐν αὐτοῖς.
\vs{8}Καὶ φάγονται τὰ κρέα τῇ νυκτὶ ταύτῃ ὀπτὰ πυρὶ, καὶ ἄζυμα ἐπὶ πικρίδων ἔδονται.
\vs{9}Οὐκ ἔδεσθε ἀπʼ αὐτῶν ὠμὸν, οὐδὲ ἡψημένον ἐν ὕδατι, ἀλλʼ ἢ ὀπτὰ πυρὶ, κεφαλὴν σὺν τοῖς ποσὶ καὶ τοῖς ἐνδοσθίοις.
\vs{10}Οὐκ ἀπολείψεται ἀπʼ αὐτοῦ ἕως πρωΐ· καὶ ὀστοῦν οὐ συντρίψετε ἀπʼ αὐτοῦ· τὰ δὲ καταλειπόμενα ἀπʼ αὐτοῦ ἕως πρωῒ ἐν πυρὶ κατακαύσετε.
\vs{11}Οὕτω δὲ φάγεσθε αὐτό· αἱ ὀσφύες ὑμῶν περιεζωσμέναι, καὶ τὰ ὑποδήματα ἐν τοῖς ποσὶν ὑμῶν, καὶ αἱ βακτηρίαι ἐν ταῖς χερσὶν ὑμῶν· καὶ ἔδεσθε αὐτὸ μετὰ σπουδῆς· Πάσχα ἐστὶ Κυρίῳ.
\vs{12}Καὶ διελεύσομαι ἐν γῇ Αἰγύπτῳ ἐν τῇ νυκτὶ ταύτῃ, καὶ πατάξω πᾶν πρωτότοκον ἐν γῇ Αἰγύπτῳ ἀπὸ ἀνθρώπου ἕως κτήνους· καὶ ἐν πᾶσι τοῖς θεοῖς τῶν Αἰγυπτίων ποιήσω τὴν ἐκδίκησιν· ἐγὼ Κύριος.
\vs{13}Καὶ ἔσται τὸ αἷμα ὑμῖν ἐν σημείῳ ἐπὶ τῶν οἰκιῶν, ἐν αἷς ὑμεῖς ἔστε ἐκεῖ· καὶ ὄψομαι τὸ αἷμα, καὶ σκεπάσω ὑμᾶς, καὶ οὐκ ἔσται ἐν ὑμῖν πληγὴ τοῦ ἐκτριβῆναι ὅταν παίω ἐν γῇ Αἰγύπτῳ.

\vs{14}Καὶ ἔσται ἡ ἡμέρα ὑμῖν αὕτη μνημόσυνον, καὶ ἑορτάσετε αὐτὴν ἑορτὴν Κυρίῳ εἰς πάσας τὰς γενεὰς ὑμῶν· νόμιμον αἰώνιον ἑορτάσετε αὐτήν.
\vs{15}Ἑπτὰ ἡμέρας ἄζυμα ἔδεσθε· ἀπὸ δὲ τῆς ἡμέρας τῆς πρώτης, ἀφανιεῖτε ζύμην ἐκ τῶν οἰκιῶν ὑμῶν· πᾶς ὃς ἂν φάγῃ ζύμην, ἐξολοθρευθήσεται ἡ ψυχὴ ἐκείνη ἐξ Ἰσραήλ, ἀπὸ τῆς ἡμέρας τῆς πρώτης ἕως τῆς ἡμέρας τῆς ἑβδόμης.
\vs{16}Καὶ ἡ ἡμέρα ἡ πρώτη, κληθήσεται ἁγία· καὶ ἡ ἡμέρα ἡ ἑβδόμη, κλητὴ ἁγία ἔσται ὑμῖν· πᾶν ἔργον λατρευτὸν οὐ ποιήσετε ἐν αὐταῖς, πλὴν ὅσα ποιηθήσεται πάσῃ ψυχῇ, τοῦτο μόνον ποιηθήσεται ὑμῖν.
\vs{17}Καὶ φυλάξετε τὴν ἐντολὴν ταύτην· ἐν γὰρ τῇ ἡμέρᾳ ταύτῃ ἐξάξω τὴν δύναμιν ὑμῶν ἐκ γῆς Αἰγύπτου, καὶ ποιήσετε τὴν ἡμέραν ταύτην εἰς γενεὰς ὑμῶν νόμιμον αἰώνιον,
\vs{18}ἐναρχόμενοι τῇ τεσσαρεσκαιδεκάτῃ ἡμέρᾳ τοῦ μηνὸς τοῦ πρώτου, ἀφʼ ἑσπέρας ἔδεσθε ἄζυμα, ἕως ἡμέρας μίας καὶ εἰκάδος τοῦ μηνὸς, ἕως ἑσπέρας.
\vs{19}Ἑπτὰ ἡμέρας ζύμη οὐχ εὑρεθήσεται ἐν ταῖς οἰκιαῖς ὑμῶν· πᾶς ὃς ἂν φάγῃ ζυμωτὸν, ἐξολοθρευθήσεται ἡ ψυχὴ ἐκείνη ἐκ συναγωγῆς Ἰσραήλ· ἔν τε τοῖς γειώραις, καὶ αὐτόχθοσι τῆς γῆς.
\vs{20}Πᾶν ζυμωτὸν οὐκ ἔδεσθε, ἐν παντὶ δὲ κατοικητηρίῳ ὑμῶν ἔδεσθε ἄζυμα.

\vs{21}Ἐκάλεσε δὲ Μωυσῆς πᾶσαν γερουσίαν υἱῶν Ἰσραὴλ, καὶ εἶπε πρὸς αὐτοὺς, ἀπελθόντες λάβετε ὑμῖν αὐτοῖς πρόβατον κατὰ συγγενείας ὑμῶν, καὶ θύσατε τὸ πάσχα.
\vs{22}Λήψεσθε δὲ δέσμην ὑσσώπου, καὶ βάψαντες ἀπὸ τοῦ αἵματος τοῦ παρὰ τὴν θύραν, καθίξετε τῆς φλιᾶς, καὶ ἐπʼ ἀμφοτέρων τῶν σταθμῶν, ἀπὸ τοῦ αἵματος ὅ ἐστι παρὰ τὴν θύραν· ὑμεῖς δὲ οὐκ ἐξελεύσεσθε ἕκαστος τὴν θύραν τοῦ οἴκου αὐτοῦ ἕως πρωΐ.
\vs{23}Καὶ παρελεύσεται Κύριος πατάξαι τοὺς Αἰγυπτίους, καὶ ὄψεται τὸ αἷμα ἐπὶ τῆς φλιᾶς, καὶ ἐπʼ ἀμφοτέρων τῶν σταθμῶν· καὶ παρελεύσεται Κύριος τὴν θύραν, καὶ οὐκ ἀφήσει τὸν ὀλοθρεύοντα εἰσελθεῖν εἰς τὰς οἰκίας ὑμῶν πατάξαι.
\vs{24}Καὶ φυλάξασθε τὸ ῥῆμα τοῦτο νόμιμον σεαυτῷ, καὶ τοῖς υἱοῖς σου, ἕως αἰῶνος.
\vs{25}Ἐὰν δὲ εἰσέλθητε εἰς τὴν γῆν, ἣν ἂν δῷ Κύριος ὑμῖν, καθότι ἐλάλησε, φυλάξασθε τὴν λατρείαν ταύτην.
\vs{26}Καὶ ἐσται ἐὰν λέγωσι πρὸς ὑμᾶς οἱ υἱοὶ ὑμῶν, τίς ἡ λατρεία αὕτη;
\vs{27}Καὶ ἐρεῖτε αὐτοῖς, θυσία τὸ πάσχα τοῦτο Κυρίῳ, ὡς ἐσκέπασε τοὺς οἴκους τῶν υἱῶν Ἰσραὴλ ἐν Αἰγύπτῳ, ἡνίκα ἐπάταξε τοὺς Αἰγυπτίους, τοὺς δὲ οἴκους ἡμῶν ἐῤῥύσατο· καὶ κύψας ὁ λαὸς προσεκύνησε.
\vs{28}Καὶ ἀπελθόντες ἐποίησαν οἱ υἱοὶ Ἰσραὴλ, καθὰ ἐνετείλατο Κύριος τῷ Μωυσῇ καὶ Ἀαρῶν, οὕτως ἐποίησαν.

\vs{29}Ἐγενήθη δὲ μεσούσης τῆς νυκτὸς, καὶ Κύριος ἐπάταξε πᾶν πρωτότοκον ἐν γῇ Αἰγύπτῳ, ἀπὸ πρωτοτόκου Φαραὼ τοῦ καθημένου ἐπὶ τοῦ θρόνου, ἕως πρωτοτόκου τῆς αἰχμαλωτίδος τῆς ἐν τῷ λάκκῳ, καὶ ἕως πρωτοτόκου παντὸς κτήνους.
\vs{30}Καὶ ἀναστὰς Φαραὼ νυκτὸς, καὶ οἱ θεράποντες αὐτοῦ, καὶ πάντες οἱ Αἰγύπτιοι, καὶ ἔγενήθη κραυγὴ μεγάλη ἐν πάσῃ γῇ Αἰγύπτῳ· οὐ γὰρ ἦν οἰκία, ἐν ᾗ οὐκ ἦν ἐν αὐτῇ τεθνηκώς.
\vs{31}Καὶ ἐκάλεσε Φαραὼ Μωυσῆν καὶ Ἀαρὼν νυκτὸς, καὶ εἶπεν αὐτοῖς, ἀνάστητε, καὶ ἐξέλθατε ἐκ τοῦ λαοῦ μου, καὶ ὑμεῖς, καὶ οἱ υἱοὶ Ἰσραήλ· βαδίζετε καὶ λατρεύσατε Κυρίῳ τῷ Θεῷ ὑμῶν, καθὰ λέγετε.
\vs{32}Καὶ τὰ πρόβατα καὶ τοὺς βόας ὑμῶν ἀναλαβόντες πορεύεσθε· εὐλογήσατε δὴ κᾀμέ.
\vs{33}Καὶ κατεβιάζοντο οἱ Αἰγύπτιοι τὸν λαὸν σπουδῇ ἐκβαλεῖν αὐτοὺς ἐν τῆς γῆς· εἶπαν γὰρ, ὅτι πάντες ἡμεῖς ἀποθνήσκομεν.
\vs{34}Ἀνέλαβε δὲ ὁ λαὸς τὸ σταῖς αὐτῶν, πρὸ τοῦ ζυμωθῆναι τὰ φυράματα αὐτῶν, ἐνδεδεμένα ἐν τοῖς ἱματίοις αὐτῶν ἐπὶ τῶν ὤμαν.
\vs{35}Οἱ δὲ υἱοὶ Ἰσραὴλ ἐποίησαν, καθὰ συνέταξεν αὐτοῖς Μωυσῆς, καὶ ᾔτησαν παρὰ τῶν Αἰγυπτίων σκεύη ἀργυρᾶ καὶ χρυσᾶ καὶ ἱματισμόν.
\vs{36}Καὶ ἔδωκε Κύριος τὴν χάριν τῷ λαῷ αὐτοῦ ἐναντίον τῶν Αἰγυπτίων, καὶ ἔχρησαν αὐτοῖς· καὶ ἐσκύλευσαν τοὺς Αἰγυπτίους.

\vs{37}Ἀπάραντες δὲ υἱοὶ Ἰσραὴλ ἐκ Ῥαμεσσῆ εἰς Σοκχὼθ εἰς ἑξακοσίας χιλιάδας πεζῶν, οἱ ἄνδρες, πλὴν τῆς ἀποσκευῆς.
\vs{38}Καὶ ἐπίμικτος πολὺς συνανέβη αὐτοῖς, καὶ πρόβατα, καὶ βόες, καὶ κτήνη πολλὰ σφόδρα.
\vs{39}Καὶ ἔπεψαν τὸ σταῖς ὃ ἐξήνεγκαν ἐξ Αἰγύπτου, ἐγκρυφίας ἀζύμους, οὐ γὰρ ἐζυμώθη· ἐξέβαλον γὰρ αὐτοὺς οἱ Αἰγύπτιοι, καὶ οὐκ ἠδυνήθησαν ἐπιμεῖναι, οὐδὲ ἐπισιτισμὸν ἐποίησαν ἑαυτοῖς εἰς τὴν ὁδόν.
\vs{40}Ἡ δὲ κατοίκησις τῶν υἱῶν Ἰσραὴλ, ἣν κατῴκησαν ἐν γῇ Αἰγύπτῳ καὶ ἐν γῇ Χαναὰν, ἔτη τετρακόσια τριάκοντα.
\vs{41}Καὶ ἐγένετο μετὰ τὰ τετρακόσια τριάκοντα ἔτη, ἐξῆλθε πᾶσα ἡ δύναμις Κυρίου ἐκ γῆς Αἰγύπτου νυκτός.
\vs{42}Προφυλακή ἐστι τῷ Κυρίῳ, ὥστε ἐξαγαγεῖν αὐτοὺς ἐκ γῆς Αἰγύπτου· ἐκείνη ἡ νὺξ αὕτη, προφυλακὴ Κυρίῳ, ὥστε πᾶσι τοῖς υἱοῖς Ἰσραὴλ εἶναι εἰς γενεὰς αὐτῶν.
\vs{43}Εἶπε δὲ Κύριος πρὸς Μωυσῆν καὶ Ἀαρὼν, οὗτος ὁ νόμος τοῦ πάσχα· πᾶς ἀλλογενὴς οὐκ ἔδεται ἀπʼ αὐτοῦ·
\vs{44}Καὶ πάντα οἰκέτην ἢ ἀργυρώνητον περιτεμεῖς αὐτόν· καὶ τότε φάγεται ἀπʼ αὐτοῦ.
\vs{45}Πάροικος ἢ μισθωτὸς οὐκ ἔδεται ἀπʼ αὐτοῦ.
\vs{46}Ἐν οἰκίᾳ μιᾷ βρωθήσεται, καὶ οὐκ ἐξοίσετε ἐκ τῆς οἰκίας τῶν κρεῶν ἔξω· καὶ ὀστοῦν οὐ συντρίψετε ἀπʼ αὐτοῦ.
\vs{47}Πᾶσα συναγωγὴ υἱῶν Ἰσραὴλ ποιήσει αὐτό.
\vs{48}Ἐὰν δέ τις προσέλθῃ πρὸς ὑμᾶς προσήλυτος ποιῆσαι τὸ πάσχα Κυρίῳ, περιτεμεῖς αὐτοῦ πᾶν ἀρσενικόν, καὶ τότε προσελεύσεται ποιῆσαι αὐτό· καὶ ἔσται ὥσπερ καὶ ὁ αὐτόχθων τῆς γῆς· πᾶς ἀπερίτμητος οὐκ ἔδεται ἀπʼ αὐτοῦ.
\vs{49}Νόμος εἷς ἔσται τῷ ἐγχωρίῳ, καὶ τῷ προσελθόντι προσηλύτῳ ἐν ὑμῖν.
\vs{50}Καὶ ἐποίησαν οἱ υἱοὶ Ἰσραὴλ καθὰ ἐνετείλατο Κύριος τῷ Μωυσῇ καὶ Ἀαρὼν πρὸς αὐτούς, οὕτως ἐποίησαν.
\vs{51}Καὶ ἐγένετο ἐν τῇ ἡμέρᾳ ἐκείνῃ, ἐξήγαγε Κύριος τοὺς υἱοὺς Ἰσραὴλ ἐκ γῆς Αἰγύπτου σὺν δυνάμει αὐτῶν.

\ch{13}
Εἶπε δὲ Κύριος πρὸς Μωυσῆν, λέγων,
\vs{2}ἁγίασόν μοι πᾶν πρωτότοκον πρωτογενὲς διανοῖγον πᾶσαν μήτραν ἐν τοῖς υἱοῖς Ἰσραὴλ ἀπὸ ἀνθρώπου ἕως κτήνους, ἐμοί ἐστιν.
\vs{3}Εἶπε δὲ Μωυσῆς πρὸς τὸν λαὸν, μνημονεύετε τὴν ἡμέραν ταύτην, ἐν ᾗ ἐξήλθατε ἐκ γῆς Αἰγύπτου, ἐξ οἴκου δουλείας· ἐν γὰρ χειρὶ κραταιᾷ ἐξήγαγεν ὑμᾶς Κύριος ἐντεῦθεν· καὶ οὐ βρωθήσεται ζύμη.
\vs{4}Ἐν γὰρ τῇ σήμερον ὑμεῖς ἐκπορεύεσθε ἐν μηνὶ τῶν νέων.
\vs{5}Καὶ ἔσται ἡνίκα ἐὰν εἰσαγάγῃ σε Κύριος ὁ Θεός σου εἰς τὴν γῆν τῶν Χαναναίων, καὶ Χετταίων, καὶ Ἀμοῤῥαίων, καὶ Εὐαίων, καὶ Ἰεβουσαίων, καὶ Γεργεσαίων, καὶ Φερεζαίων, ἣν ὤμοσε τοῖς πατράσι σου, δοῦναί σοι γῆν ῥέουσαν γάλα καὶ μέλι· καὶ ποιήσεις τὴν λατρείαν ταύτην ἐν τῷ μηνὶ τούτῳ.
\vs{6}Ἓξ ἡμέρας ἔδεσθε ἄζυμα, τῇ δὲ ἡμέρᾳ τῇ ἑβδόμῃ ἑορτὴ Κυρίου.
\vs{7}Ἄζυμα ἔδεσθε ἑπτὰ ἡμέρας· οὐκ ὀφθήσεταί σοι ζυμωτὸν, οὐδὲ ἔσται σοι ζύμη ἐν πᾶσι τοῖς ὁρίοις σου.
\vs{8}Καὶ ἀναγγελεῖς τῷ υἱῷ σου ἐν τῇ ἡμέρᾳ ἐκείνῃ, λέγων, διὰ τοῦτο ἐποίησε Κύριος ὁ Θεός μοι, ὡς ἐξεπορευόμην ἐξ Αἰγύπτου.
\vs{9}Καὶ ἔσται σοι σημεῖον ἐπὶ τῆς χειρός σου, καὶ μνημόσυνον πρὸ ὀφθαλμῶν σου, ὅπως ἂν γένηται ὁ νόμος Κυρίου ἐν τῷ στόματί σου· ἐν γὰρ χειρὶ κραταιᾷ ἐξήγαγέ σε Κύριος ὁ Θεὸς ἐξ Αἰγύπτου.
\vs{10}Καὶ φυλάξασθε τὸν νόμον τοῦτον κατὰ καιροὺς ὡρῶν, ἀφʼ ἡμερῶν εἰς ἡμέρας.

\vs{11}Καὶ ἔσται ὡς ἂν εἰσαγάγῃ σε Κύριος ὁ Θεός σου εἰς τὴν γῆν τῶν Χαναναίων, ὃν τρόπον ὤμοσε τοῖς πατράσι σου, καὶ δώσει σοι αὐτήν.
\vs{12}Καὶ ἀφελεῖς πᾶν διανοῖγον μήτραν, τὰ ἀρσενικὰ τῷ Κυρίῳ· πᾶν διανοῖγον μήτραν ἐκ βουκολίων ἢ ἐν τοῖς κτήνεσί σου, ὅσα ἐὰν γένηταί σοι, τὰ ἀρσενικὰ ἁγιάσεις τῷ Κυρίῳ.
\vs{13}Πᾶν διανοῖγον μήτραν ὄνου, ἀλλάξεις προβάτῳ· ἐὰν δὲ μὴ ἀλλάξῃς, λυτρώσῃ αὐτό· πᾶν πρωτότοκον ἀνθρώπου τῶν υἱῶν σου λυτρώσῃ.
\vs{14}Ἐὰν δὲ ἐρωτήσῃ σε ὁ υἱός σου μετὰ ταῦτα, λέγων, τί τοῦτο; καὶ ἐρεῖς αὐτῷ, ὅτι ἐν χειρὶ κραταιᾷ ἐξήγαγεν Κύριος ἡμᾶς ἐκ γῆς Αἰγύπτου, ἐξ οἴκου δουλείας.
\vs{15}Ἡνίκα δὲ ἐσκλήρυνε Φαραὼ ἐξαποστεῖλαι ἡμᾶς, ἀπέκτεινε πᾶν πρωτότοκον ἐν γῇ Αἰγύπτῳ, ἀπὸ πρωτοτόκων ἀνθρώπων ἕως πρωτοτόκων κτηνῶν· διὰ τοῦτο ἐγὼ θύω πᾶν διανοῖγον μήτραν, τὰ ἀρσενικὰ τῷ Κυρίῳ, καὶ πᾶν πρωτότοκον τῶν υἱῶν μου λυτρώσομαι.
\vs{16}Καὶ ἔσται εἰς σημεῖον ἐπὶ τῆς χειρός σου, καὶ ἀσαλευτον πρὸ ὀφθαλμων σου· ἐν γὰρ χειρὶ κραταιᾷ ἐξήγαγέ σε Κύριος ἐξ Αἰγύπτου.

\vs{17}Ὡς δὲ ἐξαπέστειλε Φαραὼ τὸν λαὸν, οὐχ ὡδήγησεν αὐτοὺς ὁ Θεὸς ὁδὸν γῆς· Φυλιστιεὶμ, ὅτι ἐγγὺς ἦν· εἶπε γὰρ ὁ Θεὸς, μήποτε μεταμελήσῃ τῷ λαῷ ἰδόντι πόλεμον, καὶ ἀποστρέψῃ εἰς Αἴγυπτον.
\vs{18}Καὶ ἐκύκλωσεν ὁ Θεὸς τὸν λαὸν ὁδὸν τὴν εἰς τὴν ἔρημον, εἰς τὴν ἐρυθρὰν θάλασσαν· πέμπτῃ δὲ γενεᾷ ἀνέβησαν οἱ υἱοὶ Ἰσραὴλ ἐκ γῆς Αἰγύπτου.
\vs{19}Καὶ ἔλαβε Μωυσῆς τὰ ὀστᾶ Ἰωσὴφ μεθʼ ἑαυτοῦ· ὅρκῳ γὰρ ὥρκισεν τοὺς υἱοὺς Ἰσραὴλ, λέγων, ἐπισκοπῇ ἐπισκέψεται ὑμᾶς Κύριος, καὶ συνανοίσετε μου τὰ ὀστᾶ ἐντεῦθεν μεθʼ ὑμῶν.

\vs{20}Ἐξάραντες δὲ οἱ υἱοὶ Ἰσραὴλ ἐκ Σοκχὼθ, ἐστρατοπέδευσαν ἐν Ὀθὼμ παρὰ τὴν ἔρημον.
\vs{21}Ὁ δὲ Θεὸς ἡγεῖτο αὐτῶν, ἡμέρας μὲν ἐν στύλῳ νεφέλης, δεῖξαι αὐτοῖς τὴν ὁδόν· τὴν δὲ νύκτα ἐν στύλῳ πυρός.
\vs{22}Οὐκ ἐξέλιπεν δὲ ὁ στύλος τῆς νεφέλης ἡμέρας, καὶ ὁ στύλος τοῦ πυρὸς νυκτὸς, ἐναντίον τοῦ λαοῦ παντός.

\ch{14}
Καὶ ἐλάλησε Κύριος πρὸς Μωυσῆν, λέγων,
\vs{2}Λάλησον τοῖς υἱοῖς Ἰσραὴλ, καὶ ἀποστρέψαντες στρατοπεδευσάτωσαν ἀπέναντι τῆς ἐπαύλεως, ἀνὰ μέσον Μαγδώλου καὶ ἀνὰ μέσον τῆς θαλάσσης, ἐξεναντίας Βεελσεπφῶν· ἐνώπιον αὐτῶν στρατοπεδεύσεις ἐπὶ τῆς θαλάσσης.
\vs{3}Καὶ ἐρεῖ Φαραὼ τῷ λαῷ αὐτοῦ, οἱ υἱοὶ Ἰσραὴλ πλανῶνται οὗτοι ἐν τῇ γῇ, συγκέκλεικε γὰρ αὐτοὺς ἡ ἔρημος.
\vs{4}Ἐγὼ δὲ σκληρυνῶ τὴν καρδίαν Φαραὼ, καὶ καταδιώξεται ὀπίσω αὐτῶν· καὶ ἐνδοξασθήσομαι ἐν Φαραῷ, καὶ ἐν πάσῃ τῇ στρατίᾳ αὐτοῦ· καὶ γνώσονται πάντες οἱ Αἰγύπτιοι ὅτι ἐγώ εἰμι Κύριος· καὶ ἐποίησαν οὕτως.
\vs{5}Καὶ ἀνηγγέλη τῷ βασιλεῖ τῶν Αἰγυπτίων ὅτι πέφευγεν ὁ λαός· καὶ μετεστράφη ἡ καρδία Φαραὼ, καὶ τῶν θεραπόντων αὐτοῦ, ἐπὶ τὸν λαὸν, καὶ εἶπαν, τί τοῦτο ἐποιήσαμεν, τοῦ ἐξαποστεῖλαι τοὺς υἱοὺς Ἰσραὴλ, τοῦ μὴ δουλεύειν ἡμῖν;
\vs{6}Ἔζευξεν οὖν Φαραὼ τὰ ἅρματα αὐτοῦ, καὶ πάντα τὸν λαὸν αὐτοῦ συναπήγαγε μεθʼ ἑαυτοῦ,
\vs{7}καὶ λαβὼν ἑξακόσια ἅρματα ἐκλεκτὰ, καὶ πᾶσαν τὴν ἵππον τῶν Αἰγυπτίων, καὶ τριστάτας ἐπὶ πάντων.
\vs{8}Καὶ ἐσκλήρυνε Κύριος τὴν καρδίαν Φαραὼ βασιλέως Αἰγύπτου, καὶ τῶν θεραπόντων αὐτοῦ, καὶ κατεδίωξεν ὀπίσω τῶν υἱῶν Ἰσραήλ· οἱ δὲ υἱοὶ Ἰσραὴλ ἐξεπορεύοντο ἐν χειρὶ ὑψηλῇ.
\vs{9}Καὶ κατεδίωξαν οἱ Αἰγύπτιοι ὀπίσω αὐτῶν, καὶ εὕροσαν αὐτοὺς παρεμβεβληκότας παρὰ τὴν θάλασσαν· καὶ πᾶσα ἡ ἵππος καὶ τὰ ἅρματα Φαραὼ, καὶ οἱ ἱππεῖς, καὶ ἡ στρατιὰ αὐτοῦ ἀπέναντι τῆς ἐπαύλεως, ἐξεναντίας Βεελσεπφῶν.
\vs{10}Καὶ Φαραὼ προσῆγε· καὶ ἀναβλέψαντες οἱ υἱοὶ Ἰσραὴλ τοῖς ὀφθαλμοῖς ὁρῶσι, καὶ οἱ Αἰγύπτιοι ἐστρατοπέδευσαν ὀπίσω αὐτῶν· καὶ ἐφοβήθησαν σφόδρα· ἀνεβόησαν δὲ οἱ υἱοὶ Ἰσραὴλ πρὸς Κύριον.
\vs{11}Καὶ εἶπαν πρὸς Μωυσῆν, παρὰ τὸ μὴ ὑπάρχειν μνήματα ἐν γῇ Αἰγύπτῳ, ἐξήγαγες ἡμᾶς θανατῶσαι ἐν τῇ ἐρήμῳ· τί τοῦτο ἐποίησας ἡμῖν, ἐξαγαγὼν ἐξ Αἰγύπτου;
\vs{12}Οὐ τοῦτο ἦν τὸ ῥῆμα, ὃ ἐλαλήσαμεν πρὸς σὲ ἐν Αἰγύπτῳ, λέγοντες, πάρες ἡμᾶς, ὅπως δουλεύσωμεν τοῖς Αἰγυπτίοις; κρεῖσσον γὰρ ἡμᾶς δουλεύειν τοῖς Αἰγυπτίοις, ἢ ἀποθανεῖν ἐν τῇ ἐρήμῳ ταύτῃ.

\vs{13}Εἶπε δὲ Μωυσῆς πρὸς τὸν λαὸν, θαρσεῖτε, στῆτε καὶ ὁρᾶτε τὴν σωτηρίαν τὴν παρὰ τοῦ Κυρίου, ἣν ποιήσει ἡμῖν σήμερον· ὃν τρόπον γὰρ ἑωράκατε τοὺς Αἰγυπτίους σήμερον, οὐ προσθήσεσθε ἔτι ἰδεῖν αὐτοὺς εἰς τὸν αἰῶνα χρόνον·
\vs{14}Κύριος πολεμήσει περὶ ὑμῶν, καὶ ὑμεῖς σιγήσετε.
\vs{15}Εἶπε δὲ Κύριος πρὸς Μωυσῆν, τί βοᾷς πρός με; λάλησον τοῖς υἱοῖς Ἰσραὴλ, καὶ ἀναζευξάτωσαν.
\vs{16}Καὶ σὺ ἔπαρον τῇ ῥάβδῳ σου, καὶ ἔκτεινον τὴν χεῖρά σου ἐπὶ τὴν θάλασσαν, καὶ ῥῆξον αὐτήν· καὶ εἰσελθάτωσαν οἱ υἱοὶ Ἰσραὴλ εἰς μέσον τῆς θαλάσσης κατὰ τὸ ξηρόν.
\vs{17}Καὶ ἰδοὺ ἐγὼ σκληρυνῶ τὴν καρδίαν Φαραὼ, καὶ τῶν Αἰγυπτίων πάντων, καὶ εἰσελεύσονται ὀπίσω αὐτῶν· καὶ ἐνδοξασθήσομαι ἐν Φαραῷ, καὶ ἐν πάσῃ τῇ στρατιᾷ αὐτοῦ, καὶ ἐν τοῖς ἅρμασι, καὶ ἐν τοῖς ἵπποις αὐτοῦ.
\vs{18}Καὶ γνώσονται πάντες οἱ Αἰγύπτιοι ὅτι ἐγώ εἰμὶ Κύριος, ἐνδοξαζομένου μου ἐν Φαραῷ, καὶ ἐν τοῖς ἅρμασι, καὶ ἵπποις αὐτοῦ.
\vs{19}Ἐξῇρε δὲ ὁ Ἄγγελος τοῦ Θεοῦ ὁ προπορευόμενος τῆς παρεμβολῆς τῶν υἱῶν Ἰσραὴλ, καὶ ἐπορεύθη ἐκ τῶν ὄπισθεν· ἐξῇρε δὲ καὶ ὁ στύλος τὴς νεφέλης ἀπὸ προσώπου αὐτῶν, καὶ ἔστη ἐκ τῶν ὀπίσω αὐτῶν.
\vs{20}Καὶ εἰσῆλθεν ἀνὰ μέσον τῆς παρεμβολῆς τῶν Αἰγυπτίων, καὶ ἀνὰ μέσον τῆς παρεμβολῆς τῶν Αἰγυπίων, καὶ ἀνὰ μέσον τῆς παρεμβολῆς Ἰσραὴλ, καὶ ἔστη· καὶ ἐγένετο σκότος καὶ γνόφος· καὶ διῆλθεν ἡ νύξ· καὶ οὐ συνέμιξαν ἀλλήλοις ὅλην τὴν νύκτα.
\vs{21}Ἐξέτεινε δὲ Μωυσῆς τὴν χεῖρα ἐπὶ τὴν θάλασσαν· καὶ ὑπήγαγε Κύριος τὴν θάλασσαν ἐν ἀνέμῳ νότῳ βιαίῳ ὅλην τὴν νύκτα, καὶ ἐποίησε τὴν θάλασσαν ξηράν· καὶ ἐσχίσθη τὸ ὕδωρ.
\vs{22}Καὶ εἰσῆλθον οἱ υἱοὶ Ἰσραὴλ εἰς μέσον τῆς θαλάσσης κατὰ τὸ ξηρόν· καὶ τὸ ὕδωρ αὐτῆς τεῖχος ἐκ δεξιῶν, καὶ τεῖχος ἐξ εὐωνύμων.

\vs{23}Καὶ κάτεδίωξαν οἱ Αἰγύπτιοι, καὶ εἰσῆλθον ὀπίσω αὐτῶν καὶ πᾶς ἵππος Φαραὼ, καὶ τὰ ἅρματα, καὶ οἱ ἀναβάται, εἰς μέσον τῆς θαλάσσης.
\vs{24}Ἐγενήθη δὲ ἐν τῇ φυλακῇ τῇ ἑωθινῇ, καὶ ἐπίβλεψε Κύριος ἐπὶ τὴν παρεμβολὴν τῶν Αἰγυπτίων ἐν στύλῳ πυρὸς καὶ νεφέλης, καὶ συνετάραξε τὴν παρεμβολὴν τῶν Αἰγυπτίων,
\vs{25}καὶ συνέδησε τοὺς ἄξονας τῶν ἁρμάτων αὐτῶν, καὶ ἤγαγεν αὐτοὺς μετὰ βίας· καὶ εἶπαν οἱ Αἰγύπτιοι, φυγωμεν ἀπὸ προσώπου Ἰσραήλ· ὁ γὰρ Κύριος πολεμεῖ περὶ αὐτῶν τοὺς Αἰγυπτίους.
\vs{26}Εἶπε δὲ Κύριος πρὸς Μωυσῆν, ἔκτεινον τὴν χεῖρά σου ἐπὶ τὴν θάλασσαν, καὶ ἀποκαταστήτω τὸ ὕδωρ, καὶ ἐπικαλυψάτω τοὺς Αἰγυπτίους, ἐπί τε τὰ ἅρματα καὶ τοὺς ἀναβάτας.
\vs{27}Ἐξέτεινε δὲ Μωυσῆς τὴν χεῖρα ἐπὶ τὴν θάλασσαν, καὶ ἀπεκατέστη τὸ ὕδωρ πρὸς ἡμέραν ἐπὶ χώρας· οἱ δὲ Αἰγύπτιοι ἔφυγον ὑπὸ τὸ ὕδωρ· καὶ ἐξετίναξε Κύριος τοὺς Αἰγυπτίους μέσον τῆς θαλάσοης.
\vs{28}Καὶ ἐπαναστραφὲν τὸ ὕδωρ ἐκάλυψε τὰ ἅρματα καὶ τοὺς ἀναβάτας, καὶ πᾶσαν τὴν δύναμιν Φαραὼ, τοὺς εἰσπεπορευμένους ὀπίσω αὐτῶν εἰς τὴν θάλασσαν· καὶ οὐ κατελείφθη ἐξ αὐτῶν οὐδὲ εἷς.
\vs{29}Οἱ δὲ υἱοὶ Ἰσραὴλ ἐπορεύθησαν διὰ ξηρᾶς ἐν μέσῳ τῆς θάλασσης· τὸ δὲ ὕδωρ αὐτοῖς τεῖχος ἐκ δεξιῶν, καὶ τεῖχος ἐξ εὐωνύμων.
\vs{30}Καὶ ἐῤῥύσατο Κύριος τὸν Ἰσραὴλ ἐν τῇ ἡμέρᾳ ἐκείνῃ ἐκ χειρὸς τῶν Αἰγυπτίων· καὶ εἶδεν Ἰσραὴλ τοὺς Αἰγυπτίους τεθνηκότας παρὰ τὸ χεῖλος τῆς θαλάσσης.
\vs{31}Εἶδε δὲ Ἰσραὴλ τὴν χεῖρα τὴν μεγάλην, ἃ ἐποίησε Κύριος τοῖς Αἰγυπτίοις· ἐφοβήθη δὲ ὁ λαὸς τὸν Κύριον, καὶ ἐπίστευσαν τῷ Θεῷ, καὶ Μωυσῇ τῷ θεράποντι αὐτοῦ.

\ch{15}
Τότε ᾖσε Μωυσῆς καὶ οἱ υἱοὶ Ἰσραὴλ τὴν ᾠδὴν ταύτην τῷ Θεῷ, καὶ εἶπαν, λέγοντες, ᾄσωμεν τῷ Κυρίῳ, ἐνδόξως γὰρ δεδόξασται· ἵππον καὶ ἀναβάτην ἔῤῥιψεν εἰς θάλασσαν.
\vs{2}Βοηθὸς καὶ σκεπαστὴς ἐγένετό μοι εἰς σωτηρίαν· οὗτός μου Θεὸς, καὶ δοξάσω αὐτόν· Θεὸς τοῦ πατρός μου, καὶ ὑψώσω αὐτόν.
\vs{3}Κύριος συντρίβων πολέμους, Κύριος ὄνομα αὐτῷ.
\vs{4}Ἅρματα Φαραὼ, καὶ τὴν δύναμιν αὐτοῦ, ἔῤῥιψεν εἰς θάλασσαν, ἐπιλέκτους ἀναβάτας τριστάτας· κατεπόθησαν ἐν ἐρυθρᾷ θαλάσσῃ·
\vs{5}Πόντῳ ἐκάλυψεν αὐτούς· κατέδυσαν εἰς βυθὸν ὡσεὶ λίθος.
\vs{6}Ἡ δεξιά σου, Κύριε, δεδόξασται ἐν ἰσχύϊ· ἡ δεξιά σου χεὶρ, Κύριε, ἔθραυσεν ἐχθρούς.
\vs{7}Καὶ τῷ πλήθει τῆς δόξης σου συνέτριψας τοὺς ὑπεναντίους· ἀπέστειλας τὴν ὀργήν σου κατέφαγεν αὐτοὺς ὡς καλάμην.
\vs{8}Καὶ διὰ πνεύματος τοῦ θυμοῦ σου διέστη τὸ ὕδωρ· ἐπάγη ὡσεὶ τεῖχος τὰ ὕδατα· ἐπάγη τὰ κύματα ἐν μέσῳ τῆς θαλάσσης.
\vs{9}Εἶπεν ὁ ἐχθρὸς, διώξας καταλήψομαι, μεριῶ σκῦλα· ἐμπλήσω ψυχήν μου, ἀνελῶ τῇ μαχαίρῃ μου, κυριεύσει ἡ χείρ μου.
\vs{10}Ἀπέστειλας τὸ πνεῦμά σου· ἐκάλυψεν αὐτοὺς θάλασσα· ἔδυσαν ὡσεὶ μόλιβος ἐν ὕδατι σφοδρῷ.
\vs{11}Τίς ὅμοιός σοι ἐν θεοῖς, Κύριε; τίς ὅμοιός σοι; δεδοξασμένος ἐν ἁγίοις, θαυμαστὸς ἐν δόξαις, ποιῶν τέρατα.
\vs{12}Ἐξέτεινας τὴν δεξιάν σου· κατέπιεν αὐτοὺς γῆ.
\vs{13}Ὡδήγησας τῇ δικαιοσύνῃ σου τὸν λαόν σου τοῦτον, ὃν ἐλυτρώσω· παρεκάλεσας τῇ ἰσχύϊ σου εἰς κατάλυμα ἅγιόν σου.
\vs{14}Ἤκουσαν ἔθνη, καὶ ὠργίσθησαν· ὠδῖνες ἔλαβον κατοικοῦντας Φυλιστιείμ.
\vs{15}Τότε ἔσπευσαν ἡγεμόνες Ἐδὼμ, καὶ ἄρχοντες Μωαβιτῶν· ἔλαβεν αὐτοὺς τρόμος· ἐτάκησαν πάντες οἱ κατοικοῦντες Χαναάν.
\vs{16}Ἐπιπέσοι ἐπʼ αὐτοὺς τρόμος καὶ φόβος· μεγέθει βραχίονός σου ἀπολιθωθήτωσαν, ἕως ἂν παρέλθῃ ὁ λαός σου, Κύριε· ἕως ἂν παρέλθῃ ὁ λαός σου οὗτος, ὃν ἐκτήσω.
\vs{17}Εἰσαγαγὼν καταφύτευσον αὐτοὺς εἰς ὄρος κληρονομίας σου, εἰς ἕτοιμον κατοικητήριόν σου, ὃ κατηρτίσω, Κύριε, ἁγίασμα, Κύριε, ὃ ἡτοίμασαν αἱ χεῖρές σου.
\vs{18}Κύριος βασιλεύων τὸν αἰῶνα, καὶ ἐπʼ αἰῶνα, καὶ ἔτι.
\vs{19}Ὅτι εἰσῆλθεν ἵππος Φαραὼ σὺν ἅρμασιν καὶ ἀναβάταις εἰς θάλασσαν, καὶ ἐπήγαγεν ἐπʼ αὐτοὺς Κύριος τὸ ὕδωρ τῆς θαλάσσης· οἱ δὲ υἱοὶ Ἰσραὴλ ἐπορεύθησαν διὰ ξηρᾶς ἐν μέσῳ τῆς θαλάσσης.

\vs{20}Λαβοῦσα δὲ Μαριὰμ ἡ προφῆτις ἡ ἀδελφὴ Ἀαρὼν τὸ τύμπανον ἐν τῇ χειρὶ αὐτῆς, καὶ ἐξήλθοσαν πᾶσαι αἱ γυναῖκες ὀπίσω αὐτῆς μετὰ τυμπάνων καὶ χορῶν.
\vs{21}Ἐξῆρχε δὲ αὐτῶν Μαριὰμ, λέγουσα, ᾄσωμεν τῷ Κυρίῳ, ἐνδόξως γὰρ δεδόξασται· ἵππον καὶ ἀναβάτην ἔῤῥιψεν εἰς θάλασσαν.
\vs{22}Ἐξῆρε δὲ Μωυσῆς τοὺς υἱοὺς Ἰσραὴλ ἀπὸ θαλάσσης ἐρυθρᾶς, καὶ ἤγαγεν αὐτοὺς εἰς τὴν ἔρημον Σούρ· καὶ ἐπορεύοντο τρεῖς ἡμέρας ἐν τῇ ἐρήμῳ, καὶ οὐχ ηὕρισκον ὕδωρ, ὥστε πιεῖν.
\vs{23}Ἦλθον δὲ εἰς Μεῤῥᾶ, καὶ οὐκ ἠδύναντο πιεῖν ἐκ Μεῤῥᾶς· πικρὸν γὰρ ἦν· διὰ τοῦτο ἐπωνόμασε τὸ ὄνομα τοῦ τόπου ἐκείνου, Πικρία.
\vs{24}Καὶ διεγόγγυζεν ὁ λαὸς ἐπὶ Μωυσῇ, λέγοντες, τί πιόμεθα;
\vs{25}Ἐβόησε δὲ Μωυσῆς πρὸς Κύριον· καὶ ἔδειξεν αὐτῷ Κύριος ξύλον, καὶ ἐνέβαλεν αὐτὸ εἰς τὸ ὕδωρ, καὶ ἐγλυκάνθη τὸ ὕδωρ· ἐκεῖ ἔθετο αὐτῷ δικαιώματα καὶ κρίσεις· καὶ ἐκεῖ αὐτὸν ἐπείρασε,
\vs{26}καὶ εἶπεν, ἐὰν ἀκοῇ ἀκούσῃς τῆς φωνῆς Κυρίου τοῦ Θεοῦ σου, καὶ τὰ ἀρεστὰ ἐναντίον αὐτοῦ ποιήσῃς, καὶ ἐνωτίσῃ ταῖς ἐντολαῖς αὐτοῦ, καὶ φυλάξῃς πάντα τὰ δικαιώματα αὐτοῦ, πᾶσαν νόσον, ἣν ἐπήγαγον τοῖς Αἰγυπτίοις, οὐκ ἐπάξω ἐπὶ σέ· ἐγὼ γάρ εἰμι Κύριος ὁ Θεός σου ὁ ἰώμενός σε.
\vs{27}Καὶ ἤλθοσαν εἰς Αἰλείμ· καὶ ἦσαν ἐκεῖ δώδεκα πηγαὶ ὑδάτων, καὶ ἑβδομήκοντα στελέχη φοινίκων· παρενέβαλον δὲ ἐκεῖ παρὰ τὰ ὕδατα.

\ch{16}
Ἀπῄραν δὲ ἐξ Αἰλεὶμ, καὶ ἤλθοσαν πᾶσα συναγωγὴ υἱῶν Ἰσραὴλ εἰς τὴν ἔρημον Σὶν, ὅ ἐστιν ἀνὰ μέσον Αἰλεὶμ, καὶ ἀνὰ μέσον Σινά. τῇ δὲ πεντεκαιδεκάτῃ ἡμέρᾳ, τῷ μηνὶ τῷ δευτέρῳ ἐξεληλυθότων αὐτῶν ἐκ γῆς Αἰγύπτου,
\vs{2}διεγόγγυζε πᾶσα συναγωγὴ υἱῶν Ἰσραὴλ ἐπὶ Μωυσὴν καὶ Ἀαρών.
\vs{3}Καὶ εἶπεν πρὸς αὐτοὺς οἱ υἱοὶ Ἰσραήλ, ὄφελον ἀπεθάνομεν πληγέντες ὑπὸ Κυρίου ἐν γῇ Αἰγύπτῳ, ὅταν ἐκαθίσαμεν ἐπὶ τῶν λεβήτων τῶν κρεῶν, καὶ ἠσθίομεν ἄρτους εἰς πλησμονήν· ὅτι ἐξηγάγετε ἡμᾶς εἰς τὴν ἔρημον ταύτην, ἀποκτεῖναι πᾶσαν τὴν συναγωγὴν ταύτην ἐν λιμῷ.
\vs{4}Εἶπε δὲ Κύριος πρὸς Μωυσῆν, ἰδοὺ ἐγὼ ὕω ὑμῖν ἄρτους ἐκ τοῦ οὐρανοῦ· καὶ ἐξελεύσεται ὁ λαὸς, καὶ συλλέξουσι τὸ τῆς ἡμέρας εἰς ἡμέραν, ὅπως πειράσω αὐτοὺς εἰ πορεύσονται τῷ νόμῳ μου, ἢ οὔ·
\vs{5}Καὶ ἔσται ἐν τῇ ἡμέρᾳ τῇ ἕκτῃ, καὶ ἑτοιμάσουσιν ὃ ἐὰν εἰσενέγκωσι· καὶ ἔσται διπλοῦν ὃ ἐὰν συναγάγωσι τὸ καθʼ ἡμέραν εἰς ἡμέραν.
\vs{6}Καὶ εἶπε Μωυσῆς καὶ Ἀαρὼν πρὸς πάσαν συναγωγὴν υἱῶν Ἰσραὴλ, ἑσπέρας γνώσεσθε, ὅτι Κύριος ἐξήγαγεν ὑμᾶς ἐκ γῆς Αἰγύπτου,
\vs{7}καὶ πρωῒ ὄψεσθε τὴν δόξαν Κυρίου ἐν τῷ εἰσακοῦσαι τὸν γογγυσμὸν ὑμῶν ἐπὶ τῷ Θεῷ· ἡμεῖς δὲ τί ἐσμεν, ὅτι διαγογγύζετε καθʼ ἡμῶν;
\vs{8}Καὶ εἶπε Μωυσῆς, ἐν τῷ διδόναι Κύριον ὑμῖν ἑσπέρας κρέα φαγεῖν, καὶ ἄρτους τὸ πρωῒ εἰς πλησμονὴν, διὰ τὸ εἰσακοῦσαι Κύριον τὸν γογγυσμὸν ὑμῶν, ὃν ὑμεῖς διαγογγύζετε καθʼ ἡμῶν· ἡμεῖς δὲ τί ἐσμεν; οὐ γὰρ καθʼ ἡμῶν ἐστιν ὁ γογγυσμὸς ὑμῶν, ἀλλʼ ἢ κατὰ τοῦ Θεοῦ.

\vs{9}Εἶπε δὲ Μωυσῆς πρὸς Ἀαρὼν, εἶπον πάσῃ συναγωγῇ υἱῶν Ἰσραὴλ, προσέλθετε ἐναντίον τοῦ Θεοῦ· εἰσακήκοε γὰρ τὸν γογγυσμὸν ὑμῶν.
\vs{10}Ἡνίκα δὲ ἐλάλει Ἀαρὼν πάσῃ συναγωγῇ υἱῶν Ἰσραὴλ, καὶ ἐπεστράφησαν εἰς τὴν ἔρημον, καὶ ἡ δόξα Κυρίου ὤφθη ἐν νεφέλῃ.
\vs{11}καὶ ἐλάλησε Κύριος πρὸς Μωυσῆν, λέγων,
\vs{12}εἰσακήκοα τὸν γογγυσμὸν τῶν υἱῶν Ἰσραήλ· λάλησον πρὸς αὐτοὺς, λέγων, τὸ πρὸς ἑσπέραν ἔδεσθε κρέα, καὶ τὸ πρωῒ πλησθήσεσθε ἄρτων· καὶ γνώσεσθε, ὅτι ἐγὼ Κύριος ὁ Θεὸς ὑμῶν.
\vs{13}Ἐγένετο δὲ ἑσπέρα· καὶ ἀνέβη ὀρτυγομήτρα, καὶ ἐκάλυψε τὴν παρεμβολήν· τὸ πρωῒ ἐγένετο καταπαυομένης τῆς δρόσου κύκλῳ τῆς παρεμβολῆς.
\vs{14}Καὶ ἰδοὺ ἐπὶ πρόσωπον τῆς ἐρήμου λεπτὸν ὡσεὶ κόριον λευκὸν, ὡσεὶ πάγος ἐπὶ τῆς γῆς.
\vs{15}Ἰδόντες δὲ αὐτὸ οἱ υἱοὶ Ἰσραὴλ, εἶπαν ἕτερος τῷ ἑτέρῳ, τί ἐστι τοῦτο; οὐ γὰρ ᾔδεισαν τί ἦν· εἶπε δὲ Μωυσῆς αὐτοῖς, οὗτος ὁ ἄρτος, ὃν ἔδωκε Κύριος ὑμῖν φαγεῖν.
\vs{16}Τοῦτο τὸ ῥῆμα ὃ συνέταξε Κύριος· συναγάγετε ἀπʼ αὐτοῦ ἕκαστος εἰς τοὺς καθήκοντας γομὸρ, κατὰ κεφαλὴν κατὰ ἀριθμὸν ψυχῶν ὑμῶν, ἕκαστος σὺν τοῖς συσκηνίοις ὑμῶν συλλέξατε.
\vs{17}Ἐποίησαν δὲ οὕτως οἱ υἱοὶ Ἰσραήλ· καὶ συνέλεξαν ὁ τὸ πολὺ καὶ ὁ τὸ ἔλαττον.
\vs{18}Καὶ μετρήσαντες γομὸρ, οὐκ ἐπλεόνασεν ὁ τὸ πόλυ, καὶ ὁ τὸ ἔλαττον οὐκ ἠλαττόνησεν· ἕκαστος εἰς τοὺς καθήκοντας παρʼ ἑαυτῷ συνέλεξαν.
\vs{19}Εἶπε δὲ Μωυσῆς πρὸς αὐτοὺς, μηδεὶς καταλειπέτω ἀπʼ αὐτοῦ εἰς τὸ πρωΐ.

\vs{20}Καὶ οὐκ εἰσήκουσαν Μωυσῆ, ἀλλὰ κατέλιπόν τινες ἀπʼ αὐτοῦ εἰς τὸ πρωΐ· καὶ ἐξέζεσε σκώληκας, καὶ ἐπώζεσε· καὶ ἐπικράνθη ἐπʼ αὐτοῖς Μωυσῆς.
\vs{21}Καὶ συνέλεξαν αὐτὸ πρωῒ πρωῒ, ἕκαστος τὸ καθῆκον αὐτῷ· ἡνίκα δὲ διεθέρμαινεν ὁ ἥλιος, ἐτήκετο.
\vs{22}Ἐγένετο δὲ τῇ ἡμέρᾳ τῇ ἕκτῃ, συνέλεξαν τὰ δέοντα διπλᾶ, δύο γομὸρ τῷ ἑνί· εἰσήλθοσαν δὲ πάντες οἱ ἄρχοντες τῆς συναγωγῆς, καὶ ἀνήγγειλαν Μωυσῇ.
\vs{23}Εἶπε δὲ Μωυσῆς πρὸς αὐτούς, οὐ τοῦτο τὸ ῥῆμά ἐστιν ὃ ἐλάλησε Κύριος; σάββατα ἀνάπαυσις ἁγία τῷ Κυρίῳ αὔριον· ὅσα ἐὰν πέσσητε, πέσσετε· καὶ ὅσα ἐὰν ἕψητε, ἕψετε· καὶ πᾶν τὸ πλεονάζον καταλείπετε αὐτὸ εἰς ἀποθήκην εἰς τὸ πρωΐ.
\vs{24}Καὶ κατελίποσαν ἀπʼ αὐτοῦ εἰς ἕως πρωῒ, καθὼς συνέταξεν αὐτοῖς Μωυσῆς· καὶ οὐκ ἐπώζεσεν, οὐδὲ σκώληξ ἐγένετο ἐν αὐτῷ.
\vs{25}Εἶπε δὲ Μωυσῆς, φάγετε σήμερον· ἔστι γὰρ σάββατα αήμερον τῷ Κυρίῳ· οὐχ εὑρεθήσεται ἐν τῷ πεδίῳ.
\vs{26}Ἓξ ἡμέρας συλλέξετε· τῇ δὲ ἡμέρᾳ τῇ ἑβδόμῃ σάββατα, ὅτι οὐκ ἔσται ἐν αὐτῇ.
\vs{27}Ἐγένετο δὲ ἐν τῇ ἡμέρᾳ τῇ ἑβδόμῃ ἐξήθλοσάν τινες ἐκ τοῦ λαοῦ συλλέξαι, καὶ οὐχ εὗρον.
\vs{28}Εἶπε δὲ Κύριος πρὸς Μωυσῆν, ἕως τίνος οὐ βούλεσθε εἰσακούειν τὰς ἐντολάς μου, καὶ τὸν νόμον μου;
\vs{29}Ἴδετε, ὁ γὰρ Κύριος ἔδωκεν ὑμῖν σάββατα τὴν ἡμέραν ταύτην· διὰ τοῦτο αὐτὸς ἔδωκεν ὑμῖν τῇ ἡμέρᾳ τῇ ἕκτῃ ἄρτους δύο ἡμερῶν· καθίσεσθε ἕκαστος εἰς τοὺς οἴκους ὑμῶν· μηδεὶς ἐκπορευέσθω ἐκ τοῦ τόπου αὐτοῦ τῇ ἡμέρᾳ τῇ ἑβδόμῃ.
\vs{30}Καὶ ἐσαββάτισεν ὁ λαὸς τῇ ἡμέρᾳ τῇ ἑβδόμῃ.
\vs{31}Καὶ ἐπωνόμασαν αὐτὸ οἱ υἱοὶ Ἰσραὴλ τὸ ὄνομα αὐτοῦ, Μάν· ἦν δὲ ὡσεὶ σπέρμα κορίου λευκόν· τὸ δὲ γεῦμα αὐτοῦ ὡς ἐγκρὶς ἐν μέλιτι.
\vs{32}Εἶπε δὲ Μωυσῆς, τοῦτο τὸ ῥῆμα, ὃ συνέταξε Κύριος, πλήσατε τὸ γομὸρ τοῦ μὰν, εἰς ἀποθήκην εἰς τὰς γενεὰς ὑμῶν· ἵνα ἴδωσι τὸν ἄρτον, ὃν ἐφάγετε ὑμεῖς ἐν τῇ ἐρήμῳ, ὡς ἐξήγαγεν ὑμᾶς Κύριος ἐκ γῆς Αἰγύπτου.
\vs{33}Καὶ εἶπε Μωυσῆς πρὸς Ἀαρὼν, λάβε στάμνον χρυσοῦν ἕνα, καὶ ἔμβαλε εἰς αὐτὸν πλῆρες τὸ γομὸρ τοῦ μὰν, καὶ ἀποθήσεις αὐτὸ ἐναντίον τοῦ Θεοῦ, εἰς διατήρησιν εἰς τὰς γενεὰς ὑμῶν,
\vs{34}ὃν τρόπον συνέταξε Κύριος τῷ Μωυσῇ· καὶ ἀπέθηκεν Ἀαρὼν ἐναντίον τοῦ μαρτυρίου εἰς διατήρησιν.
\vs{35}Οἱ δὲ υἱοὶ Ἰσραὴλ ἔφαγον τὸ μὰν ἔτη τεσσαράκοντα, ἕως ἦλθον εἰς τὴν οἰκουμένην ἐφάγοσαν τὸ μὰν, ἕως παρεγένοντο εἰς μέρος τῆς Φοινίκης.
\vs{36}Τὸ δὲ γομὸρ τὸ δέκατον τῶν τριῶν μέτρων ἦν.

\ch{17}
Καὶ ἀπῇρε πᾶσα συναγωγὴ υἱῶν Ἰσραὴλ ἐκ τῆς ἐρήμου Σὶν κατὰ παρεμβολὰς αὐτῶν, διὰ ῥήματος Κυρίου· καὶ παρενεβάλοσαν ἐν Ῥαφιδείν· οὐκ ἦν δὲ ὕδωρ τῷ λαῷ πιεῖν.
\vs{2}Καὶ ἐλοιδορεῖτο ὁ λαὸς πρὸς Μωυσῆν, λέγοντες, δὸς ἡμῖν ὕδωρ, ἵνα πίωμεν· καὶ εἶπεν αὐτοῖς Μωυσῆς, τί λοιδορεῖσθέ μοι, καὶ τί πειράζετε Κύριον;
\vs{3}Ἐδίψησε δὲ ἐκεῖ ὁ λαὸς ὕδατι· καὶ διεγόγγυσεν ἐκεῖ ὁ λαὸς πρὸς Μωυσῆν, λέγοντες, ἱνατί τοῦτο; ἀνεβίβασας ἡμᾶς ἐξ Αἰγύπτου ἀποκτεῖναι ἡμᾶς καὶ τὰ τέκνα ἡμῶν καὶ τὰ κτήνη τῷ δίψει;
\vs{4}Ἐβόησε δὲ Μωυσῆς πρὸς Κύριον, λέγων, τί ποιήσω τῷ λαῷ τούτῳ; ἔτι μικρὸν, καὶ καταλιθοβολὴσουσί με.
\vs{5}Καὶ εἶπε Κύριος πρὸς Μωυσῆν, προπορεύου τοῦ λαοῦ τούτου· λάβε δὲ σεαυτῷ ἀπὸ τῶν πρεσβυτέρων τοῦ λαοῦ· καὶ τὴν ῥάβδον, ἐν ᾗ ἐπάταξας τὸν ποταμὸν, λάβε ἐν τῇ χειρί σου, καὶ πορεύσῃ.
\vs{6}Ὅδε ἐγὼ ἕστηκα ἐκεῖ πρὸ τοῦ σὲ ἐπὶ τῆς πέτρας ἐν Χωρήβ· καὶ πατάξεις τὴν πέτραν, καὶ ἐξελεύσεται ἐξ αὐτῆς ὕδωρ, καὶ πίεται ὁ λαός. Ἐποίησε δὲ Μωυσῆς οὕτως ἐναντίον τῶν υἱῶν Ἰσραήλ.
\vs{7}Καὶ ἐπωνόμασε τὸ ὄνομα τοῦ τόπου ἐκείνου, Πειρασμὸς, καὶ Λοιδόρησις, διὰ τὴν λοιδορίαν τῶν υἱῶν Ἰσραὴλ, καὶ διὰ τὸ πειράζειν Κύριον, λέγοντας, εἰ ἔστι Κύριος ἐν ἡμῖν, ἢ οὔ;

\vs{8}Ἦλθε δὲ Ἀμαλὴκ καὶ ἐπολέμει Ἰσραὴλ ἐν Ῥαφιδείν.
\vs{9}Εἶπε δὲ Μωυσῆς τῷ Ἰησοῖ, Ἐπίλεξον σεαυτῷ ἄνδρας δυνατοὺς, καὶ ἐξελθὼν παράταξαι τῷ Ἀμαλὴκ αὔριον· καὶ ἰδοὺ ἐγὼ ἕστηκα ἐπὶ τῆς κορυφῆς τοῦ βουνοῦ, καὶ ἡ ῥάβδος τοῦ Θεοῦ ἐν τῇ χειρί μου.
\vs{10}Καὶ ἐποίησεν Ἰησοῦς καθάπερ εἶπεν αὐτῷ Μωυσῆς, καὶ ἐξελθὼν παρετάξατο τῷ Ἀμαλήκ· καὶ Μωυσῆς καὶ Ἀαρὼν καὶ Ὢρ ἀνέβησαν ἐπὶ τὴν κορυφὴν τοῦ βουνοῦ.
\vs{11}Καὶ ἐγένετο ὅταν ἐπῇρε Μωυσῆς τὰς χεῖρας, κατίσχυεν Ἰσραήλ· ὅταν δὲ καθῆκε τὰς χεῖρας, κατίσχυεν Ἀμαλήκ.
\vs{12}Αἱ δὲ χεῖρες Μωυσῆ βαρεῖαι· καὶ λαβόντες λίθον ὑπέθηκαν ὑπʼ αὐτὸν, καὶ ἐκάθητο ἐπʼ αὐτοῦ· καὶ Ἀαρὼν καὶ Ὢρ ἐστήριζον τὰς χεῖρας αὐτοῦ ἐντεῦθεν εἷς, καὶ ἐντεῦθεν εἷς· καὶ ἐγένοντο αἱ χεῖρες Μωυσῆ ἐστηριγμέναι ἕως δυσμῶν ἡλίου.
\vs{13}Καὶ ἐτρέψατο Ἰησοῦς τὸν Ἀμαλὴκ, καὶ πάντα τὸν λαὸν αὐτοῦ ἐν φόνῳ μαχαίρας.
\vs{14}Εἶπε δὲ Κύριος πρὸς Μωυσῆν, Κατάγραψον τοῦτο εἰς μνημόσυνον εἰς βιβλίον, καὶ δὸς εἰς τὰ ὦτα Ἰησοῖ· ὅτι ἀλοιφῇ ἐξαλείψω τὸ μνημόσυνον Ἀμαλὴκ ἐκ τῆς ὑπὸ τὸν οὐρανόν.
\vs{15}Καὶ ᾠκοδόμησε Μωυσῆς θυσιαστήριον Κυρίῳ· καὶ ἐπωνόμασε τὸ ὄνομα αὐτοῦ, Κύριος καταφυγή μου.
\vs{16}Ὅτι ἐν χειρὶ κρυφαίᾳ πολεμεῖ Κύριος ἐπὶ Ἀμαλὴκ ἀπὸ γενεῶν εἰς γενεάς.

\ch{18}
Ἤκουσε δὲ Ἰοθὸρ ἱερεὺς Μαδιὰμ ὁ γαμβρὸς Μωυσῆ πάντα ὅσα ἐποίησε Κύριος Ἰσραὴλ τῷ ἑαυτοῦ λαῷ· ἐξήγαγε γὰρ Κύριος τὸν Ἰσραὴλ ἐξ Αἰγύπτου.
\vs{2}Ἔλαβε δὲ Ἰοθὸρ ὁ γαμβρὸς Μωυσῆ Σεπφώραν τὴν γυναῖκα Μωυσῆ μετὰ τὴν ἄφεσιν αὐτῆς,
\vs{3}καὶ τοὺς δύο υἱοὺς αὐτῆς· ὄνομα τῷ ἑνὶ αὐτῶν Γηρσάμ, λέγων, πάροικος ἤμην ἐν γῇ ἀλλοτρίᾳ·
\vs{4}καὶ τὸ ὄνομα τοῦ δευτέρου Ἐλίεζερ, λέγων, ὁ γὰρ Θεὸς τοῦ πατρός μου βοηθός μου, καὶ ἐξείλατό με ἐκ χειρὸς Φαραώ.
\vs{5}Καὶ ἐξῆλθεν Ἰοθὸρ ὁ γαμβρὸς Μωυσῆ καὶ οἱ υἱοὶ καὶ ἡ γυνὴ πρὸς Μωυσῆν εἰς τὴν ἔρημον, οὗ παρενέβαλεν ἐπʼ ὄρους τοῦ Θεοῦ.
\vs{6}Ἀνηγγέλη δὲ Μωυσῇ, λέγοντες, ἰδοὺ ὁ γαμβρός σου Ἰοθὸρ παραγίνεται πρὸς σέ, καὶ ἡ γυνὴ καὶ οἱ δύο υἱοί σου μετʼ αὐτοῦ.
\vs{7}Ἐξῆλθε δὲ Μωυσῆς εἰς συνάντησιν τῷ γαμβρῷ, καὶ προσεκύνησεν αὐτῷ, καὶ ἐφίλησεν αυτὸν, καὶ ἠσπάσαντο ἀλλήλους, καὶ εἰσήγαγεν αὐτοὺς εἰς τὴν σκηνήν.
\vs{8}Καὶ διηγήσατο Μωυσῆς τῷ γαμβρῷ πάντα ὅσα ἐποίησε Κύριος τῷ Φαραῷ καὶ πᾶσι τοῖς Αἰγυπτίοις ἕνεκεν τοῦ Ἰσραήλ, καὶ πάντα τὸν μόχθον τὸν γενόμενον αὐτοῖς ἐν τῇ ὁδῷ, καὶ ὅτι ἐξείλατο αὐτοὺς Κύριος ἐκ χειρὸς Φαραὼ, καὶ ἐκ χειρὸς τῶν Αἰγυπτίων.
\vs{9}Ἐξέστη δὲ Ἰοθὸρ ἐπὶ πᾶσι τοῖς ἀγαθοῖς οἷς ἐποίησεν αὐτοῖς Κύριος, ὅτι ἐξείλατο αὐτοὺς ἐκ χειρὸς Αἰγυπτίων καὶ ἐκ χειρὸς Φαραώ.
\vs{10}Καὶ εἶπεν Ἰοθὸρ, εὐλογητὸς Κύριος, ὅτι ἐξείλατο αὐτοὺς ἐκ χειρὸς Αἰγυπτίων καὶ ἐκ χειρὸς Φαραώ.
\vs{11}Νῦν ἔγνων ὅτι μέγας Κύριος παρὰ πάντας τοὺς θεούς ἕνεκεν τούτου, ὅτι ἐπέθεντο αὐτοῖς.
\vs{12}Καὶ ἔλαβεν Ἰοθὸρ ὁ γαμβρὸς Μωυσῆ ὁλοκαυτώματα καὶ θυσίας τῷ Θεῷ· παρεγένετο δὲ Ἀαρὼν καὶ πάντες οἱ πρεσβύτεροι Ἰσραὴλ συμφαγεῖν ἄρτον μετὰ τοῦ γαμβροῦ Μωυσῆ, ἐναντίον τοῦ Θεοῦ.

\vs{13}Καὶ ἐγένετο μετὰ τὴν ἐπαύριον συνεκάθισε Μωυσῆς κρίνειν τὸν λαόν· παρειστήκει δὲ πᾶς ὁ λαὸς Μωυσῇ ἀπὸ πρωΐθεν ἕως δείλης.
\vs{14}Καὶ ἰδὼν Ἰοθὸρ πάντα ὅσα ποιεῖ τῷ λαῷ, λέγει, τί τοῦτο ὃ σὺ ποιεῖς τῷ λαῷ; διατί σὺ κάθησαι μόνος, πᾶς δὲ ὁ λαὸς παρέστηκέ σοι ἀπὸ πρωΐθεν ἕως δείλης;
\vs{15}Καὶ λέγει Μωυσῆς τῷ γαμβρῶ, Ὅτι παραγίνεται πρός με ὁ λαὸς ἐκζητῆσαι κρίσιν παρὰ τοῦ Θεοῦ.
\vs{16}Ὅταν γὰρ γένηται αὐτοῖς ἀντιλογία, καὶ ἔλθωσι πρός με, διακρίνω ἕκαστον, καὶ συμβιβάζω αὐτοὺς τὰ προστάγματα τοῦ Θεοῦ καὶ τὸν νόμον αὐτοῦ.
\vs{17}Εἶπε δὲ ὁ γαμβρὸς Μωυσῆ πρὸς αὐτὸν, οὐκ ὀρθῶς σὺ ποιεῖς τὸ ῥῆμα τοῦτο.
\vs{18}Φθορᾷ καταφθαρήσῃ ἀνυπομονήτῳ καὶ σὺ, καὶ πᾶς ὁ λαὸς οὗτος, ὅς ἐστι μετὰ σοῦ· βαρύ σοι τὸ ῥῆμα τοῦτο· οὐ δυνήσῃ ποιεῖν σὺ μόνος.
\vs{19}Νῦν οὖν ἄκουσόν μου, καὶ συμβουλεύσω σοι, καὶ ἔσται ὁ Θεὸς μετὰ σοῦ· γίνου σὺ τῷ λαῷ τὰ πρὸς τὸν Θεὸν, καὶ ἀνοίσεις τοὺς λόγους αὐτῶν πρὸς τὸν Θεόν.
\vs{20}Καὶ διαμαρτύρῇ αὐτοῖς τὰ προστάγματα τοῦ Θεοῦ καὶ τὸν νόμον αὐτοῦ, καὶ σημανεῖς αὐτοῖς τὰς ὁδοὺς ἐν αἷς πορεύσονται ἐν αὐταῖς, καὶ τὰ ἔργα ἃ ποιήσουσι.
\vs{21}Καὶ σὺ σεαυτῷ σκέψαι ἀπὸ παντὸς τοῦ λαοῦ ἄνδρας δυνατοὺς, θεοσεβεῖς, ἄνδρας δικαίους, μισοῦντας ὑπερηφανίαν, καὶ καταστήσεις ἐπʼ αὐτῶν χιλιάρχους καὶ ἑκατοντάρχους καὶ πεντηκοντάρχους καὶ δεκαδάρχους.
\vs{22}Καὶ κρινοῦσι τὸν λαὸν πᾶσαν ὥραν· τὸ δὲ ῥῆμα τὸ ὑπέρογκον ἀνοίσουσιν ἐπὶ σὲ· τὰ δὲ βραχέα τῶν κριμάτων κρινοῦσιν αὐτοί· καὶ κουφιοῦσιν ἀπὸ σοῦ, καὶ συναντιλήψονταί σοι.
\vs{23}Ἐὰν τὸ ῥῆμα τοῦτο ποιήσῃς, κατισχύσει σε ὁ Θεὸς, καὶ δυνήσῃ παραστῆναι, καὶ πᾶς ὁ λαὸς οὗτος εἰς τὸν ἑαυτοῦ τόπον μετʼ εἰρήνης ἥξει.
\vs{24}Ἤκουσε δὲ Μωυσῆς τῆς φωνῆς τοῦ γαμβροῦ, καὶ ἐποίησεν ὅσα εἶπεν αὐτῷ.
\vs{25}Καὶ ἐπέλεξε Μωυσῆς ἄνδρας δυνατοὺς ἀπὸ παντὸς Ἰσραὴλ, καὶ ἐποίησεν αὐτοὺς ἐπʼ αὐτῶν χιλιάρχους καὶ ἑκατοντάρχους καὶ πεντηκοντάρχους καὶ δεκαδάρχους.
\vs{26}Καὶ ἐκρίνοσαν τὸν λαὸν πᾶσαν ὥραν· πᾶν δὲ ῥῆμα ὑπέρογκον ἀνεφέροσαν ἐπὶ Μωυσῆν· πᾶν δὲ ῥῆμα ἐλαφρὸν ἐκρίνοσαν αὐτοί.
\vs{27}Ἐξαπέστειλε δὲ Μωυσῆς τὸν ἑαυτοῦ γαμβρὸν, καὶ ἀπῆλθεν εἰς τὴν γῆν αὐτοῦ.

\ch{19}
Τοῦ δὲ μηνὸς τοῦ τρίτου τῆς ἐξόδου τῶν υἱῶν Ἰσραὴλ ἐκ γῆς Αἰγύπτου τῇ ἡμέρᾳ ταύτῃ, ἤλθοσαν εἰς τὴν ἔρημον τοῦ Σινά.
\vs{2}Καὶ ἀπῆραν ἐκ Ῥαφιδεὶν, καὶ ἤλθοσαν εἰς τὴν ἔρημον τοῦ Σινὰ, καὶ παρενέβαλεν ἐκεῖ Ἰσραὴλ κατέναντι τοῦ ὄρους.
\vs{3}Καὶ Μωυσῆς ἀνέβη εἰς τὸ ὄρος τοῦ Θεοῦ· καὶ ἐκάλεσεν αὐτὸν ὁ Θεὸς ἐκ τοῦ ὄρους, λέγων, τάδε ἐρεῖς τῷ οἴκῳ Ἰακὼβ, καὶ ἀναγγελεῖς τοῖς υἱοῖς Ἰσραήλ.
\vs{4}Αὐτοὶ ἑωράκατε ὅσα πεποίηκα τοῖς Αἱγυπτίοις, καὶ ἀνέλαβον ὑμᾶς ὡσεὶ ἐπὶ πτερύγων ἀετῶν, καὶ προσηγαγόμην ὑμᾶς πρὸς ἐμαυτόν.
\vs{5}Καὶ νῦν ἐὰν ἀκοῇ ἀκούσητε τῆς ἐμῆς φωνῆς, καὶ φυλάξητε τὴν διαθήκην μου, ἔσεσθέ μοι λαὸς περιούσιος ἀπὸ πάντων τῶν ἐθνῶν· ἐμὴ γάρ ἐστι πᾶσα ἡ γῆ.
\vs{6}Ὑμεῖς δὲ ἔσεσθέ μοι βασίλειον ἱεράτευμα καὶ ἔθνος ἅγιον· ταῦτα τὰ ῥήματα ἐρεῖς τοῖς υἱοῖς Ἰσραήλ.
\vs{7}Ἦλθε δὲ Μωυσῆς, καὶ ἐκάλεσεν τοὺς πρεσβυτέρους τοῦ λαοῦ· καὶ παρέθηκεν αὐτοῖς πάντας τοὺς λόγους τούτους, οὓς συνέταξεν αὐτοῖς ὁ Θεός.
\vs{8}Ἀπεκρίθη δὲ πᾶς ὁ λαὸς ὁμοθυμαδὸν, καὶ εἶπαν, πάντα ὅσα εἶπεν ὁ Θεὸς, ποιήσομεν καὶ ἀκουσόμεθα· ἀνήνεγκε δὲ Μωυσῆς τοὺς λόγους τούτους πρὸς τὸν Θεόν.
\vs{9}Εἶπε δὲ Κύριος πρὸς Μωυσῆν, ἰδοὺ ἐγὼ παραγίνομαι πρὸς σὲ ἐν στύλῳ νεφέλης, ἵνα ἀκούσῃ ὁ λαὸς λαλοῦντός μου πρὸς σὲ, καὶ σοὶ πιστεύσωσιν εἰς τὸν αἰῶνα· ἀνήγγειλε δὲ Μωυσῆς τὰ ῥήματα τοῦ λαοῦ πρὸς Κύριον.
\vs{10}Εἶπε δὲ Κύριος πρὸς Μωυσῆν, Καταβὰς διαμάρτυραι τῷ λαῷ, καὶ ἅγνισον αὐτοὺς σήμερον καὶ αὔριον, καὶ πλυνάτωσαν τὰ ἱμάτια,
\vs{11}καὶ ἔστωσαν ἕτοιμοι εἰς τὴν ἡμέραν τὴν τρίτην· τῇ γὰρ ἡμέρᾳ τῇ τρίτῃ καταβήσεται Κύριος ἐπὶ τὸ ὄρος τὸ Σινὰ, ἐναντίον παντὸς τοῦ λαοῦ.
\vs{12}Καὶ ἀφοριεῖς τὸν λαὸν κύκλῳ, λέγων, προσέχετε ἑαυτοῖς τοῦ ἀναβῆναι εἰς τὸ ὄρος, καὶ θίγειν τι αὐτοῦ· πᾶς ὁ ἁψάμενος τοῦ ὄρους, θανάτῳ τελευτήσει.
\vs{13}Οὐχ ἅψεται αὐτοῦ χείρ· ἐν γὰρ λίθοις λιθοβοληθήσεται, ἢ βολίδι κατατοξευθήσεται· ἐάν τε κτῆνος ἐάν τε ἄνθρωπος, οὐ ζήσεται· ὅταν αἱ φωναὶ καὶ αἱ σάλπιγγες καὶ ἡ νεφέλη ἀπέλθῃ ἀπὸ τοῦ ὄρους, ἐκεῖνοι ἀναβήσονται ἐπὶ τὸ ὄρος.

\vs{14}Κατέβη δὲ Μωυσῆς ἐκ τοῦ ὄρους πρὸς τὸν λαὸν, καὶ ἡγίασεν αὐτούς· καὶ ἔπλυναν τὰ ἱμάτια.
\vs{15}Καὶ εἶπε τῷ λαῷ, γίνεσθε ἕτοιμοι, τρεῖς ἡμέρας μὴ προσέλθητε γυναικί.
\vs{16}Ἐγένετο δὲ τῇ ἡμέρᾳ τῇ τρίτῃ γενηθέντος πρὸς ὄρθρον, καὶ ἐγένοντο φωναὶ καὶ ἀστραπαὶ καὶ νεφέλη γνοφώδης ἐπʼ ὄρους Σινά· φωνὴ τῆς σάλπιγγος ἤχει μέγα· καὶ ἐπτοήθη πᾶς ὁ λαὸς ὁ ἐν τῇ παρεμβολῇ.
\vs{17}Καὶ ἐξήγαγε Μωυσῆς τὸν λαὸν εἰς συνάντησιν τοῦ Θεοῦ ἐκ τῆς παρεμβολῆς· καὶ παρέστησαν ὑπὸ τὸ ὄρος.
\vs{18}Τὸ ὄρος τὸ Σινὰ ἐκαπνίζετο ὅλον, διὰ τὸ καταβεβηκέναι ἐπʼ αὐτὸ τὸν Θεὸν ἐν πυρί· καὶ ἀνέβαινεν ὁ καπνὸς, ὡσεὶ καπνὸς καμίνου· καὶ ἐξέστη πᾶς ὁ λαὸς σφόδρα.
\vs{19}Ἐγίνοντο δὲ αἱ φωναὶ τῆς σάλπιγγος προβαίνουσαι ἰσχυρότεραι σφόδρα. Μωυσῆς ἐλάλησεν, ὁ δὲ Θεὸς ἀπεκρίνατο αὐτῷ φωνῇ.
\vs{20}Κατέβη δὲ Κύριος ἐπὶ τὸ ὄρος τὸ Σινὰ ἐπὶ τὴν κορυφὴν τοῦ ὄρους· καὶ ἐκάλεσε Κύριος Μωυσῆν ἐπὶ τὴν κορυφὴν τοῦ ὄρους· καὶ ἀνέβη Μωυσῆς.
\vs{21}Καὶ εἶπεν ὁ Θεὸς πρὸς Μωυσῆν, λέγων, καταβὰς διαμάρτυραι τῷ λαῷ, μή ποτε ἐγγίσωσι πρὸς τὸν Θεὸν κατανοῆσαι, καὶ πέσωσιν ἐξ αὐτῶν πλῆθος·
\vs{22}Καὶ οἱ ἱερεῖς οἱ ἐγγίζοντες Κυρίῳ τῷ Θεῷ ἁγιασθήτωσαν, μήποτε ἀπαλλάξῃ ἀπʼ αὐτῶν Κύριος.

\vs{23}Καὶ εἶπε Μωυσῆς πρὸς τὸν Θεὸν, οὐ δυνήσεται ὁ λαὸς προσαναβῆναι πρὸς τὸ ὄρος τὸ Σινά· σὺ γὰρ διαμεμαρτύρησαι ἡμῖν, λέγων, ἀφόρισαι τὸ ὄρος, καὶ ἁγίασαι αὐτό.
\vs{24}Εἴπε δὲ αὐτῷ Κύριος, βάδιζε, κατάβηθι, καὶ ἀνάβηθι σὺ καὶ Ἀαρὼν μετὰ σοῦ· οἱ δὲ ἱερεῖς καὶ ὁ λαὸς μὴ βιαζέσθωσαν ἀναβῆναι πρὸς τὸν Θεὸν, μὴ ποτε ἀπολέσῃ ἀπʼ αὐτῶν Κύριος.
\vs{25}Κατέβη δὲ Μωυσῆς πρὸς τὸν λαὸν, καὶ εἶπεν αὐτοῖς.

\ch{20}
Καὶ ἐλάλησε Κύριος πάντας τοὺς λόγους τούτους, λέγων,
\vs{2}ἐγώ εἰμι Κύριος ὁ Θεός σου, ὅστις ἐξήγαγόν σε ἐκ γῆς Αἰγύπτου, ἐξ οἴκου δουλείας.
\vs{3}Οὐκ ἔσονταί σοι θεοὶ ἕτεροι πλὴν ἐμοῦ.
\vs{4}Οὐ ποιήσεις σεαυτῷ εἴδωλον, οὐδὲ παντὸς ὁμοίωμα, ὅσα ἐν τῷ οὐρανῷ ἄνω, καὶ ὅσα ἐν τῇ γῇ κάτω, καὶ ὅσα ἐν τοῖς ὕδασιν ὑποκάτω τῆς γῆς.
\vs{5}Οὐ προσκυνήσεις αὐτοῖς, οὐδὲ μὴ λατρεύσεις αὐτοῖς· ἐγὼ γάρ εἰμι Κύριος ὁ Θεός σου, Θεὸς ζηλωτὴς, ἀποδιδοὺς ἁμαρτίας πατέρων ἐπὶ τέκνα, ἕως τρίτης καὶ τετάρτης γενεᾶς τοῖς μισοῦσί με,
\vs{6}καὶ ποιῶν ἔλεος εἰς χιλιάδας τοῖς ἀγαπῶσί με, καὶ τοῖς φυλάσσουσι τὰ προστάγματά μου.
\vs{7}Οὐ λήψῃ τὸ ὄνομα Κυρίου τοῦ Θεοῦ σου ἐπὶ ματαίῳ· οὐ γὰρ μὴ καθαρίσῃ Κύριος ὁ Θεός σου τὸν λαμβάνοντα τὸ ὄνομα αὐτοῦ ἐπὶ ματαίῳ.
\vs{8}Μνήσθητι τὴν ἡμέραν τῶν σαββάτων ἁγιάζειν αὐτήν.
\vs{9}Ἓξ ἡμέρας ἐργᾷ, καὶ ποιήσεις πάντα τὰ ἔργα σου.
\vs{10}Τῇ δὲ ἡμέρᾳ τῇ ἑβδόμῃ, σάββατα Κυρίῳ τῷ Θεῷ σου· οὐ ποιήσεις ἐν αὐτῇ πᾶν ἔργον σὺ, καὶ ὁ υἱός σου, καὶ ἡ θυγάτηρ σου, ὁ παῖς σου, καὶ ἡ παιδίσκη σου, ὁ βοῦς σου, καὶ τὸ ὑποζύγιόν σου, καὶ πᾶν κτῆνός σου, καὶ ὁ προσήλυτος ὁ παροικῶν ἐν σοί.
\vs{11}Ἐν γὰρ ἓξ ἡμέραις ἐποίησε Κύριος τὸν οὐρανὸν καὶ τὴν γῆν καὶ τὴν θὰλασσαν καὶ πάντα τὰ ἐν αὐτοῖς, καὶ κατέπαυσε τῇ ἡμέρᾳ τῇ ἑβδόμῃ· διὰ τοῦτο εὐλόγησε Κύριος τὴν ἡμέραν τὴν ἑβδόμην, καὶ ἡγίασεν αὐτήν.
\vs{12}Τίμα τὸν πατέρα σου, καὶ τὴν μητέρα σου, ἵνα εὖ σοι γένηται, καὶ ἵνα μακροχρόνιος γένῃ ἐπὶ τῆς γῆς τῆς ἀγαθῆς, ἧς Κύριος ὁ Θεός σου δίδωσί σοι.
\vs{13}Οὐ μοιχεύσεις.
\vs{14}Οὐ κλέψεις.
\vs{15}Οὐ φονεύσεις.
\vs{16}Οὐ ψευδομαρτυρήσεις κατὰ τοῦ πλησίον σου μαρτυρίαν ψευδῆ.
\vs{17}Οὐκ ἐπιθυμήσεις τὴν γυναῖκα τοῦ πλησίον σου· οὐκ ἐπιθυμήσεις τὴν οἰκίαν τοῦ πλησίον σου, οὔτε τὸν ἀγρὸν αὐτοῦ, οὔτε τὸν παῖδα αὐτοῦ, οὔτε τὴν παιδίσκην αὐτοῦ, οὔτε τοῦ βοὸς αὐτοῦ, οὔτε τοῦ ὑποζυγίου αὐτοῦ, οὔτε παντὸς κτήνους αὐτοῦ, οὔτε ὅσα τῷ πλησίον σου ἐστί.

\vs{18}Καὶ πᾶς ὁ λαὸς ἑώρα τὴν φωνὴν, καὶ τὰς λαμπάδας, καὶ τὴν φωνὴν τῆς σάλπιγγος, καὶ τὸ ὄρος τὸ καπνίζον· φοβηθέντες δὲ πᾶς ὁ λαὸς ἔστησαν μακρόθεν.
\vs{19}Καὶ εἶπαν πρὸς Μωυσῆν, λάλησον σὺ ἡμῖν, καὶ μὴ λαλείτω πρὸς ἡμᾶς ὁ Θεὸς, μὴ ἀποθάνωμεν.
\vs{20}Καὶ λέγει αὐτοῖς Μωυσῆς, θαρσεῖτε· ἕνεκεν γὰρ τοῦ πειράσαι ὑμᾶς παρεγενήθη ὁ Θεὸς πρὸς ὑμᾶς, ὅπως ἂν γένηται ὁ φόβος αὐτοῦ ἐν ὑμῖν, ἵνα μὴ ἁμαρτάνητε.
\vs{21}Εἱστήκει δὲ ὁ λαὸς μακρόθεν, Μωυσῆς δὲ εἰσῆλθεν εἰς τὸν γνόφον, οὗ ἦν ὁ Θεός.
\vs{22}Εἶπε δὲ Κύριος πρὸς Μωυσῆν, τάδε ἐρεῖς τῷ οἴκῳ Ἰακὼβ, καὶ ἀναγγελεῖς τοῖς υἱοῖς Ἰσραήλ· ὑμεῖς ἑωράκατε, ὅτι ἐκ τοῦ οὐρανοῦ λελάληκα πρὸς ὑμᾶς.
\vs{23}Οὐ ποιήσετε ὑμῖν αὐτοῖς θεοὺς ἀργυροῦς, καὶ θεοὺς χρυσοῦς οὐ ποιήσετε ὑμῖν αὑτοῖς.
\vs{24}Θυσιαστήριον ἐκ γῆς ποιήσετέ μοι, καὶ θύσετε ἐπʼ αὐτοῦ τὰ ὁλοκαυτώματα ὑμῶν, καὶ τὰ σωτήρια ὑμῶν, καὶ τὰ πρόβατα, καὶ τοὺς μόσχους ὑμῶν ἐν παντὶ τόπῳ, οὗ ἐὰν ἐπονομάσω τὸ ὄνομά μου ἐκεῖ, καὶ ἥξω πρὸς σὲ, καὶ εὐλογήσω σε.
\vs{25}Ἐὰν δὲ θυσιαστήριον ἐκ λίθων ποιῇς μοι, οὐκ οἰκοδομήσεις αὐτοὺς τμητούς· τὸ γὰρ ἐγχειρίδιόν σου ἐπιβέβληκας ἐπʼ αὐτοὺς, καὶ μεμίανται.
\vs{26}Οὐκ ἀναβήσῃ ἐν ἀναβαθμίσιν ἐπὶ τὸ θυσιαστήριόν μου, ὅπως ἂν μὴ ἀποκαλύψῃς τὴν ἀσχημοσύνην σου ἐπʼ αὐτοῦ.

\ch{21}
Καὶ ταῦτα τὰ δικαιώματα, ἃ παραθήσῃ ἐνώπιον αὐτῶν.
\vs{2}Ἐὰν κτήσῃ παῖδα Ἐβραῖον, ἓξ ἔτη δουλεύσει σοι· τῷ δὲ ἑβδόμῳ ἔτει ἀπελεύσεται ἐλεύθερος δωρεάν.
\vs{3}Ἐὰν αὐτὸς μόνος εἰσέλθῃ, καὶ μόνος ἐξελεύσεται· ἐὰν δὲ γυνὴ συνεισέλθῃ μετʼ αὐτοῦ, ἐξελεύσεται καὶ ἡ γυνὴ αὐτοῦ.
\vs{4}Καὶ ἐὰν δὲ ὁ κύριος δῷ αὐτῷ γυναῖκα, καὶ τέκῃ αὐτῷ υἱοὺς ἢ θυγατέρας, ἡ γυνὴ καὶ τὰ παιδία ἔσται τῷ κυρίῳ αὐτοῦ, αὐτὸς δὲ μόνος ἐξελεύσεται.
\vs{5}Ἐὰν δὲ ἀποκριθεὶς εἴπῃ ὁ παῖς, ἠγάπηκα τὸν κύριόν μου, καὶ τὴν γυναῖκα, καὶ τὰ παιδία, οὐκ ἀποτρέχω ἐλεύθερος·
\vs{6}προσάξει αὐτὸν ὁ κύριος αὐτοῦ πρὸς τὸ κριτήριον τοῦ Θεοῦ, καὶ τότε προσάξει αὐτὸν ἐπὶ τὴν θύραν ἐπὶ τὸν σταθμὸν, καὶ τρυπήσει ὁ κύριος αὐτοῦ τὸ οὖς τῷ ὀπητίῳ, καὶ δουλεύσει αὐτῷ εἰς τὸν αἰῶνα.

\vs{7}Ἐὰν δέ τις ἀποδῶται τὴν ἑαυτοῦ θυγατέρα οἰκέτιν, οὐκ ἀπελεύσεται, ὥσπερ ἀποτρέχουσιν αἱ δοῦλαι.
\vs{8}Ἐὰν μὴ εὐαρεστήσῃ τῷ κυρίῳ αὐτῆς, ἣ αὐτῷ καθωμολογήσατο, ἀπολυτρώσει αὐτήν· ἔθνει δὲ ἀλλοτρίῳ οὐ κύριός ἐστι πωλεῖν αὐτὴν, ὅτι ἠθέτησεν ἐν αὐτῇ.
\vs{9}Ἐὰν δὲ τῷ υἱῷ καθομολογήσηται αὐτὴν, κατὰ τὸ δικαίωμα τῶν θυγατέρων ποιήσει αὐτῇ.
\vs{10}Ἐὰν δὲ ἄλλην λάβῃ ἑαυτῷ, τὰ δέοντα καὶ τὸν ἱματισμὸν καὶ τὴν ὁμιλίαν αὐτῆς οὐκ ἀποστερήσει.
\vs{11}Ἐὰν δὲ τὰ τρία ταῦτα μὴ ποιήσῃ αὐτῇ, ἐξελεύσεται δωρεὰν ἄνευ ἀργυρίου.
\vs{12}Ἐὰν δὲ πατάξῃ τις τινὰ, καὶ ἀποθάνῃ, θανάτῳ θανατούσθω.
\vs{13}Ὁ δὲ οὐχ ἑκὼν, ἀλλὰ ὁ Θεὸς παρέδωκεν εἰς τὰς χεῖρας αὐτοῦ, δώσω σοι τόπον οὗ φεύξεται ἐκεῖ ὁ φονεύσας.
\vs{14}Ἐὰν δέ τις ἐπιθῆται τῷ πλησίον ἀποκτεῖναι αὐτὸν δόλῳ, καὶ καταφύγῃ, ἀπὸ τοῦ θυσιαστηρίου μου λήψῃ αὐτὸν θανατῶσαι.
\vs{15}Ὃς τύπτει πατέρα αὐτοῦ ἢ μητέρα αὐτοῦ, θανάτῳ θανατούσθω.
\vs{16}Ὁ κακολογῶν πατέρα αὐτοῦ ἢ μητέρα αὐτοῦ, τελευτήσει θανάτῳ.
\vs{17}Ὃς ἐὰν κλέψῃ τις τινὰ τῶν υἱῶν Ἰσραὴλ, καὶ καταδυναστεύσας αὐτὸν ἀποδῶται, καὶ εὑρεθῇ ἐν αὐτῷ, θανάτῳ τελευτάτω.
\vs{18}Ἐὰν δὲ λοιδορῶνται δύο ἄνδρες, καὶ πατάξωσι τὸν πλησίον λίθῳ ἢ πυγμῇ, καὶ μὴ ἀποθάνῃ, κατακλιθῇ δὲ ἐπὶ τὴν κοίτην,
\vs{19}ἐὰν ἐξαναστὰς ὁ ἄνθρωπος περιπατήσῃ ἔξω ἐπὶ ῥάβδου, ἀθῶος ἔσται ὁ πατάξας· πλὴν τῆς ἀργείας αὐτοῦ ἀποτίσει, καὶ τὰ ἰατρεῖα.
\vs{20}Ἐὰν δέ τις πατάξῃ τὸν παῖδα αὐτοῦ ἢ τὴν παιδίσκην αὐτοῦ ἐν ῥάβδῳ, καὶ ἀποθάνῃ ὑπὸ τὰς χεῖρας αὐτοῦ, δίκῃ ἐκδικηθήσεται.
\vs{21}Ἐὰν δὲ διαβιώσῃ ἡμέραν μίαν ἢ δύο, οὐκ ἐκδικηθήτω· τὸ γὰρ ἀργύριον αὐτοῦ ἐστιν.
\vs{22}Ἐὰν δὲ μάχωνται δύο ἄνδρες, καὶ πατάξωσι γυναῖκα ἐν γαστρὶ ἔχουσαν, καὶ ἐξέλθῃ τὸ παιδίον αὐτῆς μὴ ἐξεικονισμένον, ἐπιζήμιον ζημιωθήσεται· καθότι ἂν ἐπιβάλῃ ὁ ἀνὴρ τῆς γυναικὸς, δώσει μετὰ ἀξιώματος.
\vs{23}Ἐὰν δὲ ἐξεικονισμένον ᾖ, δώσει ψυχὴν ἀντὶ ψυχῆς,
\vs{24}Ὀφθαλμὸν ἀντὶ ὀφθαλμοῦ, ὀδόντα ἀντὶ ὀδόντος, χεῖρα ἀντὶ χειρὸς, πόδα ἀντὶ ποδὸς,
\vs{25}κατάκαυμα ἀντὶ κατακαύματος, τραῦμα ἀντὶ τραύματος, μώλωπα ἀντὶ μώλωπος.
\vs{26}Ἐὰν δέ τις πατάξῃ τὸν ὀφθαλμὸν τοῦ οἰκέτου αὐτοῦ, ἢ τὸν ὀφθαλμὸν τῆς θεραπαίνης αὐτοῦ, καὶ ἐκτυφλώσῃ, ἐλευθέρους ἐξαποστελεῖ αὐτοὺς ἀντὶ τοῦ ὀφθαλμοῦ αὐτῶν.
\vs{27}Ἐὰν δὲ τὸν ὀδόντα τοῦ οἰκέτου, ἢ τὸν ὀδόντα τῆς θεραπαίνης αὐτοῦ ἐκκόψῃ, ἐλευθέρους ἐξαποστελεῖ αὐτοὺς ἀντὶ τοῦ ὀδόντος αὐτῶν.
\vs{28}Ἐὰν δὲ κερατίσῃ ταῦρος ἄνδρα ἢ γυναῖκα καὶ ἀποθάνῃ, λίθοις λιθοβοληθήσεται ὁ ταῦρος, καὶ οὐ βρωθήσεται τὰ κρέα αὐτοῦ· ὁ δὲ κύριος τοῦ ταύρου ἀθῶος ἔσται.
\vs{29}Ἐὰν δὲ ὁ ταῦρος κερατιστὴς ᾖ πρὸ τῆς χθὲς καὶ πρὸ τῆς τρίτης, καὶ διαμαρτύρωνται τῷ κυρίῳ αὐτοῦ, καὶ μὴ ἀφανίσῃ αὐτὸν, ἀνέλῃ δὲ ἄνδρα ἢ γυναῖκα, ὁ ταῦρος λιθοβοληθήσεται, καὶ ὁ κύριος αὐτοῦ προσαποθανεῖται.
\vs{30}Ἐὰν δὲ λύτρα ἐπιβληθῇ αὐτῷ, δώσει λύτρα τῆς ψυχῆς αὐτοῦ ὅσα ἐὰν ἐπιβάλωσιν αὐτῶ.
\vs{31}Ἐὰν δὲ υἱὸν ἢ θυγατέρα κερατίσῃ, κατὰ τὸ δικαίωμα τοῦτο ποιήσωσιν αὐτῷ.
\vs{32}Ἐὰν δὲ παῖδα κερατίσῃ ὁ ταῦρος ἢ παιδίσκην, ἀργυρίου τριάκοντα δίδραχμα δώσει τῷ κυρίῳ αὐτῶν, καὶ ὁ ταῦρος λιθοβοληθήσεται.
\vs{33}Ἐὰν δέ τις ἀνοίξῃ λάκκον ἢ λατομήσῃ λάκκον, καὶ μὴ καλύψῃ αὐτὸν, καὶ ἐμπέσῃ ἐκεῖ μόσχος ἢ ὄνος,
\vs{34}ὁ κύριος τοῦ λάκκου ἀποτίσει, ἀργύριον δώσει τῷ κυρίῳ αὐτῶν· τὸ δὲ τετελευτηκὸς αὐτῷ ἔσται.
\vs{35}Ἐὰν δὲ κερατίσῃ τινὸς ταῦρος τόν ταῦρον τοῦ πλησίον, καὶ τελευτήσῃ, ἀποδώσονται τὸν ταῦρον τὸν ζῶντα, καὶ διελοῦνται τὸ ἀργύριον αὐτοῦ, καὶ τὸν ταῦρον τὸν τεθνηκότα διελοῦνται.
\vs{36}Ἐὰν δὲ γνωρίζηται ὁ ταῦρος ὅτι κερατιστής ἐστι πρὸ τῆς χθὲς καὶ πρὸ τῆς τρίτης ἡμέρας, καὶ διαμεμαρτυρημένοι ὦσι τῷ κυρίῳ αὐτοῦ· καὶ μὴ ἀφανίσῃ αὐτὸν, ἀποτίσει ταῦρον ἀντὶ ταύρου, ὁ δὲ τετελευτηκὼς αὐτῷ ἔσται.

\vs{37}Ἐὰν δέ τις κλέψῃ μόσχον ἢ πρόβατον, καὶ σφάξῃ ἢ ἀποδῶται, πέντε μόσχους ἀποτίσει ἀντὶ τοῦ μόσχου, καὶ τέσσερα πρόβατα ἀντὶ τοῦ προβάτου.

\ch{22}Ἐὰν δὲ ἐν τῷ διορύγματι εὑρεθῇ ὁ κλέπτης, καὶ πληγεὶς ἀποθάνῃ, οὐκ ἔστιν αὐτῷ φόνος.
\vs{2}Ἐὰν δὲ ἀνατείλῃ ὁ ἥλιος ἐπʼ αὐτῷ, ἔνοχός ἐστιν, ἀνταποθανεῖται· ἐὰν δὲ μὴ ὑπάρχῃ αὐτῷ, πραθήτω ἀντὶ τοῦ κλέμματος.
\vs{3}Ἐὰν δὲ καταλειφθῇ καὶ εὑρεθῇ ἐν τῇ χειρὶ αὐτοῦ τὸ κλέμμα ἀπό τε ὄνου ἕως προβάτου ζῶντα, διπλᾶ αὐτὰ ἀποτίσει.
\vs{4}Ἐὰν δὲ καταβοσκήσῃ τις ἀγρὸν ἢ ἀμπελῶνα, καὶ ἀφῇ τὸ κτῆνος αὐτοῦ καταβοσκῆσαι ἀγρὸν ἕτερον, ἀποτίσει ἐκ τοῦ ἀγροῦ αὐτοῦ κατὰ τὸ γέννημα αὐτοῦ· ἐὰν δὲ πάντα τὸν ἀγρὸν καταβοσκήσῃ, τὰ βέλτιστα τοῦ ἀγροῦ αὐτοῦ καὶ τὰ βέλτιστα τοῦ ἀμπελῶνος αὐτοῦ ἀποτίσει.
\vs{5}Ἐὰν δὲ ἐξελθὸν πῦρ εὕρῃ ἀκάνθας, καὶ προσεμπρήσῃ ἅλωνας ἢ στάχυς ἢ πεδίον, ἀποτίσει ὁ τὸ πῦρ ἐκκαύσας.

\vs{6}Ἐὰν δέ τις δῷ τῷ πλησίον ἀργύριον ἢ σκεύη φυλάξαι, καὶ κλαπῇ ἐκ τῆς οἰκίας τοῦ ἀνθρώπου, ἐὰν εὑρεθῇ ὁ κλέψας, ἀποτίσει τὸ διπλοῦν.
\vs{7}Ἐὰν δὲ μὴ εὑρεθῇ ὁ κλέψας, προσελεύσεται ὁ κύριος τῆς οἰκίας ἐνώπιον τοῦ Θεοῦ, καὶ ὀμεῖται ἦ μὴν μὴ αὐτὸν πεπονηρεῦσθαι ἐφʼ ὅλης τῆς παρακαταθήκης τοῦ πλησίον,
\vs{8}κατὰ πᾶν ῥητὸν ἀδίκημα, περί τε μόσχου, καὶ ὑποζυγίου, καὶ προβάτου, καὶ ἱματίου, καὶ πάσης ἀπωλίας τῆς ἐνκαλουμένης· ὅ, τι οὖν ἂν ᾖ, ἐνώπιον τοῦ Θεοῦ ἐλεύσεται ἡ κρίσις ἀμφοτέρων, καὶ ὁ ἁλοὺς διὰ τοῦ Θεοῦ, ἀποτίσει διπλοῦν τῷ πλησίον.
\vs{9}Ἐὰν δέ τις δῷ τῷ πλησίον ὑποζύγιον ἢ μόσχον ἢ πρόβατον ἢ πᾶν κτῆνος φυλάξαι, καὶ συντριβῇ ἢ τελευτήσῃ ἢ αἰχμάλωτον γένηται, καὶ μηδεὶς γνῷ,
\vs{10}ὅρκος ἔσται τοῦ Θεοῦ ἀνὰ μέσον ἀμφοτέρων, ἦ μὴν μὴ αὐτὸν πεπονηρεῦσθαι καθόλου τῆς παρακαταθήκης τοῦ πλησίον· καὶ οὕτως προσδέξεται ὁ κύριος αὐτοῦ, καὶ οὐκ ἀποτίσει.
\vs{11}Ἐὰν δὲ κλαπῇ παρʼ αὐτοῦ, ἀποτίσει τῷ κυρίῳ·
\vs{12}Ἐὰν δὲ θηριάλωτον γένηται, ἄξει αὐτὸν ἐπὶ τὴν θήραν, καὶ οὐκ ἀποτίσει.
\vs{13}Ἐὰν δὲ αἰτήσῃ τις παρὰ τοῦ πλησίον, καὶ συντριβῇ ἢ ἀποθάνῃ ἢ αἰχμάλωτον γένηται, ὁ δὲ κύριος μὴ ᾖ μετʼ αὐτοῦ, ἀποτίσει.
\vs{14}Ἐὰν δὲ ὁ κύριος ᾖ μετʼ αὐτοῦ, οὐκ ἀποτίσει· ἐὰν δὲ μισθωτὸς ᾖ, ἔσται αὐτῷ ἀντὶ τοῦ μισθοῦ αὐτοῦ.

\vs{15}Ἐὰν δὲ ἀπατήσῃ τις παρθένον ἀμνήστευτον, καὶ κοιμηθῇ μετʼ αὐτῆς, φερνῇ φερνιεῖ αὐτὴν αὐτῷ γυναῖκα.
\vs{16}Ἐὰν δὲ ἀνανεύων ἀνανεύσῃ, καὶ μὴ βούληται ὁ πατὴρ αὐτῆς δοῦναι αὐτὴν αὐτῷ γυναῖκα, ἀργύριον ἀποτίσει τῷ πατρὶ καθʼ ὅσον ἐστὶν ἡ φερνὴ τῶν παρθένων.
\vs{17}Φαρμακοὺς οὐ περιποιήσετε.
\vs{18}Πᾶν κοιμώμενον μετὰ κτήνους θανάτῳ ἀποκτενεῖτε αὐτούς.
\vs{19}Ὁ θυσιάζων θεοῖς θανάτῳ ἐξολοθρευθήσεται, πλὴν Κυρίῳ μόνῳ.

\vs{20}Καὶ προσήλυτον οὐ κακώσετε, οὐδὲ μὴ θλίψητε αὐτόν· ἦτε γάρ προσήλυτοι ἐν γῇ Αἰγύπτῳ.
\vs{21}Πᾶσαν χήραν καὶ ὀρφανὸν οὐ κακώσετε.
\vs{22}Ἐὰν δὲ κακίᾳ κακώσητε αὐτοὺς, καὶ κεκράξαντες καταβοήσωσι πρός με, ἀκοῇ εἰσακούσομαι τῆς φωνῆς αὐτῶν,
\vs{23}καὶ ὀργισθήσομαι θυμῷ, καὶ ἀποκτενῶ ὑμᾶς μαχαίρᾳ, καὶ ἔσονται αἱ γυναῖκες ὑμῶν χῆραι, καὶ τὰ παιδία ὑμῶν ὀρφανά.
\vs{24}Ἐὰν δὲ ἀργύριον ἐκδανείσῃς τῷ ἀδελφῷ τῷ πενιχρῷ παρὰ σοὶ, οὐκ ἔσῃ αὐτὸν κατεπείγων, οὐκ ἐπιθήσεις αὐτῷ τόκον.
\vs{25}Ἐὰν δὲ ἐνεχύρασμα ἐνεχυράσῃς τὸ ἱμάτιον τοῦ πλησίον, πρὸ δυσμῶν ἡλίου ἀποδώσεις αὐτῷ·
\vs{26}Ἔστι γὰρ τοῦτο περιβόλαιον αὐτοῦ, μόνον τοῦτο τὸ ἱμάτιον ἀσχημοσύνης αὐτοῦ· ἐν τίνι κοιμηθήσεται; Ἐὰν οὖν καταβοήσῃ πρός μέ, εἰσακούσομαι αὐτοῦ· ἐλεήμων γάρ εἰμι.
\vs{27}Θεοὺς οὐ κακολογήσεις, καὶ ἄρχοντα τοῦ λαοῦ σου οὐ κακῶς ἐρεῖς.
\vs{28}Ἀπαρχὰς ἅλωνος καὶ ληνοῦ σου οὐ καθυστερήσεις· τὰ πρωτότοκα τῶν υἱῶν σου δώσεις ἐμοί.
\vs{29}Οὕτω ποιήσεις τὸν μόσχον σου καὶ τὸ πρόβατόν σου καὶ τὸ ὑποζύγιόν σου· ἑπτὰ ἡμέρας ἔσται ὑπὸ τὴν μητέρα, τῇ δὲ ὀγδόῃ ἡμέρᾳ ἀποδώσεις μοι αὐτό.
\vs{30}Καὶ ἄνδρες ἅγιοι ἔσεσθέ μοι· καὶ κρέας θηριάλωτον οὐκ ἔδεσθε, τῷ κυνὶ ἀποῤῥίψατε αὐτό.

\ch{23}
Οὐ παραδέξῃ ἀκοὴν ματαίαν· οὐ συγκαταθήσῃ μετὰ τοῦ ἀδίκου γενέσθαι μάρτυς ἄδικος.
\vs{2}Οὐκ ἔσῃ μετὰ πλειόνων ἐπὶ κακίᾳ· οὐ προστεθήσῃ μετὰ πλήθους ἐκκλῖναι μετὰ τῶν πλειόνων, ὥστε ἑκκλεῖσαι κρίσιν.
\vs{3}Καὶ πένητα οὐκ ἐλεήσεις ἐν κρίσει.
\vs{4}Ἐὰν δὲ συναντήσῃς τῷ βοῒ τοῦ ἐχθροῦ σου, ἢ τῷ ὑποζυγίῳ αὐτοῦ πλανωμένοις, ἀποστρέψας ἀποδώσεις αὐτῷ.
\vs{5}Ἐὰν δὲ ἴδῃς τὸ ὑποζύγιον τοῦ ἐχθροῦ σου πεπτωκὸς ὑπὸ τὸν γόμον αὐτοῦ, οὐ παρελεύσῃ αὐτὸ, ἀλλὰ συναρεῖς αὐτὸ μετʼ αὐτοῦ.

\vs{6}Οὐ διαστρέψεις κρίμα πένητος ἐν κρίσει αὐτοῦ.
\vs{7}Ἀπὸ παντὸς ῥήματος ἀδίκου ἀποστήσῃ· ἀθῷον καὶ δίκαιον οὐκ ἀποκτενεῖς· καὶ οὐ δικαιώσεις τὸν ἀσεβῆ ἕνεκεν δώρων.
\vs{8}Καὶ δῶρα οὐ λήψῃ· τὰ γὰρ δῶρα ἐκτυφλοῖ ὀφθαλμοὺς βλεπόντων, καὶ λυμαῖνεται ῥήματα δίκαια.
\vs{9}Καὶ προσήλυτον οὐ θλίψετε· ὑμεῖς γὰρ οἴδατε τὴν ψυχὴν τοῦ προσηλύτου· αὐτοὶ γὰρ προσήλυτοι ἦτε ἐν γῇ Αἰγύπτῳ.
\vs{10}Ἓξ ἔτη σπερεῖς τὴν γῆν σου, καὶ συνάξεις τὰ γεννήματα αὐτῆς.
\vs{11}Τῷ δὲ ἑβδόμῳ ἄφεσιν ποιήσεις, καὶ ἀνήσεις αὐτὴν, καὶ ἔδονται οἱ πτωχοὶ τοῦ ἔθνους σου· τὰ δὲ ὑπολειπόμενα ἔδεται τὰ ἄγρια θηρία· οὕτως ποιήσεις τὸν ἀμπελῶνά σου, καὶ τὸν ἐλαιῶνά σου.
\vs{12}Ἓξ ἡμέρας ποιήσεις τὰ ἔργα σου, τῇ δὲ ἡμέρᾳ τῇ ἑβδόμῃ, ἀνάπαυσις· ἵνα ἀναπαύσηται ὁ βοῦς σου, καὶ τὸ ὑποζύγιόν σου, καὶ ἵνα ἀναψύξῃ ὁ υἱὸς τῆς παιδίσκης σου καὶ ὁ προσήλυτος.
\vs{13}Πάντα ὅσα εἴρηκα πρὸς ὑμᾶς, φυλάξασθε· καὶ ὄνομα θεῶν ἑτέρων οὐκ ἀναμνησθήσεσθε, οὐδὲ μὴ ἀκουσθῇ ἐκ τοῦ στόματος ὑμῶν.

\vs{14}Τρεῖς καιροὺς τοῦ ἐνιαυτοῦ ἑορτάσατέ μοι.
\vs{15}Τὴν ἑορτὴν τῶν ἀζύμων φυλάξασθε ποιεῖν· ἑπτὰ ἡμέρας ἔδεσθε ἄζυμα, καθάπερ ἐνετειλάμην σοι κατὰ τὸν καιρὸν τοῦ μηνὸς τῶν νέων· ἐν γὰρ αὐτῷ ἐξῆλθες ἐξ Αἰγύπτου· οὐκ ὀφθήσῃ ἐνώπίον μου κενός.
\vs{16}Καὶ ἑορτὴν θερισμοῦ πρωτογεννημάτων ποιήσεις τῶν ἔργων σου, ὧν ἐὰν σπείρῃς ἐν τῷ ἀγρῷ σου, καὶ ἑορτὴν συντελείας ἐπʼ ἐξόδου τοῦ ἐνιαυτοῦ ἐν τῇ συναγωγῇ τῶν ἔργων σου τῶν ἐκ τοῦ ἀγροῦ σου.
\vs{17}Τρεῖς καιροὺς τοῦ ἐνιαυτοῦ ὀφθήσεται πᾶν ἀρσενικόν σου ἐνώπιον Κυρίου τοῦ Θεοῦ σου.
\vs{18}Ὅταν γὰρ ἐκβάλω τὰ ἔθνη ἀπὸ προσώπου σου, καὶ ἐμπλατύνω τὰ ὅριά σου, οὐ θύσεις ἐπὶ ζύμῃ αἷμα θυμιάματός μου, οὐδὲ μὴ κοιμηθῇ στέαρ τῆς ἑορτῆς μου ἕως πρωΐ.
\vs{19}Τὰς ἀπαρχὰς τῶν πρωτογενημάτων τῆς γνς σου εἰσοίσεις εἰς τὸν οἶκον Κυρίου τοῦ Θεοῦ σου· οὐχ ἑψήσεις ἄρνα ἐν γάλακτι μητρὸς αὐτοῦ.
\vs{20}Καὶ ἰδοὺ ἐγὼ ἀποστέλλω τὸν ἄγγελόν μου πρὸ προσώπου σου, ἵνα φυλάξῃ σε ἐν τῇ ὁδῷ, ὅπως εἰσαγάγῃ σε εἰς τὴν γῆν, ἣν ἡτοίμασά σοι.
\vs{21}Πρόσεχε σεαυτῷ, καὶ εἰσάκουε αὐτοῦ, καὶ μὴ ἀπείθει αὐτῷ, οὐ γὰρ μὴ ὑποστείληταί σε· τὸ γὰρ ὄνομά μου ἐστὶν ἐπʼ αὐτῷ.
\vs{22}Ἐὰν ἀκοῇ ἀκούσητε τῆς ἐμῆς φωνῆς, καὶ ποιήσῃς πάντα ὅσα ἂν ἐντείλωμαί σοι, καὶ φυλάξητε τὴν διαθήκην μου, ἔσεσθέ μοι λαὸς περιούσιος ἀπὸ πάντων τῶν ἐθνῶν· ἐμὴ γάρ ἐστι πᾶσα ἡ γῆ· ὑμεῖς δὲ ἔσεσθέ μοι βασίλειον ἱεράτευμα, καὶ ἔθνος ἅγιον· ταῦτα τὰ ῥήματα ἐρεῖς τοῖς υἱοῖς Ἰσραὴλ, ἐὰν ἀκοῇ ἀκούσητε τῆς φωνῆς μου, καὶ ποιήσητε πάντα ὅσα ἂν εἴπω σοι, ἐχθρεύσω τοῖς ἐχθροῖς σου, καὶ ἀντικείσομαι τοῖς ἀντικειμένοις σοι.
\vs{23}Πορεύσεται γὰρ ὁ ἄγγελός μου ἡγούμενός σου, καὶ εἰσάξει σε πρὸς τὸν Ἀμοῤῥαῖον, καὶ Χετταῖον, καὶ Φερεζαῖον, καὶ Χαναναῖον, καὶ Γεργεσαῖον, καὶ Εὑαῖον, καὶ Ἰεβουσαῖον, καὶ ἐκτρίψω αὐτούς.
\vs{24}Οὐ προσκυνήσεις τοῖς θεοῖς αὐτῶν, οὐδὲ μὴ λατρεύσῃς αὐτοῖς· οὐ ποιήσεις κατὰ τὰ ἔργα αὐτῶν· ἀλλὰ καθαιρέσει καθελεῖς, καὶ συντρίβων συντρίψεις τὰς στήλας αὐτῶν.
\vs{25}Καὶ λατρεύσεις Κυρίῳ τῷ Θεῷ σου· καὶ εὐλογήσω τὸν ἄρτον σου καὶ τὸν οἶνόν σου καὶ τὸ ὕδωρ σου, καὶ ἀποστρέψω μαλακίαν ἀφʼ ὑμῶν.
\vs{26}Οὐκ ἔσται ἄγονος, οὐδὲ στεῖρα ἐπὶ τῆς γῆς σου· τὸν ἀριθμὸν τῶν ἡμερῶν σου ἀναπληρῶν ἀναπληρώσω.
\vs{27}Καὶ τὸν φόβον ἀποστελῶ ἡγούμενόν σου, καὶ ἐκστήσω πάντα τὰ ἔθνη, εἰς οὓς σὺ εἰσπορεύῃ εἰς αὐτούς· καὶ δώσω πάντας τοὺς ὑπεναντίους σου φυγάδας.
\vs{28}Καὶ ἀποστελῶ τὰς σφηκίας προτέρας σου· καὶ ἐκβαλεῖς τοὺς Ἀμοῤῥαίους, καὶ τοὺς Εὑαίους, καὶ τοὺς Χαναναίους, καὶ τοὺς Χετταίους ἀπὸ σοῦ.
\vs{29}Οὐκ ἐκβαλῶ αὐτοὺς ἐν ἐνιαυτῷ ἑνὶ, ἵνα μὴ γένηται ἡ γῆ ἔρημος, καὶ πολλὰ γένηται ἐπὶ σὲ τὰ θηρία τῆς γῆς.
\vs{30}Κατὰ μικρὸν ἐκβαλῶ αὐτοὺς ἀπὸ σοῦ, ἕως ἂν αὐξηθῇς καὶ κληρονομήσῃς τὴν γῆν.
\vs{31}Καὶ θήσω τὰ ὅριά σου ἀπὸ τῆς ἐρυθρᾶς θαλάσσης, ἕως τῆς θαλάσσης τῆς Φυλιστιείμ· καὶ ἀπὸ τῆς ἐρήμου, ἕως τοῦ μεγάλου ποταμοῦ Εὐφράτου· καὶ παραδώσω εἰς τὰς χεῖρας ὑμῶν τοὺς ἐγκαθημένους ἐν τῇ γῇ, καὶ ἐκβαλῶ αὐτοὺς ἀπὸ σοῦ.
\vs{32}Οὐ συγκαταθήσῃ αὐτοῖς καὶ τοῖς θεοῖς αὐτῶν διαθήκην.
\vs{33}Καὶ οὐκ ἐνκαθήσονται ἐν τῇ γῇ σου, ἵνα μὴ ἁμαρτεῖν σε ποιήσωσι πρὸς μέ· ἐὰν γὰρ δουλεύσῃς τοῖς θεοῖς αὐτῶν, οὗτοι ἔσονταί σοι πρόσκομμα.

\ch{24}
Καὶ Μωυσῇ εἶπεν, ἀνάβηθι πρὸς τὸν Κύριον σὺ καὶ Ἀαρὼν, καὶ Ναδὰβ, καὶ Ἀβιοὺδ, καὶ ἑβδομήκοντα τῶν πρεσβυτέρων Ἰσραήλ· καὶ προσκυνήσουσι μακρόθεν τῷ Κυρίῳ.
\vs{2}Καὶ ἐγγιεῖ Μωσῆς μόνος πρὸς τὸν Θεὸν, αὐτοὶ δὲ οὐκ ἐγγιοῦσιν, ὁ δὲ λαὸς οὐ συναναβήσεται μετʼ αὐτῶν.
\vs{3}Εἰσῆλθε δὲ Μωυσῆς, καὶ διηγήσατο τῷ λαῷ πάντα τὰ ῥήματα τοῦ Θεοῦ καὶ τὰ δικαιώματα· ἀπεκρίθη δὲ πᾶς ὁ λαὸς φωνῇ μιᾷ, λέγοντες, πάντας τοὺς λόγους, οὓς ἐλάλησε Κύριος, ποιήσομεν, καὶ ἀκουσόμεθα.
\vs{4}Καὶ ἔγραψε Μωυσῆς πάντα τὰ ῥήματα Κυρίου· ὀρθρίσας δὲ Μωυσῆς τὸ πρωῒ ᾠκοδόμησε θυσιαστήριον ὑπὸ τὸ ὄρος, καὶ δώδεκα λίθους εἰς τὰς δώδεκα φυλὰς τοῦ Ἰσραήλ.
\vs{5}Καὶ ἐξαπέστειλε τοὺς νεανίσκους τῶν υἱῶν Ἰσραήλ, καὶ ἀνήνεγκαν ὁλοκαυτώματα· καὶ ἔθυσαν θυσίαν σωτηρίου τῷ Θεῷ μοσχάρια.
\vs{6}Λαβὼν δὲ Μωυσῆς τὸ ἥμισυ τοῦ αἵματος, ἐνέχεεν εἰς κρατῆρας, τὸ δὲ ἥμισυ τοῦ αἵματος προσέχεε πρὸς τὸ θυσιαστήριον.
\vs{7}Καὶ λαβὼν τὸ βιβλίον τῆς διαθήκης, ἀνέγνω εἰς τὰ ὦτα τοῦ λαοῦ· καὶ εἶπαν, πάντα ὅσα ἐλάλησε Κύριος, ποιήσομεν καὶ ἀκουσόμεθα.
\vs{8}Λαβὼν δὲ Μωυσῆς τὸ αἷμα, κατεσκέδασε τοῦ λαοῦ, καὶ εἶπεν, ἰδοὺ τὸ αἷμα τῆς διαθήκης, ἧς διέθετο Κύριος πρὸς ὑμᾶς περὶ πάντων τῶν λόγων τούτων.

\vs{9}Καὶ ἀνέβη Μωυσῆς καὶ Ἀαρὼν, καὶ Ναδὰβ, καὶ Ἀβιοῦδ, καὶ ἑβδομήκοντα τῆς γερουσίας Ἰσραήλ.
\vs{10}Καὶ εἶδον τὸν τόπον οὗ εἱστήκει ὁ Θεὸς τοῦ Ἰσραήλ· καὶ τὰ ὑπὸ τοὺς πόδας αὐτοῦ, ὡσεὶ ἔργον πλίνθου σαπφείρου, καὶ ὥσπερ εἶδος στερεώματος τοῦ οὐρανοῦ τῇ καθαριότητι.
\vs{11}Καὶ τῶν ἐπιλέκτων τοῦ Ἰσραὴλ οὐ διεφώνησεν οὐδὲ εἷς· καὶ ὤφθησαν ἐν τῷ τόπῳ τοῦ Θεοῦ, καὶ ἔφαγον καὶ ἔπιον.
\vs{12}Καὶ εἶπε Κύριος πρὸς Μωυσῆν, ἀνάβηθι πρὸς με εἰς τὸ ὄρος, καὶ ἴσθι ἐκεῖ· καὶ δώσω σοι τὰ πυξία τὰ λίθινα, τὸν νόμον καὶ τὰς ἐντολας, ἃς ἔγραψα νομοθετῆσαι αὐτοῖς.
\vs{13}Καὶ ἀναστὰς Μωυσῆς καὶ Ἰησοῦς ὁ παρεστηκὼς αὐτῷ, ἀνέβησαν εἰς τὸ ὄρος τοῦ Θεοῦ.
\vs{14}Καὶ τοῖς πρεσβυτέροις εἶπαν, ἡσυχάζετε αὐτοῦ, ἕως ἀναστρέψωμεν πρὸς ὑμᾶς· καὶ ἰδοὺ Ἀαρὼν καὶ Ὢρ μεθʼ ὑμῶν· ἐάν τινι συμβῇ κρίσις, προσπορευέσθωσαν αὐτοῖς.
\vs{15}Καὶ ἀνέβη Μωυσῆς καὶ Ἰησοῦς εἰς τὸ ὄρος· καὶ ἐκάλυψεν ἡ νεφέλη τὸ ὄρος.
\vs{16}Καὶ κατέβη ἡ δόξα τοῦ Θεοῦ ἐπὶ τὸ ὄρος τὸ Σινὰ, καὶ ἐκάλυψεν αὐτὸ ἡ νεφέλη ἓξ ἡμέρας· καὶ ἐκάλεσε Κύριος τὸν Μωυσῆν τῇ ἡμέρᾳ τῇ ἑβδόμῃ ἐκ μέσου τῆς νεφέλης.
\vs{17}Τὸ δὲ εἶδος τῆς δόξης Κυρίου, ὡσεὶ πῦρ φλέγον ἐπὶ τῆς κορυφῆς τοῦ ὄρους, ἐναντίον τῶν υἱῶν Ἰσραήλ.
\vs{18}Καὶ εἰσῆλθε Μωυσῆς εἰς τὸ μέσον τῆς νεφέλης, καὶ ἀνέβη εἰς τὸ ὄρος· καὶ ἦν ἐκεῖ ἐν τῷ ὄρει τεσσεράκοντα ἡμέρας καὶ τεσσαράκοντα νύκτας.

\ch{25}
Καὶ ἐλάλησε Κύριος πρὸς Μωυσῆν, λέγων,
\vs{2}εἶπον τοῖς υἱοῖς Ἰσραὴλ, καὶ λάβετε ἀπαρχὰς παρὰ πάντων, οἷς ἂν δόξῃ τῇ καρδίᾳ, καὶ λήψεσθε τὰς ἀπαρχάς μου.
\vs{3}Καὶ αὕτη ἐστὶν ἡ ἀπαρχὴ, ἣν λήψεσθε παρʼ αὐτῶν· χρυσίον, καὶ ἀργύριον, καὶ χαλκὸν,
\vs{4}καὶ ὑάκινθον, καὶ πορφύραν, καὶ κόκκινον διπλοῦν, καὶ βύσσον κεκλωσμένην, καὶ τρίχας αἰγείας,
\vs{5}καὶ δέρματα κριῶν ἠρυθροδανωμένα, καὶ δέρματα ὑακίνθινα, καὶ ξύλα ἄσηπτα,
\vs{5a}καὶ ἔλαιον εἰς τὴν φαῦσιν, θυμιάματα εἰς τὸ ἔλαιον τῆς χρίσεως, καὶ εἰς τὴν σύνθεσιν τοῦ θυμιάματος,
\vs{7}καὶ λίθους Σαρδίου, καὶ λίθους εἰς τὴν γλυφὴν εἰς τὴν ἐπωμίδα, καὶ τὸν ποδήρη.
\vs{8}Καὶ ποιήεις μοι ἁγίασμα, καὶ ὀφθήσομαι ἐν ὑμῖν.
\vs{9}Καὶ ποιήσεις μοι κατὰ πάντα ὅσα σοι δεικνύω ἐν τῷ ὄρει, τὸ παράδειγμα τῆς σκηνῆς, καὶ τὸ παράδειγμα πάντων τῶν σκευῶν αὐτῆς· οὕτω ποιήσεις.
\vs{10}Καὶ ποιήσεις κιβωτὸν μαρτυρίου ἐκ ξύλων ἀσήπτων, δύο πήχεων καὶ ἡμίσους τὸ μῆκος, καὶ πήχεος καὶ ἡμίσους τὸ πλάτος, καὶ πήχεως καὶ ἡμίσους τὸ ὕψος.
\vs{11}Καὶ καταχρυσώσεις αὐτὴν χρυσίῳ καθαρῷ, ἔσωθεν καὶ ἔξωθεν χρυσώσεις αὐτήν· καὶ ποιήσεις αὐτῇ κυμάτια χρυσᾶ στρεπτὰ κύκλῳ.
\vs{12}Καὶ ἐλάσεις αὐτῇ τέσσαρας δακτυλίους χρυσοῦς, καὶ ἐπιθήσεις ἐπὶ τὰ τέσσαρα κλίτη· δύο δακτυλίους ἐπὶ τὸ κλίτος τὸ ἓν, καὶ δύο δακτυλίους ἐπὶ τὸ κλίτος τὸ δεύτερον.
\vs{13}Ποιήσεις δὲ ἀναφορεῖς ξύλα ἄσηπτα, καὶ καταχρυσώσεις αὐτὰ χρυσίῳ·
\vs{14}Καὶ εἰσάξεις τοὺς ἀναφορεῖς εἰς τοὺς δακτυλίους τοὺς ἐν τοῖς κλίτεσι τῆς κιβωτοῦ, αἴρειν τὴν κιβωτὸν ἐν αὐτοῖς.
\vs{15}Ἐν τοῖς δακτυλίοις τῆς κιβωτοῦ ἔσονται οἱ ἀναφορεῖς ἀκίνητοι.
\vs{16}Καὶ ἐμβαλεῖς εἰς τὴν κιβωτὸν τὰ μαρτύρια, ἃ ἂν δῶ σοι.
\vs{17}Καὶ ποιήσεις ἱλαστήριον ἐπίθεμα χρυσίου καθαροῦ, δύο πήχεων καὶ ἡμίσους τὸ μῆκος, καὶ πήχεως καὶ ἡμίσους τὸ πλάτος.
\vs{18}Καὶ ποιήσεις δύο χερουβὶμ χρυσοτορευτὰ, καὶ ἐπιθήσεις αὐτὰ ἐξ ἀμφοτέρων τῶν κλιτῶν τοῦ ἱλαστηρίου.
\vs{19}Ποιηθήσονται χεροὺβ εἷς ἐκ τοῦ κλίτους τούτου, καὶ χεροὺβ εἷς ἐκ τοῦ κλίτους τοῦ δευτέρου τοῦ ἱλαστηρίου· καὶ ποιήσεις τοὺς δύο χερουβὶμ ἐπὶ τὰ δύο κλίτη.
\vs{20}Ἔσονται οἱ χερουβὶμ ἐκτείνοντες τὰς πτέρυγας ἐπάνωθεν, συσκιάζοντες ἐν ταῖς πτέρυξιν αὐτῶν ἐπὶ τοῦ ἱλαστηρίου, καὶ τὰ πρόσωπα αὐτῶν εἰς ἄλληλα, εἰς τὸ ἱλαστήριον ἔσονται τὰ πρόσωπα τῶν χερουβίμ.
\vs{21}Καὶ ἐπιθήσεις τὸ ἱλαστήριον ἐπὶ τὴν κιβωτὸν ἄνωθεν, καὶ εἰς τὴν κιβωτὸν ἐμβαλεῖς τὰ μαρτύρια, ἃ ἂν δῶ σοι.
\vs{22}Καὶ γνωσθήσομαί σοι ἐκεῖθεν, καὶ λαλήσω σοι ἄνωθεν τοῦ ἱλαστηρίου ἀνὰ μέσον τῶν δύο χερουβὶμ, τῶν ὄντων ἐπὶ τῆς κιβωτοῦ τοῦ μαρτυρίου, καὶ κατὰ πάντα ὅσα ἐὰν ἐντείλωμαί σοι πρὸς τοὺς υἱοὺς Ἰσραήλ.
\vs{23}Καὶ ποιήσεις τράπεζαν χρυσῆν χρυσίου καθαροῦ, δύο πήχεων τὸ μῆκος, καὶ πήχεως τὸ εὖρος, καὶ πήχεως καὶ ἡμίσους τὸ ὕψος.
\vs{24}Καὶ ποιήσεις αὐτῇ στρεπτὰ κυμάτια χρυσᾶ κύκλῳ· καὶ ποιήσεις αὐτῇ στεφάνην παλαιστοῦ κύκλῳ·

\vs{25}Καὶ ποιήσεις στρεπτὸν κυμάτιον τῇ στεφάνῃ κύκλῳ.
\vs{26}Καὶ ποιήσεις τέσσαρας δακτυλίους χρυσοῦς, καὶ ἐπιθήσεις τοὺς τέσσαρας δακτυλίους ἐπὶ τὰ τέσσαρα μέρη τῶν ποδῶν αὐτῆς ὑπὸ τὴν στεφάνην.
\vs{27}Καὶ ἔσονται οἱ δακτύλιοι εἰς θήκας τοῖς ἀναφορεῦσιν, ὥστε αἴρειν ἐν αὐτοῖς τὴν τράπεζαν.
\vs{28}Καὶ ποιήσεις τοὺς ἀναφορεῖς ἐκ ξύλων ἀσήπτων, καὶ καταχρυσώσεις αὐτοὺς χρυσίῳ καθαρῷ, καὶ ἀρθήσεται ἐν αὐτοῖς ἡ τράπεζα.
\vs{29}Καὶ ποιήσεις τὰ τρυβλία αὐτῆς, καὶ τὰς θυΐσκας, καὶ τὰ σπονδεῖα, καὶ τοὺς κυάθους, ἐν οἷς σπείσεις ἐν αὐτοῖς, ἐκ χρυσίου καθαροῦ ποιήσεις αὐτά.
\vs{30}Καὶ ἐπιθήσεις ἐπὶ τὴν τράπεζαν ἄρτους ἐνωπίους ἐναντίον μου διαπαντός.

\vs{31}Καὶ ποιήσεις λυχνίαν ἐκ χρυσίου καθαροῦ, τορευτὴν ποιήσεις τὴν λυχνίαν· ὁ καυλὸς αὐτῆς, καὶ ὁ καλαμίσκοι, καὶ οἱ κρατῆρες, καὶ οἱ σφαιρωτῆρες, καὶ τὰ κρίνα ἐξ αὐτῆς ἔσται.
\vs{32}Ἓξ δὲ καλαμίσκοι ἐκπορευόμενοι ἐκ πλαγίων, τρεῖς καλαμίσκοι τῆς λυχνίας ἐκ τοῦ κλίτους τοῦ ἑνὸς αὐτῆς, καὶ τρεῖς καλαμίσκοι τῆς λυχνίας ἐκ τοῦ κλίτους τοῦ δευτέρου.
\vs{33}Καὶ τρεῖς κρατῆρες ἐκτετυπωμένοι καρυΐσκους· ἐν τῷ ἑνὶ καλαμίσκῳ σφαιρωτὴρ καὶ κρίνον· οὕτω τοῖς ἓξ καλαμίσκοις τοῖς ἐκπορευομένοις ἐκ τῆς λυχνίας.
\vs{34}Καὶ ἐν τῇ λυχνίᾳ τέσσαρες κρατῆρες ἐκτετυπωμένοι καρυΐσκους· ἐν τῷ ἑνὶ καλαμίσκῳ σφαιρωτῆρες, καὶ τὰ κρίνα αὐτῆς.
\vs{35}Ὁ σφαιρωτὴρ ὑπὸ τοὺς δύο καλαμίσκους ἐξ αὐτῆς· καὶ σφαιρωτὴρ ὑπὸ τοὺς τέσσαρας καλαμίσκους ἐξ αὐτῆς· οὕτω τοῖς ἓξ καλαμίσκοις τοῖς ἐκπορευομένοις ἐκ τῆς λυχνίας· καὶ ἐν τῇ λυχνίᾳ τέσσαρες κρατῆρες ἐκτετυπωμένοι καρυΐσκους.
\vs{36}Οἱ σφαιρωτῆρες καὶ οἱ καλαμίσκοι ἐξ αὐτῆς ἔστωσαν· ὅλη τορευτὴ ἐξ ἑνὸς χρυσίου καθαροῦ.
\vs{37}Καὶ ποιήσεις τοὺς λύχνους αὐτῆς ἑπτά· καὶ ἐπιθήσεις τοὺς λύχνους, καὶ φανοῦσιν ἐκ τοῦ ἑνὸς προσώπου.
\vs{38}Καὶ τὸν ἐπαρυστῆρα αὐτῆς, καὶ τὰ ὑποθέματα αὐτῆς ἐκ χρυσίου καθαροῦ ποιήσεις.
\vs{39}Πάντα τὰ σκεύη ταῦτα τάλαντον χρυσίου καθαροῦ.
\vs{40}Ὅρα, ποιήσεις κατὰ τὸν τύπον τὸν δεδειγμένον σοι ἐν τῷ ὄρει.

\ch{26}
Καὶ τὴν σκηνὴν ποιήσεις, δέκα αὐλαίας ἐκ βύσσου κεκλωσμένης, καὶ ὑακίνθου, καὶ πορφύρας, καὶ κοκκίνου κεκλωσμένου χερουβὶμ· ἐργασίᾳ ὑφάντου ποιήσεις αὐτάς.
\vs{2}Μῆκος τῆς αὐλαίας τῆς μιᾶς ὀκτὼ καὶ εἴκοσι πήχεων, καὶ εὖρος τεσσάρων πήχεων ἡ αὐλαία ἡ μία ἔσται· μέτρον τὸ αὐτὸ ἔσται πάσαις ταῖς αὐλαίαις.
\vs{3}Πέντε δὲ αὐλαῖαι ἔσονται ἐξ ἀλλήλων ἐχόμεναι ἡ ἑτέρα ἐκ τῆς ἑτέρας· καὶ πέντε αὐλαῖαι ἔσονται συνεχόμεναι ἑτέρα τῇ ἑτέρᾳ.
\vs{4}Καὶ ποιήσεις αὐταῖς ἀγκύλας ὑακινθίνας ἐπὶ τοῦ χείλους τῆς αὐλαίας τῆς μιᾶς, ἐκ τοῦ ἑνὸς μέρους εἰς τὴν συμβολήν· καὶ οὕτω ποιήσεις ἐπὶ τοῦ χείλους τῆς αὐλαίας τῆς ἐξωτέρας πρὸς τῇ συμβολῇ τῇ δευτέρᾳ.
\vs{5}Πεντήκοντα ἀγκύλας ποιήσεις τῇ αὐλαίᾳ τῇ μιᾷ, καὶ πεντήκοντα ἀγκύλας ποιήσεις ἐκ τοῦ μέρους τῆς αὐλαίας κατὰ τὴν συμβολὴν τῆς δευτέρας, ἀντιπρόσωποι ἀντιπίπτουσαι ἀλλήλαις εἰς ἑκάστην.
\vs{6}Καὶ ποιήσεις κρίκους πεντήκοντα χρυσοῦς· καὶ συνάψεις τὰς αὐλαίας ἑτέραν τῇ ἑτέρα τοῖς κρίκοις· καὶ ἔσται ἡ σκηνὴ μία.
\vs{7}Καὶ ποιήσεις δέῤῥεις τριχίνας σκέπην ἐπὶ τῆς σκηνῆς, ἕνδεκα δέῤῥεις ποιήσεις αὐτάς.
\vs{8}Τὸ μῆκος τῆς δέῤῥεως τῆς μιᾶς, τριάκοντα πήχεων, καὶ τεσσάρων πήχεων τὸ εὖρος τῆς δέῤῥεως τῆς μιᾶς· τὸ αὐτὸ μέτρον ἔσται ταῖς ἕνδεκα δέῤῥεσι.
\vs{9}Καὶ συνάψεις τὰς πέντε δέῤῥεις ἐπὶ τὸ αὐτὸ, καὶ τὰς ἓξ δέῤῥεις ἐπὶ τὸ αὐτό· καὶ ἐπιδιπλώσεις τὴν δέῤῥιν τὴν ἕκτην κατὰ πρόσωπον τῆς σκηνῆς.
\vs{10}Καὶ ποιήσεις ἀγκύλας πεντήκοντα ἐπὶ τοῦ χείλους τῆς δέῤῥεως τῆς μιᾶς, τῆς ἀναμέσον κατὰ συμβολήν· καὶ πεντήκοντα ἀγκύλας ποιήσεις ἐπὶ τοῦ χείλους τῆς δέῤῥεως, τῆς συναπτούσης τῆς δευτέρας.

\vs{11}Καὶ ποιήσεις κρίκους χαλκοῦς πεντήκοντα· καὶ συνάψεις τοὺς κρίκους ἐκ τῶν ἀγκυλῶν, καὶ συνάψεις τὰς δέῤῥεις, καὶ ἔσται ἕν.
\vs{12}Καὶ ὑποθήσεις τὸ πλεονάζον ἐν ταῖς δέῤῥεσι τῆς σκηνῆς· τὸ ἥμισυ τῆς δέῤῥεως τὸ ὑπολελειμμένον ὑποκαλύψεις εἰς τὸ πλεονάζον τῶν δέῤῥεων τῆς σκηνῆς, ὑποκαλύψεις ὀπίσω τῆς σκηνῆς.
\vs{13}Πῆχυν ἐκ τούτου, καὶ πῆχυν ἐκ τούτου, ἐκ τοῦ ὑπερέχοντος τῶν δέῤῥεων, ἐκ τοῦ μήκους τῶν δέῤῥεων τῆς σκηνῆς· ἔσται συγκαλύπτον ἐπὶ τὰ πλάγια τῆς σκηνῆς ἔνθεν καὶ ἔνθεν, ἵνα καλύπτῃ.
\vs{14}Καὶ ποιήσεις κατακάλυμμα τῇ σκηνῇ δέρματα κριῶν ἠρυθροδανωμένα, καὶ ἐπικαλύμματα δέρματα ὑακίνθινα ἐπάνωθεν.

\vs{15}Καὶ ποιήσεις στύλους τῆς σκηνῆς ἐκ ξύλων ἀσήπτων.
\vs{16}Δέκα πήχεων ποιήσεις τὸν στύλον τὸν ἕνα, καὶ πήχεως ἑνὸς καὶ ἡμίσους τὸ πλάτος τοῦ στύλου τοῦ ἑνός.
\vs{17}Δύο ἀγκωνίσκους τῷ στύλῳ τῷ ἑνὶ, ἀντιπίπτοντας ἕτερον τῷ ἑτέρῳ· οὕτω ποιήσεις πᾶσι τοῖς στύλοις τῆς σκηνῆς.
\vs{18}Καὶ ποιήσεις στύλους τῇ σκηνῇ, εἴκοσι στύλους ἐκ τοῦ κλίτους τοῦ πρὸς Βοῤῥᾶν.
\vs{19}Καὶ τεσσαράκοντα βάσεις ἀργυρᾶς ποιήσεις τοῖς εἴκοσι στύλοις· δύο βάσεις τῷ στύλῳ τῷ ἑνὶ εἰς ἀμφότερα τὰ μέρη αὐτοῦ· και δύο βάσεις τῷ στύλῳ τῷ ἑνὶ εἰς ἀμφοτέρα τὰ μέρη αὐτοῦ.
\vs{20}Καὶ τὸ κλίτος τὸ δεύτερον τὸ πρὸς Νότον, εἴκοσι στύλους,
\vs{21}καὶ τεσσαράκοντα βάσεις αὐτῶν ἀργυρᾶς· δύο βάσεις τῷ στύλῳ τῷ ἑνὶ εἰς ἀμφότερα τὰ μέρη αὐτοῦ, καὶ δύο βάσεις τῷ στύλῳ τῷ ἑνὶ εἰς ἀμφότερα τὰ μέρη αὐτοῦ.
\vs{22}Καὶ ἐκ τῶν ὀπίσω τῆς σκηνῆς κατὰ τὸ μέρος τὸ πρὸς θάλασσαν ποιήσεις ἓξ στύλους.
\vs{23}Καὶ δύο στύλους ποιήσεις ἐπὶ τῶν γωνιῶν τῆς σκηνῆς ἐκ τῶν ὀπισθίων.
\vs{24}Καὶ ἔσται ἐξ ἴσου κάτωθεν· κατὰ τὸ αὐτὸ ἔσονται ἴσοι ἐκ τῶν κεφαλῶν εἰς σύμβλησιν μίαν· οὕτω ποιήσεῖς ἀμφοτέραις ταῖς δυσὶ γωνίαις· ἴσαι ἔστωσαν.
\vs{25}Καὶ ἔσονται ὀκτὼ στύλοι, καὶ αἱ βάσεις αὐτῶν ἀργυραῖ δεκαέξ· δύο βάσεις τῷ ἑνὶ στύλῳ εἰς ἀμφότερα τὰ μέρη αὐτοῦ, καὶ δύο βάσεις τῷ στύλῳ τῷ ἑνί.
\vs{26}Καὶ ποιήσεις μοχλοὺς ἐκ ξύλων ἀσήπτων· πέντε τῷ ἑνὶ στύλῳ ἐκ τοῦ ἑνὸς μέρους τῆς σκηνῆς,
\vs{27}καὶ πέντε μοχλοὺς τῷ στύλῳ τῷ ἑνὶ κλίτει τῆς σκηνῆς τῷ δευτέρῳ, καὶ πέντε μοχλοὺς τῷ στύλῳ τῷ ὀπισθίῳ τῷ κλίτει τῆς σκηνῆς τῷ πρὸς θάλασσαν.
\vs{28}Καὶ ὁ μοχλὸς ὁ μέσος ἀναμέσον τῶν στύλων διϊκνείσθω ἀπὸ τοῦ ἑνὸς κλίτους εἰς τὸ ἕτερον κλίτος.
\vs{29}Καὶ τοὺς στύλους καταχρυσώσεις χρυσίῳ· καὶ τοὺς δακτυλίους ποιήσεις χρυσοῦς, εἰς οὓς εἰσάξεις τούς μοχλούς· καὶ καταχρυσώσεις τοὺς μοχλοὺς χρυσίῳ.
\vs{30}Καὶ ἀναστήσεις τὴν σκηνὴν κατὰ τὸ εἶδος τὸ δεδειγμένον σοι ἐν τῷ ὄρει.

\vs{31}Καὶ ποιήσεις καταπέτασμα ἐξ ὑακίνθου, καὶ πορφύρας, καὶ κοκκίνου κεκλωσμένου, καὶ βύσσου νενησμένης· ἔργον ὑφαντὸν ποιήσεις αὐτὸ χερουβίμ.
\vs{32}Καὶ ἐπιθήσεις αὐτὸ ἐπὶ τεσσάρων στύλων ἀσήπτων κεχρυσωμένων χρυσίῳ· καὶ αἱ κεφαλίδες αὐτῶν χρυσαῖ, καὶ αἱ βάσεις αὐτῶν τέσσαρες ἀργυραῖ.
\vs{33}Καὶ θήσεις τὸ καταπέτασμα ἐπὶ τῶν στύλων· καὶ εἰσοίσεις ἐκεῖ ἐσώτερον τοῦ καταπετάσματος τὴν κιβωτὸν τοῦ μαρτυρίου· καὶ διοριεῖ τὸ καταπέτασμα ὑμῖν ἀναμέσον τοῦ ἁγίου καὶ ἀναμέσον τοῦ ἁγίου τῶν ἁγίων.
\vs{34}Καὶ κατακαλύψεις τῷ καταπετάσματι τὴν κιβωτὸν τοῦ μαρτυρίου ἐν τῷ ἁγίῳ τῶν ἁγίων.
\vs{35}Καὶ ἐπιθήσεις τὴν τράπεζαν ἔξωθεν τοῦ καταπετάσματος, καὶ τὴν λυχνίαν ἀπέναντι τῆς τραπέζης ἐπὶ μέρους τῆς σκηνῆς τὸ πρὸς Νότον· καὶ τὴν τράπεζαν θήσεις ἐπὶ μέρους τῆς σκηνῆς τὸ πρὸς Βοῤῥᾶν.
\vs{36}Καὶ ποιήσεις ἐπίσπαστρον τῇ θύρᾳ τῆς σκηνῆς ἐξ ὑακίνθου, καὶ πορφύρας, καὶ κοκκίνου κεκλωσμένου, καὶ βύσσου κεκλωσμένης, ἔργον ποικιλτοῦ.
\vs{37}Καὶ ποιήσεις τῷ καταπετάσματι πέντε στύλους, καὶ χρυσώσεις αὐτοὺς χρυσίῳ· καὶ αἱ κεφαλίδες αὐτῶν χρυσαῖ· καὶ χωνεύσεις αὐτοῖς πέντε βάσεις χαλκᾶς.

\ch{27}
Καὶ ποιήσεις θυσιαστήριον ἐκ ξύλων ἀσήπτων, πέντε πήχεων τὸ μῆκος, καὶ πέντε πήχεων τὸ εὖρος· τετράγωνον ἔσται τὸ θυσιαστήριον, καὶ τριῶν πήχεων τὸ ὕψος αὐτοῦ.
\vs{2}Καὶ ποιήσεις τὰ κέρατα ἐπὶ τῶν τεσσάρων γωνιῶν· ἐξ αὐτοῦ ἔσται τὰ κέρατα, καὶ καλύψεις αὐτὰ χαλκῷ.
\vs{3}Καὶ ποιήσεις στεφάνην τῷ θυσιαστηρίῳ· καὶ τὸν καλυπτῆρα αὐτοῦ, καὶ τὰς φιάλας αὐτοῦ, καὶ τὰς κρεάγρας αὐτοῦ, καὶ τὸ πυρεῖον αὐτοῦ, καὶ πάντα τὰ σκεύη αὐτοῦ ποιήσεις χαλκᾶ.
\vs{4}Καὶ ποιήσεις αὐτῷ ἐσχάραν ἔργῳ δικτυωτῷ χαλκῆν· καὶ ποιήσεις τῇ ἐσχάρᾳ τέσσαρες δακτυλίους χαλκοῦς ὑπὸ τὰ τέσσαρα κλίτη.
\vs{5}Καὶ ὑποθήσεις αὐτοὺς ὑπὸ τὴν ἐσχάραν τοῦ θυσιαστήριου κάτωθεν· ἔσται δὲ ἡ ἐσχάρα ἕως τοῦ ἡμίσους τοῦ θυσιαστηρίου.
\vs{6}Καὶ ποιήσεις τῷ θυσιαστηρίῳ ἀναφορεῖς ἐκ ξύλων ἀσήπτων, καὶ περιχαλκώσεις αὐτοὺς χαλκῷ.
\vs{7}Καὶ εἰσάξεις τοὺς ἀναφορεῖς εἰς τοὺς δακτυλίους· καὶ ἔστωσαν ἀναφορεῖς κατὰ πλευρὰ τοῦ θυσιαστηρίου ἐν τῷ αἴρειν αὐτό.
\vs{8}Κοῖλον συνιδωτὸν ποιήσεις αὐτό· κατὰ τὸ παραδειχθέν σοι ἐν τῷ ὄρει, οὕτω ποιήσεις αὐτό.
\vs{9}Καὶ ποιήσεις αὐλὴν τῇ σκηνῇ· εἰς τὸ κλίτος τὸ πρὸς Λίβα ἱστία τῆς αὐλῆς ἐκ βύσσου κεκλωσμένης· μῆκος ἑκατὸν πήχεων τῷ ἑνὶ κλίτει.
\vs{10}Καὶ οἱ στύλοι αὐτῶν εἴκοσι, καὶ αἱ βάσεις αὐτῶν εἴκοσι χαλκαῖ, καὶ οἱ κρίκοι αὐτῶν καὶ αἱ ψαλίδες ἀργυραῖ.
\vs{11}Οὕτως τῷ κλίτει τῷ πρὸς ἀπηλιώτην ἱστία ἑκατὸν πήχεων μῆκος· καὶ οἱ στύλοι αὐτῶν εἴκοσι, καὶ αἱ βάσεις αὐτῶν εἴκοσι χαλκαῖ· καὶ οἱ κρίκοι καὶ αἱ ψαλίδες τῶν στύλων, καὶ αἱ βάσεις αὐτῶν περιηργυρωμέναι ἀργυρίῳ.
\vs{12}Τὸ δὲ εὖρος τῆς αὐλῆς τὸ κατὰ θάλασσαν ἱστία πεντήκοντα πήχεων· στύλοι αὐτῶν δέκα, καὶ βάσεις αὐτῶν δέκα.
\vs{13}Καὶ εὖρος τῆς αὐλῆς τῆς πρὸς Νότον ἱστία πεντήκοντα πήχεων· στύλοι αὐτῶν δέκα, καὶ βάσεις αὐτῶν δέκα.
\vs{14}Καὶ πεντεκαίδεκα πήχεων τὸ ὕψος τῶν ἱστίων τῷ κλίτει τῷ ἑνί· στύλοι αὐτῶν τρεῖς, καὶ αἱ βάσεις αὐτῶν τρεῖς.
\vs{15}Καὶ τὸ κλίτος τὸ δεύτερον δεκαπέντε πήχεων τῶν ἱστίων τὸ ὕψος· στύλοι αὐτῶν τρεῖς, καὶ αἱ βάσεις αὐτῶν τρεῖς.
\vs{16}Καὶ τῇ πύλῃ τῆς αὐλῆς κάλυμμα· εἴκοσι πήχεων τὸ ὕψος ἐξ ὑακίνθου, καὶ πορφύρας, καὶ κοκκίνου κεκλωσμένου, καὶ βύσσου κεκλωσμένης τῇ ποικιλίᾳ τοῦ ῥαφιδευτοῦ· στύλοι αὐτῶν τέσσαρες, καὶ αἱ βάσεις αὐτῶν τέσσαρες.
\vs{17}Πάντες οἱ στύλοι τῆς αὐλῆς κύκλῳ κατηργυρωμένοι ἀργυρίῳ, καὶ αἱ κεφαλίδες αὐτῶν ἀργυραῖ, καὶ αἱ βάσεις αὐτῶν χαλκαῖ.
\vs{18}Τὸ δὲ μῆκος τῆς αὐλῆς ἑκατὸν ἐφʼ ἑκατόν· καὶ εὖρος πεντήκοντα ἐπὶ πεντήκοντα· καὶ ὕψος πέντε πήχεῶν ἐκ βύσσου κεκλωσμένης, καὶ βάσεις αὐτῶν χαλκαῖ.
\vs{19}Καὶ πᾶσα ἡ κατασκευὴ καὶ πάντα τὰ ἐργαλεῖα καὶ οἱ πάσσαλοι τῆς αὐλῆς χαλκοῖ.

\vs{20}Καὶ σὺ σύνταξον τοῖς υἱοῖς Ἰσραὴλ, καὶ λαβέτωσάν σοι ἔλαιον ἐξ ἐλαιῶν ἀτρυγον καθαρὸν κεκομμένον εἰς φῶς καῦσαι, ἵνα καίηται λύχνος διαπαντός
\vs{21}ἐν τῇ σκηνῇ τοῦ μαρτυρίου· ἔξωθεν τοῦ καταπετάσματος τοῦ ἐπὶ τῆς διαθήκης καύσει αὐτὸ Ἀαρὼν καὶ οἱ υἱοὶ αὐτοῦ ἀφʼ ἑσπέρας ἕως πρωῒ, ἐναντίον Κυρίου, νόμιμον αἰώνιον εἰς τὰς γενεὰς ὑμῶν παρὰ τῶν υἱῶν Ἰσραήλ.

\ch{28}
Καὶ σὺ προσαγάγου πρὸς σεαυτὸν τόν τε Ἀαρὼν τὸν ἀδελφόν σου, καὶ τοὺς υἱοὺς αὐτοῦ, καὶ ἐκ τῶν υἱῶν Ἰσραὴλ, ἱερατεύειν μοι Ἀαρὼν, καὶ Ναδὰβ, καὶ Ἀβιοὺδ, καὶ Ἐλεάζαρ, καὶ Ἰθάμαρ, υἱοὺς Ἀαρών.
\vs{2}Καὶ ποιήσεις στολὴν ἁγίαν Ἀαρὼν τῷ ἀδελφῷ σου εἰς τιμὴν καὶ δόξαν.
\vs{3}Καὶ σύ λάλησον πᾶσι τοῖς σοφοῖς τῇ διανοίᾳ, οὓς ἐνέπλησα πνεύματος σοφίας καὶ αἰσθήσεως· καὶ ποιήσουσι τὴν στολὴν τὴν ἁγίαν Ἀαρὼν εἰς τὸ ἅγιον, ἐν ᾗ ἱερατεύσει μοι.
\vs{4}Καὶ αὗται αἱ στολαὶ, ἃς ποιησουσι· τὸ περιστήθιον, καὶ τὴν ἐπωμίδα, καὶ τὸν ποδήρη, καὶ χιτῶνα κοσυμβωτὸν, καὶ κίδαριν, καὶ ζώνην· καὶ ποιήσουσι στολὰς ἁγίας Ἀαρὼν καὶ τοῖς υἱοῖς αὐτοῦ εἰς τὸ ἱερατεύειν μοι.
\vs{5}Καὶ αὐτοὶ λήψονται τὸ χρυσίον, καὶ τὸν ὑάκινθον, καὶ τὴν πορφύραν, καὶ τὸ κόκκινον, καὶ τὴν βύσσον.
\vs{6}Καὶ ποιήσουσι τὴν ἐπωμίδα ἐκ βύσσου κεκλωσμένης, ἔργον ὑφαντὸν ποικιλτοῦ.
\vs{7}Δύο ἐπωμίδες συνέχουσαι ἔσονται αὐτῷ ἑτέρα τὴν ἑτέραν, ἐπὶ τοῖς δυσὶ μέρεσιν ἐξηρτισμέναι.
\vs{8}Καὶ τὸ ὕφασμα τῶν ἐπωμίδων ὅ ἐστιν ἐπʼ αὐτῷ, κατὰ τὴν ποίησιν ἐξ αὐτοῦ ἔσται ἐκ χρυσίου καθαροῦ, καὶ ὑακίνθου, καὶ πορφύρας, καὶ κοκκίνου διανενησμένου, καὶ βύσσου κεκλωσμένης.
\vs{9}Καὶ λήψῃ τοὺς δύο λίθους, λίθους σμαράγδου, καὶ γλύψεις ἐν αὐτοῖς τὰ ὀνόματα τῶν υἱῶν Ἰσραήλ.
\vs{10}Ἓξ ὀνόματα ἐπὶ τὸν λίθον τὸν ἕνα, καὶ τὰ ἓξ ὀνόματα τὰ λοιπὰ ἐπὶ τὸν λίθον τὸν δεύτερον κατὰ τὰς γενέσεις αὐτῶν.
\vs{11}Ἔργον λιθουργικῆς τέχνης· γλύμμα σφραγίδος διαγλύψεις τοὺς δύο λίθους ἐπὶ τοῖς ὀνόμασι τῶς υἱῶν Ἰσρσήλ.
\vs{12}Καὶ θήσεις τοὺς δύο λίθους ἐπὶ τῶς ὤμων τῆς ἐπωμίδος· λίθοι μνημοσύνου εἰσὶ τοῖς υἱοῖς Ἰσραήλ· καὶ ἀναλήψεται Ἀαρὼν τὰ ὀνόματα τῶν υἱῶν Ἰσραὴλ ἔναντι Κυρίου ἐπὶ τῶν δύο ὤμων αὐτοῦ, μνημόσυνον πεπὶ αὐτῶν.
\vs{13}Καὶ ποιήσεις ἀσπιδίσκας ἐκ χρυσίου καθαροῦ.
\vs{14}Καὶ ποιήσεις δύο κροσωτὰ ἐκ χρυσίου καθαροῦ, καταμεμιγμένα ἐν ἄνθεσιν, ἔργον πλοκῆς· καὶ ἐπιθήσεις τὰ κροσσωτὰ τὰ πεπλεγμένα ἐπὶ τὰς ἀσπιδίσκας, κατὰ τὰς παρωμίδας αὐτῶν ἐκ τῶν ἐμπροσθίων.

\vs{15}Καὶ ποιήσεις λογεῖον τῶν κρίσεων, ἔργον ποικιλτοῦ· κατὰ τὸν ῥυθμὸν τῆς ἐπωμίδος ποιήσεις αὐτὸ ἐκ χρυσίου, καὶ ὑακίνθου, καὶ πορφύρας, καὶ κοκκίνου κεκλωσμένου, καὶ βύσσου κεκλωσμένης.
\vs{16}Ποιήσεις αὐτό τετράγωνον· ἔσται διπλοῦν, σπιθαμῆς τὸ μῆκος αὐτοῦ, καὶ σπιθαμῆς τὸ εὖρος.
\vs{17}Καὶ καθυφανεῖς ἐν αὐτῷ ὕφασμα κατάλιθον τετράστιχον· στίχος λίθων ἔσται, σάρδιον, τοπάζιον, καὶ σμαράγδος, ὁ στίχος ὁ εἷς.
\vs{18}Καὶ ὁ στίχος ὁ δεύτερος, ἄνθραξ, καὶ σάπφειρος, καὶ ἴασπις.
\vs{19}Καὶ ὁ στίχος ὁ τρίτος, λιγύριον, ἀχάτης, ἀμέθυστος.
\vs{20}Καὶ ὁ στίχος ὁ τέταρτος, χρυσόλιθος, καὶ βηρύλλιον, καὶ ὀνύχιον, περικεκαλυμμένα χρυσίῳ, συνδεδεμένα ἐν χρυσίῳ· ἔστωσαν κατὰ στίχον αὐτῶν.
\vs{21}Καὶ οἱ λίθοι ἔστωσαν ἐκ τῶν ὀνομάτων τῶν υἱῶν Ἰσραὴλ δεκαδύο κατὰ τὰ ὀνόματα αὐτῶν· γλυφαὶ σφραγίδων, ἕκαστος κατὰ τὸ ὄνομα ἔστωσαν εἰς δεκαδύο φυλάς.
\vs{22}Καὶ ποιήσεις ἐπὶ τὸ λογιον κρωσσοὺς συμπεπλεγμένους, ἔργον ἁλυσιδωτὸν ἐκ χρυσίου καθαροῦ.
\vs{29}Καὶ λήψεται Ἀαρὼν τὰ ὀνόματα τῶν υἱῶν Ἰσραὴλ ἐπὶ τοῦ λογείου τῆς κρίσεως ἐπὶ τοῦ στήθους, εἰσιόντι εἰς τὸ ἅγιον μνημόσυνου ἐναντίον τοῦ Θεοῦ.
\vs{29a}Καὶ θήσεις ἐπὶ τὸ λογεῖον τῆς κρίσεως τοὺς κρωσσούς· τὰ ἁλυσιδωτὰ ἐπʼ ἀμφοτέρων τῶν κλιτῶν τοῦ λογείου ἐπιθήσεις. Καὶ τὰς δύο ἀσπιδίσκας ἐπιθήσεις ἐπʼ ἀμφοτέρους τοὺς ὤμους τῆς ἐπωμίδος κατὰ πρόσωπον.
\vs{30}Καὶ ἐπιθήσεις ἐπὶ τὸ λογεῖον τῆς κρίσεως τὴν δήλωσιν καὶ τὴν ἀλήθειαν· καὶ ἔσται ἐπὶ τοῦ στήθους Ἀαρὼν, ὃταν εἰσπορεύεται εἰς τὸ ἅγιον ἔναντὶ Κυρίου· καὶ οἴσει Ἀαρὼν τὰς κρίσεις τῶν υἱῶν Ἰσραὴλ ἐπὶ τοῦ στήθους ἔναντι Κυρίου διαπαντός.
\vs{31}Καὶ ποιήσεις ὑποδύτην ποδήρη ὅλον ὑακίνθινον.
\vs{32}Καὶ ἔσται τὸ περιστόμιον ἐξ αὐτοῦ μέσον, ὤαν ἔχον κύκλῳ τοῦ περιστομίου, ἔργον ὑφαντου, τὴν συμβολὴν συνυφασμένην ἐξ αὐτοῦ, ἵνα μὴ ῥαγῇ.
\vs{33}Καὶ ποιήσεις ὑπὸ τὸ λῶμα τοῦ ὑποδύτου κάτωθεν, ὡσεὶ ἐξανθούσης ῥόας ῥοΐσκους ἐξ ὑακίνθου, καὶ πορφύρας, καὶ κοκκίνου διανενησμένου, καὶ βύσσου κεκλωσμένης, ὑπὸ τοῦ λώματος τοῦ ὑποδύτου κύκλῳ· τὸ αὐτὸ εἶδος ῥοΐσκους χρυσοῦς, καὶ κώδωνας ἀναμέσον τούτων περικύκλῳ.
\vs{34}Παρὰ ῥοΐσκον χρυσοῦν δώδωνα, καὶ ἄνθινον ἐπὶ τοῦ λώματος τοῦ ὑποδύτου κύκλῳ·
\vs{35}Καὶ ἔσται Ἀαρὼν ἐν τῷ λειτουργεῖν ἀκουστὴ ἡ φωνὴ αὐτοῦ, εἰσιόντι εἰς τὸ ἅγιον ἔναντι Κυρίου, καὶ ἐξιόντι, ἵνα μὴ ἀποθάνῃ.
\vs{36}Καὶ ποιήσεις πέταλον χρυσοῦν καθαρόν· καὶ ἐκτυπώσεις ἐν αὐτῷ ἐκτύπωμα σφραγίδος, Ἁγίασμα Κυρίου.
\vs{37}Καὶ ἐπιθήσεις αὐτὸ ἐπὶ ὑακίνθου κεκλωσμένης· καὶ ἔσται ἐπὶ τῆς μίτρας, κατὰ πρόσωπον τῆς μίτρας ἔσται.
\vs{38}Καὶ ἔσται ἐπὶ τοῦ μετώπου Ἀαρών· καὶ ἐξαρεῖ Ἀαρὼν τὰ ἁμαρτήματα τῶν ἁγίων, ὅσα ἂν ἁγιάσωσιν οἱ υἱοὶ Ἰσραὴλ παντὸς δόματος τῶν ἁγίων αὐτῶν· καὶ ἔσται ἐπὶ τοῦ μετώπου Ἀαρὼν διαπαντὸς δεκτὸν αὐτοῖς ἔναντι Κυρίου.

\vs{39}Καὶ οἱ κοσυμβωτοὶ τῶν χιτώνων ἐκ βύσσου· καὶ ποιήσεις κίδαριν βυσσίνην· καὶ ζώνην ποιήσεις, ἔργον ποικιλτοῦ.
\vs{40}Καὶ τοῖς υἱοῖς Ἀαρὼν ποιήσεις χιτῶνας καὶ ζώνας, καὶ κιδάρεις ποιήσεις αὐτοῖς εἰς τιμὴν καὶ δόξαν.
\vs{41}Καὶ ἐνδύσεις αὐτὰ Ἀαρὼν τὸν ἀδελφόν σου, καὶ τοὺς υἱοὺς αὐτοῦ μετʼ αὐτοῦ· καὶ χρίσεις αὐτοὺς, καὶ ἐμπλήσεις αὐτῶν τὰς χεῖρας· καὶ ἁγιάσεις αὐτοὺς, ἵνα ἱερατεύωσί μοι.
\vs{42}Καὶ ποιήσεις αὐτοῖς περισκελῆ λινᾶ καλύψαι ἀσχημοσύνην χρωτὸς αὐτῶν, ἀπὸ ὀσφύος ἕως μηρῶν ἔσται.
\vs{43}Καὶ ἕξει Ἀαρὼν αὐτὰ καὶ οἱ υἱοὶ αὐτοῦ, ὅταν εἰσπορεύωνται εἰς τὴν σκηνὴν τοῦ μαρτυρίου, ἢ ὅταν προσπορεύωνται λειτουργεῖν πρὸς τὸ θυσιαστήριον τοῦ ἁγίου· καὶ οὐκ ἐπάξονται πρὸς ἑαυτοὺς ἁμαρτίαν, ἵνα μὴ ἀποθάνωσι· νόμιμον αἰώνιον αὐτῷ, καὶ τῷ σπέρματι αὐτοῦ μετʼ αὐτόν.

\ch{29}
Καὶ ταῦτά ἐστιν, ἃ ποιήσεις αὐτοῖς· ἁγιάσεις αὐτοὺς, ὥστε ἱερατεύειν μοι αὐτούς· λήψῃ δὲ μοσχάριον ἐκ βοῶν ἓν, καὶ κριοὺς ἀμώμους δύο,
\vs{2}καὶ ἄρτους ἀζύμους πεφυραμένους ἑν ἐλαίῳ, καὶ λάγανα ἄζυμα κεχρισμένα ἐν ἐλαίῳ· σεμίδαλιν ἐκ πυρῶν ποιήσεις αὐτά.
\vs{3}Καὶ ἐπιθήσεις αὐτὰ ἐπὶ κανοῦν ἕν· καὶ προσοίσεις αὐτὰ ἐπὶ τῷ κανῷ· καὶ τὸ μοσχάριον, καὶ τοὺς δύο κριούς.
\vs{4}Καὶ Ἀαρὼν καὶ τοὺς υἱοὺς αὐτοῦ προσάξεις ἐπὶ τὰς θύρας τῆς σκηνῆς τοῦ μαρτυρίου, καὶ λούσεις αὐτοὺς ἐν ὕδατι.
\vs{5}Καὶ λαβὼν τὰς στολὰς, ἐνδύσεις Ἀαρὼν τὸν ἀδελφόν σου καὶ τὸν χιτῶνα τὸν ποδήρη, καὶ τὴν ἐπωμίδα, καὶ τὸ λογεῖον· καὶ συνάψεις αὐτῷ τὸ λογεῖον πρὸς τὴν ἐπωμίδα.
\vs{6}Καὶ ἐπιθήσεις τὴν μίτραν ἐπὶ τὴν κεφαλὴν αὐτοῦ, καὶ ἐπιθήσεις τὸ πέταλον τὸ ἁγίασμα ἐπὶ τὴν μίτραν.
\vs{7}Καὶ λήψῃ τοῦ ἐλαίου τοῦ χρίσματος· καὶ ἐπιχεεῖς αὐτὸ ἐπὶ τὴν κεφαλὴν αὐτοῦ, καὶ χρίσεις αὐτόν.
\vs{8}Καὶ τοὺς υἱοὺς αὐτοῦ προσάξεις, καὶ ἐνδύσεις αὐτοὺς χιτῶνας.
\vs{9}Καὶ ζώσεις αὐτοὺς ταῖς ζωναῖς, καὶ περιθήσεις αὐτοῖς τὰς κιδάρεις· καὶ ἔσται αὐτοῖς ἱερατῖα μοι εἰς τὸν αἰῶνα· καὶ τελειώσεις Ἀαρὼν τὰς χεῖρας αὐτοῦ, καὶ τὰς χεῖρας τῶν υἱῶν αὐτοῦ.
\vs{10}Καὶ προσάξεις τὸν μόσχον ἐπὶ τὰς θύρας τῆς σκηνῆς τοῦ μαρτυρίου· καὶ ἐπιθήσουσιν Ἀαρὼν καὶ οἱ υἱοὶ αὐτοῦ τὰς χεῖρας αὐτῶν ἐπὶ τὴν κεφαλὴν τοῦ μόσχου, ἔναντι Κυρίου, παρὰ τὰς θύρας τῆς σκηνῆς τοῦ μαρτυρίου.
\vs{11}Καὶ σφάξεις τὸν μόσχον ἔναντι Κυρίου, παρὰ τὰς θύρας τῆς σκηνῆς τοῦ μαρτυρίου.
\vs{12}Καὶ λήψῃ ἀπὸ τοῦ αἵματος τοῦ μόσχου, καὶ θήσεις ἐπὶ τῶν κεράτων τοῦ θυσιαστηρίου τῷ δακτύλῳ σου· τὸ δὲ λοιπὸν πᾶν αἷμα ἐκχεεῖς παρὰ τὴν βάσιν τοῦ θυσιαστηρίου.
\vs{13}Καὶ λήψῃ πᾶν τὸ στέαρ τὸ ἐπὶ τῆς κοιλίας, καὶ τὸν λοβὸν τοῦ ἥπατος, καὶ τοὺς δύο νεφροὺς, καὶ τὸ στέαρ τὸ ἐπʼ αὐτῶν, καὶ ἐπιθήσεις ἐπὶ τὸ θυσιαστήριον.
\vs{14}Τὰ δὲ κρέατα τοῦ μόσχου, καὶ τὸ δέρμα, καὶ τὴν κόπρον κατακαύσεις πυρὶ ἔξω τῆς παρεμβολῆς· ἁμαρτίας γάρ ἐστι.

\vs{15}Καὶ τὸν κριὸν λήψῃ τὸν ἕνα, καὶ ἐπιθήσουσιν Ἀαρὼν καὶ οἱ υἱοὶ αὐτοῦ τὰς χεῖρας αὐτῶν ἐπὶ τὴν κεφαλὴν τοῦ κριοῦ.
\vs{16}Καὶ σφάξεις αὐτὸν, καὶ λαβὼν τὸ αἷμα προσχεεῖς πρὸς τὸ θυσιαστήριον κύκλῳ.
\vs{17}Καὶ τὸν κριὸν διχοτομήσεις κατὰ μέλη· καὶ πλυνεῖς τὰ ἐνδόσθια καὶ τοὺς πόδας ὕδατι, καὶ ἐπιθήσεις ἐπὶ τὰ διχοτομήματα σὺν τῇ κεφαλῇ.
\vs{18}Καὶ ἀνοίσεις ὅλον τὸν κριὸν ἐπὶ τὸ θυσιαστήριον, ὁλοκαύτωμα τῷ Κυρίῳ εἰς ὀσμὴν εὐωδίας· θυμίαμα Κυρίῳ ἐστί.
\vs{19}Καὶ λήψῃ τὸν κριὸν τὸν δεύτερον, καὶ ἐπιθήσει Ἀαρὼν καὶ οἱ υἱοὶ αὐτοῦ τὰς χεῖρας αὐτῶν ἐπὶ τὴν κεφαλὴν τοῦ κριοῦ.
\vs{20}Καὶ σφάξεις αὐτὸν, καὶ λήψῃ τοῦ αἵματος αὐτοῦ, καὶ ἐπιθήσεις ἐπὶ τὸν λοβὸν τοῦ ὠτὸς Ἀαρὼν τοῦ δεξιοῦ, καὶ ἐπὶ τὸ ἄκρον τῆς δεξιᾶς χειρὸς, καὶ ἐπὶ τὸ ἄκρον τοῦ ποδὸς τοῦ δεξιοῦ, καὶ ἐπὶ τοὺς λοβοὺς τῶν ὤτων τῶν υἱῶν αὐτοῦ τῶν δεξιῶν, καὶ ἐπὶ τὰ ἄκρα τῶν χειρῶν αὐτῶν τῶν δεξιῶν, καὶ ἐπὶ τὰ ἄκρα τῶν ποδῶν αὐτῶν τῶν δεξιῶν.
\vs{21}Καὶ λήψῃ ἀπὸ τοῦ αἵματος τοῦ ἀπὸ τοῦ θυσιαστηρίου, καὶ ἀπὸ τοῦ ἐλαίου τῆς χρίσεως, καὶ ῥανεῖς ἐπὶ Ἀαρὼν καὶ ἐπὶ τὴν στολὴν αὐτοῦ, καὶ ἐπὶ τοὺς υἱοὺς αὐτοῦ καὶ ἐπὶ τὰς στολὰς τῶν υἱῶν αὐτοῦ μετʼ αὐτοῦ· καὶ ἁγιασθήσεται αὐτὸς καὶ ἡ στολὴ αὐτοῦ, καὶ οἱ υἱοὶ αὐτοῦ καὶ αἱ στολαὶ τῶν υἱῶν αὐτοῦ μετʼ αὐτοῦ· τὸ δὲ αἷμα τοῦ κριοῦ προσχεεῖς πρὸς τὸ θυσιαστήριον κύκλῳ.
\vs{22}Καὶ λήψῃ ἀπὸ τοῦ κριοῦ τὸ στέαρ αὐτοῦ, καὶ τὸ στέαρ τὸ κατακαλύπτον τὴν κοιλίαν, καὶ τὸν λοβὸν τοῦ ἥπατος, καὶ τοὺς δύο νεφροὺς, καὶ τὸ στέαρ τὸ ἐπʼ αὐτῶν, καὶ τὸν βραχίονα τὸν δεξιόν· ἔστι γὰρ τελείωσις αὕτη.
\vs{23}Καὶ ἄρτον ἕνα ἐξ ἐλαίου, καὶ λάγανον ἓν ἀπὸ τοῦ κανοῦ τῶν ἀζύμων τῶν προτεθειμένων ἔναντι Κυρίου.
\vs{24}Καὶ ἐπιθήσεις τὰ πάντα ἐπὶ τὰς χεῖρας Ἀαρὼν, καὶ ἐπὶ τὰς χεῖρας τῶν υἱῶν αὐτοῦ· καὶ ἀφοριεῖς αὐτὰ ἀφόρισμα ἔναντι Κυρίου.
\vs{25}Καὶ λήψῃ αὐτὰ ἐκ τῶν χειρῶν αὐτῶν, καὶ ἀνοίσεις ἐπὶ τὸ θυσιαστήριον τῆς ὁλοκαυτώσεως εἰς ὀσμὴν εὐωδίας ἔναντι Κύριου· κάρπωμά ἐστι Κυρίῳ.
\vs{26}Καὶ λήψῃ τὸ στηθύνιον ἀπὸ τοῦ κριοῦ τῆς τελειώσεως, ὅ ἐστιν Ἀαρών· καὶ ἀφοριεῖς αὐτὸ ἀφόρισμα ἔναντι Κυρίου· καὶ ἔσται σοι ἐν μερίδι.
\vs{27}Καὶ ἁγιάσεις τὸ στηθύνιον ἀφόρισμα, καὶ τὸν βραχίονα τοῦ ἀφαιρέματος, ὃς ἀφώρισται, καὶ ὃς ἀφῄρηται ἀπὸ τοῦ κριοῦ τῆς τελειώσεως ἀπὸ τοῦ Ἀαρὼν, καὶ ἀπὸ τῶν υἱῶν αὐτοῦ.
\vs{28}Καὶ ἔσται Ἀαρὼν καὶ τοῖς υἱοῖς αὐτοῦ νόμιμον αἰώνιον παρὰ τῶν υἱῶν Ἰσραήλ· ἔστι γὰρ ἀφόρισμα τοῦτο· καὶ ἀφαίρεμα ἕσται παρὰ τῶν υἱῶν Ἰσραὴλ ἀπὸ τῶν θυμάτων τῶν σωτηρίων τῶν υἱῶν Ἰσραὴλ, ἀφαίρεμα Κυρίῳ.

\vs{29}Καὶ ἡ στολὴ τοῦ ἁγίου, ἥ ἐστιν Ἀαρὼν, ἔσται τοῖς υἱοῖς αὐτοῦ μετʼ αὐτὸν, χρισθῆναι αὐτοὺς ἐν αὐτοῖς, καὶ τελειῶσαι τὰς χεῖρας αὐτῶν.
\vs{30}Ἑπτὰ ἡμέρας ἐνδύσεται αὐτὰ ὁ ἱερεὺς ὁ ἀντʼ αὐτοῦ ἐκ τῶν υἱῶν αὐτοῦ, ὃς εἰσελεύσεται εἰς τὴν σκηνὴν τοῦ μαρτυρίου λειτουργεῖν ἐν τοῖς ἁγίοις.
\vs{31}Καὶ τὸν κριὸν τῆς τελειώσεως λήψῃ· καὶ ἑψήσεις τὰ κρέα ἐν τόπῳ ἁγίῳ.
\vs{32}Καὶ ἔδονται Ἀαρὼν καὶ οἱ υἱοὶ αὐτοῦ τὰ κρέα τοῦ κριοῦ, καὶ τοὺς ἄρτους τοὺς ἐν τῷ κανῷ, παρὰ τὰς θύρας τῆς σκηνῆς τοῦ μαρτυρίου.
\vs{33}Ἔδονται αὐτὰ ἐν οἷς ἡγιάσθησαν ἐν αὐτοῖς τελειῶσαι τὰς χεῖρας αὐτῶν, ἁγιάσαι αὐτούς· καὶ ἀλλογενὴς οὐκ ἔδεται ἀπʼ αὐτῶν· ἔστι γὰρ ἅγια.
\vs{34}Ἐὰν δὲ καταλειφθῇ ἀπὸ τῶν κρεῶν τῆς θυσίας τῆς τελειώσεως καὶ τῶν ἄρτων ἕως πρωῒ, κατακαύσεις τὰ λοιπὰ πυρί· οὐ βρωθήσεται· ἁγίασμα γάρ ἐστι.

\vs{35}Καὶ ποιήσεις Ἀαρὼν καὶ τοῖς υἱοῖς αὐτοῦ οὕτω κατὰ πάντα ὅσα ἐνετειλάμην σοι· ἑπτὰ ἡμέρας τελειώσεις τὰς χεῖρας αὐτῶν.
\vs{36}Καὶ τὸ μοσχάριον τῆς ἁμαρτίας ποιήσεις τῇ ἡμέρᾳ τοῦ καθαρισμοῦ· καὶ καθαριεῖς τὸ θυσιαστήριον ἐν τῷ ἁγιάζειν σε ἐπʼ αὐτῷ· καὶ χρίσεις αὐτὸ ὥστε ἁγιάσαι αὐτό.
\vs{37}Ἑπτὰ ἡμέρας καθαριεῖς τὸ θυσιαστήριον, καὶ ἁγιάσεις αὐτό· καὶ ἔσται τὸ θυσιαστήριον, ἅγιον τοῦ ἁγίου· πᾶς ὁ ἁπτόμενος τοῦ θυσιαστηρίου, ἁγιασθήσεται.
\vs{38}Καὶ ταῦτά ἐστιν, ἃ ποιήσεις ἐπὶ τοῦ θυσιαστηρίου· ἀμνοὺς ἐνιαυσίους ἀμώμους δύο τὴν ἡμέραν ἐπὶ τὸ θυσιαστήριον ἐνδελεχῶς, κάρπωμα ἐνδελεχισμοῦ.

\vs{39}Τὸν ἀμνὸν τὸν ἕνα ποιήσεις τὸ πρωῒ, καὶ τὸν ἀμνὸν τὸν δεύτερον ποιήσεις τὸ δειλινόν.
\vs{40}Καὶ δέκατον σεμιδάλεως πεφυραμένης ἐν ἐλαίῳ κεκομμένῳ τῷ τετάρτῳ τοῦ εἴν· καὶ σπονδὴν τὸ τέταρτον τοῦ εἲν οἴνου τῷ ἀμνῷ τῷ ἑνί.
\vs{41}Καὶ τὸν ἀμνὸν τὸν δεύτερον ποιήσεις τὸ δειλινὸν, κατὰ τὴν θυσίαν τὴν πρωϊνὴν, καὶ κατὰ τὴν σπονδὴν αὐτοῦ· ποιήσεις εἰς ὀσμὴν εὐωδίας κάρπωμα Κυρίῳ,
\vs{42}θυσίαν ἐνδελεχισμοῦ εἰς γενεὰς ὑμῶν, ἐπὶ θύρας τῆς σκηνῆς τοῦ μαρτυρίου ἔναντι Κυρίου, ἐν οἷς γνωσθήσομαί σοι ἐκεῖθεν, ὥστε λαλῆσαί σοι.
\vs{43}Καὶ τάξομαι ἐκεῖ τοῖς υἱοῖς Ἰσραὴλ, καὶ ἁγιασθήσομαι ἐν δόξῃ μου.
\vs{44}Καὶ ἁγιάσω τὴν σκηνὴν τοῦ μαρτυρίου, καὶ τὸ θυσιαστήριον· καὶ Ἀαρὼν καὶ τοὺς υἱοὺς αὐτοῦ ἁγιάσω, ἱερατεύειν μοι.
\vs{45}Καὶ ἐπικληθήσομαι ἐν τοῖς υἱοῖς Ἰσραὴλ, καὶ ἔσομαι αὐτῶν Θεός.
\vs{46}Καὶ γνώσονται, ὅτι ἐγώ εἰμι Κύριος ὁ Θεὸς αὐτῶν, ὁ ἐξαγαγὼν αὐτοὺς ἐκ γῆς Αἰγύπτου, ἐπικληθῆναι αὐτοῖς, καὶ εἶναι αὐτῶν Θεός.

\ch{30}
Καὶ ποιήσεις θυσιαστήριον θυμιάματος ἐκ ξύλων ἀσήπτων.
\vs{2}Καὶ ποιήσεις αὐτὸ πήχεος τὸ μῆκος, καὶ πήχεος τὸ εὖρος· τετράγωνον ἔσται· καὶ δύο πήχεων τὸ ὕψος· ἐξ αὐτοῦ ἔσται τὰ κέρατα αὐτοῦ.
\vs{3}Καὶ καταχρυσώσεις χρυσίῳ καθαρῷ τὴν ἐσχάραν αὐτοῦ, καὶ τοὺς τοίχους αὐτοῦ κύκλῳ, καὶ τὰ κέρατα αὐτοῦ· καὶ ποιήσεις αὐτῷ στρεπτὴν στεφάνην χρυσῆν κύκλῳ.
\vs{4}Καὶ δύο δακτυλίους χρυσοῦς καθαροὺς ποιήσεις ὑπὸ τὴν στρεπτὴν στεφάνην αὐτοῦ, εἰς τὰ δύο κλίτη ποιήσεις ἐν τοῖς δυσὶ πλευροῖς· καὶ ἔσονται ψαλίδες ταῖς σκυτάλαις, ὥστε αἴρειν αὐτὸ ἐν αὐταῖς.
\vs{5}Καὶ ποιήσεις σκυτάλας ἐκ ξύλων ἀσήπτων, καὶ καταχρυσώσεις αὐτὰς χρυσίῳ.
\vs{6}Καὶ θήσεις αὐτὸ ἀπέναντι τοῦ καταπετάσματος, τοῦ ὄντος ἐπὶ τῆς κιβωτοῦ τῶν μαρτυρίων, ἐν οἷς γνωσθήσομαί σοι ἐκεῖθεν.
\vs{7}Καὶ θυμιάσει ἐπʼ αὐτοῦ Ἀαρὼν θυμίαμα σύνθετον λεπτὸν τὸ πρωῒ πρωΐ· ὅταν ἐπισκευάζῃ τοὺς λύχνους, θυμιάσει ἐπʼ αὐτοῦ.
\vs{8}Καὶ ὅταν ἐξάπτῃ Ἀαρὼν τοὺς λύχνους ὀψὲ, θυμιάσει ἐπʼ αὐτοῦ. θυμίαμα ἐνδελεχισμοῦ διαπαντὸς ἔναντι Κυρίου εἰς γενεὰς αὐτῶν.
\vs{9}Καὶ οὐκ ἀνοίσει ἐπʼ αὐτοῦ θυμίαμα ἕτερον· κάρπωμα, θυσίαν, και σπονδὴν οὐ σπείσεις ἐπʼ αὐτοῦ.
\vs{10}Καὶ ἐξιλάσεται ἐπʼ αὐτοῦ Ἀαρὼν ἐπὶ τῶν κεράτων αὐτοῦ ἅπαξ τοῦ ἐνιαυτοῦ· ἀπὸ τοῦ αἵματος τοῦ καθαρισμοῦ καθαριεῖ αὐτὸ εἰς γενεὰς αὐτῶν· ἅγιον τῶν ἁγίων ἐστὶ Κυρίῳ.

\vs{11}Καὶ ἑλάλησε Κύριος πρὸς Μωυσῆν, λέγων,
\vs{12}ἐὰν λάβῃς τὸν συλλογισμὸν τῶν υἱῶν Ἰσραὴλ ἐν τῇ ἐπισκοπῇ αὐτῶν, καὶ δώσουσιν ἕκαστος λύτρα τῆς ψυχῆς αὐτοῦ Κυρίῳ, καὶ οὐκ ἔσται ἐν αὐτοῖς πτῶσις ἐν τῇ ἐπισκοπῇ αὐτῶν.
\vs{13}Καὶ τοῦτό ἐστιν ὅ δώσουσιν ὅσοι ἂν παραπορεύωνται τὴν ἐπίσκεψιν· τὸ ἥμισυ τοῦ διδράχμου ὅ ἐστιν κατὰ τὸ δίδραχμον τὸ ἅγιον, εἴκοσι ὀβολοὶ τὸ δίδραχμον, τὸ δὲ ἥμισυ τοῦ διδράχμου εἰσφορὰ Κυρίῳ.
\vs{14}Πᾶς ὁ παραπορευόμενος εἰς τὴν ἐπίσκεψιν ἀπὸ εἰκοσαετοῦς καὶ ἐπάνω, δώσουσι τὴν εἰσφορὰν Κυρίῳ.
\vs{15}Ὁ πλουτῶν οὐ προσθήσει, καὶ ὁ πενόμενος οὐκ ἐλαττονήσει ἀπὸ τοῦ ἡμίσεως τοῦ διδράχμου ἐν τῷ διδόναι τὴν εἰσφορὰν Κυρίῳ, ἐξιλάσασθαι περὶ τῶν ψυχῶν ὑμῶν.
\vs{16}Καὶ λήψῃ τὸ ἀργύριον τῆς εἰσφορᾶς παρὰ τῶν υἱῶν Ἰσραήλ, καὶ δώσεις αὐτὸ εἰς τὸ κάτεργον τῆς σκηνῆς τοῦ μαρτυρίου· καὶ ἔσται τοῖς υἱοῖς Ἰσραὴλ μνημόσυνον ἔναντι Κυρίου, ἐξιλάσασθαι περὶ τῶν ψυχῶν ὑμῶν.
\vs{17}Καὶ ἐλάλησε Κύριος πρὸς Μωυσῆν, λέγων,
\vs{18}ποίησον λουτῆρα χαλκοῦν, καὶ βάσιν αὐτῷ χαλκῆν, ὥστε νίπτεσθαι· καὶ θήσεις αὐτὸν ἀνὰ μέσον τῆς σκηνῆς τοῦ μαρτυρίου καὶ ἀνὰ μέσον τοῦ θυσιαστηρίου· καὶ ἐκχεεῖς εἰς αὐτὸν ὕδωρ.
\vs{19}Καὶ νίψεται Ἀαρὼν καὶ οἱ υἱοὶ αὐτοῦ ἑξ αὐτοῦ τὰς χεῖρας, καὶ τοὺς πόδας ὕδατι.
\vs{20}Ὅταν εἰσπορεύωνται εἰς τὴν σκηνὴν τοῦ μαρτυρίου, νίψονται ὕδατι, καὶ οὐ μὴ ἀποθάνωσιν, ὅταν προσπορεύωνται πρὸς τὸ θυσιαστήριον λειτουργεῖν καὶ ἀναφέρειν τὰ ὁλοκαυτώματα Κυρίῳ.
\vs{21}Νίψονται τὰς χεῖρας καὶ τοὺς πόδας ὕδατι, ὅταν εἰσπορεύωνται εἰς τὴν σκηνὴν τοῦ μαρτυρίου, νίψονται ὕδατι, ἵνα μὴ ἀποθάνωσι· καὶ ἔσται αὐτοῖς νόμιμον αἰώνιον, αὐτῷ καὶ ταῖς γενεαῖς αὐτοῦ μετʼ αὐτόν.
\vs{22}Καὶ ἐλάλησε Κύριος πρὸς Μωυσῆν, λέγων,
\vs{23}καὶ σὺ λάβε ἡδύσματα, τὸ ἄνθος σμύρνης ἐκλεκτῆς πεντακοσίους σίκλους, καὶ κινναμώμου εὐώδους τὸ ἥμισυ τούτου διακοσίους πεντήκοντα, καὶ καλάμου εὐώδους διακοσίους πεντήκοντα,
\vs{24}καὶ ἴρεως πεντακοσίους σίκλους τοῦ ἁγίου, καὶ ἔλαιον ἐξ ἐλαιῶν εἵν.
\vs{25}Καὶ ποιήσεις αὐτὸ ἔλαιον χρίσμα ἅγιον, μύρον μυρεψικὸν τέχνῃ μυρεψοῦ· ἔλαιον χρίσμα ἅγιον ἔσται.
\vs{26}Καὶ χρίσεις ἐξ αὐτοῦ τὴν σκηνὴν τοῦ μαρτυρίου, καὶ τὴν κιβωτὸν τῆς σκηνῆς τοῦ μαρτυρίου, καὶ πάντα τὰ σκεύη αὐτῆς,
\vs{27}καὶ τὴν λυχνίαν καὶ πάντα τὰ σκεύη αὐτῆς, καὶ τὸ θυσιαστήριον τοῦ θυμιάματος,
\vs{28}καὶ τὸ θυσιαστήριον τῶν ὁλοκαυτωμάτων καὶ πάντα αὐτοῦ τὰ σκεύη, καὶ τὴν τράπεζαν καὶ πάντα τὰ σκεύη αὐτῆς, καὶ τὸν λουτῆρα.
\vs{29}Καὶ ἁγιάσεις αὐτά· καὶ ἔσται ἅγια τῶν ἁγίων· πᾶς ὁ ἁπτόμενος αὐτῶν, ἁγιασθήσεται.
\vs{30}Καὶ Ἀαρὼν καὶ τοὺς υἱοὺς αὐτοῦ χρίσεις, καὶ ἁγιάσεις αὐτοὺς ἱερατεύειν μοι.
\vs{31}Καὶ τοῖς υἱοῖς Ἰσραὴλ λαλήσεις, λέγων, ἔλαιον ἄλειμμα χρίσεως ἅγιον ἔσται τοῦτο ὑμῖν εἰς τὰς γενεὰς ὑμῶν.
\vs{32}Ἐπὶ σάρκα ἀνθρώπου οὐ χρισθήσεται· καὶ κατὰ τὴν σύνθεσιν ταύτην οὐ ποιήσετε ὑμῖν ἑαυτοῖς ὡσαύτως· ἅγιόν ἐστιν, καὶ ἁγίασμα ἔσται ὑμῖν.
\vs{33}Ὃς ἂν ποιήσῃ ὡσαύτως, καὶ ὃς ἂν δῷ ἀπʼ αὐτοῦ ἀλλογενεῖ, ἐξολοθρευθήσεται ἐκ τοῦ λαοῦ αὐτοῦ.

\vs{34}Καὶ εἶπε Κύριος πρὸς Μωυσῆν, λάβε σεαυτῷ ἡδύσματα, στακτήν, ὄνυχα, χαλβάνην ἡδυσμοῦ καὶ λίβανον διαφανῆ· ἴσον ἴσῳ ἔσται.
\vs{35}Καὶ ποιήσουσιν ἐν αὐτῷ θυμίαμα μυρεψικὸν ἔργον μυρεψοῦ μεμιγμένον, καθαρὸν ἔργον ἅγιον.
\vs{36}Καὶ συνκόψεις ἐκ τούτων λεπτόν, καὶ θήσεις ἀπέναντι τῶν μαρτυρίων ἐν τῇ σκηνῇ τοῦ μαρτυρίου, ὅθεν γνωσθήσομαί σοι ἐκεῖθεν· ἅγιον τῶν ἁγίων ἔσται ὑμῖν θυμίαμα.
\vs{37}Κατὰ τὴν σύνθεσιν ταύτην οὐ ποιήσετε ὑμῖν ἐαυτοῖς· ἁγίασμα ἔσται ὑμῖν Κυρίῳ·
\vs{38}Ὃς ἂν ποιήσῃ ὡσαύτως ὥστε ὀσφραίνεσθαι ἐν αὐτῷ, ἀπολεῖται ἐκ τοῦ λαοῦ αὐτοῦ.

\ch{31}
Καὶ ἐλάλησε Κύριος πρὸς Μωυσῆν, λέγων,
\vs{2}ἰδοὺ ἀνακέκλημαι ἐξ ὀνόματος τὸν Βεσελεὴλ τὸν τοῦ Οὐρείου τὸν Ὣρ, ἐκ τῆς φυλῆς Ἰούδα.
\vs{3}Καὶ ἐνέπλησα αὐτὸν πνεῦμα θεῖον σοφίας καὶ συνέσεως καὶ ἐπιστήμης, ἐν παντὶ ἔργῳ διανοεῖσθαι,
\vs{4}καὶ ἀρχιτεκτονῆσαι, ἐργάζεσθαι τὸ χρυσίον, καὶ τὸ ἀργύριον, καὶ τὸν χαλκὸν, καὶ τὴν ὑάκινθον, καὶ τὴν πορφύραν, καὶ τὸ κόκκινον τὸ νηστὸν,
\vs{5}καὶ τὰ λιθουργικὰ, καὶ εἰς τὰ ἔργα τὰ τεκτονικὰ τῶν ξύλων, ἐργάζεσθαι κατὰ πάντα τὰ ἔργα.
\vs{6}Καὶ ἐγὼ ἔδωκα αὐτὸν καὶ τὸν Ἐλιὰβ τὸν τοῦ Ἀχισαμὰχ ἐκ φυλῆς Δάν· καὶ παντὶ συνετῷ καρδίᾳ δέδωκα σύνεσιν·
\vs{7}καὶ πονήσουσι πάντα ὅσα συνέταξά σοι, τὴν σκηνὴν τοῦ μαρτυρίου, καὶ τὴν κιβωτὸν τῆς διαθήκης, καὶ τὸ ἱλαστήριον τὸ ἐπʼ αὐτῆς, καὶ τὴν διασκευὴν τῆς σκηνῆς,
\vs{8}καὶ τὰ θυσιαστήρια, καὶ τὴν τράπεζαν καὶ πάντα τὰ σκεύη αὐτῆς, καὶ τὴν λυχνίαν τὴν καθαρὰν καὶ πάντα τὰ σκεύη αὐτῆς
\vs{9}καὶ τὸν λουτῆρα καὶ τὴν βάσιν αὐτοῦ,
\vs{10}καὶ τὰς στολὰς τὰς λειτουργικὰς Ἀαρὼν, καὶ τὰς στολὰς τῶν υἱῶν αὐτοῦ ἱερατεύειν μοι,
\vs{11}καὶ τὸ ἔλαιον τῆς χρίσεως, καὶ τὸ θυμίαμα τῆς συνθέσεως τοῦ ἁγίου· κατὰ πάντα ὅσα ἐγὼ ἐνετειλάμην σοι, ποιήσουσι.

\vs{12}Καὶ ἐλάλησε Κύριος πρὸς Μωυσῆν, λέγων,
\vs{13}Καὶ σὺ σύνταξον τοῖς υἱοῖς Ἰσραὴλ, λέγων, Ὁρᾶτε, καὶ τὰ σάββατά μου φυλάξεσθε· σημεῖόν ἐστι παρʼ ἐμοὶ καὶ ἐν ὑμῖν εἰς τὰς γενεὰς ὑμῶν, ἵνα γνῶτε ὅτι ἐγὼ Κύριος ὁ ἁγιάξων ὑμᾶς.
\vs{14}καὶ φυλάξεσθε τὰ σάββατα, ὅτι ἅγιον τοῦτό ἐστι Κυρίῳ ὑμῖν· ὁ βεβηλῶν αὐτὸ, θανάτῳ θανατωθήσεται· πᾶς ὃς ποιήσει ἐν αὐτῷ ἔργον, ἐξολοθρευθήσεται ἡ ψυχὴ ἐκείνη ἐκ μέσου τοῦ λαοῦ αὐτοῦ.
\vs{15}ἓξ ἡμέρας ποιήσεις ἔργα, τῇ δὲ ἡμέρᾳ τῇ ἑβδόμῃ σάββατα, ἀνάπαυσις ἁγία τῷ κυρίῳ· πᾶς ὃς ποιήσει ἔργον τῇ ἡμέρᾳ τῇ ἑβδόμῃ θανατωθήσεται.
\vs{16}Καὶ φυλάξουσιν οἱ υἱοὶ Ἰσραὴλ τὰ σάββατα, ποιεῖν αὐτὰ εἰς τὰς γενεὰς αὐτῶν·
\vs{17}Διαθήκη αἰώνιος ἐν ἐμοὶ καὶ τοῖς υἱοῖς Ἰσραὴλ, σημεῖόν ἐστιν ἐν ἐμοὶ αἰώνιον· ὅτι ἓξ ἡμεραις ἐποίησε Κύριος τὸν οὐρανὸν καὶ τὴν γῆν, καὶ τῇ ἡμέρᾳ τῇ ἑβδόμῃ κατέπαυσε, καὶ ἐπαύσατο.
\vs{18}Καὶ ἔδωκε Μωυσῇ ἡνίκα κατέπαυσε λαλῶν αὐτῷ ἐν τῷ ὄρει τῷ Σινὰ, τὰς δύο πλάκας τοῦ μαρτυρίου, πλάκας λιθίνας γεγραμμένας τῷ δακτύλῳ τοῦ Θεοῦ.

\ch{32}
Καὶ ἰδὼν ὁ λαὸς, ὅτι κεχρόνικε Μωυσῆς καταβῆναι ἐκ τοῦ ὄρους, συνέστη ὁ λαὸς ἐπὶ Ἀαρὼν, καὶ λέγουσιν αὐτῷ ἀνάστηθι, καὶ ποίησον ἡμῖν θεοὺς, οἳ προπορεύσονται ἡμῶν· ὁ γὰρ Μωυσῆς οὗτος ὁ ἄνθρωπος ὃς ἐξήγαγεν ἡμᾶς ἐκ γῆς Αἰγύπτου, οὐκ οἴδαμεν τί γέγονεν αὐτῷ.
\vs{2}Καὶ λέγει αὐτοῖς Ἀαρὼν Περιέλεσθε τὰ ἐνώτια τὰ χρυσᾶ τὰ ἐν τοῖς ὠσὶ τῶν γυναικῶν ὑμῶν καὶ θυγατέρων, καὶ ἐνέγκατε πρός με.
\vs{3}Καὶ περιείλαντο πᾶς ὁ λαὸς τὰ ἐνώτια τὰ χρυσᾶ τὰ ἐν τοῖς ὠσὶν αὐτῶν, καὶ ἤνεγκαν πρὸς Ἀαρών.
\vs{4}Καὶ ἐδέξατο ἐκ τῶν χειρῶν αὐτῶν, καὶ ἔπλασεν αὐτὰ ἐν τῇ γραφίδι· καὶ ἐποίησεν αὐτὰ μόσχον χωνευτὸν καὶ εἶπεν, Οὗτοι οἱ θεοί σου Ἰσραὴλ, οἵτινες ἀνεβίβασάν σε ἐκ γῆς Αἰγύπτου.
\vs{5}Καὶ ἰδὼν Ἀαρὼν ᾠκοδόμησε θυσιαστήριον κατέναντι αὐτοῦ· καὶ ἐκήρυξεν Ἀαρὼν λέγων, ἑορτὴ τοῦ κυρίου αὔριον.
\vs{6}Καὶ ὀρθρίσας τῇ ἐπαύριον ἀνεβίβασεν ὁλοκαυτώματα, καὶ προσήνεγκε θυσίαν σωτηρίου· καὶ ἐκάθισεν ὁ λαὸς φαγεῖν καὶ πιεῖν, καὶ ἀνέστησαν παίζειν.

\vs{7}Καὶ ἐλάλησε Κύριος πρὸς Μωυσῆν, λέγων, βάδιζε τὸ τάχος, κατάβηθι ἐντεύθεν· ἠνόμησε γὰρ ὁ λαός σου ὃν ἐξήγαγες ἐκ γῆς Αἰγύπτου.
\vs{8}Παρέβησαν ταχὺ ἐκ τῆς ὁδοῦ, ἧς ἐνετείλω αὐτοῖς· ἐποίησαν ἑαυτοῖς μόσχον, καὶ προσκεκυνήκασιν αὐτῷ, καὶ τεθύκασιν αὐτῷ, καὶ εἶπαν, Οὗτοι οἱ θεοί σου Ἰσραὴλ, οἵτινες ἀνεβίβασάν σε ἐκ γῆς Αἰγύπτου.
\vs{10}καὶ νῦν ἔασόν με, καὶ θυμωθεὶς ὀργῇ εἰς αὐτοὺς, ἐκτρίψω αὐτούς· καὶ ποιήσω σὲ εἰς ἔθνος μέγα.
\vs{11}καὶ ἐδεήθη Μωυσῆς ἔναντι Κυρίου τοῦ Θεοῦ, καὶ εἶπεν, ἱνατί, Κύριε, θυμοῖ ὀργῇ εἰς τὸν λαόν σου, οὓς ἐξήγαγες ἐκ γῆς Αἰγύπτου ἐν ἰσχύϊ μεγάλῃ, καὶ ἐν τῷ βραχίονί σου τῷ ὑψηλῷ;
\vs{12}Μή ποτε εἴπωσιν οἱ Αἰγύπτιοι λέγοντες Μετὰ πονηρίας ἐξήγαγεν αὐτοὺς ἀποκτεῖναι ἐν τοῖς ὄρεσιν καὶ ἐξαναλῶσαι αὐτοὺς ἀπὸ τῆς γῆς. παῦσαι τῆς ὀργῆς τοῦ θυμοῦ σου, καὶ ἵλεως γενοῦ ἐπὶ τῇ κακίᾳ τοῦ λαοῦ σου,
\vs{13}μνησθεὶς Ἀβραὰμ καὶ Ἰσαὰκ καὶ Ἰακὼβ τῶν σῶν οἰκετῶν, οἷς ὤμοσας κατὰ σεαυτοῦ, καὶ ἐλάλησας πρὸς αὐτοὺς, λέγων, πολυπληθυνῶ τὸ σπέρμα ὑμῶν ὡσεὶ τὰ ἄστρα τοῦ οὐρανοῦ τῷ πλήθει· καὶ πᾶσαν τὴν γῆν ταύτην ἣν εἶπας δοῦναι αὐτοῖς, καὶ καθέξουσιν αὐτὴν εἰς τὸν αἰῶνα.
\vs{14}Καὶ ἱλάσθη Κύριος περιποιῆσαι τὸν λαὸν αὐτοῦ.

\vs{15}Καὶ ἀποστρέψας Μωυσῆς, κατέβη ἀπὸ τοῦ ὄρους· καὶ αἱ δύο πλάκες τοῦ μαρτυρίου ἐν ταῖς χερσὶν αὐτοῦ, πλάκες λίθιναι καταγεγραμμέναι ἐξ ἀμφοτέρων τῶν μερῶν αὐτῶν, ἔνθεν καὶ ἔνθεν ἦσανς γεγραμμέναι·
\vs{16}καὶ αἱ πλάκες ἔργον Θεοῦ ἦσαν, καὶ ἡ γραφὴ γραφὴ Θεοῦ κεκολαμμένη ἐν ταῖς πλαξί.
\vs{17}καὶ ἀκούσας Ἰησοῦς τῆς φωνῆς τοῦ λαοῦ κραζόντων, λέγει πρὸς Μωυσὴν, Φωνὴ πολέμου ἐν τῇ παρεμβολῇ.
\vs{18}καὶ λέγει Οὐκ ἔστι φωνὴ ἐξαρχόντων κατʼ ἰσχὺν, οὐδὲ φωνὴ ἐξαρχόντων τροπῆς, ἀλλὰ φωνὴν ἐξαρχόντων οἴνου ἐγὼ ἀκούω.

\vs{19}καὶ ἡνίκα ἤγγιζε τῇ παρεμβολῇ, ὁρᾷ τὸν μόσχον καὶ τοὺς χορούς· καὶ ὀργισθεὶς θυμῷ Μωυσῆς ἔῤῥιψεν ἀπὸ τῶν χειρῶν αὐτοῦ τὰς δύο πλάκας, καὶ συνέτριψεν αὐτὰς ὑπὸ τὸ ὄρος·
\vs{20}καὶ λαβὼν τὸν μόσχον ὃν ἐποίησαν, κατέκαυσεν αὐτὸν ἐν πυρὶ, καὶ κατήλεσεν αὐτὸν λεπτὸν, καὶ ἔσπειρεν αὐτὸν ὑπὸ τὸ ὕδωρ, καὶ ἐπότισεν αὐτὸ τοὺς υἱοὺς Ἰσραήλ.
\vs{21}καὶ εἶπε Μωυσῆς τῷ Ἀαρὼν, Τί ἐποίησέ σοι ὁ λαὸς οὗτος, ὅτι ἐπήγαγες ἐπʼ αὐτοὺς ἁμαρτίαν μεγάλην;
\vs{22}καὶ εἶπεν Ἀαρὼν πρὸς Μωυσῆν, μὴ ὀργίζου, κύριε· σὺ γὰρ οἶδας τὸ ὅρμημα τοῦ λαοῦ τούτου.
\vs{23}Λέγουσι γάρ μοι, ποιήσον ἡμῖν θεοὺς, οἳ προπορεύσονται ἡμῶν· ὁ γὰρ Μωυσῆς οὗτος ὁ ἄνθρωπος, ὃς ἐξήγαγεν ἡμᾶς ἐξ Αἰγύπτου, οὐκ οἴδαμεν τί γέγονεν αὐτῷ.
\vs{24}καὶ εἶπα αὐτοῖς, εἴ τινι ὑπάρχει χρυσία, περιέλεσθε· καὶ ἔδωκάν μοι· καὶ ἔῤῥιψα εἰς τὸ πῦρ· καὶ ἐξῆλθεν ὁ μόσχος οὗτος.
\vs{25}Καὶ ἰδὼν Μωυσῆς τὸν λαὸν ὅτι διεσκέδασται· (διεσκέδασε γὰρ αὐτοὺς Ἀαρών ἐπίχαρμα τοῖς ὑπεναντίοις αὐτῶν)
\vs{26}ἔστη δὲ Μωυσῆς ἑπὶ τῆς πύλης τῆς παρεμβολῆς, καὶ εἶπε, τίς πρὸς Κύριον; ἴτω πρός με. Συνῆλθον οὖν πρὸς αὐτὸν πάντες οἱ υἱοὶ Λευί.
\vs{27}Καὶ λέγει αὐτοῖς τάδε λέγει Κύριος ὁ Θεὸς Ἰσραήλ θέσθε ἕκαστος τὴν ἑαυτοῦ ῥομφαίανν ἐπὶ τὸν μηρὸν, καὶ διέλθατε καὶ ἀνακάμψατε ἀπὸ πύλης ἐπὶ πύλην διὰ τῆς παρεμβολῆς, καὶ ἀποκτείνατε ἕκαστος τὸν ἀδελφὸν αὐτοῦ, καὶ ἕκαστος τὸν πλησίον αὐτοῦ, καὶ ἓκαστος τὸν ἔγγιστα αὐτοῦ.
\vs{28}Καὶ ἐποίησαν οἱ υἱοὶ Λευεὶ καθὰ ἐλάλησεν αὐτοῖς Μωυσῆς· καὶ ἔπεσαν ἐκ τοῦ λαοῦ ἐν ἐκείνῃ τῇ ἡμέρᾳ εἰς τρισχιλίους ἄνδρας.
\vs{29}Καὶ εἶπεν αὐτοῖς Μωυσῆς, ἐπληρώσατε τὰς χεῖρας ὑμῶν σήμερον Κυρίῳ ἕκαστος ἐν τῷ υἱῷ ἢ ἐν τῷ ἀδελφῷ αὐτοῦ, δοθῆναι ἐφʼ ὑμᾶς εὐλογίαν.

\vs{30}Καὶ ἐγένετο μετὰ τὴν αὔριον εἶπε Μωυσῆς πρὸς τὸν λαὸν, ὑμεῖς ἡμαρτήκατε ἁμαρτίαν μεγάλη· καὶ νῦν ἀναβήσομαι πρὸς τὸν Θεὸν, ἵνα ἐξιλάσωμαι περὶ τῆς ἁμαρτίας ὑμῶν.
\vs{31}Ὑπέστρεψε δὲ Μωυσῆς πρὸς Κύριον, καὶ εἶπε, δέομαι, κύριε· ἡμάρτηκεν ὁ λαὸς οὗτος ἁμαρτίαν μεγάλην, καὶ ἐποίησαν ἑαυτοῖς θεοὺς χρυσοῦς.
\vs{32}Καὶ νῦν εἰ μὲν ἀφεῖς αὐτοῖς τὴν ἁμαρτίαν αὐτῶν, ἄφες· εἰ δὲ μή, ἐξάλειψόν με ἐκ τῆς βίβλου σου, ἧς ἔγραψας.
\vs{33}Καὶ εἶπε Κύριος πρὸς Μωυσῆν, εἴ τις ἡμάρτηκεν ἐνώπιόν μου, ἐξαλείψω αὐτοὺς ἐκ τῆς βίβλου μου.
\vs{34}Νυνὶ δὲ βάδιζε, κατάβηθι, καὶ ὁδήγησον τὸν λαὸν τοῦτον εἰς τὸν τόπον, ὃν εἶπά σοι· ἰδοὺ ὁ ἄγγελός μου προπορεύσεται πρὸ προσώπου σου· ᾖ δʼ ἂν ἡμέρᾳ ἐπισκέπτωμαι, ἐπάξω ἐπʼ αὐτοὺς τὴν ἁμαρτίαν αὐτῶν
\vs{35}Καὶ ἐπάταξε Κύριος τὸν λαὸν περὶ τῆς ποιήσεως τοῦ μόσχου, οὗ ἐποίησεν Ἀαρών.

\ch{33}
Καὶ εἶπε Κύριος πρὸς Μωυσῆς, προπορεύου, ἀνάβηθι ἐντεῦθεν σὺ καὶ ὁ λαός σου, οὓς ἐξήγαγες ἐκ γῆς Αἰγύπτου, εἰς τὴν γῆν, ἣν ὤμοσα τῷ Ἀβραὰμ, καὶ Ἰσαὰκ, καὶ Ἰακὼβ, λέγων, Τῷ σπέρματι ὑμῶν δώσω αὐτήν.
\vs{2}Καὶ συναποστελῶ τὸν ἄγγελόν μου πρὸ προσώπου σου· καὶ ἐκβαλεῖ τὸν Ἀμοῤῥαῖον, καὶ Χετταῖον, καὶ Φερεζαῖον, καὶ Γεργεσαῖον, καὶ Εὐαῖον, καὶ Ἰεβουσαῖον, καὶ Χαναναῖον.
\vs{3}Καὶ εἰσάξω σε εἰς γῆν ῥέουσαν γάλα καὶ μέλι· οὐ γὰρ μὴ συναναβῶ μετὰ σου, διὰ τὸ λαὸν σκληροτράχηλόν σε εἶναι, ἵνα μὴ ἐξαναλώσω σεε ἐν τῇ ὁδῷ.
\vs{4}Καὶ ἀκούσας ὁ λαὸς τὸ ῥῆμα τὸ πονηρὸν τοῦτο, κατεπένθησεν ἐν πενθικοῖς.
\vs{5}Καὶ εἶπε Κύριος τοῖς υἱοῖς Ἰσραὴλ, ὑμεῖς λαὸς σκληροτράχηλος· ὁρᾶτε, μὴ πληγὴν ἄλλην ἐπάξω ἐγὼ ἐφʼ ὑμᾶς, καὶ ἐξαναλώσω ὑμᾶς· νῦν οὖν ἀφέλεσθε τὰς στολὰς τῶν δοξῶν ὑμῶν, καὶ τὸν κόσμον, καὶ δείξω σοι ἃ ποιήσω σοι.
\vs{6}Καὶ περιέλαντο οἱ υἱοὶ Ἰσραὴλ τὸν κόσμον αὐτῶν, καὶ τὴν περιστολὴν ἀπὸ τοῦ ὄρους τοῦ Χωρήβ.
\vs{7}Καὶ λαβὼν Μωυσῆς τὴν σκηνὴν αὐτοῦ, ἔπηξεν ἔξω τῆς παρεμβολῆς, μακρὰν ἀπὸ τῆς παρεμβολῆς· καὶ ἐκλήθν Σκηνὴ μαρτυρίου· καὶ ἐγένετο, πᾶς ὁ ζητῶν Κύριον ἐξεπορεύετο εἰς τὴν σκηνὴν τὴν ἔξω τῆς παρεμβολῆς.
\vs{8}Ἡνίκα δʼ ἂν εἰσεπορεύετο Μωυσῆς εἰς τὴν σκηνὴν ἔξω τῆς παρεμβολῆς, εἱστήκει πᾶς ὁ λαὸς σκοπεύοντες ἕκαστος παρὰ τὰς θύρας τῆς σκηνῆς αὐτοῦ· καὶ κατενοοῦσαν ἀπιόντος Μωυσῆ ἕως τοῦ εἰσελθεῖν αὐτὸν εἰς τὴν σκηνὴν.
\vs{9}Ὡς δʼ ἂν εἰσῆλθε Μωσῆς εἰς τὴν σκηνήν, κατέβαινεν ὁ στύλος τῆς νεφέλης, καὶ ἵστατο ἐπὶ τὴν θύραν τῆς σκηνῆς, καὶ ἐλάλει Μωσῇ·
\vs{10}καὶ ἐλάλει Μωυσῇ. Καὶ ἑώρα πᾶς ὁ λαὸς τὸν στύλον τῆς νεφέλης ἑστῶτα ἐπὶ τῆς θύρας τῆς σκηνῆς· καὶ στάντες πᾶς ὁ λαὸς, προσεκύνησαν ἕκαστος ἀπὸ τῆς θύρας τῆς σκηνῆς αὐτοῦ.
\vs{11}Καὶ ἐλάλησε Κύριος πρὸς Μωυσῆν, ἐνώπιος ἐνωπίῳ, ὡς εἴτις λαλήσει πρὸς τὸν ἑαυτοῦ φίλον· καὶ ἀπελύετο εἰς τὴν παρεμβολήν· ὁ δὲ θεράπων Ἰησοῦς υἱὸς Ναυὴ νέος οὐκ ἐξεπορεύετο ἐκ τῆς σκηνῆς.

\vs{12}Καὶ εἶπε Μωυσῆς, πρὸς Κύριον, Ἰδοὺ σύ μοι λέγεις, ἀνάγαγε τὸν λαὸν τοῦτον, σὺ δὲ οὐκ ἐδήλωσάς μοι, ὃν συναποστελεῖς μετʼ ἐμοῦ· σὺ δέ μοι εἶπας, Οἶδά σε παρὰ πάντας, καὶ χάριν ἔχεις παρʼ ἐμοί.
\vs{13}Εἰ οὖν εὕρηκα χάριν ἐναντίον σου, ἐμφάνισόν μοι σεαυτόν· γνωστῶς ἴδω ἴδω σε, ὅπως ἂν ὦ εὑρηκὼς χάριν ἐναντίον σου, καὶ ἵνα γνῶ, ὅτι λαός σου τὸ ἔθνος τὸ μέγα τοῦτο.
\vs{14}Καὶ λέγει, αὐτὸς προπορεύσομαί σου, καὶ καταπαύσω σε.
\vs{15}Καὶ λέγει πρὸς αὐτόν, εἰ μὴ αὐτὸς σὺ σνμπορεύῃ, μή με ἀναγάγῃς ἐντεῦθεν.
\vs{16}Καὶ πῶς γνωστὸν ἔσται ἀληθῶς, ὅτι εὕρηκα χάριν παρὰ σοί ἐγώ τε καὶ ὁ λαός σου, ἀλλʼ ἢ συμπορευομένου σου μεθʼ ἡμῶν; καὶ ἐνδοξασθήσομαι ἐγώ τε καὶ ὁ λαός σου παρὰ πάντα τὰ ἔθνη, ὅσα ἐπὶ τῆς γῆς ἐστί.
\vs{17}Καὶ εἶπε Κύριος πρὸς Μωυσῆν, Καὶ τοῦτόν σοι τὸν λόγον, ὃν εἴρηκας ποιήσω· εὕρηκας, ποιήσω· εὕρηκας γὰρ χάριν ἐνώπιον ἐμοῦ, καὶ οἶδά σε παρὰ πάντας.
\vs{18}Καὶ λέγει, ἐμφάνισόν μοι σεαυτόν.
\vs{19}Καὶ εἶπεν, ἐγὼ παρελεύσομαι πρότερός σου τῇ δόξῃ μου, καὶ καλέσω τῷ ὀνόματί μου, Κύριος ἐναντίον σου· καὶ ἐλεήσω, ὃν ἂν ἐλεῶ, καὶ οἰκτειρήσω, ὃν ἂν οἰκτείρῶ.
\vs{20}Καὶ εἶπε, οὐ δυνήσῃ ἰδεῖν τὸ πρόσωπόν μου· οὐ γὰρ μὴ ἴδῃ ἄνθρωπος τὸ πρόσωπόν μου, καὶ ζήσεται.
\vs{21}Καὶ εἶπεν Κύριος, Ἰδοὺ τόπος παρʼ ἐμοί, στήσῃ ἐπὶ τῆς πέτρας·
\vs{22}Ἡνίκα δʼ ἂν παρέλθηνᾑ ᾑ δόξα μου, καί θήσω σε εἰς ὀπὴν τῆς πέτρας, καὶ σκεπάσω τῇ χειρί μου ἐπὶ σὲ, ἕως ἂν παρέλθω.
\vs{23}Καὶ ἀφελῶ τὴν χεῖρα, καὶ τότε ὄψει τὰ ὀπίσω μου· τὸ δὲ πρόσωπόν μου οὐκ ὀφθήσεταί σοι.

\ch{34}
Καὶ εἶπε Κύριος πρὸς Μωυσῆν λαξευσον σεαυτῷ δύο πλάκας λιθίνας, καθὼς καὶ αἱ πρῶται, καὶ ἀνάβηθι πρὸς μὲ εἰς τὸ ὄρος, καὶ γράψω ἐπὶ τῶν πλακῶν τὰ ῥήματα ἃ ἦν ἐν ταῖς πλαξὶ ταῖς πρώταις, αἷς συνέτριψας.
\vs{2}Καὶ γίνου ἕτοιμος εἰς τὸ πρωί, καὶ ἀναβήσῃ ἐπὶ τὸ ὄρος τὸ Σινά, καὶ στήσῃ μοι ἐκεῖ ἐπʼ ἄκρου τοῦ ὄρους.
\vs{3}Καὶ μηδεὶς ἀναβήτω μετὰ σοῦ μηδὲ ὀφθήτω ἐν παντὶ τῷ ὄρει· καὶ τὰ πρόβατα καὶ βόες μὴ νεμέσθωσαν πλησίον τοῦ ὄρους ἐκείνου.
\vs{4}Καὶ ἐλάξευσε δύο πλάκας λιθίνας, καθάπερ καὶ αἱ πρῶται· καὶ ὀρθρίσας Μωυσῆς, ἀνέβη εἰς τὸ ὄρος τὸ Σινὰ, καθότι συνέταξεν αὐτῷ Κύριος· καὶ ἔλαβε Μωυσῆς τὰς δύο πλάκας τὰς λιθίνας.
\vs{5}Καὶ κατέβη Κύριος ἐν νεφέλῃ, καὶ παρέστη αὐτῷ ἐκεῖ, καὶ ἐκάλεσε τῷ ὀνόματι Κυρίου.
\vs{6}Καὶ παρῆλθε Κύριος πρὸ προσώπου αὐτοῦ, καὶ ἐκάλεσε κύριος ὁ Θεὸς οἰκτείρμων καὶ ἐλεήμων, μακρόθυμος καὶ πολυέλεος καὶ ἀληθινός,
\vs{7}καὶ δικαιοσύνην διατηρῶν καὶ ἔλεος εἰς χιλιάδας, ἀφαιρῶν ἀνομίας καὶ ἀδικίας καὶ ἁμαρτίας, καὶ οὐ καθαριεῖ τὸν ἔνοχον, ἐπάγων ἀνομίας πατέρων ἐπὶ τέκνα καὶ ἐπὶ τέκνα τέκνων ἐπὶ τρίτην καὶ τετάρτην γενεάν.
\vs{8}Καὶ σπεύσας Μωσῆς κύψας ἐπὶ τὴν γῆν προσεκύνησε·
\vs{9}καὶ εἶπεν, εἰ εὕρηκα χάριν ἐνώπιόν σου, συμπορευθήτω ὁ Κύριός μου μεθʼ ἡμῶν· ὁ λαὸς γὰρ σκληροτράχηλός ἐστι, καὶ ἀφελεῖς σὺ τὰς ἁμαρτίας ἡμῶν, καὶ τὰς ἀνομίας ἡμῶν, καὶ ἐσόμεθά σοι.

\vs{10}Καὶ εἶπε Κύριος πρὸς Μωυσῆν, Ἰδοὺ, ἐγὼ τίθημί σοι διαθήκην ἐνώπιον παντὸς τοῦ λαοῦ σοῦ, ποιήσω ἔνδοξα, ἃ οὐ γέγονεν ἐν πάσῃ τῇ γῇ, καὶ ἐν παντὶ ἔθνει· καὶ ὄψεται πᾶς ὁ λαὸς, ἐν οἷς εἶ σὺ, τὰ ἔργα Κυρίου, ὅτι θαυμαστά ἐστιν, ἃ ἐγὼ ποιήσω σοι.
\vs{11}Πρόσεχε σὺ πάντα ὅσα ἐγὼ ἐντέλλομαί σοι· ἰδοὺ ἐγὼ ἐκβάλλω πρὸ προσώπου ὑμῶν τὸν Ἀμοῤῥαῖον, καὶ Χαναναῖον, καὶ Φερεζαῖον, καὶ Χετταῖον, καὶ Εὑαῖον, καὶ Γεργεσαῖον, καὶ Ἰεβουσαῖον.
\vs{12}Πρόσεχε σεαυτῷ, μή ποτε θῇς διαθήκην τοῖς ἐγκαθημένοις ἐπὶ τῆς γῆς, εἰς ἣν εἰσπορεύῃ εἰς αὐτὴν, μή σοι γένηται πρόσκομμα ἐν ὑμῖν.
\vs{13}Τοὺς βωμοὺς αὐτῶν καθελεῖτε, καὶ τὰς στήλας αὐτῶν συντρίψετε, καὶ τὰ ἄλση αὐτῶν ἐκκόψετε, καὶ τὰ γλυπτὰ τῶν θεῶν αὐτῶν κατακαύσετε ἐν πυρί.
\vs{14}Οὐ γὰρ μὴ προσκυνήσητε θεοῖς ἑτέροις· ὁ γὰρ Κύριος ὁ Θεὸς, ζηλωτὸν ὄνομα, Θεὸς ζηλωτής ἐστι.
\vs{15}Μή ποτε θῇς διαθήκην τοῖς ἔγκαθημένοις ἐπὶ τῆς γῆς, καὶ ἐκπορνεύσωσιν ὀπίσω τῶν θεῶν αὐτῶν, καὶ θύσωσι τοῖς θεοῖς αὐτῶν, καὶ καλέσωσίν σε, καὶ φάγῃς τῶν αὐτῶν,
\vs{16}καὶ λάβῃς τῶν θυγατέρων αὐτῶν τοῖς υἱοῖς σου, καὶ τῶν θυγατέρων σου δῷς τοῖς υἱοῖς αὐτῶν, καὶ ἐκπορνεύσωσιν αἱ θυγατέρες σου ὀπίσω τῶν θεῶν αὐτῶν, καὶ ἐκπορνεύσωσιν οἱ υἱοί σου ὀπίσω τῶν θεῶν αὐτῶν.
\vs{17}Καὶ θεοὺς χωνευτοὺς οὐ ποιήσεις σεαυτῷ.
\vs{18}Καὶ τὴν ἑσρτὴν τῶν ἀζύμων φυλάξῃ· ἑπτὰ ἡμέρας φαγῃ ἄζυμα, καθάπερ ἐντέταλμαί σοι, εἰς τὸν καιρὸν ἐν μηνὶ τῶν νέων· ἐν γὰρ μηνὶ τῶν νέων ἐξῆλθες ἐξ Αἰγύπτου.
\vs{19}Πᾶν διανοῖγον μήτραν, ἐμοὶ τὰ ἀρσενικὰ, πᾶν πρωτότοκον μόσχου, καὶ πρωτότοκον προβάτου.
\vs{20}Καὶ πρωτότοκον ὑποζυγίου λυτρώσῃ προβάτῳ· ἐὰν δὲ μὴ λυτρώσῃ αὐτὸ, τιμὴν δώσεις. πᾶν πρωτότοκον τῶν υἱῶν σου λυτρώσῃ· οὐκ ὀφθήσῃ ἐνώπιόν μου κενός.

\vs{21}Ἓξ ἡμέρας ἐργᾷ, τῇ δὲ ἑβδόμῃ καταπαύσεις· τῷ σπόρῳ καὶ τῷ ἀμητῷ κατάπαυσις.
\vs{22}Καὶ ἑορτὴν ἑβδομάδων ποιήσεις μοι, ἀρχὴν θερισμοῦ πυροῦ· καὶ ἐορτὴν συναγωγῆς μεσοῦντος τοῦ ἐνιαυτοῦ.
\vs{23}Τρεῖς καιροὺς τοῦ ἐνιαυτοῦ ὀφθήσεται πᾶν ἀρσενικόν σου ἐνώπιον Κυρίου τοῦ Θεοῦ Ἰσραήλ.
\vs{24}Ὅταν γὰρ ἐκβάλω τὰ ἔθνη πρὸ προσώπου σου, καὶ πλατυνῶ τὰ ὅριά σου, οὐκ ἐπιθυμήσει οὐδεὶς τῆς γῆς σου, ἡνίκα ἂν ἀναβαίνῃς ὀφθῆναι ἐναντίον Κυρίου τοῦ Θεοῦ σου, τρεῖς καιροὺς τοῦ ἐνιαυτοῦ.
\vs{25}Οὐ σφάξεις ἐπὶ ζύμῃ αἷμα θυμιαμάτων μου, καὶ οὐ κοιμηθήσεται εἰς τὸ πρωῒ θύματα ἑορτῆς τοῦ πάσχα.
\vs{26}Τὰ πρωτογεννήματα τῆς γῆς σου θήσεις εἰς τὸν οἶκον Κυρίου τοῦ Θεοῦ σου· οὐχ ἑψήσεις ἄρνα ἐν γάλακτι μητρὸς αὐτοῦ.
\vs{27}Καὶ εἶπε Κύριος πρὸς Μωυσῆν, γράψον σεαυτῷ τὰ ῥήματα ταῦτα· ἐπὶ γὰρ τῶν λόγων τούτων τέθειμαί σοι διαθήκην, καὶ τῷ Ἰσραήλ.
\vs{28}Καὶ ἦν ἐκεῖ Μωυσῆς ἐναντίον Κυρίου τεσσεράκοντα ἡμέρας, καὶ τεσσεράκοντα νύκτας· ἄρτον οὐκ ἔφαγε, καὶ ὕδωρ οὐκ ἔπιε· καὶ ἔγραψεν ἐπὶ τῶν πλακῶν τὰ ῥήματα ταῦτα τῆς διαθήκης, τοὺς δέκα λόγους.

\vs{29}Ὡς δὲ κατέβαινε Μωυσῆς ἐκ τοῦ ὄρους, καὶ αἱ δύο πλάκες ἐπὶ τῶν χειρῶν Μωυσῆ· καταβαίνοντος δὲ αὐτοῦ ἐκ τοῦ ὄρους, Μωυσῆς οὐκ ᾔδει ὅτι δεδόξασται ἡ ὄψις τοῦ χρώματος τοῦ προσώπου αὐτοῦ ἐν τῷ λαλεῖν αὐτὸν αὐτῷ.
\vs{30}Καὶ εἶδεν Ἀαρὼν καὶ πάντες οἱ πρεσβύτεροι Ἰσραὴλ τὸν Μωυσῆν, καὶ ἦν δεδοξασμένη ἡ ὄψις τοῦ χρώματος τοῦ προσώπου αὐτοῦ. καὶ ἐφοβήθησαν ἐγγίσαι αὐτῷ.
\vs{31}Καὶ ἐκάλεσεν αὐτοὺς Μωυσῆς, καὶ ἐπεστράφησαν πρὸς αὐτὸν Ἀαρὼν καὶ πάντες οἱ ἄρχοντες τῆς συναγωγῆς· καὶ ἐλάλησεν αὐτοῖς Μωυσῆς.

\vs{32}Καὶ μετὰ ταῦτα προσῆλθον πρὸς αὐτὸν πάντες οἱ υἱοὶ Ἰσραήλ. καὶ ἐνετείλατο αὐτοῖς πάντα, ὅσα ἐνετείλατο Κύριος πρὸς αὐτὸν ἐν τῷ ὄρει Σινά.
\vs{33}Καὶ ἐπειδὴ κατέπαυσε λαλῶν πρὸς αὐτοὺς, ἐπέθηκεν ἐπὶ τὸ πρόσωπον αὐτοῦ κάλυμμα.
\vs{34}Ἡνίκα δʼ ἂν εἰσεπορεύετο Μωυσῆς, ἔναντι Κυρίου λαλεῖν αὐτῷ, περιῃρεῖτο τὸ κάλυμμα ἕως τοῦ ἐκπορεύεσθαι· καὶ ἐξελθὼν ἐλάλει πᾶσι τοῖς υἱοῖς Ἰσραὴλ ὅσα ἐνετείλατο αὐτῷ Κύριος.
\vs{35}Καὶ εἶδον οἱ υἱοὶ Ἰσραὴλ τὸ πρόσωπον Μωυσέως, ὅτι δεδόξασται· καὶ περιέθηκε Μωυσῆς κάλυμμα ἐπὶ τὸ πρόσωπον ἑαυτοῦ, ἕως ἂν εἰσέλθῃ συλλαλεῖν αὐτῷ.

\ch{35}
Καὶ συνήθροισε Μωυσῆς πᾶσαν συναγωγὴν υἱῶν Ἰσραὴλ, καὶ εἶπεν, οὗτοι οἱ λόγοι, οὓς εἶπε Κύριος ποιῆσαι αὐτούς.
\vs{2}Ἓξ ἡμέρας ποιήσεις ἔργα, τῇ δὲ ἡμέρᾳ τῇ ἑβδόμῃ κατάπαυσις· ἅγια, σάββατα· ἀνάπαυσις Κυρίῳ· πᾶς ὁ ποιῶν ἔργον ἐν αὐτῇ, τελευτάτω.
\vs{3}Οὐ καύσετε πῦρ ἐν πάσῃ κατοικίᾳ ὑμῶν τῇ ἡμέρᾳ τῶν σαββάτων· ἐγὼ Κύριος.
\vs{4}Καὶ εἶπε Μωυσῆς πρὸς πᾶσαν συναγωγὴν υἱῶν Ἰσραὴλ, λέγων, τοῦτο τὸ ῥῆμα, ὃ συνέταξε Κύριος, λέγων,
\vs{5}λάβετε παρʼ ὑμῶν αὐτῶν ἀφαίρεμα Κυρίῳ· πᾶς ὁ καταδεχόμενος τῇ καρδίᾳ, οἴσουσι τὰς ἀπαρχὰς Κυρίῳ, χρυσίον, ἀργύριον, χαλκὸν,
\vs{6}ὑάκινθον, πορφύραν, κόκκινον διπλοῦν διανενησμένον, καὶ βύσσον κεκλωσμένην, καὶ τρίχας αἰγείας,
\vs{7}καὶ δέρματα κριῶν ἠρυθροδανωμένα, καὶ δέρματα ὑακινθινα, καὶ ξύλα ἄσηπτα,
\vs{9}καὶ λίθους σαρδίου, καὶ λίθους εἰς τὴν γλυφὴν εἰς τὴν ἐπωμίδα καὶ τὸν ποδήρη.
\vs{10}Καὶ πᾶς σοφὸς τῇ καρδίᾳ ἐν ὑμῖν, ἐλθὼν ἐργαζέσθω πάντα ὅσα συνέταξε Κύριος·
\vs{11}Τὴν σκηνὴν, καὶ τὰ παραρύματα, καὶ τὰ κατακαλύμματα, καὶ τὰ διατόνια, καὶ τοὺς μοχλοὺς, καὶ τοὺς στύλους,
\vs{12}καὶ τὴν κιβωτὸν τοῦ μαρτυρίου, καὶ τοὺς ἀναφορεῖς αὐτῆς, καὶ τὸ ἱλαστήριον αὐτῆς, καὶ τὸ καταπέτασμα,
\vs{12a}καὶ τὰ ἱστία τῆς αὐλῆς, καὶ τοὺς στύλους αὐτῆς, καὶ τοὺς λίθους τοὺς τῆς σμαράγδου, καὶ τὸ θυμίαμα, καὶ τὸ ἔλαιον τοῦ χρίσματος,
\vs{13}καὶ τὴν τράπεζαν καὶ πάντα τὰ σκεύη αὐτῆς,
\vs{14}καὶ τὴν λυχνίαν τοῦ φωτὸς καὶ πάντα τὰ σκεύη αὐτῆς,
\vs{16}καὶ τὸ θυσιαστήριον καὶ πάντα τὰ σκεύη αὐτοῦ,
\vs{19}καὶ τὰς στολὰς τὰς ἁγίας Ἀαρὼν τοῦ ἱερέως, καὶ τὰς στολὰς ἐν αἷς λειτουργήσουσιν ἐν αὐταῖς, καὶ τοὺς χιτῶνας τοῖς υἱοῖς Ἀαρὼν τῆς ἱερατίας, καὶ τὸ ἔλαιον τοῦ χρίσματος, καὶ τὸ θυμίαμα τῆς συνθέσεως.

\vs{20}Καὶ ἐξῆλθε πᾶσα συναγωγὴ υἱῶν Ἰσραὴλ ἀπὸ Μωυσῆ.
\vs{21}Καὶ ἤνεγκαν ἕκαστος, ὧν ἔφερεν ἡ καρδία αὐτῶν, καὶ ὅσοις ἔδοξε τῇ ψυχῇ αὐτῶν, ἀφαίρεμα· καὶ ἤνεγκαν ἀφαίρεμα Κυρίῳ εἰς πάντα τὰ ἔργα τῆς σκηνῆς τοῦ μαρτυρίου, καὶ εἰς πάντα τὰ κάτεργα αὐτῆς καὶ εἰς πάσας τὰς στολὰς τοῦ ἁγίου.
\vs{22}Καὶ ἤνεγκαν οἱ ἄνδρες παρὰ τῶν γυναικῶν, πᾶς ᾧ ἔδοξε τῇ διανοίᾳ, ἤνεγκαν σφραγίδας, καὶ ἐνώτια, καὶ δακτυλίους, καὶ ἐμπλόκια, καὶ περιδέξια, πᾶν σκεῦος χρυσοῦν.
\vs{23}Καὶ πάντες ὅσοι ἤνεγκαν ἀφαιρέματα χρυσίου Κυρίῳ, καὶ παρʼ ᾧ εὑρέθη βύσσος· καὶ δέρματα ὑακίνθινα καὶ δέρματα κριῶν ἠρυθροδανωμένα ἤνεγκαν.
\vs{24}Καὶ πᾶς ὁ ἀφαιρῶν ἀφαίρεμα, ἤνεγκαν ἀργύριον καὶ χαλκὸν, τὰ ἀφαιρέματα Κυρίῳ· καὶ παρʼ οἷς εὑρέθη ξύλα ἄσηπτα· καὶ εἰς πάντα τὰ ἔργα τῆς παρασκευῆς ἤνεγκαν.
\vs{25}Καὶ πᾶσα γυνὴ σοφὴ τῇ διανοίᾳ ταῖς χερσὶ νήθειν, ἤνεγκαν νενησμένα, τὴν ὑάκινθον, καὶ τὴν πορφύραν, καὶ τὸ κόκκινον, καὶ τὴν βύσσον.
\vs{26}Καὶ πᾶσαι αἱ γυναῖκες, αἷς ἔδοξε τῇ διανοίᾳ αὐτῶν ἐν σοφίᾳ, ἔνησαν τὰς τρίχας τὰς αἰγείας.
\vs{27}Καὶ οἱ ἄρχοντες ἤνεγκαν τοὺς λίθους τῆς σμαράγδου, καὶ τοὺς λίθους τῆς πληρώσεως εἰς τὴν ἐπωμίδα, καὶ τὸ λογεῖον,
\vs{28}καὶ τὰς συνθέσεις, καὶ εἰς τὸ ἔλαιον τῆς χρίσεως, καὶ τὴν συνθεσιν τοῦ θυμιάματος.
\vs{29}Καὶ πᾶς ἀνὴρ καὶ γυνὴ, ὧν ἔφερεν ἡ διάνοια αὐτῶν εἰσελθόντας ποιεῖν πάντα τὰ ἔργα, ὅσα συνέταξε Κύριος ποιῆσαι αὐτὰ διὰ Μωυσῆ, ἤνεγκαν οἱ υἱοὶ Ἰσραὴλ, ἀφαίρεμα Κυρίῳ.
\vs{30}Καὶ εἶπε Μωυσῆς τοῖς υἱοῖς Ἰσραὴλ, Ἰδοὺ ἀνακέκληκεν ὁ Θεὸς ἐξ ὀνόματος τὸν Βεσελεὴλ τὸν τοῦ Οὐρείου τὸν Ὢρ, ἐκ τῆς φυλῆς Ἰούδα,
\vs{31}καὶ ἐνέπλησεν αὐτὸν πνεῦμα θεῖον σοφίας καὶ συνέσεως, καὶ ἐπιστήμης πάντων,
\vs{32}ἀρχιτεκτονεῖν κατὰ πάντα τὰ ἔργα τῆς ἀρχιτεκτονίας, ποιεῖν τὸ χρυσίον καὶ τὸ ἀργύριον καὶ τὸν χαλκόν,
\vs{33}καὶ λιθουργῆσαι τὸν λίθον, καὶ κατεργάζεσθαι τὰ ξύλα, καὶ ποιεῖν ἐν παντὶ ἔργῳ σοφίας.
\vs{34}Καὶ προβιβάσαι γε ἔδωκεν ἐν τῇ διανοίᾳ αὐτῷ τε, καὶ τῷ Ἐλιὰβ τῷ τοῦ Ἀχισαμὰκ, ἐκ φυλῆς Δάν·
\vs{35}Καὶ ἐνέπλησεν αὐτοὺς σοφίας, συνέσεως, διανοίας, πάντα συνιέναι ποιῆσαι τὰ ἔργα τοῦ ἁγίου, καὶ τὰ ὑφαντὰ καὶ ποικιλτὰ ὑφάναι τῷ κοκκίνῳ, καὶ τῇ βύσσῳ, ποιεῖν πᾶν ἔργον ἀρχιτεκτονίας, ποικιλίας.

\ch{36}
Καὶ ἐποίησε Βεσελεὴλ καὶ Ἐλιὰβ, καὶ πᾶς σοφὸς τῇ διανοὶᾳ, ᾧ ἐδόθη σοφία καὶ ἐπιστήμη ἐν αὑτοῖς, συνιέναι ποιεῖν πάντα τὰ ἔργα, κατὰ τὰ ἅγια καθήκοντα, κατὰ πάντα ὅσα συνέταξε Κύριος.
\vs{2}Καὶ ἐκάλεσε Μωυσῆς Βεσελεὴλ καὶ Ἐλιὰβ, καὶ πάντας τοὺς ἔχοντας τὴν σοφίαν, ᾧ ἔδωκεν ὁ Θεὸς ἐπιστήμην ἐν τῇ καρδίᾳ, καὶ πάντας τοὺς ἑκουσίως βουλομένους προσπορεύεσθαι πρὸς τὰ ἔργα, ὥστε συντελεῖν αὐτά.
\vs{3}Καὶ ἔλαβον παρὰ Μωσῆ πάντα τὰ ἀφαιρέματα, ἃ ἤνεγκαν οἱ υἱοὶ Ἰσραὴλ εἰς πάντα τὰ ἔργα τοῦ ἁγίου ποιεῖν αὐτά· καὶ αὐτοὶ προσεδέχοντο ἔτι τὰ προσφερόμενα παρὰ τῶν φερόντων τὸ πρωΐ.
\vs{4}Καὶ παρεγίνοντο πάντες οἱ σοφοὶ οἱ ποιοῦντες τὰ ἔργα τοῦ ἁγίου, ἕκαστος κατὰ τὸ αὐτοῦ ἔργον, ὃ εἰργάζοντο αὐτοί.
\vs{5}Καὶ εἶπε πρὸς Μωυσῆν, ὅτι πλῆθος φέρει ὁ λαὸς κατὰ τὰ ἔργα ὅσα συνέταξε Κύριος ποιῆσαι.
\vs{6}Καὶ προσέταξε Μωυσῆς, καὶ ἐκήρυξεν ἐν τῇ παρεμβολῇ, λέγων, ἀνὴρ καὶ γυνὴ μηκέτι ἐργαζέσθωσαν εἰς τάς ἀπαρχὰς τοῦ ἁγίου· καὶ ἐκωλύθη ὁ λαὸς ἔτι προσφέρειν.
\vs{7}Καὶ τὰ ἔργα ἦν αὐτοῖς ἱκανὰ εἰς τὴν κατασκευὴν ποιῆσαι, καὶ προσκατέλιπον.
\vs{8}Καὶ ἐποίησε πᾶς σοφὸς ἐν τοῖς ἐργαζομένοις τὰς στολὰς τῶν ἁγίων, αἵ εἰσιν Ἀαρὼν τῷ ἱερεῖ, καθὰ συνέταξε Κύριος τῷ Μωυσῇ.
\vs{9}Καὶ ἐποίησε τὴν ἐπωμίδα ἐκ χρυσίου, καὶ ὑακίνθου, καὶ πορφύρας, καὶ κοκκίνου νενησμένου, καὶ βύσσου κεκλωσμένης·
\vs{10}καὶ ἐτμήθη τὰ πέταλα τοῦ χρυσίου τρίχες, ὥστε συνυφάναι σὺν τῇ ὑακίνθῳ, καὶ τῇ πορφύρᾳ, καὶ σὺν τῷ κοκκίνῳ τῷ διανενησμένῳ, καὶ τῇ βύσσῳ τῇ κεκλωσμένῃ· ἔργον ὑφαντὸν ἐποίησαν αὐτό·
\vs{11}ἐπωμίδας συνεχούσας ἐξ ἀμφοτέρων τῶν μερῶν, ἔργον ὑφαντὸν εἰς ἄλληλα συμπεπλεγμένα καθʼ ἑαυτό.
\vs{12}Ἐξ αὐτοῦ ἐποίησαν αὐτὸ κατὰ τὴν αὐτοῦ ποίησιν, ἐκ χρυσίου, καὶ ὑακίνθου, καὶ πορφύρας, καὶ κοκκίνου διανενησμένου, καὶ βύσσου κεκλωσμένης, καθὰ συνέταξε Κύριος τῷ Μωυσῇ·
\vs{13}καὶ ἐποίησαν ἀμφοτέρους τοὺς λίθους τῆς σμαράγδου συνπεπορπημένους καὶ περισεσιαλωμένους χρυσίῳ, γεγλυμμένους καὶ ἐκκεκολαμμένους ἐγκόλαμμα σφραγίδος ἐκ τῶν ὀνομάτων τῶν υἱῶν Ἰσραήλ·
\vs{14}καὶ ἐπέθηκεν αὐτοὺς ἐπὶ τοὺς ὤμους τῆς ἐπωμίδος, λίθους μνημοσύνου τῶν υἱῶν Ἰσραὴλ, καθὰ συνέταξε Κύριος τῷ Μωυσῇ.

\vs{15}Καὶ ἐποίησαν λογεῖον, ἔργον ὑφαντὸν ποικιλίᾳ κατὰ τὸ ἔργον τῆς ἐπωμίδος, ἐκ χρυσίου, καὶ ὑακίνθου, καὶ πορφύρας, καὶ κοκκίνου διανενησμένου, καὶ βύσσου κεκλωσμένης·
\vs{16}τετράγωνον διπλοῦν ἐπόησαν τὸ λογεῖον· σπιθαμῆς τὸ μῆκος, καὶ σπιθαμῆς τὸ εὖρος διπλοῦν.
\vs{17}Καὶ συνυφάνθη ἐν αὐτῷ ὕφασμα κατάλιθον τετράστιχον· στίχος λίθων, σάρδιον καὶ τοπάζιον καὶ σμάραγδος, ὁ στίχος ὁ εἷς·
\vs{18}καὶ ὁ στίχος ὁ δεύτερος, ἄνθραξ καὶ σάπφειρος καὶ ἴασπις·
\vs{19}καὶ ὁ στίχος ὁ τρίτος, λιγύριον καὶ ἀχάτης καὶ ἀμέθυστος·
\vs{20}καὶ ὁ στίχος ὁ τέταρτος, χρυσόλιθος καὶ βηρύλλιον καὶ ὀνύχιον περικεκυκλωμένα χρυσίῳ, καὶ συνδεδεμένα χρυσίῳ.
\vs{21}Καὶ οἱ λίθοι ἦσαν ἐκ τῶν ὀνομάτων τῶν υἱῶν Ἰσραὴλ δώδεκα, ἐκ τῶν ὀνομάτων αὐτῶν ἐγγεγλυμμένα εἰς σφραγίδας, ἕκαστος ἐκ τοῦ ἑαυτοῦ ὀνόματος εἰς τὰς δώδεκα φυλάς.
\vs{22}Καὶ ἐποίησαν ἐπὶ τὸ λογεῖον κρωσσοὺς συμπεπλεγμένους, ἔργον ἐμπλοκίου, ἐκ χρυσίου καθαροῦ.
\vs{23}Καὶ ἐποίησαν δύο ἀσπιδίσκας χρυσᾶς, καὶ δύο δακτυλίους χρυσοῦς· καὶ ἐπέθηκαν τοὺς δύο δακτυλίους τοὺς χρυσοῦς ἐπʼ ἀμφοτέρας τὰς ἀρχὰς τοῦ λογείου.
\vs{24}Καὶ ἐπέθηκαν τὰ ἐμπλόκια ἐκ χρυσίου ἐπὶ τοὺς δακτυλίους ἐπʼ ἀμφοτέρων τῶν μερῶν τοῦ λογείου·
\vs{25}καὶ εἰς τὰς δύο συμβολὰς τὰ δύο ἐμπλόκια. Καὶ ἐπέθηκαν ἐπὶ τὰς δύο ἀσπιδίσκας· καὶ ἐπέθηκαν ἐπὶ τοὺς ὤμους τῆς ἐπωμίδος ἐξεναντίας κατὰ πρόσωπον.
\vs{26}Καὶ ἐποίησαν δύο δακτυλίους χρυσοῦς, καὶ ἐπέθηκαν ἐπὶ τὰ δύο πτερύγια ἐπʼ ἄκρου τοῦ λογείου, καὶ ἐπὶ τὸ ἄκρον τοῦ ὀπισθίου τῆς ἐπωμίδος ἔσωθεν·
\vs{27}Καὶ ἐποίησαν δύο δακτυλίους χρυσοῦς, καὶ ἐπέθηκαν ἐπʼ ἀμφοτέρους τοὺς ὤμους τῆς ἐπωμίδος κάτωθεν αὐτοῦ, κατὰ πρόσωπον κατὰ τὴν συμβολὴν ἄνωθεν τῆς συνυφῆς τῆς ἐπωμίδος·
\vs{28}καὶ συνέσφιγξε τὸ λογεῖον ἀπὸ τῶν δακτυλίων τῶν ἐπʼ αὐτοῦ εἰς τοὺς δακτυλίους τῆς ἐπωμίδος, συνεχομένους ἐκ τῆς ὑακίνθου, συμπεπλεγμένους εἰς τὸ ὕφασμα τῆς ἐπωμίδος, ἵνα μὴ χαλᾶται τὸ λογεῖον ἀπὸ τῆς ἐπωμίδος, καθὰ συνέταξε Κύριος τῷ Μωυσῇ.
\vs{29}Καὶ ἐποίησαν τὸν ὑποδύτην ὑπὸ τὴν ἐπωμίδα, ἔργον ὑφαντὸν, ὅλον ὑακίνθινον·
\vs{30}τὸ δὲ περιστόμιον τοῦ ὑποδύτου ἐν τῷ μέσῳ διυφασμένον συμπλεκτὸν, ὤαν ἔχον κύκλῳ τὸ περιστόμιονν ἀδιάλυτον·
\vs{31}Καὶ ἐποίησαν ἐπὶ τοῦ λώματος τοῦ ὑποδύτου κάτωθεν ὡς ἐξανθούσης ῥόας ῥοΐσκους, ἐξ ὑακίνθου, καὶ πορφύρας, καὶ κοκκίνου νενησμένου, καὶ βύσσου κεκλωσμένης.
\vs{32}Καὶ ἐποίησαν κώδωνας χρυσοῦς, καὶ ἐπέθηκαν τοὺς κώδωνας ἐπὶ τὸ λῶμα τοῦ ὑποδύτου κύκλῳ ἀνὰ μέσον τῶν ῥοΐσκων·
\vs{33}κώδων χρυσοῦς καὶ ῥοΐσκος ἐπὶ τοῦ λώματος τοῦ ὑποδύτου κύκλῳ, εἰς τὸ λειτουργεῖν, καθὰ συνέταξε Κύριος τῷ Μωυσῇ.
\vs{34}Καὶ ἐποίησαν χιτῶνας βυσσίνους, ἔργον ὑφαντὸν, Ἀαρὼν καὶ τοῖς υἱοῖς αὐτοῦ,
\vs{35}καὶ τὰς κιδάρεις ἐκ βύσσου, καὶ τὴν μίτραν ἐκ βύσσου, καὶ τὰ περισκελῆ ἐκ βύσσου κεκλωσμένης,
\vs{36}καὶ τὰς ζώνας αὐτῶν ἐκ βύσσου, καὶ ὑακίνθου, καὶ πορφύρας, καὶ κοκκίνου νενησμένου, ἔργον ποικιλτοῦ, ὃν τρόπον συνέταξε Κύριος τῷ Μωυσῇ.
\vs{37}Καὶ ἐποίησαν τὸ πέταλον τὸ χρυσοῦν, ἀφόρισμα τοῦ ἁγίου, χρυσίου καθαροῦ· καὶ ἔγραψεν ἐπʼ αὐτοῦ γράμματα ἐκτετυπωμένα, σφραγίδος, Ἁγίασμα Κυρίῳ·
\vs{38}Καὶ ἐπέθηκαν ἐπὶ τὸ λῶμα ὑακίνθινον, ὥστε ἐπικεῖσθαι ἐπὶ τὴν μίτραν ἄνωθεν, ὅν τρόπον συνέταξε Κύριος τῷ Μωυσῇ.

\ch{37}
Καὶ ἐποίησαν τῇ σκηνῇ δέκα αὐλαίας·
\vs{2}ὀκτὼ καὶ εἴκοσι πήχεων μῆκος τῆς αὐλαίας τῆς μιᾶς· τὸ αὐτὸ ἦν πάσαις· καὶ τεσσάρων πήχεων τὸ εὖρος τῆς αὐλαίας τῆς μιᾶς.
\vs{3}καὶ ἐποίησαν τὸ καταπέτασμα ἐξ ὑακίνθου, καὶ πορφύρας, καὶ κοκκίνου νενησμένου, καὶ βύσσου κεκλωσμένης, ἔργον ὑφαντὸυ χερουβείμ·
\vs{4}καὶ ἐπέθηκαν αὐτὸ ἐπὶ τέσσαρας στύλους ἀσήπτους κατακεχρυσωμένους ἐν χρυσίῳ· καὶ αἱ κεφαλίδες αὐτῶν χρυσαῖ, καὶ αἱ βάσεις αὐτῶν τέσσαρες ἀργυραῖ.
\vs{5}καὶ ἐποίησαν τὸ καταπέτασμα τῆς θύρας τῆς σκηνῆς τοῦ μαρτυρίου ἐξ ὑακίνθου, καὶ πορφύρας, καὶ κοκκίνου νενησμένου, καὶ βύσσου κεκλωσμένης, ἔργον ὑφαντὸντοῦ χερουβείμ·
\vs{6}καὶ τοὺς στύλους αὐτῶν πέντε, καὶ τοὺς κρίκους· καὶ τὰς κεφαλίδας αὐτῶν, καὶ τὰς ψαλίδας αὐτῶν κατεχρύσωσαν χρυσίῳ· καὶ αἱ βάσεις αὐτῶν πέντε χαλκαῖ.

\vs{7}Καὶ ἐποίησαν τὴν αὐλῆν τὰ πρὸς Λίβα, ἱστία τῆς αὐλῆς ἐκ βύσσου κεκλωσμένης ἑκατὸν ἐφʼ ἑκατόν·
\vs{8}καὶ οἱ στύλοι αὐτῶν εἴκοσι, καὶ αἱ βάσεις αὐτῶν εἴκοσι.
\vs{9}καὶ τὸ κλίτος τὸ πρὸς Βοῤῥᾶν, ἑκατὸν ἐφʼ ἑκατόν· καὶ τὸ κλίτος τὸ πρὸς Νότον, ἑκατὸν ἐφʼ ἑκατόν· καὶ οἱ στύλοι αὐτῶν εἴκοσι, καὶ αἱ βάσεις αὐτῶν εἴκοσι·
\vs{10}Καὶ τὸ κλίτος τὸ πρὸς θάλασσαν αὐλαῖαι πεντήκοντα πήχεων· στύλοι αὐτῶν δέκα, καὶ αἱ βάσεις αὐτῶν δέκα·
\vs{11}Καὶ τὸ κλίτος τὸ πρὸς ἀνατολὰς πεντήκοντα πήχεων ἱστία, πεντεαίδεκα πήχεων τὸ κατὰ νώτου·
\vs{12}καὶ οἱ στύλοι αὐτῶν τρεῖς, καὶ αἱ βάσεις αὐτῶν τρεῖς·
\vs{13}Καὶ ἐπὶ τοῦ νώτου τοῦ δευτέρου ἔνθεν καὶ ἔνθεν κατὰ τὴν πύλην τῆς αὐλῆς, αὐλαῖαι πεντεκαίδεκα πήχεων· στύλοι αὐτῶν τρεῖς, καὶ αἱ βάσεις αὐτῶν τρεῖς·
\vs{14}πᾶσαι αἱ αὐλαῖαι τῆς σκηνῆς ἐκ βύσσου κεκλωσμένης.
\vs{15}Καὶ αἱ βάσεις τῶν στύλων αὐτῶν χαλκαῖ, καὶ αἱ ἀγκύλαι αὐτῶν ἀργυραῖ, καὶ αἱ κεφαλίδες αὐτῶν περιηργυρωμέναι ἀργυρίῳ, καὶ οἱ στύλοι περιηργυρωμένοι ἀργυρίῳ πάντες οἱ στύλοι τῆς αὐλῆς·
\vs{16}καὶ τὸ καταπέτασμα τῆς πύλης τῆς αὐλῆς ἔργον ποικιλτοῦ ἐξ ὑακίνθου, καὶ πορφύρας, καὶ κοκκίνου νενησμένου, καὶ βύσσου κεκλωσμένης· εἴκοσι πήχεων τὸ μῆκος, καὶ τὸ ὕψος καὶ τὸ εὖρος πέντε πήχεων ἐξισούμενον τοῖς ἱστίοις τῆς αὐλῆς·
\vs{17}καὶ οἱ στύλοι αὐτῶν τέσσαρες, καὶ αἱ βάσεις αὐτῶν τέσσαρες χαλκαῖ, καὶ αἱ ἀγκύλαι αὐτῶν ἀργυραῖ, καὶ αἱ κεφαλίδες αὐτῶν περιηργυρωμέναι ἀργυρίῳ.
\vs{18}Καὶ πάντες οἱ πάσσαλοι τῆς αὐλῆς κύκλῳ χαλκοῖ, καὶ αὐτοὶ περιηργυρωμένοι ἀργυρίῳ.
\vs{19}Καὶ αὕτη ἡ σύνταξις τῆς σκηνῆς τοῦ μαρτυρίου, καθὰ συνετάγη Μωυσῇ, τὴν λειτουργίαν εἶναι τῶν Λευιτῶν διὰ Ἰθάμαρ τοῦ υἱοῦ Ἀαρὼν τοῦ ἱερέως.

\vs{20}Καὶ Βεσελεὴλ ὁ τοῦ Οὐρείου, ἐκ φυλῆς Ἰούδα, ἐποίησε καθὰ συνέταξε Κύριος τῷ Μωυσῇ,
\vs{21}καὶ Ἐλιὰβ ὁ τοῦ Ἀχισαμὰχ ἐκ φυλῆς Δὰν, ὅς ἠρχιτεκτόνησε τὰ ὑφαντὰ καὶ τὰ ῥαφιδευτὰ καὶ ποικιλτικά, ὑφάναι τῷ κοκκίνῳ καὶ τῇ βύσσῳ.

\ch{38}
Καὶ ἐποίησε Βεσελεὴλ τὴν κιβωτόν,
\vs{2}καὶ κατεχρύσωσεν αὐτὴν χρωσίῳ καθαρῷ ἔσωθεν καὶ ἔξωθεν·
\vs{3}καὶ ἐχώνευσεν αὐτῇ τέσσαρας δακτυλίους χρυσοῦς· δύο ἐπὶ τὸ κλίτος τὸ ἓν, καὶ δύο ἐπὶ τὸ κλίτος τὸ δεύτερον,
\vs{4}εὐρεῖς τοῖς διωστῆρσιν, ὥστε αἴρειν αὐτὴν ἐν αὐτοῖς.
\vs{5}Καὶ ἐποίησε τὸ ἱλαστήριον ἐπάνωθεν τῆς κιβωτοῦ ἐκ χρυσίου καθαροῦ,
\vs{6}καὶ τοὺς δύο χερουβεὶμ χρυσοῦς·
\vs{7}χεροὺβ ἕνα ἐπὶ τὸ ἄκρον τοῦ ἱλαστηρίου τὸ ἓν, καὶ χεροὺβ ἕνα ἐπὶ τὸ ἄκρον τοῦ ἱλαστηρίου τὸ δεύτερον,
\vs{8}σκιάζοντα ταῖς πτέρυξιν αὐτῶν ἐπὶ τὸ ἱλαστήριον.
\vs{9}Καὶ ἐποίησε τὴν τράπεζαν τὴν προκειμένην ἐκ χρυσίου καθαροῦ,
\vs{10}καὶ ἐχώνευσεν αὐτῇ τέσσαρας δακτυλίους, δύο ἐπὶ τοῦ κλίτους τοῦ ἑνὸς, καὶ δύο ἐπὶ τοῦ κλίτους τοῦ δευτέρου, εὐρεῖς, ὥστε αἴρειν τοῖς διωστῆρσιν ἐν αὐτοῖς.
\vs{11}Καὶ τοὺς διωστῆρας τῆς κιβωτοῦ καὶ τῆς τραπέζης ἐποίησε, καὶ κατεχρύσωσεν αὐτοὺς χρυσίῳ.
\vs{12}Καὶ ἐποίησε τὰ σκεύη τῆς τραπέζης, τά τε τρυβλία, καὶ τὰς θυίσκας, καὶ τοὺς κυάθους, καὶ τὰ σπονδεῖα, ἐν οἷς σπείσει ἐν αὐτοῖς, χρυσᾶ.
\vs{13}Καὶ ἐποίησε τὴν λυχνίαν ἣ φωτίζει, χρυσῆν,
\vs{14}στερεὰν τὸν καυλόν, καὶ τοὺς καλαμίσκους ἐξ ἀμφοτέρων τῶν μερῶν αὐτῆς·
\vs{15}ἐκ τῶν καλαμίσκων αὐτῆς οἱ βλαστοὶ ἐκ τῶν καλαμίσκων αὐτῆς οἱ βλαστοὶ ἐξέχοντες· τρεῖς ἐκ τούτου, καὶ τρεῖς ἐκ τούτου, ἐξισούμενοι ἀλλήλοις.
\vs{16}Καὶ τὰ λαμπάδια αὐτῶν, ἅ ἐστιν ἐπὶ τῶν ἄκρων, καρυωτὰ ἐξ αὐτῶν· καὶ τὰ ἐνθέμια ἐξ αὐτῶν, ἵνα ὦσιν οἱ λύχνοι ἐπʼ αὐτῶν· καὶ τὸ ἐνθέμιον τὸ ἕβδομον, τὸ ἐπʼ ἄκρου τοῦ λαμπαδίου, ἐπὶ τῆς κορυφῆς ἄνωθεν, στερεὸν ὅλον χρυσοῦν·
\vs{17}Καὶ ἑπτὰ λύχνους ἐπʼ αὐτῆς χρυσοῦς, καὶ τὰς λαβίδας αὐτῆς χρυσᾶς, καὶ τὰς ἐπαρυστρίδας αὐτῶν χρυσᾶς.
\vs{18}Οὗτος περιηργύρωσε τοὺς στύλους, καὶ ἐχώνευσε τῷ στύλῳ δακτυλίους χρυσοῦς, καὶ ἐχρύσωσε τοὺς μοχλοὺς χρυσίῳ, καὶ κατεχρύσωσε τοὺς στύλους τοῦ καταπετάσματος χρυσίῳ, καὶ ἐποίησε τὰς ἀγκύλας χρυσᾶς.
\vs{19}Οὗτος ἐποίησε καὶ τοὺς κρίκους τῆς σκηνῆς χρυσοῦς, καὶ τοὺς κρίκους τῆς αὐλῆς, καὶ κρίκους εἰς τὸ ἐκτείνειν τὸ κατακάλυμμα ἄνωθεν χαλκοῦς·
\vs{20}Οὗτος ἐχώνευσε τὰς κεφαλίδας τὰς ἀργυρᾶς τῆς σκηνῆς, καὶ τὰς κεφαλίδας τὰς χαλκᾶς τῆς θύρας τῆς σκηνῆς, καὶ τὴν πύλην τῆς αὐλῆς· καὶ ἀγκύλας ἐποίησε τοῖς στύλοις ἀργυρᾶς, ἐπὶ τῶν στύλων οὗτος περιηργύρωσεν αὐτάς·
\vs{21}Οὗτος ἐποίησε τοὺς πασσάλους τῆς σκηνῆς, καὶ τοὺς πασσάλους τῆς αὐλῆς χαλκοῦς·
\vs{22}Οὗτος ἐποίησε τὸ θυσιαστήριον τὸ χαλκοῦν ἐκ τῶν πυρείων τῶν χαλκῶν, ἃ ἦσαν τοῖς ἀνδράσιν τοῖς καταστασιάσασι μετὰ τῆς Κορὲ συναγωγῆς·
\vs{23}Οὗτος ἐποίησε πάντα τὰ σκεύη τοῦ θυσιαστηρίου, καὶ τὸ πυρεῖον αὐτοῦ, καὶ τὴν βάσιν, καὶ τὰς φιάλας, καὶ τὰς κρεάγρας τὰς χαλκᾶς·
\vs{24}Οὗτος ἐποίησε θυσιαστηρίῳ παράθεμα, ἔργον δικτυωτὸν κάτωθεν τοῦ πυρείου ὑπὸ αὐτὸ ἕως τοῦ ἡμίσους αὐτοῦ· καὶ ἐπέθηκεν αὐτῷ τέσσαρας δακτυλίους ἐκ τῶν τεσσάρων μερῶν τοῦ παραθέματος τοῦ θυσιαστηρίου χαλκοῦς, εὐρεῖς τοῖς μοχλοῖς, ὥστε αἴρειν ἐν αὐτοῖς τὸ θυσιαστήριον·
\vs{25}Οὗτος ἐποίησε τὸ ἔλαιον τῆς χρίσεως τὸ ἅγιον, καὶ τὴν σύνθεσιν τοῦ θυμιάματος καθαρὸν ἔργον μυρεψοῦ·
\vs{26}Οὗτος ἐποίησε τὸν λουτῆρα τὸν χαλκοῦν, καὶ τὴν βάσιν αὐτοῦ χαλκῆν ἐκ τῶν κατόπτρων τῶν νηστευσασῶν, αἳ ἐνήστευσαν παρὰ τὰς θύρας τῆς σκηνῆς τοῦ μαρτυρίου, ἐν ᾗ ἡμέρᾳ ἔπηξεν αὐτήν.

\vs{27}Καὶ ἐποίησε τὸν λουτῆρα, ἵνα νίπτωνται ἐξ αὐτοῦ Μωυσῆς καὶ Ἀαρὼν καὶ οἱ υἱοὶ αὐτοῦ τὰς χεῖρας αὐτῶν καὶ τοὺς πόδας, εἰσπορευομένων αὐτῶν εἰς τὴν σκηνὴν τοῦ μαρτυρίου, ἢ ὅταν προσπορεύωνται πρὸς τὸ θυσιαστήριον λειτουργεῖν, ἐνίπτοντο ἐξ αὐτοῦ, καθάπερ συνέταξε Κύριος τῷ Μωυσῇ.

\ch{39}
Πᾶν τὸ χρυσίον, ὃ κατειργάσθη εἰς τὰ ἔργα κατὰ πᾶσαν τὴν ἐργασίαν τῶν ἁγίων, ἐγένετο χρυσίου τοῦ τῆς ἀπαρχῆς, ἐννέα καὶ εἴκοσι τάλαντα, καὶ ἑπτακόσιοι εἴκοσι σίκλοι κατὰ τὸν σίκλον τὸν ἅγιον·
\vs{2}Καὶ ἀργυρίου ἀφαίρεμα παρὰ τῶν ἐπεσκεμμένων ἀνδρῶν τῆς συναγωγῆς ἑκατὸν τάλαντα, καὶ χίλιοι ἑπτακόσιοι ἑβδομηκονταπέντε σίκλοι· δραχμὴ μία τῇ κεφαλῇ τὸ ἥμισυ τοῦ σίκλου, κατὰ τὸν σίκλον τὸν ἅγιον·
\vs{3}Πᾶς ὁ παραπορευόμενος τὴν ἐπίσκεψιν ἀπὸ εἰκοσαετοῦς καὶ ἐπάνω εἰς τὰς ἑξήκοντα μυριάδας, καὶ τρισχίλιοι πεντακόσιοι καὶ πεντήκοντα.
\vs{4}Καὶ ἐγενήθη τὰ ἑκατὸν τάλαντα τοῦ ἀργυρίου εἰς τὴν χώνευσιν τῶν ἑκατὸν κεφαλίδων τῆς σκηνῆς, καὶ εἰς τὰς κεφαλίδας τοῦ καταπετάσματος,
\vs{5}ἑκατὸν κεφαλίδες εἰς τὰ ἑκατὸν τάλαντα, τάλαντον τῇ κεφαλίδι·
\vs{6}Καὶ τοὺς χιλίους ἑπτακοσίους ἑβδομηκοντα πέντε σίκλους ἐποίησεν εἰς τὰς ἀγκύλας τοῖς στύλοις· καὶ κατεχρύσωσε τὰς κεφαλίδας αὐτῶν, καὶ κατεκόσμησεν αὐτούς.

\vs{7}Καὶ ὁ χαλκὸς τοῦ ἀφαιρέματος ἑβδομήκοντα τάλαντα, καὶ χίλιοι πεντακόσιοι σίκλοι·
\vs{8}Καὶ ἐποίησαν ἐξ αὐτον τὰς βάσεις τῆς θύρας τῆς σκηνῆς τοῦ μαρτυρίου,
\vs{9}καὶ τὰς βάσεις τῆς αὐλῆς κύκλῳ, καὶ τὰς βάσεις τῆς πύλης τῆς αὐλῆς, καὶ τοὺς πασσάλους τῆς σκηνῆς, καὶ τοὺς πασσάλους τῆς αὐλῆς κύκλῳ,
\vs{10}καὶ τὸ παράθεμα τὸ χαλκοῦν τοῦ θυσιαστηρίου, καὶ πάντα τὰ σκεύη τοῦ θυσιαστηρίου, καὶ πάντα τὰ ἐργαλεῖα τῆς σκηνῆς τοῦ μαρτυρίου·
\vs{11}Καὶ ἐποίησαν οἱ υἱοὶ Ἰσραὴλ, καθὰ συνέταξε Κύριος τῷ Μωυσῇ, οὕτως ἐποίησαν·
\vs{12}Τὸ δὲ λοιπὸν χρυσίον τοῦ ἀφαιρέματος ἐποίησαν σκεύη εἰς τὸ λειτουργεῖν ἐν αὐτοῖς ἔναντι Κυρίου·
\vs{13}Καὶ τὴν καταλειφθεῖσαν ὑάκινθον, καὶ πορφύραν, καὶ τὸ κόκκινον ἐποίησαν στολὰς λειτουργικὰς Ἀαρών, ὥστε λειτουργεῖν ἐν αὐταῖς ἐν τῷ ἁγίῳ·
\vs{14}Καὶ ἤνεγκαν τὰς στολὰς πρὸς Μωυσῆν, καὶ τὴν σκηνὴν, καὶ τὰ σκεύη αὐτῆς, τὰς βάσεις καὶ τοὺς μοχλοὺς αὐτῆς, καὶ τοὺς στύλους·
\vs{15}καὶ τὸ θυσιαστήριον, καὶ πάντα τὰ σκεύη αὐτοῦ.

\vs{16}Καὶ τὸ ἔλαιον τῆς χρίσεως, καὶ τὸ θυμίαμα τῆς συνθέσεως, καὶ τὴν λυχνίαν τὴν καθαρὰν,
\vs{17}καὶ τοὺς λύχνους αὐτῆς, λύχνους τῆς καύσεως, καὶ τὸ ἔλαιον τοῦ φωτός·
\vs{18}Καὶ τὴν τράπεζαν τῆς προθέσεως, καὶ πάντα τὰ σκεύη αὐτῆς· καὶ τοὺς ἄρτους τοὺς προκειμένους·
\vs{19}Καὶ τὰς στολὰς τοῦ ἁγίου, αἵ εἰσιν Ἀαρών, καὶ τὰς στολὰς τῶν υἱῶν αὐτοῦ, εἰς τὴν ἱερατείαν·
\vs{20}Καὶ τὰ ἱστία τῆς αὐλῆς, καὶ τοὺς στύλους· καὶ τὸ καταπέτασμα τῆς θύρας τῆς σκηνῆς, καὶ τῆς πύλης τῆς αὐλῆς·
\vs{21}Καὶ πάντα τὰ σκεύη τῆς σκηνῆς, καὶ πάντα τὰ ἐργαλεῖα αὐτῆς· καὶ τὰς διφθέρας δέρματα κριῶν ἠρυθροδανωμένα, καὶ τὰ καλύμματα ὑακίνθινα, καὶ τῶν λοιπῶν τὰ ἐπικαλύμματα· καὶ τοὺς πασσάλους, καὶ πάντα τὰ ἐργαλεῖα τὰ εἰς τὰ ἔργα τῆς σκηνῆς τοῦ μαρτυρίου·
\vs{22}Ὃσα συνέταξε Κύριος τῷ Μωυσῇ, οὕτως ἐποίησαν οἱ υἱοὶ Ἰσραὴλ πᾶσαν τὴν ἀποσκευήν·
\vs{23}Καὶ εἶδε Μωυσῆς πάντα τὰ ἔργα, καὶ ἦσαν πεποιηκότες αὐτὰ ὃν τρόπον συνέταξε Κύριος τῷ Μωυσῇ, οὕτως ἐποίησαν αὐτὰ, καὶ εὐλόγησεν αὐτοὺς Μωυσῆς.

\ch{40}
Καὶ ἐλάλησε Κύριος πρὸς Μωυσῆν, λέγων,
\vs{2}ἐν ἡμέρᾳ μιᾷ τοῦ μηνὸς τοῦ πρώτου νουμηνίᾳ, στήσεις τὴν σκηνὴν τοῦ μαρτυρίου.
\vs{3}Καὶ θήσεις τὴν κιβωτὸν τοῦ μαρτυρίου, καὶ σκεπάσεις τὴν κιβωτὸν τῷ καταπετάσματι.
\vs{4}Καὶ εἰσοίσεις τὴν τράπεζαν, καὶ προθήσεις τὴν πρόθεσιν αὐτῆς· καὶ εἰσοίσεις τὴν λυχνίαν, καὶ ἐπιθήσεις τοὺς λύχνους αὐτῆς.
\vs{5}Καὶ θήσεις τὸ θυσιαστήριον τὸ χρυσοῦν, εἰς τὸ θυμιᾷν ἐναντίον τῆς κιβωτοῦ· καὶ ἐπιθήσεις κάλυμμα καταπετάσματος ἐπὶ τὴν θύραν τῆς σκηνῆς τοῦ μαρτυρίου.
\vs{6}Καὶ τὸ θυσιαστήριον τῶν καρπωμάτων θήσεις παρὰ τὰς θύρας τῆς σκηνῆς τοῦ μαρτυρίου·
\vs{8}καὶ περιθήσεις τὴν σκηνήν, καὶ πάντα τὰ αὐτῆς ἁγιάσεις κύκλῳ.
\vs{9}Καὶ λήψῃ τὸ ἔλαιον τοῦ χρίσματος, καὶ χρίσεις τὴν σκηνὴν, καὶ πάντα τὰ ἐν αὐτῇ, καὶ ἁγιάσεις αὐτὴν, καὶ πάντα τὰ σκεύη αὐτῆς, καὶ ἔσται ἁγία.
\vs{10}Καὶ χρίσεις τὸ θυσιαστήριον τῶν καρπωμάτων, καὶ πάντα τὰ σκεύη αὐτοῦ· καὶ ἁγιάσεις τὸ θυσιαστήριον, καὶ ἔσται τὸ θυσιαστήριον ἅγιον τῶν ἁγίων.
\vs{12}Καὶ προσάξεις Ἀαρὼν καὶ τοὺς υἱοὺς αὐτοῦ ἐπὶ τὰς θύρας τῆς σκηνῆς τοῦ μαρτυρίου, καὶ λούσεις αὐτοὺς ὕδατι.
\vs{13}Καὶ ἐνδύσεις Ἀαρὼν τὰς στολὰς τὰς ἁγίας, καὶ χρίσεις αὐτὸν, καὶ ἁγιάσεις αὐτὸν, καὶ ἱερατεύει μοι.
\vs{14}Καὶ τοὺς υἱοὺς αὐτοῦ προσάξεις, καὶ ἐνδύσεις αὐτοὺς χιτῶνας.
\vs{15}Καὶ ἀλείψεις αὐτοὺς ὃν τρόπον ἤλειψας τὸν πατέρα αὐτῶν, καὶ ἱερατεύσουσί μοι· καὶ ἔσται, ὥστε εἶναι αὐτοῖς χρίσμα ἱερατείας εἰς τὸν αἰῶνα, εἰς τὰς γενεὰς αὐτῶν.
\vs{16}Καὶ ἐποίησε Μωυσῆς πάντα, ὅσα ἐνετείλατο αὐτῷ Κύριος, οὕτως ἐποίησε.

\vs{17}Καὶ ἐγένετο ἐν τῷ μηνὶ τῷ πρώτῳ, τῷ δευτέρῳ ἔτει, ἐκπορευομένων αὐτῶν ἐξ Αἰγύπτου, νουμηνίᾳ ἐστάθη ἡ σκηνή.
\vs{18}Καὶ ἔστησε Μωυσῆς τὴν σκηνὴν, καὶ ἐπέθηκε τὰς κεφαλιδας, καὶ διενέβαλε τοὺς μοχλοὺς, καὶ ἔστησε τοὺς στύλους.
\vs{19}Καὶ ἐξέτεινε τὰς αὐλαίας ἐπὶ τὴν σκηνὴν, καὶ ἐπέθηκε τὸ κατακάλυμμα τῆς σκηνῆς ἐπʼ αὐτὴν ἄνωθεν, καθὰ συνέταξε Κύριος τῷ Μωυσῇ.
\vs{20}Καὶ λαβὼν τὰ μαρτύρια ἐνέβαλεν εἰς τὴν κιβωτόν· καὶ ὑπέθηκε τοὺς διωστῆρας ὑπὸ τὴν κιβωτὸν,
\vs{21}καὶ εἰσήνεγκε τὴν κιβωτὸν εἰς τὴν σκηνὴν, καὶ ἐπέθηκε τὸ κατακάλυμμα τοῦ καταπετάσματος, καὶ ἐσκέπασε τὴν κιβωτὸν τοῦ μαρτυρίου, ὃν τρόπον συνέταξε Κύριος τῷ Μωυσῇ·
\vs{22}Καὶ ἐπέθηκε τὴν τράπεζαν εἰς τὴν σκηνὴν τοῦ μαρτυρίου, τὸ πρὸς Βοῤῥᾶν ἔξωθεν τοῦ καταπετάσματος τῆς σκηνῆς.
\vs{23}Καὶ προσέθηκεν ἐπʼ αὐτῆς ἄρτους τῆς προθέσεως ἔναντι Κυρίου, ὃν τρόπον συνέταξε Κύριος τῷ Μωυσῇ.
\vs{24}Καὶ ἔθηκε τὴν λυχνίαν εἰς τὴν σκηνὴν τοῦ μαρτυρίου, εἰς τὸ κλίτος τῆς σκηνῆς τὸ πρὸς Νότον.
\vs{25}Καὶ ἐπέθηκε τοὺς λύχνους αὐτῆς ἔναντι Κυρίου, ὃν τρόπον συνέταξε Κύριος τῷ Μωυσῇ.
\vs{26}Καὶ ἔθηκε τὸ θυσιαστήριον τὸ χρυσοῦν ἐν τῇ σκηνῇ τοῦ μαρτυρίου ἀπέναντι τοῦ καταπετάσματος,
\vs{27}καὶ ἐθυμίασεν ἐνʼ αὐτοῦ θυμίαμα τῆς συνθέσεως, καθάπερ συνέταξε Κύριος τῷ Μωυσῇ.
\vs{29}Καὶ τὸ θυσιαστήριον τῶν καρπωμάτων ἔθηκε παρὰ τὰς θύρας τῆς σκηνῆς.
\vs{33}Καὶ ἔστησε τὴν αὐλὴν κύκλῳ τῆς σκηνῆς, και τοῦ θυσιαστηρίου· καὶ συνετέλεσε Μωυσῆς πάντα τὰ ἔργα.

\vs{34}Καὶ ἐκάλυψεν ἡ νεφέλη τὴν σκηνὴν τοῦ μαρτυρίου· καὶ δόξης Κυρίου ἐπλήσθη ἡ σκηνή.
\vs{35}Καὶ οὐκ ἠδυνάσθη Μωυσῆς εἰσελθεῖν εἰς τὴν σκηνὴν τοῦ μαρτυρίου, ὅτι ἐπεσκίαζεν ἐπʼ αὐτὴν ἡ νεφέλη, καὶ δόξης Κυρίου ἐνεπλήσθη ἡ σκηνή.
\vs{36}Ἡνίκα δʼ ἂν ἀνέβη ἡ νεφέλη ἀπὸ τῆς σκηνῆς, ἀνεζεύγνυσαν οἱ υἱοὶ Ἰσραὴλ σὺν τῇ ἀπαρτίᾳ αὐτῶν.
\vs{37}Εἰ δὲ μὴ ἀνέβη ἡ νεφέλη, οὐκ ἀνεζεύγνυσαν ἕως ἡμέρας, ἧς ἀνέβη ἡ νεφέλη.
\vs{38}Νεφέλη γὰρ ἦν ἐπὶ τῆς σκηνῆς ἡμέρας, καὶ πῦρ ἦν ἐπʼ αὐτῆς νυκτὸς ἐναντίον παντὸς Ἰσραὴλ, ἐν πάσαις ταῖς ἀναζυγαῖς αὐτῶν.


\def\book{ΛΕΥΙΤΙΚΟΝ}
\biblebook{ΛΕΥΙΤΙΚΟΝ}


\lettrine[lines=2, loversize=0.2, nindent=0em, findent=.25em]{\textcolor{bookheadingcolor}{Κ}}{ΑΙ} ἀνεκάλεσε Μωυσῆν, καὶ ἐλάλησε Κύριος αὐτῷ ἐκ τῆς σκηνῆς τοῦ μαρτυρίου, λέγων,
\vs{2}λάλησον τοῖς υἱοῖς Ἰσραὴλ, καὶ ἐρεῖς πρὸς αὐτοὺς, ἄνθρωπος ἐξ ὑμῶν ἐὰν προσαγάγῃ δῶρα τῷ Κυρίῳ, ἀπὸ τῶν κτηνῶν καὶ ἀπὸ τῶν βοῶν καὶ ἀπὸ τῶν προβάτων προσοίσετε τὰ δῶρα ὑμῶν.
\vs{3}Ἐὰν ὁλοκαύτωμα τὸ δῶρον αὐτοῦ, ἐκ τῶν βοῶν ἄρσεν ἄμωμον προσάξει πρὸς τὴν θύραν τῆς σκηνῆς τοῦ μαρτυρίου, προσοίσει αὐτὸ δεκτὸν ἐναντίον Κυρίου.
\vs{4}Καὶ ἐπιθήσει τὴν χεῖρα ἐπὶ τὴν κεφαλὴν τοῦ καρπώματος δεκτὸν αὐτῷ, ἐξιλάσασθαι περὶ αὐτοῦ.
\vs{5}Καὶ σφάξουσι τὸν μόσχον ἔναντι Κυρίου· καὶ προσοίσουσιν οἱ υἱοὶ Ἀαρὼν οἱ ἱερεῖς τὸ αἷμα, καὶ προσχεοῦσι τὸ αἷμα ἐπὶ τὸ θυσιαστήριον κύκλῳ τὸ ἐπὶ τῶν θυρῶν τῆς σκηνῆς τοῦ μαρτυρίου·
\vs{6}καὶ ἐκδείραντες τὸ ὁλοκαύτωμα, μελιοῦσιν αὐτὸ κατὰ μέλη.
\vs{7}Καὶ ἐπιθήσουσιν οἱ υἱοὶ Ἀαρὼν οἱ ἱερεῖς πῦρ ἐπὶ τὸ θυσιαστήριον, καὶ ἐπιστοιβάσουσι ξύλα ἐπὶ τὸ πῦρ.
\vs{8}Καὶ ἐπιστοιβάσουσιν οἱ υἱοὶ Ἀαρὼν οἱ ἱερεῖς τὰ διχοτομήματα, καὶ τὴν κεφαλὴν, καὶ τὸ στέαρ ἐπὶ τὰ ξύλα τὰ ἐπὶ τοῦ πυρὸς τὰ ὄντα ἐπὶ τοῦ θυσιαστηρίου.
\vs{9}Τὰ δὲ ἐγκοίλια καὶ τοὺς πόδας πλυνοῦσιν ὕδατι· καὶ ἐπιθήσουσιν οἱ ἱερεῖς τὰ πάντα ἐπὶ τὸ θυσιαστήριον· κάρπωμά ἐστι θυσία ὀσμὴ εὐωδίας τῷ Κυρίῳ.
\vs{10}Ἐὰν δὲ ἀπὸ τῶν προβάτων τὸ δῶρον αὐτοῦ τῷ Κυρίῳ, ἀπό τε τῶν ἀρνῶν, καὶ τῶν ἐρίφων εἰς ὁλοκαυτώματα, ἄρσεν ἄμωμον προσάξει αὐτό.
\vs{11}Καὶ ἐπιθήσει τὴν χεῖρα ἐπὶ τὴν κεφαλὴν αὐτοῦ· καὶ σφάξουσιν αὐτὸ ἐκ πλαγίων τοῦ θυσιαστηρίου πρὸς βοῤῥᾶν ἔναντι Κυρίου· καὶ προσχεοῦσιν οἱ υἱοὶ Ἀαρὼν οἱ ἱερεῖς τὸ αἷμα αὐτοῦ ἐπὶ τὸ θυσιαστήριον κύκλῳ.
\vs{12}Καὶ διελουσιν αὐτὸ κατὰ μέλη, καὶ τὴν κεφαλὴν, καὶ τὸ στέαρ· καὶ ἐπιστοιβάσουσιν οἱ ἱερεῖς αὐτὰ ἐπὶ τὰ ξύλα τὰ ἐπὶ τοῦ πυρὸς τὰ ἐπὶ τοῦ θυσιαστηρίου.
\vs{13}Καὶ τὰ ἐγκοίλια, καὶ τοὺς πόδας πλυνοῦσιν ὕδατι· καὶ προσοίσει ὁ ἱερεὺς τὰ πάντα, καὶ ἐπιθήσει ἐπὶ τὸ θυσιαστήριον· κάρπωμά ἐστι θυσία ὀσμὴ εὐωδίας τῷ Κυρίῳ.
\vs{14}Ἐὰν δὲ ἀπὸ τῶν πετεινῶν κάρπωμα προσφέρει δῶρον αὐτοῦ τῷ κυρίῳ, καὶ προσοίσει ἀπὸ τῶν τρυγόνων, ἢ ἀπὸ τῶν περιστερῶν τὸ δῶρον αὐτοῦ.
\vs{15}Καὶ προσοίσει αὐτὸ ὁ ἱερεὺς πρὸς τὸ θυσιαστήριον, καὶ ἀποκνίσει τὴν κεφαλήν, καὶ ἐπιθήσει ὁ ἱερεὺς ἐπὶ τὸ θυσιαστήριον, καὶ στραγγιεῖ τὸ αἷμα πρὸς τὴν βάσιν τοῦ θυσιαστηρίου.
\vs{16}Καὶ ἀφελεῖ τὸν πρόλοβον σὺν τοῖς πτεροῖς, καὶ ἐκβαλεῖ αὐτὸ παρὰ τὸ θυσιαστήριον κατʼ ἀνατολὰς εἰς τὸν τόπον τῆς σποδοῦ·
\vs{17}Καὶ ἐκκλάσει αὐτὸ ἐκ τῶν πτερύγων, καὶ οὐ διελεῖ, καὶ ἐπιθήσει αὐτὸ ὁ ἱερεὺς ἐπὶ τὸ θυσιαστήριον ἐπὶ τὰ ξύλα τὰ ἐπὶ τοῦ πυρός· κάρπωμά ἐστι θυσία ὀσμὴ εὐωδίας τῷ Κυρίῳ.

\ch{2}
Ἐὰν δὲ ψυχὴ προσφέρῃ δῶρον θυσίαν τῷ Κυρίῳ, σεμίδαλις ἔσται τὸ δῶρον αὐτοῦ, καὶ ἐπιχεεῖ ἐπʼ αὐτὸ ἔλαιον, καὶ ἐπιθήσει ἐπʼ αὐτὸ λίβανον· θυσία ἐστί.
\vs{2}Καὶ οἴσει πρὸς τοὺς υἱοὺς Ἀαρὼν τοὺς ἱερεῖς· καὶ δραξάμενος ἀπʼ αὐτῆς πλήρη τὴν δράκα ἀπὸ τῆς σεμιδάλεως σὺν τῷ ἐλαίῳ, καὶ πάντα τὸν λίβανον αὐτῆς, καὶ ἐπιθήσει ὁ ἱερεὺς τὸ μνημόσυνον αὐτῆς ἐπὶ τὸ θυσιαστήριον· θυσία ὀσμὴ εὐωδίας τῷ Κυρίῳ.
\vs{3}Καὶ τὸ λοιπὸν ἀπὸ τῆς θυσίας Ἀαρὼν καὶ τοῖς υἱοῖς αὐτοῦ, ἅγιον τῶν ἁγίων ἀπὸ τῶν θυσιῶν Κυρίου.
\vs{4}Ἐὰν δὲ προσφέρῃ δῶρον θυσίαν πεπεμμένην ἐκ κλιβάνου δῶρον Κυρίῳ ἐκ σεμιδάλεως, ἄρτους ἀζύμους πεφυραμένους ἐν ἐλαίῳ, καὶ λάγανα ἄζυμα διακεχρισμένα ἐν ἐλαίῳ.
\vs{5}Ἐὰν δὲ θυσία ἀπὸ τηγάνου τὸ δῶρόν σου, σεμίδαλις πεφυραμένη ἐν ἐλαίῳ ἄζυμά ἐστι.
\vs{6}Καὶ διαθρύψεις αὐτὰ κλάσματα, καὶ ἐπιχεεῖς ἐπʼ αὐτὰ ἔλαιον· θυσία ἐστὶ Κυρίῳ.
\vs{7}Ἐὰν δὲ θυσία ἀπὸ ἐσχάρας τὸ δῶρόν σου, σεμίδαλις ἐν ἐλαίῳ ποιηθήσεται.
\vs{8}Καὶ προσοίσει τὴν θυσίαν ἣν ἂν ποιήσῃ ἐκ τούτων τῷ Κυρίῳ, καὶ προσοίσει πρὸς τὸν ἱερέα.
\vs{9}Καὶ προσεγγίσας πρὸς τὸ θυσιαστήριον, ἀφελεῖ ὁ ἱερεὺς ἀπὸ τῆς θυσίας τὸ μνημόσυνον αὐτῆς, καὶ ἐπιθήσει ὁ ἱερεὺς ἐπὶ τὸ θυσιαστήριον, κάρπωμα· ὀσμὴ εὐωδίας Κυρίῳ.
\vs{10}Τὸ δὲ καταλειφθὲν ἀπὸ τῆς θυσίας, Ἀαρὼν καὶ τοῖς υἱοῖς αὐτοῦ, ἅγια τῶν ἁγίων ἀπὸ τῶν καρπωμάτων Κυρίου.
\vs{11}Πᾶσαν θυσίαν, ἣν ἂν προσφέρητε Κυρίῳ, οὐ ποιήσετε ζυμωτόν· πᾶσαν γὰρ ζύμην, καὶ πᾶν μέλι οὐ προσοίσετε ἀπʼ αὐτοῦ, καρπῶσαι Κυρίῳ δῶρον.
\vs{12}Ἀπαρχῆς προσοίσετε αὐτὰ Κυρίῳ, ἐπὶ δὲ τὸ θυσιαστήριον οὐκ ἀναβιβασθήσεται εἰς ὀσμὴν εὐωδίας Κυρίῳ.
\vs{13}Καὶ πᾶν δῶρον θυσίας ὑμῶν ἁλὶ ἁλισθήσεται· οὐ διαπαύσατε ἅλας διαθήκης Κυρίου ἀπὸ θυσιασμάτων ὑμῶν· ἐπὶ παντὸς δώρου ὑμῶν προσοίσετε Κυρίῳ τῷ Θεῷ ὑμῶν ἅλας.
\vs{14}Ἐὰν δὲ προσφέρῃς θυσίαν πρωτογεννημάτων τῷ Κυρίῳ, νέα πεφρυγμένα χίδρα ἐρικτὰ τῷ Κυρίῳ· καὶ προσοίσεις τὴν θυσίαν τῶν πρωτογεννημάτων.
\vs{15}Καὶ ἐπιχεεῖς ἐπʼ αὐτὴν ἔλαιον, καὶ ἐπιθήσεις ἐπʼ αὐτὴν λίβανον· θυσία ἐστί.
\vs{16}Καὶ ἀνοίσει ὁ ἱερεὺς τὸ μνημόσυνον αὐτῆς ἀπὸ τῶν χίδρων σὺν τῷ ἐλαίῳ, καὶ πάντα τὸν λίβανον αὐτῆς· κάρπωμά ἐστι Κυρίῳ.

\ch{3}
Ἐὰν δὲ θυσία σωτηρίου τὸ δῶρον αὐτοῦ τῷ Κυρίῳ, ἐὰν μὲν ἐκ τῶν βοῶν αὐτὸ προσαγάγῃ, ἐάν τω ἄρσεν, ἐάν τε θῆλυ, ἄμωμον προσάξει αὐτὸ ἔναντι Κυρίου.
\vs{2}Καὶ ἐπιθήσει τὰς χεῖρας ἐπὶ τὴν κεφαλὴν τοῦ δώρου, καὶ σφάξει αὐτὸ ἐναντίον Κυρίου παρὰ τὰς θύρας τῆς σκηνῆς τοῦ μαρτυρίου· καὶ προσχεοῦσιν οἱ υἱοὶ Ἀαρὼν οἱ ἱερεῖς τὸ αἷμα ἐπὶ τὸ θυσιαστήριον τῶν ὁλοκαυτωμάτων κύκλῳ.
\vs{3}Καὶ προσάξουσιν ἀπὸ τῆς θυσίας τοῦ σωτηρίου κάρπωμα Κυρίῳ, τὸ στέαρ τὸ κατακαλύπτον τὴν κοιλίαν, καὶ πᾶν τὸ στέαρ τὸ ἐπὶ τῆς κοιλίας.
\vs{4}Καὶ τοὺς δύο νεφροὺς, καὶ τὸ στέαρ τὸ ἐπʼ αὐτῶν, τὸ ἐπὶ τῶν μηρίων, καὶ τὸν λοβὸν τὸν ἐπὶ τοῦ ἥπατος σὺν τοῖς νεφροῖς περιελεῖ.
\vs{5}Καὶ ἀνοίσουσιν αὐτὰ οἱ υἱοὶ Ἀαρὼν οἱ ἱερεῖς ἐπὶ τὸ θυσιαστήριον ἐπὶ τὰ ὁλοκαυτώματα ἐπὶ τὰ ξύλα, τὰ ἐπὶ τοῦ πυρὸς ἐπὶ τοῦ θυσιαστηρίου· κάρπωμα ὀσμὴ εὐωδίας Κυρίῳ.
\vs{6}Ἐὰν δὲ ἀπὸ τῶν προβάτων τὸ δῶρον αὐτοῦ θυσία σωτηρίου τῷ Κυρίῳ, ἄρσεν ἢ θῆλυ, ἄμωμον προσοίσει αὐτό.
\vs{7}Ἐὰν ἄρνα προσαγάγῃ τὸ δῶρον αὐτοῦ, προσάξει αὐτὸ ἔναντι Κυρίου.
\vs{8}Καὶ ἐπιθήσει τὰς χεῖρας ἐπὶ τὴν κεφαλὴν τοῦ δώρου αὐτοῦ, καὶ σφάξει αὐτὸ παρὰ τὰς θύρας τῆς σκηνῆς τοῦ μαρτυρίου· καὶ προσχεοῦσιν οἱ υἱοὶ Ἀαρὼν οἱ ἱερεῖς τὸ αἷμα ἐπὶ τὸ θυσιαστήριον κύκλῳ.
\vs{9}Καὶ προσοίσει ἀπὸ τῆς θυσίας τοῦ σωτηρίου κάρπωμα τῷ Κυρίῷ· τὸ στέαρ καὶ τὴν ὀσφὺν ἄμωμον σὺν ταῖς ψόαις περιελεῖ αὐτό· καὶ πᾶν τὸ στέαρ τὸ κατακαλύπτον τὴν κοιλίαν, καὶ πᾶν τὸ στέαρ τὸ ἐπὶ τῆς κοιλίας.
\vs{10}Καὶ ἀμφοτέρους τοὺς νεφροὺς, καὶ τὸ στέαρ τὸ ἐπʼ αὐτῶν, τὸ ἐπὶ τῶν μηρίων, καὶ τὸν λοβὸν τὸν ἐπὶ τοῦ ἥπατος σὺν τοῖς νεφροῖς περιελὼν,
\vs{11}ἀνοίσει ὁ ἱερεὺς ἐπὶ τὸ θυσιαστήριον· ὀσμὴ εὐωδίας κάρπωμα Κυρίῳ.

\vs{12}Ἐὰν δὲ ἀπὸ τῶν αἰγῶν τὸ δῶρον αὐτοῦ, καὶ προσάξει ἔναντι Κυρίου.
\vs{13}Καὶ ἐπιθήσει τὰς χεῖρας ἐπὶ τὴν κεφαλὴν αὐτοῦ, καὶ σφάξουσιν αὐτὸ ἔναντι Κυρίου παρὰ τὰς θύρας τῆς σκηνῆς τοῦ μαρτυρίου· καὶ προσχεοῦσιν οἱ υἱοὶ Ἀαρὼν οἱ ἱερεῖς τὸ αἷμα ἐπὶ τὸ θυσιαστήριον κύκλῳ.
\vs{14}Καὶ ἀνοίσει ἀπʼ αὐτοῦ κάρπωμα Κυρίῳ τὸ στέαρ τὸ κατακαλύπτον τὴν κοιλίαν, καὶ πᾶν τὸ στέαρ τὸ ἐπὶ τῆς κοιλίας.
\vs{15}Καὶ ἀμφοτέρους τοὺς νεφροὺς, καὶ πᾶν τὸ στέαρ τὸ ἐπʼ αὐτῶν, τὸ ἐπὶ τῶν μηρίων, καὶ τὸν λοβὸν τοῦ ἥπατος σὺν τοῖς νεφροῖς περιελεῖ,
\vs{16}καὶ ἀνοίσει ὁ ἱερεὺς ἐπὶ τὸ θυσιαστήριον· κάρπωμα ὀσμὴ εὐωδίας τῷ Κυρίῳ· πᾶν τὸ στέαρ τῷ Κυρίῳ.
\vs{17}Νόμιμον εἰς τὸν αἰῶνα εἰς τὰς γενεὰς ὑμῶν, ἐν πάσῃ κατοικίᾳ ὑμῶν· πᾶν στέαρ καὶ πᾶν αἷμα οὐκ ἔδεσθε.

\ch{4}
Καὶ ἐλάλησε Κύριος πρὸς Μωυσῆν, λέγων,
\vs{2}Λάλησον πρὸς τοὺς υἱοὺς Ἰσραὴλ, λέγων, ψυχὴ ἐὰν ἁμάρτῃ ἔναντι Κυρίου ἀκουσίως ἀπὸ πάντων τῶν προσταγμάτων Κυρίου, ὧν οὐ δεῖ ποιεῖν, καὶ ποιήσῃ ἕν τι ἀπʼ αὐτῶν·
\vs{3}Ἐὰν μὲν ὁ ἀρχιερεὺς ὁ κεχρισμένος ἁμάρτῃ τοῦ τὸν λαὸν ἁμαρτεῖν, καὶ προσάξει περὶ τῆς ἁμαρτίας αὐτοῦ, ἧς ἥμαρτε, μόσχον ἐκ βοῶν ἄμωμον τῷ Κυρίῳ περὶ τῆς ἁμαρτίας.
\vs{4}Καὶ προσάξει τὸν μόσχον παρὰ τὴν θύραν τῆς σκηνῆς τοῦ μαρτυρίου ἔναντι Κυρίου, καὶ ἐπιθήσει τὴν χεῖρα αὐτοῦ ἐπὶ τὴν κεφαλὴν τοῦ μόσχου ἔναντι Κυρίου, καὶ σφάξει τὸν μόσχον ἐνώπιον Κυρίου.
\vs{5}Καὶ λαβὼν ὁ ἱερεὺς ὁ χριστὸς ὁ τετελειωμένος τὰς χεῖρας ἀπὸ τοῦ αἵματος τοῦ μόσχου, καὶ εἰσοίσει αὐτὸ εἰς τὴν σκηνὴν τοῦ μαρτυρίου.
\vs{6}Καὶ βάψει ὁ ἱερεὺς τὸν δάκτυλον εἰς τὸ αἷμα, καὶ προσρανεῖ ἀπὸ τοῦ αἵματος ἑπτάκις ἔναντι Κυρίου, κατὰ τὸ καταπέτασμα τὸ ἅγιον.
\vs{7}Καὶ ἐπιθήσει ὁ ἱερεὺς ἀπὸ τοῦ αἵματος τοῦ μόσχου ἐπὶ τὰ κέρατα τοῦ θυσιαστηρίου τοῦ θυμιάματος τῆς συνθέσεως τοῦ ἐναντίον Κυρίου, ὅ ἐστιν ἐν τῇ σκηνῇ τοῦ μαρτυρίου· καὶ πᾶν τὸ αἷμα τοῦ μόσχου ἐκχεεῖ παρὰ τὴν βάσιν τοῦ θυσιαστηρίου τῶν ὁλοκαυτωμάτων, ὅ ἐστι παρὰ τὰς θύρας τῆς σκηνῆς τοῦ μαρτυρίου.
\vs{8}Καὶ πὰν τὸ στέαρ τοῦ μόσχου τοῦ τῆς ἁμαρτίας περιελεῖ ἀπʼ αὐτοῦ, τὸ στέαρ τὸ κατακαλύπτον τὰ ἐνδόσθια, καὶ πᾶν τὸ στέαρ τὸ ἐπὶ τῶν ἐνδοσθίων,
\vs{9}καὶ τοὺς δύο νεφροὺς, καὶ τὸ στέαρ τὸ ἐπʼ αὐτῶν, ὅ ἐστιν ἐπὶ τῶν μηρίων, καὶ τὸν λοβὸν τὸν ἐπὶ τοῦ ἥπατος σὺν τοῖς νεφροῖς περιελεῖ αὐτό,
\vs{10}ὃν τρόπον ἀφαιρεῖται αὐτὸ ἀπὸ τοῦ μόσχου τοῦ τῆς θυσίας τοῦ σωτηρίου, καὶ ἀνοίσει ὁ ἱερεὺς ἐπὶ τὸ θυσιαστήριον τῆς καρπώσεως.
\vs{11}Καὶ τὸ δέρμα τοῦ μόσχου, καὶ πᾶσαν αὐτοῦ τὴν σάρκα σὺν τῇ κεφαλῇ καὶ τοῖς ἀκρωτηρίοις καὶ τῇ κοιλίᾳ καὶ τῇ κόπρῳ·
\vs{12}καὶ ἐξοίσουσιν ὅλον τὸν μόσχον ἔξω τῆς παρεμβολῆς εἰς τόπον καθαρὸν, οὗ ἐκχεοῦσι τὴν σποδιὰν, καὶ κατακαύσουσιν αὐτὸν ἐπὶ ξύλων ἐν πυρί· ἐπὶ τῆς ἐκχύσεως τῆς σποδιᾶς καυθήσεται.

\vs{13}Ἐὰν δὲ πᾶσα συναγωγὴ Ἰσραὴλ ἀγνοήσῃ ἀκουσίως, καὶ λάθῃ ῥῆμα ἐξ ὀφθαλμῶν τῆς συναγωγῆς, καὶ ποιήσωσι μίαν ἀπὸ πασῶν τῶν ἐντολῶν Κυρίου, ἣ οὐ ποιηθήσεται, καὶ πλήμμελήσωσι,
\vs{14}καὶ γνωσθῇ αὐτοῖς ἡ ἁμαρτία, ἣν ἥμαρτον ἐν αὐτῇ, καὶ προσάξει ἡ συναγωγὴ μόσχον ἐκ βοῶν ἄμωμον περὶ τῆς ἁμαρτίας, καὶ προσάξει αὐτὸν παρὰ τὰς θύρας τῆς σκηνῆς τοῦ μαρτυρίου.
\vs{15}Καὶ ἐπιθήσουσιν οἱ πρεσβύτεροι τῆς συναγωγῆς τὰς χείρας αὐτῶν ἐπὶ τὴν κεφαλὴν τοῦ μόσχου ἔναντι Κυρίου, καὶ σφάξουσιν τὸν μόσχον ἔναντι Κυρίου.
\vs{16}Καὶ εἰσοίσει ὁ ἱερεὺς ὁ χριστὸς ἀπὸ τοῦ αἵματος τοὺ μόσχου εἰς τὴν σκηνὴν τοῦ μαρτυρίου.
\vs{17}Καὶ βάψει ὁ ἱερεὺς τὸν δάκτυλον ἀπὸ τοῦ αἵματος τοῦ μόσχου, καὶ ῥανεῖ ἑπτάκις ἔναντι Κυρίου, κατενώπιον τοῦ καταπετάσματος τοῦ ἁγίου.
\vs{18}Καὶ ἀπὸ τοῦ αἵματος ἐπιθήσει ὁ ἱερεὺς ἐπὶ τὰ κέρατα τοῦ θυσιαστηρίου τῶν θυμιαμάτων τῆς συνθέσεως, ὅ ἐστιν ἐνώπιον Κυρίου, ὅ ἐστιν ἐν τῇ σκηνῇ τοῦ μαρτυρίου· καὶ τὸ πᾶν αἷμα ἐκχεεῖ πρὸς τὴν βάσιν τοῦ θυσιαστηρίου τῶν καρπώσεων, τοῦ πρὸς τῇ θύρᾳ τῆς σκηνῆς τοῦ μαρτυρίου.
\vs{19}Καὶ τὸ πᾶν στέαρ περιελεῖ ἀπʼ αὐτοῦ, καὶ ἀνοίσει ἐπὶ τὸ θυσιαστήριον.
\vs{20}Καὶ ποιήσει τὸν μόσχον, ὃν τρόπον ἐποίησε τὸν μόσχον τὸν τῆς ἁμαρτίας, οὕτω ποιηθήσεται· καὶ ἐξιλάσεται περὶ αὐτῶν ὁ ἱερεὺς, καὶ ἀφεθήσεται αὐτοῖς ἡ ἁμαρτία.
\vs{21}Καὶ ἐξοίσουσι τὸν μόσχον ὅλον ἔξω τῆς παρεμβολῆς, καὶ κατακαύσουσι τὸν μόσχον, ὃν τρόπον κατέκαυσαν τὸν μόσχον τὸν πρότερον· ἁμαρτία συναγωγῆς ἐστιν.

\vs{22}Ἐὰν δὲ ὁ ἄρχων ἁμάρτῃ, καὶ ποιήσῃ μίαν ἀπὸ πασῶν τῶν ἐντολῶν Κυρίου τοῦ Θεοῦ αὐτοῦ, ἣ οὐ ποιηθήσεται, ἀκουσίως, καὶ ἁμάρτῃ καὶ πλημμελήσῃ,
\vs{23}καὶ γνωσθῇ αὐτῷ ἡ ἁμαρτία, ἣν ἥμαρτεν ἐν αὐτῇ, καὶ προσοίσει τὸ δῶρον αὐτοῦ χίμαρον ἐξ αἰγῶν, ἄρσεν ἄμωμον.
\vs{24}Καὶ ἐπιθήσει τὴν χεῖρα ἐπὶ τὴν κεφαλὴν τοῦ χιμάρου· καὶ σφάξουσιν αὐτὸν ἐν τόπῳ οὗ σφάζουσι τὰ ὁλοκαυτώματα ἐνώπιον Κυρίου· ἁμαρτία ἐστί.
\vs{25}Καὶ ἐπιθήσει ὁ ἱερεὺς ἀπὸ τοῦ αἵματος τοῦ τῆς ἁμαρτίας τῷ δακτύλῳ ἐπὶ τὰ κέρατα τοῦ θυσιαστηρίου τῶν ὁλοκαυτωμάτων· καὶ τὸ πᾶν αἷμα αὐτοῦ ἐκχεεῖ παρὰ τὴν βάσιν τοῦ θυσιαστηρίου τῶν ὁλοκαυτωμάτων.
\vs{26}Καὶ τὸ πᾶν στέαρ αὐτοῦ ἀνοίσει ἐπὶ τὸ θυσιαστήριον, ὥσπερ τὸ στέαρ θυσίας σωτηρίου· καὶ ἐξιλάσεται περὶ αὐτοῦ ὁ ἱερεὺς ἀπὸ τῆς ἁμαρτίας αὐτοῦ, καὶ ἀφεθήσεται αὐτῷ.

\vs{27}Ἐὰν δὲ ψυχὴ μία ἁμάρτῃ ἀκουσίως ἐκ τοῦ λαοῦ τῆς γῆς, ἐν τῷ ποιῆσαι μίαν ἀπὸ πασῶν τῶν ἐντολῶν Κυρίου, ἣ οὐ ποιηθήσεται, καὶ πλημμελήσῃ·
\vs{28}καὶ γνωσθῇ αὐτῷ ἡ ἁμαρτία, ἣν ἥμαρτεν ἐν αὐτῇ, καὶ οἴσει χίμαιραν ἐξ αἰγῶν, θήλειαν ἄμωμον οἴσει περὶ τῆς ἁμαρτίας, ἧς ἥμαρτε.
\vs{29}Καὶ ἐπιθήσει τὴν χεῖρα ἐπὶ τὴν κεφαλὴν τοῦ ἁμαρτήματος αὐτοῦ· καὶ σφάξουσιν τὴν χίμαιραν τὴν τῆς ἁμαρτίας ἐν τῷ τόπῳ, οὗ σφάζουσι τὰ ὁλοκαυτώματα.
\vs{30}Καὶ λήψεται ὁ ἱερεὺς ἀπὸ τοῦ αἵματος αὐτῆς τῷ δακτύλῳ, καὶ ἐπιθήσει ἐπὶ τὰ κέρατα τοῦ θυσιαστηρίου τῶν ὁλοκαυτωμάτων· καὶ πᾶν τὸ αἷμα αὐτῆς ἐκχεεῖ παρὰ τὴν βάσιν τοῦ θυσιαστηρίου.
\vs{31}Καὶ πᾶν τὸ στέαρ περιελεῖ, ὃν τρόπον περιαιρεῖται στέαρ ἀπὸ θυσίας σωτηρίου· καὶ ἀνοίσει ὁ ἱερεὺς ἐπὶ τὸ θυσιαστήριον εἰς ὀσμὴν εὐωδίας Κυρίῳ· καὶ ἐξιλάσεται περὶ αὐτοῦ ὁ ἱερεὺς, καὶ ἀφεθήσεται αὐτῷ.

\vs{32}Εὰν δὲ πρόβατον προσενέγκῃ τὸ δῶρον αὐτοῦ περὶ τῆς ἁμαρτίας, θῆλυ ἄμωμον προσοίσει αὐτό.
\vs{33}Καὶ ἐπιθήσει τὴν χεῖρα ἐπὶ τῆν κεφαλὴν τοῦ τῆς ἁμαρτίας· καὶ σφάξουσιν αὐτὸ ἐν τόπῳ, οὗ σφάζουσι τὰ ὁλοκαυτώματα.
\vs{34}Καὶ λαβὼν ὁ ἱερεὺς ἀπὸ τοῦ αἵματος τοῦ τῆς ἁμαρτίας τῷ δακτύλῳ, ἐπιθήσει ἐπὶ τὰ κέρατα τοῦ θυσιαστηρίου τῆς ὁλοκαρπώσεως· καὶ πᾶν αὐτοῦ τὸ αἷμα ἐκχεεῖ παρὰ τὴν βάσιν τοῦ θυσιαστηρίου τῆς ὁλοκαυτώσεως.
\vs{35}Καὶ πᾶν αὐτοῦ τὸ στέαρ περιελεῖ, ὃν τρόπον περιαιρεῖται στέαρ προβάτου ἐκ τῆς θυσίας τοῦ σωτηρίου· καὶ ἐπιθήσει αὐτὸ ὁ ἱερεὺς ἐπὶ τὸ θυσιαστήριον ἐπὶ τὸ ὁλοκαύτωμα Κυρίου· καὶ ἐξιλάσεται περὶ αὐτοῦ ὁ ἱερεὺς περὶ τῆς ἁμαρτίας ἧς ἥμαρτε, καὶ ἀφεθήσεται αὐτῷ.

\ch{5}
Ἐὰν δὲ ψυχὴ ἁμάρτῃ, καὶ ἀκούσῃ φωνὴν ὁρκισμοῦ, καὶ οὗτος μάρτυς ἢ ἑώρακεν ἢ σύνοιδεν, ἐὰν μὴ ἀπαγγείλῃ, λήψεται τὴν ἁμαρτίαν.
\vs{2}Ἡ ψυχὴ ἐκείνη ἥτις ἐὰν ἅψηται παντὸς πράγματος ἀκαθάρτου, ἢ θνησιμαίου, ἢ θηριαλώτου ἀκαθάρτου, ἢ τῶν θνησιμαίων βδελυγμάτων τῶν ἀκαθάρτων, ἢ τῶν θνησιμαίων κτηνῶν τῶν ἀκαθάρτων,
\vs{3}ἢ ἅψηται ἀπὸ ἀκαθαρσίας ἀνθρώπου, ἀπὸ πάσης ἀκαθαρσίας αὐτοῦ, ἧς ἂν ἁψάμενος μιανθῇ καὶ ἔλαθεν αὐτόν, μετὰ τοῦτο δὲ γνῷ, καὶ πλημμελήσῃ.
\vs{4}Ἡ ψυχὴ ἡ ἄνομος, ἡ διαστέλλουσα τοῖς χείλεσι κακοποιῆσαι ἢ καλῶς ποιῆσαι κατὰ πάντα ὅσα ἐὰν διαστείλῃ ὁ ἄνθρωπος μεθʼ ὅρκου, καὶ λάθῃ αὐτὸν πρὸ ὀφθαλμῶν, καὶ οὗτος γνῷ, καὶ ἁμάρτῃ ἕν τι τούτων.
\vs{5}Καὶ ἐξαγορεύσει τὴν ἁμαρτίαν περὶ ὧν ἡμάρτηκε κατʼ αὐτῆς.
\vs{6}Καὶ οἴσει περὶ ὧν ἐπλημμέλησε Κυρίῳ, περὶ τῆς ἁμαρτίας ἧς ἥμαρτε, θῆλυ ἀπὸ τῶν προβάτων ἀμνάδα, ἢ χίμαιραν ἐξ αἰγῶν, περὶ ἁμαρτίας· καὶ ἐξιλάσεται περὶ αὐτοῦ ὁ ἱερεὺς περὶ τῆς ἁμαρτίας αὐτοῦ, ἧς ἥμαρτε, καὶ ἀφεθήσεται αὐτῷ ἡ ἁμαρτία.
\vs{7}Ἐὰν δὲ μὴ ἰσχύῃ ἡ χεὶρ αὐτοῦ τὸ ἱκανὸν εἰς τὸ πρόβατον, οἴσει περὶ τῆς ἁμαρτίας αὐτοῦ, ἧς ἥμαρτε, δύο τρυγόνας, ἢ δύο νοσσοὺς περιστερῶν Κυρίῳ, ἕνα περὶ ἁμαρτίας, καὶ ἕνα εἰς ὁλοκαύτωμα.
\vs{8}Καὶ οἴσει αὐτὰ πρὸς τὸν ἱερέα· καὶ προσάξει ὁ ἱερεὺς τὸ περὶ τῆς ἁμαρτίας πρότερον· καὶ ἀποκνίσει ὁ ἱερεὺς τὴν κεφαλὴν αὐτοῦ ἀπὸ τοῦ σφονδύλου, καὶ οὐ διελεῖ.
\vs{9}Καὶ ῥανεῖ ἀπὸ τοῦ αἵματος τοῦ περὶ τῆς ἁμαρτίας ἐπὶ τὸν τοῖχον τοῦ θυσιαστηρίου· τὸ δὲ κατάλοιπον τοῦ αἵματος καταστραγγιεῖ ἐπὶ τὴν βάσιν τοῦ θυσιαστηρίου· ἁμαρτία γάρ ἐστι·
\vs{10}Καὶ τὸ δεύτερον ποιήσει ὁλοκάρπωμα, ὡς καθήκει· καὶ ἐξιλάσεται ὁ ἱερεὺς περὶ τῆς ἁμαρτίας αὐτοῦ, ἧς ἥμαρτε, καὶ ἀφεθήσεται αὐτῷ.

\vs{11}Ἐὰν δὲ μὴ εὑρίσκῃ ἡ χεὶρ αὐτοῦ ζεῦγος τρυγόνων, ἢ δύο νοσσούς περιστερῶν, καὶ οἴσει τὸ δῶρον αὐτοῦ, περὶ οὗ ἥμαρτε, τὸ δέκατον τοῦ οἰφὶ σεμιδάλεως περὶ ἁμαρτίας· οὐκ ἐπιχεεῖ ἐπʼ αὐτὸ ἔλαιον, οὐδὲ ἐπιθήσει ἐπʼ αὐτῷ λίβανον, ὅτι περὶ ἁμαρτίας ἐστί.
\vs{12}Καὶ οἴσει αὐτὸ πρὸς τὸν ἱερέα· καὶ δραξάμενος ὁ ἱερεὺς ἀπʼ αὐτῆς πλήρη τὴν δράκα, τὸ μνημόσυνον αὐτῆς ἐπιθήσει ἐπὶ τὸ θυσιαστήριον τῶν ὁλοκαυτωμάτων Κυρίῳ· ἁμαρτία ἐστί.
\vs{13}Καὶ ἐξιλάσεται περὶ αὐτοῦ ὁ ἱερεὺς περὶ τῆς ἁμαρτίας αὐτοῦ, ἧς ἥμαρτεν ἀφʼ ἑνὸς τούτων, καὶ ἀφεθήσεται αὐτῷ. τὸ δὲ καταλειφθὲν ἔσται τῷ ἱερεῖ, ὡς θυσία τῆς σεμιδάλεως.

\vs{14}Καὶ ἐλάλησε Κύριος πρὸς Μωυσῆν, λέγων,
\vs{15}ψυχὴ ἣ ἂν λάθῃ αὐτὸν λήθῃ, καὶ ἁμάρτῃ ἀκουσίως ἀπὸ τῶν ἁγίων Κυρίου, καὶ οἴσει τῆς πλημμελείας αὐτοῦ τῷ Κυρίῳ κριὸν ἄμωμον ἐκ τῶν προβάτων, τιμῆς ἀργυρίου σίκλων, τῷ σίκλῳ τῶν ἁγίων, περὶ οὗ ἐπλημμέλησε.
\vs{16}Καὶ ὃ ἥμαρτεν ἀπὸ τῶν ἁγίων ἀποτίσει αὐτὸ, καὶ τὸ ἐπίπεμπτον προσθήσει ἐπʼ αὐτὸ, καὶ δώσει αὐτὸ τῷ ἱερεῖ· καὶ ὁ ἱερεὺς ἐξιλάσεται περὶ αὐτοῦ ἐν τῷ κριῷ τῆς πλημμελείας, καὶ ἀφεθήσεται αὐτῷ.
\vs{17}Καὶ ἡ ψυχὴ ἣ ἂν ἁμάρτῃ, καὶ ποιήσει μίαν ἀπὸ πασῶν τῶν ἐντολῶν Κυρίου, ὧν οὐ δεῖ ποιεῖν, καὶ οὐκ ἔγνω, καὶ πλημμελήσῃ, καὶ λάβῃ τὴν ἁμαρτίαν,
\vs{18}καὶ οἴσει κριὸν ἄμωμον ἐκ τῶν προβάτων, τιμῆς ἀργυρίου εἰς πλημμέλειαν πρὸς τὸν ἱερέα· καὶ ἐξιλάσεται περὶ αὐτοῦ ὁ ἱερεὺς περὶ τῆς ἀγνοίας αὐτοῦ, ἧς ἠγνόησε, καὶ αὐτὸς οὐκ ᾔδει, καὶ ἀφεθήσεται αὐτῷ.
\vs{19}Ἐπλημμέλησε γὰρ πλημμελείᾳ ἔναντι Κυρίου.

\vs{20}Καὶ ἐλάλησε Κύριος πρὸς Μωυσῆν, λέγων,
\vs{21}ψυχὴ ἣ ἂν ἁμάρτῃ, καὶ παριδὼν παρίδῃ τὰς ἐντολὰς Κυρίου, καὶ ψεύσηται τὰ πρὸς τὸν πλησίον ἐν παραθήκῃ, ἢ περὶ κοινωνίας, ἢ περὶ ἁρπαγῆς, ἢ ἠδίκησέ τι τὸν πλησίον,
\vs{22}ἢ εὗρεν ἀπωλίαν, καὶ ψεύσηται περὶ αὐτῆς, καὶ ὀμόσῃ ἀδίκως περὶ ἑνὸς ἀπὸ πάντων, ὧν ἐὰν ποιήσῃ ὁ ἄνθρωπος, ὥστε ἁμαρτεῖν ἐν τούτοις·
\vs{23}Καὶ ἔσται ἡνίκα ἐὰν ἁμάρτῃ, καὶ πλημμελήσῃ, καὶ ἀποδῷ τὸ ἅρπαγμα, ὃ ἥρπασεν, ἢ τὸ ἀδίκημα, ὃ ἠδίκησεν, ἢ τὴν παραθήκην, ἥτις παρετέθη αὐτῷ, ἢ τὴν ἀπώλειαν, ἣν εὗρεν
\vs{24}ἀπὸ παντὸς πράγματος, οὗ ὤμοσε περὶ αὐτοῦ ἀδίκως, καὶ ἀποτίσει αὐτὸ τὸ κεφάλαιον, καὶ τὸ ἐπίπεμπτον προσθήσει ἐπʼ αὐτὸ, τίνος ἐστίν, αὐτῷ ἀποδώσει ᾗ ἡμέρᾳ ἐλεγχθῇ.
\vs{25}Καὶ τῆς πλημμελείας αὐτου οἴσει τῷ Κυρίῳ κριὸν ἀπὸ τῶν προβάτων ἄμωμον, τιμῆς, εἰς ὃ ἐπλημμέλησε.
\vs{26}Καὶ ἐξιλάσεται περὶ αὐτοῦ ὁ ἱερεὺς ἔναντι Κυρίου, καὶ ἀφεθήσεται αὐτῷ περὶ ἑνὸς ἀπὸ πάντων ὧν ἐποίησε καὶ ἐπλημμέλησεν αὐτῷ.

\ch{6}
Καὶ ἐλάλησε Κύριος πρὸς Μωυσῆν, λέγων,
\vs{2}ἔντειλαι τῷ Ἀαρὼν καὶ τοῖς υἱοῖς αὐτοῦ, λέγων, οὗτος ὁ νόμος τῆς ὁλοκαυτώσεως· αὕτη ἡ ὁλοκαύτωσις ἐπὶ τῆς καύσεως αὐτῆς ἐπὶ τοῦ θυσιαστηρίου ὅλην τὴν νύκτα ἕως τοπρωῒ, καὶ τὸ πῦρ τοῦ θυσιαστηρίου καυθήσεται ἐπʼ αὐτοῦ, οὐ σβεσθήσεται.
\vs{3}Καὶ ἐνδύσεται ὁ ἱερεὺς χιτῶνα λινοῦν, καὶ περισκελὲς λινοῦν ἐνδύσεται περὶ τὸ σῶμα αὐτοῦ, καὶ ἀφελεῖ τὴν κατακάρπωσιν, ἣν ἂν καταναλώσῃ τὸ πῦρ, τὴν ὁλοκαύτωσιν ἀπὸ τοῦ θυσιαστηρίου· καὶ παραθήσει αὐτὸ ἐχόμενον τοῦ θυσιαστηρίου.
\vs{4}Καὶ ἐκδύσεται τὴν στολὴν αὐτοῦ, καὶ ἐνδύσεται στολὴν ἄλλην· καὶ ἐξοίσει τὴν κατακάρπωσιν ἔξω τῆς παρεμβολῆς εἰς τόπον καθαρόν.
\vs{5}Καὶ πῦρ ἐπὶ τὸ θυσιαστήριον καυθήσεται ἀπʼ αὐτοῦ, καὶ οὐ σβεσθήσεται· καὶ καύσει ἐπʼ αὐτοῦ ὁ ἱερεὺς ξύλα τοπρωῒ πρωῒ, καὶ στοιβάσει ἐπʼ αὐτοῦ τὴν ὁλοκαύτωσιν, καὶ ἐπιθήσει ἐπʼ αὐτὸ τὸ στέαρ τοῦ σωτηρίου.
\vs{6}Καὶ πῦρ διαπαντὸς καυθήσεται ἐπὶ τὸ θυσιαστήριον, οὐ σβεσθήσεται.
\vs{7}Οὗτος ὁ νόμος τῆς θυσίας, ἣν προσάξουσιν αὐτὴν οἱ υἱοὶ Ἀαρὼν ἔναντι Κυρίου, ἀπέναντι τοῦ θυσιαστηρίου.
\vs{8}Καὶ ἀφελεῖ ἀπʼ αὐτοῦ τῇ δρακὶ ἀπὸ τῆς σεμιδάλεως τῆς θυσίας σὺν τῷ ἐλαίῳ αὐτῆς, καὶ σὺν παντὶ τῷ λιβάνῳ αὐτῆς, τὰ ὄντα ἐπὶ τῆς θυσίας· καὶ ἀνοίσει ἐπὶ τὸ θυσιαστήριον κάρπωμα ὀσμὴν εὐωδίας, τὸ μνημόσυνον αὐτῆς τῷ Κυρίῳ.
\vs{9}Τὸ δὲ καταλειφθὲν ἀπʼ αὐτῆς ἔδεται Ἀαρὼν καὶ οἱ υἱοὶ αὐτοῦ· ἄζυμα βρωθήσεται ἐν τόπῳ ἁγίῳ· ἐν αὐλῇ τῆς σκηνῆς τοῦ μαρτυρίου ἔδονται αὐτήν.
\vs{10}Οὐ πεφθήσεται ἐζυμωμένη· μερίδα αὐτὴν ἔδωκα αὐτοῖς ἀπὸ τῶν καρπωμάτων Κυρίου· ἅγια ἁγίων ἐστὶν, ὥσπερ τὸ τῆς ἁμαρτίας, καὶ ὥσπερ τὸ τῆς πλημμελείας.
\vs{11}Πᾶν ἀρσενικὸν τῶν ἱερέων ἔδονται αὐτήν· νόμιμον αἰώνιον εἰς τὰς γενεὰς ὑμῶν ἀπὸ τῶν καρπωμάτων Κυρίου· πᾶς ὃς ἐὰν ἅψηται αὐτῶν, ἁγιασθήσεται.

\vs{12}Καὶ ἐλάλησε Κύριος πρὸς Μωυσῆν, λέγων,
\vs{13}τοῦτο τὸ δῶρον Ἀαρὼν καὶ τῶν υἱῶν αὐτοῦ, ὃ προσοίσουσι Κυρίῳ ἐν τῇ ἡμέρᾳ, ᾗ ἂν χρίσῃς αὐτόν· τὸ δέκατον τοῦ οἰφὶ σεμιδάλεως εἰς θυσίαν διαπαντὸς, τὸ ἥμισυ αὐτῆς τὸπρωῒ, καὶ τὸ ἥμισυ αὐτῆς τοδειλινόν.
\vs{14}Ἐπὶ τηγάνου ἐν ἐλαίῳ ποιηθήσεται, πεφυραμένην οἴσει αὐτήν ἑλικτά, θυσίαν ἐκ κλασμάτων, θυσίαν εἰς ὀσμὴν εὐωδίας Κυρίῳ.
\vs{15}Ὁ ἱερεὺς ὁ χριστὸς ὁ ἀντʼ αὐτοῦ ἐκ τῶν υἱῶν αὐτοῦ ποιήσει αὐτήν· νόμος αἰώνιος· ἅπαν ἐπιτελεσθήσεται.
\vs{16}Καὶ πᾶσα θυσία ἱερέως ὁλόκαυτος ἔσται, καὶ οὐ βρωθήσεται.
\vs{17}Καὶ ἐλάλησε Κύριος πρὸς Μωυσῆν, λέγων,
\vs{18}λάλησον τῷ Ἀαρὼν καὶ τοῖς υἱοῖς αὐτοῦ, λέγων, οὗτος ὁ νόμος τῆς ἁμαρτίας· ἐν τόπῳ οὗ σφάζουσι τὸ ὁλοκαύτωμα, σφάξουσι τὰ περὶ τῆς ἁμαρτίας ἔναντι Κυρίου· ἅγια ἁγίων ἐστίν.
\vs{19}Ὁ ἱερεὺς ὁ ἀναφέρων αὐτὴν, ἔδεται αὐτήν· ἐν τόπῳ ἁγίῳ βρωθήσεται, ἐν αὐλῇ τῆς σκηνῆς τοῦ μαρτυρίου.
\vs{20}Πᾶς ὁ ἁπτόμενος τῶν κρεῶν αὐτῆς, ἁγιασθήσεται· καὶ ᾧ ἐὰν ἐπιῤῥαντισθῇ ἀπὸ τοῦ αἵματος αὐτῆς ἐπὶ τὸ ἱμάτιον, ὃς ἐὰν ῥαντισθῇ ἐπʼ αὐτὸ, πλυθήσεται ἐν τόπῳ ἁγίῳ.
\vs{21}Καὶ σκεῦος ὀστράκινον, οὗ ἐὰν ἑψεθῇ ἐν αὐτῷ, συντριβήσεται· ἐὰν δὲ ἐν σκεύει χαλκῷ ἑψηθῇ, ἐκτρίψει αὐτὸ, καὶ ἐκκλύσει ὕδατι.
\vs{22}Πᾶς ἄρσην ἐν τοῖς ἱερεῦσιν φάγεται αὐτά ἅγια ἁγίων ἐστὶ Κυρίῳ.
\vs{23}Καὶ πάντα τὰ περὶ τῆς ἁμαρτίας, ὧν ἐὰν εἰσενεχθῇ ἀπὸ τοῦ αἵματος αὐτῶν εἰς τὴν σκηνὴν τοῦ μαρτυρίου ἐξιλάσασθαι ἐν τῷ ἁγίῳ, οὐ βρωθήσεται· ἐν πυρὶ κατακαυθήσεται.

\ch{7}
Καὶ οὗτος ὁ νόμος τοῦ κριοῦ τοῦ περὶ τῆς πλημμελείας· ἅγια ἁγίων ἐστίν.
\vs{2}Ἐν τόπῳ οὗ σφάζουσι τὸ ὁλοκαύτωμα, σφάξουσι τὸν κριὸν τῆς πλημμελείας ἔναντι Κυρίου· καὶ τὸ αἷμα προσχεεῖ ἐπὶ τὴν βάσιν τοῦ θυσιαστηρίου κύκλῳ·
\vs{3}Καὶ πᾶν τὸ στέαρ αὐτοῦ προσοίσει ἀπʼ αὐτοῦ, καὶ τὴν ὀσφὺν, καὶ πᾶν τὸ στέαρ τὸ κατακαλύπτον τὰ ἐνδόσθια, καὶ πᾶν τὸ στέαρ τὸ ἐπὶ τῶν ἐνδοσθίων,
\vs{4}καὶ τοὺς δύο νεφροὺς, καὶ τὸ στέαρ τὸ ἐπʼ αὐτῶν, τὸ ἐπὶ τῶν μηρίων, καὶ τὸν λοβὸν τὸν ἐπὶ τοῦ ἥπατος σὺν τοῖς νεφροῖς, περιελεῖ αὐτά.
\vs{5}Καὶ ἀνοίσει αὐτὰ ὁ ἱερεὺς ἐπὶ τὸ θυσιαστήριον κάρπωμα τῷ Κυρίῳ· περὶ πλημμελείας ἐστί.
\vs{6}Πᾶς ἄρσην ἐκ τῶν ἱερέων ἔδεται αὐτά· ἐν τόπῳ ἁγίῳ ἔδονται αὐτά ἅγια ἁγίων ἐστίν.
\vs{7}Ὥσπερ τὸ περὶ τῆς ἁμαρτίας, οὕτω καὶ τὸ τῆς πλημμελείας· νόμος εἷς αὐτῶν· ὁ ἱερεὺς ὅστις ἐξιλάσεται ἐν αὐτῷ, αὐτῷ ἔσται.
\vs{8}Καὶ ὁ ἱερεὺς ὁ προσάγων ὁλοκαύτωμα ἀνθρώπου, τὸ δέρμα τῆς ὁλοκαυτώσεως, ἧς προσφέρει αὐτὸς, αὐτῷ ἔσται.
\vs{9}Καὶ πᾶσα θυσία ἥτις ποιηθήσεται ἐν τῷ κλιβάνῳ, καὶ πᾶσα ἥτις ποιηθήσεται ἐπʼ ἐσχάρας, ἢ ἐπὶ τηγάνου, τοῦ ἱερέως τοῦ προσφέροντος αὐτὴν, αὐτῷ ἔσται.
\vs{10}Καὶ πᾶσα θυσία ἀναπεποιημένη ἐν ἐλαίῳ, καὶ μὴ ἀναπεποιημένη, πᾶσι τοῖς υἱοῖς Ἀαρὼν ἔσται, ἑκάστῳ τὸ ἶσον.

\vs{11}Οὗτος ὁ νόμος θυσίας σωτηρίου, ἣν προσοίσουσι Κυρίῳ.
\vs{12}Ἐὰν μὲν περὶ αἰνέσεως προσφέρῃ αὐτήν, καὶ προσοίσει ἐπὶ τῆς θυσίας τῆς αἰνέσεως ἄρτους ἐκ σεμιδάλεως ἀναπεποιημένους ἐν ἐλαίῳ, καὶ λάγανα ἄζυμα διακεχρισμένα ἐν ἐλαίῳ, καὶ σεμίδαλιν πεφυραμένην ἐν ἐλαίῳ.
\vs{13}Ἐπʼ ἄρτοις ζυμίταις προσοίσει τὰ δῶρα αὐτοῦ ἐπὶ θυσίᾳ αἰνέσεως σωτηρίου.
\vs{14}Καὶ προσάξει ἓν ἀπὸ πάντων τῶν δώρων αὐτοῦ, ἀφαίρεμα Κυρίῳ· τῷ ἱερεῖ τῷ προσχέοντι τὸ αἷμα τοῦ σωτηρίου, αὐτῷ ἔσται.
\vs{15}Καὶ τὰ κρέα θυαίας αἰνέσεως σωτηρίου αὐτῷ ἔσται· καὶ ἐν ᾗ ἡμέρᾳ δωρεῖται, βρωθήσεται· οὐ καταλείψουσιν ἀπʼ αὐτοῦ εἰς τὸ πρωΐ.
\vs{16}Καὶ ἐὰν εὐχὴ ᾖ, ἢ ἑκούσιον θυσιάζῃ τὸ δῶρον αὐτοῦ, ᾗ ἂν ἡμέρᾳ προσαγάγῃ τὴν θυσίαν αὐτοῦ, βρωθήσεται, καὶ τῇ αὔριον.
\vs{17}Καὶ τὸ καταλειφθὲν ἀπὸ τῶν κρεῶν τῆς θυσίας ἕως ἡμέρας τρίτης, ἐν πυρὶ κατακαυθήσεται.
\vs{18}Ἐὰν δὲ φαγὼν φάγῃ ἀπὸ τῶν κρεῶν τῇ ἡμέρᾳ τῇ τρίτῃ, οὐ δεχθήσεται αὐτῷ τῷ προσφέροντι αὐτό· οὐ λογισθήσεται αὐτῷ, μίασμά ἐστιν· ἡ δὲ ψυχὴ ἥτις ἐὰν φάγῃ ἀπʼ αὐτοῦ, τὴν ἁμαρτίαν λήψεται.
\vs{19}Καὶ κρέα ὅσα ἐὰν ἅψηται παντὸς ἀκαθάρτου, οὐ βρωθήσεται, ἐν πυρὶ κατακαυθήσεται· πᾶς καθαρὸς φάγεται κρέα.
\vs{20}Ἡ δὲ ψυχὴ ἥτις ἐὰν φάγῃ ἀπὸ τῶν κρεῶν τῆς θυσίας τοῦ σωτηρίου, ὅ ἐστι Κυρίου, καὶ ἡ ἀκαθαρσία αὐτοῦ ἐπʼ αὐτῷ, ἀπολεῖται ἡ ψυχὴ ἐκείνη ἐκ τοῦ λαοῦ αὐτῆς.
\vs{21}Καὶ ἡ ψυχὴ ἣ ἂν ἅψηται παντὸς πράγματος ἀκαθάρτου, ἢ ἀπὸ ἀκαθαρσίας ἀνθρώπου, ἢ τῶν τετραπόδων τῶν ἀκαθάρτω ἢ παντὸς βδελύγματος ἀκαθάρτου, καὶ φάγῃ ἀπὸ τῶν κρεῶν τῆς θυσίας τοῦ σωτηρίου, ὅ ἐστι Κυρίου, ἀπολεῖται ἡ ψυχὴ ἐκείνη ἐκ τοῦ λαοῦ αὐτῆς.

\vs{22}Καὶ ἐλάλησε Κύριος πρὸς Μωυσῆν, λέγων,
\vs{23}λάλησον τοῖς υἱοῖς Ἰσραὴλ, λέγων, πᾶν στέαρ βοῶν, καὶ προβάτων, καὶ αἰγῶν οὐκ ἔδεσθε.
\vs{24}Καὶ στέαρ θνησιμαίων καὶ θηριαλώτων ποιηθήσεται εἰς πᾶν ἔργον, καὶ εἰς βρῶσιν οὐ βρωθήσεται.
\vs{25}Πᾶς ὁ ἔσθων στέαρ ἀπὸ τῶν κτηνῶν, ὧν προσάξει ἀπʼ αὐτῶν κάρπωμα Κυρίῳ, ἀπολεῖται ἡ ψυχὴ ἐκείνη ἀπὸ τοῦ λαοῦ αὐτῆς.
\vs{26}Πᾶν αἷμα οὐκ ἔδεσθε ἐν πάσῃ τῇ κατοικίᾳ ὑμῶν, ἀπό τε τῶν κτηνῶν καὶ ἀπὸ τῶν πετεινῶν.
\vs{27}Πᾶσα ψυχὴ ἣ ἂν φάγῃ αἷμα, ἀπολεῖται ἡ ψυχὴ ἐκείνη ἀπὸ τοῦ λαοῦ αὐτῆς.

\vs{28}Καὶ ἐλάλησε Κύριος πρὸς Μωυσῆν, λέγων,
\vs{29}καὶ τοῖς υἱοῖς Ἰσραὴλ λαλήσεις, λέγων, ὁ προσφέρων θυσίαν σωτηρίου, οἴσει τὸ δῶρον αὐτοῦ Κυρίῳ καὶ ἀπὸ τῆς θυσίας τοῦ σωτηρίου.
\vs{30}Αἱ χεῖρες αὐτοῦ προσοίσουσι τὰ καρπώματα Κυρίῳ· τὸ στέαρ τὸ ἐπὶ τοῦ στηθυνίου, καὶ τὸν λοβὸν τοῦ ἥπατος προσοίσει αὐτὰ, ὥστε ἐπιτιθέναι δόμα ἔναντι Κυρίου.
\vs{31}Καὶ ἀνοίσει ὁ ἱερεὺς τὸ στέαρ ἐπὶ τοῦ θυσιαστηρίου· καὶ ἔσται τὸ στηθύνιον Ἀαρὼν καὶ τοῖς υἱοῖς αὐτοῦ.
\vs{32}Καὶ τὸν βραχίονα τὸν δεξιὸν δώσετε ἀφαίρεμα τῷ ἱερεῖ ἀπὸ τῶν θυσιῶν τοῦ σωτηρίου ὑμῶν.
\vs{33}Ὁ προσφέρων τὸ αἷμα τοῦ σωτηρίου, καὶ τὸ στέαρ τὸ ἀπὸ τῶν υἱῶν Ἀαρὼν, αὐτῷ ἔσται ὁ βραχίων ὁ δεξιὸς ἐν μερίδι.
\vs{34}Τὸ γὰρ στηθύνιον τοῦ ἐπιθέματος καὶ τὸν βραχίονα τοῦ ἀφαιρέματος εἴληφα παρὰ τῶν υἱῶν Ἰσραὴλ ἀπὸ τῶν θυσιῶν τοῦ σωτηρίου ὑμῶν, καὶ ἔδωκα αὐτὰ Ἀαρὼν τῷ ἱερεῖ καὶ τοῖς υἱοῖς αὐτοῦ, νόμιμον αἰώνιον παρὰ τῶν υἱῶν Ἰσραήλ.
\vs{35}Αὕτη ἡ χρίσις Ἀαρὼν, καὶ ἡ χρίσις τῶν υἱῶν αὐτοῦ ἀπὸ τῶν καρπωμάτων Κυρίου, ἐν ᾗ ἡμέρᾳ προσηγάγετο αὐτοὺς τοῦ ἱερατεύειν τῷ Κυρίῳ,
\vs{36}καθὰ ἐνετείλατο Κύριος δοῦναι αὐτοῖς ᾗ ἡμέρᾳ ἔχρισεν αὐτοὺς παρὰ τῶν υἱῶν Ἰσραὴλ, νόμιμον αἰώνιον εἰς τὰς γενεὰς αὐτῶν.
\vs{37}Οὗτος ὁ νόμος τῶν ὁλοκαυτωμάτων, καὶ θυσίας, καὶ περὶ ἁμαρτίας, καὶ τῆς πλημμελείας καὶ τῆς τελειώσεως, καὶ τῆς θυσίας τοῦ σωτηρίου,
\vs{38}ὃν τρόπον ἐνετείλατο Κύριος τῷ Μωυσῇ ἐν τῷ ὄρει Σινᾷ, ᾗ ἡμέρᾳ ἐνετείλατο τοῖς υἱοῖς Ἰσραὴλ προσφέρειν τὰ δῶρα αὐτῶν ἔναντι Κυρίου ἐν τῇ ἐρήμῳ Σινᾷ.

\ch{8}
Καὶ ἐλάλησε Κύριος πρὸς Μωυσῆν, λέγων,
\vs{2}λάβε Ἀαρὼν καὶ τοὺς υἱοὺς αὐτοῦ, καὶ τὰς στολὰς αὐτοῦ, καὶ τὸ ἔλαιον τῆς χρίσεως, καὶ τὸν μόσχον τὸν περὶ τῆς ἁμαρτίας, καὶ τοὺς δύο κριοὺς, καὶ τὸ κανοῦν τῶν ἀζύμων,
\vs{3}καὶ πᾶσαν τὴν συναγωγὴν ἐκκλησίασον ἐπὶ τὴν θύραν τῆς σκηνῆς τοῦ μαρτυρίου.
\vs{4}Καὶ ἐποίησε Μωυσῆς ὃν τρόπον συνέταξεν αὐτῷ Κύριος· καὶ ἐξεκκλησίασε τὴν συναγωγὴν ἐπὶ τὴν θύραν τῆς σκηνῆς τοῦ μαρτυρίου.
\vs{5}Καὶ εἶπε Μωυσῆς τῇ συναγωγῇ, τοῦτό ἐστι τὸ ῥῆμα, ὃ ἐνετείλατο Κύριος ποιῆσαι.
\vs{6}Καὶ προσήνεγκε Μωυσῆς τὸν Ἀαρὼν, καὶ τοὺς υἱοὺς αὐτοῦ, καὶ ἔλουσεν αὐτοὺς ὕδατι.
\vs{7}Καὶ ἐνέδυσεν αὐτὸν τὸν χιτῶνα, καὶ ἔζωσεν αὐτὸν τὴν ζώνην, καὶ ἐνέδυσεν αὐτὸν τὸν ὑποδύτην, καὶ ἐπέθηκεν ἐπʼ αὐτὸν τὴν ἐπωμίδα.
\vs{8}Καὶ συνέζωσεν αὐτὸν κατὰ τὴν ποίησιν τῆς ἐπωμίδος, καὶ συνέσφιγξεν αὐτὸν ἐν αὐτῇ· καὶ ἐπέθηκεν ἐπʼ αὐτὴν τὸ λογεῖον, καὶ ἐπέθηκεν ἐπὶ τὸ λογεῖον τὴν δήλωσιν καὶ τὴν ἀλήθειαν.
\vs{9}Καὶ ἐπέθηκε τὴν μίτραν ἐπὶ τὴν κεφαλὴν αὐτοῦ, καὶ ἐπέθηκεν ἐπὶ τὴν μίτραν κατὰ πρόσωπον αὐτοῦ τὸ πέταλον τὸ χρυσοῦν τὸ καθηγιασμένον ἅγιον, ὃν τρόπον συνέταξε Κύριος τῷ Μωυσῇ.

\vs{10}Καὶ ἔλαβε Μωυσῆς ἀπὸ τοῦ ἐλαίου τῆς χρίσεως,
\vs{11}καὶ ἔῤῥανεν ἀπʼ αὐτοῦ ἐπὶ τὸ θυσιαστήριον ἑπτάκις· καὶ ἔχρισε τὸ θυσιαστήριον, καὶ ἡγίασεν αὐτὸ, καὶ πάντα τὰ ἐν αὐτῷ, καὶ τὸν λουτῆρα, καὶ τὴν βάσιν αὐτοῦ, καὶ ἡγίασεν αὐτά· καὶ ἔχρισε τὴν σκηνὴν, καὶ πάντα τὰ σκεύη αὐτῆς, καὶ ἡγίασεν αὐτήν.
\vs{12}Καὶ ἐπέχεε Μωυσῆς ἀπὸ τοῦ ἐλαίου τῆς χρίσεως ἐπὶ τὴν κεφαλὴν Ἀαρών· καὶ ἔχρισεν αὐτὸν, καὶ ἡγίασεν αὐτόν.
\vs{13}Καὶ προσήγαγε Μωυσῆς τοὺς υἱοὺς Ἀαρὼν, καὶ ἐνέδυσεν αὐτοὺς χιτῶνας, καὶ ἔζωσεν αὐτοὺς ζώνας, καὶ περιέθηκεν αὐτοῖς κιδάρεις, καθάπερ συνέταξε Κύριος τῷ Μωυσῇ.

\vs{14}Καὶ προσήγαγε Μωυσῆς τὸν μόσχον τὸν περὶ τῆς ἁμαρτίας· καὶ ἐπέθηκεν Ἀαρὼν καὶ οἱ υἱοὶ αὐτοῦ τὰς χεῖρας ἐπὶ τὴν κεφαλὴν τοῦ μόσχου τοῦ τῆς ἁμαρτίας.
\vs{15}Καὶ ἔσφαξεν αὐτόν· καὶ ἔλαβε Μωυσῆς ἀπὸ τοῦ αἵματος, καὶ ἐπέθηκεν ἐπὶ τὰ κέρατα τοῦ θυσιαστηρίου κύκλῳ τῷ δακτύλῳ, καὶ ἐκαθάρισε τὸ θυσιαστήριον· καὶ τὸ αἷμα ἐξέχεεν ἐπὶ τὴν βάσιν τοῦ θυσιαστηρίου, καὶ ἡγίασεν αὐτὸ, τοῦ ἐξιλάσασθαι ἐπʼ αὐτοῦ.
\vs{16}Καὶ ἔλαβε Μωυσῆς πᾶν τὸ στέαρ τὸ ἐπὶ τῶν ἐνδοσθίων, καὶ τὸν λοβὸν τὸν ἐπὶ τοῦ ἥπατος, καὶ ἀμφοτέρους τοὺς νεφροὺς, καὶ τὸ στέαρ τὸ ἐπʼ αὐτῶν, καὶ ἀνήνεγκε Μωυσῆς ἐπὶ τὸ θυσιαστήριον.
\vs{17}Καὶ τὸν μόσχον, καὶ τὴν βύρσαν αὐτοῦ, καὶ τὰ κρέα αὐτοῦ, καὶ τὴν κόπρον αὐτοῦ, κατέκαυσεν αὐτὰ πυρὶ ἔξω τῆς παρεμβολῆς, ὃν τρόπον συνέταξε Κύριος τῷ Μωυσῇ.
\vs{18}Καὶ προσήγαγε Μωυσῆς τὸν κριὸν τὸν εἰς ὁλοκαύτωμα· καὶ ἐπέθηκεν Ἀαρὼν καὶ υἱοὶ αὐτοῦ τὰς χεῖρας αὐτῶν ἐπὶ τὴν κεφαλὴν τοῦ κριοῦ. Καὶ ἔσφαξε Μωυσῆς τὸν κριόν· καὶ προσέχεε Μωυσῆς τὸ αἷμα ἐπὶ τὸ θυσιαστήριον κύκλῳ.
\vs{19}Καὶ τὸν κριὸν ἐκρεανόμησε κατὰ μέλη· καὶ ἀνήνεγκε Μωυσῆς τὴν κεφαλὴν, καὶ τὰ μέλη, καὶ τὸ στέαρ· καὶ τὴν κοιλίαν, καὶ τοὺς πόδας ἔπλυνεν ὕδατι.
\vs{20}Καὶ ἀνήνεγκε Μωυσῆς ὅλον τὸν κριὸν ἐπὶ τὸ θυσιαστήριον· ὁλοκαύτωμά ἐστιν εἰς ὀσμὴν εὐωδίας· κάρπωμά ἐστι τῷ Κυρίῳ, καθάπερ ἐνετείλατο Κύριος τῷ Μωυσῇ.

\vs{21}Καὶ προσήγαγε Μωυσῆς τὸν κριὸν τὸν δεύτερον, κριὸν τελειώσεως· καὶ ἐπέθηκεν Ἀαρὼν καὶ οἱ υἱοὶ αὐτοῦ τὰς χεῖρας αὐτῶν ἐπὶ τὴν κεφαλὴν τοῦ κριοῦ.
\vs{22}Καὶ ἔσφαξεν αὐτόν· καὶ ἔλαβε Μωυσῆς ἀπὸ τοῦ αἵματος αὐτοῦ, καὶ ἐπέθηκεν ἐπὶ τὸν λοβὸν τοῦ ὠτὸς Ἀαρὼν τοῦ δεξιοῦ, καὶ ἐπὶ τὸ ἄκρον τῆς χειρὸς τῆς δεξιᾶς, καὶ ἐπὶ τὸ ἄκρον τοῦ ποδὸς τοῦ δεξιοῦ.
\vs{23}Καὶ προσήγαγε Μωυσῆς τοὺς υἱοὺς Ἀαρών· καὶ ἐπέθηκε Μωυσῆς ἀπὸ τοῦ αἵματος ἐπὶ τοὺς λοβοὺς τῶν ὤτων τῶν δεξιῶν, καὶ ἐπὶ τὰ ἄκρα τῶν χειρῶν αὐτῶν τῶν δεξιῶν· καὶ ἐπὶ τὰ ἄκρα τῶν ποδῶν αὐτῶν τῶν δεξιῶν· καὶ προσέχεε Μωυσῆς τὸ αἷμα ἐπὶ τὸ θυσιαστήριον κύκλῳ.
\vs{24}Καὶ ἔλαβε τὸ στέαρ, καὶ τὴν ὀσφὺν, καὶ τὸ στέαρ τὸ ἐπὶ τῆς κοιλίας, καὶ τὸν λοβὸν τοῦ ἥπατος, καὶ τοὺς δύο νεφροὺς, καὶ τὸ στέαρ τὸ ἐπʼ αὐτῶν, καὶ τὸν βραχίονα τὸν δεξιόν.
\vs{25}Καὶ ἀπὸ τοῦ κανοῦ τῆς τελειώσεως, τοῦ ὄντος ἔναντι Κυρίου, καὶ ἔλαβεν ἄρτον ἕνα ἄζυμον, καὶ ἄρτον ἐξ ἐλαίου ἕνα, καὶ λάγανον ἓν, καὶ ἐπέθηκεν ἐπὶ τὸ στέαρ, καὶ τὸν βραχίονα τὸν δεξιόν.
\vs{26}Καὶ ἐπέθηκεν ἅπαντα ἐπὶ τὰς χεῖρας Ἀαρὼν, καὶ ἐπὶ τὰς χεῖρας τῶν υἱῶν αὐτοῦ, καὶ ἀνήνεγκεν αὐτὰ ἀφαίρεμα ἔναντι Κυρίου.
\vs{27}Καὶ ἔλαβε Μωυσῆς ἀπὸ τῶν χειρῶν αὐτῶν, καὶ ἀνήνεγκεν αὐτὰ Μωυσῆς ἐπὶ τὸ θυσιαστήριον, ἐπὶ τὸ ὁλοκαύτωμα τῆς τελειώσεως, ὅ ἐστι ὀσμὴ εὐωδίας· κάρπωμά ἐστιν τῷ Κυρίῳ.
\vs{28}Καὶ λαβὼν Μωυσῆς τὸ στηθύνιον, ἀφεῖλεν αὐτὸ ἐπίθεμα ἔναντι Κυρίου, ἀπὸ τοῦ κριοῦ τῆς τελειώσεως· καὶ ἐγένετο Μωυσῇ ἐν μεριδι, καθὰ ἐνετείλατο Κύριος τῷ Μωυσῇ.

\vs{29}Καὶ ἔλαβε Μωυσῆς ἀπὸ τοῦ ἐλαίου τῆς χρίσεως, καὶ ἀπὸ τοῦ αἵματος τοῦ ἐπὶ τοῦ θυσιαστηρίου, καὶ προσέῤῥνεν ἐπὶ Ἀαρὼν, καὶ τὰς στολὰς αὐτοῦ, καὶ τοὺς υἱοὺς αὐτοῦ, καὶ τὰς στολὰς τῶν υἱῶν αὐτοῦ μετʼ αὐτοῦ.
\vs{30}Καὶ ἡγίασεν Ἀαρὼν, καὶ τὰς στολάς αὐτοῦ, καὶ τοὺς υἱοὺς αὐτοῦ, καὶ τὰς στολὰς τῶν υἱῶν αὐτοῦ μετʼ αὐτοῦ.
\vs{31}Καὶ εἶπε Μωυσῆς πρὸς Ἀαρὼν, καὶ τοὺς υἱοὺς αὐτοῦ, ἑψήσατε τὰ κρέα ἐν τῇ αὐλῇ τῆς σκηνῆς τοῦ μαρτυρίου ἐν τόπῳ ἁγίῳ· καὶ ἐκεῖ φάγεσθε αὐτὰ, καὶ τοὺς ἄρτους τοὺς ἐν τῷ κανῷ τῆς τελειώσεως, ὃν τρόπον συντέτακταί μοι, λέγων, Ἀαρὼν καὶ οἱ υἱοὶ αὐτοῦ φάγονται αὐτά.
\vs{32}Καὶ τὸ καταλειφθὲν τῶν κρεῶν καὶ τῶν ἄρτων ἐν πυρὶ κατακαύσατε.
\vs{33}Καὶ ἀπὸ τῆς θύρας τῆς σκηνῆς τοῦ μαρτυρίου οὐκ ἐξελεύσεσθε ἑπτὰ ἡμέρας, ἕως ἡμέρα πληρωθῇ, ἡμέρα τελειώσεως ὑμῶν· ἑπτὰ γὰρ ἡμέρας τελειώσει τὰς χεῖρας ὑμῶν.
\vs{34}Καθάπερ ἐποίησεν ἐν τῇ ἡμέρᾳ ταύτῃ, ᾗ ἐνετείλατο Κύριος τοῦ ποιῆσαι, ὥστε ἐξιλάσασθαι περὶ ὑμῶν.
\vs{35}Καὶ ἐπὶ τὴν θύραν τῆς σκηνῆς τοῦ μαρτυρίου καθήσεσθε ἑπτὰ ἡμέρας, ἡμέραν καὶ νύκτα· φυλάξεσθε τὰ φυλάγματα Κυρίου, ἵνα μὴ ἀποθάνητε· οὕτω γὰρ ἐνετείλατό μοι Κύριος ὁ Θεός.
\vs{36}Καὶ ἐποίησεν Ἀαρὼν καὶ οἱ υἱοὶ αὐτοῦ πάντας τοὺς λόγους, οὓς συνέταξε Κύριος τῷ Μωυσῇ.

\ch{9}
Καὶ ἐγενήθη τῇ ἡμέρᾳ τῇ ὀγδόῃ, ἐκάλεσε Μωυσῆς Ἀαρὼν, καὶ τοὺς υἱοὺς αὐτοῦ, καὶ τὴν γερουσίαν Ἰσραὴλ,
\vs{2}καὶ εἶπε Μωυσῆς πρὸς Ἀαρών, λάβε σεαυτῷ μοσχάριον ἐκ βοῶν περὶ ἁμαρτίας, καὶ κριὸν εἰς ὁλοκαύτωμα, ἄμωμα, καὶ προσένεγκε αὐτὰ ἔναντι Κυρίου.
\vs{3}Καὶ τῇ γερουσίᾳ Ἰσραὴλ λάλησον, λέγων, λάβετε χίμαρον ἐξ αἰγῶν ἕνα περὶ ἁμαρτίας, καὶ μοσχάριον, καὶ ἀμνὸν ἐνιαύσιον εἰς ὁλοκάρπωσιν, ἄμωμα,
\vs{4}καὶ μόσχον, καὶ κριὸν εἰς θυσίαν σωτηρίου ἔναντι Κυρίου, καὶ σεμίδαλιν πεφυραμένην ἐν ἐλαίῳ· ὅτι σήμερον Κύριος ὀφθήσεται ἐν ὑμῖν.
\vs{5}Καὶ ἔλαβον καθὸ ἐνετείλατο Μωυσῆς ἀπέναντι τῆς σκηνῆς τοῦ μαρτυρίου· καὶ προσῆλθε πᾶσα συναγωγὴ, καὶ ἔστησαν ἔναντι Κυρίου.
\vs{6}Καὶ εἶπε Μωυσῆς, τοῦτο τὸ ῥῆμα, ὃ εἶπε Κύριος, ποιήσατε, καὶ ὀφθήσεται ἐν ὑμῖν ἡ δόξα Κυρίου.
\vs{7}Καὶ εἶπε Μωυσῆς τῷ Ἀαρὼν, πρόσελθε πρὸς τὸ θυσιαστήριον, καὶ ποίησον τὸ περὶ τῆς ἁμαρτίας σου, καὶ τὸ ὁλοκαύτωμά σου, καὶ ἐξίλασαι περὶ σεαυτοῦ, καὶ τοῦ οἴκου σου· καὶ ποίησον τὰ δῶρα τοῦ λαοῦ, καὶ ἐξίλασαι περὶ αὐτῶν, καθάπερ ἐνετείλατο Κύριος τῷ Μωυσῇ.
\vs{8}Καὶ προσῆλθεν Ἀαρὼν πρὸς τὸ θυσιαστήριον, καὶ ἔσφαξε τὸ μοσχάριον τὸ περὶ τῆς ἁμαρτίας αὐτοῦ.
\vs{9}Καὶ προσήνεγκαν οἱ υἱοὶ Ἀαρὼν τὸ αἷμα πρὸς αὐτόν· καὶ ἔβαψε τὸν δάκτυλον εἰς τὸ αἷμα, καὶ ἐπέθηκεν ἐπὶ τὰ κέρατα τοῦ θυσιαστηρίου· καὶ τὸ αἷμα ἐξέχεεν ἐπὶ τὴν βάσιν τοῦ θυσιαστηρίου.
\vs{10}Καὶ τὸ στέαρ καὶ τοὺς νεφροὺς καὶ τὸν λοβὸν τοῦ ἥπατος τοῦ περὶ τῆς ἁμαρτίας ἀνήνεγκεν ἐπὶ τὸ θυσιαστήριον, ὃν τρόπον ἐνετείλατο Κύριος τῷ Μωυσῇ.
\vs{11}Καὶ τὰ κρέα καὶ τὴν βύρσαν κατέκαυσεν αὐτὰ πυρὶ, ἔξω τῆς παρεμβολῆς.
\vs{12}Καὶ ἔσφαξε τὸ ὁλοκαύτωμα· καὶ προσήνεγκαν οἱ υἱοὶ Ἀαρὼν τὸ αἷμα πρὸς αὐτόν· καὶ προσέχεεν ἐπὶ τὸ θυσιαστήριον κύκλῳ.
\vs{13}Καὶ τὸ ὁλοκαύτωμα προσήνεγκαν αὐτὸ κατὰ μέλη· αὐτὰ καὶ τὴν κεφαλὴν ἐπέθηκεν ἐπὶ τὸ θυσιαστήριον.
\vs{14}Καὶ ἔπλυνε τὴν κοιλίαν καὶ τοὺς πόδας ὕδατι· καὶ ἐπέθηκεν ἐπὶ τὸ ὁλοκαύτωμα ἐπὶ τὸ θυσιαστήριον.

\vs{15}Καὶ προσήνεγκε τὸ δῶρον τοῦ λαοῦ, καὶ ἔλαβε τὸν χίμαρον τὸν περὶ τῆς ἁμαρτίας τοῦ λαοῦ, καὶ ἔσφαξεν αὐτὸν, καὶ ἐκαθάρισεν αὐτὸν, καθὰ καὶ τὸν πρῶτον.
\vs{16}Καὶ προσήνεγκε τὸ ὁλοκαύτωμα, καὶ ἐποίησεν αὐτὸ ὡς καθήκει.
\vs{17}Καὶ προσήνεγκε τὴν θυσίαν, καὶ ἔπλησε τὰς χεῖρας ἀπʼ αὐτῆς, καὶ ἐπέθηκεν ἐπὶ τὸ θυσιαστήριον χωρὶς τοῦ ὁλοκαυτώματος τοῦ πρωϊνοῦ.
\vs{18}Καὶ ἔσφαξε τὸν μόσχον, καὶ τὸν κριὸν τῆς θυσίας τοῦ σωτηρίου τῆς τοῦ λαοῦ· καὶ προσήνεγκαν οἱ υἱοὶ Ἀαρὼν τὸ αἷμα πρὸς αὐτόν, καὶ προσέχεε πρὸς τὸ θυσιαστήριον κύκλῳ,
\vs{19}καὶ τὸ στέαρ τὸ ἀπὸ τοῦ μόσχου, καὶ τοῦ κριοῦ τὴν ὀσφὺν, καὶ τὸ στέαρ τὸ κατακαλύπτον ἐπὶ τῆς κοιλίας, καὶ τοὺς δύο νεφροὺς, καὶ τὸ στέαρ τὸ ἐπʼ αὐτῶν, καὶ τὸν λοβὸν τὸν ἐπὶ τοῦ ἥπατος.
\vs{20}Καὶ ἐπέθηκε τὰ στέατα ἐπὶ τὰ στηθύνια καὶ ἀνήνεγκε τὰ στέατα ἐπὶ τὸ θυσιαστήριον.
\vs{21}Καὶ τὸ στηθύνιον, καὶ τὸν βραχίονα τὸν δεξιὸν ἀφεῖλεν Ἀαρὼν ἀφαίρεμα ἔναντι Κυρίου, ὃν τρόπον συνέταξε Κύριος τῷ Μωυσῇ.
\vs{22}Καὶ ἐξάρας Ἀαρὼν τὰς χεῖρας ἐπὶ τὸν λαὸν, εὐλόγησεν αὐτούς· καὶ κατέβη ποιήσας τὸ περὶ τῆς ἁμαρτίας, καὶ τὰ ὁλοκαυτώματα, καὶ τὰ τοῦ σωτηρίου.
\vs{23}Καὶ εἰσῆλθε Μωυσῆς καὶ Ἀαρὼν εἰς τὴν σκηνὴν τοῦ μαρτυρίου· καὶ ἐξελθόντες εὐλόγησαν πάντα τὸν λαὸν· καὶ ὤφθη δόξα Κυρίου παντὶ τῷ λαῷ.
\vs{24}Καὶ ἐξῆλθε πῦρ παρὰ Κυρίου, καὶ κατέφαγε τὰ ἐπὶ τοῦ θυσιαστηρίου, τά τε ὁλοκαυτώματα, καὶ τὰ στέατα· καὶ εἶδε πᾶς ὁ λαὸς, καὶ ἐξέστη, καὶ ἔπεσαν ἐπὶ πρόσωπον.

\ch{10}
Καὶ λαβόντες οἱ δύο υἱοὶ Ἀαρὼν Ναδὰβ καὶ Ἀβιοὺδ, ἕκαστος τὸ πυρεῖον αὐτοῦ, ἐπέθηκαν ἐπʼ αὐτὸ πῦρ, καὶ ἐπέβαλον ἐπʼ αὐτὸ θυμίαμα, καὶ προσήνεγκαν ἔναντι Κυρίου πῦρ ἀλλότριον, ὃ οὐ προσέταξε Κύριος αὐτοῖς.
\vs{2}Καὶ ἐξῆλθε πῦρ παρὰ Κυρίου, καὶ κατέφαγεν αὐτοὺς, καὶ ἀπέθανον ἔναντι Κυρίου.
\vs{3}Καὶ εἶπε Μωυσῆς πρὸς Ἀαρὼν, τοῦτό ἐστιν, ὃ εἶπε Κύριος, λέγων, ἐν τοῖς ἐγγίζουσί μοι ἁγιασθήσομαι, καὶ ἐν πάσῃ τῇ συναγωγῇ δοξασθήοσμαι· καὶ κατενύχθη Ἀαρών.
\vs{4}Καὶ ἐκάλεσε Μωυσῆς τὸν Μισαδάη, καὶ τὸν Ἐλισαφὰν, υἱοὺς Ὀζιὴλ, υἱοὺς τοῦ ἀδελφοῦ τοῦ πατρὸς Ἀαρὼν, καὶ εἶπεν αὐτοῖς, προσέλθατε καὶ ἄρατε τοὺς ἀδελφοὺς ὑμῶν ἐκ προσώπου τῶν ἁγίων ἔξω τῆς παρεμβολῆς.
\vs{5}Καὶ προσῆλθον, καὶ ᾖραν αὐτοὺς ἐν τοῖς χιτῶσιν αὐτῶν ἔξω τῆς παρεμβολῆς, ὃν τρόπον εἶπε Μωυσῆς.
\vs{6}Καὶ εἶπε Μωυσῆς πρὸς Ἀαρὼν καὶ Ἐλεάζαρ καὶ Ἰθάμαρ τοὺς υἱοὺς αὐτοῦ τοὺς καταλελειμμένους, τὴν κεφαλὴν ὑμῶν οὐκ ἀποκιδαρώσετε, καὶ τὰ ἱμάτια ὑμῶν οὐ διαῤῥήξετε, ἵνα μὴ ἀποθάνητε, καὶ ἐπὶ πᾶσαν τὴν συναγωγὴν ἔσται θυμός· οἱ δὲ ἀδελφοὶ ὑμῶν, πᾶς ὁ οἶκος Ἰσραὴλ, κλαύσονται τὸν ἐμπυρισμὸν, ὃν ἐνεπυρίσθησαν ὑπὸ Κυρίου.
\vs{7}Καὶ ἀπὸ τὴς θύρας τῆς σκηνῆς τοῦ μαρτυρίου οὐκ ἐξελεύσεσθε, ἵνα μὴ ἀποθάνητε· τὸ ἔλαιον γὰρ τῆς χρίσεως, τὸ παρὰ Κυρίου, ἐφʼ ὑμῖν, καὶ ἐποίησαν κατὰ τὸ ῥῆμα Μωυσῆ.

\vs{8}Καὶ ἐλάλησε Κύριος τῷ Ἀαρὼν, λέγων,
\vs{9}οἶνον καὶ σίκερα οὐ πίεσθε σὺ καὶ οἱ υἱοί σου μετὰ σοῦ, ἡνίκα ἐὰν εἰσπορεύησθε εἰς τὴν σκηνὴν τοῦ μαρτυρίου, ἢ προσπορευομένων ὑμῶν πρὸς τὸ θυσιαστήριον, καὶ οὐ μὴ ἀποθάνητε· νόμιμον αἰώνιον εἰς τὰς γενεὰς ὑμῶν,
\vs{10}διαστεῖλαι ἀναμέσον τῶν ἁγίων καὶ τῶν βεβήλων, καὶ ἀναμέσον τῶν ἀκαθάρτων καὶ τῶν καθαρῶν,
\vs{11}καὶ συμβιβάξειν τοὺς υἱοὺς Ἰσραὴλ ἅπαντα τὰ νόμιμα, ἃ ἐλάλησε Κύριος πρὸς αὐτοὺς διὰ χειρὸς Μωυσῆ.
\vs{12}Καὶ εἶπε Μωυσῆς πρὸς Ἀαρὼν καὶ πρὸς Ἐλεάζαρ καὶ Ἰθάμαρ τοὺς υἱοὺς Ἀαρὼν τοὺς καταλειφθέντας, λάβετε τὴν θυσίαν τὴν καταλειφθεῖσαν ἀπὸ τῶν καρπωμάτων Κυρίου, καὶ φάγεσθε ἄζυμα παρὰ τὸ θυσιαστήριον· ἅγια ἁγίων ἐστί.
\vs{13}Καὶ φάγεσθε αὐτὴν ἐν τόπῳ ἁγίῳ· νόμιμον γάρ σοι ἐστὶ, καὶ νόμιμον τοῖς υἱοῖς σου τοῦτο ἀπὸ τῶν καρπωμάτων Κυρίου· οὕτω γὰρ ἐντέταλταί μοι.
\vs{14}Καὶ τὸ στηθύνιον τοῦ ἀφορίσματος, καὶ τὸν βραχίονα τοῦ ἀφαιρέματος φάγεσθε ἐν τόπῳ ἁγίῳ, σὺ καὶ οἱ υἱοί σου καὶ ὁ οἶκός σου μετὰ σοῦ· νόμιμον γὰρ σοι, καὶ νόμιμον τοῖς υἱοῖς σου ἐδόθη ἀπὸ τῶν θυσιῶν τοῦ σωτηρίου τῶν υἱῶν Ἰσραήλ.
\vs{15}Τὸν βραχίονα τοῦ ἀφαιρέματος, καὶ τὸ στηθύνιον τοῦ ἀφορίσματος ἐπὶ τῶν καρπωμάτων τῶν στεάτωι· προσοίσουσιν ἀφόρισμα ἀφορίσαι ἔναντι Κυρίου· καὶ ἔσται σοι καὶ τοῖς υἱοῖς σου καὶ ταῖς θυγατράσι σου μετὰ σοῦ νόμιμον αἰώνιον, ὃν τρόπον συνέταξε Κύριος τῷ Μωυσῇ.

\vs{16}Καὶ τὸν χίμαρον τὸν περὶ τῆς ἁμαρτίας ζητῶν ἐξεζήτησε Μωυσῆς· καὶ ὁ δὲ ἐνπεπύριστο· καὶ ἐθυμώθη Μωυσῆς ἐπὶ Ἐλεάζαρ καὶ Ἰθάμαρ τοὺς υἱοὺς Ἀαρὼν τοὺς καταλελειμμένους, λέγων,
\vs{17}διατί οὐκ ἐφάγετε τὸ περὶ τῆς ἁμαρτίας ἐν τόπῳ ἁγίῳ; ὅτι γὰρ ἅγια ἁγίων ἐστι, τοῦτο ἔδωκεν ὑμῖν φαγεῖν, ἵνα ἀφέλητε τὴν ἁμαρτίαν τῆς συναγωγῆς, καὶ ἐξιλάσησθε περὶ αὐτῶν ἔναντι Κυρίου.
\vs{18}Οὐ γὰρ εἰσήχθη τοῦ αἵματος αὐτοῦ εἰς τὸ ἅγιον· κατὰ πρόσωπον ἔσω φάγεσθε αὐτὸ ἐν τόπῳ ἁγίῳ, ὃν τρόπον μοι συνέταξε Κύριος.
\vs{19}Καὶ ἐλάλησεν Ἀαρὼν πρὸς Μωυσῆν, λέγων, εἰ σήμερον προσαγιόχασι τὰ περὶ τῆς ἁμαρτίας αὐτῶν, καὶ τὰ ὁλοκαυτώματα αὐτῶν ἔναντι Κυρίου, καὶ συμβέβηκέ μοι τοιαῦτα, καὶ φάγομαι τὰ περὶ τῆς ἁμαρτίας σήμερον, μὴ ἀρεστὸν ἔται Κυρίῳ;
\vs{20}Καὶ ἤκουσε Μωυσῆς, καὶ ἤρεσεν αὐτῷ.

\ch{11}
Καὶ ἐλάλησε Κύριος πρὸς Μωυσῆν καὶ Ἀαρὼν, λέγων,
\vs{2}λαλήσατε τοῖς υἱοῖς Ἰσραὴλ, λέγοντες, ταῦτα τὰ κτήνη, ἃ φάγεσθε ἀπὸ πάντων τῶν κτηνῶν τῶν ἐπὶ τῆς γῆς.
\vs{3}Πᾶν κτῆνος διχηλοῦν ὁπλὴν καὶ ὀνυχιστῆρας ὀνυχίζον δύο χηλῶν, καὶ ἀνάγον μηρυκισμὸν ἐν τοῖς κτήνεσι, ταῦτα φάγεσθε.
\vs{4}Πλὴν ἀπὸ τούτων οὐ φάγεσθε, ἀπὸ τῶν ἀναγόντων μηρυκισμὸν, καὶ ἀπὸ τῶν διχηλούντων τὰς ὁπλὰς, καὶ ὀνυχιζόντων ὀνυχιστῆρας· τὸν κάμηλον, ὅτι ἀνάγει μηρυκισμὸν τοῦτο, ὁπλὴν δὲ οὐ διχηλεῖ, ἀκάθαρτον τοῦτο ὑμῖν.
\vs{5}Καὶ τὸν δασύποδα, ὅτι ἀνάγει μηρυκισμὸν τοῦτο, καὶ ὁπλὴν οὐ διχηλεῖ, ἀκάθαρτον τοῦτο ὑμῖν.
\vs{6}Καὶ τὸν χοιρογρύλλιον, ὅτι οὐκ ἀνάγει μηρυκισμὸν τοῦτο, καὶ ὁπλὴν οὐ διχηλεῖ, ἀκάθαρτον τοῦτο ὑμῖν.
\vs{7}Καὶ τὸν ὗν, ὅτι διχηλεῖ ὁπλὴν τοῦτο, καὶ ὀνυχίζει ὄνυχας ὁπλῆς, καὶ τοῦτο οὐκ ἀνάγει μηρυκισμὸν, ἀκάθαρτον τοῦτο ὑμῖν.
\vs{8}Ἀπὸ τῶν κρεῶν αὐτῶν οὐ φάγεσθε, καὶ τῶν θνησιμαίων αὐτῶν οὐχ ἅψεσθε· ἀκάθαρτα ταῦτα ὑμῖν.

\vs{9}Καὶ ταῦτα, ἃ φάγεσθε ἀπὸ πάντων τῶν ἐν τοῖς ὕδασι· πάντα ὅσα ἐστὶν αὐτοῖς πτερύγια καὶ λεπίδες ἐν τοῖς ὕδασι, καὶ ἐν ταῖς θαλάσσαις, καὶ ἐν τοῖς χειμάῤῥοις, ταῦτα φάγεσθε.
\vs{10}Καὶ πάντα ὅσα οὐκ ἔστιν αὐτοῖς πτερύγια, οὐδὲ λεπίδες ἐν τῷ ὕδατι, ἢ ἐν ταῖς θαλάσσαις, καὶ ἐν τοῖς χειμάῤῥοις, ἀπὸ πάντων ὧν ἐρεύγεται τὰ ὕδατα, καὶ ἀπὸ πάσης ψυχῆς τῆς ζώσης ἐν τῷ ὕδατι, βδέλυγμά ἐστι, καὶ βδελύγματα ἔσονται ὑμῖν.
\vs{11}Ἀπὸ τῶν κρεῶν αὐτῶν οὐκ ἔδεσθε, καὶ τὰ θνησιμαῖα αὐτῶν βδελύξεσθε.
\vs{12}Καὶ πάντα ὅσα οὐκ ἔστιν αὐτοῖς πτερύγια, οὐδὲ λεπίδες τῶν ἐν τοῖς ὕδασι, βδέλυγμα τοῦτό ἐστιν ὑμῖν.
\vs{13}Καὶ ταῦτα, ἃ βδελύξεσθε ἀπὸ τῶν πετεινῶν, καὶ οὐ βρωθήσεται, βδέλυγμά ἐστι· τὸν ἀετὸν, καὶ τὸν γρύπα, καὶ τὸν ἁλιαίετον,
\vs{14}καὶ τὸν γύπα, καὶ τὸν ἴκτινον καὶ τὰ ὅμοια αὐτῷ.
\vs{15}Καὶ στρουθὸν, καὶ γλαῦκα, καὶ λάρον, καὶ τὰ ὅμοια αὐτῷ·
\vs{16}Καὶ πάντα κόρακα, καὶ τὰ ὅμοια αὐτῷ· καὶ ἱέρακα, καὶ τὰ ὅμοια αὐτῷ·
\vs{17}καὶ νυκτικόρακα, καὶ καταράκτην, καὶ ἴβιν,
\vs{18}καὶ πορφυρίωνα, καὶ πελεκᾶνα, καὶ κύκνον,
\vs{19}καὶ ἐρωδιὸν, καὶ χαράδριον, καὶ τὰ ὅμοια αὐτῷ· καὶ ἔποπα, καὶ νυκτερίδα.
\vs{20}Καὶ πάντα τὰ ἑρπετὰ τῶν πετεινῶν, ἃ πορεύεται ἐπὶ τέσσαρα, βδελύγματά ἐστιν ὑμῖν.
\vs{21}Ἀλλὰ ταῦτα φάγεσθε ἀπὸ τῶν ἑρπετῶν τῶν πετεινῶν, ἃ πορεύεται ἐπὶ τέσσαρα, ἃ ἔχει σκέλη ἀνώτερον τῶν ποδῶν αὐτοῦ, πηδᾷν ἐν αὐτοῖς ἐπὶ τῆς γῆς.
\vs{22}Καὶ ταῦτα φάγεσθε ἀπʼ αὐτῶν· τὸν βροῦχον, καὶ τὰ ὅμοια αὐτῷ· καὶ τὸν ἀττάκην, καὶ τὰ ὅμοια αὐτῷ· καὶ ὀφιομάχην, καὶ τὰ ὅμοια αὐτῷ· καὶ τὴν ἀκρίδα, καὶ τὰ ὅμοια αὐτῇ·
\vs{23}Πᾶν ἑρπετὸν ἀπὸ τῶν πετεινῶν, οἷς εἰσι τέσσαρες πόδες, βδελύγματά ἐστιν ὑμῖν,
\vs{24}καὶ ἐν τούτοις μιανθήσεσθε· πᾶς ὁ ἁπτόμενος τῶν θνησιμαίων αὐτῶν, ἀκάθαρτος ἔσται ἕως ἑσπέρας.
\vs{25}Καὶ πᾶς ὁ αἴρων τῶν θνησιμαίων αὐτῶν, πλυνεῖ τὰ ἱμάτια αὐτοῦ, καὶ ἀκάθαρτος ἔσται ἕως ἑσπέρας.
\vs{26}Καὶ ἐν πᾶσι τοῖς κτήνεσιν ὅ ἐστι διχηλοῦν ὁπλὴν, καὶ ὀνυχιστῆρας ὀνυχίζει, καὶ μηρυκισμὸν οὐ μηρυκᾶται, ἀκάθαρτα ἔσονται ὑμῖν· πᾶς ὁ ἁπτόμενος τῶν θνησιμαίων αὐτῶν, ἀκάθαρτος ἔσται ἕως ἑσπέρας.
\vs{27}Καὶ πᾶς ὃς πορεύεται ἐπὶ χειρῶν ἐν πᾶσι τοῖς θηρίοις, ἃ πορεύεται ἐπὶ τέσσαρα, ἀκάθαρτά ἐστιν ὑμῖν· πᾶς ὁ ἁπτόμενος τῶν θνησιμαίων αὐτῶν, ἀκάθαρτος ἔσται ἕως ἑσπέρας.
\vs{28}Καὶ ὁ αἴρων τῶν θνησιμαίων αὐτῶν, πλυνεῖ τὰ ἱμάτια αὐτοῦ, καὶ ἀκάθαρτος ἔσται ἕως ἑσπέρας· ἀκάθαρτα ταῦτά ἔστιν ὑμῖν.

\vs{29}Καὶ ταῦτα ὑμῖν ἀκάθαρτα ἀπὸ τῶν ἑρπετῶν τῶν ἐπὶ τῆς γῆς· ἡ γαλὴ, καὶ ὁ μῦς, καὶ ὁ κροκόδειλος ὁ χερσαῖος,
\vs{30}μυγάλη, καὶ χαμαιλέων, καὶ χαλαβώτης, καὶ σαῦρα, καὶ ἀσπάλαξ.
\vs{31}Ταῦτα ἀκάθαρτα ὑμῖν ἀπὸ πάντων τῶν ἑρπετῶν τῶν ἐπὶ τῆς γῆς· πᾶς ὁ ἁπτόμενος αὐτῶν τεθνηκότων, ἀκάθαρτος ἔσται ἕως ἑσπέρας.
\vs{32}Καὶ πᾶν ἐφʼ ὃ ἂν ἐπιπέσῃ ἀπʼ αὐτῶν ἐπʼ αὐτὸ τεθνηκότων αὐτῶν, ἀκάθαρτον ἔσται ἀπὸ παντὸς σκεύους ξυλίνου ἢ ἱματίου ἢ δέρματος ἢ σάκκου· πᾶν σκεῦος ὃ ἂν ποιηθῇ ἔργον ἐν αὐτῷ, εἰς ὕδωρ βαφήσεται, καὶ ἀκάθαρτον ἔσται ἕως ἑσπέρας· καὶ καθαρὸν ἔσται.
\vs{33}Καὶ πᾶν σκεῦος ὀστράκινον εἰς ὃ ἐὰν πέσῃ ἀπὸ τούτων ἔνδον, ὅσα ἐὰν ἔνδον ᾖ, ἀκάθαρτα ἔσται, καὶ αὐτὸ συντριβήσεται.
\vs{34}Καὶ πᾶν βρῶμα, ὃ ἔσθεται, εἰς ὃ ἂν ἐπέλθῃ ἐπʼ αὐτὸ ὕδωρ, ἀκάθαρτον ἔσται· καὶ πᾶν ποτὸν, ὃ πίνεται ἐν παντὶ ἀγγείῳ, ἀκάθαρτον ἔσται.
\vs{35}Καὶ πᾶν ὃ ἐὰν ἐπιπέσῃ ἀπὸ τῶν θνησιμαίων αὐτῶν ἐπʼ αὐτό, ἀκάθαρτον ἔσται· κλίβανοι καὶ χυτρόποδες καθαιρεθήσονται· ἀκάθαρτα ταῦτά ἐστι, καὶ ἀκάθαρτα ταῦτα ὑμῖν ἔσονται.
\vs{36}Πλὴν πηγῶν ὑδάτων καὶ λάκκου καὶ συναγωγῆς ὕδατος, ἔσται καθαρόν· ὁ δὲ ἁπτόμενος τῶν θνησιμαίων αὐτῶν, ἀκάθαρτος ἔσται.
\vs{37}Ἐὰν δὲ ἐπιπέσῃ ἀπὸ τῶν θνησιμαίων αὐτῶν ἐπὶ πᾶν σπέρμα σπόριμον, ὃ σπαρήσεται, καθαρὸν ἔσται.
\vs{38}Ἐὰν δὲ ἐπιχυθῇ ὕδωρ ἐπὶ πᾶν σπέρμα, καὶ ἐπιπέσῃ τῶν θνησιμαίων αὐτῶν ἐπʼ αὐτό, ἀκάθαρτόν ἐστιν ὑμῖν.
\vs{39}Ἐὰν δὲ ἀποθάνῃ τῶν κτηνῶν, ὅ ἐστιν ὑμῖν φαγεῖν τοῦτο, ὁ ἁπτόμενος τῶν θνησιμαίων αὐτῶν, ἀκάθαρτος ἔσται ἕως ἑσπέρας·
\vs{40}καὶ ὁ ἐσθίων ἀπὸ τῶν θνησιμαίων τούτων, πλυνεῖ τὰ ἱμάτια, καὶ ἀκάθαρτος ἔσται ἕως ἑσπέρας· καὶ ὁ αἴρων ἀπὸ θνησιμαίων αὐτῶν, πλυνεῖ τὰ ἱμάτια, καὶ λούσεται ὕδατι, καὶ ἀκάθαρτος ἔσται ἕως ἑσπέρας.
\vs{41}Καὶ πᾶν ἑρπετὸν, ὃ ἕρπει ἐπὶ τῆς γῆς, βδέλυγμα ἔσται τοῦτο ὑμῖν· οὐ βρωθήσεται.
\vs{42}Καὶ πᾶς ὁ πορευόμενος ἐπὶ κοιλίας, καὶ πᾶς ὁ πορευόμενος ἐπὶ τέσσαρα διαπαντός, ὃ πολυπληθεῖ ποσὶν ἐν πᾶσι τοῖς ἑρπετοῖς τοῖς ἕρπουσιν ἐπὶ τῆς γῆς, οὐ φάγεσθε αὐτὸ, ὅτι βδέλυγμα ὑμῖν ἐστι.
\vs{43}Καὶ οὐ μὴ βδελύξητε τὰς ψυχὰς ὑμῶν ἐν πᾶσι τοῖς ἑρπετοῖς τοῖς ἕρπουσιν ἐπὶ τῆς γῆς, καὶ οὐ μιανθήσεσθε ἐν τούτοις, καὶ οὐκ ἀκάθαρτοι ἔσεσθε ἐν αὐτοῖς,
\vs{44}ὅτι ἐγώ εἰμι Κύριος ὁ Θεὸς ὑμῶν· καὶ ἁγιασθήσεσθε, καὶ ἅγιοι ἔσεσθε, ὅτι ἅγιός εἰμι ἐγὼ Κύριος ὁ Θεὸς ὑμῶν· καὶ οὐ μιανεῖτε τὰς ψυχὰς ὑμῶν ἐν πᾶσι τοῖς ἑρπετοῖς τοῖς κινουμένοις ἐπὶ τῆς γῆς,
\vs{45}ὅτι ἐγώ εἰμι Κύριος ὁ ἀναγαγὼν ὑμᾶς ἐκ γῆς Αἰγύπτου εἶναι ὑμῶν Θεός· καὶ ἔσεσθε ἅγιοι, ὅτι ἅγιός εἰμι ἐγὼ Κύριος.
\vs{46}Οὗτος ὁ νόμος περὶ τῶν κτηνῶν καὶ τῶν πετεινῶν καὶ πάσης ψυχῆς τῆς κινουμένης ἐν τῷ ὕδατι, καὶ πάσης ψυχῆς ἑρπούσης ἐπὶ τῆς γῆς,
\vs{47}διαστεῖλαι ἀναμέσον τῶν ἀκαθάρτων καὶ ἀναμέσον τῶν καθαρῶν, καὶ ἀναμέσον τῶν ζωογονούντων τὰ ἐσθιόμενα καὶ ἀναμέσον τῶν ζωογονούντων τὰ μὴ ἐσθιόμενα.

\ch{12}
Καὶ ἐλάλησε Κύριος πρὸς Μωυσῆν, λέγων,
\vs{2}λάλησον τοῖς υἱοῖς Ἰσραὴλ, καὶ ἐρεῖς πρὸς αὐτοὺς, γυνὴ ἥτις ἐὰν σπερματισθῇ, καὶ τέκῃ ἄρσεν, καὶ ἀκάθαρτος ἔσται ἑπτὰ ἡμέρας· κατὰ τὰς ἡμέρας τοῦ χωρισμοῦ τῆς ἀφέδρου αὐτῆς, ἀκάθαρτος ἔσται.
\vs{3}Καὶ τῇ ἡμέρᾳ τῇ ὀγδόῃ περιτεμεῖ τὴν σάρκα τῆς ἀκροβυστίας αὐτοῦ.
\vs{4}Καὶ τριάκοντα καὶ τρεῖς ἡμέρας καθήσεται ἐν αἵματι ἀκαθάρτῳ αὐτῆς· παντὸς ἁγίου οὐξ ἅψεται, καὶ εἰς τὸ ἁγιαστήριον οὐκ εἰσελεύσεται, ἕως ἂν πληρωθῶσιν αἱ ἡμέραι καθάρσεως αὐτῆς.
\vs{5}Ἐὰν δὲ θῆλυ τέκῃ, καὶ ἀκάθαρτος ἔσται δὶς ἑπτὰ ἡμέρας, κατὰ τὴν ἄφεδρον αὐτῆς· καὶ ἑξήκοντα ἡμέρας καὶ ἓξ καθεσθήσεται ἐν αἵματι ἀκαθάρτῳ αὐτῆς.

\vs{6}Καὶ ὅταν ἀναπληρωθῶσιν αἱ ἡμέραι καθάρσεως αὐτῆς ἐφʼ υἱῷ ἢ ἐπὶ θυγατρι, προσοίσει ἀμνὸν ἐνιαύσιον ἄμωμον εἰς ὁλοκαύτωμα, καὶ νοσσὸν περιστερᾶς ἢ τρυγόνα περὶ ἁμαρτίας ἐπὶ τὴν θύραν τῆς σκηνῆς τοῦ μαρτυρίου, πρὸς τὸν ἱερέα.
\vs{7}Καὶ προσοίσει αὐτὸν ἔναντι Κυρίου· καὶ ἐξιλάσεται περὶ αὐτῆς ὁ ἱερεὺς, καὶ καθαριεῖ αὐτὴν ἀπὸ τῆς πηγῆς τοῦ αἵματος αὐτῆς· οὗτος ὁ νόμος τῆς τικτούσης ἄρσεν ἢ θῆλυ.
\vs{8}Ἐὰν δὲ μὴ εὑρίσκῃ ἡ χεὶρ αὐτῆς τὸ ἱκανὸν εἰς ἀμνὸν, καὶ λήψεται δύο τρυγόνας ἢ δύο νοσσοὺς περιστερῶν, μίαν εἰς ὁλοκαύτωμα, καὶ μίαν περὶ ἁμαρτίας· καὶ ἐξιλάσεται περὶ αὐτῆς ὁ ἱερεὺς, καὶ καθαρισθήσεται.

\ch{13}
Καὶ ἐλάλησε Κύριος πρὸς Μωυσῆν καὶ Ἀαρὼν, λέγων,
\vs{2}ἀνθρώπῳ ἐάν τινι γένηται ἐν δέρματι χρωτὸς αὐτοῦ οὐλὴ σημασίας τηλαυγὴς, καὶ γένηται ἐν δέρματι χρωτὸς αὐτοῦ ἁφὴ λέπρας, ἀχθήσεται πρὸς Ἀαρὼν τὸν ἱερέα, ἢ ἕνα τῶν υἱῶν αὐτοῦ τῶν ἱερέων.
\vs{3}Καὶ ὄψεται ὁ ἱερεὺς τὴν ἁφὴν ἐν δέρματι τοῦ χρωτὸς αὐτοῦ, καὶ ἡ θρὶξ ἐν τῇ ἁφῇ μεταβάλῃ λευκὴ, καὶ ἡ ὄψις τῆς ἁφῆς ταπεινὴ ἀπὸ τοῦ δέρματος τοῦ χρωτὸς, ἁφὴ λέπρας ἐστί· καὶ ὄψεται ὁ ἱερεὺς, καὶ μιανεῖ αὐτόν.
\vs{4}Ἐὰν δὲ καὶ τηλαυγὴς λευκὴ ἦ ἐν τῷ δέρματι τοῦ χρωτὸς αὐτοῦ, καὶ ταπεινὴ μὴ ἦ ἡ ὄψις αὐτῆς ἀπὸ τοῦ δέρματος, καὶ ἡ θρὶξ αὐτοῦ οὐ μετέβαλε τρίχα λευκὴν, αὐτὴ δέ ἐστιν ἀμαυρὰ, καὶ ἀφοριεῖ ὁ ἱερεὺς τὴν ἁφὴν ἑπτὰ ἡμέρας.
\vs{5}Καὶ ὄψεται ὁ ἱερεὺς τὴν ἁφὴν τῇ ἡμέρᾳ τῇ ἑβδόμῃ· καὶ ἰδοὺ ἡ ἁφὴ μένει ἐναντίον αὐτοῦ, οὐ μετέπεσεν ἡ ἁφὴ ἐν τῷ δέρματι, καὶ ἀφοριεῖ αὐτὸν ὁ ἱερεὺς ἑπτὰ ἡμέρας τοδεύτερον.
\vs{6}Καὶ ὄψεται ὁ ἱερεὺς αὐτὸν τῇ ἡμέρᾳ τῇ ἑβδόμῃ τοδεύτερον· καὶ ἰδοὺ ἀμαυρὰ ἡ ἁφή, οὐ μετέπεσεν ἡ ἁφὴ ἐν τῷ δέρματι· καὶ καθαριεῖ αὐτὸν ὁ ἱερεύς, σημασία γάρ ἐστι· καὶ πλυνάμενος τὰ ἱμάτια αὐτοῦ, καθαρὸς ἔσται.
\vs{7}Ἐὰν δὲ μεταβαλοῦσα μεταπέσῃ ἡ σημασία ἐν τῷ δέρματι, μετὰ τὸ ἰδεῖν αὐτὸν τὸν ἱερέα τοῦ καθαρίσαι αὐτόν, καὶ ὀφθήσεται τοδεύτερον τῷ ἱερεῖ.
\vs{8}Καὶ ὄψεται αὐτὸν ὁ ἱερεύς, καὶ ἰδοὺ μετέπεσεν ἡ σημασία ἐν τῷ δέρματι, καὶ μιανεῖ αὐτὸν ὁ ἱερεύς· λέπρα ἐστί.

\vs{9}Καὶ ἁφὴ λέπρας ἐὰν γένηται ἐν ἀνθρώπῳ, και ἥξει πρὸς τὸν ἱερέα·
\vs{10}Καὶ ὄψεται ὁ ἱερεὺς, καὶ ἰδοὺ οὐλὴ λευκὴ ἐν τῷ δέρματι, καὶ αὕτη μετέβαλε τρίχα λευκὴν, καὶ ἀπὸ τοῦ ὑγιοῦς τῆς σαρκὸς τῆς ζώσης ἐν τῇ οὐλῇ.
\vs{11}Λέπρα παλαιουμένη ἐστὶν ἐν τῷ δέρματι τοῦ χρωτός, καὶ μιανεῖ αὐτὸν ὁ ἱερεὺς, καὶ ἀφοριεῖ αὐτὸν, ὅτι ἀκάθαρτός ἐστιν.

\vs{12}Ἐὰν δὲ ἀνθοῦσα ἐξανθήσῃ λέπρα ἐν τῷ δέρματι, καὶ καλύψῃ ἡ λέπρα πᾶν τὸ δέρμα τῆς ἁφῆς ἀπὸ κεφαλῆς ἕως ποδῶν, καθʼ ὅλην τὴν ὅρασιν τοῦ ἱερέως·
\vs{13}Καὶ ὄψεται ὁ ἱερεὺς, καὶ ἰδοὺ ἐκάλυψεν ἡ λέπρα πᾶν τὸ δέρμα τοῦ χρωτός· καὶ καθαριεῖ αὐτὸν ὁ ἱερεὺς τὴν ἁφήν, ὅτι πᾶν μετέβαλε λευκὸν, καθαρόν ἐστι.
\vs{14}Καὶ ᾗ ἂν ἡμέρᾳ ὀφθῇ ἐν αὐτῷ χρὼς ζῶν, μιανθήσεται.
\vs{15}Καὶ ὅψεται ὁ ἱερεὺς τὸν χρῶτα τὸν ὑγιῆ, καὶ μιανεῖ αὐτὸν ὁ χρὼς ὁ ὑγιὴς, ὅτι ἀκάθαρτός ἐστι· λέπρα ἐστίν.
\vs{16}Ἐὰν δὲ ἀποκαταστῇ ὁ χρὼς ὁ ὑγιὴς, καὶ μεταβάλῃ λευκὴ, καὶ ἐλεύσεται πρὸς τὸν ἱερέα·
\vs{17}Καὶ ὄψεται ὁ ἱερεὺς, καὶ ἰδοὺ μετέβαλεν ἡ ἁφὴ εἰς τὸ λευκόν, καὶ καθαριεῖ ὁ ἱερεὺς τὴν ἁφήν· καθαρός ἐστι.

\vs{18}Καὶ σὰρξ ἐὰν γένηται ἐν τῷ δέρματι αὐτοῦ ἕλκος, καὶ ὑγιασθῇ,
\vs{19}καὶ γένηται ἐν τῷ τόπῳ τοῦ ἕλκους οὐλὴ λευκὴ, ἢ τηλαυγὴς λευκαίνουσα, ἢ πυῤῥίζουσα, καὶ ὀφθήσεται τῷ ἱερεῖ·
\vs{20}Καὶ ὄψεται ὁ ἱερεὺς, καὶ ἰδοὺ ἡ ὄψις ταπεινοτέρα τοῦ δέρματος, καὶ ἡ θρὶξ αὐτῆς μετέβαλεν εἰς λευκὴν, καὶ μιανεῖ αὐτὸν ὁ ἱερεὺς, ὅτι λέπρα ἐστίν· ἐν τῷ ἕλκει ἐξήνθησεν.
\vs{21}Ἐὰν δὲ ἴδῃ ὁ ἱερεὺς, καὶ ἰδοὺ οὐκ ἔστιν ἐν αὐτῷ θρὶξ λευκὴ, καὶ ταπεινὸν μὴ ᾖ ἀπὸ τοῦ δέρματος τοῦ χρωτὸς, καὶ αὐτὴ ᾖ ἀμαυρά, καὶ ἀφοριεῖ αὐτὸν ὁ ἱερεὺς ἑπτὰ ἡμέρας.
\vs{22}Ἐὰν δὲ διαχύσει διαχέηται ἐν τῷ δέρματι, καὶ μιανεῖ αὐτὸν ὁ ἱερεὺς, ἁφὴ λέπρας ἐστίν· ἐν τῷ ἕλκει ἐξήνθησεν·
\vs{23}Ἐὰν δὲ κατὰ χώραν μείνῃ τὸ τηλαύγημα καὶ μὴ διαχέηται, οὐλὴ τοῦ ἕλκους ἐστὶ, καὶ καθαριεῖ αὐτὸν ὁ ἱερεύς.

\vs{24}Καὶ σὰρξ ἐὰν γένηται ἐν τῷ δέρματι αὐτοῦ κατάκαυμα πυρὸς, καὶ γένηται ἐν τῷ δέρματι αὐτοῦ τὸ ὑγιασθὲν τοῦ κατακαύματος αὐγάζον τηλαυγὲς λευκὸν, ὑποπυῤῥίζον, ἢ ἔκλευκον·
\vs{25}Καὶ ὄψεται αὐτὸν ὁ ἱερεὺς, καὶ ἰδοὺ μετέβαλε θρὶξ λευκὴ εἰς τὸ αὐγάζον, καὶ ἡ ὄψις αὐτοῦ ταπεινὴ ἀπὸ τοῦ δέρματος, λέπρα ἐστίν· ἐν τῷ κατακαύματι ἐξήνθησε· καὶ μιανεῖ αὐτὸν ὁ ἱερεὺς, ἁφὴ λέπρας ἐστίν.
\vs{26}Ἐὰν δὲ ἴδῃ ὁ ἱερεὺς, καὶ ἰδοὺ οὐκ ἔστιν ἐν τῷ· αὐγάζοντι θρὶξ λευκή, καὶ ταπεινὸν μὴ ᾖ ἀπὸ τοῦ δέρματος, αὐτὸ δὲ ἀμαυρὸν, καὶ ἀφοριεῖ αὐτὸν ὁ ἱερεὺς ἑπτὰ ἡμέρας.
\vs{27}Καὶ ὄψεται αὐτὸν ὁ ἱερεὺς τῇ ἡμέρᾳ τῇ ἑβδόμῃ· ἐὰν δὲ διαχύσει διαχέηται ἐν τῷ δέρματι, καὶ μιανεῖ αὐτὸν ὁ ἱερεὺς, ἁφὴ λέπρας ἐστίν· ἐν τῷ ἕλκει ἐξήνθησεν.
\vs{28}Ἐὰν δὲ κατὰ χώραν μείνῃ τὸ αὐγάζον, καὶ μὴ διαχυθῇ ἐν τῷ δέρματι, αὐτὴ δὲ ἀμαυρὰ ᾖ, οὐλὴ τοῦ κατακαύματός ἐστι, καὶ καθαριεῖ αὐτὸν ὁ ἱερεύς· ὁ γὰρ χαρακτὴρ τοῦ κατακαύματός ἐστι.

\vs{29}Καὶ ἀνδρὶ ἢ γυναικὶ ἐὰν γένηται ἐν αὐτοῖς ἁφὴ λέπρας ἐν τῇ κεφαλῇ ἢ ἐν τῷ πώγωνι·
\vs{30}Καὶ ὄψεται ὁ ἱερεὺς τὴν ἁφὴν, καὶ ἰδοὺ ἡ ὄψις αὐτῆς ἐγκοιλοτέρα τοῦ δέρματος, ἐν αὐτῇ δὲ θρὶξ ξανθίζουσα λεπτὴ, καὶ μιανεῖ αὐτὸν ὁ ἱερεύς· θραῦσμά ἐστι, λέπρα τῆς κεφαλῆς ἢ λέπρα τοῦ πώγωνός ἐστι.
\vs{31}Καὶ ἐὰν ἴδῃ ὁ ἱερεὺς τὴν ἁφὴν τοῦ θραύσματος, καὶ ἰδοὺ οὐχ ἡ ὄψις ἐγκοιλοτέρα τοῦ δέρματος, καὶ θρὶξ ξανθίζουσα οὐκ ἔστιν ἐν αὐτῇ, καὶ ἀφοριεῖ ὁ ἱερεὺς τὴν ἁφὴν τοῦ θραύσματος ἑπτὰ ἡμέρας.
\vs{32}Καὶ ὄψεται ὁ ἱερεὺς τὴν ἁφὴν τῇ ἡμέρᾳ τῇ ἑβδόμῃ, καὶ ἰδοὺ οὐ διεχύθη τὸ θραῦσμα, καὶ θρὶξ ξανθίζουσα οὐκ ἔστιν ἐν αὐτῇ, καὶ ἡ ὄψις τοῦ θραύσματος οὐκ ἔστι κοίλη ἀπὸ τοῦ δέρματος·
\vs{33}Καὶ ξυρηθήσεται τὸ δέρμα, τὸ δὲ θραῦσμα οὐ ξυρηθήσεται, καὶ ἀφοριεῖ ὁ ἱερεὺς τὸ θραῦσμα ἑπτὰ ἡμέρας τὸ δεύτερον.
\vs{34}Καὶ ὄψεται ὁ ἱερεὺς τὸ θραῦσμα τῇ ἡμέρᾳ τῇ ἑβδόμῃ, καὶ ἰδοὺ οὐ διεχύθη τὸ θραῦσμα ἐν τῷ δέρματι μετὰ τὸ ξυρηθῆναι αὐτόν, καὶ ἡ ὄψις τοῦ θραύσματος οὐκ ἔστιν κοίλη ἀπὸ τοῦ δέρματος, καὶ καθαριεῖ αὐτὸν ὁ ἱερεὺς, καὶ πλυνάμενος τὰ ἱμάτια, καθαρὸς ἔσται.
\vs{35}Ἐὰν δὲ διαχύσει διαχέηται τὸ θραῦσμα ἐν τῷ δέρματι μετὰ τὸ καθαρισθῆναι αὐτόν·
\vs{36}Καὶ ὄψεται ὁ ἱερεὺς, καὶ ἰδοὺ διακέχυται τὸ θραῦσμα ἐν τῷ δέρματι, οὐκ ἐπισκέψεται ὁ ἱερεὺς περὶ τῆς τριχὸς τῆς ξανθῆς, ὅτι ἀκάθαρτός ἐστιν.
\vs{37}Ἐὰν δὲ ἐνώπιον μείνῃ ἐπὶ χώρας τὸ θραῦσμα, καὶ θρὶξ μέλαινα ἀνατείλῃ ἐν αὐτῷ, ὑγίακε τὸ θραῦσμα, καθαρός ἐστι, καὶ καθαριεῖ αὐτὸν ὁ ἱερεύς.
\vs{38}Καὶ ἀνδρὶ ἢ γυναικὶ ἐὰν γένηται ἐν δέρματι τῆς σαρκὸς αὐτοῦ αὐγάματα αὐγάζοντα λευκανθίζοντα·
\vs{39}καὶ ὄψεται ὁ ἱερεὺς, καὶ ἰδοὺ ἐν δέρματι τῆς σαρκὸς αὐτοῦ αὐγάσματα αὐγάζοντα λευκαθίζοντα, ἀλφός ἐστιν· ἐξανθεῖ ἐν τῷ δέρματι τῆς σαρκὸς αὐτοῦ, καθαρός ἐστιν.
\vs{40}Ἐὰν δέ τινι μαδήσῃ ἡ κεφαλὴ αὐτοῦ, φαλακρός ἐστι, καθαρός ἐστιν.
\vs{41}Ἐὰν δὲ κατὰ πρόσωπον μαδήσῃ ἡ κεφαλὴ αὐτοῦ, ἀναφάλαντός ἐστι, καθαρός ἐστιν.
\vs{42}Ἐὰν δὲ γένηται ἐν τῷ φαλακρώματι αὐτοῦ ἢ ἐν τῷ ἀναφαλαντώματι αὐτοῦ ἁφὴ λευκὴ ἢ πυῤῥίζουσα, λέπρα ἐστὶν ἐν τῷ φαλακρώματι αὐτοὺ, ἢ ἐν τῷ ἀναφαλαντώματι αὐτοῦ·
\vs{43}καὶ ὄψεται αὐτὸν ὁ ἱερεύς, καὶ ἰδοὺ ἡ ὄψις τῆς ἁφῆς λευκὴ ἢ πυῤῥίζουσα ἐν τῷ φαλακρώματι αὐτοῦ ἢ ἐν τῷ ἀναφαλαντώματι αὐτοῦ, ὡς ωἶδος λέπρας ἐν δέρματι τῆς σαρκὸς αὐτοῦ·
\vs{44}Ἄνθρωπος λεπρός ἐστι· μιάνσει μιανεῖ αὐτὸν ὁ ἱερεὺς, ἐν τῇ κεφαλῇ αὐτοῦ ἡ ἁφὴ αὐτοῦ.
\vs{45}Καὶ ὁ λεπρὸς ἐν ᾧ ἐστιν ἡ ἁφὴ, τὰ ἱμάτια αὐτοῦ ἔστω παραλελυμένα, καὶ ἡ κεφαλὴ αὐτοῦ ἀκάλυπτος, καὶ περὶ τὸ στόμα αὐτοῦ περιβαλέσθω, καὶ ἀκάθαρτος κεκλήσεται.
\vs{46}Πάσας τὰς ἡμέρας, ὅσας ἐὰν ᾖ ἐπʼ αὐτὸν ἡ ἁφὴ, ἀκάθαρτος ὢν ἀκάθαρτος ἔσται· κεχωρισμένος καθήσεται, ἔξω τῆς παρεμβολῆς αὐτοῦ ἔσται ἡ διατριβή.

\vs{47}Καὶ ἱματίῳ ἐὰν γένηται ἁφὴ ἐν αὐτῷ λέπρας, ἐν ἱματίῳ ἐρέῳ, ἢ ἐν ἱματίῳ στυππυίνῳ,
\vs{48}ἢ ἐν στήμονι, ἢ ἐν κρόκῃ, ἢ ἐν τοῖς λινοῖς, ἢ ἐν τοῖς ἐρέοις, ἢ ἐν δέρματι, ἢ ἐν παντὶ ἐργασίμῳ δέρματι,
\vs{49}καὶ γένηται ἡ ἁφὴ χλωρίζουσα ἢ πυῤῥίζουσα ἐν τῷ δέρματι, ἢ ἐν τῷ ἱματίῳ, ἢ ἐν τῷ στήμονι, ἢ ἐν τῇ κρόκῃ, ἢ ἐν παντὶ σκεύει ἐργασίμῳ δέρματος, ἁφὴ λέπρας ἐστί· καὶ δείξει τῷ ἱερεῖ·
\vs{50}Καὶ ὄψεται ὁ ἱερεὺς τὴν ἁφήν, καὶ ἀφοριεῖ ὁ ἱερεὺς τὴν ἁφὴν ἑπτὰ ἡμέρας.
\vs{51}Καὶ ὄψεται ὁ ἱερεὺς τὴν ἁφὴν τῇ ἡμέρᾳ τῇ ἑβδόμῃ· ἐὰν δὲ διαχέηται ἡ ἁφὴ ἐν τῷ ἱματίῳ, ἢ ἐν τῷ στήμονι, ἢ ἐν τῇ κρόκῃ, ἢ ἐν τῷ δέρματι, κατὰ πάντα ὅσα ἐὰν ποιηθῇ δέρματα ἐν τῇ ἐργασίᾳ, λέπρα ἔμμονός ἐστιν ἡ ἁφὴ, ἀκάθαρτός ἐστι.
\vs{52}Κατακαύσει τὸ ἱμάτιον, ἢ τὸν στήμονα, ἢ τὴν κρόκην ἐν τοῖς ἐρέοις, ἢ ἐν τοῖς λινοῖς, ἢ ἐν παντὶ σκεύει δερματίνῳ, ἐν ᾧ ἂν ᾖ ἐν αὐτῷ ἡ ἁφὴ, ὅτι λέπρα ἔμμονός ἐστιν, ἐν πυρὶ κατακαυθήσεται.

\vs{53}Ἐὰν δὲ ἴδῃ ὁ ἱερεὺς, καὶ μὴ διαχέηται ἡ ἁφὴ ἐν τῷ ἱματίῳ, ἢ ἐν τῷ στήμονι, ἢ ἐν τῇ κρόκῃ, ἢ ἐν παντὶ σκεύει δερματίνῳ·
\vs{54}Καὶ συντάξει ὁ ἱερεύς, καὶ πλυνεῖ ἐφʼ οὗ ἐὰν ᾖ ἐπʼ αὐτοῦ ἡ ἁφὴ, καὶ ἀφοριεῖ ὁ ἱερεὺς τὴν ἁφὴν ἑπτὰ ἡμέρας τοδέυτερον.
\vs{55}Καὶ ὄψεται ὁ ἱερεὺς μετὰ τὸ πλυθῆναι αὐτὸ τὴν ἁφὴν, καὶ ἥδε οὐ μὴ μετέβαλεν ἡ ἁφὴ τὴν ὄψιν, καὶ ἡ ἁφὴ οὐ διαχεῖται, ἀκάθαρτόν ἐστιν, ἐν πυρὶ κατακαυθήσεται· ἐστήρικται ἐν τῷ ἱματίῳ, ἢ ἐν τῷ στήμονι, ἢ ἐν τῇ κρόκη.
\vs{56}Καὶ ἐὰν ἴδῃ ὁ ἱερεὺς, καὶ ᾖ ἀμαυρὰ ἡ ἁφὴ μετὰ τὸ πλυθήναι αὐτὸ, ἀποῤῥήξει αὐτὸ ἀπὸ τοῦ ἱματίου, ἢ ἀπὸ τοῦ στήμονος, ἢ ἀπὸ τῆς κρόκης, ἢ ἀπὸ τοῦ δέρματος.
\vs{57}Ἐὰν δὲ ὀφθῇ ἔτι ἐν τῷ ἱματίῳ, ἢ ἐν τῷ στήμονι, ἢ ἐν τῇ κρόκῃ, ἢ ἐν παντὶ σκεύει δερματίνῳ, λέπρα ἐξανθοῦσά ἐστιν ἐν πυρὶ κατακαυθήσεται ἐν ᾧ ἐστιν ἡ ἁφή.
\vs{58}Καὶ τὸ ἱμάτιον, ἢ ὁ στήμων, ἢ ἡ κρόκη, ἢ πᾶν σκεῦος δερμάτινον, ὃ πλυθήσεται, καὶ ἀποστήσεται ἀπʼ αὐτοῦ ἡ ἁφὴ, καὶ πλυθήσεται τὸ δέυτερον, καὶ καθαρὸν ἔσται.
\vs{59}Οὗτος ὁ νόμος ἁφῆς λέπρας ἱματίου ἐρέου, ἢ στυππυίνου, ἢ στήμονος, ἢ κρόκης, ἢ παντὸς σκεύους δερματίνου, εἰς τὸ καθαρίσαι αὐτὸ, ἢ μιᾶναι αὐτό.

\ch{14}
Καὶ ἐλάλησε Κύριος πρὸς Μωυσῆν, λέγων,
\vs{2}οὗτος ὁ νὁμος τοῦ λεπροῦ· ᾗ ἂν ἡμέρᾳ καθαρισθῇ, καὶ προσαχθήσεται πρὸς τὸν ἱερέα.
\vs{3}Καὶ ἐξελεύσεται ὁ ἱερεὺς ἔξω τῆς παρεμβολῆς, καὶ ὄψεται ὁ ἱερεὺς, καὶ ἰδοὺ ἰᾶται ἡ ἁφὴ τῆς λέπρας ἀπὸ τοῦ λεπροῦ.
\vs{4}Καὶ προστάξει ὁ ἱερεὺς, καὶ λήψονται τῷ κεκαθαρισμένῳ δύο ὀρνίθια ζῶντα καθαρὰ, καὶ ξύλον κέδρινον, καὶ κεκλωσμένον κόκκινον, καὶ ὕσσωπον.
\vs{5}Καὶ προστάξει ὁ ἱερεὺς, καὶ σφάξουσι τὸ ὀρνίθιον τὸ ἓν εἰς ἀγγεῖον ὀστράκινον ἐφʼ ὕδατι ζῶντι.
\vs{6}Καὶ τὸ ὀρνίθιον τὸ ζῶν λήμψεται αὐτὸ, καὶ τὸ ξύλον τὸ κέδρινον, καὶ τὸ κλωστὸν κόκκινον, καὶ τὸν ὕσσωπον, καὶ βάψει αὐτὰ καὶ τὸ ὀρνίθιον τὸ ζῶν εἰς τὸ αἷμα τοῦ ὀρνιθίου τοῦ σφαγέντος ἐφʼ ὕδατι ζῶντι.
\vs{7}Καὶ περιῤῥανεῖ ἐπὶ τὸν καθαρισθέντα ἀπὸ τῆς λέπρας ἑπτάκις, καὶ καθαρὸς ἔσται· καὶ ἐξαποστελεῖ τὸ ὀρνίθιον τὸ ζῶν εἰς τὸ πεδίον.
\vs{8}Καὶ πλυνεῖ ὁ καθαρισθεὶς τὰ ἱμάτια αὐτοῦ, καὶ ξυρηθήσεται αὐτοῦ πᾶσαν τὴν τρίχα, καὶ λούσεται ἐν ὕδατι, καὶ καθαρὸς ἔσται· καὶ μετὰ ταῦτα εἰσελεύσεται εἰς τὴν παρεμβολὴν, καὶ διατρίψει ἔξω τοῦ οἴκου αὐτοῦ ἑπτὰ ἡμέρας.
\vs{9}Καὶ ἔσται τῇ ἡμέρᾳ τῇ ἑβδόμῃ, ξυρηθήσεται πᾶσαν τὴν τρίχα αὐτοῦ, τὴν κεφαλὴν αὐτοῦ, καὶ τὸν πώγωνα, καὶ τὰς ὀφρῦς, καὶ πᾶσαν τὴς τρίχα αὐτοῦ ξυρηθήσεται· καὶ πλυνεῖ τὰ ἱμάτια, καὶ λούσεται τὸ σῶμα αὐτοῦ ὕδατι, καὶ καθαρὸς ἔσται.
\vs{10}Καὶ τῇ ἡμέρᾳ τῇ ὀγδόῃ λήψεται δύο ἀμνοὺς ἀμώμους ἑνιαυσίους, καὶ πρόβατον ἄμωμον ἑνιαύσιον, καὶ τρία δέκατα σεμιδάλεως εἰς θυσίαν πεφυραμένης ἐν ἐλαίῳ, καὶ κοτύλην ἐλαίου μίαν.
\vs{11}Καὶ στήσει ὁ ἱερεὺς ὁ καθαρίζων, τὸν ἄνθρωπον τὸν καθαριζόμενον, καὶ ταῦτα ἔναντι Κυρίου, ἐπὶ τὴν θύραν τῆς σκηνῆς τοῦ μαρτυρίου.
\vs{12}Καὶ λήψεται ὁ ἱερεὺς τὸν ἀμνὸν τὸν ἕνα, καὶ προσάξει αὐτὸν τῆς πλημμελείας, καὶ τὴν κοτύλην τοῦ ἐλαίου, καὶ ἀφοριεῖ αὐτὰ ἀφόρισμα ἔναντι Κυρίου.
\vs{13}Καὶ σφάξουσι τὸν ἀμνὸν ἐν τόπῳ, οὗ σφάζουσι τὰ ὁλοκαυτώματα, καὶ τὰ περὶ ἁμαρτίας, ἐν τόπῳ ἁγίῳ· ἔστι γὰρ τὸ περὶ ἁμαρτίας, ὥσπερ τὸ τῆς πλημμελείας ἐστὶ τῷ ἱερει· ἅγια ἁγίων ἐστί.
\vs{14}Καὶ λήψεται ὁ ἱερεὺς ἀπὸ τοῦ αἵματος τοῦ τῆς πλημμελείας, καὶ ἐπιθήσει ὁ ἱερεὺς ἐπὶ τὸν λοβὸν τοῦ ὠτὸς τοῦ καθαριζομένου τοῦ δεξιοῦ, καὶ ἐπὶ τὸ ἄκρον τῆς χειρὸς τῆς δεξιᾶς, καὶ ἐπὶ τὸ ἄκρον τοῦ ποδὸς τοῦ δεξιοῦ.
\vs{15}Καὶ λαβὼν ὁ ἱερεὺς ἀπὸ τῆς κοτύλης τοῦ ἐλαίου, ἐπιχεεῖ ἐπὶ τὴν χεῖρα τοῦ ἱερέως τὴν ἀριστεράν.
\vs{16}Καὶ βάψει τὸν δάκτυλον τὸν δεξιὸν ἀπὸ τοὺ ἐλαίου τοῦ ὄντος ἐπὶ τῆς χειρὸς αὐτοῦ τῆς ἀριστερᾶς· καὶ ῥανεῖ τῷ δακτύλῳ ἑπτάκις ἔναντι Κυρίου.
\vs{17}Τὸ δὲ καταλειφθὲν ἔλαιον τὸ ὂν ἐν τῇ χειρὶ, ἐπιθήσει ὁ ἱερεὺς ἐπὶ τὸν λοβὸν τοῦ ὠτὸς τοῦ καθαριζομένου τοῦ δεξιοῦ, καὶ ἐπὶ τὸ ἄκρον τῆς χειρὸς τῆς δεξιᾶς, καὶ ἐπὶ τὸ ἄκρον τοῦ ποδὸς τοῦ δεξιοῦ, ἐπὶ τὸν τόπον τοῦ αἵματος τοῦ τῆς πλημμελείας.
\vs{18}Τὸ δὲ καταλειφθὲν ἔλαιον τὸ ἐπὶ τῆς χειρὸς τοῦ ἱερέως, ἐπιθήσει ὁ ἱερεὺς ἐπὶ τὴν κεφαλὴν τοῦ καθαρισθέντος· καὶ ἐξιλάσεται περὶ αὐτοῦ ὁ ἱερεὺς ἔναντι Κυρίου.
\vs{19}Καὶ ποιήσει ὁ ἱερεὺς τὸ περὶ τῆς ἁμαρτίας, καὶ ἐξιλάσεται ὁ ἱερεὺς περὶ τοῦ καθαριζομένου ἀπὸ τῆς ἁμαρτίας αὐτοῦ· καὶ μετὰ τοῦτο σφάξει ὁ ἱερεὺς τὸ ὁλοκαύτωμα.
\vs{20}Καὶ ἀνοίσει ὁ ἱερεὺς τὸ ὁλοκαύτωμα, καὶ τὴν θυσίαν ἐπὶ τὸ θυσιαστήριον ἔναντι κυρίου· καὶ ἐξιλάσεται περὶ αὐτοῦ ὁ ἱερεὺς, καὶ καθαρισθήσεται.
\vs{21}Ἐὰν δὲ πένηται, καὶ ἡ χεὶρ αὐτοῦ μὴ εὑρίσκῃ, λήψεται ἀμνὸν ἕνα εἰς ὃ ἐπλημμέλησεν εἰς ἀφαίρεμα, ὥστε ἐξιλάσασθαι περὶ αὐτοῦ, καὶ δέκατον σεμιδάλεως πεφυραμένης ἐν ἐλαίῳ εἰς θυσίαν, καὶ κοτύλην ἐλαίου μίαν,
\vs{22}καὶ δύο τρυγόνας, ἢ δύο νοσσοὺς περιστερῶν, ὅσα εὗρεν ἡ χεὶρ αὐτοῦ, καὶ ἔσται ἡ μία περὶ ἁμαρτίας, καὶ ἡ μία εἰς ὁλοκαύτωμα.
\vs{23}Καὶ προσοίσει αὐτὰ τῇ ἡμέρᾳ τῇ ὀγδόῃ, εἰς τὸ καθαρίσαι αὐτὸν, πρὸς τὸν ἱερέα, ἐπὶ τὴν θύραν τῆς σκηνῆς τοῦ μαρτυρίου ἔναντι Κυρίου.
\vs{24}Καὶ λαβὼν ὁ ἱερεὺς τὸν ἀμνὸν τῆς πλημμελείας, καὶ τὴν κοτύλην τοῦ ἐλαίου, ἐπιθήσει αὐτὰ ἐπίθεμα ἔναντι Κυρίου.
\vs{25}Καὶ σφάξει τὸν ἀμνὸν τὸν τῆς πλημμελείας, καὶ λήψεται ὁ ἱερεὺς ἀπὸ τοῦ αἵματος τοῦ τῆς πλημμελείας, καὶ ἐπιθήσει ἐπὶ τὸν λοβὸν τοῦ ὠτὸς τοῦ καθαριζομένου τοῦ δεξιοῦ, καὶ ἐπὶ τὸ ἄκρον τῆς χειρὸς τῆς δεξιᾶς, καὶ ἐπὶ τὸ ἄκρον τοῦ ποδὸς τοῦ δεξιοῦ.
\vs{26}Καὶ ἀπὸ τοῦ ἐλαίου ἐπιχεεῖ ὁ ἱερεὺς ἐπὶ τὴν χεῖρα τοῦ ἱερέως τὴν ἀριστεράν.
\vs{27}Καὶ ῥανεῖ ὁ ἱερεὺς τῷ δακτύλῳ τῷ δεξιῷ ἀπὸ τοῦ ἐλαίου τοῦ ἐν τῇ χειρὶ αὐτοῦ τῇ ἀριστερᾷ ἑπτάκις ἔναντι Κυρίου.
\vs{28}Καὶ ἐπιθήσει ὁ ἱερεὺς ἀπὸ τοῦ ἐλαίου τοῦ ἐπὶ τῆς χειρὸς αὐτοῦ ἐπὶ τὸν λοβὸν τοῦ ὠτὸς τοῦ καθαριζομένου τοῦ δεξιοῦ, καὶ ἐπὶ τὸ ἄκρον τῆς χειρὸς αὐτοῦ τῆς δεξιᾶς, καὶ ἐπὶ τὸ ἄκρον τοῦ ποδὸς αὐτοῦ τοῦ δεξιοῦ, ἐπὶ τὸν τόπον τοῦ αἵματος τοῦ τῆς πλημμελείας.
\vs{29}Τὸ δὲ καταλειφθὲν ἀπὸ τοῦ ἐλαίου τὸ ὂν ἐπὶ τῆς χειρὸς τοῦ ἱερέως, ἐπιθήσει ἐπὶ τὴν κεφαλὴν τοῦ καθαρισθέντος· καὶ ἐξιλάσεται περὶ αὐτοῦ ὁ ἱερεὺς ἔναντι Κυρίου.

\vs{30}Καὶ ποιήσει μίαν ἀπὸ τῶν τρυγόνων ἢ ἀπὸ τῶν νοσσῶν τῶν περιστερῶν, καθότι εὗρεν αὐτοῦ ἡ χεὶρ,
\vs{31}τὴν μίαν περὶ ἁμαρτίας, καὶ τὴν μίαν εἰς ὁλοκαύτωμα σὺν τῇ θυσίᾳ· καὶ ἐξιλάσεται ὁ ἱερεὺς περὶ τοῦ καθαριζομένου ἔναντι Κυρίου.
\vs{32}Οὗτος ὁ νόμος ἐν ᾧ ἐστιν ἡ ἁφὴ τῆς λέπρας, καὶ τοῦ μὴ εὑρίσκοντος τῇ χειρὶ εἰς τὸν καθαρισμὸν αὐτοῦ.

\vs{33}Καὶ ἐλάλησε Κύριος πρὸς Μωυσῆν καὶ Ἀαρὼν, λέγων,
\vs{34}ὡς ἂν εἰσέλθητε εἰς τὴν γῆν τῶν Χαναναίων, ἣν ἐγὼ δίδωμι ὑμῖν ἐν κτήσει, καὶ δώσω ἁφὴν λέπρας ἐν ταῖς οἰκίαις τῆς γῆς τῆς ἐγκτήτου ὑμῖν·
\vs{35}καὶ ἥξει τίνος αὐτοῦ ἡ οἰκία, καὶ ἀναγγελεῖ τῷ ἱερεῖ, λέγων, ὥσπερ ἁφὴ ἑώραταί μοι ἐν τῇ οἰκίᾳ.
\vs{36}Καὶ προστάξει ὁ ἱερεὺς ἀποσκευάσαι τὴν οἰκίαν, πρὸ τοῦ εἰσελθόντα τὸν ἱερέα ἰδεῖν τὴν ἁφὴν, καὶ οὐ μὴ ἀκάθαρτα γένηται ὅσα ἂν ᾖ ἐν τῇ οἰκίᾳ· καὶ μετὰ ταῦτα εἰσελεύσεται ὁ ἱερεὺς καταμαθεῖν τὴν οἰκίαν.
\vs{37}Καὶ ὄψεται τὴν ἁφὴν, καὶ ἰδοὺ ἡ ἁφὴ ἐν τοῖς τοίχοις τῆς οἰκίας, κοιλάδας χλωριζούσας, ἢ πυῤῥιζούσας, καὶ ἡ ὄψις αὐτῶν ταπεινοτέρα τῶν τοίχων.
\vs{38}Καὶ ἐξελθὼν ὁ ἱερεὺς ἐκ τῆς οἰκίας ἐπὶ τὴν θύραν τῆς οἰκίας, καὶ ἀφοριεῖ ὁ ἱερεὺς τὴν οἰκίαν ἑπτὰ ἡμέρας.
\vs{39}Καὶ ἐπανήξει ὁ ἱερεὺς τῇ ἡμέρᾳ τῇ ἑβδόμῃ, καὶ ὄψεται τὴν οἰκίαν, καὶ ἰδοὺ διεχύθη ἡ ἁφὴ ἐν τοῖς τοίχοις τῆς οἰκίας.
\vs{40}Καὶ προστάξει ὁ ἱερεὺς, καὶ ἐξελοῦσι τοὺς λίθους ἐν οἷς ἐστιν ἡ ἁφὴ, καὶ ἐκβαλοῦσιν αὐτοὺς ἔξω τῆς πόλεως εἰς τόπον ἀκάθαρτον.
\vs{41}Καὶ τὴν οἰκίαν ἀποξύσουσιν ἔσωθεν κύκλῳ, καὶ ἐκχεοῦσι τὸν χοῦν τὸν ἀπεξυσμένον ἔξω τῆς πόλεως εἰς τόπον ἀκάθαρτον.
\vs{42}Καὶ λήψονται λίθους ἀπεξυσμένους ἑτέρους, καὶ ἀντιθήσουσιν ἀντὶ τῶν λίθων· καὶ χοῦν ἕτερον λήψονται, καὶ ἐξαλείψουσι τὴν οἰκίαν.
\vs{43}Ἐὰν δὲ ἐπέλθῃ πάλιν ἡ ἁφὴ, καὶ ἀνατείλῃ ἐν τῇ οἰκίᾳ μετὰ τὸ ἐξελεῖν τοὺς λίθους, καὶ μετὰ τὸ ἀποξυσθῆναι τὴν οἰκίαν, καὶ μετὰ τὸ ἐξαλειφθῆναι,
\vs{44}καὶ εἰσελεύσεται ὁ ἱερεὺς, καὶ ὄψεται εἰ διακέχυται ἡ ἁφὴ ἐν τῇ οἰκίᾳ, λέπρα ἔμμονός ἐστιν ἐν τῇ οἰκίᾳ, ἀκάθαρτός ἐστι.
\vs{45}Καὶ καθελοῦσι τὴν οἰκίαν, καὶ τὰ ξύλα αὐτῆς, καὶ τοὺς λίθους αὐτῆς, καὶ πάντα τὸν χοῦν ἐξοίσουσιν ἔξω τῆς πόλεως εἰς τόπον ἀκάθαρτον.
\vs{46}Καὶ ὁ εἰσπορευόμενος εἰς τὴν οἰκίαν πάσας τὰς ἡμέρας, ἃς ἀφωρισμένη ἐστὶν, ἀκάθαρτος ἔσται ἕως ἑσπέρας·
\vs{47}Καὶ ὁ κοιμώμενος ἐν τῇ οἰκίᾳ, πλυνεῖ τὰ ἱμάτια αὐτοῦ, καὶ ἀκάθαρτος ἔσται ἕως ἑσπέρας· καὶ ὁ ἔσθων ἐν τῇ οἰκίᾳ, πλυνεῖ τὰ ἱμάτια αὐτοῦ, καὶ ἀκάθαρτος ἔσται ἕως ἑσπέρας.

\vs{48}Ἐὰν δὲ παραγενόμενος εἰσέλθῃ ὁ ἱερεὺς καὶ ἴδῃ, καὶ ἰδοὺ οὐ διαχύσει οὐ διαχεῖται ἡ ἁφὴ ἐν τῇ οἰκίᾳ μετὰ τὸ ἐξαλειφθῆναι τὴν οἰκίαν, καὶ καθαριεῖ ὁ ἱερεὺς τὴν οἰκίαν, ὅτι ἰάθη ἡ ἁφή.
\vs{49}Καὶ λήψεται ἀφαγνίσαι τὴν οἰκίαν, δύο ὀρνίθια ζῶντα καθαρὰ, καὶ ξύλον κέδρινον, καὶ κεκλωσμένον κόκκινον, καὶ ὕσσωπον.
\vs{50}Καὶ σφάξει τὸ ὀρνίθιον τὸ ἓν εἰς σκεῦος ὀστράκινον ἐφʼ ὕδατι ζῶντι·
\vs{51}Καὶ λήψεται τὸ ξύλον τὸ κέδρινον, καὶ τὸ κεκλωσμένον κόκκινον, καὶ τὸν ὕσσωπον, καὶ τὸ ὀρνίθιον τὸ ζῶν· καὶ βάψει αὐτὸ εἰς τὸ αἷμα τοῦ ὀρνιθίου τοῦ ἐσφαγμενου ἐφʼ ὕδατι ζῶντι· καὶ περιῤῥανεῖ ἐν αὐτοῖς ἐπὶ τὴν οἰκίαν ἑπτάκις.
\vs{52}Καὶ ἀφαγνιεῖ τὴν οἰκίαν ἐν τῷ αἵματι τοῦ ὀρνιθίου, καὶ ἐν τῷ ὕδατι τῷ ζῶντι, καὶ ἐν τῷ ὀρνιθίῳ τῷ ζῶντι, καὶ ἐν τῷ ξύλῳ τῷ κεδρίνῳ, καὶ ἐν τῷ ὑσσώπῳ, καὶ ἐν τῷ κεκλωσμένῳ κοκκίνῳ.
\vs{53}Καὶ ἐξαποστελεῖ τὸ ὀρνίθιον τὸ ζῶν ἔξω τῆς πόλεως εἰς τὸ πεδίον, καὶ ἐξιλάσεται περὶ τῆς οἰκίας, καὶ καθαρὰ ἔσται.
\vs{54}Οὗτος ὁ νόμος κατὰ πᾶσαν ἁφὴν λέπρας, καὶ θραύσματος,
\vs{55}καὶ τῆς λέπρας ἱματίου, καὶ οἰκίας,
\vs{56}καὶ οὐλῆς, καὶ σημασίας, καὶ τοῦ αὐγάζοντος,
\vs{57}καὶ τοῦ ἐξηγήσασθαι ᾗ ἡμέρᾳ ἀκάθαρτον, καὶ ᾗ ἡμέρᾳ καθαρισθήσεται· οὗτος ὁ νόμος τῆς λέπρας.

\ch{15}
Καὶ ἐλάλησε Κύριος πρὸς Μωυσῆν καὶ Ἀαρὼν, λέγων,
\vs{2}λάλησον τοῖς υἱοῖς Ἰσραὴλ, καὶ ἐρεῖς αὐτοῖς, ἀνδρὶ ἀνδρὶ ᾧ ἐὰν γένηται ῥύσις ἐκ τοῦ σώματος αὐτοῦ, ἡ ῥύσις αὐτοῦ ἀκάθαρτός ἐστι.
\vs{3}Καὶ οὗτος ὁ νόμος τῆς ἀκαθαρσίας αὐτοῦ· ῥέων γόνον ἐκ σώματος αὐτοῦ, ἐκ τῆς ῥύσεως, ἧς συνέστηκε τὸ σῶμα αὐτοῦ διὰ τῆς ῥύσεως, αὕτη ἡ ἀκαθαρσία αὐτοῦ ἐν αὐτῷ· πᾶσαι αἱ ἡμέραι ῥύσεως σώματος αὐτοῦ, ᾗ συνέστηκε τὸ σῶμα αὐτοῦ διὰ τῆς ῥύσεως, ἀκαθαρσία αὐτοῦ ἐστι.
\vs{4}Πᾶσα κοίτη ἐφʼ ἧς ἂν κοιμηθῇ ἐπʼ αὐτῆς ὁ γονοῤῥυὴς, ἀκάθαρτός ἐστι, καὶ πᾶν σκεῦος ἐφʼ ὃ ἂν καθίσῃ ἐπʼ αὐτὸ ὁ γονοῤῤυὴς, ἀκάθαρτον ἔσται.
\vs{5}Καὶ ἄνθρωπος, ὃς ἐὰν ἅψηται τῆς κοίτης αὐτοῦ, πλυνεῖ τὰ ἱμάτια αὐτοῦ, καὶ λούσεται ὕδατι, καὶ ἀκάθαρτος ἔσται ἕως ἑσπέρας.
\vs{6}Καὶ ὁ καθήμενος ἐπὶ τοῦ σκεύους ἐφʼ ὃ ἂν καθίσῃ ὁ γονοῤῥυὴς, πλυνεῖ τὰ ἱμάτια αὐτοῦ, καὶ λούσεται ὕδατι, καὶ ἀκάθαρτος ἔσται ἕως ἑσπέρας.
\vs{7}Καὶ ὁ ἁπτόμενος τοῦ χρωτὸς τοῦ γονοῤῥυοῦς, πλυνεῖ τὰ ἱμάτια, καὶ λούσεται ὕδατι, καὶ ἀκάθαρτος ἔσται ἕως ἑσπέρας.
\vs{8}Ἐὰν δὲ προσσιελίσῃ ὁ γονοῤῥυὴς ἐπὶ τὸν καθαρὸν, πλυνεῖ τὰ ἱμάτια αὐτοῦ, καὶ λούσεται ὕδατι, καὶ ἀκάθαρτος ἔσται ἕως ἑσπέρας.
\vs{9}Καὶ πᾶν ἐπίσαγμα ὄνου, ἐφʼ ὃ ἂν ἐπιβῇ ἐπʼ αὐτὸ ὁ γονοῤῥυὴς, ἀκάθαρτον ἔσται ἕως ἑσπέρας.
\vs{10}Καὶ πᾶς ὁ ἁπτόμενος ὅσα ἂν ᾖ ὑποκάτω αὐτοῦ, ἀκάθαρτος ἔσται ἕως ἑσπέρας· καὶ ὁ αἴρων αὐτὰ, πλυνεῖ τὰ ἱμάτια αὐτοῦ, καὶ λούσεται ὕδατι, καὶ ἀκάθαρτος ἔσται ἕως ἑσπέρας.
\vs{11}Καὶ ὅσων ἐὰν ἅψηται ὁ γονοῤῥυὴς, καὶ τὰς χεῖρας οὐ νένιπται ὕδατι, πλυνεῖ τὰ ἱμάτια, καὶ λούσεται τὸ σῶμα ὕδατι, καὶ ἀκάθαρτος ἔσται ἕως ἑσπέρας.
\vs{12}Καὶ σκεῦος ὀστράκινον οὗ ἂν ἅψηται ὁ γονοῤῥυὴς, συντριβήσεται· καὶ σκεῦος ξύλινον νιφήσεται ὕδατι, καὶ καθαρὸν ἔσται.
\vs{13}Ἐὰν δὲ καθαρισθῇ ὁ γονοῤῥυὴς ἐκ τῆς ῥύσεως αὐτοῦ, καὶ ἐξαριθμηθήσεται αὐτῷ ἑπτὰ ἡμέρας εἰς τὸν καθαρισμὸν αὐτοῦ, καὶ πλυνεῖ τὰ ἱμάτια αὐτοῦ, καὶ λούσεται τὸ σῶμα ὕδατι, καὶ καθαρὸς ἔσται.
\vs{14}Καὶ τῇ ἡμέρᾳ τῇ ὀγδόῃ λήψεται ἑαυτῷ δύο τρυγόνας, ἢ δύο νοσσοὺς περιστερῶν, καὶ οἴσει αὐτὰ ἔναντι Κυρίου ἐπὶ τὰς θύρας τῆς σκηνῆς τοῦ μαρτυρίου, καὶ δώσει αὐτὰ τῷ ἱερεῖ.
\vs{15}Καὶ ποιήσει αὐτὰ ὁ ἱερεὺς μίαν περὶ ἁμαρτίας, καὶ μίαν εἰς ὁλοκαύτωμα· καὶ ἐξιλάσεται περὶ αὐτοῦ ὁ ἱερεὺς ἔναντι Κυρίου ἀπὸ τῆς ῥύσεως αὐτοῦ.

\vs{16}Καὶ ἄνθρωπος ᾧ ἂν ἐξέλθῃ ἐξ αὐτοῦ κοίτη σπέρματος, καὶ λούσεται ὕδατι πᾶν τὸ σῶμα αὐτοῦ, καὶ ἀκάθαρτος ἔσται ἕως ἑσπέρας.
\vs{17}Καὶ πᾶν ἱμάτιον, καὶ πᾶν δέρμα ἐφʼ ὃ ἂν ᾖ ἐπʼ αὐτὸ κοίτη σπέρματος, καὶ πλυθήσεται ὕδατι, καὶ ἀκάθαρτον ἔσται ἕως ἑσπέρας.
\vs{18}Καὶ γυνὴ ἐὰν κοιμηθῇ ἀνὴρ μετʼ αὐτῆς κοίτην σπέρματος, καὶ λούσονται ὕδατι, καὶ ἀκάθαρτοι ἔσονται ἕως ἑσπέρας.
\vs{19}Καὶ γυνὴ ἥτις ἂν ᾖ ῥέουσα αἵματι, καὶ ἔσται ἡ ῥύσις αὐτῆς ἐν τῷ σώματι αὐτῆς, ἑπτὰ ἡμέρας ἔσται ἐν τῇ ἀφέδρῳ αὐτῆς· πᾶς ὁ ἁπτόμενος αὐτῆς, ἀκάθαρτος ἔσται ἕως ἑσπέρας.
\vs{20}Καὶ πᾶν ἐφʼ ὃ ἂν κοιτάζηται ἐπʼ αὐτὸ ἐν τῇ ἀφέδρῳ αὐτῆς, ἀκάθαρτον ἔσται· καὶ πᾶν ἐφʼ ὃ ἂν ἐπικαθίσῃ ἐπʼ αὐτὸ, ἀκάθαρτον ἔσται.
\vs{21}Καὶ πᾶς ὃς ἂν ἅψηται τῆς κοίτης αὐτῆς, πλυνεῖ τὰ ἱμάτια αὐτοῦ, καὶ λούσεται τὸ σῶμα αὐτοῦ ὕδατι, καὶ ἀκάθαρτος ἔσται ἕως ἑσπέρας.
\vs{22}Καὶ πᾶς ὁ ἁπτόμενος παντὸς σκεύους οὗ ἐὰν καθίσῃ ἐπʼ αὐτὸ, πλυνεῖ τὰ ἱμάτια αὐτοῦ, καὶ λούσεται ὕδατι, καὶ ἀκάθαρτος ἔσται ἕως ἑσπέρας.
\vs{23}Ἐὰν δὲ ἐν τῇ κοίτῃ αὐτῆς οὔσης, ἢ ἐπὶ τοῦ σκεύους οὗ ἐὰν καθίσῃ ἐπʼ αὐτῷ ἐν τῷ ἅπτεσθαι αὐτὸν αὐτῆς, ἀκάθαρτος ἔσται ἕως ἑσπέρας.

\vs{24}Ἐὰν δὲ κοίτῃ κοιμηθῇ τις μετʼ αὐτῆς, καὶ γένηται ἡ ἀκαθαρσία αὐτῆς ἐπʼ αὐτῷ, ἀκάθαρτος ἔσται ἑπτὰ ἡμέρας· καὶ πᾶσα κοίτη ἐφʼ ᾗ ἂν κοιμηθῇ ἐπʼ αὐτῇ, ἀκάθαρτος ἔσται.
\vs{25}Καὶ γυνὴ ἐὰν ῥέῃ ῥύσει αἵματος ἡμέρας πλείους, οὐκ ἐν καιρῷ τῆς ἀφέδρου αὐτῆς, ἐὰν καὶ ῥέῃ μετὰ τὴν ἄφεδρον αὐτῆς, πᾶσαι αἱ ἡμέραι ῥύσεως ἀκαθαρσίας αὐτῆς, καθάπερ αἱ ἡμέραι τῆς ἀφέδρου αὐτῆς, ἔσται ἀκάθαρτος.
\vs{26}Καὶ πᾶσα κοίτη ἐφʼ ἧς ἂν κοιμηθῇ ἐπʼ αὐτῆς πάσας τὰς ἡμέρας τῆς ῥύσεως, κατὰ τὴν κοίτην τῆς ἀφέδρου, ἔσται αὐτῇ· καὶ πᾶν σκεῦος ἐφʼ ὃ ἂν καθίσῃ ἐπʼ αὐτὸ, ἀκάθαρτον ἔσται κατὰ τὴν ἀκαθαρσίαν τῆς ἀφέδρου.
\vs{27}Πᾶς ὁ ἁπτόμενος αὐτῆς ἀκάθαρτος ἔσται, καὶ πλυνεῖ τὰ ἱμάτια καὶ λούσεται τὸ σῶμα ὕδατι, καὶ ἀκάθαρτος ἔσται ἕως ἑσπέρας.
\vs{28}Ἐὰν δὲ καθαρισθῇ ἀπὸ τῆς ῥύσεως, καὶ ἐξαριθμήσεται αὐτῇ ἑπτὰ ἡμέρας, καὶ μετὰ ταῦτα καθαρισθήσεται.
\vs{29}Καὶ τῇ ἡμέρᾳ τῇ ὀγδόῃ λήψεται αὑτῇ δύο τρυγόνας, ἢ δύο νοσσοὺς περιστερῶν, καὶ οἴσει αὐτὰ πρὸς τὸν ἱερέα ἐπὶ τὴν θύραν τῆς σκηνῆς τοῦ μαρτυρίου.
\vs{30}Καὶ ποιήσει ὁ ἱερεὺς τὴν μίαν περὶ ἁμαρτίας, καὶ τὴν μίαν εἰς ὁλοκαύτωμα· καὶ ἐξιλάσεται περὶ αὐτῆς ὁ ἱερεὺς ἔναντι Κυρίου ἀπὸ ῥύσεως ἀκαθαρσίας αὐτῆς.

\vs{31}Καὶ εὐλαβεῖς ποιήσετε τοὺς υἱοὺς Ἰσραὴλ ἀπὸ τῶν ἀκαθαρσιῶν αὐτῶν· καὶ οὐκ ἀποθανοῦνται διὰ τὴν ἀκαθαρσίαν αὐτῶν, ἐν τῷ μιαίνειν αὐτοὺς τὴν σκηνήν μου τὴν ἐν αὐτοῖς.
\vs{32}Οὗτος ὁ νόμος τοῦ γονοῤῥυοῦς· καὶ ἐάν τινι ἐξέλθῃ ἐξ αὐτοῦ κοίτη σπέρματος, ὥστε μιανθῆναι ἐν αὐτῇ,
\vs{33}καὶ τῇ αἱμοῤῥοούσῃ ἐν τῇ ἀφέδρῳ αὐτῆς, καὶ ὁ γονοῤῥυὴς ἐν τῇ ῥύσει αὐτοῦ τῷ ἄρσενι ἢ τῇ θηλείᾳ, καὶ τῷ ἀνδρὶ, ὃς ἂν κοιμηθῇ μετὰ ἀποκαθημένης.

\ch{16}
Καὶ ἐλάλησε Κύριος πρὸς Μωυσῆν, μετὰ τὸ τελευτῆσαι τοὺς δύο υἱοὺς Ἀαρὼν ἐν τῷ προσάγειν αὐτοὺς πῦρ ἀλλότριον ἔναντι Κυρίου, καὶ ἐτελεύτησαν.
\vs{2}Καὶ εἶπε Κύριος πρὸς Μωυσῆν, λάλησον πρὸς Ἀαρὼν τὸν ἀδελφόν σου, καὶ μὴ εἰσπορευέσθω πᾶσαν ὥραν εἰς τὸ ἅγιον ἐσώτερον τοῦ καταπετάσματος εἰς πρόσωπον τοῦ ἱλαστηρίου, ὅ ἐστιν ἐπὶ τῆς κιβωτοῦ τοῦ μαρτυρίου, καὶ οὐκ ἀποθανεῖται· ἐν γὰρ νεφέλῃ ὀφθήσομαι ἐπὶ τοῦ ἱλαστηρίου.
\vs{3}Οὕτως εἰσελεύσεται Ἀαρὼν εἰς τὸ ἅγιον· ἐν μόσχῳ ἐκ βοῶν περὶ ἁμαρτίας, καὶ κριὸν εἰς ὁλοκαύτωμα.
\vs{4}Καὶ χιτῶνα λινοῦν ἡγιασμένον ἐνδύσεται, καὶ περισκελὲς λινοῦν ἔσται ἐπὶ τοῦ χρωτὸς αὐτοῦ, καὶ ζώνῃ λινῇ ζώσεται, καὶ κίδαριν λινῆν περιθήσεται, ἱμάτια ἅγιά ἐστι· καὶ λούσεται ὕδατι πᾶν τὸ σῶμα αὐτοῦ καὶ ἐνδύσεται αὐτά.
\vs{5}Καὶ παρὰ τῆς συναγωγῆς τῶν υἱῶν Ἰσραὴλ λήψεται δύο χιμάρους ἐξ αἰγῶν περὶ ἁμαρτίας, καὶ κριὸν ἕνα εἰς ὁλοκαύτωμα.
\vs{6}Καὶ προσάξει Ἀαρὼν τὸν μόσχον τὸν περὶ τῆς ἁμαρτίας αὐτοῦ, καὶ ἐξιλάσεται περὶ αὐτοῦ, καὶ τοῦ οἴκου αὐτοῦ.
\vs{7}Καὶ λήψεται τοὺς δύο χιμάρους, καὶ στήσει αὐτοὺς ἔναντι Κυρίου παρὰ τὴν θύραν τῆς σκηνῆς τοῦ μαρτυρίου.
\vs{8}Καὶ ἐπιθήσει Ἀαρὼν ἐπὶ τοὺς δύο χιμάρους κλῆρους· κλῆρον ἕνα τῷ Κυρίῳ, καὶ κλῆρον ἕνα τῷ ἀποπομπαίῳ.
\vs{9}Καὶ προσάξει Ἀαρὼν τὸν χίμαρον ἐφʼ ὃν ἐπῆλθεν ἐπʼ αὐτὸν ὁ κλῆρος τῷ Κυρίῳ, καὶ προσοίσει περὶ ἁμαρτίας.
\vs{10}Καὶ τὸν χίμαρον, ἐφʼ ὃν ἐπῆλθεν ἐπʼ αὐτὸν ὁ κλῆρος τοῦ ἀποπομπαίου, στήσει αὐτὸν ζῶντα ἔναντι Κυρίου, τοῦ ἐξιλάσασθαι ἐπʼ αὐτοῦ, ὥστε ἀποστεῖλαι αὐτὸν εἰς τὴν ἀποπομπήν, καὶ ἀφήσει αὐτὸν εἰς τὴν ἔρημον.
\vs{11}Καὶ προσάξει Ἀαρὼν τὸν μόσχον τὸν περὶ τῆς ἁμαρτίας αὑτοῦ, καὶ ἐξιλάσεται περὶ ἑαυτοῦ, καὶ τοῦ οἴκου· καὶ σφάξει τὸν μόσχον περὶ τῆς ἁμαρτίας αὐτοῦ.
\vs{12}Καὶ λήψεται τὸ πυρεῖον πλῆρες ἀνθράκων πυρὸς ἀπὸ τοῦ θυσιαστηρίου, τοῦ ἀπέναντι Κυρίου· καὶ πλήσει τὰς χεῖρας θυμιάματος συνθέσεως λεπτῆς, καὶ εἰσοίσει ἐσώτερον τοῦ καταπετάσματος.
\vs{13}Καὶ ἐπιθήσει τὸ θυμίαμα ἐπὶ τὸ πῦρ ἔναντι Κυρίου· καὶ καλύψει ἡ ἀτμὶς τοῦ θυμιάματος τὸ ἱλαστήριον τὸ ἐπὶ τῶν μαρτυρίων, καὶ οὐκ ἀποθανεῖται.
\vs{14}Καὶ λήψεται ἀπὸ τοῦ αἵματος τοῦ μόσχου, καὶ ῥανεῖ τῷ δακτύλῳ ἐπὶ τὸ ἱλαστήριον κατὰ ἀνατολάς· κατὰ πρόσωπον τοῦ ἱλαστηρίου ῥανεῖ ἑπτάκις ἀπὸ τοῦ αἵματος τῷ δακτύλῳ.

\vs{15}Καὶ σφάξει τὸν χίμαρον τὸν περὶ ἁμαρτίας, τὸν περὶ τοῦ λαοῦ, ἔναντι Κυρίου· καὶ εἰσοίσει τοῦ αἵματος αὐτοῦ ἐσώτερον τοῦ καταπετάσματος, καὶ ποιήσει τὸ αἷμα αὐτοῦ, ὃν τρόπον ἐποίησε τὸ αἷμα τοῦ μόσχου· καὶ ῥανεῖ τὸ αἷμα αὐτοῦ ἐπὶ τὸ ἱλαστήριον, κατὰ πρόσωπον τοῦ ἱλαστηρίου.
\vs{16}Καὶ ἐξιλάσεται τὸ ἅγιον ἀπὸ τῶν ἀκαθαρσιῶν τῶν υἱῶν Ἰσραὴλ, καὶ ἀπὸ τῶν ἀδικημάτων αὐτῶν περὶ πασῶν τῶν ἁμαρτιῶς αὐτῶν· καὶ οὕτω ποιήσει τῇ σκηνῇ τοῦ μαρτυρίου τῇ ἐκτισμένῃ ἐν αὐτοῖς ἐν μέσῳ τῆς ἀκαθαρσίας αὐτῶν.
\vs{17}Καὶ πᾶς ἄνθρώπος οὐκ ἔσται ἐν τῇ σκηνῇ τοῦ μαρτυρίου, εἰσπορευομένου αὐτοῦ ἐξιλάσασθαι ἐν τῷ ἁγίῳ, ἕως ἂν ἐξέλθῃ· καὶ ἐξιλάσεται περὶ ἑαυτοῦ, καὶ τοῦ οἴκου αὐτοῦ, καὶ περὶ πάσης συναγωγῆς υἱῶν Ἰσραήλ.
\vs{18}Καὶ ἐξελεύσεται ἐπὶ τὸ θυσιαστήριον τὸ ὂν ἀπέναντι Κυρίου, καὶ ἐξιλάσεται ἐπʼ αὐτοῦ· καὶ λήψεται ἀπὸ τοῦ αἵματος τοῦ μόσχου, καὶ ἀπὸ τοῦ αἵματος τοῦ χιμάρου, καὶ ἐπιθήσει ἐπὶ τὰ κέρατα τοῦ θυσιαστηρίου κύκλῳ.
\vs{19}Καὶ ῥανεῖ ἐπʼ αὐτὸ ἀπὸ τοῦ αἵματος τῷ δακτύλῳ ἑπτάκις, καὶ καθαριεῖ αὐτό, καὶ ἁγιάσει αὐτὸ ἀπὸ τῶν ἀκαθαρσιῶν τῶν υἱῶν Ἰσραήλ.
\vs{20}Καὶ συντελέσει ἐξιλασκόμενος τὸ ἅγιον, καὶ τὴν σκηνὴν τοῦ μαρτυρίου, καὶ τὸ θυσιαστήριον, καὶ περὶ τῶν ἱερέων καθαριεῖ· καὶ προσάξει τὸν χίμαρον τὸν ζῶντα.
\vs{21}Καὶ ἐπιθήσει Ἀαρὼν τὰς χεῖρας αὐτοῦ ἐπὶ τὴν κεφαλὴν τοῦ χιμάρου τοῦ ζῶντος, καὶ ἐξαγορεύσει ἐπʼ αὐτοῦ πάσας τὰς ἀνομίας τῶν υἱῶν Ἰσραὴλ, καὶ πάσας τὰς ἀδικίας αὐτῶν, καὶ πάσας τὰς ἁμαρτίας αὐτῶν· καὶ ἐπιθήσει αὐτὰς ἐπὶ τὴν κεφαλὴν τοῦ χιμάρου τοῦ ζῶντος· καὶ ἐξαποστελεῖ ἐν χειρὶ ἀνθρώπου ἑτοιμου εἰς τὴν ἔρημον.
\vs{22}Καὶ λήψεται ὁ χίμαρος ἐφʼ ἑαυτῷ τὰς ἀδικίας αὐτῶν εἰς γῆν ἄβατον· καὶ ἐξαποστελεῖ τὸν χίμαρον εἰς τὴν ἔρημον.
\vs{23}Καὶ εἰσελεύσεται Ἀαρὼν εἰς τὴν σκηνὴν τοῦ μαρτυρίου, καὶ ἐκδύσεται τὴν στολὴν τὴν λινῆν, ἣν ἐνδεδύκει, εἰσπορευομένου αὐτοῦ εἰς τὸ ἅγιον, καὶ ἀποθήσει αὐτὴν ἐκεῖ.
\vs{24}Καὶ λούσεται τὸ σῶμα αὐτοῦ ὕδατι ἐν τόπῳ ἁγίῳ, καὶ ἐνδύσεται τὴν στολὴν αὐτοῦ, καὶ ἐξελθὼν ποιήσει τὸ ὁλοκαύτωμα αὐτοῦ καὶ τὸ ὁλοκάρπωμα τοῦ λαοῦ, καὶ ἐξιλάσεται περὶ αὐτοῦ, καὶ περὶ τοῦ οἴκου αὐτοῦ, καὶ περὶ τοῦ λαοῦ, ὡς περὶ τῶν ἱερέων.
\vs{25}Καὶ τὸ στέαρ τὸ περὶ τῶν ἁμαρτιῶν ἀνοίσει ἐπὶ τὸ θυσιαστήριον.

\vs{26}Καὶ ὁ ἐξαποστέλλων τὸν χίμαρον τὸν διεσταλμένον εἰς ἄφεσιν, πλυνεῖ τὰ ἱμάτια, καὶ λούσεται τὸ σῶμα αὐτοῦ ὕδατι, καὶ μετὰ ταῦτα εἰσελεύσεται εἰς τὴν παρεμβολήν.
\vs{27}Καὶ τὸν μόσχον τὸν περὶ τῆς ἁμαρτίας, καὶ τὸν χίμαρον τὸν περὶ τῆς ἁμαρτίας, ὧν τὸ αἷμα εἰσηνέχθη ἐξιλάσασθαι ἐν τῷ ἁγίῳ, ἐξοίσουσιν αὐτὰ ἔξω τῆς παρεμβολῆς, ταὶ κατακαύσουσιν αὐτὰ ἐν πυρὶ, καὶ τὰ δέρματα αὐτῶν καὶ τὰ κρέα αὐτῶν καὶ τὴν κόπρον αὐτῶν.
\vs{28}Ὁ δὲ κατακαίων αὐτὰ, πλυνεῖ τὰ ἱμάτια, καὶ λούσεται τὸ σῶμα αὐτοῦ ὕδατι, καὶ μετὰ ταῦτα εἰσελεύσεται εἰς τὴν παρεμβολήν.

\vs{29}Καὶ ἔσται τοῦτο ὑμῖν νόμιμον αἰώνιον· ἐν τῷ μηνὶ τῷ ἑβδόμῳ, δεκάτῃ τοῦ μηνὸς, ταπεινώσετε τὰς ψυχὰς ὑμῶν, καὶ πᾶν ἔργον οὐ ποιήσετε ὁ αὐτόχθων, καὶ ὁ προσήλυτος ὁ προσκείμενος ἐν ὑμῖν.
\vs{30}Ἐν γὰρ τῇ ἡμέρᾳ ταύτῃ ἐξιλάσεται περὶ ὑμῶν, καθαρίσαι ὑμᾶς ἀπὸ πασῶν τῶν ἁμαρτιῶν ὑμῶν ἔναντι Κυρίου, καὶ καθαρισθήσεσθε.
\vs{31}Σάββατα σαββάτων ἀνάπαυσις αὕτη ἔσται ὑμῖν· καὶ ταπεινώσετε τὰς ψυχὰς ὑμῶν, νόμιμον αἰώνιον.
\vs{32}Ἐξιλάσεται ὁ ἱερεὺς, ὃν ἂν χρίσωσιν αὐτὸν, καὶ ὃν ἂν τελειώσωσι τὰς χεῖρας αὐτοῦ ἱερατεύειν μετὰ τὸν πατέρα αὐτοῦ· καὶ ἐνδύσεται τὴν στολὴν τὴν λινῆν, στολὴν ἁγίαν.
\vs{33}Καὶ ἐξιλάσεται τὸ ἅγιον τοῦ ἁγίου, καὶ τὴν σκηνὴν τοῦ μαρτυρίου, καὶ τὸ θυσιαστήριον ἐξιλάσεται, καὶ περὶ τῶν ἱερέων, καὶ περὶ πάσης συναγωγῆς ἐξιλάσεται.
\vs{34}Καὶ ἔσται τοῦτο ὑμῖν νόμιμον αἰώνιον ἐξιλάσκεσθαι περὶ τῶν νἱῶν Ἰσραὴλ ἀπὸ πασῶν τῶν ἁμαρτιῶν αὐτῶν· ἅπαξ τοῦ ἐνιαυτοῦ ποιηθήσεται, καθὰ συνέταξε Κύριος τῷ Μωυσῇ.

\ch{17}
Καὶ ἐλάλησε Κύριος πρὸς Μωυσῆν, λέγων,
\vs{2}λάλησον πρὸς Ἀαρὼν καὶ πρὸς τοὺς υἱοὺς αὐτοῦ, καὶ πρὸς πάντας υἱοὺς Ἰσραὴλ, καὶ ἐρεῖς πρὸς αὐτοὺς, τοῦτο τὸ ῥῆμα ὃ ἐνετείλατο Κύριος, λέγων,
\vs{3}ἄνθρωπος ἄνθρωπος τῶν υἱῶν Ἰσραὴλ, ἢ τῶν προσηλύτων τῶν προσκειμένων ἐν ὑμῖν, ὃς ἐὰν σφάξῃ μόσχον, ἢ πρόβατον, ἢ αἶγα ἐν τῇ παρεμβολῇ, καὶ ὃς ἂν σφάξῃ ἔξω τῆς παρεμβολῆς,
\vs{4}καὶ ἐπὶ τὴν θύραν τῆς σκηνῆς τοῦ μαρτυρίου μὴ ἐνέγκῃ, ὥστε ποιῆσαι αὐτὸ εἰς ὁλοκαύτωμα ἢ σωτήριον Κυρίῳ δεκτὸν εἰς ὀσμὴν εὐωδίας· καὶ ὃς ἂν σφάξῃ ἔξω, καὶ ἐπὶ τὴν θύραν τῆς σκηνῆς τοῦ μαρτυρίου μὴ ἐνέγκῃ αὐτὸ, ὥστε προσενέγκαι δῶρον τῷ Κυρίῳ ἀπέναντι τῆς σκηνῆς Κυρίου· καὶ λογισθήσεται τῷ ἀνθρώπῳ ἐκείνῳ αἷμα· αὶμα ἐξέχεεν· ἐξολοθρευθήσεται ἡ ψυχὴ ἐκείνη ἐκ τοῦ λαοῦ αὐτῆς.
\vs{5}Ὅπως ἀναφέρωσιν οἱ υἱοὶ Ἰσραὴλ τὰς θυσίας αὐτῶν, ὅσας ἂν αὐτοὶ σφάξουσιν ἐν τοῖς πεδίοις, καὶ οἴσουσι τῷ Κυρίῳ ἐπὶ τὰς θύρας τῆς σκηνῆς τοῦ μαρτυρίου πρὸς τὸν ἱερέα· καὶ θύσουσι θυσίαν σωτηρίου τῷ Κυρίῳ αὐτά.
\vs{6}Καὶ προσχεεῖ ὁ ἱερεὺς τὸ αἷμα ἐπὶ τὸ θυσιαστήριον κύκλῳ ἀπέναντι Κυρίου παρὰ τὰς θύρας τῆς σκηνῆς τοῦ μαρτυρίου· καὶ ἀνοίσει τὸ στέαρ εἰς ὀσμὴν εὐωδίας Κυρίῳ.

\vs{7}Καὶ οὐ θύσουσιν ἔτι τὰς θυσίας αὐτῶν τοῖς ματαίοις, οἷς αὐτοὶ ἐκπορνεύουσιν ὀπίσω αὐτῶν· νόμιμον αἰώνιον ἔσται ὑμῖν εἰς τὰς γενεὰς ὑμῶν.
\vs{8}Καὶ ἐρεῖς πρὸς αὐτοὺς, ἄνθρωπος ἄνθρωπος τῶν υἱῶν Ἰσραὴλ, ἢ ἀπὸ τῶν υἱῶν τῶν προσηλύτων τῶν προσκειμένων ἐν ὑμῖν, ὃς ἂν ποιήσῃ ὁλοκαύτωμα ἢ θυσίαν,
\vs{9}καὶ ἐπὶ τὴν θύραν τῆς σκηνῆς τοῦ μαρτυρίου μὴ ἐνέγκῃ ποιῆσαι αὐτὸ τῷ Κυρίῳ, ἐξολοθρευθήσεται ὁ ἄνθρωπος ἐκεῖνος ἐκ τοῦ λαοῦ αὐτοῦ.
\vs{10}Καὶ ἄνθρωπος ἄνθρωπος τῶν υἱῶν Ἰσραὴλ, ἢ τῶν προσηλύτων τῶν προσκειμένων ἐν ὑμῖν, ὃς ἂν φάγῃ πᾶν αἷμα· καὶ ἐπιστήσω τὸ πρόσωπόν μου ἐπὶ τὴν ψυχὴν τὴν ἔσθουσαν τὸ αἷμα, καὶ ἀπολῶ αὐτὴν ἐκ τοῦ λαοῦ αὐτῆς.
\vs{11}Ἡ γὰρ ψυχὴ πάσης σαρκὸς αἷμα αὐτοῦ ἐστι· καὶ ἐγὼ δέδωκα αὐτὸ ὑμῖν ἐπὶ τοῦ θυσιαστηρίου ἐξιλάσκεσθαι περὶ τῶν ψυχῶν ὑμῶν· τὸ γὰρ αἷμα αὐτοῦ ἀντὶ ψυχῆς ἐξιλάσεται.
\vs{12}Διὰ τοῦτο εἴρηκα τοῖς υἱοῖς Ἰσραὴλ, πᾶσα ψυχὴ ἐξ ὑμῶν οὐ φάγεται αἷμα· καὶ ὁ προσήλυτος ὁ προσκείμενος ἐν ὑμῖν οὐ φάγεται αἷμα.
\vs{13}Καὶ ἄνθρωπος ἄνθρωπος τῶν υἱῶν Ἰσραὴλ, ἢ τῶν προσηλύτων τῶν προσκειμένων ἐν ὑμῖν, ὃς ἂν θηρεύσῃ θήρευμα θηρίον ἢ πετεινὸν, ὃ ἔσθεται, καὶ ἐκχεεῖ τὸ αἷμα, καὶ καλύψει αὐτὸ τῇ γῇ.
\vs{14}Ἡ γὰρ ψυχὴ πάσης σαρκὸς αἷμα αὐτοῦ ἐστι· καὶ εἶπα τοῖς υἱοῖς Ἰσραὴλ, αἷμα πάσης σαρκὸς οὐ φάγεσθε, ὅτι ἡ ψυχὴ πάσης σαρκὸς αἷμα αὐτοῦ ἐστί· πᾶς ὁ ἔσθων αὐτὸ, ἐξολοθρευθήσεται.
\vs{15}Καὶ πᾶσα ψυχὴ, ἥτις φάγεται θνησιμαῖον, ἢ θηριάλωτον ἐν τοῖς αὐτόχθοσιν, ἢ ἐν τοῖς προσηλύτοις, πλυνεῖ τὰ ἱμάτια αὐτοῦ, καὶ λούσεται ὕδατι, καὶ ἀκάθαρτος ἔσται ἕως ἑσπέρας, καὶ καθαρὸς ἔσται.
\vs{16}Ἐὰν δὲ μὴ πλύνῃ τὰ ἱμάτια, καὶ τὸ σῶμα μὴ λούσηται ὕδατι, καὶ λήψεται ἀνόμημα αὐτοῦ.

\ch{18}
Καὶ εἶπε Κύριος πρὸς Μωυσῆν, λέγων,
\vs{2}λάλησον τοῖς υἱοῖς Ἰσραὴλ, καὶ ἐρεῖς πρὸς αὐτοὺς, ἐγὼ Κύριος ὁ Θεὸς ὑμῶν.
\vs{3}Κατὰ τὰ ἐπιτηδεύματα Αἰγύπτου, ἐν ᾗ κατῳκήσατε ἐπʼ αὐτῇ, οὐ ποιήσετε· καὶ κατὰ τὰ ἐπιτηδεύματα γῆς Χαναὰν, εἰς ἣν ἐγὼ εἰσάγω ὑμᾶς ἐκεῖ, οὐ ποιήσετε, καὶ τοῖς νομίμοις αὐτῶν οὐ πορεύσεσθε.
\vs{4}Τὰ κρίματά μου ποιήσετε, καὶ τὰ προστάγματά μου φυλάξεσθε, καὶ πορεύεσθε ἐν αὐτοῖς· ἐγὼ Κύριος ὁ Θεὸς ὑμῶν.
\vs{5}Καὶ φυλάξεσθε πάντα τὰ προστάγματά μου, καὶ πάντα τὰ κρίματά μου, καὶ ποιήσετε αὐτά· ἃ ποιήσας αὐτὰ ἄνθρωπος, ζήσεται ἐν αὐτοῖς· ἐγὼ Κύριος ὁ Θεὸς ὑμῶν.
\vs{6}Ἄνθρωπος ἄνθρωπος πρὸς πάντα οἰκεῖα σαρκὸς αὐτοῦ οὐ προσελεύσεται ἀποκαλύψαι ἀσχημοσύνην· ἐγὼ Κύριος.
\vs{7}Ἀσχημοσύνην πατρός σου καὶ ἀσχημοσύνην μητρός σου οὐκ ἀποκαλύψεις, μήτηρ γάρ σου ἐστὶν, οὐκ ἀποκαλύψεις τὴν ἀσχημοσύνην αὐτῆς.
\vs{8}Ἀσχημοσύνην γυναικὸς πατρός σου οὐκ ἀποκαλύψεις, ἀσχημοσύνη πατρός σου ἐστίν.
\vs{9}Ἀσχημοσύνην τῆς ἀδελφῆς σου ἐκ πατρός σου ἢ ἐκ μητρός σου, ἐνδογενοῦς ἢ γεγεννημένης ἔξω, οὐκ ἀποκαλύψεις ἀσχημοσύνην αὐτῶν.
\vs{10}Ἀσχημοσύνην θυγατρὸς υἱοῦ σου, ἢ θυγατρὸς θυγατρός σου, οὐκ ἀποκαλύψεις τὴν ἀσχημοσύνην αὐτῶν, ὅτι σὴ ἀσχημοσύνη ἐστίν.
\vs{11}Ἀσχημοσύνην θυγατρὸς γυναικὸς πατρός σου οὐκ ἀποκαλύψεις, ὁμοπατρία ἀδελφή σου ἐστὶν, οὐκ ἀποκαλύψεις τὴν ἀσχημοσύνην αὐτῆς.
\vs{12}Ἀσχημοσύνην ἀδελφῆς πατρός σου οὐκ ἀποκαλύψεις, οἰκεία γὰρ πατρός σου ἐστιν.
\vs{13}Ἀσχημοσύνην ἀδελφῆς μητρός σου οὐκ ἀποκαλύψεις, οἰκεία γὰρ μητρός σου ἐστίν.
\vs{14}Ἀσχημοσύνην ἀδελφοῦ τοῦ πατρός σου οὐκ ἀποκαλύψεις, καὶ πρὸς τὴν γυναῖκα αὐτοῦ οὐκ εἰσελεύσῃ, συγγενὴς γάρ σου ἐστίν.
\vs{15}Ἀσχημοσύνην νύμφης σου οὐκ ἀποκαλύψεις, γυνὴ γὰρ υἱοῦ σου ἐστὶν, οὐκ ἀποκαλύψεις τὴν ἀσχημοσύνην αὐτῆς.
\vs{16}Ἀσχημοσύνην γυναικὸς ἀδελφοῦ σου οὐκ ἀποκαλύψεις, ἀσχημοσύνη ἀδελφοῦ σου ἐστίν.
\vs{17}Ἀσχημοσύνην γυναικὸς καὶ θυγατρὸς αὐτῆς οὐκ ἀποκαλύψεις· τὴν θυγατέρα τοῦ υἱοῦ αὐτῆς, καὶ τὴν θυγατέρα τῆς θυγατρὸς αὐτῆς οὐ λήψῃ ἀποκαλύψαι τὴν ἀσχημοσύνην αὐτῶν, οἰκεῖαι γάρ σου εἰσίν· ἀσέβημα ἐστι.
\vs{18}Γυναῖκα ἐπʼ ἀδελφῇ αὐτῆς οὐ λήψῃ ἀντίζηλον ἀποκαλύψαι τὴν ἀσχημοσύνην αὐτῆς ἐπʼ αὐτῇ, ἔτι ζώσης αὐτῆς.

\vs{19}Καὶ πρὸς γυναῖκα ἐν χωρισμῷ ἀκαθαρσίας αὐτῆς οὐκ εἰσελεύσῃ ἀποκαλύψαι τὴν ἀσχημοσύνην αὐτῆς.
\vs{20}Καὶ πρὸς τὴν γυναῖκα τοῦ πλησίον σου οὐ δώσεις κοίτην σπέρματός σου, ἐκμιανθῆναι πρὸς αὐτήν.
\vs{21}Καὶ ἀπὸ τοῦ σπέρματός σου οὐ δώσεις λατρεύειν ἄρχοντι· καὶ οὐ βεβηλώσεις τὸ ὄνομα τὸ ἅγιον· ἐγὼ Κύριος.
\vs{22}Καὶ μετὰ ἄρσενος οὐ κοιμηθήσῃ κοίτην γυναικείαν, βδέλυγμα γάρ ἐστι.
\vs{23}Καὶ πρὸς πᾶν τετράπουν οὐ δώσεις τὴν κοίτην σου εἰς σπερματισμὸν, ἐκμιανθῆναι πρὸς αὐτό· καὶ γυνὴ οὐ στήσεται πρὸς πᾶν τετράπουν βιβασθῆναι· μυσαρὸν γάρ ἐστι.
\vs{24}Μὴ μιαίνεσθε ἐν πᾶσι τούτοις· ἐν πᾶσι γὰρ τούτοις ἐμίανθησαν τὰ ἔθνη, ἃ ἐγὼ ἐξαποστέλλω πρὸ προσώπου ὑμῶν,
\vs{25}καὶ ἐξεμιάνθη ἡ γῆ καὶ ἀνταπέδωκα ἀδικίαν αὐτοῖς διʼ αὐτὴν, καὶ προσώχθισεν ἡ γῆ τοῖς ἐγκαθημένοις ἐπʼ αὐτῆς.
\vs{26}Καὶ φυλάξεσθε πάντα τὰ νόμιμά μου, καὶ πάντα τὰ προστάγματά μου, καὶ οὐ ποιήσετε ἀπὸ πάντων τῶν βδελυγμάτων τούτων ὁ ἐγχώριος, καὶ ὁ προσγενόμενος προσήλυτος ἐν ὑμῖν·
\vs{27}(Πάντα γὰρ τὰ βδελύγματα ταῦτα ἐποίησαν οἱ ἄνθρωποι τῆς γῆς, οἱ ὄντες πρότερον ὑμῶν, καὶ ἐμιάνθη ἡ γῆ·)
\vs{28}καὶ ἵνα μὴ προσοχθίσῃ ὑμῖν ἡ γῆ ἐν τῷ μιαίνειν ὑμᾶς αὐτὴν, ὃν τρόπον προσώχθισε τοῖς ἔθνεσι τοῖς πρὸ ὑμῶν.
\vs{29}Ὅτι πᾶς ὃς ἐὰν ποιήσῃ ἀπὸ πάντων τῶν βδελυγμάτων τούτων, ἐξολοθρευθήσονται αἱ ψυχαὶ αἱ ποιοῦσαι ἐκ τοῦ λαοῦ αὐτῶν.
\vs{30}Καὶ φυλάξετε τὰ προστάγματά μου, ὅπως μὴ ποιήσητε ἀπὸ πάντων τῶν νομίμων τῶν ἐβδελυγμένων, ἃ γέγονε πρὸ τοῦ ὑμᾶς· καὶ οὐ μιανθήσεσθε ἐν αὐτοῖς, ὅτι ἐγὼ Κύριος ὁ Θεὸς ὑμῶν.

\ch{19}
Καὶ ἐλάλησε Κύριος πρὸς Μωυσῆν, λέγων,
\vs{2}λάλησον τῇ συναγωγῇ τῶν υἱῶν Ἰσραὴλ, καὶ ἐρεῖς πρὸς αὐτοῦς, ἅγιοι ἔσεσθε, ὅτι ἅγιος ἐγὼ Κύριος ὁ Θεὸς ὑμῶν.
\vs{3}Ἕκαστος πατέρα αὐτοῦ καὶ μητέρα αὐτοῦ φοβείσθω, καὶ τὰ σάββατά μου φυλάξεσθε· ἐγὼ Κύριος ὁ Θεὸς ὑμῶν.
\vs{4}Οὐκ ἐπακολουθήσετε εἰδώλοις, καὶ θεοὺς χωνευτοὺς οὐ ποιήσετε ὑμῖν· ἐγὼ Κύριος ὁ Θεὸς ὑμῶν.
\vs{5}Καὶ ἐὰν θύσητε θυσίαν σωτηρίου τῷ Κυρίῳ, δεκτὴν ὑμῶν θύσετε.
\vs{6}ᾟ ἂν ἡμέρᾳ θύσετε, βρωθήσεται, καὶ τῇ αὔριον· καὶ ἐὰν καταλειφθῇ ἕως ἡμέρας τρίτης, ἐν πυρὶ κατακαυθήσεται.
\vs{7}Ἐὰν δὲ βρώσει βρωθῇ τῇ ἡμέρᾳ τῇ τρίτῃ, ἄθυτόν ἐστιν, οὐ δεχθήσεται.
\vs{8}Ὁ δὲ ἔσθων αὐτὸ, ἁμαρτίαν λήψεται, ὅτι τὰ ἅγια Κυρίου ἐβεβήλωσε· καὶ ἐξολοθρευθήσονται αἱ ψυχαὶ αἱ ἔσθουσαι ἐκ τοῦ λαοῦ αὐτῶν.

\vs{9}Καὶ ἐκθεριζόντων ὑμῶν τὸν θερισμὸν τῆς γῆς ὑμῶν, οὐ συντελέσετε τὸν θερισμὸν ὑμῶν τοῦ ἀγροῦ σου ἐκθερίσαι· καὶ τὰ ἀποπίπτοντα τοῦ θερισμοῦ σου οὐ συλλέξεις,
\vs{10}καὶ τὸν ἀμπελῶνά σου οὐκ ἐπανατρυγήσεις, οὐδὲ τὰς ῥῶγας τοῦ ἀμπελῶνός σου συλλέξεις· τῷ πτωχῷ καὶ τῷ προσηλύτῳ καταλείψεις αὐτά. ἐγώ εἰμι Κύριος ὁ Θεὸς ὑμῶν.
\vs{11}Οὐ κλέψετε, οὐ ψεύσεσθε, οὐδὲ συκοφαντήσει ἕκαστος τὸν πλησίον.
\vs{12}Καὶ οὐκ ὀμεῖσθε τῷ ὀνόματί μου ἐπʼ ἀδίκῳ, καὶ οὐ βεβηλώσετε τὸ ὄνομα τὸ ἅγιον τοῦ Θεοῦ ὑμῶν· ἐγώ εἰμι Κύριος ὁ Θεὸς ὑμῶν.
\vs{13}Οὐκ ἀδικήσεις τὸν πλησίον, καὶ οὐχ ἁρπᾷ· καὶ οὐ μὴ κοιμηθήσεται ὁ μισθὸς τοῦ μισθωτοῦ σου παρὰ σοὶ ἕως πρωΐ.

\vs{14}Οὐ κακῶς ἐρεῖς κωφόν, καὶ ἀπέναντι τυφλοῦ οὐ προσθήσεις σκάνδαλον· καὶ φοβηθήσῃ Κύριον τὸν Θεόν σου· ἐγώ εἰμι Κύριος ὁ Θεὸς ὑμῶν.
\vs{15}Οὐ ποιήσετε ἄδικον ἐν κρίσει· οὐ λήψῃ πρόσωπον πτωχοῦ, οὐδὲ μὴ θαυμάσῃς πρόσωπον δυνάστου· ἐν δικαιοσύνῃ κρίνεις τὸν πλησίον σου.
\vs{16}Οὐ πορεύσῃ δόλῳ ἐν τῷ ἔθνει σου· οὐκ ἐπιστήσῃ ἐφʼ αἷμα τοῦ πλησίον σου· ἐγώ εἰμι Κύριος ὁ Θεὸς ὑμῶν.
\vs{17}Οὐ μισήσεις τὸν ἀδελφόν σου τῇ διανοίᾳ σου· ἐλεγμῷ ἐλέγξεις τὸν πλησίον σου, καὶ οὐ λήψῃ διʼ αὐτὸν ἁμαρτίαν.
\vs{18}Καὶ οὐκ ἐκδικᾶταί σου ἡ χείρ· καὶ οὐ μηνιεῖς τοῖς υἱοῖς τοῦ λαοῦ σου· καὶ ἀγαπήσεις τὸν πλησίον σου ὡς σεαυτόν· ἐγώ εἰμι Κύριος.

\vs{19}Τὸν νόμον μου φυλάξεσθε· τὰ κτήνη σου οὐ κατοχεύσεις ἑτεροζύγῳ· καὶ τὸν ἀμπελῶνά σου οὐ κατασπερεῖς διάφορον· καὶ ἱμάτιον ἐκ δύο ὑφασμένον κίβδηλον οὐκ ἐπιβαλεῖς σεαυτῷ.
\vs{20}Καὶ ἐάν τις κοιμηθῇ μετὰ γυναικὸς κοίτην σπέρματος, καὶ αὕτη ᾖ οἰκέτις διαπεφυλαγμένη ἀνθρώπῳ, καὶ αὕτη λύτροις οὐ λελύτρωται, ἢ ἐλευθερία οὐκ ἐδόθη αὐτῇ, ἐπισκοπὴ ἔσται αὐτοῖς· οὐκ ἀποθανοῦνται, ὅτι οὐκ ἀπηλευθερώθη.
\vs{21}Καὶ προσάξει τῆς πλημμελείας αὐτοῦ τῷ Κυρίῳ παρὰ τὴν θύραν τῆς σκηνῆς τοῦ μαρτυρίου κριὸν πλημμελείας.
\vs{22}Καὶ ἐξιλάσεται περὶ αὐτοῦ ὁ ἱερεὺς ἐν τῷ κριῷ τῆς πλημμελείας ἔναντι Κυρίου περὶ τῆς ἁμαρτίας ἧς ἥμαρτε, καὶ ἀφεθήσεται αὐτῷ ἡ ἁμαρτία ἣν ἥμαρτεν.
\vs{23}Ὅταν δὲ εἰσέλθητε εἰς τὴν γῆν, ἣν Κύριος ὁ Θεὸς ὑμῶν δίδωσιν ὑμῖν, καὶ καταφυτεύσετε πᾶν ξύλον βρώσιμον, καὶ περικαθαριεῖτε τὴν ἀκαθαρσίαν αὐτοῦ· ὁ καρπὸς αὐτοῦ τρία ἔτη ἔσται ὑμῖν ἀπερικάθαρτος, οὐ βρωθήσεται.
\vs{24}Καὶ τῷ ἔτει τῷ τετάρτῳ ἔσται πᾶς ὁ καρπὸς αὐτοῦ ἅγιος αἰνετὸς τῷ Κυρίῳ.
\vs{25}Ἐν δὲ τῷ ἔτει τῷ πέμπτῳ φάγεσθε τὸν καρπόν, πρόσθεμα ὑμῖν τὰ γεννήματα αὐτοῦ· ἐγώ εἰμι Κύριος ὁ Θεὸς ὑμῶν.

\vs{26}Μὴ ἔσθετε ἐπὶ τῶν ὀρέων, καὶ οὐκ οἰωνιεῖσθε, οὐδὲ ὀρνιθοσκοπήσεσθε.
\vs{27}Οὐ ποιήσετε σισόην ἐκ τῆς κόμης τῆς κεφαλῆς ὑμῶν, οὐδὲ φθερεῖτε τὴν ὄψιν τοῦ πώγωνος ὑμῶν.
\vs{28}Καὶ ἐντομίδας οὐ ποιήσετε ἐπὶ ψυχῇ ἐν τῷ σώματι ὑμῶν· καὶ γράμματα στικτὰ οὐ ποιήσετε ἐν ὑμῖν· ἐγώ εἰμι Κύριος ὁ Θεὸς ὑμῶν.
\vs{29}Οὐ βεβηλώσεις τὴν θυγατέρα σου ἐκπορνεῦσαι αὐτήν· καὶ οὐκ ἐκπορνεύσει ἡ γῆ, καὶ ἡ γῆ πλησθήσεται ἀνομίας.
\vs{30}Τὰ σάββατά μον φυλάξεσθε, καὶ ἀπὸ τῶν ἁγίων μου φοβηθήσεσθε· ἐγώ εἰμι Κύριος.
\vs{31}Οὐκ ἐπακολουθήσετε ἐγγαστριμύθοις, καὶ τοῖς ἐπαοιδοῖς οὐ προσκολληθήσεσθε, ἐκμιανθῆναι ἐν αὐτοῖς· ἐγώ εἰμι Κύριος ὁ Θεὸς ὑμῶν.
\vs{32}Ἀπὸ προσώπου πολιοῦ ἐξαναστήσῃ, καὶ τιμήσεις πρόσωπον πρεσβυτέρου, καὶ φοβηθήσῃ τὸν Θεόν σου· ἐγώ εἰμι Κύριος ὁ Θεὸς ὑμῶν.
\vs{33}Ἐὰν δέ τις προσέλθῃ ὑμῖν προσήλυτος ἐν τῇ γῇ ὑμῶν, οὐ θλίψετε αὐτόν.
\vs{34}Ὡς ὁ αὐτόχθων ἐν ὑμῖν ἔσται ὁ προσήλυτος ὁ προσπορευόμενος πρὸς ὑμᾶς, καὶ ἀγαπήσεις αὐτὸν ὡς σεαυτόν· ὅτι προσήλυτοι ἐγενήθητε ἐν γῇ Αἰγύπτῳ· ἐγώ εἰμι Κύριος ὁ Θεὸς ὑμῶν.
\vs{35}Οὐ ποιήσετε ἄδικον ἐν κρίσει, ἐν μέτροις καὶ ἐν σταθμίοις καὶ ἐν ζυγοῖς.
\vs{36}Ζυγὰ δίκαια καὶ σταθμία δίκαια καὶ χοῦς δίκαιος ἔσται ἐν ὑμιν· ἐγώ εἰμι Κύριος ὁ Θεὸς ὑμῶν, ὁ ἐξαγαγὼν ὑμᾶς ἐκ γῆς Αἰγύπτου.
\vs{37}Καὶ φυλάξεσθε πάντα τὸν νόμον μου, καὶ πάντα τὰ προστάγματά μου, καὶ ποιήσετε αὐτά· ἐγώ εἰμι Κύριος ὁ Θεὸς ὑμῶν.

\ch{20}
Καὶ ἐλάλησε Κύριος πρὸς Μωυσῆν, λέγων, καὶ τοῖς υἱοῖς Ἰσραὴλ λαλήσεις,
\vs{2}ἐάν τις ἀπὸ τῶν υἱῶν Ἰσραὴλ, ἢ ἀπὸ τῶν γεγενημένων προσηλύτων ἐν Ἰσραὴλ, ὃς ἂν δῷ τοῦ σπέρματος αὐτοῦ ἄρχοντι, θανάτῳ θανατούσθω· τὸ ἔθνος τὸ ἐπὶ τῆς γῆς λιθοβολήσουσιν αὐτὸν ἐν λίθοις.
\vs{3}Καὶ ἐγὼ ἐπιστήσω τὸ πρόσωπόν μου ἐπὶ τὸν ἄνθρωπον ἐκεῖνον, καὶ ἀπολῶ αὐτὸν ἐκ τοῦ λαοῦ αὐτοῦ, ὅτι τοῦ σπέρματος αὐτοῦ ἔδωκεν ἄρχοντι, ἵνα μιάνῃ τὰ ἅγιά μου, καὶ βεβηλώσῃ τὸ ὄνομα τῶν ἡγιασμένων μοι.
\vs{4}Ἐὰν δὲ ὑπερόψει ὑπερίδωσιν οἱ αὐτόχθονες τῆς γῆς τοῖς ὀφθαλμοῖς αὐτῶν ἀπὸ τοῦ ἀνθρώπου ἐκείνου, ἐν τῷ δοῦναι αὐτὸν τοῦ σπέρματος αὐτοῦ ἄρχοντι, τοῦ μὴ ἀποκτεῖναι αὐτόν·
\vs{5}καὶ ἐπιστήσω τὸ πρόσωπόν μου ἐπὶ τὸν ἄνθρωπον ἐκεῖνον, καὶ τὴν συγγένειαν αὐτοῦ, καὶ ἀπολῶ αὐτὸν, καὶ πάντας τοὺς ὁμονοοῦντας αὐτῷ, ὥστε ἐκπορνεύειν αὐτὸν εἰς τοὺς ἄρχοντας, ἐκ τοῦ λαοῦ αὐτῶν.

\vs{6}Καὶ ψυχὴ ἣ ἂν ἐπακολουθήσῃ ἐγγαστριμύθοις ἢ ἐπαοιδοῖς, ὥστε ἐκπορνεῦσαι ὀπίσω αὐτῶν, ἐπιστήσω τὸ πρόσωπόν μου ἐπὶ τὴν ψυχὴν ἐκείνην, καὶ ἀπολῶ αὐτῆν ἐκ τοῦ λαοῦ αὐτῆς.
\vs{7}Καὶ ἔσεσθε ἅγιοι, ὅτι ἅγιος ἐγὼ Κύριος ὁ Θεὸς ὑμῶν.
\vs{8}Καὶ φυλάξεσθε τὰ προστάγματά μου, καὶ ποιήσετε αὐτά· ἐγὼ Κύριος ὁ ἁγιάζεν ὑμᾶς.
\vs{9}Ἄνθρωπος ἄνθρωπος, ὃς ἂν κακῶς εἴπῃ τὸν πατέρα αὐτοῦ ἢ τὴν μητέρα αὐτοῦ, θανάτῳ θανατούσθω· πατέρα αὐτοῦ ἢ μητέρα αὐτοῦ κακῶς εἶπεν; ἔνοχος ἔσται.

\vs{10}Ἄνθρωπος ὃς ἂν μοιχεύσηται γυναῖκα ἀνδρός, ἢ ὃς ἂν μοιχεύσηται γυναῖκα τοῦ πλησίον, θανάτῳ θανατούσθωσαν, ὁ μοιχεύων καὶ ἡ μοιχευομένη.
\vs{11}Καὶ ἐάν τις κοιμηθῇ μετὰ γυναικὸς τοῦ πατρὸς αὐτοῦ, ἀσχημοσύνην τοῦ πατρὸς αὐτοῦ ἀπεκάλυσε· θανάτῳ θανατούσθωσαν ἀμφότεροι, ἔνοχοί εἰσι.
\vs{12}Καὶ ἐάν τις κοιμηθῇ μετὰ νύμφης αὐτοῦ, θανάτῳ θανατούσθωσαν ἀμφότεροι· ἠσεβήκασι γάρ, ἔνοχοί εἰσι.
\vs{13}Καὶ ὃς ἂν κοιμηθῇ μετὰ ἄρσενος κοίτην γυναικὸς, βδέλυγμα ἐποίησαν ἀμφότεροι· θανάτῳ θανατούσθωσαν, ἔνοχοι εἰσιν.
\vs{14}Ὃς ἂν λάβῃ γυναῖκα καὶ τὴν μητέρα αὐτῆς, ἀνόμημά ἐστιν· ἐν πυρὶ κατακαύσουσιν αὐτὸν καὶ αὐτὰς, καὶ οὐκ ἔσται ἀνομία ἐν ὑμῖν.
\vs{15}Καὶ ὃς ἂν δῷ κοιτασίαν αὐτοῦ ἐν τετράποδι, θανάτῳ θανάτούσθω, καὶ τὸ τετράπουν ἀποκτενεῖτε.
\vs{16}Καὶ γυνὴ ἥτις προσελεύσεται πρὸς πᾶν κτῆνος βιβασθῆναι αὐτὴν ὑπʼ αὐτοὺ, ἀποκτενεῖτε τὴν γυναῖκα καὶ τὸ κτῆνος· θανάτῳ θανατούσθωσαν, ἔνοχοί εἰσιν.
\vs{17}Ὃς ἂν λάβῃ τὴν ἀδελφὴν αὐτοῦ ἐκ πατρὸς αὐτοῦ ἢ ἐκ μητρὸς αὐτοῦ, καὶ ἴδῃ τὴν ἀσχημοσύνην αὐτῆς, καὶ αὕτη ἴδῃ τὴν ἀσχημοσύνην αὐτοῦ, ὄνειδός ἐστιν, ἐξολοθρευθήσονται ἐνωπιον υἱῶν γένους αὐτῶν· ἀσχημοσύνην ἀδελφῆς αὐτοῦ ἀπεκάλυψεν, ἁμαρτίαν κομιοῦνται.
\vs{18}Καὶ ἀνὴρ ὃς ἂν κοιμηθῇ μετὰ γυναικὸς ἀποκαθημένης, καὶ ἀποκαλύψῃ τὴν ἀσχημοσύνην αὐτῆς, τὴν πηγὴν αὐτῆς ἀπεκάλυψε, καὶ αὕτη ἀπεκάλυψε τὴν ῥύσιν τοῦ αἵματος αὐτῆς· ἐξολοθρευθήσονται ἀμφότεροι ἐκ τῆς γενεᾶς αὐτῶν.
\vs{19}Καὶ ἀσχημοσύνην ἀδελφῆς πατρὸς σου, καὶ ἀδελφῆς μητρός σου οὐκ ἀποκαλύψεις· τὴν γὰρ οἰκειότητα ἀπεκάλυψεν, ἁμαρτίαν ἀποίσονται.
\vs{20}Ὃς ἂν κοιμηθῇ μετὰ τῆς συγγενοῦς αὐτοῦ, ἀσχημοσύνην τῆς συγγενείας αὐτοῦ ἀπεκάλυψεν, ἄτεκνοι ἀποθανοῦνται.
\vs{21}Ὃς ἐὰς λάβῃ γυναῖκα τοῦ ἀδελφοῦ αὐτοῦ, ἀκαθαρσία ἐστίν· ἀσχημοσύνην τοῦ ἀδελφοῦ αὐτοῦ ἀπεκάλυψεν, ἄτεκνοι ἀποθανοῦνται.

\vs{22}Καὶ φυλάξασθε πάντα τὰ προστάγματά μου, καὶ τὰ κρίματά μου, καὶ ποιήσετε αὐτὰ, καὶ οὐ μὴ προσοχθίσῃ ὑμῖν ἡ γῆ, εἰς ἣν ἐγὼ εἰσάγω ὑμᾶς ἐκεῖ κατοικεῖν ἐπʼ αὐτῆς.
\vs{23}Καὶ οὐχὶ πορεύεσθε τοῖς νομίμοις τῶν ἐθνῶν, οὓς ἐξαποστέλλω ἀφʼ ὑμῶν· ὅτι ταῦτα πάντα ἐποίησαν, καὶ ἐβδελυξάμην αὐτούς.
\vs{24}Καὶ εἶπα ὑμῖν, ὑμεῖς κληρονομήσετε τὴν γῆν αὐτῶν, καὶ ἐγὼ δώσω ὑμῖν αὐτὴν ἐν κτήσει, γῆν ῥέουσαν γάλα καὶ μέλι· ἐγὼ Κύριος ὁ Θεὸς ὑμῶν, ὃς διώρισα ὑμᾶς ἀπὸ πάντων τῶν ἐθνῶν.
\vs{25}Καὶ ἀφοριεῖτε αὐτοὺς ἀναμέσον τῶν κτηνῶν τῶν καθαρῶν καὶ ἀναμέσον τῶν κτηνῶν τῶν ἀκαθάρτων, καὶ ἀναμέσον τῶν πετεινῶν τῶν καθαρῶν καὶ τῶν ἀκαθάρτων· καὶ οὐ βδελύξετε τὰς ψυχὰς ὑμῶν ἐν τοῖς κτήνεσι, καὶ ἐν τοῖς πετεινοῖς, καὶ ἐν πᾶσι τοῖς ἑρπετοῖς τῆς γῆς ἃ ἐγὼ ἀφώρισα ὑμῖν ἐν ἀκαθαρσίᾳ.
\vs{26}Καὶ ἔσεσθέ μοι ἅγιοι, ὅτι ἐγὼ ἅγιός εἶμι Κύριος ὁ Θεὸς ὑμῶν, ὁ ἀφορίσας ὑμᾶς ἀπὸ πάντων τῶν ἐθνῶν, εἶναί μοι.

\vs{27}Καὶ ἀνὴρ ἢ γυνὴ ὃς ἂν γένηται αὐτῶν ἐγγαστρίμυθος ἢ ἐπαοιδός, θανάτῳ θανατούσθωσαν ἀμφότεροι· λίθοις λιθοβολήσετε αὐτοὺς, ἔνοχοί εἰσι.

\ch{21}
Καὶ εἶπε Κύριος πρὸς Μωυσῆν, λέγων, εἶπον τοῖς ἱερεύσι τοῖς υἱοῖς Ἀαρὼν, καὶ ἐρεῖς πρὸς αὐτούς, ἐν ταῖς ψυχαῖς οὐ μιανθήσονται ἐν τῷ ἔθνει αὐτῶν,
\vs{2}ἀλλʼ ἢ ἐν τῷ οἰκείῳ τῷ ἔγγιστα αὐτῶν, ἐπὶ πατρὶ καὶ μητρὶ, καὶ υἱοῖς, καὶ θυγατράσιν, ἐπʼ ἀδελφῷ,
\vs{3}καὶ ἐπʼ ἀδελφῇ παρθένῳ τῇ ἐγγιζούσῃ αὐτῷ, τῇ μὴ ἐκδεδομένῃ ἀνδρί, ἐπὶ τούτοις μιανθήσεται.
\vs{4}Οὐ μιανθήσεται ἐξάπινα ἐν τῷ λαῷ αὐτοῦ εἰς βεβήλωσιν αὐτοῦ.
\vs{5}Καὶ φαλάκρωμα οὐ ξυρηθήσεσθε τὴν κεφαλὴν ἐπὶ νεκρῷ· καὶ τὴν ὄψιν τοῦ πώγωνος οὐ ξυρήσονται· καὶ ἐπὶ τὰς σάρκας αὐτῶν οὐ κατατεμοῦσιν ἐντομίδας.
\vs{6}Ἅγιοι ἔσονται τῷ Θεῷ αὐτῶν, καὶ οὐ βεβηλώσουσι τὸ ὄνομα τοῦ Θεοῦ αὐτῶν· τὰς γὰρ θυσίας Κυρίου δῶρα τοῦ Θεοῦ αὐτῶν αὐτοὶ προσφέρουσι, καὶ ἔσονται ἅγιοι.
\vs{7}Γυναῖκα πόρνην καὶ βεβηλωμένην οὐ λήψονται, καὶ γυναῖκα ἐκβεβλημένην ἀπὸ ἀνδρὸς αὐτῆς, ὅτι ἅγιός ἐστι Κυρίῳ τῷ Θεῷ αὐτοῦ.
\vs{8}Καὶ ἁγιάσεις αὐτόν· τὰ δῶρα Κυρίου τοῦ Θεοῦ ὑμῶν οὗτος προσφέρει, ἅγιος ἔσται· ὅτι ἅγιος ἐγὼ Κύριος ὁ ἁγιάζων αὐτούς.
\vs{9}Καὶ θυγάτηρ ἀνθρώπου ἱερέως ἐὰν βεβηλωθῇ τοῦ ἐκπορνεύσαι, τὸ ὄνομα τοῦ πατρὸς αὐτῆς αὐτὴ βεβηλοῖ· ἐπὶ πυρὸς κατακαυθήσεται.

\vs{10}Καὶ ὁ ἱερεὺς ὁ μέγας ἀπὸ τῶν ἀδελφῶν αὐτοῦ, τοῦ ἐπικεχυμένου ἐπὶ τὴν κεφαλὴν τοῦ ἐλαίου τοῦ χριστοῦ, καὶ τετελειωμένου ἐνδύσασθαι τὰ ἱμάτια, τὴν κεφαλὴν οὐκ ἀποκιδαρώσει, καὶ τὰ ἱμάτια οὐ διαῤῥήξει,
\vs{11}καὶ ἐπὶ πάσῃ ψυχῇ τετελευτηκυίᾳ οὐκ εἰσελεύσεται, ἐπὶ πατρὶ αὐτοῦ οὐδὲ ἐπὶ μητρὶ αὐτοῦ οὐ μιανθήσεται.
\vs{12}Καὶ ἐκ τῶν ἁγίων οὐκ ἐξελεύσεται, καὶ οὐ βεβηλώσει τὸ ἡγιασμένον τοῦ Θεοῦ αὐτοῦ, ὅτι τὸ ἅγιον ἔλαιον τὸ χριστὸν τοῦ Θεοῦ ἐπʼ αὐτῷ· ἐγὼ Κύριος.
\vs{13}Οὗτος γυναῖκα παρθένον ἐκ τοῦ γένους αὐτοῦ λήψεται.
\vs{14}Χήραν δὲ καὶ ἐκβεβλημένην καὶ βεβηλωμένην καὶ πόρνην, ταύτας οὐ λήψεται, ἀλλʼ ἢ παρθένον ἐκ τοῦ λαοῦ αὐτοῦ λήψεται γυναῖκα.
\vs{15}Καὶ οὐ βεβηλώσει τὸ σπέρμα αὐτοῦ ἐν τῷ λαῷ αὐτοῦ· ἐγὼ Κύριος ὁ ἁγιάζων αὐτόν.
\vs{16}Καὶ ἐλάλησε Κύριος πρὸς Μωυσῆν, λέγων,
\vs{17}εἶπον Ἀαρὼν, ἄνθρωπος ἐκ τοῦ γένους σου εἰς τὰς γενεὰς ὑμῶν, τινὶ ἐὰν ᾖ ἐν αὐτῷ μῶμος, οὐ προσελεύσεται προσφέρειν τὰ δῶρα τοῦ Θεοῦ αὐτοῦ.
\vs{18}πᾶς ἄνθρωπος ᾧ ἐστιν ἐν αὐτῷ μῶμος, οὐ προσελεύσεται· ἄνθρωπος τυφλὸς.
\vs{19}ἢ χωλὸς, ἢ κολοβόριν, ἢ ὠτότμητος, ἢ ἄνθρωπος ᾧ ἂν ᾖ ἐν αὐτῳ σύντριμμα χειρὸς, ἢ σύντριμμα ποδὸς,
\vs{20}ἢ κυρτὸς, ἢ ἔφηλος, ἢ πτίλλος τοὺς ὀφθαλμούς, ἢ ἄνθρωπος ᾧ ἂν ᾖ ἐν αὐτῷ ψώρα ἀγρία, ἢ λειχὴν, ἢ μονόρχις.
\vs{21}Πᾶς ᾧ ἐστιν ἐν αὐτῷ μῶμος, ἐκ τοῦ σπέρματος Ἀαρὼν τοῦ ἱερέως, οὐκ ἐγγιεῖ τοῦ προσενεγκεῖν τὰς θυσίας τῷ Θεῷ σου, ὅτι μῶμος ἐν αὐτῷ· τὰ δῶρα τοῦ Θεοῦ οὐ προσελεύσεται προσενεγκεῖν.
\vs{22}Τὰ δῶρα τοῦ Θεοῦ τὰ ἅγια τῶν ἁγίων, καὶ ἀπὸ τῶν ἁγίων φάγεται.
\vs{23}Πλὴν πρὸς τὸ καταπέτασμα οὐ προσελεύσεται, καὶ πρὸς τὸ θυσιαστήριον οὐκ ἐγγιεῖ, ὅτι μῶμον ἔχει· καὶ οὐ βεβηλώσει τὸ ἅγιον τοῦ Θεοῦ αὐτοῦ, ὅτι ἐγώ εἰμι Κύριος ὁ ἁγιάζων αὐτούς.
\vs{24}Καὶ ἐλάλησε Μωυσῆς πρὸς Ἀαρὼν καὶ τοὺς υἱοὺς αὐτοῦ, καὶ πρὸς πάντας υἱοὺς Ἱσραήλ.

\ch{22}
Καὶ ἐλάλησε Κύριος πρὸς Μωυσῆν, λέγων,
\vs{2}εἶπον Ἀαρὼν καὶ τοῖς υἱοῖς αὐτοῦ· καὶ προσεχέτωσαν ἀπὸ τῶν ἁγίων τῶν υἱῶν Ἰσραὴλ, καὶ οὐ βεβηλώσουσι τὸ ὄνομα τὸ ἅγιόν μου, ὅσα αὐτοὶ ἁγιάζουσί μοι· ἐγὼ Κύριος.
\vs{3}Εἶπον αὐτοῖς, εἰς τὰς γενεὰς ὑμῶν πᾶς ἄνθρωπος, ὃς ἂν προσέλθῃ ἀπὸ παντὸς τοῦ σπέρματος ὑμῶν πρὸς τὰ ἅγια, ὅσα ἂν ἁγιάζωσιν οἱ υἱοὶ Ἰσραὴλ τῷ Κυρίῳ, καὶ ἡ ἀκαθαρσία αὐτοῦ ἐπʼ αὐτῷ ᾖ, ἐξολοθρευθήσεται ἡ ψυχὴ ἐκείνη ἀπʼ ἐμοῦ· ἐγὼ Κύριος ὁ Θεὸς ὑμῶν.
\vs{4}Καὶ ἄνθρωπος ἐκ τοῦ σπέρματος Ἀαρὼν τοῦ ἱερέως, καὶ οὗτος λεπρᾷ ἢ γονοῤῥυεῖ, τῶν ἁγίων οὐκ ἔδεται, ἕως ἂν καθαρισθῇ· καὶ ὁ ἁπτόμενος πάσης ἀκαθαρσίας ψυχῆς, ἢ ἄνθρωπος ᾧ ἂν ἐξέλθῃ ἐξ αὐτοῦ κοίτη σπέρματος,
\vs{5}ἢ ὅστις ἂν ἅψηται παντὸς ἑρπετοῦ ἀκαθάρτου, ὃ μιανεῖ αὐτὸν, ἢ ἐπʼ ἀνθρώπῳ, ἐν ᾧ μιανεῖ αὐτὸν κατὰ πᾶσαν ἀκαθαρσίαν αὐτοῦ·
\vs{6}Ψυχὴ ἥτις ἐὰν ἅψηται αὐτῶν, ἀκάθαρτος ἔσται ἕως ἑσπέρας· οὐκ ἔδεται ἀπὸ τῶν ἁγίων, ἐὰν μὴ λούσηται τὸ σῶμα αὐτοῦ ὕδατι.
\vs{7}Καὶ δύῃ ὁ ἥλιος, καὶ καθαρὸς ἔσται· καὶ τότε φάγεται τῶν ἁγίων, ὅτι ἄρτος αὐτοῦ ἐστι.
\vs{8}Θνησιμαῖον καὶ θηριάλωτον οὐ φάγεται, μιανθῆναι αὐτὸν ἐν αὐτοῖς· ἐγὼ Κύριος.
\vs{9}Καὶ φυλάξονται τὰ φυλάγματά μου, ἵνα μὴ λάβωσι διʼ αὐτὰ ἁμαρτίαν, καὶ ἀποθάνωσι διʼ αὐτὰ, ἐὰν βεβηλώσουσιν αὐτά· ἐγὼ Κύριος ὁ Θεὸς ὁ ἁγιάζων αὐτούς.
\vs{10}Καὶ πᾶς ἀλλογενὴς οὐ φάγεται ἅγια· πάροικος ἱερέως, ἢ μισθωτὸς, οὐ φάγεται ἅγια.
\vs{11}Ἐὰν δὲ ἱερεὺς κτήσηται ψυχὴν ἔγκτητον ἀργυρίου, οὗτος φάγεται ἐκ τῶν ἄρτων αὐτοῦ· καὶ οἱ οἰκογενεῖς αὐτοῦ, καὶ οὗτοι φάγονται τῶν ἄρτων αὐτοῦ.
\vs{12}Καὶ θυγάτηρ ἀνθρώπου ἱερέως ἐὰν γένηται ἀνδρὶ ἀλλογενεῖ, αὐτὴ τῶν ἀπαρχῶν ἁγίου οὐ φάγεται.
\vs{13}Καὶ θυγάτηρ ἱερέως ἐὰν γένηται χήρα ἢ ἐκβεβλημένη, σπέρμα δὲ μὴ ᾖ αὐτῇ, ἐπαναστρέψει ἐπὶ τὸν οἶκον τὸν πατρικὸν κατὰ τὴν νεότητα αὐτῆς· ἀπὸ τῶν ἄρτων τοῦ πατρὸς αὐτῆς φάγεται· καὶ πᾶς ἀλλογενὴς οὐ φάγεται ἀπʼ αὐτῶν.
\vs{14}Καὶ ἄνθρωπος ὃς ἂν φάγῃ ἅγια κατʼ ἄγνοιαν, καὶ προσθήσει τὸ ἐπίπεμπτον αὐτοῦ ἐπʼ αὐτὸ, καὶ δώσει τῷ ἱερεῖ τὸ ἅγιον.
\vs{15}Καὶ οὐ βεβηλώσουσι τὰ ἅγια τῶν υἱῶν Ἰσραὴλ, ἃ αὐτοὶ ἀφαιροῦσι τῷ Κυρίῳ,
\vs{16}καὶ ἐπάξουσιν ἐφʼ ἑαυτοὺς ἀνομίαν πλημμελείας ἐν τῷ ἐσθίειν αὐτοὺς τὰ ἅγια αὐτῶν, ὅτι ἐγὼ Κύριος ὁ ἁγιάζων αὐτούς.

\vs{17}Καὶ ἐλάλησε Κύριος πρὸς Μωυσῆν, λέγων,
\vs{18}λάλησον Ἀαρὼν καὶ τοῖς υἱοῖς αὐτοῦ, καὶ πάσῃ συναγωγῇ Ἰσραὴλ, καὶ ἐρεῖς πρὸς αὐτοὺς, ἄνθρωπος ἄνθρωπος ἀπὸ τῶν υἱῶν Ἰσραὴλ, ἢ τῶν προσηλύτων τῶν προσκειμένων πρὸς αὐτοὺς ἐν Ἰσραὴλ, ὃς ἂν προσενέγκῃ τὰ δῶρα αὐτοῦ κατὰ πᾶσαν ὁμολογίαν αὐτῶν, ἢ κατὰ πᾶσαν αἵρεσιν αὐτῶν, ὅσα ἂν προσενέγκωσι τῷ Θεῷ εἰς ὁλοκαύτωμα.
\vs{19}Δεκτὰ ὑμῖν ἄμωμα ἄρσενα ἐκ τῶν βουκολίων, ἢ ἐκ τῶν προβάτων, καὶ ἐκ τῶν αἰγῶν.
\vs{20}Πάντα ὅσα ἂν ἔχῃ μῶμον ἐν αὐτῷ οὐ προσάξουσι Κυρίῳ, διότι οὐ δεκτὸν ἔσται ὑμῖν.
\vs{21}Καὶ ἄνθρωπος ὃς ἂν προσενέγκῃ θυσίαν σωτηρίου τῷ Κυρίῳ, διαστείλας εὐχὴν ἢ κατὰ αἵρεσιν ἢ ἐν ταῖς ἑορταῖς ὑμῶν, ἐκ τῶν βουκολίων ἢ ἐκ τῶν προβάτων, ἄμωμον ἔσται εἰσδεκτὸν, πᾶς μῶμος οὐκ ἔσται ἐν αὐτῷ.
\vs{22}Τυφλὸν ἢ συντετριμμένον ἢ γλωσσότμητον ἢ μυρμηκιῶντα ἢ ψωραγριῶντα ἢ λειχῆνας ἔχοντα, οὐ προσάξουσι ταῦτα τῷ Κυρίῳ, καὶ εἰς κάρπωσιν οὐ δώσετε ἀπʼ αὐτῶν ἐπὶ τὸ θυσιαστήριον τῷ Κυρίῳ.
\vs{23}Καὶ μόσχον ἢ πρόβατον ὠτότμητον ἢ κολοβόκερκον, σφάγια ποιήσεις αὐτὰ σεαυτῷ, εἰς δὲ εὐχήν σου οὐ δεχθήσεται.
\vs{24}Θλαδίαν καὶ ἐκτεθλιμμένον καὶ ἐκτομίαν καὶ ἀπεσπασμένον, οὐ προσάξεις αὐτὰ τῷ Κυρίῳ, καὶ ἐπὶ τῆς γῆς ὑμῶν οὐ ποιήσετε.
\vs{25}Καὶ ἐκ χειρὸς ἀλλογενοῦς οὐ προσοίσετε τὰ δῶρα τοῦ Θεοῦ ὑμῶν ἀπὸ πάντων τούτων· ὅτι φθάρματά ἐστιν ἐν αὐτοῖς, μῶμος ἐν αὐτοῖς οὐ δεχθήσεται ταῦτα ὑμῖν.
\vs{26}Καὶ ἐλάλησε Κύριος πρὸς Μωυσῆν, λέγων,
\vs{27}μόσχον ἢ πρόβατον ἢ αἶγα, ὡς ἂν τεχθῇ, καὶ ἔσται ἑπτὰ ἡμέρας ὑπὸ τὴν μητέρα, τῇ δὲ ἡμέρᾳ τῇ ὀγδόῃ καὶ ἐπέκεινα δεχθήσεται εἰς δῶρα, κάρπωμα Κυρίῳ.
\vs{28}Καὶ μόσχον καὶ πρόβατον, αὐτὴν καὶ τὰ παιδία αὐτῆς, οὐ σφάξεις ἐν ἡμέρᾳ μιᾷ.

\vs{29}Ἐὰν δὲ θύσῃς θυσίαν εὐχὴν χαρμοσύνης Κυρίῳ, εἰσδεκτὸν ὑμῖν θύσετε αὐτό.
\vs{30}Αὐτῇ τῇ ἡμέρᾳ ἐκείνῃ βρωθήσεται· οὐκ ἀπολείψετε ἀπὸ τῶν κρεῶν εἰς τοπρωΐ· ἐγώ εἰμι Κύριος.
\vs{31}Καὶ φυλάξετε τὰς ἐντολάς μου, καὶ ποιήσετε αὐτάς.
\vs{32}Καὶ οὐ βεβηλώσετε τὸ ὄνομα τοῦ ἁγίου, καὶ ἁγιασθήσομαι ἐν μέσῳ τῶν υἱῶν Ἰσραήλ· ἐγὼ Κύριος ὁ ἁγιάζων ὑμᾶς,
\vs{33}ὁ ἐξαγαγὼν ὑμᾶς ἐκ γῆς Αἰγύπτου, ὥστε εἶναι ὑμῶν Θεός· ἐγὼ Κύριος.

\ch{23}
Καὶ εἶπε Κύριος πρὸς Μωυσῆν, λέγων,
\vs{2}λάλησον τοῖς υἱοῖς Ἰσραὴλ, καὶ ἐρεῖς πρὸς αὐτούς, αἱ ἑορταὶ Κυρίου ἃς καλέσετε αὐτὰς κλητὰς ἁγίας, αὗταί εἰσιν αἱ ἑορταί μου.
\vs{3}Ἓξ ἡμέρας ποιήσεις ἔργα, τῇ δὲ ἡμέρᾳ τῇ ἑβδόμῃ σάββατα, ἀνάπαυσις, κλητὴ ἁγία τῷ Κυρίῳ· πᾶν ἔργον οὐ ποιήσεις· σάββατά ἐστι τῷ Κυρίῳ ἐν πάσῃ κατοικίᾳ ὑμῶν.

\vs{4}Αὗται αἱ ἑορταὶ τῷ Κυρίῳ κληταὶ ἅγιαι, ἃς καλέσετε αὐτὰς ἐν τοῖς καιροῖς αὐτῶν.
\vs{5}Ἐν τῷ πρώτῳ μηνὶ, ἐν τῇ τεσσαρεσκαιδεκάτῃ ἡμέρᾳ τοῦ μηνὸς, ἀναμέσον τῶν ἑσπερινῶν πάσχα τῷ Κυρίῳ.
\vs{6}Καὶ ἐν τῇ πεντεκαιδεκάτῃ ἡμέρᾳ τοῦ μηνὸς τούτου ἑορτὴ τῶν ἀζύμων τῷ Κυρίῳ· ἑπτὰ ἡμέρας ἄζυμα ἔδεσθε.
\vs{7}Καὶ ἡμέρα ἡ πρώτη κλητὴ ἁγία ἔσται ὑμῖν· πᾶν ἔργον λατρευτὸν οὐ ποιήσετε.
\vs{8}Καὶ προσάξετε ὁλοκαυτώματα τῷ Κυρίῳ ἑπτὰ ἡμέρας· καὶ ἡ ἡμέρα ἡ ἑβδόμη κλητὴ ἁγία ἔσται ὑμῖν· πᾶν ἔργον λατρευτὸν οὐ ποιήσετε.
\vs{9}Καὶ ἐλάλησε Κύριος πρὸς Μωυσῆν, λέγων,
\vs{10}εἶπον τοῖς υἱοῖς Ἰσραὴλ, καὶ ἐρεῖς πρὸς αὐτοὺς, ὅταν εἰσέλθητε εἰς τὴν γῆν, ἣν ἐγὼ δίδωμι ὑμῖν, καὶ θερίζητε τὸν θερισμὸν αὐτῆς, καὶ οἴσετε τὸ δράγμα ἀπαρχὴν τοῦ θερισμοῦ ὑμῶν πρὸς τὸν ἱερέα·
\vs{11}Καὶ ἀνοίσει τὸ δράγμα ἔναντι Κυρίου δεκτὸν ὑμῖν· τῇ ἐπαύριον τῆς πρώτης ἀνοίσει αὐτό ὁ ἱερεύς.
\vs{12}Καὶ ποιήσετε ἐν τῇ ἡμέρᾳ ἐν ᾗ ἂν φέρητε τὸ δράγμα, πρόβατον ἄμωμον ἐνιαύσιον εἰς ὁλοκαύτωμα τῷ Κυρίῳ.
\vs{13}Καὶ τὴν θυσίαν αὐτοῦ δύο δέκατα σεμιδάλεως ἀναπεποιημένης ἐν ἐλαίῳ· θυσία τῷ Κυρίῳ, ὀσμὴ εὐωδίας Κυρίῳ· καὶ σπονδὴν αὐτοῦ τὸ τέταρτον τοῦ ἳν οἴνου.
\vs{14}Καὶ ἄρτον, καὶ πεφρυγμένα χίδρα νέα οὐ φάγεσθε ἕως εἰς αὐτὴν τὴν ἡμέραν ταύτην, ἕως ἂν προσενέγκητε ὑμεῖς τὰ δῶρα τῷ Θεῷ ὑμῶν· νόμιμον αἰώνιον εἰς τὰς γενεὰς ὑμῶν ἐν πάσῃ κατοικίᾳ ὑμῶν.

\vs{15}Καὶ ἀριθμήσετε ὑμῖν ἀπὸ τῆς ἐπαύριον τῶν σαββάτων, ἀπὸ τῆς ἡμέρας ἧς ἂν προσενέγκητε τὸ δράγμα τοῦ ἐπιθέματος, ἑπτὰ ἑβδομάδας ὁλοκλήρους,
\vs{16}ἕως τῆς ἐπαύριον τῆς ἐσχάτης ἑβδομάδος ἀριθμήσετε πεντήκοντα ἡμέρας, καὶ προσοίσετε θυσίαν νέαν τῷ Κυρίῳ.
\vs{17}Ἀπὸ τῆς κατοικίας ὑμῶν προσοίσετε ἄρτους ἐπίθεμα, δύο ἄρτους· ἐκ δύο δεκάτων σεμιδάλεως ἔσονται, ἐζυμωμένοι πεφθήσονται πρωτογεννημάτων τῷ Κυρίῳ.
\vs{18}Καὶ προσάξετε μετὰ τῶν ἄρτων ἑπτὰ ἀμνοὺς ἀμώμους ἐνιαυσίους, καὶ μόσχον ἕνα ἐκ βουκολίου, καὶ κριοὺς δύο ἀμώμους, καὶ ἔσονται ὁλοκαύτωμα τῷ Κυρίῳ· καὶ αἱ θυσίαι αὐτῶν καὶ αἱ σπονδαὶ αὐτῶν θυσία ὀσμὴ εὐωδίας τῷ Κυρίῳ.
\vs{19}Καὶ ποιήσουσι χίμαρον ἐξ αἰγῶν ἕνα περὶ ἁμαρτίας, καὶ δύο ἀμνοὺς ἐνιαυσίους εἰς θυσίαν σωτηρίου μετὰ τῶν ἄρτων τοῦ πρωτογεννήματος.
\vs{20}Καὶ ἐπιθήσει αὐτὰ ὁ ἱερεὺς μετὰ τῶν ἄρτων τοῦ πρωτογεννήματος ἐπίθεμα ἐναντίον Κυρίου μετὰ τῶν δύο ἀμνῶν, ἅγια ἔσονται τῷ Κυρίῳ· τῷ ἱερεῖ τῷ προσφέροντι αὐτὰ αὐτῷ ἔσται.
\vs{21}Καὶ καλέσετε ταύτην τὴν ἡμέραν κλητήν· ἁγία ἔσται ὑμῖν· πᾶν ἔργον λατρευτὸν οὐ ποιήσετε ἐν αὐτῇ· νόμιμον αἰώνιον εἰς τὰς γενεὰς ὑμῶν ἐν πάσῃ τῇ κατοικίᾳ ὑμῶν.
\vs{22}Καὶ ὅταν θερίζητε τὸν θερισμὸν τῆς γῆς ὑμῶν, οὐ συντελέσετε τὸ λοιπὸν τοῦ θερισμοῦ τοῦ ἀγροῦ σου ἐν τῷ θερίζειν σε, καὶ τὰ ἀποπίπτοντα τοῦ θερισμοῦ σου οὐ συλλέξεις· τῷ πτωχῷ καὶ τῷ προσηλύτῳ ὑπολείψεις αὐτά· ἐγὼ Κύριος ὁ Θεὸς ὑμῶν.

\vs{23}Καὶ ἐλάλησε Κύριος πρὸς Μωυσῆν, λέγων,
\vs{24}λάλησον τοῖς υἱοῖς Ἰσραὴλ, λέγων, τοῦ μηνὸς τοῦ ἐβδόμου μιᾷ τοῦ μηνὸς ἔσται ὑμῖν ἀνάπαυσις, μνημόσυνον σαλπίγγων· κλητὴ ἁγία ἔσται ὑμῖν·
\vs{25}Πᾶν ἔργον λατρευτὸν οὐ ποιήσετε· καὶ προσάξετε ὁλοκαύτωμα Κυρίῳ.

\vs{26}Καὶ ἐλάλησε Κύριος πρὸς Μωυσῆν, λέγων,
\vs{27}καὶ τῇ δεκάτῃ τοῦ μηνὸς τοῦ ἐβδόμου τούτου, ἡμέρα ἐξιλασμοῦ, κλητὴ ἁγία ἔσται ὑμῖν· καὶ ταπεινώσετε τὰς ψυχὰς ὑμῶν, καὶ προσάξετε ὁλοκαύτωμα τῷ Κυρίῳ.
\vs{28}Πᾶν ἔργον οὐ ποιήσετε ἐν αὐτῇ τῇ ἡμέρᾳ ταύτῃ· ἔστι γὰρ ἡμέρα ἐξιλασμοῦ αὕτη ὑμῖν, ἐξιλάσασθαι περὶ ὑμῶν ἔναντι Κυρίου τοῦ Θεοῦ ὑμῶν.
\vs{29}Πᾶσα ψυχὴ, ἥτις μὴ ταπεινωθήσεται ἐν αὐτῇ τῇ ἡμέρᾳ ταύτῃ, ἐξολοθρευθήσεται ἐκ τοῦ λαοῦ αὐτῆς.
\vs{30}Καὶ πᾶσα ψυχὴ, ἥτις ποιήσει ἔργον ἐν αὐτῇ τῇ ἡμέρᾳ ταύτῃ, ἀπολεῖται ἡ ψυχὴ ἐκείνη ἐκ τοῦ λαοῦ αὐτῆς.
\vs{31}Πᾶν ἔργον οὐ ποιήσετε· νόμιμον αἰώνιον εἰς τὰς γενεὰς ὑμῶν ἐν πάσαις κατοικίαις ὑμῶν.
\vs{32}Σάββατα σαββάτων ἔσται ὑμῖν· καὶ ταπεινώσετε τὰς ψυχὰς ὑμῶν· ἀπὸ ἐνάτης τοῦ μηνὸς, ἀπὸ ἑσπέρας ἕως ἑσπέρας σαββατιεῖτε τὰ σάββατα ὑμῶν.

\vs{33}Καὶ ἐλάλησε Κύριος πρὸς Μωυσῆν, λέγων,
\vs{34}λάλησον τοῖς υἱοῖς Ἰσραὴλ, λέγων, τῇ πεντεκαιδεκάτῃ τοῦ μηνὸς τοῦ ἐβδόμου τούτου, ἑορτὴ σκηνῶν ἑπτὰ ἡμέρας τῷ Κυρίῳ.
\vs{35}Καὶ ἡ ἡμέρα ἡ πρώτη κλητὴ ἁγία· πᾶν ἔργον λατρευτὸν οὐ ποιήσετε.
\vs{36}Ἑπτὰ ἡμέρας προσάξετε ὁλοκαυτώματα τῷ Κυρίῳ, καὶ ἡ ἡμέρα ἡ ὀγδόη κλητὴ ἁγία ἔσται ὑμῖν· καὶ προσάξετε ὁλοκαυτώματα Κυρίῳ· ἐξόδιόν ἐστι· πᾶν ἔργον λατρευτὸν οὐ ποιήσετε.
\vs{37}Αὗται ἑορταὶ Κυρὶῳ, ἃς καλέσετε κλητὰς ἁγίας, ὥστε προσενέγκαι καρπώματα τῷ Κυρίῳ, ὁλοκαυτώματα καὶ θυσίας αὐτῶν, καὶ σπονδὰς αὐτῶν τὸ καθʼ ἡμέραν εἰς ἡμέραν·
\vs{38}πλὴν τῶν σαββάτων Κυρίου, καὶ πλὴν τῶν δομάτων ὑμῶν, καὶ πλὴν πασῶν τῶν εὐχῶν ὑμῶν, καὶ πλὴν τῶν ἐκουσίων ὑμῶν, ἃ ἂν δώτε τῷ Κυρίῳ.
\vs{39}Καὶ ἐν τῇ πεντεκαιδεκάτῃ ἡμέρᾳ τοῦ μηνὸς τοῦ ἑβδόμου τούτου, ὅταν συντελέσητε τὰ γεννήματα τῆς γῆς, ἑορτάσατε τῷ Κυρίῳ ἑπτὰ ἡμέρας· τῇ ἡμέρᾳ τῇ πρώτῃ ἀνάπαυσις, καὶ τῇ ἡμέρᾳ τῇ ὀγδόῃ ἀνάπαυσις.
\vs{40}Καὶ λήψεσθε τῇ ἡμέρᾳ τῇ πρώτῃ καρπὸν ξύλου ὡραῖου, καὶ κάλλυνθρα φοινίκων, καὶ κλάδους ξύλου δασεῖς, καὶ ἰτέας, καὶ ἄγνου κλάδους ἐκ χειμάῤῥου, εὐφρανθῆναι ἔναντι Κυρίου τοῦ Θεοῦ ὑμῶν ἑπτὰ ἡμέρας τοῦ ἐνιαυτοῦ.
\vs{41}Νόμιμον αἰώνιον εἰς τὰς γενεὰς ὑμῶν· ἐν τῷ μηνὶ τῷ ἑβδόμῳ ἑορτάσετε αὐτήν.
\vs{42}Ἐν σκηναῖς κατοικήσετε ἑπτὰ ἡμέρας· πᾶς ὁ αὐτόχθων ἐν Ἰσραὴλ κατοικήσει ἐν σκηναῖς,
\vs{43}ὅπως ἴδωσιν αἱ γενεαὶ ὑμῶν, ὅτι ἐν σκηναῖς κατῴκισα τοὺς υἱοὺς Ἰσραὴλ, ἐν τῷ ἐξαγαγεῖν με αὐτοὺς ἐκ γῆς Αἰγύπτου· ἐγὼ Κύριος ὁ Θεὸς ὑμῶν.
\vs{44}Καὶ ἐλάλησε Μωυσῆς τὰς ἑορτὰς Κυρίου τοῖς υἱοῖς Ἰσραήλ.

\ch{24}
Καὶ ἐλάλησε Κύριος πρὸς Μωυσῆν, λέγων,
\vs{2}ἔντειλαι τοῖς υἱοῖς Ἰσραὴλ, καὶ λαβέτωσάν σοι ἔλαιον ἐλάϊνον καθαρὸν κεκομμένον εἰς φῶς, καῦσαι λύχνον διαπαντὸς,
\vs{3}ἔξωθεν τοῦ καταπετάσματος ἐν τῇ σκηνῇ τοῦ μαρτυρίου· καὶ καύσουσιν αὐτὸ Ἀαρὼν καὶ οἱ υἱοὶ αὐτοῦ ἀπὸ ἑσπέρας ἕως πρωῒ ἐνώπιον Κυρίου ἐνδελεχῶς, νόμιμον αἰώνιον εἰς τὰς γενεὰς ὑμῶν.
\vs{4}Ἐπὶ τῆς λυχνίας τῆς καθαρᾶς καύσετε τοὺς λύχνους ἐναντίον Κυρίου ἕως εἰς τοπρωΐ.
\vs{5}Καὶ λήμψεσθε σεμίδαλιν, καὶ ποιήσετε αὐτὴν δώδεκα ἄρτους· δύο δεκάτων ἔσται ὁ ἄρτος ὁ εἷς.
\vs{6}Καὶ ἐπιθήσετε αὐτοὺς δύο θέματα, ἓξ ἄρτους τὸ ἓν θέμα ἐπὶ τὴν τράπεζαν τὴν καθαρὰν ἔναντι Κυρίου.
\vs{7}Καὶ ἐπιθήσετε ἐπὶ τὸ θέμα λίβανον καθαρὸν καὶ ἅλα, καὶ ἔσονται εἰς ἄρτους εἰς ἀνάμνησιν προκείμενα τῷ Κυρίῳ.
\vs{8}Τῇ ἡμέρᾳ τῶν σαββάτων προσθήσεται ἔναντι Κυρίου διαπαντὸς ἐνώπιον τῶν υἱῶν Ἰσραὴλ, διαθήκην αἰώνιον.
\vs{9}Καὶ ἔσται Ἀαρὼν καὶ τοῖς υἱοῖς αὐτοῦ· καὶ φάγονται αὐτὰ ἐν τόπῳ ἁγίῳ· ἔστι γὰρ ἅγια τῶν ἁγίῳν τοῦτο αὐτῶν ἀπὸ τῶν θυσιαζομένων τῷ Κυρίῳ, νόμιμον αἰώνιον.

\vs{10}Καὶ ἐξῆλθεν υἱὸς γυναικὸς Ἰσραηλίτιδος, καὶ οὗτος ἦν υἱὸς Αἰγυπτίου ἐν τοῖς υἱοῖς Ἰσραήλ· καὶ ἐμαχέσαντο ἐν τῇ παρεμβολῇ ὁ ἐκ τῆς Ἰσραηλίτιδος, καὶ ὁ ἄνθρωπος ὁ Ἰσραηλίτης.
\vs{11}Καὶ ἐπονομάσας ὁ υἱὸς τῆς γυναικὸς τῆς Ἰσραηλίτιδος τὸ ὄνομα κατηράσατο· καὶ ἤγαγον αὐτὸν πρὸς Μωυσῆν· καὶ τὸ ὄνομα τῆς μητρὸς αὐτοῦ Σαλωμεὶθ θυγάτηρ Δαβρεὶ ἐκ τῆς φυλῆς Δάν.
\vs{12}Καὶ ἀπέθεντο αὐτὸν εἰς φυλακὴν διακρῖναι αὐτὸν διὰ προστάγματος Κυρίου.
\vs{13}Καὶ ἐλάλησε Κύριος πρὸς Μωυσῆν, λέγων,
\vs{14}ἐξάγαγε τὸν καταρασάμενον ἔξω τῆς παρεμβολῆς, καὶ ἐπιθήσουσι πάντες οἱ ἀκούσαντες τὰς χεῖρας αὐτῶν ἐπὶ τὴν κεφαλὴν αὐτοῦ, καὶ λιθοβολήσουσιν αὐτὸν πᾶσα ἡ συναγωγή.
\vs{15}Καὶ τοῖς υἱοῖς Ἰσραὴλ λάλησον, καὶ ἐρεῖς πρὸς αὐτοὺς, ἄνθρωπος ὃς ἐὰν καταράσηται Θεὸν, ἁμαρτίαν λήψεται.
\vs{16}Ὀνομάζων δὲ τὸ ὄνομα Κυρίου, θανάτῳ θανατούσθω· λίθοις λιθοβολείτω αὐτὸν πᾶσα ἡ συναγωγὴ Ἰσραήλ· ἐάν τε προσήλυτος ἐάν τε αὐτόχθων, ἐν τῷ ὀνομάσαι αὐτὸν τὸ ὄνομα Κυρίου, τελευτάτω.
\vs{17}Καὶ ἄνθρωπος ὃς ἂν πατάξῃ ψυχὴν ἀνθρώπου, καὶ ἀποθάνῃ, θανάτῳ θανατούσθω.
\vs{18}Καὶ ὃς ἂν πατάξῃ κτῆνος, καὶ ἀποθάνῃ, ἀποτισάτω ψυχὴν ἀντὶ ψυχῆς.
\vs{19}Καὶ ἐάν τις δῷ μῶμον τῷ πλησίον, ὡς ἐποίησεν αὐτῷ, ὡσαύτως ἀντιποιηθήσεται αὐτῷ·
\vs{20}Σύντριμμα ἀντὶ συντρίμματος, ὀφθαλμὸν ἀντὶ ὀφθαλμοῦ, ὀδόντα ἀντὶ ὀδόντος, καθότι ἂν δῷ μῶμον τῷ ἀνθρώπῳ, οὕτω δοθήσεται αὐτῷ.
\vs{21}Ὃς ἂν πατάξῃ ἄνθρωπον, καὶ ἀποθάνῃ, θανάτῳ θανατούσθω.
\vs{22}Δικαίωσις μία ἔσται τῷ προσηλύτῳ καὶ τῷ ἐγχωρίῳ, ὅτι ἐγώ εἰμι Κύριος ὁ Θεὸς ὑμῶν.
\vs{23}Καὶ ἐλάλησε Μωυσῆς τοῖς υἱοῖς Ἰσραήλ· καὶ ἐξήγαγον τὸν καταρασάμενον ἔξω τῆς παρεμβολῆς, καὶ ἐλιθοβόλησαν αὐτὸν ἐν λίθοις· καὶ οἱ υἱοὶ Ἰσραὴλ ἐποίησαν καθάπερ συνέταξε Κύριος τῷ Μωυσῇ.

\ch{25}
Καὶ ἐλάλησε Κύριος πρὸς Μωυσῆν ἐν τῷ ὄρει Σινᾷ, λέγων,
\vs{2}λάλησον τοῖς υἱοῖς Ἰσραὴλ, καὶ ἐρεῖς πρὸς αὐτοὺς, ὅταν εἰσέλθητε εἰς τὴν γῆν, ἣν ἐγὼ δίδωμι ὑμῖν, καὶ ἀναπαύσεται ἡ γῆ, ἣν ἐγὼ δίδωμι ὑμῖν, σάββατα τῷ Κυρίῳ.
\vs{3}Ἓξ ἔτη σπερεῖς τὸν ἀγρόν σου, καὶ ἓξ ἔτη τεμεῖς τὴν ἄμπελόν σου, καὶ συνάξεις τὸν καρπὸν αὐτῆς.
\vs{4}Τῷ δὲ ἔτει τῷ ἑβδόμῳ σάββατα· ἀνάπαυσις ἔσται τῇ γῇ, σάββατα τῷ Κυρίῳ· τὸν ἀγρόν σου οὐ σπερεῖς, καὶ τὴν ἄμπελόν σου οὐ τεμεῖς,
\vs{5}καὶ τὰ αὐτόματα ἀναβαίνοντα τοῦ ἀγροῦ σου οὐκ ἐκθερίσεις, καἰ τὴν σταφυλὴν τοῦ ἁγιάσματός σου οὐκ ἐκτρυγήσεις· ἐνιαυτὸς ἀναπαύσεως ἔσται τῇ γῇ.
\vs{6}Καὶ ἔσται τὰ σάββατα τῆς γῆς βρώματά σοι, καὶ τῷ παιδί σου, καὶ τῇ παιδίσκῃ σου, καὶ τῷ μισθωτῷ σου, καὶ τῷ παροίκῳ τῷ προσκειμένῳ πρὸς σέ.
\vs{7}Καὶ τοῖς κτήνεσί σου, καὶ τοῖς θηρίοις τοῖς ἐν τῇ γῇ σου ἔσται πᾶν τὸ γέννημα αὐτοῦ εἰς βρῶσιν.

\vs{8}Καὶ ἐξαριθμήσεις σεαυτῷ ἑπτὰ ἀναπαύσεις ἐτῶν, ἑπτὰ ἔτη ἑπτάκις· καὶ ἔσονταί σοι ἑπτὰ ἑβδομάδες ἐτῶν ἐννέα καὶ τεσσαράκοντα ἔτη.
\vs{9}Διαγγελεῖτε σάλπιγγος φωνῇ ἐν πάσῃ τῇ γῇ ὑμῶν ἐν τῷ μηνὶ τῷ ἑβδόμῳ τῇ δεκάτῃ τοῦ μηνός· τῇ ἡμέρᾳ τοῦ ἱλασμοῦ διαγγελεῖτε σάλπιγγι ἐν πάσῃ τῇ γῇ ὑμῶν.
\vs{10}Καὶ ἁγιάσετε τὸ ἔτος τὸν πεντηκοστὸν ἐνιαυτὸν, καὶ διαβοήσετε ἄφεσιν ἐπὶ τῆς γῆς πᾶσι τοῖς κατοικοῦσιν αὐτήν· ἐνιαυτὸς ἀφέσεως σημασία αὕτη ἔσται ὑμῖν· καὶ ἀπελεύσεται εἷς ἕκαστος εἰς τὴν κτῆσιν αὐτοῦ, καὶ ἕκαστος εἰς τὴν πατριὰν αὐτοῦ ἀπελεύσεσθε.
\vs{11}Ἀφέσεως σημασία αὕτη, τὸ ἔτος τὸ πεντηκοστὸν ἐνιαυτὸς ἔσται ὑμῖν· οὐ σπερεῖτε, οὐδὲ ἀμήσετε τὰ αὐτόματα ἀναβαίνοντα αὐτῆς, καὶ οὐ τρυγήσετε τὰ ἡγιασμένα αὐτῆς,
\vs{12}ὅτι ἀφέσεως σημασία ἐστίν· ἅγιον ἔσται ὑμῖν· ἀπὸ τῶν πεδίων φάγεσθε τὰ γεννήματα αὐτῆς.
\vs{13}Ἐν τῷ ἔτει τῆς ἀφέσεως σημασίας αὐτῆς ἐπανελεύσεται εἰς τὴν ἔγκτησιν αὐτοῦ.
\vs{14}Ἐὰν δὲ ἀποδῷ πράσιν τῷ πλησίον σου, ἐὰν δὲ καὶ κτήσῃ παρὰ τοῦ πλησίον σου, μὴ θλιβέτω ἄνθρωπος τὸν πλησίον.
\vs{15}Κατὰ ἀριθμὸν ἐτῶν μετὰ τὴν σημασίαν κτήσῃ παρὰ τοῦ πλησίον, κατὰ ἀριθμὸν ἐνιαυτῶν γεννημάτων ἀποδώσεταί σοι.
\vs{16}Καθότι ἂν πλεῖον τῶν ἐτῶν πληθυνεῖ τὴν ἔγκτησιν αὐτοῦ, καὶ καθότι ἂν ἔλαττον τῶν ἐτῶν ἐλαττονώσει τὴν ἔγκτησιν αὐτοῦ· ὅτι ἀριθμὸν γεννημάτων αὐτοῦ, οὕτως ἀποδώσεταί σοι.
\vs{17}Μὴ θλιβέτω ἄνθρωπος τὸν πλησίον· καὶ φοβηθήσῃ Κύριον τὸν Θεόν σου· ἐγώ εἰμι Κύριος ὁ Θεὸς ὑμῶν.

\vs{18}Καὶ ποιήσετε πάντα τὰ δικαιώματά μου, καὶ πάσας τὰς κρίσεις μου, καὶ φυλάξασθε, καὶ ποιήσετε αὐτὰ, καὶ κατοικήσετε ἐπὶ τῆς γῆς πεποιθότες.
\vs{19}Καὶ δώσει ἡ γῆ τὰ ἐκφόρια αὐτῆς, καὶ φάγεσθε εἰς πλησμονὴν, καὶ κατοικήσετε πεποιθότες ἐπʼ αὐτῆς.
\vs{20}Ἐὰν δὲ λέγητε, τί φαγόμεθα ἐν τῷ ἔτει τῷ ἑβδόμῳ τούτῳ, ἐὰν μὴ σπείρωμεν μηδὲ συναγάγωμεν τὰ γεννήματα ἡμῶν;
\vs{21}Καὶ ἀποστέλλω τὴν εὐλογίαν μου ὑμῖν ἐν τῷ ἔτει τῷ ἕκτῳ, καὶ ποιήσει τὰ γεννήματα αὐτῆς εἰς τὰ τρία ἔτη.
\vs{22}Καὶ σπερεῖτε τὸ ἔτος τὸ ὄγδοον, καὶ φάγεσθε ἀπὸ τῶν γεννημάτων παλαιὰ ἕως τοῦ ἔτους τοῦ ἐνάτου· ἕως ἂν ἔλθῃ τὸ γέννημα αὐτῆς, φάγεσθε παλαιὰ παλαιῶν.
\vs{23}Καὶ ἡ γῆ οὐ πραθήσεται εἰς βεβαίωσιν· ἐμὴ γάρ ἐστιν ἡ γῆ, διότι προσήλυτοι καὶ πάροικοι ὑμεῖς ἐστε ἐναντίον μου.
\vs{24}Καὶ κατὰ πᾶσαν γῆν κατασχέσεως ὑμῶν, λύτρα δώσετε τῆς γῆς.
\vs{25}Ἐὰν δὲ πένηται ὁ ἀδελφός σου ὁ μετὰ σοῦ, καὶ ἀποδῶται ἀπὸ τῆς κατασχέσεως αὐτοῦ, καὶ ἔλθῃ ὁ ἀγχιστεύων ὁ ἐγγίζων αὐτῷ, καὶ λυτρώσεται τὴν πρᾶσιν τοῦ ἀδελφοῦ αὐτοῦ.
\vs{26}Ἐὰν δὲ μὴ ᾖ τινι ὁ ἀγχιστεύων, καὶ εὐπορηθῇ τῇ χειρὶ, καὶ εὑρεθῇ αὐτῷ τὸ ἱκανὸν, λύτρα αὐτοῦ·
\vs{27}καὶ συλλογιεῖται τὰ ἔτη τῆς πράσεως αὐτοῦ, καὶ ἀποδώσει ὅ ὑπερέχει τῷ ἀνθρώπῳ, ᾧ ἀπέδοτο αὐτὸ αὐτῷ, καὶ ἀπελεύσεται εἰς τὴν κατάσχεσιν αὐτοῦ.
\vs{28}Ἐὰν δὲ μὴ εὑπορηθῇ αὐτοῦ ἡ χεὶρ τὸ ἱκανὸν, ὥστε ἀποδοῦναι αὐτῷ, καὶ ἔσται ἡ πράσις τῷ κτησαμένῳ αὐτὰ ἕως τοῦ ἕκτου ἔτους τῆς ἀφέσεως, καὶ ἐξελεύσεται ἐν τῇ ἀφέσει, καὶ ἀπελεύσεται εἰς τὴν κατάσχεσιν αὐτοῦ.

\vs{29}Ἐὰν δέ τις ἀποδῶται οἰκίαν οἰκητὴν ἐν πόλει τετειχισμένῃ, καὶ ἔσται ἡ λύτρωσις αὐτῆς, ἕως πληρωθῇ· ἐνιαυτὸς ἡμερῶν ἔσται ἡ λύτρωσις αὐτῆς.
\vs{30}Ἐὰν δὲ μὴ λυτρωθῇ ἕως ἂν πληρωθῇ αὐτῆς ἐνιαυτὸς ὅλος, κυρωθήσεται ἡ οἰκία ἡ οὖσα ἐν πόλει τῇ ἐχούσῃ τεῖχος, βεβαίως τῷ κτησαμένῳ αὐτὴν εἰς τὰς γενεὰς αὐτοῦ, καὶ οὐκ ἐξελεύσεται ἐν τῇ ἀφέσει.
\vs{31}Αἱ δὲ οἰκίαι αἱ ἐν ἐπαύλεσιν, αἷς οὐκ ἔστιν ἐν αὐταῖς τεῖχος κύκλῳ, πρὸς τὸν ἀγρὸν τῆς γῆς λογισθήσονται· λυτρωταὶ διαπαντὸς ἔσονται, καὶ ἐν τῇ ἀφέσει ἐξελεύσονται.
\vs{32}Καὶ αἱ πόλεις τῶν Λευιτῶν, οἰκίαι τῶν πόλεων κατασχέσεως αὐτῶν, λυτρωταὶ διαπαντὸς ἔσονται τοῖς Λευίταις.
\vs{33}Καὶ ὃς ἂν λυτρώσηται παρὰ τῶν Λευιτῶν, καὶ ἐξελεύσεται ἡ διάπρασις αὐτῶν οἰκιῶν πόλεως κατασχέσεως αὐτῶν ἐν τῇ ἀφέσει, ὅτι οἰκίαι τῶν πόλεων τῶν Λευιτῶν κατάσχεσις αὐτῶν ἐν μέσῳ υἱῶν Ἰσραήλ.
\vs{34}Καὶ οἱ ἀγροὶ ἀφωρισμένοι ταῖς πόλεσιν αὐτῶν οὐ πραθήσονται, ὅτι κατάσχεσις αἰωνία τοῦτο αὐτῶν ἐστον.

\vs{35}Ἐὰν δὲ πένηται ὁ ἀδελφός σοῦ ὁ μετὰ σοῦ, καὶ ἀδυνατήσῃ ταῖς χερσὶ παρὰ σοὺ, ἀντιλήψῃ αὐτοῦ ὡς προσηλύτου καὶ παροίκου, καὶ ζήσεται ὁ ἀδελφός σου μετὰ σοῦ.
\vs{36}Οὐ λήψῃ παρʼ αὐτοῦ τόκον, οὐδὲ ἐπὶ πλήθει· καὶ φοβηθήσῃ τὸν Θεόν σου· ἐγὼ Κύριος· καὶ ζήσεται ὁ ἀδελφός σου μετὰ σοῦ.
\vs{37}Τὸ ἀργύριόν σου οὐ δώσεις αὐτῷ ἐπὶ τὸκῳ, καὶ ἐπὶ πλεονασμῷ οὐ δώσεις αὐτῷ τὰ βρώματά σου.
\vs{38}Ἐγὼ Κύριος ὁ Θεὸς ὑμῶν, ὁ ἐξαγαγὼν ὑμᾶς ἐκ γῆς Αἰγύπτου, δοῦναι ὑμῖν τὴν γῆν Χαναὰν, ὥστε εἶναι ὑμῶν Θεός.

\vs{39}Ἐὰν δὲ ταπεινωθῇ ὁ ἀδελφός σου παρὰ σοὶ, καὶ πραθῇ σοι, οὐ δουλεύσει σοι δουλείαν οἰκέτου.
\vs{40}Ὡς μισθωτὸς ἢ πάροικος ἔσται σοι· ἕως τοῦ ἔτους τῆς ἀφέσεως ἐργᾶται παρὰ σοί,
\vs{41}καὶ ἐξελεύσεται τῇ ἀφέσει, καὶ τὰ τέκνα αὐτοῦ μετʼ αὐτοῦ, καὶ ἀπελεύσεται εἰς τὴν γενεὰν αὐτοῦ, εἰς τὴν κατάσχεσιν τὴν πατρικὴν ἀποδραμεῖται.
\vs{42}Διότι οἰκέται μου εἰσὶν οὗτοι, οὓς ἐξήγαγον ἐκ γῆς Αἰγύπτου· οὐ πραθήσεται ἐν πράσει οἰκέτου.
\vs{43}Οὐ κατατενεῖς αὐτὸν ἐν τῷ μόχθῳ, καὶ φοβηθήσῃ Κύριον τὸν Θεόν σου,
\vs{44}καὶ παῖς καὶ παιδίσκη ὅσοι ἂν γένωνταί σοι, ἀπὸ τῶν ἐθνῶν ὅσοι κύκλῳ σου εἰσὶν, ἀπʼ αὐτῶν κτήσεσθε δοῦλον καὶ δούλην,
\vs{45}καὶ ἀπὸ τῶν υἱῶν τῶν παροίκων τῶν ὄντων ἐν ὑμῖν, ἀπὸ τούτων κτήσεσθε καὶ ἀπὸ τῶν συγγενῶν αὐτῶν, ὅσοι ἂν γένωνται ἐν τῇ γῇ ὑμῶν, ἔστωσαν ὑμῖν εἰς κατάσχεσιν.
\vs{46}Καὶ καταμεριεῖτε αὐτοὺς τοῖς τέκνοις ὑμῶν μεθʼ ὑμᾶς· καὶ ἔσονται ὑμῖν κατόχιμοι εἰς τὸν αἰῶνα· τῶν δὲ ἀδελφῶν ὑμῶν τῶν υἱῶν Ἰσραὴλ, ἕκαστος τὸν ἀδελφὸν αὐτοῦ οὐ κατατενεῖ αὐτὸν ἐν τοῖς μόχθοις.

\vs{47}Ἐὰν δὲ εὕρῃ ἡ χεὶρ τοὺ προσηλύτου ἢ τοῦ παροίκου τοῦ παρὰ σοὶ, καὶ ἀπορηθεὶς ὁ ἀδελφός σου πραθῇ τῷ προσηλύτῳ ἢ τῷ παροίκῳ τῷ παρὰ σοὶ, ἢ ἐκ γενετῆς προσηλύτῳ,
\vs{48}μετὰ τὸ πραθῆναι αὐτῷ, λύτρωσις ἔσται αὐτοῦ· εἷς τῶν ἀδελφῶν αὐτοῦ λυτρώσεται αὐτόν.
\vs{49}Ἀδελφὸς πατρὸς αὐτοῦ, ἢ υἱὸς ἀδελφοῦ πατρὸς λυτρώσεται αὐτόν, ἢ ἀπὸ τῶν οἰκείων τῶν σαρκῶν αὐτοῦ ἐκ τῆς φυλῆς αὐτοῦ λυτρῶται αὐτόν· ἐὰν δὲ εὐπορηθεὶς ταῖς χερσὶ λυτρῶται ἑαυτὸν,
\vs{50}καὶ συλλογιεῖται πρὸς τὸν κεκτημένον αὐτὸν ἀπὸ τοῦ ἔτους οὗ ἀπέδοτο ἑαυτὸν αὐτῷ ἕως τοῦ ἐνιαυτοῦ τῆς ἀφέσεως· καὶ ἔσται τὸ ἀργύριον τῆς πράσεως αὐτοῦ ὡς μισθίου· ἔτος ἐξ ἔτους ἔσται μετʼ αὐτοῦ.
\vs{51}Ἐὰν δέ τινι πλεῖον τῶν ἐτῶν ᾖ, πρὸς ταῦτα ἀποδώσει τὰ λύτρα αὐτοῦ ἀπὸ τοῦ ἀργυρίου τῆς πράσεως αὐτοῦ.
\vs{52}Ἐὰν δὲ ὀλίγον καταλειφθῇ ἀπὸ τῶν ἐτῶν εἰς τὸν ἐνιαυτὸν τῆς ἀφέσεως, καὶ συλλογιεῖται αὐτῷ κατὰ τὰ ἔτη αὐτοῦ, καὶ ἀποδώσει τὰ λύτρα αὐτοῦ ὡς μισθωτός·
\vs{53}ἐνιαυτὸς ἐξ ἐνιαυτοῦ ἔσται μετʼ αὐτοῦ· οὐ κατατενεῖς αὐτὸν ἐν τῷ μόχθῳ ἐνώπιόν σου.
\vs{54}Ἐὰν δὲ μὴ λυτρῶται κατὰ ταῦτα, ἐξελεύσεται ἐν τῷ ἔτει τῆς ἀφέσεως αὐτὸς καὶ τὰ παιδία αὐτοῦ μετʼ αὐτοῦ.
\vs{55}Ὅτι ἐμοὶ οἱ υἱοὶ Ἰσραὴλ οἰκέται εἰσὶ, παῖδές μου οὗτοί εἰσιν, οὓς ἐξήγαγον ἐκ γῆς Αἰγύπτου.

\ch{26}
Ἐγὼ Κύριος ὁ Θεὸς ὑμῶν· οὐ ποιήσετε ὑμῖν αὐτοῖς χειροποίητα, οὐδὲ γλυπτὰ, οὐδὲ στήλην ἀναστήσετε ὑμῖν, οὐδὲ λίθον σκοπὸν θήσετε ἐν τῇ γῇ ὑμῶν προσκυνῆσαι αὐτῷ· ἐγώ εἰμι Κύριος ὁ Θεὸς ὑμῶν.
\vs{2}Τὰ σάββατά μου φυλάξεσθε, καὶ ἀπὸ τῶν ἁγίων μου φοβηθήσεσθε· ἐγώ εἰμι Κύριος.
\vs{3}Ἐὰν τοῖς προστάγμασί μου πορεύησθε, καὶ τὰς ἐντολάς μου φυλάσσησθε, καὶ ποιήσητε αὐτὰς,
\vs{4}καὶ δώσω τὸν ὑετὸν ὑμῖν ἐν καιρῷ αὐτοῦ, καὶ ἡ γῆ δώσει τὰ γεννήματα αὐτῆς, καὶ τὰ ξύλα τῶν πεδίων ἀποδώσει τὸν καρπὸν αὐτῶν·
\vs{5}Καὶ καταλήψεται ὑμῖν ὁ ἁλοητὸς τὸν τρυγητὸν, καὶ ὁ τρυγητὸς καταλήψεται τὸν σπόρον· καὶ φάγεσθε τὸν ἄρτον ὑμῶν εἰς πλησμονήν· καὶ κατοικήσετε μετὰ ἀσφαλείας ἐπὶ τῆς γῆς ὑμῶν, καὶ πόλεμος οὐ διελεύσεται διὰ τῆς γῆς ὑμῶν·
\vs{6}Καὶ δώσω εἰρήνην ἐν τῇ γῇ ὑμῶν· καὶ κοιμηθήσεσθε, καὶ οὐκ ἔσται ὑμᾶς ὁ ἐκφοβῶν· καὶ ἀπολῶ θηρία πονηρὰ ἐκ τῆς γῆς ὑμῶν.
\vs{7}Καὶ διώξεσθε τοὺς ἐχθροὺς ὑμῶν, καὶ πεσοῦνται ἐναντίον ὑμῶν φόνῳ.
\vs{8}Καὶ διώξονται ἐξ ὑμῶν πέντε ἑκατὸν, καὶ ἑκατὸν ὑμῶν διώξονται μυριάδας· καὶ πεσοῦνται οἱ ἐχθροὶ ὑμῶν ἐναντίον ὑμῶν μαχαίρᾳ.
\vs{9}Καὶ ἐπιβλέψω ἐφʼ ὑμᾶς, καὶ αὐξανῶ ὑμᾶς, καὶ πληθυνῶ ὑμᾶς, καὶ στήσω τὴν διαθήκην μου μεθʼ ὑμῶν·
\vs{10}Καὶ φάγεσθε παλαιὰ καὶ παλαιὰ παλαιῶν, καὶ παλαιὰ ἐκ προσώπου νέων ἐξοίσετε.
\vs{11}Καὶ θήσω τὴν σκηνήν μου ἐν ὑμῖν, καὶ οὐ βδελύξεται ἡ ψυχή μου ὑμᾶς,
\vs{12}καὶ ἐμπεριπατήσω ἐν ὑμῖν· καὶ ἔσομαι ὑμῶν Θεὸς, καὶ ὑμεῖς ἔσεσθέ μοι λαός.
\vs{13}Ἐγώ εἰμι Κύριος ὁ Θεὸς ὑμῶν, ὁ ἐξαγαγὼν ὑμᾶς ἐκ γῆς Αἰγύπτου, ὄντων ὑμῶν δούλων· καὶ συνέτριψα τὸν δεσμὸν τοῦ ζυγοῦ ὑμῶν, καὶ ἤγαγον ὑμᾶς μετὰ παῤῥησίας.

\vs{14}Ἐὰν δὲ μὴ ὑπακούσητέ μου, μηδὲ ποιήσητε τὰ προστάγματά μου ταῦτα,
\vs{15}ἀλλὰ ἀπειθήσητε αὐτοῖς, καὶ τοῖς κρίμασί μου προσοχθίσῃ ἡ ψυχὴ ὑμῶν, ὥστε ὑμᾶς μὴ ποιεῖν πάσας τὰς ἐντολάς μου, ὥστε διασκεδάσαι τὴν διαθήκην μου,
\vs{16}καὶ ἐγὼ ποιήσω οὕτως ὑμῖν· καὶ ἐπιστήσω ἐφʼ ὑμᾶς τὴν ἀπορίαν, τήν τε ψώραν, καὶ τὸν ἴκτερα σφακελίζοντα τοὺς ὀφθαλμοὺς ὑμῶν, καὶ τὴν ψυχὴν ὑμῶν ἐκτήκουσαν· καὶ σπερεῖτε διακενῆς τὰ σπέρματα ὑμῶν, καὶ ἔδονται οἱ ὑπεναντίοι ὑμῶν.
\vs{17}Καὶ ἐπιστήσω τὸ πρόσωπόν μου ἐφʼ ὑμᾶς, καὶ πεσεῖσθε ἐναντίον τῶν ἐχθρῶν ὑμῶν, καὶ διώξονται ὑμᾶς οἱ μισοῦντες ὑμᾶς, καὶ φεύξεσθε οὐδενὸς διώκοντος ὑμᾶς.
\vs{18}Καὶ ἐὰν ἕως τούτου μὴ ὑπακούσητέ μου, καὶ προσθήσω τοῦ παιδεῦσαι ὑμᾶς ἑπτάκις ἐπὶ ταῖς ἁμαρτίαις ὑμῶν.
\vs{19}Καὶ συντρίψω τὴν ὕβριν τῆς ὑπερηφανίας ὑμῶν· καὶ θήσω τὸν οὐρανὸν ὑμῖν σιδηροῦν, καὶ τὴν γῆν ὑμῶν ὡσεὶ χαλκῆν.
\vs{20}Καὶ ἔσται εἰς κενὸν ἡ ἰσχὺς ὑμῶν· καὶ οὐ δώσει ἡ γῆ ὑμῶν τὸν σπόρον αὐτῆς, καὶ τὸ ξύλον τοῦ ἀγρου ὑμῶν οὐ δώσει τὸν καρπὸν αὐτοῦ.

\vs{21}Καὶ ἐὰν μετὰ ταῦτα πορεύησθε πλάγιοι, καὶ μὴ βούλησθε ὑπακούειν μου, προσθήσω ὑμῖν πληγὰς ἑπτὰ κατὰ τὰς ἁμαρτίας ὑμῶν.
\vs{22}Καὶ ἀποστέλλω ἐφʼ ὑμᾶς τὰ θηρία τὰ ἄγρια τῆς γῆς, καὶ κατέδεται ὑμᾶς, καὶ ἐξαναλώσει τὰ κτήνη ὑμῶν, καὶ ὀλιγοστοὺς ποιήσω ὑμᾶς, καὶ ἐρημωθήσονται αἱ ὁδοὶ ὑμῶν.
\vs{23}Καὶ ἐπὶ τούτοις ἐὰν μὴ παιδευθῆτε, ἀλλὰ πορεύησθε πρός με πλάγιοι,
\vs{24}πορεύσομαι κᾀγὼ μεθʼ ὑμῶν θυμῷ πλαγίῳ, καὶ πατάξω ὑμᾶς κᾀγὼ ἑπτάκις ἀντὶ τῶν ἁμαρτιῶν ὑμῶν.
\vs{25}Καὶ ἐπάξω ἐφʼ ὑμᾶς μάχαιραν ἐκδικοῦσαν δίκην διαθήκης, καὶ καταφεύξεσθε εἰς τὰς πόλεις ὑμῶν· καὶ ἐξαποστελῶ θάνατον εἰς ὑμᾶς, καὶ παραδοθήσεσθε εἰς χεῖρας τῶν ἐχθρῶν.
\vs{26}Ἐν τῷ θλίψαι ὑμᾶς σιτοδείᾳ ἄρτων, καὶ πέψουσι δέκα γυναῖκες τοὺς ἄρτους ὑμῶν ἐν κλιβάνῳ ἑνὶ, καὶ ἀποδώσουσι τοὺς ἄρτους ὑμῶν ἐν σταθμῷ, καὶ φάγεσθε, καὶ οὐ μὴ ἐμπλησθῆτε.

\vs{27}Ἐὰν δὲ ἐπὶ τούτοις μὴ ὑπακούσητέ μου, καὶ πορεύησθε πρός με πλάγιοι,
\vs{28}καὶ αὐτὸς πορεύσομαι μεθʼ ὑμῶν ἐν θυμῷ πλαγίῳ, καὶ παιδεύσω ὑμᾶς ἐγὼ ἑπτάκις κατὰ τὰς ἁμαρτίας ὑμῶν.
\vs{29}Καὶ φάγεσθε τὰς σάρκας τῶν υἱῶν ὑμῶν, καὶ τὰς σάρκας τῶν θυγατέρων ὑμῶν φάγεσθε.
\vs{30}Καὶ ἐρημώσω τὰς στήλας ὑμῶν, καὶ ἐξολοθρεύσω τὰ ξύλινα χειροποίητα ὑμῶν, καὶ θήσω τὰ κῶλα ὑμῶν ἐπὶ τὰ κῶλα τῶν εἰδώλων ὑμῶν, καὶ προσοχθιεῖ ἡ ψυχή μου ὑμῖν.
\vs{31}Καὶ θήσω τὰς πόλεις ὑμῶν ἐρήμους, καὶ ἐξερημώσω τὰ ἅγια ὑμῶν, καὶ οὐ μὴ ὀσφρανθῶ τῆς ὀσμῆς τῶν θυσιῶν ὑμῶν.
\vs{32}Καὶ ἐξερημώσω ἐγὼ τὴν γῆν ὑμῶν, καὶ θαυμάσονται ἐπʼ αὐτῇ οἱ ἐχθροὶ ὑμῶν, οἱ ἐνοικοῦντες ἐν αὐτῇ.
\vs{33}Καὶ διασπερῶ ὑμᾶς εἰς τὰ ἔθνη, καὶ ἐξαναλώσει ὑμᾶς ἐπιπορευομένη ἡ μάχαιρα, καὶ ἔσται ἡ γῆ ὑμῶν ἔρημος, καὶ αἱ πόλεις ὑμῶν ἔσονται ἔρημοι.
\vs{34}Τότε εὐδοκήσει ἡ γῆ τὰ σάββατα αὐτῆς πάσας τὰς ἡμέρας τῆς ἐρημώσεως αὐτῆς,
\vs{35}καὶ ὑμεῖς ἔσεσθε ἐν τῇ γῇ τῶν ἐχθρῶν ὑμῶν· τότε σαββατιεῖ ἡ γη, καὶ εὐδοκήσει ἡ γῆ τὰ σάββατα αὐτῆς πάσας τὰς ἡμέρας τῆς ἐρημώσεως αὐτῆς· σαββατιεῖ ἃ οὐκ ἐσαββάτισεν ἐν τοῖς σαββάτοις ὑμῶν, ἡνίκα κατῳκεῖτε αὐτήν.
\vs{36}Καὶ τοῖς καταλειφθεῖσιν ἐξ ὑμῶν ἐπάξω δουλείαν εἰς τὴν καρδίαν αὐτῶν ἐν τῇ γῇ τῶν ἐχθρῶν αὐτῶν· καὶ διώξεται αὐτοὺς φωνὴ φύλλου φερομένου, καὶ φεύξονται ὡς φεύγοντες ἀπὸ πολέμου, καὶ πεσοῦνται οὐθενὸς διώκοντος.
\vs{37}Καὶ ὑπερόψεται ὁ ἀδελφὸς τὸν ἀδελφὸν ὡσεὶ ἐν πολέμῳ, οὐθενὸς κατατρέχοντος· καὶ οὐ δυνήσεσθε ἀντιστῆναι τοῖς ἐχθροῖς ὑμῶν.
\vs{38}Καὶ ἀπολεῖσθε ἐν τοῖς ἔθνεσι, καὶ κατέδεται ὑμᾶς ἡ γῆ τῶν ἐχθρῶν ὑμῶν.
\vs{39}Καὶ οἱ καταλειφθέντες ἀφʼ ὑμῶν, καταφθαρήσονται διὰ τὰς ἁμαρτίας αὐτῶν, καὶ διὰ τὰς ἁμαρτίας τῶν πατέρων αὐτῶν· ἐν τῇ γῇ τῶν ἐχθρῶν αὐτῶν τακήσονται.

\vs{40}Καὶ ἐξαγορεύσουσι τὰς ἁμαρτίας αὐτῶν, καὶ τὰς ἁμαρτίας τῶν πατέρων αὐτῶν, ὅτι παρέβησαν καὶ ὑπερεῖδόν με, καὶ ὅτι ἐπορεύθησαν ἐναντίον μου πλάγιοι,
\vs{41}καὶ ἐγὼ ἐπορεύθην μετʼ αὐτῶν ἐν θυμῷ πλαγίῳ· καὶ ἀπολῶ αὐτοὺς ἐν τῇ γῇ τῶν ἐχθρῶν αὐτῶν· τότε ἐντραπήσεται ἡ καρδία αὐτῶν ἡ ἀπερίτμητος, καὶ τότε εὐδοκήσουσι τὰς ἁμαρτίας αὐτῶν.
\vs{42}Καὶ μνησθήσομαι τῆς διαθήκης Ἰακὼβ, καὶ τῆς διαθήκης Ἰσαὰκ, καὶ τῆς διαθήκης Ἁβραὰμ μνησθήσομαι.

\vs{43}Καὶ τῆς γῆς μνησθήσομαι, καὶ ἡ γῆ ἐγκαταλειφθήσεται ἀπʼ αὐτῶν· τότε προσδέξεται ἡ γῆ τὰ σάββατα αὐτῆς, ἐν τῷ ἐρημωθῆναι αὐτὴν διʼ αὐτούς· καὶ αὐτοὶ προσδέξονται τὰς αὐτῶν ἀνομίας, ἀνθʼ ὧν τὰ κρίματά μου ὑπερεῖδον, καὶ τοῖς προστάγμασί μου προσώχθισαν τῇ ψυχῇ αὐτῶν.
\vs{44}Καὶ οὐδʼ ὡς ὄντων αὐτῶν ἐν τῇ γῇ τῶν ἐχθρῶν αὐτῶν, οὐχ ὑπερεῖδον αὐτούς, οὐδὲ προσώχθισα αὐτοῖς ὥστε ἐξαναλῶσαι αὐτούς τοῦ διασκεδάσαι τὴν διαθήκην μου τὴν πρὸς αὐτούς· ἐγὼ γάρ εἰμι Κύριος ὁ Θεὸς αὐτῶν.
\vs{45}Καὶ μνησθήσομαι διαθήκης αὐτῶν τῆς προτέρας, ὅτε ἐξήγαγον αὐτοὺς ἐκ γῆς Αἰγύπτου, ἐξ οἴκου δουλείας ἔναντι τῶν ἐθνῶν, τοῦ εἶναι αὐτῶν Θεός· ἐγώ εἰμι Κύριος.
\vs{46}Ταῦτα τὰ κρίματά μου, καὶ τὰ προστάγματά μου, καὶ ὁ νόμος ὃν ἔδωκε Κύριος ἀναμέσον αὐτοῦ καὶ ἀναμέσον τῶν υἱῶν Ἰσραὴλ, ἐν τῷ ὄρει Σινᾷ ἐν χειρὶ Μωυσῆ.

\ch{27}
Καὶ ἐλάλησε Κύριος πρὸς Μωυσῆν, λέγων,
\vs{2}λάλησον τοῖς υἱοῖς Ἰσραὴλ, καὶ ἐρεῖς αὐτοῖς, ὃς ἂν εὔξηται εὐχὴν ὥστε τιμὴν τῆς ψυχῆς αὐτοῦ τῷ Κυρίῳ,
\vs{3}ἔσται ἡ τιμὴ τοῦ ἄρσενος ἀπὸ εἰκοσαετοῦς, ἕως ἑξηκονταετοῦς, ἔσται αὐτοῦ ἡ τιμὴ πεντήκοντα δίδραχμα ἀργυρίου τῷ σταθμῷ τῷ ἁγίῳ·
\vs{4}Τῆς δὲ θηλείας ἔσται ἡ συντίμησις τριάκοντα δίδραχμα.
\vs{5}Ἐὰν δὲ ἀπὸ πενταετοῦς ἕως εἴκοσι ἐτῶν, ἔσται ἡ τιμὴ τοῦ ἄρσενος εἴκοσι δίδραχμα· τῆς δὲ θηλείας, δέκα δίδραχμα.
\vs{6}Ἀπὸ δὲ μηνιαίου ἕως πενταετοῦς, ἔσται ἡ τιμὴ τοῦ ἄρσενος πέντε δίδραχμα· τῆς δὲ θηλείας, τρία δίδραχμα ἀργυρίου.
\vs{7}Ἐὰν δὲ ἀπὸ ἑξήκοντα ἐτῶν καὶ ἐπάνω, ἐὰν μὲν ἄρσεν ᾖ, ἔσται ἡ τιμὴ αὐτοῦ πεντεκαίδεκα δίδραχμα ἀργυρίου· ἐὰν δὲ θήλεια, δέκα δίδραχμα.
\vs{8}Ἐὰν δὲ ταπεινὸς ᾖ τῇ τιμῇ, στήσεται ἐναντίον τοῦ ἱερέως· καὶ τιμήσεται αὐτὸν ὁ ἱερεύς· καθάπερ ἰσχύει ἡ χεὶρ τοῦ εὐξαμένου, τιμήσεται αὐτὸν ὁ ἱερεύς.

\vs{9}Ἐὰν δὲ ἀπὸ τῶν κτηνῶν τῶν προσφερομένων ἀπʼ αὐτῶν δῶρον τῷ Κυρίῳ, ὃς ἂν δῷ ἀπὸ τούτων τῷ Κυρίῳ, ἔσται ἅγιον.
\vs{10}Οὐκ ἀλλάξει αὐτὸ καλὸν πονηρῷ, οὐδὲ πονηρὸν καλῷ· ἐὰν δὲ ἀλλάσσων ἀλλάξῃ αὐτὸ κτῆνος κτήνει, ἔσται αὐτὸ καὶ τὸ ἄλλαγμα ἅγια.
\vs{11}Ἐὰν δὲ πᾶν κτῆνος ἀκάθαρτον, ἀφʼ ὧν οὐ προσφέρεται ἀπʼ αὐτῶν δῶρον τῷ Κυρίῳ, στήσει τὸ κτῆνος ἔναντι τοῦ ἱερέως,
\vs{12}καὶ τιμήσεται αὐτὸ ὁ ἱερεὺς ἀναμέσον καλοῦ καὶ ἀναμέσον πονηροῦ· καὶ καθότι ἂν τιμήσηται αὐτὸ ὁ ἱερεὺς, οὕτω στήσεται.
\vs{13}Ἐὰν δὲ λυτρούμενος λυτρώσηται αὐτὸ, προσθήσει τὸ ἐπίπεμπτον πρὸς τὴν τιμὴν αὐτοῦ.
\vs{14}Καὶ ἄνθρωπος ὃς ἂν ἁγιάσῃ τὴν οἰκίαν αὐτοῦ ἁγίαν τῷ Κυρίῳ, καὶ τιμήσεται αὐτὴν ὁ ἱερεὺς ἀναμέσον καλῆς καὶ ἀναμέσον πονηρᾶς· ὡς ἂν τιμήσηται αὐτὴν ὁ ἱερεὺς, οὕτω σταθήσεται.
\vs{15}Ἐὰν δὲ ὁ ἁγιάσας αὐτὴν λυτρῶται τὴν οἰκίαν αὐτοῦ, προσθήσει ἐπʼ αὐτὸ τὸ ἐπίπεμπτον τοῦ ἀργυρίου τῆς τιμῆς, καὶ ἔσται αὐτῷ.

\vs{16}Ἐὰν δὲ ἀπὸ τοῦ ἀγροῦ τῆς κατασχέσεως αὐτοῦ ἁγιάσῃ ἄνθρωπος τῷ Κυρίῳ, καὶ ἔσται ἡ τιμὴ κατὰ τὸν σπόρον αὐτοῦ, κόρου κριθῶν πεντήκοντα δίδραχμα ἀργυρίου.
\vs{17}Ἐὰν δὲ ἀπὸ τοῦ ἐνιαυτοῦ τῆς ἀφέσεως ἁγιάσῃ τὸν ἀγρὸν αὐτοῦ, κατὰ τὴν τιμὴν αὐτοῦ στήσεται.
\vs{18}Ἐὰν δὲ ἔσχατον μετὰ τὴν ἄφεσιν ἁγιάσῃ τὸν ἀγρὸν αὐτοῦ, προσλογιεῖται αὐτῷ ὁ ἱερεὺς τὸ ἀργύριον ἐπὶ τὰ ἔτη τὰ ἐπίλοιπα, ἕως εἰς τὸν ἐνιαυτὸν τῆς ἀφέσεως, καὶ ἀνθυφαιρεθήσεται ἀπὸ τῆς συντιμήσεως αὐτοῦ.
\vs{19}Ἐὰν δὲ λυτρῶται τὸν ἀγρὸν ὁ ἁγιάσας αὐτὸν, προσθήσει τὸ ἐπίπεμπτον τοῦ ἀργυρίου πρὸς τὴν τιμὴν αὐτοῦ, καὶ ἔσται αὐτῷ.
\vs{20}Ἐὰν δὲ μὴ λυτρῶται τὸν ἀγρὸν, καὶ ἀποδῶται τὸν ἀγρὸν ἀνθρώπῳ ἑτέρῳ, οὐκέτι μὴ λυτρώσηται αὐτόν.
\vs{21}Ἀλλʼ ἔσται ὁ ἀγρὸς ἐξεληλυθυίας τῆς ἀφέσεως ἅγιος τῷ Κυρίῳ, ὥσπερ ἡ γῆ ἡ ἀφωρισμένη τῷ ἱερεῖ ἔσται κατάσχεσις αὐτοῦ.
\vs{22}Ἐὰν δὲ ἀπὸ τοῦ ἀγροῦ οὗ κέκτηται, ὃς οὐκ ἔστιν ἀπὸ τοῦ ἀγροῦ τῆς κατασχέσεως αὐτοῦ, ἁγιάσῃ τῷ Κυρίῳ,
\vs{23}λογιεῖται πρὸς αὐτὸν ὁ ἱερεὺς τὸ τέλος τῆς τιμῆς ἐκ τοῦ ἐνιαυτοῦ τῆς ἀφέσεως, καὶ ἀποδώσει τὴν τιμὴν ἐν τῇ ἡμέρᾳ ἐκείνῃ ἁγίαν τῷ Κυρίῳ·
\vs{24}Καὶ ἐν τῷ ἐνιαυτῷ τῆς ἀφέσεως ἀποδοθήσεται ὁ ἀγρὸς τῷ ἀνθρώπῳ παρʼ οὗ κέκτηται αὐτὸν, οὗ ἦν ἡ κατάσχεσις τῆς γῆς.
\vs{25}Καὶ πᾶσα τιμὴ ἔσται σταθμίοις ἁγίοις· εἴκοσι ὀβολοὶ ἔσται τὸ δίδραχμον.
\vs{26}Καὶ πᾶν πρωτότοκον ὃ ἐὰν γένηται ἐν τοῖς κτήνεσί σου, ἔσται τῷ Κυρίῳ, καὶ οὐ καθαγιάσει αὐτὸ οὐδείς· ἐάν τε μόσχον, ἐάν τε πρόβατον, τῷ Κυρίῳ ἐστίν.
\vs{27}Ἐὰν δὲ τῶν τετραπόδων τῶν ἀκαθάρτων ἀλλάξῃ κατὰ τὴν τιμὴν αὐτοῦ, καὶ προσθήσει τὸ ἐπίπεμπτον πρὸς αὐτὸ, καὶ ἔσται αὐτῷ· ἐὰν δὲ μὴ λυτρῶται, πραθήσεται κατὰ τὸ τίμημα αὐτοῦ.

\vs{28}Πᾶν δὲ ἀνάθεμα, ὃ ἂν ἀναθῇ ἄνθρωπος τῷ Κυρίῳ ἀπὸ πάντων, ὅσα αὐτῷ ἐστιν, ἀπὸ ἀνθρώπου ἕως κτήνους, καὶ ἀπὸ ἀγροῦ κατασχέσεως αὐτοῦ, οὐκ ἀποδώσεται οὐδὲ λυτρώσεται· πᾶν ἀνάθεμα ἅγιον ἁγίων ἔσται τῷ Κυρίῳ.
\vs{29}Καὶ πᾶν ὃ ἐὰν ἀνατεθῇ ἀπὸ τῶν ἀνθρώπων, οὐ λυτρωθήσεται, ἀλλὰ θανάτῳ θανατωθήσεται.
\vs{30}Πᾶσα δεκάτη τῆς γῆς, ἀπὸ τοῦ σπέρματος τῆς γῆς, καὶ τοῦ καρποῦ τοῦ ξυλίνου, τῷ. Κυρίῳ ἐστίν, ἅγιον τῷ Κυρίῳ.
\vs{31}Ἐὰν δὲ λυτρῶται λύτρῳ ἄνθρωπος τὴν δεκάτην αὐτοῦ, τὸ ἐπίπεμπτον προσθήσει πρὸς αὐτὸν, καὶ ἔσται αὐτῷ.
\vs{32}Καὶ πᾶσα δεκάτη βοῶν, καὶ προβάτων, καὶ πᾶν ὃ ἂν ἔλθῃ ἐν τῷ ἀριθμῷ ὑπὸ τὴν ῥάβδον, τὸ δέκατον ἔσται ἅγιον τῷ Κυρίῳ.
\vs{33}Οὐκ ἀλλάξεις καλὸν πονηρῷ, οὐδὲ πονηρὸν καλῷ· ἐὰν δὲ ἀλλάσσων ἀλλάξῃς αὐτό, καὶ τὸ ἄλλαγμα αὐτοῦ ἔσται ἅγιον, οὐ λυτρωθήσεται.

\vs{34}Αὗταί εἰσιν αἱ ἐντολαὶ ἃς ἐνετείλατο Κύριος τῷ Μωυσῇ πρὸς τοὺς υἱοὺς Ἰσραὴλ ἐν τῷ ὄρει Σινᾷ.


\def\book{ΑΡΙΘΜΟΙ}
\biblebook{ΑΡΙΘΜΟΙ}


\lettrine[lines=2, loversize=0.2, nindent=0em, findent=.25em]{\textcolor{bookheadingcolor}{Κ}}{ΑΙ} ἐλάλησε Κύριος πρὸς Μωυσῆν ἐν τῇ ἐρήμῳ τῇ Σινᾷ, ἐν τῇ σκηνῇ τοῦ μαρτυρίου, ἐν μιᾷ τοῦ μηνὸς τοῦ δευτέρου, ἔτους δευτέρου ἐξελθόντων αὐτῶν ἐκ γῆς Αἰγύπτου, λέγων,
\vs{2}λάβετε ἀρχὴν πάσης συναγωγῆς Ἰσραὴλ κατὰ συγγενείας, κατʼ οἴκους πατριῶν αὐτῶν, κατὰ ἀριθμὸν ἐξ ὀνόματος αὐτῶν, κατὰ κεφαλὴν αὐτῶν·
\vs{3}πᾶς ἄρσην ἀπὸ εἰκοσαετοῦς καὶ ἐπάνω, πᾶς ὁ ἐκπορευόμενος ἐν δυνάμει Ἰσραήλ, ἐπισκέψασθε αὐτοὺς σὺν δυνάμει αὐτῶν· σὺ καὶ Ἀαρὼν ἐπισκέψασθε αὐτούς.
\vs{4}Καὶ μεθʼ ὑμῶν ἔσονται ἕκαστος κατὰ φυλὴν ἑκάστου ἀρχόντων, κατʼ οἴκους πατριῶν ἔσονται.

\vs{5}Καὶ ταῦτα τὰ ὀνόματα τῶν ἀνδρῶν, οἵτινες παραστήσονται μεθʼ ὑμῶν· τῶν Ῥουβὴν, Ἐλισοὺρ υἱὸς Σεδιούρ·
\vs{6}Τῶν Συμεὼν, Σαλαμιὴλ υἱὸς Σουρισαδαί.
\vs{7}Τῶν Ἰούδα, Ναασσὼν υἱὸς Ἀμιναδάβ.
\vs{8}Τῶν Ἰσσάχαρ, Ναθαναὴλ υἱὸς Σωγάρ·
\vs{9}Τῶν Ζαβουλὼν, Ἐλιὰβ υἱὸς Χαιλών· Τῶν υἱῶν Ἰωσὴφ τῶν Ἐφραὶμ, Ἐλισαμὰ υἱὸς Ἐμιούδ·
\vs{10}τῶν Μανασσῆ, Γαμαλιὴλ υἱὸς Φαδασούρ.
\vs{11}Τῶν Βενιαμὶν, Ἀβιδὰν υἱὸς Γαδεωνί.
\vs{12}Τῶν Δὰν, Ἀχιέζερ υἱὸς Ἀμισαδαΐ.
\vs{13}Τῶν Ἀσὴρ, Φαγαϊὴλ, υἱὸς Ἐχράν.
\vs{14}Τῶν Γὰδ, Ἐλισὰφ υἱὸς Ῥαγουήλ.
\vs{15}Τῶν Νεφθαλὶ, Ἀχιρὲ υἱὸς Αἰνάν.
\vs{16}Οὗτοι ἐπίκλητοι τῆς συναγωγῆς, ἄρχοντες τῶν φυλῶν κατὰ πατριὰς αὐτῶν, χιλίαρχοι Ἰσραήλ εἰσι.

\vs{17}Καὶ ἔλαβε Μωυσῆς καὶ Ἀαρὼν τοὺς ἄνδρας τούτους τοὺς ἀνακληθέντας ἐξ ὀνόματος.
\vs{18}Καὶ πᾶσαν τὴν συναγωγὴν συνήγαγον ἐν μιᾷ τοῦ μηνὸς τοῦ δευτέρου ἔτους· καὶ ἐπηξονοῦσαν κατὰ γενέσεις αὐτῶν, κατὰ πατριὰς αὐτῶν, κατὰ ἀριθμὸν ὀνομάτων αὐτῶν, ἀπὸ εἰκοσαετοῦς καὶ ἐπάνω, πᾶν ἀρσενικὸν κατὰ κεφαλὴν αὐτῶν,
\vs{19}ὃν τρόπον συνέταξε Κύριος τῷ Μωυσῇ· καὶ ἐπεσκέπησαν ἐν τῇ ἐρήμῳ τοῦ Σινά.

\vs{20}Καὶ ἐγένοντο οἱ υἱοὶ Ῥουβὴν πρωτοτόκου Ἰσραὴλ κατὰ συγγενείας αὐτῶν, κατὰ δήμους αὐτῶν, κατʼ οἴκους πατριῶν αὐτῶν, κατὰ ἀριθμὸν ὀνομάτων αὐτῶν, κατὰ κεφαλὴν αὐτῶν, πάντα ἀρσενικὰ ἀπὸ εἰκοσαετοῦς καὶ ἐπάνω, πᾶς ὁ ἐκπορευόμενος ἐν τῇ δυνάμει,
\vs{21}ἡ ἐπίσκεψις αὐτῶν ἐκ τῆς φυλῆς Ῥουβὴν, ἓξ καὶ τεσσαράκοντα χιλιάδες καὶ πεντακόσιοι.
\vs{22}Τοῖς υἱοῖς Συμεὼν κατὰ συγγενείας αὐτῶν, κατὰ δήμους αὐτῶν, κατʼ οἴκους πατριῶν αὐτῶν, κατὰ ἀριθμὸν ὀνομάτων αὐτῶν, κατὰ κεφαλὴν αὐτῶν, πάντα ἀρσενικὰ ἀπὸ εἰκοσαετοῦς καὶ ἐπάνω, πᾶς ὁ ἐκπορευόμενος ἐν τῇ δυνάμει,
\vs{23}ἡ ἐπίσκεψις αὐτῶν ἐκ τῆς φυλῆς Συμεὼν, ἐννέα καὶ πεντήκοντα χιλιάδες καὶ τριακόσιοι.

\vs{24}Τοῖς υἱοῖς Ἰούδα κατὰ συγγενείας αὐτῶν, κατὰ δήμους αὐτῶν, κατʼ οἴκους πατριῶν αὐτῶν, κατὰ ἀριθμὸν ὀνομάτων αὐτῶν, κατὰ κεφαλὴν αὐτῶν, πάντα ἀρσενικὰ ἀπὸ εἰκοσαετοῦς καὶ ἐπάνω, πᾶς ὁ ἐκπορευόμενος ἐν τῇ δυνάμει,
\vs{25}ἡ ἐπίσκεψις αὐτῶν ἐκ τῆς φυλῆς Ἰούδα, τέσσαρες καὶ ἑβδομήκοντα χιλιάδες καὶ ἑξακόσιοι.

\vs{26}Τοῖς υἱοῖς Ἰσσάχαρ κατὰ συγγενίας αὐτῶν, κατὰ δήμους αὐτῶν, κατʼ οἴκους πατριῶν αὐτῶν, κατὰ ἀριθμὸν ὀνομάτων αὐτῶν, κατὰ κεφαλὴν αὐτῶν, πάντα ἀρσενικὰ ἀπὸ εἰκοσαετοῦς καὶ ἐπάνω, πᾶς ὁ ἐκπορευόμενος ἐν τῇ δυνάμει,
\vs{27}ἡ ἐπίσκεψις αὐτῶν ἐκ τῆς φυλῆς Ἰσσάχαρ, τέσσαρες καὶ πεντήκοντα χιλιάδες καὶ τετρακόσιοι.
\vs{28}Τοῖς υἱοῖς Ζαβουλὼν κατὰ συγγενείας αὐτῶν, κατὰ δήμους αὐτῶν, κατʼ οἴκους πατριῶν αὐτῶν, κατὰ ἀριθμὸν ὀνομάτων αὐτῶν, κατὰ κεφαλὴν αὐτῶν, πάντα ἀρσενικὰ ἀπὸ εἰκοσαετοῦς καὶ ἐπάνω, πᾶς ὁ ἐκπορευόμενος ἐν τῇ δυνάμει,
\vs{29}ἡ ἐπίσκεψις αὐτῶν ἐκ τῆς φυλῆς Ζαβουλὼν, ἑπτὰ καὶ πεντήκοντα χιλιάδες καὶ τετρακόσιοι.

\vs{30}Τοῖς υἱοῖς Ἰωσὴφ υἱοῖς Ἐφραὶμ κατὰ συγγενείας αὐτῶν, κατὰ δήμους αὐτῶν, κατʼ οἴκους πατριῶν αὐτῶν, κατὰ ἀριθμὸν ὀνομάτων αὐτῶν, κατὰ κεφαλὴν αὐτῶν, πάντα ἀρσενικὰ ἀπὸ εἰκοσαετοῦς καὶ ἐπάνω, πᾶς ὁ ἐκπορευόμενος ἐν τῇ δυνάμει,
\vs{31}ἡ ἐπίσκέψις αὐτῶν ἐκ τῆς φυλῆς Ἐφραὶμ, τεσσαράκοντα χιλιάδες καὶ πεντακόσιοι.
\vs{32}Τοῖς υἱοῖς Μανασσὴ κατὰ συγγενείας αὐτῶν, κατὰ δήμους αὐτῶν, κατʼ οἴκους πατριῶν αὐτῶν, κατὰ ἀριθμὸν ὀνομάτων αὐτῶν, κατὰ κεφαλὴν αὐτῶν, πάντα ἀρσενικὰ, ἀπὸ εἰκοσαετοῦς καὶ ἐπάνω, πᾶς ὁ ἐκπορευόμενος ἐν τῇ δυνάμει,
\vs{33}ἡ ἐπίσκεψις αὐτῶν ἐκ τῆς φυλῆς Μανασσὴ, δύο καὶ τριάκοντα χιλιάδες καὶ διακόσιοι.
\vs{34}Τοῖς υἱοῖς Βενιαμὶν κατὰ συγγενείας αὐτῶν, κατὰ δήμους αὐτῶν, κατʼ οἴκους πατριῶν αὐτῶν, κατὰ ἀριθμὸν ὀνομάτων αὐτῶν, κατὰ κεφαλὴν αὐτῶν, πάντα ἀρσενικὰ ἀπὸ εἰκοσαετοῦς καὶ ἐπάνω, πᾶς ὁ ἐκπορευόμενος ἐν τῇ δυνάμει,
\vs{35}ἡ ἐπίσκεψις αὐτῶν ἐκ τῆς φυλῆς Βενιαμὶν, πέντε καὶ τριάκοντα χιλιάδες καὶ τετρακόσιοι.
\vs{36}Τοῖς υἱοῖς Γὰδ κατὰ συγγενείας αὐτῶν, κατὰ δήμους αὐτῶν, κατʼ οἴκους πατριῶν αὐτῶν, κατὰ ἀριθμὸν ὀνομάτων αὐτῶν, κατὰ κεφαλὴν αὐτῶν, πάντα ἀρσενικὰ ἀπὸ εἰκοσαετοῦς καὶ ἐπάνω, πᾶς ὁ ἐκπορευόμενος ἐν τῇ δυνάμει,
\vs{37}ἡ ἐπίσκεψις αὐτῶν· ἐκ τῆς φυλῆς Γὰδ, πέντε καὶ τεσσαράκοντα χιλιάδες καὶ ἑξακόσιοι καὶ πεντήκοντα.

\vs{38}Τοῖς υἱοῖς Δὰν κατὰ συγγενείας αὐτῶν, κατὰ δήμους αὐτῶν, κατʼ οἴκους πατριῶν αὐτῶν, κατὰ ἀριθμὸν ὀνομάτων αὐτῶν, κατὰ κεφαλὴν αὐτῶν, πάντα ἀρσενικὰ ἀπὸ εἰκοσαετοῦς καὶ ἐπάνω, πᾶς ὁ ἐκπορευόμενος ἐν τῇ δυνάμει,
\vs{39}ἡ ἐπίσκεψις αὐτῶν ἐκ τῆς φυλῆς Δὰν, δύο καὶ ἑξήκοντα χιλιάδες καὶ ἑπτακόσιοι.
\vs{40}Τοῖς υἱοῖς Ἀσὴρ κατὰ συγγενείας αὐτῶν, κατὰ δήμους αὐτῶν, κατʼ οἴκους πατριῶν αὐτῶν, κατὰ ἀριθμὸν ὀνομάτων αὐτῶν, κατὰ κεφαλὴν αὐτῶν, πάντα ἀρσενικὰ ἀπὸ εἰκοσαετοῦς καὶ ἐπάνω, πᾶς ὁ ἐκπορευόμενος ἐν τῇ δυνάμει,
\vs{41}ἡ ἐπίσκεψις αὐτῶν ἐκ τῆς φυλῆς Ἀσὴρ, μία καὶ τεσσαράκοντα χιλιάδες καὶ πεντακόσιοι.

\vs{42}Τοῖς υἱοῖς Νεφθαλὶ κατὰ συγγενείας αὐτῶν, κατὰ δήμους αὐτῶν, κατʼ οἴκους πατριῶν αὐτῶν, κατὰ ἀριθμὸν ὀνομάτων αὐτῶν, κατὰ κεφαλὴν αὐτῶν, πάντα ἀρσενικὰ ἀπὸ εἰκοσαετοῦς καὶ ἐπάνω, πᾶς ὁ ἐκπορευόμενος ἐν τῇ δυνάμει,
\vs{43}ἡ ἐπίσκεψις αὐτῶν ἐκ τῆς φυλῆς Νεφθαλὶ, τρεῖς καὶ πεντήκοντα χιλιάδες καὶ τετρακόσιοι.

\vs{44}Αὕτη ἡ ἐπίσκεψις, ἣν ἐπεσκέψαντο Μωυσῆς καὶ Ἀαρὼν καὶ οἱ ἄρχοντες Ἰσραὴλ δώδεκα ἄνδρες· ἀνὴρ εἷς κατὰ φυλὴν μίαν, κατὰ φυλὴν οἴκων πατριᾶς ἦσαν.
\vs{45}Καὶ ἐγένετο πᾶσα ἡ ἐπίσκεψις υἱῶν Ἰσραὴλ σὺν δυνάμει αὐτῶν ἀπὸ εἰκοσαετοῦς καὶ ἐπάνω, πᾶς ὁ ἐκπορευόμενος παρατάξασθαι ἐν Ἰσραὴλ
\vs{46}ἑξακόσιαι χιλιάδες καὶ τρισχίλιοι καὶ πεντακόσιοι καὶ πεντήκοντα.

\vs{47}Οἱ δὲ Λευῖται ἐκ τῆς φυλῆς πατριᾶς αὐτῶν οὐκ ἐπεσκέπησαν ἐν τοῖς υἱοῖς Ἰσραήλ.
\vs{48}καὶ ἐλάλησε Κύριος πρὸς Μωυσῆν, λέγων,
\vs{49}ὅρα, τὴν φυλὴν Λευὶ οὐ συνεπισκέψῃ, καὶ τὸν ἀριθμὸν αὐτῶν οὐ λήμῃ, ἐν μέσῳ υἱῶν Ἰσραήλ.
\vs{50}Καὶ σὺ ἐπίστησον τοὺς Λευίτας ἐπὶ τὴν σκηνὴν τοῦ μαρτυρίου, καὶ ἐπὶ πάντα τὰ σκεύη αὐτῆς, καὶ ἐπὶ πάντα ὅσα ἐστὶν ἐν αὐτῇ· ἀροῦσιν αὐτοὶ τὴν σκηνὴν, καὶ πάντα τὰ σκεύη αὐτῆς· καὶ αὐτοὶ λειτουργήσουσιν ἐν αὐτῇ, καὶ κύκλῳ τῆς σκηνῆς παρεμβαλοῦσι.
\vs{51}Καὶ ἐν τῷ ἐξαίρειν τὴν σκηνὴν, καθελοῦσιν αὐτὴν οἱ Λευῖται, καὶ ἐν τῷ παρεμβάλλειν τὴν σκηνὴν, ἀναστήσουσι· καὶ ὁ ἀλλογενῆς ὁ προσπορευόμενος ἀποθανέτω.
\vs{52}Καὶ παρεμβαλοῦσιν οἱ υἱοὶ Ἰσραὴλ, ἀνὴρ ἐν τῇ ἑαυτοῦ τάξει, καὶ ἀνὴρ κατὰ τὴν ἑαυτοῦ ἡγεμονίαν, σὺν δυνάμει αὐτῶν.
\vs{53}Οἱ δὲ Λευῖται παρεμβαλλέτωσαν ἐναντίοι κύκλῳ τῆς σκηνῆς τοῦ μαρτυρίου, καὶ οὐκ ἔσται ἁμάρτημα ἐν υἱοῖς Ἰσραήλ. καὶ φυλάξουσιν οἱ Λευῖται αὐτοὶ τὴν φυλακὴν τῆς σκηνῆς τοῦ μαρτυρίου.
\vs{54}Καὶ ἐποίησαν οἱ υἱοὶ Ἰσραὴλ, κατὰ πάντα ἃ ἐνετείλατο Κύριος τῷ Μωυσῇ καὶ Ἀαρὼν, οὕτως ἐποίησαν.

\ch{2}
Καὶ ἐλάλησε Κύριος πρὸς Μωυσῆν καὶ Ἀαρὼν, λέγων,
\vs{2}ἄνθρωπος ἐχόμενος αὐτοῦ κατὰ τάγμα, κατὰ σημαίας, κατʼ οἴκους πατριῶν αὐτῶν, παρεμβαλλέτωσαν οἱ υἱοὶ Ἰσραήλ ἐναντίοι· κύκλῳ τῆς σκηνῆς τοῦ μαρτυρίου παρεμβαλοῦσιν οἱ υἱοὶ Ἰσραήλ.
\vs{3}Καὶ οἱ παρεμβάλλοντες πρῶτοι κατὰ ἀνατολὰς, τάγμα παρεμβολῆς Ἰούδα σὺν δυνάμει αὐτῶν, καὶ ὁ ἄρχων τῶν υἱῶν Ἰούδα, Ναασσὼν υἱὸς Ἀμιναδάβ.
\vs{4}Δύναμις αὐτοῦ οἱ ἐπεσκεμμένοι, τέσσαρες καὶ ἑβδομήκοντα χιλιάδες καὶ ἑξακόσιοι.
\vs{5}Καὶ οἱ παρεμβάλλοντες ἐχόμενοι φυλῆς Ἰσσάχαρ, καὶ ὁ ἄρχων τῶν υἱῶν Ἰσσάχαρ, Ναθαναὴλ υἱὸς Σωγάρ.
\vs{6}Δύναμις αὐτοῦ οἱ ἐπεσκεμμένοι, τέσσαρες καὶ πεντήκοντα χιλιάδες καὶ τετρακόσιοι.
\vs{7}Καὶ οἱ παρεμβάλλοντες ἐχόμενοι φυλῆς Ζαβουλὼν, καὶ ὁ ἄρχων τῶν υἱῶν Ζαβουλὼν, Ἑλιὰβ υἱὸς Χαιλών.
\vs{8}Δύναμις αὐτοῦ οἱ ἐπεσκεμμένοι, ἑπτὰ καὶ πεντήκοντα χιλιάδες καὶ τετρακόσιοι.
\vs{9}Πάντες οἱ ἐπεσκεμμένοι ἐκ τῆς παρεμβολῆς Ἰούδα, ἑκατὸν ὀγδοήκοντα χιλιάδες καὶ ἑξακισχίλιοι καὶ τετρακόσιοι, σὺν δυνάμει αὐτῶν πρῶτον ἐξαροῦσι.
\vs{10}Τάγματα παρεμβολῆς Ῥουβὴν, πρὸς λίβα δύναμις αὐτῶν, καὶ ὁ ἄρχων τῶν υἱῶν ʼΡουβὴν, Ἐλισοὺρ υἱὸς Σεδιούρ.
\vs{11}Δύναμις αὐτοῦ οἱ ἐπεσκεμμένοι, ἓξ καὶ τεσσαράκοντα χιλιάδες καὶ πεντακόσιοι.
\vs{12}Καὶ οἱ παρεμβάλλοντες ἐχόμενοι αὐτοῦ φυλῆς Συμεὼν, καὶ ὁ ἄρχων τῶν υἱῶν Συμεὼν, Σαλαμιὴλ υἱὸς Σουρισαδαί.
\vs{13}Δύναμις αὐτοῦ οἱ ἐπεσκεμμένοι, ἐννέα καὶ πεντήκοντα χιλιάδες καὶ τριακόσιοι.
\vs{14}Καὶ οἱ παρεμβάλοντες ἐχόμενοι αὐτοῦ φυλὴ Γὰδ, καὶ ὁ ἄρχων τῶν υἱῶν Γὰδ, Ἑλισὰφ υἱὸς Ῥαγουήλ.
\vs{15}Δύναμις αὐτοῦ οἱ ἐπεσκεμμένοι, πέντε καὶ τεσσαράκοντα χιλιάδες καὶ ἑξακόσιοι καὶ πεντήκοντα.
\vs{16}Πάντες οἱ ἐπεσκεμμένοι τῆς παρεμβολῆς Ῥουβὴν, ἑκατὸν πεντήκοντα μία χιλιάδες καὶ τετρακόσιοι καὶ πεντήκοντα, σὺν δυνάμει αὐτῶν δεύτεροι ἐξαροῦσι.

\vs{17}Καὶ ἀρθήσεται ἡ σκηνὴ τοῦ μαρτυρίου, καὶ ἡ παρεμβολὴ τῶν Λευιτῶν μέσον τῶν παρεμβολῶν· ὡς καὶ παρεμβαλοῦσιν, οὕτω καὶ ἐξαροῦσιν ἕκαστος ἐχόμενος καθʼ ἡγεμονίας.
\vs{18}Τάγμα παρεμβολῆς Ἐφραὶμ παρὰ θάλασσαν σὺν δυνάμει αὐτῶν, καὶ ὁ ἄρχων τῶν υἱῶν Ἐφραὶμ, Ἐλισαμὰ υἱὸς Ἐμιούδ.
\vs{19}Δύναμις αὐτοῦ οἱ ἐπεσκεμμένοι, τεσσαράκοντα χιλιάδες καὶ πεντακόσιοι.

\vs{20}Καὶ οἱ παρεμβάλλοντες ἐχόμενοι φυλῆς Μανασσῆ, καὶ ὁ ἄρχων τῶν υἱῶν Μανασσῆ, Γαμαλιὴλ υἱὸς Φαδασσούρ.
\vs{21}Δύναμις αὐτοῦ οἱ ἐπεσκεμμένοι, δύο καὶ τριάκοντα χιλιάδες καὶ διακόσιοι.
\vs{22}Καὶ οἱ παρεμβάλλοντες ἐχόμενοι φυλῆς Βενιαμὶν, καὶ ὁ ἄρχων τῶν υἱῶν Βενιαμὶν, Ἀβιδὰν υἱὸς Γαδεωνί.
\vs{23}Δύναμις αὐτοῦ οἱ ἐπεσκεμμένοι, πέντε καὶ τριάκοντα χιλιάδες καὶ τετρακόσιοι.
\vs{24}Πάντες οἱ ἐπεσκεμμένοι τῆς παρεμβολῆς Ἐφραὶμ, ἑκατὸν χιλιάδες καὶ ὀκτακισχίλιοι καὶ ἑκατὸν· σὺν δυνάμει αὐτῶν τρίτοι ἐξαροῦσι.

\vs{25}Τάγμα παρεμβολῆς Δὰν πρὸς βοῤῥᾶν σὺν δυνάμει αὐτῶν, καὶ ὁ ἄρχων τῶν υἱῶν Δὰν, Ἀχιέζερ υἱὸς Ἀμισαδαί.
\vs{26}Δύναμις αὐτοῦ οἱ ἐπεσκεμμένοι, δύο καὶ ἑξήκοντα χιλιάδες καὶ ἑπτακόσιοι.
\vs{27}Καὶ οἱ παρεμβάλλοντες ἐχόμενοι αὐτοῦ φυλὴ Ἀσὴρ, καὶ ὁ ἄρχων τῶν υἱῶν Ἀσὴρ, Φαγεὴλ υἱὸς Ἐχράν.
\vs{28}Δύναμις αὐτοῦ οἱ ἐπεσκεμμένοι, μία καὶ τεσσαράκοντα χιλιάδες καὶ πεντακόσιοι.
\vs{29}Καὶ οἱ παρεμβάλλοντες ἐχόμενοι φυλῆς Νεφθαλί, καὶ ὁ ἄρχων τῶν υἱῶν Νεφθαλὶ, Ἀχιρὲ υἱὸς Αἰνάν.
\vs{30}Δύναμις αὐτοῦ οἱ ἐπεσκεμμένοι, τρεῖς καὶ πεντήκοντα χιλιάδες καὶ τετρακόσιοι.
\vs{31}Πάντες οἱ ἐπεσκεμμένοι τῆς παρεμβολῆς Δὰν, ἑκατὸν καὶ πεντηκονταεπτὰ χιλιάδες καὶ ἑξακόσιοι· ἔσχατοι ἐξαροῦσι κατὰ τάγμα αὐτῶν.

\vs{32}Αὕτη ἡ ἐπίσκεψις τῶν υἱῶν Ἰσραὴλ κατʼ οἴκους πατριῶν αὐτῶν· πᾶσα ἡ ἐπίσκεψις τῶν παρεμβολῶν σὺν ταῖς δυνάμεσιν αὐτῶν, ἑξακόσιαι χιλιάδες καὶ τρισχίλιοι πεντακόσιοι πεντήκοντα.
\vs{33}Οἱ δὲ Λευῖται οὐ συνεπεσκέπησαν ἐν αὐτοῖς, καθὰ ἐνετείλατο Κύριος τῷ Μωυσῇ.
\vs{34}Καὶ ἐποίησαν οἱ υἱοὶ Ἰσραὴλ πάντα ὅσα συνέταξε Κύριος τῷ Μωυσῇ· οὕτω παρενέβαλον κατὰ τάγμα αὐτῶν, καὶ οὕτως ἐξῇρον ἕκαστος ἐχόμενοι κατὰ δήμους αὐτῶν, κατʼ οἴκους πατριῶν αὐτῶν.

\ch{3}
Καὶ αὗται αἱ γενέσεις Ἀαρὼν καὶ Μωυσῆ, ἐν ᾗ ἡμέρᾳ ἐλάλησε Κύριος τῷ Μωυσῇ ἐν ὄρει Σινᾷ.
\vs{2}Καὶ ταῦτα τὰ ὀνόματα τῶν υἱῶν Ἀαρών· πρωτότοκος Ναδάβ, καὶ Ἀβιοὺδ, Ἐλεάζαρ, καὶ Ἰθάμαρ.
\vs{3}Ταῦτα τὰ ὀνόματα τῶν υἱῶν Ἀαρὼν, οἱ ἱερεῖς οἱ ἠλειμμένοι, οὓς ἐτελείωσαν τὰς χεῖρας αὐτῶν ἱερατεύειν.
\vs{4}Καὶ ἐτελεύτησε Ναδὰβ καὶ Ἀβιοὺδ ἔναντι Κυρίου, προσφερόντων αὐτῶν πῦρ ἀλλότριον ἔναντι Κυρίου, ἐν τῇ ἐρήμῳ Σινᾷ, καὶ παιδία οὐκ ἦν αὐτοῖς· καὶ ἱεράτευσεν Ἐλεάζαρ καὶ Ἰθάμαρ μετὰ Ἀαρὼν τοῦ πατρὸς αὐτῶν.

\vs{5}Καὶ ἐλάλησε Κύριος πρὸς Μωυσῆν, λέγων,
\vs{6}λάβε τὴν φυλὴν Λευὶ, καὶ στήσεις αὐτοὺς ἐναντίον Ἀαρὼν τοῦ ἱερέως, καὶ λειτουργήσουσιν αὐτῷ,
\vs{7}καὶ φυλάξουσι τὰς φυλακὰς αὐτοῦ, καὶ τὰς φυλακὰς τῶν υἱῶν Ἰσραὴλ ἔναντι τῆς σκηνῆς τοῦ μαρτυρίου, ἐργάζεσθαι τὰ ἔργα τῆς σκηνῆς.
\vs{8}Καὶ φυλάξουσι πάντα τὰ σκεύη τῆς σκηνῆς τοῦ μαρτυρίου, καὶ τὰς φυλακὰς τῶν υἱῶν Ἰσραὴλ κατὰ πάντα τὰ ἔργα τῆς σκηνῆς.
\vs{9}Καὶ δώσεις τοὺς Λευίτας Ἀαρὼν, καὶ τοῖς υἱοῖς αὐτοῦ τοῖς ἱερεῦσι· δεδομένοι δόμα οὗτοί μοι εἰσὶν ἀπὸ τῶν υἱῶν Ἰσραήλ.
\vs{10}Καὶ Ἀαρὼν καὶ τοὺς υἱοὺς αὐτοῦ καταστήσεις ἐπὶ τῆς σκηνῆς τοῦ μαρτυρίου· καὶ φυλάξουσι τὴν ἱερατείαν αὐτῶν, καὶ πάντα τὰ κατὰ τὸν βωμὸν, καὶ ἔσω τοῦ καταπετάσματος· καὶ ὁ ἀλλογενὴς ὁ ἁπτόμενος ἀποθανεῖται.
\vs{11}Καὶ ἐλάλησε Κύριος πρὸς Μωυσῆν, λέγων,
\vs{12}καὶ ἰδοὺ ἐγὼ εἴληφα τοὺς Λευίτας ἐκ μέσου τῶν υἱῶν Ἰσραὴλ ἀντὶ παντὸς πρωτοτόκου διανοίγοντος μήτραν παρὰ τῶν υἱῶν Ἰσραήλ· λύτρα αὐτῶν ἔσονται, καὶ ἔσονται ἐμοὶ οἱ Λευῖται.
\vs{13}Ἐμοὶ γὰρ πᾶν προτότοκον· ἐν ᾗ ἡμέρᾳ ἐπάταξα πᾶν πρωτότοκον ἐν γῇ Αἰγύπτου, ἡγίασα ἐμοὶ πᾶν πρωτότοκον ἐν Ἰσραήλ· ἀπὸ ἀνθρώπου ἕως κτήνους ἐμοὶ ἔσονται· ἐγὼ Κύριος.

\vs{14}Καὶ ἐλάλησε Κύριος πρὸς Μωυσῆν ἐν τῇ ἐρήμῳ Σινᾷ, λέγων,
\vs{15}ἐπίσκεψαι τοὺς υἱοὺς Λευὶ κατʼ οἴκους πατριῶν αὐτῶν, κατὰ δήμους αὐτῶν· πᾶν ἀρσενικὸν ἀπὸ μηνιαίου καὶ ἐπάνω, ἐπισκέψασθε αὐοτύς.
\vs{16}Καὶ ἐπεσκέψαντο αὐτοὺς Μωυσῆς καὶ Ἀαρὼν διὰ φωνῆς Κυρίου, ὃν τρόπον συνέταξεν αὐτοῖς Κύριος.

\vs{17}Καὶ ἦσαν οὗτοι οἱ υἱοὶ Λευὶ ἐξ ὀνομάτων αὐτῶν· Γεδσὼν, Καὰθ, καὶ Μεραρί.
\vs{18}Καὶ ταῦτα τὰ ὀνόματα τῶν υἱῶν Γεδσὼν κατὰ δήμους αὐτῶν· Λοβενὶ καὶ Σεμεΐ.
\vs{19}Καὶ υἱοὶ Καὰθ κατὰ δήμους αὐτῶν· Ἀμρὰμ καὶ Ἰσσαὰρ, Χεβρὼν καὶ Ὀζιήλ.
\vs{20}Καὶ υἱοὶ Μεραρὶ κατὰ δήμους αὐτῶν· Μοολὶ καὶ Μουσί· οὗτοί εἰσι δῆμοι τῶν Λευιτῶν κατʼ οἴκους πατριῶν αὐτῶν.
\vs{21}Τῷ Γεδσὼν δῆμος τοῦ Λοβενὶ, καὶ δῆμος τοῦ Σεμεΐ· οὗτοι δῆμοι τοῦ Γεδσών.
\vs{22}Ἡ ἐπίσκεψις αὐτῶν κατὰ ἀριθμὸν παντὸς ἀρσενικοῦ ἀπὸ μηνιαίου καὶ ἐπάνω, ἡ ἐπίσκεψις αὐτῶν, ἑπτακισχίλιοι καὶ πεντακόσιοι.
\vs{23}Καὶ οἱ υἱοὶ Γεδσὼν ὀπίσω τῆς σκηνῆς παρεμβαλοῦσι παρὰ θάλασσαν.
\vs{24}Καὶ ὁ ἄρχων οἴκου πατριᾶς τοῦ δήμου τοῦ Γεδσὼν, Ἑλισὰφ υἱὸς Δαήλ.
\vs{25}Καὶ ἡ φυλακὴ υἱῶν Γεδσὼν ἐν τῇ σκηνῇ τοῦ μαρτυρίου, ἡ σκηνὴ καὶ τὸ κάλυμμα, καὶ τὸ κατακάλυμμα τῆς θύρας τῆς σκηνῆς τοῦ μαρτυρίου,
\vs{26}καὶ τὰ ἱστία τῆς αὐλῆς, καὶ τὸ καταπέτασμα τῆς πύλης τῆς αὐλῆς τῆς οὔσης ἐπὶ τῆς σκηνῆς, καὶ τὰ κατάλοιπα πάντων τῶν ἔργων αὐτοῦ.

\vs{27}Τῷ Καὰθ δῆμος ὁ Ἀμρὰμ εἷς, καὶ δῆμος ὁ Ἰσσαὰρ εἷς, καὶ δῆμος ὁ Χεβρων εἷς, καὶ δῆμος ὁ Ὀζιὴλ εἷς· οὗτοί εἰσιν οἱ δῆμοι τοῦ Καὰθ, κατὰ ἀριθμόν.
\vs{28}Πᾶν ἀρσενικὸν ἀπὸ μηνιαίου καὶ ἐπάνω, ὀκτακισχίλιοι καὶ ἑξακόσιοι, φυλάσσοντες τὰς φυλακὰς τῶν ἁγίων.
\vs{29}Οἱ δῆμοι τῶν υἱῶν Καὰθ παρεμβαλοῦσιν ἐκ πλαγίων τῆς σκηνῆς κατὰ Λίβα.
\vs{30}Καὶ ὁ ἄρχων οἴκου πατριῶν τῶν δήμων τοῦ Καὰθ, Ἑλισαφὰν υἱὸς Ὀζιήλ.

\vs{31}Καὶ ἡ φυλακὴ αὐτῶν ἡ κιβωτὸς, καὶ ἡ τράπεζα, καὶ ἡ λυχνία, καὶ τὰ θυσιαστήρια, καὶ τὰ σκεύη τοῦ ἁγίου ὅσα λειτουργοῦσιν ἐν αὐτοῖς, καὶ τὸ κατακάλυμμα, καὶ πάντα τὰ ἔργα αὐτῶν.
\vs{32}Καὶ ὁ ἄρχων ἐπὶ τῶν ἀρχόντων τῶν Λευιτῶν, Ἐλεάζαρ ὁ υἱὸς Ἀαρὼν τοῦ ἱερέως, καθεσταμένος φυλάσσειν τὰς φυλακὰς τῶν ἁγίων.
\vs{33}Τῷ Μεραρὶ δῆμος ὁ Μοολὶ, καὶ δῆμος ὁ Μουσί· οὗτοί εἰσι δῆμοι τοῦ Μεραρί.
\vs{34}Ἡ ἐπίσκεψις αὐτῶν κατὰ ἀριθμὸν, πᾶν ἀρσενικὸν ἀπὸ μηνιαίου καὶ ἐπάνω, ἑξακισχίλιοι καὶ πεντήκοντα.
\vs{35}Καὶ ὁ ἄρχων οἴκου πατριῶν τοῦ δήμου τοῦ Μεραρὶ, Σουριὴλ υἱὸς Ἀβιχαίλ· ἐκ πλαγίων τῆς σκηνῆς παρεμβαλοῦσι πρὸς βοῤῥᾶν.
\vs{36}Ἡ ἐπίσκεψις τῆς φυλακῆς υἱῶν Μεραρὶ, τὰς κεφαλίδας τῆς σκηνῆς, καὶ τοὺς μοχλοὺς αὐτῆς, καὶ τοὺς στύλους αὐτῆς, καὶ τὰς βάσεις αὐτῆς, καὶ πάντα τὰ σκεύη αὐτῶν, καὶ τὰ ἔργα αὐτῶν,
\vs{37}καὶ τοὺς στύλους τῆς αὐλῆς κύκλῳ, καὶ τὰς βάσεις αὐτῶν, καὶ τοὺς πασσάλους, καὶ τοὺς κάλους αὐτῶν.

\vs{38}Οἱ παρεμβάλλοντες κατὰ πρόσωπον τῆς σκηνῆς τοῦ μαρτυρίου ἀπὸ ἀνατολῆς, Μωυσῆς καὶ Ἀαρὼν καὶ οἱ υἱοὶ αὐτοῦ, φυλάσσοντες τὰς φυλακὰς τοῦ ἁγίου εἰς τὰς φυλακὰς τῶν υἱῶν Ἰηραήλ· καὶ ὁ ἀλλογενὴς ὁ ἁπτόμενος, ἀποθανεῖται.
\vs{39}Πᾶσα ἡ ἐπίσκεψις τῶν Λευιτῶν, οὓς ἐπεσκέψατο Μωυσῆς καὶ Ἀαρὼν διὰ φωνῆς Κυρίου κατὰ δήμους αὐτῶν, πᾶν ἀρσενικὸν ἀπὸ μηνιαίου καὶ ἐπάνω, δύο καὶ εἴκοσι χιλιάδες.

\vs{40}Καὶ εἶπε Κύριος πρὸς Μωυσῆν, λέγων, ἐπίσκεψαι πᾶν πρωτότοκον ἄρσεν τῶν υἱῶν Ἰσραὴλ ἀπὸ μηνιαίου καὶ ἐπάνω· καὶ λάβετε τὸν ἀριθμὸν ἐξ ὀνόματος.
\vs{41}Καὶ λήψῃ τοὺς Λευίτας ἐμοί, ἐγὼ Κύριος, ἀντὶ πάντων τῶν πρωτοτόκων τῶν υἱῶν Ἰσραήλ, καὶ τὰ κτήνη τῶν Λευιτῶν ἀντὶ πάντων τῶν πρωτοτόκων ἐν τοῖς κτήνεσι τῶν υἱῶν Ἰσραήλ.
\vs{42}Καὶ ἐπεσκέψατο Μωυσῆς ὃν τρόπον ἐνετείλατο Κύριος πᾶν πρωτότοκον ἐν τοῖς υἱοῖς Ἰσραήλ.
\vs{43}Καὶ ἐγένοντο πάντα τὰ πρωτότοκα τὰ ἀρσενικὰ κατὰ ἀριθμὸν ἐξ ὀνόματος ἀπὸ μηνιαίου καὶ ἐπάνω ἐκ τῆς ἐπισκέψεως αὐτῶν, δύο καὶ εἴκοσι χιλιάδες καὶ τρεῖς καὶ ἑβδομήκοντα καὶ διακόσιοι.
\vs{44}Καὶ ἐλάλησε Κύριος πρὸς Μωυσῆν, λέγων,
\vs{45}λάβε τοὺς Λευίτας ἀντὶ πάντων τῶν πρωτοτόκων υἱῶν Ἰσραὴλ, καὶ τὰ κτήνη τῶν Λευιτῶν ἀντὶ τῶν κτηνῶν αὐτῶν, καὶ ἔσονται ἐμοὶ οἱ Λευῖται· ἐγὼ Κύριος.
\vs{46}Καὶ τὰ λύτρα τριῶν καὶ ἑβδομήκοντα καὶ διακοσίων οἱ πλεονάζοντες παρὰ τοὺς Λευίτας ἀπὸ τῶν πρωτοτόκων τῶν υἱῶν Ἰσραήλ·
\vs{47}Καὶ λήψῃ πέντε σίκλους κατὰ κεφαλὴν, κατὰ τὸ δίδραχμον τὸ ἅγιον λήψῃ, εἴκοσι ὀβολοὺς τοῦ σίκλου.
\vs{48}Καὶ δώσεις τὸ ἀργύριον Ἀαρὼν καὶ τοῖς υἱοῖς αὐτοῦ, λύτρα τῶν πλεοναζόντων ἐν αὐτοῖς.
\vs{49}Καὶ ἔλαβε Μωυσῆς τὸ ἀργύριον τὰ λύτρα τῶν πλεοναζόντων εἰς τὴν ἐκλύτρωσιν τῶν Λευιτῶν.
\vs{50}Παρὰ τῶν πρωτοτόκων τῶν υἱῶν Ἰσραὴλ ἔλαβε τὸ ἀργύριον, χιλίους τριακοσίους ἑξηκονταπέντε σίκλους, κατὰ τὸν σίκλον τὸν ἅγιον.
\vs{51}Καὶ ἔδωκε Μωυσῆς τὰ λύτρα τῶν πλεοναζόντων Ἀαρὼν καὶ τοῖς υἱοῖς αὐτοῦ, διὰ φωνῆς Κυρίου, ὃν τρόπον συνέταξε Κύριος τῷ Μωυσῇ.

\ch{4}
Καὶ ἐλάλησε Κύριος πρὸς Μωυσῆν καὶ Ἀαρὼν, λέγεν,
\vs{2}λάβε τὸ κεφάλαιον τῶν υἱῶν Καὰθ ἐκ μέσου υἱῶν Λευὶ, κατὰ δήμους αὐτῶν, κατʼ οἴκους πατριῶν αὐτῶν,
\vs{3}ἀπὸ εἴκοσι καὶ πέντε ἐτῶν καὶ ἐπάνω ἕως πεντήκοντα ἐτῶν, πᾶς ὁ εἰσπορευόμενος λειτουργεῖν, ποιῆσαι πάντα τὰ ἔργα ἐν τῇ σκηνῇ τοῦ μαρτυρίου.

\vs{4}Καὶ ταῦτα τὰ ἔργα τῶν υἱῶν Καὰθ ἐν τῇ σκηνῇ τοῦ μαρτυρίου· ἅγιον τῶν ἁγίων.
\vs{5}Καὶ εἰσελεύσεται Ἀαρὼν καὶ υἱοὶ αὐτοῦ, ὅταν ἐξαίρῃ ἡ παρεμβολὴ, καὶ καθελοῦσι τὸ καταπέτασμα τὸ συσκιάζον, καὶ κατακαλύψουσιν ἐν αὐτῷ τὴν κιβωτὸν τοῦ μαρτυρίου,
\vs{6}καὶ ἐπιθήσουσιν ἐπʼ αὐτὸ κατακάλυμμα δέρμα ὑακίνθινον, καὶ ἐπιβαλοῦσιν ἐπʼ αὐτὴν ἱμάτιον ὅλον ὑακίνθινον ἄνωθεν, καὶ διεμβαλοῦσι τοὺς ἀναφορεῖς.

\vs{7}Καὶ ἐπί τὴν τράπεζαν τὴν προκειμένην ἐπιβαλοῦσιν ἐπʼ αὐτὴν ἱμάτιον ὁλοπόρφυρον, καὶ τὰ τρυβλία, καὶ τὰς θυΐσκας, καὶ τοὺς κυάθους, καὶ τὰ σπονδεῖα ἐν οἷς σπένδει, καὶ οἱ ἄρτοι οἳ διαπαντὸς ἐπʼ αὐτῆς ἔσονται.
\vs{8}Καὶ ἐπιβαλοῦσιν ἐπʼ αὐτὴν ἱμάτιον κόκκινον, καὶ καλύψουσιν αὐτὴν καλύμματι δερματίνῳ ὑακινθίνῳ, καὶ διεμβαλοῦσι διʼ αὐτῆς τοὺς ἀναφορεῖς.
\vs{9}Καὶ λήψονται ἱμάτιον ὑακίνθινον, καὶ καλύψουσι τὴν λυχνίαν τὴν φωτίζουσαν, καὶ τοὺς λύχνους αὐτῆς, καὶ τὰς λαβίδας αὐτῆς, καὶ τὰς ἐπαρυστρίδας αὐτῆς, καὶ πάντα τὰ ἀγγεῖα τοῦ ἐλαίου οἷς λειτουργοῦσιν ἐν αὐτοῖς.
\vs{10}Καὶ ἐμβαλοῦσιν αὐτὴν, καὶ πάντα τὰ σκεύη αὐτῆς, εἰς κάλυμμα δερμάτινον ὑακίνθινον, καὶ ἐπιθήσουσιν αὐτὴν ἐπʼ ἀναφορέων.
\vs{11}Καὶ ἐπὶ τὸ θυσιαστήριον τὸ χρυσοῦν ἐπικαλύψουσιν ἱμάτιον ὑακίνθινον, καὶ καλύψουσιν αὐτὸ καλύμματι δερματίνῳ ὑακινθίνῳ, καὶ διεμβαλοῦσι τοὺς ἀναφορεῖς αὐτοῦ.

\vs{12}καὶ λήψονται πάντα τὰ σκεύη τὰ λειτουργικὰ ὅσα λειτουργοῦσιν ἐν αὐτοῖς ἐν τοῖς ἁγίοις· καὶ ἐμβαλοῦσιν εἰς ἱμάτιον ὑακίνθινον, καὶ καλύψουσιν αὐτὰ καλύμματι δερματίνῳ ὑακινθίνῳ, καὶ ἐπιθήσουσιν ἐπὶ ἀναφορεῖς.
\vs{13}Καὶ τὸν καλυπτῆρα ἐπιθήσει ἐπὶ τὸ θυσιαστήριον, καὶ ἐπικαλύψουσιν ἐπʼ αὐτὸ ἱμάτιον ὁλοπόρφυρον.
\vs{14}Καὶ ἐπιθήσουσιν ἐπʼ αὐτὸ πάντα τὰ σκεύη ὅσοις λειτουργοῦσιν ἐπʼ αὐτῷ ἐν αὐτοῖς, καὶ τὰ πυρεῖα, καὶ τὰς κρεάγρας, καὶ τὰς φιάλας, καὶ τὸν καλυπτῆρα, καὶ πάντα τὰ σκεύη τοῦ θυσιαστηρίου· καὶ ἐπιβαλοῦσιν ἐπʼ αὐτὸ κάλυμμα δερμάτινον ὑακίνθινον, καὶ διεμβαλοῦσι τοὺς ἀναφορεῖς αὐτοῦ· καὶ λήψονται ἱμάτιον πορφυροῦν, καὶ συγκαλύψουσι τὸν λουτῆρα καὶ τὴν βάσιν αὐτοῦ, καὶ ἐμβαλοῦσιν αὐτὸ εἰς κάλυμμα δερμάτινον ὑακίνθινον, καὶ ἐπιθήσουσιν ἐπὶ ἀναφορεῖς,
\vs{15}καὶ συντελέσουσιν Ἀαρὼν καὶ οἱ υἱοὶ αὐτοῦ, καλύπτοντες τὰ ἅγια, καὶ πάντα τὰ σκεύη τὰ ἅγια, ἐν τῷ ἐξαίρειν τὴν παρεμβολήν· καὶ μετὰ ταῦτα εἰσελεύσονται υἱοὶ Καὰθ αἴρειν, καὶ οὐχ ἅψονται τῶν ἁγίων, ἵνα μὴ ἀποθάνωσι· ταῦτα ἀροῦσιν οἱ υἱοὶ Καὰθ ἐν τῇ σκηνῇ τοῦ μαρτυρίου.

\vs{16}Ἐπίσκοπος Ἐλεάζαρ υἱὸς Ἀαρὼν τοῦ ἱερέως, τὸ ἔλαιον τοῦ φωτὸς, καὶ τὸ θυμίαμα τῆς συνθέσεως, καὶ ἡ θυσία ἡ καθʼ ἡμέραν, καὶ τὸ ἔλαιον τῆς χρίσεως, ἡ ἐπισκοπὴ ὅλης τῆς σκηνῆς, καὶ ὅσα ἐστὶν ἐν αὐτῇ ἐν τῷ ἁγίῳ, ἐν πᾶσι τοῖς ἔργοις.

\vs{17}Καὶ ἐλάλησε Κύριος πρὸς Μωυσῆν καὶ Ἀαρὼν, λέγων,
\vs{18}μὴ ὀλοθρεύσητε τῆς φυλῆς τὸν δῆμον τὸν Καὰθ ἐκ μέσου τῶν Λευιτῶν.
\vs{19}Τοῦτο ποιήσατε αὐτοῖς, καὶ ζήσονται καὶ οὐ μὴ ἀποθάνωσι, προσπορευομένων αὐτῶν πρὸς τὰ ἅγια τῶν ἁγίων· Ἀαρὼν καὶ οἱ υἱοὶ αὐτοῦ προσπορεύεσθωσαν, καὶ καταστήσουσιν αὐτοὺς ἕκαστον κατὰ τὴν ἀναφορὰν αὐτοῦ,
\vs{20}καὶ οὐ μὴ εἰσέλθωσιν ἰδεῖν ἐξάπινα τὰ ἅγια, καὶ ἀποθανοῦνται.

\vs{21}Καὶ ἐλάλησε Κύριος πρὸς Μωυσῆν, λέγων,
\vs{22}λάβε τὴν ἀρχὴν τῶν υἱῶν Γεδσὼν, καὶ τούτους κατʼ οἴκους πατριῶν αὐτῶν, κατὰ δήμους αὐτῶν,
\vs{23}ἀπὸ πέντε καὶ εἰκοσαετοῦς καὶ ἐπάνω ἕως πεντηκονταετοῦς ἐπίσκεψαι αὐτοὺς, πᾶς ὁ εἰσπορευόμενος λειτουργεῖν, ποιεῖν τὰ ἔργα αὐτοῦ ἐν τῇ σκηνῇ τοῦ μαρτυρίου.
\vs{24}Αὕτη ἡ λειτουργία τοῦ δήμου τοῦ Γεδσὼν, λειτουργεῖν καὶ αἴρειν.
\vs{25}Καὶ ἀρεῖ τὰς δέῤῥεις τῆς σκηνῆς, καὶ τὴν σκηνὴν τοῦ μαρτυρίου, καὶ τὸ κάλυμμα αὐτῆς, καὶ τὸ κατακάλυμμα τὸ ὑακίνθινον τὸ ὂν ἐπʼ αὐτῆς ἄνωθεν, καὶ τὸ κάλυμμα τῆς θύρας τῆς σκηνῆς τοῦ μαρτυρίου,
\vs{26}καὶ τὰ ἱστία τῆς αὐλῆς, ὅσα ἐπὶ τῆς σκηνῆς τοῦ μαρτυρίου, καὶ τὰ περισσὰ, καὶ πάντα τὰ σκεύη τὰ λειτουργικὰ ὅσα λειτουργοῦσιν ἐν αὐτοῖς ποιήσουσι.
\vs{27}Κατὰ στόμα Ἀαρὼν καὶ τῶν υἱῶν αὐτοῦ ἔσται ἡ λειτουργία τῶν υἱῶν Γεδσὼν κατὰ πάσας τὰς λειτουργίας αὐτῶν, καὶ κατὰ πάντα τὰ ἔργα αὐτῶν· καὶ ἐπισκέψῃ αὐτοὺς ἐξ ὀνόματος πάντα τὰ ἀρτὰ ὑπʼ αὐτῶν.
\vs{28}Αὕτη ἡ λειτουργία τῶν υἱῶν Γεδσὼν ἐν τῇ σκηνῇ τοῦ μαρτυρίου, καὶ ἡ φυλακὴ αὐτῶν ἐν χειρὶ Ἰθάμαρ τοῦ υἱοῦ Ἀαρὼν τοῦ ἱερέως.

\vs{29}Οἱ υἱοὶ Μεραρὶ κατὰ δήμους αὐτῶν, κατʼ οἴκους πατριῶν αὐτῶν, ἐπισκέψασθε αὐτοὺς,
\vs{30}ἀπὸ πέντε καὶ εἰκοσαετοῦς καὶ ἐπάνω ἕως πεντηκονταετοῦς ἐπισκέψασθε αὐτοὺς, πᾶς ὁ εἰσπορευόμενος λειτουργεῖν τὰ ἔργα τῆς σκηνῆς τοῦ μαρτυρίου.
\vs{31}Καὶ ταῦτα τὰ φυλάγματα τῶν αἰρομένων ὑπʼ αὐτῶν κατὰ πάντα τὰ ἔργα αὐτῶν ἐν τῇ σκηνῇ τοῦ μαρτυρίου· τὰς κεφαλίδας τῆς σκηνῆς, καὶ τοὺς μοχλοὺς, καὶ τοὺς στύλους αὐτῆς, καὶ τὰς βάσεις αὐτῆς, καὶ τὸ κατακάλυμμα, καὶ αἱ βάσεις αὐτῶν, καὶ οἱ στύλοι αὐτῶν, καὶ τὸ κατακάλυμμα τῆς θύρας τῆς σκηνῆς,
\vs{32}καὶ τοὺς στύλους τῆς αὐλῆς κύκλῳ, καὶ αἱ βάσεις αὐτῶν, καὶ τοὺς στύλους τοῦ καταπετάσματος τῆς πύλης τῆς αὐλῆς, καὶ τὰς βάσεις αὐτῶν, καὶ τοὺς πασσάλους αὐτῶν, καὶ τοὺς κάλους αὐτῶν, καὶ πάντα τὰ σκεύη αὐτῶν, καὶ πάντα τὰ λειτουργήματα αὐτῶν· ἐξ ὀνομάτων ἐπισκέψασθε αὐτοὺς, καὶ πάντα τὰ σκεύη τῆς φυλακῆς τῶν αἰρομένων ὑπʼ αὐτῶν.
\vs{33}Αὕτη ἡ λειτουργία δήμου υἱῶν Μεραρὶ ἐν πᾶσι τοῖς ἔργοις αὐτῶν ἐν τῇ σκηνῇ τοῦ μαρτυρίου ἐν χειρὶ Ἰθάμαρ τοῦ υἱοῦ Ἀαρὼν τοῦ ἱερέως.

\vs{34}Καὶ ἐπεσκέψατο Μωυσῆς καὶ Ἀαρὼν καὶ οἱ ἄρχοντες Ἰσραὴλ τοὺς υἱοὺς Καὰθ κατὰ δήμους αὐτῶν, κατʼ οἴκους πατριῶν αὐτῶν,
\vs{35}ἀπὸ πέντε καὶ εἰκοσαετοῦς καὶ ἐπάνω ἕως πεντηκονταετοῦς, πᾶς ὁ εἰσπορευόμενος λειτουργεῖν καὶ ποιεῖν ἐν τῇ σκηνῇ τοῦ μαρτυρίου.
\vs{36}Καὶ ἐγένετο ἡ ἐπίσκεψις αὐτῶν κατὰ δήμους αὐτῶν, δισχίλιοι ἑπτακόσιοι πεντήκοντα.
\vs{37}Αὕτη ἡ ἐπίσκεψις δήμου Καάθ, πᾶς ὁ λειτουργῶν ἐν τῇ σκηνῇ τοῦ μαρτυρίου, καθὰ ἐπεσκέψατο Μωυσῆς καὶ Ἀαρὼν διὰ φωνῆς Κυρίου, ἐν χειρὶ Μωυσῆ.

\vs{38}Καὶ ἐπεσκέπησαν υἱοὶ Γεδσὼν κατὰ δήμους αὐτῶν, κατʼ οἴκους πατριῶν αὐτῶν,
\vs{39}ἀπὸ πέντε καὶ εἰκοσαετοῦς καὶ ἐπάνω ἕως πεντηκονταετοῦς, πᾶς ὁ εἰσπορευόμενος λειτουργεῖν καὶ ποιεῖν τὰ ἔργα ἐν τῇ σκηνῇ τοῦ μαρτυρίου.
\vs{40}Καὶ ἐγένετο ἡ ἐπίσκεψις αὐτῶν, κατὰ δήμους αὐτῶν, κατʼ οἴκους πατριῶν αὐτῶν, δισχίλιοι ἑξακόσιοι τριάκοντα.
\vs{41}Αὕτη ἡ ἐπίσκεψις δήμου υἱῶν Γεδσὼν, πᾶς ὁ λειτουργῶν ἐν τῇ σκηνῇ τοῦ μαρτυρίου, οὓς ἐπεσκέψατο Μωυσῆς καὶ Ἀαρὼν διὰ φωνῆς Κυρίου, ἐν χειρὶ Μωυσῆ.

\vs{42}Ἐπεσκέπησαν δὲ καὶ δῆμος υἱῶν Μεραρὶ κατὰ δήμους αὐτῶν, κατʼ οἴκους πατριῶν αὐτῶν,
\vs{43}ἀπὸ πέντε καὶ εἰκοσαετοῦς καὶ ἐπάνω ἕως πεντηκονταετοῦς, πᾶς ὁ εἰσπορευόμενος λειτουργεῖν πρὸς τὰ ἔργα τῆς σκηνῆς τοῦ μαρτυρίου.
\vs{44}Καὶ ἐγενήθη ἡ ἐπίσκεψις αὐτῶν κατὰ δήμους αὐτῶν, κατʼ οἴκους πατριῶν αὐτῶν, τρισχίλιοι καὶ διακόσιοι.
\vs{45}Αὕτη ἡ ἐπίσκεψις δήμου υἱῶν Μεραρὶ, οὓς ἐπεσκέψατο Μωυσῆς καὶ Ἀαρὼν διὰ φωνῆς Κυρίου, ἐν χειρὶ Μωυσῆ.
\vs{46}Πάντες οἱ ἐπεσκεμμένοι, οὓς ἐπεσκέψατο Μωυσῆς καὶ Ἀαρὼν καὶ οἱ ἄρχοντες Ἰσραὴλ τοὺς Λευίτας, κατὰ δήμους καὶ κατʼ οἴκους πατριῶν αὐτῶν,
\vs{47}ἀπὸ πέντε καὶ εἰκοσαετοῦς καὶ ἐπάνω ἕως πεντηκονταετοῦς, πᾶς ὁ εἰσπορευόμενος πρὸς τὸ ἔργον τῶν ἔργων, καὶ τὰ ἔργα τὰ αἰρόμενα ἐν τῇ σκηνῇ τοῦ μαρτυρίου.
\vs{48}Καὶ ἐγενήθησαν οἱ ἐπισκεπέντες, ὀκτακισχίλιοι πεντακόσιοι ὀγδοήκοντα.
\vs{49}Διὰ φωνῆς Κυρίου ἐπεσκέψατο αὐτοὺς ἐν χειρὶ Μωυσῆ, ἄνδρα κατὰ ἄνδρα ἐπὶ τῶν ἔργων αὐτῶν, καὶ ἐπὶ ὧν αἴρουσιν αὐτοί· καὶ ἐπεσκέπησαν, ὃν τρόπον συνέταξε Κύριος τῷ Μωυσῇ.

\ch{5}
Καὶ ἐλάλησε Κύριος πρὸς Μωυσῆν, λέγων,
\vs{2}πρόσταξον τοῖς υἱοῖς Ἰσραὴλ, καὶ ἐξαποστειλάτωσαν ἐκ τῆς παρεμβολῆς πάντα λεπρὸν, καὶ πάντα γονοῤῥυῆ, καὶ πάντα ἀκάθαρτον ἐπὶ ψυχῇ.
\vs{3}Ἀπὸ ἀρσενικοῦ ἕως θηλυκοῦ, ἐξαποστείλατε ἔξω τῆς παρεμβολῆς, καὶ οὐ μὴ μιανοῦσι τὰς παρεμβολὰς αὐτῶν, ἐν οἷς ἐγὼ καταγίνομαι ἐν αὐτοῖς.
\vs{4}Καὶ ἐποίησαν οὕτως οἱ υἱοὶ Ἰσραὴλ, καὶ ἐξαπέστειλαν αὐτοὺς ἔξω τῆς παρεμβολῆς· καθὰ ἐλάλησε Κύριος Μωυσῇ, οὕτως ἐποίησαν οἱ υἱοὶ Ἰσραήλ.

\vs{5}Καὶ ἐλάλησε Κύριος πρὸς Μωυσῆν, λέγων,
\vs{6}λάλησον τοῖς υἱοῖς Ἰσραὴλ, λέγων, ἀνὴρ ἢ γυνὴ, ὅστις ἂν ποιήσῃ ἀπὸ πασῶν τῶν ἁμαρτιῶν τῶν ἀνθρωπίνων, καὶ παριδὼν παρίδῃ καὶ πλημμελήσῃ ἡ ψυχὴ ἐκείνη,
\vs{7}ἐξαγορεύσει τὴν ἁμαρτίαν, ἣν ἐποίησε, καὶ ἀποδώσει τὴν πλημμέλειαν· τὸ κεφάλαιον, καὶ τὸ ἐπίπεμπτον αὐτοῦ προσθήσει ἐπʼ αὐτὸ, καὶ ἀποδώσει τίνι ἐπλημμέλησεν αὐτῷ.
\vs{8}Ἐὰν δὲ μὴ ᾖ τῷ ἀνθρώπῳ ὁ ἀγχιστεύων, ὥστε ἀποδοῦναι αὐτῷ τὸ πλημμέλημα πρὸς αὐτὸν, τὸ πλημμέλημα τὸ ἀποδιδόμενον Κυρίῳ, τῷ ἱερεῖ ἔσται, πλὴν τοῦ κριοῦ τοῦ ἱλασμοῦ, διʼ οὗ ἐξιλάσεται ἐν αὐτῷ περὶ αὐτοῦ.

\vs{9}Καὶ πᾶσα ἀπαρχὴ κατὰ πάντα τὰ ἁγιαζόμενα ἐν υἱοῖς Ἰσραὴλ, ὅσα ἐὰν προσφέρωσι κυρίῳ, τῷ ἱερεῖ αὐτῷ ἔσται·
\vs{10}Καὶ ἑκάστου τὰ ἡγιασμένα, αὐτοῦ ἔσται· καὶ ἀνὴρ, ὃς ἂν δῷ τῷ ἱερεῖ, αὐτῷ ἔσται.

\vs{11}Καὶ ἐλάλησε Κύριος πρὸς Μωυσῆν, λέγων,
\vs{12}λάλησον τοῖς υἱοῖς Ἰσραὴλ, καὶ ἐρεῖς πρὸς αὐτοὺς, ἀνδρὸς ἀνδρὸς ἐὰν παραβῇ ἡ γυνὴ αὐτοῦ, καὶ ὑπεριδοῦσα παρίδῃ αὐτὸν,
\vs{13}καὶ κοιμηθῇ τις μετʼ αὐτῆς κοίτην σπέρματος, καὶ λάθῃ ἐξ ὀφθαλμῶν τοῦ ἀνδρὸς αὐτῆς, καὶ κρύψῃ, αὐτὴ δὲ ᾖ μεμιασμένη, καὶ μάρτυς μὴ ἦν μετʼ αὐτῆς, καὶ αὐτὴ μὴ ᾖ συνειλημμένη,
\vs{14}καὶ ἐπέλθῃ αὐτῷ πνεῦμα ζηλώσεως, καὶ ζηλώσῃ τὴν γυναῖκα αὐτοῦ, αὐτὴ δὲ μεμίανται, ἢ ἐπέλθῃ αὐτῷ πνεῦμα ζηλώσεως, καὶ ζηλώσῃ τὴν γυναῖκα αὐτοῦ, αὐτὴ δὲ μὴ ᾖ μεμιασμένη,
\vs{15}καὶ ἄξει ὁ ἄνθρωπος τὴν γυναῖκα αὐτοῦ πρὸς τὸν ἱερέα, καὶ προσοίσει τὸ δῶρον περὶ αὐτῆς, τὸ δέκατον τοῦ οἰφὶ ἄλευρον κρίθινον· οὐκ ἐπιχεεῖ ἐπʼ αὐτὸ ἔλαιον, οὐδὲ ἐπιθήσει ἐπʼ αὐτὸ λίβανον· ἔστι γὰρ θυσία ζηλοτυπίας, θυσία μνημοσύνου, ἀναμιμνήσκουσα ἁμαρτίαν.

\vs{16}Καὶ προσάξει αὐτὴν ὁ ἱερεὺς, καὶ στήσει αὐτὴν ἔναντι Κυρίου.
\vs{17}Καὶ λήψεται ὁ ἱερεὺς ὕδωρ καθαρὸν ζῶν ἐν ἀγγείῳ ὀστρακίνῳ, καὶ τῆς γῆς τῆς οὔσης ἐπὶ τοῦ ἐδάφους τῆς σκηνῆς τοῦ μαρτυρίου, καὶ λαβὼν ὁ ἱερεὺς ἐμβαλεῖ εἰς τὸ ὕδωρ.
\vs{18}Καὶ στήσει ὁ ἱερεὺς τὴν γυναῖκα ἔναντι Κυρίου, καὶ ἀποκαλύψει τὴν κεφαλὴν τῆς γυναικὸς, καὶ δώσει ἐπὶ τὰς χεῖρας αὐτῆς τὴν θυσίαν τοῦ μνημοσύνου, τὴν θυσίαν τῆς ζηλοτυπίας· ἐν δὲ τῇ χειρὶ τοῦ ἱερέως ἔσται τὸ ὕδωρ τοῦ ἐλεγμοῦ τοῦ ἐπικαταρωμένου τούτου.
\vs{19}Καὶ ὁρκιεῖ αὐτὴν ὁ ἱερεὺς, καὶ ἐρεῖ τῇ γυναικὶ, εἰ μὴ κεκοίμηταί τις μετὰ σοῦ, εἰ μὴ παραβέβηκας μιανθῆναι ὑπὸ τὸν ἄνδρα τὸν σεαυτῆς, ἀθῶα ἴσθι ἀπὸ τοῦ ὕδατος τοῦ ἐλεγμοῦ τοῦ ἐπικαταρωμένου τούτου.
\vs{20}Εἰ δὲ σὺ παραβέβηκας ὕπανδρος οὖσα, ἢ μεμίανσαι, καὶ ἔδωκέ τις τὴν κοίτην αὐτοῦ ἐν σοὶ, πλὴν τοῦ ἀνδρός σου·
\vs{21}Καὶ ὁρκιεῖ ὁ ἱερεὺς τὴν γυναῖκα ἐν τοῖς ὅρκοις τῆς ἀρᾶς ταύτης, καὶ ἐρεῖ ὁ ἱερεὺς τῇ γυναικὶ, δῴη σε Κύριος ἐν ἀρᾷ καὶ ἐνόρκιον ἐν μέσῳ τοῦ λαοῦ σου, ἐν τῷ δοῦναι Κύριον τὸν μηρόν σου διαπεπτωκότα, καὶ τὴν κοιλίαν σου πεπρησμένην.
\vs{22}Καὶ εἰσελεύσεται τὸ ὕδωρ τὸ ἐπικαταρώμενον τοῦτο εἰς τὴν κοιλίαν σου πρῆσαι γαστέρα, καὶ διαπεσεῖν μηρόν σου· καὶ ἐρεῖ ἡ γυνὴ, γένοιτο, γένοιτο.

\vs{23}Καὶ γράψει ὁ ἱερεὺς τὰς ἀρὰς ταύτας εἰς βιβλίον, καὶ ἐξαλείψει εἰς τὸ ὕδωρ τοῦ ἐλεγμοῦ τοῦ ἐπικαταρωμένου.
\vs{24}Καὶ ποτιεῖ τὴν γυναῖκα τὸ ὕδωρ τοῦ ἐλεγμοῦ τοῦ ἐπικαταρωμένου· καὶ εἰσελεύσεται εἰς αὐτὴν τὸ ὕδωρ τὸ ἐπικαταρώμενον τοῦ ἐλεγμοῦ.

\vs{25}Καὶ λήψεται ὁ ἱερεὺς ἐκ χειρὸς τῆς γυναικὸς τὴν θυσίαν τῆς ζηλοτυπίας, καὶ ἐπιθήσει τὴν θυσίαν ἔναντι Κυρίου, καὶ προσοίσει αὐτὴν πρὸς τὸ θυσιαστήριον.
\vs{26}Καὶ δράξεται ὁ ἱερεὺς ἀπὸ τῆς θυσίας τὸ μνημόσυνον αὐτῆς, καὶ ἀνοίσεται αὐτὸ ἐπὶ τὸ θυσιαστήριον, καὶ μετὰ. ταῦτα ποτιεῖ τὴν γυναῖκα τὸ ὕδωρ.
\vs{27}Καὶ ἔσται ἐὰν ᾖ μεμιασμένη καὶ λήθῃ λάθῃ τὸν ἄνδρα αὐτῆς, καὶ εἰσελεύσεται εἰς αὐτὴν τὸ ὕδωρ τοῦ ἐλεγμοῦ τὸ ἐπικαταρώμενον, καὶ πρησθήσεται τὴν κοιλίαν, καὶ διαπεσεῖται ὁ μηρὸς αὐτῆς, καὶ ἔσται ἡ γυνὴ εἰς ἀρὰν τῷ λαῷ αὐτῆς.
\vs{28}Ἐὰν δὲ μὴ μιανθῇ ἡ γυνὴ, καὶ καθαρὰ ᾖ, καὶ ἀθῶα ἔσται καὶ ἐκσπερματιεῖ σπέρμα.
\vs{29}Οὗτος ὁ νόμος τῆς ζηλοτυπίας, ᾧ ἂν παραβῇ ἡ γυνὴ ὕπανδρος οὖσα, καὶ μιανθῇ.
\vs{30}Ἢ ἄνθρωπος ὃς ἐὰν ἐπέλθῃ ἐπʼ αὐτὸν πνεῦμα ζηλώσεως, καὶ ζηλώσῃ τὴν γυναῖκα αὐτοῦ, καὶ στήσῃ τὴν γυναῖκα αὐτοῦ ἔναντι Κυρίου, καὶ ποιήσει αὐτῇ ὁ ἱερεὺς πάντα τὸν νόμον τοῦτον,
\vs{31}καὶ ἀθῶος ἔσται ὁ ἄνθρωπος ἀπὸ ἁμαρτίας· καὶ γυνὴ ἐκείνη λήψεται τὴν ἁμαρτίαν αὐτῆς.

\ch{6}
Καὶ ἐλάλησε Κύριος πρὸς Μωυσῆν, λέγων, λάλησον τοῖς υἱοῖς Ἰσραὴλ,
\vs{2}καὶ ἐρεῖς πρὸς αὐτοὺς, ἀνὴρ ἢ γυνὴ, ὃς ἂν μεγάλως εὔξηται εὐχὴν ἀφαγνίσασθαι ἁγνείαν Κυρίῳ,
\vs{3}ἀπὸ οἴνου καὶ σίκερα ἁγνισθήσεται· καὶ ὄξος ἐξ οἴνου καὶ ὄξος ἐκ σίκερα οὐ πίεται· καὶ ὅσα κατεργάζεται ἐκ σταφυλῆς οὐ πίεται· καὶ σταφυλὴν πρόσφατον καὶ σταφίδα οὐ φάγεται πάσας τὰς ἡμέρας τῆς εὐχῆς αὐτοῦ·
\vs{4}ἀπὸ πάντων ὅσα γίνεται ἐξ ἀμπέλου, οἶνον ἀπὸ στεμφύλων ἕως γιγάρτου οὐ φάγεται πάσας τὰς ἡμέρας τοῦ ἁγνισμοῦ·
\vs{5}ξυρὸν οὐκ ἐπελεύσεται ἐπὶ τὴν κεφαλὴν αὐτοῦ, ἓως ἄν πληρωθῶσιν αἱ ἡμέραι, ὅσας ηὐξατο Κυρίῳ· ἅγιος ἔσται τρέφων κόμην τρίχα κεφαλῆς πάσας τὰς ἡμέρας τῆς εὐχῆς Κυρίῳ·
\vs{6}ἐπὶ πάσῃ ψυχῇ τετελευτηκυίᾳ οὐκ εἰσελεύσεται ἐπὶ πατρὶ καὶ μητρὶ,
\vs{7}καὶ ἐπʼ ἀδελφῷ καὶ ἐπʼ ἀδελφῇ, οὐ μιανφήσεται ἐπʼ αὐτοὶς ἀποφανόντων αὐτῶν, ὅτι εὐχὴ Θεοῦ αὐτοῦ ἐπʼ αὐτῷ ἐπὶ κεφαλῆς αὐτοῦ·

\vs{8}Πάσας τὰς ἡμέρας τῆς εὐχῆς αὐτοῦ ἅγιος ἔσται Κυρίῳ.
\vs{9}Ἐὰν δέ τις ἀποθάνῃ ἐπʼ αὐτῷ ἐξάπινα, παραχρῆμα μιανθήσεται ἡ κεφαλὴ εὐχῆς αὐτοῦ· καὶ ξυρήσεται τὴν κεφαλὴν αὐτοῦ ᾗ ἂν ἡμέρᾳ καθαρισθῇ· τῇ ἡμέρᾳ τῇ ἑβδόμῃ ξυρηθήσεται.
\vs{10}Καὶ τῇ ἡμέρᾳ τῇ ὀγδόῃ οἴσει δύο τρυγόνας, ἢ δύο νοσσοὺς περιστερῶν πρὸς τὸν ἱερέα, ἐπὶ τὰς θύρας τῆς σκηνῆς τοῦ μαρτυρίου.

\vs{11}Καὶ ποιήσει ὁ ἱερεὺς μίαν περὶ ἁμαρτίας, καὶ μίαν εἰς ὁλοκαύτωμα· καὶ ἐξιλάσεται περὶ αὐτοῦ ὁ ἱερεὺς περὶ ὧν ἥμαρτε περὶ τῆς ψυχῆς· καὶ ἁγιάσει τὴν κεφαλὴν αὐτοῦ ἐν ἐκείνῃ τῇ ἡμέρᾳ, ᾗ ἡγιάσθη Κυρίῳ, τὰς ἡμέρας τῆς εὐχῆς·
\vs{12}καὶ προσάξει ἀμνὸν ἐνιαύσιον εἰς πλημμέλειαν· καὶ αἱ ἡμέραι αἱ πρότεραι ἄλογοι ἔσονται, ὅτι ἐμιάνθη ἡ κεφαλὴ εὐχῆς αὐτοῦ.

\vs{13}Καὶ οὗτος ὁ νόμος τοῦ εὐξαμένου· ᾗ ἂν ἡμέρᾳ πληρώσῃ ἡμέρας εὐχῆς αὐτοῦ, προσοίσει αὐτὸς παρὰ τὰς θύρας τῆς σκηνῆς τοῦ μαρτυρίου.
\vs{14}Καὶ προσάξει τὸ δῶρον αὐτοῦ Κυρίῳ ἀμνὸν ἐνιαύσιον ἄμωμον ἕνα εἰς ὁλοκαύτωσιν, καὶ ἀμνάδα ἐνιαυσίαν μίαν ἄμωμον εἰς ἁμαρτίαν, καὶ κριὸν ἕνα ἄμωμον εἰς σωτήριον,
\vs{15}καὶ κανοῦν ἀζύμων σεμιδάλεως ἄρτους ἀναπεποιημένους ἐν ἐλαίῳ, καὶ λάγανα ἄζυμα κεχρισμένα ἐν ἐλαίῳ, καὶ θυσίαν αὐτῶν, καὶ σπονδὴν αὐτῶν.
\vs{16}Καὶ προσοίσει ὁ ἱερεὺς ἔναντι Κυρίου, καὶ ποιήσει τὸ περὶ ἁμαρτίας αὐτοῦ, καὶ τὸ ὁλοκαύτωμα αὐτοῦ.
\vs{17}Καὶ τὸν κριὸν ποιήσει θυσίαν σωτηρίου τῷ Κυρίῳ ἐπὶ τῷ κανῷ τῶν ἀζύμων· καὶ ποιήσει ὁ ἱερεὺς τὴν θυσίαν αὐτοῦ, καὶ τὴν σπονδὴν αὐτοῦ.
\vs{18}Καὶ ξυρήσεται ὁ ηὐγμένος παρὰ τὰς θύρας τῆς σκηνῆς τοῦ μαρτυρίου τὴν κεφαλὴν τῆς εὐχῆς αὐτοῦ, καὶ ἐπιθήσει τὰς τρίχας ἐπὶ τὸ πῦρ, ὅ ἐστιν ὑπὸ τὴν θυσίαν τοῦ σωτηρίου.

\vs{19}Καὶ λήψεται ὁ ἱερεὺς τὸν βραχίονα ἐφθὸν ἀπὸ τοῦ κριοῦ, καὶ ἄρτον ἕνα ἄζυμον ἀπὸ τοῦ κανοῦ, καὶ λάγανον ἄζυμον ἕν, καὶ ἐπιθήσει ἐπὶ τὰς χεῖρας τοῦ ηὐγμένου μετὰ τὸ ξυρήσασθαι αὐτὸν τὴν εὐχὴν αὐτοῦ,
\vs{20}καὶ προσοίσει αὐτὰ ὁ ἱερεὺς ἐπίθεμα ἔναντι Κυρίου· ἅγιον ἔσται τῷ ἱερεῖ ἐπὶ τοῦ στηθηνίου τοῦ ἐπιθέματος, καὶ ἐπὶ τοῦ βραχίονος τοῦ ἀφαιρέματος· καὶ μετὰ ταῦτα πίεται ὁ ηὐγμένος οἶνον.
\vs{21}Οὗτος ὁ νόμος τοῦ εὐξαμένου, ὃς ἂν εὔξηται Κυρίῳ δῶρον αὐτοῦ Κυρίῳ περὶ τῆς εὐχῆς, χωρὶς ὧν ἂν εὕρῃ ἡ χεὶρ αὐτοῦ, κατὰ δύναμιν τῆς εὐχῆς αὐτοῦ, ἣν ἂν εὔξηται κατὰ νόμον ἁγνείας.

\vs{22}Καὶ ἐλάλησε Κύριος πρὸς Μωυσῆν, λέγων,
\vs{23}λάλησον Ἀαρὼν καὶ τοῖς υἱοῖς αὐτοῦ, λέγων, οὕτως εὐλογήσετε τοὺς υἱοὺς Ἰσραὴλ, λέγοντες αὐτοῖς,
\vs{24}εὐλογήσαι σε Κύριος, καὶ φυλάξαι, σε.
\vs{25}Ἐπιφάναι Κύριος τὸ πρόσωπον αὐτοῦ ἐπὶ σὲ, καὶ ἐλεήσαι σε.
\vs{26}Ἐπάραι Κύριος τὸ πρόσωπον αὐτοῦ ἐπὶ σὲ, καὶ δῴη σοι εἰρήνην.
\vs{27}Καὶ ἐπιθήσουσι τὸ ὄνομά μου ἐπὶ τοὺς υἱοὺς Ἰσραὴλ, καὶ ἐγὼ Κύριος εὐλογήσω αὐτούς.

\ch{7}
Καὶ ἐγένετο ᾗ ἡμέρᾳ συνετέλεσε Μωυσῆς, ὥστε ἀναστῆσαι τὴν σκηνήν, καὶ ἔχρισεν αὐτὴν, καὶ ἡγίασεν αὐτὴν, καὶ πάντα τὰ σκεύη αὐτῆς, καὶ τὸ θυσιαστήριον, καὶ πάντα τὰ σκεύη αὐτοῦ, καὶ ἔχρισεν αὐτὰ, καὶ ἡγίασεν αὐτά.
\vs{2}Καὶ προσήνεγκαν οἱ ἄρχοντες Ἰσραὴλ, δώδεκα ἄρχοντες οἴκων πατριῶν αὐτῶν· οὗτοι οἱ ἄρχοντες φυλῶν, οὗτοι οἱ παρεστηκότες ἐπὶ τῆς ἐπισκοπῆς.
\vs{3}Καὶ ἤνεγκαν τὸ δῶρον αὐτῶν ἔναντι Κυρίου, ἓξ ἁμάξας λαμπηνίκας, καὶ δώδεκα βόας· ἅμαξαν παρὰ δύο ἀρχόντων, καὶ μόσχον παρὰ ἑκάστου· καὶ προσήγαγον ἐναντίον τῆς σκηνῆς.
\vs{4}Καὶ εἶπε Κύριος πρὸς Μωυσῆν, λέγων,
\vs{5}λάβε παρʼ αὐτῶν, καὶ ἔσονται πρὸς τὰ ἔργα τὰ λειτουργικὰ τῆς σκηνῆς τοῦ μαρτυρίου· καὶ δώσεις αὐτὰ τοῖς Λευίταις, ἑκάστῳ κατὰ τὴν αὐτοῦ λειτουργίαν.
\vs{6}Καὶ λαβὼν Μωυσῆς τὰς ἁμάξας καὶ τοὺς βόας, ἔδωκεν αὐτὰ τοῖς Λευίταις.
\vs{7}Καὶ τὰς δύο ἁμάξας καὶ τοὺς τέσσαρας βόας ἔδωκε τοῖς υἱοῖς Γεδσὼν κατὰ τὰς λειτουργίας αὐτῶν.
\vs{8}Καὶ τὰς τέσσαρας ἁμάξας καὶ τοὺς ὀκτὼ βόας ἔδωκε τοῖς υἱοῖς Μεραρὶ κατὰ τὰς λειτουργίας αὐτῶν, διὰ Ἰθάμαρ υἱοῦ Ἀαρὼν τοῦ ἱερέως.
\vs{9}Καὶ τοῖς υἱοῖς Καὰθ οὐ δέδωκεν, ὅτι τὰ λειτουργήματα τοῦ ἁγίου ἔχουσιν· ἐπʼ ὤμων ἀροῦσι.

\vs{10}Καὶ προσήνεγκαν οἱ ἄρχοντες εἰς τὸν ἐγκαινισμὸν τοῦ θυσιαστηρίου, ἐν τῇ ἡμέρᾳ ᾗ ἔχρισεν αὐτὸ, καὶ προσήνεγκαν οἱ ἄρχοντες τὰ δῶρα αὐτῶν ἀπέναντι τοῦ θυσιαστηρίου.
\vs{11}Καὶ εἶπε Κύριος πρὸς Μωυσῆν, ἄρχων εἷς καθʼ ἡμέραν, ἄρχων καθʼ ἡμέραν προσοίσουσι τὰ δῶρα αὐτῶν εἰς τὸν ἐγκαινισμὸν τοῦ θυσιαστηρίου.

\vs{12}Καὶ ἦν ὁ προσφέρων ἐν τῇ ἡμέρᾳ τῇ πρώτῃ τὸ δῶρον αὐτοῦ, Ναασσὼν υἱὸς Ἀμιναδὰβ, ἄρχων τῆς φυλῆς Ἰούδα.
\vs{13}Καὶ προσήνεγκε τὸ δῶρον αὐτοῦ, τρυβλίον ἀργυροῦν ἕν, τριάκοντα καὶ ἑκατὸν ὁλκὴ αὐτοῦ· φιάλην μίαν ἀργυρᾶν, ἑβδομήκοντα σίκλων κατὰ τὸν σίκλον τὸν ἅγιον· ἀμφοτέρα πλήρη σεμιδάλεως ἀναπεποιημένης ἐν ἐλαίῳ εἰς θυσίαν.
\vs{14}Θυΐσκην μίαν δέκα χρυσῶν, πλήρη θυμιάματος.
\vs{15}Μόσχον ἕνα ἐκ βοῶν, κριὸν ἕνα, ἀμνὸν ἕνα ἐνιαύσιον εἰς ὁλοκαύτωμα,
\vs{16}καὶ χίμαρον ἐξ αἰγῶν ἕνα περὶ ἁμαρτίας.
\vs{17}Καὶ εἰς θυσίαν σωτηρίου δαμάλεις δύο, κριοὺς πέντε, τράγους πέντε, ἀμνάδας ἐνιαυσίας πέντε· τοῦτο δῶρον Ναασσὼν υἱοῦ Ἀμιναδάβ.

\vs{18}Τῇ ἡμέρᾳ τῇ δευτέρᾳ προσήνεγκε Ναθαναὴλ υἱὸς Σωγὰρ, ὅ ἄρχων τῆς φυλῆς Ἰσσάχαρ.
\vs{19}Καὶ προσήνεγκε τὸ δῶρον αὐτοῦ, τρυβλίον ἀργυροῦν ἓν, τριάκοντα καὶ ἑκατὸν ὁλκὴ αὐτοῦ· φιάλην μίαν ἀργυρᾶν, ἑβδομήκοντα σίκλων κατὰ τὸν σίκλον τὸν ἅγιον· ἀμφότερα πλήρη σεμιδάλεως ἀναπεποιημένης ἐν ἐλαίῳ εἰς θυσίαν.
\vs{20}Θυΐσκην μίαν δέκα χρυσῶν, πλήρη θυμιάματος.
\vs{21}Μόσχον ἔνα ἐκ βοῶν, κριὸν ἕνα, ἀμνὸν ἕνα ἐνιαύσιον εἰς ὁλοκαύτωμα,
\vs{22}καὶ χίμαρον ἐξ αἰγῶν ἕνα περὶ ἁμαρτίας.
\vs{23}Καὶ εἰς θυσίαν σωτηρίου δαμάλεις δύο, κριοὺς πέντε, τράγους πέντε, ἀμνάδας ἐνιαυσίας πέντε· τοῦτο τὸ δῶρον Ναθαναὴλ υἱοῦ Σωγάρ.

\vs{24}Τῇ ἡμέρᾳ τῇ τρίτῃ ἄρχων τῶν υἱῶν Σαβουλὼν, Ἑλιὰβ υἱὸς Χαιλών.
\vs{25}Τὸ δῶρον αὐτοῦ, τρυβλίον ἀργυροῦν ἕν, τριάκοντα καὶ ἑκατὸν ὁλκὴ αὐτοῦ· φιάλην μίαν ἀργυρᾶν, ἑβδομήκοντα σίκλων κατὰ τὸν σίκλον τὸν ἅγιον· ἀμφότερα πλήρη σεμιδάλεως ἀναπεποιημένης ἐν ἐλαίῳ εἰς θυσίαν·
\vs{26}Θυΐσκην μίαν δέκα χρυσῶν, πλήρη θυμιάματος.
\vs{27}Μόσχον ἕνα ἐκ βοῶν, κριὸν ἕνα, ἀμνὸν ἕνα ἐνιαύσιον εἰς ὁλοκαύτωμα,
\vs{28}καὶ χίμαρον ἐξ αἰγῶν ἕνα περὶ ἁμαρτίας.
\vs{29}Καὶ εἰς θυσίαν σωτηρίου δαμάλεις δύο, κριοὺς πέντε, τράγους πέντε, ἀμνάδας ἐνιαυσίας πέντε· τοῦτο τὸ δῶρον Ἑλιὰβ υἱοῦ Χαιλών.

\vs{30}Τῇ ἡμέρᾳ τῇ τετάρτῃ ἄρχων τῶν υἱῶν Ῥουβὴν, Ἑλισοῦρ υἱὸς Σεδιούρ.
\vs{31}Τὸ δῶρον αὐτοῦ, τρυβλίον ἀργυροῦν ἓν, τριάκοντα καὶ ἑκατὸν ὁλκὴ αὐτοῦ· φιάλην μίαν ἀργυρᾶν, ἑβδομήκοντα σίκλων κατὰ τὸν σίκλον τὸν ἅγιον· ἀμφότερα πλήρη σεμιδάλεως ἀναπεποιημένης ἐν ἐλαίῳ εἰς θυσίαν·
\vs{32}Θυΐσκην μίαν δέκα χρυσῶν, πλήρη θυμιάματος.
\vs{33}Μόσχον ἕνα ἐκ βοῶν, κριὸν ἕνα, ἀμνὸν ἕνα ἐνιαύσιον εἰς ὁλοκαύτωμα,
\vs{34}καὶ χίμαρον ἐξ αἰγῶν ἕνα περὶ ἁμαρτίας.
\vs{35}Καὶ εἰς θυσίαν σωτηρίου δαμάλεις δύο, κριοὺς πέντε, τράγους πέντε, ἀμνάδας ἐνιαυσίας πέντε· τοῦτο τὸ δῶρον Ἑλισοὺρ υἱοῦ Σεδιούρ.

\vs{36}Τῇ ἡμέρᾳ τῇ πέμπτῃ ἄρχων τῶν υἱῶν Συμεὼν, Σαλαμιὴλ υἱὸς Σουρισαδαί.
\vs{37}Τὸ δῶρον αὐτοῦ, τρυβλίον ἀργυροῦν ἓν, τριάκοντα καὶ ἑκατὸν ὁλκὴ αὐτοῦ· φιάλην μίαν ἀργυρᾶν, ἑβδομήκοντα σίκλων κατὰ τὸν σίκλον τὸν ἅγιον· ἀμφότερα πλήρη σεμιδάλεως ἐναπεποιημένης ἐν ἐλαίῳ εἰς θυσίαν.
\vs{38}Θυϊσκην μίαν δέκα χρυσῶν, πλήρη θυμιάματος.
\vs{39}Μόσχον ἕνα ἐκ βοῶν, κριὸν ἕνα, ἀμνὸν ἕνα ἐνιαύσιον εἰς ὁλοκαύτωμα,
\vs{40}καὶ χίμαρον ἐξ αἰγῶν ἕνα περὶ ἁμαρτίας.
\vs{41}Καὶ εἰς θυσίαν σωτηρίου δαμάλεις δύο, κριοὺς πέντε, τράγους πέντε, ἀμνάδας ἐνιαυσίας πέντε· τοῦτο τὸ δῶρον Σαλαμιὴλ υἱοῦ Σουρισαδαί.

\vs{42}Τῇ ἡμέρᾳ τῇ ἕκτῃ ἄρχων τῶν υἱῶν Γὰδ, Ἐλεισὰφ υἱὸς Ῥαγουήλ.
\vs{43}Τὸ δῶρον αὐτοῦ, τρυβλίον ἀργυροῦν ἕν, τριάκοντα καὶ ἑκατὸν ὁλκὴ αὐτοῦ· φιάλην μίαν ἀργυρᾶν, ἑβδομήκοντα σίκλων κατὰ τὸν σίκλον τὸν ἅγιον· ἀμφότερα πλήρη σεμιδάλεως ἀναπεποιημένης ἐν ἐλαίῳ εἰς θυσίαν.
\vs{44}Θυΐσκην μίαν δέκα χρυσῶν, πλήρη θυμιάματος.
\vs{45}Μόσχον ἕνα ἐκ βοῶν, κριὸν ἕνα, ἀμνὸν ἕνα ἐνιαύσιον εἰς ὁλοκαύτωμα,
\vs{46}καὶ χίμαρον ἐξ αἰγῶν ἕνα περὶ ἁμαρτίας.
\vs{47}Καὶ εἰς θυσίαν σωτηρίου δαμάλεις δύο, κριοὺς πέντε, τράγους πέντε, ἀμνάδας ἐνιαυσίας πέντε· τοῦτο τὸ δῶρον Ἐλισὰφ υἱοῦ Ῥαγουήλ.

\vs{48}Τῇ ἡμέρᾳ τῇ ἑβδόμῃ ἄρχων τῶν υἱῶν Ἐφραὶμ, Ἐλισαμὰ υἱὸς Ἐμιούδ.
\vs{49}Τὸ δῶρον αὐτοῦ, τρυβλίον ἀργυροῦν ἓν, τριάκοντα καὶ ἑκατὸν ὁλκὴ αὐτοῦ· φιάλην μίαν ἀργυρᾶν, ἑβδομήκοντα σίκλων κατὰ τὸν σίκλον τὸν ἅγιον· ἀμφότερα πλήρη σεμιδάλεως ἀναπεποιημένης ἐν ἐλαίῳ εἰς θυσίαν.
\vs{50}Θυΐσκην μίαν δέκα χρυσῶν, πλήρη θυμιάματος.
\vs{51}Μόσχον ἕνα ἐκ βοῶν, κριὸν ἕνα, ἀμνὸν ἕνα ἐνιαύσιον εἰς ὁλοκαύτωμα,
\vs{52}καὶ χίμαρον ἐξ αἰγῶν ἕνα περὶ ἁμαρτίας.
\vs{53}Καὶ εἰς θυσίαν σωτηρίου δαμάλεις δύο, κριοὺς πέντε, τράγους πέντε, ἀμνάδας ἐνιαυσίας πέντε· τοῦτο τὸ δῶρον Ἐλισαμὰ υἱοῦ Ἐμιούδ.

\vs{54}Τῇ ἡμέρᾳ τῇ ὀγδόῃ ἄρχων τῶν υἱῶν Μανασσῆ, Γαμαλιὴλ υἱὸς Φαδασσούρ.
\vs{55}Τὸ δῶρον αὐτοῦ, τρυβλίον ἀργυροῦν ἓν, τριάκοντα καὶ ἑκατὸν ὁλκὴ αὐτοῦ· φιάλην μίαν ἀργυρᾶν, ἑβδομήκοντα σίκλων κατὰ τὸν σίκλον τὸν ἅγιον· ἀμφότερα πλήρη σεμιδάλεως ἀναπεποιημένης ἐν ἐλαίῳ εἰς θυσίαν.
\vs{56}Θυΐσκην μίαν δέκα χρυσῶν, πλήρη θυμιάματος.
\vs{57}Μόσχον ἕνα ἐκ βοῶν, κριὸν ἕνα, ἀμνὸν ἕνα ἐνιαύσιον εἰς ὁλοκαύτωμα,
\vs{58}καὶ χίμαρον ἐξ αἰγῶν ἕνα περὶ ἁμαρτίας·
\vs{59}Καὶ εἰς θυσίαν σωτηρίου δαμάλεις δύο, κριοὺς πέντε, τράγους πέντε, ἀμνάδας ἐνιαυσίας πέντε· τοῦτο τὸ δῶρον Γαμαλιὴλ υἱοὺ Φαδασσούρ.

\vs{60}Τῇ ἡμέρᾳ τῇ ἐνάτῃ ἄρχων τῶν υἱῶν Βενιαμὶν, Ἀβιδὰν υἱὸς Γαδεωνί.
\vs{61}Τὸ δῶρον αὐτοῦ, τρυβλίον ἀργυροῦν ἓν, τριάκοντα καὶ ἑκατὸν ὁλκὴ αὐτοῦ· φιάλην, μίαν ἀργυρᾶν, ἑβδομήκοντα σίκλων κατὰ τὸν σίκλον τὸν ἅγιον· ἀμφότερα πλήρη σεμιδάλεως ἀναπεποιημένης ἐν ἐλαίῳ εἰς θυσίαν·
\vs{62}Θυΐσκην μίαν δέκα χρυσῶν, πλήρη θυμιάματος.
\vs{63}Μόσχον ἕνα ἐκ βοῶν, κριὸν ἕνα, ἀμνὸν ἕνα ἐνιαύσιον εἰς ὁλοκαύτωμα,
\vs{64}καὶ χίμαρον ἐξ αἰγῶν ἕνα περὶ ἁμαρτίας.
\vs{65}Καὶ εἰς θυσίαν σωτηρίου δαμάλεις δύο, κριοὺς πέντε, τράγους πέντε, ἀμνάδας ἐνιαυσίας πέντε· τοῦτο τὸ δῶρον Ἀβιδὰν υἱοῦ Γαδεωνί.

\vs{66}Τῇ ἡμέρᾳ τῇ δεκάτῃ ἄρχων τῶν υἱῶν Δὰν, Ἀχιέζερ υἱὸς Ἀμισαδαί.
\vs{67}Τὸ δῶρον αὐτοῦ, τρυβλίον ἀργυροῦν ἕν, τριάκοντα καὶ ἑκατὸν ὁλκὴ αὐτοῦ· φιάλην μίαν ἀργυρᾶν, ἑβδομήκοντα σίκλων κατὰ τὸν σίκλον τὸν ἅγιον· ἀμφότερα πλήρη σεμιδάλεως ἀναπεποιημένης ἐν ἐλαίῳ εἰς θυσίαν.
\vs{68}Θυΐσκην μίαν δέκα χρυσῶν, πλήρη θυμιάματος.
\vs{69}Μόσχον ἕνα ἐκ βοῶν, κριὸν ἕνα, ἀμνὸν ἕνα ἐνιαύσιον εἰς ὁλοκαύτωμα,
\vs{70}καὶ χίμαρον ἐξ αἰγῶν ἕνα περὶ ἁμαρτίας.
\vs{71}Καὶ εἰς θυσίαν σωτηρίου δαμάλεις δύο, κριοὺς πέντε, τράγους πέντε, ἀμνάδας ἐνιαυσίας πέντε· τοῦτο τὸ δῶρον Ἀχιέζερ υἱοῦ Ἀμισαδαί.

\vs{72}Τῇ ἡμέρᾳ τῇ ἑνδέκατῃ ἄρχων τῶν υἱῶν Ἀσὴρ, Φαγεὴλ υἱὸς Ἐχράν.
\vs{73}Τὸ δῶρον αὐτοῦ, τρυβλίον ἀργυροῦν ἓν, τριάκοντα καὶ ἑκατὸν ὁλκὴ αὐτοῦ· φιάλην μίαν ἀργυρᾶν, ἑβδομήκοντα σίκλων κατὰ τὸν σίκλον τὸν ἅγιον· ἀμφότερα πλήρη σεμιδάλεως ἀναπεποιημένης ἐν ἐλαίῳ εἰς θυσίαν.
\vs{74}Θυΐσκην μίαν δέκα χρυσῶν, πλήρη θυμιάματος.
\vs{75}Μόσχον ἕνα ἐκ βοῶν, κριὸν ἕνα, ἀμνὸν ἐνιαύσιον ἕνα εἰς ὁλοκαύτωμα,
\vs{76}καὶ χίμαρον ἐξ αἰγῶν ἕνα περὶ ἁμαρτίας.
\vs{77}Καὶ εἰς θυσίαν σωτηρίου δαμάλεις δύο, κριοὺς πέντε, τράγους πέντε, ἀμνάδας ἐνιαυσίας πέντε· τοῦτο τὸ δῶρον Φαγεὴλ υἱοῦ Ἐχράν.

\vs{78}Τῇ ἡμέρᾳ τῇ δωδεκάτῃ ἄρχων τῶν υἱῶν Νεφθαλί, Ἀχιρὲ υἱὸς Αἰνάν.
\vs{79}Τὸ δῶρον αὐτοῦ, τρυβλίον ἀργυροῦν ἓν, τριάκοντα καὶ ἑκατὸν ὁλκὴ αὐτοῦ· φιάλην μίαν ἀργυρᾶν, ἑβδομήκοντα σίκλων κατὰ τὸν σίκλον τὸν ἅγιον· ἀμφότερα πλήρη σεμιδάλεως ἀναπεποιημένης ἐν ἐλαίῳ εἰς θυσίαν.
\vs{80}Θυΐσκην μίαν δέκα χρυσῶν, πλήρη θυμιάματος.
\vs{81}Μόσχον ἕνα ἐκ βοῶν, κριὸν ἕνα, ἀμνὸν ἕνα ἐνιαύσιον εἰς ὁλόκαύτωμα,
\vs{82}καὶ χίμαρον ἐξ αἰγῶν ἕνα περὶ ἁμαρτίας.
\vs{83}Καὶ εἰς θυσίαν σωτηρίου δαμάλεις δύο, κριοὺς πέντε, τράγους πέντε, ἀμνάδας ἐνιαυσίας πέντε· τοῦτο τὸ δῶρον Ἀχιρὲ υἱοῦ Αἰνάν.

\vs{84}Οὗτος ὁ ἐγκαινισμὸς τοῦ θυσιαστηρίου ᾗ ἡμέρᾳ ἔχρισεν αὐτὸ, παρὰ τῶν ἀρχόντων τῶν υἱῶν Ἰσραὴλ· τρυβλία ἀργυρᾶ δώδεκα, φιάλαι ἀργυραῖ δώδεκα, θυΐσκαι χρυσαῖ δώδεκα.
\vs{85}Τριάκοντα καὶ ἑκατὸν σίκλων, τὸ τρυβλίον τὸ ἕν, καὶ ἑβδομήκοντα σίκλων ἡ φιάλη ἡ μία· πᾶν τὸ ἀργύριον τῶν σκευῶν, δισχίλιοι καὶ τετρακόσιοι σίκλοι· σίκλοι, ἐν τῷ σίκλῳ τῷ ἁγίῳ.
\vs{86}Θυΐσκαι χρυσαῖ δώδεκα πλήρεις θυμιάματος· πᾶν τὸ χρυσίον τῶν θυϊσκῶν, εἴκοσι καὶ ἑκατὸν χρυσοῖ.
\vs{87}Πᾶσαι αἱ βόες αἱ εἰς ὁλοκαύτωσιν, μόσχοι δώδεκα, κριοὶ δώδεκα, ἀμνοὶ ἐνιαύσιοι δώδεκα, καὶ αἱ θυσίαι αὐτῶν, καὶ αἱ σπονδαὶ αὐτῶν· καὶ χίμαροι ἐξ αἰγῶν δώδεκα περὶ ἁμαρτίας.
\vs{88}Πᾶσαι αἱ βόες εἰς θυσίαν σωτηρίου, δαμάλεις εἰκοσιτέσσαρες, κριοὶ ἑξήκοντα, τράγοι ἑξήκοντα ἐνιαύσιαι, ἀμνάδες ἑξήκοντα ἐνιαύσιοι ἄμωμοι· αὕτη ἡ ἐγκαίνωσις τοῦ θυσιαστηρίου, μετὰ τὸ πληρῶσαι τὰς χεῖρας αὐτοῦ, καὶ μετὰ τὸ χρίσαι αὐτόν.

\vs{89}Ἐν τῷ εἰσπορεύεσθαι Μωυσῆν εἰς τὴν σκηνὴν τοῦ μαρτυρίου λαλῆσαι αὐτῷ, καὶ ἤκουσε τὴν φωνὴν Κυρίου λαλοῦντος πρὸς αὐτὸν ἄνωθεν τοῦ ἱλαστηρίου, ὅ ἐστιν ἐπὶ τῆς κιβωτοῦ τοῦ μαρτυρίου, ἀναμέσον τῶν δύο χερουβίμ· καὶ ἐλάλει πρὸς αὐτόν.

\ch{8}
Καὶ ἐλάλησε Κύριος πρὸς Μωυσῆν, λέγων, λάλησον τῷ Ἀαρὼν,
\vs{2}καὶ ἐρεῖς πρὸς αὐτὸν, ὅταν ἐπιτιθῇς τοὺς λύχνους ἐκ μέρους, κατὰ πρόσωπον τῆς λυχνίας φωτιοῦσιν οἱ ἑπτὰ λύχνοι.
\vs{3}Καὶ ἐποίησεν οὕτως Ἀαρών· ἐκ τοῦ ἑνὸς μέρους κατὰ πρόσωπον τῆς λυχνίας ἐξῆψε τοὺς λύχνους αὐτῆς, καθὰ συνέταξε Κύριος τῷ Μωυσῇ·
\vs{4}Καὶ αὕτη ἡ κατασκευὴ τῆς λυχνίας· στερεὰ, χρυσῆ, ὁ καυλὸς αὐτῆς, καὶ τὰ κρίνα αὐτῆς, στερεὰ ὅλη· κατὰ τὸ εἶδος ὃ ἔδειξε Κύριος τῷ Μωυσῇ, οὕτως ἐποίησε τὴν λυχνίαν.

\vs{5}Καὶ ἐλάλησε Κύριος πρὸς Μωυσῆν λέγων,
\vs{6}λάβε τοὺς Λευίτας ἐκ μέσου υἱῶν Ἰσραὴλ, καὶ ἀφαγνιεῖς αὐτούς.
\vs{7}Καὶ οὕτω ποιήσεις αὐτοῖς τὸν ἁγνισμὸν αὐτῶν· περιῤῥανεῖς αὐτοὺς ὕδωρ ἁγνισμοῦ· καὶ ἐπελεύσεται ξυρὸν ἐπὶ πᾶν τὸ σῶμα αὐτῶν, καὶ πλυνοῦσι τὰ ἱμάτια αὐτῶν, καὶ καθαροὶ ἔσονται.

\vs{8}Καὶ λήμψονται μόσχον ἕνα ἐκ βοῶν, καὶ τούτου θυσίαν σεμίδαλιν ἀναπεποιημένην ἐν ἐλαίῳ· καὶ μόσχον ἐνιαύσιον ἐκ βοῶν λήψῃ περὶ ἁμαρτίας.
\vs{9}Καὶ προσάξεις τοὺς Λευίτας ἔναντι τῆς σκηνῆς τοῦ μαρτυρίου· καὶ συνάξεις πᾶσαν συναγωγὴν υἱῶν Ἰσραήλ·
\vs{10}Καὶ προσάξεις τοὺς Λευίτας ἔναντι Κυρίου, καὶ ἐπιθήσουσιν οἱ υἱοὶ Ἰσραὴλ τὰς χεῖρας αὐτῶν ἐπὶ τοὺς Λευίτας·
\vs{11}Καὶ ἀφοριεῖ Ἀαρὼν τοὺς Λευίτας ἀπόδομα ἔναντι Κυρίου παρὰ τῶν υἱῶν Ἰσραήλ· καὶ ἔσονται ὥστε ἐργάζεσθαι τὰ ἔργα Κυρίου.
\vs{12}Οἱ δὲ Λευῖται ἐπιθήσονσι τὰς χεῖρας ἐπὶ τὰς κεφαλὰς τῶν μόσχων· καὶ ποιήσεις τὸν ἕνα περὶ ἁμαρτίας, καὶ τὸν ἕνα εἰς ὁλοκαύτωμα Κυρίῳ, ἐξιλάσασθαι περὶ αὐτῶν.

\vs{13}Καὶ στήσεις τοὺς Λευίτας ἔναντι Κυρίου, καὶ ἔναντι Ἀαρὼν, καὶ ἔναντι τῶν υἱῶν αὐτοῦ, καὶ ἀποδώσεις αὐτοὺς ἀπόδομα ἔναντι Κυρίου·
\vs{14}Καὶ διαστελεῖς τοὺς Λευίτας ἐκ μέσου υἱῶν Ἰσραήλ· καὶ ἔσονταί μοι.
\vs{15}Καὶ μετὰ ταῦτα εἰσελεύσονται οἱ Λευῖται ἐργάζεσθαι τὰ ἔργα τῆς σκηνῆς τοῦ μαρτυρίον· καὶ καθαριεῖς αὐτοὺς, καὶ ἀποδώσεις αὐτοὺς ἔναντι Κυρίου.
\vs{16}ὅτι ἀπόδομα ἀποδεδομένοι οὗτοί μοι εἰσὶν ἐκ μέσου υἱῶν Ἰσραήλ· ἀντὶ τῶν διανοιγόντων πᾶσαν μήτραν πρωτοτόκων πάντων ἐκ τῶν υἱῶν Ἰσραὴλ εἴληφα αὐτοὺς ἐμοί.
\vs{17}Ὅτι ἐμοὶ πᾶν πρωτότοκον ἐν υἱοῖς Ἰσραὴλ ἀπὸ ἀνθρώπων ἕως κτήνους· ᾗ ἡμέρᾳ ἐπάταξα πᾶν πρωτότοκον ἐν γῇ Αἰγύπτου, ἡγίασα αὐτοὺς ἐμοὶ,
\vs{18}καὶ ἔλαβον τοὺς Λευίτας ἀντὶ παντὸς πρωτοτόκου ἐν υἱοῖς Ἰσραήλ.
\vs{19}Καὶ ἀπέδωκα τοὺς Λευίτας ἀπόδομα δεδομένους Ἀαρὼν καὶ τοῖς υἱοῖς αὐτοῦ ἐκ μέσου υἱῶν Ἰσραὴλ, ἐργάζεσθαι τὰ ἔργα τῶν υἱῶν Ἰσραὴλ ἐν τῇ σκηνῇ τοῦ μαρτυρίου, καὶ ἐξιλάσκεσθαι περὶ τῶν υἱῶν Ἰσραήλ· καὶ οὐκ ἔσται ἐν τοῖς υἱοῖς Ἰσραὴλ προσεγγίζων πρὸς τὰ ἅγια.

\vs{20}Καὶ ἐποίησε Μωυσῆς καὶ Ἀαρὼν καὶ πᾶσα ἡ συναγωγὴ υἱῶν Ἰσραὴλ τοῖς Λευίταις καθὰ ἐνετείλατο Κύριος τῷ Μωυσῇ περὶ τῶν Λευιτῶν, οὕτως ἐποίησαν αὐτοῖς οἱ υἱοὶ Ἰσραήλ.
\vs{21}Καὶ ἡγνίσαντο οἱ Λευῖται, καὶ ἐπλύναντο τὰ ἱμάτια· καὶ ἀπέδωκεν αὐτοὺς Ἀαρὼν ἀπόδομα ἔναντι Κυρίου, καὶ ἐξιλάσατο περὶ αὐτῶν Ἀαρὼν ἀφαγνίσασθαι αὐτούς.
\vs{22}Καὶ μετὰ ταῦτα εἰσῆλθον οἱ Λευῖται λειτουργεῖν τὴν λειτουργίαν αὐτῶν ἐν τῇ σκηνῇ τοῦ μαρτυρίου ἔναντι Ἀαρὼν, καὶ ἔναντι τῶν υἱῶν αὐτοῦ· καθὰ συνέταξε Κύριος τῷ Μωυσῇ περὶ τῶν Λευιτῶν, οὕτως ἐποίησαν αὐτοῖς.

\vs{23}Καὶ ἐλάλησε Κύριος πρὸς Μωυσῆν, λέγων,
\vs{24}τοῦτό ἐστι τὸ περὶ τῶν Λευιτῶν· ἀπὸ πέντε καὶ εἰκοσαετοῦς καὶ ἐπάνω, εἰσελεύσονται ἐνεργεῖν ἐν τῇ σκηνῇ τοῦ μαρτυρίου·
\vs{25}Καὶ ἀπὸ πεντηκονταετοῦς ἀποστήσεται ἀπὸ τῆς λειτουργίας, καὶ οὐκ ἐργᾶται ἔτι.
\vs{26}Καὶ λειτουργήσει ὁ ἀδελφὸς αὑτοῦ ἐν τῇ σκηνῇ τοῦ μαρτυρίου φυλάσσειν φυλακὰς, ἔργα δὲ οὐκ ἐργᾶται· οὕτως ποιήσεις τοῖς Λευίταις ἐν ταῖς φυλακαῖς αὐτῶν.

\ch{9}
Καὶ ἐλάλησε Κύριος πρὸς Μωυσῆν ἐν τῇ ἐρήμῳ Σινᾷ ἐν τῷ ἔτει τῷ δευτέρῳ, ἐξελθόντων αὐτῶν ἐκ γῆς Αἰγύπτου, ἐν τῷ μηνὶ τῷ πρώτῳ, λέγων,
\vs{2}εἶπον, καὶ ποιείτωσαν οἱ υἱοὶ Ἰσραὴλ τὸ πάσχα καθʼ ὥραν αὐτοῦ,
\vs{3}τῇ τεσσαρεσκαιδεκάτῃ ἡμέρᾳ τοῦ μηνὸς τοῦ πρώτου πρὸς ἑσπέραν, ποιήσεις αὐτὸ κατὰ καιρούς· κατὰ τὸν νόμον αὐτοῦ, καὶ κατὰ τὴν σύγκρισιν αὐτοῦ ποιήσεις αὐτό.
\vs{4}Καὶ ἐλάλησε Μωυσῆς τοῖς υἱοῖς Ἰσραὴλ ποιῆσαι τὸ πάσχα ἐναρχομένου τῇ τεσσαρεσκαιδεκάτῃ ἡμέρᾳ τοῦ μηνὸς ἐν τῇ ἐρήμῳ τοῦ Σινᾶ.
\vs{5}καθὰ συνέταξε Κύριος τῷ Μωυσῇ, οὕτως ἐποίησαν οἱ υἱοὶ Ἰσραήλ.

\vs{6}Καὶ παρεγένοντο οἱ ἄνδρες οἳ ἦσαν ἀκάθαρτοι ἐπὶ ψυχῇ ἀνθρώπου, καὶ οὐκ ἠδύναντο ποιῆσαι τὸ πάσχα ἐν τῇ ἡμέρᾳ ἐκείνῃ· καὶ προσῆλθον ἐναντίον Μωυσῆ καὶ Ἀαρὼν ἐν ἐκείνῃ τῇ ἡμέρᾳ·
\vs{7}Καὶ εἶπαν οἱ ἄνδρες ἐκεῖνοι πρὸς αὐτὸν, ἡμεῖς ἀκάθαρτοι ἐπὶ ψυχῇ ἀνθρώπου· μὴ οὖν ὑστερήσωμεν προσενέγκαι τὸ δῶρον Κυρίῳ κατὰ καιρὸν αὐτοῦ ἐν μέσῳ υἱῶν Ἰσραήλ;
\vs{8}Καὶ εἶπε πρὸς αὐτοὺς Μωυσῆς, στῆτε αὐτοῦ, καὶ ἀκούσομαι τί ἐντελεῖται Κύριος περὶ ὑμῶν.
\vs{9}Καὶ ἐλάλησε Κύριος πρὸς Μωυσῆν, λέγων,
\vs{10}λάλησον τοῖς υἱοῖς Ἰσραὴλ, λέγων, ἄνθρωπος ἄνθρωπος, ὃν ἐὰν γένηται ἀκάθαρτος ἐπὶ ψυχῇ, ἀνθρώπου, ἢ ἐν ὁδῷ μακρὰν ὑμῖν, ἢ ἐν ταῖς γενεαῖς ὑμῶν, καὶ ποιήσει τὸ πάσχα Κυρίῳ ἐν τῷ μηνὶ τῷ δευτέρῳ ἐν τῇ τεσσαρεσκαιδεκάτῃ ἡμέρᾳ·
\vs{11}τὸ πρὸς ἑσπέραν ποιήσουσιν αὐτὸ, ἐπʼ ἀζύμων καὶ πικρίδων φάγονται αὐτό.
\vs{12}Οὐ καταλείψουσιν ἀπʼ αὐτοῦ εἰς τὸ πρωῒ, καὶ ὀστοῦν οὐ συντρίψουσιν ἀπʼ αὐτοῦ· κατὰ τὸν νόμον τοῦ πάσχα ποιήσουσιν αὐτό.
\vs{13}Καὶ ἄνθρωπος ὃς ἐὰν καθαρὸς ᾖ, καὶ ἐν ὁδῷ μακρὰν οὐκ ἔστι, καὶ ὑστερήσῃ ποιῆσαι τὸ πάσχα, ἐξολοθρευθήσεται ἡ ψυχὴ ἐκείνη ἐκ τοῦ λαοῦ αὐτῆς, ὅτι τὸ δῶρον Κυρίῳ οὐ προσήνεγκε κατὰ τὸν καιρὸν αὐτοῦ· ἁμαρτίαν αὐτοῦ λήψεται ὁ ἄνθρωπος ἐκεῖνος.
\vs{14}Ἐὰν δὲ προσέλθῃ πρὸς ὑμᾶς προσήλυτος ἐν τῇ γῇ ὑμῶν, καὶ ποιήσῃ τὸ πάσχα Κυρίῳ, κατὰ τὸν νόμον τοῦ πάσχα, καὶ κατὰ τὴν σύνταξιν αὐτοῦ ποιήσει αὐτό· νόμος εἷς ἔσται ὑμῖν καὶ τῷ προσηλύτῳ καὶ τῷ αὐτόχθονι τῆς γῆς.

\vs{15}Καὶ τῇ ἡμέρᾳ ᾗ ἐστάθη ἡ σκηνὴ, ἐκάλυψεν ἡ νεφελη τὴν σκηνὴν, τὸν οἶκον τοῦ μαρτυρίου· καὶ τὸ ἑσπέρας ἦν ἐπὶ τῆς σκηνῆς ὡς εἶδος πυρὸς ἕως πρωΐ.
\vs{16}Οὕτως ἐγίνετο διαπαντός· ἡ νεφέλη ἐκάλυπτεν αὐτὴν ἡμέρας, καὶ εἶδος πυρὸς τὴν νύκτα.
\vs{17}Καὶ ἡνίκα ἀνέβη ἡ νεφέλη ἀπὸ τῆς σκηνῆς, καὶ μετὰ ταῦτα ἀπῇραν οἱ υἱοὶ Ἰσραήλ· καὶ ἐν τῷ τόπῳ οὗ ἂν ἔστη ἡ νεφέλη, ἐκεῖ παρενέβαλον οἱ υἱοὶ Ἰσραήλ.
\vs{18}Διὰ προστάγματος Κυρίου παρεμβαλοῦσιν οἱ υἱοὶ Ἰσραὴλ, καὶ διὰ προστάγματος Κυρίου ἀπαροῦσι· πάσας τὰς ἡμέρας ἐν αἷς σκιάζει ἡ νεφέλη ἐπὶ τῆς σκηνῆς, παρεμβαλοῦσιν οἱ υἱοὶ Ἰσραήλ.
\vs{19}Καὶ ὅταν ἐφέλκηται ἡ νεφέλη ἐπὶ τῆς σκηνῆς ἡμέρας πλείους, καὶ φυλάξονται οἱ υἱοὶ Ἰσραὴλ τὴν φυλακὴν τοῦ Θεοῦ, καὶ οὐ μὴ ἐξάρωσι.
\vs{20}Καὶ ἔσται ὅταν σκεπάζῃ ἡ νεφέλη ἡμέρας ἀριθμῷ ἐπὶ τῆς σκηνῆς, διὰ φωνῆς Κυρίου παρεμβαλοῦσι, καὶ διὰ προστάγματος Κυρίου ἀπαροῦσι.
\vs{21}Καὶ ἔσται ὅταν γένηται ἡ νεφέλη ἀφʼ ἑσπέρας ἕως πρωῒ, καὶ ἀναβῇ ἡ νεφέλη τοπρωῒ, καὶ ἀπαροῦσιν ἡμέρας ἢ νυκτός.
\vs{22}Μηνὸς ἡμέρας πλεοναζούσης τῆς νεφέλης σκιαζούσης ἐπʼ αὐτῆς, παρεμβαλοῦσιν οἱ υἱοὶ Ἰσραὴλ, καὶ οὐ μὴ ἀπάρωσιν.
\vs{23}Ὅτι διὰ προστάγματος Κυρίου ἀπαροῦσι· τὴν φυλακὴν Κυρίου ἐφυλάξαντο διὰ προστάγματος Κυρίου ἐν χειρὶ Μωυσῆ.

\ch{10}
Καὶ ἐλάλησε Κύριος πρὸς Μωυσῆν, λέγων, ποίησον σεαυτῷ δύο σάλπιγγας ἀγυρᾶς·
\vs{2}ἐλατὰς ποιήσεις αὐτάς· καὶ ἔσονταί σοι ἀνακαλεῖν τὴν συναγωγὴν, καὶ ἐξαίρειν τὰς παρεμβολάς.
\vs{3}Καὶ σαλπιεῖς ἐν αὐταῖς, καὶ συναχθήσεται πᾶσα ἡ συναγωγὴ ἐπὶ τὴν θύραν τῆς σκηνῆς τοῦ μαρτυρίου·
\vs{4}Ἐὰν δὲ ἐν μιᾷ σαλπίσωσι, προσελεύσονται πρὸς σὲ πάντες οἱ ἄρχοντες ἀρχηγοὶ Ἰσραήλ.
\vs{5}καὶ σαλπιεῖτε σημασίαν, καὶ ἐξαροῦσιν αἱ παρεμβάλλουσαι αἱ παρεμβαλοῦσαι ἀνατολάς·
\vs{6}καὶ σαλπιεῖτε σημασίαν δευτέραν, καὶ ἐξαροῦσιν αἱ παρεμβολαὶ αἱ παρεμβάλλουσαι Λίβα· καὶ σαλπιεῖτε σημασίαν τρίτην, καὶ ἐξαροῦσιν αἱ παρεμβολαὶ αἱ παρεμβάλλουσαι παρὰ θάλασσαν· καὶ σαλπιεῖτε σημασίαν τετάρτην, καὶ ἐξαροῦσιν αἱ παρεμβολαὶ αἱ παρεμβάλλουσαι πρὸς Βοῤῥᾶν· σημασίᾳ σαλπιοῦσιν ἐν τῇ ἐξάρσει αὐτῶν.
\vs{7}καὶ ὅταν συναγάγητε τὴν συναγωγὴν, σαλπιεῖτε, καὶ οὐ σημασίᾳ.
\vs{8}Καὶ οἱ υἱοὶ Ἀαρὼν οἱ ἱερεῖς σαλπιοῦσι ταῖς σάλπιγξι· καὶ ἔσται ὑμῖν νόμιμον αἰώνιον εἰς τὰς γενεὰς ὑμῶν.
\vs{9}Ἐὰν δὲ ἐξέλθητε εἰς πόλεμον ἐν τῇ γῇ ὑμῶν πρὸς τοὺς ὑπεναντίους τοὺς ἀνθεστηκότας ὑμῖ, καὶ σημανεῖτε ταῖς σάλπιγξιν, καὶ ἀναμνησθήσεσθε ἔναντι Κυρίου, καὶ διασωθήσεσθε ἀπὸ τῶν ἐχθρῶν ὑμῶν.
\vs{10}Καὶ ἐν ταῖς ἡμέραις τῆς εὐφροσύνης ὑμῶν, καὶ ἐν ταῖς ἑορταῖς ὑμῶν, καὶ ἐν ταῖς νουμηνίαις ὑμῶν, σαλπιεῖτε ταῖς σάλπιγξιν ἐπὶ τοῖς ὁλοκαυτώμασι, καὶ ἐπὶ ταῖς θυσίαις τῶν σωτηρίων ὑμῶν· καὶ ἔσται ὑμῖν ἀνάμνησις ἔναντι τοῦ Θεοῦ ὑμῶν· ἐγὼ Κύριος ὁ Θεὸς ὑμῶν.

\vs{11}Καὶ ἐγένετο ἐν τῷ ἐνιαυτῷ τῷ δευτέρῳ ἐν τῷ μηνὶ τῷ δευτέρῳ εἰκάδι τοῦ μηνὸς, ἀνέβη ἡ νεφέλη ἀπὸ τῆς σκηνῆς τοῦ μαρτυρίου.
\vs{12}καὶ ἐξῇραν οἱ υἱοὶ Ἰσραὴλ σὺν ἀπαρτίαις αὐτῶν ἐν τῇ ἐρήμῳ Σινᾶ· καὶ ἔστη ἡ νεφέλη ἐν τῇ ἐρήμῳ τοῦ Φαράν.
\vs{13}Καὶ ἐξῆραν πρῶτοι διὰ φωνῆς Κυρίου ἐν χειρὶ Μωυσῆ.

\vs{14}Καὶ ἐξῇραν τάγμα παρεμβολῆς υἱῶν Ἰούδα πρῶτοι σὺν δυνάμει αὐτῶν· καὶ ἐπὶ τῆς δυνάμεως αὐτῶν, Ναασσὼν υἱὸς Ἀμιναδάβ.
\vs{15}καὶ ἐπὶ τῆς δυνάμεως φυλῆς υἱῶν Ἰσσάχαρ, Ναθαναὴλ υἱὸς Σωγάρ.
\vs{16}καὶ ἐπὶ τῆς δυνάμεως φυλῆς υἱῶν Ζαβουλὼν, Ἐλιὰβ, υἱὸς Χαιλών.
\vs{17}Καὶ καθελοῦσι τὴν σκηνὴν, καὶ ἐξαροῦσιν οἱ υἱοὶ Γεδσὼν, καὶ οἱ υἱοὶ Μεραρὶ, οἱ αἴροντες τὴν σκηνήν.

\vs{18}Καὶ ἐξῇραν τάγμα παρεμβολῆς Ῥουβὴν σὺν δυνάμει αὐτῶν· καὶ ἐπὶ τῆς δυνάμεως αὐτῶν, Ἐλισοὺρ υἱὸς Σεδιούρ·
\vs{19}Καὶ ἐπὶ τῆς δυνάμεως φυλῆς υἱῶν Συμεὼν, Σαλαμιὴλ υἱὸς Σουρισαδαΐ.
\vs{20}Καὶ ἐπὶ τῆς δυνάμεως φυλῆς υἱῶν Γὰδ, Ἐλισὰφ ὁ τοῦ Ῥαγουήλ.
\vs{21}Καὶ ἐξαροῦσιν οἱ υἱοὶ Καὰθ αἴροντες τὰ ἅγια· καὶ στήσουσι τὴν σκηνὴν ἕως παραγένωνται.
\vs{22}Καὶ ἐξαροῦσι τάγμα παρεμβολῆς Ἐφραὶμ σὺν δυνάμει αὐτῶν· καὶ ἐπὶ τῆς δυνάμεως αὐτῶν, Ἐλισαμὰ υἱὸς Σμιούδ.

\vs{23}Καὶ ἐπὶ τῆς δυνάμεως φυλῆς υἱῶν Μανασσῆ, Γαμαλιὴλ ὁ τοῦ Φαδασσούρ.
\vs{24}Καὶ ἐπὶ τῆς δυνάμεως φυλῆς υἱῶν Βενιαμὶν, Ἀβιδὰν ὁ τοῦ Γαδεωνί.
\vs{25}Καὶ ἐξαροῦσι τάγμα παρεμβολῆς υἱῶν Δὰν, ἔσχατοι πασῶν τῶν παρεμβολῶν, σὺν δυνάμει αὐτῶν· καὶ ἐπὶ τῆς δυνάμεως αὐτῶν, Ἀχιέζερ ὁ τοῦ Ἀμισαδαΐ.
\vs{26}Καὶ ἐπὶ τῆς δυνάμεως φυλῆς υἱῶν Ἀσὴρ, Φαγεὴλ υἱὸς Ἐχράν.
\vs{27}Καὶ ἐπὶ τῆς δυνάμεως φυλῆς υἱῶν Νεφθαλὶ, Ἀχιρὲ υἱὸς Αἰνάν.
\vs{28}Αὗται αἱ στρατιαὶ υἱῶν Ἰσραήλ· καὶ ἐξῇραν σὺν δυνάμει αὐτῶν.

\vs{29}Καὶ εἶπε Μωυσῆς τῷ Ὀβὰβ υἱῷ Ῥαγουὴλ τῷ Μαδιανίτῃ τῷ γαμβρῷ Μωυσῆ, ἐξαίρομεν ἡμεῖς εἰς τὸν τόπον ὃν εἶπε Κύριος, τοῦτον δώσω ὑμῖν· δεῦρο μεθʼ ἡμῶν, καὶ εὖ σε ποιήσομεν, ὅτι Κύριος ἐλάλησε καλὰ περὶ Ἰσραήλ.
\vs{30}Καὶ εἶπε πρὸς αὐτόν, οὐ πορεύσομαι, ἀλλὰ εἰς τὴν γῆν μου, καὶ εἰς τὴν γενεάν μου.
\vs{31}Καὶ εἶπε, μὴ ἐγκαταλίπῃς ἡμᾶς, οὗ ἕνεκεν ἦσθα μεθʼ ἡμῶν ἐν τῇ ἐρήμῳ, καὶ ἔσῃ ἐν ἡμῖν πρεσβύτης.
\vs{32}Καὶ ἔσται ἐὰν πορευθῇς μεθʼ ἡμῶν, καὶ ἔσται τὰ ἀγαθὰ ἐκεῖνα ὅσα ἂν ἀγαθοποιήσῃ Κύριος ἡμᾶς, καὶ εὖ σε ποιήσομεν.

\vs{33}Καὶ ἐξῇραν ἐκ τοῦ ὄρους Κυρίου ὁδὸν τριῶν ἡμερῶν· καὶ ἡ κιβωτὸς τῆς διαθήκης Κυρίου προεπορεύετο προτέρα αὐτῶν ὁδὸν τριῶν ἡμερῶν κατασκέψασθαι αὐτοῖς ἀνάπαυσιν.
\vs{34}Καὶ ἐγένετο ἐν τῷ ἐξαίρειν τὴν κιβωτὸν, καὶ εἶπε Μωυσῆς, ἐξεγέρθητι Κύριε, καὶ διασκορπισθήτωσαν οἱ ἐχθροί σου, φυγέτωσαν πάντες οἱ μισοῦντές σε.
\vs{35}Καὶ ἐν τῇ καταπαύσει εἶπεν, ἐπίστρεφε Κύριε χιλιάδας μυριάδας ἐν τῷ Ἰσραήλ.
\vs{36}Καὶ ἡ νεφέλη ἐγένετο σκιάζουσα ἐπʼ αὐτοῖς ἡμέρας, ἐν τῷ ἐξαίρειν αὐτοὺς ἐκ τῆς παρεμβολῆς.

\ch{11}
Καὶ ἦν ὁ λαὸς γογγύζων πονηρὰ ἔναντι Κυρίου· καὶ ἤκουσε Κύριος, καὶ ἐθυμώθη ὀργῇ· καὶ ἐξεκαύθη ἐν αὐτοῖς πῦρ παρὰ Κυρίου, καὶ κατέφαγε μέρος τι τῆς παρεμβολῆς.
\vs{2}Καὶ ἐκέκραξεν ὁ λαὸς πρὸς Μωυσῆν· καὶ ηὔξατο Μωυσῆς πρὸς Κύριον, καὶ ἐκόπασε τὸ πῦρ.
\vs{3}Καὶ ἐκλήθη τὸ ὄνομα τοῦ τόπου ἐκείνου, Ἐμπυρισμός· ὅτι ἐξεκαύθη ἐν αὐτοῖς παρὰ Κυρίου.
\vs{4}Καὶ ὁ ἐπίμικτος ὁ ἐν αὐτοῖς ἐπεθύμησεν ἐπιθυμίαν· καὶ καθίσαντες ἔκλαιον καὶ οἱ υἱοὶ Ἰσραὴλ, καὶ εἶπαν, τίς ἡμᾶς ψωμιεῖ κρέα;
\vs{5}Ἐμνήσθημεν τοὺς ἰχθύας, οὓς ἠσθίομεν ἐν Αἰγύπτῳ δωρεὰν. καὶ τοὺς σικύας, καὶ τοὺς πέπονας, καὶ τὰ πράσα, καὶ τὰ κρόμμυα, καὶ τὰ σκόρδα.
\vs{6}Νυνὶ δὲ ἡ ψυχὴ ἡμῶν κατάξηρος· οὐδὲν πλὴν εἰς τὸ μάννα οἱ ὀφθαλμοὶ ἡμῶν.
\vs{7}Τὸ δὲ μάννα ὡσεὶ σπέρμα κορίου ἐστὶ, καὶ τὸ εἶδος αὐτοῦ εἶδος κρυστάλλου.
\vs{8}Καὶ διεπορεύετο ὁ λαὸς, καὶ συνέλεγον, καὶ ἤληθον αὐτὸ ἐν τῷ μύλῳ, καὶ ἔτριβον ἐν τῇ θυΐᾳ, καὶ ἥψουν αὐτὸ ἐν τῇ χύτρᾳ, καὶ ἐποίουν αὐτὸ ἐνκρυφίας· καὶ ἦν ἡ ἡδονὴ αὐτοῦ ὡσεὶ γεῦμα ἐγκρὶς ἐξ ἐλαίου.
\vs{9}Καὶ ὅταν κατέβη ἡ δρόσος ἐπὶ τὴν παρεμβολὴν νυκτὸς, κατέβαινε τὸ μάννα ἐπʼ αὐτῆς.

\vs{10}Καὶ ἤκουσε Μωυσῆς κλαίοντων αὐτῶν κατά δήμους αὐτῶν, ἕκαστον ἐπὶ τῆς θύρας αὐτοῦ· καὶ ἐθυμώθη ὀργῇ Κύριος σφόδρα· καὶ ἔναντι Μωυσῆ ἦν πονηρόν.
\vs{11}Καὶ εἶπε Μωυσῆς πρὸς Κύριον, ἱνατί ἐκάκωσας τὸν θεράποντά σου, καὶ διατί οὐχ εὕρηκα χάριν ἐναντίον σου, ἐπιθεῖναι τὴν ὁρμὴν τοῦ λαοῦ τούτου ἐπʼ ἐμέ;
\vs{12}Μὴ ἐγὼ ἐν γαστρὶ ἔλαβον πάντα τὸν λαὸν τοῦτον, ἢ ἐγὼ ἔτεκον αὐτούς; ὅτι λέγεις μοι, λάβε αὐτὸν εἰς τὸν κόλπον σου, ὡσεὶ ἄραι τιθηνὸς τὸν θηλάζοντα, εἰς τὴν γῆν ἣν ὤμοσας τοῖς πατράσιν αὐτῶν;
\vs{13}Πόθεν μοι κρέα, δοῦναι παντὶ τῷ λαῷ τούτῳ; ὅτι κλαίουσιν ἐπʼ ἐμοὶ, λέγοντες, δὸς ἡμῖν κρέα, ἵνα φάγωμεν.
\vs{14}Οὐ δυνήσομαι ἐγὼ μόνος φέρειν τὸν λαὸν τοῦτον, ὅτι βαρύτερόν μοι ἐστὶ τὸ ῥῆμα τοῦτο.
\vs{15}Εἰ δʼ οὕτω σὺ ποιεῖς μοι, ἀπόκτεινόν με ἀναιρέσει, εἰ εὕρηκα ἔλεος παρὰ σοὶ, ἵνα μὴ ἴδω τὴν κάκωσίν μου.

\vs{16}Καὶ εἶπε Κύριος πρὸς Μωυσῆς, συνάγαγέ μοι ἑβδομήκοντα ἄνδρας ἀπὸ τῶν πρεσβυτέρων Ἰσραήλ, οὓς αὐτὸς σὺ οἶδας, ὅτι οὗτοί εἰσι πρεσβύτεροι τοῦ λαοῦ καὶ γραμματεῖς αὐτῶς· καὶ ἄξεις αὐτοὺς πρὸς τὴν σκηνὴν τοῦ μαρτυρίου, καὶ στήσονται ἐκεῖ μετὰ σοῦ.
\vs{17}Καὶ καταβήσομαι, καὶ λαλήσω ἐκεῖ μετὰ σοῦ· καὶ ἀφελῶ ἀπὸ τοῦ πνεύματος τοῦ ἐπὶ σοὶ, καὶ ἐπιθήσω ἐπʼ αὐτούς· καὶ συναντιλήψονται μετὰ σοῦ τὴν ὁρμὴν τοῦ λαοῦ, καὶ οὐκ οἴσεις αὐτὸς σὺ μόνος.
\vs{18}Καὶ τῷ λαῷ ἐρεῖς, ἁγνίσασθε εἰς αὔριον, καὶ φάγεσθε κρέα· ὅτι ἐκλαύσατε ἔναντι Κυρίου, λέγοντες, τίς ἡμᾶς ψωμιεῖ κρέα; ὅτι καλὸν ἡμῖν ἐστιν ἐν Αἰγυπτῳ· καὶ δώσει Κύριος ὑμῖν φαγεῖν κρέα, καὶ φάγεσθε κρέα.
\vs{19}Οὐχ ἡμέραν μίαν φάγεσθε, οὐ δὲ δύο, οὐ δὲ πέντε ἡμέρας, οὐ δὲ δέκα ἡμέρας, οὐ δὲ εἴκοσι ἡμέρας,
\vs{20}ἕως μηνὸς ἡμερῶν φάγεσθε, ἕως ἂν ἐξέλθῃ ἐκ τῶν μυκτήρων ὑμῶν· καὶ ἔσται ὑμῖν εἰς χολέραν, ὅτι ἠπειθήσατε Κυρίῳ, ὅς ἐστιν ἐν ὑμῖν, καὶ ἐκλαύσατε ἐναντίον αὐτοῦ, λέγοντες, ἱνατί ἡμῖν ἐξελθεῖν ἐξ Αἰγύπτου;
\vs{21}Καὶ εἶπε Μωυσῆς, ἑξακόσιαι χιλιάδες πεζῶν ὁ λαὸς, ἐν οἷς εἰμι ἐν αὐτοῖς· καὶ σὺ εἶπας, κρέα δώσω αὐτοῖς φαγεῖν, καὶ φάγονται μῆνα ἡμερῶν·
\vs{22}Μὴ πρόβατα καὶ βόες σφαγήσονται αὐτοῖς, καὶ ἀρκέσει αὐτοῖς; ἢ πᾶν τὸ ὄψος τῆς θαλάσσης συναχθήσεται αὐτοῖς, καὶ ἀρκέσει αὐτοῖς;
\vs{23}Καὶ εἶπε Κύριος πρὸς Μωυσῆν, μὴ χεὶρ Κυρίου οὐκ ἐξαρκέσει; ἤδη γνώσῃ εἰ ἐπικαταλήμψεταί σε ὁ λόγος μου ἢ οὔ.

\vs{24}Καὶ ἐξῆλθε Μωυσῆς, καὶ ἐλάλησε πρὸς τὸν λαὸν τὰ ῥήματα Κυρίου· καὶ συνήγαγεν ἑβδομήκοντα ἄνδρας ἀπὸ τῶν πρεσβυτέρων τοῦ λαοῦ, καὶ ἔστησεν αὐτὸς κύκλῳ τῆς σκηνῆς.
\vs{25}Καὶ κατέβη Κύριος ἐν νεφέλῃ, καὶ ἐλάλησε πρὸς αὐτόν· καὶ παρείλατο ἀπὸ τοῦ πνεύματος τοῦ ἐπʼ αὐτῷ, καὶ ἐπέθηκεν ἐπὶ τοὺς ἑβδομήκοντα ἄνδρας τοὺς πρεσβυτέρους· ὡς δὲ ἐπανεπαύσατο πνεῦμα ἐπʼ αὐτοὺς, καὶ ἐπροφήτευσαν, καὶ οὐκ ἔτι προσέθεντο.
\vs{26}Καὶ κατελείφθησαν δύο ἄνδρες ἐν τῇ παρεμβολῇ, ὄνομα τῷ ἑνὶ Ἑλδὰδ, καὶ ὄνομα τῷ δευτέρῳ Μωδάδ· καὶ ἐπανεπαύσατο ἐπʼ αὐτοὺς πνεῦμα· καὶ οὗτοι ἦσαν τῶν καταγεγραμμένων, καὶ οὐκ ἦλθον πρὸς τὴν σκηνήν· καὶ ἐπροφήτευσαν ἐν τῇ παρεμβολῇ.
\vs{27}Καὶ προσδραμὼν ὁ νεανίσκος, ἀπήγγειλε Μωυσῇ· καὶ εἶπε, λέγων, Ἑλδὰδ καὶ Μωδὰδ προφητεύουσιν ἐν τῇ παρεμβολῇ.
\vs{28}Καὶ ἀποκριθεὶς Ἰησοῦς ὁ τοῦ Ναυὴ, ὁ παρεστηκὼς Μωυσῇ, ὁ ἐκλεκτὸς, εἶπε, κύριε Μωυσῆ, κώλυσον αὐτούς.
\vs{29}Καὶ εἶπε Μωυσῆς αὐτῷ, μὴ ζηλοῖς ἐμέ; καὶ τίς δῴη πάντα τὸν λαὸν Κυρίου προφήτας, ὅταν δῷ Κύριος τὸ πνεῦμα αὐτοῦ ἐπʼ αὐτούς;
\vs{30}Καὶ ἀπῆλθε Μωυσῆς εἰς τὴν παρεμβολὴν αὐτὸς καὶ οἱ πρεσβύτεροι Ἰσραήλ.

\vs{31}Καὶ πνεῦμα ἐξῆλθε παρὰ Κυρίου, καὶ ἐξεπέρασεν ὀρτυγομήτραν ἀπὸ τῆς θαλάσσης· καὶ ἐπέβαλεν ἐπὶ τὴν παρεμβολὴν ὁδὸν ἡμέρας ἐντεῦθεν, καὶ ὁδὸν ἡμέρας ἐντεῦθεν, κύκλῳ τῆς παρεμβολῆς, ὡσεὶ δίπηχυ ἀπὸ τῆς γῆς.
\vs{32}Καὶ ἀναστὰς ὁ λαὸς ὅλην τὴν ἡμέραν, καὶ ὅλην τὴν νύκτα, καὶ ὅλην τὴν ἡμέραν τὴν ἐπαύριον, καὶ συνήγαγον τὴν ὀρτυγομήτραν· ὁ τὸ ὀλίγον, συνήγαγε δέκα κόρους· καὶ ἔψυξαν ἑαυτοῖς ψυγμοὺς κύκλῳ τῆς παρεμβολῆς.
\vs{33}Τὰ κρέα ἔτι ἦν ἐν τοῖς ὀδοῦσιν αὐτῶν πρινὴ ἐκλείπειν, καὶ Κύριος ἐθυμώθη εἰς τὸν λαὸν, καὶ ἐπάταξε Κύριος τὸν λαὸν πληγὴν μεγάλην σφόδρα.
\vs{34}Καὶ ἐκλήθη τὸ ὄνομα τοῦ τόπου ἐκείνου, Μνήματα τῆς ἐπιθυμίας· ὅτι ἐκεῖ ἔθαψαν τὸν λαὸν τὸν ἐπιθυμητήν.
\vs{35}Ἀπὸ Μνημάτων ἐπιθυμίας ἐξῇρεν ὁ λαὸς εἰς Ἀσηρώθ· καὶ ἐγένετο ὁ λαὸς ἐν Ἀσηρώθ.

\ch{12}
Καὶ ἐλάλησε Μαριὰμ καὶ Ἀαρὼν κατὰ Μωυσῆ, ἕνεκεν τῆς γυναικὸς τῆς Αἰθιοπίσσης ἣν ἔλαβε Μωυσῆς, ὅτι γυναῖκα Αἰθιόπισσαν ἔλαβε,
\vs{2}καὶ εἶπαν, μὴ Μωυσῇ μόνῳ λελάληκε Κύριος; οὐχὶ καὶ ἡμῖν ἐλάλησε; καὶ ἤκουσε Κύριος.
\vs{3}Καὶ ὁ ἄνθρωπος Μωυσῆς πραῢς σφόδρα παρὰ πάντας τοὺς ἀνθρώπους τοὺς ὄντας ἐπὶ τῆς γῆς.
\vs{4}Καὶ εἶπε Κύριος παραχρῆμα πρὸς Μωυσῆν καὶ Ἀαρὼν καὶ Μαριὰμ, ἐξέλθετε ὑμεῖς οἱ τρεῖς εἰς τὴν σκηνὴν τοῦ μαρτυρίου.
\vs{5}Καὶ ἐξῆλθον οἱ τρεῖν εἰν τὴν σκηνὴν τοῦ μαρτυρίου· καὶ κατέβη Κύριος ἐν στύλῳ νεφέλης, καὶ ἔστη ἐπὶ τῆς θύρας τῆς σκηνῆς τοῦ μαρτυρίου· καὶ ἐκλήθησαν Ἀαρὼν καὶ Μαριάμ· καὶ ἐξήλθοσαν ἀμφότεροι.
\vs{6}Καὶ εἶπε πρὸς αὐτοὺς, ἀκούσατε τῶν λόγων μου· ἐὰν γένηται προφήτης ὑμῶν Κυρίῳ, ἐν ὁράματι αὐτῷ γυωσθήσομαι, καὶ ἐν ὕπνῳ λαλήσω αὐτῷ.
\vs{7}Οὐχ οὕτως ὁ θεράπων μου Μωυσῆς, ἐν ὅλῳ τῷ οἴκῳ μου πιστός ἐστι·
\vs{8}Στόμα κατὰ στόμα λαλήσω αὐτῶ, ἐν εἴδει, καὶ οὐ διʼ αἰνιγμάτων, καὶ τὴν δόξαν Κυρίου εἶδε· καὶ διατί οὐκ ἐφοβήθητε καταλαλῆσαι κατὰ τοῦ θεράποντός μου Μωυσῆ;
\vs{9}Καὶ ὀργὴ θυμοῦ Κυρίου ἐπʼ αὐτοῖς, καὶ ἀπῆλθε.
\vs{10}Καὶ ἡ νεφέλη ἀπέστη ἀπὸ τῆς σκηνῆς· καὶ ἰδοὺ Μαριὰμ λεπρῶσα ὡσεὶ χιών· καὶ ἐπέβλεψεν Ἀαρὼν ἐπὶ Μαριὰμ, καὶ ἰδοὺ λεπρῶσα.
\vs{11}Καὶ εἶπεν Ἀαρὼν πρὸς Μωυσῆν, δέομαι κύριε, μὴ συνεπιθῇ ἡμῖν ἁμαρτίαν, διότι ἠγνοήσαμεν καθʼ ὅτι ἡμάρτομεν·
\vs{12}Μὴ γένηται ὡσεὶ ἶσον θανάτῳ, ὡσεὶ ἔκτρωμα ἐκπορευόμενον ἐκ μήτρας μητρὸς, καὶ κατεσθίει τὸ ἥμισυ τῶν σαρκῶν αὐτῆς.
\vs{13}Καὶ ἐβόησε Μωυσῆς πρὸς Κύριον, λέγων, ὁ Θεὸς δέομαί σου, ἴασαι αὐτήν.
\vs{14}Καὶ εἶπε Κύριος πρὸς Μωυσῆν, εἰ ὁ πατὴρ αὐτῆς πτύων ἐνέπτυσεν εἰς τὸ πρόσωπον αὐτῆς οὐκ ἐντραπήσεται ἑπτὰ ἡμέρας; ἀφορισθήτω ἑπτὰ ἡμέρας ἔξω τῆς παρεμβολῆς, καὶ μετὰ ταῦτα εἰσελεύσεται.

\vs{15}Καὶ ἀφωρίσθη Μαριὰμ ἔξω τῆς παρεμβολῆς ἑπτὰ ἡμέρας· καὶ ὁ λαὸς οὐκ ἐξῇρεν, ἕως ἐκαθαρίσθη Μαριάμ.

\vs{16}Καὶ μετὰ ταῦτα ἐξῇρεν ὁ λαὸς ἐξ Ἀσηρὼθ, καὶ παρενέβαλον ἐν τῇ ἐρήμῳ τοῦ Φαράν.

\ch{13}Καὶ ἐλάλησε Κύριος πρὸς Μωυσῆν, λέγων, ἀπόστειλον σεαυτῷ ἄνδρας,
\vs{2}καὶ κατασκεψάσθωσαν τὴν γῆν τῶν Χαναναίων, ἣν ἐγὼ δίδωμι τοῖς υἱοῖς Ἰσραὴλ εἰς κατάσχεσιν· ἄνδρα ἕνα κατὰ φυλήν, κατὰ δήμους πατριῶν αὐτῶν ἀποστελεῖς αὐτοὺς, πάντα ἀρχηγὸν ἐξ αὐτῶν.

\vs{3}Καὶ ἐξαπέστειλεν αὐτοὺς Μωυσῆς ἐκ τῆς ἐρήμου Φαρὰν διὰ φωνῆς Κυρίου· πάντες ἄνδρες ἀρχηγοὶ υἱῶν Ἰσραὴλ οὗτοι.
\vs{4}Καὶ ταῦτα τὰ ὀνόματα αὐτῶν· τῆς φυλῆς Ῥουβὴν, Σαμουὴλ υἱὸς Ζαχούρ.
\vs{5}Τῆς φυλῆς Συμεὼν, Σαφὰτ υἱὸς Σουρί.
\vs{6}Τῆς φυλῆς Ἰούδα, Χάλεβ υἱὸς Ἰεφοννή.
\vs{7}Τῆς φυλῆς Ἰσσάχαρ, Ἰλαὰλ υἱὸς Ἰωσήφ.
\vs{8}Τῆς φυλῆς Ἐφραὶμ, Αὐσὴ υἱὸς Ναυή.
\vs{9}Τῆς φυλῆς Βενιαμὶν, Φαλτὶ υἱὸς Ῥαφοῦ.
\vs{10}Τῆς φυλῆς Ζαβουλὼν, Γουδιὴλ υἱὸς Σουδί.
\vs{11}Τῆς φυλῆς Ἰωσὴφ τῶν υἱῶν Μανασσῆ, Γαδδὶ υἱὸς Σουσί.
\vs{12}Τῆς φυλῆς Δὰν, Ἀμιὴλ υἱὸς Γαμαλί.
\vs{13}Τῆς φυλῆς Ἀσὴρ, Σαθοὺρ υἱὸς Μιχαήλ.
\vs{14}Τῆς φυλῆς Νεφθαλὶ, Ναβὶ υἱὸς Σαβί.
\vs{15}Τῆς φυλῆς Γὰδ, Γουδιὴλ υἱὸς Μακχί.
\vs{16}Ταῦτα τὰ ὀνόματα τῶν ἀνδρῶν, οὓς ἀπέστειλε Μωυσῆς κατασκέψασθαι τὴν γῆν· καὶ ἐπωνόμασε Μωυσῆς τὸν Αὐσὴ υἱὸν Ναυὴ, Ἰησοῦν.

\vs{17}Καὶ ἀπέστειλεν αὐτοὺς Μωυσῆς κατασκέψασθαι τὴν γῆν Χαναὰν, καὶ εἶπε πρὸς αὐτοὺς, ἀνάβητε ταύτῃ τῇ ἐρήμῳ, καὶ ἀναβήσεσθε εἰς τὸ ὄρος,
\vs{18}καὶ ὄψεσθε τὴν γῆν τίς ἐστι, καὶ τὸν λαὸν τὸν ἐγκαθήμενον ἐπʼ αὐτῆς, εἰ ἰσχυρός ἐστιν ἢ ἀσθενὴς, ἢ ὀλίγοι εἰσὶν ἢ πολλοί.
\vs{19}Καὶ τίς ἡ γῆ εἰς ἣν οὗτοι ἐγκάθηνται ἐπʼ αὐτῆς, ἢ καλή ἐστιν ἢ πονηρά· καὶ τίνες αἱ πόλεις ἃς οὗτοι κατοικοῦσιν ἐν αὐταῖς, εἰ ἐν τειχήρεσιν ἢ ἐν ἀτειχίστοις.
\vs{20}Καὶ τίς ἡ γῆ, ἢ πίων ἢ παρειμένη· εἰ ἔστιν ἐν αὐτῇ δένδρα, ἢ οὔ· καὶ προσκαρτερήσαντες λήψεσθε ἀπὸ τῶν καρπῶν τῆς γῆς. καὶ αἱ ἡμέραι, ἡμέραι ἔαρος, πρόδρομοι σταφυλῆς.

\vs{21}Καὶ ἀναβάντες κατεσκέψαντο τὴν γῆν ἀπὸ τῆς ἐρήμου Σὶν ἕως Ῥοὸβ, εἰσπορευομένων Αἰμάθ.
\vs{22}Καὶ ἀνέβησαν κατὰ τὴν ἔρημον, καὶ ἀπῆλθον ἕως Χεβρὼν, καὶ ἐκεῖ Ἀχιμὰν, καὶ Σεσσὶ, καὶ Θελαμὶ, γενεαὶ Ἐνάχ· καὶ Χεβρὼν ἐπτὰ ἔτεσιν ᾠκοδομήθη πρὸ τοῦ Τανὶν Αἰγύπτου.
\vs{23}Καὶ ἤλθοσαν ἕως φάραγγος βότρυος, καὶ κατεσκέψαντο αὐτήν· καὶ ἔκοψαν ἐκεῖθεν κλῆμα καὶ βότρυν σταφυλῆς ἕνα ἐπʼ αὐτοῦ, καὶ ᾖραν αὐτὸν ἐπʼ ἀναφορεύσι, καὶ ἀπὸ τῶν ῥοῶν, καὶ ἀπὸ τῶν συκῶν.
\vs{24}Καὶ τὸν τόπον ἐκεῖνον ἐπωνόμασαν Φάραγξ βότρυος, διὰ τὸν βότρυν, ὃν ἔκοψαν ἐκεῖθεν οἱ υἱοὶ Ἰσραήλ.
\vs{25}Καὶ ἀπέστρεψαν ἐκεῖθεν κατασκεψάμενοι τὴν γῆν μετὰ τεσσαράκοντα ἡμέρας.

\vs{26}Καὶ πορευθέντες ἦλθον πρὸς Μωυσῆν καὶ Ἀαρὼν καὶ πρὸς πᾶσαν συναγωγὴν υἱῶν Ἰσραὴλ, εἰς τὴν ἔρημον Φαρὰν Κάδης· καὶ ἀπεκρίθησαν αὐτοῖς ῥῆμα καὶ πάσῃ συναγωγῇ, καὶ ἔδειξαν τὸν καρπὸν τῆς γῆς,
\vs{27}καὶ διηγήσαντο αὐτῷ, καὶ εἶπαν, ἤλθαμεν εἰς τὴν γῆν εἰς ἣν ἀπέστειλας ἡμᾶς, γῆν ῥέουσαν γάλα καὶ μέλι· καὶ οὗτος ὁ καρπὸς αὐτῆς.
\vs{28}Ἀλλʼ ἢ ὅτι θρασὺ τὸ ἔθνος τὸ κατοικοῦν ἐπʼ αὐτῆς, καὶ πόλεις ὀχυραὶ τετειχισμέναι μεγάλαι σφόδρα· καὶ τὴν γενεὰν Ἐνὰχ ἑωράκαμεν ἐκεῖ.
\vs{29}Καὶ Ἀμαλὴκ κατοικεῖ ἐν τῇ γῇ τῇ πρὸς Νότον· καὶ ὁ Χετταῖος, καὶ ὁ Εὐαῖος, καὶ ὁ Ἰεβουσαῖος, καὶ ὁ Ἀμοῤῥαῖος κατοικεῖ ἐν τῇ ὀρεινῇ· καὶ ὁ Χαναναῖος κατοικεῖ παρὰ θάλασσαν, καὶ παρὰ τὸν Ἰορδάνην ποταμόν.
\vs{30}Καὶ κατεσιώπησε Χάλεβ τὸν λαὸν πρὸς Μωυσῆν, καὶ εἶπεν αὐτῷ, οὐχὶ, ἀλλὰ ἀναβάντες ἀναβησόμεθα, καὶ κατακληρονομήσομεν αὐτὴν, ὅτι δυνατοὶ δυνησόμεθα πρὸς αὐτούς.
\vs{31}Καὶ οἱ ἄνθρωποι οἱ συναναβάντες μετʼ αὐτοῦ, εἶπαν, οὐκ ἀναβαίνομεν, ὅτι οὐ μὴ δυνώμεθα ἀναβῆναι πρὸς τὸ ἔθνος, ὅτι ἰσχυρότερον ἡμῶν ἐστι μᾶλλον.
\vs{32}Καὶ ἐξήνεγκαν ἔκστασιν τῆς γῆς ἣν κατεσκέψαντο αὐτὴν πρὸς τοὺς υἱοὺς Ἰσραὴλ, λέγοντες, τὴν γῆς ἣν παρήλθομεν αὐτὴν κατασκέψασθαι, γῆ κατέσθουσα τοὺς κατοικοῦντας ἐπʼ αὐτῆς ἐστι· καὶ πᾶς ὁ λαὸς ὃν ἑωράκαμεν ἐν αὐτῇ, ἄνδρες ὑπερμήκεις.
\vs{33}Καὶ ἐκεῖ ἑωράκαμεν τοὺς γίγαντας, καὶ ἦμεν ἐνώπιον αὐτῶν ὡσεὶ ἀκρίδες· ἀλλὰ καὶ οὕτως ἦμεν ἐνώπιον αὐτῶν.

\ch{14}
Καὶ ἀναλαβοῦσα πᾶσα ἡ συναγωγὴ, ἐνέδωκε φωνήν· καὶ ἔκλαιεν ὁ λαὸς ὅλην τὴν νύκτα ἐκείνην.
\vs{2}Καὶ διεγόγγυζον ἐπὶ Μωυσῆν καὶ Ἀαρὼν πάντες οἱ υἱοὶ Ἰσραήλ· καὶ εἶπαν πρὸς αὐτοὺς πᾶσα ἡ συναγωγὴ,
\vs{3}Ὄφελον ἀπεθάνομεν ἐν γῇ Αἰγύπτῳ, ἢ ἐν τῇ ἐρήμῳ ταύτῃ, εἰ ἀπεθάνομεν· καὶ ἱνατί Κύριος εἰσάγει ἡμᾶς εἰς τὴν γῆν ταύτην πεσεῖν ἐν πολέμῳ; αἱ γυναῖκες ἡμῶν καὶ τὰ παιδία ἔσονται εἰς διαρπαγήν· νῦν οὖν βέλτιόν ἐστιν ἀποστραφῆναι εἰς Αἴγυπτον.
\vs{4}Καὶ εἶπαν ἕτερος τῷ ἑτέρῳ, δῶμεν ἀρχηγὸν, καὶ ἀποστρέψωμεν εἰς Αἴγυπτον.
\vs{5}Καὶ ἔπεσε Μωυσῆς καὶ Ἀαρὼν ἐπὶ πρόσωπον ἐναντίον πάσης συναγωγῆς υἱῶν Ἰσραήλ.

\vs{6}Ἰησοῦς δὲ ὁ τοῦ Ναυὴ, καὶ Χάλεβ ὁ τοῦ Ἰεφοννὴ τῶν κατασκεψαμένων τὴν γῆν, διέῤῥηξαν τὰ ἱμάτια αὐτῶν,
\vs{7}καὶ εἶπαν πρὸς πᾶσαν συναγωγὴν υἱῶν Ἰσραὴλ, λέγοντες, ἡ γῆ ἣν κατεσκεψάμεθα αὐτὴν, ἀγαθή ἐστι σφόδρα σφόδρα.
\vs{8}Εἰ αἱρετίζει ἡμᾶς Κύριος, εἰσάξει ἡμᾶς εἰς τὴν γῆν ταύτην, καὶ δώσει αὐτὴν ἡμῖν· γῆ ἥτις ἐστὶ ῥέουσα γάλα καὶ μέλι.
\vs{9}Ἀλλὰ ἀπὸ τοῦ Κυρίου μὴ ἀποστάται γίνεσθε· ὑμεῖς δὲ μὴ φοβηθῆτε τὸν λαὸν τῆς γῆς, ὅτι κατάβρωμα ὑμῖν ἐστιν· ἀφέστηκε γὰρ ὁ καιρὸς ἀπʼ αὐτῶν· ὁ δὲ Κύριος ἐν ἡμῖν· μὴ φοβηθήτε αὐτούς.

\vs{10}Καὶ εἶπε πᾶσα ἡ συναγωγὴ καταλιθοβολῆσαι αὐτοὺς ἐν λίθοις· καὶ ἡ δόξα Κυρίου ὤφθη ἐν τῇ νεφέλῃ ἐπὶ τῆς σκηνῆς τοῦ μαρτυρίου πᾶσι τοῖς υἱοῖς Ἰσραήλ.
\vs{11}Καὶ εἶπε Κύριος πρὸς Μωυσῆν, ἕως τίνος παροξύνει με ὁ λαὸς οὗτος; καὶ ἕως τίνος οὐ πιστεύουσί μοι ἐπὶ πᾶσι τοῖς σημείοις, οἷς ἐποίησα ἐν αὐτοῖς;
\vs{12}Πατάξω αὐτοὺς θανάτῳ, καὶ ἀπολῶ αὐτούς· καὶ ποιήσω σε καὶ τὸν οἶκον τοῦ πατρός σου εἰς ἔθνος μέγα, καὶ πολὺ μᾶλλον ἢ τοῦτο.
\vs{13}Καὶ εἶπε Μωυσῆς πρὸς Κύριον, καὶ ἀκούσεται Αἴγυπτος, ὅτι ἀνήγαγες τῇ ἰσχύϊ σου τὸν λαὸν τοῦτον ἐξ αὐτῶν.
\vs{14}Ἀλλὰ καὶ πάντες οἱ κατοικοῦντες ἐπὶ τῆς γῆς ταύτης ἀκηκόασιν, ὅτι σὺ εἶ Κύριος ἐν τῷ λαῷ τούτῷ, ὅστις ὀφθαλμοῖς κατʼ ὀφθαλμοὺς ὀπτάζῃ Κύριε, καὶ ἡ νεφέλη σου ἐφέστηκεν ἐπʼ αὐτῶν, καὶ ἐν στύλῳ νεφέλης σὺ πορεύῃ πρότερος αὐτῶν τὴν ἡμέραν, καὶ ἐν στύλῳ πυρὸς τὴν νύκτα.
\vs{15}Καὶ ἐκτρίψεις τὸν λαὸν τοῦτον ὡσεὶ ἄνθρωπον ἕνα· καὶ ἐροῦσι τὰ ἔθνη ὅσοι ἀκηκόασι τὸ ὄνομά σου, λέγοντες,
\vs{16}παρὰ τὸ μὴ δύνασθαι Κύριον εἰσαγαγεῖν τὸν λαὸν τοῦτον εἰς τὴν γῆν ἣν ὤμοσεν αὐτοῖς, κατέστρωσεν αὐτοὺς ἐν τῇ ἐρήμῳ.
\vs{17}Καὶ νῦν ὑψωθήτω ἡ ἰσχύς σου Κύριε, ὃν τρόπον εἶπας, λέγων,
\vs{18}Κύριος μακρόθυμος, καὶ πολυέλεος, καὶ ἀληθινὸς, ἀφαιρῶν ἀνομίας καὶ ἀδικίας καὶ ἁμαρτίας, καὶ καθαρισμῷ οὐ καθαριεῖ τὸν ἔνοχον, ἀποδιδοὺς ἁμαρτίας πατέρων ἐπὶ τέκνα ἕως τρίτης καὶ τετάρτης γενεᾶς.
\vs{19}Ἄφες τὴν ἁμαρτίαν τῷ λαῷ τούτῳ κατὰ τὸ μέγα ἔλεός σου, καθάπερ ἵλεως ἐγένου αὐτοῖς ἀπʼ Αἰγύπτου ἕως τοῦ νῦν.

\vs{20}Καὶ εἶπε Κύριος πρὸς Μωυσῆν, ἵλεως αὐτοῖς εἰμι κατὰ τὸ ῥῆμά σου.
\vs{21}Ἀλλὰ ζῶ ἐγὼ καὶ ζῶν τὸ ὄνομά μου, καὶ ἐμπλήσει ἡ δόξα Κυρίου πᾶσαν τὴν γῆν.
\vs{22}Ὅτι πάντες οἱ ἄνδρες οἱ ὁρώντες τὴν δόξαν μου, καὶ τὰ σημεῖα ἃ ἐποίησα ἐν Αἰγύπτῳ, καὶ ἐν τῇ ἐρήμῳ, καὶ ἐπείρασάν με τοῦτο δέκατον, καὶ οὐκ εἰσήκουσαν τῆς φωνῆς μου,
\vs{23}ἦ μὴν οὐκ ὄψονται τὴν γῆν, ἣν ὤμοσα τοῖς πατράσιν αὐτῶν· ἀλλʼ ἢ τὰ τέκνα αὐτῶν ἅ ἐστι μετʼ ἐμοῦ ὧδε, ὅσοι οὐκ οἴδασιν ἀγαθὸν οὐδὲ κακὸν, πᾶς νεώτερος ἄπειρος, τούτοις δώσω τὴν γῆν· πάντες δὲ οἱ παροξύναντές με, οὐκ ὄψονται αὐτήν.
\vs{24}Ὁ δὲ παῖς μου Χάλεβ, ὅτι πνεῦμα ἕτερον ἐν αὐτῷ, καὶ ἐπηκολούθησέ μοι, εἰσάξω αὐτὸν εἰς τὴν γῆν εἰς ἣν εἰσῆλθεν ἐκεῖ, καὶ τὸ σπέρμα αὐτοῦ κληρονομήσει αὐτήν.
\vs{25}Ὁ δὲ Ἀμαλὴκ καὶ ὁ Χαναναῖος κατοικοῦσιν ἐν τῇ κοιλάδι· αὔριον ἐπιστράφητε καὶ ἀπάρατε ὑμεῖς εἰς τὴν ἔρημον, ὁδὸν θάλασσαν ἐρυθρᾶν.

\vs{26}Καὶ εἶπε Κύριος πρὸς Μωυσῆν καὶ Ἀαρὼν, λέγων,
\vs{27}ἕως τίνος τὴν συναγωγὴν τὴν πονηρὰν ταύτην; ἅ αὐτοὶ γογγύζουσιν ἐναντίον μου, τὴν γόγγυσιν τῶν υἱῶν Ἰσραὴλ, ἣν ἐγόγγυσαν περὶ ὑμῶν, ἀκήκοα.
\vs{28}Εἶπον αὐτοῖς, ζῶ ἐγὼ, λέγει Κύριος· ἦ μὴν ὃν τρόπον λελαλήκατε εἰς τὰ ὦτά μου, οὕτω ποιήσω ὑμῖν.
\vs{29}Ἐν τῇ ἐρήμῳ ταύτῃ πεσεῖται τὰ κῶλα ὑμῶν· καὶ πᾶσα ἡ ἐπισκοπὴ ὑμῶν, καὶ οἱ κατηριθμημένοι ὑμῶν ἀπὸ εἰκοσαετοῦς καὶ ἐπάνω, ὅσοι ἐγόγγυσαν ἐπʼ ἐμοί·
\vs{30}εἰ ὑμεῖς εἰσελεύσεσθε εἰς τὴν γῆν ἐφʼ ἣν ἐξέτεινα τὴν χεῖρά μου κατασκηνῶσαι ὑμᾶς ἐπʼ αὐτῆς· ἀλλʼ ἢ Χάλεβ υἱὸς Ἰεφοννὴ, καὶ Ἰησοῦς ὁ τοῦ Ναυὴ.
\vs{31}Καὶ τὰ παιδία, ἃ εἴπατε ἐν διαρπαγῇ ἔσεσθαι, εἰσάξω αὐτοὺς εἰς τὴν γῆν· καὶ κληρονομήσουσι τὴν γῆν, ἣν ὑμεῖς ἁπέστητε ἀπʼ αὐτῆς.
\vs{32}Καὶ τὰ κῶλα ὑμῶν πεσεῖται ἐν τῇ ἐρήμῳ ταύτῃ.
\vs{33}Οἱ δὲ υἱοὶ ὑμῶν ἔσονται νεμόμενοι ἐν τῇ ἐρήμῳ τεσσαράκοντα ἔτη· καὶ ἀνοίσουσι τὴν πορνείαν ὑμῶν, ἕως ἂν ἀναλωθῇ τὰ κῶλα ὑμῶν ἐν τῇ ἐρήμῳ,
\vs{34}κατὰ τὸν ἀριθμὸν τῶν ἡμερῶν ὅσας κατεσκέψασθε τὴν γῆν, τεσσαράκοντα ἡμέρας, ἡμέραν τοῦ ἐνιαυτοῦ, λήμψεσθε τὰς ἁμαρτίας ὑμῶν τεσσαράκοντα ἔτη· καὶ γνώσεσθε τὸν θυμὸν τῆς ὀργῆς μου.
\vs{35}Ἐγὼ Κύριος ἐλάλησα, ἦ μὴν οὕτω ποιήσω τῇ συναγωγῇ τῇ πονηρᾷ ταύτῃ, τῇ ἐπισυνισταμένῃ ἐπʼ ἐμέ· ἐν τῇ ἐρήμῳ ταύτη ἐξαναλωθήσονται, καὶ ἐκεῖ ἀποθανοῦνται.

\vs{36}Καὶ οἱ ἄνθρωποι, οὓς ἀπέστειλεν Μωυσῆς κατασκέψασθαι τὴν γῆν, καὶ παραγενηθέντες διεγόγγυσαν κατʼ αὐτῆς πρὸς τὴν συναγωγὴν ἐξενέγκαι ῥήματα πονηρὰ περὶ τῆς γῆς,
\vs{37}καὶ ἀπέθανον οἱ ἄνθρωποι οἱ κατείπαντες πονηρὰ κατὰ τῆς γῆς ἐν τῇ πληγῇ ἔναντι Κυρίου.
\vs{38}Καὶ Ἰησοῦς υἱὸς Ναυὴ καὶ Χάλεβ υἱὸς Ἰεφοννὴ ἔζησαν ἀπὸ τῶν ἀνθρώπων ἐκείνων τῶν πεπορευμένων κατασκέψασθαι τὴν γῆν.
\vs{39}Καὶ ἐλάλησε Μωυσῆς τὰ ῥήματα ταῦτα πρὸς πάντας υἱοὺς Ἰσραήλ· καὶ ἐπενθησεν ὁ λαὸς σφόδρα.

\vs{40}Καὶ ὀρθρίσαντες τοπρωῒ ἀνέβησαν εἰς τὴν κορυφὴν τοῦ ὄρους, λέγοντες, ἰδοὺ, οἵδε ἡμεῖς ἀναβησόμεθα εἰς τὸν τόπον ὃν εἶπε Κύριος, ὅτι ἡμάρτομεν.
\vs{41}Καὶ εἶπε Μωυσῆς, ἱνατί ὑμεῖς παραβαίνετε τὸ ῥῆμα Κυρίου; οὐκ εὔοδα ἔσται ὑμῖν.
\vs{42}Μὴ ἀναβαίνετε, οὐ γάρ ἐστι Κύριος μεθʼ ὑμῶν· καὶ πεσεῖσθε πρὸ προσώπου τῶν ἐχθρῶν ὑμῶν.
\vs{43}Ὅτι ὁ Ἀμαλὴκ καὶ ὁ Χαναναῖος ἐκεῖ ἔμπροσθεν ὑμῶν, καὶ πεσεῖσθε μαχαίρᾳ, οὗ εἵνεκεν ἀπεστράφητε ἀπειθοῦντες Κυρίῳ, καὶ οὐκ ἔσται Κύριος ἐν ὑμῖν.
\vs{44}Καὶ διαβιασάμενοι, ἀνέβησαν ἐπὶ τὴν κορυφὴν τοῦ ὄρους· ἡ δὲ κιβωτὸς τῆς διαθήκης Κυρίου καὶ Μωυσῆς οὐκ ἐκινήθησαν ἐκ τῆς παρεμβολῆς.
\vs{45}Καὶ κατέβη ὁ Ἀμαλὴκ καὶ ὁ Χαναναῖος ὁ ἐνκαθήμενος ἐν τῷ ὄρει ἐκείνῳ, καὶ ἐτρέψαντο αὐτοὺς, καὶ κατέκοψαν αὐτοὺς ἕως Ἑρμάν· καὶ ἀπεστράφησαν εἰς τὴν παρεμβολήν.

\ch{15}
Καὶ εἶπε Κύριος πρὸς Μωυσῆν, λέγων, λάλησον τοῖς ὑοῖς Ἰσραὴλ,
\vs{2}καὶ ἐρεῖς πρὸς αὐτούς, ὅταν εἰσέλθητε εἰς τὴν γῆν τῆς κατοικήσεως ὑμῶν, ἣν ἐγὼ δίδωμι ὑμῖν,
\vs{3}καὶ ποιήσεις ὁλοκαυτώματα Κυρίῳ, ὁλοκάρπωμα ἢ θυσίαν, μεγαλῦναι εὐχὴν, ἢ καθʼ ἑκούσιον, ἢ ἐν ταῖς ἑορταῖς ὑμῶν ποιῆσαι ὀσμὴν εὐωδίας τῷ Κυρίῳ, εἰ μὲν ἀπὸ τῶν βοῶν ἢ ἀπὸ τῶν προβάτων.
\vs{4}Καὶ προσοίσει ὁ προσφέρων τὸ δῶρον αὐτοῦ Κυρίῳ, θυσίαν σεμιδάλεως δέκατον τοῦ οἰφί ἀναπεποιημένης ἐν ἐλαίῳ ἐν τετάρτῳ τοῦ ἴν.
\vs{5}Καὶ οἶνον εἰς σπονδὴν τὸ τέταρτον τοῦ ἴν ποιήσετε ἐπὶ τῆς ὁλοκαυτώσεως, ἢ ἐπὶ τῆς θυσίας· τῷ ἀμνῷ τῷ ἑνὶ ποιήσεις τοσοῦτο, κάρπωμα ὀσμὴν εὐωδίας τῷ Κυρίῳ.
\vs{6}Καὶ τῷ κριῷ, ὅταν ποιῆτε αὐτὸν εἰς ὁλοκαύτωμα ἢ εἰς θυσίαν, ποιήσεις θυσίαν σεμιδάλεως δύο δέκατα ἀναπεποιημένης ἐν ἐλαίῳ τὸ τρίτον τοῦ ἴν·
\vs{7}Καὶ οἶνον εἰς σπονδὴν τὸ τρίτον τοῦ ἴν προσοίσετε εἰς ὀσμὴν εὐωδίς Κυρίῳ.

\vs{8}Ἐὰν δὲ ποιῆτε ἀπὸ τῶν βοῶν εἰς ὁλοκαύτωσιν ἤ εἰς θυσίαν μεγαλῦναι εὐχὴν, ἢ εἰς σωτήριον Κυρίῳ,
\vs{9}καὶ προσοίσει ἐπὶ τοῦ μόσχου θυσίαν σεμιδάλεως τρία δέκατα ἀναπεποιημένης ἐν ἐλαίῳ ἥμισυ τοῦ ἴν.
\vs{10}Καὶ οἶνον εἰς σπονδὴν τὸ ἥμιου τοῦ ἴν, κάρπωμα ὀσμὴν εὐωδίας Κυρίῳ.

\vs{11}Οὕτω ποιήσεις τῷ μόσχῳ τῷ ἑνὶ, ἢ τῷ κριῷ τῷ ἑνὶ, ἢ τῷ ἀμνῷ τῷ ἑνὶ ἐκ τῶν προβάτων ἢ ἐκ τῶν αἰγῶν·
\vs{12}Κατὰ τὸν ἀριθμὸν ὧν ἐὰν ποιήσντε, οὕτως ποιήσετε τῷ ἑνὶ, κατὰ τὸν ἀριθμὸν αὐτῶν.

\vs{13}Πᾶς ὁ αὐτόχθων ποιήσει οὕτως τοιαῦτα προσενέγκαι καρπώματα εἰς ὀσμὴν εὐωδίας Κυρίῳ.
\vs{14}Ἐὰν δὲ προσήλυτος ἐν ὑμῖν προσγένηται ἐν τῇ γῇ ὑμῶν, ἢ ὂς ἂν γένηται ἐν ὑμῖν ἐν ταῖς γενεαῖς ὑμῶν, καὶ ποιήσει κάρπωμα ὀσμὴν εὐωδίας Κυρίῳ, ὃν τρόπον ποιεῖτε ὑμεῖς, οὕτω ποιήσει ἡ συναγωγὴ Κυρίῳ.

\vs{15}Νόμος εἷς ἔσται ὑμῖν καὶ τοῖς προσηλύτοις τοῖς προσκειμένοις ἐν ὑμῖν, νὀμος αἰώνιος εἰς τὰς γενεὰς ὑμῶν· ὡς ὑμεῖς, καὶ ὁ προσήλυτος ἔσται ἔναντι Κυρίου.
\vs{16}Νόμος εἷς ἔσται καὶ δικαίωμα ἓν ἔσται ὑμῖν καὶ τῷ προσηλύτῳ τῷ προσκειμένῳ ἐν ὑμῖν.

\vs{17}Καὶ ἐλάλησε Κύριος πρὸς Μωυσῆν, λέγων,
\vs{18}λάλησον τοῖς υἱοῖς Ἰσραὴλ, καὶ ἐρεῖς πρὸς αὐτοὺς, ἐν τῷ εἰσπορεύεσθαι ὑμᾶς εἰς τὴν γῆν, εἰς ἣν ἐγὼ εἰσάγω ὑμᾶς ἐκεῖ,
\vs{19}καὶ ἔσται ὅταν ἔσθητε ὑμεῖς ἀπὸ τῶν ἄρτων τῆς γῆς, ἀφελεῖτε ἀφαίρεμα ἀφόρισμα Κυρίῳ, ἀπαρχὴν φυράματον ὑμῶν.
\vs{20}Ἄρτον ἀφοριεῖτε ἀφαίρεμα αὐτό· ὡς ἀφαίρεμα ἀπὸ ἅλω, οὕτως ἀφελεῖτε αὐτὸν,
\vs{21}ἀπαρχὴν φυράματος ὑμῶν, καὶ δώσετε Κυρίῳ ἀφαίρεμα εἰς τὰς γενεὰς ὑμῶν.

\vs{22}Ὅταν δὲ διαμάρτητε καὶ μὴ ποιήσητε πάσας τὰς ἐντολὰς ταύτας, ἅς ἐλάλησε Κύριος πρὸς Μωυσῆν,
\vs{23}καθὰ συνέταξε Κύριος πρὸς ὑμᾶς ἐν χειρὶ Μωυσῆ, ἀπὸ τῆς ἡμέρας ᾗ συνέταξε Κύριος πρὸς ὑμᾶς καὶ ἐπέκεινα εἰς τὰς γενεὰς ὑμῶν,
\vs{24}καὶ ἔσται ἐὰν ἐξ ὀφθαλμῶν τῆς συναγωγῆς γενηθῇ ἀκουσίως, καὶ ποιήσει πᾶσα ἡ συναγωγὴ μόσχον ἕνα ἐκ βοῶν ἄμωμον εἰς ὁλοκαύτωμα εἰς ὀσμὴν εὐωδίας Κυρίῳ, καὶ θυσίαν τούτου καὶ σπονδὴν αὐτοῦ κατὰ τὴν σύνταξιν, καὶ χίμαρον ἐξ αἰγῶν ἕνα περὶ ἁμαρτίας.
\vs{25}Καὶ ἐξιλάσεται ὁ ἱερεὺς περὶ πάσης συναγωγῆς υἱῶν Ἰσραὴλ, καὶ ἀφεθήσεται αὐτοῖς, ὅτι ἀκούσιόν ἐστι· καὶ αὐτοὶ ἤνεγκαν τὸ δῶρον αὐτῶν κάρπωμα Κυρίῳ περὶ τῆς ἁμαρτίας αὐτῶν ἔνατι Κυρίου, περὶ τῶν ἀκουσίων αὐτῶν.
\vs{26}Καὶ ἀφεθήσεται κατὰ πᾶσαν συναγωγην υἱῶν Ἰσραὴλ, καὶ τῷ προσηλύτῳ τῷ προσκειμένῳ πρὸς ὑμᾶς, ὅτι παντὶ τῷ λαῷ ἀκούσιον.

\vs{27}Ἐάν τε ψυχὴ μία ἁμάρτῃ ἀκουσίως, προσάξει αἶγα μίαν ἐνιαυσίαν περὶ ἁμαρτίας.
\vs{28}Καὶ ἐξιλάσεται ὁ ἱερεὺς περὶ τῆς ψυχῆς τῆς ἀκουσιασθείσης, καὶ ἁμαρτούσης ἀκουσίως ἔναντι Κυρίου, ἐξιλάσασθαι περὶ αὐτοῦ.
\vs{29}Τῷ ἐγχωρίῳ ἐν υἱοῖς Ἰσραὴλ, καὶ τῷ προσηλύτῳ τῷ προσκειμένῳ ἐν αὐτοῖς νόμος εἷς ἔσται αὐτοῖς, ὃς ἐὰν ποιήσῃ ἀκουσίως.

\vs{30}Καὶ ψυχὴ ἥτις ποιήσῃ ἐν χειρὶ ὑπερηφανίας ἀπὸ τῶν αὐτοχθόνων ἢ ἀπὸ τῶν προσηλύτων, τὸν Θεὸν οὗτος παροξυνεῖ, ἐξόλοθρευθήσεται ἡ ψυχὴ ἐκείνη ἐκ τοῦ λαοῦ αὐτῆς,
\vs{31}ὅτι τὸ ῥῆμα Κυρίου ἐφαύλισε, καὶ τὰς ἐντολὰς αὐτοῦ διεσκέδασεν· ἐκτρίψει ἐκτριβήσεται ἡ ψυχὴ ἐκείνη, ἡ ἁμαρτία αὐτῆς ἐν αὐτῇ.

\vs{32}Καὶ ἦσαν οἱ υἱοὶ Ἰσραὴλ ἑν τῇ ἐρήμῳ, καὶ εὗρον ἄνδρα συλλέγοντα ξύλα τῇ ἡμέρᾳ τῶν σαββάτων.
\vs{33}Καὶ προσήγαγον αὐτὸν οἱ εὐρόντες συλλέγοντα ξύλα τῇ ἡμέρᾳ τῶν σαββάτων πρὸς Μωυσῆν καὶ Ἀαρὼν, καὶ πρὸς πᾶσαν συναγωγὴν υἱῶν Ἰσραήλ.
\vs{34}Καὶ ἀπέθεντο αὐτὸν εἰς φυλακὴν, οὐ γὰρ συνέκριναν τί ποιήσωσιν αὐτόν.
\vs{35}Καὶ ἐλάλησε Κύριος πρὸς Μωυσῆν, λέγων, θανάτῳ θανατούσθω ὁ ἄνθρωπος· λιθοβολήσατε αὐτὸν λίθοις πᾶσα ἡ συναγωγή.
\vs{36}Καὶ ἐξήγαγον αὐτὸν πᾶσα ἡ συναγωγὴ ἔξω τῆς παρεμβολῆς· καὶ ἐλιθοβόλησεν αὐτὸν πᾶσα ἡ συναγωγὴ λίθοις ἔξω τῆς παρεμβολῆς, καθὰ συνέταξε Κύπιος τῷ Μωυσῇ.

\vs{37}Καὶ εἶπε Κύριος πρὸς Μωυσῆν, λέγων,
\vs{38}λάλησον τοῖς υἱοῖς Ἰσραὴλ, καὶ ἐρεῖς πρὸς αὐτοὺς, καὶ ποιησάτωσαν ἐαυτοῖς κράσπεδα ἐπὶ τὰ πτερύγια τῶν ἱματίων αὐτῶν εἰς τὰς γενεὰς αὐτῶν· καὶ ἐπιθήσετε ἐπὶ τὰ κράσπεδα τῶν πτερυγίων κλῶσμα ὑακίνθινον.
\vs{39}Καὶ ἔσται ὑμῖν ἐν τοῖς κρασπέδοις, καὶ ὄψεσθε αὐτά· καὶ μνησθήσεσθε πασῶν τῶν ἐντολῶν Κυρίου, καὶ ποιήσετε αὐτάς· καὶ οὐ διαστραφήσεσθε ὀπίσω τῶν διανοιῶν ὑμῶν, καὶ τῶν ὀφθαλμῶν ἐν οἷς ὑμεῖς ἐκπορνεύετε ὀπίσω αὐτῶν, ὅπως ἂν μνησθῆτε καὶ ποιήσητε πάσας τὰς ἐντολάς μου,
\vs{40}καὶ ἔσεσθε ἅγιοι τῷ Θεῷ ὑμῶν.
\vs{41}Ἐγὼ Κύριος ὁ Θεὸς ὑμῶν ὁ ἐξαγαγὼν ὑμᾶς ἐκ γῆς Αἰγύπτου, εἶναι ὑμῶν Θεός· ἐγὼ Κύριος ὁ Θεὸς ὑμῶν.

\ch{16}
Καὶ ἐλάλησε Κορὲ υἱὸς Ἰσαὰρ υἱοῦ Καὰθ υἱοῦ Λευί, καὶ Δαθὰν καὶ Ἀβειρὼν υἱοὶ Ἐλιάβ, καὶ Αὒν υἱὸς Φαλὲθ υἱοῦ Ῥουβήν·
\vs{2}καὶ ἀνέστησαν ἔναντι Μωυσῆ, καὶ ἄνδρες τῶν υἱῶν Ἰσραὴλ πεντήκοντα καὶ διακόσιοι, ἀρχηγοὶ συναγωγῆς, σύγκλητοι βουλῆς, καὶ ἄνδρες ὀνομαστοί.
\vs{3}Συνέστησαν ἐπὶ Μωυσῆν καὶ Ἀαρὼν, καὶ εἶπαν, ἐχέτω ὑμῖν ὅτι πᾶσα ἡ συναγωγὴ πάντες ἅγιοι, καὶ ἐν αὐτοῖς Κύριος· καὶ διατί κατανίστασθε ἐπὶ τὴν συναγωγὴν Κυρίου;
\vs{4}Καὶ ἀκούσας Μωυσῆς, ἔπεσεν ἐπὶ πρόσωπον.
\vs{5}Καὶ ἐλάλησε πρὸς Κόρε καὶ πρὸς πᾶσαν αὐτοῦ τὴν συναγωγὴν, λέγων, ἐπέσκεπται καὶ ἔγνω ὁ Θεὸς τοὺς ὄντας αὐτοῦ καὶ τοὺς ἁγίους, καὶ προσηγάγετο πρὸς ἑαυτόν· καὶ οὓς ἐξελέξατο ἑαυτῷ, προσηγάγετο πρὸς ἑαυτόν.
\vs{6}Τοῦτο ποιήσατε· λάβετε ὑμῖν αὐτοῖς πυρεῖα Κορὲ, καὶ πᾶσα ἡ συναγωγὴ αὐτοῦ,
\vs{7}καὶ ἐπίθετε ἐπʼ αὐτὰ πῦρ, καὶ ἐπίθετε ἐπʼ αὐτὰ θυμίαμα ἔναντι Κυρίου αὔριον· καὶ ἔσται ὁ ἀνὴρ ὃν ἐκλέλεκται Κύριος, οὗτος ἅγιος· ἱκανούσθω ὑμῖν υἱοὶ Λευί.
\vs{8}Καὶ εἶπε Μωυσῆς πρὸς Κορὲ, εἰσακούσατέ μου υἱοὶ Λευί.
\vs{9}Μὴ μικρόν ἐστι τοῦτο ὑμῖν, ὅτι διέστειλεν ὁ Θεὸς Ἰσραὴλ ὑμᾶς ἐκ συναγωγῆς Ἰσραὴλ, καὶ προσηγάγετο ὑμᾶς πρὸς ἑαυτὸν λειτουργεῖν τᾶς λειτουργίας τῆς σκηνῆς Κυρίου, καὶ παρίστασθαι ἔναντι τῆς σκηνῆς λατρεύειν αὐτοῖς;
\vs{10}καὶ προσηγάγετό σε καὶ πάντας τοὺς ἀδελφούς σου υἱοὺς Λευὶ μετὰ σοῦ, καὶ ζητεῖτε καὶ ἱερατεύειν;
\vs{11}Οὕτως σὺ καὶ πᾶσα ἡ συναγωγή σου ἡ συνηθροισμένη πρὸς τὸν Θεόν· καὶ Ἀαρὼν τίς ἐστιν, ὅτι διαγογγύζετε κατʼ αὐτοῦ;

\vs{12}Καὶ ἀπέστειλε Μωυσῆς καλέσαι Δαθὰν καὶ Ἀβειρὼν υἱοὺς Ἑλιάβ· καὶ εἶπαν, οὐκ ἀναβαίνομεν.
\vs{13}Μὴ μικρὸν τοῦτο, ὅτι ἀνήγαγες ἡμᾶς εἰς γῆν ῥέουσαν γάλα καὶ μέλι, ἀποκτεῖναι ἡμᾶς ἐν τῇ ἐρήμῳ, ὅτι κατάρχεις ἡμῶν;
\vs{14}Ἀρχων εἶ· καὶ σὺ εἰς γῆν ῥέουσαν γάλα καὶ μέλι εἰσήγαγες ἡμᾶς, καὶ ἔδωκας ἡμῖν κλῆρον ἀγροῦ καὶ ἀμπελῶνας; τοὺς ὀφθαλμοὺς τῶν ἀνθρώπων ἐκείνων ἄν ἐξέκοψας; οὐκ ἀναβαίνομεν.
\vs{15}Καὶ ἐβαρυθύμησε Μωυσῆς σφόδρα, καὶ εἶπε πρὸς Κύριον, μὴ πρόσχῃς εἰς τὴν θυσίαν αὐτῶν· οὐκ ἐπιθύμημα οὐδενὸς αὐτῶν εἴληφα, οὐδὲ ἐκάκωσα οὐδένα αὐτῶν.
\vs{16}Καὶ εἶπε Μωυσῆς πρὸς Κορὲ, ἁγίασον τὴν συναγωγήν σου, καὶ γίνεσθε ἕτοιμοι ἔναντι Κυρίου σὺ καὶ Ἀαρὼν καὶ αὐτοὶ αὔριον·
\vs{17}Καὶ λάβετε ἕκαστος τὸ πυρεῖον αὐτοῦ, καὶ ἐπιθήσετε ἐπʼ αὐτὰ θυμίαμα, καὶ προσάξετε ἔναντι Κυρίου ἕκαστος τὸ πυρεῖον αὐτοῦ, πεντήκοντα καὶ διακόσια πυρεῖα, καὶ σὺ καὶ Ἀαρὼν ἕκαστος τὸ πυρεῖον αὐτοῦ.

\vs{18}Καὶ ἔλαβεν ἕκαστος τὸ πυρεῖον αὐτοῦ, καὶ ἐπέθηκαν ἐπʼ αὐτὰ πῦρ, καὶ ἐπέβαλον ἐπʼ αὐτὰ θυμίαμα· καὶ ἔστησαν παρὰ τὰς θύρας τῆς σκηνῆς τοῦ μαρτυρίου Μωυσῆς καὶ Ἀαρών.
\vs{19}Καὶ ἐπισυνέστησεν ἐπʼ αὐτοὺς Κορὲ τὴν πᾶσαν αὐτοῦ συναγωγὴν παρὰ τὴν θύραν τῆς σκηνῆς τοῦ μαρτυρίου· καὶ ὤφθη ἡ δόξα Κυρίου πάσῃ τῇ συναγωγῇ.
\vs{20}Καὶ ἐλάλησε Κύριος πρὸς Μωυσῆν καὶ Ἀαρὼν, λέγων,
\vs{21}ἀποσχίσθητε ἐκ μέσου τῆς συναγωγῆς ταύτης, καὶ ἐξαναλώσω αὐτοὺς εἰσάπαξ.
\vs{22}Καὶ ἔπεσαν ἐπὶ πρόσωπον αὐτῶν, καὶ εἶπαν, Θεὸς, Θεὸς τῶν πνευμάτων καὶ πάσης σαρκὸς, εἰ ἄνθρωπος εἷς ἥμαρτεν, ἐπὶ πᾶσαν τὴν συναγωγὴν ὀργὴ Κυρίου;
\vs{23}Καὶ ἐλάλησε Κύριος πρὸς Μωυσῆν, λέγων,
\vs{24}λάλησον τῇ συναγωγῇ, λέγων, ἀναχωρήσατε κύκγῳ ἀπὸ τῆς συναγωγῆς Κορὲ.

\vs{25}Καὶ ἀνέστη Μωυσῆς, καὶ ἐπορεύθη πρὸς Δαθὰν καὶ Ἀβειρὼν, καὶ συνεπορεύθησαν μετʼ αὐτοῦ πάντες οἱ πρεσβύτεροι Ἰσραήλ.
\vs{26}Καὶ ἐλάλησε πρὸς τὴν συναγωγὴν, λέγων, ἀποσχίσθητε ἀπὸ τῶν σκηνῶν τῶν ἀνθρώπων τῶν σκληρῶν τούτων, καὶ μὴ ἅπτεσθε ἀπὸ πάντων ὧν ἐστιν αὐτοῖς, μὴ συναπόλησθε ἐν πάσῃ τῇ ἁμαρτίᾳ αὐτῶν.
\vs{27}Καὶ ἀπέστησαν ἀπὸ τῆς σκηνῆς Κορὲ κύκλῳ· καὶ Δαθὰν καὶ Ἀβειρὼν ἐξῆλθον, καὶ εἱστήκεισαν παρὰ τὰς θύρας τῶν σκηνῶν αὐτῶν, καὶ αἱ γυναῖκες αὐτῶν, καὶ τὰ τέκνα αὐτῶν, καὶ ἡ ἀποσκευὴ αὐτῶν.

\vs{28}Καὶ εἶπε Μωυσῆς, ἐν τούτῳ γνώσεσθε ὅτι Κύριος ἀπέστειλέ με ποιῆσαι πάντα τὰ ἔργα ταῦτα, ὅτι οὐκ ἀπʼ ἐμαυτοῦ.
\vs{29}Εἰ κατὰ θάνατον πάντων ἀνθρώπων ἀποθανοῦνται οὗτοι, εἰ καὶ κατʼ ἐπίσκεψιν πάντων ἀνθρώπων ἐπισκοπὴ ἔσται αὐτῶν, οὐχὶ Κύριος ἀπέσταλκέ με.
\vs{30}Ἀλλʼ ἢ ἐν φάσματι δείξει Κύριος, καὶ ἀνοίξασα ἡ γῆ τὸ στόμα αὐτῆς καταπίεται αὐτοὺς, καὶ τοὺς οἴκους αὐτῶν, καὶ τὰς σκηνὰς αὐτῶν, καὶ πάντα ὅσα ἐστὶν αὐτοῖς, καὶ καταβήσονται ζῶντες εἰς ᾅδου, καὶ γνώσεσθε, ὅτι παρώξυναν οἱ ἄνθρωποι οὗτοι τὸν Κύριον.

\vs{31}Ὡς δὲ ἐπαύσατο λαλῶν πάντας τοὺς λόγους τούτους, ἐῤῥάγη ἡ γῆ ὑποκάτω αὐτῶν.
\vs{32}Καὶ ἠνοίχθη ἡ γῆ, καὶ κατέπιεν αὐτοὺς, καὶ τοὺς οἴκους αὐτῶν, καὶ πάντας τοὺς ἀνθρώπους τοὺς ὄντας μετὰ Κορὲ, καὶ τὰ κτήνη αὐτῶν.
\vs{33}Καὶ κατέβησαν αὐτοὶ, καὶ ὅσα ἐστὶν αὐτῶν ζῶντα εἰς ᾅδου, καὶ ἐκάλυψεν αὐτοὺς ἡ γῆ, καὶ ἀπώλοντο ἐκ μέσου τῆς συναγωγῆς.
\vs{34}Καὶ πᾶς Ἰσραὴλ οἱ κύκλῳ αὐτῶν ἔφυγον ἀπὸ τῆς φωνῆς αὐτῶν, ὅτι λέγοντες, μή ποτε καταπίῃ ἡμᾶς ἡ γῆ.
\vs{35}Καὶ πῦρ ἐξῆλθε παρὰ Κυρίου, καὶ κατέφαγε τοὺς πεντήκοντα καὶ διακοσίους ἄνδρας τοὺς προσφέροντας τὸ θυμίαμα.

\ch{17}
Καὶ εἶπε Κύριος πρὸς Μωυσῆν,
\vs{2}καὶ πρὸς Ἐλεάζαρ τὸν υἱὸν Ἀαρὼν τὸν ἱερέα, ἀνέλεσθε τὰ πυρεῖα τὰ χαλκᾶ ἐκ μέσου τῶν κατακεκαυμένων, καὶ τὸ πῦρ τὸ ἀλλότριον τοῦτο σπεῖρον ἐκεῖ, ὅτι ἡγίασαν τὰ πυρεῖα τῶν ἁμαρτωλῶν τούτων ἐν ταῖς ψυχαῖς αὐτῶν,
\vs{3}καὶ ποίησον αὐτὰ λεπίδας ἐλατὰς περίθεμα τῷ θυσιαστηρίῳ, ὅτι προσηνέχθησαν ἔναντι Κυρίου καὶ ἡγιάσθησαν· καὶ ἐγένοντο εἰς σημεῖον τοῖς υἱοῖς Ἰσραήλ.
\vs{4}Καὶ ἔλαβεν Ἐλεάζαρ υἱοὶς Ἀαρὼν τοῦ ἱερέως τὰ πυρεῖα τὰ χαλκᾶ, ὅσα προσήνεγκαν οἱ κατακεκαυμένοι, καὶ προσέθηκαν αὐτὰ περίθεμα τῷ θυσιαστηρίῳ,
\vs{5}μνημόσυνον τοῖς υἱοῖς Ἰσραήλ, ὅπως ἂν μὴ προσέλθῃ μηδὶς ἀλλογενής, ὃς οὐκ ἔστιν ἐκ τοῦ σπέρματος Ἀαρὼν, ἐπιθεῖναι θυμίαμα ἔναντι Κυρίου· καὶ οὐκ ἔσται ὥσπερ Κορὲ, καὶ ἡ ἐπισύστασις αὐτοῦ, καθὰ ἐλάλησε Κύριος ἐν χειρὶ Μωυσῆ αὐτῷ.

\vs{6}Καὶ ἐγόγγυσαν οἱ υἱοὶ Ἰσραὴλ τῇ ἐπαύριον ἐπὶ Μωυσῆν καὶ Ἀαρὼν, λέγοντες, ὑμεῖς ἀπεκτάγκατε τὸν λαὸν Κυρίου.
\vs{7}Καὶ ἐγένετο ἐν τῷ ἐπισυστρέφεσθαι τὴν συναγωγὴν ἐπὶ Μωυσῆν καὶ Ἀαρὼν, καὶ ὥρμησαν ἐπὶ τὴν σκηνὴν τοῦ μαρτυρίου· καὶ τήνδε ἐκάλυψεν αὐτὴν ἡ νεφέλη, καὶ ὤφθη ἡ δόξα Κυρίου.
\vs{8}Καὶ εἰσῆλθε Μωυσῆς καὶ Ἀαρὼν κατὰ πρόσωπον τῆς σκηνῆς τοῦ μαρτυρίου.
\vs{9}Καὶ ἐλάλησε Κύριος πρὸς Μωυσῆν καὶ Ἀαρὼν, λέγων,
\vs{10}ἐκχωρήσατε ἐκ μέσου τῆς συναγωγῆς ταύτης, καὶ ἐξαναλώσω αὐτοὺς εἰσάπαξ· καὶ ἔπεσον ἐπὶ πρόσωπον αὐτῶν.
\vs{11}Καὶ εἶπε Μωυσῆς πρὸς Ἀαρὼν, λάβε τὸ πυρεῖον, καὶ ἐπίθες ἐπʼ αὐτὸ πῦρ ἀπὸ τοῦ θυσιαστηρίου, καὶ ἐπίβαλε ἐπʼ αὐτὸ θυμίαμα, καὶ ἀπένεγκε τοτάχος εἰς τὴν παρεμβολὴν, καὶ ἐξίλασαι περὶ αὐτῶν· ἐξῆλθε γὰρ ὀργὴ ἀπὸ προσώπου Κυρίου, ἦρκται θραύειν τὸν λαόν.
\vs{12}Καὶ ἔλαβεν Ἀαρὼν καθάπερ ἐλάλησεν αὐτῷ Μωυσῆς, καὶ ἔδραμεν εἰς τὴν συναγωγήν· καὶ ἤδη ἐνῆρκτο ἡ θραῦσις ἐν τῷ λαῷ· καὶ ἐπέβαλε τὸ θυμίαμα καὶ ἐξιλάσατο περὶ τοῦ λαοῦ.
\vs{13}Καὶ ἔστη ἀναμέσον τῶν τεθνηκότων καὶ τῶν ζώντων, καὶ ἐκόπασεν ἡ θραῦσις.
\vs{14}Καὶ ἐγένοντο οἱ τεθνηκότες ἐν τῇ θραύσει τεσσαρεσκαίδεκα χιλιάδες καὶ ἑπτακόσιοι, χωρὶς τῶν τεθνηκότων ἕνεκεν Κορέ.
\vs{15}Καὶ ἐπέστρεψεν Ἀαρὼν πρὸς Μωυσῆν ἐπὶ τὴν θύραν τῆς σκηνῆς τοῦ μαρτυρίου, καὶ ἐκόπασεν ἡ θραῦσις.

\vs{16}Καὶ ἐλάλησε Κύριος πρὸς Μωυσῆν, λέγων,
\vs{17}λάλησον τοῖς υἱοῖς Ἰσραὴλ, καὶ λάβε παρʼ αὐτῶν ῥάβδον, ῥάβδον κατʼ οἴκους πατριῶν παρὰ πάντων τῶν ἀρχόντων αὐτῶν, κατʼ οἴκους πατριῶν αὐτῶν, δώδεκα ῥάβδους, καὶ ἑκάστου τὸ ὄνομα αὐτοῦ ἐπὶγραψον ἐπὶ τῆς ῥάβδου.
\vs{18}Καὶ τὸ ὄνομα Ἀαρὼν ἐπίγραψον ἐπὶ τῆς ῥάβδου Λευί· ἔστι γὰρ ῥάβδος μία· κατὰ φυλὴν οἴκου πατριῶν αὐτῶν δώσουσι.
\vs{19}Καὶ θήσεις αὐτὰς ἐν τῇ σκηνῇ τοῦ μαρτυρίου, κατέναντι τοῦ μαρτυρίου, ἐν οἷς γνωσθήσομαί σοι ἐκεῖ.
\vs{20}Καὶ ἔσται ὁ ἄνθρωπος ὃν ἂν ἐκλέξωμαι αὐτὸν, ἡ ῥάβδος αὐτοῦ ἐκβλαστήσει· καὶ περιελῶ ἀπʼ ἐμοῦ τὸν γογγυσμὸν υἱῶν Ἰσραήλ, ἃ αὐτοὶ γογγύζουσιν ἐφʼ ὑμῖν.

\vs{21}Καὶ ἐλάλησε Μωυσῆς τοῖς υἱοῖς Ἰσραήλ· καὶ ἔδωκαν αὐτῷ πάντες οἱ ἄρχοντες αὐτῶν ῥάβδον· τῷ ἄρχοντι τῷ ἑνὶ ῥάβδον κατʼ ἄρχοντα, κατʼ οἴκους πατριῶν αὐτῶν, δώδεκα ῥάβδους· καὶ ἡ ῥάβδος Ἀαρὼν ἀναμέσον τῶν ῥάβδων αὐτῶν.
\vs{22}Καὶ ἀπέθηκε Μωυσῆς τὰς ῥάβδους ἔναντι Κυρίου ἐν τῇ σκηνῇ τοῦ μαρτυρίου.
\vs{23}Καὶ ἐγενετο τῇ ἐπαύριον, καὶ εἰσῆλθε Μωυσῆς καὶ Ἀαρὼν ἐν τῇ σκηνῇ τοῦ μαρτυρίου· καὶ ἰδοὺ ἐβλάστησεν ἡ ῥάβδος Ἀαρὼν εἰς οἶκον Λευί, καὶ ἐξήνεγκε βλαστὸν, καὶ ἐξήνθησεν ἄνθη, καὶ ἐβλάστησε κάρυα.
\vs{24}Καὶ ἐξήνεγκε Μωυσῆς πάσας τὰς ῥάβδους ἀπὸ προσώπου Κυρίου πρὸς πάντας υἱοὺς Ἰσραήλ· καὶ εἶδον, καὶ ἔλαβον ἕκαστος τὴν ῥάβδον αὐτοῦ.

\vs{25}Καὶ εἶπε Κύριος πρὸς Μωυσῆν, ἀπόθες τὴν ῥάβδον Ἀαρὼν ἐνώπιον τῶν μαρτυρίων εἰς διατήρησιν, σημεῖον τοῖς υἱοῖς τῶν ἀνηκόων· καὶ παυσάσθω ὁ γογγυσμὸς αὐτῶν ἀπʼ ἐμοῦ, καὶ οὐ μὴ ἀποθάνωσι.
\vs{26}Καὶ ἐποίησε Μωυσῆς καὶ Ἀαρὼν, καθὰ συνέταξε κύριος τῷ Μωυσῇ, οὕτως ἐποίησαν.
\vs{27}Καὶ εἶπαν οἱ υἱοὶ Ἰσραὴλ πρὸς Μωυσῆν, λέγοντες, ἰδοὺ ἐξανηλώμεθα, ἀπολώλαμεν, παρανηλώμεθα.
\vs{28}Πᾶς ὁ ἁπτόμενος τῆς σκηυῆς Κυρίου, ἀποθνήσκει· ἕως εἰς τέλος ἀποθάνωμεν;

\ch{18}
Καὶ εἶπε Κύριος πρὸς Ἀαρὼν, λέγων, σὺ καὶ οἱ υἱοί σου καὶ ὁ οἶκος τοῦ πατρός σου λήψεσθε τὰς ἀμαρτίας τῶν ἁγίων, καὶ σὺ καὶ οἱ υἱοί σου λήψεσθε τὰς ἁμαρτίας τῆς ἱερατείας ὑμῶν.
\vs{2}Καὶ τοὺς ἀδελφούς σου φυλὴν Λευὶ, δῆμον τοῦ πατρός σου προσαγάγου πρὸς σεαυτὸν, καὶ προστεθήτωσάν σοι, καὶ λειτουργείτωσάν σοι· καὶ σὺ καὶ οἱ υἱοί σου μετὰ σοῦ ἀπέναντι τῆς σκηνῆς τοῦ μαρτυρίου.
\vs{3}Καὶ φυλάξονται τὰς φυλακάς σου, καὶ τὰς φυλακὰς τῆς σκηνῆς. πλὴν πρὸς τὰ σκεύη τὰ ἅγια, καὶ πρὸς τὸ θυσιαστήριον οὐ προσελεύσονται, καὶ οὐκ ἀποθανοῦνται καὶ οὗτοι καὶ ὑμεῖς.
\vs{4}Καὶ προστεθήσονται πρὸς σέ, καὶ φυλάξονται τὰς φυλακὰς τῆς σκηνῆς τοῦ μαρτυρίου, κατὰ πάσας τὰς λειτουργίας τῆς σκηνῆς· καὶ ὁ ἀλλογενὴς οὐ προσελεύσεται πρὸς σέ.
\vs{5}Καὶ φυλάξεσθε τὰς φυλακὰς τῶν ἁγίων, καὶ τὰς φυλακὰς τοῦ θυσιαστηρίου, καὶ οὐκ ἔσται θυμὸς ἐν τοῖς υἱοῖς Ἰσραήλ.
\vs{6}Καὶ ἐγὼ εἴληφα τοὺς ἀδελφοὺς ὑμῶν τοὺς Λευίτας ἐκ μέσου τῶν υἱῶν Ἰσραὴλ δόμα δεδομένον Κυρίῳ, λειτουργεῖν τὰς λειτουργίας τῆς σκηνῆς τοῦ μαρτυρίου.
\vs{7}Καὶ σὺ καὶ οἱ υἱοί σου μετὰ σοῦ διατηρήσετε τὴν ἱερατείαν ὑμῶν, κατὰ πάντα τρόπον τοῦ θυσιαστηρίου, καὶ τὸ ἔνδοθεν τοῦ καταπετάσματος· καὶ λειτουργήσετε τὰς λειτουργίας δόμα τῆς ἱερατείας ὑμῶν· καὶ ὁ ἀλλογενὴς ὁ προσπορευόμενος ἀποθανεῖται.

\vs{8}Καὶ ἐλάλησε Κύριος πρὸς Ἀαρὼν, καὶ ἰδοὺ ἐγὼ δέδωκα ὑμῖν τὴν διατήρησιν τῶν ἀπαρχῶν ἀπὸ πάντων τῶν ἡγιασμένων μοι παρὰ τῶν υἱῶν Ἰσραήλ· σοὶ δέδωκα αὐτὰ εἰς γέρας, καὶ τοῖς υἱοῖς σου μετὰ σὲ νόμιμον αἰώνιον.
\vs{9}Καὶ τοῦτο ἔστω ὑμῖν ἀπὸ τῶν ἡγιασμένων ἁγίων τῶν καρπωμάτων, ἀπὸ πάντων τῶν δώρων αὐτῶν, καὶ ἀπὸ πάντων τῶν θυσιασμάτων αὐτῶν, καὶ ἀπὸ πάσης πλημμελείας αὐτῶν, καὶ ἀπὸ πασῶν τῶν ἁμαρτιῶν αὐτῶν, ὅσα ἀποδιδόασί μοι ἀπὸ πάντων τῶν ἁγίων, σοὶ ἔσται καὶ τοῖς υἱοῖς σου.
\vs{10}Ἐν τῷ ἁγίῳ τῶν ἁγίων φάγεσθε αὐτά· πᾶν ἀρσενικὸν φάγεται αὐτά· σὺ καὶ οἱ υἱοί σου· ἅγια ἔσται σοι.

\vs{11}Καὶ τοῦτο ἔσται ὑμῖν ἀπαρχῶν δομάτων αὐτῶν, ἀπὸ πάντων τῶν ἐπιθεμάτων τῶν υἱῶν Ἰσραήλ· σοὶ δέδωκα αὐτὰ καὶ τοῖς υἱοῖς σου καὶ ταῖς θυγατράσι σου μετὰ σοῦ, νόμιμον αἰώνιον· πᾶς καθαρὸς ἐν τῷ οἴκῳ σου ἔδεται αὐτά.

\vs{12}Πᾶσα ἀπαρχὴ ἐλαίου, καὶ πᾶσα ἀπαρχὴ οἴνου, σίτου ἀπαρχὴ αὐτῶν ὅσα ἂν δῶσι τῷ Κυρίῳ, σοὶ δέδωκα αὐτά.
\vs{13}Τὰ πρωτογεννήματα πάντα ὅσα ἐν τῇ γῇ αὐτῶν, ὅσα ἂν ἐνέγκωσι Κυρίῳ, σοὶ ἔσται· πᾶς καθαρὸς ἐν τῷ οἴκῳ σου ἔδεται αὐτά.

\vs{14}Πᾶν ἀνατεθεματισμένον ἐν υἱοῖς Ἰσραὴλ, σοὶ ἔσται.
\vs{15}Καὶ πᾶν διανοῖγον μήτραν ἀπὸ πάσης σαρκὸς, ὅσα προσφέρουσι Κυρίῳ ἀπὸ ἀνθρώπου ἕως κτήνους, σοὶ ἔσται· ἀλλʼ ἢ λύτροις λυτρωθήσεται τὰ πρωτότοκα τῶν ἀνθρώπων, καὶ τὰ πρωτότοκα τῶν κτηνῶν τῶν ἀκαθάρτων λυτρώσῃ.
\vs{16}Καὶ ἡ λύτρωσις αὐτοῦ, ἀπὸ μηνιαίου· ἡ συντίμησις πέντε σίκλων, κατὰ τὸν σίκλον τὸν ἅγιον εἴκοσι ὀβολοί εἰσι.
\vs{17}Πλὴν πρωτότοκα μόσχων καὶ πρωτότοκα προβάτων, καὶ πρωτότοκα αἰγῶν οὐ λυτρώσῃ· ἅγιά ἐστι· καὶ τὸ αἷμα αὐτῶν προσχεεῖς πρὸς τὸ θυσιαστήριον, καὶ τὸ στέαρ ἀνοίσεις κάρπωμα εἰς ὀσμὴν εὐωδίας Κυρίῳ.

\vs{18}Καὶ τὰ κρέα ἔσται σοι, καθὰ καὶ τὸ στηθύνιον τοῦ ἐπιθέματος· καὶ κατὰ τὸν βραχίονα τὸν δεξιὸν, σοὶ ἔσται.
\vs{19}Πᾶν ἀφαίρεμα τῶν ἁγίων, ὅσα ἐὰν ἀφέλωσιν οἱ υἱοὶ Ἰσραὴλ Κυρίῳ, δέδωκά σοι καὶ τοῖς υἱοῖς σου καὶ ταῖς θυγατράσι σου μετὰ σοῦ, νόμιμον αἰώνιον· διαθήκη ἁλὸς αἰωνίου ἔστιν ἔναντι Κυρίου, σοὶ καὶ τῷ σπέρματί σου μετὰ σέ.

\vs{20}Καὶ ἐλάλησε Κύριος πρὸς Ἀαρὼν, ἐν τῇ γῇ αὐτῶν οὐ κληρονομήσεις, καὶ μερὶς οὐκ ἔσται σοι ἐν αὐτοῖς, ὅτι ἐγὼ μερίς σου καὶ κληρονομία σου ἐν μέσῳ τῶν υἱῶν Ἰσραήλ.

\vs{21}Καὶ τοῖς υἱοῖς Λευὶ ἰδοὺ δέδωκα πᾶν ἐπιδέκατον ἐν Ἰσραὴλ ἐν κλήρῳ ἀντὶ τῶν λειτουργιῶν αὐτῶν, ὅσα αὐτοὶ λειτουργοῦσι λειτουργίαν ἐν τῇ σκηνῇ τοῦ μαρτυρίου.
\vs{22}Καὶ οὐ προσελεύσονται ἔτι οἱ υἱοὶ Ἰσραὴλ εἰς τὴν σκηνὴν τοῦ μαρτυρίου λαβεῖν ἁμαρτίαν θανατηφόρον.
\vs{23}Καὶ λειτουργήσει ὁ Λενίτης αὐτὸς τὴν λειτουργίαν τῆς σκηνῆς τοῦ μαρτυρίου· καὶ αὐτοὶ λήψονται τὰ ἁμαρτήματα αὐτῶν, νόμιμον αἰώνιον εἰς τὰς γενεὰς αὐτῶν· καὶ ἐν μέσῳ υἱῶν Ἰσραὴλ οὐ κληρονομήσουσι κληρονομίαν.
\vs{24}Ὅτι τὰ ἐπιδέκατα τῶν υἱῶν Ἰσραὴλ ὅσα ἐὰν ἀφορίσωσι Κυρίῳ, ἀφαίρεμα δέδωκα τοῖς Λευίταις ἐν κλήρῳ· διὰ τοῦτο εἴρηκα αὐτοῖς, ἐν μέσῳ υἱῶν Ἰσραὴλ οὐ κληρονομήσουσι κλῆρον.

\vs{25}Καὶ ἐλάλησε Κύριος πρὸς Μωυσῆν, λέγων,
\vs{26}καὶ τοῖς Λενίταις λαλήσεις, καὶ ἐρεῖς πρὸς αὐτοὺς, ἐὰν λάβητε παρὰ τῶν υἱῶν Ἰσραὴλ τὸ ἐπιδέκατον, ὃ δέδωκα ὑμῖν παρʼ αὐτῶν ἐν κλήρῳ, καὶ ἀφελεῖτε ὑμεῖς ἀπʼ αὐτοῦ ἀφαίρεμα Κυρίῳ, ἐπιδέκατον ἀπὸ τοῦ ἐπιδεκάτου.
\vs{27}Καὶ λογισθήσεται ὑμῖν τὰ ἀφαιρέματα ὑμῶν ὡς σῖτος ἀπὸ ἅλω, καὶ ἀφαίρεμα ἀπὸ ληνοῦ.
\vs{28}Οὕτως ἀφελεῖτε αὐτοὺς καὶ ὑμεῖς ἀπὸ πάντων τῶν ἀφαιρεμάτων Κυρίου ἀπὸ πάντων τῶν ἐπιδεκάτων ὑμῶν, ὅσα ἐὰν λάβητε παρὰ τῶν υἱῶν Ἰσραήλ· καὶ δώσετε ἀπʼ αὐτῶν ἀφαίρεμα Κυρίῳ Ἀαρὼν τῷ ἱερεῖ.
\vs{29}Ἀπὸ πάντων τῶν δομάτων ὑμῶν ἀφελεῖτε ἀφαιρεμα Κυρίῳ, ἢ ἀπὸ πάντων τῶν ἀπαρχῶν τὸ ἡγιασμένον ἀπʼ αὐτοῦ.
\vs{30}Καὶ ἐρεῖς πρὸς αὐτοὺς, ὅταν ἀφαιρῆτε τὴν ἀπαρχὴν ἀπʼ αὐτοῦ, καὶ λογισθήσεται τοῖς Λευίταις ὡς γέννημα ἀπὸ ἅλω, καὶ ὡς γέννημα ἀπὸ ληνοῦ.
\vs{31}Καὶ ἔδεσθε αὐτὸ ἐν παντὶ τόπῳ ὑμεῖς καὶ οἱ οἶκοι ὑμῶν, ὅτι μισθὸς οὗτος ὑμῖν ἐστιν ἀντὶ τῶν λειτουργιῶν ὑμῶν τῶν ἐν τῇ σκηνῇ τοῦ μαρτυρίου.
\vs{32}Καὶ οὐ λήψεσθε διʼ αὐτὸ ἁμαρτίαν, ὅτι ἂν ἀφαιρῆτε τὴν ἀπαρχὴν ἀπʼ αὐτοῦ· καὶ τὰ ἅγια τῶν υἱῶν Ἰσραὴλ οὐ βεβηλώσετε, ἵνα μὴ ἀποθάνητε.

\ch{19}
Καὶ ἐλάλησε Κύριος πρὸς Μωυσῆν καὶ Ἀαρὼν, λέγων,
\vs{2}αὕτη ἡ διαστολὴ τοῦ νόμου, ὅσα συνέταξε Κύριος, λέγων, λάλησον τοῖς υἱοῖς Ἰσραήλ· καὶ λαβέτωσαν πρὸς σὲ δάμαλιν πυῤῥὰν ἄμωμου, ἥτις οὐκ ἔχει ἐν αὐτῇ μῶμον, καὶ ᾗ οὐκ ἐπεβλήθη ἐπʼ αὐτὴν ζυγός.
\vs{3}Καὶ δώσεις αὐτὴν πρὸς Ἐλεάζαρ τὸν ἱερέα· καὶ ἐξάξουσιν αὐτὴν ἔξω τῆς παρεμβολῆς εἰς τόπον καθαρὸν, καὶ σφάξουσιν αὐτὴν ἐνώπιον αὐτοῦ.
\vs{4}Καὶ λήψεται Ἐλεάζαρ ἀπὸ τοῦ αἵματος αὐτῆς, καὶ ῥανεῖ ἀπέναντι τοῦ προσώπου τῆς σκηνῆς τοῦ μαρτυρίου ἀπὸ τοῦ αἵματος αὐτῆς ἑπτάκις.
\vs{5}Καὶ κατακαύσουσιν αὐτὴν ἐναντίον αὐτοῦ· καὶ τὸ δέρμα καὶ τὰ κρέα αὐτῆς καὶ τὸ αἷμα αὐτῆς σὺν τῇ κόπρῳ αὐτῆς κατακαυθήσεται.
\vs{6}Καὶ λήψεται ὁ ἱερεὺς ξύλον κέδρινον καὶ ὕσσωπον καὶ κόκκινον, καὶ ἐμβαλοῦσιν εἰς μέσον τοῦ κατακαύματος τῆς δαμάλεως.

\vs{7}Καὶ πλυνεῖ τὰ ἱμάτια αὐτοῦ ὁ ἱερεὺς, καὶ λούσεται τὸ σῶμα αὐτοῦ ὕδατι, καὶ μετὰ ταῦτα εἰσελεύσεται εἰς τὴν παρεμβολὴν, καὶ ἀκάθαρτος ἔσται ὁ ἱερεὺς ἕως ἑσπέρας.
\vs{8}Καὶ ὁ κατακαίων αὐτὴν πλυνεῖ τὰ ἱμάτια αὐτοῦ, καὶ λούσεται τὸ σῶμα αὐτοῦ, καὶ ἀκάθαρτος ἔσται ἕως ἑσπέρας.
\vs{9}Καὶ συνάξει ἄνθρωπος καθαρὸς τὴν σποδὸν τῆς δαμάλεως, καὶ ἀποθήσει ἔξω τῆς παρεμβολῆς εἰς τόπον καθαρόν· καὶ ἔσται τῇ συναγωγῇ υἱῶν Ἰσραὴλ εἰς διατήρησιν· ὕδωρ ῥαντισμοῦ ἅγνισμά ἐστι.
\vs{10}Καὶ ὁ συνάγων τὴν σποδιὰν τῆς δαμάλεως, πλυνεῖ τὰ ἱμάτια αὐτου, καὶ ἀκάθαρτος ἔσται ἕως ἑσπέρας· καὶ ἔσται τοῖς υἱοῖς Ἰσραὴλ καὶ τοῖς προσηλύτοις προσκειμένοις νόμιμον αἰώνιον.

\vs{11}Ὁ ἁπτόμενος τοῦ τεθνηκότος πάσης ψυχῆς ἀνθρώπου, ἀκάθαρτος ἔσται ἑπτὰ ἡμέρας.
\vs{12}Οὗτος ἁγνισθήσεται τῇ ἡμέρᾳ τῇ τρίτῃ καὶ τῇ ἡμέρᾳ τῇ ἑβδόμῃ, καὶ καθαρὸς ἔσται· ἐὰν δὲ μὴ ἀφαγνισθῇ τῇ ἡμέρᾳ τῇ τρίτῃ καὶ τῇ ἡμέρᾳ τῇ ἑβδόμῃ, οὐ καθαρὸς ἔσται.
\vs{13}Πᾶς ὁ ἁπτόμενος τοῦ τεθνηκότος ἀπὸ ψυχῆς ἀνθρώπου, ἐὰν ἀποθάνῃ, καὶ μὴ ἀφαγνισθῇ, τὴν σκηνὴν Κυρίου ἐμίανεν· ἐκτριβήσεται ἡ ψυχὴ ἐκείνη ἐξ Ἰσραὴλ, ὅτι ὕδωρ ῥαντισμοῦ οὐ περιεῤῥαντίσθη ἐπʼ αὐτόν· ἀκάθαρτός ἐστιν· ἔτι ἡ ἀκαθαρσία αὐτοῦ ἐν αὐτῷ ἐστι.
\vs{14}Καὶ οὗτος ὁ νόμος· ἄνθρωπος ἐὰν ἀποθάνῃ ἐν οἰκίᾳ, πᾶς ὁ εἰσπορευόμενος εἰς τὴν οἰκίαν, καὶ ὅσα ἐστὶν ἐν τῇ οἰκίᾳ, ἀκάθαρτα ἔσται ἑπτὰ ἡμέρας.
\vs{15}Καὶ πᾶν σκεῦος ἀνεῳγμένον ὅσα οὐχὶ δεσμὸν καταδέδεται ἐπʼ αὐτῷ, ἀκάθαρτά ἐστι.
\vs{16}Καὶ πᾶς ὃς ἂν ἅψηται ἐπὶ προσώπου τοῦ πεδίου τραυματίου ἢ νεκροῦ ἢ ὀστέου ἀνθρωπίνου ἢ μνήματος, ἑπτὰ ἡμέρας ἀκάθαρτος ἔσται.

\vs{17}Καὶ λήψονται τῷ ἀκαθάρτῳ ἀπὸ τῆς σποδιᾶς τῆς κατακεκαυμένης τοῦ ἁγνισμοῦ, καὶ κέχεοῦσιν ἐπʼ αὐτὴν ὕδωρ ζῶν εἰς σκεῦος.
\vs{18}Καὶ λήψεται ὕσσωπον, καὶ βάψει εἰς τὸ ὕδωρ ἀνὴρ καθαρὸς, καὶ περιῤῥανεῖ ἐπὶ τὸν οἶκον, καὶ ἐπὶ τὰ σκεύη, καὶ ἐπὶ τὰς ψυχὰς, ὅσαι ἂν ὦσιν ἐκεῖ, καὶ ἐπὶ τὸν ἡμμένον τοῦ ὀστέου τοῦ ἀνθρωπίνου, ἢ τοῦ τραυματίου, ἢ τοῦ τεθνηκότος, ἢ τοῦ μνήματος.
\vs{19}Καὶ περιῤῥανεῖ ὁ καθαρὸς ἐπὶ τὸν ἀκάθαρτον ἐν τῇ ἡμέρᾳ τῇ τρίτῃ καὶ ἐν τῇ ἡμέρᾳ τῇ ἑβδόμῃ, καὶ ἀφαγνισθήσεται τῇ ἡμέρᾳ τῇ ἑβδόμῃ· καὶ πλυνεῖ τὰ ἱμάτια αὐτοῦ, καὶ λούσεται ὕδατι, καὶ ἀκάθαρτος ἔσται ἕως ἑσπέρας.
\vs{20}Καὶ ἄνθρωπος ὃς ἂν μιανθῇ, καὶ μὴ ἀφαγνισθῇ, ἐξολοθρευθήσεται ἡ ψυχὴ ἐκείνη ἐκ μέσου τῆς συναγωγῆς, ὅτι τὰ ἅγια Κυρίου ἐμίανεν, ὅτι ὕδωρ ῥαντισμοῦ οὐ περιεῤῥαντίσθη ἐπʼ αὐτόν· ἀκάθαρτός ἐστι.
\vs{21}Καὶ ἔσται ὑμῖν νόμιμον αἰώνιον· καὶ ὁ περιῤῥαίνων ὕδωρ ῥαντισμοῦ, πλυνεῖ τὰ ἱμάτια αὐτοῦ· καὶ ὁ ἁπτόμενος τοῦ ὕδατος τοῦ ῥαντισμοῦ, ἀκάθαρτος ἔσται ἕως ἑσπέρας.
\vs{22}Καὶ παντὸς οὗ ἐὰν ἅψηται αὐτοῦ ὁ ἀκάθαρτος, ἀκάθαρτον ἔσται· καὶ ψυχὴ ἡ ἁπτομένη, ἀκάθαρτος ἔσται ἕως ἑσπέρας.

\ch{20}
Καὶ ἦλθον οἱ υἱοὶ Ἰσραὴλ, πᾶσα ἡ συναγωγὴ, εἰς τὴν ἔρημον Σὶν, ἐν τῷ μηνὶ τῷ πρώτῳ, καὶ κατέμεινεν ὁ λαὸς ἐν Κάδης· καὶ ἐτελεύτησεν ἐκεῖ Μαριὰμ, καὶ ἐτάφη ἐκεῖ.
\vs{2}Καὶ οὐκ ἦν ὕδωρ τῇ συναγωγῇ· καὶ ἠθροίσθησαν ἐπὶ Μωυσῆν καὶ Ἀαρών.
\vs{3}Καὶ ἐλοιδορεῖτο ὁ λαὸς πρὸς Μωυσῆν, λέγοντες, ὄφελον ἀπεθάνομεν ἐν τῇ ἀπωλείᾳ τῶν ἀδελφῶν ἡμῶν ἔναντι Κυρίου.
\vs{4}Καὶ ἱνατί ἀνηγάγετε τὴν συναγωγὴν Κυρίου εἰς τὴν ἔρημον ταύτην ἀποκτεῖναι ἡμᾶς, καὶ τα κτήνη ἡμῶν;
\vs{5}Καὶ ἱνατί τοῦτο; ἀνηγάγετε ἡμᾶς ἐξ Αἰγύπτου, παραγενέσθαι εἰς τὸν τόπον τὸν πονηρὸν τοῦτον· τόπος οὗ οὐ σπείρεται, οὐδὲ συκαῖ, οὐδὲ ἄμπελοι, οὔτε ῥοαὶ, οὔτε ὕδωρ ἐστὶ πιεῖν.

\vs{6}Καὶ ἦλθε Μωυσῆς καὶ Ἀαρὼν ἀπὸ προσώπου τῆς συναγωγῆς ἐπὶ τὴν θύραν τῆς σκηνῆς τοῦ μαρτυρίου, καὶ ἔπεσον ἐπὶ πρόσωπον· καὶ ὤφθη ἡ δόξα Κυρίου πρὸς αὐτοὺς.
\vs{7}Καὶ ἐλάλησε Κύριος πρὸς Μωυσῆν, λέγων,
\vs{8}λάβε τὴν ῥάβδον σου, καὶ ἐκκλησίασον τὴν συναγωγὴν σὺ καὶ Ἀαρὼν ὁ ἀδελφός σου, καὶ λαλήσατε πρὸς τὴν πέτραν ἐναντίον αὐτῶν, καὶ δώσει τὰ ὕδατα αὐτῆς, καὶ ἐξοίσετε αὐτοῖς ὕδωρ ἐκ τῆς πέτρας, καὶ ποτιεῖτε τὴν συναγωγὴν, καὶ τὰ κτήνη αὐτῶν.
\vs{9}Καὶ ἔλαβε Μωυσῆς τὴν ῥάβδον τὴν ἀπέναντι Κυρίου, καθὰ συνέταξε Κύριος.
\vs{10}Καὶ ἐξεκκλησίασε Μωυσῆς καὶ Ἀαρὼν τὴν συναγωγὴν ἀπέναντι τῆς πέτρας, καὶ εἶπε πρὸς αὐτοὺς, ἀκούσατέ μου οἱ ἀπειθεῖς· μὴ ἐκ τῆς πέτρας ταύτης ἐξάξομεν ὑμῖν ὕδωρ;
\vs{11}Καὶ ἐπάρας Μωυσῆς τὴν χεῖρα αὐτοῦ, ἐπάταξε τὴν πέτραν τῇ ῥάβδῳ δίς· καὶ ἐξῆλθεν ὕδωρ πολὺ, καὶ ἔπιεν ἡ συναγωγὴ, καὶ τὰ κτήνη αὐτῶν.
\vs{12}Καὶ εἶπε Κύριος πρὸς Μωυσῆν καὶ Ἀαρὼν, ὅτι οὐκ ἐπιστεύσατε ἁγιάσαι με ἐναντίον τῶν υἱῶν Ἰσραὴλ, διὰ τοῦτο οὐκ εἰσάξετε ὑμεῖς τὴν συναγωγὴν ταύτην εἰς τὴν γῆν ἣν δέδωκα αὐτοῖς.
\vs{13}Τοῦτο τὸ ὕδωρ Ἀντιλογίας, ὅτι ἐλοιδορήθησαν οἱ υἱοὶ Ἰσραὴλ ἔναντι Κυρίου, καὶ ἡγιάσθη ἐν αὐτοῖς.

\vs{14}Καὶ ἀπέστειλε Μωυσῆς ἀγγέλους ἐκ Κάδης πρὸς βασιλέα Ἐδὼμ, λέγων, τάδε λέγει ὁ ἀδελφός σου Ἰσραήλ· σὺ ἐπίστῃ πάντα τὸν μόχθον τὸν εὑρόντα ἡμᾶς.
\vs{15}Καὶ κατέβησαν οἱ πατέρες ἡμῶν εἰς Αἴγυπτον, καὶ παρῳκήσαμεν ἐν Αἰγύπτῳ ἡμέρας πλείους, καὶ ἐκάκωσαν ἡμᾶς οἱ Αἰγύπτιοι καὶ τοὺς πατέρας ἡμῶν.
\vs{16}Καὶ ἀνεβοήσαμεν πρὸς Κύριον, καὶ εἰσήκουσε Κύριος τῆς φωνῆς ἡμῶν, καὶ ἀποστείλας ἄγγελον, ἐξήγαγεν ἡμᾶς ἐξ Αἰγύπτου· καὶ νῦν ἐσμὲν ἐν Κάδης πόλει, ἐκ μέρους τῶν ὁρίων σου.
\vs{17}Παρελευσόμεθα διὰ τῆς γῆς σου· οὐ διελευσόμεθα διʼ ἀγρῶν, οὐδὲ διʼ ἀμπελώνων, οὐδὲ πιόμεθα ὕδωρ ἐκ λάκκου σου· ὁδῷ βασιλικῇ πορευσόμεθα· οὐκ ἐκκλινοῦμεν δεξιὰ οὐδὲ εὐώνυμα, ἕως ἂν παρέλθωμεν τὰ ὅριά σου.
\vs{18}Καὶ εἶπε πρὸς αὐτὸν Ἐδὼμ, Οὐ διελεύσῃ διʼ ἐμοῦ· εἰ δὲ μὴ, ἐν πολέμῳ ἐξελεύσομαι εἰς συνάντησίν σοι.
\vs{19}Καὶ λέγουσιν αὐτῷ οἱ υἱοὶ Ἰσραὴλ, παρὰ τὸ ὄρος παρελευσόμεθα· ἐὰν δὲ τοῦ ὕδατός σου πίωμεν ἐγώ τε καὶ τὰ κτήνη μου, δώσω τιμήν σοι· ἀλλὰ τὸ πρᾶγμα οὐδέν ἐστι· παρὰ τὸ ὄρος παρελευσόμεθα.
\vs{20}Ὁ δὲ εἶπεν, οὐ διελεύσῃ διʼ ἐμοῦ· καὶ ἐξῆλθεν Ἐδὼμ εἰς συνάντησιν αὐτῷ ἐν ὄχλῳ βαρεῖ, καὶ ἐν χειρὶ ἰσχυρᾷ.
\vs{21}Καὶ οὐκ ἠθέλησεν Ἐδὼμ δοῦναι τῷ Ἰσραὴλ παρελθεῖν διὰ τῶν ὁρίων αὐτοῦ· καὶ ἐξέκλινεν Ἰσραὴλ ἀπʼ αὐτοῦ.
\vs{22}Καὶ ἀπῇραν ἐκ Καδης· καὶ παρεγένοντο οἱ υἱοὶ Ἰσραὴλ πᾶσα ἡ συναγωγὴ εἰς Ὢρ τὸ ὄρος.

\vs{23}Καὶ εἶπε Κύριος πρὸς Μωυσῆν καὶ Ἀαρὼν ἐν Ὢρ τῷ ὄρει ἐπὶ τῶν ὁρίων τῆς γῆς Ἐδὼμ, λέγων,
\vs{24}προστεθήτω Ἀαρὼν πρὸς τὸν λαὸν αὐτοῦ, ὅτι οὐ μὴ εἰσέλθητε εἰς τὴν γῆν ἣν δέδωκα τοῖς υἱοῖς Ἰσραήλ, διότι παρωξύνατέ με ἐπὶ τοῦ ὕδατος τῆς λοιδορίας.
\vs{25}Λάβε τὸν Ἀαρὼν, καὶ Ἐλεάζαρ τὸν υἱὸν αὐτοῦ, καὶ ἀναβίβασον αὐτοὺς εἰς Ὢρ τὸ ὄρος, ἔναντι πάσης τῆς συναγωγῆς,
\vs{26}καὶ ἔκδυσον Ἀαρὼν τὴν στολὴν αὐτοῦ, καὶ ἔνδυσον Ἐλεάζαρ τὸν υἱὸν αὐτοῦ· καὶ Ἀαρὼν προστεθεὶς ἀποθανέτω ἐκεῖ.
\vs{27}Καὶ ἐποίησε Μωυσῆς καθὰ συνέταξε Κύριος αὐτῷ, καὶ ἀνεβίβασεν αὐτὸν εἰς Ὢρ τὸ ὄρος, ἐναντίον πάσης τῆς συναγωγῆς,
\vs{28}καὶ ἐξέδυσε τὸν Ἀαρὼν τὰ ἱμάτια αὐτοῦ, καὶ ἐνέδυσεν αὐτὰ Ἐλεάζαρ τὸν υἱὸν αὐτοῦ· καὶ ἀπέθανεν Ἀαρὼν ἐπὶ τῆς κορυφῆς τοῦ ὄρους· καὶ κατέβη Μωυσῆς καὶ Ἐλεάζαρ ἐκ τοῦ ὄρους.
\vs{29}Καὶ εἶδε πᾶσα ἡ συναγωγὴ ὅτι ἀπελύθη Ἀαρὼν, καὶ ἔκλαυσαν τὸν Ἀαρὼν τριάκοντα ἡμέρας πᾶς οἶκος Ἰσραήλ.

\ch{21}
Καὶ ἤκουσεν ὁ Χανανεὶς βασιλεὺς Ἀρὰδ ὁ κατοικῶν κατὰ τὴν ἔρημον, ὅτι ἦλθεν Ἰσραὴλ ὁδὸν Ἀθαρεὶν, καὶ ἐπολέμησε πρὸς Ἰσραὴλ, καὶ κατεπροενόμευσεν ἐξ αὐτῶν αἰχμαλωσίαν.
\vs{2}Καὶ ηὔξατο Ἰσραὴλ εὐχὴν Κυρίῳ, καὶ εἶπεν, ἐάν μοι παραδῷς τὸν λαὸν τοῦτον ὑποχείριον, ἀναθεματιῶ αὐτὸν καὶ τὰς πόλεις αὐτοῦ.
\vs{3}Καὶ εἰσήκουσε Κύριος τῆς φωνῆς Ἰσραὴλ, καὶ παρέδωκε τὸν Χανανεὶν ὑποχείριον αὐτοῦ· καὶ ἀνεθεμάτισεν αὐτὸν, καὶ τὰς πόλεις αὐτοῦ· καὶ ἐπεκάλεσαν τὸ ὄνομα τοῦ τόπου ἐκείνου, Ἀνάθεμα.

\vs{4}Καὶ ἀπάραντες ἐξ Ὢρ τοῦ ὄρους ὁδὸν ἐπὶ θάλασσαν ἐρυθρᾶν, περιεκύκλωσαν γῆν Ἐδώμ· καὶ ὠλιγοψύχησεν ὁ λαὸς ἐν τῇ ὁδῷ.
\vs{5}Καὶ κατελάλει ὁ λαὸς πρὸς τὸν Θεὸν καὶ κατὰ Μωυσῆ, λέγοντες, ἱνατί τοῦτο; ἐξήγαγες ἡμᾶς ἐξ Αἰγύπτου ἀποκτεῖναι ἐν τῇ ἐρήμῳ; ὅτι οὐκ ἔστιν ἄρτος, οὐδὲ ὕδωρ· ἡ δὲ ψυχὴ ἡμῶν προσώχθισεν ἐν τῷ ἄρτῳ τῷ διακένῳ τούτῳ.
\vs{6}Καὶ ἀπέστειλε Κύριος εἰς τὸν λαὸν τοὺς ὄφεις τοὺς θανατοῦντας, καὶ ἔδακνον τὸν λαόν, καὶ ἀπέθανε λαὸς πολὺς τῶν υἱῶν Ἰσραήλ.
\vs{7}Καὶ παραγενόμενος ὁ λαὸς πρὸς Μωυσῆν, ἔλεγον, ὅτι ἡμάρτομεν, ὅτι κατελαλήσαμεν κατὰ τοῦ Κυρίου, καὶ κατὰ σοῦ· εὔξαι οὖν πρὸς Κύριον, καὶ ἀφελέτω ἀφʼ ἡμῶν τὸν ὄφιν.
\vs{8}Καὶ ηὔξατο Μωυσῆς πρὸς Κύριον περὶ τοῦ λαοῦ· καὶ εἰπε Κύριος πρὸς Μωυσῆν, ποίησον σεαυτῷ ὄφιν, καὶ θὲς αὐτὸν ἐπὶ σημείου, καὶ ἔσται ἐὰν δάκῃ ὄφις ἄνθρωπον, πᾶς ὁ δεδηγμένος ἰδὼν αὐτὸν ζήσεται.
\vs{9}Καὶ ἐποίησε Μωυσῆς ὄφιν χαλκοῦν, καὶ ἔστησεν αὐτὸν ἐπὶ σημείου· καὶ ἐγένετο ὅταν ἔδακνεν ὄφις ἄνθρωπον, καὶ ἐπέβλεψεν ἐπὶ τὸν ὄφιν τὸν χαλκοῦν, καὶ ἔζη.

\vs{10}Καὶ ἀπῇραν οἱ υἱοὶ Ἰσραὴλ, καὶ παρενέβαλον ἐν Ὠβώθ.
\vs{11}Καὶ ἐξάραντες ἐξ Ὠβὼθ, καὶ παρενέβαλον ἐν Ἀχαλγαὶ ἐκ τοῦ πέραν ἐν τῇ ἐρήμῳ, ἥ ἐστι κατὰ πρόσωπον Μωὰβ, κατʼ ἀνατολὰς ἡλίου.
\vs{12}Καὶ ἐκεῖθεν ἀπῇραν, καὶ παρενέβαλον εἰς φάραγγα Ζαρέδ.
\vs{13}Καὶ ἐκεῖθεν ἀπάραντες παρενέβαλον εἰς τὸ πέραν Ἀρνῶν ἐν τῇ ἐρήμῳ, τὸ ἐξέχον ἀπὸ τῶν ὁρίων τῶν Ἀμοῤῥαίων· ἔστι γὰρ Ἀρνῶν ὅρια Μωὰβ, ἀναμέσον Μωὰβ καὶ ἀναμέσον τοῦ Ἀμοῤῥαίου.
\vs{14}Διὰ τοῦτο λέγεται ἐν βιβλίῳ, πόλεμος τοῦ Κυρίου τὴν Ζωὸβ ἐφλόγισε, καὶ τοὺς χιμάῤῥους Ἀρνῶν.
\vs{15}Καὶ τοὺς χιμάῤῥους κατέστησε κατοικίσαι Ἤρ· καὶ πρόσκειται τοῖς ὁρίοις Μωάβ.

\vs{16}Καὶ ἐκεῖθεν τὸ φρέαρ· τοῦτο φρέαρ, ὃ εἶπε Κύριος πρὸς Μωυσῆν, συνάγαγε τὸν λαὸν, καὶ δώσω αὐτοῖς ὕδωρ πιεῖν.
\vs{17}Τότε ᾖσεν Ἰσραὴλ τὸ ἆσμα τοῦτο ἐπὶ τοῦ φρέατος, ἐξάρχετε αὐτῷ φρέαρ,
\vs{18}ὤρυξαν αὐτὸ ἄρχοντες, ἐξελατόμησαν αὐτὸ βασιλεῖς ἐθνῶν ἐν τῇ βασιλείᾳ αὐτῶν, ἐν τῷ κυριεῦσαι αὐτῶν·
\vs{19}καὶ ἀπὸ φρέατος εἰς Μανθαναεὶν, καὶ ἀπὸ Μανθαναεὶν εἰς Νααλιὴλ, καὶ ἀπὸ Νααλιὴλ εἰς Βαμὼθ, καὶ ἀπὸ Βαμὼθ εἰς Ἰανὴν, ἥ ἐστιν ἐν τῷ πεδίῳ Μωὰβ, ἀπὸ κορυφῆς τοῦ λελαξευμένου, τὸ βλέπον κατὰ πρόσωπον τῆς ἐρήμου.

\vs{20}Καὶ ἀπέστειλε Μωυσῆς πρέαβεις πρὸς Σηὼν βασιλέα Ἀμοῤῥαίων, λόγοις εἰρηνικοῖς, λέγων,
\vs{21}παρελευσόμεθα διὰ τῆς γῆς σου, τῇ ὁδῷ πορευσόμεθα· οὐκ ἐκκλινοῦμεν οὔτε εἰς ἀγρὸν, οὔτε εἰς ἀμπελῶνα·
\vs{22}Οὐ πιόμεθα ὕδωρ ἐκ φρέατός σου· ὁδῷ βασιλικῇ παρευσόμεθα, ἕως παρέλθωμεν τὰ ὅριά σου.
\vs{23}Καὶ οὐκ ἔδωκε Σηὼν τῷ Ἰσραὴλ παρελθεῖν διὰ τῶν ὁρίων αὐτοῦ· καὶ συνήγαγε Σηὼν πάντα τὸν λαὸν αὐτοῦ, καὶ ἐξῆλθε παρατάξασθαι τῷ Ἰσραὴλ εἰς τὴν ἔρημον· καὶ ἦλθεν εἰς Ἰασσὰ, καὶ παρετάξατο τῷ Ἰσραήλ.
\vs{24}Καὶ ἐπάταξεν αὐτὸν Ἰσραὴλ φόνῳ μαχαίρης, καὶ κατεκυρίευσαν τῆς γῆς αὐτοῦ, ἀπὸ Ἀρνῶν ἕως Ἰαβὸκ, ἕως υἱῶν Ἀμμὰν, ὅτι Ἰαζὴρ ὅρια υἱῶν Ἀμμάν ἐστι.
\vs{25}Καὶ ἔλαβεν Ἰσραὴλ πάσας τὰς πόλεις ταύτας, καὶ κατῷκησεν Ἰσραὴλ ἐν πάσαις ταῖς πόλεσι τῶν Ἀμοῤῥαίων, ἐν Ἑσεβὼν, καὶ ἐν πάσαις ταῖς συγκυρούσαις αὐτῇ.
\vs{26}Ἔστι γὰρ Ἐσεβὼν, πόλις Σηὼν τοῦ βασιλέως τῶν Ἀμοῤῥαίων ἐστίν· καὶ οὗτος ἐπολέμησε βασιλέα Μωὰβ τὸ πρότερον· καὶ ἔλαβον πᾶσαν τὴν γῆν αὐτοῦ, ἀπὸ Ἀροὴρ ἕως Ἀρνῶν.
\vs{27}Διὰ τοῦτο ἐροῦσιν οἱ αἰνιγματισταὶ, ἔλθετε εἰς Ἐσεβὼν, ἵνα οἰκοδομηθῇ καὶ κατασκευασθῇ πόλις Σηών·
\vs{28}ὅτι πῦρ ἐξῆλθεν ἐξ Ἑσεβὼν, φλὸξ ἐκ πόλεως Σηὼν, καὶ κατέφαγεν ἕως Μωὰβ, καὶ κατέπιε στήλας Ἀρνῶν.
\vs{29}Οὐαί σοι Μωὰβ, ἀπώλου λαὸς Χαμώς· ἀπεδόθησαν οἱ υἱοὶ αὐτῶν διασώζεσθαι, καὶ αἱ θυγατέρες αὐτῶν αἰχμάλωτοι τῷ βασιλεῖ τῶν Ἀμοῤῥαίων Σηὼν,
\vs{30}καὶ τὸ σπέρμα αὐτῶν ἀπολεῖται, Ἑσεβὼν ἕως Δαιβών· καὶ αἱ γυναῖκες ἔτι προσεξέκαυσαν πῦρ ἐπὶ Μωάβ.

\vs{31}Κατῴκησε δὲ Ἰσραὴλ ἐν πάσαις ταῖς πόλεσι τῶν Ἀμοῤῥαίων.
\vs{32}Καὶ ἀπέστειλε Μωυσῆς κατασκέψασθαι τὴν Ἰαζήρ· καὶ κατελάβοντο αὐτὴν, καὶ τὰς κώμας αὐτῆς, καὶ ἐξέβαλον τὸν Ἀμοῤῥαῖον τὸν κατοικοῦντα ἐκεῖ.
\vs{33}Καὶ ἐπιστρέψαντες, ἀνέβησαν ὁδὸν τὴν εἰς Βασάν· καὶ ἐξῆλθεν Ὢγ βασιλεὺς τῆς Βασὰν εἰς συνάντησιν αὐτοῖς, καὶ πᾶς ὁ λαὸς αὐτοῦ εἰς πόλεμον εἰς Ἐδραείν.
\vs{34}Καὶ εἶπε Κύριος πρὸς Μωυσῆν, μὴ φοβηθῇς αὐτὸν, ὅτι εἰς τὰς χεῖράς σου παραδέδωκα αὐτὸν, καὶ πάντα τὸν λαὸν αὐτοῦ, καὶ πᾶσαν τὴν γῆν αὐτοῦ· καὶ ποιήσεις αὐτῷ καθὼς ἐποίησας τῷ Σηὼν βασιλεῖ τῶν Ἀμοῤῥαίων, ὃς κατῴκει ἐν Ἑσεβών.
\vs{35}Καὶ ἐπάταξεν αὐτὸν καὶ τοὺς υἱοὺς αὐτοῦ, καὶ πάντα τὸν λαὸν αὐτοῦ, ἕως τοῦ μὴ καταλιπεῖν αὐτοῦ ζωγρείαν· καὶ ἐκληρονόμησαν τὴν γῆν αὐτον.

\ch{22}
Καὶ ἀπάραντες οἱ υἱοὶ Ἰσραὴλ παρενέβαλον ἐπὶ δυσμῶν Μωὰβ παρὰ τὸν Ἰορδάνην κατὰ Ἰεριχώ.
\vs{2}Καὶ ἰδὼν Βαλὰκ υἱὸς Σεπφὼρ πάντα ὅσα ἐποίησεν Ἰσραὴλ τῷ Ἀμοῤῥαίῳ,
\vs{3}καὶ ἐφοβήθη Μωὰβ τὸν λαὸν σφόδρα ὅτι πολλοὶ ἦσαν· καὶ προσώχθισε Μωὰβ ἀπὸ προσώπου υἱῶν Ἰσραήλ.
\vs{4}Καὶ εἶπε Μωὰβ τῇ γερουσίᾳ Μαδιὰμ, νῦν ἐκλείξει ἡ συναγωγὴ αὕτη πάντας τοὺς κύκλῳ ἡμῶν, ὡσεὶ ἐκλείξαι ὁ μόσχος τὰ χλωρὰ ἐκ τοῦ πεδίου· καὶ Βαλὰκ υἱὸς Σεπφὼρ βασιλεὺς Μωὰβ ἦν κατὰ τὸν καιρὸν ἐκεῖνον.
\vs{5}Καὶ ἀπέστειλε πρέσβεις πρὸς Βαλαὰμ υἱὸν Βεὼρ Φαθουρὰ, ὅ ἐστιν ἐπὶ τοῦ ποταμοῦ γῆς υἱῶν λαοῦ αὐτοῦ, καλέσαι αὐτὸν, λέγων, ἰδοὺ λαὸς ἐξελήλυθεν ἐξ Αἰγύπτου, καὶ ἰδοὺ κατεκάλυψε τὴν ὄψιν τῆς γῆς, καὶ οὗτος ἐγκάθηται ἐχόμενός μου.
\vs{6}Καὶ νῦν δεῦρο ἄρασαί μοι τὸν λαὸν τοῦτον, ὅτι ἰσχύει οὗτος ἢ ἡμεῖς, ἐὰν δυνώμεθα πατάξαι ἐξ αὐτῶν, καὶ ἐκβαλῶ αὐτοὺς ἐκ τῆς γῆς· ὅτι οἶδα οὓς ἐὰν εὐλογήσῃς σὺ, εὐλόγηνται, καὶ οὓς ἂν καταράσῃ σὺ, κεκατήρανται.
\vs{7}Καὶ ἐπορεύθη ἡ γερουσία Μωὰβ, καὶ ἡ γερουσία Μαδιὰμ, καὶ τὰ μαντεῖα ἐν ταῖς χερσὶν αὐτῶν· καὶ ἦλθον πρὸς Βαλαὰμ, καὶ εἶπαν αὐτῷ τὰ ῥήματα Βαλάκ.
\vs{8}Καὶ εἶπε πρὸς αὐτοὺς, καταλύσατε αὐτοῦ τὴν νύκτα, καὶ ἀποκριθήσομαι ὑμῖν πράγματα ἃ ἂν λαλήσῃ Κύριος πρὸς μέ· καὶ κατέμειναν οἱ ἄρχοντες Μωὰβ παρὰ Βαλαάμ.

\vs{9}Καὶ ἦλθεν ὁ Θεὸς πρὸς Βαλαὰμ, καὶ εἶπεν αὐτῷ, τί οἱ ἄνθρωποι οὗτοι παρὰ σοί;
\vs{10}Καὶ εἶπε Βαλαὰμ πρὸς τὸν Θεὸν, Βαλὰκ υἱὸς Σεπφὼρ, βασιλεὺς Μωὰβ, ἀπέστειλεν αὐτοὺς πρὸς μὲ, λέγων,
\vs{11}ἰδοὺ λαὸς ἐξελήλυθεν ἐξ Αἰγύπτου, καὶ κεκάλυφεν τὴν ὄψιν τῆς γῆς, καὶ οὗτος ἐγκάθηται ἐχόμενός μου, καὶ νῦν δεῦρο ἄρασαί μοι αὐτὸν, εἰ ἄρα δυνήσομαι πατάξαι αὐτὸν, καὶ ἐκβαλῶ αὐτὸν ἀπὸ τῆς γῆς.
\vs{12}Καὶ εἶπεν ὁ Θεὸς πρὸς Βαλαὰμ, οὐ πορεύσῃ μετʼ αὐτῶν, οὐδὲ καταράσῃ τὸν λαόν· ἔστι γὰρ εὐλογημένος.
\vs{13}Καὶ ἀναστὰς Βαλαὰμ τοπρωῒ, εἶπε τοῖς ἄρχουσι Βαλὰκ, ἀποτρέχετε πρὸς τὸν κύριον ὑμῶν, οὐκ ἀφίησί με ὁ Θεὸς πορεύεσθαι μεθʼ ὑμῶν.
\vs{14}Καὶ ἀναστάντες οἱ ἄρχοντες Μωὰβ, ἦλθον πρὸς Βαλὰκ, καὶ εἶπαν, οὐ θέλει Βαλαὰμ πορευθῆναι μεθʼ ἡμῶν.

\vs{15}Καὶ προσέθετο Βαλὰκ ἔτι ἀποστεῖλαι ἄρχοντας πλείους, καὶ ἐντιμοτέρους τούτων.
\vs{16}Καὶ ἦλθον πρὸς Βαλαὰμ, καὶ λέγουσιν αὐτῷ, τάδε λέγει Βαλὰκ ὁ τοῦ Σεπφώρ· ἀξιῶ σε μὴ ὀκνήσῃς ἐλθεῖν πρὸς μέ.
\vs{17}ἐντίμως γὰρ τιμήσω σε, καὶ ὅσα ἐὰν εἴπῃς ποιήσω σοι· καὶ δεῦρο ἐπικατάρασαί μοι τὸν λαὸν τοῦτον.
\vs{18}Καὶ ἀπεκρίθη Βαλαὰμ, καὶ εἶπε τοῖς ἄρχουσι Βαλὰκ, ἐὰν δῷ μοι Βαλὰκ πλήρη τὸν οἶκον αὐτοῦ ἀργυρίου καὶ χρυσίου, οὐ δυνήσομαι παραβῆναι τὸ ῥῆμα Κυρίου τοῦ Θεοῦ, ποιῆσαι αὐτὸ μικρὸν ἢ μέγα ἐν τῇ διανοίᾳ μου.
\vs{19}Καὶ νῦν ὑπομείνατε αὐτοῦ καὶ ὑμεῖς τὴν νύκτα ταύτην, καὶ γνώσομαι τί προσθήσει Κύριος λαλῆσαι πρὸς μέ.
\vs{20}Καὶ ἦλθεν ὁ Θεὸς πρὸς Βαλαὰμ νυκτὸς, καὶ εἶπεν αὐτῷ, εἰ καλέσαι σε πάρεισιν οἱ ἄνθρωποι οὗτοι, ἀναστὰς ἀκολούθησον αὐτοῖς· ἀλλὰ τὸ ῥῆμα ὃ ἐὰν λαλήσω πρὸς σὲ, τοῦτο ποιήσεις.

\vs{21}Καὶ ἀναστὰς Βαλαὰμ τοπρωῒ, ἐπέσαξε τὴν ὄνον αὐτοῦ, καὶ ἐπορεύθη μετὰ τῶν ἀρχόντων Μωάβ.
\vs{22}Καὶ ὠργίσθη θυμῷ ὁ Θεὸς ὅτι ἐπορεύθη αὐτός· καὶ ἀνέστη ὁ ἄγγελος τοῦ Θεοῦ διαβαλεῖν αὐτόν· καὶ αὐτὸς ἐπιβεβήκει ἐπὶ τῆς ὄνου αὐτοῦ, καὶ οἱ δύο παῖδες αὐτοῦ μετʼ αὐτοῦ.
\vs{23}Καὶ ἰδοῦσα ἡ ὄνος τὸν ἄγγελον τοῦ Θεοῦ ἀνθεστηκότα ἐν τῇ ὁδῷ, καὶ τὴν ῥομφαίαν ἐσπασμένην ἐν τῇ χειρὶ αὐτοῦ, καὶ ἐξέκλινεν ἡ ὄνος ἐκ τῆς ὁδοῦ, καὶ ἐπορεύετο εἰς τὸ πεδίον· καὶ ἐπάταξε τὴν ὄνον ἐν τῇ ῥάβδῳ αὐτοῦ τοῦ εὐθῦναι αὐτὴν ἐν τῇ ὁδῷ.

\vs{24}Καὶ ἔστη ὁ ἄγγελος τοῦ Θεοῦ ἐν ταῖς αὔλαξι τῶν ἀμπέλων, φραγμὸς ἐντεῦθεν καὶ φραγμὸς ἐντεῦθεν·
\vs{25}Καὶ ἰδοῦσα ἡ ὄνος τὸν ἄγγελον τοῦ Θεοῦ, προσέθλιψεν ἑαυτὴν πρὸς τὸν τοῖχον, καὶ ἀπέθλιψεν τὸν πόδα Βαλαὰμ πρὸς τὸν τοῖχον· καὶ προσέθετο ἔτι μαστίξαι αὐτήν.

\vs{26}Καὶ προσέθετο ὁ ἄγγελος τοῦ Θεοῦ, καὶ ἀπελθὼν ὑπέστη ἐν τόπῳ στενῷ, εἰς ὃν οὐκ ἦν ἐκκλῖναι δεξιὰν ἢ ἀριστεράν.
\vs{27}Καὶ ἰδοῦσα ἡ ὄνος τὸν ἄγγελον τοῦ Θεοῦ, συνεκάθισεν ὑποκάτω Βαλαάμ· καὶ ἐθυμώθη Βαλαὰμ, καὶ ἔτυπτε τὴν ὄνον τῇ ῥάβδῳ.
\vs{28}Καὶ ἤνοιξεν ὁ Θεὸς τὸ στόμα τῆς ὄνου, καὶ λέγει τῷ Βαλαὰμ, τί ἐποίησά σοι, ὅτι πέπαικάς με τρίτον τοῦτο;
\vs{29}Καὶ εἶπε Βαλαὰμ τῇ ὄνῳ, ὅτι ἐμπέπαιχάς μοι, καὶ εἰ εἶχον μάχαιραν ἐν τῇ χειρὶ ἤδη ἂν ἐξεκέντησά σε.
\vs{30}Καὶ λέγει ἡ ὄνος τῷ Βαλαὰμ, οὐκ ἐγὼ ἡ ὄνος σου ἐφʼ ἧς ἐπέβαινες ἀπὸ νεότητός σου, ἕως τῆς σήμερον ἡμέρας; μὴ ὑπεροράσει ὑπεριδούσα ἐποίησά σοι οὕτως; ὁ δὲ εἶπεν, οὐχί.
\vs{31}Ἀπεκάλυψε δὲ ὁ Θεὸς τοὺς ὀφθαλμοὺς Βαλαὰμ, καὶ ὁρᾷ τὸν ἄγγελον Κυρίου ἀνθεστηκότα ἐν τῇ ὁδῷ, καὶ τὴν μάχαιραν ἐσπασμένην ἐν τῇ χειρὶ αὐτοῦ, καὶ κύψας προσεκύνησε τῷ προσώπῳ αὐτοῦ.
\vs{32}Καὶ εἶπεν αὐτῷ ὁ ἄγγελος τοῦ Θεοῦ, διατί ἐπάταξας τὴν ὄνον σου τοῦτο τρίτον; καὶ ἰδοὺ ἐγὼ ἐξῆλθον εἰς διαβολήν σου, ὅτι οὐκ ἀστεία ἡ ὁδός σου ἐναντίον μου, καὶ ἰδοῦσά με ἡ ὄνος, ἐξέκλινεν ἀπʼ ἐμοῦ τρίτον τοῦτο.
\vs{33}Καὶ εἰ μὴ ἐξέκλινεν, νῦν οὖν σὲ μὲν ἀπέκτεινα, ἐκείνην δʼ ἂν περιεποιησάμην.
\vs{34}Καὶ εἶπε Βαλαὰμ τῷ ἀγγέλῳ Κυρίου, ἡμάρτηκα, οὐ γὰρ ἠπιστάμην ὅτι σύ μοι ἀνθέστηκας ἐν τῇ ὁδῷ εἰς συνάντησιν· καὶ νῦν εἰ μή σοι ἀρκέσει, ἀποστραφήσομαι.
\vs{35}Καὶ εἶπεν ὁ ἄγγελος τοῦ Θεοῦ πρὸς Βαλαὰμ, συμπορεύθητι μετὰ τῶν ἀνθρώπων· πλὴν τὸ ῥῆμα ὃ ἐὰν εἴπω πρὸς σὲ, τοῦτο φυλάξῃ λαλῆσαι. Καὶ ἐπορεύθη Βαλαὰμ μετὰ τῶν ἀρχόντων Βαλάκ.

\vs{36}Καὶ ἀκούσας Βαλὰκ ὅτι ἥκει Βαλαὰμ, ἐξῆλθεν εἰς συνάντησιν αὐτῷ, εἰς πόλιν Μωάβ, ἥ ἐστιν ἐπὶ τῶν ὁρίων Ἀρνῶν, ἥ ἐστιν ἐκ μέρους τῶν ὁρίων.
\vs{37}Καὶ εἶπε Βαλὰκ πρὸς Βαλαὰμ, οὐχὶ ἀπέστειλα πρὸς σὲ καλέσαι σε; διατί οὐκ ἤρχου πρὸς μέ; ὄντως οὐ δυνήσομαι τιμῆσαί σε;
\vs{38}Καὶ εἶπε Βαλαὰμ πρὸς Βαλὰκ, ἰδοὺ ἥκω πρὸς σὲ νῦν· δυνατὸς ἔσομαι λαλῆσαί τι; τὸ ῥῆμα ὃ ἐὰν ἐμβάλῃ ὁ Θεὸς εἰς τὸ στόμα μου, τοῦτο λαλήσω.
\vs{39}Καὶ ἐπορεύθη Βαλαὰμ μετὰ Βαλὰκ, καὶ ἦλθον εἰς πόλεις ἐπαύλεων.
\vs{40}Καὶ ἔθυσε Βαλὰκ πρόβατα καὶ μόσχους, καὶ ἀπέστειλε τῷ Βαλαὰμ καὶ τοῖς ἄρχουσι τοῖς μετʼ αὐτοῦ.
\vs{41}Καὶ ἐγενήθη πρωΐ· καὶ παραλαβὼν Βαλὰκ τὸν Βαλαάμ, ἀνεβίβασεν αὐτὸν ἐπὶ τὴν στήλην τοῦ Βαὰλ, καὶ ἔδειξεν αὐτῷ ἐκεῖθεν μέρος τι τοῦ λαοῦ.

\ch{23}
Καὶ εἶπε Βαλαὰμ τῷ Βαλὰκ οἰκοδόμησόν μοι ἐνταῦθα ἑπτὰ βωμοὺς, καὶ ἐτοίμασόν μοι ἐνταῦθα ἑπτὰ μόσχους, καὶ ἑπτὰ κριούς.
\vs{2}Καὶ ἐποίησε Βαλὰκ ὃν τρόπον εἶπεν αὐτῷ Βαλαάμ· καὶ ἀνήνεγκε μόσχον καὶ κριὸν ἐπὶ τὸν βωμόν.
\vs{3}Καὶ εἶπε Βαλαὰμ πρὸς Βαλὰκ, παράστηθι ἐπὶ τῆς θυσίας σου, καὶ πορεύσομαι εἴ μοι φανεῖται ὁ Θεὸς ἐν συναντήσει, καὶ ῥῆμα ὃ ἐάν μοι δείξῃ, ἀναγγελῶ σοι· καὶ παρέστη Βαλὰκ ἐπὶ τῆς θυσίας αὐτοῦ.
\vs{4}Καὶ Βαλαὰμ ἐπορεύθη ἐπερωτῆσαι τὸν Θεόν· καὶ ἐπορεύθη εὐθεῖαν· καὶ ἐφάνη ὁ Θεὸς τῷ Βαλαὰμ· καὶ εἶπε πρὸς αὐτὸν Βαλαάμ, τοὺς ἑπτὰ βωμοὺς ἡτοίμασα, καὶ ἀνεβίβασα μόσχον καὶ κριὸν ἐπὶ τὸν βωμόν.
\vs{5}Καὶ ἐνέβαλεν ὁ Θεὸς ῥῆμα εἰς τὸ στόμα Βαλαὰμ, καὶ εἶπεν, ἐπιστραφεὶς πρὸς Βαλὰκ, οὕτω λαλήσεις.
\vs{6}Καὶ ἀπεστράφη πρὸς αὐτόν· καὶ ὁ δὲ ἐφιστήκει ἐπὶ τῶν ὁλοκαυτωμάτων αὐτοῦ, καὶ πάντες οἱ ἄρχοντες Μωὰβ μετʼ αὐτοῦ· καὶ ἐγενήθη πνεῦμα Θεοῦ ἐπʼ αὐτῷ.
\vs{7}Καὶ ἀναλαβὼν τὴν παραβολὴν αὐτοῦ, εἶπεν ἐκ Μεσοποταμίας μετεπέμψατό με Βαλὰκ βασιλεὺς Μωὰβ ἐξ ὀρέων ἀπʼ ἀνατολῶν, λέγων, δεῦρο ἄρασαί μοι τὸν Ἰακώβ, καὶ δεῦρο ἐπικατάρασαί μοι τὸν Ἰσραήλ.
\vs{8}Τί ἀράσωμαι ὃν μὴ ἀρᾶται Κύριος; ἢ τί καταράσωμαι ὃν μὴ καταρᾶται ὁ Θεός;
\vs{9}Ὅτι ἀπὸ κορυφῆς ὀρέων ὄψομαι αὐτὸν, καὶ ἀπὸ βουνῶν προσνοήσω αὐτόν· ἰδοὺ λαὸς μόνος κατοικήσει, καὶ ἐν ἔθνεσιν οὐ συλλογισθήσεται.
\vs{10}Τίς ἐξηκριβάσατο τὸ σπέρμα Ἰακὼβ, καὶ τίς ἐξαριθμήσεται δήμους Ἰσραήλ; ἀποθάνοι ἡ ψυχή μου ἐν ψυχαῖς δικαίων, καὶ γένοιτο τὸ σπέρμα μου ὡς τὸ σπέρμα τούτων.

\vs{11}Καὶ εἶπε Βαλὰκ πρὸς Βαλαὰμ, τί πεποίηκάς μοι; εἰς κατάρασιν ἐχθρῶν μου κέκληκά σε, καὶ ἰδοὺ εὐλόγηκας εὐλογίαν.
\vs{12}Καὶ εἶπε Βαλαὰμ πρὸς Βαλὰκ, οὐχὶ ὅσα ἂν ἐμβάλῃ ὁ Θεὸς εἰς τὸ στόμα μου, τοῦτο φυλάξω λαλῆσαι;
\vs{13}Καὶ εἶπε πρὸς αὐτὸν Βαλὰκ, δεῦρο ἔτι μετʼ ἐμοῦ εἰς τόπον ἄλλον ἐξ οὗ οὐκ ὄψει αὐτὸν ἐκεῖθεν, ἀλλʼ ἢ μέρος τι αὐτοῦ ὄψει, πάντας δὲ οὐ μὴ ἴδῃς, καὶ κατάρασαί μοι αὐτὸν ἐκεῖθεν.

\vs{14}Καὶ παρέλαβεν αὐτὸν εἰς ἀγροῦ σκοπιὰν ἐπὶ κορυφὴν λελαξευμένου· καὶ ᾠκοδόμησεν ἐκεῖ ἑπτὰ βωμοὺς, καὶ ἀνεβίβασε μόσχον καὶ κριὸν ἐπὶ τὸν βωμόν.
\vs{15}Καὶ εἶπε Βαλαὰμ πρὸς Βαλὰκ, παράστηθι ἐπὶ τῆς θυσίας σου, ἐγὼ δὲ πορεύσομαι ἐπερωτῆσαι τὸν Θεόν.
\vs{16}Καὶ συνήντησεν ὁ Θεὸς τῷ Βαλαὰμ, καὶ ἐνέβαλε ῥῆμα εἰς τὸ στόμα αὐτοῦ, καὶ εἶπεν, ἀποστράφηθι πρὸς Βαλὰκ, καὶ τάδε λαλήσεις.
\vs{17}Καὶ ἀπεστράφη πρὸς αὐτόν· καὶ ὁ δὲ ἐφειστήκει ἐπὶ τῆς ὁλοκαυτώσεως αὐτοῦ, καὶ πάντες οἱ ἄρχοντες Μωὰβ μετʼ αὐτοῦ· καὶ εἶπεν αὐτῷ Βαλὰκ, τί ἐλάλησε Κύριος;
\vs{18}Καὶ ἀναλαβὼν τὴν παραβολὴν αὐτοῦ, εἶπεν, ἀνάστηθι Βαλὰκ, καὶ ἄκουε, ἐνώτισαι μάρτυς υἱὸς Σεπφώρ.
\vs{19}Οὐχ ὡς ἄνθρωπος ὁ Θεὸς διαρτηθῆναι, οὐδʼ ὡς υἱὸς ἀνθρώπου ἀπειληθῆναι, αὐτὸς εἴπας, οὐχὶ ποιήσει; λαλήσει, καὶ οὐχὶ ἐμμενεῖ;
\vs{20}Ἰδοὺ εὐλογεῖν παρείλημμαι· εὐλογήσω, καὶ οὐ μὴ ἀποστρέψω.
\vs{21}Οὐκ ἔσται μόχθος ἐν Ἰακὼβ, οὐδὲ ὀφθήσεται πόνος ἐν Ἰσραήλ· Κύριος ὁ θεὸς αὐτοῦ μετʼ αὐτοῦ, τὰ ἔνδοξα ἀρχόντων ἐν αὐτῷ.
\vs{22}Θεὸς ὁ ἐξαγαγὼν αὐτὸν ἐξ Αἰγύπτου, ὡς δόξα μονοκέρωτος αὐτῷ.
\vs{23}Οὐ γάρ ἐστιν οἰωνισμὸς ἐν Ἰακὼβ, οὐδὲ μαντεία ἐν Ἰσραήλ· κατὰ καιρὸν ῥηθήσεται Ἰακὼβ, καὶ τῷ Ἰσραὴλ, τί ἐπιτελέσει ὁ Θεός;
\vs{24}Ἰδοὺ λαὸς ὡς σκύμνος ἀναστήσεται, καὶ ὡς λέων γαυρωθήσεται· οὐ κοιμηθήσεται ἕως φάγῃ θήραν, καὶ αἷμα τραυματιῶν πίεται.

\vs{25}Καὶ εἶπε Βαλὰκ πρὸς Βαλαὰμ, οὔτε κατάραις καταράσῃ μοι αὐτὸν, οὔτε εὐλογῶν μὴ εὐλογήσῃς αὐτόν.
\vs{26}Καὶ ἀποκριθεὶς Βαλαὰμ, εἶπε τῷ Βαλὰκ, οὐκ ἐλάλησά σοι, λέγων, τὸ ῥῆμα ὃ ἐὰν λαλήσῃ ὁ Θεὸς, τοῦτο ποιήσω;
\vs{27}Καὶ εἶπε Βαλὰκ πρὸς Βαλαὰμ, δεῦρο παραλάβω σε εἰς τόπον ἄλλον, εἰ ἀρέσει τῷ Θεῷ, καὶ κατάρασαί μοι αὐτὸν ἐκεῖθεν.
\vs{28}Καὶ παρέλαβε Βαλὰκ τὸν Βαλαὰμ ἐπὶ κορυφὴν τοῦ Φογὼρ, τὸ παρατεῖνον εἰς τὴν ἔρημον.
\vs{29}Καὶ εἶπε Βαλαὰμ πρὸς Βαλὰκ, οἰκοδόμησόν μοι ὧδε ἑπτὰ βωμοὺς, καὶ ἑτοίμασόν μοι ὧδε ἑπτὰ μόσχους, καὶ ἑπτὰ κριούς.
\vs{30}Καὶ ἐποίησε Βαλὰκ καθάπερ εἶπεν αὐτῷ Βαλαάμ, καί ἀνήνεγκε μόσχον καὶ κριὸν ἐπὶ τὸν βωμόν.

\ch{24}
Καὶ ἰδὼν Βαλαὰμ ὅτι καλόν ἐστιν ἐναντίον Κυρίου εὐλογεῖν τὸν Ἰσραὴλ, οὐκ ἐπορεύθη κατὰ τὸ εἰωθὸς αὐτῷ εἰς συνάντησιν τοῖς οἰωνοῖς, καὶ ἀπέστρεψε τὸ πρόσωπον αὐτοῦ εἰς τὴν ἔρημον.
\vs{2}Καὶ ἐξάρας Βαλαὰμ τοὺς ὀφθαλμοὺς αὐτοῦ, καθορᾷ τὸν Ἰσραὴλ ἐστρατοπεδευκότα κατὰ φυλάς· καὶ ἐγένετο ἐπʼ αὐτῷ πνεῦμα Θεου.
\vs{3}Καὶ ἀναλαβὼν τὴν παραβολὴν αὐτοῦ, εἶπε, φησὶ Βαλαὰμ υἱοῖς Βεὼρ, φησὶν ὁ ἄνθρωπος ὁ ἀληθινῶς ὁρῶν,
\vs{4}φησὶν ἀκούων λόγια ἰσχυροῦ, ὅστις ὅρασιν Θεοῦ εἶδεν ἐν ὕπνῳ· ἀποκεκαλυμμένοι οἱ ὀφθαλμοὶ αὐτοῦ.
\vs{5}Ὡς καλοὶ οἱ οἶκοί σου Ἰακὼβ, αἱ σκηναί σου Ἰσραήλ·
\vs{6}Ὡσεὶ νάπαι σκιάζουσαι, καὶ ὡσεὶ παράδεισοι ἐπὶ ποταμῷ, καὶ ὡσεὶ, σκηναὶ, ἃς ἔπηξε Κύριος, καί ὡσεὶ κέδροι παρʼ ὕδατα.
\vs{7}Ἐξελεύσεται ἄνθρωπος ἐκ τοῦ σπέρματος αὐτοῦ, καὶ κυριεύσει ἐθνῶν πολλῶν· καὶ ὑψωθήσεται ἡ Γὼγ βασιλεία, καὶ αὐξηθήσεται βασιλεία αὐτοῦ.
\vs{8}Θεὸς ὡδήγησεν αὐτὸν ἐξ Αἰγύπτου· ὡς δόξα μονοκέρωτος αὐτῷ· ἔδεται ἔθνη ἐχθρῶν αὐτοῦ, καὶ τὰ πάχη αὐτῶν ἐκμυελιεῖ, καὶ ταῖς βολίσιν αὐτοῦ κατατοξεύσει ἐχθρόν.
\vs{9}Κατακλιθεὶς ἀνεπαύσατο ὡς λέων, καὶ ὡς σκύμνος· τίς ἀναστήσει αὐτόν· οἱ εὐλογοῦντές σε, εὐλόγηνται· καὶ οἱ καταρώμενοί σε, κεκατήρανται.

\vs{10}Καὶ ἐθυμώθη Βαλὰκ ἐπὶ Βαλαὰμ, καὶ συνεκρότησε ταῖς χερσὶν αὐτοῦ, καὶ εἶπε Βαλὰκ πρὸς Βαλαὰμ, καταρᾶσθαι τὸν ἐχθρόν μου κέκληκά σε, καὶ ἰδοὺ εὐλογῶν εὐλόγησας τρίτον τοῦτο.
\vs{11}Νῦν οὖν φεῦγε εἰς τὸν τόπον σου· εἶπα, τιμήσω σε, καὶ νῦν ἐστέρησέ σε Κύριος τῆς δόξης.
\vs{12}Καὶ εἶπε Βαλαὰμ πρὸς Βαλὰκ, οὐχὶ καὶ τοῖς ἀγγέλοις σου οὓς ἀπέστειλας πρὸς με ἐλάλησα, λέγων,
\vs{13}ἐάν μοι δῷ Βαλὰκ πλήρη τὸν οἶκον αὐτοῦ ἀργυρίου καὶ χρυσίου, οὐ δυνήσομαι παραβῆναι τὸ ῥῆμα Κυρίου ποιῆσαι αὐτὸ καλὸν ἤ πονηρὸν παρʼ ἐμαυτοῦ· ὅσα ἂν εἴπῃ ὁ Θεός, ταῦτα ἐρῶ.
\vs{14}Καὶ νῦν ἰδοὺ ἀποτρέχω εἰς τὸν τόπον μου· δεῦρο, συμβουλεύσω σοι, τί ποιήσει ὁ λαὸς οὗτος τὸν λαόν σου ἐπʼ ἐσχάτου τῶν ἡμερῶν.

\vs{15}Καὶ ἀναλαβὼν τὴν παραβολὴν αὐτοῦ, εἶπε,

\vs{16}Φησὶ Βαλαὰμ υἱὸς Βεὼρ, φησὶν ὁ ἄνθρωπος ὁ ἀληθινῶς ὁρῶν, ἀκούων λόγια Θεοῦ, ἐπιστάμενος ἐπιστήμην παρὰ ὑψίστου, καὶ ὅρασιν Θεοῦ ἰδὼν ἐν ὕπνῳ· ἀποκεκαλυμμένοι οἱ ὀφθαλμοὶ αὐτοῦ.
\vs{17}Δείξω αὐτῷ, καὶ οὐχὶ νῦν· μακαρίζω, καὶ οὐκ ἐγγίζει· ἀνατελεῖ ἄστρον ἐξ Ἰακὼβ, ἀναστήσεται ἄνθρωπος ἐξ Ἰσραήλ· καὶ θραύσει τοὺς ἀρχηγοὺς Μωὰβ, καῒ προνομεύσει πάντας υἱοὺς Σήθ.
\vs{18}Καὶ ἔσται Ἐδὼμ κληρονομία, καὶ ἔσται κληρονομία Ἡσαῦ ὁ ἐχθρὸς αὐτοῦ· καὶ Ἰσραὴλ ἐποίησεν ἐν ἰσχύϊ.
\vs{19}Καὶ ἐξεγερθήσεται ἐξ Ἰακὼβ, καὶ ἀπολεῖ σωζόμενον ἐκ πόλεως.
\vs{20}Καὶ ἰδὼν τὸν Ἀμαλὴκ, καὶ ἀναλαβὼν τὴν παραβολὴν αὐτοῦ, εἶπεν, ἀρχὴ ἐθνῶν Ἀμαλὴκ, καὶ τὸ σπέρμα αὐτῶν ἀπολεῖται.
\vs{21}Καὶ ἰδὼν τὸν Κεναῖον, καὶ ἀναλαβὼν τὴν παραβολὴν αὐτοῦ, εἶπεν, ἰσχυρὰ ἡ κατοικία σου· καὶ ἐὰν θῇς ἐν πέτρᾳ τὴν νοσσίαν σου,
\vs{22}καὶ ἐὰν γένηται τῷ Βεὼρ νοσσιὰ πανουργίας, Ἀσσύριοι αἰχμαλωτεύσουσί σε.
\vs{23}Καὶ ἰδὼν τὸν Ὢγ, καὶ ἀναλαβὼν τὴν παραβολὴν αὐτοῦ, εἶπεν, ὢ ὢ, τίς ζήσεται, ὅταν θῇ ταῦτα ὁ Θεός;
\vs{24}Καὶ ἐξελεύσεται ἐκ χειρῶς Κιτιαίων, καὶ κακώσουσιν Ἀσσοὺρ, καὶ κακώσουσιν Ἐβραίους, καὶ αὐτοὶ ὁμοθυμαδὸν ἀπολοῦνται.
\vs{25}Καὶ ἀναστὰς Βαλαὰμ ἀπῆλθεν, ἀποστραφεὶς εἰς τὸν τόπον αὐτοῦ· καὶ Βαλὰκ ἀπῆλθε πρὸς ἑαυτόν.

\ch{25}
Καὶ κατέλυσεν Ἰσραὴλ ἐν Σαττεὶν, καὶ ἐβεβηλώθη ὁ λαὸς ἐκπορνεῦσαι εἰς τὰς θυγατέρας Μωάβ.
\vs{2}Καὶ ἐκάλεσαν αὐτοὺς εἰς τὰς θυσίας τῶν εἰδώλων αὐτῶν· καὶ ἔφαγεν ὁ λαὸς τῶν θυσιῶν αὐτῶν, καὶ προσεκύνησαν τοῖς εἰδώλοις αὐτῶν·
\vs{3}Καὶ ἐτελέσθη Ἰσραὴλ τῷ Βεελφεγώρ· καὶ ὠργίσθη θυμῷ Κύριος τῷ Ἰσραήλ.
\vs{4}Καὶ εἶπε Κύριος τῷ Μωυσῇ, λάβε πάντας τοὺς ἀρχηγοὺς τοῦ λαοῦ, καὶ παραδειγμάτισον αὐτοὺς Κυρίῳ κατέναντι τοῦ ἡλίου, καὶ ἀποστραφήσεται ὀργὴ θυμοῦ Κυρίου ἀπὸ Ἰσραήλ.
\vs{5}Καὶ εἶπε Μωυσῆς ταῖς φυλαῖς Ἰσραήλ, ἀποκτείνατε ἕκαστος τὸν οἰκεῖον αὐτοῦ τὸν τετελεσμένον τῷ Βεελφεγώρ.
\vs{6}Καὶ ἰδοὺ ἄνθρωπος τῶν υἱῶν Ἰσραὴλ ἐλθὼν προσήγαγε τὸν ἀδελφὸν αὐτοῦ πρὸς τὴν Μαδιανίτιν ἐναντίον Μωυσῆ, καὶ ἐναντίον πάσης συναγωγῆς υἱῶν Ἰσραήλ· αὐτοὶ δὲ ἔκλαιον παρὰ τὴν θύραν τῆς σκηνῆς τοῦ μαρτυρίου.
\vs{7}Καὶ ἰδὼν Φινεὲς υἱὸς Ἐλεάζαρ υἱοῦ Ἀαρὼν τοῦ ἱερέως, ἐξανέστη ἐκ μέσου τῆς συναγωγῆς, καὶ λαβὼν σειρομάστην ἐν τῇ χειρὶ,
\vs{8}εἰσῆλθεν ὀπίσω τοῦ ἀνθρώπου τοῦ Ἰσραηλίτου εἰς τὴν κάμινον, καὶ ἀπεκέντησεν ἀμφοτέρους, τόν τε ἄνθρωπον τὸν Ἰσραηλίτην, καὶ τὴν γυναῖκα διὰ τῆς μήτρας αὐτῆς· καὶ ἐπαύσατο ἡ πληγὴ ἀπὸ υἱῶν Ἰσραήλ.
\vs{9}Καὶ ἐγένοντο οἱ τεθνηκότες ἐν τῇ πληγῇ, τέσσαρες καὶ εἴκοσι χιλιάδες.

\vs{10}Καὶ ἐλάλησε Κύριος πρὸς Μωυσῆν, λέγων,
\vs{11}Φινεὲς υἱὸς Ἐλεάζαρ υἱοῦ Ἀαρὼν τοῦ ἱερέως κατέπαυσε τὸν θυμόν μου ἀπὸ υἱῶν Ἰσραὴλ, ἐν τῷ ζηλῶσαί μου τὸν ζῆλον ἐν αὐτοῖς, καὶ οὐκ ἐξανήλωσα τοὺς υἱοὺς Ἰσραὴλ ἐν τῷ ζήλῳ μου.
\vs{12}Οὕτως εἰπον, ἰδοὺ ἐγὼ δίδωμι αὐτῷ διαθήκην εἰρήνης,
\vs{13}καὶ ἔσται αὐτῷ καὶ τῷ σπέρματι αὐτοῦ μετʼ αὐτὸν διαθήκη ἱερατείας αἰωνία, ἀνθʼ ὧν ἐζήλωσεν τῷ Θεῷ αὐτοῦ, καὶ ἐξιλάσατο περὶ τῶν υἱῶν Ἰσραήλ.
\vs{14}Τὸ δὲ ὄνομα τοῦ ἀνθρώπου τοῦ Ἰσραηλείτου τοῦ πεπληγότος, ὃς ἐπλήγη μετὰ τῆς Μαδιανίτιδος, Ζαμβρὶ, υἱὸς Σαλμὼν, ἄρχων οἴκου πατριᾶς τῶν Συμεών.
\vs{15}Καὶ ὄνομα τῇ γυναικὶ τῇ Μαδιανίτιδι τῇ πεπληγυίᾳ, Χασβὶ, θυγάτηρ Σοὺρ, ἄρχοντος ἔθνους Ὀμμὼθ· οἴκου πατριᾶς ἐστι τῶν Μαδιάμ.

\vs{16}Καὶ ἐλάλησε Κύριος πρὸς Μωυσῆν, λέγων, λάλησον τοῖς υἱοῖς Ἰσραὴλ, λέγων,
\vs{17}ἐχθραίνετε τοῖς Μαδιηναίοις καὶ πατάξατε αὐτοὺς,
\vs{18}ὅτι ἐχθραίνουσιν αὐτοὶ ὑμῖν ἐν δολιότητι, ὅσα δολιοῦσιν ὑμᾶς διὰ Φογὼρ, καὶ διὰ Χασβὶ θυγατέρα ἄρχοντος Μαδιὰμ ἀδελφὴν αὐτῶν, τὴν πεπληγυῖαν ἐν τῇ ἡμέρᾳ τῆς πληγῆς διὰ Φογώρ.

\ch{26}
Καὶ ἐγένετο μετὰ τὴν πληγὴν, καὶ ἐλάλησε Κύριος πρὸς Μωυσῆν καὶ Ἐλεάζαρ τὸν ἱερέα, λέγων,
\vs{2}λάβε τὴν ἀρχὴν πάσης συναγωγῆς υἱῶν Ἰσραὴλ ἀπὸ εἰκοσαετοῦς καὶ ἐπάνω κατʼ οἴκους πατριῶν αὐτῶν, πᾶς ὁ ἐκπορευόμενος παρατάξασθαι ἐν Ἰσραήλ.

\vs{3}Καὶ ἐλάλησε Μωυσῆς καὶ Ἐλεάζαρ ὁ ἱερεὺς ἐν Ἀραβὼθ Μωἀβ ἐπὶ τοῦ Ἰορδάνου κατὰ Ἱεριχὼ, λέγων,
\vs{4}ἀπὸ εἰκοσαετοῦς καὶ ἐπάνω, ὃν τρόπον συνέταξε Κύριος τῷ Μωυσῇ· καὶ οἱ υἱοὶ Ἰσραὴλ οἱ ἐξελθόντες ἐξ Αἰγύπτου,
\vs{5}Ῥουβὴν πρωτότοκος Ἰσραήλ· υἱοὶ δὲ Ῥουβὴν Ἐνὼχ, καὶ δῆμος τοῦ Ἐνώχ· τῷ Φαλλοῦ, δῆμος τοῦ Φαλλουί.
\vs{6}Τῷ Ἀσρὼν, δῆμος τοῦ Ἀσρωνεί· τῷ Χαρμὶ, δῆμος τοῦ Χαρμί.
\vs{7}Οὗτοι δῆμοι Ῥουβήν· καὶ ἐγένετο ἡ ἐπίσκεψις αὐτῶν, τρεῖς καὶ τεσσαράκοντα χιλιάδες καὶ ἑπτακόσιοι καὶ τριάκοντα.

\vs{8}Καὶ υἱοῖ Φαλλοὺ, Ἑλιάβ.
\vs{9}Καὶ υἱοὶ Ἑλιὰβ, Ναμουὴλ, καὶ Δαθὰν, καὶ Ἀβειρών· οὗτοι ἐπίκλητοι τῆς συναγωγῆς· οὗτοί εἰσιν οἱ ἐπισυστάντες ἐπὶ Μωυσῆν καὶ Ἀαρὼν ἐν τῇ συναγωγῇ Κορὲ, ἐν τῇ ἐπισυστάσει Κυρίου.
\vs{10}Καὶ ἀνοίξασα ἡ γῆ τὸ στόμα αὐτῆς, κατέπιεν αὐτοὺς καὶ Κορὲ, ἐν τῷ θανάτῳ τῆς συναγωγῆ αὐτοῦ, ὅτε κατέφαγε τὸ πῦρ τοὺς πεντήκοντα καὶ διακοσίους, καὶ ἐγενήθησαν ἐν σημείῳ·
\vs{11}οἱ δὲ υἱοὶ Κορὲ οὐκ ἀπέθανον.

\vs{12}Καὶ οἱ υἱοὶ Συμεὼν, ὁ δῆμος τῶν υἱῶν Συμεών· τῷ Ναμουὴλ, δῆμος ὁ Ναμουηλί· τῷ Ἰαμὶν, δῆμος ὁ Ἰαμινί· τῷ Ἰαχὶν, δῆμος Ἰαχινί·
\vs{13}Τῷ Ζαρὰ δῆμος ὁ Ζαραΐ· τῷ Σαοὺλ, δῆμος ὁ Σαουλί.
\vs{14}Οὗτοι δῆμοι Συμεὼν ἐκ τῆς ἐπισκέψεως αὐτῶν, δύο καὶ εἴκοσι χιλιάδες καὶ διακόσιοι.

\vs{15}Υἱοὶ δὲ Ἰούδα, Ἢρ καὶ Αὐνάν· καὶ ἀπέθανον Ἤρ καὶ Αὐνὰν ἐν γῇ Χαναάν.
\vs{16}Καὶ ἐγένοντο οἱ υἱοὶ Ἰούδα κατὰ δήμους αὐτῶν· τῷ Σηλὼμ, δῆμος ὁ Σηλωνί· τῷ Φαρὲς, δῆμος ὁ Φαρεσί· τῷ Ζαρὰ, δῆμος ὁ Ζαραΐ.
\vs{17}Καὶ ἐγένοντο οἱ υἱοὶ Φαρὲς, τῷ Ἀσρὼν, δῆμος ὁ Ἀσρωνί· τῷ Ἰαμοῦν, δῆμος ὁ Ἰαμουνί.
\vs{18}Οὗτοι δῆμοι τοῦ Ἰούδα κατὰ τὴν ἐπίσκεψιν αὐτῶν, ἓξ καὶ ἑβδομήκοντα χιλιάδες καὶ πεντακόσιοι.

\vs{19}Καὶ υἱοὶ Ἰσσάχαρ κατὰ δήμους αὐτῶν· τῷ Θωλᾷ, δῆμος ὁ Θωλαΐ· τῷ Φουᾷ, δῆμος ὁ Φουαΐ·
\vs{20}Τῷ Ἰασοὺβ, δῆμος ὁ Ἰασουβί· τῷ Σαμρὰμ, δῆμος ὁ Σαμραμί.
\vs{21}Οὗτοι δῆμοι Ἰσσάχαρ ἐξ ἐπισκέψεως αὐτῶν, τέσσαρες καὶ ἑξήκοντα χιλιάδες καὶ τετρακόσιοι.

\vs{22}Υἱοὶ Ζαβουλὼν κατὰ δήμους αὐτῶν· τῷ Σαρὲδ, δῆμος ὁ Σαρεδί· τῷ Ἀλλὼν, δῆμος ὁ Ἀλλωνί· τῷ Ἀλλὴλ, δῆμος ὁ Ἀλληλί.
\vs{23}Οὗτοι δῆμοι Ζαβουλὼν ἐξ ἐπισκέψεως αὐτῶν, ἐξήκοντα χιλιάδες καὶ πεντακόσιοι.

\vs{24}Υἱοὶ Γὰδ κατὰ δήμους αὐτῶν· τῷ Σαφὼν, δῆμος ὁ Σαφωνί· τῷ Ἀγγὶ, δῆμος ὁ Ἀγγί· τῷ Σουνὶ, δῆμος ὁ Σουνί·
\vs{25}τῷ Ἀζενὶ, δῆμος ὁ Ἀζενί· τῷ Ἀδδὶ, δῆμος ὁ Ἀδδί·
\vs{26}τῷ Ἀροδὶ, δῆμος ὁ Αροαδί· τῷ Ἀριὴλ, δῆμος ὁ Ἀριηλί.
\vs{27}Οὗτοι δῆμοι υἱῶν Γὰδ ἐξ ἑπισκέψεως αὐτῶν, τέσσαρες καὶ τεσσαράκοντα χιλιάδες καὶ πεντακόσιοι.

\vs{28}Υἱοὶ Ἀσὴρ κατὰ δήμους αὐτῶν· τῷ Ἰαμὶν, δῆμος ὁ Ἰαμινί· τῷ Ἰεσοὺ, δῆμος ὁ Ἰεσουΐ· τῷ Βαριὰ, δῆμος ὁ Βαριαΐ.
\vs{29}Τῷ Χοβὲρ, δῆμος ὁ Χοβερί· τῷ Μελχιὴλ, δῆμος ὁ Μελχιηλί.
\vs{30}Καὶ τὸ ὄνομα θυγατρὸς Ἀσὴρ, Σάρα.
\vs{31}Οὗτοι δῆμοι Ἀσὴρ ἐξ ἐπισκέψεως αὐτῶν, τρεῖς καὶ τεσσαράκοντα χιλιάδες καὶ τετρακόσιοι.

\vs{32}Υἱοὶ Ἰωσὴφ κατὰ δήμους αὐτῶν, Μανασσῆ καὶ Ἐφραίμ.

\vs{33}Υἱοὶ Μανασσῆ. Τῷ Μαχὶρ, δῆμος ὁ Μαχιρί· καὶ Μαχὶρ ἐγέννησε τὸν Γαλαάδ· τῷ Γαλαὰδ, δῆμος ὁ Γαλααδί.
\vs{34}Καὶ οὗτοι υἱοὶ Γαλαάδ· Ἀχιέζερ, δῆμος ὁ Ἀχιεζερί· τῷ Χελὲγ, δῆμος ὁ Χελεγί.
\vs{35}Τῷ Ἐσριὴλ, δῆμος ὁ Ἐσριηλί· τῷ Συχὲμ, δῆμος ὁ Συχεμί·
\vs{36}Τῷ Συμαὲρ, δῆμος ὁ Συμαερί· καὶ τῷ Ὀφὲρ, δῆμος ὁ Ὀφερὶ.
\vs{37}Καὶ τῷ Σαλπαὰδ, υἱῷ Ὀφὲρ, οὐκ ἐγένοντο αὐτῷ υἱοὶ, ἀλλʼ ἢ θυγατέρες· καὶ ταῦτα τὰ ὀνόματα τῶν θυγατέρων Σαλπαάδ· Μαλὰ, καὶ Νουὰ, καὶ Ἐγλὰ, καὶ Μελχὰ, καὶ Θερσά.
\vs{38}Οὗτοι δῆμοι Μανασσὴ ἐξ ἐπισκέψεως αὐτῶν, δύο καὶ πεντήκοντα χιλιαδες καὶ ἑπτακόσιοι.

\vs{39}Καὶ οὗτοι υἱοὶ Ἐφραίμ· τῷ Σουθαλὰ, δῆμος ὁ Σουθαλάν· τῷ Τανὰχ, δῆμος ὁ Ταναχί.
\vs{40}Οὗτοι υἱοὶ Σουθαλά· τῷ Ἐδὲν, δῆμος ὁ Ἐδενί.
\vs{41}Οὗτοι δῆμοι Ἐφραὶμ ἐξ ἐπισκέψεως αὐτῶν, δύο καὶ τριάκοντα χιλιάδες καὶ πεντακόσιοι· οὗτοι δῆμοι υἱῶν Ἰωσὴφ κατὰ δήμους αὐτῶν.

\vs{42}Υἱοὶ Βενιαμὶν κατὰ δήμους αὐτῶν· τῷ Βαλὲ, δῆμος ὁ Βαλί· τῷ Ἀσυβὴρ, δῆμος ὁ Ἀσυβηρί· τῷ Ἰαχιρὰν, δῆμος ὁ Ἰαχιρανί.

\vs{43}Τῷ Σωφὰν, δῆμος ὁ Σωφανι.
\vs{44}Καὶ ἐγένοντο οἱ υἱοὶ Βαλὲ, Ἀδὰρ, καὶ Νοεμάν· τῷ Ἀδὰρ, δῆμος ὁ Αδαρί· καὶ τῷ Νοεμὰν, δῆμος ὁ Νοεμανί.
\vs{45}Οὗτοι υἱοὶ Βενιαμὶν κατὰ δήμους αὐτῶν ἐξ ἐπισκέψεως αὐτῶν, πέντε καὶ τριάκοντα χιλιάδες καὶ πεντακόσιοι.

\vs{46}Καὶ υἱοὶ Δὰν κατὰ δήμους αὐτῶν· τῷ Σαμὲ, δῆμος ὁ Σαμεΐ, οὗτοι δῆμοι Δὰν κατὰ δήμους αὐτῶν.
\vs{47}Πάντες οἱ δῆμοι Σαμεῒ κατʼ ἐπισκοπὴν αὐτῶν, τέσσαρες καὶ ἑξήκοντα χιλιάδες καὶ τετρακόσιοι.

\vs{48}Υἱοὶ Νεφθαλὶ κατὰ δήμους αὐτῶν· τῷ Ἀσιὴλ· δῆμος ὁ Ἀσιηλί· τῷ Γαυνὶ, δῆμος ὁ Γαυνί.
\vs{49}Τῷ Ἰεσὲρ, δῆμος ὁ Ἰεσερί· τῷ Σελλὴμ, δῆμος ὁ Σελλημί.
\vs{50}Οὗτοι δῆμοι Νεφθαλὶ ἐξ ἐπισκέψεως αὐτῶν, τεσσαράκοντα χιλιάδες καὶ τριακόσιοι.

\vs{51}Αὕτη ἡ ἐπίσκεψις υἱῶν Ἰσραὴλ, ἑξακόσιαι χιλιάδες καὶ χίλιοι καὶ ἑπτακόσιοι καὶ τριάκοντα.

\vs{52}Καὶ ἐλάλησε Κύριος πρὸς Μωυσῆν, λέγων,
\vs{53}τούτοις μερισθήσεται ἡ γῆ, κληρονομεῖν ἐξ ἀριθμοῦ ὀνομάτων.
\vs{54}Τοῖς πλείοσι πλεονάσεις τὴν κληρονομίαν, καὶ τοῖς ἐλάττοσιν ἐλαττώσεις τὴν κληρονομίαν αὐτῶν· ἑκάστῳ, καθὼς ἐπεσκέπησαν, δοθήσεται ἡ κληρονομία αὐτῶν.
\vs{55}Διὰ κλήρων μερισθήσεται ἡ γῆ τοῖς ὀνόμασι· κατὰ φυλὰς πατριῶν αὐτῶν κληρονομήσουσιν.
\vs{56}Ἐκ τοῦ κλήρου μεριεῖς τὴν κληρονομίαν αὐτῶν ἀναμέσον πολλῶν καὶ ὀλίγων.

\vs{57}Καὶ υἱοὶ Λευὶ κατὰ δήμους αὐτῶν· τῷ Γεδσὼν, δῆμος ὁ Γεδσωνί· τῷ Καὰθ, δῆμος ὁ Κααθί· τῷ Μεραρὶ, δῆμος ὁ Μεραρί.
\vs{58}Οὗτοι δῆμοι υἱῶν Λευί· δῆμος ὁ Λοβενί, δῆμος ὁ Χεβρωνὶ, δῆμος ὁ Κορὲ, καὶ δῆμος ὁ Μουσί· καὶ Καὰθ ἐγέννησε τὸν Ἁμράμ.
\vs{59}Τὸ δὲ ὄνομα τῆς γυναικὸς αὐτοῦ Ἰωχαβὲδ, θυγάτηρ Λευὶ, ἣ ἔτεκε τούτους τῷ Λευὶ ἐν Αἰγύπτῳ, καὶ ἔτεκε τῷ Ἀμρὰμ τὸν Ἀαρὼν καὶ Μωυσῆν, καὶ Μαριὰμ τὴν ἀδελφὴν αὐτῶν.
\vs{60}Καὶ ἐγενήθησαν τῷ Ἀαρὼν, ὅ, τε Ναδὰβ, καὶ Ἀβιοὺδ, καὶ Ἐλεάζαρ, καὶ Ἰθάμαρ.
\vs{61}Καὶ ἀπέθανε Ναδὰβ καὶ Ἀβιοὺδ ἐν τῷ προσφέρειν αὐτοὺς πῦρ ἀλλότριον ἔναντι Κυρίου ἐν τῇ ἐρήμῳ Σινᾷ.
\vs{62}Καὶ ἐγενήθησαν ἐξ ἐπισκέψεως αὐτῶν τρεῖς καὶ εἴκοσι χιλιάδες, πᾶν ἀρσενικὸν ἀπὸ μηνιαίου καὶ ἐπάνω· οὐ γὰρ συνεπεσκέπησαν ἐν μέσῳ υἱῶν Ἰσραὴλ, ὅτι οὐ δίδοται αὐτοῖς κλῆρος ἐν μέσῳ υἱῶν Ἰσραήλ.

\vs{63}Καὶ αὕτη ἡ ἐπίσκεψις Μωυσῆ καὶ Ἐλεάζαρ τοῦ ἱερέως, οἱ ἐπεσκέψαντο τοὺς υἱοὺς Ἰσραὴλ ἐν Ἀραβὼθ Μωὰβ, ἐπὶ τοῦ Ἰορδάνου κατὰ Ἱεριχώ.
\vs{64}Καὶ ἐν τούτοις οὐκ ἦν ἄνθρωπος τῶν ἐπεσκεμμένων ὑπὸ Μωυσῆ καὶ Ἀαρὼν, οὓς ἐπεσκέψαντο τοὺς υἱοὺς Ἰσραὴλ ἐν τῇ ἐρήμῳ Σινᾷ.
\vs{65}Ὅτι εἶπε Κύριος αὐτοῖς, θανάτῳ ἀποθανοῦνται ἐν τῇ ἐρήμῳ· καὶ οὐ κατελείφθη ἐξ αὐτῶν οὐδὲ εἷς, πλὴν Χάλεβ υἱὸς Ἰεφοννὴ, καὶ Ἰησοῦς ὁ τοῦ Ναυή.

\ch{27}
Καὶ προσελθοῦσαι αἱ θυγατέρες Σαλπαὰδ υἱοῦ Ὀφὲρ, υἱοῦ Γαλαὰδ, υἱοῦ Μαχιρ, τοῦ δήμου Μανασσῆ, τῶν υἱῶν Ἰωσὴφ, καὶ ταῦτα τὰ ὀνόματα αὐτῶν, Μααλὰ, καὶ Νουὰ, καὶ Ἐγλὰ, καὶ Μελχὰ, καὶ Θερσὰ,
\vs{2}καὶ στᾶσαι ἔναντι Μωυσῆ, καὶ ἔναντι Ἐλεάζαρ τοῦ ἱερέως, καὶ ἔναντι τῶν ἀρχόντων, καὶ ἔναντι πάσης συναγωγῆς ἐπὶ τῆς θύρας τῆς σκηνῆς τοῦ μαρτυρίου, λέγουσιν,
\vs{3}ὁ πατὴρ ἡμῶν ἀπέθανεν ἐν τῇ ἐρήμῳ, καὶ αὐτὸς οὐκ ἦν ἐν μέσῳ τῆς συναγωγῆς τῆς ἐπισυστάσης ἔναντι Κυρίου ἐν τῇ συναγωγῇ Κορὲ, ὅτι διʼ ἁμαρτίαν αὐτοῦ ἀπέθανε, καὶ υἱοὶ οὐκ ἐγένοντο αὐτῷ· μὴ ἐξαλειφθήτω τὸ ὄνομα τοῦ πατρὸς ἡμῶν ἐκ μέσου τοῦ δήμου αὐτοῦ, ὅτι οὐκ ἔστιν αὐτῷ υἱός· δότε ἡμῖν κατάσχεσιν ἐν μέσῳ ἀδελφῶν πατρὸς ἡμῶν.
\vs{4}Καὶ προσήγαγε Μωυσῆν τὴν κρίσιν αὐτῶν ἔναντι Κυρίου.

\vs{5}Καὶ ἐλάλησε Κύριος πρὸς Μωυσῆν,
\vs{6}λέγεν, ὀρθῶς θυγατέρες Σαλπαὰδ λελαλήκασι· δόμα δώσεις αὐταῖς κατάσχεσιν κληρονομίας ἐν μέσῳ ἀδελφῶν πατρὸς αὐτῶν, καὶ περιθήσεις τὸν κλῆρον τοῦ πατρὸς αὐτῶν αὐταῖς.
\vs{7}Καὶ τοῖς υἱοῖς Ἰσραὴλ λαλήσεις,
\vs{8}λέγων, ἄνθρωπος ἐὰν ἀποθάνῃ, καὶ υἱὸς μὴ ᾖ αὐτῷ, περιθήσετε τὴν κληρονομίαν αὐτοῦ τῇ θυγατρὶ αὐτοῦ·
\vs{9}Ἐὰν δὲ μὴ ᾖ θυγάτηρ αὐτῷ, δώσετε τὴν κληρονομίαν τῷ ἀδελφῷ αὐτοῦ.
\vs{10}Ἐὰν δὲ μὴ ὦσιν αὐτῷ ἀδελφοὶ, δώσετε τὴν κληρονομίαν τῷ ἀδελφῷ τοῦ πατρὸς αὐτοῦ.
\vs{11}Ἐὰν δὲ μὴ ὦσιν ἀδελφοὶ τοῦ πατρὸς αὐτοῦ, δώσετε τὴν κληρονομίαν τῷ οἰκείῳ τῷ ἔγγιστα αὐτοῦ ἐκ τῆς φυλῆς αὐτοῦ, κληρονομῆσαι τὰ αὐτοῦ· καὶ ἔσται τοῦτο τοῖς υἱοῖς Ἰσραὴλ δικαίωμα κρίσεως, καθὰ συνέταξε Κύριος τῷ Μωυσῇ.

\vs{12}Καὶ εἶπε Κύριος πρὸς Μωυσῆν, ἀνάβηθι εἰς τὸ ὄρος τὸ ἐν τῷ πέραν τοῦ Ἰορδάνου, τοῦτο τὸ ὄρος Ναβαὺ, καὶ ἴδε τὴν γῆν Χαναὰν, ἣν ἐγὼ δίδωμι τοῖς υἱοῖς Ἰσραὴλ ἐν κατασχέσει.
\vs{13}Καὶ ὄψῃ αὐτὴν, καὶ προστεθήσῃ πρὸς τὸν λαόν σου καὶ σὺ, καθὰ προσετέθη Ἀαρὼν ὁ ἀδελφός σου ἐν Ὢρ τῷ ὄρει.
\vs{14}Διότι παρέβητε τὸ ῥῆμά μου ἐν τῇ ἐρήμῳ Σὶν, ἐν τῷ ἀντιπίπτειν τὴν συναγωγὴν ἁγιάσαι με, οὐχ ἡγιάσατέ με ἐπὶ τῷ ὕδατι ἔναντι αὐτῶν· τοῦτʼ ἔστι τὸ ὕδωρ ἀντιλογίας ἐν Κάδης ἐν τῇ ἐρήμῳ Σίν.
\vs{15}Καὶ εἶπε Μωυσῆς πρὸς Κύριον,
\vs{16}ἐπισκεψάσθω Κύριος ὁ Θεὸς τῶν πνευμάτων καὶ πάσης σαρκὸς ἄνθρωπον ἐπὶ τῆς συναγωγῆς ταύτης,
\vs{17}ὅστις ἐξελεύσεται πρὸ προσώπου αὐτῶν, καὶ ὅστις εἰσελεύσεται πρὸ προσώπου αὐτῶν, καὶ ὅστις ἐξάξει αὐτοὺς, καὶ ὃστις εἰσάξει αὐτοὺς, καὶ οὐκ ἔσται ἡ συναγωγὴ Κυρίου ὡσεὶ πρόβατα οἷς οὐκ ἔστι ποιμήν.
\vs{18}Καὶ ἐλάλησε Κύριος πρὸς Μωυσῆν, λέγων, λάβε πρὸς σεαυτὸν Ἰησοῦν υἱὸν Ναυὴ, ἄνθρωπον ὃς ἔχει πνεῦμα ἐν ἑαυτῷ, καὶ ἐπιθήσεις τὰς χεῖράς σου ἐπʼ αὐτόν·
\vs{19}Καὶ στήσεις αὐτὸν ἔναντι Ἐλεάζαρ τοῦ ἱερέως, καὶ ἐντελῇ αὐτῷ ἔναντι πάσης συναγωγῆς, καὶ ἐντελῇ περὶ αὐτοῦ ἐναντίον αὐτῶν.
\vs{20}Καὶ δώσεις τῆς δόξης σου ἐπʼ αὐτὸν, ὅπως ἂν εἰσακούσωσιν αὐτοῦ οἱ υἱοὶ Ἰσραήλ.
\vs{21}Καὶ ἔναντι Ἐλεάζαρ τοῦ ἱερέως στήσεται, καὶ ἐπερωτήσουσιν αὐτὸν τὴν κρίσιν τῶν δήλων ἔναντι Κυρίου· ἐπὶ τῷ στόματι αὐτοῦ ἐξελεύσονται, καὶ ἐπὶ τῷ στόματι αὐτοῦ εἰσελεύσονται αὐτὸς καὶ οἱ υἱοὶ Ἰσραὴλ ὁμοθυμαδὸν, καὶ πᾶσα ἡ συναγωγή.

\vs{22}Καὶ ἐποίησε Μωυσῆς καθὰ ἐνετείλατο αὐτῷ Κύριος· καὶ λαβὼν τὸν Ἰησοῦν, ἔστησεν αὐτὸν ἐναντίον Ἐλεάζαρ τοῦ ἱερέως, καὶ ἐναντίον πάσης συναγωγῆς,
\vs{23}καὶ ἐπέθηκε τὰς χεῖρας αὐτοῦ ἐπʼ αὐτὸν, καὶ συνέστησεν αὐτὸν καθάπερ συνέταξε Κύριος τῷ Μωυσῇ.

\ch{28}
Καὶ ἐλάλησε Κύριος πρὸς Μωυσῆν, λέγων,
\vs{2}ἔντειλαι τοῖς υἱοῖς Ἰσραὴλ, καὶ ἐρεῖς πρὸς αὐτοὺς, λέγων, τὰ δῶρά μου δόματά μου καρπώματά μου εἰς ὀσμὴν εὐωδίας διατηρήσετε προσφέρειν ἐμοὶ ἐν ταῖς ἑορταῖς μου.
\vs{3}Καὶ ἐρεῖς πρὸς αὐτοὺς, ταῦτα τὰ καρπώματα ὅσα προσάξετε Κυρίῳ, ἀμνοὺς ἐνιαυσίους ἀμώμους δύο τὴν ἡμέραν εἰς ὁλοκαύτωσιν ἐνδελεχῶς.
\vs{4}Τὸν ἀμνὸν τὸν ἕνα ποιήσεις τὸ τοπρωῒ, καὶ τὸν ἀμνὸν τὸν δεύτερον ποιήσεις τὸ πρὸς ἑσπέραν.

\vs{5}Καὶ ποιήσεις τὸ δέκατον τοῦ οἰφὶ σεμίδαλιν εἰς θυσίαν ἀναπεποιημένην ἐν ἐλαίῳ ἐν τετάρτῳ τοῦ ἴν.
\vs{6}Ὁλοκαύτωμα ἐνδελεχισμοῦ, ἡ γενομένη ἐν τῷ ὄρει Σινᾷ εἰς ὀσμὴν εὐωδίας Κυρίῳ.
\vs{7}Καὶ σπονδὴν αὐτοῦ τὸ τέταρτον τοῦ ἴν τῷ ἀμνῷ τῷ ἑνί· ἐν τῷ ἁγίῳ σπείσεις σπονδὴν σίκερα Κυρίῳ·
\vs{8}καὶ τὸν ἀμνὸν τὸν δεύτερον ποιήσεις τὸ πρὸς ἑσπέραν· κατὰ τὴν θυσίαν αὐτοῦ καὶ κατὰ τὴν σπονδὴν αὐτοῦ ποιήσετε εἰς ὀσμὴν εὐωδίας Κυρίῳ.
\vs{9}Καὶ τῇ ἡμέρᾳ τῶν σαββάτων προσάξετε δύο ἀμνοὺς ἐνιαυσίους ἀμώμους, καὶ δύο δέκατα σεμιδάλεως ἀναπεποιημένης ἐν ἐλαίῳ εἰς θυσίαν καὶ σπονδὴν,
\vs{10}ὁλοκαύτωμα σαββάτων ἐν τοῖς σαββάτοις ἐπὶ τὴς ὁλοκαυτώσεως τῆς διαπαντὸς, καὶ τὴν σπονδὴν αὐτοῦ.

\vs{11}Καὶ ἐν ταῖς νεομηνίαις προσάξετε ὁλοκαύτωμα τῷ Κυρίῳ, μόσχους ἐκ βοῶν δύο, καὶ κριὸν ἕνα, ἀμνοὺς ἐνιαυσίους ἑπτὰ ἀμώμους·
\vs{12}Τρία δέκατα σεμιδάλεως ἀναπεποιημένης ἐν ἐλαίῳ τῷ μόσχῳ τῷ ἑνὶ, καὶ δύο δέκατα σεμιδάλεως ἀναπεποιημένης ἐν ἐλαίῳ τῷ κριῷ τῷ ἑνί·
\vs{13}Δέκατον δέκατον σεμιδάλεως ἀναπεποιημένης ἐν ἐλαίῳ τῷ ἀμνῷ τῷ ἑνὶ, θυσίαν ὀσμὴν εὐωδίας κάρπωμα Κυρίῳ.
\vs{14}Ἡ σπονδὴ αὐτῶν τὸ ἥμισυ τοῦ ἴν ἔσται τῷ μόσχῳ τῷ ἑνί· καὶ τὸ τρίτον τοῦ ἴν ἔσται τῷ κριῷ τῷ ἑνί. Καὶ τὸ τέταρτον τοῦ ἴν ἔσται τῷ ἀμνῷ τῷ ἑνὶ οἴνου· τοῦτο τὸ ὁλοκαύτωμα μῆνα ἐκ μηνὸς εἰς τοὺς μῆνας τοῦ ἐνιαυτοῦ.

\vs{15}Καὶ χίμαρον ἐξ αἰγῶν ἕνα περὶ ἁμαρτίας Κυρίῳ, ἐπὶ τῆς ὁλοκαυτώσεως τῆς διαπαντὸς ποιηθήσεται, καὶ ἡ σπονδὴ αὐτοῦ.

\vs{16}Καὶ ἐν τῷ μηνὶ τῷ πρώτῳ τεσσαρεσκαιδεκάτῃ ἡμέρᾳ τοῦ μηνὸς πάσχα Κυρίῳ.
\vs{17}Καὶ τῇ πεντεκαιδεκάτῃ ἡμέρᾳ τοῦ μηνὸς τούτου ἑορτή· ἑπτὰ ἡμέρας ἄζυμα ἔδεσθε.
\vs{18}Καὶ ἡ ἡμέρα ἡ πρώτη ἐπίκλητος ἁγία ἔσται ὑμῖν· πᾶν ἔργον λατρευτὸν οὐ ποιήσετε.
\vs{19}Καὶ προσάξετε ὁλοκαυτώματα κάρπωμα Κυρίῳ, μόσχους ἐκ βοῶν δύο, κριὸν ἕνα, ἀμνοὺς ἐνιαυσίους ἑπτά· ἄμωμοι ἔσονται ὑμῖν.
\vs{20}Καὶ θυσία αὐτῶν σεμίδαλις ἀναπεποιημένη ἐν ἐλαίῳ· τρία δέκατα τῷ μόσχῳ τῷ ἑνὶ, καὶ δύο δέκατα τῷ κριῷ τῷ ἑνί.
\vs{21}Δέκατον δέκατον ποιήσεις τῷ ἀμνῷ τῷ ἑνὶ, τοῖς ἑπτὰ ἀμνοῖς.
\vs{22}Καὶ χίμαρον ἐξ αἰγῶν ἕνα περὶ ἁμαρτίας, ἐξιλάσασθαι περὶ ὑμῶν·
\vs{23}Πλὴν τῆς ὁλοκαυτώσεως τῆς διαπαντὸς τῆς πρωϊνῆς, ὅ ἐστιν ὁλοκαύτωμα ἐνδελεχισμοῦ.
\vs{24}Ταῦτα κατὰ ταῦτα ποιήσετε τὴν ἡμέραν εἰς τὰς ἑπτὰ ἡμέρας, δῶρον κάρπωμα εἰς ὀσμὴν εὐωδίας Κυρίῳ, ἐπὶ τοῦ ὁλοκαυτώματος τοῦ διαπαντὸς ποιήσεις τὴν σπονδὴν αὐτοῦ.
\vs{25}Καὶ ἡμέρα ἡ ἑβδόμη κλητὴ ἁγία ἔσται ὑμῖν· πᾶν ἔργον λατρευτὸν οὐ ποιήσετε ἐν αὐτῇ.

\vs{26}Καὶ τῇ ἡμέρᾳ τῶν νέων, ὅταν προσφέρητε θυσίαν νέαν Κυρίῳ τῶν ἑβδομάδων, ἐπίκλητος ἁγία ἔσται ὑμῖν· πᾶν ἔργον λατρευτὸν οὐ ποιήσετε.
\vs{27}Καὶ προσάξετε ὁλοκαυτώματα εἰς ὀσμὴν εὐωδίας Κυρίῳ, μόσχους ἐκ βοῶν δύο, κριὸν ἕνα, ἀμνοὺς ἐνιαυσίους ἑπτὰ ἀμώμους.
\vs{28}Ἡ θυσία αὐτῶν σεμίδαλις ἀναπεποιημένη ἐν ἐλαίῳ· τρία δέκατα τῷ μόσχῳ τῷ ἑνὶ, καὶ δύο δέκατα τῷ κριῷ τῷ ἑνί.
\vs{29}Δέκατον δέκατον τῷ ἀμνῷ τῷ ἑνὶ, τοῖς ἑπτὰ ἀμνοῖς·
\vs{30}καὶ χίμαρον ἐξ αἰγῶν ἕνα περὶ ἁμαρτίας, ἐξιλάσασθαι περὶ ὑμῶν· πλὴν τοῦ ὁλοκαυτώματος τοῦ διαπαντός·
\vs{31}καὶ τὴν θυσίαν αὐτῶν ποιήσετέ μοι, ἄμωμοι ἔσονται ὑμῖν, καὶ τὰς σπονδὰς αὐτῶν.

\ch{29}
Καὶ τῷ μηνὶ τῷ ἑβδόμῳ, μιᾷ τοῦ μηνὸς, ἐπίκλητος ἁγία ἔσται ὑμῖν· πᾶν ἔργον λατρευτὸν οὐ ποιήσετε· ἡμέρα σημασίας ἔσται ὑμῖν.
\vs{2}Καὶ ποιήσετε ὁλοκαυτώματα εἰς ὀσμὴν εὐωδίας Κυρίῳ, μόσχον ἕνα ἐκ βοῶν, κριὸν ἕνα, ἀμνοὺς ἐνιαυσίους ἑπτὰ ἀμώμους.
\vs{3}Ἡ θυσία αὐτῶν σεμίδαλις ἀναπεποιημένη ἐν ἐλαίῳ· τρία δέκατα τῷ μόσχῳ τῷ ἑνὶ, καὶ δύο δέκατα τῷ κρίῳ τῷ ἑνί·
\vs{4}Δέκατον δέκατον τῷ ἀμνῷ τῷ ἑνὶ, τοῖς ἑπτὰ ἀμνοῖς·
\vs{5}Καὶ χίμαρον ἐξ αἰγῶν ἕνα περὶ ἁμαρτίας, ἐξιλάσασθαι περὶ ὑμῶν·
\vs{6}Πλὴν τῶν ὁλοκαυτωμάτων τῆς νουμηνίας· καὶ αἱ θυσίαι αὐτῶν, καὶ αἱ σπονδαὶ αὐτῶν, καὶ τὸ ὁλοκαύτωμα τὸ διαπαντός· καί αἱ θυσίαι αὐτῶν καὶ αἱ σπονδαὶ αὐτῶν κατὰ τὴν σύγκρισιν αὐτῶν εἰς ὀσμὴν εὐωδίας Κυρίῳ.

\vs{7}Καὶ τῇ δεκάτῃ τοῦ μηνὸς τούτου ἐπίκλητος ἁγία ἔσται ὑμῖν· καὶ κακώσετε τὰς ψυχὰς ὑμῶν, καὶ πᾶν ἔργον οὐ ποιήσετε.
\vs{8}Καὶ προσοίσετε ὁλοκαυτώματα εἰς ὀσμὴν εὐωδίας Κυρίῳ, καρπώματα Κυρίῳ, μόσχον ἐκ βοῶν ἕνα, κριὸν ἕνα, ἀμνους ἐνιαυσίους ἑπτά· ἄμωμοι ἔσονται ὑμῖν.
\vs{9}Ἡ θυσία αὐτῶν σεμίδαλις ἀναπεποιημένη ἐν ἐλαίῳ· τρία δέκατα τῷ μόσχῳ τῷ ἑνὶ, καὶ δύο δέκατα τῷ κριῷ τῷ ἑνί·
\vs{10}Δέκατον δέκατον τῷ ἀμνῷ τῷ ἑνὶ, εἰς τοὺς ἑπτὰ ἀμνούς·
\vs{11}Καὶ χίμαρον ἐξ αἰγῶν ἕνα περὶ ἁμαρτίας, ἐξιλάσασθαι περὶ ὑμῶν· πλὴν τὸ περὶ τῆς ἁμαρτίας τῆς ἐξιλάσεως, καὶ ἡ ὁλοκαύτωσις ἡ διαπαντός· ἡ θυσία αὐτῆς, καὶ ἡ σπονδὴ αὐτῆς κατὰ τὴν σύγκρισιν εἰς ὀσμὴν εὐωδίας κάρπωμα Κυρίῳ.

\vs{12}Καὶ τῇ πεντεκαιδεκάτῃ ἡμέρᾳ τοῦ μηνὸς τοῦ ἑβδόμου τούτου ἐπίκλητος ἁγία ἔσται ὑμῖν· πᾶν ἔργον λατρευτὸν οὐ ποιήσετε· καὶ ἑορτάσατε αὐτὴν ἑορτὴν Κυρίῳ ἑπτὰ ἡμέρας.
\vs{13}Καὶ προσάξετε ὁλοκαυτώματα κάρπωμα εἰς ὀσμὴν εὐωδίας Κυρίῳ, τῇ ἡμέρᾳ τῇ πρώτῃ μόσχους ἐκ βοῶν τρεῖς καὶ δέκα, κριοὺς δυο, ἀμνοὺς ἐνιαυσίους δεκατέσσαρας· ἄμωμοι ἔσονται.
\vs{14}Αἱ θυσίαι αὐτῶν σεμίδαλις ἀναπεποιημένη ἐν ἐλαίῳ· τρία δέκατα τῷ μόσχῳ τῷ ἑνὶ, τοῖς τρισκαίδεκα μόσχοις· καὶ δύο δέκατα τῷ κριῷ τῷ ἑνὶ, ἐπὶ τοὺς δύο κριούς·
\vs{15}Δέκατον δέκατον τῷ ἀμνῷ τῷ ἑνὶ, ἐπὶ τοὺς τέσσαρας καὶ δέκα ἀμνούς·
\vs{16}Καὶ χίμαρον ἐξ αἰγῶν ἕνα περὶ ἁμαρτίας· πλὴν τῆς ὁλοκαυτώσεως τῆς διαπαντός· αἱ θυσίαι αὐτῶν καὶ αἱ σπονδαὶ αὐτῶν.

\vs{17}Καὶ τῇ ἡμέρᾳ τῇ δευτέρᾳ μόσχους δώδεκα, κριοὺς δύο, ἀμνοὺς ἐνιαυσίους τέσσαρας καὶ δέκα ἀμώμους.
\vs{18}Ἡ θυσία αὐτῶν καὶ ἡ σπονδὴ αὐτῶν τοῖς μόσχοις καὶ τοῖς κριοῖς καὶ τοῖς ἀμνοῖς κατὰ ἀριθμὸν αὐτῶν, κατὰ τὴν σύγκρισιν αὐτῶν.
\vs{19}καὶ χίμαρον ἐξ αἰγῶν ἕνα περὶ ἁμαρτίας· πλὴν τῆς ὁλοκαυτώσεως τῆς διαπαντός· αἱ θυσίαι αὐτῶν καὶ αἱ σπονδαὶ αὐτῶν.

\vs{20}Τῇ ἡμέρᾳ τῇ τρίτῃ μόσχους ἕνδεκα, κριοὺς δύο, ἀμνοὺς ἐνιαυσίους τέσσαρας καὶ δέκα ἀμώμους.
\vs{21}Ἡ θυσία αὐτῶν καὶ ἡ σπονδὴ αὐτῶν τοῖς μόσχοις καὶ τοῖς κριοῖς καὶ τοῖς ἀμνοῖς κατὰ ἀριθμὸν αὐτῶν, κατὰ τὴν σύγκρισιν αὐτῶν.
\vs{22}Καὶ χίμαρον ἐξ αἰγῶν ἕνα περὶ ἁμαρτίας· πλὴν τῆς ὁλοκαυτώσεως τῆς διαπαντός· αἱ θυσίαι αὐτῶν καὶ αἱ σπονδαὶ αὐτῶν.

\vs{23}Τῇ ἡμέρᾳ τῇ τετάρτῃ μόσχους δέκα, κριοὺς δύο, ἀμνοὺς ἐνιαυσίους τέσσαρας καὶ δέκα ἀμώμους.
\vs{24}Αἱ θυσίαι αὐτῶν καὶ αἱ σπονδαὶ αὐτῶν τοῖς μόσχοις καὶ τοῖς κριοῖς καὶ τοῖς ἀμνοῖς κατὰ ἀριθμὸν αὐτῶν, κατὰ τὴν σύγκρισιν αὐτῶν.
\vs{25}Καὶ χίμαρον ἐξ αἰγῶν ἕνα περὶ ἁμαρτίας· πλὴν τῆς ὁλοκαυτώσεως τῆς διαπαντός· αἱ θυσίαι αὐτῶν καὶ αἱ σπονδαὶ αὐτῶν.

\vs{26}Τῇ ἡμέρᾳ τῇ πέμπτῃ μόσχους ἐννέα, κριοὺς δύο, ἀμνοὺς ἐνιαυσίους τέσσαρας καὶ δέκα ἀμώμους.
\vs{27}Αἱ θυσίαι αὐτῶν καὶ αἱ σπονδαὶ αὐτῶν τοῖς μόσχοις καὶ τοῖς κριοῖς καὶ τοῖς ἀμνοῖς κατὰ ἀριθμὸν αὐτῶν, κατὰ τὴν σύγκρισιν αὐτῶν.
\vs{28}Καὶ χίμαρον ἐξ αἰγῶν ἕνα περὶ ἁμαρτίας πλὴν τῆς ὁλοκαυτώσεως τῆς διαπαντός· αἱ θυσίαι αὐτῶν καὶ αἱ σπονδαὶ αὐτῶν.

\vs{29}Τῇ ἡμέρᾳ τῇ ἕκτῃ μόσχους ὀκτὼ, κριοὺς δύο, ἀμνοὺς ἐνιαυσίους δεκατέσσαρας ἀμώμους.
\vs{30}Αἱ θυσίαι αὐτῶν καὶ αἱ σπονδαὶ αὐτῶν τοῖς μόσχοις καὶ τοῖς κριοῖς καὶ τοῖς ἀμνοῖς κατὰ ἀριθμὸν αὐτῶν, κατὰ τὴν σύγκρισιν αὐτῶν.
\vs{31}Καὶ χίμαρον ἐξ αἰγῶν ἕνα περὶ ἁμαρτίας· πλὴν τῆς ὁλοκαυτώσεως τῆς διαπαντός· αἱ θυσίαι αὐτῶν καὶ αἱ σπονδαὶ αὐτῶν.

\vs{32}Τῇ ἡμέρᾳ τῇ ἑβδομῃ μόσχους ἑπτὰ, κριοὺς δύο, ἀμνοὺς ἐνιαυσίους δεκατέσσαρας ἀμώμους.
\vs{33}Αἱ θυσίαι αὐτῶν καὶ αἱ σπονδαὶ αὐτῶν τοῖς μόσχοις καὶ τοῖς κριοῖς καὶ τοῖς ἀμνοῖς κατὰ ἀριθμὸν αὐτῶν, κατὰ τὴν σύγκρισιν αὐτῶν.
\vs{34}Καὶ χίμαρον ἐξ αἰγῶν ἕνα περὶ ἁμαρτίας· πλὴν τῆς ὁλοκαυτώσεως τῆς διαπαντός· αἱ θυσίαι αὐτῶν καὶ αἱ σπονδαὶ αὐτῶν.
\vs{35}Καὶ τῇ ἡμέρᾳ τῇ ὀγδόῃ ἐξόδιον ἔσται ὑμῖν· πᾶν ἔργον λατρευτὸν οὐ ποιησετε ἐν αὐτῇ.
\vs{36}Καὶ προσάξετε ὁλοκαυτώματα εἰς ὀσμὴν εὐωδίας καρπώματα τῷ Κυρίῳ, μόσχον ἕνα, κριὸν ἕνα, ἀμνοὺς ἐνιαυσίους ἑπτὰ ἀμώμους.
\vs{37}Αἱ θυσίαι αὐτῶν καὶ αἱ σπονδαὶ αὐτῶν τῷ μόσχῳ καὶ τῷ κριῷ καὶ τοῖς ἀμνοῖς κατὰ ἀριθμὸν αὐτῶν, κατὰ τὴν σύνκρισιν αὐτῶν.
\vs{38}Καὶ χίμαρον ἐξ αἰγῶν ἕνα περὶ ἁμαρτίας· πλὴν τῆς ὁλοκαυτώσεως τῆς διαπαντός· αἱ θυσίαι αὐτῶν καὶ αἱ σπονδαὶ αὐτῶν.

\vs{39}Ταῦτα ποιήσετε Κυρίῳ ἐν ταῖς ἑορταῖς ὑμῶν, πλὴν τῶν εὐχῶν ὑμῶν, καὶ τὰ ἑκούσια ὑμῶν, καὶ τὰ ὁλοκαυτώματα ὑμῶν, καὶ τὰς θυσίας ὑμῶν, καὶ τὰς σπονδὰς ὑμῶν, καὶ τὰ σωτήρια ὑμῶν.

\ch{30}
Καὶ ἐλάλησε Μωυσῆς τοῖς υἱοῖς Ἰσραὴλ κατὰ πάντα ὅσα ἐνετείλατο Κύριος τῷ Μωυσῇ.
\vs{2}Καὶ ἐλάλησε Μωυσῆς πρὸς τοὺς ἄρχοντας τῶν φυλῶν υἱῶν Ἰσραὴλ, λέγων, τοῦτο τὸ ῥῆμα ὃ συνέταξε Κύριος.
\vs{3}Ἄνθρωπος ἄνθρωπος ὃς ἂν εὔξηται εὐχὴν Κυρίῳ, ἢ ὀμόσῃ ὅρκον, ἢ ὁρίσηται ὁρισμῷ περὶ τῆς ψυχῆς αὐτοῦ, οὐ βεβηλώσει τὸ ῥῆμα αὐτοῦ· πάντα ὅσα ἂν ἐξέλθῃ ἐκ τοῦ στόματος αὐτοῦ, ποιήσει.
\vs{4}Ἐὰν δὲ εὔξηται γυνὴ εὐχὴν Κυρίῳ, ἢ ὁρίσηται ὁρισμὸν ἐν τῷ οἴκῳ τοῦ πατρὸς αὐτῆς ἐν τῇ νεότητι αὐτῆς,
\vs{5}καὶ ἀκούσῃ ὁ πατὴρ αὐτῆς τὰς εὐχὰς αὐτῆς, καὶ τοὺς ὁρισμοὺς αὐτῆς, οὓς ὡρίσατο κατὰ τῆς ψυχῆς αὐτῆς, καὶ παρασιωπήσῃ αὐτῆς ὁ πατὴρ, καὶ στήσονται πᾶσαι αἱ εὐχαὶ αὐτῆς, καὶ πάντες οἱ ὁρισμοὶ οὓς ὡρίσατο κατὰ τῆς ψυχῆς αὐτῆς, μενοῦσιν αὐτῇ·
\vs{6}Ἐὰν δὲ ἀνανεύων ἀνανεύσῃ ὁ πατὴρ αὐτῆς, ᾗ ἂν ἡμέρᾳ ἀκούσῃ πάσας τὰς εὐχὰς αὐτῆς καὶ τοὺς ὁρισμοὺς, οὓς ὡρίσατο κατὰ τῆς ψυχῆς αὐτῆς, οὐ στήσονται· καὶ Κύριος καθαριεῖ αὐτὴν, ὅτι ἀνένευσεν ὁ πατὴρ αὐτῆς.

\vs{7}Ἐὰν δὲ γενομένη γένηται ἀνδρὶ, καὶ αἱ εὐχαὶ αὐτῆς ἐπʼ αὐτῇ κατὰ τὴν διαστολὴν τῶν χειλέων αὐτῆς, οὓς ὡρίσατο κατὰ τῆς ψυχῆς αὐτῆς,
\vs{8}καὶ ἀκούσῃ ὁ ἀνὴρ αὐτῆς, καὶ παρασιωπήσῃ αὐτῇ ᾗ ἂν ἡμέρᾳ ἀκούσῃ, καὶ οὕτω στήσονται πᾶσαι αἱ εὐχαὶ αὐτῆς, καὶ οἱ ὁρισμοὶ αὐτῆς, οὓς ὡρίσατο κατὰ τῆς ψυχῆς αὐτῆς, στήσονται.
\vs{9}Ἐὰν δὲ ἀνανεύων ἀνανεύσῃ ὁ ἀνὴρ αὐτῆς ᾗ ἐὰν ἡμέρᾳ ἀκούσῃ, πᾶσαι αἱ εὐχαὶ αὐτῆς, καὶ οἱ ὁρισμοὶ αὐτῆς οὓς ὡρίσατο κατὰ τῆς ψυχῆς αὐτῆς, οὐ μενοῦσιν, ὅτι ὁ ἀνὴρ ἀνένευσεν ἀπʼ αὐτῆς· καὶ Κύριος καθαριεῖ αὐτήν.

\vs{10}Καὶ εὐχὴ χήρας καὶ ἐκβεβλημένης ὅσα ἐὰν εὔξηται κατὰ τῆς ψυχῆς αὐτῆς, μενοῦσιν αὐτῇ.
\vs{11}Ἐὰν δὲ ἐν τῷ οἴκῳ τοῦ ἀνδρὸς αὐτῆς ἡ εὐχὴ αὐτῆς, ἢ ὁ ὁρισμὸς κατὰ τῆς ψυχῆς αὐτῆς μεθʼ ὅρκου,
\vs{12}καὶ ἀκούσῃ ὁ ἀνὴρ αὐτῆς, καὶ παρασιωπήσῃ αὐτῇ, καὶ μὴ ἀνανεύσῃ αὐτῇ, καὶ στήσονται πᾶσαι αἱ εὐχαὶ αὐτῆς, καὶ πάντες οἱ ὁρισμοὶ αὐτῆς οὓς ὡρίσατο κατὰ τῆς ψυχῆς αὐτῆς, στήσονται κατʼ αὐτῆς.
\vs{13}Ἐὰν δὲ περιελὼν περιέλῃ ὁ ἀνὴρ αὐτῆς ᾗ ἂν ἡμέρᾳ ἀκούσῃ, πάντα ὅσα ἐὰν ἐξέλθῃ ἐκ τῶν χειλέων αὐτῆς κατὰ τὰς εὐχὰς αὐτῆς, καὶ κατὰ τοὺς ὁρισμοὺς τοὺς κατὰ τῆς ψυχῆς αὐτῆς, οὐ μενεῖ αὐτῇ· ὁ ἀνὴρ αὐτῆς περιεῖλε, καὶ Κύριος καθαριεῖ αὐτήν.
\vs{14}Πᾶσα εὐχὴ καὶ πᾶς ὅρκος δεσμοῦ κακῶσαι ψυχὴν, ὁ ἀνὴρ αὐτῆς στήσει αὐτῇ, καὶ ὁ ἀνὴρ αὐτῆς περιελεῖ
\vs{15}Ἐὰν δὲ σιωπῶν παρασιωπήσῃ αὐτῇ ἡμέραν ἐξ ἡμέρας, καὶ στήσει αὐτῇ πάσας τὰς εὐχὰς αὐτῆς, καὶ τοὺς ὁρισμοὺς τοὺς ἐπʼ αὐτῆς στήσει αὐτῇ, ὅτι ἐσιώπησεν αὐτῇ τῇ ἡμέρᾳ ᾗ ἤκουσεν.
\vs{16}Ἐὰν δὲ περιελὼν περιέλῃ ὁ ἀνὴρ αὐτῆς μετὰ τὴν ἡμέραν ἣν ἤκουσε, καὶ λήψεται τὴν ἁμαρτίαν αὐτοῦ.
\vs{17}Ταῦτα τὰ δικαιώματα ὅσα ἐνετείλατο Κύριος τῷ Μωυσῇ, ἀναμέσον ἀνδρὸς καὶ γυναικὸς αὐτοῦ, καὶ ἀναμέσον πατρὸς καὶ θυγατρὸς ἐν νεότητι ἐ οἴκῳ πατρός.

\ch{31}
Καὶ ἐλάλησε Κύριος πρὸς Μωυσῆν, λέγων,
\vs{2}ἐκδίκει τὴν ἐκδίκησιν υἱῶν Ἰσραὴλ ἐκ τῶν Μαδιανιτῶν, καὶ ἔσχατον προστεθήσῃ πρὸς τὸν λαόν σου.
\vs{3}Καὶ ἐλάλησε Μωυσῆς πρὸς τὸν λαὸν, λέγων, ἐξοπλίσατε ἐξ ὑμῶν ἄνδρας, καὶ παρατάξασθε ἔναντι Κυρίου ἐπὶ Μαδιὰν, ἀποδοῦναι ἐκδίκησιν παρὰ τοῦ Κυρίου τῇ Μαδιάν.
\vs{4}Χιλίους ἐκ φυλῆς, χιλίους ἐκ φυλῆς, ἐκ πασῶν φυλῶν υἱῶν Ἰσραὴλ, ἀποστείλατε παρατάξασθαι.
\vs{5}Καὶ ἐξηρίθμησαν ἐκ τῶν χιλιάδων Ἰσραὴλ χιλίους ἐκ φυλῆς, δώδεκα χιλιάδας ἐνωπλισμένοι εἰς παράταξιν.
\vs{6}Καὶ ἀπέστειλεν αὐτοὺς Μωυσῆς χιλίους ἐκ φυλῆς, χιλίους ἐκ φυλῆς σὺν δυνάμει αὐτῶν, καὶ Φινεὲς υἱὸν Ἐλεάζαρ υἱοῦ Ἀαρὼν τοῦ ἱερέως· καὶ τὰ σκεύη τὰ ἅγια, καὶ αἱ σάλπιγγες τῶν σημασιῶν ἐν ταῖς χερσὶν αὐτῶν.

\vs{7}Καὶ παρετάξαντο ἐπὶ Μαδιὰν, καθὰ ἐνετείλατο Κύριος Μωυσῇ· καὶ ἀπέκτειναν πᾶν ἀρσενικόν.
\vs{8}Καὶ τοὺς βασιλεῖς Μαδιὰν ἀπέκτειναν ἅμα τοῖς τραυματίαις αὐτῶν· καὶ τὸν Εὐὶν, καὶ τὸν Ῥοκὸν, καὶ τὸν Σοὺρ, καὶ τὸν Οὒρ, καὶ τὸν Ῥοβὸκ, πέντε βασιλεῖς Μαδιάν· καὶ τὸν Βαλαὰμ υἱὸν Βεὼρ ἀπέκτειναν ἐν ῥομφαίᾳ σὺν τοῖς τραυματίαις αὐτῶν·
\vs{9}Καὶ ἐπρονόμευσαν τὰς γυναῖκας Μαδιὰν, καὶ τὴν ἀποσκευὴν αὐτῶν, καὶ τὰ κτήνη αὐτῶν, καὶ πάντα τὰ ἔγκτητα αὐτῶν· καὶ τὴν δύναμιν αὐτῶν ἐπρονόμευσαν·
\vs{10}Καὶ πάσας τὰς πόλεις αὐτῶν τὰς ἐν ταῖς κατοικίαις αὐτῶν, καὶ τὰς ἐπαύλεις αὐτῶν ἐνέπρησαν ἐν πυρί.
\vs{11}Καὶ ἔλαβον πᾶσαν τὴν προνομὴν αὐτῶν, καὶ πάντα τὰ σκῦλα αὐτῶν ἀπὸ ἀνθρώπου ἕως κτήνους.
\vs{12}Καὶ ἤγαγον πρὸς Μωυσῆν καὶ πρὸς Ἐλεάζαρ τὸν ἱερέα, καὶ πρὸς πάντας υἱοὺς Ἰσραὴλ, τὴν αἰχμαλωσίαν, καὶ τὰ σκύλα, καὶ τὴν προνομὴν εἰς τὴν παρεμβολὴν εἰς Ἀραβὼθ Μωὰβ, ἥ ἐστιν ἐπὶ τοῦ Ἰορδάνου κατὰ Ἰεριχώ.
\vs{13}Καὶ ἐξῆλθε Μωυσῆς καὶ Ἐλεάζαρ ὁ ἱερεὺς καὶ πάντες οἱ ἄρχοντες τῆς συναγωγῆς εἰς συνάντησιν αὐτοῖς ἔξω τῆς παρεμβολῆς.
\vs{14}Καὶ ὠργίσθη Μωυσῆς ἐπὶ τοῖς ἐπισκόποις τῆς δυνάμεως, χιλιάρχοις καὶ ἑκατοντάρχοις τοῖς ἐρχομένοις ἐκ τῆς παρατάξεως τοῦ πολέμου.
\vs{15}Καὶ εἶπεν αὐτοῖς Μωυσῆς, ἱνατί ἐζωγρήσατε πᾶν θῆλυ;
\vs{16}Αὗται γὰρ ἦσαν τοῖς υἱοῖς Ἰσραὴλ κατὰ τὸ ῥῆμα Βαλαὰμ τοῦ ἀποστῆσαι καὶ ὑπεριδεῖν τὸ ῥῆμα Κυρίου, ἕνεκεν Φογώρ· καὶ ἐγένετο ἡ πληγὴ ἐν τῇ συναγωγῇ Κυρίου.
\vs{17}Καὶ νῦν ἀποκτείνατε πᾶν ἀρσενικὸν ἐν πάσῃ τῇ ἀπαρτίᾳ, πᾶσαν γυναῖκα, ἥτις ἔγυω κοίτην ἄρσενος, ἀποκτείνατε.
\vs{18}Καὶ πᾶσαν τὴν ἀπαρτίαν τῶν γυναικῶν, ἥτις οὐκ οἶδε κοίτην ἄρσενος, ζωγρήσατε αὐτάς.
\vs{19}Καὶ ὑμεῖς παρεμβάλετε ἔξω τῆς παρεμβολῆς ἐπτὰ ἡμέρας· πᾶς ὁ ἀνελὼν καὶ ὁ ἁπτόμενος τοῦ τετρωμένου ἁγνισθήσεται τῇ ἡμέρᾳ τῇ τρίτῃ, καὶ τῇ ἡμέρᾳ τῇ ἑβδόμῃ ὑμεῖς καὶ ἡ αἰχμαλωσία ὑμῶν.
\vs{20}Καὶ πᾶν περίβλημα καὶ πᾶν σκεῦος δερμάτινον, καὶ πᾶσαν ἐργασίαν ἐξ αἰγείας, καὶ πᾶν σκεῦος ξύλινον ἀφαγνιεῖτε.

\vs{21}Καὶ εἶπεν Ἐλεάζαρ ὁ ἱερεὺς πρὸς τοὺς ἄνδρας τῆς δυνάμεως τοὺς ἐρχομένους ἐκ τῆς παρατάξεως τοῦ πολέμου, τοῦτο τὸ δικαίωμα τοῦ νόμου ὃ συνέταξε Κύριος τῷ Μωυσῇ.
\vs{22}Πλὴν τοῦ χρυσίου καὶ τοῦ ἀργυρίου καὶ χαλκοῦ καὶ σιδήρου καὶ μολίβου καὶ κασσιτέρου,
\vs{23}πᾶν πρᾶγμα ὃ διελεύσεται ἐν πυρὶ, καὶ καθαρισθήσεται, ἀλλʼ ἢ τῷ ὕδατι τοῦ ἁγνισμοῦ ἁγνισθήσεται· καὶ πάντα ὅσα ἐὰν μὴ διαπορεύηται διὰ πυρὸς, διελεύσεται διʼ ὕδατος.
\vs{24}Καὶ πλυνεῖσθε τὰ ἱμάτια τῇ ἡμέρᾳ τῇ ἑβδόμῃ, καὶ καθαρισθήσεσθε· καὶ μετὰ ταῦτα εἰσελεύσεσθε εἰς τὴν παρεμβολήν.

\vs{25}Καὶ ἐλάλησε Κύριος πρὸς Μωυσῆν, λέγων,
\vs{26}λάβε τὸ κεφάλαιον τῶν σκύλων τῆς αἰχμαλωσίας ἀπὸ ἀνθρώπου ἕως κτήνους σὺ καὶ Ἐλεάζαρ ὁ ἱερεὺς καὶ οἱ ἄρχοντες τῶν πατριῶν τῆς συναγωγῆς.
\vs{27}Καὶ διελεῖτε τὰ σκῦλα ἀναμέσον τῶν πολεμιστῶν τῶν ἐκπεπορευμένων εἰς τὴν παράταξιν, καὶ ἀναμέσον πάσης συναγωγῆς.
\vs{28}Καὶ ἀφελεῖτε τέλος Κυρίῳ παρὰ τῶν ἀνθρώπων τῶν πολεμιστῶν τῶν ἐκπεπορευμένων εἰς τὴν παράταξιν, μίαν ψυχὴν ἀπὸ πεντακοσίων, ἀπὸ τῶν ἀνθρώπων, καὶ ἀπὸ τῶν κτηνῶν, καὶ ἀπὸ τῶν βοῶν, καὶ ἀπὸ τῶν προβάτων, καὶ ἀπὸ τῶν ὄνων·
\vs{29}καὶ ἀπὸ τοῦ ἡμίσους αὐτῶν λήψεσθε. Καὶ δώσεις Ἐλεάζαρ τῷ ἱερεῖ τὰς ἀπαρχὰς Κυρίου.
\vs{30}Καὶ ἀπὸ τοῦ ἡμίσους τοῦ τῶν υἱῶν Ἰσραὴλ λήψῃ ἕνα ἀπὸ πεντήκοντα ἀπὸ τῶν ἀνθρώπων, καὶ ἀπὸ τῶν βοῶν, καὶ ἀπὸ τῶν προβάτων, καὶ ἀπὸ τῶν ὄνων, καὶ ἀπὸ πάντων τῶν κτηνῶν· καὶ δώσεις αὐτὰ τοῖς Λευίταις τοῖς φυλάσσουσι τὰς φυλακὰς ἐν τῇ σκηνῇ Κυρίου.

\vs{31}Καὶ ἐποίησε Μωυσῆς καὶ Ἐλεάζαρ ὁ ἱερεὺς, καθὰ συνέταξε Κύριος τῷ Μωυσῇ.
\vs{32}Καὶ ἐγενήθη τὸ πλεόνασμα τῆς προνομῆς ὁ προενόμευσαν οἱ ἄνδρες οἱ πολεμισταὶ, ἀπὸ τῶν προβάτων, ἑξακόσιαι χιλιάδες καὶ ἑβδομήκοντα καὶ πέντε χιλιάδες·
\vs{33}Καὶ βόες, δύο καὶ ἑβδομήκοντα χιλιάδες·
\vs{34}Καὶ ὄνοι, μία καὶ ἑξήκοντα χιλιάδες·
\vs{35}Καὶ ψυχαὶ ἀνθρώπων ἀπὸ τῶν γυναικῶν αἳ οὐκ ἔγνωσαν κοίτην ἀνδρὸς, πᾶσαι ψυχαὶ, δύο καὶ τριάκοντα χιλιάδες.
\vs{36}Καὶ ἐγενήθη τὸ ἡμίσευμα ἡ μερὶς τῶν ἐκπεπορευμένων εἰς τὸν πόλεμον ἐκ τοῦ ἀριθμοῦ τῶν προβάτων, τριακόσιαι καὶ τριάκοντα χιλιάδες καὶ ἑπτακισχίλια καὶ πεντακόσια.
\vs{37}Καὶ ἐγένετο τὸ τέλος Κυρίῳ ἀπὸ τῶν προβάτων, ἑξακόσιαι ἑβδομήκοντα πέντε·
\vs{38}Καὶ βόες, ἓξ καὶ τριάκοντα χιλιάδες, καὶ τὸ τέλος Κυρίῳ, δύο καὶ ἑβδομήκοντα·
\vs{39}Καὶ ὄνοι, τριάκοντα χιλιάδες καὶ πεντακόσιοι, καὶ τὸ τέλος Κυρίῳ, εἷς καὶ ἑξήκοντα·
\vs{40}Καὶ ψυχαὶ ἀνθρώπων, ἑκκαίδεκα χιλιάδες, καὶ τὸ τέλος αὐτῶν Κυρίῳ, δύο καὶ τριάκοντα ψυχαί.

\vs{41}Καὶ ἔδωκε Μωυσῆς τὸ τέλος Κυρίῳ τὸ ἀφαίρεμα τοῦ Θεοῦ Ἐλεάζαρ τῷ ἱερεῖ, καθὰ συνέταξε Κύριος τῷ Μωυσῇ·
\vs{42}ἀπὸ τοῦ ἡμισεύματος τῶν υἱῶν Ἰσραὴλ, οὓς διεῖλε Μωυσῆς ἀπὸ τῶν ἀνδρῶν τῶν πολεμιστῶν.
\vs{43}Καὶ ἐγένετο τὸ ἡμίσευμα ἀπὸ τῆς συναγωγῆς ἀπὸ τῶν προβάων, τριακόσιαι καὶ τριάκοντα χιλιάδες καὶ ἑπτακισχίλια καὶ πεντακόσια·
\vs{44}Καὶ βόες, ἓξ καὶ τριάκοντα χιλιάδες·
\vs{45}Ὄνοι, τριάκοντα χιλιάδες καὶ πεντακόσιοι·
\vs{46}Καὶ ψυχαὶ ἀνθρώπων, ἓξ καὶ δέκα χιλιάδες.
\vs{47}Καὶ ἔλαβε Μωυσῆς ἀπὸ τοῦ ἡμισεύματος τῶν υἱῶν Ἰσραὴλ τὸ ἓν ἀπὸ τῶν πεντήκοντα, ἀπὸ τῶν ἀνθρώπων καὶ ἀπὸ τῶν κτηνῶν, καὶ ἔδωκεν αὐτὰ τοῖς Λευίταις τοῖς φυλάσσουσι τὰς φυλακὰς τῆς σκηνῆς Κυρίου, ὃν τρόπον συνέταξε Κύριος τῷ Μωυσῇ.

\vs{48}Καὶ προσῆλθον πρὸς Μωυσῆν πάντες οἱ καθεσταμένοι εἰς τὰς χιλιαρχίας τῆς δυνάμεως, χιλίαρχοι καὶ ἑκατόνταρχοι,
\vs{49}καὶ εἶπαν πρὸς Μωυσῆν, Οἱ παῖδές σου εἰλήφασι τὸ κεφάλαιον τῶν ἀνδρῶν τῶν πολεμιστῶν τῶν παρʼ ἡμῖν, καὶ οὐ διαπεφώνηκεν ἀπʼ αὐτῶν οὐδὲ εἷς.
\vs{50}Καὶ προσενηνόχαμεν τὸ δῶρον Κυρίῳ, ἀνὴρ ὃ εὗρε σκεῦος χρυσοῦν καὶ χλιδῶνα καὶ ψέλλιον καὶ δακτύλιον καὶ περιδέξιον καὶ ἐμπλοκιον, ἐξιλάσασθαι περὶ ἡμῶν ἔναντι Κυρίου.
\vs{51}Καὶ ἔλαβε Μωυσῆς καὶ Ἐλεάζαρ ὁ ἱερεὺς τὸ χρυσίον παρʼ αὐτῶν πᾶν σκεῦος εἰργασμένον.
\vs{52}Καὶ ἐγένετο πᾶν τὸ χρυσίον τὸ ἀφαίρεμα ὃ ἀφεῖλον Κυρίῳ, ἑκκίδεκα χιλιάδες καὶ ἑπτακόσιοι καὶ πεντήκοντα σίκλοι παρὰ τῶν χιλιάρχων καὶ παρὰ τῶν ἑκατοντάρχων.
\vs{53}Καὶ οἱ ἄνδρες οἱ πολεμισταὶ ἐπρονόμευσαν ἕκαστος ἑαυτῷ.
\vs{54}Καὶ ἔλαβε Μωυσῆς καὶ Ἐλεάζαρ ὁ ἱερεὺς τὸ χρυσίον παρὰ τῶν χιλιάρχων καὶ παρὰ τῶν ἑκατοντάρχων, καὶ εἰσήνεγκεν αὐτὰ εἰς τὴν σκηνὴν τοῦ μαρτυρίου, μνημόσυνον τῶν υἱῶν Ἰσραὴλ ἔναντι Κυρίου.

\ch{32}
Καὶ κτήνη πλῆθος ἦν τοῖς υἱοῖς Ῥουβὴν καὶ τοῖς υἱοῖς Γὰδ, πλῆθος σφόδρα· καὶ εἶδον τὴν χώραν Ἰαζὴρ, καὶ τὴν χώραν Γαλαάδ· καὶ ἦν ὁ τόπος, τόπος κτήνεσι·
\vs{2}Καὶ προσελθόντες οἱ υἱοὶ Ῥουβὴν καὶ οἱ υἱοὶ Γὰδ, εἶπαν πρὸς Μωυσῆν καὶ πρὸς Ἐλεάζαρ τὸν ἱερέα καὶ πρὸς τοὺς ἄρχοντας τῆς συναγωγῆς, λέγοντες,
\vs{3}Ἀταρὼθ, καὶ Δαιβὼν, καὶ Ἰαζὴρ, καὶ Ναμρὰ, καὶ Ἐσεβὼν, καὶ Ἐλεαλὴ, καὶ Σεβαμὰ, καὶ Ναβαὺ, καὶ Βαιάν,
\vs{4}τὴν γῆν ἣν παραδέδωκε Κύριος ἐνώπιον τῶν υἱῶν Ἰσραήλ, γῆ κτηνοτρόφος ἐστὶ, καὶ τοῖς παισί σου κτήνη ὑπάρχει.
\vs{5}Καὶ ἔλεγον, εἰ εὕρομεν χάριν ἐνώπιόν σου, δοθήτω ἡ γῆ αὕτη τοῖς οἰκέταις σου ἐν κατασχέσει, καὶ μὴ διαβιβάσῃς ἡμᾶς τὸν Ἰορδάνην.

\vs{6}Καὶ εἶπε Μωυσῆς τοῖς υἱοῖς Γὰδ καὶ τοῖς υἱοῖς Ῥουβὴν, οἱ ἀδελφοὶ ὑμῶν πορεύονται εἰς τὸν πόλεμον, καὶ ὑμεῖς καθήσεσθε αὐτοῦ;
\vs{7}Καὶ ἱνατί διαστρέφετε τὰς διανοίας τῶν υἱῶν Ἰσραὴλ μὴ διαβῆναι εἰς τὴν γῆν, ἣν Κύριος δίδωσιν αὐτοῖς;
\vs{8}Οὐχ οὕτως ἐποίησαν οἱ πατέρες ὑμῶν, ὅτε ἀπέστειλα αὐτοὺς ἐκ Κάδης Βαρνὴ κατανοῆσαι τὴν γῆν;
\vs{9}καὶ ἀνέβησαν φάραγγα βότρυος, καὶ κατενόησαν τὴν γῆν, καὶ ἀπέστησαν τὴν καρδίαν τῶν υἱῶν Ἰσραὴλ, ὅπως μὴ εἰσέλθωσιν εἰς τὴν γῆν, ἣν ἔδωκε Κύριος αὐτοῖς.
\vs{10}Καὶ ὠργίσθη θυμῷ Κύριος ἐν τῇ ἡμέρᾳ ἐκείνῃ, καὶ ὤμοσε, λέγων,
\vs{11}εἰ ὄψονται οἱ ἄνθρωποι οὗτοι οἱ ἀναβάντες ἐξ Αἰγύπτου ἀπὸ εἰκοσαετοῦς καὶ ἐπάνω, οἱ ἐπιστάμενοι τὸ ἀγαθὸν καὶ τὸ κακὸν, τὴν γῆν ἣν ὤμοσα τῷ Ἁβραὰμ καὶ Ἰσαὰκ καὶ Ἰακὼβ, οὐ γὰρ συνεπηκολούθησαν ὀπίσω μου·

\vs{12}πλὴν Χάλεβ υἱὸς Ἰεφοννὴ ὁ διακεχωρισμένος, καὶ Ἰησοῦς ὁ τοῦ Ναυὴ, ὅτι συνεπηκολούθησεν ὀπίσω Κυρίου.
\vs{13}Καὶ ὠργίσθη θυμῷ Κύριος ἐπὶ τὸν Ἰσραήλ, καὶ κατερόμβευσεν αὐτοὺς ἐν τῇ ἐρήμῳ τεσσεράκοντα ἔτη, ἕως ἐξανηλώθη πᾶσα ἡ γενεὰ, οἱ ποιοῦντες τὰ πονηρὰ ἔναντι Κυρίου.
\vs{14}Ἰδοὺ ἀνέστητε ἀντὶ τῶν πατέρων ὑμῶν, σύντριμμα ἀνθρώπων ἁμαρτωλῶν, προσθεῖναι ἔτι ἐπὶ τὸν θυμὸν τῆς ὀργῆς Κυρίου ἐπὶ Ἰσραήλ.
\vs{15}Ὅτι ἀποστραφήσεσθε ἀπʼ αὐτοῦ προσθεῖναι ἔτι καταλιπεῖν αὐτὸν ἐν τῇ ἐρήμῳ, καὶ ἀνομήσετε εἰς ὅλην τὴν συναγωγὴν ταύτην.

\vs{16}Καὶ προσῆλθον αὐτῷ, καὶ ἔλεγον, ἐπαύλεις προβάτων οἰκοδομήσομεν ὧδε τοῖς κτήνεσιν ἡμῶν, καὶ πόλεις ταῖς ἀποσκευαῖς ἡμῶν.
\vs{17}Καὶ ἡμεῖς ἐνοπλισάμενοι προφυλακὴν πρότεροι τῶν υἱῶν Ἰσραὴλ, ἕως ἂν ἀγάγωμεν αὐτοὺς εἰς τὸν ἑαυτῶν τόπον· καὶ κατοικήσει ἡ ἀποσκευὴ ἡμῶν ἐν πόλεσι τετειχισμέναις διὰ τοὺς κατοικοῦντας τὴν γῆν.
\vs{18}Οὐ μὴ ἀποστραφῶμεν εἰς τὰς οἰκίας ἡμῶν ἕως ἂν καταμερισθῶσιν οἱ υἱοὶ Ἰσραὴλ, ἕκαστος εἰς τὴν κληρονομίαν αὐτοῦ.
\vs{19}Καὶ οὐκέτι κληρονομήσομεν ἐν αὐτοῖς ἀπὸ τοῦ πέραν τοῦ Ἰορδάνου καὶ ἐπέκεινα, ὅτι ἀπέχομεν τοὺς κλήρους ἡμῶν ἐν τῷ πέραν τοῦ Ἰορδάνου ἐν ἀνατολαῖς.

\vs{20}Καὶ εἶπε πρὸς αὐτοὺς Μωυσῆς, ἐὰν ποιήσητε κατὰ τὸ ῥῆμα τοῦτο, ἐὰν ἐξοπλίσησθε ἔναντι Κυρίου εἰς πόλεμον,
\vs{21}καὶ παρελεύσεται ὑμῶν πᾶς ὁπλίτης τὸν Ἰορδάνην ἔναντι Κυρίου, ἕως ἂν ἐκτριβῇ ὁ ἐχθρὸς αὐτοῦ ἀπὸ προσώπου αὐτοῦ,
\vs{22}καὶ κατακυριευθῇ ἡ γῆ ἔναντι Κυρίου, καὶ μετὰ ταῦτα ἀποστραφήσεσθε, καὶ ἔσεσθε ἀθῶοι ἔναντι Κυρίου, καὶ ἀπὸ Ἰσραήλ· καὶ ἔσται ἡ γῆ αὕτη ὑμῖν ἐν κατασχέσει ἔναντι Κυρίου.
\vs{23}Ἐὰν δὲ μὴ ποιήσητε οὕτως, ἁμαρτήσεσθε ἔναντι Κυρίου· καὶ γνώσεσθε τὴν ἁμαρτίαν ὑμῶν, ὅταν ὑμᾶς καταλάβῃ τὰ κακά.
\vs{24}Καὶ οἰκοδομήσετε ὑμῖν ἑαυτοῖς πόλεις τῇ ἀποσκευῇ ὑμῶν, καὶ ἐπαύλεις τοῖς κτήνεσιν ὑμῶν· καὶ τὸ ἐκπορευόμενον ἐκ τοῦ στόματος ὑμῶν ποιήσετε.

\vs{25}Καὶ εἶπαν υἱοὶ Ῥουβὴν καὶ υἱοὶ Γὰδ πρὸς Μωυσῆν, λέγοντες, οἱ παῖδές σου ποιήσουσι καθὰ ὁ Κύριος ἡμῶν ἐντέλλεται.
\vs{26}Ἡ ἀποσκευὴ ἡμῶν, καὶ αἱ γυναῖκες ἡμῶν, καὶ πάντα τὰ κτήνη ἡμῶν ἔσονται ἐν ταῖς πόλεσιν Γαλαάδ.
\vs{27}Οἱ δὲ παῖδές σου παρελεύσονται πάντες ἐνωπλισμένοι καὶ ἐκτεταγμένοι ἔναντι Κυρίου εἰς τὸν πόλεμον, ὃν τρόπον ὁ κύριος λέγει.

\vs{28}Καὶ συνέστησεν αὐτοῖς Μωυσῆς Ἐλεάζαρ τὸν ἱερέα, καὶ Ἰησοῦν υἱὸν Ναυὴ, καὶ τοὺς ἄρχοντας πατριῶν τῶν φυλῶν Ἰσραήλ.
\vs{29}Καὶ εἶπε πρὸς αὐτοὺς Μωυσῆς, ἐὰν διαβῶσιν οἱ υἱοὶ Ῥουβὴν καὶ οἱ υἱοὶ Γὰδ μεθʼ ὑμῶν τὸν Ἰορδάνην, πᾶς ἐνωπλισμένος εἰς πόλεμον ἔναντι Κυρίου, καὶ κατακυριεύσητε τῆς γῆς ἀπέναντι ὑμῶν, καὶ δώσετε αὐτοῖς τὴν γῆν Γαλαὰδ ἐν κατασχέσει.
\vs{30}Ἐὰν δὲ μὴ διαβῶσιν ἐνωπλισμένοι μεθʼ ὑμῶν εἰς τὸν πόλεμον ἔναντι Κυρίου, καὶ διαβιβάσετε τὴν ἀποσκευὴν αὐτῶν, καὶ τὰς γυναῖκας αὐτῶν, καὶ τὰ κτήνη αὐτῶν πρότερα ὑμῶν εἰς γῆν Χαναὰν, καὶ συγκατακληρονομηθήσονται ἐν ὑμῖν ἐν τῇ γῇ Χαναάν.
\vs{31}Καὶ ἀπεκρίθησαν οἱ υἱοὶ Ῥουβὴν καὶ οἱ υἱοὶ Γὰδ λέγοντες, ὅσα ὁ Κύριος λέγει τοῖς θεράπουσιν, οὕτω ποιήσομεν ἡμεῖς.
\vs{32}Διαβησόμεθα ἐνωπλισμένοι ἔναντι Κυρίου εἰς γῆν Χαναὰν, καὶ δώσετε τὴν κατάσχεσιν ἡμῖν ἐν τῷ πέραν τοῦ Ἰορδάνου.

\vs{33}Καὶ ἔδωκεν αὐτοῖς Μωυσῆς τοῖς υἱοῖς Γὰδ, καὶ τοῖς υἱοῖς Ῥουβὴν, καὶ τῷ ἡμίσει φυλῆς Μανασσῆ υἱῶν Ἰωσὴφ, τὴν βασιλείαν Σηὼν βασιλέως Ἀμοῤῥαίων, καὶ τὴν βασιλείαν Ὢγ βασιλέως τῆς Βασὰν, τὴν γῆν καὶ τὰς πόλεις σὺν τοῖς ὁρίοις αὐτῆς, πόλεις τῆς γῆς κύκλῳ.
\vs{34}Καὶ ᾠκοδόμησαν οἱ υἱοὶ Γὰδ τὴν Δαιβὼν, καὶ τὴν Ἀταρὼθ, καὶ τὴν Ἀροὴρ,
\vs{35}καὶ τὴν Σοφὰρ, καὶ τὴν Ἰαζὴρ, καὶ ὕψωσαν αὐτὰς,
\vs{36}καὶ τὴν Ναμρὰμ, καὶ τὴν Βαιθαρὰν, πόλεις ὀχυρὰς, καὶ ἐπαύλεις προβάτων.
\vs{37}Καὶ οἱ υἱοὶ Ῥουβὴν ᾠκοδόμησαν τὴν Ἐσεβὼν, καὶ Ἐλεάλην, καὶ Καριαθὰμ,
\vs{38}καὶ τὴν Βεελμεὼν, περικεκυκλωμένας, καὶ τὴν Σεβαμά· καὶ ἐπωνόμασαν κατὰ τὰ ὀνόματα αὐτῶν τὰ ὀνόματα τῶν πόλεων, ἃς ᾠκοδόμησαν.
\vs{39}Καὶ ἐπορεύθη υἱὸς Μαχὶρ υἱοῦ Μανασσὴ Γαλαὰδ, καὶ ἔλαβεν αὐτὴν, καὶ ἀπώλεσε τὸν Ἀμοῤῥαῖον τὸν κατοικοῦντα ἐν αὐτῇ.
\vs{40}Καὶ ἔδωκε Μωυσῆς τὴν Γαλαὰδ τῷ Μαχὶρ υἱῷ Μανασσῆ, καὶ κατῴκησεν ἐκεῖ.
\vs{41}Καὶ Ἰαῒρ ὁ τοῦ Μανασσῆ ἐπορεύθη, καὶ ἔλαβε τὰς ἐπαύλεις αὐτῶν, καὶ ἐπωνόμασεν αὐτὰς ἐπαύλεις Ἰαΐρ.
\vs{42}Καὶ Ναβαὺ ἐπορεύθη, καὶ ἔλαβε τὴν Καὰθ καὶ τὰς κώμας αὐτῆς, καὶ ἐπωνόμασεν αὐτὰς Ναβὼθ ἐκ τοῦ ὀνόματος αὐτοῦ.

\ch{33}
Καὶ οὗτοι οἱ σταθμοὶ τῶν υἱῶν Ἰσραὴλ, ὡς ἐξῆλθον ἐκ γῆς Αἰγύπτου σὺν δυνάμει αὐτῶν ἐν χειρὶ Μωυσῆ καὶ Ἀαρών.
\vs{2}Καὶ ἔγραψε Μωυσῆς τὰς ἀπάρσεις αὐτῶν, καὶ τοὺς σταθμοὺς αὐτῶν, διὰ ῥήματος Κυρίου· καὶ οὗτοι σταθμοὶ τῆς πορείας αὐτῶν.
\vs{3}Ἀπῆραν ἐκ Ῥαμεσσὴ τῷ μηνὶ τῷ πρώτῳ, τῇ πεντεκαιδεκάτῃ ἡμέρᾳ τοῦ μηνὸς τοῦ πρώτου· τῇ ἐπαύριον τοῦ πάσχα ἐξῆλθον οἱ υἱοὶ Ἰσραὴλ ἐν χειρὶ ὑψηλῇ ἐναντίον πάντων τῶν Αἰγυπτίων.
\vs{4}Καὶ οἱ Αἰγύπτιοι ἔθαπτον ἐξ αὐτῶν τοὺς τεθνηκότας πάντας οὓς ἐπάταξε Κύριος, πᾶν πρωτότοκον ἐν γῇ Αἰγύπτῳ· καὶ ἐν τοῖς θεοῖς αὐτῶν ἐποίησε τὴν ἐκδίκησιν Κύριος.
\vs{5}Καὶ ἀπάραντες οἱ υἱοὶ Ἰσραὴλ ἐκ Ῥαμεσσῆ, παρενέβαλον εἰς Σοκχώθ.
\vs{6}Καὶ ἀπάραντες ἐκ Σοκχὼθ, παρενέβαλον εἰς Βουθὰν, ὅ ἐστι μέρος τι τῆς ἐρήμου·
\vs{7}Καὶ ἀπῇραν ἐκ Βουθὰν, καὶ παρενέβαλον ἐπὶ τὸ στόμα Εἰρὼθ, ὅ ἐστιν ἀπέναντι Βεελσεπφὼν, καὶ παρενέβαλον ἀπέναντι Μαγδώλου.
\vs{8}Καὶ ἀπῇραν ἀπέναντι Εἰρὼθ, καὶ διέβησαν μέσον τῆς θαλάσσης εἰς τὴν ἔρημον· καὶ ἐπορεύθησαν ὁδὸν τριῶν ἡμερῶν διὰ τῆς ἐρήμου αὐτοὶ, καὶ παρενέβαλον ἐν Πικρίαις.
\vs{9}Καὶ ἀπῇραν ἐκ Πικριῶν, καὶ ἦλθον εἰς Αἰλίμ· καὶ ἐν Αἰλὶμ δώδεκα πηγαὶ ὑδάτων, καὶ ἑβδομήκοντα στελέχη φοινίκων, καὶ παρενέβαλον ἐκεῖ παρὰ τὸ ὕδωρ.
\vs{10}Καὶ ἀπῇραν ἐξ Αἰλὶμ, καὶ παρενέβαλον ἐπὶ θάλασσαν ἐρυθράν.
\vs{11}Καὶ ἀπῇραν ἀπὸ θαλάσσης ἐρυθρᾶς, καὶ παρενέβαλον εἰς τὴν ἔρημον Σίν.

\vs{12}Καὶ ἀπῇραν ἐκ τῆς ἐρήμου Σὶν, καὶ παρενέβαλον εἰς Ῥαφακά.
\vs{13}Καὶ ἀπῇραν ἐκ Ῥαφακὰ, καὶ παρενέβαλον ἐν Αἰλούς.
\vs{14}Καὶ ἀπῇραν ἐξ Αἰλούς, καὶ παρενέβαλον ἐν Ῥαφιδίν· καὶ οὐκ ἦν ἐκεῖ ὕδωρ τῷ λαῷ πιεῖν.
\vs{15}Καὶ ἀπῇραν ἐκ Ῥαφιδὶν, καὶ παρενέβαλον ἐν τῇ ἐρήμῳ Σινᾶ.
\vs{16}Καὶ ἀπῇραν ἐκ τῆς ἐρήμου Σινᾶ, καὶ παρενέβαλον ἐν μνήμασι τῆς ἐπιθυμίας.
\vs{17}Καὶ ἀπῇραν ἐκ μνημάτων τῆς ἐπιθυμίας, καὶ παρενέβαλον ἐν Ἀσηρώθ.
\vs{18}Καὶ ἀπῇραν ἐξ Ἀσηρὼθ, καὶ παρενέβαλον ἐν Ῥαθαμᾷ.

\vs{19}Καὶ ἀπῇραν ἐκ Ῥαθαμᾶ, καὶ παρενέβαλον ἐν Ῥεμμὼν Φαρές.
\vs{20}Καὶ ἀπῇραν ἐκ Ῥεμμὼν Φαρὲς, καὶ παρενέβαλον εἰς Λεβωνᾶ.
\vs{21}Καὶ ἀπῇραν ἐκ Λεβῶνα, καὶ παρενέβαλον εἰς Ῥεσσάν.
\vs{22}Καὶ ἀπῇραν ἐκ Ῥεσσὰν, καὶ παρενέβαλον εἰς Μακελλάθ.
\vs{23}Καὶ ἀπῇραν ἐκ Μακελλὰθ, καὶ παρενέβαλον εἰς Σαφάρ.
\vs{24}Καὶ ἀπῇραν ἑκ Σαφὰρ, καὶ παρενέβαλον εἰς Χαραδάθ.
\vs{25}Καὶ ἀπῇραν ἐκ Χαραδὰθ, καὶ παρενέβαλον εἰς Μακηλώθ.
\vs{26}Καὶ ἀπῇραν ἐκ Μακηλὼθ, καὶ παρενέβαλον εἰς Καταάθ.
\vs{27}Καὶ ἀπῇραν ἐκ Καταὰθ, καὶ παρενέβαλον εἰς Ταράθ.
\vs{28}Καὶ ἀπῇραν ἐκ Ταρὰθ, καὶ παρενέβαλον εἰς Μαθεκκά.
\vs{29}Καὶ ἀπῇραν ἐκ Μαθεκκὰ, καὶ παρενέβαλον εἰς Σελμωνᾶ.
\vs{30}Καὶ ἀπῇραν ἐκ Σελμωνᾶ, καὶ παρενέβαλον εἰς Μασουρούθ.
\vs{31}Καὶ ἀπῇραν ἐκ Μασουροὺθ, καὶ παρενέβαλον εἰς Βαναία.
\vs{32}Καὶ ἀπῇραν ἐκ Βαναία, καὶ παρενέβαλον εἰς τὸ ὄρος Γαδγάδ.

\vs{33}Καὶ ἀπῇραν ἐκ τοῦ ὄρους Γαδγὰδ, καὶ παρενέβαλον εἰς Ἐτεβαθά.
\vs{34}Καὶ ἀπῇραν ἐξ Ἐτεβαθὰ, καὶ παρενέβαλον εἰς Ἐβρωνά.
\vs{35}Καὶ ἀπῇραν ἐξ Ἐβρωνὰ, καὶ παρενέβαλον εἰς Γεσιὼν Γάβερ.
\vs{36}Καὶ ἀπῇραν ἐκ Γεσιὼν Γάβερ, καὶ παρενέβαλον ἐν τῇ ἐρήμῳ Σίν· καὶ ἀπῇραν ἐκ τῆς ἐρήμου Σὶν, καὶ παρενέβαλον εἰς τὴν ἔρημον Φαράν· αὕτη ἐστὶ Κάδης.
\vs{37}Καὶ ἀπῇραν ἐκ Κάδης, καὶ παρενέβαλον εἰς Ὢρ τὸ ὄρος πλησίον γῆς Ἐδώμ.

\vs{38}Καὶ ἀνέβη Ἀαρὼν ὁ ἱερεὺς διὰ προστάγματος Κυρίου, καὶ ἀπέθανεν ἐκεῖ ἐν τῷ τεσσαρακοστῷ ἔτει τῆς ἐξόδου τῶν υἱῶν Ἰσραὴλ ἐκ γῆς Αἰγύπτου, τῷ μηνὶ τῷ πέμπτῳ μιᾷ τοῦ μηνός.
\vs{39}Καὶ Ἀαρὼν ἦν τριῶν καὶ εἴκοσι καὶ ἑκατὸν ἐτῶν, ὅτε ἀπέθνησκεν ἐν Ὢρ τῷ ὄρει.
\vs{40}Καὶ ἀκούσας ὁ Χανανὶς βασιλεὺς Ἀρὰδ, καὶ οὗτος κατῴκει ἐν γῇ Χαναὰν, ὅτε εἰσεπορεύοντο οἱ υἱοὶ Ἰσραήλ·
\vs{41}καὶ ἀπῇραν ἐξ Ὢρ τοῦ ὄρους, καὶ παρενέβαλον εἰς Σελμωνᾶ.
\vs{42}Καὶ ἀπῇραν ἐκ Σελμωνᾶ, καὶ παρενέβαλον εἰς Φινώ.
\vs{43}Καὶ ἀπῇραν ἐκ Φινὼ, καὶ παρενέβαλον ἐν Ὠβώθ.

\vs{44}Καὶ ἀπῇραν ἐξ Ὠβὼθ, καὶ παρενέβαλον ἐν Γαῒ, ἐν τῷ πέραν ἐπὶ τῶν ὁρίων Μωάβ.
\vs{45}Καὶ ἀπῇραν ἐκ Γαῒ, καὶ παρενέβαλον εἰς Δαιβὼν Γάδ.
\vs{46}Καὶ ἀπῇραν ἐκ Δαιβὼν Γὰδ, καὶ παρενέβαλον ἐν Γελμὼν Δεβλαθαίμ.
\vs{47}Καὶ ἀπῇραν ἐκ Γελμὼν Δεβλαθαὶμ, καὶ παρενέβαλον ἐπὶ τὰ ὄρη τὰ Ἀβαρὶμ, ἀπέναντι Ναβαῦ.
\vs{48}Καὶ ἀπῇραν ἀπὸ ὀρέων Ἀβαρὶμ, καὶ παρενέβαλον ἐπὶ δυσμῶν Μωὰβ, ἐπὶ τοῦ Ἰορδάνου κατὰ Ἱεριχώ.
\vs{49}Καὶ παρενέβαλον παρὰ τὸν Ἰορδάνην ἀναμέσον Αἰσιμώθ, ἕως Βελσᾶ τὸ κατὰ δυσμὰς Μωάβ.

\vs{50}Καὶ ἐλάλησε Κύριος πρὸ; Μωυσῆν ἐπὶ δυσμῶν Μωὰβ παρὰ τὸν Ἰορδάνην κατὰ Ἱεριχὼ, λέγων,
\vs{51}λάλησον τοῖς υἱοῖς Ἰσραὴλ, καὶ ἐρεῖς πρὸς αὐτοὺς, ὑμεῖς διαβαίνετε τὸν Ἰορδάνην εἰς γῆν Χαναάν.
\vs{52}Καὶ ἀπολεῖτε πάντας τοὺς κατοικοῦντας ἐν τῇ γῇ πρὸ προσώπου ὑμῶν, καὶ ἐξαρεῖτε τὰς σκοπιὰς αὐτῶν, καὶ πάντα τὰ εἴδωλα τὰ χωνευτὰ αὐτῶν ἀπολεῖτε αὐτὰ, καὶ πάσας τὰς στήλας αὐτῶν ἐξαρεῖτε.
\vs{53}Καὶ ἀπολεῖτε πάντας τοὺς κατοικοῦντας τὴν γῆν, καὶ κατοικήσετε ἐν αὐτῇ, ὑμῖν γὰρ δέδωκα τὴν γῆν αὐτῶν ἐν κλήρῳ.
\vs{54}Καὶ κατακληρονομήσετε τὴν γῆν αὐτῶν ἐν κλήρῳ κατὰ φυλὰς ὑμῶν· τοῖς πλείοσι πληθυνεῖτε τὴν κατάσχεσιν αὐτῶν, καὶ τοῖς ἐλάττοσιν ἐλαττώσετε τὴν κατάσχεσιν αὐτῶν· εἰς ὃ ἂν ἐξέλθῃ τὸ ὄνομα αὐτοῦ, ἐκεῖ αὐτοῦ ἔσται· κατὰ φυλὰς πατριῶν ὑμῶν κληρονομήσετε.
\vs{55}Ἐὰν δὲ μὴ ἀπολέσητε τοὺς κατοικοῦντας ἐπὶ τῆς γῆς ἀπὸ προσώπου ὑμῶν, καὶ ἔσται οὓς ἐὰν καταλίπητε ἐξ αὐτῶν, σκόλοπες ἐν τοῖς ὀφθαλμοῖς ὑμῶν, καὶ βολίδες ἐν ταῖς πλευραῖς ὑμῶν, καὶ ἐχθρεύσουσιν ὑμῖν ἐπὶ τῆς γῆς, ἐφʼ ἣν ὑμεῖς κατοικήσετε.
\vs{56}Καὶ ἔσται καθότι διεγνώκειν ποιῆσαι αὐτοὺς, ποιήσω ὑμᾶς.

\ch{34}
Καὶ ἐλάλησε Κύριος πρὸς Μωυσῆν, λέγων,
\vs{2}ἔντειλαι τοῖς υἱοῖς Ἰσραὴλ, καὶ ἐρεῖς πρὸς αὐτοὺς, ὑμεῖς εἰσπορεύεσθε εἰς τὴν γῆν Χαναὰν· αὕτη ἔσται ὑμῖν εἰς κληρονομίαν, γῆ Χαναὰν σὺν τοῖς ὁρίοις αὐτῆς.
\vs{3}Καὶ ἔσται ὑμῖν τὸ κλίτος τὸ πρὸς Λίβα ἀπὸ ἐρήμου Σὶν ἕως ἐχόμενον Ἐδὼμ, καὶ ἔσται ὑμῖν τὰ ὅρια πρὸς λίβα ἀπὸ μέρους τῆς θαλάσσης τῆς ἁλυκῆς ἀπὸ ἀνατολῶν.
\vs{4}Καὶ κυκλώσει ὑμᾶς τὰ ὅρια ἀπὸ Λιβὸς πρὸς ἀνάβασιν Ἀκραβὶν, καὶ παρελεύσεται Ἐννὰκ, καὶ ἔσται ἡ διέξοδος αὐτοῦ πρὸς Λίβα Κάδης τοῦ Βαρνὴ, καὶ ἐξελεύσεται εἰς ἔπαυλιν Ἀρὰδ, καὶ παρελεύσεται Ἀσεμωνᾶ.
\vs{5}Καὶ κυκλώσει τὰ ὅρια ἀπὸ Ἀσεμωνᾶ χειμάῤῥουν Αἰγύπτου, καὶ ἔσται ἡ διέξοδος ἡ θάλασσα.
\vs{6}Καὶ τὰ ὅρια τῆς θαλάσσης ἔσται ὑμῖν, ἡ θάλασσα ἡ μεγάλη ὁριεῖ, τοῦτο ἔσται ὑμῖν τὰ ὅρια τῆς θαλάσσης.

\vs{7}Καὶ τοῦτο ἔσται ὑμῖν τὰ ὅρια πρὸς βοῤῥᾶν· ἀπὸ τῆς θαλάσσης τῆς μεγάλης καταμετρήσετε ὑμῖν αὐτοῖς παρὰ τὸ ὄρος τὸ ὄρος.
\vs{8}Καὶ ἀπὸ τοῦ ὄρους τὸ ὄρος καταμετρήσετε αὐτοῖς, εἰσπορευομένων εἰς Ἐμὰθ, καὶ ἔσται ἡ διέξοδος αὐτοῦ τὰ ὅρια Σαραδάκ.
\vs{9}Καὶ ἐξελεύσεται τὰ ὅρια Δεφρωνὰ, καὶ ἔσται ἡ διέξοδος αὐτοῦ Ἀρσεναΐν· τοῦτο ἔσται ὑμῖν ὅρια ἀπὸ Βοῤῥᾶ.
\vs{10}Καὶ καταβήσεται τὰ ὅρια ἀπὸ Σεπεφαμὰρ Βηλὰ ἀπὸ ἀνατολῶν ἐπὶ πηγὰς, και καταβήσεται τὰ ὅρια Βηλὰ ἐπὶ νώτου θαλάσσης Χενερὲθ ἀπὸ ἀνατολῶν.
\vs{11}Καὶ καταβήσεται τὰ ὅρια ἀπὸ Σεπφαμὰρ Βηλὰ ἀπὸ ἀνατολῶν ἐπὶ πηγάς, καὶ καταβήσεται τὰ ὅρια Βηλὰ ἐπὶ νώτου θαλάσσης Χενερὲθ ἀπὸ ἀνατολῶν.
\vs{12}Καὶ καταβήσεται τὰ ὅρια ἐπὶ τὸν Ἰορδάνην, καὶ ἔσται ἡ διέξοδος θάλασσα ἡ ἁλυκή· αὕτη ἔσται ὑμῖν ἡ γῆ καὶ τὰ ὅρια αὐτῆς κύκλῳ.

\vs{13}Καὶ ἐνετείλατο Μωυσῆς τοῖς υἱοῖς Ἰσραὴλ, λέγων, αὕτη ἡ γῆ ἣν κατακληρονομήσετε αὐτὴν μετὰ κλήρου, ὃν τρόπον συνέταξε Κύριος δοῦναι αὐτὴν ταῖς ἐννέα φυλαῖς καὶ τῷ ἡμίσει φυλῆς Μανασσῆ.
\vs{14}Ὅτι ἔλαβε φυλὴ υἱῶν Ῥουβὴν, καὶ φυλὴ υἱῶν Γὰδ κατʼ οἴκους πατριῶν αὐτῶν· καὶ τὸ ἥμισυ φυλῆς Μανασσῆ ἀπέλαβον τοὺς κλήρους αὐτῶν.
\vs{15}Δύο φυλαὶ καὶ ἥμισυ φυλῆς ἔλαβον τοὺς κλήρους αὐτῶν πέραν τοῦ Ἰορδάνου κατὰ Ἱεριχὼ ἀπὸ Νότου κατʼ ἀνατολάς.

\vs{16}Καὶ ἐλάλησε Κύριος πρὸς Μωυσῆν, λέγων,
\vs{17}ταῦτα τὰ ὀνόματα τῶν ἀνδρῶν, οἳ κληρονομήσουσιν ὑμῖν τὴν γῆν· Ἐλεάζαρ ὁ ἱερεὺς καὶ Ἰησοῦς ὁ τοῦ Ναυή.
\vs{18}Καὶ ἄρχοντα ἕνα ἐκ φυλῆς λήψεσθε κατακληρονομῆσαι ὑμῖν τὴν γῆν.

\vs{19}Καὶ ταῦτα τὰ ὀνόματα τῶν ἀνδρῶν· τῆς φυλῆς Ἰούδα, Χάλεβ υἱὸς Ἰεφοννή.
\vs{20}Τῆς φυλῆς Συμεὼν, Σαλαμιὴλ υἱὸς Σεμιούδ.
\vs{21}Τῆς φυλῆς Βενιαμὶν, Ἐλδὰδ υἱὸς Χασλών·
\vs{22}Τῆς φυλῆς Δὰν, ἄρχων Βακχὶρ υἱὸς Ἐγλί.
\vs{23}Τῶν υἱῶν Ἰωσὴφ φυλῆς υἱῶν Μανασσῆ, ἄρχων Ἀνιὴλ υἱὸς Σουφί.
\vs{24}Τῆς φυλῆς υἱῶν Ἐφραὶμ, ἄρχων Καμουὴλ υἱὸς Σαβαθᾶν.
\vs{25}Τῆς φυλῆς Ζαβουλὼν, ἄρχων Ἐλισαφὰν υἱὸς Φαρνάχ.
\vs{26}Τῆς φυλῆς υἱῶν Ἰσσάχαρ, ἄρχων Φαλτιὴλ υἱὸς Ὀζᾶ.
\vs{27}Τῆς φυλῆς υἱῶν Ἀσὴρ, ἄρχων Ἀχιὼρ υἱὸς Σελεμί.
\vs{28}Τῆς φυλῆς Νεφθαλὶ, ἄρχων Φαδαὴλ υἱὸς Ἰαμιούδ.

\vs{29}Τούτοις ἐνετείλατο Κύριος καταμερίσαι τοῖς υἱοῖς Ἰσραὴλ ἐν γῇ Χαναάν.

\ch{35}
Καὶ ἐλάλησε Κύριος πρὸς Μωυσῆν ἐπὶ δυσμῶν Μωὰβ παρὰ τὸν Ἰορδάνην κατὰ Ἱεριχὼ, λέγων,
\vs{2}σύνταξον τοῖς υἱοῖς Ἰσραὴλ, καὶ δώσουσι τοῖς Λευίταις ἀπὸ τῶν κλήρων κατασχέσεως αὐτῶν πόλεις κατοικεῖν· καὶ τὰ προάστεια τῶν πόλεων κύκλῳ αὐτῶν δώσουσι τοῖς Λευίταις.
\vs{3}Καὶ ἔσονται αὐτοῖς αἱ πόλεις κατοικεῖν, καὶ τὰ ἀφορίσματα αὐτῶν ἔσται τοῖς κτήνεσιν αὐτῶν, καὶ πᾶσι τοῖς τετράποσιν αὐτῶν.
\vs{4}Καὶ τὰ συνκυροῦντα τῶν πόλεων, ἃς δώσετε τοῖς Λευίταις, ἀπὸ τείχους τῆς πόλεως καὶ ἔξω δισχιλίους πήχεις κύκλῳ.
\vs{5}Καὶ μετρήσεις ἔξω τῆς πόλεως τὸ κλίτος τὸ πρὸς ἀνατολὰς δισχιλίους πήχεις, καὶ τὸ κλίτος τὸ πρὸς Λίβα δισχιλίους πήχεις, καὶ τὸ κλίτος τὸ πρὸς θάλασσαν διαχιλίους πήχεις, καὶ τὸ κλίτος τὸ πρὸς Βοῤῥᾶν δισχιλίους πήχεις· καὶ ἡ πόλις μέσον τούτου ἔσται ὑμῖν, καὶ τὰ ὅμορα τῶν πόλεων.
\vs{6}Καὶ τὰς πόλεις δώσετε τοῖς Λευίταις, τὰς ἓξ πόλεις τῶν φυγαδευτηρίων ἃς δώσετε φυγεῖν ἐκεῖ τῷ φονεύσαντι, καὶ πρὸς ταύταις τεσσαράκοντα καὶ δύο πόλεις.
\vs{7}Πάσας τὰς πόλεις δώσετε τοῖς Λευίταις τεσσαράκοντα καὶ ὀκτὼ πόλεις· ταύτας, καὶ τὰ προάστεια αὐτῶν.
\vs{8}Καὶ τὰς πόλεις ἃς δώσετε ἀπὸ τῆς κατασχέσεως υἱῶν Ἰσραὴλ, ἀπὸ τῶν τὰ πολλὰ, πολλά· καὶ ἀπὸ τῶν ἐλαττόνων, ἐλάττω· ἕκαστος κατὰ τὴν κληρονομίαν αὐτοῦ ἣν κατακληρονομήσουσι, δώσουσιν ἀπὸ τῶν πόλεων τοῖς Λευίταις.

\vs{9}Καὶ ἐλάλησε Κύριος πρὸς Μωυσῆν, λέγων,
\vs{10}λάλησον τοῖς υἱοῖς Ἰσραὴλ, καὶ ἐρεῖς πρὸς αὐτοὺς, ὑμεῖς διαβαίνετε τὸν Ἰορδάνην εἰς γῆν Χαναάν.
\vs{11}Καὶ διαστελεῖτε ὑμῖν αὐτοῖς πόλεις· φυγαδευτήρια ἔσται ὑμῖν φυγεῖν ἐκεῖ τὸν φονευτὴν, πᾶς ὁ πατάξας ψυχὴν ἀκουσίως.
\vs{12}Καὶ ἔσονται αἱ πόλεις ὑμῖν φυγαδευτήρια ἀπὸ τοῦ ἀγχιστεύοντος τὸ αἷμα, καὶ οὐ μὴ ἀποθάνῃ ὁ φονεύων ἕως ἂν στῇ ἔναντι τῆς συναγωγῆς εἰς κρίσιν.
\vs{13}Καὶ αἱ πόλεις ἃς δώσετε τὰς ἓξ πόλεις, φυγαδευτήρια ἔσονται ὑμῖν.
\vs{14}Τὰς τρεῖς πόλεις δώσετε πέραν τοῦ Ἰορδάνου, καὶ τὰς τρεῖς πόλεις δώσετε ἐν γῇ Χαναάν.

\vs{15}Φυγαδεῖον ἔσται τοῖς υἱοῖς Ἰσραήλ, καὶ τῷ προσηλύτῳ, καὶ τῷ παροίκῳ τῷ ἐν ὑμῖν· ἔσονται αἱ πόλεις αὗται εἰς φυγαδευτήριον, φυγεῖν ἐκεῖ παντὶ πατάξαντι ψυχὴν ἀκουσίως.

\vs{16}Ἐὰν δὲ ἐν σκεύει σιδήρου πατάξῃ αὐτὸν, καὶ τελευτήσῃ, φονευτής ἐστι· θανάτῳ θανατούσθω ὁ φονευτής.
\vs{17}Ἐὰν δὲ ἐν λίθῳ ἐκ χειρὸς ἐν ᾧ ἀποθανεῖται ἐν αὐτῷ, πατάξῃ αὐτὸν, καὶ ἀποθάνῃ, φονευτής ἐστι· θανάτῳ θανατούσθω ὁ φονευτής.
\vs{18}Ἐὰν δὲ ἐν σκεύει ξυλίνῳ ἐκ χειρὸς ἐξ οὗ ἀποθανεῖται ἐν αὐτῷ, πατάξῃ αὐτόν, καὶ ἀποθάνῃ, φονευτής ἐστι· θανάτῳ θανατούσθω ὁ φονευτής.

\vs{19}Ὁ ἀγχιστεύων τὸ αἷμα, οὗτος ἀποκτενεῖ τὸν φονεύσαντα· ὅταν συναντήσῃ αὐτῷ οὗτος, ἀποκτενεῖ αὐτόν.
\vs{20}Ἐὰν δὲ διʼ ἔχθραν ὤσῃ αὐτὸν, καὶ ἐπιῤῥίψῃ ἐπʼ αὐτὸν πᾶν σκεῦος ἐξ ἐνέδρου, καὶ ἀποθάνῃ,
\vs{21}ἢ διὰ μῆνιν ἐπάταξεν αὐτὸν τῇ χειρί, καὶ ἀποθάνῃ, θανάτῳ θανατούσθω ὁ πατάξας, φονευτής ἐστι· θανάτῳ θανατούσθω ὁ φονεύων· ὁ ἀγχιστεύων τὸ αἷμα ἀποκτενεῖ τὸν φονεύσαντα ἐν τῷ συναντῆσαι αὐτῷ.

\vs{22}Ἐὰν δὲ ἐξάπινα, οὐ διʼ ἔχθραν ὤσῃ αὐτὸν, ἢ ἐπιῤῥίψῃ ἐπʼ αὐτὸν πᾶν σκεῦος, οὐκ ἐξ ἐνέδρου,
\vs{23}ἢ παντὶ λίθῳ, ἐν ᾧ ἀποθανεῖται ἐν αὐτῷ, οὐκ εἰδὼς, καὶ ἐπιπέσῃ ἐπʼ αὐτὸν, καὶ ἀποθάνῃ, αὐτὸς δὲ οὐκ ἐχθρὸς αὐτοῦ ἦν, οὐδὲ ζητῶν κακοποιῆσαι αὐτόν·
\vs{24}καὶ κρινεῖ ἡ συναγωγὴ ἀναμέσον τοῦ πατάξαντος καὶ ἀναμέσον τοῦ ἀγχιστεύοντος τὸ αἷμα, κατὰ τὰ κρίματα ταῦτα.
\vs{25}Καὶ ἐξελεῖται ἡ συναγωγὴ τὸν φονεύσαντα ἀπὸ τοῦ ἀγχιστεύοντος τὸ αἷμα, καὶ ἀποκαταστήσουσιν αὐτὸν ἡ συναγωγὴ εἰς τὴν πόλιν τοῦ φυγαδευτηρίου αὐτοῦ, οὗ κατέφυγε, καὶ κατοικήσει ἐκεῖ ἕως ἂν ἀποθάνῃ ὁ ἱερεὺς ὁ μέγας, ὃν ἔχρισαν αὐτὸν τῷ ἐλαίῳ τῷ ἁγίῳ.

\vs{26}Ἐὰν δὲ ἐξόδῳ ἐξέλθῃ ὁ φονεύασς τὰ ὅρια τῆς πόλεως εἰς ἣν κατέφυγεν ἐκεῖ,
\vs{27}καὶ εὕρῃ αὐτὸν ὁ ἀγχιστεύων τὸ αἷμα ἔξω τῶν ὁρίων τῆς πόλεως καταφυγῆς αὐτοῦ, καὶ φονεύσῃ ὁ ἀγχιστεύων τὸ αἷμα τὸν φονεύσαντα, οὐκ ἔνοχός ἐστιν.
\vs{28}Ἐν γὰρ τῇ πόλει τῆς καταφυγῆς κατοικείτω ἕως ἂν ἀποθάνῃ ὁ ἱερεὺς ὁ μέγας· καὶ μετὰ τὸ ἀποθανεῖν τὸν ἱερέα τὸν μέγαν, ἐπαναστραφήσεται ὁ φονεύσας εἰς τὴν γῆν τῆς κατασχέσεως αὐτοῦ.

\vs{29}Καὶ ἔσται ταῦτα ὑμῖν εἰς δικαίωμα κρίματος εἰς τὰς γενεὰς ὑμῶν ἐν πάσαις ταῖς κατοικίαις ὑμῶν.
\vs{30}Πᾶς πατάξας ψυχήν, διὰ μαρτύρων φονεύσεις τὸν φονεύσαντα· καὶ μάρτυς εἷς οὐ μαρτυρήσει ἐπὶ ψυχὴν ἀποθανεῖν.
\vs{31}Καὶ οὐ λήψεσθε λύτρα περὶ ψυχῆς παρὰ τοῦ φονεύσαντος τοῦ ἐνόχου ὄντος ἀναιρεθῆναι· θανάτῳ γὰρ θανατωθήσεται.
\vs{32}Οὐ λήψεσθε λύτρα τοῦ φυγεῖν εἰς πόλιν τῶν φυγαδευτηρίων, τοῦ πάλιν κατοικεῖν ἐπὶ τῆς γῆς, ἕως ἂν ἀποθάνῃ ὁ ἱερεὺς ὁ μέγας.
\vs{33}Καὶ οὐ μὴ φονοκτονήσητε τὴν γῆν εἰς ἣν ὑμεῖς κατοικεῖτε· τὸ γὰρ αἷμα τοῦτο φονοκτονεῖ τὴν γῆν, καὶ οὐκ ἐξιλασθήσεται ἡ γῆ ἀπὸ τοῦ αἵματος τοῦ ἐκχυθέντος ἐπʼ αὐτῆς, ἀλλʼ ἐπὶ τοῦ αἵματος τοῦ ἐκχέοντος.
\vs{34}Καὶ οὐ μιανεῖτε τὴν γῆν ἐφʼ ἧς κατοικεῖτε ἐπʼ αὐτῆς, ἐφʼ ἧς ἐγὼ κατασκηνώ ἐν ὑμῖν· ἐγὼ γάρ εἰμι Κύριος κατασκηνῶν ἐν μέσῳ τῶν υἱῶν Ἰσραήλ.

\ch{36}
Καὶ προσῆλθον οἱ ἄρχοντες φυλῆς υἱῶν Γαλαὰδ υἱοῦ Μαχὶρ υἱοῦ Μανασσῆ ἐκ τῆς φυλῆς υἱῶν Ἰωσὴφ, καὶ ἐλάλησαν ἔναντι Μωυσῆ, καὶ ἔναντι Ἐλεάζαρ τοῦ ἱερέως, καὶ ἔναντι τῶν ἀρχόντων οἴκων πατριῶν τῶν υἱῶν Ἰσραὴλ,
\vs{2}καὶ εἶπαν, τῷ κυρῖῳ ἡμῶν ἐνετείλατο Κύριος ἀποδοῦναι τὴν γῆν τῆς κληρονομίας ἐν κλήρῳ τοῖς υἱοῖς Ἰσραήλ· καὶ τῷ κυρίῳ συνέταξε Κύριος δοῦναι τὴν κληρονομίαν Σαλπαὰδ τοῦ ἀδελφοῦ ἡμῶν ταῖς θυγατράσιν αὐτοῦ.
\vs{3}Καὶ ἔσονται ἑνὶ τῶν φυλῶν υἱῶν Ἰσραὴλ γυναῖκες· καὶ ἀφαιρεθήσεται ὁ κλῆρος αὐτῶν ἐκ τῆς κατασχέσεως τῶν πατέρεν ἡμῶν, καὶ προστεθήσεται εἰς κληρονομίαν τῆς φυλῆς, οἷς ἂν γένωνται γυναῖκες, καὶ ἐκ τοῦ κλήρου τῆς κληρονομίας ἡμῶν ἀφαιρεθήσεται.
\vs{4}Ἐὰν δὲ γένηται ἡ ἄφεσις τῶν υἱῶν Ἰσραὴλ, καὶ προστεθήσεται ἡ κληρονομία αὐτῶν ἐπὶ τὴν κληρονομίαν τῆς φυλῆς, οἷς ἂν γένωνται γυναῖκες, καὶ ἀπὸ τῆς κληρονομίας φυλῆς πατριᾶς ἡμῶν ἀφαιρεθήσεται ἡ κληρονομία αὐτῶν.

\vs{5}Καὶ ἐνετείλατο Μωυσῆς τοῖς υἱοῖς Ἰσραὴλ διὰ προστάγματος Κυρίου, λέγων, οὕτως φυλὴ υἱῶν Ἰωσὴφ λέγουσι.
\vs{6}Τοῦτο τὸ ῥῆμα ὃ συνέταξε Κύριος ταῖς θυγατράσι Σαλπαὰδ, λέγων, οὗ ἀρέσκῃ ἐναντίον αὐτῶν, ἔστωσαν γυναῖκες, πλὴν ἐκ τοῦ δήμου τοῦ πατρὸς αὐτῶν ἔστωσαν γυναῖκες.
\vs{7}Καὶ οὐχὶ περιστραφήσεται κληρονομία τοῖς υἱοῖς Ἰσραὴλ ἀπὸ φυλῆς ἐπὶ φύλην, ὅτι ἕκαστος ἐν τῇ κληρονομίᾳ τῆς φυλῆς τῆς πατριᾶς αὐτοῦ προσκολληθήσονται οἱ υἱοὶ Ἰσραήλ.
\vs{8}Καὶ πᾶσα θυγάτηρ ἀγχιατεύουσα κληρονομίαν ἐκ τῶν φυλῶν υἱῶν Ἰσραὴλ, ἑνὶ τῶν ἐκ τοῦ δήμου τοῦ πατρὸς αὐτῆς ἔσονται γυναῖκες, ἵνα ἀγχιστεύσωσιν οἱ υἱοὶ Ἰσραὴλ ἕκαστος τὴν κληρονομίαν τὴν πατρικὴν αὐτοῦ.
\vs{9}Καὶ οὐ περιστραφήσεται ὁ κλῆρος ἐκ φυλῆς ἐπὶ φυλὴν ἑτέραν, ἀλλʼ ἕκαστος ἐν τῇ κληρονομίᾳ αὐτοῦ προσκολληθήσονται οἱ υἱοὶ Ἰσραήλ.

\vs{10}Ὃν τρόπον συνέταξε Κύριος Μωυσῇ, οὕτως ἐποίησαν θυγατράσι Σαλπαάδ.
\vs{11}Καὶ ἐγένοντο Θερσὰ καὶ Ἐγλὰ καὶ Μελχὰ καὶ Νούα καὶ Μαλαὰ θυγατέρες Σαλπαὰδ, τοῖς ἀνεψιοῖς αὐτῶν,
\vs{12}ἐκ τοῦ δήμου τοῦ Μανασσῆ υἱῶν Ἰωσὴφ ἐγενήθησαν γυναῖκες· καὶ ἐγενήθη ἡ κληρονομία αὐτῶν ἐπὶ τὴν φυλὴν δήμου τοῦ πατρὸς αὐτῶν.
\vs{13}Αὗται αἱ ἐντολαὶ καὶ τὰ δικαιώματα καὶ τὰ κρίματα, ἃ ἐνετείλατο Κύριος ἐν χειρὶ Μωυσῆ ἐπὶ δυσμῶν Μωὰβ ἐπὶ τοῦ Ἰορδάνου κατὰ Ἰεριχώ.


\def\book{ΔΕΥΤΕΡΟΝΟΜΙΟΝ}
\biblebook{ΔΕΥΤΕΡΟΝΟΜΙΟΝ}


\lettrine[lines=2, loversize=0.2, nindent=0em, findent=.25em]{\textcolor{bookheadingcolor}{Ο}}{ΥΤΟΙ} οἱ λόγοι οὓς ἐλάλησε Μωυσῆς παντὶ Ἰσραὴλ πέραν τοῦ Ἰορδάνου ἐν τῇ ἐρήμῳ πρὸς δυσμαῖς πλησίον τῆς ἐρυθρᾶς θαλάσσης ἀναμέσον Φαρὰν Τοφὸλ, καὶ Λοβὸν, καὶ Αὐλῶν, καὶ καταχρύσεα.
\vs{2}Ἕνδεκα ἡμερῶν ἐκ Χωρὴβ ὁδὸς ἐπʼ ὄρος
\vs{3}Σηεὶρ ἕως Κάδης Βαρνή. Καὶ ἐγενήθη ἐν τῷ τεσσαρακοστῷ ἔτει ἐν τῷ ἑνδεκάτῳ μηνὶ μιᾷ τοῦ μηνὸς, ἐλάλησε Μωυσῆς πρὸς πάντας υἱοὺς Ἰσραὴλ, κατὰ πάντα ὅσα ἐνετείλατο Κύριος αὐτῷ πρὸς αὐτούς· μετὰ τὸ πατάξαι
\vs{4}Σηὼν βασιλέα Ἀμοῤῥαίων τὸν κατοικήσαντα ἐν Ἐσεβὼν, καὶ τὸν Ὢγ βασιλέα τῆς Βασὰν τὸν κατοικήσαντα ἐν Ἀσταρὼθ καὶ ἐν
\vs{5}Ἐδραῒν, ἐν τῷ πέραν τοῦ Ἰορδάνου ἐν γῇ Μωὰβ, ἤρξατο Μωυσῆς διασαφῆσαι τὸν νόμον τοῦτον, λέγων,
\vs{6}Κύριος ὁ Θεὸς ἡμῶν ἐλάλησεν ἡμῖν ἐν Χωρὴβ, λέγων, ἰκανούσθω ὑμῖν κατοικεῖν ἐν τῷ ὄρει τούτῳ.
\vs{7}Ἐπιστράφητε καὶ ἀπάρατε ὑμεῖς καὶ εἰσπορεύεσθε εἰς ὄρος Ἀμοῤῥαίων, καὶ πρὸς πάντας τοὺς περιοίκους Ἄραβα, εἰς ὄρος καὶ πεδίον, καὶ πρὸς Λίβα, καὶ παραλίαν γῆν Χαναναίων, καὶ Ἀντιλίβανον ἕως τοῦ ποταμοῦ τοῦ μεγάλου, ποταμοῦ Εὐφράτου.
\vs{8}Ἴδετε, παραδέδωκεν ἐνώπιον ὑμῶν τὴν γῆν· εἰσπορευθέντες κληρονομήσατε τὴν γῆν, ἣν ὤμοσα τοῖς πατράσιν ὑμῶν τῷ Ἁβραὰμ, καὶ Ἰσαὰκ, καὶ Ἰσκὼβ, δοῦναι αὐτοῖς καὶ τῷ σπέρματι αὐτῶν μετʼ αὐτούς.

\vs{9}Καὶ εἶπα πρὸς ὑμᾶς ἐν τῷ καιρῷ ἐκείνῳ, λέγων, οὐ δυνήσομαι μόνος φέρειν ὑμᾶς.
\vs{10}Κύριος ὁ Θεὸς ὑμῶν ἐπλήθυνεν ὑμᾶς, καὶ ἰδού ἐστε σήμερον ὡσεὶ τὰ ἄστρα τοῦ οὐρανοῦ τῷ πλήθει.
\vs{11}Κύριος ὁ Θεὸς τῶν πατέρων ὑμῶν προσθείῃ ὑμῖν ὡς ἐστὲ χιλιοπλασίως, καὶ εὐλογήσαι ὑμᾶς καθότι ἐλάλησεν ὑμῖν.
\vs{12}Πῶς δυνήσομαι μόνος φέρειν τὸν κόπον ὑμῶν καὶ τὴν ὑπόστασιν ὑμῶν καὶ τὰς ἀντιλογίας ὑμῶν;
\vs{13}Δότε ἑαυτοῖς ἄνδρας σοφοὺς καὶ ἐπιστήμονας καὶ συνετοὺς εἰς τὰς φυλὰς ὑμῶν, καὶ καταστήσω ἐφʼ ὑμῶν, ἡγουμένους ὑμῶν.
\vs{14}Καὶ ἀπεκρίθητέ μοι, καὶ εἴπατε, καλὸν τὸ ῥῆμα ὃ ἐλάλησας ποιῆσαι.
\vs{15}Καὶ ἔλαβον ἐξ ὑμῶν ἄνδρας σοφοὺς καὶ ἐπιστήμονας καὶ συνετούς, καὶ κατέστησα αὐτοὺς ἡγεῖσθαι ἐφʼ ὑμῶν χιλιάρχους, καὶ ἑκατοντάρχους, καὶ πεντηκοντάρχους, καὶ δεκάρχους, καὶ γραμματοεισαγωγεῖς τοῖς κριταῖς ὑμῶν·
\vs{16}Καὶ ἐνετειλάμην τοῖς κριταῖς ὑμῶν ἐν τῷ καιρῷ ἐκείνῳ, λέγων, διακούετε ἀναμέσον τῶν ἀδελφῶν ὑμῶν, καὶ κρίνατε δικαίως ἀναμέσον ἀνδρὸς, καὶ ἀναμέσον ἀδελφοῦ, καὶ ἀναμέσον προσηλύτου αὐτοῦ.
\vs{17}Οὐκ ἐπιγνώσῃ πρόσωπον ἐν κρίσει, κατὰ τὸν μικρὸν καὶ κατὰ τὸν μέγαν κρινεῖς, οὐ μὴ ὑποστείλῃ πρόσωπον ἀνθρώπου, ὅτι ἡ κρίσις τοῦ Θεοῦ ἐστι· καὶ τὸ ῥῆμα ὃ ἐὰν σκληρὸν ἠ· ἀφʼ ὑμῶν, ἀνοίσετε αὐτὸ ἐπʼ ἐμὲ, καὶ ἀκούσομαι αὐτό.
\vs{18}Καὶ ἐνετειλάμην ὑμῖν ἐν τῷ καιρῷ ἐκείνῳ πάντας τοὺς λόγους, οὓς ποιήσετε.

\vs{19}Καὶ ἀπάραντες ἐκ Χωρὴβ ἐπορεύθημεν πᾶσαν τὴν ἔρημον τὴν μεγάλην καὶ τὴν φοβερὰν ἐκείνην, ἣν εἴδετε, ὁδὸν ὄρους τοῦ Ἀμοῤῥαίου, καθότι ἐνετείλατο Κύριος ὁ Θεὸς ἡμῶν ἡμῖν, καὶ ἤλθομεν ἕως Κάδης Βαρνή.
\vs{20}Καὶ εἶπα πρὸς ὑμᾶς, ἤλθατε ἕως τοῦ ὄρους τοῦ Ἀμοῤῥαίου, ὃ Κύριος ὁ Θεὸς ἡμῶν δίδωσιν ὑμῖν.
\vs{21}Ἴδετε, παραδέδωκεν ὑμῖν Κύριος ὁ Θεὸς ὑμῶν πρὸ προσώπου ὑμῶν τὴν γῆν· ἀναβάντες κληρονομήσατε ὃν τρόπον εἶπε Κύριος ὁ Θεὸς τῶν πατέρων ὑμῶν ὑμῖν· μὴ φοβεῖσθε, μηδὲ δειλιάσητε.
\vs{22}Καὶ προσήλθατέ μοι πάντες, καὶ εἴπατε, Ἀποστείλωμεν ἄνδρας προτέρους ἡμῶν, καὶ ἐφοδευσάτωσαν ἡμῖν τὴν γῆν, καὶ ἀναγγειλάτωσαν ἡμῖν ἀπόκρισιν τὴν ὁδὸν διʼ ἧς ἀναβησόμεθα ἐν αὐτῇ, καὶ τὰς πόλεις εἰς ἃς εἰσπορευσόμεθα εἰς αὐτάς.
\vs{23}Καὶ ἤρεσεν ἐναντίον μου τὸ ῥῆμα· καὶ ἔλαβον ἐξ ὑμῶν δώδεκα ἄνδρας, ἄνδρα ἕνα κατά φυλήν.
\vs{24}Καὶ ἐπιστραφέντες ἀνέβησαν εἰς τὸ ὄρος, καὶ ἤλθοσαν ἕως φάραγγος βότρυος, καὶ κατεσκόπευσαν αὐτήν.
\vs{25}Καὶ ἐλάβοσαν ἐν ταῖς χερσὶν αὐτῶν ἀπὸ τοῦ καρποῦ τῆς γῆς, καὶ κατήνεγκαν πρὸς ὑμᾶς, καὶ ἔλεγον, Ἀγαθὴ ἡ γῆ, ἣν Κύριος ὁ Θεὸς ἡμῶν δίδωσιν ἡμῖν.

\vs{26}Καὶ οὐκ ἠθελήσατε ἀναβῆναι, ἀλλʼ ἠπειθήσατε τῷ ῥήματι Κυρίου τοῦ Θεοῦ ἡμῶν.
\vs{27}Καὶ διεγογγύζετε ἐν ταῖς σκηναῖς ὑμῶν, καὶ εἴπατε, διὰ τὸ μισεῖν Κύριον ἡμᾶς, ἐξήγαγεν ἡμᾶς ἐκ γῆς Αἰγύπτου παραδοῦναι ἡμᾶς εἰς χεῖρας Ἀμοῤῥαίων, ἐξολοθρεῦσαι ἡμᾶς.
\vs{28}Ποῦ ἡμεῖς ἀναβαίνομεν; οἱ δὲ ἀδελφοὶ ὑμῶν ἀπέστησαν τὴν καρδίαν ὑμῶν, λέγοντες, ἔθνος μέγα καὶ πολὺ καὶ δυνατώτερον ἡμῶν, καὶ πόλεις μεγάλαι καὶ τετειχισμέναι ἕως τοῦ οὐρανοῦ· ἀλλὰ καὶ υἱοὺς γιγάντων ἑωράκαμεν ἐκεῖ.
\vs{29}Καὶ εἶπα πρὸς ὑμᾶς, μὴ πτήξητε, μηδὲ φοβηθῆτε ἀπʼ αὐτῶν.
\vs{30}Κύριος ὁ Θεὸς ὑμῶν ὁ προπορευόμενος πρὸ προσώπου ὑμῶν, αὐτὸς συνεκπολεμήσει αὐτοὺς μεθʼ ὑμῶν κατὰ πάντα ὅσα ἐποίησεν ὑμῖν ἐν γῇ Αἰγύπτῳ, καὶ ἐν τῇ ἐρήμῳ ταύτῃ ἣν εἴδετε, ὁδὸν ὄρους τοῦ Ἀμοῤῥαίου· ὡς τροφοφορήσαι σε Κύριος ὁ
\vs{31}Θεός σου, ὡς εἴτις τροφοφορήσαι ἄνθρωπος τὸν υἱὸν αὐτοῦ, κατὰ πᾶσαν τὴν ὁδὸν εἰς ἣν ἐπορεύθητε ἕως ῆλθετε εἰς τὸν τόπον τοῦτον.

\vs{32}Καὶ ἐν τῷ λόγῳ τούτῳ οὐκ ἐνεπιστεύσατε Κυρίῳ τῷ Θεῷ ἡμῶν,
\vs{33}ὃς προπορεύεται πρότερος ὑμῶν ἐν τῇ ὁδῷ ἐκλέγεσθαι ὑμῖν τόπον, ὁδηγῶν ὑμᾶς ἐν πυρὶ νυκτὸς, δεικνύων ὑμῖν τὴν ὁδὸν καθʼ ἣν πορεύεσθε ἐπʼ αὐτῆς, καὶ ἐν νεφέλῃ ἡμέρας.

\vs{34}Καὶ ἤκουσε Κύριος τὴν φωνὴν τῶν λόγων ὑμῶν, καὶ παροξυνθεὶς ὤμοσε, λέγων,
\vs{35}εἰ ὄψεταί τις τῶν ἀνδρῶν τούτων τὴν γῆν ἀγαθὴν ταύτην, ἣν ὤμοσα τοῖς πατράσιν αὐτῶν, πλὴν
\vs{36}Χάλεβ υἱὸς Ἰεφοννὴ, οὗτος ὄψεται αὐτὴν, καὶ τούτῳ δώσω τὴν γῆν ἐφʼ ἣν ἐπέβη, καὶ τοῖς υἱοῖς αὐτοῦ, διὰ τὸ προσκεῖσθαι αὐτὸν τὰ πρὸς Κύριον.
\vs{37}Καὶ ἐμοὶ ἐθυμώθη Κύριος διʼ ὑμᾶς, λέγων, οὐδὲ σὺ οὐ μὴ εἰσέλθῃς ἐκεῖ.
\vs{38}Ἰησοῦς υἱὸς Ναυὴ ὁ παρεστηκώς σοι, οὗτος εἰσελεύσεται ἐκεῖ· αὐτὸν κατίσχυσον, ὅτι αὐτὸς κατακληρονομήσει αὐτὴν τῷ Ἰσραήλ.
\vs{39}Καὶ πᾶν παιδίον νέον ὅστις οὐκ οἶδε σήμερον ἀγαθὸν ἢ κακόν, οὗτοι εἰσελεύσονται ἐκεῖ, καὶ τούτοις δώσω αὐτήν, καὶ αὐτοὶ κληρονομήσουσιν αὐτήν.
\vs{40}Καὶ ὑμεῖς ἐπιστράφεντες ἐστρατοπεδεύσατε εἰς τὴν ἔρημον, ὁδὸν τὴν ἐπὶ τῆς ἐρυθρᾶς θαλάσσης.

\vs{41}Καὶ ἀπεκρίθητε, καὶ εἴπατε, ἡμάρτομεν ἔναντι Κυρίου τοῦ Θεοῦ ἡμῶν· ἡμεῖς ἀναβάντες πολεμήσομεν κατὰ πάντα ὅσα ἐνετείλατο Κύριος ὁ Θεὸς ἡμῶν ἡμῖν· καὶ ἀναλαβόντες ἕκαστος τὰ σκεύη τὰ πολεμικὰ αὐτοῦ, καὶ συναθροισθέντες ἀναβαίνετε εἰς τὸ ὄρος.
\vs{42}Καὶ εἶπε Κύριος πρὸς μὲ, εἶπον αὐτοῖς, οὐκ ἀναβήσεσθε οὐδὲ μὴ πολεμήσετε, οὐ γάρ εἰμι μεθʼ ὑμῶν, καὶ οὐ μὴ συντριβῆτε ἐνώπιον τῶν ἐχθρῶν ὑμῶν.
\vs{43}Καὶ ἐλάλησα ὑμῖν, καὶ οὐκ εἰσηκούσατέ μου· καὶ παρέβητε τὸ ῥῆμα Κυρίου· καὶ παραβιασάμενοι ἀνέβητε εἰς τὸ ὄρος.
\vs{44}Καὶ ἐξῆλθεν ὁ Ἀμοῤῥαῖος ὁ κατοικῶν ἐν τῷ ὄρει ἐκείνῳ εἰς συνάντησιν ὑμῖν, καὶ κατεδίωξεν ὑμᾶς ὡσεὶ ποιήσαισαν αἱ μέλισσαι, καὶ ἐτίτρωσκον ὑμᾶς ἀπὸ Σηεὶρ ἕως Ἑρμᾶ.
\vs{45}Καὶ καθίσαντες ἐκλαίετε ἐναντίον Κυρίου τοῦ Θεοῦ ἡμῶν, καὶ οὐκ εἰσήκουσε Κύριος τῆς φωνῆς ὑμῶν, οὐδὲ προσέσχεν ὑμῖν.

\vs{46}Καὶ ἐνεκάθησθε ἐν Κάδης ἡμέρας πολλάς, ὅσας ποτὲ ἡμέρας ἐνεκάθησθε.

\ch{2}
Καὶ ἐπιστραφέντες ἀπῄραμεν εἰς τὴν ἔρημον, ὁδὸν θάλασσαν ἐρυθράν, ὃν τρόπον ἐλάλησε Κύριος πρὸς μὲ, καὶ ἐκυκλώσαμεν τὸ ὄρος τὸ Σηεὶρ ἡμέρας πολλάς.
\vs{2}Καὶ εἶπε Κύριος πρὸς μέ,
\vs{3}ἱκανούσθω ὑμῖν κυκλούν τὸ ὄρος τοῦτο· ἐπιστράφητε οὖν ἐπὶ Βοῤῥᾶν·
\vs{4}Καὶ τῷ λαῷ ἔντειλαι, λέγων, ὑμεῖς παραπορεύεσθε διὰ τῶν ὁρίων τῶν ἀδελφῶν ὑμῶν υἱῶν Ἡσαύ, οἳ κατοικοῦσιν ἐν Σηεὶρ, καὶ φοβηθήσονται ὑμᾶς, καὶ εὐλαβηθήσονται ὑμᾶς σφόδρα.
\vs{5}Μὴ συνάψητε πρὸς αὐτοὺς πόλεμον, οὐ γὰρ δῶ ὑμῖν ἀπὸ τῆς γῆς αὐτῶν οὐδὲ βῆμα ποδός, ὅτι ἐν κλήρῳ δέδωκα τοῖς υἱοῖς Ἡσαὺ τὸ ὄρος τὸ Σηείρ.
\vs{6}Ἀργυρίου βρώματα ἀγοράσατε παρʼ αὐτῶν καὶ φάγεσθε, καὶ ὕδωρ μέτρῳ λήψεσθε παρʼ αὐτῶν ἀργυρίου καὶ πίεσθε.
\vs{7}Ὁ γὰρ Κύριος ὁ Θεὸς ἡμῶν εὐλόγησέ σε ἐν παντὶ ἔργῳ τῶν χειρῶν σου· διάγνωθι πῶς διῆλθες τὴν ἔρημον τὴν μεγάλην καὶ τὴν φοβερὰν ἐκείνην· ἰδοὺ τεσσαράκοντα ἔτη Κύριος ὁ Θεός σου μετὰ σοῦ· οὐκ ἐπεδεήθης ῥήματος.

\vs{8}Καὶ παρήλθομεν τοὺς ἀδελφοὺς ἡμῶν υἱοὺς Ἡσαῦ, τοὺς κατοικοῦντας ἐν Σηεὶρ, παρὰ τὴν ὁδὸν τὴν Ἄραβα ἀπὸ Αἰλὼν καὶ ἀπὸ Γεσιὼν Γάβερ· καὶ ἐπιστρέψαντες παρήλθομεν ὁδὸν ἔρημον Μωάβ.
\vs{9}Καὶ εἶπε Κύριος πρὸς μέ, μὴ ἐχθραίνετε τοῖς Μωαβιταῖς, καὶ μὴ συνάψητε πρὸς αὐτοὺς πόλεμον· οὐ γὰρ μὴ δῶ ἀπὸ τῆς γῆς αὐτῶν ὑμῖν ἐν κλήρῳ, τοῖς γὰρ υἱοῖς Λὼτ δέδωκα τὴν Ἀροὴρ κληρονομεῖν.
\vs{10}Οἱ Ὀμμὶν πρότεροι ἐνεκάθηντο ἐπʼ αὐτῆς, ἔθνος μέγα καὶ πολὺ καὶ ἰσχύοντες, ὥσπερ οἱ Ἐνακίμ.
\vs{11}Ῥαφαῒν λογισθήσονται καὶ οὗτοι ὥσπερ καὶ οἱ Ἐνακίμ· καὶ οἱ Μωαβῖται ἐπονομάζουσιν αὐτοὺς Ὀμμείν.
\vs{12}Καὶ ἐν Σηεὶρ ἐνεκάθητο ὁ Χοῤῥαῖος τὸ πρότερον, καὶ υἱοὶ Ἡσαὺ ἀπώλεσαν αὐτοὺς, καὶ ἐξέτριψαν αὐτοὺς ἀπὸ προσώπου αὐτῶν. καὶ κατῳκίσθησαν ἀντʼ αὐτῶν, ὃν τρόπον ἐποίησεν Ἰσραὴλ τὴν γῆν τῆς κληρονομίας αὐτοῦ, ἣν δέδωκε Κύριος αὐτοῖς.
\vs{13}Νῦν οὖν ἀνάστητε καὶ ἀπάρατε ὑμεῖς, καὶ παραπορεύεσθε τὴν φάραγγα Ζάρετ.

\vs{14}Καὶ αἱ ἡμέραι ἃς παρεπορεύθημεν ἀπὸ Κάδης Βαρνὴ ἕως οὗ παρήλθομεν τὴν φάραγγα Ζαρὲτ, τριάκοντα καὶ ὀκτὼ ἔτη, ἕως οὗ διέπεσε πᾶσα γενεὰ ἀνδρῶν πολεμιστῶν ἀποθνήσκοντες ἐκ τῆς παρεμβολῆς, καθότι ὤμοσεν Κύριος ὁ
\vs{15}Θεός αὐτοὺς. Καὶ ἡ χεὶρ τοῦ Θεοῦ ἦν ἐπʼ αὐτοῖς ἐξαναλῶσαι αὐτοὺς ἐκ μέσου τῆς παρεμβολῆς ἕως οὗ διέπεσαν.

\vs{16}Καὶ ἐγενήθη ἐπειδὰν ἔπεσαν πάντες οἱ ἄνδρες οἱ πολεμισταὶ ἀποθνήσκοντες ἐκ μέσου τοῦ λαοῦ, καὶ ἐλάλησε
\vs{17}Κύριος πρὸς μὲ, λέγων, σὺ παραπορεύσῃ σήμερον τὰ ὅρια
\vs{18}Μωὰβ τὴν Ἀροὴρ,
\vs{19}καὶ προσάξετε ἐγγὺς υἱῶν Ἀμμάν· μὴ ἐχθραίνετε αὐτοῖς, μηδὲ συνάψετε αὐτοῖς εἰς πόλεμον· οὐ γὰρ μὴ δῶ ἀπὸ τῆς γῆς υἱῶν Ἀμμάν σοι ἐν κλήρῳ, ὅτι τοῖς υἱοῖς Λὼτ δέδωκα αὐτὴν ἐν κλήρῳ.
\vs{20}Γὴ Ῥαφαῒν λογισθήσεται, καὶ γὰρ ἐπʼ αὐτῆς κατῴκουν οἱ Ῥαφαῒν τοπρότερον καὶ οἱ Ἀμμανῖται ἐπονομάζουσιν αὐτοὺς Ζοχομμίν.
\vs{21}Ἔθνος μέγα καὶ πολὺ καὶ δυνατώτερον ὑμῶν, ὥσπερ καὶ οἱ Ἐνακείμ· καὶ ἀπώλεσεν αὐτοὺς Κύριος πρὸ πρόσωπου αὐτῶν, καὶ κατεκληρονόμησαν καὶ κατῳκίσθησαν, ἀντʼ αὐτῶν ἕως τῆς ἡμέρας ταύτης.
\vs{22}Ὥσπερ ἐποίησαν τοῖς υἱοῖς Ἡσαὺ τοῖς κατοικοῦσιν ἐν Σηείρ, ὃν τρόπον ἐξέτριψαν τὸν Χοῤῥαῖον ἀπὸ προσώπου αὐτῶν, καὶ κατεκληρονόμησαν αὐτοὺς, καὶ κατῳκίσθησαν ἀντʼ αὐτῶν ἕως τῆς ἡμέρας ταύτης.
\vs{23}Καὶ οἱ Εὐαῖοι οἱ κατοικοῦντες ἐν Ἀσηδὼθ ἕως Γάζης, καὶ οἱ Καππάδοκες οἱ ἐξελθόντες ἐκ Καππαδοκίας, ἐξέτριψαν αὐτοὺς, καὶ κατῳκίσθησαν ἀντʼ αὐτῶν.

\vs{24}Νῦν οὖν ἀνάστητε καὶ ἀπάρατε, καὶ παρέλθετε ὑμεῖς τὴν φάραγγα Ἀρνών· ἰδοὺ παραδέδωκα εἰς χεῖράς σου τὸν Σηὼν βασιλέα Ἐσεβὼν τὸν Ἀμοῤῥαῖον καὶ τὴν γῆν αὐτοῦ· ἐνάρχου κληρονομεῖν· σύναπτε πρὸς αὐτὸν πόλεμον ἐν τῇ ἡμέρᾳ ταύτῃ.
\vs{25}Ἐνάρχου δοῦναι τὸν τρόμον σου καὶ τὸν φόβον σου ἐπὶ προσώπου πάντων τῶν ἐθνῶν τῶν ὑποκάτω τοῦ οὐρανοῦ, οἵτινες ἀκούσαντες τὸ ὄνομά σου ταραχθήσονται, καὶ ὠδῖνας ἕξουσιν ἀπὸ προσώπου σου.

\vs{26}Καὶ ἀπέστειλα πρέσβεις ἐκ τῆς ἐρήμου Κεδαμὼθ πρὸς Σηὼν βασιλέα Ἐσεβὼν λόγοις εἰρηνικοῖς, λέγων,
\vs{27}παρελεύσομαι διὰ τῆς γῆς σου· ἐν τῇ ὁδῷ παρεύσομαι, οὐκ ἐκκλινῶ δεξιὰ οὐδʼ ἀριστερά.
\vs{28}Βρώματα ἀργυρίου ἀποδώσῃ μοι, καὶ φάγομαι· καὶ ὕδωρ ἀργυρίου ἀποδώσῃ μοι, καὶ πίομαι· πλὴν ὅτι παρελεύσομαι τοῖς ποσί·
\vs{29}Καθὼς ἐποίησάν μοι οἱ υἱοὶ Ἠσαῦ οἱ κατοικοῦντες ἐν Σηεὶρ, καὶ οἱ Μωαβῖται οἱ κατοικοῦντες ἐν Ἀροήρ· ἕως παρέλθω τὸν Ἰορδάνην εἰς τὴν γῆν, ἣν Κύριος ὁ Θεὸς ἡμῶν δίδωσιν ἡμῖν.
\vs{30}Καὶ οὐκ ἠθέλησεν Σηὼν βασιλεὺς Ἐσεβὼν παρελθεῖν ἡμᾶς διʼ αὐτοῦ, ὅτι ἐσκλήρυνε Κύριος ὁ Θεὸς ἡμῶν τὸ πνεῦμα αὐτοῦ, καὶ κατίσχυσε τὴν καρδίαν αὐτοῦ, ἵνα παραδοθῇ εἰς τὰς χεῖράς σου ὡς ἐν τῇ ἡμέρᾳ ταύτῃ.

\vs{31}Καὶ εἶπε Κύριος πρὸς μέ ἰδοὺ ἦργμαι παραδοῦναι πρὸ προσώπου σου τὸν Σηὼν βασιλέα Ἐσεβὼν τὸν Ἀμοῤῥαῖον, καὶ τὴν γῆν αὐτοῦ, καὶ ἔναρξαι κληρονομῆσαι τὴν γῆν αὐτοῦ.
\vs{32}Καὶ ἐξῆλθε Σηὼν βασιλεὺς Ἐσεβὼν εἰς συνάντησιν ἡμῖν, αὐτὸς καὶ πᾶς ὁ λαὸς αὐτοῦ, εἰς πόλεμον εἰς Ἰασσά.
\vs{33}Καὶ παρέδωκεν αὐτὸν Κύριος ὁ Θεὸς ἡμῶν πρὸ προσώπου ἡμῶν· καὶ ἐπατάξαμεν αὐτὸν καὶ τοὺς υἱοὺς αὐτοῦ καὶ πάντα τὸν λαὸν αὐτοῦ.
\vs{34}Καὶ ἐκρατήσαμεν πασῶν τῶν πόλεων αὐτοῦ ἐν τῷ καιρῷ ἐκείνῳ, καὶ ἐξωλοθρεύσαμεν πᾶσαν πόλιν ἑξῆς, καὶ τὰς γυναῖκας αὐτῶν καὶ τὰ τέκνα αὐτῶν· οὐ κατελίπομεν ζωγρίαν.
\vs{35}Πλὴν τὰ κτήνη ἐπρονομεύσαμεν, καὶ τὰ σκῦλα τῶν πόλεων ἐλάβομεν
\vs{36}ἐξ Ἀροὴρ, ἥ ἐστι παρὰ τὸ χεῖλος χειμάῤῥου Ἀρνών, καὶ τὴν πόλιν τὴν οὖσαν ἐν τῇ φάραγγι, καὶ ἕως ὄρους τοῦ Γαλαάδ· οὐκ ἐγενήθη πόλις ἥτις διέφυγεν ἡμᾶς. τὰς πάσας παρέδωκε Κύριος ὁ Θεὸς ἡμῶν εἰς τὰς χεῖρας ἡμῶν.
\vs{37}Πλὴν ἐγγὺς υἱῶν Ἀμμὰν οὐ προσήλθομεν πάντα τὰ συγκυροῦντα χειμάῤῥου Ἰαβὸκ, καὶ τὰς πόλεις τὰς ἐν τῇ ὀρεινῇ, καθότι ἐνετείλατο Κύριος ὁ Θεὸς ἡμῶν ἡμῖν.

\ch{3}
Καὶ ἐπιστραφέντες, ἀνέβημεν ὁδὸν τὴν εἰς Βασάν· καὶ ἐξῆλθεν Ὤγ βασιλεὺς τῆς Βασὰν εἰς συνάντησιν ἡμῖν, αὐτὸς καὶ πᾶς ὁ λαὸς αὐτου εἰς πόλεμον εἰς Ἑδραΐμ.
\vs{2}Καὶ εἶπε Κύριος πρὸς μέ, μὴ φοβηθῇς αὐτόν, ὅτι εἰς τὰς χεῖράς σου παραδέδωκα αὐτὸν, καὶ πάντα τὸν λαὸν αὐτοῦ, καὶ πᾶσαν τὴν γῆν αὐτοῦ· καὶ ποιήσεις αὐτῷ, ὥσπερ ἐποίησας Σηὼν βασιλεῖ τῶν Ἀμοῤῥαίων, ὃς κατῴκει ἐν Ἐσεβών.
\vs{3}Καὶ παρέδωκεν αὐτὸν Κύριος ὁ Θεὸς ἡμῶν εἰς τὰς χεῖρας ἡμῶν, καὶ τὸν Ὢγ βασιλέα τῆς Βασὰν, καὶ πάντα τὸν λαὸν αὐτοῦ· καὶ ἐπατάξαμεν αὐτὸν, ἕως τοῦ μὴ καταλιπεῖν αὐτοῦ σπέρμα.

\vs{4}Καὶ ἐκρατήσαμεν πασῶν τῶν πόλεων αὐτοῦ ἐν τῷ καιρῷ ἐκείνῷ· οὐκ ἦν πόλις, ἣν οὐκ ἐλάβομεν παρʼ αὐτῶν· ἑξήκοντα πόλεις, πάντα τὰ περίχωρα Ἀργὸβ βασιλέως Ὢγ ἐν Βασάν·
\vs{5}Πᾶσαι πόλεις ὀχυραί, τείχη ὑψηλά, πύλαι καὶ μοχλοί· πλὴν τῶν πόλεων τῶν Φερεζαίων τῶν πολλῶν σφόδρα.
\vs{6}Ἐξωλοθρεύσαμεν, ὥσπερ ἐποιήσαμεν τὸν Σηὼν βασιλέα Ἐσεβών, καὶ ἐξωλοθρεύσαμεν πᾶσαν πόλιν ἑξῆς, καὶ τὰς γυναῖκας, καὶ τὰ παιδία,
\vs{7}καὶ πάντα τὰ κτήνη· καὶ τὰ σκῦλα τῶν πόλεων ἐπρονομεύσαμεν ἑαυτοῖς.

\vs{8}Καὶ ἐλάβομεν ἐν τῷ καιρῷ ἐκείνῳ τὴν γῆν ἐκ χειρῶν δύο βασιλέων τῶν Ἀμοῤῥαίων, οἳ ἦσαν πέραν τοῦ Ἰορδάνου ἀπὸ τοῦ χειμάῤῥου Ἀρνὼν καὶ ἕως Ἀερμών·
\vs{9}Οἱ Φοίνικες ἐπονομάζουσιν τὸ Ἀερμὼν Σανιώρ, καὶ ὁ Ἀμοῤῥαῖος ἐπωνόμασεν αὐτὸ Σανίρ·
\vs{10}Πᾶσαι πόλεις Μισὼρ, καὶ πᾶσα Γαλαάδ, καὶ πᾶσα Βασὰν ἕως Ἑλχᾶ καὶ Ἑδραῒμ, πόλεις βασιλείας τοῦ Ὢγ ἐν τῇ Βασάν·
\vs{11}Ὅτι πλὴν Ὢγ βασιλεὺς Βασὰν κατελείφθη ἀπὸ τῶν Ῥαφαΐν· ἰδοὺ ἡ κλὶνη αὐτοῦ κλὶνη σιδηρᾶ, ἰδοὺ αὕτη ἐν τῇ ἄκρᾳ τῶν υἱῶν Ἀμμών· ἐννέα πήχεων τὸ μηκος αὐτῆς, καὶ τεσσάρων πήχεων τὸ εὖρος αὐτῆς ἐν πήχει ἀνδρός.
\vs{12}Καὶ τὴν γῆν ἐκείνην ἐκληρονομήσαμεν ἐν τῷ καιρῷ ἐκείνῳ ἀπὸ Ἀροήρ, ἥ ἐστι παρὰ τὸ χείλος χειμάῤρου Ἀρνών, καὶ τὸ ἥμισυ τοῦ ὄρους Γαλαάδ· καὶ τὰς πόλεις αὐτοῦ ἔδωκα τῷ Ῥουβὴν καὶ τῷ Γάδ.
\vs{13}Καὶ τὸ κατάλοιπον τοῦ Γαλαὰδ, καὶ πᾶσαν τὴν Βασάν βασιλείαν Ὤγ ἔδωκα τῷ ἡμίσει φυλῆς Μανασσή, καὶ πᾶσαν περίχωρον Ἀργόβ, πᾶσαν Βασὰν ἐκείνην· γῆ Ῥαφαῒν λογισθήσεται.
\vs{14}Καὶ Ἰαῒρ υἱὸς Μανασσὴ ἔλαβε πᾶσαν τὴν περίχωρον Ἀργὸβ ἕως τῶν ὁρίων Γαργασὶ καὶ Μαχαθί· ἐπωνόμασεν αὐτὰς ἐπὶ τῷ ὀνόματι αὐτοῦ τὴν Βασὰν Θαυὼθ Ἰαεῒρ ἕως τῆς ἡμέρας ταύτης.
\vs{15}Καὶ τῷ Μαχὶρ ἔδωκα τὴν Γαλαάδ.
\vs{16}Καὶ τῷ Ῥουβὴν καὶ τῷ Γάδ δέδωκα ὑπὸ τῆς Γαλαὰδ ἕως χειμάῤῥου Ἀρνών μέσον τοῦ χειμάῤῥου ὅριον καὶ ἕως τοῦ Ἰαβόκ· ὁ χειμάῤῥους ὅριον τοῖς υἱοῖς Ἀμμάν·
\vs{17}Καὶ ἡ Ἄραβα καὶ ὁ Ἰορδάνης ὅριον Μαχαναρὲθ, καὶ ἕως θαλάσσης Ἄραβά, θαλάσσης ἁλυκῆς ὑπὸ Ἀσηδὼθ τὴν Φασγὰ ἀνατολῶν.

\vs{18}Καὶ ἐνετειλάμην ὑμῖν ἐν τῷ καιρῷ ἐκείνῳ, λέγων, Κύριος ὁ Θεὸς ὑμῶν ἔδωκεν ὑμῖν τὴν γῆν ταύτην ἐν κλήρῳ· ἐνοπλισάμενοι προπορεύεσθε πρὸ προσώπου τῶν ἀδελφῶν ὑμῶν υἱῶν Ἰσραήλ πᾶς δυνατός.
\vs{19}Πλὴν αἱ γυναῖκες ὑμῶν καὶ τὰ τέκνα ὑμῶν καὶ τὰ κτήνη ὑμῶν, οἶδα ὅτι πολλὰ κτήνη ὑμῖν, κατοικείτωσαν ἐν ταῖς πόλεσιν ὑμῶν, αἷς ἔδωκα ὑμῖν, ἕως ἂν καταπαύσῃ Κύριος ὁ
\vs{20}Θεὸς ὑμῶν τοὺς ἀδελφοὺς ὑμῶν, ὥσπερ καὶ ὑμᾶς, καὶ κατακληρονομήσωσι καὶ οὗτοι τὴν γῆν, ἣν Κύριος ὁ Θεὸς ἡμῶν δίδωσιν αὐτοῖς ἐν τῷ πέραν τοῦ Ἰορδάνου· καὶ ἐπαναστραφήσεσθε ἕκαστος εἰς τὴν κληρονομίαν αὐτοῦ, ἣν ἔδωκα ὑμῖν.

\vs{21}Καὶ τῷ Ἰησοῖ ἐνετειλάμην ἐν τῷ καιρῷ ἐκείνῳ, λέγων, οἱ ὀφθαλμοὶ ὑμῶν ἑωράκασιν πάντα, ὅσα ἐποίησε Κύριος ὁ Θεὸς ἡμῶν τοῖς δυσὶ βασιλεῦσι τούτοις· οὕτως ποιήσει Κύριος ὁ Θεὸς ἡμῶν πάσας τὰς βασιλείας ἐφʼ ἃς σὺ διαβαίνεις ἐκεῖ.
\vs{22}Οὐ φοβηθήσεσθε ἀπʼ αὐτῶν, ὅτι Κύριος ὁ Θεὸς ἡμῶν αὐτὸς πολεμήσει περὶ ὑμῶν.

\vs{23}Καὶ ἐδεήθην Κυρίου ἐν τῷ καιρῷ ἐκείνῳ, λέγων,
\vs{24}Κύριε Θεὲ, σὺ ἤρξω δεῖξαι τῷ σῷ θεράποντι τὴν ἰσχύν σου, καὶ τὴν δύναμίν σου, καὶ τὴν χεῖρα τὴν κραταιὰν, καὶ τὸν βραχίονα τὸν ὑψηλόν· τίς γάρ ἐστι Θεὸς ἐν τῷ οὐρανῷ ἢ ἐπὶ τῆς γῆς, ὅστις ποιήσει καθὰ ἐποίησας σὺ, καὶ κατὰ τὴν ἰσχύν σου;
\vs{25}Διαβὰς οὖν ὄψομαι τὴν γῆν τὴν ἀγαθὴν ταύην τὴν οὖσαν πέραν τοῦ Ἰορδάνου, τὸ ὄρος τοῦτο τὸ ἀγαθὸν καὶ τὸν Ἀντιλίβανον.

\vs{26}Καὶ ὑπερεῖδε Κύριος ἐμὲ ἕνεκεν ὑμῶν, καὶ οὐκ εἰσήκουσέ μου· καὶ εἶπε Κύριος πρὸς μέ, ἱκανούσθω σοι, μὴ προσθῇς ἔτι λαλῆσαι τὸν λόγον τοῦτον.
\vs{27}Ἀνάβηθι ἐπὶ τὴν κορυφὴν τοῦ λελαξευμένου, καὶ ἀναβλέψας τοὶς ἀφθαλμοῖς σου κατὰ θάλασσαν καὶ Βοῤῥᾶν καὶ Λίβα καὶ ἀνατολάς, καὶ ἴδε τοῖς ὀφθαλμοῖς σου, ὅτι οὐ διαβήσῃ τὸν Ἰορδάνην τοῦτον.
\vs{28}Καὶ ἔντειλαι Ἰησοῖ καὶ κατίσχυσον αὐτὸν καὶ παρακάλεσον αὐτὸν, ὅτι οὗτος διαβήσεται πρὸ προσώπου τοῦ λαοῦ τούτου, καὶ οὗτος κατακληρονομήσει αὐτοῖς πᾶσαν τὴν γῆν ἣν ἑώρακας.
\vs{29}Καὶ ἐνεκαθήμεθα ἐν νάπῃ σύνεγγυς οἴκου Φογώρ.

\ch{4}
Καὶ νῦν Ἰσραὴλ ἄκουε τῶν δικαιωμάτων καὶ τῶν κριμάτων, ὅσα ἐγὼ διδάσκω ὑμᾶς σήμερον ποιεῖν, ἵνα ζῆτε, καὶ πολυπλασιασθῆτε, καὶ εἰσελθόντες κληρονομήσητε τὴν γῆν, ἣν Κύριος ὁ Θεὸς τῶν πατέρων ὑμῶν δίδωσιν ὑμῖν.
\vs{2}Οὐ προσθήσετε πρὸς τὸ ῥῆμα ὃ ἐγὼ ἐντέλλομαι ὑμῖν, καὶ οὐκ ἀφελεῖτε ἀπʼ αὐτοῦ· φυλάσσεσθε τὰς ἐντολὰς Κυρίου τοῦ Θεοῦ ἡμῶν, ὅσα ἐγὼ ἐντέλλομαι ὑμῖν σήμερον.
\vs{3}Οἱ ὀφθαλμοὶ ὑμῶν ἑωράκασι πάντα ὅσα ἐποίησε Κύριος ὁ Θεὸς ἡμῶν τῷ Βεελφεγὼρ, ὅτι πᾶς ἄνθρωπος ὅστις ἐπορεύθη ὀπίσω Βεελφεγὼρ, ἐξέτριψεν αὐτὸν Κύριος ὁ Θεὸς ὑμῶν ἐξ ὑμῶν.
\vs{4}Ὑμεῖς δὲ οἱ προσκείμενοι Κυρίῳ τῷ Θεῷ ὑμῶν, ζῆτε πάντες ἐν τῇ σήμερον.

\vs{5}Ἴδετε, δέδειχα ὑμῖν δικαιώματα καὶ κρίσεις καθὰ ἐνετείλατό μοι Κύριος, ποιῆσαι οὕτως ἐν τῇ γῇ εἰς ἣν ὑμεῖς εἰσπορεύεσθε ἐκεῖ κληρονομεῖν αὐτήν.
\vs{6}Καὶ φυλάξεσθε καὶ ποιήσετε· ὅτι αὕτη ἡ σοφία ὑμῶν καὶ ἡ σύνεσις ἐναντίον πάντων τῶν ἐθνῶν, ὅσοι ἂν ἀκούσωσι πάντα τὰ δικαιώματα ταῦτα· καὶ ἐροῦσιν, ἰδοὺ λαὸς σοφὸς καὶ ἐπιστήμων τὸ ἔθνος τὸ μέγα τοῦτο.
\vs{7}Ὅτι ποῖον ἔθνος μέγα, ᾧ ἐστιν αὐτῷ Θεὸς ἐγγίζων αὐτοῖς ὡς Κύριος ὁ Θεὸς ἡμῶν ἐν πᾶσιν οἷς ἐὰν αὐτὸν ἐπικαλεσώμεθα;
\vs{8}Καὶ ποῖον ἔθνος μέγα, ᾧ ἐστιν αὐτῷ δικαιώματα καὶ κρίματα δίκαια κατὰ πάντα τὸν νόμον τοῦτον, ὃν ἐγὼ δίδωμι ἐνώπιον ὑμῶν σήμερον;

\vs{9}Πρόσεχε σεαυτῷ, καὶ φύλαξον τὴν ψυχήν σου σφόδρα· μὴ ἐπιλάθῃ πάντας τοὺς λόγους, οὓς ἑωράκασιν οἱ ὀφθαλμοί σου, καὶ μὴ ἀποστήτωσαν ἀπὸ τῆς καρδίας σου πάσας τὰς ἡμέρας τῆς ζωῆς σου· καὶ συμβιβάσεις τοὺς υἱούς σου καὶ τοὺς υἱοὺς τῶν υἱῶν σου, ἡμέραν ἣν ἔστητε ἐνώπιον Κυρίου τοῦ Θεοῦ ἡμῶν ἐν Χωρὴβ τῇ ἡμέρᾳ τῆς ἐκκλησίας· ὅτε εἶπε
\vs{10}Κύριος πρὸς μὲ, ἐκκλησίασον πρὸς μὲ τὸν λαὸν, καὶ ἀκουσάτωσαν τὰ ῥήματά μου, ὅπως μάθωσι φοβεῖσθαί με πάσας τὰς ἡμέρας ἃς αὐτοὶ ζῶσιν ἐπὶ τῆς γῆς, καὶ τοὺς υἱοὺς αὐτῶν διδάξουσι.
\vs{11}Καὶ προσήλθετε καὶ ἔστητε ὑπὸ τὸ ὄρος· καὶ τὸ ὄρος ἐκαίετο πυρὶ ἕως τοῦ οὐρανοῦ· σκότος, γνόφος, θύελλα.
\vs{12}Καὶ ἐλάλησε Κύριος πρὸς ὑμᾶς ἐκ μέσου τοῦ πυρὸς φωνὴν ῥημάτων, ἣν ὑμεῖς ἠκούσατε· καὶ ὁμοίωμα οὐκ εἴδετε, ἀλλʼ ἢ φωνήν·
\vs{13}Καὶ ἀνήγγειλεν ὑμῖν τὴν διαθήκην αὐτοῦ, ἣν ἑνετείλατο ὑμῖν ποιεῖν, τὰ δέκα ῥήματα, καὶ ἔγραψεν αὐτὰ ἐπὶ δύο πλάκας λιθίνας.

\vs{14}Καὶ ἐμοὶ ἐνετείλατο Κύριος ἐν τῷ καιρῷ ἐκείνῳ, διδάξαι ὑμᾶς δικαιώματα καὶ κρίσεις, ποιεῖν ὑμᾶς αὐτὰ ἐπὶ τῆς γῆς, εἰς ἣν ὑμεῖς εἰσπορεύεσθε ἐκεῖ κληρονομῆσαι αὐτήν.
\vs{15}Καὶ φυλάξεσθε σφόδρα τὰς ψυχὰς ὑμῶν, ὅτι οὐκ εἴδετε ὁμοίωμα ἐν τῇ ἡμέρᾳ ᾗ ἐλάλησε Κύριος πρὸς ὑμᾶς ἐν Χωρὴβ ἐν τῷ ὄρει ἐκ μέσου τοῦ πυρός.
\vs{16}Μὴ ἀνομήσητε καὶ ποιήσητε ὑμῖν ἑαυτοῖς γλυπτὸν ὁμοίωμα, πᾶσαν εἰκόνα ὁμοίωμα ἀρσενικοῦ ἢ θηλυκοῦ,
\vs{17}ὁμοίωμα παντὸς κτήνους τῶν ὄντων ἐπὶ τῆς γῆς, ὁμοίωμα παντὸς ὀρνέου πτερωτοῦ ὃ πέταται ὑπὸ τὸν οὐρανὸν,
\vs{18}ὁμοίωμα παντὸς ἑρπετοῦ ὃ ἕρπει ἐπὶ τῆς γῆς, ὁμοίωμα παντὸς ἰχθύος, ὅσα ἐστὶν ἐν τοῖς ὕδασιν ὑποκάτω τῆς γῆς.
\vs{19}Καὶ μὴ ἀναβλέψας εἰς τὸν οὐρανὸν, καὶ ἰδὼν τὸν ἥλιον καὶ τὴν σελήνην καὶ τοὺς ἀστέρας, καὶ πάντα τὸν κόσμον τοῦ οὐρανοῦ, πλανηθεὶς προσκυνήσῃς αὐτοῖς, καὶ λατρεύσῃς αὐτοῖς, ἃ ἀπένειμε Κύριος ὁ Θεός σου αὐτὰ πᾶσι τοῖς ἔθνεσι τοῖς ὑποκάτω τοῦ οὐρανοῦ.
\vs{20}Ὑμᾶς δὲ ἔλαβεν ὁ Θεὸς, καὶ ἐξήγαγεν ὑμᾶς ἐκ γῆς Αἰγύπτου, ἐκ τῆς καμίνου τῆς σιδηρᾶς, ἐξ Αἰγύπτου, εἶναι αὐτῷ λαὸν ἔγκληρον, ὡς ἐν τῇ ἡμέρᾳ ταύτῃ.

\vs{21}Καὶ Κύριος ὁ Θεὸς ἐθυμώθη μοι περὶ τῶν λεγομένων ὑφʼ ὑμῶν, καὶ ὤμοσεν ἵνα μὴ διαβῶ τὸν Ἰορδάνην τοῦτον, καὶ ἵνα μὴ εἰσέλθω εἰς τὴν γῆν, ἣν Κύριος ὁ Θεός σου δίδωσί σοι ἐν κλήρῳ.
\vs{22}Ἐγὼ γὰρ ἀποθνήσκω ἐν τῇ γῇ ταύτῃ, καὶ οὐ διαβαίνω τὸν Ἰορδάνην τοῦτον· ὑμεῖς δὲ διαβαίνετε, καὶ κληρονομήσετε τὴν γῆν τὴν ἀγαθὴν ταύτην.
\vs{23}Προσέχετε ὑμῖν, μὴ ἐπιλάθησθε τὴν διαθήκην Κύριου τοῦ Θεοῦ ἡμῶν, ἣν διέθετο πρὸς ὑμᾶς, καὶ ἀνομήσητε, καὶ ποιήσητε ὑμῖν ἑαυτοῖς γλυπτὸν ὁμοίωμα πάντων ὧν συνέταξέ σοι Κύριος ὁ Θεός σου.
\vs{24}Ὅτι Κύριος ὁ Θεός σου πῦρ καταναλίσκον ἐστί, Θεὸς ζηλωτής.

\vs{25}Ἐὰν δὲ γεννήσῃς υἱοὺς καὶ υἱοὺς τῶν υἱῶν σου, καὶ χρονίσητε ἐπὶ τῆς γῆς, καὶ ἀνομήσητε, καὶ ποιήσετε γλυπτὸν ὁμοίωμα παντός, καὶ ποιήσητε τὸ πονηρὸν ἐνώπιου Κυρίου τοῦ
\vs{26}Θεοῦ ὑμῶν παροργίσαι αὐτόν, διαμαρτύρομαι ὑμῖν σήμερον τὸν τε οὐρανὸν καὶ τὴν γῆν, ὅτι ἀπωλίᾳ ἀπολεῖσθε ἀπὸ τῆς γῆς, εἰς ἣν ὑμεῖς διαβαίνετε τὸν Ἰορδάνην ἐκεῖ κληρονομῆσαι· οὐχὶ πολυχρονιεῖτε ἡμέρας ἐπʼ αὐτῆς, ἀλλʼ ἢ ἐκτριβῇ ἐκτριβήσεσθε.
\vs{27}Καὶ διασπερεῖ Κύριος ὑμᾶς ἐν πᾶσι τοῖς ἔθνεσι, καὶ καταλειφθήσεσθε ὀλίγοι ἀριθμῷ ἐν πᾶσι τοῖς ἔθνεσιν, εἰς οὓς εἰσάξει Κύριος ὑμᾶς ἐκεῖ.
\vs{28}Καὶ λατρεύσετε ἐκεῖ θεοῖς ἑτέροις ἔργοις χειρῶν ἀνθρώπων, ξύλοις καὶ λίθοις, οἳ οὐκ ὄψονται, οὔτε μὴ ἀκούσωσιν, οὔτε μὴ φάγωσιν, οὔτε μὴ ὀσφρανθῶσι.
\vs{29}Καὶ ζητήσετε ἐκεῖ Κύριον τὸν Θεὸν ὑμῶν, καὶ εὑρήσετε αὐτὸν ὅταν ἐκζητήσητε αὐτὸν ἐξ ὅλης τῆς καρδίας σου, καὶ ἐξ ὅλης τῆς ψυχῆς σου ἐν τῇ θλίψει σου·
\vs{30}Καὶ εὑρήσουσί σε πάντες οἱ λόγοι οὗτοι ἐπʼ ἐσχάτῳ τῶν ἡμερῶν, καὶ ἐπιστραφήσῃ πρὸς Κύριον τὸν Θεόν σου, καὶ εἰσακούσῃ τῆς φωνῆς σὐτοῦ·
\vs{31}Ὅτι Θεὸς οἰκτίρμων Κύριος ὁ Θεός σου· οὐκ ἐγκαταλείψει σε, οὐδὲ μὴ ἐκτρίψει σε· οὐκ ἐπιλήσεται τὴν διαθήκην τῶν πατέρων σου, ἣν ὤμοσεν αὐτοῖς Κύριος.

\vs{32}Ἐπερωτήσατε ἡμέρας προτέρας τὰς γενομένας προτέρας σου ἀπὸ τῆς ἡμέρας ἧς ἔκτισεν ὁ Θεὸς ἄνθρωπον ἐπὶ τῆς γῆς, καὶ ἐπὶ τὸ ἄκρον τοῦ οὐρανοῦ ἕως τοῦ ἄκρου τοῦ οὐρανοῦ, εἰ γέγονε κατὰ τὸ ῥῆμα τὸ μέγα τοῦτο, εἰ ἤκουσται τοιοῦτο· εἰ ἀκήκοεν ἔθνος φωνὴν
\vs{33}Θεοῦ ζῶντος λαλοῦντος ἐκ μέσου τοῦ πυρός, ὃν τρόπον ἀκήκοας σὺ καὶ ἔζησας·
\vs{34}εἰ ἐπείρασεν ὁ Θεὸς εἰσελθὼν λαβεῖν ἑαυτῷ ἔθνος ἐκ μέσου ἔθνους ἐν πειρασμῷ, καὶ ἐν σημείοις, καὶ ἐν τέρασι, καὶ ἐν πολέμῳ, καὶ ἐν χειρὶ κραταιᾷ, καὶ ἐν βραχίονι ὑψηλῷ, καὶ ἐν ὁράμασιν μεγάλοις, κατὰ πάντα ὅσα ἐποίησε Κύριος ὁ Θεὸς ἡμῶν ἐν Αἰγύπτῳ ἐνώπιόν σου βλέποντος·
\vs{35}ὥστε εἰδῆσαί σε ὅτι Κύριος ὁ Θεός σου οὗτος Θεός ἐστι, καὶ οὐκ ἔστιν ἔτι πλὴν αὐτοῦ.
\vs{36}Ἐκ τοῦ οὐρανοῦ ἀκουστὴ ἐγένετο ἡ φωνὴ αὐτοῦ παιδεῦσαί σε, καὶ ἐπὶ τῆς γῆς ἔδειξέ σοι τὸ πῦρ αὐτοῦ τὸ μέγα, καὶ τὰ ῥήματα αὐτοῦ ἤκουσας ἐκ μέσου τοῦ πυρός.

\vs{37}Διὰ τὸ ἀγαπῆσαι αὐτὸν τοὺς πατέρας σου, καὶ ἐξελέξατο τὸ σπέρμα αὐτῶν μετʼ αὐτοὺς ὑμᾶς, καὶ ἐξήγαγέ σε αὐτὸς ἐν τῇ ἰσχύϊ αὐτοῦ τῇ μεγάλῃ ἐξ Αἰγύπτου,
\vs{38}ἐξολοθρεῦσαι ἔθνη μεγάλα καὶ ἰσχυρότερά σου πρὸ προσώπου σου, εἰσαγαγεῖν σε δοῦναί σοι τὴν γῆν αὐτῶν κληρονομεῖν, καθὼς ἔχεις σήμερον.

\vs{39}Καὶ γνώσῃ σήμερον, καὶ ἐπιστραφήσῃ τῇ διανοίᾳ, ὅτι Κύριος ὁ Θεός σου οὗτος Θεὸς ἐν τῷ οὐρανῷ ἄνω καὶ ἐπὶ τῆς γῆς κάτω, καὶ οὐκ ἔστιν ἔτι πλὴν αὐτοῦ.
\vs{40}Καὶ φυλάξασθε τὰς ἐντολὰς αὐτοῦ, καὶ τὰ δικαιώματα αὐτοῦ, ὅσα ἐγὼ ἀντέλλομαί σοι σήμερον, ἵνα εὖ σοι γένηται καὶ τοῖς υἱοῖς σου μετὰ σὲ, ὅπως μακροήμεροι γένησθε ἐπὶ τῆς γῆς, ἧς Κύριος ὁ Θεός σου δίδωσί σοι πάσας τὰς ἡμέρας.
\vs{41}Τότε ἀφώρισε Μωυσῆς τρεῖς πόλεις πέραν τοῦ Ἰορδάνου ἀπὸ ἀνατολῶν ἡλίου,
\vs{42}φυγεῖν ἐκεῖ τὸν φονευτὴν ὃς ἂν φονεύσῃ τὸν πλησίον οὐκ εἰδὼς, καὶ οὗτος οὐ μισῶν αὐτὸν πρὸ τῆς χθὲς καὶ τῆς τρίτης, καὶ καταφεύξεται εἰς μίαν τῶν πόλεων τούτων, καὶ ζήσεται·
\vs{43}τὴν Βοσὸρ ἐν τῇ ἐρήμῳ ἐν τῇ γῇ τῇ πεδινῇ τῷ Ῥουβήν, καὶ τὴν Ῥαμὼθ ἐν Γαλαὰδ τῷ Γαδδί, καὶ τὴν Γαυλὼν ἐν Βασὰν τῷ Μανασσῄ.

\vs{44}Οὗτος ὁ νόμος, ὃν παρέθετο Μωυσῆς ἐνώπιον υἱῶν Ἰσραήλ.
\vs{45}Ταῦτα τὰ μαρτύρια, καὶ τὰ δικαιώματα, καὶ τὰ κρίματα, ὅσα ἐλάλησε Μωυσῆς τοῖς υἱοῖς Ἰσραήλ, ἐξελθόντων αὐτῶν ἐκ γῆς Αἰγύπτου,
\vs{46}ἐν τῷ πέραν τοῦ Ἰορδάνου, ἐν φάραγγι, ἐγγὺς οἴκου Φογώρ, ἐν γῇ Σηὼν βασιλέως τῶν Ἀμοῤῥαίων, ὃς κατῴκει ἐν Ἐσεβὼν, ὅν ἐπάταξε Μωυσῆς, καὶ οἱ υἱοὶ Ἰσραήλ, ἐξελθόντων αὐτῶν ἐκ γῆς Αἰγύπτου.
\vs{47}Καὶ ἐκληρονόμησαν τὴν γῆν αὐτοῦ, καὶ τὴν γῆν Ὢγ βασιλέως τῆς Βασὰν, δύο βασιλέων τῶν Ἀμοῤῥαίων, οἳ ἦσαν πέραν τοῦ Ἰορδάνου κατὰ ἀνατολὰς ἡλίου,
\vs{48}ἀπὸ Ἀροὴρ, ἥ ἐστιν ἐπὶ τοῦ χείλους χειμάῤῥου Ἀρνών, καὶ ἐπὶ τοῦ ὄρους τοῦ Σηὼν, ὅ ἐστιν Ἀερμὼν,
\vs{49}πᾶσαν τὴν Ἄραβα πέραν τοῦ Ἰορδάνου κατὰ ἀνατολὰς ἡλίου ὑπὸ Ἀσηδὼθ τὴν λαξευτήν.

\ch{5}
Καὶ ἐκάλεσε Μωυσῆς πάντα Ἰσραὴλ, καὶ εἶπε πρὸς αὐτούς, ἄκουε Ἰσραὴλ τὰ δικαιώματα καὶ τὰ κρίματα, ὅσα ἐγὼ λαλῶ ἐν τοῖς ὠσὶν ὑμῶν ἐν τῇ ἡμέρᾳ ταύτῃ, καὶ μαθήσεσθε αὐτὰ, καὶ φυλάξεσθε ποιεῖν αὐτά.
\vs{2}Κύριος ὁ Θεὸς ὑμῶν διέθετο πρὸς ὑμᾶς διαθήκην ἐν Χωρήβ.
\vs{3}Οὐχὶ τοῖς πατράσιν ὑμῶν διέθετο Κύριος τὴν διαθήκην ταύτην, ἀλλʼ ἢ πρὸς ὑμᾶς· ὑμεῖς ὧδε πάντες ζῶντες σήμερον.
\vs{4}Πρόσωπον κατὰ πρόσωπον ἐλάλησε Κύριος πρὸς ὑμᾶς ἐν τῷ ὄρει ἐκ μέσου τοῦ πυρός.
\vs{5}Κᾀγὼ εἱστήκειν ἀναμέσον Κυρίου καὶ ὑμῶν ἐν τῷ καιρῷ ἐκείνῳ ἀναγγεῖλαι ὑμῖν τὰ ῥήματα Κυρίου, ὅτι ἐφοβήθητε ἀπὸ προσώπου τοῦ πυρὸς, καὶ οὐκ ἀνέβητε εἰς τὸ ὄρος, λέγων
\vs{6}ἐγώ εἰμι Κύριος ὁ Θεός σου ὁ ἐξαγαγών σε ἐκ γῆς Αἰγύπτοῦ, ἐξ οἴκου δουλείας.

\vs{7}Οὐκ ἔσονταί σοι θεοὶ ἕτεροι πρὸ προσώπου μου.

\vs{8}Οὐ ποιήσεις σεαυτῷ εἴδωλον, οὐδὲ παντὸς ὁμοίωμα, ὅσα ἐν τῷ οὐρανῷ ἄνω, καὶ ὅσα ἐν τῇ γῇ κάτω, καὶ ὅσα ἐν τοῖς ὕδασιν ὑποκάτω τῆς γῆς.
\vs{9}Οὐ πρόσκυνήσεις αὐτοῖς, οὐδὲ μὴ λατρεύσῃς αὐτοῖς· ὅτι ἐγώ εἰμι Κύριος ὁ Θεός σου, Θεὸς ζηλωτὴς, ἀποδιδοὺς ἁμαρτίας πατέρων ἐπὶ τέκνα ἐπὶ τρίτην καὶ τετάρτην γενεὰν τοῖς μισοῦσί με,
\vs{10}καὶ ποιῶν ἔλεος εἰς χιλιάδας τοῖς ἀγαπῶσί με, καὶ τοῖς φυλάσσουσι τὰ προστάγματά μου.
\vs{11}Οὐ λήψῃ τὸ ὄνομα Κυρίου τοῦ Θεοῦ σου ἐπὶ ματαίῳ· οὐ γὰρ μὴ καθαρίσῃ Κύριος ὁ Θεός σου τὸν λαμβάνοντα τὸ ὄνομα αὐτοῦ ἐπὶ ματαίῳ.

\vs{12}Φύλαξαι τὴν ἡμέραν τῶν σαββάτων ἁγιάζειν αὐτὴν, ὃν τρόπον ἐνετείλατό σοι Κύριος ὁ Θεός σου.
\vs{13}Ἓξ ἡμέρας ἐργᾷ καὶ ποιήσεις πάντα τὰ ἔργα σου·
\vs{14}τῇ δὲ ἡμέρᾳ τῇ ἑβδόμῃ σάββατα Κυρίῳ τῷ Θεῷ σου· οὐ ποιήσεις ἐν αὐτῇ πᾶν ἔργον σὺ καὶ ὁ υἱός σου καὶ ἡ θυγάτηρ σου, ὁ παῖς σου καὶ ἡ παιδίσκη σου, ὁ βοῦς σου καὶ τὸ ὑποζύγιόνσου, καὶ πᾶν κτῆνός σου, καὶ προσήλυτος ὁ παροικῶν ἐν σοὶ· ἵνα ἀναπαύσηται ὁ παῖς σου, καὶ ἡ παιδίσκη σου, καὶ τὸ ὑποζύγίον σου, ὥσπερ καὶ σύ.
\vs{15}Καὶ μνησθήσῃ ὅτι οἰκέτης ἦσθα ἐν γῇ Αἰγύπτῳ, καὶ ἐξήγαγέ σε Κύριος ὁ Θεός σου ἐκεῖθεν ἐν χειρὶ κραταιᾷ, καὶ ἐν βραχίονι ὑψηλῷ· διὰ τοῦτο συνέταξέ σοι Κύριος ὁ Θεός σου ὥστε φυλάσσεσθαι τὴν ἡμέραν τῶν σαββάτων καὶ ἁγιάζειν αὐτήν.
\vs{16}Τίμα τὸν πατέρα σου καὶ τὴν μητέρα σου, ὃν τρόπον ἐνετείλατό σοι Κύριος ὁ Θεός σου, ἵνα εὖ σοι γένηται, καὶ ἵνα μακροχρόνιος γένῃ ἐπὶ τῆς γῆς, ἧς Κύριος ὁ Θεός σου δίδωσί σοι.
\vs{17}Οὐ φονεύσεις.
\vs{18}Οὐ μοιχεύσεις.
\vs{19}Οὐ κλέψεις.
\vs{20}Οὐ ψευδομαρτυρήσεις κατὰ τοῦ πλησίον σου μαρτυρίαν ψευδῆ.
\vs{21}Οὐκ ἐπιθυμήσεις τὴν γυναῖκα τοῦ πλησίον σου· οὐκ ἐπιθυμήσεις τὴν οἰκίαν τοῦ πλησίον σου, οὔτε τὸν ἀγρὸν αὐτοῦ, οὔτε τὸν παῖδα αὐτοῦ, οὔτε τὴν παιδίσκην αὐτοῦ, οὔτε τοῦ βοὸς αὐτοῦ, οὔτε τοῦ ὑποζυγίου αὐτοῦ, οὔτε παντὸς κτήνους αὐτοῦ, οὔτε πάντα ὅσα τῷ πλησίον σου ἐστί.

\vs{22}Ταῦτα τὰ ῥήματα ἐλάλησε Κύριος πρὸς πᾶσαν συναγωγὴν ὑμῶν ἐν τῷ ὄρει ἐκ μέσου τοῦ πυρός· σκότος, γνόφος, θύελλα, φωνὴ μεγάλη· καὶ οὐ προσέθηκε· καὶ ἔγραψεν αὐτὰ ἐπὶ δύο πλάκας λιθίνας, καὶ ἔδωκέ μοι.
\vs{23}Καὶ ἐγένετο ὡς ἠκούσατε τὴν φωνὴν ἐκ μέσου τοῦ πυρός, καὶ τὸ ὄρος ἐκαίετο πυρί, καὶ προσήλθετε πρὸς μὲ πάντες οἱ ἡγούμενοι τῶν φυλῶν ὑμῶν, καὶ ἡ γερουσία ὑμῶν,
\vs{24}καὶ ἐλέγετε, ἰδοὺ ἔδειξεν ἡμῖν Κύριος ὁ Θεὸς ἡμῶν τὴν δόξαν αὐτοῦ, καὶ τὴν φωνὴν αὐτοῦ ἠκούσαμεν ἐκ μέσου τοῦ πυρός· ἐν τῇ ἡμέρᾳ ταύτῃ εἴδομεν ὅτι λαλήσει ὁ Θεὸς πρὸς ἄνθρωπον, καὶ ζήσεται.
\vs{25}Καὶ νῦν μὴ ἀποθάνωμεν, ὅτι ἐξαναλώσει ἡμᾶς τὸ πῦρ τὸ μέγα τοῦτο, ἐὰν προσθώμεθα ἡμεῖς ἀκοῦσαι τὴν φωνὴν Κυρίου τοῦ Θεοῦ ἡμῶν ἔτι, καὶ ἀποθανούμεθα.
\vs{26}Τίς γὰρ σὰρξ ἥτις ἤκουσε φωνὴν Θεοῦ ζῶντος, λαλοῦντος ἐκ μέσου τοῦ πυρὸς, ὡς ἡμεῖς, καὶ ζήσεται;
\vs{27}Πρόσελθε σὺ, καὶ ἄκουσον πάντα ὅσα ἂν εἴπῃ Κύριος ὁ Θεὸς ἡμῶν, καὶ σὺ λαλήσεις πρὸς ἡμᾶς πάντα ὅσα ἂν λαλήσει Κύριος ὁ Θεὸς ἡμῶν πρὸς σὲ, καὶ ἀκουσόμεθα, καὶ ποιήσομεν.

\vs{28}Καὶ ἤκουσε Κύριος τὴν φωνὴν τῶν λόγων ὑμῶν λαλούντων πρὸς μέ· καὶ εἶπε Κύριος πρὸς μέ, ἤκουσα τὴν φωνὴν τῶν λόγων τοῦ λαοῦ τούτου ὅσα ἐλάλησαν πρὸς σέ· ὀρθῶς πάντα ὅσα ἐλάλησαν.
\vs{29}Τίς δώσει εἶναι οὕτω τὴν καρδίαν αὐτῶν ἐν αὐτοῖς, ὥστε φοβεῖσθαί με καὶ φυλάσσεσθαι τὰς ἐντολάς μου πάσας τὰς ἡμέρας, ἵνα εὖ ἠ· αὐτοῖς, καὶ τοῖς υἱοῖς αὐτῶν διʼ αἰῶνος;
\vs{30}Βάδισον, εἶπον αὐτοῖς, ἀποστράφητε ὑμεῖς εἰς τοὺς οἴκους ὑμῶν·
\vs{31}σὺ δὲ αὐτοῦ στῆθι μετʼ ἐμοῦ, καὶ λαλήσω πρὸς σὲ τὰς ἐντολὰς καὶ τὰ δικαιώματα καὶ τὰ κρίματα ὅσα διδάξεις αὐτοὺς, καὶ ποιείτωσαν οὕτως ἐν τῇ γῇ ἣν ἐγὼ δίδωμι αὐτοῖς ἐν κλήρῳ.
\vs{32}Καὶ φυλάξεσθε ποιεῖν ὃν τρόπον ἐνετειλατό σοι Κύριος ὁ Θεός σου· οὐκ ἐκκλινεῖτε εἰς δεξιὰ οὐδὲ εἰς ἀριστερά,
\vs{33}κατὰ πᾶσαν τὴν ὁδὸν, ἣν ἐνετείλατό σοι Κύριος ὁ Θεός σου πορεύεσθαι ἐν αὐτῇ, ὅπως καταπαύσῃ σε, καὶ εὖ σοι ἠ·, καὶ μακροημερεύσητε ἐπὶ τῆς γῆς ἣν κληρονομήσετε.

\ch{6}
Καὶ αὗται αἱ ἐντολαὶ καὶ τὰ δικαιώματα καὶ τὰ κρίματα ὅσα ἐνετείλατο Κύριος ὁ Θεὸς ἡμῶν διδάξαι ὑμᾶς ποιεῖν οὕτως ἐν τῇ γῇ, εἰς ἣν ὑμεῖς εἰσπορεύεσθε ἐκεῖ κληρονομῆσαι αὐτήν.
\vs{2}Ἵνα φοβῆσθε Κύριον τὸν Θεὸν ὑμῶν, φυλάσσεσθε πάντα τὰ δικαιώματα αὐτοῦ, καὶ τὰς ἐντολὰς αὐτοῦ, ἃς ἐγὼ ἐντέλλομαί σοι σήμερον, σὺ καὶ οἱ υἱοί σου, καὶ οἱ υἱοὶ τῶν υἱῶν σου πάσας τὰς ἡμέρας τῆς ζωῆς σου, ἵνα μακροημερεύσητε.

\vs{3}Καὶ ἄκουσον Ἰσραὴλ, καὶ φύλαξον ποιεῖν, ὅπως εὖ σοι ἠ·, καὶ ἵνα πληθυνθῆτε σφόδρα, καθάπερ ἐλάλησε Κύριος ὁ Θεὸς τῶν πατέρων σου δοῦναί σοι γῆν ῥέουσαν γάλα καὶ μέλι·
\vs{4}καὶ ταῦτα τὰ δικαιώματα καὶ τὰ κρίματα, ὅσα ἐνετείλατο Κύριος τοῖς υἱοῖς Ἰσραὴλ ἐν τῇ ἐρήμῳ, ἐξελθόντων αὐτῶν ἐκ γῆν Αἰγύπτου. Ἄκουε Ἰσραὴλ, Κύριος ὁ Θεὸς ἡμῶν, Κύριος εἷς ἐστι.
\vs{5}Καὶ ἀγαπήσεις Κύριον τὸν Θεόν σου ἐξ ὅλης τῆς διανοίας σου, καὶ ἐξ ὅλης τῆς ψυχῆς σου, καὶ ἐξ ὅλης τῆς δυνάμεώς σου.
\vs{6}Καὶ ἔσται τὰ ῥήματα ταῦτα, ὅσα ἐγὼ ἐντέλλομαί σοι σήμερον, ἐν τῇ καρδίᾳ σου, καὶ ἐν τῇ ψυχῇ σου.
\vs{7}Καὶ προβιβάσεις αὐτὰ τοὺς υἱούς σου, καὶ λαλήσεις ἐν αὐτοῖς καθήμενος ἐν οἴκῳ, καὶ πορευόμενος ἐν ὁδῷ, καὶ κοιταζόμενος, καὶ διανιστάμενος.
\vs{8}Καὶ ἀφάψεις αὐτὰ εἰς σημεῖον ἐπὶ τῆς χειρός σου, καὶ ἔσται ἀσάλευτον πρὸ ὀφθαλμῶν σου.
\vs{9}Καὶ γράψετε αὐτὰ ἐπὶ τὰς φλιὰς τῶν οἰκιῶν ὑμῶν, καὶ τῶν πυλῶν ὑμῶν.

\vs{10}Καὶ ἔσται ὅταν εἰσαγάγῃ σε Κύριος ὁ Θεός σου εἰς τὴν γῆν ἣν ὤμοσε τοῖς πατράσι σου, τῷ Ἁβραὰμ, καὶ τῷ Ἰσαακ, καὶ τῷ Ἰακώβ, δοῦναί σοι πόλεις μεγάλας καὶ καλὰς ἃς οὐκ ᾠκοδόμησας,
\vs{11}οἰκίας πλήρεις πάντων ἀγαθῶν ἃς οὐκ ἐνέπλησας, λάκκους λελατομημένους οὓς οὐκ ἐξελατόμησας, ἀμπελῶνας καὶ ἐλαιῶνας οὓς οὐ κατεφύτευσας, καὶ φαγὼν καὶ ἐμπλησθεὶς, πρόσεχε σεαυτῷ μὴ ἐπιλάθῃ
\vs{12}Κυρίου τοῦ Θεοῦ σου τοῦ ἐξαγαγόντος σε ἐκ γῆς Αἰγύπτου, ἐξ οἴκου δουλείας.
\vs{13}Κύριον τὸν Θεόν σου φοβήθήσῃ, καὶ αὐτῷ μόνῳ λατρεύσεις, καὶ πρὸς αὐτὸν κολληθήσῃ, καὶ ἐπὶ τῷ ὀνόματι αὐτοῦ ὀμῇ.

\vs{14}Οὐ πορεύεσθε ὀπίσω θεῶν ἑτέρων ἀπὸ τῶν θεῶν τῶν ἐθνῶν τῶν περικύκλῳ ὑμῶν,
\vs{15}ὅτι ὁ Θεὸς ζηλωτὴς Κύριος ὁ Θεός σου ἐν σοί· μὴ ὀργισθεὶς θυμῷ Κύριος ὁ Θεός σου σοὶ, ἐξολοθρεύσῃ σε ἀπὸ προσώπου τῆς γῆς.

\vs{16}Οὐκ ἐκπειράσεις Κύριον τὸν Θεόν σου, ὃν τρόπον ἐξεπειράσατε ἐν τῷ πειρασμῷ.
\vs{17}Φυλάσσων φυλάξῃ τὰς ἐντολὰς Κυρίου τοῦ Θεοῦ σου, τὰ μαρτύρια, καὶ τὰ δικαιώματα, ὅσα ἐνετείλατό σοι.
\vs{18}Καὶ ποιήσεις τὸ ἀρεστὸν καὶ τὸ καλὸν ἔναντι Κυρίου τοῦ Θεοῦ σου, ἵνα εὖ σοι γένηται, καὶ εἰσέλθῃς καὶ κληρονομήσῃς τὴν γῆν τὴν ἀγαθὴν, ἣν ὤμοσε Κύριος τοῖς πατράσιν ὑμῶν,
\vs{19}ἐκδιῶξαι πάντας τοὺς ἐχθρούς σου πρὸ προσώπου σου, καθὰ ἐλάλησε Κύριος.

\vs{20}Καὶ ἔσται ὅταν ἐρωτήσῃ σε ὁ υἱός σου αὔριον, λέγων, τί ἐστι τὰ μαρτύρια, καὶ τὰ δικαιώματα καὶ τὰ κρίματα, ὅσα ἐνετείλατο Κύριος ὁ Θεὸς ἡμῶν ἡμῖν;
\vs{21}Καὶ ἐρεῖς τῷ υἱῷ σου, οἰκέται ἦμεν τῷ Φαραὼ ἐν γῇ Αἰγύπτῳ, καὶ ἐξήγαγεν ἡμᾶς Κύριος ἐκεῖθεν ἐν χειρὶ κραταιᾷ, καὶ ἐν βραχίονι ὑψηλῷ.
\vs{22}Καὶ ἔδωκε Κύριος σημεῖα καὶ τέρατα μεγάλα καὶ πονηρὰ ἐν Αἰγύπτῳ ἐν Φαραὼ καὶ ἐν τῷ οἴκῳ αὐτοῦ ἐνώπιον ἡμῶν,
\vs{23}καὶ ἡμᾶς ἐξήγαγεν ἐκεῖθεν δαῦναι ἡμῖν τὴν γῆν ταύτην, ἣν ὤμοσε δοῦναι τοῖς πατράσιν ἡμῶν.
\vs{24}Καὶ ἐνετείλατο ἡμῖν Κύριος ποιεῖν πάντα τὰ δικαιώματα ταῦτα· φοβεῖσθαι Κύριον τὸν Θεὸν ἡμῶν, ἵνα εὖ ἠ· ἡμῖν πάσας τὰς ἡμέρας, ἵνα ζῶμεν ὥσπερ καὶ σήμερον.
\vs{25}Καὶ ἐλεημοσύνη ἔσται ἡμῖν, ἐὰν φυλασσώμεθα ποιεῖν πάσας τὰς ἐντολὰς ταύτας ἐναντίον Κυρίου τοῦ Θεοῦ ἡμῶν, καθὰ ἐνετείλατο ἡμῖν.

\ch{7}
Ἐὰν δὲ εἰσάγῃ σε Κύριος ὁ Θεός σου εἰς τὴν γῆν, εἰς ἣν εἰσπορεύῃ ἐκεῖ κληρονομῆσαι αὐτὴν, καὶ ἐξάρῃ ἔθνη μεγάλα ἀπὸ προσώπου σου, τὸν Χετταῖον καὶ Γεργεσαῖον καὶ Ἀμοῤῥαῖον καὶ Χαναναῖον καὶ Φερεζαῖον καὶ Εὐαῖον καὶ Ἰεβουσαῖον, ἐπτὰ ἔθην πολλὰ καὶ ἰσχυρότερα ὑμῶν·
\vs{2}Καὶ παραδώσει αὐτοὺς Κύριος ὁ Θεός σου εἰς τὰς χεῖράς σου, καὶ πατάξεις αὐτούς· ἀφανισμῷ ἀφανιεῖς αὐτούς· οὐ διαθήσῃ πρὸς αὐτοὺς διαθήκην, οὐδὲ μὴ ἐλεήσητε αὐτούς,
\vs{3}οὐδὲ μὴ γαμβρεύσητε πρὸς αὐτούς· τὴν θυγατέρα σου οὐ δώσεις τῷ υἱῷ αὐτοῦ, καὶ τὴν θυγατέρα αὐτοῦ οὐ λήψῃ τῷ υἱῷ σου.
\vs{4}Ἀποστήσει γὰρ τὸν υἱόν σου ἀπʼ ἐμοῦ, καὶ λατρεύσει θεοῖς ἑτέροις· καὶ ὀργισθήσεται θυμῷ Κύριος εἰς ὑμᾶς, καὶ ἐξολοθρεύσει σε τοτάχος.
\vs{5}Ἀλλʼ οὕτω ποιήσετε αὐτοῖς· τοὺς βωμοὺς αὐτῶν καθελεῖτε, καὶ τὰς στήλας αὐτῶν συντρίψετε, καὶ τὰ ἄλση αὐτῶν ἐκκόψετε, καὶ τὰ γλυπτὰ τῶν θεῶν αὐτῶν κατακαύσετε πυρί.
\vs{6}Ὅτι λαὸς ἅγιος εἶ Κυρίῳ τῷ Θεῷ σου· καὶ σὲ προείλετο Κύριος ὁ Θεός σου εἶναι αὐτῷ λαὸν περιούσιον παρὰ πάντα τὰ ἔθνη, ὅσα ἐπὶ προσώπου τῆς γῆς.

\vs{7}Οὐχ ὅτι πολυπληθεῖτε παρὰ πάντα τὰ ἔθνη, προείλετο Κύριος ὑμᾶς, καὶ ἐξελέξατο Κύριος ὑμᾶς· ὑμεῖς γάρ ἐστε ὀλιγοστοὶ παρὰ πάντα τὰ ἔθνη.
\vs{8}Ἀλλὰ παρὰ τὸ ἀγαπᾷν Κύριον ὑμᾶς, καὶ διατηρῶν τὸν ὅρκον ὃν ὤμοσε τοῖς πατράσιν ὑμῶν, ἐξήγαγεν ὑμᾶς Κύριος ἐν χειρὶ κραταιᾷ, καὶ ἐλυτρώσατό σε Κύριος ἐξ οἴκου δουελίας, ἐκ χειρὸς Φαραὼ βασιλέως Αἰγύπτου.
\vs{9}Καὶ γνώσῃ, ὅτι Κύριος ὁ Θεός σου, οὗτος Θεός· Θεὸς πιστός, ὁ φυλάσσων διαθήκην καὶ ἔλεος τοῖς ἀγαπῶσιν αὐτὸν καὶ τοῖς φυλάσσουσι τὰς ἐντολὰς αὐτοῦ εἰς χιλίας γενεάς,
\vs{10}καὶ ἀποδιδοὺς τοῖς μισοῦσι κατὰ πρόσωπον ἐξολοθρεῦσαι αὐτούς· καὶ οὐχὶ βραδυνεῖ τοίς μισοῦσι· κατὰ πρόσωπον ἀποδώσει αὐτοῖς.

\vs{11}Καὶ φυλάξῃ τὰς ἐντολὰς, καὶ τὰ δικαιώματα, καὶ τὰ κρίματα ταῦτα, ὅσα ἐγὼ ἐντέλλομαί σοι σήμερον ποιεῖν.
\vs{12}Καὶ ἔσται ἡνίκα ἂν ἀκούσητε τὰ δικαιώματα ταῦτα, καὶ φυλάξητε καὶ ποιήσητε αὐτὰ, καὶ διαφυλάξει Κύριος ὁ Θεός σου σοὶ τὴν διαθήκην καὶ τὸ ἔλεος, ὃ ὤμοσε τοῖς πατράσιν ὑμῶν.
\vs{13}Καὶ ἀγαπήσει σε, καὶ εὐλογήσει σε, καὶ πληθυνεῖ σε, καὶ εὐλογήσει τὰ ἔγγονα τῆς κοιλίας σου, καὶ τὸν καρπὸν τῆς γῆς σου, τὸν σῖτόν σου, καὶ τὸν οἶνόν σου, καὶ τὸ ἔλαιόν σου, τὰ βουκόλια τῶν βοῶν σου, καὶ τὰ ποίμνια τῶν προβάτων σου ἐπὶ τῆς γῆς, ἧς ὤμοσε Κύριος τοῖς πατράσι σου δοῦναί σοι.
\vs{14}Εὐλογητὸς ἔσῃ παρὰ πάντα τὰ ἔθνη· οὐκ ἔσται ἐν ὑμῖν ἄγονος, οὐδὲ στεῖρα, καὶ ἐν τοῖς κτήνεσί σου.
\vs{15}Καὶ περιελεῖ Κύριος ὁ Θεός σου ἀπὸ σοῦ πᾶσαν μαλακίαν, καὶ πάσας νόσους Αἰγύπτου τὰς πονηρὰς, ἃς ἑώρακας, καὶ ὅσα ἔγνως, οὐκ ἐπιθήσει ἐπὶ σὲ· καὶ ἐπιθήσει αὐτὰ ἐπὶ πάντας τοὺς μισοῦντὰς σε.

\vs{16}Καὶ φαγῇ πάντα τὰ σκῦλα τῶν ἐθνῶν, ἃ Κύριος ὁ Θεός σου δίδωσί σοι· οὐ φείσεται ὁ ὀφθαλμός σου ἐπʼ αὐτοῖς, καὶ οὐ μὴ λατρεύσῃς τοῖς θεοῖς αὐτῶν· ὅτι σκῶλον τοῦτό ἐστί σοί.

\vs{17}Ἐὰν δὲ λέγῃς ἐν τῇ διανοίᾳ σου, ὅτι πολὺ τὸ ἔθνος τοῦτο ἢ ἐγώ, πῶς δυνήσομαι ἐξολοθρεῦσαι αὐτούς;
\vs{18}Οὐ φοβηθήσῃ αὐτούς· μνείᾳ μνησθήσῃ, ὅσα ἐποίησε Κύριος ὁ Θεός σου τῷ Φαραὼ καὶ πᾶσι τοῖς Αἰγυπτίοις·
\vs{19}Τοὺς πειρασμοὺς τοὺς μεγάλους, οὓς ἴδοσαν οἱ ὀφθαλμοί σου, τὰ σημεῖα καὶ τὰ τέρατα τὰ μεγάλα ἐκεῖνα, τὴν χεῖρα τὴν κραταιὰν, καὶ τὸν βραχίονα τὸν ὑψηλόν· ὡς ἐξήγαγέ σε Κύριος ὁ Θεός σου, οὕτως ποιήσει Κύριος ὁ Θεὸς ὑμῶν πᾶσιν τοῖς ἔθνεσιν, οὓς σὺ φοβῇ ἀπὸ προσώπου αὐτῶν.
\vs{20}Καὶ τὰς σφηκίας ἀποστελεῖ Κύριος ὁ Θεός σου εἰς αὐτοὺς, ἕως ἂν ἐκτριβῶσιν οἱ καταλελειμμένοι καὶ οἱ κεκρυμμένοι ἀπὸ σοῦ·
\vs{21}Οὐ τρωθήσῃ ἀπὸ προσώπου αὐτῶν, ὅτι Κύριος ὁ Θεός σου ἐν σοί, Θεὸς μέγας καὶ κραταιός.
\vs{22}Καὶ καταναλώσει Κύριος ὁ Θεός σου τὰ ἔθνη ταῦτα ἀπὸ προσώπου σου κατὰ μικρὸν μικρόν· οὐ δυνήσῃ ἐξαναλῶσαι αὐτοὺς τοτάχος, ἵνα μὴ γένηται ἡ γῆ ἔρημος, καὶ πληθυνθῇ ἐπὶ σὲ τὰ θηρία τὰ ἄγρια.
\vs{23}Καὶ παραδώσει αὐτοὺς Κύριος ὁ Θεός σου εἰς τὰς χεῖράς σου, καὶ ἀπολεῖς αὐτοὺς ἀπωλείᾳ μεγάλῃ, ἕως ἂν ἐξολοθρεύσητε αὐτούς·
\vs{24}Καὶ παραδώσει τοὺς βασιλεῖς αὐτῶν εἰς τὰς χεῖρας ὑμῶν, καὶ ἀπολεῖτε τὸ ὄνομα αὐτῶν ἐκ τοῦ τόπου ἐκείνου· οὐκ ἀντιστήσεται οὐθεὶς κατὰ πρόσωπόν σου, ἕως ἂν ἐξολοθρεύσῃς αὐτούς.

\vs{25}Τὰ γλυπτὰ τῶν Θεῶν αὐτῶν καύσετε πυρί· οὐκ ἐπιθυμήσεις ἀργύριον, οὐδὲ χρυσίον ἀπʼ αὐτῶν οὐ λήψῃ σεαυτῷ, μὴ πταίσῃς διʼ αὐτὸ, ὅτι βδέλυγμα Κυρίῳ τῷ Θεῷ σού ἐστί.
\vs{26}Καὶ οὐκ εἰσοίσεις βδέλυγμα εἰς τὸν οἶκόν σου, καὶ ἀνάθεμα ἔσῃ ὥσπερ τοῦτο· προσοχθίσματι προσοχθιεῖς, καὶ βδελύγματι βδελύξῃ, ὅτι ἀνάθεμά ἐστι.

\ch{8}
Πάσας τὰς ἐντολὰς, ἃς ἐγὼ ἐντέλλομαι ὑμῖν σήμερον, φυλάξεσθε ποιεῖν, ἵνα ζῆτε καὶ πολυπλασιασθῆτε, καὶ εἰσέλθητε καὶ κληρονομήσητε τὴν γῆν, ἣν ὤμοσε Κύριος ὁ Θεὸς ὑμῶν τοῖς πατράσιν ὑμῶν.
\vs{2}Καὶ μνησθήσῃ πᾶσαν τὴν ὁδὸν, ἣν ἤγαγέ σε Κύριος ὁ Θεός σου ἐν τῇ ἐρήμῳ, ὅπως ἂν κακώσῃ σε καὶ πειράσῃ σε, καὶ διαγνωσθῇ τὰ ἐν τῇ καρδίᾳ σου, εἰ φυλάξῃ τὰς ἐντολὰς αὐτοῦ ἢ οὔ.
\vs{3}Καὶ ἐκάκωσέ σε, καὶ ἐλιμαγχόνησέ σε, καὶ ἐψώμισέ σε τὸ μάννα, ὃ οὐκ ᾔδεισαν οἱ πατέρες σου· ἵνα ἀναγγείλῃ σοι, ὅτι οὐκ ἐπʼ ἄρτῳ μόνῳ ζήσεται ὁ ἄνθρωπος, ἀλλʼ ἐπὶ παντὶ ῥήματι τῷ ἐκπορευομένῳ διὰ στόματος Θεοῦ ζήσεται ὁ ἄνθρωπος.
\vs{4}Τὰ ἱμάτιά σου οὐκ ἐπαλαιώθη ἀπὸ σοῦ, τὰ ὑποδήματά σου οὐ κατετρίβη ἀπὸ σοῦ· οἱ πόδες σου οὐκ ἐτυλώθησαν, ἰδοὺ τεσσαράκοντα ἔτη.

\vs{5}Καὶ γνώσῃ τῇ καρδίᾳ σου, ὅτι ὡς εἴτις ἄνθρωπος παιδεύσῃ τὸν υἱὸν αὐτοῦ, οὕτως Κύριος ὁ Θεός σου παιδεύσει σε.
\vs{6}Καὶ φυλάξῃ τὰς ἐντολὰς Κυρίου τοῦ Θεοῦ σου πορεύεσθαι ἐν ταῖς ὁδοῖς αὐτοῦ, καὶ φοβεῖσθαι αὐτόν.

\vs{7}Ὁ γὰρ Κύριος ὁ Θεός σου εἰσάξει σε εἰς γῆν ἀγαθὴν καὶ πολλὴν, οὗ χείμαῤῥοι ὑδάτων, καὶ πηγαὶ ἀβύσσων ἐκπορευόμεναι διὰ τῶν πεδίων καὶ διὰ τῶν ὀρέων·
\vs{8}Γῆ πυροῦ καὶ κριθῆς, ἄμπελοι, συκαῖ, ῥοαί· γῆ ἐλαίας ἐλαίου καὶ μέλιτος·
\vs{9}γῆ ἐφʼ ἧς οὐ μετὰ πτωχείας φαγῇ τὸν ἄρτον σου, καὶ οὐκ ἐνδεηθήσῃ ἐπʼ αὐτῆς οὐδέν· γῆ ἧς οἱ λίθοι σίδηρος, καὶ ἐκ τῶν ὀρέων αὐτῆς μεταλλεύσεις χαλκόν.

\vs{10}Καὶ φαγῇ καὶ ἐμπλησθήσῃ, καὶ εὐλογήσεις Κύριον τὸν Θεόν σου ἐπὶ τῆς γῆς τῆς ἀγαθῆς, ἧς δέδωκέ σοι.
\vs{11}Πρόσεχε σεαυτῷ μὴ ἐπιλάθῃ Κυρίου τοῦ Θεοῦ σου, τοῦ μὴ φυλάξαι τὰς ἐντολὰς αὐτοῦ, καὶ τὰ κρίματα καὶ τὰ δικαιώματα αὐτοῦ, ὅσα ἐγὼ ἐντέλλομαί σοι σήμερον·
\vs{12}Μὴ φαγὼν καὶ ἐμπλησθεὶς, καὶ οἰκίας καλὰς οἰκοδομήσας καὶ κατοικήσας ἐν αὐταῖς,
\vs{13}καὶ τῶν βοῶν σου καὶ τῶν προβάτων σου πληθυνθέντων σοι, ἀργυρίου καὶ χρυσίου πληθυνθέντος σοι, καὶ πάντων ὅσων σοι ἔσται πληθυνθέντων σοι,
\vs{14}ὑψωθῇς τῇ καρδίᾳ, καὶ ἐπιλάθῃ Κυρίου τοῦ Θεοῦ σου, τοῦ ἐξαγαγόντος σε ἐκ γῆς Αἰγύπτου, ἐξ οἴκου δουλείας·
\vs{15}τοῦ ἀγαγόντος σε διὰ τῆς ἐρήμου τῆς μεγάλης καὶ τῆς φοβερᾶς ἐκείνης, οὗ ὄφις δάκνων, καὶ σκορπίος, καὶ δίψα, οὗ οὐκ ἦν ὕδωρ· τοῦ ἐξαγαγόντος σοι ἐκ πέτρας ἀκροτόμου πηγὴν ὕδατος·
\vs{16}τοῦ ψωμίσαντός σε τὸ μάννα ἐν τῇ ἐρήμῳ ὃ οὐκ ᾔδεις σὺ, καὶ οὐκ ᾔδεισαν οἱ πατέρες σου, ἵνα κακώσῃ σε, καὶ ἐκπειράσῃ σε, καὶ εὖ σε ποιήσῃ ἐπʼ ἐσχάτων τῶν ἡμερῶν σου.
\vs{17}Μὴ εἴπῃς ἐν τῇ καρδίᾳ σου, ἡ ἰσχύς μου, καὶ τὸ κράτος τῆς χειρός μου ἐποίησέ μοι τὴν δύναμιν τὴν μεγάλην ταύτην.
\vs{18}Καὶ μνησθήσῃ Κυρίου τοῦ Θεοῦ σου, ὅτι αὐτός σοι δίδωσιν ἰσχὺν τοῦ ποιῆσαι δύναμιν, καὶ ἵνα στήσῃ τὴν διαθήκην αὐτοῦ ἣν ὤμοσε Κύριος τοῖς πατράσι σου, ὡς σήμερον.

\vs{19}Καὶ ἔσται ἐὰν λήθῃ ἐπιλάθῃ Κυρίου τοῦ Θεοῦ σου, καὶ πορευθῇς, ὀπίσω θεῶν ἑτέρων, καὶ λατρεύσῃς αὐτοῖς, καὶ προσκυνήσῃς αὐτοῖς, διαμαρτύρομαι ὑμῖν σήμερον τόν τε οὐρανὸν καὶ τὴν γῆν, ὅτι ἀπωλείᾳ ἀπολεῖσθε.
\vs{20}Καθὰ καὶ τὰ λοιπὰ ἔθνη ὅσα Κύριος ὁ Θεὸς ἀπολλύει πρὸ προσώπου ὑμῶν, οὕτως ἀπολεῖσθε, ἀνθʼ ὧν οὐκ ἠκούσατε τῆς φωνῆς Κυρίου τοῦ Θεοῦ ὑμῶν.

\ch{9}
Ἄκουε Ἰσραήλ· σὺ διαβαίνεις σήμερον τὸν Ἰορδάνην εἰσελθεῖν κληρονομῆσαι ἔθνη μεγάλα καὶ ἰσχυρότερα μᾶλλον ἢ ὑμεῖς, πόλεις μεγάλας καὶ τειχήρεις ἕως τοῦ οὐρανοῦ,
\vs{2}λαὸν μέγαν καὶ πολὺν καὶ εὐμήκη, υἱοὺς Ἐνάκ, οὓς σὺ οἶσθα, καὶ σὺ ἀκήκοας, τίς ἀντιστήσεται κατὰ πρόσωπον υἱῶν Ἐνάκ;
\vs{3}Καὶ γνώσῃ σήμερον, ὅτι Κύριος ὁ Θεός σου οὗτος προπορεύσεται πρὸ προσώπου σου· πῦρ καταναλίσκον ἐστίν· οὗτος ἐξολοθρεύσει αὐτούς, καὶ οὗτος ἀποστρέψει αὐτοὺς ἀπὸ προσώπου σου, καὶ ἀπολεῖ αὐτοὺς ἐν τάχει, καθάπερ εἶπέ σοι Κύριος.
\vs{4}Μὴ εἴπῃς ἐν τῇ καρδίᾳ σου ἐν τῷ ἐξαναλῶσαι Κύριον τὸν Θεόν σου τὰ ἔθνη ταῦτα πρὸ προσώπου σου, λέγων, διὰ τὴν δικαιοσύνην μου εἰσήγαγέ με Κύριος κληρονομῆσαι τὴν γὴν τὴν ἀγαθὴν ταύτην.
\vs{5}Οὐχὶ διὰ τὴν δικαιοσύνην σου, οὐδὲ διὰ τὴν ὁσιότητα τῆς καρδίας σου σὺ εἰσπορεύῃ κληρονομῆσαι τὴν γῆν αὐτῶν, ἀλλὰ διὰ τὴν ἀσέβειαν τῶν ἐθνῶν τούτων Κύριος ἐξολοθρεύσει αὐτοὺς ἀπὸ προσώπου σου, καὶ ἵνα στήσῃ τὴν διαθήκην, ἥ ὤμοσε Κύριος τοῖς πατράσιν ἡμῶν τῷ Ἁβραὰμ καὶ τῷ Ἰσαὰκ καὶ τῷ Ἰακώβ.

\vs{6}Καὶ γνώσῃ σήμερον, ὅτι οὐχὶ διὰ τὰς δικαιοσύνας σου Κύριος ὁ Θεός σου δίδωσί σοι τὴν γῆν τὴν ἀγαθὴν ταύτην κληρονομῆσαι, ὅτι λαὸς σκληροτράχηλος εἶ.
\vs{7}Μνήσθητι, μὴ ἐπιλάθῃ ὅσα παρώξυνας Κύριον τὸν Θεόν σου ἐν τῇ ἐρήμῳ· ἀφʼ ἧς ἡμέρας ἐξήλθετε ἐξ Αἰγύπτου, καὶ ἤλθετε εἰς τὸν τόπον τοῦτον, ἀπειθοῦντες διετελεῖτε τὰ πρὸς Κύριον.

\vs{8}Καὶ ἐν Χωρὴβ παρωξύνατε Κύριον, καὶ ἐθυμώθη Κύριος ἐφʼ ὑμῖν ἐξολεθρεῦσαι ὑμᾶς,
\vs{9}ἀναβαίνοντός μου εἰς τὸ ὄρος λαβεῖν τὰς πλάκας τὰς λιθίνας, πλάκας διαθήκης ἃς διέθετο Κύριος πρὸς ὑμᾶς, καὶ κατεγενόμην ἐν τῷ ὄρει τεσσαράκοντα ἡμέρας καὶ τεσσεράκοντα νύκτας, ἄρτον οὐκ ἔφαγον καὶ ὕδωρ οὐκ ἔπιον.
\vs{10}Καὶ ἔδωκέ μοι Κύριος τὰς δύο πλάκας τὰς λιθίνας γεγραμμένας ἐν τῷ δακτύλῳ τοῦ Θεοῦ, καὶ ἐπʼ αὐταῖς ἐγέγραπτο πάντες οἱ λόγοι οὓς ἐλάλησε Κύριος πρὸς ὑμᾶς ἐν τῷ ὄρει ἡμέρᾳ ἐκκλησίας·
\vs{11}Καὶ ἐγένετο διὰ τεσσεράκοντα ἡμερῶν καὶ διὰ τεσσαράκοντα νυκτῶν ἔδωκε Κύριος ἐμοὶ τὰς δύο πλάκας τὰς λιθίνας, πλάκας διαθήκης.
\vs{12}Καὶ εἶπε Κύριος πρὸς μέ, ἀνάστηθι, κατάβηθι τοτάχος ἐντεῦθεν, ὅτι ἠνόμησεν ὁ λαός σου, οὓς ἐξήγαγες ἐκ γῆς Αἰγύπτου· παρέβησαν ταχὺ ἐκ τῆς ὁδοῦ ἧς ἐνετείλω αὐτοῖς, καὶ ἐποίησαν ἑαυτοῖς χώνευμα.

\vs{13}Καὶ εἶπε Κύριος πρὸς μὲ, λέγων, λελάληκα πρὸς σὲ ἅπαξ καὶ δὶς, λέγων, ἑώρακα τὸν λαὸν τοὗτον, καὶ ἰδοὺ λαὸς σκληροτράχηλός ἐστι·
\vs{14}Καὶ νῦν ἔασόν με ἐξολοθρεῦσαι αὐτούς, καὶ ἐξαλείψω τὸ ὄνομα αὐτῶν ὑποκάτωθεν τοῦ οὐρανοῦ, καὶ ποιήσω σε εἰς ἔθνος μέγα, καὶ ἰσχυρὸν, καὶ πολὺ μᾶλλον ἢ τοῦτο.
\vs{15}Καὶ ἐπιστρέψας, κατέβην ἐκ τοῦ ὄρους· καὶ τὸ ὄρος ἐκαίετο πυρὶ ἕως τοῦ οὐρανοῦ· καὶ αἱ δύο πλάκες τῶν μαρτυρίων ἐπὶ ταῖς δυσὶ χερσί μου.
\vs{16}Καὶ ἰδὼν ὅτι ἡμάρτετε ἐναντίον Κυρίου τοῦ Θεοῦ ὑμῶν, καὶ ἐποιήσατε ὑμῖν αὐτοῖς χωνευτόν, καὶ παρέβητε ἀπὸ τῆς ὁδοῦ, ἧς ἐνετείλατο Κύριος ὑμῖν ποιεῖν·
\vs{17}καὶ ἐπιλαβόμενος τῶν δύο πλακῶν, ἔῤῥιψα αὐτὰς ἀπὸ τῶν δύο χειρῶν μου, καὶ συνέτριψα ἐναντίον ὑμῶν.
\vs{18}Καὶ ἐδεήθην ἐναντίον Κυρίου δεύτερον καθάπερ καὶ τὸ πρότερον τεσσαράκοντα ἡμέρας καὶ τεσσαράκοντα νύκτας, ἄρτον οὐκ ἔφαγον καὶ ὕδωρ οὐκ ἔπιον, περὶ πασῶν τῶν ἁμαρτιῶν ὑμῶν ὧν ἡμάρτετε ποιῆσαι τὸ πονηρὸν ἐναντίον Κυρίου τοῦ Θεοῦ παροξῦναι αὐτόν.
\vs{19}Καὶ ἔκφοβός εἰμι διὰ τὸν θυμόν καὶ τὴν ὀργὴν, ὅτι παρωξύνθη Κύριος ἐφʼ ὑμῖν τοῦ ἐξολοθρεῦσαι ὑμᾶς· καὶ εἰσήκουσε Κύριος ἐμοῦ καὶ ἐν τῷ καιρῷ τούτῳ.
\vs{20}Καὶ ἐπὶ Ἀαρὼν ἐθυμώθη ἐξολοθρεῦσαι αὐτόν, καὶ ηὐξάμην καὶ περὶ Ἀαρὼν ἐν τῷ καιρῷ ἐκείνῳ.
\vs{21}Καὶ τὴν ἁμαρτίαν ὑμῶν, ἣν ἐποιήσατε, τὸν μόσχον ἔλαβον αὐτὸν, καὶ κατέκαυσα αὐτὸν ἐν πυρί, καὶ συνέκοψα αὐτὸν καταλέσας σφόδρα ἕως ἐγένετο λεπτόν, καὶ ἐγένετο ὡσεὶ κονιορτός· καὶ ἔῤῥιψα τὸν κονιορτὸν εἰς τὸν χειμάῤῥουν τὸν καταβαίνοντα ἐκ τοῦ ὄρους.

\vs{22}Καὶ ἐν τῷ ἐμπυρισμῷ, καὶ ἐν τῷ πειρασμῷ, καὶ ἐν τοῖς μνήμασι τῆς ἐπιθυμίας παροξύναντες ἦτε Κύριον.
\vs{23}Καὶ ὅτε ἐξαπέστειλεν ὑμᾶς Κύριος ἐκ Κάδης Βαρνὴ, λέγων, ἀνάβητε καὶ κληρονομήσατε τὴν γῆν, ἣν δίδωμι ὑμῖν, καὶ ἠπειθήσατε τῷ ῥήματι Κυρίου τοῦ Θεοῦ ὑμῶν, καὶ οὐκ ἐπιστεύσατε αὐτῷ, καὶ οὐκ εἰσηκούσατε τῆς φωνῆς αὐτοῦ.
\vs{24}Ἀπειθοῦντες ἦτε τὰ πρὸς Κύριον ἀπὸ τῆς ἡμέρας ἧς ἐγνώσθη ὑμῖν.
\vs{25}Καὶ ἐδεήθην ἔναντι Κυρίου τεσσαράκοντα ἡμέρας καὶ τεσσαράκοντα νύκτας, ὅσας ἐδεήθην· εἶπε γὰρ Κύριος ἐξολοθρεῦσαι ὑμᾶς.
\vs{26}Καὶ ηὐξάμην πρὸς τὸν Θεὸν, καὶ εἶπα, Κύριε βασιλεῦ τῶν θεῶν, μὴ ἐξολοθρεύσῃς τὸν λαόν σου καὶ τὴν μερίδα σου, ἣν ἐλυτρώσω, οὓς ἐξήγαγες ἐκ γῆς Αἰγύπτου ἐν τῇ ἰσχύϊ σου τῇ μεγάλῃ, καὶ ἐν τῇ χειρί σου τῇ κραταιᾷ, καὶ ἐν τῷ βραχίονί σου τῷ ὑψηλῷ.
\vs{27}Μνήσθητι Ἁβραὰμ καὶ Ἰσαὰκ καὶ Ἰακὼβ τῶν θεραπόντων σου, οἷς ὤμοσας κατὰ σεαυτοῦ· μὴ ἐπιβλέψῃς ἐπὶ τὴν σκληρότητα τοῦ λαοῦ τούτου, καὶ τὰ ἀσεβήματα, καὶ ἐπὶ τὰ ἁμαρτήματα αὐτῶν.
\vs{28}Μὴ εἴπωσιν οἱ κατοικοῦντες τὴν γῆν ὅθεν ἐξήγαγες ἡμᾶς ἐκεῖθεν, λέγοντες, παρὰ τὸ μὴ δύνασθαι Κύριον εἰσαγαγεῖν αὐτοὺς εἰς τὴν γῆν ἣν εἶπεν αὐτοῖς, καὶ παρὰ τὸ μισῆσαι αὐτοὺς, ἐξήγαγεν αὐτοὺς ἐν τῇ ἐρήμῳ ἀποκτεῖναι αὐτούς.
\vs{29}Καὶ οὗτοι λαός σου καὶ κλῆρός σου, οὓς ἐξήγαγες ἐκ γῆς Αἰγύπτου ἐν τῇ ἰσχύϊ σου τῇ μεγάλῃ, καὶ ἐν τῇ χειρί σου τῇ κραταιᾷ, καὶ ἐν τῷ βραχίονί σου τῷ ὑψηλῷ.

\ch{10}
Ἐν ἐκείνῳ τῷ καιρῷ εἶπε Κύριος πρὸς μὲ, λάξευσον σεαυτῷ δύο πλάκας λιθίνας ὥσπερ τὰς πρώτας, καὶ ἀνάβηθι πρὸς μὲ εἰς τὸ ὄρος, καὶ ποιήσεις σεαυτῷ κιβωτὸν ξυλίνην.
\vs{2}Καὶ γράψεις ἐπὶ τὰς πλάκας τὰ ῥήματα, ἃ ἦν ἐν ταῖς πλαξὶν ταῖς πρώταις ἃς συνέτριψας, καὶ ἐμβαλεῖς αὐτὰ εἰς τὴν κιβωτόν.
\vs{3}Καὶ ἐποίησα κιβωτὸν ἐκ ξύλων ἀσήπτων, καὶ ἐλάξευσα τὰς πλάκας λιθίνας ὡς αἱ πρῶται, καὶ ἀνέβην εἰς τὸ ὄρος καὶ αἱ δύο πλάκες ἐπὶ ταῖς χερσί μου.
\vs{4}Καὶ ἔγραψεν ἐπὶ τὰς πλάκας κατὰ τὴν γραφὴν τὴν πρώτην τοὺς δέκα λόγους, οὓς ἐλάλησε Κύριος πρὸς ὑμᾶς ἐν τῷ ὄρει ἐκ μέσου τοῦ πυρός, καὶ ἔδωκεν αὐτὰς Κύριος ἐμοί.
\vs{5}Καὶ ἐπιστρέψας κατέβην ἐκ τοῦ ὄρους, καὶ ἐνέβαλον τὰς πλάκας εἰς τὴν κιβωτὸν ἣν ἐποίησα· καὶ ἦσαν ἐκεῖ, καθὰ ἐνετείλατό μοι Κύριος.
\vs{6}Καὶ οἱ υἱοὶ Ἰσραὴλ ἀπῇραν ἐκ Βηρὼθ υἱῶν Ἰακεὶμ Μισαδαΐ· ἐκεῖ ἀπέθανεν Ἀαρὼν, καὶ ἐτάφη ἐκεῖ, καὶ ἱεράτευσεν Ἐλεάζὰρ υἱὸς αὐτοῦ ἀντʼ αὐτοῦ.
\vs{7}Ἐκεῖθεν ἀπῇραν εἰς Γαδγάδ· καὶ ἀπὸ Γαδγὰδ εἰς Ἐτεβαθᾶ, γῆ χείμαῤῥοι ὑδάτων.

\vs{8}Ἐν ἐκείνῳ τῷ καιρῷ διέστειλε Κύριος τὴν φυλὴν τὴν Λευὶ, αἴρειν τὴν κιβωτὸν τῆς διαθήκης Κυρίου, παρεστάναι ἔναντι Κυρίου, λειτουργεῖν καὶ ἐπεύχεσθαι ἐπὶ τῷ ὀνόματι αὐτοῦ ἕως τῆς ἡμέρας ταύτης.
\vs{9}Διὰ τοῦτο οὐκ ἔστιν τοῖς Λευίταις μερὶς καὶ κλῆρος ἐν τοῖς ἀδελφοῖς αὐτῶν· Κύριος αὐτὸς κλῆρος αὐτοῦ, καθότι εἶπεν αὐτῷ.
\vs{10}Κᾀγὼ εἱστήκειν ἐν τῷ ὄρει τεσσαράκοντα ἡμέρας καὶ τεσσαράκοντα νύκτας. καὶ εἰσήκουσε Κύριος ἐμοῦ καὶ ἐν τῷ καιρῷ τούτῳ, καὶ οὐκ ἠθέλησε Κύριος ἐξολοθρεῦσαι ὑμᾶς.
\vs{11}Καὶ εἶπε Κύριος πρὸς μέ, Βάδιζε, ἄπαρον ἐναντίον τοῦ λαοῦ τούτου, καὶ εἰσπορευέσθωσαν καὶ κληρονομείτωσαν τὴν γῆν, ἣν ὤμοσα τοῖς πατράσιν αὐτῶν δοῦναι αὐτοῖς.

\vs{12}Καὶ νῦν, Ἰσραήλ, τί Κύριος ὁ Θεός σου αἰτεῖται παρὰ σοῦ, ἀλλʼ ἢ φοβεῖσθαι Κύριον τὸν Θεόν σου, καὶ πορεύεσθαι ἐν πάσαις ταῖς ὁδοῖς αὐτοῦ, καὶ ἀγαπᾷν αὐτόν, καὶ λατρεύειν Κυρίῳ τῷ Θεῷ σου ἐξ ὅλης τῆς καρδίας σου, καὶ ἐξ ὅλης τῆς ψυχῆς σου,
\vs{13}φυλάσσεσθαι τὰς ἐντολὰς Κυρίου τοῦ Θεοῦ σου, καὶ τὰ δικαιώματα αὐτοῦ, ὅσα ἐγὼ ἐντέλλομαί σοι σήμερον, ἵνα εὖ σοι ἠ·;
\vs{14}Ἰδοὺ Κυρίου τοῦ Θεοῦ σου ὁ οὐρανὸς καὶ ὁ οὐρανὸς τοῦ οὐρανοῦ, ἡ γῆ καὶ πάντα ὅσα ἐστὶν ἐν αὐτῇ.
\vs{15}Πλὴν τοὺς πατέρας ὑμῶν προείλετο Κύριος ἀγαπᾷν αὐτούς, καὶ ἐξελέξατο τὸ σπέρμα αὐτῶν μετʼ αὐτοὺς, ὑμᾶς, παρὰ πάντα τὰ ἔθνη, κατὰ τὴν ἡμέραν ταύτην.
\vs{16}Καὶ περιτεμεῖσθε τὴν σκληροκαρδίαν ὑμῶν, καὶ τὸν τράχηλον ὑμῶν οὐ σκληρυνεῖτε.
\vs{17}Ὁ γὰρ Κύριος ὁ Θεὸς ὑμῶν, οὗτος Θεὸς τῶν θεῶν, καὶ κύριος τῶν κυρίων, ὁ Θεὸς ὁ μέγας, καὶ ἰσχυρὸς, καὶ φοβερὸς, ὅστις οὐ θαυμάζει πρόσωπον, οὐδὲ οὐ μὴ λάβῃ δῶρον·
\vs{18}ποιῶν κρίσιν προσηλύτῳ καὶ ὀρφανῷ καὶ χήρᾳ, καὶ ἀγαπᾷ τὸν προσήλυτον δοῦναι αὐτῷ ἄρτον καὶ ἱμάτιον.
\vs{19}Καὶ ἀγαπήσετε τὸν προσήλυτον· προσήλυτοι γὰρ ἦτε ἐν γῇ Αἰγύπτῳ.

\vs{20}Κύριον τὸν Θεόν σου φοβηθήσῃ, καὶ αὐτῷ λατρεύσεις, καὶ πρὸς αὐτὸν κολληθήσῃ, καὶ ἐπὶ τῷ ὀνόματι αὐτοῦ ὀμῇ·
\vs{21}Οὗτος καύχημά σου, καὶ οὗτος Θεός σου, ὅστις ἐποίησεν ἐν σοὶ τὰ μεγάλα καὶ τὰ ἔνδοξα ταῦτα, ἃ ἴδοσαν οἱ ὀφθαλμοί σου.
\vs{22}Ἐν ἑβδομήκοντα ψυχαῖς κατέβησαν οἱ πατέρες σου εἰς Αἴγυπτον· νυνὶ δὲ ἐποίησέ σε Κύριος ὁ Θεός σου ὡσεὶ τὰ ἄστρα τοῦ οὐρανοῦ τῷ πλήθει.

\ch{11}
Καὶ ἀγαπήσεις Κύριον τὸν Θεόν σου, καὶ φυλάξῃ τὰ φυλάγματα αὐτοῦ, καὶ τὰ δικαιώματα αὐτοῦ, καὶ τὰς ἐντολὰς αὐτοῦ, καὶ τὰς κρίσεις αὐτοῦ πάσας τὰς ἡμέρας.
\vs{2}Καὶ γνώσεσθε σήμερον, ὅτι οὐχὶ τὰ παιδία ὑμῶν, ὅσοι οὐκ οἴδασιν οὐδὲ ἴδοσαν τὴν παιδείαν Κυρίου τοῦ Θεοῦ σου, καὶ τὰ μεγαλεῖα αὐτοῦ, καὶ τὴν χεῖρα τὴν κραταιὰν, καὶ τὸν βραχίονα τὸν ὑψηλόν,
\vs{3}καὶ τὰ σημεῖα αὐτοῦ, καὶ τὰ τέρατα αὐτοῦ, ὅσα ἐποίησεν ἐν μέσῳ Αἰγύπτου Φαραὼ βασιλεῖ Αἰγύπτου, καὶ πάσῃ τῇ γῇ αὐτοῦ,
\vs{4}καὶ ὅσα ἐποίησε τὴν δύναμιν τῶν Αἰγυπτίων, καὶ τὰ ἅρματα αὐτῶν, καὶ τὴν ἵππον αὐτῶν, καὶ τὴν δύναμιν αὐτῶν, ὡς ἐπέκλυσε τὸ ὕδωρ τῆς θαλάσσης τῆς ἐρυθρᾶς ἐπὶ προσώπου αὐτῶν καταδιωκόντων αὐτῶν ἐκ τῶν ὀπίσω ὑμῶν, καὶ ἀπώλεσεν αὐτοὺς Κύριος ἕως τῆς σήμερον ἡμέρας,
\vs{5}καὶ ὅσα ἐποίησεν ἡμῖν ἐν τῇ ἐρήμῳ ἕως ἤλθετε εἰς τὸν τόπον τοῦτον,
\vs{6}καὶ ὅσα ἐποίησε τῷ Δαθὰν καὶ Ἀβειρὼν υἱοῖς Ἑλιὰβ υἱοῦ Ῥουβὴν, οὓς ἀνοίξασα ἡ γῆ τὸ στόμα αὐτῆς κατέπιεν αὐτοὺς, καὶ τοὺς οἴκους αὐτῶν, καὶ τὰς σκηνὰς αὐτῶν, καὶ πᾶσαν αὐτῶν τὴν ὑπόστασιν τὴν μετʼ αὐτῶν ἐν μέσῳ παντὸς Ἰσραήλ·
\vs{7}Ὅτι οἱ ὀφθαλμοὶ ὑμῶν ἑώρακαν πάντα τὰ ἔργα Κυρίου τὰ μεγάλα, ὅσα ἐποίησεν ἐν ὑμῖν σήμερον.

\vs{8}Καὶ φυλάξεσθε πάσας τὰς ἐντολὰς αὐτοῦ ὅσας ἐγὼ ἐντέλλομαί σοι σήμερον, ἵνα ζῆτε, καὶ πολυπλασιασθῆτε, καὶ εἰσελθόντες κληρονομήσητε τὴν γῆν, εἰς ἣν ὑμεῖς διαβαίνετε τὸν Ἰορδάνην ἐκεῖ κληρονομῆσαι αὐτήν·
\vs{9}Ἵνα μακροημερεύσητε ἐπὶ τῆς γῆς, ἧς ὤμοσε Κύριος τοῖς πατράσιν ὑμῶν δοῦναι αὐτοῖς καὶ τῷ σπέρματι αὐτῶν μετʼ αὐτούς, γῆν ῥέουσαν γάλα καὶ μέλι.
\vs{10}Ἔστι γὰρ ἡ γῆ εἰς ἣν εἰσπορεύῃ ἐκεῖ κληρονομῆσαι αὐτήν, οὐχ ὥσπερ γῆ Αἰγύπτου ἐστίν, ὅθεν ἐκπεπόρευσθε ἐκεῖθεν, ὅταν σπείρωσι τὸν σπόρον, καὶ ποτίζωσι τοῖς ποσὶν αὐτῶν, ὡσεὶ κῆπον λαχανείας·
\vs{11}Ἡ δὲ γῆ εἰς ἣν εἰσπορεύῃ ἐκεῖ κληρονομῆσαι αὐτὴν, γῆ ὀρεινὴ καὶ πεδινὴ· ἐκ τοῦ ὑετοῦ τοῦ οὐρανοῦ πίεται ὕδωρ.
\vs{12}Γῆ, ἣν Κύριος ὁ Θεός σου ἐπισκοπεῖται αὐτὴν διαπαντός, οἱ ὀφθαλμοὶ Κυρίου τοῦ Θεοῦ σου ἐπʼ αὐτῆς ἀπʼ ἀρχῆς τοῦ ἐνιαυτοῦ καὶ ἕως συντελείας τοῦ ἐνιαυτοῦ.

\vs{13}Ἐὰν δὲ ἀκοῇ ἀκούσητε πάσας τὰς ἐντολὰς, ἃς ἐγὼ ἐντελλομαί σοι σήμερον, ἀγαπᾷν Κύριον τὸν Θεόν σου, καὶ λατρεύειν αὐτῷ ἐξ ὅλης τῆς καρδίας σου, καὶ ἐξ ὅλης τῆς ψυχῆς σου,
\vs{14}καὶ δώσει τὸν ὑετὸν τῇ γῇ σου καθʼ ὥραν πρώϊμον καὶ ὄψιμον, καὶ εἰσοίσεις τὸν σῖτόν σου, καὶ τὸν οἶνόν σου, καὶ τὸ ἔλαιόν σου,
\vs{15}καὶ δώσει χορτάσματα ἐν τοῖς ἀγροῖς σου τοῖς κτήνεσί σου· καὶ φαγὼν, καὶ ἐμπλησθεὶς,
\vs{16}πρόσεχε σεαυτῷ μὴ πλατυνθῇ ἡ καρδία σου, καὶ παραβῆτε, καὶ λατρεύσητε θεοῖς ἑτέροις, καὶ προσκυνήσητε αὐτοῖς,
\vs{17}καὶ θυμωθεὶς ὀργῇ Κύριος ἐφʼ ὑμῖν, καὶ συσχῇ τὸν οὐρανὸν, καὶ οὐκ ἔσται ὑετὸς, καὶ ἡ γῆ οὐ δώσει τὸν καρπὸν αὐτῆς, καὶ ἀπολεῖσθε ἐν τάχει ἀπὸ τῆς γῆς τῆς ἀγαθῆς, ἧς Κύριος ἔδωκεν ὑμῖν.

\vs{18}Καὶ ἐμβαλεῖτε τὰ ῥήματα ταῦτα εἰς τὴν καρδίαν ὑμῶν καὶ εἰς τὴν ψυχὴν ὑμῶν, καὶ ἀφάψετε αὐτὰ εἰς σημεῖον ἐπὶ τῆς χειρὸς ὑμῶν, καὶ ἔσται ἀσάλευτον πρὸ ὀφθαλμῶν ὑμῶν·
\vs{19}Καὶ διδάξετε αὐτὰ τὰ τέκνα ὑμῶν λαλεῖν ἐν αὐτοῖς καθημένου σου ἐν οἴκῳ, καὶ πορευομένου σου ἐν ὁδῷ, καὶ καθεύδοντός σου, καὶ διανισταμένου σου.
\vs{20}Καὶ γράψετε αὐτὰ ἐπὶ τὰς φλιὰς τῶν οἰκιῶν ὑμῶν, καὶ τῶν πυλῶν ὑμῶν,
\vs{21}ἵνα μακροημερεύσητε, καὶ αἱ ἡμέραι τῶν υἱῶν ὑμῶν ἐπὶ τῆς γῆς, ἧς ὤμοσε Κύριος τοῖς πατράσιν ὑμῶν δοῦναι αὐτοῖς, καθὼς αἱ ἡμέραι τοῦ οὐρανοῦ ἐπὶ τῆς γῆς.
\vs{22}Καὶ ἔσται ἐὰν ἀκοῇ ἀκούσητε πάσας τὰς ἐντολὰς ταύτας ἃς ἐγὼ ἐντέλλομαί σοι σήμερον ποιεῖν, ἀγαπᾷν Κύριον τὸν Θεὸν ἡμῶν, καὶ πορεύεσθαι ἐν πάσαις ταῖς ὁδοῖς αὐτοῦ, καὶ προσκολλᾶσθαι αὐτῷ,
\vs{23}καὶ ἐκβαλεῖ Κύριος πάντα τὰ ἔθνη ταῦτα ἀπὸ προσώπου ὑμῶν, καὶ κληρονομήσετε ἔθνη μεγάλα καὶ ἰσχυρὰ μᾶλλον ἢ ὑμεῖς.
\vs{24}Πάντα τὸν τόπον οὗ ἐὰν πατήσῃ τὸ ἴχνος τοῦ ποδὸς ὑμῶν, ὑμῖν ἔσται· ἀπὸ τῆς ἐρήμου καὶ Ἀντιλιβάνου, καὶ ἀπὸ τοῦ ποταμοῦ τοῦ μελάλου, ποταμοῦ Εὐφράτου, καὶ ἕως τῆς θαλάσσης τῆς ἐπὶ δυσμῶν ἔσται τὰ ὅριά σου.
\vs{25}Οὐκ ἀντιστήσεται οὐδεὶς κατὰ πρόσωπον ὑμῶν· καὶ τὸν φόβον ὑμῶν καὶ τὸν τρόμον ὑμῶν ἐπιθήσει Κύριος ὁ Θεὸς ὑμῶν ἐπὶ πρόσωπον πάσης τῆς γῆς, ἐφʼ ἧς ἂν ἐπιβῆτε ἐπʼ αὐτῆς, ὃν τρόπον ἐλάλησε πρὸς ὑμᾶς.

\vs{26}Ἰδοὺ ἐγὼ δίδωμι ἐνώπιον ὑμῶν σήμερον τὴν εὐλογίαν καὶ τὴν κατάραν·
\vs{27}Τὴν εὐλογίαν, ἐὰν ἀκούσητε τὰς ἐντολὰς Κυρίου τοῦ Θεοῦ ὑμῶν, ὅσας ἐγὼ ἐντέλλομαι ὑμῖν σήμερον·
\vs{28}καὶ τὴν κατάραν, ἐὰν μὴ ἀκούσητε τὰς ἐντολὰς Κυρίου τοῦ Θεοῦ ὑμῶν, ὅσα ἐγὼ ἐντέλλομαι ὑμῖν σήμερον, καὶ πλανηθῆτε ἀπὸ τῆς ὁδοῦ ἧς ἐνετειλάμην ὑμῖν, πορευθέντες λατρεύειν θεοῖς ἑτέροις, οὓς οὐκ οἴδατε.
\vs{29}Καὶ ἔσται ὅταν εἰσαγάγῃ σε Κύριος ὁ Θεός σου εἰς τὴν γῆν εἰς ἣν διαβαίνεις ἐκεῖ κληρονομῆσαι αὐτήν, καὶ δώσεις εὐλογίαν ἐπʼ ὄρος Γαριζὶν, καὶ τὴν κατάραν ἐπʼ ὄρος Γαιβάλ.
\vs{30}Οὐκ ἰδοὺ ταῦτα πέραν τοῦ Ἰορδάνου, ὀπίσω ὁδὸν δυσμῶν ἡλίου ἐν γῇ Χαναὰν, τὸ κατοικοῦν ἐπὶ δυσμῶν ἐχόμενον τοῦ Γολγὸλ πλησίον τῆς δρυὸς τῆς ὑψηλῆς;
\vs{31}Ὑμεῖς γὰρ διαβαίνετε τὸν Ἰορδάνην, εἰσελθόντες κληρονομῆσαι τὴν γῆν, ἣν Κύριος ὁ Θεὸς ἡμῶν δίδωσιν ὑμῖν ἐν κλήρῳ πάσας τὰς ἡμέρας, καὶ κατοικήσετε ἐν αὐτῇ.

\vs{32}Καὶ φυλάξεσθε τοῦ ποιεῖν πάντα τὰ προστάγματα αὐτοῦ, καὶ τὰς κρίσεις ταύτας, ὅσας ἐγὼ δίδωμι ἐνώπιον ὑμῶν σήμερον.

\ch{12}
Καὶ ταῦτα τὰ προστάγματα καὶ αἱ κρίσεις, ἃς φυλάξετε τοῦ ποιεῖν ἐν τῇ γῇ, ἣν Κύριος ὁ Θεὸς τῶν πατέρων ὑμῶν δίδωσιν ὑμῖν ἐν κλήρῳ, πάσας τὰς ἡμέρας, ἃς ὑμεῖς ζῆτε ἐπὶ τῆς γῆς.
\vs{2}Ἀπωλίᾳ ἀπολεῖτε πάντας τοὺς τόπους ἐν οἷς ἐλάτρευσαν ἐκεῖ τοῖς θεοῖς αὐτῶν, οὓς ὑμεῖς κληρονομεῖτε αὐτοὺς, ἐπὶ τῶν ὀρέων τῶν ὑψηλῶν, καὶ ἐπὶ τῶν θινῶν, καὶ ὑποκάτω δένδρου δασέος·
\vs{3}Καὶ κατασκάψετε τοὺς βωμοὺς αὐτῶν, καὶ συντρίψετε τὰς στήλας αὐτῶν, καὶ τὰ ἄλση αὐτῶν ἐκκόψετε, καὶ τὰ γλυπτὰ τῶν θεῶν αὐτῶν κατακαύσετε πυρί, καὶ ἀπολεῖτε τὸ ὄνομα αὐτῶν ἐκ τοῦ τόπου ἐκείνου.
\vs{4}Οὐ ποιήσετε οὕτω Κυρίῳ τῷ Θεῷ ὑμῶν.
\vs{5}Ἀλλʼ ἢ εἰς τὸν τόπον, ὃν ἂν ἐκλέξηται Κύριος ὁ Θεός σου ἐν μιᾷ τῶν πόλεων ὑμῶν ἐπονομάσαι τὸ ὄνομα αὐτοῦ ἐκεῖ καὶ ἐπικληθῆναι, καὶ ἐκζητήσετε καὶ ἐλεύσεσθε ἐκεῖ.
\vs{6}Καὶ οἴσετε ἐκεῖ τὰ ὁλοκαυτώματα ὑμων, καὶ τὰ θυσιάσματα ὑμῶν, καὶ τὰς ἀπαρχὰς ὑμῶν, καὶ τὰς εὐχὰς ὑμῶν, καὶ τὰ ἐκούσια ὑμῶν, καὶ τὰς ὁμολογίας ὑμῶν, τὰ πρωτότοκα τῶν βοῶν ὑμῶν, καὶ τῶν προβάτων ὑμῶν.
\vs{7}Καὶ φάγεσθε ἐκεῖ ἐναντίον Κυρίου τοῦ Θεοῦ ὑμῶν, καὶ εὐφρανθήσεσθε ἐπὶ πᾶσιν, οὗ ἐὰν ἐπιβάλητε τὴν χεῖρα ὑμεῖς, καὶ οἱ οἶκοι ὑμῶν, καθότι εὐλόγησέ σε Κύριος ὁ Θεός σου.

\vs{8}Οὐ ποιήσετε πάντα ὅσα ἡμεῖς ποιοῦμεν ὧδε σήμερον, ἕκαστος τὸ ἀρεστὸν ἐνώπιον αὐτοῦ·
\vs{9}Οὐ γὰρ ἥκατε ἕως τοῦ νῦν εἰς τὴν κατάπαυσιν, καὶ εἰς τὴν κληρονομίαν, ἣν Κύριος ὁ Θεὸς ἡμῶν δίδωσιν ὑμῖν.
\vs{10}Καὶ διαβήσεσθε τὸν Ἰορδάνην, καὶ κατοικήσετε ἐπὶ τῆς γῆς, ἧς Κύριος ὁ Θεὸς ἡμῶν κατακληρονομεῖ ὑμῖν, καὶ καταπαύσει ὑμᾶς ἀπὸ πάντων τῶν ἐχθρῶν ὑμῶν τῶν κύκλῳ, καὶ κατοικήσετε μετὰ ἀσφαλείας.
\vs{11}Καὶ ἔσται ὁ τόπος, ὃν ἂν ἐκλέξηται Κύριος ὁ Θεός σου ἐπικληθῆναι τὸ ὄνομα αὐτοῦ ἐκεῖ, ἐκεῖ οἴσετε πάντα ὅσα ἐγὼ ἐντέλλομαι ὑμῖν σήμερον· τὰ ὁλοκαυτώματα ὑμῶν, καὶ τὰ θυσιάσματα ὑμῶν, καὶ τὰ ἐπιδέκατα ὑμῶν, καὶ τὰς ἀπαρχὰς τῶν χειρῶν ὑμῶν, καὶ πᾶν ἐκλεκτὸν τῶν δώρων ὑμῶν, ὅσα ἂν εὔξησθε Κυρίῳ τῷ Θεῷ ὑμῶν.
\vs{12}Καὶ εὐφρανθήσεσθε ἐναντίον Κυρίου τοῦ Θεοῦ ὑμῶν, ὑμεῖς καὶ οἱ υἱοὶ ὑμῶν, καὶ αἱ θυγατέρες ὑμῶν, καὶ οἱ παῖδες ὑμῶν, καὶ αἱ παιδίσκαι ὑμῶν, καὶ ὁ Λευίτης ὁ ἐπὶ τῶν πυλῶν ὑμῶν· ὅτι οὐκ ἔστιν αὐτῷ μερὶς οὐδὲ κλῆρος μεθʼ ὑμῶν.
\vs{13}Πρόσεχε σεαυτῷ, μὴ ἀνενέγκῃς τὰ ὁλοκαυτώματά σου ἐν παντὶ τόπῳ οὗ ἐὰν ἴδῃς.
\vs{14}Ἀλλʼ ἢ εἰς τὸν τόπον, ὃν ἂν ἐκλέξεται Κύριος ὁ Θεός σου αὐτὸν, ἐν μιᾷ τῶν φυλῶν σου, ἐκεῖ ἀνοίσετε τὰ ὁλοκαυτώματα ὑμῶν, καὶ ἐκεῖ ποιήσεις πάντα ὅσα ἐγὼ ἐντέλλομαί σοι σήμερον.
\vs{15}Ἀλλʼ ἢ ἐν πάσῃ ἐπιθυμίᾳ σου θύσεις, καὶ φαγῇ κρέα κατὰ τὴν εὐλογίαν Κυρίου τοῦ Θεοῦ σου, ἣν ἔδωκέ σοι ἐν πάσῃ πόλει· ὁ ἀκάθαρτος ἐν σοὶ καὶ ὁ καθαρὸς ἐπὶ τὸ αὐτό φάγεται αὐτό, ὡς δορκάδα ἢ ἔλαφον.
\vs{16}Πλὴν τὸ αἷμα οὐ φάγεσθε· ἐπὶ τὴν γῆν ἐκχεεῖτε αὐτὸ, ὡς ὕδωρ.

\vs{17}Οὐ δυνήσῃ φαγεῖν ἐν ταῖς πόλεσί σου τὸ ἐπιδέκατον τοῦ σίτου σου, καὶ τοῦ οἴνου σου, καὶ τοῦ ἐλαίου σου, τὰ πρωτότοκα τῶν βοῶν σου, καὶ τῶν προβάτων σαυ, καὶ πάσας τὰς εὐχὰς, ὅσας ἂν εὔξησθε, καὶ τὰς ὁμολογίας ὑμῶν, καὶ τὰς ἀπαρχὰς τῶν χειρῶν σου.
\vs{18}Ἀλλʼ ἢ ἐναντίον Κυρίου τοῦ Θεοῦ σου φάγῃ αὐτὸ ἐν τῷ τόπῳ, ᾧ ἂν ἐκλέξηται Κύριος ὁ Θεός σου αὐτῷ, σὺ καὶ ὁ υἱός σου, καὶ ἡ θυγάτηρ σου, ὁ παῖς σου, καὶ ἡ παιδίσκη σου, καὶ ὁ προσήλυτος ὁ ἐν ταῖς πόλεσιν ὑμῶν· καὶ εὐφρανθήσῃ ἐναντίον Κυρίου τοῦ Θεοῦ σου ἐπὶ πάντα, οὗ ἐὰν ἐπιβάλῃς τὴν χεῖρά σου.

\vs{19}Πρόσεχε σεαυτῷ μὴ ἐγκαταλίπῃς τὸν Λευίτην πάντα τὸν χρόνον ὅσον ἂ ζῇς ἐπὶ τῆς γῆς.

\vs{20}Ἐὰν δὲ ἐμπλατύνῃ Κύριος ὁ Θεός σου τὰ ὅριά σου, καθάπερ ἐλάλησέ σοι, καὶ ἐρεῖς, φάγομαι κρέα, ἐὰν ἐπιθυμήσῃ ἡ ψυχή σου ὥστε φαγεῖν κρέα, ἐν πάσῃ ἐπιθυμίᾳ τῆς ψυχῆς σου φάγῃ κρέα.
\vs{21}Ἐὰν δὲ μακρὰν ἀπέχῃ σου ὁ τόπος, ὃν ἂν ἐκλέξηται Κύριος ὁ Θεός σου ἐκεῖ ἐπικληθῆναι τὸ ὄνομα αὐτοῦ ἐκεῖ, καὶ θύσεις ἀπὸ τῶν βοῶν σου, καὶ ἀπὸ τῶν προβάτων σου ὧν ἂν δῷ ὁ Θεός σοι, ὃν τρόπον ἐνετειλάμην σοι, καὶ φαγῃ ἐν ταῖς πόλεσί σου κατὰ τὴν ἐπιθυμίαν τῆς ψυχῆς σου.
\vs{22}Ὡς ἔσθεται ἡ δορκὰς καὶ ἡ ἔλαφος, οὕτω φαγῃ αὐτό· ὁ ἀκάθαρτος ἐν σοὶ καὶ ὁ καθαρὸς ὡσαύτως ἔδεται.
\vs{23}Πρόσεχε ἰσχυρῶς τοῦ μὴ φαγεῖν αἷμα, ὅτι αἷμα αὐτοῦ ψυχή· οὐ βρωθήσεται ψυχὴ μετὰ τῶν κρεῶν.
\vs{24}Οὐ φάγεσθε· ἐπὶ τὴν γῆν ἐκχεεῖτε αὐτὸ ὡς ὕδωρ.
\vs{25}Οὐ φαγῇ αὐτὸ, ἵνα εὖ σοι γένηται καὶ τοῖς υἱοῖς σου μετὰ σὲ, ἐὰν ποιήσῃς τὸ καλὸν καὶ τὸ ἀρεστὸν ἐναντίον Κυρίου τοῦ Θεοῦ σου.
\vs{26}Πλὴν τὰ ἅγιά σου ἐὰν γένηταί σοι, καὶ τὰς εὐχάς σου λαβὼν ἥξεις εἰς τὸν τόπον, ὃν ἂν ἐκλέξηται Κύριος ὁ Θεός σου ἐπικληθῆναι τὸ ὄνομα αὐτοῦ ἐκεῖ.
\vs{27}Καὶ ποιήσεις τὰ ὁλοκαυτώματά σου, τὰ κρέα ἀνοίσεις ἐπὶ τὸ θυσιαστήριον Κυρίου τοῦ Θεοῦ σου· τὸ δὲ αἷμα τῶν θυσιῶν σου προσχεεῖς πρὸς τὴν βάσιν τοῦ θυσιαστηρίου Κυρίου τοῦ Θεοῦ σου, τὰ δὲ κρέα φαγῇ.
\vs{28}Φυλάσσου καὶ ἄκουε καὶ ποιήσεις πάντας τοὺς λόγους οὓς ἐγὼ ἐντέλλομαί σοι, ἵνα εὖ σοι γένηται καὶ τοῖς υἱοῖς σου διʼ αἰῶνος, ἐὰν ποιήσῃς τὸ ἀρεστὸν καὶ τὸ καλὸν ἐναντίον Κυρίου τοῦ Θεοῦ σου.

\vs{29}Ἐὰν δὲ ἐξολοθρεύσῃ Κύριος ὁ Θεός σου τὰ ἔθνη, εἰς οὓς εἰσπορεύῃ ἐκεῖ κληρονομῆσαι τὴν γῆν αὐτῶν, ἀπὸ προσώπου σου, καὶ κατακληρονομήσῃς αὐτὴν, καὶ κατοικήσῃς ἐν τῇ γῇ αὐτῶν,
\vs{30}πρόσεχε σεαυτῷ μὴ ἐκζητήσῃς ἐπακολουθῆσαι αὐτοῖς μετὰ τὸ ἐξολοθρευθῆναι αὐτοὺς ἀπὸ προσώπου σου, λέγων, πῶς ποιοῦσι τὰ ἔθνη ταῦτα τοῖς θεοῖς αὐτῶν; ποιήσω κᾀγώ.
\vs{31}Οὐ ποιήσεις οὕτω τῷ Θεῷ σου· τὰ γὰρ βδελύματα Κυρίου ἃ ἐμίσησεν, ἐποίησαν ἐν τοῖς θεοῖς αὐτῶν, ὅτι τοὺς υἱοὺς αὐτῶν, καὶ τὰς θυγατέρας αὐτῶν κατακαίουσιν ἐν πυρὶ τοῖς θεοῖς αὐτῶν.

\ch{13}Πᾶν ῥῆμα ὃ ἐγὼ ἐντλλομαι ὑμῖν σήμερον, τοῦτο φυλάξῃ ποιεῖν· οὐ προσθήσεις ἐπʼ αὐτό, οὐδὲ ἀφελεῖς ἀπʼ αὐτοῦ.

\vs{2}Ἐὰν δὲ ἀναστῇ ἐν σοὶ προφήτης ἢ ἐνυπνιαζόμενος τὸ ἐνύπνιον, καὶ δῷ σοι σημεῖον ἢ τέρας,
\vs{3}καὶ ἔλθῃ τὸ σημεῖον ἢ τὸ τέρας ὃ ἐλάλησε πρὸς σὲ, λέγων, πορευθῶμεν καὶ λατρεύσωμεν θεοῖς ἑτέροις οὓς οὐκ οἴδατε,
\vs{4}οὐκ ἀκούσεσθε τῶν λόγων τοῦ προφήτου ἐκείνου ἢ τοῦ ἐνυπνιαζομένου τὸ ἐνύπνιον ἐκεῖνο· ὅτι πειράζει Κύριος ὁ Θεός σου ὑμᾶς, εἰδέναι εἰ ἀγαπᾶτε τὸν Θεὸν ὑμῶν ἐξ ὅλης τῆς καρδίας ὑμῶν, καὶ ἐξ ὅλης τῆς ψυχῆς ὑμῶν.
\vs{5}Ὀπίσω Κυρίου τοῦ Θεοῦ ὑμῶν πορεύεσθε, καὶ τοῦτον φοβηθήσεσθε, καὶ τῆς φωνῆς αὐτοῦ ἀκούσεσθε, καὶ αὐτῷ προστεθήσεσθε.
\vs{6}Καὶ ὁ προφήτης ἐκεῖνος ἢ ὁ τὸ ἐνύπνιον ἐνυπνιαζόμενος ἐκεῖνος, ἀποθανεῖται· ἐλάλησε γὰρ πλανῆσαί σε ἀπὸ Κυρίου τοῦ Θεοῦ σου ἐξαγαγόντος σε ἐκ γῆς Αἰγύπτου, τοῦ λυτρωσαμένου σε ἐκ τῆς δουλείας, ἐξῶσαί σε ἀπὸ τῆς ὁδοῦ, ἧς ἐνετείλατό σοι Κύριος ὁ Θεός σου πορεύεσθαι ἐν αὐτῇ· καὶ ἀφανιεῖς τὸν πονηρὸν ἐξ ὑμῶν αὐτῶν.

\vs{7}Ἐὰν δὲ παρακαλέσῃ σε ὁ ἀδελφός σου ἐκ πατρός σου ἢ ἐκ μητρός σου, ἢ ὁ υἱός σου, ἢ ἡ θυγάτηρ, ἢ ἡ γυνή σου ἡ ἐν κόλπῳ σου, ἢ φίλος ἴσος τῇ ψυχῇς σου λάθρα, λέγων, βαδίσωμεν καὶ λατρεύσωμεν θεοῖς ἑτέροις, οὓς οὐκ ᾔδεις σὺ καὶ οἱ πατέρες σου,
\vs{8}ἀπὸ τῶν θεῶν τῶν ἐθνῶν τῶν περὶ κύκλῳ ὑμῶν, τῶν ἐγγιζόντων σοι ἢ τῶν μακρὰν ἀπὸ σοῦ, ἀπʼ ἄκρου τῆς γῆς ἕως ἄκρου τῆς γῆς,
\vs{9}οὐ συνθελήσεις αὐτῷ, καὶ οὐκ εἰσακούσῃ αὐτοῦ, καὶ οὐ φείσεται ὁ ὀφθαλμός σου ἐπʼ αὐτῷ, οὐκ ἐπιποθήσεις ἐπʼ αὐτῷ, οὐδʼ οὐ μὴ σκεπάσῃς αὐτόν·
\vs{10}ἀναγγέλλων ἀναγγελεῖς περὶ αὐτοῦ, καὶ αἱ χεῖρές σου ἔσονται ἐπʼ αὐτὸν ἐν πρώτοις ἀποκτεῖναι αὐτὸν, καὶ αἱ χεῖρες παντὸς τοῦ λαοῦ ἐπʼ ἐσχάτῳ.
\vs{11}Καὶ λιθοβολήσουσιν αὐτὸν ἐν λίθοις, καὶ ἀποθανεῖται, ὅτι ἐζήτησεν ἀποστῆσαί σε ἀπὸ Κυρίου τοῦ Θεοῦ σου τοῦ ἐξαγαγόντος σε ἐκ γῆς Αἰγύπτου, ἐξ οἴκου δουλείας.
\vs{12}Καὶ πᾶς Ἰσραὴλ ἀκούσας φοβηθήσεται, καὶ οὐ προσθήσει ποιῆσαι ἔτι κατὰ τὸ ῥῆμα τὸ πονηρὸν τοῦτο ἐν ὑμῖν.

\vs{13}Ἐὰν δὲ ἀκούσῃς ἐν μιᾷ τῶν πόλεών σου, ὧν Κύριος ὁ Θεός σου δίδωσί σοι κατοικεῖν σε ἐκεῖ, λεγόντων,
\vs{14}ἐξήλθοσαν ἄνδρες παράνομοι ἐξ ὑμῶν, καὶ ἀπέστησαν πάντας τοὺς κατοικοῦντας τὴν γῆν αὐτῶν, λέγοντες, πορευθῶμεν καὶ λατρεύσωμεν θεοῖς ἑτέροις, οὓς οὐκ ᾔδειτε,
\vs{15}καὶ ἐτάσεις καὶ ἐρωτήσεις, καὶ ἐρευνήσεις σφόδρα, καὶ ἰδοὺ ἀληθὴς σαφῶς ὁ λόγος, γεγένηται τὸ βδέλυγμα τοῦτο ἐν ὑμῖν·
\vs{16}ἀναιρῶν ἀνελεῖς πάντας τοὺς κατοικοῦντας ἐν τῇ γῇ ἐκείνῃ ἐν φόνῳ μαχαίρας, ἀναθέματι ἀναθεματιεῖτε αὐτὴν, καὶ πάντα τὰ ἐν αὐτῇ.
\vs{17}Καὶ πάντα τὰ σκῦλα αὐτῆς συνάξεις εἰς τὰς διόδους αὐτῆς, καὶ ἐμπρήσεις τὴν πόλιν ἐν πυρὶ, καὶ πάντα τὰ σκῦλα αὐτῆς πανδημεὶ ἐναντίον Κυρίου τοῦ Θεοῦ σου· καὶ ἔσται ἀοίκητος εἰς τὸν αἰῶνα, οὐκ ἀνοικοδομηθήσεται ἔτι.
\vs{18}Καὶ οὐ προσκολληθήσεται οὐδὲν ἀπὸ τοῦ ἀναθέματος ἐν τῇ χειρί σου, ἵνα ἀποστραφῇ Κύριος ἀπὸ θυμοῦ τῆς ὀργῆς αὐτοῦ, καὶ δώσῃ σοι ἔλεος, καὶ ἐλεήσῇ σε, καὶ πληθύνῃ σε, ὃν τρόπον ὤμοσε τοῖς πατράσι σου,
\vs{19}ἐὰν ἀκούσῃς τῆς φωνῆς Κυρίου τοῦ Θεοῦ σου, φυλάσσειν τὰς ἐντολὰς αὐτοῦ, ὅσας ἐγὼ ἐντέλλομαί σοι σήμερον, ποιεῖν τὸ καλὸν καὶ τὸ ἀρεστὸν ἐναντίον Κυρίου τοῦ Θεοῦ σου.

\ch{14}
Υἱοί ἐστε Κυρίου τοῦ Θεοῦ ὑμῶν· οὐκ ἐπιθήσετε φαλάκρωμα ἀναμέσον τῶν ὀφθαλμῶν ὑμῶν ἐπὶ νεκρῷ.
\vs{2}Ὅτι λαὸς ἅγιος εἶ Κυρίῳ τῷ Θεῷ σου, καὶ σε ἐξελέξατο Κύριος ὁ Θεός σου γενέσθαι σε λαὸν αὐτῷ περιούσιον ἀπὸ πάντων τῶν ἐθνῶν τῶν ἐπὶ προσώπου τῆς γῆς.
\vs{3}Οὐ φάγεσθε πᾶν βδέλυγμα.
\vs{4}Ταῦτα κτήνη ἃ φάγεσθε· μόσχον ἐκ βοῶν, καὶ ἀμνὸν ἐκ προβάτων, καὶ χίμαρον ἐξ αἰγῶν·
\vs{5}ἔλαφον, καὶ δορκάδα, καὶ πύγαργον, ὄρυγα, καὶ καμηλοπάρδαλιν.
\vs{6}Πᾶν κτῆνος διχηλοῦν ὁπλὴν, καὶ ὀνυχιστῆρας ὀνυχίζον δύο χηλῶν, καὶ ἀνάγον μηρυκισμὸν ἐν τοῖς κτήνεσι, ταῦτα φάγεσθε.
\vs{7}Καὶ ταῦτα οὐ φάγεσθε ἀπὸ τῶν ἀναγόντων μηρυκισμὸν, καὶ ἀπὸ τῶν διχηλούντων τὰς ὁπλὰς, καὶ ὀνυχιζόντων ὀνυχιστῆρας· τὸν κάμηλον, καὶ δασύποδα, καὶ χοιρογρύλλιον· ὅτι ἀνάγουσι μηρυκισμὸν, καὶ ὁπλὴν οὐ διχηλοῦσιν, ἀκάθαρτα ταῦτα ὑμῖν ἐστι.
\vs{8}Καὶ τὸν ὗν, ὅτι διχηλεῖ ὁπλὴν τοῦτο, καὶ ὀνυχίζει ὀνυχιστῆρας ὁπλῆς, καὶ τοῦτο μηρυκισμὸν οὐ μηρυκᾶται ἀκάθαρτον τοῦτο ὑμῖν· ἀπὸ τῶν κρεῶν αὐτῶν οὐ φάγεσθε, τῶν θνησιμαίων αὐτῶν οὐχ ἅψεσθε.

\vs{9}Καὶ ταῦτα φάγεσθε ἀπὸ πάντων τῶν ἐν τῷ ὕδατι, πάντα ὅσα ἐστὶν ἐν αὐτοῖς πτερύγια καὶ λεπίδες, φάγεσθε.
\vs{10}Καὶ πάντα ὅσα οὐκ ἔστιν αὐτοῖς πτερύγια καὶ λεπίδες, οὐ φάγεσθε· ἀκάθαρτα ὑμῖν ἐστι.
\vs{11}Πᾶν ὄρνεον καθαρὸν φάγεσθε.
\vs{12}Καὶ ταῦτα οὐ φάγεσθε ἀπʼ αὐτῶν· τὸν ἀετὸν, καὶ τὸν γρύπα, καὶ τὸν ἁλιάετον,
\vs{13}καὶ τὸν γύπα, καὶ τὸν ἴκτινον, καὶ τὰ ὅμοια αὐτῷ,
\vs{15}καὶ στρουθὸν, καὶ γλαῦκα, καὶ λάρον,
\vs{16}καὶ ἐρωδιὸν, καὶ κύκνον, καὶ ἶβιν,
\vs{17}καὶ καταράκτην, καὶ ἱέρακα, καὶ τὰ ὅμοια αὐτῷ, καὶ ἔποπα, καὶ νυκτικόρακα,
\vs{18}καὶ πελακᾶνα, καὶ χαραδριὸν, καὶ τὰ ὅμοια αὐτῷ, καὶ πορφυρίωνα, καὶ νυκτερίδα.
\vs{19}Πάντα τὰ ἑρπετὰ τῶν πετεινῶν ἀκάθαρτά ἐστιν ὑμῖν· οὐ φάγεσθε ἀπʼ αὐτῶν.
\vs{20}Πᾶν πετεινὸν καθαρὸν φάγεσθε.
\vs{21}Πᾶν θνησιμαῖον οὐ φάγεσθε· τῷ παροίκῳ τῷ ἐν ταῖς πόλεσί σου δοθήσεται καὶ φάγεται, ἢ ἀποδώσῃ τῷ ἀλλοτρίῳ, ὅτι λαὸς ἅγιος εἶ Κυρίῳ τῷ Θεῷ σου. Οὐχ ἑψήσεις ἄρνα ἐν γάλακτι μητρὸς αὐτοῦ.

\vs{22}Δεκάτην ἀποδεκατώσεις παντὸς γεννήματος τοῦ σπέρματός σου, τὸ γέννημα τοῦ ἀγροῦ σου ἐνιαυτὸν κατʼ ἐνιαυτόν.
\vs{23}Καὶ φαγῇ αὐτὸ ἐν τῷ τόπῳ, ᾧ ἐὰν ἐκλέξηται Κύριος ὁ Θεός σου ἐπικληθῆναι τὸ ὄνομα αὐτοῦ ἐκεῖ· οἴσετε τὰ ἐπιδέκατα τοῦ σίτου σου, καὶ τοῦ οἴνου σου, καὶ τοῦ ἐλαίου σου, τὰ πρωτότοκα τῶν βοῶν σου, καὶ τῶν προβάτων σου, ἵνα μάθῃς φοβεῖσθαι Κύριον τὸν Θεόν σου πάσας τὰς ἡμέρας.
\vs{24}Ἐὰν δὲ μακρὰν γένηται ἡ ὁδὸς ἀπὸ σοῦ, καὶ μὴ δύνῃ ἀναφέρειν αὐτὰ, ὅτι μακρὰν ἀπὸ σοῦ ὁ τόπος, ὃν ἂν ἐκλέξηται Κύριος ὁ Θεός σου ἐπικληθῆναι τὸ ὄνομα αὐτοῦ ἐκεῖ, ὅτι εὐλογήσει σε Κύριος ὁ Θεός σοῦ,
\vs{25}καὶ ἀποδώσῃ αὐτὰ ἀργυρίου, καὶ λήψῃ τὸ ἀργύριον ἐν ταῖς χερσί σου, καὶ πορεύσῃ εἰς τὸν τόπον ὃν ἂν ἐκλέξηται Κύριος ὁ Θεός σου αὐτόν.
\vs{26}Καὶ δώσεις ἀργύριον ἐπὶ παντὸς οὗ ἂν ἐπιθυμεῇ ἡ ψυχή σου, ἐπὶ βουσὶν ἤ ἐπὶ προβάτοις, ἢ ἐτʼ οἴνῳ ἢ ἐπὶ σίκερα, ἢ ἐπὶ σίκερα, ἢ ἐπὶ παντὸς οὗ ἂν ἐπιθυμῇ ἡ φυχή σου, καὶ φαγῇ ἐκεῖ ἐναντίον Κυρίου τοῦ Θεοῦ σου, καὶ εὐφρανθήσῃ σὺ καὶ ὁ οἶκός σου,
\vs{27}καὶ ὁ Λευίτης ὁ ἐν ταῖς πόλεσί σου, ὅτι οὐκ ἔστιν αὐτῷ μερὶς οὐδὲ κλῆρος μετὰ σοῦ.

\vs{28}Μετὰ τρία ἔτη ἐξοίσεις πᾶν τὸ ἐπιδέκατον τῶν γεννημάτων σου, ἐν τῷ ἐνιαυτῷ ἐκείνῳ θήσεις αὐτὸ ἐν ταῖς πόλεσί σου.
\vs{29}Καὶ ἐλεύσεται ὁ Λευίτης, ὅτι οὐκ ἔστιν αὐτῷ μερὶς οὐδὲ κλῆρος μετὰ σοῦ, καὶ ὁ προσήλυτος καὶ ὁ ὀρφανὸς καὶ ἡ χήρα ἡ ἐν ταῖς πόλεσί σου, καὶ φάγονται καὶ ἐμπλησθήσονται, ἵνα εὐλογήσῃ σε Κύριος ὁ Θεός σου ἐν πᾶσι τοῖς ἔργοις οἷς ἐὰν ποιῇς.

\ch{15}
Διʼ ἑπτὰ ἐτῶν ποιήσεις ἄφεσιν.
\vs{2}Καὶ οὕτω τὸ πρόσταγμα τῆς ἀφέσεως· ἀφήσεις πᾶν χρέος ἴδιον, ὃ ὀφείλει σοι ὁ πλησίον, καὶ τὸν ἀδελφόν σου οὐκ ἀπαιτήσεις· ἐπικέκληται γὰρ ἄφεσις Κυρίῳ τῷ Θεῷ σου.
\vs{3}Τὸν ἀλλότριον ἀπαιτήσεις ὅσα ἐὰν ἠ· σοι παρʼ αὐτῷ, τῷ δὲ ἀδελφῷ σου ἄφεσιν ποιήσεις τοῦ χρέους σου.
\vs{4}Ὅτι οὐκ ἔσται ἐν σοὶ ἐνδεὴς, ὅτι εὐλογῶν εὐλογήσει σε Κύριος ὁ Θεός σου ἐν τῇ γῇ, ᾗ Κύριος ὁ Θεός σου δίδωσί σοι ἐν κλήρῳ κατακληρονομεῖν σε αὐτήν.

\vs{5}Ἐὰν δὲ ἀκοῇ εἰσακούσητε τῆς φωνῆς Κυρίου τοῦ Θεοῦ ὑμῶν φυλάσσειν καὶ ποιεῖν πάσας τὰς ἐντολὰς ταύτας, ὅσας ἐγὼ ἐντέλλομαί σοι σήμερον,
\vs{6}ὅτι Κύριος ὁ Θεός σου εὐλόγησέ σε ὃν τρόπον ἐλάλησέ σοι· καὶ δανειεῖς ἔθνεσι πολλοῖς, σὺ δὲ οὐ δανειῇ· καὶ ἄρξεις ἐθνῶν πολλῶν, σοῦ δὲ οὐκ ἄρξουσιν.

\vs{7}Ἐὰν δὲ γένηται ἐν σοὶ ἐνδεὴς ἐκ τῶν ἀδελφῶν σου ἐν μιᾷ τῶν πόλεών σου ἐν τῇ γῇ, ᾗ Κύριος ὁ Θεός σου δίδωσί σοι, οὐκ ἀποστέρξεις τὴν καρδίαν σου, οὐδʼ οὐ μὴ συσφίγξεις τὴν χεῖρά σου ἀπὸ τοῦ ἀδελφοῦ σου τοῦ ἐπιδεομένου.
\vs{8}Ἀνοίγων ἀνοίξεις τὰς χεῖράς σου αὐτῷ, καὶ δάνειον δανειεῖς αὐτῷ ὅσον ἐπιδέεται, καθότι ἐνδεεῖται.
\vs{9}Πρόσεχε σεαυτῷ μὴ γένηται ῥῆμα κρυπτὸν ἐν τῇ καρδίᾳ σου ἀνόμημα, λέγων, Ἐγγίζει τὸ ἔτος τὸ ἕβδομον, ἔτος τῆς ἀφέσεως, καὶ πονηρεύσηται ὁ ὀφθαλμός σου τῷ ἀδελφῷ σου τῷ ἐπιδεομένῳ, καὶ οὐ δώσεις αὐτῷ, καὶ καταβοήσεται κατὰ σοῦ πρὸς Κύριον, καὶ ἔσται ἐν σοὶ ἁμαρτία μεγάλη.
\vs{10}Διδοὺς δώσεις αὐτῷ, καὶ δάνειον δανειεῖς αὐτῷ ὅσον ἐπιδέεται, καθότι ἐνδεεῖται· καὶ οὐ λυπηθήσῃ τῇ καρδίᾳ σου διδόντος σου αὐτῷ, ὅτι διὰ τὸ ῥῆμα τοῦτο εὐλογήσει σε Κύριος ὁ Θεός σου ἐν πᾶσιν τοῖς ἔργοις καὶ ἐν πᾶσι οὗ ἂν ἐπιβάλῃς τὴν χεῖρά σου.
\vs{11}Οὐ γὰρ μὴ ἐκλίπῃ ἐνδεὴς ἀπὸ τῆς γῆς σου· διὰ τοῦτο ἐγώ σοι ἐντέλλομαι ποιεῖν τὸ ῥῆμα τοῦτο, λέγων, ἀνοίγων ἀνοίξεις τὰς χεῖράς σου τῷ ἀδελφῷ σου τῷ πένητι καὶ τῷ ἐπιδεομένῳ τῷ ἐπὶ τῆς γῆς σου.

\vs{12}Ἐὰν δὲ πραθῇ σοι ὁ ἀδελφός σου ὁ Ἑβραῖος ἢ Ἐβραία, δουλεύσει σοι ἓξ ἔτη, καὶ τῷ ἑβδόμῳ ἐξαποστελεῖς αὐτὸν ἐλεύθερον ἀπὸ σοῦ.
\vs{13}Ὅταν δὲ ἐξαποστέλλῃς αὐτὸν ἐλεύθερον ἀπὸ σοῦ, οὐκ ἐξαποστελεῖς αὐτὸν κενόν.
\vs{14}Ἐφόδιον ἐφοδιάσεις αὐτὸν ἀπὸ τῶν προβάτων σου, καὶ ἀπὸ τοῦ σίτου σου, καὶ ἀπὸ τοῦ οἴνου σου· καθὰ εὐλόγησέ σε Κύριος ὁ Θεός σου, δώσεις αὐτῷ.

\vs{15}Καὶ μνησθήσῃ ὅτι οἰκέτης ἦσθα ἐν γῇ Αἰγύπτου, καὶ ἐλυτρώσατό σε Κύριος ὁ Θεός σοῦ ἐκεῖθεν· διὰ τοῦτο ἐγώ σοι ἐντέλλομαι ποιεῖν τὸ ῥῆμα τοῦτο.
\vs{16}Ἐὰν δὲ λέγῃ πρὸς σὲ, οὐκ ἐξελεύσομαι ἀπὸ σοῦ, ὅτι ἠγάπηκέ σε καὶ τὴν οἰκίαν σου, ὅτι εὖ ἐστιν αὐτῷ παρὰ σοί.
\vs{17}Καὶ λήψῃ τὸ ὀπήτιον, καὶ τρυπήσεις τὸ ὠτίον αὐτοῦ πρὸς τὴν θύραν, καὶ ἔσται σοι οἰκέτης εἰς τὸν αἰῶνα· καὶ τὴν παιδίσκην σου ὡσαύτως ποιήσεις.
\vs{18}Οὐ σκληρὸν ἔσται ἐναντίον σου ἐξαποστελλομένων αὐτῶν ἐλευθέρων ἀπὸ σου, ὅτι ἐπέτειον μισθὸν τοῦ μισθωτοῦ ἐδούλευσέ σοι ἓξ ἔτη· καὶ εὐλογήσει σε Κύριος ὁ Θεός σου ἐν πᾶσιν οἷς ἐὰν ποιῇς.

\vs{19}Πᾶν πρωτότοκον ὃ ἐὰν τεχθῇ ἐν ταῖς βουσί σου, καὶ ἐν τοῖς προβάτοις σου, τὰ ἀρσενικὰ ἁγιάσεις Κυρίῳ τῷ Θεῷ σου· οὐκ ἐργᾷ ἐν τῷ πρωτοτόκῳ μόσχῳ σου, καὶ οὐ μὴ κείρῃς τὰ πρωτότοκα τῶν προβάτων σου.
\vs{20}Ἔναντι Κυρίου φαγῇ αὐτὸ ἐνιαυτὸν ἐξ ἐνιαυτοῦ ἐν τῷ τόπῳ ᾧ ἐὰν ἐκλέξηται Κύριος ὁ Θεός σου, σὺ καὶ ὁ οἶκός σου.
\vs{21}Ἐὰν δὲ ἠ· ἐν αὐτῷ μῶμος, χωλὸν ἢ τυφλον, μῶμον πονηρὸν, οὐ θύσεις αὐτὸ Κυρίῳ τῷ Θεῷ σου.

\vs{22}Ἐν ταῖς πόλεσί σου φαγῇ αὐτό· ὁ ἀκάθαρτος ἐν σοι, καὶ ὁ καθαρὸς ὡσαύτως ἔδεται ὡς δορκάδα ἢ ἔλαφον.
\vs{23}Πλὴν αἷμα οὐ φἀγεσθε· ἐπὶ τὴν γῆν ἐκχεεῖς αὐτὸ ὡς ὕδωρ.

\ch{16}
Φυλάξαι τὸν μῆνα τῶν νέων, καὶ ποιήσεις τὸ πάσχα Κυρίῳ τῷ Θεῷ σου, ὅτι ἐν τῷ μηνὶ τῶν νέων ἐξῆλθες ἐξ Αἰγύπτου νυκτός.
\vs{2}Καὶ θύσεις τὸ πάσχα Κυρίῳ τῷ Θεῷ σου πρόβατα καὶ βόας ἐν τῷ τόπῳ, ᾧ ἐὰν ἐκλέξηται Κύριος ὁ Θεός σου αὐτὸν, ἐπικληθῆναι τὸ ὄνομα αὐτοῦ ἐκεῖ.
\vs{3}Οὐ φαγῇ ἐπʼ αὐτοῦ ζύμην· ἐπτὰ ἡμέρας φαγῇ ἐπʼ αὐτοῦ ἄζυμα, ἄρτον κακώσεως, ὅτι ἐν σπουδῇ ἐξήλθετε ἐξ Αἰγύπτου, ἵνα μνησθῆτε τὴν ἡμέραν τῆς ἐξοδίας ὑμῶν ἐκ γῆς Αἰγύπτου πάσας τὰς ἡμέρας τῆς ζωῆς ὑμῶν.
\vs{4}Οὐκ ὀφθήσεται σοι ζύμη ἐν πᾶσι τοῖς ὁρίοις σου ἑπτὰ ἡμέρας, καὶ οὐ κοιμήθησεται ἀπὸ τῶν κρεῶν ὧν ἐὰν θύσῃς τὸ ἑσπέρας τῇ ἡμέρᾳ τῇ πρώτῃ εἰς τοπρωΐ.
\vs{5}Οὐ δυνήσῃ θῦσαι τὸ πάσχα ἐν οὐδεμιᾷ τῶν πόλεών σου, ὧν Κύριος ὁ Θεός σου δίδωσί σοι·
\vs{6}ἀλλʼ ἢ εἰς τὸν τόπον, ὃν ἂν ἐκλέξηται Κύριος ὁ Θεός σου, ἐπικληθῆναι τὸ ὄνομα αὐτοῦ ἐκεῖ, θύσεις τὸ πάσχα ἑσπέρας πρὸς δυσμὰς ἡλίου, ἐν τῷ καιρῷ ᾧ ἐξῆλθες ἐξ Αἰγύπτου.
\vs{7}Καὶ ἑψήσεις καὶ ὀπτήσεις καὶ φαγῇ ἐν τῷ τόπῳ, οὗ ἐὰν ἐκλέξηται Κύριος ὁ Θεός σου αὐτόν· καὶ ἀποστραφήσῃ τοπρωΐ, καὶ ἐλεύσῃ εἰς τοὺς οἴκους σου.
\vs{8}Ἓξ ἡμέρας φαγῇ ἄζυμα, καὶ τῇ ἡμέρᾳ τῇ ἑβδόμῃ ἐξόδιον ἑορτὴ Κυρίῳ τῷ Θεῷ σου· οὐ ποιήσεις ἐν αὐτῇ πᾶν ἔργον, πλὴν ὅσα ποιηθήσεται ψυχῇ.

\vs{9}Ἑπτὰ ἑβδομάδας ἐξαριθμήσεις σεαυτῷ· ἀρξαμένου σου δρέπανον ἐπʼ ἀμητὸν, ἄρξῃ ἐξαριθμῆσαι ἑπτὰ ἑβδομάδας.
\vs{10}Καὶ ποιήσεις ἑορτὴν ἑβδομάδων Κυρίῳ τῷ Θεῷ σου, καθὼς ἡ χείρ σου ἰσχύει, ὅσα ἂν δῷ Κύριος ὁ Θεός σου.

\vs{11}Καὶ εὐφρανθήσῃ ἐναντίον Κυρίου τοῦ Θεοῦ σου, σὺ καὶ ὁ υἱός σου, καὶ ἡ θυγάτηρ σου, ὁ παῖς σου, καὶ ἡ παιδίσκη σου, καὶ ὁ Λευίτης, καὶ ὁ προσήλυτος, καὶ ὁ ὀρφανὸς, καὶ ἡ χήρα ἡ οὖσα ἐν ὑμῖν, ἐν τῷ τόπῳ, ᾧ ἐὰν ἐκλέξηται Κύριος ὁ Θεός σου αὐτὸν, ἐπικληθῆναι τὸ ὄνομα αὐτοῦ ἐκεῖ.

\vs{12}Καὶ μνησθήσῃ ὅτι οἰκέτης ἐγένου ἐν γῇ Αἰγύπτῳ, καὶ φυλάξῃ καὶ ποιήσεις τὰς ἐντολὰς ταύτας.
\vs{13}Ἐορτὴν σκηνῶν ποιήσεις σεαυτῷ ἑπτὰ ἡμέρας ἐν τῷ συναγαγεῖν σε ἐκ τῆς ἅλωνός σου καὶ ἀπὸ τῆς ληνοῦ σου.
\vs{14}Καὶ εὐφρανθήσῃ ἐν τῇ ἑορτῇ σου, σὺ καὶ ὁ υἱός σου, καὶ ἡ θυγάτηρ σου, ὁ παῖς σου, καὶ ἡ παιδίσκη σου, καὶ ὁ Λευίτης, καὶ ὁ προσήλυτος, καὶ ὁ ὀρφανὸς, καὶ ἡ χήρα ἡ οὖσα ἐν ταῖς πόλεσί σου.
\vs{15}Ἑπτὰ ἡμέρας ἑορτάσεις Κυρἱω τῷ Θεῷ σου ἐν τῷ τόπῳ, ᾧ ἂν ἐκλέξηται Κύριος ὁ Θεός σου αὐτῷ· ἐὰν δὲ εὐλογήσῃ σε Κύριος ὁ Θεός σου ἐν πᾶσι γεννήμασί σου, καὶ ἐν πάντι ἔργῳ τῶν χειρῶν σου, καὶ ἔσῃ εὐφραινόμενος.

\vs{16}Τρεῖς καιροὺς τοῦ ἐνιαυτοῦ ὀφθήσεται πᾶν ἀρσενικόν σου ἐναντίον Κυρίου τοῦ Θεοῦ σου ἐν τῷ τόπῳ, ᾧ ἐὰν ἐκλέξηται αὐτὸν Κύριος· ἐν τῇ ἑορτῇ τῶν ἀζύμων, καὶ ἐν τῇ ἑορτῇ τῶν ἑβδομάδων, καὶ ἐν τῇ ἑορτῇ τῆς σκηνοπηγίας· οὐκ ὀφθήσῃ ἐνώπιον Κυρίου τοῦ Θεοῦ σου κενός.
\vs{17}Ἕκαστος κατὰ δύναμιν τῶν χειρῶν ὑμῶν, κατὰ τὴν εὐλογίαν Κυρίου τοῦ Θεοῦ σου ἣν ἔδωκέ σοι.

\vs{18}Κριτὰς καὶ γραμματοεισαγωγεῖς ποιήσεις σεαυτῷ ἐν ταῖς πόλεσί σου, αἷς Κύριος ὁ Θεός σου δίδωσί σοι κατὰ φυλάς· καὶ κρινοῦσι τὸν λαὸν κρίσιν δικαίαν.
\vs{19}Οὐκ ἐκκλινοῦσι κρίσιν, οὐδὲ λήψονται δῶρον· τὰ γὰρ δῶρα ἀποτυφλοῖ ὀφθαλμοὺς σοφῶν, καὶ ἐξαίρει λόγους δικαίων.
\vs{20}Δικαίως τὸ δίκαιον διώξῃ, ἵνα ζῆτε, καὶ εἰσελθόντες κληρονομήσητε τὴν γῆν ἣν Κύριος ὁ Θεός σου δίδωσί σοι.

\vs{21}Οὐ φυτεύσεις σεαυτῷ ἄλσος· πᾶν ξύλον παρὰ τὸ θυσιαστήριον τοῦ Θεοῦ σου οὐ ποιήσεις σεαυτῷ.
\vs{22}Οὐ στήσεις σεαυτῷ στήλην, ἃ ἐμίσησε Κύριος ὁ Θεός σου.

\ch{17}
Οὐ θύσεις Κυρίῳ τῷ Θεῷ σου μόσχον ἢ πρόβατον, ἐν ᾧ ἐστιν ἐν αὐτῷ μῶμος, πᾶν ῥῆμα πονηρόν· ὅτι βδέλυγμα Κυρίῳ τῷ Θεῷ σου ἐστίν.

\vs{2}Ἐὰν δὲ εὑρεθῇ ἐν μιᾷ τῶν πόλεών σου, ὧν Κύριος ὁ Θεός σου δίδωσί σοι, ἀνὴρ ἢ γυνὴ ὃς ποιήσει τὸ πονηρὸν ἐναντίον Κυρίου τοῦ Θεοῦ σου, παρελθεῖν τὴν διαθήκην αὐτοῦ,
\vs{3}καὶ ἐλθόντες λατρεύσωσι θεοῖς ἑτέροις, καὶ προσκυνήσωσιν αὐτοῖς, τῷ ἡλίῳ, ἢ τῇ σελήνῃ, ἢ παντὶ τῶν ἐκ τοῦ κόσμου τοῦ οὐρανοῦ, ἂ οὐ προσέταξέ σοι,
\vs{4}καὶ ἀναγγελῇ σοι καὶ ἐκζητήσῃς σφόδρα, καὶ ἰδοὺ ἀληθῶς γέγονε τὸ ῥῆμα, γεγένηται τὸ βδέλυγμα τοῦτο ἐν Ἰσραήλ·
\vs{5}Καὶ ἐξάξεις τὸν ἄνθρωπον ἐκεῖνον, ἢ τὴν γυναῖκα ἐκείνην, καὶ λιθοβολήσετε αὐτοὺς ἐν λίθοις, καὶ τελευτήσουσιν.
\vs{6}Ἐπὶ δυσὶ μάρτυσιν ἢ ἐπὶ τρισὶ μάρτυσιν ἀποθανεῖται· ὁ ἀποθνήσκων οὐκ ἀποθανεῖται ἐφʼ ἑνὶ μάρτυρι.
\vs{7}Καὶ ἡ χεὶρ τῶν μαρτύρων ἔσται ἐπʼ αὐτῷ ἐν πρώτοις θανατῶσαι αὐτὸν, καὶ ἡ χεὶρ τοῦ λαοῦ ἐπʼ ἐσχάτων· καὶ ἐξαρεῖς τὸν πονηρὸν ἐξ ὑμῶν αὐτῶν.

\vs{8}Ἐὰν δὲ ἀδυνατήσῃ ἀπὸ σοῦ ῥῆμα ἐν κρίσει ἀναμέσον αἷμα αἵματος, καὶ ἀναμέσον κρίσις κρίσεως, καὶ ἀναμέσον ἁφὴ ἁφῆς, καὶ ἀναμέσον ἀντιλογία ἀντιλογίας, ῥήματα κρίσεως ἐν ταῖς πόλεσιν ὑμῶν, καὶ ἀναστὰς ἀναβήσῃ εἰς τὸν τόπον ὃν ἂν ἐκλέξηται Κύριος ὁ Θεός σου ἐκεῖ,
\vs{9}καὶ ἐγεύσῃ πρὸς τὸνς ἱερεῖς τοὺς Λευίτας, καὶ πρὸς τὸν κριτὴν ὃς ἂν γένηται ἐν ταῖς ἡμέραις ἐκείναις, καὶ ἐκζητήσαντες ἀναγγελοῦσί σοι τὴν κρίσιν.
\vs{10}Καὶ ποιήσεις κατὰ τὸ πρᾶγμα ὃ ἂν ἀναγγείλωσί σοι ἐκ τοῦ τόπου, οὗ ἐὰν ἐκλέξηται Κύριος ὁ Θεός σου, καὶ φυλάξῃ ποιῆσαι πάντα ὅσα ἂν νομοθετηθῇ σοι.
\vs{11}Κατὰ τὸν νόμον καὶ κατὰ τὴν κρίσιν ἣν ἂν εἴπωσί σοι, ποιήσεις· οὐκ ἐκκλινεῖς ἀπὸ τοῦ ῥήματος οὗ ἐὰν ἀναγγείλωσί σοι δεξιὰ οὐδὲ ἀριστερά.

\vs{12}Καὶ ὁ ἄνθρωπος ὃς ἐὰν ποιήσῃ ἐν ὑπερηφανίᾳ, ὥστε μὴ ὑπακοῦσαι τοῦ ἱερέως τοῦ παρεστηκότος λειτουργεῖν ἐπὶ τῷ ὀνόματι Κυρίου τοῦ Θεοῦ σου, ἢ τοῦ κριτοῦ ὃς ἂν ἠ· ἐν ταῖς ἡμέραις ἐκείναις, καὶ ἀποθανεῖται ὁ ἄνθρωπος ἐκεῖνος, καὶ ἐξαρεῖς τὸν πονηρὸν ἐξ Ἰσραήλ.
\vs{13}Καὶ πᾶς ὁ λαὸς ἀκούσας φοβηθήσεται, καὶ οὐκ ἀσεβήσει ἔτι.

\vs{14}Ἐὰν δὲ εἰσέλθῃς εἰς τὴν γῆν ἣν Κύριος ὁ Θεός σου δίδωσί σοι, καὶ κληρονομήσῃς αὐτὴν, καὶ κατοικήσῃς ἐπʼ αὐτὴν, καὶ εἴπῃς, καταστήσω ἐπʼ ἐμαυτὸν ἄρχοντα, καθὰ καὶ τὰ λοιπὰ ἔθνη τὰ κύκλῳ μου·
\vs{15}καθιστῶν καταστήσεις ἐπὶ σεαυτὸν ἄρχοντα, ὃν ἂν ἐκλέξηται Κύριος ὁ Θεὸς αὐτόν· ἐκ τῶν ἀδελφῶν σου καταστήσεις ἐπὶ σεαυτὸν ἄρχοντα· οὐ δυνήσῃ καταστῆσαι ἐπὶ σεαυτὸν ἄνθρωπον ἀλλότριον, ὅτι οὐκ ἀδελφός σου ἐστί.
\vs{16}Διότι οὐ πληθυνεῖ ἑαυτῷ ἵππον, οὐδὲ μὴ ἀποστρέψῃ τὸν λαὸν εἰς Αἴγυπτον, ὅπως μὴ πληθύνῃ αὑτῷ ἵππον· ὁ δὲ Κύριος εἶπεν, οὐ προσθήσεσθε ἀποστρέψαι τῇ ὁδῷ ταύτῃ ἔτι.
\vs{17}Καὶ οὐ πληθυνεῖ ἑαυτῷ γυναῖκας, ἵνα μὴ μεταστῇ αὐτοῦ καρδία· καὶ ἀργύριον καὶ χρυσίον οὐ πληθυνεῖ ἑαυτῷ σφόδρα.

\vs{18}Καὶ ὅταν καθίσῃ ἐπὶ τῆς ἀρχῆς αὐτοῦ, καὶ γράψει αὑτῷ τὸ δευτερονόμιον τοῦτο εἰς βιβλίον παρὰ τῶν ἱερέων τῶν Λευιτῶν,
\vs{19}καὶ ἔσται μετʼ αὐτοῦ, καὶ ἀναγνώσεται ἐν αὐτῷ πάσας τὰς ἡμέρας τῆς ζωῆς αὐτοῦ, ἵνα μάθῃ φοβεῖσθαι Κύριον τὸν Θεόν σου, καὶ φυλάσσεσθαι πάσας τὰς ἐντολὰς ταύτας, καὶ τὰ δικαιώματα ταῦτα ποιεῖν·
\vs{20}ἵνα μὴ ὑψωθῇ ἡ καρδία αὐτοῦ ἀπὸ τῶν ἀδελφῶν αὐτοῦ, ἵνα μὴ παραβῇ ἀπὸ τῶν ἐντολῶν δεξιὰ ἢ ἀριστερὰ, ὅπως ἂν μακροχρονίσῃ ἐπὶ τῆς ἀρχῆς αὐτοῦ, αὐτὸς καὶ οἱ υἱοὶ αὐτοῦ ἐν τοῖς υἱοῖς Ἰσραήλ.

\ch{18}
Οὐκ ἔσται τοῖς ἱερεῦσι τοῖς Λευίταις ὅλῃ φυλῇ Λευὶ μερὶς οὐδὲ κλῆρος μετὰ Ἰσραήλ· καρπώματα Κυρίου ὁ κλῆρος αὐτῶν, φάγονται αὐτά.
\vs{2}Κλῆρος δὲ οὐκ ἔσται αὐτοῖς ἐν τοῖς ἀδελφοῖς αὐτῶν· Κύριος αὐτὸς κλῆρος αὐτοῦ, καθότι εἶπεν αὐτῷ.
\vs{3}Καὶ αὕτη ἡ κρίσις τῶν ἱερέων τὰ παρὰ τοῦ λαοῦ παρὰ τῶν θυόντων τὰ θύματα, ἐάν τε μόσχον, ἐάν τε πρόβατον· καὶ δώσεις τὸν βραχίονα τῷ ἱερεῖ, καὶ τὰ σιαγόνια, καὶ τὸ ἔνυστρον,
\vs{4}καὶ τὰς ἀπαρχὰς τοῦ σίτου σου, καὶ τοῦ οἴνου σου, καὶ τοῦ ἐλαίου σου· καὶ τὴν ἀπαρχὴν τῶν κουρῶν τῶν προβάτων σου δώσεις αὐτῷ.
\vs{5}Ὅτι αὐτὸν ἐξελέξατο Κύριος ἐκ πασῶν τῶν φυλῶν σου, παρεστάναι ἔναντι Κυρίου τοῦ Θεοῦ, λειτουργεῖν καὶ εὐλογεῖν ἐπὶ τῷ ὀνόματι αὐτοῦ, αὐτὸς καὶ οἱ υἱοὶ αὐτοῦ ἐν τοῖς υἱοῖς Ἰσραήλ.

\vs{6}Ἐὰν δὲ παραγένηται ὁ Λευίτης ἐκ μιᾶς τῶν πόλεων ἐκ πάντων τῶν υἱῶν Ἰσραὴλ, οὗ αὐτὸς παροικεῖ, καθʼ ὅτι ἐπιθυμεῖ ἡ ψυχὴ αὐτοῦ, εἰς τὸν τόπον ὃν ἂν ἐκλέξηται,
\vs{7}λειτουργήσει τῷ ὀνόματι Κυρίου τοῦ Θεοῦ αὐτοῦ, ὥσπερ πάντες οἱ ἀδελφοὶ αὐτοῦ οἱ Λευῖται οἱ παρεστηκότες ἐκεῖ ἐναντίον Κυρίου τοῦ Θεοῦ σου.
\vs{8}Μερίδα μεμερισμένην φάγεται, πλὴν τῆς πράσεως τῆς κατὰ πατριάν.
\vs{9}Ἐὰν δὲ εἰσέλθῃς εἰς τὴν γῆν ἣν Κύριος ὁ Θεός σου δίδωσί σοι, οὐ μαθήσῃ ποιεῖν κατὰ τὰ βδελύγματα τῶν ἐθνῶν ἐκείνων.

\vs{10}Οὐχ εὑρεθήσεται ἐν σοὶ περικαθαίρων τὸν υἱὸν αὐτοῦ καὶ τὴν θυγατέρα αὐτοῦ ἐν πυρὶ, μαντευόμενος μαντείαν, κληδονιζόμενος, καὶ οἰωνιζόμενος, φαρμακοῖς
\vs{11}ἐπαείδων ἐπαοιδὴν, ἐγγαστρίμυθος, καὶ τερατοσκόπος, ἐπερωτῶν τοὺς νεκρούς.
\vs{12}Ἔστι γὰρ βδέλυγμα Κυρίῳ τῷ Θεῷ σου πᾶς ποιῶν ταῦτα· ἕνεκεν γὰρ τῶν βδελυγμάτων τούτων Κύριος ἐξολοθρεύσει αὐτοὺς ἀπὸ προσώπου σου.
\vs{13}Τέλειος ἔσῃ ἐναντίον Κυρίου τοῦ Θεοῦ σου.
\vs{14}Τὰ γὰρ ἔθνη ταῦτα, οὓς σὺ κατακληρονομεῖς αὐτοὺς, οὗτοι κληδόνων καὶ μαντειῶν ἀκούσονται· καί σοι οὐχ οὕτως ἔδωκε Κύριος ὁ Θεός σου.

\vs{15}Προφήτην ἐκ τῶν ἀδελφῶν σου, ὡς ἐμὲ, ἀναστήσει σοι Κύριος ὁ Θεός σου· αὐτοῦ ἀκούσεσθε·
\vs{16}Κατὰ πάντα ὅσα ᾐτήσω παρὰ Κυρίου τοῦ Θεοῦ σου ἐν Χωρὴβ τῇ ἡμέρᾳ τῆς ἐκκλησίας, λέγοντες, οὐ προσθήσομεν ἀκοῦσαι τὴν φωνὴν Κυρίου τοῦ Θεοῦ σου, καὶ τὸ πῦρ τοῦτο τὸ μὲγα οὐκ ὀψόμεθα ἔτι, οὐδὲ μὴ ἀποθάνωμεν.
\vs{17}Καὶ εἶπε Κύριος πρὸς μὲ, ὀρθῶς πάντα ὅσα ἐλάλησαν πρὸς σέ.
\vs{18}Προφήτην ἀναστήσω αὐτοῖς ἐκ τῶν ἀδελφῶν αὐτῶν, ὥσπερ σέ· καὶ δώσω τὰ ῥήματα ἐν τῷ στόματι αὐτοῦ, καὶ λαλήσει αὐτοῖς καθʼ ὅτι ἂν ἐντείλωμαι αὐτῷ.
\vs{19}Καὶ ὁ ἄνθρωπος ὃς ἐὰν μὴ ἀκούσῃ ὅσα ἂν λαλήσῃ ὁ προφήτης ἐκεῖνος ἐπὶ τῷ ὀνόματί μου, ἐγὼ ἐκδικήσω ἐξ αὐτοῦ.
\vs{20}Πλὴν ὁ προφήτης ὃς ἂν ἀσεβήσῃ λαλῆσαι ἐπὶ τῷ ὀνόματί μου ῥῆμα ὃ οὐ προσέταξα λαλῆσαι, καὶ ὃς ἂν λαλήσῃ ἐν ὀνόματι θεῶν ἑτέρων, ἀποθανεῖται ὁ προφήτης ἐκεῖνος.
\vs{21}Ἐὰν δὲ εἴπῃς ἐν τῇ καρδίᾳ σου, πῶς γνωσόμεθα τὸ ῥῆμα ὃ οὐκ ἐλάλησε Κύριος;
\vs{22}Ὅσα ἐὰν λαλήσῃ ὁ προφήτης ἐκεῖνος τῷ ὀνόματι Κυρίου, καὶ μὴ γένηται, καὶ μὴ συμβῇ, τοῦτο τὸ ῥῆμα ὃ οὐκ ἐλάλησε Κύριος, ἐν ἀσεβείᾳ ἐλάλησεν ὁ προφήτης ἐκεῖνος· οὐκ ἀφέξεσθε αὐτοῦ.

\ch{19}
Ἐὰν δὲ ἀφανίσῃ Κύριος ὁ Θεός σου τὰ ἔθνη, ἃ ὁ Θεὸς δίδωσί σοι τὴν γῆν, καὶ κατακληρονομήσητε αὐτοὺς, καὶ κατοικήσητε ἐν ταῖς πόλεσιν αὐτῶν, καὶ ἐν τοῖς οἴκοις αὐτῶν,
\vs{2}τρεῖς πόλεις διαστελεῖς σεαυτῷ ἐν μέσῳ τῆς γῆς σου, ἧς Κύριος ὁ Θεός σου δίδωσί σοι.
\vs{3}Στόχασαί σοι τὴν ὁδὸν, καὶ τριμεριεῖς τὰ ὅρια τῆς γῆς σου, ἣν καταμερίζει σοι Κύριος ὁ Θεός σου, καὶ ἔσται καταφυγὴ παντὶ φονευτῇ.
\vs{4}Τοῦτο δὲ ἔσται τὸ πρόσταγμα τοῦ φονευτοῦ, ὃς ἂν φύγῃ ἐκεῖ, καὶ ζήσεται, ὃς ἂν πατάξῃ τὸν πλησίον αὐτοῦ οὐκ εἰδῶς, καὶ οὗτος οὐ μισῶν αὐτὸν πρὸ τῆς χθὲς καὶ τρίτης.
\vs{5}Καὶ ὅς ἐὰν εἰσέλθῃ μετὰ τοῦ πλησίον εἰς τὸν δρυμὸν συναγαγεῖν ξύλα, καὶ ἐκκρουσθῇ ἡ χεὶρ αὐτοῦ τῇ ἀξίνῃ κόπτοντος τὸ ξύλον, καὶ ἐκπεσὸν τὸ σιδήριον ἀπὸ τοῦ ξύλου τύχῃ τοῦ πλησίον, καὶ ἀποθάνῃ, οὗτος καταφεύξεται εἰς μίαν τῶν πόλεων τούτων, καὶ ζήσεται.
\vs{6}Ἵνα μὴ διώξας ὁ ἀγχιστεύων τοῦ αἵματος ὀπίσω τοῦ φονεύσαντος, ὅτι παρατεθέρμανται τῇ καρδίᾳ, καὶ καταλάβῃ αὐτὸν, ἐὰν μακροτέρα ἠ· ἡ ὁδὸς, καὶ πατάξῃ αὐτοῦ ψυχήν· καὶ τούτῳ οὐκ ἔστι κρίσις θανάτου, ὅτι οὐ μισῶν ἦν αὐτὸν πρὸ τῆς χθὲς, οὐδὲ πρὸ τῆς τρίτης.
\vs{7}Διὰ τοῦτο ἐγώ σοι ἐντέλλομαι τὸ ῥῆμα τοῦτο, λέγων, τρεῖς πόλεις διαστελεῖς σεαυτῷ.

\vs{8}Ἐὰν δὲ ἐμπλατύνῃ Κύριος ὁ Θεός σου τὰ ὅριά σου, ὃν τρόπον ὤμοσε τοῖς πατράσι σου, καὶ δῷ σοι Κύριος πᾶσαν τὴν γῆν, ἣν εἶπε δοῦναι τοῖς πατράσι σου,
\vs{9}ἐὰν ἀκούσῃς ποιεῖν πάσας τὰς ἐντολὰς ταύτας, ἃς ἐγὼ ἐντέλλομαί σοι σήμερον, ἀγαπᾷν Κύριον τὸν Θεόν σου, πορεύεσθαι ἐν πάσαις ταῖς ὁδοῖς αὐτοῦ πάσας τὰς ἡμέρας· προσθήσεις σεαυτῷ ἔτι τρεῖς πόλεις πρὸς τὰς τρεῖς ταύτας.
\vs{10}Καὶ οὐκ ἐκχυθήσεται αἷμα ἀναίτιον ἐν τῇ γῇ, ἡ· Κύριος ὁ Θεός σου δίδωσί σοι ἐν κλήρῳ, καὶ οὐκ ἔσται ἐν σοὶ αἵματι ἔνοχος.

\vs{11}Ἐὰν δὲ γένηται ἐν σοὶ ἄνθρωπος μισῶν τὸν πλησίον, καὶ ἐνεδρεύσῃ αὐτὸν, καὶ ἐπαναστῇ ἐπʼ αὐτὸν, καὶ πατάξῃ αὐτοῦ ψυχὴν, καὶ ἀποθάνῃ, καὶ φύγῃ εἰς μίαν τῶν πόλεων τούτων·
\vs{12}καὶ ἀποστελοῦσιν ἡ γερουσία τῆς πόλεως αὐτοῦ, καὶ λήμψονται αὐτὸν ἐκεῖθεν, καὶ παραδώσουσιν αὐτὸν εἰς χεῖρας τῶν ἀγχιστευόντων τοῦ αἵματος, καὶ ἀποθανεῖται.
\vs{13}Οὐ φείσεται ὁ ὀφθαλμός σου ἐπʼ αὐτῷ, καὶ καθαριεῖς τὸ αἷμα τὸ ἀναίτιον ἐξ Ἰσραὴλ, καὶ εὖ σοι ἔσται.

\vs{14}Οὐ μετακινήσεις ὅρια τοῦ πλησίον, ἃ ἔστησαν οἱ πατέρες σου ἐν τῇ κληρονομίᾳ, ᾗ κατεκληρονομήθης ἐν τῇ γῇ, ἣν Κύριος ὁ Θεός σου δίδωσί σοι ἐν κλήρῳ.
\vs{15}Οὐκ ἐμμενεῖ μάρτυς εἷς μαρτυρῆσαι κατὰ ἀνθρώπου κατὰ πᾶσαν ἀδικίαν, καὶ κατὰ πᾶν ἁμάρτημα, καὶ κατὰ πᾶσαν ἁμαρτίαν, ἣν ἐὰν ἁμάρτῃ· ἐπὶ στόματος δύο μαρτύρων, καὶ ἐπὶ στόματος τριῶν μαρτύρων, στήσεται πᾶν ῥῆμα.
\vs{16}Ἐὰν δὲ καταστῇ μάρτυς ἄδικος κατὰ ἀνθρώπου, καταλέγων αὐτοῦ ἀσέβειαν·
\vs{17}καὶ στήσονται οἱ δύο ἄνθρωποι οἷς ἐστιν αὐτοῖς ἡ ἀντιλογία, ἔναντι Κυρίου, καὶ ἔναντι τῶν ἱερέων, καὶ ἔναντι τῶν κριτῶν, οἳ ἂν ὦσιν ἐν ταῖς ἡμέραις ἐκείναις·
\vs{18}Καὶ ἐξετάσωσιν οἱ κριταὶ ἀκριβῶς, καὶ ἰδοὺ μάρτυς ἄδικος ἐμαρτύρησεν ἄδικα, ἀντέστη κατὰ τοῦ ἀδελφοῦ αὐτοῦ·
\vs{19}Καὶ ποιήσετε αὐτῷ ὃν τρόπον ἐπονηρεύσατο ποιῆσαι κατὰ τοῦ ἀδελφοῦ αὐτοῦ, καὶ ἐξαρεῖς τὸ πονηρὸν ἐξ ὑμῶν αὐτῶν.
\vs{20}Καὶ οἱ ἐπίλοιποι ἀκούσαντες φοβηθήσονται, καὶ οὐ προσθήσουσιν ἔτι ποιῆσαι κατὰ τὸ ῥῆμα τὸ πονηρὸν τοῦτο ἐν ὑμῖν.
\vs{21}Οὐ φείσεται ὁ ὀφθαλμός σου ἐπʼ αὐτῷ· ψυχὴν ἀντὶ ψυχῆς, ὀφθαλμὸν ἀντὶ ὀφθαλμοῦ, ὀδόντα ἀντὶ ὀδόντος, χεῖρα ἀντὶ χειρὸς, πόδα ἀντὶ ποδός.

\ch{20}
Ἐὰν δὲ ἐξέλθῃς εἰς πόλεμον ἐπὶ τοὺς ἐχθρούς σου, καὶ ἴδῃς ἵππον καὶ ἀναβάτην καὶ λαὸν πλείονά σου, οὐ φοβηθήσῃ ἀπʼ αὐτῶν, ὅτι Κύριος ὁ Θεός σου μετὰ σοῦ, ὁ ἀναβιβάσας σε ἐκ γῆς Αἰγύπτου.
\vs{2}Καὶ ἔσται ὅταν ἐγγίσῃς τῷ πολέμῳ, καὶ προσεγγίσας ὁ ἱερεὺς λαλήσει τῷ λαῷ,
\vs{3}καὶ ἐρεῖ πρὸς αὐτοὺς, ἄκουε Ἰσραήλ· ὑμεῖς πορεύεσθε σήμερον εἰς τὸν πόλεμον ἐπὶ τοὺς ἐχθροὺς ὑμῶν· μὴ ἐκλυέσθω ἡ καρδία ὑμῶν, μὴ φοβεῖσθε, μηδὲ θραύεσθε, μηδὲ ἐκκλίνετε ἀπὸ προσώπου αὐτῶν.
\vs{4}Ὅτι Κύριος ὁ Θεὸς ὑμῶν ὁ προπορευόμενος μεθʼ ὑμῶν, συνεκπολεμῆσαι ὑμῖν τοὺς ἐχθροὺς ὑμῶν διασῶσαι ὑμᾶς.

\vs{5}Καὶ λαλήσουσιν οἱ γραμματεῖς πρὸς τὸν λαὸν, λέγοντες, τίς ὁ ἄνθρωπος ὁ οἰκοδομήσας οἰκίαν καινὴν, καὶ οὐκ ἐνεκαίνισεν αὐτήν; πορευέσθω καὶ ἀποστραφήτω εἰς τὴν· οἰκίαν αὐτοῦ, μὴ ἀποθάνῃ ἐν τῷ πολέμῳ, καὶ ἄνθρωπος ἕτερος ἐγκαινιεῖ αὐτήν.
\vs{6}Καὶ τίς ὁ ἄνθρωπος ὅστις ἐφύτευσεν ἀμπελῶνα, καὶ οὐκ εὐφράνθη ἐξ αὐτοῦ; πορευέσθω καὶ ἀποστραφήτω εἰς τὴν οἰκίαν αὐτοῦ, μὴ ἀποθάνῃ ἐν τῷ πολέμῳ, καὶ ἄνθρωπος ἕτερος εὐφρανθήσεται ἐξ αὐτοῦ.
\vs{7}Καὶ τίς ὁ ἄνθρωπος ὅστις μεμνήστευται γυναῖκα, καὶ οὐκ ἔλαβεν αὐτήν; πορευέσθω καὶ ἀποστραφήτω εἰς τὴν οἰκίαν αὐτοῦ, μὴ ἀποθάνῃ ἐν τῷ πολέμῳ, καὶ ἄνθρωπος ἕτερος λήψεται αὐτήν.
\vs{8}Καὶ προσθήσουσιν οἱ γραμματεῖς λαλῆσαι πρὸς τὸν λαὸν, καὶ ἐροῦσι, τίς ὁ ἄνθρωπος ὁ φοβούμενος καὶ δειλὸς τῇ καρδίᾳ; πορευέσθω καὶ ἀποστραφήτω εἰς τὴν οἰκίαν αὐτοῦ, ἵνα μὴ δειλιάνῃ τὴν καρδίαν τοῦ ἀδελφοῦ αὐτοῦ, ὥσπερ ἡ αὐτοῦ.
\vs{9}Καὶ ἔσται ὅταν παύσωνται οἱ γραμματεῖς λαλοῦντες πρὸς τὸν λαὸν, καὶ καταστήσουσιν ἄρχοντας τῆς στρατιᾶς προηγουμένους τοῦ λαοῦ.

\vs{10}Ἐὰν δὲ προσέλθῃς πρὸς πόλιν ἐκπολεμῆσαι αὐτοὺς, καὶ ἐκκαλέσαι αὐτοὺς μετʼ εἰρήνης.
\vs{11}Ἐὰν μὲν εἰρηνικὰ ἀποκριθῶσί σοι, καὶ ἀνοίξωσί σοι, ἔσται πᾶς ὁ λαὸς οἱ εὑρεθέντες ἐν αὐτῇ ἔσονταί σοι φορολόγητοι καὶ ὑπήκοοί σου.
\vs{12}Ἐὰν δὲ μὴ ὑπακούσωσί σοι, καὶ ποιῶσι πρὸς σὲ πόλεμον, περικαθιεῖς αὐτὴν,
\vs{13}ἕως ἂν παραδῷ σοι αὐτὴν Κύριος ὁ Θεός σου εἰς τὰς χεῖράς σου, καὶ πατάξεις πᾶν ἀρσενικὸν αὐτῆς ἐν φόνῳ μαχαίρας,
\vs{14}πλὴν τῶν γυναικῶν καὶ τῆς ἀποσκευῆς· καὶ πάντα τὰ κτήνη, καὶ πάντα ὅσα ἂν ὑπάρχῃ ἐν τῇ πόλει, καὶ πᾶσαν τὴν ἀπαρτίαν προνομεύσεις σεαυτῷ, καὶ φαγῇ πᾶσαν τὴν προνομὴν τῶν ἐχθρῶν σου, ὧν Κύριος ὁ Θεός σου δίδωσί σοι.
\vs{15}Οὕτως ποιήσεις πάσας τὰς πόλεις τὰς μακρὰν οὔσας σου σφόδρα, οὐχὶ ἐκ τῶν πόλεων τῶν ἐθνῶν τούτων,
\vs{16}ὧν Κύριος ὁ Θεός σου δίδωσί σοι κληρονομεῖν τὴν γῆν αὐτῶν. Οὐ ζωγρήσετε πᾶν ἐμπνέον,
\vs{17}ἀλλʼ ἢ ἀναθέματι ἀναθεματιεῖτε αὐτοὺς, τὸν Χετταῖον, καὶ Ἀμοῤῥαῖον, καὶ Χαναναῖον, καὶ Φερεζαῖον, καὶ Εὐαῖον, καὶ Ἰεβουσαῖον, καὶ Γεργεσαῖον, ὃν τρόπον ἐνετείλατό σοι Κύριος ὁ Θεός σου,
\vs{18}ἵνα μὴ διδάξωσι ποιεῖν ὑμᾶς πάντα τὰ βδελύγματα αὐτῶν, ὅσα ἐποίησαν τοῖς θεοῖς αὐτῶν, καὶ ἁμαρτήσεσθε ἐναντίον Κυρίου τοῦ Θεοῦ ὑμῶν.

\vs{19}Ἐὰν δὲ περικαθίσῃς περὶ πόλιν μίαν ἡμέρας πλείους ἐκπολεμῆσαι αὐτὴν εἰς κατάληψιν αὐτῆς, οὐκ ἐξολεθρεύσεις τὰ δένδρα αὐτῆς, ἐπιβαλεῖν ἐπʼ αὐτὰ σίδηρον, ἀλλʼ ἢ ἀπʼ αὐτοῦ φαγῇ, αὐτὸ δὲ οὐκ ἐκκόψεις· μὴ ἄνθρωπος τὸ ξύλον τὸ ἐν τῷ ἀγρῷ, εἰσελθεῖν ἀπὸ προσώπου σου εἰς τὸν χάρακα;
\vs{20}Ἀλλὰ ξύλον ὃ ἐπίστασαι ὅτι οὐ καρπόβρωτόν ἐστι, τοῦτο ὀλοθρεύσεις καὶ ἐκκόψεις· καὶ οἰκοδομήσεις χαράκωσιν ἐπὶ τὴν πόλιν, ἥτις ποιεῖ πρὸς σὲ τὸν πόλεμον, ἕως ἂν παραδοθῇ.

\ch{21}
Ἐὰν δὲ εὑρεθῇ τραυματίας ἐν τῇ γῇ, ᾗ Κύριος ὁ Θεός σου δίδωσί σοι κληρονομῆσαι, πεπτωκὼς ἐν τῷ πεδίῳ, καὶ οὐκ οἴδασι τὸν πατάξαντα,
\vs{2}ἐξελεύσεται ἡ γερουσία σου καὶ οἱ κριταί σου, καὶ ἐκμετρήσουσιν ἐπὶ τὰς πόλεις τὰς κύκλῳ τοῦ τραυματίου·
\vs{3}Καὶ ἔσται ἡ πόλις ἡ ἐγγίζουσα τῷ τραυματίᾳ, καὶ λήψεται ἡ γερουσία τῆς πόλεως ἐκείνης δάμαλιν ἐκ βοῶν, ἥτις οὐκ εἴργασται, καὶ ἥτις οὐχ εἵλκυσε ζυγόν·
\vs{4}Καὶ καταβιβάσουσιν ἡ γερουσία τῆς πόλεως ἐκείνης δάμαλιν εἰς φάραγγα τραχεῖαν, ἥτις οὐκ εἴργασται οὐδὲ σπείρεται, καὶ νευροκοπήσουσι τὴν δάμαλιν ἐν τῇ φάραγγι.
\vs{5}Καὶ προσελεύσονται οἱ ἱερεῖς οἱ Λευῖται, ὅτι αὐτοὺς ἐπέλεξε Κύριος ὁ Θεὸς παρεστηκέναι αὐτῷ, καὶ εὐλογεῖν ἐπὶ τῷ ὀνόματι αὐτοῦ· καὶ ἐπὶ τῷ στόματι αὐτῶν ἔσται πᾶσα ἀντιλογία, καὶ πᾶσα ἁφή.
\vs{6}Καὶ πᾶσα ἡ γερουσία τῆς πόλεως ἐκείνης οἱ ἐγγίζοντες τῷ τραυματίᾳ νίψονται τὰς χεῖρας ἐπὶ τὴν κεφαλὴν τῆς δαμάλεως τῆς νενευροκοπημένης ἐν τῇ φάραγγι·
\vs{7}καὶ ἀποκριθέντες, ἐροῦσιν, αἱ χεῖρες ἡμῶν οὐκ ἐξέχεαν τὸ αἷμα τοῦτο, καὶ οἱ ὀφθαλμοὶ ἡμῶν οὐκ ἑωράκασιν.
\vs{8}Ἵλεως γένου τῷ λαῷ σου Ἰσραὴλ, οὓς ἐλυτρώσω Κύριε, ἵνα μὴ γένηται αἷμα ἀναίτιον ἐν τῷ λαῷ σου Ἰσραήλ· καὶ ἐξιλασθήσεται αὐτοῖς τὸ αἷμα.
\vs{9}Σὺ δὲ ἐξαρεῖς τὸ αἷμα τὸ ἀναίτιον ἐξ ὑμῶν αὐτῶν, ἐὰν ποιήσῃς τὸ καλὸν καὶ τὸ ἀρεστὸν ἔναντι Κυρίου τοῦ Θεοῦ σου.

\vs{10}Ἐὰν δὲ ἐξελθὼν εἰς πόλεμον ἐπὶ τοὺς ἐχθρούς σου, καὶ παραδῷ σοι Κύριος ὁ Θεός σου εἰς τὰς χεῖράς σου, καὶ προνομεύσῃς τὴν προνομὴν αὐτῶν,
\vs{11}καὶ ἴδῃς ἐν τῇ προνομῇ γυναῖκα καλὴν τῷ εἴδει, καὶ ἐνθυμηθῇς αὐτῆς, καὶ λάβῃς αὐτὴν σεαυτῷ γυναῖκα,
\vs{12}καὶ εἰσάξῃς αὐτὴν ἔνδον εἰς τὴν οἰκίαν σου, καὶ ξυρήσεις τὴν κεφαλὴν αὐτῆς, καὶ περιονυχιεῖς αὐτὴν,
\vs{13}καὶ περιελεῖς τὰ ἱμάτια τῆς αἰχμαλωσίας ἀπʼ αὐτῆς, καὶ καθιεῖται ἐν τῇ οἰκίᾳ σου, καὶ κλαύσεται τὸν πατέρα καὶ τὴν μητέρα μηνὸς ἡμέρας· καὶ μετὰ ταῦτα εἰσελεύσῃ πρὸς αὐτὴν καὶ συνοικισθήσῃ αὐτῇ, καὶ ἔσται σου γυνή.

\vs{14}Καὶ ἔσται ἐὰν μὴ θέλῃς αὐτὴν, ἐξαποστελεῖς αὐτὴν ἐλευθέραν, καὶ πράσει οὐ πραθήσεται ἀργυρίου· οὐκ ἀθετήσεις αὐτὴν, διότι ἐταπείνωσας αὐτήν.

\vs{15}Ἐὰν δὲ γένωνται ἀνθρώπῳ δύο γυναῖκες, μία αὐτῶν ἠγαπημένη, καὶ μία αὐτῶν μισουμένη, καὶ τέκωσιν αὐτῷ ἡ ἠγαπημένη καὶ ἡ μισουμένη, καὶ γένηται υἱὸς πρωτότοκος τῆς μισουμένης·
\vs{16}Καὶ ἔσται ἡ· ἂν ἡμέρᾳ κατακληρονομῇ τοῖς υἱοῖς αὐτοῦ τὰ ὑπάρχοντα αὐτοῦ, οὐ δυνήσεται πρωτοτοκεῦσαι τῷ υἱῷ τῆς ἠγαπημένης, ὑπεριδὼν τὸν υἱὸν τῆς μισουμένης τὸν πρωτότοκον·
\vs{17}Ἀλλὰ τὸν πρωτότοκον υἱὸν τῆς μισουμένης ἐπιγνώσεται δοῦναι αὐτῷ διπλᾶ ἀπὸ πάντων ὧν ἂν εὑρεθῇ αὐτῷ, ὅτι οὗτός ἐστιν ἀρχὴ τέκνων αὐτοῦ, καὶ τούτῳ καθήκει τὰ πρωτοτοκεῖα.
\vs{18}Ἐὰν δέ τινι ἠ· υἱὸς ἀπειθῆς καὶ ἐρεθιστὴς οὐχ ὑπακούων φωνὴν πατρὸς καὶ φωνὴν μητρὸς, καὶ παιδεύωσιν αὐτὸν, καὶ μὴ εἰσακούῃ αὐτῶν·
\vs{19}Καὶ συλλαβόντες αὐτὸν ὁ πατὴρ αὐτοῦ καὶ ἡ μήτηρ αὐτοῦ, καὶ ἐξάξουσιν αὐτὸν ἐπὶ τὴν γερουσίαν τῆς πόλεως αὐτοῦ, καὶ ἐπὶ τὴν πύλην τοῦ τόπου·
\vs{20}Καὶ ἐροῦσι τοῖς ἀνδράσι τῆς πόλεως αὐτῶν, ὁ υἱὸς ἡμῶν οὗτος ἀπειθεῖ καὶ ἐρεθίζει, οὐχ ὑπακούει τῆς φωνῆς ἡμῶν, συμβολοκοπῶν οἰνοφλυγεῖ.
\vs{21}Καὶ λιθοβολήσουσιν αὐτὸν οἱ ἄνδρες τῆς πόλεως αὐτοῦ ἐν λίθοις, καὶ ἀποθανεῖται· καὶ ἐξαρεῖς τὸν πονηρὸν ἐξ ὑμῶν αὐτῶν· καὶ οἱ ἐπίλοιποι ἀκούσαντες φοβηθήσονται.

\vs{22}Ἐὰν δὲ γένηται ἔν τινι ἁμαρτία, κρίμα θανάτου, καὶ ἀποθάνῃ, καὶ κρεμάσητε αὐτὸν ἐπὶ ξύλου·
\vs{23}οὐ κοιμηθήσεται τὸ σῶμα αὐτοῦ ἐπὶ τοῦ ξύλου, ἀλλὰ ταφῇ θάψετε αὐτὸ ἐν τῇ ἡμέρᾳ ἐκείνῃ, ὅτι κεκατηραμένος ὑπὸ Θεοῦ πᾶς κρεμάμενος ἐπὶ ξύλου· καὶ οὐ μὴ μιανεῖτε τὴν γῆν, ἣν Κύριος ὁ Θεός σου δίδωσί σοι ἐν κλήρῳ.

\ch{22}
Μὴ ἰδὼν τὸν μόσχον τοῦ ἀδελφοῦ σου, ἢ τὸ πρόβατον αὐτοῦ, πλανώμενα ἐν τῇ ὁδῷ, ὑπερίδῃς αὐτά· ἀποστροφῇ ἀποστρέψεις αὐτὰ τῷ ἀδελφῷ σου, καὶ ἀποδώσεις αὐτῷ.
\vs{2}Ἐὰν δὲ μὴ ἐγγίζῃ ὁ ἀδελφός σου πρὸς σὲ, μηδὲ ἐπίοτῃ αὐτὸν, συνάξεις αὐτὸν ἔνδον εἰς τὴν οἰκίαν σου, καὶ ἔσται μετὰ σοῦ ἕως ἂν ζητήσῃ αὐτὰ ὁ ἀδελφός σου, καὶ ἀποδώσεις αὐτῷ.
\vs{3}Οὕτω ποιήσεις τὸν ὄνον αὐτοῦ, καὶ οὕτω ποιήσεις τὸ ἱμάτιον αὐτοῦ, καὶ οὕτω ποιήσεις κατὰ πᾶσαν ἀπώλειαν τοῦ ἀδελφοῦ σου· ὅσα ἐὰν ἀπολῆται παρʼ αὐτοῦ, καὶ εὕρῃς, οὐ δυνήσῃ ὑπεριδεῖν.
\vs{4}Οὐκ ὄψῃ τὸν ὄνον τοῦ ἀδελφοῦ σου ἢ τὸν μόσχον αὐτοῦ πεπτωκότας ἐν τῇ ὁδῷ, μὴ ὑπερίδῃς αὐτοὺς, ἀνιστῶν ἀναστήσεις μετʼ αὐτοῦ.

\vs{5}Οὐκ ἔσται σκεύη ἀνδρὸς ἐπὶ γυναικὶ, οὐδὲ μὴ ἐνδύσηται ἀνὴρ στολὴν γυναικίαν, ὅτι βδέλυγμα Κυρίῳ τῷ Θεῷ σου ἐστὶ πᾶς ποιῶν ταῦτα.
\vs{6}Ἐὰν δὲ συναντήσῃς νοσσιᾷ ὀρνέων πρὸ προσώπου σου ἐν τῇ ὁδῷ ἢ ἐπὶ παντὶ δένδρῳ, ἢ ἐπὶ τῆς γῆς, νοσσοῖς ἢ ὠοῖς, καὶ ἡ μήτηρ θάλπῃ ἐπὶ τῶν νοσσῶν ἢ ἐπὶ τῶν ὠῶν, οὐ λήψῃ τὴν μητέρα μετὰ τῶν τέκνων.
\vs{7}Ἀποστολῇ ἀποστελεῖς τὴν μητέρα, τὰ δὲ παιδία λήψῃ σεαυτῷ, ἵνα εὖ σοι γένηται καὶ πολυήμερος γένῃ.

\vs{8}Ἐὰν οἰκοδομήσῃς οἰκίαν καινήν, καὶ ποιήσεις στεφάνην τῷ δώματί σου, καὶ οὐ ποιήσεις φόνον ἐν τῇ οἰκίᾳ σου, ἐὰν πέσῃ ὁ πεσὼν ἀπʼ αὐτοῦ.
\vs{9}Οὐ κατασπερεῖς τὸν ἀμπελῶνά σου διάφορον, ἵνα μὴ ἁγιασθῇ τὸ γέννημα, καὶ τὸ σπέρμα ὃ ἐὰν σπείρῃς μετὰ τοῦ γεννήματος τοῦ ἀμπελῶνός σου.
\vs{10}Οὐκ ἀροτριάσεις ἐν μόσχῳ καὶ ὄνῳ ἐπὶ τὸ αὐτό.
\vs{11}Οὐκ ἐνδύσῃ κίβδηλον, ἔρια καὶ λίνον ἐν τῷ αὐτῷ.
\vs{12}Στρεπτὰ ποιήσεις σεαυτῷ ἐπὶ τῶν τεσσάρων κρασπέδων τῶν περιβολαίων σου, ἃ ἐὰν περιβάλῃ ἐν αὐτοῖς.

\vs{13}Ἐὰν δέ τις λάβῃ γυναῖκα καὶ συνοικήσῃ αὐτῇ, καὶ μισήσῃ αὐτήν,
\vs{14}καὶ ἐπιθῇ αὐτῇ προφασιστικοὺς λόγους, καὶ κατενέγκῃ αὐτῆς ὄνομα πονηρὸν, καὶ λέγῃ, τὴν γυναῖκα ταύτην εἴληφα, καὶ προσελθὼν αὐτῇ οὐχ εὕρηκα αὐτῆς τὰ παρθένια·
\vs{15}Καὶ λαβὼν ὁ πατὴρ τῆς παιδὸς καὶ ἡ μήτηρ ἐξοίσουσι τὰ παρθένια τῆς παιδὸς πρὸς τὴν γερουσίαν ἐπὶ τὴν πύλην.
\vs{16}Καὶ ἐρεῖ ὁ πατὴρ τῆς παιδὸς τῇ γερουσίᾳ, τὴν θυγατέρα μου ταύτην δέδωκα τῷ ἀνθρώπῳ τούτῳ γυναῖκα, καὶ μισήσας αὐτὴν
\vs{17}νῦν οὗτος, ἐπιτίθησιν αὐτῇ προφασιστικοὺς λόγους, λέγων, οὐχ εὕρηκα τῇ θυγατρί σου παρθένια· καὶ ταῦτα τὰ παρθένια τῆς θυγατρός μου. Καὶ ἀναπτύξουσι τὸ ἱμάτιον ἐναντίον τῆς γερουσίας τῆς πόλεως.
\vs{18}Καὶ λήψεται ἡ γερουσία τῆς πόλεως ἐκείνης τὸν ἄνθρωπον ἐκεῖνον, καὶ παιδεύσουσιν αὐτόν,
\vs{19}καὶ ξημιώσουσιν αὐτὸν ἑκατὸν σίκλους, καὶ δώσουσι τῷ πατρὶ τῆς νεάνιδος, ὅτι ἐξήνεγκεν ὄνομα πονηρὸν ἐπὶ παρθένον Ἰσραηλίτιν, καὶ αὐτοῦ ἔσται γυνή· οὐ δυνήσεται ἐξαποστεῖλαι αὐτὴν τὸν ἅπαντα χρόνον.

\vs{20}Ἐὰν δὲ ἐπʼ ἀληθείας γένηται ὁ λόγος οὗτος, καὶ μὴ εὑρεθῇ παρθένια τῇ νεάνιδι,
\vs{21}καὶ ἐξάξουσι τὴν νεᾶνιν ἐπὶ τὰς θύρας τοῦ οἴκου τοῦ πατρὸς αὐτῆς, καὶ λιθοβολήσουσιν αὐτὴν ἐν λίθοις, καὶ ἀποθανεῖται, ὅτι ἐποίησεν ἀφροσύνην ἐν υἱοῖς Ἰσραήλ ἐκπορνεῦσαι τὸν οἶκον τοῦ πατρὸς αὐτῆς· καὶ ἐξαρεῖς τὸν πονηρὸν ἐξ ὑμῶν αὐτῶν.

\vs{22}Ἐὰν δὲ εὑρεθῇ ἄνθρωπος κοιμώμενος μετὰ γυναικὸς συνωκισμένης ἀνδρί, ἀποκτενεῖτε ἀμφοτέρους, τὸν ἄνδρα τὸν κοιμώμενον μετὰ τῆς γυναικὸς, καὶ τὴν γυναῖκα· καὶ ἐξαρεῖς τὸν πονηρὸν ἐξ Ἰσραήλ.
\vs{23}Ἐὰν δὲ γένηται παῖς παρθένος μεμνηστευμένη ἀνδρί, καὶ εὑρὼν αὐτὴν ἄνθρωπος ἐν πόλει κοιμηθῇ μετʼ αὐτῆς,
\vs{24}ἐξάξετε ἀμφοτέρους ἐπὶ τὴν πυλὴν τῆς πόλεως αὐτῶν, καὶ λιθοβοληθήσονται ἐν λίθοις, καὶ ἀποθανοῦνται· τὴν νεᾶνιν, ὅτι οὐκ ἐβόησεν ἐν τῇ πόλει· καὶ τὸν ἄνθρωπον, ὅτι ἐταπείνωσε τὴν γυναῖκα τοῦ πλησίον· καὶ ἐξαρεῖς τὸν πονηρὸν ἐξ ὑμῶν αὐτῶν.
\vs{25}Ἐὰν δὲ ἐν πεδίῳ εὕρῃ ἄνθρωπος τὴν παῖδα τὴν μεμνηστευμένῃν, καὶ βιασάμενος κοιμηθῇ μετʼ αὐτῆς, ἀποκτενεῖτε τὸν κοιμώμενον μετʼ αὐτῆς μόνον.
\vs{26}Καὶ τῇ νεάνιδι οὐκ ἔστιν ἁμάρτημα θανάτου· ὡς εἴ τις ἐπαναστῇ ἄνθρωπος ἐπὶ τὸν πλησίον, καὶ φονεύσῃ αὐτοῦ ψυχήν, οὕτω τὸ πρᾶγμα τοῦτο,
\vs{27}ὅτι ἐν τῷ ἀγρῷ εὗρεν αὐτήν· ἐβόησεν ἡ νεᾶνις ἡ μεμνηστευμένη, καὶ οὐκ ἦν ὁ βοηθήσων αὐτῇ.

\vs{28}Ἐὰν δέ τις εὕρῃ τὴν παῖδα τὴν παρθένον, ἥτις οὐ μεμνήστευται, καὶ βιασάμενος κοιμηθῇ μετʼ αὐτῆς, καὶ εὑρεθῇ,
\vs{29}δώσει ὁ ἄνθρωπος ὁ κοιμηθεὶς μετʼ αὐτῆς τῷ πατρὶ τῆς νεάνιδος πεντήκοντα δίδραχμα ἀργυρίου, καὶ αὐτοῦ ἔσται γυνή, ἀνθʼ ὧν ἐταπείνωσεν αὐτήν· οὐ δυνήσεται ἐξαποστεῖλαι αὐτὴν τὸν ἅπαντα χρόνον.

\ch{23}Οὐ λήψεται ἄνθρωπος τὴν γυναῖκα τοῦ πατρὸς αὐτοῦ, καὶ οὐκ ἀποκαλύψει συγκάλυμμα τοῦ πατρὸς αὐτοῦ.

\vs{2}Οὐκ εἰσελεύσεται θλαδίας, οὐδὲ ἀποκεκομμένος, εἰς ἐκκλησίαν Κυρίου.
\vs{3}Οὐκ εἰσελεύσεται ἐκ πόρνης εἰς ἐκκλησίαν Κύριου.

\vs{4}Οὐκ εἰσελεύσεται Ἀμμανίτης καὶ Μωαβίτης εἰς ἐκκλησίαν Κυρίου, καὶ ἕως δεκάτης γενεᾶς οὐκ εἰσελεύσεται εἰς ἐκκλησίαν Κυρίου, καὶ ἕως εἰς τὸν αἰῶνα·
\vs{5}παρὰ τὸ μὴ συναντῆσαι αὐτοὺς ὑμῖν μετὰ ἄρτων καὶ ὕδατος ἐν τῇ ὁδῷ, ἐκπορευομένων ὑμῶν ἐξ Αἰγύπτου, καὶ ὅτι ἐμισθώσαντο ἐπὶ σὲ τὸν Βαλαὰμ υἱὸν Βεὼρ ἐκ τῆς Μεσοποταμίας καταρᾶσθαί σε.
\vs{6}Καὶ οὐκ ἠθέλησε Κύριος ὁ Θεός σου εἰσακοῦσαι τοῦ Βαλαάμ· καὶ μετέστρεψε Κύριος ὁ Θεός σου τὰς κατάρας εἰς εὐλογίαν, ὅτι ἠγάπησέ σε Κύριος ὁ Θεός σου.
\vs{7}Οὐ προσαγορεύσεις εἰρηνικὰ αὐτοῖς καὶ συμφέροντα αὐτοῖς πάσας τὰς ἡμέρας σου εἰς τὸν αἰῶνα.
\vs{8}Οὐ βδελύξῃ Ἰδουμαῖον, ὅτι ἀδελφός σού ἐστίν· οὐ βδελύξῃ Αἰγύπτιον, ὅτι πάροικος ἐγένου ἐν τῇ γῇ αὐτοῦ.
\vs{9}Υἱοὶ ἐὰν γεννηθῶσιν αὐτοῖς, γενεᾷ τρίτῃ εἰσελεύσονται εἰς ἐκκλησίαν Κυρίου.

\vs{10}Ἐὰν δὲ ἐξέλθῃς παρεμβαλεῖν ἐπὶ τοὺς ἐχθρούς σου, καὶ φυλάξῃ ἀπὸ παντὸς ῥήματος πονηροῦ.
\vs{11}Ἐὰν ἠ· ἐν σοὶ ἄνθρωπος ὃς οὐκ ἔσται καθαρὸς ἐκ ῥύσεως αὐτοῦ νυκτὸς, καὶ ἐξελεύσεται ἔξω τῆν παρεμβολῆς, καὶ οὐκ εἰσελεύσεται εἰς τὴν παρεμβολήν.
\vs{12}Καὶ ἔσται τὸ πρὸς ἐσπέραν λούσεται τὸ σῶμα αὐτοῦ ὕδατι, καὶ δεδυκότος ἡλίου εἰσελεύσεται εἰς τὴν παρεμβολήν.
\vs{13}Καὶ τόπος ἔσται σοι ἔξω τῆς παρεμβολῆς, καὶ ἐξελεύσῃ ἐκεῖ ἔξω.
\vs{14}Καὶ πάσσαλος ἔσται σοι ἐπὶ τῆς ζώνης σου· καὶ ἔσται ὅταν διακαθιζάνῃς ἔξω, καὶ ὀρύξεις ἐν αὐτῷ, καὶ ἐπαγαγὼν καλύψεις τὴν ἀσχημοσύνην σου·
\vs{15}Ὅτι Κύριος ὁ Θεός σου ἐμπεριπατεῖ ἐν τῇ παρεμβολῇ σου ἐξελέσθαι σε καὶ παραδοῦναι τὸν ἐχθρόν σου πρὸ προσώπου σου· καὶ ἔσται ἡ παρεμβολή σου ἁγία, καὶ οὐκ ὀφθήσεται ἐν σοὶ ἀσχημοσύνη πράγματος, καὶ ἀποστρέψει ἀπὸ σοῦ.

\vs{16}Οὐ παραδώσεις παῖδα τῷ κυρίῳ αὐτοῦ, ὃς προστέθειταί σοι παρὰ τοῦ κυρίου αὐτοῦ.
\vs{17}Μετὰ σοῦ κατοικήσει, ἐν ὑμῖν κατοικήσει οὗ ἂν ἀρέσῃ αὐτῷ· οὐ θλίψεις αὐτόν.
\vs{18}Οὐκ ἔσται πόρνη ἀπὸ θυγατέρων Ἰσραὴλ, καὶ οὐκ ἔσται πορνεύων ἀπὸ υἱῶν Ἰσραήλ· οὐκ ἔσται τελεσφόρος ἀπὸ θυγατέρων Ἰσραήλ, καὶ οὐκ ἔσται τελισκόμενος ἀπὸ υἱῶν Ἰσραήλ.
\vs{19}Οὐ προσοίσεις μίσθωμα πόρνης, οὐδὲ ἄλλαγμα κυνὸς εἰς τὸν οἶκον Κυρίου τοῦ Θεοῦ σου πρὸς πᾶσαν εὐχὴν, ὅτι βδέλυγμα Κυρίῳ τῷ Θεῷ σου ἐστὶ καὶ ἀμφότερα.

\vs{20}Οὐκ ἐκτοκιεῖς τῷ ἀδελφῷ σου τόκον ἀργυρίου, καὶ τόκον βρωμάτων, καὶ τοκον παντὸς πράγματος, οὗ ἐὰν ἐκδανίσῃς.
\vs{21}Τῷ ἀλλοτρίῳ ἐκτοκιεῖς, τῷ δὲ ἀδελφῷ σου οὐκ ἐκτοκιεῖς, ἵνα εὐλογήσῃ σε Κύριος ὁ Θεός σου ἐν πᾶσι τοῖς ἔργοις σου ἐπὶ τῆς γῆς, εἰς ἣν εἰσπορεύῃ ἐκεῖ κληρονομῆσαι αὐτήν.

\vs{22}Ἐὰν δὲ εὔξῃ εὐχὴν Κυρίῳ τῷ Θεῷ σου, οὐ χρονιεῖς ἀποδοῦναι αὐτὴν, ὅτι ἐκζητῶν ἐκζητήσει Κύριος ὁ Θεός σου παρὰ σοῦ, καὶ ἔσται ἐν σοὶ ἁμαρτία.
\vs{23}Ἐὰν δὲ μὴ θέλῃς εὔξασθαι, οὐκ ἔστιν ἐν σοὶ ἁμαρτία.
\vs{24}Τὰ ἐκπορευόμενα διὰ τῶν χειλέων σου φυλάξῃ, καὶ ποιήσεις ὃν τρόπον ηὔξω Κυρίῳ τῷ Θεῷ δόμα, ὃ ἐλάλησας τῷ στόματί σου.

\vs{25}Ἐὰν δὲ εἰσέλθῃς εἰς ἀμητὸν τοῦ πλησίον σου, καὶ συλλέξῃς ἐν ταῖς χερσί σου στάχυς, καὶ δρέπανον οὐ μὴ ἐπιβάλῃς ἐπʼ ἀμητὸν τοῦ πλησίον σου.
\vs{26}Ἐὰν δὲ εἰσέλθῃς εἰς τὸν ἀμπελῶνα τοῦ πλησίον σου, φαγῇ σταφυλὴν, ὅσον ψυχήν σου ἐμπλησθῆναι, εἰς δὲ ἄγγος οὐκ ἐμβάλῃς.

\ch{24}Ἐὰν δέ τις λάβῃ γυναῖκα, καὶ συνοικήσῃ αὐτῇ, καὶ ἔσται ἐὰν μὴ εὕρῃ χάριν ἐναντίον αὐτοῦ, ὅτι εὗρεν ἐν αὐτῇ ἄσχημον πρᾶγμα, καὶ γράψει αὐτῇ βιβλίον ἀποστασίου, καὶ δώσει εἰς τὰς χεῖρας αὐτῆς, καὶ ἐξαποστελεῖ αὐτὴν ἐκ τῆς οἰκίας αὐτοῦ,
\vs{2}καὶ ἀπελθοῦσα γένηται ἀνδρὶ ἑτέρῳ,
\vs{3}καὶ μισήσῃ αὐτὴν ὁ ἀνὴρ ὁ ἔσχατος, καὶ γράψῃ αὐτῇ βιβλίον ἀποστασίου, καὶ δώσει εἰς τὰς χεῖρας αὐτῆς, καὶ ἐξαποστελεῖ αὐτὴν ἐκ τῆς οἰκίας αὐτοῦ, καὶ ἀποθάνῃ ὁ ἀνὴρ ὁ ἔσχατος, ὃς ἔλαβεν αὐτὴν ἐαυτῷ γυναῖκα,
\vs{4}οὐ δυνήσεται ὁ ἀνὴρ ὁ πρότερος ὁ ἐξαποστείλας αὐτὴν, ἐπαναστρέψας λαβεῖν αὐτὴν ἑαυτῷ γυναῖκα, μετὰ τὸ μιανθῆναι αὐτὴν, ὅτι βδέλυγμά ἐστιν ἐναντίον Κυρίου τοῦ Θεοῦ σου, καὶ οὐ μιανεῖτε τὴν γῆν, ἣν Κύριος ὁ Θεός σου δίδωσί σοι ἐν κλήρῳ.

\vs{5}Ἐὰν δέ τις λάβῃ γυναῖκα προσφάτως, οὐκ ἐξελεύσεται εἰς πόλεμον, καὶ οὐκ ἐπιβληθήσεται αὐτῷ οὐδὲν πρᾶγμα· ἀθῶος ἔσται ἐν τῇ οἰκίᾳ αὐτοῦ, ἐνιαυτὸν ἕνα εὐφρανεῖ τὴν γυναῖκα αὐτοῦ ἣν ἔλαβεν.

\vs{6}Οὐκ ἐνεχυράσεις μύλον, οὐδὲ ἐπιμύλιον, ὅτι ψυχὴν οὗτος ἐνεχυράζει.
\vs{7}Ἐὰν δὲ ἁλῷ ἄνθρωπος κλέπτων ψυχὴν ἐκ τῶν ἀδελφῶν αὐτοῦ τῶν υἱῶν Ἰσραὴλ, καὶ καταδυναστεύσας αὐτὸν ἀποδῶται, ἀποθανεῖται ὁ κλέπτης ἐκεῖνος· καὶ ἐξαρεῖς τὸν πονηρὸν ἐξ ὑμῶν αὐτῶν.
\vs{8}Πρόσεχε σεαυτῷ ἐν τῇ ἁφῇ τῆς λέπρας· φυλάξῃ σφόδρα ποιεῖν κατὰ πάντα τὸν νόμον, ὃν ἂν ἀναγγείλωσιν ὑμῖν οἱ ἱερεῖς οἱ Λευῖται· ὃν τρόπον ἐνετειλάμην ὑμῖν, φυλάξασθε ποιεῖν.
\vs{9}Μνήσθητι ὅσα ἐποίησε Κύριος ὁ Θεός σου τῇ Μαριὰμ ἐν τῇ ὁδῷ, ἐκπορευομένων ὑμῶν ἐξ Αἰγύπτου.

\vs{10}Ἐὰν ὀφείλημα ἠ· ἐν τῷ πλησίον σου, ὀφείλημα ὁτιοῦν, οὐκ εἰσελεύσῃ εἰς τὴν οἰκίαν αὐτοῦ ἐνεχυράσαι τὸ ἐνέχυρον αὐτοῦ.
\vs{11}Ἔξω στήσῃ, καὶ ὁ ἄνθρωπος οὗ τὸ δάνειόν σού ἐστιν ἐν αὐτῷ, ἐξοίσει σοι τὸ ἐνέχυρον ἔξω.
\vs{12}Ἐὰν δὲ ὁ ἄνθρωπος πένηται, οὐ κοιμηθήσῃ ἐν τῷ ἐνεχύρῳ αὐτοῦ.
\vs{13}Ἀποδόσει ἀποδώσεις τὸ ἐνέχυρον αὐτοῦ πρὸς δυσμαῖς ἡλίου, καὶ κοιμηθήσεται ἐν τῷ ἱματίῳ αὐτοῦ, καὶ εὐλογήσει σε, καὶ ἔσται σοι ἐλεημοσύνη ἐναντίον Κυρίου τοῦ Θεοῦ σου.
\vs{14}Οὐκ ἀπαδικήσεις μισθὸν πένητος καὶ ἐνδεοῦς ἐκ τῶν ἀδελφῶν σου, ἢ ἐκ τῶν προσηλύτων τῶν ἐν ταῖς πόλεσίν σου.
\vs{15}Αὐθημερὸν ἀποδώσεις τὸν μισθὸν αὐτοῦ, οὐκ ἐπιδύσεται ὁ ἥλιος ἐπʼ αὐτῷ, ὅτι πένης ἐστὶ, καὶ ἐν αὐτῷ ἔχει τὴν ἐλπίδα, καὶ καταβοήσεται κατὰ σοῦ πρὸς Κύριον, καὶ ἔσται ἐν σοὶ ἁμαρτία.
\vs{16}Οὐκ ἀποθανοῦνται πατέρες ὑπὲρ τέκνων, καὶ οἱ υἱοὶ οὐκ ἀποθανοῦνται ὑπὲρ πατέρων· ἕκαστος ἐν τῇ ἑαυτοῦ ἁματρίᾳ ἀποθανεῖται.
\vs{17}Οὐκ ἐκκλινεῖς κρίσιν προσηλύτου καὶ ὀρφανοῦ καὶ χήρας·
\vs{18}οὐκ ἐνεχυράσεις ἱμαίτιον χήρας, καὶ μνησθήσῃ ὅτι οἰκέτης ἦσθα ἐν γῇ Αἰγύπτῳ, καὶ ἐλυτρώσατό σε Κύριος ὁ Θεός σου ἐκεῖθεν· διὰ τοῦτο ἐγώ σοι ἐντέλλομαι ποιεῖν τὸ ῥῆμα τοῦτο.

\vs{19}Ἐὰν δὲ ἀμήσῃς ἀμητὸν ἐν τῷ ἀγρῷ σου, καὶ ἐπιλάθῃ δράγμα ἐν τῷ ἀγρῷ σου, οὐκ ἀναστραφήσῃ λαβεῖν αὐτό· τῷ προσηλύτῳ καὶ τῷ ὀρφανῷ καὶ τῇ χήρᾳ ἔσται, ἵνα εὐλογήσῃ σε Κύριος ὁ Θεός σου ἐν πᾶσι τοῖς ἔργοις τῶν χειρῶν σου.
\vs{20}Ἐὰν δὲ ἐλαιολογῇς, οὐκ ἐπαναστρέψεις καλαμήσασθαι τὰ ὀπίσω σου· τῷ προσηλύτῳ καὶ τῷ ὀρφανῷ καὶ τῇ χήρᾳ ἔσται· καὶ μνησθήσῃ ὅτι οἰκέτης ἦσθα ἐν γῇ Αἰγύπτῳ· διὰ τοῦτο ἐγώ σοι ἐντέλλομαι ποιεῖν τὸ ῥῆμα τοῦτο.
\vs{21}Ἐὰν δὲ τρυγήσῃς τὸν ἀμπελῶνά σου, οὐκ ἐπανατρυγήσεις αὐτὸν τὰ ὀπίσω σου· τῷ προσηλύτῳ καὶ τῷ ὀρφανῷ καὶ τῇ χήρᾳ ἔσται·
\vs{22}καὶ μνησθήσῃ ὅτι οἰκέτης ἦσθα ἐν γῇ Αἰγύπτῳ· διὰ τοῦτο ἐγώ σοι ἐντέλλομαι ποιεῖν τὸ ῥῆμα τοῦτο.

\ch{25}
Ἐὰν δὲ γένηται ἀντιλογία ἀναμέσον ἀνθρώπων, καὶ προσέλθωσιν εἰς κρίσιν, καὶ κρίνωσι, καὶ δικαιώσωσι τὸ δίκαιον, καὶ καταγνῶσι τοῦ ἀσεβοῦς·
\vs{2}Καὶ ἔσται, ἐὰν ἄξιος ἠ· πληγῶν ὁ ἀσεβῶν, καθιεῖς αὐτὸν ἔναντι τῶν κριτῶν, καὶ μαστιγώσουσιν αὐτὸν ἐναντίον αὐτῶν κατὰ τὴν ἀσέβειαν αὐτοῦ.
\vs{3}Καὶ ἀριθμῷ τεσσαράκοντα μαστιγώσουσιν αὐτόν οὐ προσθήσουσιν· ἐὰν δὲ προσθῇς μαστιγῶσαι ὑπὲρ ταύτας τὰς πληγὰς πλείους, ἀσχημονήσει ὁ ἀδελφός σου ἐναντίον σου.
\vs{4}Οὐ φιμώσεις βοῦν ἀλοῶντα.

\vs{5}Ἐὰν δὲ κατοικῶσιν ἀδελφοὶ ἐπὶ τὸ αὐτό, καὶ ἀποθάνῃ εἷς ἐξ αὐτῶν, σπέρμα δὲ μὴ ἠν αὐτῷ, οὐκ ἔσται ἡ γυνὴ τοῦ τεθνηκότος ἔξω ἀνδρὶ μὴ ἐγγίζοντι· ὁ ἀδελφὸς τοῦ ἀνδρὸς αὐτῆς εἰσελεύσεται πρὸς αὐτὴν, καὶ λήψεται αὐτὴν ἑαυτῷ γυναῖκα, καὶ συνοικήσει αὐτῇ.
\vs{6}Καὶ ἔσται τὸ παιδίον ὃ ἐὰν τέκῃ, κατασταθήσεται ἐκ τοῦ ὀνόματος τοῦ τετελευτηκότος, καὶ οὐκ ἐξαλειφθήσεται τὸ ὄνομα αὐτοῦ ἐξ Ἰσραήλ.

\vs{7}Ἐὰν δὲ μὴ βούληται ὁ ἄνθρωπος λαβεῖν τὴν γυναῖκα τοῦ ἀδελφοῦ αὐτοῦ, καὶ ἀναβήσεται ἡ γυνὴ ἐπὶ τὴν πύλην ἐπὶ τὴν γερουσίαν, καὶ ἐρεῖ, οὐ θέλει ὁ ἀδελφὸς τοῦ ἀνδρός μου ἀναστῆσαι τὸ ὄνομα τοῦ ἀδελφοῦ αὐτοῦ ἐν Ἰσραήλ, οὐκ ἠθέλησεν ὁ ἀδελφὸς τοῦ ἀνδρός μου.
\vs{8}Καὶ καλέσουσιν αὐτὸν ἡ γερουσία τῆς πόλεως αὐτοῦ, καὶ ἐροῦσιν αὐτῷ· καὶ στὰς εἴπῃ, οὐ βούλομαι λαβεῖν αὐτήν·
\vs{9}καὶ προσελθοῦσα ἡ γυνὴ τοῦ ἀδελφοῦ αὐτοῦ ἔναντι τῆς γερουσίας, καὶ ὑπολύσει τὸ ὑπόδημα αὐτοῦ τὸ ἓν ἀπὸ τοῦ ποδὸς αὐτοῦ, καὶ ἐμπτύσεται κατὰ πρόσωπον αὐτοῦ, καὶ ἀποκριθεῖσα ἐρεῖ, οὕτως ποιήσουσι τῷ ἀνθρώπῳ, ὃς οὐκ οἰκοδομήσει τὸν οἴκον τοῦ ἀδελφοῦ αὐτοῦ ἐν Ἰσραήλ.
\vs{10}Καὶ κληθήσεται τὸ ὄνομα αὐτοῦ ἐν Ἰσραὴλ, οἶκος τοῦ ὑπολυθέντος τὸ ὑπόδημα.

\vs{11}Ἐὰν δὲ μάχωνται ἄνθρωποι ἐπὶ τὸ αὐτό, ἄνθρωπος μετὰ τοῦ ἀδελφοῦ αὐτοῦ, καὶ προσέλθῃ ἡ γυνὴ ἑνὸς αὐτῶν ἐξελέσθαι τὸν ἄνδρα αὐτῆς ἐκ χειρὸς τοῦ τύπτοντος αὐτόν, καὶ ἐκτείνασα τὴν χεῖρα ἐπιλάβηται τῶν διδύμων αὐτοῦ,
\vs{12}ἀποκόψεις τὴν χεῖρα· οὐ φείσεται ὁ ὀφθαλμός σου ἐπʼ αὐτῇ.

\vs{13}Οὐκ ἔσται σοι ἐν τῷ μαρσίππῳ σου στάθμιον καὶ στάθμιον, μέγα ἢ μικρόν.
\vs{14}Οὐκ ἔσται ἐν τῇ οἰκίᾳ σου μέτρον καὶ μέτρον, μέγα ἢ μικρόν.
\vs{15}Στάθμιον ἀληθινὸν καὶ δίκαιον ἔσται σοι, καὶ μέτρον ἀληθινὸν καὶ δίκαιον ἔσται σοι, ἵνα πολυήμερος γένῃ ἐπὶ τῆς γῆς, ἧς Κύριος ὁ Θεός σου δίδωσί σοι ἐν κλήρῳ.
\vs{16}Ὅτι βδέλυγμα Κυρίῳ τῷ Θεῷ σου πας ποιῶν ταῦτα, πας ποιῶν ἄδικον.

\vs{17}Μνήσθητι ὅσα ἐποίησέ σοι Ἀμαλὴκ ἐν τῇ ὁδῷ, ἐκπορευομένου σου ἐκ γῆς Αἰγύπτου,
\vs{18}πῶς ἀντέστη σοι ἐν τῇ ὁδῷ, καὶ ἔκοψέ σου τὴν οὐραγίαν, τοὺς κοπιῶντας ὀπίσω σου, σὺ δὲ ἐπείνας καὶ ἐκοπίας· καὶ οὐκ ἐφοβήθη τὸν Θεόν.
\vs{19}Καὶ ἔσται ἡνίκα ἐὰν καταπαύσῃ σε Κύριος ὁ Θεός σου ἀπὸ πάντων τῶν ἐχθρῶν σου τῶν κύκλῳ σου ἐν τῇ γῇ, ᾗ Κύριος ὁ Θεός σου δίδωσί σοι κληρονομῆσαι, ἐξαλείψεις τὸ ὄνομα Ἀμαλὴκ ἐκ τῆς ὑπὸ τὸν οὐρανόν, καὶ οὐ μὴ ἐπιλάθῃ.

\ch{26}
Καὶ ἔσται ἐὰν εἰσέλθῆς εἰς τὴν γῆν, ἣν Κύριος ὁ Θεός σου δίδωσί σοι κληρονομῆσαι, καὶ κατακληρονομήσῃς αὐτὴν, καὶ κατοικήσῃς ἐπʼ αὐτὴν,
\vs{2}καὶ λήμψῃ ἀπὸ τῆς ἀπαρχῆς τῶν καρπῶν τῆς γῆς σου, ἧς Κύριος ὁ Θεός σου δίδωσί σοι, καὶ ἐμβαλεῖς εἰς κάρταλλον, καὶ πορεύσῃ εἰς τὸν τόπον, ὃν ἂν ἐκλέξηται Κύριος ὁ Θεός σου ἐπικληθῆναι τὸ ὄνομα αὐτοῦ ἐκεῖ.
\vs{3}Καὶ ἐλεύσῃ πρὸς τὸν ἱερέα ὃς ἔσται ἐν ταῖς ἡμέραις ἐκείναις, καὶ ἐρεῖς πρὸς αὐτὸν, ἀναγγέλλω σήμερον Κυρίῳ τῷ Θεῷ μου, ὅτι εἰσελήλυθα εἰς τὴν γῆν, ἣν ὤμοσε Κύριος τοῖς πατράσιν ἡμῶν δοῦναι ἡμῖν.
\vs{4}Καὶ λήψεται ὁ ἱερεὺς τὸν κάρταλλον ἐκ τῶν χειρῶν σου, καὶ θήσει αὐτὸν ἀπέναντι τοῦ θυσιαστηρίου Κυρίου τοῦ Θεοῦ σου.
\vs{5}Καὶ ἀποκριθεὶς ἐρεῖ ἔναντι Κυρίου τοῦ Θεοῦ σου, Συρίαν ἀπέβαλεν ὁ πατήρ μου, καὶ κατέβη εἰς Αἴγυπτον, καὶ παρῴκησεν ἐκεῖ ἐν ἀριθμῷ βραχεῖ, καὶ ἐγένετο ἐκεῖ εἰς ἔθνος μέγα καὶ πλῆθος πολύ.
\vs{6}Καὶ ἐκάκωσαν ἡμᾶς οἱ Αἰγύπτιοι, καὶ ἐταπείνωσαν ἡμᾶς, καὶ ἐπέθηκαν ἡμῖν ἔργα σκληρά·
\vs{7}Καὶ ἀνεβοήσαμεν πρὸς Κύριον τὸν Θεὸν ἡμῶν, καὶ εἰσήκουσε Κύριος τῆς φωνῆς ἡμῶν, καὶ εἴδε τὴν ταπείνωσιν ἡμῶν, καὶ τὸν μόχθον ἡμῶν, καὶ τὸν θλιμμὸν ἡμῶν.
\vs{8}Καὶ ἐξήγαγεν ἡμᾶς Κύριος ἐξ Αἰγύπτου αὐτὸς ἐν ἰσχύϊ αὐτοῦ τῇ μεγάλῃ, καὶ ἐν χειρὶ κραταιᾷ, καὶ βραχίονι ὑψηλῷ, καὶ ἐν ὁράμασι μεγάλοις, καὶ ἐν σημείοις, καὶ ἐν τέρασι.
\vs{9}Καὶ εἰσήγαγεν ἡμᾶς εἰς τὸν τόπον τοῦτον, καὶ ἔδωκεν ἡμῖν τὴν γῆν ταύτην, γῆν ῥέουσαν γάλα καὶ μέλι.
\vs{10}Καὶ νῦν ἰδοὺ ἐνήνοχα τὴν ἀπαρχὴν τῶν γεννμάτων τῆς γῆς, ἧς ἔδωκάς μοι Κύριε, γῆν ῥέουσαν γάλα καὶ μέλι· καὶ ἀφήσεις αὐτὸ ἀπέναντι Κυρίου τοῦ Θεοῦ σου, καὶ προσκυνήσεις ἔναντι Κυρίου τοῦ Θεοῦ σου,
\vs{11}καὶ εὐφρανθήσῃ ἐν πᾶσι τοῖς ἀγαθοῖς, οἷς ἔδωκέ σοι Κύριος ὁ Θεός σου, καὶ ἡ οἰκία σου, καὶ ὁ Λευίτης, καὶ ὁ προσήλυτος ὁ ἐν σοί.

\vs{12}Ἐὰν δὲ συντελέσῃς ἀποδεκατῶσαι πᾶν τὸ ἐπιδέκατον τῶν γενημάτων σου ἐν τῷ ἔτει τῷ τρίτῳ, τὸ δεύτερον ἐπιδέκατον δώσεις τῷ Λευίτῃ καὶ τῷ προσηλύτῳ καὶ τῷ ὀρφανῷ καὶ τῇ χήρᾳ, καὶ φάγονται ἐν ταῖς πόλεσί σου, καὶ εὐφρανθήσονται.

\vs{13}Καὶ ἐρεῖς ἔναντι Κυρίου τοῦ Θεοῦ σου, ἐξεκάθαρα τὰ ἅγια ἐκ τῆς οἰκίας μου, καὶ ἔδωκα αὐτὰ τῷ Λευίτῃ καὶ τῷ προσηλύτῳ καὶ τῷ ὀρφανῷ καὶ τῇ χήρᾳ, κατὰ πάσας τὰς ἐντολὰς ἃς ἐνετείλω μοι· οὐ παρῆλθον τὴν ἐντολήν σου, καὶ οὐκ ἐπελαθόμην.
\vs{14}Καὶ οὐκ ἔφαγον ἐν ὀδύνῃ μου ἀπʼ αὐτῶν, οὐκ ἐκάρπωσα ἀπʼ αὐτῶν εἰς ἀκάθαρτον, οὐκ ἔδωκα ἀπʼ αὐτῶν τῷ τεθνηκότι· ὑπήκουσα τῆς φωνῆς Κυρίου τοῦ Θεοῦ ἡμῶν, ἐποίησα καθὰ ἐνετείλω μοι.
\vs{15}Κάτιδε ἐκ τοῦ οἴκου τοῦ ἁγίου σου ἐκ τοῦ οὐρανοῦ, καὶ εὐλόγησον τὸν λαόν σου τὸν Ἰσραὴλ, καὶ τὴν γῆν ἣν ἔδωκας αὐτοῖς, καθὰ ὤμοσας τοῖς πατράσιν ἡμῶν, δοῦναι ἡμῖν γῆν ῥέουσαν γάλα καὶ μέλι.

\vs{16}Ἐν τῇ ἡμέρᾳ ταύτῃ Κύριος ὁ Θεός σου ἐνετείλατό σοι ποιῆσαι πάντα τὰ δικαιώματα καὶ τὰ κρίματα· καὶ φυλάξεσθε καὶ ποιήσετε αὐτὰ ἐξ ὅλης τῆς καρδίας ὑμῶν, καὶ ἐξ ὅλης τῆς ψυχῆς ὑμῶν.
\vs{17}Τὸν Θεὸν εἵλου σήμερον εἶναί σου Θεόν, καὶ πορεύεσθαι ἐν πάσαις ταῖς ὁδοῖς αὐτοῦ, καὶ φυλάσσεσθαι τὰ δικαιώματα καὶ τὰ κρίματα, καὶ ὑπακούειν τῆς φωνῆς αὐτοῦ.
\vs{18}Καὶ Κύριος εἵλατό σε σήμερον γενέσθαι σε αὐτῷ λαὸν περιούσιον, καθάπερ εἶπε, φυλάττειν τὰς ἐντολὰς αὐτοῦ,
\vs{19}καὶ εἶναί σε ὑπεράνω πάντων τῶν ἐθνῶν, ὡς ἐποίησέ σε ὀνομαστὸν καὶ καύχημα καὶ δοξαστόν, εἶναί σε λαὸν ἅγιον Κυρίῳ τῷ Θεῷ σου, καθὼς ἐλάλησε.

\ch{27}
Καὶ προσέταξε Μωυσῆς καὶ ἡ γερουσία Ἰσραὴλ, λέγων, φυλάσσεσθε πάσας τὰς ἐντολάς ταύτας, ὅσας ἐγὼ ἐντέλλομαι ὑμῖν σήμερον.
\vs{2}Καὶ ἔσται ᾗ ἂν ἡμέρᾳ διαβῆτε τὸν Ἰορδάνην εἰς τὴν γῆν, ἣν Κύριος ὁ Θεός σου δίδωσί σοι, καὶ στήσεις σεαυτῷ λίθους μεγάλους, καὶ κονιάσεις αὐτοὺς κονίᾳ.
\vs{3}Καὶ γράψεις ἐπὶ τῶν λίθων τούτων πάντας τούς λόγους τοῦ νόμου τούτου, ὠς ἂν διαβῆτε τὸν Ἰορδάνην, ἡνίκα ἂν εἰσέλθητε εἰς τὴν γῆν, ἣν Κύριος ὁ Θεὸς τῶν πατέρων σου δίδωσί σοι, γῆν ῥέουσαν γάλα καὶ μέλι, ὃν τρόπον εἶπε Κύριος ὁ Θεὸς τῶν πατέρων σού σοί.
\vs{4}Καὶ ἔσται ὡς ἂν διαβῆτε τὸν Ἰορδάνην, στήσετε τοὺς λίθους τούτους, οὓς ἐγὼ ἐντέλλομαί σοι σήμερον, ἐν ὄρει Γαιβάλ, καὶ κονιάσεις αὐτοὺς κονίᾳ.
\vs{5}Καὶ οἰκοδομήσεις ἐκεῖ θυσιαστήριον Κυρίῳ τῷ Θεῷ σου, θυσιαστήριον ἐκ λίθων· οὐκ ἐπιβαλεῖς ἐπʼ αὐτὸ σίδηρον·
\vs{6}λίθους ὁλοκλήρους οἰκοδομήσεις θυσιαστήριον Κυρίῳ τῷ Θεῷ σου, καὶ ἀνοίσεις ἐπʼ αὐτὸ ὁλοκαυτώματα Κυρίῳ τῷ Θεῷ σου.
\vs{7}Καὶ θύσεις ἐκεῖ θυσίαν σωτηρίου· καὶ φάγῃ, καὶ ἐμπλησθήσῃ, καὶ εὐφρανθήσῃ ἔναντι Κυρίου τοῦ Θεοῦ σου.
\vs{8}Καὶ γράψεις ἐπὶ τῶν λίθων πάντα τὸν νόμον τοῦτον σαφῶς σφόδρα.

\vs{9}Καὶ ἐλάλησε Μωυσῆς καὶ οἱ ἱερεῖς οἱ Λευῖται παντὶ Ἰσραὴλ, λέγοντες, σιώπα καὶ ἄκουε Ἰσραήλ· ἐν τῇ ἡμέρᾳ ταύτῃ γέγονας εἰς λαὸν Κυρίῳ τῷ Θεῷ σου,
\vs{10}καὶ εἰσακούσῃ τῆς φωνῆς Κυρίου τοῦ Θεοῦ σου, καὶ ποιήσεις πάσας τὰς ἐντολὰς αὐτοῦ, καὶ τὰ δικαιώματα αὐτοῦ, ὅσα ἐγὼ ἐντέλλομαί σοι σήμερον.

\vs{11}Καὶ ἐνετείλατο Μωυσῆς τῷ λαῷ ἐν τῇ ἡμέρᾳ ἐκείνῃ, λέγων,
\vs{12}οὗτοι στήσονται εὐλογεῖν τὸν λαὸν ἐν ὄρει Γαριζίν διαβάντες τὸν Ἰορδάνην, Συμεὼν, Λευὶ, Ἰουδας, Ἰσσάχαρ, Ἰωσὴφ, καὶ Βενιαμίν.
\vs{13}Καὶ οὗτοι στήσονται ἐπὶ τῆς κατάρας ἐν ὄρει Γαιβάλ, Ῥουβὴν, Γὰδ, καὶ Ἀσὴρ, Ζαβουλὼν, Δὰν, καὶ Νεφθαλί.

\vs{14}Καὶ ἀποκριθέντες ἐροῦσιν οἱ Λευῖται παντὶ Ἰσραὴλ φωνῇ μεγάλῃ,
\vs{15}ἐπικατάρατος ἄνθρωπος ὅστις ποιήσει γλυπτὸν καὶ χωνευτὸν, βδέλυγμα Κυρίῳ, ἔργον χειρῶν τεχνιτῶν, καὶ θήσει αὐτὸ ἐν ἀποκρύφῳ· καὶ ἀποκριθεὶς πᾶς ὁ λαὸς, ἐροῦσι, γένοιτο.
\vs{16}Ἐπικατάρατος ὁ ἀτιμάζων πατέρα αὐτοῦ ἢ μητέρα αὐτοῦ· καὶ ἐροῦσι πᾶς ὁ λαὸς, γένοτιο.
\vs{17}Ἐπικατάρατος ὁ μετατιθεὶς ὅρια τοῦ πλησὶον· καὶ ἐροῦσι πᾶς ὁ λαός, γένοιτο.
\vs{18}Ἐπικατάρατος ὁ πλανῶν τυφλὸν ἐν ὁδῷ· καὶ ἐροῦσι πᾶς ὁ λαός, γένοιτο.
\vs{19}Ἐπικατάρατος ὃς ἂν ἐκκλίνῃ κρίσιν προσηλύτου καὶ ὀρφανοῦ καὶ χήρας· καὶ ἐροῦσι πᾶς ὁ λαός, γένοιτο.
\vs{20}Ἐπικατάρατος ὁ κοιμώμενος μετὰ γυναικὸς τοῦ πατρὸς αὐτοῦ, ὅτι ἀπεκάλυψε συγκάλυμμα τοῦ πατρὸς αὐτοῦ· καὶ ἐροῦσι πᾶς ὁ λαὸς, γένοιτο.
\vs{21}Ἐπικατάρατος ὁ κοιμώμενος μετὰ παντὸς κτήνους· καὶ ἐροῦσι πᾶς ὁ λαὸς, γένοιτο.
\vs{22}Ἐπικατάρατος ὁ κοιμώμενος μετὰ ἀδελφῆς ἐκ πατρὸς ἢ μητρὸς αὐτοῦ. καὶ ἐροῦσι πᾶς ὁ λαός, γένοιτο.
\vs{23}Ἐπικατάρατος ὁ κοιμώμενος μετὰ νύμφης αὐτοῦ· καὶ ἐροῦσι πᾶς ὁ λαός, γένοιτο· ἐπικατάρατος ὁ κοιμώμενος μετὰ τῆς ἀδελφῆς τῆς γυναικὸς αὐτοῦ· καὶ ἐροῦσι πᾶς ὁ λαὸς, γένοιτο.
\vs{24}Ἐπικατάρατος ὁ τύπτων τὸν πλησίον δόλῳ· καὶ ἐροῦσι πᾶς ὁ λαός, γένοιτο.
\vs{25}Ἐπικατάρατος ὃς ἂν λάβῃ δῶρα πατάξαι ψυχὴν αἵματος ἀθῴου· καὶ ἐροῦσι πᾶς ὁ λαὸς, γένοιτο.
\vs{26}Ἐπικατάρατος πᾶς ἄνθρωπος ὃς οὐκ ἐμμένει ἐν πᾶσι τοῖς λόγοις τοῦ νόμου τούτου ποιῆσαι αὐτούς· καὶ ἐροῦσι πᾶς ὁ λαὸς, γένοιτο.

\ch{28}
Καὶ ἔσται ἐὰν ἀκοῇ ἀκούσῃς τῆς φωνῆς Κυρίου τοῦ Θεοῦ σου, φυλάσσειν καὶ ποιεῖν πάσας τὰς ἐντολὰς ταύτας, ἃς ἐγὼ ἐντέλλομαί σοι σήμερον, καὶ δώσει σε Κύριος ὁ Θεός σου ὑπεράνω ἐπὶ πάντα τὰ ἔθνη τῆς γῆς,
\vs{2}καὶ ἥξουσιν ἐπὶ σὲ πᾶσαι αἱ εὐλογίαι αὗται, καὶ εὑρήσουσί σε· ἐὰν ἀκοῇ ἀκούσῃς τῆς φωνῆς Κυρίου τοῦ Θεοῦ σου,
\vs{3}εὐλογημένος σὺ ἐν πόλει, καὶ εὐλογημένος σὺ ἐν ἀγρῶ.
\vs{4}Εὐλογημένα τὰ ἔκγονα τῆς κοιλίας σου, καὶ τὰ γεννήματα τῆς γης σου, καὶ τὰ βουκόλια τῶν βοῶν σου, καὶ τὰ ποίμνια τῶν προβάτων σου.
\vs{5}Εὐλογημέναι αἱ ἀποθῆκαί σου, καὶ τὰ ἐγκαταλείμματά σου.
\vs{6}Εὐλογημένος σὺ ἐν τῷ εἰσπορεύεσθαί σε, καὶ εὐλογημένος σὺ ἐν τῷ ἐκπορεύεσθαί σε.

\vs{7}Παραδῷ Κύριος ὁ Θεός σου τοὺς ἐχθρούς σου τοὺς ἀνθεστηκότας σοι συντετριμμένους πρὸ προσώπου σου· ὁδῷ μιᾷ ἐξελεύσονται πρὸς σέ, καὶ ἐν ἑπτὰ ὁδοῖς φεύξονται ἀπὸ προσώπου σου.
\vs{8}Ἀποστείλαι Κύριος ἐπὶ σὲ τὴν εὐλογίαν ἐν τοῖς ταμείοις σου, καὶ ἐπὶ πάντα οὗ ἂν ἐπιβάλῃς τὴν χεῖρά σου, ἐπὶ τῆς γῆς, ἧς Κύριος ὁ Θεός σου δίδωσί σοι.
\vs{9}Ἀναστήσαι σε Κύριος ἑαυτῷ λαὸν ἅγιον, ὃν τρόπον ὤμοσε τοῖς πατράσι σου· ἐὰν ἀκούσῃς τῆς φωνῆς Κυρίου τοῦ Θεοῦ σου, καὶ πορευθῇς ἐν πάσαις ταῖς ὁδοῖς αὐτοῦ,
\vs{10}καὶ ὄψονταί σε πάντα τὰ ἔθνη τῆς γῆς, ὅτι τὸ ὄνομα Κυρίου ἐπικέκληταί σοι, καὶ φοβηθήσονταί σε.
\vs{11}Καὶ πληθυνεῖ σε Κύριος ὁ Θεός σου εἰς ἀγαθὰ ἐν τοῖς ἐκγόνοις τῆς κοιλίας σου, καὶ ἐπὶ τοῖς ἐκγόνοις τῶν κτηνῶν σου, καὶ ἐπὶ τοῖς γεννήμασι τῆς γῆς σου, ἐπὶ τῆς γῆς σου ἧς ὤμοσε Κύριος τοῖς πατράσι σου δοῦναί σοι.

\vs{12}Ἀνοίξαι σοι Κύριος τὸν θησαυρὸν αὐτοῦ τὸν ἀγαθὸν, τὸν οὐρανὸν, δοῦναι τὸν ὑετὸν τῇ γῇ σου ἐπὶ καιροῦ· εὐλογῆσαι πάντα τὰ ἔργα τῶν χειρῶν σου· καὶ δανειεῖς ἔθνεσι πολλοῖς, σὺ δὲ οὐ δανειῇ. καὶ ἄρξεις σὺ ἐθνῶν πολλῶν, σοῦ δὲ οὐκ ἄρξουσι.
\vs{13}Καταστήσαι σε Κύριος ὁ Θεός σου εἰς κεφαλὴν καὶ μὴ εἰς οὐρὰν, καὶ ἔσῃ τότε ἐπάνω καὶ οὐκ ἔσῃ ὑποκάτω, ἐὰν ἀκούσῃς τῆς φωνῆς Κυρίου τοῦ Θεοῦ σου, ὅσα ἐγὼ ἐντέλλομαί σοι σήμερον φυλάσσειν.
\vs{14}Οὐ παραβήσῃ ἀπὸ πασῶν τῶν ἐντολῶν, ὧν ἐγὼ ἐντέλλομαί σοι σήμερον δεξιὰ οὐδὲ ἀριστερὰ, πορεύσεθαι ὀπίσω θεῶν ἑτέρων λατρεύειν αὐτοῖς.

\vs{15}Καὶ ἔσται ἐὰν μὴ εἰσακούσῃς τῆς φωνῆς Κυρίου τοῦ Θεοῦ σου, φυλάσσεσθαι πάσας τὰς ἐντολὰς αὐτοῦ, ὅσας ἐγὼ ἐντέλλομαί σοι σήμερον, καὶ ἐλεύσονται ἐπὶ σὲ πᾶσαι αἱ κατάραι αὗται, καὶ καταλήψονταί σε.
\vs{16}Ἐπικατάρατος σὺ ἐν πόλει, καὶ ἐπικατάρατος σὺ ἐν ἀγρῷ.
\vs{17}Ἐπικατάρατοι αἱ ἀποθῆκαί σου, καὶ τὰ ἐγκαταλείμματά σου.
\vs{18}Ἐπικατάρατα τὰ ἔκγονα τῆς κοιλίας σου, καὶ τὰ γεννήματα τῆς γῆς σου, τὰ βουκόλια τῶν βοῶν σου, καὶ τὰ ποίμνια τῶν προβάτων σου·
\vs{19}Ἐπικατάρατος σὺ ἐν τῷ εἰσπορεύεσθαί σε, καὶ ἐπικατάρατος σὺ ἐν τῷ ἐκπορεύεσθαί σε.

\vs{20}Ἀποστείλαι Κύριος ἐπὶ σὲ τὴν ἔνδειαν καὶ τὴν ἐκλιμίαν καὶ τὴν ἀνάλωσιν ἐπὶ πάντα οὗ ἐὰν ἐπιβάλῃς τὴν χεῖρά σου, ἕως ἂν ἐξολοθρεύσῃ σε, καὶ ἕως ἂν ἀπολέσῃ σε ἐν τάχει διὰ τὰ πονηρὰ ἐπιτηδεύματά σου, διότι ἐνκατέλιπές με.
\vs{21}Προσκολλήσαι Κύριος εἰς σὲ τὸν θάνατον, ἕως ἂν ἐξαναλώσῃ σε ἀπὸ τῆς γῆς, εἰς ἣν εἰσπορεύῃ ἐκεῖ κληρονομῆσαι αὐτήν.
\vs{22}Πατάξαι σε Κύριος ἐν ἀπορίᾳ, καὶ πυρετῷ, καὶ ῥίγει, καὶ ἐρεθισμῷ, καὶ ἀνεμοφθορίᾳ, καὶ τῇ ὤχρᾳ, καὶ καταδιώξονταί σε ἕως ἂν ἀπολέσωσί σε.
\vs{23}Καὶ ἔσται σοι ὁ οὐρανὸς ὁ ὑπὲρ κεφαλῆς σου χαλκοῦς, καὶ ἡ γῆ ἡ ὑποκάτω σου σιδηρᾶ.
\vs{24}Δῴη Κύριος ὁ Θεός σου τὸν ὑετὸν τῇς γῇς σου κονιορτὸν, καὶ χοῦς ἐκ τοῦ οὐρανοῦ καταβήσεται, ἕως ἂν ἐκτρίψῃ σε, καὶ ἕως ἂν ἀπολέσῃ σε ἐν τάχει.
\vs{25}Δῴη σε Κύριος ἐπὶ κοπὴν ἐναντίον τῶν ἐχθρῶν· ἐν ὁδῷ μιᾷ ἐξελεύσῃ πρὸς αὐτοὺς, καὶ ἐν ἑπτὰ ὁδοῖς φεύξῃ ἀπὸ προσώπου αὐτῶν· καὶ ἔσῃ διασπορὰ ἐν πάσαις βασιλείαις τῆς γῆς.
\vs{26}Καὶ ἔσονται οἱ νεκροὶ ὑμῶν κατάβρωμα τοῖς πετεινοῖς τοῦ οὐρανοῦ, καὶ τοῖς θηρίοις τῆς γῆς, καὶ οὐκ ἔσται ὁ ἐκφοβῶν.
\vs{27}Πατάξαι σε Κύριος ἕλκει Αἰγυπτίῳ εἰς τὴν ἕδραν, καὶ ψώρᾳ ἀγρίᾳ, καὶ κνήφῃ, ὥστε μὴ δύνασθαί σε ἰαθῆναι.
\vs{28}Πατάξαι σε Κύριος παραπληξίᾳ, καὶ ἀορασίᾳ, καὶ ἐκστάσει διανοίας.
\vs{29}Καὶ ἔσῃ ψηλαφῶν μεσημβρίας, ὡσεί τις ψηλαφήσαι τυφλὸς ἐν τῷ σκότει, καὶ οὐκ εὐοδώσει τὰς ὁδούς σου· καὶ ἔσῃ τότε ἀδικούμενος, καὶ διαρπαζόμενος πάσας τὰς ἡμέρας, καὶ οὐκ ἔσται ὁ βοηθῶν.

\vs{30}Γυναῖκα λήψῃ, καὶ ἀνὴρ ἕτερος ἕξει αὐτήν· οἰκίαν οἰκοδομήσεις, καὶ οὐκ οἰκήσεις ἐν αὐτῇ· ἀμπελῶνα φυτεύσεις, καὶ οὐ μὴ τρυγήσεις αὐτόν.
\vs{31}Ὁ μόσχος σου ἐσφαγμένος ἐναντίον σου, καὶ οὐ φάγῃ ἐξ αὐτοῦ· ὁ ὄνος σου ἡρπασμένος ἀπὸ σοῦ, καὶ οὐκ ἀποδοθήσεταί σοι· τὰ πρόβατά σου δεδομένα τοῖς ἐχθροῖς σου,
\vs{32}καὶ οὐκ ἔσται σοι ὁ βοηθῶν. Οἱ υἱοί σου καὶ αἱ θυγατέρες σου δεδομέναι ἔθνει ἑτέρῳ, καὶ οἱ ὀφθαλμοί σου βλέψονται σφακελίζοντες εἰς αὐτὰ· οὐκ ἰσχύσει ἡ χείρ σου.
\vs{33}Τὰ ἐκφόρια τῆς γῆς σου, καὶ πάντας τοὺς πόνους σου φάγεται ἔθνος, ὃ οὐκ ἐπίστασαι· καὶ ἔσῃ ἀδικούμενος καὶ τεθραυσμένος πάσας τὰς ἡμέρας.
\vs{34}Καὶ ἔσῃ παράπληκτος διὰ τὰ ὁράματα τῶν ὀφθαλμῶν σου, ἃ βλέψῃ.

\vs{35}Πατάξαι σε Κύριος ἐν ἕλκει πονηρῷ ἐπὶ τὰ γόνατα καὶ ἐπὶ τὰς κνήμας, ὥστε μὴ δύνασθαι ἰαθῆναί σε ἀπὸ ἴχνους τῶν ποδῶν σου ἕως τῆς κορυφῆς σου.

\vs{36}Ἀπαγάγοι Κύριός σε καὶ τοὺς ἄρχοντάς σου, οὓς ἂν καταστήσῃς ἐπὶ σεαυτὸν, ἐπʼ ἔθνος ὃ οὐκ ἐπίστασαι σὺ καὶ οἱ πατέρες σου, καὶ λατρεύσεις ἐκεῖ θεοῖς ἑτέροις ξύλοις καὶ λίθοις.
\vs{37}Καὶ ἔσῃ ἐκεῖ ἐν αἰνίγματι καὶ παραβολῇ καὶ διηγήματι ἐν πᾶσι τοῖς ἔθνεσιν, εἰς οὓς ἂν ἀπαγάγῃ σε Κύριος ἐκεῖ.

\vs{38}Σπέρμα πολὺ ἐξοίσεις εἰς τὸ πεδίον, καὶ ὀλίγα εἰσοίσεις, ὅτι κατέδεται αὐτὰ ἡ ἀκρίς·
\vs{39}Ἀμπελῶνα φυτεύσεις καὶ κατεργᾷ, καὶ οἶνον οὐ πίεσαι οὐδὲ εὐφρανθήσῃ ἐξ αὐτοῦ, ὅτι καταφάγεται αὐτὰ ὁ σκώληξ.
\vs{40}Ἐλαῖαι ἔσονταί σοι ἐν πᾶσι τοὶς ὁρίοις σου, καὶ ἔλαιον οὐ χρίσῃ, ὅτι ἐκρυήσεται ἡ ἐλαία σου.
\vs{41}Υἱοὺς καὶ θυγατέρας γεννήσεις καὶ οὐκ ἔσονται· ἀπελεύσονται γὰρ ἐν αἰχμαλωσίᾳ.
\vs{42}Πάντα τὰ ξύλινά σου, καὶ τὰ γεννήματα τῆς γῆς σου ἐξαναλώσει ἡ ἐρισύβη.
\vs{43}Ὁ προσήλυτος ὅς ἐστιν ἐν σοὶ, ἀναβήσεται ἄνω ἄνω, σὺ δὲ καταβήσῃ κάτω κάτω.
\vs{44}Οὗτος δανειεῖ σοι, σὺ δὲ τούτῳ οὐ δανειεῖς· οὗτος ἔσται κεφαλὴ, σὺ δὲ ἔσῃ οὐρά.

\vs{45}Καὶ ἐλεύσονται ἐπὶ σὲ πᾶσαι αἱ κατάραι αὗται, καὶ καταδιώξονταί σε, καὶ καταλήψονταί σε, ἔως ἂν ἐξολοθρεύσῃ σε, καὶ ἕως ἂν ἀπολέσῃ σε· ὅτι οὐκ εἰσήκουσας τῆς φωνῆς Κυρίου τοῦ Θεοῦ σου, φυλάξαι τὰς ἐντολὰς αὐτοῦ, καὶ τὰ δικαιώματα ὅσα ἐνετείλατό σοι.
\vs{46}Καὶ ἔσται ἐν σοὶ σημεῖα, καὶ τέρατα ἐν τῷ σπέρματί σου ἕως τοῦ αἰῶνος,
\vs{47}ἀνθʼ ὧν οὐκ ἐλάτρευσας Κυρίῳ τῷ Θεῷ σου ἐν εὐφροσύνῃ καὶ ἀγαθῇ διανοίᾳ διὰ τὸ πλῆθος πάντων.

\vs{48}Καὶ λατρεύσεις τοῖς ἐχθροῖς σου, οὓς ἐπαποστελεῖ Κύριος ἐπὶ σέ, ἐν λιμῷ, καὶ ἐν δίψει, καὶ ἐν γυμνότητι, καὶ ἐν ἐκλείψει πάντων· καὶ ἐπιθήσῃ κλοιὸν σιδηροῦν ἐπὶ τὸν τράχηλόν σου, ἕως ἂν ἐξολοθρεύσῃ σε.
\vs{49}Ἐπάξει ἐπὶ σὲ Κύριος ἔθνος μακρόθεν ἀπʼ ἐσχάτου τῆς γῆς ὡσεὶ ὅρμημα ἀετοῦ, ἔθνος ὃ οὐκ ἀκούσῃ τῆς φωνῆς αὐτοῦ,
\vs{50}ἔθνος ἀναιδὲς προσώπῳ, ὅστις οὐ θαυμάσει πρόσωπον πρεσβύτου, καὶ νέον οὐκ ἐλεήσει.
\vs{51}Καὶ κατέδεται τὰ ἔγκονα τῶν κτηνῶν σου, καὶ τὰ γεννήματα τῆς γῆς σου, ὥστε μὴ καταλιπεῖν σοι σῖτον, οἶνον, ἔλαιον, τὰ βουκόλια τῶν βοῶν σου, καὶ τὰ ποίμνια τῶν προβάτων σου, ἕως ἂν ἀπολέσῃ σε.
\vs{52}Καὶ ἐκτρίψῃ σε ἐν ταῖς πόλεσί σου, ἕως ἂν καθαιρεθῶσι τὰ τείχη τὰ ὑψηλὰ καὶ τὰ ὀχυρὰ, ἐφʼ οἷς σὺ πέποιθας ἐπʼ αὐτοῖς, ἐν πάσῃ τῇ γῇ σου· καὶ θλίψει σε ἐν ταῖς πόλεσί σου, αἷς ἔδωκέ σοι.
\vs{53}Καὶ φαγῇ τὰ ἔκγονα τῆς κοιλίας σου, κρέα υἱῶν σου καὶ θυγατέρων σου, ὅσα ἔδωκέ σοι, ἐν τῇ στενοχωρίᾳ σου καὶ ἐν τῇ θλίψει σου, ἡ· θλίψει σε ὁ ἐχθρός σου.

\vs{54}Ὁ ἁπαλὸς ὁ ἐν σοὶ καὶ ὁ τρυφερὸς σφόδρα, βασκανεῖ τῷ ὀφθαλμῷ αὐτοῦ τὸν ἀδελφὸν αὐτοῦ, καὶ τὴν γυναῖκα τὴν ἐν τῷ κόλπῳ αὐτοῦ, καὶ τὰ καταλελειμμένα τέκνα, ἃ ἂν καταλειφθῇ
\vs{55}αὐτῷ, ὥστε δοῦναι ἑνὶ αὐτῶν ἀπὸ τῶν σαρκῶν τῶν τέκνων αὐτοῦ, ὧν ἂν κατέσθῃ διὰ τὸ μὴ καταλειφθῆναι αὐτῷ οὐδὲν ἐν τῇ στενοχωρίᾳ σου, καὶ ἐν τῇ θλίψει σου, ἡ· ἂν θλίψωσί σε οἱ ἐχθροί σου ἐν πάσαις ταῖς πόλεσί σου.

\vs{56}Καὶ ἡ ἁπαλὴ ἐν ὑμῖν καὶ ἡ τρυφερά, ἧς οὐχὶ πεῖραν ἔλαβεν ὁ ποὺς αὐτῆς βαίνειν ἐπὶ τῆς γῆς διὰ τὴν τρυφερότητα καὶ διὰ τὴν ἁπαλότητα, βασκανεῖ τῷ ὀφθαλμῷ αὐτῆς τὸν ἄνδρα αὐτῆς τὸν ἐν κόλπῳ αὐτῆς, καὶ τὸν υἱὸν καὶ τὴν θυγατέρα αὐτῆς,
\vs{57}καὶ τὸ κόριον αὐτῆς τὸ ἐξελθὸν διὰ τῶν μηρῶν αὐτῆς, καὶ τὸ τέκνον αὐτῆς ὃ ἐὰν τέκῃ· καταφάγεται γὰρ αὐτὰ διὰ τὴν ἔνδειαν πάντων κρυφῇ ἐν τῇ στενοχωρίᾳ σου, καὶ ἐν τῇ θλίψει σου, ἡ· θλίψει σε ὁ ἐχθρός σου ἐν ταῖς πόλεσί σου,
\vs{58}ἐὰν μὴ εἰσακούσῃς ποιεῖν πάντα τὰ ῥήματα τοῦ νόμου τούτου, τὰ γεγραμμένα ἐν τῷ βιβλίῳ τούτῳ, φοβεῖσθαι τὸ ὄνομα τὸ ἔντιμον τὸ θαυμαστὸν τοῦτο, ΚΥΡΙΟΝ τὸν ΘΕΟΝ σου.
\vs{59}Καὶ παραδοξάσει Κύριος τὰς πληγάς σου, καὶ τὰς πληγὰς τοῦ σπέρματός σου, πληγὰς μεγάλας καὶ θαυμαστὰς, καὶ νόσους πονηρὰς καὶ πιστάς.
\vs{60}Καὶ ἐπιστρέψει πᾶσαν τὴν ὀδύνην Αἰγύπτου τὴν πονηρὰν, ἣν διευλαβοῦ ἀπὸ προσώπου αὐτῶν, καὶ κολληθήσονται ἐν σοί.
\vs{61}Καὶ πᾶσαν μαλακίαν, καὶ πᾶσαν πληγὴν τὴν μὴ γεγραμμένην, καὶ πᾶσαν τὴν γεγραμμένην ἐν τῷ βιβλίῳ τοῦ νόμου τούτου, ἐπάξει Κύριος ἐπὶ σέ, ἕως ἂν ἐξολοθρεύσῃ σε.
\vs{62}Καὶ καταλειφθήσεσθε ἐν ἀριθμῷ βραχεῖ, ἀνθʼ ὧν ὅτι ἦτε ὡσεὶ τὰ ἄστρα τοῦ οὐρανοῦ τῷ πλήθει, ὅτι οὐκ εἰσήκουσας τῆς φωνῆς Κυρίου τοῦ Θεοῦ σου.

\vs{63}Καὶ ἔσται ὃν τρόπον εὐφράνθη Κύριος ἐφʼ ὑμῖν εὖ ποιῆσαι ὑμᾶς, καὶ πληθῦναι ὑμᾶς, οὕτως εὐφρανθήσεται Κύριος ἐφʼ ὑμῖν ἐξολοθρεῦσαι ὑμᾶς· καὶ ἐξαρθήσεσθε ἐν τάχει ἀπὸ τῆς γῆς, εἰς ἣν εἰσπορεύῃ ἐκεῖ κληρονομῆσαι αὐτήν.
\vs{64}Καὶ διασπερεῖ σε Κύριος ὁ Θεός σου εἰς πάντα τὰ ἔθνη, ἀπʼ ἄκρου τῆς γῆς ἕως ἄκρου τῆς γῆς, καὶ δουλεύσεις ἐκεῖ θεοῖς ἑτέροις, ξύλοις καὶ λίθοις, οὓς οὐκ ἠπίστω σὺ καὶ οἱ πατέρες σου.
\vs{65}Ἀλλὰ καὶ ἐν τοῖς ἔθνεσιν ἐκείνοις οὐκ ἀναπαύσει σε, οὐδʼ οὐ μὴ γένηται στάσις τῷ ἴχνει τοῦ ποδός σου· καὶ δώσει σοι Κύριος ἐκεῖ καρδίαν ἑτέραν ἀπειθοῦσαν, καὶ ἐκλείποντας ὀφθαλμοὺς, καὶ τηκομένην ψυχήν.
\vs{66}Καὶ ἔσται ἡ ζωή σου κρεμαμένη ἀπέναντι τῶν ὀφθαλμῶν σου· καὶ φοβηθήσῃ ἡμέρας καὶ νυκτὸς, καὶ οὐ πιστεύσεις τῇ ζωῇ σου.
\vs{67}Τὸ πρωῒ ἐρεῖς, πῶς ἂν γένοιτο ἑσπέρα· καὶ τὸ ἑσπέρας ἐρεῖς, πῶς ἂν γένοιτο πρωΐ· ἀπὸ τοῦ φόβου τῆς καρδίας σου ἃ φοβηθήσῃ, καὶ ἀπὸ τῶν ὁραμάτων τῶν ὀφθαλμῶν σου ὧν ὄψῃ.
\vs{68}Καὶ ἀποστρέψει σε Κύριος εἰς Αἴγυπτον ἐν πλοίοις, ἐν τῇ ὁδῷ ἡ· εἶπα, οὐ προσθήσῃ ἔτι ἰδεῖν αὐτήν· καὶ πραθήσεσθε ἐκεῖ τοῖς ἐχθροῖς ὑμῶν εἰς παῖδας καὶ παιδίσκας, καὶ οὐκ ἔσται ὁ κτώμενος.

\vs{69}Οὗτοι οἱ λόγοι τῆς διαθήκης, οὓς ἐνετείλατο Κύριος Μωυσῇ στῆσαι τοῖς υἱοῖς Ἰσραὴλ ἐν γῇ Μωὰβ, πλὴν τῆς διαθήκης ἧς διέθετο αὐτοῖς ἐν Χωρήβ.

\ch{29}
Καὶ ἐκάλεσε Μωυσῆς πάντας τοὺς υἱοὺς Ἰσραὴλ, καὶ εἶπε πρὸς αὐτοὺς, ὑμεῖς ἑωράκατε πάντα ὅσα ἐποίησε Κύριος ἐν γῇ Αἰγύπτῳ ἐνώπιον ὑμῶν Φαραὼ καὶ τοῖς θεράπουσιν αὐτοῦ, καὶ πάσῃ τῇ γῇ αὐτοῦ,
\vs{2}τοὺς πειρασμοὺς τοὺς μεγάλους οὓς ἑωράκασιν οἱ ὀφθαλμοί σου, τὰ σημεῖα καὶ τὰ τέρατα τὰ μεγάλα ἐκεῖνα.
\vs{3}Καὶ οὐκ ἔδωκε Κύριος ὁ Θεὸς ὑμῖν καρδίαν εἰδέναι, καὶ ὀφθαλμοὺς βλέπειν, καὶ ὦτα ἀκούειν ἕως τῆς ἡμέρας ταύτης.
\vs{4}Καὶ ἤγαγεν ὑμᾶς τεσσαράκοντα ἔτη ἐν τῇ ἐρήμῳ· οὐκ ἐπαλαιώθη τὰ ἱμάτια ὑμῶν, καὶ τὰ ὑποδήματα ὑμῶν οὐ κατετρίβη ἀπὸ τῶν ποδῶν ὑμῶν.
\vs{5}Ἄρτον οὐκ ἐφάγετε, οἶνον καὶ σίκερα οὐκ ἐπίετε, ἵνα γνῶτε ὃτι Κύριος ὁ Θεὸς ὑμῶν ἐγώ.
\vs{6}Καὶ ἤλθετε ἕως τοῦ τόπου τούτου· καὶ ἑξῆλθε Σηὼν βασιλεὺς Ἐσεβὼν, καὶ Ὢγ βασιλεὺς Βασὰν εἰς συνάντησιν ἡμῖν ἐν πολέμῳ.
\vs{7}Καὶ ἐπατάξαμεν αὐτοὺς, καὶ ἐλάβομεν τὴν γῆν αὐτῶν, καὶ ἔδωκα αὐτὴν ἐν κλήρῳ τῷ Ῥουβὴν, καὶ τῷ Γαδδὶ, καὶ τῷ ἡμίσει φυλῆς Μανασσῆ.
\vs{8}Καὶ φυλάξεσθε ποιεῖν πάντας τοὺς λόγους τῆς διαθήκης ταύτης, ἵνα συνῆτε πάντα ὅσα ποιήσετε.

\vs{9}Ὑμεῖς ἑστήκατε πάντες σήμερον ἐναντίον Κυρίου τοῦ Θεοῦ ὑμῶν, οἱ ἀρχίφυλοι ὑμῶν, καὶ ἡ γερουσία ὑμῶν, καὶ οἱ κριταὶ ὑμῶν, καὶ οἱ γραμματοεισαγωγεῖς ὑμῶν, πᾶς ἀνὴρ Ἰσραήλ,
\vs{10}αἱ γυναῖκες ὑμῶν, καὶ τὰ ἔκγονα ὑμῶν καὶ ὁ προσήλυτος ὁ ἐν μέσῳ τῆς παρεμβολῆς ὑμῶν, ἀπὸ ξυλοκόπου ὑμῶν καὶ ἕως ὑδροφόρου ὑμῶν,
\vs{11}παρελθεῖν ἐν τῇ διαθήκῃ Κυρίου τοῦ Θεοῦ ὑμῶν, καὶ ἐν ταῖς ἀραῖς αὐτοῦ, ὅσα Κύριος ὁ Θεός σου διατίθεται πρὸς σὲ σήμερον·
\vs{12}ἵνα στήσῃ σε αὐτῷ εἰς λαὸν, καὶ αὐτὸς ἔσται σου Θεὸς, ὃν τρόπον εἶπέ σοι, καὶ ὃν τρόπον ὤμοσε τοῖς πατράσι σου Ἁβραὰμ καὶ Ἰσαὰκ καὶ Ἰακώβ.
\vs{13}Καὶ οὐχ ὑμῖν μόνοις ἐγὼ διατίθεμαι τὴν διαθήκην ταύτην καὶ τὴν ἀρὰν ταύτην,
\vs{14}ἀλλὰ καὶ τοῖς ὧδε οὖσι μεθʼ ὑμῶν σήμερον ἐναντίον Κυρίου τοῦ Θεοῦ ὑμῶν, καὶ τοῖς μὴ οὖσι μεθʼ ὑμῶν ὧδε σήμερον.

\vs{15}Ὅτι ὑμεῖς οἴδατε πῶς κατῳκήσαμεν ἐν γῇ Αἰγύπτῳ, ὡς παρήλθομεν ἐν μέσῳ τῶν ἐθνῶν οὓς παρήλθετε.
\vs{16}Καὶ ἴδετε τὰ βδελύγματα αὐτῶν, καὶ τὰ εἴδωλα αὐτῶν, ξύλον καὶ λίθον, ἀργύριον καὶ χρυσίον, ἅ ἐστι παρʼ αὐτοῖς.
\vs{17}Μή τις ἐστὶν ἐν ὑμῖν ἀνὴρ, ἢ γυνὴ, ἢ πατριὰ, ἢ φυλὴ, τινὸς ἡ διάνοια ἐξέκλινεν ἀπὸ Κυρίου τοῦ Θεου ὑμῶν, πορευθέντες λατρεύειν τοῖς θεοῖς τῶν ἐθνῶν ἐκείνων· μή τις ἐστὶν ἐν ὑμῖν ῥίζα ἄνω φύουσα ἐν χολῇ καὶ πικρίᾳ·
\vs{18}Καὶ ἔσται ἐὰν ἀκούσῃ τὰ ῥήματα τῆς ἀρᾶς ταύτης, καὶ ἐπιφημίσηται ἐν τῇ καρδίᾳ αὐτοῦ, λέγων, ὅσιά μοι γένοιτο, ὅτι ἐν τῇ ἀποπλανήσει τῆς καρδίας μου πορεύσομαι, ἵνα μὴ συναπολέσῃ ὁ ἁμαρτωλὸς τὸν ἀναμάρτητον·
\vs{19}Οὐ μὴ θελήσει ὁ Θεὸς εὐϊλατεῦσαι αὐτῷ, ἀλλʼ ἢ τότε ἐκκαυθήσεται ὀργὴ Κυρίου καὶ ὁ ζῆλος αὐτοῦ ἐν τῷ ἀνθρώπῳ ἐκείνῳ· καὶ κολληθήσονται ἐν αὐτῷ πᾶσαι αἱ ἀραὶ τῆς διαθήκης ταύτης, αἱ γεγραμμέναι ἐν τῷ βιβλίῳ τούτῳ· καὶ ἐξαλείψει Κύριος τὸ ὄνομα αὐτοῦ ἐκ τῆς ὑπὸ τὸν οὐρανόν.
\vs{20}Καὶ διαστελεῖ αὐτὸν Κύριος εἰς κακὰ ἐκ πάντων υἱῶν Ἰσραὴλ, κατὰ πάσας τὰς ἀρὰς τῆς διαθήκης τὰς γεγραμμένας ἐν τῷ βιβλίῳ τοῦ νόμου τούτου.

\vs{21}Καὶ ἐροῦσιν ἡ γενεὰ ἡ ἑτέρα οἱ υἱοὶ ὑμῶν, οἳ ἀναστήσονται μεθʼ ὑμᾶς, καὶ ὁ ἀλλότριος ὃς ἂν ἔλθῃ ἐκ γῆς μακρόθεν, καὶ ὄψονται τὰς πληγὰς τῆς γῆς ἐκείνης καὶ τὰς νόσους αὐτῆς, ἃς ἀπέστειλε Κύριος ἐπʼ αὐτὴν,
\vs{22}θεῖον καὶ ἅλα κατακεκαυμένον· πᾶσα ἡ γῆ αὐτῆς οὐ σπαρήσεται, οὐδὲ ἀνατελεῖ, οὐδὲ μὴ ἀναβῇ ἐπʼ αὐτὴν πᾶν χλωρόν. ὥσπερ κατεστράφη Σόδομα καὶ Γόμοῤῥα, Ἀδαμὰ καὶ Σεβωῒμ, ἃς κατέστρεψε Κύριος ἐν θυμῷ καὶ ὀργῇ·
\vs{23}Καὶ ἐροῦσι πάντα τὰ ἔθνη, διατί ἐποίησε Κύριος οὕτω τῇ γῇ ταύτῃ; τίς ὁ θυμὸς τῆς ὀργῆς ὁ μέγας οὗτος;
\vs{24}Καὶ ἐροῦσιν, ὅτι κατέλιπον τὴν διαθήκην Κυρίου τοῦ Θεοῦ τῶν πατέρων αὐτῶν, ἃ διέθετο τοῖς πατράσιν αὐτῶν, ὅτε ἐξήγαγεν αὐτοὺς ἐκ γῆς Αἰγύπτου,
\vs{25}καὶ πορεύθεντες ἐλάτρευσαν θεοῖς ἑτέροις, οὓς οὐκ ἠπίσταντο, οὐδὲ διένειμεν αὐτοῖς·
\vs{26}Καὶ ὠργίσθη θυμῷ Κύριος ἐπὶ τὴν γῆν ἐκείνην ἐπαγαγεῖν ἐπʼ αὐτὴν κατὰ πάσας τὰς κατάρας τὰς γεγραμμένας ἐν τῷ βιβλίῳ τοῦ νόμου τούτου.
\vs{27}Καὶ ἐξῇρεν αὐτοὺς Κύριος ἀπὸ τῆς γῆς αὐτῶν ἐν θυμῷ καὶ ὀργῇ καὶ παροξυσμῷ μεγάλῳ σφόδρα, καὶ ἐξέβαλεν αὐτοὺς εἰς γῆν ἑτέραν ὡσεὶ νῦν.

\vs{28}Τὰ κρυπτὰ Κυρίῳ τῷ Θεῷ ἡμῶν, τὰ δὲ φανερὰ ἡμῖν καὶ τοῖς τέκνοις ἡμῶν εἰς τὸν αἰῶνα, ποιεῖν πάντα τὰ ῥήματα τοῦ νόμου τούτου.

\ch{30}
Καὶ ἔσται ὡς ἂν ἔλθωσιν ἐπὶ σὲ πάντα τὰ ῥήματα ταῦτα, ἡ εὐλογία καὶ ἡ κατάρα, ἣν ἔδωκα πρὸ προσώπου σου; καὶ δέξῃ εἰς τὴν καρδίαν σου ἐν πᾶσι τοῖς ἔθνεσιν, οὗ ἐὰν σε διασκορπίσῃ σε Κύριος ἐκεῖ,
\vs{2}καὶ ἐπιστραφήσῃ ἐπὶ Κύριον τὸν Θεόν σου, καὶ εἰσακούσῃ τῆς φωνῆς αὐτοῦ κατὰ πάντα ὅσα ἐγὼ ἐντέλλομαί σοι σήμερον, ἐξ ὅλης τῆς καρδίας σου, καὶ ἐξ ὅλης τῆς ψυχῆς σου,
\vs{3}καὶ ἰάσεται Κύριος τὰς ἁμαρτίας σου, καὶ ἐλεήσει σε, καὶ πάλιν συνάξει σε ἐκ πάντων τῶν ἐθνῶν, εἰς οὓς διεσκόρπισέ σε Κύριος ἐκεῖ.
\vs{4}Ἐὰν ἠ· ἡ διασπορά σου ἀπʼ ἄκρου τοῦ οὐρανοῦ ἕως ἄκρου τοῦ οὐρανοῦ, ἐκεῖθεν συνάξει σε Κύριος ὁ Θεός σου, καὶ ἐκεῖθεν λήψεταί σε Κύριος ὁ Θεός σου.
\vs{5}Καὶ εἰσάξει σε ὁ Θεός σου ἐκεῖθεν εἰς τὴν γῆν ἣν ἐκληρονόμησαν οἱ πατέρες σου, καὶ κληρονομήσεις αὐτήν· καὶ εὖ σε ποιήσει, καὶ πλεοναστόν σε ποιήσει ὑπὲρ τοὺς πατέρας σου.
\vs{6}Καὶ περικαθαριεῖ Κύριος τὴν καρδίαν σου, καὶ τὴν καρδίαν τοῦ σπέρματός σου, ἀγαπᾷν Κύριον τὸν Θεόν σου ἐξ ὅλης τῆς καρδίας σου, καὶ ἐξ ὅλης τῆς ψυχῆς σου, ἵνα ζῇς σύ.

\vs{7}Καὶ δώσει Κύριος ὁ Θεός σου τὰς ἀρὰς ταύτας ἐπὶ τοὺς ἐχθρούς σου, καὶ ἐπὶ τοὺς μισοῦντάς σε, οἳ ἐδίωξάν σε.
\vs{8}Καὶ σὺ ἐπιστραφήσῃ καὶ εἰσακούσῃ τῆς φωνῆς Κυρίου τοῦ Θεοῦ σου, καὶ ποιήσεις τὰς ἐντολὰς αὐτοῦ, ὅσας ἐγὼ ἐντέλλομαί σοι σήμερον.
\vs{9}Καὶ εὐλογήσει σε Κύριος ὁ Θεός σου ἐν παντὶ ἔργῳ τῶν χειρῶν σου, ἐν τοῖς ἐκγόνοις τῆς κοιλίας σου, καὶ ἐν τοῖς ἐκγόνοις τῶν κτηνῶν σου, καὶ ἐν τοῖς γεννήμασι τῆς γῆς σου, ὅτι ἐπιστρέψει Κύριος ὁ Θεός σου εὐφρανθῆναὶ ἐπὶ σοὶ εἰς ἀγαθὰ, καθότι εὐφράνθη ἐπὶ τοῖς πατράσι σου·
\vs{10}Ἐὰν εἰσακούσῃς τῆς φωνῆς Κυρίου τοῦ Θεοῦ σου, φυλάσσεσθαι τὰς ἐντολὰς αὐτοῦ, καὶ τὰ δικαιώματα αὐτοῦ, καὶ τὰς κρίσεις αὐτοῦ τὰς γεγραμμένας ἐν τῷ βιβλίῳ τοῦ νόμου τούτου· ἐὰν ἐπιστραφῇς ἐπὶ Κύριον τὸν Θεόν σου ἐξ ὅλης τῆς καρδίας σου, καὶ ἐξ ὅλης τῆς ψυχῆς σου.
\vs{11}Ὅτι ἡ ἐντολὴ αὕτη ἣν ἐγὼ ἐντέλλομαί σοι σήμερον, οὐχ ὑπέρογκός ἐστιν, οὐδέ μακρὰν ἀπὸ σοῦ ἐστιν.
\vs{12}Οὐκ ἐν τῷ οὐρανῷ ἄνω ἐστὶ, λέγων, τίς ἀναβήσεται ἡμῖν εἰς τὸν οὐρανὸν, καὶ λήψεται ἡμῖν αὐτὴν, καὶ ἀκούσαντες αὐτὴν ποιήσομεν;
\vs{13}Οὐδὲ πέραν τῆς θαλάσσης ἐστί, λέγων, τίς διαπεράσει ἡμῖν εἰς τὸ πέραν τῆς θαλάσσης, καὶ λάβῃ ἡμῖν αὐτὴν, καὶ ἀκουστὴν ἡμῖν ποιήσῃ αὐτὴν, καὶ ποιήσομεν;
\vs{14}Ἐγγύς σου ἐστὶ τὸ ῥῆμα σφόδρα ἐν τῷ στόματί σου, καὶ ἐν τῇ καρδίᾳ σου, καὶ ἐν ταῖς χερσί σου ποιεῖν αὐτό.

\vs{15}Ἰδοὺ δέδωκα πρὸ προσώπου σου σήμερον τὴν ζωὴν καὶ τὸν θάνατον, τὸ ἀγαθὸν καὶ τὸ κακόν.
\vs{16}Ἐὰν εἰσακούσῃς τὰς ἐντολὰς Κυρίου τοῦ Θεοῦ σου, ἃς ἐγὼ ἐντέλλομαί σοι σήμερον, ἀγαπᾷν Κύριον τὸν Θεόν σου, πορεύεσθαι ἐν πάσαις ταῖς ὁδοῖς αὐτοῦ, καὶ φυλάσσεσθαι τὰ δικαιώματα αὐτοῦ, καὶ τὰς κρίσεις αὐτοῦ, καὶ ζήσεσθε, καὶ πολλοὶ ἔσεσθε, καὶ εὐλογήσει σε Κύριος ὁ Θεός σου ἐν πάσῃ τῇ γῇ, εἰς ἣν εἰσπορεύῃ ἐκεῖ κληρονομῆσαι αὐτήν.
\vs{17}Καὶ ἐὰν μεταστῇ ἡ καρδία σου, καὶ μὴ εἰσακούσῃς, καὶ πλανηθεὶς προσκυνήσῃς θεοῖς ἑτέροις καὶ λατρεύσῃς αὐτοῖς,
\vs{18}ἀναγγέλλω σοι σήμερον, ὅτι ἀπωλείᾳ ἀπολεῖσθε, καὶ οὐ μὴ πολυήμεροι γένησθε ἐπὶ τῆς γῆς, εἰς ἣν ὑμεῖς διαβαίνετε τὸν Ἰορδάνην ἐκεῖ κληρονομῆσαι αὐτήν.

\vs{19}Διαμαρτύρομαι ὑμῖν σήμερον τόν τε οὐρανὸν καὶ τὴν γῆν, τὴν ζωὴν καὶ τὸν θάνατον δέδωκα πρὸ προσώπου ὑμῶν, τὴν εὐλογίαν καὶ τὴν κατάραν· ἔκλεξαι τὴν ζωὴν σὺ, ἵνα ζήσῃς σὺ καὶ τὸ σπέρμα σου,
\vs{20}ἀγαπᾷν Κύριον τὸν Θεόν σου, εἰσακούειν τῆς φωνῆς αὐτοῦ, καὶ ἔχεσθαι αὐτοῦ· ὅτι τοῦτο ἡ ζωή σου καὶ ἡ μακρότης τῶν ἡμερῶν σου, τὸ κατοικεῖν ἐπὶ τῆς γῆς, ἧς ὤμοσε Κύριος τοῖς πατράσι σου Ἁβραὰμ καὶ Ἰσαὰκ καὶ Ἰακὼβ δοῦναι αὐτοῖς.

\ch{31}
Καὶ συνετέλεσε Μωυσῆς λαλῶν πάντας τοὺς λόγους τούτους πρὸς πάντας υἱοὺς Ἰσραὴλ,
\vs{2}καὶ εἶπε πρὸς αὐτοὺς, ἐκατὸν καὶ εἴκοσι ἐτῶν ἐγώ εἰμι σήμερον· οὐ δυνήσομαι ἔτι εἰσπορεύεσθαι καὶ ἐκπορεύεσθαι· Κύριος δὲ εἶπε πρὸς μὲ, οὐ διαβήσῃ τὸν Ιορδάνην τοῦτον.
\vs{3}Κύριος ὁ Θεός σου ὁ προπορευόμενος πρὸ προσώπου σου, οὗτος ἐξολοθρεύσει τὰ ἔθνη ταῦτα ἀπὸ προσώπου σου, καὶ κατακληρονομήσεις αὐτούς· καὶ Ἰησοῦς ὁ προπορευόμενος πρὸ προσώπου σου, καθὰ ἐλάλησε Κύριος.
\vs{4}Καὶ ποιήσει Κύριος ὁ Θεός σου αὐτοῖς καθὼς ἐποίησε Σηὼν καὶ Ὢγ δυσὶ βασιλεῦσι τῶν Ἀμοῤῥαίων, οἳ ἦσαν πέραν τοῦ Ἰορδάνου, καὶ τῇ γῇ αὐτῶν, καθότι ἐξωλόθρευσεν αὐτούς.
\vs{5}Καὶ παρέδωκεν αὐτοὺς Κύριος ὑμῖν· καὶ ποιήσετε αὐτοῖς, καθότι ἐνετειλάμην ὑμῖν.
\vs{6}Ἀνδρίζου καὶ ἴσχυε, μὴ φοβοῦ, μηδὲ δειλιάσῃς, μηδὲ πτοηθῇς ἀπὸ προσώπου αὐτῶν· ὅτι Κύριος ὁ Θεός σου ὁ προπορευόμενος μεθʼ ὑμῶν ἐν ὑμῖν, οὔτε μή σε ἀνῇ, οὔτε μή σε ἐγκαταλίπῃ.
\vs{7}Καὶ ἐκάλεσε Μωυσῆς Ἰησοῦν, καὶ εἶπεν αὐτῷ ἔναντι παντὸς Ἰσραὴλ, ἀνδρίζου καὶ ἴσχυε, σὺ γὰρ εἰσελεύσῃ πρὸ προσώπου τοῦ λαοῦ τούτου εἰς τὴν γῆν ἣν ὤμοσε Κύριος τοῖς πατράσιν ὑμῶν δοῦναι αὐτοῖς, καὶ σὺ κατακληρονομήσεις αὐτοῖς.
\vs{8}Καὶ Κύριος ὁ συμπορευόμενος μετὰ σοῦ, οὐκ ἀνήσει σε, οὐδὲ μὴ σε ἐγκαταλίπῃ· μή φοβοῦ, μηδὲ δειλία.

\vs{9}Καὶ ἔγραψε Μωυσῆς τὰ ῥήματα τοῦ νόμου τούτου εἰς βιβλίον, καὶ ἔδωκε τοῖς ἱερεῦσι τοῖς υἱοῖς Λευὶ τοῖς αἴρουσι τὴν κιβωτὸν τῆς διαθήκης Κυρίου, καὶ τοῖς πρεσβυτέροις τῶν υἱῶν Ἰσραήλ.

\vs{10}Καὶ ἐνετείλατο Μωυσῆς αὐτοῖς ἐν τῇ ἡμέρᾳ ἐκείνῃ, λέγων, μετὰ ἑπτὰ ἔτη ἐν καιρῷ ἐνιαυτοῦ ἀφέσεως ἐν ἑορτῇ σκηνοπηγίας,
\vs{11}ἐν τῷ συμπορεύεσθαι πάντα Ἰσραὴλ ὀφθῆναι ἐνώπιον Κύριου τοῦ Θεοῦ ὑμῶν, ἐν τῷ τόπῳ ᾧ ἂν ἐκλέξηται Κύριος, ἀναγνώσεσθε τὸν νόμον τοῦτον ἐναντίον παντὸς Ἰσραὴλ εἰς τὰ ὦτα αὐτῶν,
\vs{12}ἐκκλησιάσας τὸν λαὸν, τοὺς ἄνδρας καὶ τὰς γυναῖκας καὶ τὰ ἔκγονα καὶ τὸν προσήλυτον τὸν ἐν ταῖς πόλεσιν ὑμῶν, ἵνʼ ἀκούσωσι, καὶ ἵνα μάθωσι φοβεῖσθαι Κύριον τὸν Θεὸν ὑμῶν· καὶ ἀκούσονται ποιεῖν πάντας τοὺς λόγους τοῦ νόμου τούτου.
\vs{13}Καὶ οἱ υἱοὶ αὐτῶν οἳ οὐκ οἴδασιν, ἀκούσονται, καὶ μαθήσονται φοβεῖσθαι Κύριον τὸν Θεόν σου πάσας τὰς ἡμέρας ὅσας αὐτοὶ ζῶσιν ἐπὶ τῆς γῆς, εἰς ἣν ὑμεῖς διαβαίνετε τὸν Ἰορδάνην ἐκεῖ κληρονομῆσαι αὐτήν.

\vs{14}Καὶ εἶπε Κύριος πρὸς Μωυσῆν, ἰδοὺ ἐγγίκασιν αἱ ἡμέραι τοῦ θανάτου σου· κάλεσον Ἰησοῦν, καὶ στῆτε παρὰ τὰς θύρας τῆς σκηνῆς τοῦ μαρτυρίου, καὶ ἐντελοῦμαι αὐτῷ· καὶ ἐπορεύθη Μωυσῆς καὶ Ἰησοῦς εἰς τὴν σκηνὴν τοῦ μαρτυρίου, καὶ ἔστησαν παρὰ τὰς θύρας τῆς σκηνῆς τοῦ μαρτυρίου.
\vs{15}Καὶ κατέβη Κύριος ἐν νεφέλῃ, καὶ ἔστη παρὰ τὰς θύρας τῆς σκηνῆς τοῦ μαρτυρίου· καὶ ἔστη ὁ στύλος τῆς νεφέλης παρὰ τὰς θύρας τῆς σκηνῆς τοῦ μαρτυρίου.
\vs{16}Καὶ εἶπε Κύριος πρὸς Μωυσῆν, ἰδοὺ σὺ κοιμᾷ μετὰ τῶν πατέρων σου, καὶ ἀναστὰς οὗτος ὁ λαὸς ἐκπορνεύσει ὀπίσω θεῶν ἀλλοτρίων τῆς γῆς, εἰς ἣν οὗτος εἰσπορεύεται, καὶ καταλείψουσί με, καὶ διασκεδάσουσι τὴν διαθήκην μου, ἣν διεθέμην αὐτοῖς.
\vs{17}Καὶ ὀργισθήσομαι θυμῷ εἰς αὐτοὺς ἐν τῇ ἡμέρᾳ ἐκείνῃ, καὶ καταλείψω αὐτοὺς, καὶ ἀποστρέψω τὸ πρόσωπόν μου ἀπʼ αὐτῶν, καὶ ἔσται κατάβρωμα· καὶ εὑρήσουσιν αὐτὸν κακὰ πολλὰ καὶ θλίψεις· καὶ ἐρεῖ ἐν τῇ ἡμέρᾳ ἐκείνῃ, διότι οὐκ ἔστι Κύριος ὁ Θεός μου ἐν ἐμοὶ, εὕροσάν με τὰ κακὰ ταῦτα.
\vs{18}Ἐγὼ δὲ ἀποστροφῇ ἀποστρέψω τὸ πρόσωπόν μου ἀπʼ αὐτῶν ἐν τῇ ἡμέρᾳ ἐκείνῃ, διὰ πάσας τὰς κακίας ἃς ἐποίησαν, ὅτι ἀπέστρεψαν ἐπὶ θεοὺς ἀλλοτρίους.

\vs{19}Καὶ νῦν γράψατε τὰ ῥήματα τῆς ᾠδῆς ταύτης, καὶ διδάξατε αὐτὴν τοὺς υἱοὺς Ἰσραὴλ, καὶ ἐμβαλεῖτε αὐτὴν εἰς τὸ στόμα αὐτῶν, ἵνα γένηταί μοι ἡ ᾠδὴ αὕτη κατὰ πρόσωπον μαρτυροῦσα ἐν υἱοῖς Ἰσραήλ.
\vs{20}Εἰσάξω γὰρ αὐτοὺς εἰς τὴν γῆν τὴν ἀγαθὴν, ἣν ὤμοσα τοῖς πατράσιν αὐτῶν, δοῦναι αὐτοῖς γῆν ῥέουσαν γάλα καὶ μέλι, καὶ φάγονται, καὶ ἐμπλησθέντες κορήσουσι, καὶ ἐπιστραφήσονται ἐπὶ θεοὺς ἀλλοτρίους, καὶ λατρεύσουσιν αὐτοῖς, καὶ παροξυνοῦσί με, καὶ διασκεδάσουσι τὴν διαθήκην μου.
\vs{21}Καὶ ἀντικαταστήσεται ἡ ᾠδὴ αὕτη κατὰ πρόσωπον μαρτυροῦσα· οὐ γὰρ μὴ ἐπιλησθῇ ἀπὸ στόματος αὐτῶν, καὶ ἀπὸ στόματος τοῦ σπέρματος αὐτῶν· ἐγὼ γὰρ οἶδα τὴν πονηρίαν αὐτῶν, ὅσα ποιοῦσιν ὧδε σήμερον, πρὸ τοῦ εἰσαγαγεῖν με αὐτοὺς εἰς τὴν γῆν τὴν ἀγαθὴν, ἣν ὤμοσα τοῖς πατράσιν αὐτῶν.

\vs{22}Καὶ ἔγραψε Μωυσῆς τὴν ᾠδὴν ταύτην ἐν ἐκείνῃ τῇ ἡμέρᾳ, καὶ ἐδίδαξεν αὐτὴν τοὺς υἱοὺς Ἰσραήλ.
\vs{23}Καὶ ἐνετείλατο Ἰησοῖ, καὶ εἶπεν, ἀνδρίζου καὶ ἴσχυε, σὺ γὰρ εἰσάξεις τοὺς υἱοὺς Ἰσραὴλ εἰς τὴν γῆν, ἣν ὤμοσεν αὐτοῖς Κύριος, καὶ αὐτὸς ἔσται μετὰ σοῦ.

\vs{24}Ἡνίκα δὲ συνετέλεσε Μωυσῆς γράφων πάντας τοὺς λόγους τοῦ νόμου τούτου εἰς βιβλίον ἕως εἰς τέλος,
\vs{25}καὶ ἐνετείλατο τοῖς Λευίταις τοῖς αἴρουσι τὴν κιβωτὸν τῆς διαθήκης Κυρίου, λέγων,
\vs{26}λαβόντες τὸ βιβλίον τοῦ νόμου τούτου, θήσετε αὐτὸ ἐκ πλαγίων τῆς κιβωτοῦ τῆς διαθήκης Κυρίου τοῦ Θεοῦ ὑμῶν· καὶ ἔσται ἐκεῖ ἐν σοὶ εἰς μαρτύριον.
\vs{27}Ὅτι ἐγὼ ἐπίσταμαι τὸν ἐρεθισμόν σου, καὶ τὸν τράχηλόν σου τὸν σκληρόν· ἔτι γὰρ ἐμοῦ ζῶντος μεθʼ ὑμῶν σήμερον, παραπικραίνοντες ἦτε τὰ πρὸς τὸν Θεόν· πῶς οὐχὶ καὶ ἔσχατον τοῦ θανάτου μου;
\vs{28}Ἐκκλησιάσατε πρὸς μὲ τοὺς φυλάρχους ὑμῶν, καὶ τοὺς πρεσβυτέρους ὑμῶν, καὶ τοὺς κριτὰς ὑμῶν, καὶ τοὺς γραμματοεισαγωγεῖς ὑμῶν, ἵνα λαλήσω εἰς τὰ ὦτα αὐτῶν πάντας τοὺς λόγους τούτους· καὶ διαμαρτύρομαι αὐτοῖς τόν τε οὐρανὸν καὶ τὴν γῆν.
\vs{29}Οἶδα γὰρ ὅτι ἔσχατον τῆς τελευτῆς μου ἀνομίᾳ ἀνομήσετε, καὶ ἐκκλινεῖτε ἐκ τῆς ὁδοῦ ἧς ἐνετειλάμην ὑμῖν, καὶ συναντήσεται ὑμῖν τὰ κακὰ ἔσχατον τῶν ἡμερῶν, ὅτι ποιήσετε τὰ πονηρὰ ἐναντίον Κυρίου, παροργίσαι αὐτὸν ἐν τοῖς ἔργοις τῶν χειρῶν ὑμῶν.

\vs{30}Καὶ ἐλάλησε Μωυσῆς εἰς τὰ ὦτα πάσης ἐκκλησίας τὰ ῥήματα τῆς ᾠδῆς ταύτης ἕως εἰς τέλος.

\ch{32}
Πρόσεχε οὐρανὲ, καὶ λαλήσω, καὶ ἀκουέτω ἡ γῆ ῥήματα ἐκ στόματός μου.
\vs{2}Προσδοκάσθω ὡς ὑετὸς τὸ ἀπόφθεγμά μου, καὶ καταβήτω ὡς δρόσος τὰ ῥήματά μου, ὡσεὶ ὄμβρος ἐπʼ ἄγρωστιν, καὶ ὡσεὶ νιφετὸς ἐπὶ χόρτον.
\vs{3}Ὅτι τὸ ὄνομα Κυρίου ἐκάλεσα· δότε μεγαλωσύνην τῷ Θεῷ ἡμῶν.
\vs{4}Θεὸς, ἀληθινὰ τὰ ἔργα αὐτοῦ, καὶ πᾶσαι αἱ ὁδοὶ αὐτοῦ κρίσεις· Θεὸς πιστὸς, καὶ οὐκ ἔστιν ἀδικία· δίκαιος καὶ ὅσιος Κύριος.
\vs{5}Ἡμάρτοσαν οὐκ αὐτῷ τέκνα μωμητά· γενεὰ σκολιὰ καὶ διεστραμμένη.
\vs{6}Ταῦτα Κυρίῳ ἀνταποδίδοτε; οὕτω λαὸς μωρὸς καὶ οὐχὶ σοφός; οὐκ αὐτὸς οὗτός σου πατὴρ ἐκτήσατό σε καὶ ἐποίησέ σε καὶ ἔπλασέ σε;
\vs{7}Μνήσθητε ἡμέρας αἰῶνος, σύνετε ἔτη γενεῶν γενεαῖς. ἐπερώτησον τὸν πατέρα σου καὶ ἀναγγελεῖ σοι, τοὺς πρεσβυτέρους σου καὶ ἐροῦσί σοι.

\vs{8}Ὅτε διεμέριζεν ὁ ὕψιστος ἔθνη, ὡς διέσπειρεν υἱοὺς Αδὰμ, ἔστησεν ὅρια ἐθνῶν κατὰ ἀριθμὸν ἀγγέλων Θεοῦ.
\vs{9}Καὶ ἐγενήθη μερὶς Κυρίου λαὸς αὐτοῦ Ἰακώβ· σχοίνισμα κληρονομίας αὐτοῦ Ἰσραήλ.
\vs{10}Αὐτάρκησεν αὐτὸν ἐν τῇ ἐρήμῳ, ἐν δίψει καύματος ἐν γῇ ἀνύδρῳ· ἐκύκλωσεν αὐτὸν καὶ ἐπαίδευσεν αὐτὸν, καὶ διεφύλαξεν αὐτὸν, ὡς κόρην ὀφθαλμοῦ·
\vs{11}Ὡς ἀετὸς σκεπάσαι νοσσιὰν αὐτοῦ, καὶ ἐπὶ τοῖς νοσσοῖς αὐτοῦ ἐπεπόθησε, διεὶς τὰς πτέρυγας αὐτοῦ ἐδέξατο αὐτοὺς, καὶ ἀνέλαβεν αὐτοὺς ἐπὶ τῶν μεταφρένων αὐτοῦ.
\vs{12}Κύριος μόνος ἦγεν αὐτοὺς, οὐκ ἦν μετʼ αὐτῶν θεὸς ἀλλότριος.
\vs{13}Ἀνεβίβασεν αὐτοὺς ἐπὶ τὴν ἰσχὺν τῆς γῆς· ἐψώμισεν αὐτοὺς γεννήματα ἀγρῶν· ἐθήλασαν μέλι ἐκ πέτρας, καὶ ἔλαιον ἐκ στερεᾶς πέτρας.
\vs{14}Βούτυρον βοῶν, καὶ γάλα προβάτων, μετὰ στέατος ἀρνῶν καὶ κριῶν, υἱῶν ταύρων καὶ τράγων, μετὰ στέατος νεφρῶν πυροῦ, καὶ αἷμα σταφυλῆς ἔπιεν οἶνον.
\vs{15}Καὶ ἔφαγεν Ἰακὼβ καὶ ἐνεπλήσθη, καὶ ἀπελάκτισεν ὁ ἠγαπημένος· ἐλιπάνθη, ἐπαχύνθη, ἐπλατύνθη, καὶ ἐγκατέλιπε τὸν Θεὸν τὸν ποιήσαντα αὐτὸν, καὶ ἀπέστη ἀπὸ Θεοῦ σωτῆρος αὐτοῦ.

\vs{16}Παρώξυνάν με ἐπʼ ἀλλοτρίοις· ἐν βδελύγμασιν αὐτῶν παρεπίκρανάν με.
\vs{17}Εθυσαν δαιμονίοις, καὶ οὐ Θεῷ· θεοῖς οἷς οὐκ ᾔδεισαν· καινοὶ πρόσφατοι ἥκασιν, οὓς οὐκ ᾔδεισαν οἱ πατέρες αὐτῶν.
\vs{18}Θεὸν τὸν γεννήσαντά σε ἐγκατέλιπες, καὶ ἐπελάθου Θεοῦ τοῦ τρέφοντός σε.

\vs{19}Καὶ εἶδε Κύριος, καὶ ἐζήλωσε· καὶ παρωξύνθη διʼ ὀργὴν υἱῶν αὐτοῦ καὶ θυγατέρων,
\vs{20}καὶ εἶπεν, ἀποστρέψω τὸ πρόσωπόν μου ἀπʼ αὐτῶν, καὶ δείξω τί ἔσται αὐτοῖς ἐπʼ ἐσχάτων ἡμερῶν· ὅτι γενεὰ ἐξεστραμμένη ἐστίν, υἱοὶ οἷς οὐκ ἔστι πίστις ἐν αὐτοῖς.

\vs{21}Αὐτοὶ παρεζήλωσάν με ἐπʼ οὐ Θεῷ, παρώξυνάν με ἐν τοῖς εἰδώλοις αὐτῶν· κᾀγὼ παραζηλώσω αὐτοὺς ἐπʼ οὐκ ἔθνει, ἐπὶ ἔθνει ἀσυνέτῳ παροργιῶ αὐτούς.
\vs{22}Ὅτι πῦρ ἐκκέκαυται ἐκ τοῦ θυμοῦ μου, καυθήσεται ἕως ᾄδου κάτω· καταφάγεται γῆν καὶ τὰ γεννήματα αὐτῆς· φλέξει θεμέλια ὀρέων.
\vs{23}Συνάξω εἰς αὐτοὺς κακὰ, καὶ τὰ βέλη μου συμπολεμήσω εἰς αὐτούς.
\vs{24}Τηκόμενοι λιμῷ καὶ βρώσει ὀρνέων, καὶ ὀπισθότονος ἀνίατος· ὀδόντας θηρίων ἐπαποστελῶ εἰς αὐτοὺς, μετὰ θυμοῦ συρόντων ἐπὶ γῆν.
\vs{25}Ἔξωθεν ἀτεκνώσει αὐτοὺς μάχαιρα, καὶ ἐκ τῶν ταμιείων, φόβος· νεανίσκος σὺν παρθένῳ, θηλάζων μετὰ καθεστηκότος πρεσβύτου.
\vs{26}Εἶπα, διασπερῶ αὐτοὺς, παύσω δὲ ἐξ ἀνθρώπων τὸ μνημόσυνον αὐτῶν.
\vs{27}Εἰ μὴ διʼ ὀργὴν ἐχθρῶν, ἵνα μὴ μακροχρονίσωσι, ἵνα μὴ συνεπιθῶνται οἱ ὑπεναντίοι· μὴ εἴπωσιν, ἡ χεὶρ ἡμῶν ἡ ὑψηλὴ, καὶ οὐχὶ Κύριος, ἐποίησε ταῦτα πάντα.

\vs{28}Ἔθνος ἀπολωλεκὸς βουλήν ἐστι, καὶ οὐκ ἔστιν ἐν αὐτοῖς ἐπιστήμη.
\vs{29}Οὐκ ἐφρόνησαν συνιέναι· ταῦτα καταδεξάσθωσαν εἰς τὸν ἐπιόντα χρόνον.
\vs{30}Πῶς διώξεται εἷς χιλίους, καὶ δύο μετακινήσουσι μυριάδας, εἰ μὴ ὁ Θεὸς ἀπέδοτο αὐτοὺς, καὶ Κύριος παρέδωκεν αὐτούς;
\vs{31}Ὅτι οὐκ εἰσὶν ὡς ὁ Θεὸς ἡμῶν οἱ θεοὶ αὐτῶν· οἱ δὲ ἐχθροὶ ἡμῶν ἀνόητοι.
\vs{32}Ἐκ γὰρ ἀμπέλου Σοδόμων ἡ ἄμπελος αὐτῶν, καὶ ἡ κληματὶς αὐτῶν ἐκ Γομόῤῥας· σταφυλὴ αὐτῶν σταφυλὴ χολῆς, βότρυς πικρίας αὐτοῖς.
\vs{33}Θυμὸς δρακόντων ὁ οἶνος αὐτῶν, καὶ θυμὸς ἀσπίδων ἀνίατος.
\vs{34}Οὐκ ἰδοὺ ταῦτα συνῆκται παρʼ ἐμοὶ, καὶ ἐσφράγισται ἐν τοῖς θησαυροῖς μου;
\vs{35}Ἐν ἡμέρᾳ ἐκδικήσεως ἀνταποδώσω, ὅταν σφαλῇ ὁ ποὺς αὐτῶν· ὅτι ἐγγὺς ἡμέρα ἀπωλίας αὐτοῖς, καὶ πάρεστιν ἕτοιμα ὑμῖν.
\vs{36}Ὅτι κρινεῖ Κύριος τὸν λαὸν αὐτοῦ, καὶ ἐπὶ τοῖς δούλοις αὐτοῦ παρακληθήσεται· εἶδε γὰρ παραλελυμένους αὐτοὺς, καὶ ἐκλελοιπότας ἐν ἐπαγωγῇ, καὶ παρειμένους·
\vs{37}Καὶ εἶπε Κύριος, ποῦ εἰσιν οἱ θεοὶ αὐτῶν, ἐφʼ οἷς ἐπεποίθεισαν ἐπʼ αὐτοῖς,
\vs{38}ὧν τὸ στέαρ τῶν θυσιῶν αὐτῶν ἠσθίετε, καὶ ἐπίνετε τὸν οἶνον τῶν σπονδῶν αὐτῶν; ἀναστήτωσαν καὶ βοηθησάτωσαν ὑμῖν καὶ γενηθήτωσαν ὑμῖν σκεπασταί.
\vs{39}Ἴδετε ἴδετε ὅτι ἐγώ εἰμι, καὶ οὐκ ἔστι Θεὸς πλὴν ἐμοῦ· ἐγὼ ἀποκτείνω, καὶ ζῇν ποιήσω· πατάξω, κἀγὼ ἰάσομαι· καὶ οὐκ ἔστιν ὃς ἐξελεῖται ἐκ τῶν χειρῶν μου.
\vs{40}Ὅτι ἀρῶ εἰς τὸν οὐρανὸν τὴν χεῖρά μου, καὶ ὀμοῦμαι τὴν δεξιάν μου· καὶ ἐρῶ, ζῶ ἐγὼ εἰς τὸν αἰῶνα·
\vs{41}Ὅτι παροξυνῶ ὡς ἀστραπὴν τὴν μάχαιράν μου, καὶ ἀνθέξεται κρίματος ἡ χείρ μου, καὶ ἀποδώσω δίκην τοῖς ἐχθροῖς, καὶ τοῖς μισοῦσί με ἀνταποδώσω.
\vs{42}Μεθύσω τὰ βέλη μου ἀφʼ αἵματος, καὶ ἡ μάχαιρά μου φάγεται κρέα ἀφʼ αἵματος τραυματιῶν καὶ αἰχμαλωσίας ἀπὸ κεφαλῆς ἀρχόντων ἐχθρῶν.

\vs{43}Εὐφράνθητε οὐρανοὶ ἅμα αὐτῷ, καὶ προσκυνησάτωσαν αὐτῷ πάντες ἄγγελοι Θεοῦ· εὐφράνθητε ἔθνη μετὰ τοῦ λαοῦ αὐτοῦ, καὶ ἐνισχυσάτωσαν αὐτῷ πάντες υἱοὶ Θεοῦ, ὅτι τὸ αἷμα τῶν υἱῶν αὐτοῦ ἐκδικᾶται· καὶ ἐκδικήσει καὶ ἀνταποδώσει δίκην τοῖς ἐχθροῖς, καὶ τοῖς μισοῦσιν ἀνταποδώσει· καὶ ἐκκαθαριεῖ Κύριος τὴν γῆν τοῦ λαοῦ αὐτοῦ.

\vs{44}Καὶ ἔγραψε Μωυσῆς τὴν ᾠδὴν ταύτην ἐν τῇ ἡμέρᾳ ἐκείνῃ, καὶ ἐδίδαξεν αὐτὴν τοὺς υἱοὺς Ἰσραήλ· καὶ εἰσῆλθε Μωυσῆς, καὶ ἐλάλησε πάντας τοὺς λόγους τοῦ νόμου τούτου εἰς τὰ ὦτα τοῦ λαοῦ, αὐτὸς καὶ Ἰησοῦς ὁ τοῦ Ναυή.
\vs{45}Καὶ ἐξετέλεσε Μωυσῆς λαλῶν παντὶ Ἰσραήλ.
\vs{46}Καὶ εἶπε πρὸς αὐτοὺς, προσέχετε τῇ καρδίᾳ ἐπὶ πάντας τοὺς λόγους τούτους, οὓς ἐγὼ διαμαρτύρομαι ὑμῖν σήμερον, ἃ ἐντελεῖσθε τοῖς υἱοῖς ὑμῶν, φυλάσσειν καὶ ποιεῖν πάντας τοὺς λόγους τοῦ νόμου τούτου.
\vs{47}Ὅτι οὐχὶ λόγος κενὸς οὗτος ὑμῖν· ὅτι αὕτη ἡ ζωὴ ὑμῶν, καὶ ἕνεκεν τοῦ λόγου τούτου μακροημερεύσετε ἐπὶ τῆς γῆς, εἰς ἣν ὑμεῖς διαβαίνετε τὸν Ἰορδάνην ἐκεῖ κληρονομῆσαι.
\vs{48}Καὶ ἐλάλησε Κύριος πρὸς Μωυσῆν ἐν τῇ ἡμέρᾳ ταύτῃ, λέγων,
\vs{49}ἀνάβηθι εἰς τὸ ὄρος τὸ Ἀβαρίμ, τοῦτο ὄρος Ναβαῦ ὅ ἐστιν ἐν γῇ Μωὰβ κατὰ πρόσωπον Ἰεριχὼ, καὶ ἴδε τὴν γῆν Χαναάν, ἣν ἐγὼ δίδωμι τοῖς υἱοῖς Ἰσραὴλ,
\vs{50}καὶ τελεύτα ἐν τῷ ὄρει εἰς ὃ ἀναβαίνεις ἐκεῖ, καὶ προστέθητι πρὸς τὸν λαόν σου· ὃν τρόπον ἀπέθανεν Ἀαρὼν ὁ ἀδελφός σου ἐν Ὢρ τῷ ὄρει, καὶ προσετέθη πρὸς τὸν λαὸν αὐτοῦ.
\vs{51}Ὅτι ἠπειθήσατε τῷ ῥήματί μου ἐν τοῖς υἱοῖς Ἰσραὴλ ἐπὶ τοῦ ὕδατος ἀντιλογίας Κάδης ἐν τῇ ἐρήμῳ Σὶν, διότι οὐχ ἡγιάσατέ με ἐν τοῖς υἱοῖς Ἰσραήλ.
\vs{52}Ἀπέναντι ὄψει τὴν γῆν, καὶ ἐκεῖ οὐκ εἰσελεύσῃ.

\ch{33}
Καὶ αὕτη ἡ εὐλογία ἣν ηὐλόγησε Μωυσῆς ἄνθρωπος τοῦ Θεοῦ τοὺς υἱοὺς Ἰσραὴλ πρὸ τῆς τελευτῆς αὐτοῦ.
\vs{2}Καὶ εἶπε, Κύριος ἐκ Σινὰ ἥκει, καὶ ἐπέφανεν ἐκ Σηεὶρ ἡμῖν, καὶ κατέσπευσεν ἐξ ὄρους Φαρὰν, σὺν μυρίασι Κάδης, ἐκ δεξιῶν αὐτοῦ ἄγγελοι μετʼ αὐτοῦ.
\vs{3}Καὶ ἐφείσατο τοῦ λαοῦ αὐτοῦ, καὶ πάντες οἱ ἡγιασμένοι ὑπὸ τὰς χεῖράς σου· καὶ οὗτοι ὑπὸ σέ εἰσί· καὶ ἐδέξατο ἀπὸ τῶν λόγων αὐτοῦ
\vs{4}νόμον, ὃν ἐνετείλατο ἡμῖν Μωυσῆς, κληρονομίαν συναγωγαῖς Ἰακώβ.
\vs{5}Καὶ ἔσται ἐν τῷ ἠγαπημένῳ ἄρχων, συναχθέντων ἀρχόντων λαῶν ἅμα φυλαῖς Ἰσραήλ.
\vs{6}Ζήτω Ῥουβὴν, καὶ μὴ ἀποθανέτω, καὶ ἔστω πολὺς ἐν ἀριθμῷ.

\vs{7}Καὶ αὕτη Ἰούδα· εἰσάκουσον Κύριε φωνῆς Ἰούδα, καὶ εἰς τὸν λαὸν αὐτοῦ ἔλθοις ἄν· αἱ χεῖρες αὐτοῦ διακρινοῦσιν αὐτῷ, καὶ βοηθὸς ἐκ τῶν ἐχθρῶν ἔσῃ.

\vs{8}Καὶ τῷ Λευὶ εἶπε, δότε Λευὶ δήλους αὐτοῦ, καὶ ἀλήθειαν αὐτοῦ τῷ ἀνδρὶ τῷ ὁσίῳ, ὃν ἐπείρασαν αὐτὸν ἐν πείρᾳ· ἐλοιδόρησαν αὐτὸν ἐφʼ ὕδατος ἀντιλογίας·
\vs{9}Ὁ λέγων τῷ πατρὶ καὶ τῇ μητρὶ, οὐχ ἑώρακά σε, καὶ τοὺς ἀδελφοὺς αὐτοῦ οὐκ ἐπέγνω, καὶ τοὺς υἱοὺς αὐτοῦ ἀπέγνω· ἐφύλαξε τὰ λόγιά σου, καὶ τὴν διαθήκην σου διετήρησε.
\vs{10}Δηλώσουσι τὰ δικαιώματά σου τῷ Ἰακὼβ, καὶ τὸν νόμον σου τῷ Ἰσραήλ· ἐπιθήσουσι θυμίαμα ἐν ὀργῇ σου διαπαντὸς ἐπὶ τὸ θυσιαστήριόν σου.
\vs{11}Εὐλόγησον, Κύριε, τὴν ἰσχὺν αὐτοῦ, καὶ τὰ ἔργα τῶν χειρῶν αὐτοῦ δέξαι· κάταξον ὀσφῦν ἐχθρῶν ἐπανεστηκότων αὐτῷ, καὶ οἱ μισοῦντες αὐτὸν μὴ ἀναστήτωσαν.
\vs{12}Καὶ τῷ Βενιαμὶν εἶπεν, ἠγαπημένος ὑπὸ Κυρίου κατασκηνώσει πεποιθὼς, καὶ ὁ Θεὸς σκιάζει ἐπʼ αὐτῷ πάσας τὰς ἡμέρας, καὶ ἀναμέσον τῶν ὤμων αὐτοῦ κατέπαυσε.

\vs{13}Καὶ τῷ Ἰωσὴφ εἶπεν, ἀπʼ εὐλογίας Κύριου ἡ γῆ αὐτοῦ, ἀπὸ ὡρῶν οὐρανοῦ, καὶ δρόσου, καὶ ἀπὸ ἀβύσσων πηγῶν κάτωθεν,
\vs{14}καὶ καθʼ ὥραν γεννημάτων ἡλίου τροπῶν, καὶ ἀπὸ συνόδων μηνῶν,
\vs{15}ἀπὸ κορυφῆς ὀρέων ἀρχῆς, καὶ ἀπὸ κορυφῆς βουνῶν ἀενάων,
\vs{16}καὶ καθʼ ὥραν γῆς πληρώσεως· καὶ τὰ δεκτὰ τῷ ὀφθέντι ἐν τῇ βάτῳ ἔλθοισαν ἐπὶ κεφαλὴν Ἰωσὴφ, καὶ ἐπὶ κορυφῆς δοξασθεὶς ἐπʼ ἀδελφοῖς.
\vs{17}Πρωτότοκος ταύρου τὸ κάλλος αὐτοῦ, κέρατα μονοκέρωτος τὰ κέρατα αὐτοῦ· ἐν αὐτοῖς ἔθνη κερατιεῖ ἅμα, ἕως ἀπʼ ἄκρου γῆς· αὗται μυρίαδες Ἐφραὶμ, καὶ αὗται χιλιάδες Μανασσῆ.
\vs{18}Καὶ τῷ Ζαβουλὼν εἶπεν, εὐφράνθητι Ζαβουλὼν ἐν ἐξοδίᾳ σου, καὶ Ἰσσάχαρ ἐν τοῖς σκηνώμασιν αὐτοῦ.
\vs{19}Ἔθνη ἐξολοθρεύσουσι· καὶ ἐπικαλέσεσθε ἐκεῖ, καὶ θύσετε ἐκεῖ θυσίαν δικαιοσύνης· ὅτι πλοῦτος θαλάσσης θηλάσει σε, καὶ ἐμπόρια παράλιον κατοικούντων.

\vs{20}Καὶ τῷ Γὰδ εἶπεν, εὐλογημένος ἐμπλατύνων Γάδ· ὡς λέων ἐνεπαύσατο, συντρίψας βραχίονα καὶ ἄρχοντα.
\vs{21}Καὶ εἶδεν ἀπαρχὴν αὐτοῦ, ὅτι ἐκεῖ ἐμερίσθη γῆ ἀρχόντων συνηγμένων ἅμα ἀρχηγοῖς λαῶν· δικαιοσύνην Κύριος ἐποίησε, καὶ κρίσιν αὐτοῦ μετὰ Ἰσραήλ.

\vs{22}Καὶ τῷ Δὰν εἶπε, Δὰν σκύμνος λέοντος, καὶ ἐκπηδήσεται ἐκ τοῦ Βασάν.
\vs{23}Καὶ τῷ Νεφθαλὶ εἶπε, Νεφθαλὶ πλησμονὴ δεκτῶν· καὶ ἐμπλησθήτω εὐλογία παρὰ Κυρίου· θάλασσαν καὶ Λίβα κληρονομήσει.
\vs{24}Καὶ τῷ Ἀσὴρ εἶπεν, εὐλογημένος ἀπὸ τέκνων Ἀσήρ, καὶ ἔσται δεκτὸς τοῖς ἀδελφοῖς αὐτοῦ· βάψει ἐν ἐλαίῳ τὸν πόδα αὐτοῦ.
\vs{25}Σίδηρος καὶ χαλκὸς τὸ ὑπόδημα αὐτοῦ ἔσται· ὡς αἱ ἡμέραι σου, ἡ ἰσχύς σου.

\vs{26}Οὐκ ἔστιν ὥσπερ ὁ Θεὸς τοῦ ἠγαπημένου, ὁ ἐπιβαίνων ἐπὶ τὸν οὐρανὸν βοηθός σου, καὶ ὁ μεγαλοπρεπὴς τοῦ στερεώματος.
\vs{27}Καὶ σκεπάσει σε Θεοῦ ἀρχὴ, καὶ ὑπὸ ἰσχὺν βραχιόνων ἀενάων· καὶ ἐκβαλεῖ ἀπὸ προσώπου σου ἐχθρὸν, λέγων, ἀπόλοιο.
\vs{28}Καὶ κατασκηνώσει Ἰσραὴλ πεποιθὼς, μόνος ἐπὶ γῆς Ἰακὼβ, ἐπὶ σίτῳ καὶ οἴνῳ· καὶ ὁ οὐρανός σοι συννεφὴς δρόσῳ.
\vs{29}Μακάριος σὺ Ἰσραήλ· τίς ὅμοιός σοι λαὸς σωζόμενος ὑπὸ Κυρίου; ὑπερασπιεῖ ὁ βοηθός σου, καὶ ἡ μάχαιρα καύχημά σου· καὶ ψεύσονταί σε οἱ ἐχθροί σου· καὶ σὺ ἐπὶ τὸν τράχηλον αὐτῶν ἐπιβήσῃ.

\ch{34}
Καὶ ἀνεβη Μωυσῆς ἀπὸ Ἀραβὼθ Μωὰβ ἐπὶ τὸ ὄρος Ναβαῦ, ἐπὶ κορυφὴν Φασγὰ, ἥ ἐστιν ἐπὶ προσώπου Ἱεριχώ· καὶ ἔδειξεν αὐτῷ Κύριος πᾶσαν τὴν γῆν Γαλαὰδ ἕως Δὰν,
\vs{2}καὶ πᾶσαν τὴν γῆν Νεφθαλὶ, καὶ πᾶσαν τὴν γῆν Ἐφραὶμ, καὶ Μανασσῆ, καὶ πᾶσαν τὴν γῆν Ἰούδα ἕως τῆς θαλάσσης τῆς ἐσχάτης,
\vs{3}καὶ τὴν ἔρημον, καὶ τὰ περίχωρα Ἱεριχὼ, πόλιν φοινίκων ἕως Σηγώρ.
\vs{4}Καὶ εἶπε Κύριος πρὸς Μωυσῆν, αὕτη ἡ γῆ ἣν ὤμοσα τῷ Ἁβραὰμ καὶ Ἰσαὰκ καὶ Ἰακὼβ, λέγων, τῷ σπέρματι ὑμῶν δώσω αὐτήν· καὶ ἔδειξα τοῖς ὀφθαλμοῖς σου, καὶ ἐκεῖ οὐκ εἰσελεύσῃ.

\vs{5}Καὶ ἐτελεύτησε Μωυσῆς ὁ οἰκέτης Κυρίου ἐν γῇ Μωὰβ διὰ ῥήματος Κυρίου.
\vs{6}Καὶ ἔθαψαν αὐτὸν ἐν Γαὶ ἐγγὺς οἴκου Φογώρ· καὶ οὐκ εἶδεν οὐδεὶς τὴν ταφὴν αὐτοῦ ἕως τῆς ἡμέρας ταύτης.
\vs{7}Μωυσῆς δὲ ἦν ἑκατὸν καὶ εἴκοσι ἐτῶν ἐν τῷ τελευτᾷν αὐτόν· οὐκ ἠμαυρώθησαν οἱ ὀφθαλμοὶ αὐτοῦ, οὐδὲ ἐφθάρησαν τὰ χελώνια αὐτοῦ.

\vs{8}Καὶ ἔκλαυσαν οἱ υἱοὶ Ἰσραὴλ Μωυσῆν ἐν Ἀραβὼθ Μωὰβ ἐπὶ τοῦ Ἰορδάνου κατὰ Ἱεριχὼ τριάκοντα ἡμέρας, καὶ συνετελέσθησαν αἱ ἡμέραι πένθους κλαυθμοῦ Μωυσῆ.
\vs{9}Καὶ Ἰησοῦς υἱὸς Ναυῆ ἐνεπλήσθη πνεύματος συνέσεως, ἐπέθηκε γὰρ Μωυσῆς τὰς χεῖρας αὐτοὺ ἐπʼ αὐτόν· καὶ εἰσήκουσαν αὐτοῦ οἱ υἱοὶ Ἰσραήλ· καὶ ἐποίησαν καθότι ἐνετείλατο Κύριος τῷ Μωυσῇ.

\vs{10}Καὶ οὐκ ἀνέστη ἔτι προφήτης ἐν Ἰσραὴλ ὡς Μωσῆς· ὃν ἔγνω Κύριος αὐτὸν πρόσωπον κατὰ πρόσωπον
\vs{11}ἐν πᾶσι τοῖς σημείοις καὶ τέρασιν, ὃν ἀπέστειλεν αὐτὸν Κύριος ποιῆσαι αὐτὰ ἐν γῇ Αἰγύπτῳ Φαραὼ, καὶ τοῖς θεράπουσιν αὐτοῦ, καὶ πάσῃ τῇ γῇ αὐτοῦ·
\vs{12}τὰ θαυμάσια τὰ μεγάλα, καὶ τὴν χεῖρα τὴν κραταιὰν, ἃ ἐποίησε Μωυσῆς ἔναντι παντὸς Ἰσραήλ.


\def\book{ΙΗΣΟΥΣ ΝΑΥΗ}
\biblebook{ΙΗΣΟΥΣ ΝΑΥΗ}


\lettrine[lines=2, loversize=0.2, nindent=0em, findent=.25em]{\textcolor{bookheadingcolor}{Κ}}{ΑΙ} ἐγένετο μετὰ τὴν τελευτὴν Μωυσῆ, εἶπε Κύριος τῷ Ἰησοῖ υἱῷ Ναυῆ τῷ ὑπουργῷ Μωυσῆ, λέγων,
\vs{2}Μωυσῆς ὁ θεράπωνμου τετελεύτηκε. νῦν οὖν ἀναστὰς, διάβηθι τὸν Ἰορδάνην σὺ καὶ πᾶς ὁ λαὸς οὗτος εἰς τὴν γῆν, ἣν ἐγὼ δίδωμι αὐτοῖς.
\vs{3}πᾶς ὁ τόπος ἐφʼ ὃν ἂν ἐπιβῆτε τῷ ἴχνει τῶν ποδῶν ὑμῶν, ὑμῖν δώσω αὐτόν, ὃν τρόπον εἴρηκα τῷ Μωυσῇ·
\vs{4}Τὴν ἔρημον καὶ τὸν Ἀντιλίβανον, ἕως τοῦ ποταμοῦ τοῦ μεγάλου, ποταμοῦ Εὐφράτου, καὶ ἕως τῆς θαλάσσης τῆς ἐσχάτης· ἀφʼ ἡλίου δυσμῶν ἔσται τὰ ὅρια ὑμῶν.
\vs{5}οὐκ ἀντιστήσεται ἄνθρωπος κατενώπιον ὑμῶν πάσας τὰς ἡμέρας τῆς ζωῆς σου· καὶ ὥσπερ ἤμην μετὰ Μωυσῆ, οὕτως ἔσομαι καὶ μετὰ σοῦ, καὶ οὐκ ἐγκαταλείψω σε οὐδʼ ὑπερόψομαί σε.
\vs{6}Ἴσχυε καὶ ἀνδρίζου· σὺ γὰρ ἀποδιεῖς τῷ λαῷ τούτῳ τὴν γῆν, ἣν ὤμοσα τοῖς πατράσινν ὑμῶν δοῦναι αὐτοῖς.
\vs{7}Ἴσχυε οὖν καὶ ἀνδρίζου, φυλάσσεσθαι καὶ ποιεῖν καθότι ἐνετείλατό σοι Μωυσῆς ὁ παῖς μου· καὶ οὐκ ἐκκλινεῖς ἀπʼ αὐτῶν εἰς δεξιὰ οὐδὲ εἰς ἀριστερὰ, ἵνα συνῇς ἐν πᾶσιν οἷς ἐὰν πράσσῃς.
\vs{8}Καὶ οὐκ ἀποστήσεται ἡ βίβλος τοῦ νόμου τούτου ἐκ τοῦ στόματός σου, καὶ μελετήσεις ἐν αὐτῇ ἡμέρας καὶ νυκτὸς, ἵνα εἰδῇς ποιεῖν πάντα τὰ γεγραμμένα· τότε εὐοδωθήσῃ, καὶ εὐοδώσεις τὰς ὁδούς σου, καὶ τότε συνήσεις.
\vs{9}Ἰδοὺ ἐντέαλμαί σοι· ἴσχυε καὶ ἀνδρίζου, μὴ δειλιάσῃς, μηδὲ φοβηθῇ, ὅτι μετὰ σοῦ Κύριος ὁ Θεός σου εἰς πάντα οὗ ἐὰν πορεύῃ·
\vs{10}Καὶ ἐνετείλατο Ἰησοῦς τοῖς γραμματεῦσιν τοῦ λαοῦ, λέγων,
\vs{11}Εἰσέλθατε κατὰ μέσον τῆς παρεμβολῆς τοῦ λαοῦ, καὶ ἐντείλασθε τῷ λαῷ, λέγοντες, ἐτοιμάζεσθε ἐπισιτισμόν, ὅτι ἔτι τρεῖς ἡμέραι καὶ ὑμεῖς διαβαίνετε τὸν Ἰορδάνην τοῦτον, εἰσελθόντες κατασχεῖν τὴν γῆν, ἣν Κύριος ὁ Θεὸς τῶν πατέρων ὑμῶν δίδωσιν ὑμῖν.

\vs{12}Καὶ τῷ Ῥουβὴν, καὶ τῷ Γὰδ, καὶ τῷ ἡμίσει φυλῆς Μανασσῆ εἶπεν Ἰησοῦς,
\vs{13}μνήσθητε τὸ ῥῆμα ὃ ἐνετείλατο ὑμῖν Μωυσῆς ὁ παῖς Κυρίου, λέγων, Κύριος ὁ Θεὸς ὑμῶν κατέπαυσεν ὑμᾶς, καὶ ἔδωκεν ὑμῖν τὴν γῆν ταύτην.
\vs{14}Αἱ γυναῖκες ὑμῶν καὶ τὰ παιδία ὑμῶν καὶ τὰ κτήνη ὑμῶν κατοικείτωσαν ἐν τῇ γῇ, ᾗ ἔδωκεν ὑμῖν· ὑμεῖς δὲ διαβήσεσθε εὔζωνοι πρότεροι τῶν ἀδελφῶν ὑμῶν πᾶς ὁ ἰσχύων· καὶ συμμαχήσετε αὐτοῖς,
\vs{15}ἕως ἂν καταπαύσῃ Κύριος ὁ Θεὸς ἡμῶν τοὺς ἀδελφοὺς ὑμῶν, ὥσπερ καὶ ὑμᾶς, καὶ κληρονομήσωσι καὶ οὗτοι τὴν γῆν, ἣν Κύριος ὁ Θεὸς ἡμῶν δίδωσιν αὐτοῖς· καὶ ἀπελεύσεσθε ἕκαστος εἰς τὴν κληρονομίαν αὐτοῦ, ἣν ἔδωκεν ὑμῖν Μωυσῆς εἰς τὸ πέραν τοῦ Ἰορδάνου ἐπʼ ἀνατολῶν ἡλίου.
\vs{16}Καὶ ἀποκριθέντες τῷ Ἰησοῦ εἶπαν, πάντα ὅσα ἐὰν ἐντείλῃ ἡμῖν, ποιήσομεν, καὶ εἰς πάντα τόπον οὗ ἐὰν ἀποστείλῃς ἡμᾶς, πορευσόμεθα.
\vs{17}Κατὰ πάντα ὅσα ἠκούσαμεν Μωυσῆ, ἀκουσόμεθά σοῦ· πλὴν ἔστω Κύριος ὁ Θεὸς ἡμῶν μετὰ σοῦ, ὃν τρόπον ἦν μετὰ Μωυσῆ.
\vs{18}Ὁ δὲ ἄνθρωπος ὃς ἂν ἀπειθήσῃ σοι, καὶ ὅστις μὴ ἀκούσῃ τῶν ῥημάτων σου καθότι ἐὰν ἐντείλῃ αὐτῷ, ἀποθανέτω· ἀλλὰ ἴσχυε καὶ ἀνδρίζου.

\ch{2}
Καὶ ἀπέστειλεν Ἰησοῦς υἱὸς Ναυῆ ἐκ Σαττὶν δύο νεανίσκους κατασκοπεῦσαι, λέγων, ἀνάβητε καὶ ἴδετε τὴν γῆν καὶ τὴν Ἱεριχώ· καὶ πορευθέντες οἱ δύο νεανίσκοι εἰσήλθοσαν εἰς Ἱεριχώ· καὶ εἰσήλθοσαν εἰς οἰκίαν γυναικὸς πόρνης, ᾗ ὄνομα Ῥαάβ· Καὶ κατέλυσαν ἐκεῖ.

\vs{2}Καὶ ἀπηγγέλη τῷ βασιλεῖ Ἱεριχὼ, λέγοντες, εἰσπεπόρευνται ὧδε ἄνδρες τῶν υἱῶν Ἰσραὴλ κατασκοπεῦσαι τὴν γῆν.
\vs{3}Καὶ ἀπέστειλεν ὁ βασιλεὺς Ἱεριχὼ, καὶ εἶπε πρὸς Ῥαὰβ, λέγων, ἐξάγαγε τοὺς ἄνδρας τοὺς εἰσπεπορευμένους εἰς τὴν οἰκίαν σου τὴν νύκτα, κατασκοπεῦσαι γὰρ τὴν γῆν ἥκασι.
\vs{4}Καὶ λαβοῦσα ἡ γυνὴ τοὺς δύο ἄνδρας, ἔκρυψεν αὐτούς· καὶ εἶπεν αὐτοῖς, λέγουσα, εἰσεληλύθασι πρὸς μὲ οἱ ἄνδρες,
\vs{5}ὡς δὲ ἡ πύλη ἐκλείετο ἐν τῷ σκότει, καὶ οἱ ἄνδρες ἐξῆλθον· οὐκ ἐπίσταμαι ποῦ πεπόρευνται· καταδιώξατε ὀπίσω αὐτῶν, εἰ καταλήψεσθε αὐτούς.
\vs{6}Αὕτη δὲ ἀνεβίβασεν αὐτοὺς ἐπὶ τὸ δῶμα, καὶ ἔκρυψεν αὐτοὺς ἐν τῇ λινοκαλάμῃ τῇ ἐστοιβασμένῃ αὐτῇ ἐπὶ τοῦ δώματος.
\vs{7}Καὶ οἱ ἄνδρες κατεδίωξαν ὀπίσω αὐτῶν ὁδὸν τὴν ἐπὶ τοῦ Ἰορδάνου ἐπὶ τὰς διαβάσεις, καὶ ἡ πύλη ἐκλείσθη.

\vs{8}Καὶ ἐγένετο ὡς ἐξήλθοσαν οἱ διώκοντες ὀπίσω αὐτῶν, καὶ αὐτοὶ δὲ πρὶν ἢ κοιμηθῆναι αὐτοὺς, αὕτη δὲ ἀνέβη πρὸς αὐτοὺς ἐπὶ τὸ δῶμα,
\vs{9}καὶ εἶπε πρὸς αὐτοὺς, ἐπίσταμαι ὅτι ἔδωκεν ὑμῖν Κύριος τὴν γῆν· ἐπιπέπτωκε γὰρ ὁ φόβος ὑμῶν ἐφʼ ἡμᾶς.
\vs{10}Ἀκηκόαμεν γὰρ ὅτι κατεξήρανε Κύριος ὁ Θεὸς τὴν ἐρυθρὰν θάλασσαν ἀπὸ προσώπου ὑμῶν, ὅτε ἐξεπορεύεσθε ἐκ γῆς Αἰγυπτου, καὶ ὅσα ἐποίησε τοῖς δυσὶ βασιλεῦσι τῶν Ἀμοῤῥαίων, οἳ ἦσαν πέραν τοῦ Ἰορδάνου, τῷ Σηὼν καὶ Ὢγ, οὓς ἐξωλοθρεύσατε αὐτούς.
\vs{11}Καὶ ἀκούσαντες ἡμεῖς ἐξέστημεν τῇ καρδίᾳ ἡμῶν, καὶ οὐκ ἔστη ἔτι πνεῦμα ἐν οὐδενὶ ἡμῶν ἀπὸ προσώπου ὑμῶν· ὅτι Κύριος ὁ Θεὸς ὑμῶν, Θεὸς ἐν οὐρανῷ ἄνω καὶ ἐπὶ τῆς γῆς κάτω.
\vs{12}Καὶ νῦν ὀμόσατέ μοι Κύριον τὸν Θεὸν, ὅτι ποιῶ ὑμῖν ἔλεος, καὶ ποιήσατε Καὶ ὑμεῖς ἔλεος ἐν τῷ οἴκῳ τοῦ πατρός μου,
\vs{13}καὶ ζωγρήσατε τὸν οἶκον τοῦ πατρὸς μου, τὴν μητέρα μου, καὶ τοὺς ἀδελφούς μου, καὶ πάντα τὸν οἶκόν μου, καὶ πάντα ὅσα ἐστὶν αὐτοῖς, καὶ ἐξελεῖσθε τὴν ψυχήν μου ἐκ θανάτου.

\vs{14}Καὶ εἶπαν αὐτῇ οἱ ἄνδρες, ἡ ψυχὴ ἡμῶν ἀνθʼ ὑμῶν εἰς θάνατον· καὶ αὐτὴ εἶπεν, ὡς ἂν παραδῷ Κύριος ὑμῖν τὴν πόλιν, ποιήσετε εἰς ἐμὲ ἔλεος καὶ ἀλήθειαν.
\vs{15}Καὶ κατεχάλασεν αὐτοὺς διὰ τῆς θυρίδος,
\vs{16}καὶ εἶπεν αὐτοῖς, εἰς τὴν ὀρεινὴν ἀπέλθετε, μὴ συναντήσωσιν ὑμῖν οἱ καταδιώκοντες, καὶ κρυβήσεσθε ἐκεῖ τρεῖς ἡμέρας ἕως ἂν ἀποστρέψωσιν οἱ καταδιώκοντες ὀπίσω ὑμῶν, καὶ μετὰ ταῦτα ἀπελεύσεσθε εἰς τὴν ὁδὸν ὑμῶν.

\vs{17}Καὶ εἶπαν πρὸς αὐτὴν οἱ ἄνδρες, ἀθῶοι ἐσμὲν τῷ ὅρκῳ σου τούτῳ.
\vs{18}Ἰδοὺ ἡμεῖς εἰσπορευόμεθα εἰς μέρος τῆς πόλεως, καὶ θήσεις τὸ σημεῖον, τὸ σπαρτίον τὸ κόκκινον τοῦτο ἐκδήσεις εἰς τὴν θυρίδα διʼ ἧς κατεβίβασας ἡμᾶς διʼ αὐτῆς· τὸν δὲ πατέρα σου, καὶ τὴν μητέρα σου, καὶ τοὺς ἀδελφούς σου, καὶ πάντα τὸν οἶκον τοῦ πατρός σου συνάξεις πρὸς σεαυτὴν εἰς τὴν οἰκίαν σου.
\vs{19}Καὶ ἔσται πᾶς ὃς ἂν ἐξέλθῃ τὴν θύραν τῆς οἰκίας σου ἔξω, ἔνοχος ἑαυτῷ ἔσται, ἡμεῖς δὲ ἀθῶοι τῷ ὅρκῳ σου τούτῳ· καὶ ὅσοι ἐὰν γένωνται μετὰ σοῦ ἐν τῇ οἰκίᾳ σου, ἡμεῖς ἔνοχοι ἐσόμεθα.
\vs{20}Ἐὰν δέ τις ἡμᾶς ἀδικήσῃ ἢ καὶ ἀποκαλύψῃ τοὺς λόγους ἡμῶν τούτους, ἐσόμεθα ἀθῶοι τῷ ὅρκῳ σου τούτῳ.
\vs{21}Καὶ εἶπεν αὐτοῖς, κατὰ τὸ ῥῆμα ὑμῶν ἔστω· καὶ ἐξαπέστειλεν αὐτοὺς, καὶ ἐπορεύθησαν.
\vs{22}Καὶ ἤλθοσαν εἰς τὴν ὀρεινήν, καὶ κατέμειναν ἐκεῖ τρεῖς ἡμέρας· καὶ ἐξεζήτησαν οἱ καταδιώκοντες πάσας τὰς ὁδοὺς, καὶ οὐχ εὕροσαν.

\vs{23}Καὶ ὑπέστρεψαν οἱ δύο νεανίσκοι, καὶ κατέβησαν ἐκ τοῦ ὄρους· καὶ διέβησαν πρὸς Ἰησοῦν υἱὸν Ναυὴ, καὶ διηγήσαντο αὐτῷ πάντα τὰ συμβεβηκότα αὐτοῖς.
\vs{24}Καὶ εἶπαν πρὸς Ἰησοῦν, ὅτι παραδέδωκε Κύριος πᾶσαν τὴν γῆν ἐν χειρὶ ἡμῶν, καὶ κατέπτηχε πᾶς ὁ κατοικῶν τὴν γῆν ἐκείνην ἀφʼ ἡμῶν.

\ch{3}
Καὶ ὤρθρισεν Ἰησοῦς τοπρωῒ, καὶ ἀπῇρεν ἐκ Σαττὶν, καὶ ἤλθοσαν ἕως τοῦ Ἰορδάνου, καὶ κατέλυσαν ἐκεῖ πρὸ τοῦ διαβῆναι.
\vs{2}Καὶ ἐγένετο μετὰ τρεῖς ἡμέρας διῆλθον οἱ γραμματεῖς διὰ τῆς παρεμβολῆς,
\vs{3}καὶ ἐνετείλαντο τῷ λαῷ, λέγοντες, ὅταν ἴδητε τὴν κιβωτὸν τῆς διαθήκης Κυρίου τοῦ Θεοῦ ἡμῶν, καὶ τοὺς ἱερεῖς ἡμῶν καὶ τοὺς Λευίτας αἴροντας αὐτὴν, ἀπαρεῖτε ἀπὸ τῶν τόπων ὑμῶν, καὶ πορεύσεσθε ὀπίσω αὐτῆς.
\vs{4}Ἀλλὰ μακρὰν ἔστω ἀναμέσον ὑμῶν καὶ ἐκείνης, ὅσον δισχιλίους πήχεις στήσεσθε· μὴ προσεγγίσητε αὐτῇ, ἵνα ἐπίστησθε τὴν ὁδὸν, ἣν πορεύσεσθε αὐτήν· οὐ γὰρ πεπόρευσθε τὴν ὁδὸν ἀπʼ ἐχθὲς καὶ τρίτης ἡμέρας.

\vs{5}Καὶ εἶπεν Ἰησοῦς τῷ λαῷ, ἁγνίσασθε εἰς αὔριον, ὅτι αὔριον ποιήσει Κύριος ἐν ὑμῖν θαυμαστά.
\vs{6}Καὶ εἶπεν Ἰησοῦς τοῖς ἱερεῦσιν, ἄρατε τὴν κιβωτὸν τῆς διαθήκης Κυρίου, καὶ προπορεύεσθε τοῦ λαοῦ· καὶ ᾖραν οἱ ἱερεῖς τὴν κιβωτὸν τῆς διαθήκης Κυρίου, καὶ ἐπορεύοντο ἔμπροσθεν τοῦ λαοῦ.
\vs{7}Καὶ εἶπε Κύριος πρὸς Ἰησοῦν, ἐν τῇ ἡμέρᾳ ταύτῃ ἄρχομαι ὑψῶσαί σε κατενώπιον πάντων υἱῶν Ἰσραὴλ, ἵνα γνῶσιν ὅτι καθότι ἤμην μετὰ Μωυσῆ, οὕτως ἔσομαι καὶ μετὰ σοῦ.
\vs{8}Καὶ νῦν ἔντειλαι τοῖς ἱερεῦσι τοῖς αἴρουσι τὴν κιβωτὸν τῆς διαθήκης, λέγων, ὠς ἂν εἰσέλθητε ἐπὶ μέρους τοῦ ὕδατος τοῦ Ἰορδάνου, καὶ ἐν τῷ Ἰορδάνῃ στήσεσθε.

\vs{9}Καὶ εἶπεν Ἰησοῦς τοῖς υἱοῖς Ἰσραὴλ, προσαγάγετε ὧδε, καὶ ἀκούσατε τὸ ῥῆμα Κυρίου τοῦ Θεοῦ ἡμῶν.
\vs{10}Ἐν τούτῳ γνώσεσθε, ὅτι Θεὸς ζῶν ἐν ὑμῖν, καὶ ὀλοθρεύων ὀλοθρεύσει ἀπὸ προσώπου ἡμῶν τὸν Χαναναῖον, καὶ τὸν Χετταῖον, καὶ τὸν Φερεζαῖον, καὶ τὸν Εὐαῖον, καὶ τὸν Ἀμοῤῥαῖον, καὶ τὸν Γεργεσαῖον, καὶ τὸν Ἰεβουσαῖον.
\vs{11}Ἰδοὺ ἡ κιβωτὸς διαθήκης Κυρίου πάσης τῆς γῆς διαβαίνει τὸν Ἰορδάνην.
\vs{12}Προχειρίσασθε ὑμῖν δώδεκα ἄνδρας ἀπὸ τῶν υἱῶν Ἰσραὴλ, ἕνα ἀφʼ ἑκάστης φυλῆς.
\vs{13}Καὶ ἔσται, ὡς ἂν καταπαύσωσιν οἱ πόδες τῶν ἱερέων τῶν αἰρόντων τὴν κιβωτὸν τῆς διαθήκης Κυρίου πάσης τῆς γῆς ἐν τῷ ὕδατι τοῦ Ἰορδάνου, τὸ ὕδωρ τοῦ Ἰορδάνου ἐκλείψει, τὸ δὲ ὕδωρ τὸ καταβαῖνον στήσεται.

\vs{14}Καὶ ἀπῇρεν ὁ λαὸς ἐκ τῶν σκηνωμάτων αὐτῶν διαβῆναι τὸν Ἰορδάνην, οἱ δὲ ἱερεῖς ᾔροσαν τὴν κιβωτὸν τῆς διαθήκης Κυρίου πρότεροι τοῦ λαοῦ.
\vs{15}Ὡς δὲ εἰσεπορεύοντο οἱ ἱερεῖς οἱ αἴροντες τὴν κιβωτὸν τῆς διαθήκης ἐπὶ τὸν Ἰορδάνην, καὶ οἱ πόδες τῶν ἱερέων τῶν αἰρόντων τὴν κιβωτὸν τῆς διαθήκης Κυρίου ἐβάφησαν εἰς μέρος τοῦ ὕδατος τοῦ Ἰορδάνου· ὁ δὲ Ἰορδάνης ἐπληροῦτο καθʼ ὅλην τὴν κρηπίδα αὐτοῦ, ὡσεὶ ἡμέραι θερισμοῦ πυρῶν·
\vs{16}Καὶ ἔστη τὰ ὕδατα τὰ καταβαίνοντα ἄνωθεν, ἔστη πῆγμα ἓν ἀφεστηκὸς μακρὰν σφόδρα σφοδρῶς ἕως μέρους Καριαθιαρίμ· τὸ δὲ καταβαῖνον κατέβη εἰς τὴν θάλασσαν Ἄραβα θάλασσαν ἁλὸς, ἕως εἰς τὸ τέλος ἐξέλιπε· καὶ ὁ λαὸς εἱστήκει ἀπέναντι Ἱεριχώ.
\vs{17}Καὶ ἔστησαν οἱ ἱερεῖς οἱ αἴροντες τὴν κιβωτὸν τῆς διαθήκης Κυρίου ἐπὶ ξηρᾶς ἐν μέσῳ τοῦ Ἰορδάνου· καὶ πάντες οἱ υἱοὶ Ἰσραὴλ διέβαινον διὰ ξηρᾶς, ἕως συνετέλεσε πᾶς ὁ λαὸς διαβαίνων τὸν Ἰορδάνην.

\ch{4}
Καὶ ἐπεὶ συνετέλεσε πᾶς ὁ λαὸς διαβαίνων τὸν Ἰορδανην, καὶ εἶπε, Κύριος τῷ Ἰησοῖ, λέγων,
\vs{2}ταραλαβὼν ἄνδρας ἀπὸ τοῦ λαοῦ, ἕνα ἀφʼ ἑκάστης φυλῆς,
\vs{3}σύνταξον αὐτοῖς· καὶ ἀνέλεσθε ἐκ μέσου τοῦ Ἰορδάνου ἑτοίμους δώδεκα λίθους, καὶ τούτους διακομίσαντες ἅμα ὑμῖν αὐτοῖς, θέτε αὐτοὺς ἐν τῇ στρατοπεδίᾳ ὑμῶν, οὗ ἐὰν παρεμβάλητε ἐκεῖ τὴν νύκτα.

\vs{4}Καὶ ἀνακαλεσάμενος Ἰησοῦς δώδεκα ἄνδρας τῶν ἐνδόξων ἀπὸ τῶν υἱῶν Ἰσραὴλ, ἕνα ἀφʼ ἑκάστης φυλῆς,
\vs{5}εἶπεν αὐτοῖς, προσαγάγετε ἔμπροσθέν μου πρὸ προσώπου Κυρίου εἰς μέσον τοῦ Ἰροδάνου· καὶ ἀνελόμενος ἐκεῖθεν ἕκαστος λίθον, ἀράτω ἐπὶ τῶν ὤμων αὐτοῦ κατὰ τὸν ἀριθμὸν τῶν δώδεκα φυλῶν τοῦ Ἰσραὴλ,
\vs{6}ἵνα ὑπάρχωσιν ὑμῖν οὗτοι εἰς σημεῖον κείμενον διαπαντός· ἵνα ὅταν ἐρωτᾷ σε ὁ υἱός σου αὔριον λέγων, τί εἰσιν οἱ λιθοι οὗτοι ἡμῖν;
\vs{7}Καὶ σὺ δηλώσεις τῷ υἱῷ σου, λέγων, ὅτι ἐξέλιπεν ὁ Ἰορδάνης ποταμὸς ἀπὸ προσώπου κιβωτοῦ διαθήκης Κυρίου πάσης τῆς γῆς, ὡς διέβαινεν αὐτόν· καὶ ἔσονται οἱ λίθοι οὗτοι ὑμῖν μνημόσυνον τοῖς υἱοῖς Ἰσραὴλ ἕως τοῦ αἰῶνος.

\vs{8}Καὶ ἐποίησαν οὕτως οἱ υἱοὶ Ἰσραὴλ, καθότι ἐνετείλατο Κύριος τῷ Ἰησοῖ· καὶ ἀναλαβόντες δώδεκα λίθους ἐκ μέσου τοῦ Ἰορδάνου, καθάπερ συνέταξε Κύριος τῷ Ἰησοῖ ἐν τῇ συντελείᾳ τῆς διαβάσεως τῶν υἱῶν Ἰσραὴλ, καὶ διεκόμισαν ἅμα ἑαυτοῖς εἰς τὴν παρεμβολὴν, καὶ ἀπέθηκαν ἐκεῖ.
\vs{9}Ἔστησε δὲ Ἰησοῦς καὶ ἄλλους δώδεκα λίθους ἐν αὐτῷ τῷ Ἰορδάνῃ, ἐν τῷ γενομένῳ τόπῳ ὑπὸ τοὺς πόδας τῶν ἱερέων τῶν αἰρόντων τὴν κιβωτὸν τῆς διαθήκης Κυρίου· καὶ εἰσὶν ἐκεῖ ἕως τῆς σήμερον ἡμέρας.

\vs{10}Εἱστήκεισαν δὲ οἱ ἱερεῖς οἱ αἴροντες τὴν κιβωτὸν τῆς διαθήκης ἐν τῷ Ἰορδάνῃ, ἕως οὗ συνετέλεσεν Ἰησοῦς πάντα ἃ ἐνετείλατο Κύριος ἀναγγεῖλαι τῷ λαῷ· καὶ ἔσπευσεν ὁ λαὸς, καὶ διέβησαν.
\vs{11}Καὶ ἐγένετο ὡς συνετέλεσε πᾶς ὁ λαὸς διαβῆναι, καὶ διέβη ἡ κιβωτὸς τῆς διαθήκης Κυρίου, καὶ οἱ λίθοι ἔμπροσθεν αὐτῶν.
\vs{12}Καὶ διέβησαν οἱ υἱοὶ Ῥουβὴν, καὶ οἱ υἱοὶ Γὰδ, καὶ οἱ ἡμίσεις φυλῆς Μανασσῆ διεσκευασμένοι ἔμπροσθεν τῶν υἱῶν Ἰσραὴλ, καθάπερ ἐνετείλατο αὐτοῖς Μωυσῆς.
\vs{13}Τετρακισμύριοι εὔζωνοι εἰς μάχην διέβησαν ἐναντίον Κυρίου εἰς πόλεμον πρὸς τὴν Ἱεριχὼ πόλιν.
\vs{14}Ἐν ἐκείνῃ τῇ ἡμέρᾳ ηὔξησε Κύριος τὸν Ἰησοῦν ἐναντίον τοῦ παντὸς γένους Ἰσραὴλ· καὶ ἐφοβοῦντο αὐτὸν, ὥσπερ Μωυσῆν, ὅσον χρόνον ἔζη.

\vs{15}Καὶ εἶπε Κύριος τῷ Ἰησοῖ, λέγων,
\vs{16}ἔντειλαι τοῖς ἱερεῦσι τοῖς αἴρουσι τὴν κιβωτὸν τῆς διαθήκης τοῦ μαρτυρίου Κυρίου, ἐκβῆναι ἐκ τοῦ Ἰορδάνου.
\vs{17}Καὶ ἐνετείλατο Ἰησοῦς τοῖς ἱερεῦσι, λέγων, ἔκβητε ἐκ τοῦ Ἰορδάνου.
\vs{18}Καὶ ἐγένετο ὡς ἐξέβησαν οἱ ἱερεῖς οἱ αἴροντες τὴν κιβωτὸν τῆς διαθήκης Κυρίου ἐκ τοῦ Ἰορδάνου, καὶ ἔθηκαν τοὺς πόδας ἐπὶ τῆς γῆς, ὥρμησε τὸ ὕδωρ τοῦ Ἰορδάνου κατὰ χώραν, καὶ ἐπορεύετο καθὰ χθὲς καὶ τρίτην ἡμέραν διʼ ὅλης τῆς κρηπίδος.

\vs{19}Καὶ ὁ λαὸς ἀνέβη ἐκ τοῦ Ἰορδάνου δεκάτῃ τοῦ μηνὸς τοῦ πρώτου· καὶ κατεστρατοπέδευσαν οἱ υἱοὶ Ἰσραὴλ ἑν Γαλγάλοις κατὰ μέρος τὸ πρὸς ἡλίου ἀνατολὰς ἀπὸ τῆς Ἱεριχώ.
\vs{20}Καὶ τοὺς δώδεκα λίθους τούτους, οὓς ἔλαβεν ἐκ τοῦ Ἰορδάνου, ἔστησεν Ἰησοῦς ἐν Γαλγάλοις,
\vs{21}λέγων, ὅταν ἐρωτῶσιν ὑμᾶς οἱ υἱοι ὑμῶν λέγοντες, τί εἰσιν οἱ λίθοι οὗτοι;
\vs{22}Ἀναγγείλατε τοῖς υἱοῖς ὑμῶν, ὅτι ἐπὶ ξηρᾶς διέβη Ἰσραὴλ τὸν Ἰορδάνην τοῦτον,
\vs{23}ἀποξηράναντος Κυρίου τοῦ Θεοῦ ἡμῶν τὸ ὕδωρ τοῦ Ἰορδάνου ἐκ τῶν ἔμπροσθεν αὐτῶν, μέχρις οὗ διέβησαν· καθάπερ ἐποίησε Κύριος ὁ Θεὸς ἡμῶν τὴν ἐρυθρὰν θάλασσαν, ἣν ἀπεξήρανε Κύριος ὁ Θεὸς ἡμῶν ἔμπροσθεν ἡμῶν, ἕως παρήλθομεν·
\vs{24}Ὅπως γνῶσι πάντα τὰ ἔθνη τῆς γῆς, ὅτι ἡ δύναμις τοῦ κυρίου ἰσχυρά ἐστι, καὶ ἵνα ὑμεῖς σέβησθε Κύριον τὸν Θεὸν ἡμῶν ἐν παντὶ ἔργῳ.

\ch{5}
Καὶ ἐγένετο ὡς ἤκουσαν οἱ βασιλεῖς τῶν Ἀμοῤῥαίων οἳ ἦσαν πέραν τοῦ Ἰορδάνου, καὶ οἱ βασιλεῖς τῆς Φοινίκης οἳ παρὰ τὴν θάλασσαν, ὅτι ἀπεξήρανε Κύριος ὁ Θεὸς τὸν Ἰορδάνην ποταμὸν ἐκ τῶν ἔμπροσθεν τῶν νἱῶν Ἰσραὴλ ἐν τῷ διαβαίνειν αὐτοὺς καὶ ἐτάκησαν αὐτῶν αἱ διάνοιαι, καὶ κατεπλάγησαν, καὶ οὐκ ἦν ἐν αὐτοῖς φρόνησις οὐδεμία ἀπὸ προσώπου τῶν υἱῶν Ἰσραήλ.

\vs{2}Ὑπὸ δὲ τοῦτον τὸν καιρὸν εἶπε Κύριος τῷ Ἰησοῖ, ποίησον σεαυτῷ μαχαίρας πετρίνας ἐκ πέτρας ἀκροτόμου, καὶ καθίσας περίτεμε τοὺς υἱοὺς Ἰσραὴλ ἐκ δευτέρον.
\vs{3}Καὶ ἐποίησεν Ἰησοῦς μαχαίρας πετρίνας ἀκροτόμους, καὶ περιέτεμε τοὺς υἱοὺς Ἰσραὴλ ἐπὶ τοῦ καλουμένου τόπου, Βουνὸς τῶν ἀκροβυστιῶν.
\vs{4}Ὃν δὲ τρόπον περιεκάθαρεν Ἰησοῦς τοὺς υἱοὺς Ἰσραήλ· ὅσοι ποτὲ ἐγένοντο ἐν τῇ ὁδῷ, καὶ ὅσοι ποτὲ ἀπερίτμητοι ἦσαν τῶν ἐξεληλυθότων ἐξ Αἰγύπτου,
\vs{5}πάντας τούτους περιέτεμεν Ἰησοῦς· τεσσαράκοντα γὰρ καὶ δύο ἔτη ἀνέστραπται Ἰσραὴλ ἐν τῇ ἐρήμῳ τῇ Μαβδαρίτιδτ.
\vs{6}Διὸ ἀπερίτμητοι ἦσαν οἱ πλεῖστοι αὐτῶν τῶν μαχίμων τῶν ἐξεληλυθότων ἐκ γῆς Αἰγύπτου, οἱ ἀπειθήσαντες τῶν ἐντολῶν τοῦ Θεοῦ, οἷς καὶ διώρισε μὴ ἰδεῖν αὐτοὺς τὴν γῆν, ἣν ὤμοσε Κύριος τοῖς πατράσιν αὐτῶν δοῦναι γῆν ῥέουσαν γάλα καὶ μέλι.
\vs{7}Ἀντὶ δὲ τούτων ἀντικατέστησε τοὺς υἱοὺς αὐτῶν, οὓς Ἰησοῦς περιέτεμε, διὰ τὸ αὐτοὺς γεγεννῆσθαι κατὰ τὴν ὁδὸν ἀπεριτμήτους.
\vs{8}Περιτμηθέντες δὲ ἡσυχίαν εἶχον αὐτόθι καθήμενοι ἐν τῇ παρεμβολῇ ἕως ὑγιάσθησαν.
\vs{9}Καὶ εἶπε Κύριος τῷ Ἰησοῖ υἱῷ Ναυή, ἐν τῇ σήμερον ἡμέρᾳ ἀφεῖλον τὸν ὀνειδισμὸν Αἰγύπτου ἀφʼ ὑμῶν· καὶ ἐκάλεσε τὸ ὄνομα τοῦ τόπου ἐκείνου, Γάλγαλα.

\vs{10}Καὶ ἐποίησαν οἱ υἱοὶ Ἰσραὴλ τὸ πάσχα τῇ τεσσαρεσκαιδεκάτῃ ἡμέρᾳ τοῦ μηνὸς ἀφʼ ἑσπέρας ἐπὶ δυσμῶν Ἱεριχὼ ἐν τῷ πέραν τοῦ Ἰορδάνου ἐν τῷ πεδίῳ.
\vs{11}Καὶ ἐφάγοσαν ἀπὸ τοῦ σίτου τῆς γῆς ἄζυμα καὶ νέα.
\vs{12}Ἐν ταύτῃ τῇ ἡμέρᾳ ἐξέλιπε τὸ μάννα μετὰ τὸ βεβρωκέναι αὐτοὺς ἐκ τοῦ σίτου τῆς γῆς, καὶ οὐκέτι ὑπῆρχε τοῖς υἱοῖς Ἰσραὴλ μάννα· ἐκαρπίσαντο δὲ τὴν χώραν τῶν Φοινίκων ἐν τῷ ἐνιαυτῷ ἐκείνῳ.

\vs{13}Καὶ ἐγένετο ὡς ἦν Ἰησοῦς ἐν Ἱεριχὼ, καὶ ἀναβλέψας τοῖς ὀφθαλμοῖς εἶδεν ἄνθρωπον ἑστηκότα ἐναντίον αὐτοῦ, καὶ ἡ ῥομφαία ἐσπασμένη ἐν τῇ χειρὶ αὐτοῦ· καὶ προσελθὼν Ἰησοῦς, εἶπεν αὐτῷ, ἡμέτερος εἶ, ἢ τῶν ὑπεναντίων;
\vs{14}Ὁ δὲ εἶπεν αὐτῷ, ἐγὼ ἀρχιστράτηγος δυνάμεως Κυρίου, νυνὶ παραγέγονα. Καὶ Ἰησοῦς ἔπεσεν ἐπὶ πρόσωπον ἐπὶ τὴν γῆν, καὶ εἶπεν αὐτῷ, δέσποτα, τί προστάσσεις τῷ σῷ οἰκέτῃ;
\vs{15}Καὶ λέγει ὁ ἀρχιστράτηγος Κυρίου πρὸς Ἰησοῦν, Λῦσαι τὸ ὑπόδημα ἐκ τῶν ποδῶν σου, ὁ γὰρ τόπος ἐφʼ ᾧ νῦν ἔστηκας ἐπʼ αὐτοῦ, ἅγιός ἐστι.

\ch{6}
Καὶ Ἱεριχὼ συρκεκλεισμένη καὶ ὠχυρωμένη, καὶ οὐδεὶς ἐξεπορεύετο ἐξ αὐτῆς, οὐδὲ εἰσεπορεύετο.
\vs{2}Καὶ εἶπε Κύριος πρὸς Ἰησοῦν, ἰδον ἐγὼ παραδίδωμι ὑποχείριόν σοι τὴν Ἱεριχὼ, καὶ τὸν βασιλέα αὐτῆς τὸν ἐν αὐτῇ, δυνατοὺς ὄντας ἐν ἰσχύϊ.
\vs{3}Σὺ δὲ περίστησον αὐτῇ τοὺς μαχίμους κύκλῳ.
\vs{5}Καὶ ἔσται ὡς ἂν σαλπίσητε τῇ σάλπιγγι, ἀνακραγέτω πᾶς ὁ λαὸς ἅμα, καὶ ἀνακραγόντων αὐτῶν πεσεῖται αὐτόματα τὰ τείχη τῆς πόλεως, καὶ εἰσελεύσεται πᾶς ὁ λαὸς ὁρμήσας ἕκαστος κατὰ πρόσωπον εἰς τὴν πόλιν.

\vs{6}Καὶ εἰσῆλθεν Ἰησοῦς ὁ τοῦ Ναυῆ πρὸς τοὺς ἱερεῖς,
\vs{7}καὶ εἶπεν αὐτοῖς, λέγων, παραγγείλατε τῷ λαῷ περιελθεῖν, καὶ κυκλῶσαι τὴν πόλιν· καὶ οἱ μάχιμοι παραπορευέσθωσαν ἐνωπλισμένοι ἐναντίον Κυρίου.
\vs{8}Καὶ ἑπτὰ ἱερεῖς ἔχοντες ἑπτὰ σάλπιγγας ἱερὰς παρελθέτωσαν ὡσαύτως ἐναντίον τοῦ Κυρίου, καὶ σημαινέτωσαν εὐτόνως· καὶ ἡ κιβωτὸς τῆς διαθήκης Κυρίου ἐπακολουθείτω.
\vs{9}Οἱ δὲ μάχιμοι παραπορευέσθωσαν ἔμπροσθεν, καὶ οἱ ἱερεῖς οἱ οὐραγοῦντες ὀπίσω τῆς κιβωτοῦ τῆς διαθήκης Κυρίου σαλπίζοντες.
\vs{10}Τῷ δὲ λαῷ ἐνετείλατο Ἰησοῦς, λέγων, μὴ βοᾶτε, μηδὲ ἀκουσάτω μηδεὶς τὴν φωνὴν ὑμῶν, ἕως ἂν ἡμέραν διαγγείλῃ αὐτὸς ἀναβοῆσαι, καὶ τότε ἀναβοήσετε·
\vs{11}Καὶ περιελθοῦσα ἡ κιβωτὸς τῆς διαθήκης τοῦ Θεοῦ εὐθέως ἀπῆλθεν εἰς τὴν παρεμβολὴν, καὶ ἐκοιμήθη ἐκεῖ.

\vs{12}Καὶ τῇ ἡμέρᾳ τῇ δευτέρᾳ ἀνέστη Ἰησοῦς τοπρωῒ, καὶ ᾖραν οἱ ἱερεῖς τὴν κιβωτὸν τῆς διαθήκης Κυρίου.
\vs{13}Καὶ οἱ ἑπτὰ ἱερεῖς οἱ φέροντες τὰς σάλπιγγας τὰς ἑπτὰ προεπορεύοντο ἐναντίον Κυρίου· καὶ μετὰ ταῦτα εἰσεπορεύοντο οἱ μάχιμοι, καὶ ὁ λοιπὸς ὄχλος ὄπισθεν τῆς κιβωτοῦ τῆς διαθήκης Κυρίου·
\vs{14}καὶ οἱ ἱερεῖς ἐσάλπισαν ταῖς σάλπιγξι, καὶ ὁ λοιπὸς ὄχλος ἅπας περιεκύκλωσε τὴν πόλιν ἑξάκις ἐγγύθεν, καὶ ἀπῆλθε πάλιν εἰς τὴν παρεμβολήν· οὕτως ἐποίει ἐπὶ ἓξ ἡμέρας.

\vs{15}Καὶ τῇ ἡμέρᾳ τῇ ἑβδόμῃ ἀνέστησαν ὄρθρου, καὶ περιήλθοσαν τὴν πόλιν ἐν τῇ ἡμέρᾳ ἐκείνῃ ἑπτάκις.
\vs{16}Καὶ ἐγένετο τῇ περιόδῳ τῇ ἑβδόμῃ ἐσάλπισαν οἱ ἱερεῖς· καὶ εἶπεν Ἰησοῦς τοῖς υἱοῖς Ἰσραὴλ, κεκράξατε, παρέδωκε γὰρ Κύριος ὑμῖν τὴν πόλιν.
\vs{17}Καὶ ἔσται ἡ πόλις ἀνάθεμα, αὐτὴ καὶ πάντα ὅσα ἐστὶν ἐν αὐτῇ, Κυρὶῳ σαβαώθ· πλὴν Ῥαὰβ τὴν πόρνην περιποιήσασθε αὐτὴν, καὶ πάντα ὅσα ἐστὶν ἐν τῷ οἴκῳ αὐτῆς.
\vs{18}Ἀγγὰ ὑμεῖς φυλάξεσθε σφόδρα ἀπὸ τοῦ ἀναθέματος, μήποτε ἐνθυμηθέντες ὑμεῖς αὐτοὶ λάβητε ἀπὸ τοῦ ἀναθέματος, καὶ ποιήσητε τὴν παρεμβολὴν τῶν υἱῶν Ἰσραὴλ ἀνάθεμα, καὶ ἐκτρίψητε ἡμᾶς.
\vs{19}Καὶ πᾶν ἀργύριον ἢ χρυσίον, ἢ χαλκὸς ἢ σίδηρος, ἅγιον ἔσται τῷ Κυρὶῳ· εἰς θησαυρὸν Κυρίου εἰσενεχθήσεται.

\vs{20}Καὶ ἐσάλπισαν ταῖς σάλπιγξιν οἱ ἱερεῖς· ὡς δὲ ἤκουσεν ὁ λαὸς τῶν σαλπίγγων, ἠλάλαξε πᾶς ὁ λαὸς ἅμα ἀλαλαγμῷ μεγάλῳ καὶ ἰσχυρῷ· καὶ ἔπεσεν ἅπαν τὸ τεῖχος κύκλῳ· καὶ ἀνέβη πᾶς ὁ λαὸς εἰς τὴν πόλιν.
\vs{21}Καὶ ἀνεθεμάτισεν αὐτὴν Ἰησοῦς, καὶ ὅσα ἦν ἐν τῇ πόλει ἀπὸ ἀνδρὸς καὶ ἕως γυναικὸς, ἀπὸ νεανίσκου καὶ ἕως πρεσβύτου, καὶ ἕως μόσχου καὶ ὑποζυγίου, ἐν στόματι ῥομφαίας.

\vs{22}Καὶ τοῖς δυσὶ νεανίσκοις τοῖς κατασκοπεύσασιν εἶπεν Ἰησοῦς, εἰσέλθατε εἰς τὴν οἰκίαν τῆς γυναικὸς, καὶ ἐξαγάγετε αὐτὴν ἐκεῖθεν, καὶ ὅσα ἐστὶν αὐτῇ.
\vs{23}Καὶ εἰσῆλθον οἱ δύο νεανίσκοι οἱ κατασκοπεύσαντες τὴν πόλιν, εἰς τὴν οἰκίαν τῆς γυναικὸς, καὶ ἐξηγάγοσαν Ῥαὰβ τὴν πόρνην, καὶ τὸν πατέρα αὐτῆς, καὶ τὴν μητέρα αὐτῆς, καὶ τοὺς ἀδελφοὺς αὐτῆς, καὶ τὴν συγγένειαν αὐτῆς, καὶ πάντα ὅσα ἦν αὐτῆ· καὶ κατέστησαν αὐτὴν ἔξω τῆς παρεμβολῆς Ἰσραήλ.
\vs{24}Καὶ ἡ πόλις ἐνεπρήσθη ἐν πυρισμῷ σὺν πᾶσι τοῖς ἐν αὐτῇ· πλὴν ἀργυρίου καὶ χρυσίου καὶ χαλκοῦ καὶ σιδήρου ἔδωκαν εἰς θησαυρὸν Κυρίου εἰσενεχθῆναι.

\vs{25}Καὶ Ῥαὰβ τὴν πόρνην, καὶ πάντα τὸν οἶκον αὐτῆς τὸν πατρικὸν ἐζώγρησεν Ἰησοῦς· καὶ κατῴκισεν ἐν τῷ Ἰσραὴλ ἕως τῆς σήμερον ἡμέρας, διότι ἔκρυψε τοὺς κατασκοπεύαντας, οὓς ἀπέστειλεν Ἰησοῦς κατασκοπεῦσαι τὴν Ἱεριχώ.
\vs{26}Καὶ ὥρκισεν Ἰησοῦς ἐν τῇ ἡμέρᾳ ἐκείνῃ ἐναντίον Κυρίου, λέγων, ἐπικατάρατος ὁ ἄνθρωπος, ὃς οἰκοδομήσει τὴν πόλιν ἐκείνην· ἐν τῷ πρωτοτόκῳ αὐτοῦ θεμελιώσει αὐτὴν, καὶ ἐν τῷ ἐλαχίστῳ αὐτοῦ ἐπιστήσει τὰς πύλας αὐτῆς. Καὶ οὕτως ἐποίησεν Ὁζᾶν ὁ ἐκ Βαιθήλ· ἐν τῷ Ἀβιρὼν τῷ πρωτοτόκῳ ἐθεμελίωσεν αὐτὴν, καὶ ἐν τῷ ἐλαχίστῳ διασωθέντι ἐπέστησε τὰς πύλας αὐτῆς.

\vs{27}Καὶ ἦν Κύριος μετὰ Ἰησοῦ, καὶ ἦν τὸ ὄνομα αὐτοῦ κατὰ πᾶσαν τὴν γῆν.

\ch{7}
Καὶ ἐπλημμέλησαν οἱ υἱοὶ Ἰσραὴλ πλημμέλιαν μεγάλην, καὶ ἐνοσφίσαντο ἀπὸ τοῦ ἀναθέματος· καὶ ἔλαβεν Ἄχαρ υἱὸς Χαρμὶ υἱοῦ Ζαμβρὶ υἱοῦ Ζαρὰ ἐκ τῆς φυλῆς Ἰούδα ἀπὸ τοῦ ἀναθέματος· καὶ ἐθυμώθη Κύριος ὀργῇ τοῖς υἱοῖς Ἰσραήλ.

\vs{2}Καὶ ἀπέστειλεν Ἰησοῦς ἄνδρας εἰς Γαὶ, ἥ ἐστι κατὰ Βαιθὴλ, λέγων, κατασκέψασθε τὴν Γαί
\vs{3}καὶ ἀνέβησαν οἱ ἄνδρες καὶ κατεσκέψαντο τὴν Γαί· Καὶ ἀνέστρεψαν πρὸς Ἰησοῦν, καὶ εἶπαν πρὸς αὐτὸν, μὴ ἀναβήτω πᾶς ὁ λαὸς, ἀλλʼ ὡσεὶ δισχίλιοι ἢ τρισχίλιοι ἄνδρες ἀναβήτωσαν καὶ ἐκπολιορκησάτωσαν τὴν πόλιν· μὴ ἀναγάγῃς ἐκεῖ τὸν λαὸν ἅπαντα, ὀλίγοι γάρ εἰσι.
\vs{4}Καὶ ἀνέβησαν ὡσεὶ τρισχίλιοι ἄνδρες, καὶ ἔφυγον ἀπὸ προσώπου ἀνδρῶν Γαί.
\vs{5}Καὶ ἀπέκτειναν ἀπʼ αὐτῶν ἄνδρες Γαὶ εἰς τριακονταὲξ ἄνδρας, καὶ κατεδίωξαν αὐτοὺς ἀπὸ τῆς πύλης, καὶ συνέτριψαν αὐτοὺς ἀπὸ τοῦ καταφεροῦς· καὶ ἐπτοήθη ἡ καρδία τοῦ λαοῦ, καὶ ἐγένετο ὥσπερ ὕδωρ.

\vs{6}Καὶ διέῤῥηξεν Ἰησοῦς τὰ ἱμάτια αὐτοῦ· καὶ ἔπεσεν Ἰησοῦς ἐπὶ τὴν γῆν ἐπὶ πρόσωπον ἐναντίον Κυρίου ἕως ἑσπέρας, αὐτὸς καὶ οἱ πρεσβύτεροι Ἰσραήλ· καὶ ἐπεβάλοντο χοῦν ἐπὶ τὰς κεφαλὰς αὐτῶν.
\vs{7}Καὶ εἶπεν Ἰησοῦς, δέομαι Κύριε· ἱνατί διεβίβασεν ὁ παῖς σου τὸν λαὸν τοῦτον τὸν Ἰορδάνην παραδοῦναι αὐτὸν τῷ Ἀμοῤῥαίῳ, ἀπολέσαι ἡμᾶς; καὶ εἰ κατεμείναμεν καὶ κατῳκίσθημεν παρὰ τὸν Ἰορδάνην.
\vs{8}Καὶ τί ἐρῶ ἐπεὶ μετέβαλεν Ἰσραὴλ αὐχένα ἀπέναντι τοῦ ἐχθροῦ αὐτοῦ;
\vs{9}Καὶ ἀκούσας ὁ Χαναναῖος καὶ πάντες οἱ κατοικοῦντες τὴν γῆν, περικυκλώσουσιν ἡμᾶς, καὶ ἐκτρίψουσιν ἡμᾶς ἀπὸ τῆς γῆς· καὶ τί ποιήσεις τὸ ὄνομά σου τὸ μέγα;

\vs{10}Καὶ εἶπε Κύριος πρὸς Ἰησοῦν, ἀνάστηθι, ἱνατί τοῦτο σὺ πέπτωκας ἐπὶ πρόσωπόν σου;
\vs{11}Ἡμάρτηκεν ὁ λαὸς καὶ παρέβη τὴν διαθήκην, ἣν διεθέμην πρὸς αὐτοὺς, κλέψαντες ἀπὸ τοῦ ἀναθέματος ἐνέβαλον εἰς τὰ σκεύη αὐτῶν.
\vs{12}Καὶ οὐ μὴ δύνωνται οἱ υἱοὶ Ἰσραὴλ ὑποστῆναι κατὰ πρόσωπον τῶν ἐχθρῶν αὐτῶν· αὐχένα ὑποστρέψουσιν ἔναντι τῶν ἐχθρῶν αὐτῶν, ὅτι ἐγενήθησαν ἀνάθεμα· οὐ προσθήσω ἔτι εἶναι μεθʼ ὑμῶν, ἐὰν μὴ ἐξάρητε τὸ ἀνάθεμα ἐξ ὑμῶν αὐτῶν.
\vs{13}Ἀναστὰς ἁγίασον τὸν λαὸν, καὶ εἶπον ἁγιασθῆναι εἰσαύριον· τάδε λέγει Κύριος ὁ Θεὸς Ἰσραὴλ, τὸ ἀνάθεμά ἐστιν ἐν ὑμῖν· οὐ δυνήσεσθε ἀντιστῆναι ἀπέναντι τῶν ἐχθρῶν ὑμῶν, ἕως ἂν ἐξάρητε τὸ ἀνάθεμα ἐξ ὑμῶν.
\vs{14}Καὶ συναχθήσεσθε πάντες τοπρωῒ κατὰ φυλὰς, καὶ ἔσται ἡ φυλὴ ἣν ἂν δείξῃ Κύριος, προσάξετε κατὰ δήμους· καὶ τὸν δῆμον ὃν ἂν δείξῃ Κύριος, προσάξετε κατʼ οἶκον· καὶ τὸν οἶκον ὃν ἂν δέξῃ Κύριος, προσάξετε κατʼ ἄνδρα.
\vs{15}Καὶ ὃς ἂν ἐνδειχθῇ, κατακαυθήσεται ἐν πυρὶ, καὶ πάντα ὅσα ἐστὶν αὐτῷ· ὅτι παρέβη τὴν διαθήκην Κυρίου, καὶ ἐποίησεν ἀνόμημα ἐν Ἰσραήλ.

\vs{16}Καὶ ὤρθρισεν Ἰησοῦς, καὶ προσήγαγε τὸν λαὸν κατὰ φυλάς· καὶ ἐνεδείχθη ἡ φυλὴ Ἰούδα.
\vs{17}Καὶ προσήχθη κατὰ δήμους, καὶ ἐνεδείχθη δῆμος Ζαραΐ. Καὶ προσήχθη κατʼ ἄνδρα,
\vs{18}καὶ ἐνεδείχθη Ἄχαρ υἱὸς Ζαμβρὶ υἱοῦ Ζάρά.

\vs{19}Καὶ εἶπεν Ἰησοῦς τῷ Ἄχαρ, δὸς δόξαν σήμερον τῷ Κυρίῳ Θεῷ Ἰσραὴλ, καὶ δὸς τὴν ἐξομολόγησιν, καὶ ἀνάγγειλόν μοι τί ἐποίησας, καὶ μὴ κρύψῃς ἀπʼ ἐμοῦ.
\vs{20}Καὶ ἀπεκρίθη Ἄχαρ τῷ Ἰησοῖ, καὶ εἶπεν, ἀληθῶς ἥμαρτον ἐναντίον Κυρίου τοῦ Θεοῦ Ἰσραήλ· οὕτως καὶ οὕτως ἐποίησα.
\vs{21}Εἶδον ἐν τῇ προνομῇ ψιλὴν ποικίλην, καὶ διακόσια δίδραχμα ἀργυρίου, καὶ γλῶσσαν μίαν χρυσῆν πεντήκοντα διδράχμων, καὶ ἐνθυμηθεὶς αὐτῶν ἔλαβον· καὶ ἰδοὺ αὐτὰ ἐγκέκρυπται ἐν τῇ σκηνῇ μου, καὶ τὸ ἀργύριον κέκρυπται ὑποκάτω αὐτῶν.
\vs{22}Καὶ ἀπέστειλεν Ἰησοῦς ἀγγέλους, καὶ ἔδραμον εἰς τὴν σκηνὴν εἰς τὴν παρεμβολήν· καὶ ταῦτα ἦν κεκρυμμένα εἰς τὴν σκηνὴν αὐτοῦ, καὶ τὸ ἀργύριον ὑποκάτω αὐτῶν.
\vs{23}Καὶ ἐξήνεγκαν αὐτὰ ἐκ τῆς σκηνῆς, καὶ ἤνεγκαν πρὸς Ἰησοῦν καὶ τοὺς πρεσβυτέρους Ἰσραὴλ, καὶ ἔθηκαν αὐτὰ ἔναντι Κυρίου.

\vs{24}Καὶ ἔλαβεν Ἰησοῦς τὸν Ἄχαρ υἱὸν Ζαρὰ, καὶ ἀνήγαγεν αὐτὸν εἰς φάραγγα Ἀχὼρ, καὶ τοὺς υἱοὺς αὐτοῦ, καὶ τὰς θυγατέρας αὐτοῦ, καὶ τοὺς μόσχους αὐτοῦ, καὶ τὰ ὑποζύγια αὐτοῦ, καὶ πάντα τὰ πρόβατα αὐτοῦ, καὶ τὴν σκηνὴν αὐτοῦ, καὶ πάντα τὰ ὑπάρχοντα αὐτοῦ, καὶ πᾶς ὁ λαὸς μετʼ αὐτοῦ· καὶ ἀνήγαγεν αὐτοὺς εἰς Ἐμεκαχώρ.
\vs{25}Καὶ εἶπεν Ἰησοῦς τῷ Ἄχαρ, τί ὠλόθρευσας ἡμᾶς; ἐξολοθρεύσαι σε Κύριος καθὰ καὶ σήμερον.
\vs{26}Καὶ ἐλιθοβόλησαν αὐτὸν λίθοις πᾶς Ἰσραὴλ, καὶ ἐπέστησαν αὐτῷ σωρὸν λίθων μέγαν· καὶ ἐπαύσατο Κύριος τοῦ θυμοῦ τῆς ὀργῆς. Διὰ τοῦτο ἐπωνόμασεν αὐτὸ Ἐμεκαχὼρ ἕως τῆς ἡμέρας ταύτης.

\ch{8}
Καὶ εἶπε Κύριος πρὸς Ἰησοῦν, μὴ φοβηθῇς, μηδὲ δειλιάσῃς· λάβε μετὰ σοῦ πάντας τοὺς ἄνδρας τοὺς πολεμιστάς, καὶ ἀναστὰς ἀνάβηθι εἰς Γαί· ἰδοὺ δέδωκα εἰς τὰς χεῖράς σου τὸν βασιλέα Γαὶ, καὶ τὴν γῆν αὐτοῦ.
\vs{2}Καὶ ποιήσεις τὴν Γαὶ, ὃν τρόπον ἐποίησας τὴν Ἱεριχὼ, καὶ τὸν βασιλέα αὐτῆς· καὶ τὴν προνομὴν τῶν κτηνῶν προνομεύσεις σεαυτῷ· κατάστησον δὲ σεαυτῷ ἔνεδρα τῇ πόλει εἰς τὰ ὀπίσω.

\vs{3}Καὶ ἀνέστη Ἰησοῦς καὶ πᾶς ὁ λαὸς ὁ πολεμιστὴς ὥστε ἀναβῆναι εἰς Γαί. ἐπέλεξε δὲ Ἰησοῦς τριάκοντα χιλιάδας ἀνδρῶν δυνατοὺς ἐν ἰσχύϊ, καὶ ἀπέστειλεν αὐτοὺς νυκτός.
\vs{4}Καὶ ἐνετείλατο αὐτοῖς, λέγων, ὑμεῖς ἐνεδρεύσατε ὀπίσω τῆς πόλεως· μὴ μακρὰν γίνεσθε ἀπὸ τῆς πόλεως, καὶ ἔσεσθε πάντες ἕτοιμοι.
\vs{5}Καὶ ἐγὼ καὶ πάντες οἱ μετʼ ἐμοῦ προσάξομεν πρὸς τὴν πόλιν· καὶ ἔσται ὡς ἂν ἐξέλθωσιν οἱ κατοικοῦντες Γαὶ εἰς συνάντησιν ἡμῖν, καθάπερ καὶ πρώην, καὶ φευξόμεθα ἀπὸ προσώπου αὐτῶν.
\vs{6}Καὶ ὡς ἂν ἐξέλθωσιν ὀπίσω ἡμῶν, ἀποσπάσομεν αὐτοὺς ἀπὸ τῆς πόλεως· καὶ ἐροῦσι, φεύγουσιν οὗτοι ἀπὸ προσώπου ἡμῶν, ὃν τρόπον καὶ ἔμπροσθεν.
\vs{7}Ὑμεῖς δὲ ἐξαναστήσεσθε ἐκ τῆς ἐνέδρας, καὶ πορεύσεσθε εἰς τὴν πόλιν.
\vs{8}Κατὰ τὸ ῥῆμα τοῦτο ποιήσετε· ἰδοὺ ἐντέταλμαι ὑμῖν.
\vs{9}Καὶ ἀπέστειλεν αὐτοὺς Ἰησοῦς, καὶ ἐπορεύθησαν εἰς τὴν ἔνεδραν· καὶ ἐνεκάθισαν ἀναμέσον Βαιθὴλ καὶ ἀναμέσον Γαὶ, ἀπὸ θαλάσσης τῆς Γαί.

\vs{10}Καὶ ὀρθρίσας Ἰησοῦς τοπρωῒ, ἐπεσκέψατο τὸν λαόν· καὶ ἀνέβησαν αὐτὸς καὶ οἱ πρεσβύτεροι κατὰ πρόσωπον τοῦ λαοῦ ἐπὶ Γαί.
\vs{11}Καὶ πᾶς ὁ λαὸς ὁ πολεμιστὴς μετʼ αὐτοῦ ἀνέβησαν· καὶ πορευόμενοι ἦλθον ἐξεναντίας τῆς πόλεως ἀπὸ ἀνατολῶν.
\vs{12}Καὶ τὰ ἔνεδρα τῆς πόλεως ἀπὸ θαλάσσης·
\vs{14}Καὶ ἐγένετο ὡς εἶδε βασιλεὺς Γαὶ, ἔσπευσε καὶ ἐξῆλθεν εἰς συνάντησιν αὐτοῖς ἐπʼ εὐθείας εἰς τὸν πόλεμον, αὐτὸς καὶ πᾶς ὁ λαὸς ὁ μετʼ αὐτοῦ· καὶ αὐτὸς οὐκ ᾔδει ὅτι ἔνεδρα αὐτῷ ἐστιν ὀπίσω τῆς πόλεως.
\vs{15}Καὶ εἶδε, καὶ ἀνεχώρησεν Ἰησοῦς καὶ Ἰσραὴλ ἀπὸ προσώπου αὐτῶν.
\vs{16}Καὶ κατεδίωξαν ὀπίσω τῶν υἱῶν Ἰσραήλ· καὶ αὐτοὶ ἀπέστησαν ἀπὸ τῆς πόλεως.
\vs{17}Οὐ κατελείφθη οὐδεὶς ἐν τῇ Γαὶ, ὃς οὐ κατεδίωξεν ὀπίσω Ἰσραήλ· καὶ κατέλιπον τὴν πόλιν ἠνεῳγμένην, καὶ κατεδίωξαν ὀπίσω Ἰσραήλ·

\vs{18}Καὶ εἶπε Κύριος πρὸς Ἰησοῦν, ἔκτεινον τὴν χεῖρά σου ἐν τῷ γαισῷ τῷ ἐν τῇ χειρί σου ἐπὶ τὴν πόλιν, εἰς γὰρ τὰς χεῖράς σου παραδέδωκα αὐτήν· καὶ τὰ ἔνεδρα ἐξαναστήσονται ἐν τάχει ἐκ τοῦ τόπου αὐτῶν. Καὶ ἐξέτεινεν Ἰησοῦς τὴν χεῖρα αὐτοῦ τὸν γαισὸν ἐπὶ τὴν πόλιν·
\vs{19}καὶ τὰ ἔνεδρα ἐξανέστησαν ἐν τάχει ἐκ τοῦ τόπου αὐτῶν· καὶ ἐξήλθοσαν ὅτε ἐξέτεινε τὴν χεῖρα, καὶ εἰσήλθοσαν καὶτὴν πόλιν, καὶ κατελάβοντο αὐτήν· καὶ σπεύσαντες ἐνέπρησαν τὴν πόλιν ἐν πυρί.
\vs{20}Καὶ περιβλέψαντες οἱ κάτοικοι Γαὶ εἰς τὰ ὀπίσω αὐτῶν, καὶ ἐθεώρουν καπνὸν ἀναβαίνοντα ἐκ τῆς πόλεως εἰς τὸν οὐρανόν· καὶ οὐκ ἔτι εἶχον ποῦ φύγωσιν ὧδε ἢ ὧδε.
\vs{21}Καὶ Ἰησοῦς καὶ πᾶς Ἰσραὴλ εἶδον, ὅτι ἔλαβον τὰ ἔνεδρα τὴν πόλιν, καὶ ὅτι ἀνέβη ὁ καπνὸς τῆς πόλεως εἰς τὸν οὐρανόν· καὶ μεταβαλλόμενοι, ἐπάταξαν τοὺς ἄνδρας τῆς Γαί.
\vs{22}Καὶ οὗτοι ἐξήλθοσαν ἐκ τῆς πόλεως εἰς συνάντησιν· καὶ ἐκενήθησαν ἀναμέσον τῆς παρεμβολῆς, οὗτοι ἐντεῦθεν καὶ οὗτοι ἐντεῦθεν· καὶ ἐπάταξαν αὐτοὺς ἕως τοῦ μὴ καταλειφθῆναι αὐτῶν σεσωσμένον καὶ διαπεφευγότα.
\vs{23}Καὶ τὸν βασιλέα τῆς Γαὶ συνέλαβον ζῶντα, καὶ προσήγαγον αὐτὸν πρὸς Ἰησοῦν.

\vs{24}Καὶ ὡς ἐπαύσαντο οἱ υἱοὶ Ἰσραὴλ ἀποκτείνοντες πάντας τοὺς ἐν τῇ Γαὶ, καὶ τοὺς ἐν τοῖς πεδίοις, καὶ ἐν τῷ ὄρει ἐπὶ τῆς καταβάσεως, οὗ κατεδίωξαν αὐτοὺς ἀπʼ αὐτῆς εἰς τέλος, καὶ ἐπέστρεψεν Ἰησοῦς εἰς Γαὶ, καὶ ἐπάταξεν αὐτὴν ἐν στόματι ῥομφαίας.
\vs{25}Καὶ ἐγενήθησαν οἱ πεσόντες ἐν τῇ ἡμέρᾳ ἐκείνῃ ἀπὸ ἄνδρος καὶ ἕως γυναικὸς, δώδεκα χιλιάδες, πάντας τοὺς κατοικοῦντας Γαί.
\vs{27}Πλὴν τῶν σκύλων τῶν ἐν τῇ πόλει πάντα, ἃ ἐπρονόμευσαν ἑαυτοῖς οἱ υἱοὶ Ἰσραήλ κατὰ πρόσταγμα Κυρίου, ὃν τρόπον συνέταξε Κύριος τῷ Ἰησοῖ.

\vs{28}Καὶ ἐνεπύρισεν Ἰησοῦς τὴν πόλιν ἐν πυρί· χῶμα ἀοίκητον εἰς τὸν αἰῶνα ἔθηκεν αὐτὴν ἕως τῆς ἡμέρας ταύτης.
\vs{29}Καὶ τὸν βασιλέα τῆς Γαὶ ἐκρέμασεν ἐπὶ ξύλου διδύμου· καὶ ἦν ἐπὶ τοῦ ξύλου ἕως ἑσπέρας· καὶ ἐπιδύνοντος τοῦ ἡλίου συνέταξεν Ἰησοῦς, καὶ καθείλοσαν τὸ σῶμα αὐτοῦ ἀπὸ τοῦ ξύλου, καὶ ἔῤῥιψαν αὐτὸ εἰς τὸν βόθρον· καὶ ἐπέστησαν αὐτῷ σωρὸν λίθων, ἕως τῆς ἡμέρας ταύτης.

\ch{9}
Ὡς δὲ ἤκουσαν οἱ βασιλεῖς τῶν Ἀμοῤῥαίων οἱ ἐν τῷ πέραν τοῦ Ἰορδάνου, οἱ ἐν τῇ ὀρεινῇ, καὶ οἱ ἐν τῇ πεδινῇ, καὶ οἱ ἐν πάσῃ τῇ παραλίᾳ τῆς θαλάσσης τῆς μεγάλης, καὶ οἱ πρὸς τῷ Ἀντιλιβάνῳ, καὶ οἱ Χετταῖοι, καὶ οἱ Χαναναῖοι, καὶ οἱ Φερεζαῖοι, καὶ οἱ Εὐαῖοι, καὶ οἱ Ἀμοῤῥαῖοι, καὶ οἱ Γεργεσαῖοι, καὶ οἱ Ἰεβουσαῖοι,
\vs{2}συνήλθοσαν ἐπὶ τὸ αὐτὸ ἐκπολεμῆσαι Ἰησοῦν καὶ Ἰσραὴλ ἅμα πάντες.

\vs{2a}Τότε ᾠκοδόμησεν Ἰησοῦς θυσιαστήριον Κυρίῳ τῷ Θεῷ Ἰσραὴλ ἐν ὄρει Γαιβὰλ,
\vs{2b}καθότι ἐνετείλατο Μωυσῆς ὁ θεράπων Κυρίου τοῖς υἱοῖς Ἰσραῆλ, καθὰ γέγραπται ἐν τῷ νόμῳ Μωυσῆ, θυσιαστήριον λίθων ὁλοκλήρων, ἐφʼ οὓς οὐκ ἐπεβλήθη σίδηρος· καὶ ἀνεβίβασεν ἐκεῖ ὁλοκαυτώματα Κυρίῳ, καὶ θυσίαν σωτηρίου.
\vs{2c}Καὶ ἔγραψεν Ἰησοῦς ἐπὶ τῶν λίθων τὸ δευτερονόμιον, νόμον Μωυσῆ, ἐνώπιον τῶν υἱῶν Ἰσραήλ.
\vs{2d}Καὶ πᾶς Ἰσραὴλ, καὶ οἱ πρεσβύτεροι αὐτῶν, καὶ οἱ δικασταὶ, καὶ οἱ γραμματεῖς αὐτῶν, παρεπορεύοντο ἔνθεν καὶ ἔνθεν τῆς κιβωτοῦ ἀπέναντι· καὶ οἱ ἱερεῖς καὶ οἱ Λευῖται ᾖρν τὴν κιβωτὸν τῆς διαθήκης Κυρίου· καὶ ὁ προσήλυτος καὶ ὁ αὐτόχθων, οἳ ἦσαν ἥμισυ πλησίον ὄρους Γαριζὶν, καὶ οἳ ἦσαν ἥμισυ πλησίον ὄρους Γαιβὰλ, καθότι ἐνετείλατο Μωυσῆς ὁ θεράπων Κυρίου εὐλογῆσαι τὸν λαὸν ἐν πρώτοις.

\vs{2e}Καὶ μετὰ ταῦτα οὕτως ἀνέγνω Ἰησοῦς πάντα τὰ ῥήματα τοῦ νόμου τούτου, τὰς εὐλογίας καὶ τὰς κατάρας, κατὰ πάντα τὰ γεγραμμένα ἐν τῷ νόμῳ Μωυσῆ.
\vs{2f}Οὐκ ἦν ῥῆμα ἀπὸ πάντων ὧν ἐνετείλατο Μωυσῆς τῷ Ἰησοῖ, ὃ οὐκ ἀνέγνω Ἰησοῦς εἰς τὰ ὦτα πάσης ἐκκλησίας υἱῶν Ἰσραὴλ, τοῖς ἀνδράσι καὶ ταῖς γυναιξὶ καὶ τοῖς παιδίοις καὶ τοῖς προσηλύτοις τοῖς προσπορευομένοις τῷ Ἰσραήλ.

\vs{3}Καὶ οἱ κατοικοῦντες Γαβαὼν ἤκουσαν πάντα ὅσα ἐποίησε Κύριος τῇ Ἱεριχὼ καὶ τῇ Γαί.
\vs{4}Καὶ ἐποίησαν καί γε αὐτοὶ μετὰ πανουργίας· καὶ ἐλθόντες ἐπεσιτίσαντο καὶ ἡτοιμάσαντο· καὶ λαβόντες σάκκους παλαιοὺς ἐπὶ τῶν ὤμων αὐτῶν, καὶ ἀσκοὺς οἴνου παλαιοὺς καὶ κατεῤῥωγότας ἀποδεδεμένους,
\vs{5}καὶ τὰ κοῖλα τῶν ὑποδημάτων αὐτῶν, καὶ τὰ σανδάλια αὐτῶν παλαιὰ καὶ καταπεπελματωμένα ἐν τοῖς ποσὶν αὐτῶν, καὶ τὰ ἱμάτια αὐτῶν πεπαλαιωμένα ἐπάνω αὐτῶν, καὶ ὁ ἄρτος αὐτῶν τοῦ ἐπισιτισμοῦ ξηρὸς καὶ εὐρωτιῶν καὶ βεβρωμένος.

\vs{6}Καὶ ἤλθοσαν πρὸς Ἰησοῦν εἰς τὴν παρεμβολὴν Ἰσραὴλ εἰς Γάλγαλα, καὶ εἶπαν πρὸς Ἰησοῦν καὶ Ἰσραήλ, ἐκ γῆς μακρόθεν ἥκαμεν· καὶ νῦν διάθεσθε ἡμῖν διαθήκην.
\vs{7}Καὶ εἶπαν οἱ υἱοὶ Ἰσραὴλ πρὸς τὸν Χοῤῥαῖον, ὅρα μὴ ἐν ἐμοὶ κατοικεῖς· καὶ πῶς σοι διαθῶμαι διαθήκην;
\vs{8}Καὶ εἶπαν πρὸς Ἰησοῦν, οἰκεται σου ἐσμέν· καὶ εἶπε πρὸς αὐτοὺς Ἰησοῦς, πόθεν ἐστὲ, καὶ πόθεν παραγεγόνατε;
\vs{9}Καὶ εἶπαν, ἐκ γῆς μακρόθεν σφόδρα ἥκασιν οἱ παῖδές σου ἐν ὀνόματι Κυρίου τοῦ Θεοῦ σου· ἀκηκόαμεν γὰρ τὸ ὄνομα αὐτοῦ, καὶ ὅσα ἐποίησεν ἐν Αἰγύπτῳ,
\vs{10}καὶ ὅσα ἐποίησε τοῖς βασιλεῦσι τῶν Ἀμοῤῥαίων, οἳ ἦσαν πέραν τοῦ Ἰορδάνου, τῷ Σηὼν βασιλεῖ τῶν Ἀμοῤῥαίων, καὶ τῷ Ὢγ βασιλεῖ τῆς Βασὰν, ὃς κατῴκει ἐν Ἀσταρὼθ καὶ ἐν Ἐδραΐν.
\vs{11}Καὶ ἀκούσαντες εἶπαν πρὸς ἡμᾶς οἱ πρεσβύτεροι ἡμῶν καὶ πάντες οἱ κατοικοῦντες τὴν γῆν ἡμῶν, λέγοντες, λάβετε ἑαυτοῖς ἐπισιτισμὸν εἰς τὴν ὁδὸν, καὶ πορεύθητε εἰς συνάντησιν αὐτῶν, καὶ ἐρεῖτε πρὸς αὐτοὺς, οἰκέται σου ἐσμὲν, καὶ νῦν διάθεσθε ἡμῖν τὴν διαθήκην.
\vs{12}Οὗτοι οἱ ἄρτοι, θερμοὺς ἐφωδιάσθημεν αὐτοὺς ἐν τῇ ἡμέρᾳ ᾗ ἐξήλθομεν παραγενέσθαι πρὸς ὑμᾶς· νῦν δὲ ἐξηράνθησαν καὶ γεγόνασι βεβρωμένοι.
\vs{13}Καὶ οὗτοι οἱ ἀσκοὶ τοῦ οἴνου οὓς ἐπλήσαμεν καινοὺς, καὶ οὗτοι ἐῤῥώγασι· καὶ τὰ ἱμάτια ἡμῶν, καὶ τὰ ὑποδήματα ἡμῶν πεπαλαίωται ἀπὸ τῆς πολλῆς ὁδοῦ σφόδρα.

\vs{14}Καὶ ἔλαβον οἱ ἄρχοντες τοῦ ἐπισιτισμοῦ αὐτῶν, καὶ Κύριον οὐκ ἐπηρώτησαν.
\vs{15}Καὶ ἐποίησεν Ἰησοῦς πρὸς αὐτοὺς εἰρήνην, καὶ διέθεντο πρὸς αὐτοὺς διαθήκην τοῦ διασῶσαι αὐτούς· καὶ ὤμοσαν αὐτοῖς οἱ ἄρχοντες τῆς συναγωγῆς.

\vs{16}Καὶ ἐγένετο μετὰ τρεῖς ἡμέρας μετὰ τὸ διαθέσθαι πρὸς αὐτοὺς διαθήκην, ἤκουσαν ὅτι ἐγγύθεν αὐτῶν εἰσι, καὶ ὅτι ἐν αὐτοῖς κατοικοῦσι.
\vs{17}Καὶ ἀπῇραν οἱ υἱοὶ Ἰσραὴλ, καὶ ἦλθον εἰς τὰς πόλεις αὐτῶν· αἱ δὲ πόλεις αὐτῶν Γαβαὼν καὶ Κεφιρὰ καὶ Βηρὼτ, καὶ πόλεις Ἰαρίν.
\vs{18}Καὶ οὐκ ἐμαχέσαντο αὐτοῖς οἱ υἱοὶ Ἰσραὴλ, ὅτι ὤμοσαν αὐτοῖς πάντες οἱ ἄρχοντες Κύριον τὸν Θεὸν Ἰσραήλ· καὶ διεγόγγυσαν πᾶσα ἡ συναγωγὴ ἐπὶ τοῖς ἄρχουσι.

\vs{19}Καὶ εἶπαν οἱ ἄρχοντες πάσῃ τῇ συναγωγῇ, ἡμεῖς ὠμόσαμεν αὐτοῖς Κύριον τὸν Θεὸν Ἰσραὴλ, καὶ νῦν οὐ δυνησόμεθα ἅψασθαι αὐτῶν.
\vs{20}Τοῦτο ποιήσομεν, ζωγρῆσαι αὐτοὺς, καὶ περιποιησόμεθα αὐτούς· καὶ οὐκ ἔσται καθʼ ἡμῶν ὀργὴ διὰ τὸν ὅρκον, ὃν ὠμόσαμεν αὐτοῖς.
\vs{21}Ζήσονται, καὶ ἔσονται ξυλοκόποι καὶ ὑδροφόροι πάσῃ τῇ συναγωγῇ, καθάπερ εἶπαν αὐτοῖς οἱ ἄρχοντες.

\vs{22}Καὶ συνεκάλεσεν αὐτοὺς Ἰησοῦς, καὶ εἶπεν αὐτοῖς, διατί παρελογίσασθέ με, λέγοντες, μακρὰν ἀπὸ σοῦ ἐσμὲν σφόδρα· ὑμεῖς δὲ ἐγχώριοί ἐστε τῶν κατοικούντων ἐν ἡμῖν;
\vs{23}Καὶ νῦν ἐπικατάρατοί ἐστε· οὐ μὴ ἐκλείπῃ ἐξ ὑμῶν δοῦλος, οὐδὲ ξυλοκόπος, οὐδὲ ὑδροφόρος ἐμοὶ καὶ τῷ Θεῷ μου.
\vs{24}Καὶ ἀπεκρίθησαν τῷ Ἰησοῖ, λέγοντες, ἀνηγγέλη ἡμῖν ὅσα συνέταξε Κύριος ὁ Θεός σου Μωυσῇ τῷ παιδὶ αὐτοῦ, δοῦναι ὑμῖν τὴν γῆν ταύτην, καὶ ἐξολοθρεῦσαι ἡμᾶς καὶ πάντας τοὺς κατοικοῦντας ἐπʼ αὐτῆς ἀπὸ προσώπου ὑμῶν· καὶ ἐφοβήθημεν σφόδρα περὶ τῶν ψυχῶν ἡμῶν ἀπὸ προσώπου ὑμῶν, καὶ ἐποιήσαμεν τὸ πρᾶγμα τοῦτο.
\vs{25}Καὶ νῦν ἰδοὺ ἡμεῖς ὑποχείριοι ὑμῖν· ὡς ἀρέσκει ὑμῖν καὶ ὡς δοκεῖ ὑμῖν, ποιήσατε ἡμῖν.

\vs{26}Καὶ ἐποίησαν αὐτοῖς οὕτως· καὶ ἐξείλατο αὐτοὺς Ἰησοῦς ἐν τῇ ἡμέρᾳ ἐκείνῃ ἐκ χειρῶν υἱῶν Ἰσραὴλ, καὶ οὐκ ἀνεῖλον αὐτούς.
\vs{27}Καὶ κατέστησεν αὐτοὺς Ἰησοῦς ἐν τῇ ἡμέρᾳ ἐκείνῃ ξυλοκόπους καὶ ὑδροφόρους πάσῃ τῇ συναγωγῇ, καὶ τῷ θυσιαστηρίῳ τοῦ Θεοῦ· διὰ τοῦτο ἐγένοντο οἱ κατοικοῦντες Γαβαὼν ξυλοκόποι καὶ ὑδροφόροι τοῦ θυσιαστηρίου τοῦ θεοῦ ἕως τῆς σήμερον ἡμέρας, καὶ εἰς τὸν τόπον ὃν ἂν ἐκλέξηται Κύριος.

\ch{10}
Ὡς δὲ ἤκουσεν Ἀδωνιβεζέκ βασιλεὺς Ἱερουσαλὴμ ὅτι ἔλαβεν Ἰησοῦς τὴν Γαὶ, καὶ ἐξωλόθρευσεν αὐτὴν, ὃν τρόπον ἐποίησαν τὴν Ἱεριχὼ καὶ τὸν βασιλέα αὐτῆς, οὕτως ἐποίησαν καὶ τὴν Γαὶ καὶ τὸν βασιλέα αὐτῆς, καὶ ὅτι ηὐτομόλησαν οἱ κατοικοῦντες Γαβαὼν πρὸς Ἰησοῦν καὶ πρὸς Ἰσραὴλ,
\vs{2}καὶ ἐφοβήθησαν ἀπʼ αὐτῶν σφόδρα· ᾔδει γὰρ ὅτι πόλις μεγάλη Γαβαὼν, ὡσεὶ μία τῶν μητροπόλεων, καὶ πάντες οἱ ἄνδρες αὐτῆς ἰσχυροί.
\vs{3}Καὶ ἀπέστειλεν Ἀδωνιβεζὲκ βασιλεὺς Ἰερουσαλὴμ πρὸς Ἐλὰμ βασιλέα Χεβρὼν, καὶ πρὸς Φιδὼν βασιλέα Ἱερειμοὺθ, καὶ πρὸς Ἰεφθα βασιλέα Λαχὶς καὶ πρὸς Δαβεὶν βασιλέα Ὀδολλὰμ, λέγων,
\vs{4}δεῦτε, ἀνάβητε πρός με, καὶ βοηθήσατέ μοι, καὶ ἐκπολεμήσωμεν Γαβαών· ηὐτομόλησαν γὰρ πρὸς Ἰησοῦν καὶ πρὸς τοὺς υἱοὺς Ἰσραήλ.
\vs{5}Καὶ ἀνέβησαν οἱ πέντε βασιλεῖς τῶν Ἰεβουσαίων, βασιλεὺς Ἱερουσαλὴμ, καὶ βασιλεὺς Χεβρὼν, καὶ βασιλεὺς Ἱεριμοὺθ, καὶ βασιλεὺς Λαχὶς, καὶ βασιλεὺς Ὀδολλάμ, αὐτοὶ καὶ πᾶς ὁ λαὸς αὐτῶν. καὶ περιεκάθισαν τὴν Γαβαὼν, καὶ ἐξεπολιόρκουν αὐτήν.

\vs{6}Καὶ ἀπέστειλαν οἱ κατοικοῦντες Γαβαὼν πρὸς Ἰησοῦν εἰς τὴν παρεμβολὴν Ἰσραὴλ εἰς Γάλγαλα, λέγοντες, μὴ ἐκλύσῃς τὰς χεῖράς σου ἀπὸ τῶν παίδων σου· ἀνάβηθι πρὸς ἡμᾶς, τοτάχος, και βοήθησον ἡμῖν, καὶ ἐξελοῦ ἡμᾶς· ὅτι συνηγμένοι εἰσὶν ἐφʼ ἡμᾶς πάντες οἱ βασιλεῖς τῶν Ἀμοῤῥαίων, οἱ κατοικοῦντες τὴν ὀρεινὴν.
\vs{7}Καὶ ἀνέβη Ἰησοῦς ἐκ Γαλγάλων, αὐτὸς καὶ πᾶς ὁ λαὸς ὁ πολεμιστὴς μετʼ αὐτοῦ, πᾶς δυνατὸς ἐν ἰσχύϊ.

\vs{8}Καὶ εἶπε Κύριος πρὸς Ἰησοῦν, μὴ φοβηθῇς αὐτούς, εἰς γὰρ τὰς χεῖράς σου παραδέδωκα αὐτούς· οὐχ ὑπολειφθήσεται ἐξ αὐτῶν οὐδεὶς ἐνώπιον ὑμῶν.

\vs{9}Καὶ ἐπεὶ παρεγένετο Ἰησοῦς ἐπʼ αὐτοὺς ἄφνω, ὅλην τὴν νύκτα εἰσεπορεύθη ἐκ Γαλγάλων.
\vs{10}Καὶ ἐξέστησεν αὐτοὺς Κύριος ἀπὸ προσώπου τῶν υἱῶν Ἰσραήλ· καὶ συνέτριψεν αὐτοὺς Κύριος συντρίψει μεγάλῃ ἐν Γαβαών· καὶ κατεδίωξαν αὐτοὺς ὁδὸν ἀναβάσεως Ὠρωνείν, καὶ κατέκοπτον αὐτοὺς ἕως Ἀζηκὰ καὶ ἕως Μακηδά.
\vs{11}Ἐν δὲ τῷ φεύγειν αὐτοὺς ἀπὸ προσώπου τῶν υἱῶν Ἰσραὴλ ἐπὶ τὴς καταβάσεως Ὡρωνεὶν, καὶ Κύριος ἐπέῤῥιψεν αὐτοῖς λίθους χαλάζης ἐκ τοῦ οὐρανοῦ ἕως Ἀζηκά· καὶ ἐγένοντο πλείους οἱ ἀποθανόντες διὰ τοὺς λίθους τῆς χαλάζης, ἢ οὓς ἀπέκτειναν οἱ υἱοὶ Ἰσραὴλ μαχαίρᾳ ἐν τῷ πολέμῳ.

\vs{12}Τότε ἐλάλησεν Ἰησοῦς πρὸς Κύριον, ᾗ ἡμέρᾳ παρέδωκεν ὁ Θεὸς τὸν Ἀμοῤῥαῖον ὑποχείριον Ἰσραὴλ, ἡνίκα συνέτριψεν αὐτοὺς ἐν Γαβαὼν, καὶ συνετρίβησαν ἀπὸ προσώπου υἱῶν Ἰσραήλ· καὶ εἶπεν Ἰησοῦς, στήτω ὁ ἥλιος κατὰ Γαβαὼν, καὶ ἡ σελήνη κατὰ φάραγγα Αἰλών.
\vs{13}Καὶ ἔστη ὁ ἥλιος καὶ ἡ σελήνη ἐν στάσει, ἕως ἠμύνατο ὁ Θεὸς τοὺς ἐχθροὺς αὐτῶν· καὶ ἔστη ὁ ἥλιος κατὰ μέσον τοῦ οὐρανοῦ· οὐ προεπορεύετο εἰς δυσμὰς εἰς τέλος ἡμέρας μιᾶς.
\vs{14}Καὶ οὐκ ἐγένετο ἡμέρα τοιαύτη οὐδὲ τὸ πρότερον οὐδὲ τὸ ἔσχατον, ὥστε ἐπακοῦσαι Θεὸν ἀνθρώπου, ὅτι Κύριος συνεξεπολέμησε τῷ Ἰσραήλ.

\vs{16}Καὶ ἔφυγον οἱ πέντε βασιλεῖς οὗτοι, καὶ κατεκρύβησαν εἰς τὸ σπήλαιον τὸ ἐν Μακηδά.
\vs{17}Καὶ ἀπηγγέλη τῷ Ἰησοῦ, λέγοντες, εὕρηνται οἱ πέντε βασιλεῖς κεκρυμμένοι ἐν τῷ σπηλαίῳ τῷ ἐν Μακηδά.
\vs{18}Καὶ εἶπεν Ἰησοῦς, κυλίσατε λίθους ἐπὶ τὸ στόμα τοῦ σπηλαίου, καὶ καταστήσατε ἄνδρας φυλάσσειν ἐπʼ αὐτούς.
\vs{19}Ὑμεῖς δὲ μὴ ἑστήκατε, καταδιώκοντες ὀπίσω τῶν ἐχθρῶν ὑμῶν, καὶ καταλάβετε τὴν οὐραγίαν αὐτῶν, καὶ μὴ ἀφῆτε εἰσελθεῖν εἰς τὰς πόλεις αὐτῶν· παρέδωκε γὰρ αὐτοὺς Κύριος ὁ Θεὸς ἡμῶν εἰς τὰς χεῖρας ἡμῶν.
\vs{20}Καὶ ἐγένετο ὡς κατέπαυσεν Ἰησοῦς καὶ πᾶς υἱὸς Ἰσραὴλ κόπτοντες αὐτοὺς κοπὴν μεγάλην σφόδρα ἕως εἰς τέλος, καὶ οἱ διασωζόμενοι διεσώθησαν εἰς τὰς πόλεις τὰς ὀχυράς.

\vs{21}Καὶ ἀπεστράφη πᾶς ὁ λαὸς πρὸς Ἰησοῦν εἰς Μακηδὰ ὑγιεῖς· καὶ οὐκ ἔγρυξεν οὐδεὶς τῶν υἱῶν Ἰηραὴλ τῇ γλώσσῃ αὐτοῦ.

\vs{22}Καὶ εἶπεν Ἰησοῦς, ἀνοίξατε τὸ σπήλαιον, καὶ ἐξαγάγετε τοὺς πέντε βασιλεῖς τούτους ἐκ τοῦ σπηλαίου.
\vs{23}Καὶ ἐξηγάγοσαν τοὺς πέντε βασιλεῖς ἐκ τοῦ σπηλαίου, τὸν βασιλέα Ἱερουσαλὴμ, καὶ τὸν βασιλέα Χεβρὼν, καὶ τὸν βασιλέα Ἱεριμοὺθ, καὶ τὸν βασιλέα Λαχὶς, καὶ τὸν βασιλέα Ὀδολλάμ.
\vs{24}Καὶ ἐπεὶ ἐξήγαγον αὐτοὺς πρὸς Ἰησοῦν, καὶ συνεκάλεσεν Ἰησοῦς πάντα Ἰσραὴλ, καὶ τοὺς ἐναρχομένους τοῦ πολέμου τοὺς συμπορευομένους αὐτῷ, λέγων αὐτοῖς, προπορεύεσθε καὶ ἐπίθετε τοὺς πόδας ὑμῶν ἐπὶ τοὺς τραχήλους αὐτῶν· καὶ προσελθόντες ἐπέθηκαν τοὺς πόδας αὐτῶν ἐπὶ τοὺς τραχήλους αὐτῶν.
\vs{25}Καὶ εἶπεν Ἰησοῦς πρὸς αὐτοὺς, μὴ φοβηθῆτε αὐτοὺς, μηδὲ δειλιάσητε, ἀνδρίζεσθε καὶ ἰσχύετε, ὅτι οὕτω ποιήσει Κύριος πᾶσι τοῖς ἐχθροῖς ὑμῶν, οὓς ὑμεῖς καταπολεμεῖτε αὐτούς.
\vs{26}Καὶ ἀπέκτεινεν αὐτοὺς Ἰησοῦς, καὶ ἐκρέμασεν αὐτοὺς ἐπὶ πέντε ξύλων· καὶ ἦσαν κρεμάμενοι ἐπὶ τῶν ξύλων ἕως ἑσπέρας.
\vs{27}Καὶ ἐγενήθη πρὸς ἡλίου δυσμὰς, ἐνετείλατο Ἰησοῦς, καὶ καθεῖλον αὐτοὺς ἀπὸ τῶν ξύλων, καὶ ἔῤῥιψαν αὐτοὺς εἰς τὸ σπήλαιον, εἰς ὃ κατεφύγοσαν ἐκεῖ, καὶ ἐπεκύλισαν λίθους ἐπὶ τὸ σπήλαιον ἕως τῆς σήμερον ἡμέρας.

\vs{28}Καὶ τὴν Μακηδὰ ἐλάβοσαν ἐν τῇ ἡμέρᾳ ἐκείνῃ, καὶ ἐφόνευσαν αὐτὴν ἐν στόματι ξίφους, καὶ ἐξωλόθρευσαν πᾶν ἐμπνέον ὃ ἦν ἐν αὐτῇ· καὶ οὐ κατελείφθη οὐδεὶς ἐν αὐτῇ διασεσωσμένος καὶ διαπεφευγώς· καὶ ἐποίησαν τῷ βασιλεῖ Μακηδὰ, ὃν τρόπον ἐποίησαν τῷ βασιλεῖ Ἱεριχώ.

\vs{29}Καὶ ἀπῆλθεν Ἰησοῦς καὶ πᾶς Ἰσραὴλ μετʼ αὐτοῦ ἐκ Μακηδὰ εἰς Λεβνὰ, καὶ ἐπολιόρκει Λεβνά.
\vs{30}Καὶ παρέδωκεν αὐτὴν Κύριος εἰς χεῖρας Ἰσραήλ· καὶ ἔλαβον αὐτὴν, καὶ τὸν βασιλέα αὐτῆς, καὶ ἐφόνευσαν αὐτὴν ἐν στόματι ξίφους, καὶ πᾶν ἐμπνέον ἐν αὐτῇ· καὶ οὐ κατελείφθη ἐν αὐτῇ διασεσωσμένος καὶ διαπεφευγώς· καὶ ἐποίησαν τῷ βασιλεῖ αὐτῆς, ὃν τρόπον ἐποίησαν τῷ βασιλεῖ Ἱεριχώ.

\vs{31}Καὶ ἀπῆλθεν Ἰησοῦς καὶ πᾶς Ἰσραὴλ μετʼ αὐτοῦ ἐκ Λεβνὰ εἰς Λαχὶς, καὶ περιεκάθισεν αὐτὴν, καὶ ἐπολιόρκει αὐτήν.
\vs{32}Καὶ παρέδωκε Κύριος τὴν Λαχὶς εἰς τὰς χεῖρας Ἰσραήλ. καὶ ἔλαβεν αὐτὴν ἐν τῇ ἡμέρᾳ τῇ δευτέρᾳ, καὶ ἐφόνευσαν αὐτὴν ἐν στόματι ξίφους, καὶ ἐξωλόθρευσαν αὐτὴν, ὃν τρόπον ἐποίησαν τὴν Λεβνά.
\vs{33}Τότε ἀνέβη Ἐλὰμ βασιλεὺς Γαζὲρ βοηθήσων τῇ Λαχείς· καὶ ἐπάταξεν αὐτὸν Ἰησοῦς ἐν στόματι ξίφους, καὶ τὸν λαὸν αὐτοῦ, ἕως τοῦ μὴ καταλειφθῆναι αὐτῶν σεσωσμένον καὶ διαπεφευγότα.

\vs{34}Καὶ ἀπῆλθεν Ἰησοῦς καὶ πᾶς Ἰσραὴλ μετʼ αὐτοῦ ἐκ Λαχὶς εἰς Ὀδολλὰμ, καὶ περιεκάθισεν αὐτὴν καὶ ἐξεπολιόρκησεν αὐτήν.
\vs{35}Καὶ παρέδωκεν αὐτὴν Κύριος ἐν χειρὶ Ἰσραήλ· καὶ ἔλαβεν αὐτὴν ἐν τῇ ἡμέρᾳ ἐκείνῃ, καὶ ἐφόνευσεν αὐτὴν ἐν στόματι ξίφους, καὶ πᾶν ἐμπνέον ἐν ἀτῇ ἐφόνευσαν, ὃν τρόπον ἐποίησαν τῇ Λαχίς.

\vs{36}Καὶ ἀπῆλθεν Ἰησοῦς καὶ πᾶς Ἰσραὴλ μετʼ αὐτοῦ εἰς Χεβρὼν, καὶ περιεκάθισεν αὐτήν.
\vs{37}Καὶ ἐπάταξεν αὐτὴν ἐν στόματι ξίφους, καὶ πᾶν τὸ ἐμπνέον ὅσα ἦν ἐν αὐτῇ· οὐκ ἦν διασεσωσμένος· ὃν τρόπον ἐποίησαν τὴν Ὀδολλὰμ, ἐξωλόθρευσαν αὐτὴν, καὶ ὅσα ἦν ἐν αὐτῇ.

\vs{38}Καὶ ἀπέστρεψεν Ἰησοῦς καὶ πᾶς Ἰσραὴλ εἰς Δαβίρ· καὶ περικαθίσαντες αὐτὴν,
\vs{39}ἔλαβον αὐτὴν, καὶ τὸν βασιλέα αὐτῆς, καὶ τὰς κώμας αὐτῆς· καὶ ἐπάταξεν αὐτὴν ἐν στόματι ξίφους, καὶ ἐξωλόθρευσαν αὐτὴν, καὶ πᾶν ἐμπνέον ἐν αὐτῇ· καὶ οὐ κατέλιπον αὐτῇ οὐδένα διασεσωσμένον· ὃν τρόπον ἐποίησαν τῇ Χεβρὼν καὶ τῷ βασιλεῖ αὐτῆς, οὕτως ἐποίησαν τῇ Δαβὶρ καὶ τῷ βασιλεῖ αὐτῆς.

\vs{40}Καὶ ἐπάταξεν Ἰησοῦς πᾶσαν τὴν γῆν τῆς ὀρεινῆς, καὶ τὴν Ναγὲβ, καὶ τὴν πεδινὴν, καὶ τὴν Ἀσηδὼθ, καὶ τοὺς βασιλεῖς αὐτῆς· οὐ κατέλιπον αὐτῶν σεσωσμένον· καὶ πᾶν ἐμπνέον ζωῆς ἐξωλόθρευσεν, ὃν τρόπον ἐνετείλατο Κύριος ὁ Θεὸς Ἰσραὴλ,
\vs{41}ἀπὸ Κάδης Βαρνῆ ἕως Γάζης πᾶσαν τὴν Γοσὸμ ἕως τῆς Γαβαών.
\vs{42}Καὶ πάντας τοὺς βασιλεῖς αὐτῶν, καὶ τὴν γῆν αὐτῶν ἐπάταξεν Ἰησοῦς εἰσάπαξ· ὅτι Κύριος ὁ Θεὸς Ἰσραὴλ συνεπολέμει τῷ Ἰσραήλ.

\ch{11}
Ὡς δὲ ἤκουσεν Ἰαβὶς βασιλεὺς Ἀσὼρ, ἀπέστειλε πρὸς Ἰωβὰβ βασιλέα Μαρῶν, καὶ πρὸς βασιλέα Συμοὼν, καὶ πρὸς βασιλέα Ἀζὶφ,
\vs{2}καὶ πρὸς βασιλεῖς τοὺς κατὰ Σιδῶνα τὴν μεγάλην, εἰς τὴν ὀρεινὴν καὶ εἰς Ἄραβα ἀπέναντι Κενερὼθ, καὶ εἰς τὸ πεδίον, καὶ εἰς Φεναεδδὼρ,
\vs{3}καὶ εἰς τοὺς παραλίους Χαναναίους ἀπὸ ἀνατολῶν, καὶ εἰς τοὺς παραλίους Ἀμοῤῥαίους, καὶ τοὺς Χετταίους, καὶ Φερεζαίους, καὶ Ἰεβουσαίους τοὺς ἐν τῷ ὄρει, καὶ τοὺς Εὐαίους, καὶ τοὺς ὑπὸ τὴν Ἀερμὼν εἰς γῆν Μασσύμα.
\vs{4}Καὶ ἐξῆλθον αὐτοὶ καὶ οἱ βασιλεῖς αὐτῶν μετʼ αὐτῶν, ὥσπερ ἡ ἄμμος τῆς θαλάσσης τῷ πλήθει, καὶ ἵπποι καὶ ἅρματα πολλὰ σφόδρα.
\vs{5}Καὶ συνῆλθον πάντες οἱ βασιλεῖς αὐτοὶ καὶ παρεγένοντο ἐπὶ τὸ αὐτὸ, καὶ παρενέβαλον ἐπὶ τοῦ ὕδατος Μαρὼν πολεμῆσαι τὸν Ἰσραήλ.

\vs{6}Καὶ εἶπε Κύριος πρὸς Ἰησοῦν, μὴ φοβηθῇς ἀπὸ προσώπου αὐτῶν, ὅτι αὔριον ταύτην τὴν ὥραν ἐγὼ παραδίδωμι τετροπωμένους αὐτοὺς ἐναντίον τοῦ Ἰσραήλ· τοὺς ἵππους αὐτῶν νευροκοπήσεις, καὶ τὰ ἅρματα αὐτῶν κατακαύσεις ἐν πυρί.
\vs{7}Καὶ ἦλθεν Ἰησοῦς καὶ πᾶς ὁ λαὸς ὁ πολεμιστὴς ἐπʼ αὐτοὺς ἐπὶ τὸ ὕδωρ Μαρὼν ἐξάπινα· καὶ ἐπέπεσαν ἐπʼ αὐτοὺς ἐν τῇ ὀρεινῇ.
\vs{8}Καὶ παρέδωκεν αὐτοὺς Κύριος ὑποχειρίους Ἰσραήλ· καὶ κόπτοντες αὐτοὺς κατεδίωκον ἕως Σιδῶνος τῆς μεγάλης, καὶ ἕως Μασερὼν, καὶ ἕως τῶν πεδίων Μασσὼχ κατʼ ἀνατολάς· καὶ κατέκοψαν αὐτοὺς ἕως τοῦ μὴ καταλειφθῆναι αὐτῶν διασεσωσμένον.
\vs{9}Καὶ ἐποίησεν αὐτοῖς Ἰησοῦς, ὃν τρόπον ἐνετείλατο αὐτῷ Κύριος· τοὺς ἵππους αὐτῶν ἐνευροκόπησε, καὶ τὰ ἅρματα αὐτῶν ἐνέπρησε πυρί.

\vs{10}Καὶ ἐπεστράφη Ἰησοῦς ἐν τῷ καιρῷ ἐκείνῳ, καὶ κατελάβετο Ἀσὼρ, καὶ τὸν βασιλέα αὐτῆς· ἦν δὲ Ἀσὼρ τοπρότερον ἄρχουσα πασῶν τῶν βασιλειῶν τούτων.
\vs{11}Καὶ ἀπέκτειναν πᾶν ἐμπνέον ἐν αὐτῇ ἐν ξίφει, καὶ ἐξωλόθρευσαν πάντας, καὶ οὐ κατελείφθη ἐν αὐτῇ ἐμπνέον· καὶ τὴν Ἀσὼρ ἐνέπρησαν ἐν πυρί.
\vs{12}Καὶ πάσας τὰς πόλεις τῶν βασιλειῶν, καὶ τοὺς βασιλεῖς αὐτῶν ἔλαβεν Ἰησοῦς, καὶ ἀνεῖλεν αὐτοὺς ἐν στόματι ξίφους· καὶ ἐξωλόθρευσαν αὐτοὺς, ὃν τρόπον συνέταξε Μωυσῆς ὁ παῖς Κυρίου.
\vs{13}Ἀλλὰ πάσας τὰς πόλεις τὰς κεχωματισμένας οὐκ ἐνέπρησεν Ἰσραήλ· πλὴν Ἀσὼρ μόνην ἐνέπρησεν Ἰσραὴλ,
\vs{14}καὶ πάντα τὰ σκῦλα αὐτῆς ἐπρονόμευσαν ἑαυτοῖς οἱ υἱοὶ Ἰσραήλ· αὐτοὺς δὲ πάντας ἐξωλόθρευσαν ἐν στόματι ξίφους, ἕως ἀπώλεσεν αὐτοὺς· οὐ κατέλιπον ἐξ αὐτῶν οὐδὲν ἐμπνέον.
\vs{15}Ὃν τρόπον συνέταξε Κύριος τῷ Μωυσῇ τῷ παιδὶ αὐτοῦ, καὶ Μωυσῆς ὡσαύτως ἐνετείλατο τῷ Ἰησοῖ· καὶ οὕτως ἐποίησεν Ἰησοῦς, οὐ παρέβη οὐδὲν ἀπὸ πάντων ὧν συνέταξεν αὐτῷ Μωυσῆς.

\vs{16}Καὶ ἔλαβεν Ἰησοῦς πᾶσαν τὴν γῆν τὴν ὀρεινὴν, καὶ πᾶσαν τὴν γῆν Ναγὲβ, καὶ πᾶσαν τὴν γῆν Γοσὸμ, καὶ τὴν πεδινὴν, καὶ τὴν πρὸς δυσμαῖς, καὶ τὸ ὄρος Ἰσραὴλ, καὶ τὰ ταπεινὰ
\vs{17}τὰ πρὸς τῷ ὄρει ἀπὸ ὄρους Χελχὰ, καὶ ὃ προσαναβαίνει εἰς Σηεὶρ, καὶ ἕως Βαλαγὰδ, καὶ τὰ πεδία τοῦ Λιβάνου ὑπὸ τὸ ὄρος τὸ Ἀερμών· καὶ πάντας τοὺς βασιλεῖς αὐτῶν ἔλαβε, καὶ ἀνεῖλε, καὶ ἀπέκτεινε.
\vs{18}Καὶ ἡμέρας πλείους ἐποίησεν Ἰησοῦς πρὸς τοὺς βασιλεῖς τούτους τὸν πόλεμον.

\vs{19}Καὶ οὐκ ἦν πόλις, ἣν οὐκ ἔλαβεν Ἰσραήλ· πάντα ἐλάβοσαν ἐν πολέμῳ.
\vs{20}Ὅτι διὰ Κυρίου ἐγένετο κατισχύσαι αὐτῶν τὴν καρδίαν συναντᾷν εἰς πόλεμον πρὸς Ἰσραὴλ, ἵνα ἐξολοθρευθῶσιν, ὅπως μὴ δοθῇ αὐτοῖς ἔλεος, ἀλλʼ ἵνα ἐξολοθρευθῶσιν, ὃν τρόπον εἶπε Κύριος πρὸς Μωυσῆν.

\vs{21}Καὶ ἦλθεν Ἰησοῦς ἐν τῷ καιρῷ ἐκείνῳ, καὶ ἐξωλόθρευσε τοὺς Ἐνακὶμ ἐκ τῆς ὀρεινῆς, ἐκ Χεβρὼν, καὶ ἐκ Δαβὶρ, καὶ ἐξ Ἀναβὼθ, καὶ ἐκ παντὸς Ἰούδα σὺν ταῖς πόλεσιν αὐτῶν· καὶ ἐξωλόθρευσεν αὐτοὺς Ἰησοῦς.
\vs{22}Οὐ κατελείφθη τῶν Ἐνακὶμ ἀπὸ τῶν υἱῶν Ἰσραὴλ, ἀλλὰ πλὴν ἐν Γάζῃ, καὶ ἐν Γὲθ, καὶ ἐν Ἀσελδὼ κατελείφθη.

\vs{23}Καὶ ἔλαβεν Ἰησοῦς πᾶσαν τὴν γῆν, καθότι ἐνετείλατο Κύριος τῷ Μωυσῇ· καὶ ἔδωκεν αὐτοὺς Ἰησοῦς ἐν κληρονομίᾳ Ἰσραὴλ ἐν μερισμῷ κατὰ φυλὰς αὐτῶν· καὶ ἡ γῆ κατέπαυσε πολεμουμένη.

\ch{12}
Καὶ οὗτοι οἱ βασιλεῖς τῆς γῆς, οὓς ἀνεῖλον οἱ υἱοὶ Ἰσραὴλ, καὶ κατεκληρονόμησαν τὴν γῆν αὐτῶν πέραν τοῦ Ἰορδάνου ἀφʼ ἡλίου ἀνατολῶν ἀπὸ φάραγγος Ἀρνῶν ἕως τοῦ ὄρους Ἀερμὼν, καὶ πᾶσαν τὴν γῆν Ἄραβα ἀπʼ ἀνατολῶν.
\vs{2}Σηὼν τὸν βασιλέα τῶν Ἀμοῤῥαίων, ὃς κατῴκει ἐν Ἐσεβὼν, κυριεύων ἀπὸ Ἀρνῶν, ἥ ἐστιν ἐν τῇ φάραγγι κατὰ μέρος τῆς φάραγγος, καὶ τὸ ἥμισυ τῆς Γαλαὰδ ἕως Ἰαβὸκ, ὅρια υἱῶν Ἀμμών.
\vs{3}Καὶ Ἄραβ ἕως τῆς θαλάσσης Χενερὲθ κατʼ ἀνατολὰς, καὶ ἕως τῆς θαλάσσης Ἄραβα, θάλασσαν τῶν ἁλῶν ἀπὸ ἀνατολῶν ὁδὸν τὴν κατὰ Ἀσειμὼθ, ἀπὸ Θαιμὰν τὴν ὑπὸ Ἀσηδὼθ Φασγά.
\vs{4}Καὶ Ὢγ βασιλεὺς Βασὰν ὑπελείφθη ἐκ τῶν γιγάντων, ὁ κατοικῶν ἐν Ἀσταρὼθ καὶ ἐν Ἐδραῒν,
\vs{5}ἄρχων ἀπὸ ὄρους Ἀερμὼν καὶ ἀπὸ Σεκχαὶ, καὶ πᾶσαν τὴν Βασὰν ἕως ὁρίων Γεργεσὶ, καὶ τὴν Μαχὶ, καὶ τὸ ἥμισυ Γαλαδ ὁρίων Σηὼν βασιλέως Ἐσεβών.
\vs{6}Μωυσῆς ὁ παῖς Κυρίου καὶ οἱ υἱοὶ Ἰσραὴλ ἐπάταξαν αὐτούς· καὶ ἔδωκεν αὐτὴν Μωυσῆς ἐν κληρονομίᾳ Ῥουβὴν, καὶ Γὰδ, καὶ τῷ ἡμίσει φυλῆς Μανασσῆ.

\vs{7}Καὶ οὗτοι οἱ βασιλεῖς τῶν Ἀμοῤῥαίων, οὓς ἀνεῖλεν Ἰησοῦς καὶ υἱοὶ Ἰσραὴλ ἐν τῷ πέραν τοῦ Ἰορδάνου παρὰ θάλασσαν Βαλαγὰδ ἐν τῷ πεδίῳ τοῦ Λιβάνου, καὶ ἕως ὄρους τοῦ Χελχὰ ἀναβαινόντων εἰς Σηείρ· καὶ ἔδωκεν αὐτὴν Ἰησοῦς ταῖς φυλαῖς Ἰσραὴλ κληρονομεῖν κατὰ κλῆρον αὐτῶν,
\vs{8}ἐν τῷ ὄρει, καὶ ἐν τῷ πεδίῳ, καὶ ἐν Ἄραβ, καὶ ἐν Ἀσηδὼθ, καὶ ἐν τῇ ἐρήμῳ, καὶ Ναγὲβ· τὸν Χετταῖον, καὶ τὸν Ἀμοῤῥαῖον, καὶ τὸν Χαναναῖον, καὶ τὸν Φερεζαῖον, καὶ τὸν Εὐαῖον, καὶ τὸν Ἰεβουσαῖον.

\vs{9}Τὸν βασιλέα Ἱεριχὼ, καὶ τὸν βασιλέα τῆς Γαὶ, ἥ ἐστι πλησίον Βαιθὴλ,
\vs{10}βασιλέα Ἱερουσαλὴμ, βασιλέα Χεβρὼν,
\vs{11}βασιλέα Ἱεριμοὺθ, βασιλέα Λαχὶς,
\vs{12}βασιλέα Αἰλὰμ, βασιλέα Γαζὲρ,
\vs{13}βασιλέα Δαβὶρ, βασιλέα Γαδὲρ,
\vs{14}βασιλέα Ἑρμὰθ, βασιλέα Ἀδὲρ,
\vs{15}βασιλέα Λεβνὰ, βασιλέα Ὀδολλὰμ,
\vs{16}βασιλέα Ἠλὰθ,
\vs{17}βασιλέα Ταφοὺτ, βασιλέα Ὀφὲρ,
\vs{18}βασιλέα Ὀφὲκ τῆς Ἀρὼκ,
\vs{19}βασιλέα Ἀσὼμ,
\vs{20}βασιλέα Συμοὼν, βασιλέα Μαμβρὼθ, βασιλέα Ἀζὶφ,
\vs{21}βασιλέα Κάδης, βασιλέα Ζαχὰκ,
\vs{22}βασιλέα Μαρεδὼθ, βασιλέα Ἰεκὸμ τοῦ Χερμὲλ,
\vs{23}βασιλέα Ὀδολλὰμ τοῦ Φεννεαδὼρ, βασιλέα Γεῒ τῆς Γαλιλαίας,
\vs{24}βασιλέα Θερσά· πάντες οὗτοι βασιλεῖς εἰκοσιεννέα.

\ch{13}
Καὶ Ἰησοῦς πρεσβύτερος προβεβηκὼς τῶν ἡμερῶν· καὶ εἶπε Κύριος πρὸς Ἰησοῦν, σὺ προβέβηκας τῶν ἡμερῶν, καὶ ἡ γῆ ὑπολέλειπται πολλὴ εἰς κληρονομίαν.
\vs{2}Καὶ αὕτη ἡ γῆ καταλελειμμένη· ὅρια Φυλιστιεὶμ, ὁ Γεσιρὶ, καὶ ὁ Χαναναῖος,
\vs{3}ἀπὸ τῆς ἀοικήτου τῆς κατὰ πρόσωπον Αἰγύπτου ἕως τῶν ὁρίων Ἀκκαρὼν ἐξ εὐωνύμων τῶν Χαναναίων προσλογίζεται ταῖς πέντε σατραπείαις τῶν Φυλιστιείμ, τῷ Γαζαίῳ, καὶ τῷ Ἀζωτίῳ, καὶ τῷ Ἀσκαλωνίτῃ, καὶ τῷ Γετθαίῳ, καὶ τῷ Ἀκκαρωνίτῃ, καὶ τῷ Εὐαίῳ,
\vs{4}ἐκ Θαιμὰν καὶ πάσῃ γῇ Χαναὰν ἐναντίον Γάζης, καὶ οἱ Σιδώνιοι ἕως Ἀφὲκ, ἕως τῶν ὁρίων τῶν Ἀμοῤῥαίων,
\vs{5}καὶ πᾶσαν τὴν γῆν Γαλιὰθ Φυλιστιεὶμ, καὶ πάντα τὸν Λίβανον ἀπὸ ἀνατολῶν ἡλίου ἀπὸ Γαλγὰλ ὑπὸ τὸ ὄρος τὸ Ἀερμὼν ἕως τῆς εἰσόδου Ἐμὰθ,
\vs{6}πᾶς ὁ κατοικῶν τὴν ὀρεινὴν ἀπὸ τοῦ Λιβάνου ἕως τῆς Μασερὲθ Μεμφωμαὶμ. Πάντας τοὺς Σιδωνίους, ἐγὼ αὐτοὺς ἐξολοθρεύσω ἀπὸ προσώπου Ἰσραήλ· ἀλλὰ διάδος αὐτὴν ἐν κλήρῳ τῷ Ἰσραὴλ, ὃν τρόπον σοὶ ἐνετειλάμην.

\vs{7}Καὶ νῦν μέρισον τὴν γῆν ταύτην ἐν κληρονομίᾳ ταῖς ἐννέα φυλαῖς, καὶ τῷ ἡμίσει φυλῆς Μανασσή. Ἀπὸ τοῦ Ἰορδάνου ἕως τῆς θαλάσσης τῆς μεγάλης κατὰ δυσμὰς ἡλίου δώσεις αὐτήν· ἡ θάλασσα ἡ μεγάλη ὁριεῖ.
\vs{8}ταῖς δυσὶ φυλαῖς, καὶ τῷ ἡμίσει φυλῆς Μανασσή, τῷ Ῥουβὴν, καὶ τῷ Γὰδ ἔδωκε Μωυσῆς ἐν τῷ πέραν τοῦ Ἰορδάνου· κατʼ ἀνατολὰς ἡλίου δέδωκεν αὐτῷ Μωυσῆς ὁ παῖς Κυρίου,
\vs{9}ἀπὸ Ἀροὴρ, ἥ ἐστιν ἐπὶ τοῦ χείλους χειμάῤῥου Ἀρνών, καὶ τὴν πόλιν τὴν ἐν μέσῳ τῆς φάραγγος, καὶ πᾶσαν τὴν Μισὼρ ἀπὸ Μαιδαβάν·
\vs{10}Πάσας τὰς πόλεις Σηὼν βασιλέως Ἀμοῤῥαίων, ὃς ἐβασίλευσεν ἐν Ἐσεβὼν ἕως τῶν ὁρίων υἱῶν Ἀμμῶν·
\vs{11}Καὶ τὴν Γαλααδίτιδα, καὶ τὰ ὅρια Γεσιρὶ, καὶ τοὺς Μαχατὶ, πᾶν ὄρος Ἀερμὼν, καὶ πᾶσαν τὴν Βασανίτιν ἕως Ἀχά·
\vs{12}Πᾶσαν τὴν βασιλείαν Ὢγ ἐν τῇ Βασανίτιδι, ὃς ἐβασίλευσεν ἐν Ἀσταρὼθ καὶ ἐν Ἐδραΐν· οὗτος κατελείφθη ἀπὸ τῶν γιγάντων, καὶ ἐπάταξεν αὐτὸν Μωυσῆς, καὶ ἐξωλόθρευσε.
\vs{13}Καὶ οὐκ ἐξωλόθρευσαν οἱ υἱοὶ Ἰσραὴλ τὸν Γεσιρὶ, καὶ τὸν Μαχατὶ, καὶ τὸν Χαναναῖον· καὶ κατῷκει βασιλεὺς Γεσιρὶ καὶ ὁ Μαχατὶ ἐν τοῖς υἱοῖς Ἰσραὴλ ἕως τῆς σήμερον ἡμέρας.

\vs{14}Πλὴν τῆς φυλῆς Λευὶ οὐκ ἐδόθη κληρονομία· Κύριος ὁ Θεὸς Ἰσραὴλ, οὗτος κληρονομία αὐτῶν, καθὰ εἶπεν αὐτοῖς Κύριος· καὶ οὗτος ὁ καταμερισμὸς, ὃν κατεμέρισε Μωυσῆς τοῖς υἱοῖς Ἰσραὴλ ἐν Ἀραβὼθ Μωὰβ ἐν τῷ πέραν τοῦ Ἰορδάνου κατὰ Ἱεριχώ.

\vs{15}Καὶ ἔδωκε Μωυσῆς τῇ φυλῇ Ῥουβὴν κατὰ δήμους αὐτῶν.
\vs{16}Καὶ ἐγενήθη αὐτῶν τὰ ὅρια ἀπὸ Ἀροὴρ, ἥ ἐστι κατὰ πρόσωπον φάραγγος Ἀρνῶν, καὶ ἡ πόλις ἡ ἐν τῇ φάραγγι Ἀρνῶν· καὶ πᾶσαν τὴν Μισὼρ,
\vs{17}ἕως Ἐσεβὼν, καὶ πάσας τὰς πόλεις τὰς οὔσας ἐν τῇ Μισὼρ, καὶ Δαιβὼν, καὶ Βαιμὼν Βαὰλ, καὶ οἴκου Μεελβὼθ,
\vs{18}καὶ Βασὰν, καὶ Βακεδμὼθ, καὶ Μαιφαὰδ, καὶ, Καριθαὶμ,
\vs{19}καὶ Σεβαμὰ, καὶ Σεραδὰ, καὶ Σιὼν ἐν τῷ ὄρει Ἐνὰβ,
\vs{20}καὶ Βαιθφογὼρ, καὶ Ἀσηδὼθ Φασγὰ, καὶ Βαιτθασεινὼθ,
\vs{21}καὶ πάσας τὰς πόλεις τοῦ Μισὼρ, καὶ πᾶσαν τὴν βασιλείαν τοῦ Σηὼν βασιλέως τῶν Ἀμοῤῥαίων, ὃν ἐπάταξε Μωυσῆς αὐτὸν καὶ τοὺς ἡγουμένους Μαδιὰμ, καὶ τὸν Εὐὶ, καὶ τὸν Ῥοβὸκ, καὶ τὸν Σοὺρ, καὶ τὸν Οὒρ, καὶ τὸν Ῥοβὲ ἄρχοντα ἔναρα Σιὼν, καὶ τοὺς κατοικοῦντας Σιών.

\vs{22}Καὶ τὸν Βαλαὰμ τὸν τοῦ Βαιὼρ τὸν μάντιν ἀπέκτειναν ἐν τῇ ῥοπῇ.

\vs{23}Ἐγένετο δὲ τὰ ὅρια Ῥουβὴν, Ἰορδάνης ὅριον· αὑτη ἡ κληρονομία υἱῶν Ῥουβὴν κατὰ δήμους αὐτῶν, αἱ πόλεις αὐτῶν καὶ αἱ ἐπαύλεις αὐτῶν.

\vs{24}Ἔδωκε δὲ Μωυσῆς τοῖς υἱοῖς Γὰδ κατὰ δήμους αὐτῶν.
\vs{25}Καὶ ἐγένετο τὰ ὅρια αὐτῶν Ἰαζήρ· πᾶσαι πόλεις Γαλαὰδ, καὶ τὸ ἥμισυ γῆς υἱῶν Ἀμμῶν ἕως Ἄραβα, ἥ ἐστι κατὰ προσωπον Ἀράδ.
\vs{26}Καὶ ἀπὸ Ἐσεβὼν ἕως Ἀραβὼθ κατὰ τὴν Μασσηφὰ, καὶ Βοτανὶμ, καὶ Μαὰν ἕως τῶν ὁρίων Δαιβὼν,
\vs{27}καὶ Ἐναδὼμ καὶ Ὀθαργαῒ καὶ Βαινθαναβρὰ καὶ Σοκχωθὰ καὶ Σαφὰν καὶ τὴν λοιπὴν βασιλείαν Σηὼν βασιλέως Ἐσεβών· καὶ ὁ Ἰορδάνης ὁριεῖ ἕως μέρους τῆς θαλάσσης Χενερὲθ πέραν τοῦ Ἰορδάνου ἀπʼ ἀνατολῶν.
\vs{28}Αὕτη ἡ κληρονομία υἱῶν Γὰδ κατὰ δήμους αὐτῶν καὶ κατὰ πόλεις αὐτῶν· κατὰ δήμους αὐτῶν αὐχένα ἐπιστρέψουσιν ἐναντίον τῶν ἐχθρῶν αὐτῶν, ὅτι ἐγενήθη κατὰ δήμους αὐτῶν αἱ πόλεις αὐτῶν, καὶ αἱ ἐπαύλεις αὐτῶν.

\vs{29}Καὶ ἔδωκε Μωυσῆς τῷ ἡμίσει φυλῆς Μανασσῆ κατὰ δήμους αὐτῶν.
\vs{30}Καὶ ἐγένετο τὰ ὅρια αὐτῶν ἀπὸ Μαὰν, καὶ πᾶσα βασιλεία Βασὰν, καὶ πᾶσα βασιλεία Ὢγ βασιλέως τῆς Βασὰν, καὶ πάσας τὰς κώμας Ἰαῒρ, αἵ εἰσιν ἐν τῇ Βασανίτιδι, ἑξήκοντα πόλεις.
\vs{31}Καὶ τὸ ἥμισυ τῆς Γαλαάδ· καὶ ἐν Ἀσταρὼθ, καὶ ἐν Ἐδραῒν πόλεις βασιλείας Ὢγ ἐν τῇ Βασανίτιδι, τοῖς υἱοῖς Μαχὶρ υἱοῖς Μανασσῆ, καὶ τοῖς ἡμίσεσιν υἱοις Μαχὶρ υἱοῖς Μανασσῆ, κατὰ δήμους αὐτῶν.
\vs{32}Οὗτοι οὓς κατεκληρονόμησε Μωυσῆς πέραν τοῦ Ἰορδάνου ἐν Ἀραβὼθ Μωὰβ ἐν τῷ πέραν τοῦ Ἰορδάνου τοῦ κατὰ Ἱεριχὼ ἀπʼ ἀνατολῶν.

\ch{14}
Καὶ οὗτοι οἱ κατακληρονομήσαντες υἱῶν Ἰσραὴλ ἐν τῇ γῇ Χαναὰν, οἷς κατεκληρονόμησαν αὐτοῖς Ἐλεάζαρ ὁ ἱερεὺς, καὶ Ἰησοῦς ὁ τοῦ Ναυῆ, καὶ οἱ ἄρχοντες πατριῶν φυλῶν τῶν υἱῶν Ἰσραήλ.
\vs{2}Κατὰ κλήρους ἐκληρονόμησαν, ὃν τρόπον ἐνετείλατο Κύριος ἐν χειρὶ Ἰησοῦ ταῖς ἐννέα φυλαῖς, καὶ τῷ ἡμίσει φυλῆς
\vs{3}ἀπὸ τοῦ πέραν τοῦ Ἰορδάνου. Καὶ τοῖς Λευίταις οὐκ ἔδωκε κλῆρον ἐν αὐτοῖς,
\vs{4}ὅτι ἦσαν οἱ υἱοὶ Ἰωσὴφ δύο φυλαὶ Μανασσὴ καὶ Ἐφραΐμ· καὶ οὐκ ἐδόθη μερὶς ἐν τῇ γῇ τοῖς Λευίταις, ἀλλʼ ἤ πόλεις κατοικεῖν, καὶ τὰ ἀφωρισμένα αὐτῶν τοῖς κτήνεσι, καὶ τὰ κτήνη αὐτῶν.
\vs{5}Ὃν τρόπον ἐνετείλατο Κύριος τῷ Μωυσῇ, οὕτως ἐποίησαν οἱ υἱοὶ Ἰσραὴλ, καὶ ἐμέρισαν τὴν γῆν.

\vs{6}Καὶ προσήλθοσαν οἱ υἱοῖ Ἰούδα πρὸς Ἰησοῦν ἐν Γαλγάλ· καὶ εἶπε πρὸς αὐτὸν Χάλεβ ὁ τοῦ Ἰεφονὴ ὁ Κενεζαῖος, σὺ ἐπίστῃ τὸ ῥῆμα, ὃ ἐλάλησε Κύριος πρὸς Μωυσῆν ἄνθρωπον τοῦ Θεοῦ περὶ ἐμοῦ καὶ σοῦ ἐν Κάδης Βαρνῆ.
\vs{7}Τεσσαράκοντα γὰρ ἐτῶν ἤμην ὅτε ἀπέστειλέ με Μωυσῆς ὁ παῖς τοῦ Θεοῦ ἐκ Κάδης Βαρνῆ κατασκοπεῦσαι τὴν γῆν· καὶ ἀπεκρίθην αὐτῷ λόγον κατὰ τὸν νοῦν αὐτοῦ.
\vs{8}Οἱ ἀδελφοί μου οἱ ἀναβάντες μετʼ ἐμοῦ μετέστησαν τὴν καρδίαν τοῦ λαοῦ, ἐγὼ δὲ προσετέθην ἐπακολουθῆσαι Κυρίῳ τῷ Θεῷ μου.
\vs{9}Καὶ ὤμοσε Μωυσῆς ἐν ἐκείνῃ τῇ ἡμέρᾳ, λέγων, ἡ γῆ ἐφʼ ἣν ἐπέβης, σοὶ ἔσται ἐν κλήρῳ καὶ τοῖς τέκνοις σου εἰς τὸν αἰῶνα, ὅτι προσετέθης ἐπακολουθῆσαι ὀπίσω Κυρίου τοῦ Θεοῦ ἡμῶν.
\vs{10}Καὶ νῦν διέθρεψέ με Κύριος ὃν τρόπον εἶπε· τοῦτο τεσσαρακοστὸν καὶ πέμπτον ἔτος, ἀφʼ οὗ ἐλάλησε Κύριος τὸ ῥῆμα τοῦτο πρὸς Μωυσῆν· καὶ ἐπορεύθη Ἰσραὴλ ἐν τῇ ἐρήμῳ· καὶ νῦν ἰδοὺ ἐγὼ σήμερον ὀγδοήκοντα καὶ πέντε ἐτῶν·
\vs{11}Ἔτι εἰμὶ σήμερον ἰσχύων, ὡσεὶ ὅτε ἀπέστειλέ με Μωυσῆς, ὡσαύτως ἰσχύω νῦν ἐζελθεῖν καὶ εἰσελθεῖν εἰς τὸν πόλεμον.
\vs{12}Καὶ νῦν αἰτοῦμαί σε τὸ ὄρος τοῦτο, καθὰ εἰπε Κύριος τῇ ἡμέρᾳ ἐκείνῃ, ὅτι σὺ ἀκήκοας τὸ ῥῆμα τοῦτο ἐν τῇ ἡμέρᾳ ἐκείνῃ· νῦν δὲ οἱ Ἐνακὶμ ἐκεῖ εἰσι, πόλεις ὀχυραὶ καὶ μεγάλαι· ἐὰν οὖν Κύριος μετʼ ἐμοῦ ᾖ, ἐξολοθρεύσω αὐτοὺς, ὃν τρόπον εἶπέ μοι Κύριος.

\vs{13}Καὶ εὐλόγησεν αὐτὸν Ἰησοῦς, καὶ ἔδωκε Χεβρὼν τῷ Χάλεβ υἱῷ Ἰεφονῆ υἱῷ Κενὲζ ἐν κλήρῳ.
\vs{14}Διὰ τοῦτο ἐγενήθη ἡ Χεβρὼν τῷ Χάλεβ τῷ τοῦ Ἰεφονῆ τοῦ Κενεζαίου ἐν κλήρῳ ἕως τῆς ἡμέρας ταύτης, διὰ τὸ αὐτὸν ἐπακολουθῆσαι τῷ προστάγματι Κυρίου Θεοῦ Ἰσραήλ.
\vs{15}Τὸ δὲ ὄνομα τῆς Χεβρὼν ἦν τὸ πρότερον πόλις Ἀργὸβ, μητρόπολις τῶν Ἐνακὶμ αὕτη· καὶ ἡ γῆ ἐκόπασε τοῦ πολέμου.

\ch{15}
Καὶ ἐγένετο τὰ ὅρια φυλῆς Ἰούδα κατὰ δήμους αὐτῶν ἀπὸ τῶν ὁρίων τῆς Ἰδουμαίας ἀπὸ τῆς ἐρήμου Σὶν ἕως Κάδης πρὸς Λίβα.

\vs{2}Καὶ ἐγενήθη αὐτῶν τὰ ὅρια ἀπὸ Λιβὸς ἕως μέρους θαλάσσης τῆς ἁλυκῆς ἀπὸ τῆς λοφιᾶς τῆς φερούσης ἐπὶ Λίβα.
\vs{3}Καὶ διαπορεύεται ἀπέναντι τῆς προσαναβάσεως Ἀκραβίν· καὶ ἐκπεριπορεύεται Σενὰ, καὶ ἀναβαίνει ἀπὸ Λιβὸς ἐπὶ Κάδης Βαρνῆ· καὶ ἐκπορεύεται Ἀσωρὼν, καὶ προσαναβαίνει εἰς Σάραδα· καὶ ἐκπορεύεται τὴν κατὰ δυσμὰς Κάδης,
\vs{4}καὶ ἐκπορεύεται ἐπὶ Σελμωνὰν, καὶ διεκβάλλει ἕως φάραγγος Αἰγύπτου· καὶ ἔσται αὐτοῦ ἡ διέξοδος τῶν ὁρίων ἐπὶ τὴν θάλασσαν· τοῦτό ἐστιν αὐτῶν ὅρια ἀπὸ Λιβός.

\vs{5}Καὶ τὰ ὅρια ἀπὸ ἀνατολῶν πᾶσα ἡ θάλασσα ἡ ἁλυκὴ ἕως τοῦ Ἰορδάνου. Καὶ τὰ ὅρια αὐτῶν ἀπὸ Βοῤῥᾶ, καὶ ἀπὸ τῆς λοφιᾶς τῆς θαλάσσης καὶ ἀπὸ τοῦ μέρους τοῦ Ἰορδάνου.
\vs{6}Ἐπιβαίνει τὰ ὅρια ἐπὶ Βαιθαγλαάμ· καὶ παραπορεύεται ἀπὸ Βοῤῥᾶ ἐπὶ Βαιθάραβα, καὶ προσαναβαίνει τὰ ὅρια ἐπὶ λίθον Βαιὼν υἱοῦ Ῥουβήν·
\vs{7}Καὶ προσαναβαίνει τὰ ὅρια ἐπὶ τὸ τέταρτον τῆς φάραγγος Ἀχὼρ, καὶ καταβαίνει ἐπὶ Γαλγὰλ, ἥ ἐστιν ἀπέναντι τῆς προσβάσεως Ἀδαμμὶν, ἥ ἐστι κατὰ Λίβα τῇ φάραγγι, καὶ διεκβάλλεῖ ἐπὶ τὸ ὕδωρ τῆς πηγῆς τοῦ ἡλίου· καὶ ἔσται αὐτοῦ ἡ διέξοδος πηγὴ Ῥωγήλ·
\vs{8}Καὶ ἀναβαίνει τὰ ὅρια εἰς φάραγγα Ἐννὸμ, ἐπὶ νώτου τοῦ Ἰεβοῦις ἀπὸ Λιβός· αὕτη ἐστιν Ἱερουσαλήμ· καὶ διεκβάλλει τὰ ὅρια ἐπὶ κορυφὴν ὄρους, ἥ ἐστι κατὰ πρόσωπον φάραγγος Ἐννὸμ πρὸς θαλάσσης, ἥ ἐστιν ἐκ μέρους γῆς Ῥαφαῒν ἐπὶ βοῤῥᾶ·
\vs{9}Καὶ διεκβάλλει τὸ ὅριον ἀπὸ κορυφῆς τοῦ ὄρους ἐπὶ πηγὴν ὕδατος Ναφθὼ, καὶ διεκβάλλει εἰς τὸ ὄρος Ἐφρών· καὶ ἄξει τὸ ὅριον εἰς Βαάλ· αὕτη ἐστὶ πόλις Ἰαρίμ.
\vs{10}Καὶ περιελεύσεται ὅριον ἀπὸ Βαὰλ ἐπὶ θάλασσαν, καὶ παρελεύσεται εἰς ὄρος Ἀσσὰρ ἐπὶ νώτου πόλιν Ἰαρὶν ἀπὸ Βοῤῥᾶ· αὕτη ἐστὶ Χασλών· καὶ καταβήσεται ἐπὶ πόλιν ἡλίου, καὶ παρελεύσεται ἐπὶ Λίβα·
\vs{11}Καὶ διεκβάλλει τὸ ὅριον κατὰ νώτου Ἀκκαρὼν ἐπὶ βοῤῥᾶν, καὶ διεκβαλεῖ τὰ ὅρια εἰς Σοκχὼθ, καὶ παρελεύσεται ὅρια ἐπὶ Λίβα, καὶ διεκβαλεῖ ἐπὶ Λεβνὰ, καὶ ἔσται ἡ διέξοδος τῶν ὁρίων ἐπὶ θάλασσαν·
\vs{12}Καὶ τὰ ὅρια αὐτῶν ἀπὸ θαλάσσης, ἡ θάλασσα ἡ μεγάλη ὁριεῖ. Ταῦτα τὰ ὅρια υἱῶν Ἰούδα κύκλῳ κατὰ δήμους αὐτῶν.

\vs{13}Καὶ τῷ Χάλεβ υἱῷ Ἰεφονῆ ἔδωκε μερίδα ἐν μέσῳ υἱῶν Ἰούδα διὰ προστάγματος τοῦ Θεοῦ· καὶ ἔδωκεν αὐτῷ Ἰησοῦς τὴν πόλιν Ἀρβὸκ μητρόπολιν Ἐνάκ· αὕτη ἐστὶ Χεβρών.
\vs{14}Καὶ ἐξωλόθρευσεν ἐκεῖθεν Χάλεβ υἱὸς Ἰεφονὴ τοὺς τρεῖς υἱοὺς Ἐνὰκ, τὸν Σουσὶ καὶ Θολαμὶ καὶ τὸν Ἀχιμᾶ.
\vs{15}Καὶ ἀνέβη ἐκεῖθεν Χάλεβ ἐπὶ τοὺς κατοικοῦντας Δαβίρ· τὸ δὲ ὄνομα Δαβὶρ ἦν τὸ πρότερον πόλις Γραμμάτων.

\vs{16}Καὶ εἶπε Χάλεβ, ὃς ἂν λάβῃ καὶ ἐκκόψῃ τὴν πόλιν τῶν Γραμμάτων καὶ κυριεύσῃ αὐτῆς, δώσω αὐτῷ τὴν Ἀσχὰν θυγατέρα μου εἰς γυναῖκα.
\vs{17}Καὶ ἔλαβεν αὐτὴν Γοθονιὴλ υἱὸς Χενὲζ ἀδελφοῦ Χάλεβ· καὶ ἔδωκεν αὐτῷ τὴν Ἀσχὰν θυγατέρα αὐτοῦ γυναῖκα.
\vs{18}Καὶ ἐγένετο ἐν τῷ ἐκπορεύέσθαι αὐτὴν καὶ συνέβουλεύσατο αὐτῷ, λέγουσα, αἰτήσομαι τὸν πατέρα μου ἀγρόν· καὶ ἐβόησεν ἐκ τοῦ ὄνου· καὶ εἶπεν αὐτῇ Χάλεβ, τί ἐστί σοι;
\vs{19}Καὶ εἶπεν αὐτῷ, δός μοι εὐλογίαν, ὅτι εἰς γῆν Ναγὲβ δέδωκάς με· δός μοι τὴν Βοτθανίς· καὶ ἔδωκεν αὐτῇ τὴν Γοναιθλὰν τὴν ἄνω καὶ τὴν Γοναιθλὰν τὴν κάτω.

\vs{20}Αὕτη ἡ κληρονομία φυλῆς υἱῶν Ἰούδα.
\vs{21}Ἐγενήθησαν δὲ πόλεις αὐτῶν πόλεις πρὸς τῇ φυλῇ υἱῶν Ἰούδα ἐφʼ ὁρίων Ἐδὼμ ἐπὶ τῆς ἐρήμου, καὶ Βαισελεὴλ, καὶ Ἀρὰ, καὶ Ἀσὼρ,
\vs{22}καὶ Ἰκὰμ, καὶ Ῥεγμὰ, καὶ Ἀρουὴλ,
\vs{23}καὶ Κάδης, καὶ Ἀσοριωναὶν, καὶ Μαινὰμ,
\vs{24}καὶ Βαλμαινὰν, καὶ αἱ κῶμαι αὐτῶν,
\vs{25}καὶ αἱ πόλεις Ἀσερὼν, αὕτη Ἀσὼρ,
\vs{26}καὶ Σὴν, καὶ Σαλμαὰ, καὶ Μωλαδὰ,
\vs{27}καὶ Σερὶ, καὶ Βαιφαλὰθ,
\vs{28}καὶ Χολασεωλὰ, καὶ Βηρσαβεέ· καὶ αἱ κῶμαι αὐτῶν, καὶ αἱ ἐπαύλεις αὐτῶν,
\vs{29}Βαλὰ, καὶ Βακὼκ, καὶ Ἀσὸμ,
\vs{30}καὶ Ἐλβωϋδὰδ, καὶ Βαιθὴλ, καὶ Ἑρμὰ,
\vs{31}καὶ Σεκελὰκ, καὶ Μαχαρὶμ, καὶ Σεθεννὰκ,
\vs{32}καὶ Λαβὼς, καὶ Σαλὴ, καὶ Ἐρωμώθ· πόλεις εἰκοσιεννέα, καὶ αἱ κῶμαι αὐτῶν.

\vs{33}Ἐν τῇ πεδινῇ Ἀσταὼλ, καὶ Ῥάα, καὶ Ἄσσα.
\vs{34}Καὶ Ῥάμεν, καὶ Τανὼ, καὶ Ἰλουθὼθ, καὶ Μαιανὶ,
\vs{35}καὶ Ἰερμοὺθ, καὶ Ὀδολλὰμ, καὶ Μεμβρὰ, καὶ Σαωχὼ, καὶ Ἰαζηκὰ,
\vs{36}καὶ Σακαρὶμ, καὶ Γάδηρα, καὶ αἱ ἐπαύλεις αὐτῆς· πόλεις δεκατέσσαρες, καὶ αἱ κῶμαι αὐτῶν.
\vs{37}Σεννὰ, καὶ Ἀδασὰν, καὶ Μαγαδαλγὰδ,
\vs{38}καὶ Δαλὰδ, καὶ Μασφὰ, καὶ Ἰαχαρεὴλ,
\vs{39}καὶ Βασηδὼθ, καὶ Ἰδεαδαλέα, καὶ
\vs{40}Χαβρά, καὶ Μαχὲς, καὶ Μααχὼς,
\vs{41}καὶ Γεδδὼρ, καὶ Βαγαδιὴλ, καὶ Νωμὰν, καὶ Μαχηδάν· πόλεις ἑκκαίδεκα, καὶ αἱ κῶμαι αὐτῶν·
\vs{42}Λεβνὰ, καὶ Ἰθὰκ, καὶ Ἀνὼχ,
\vs{43}καὶ Ἰανὰ, καὶ Νασὶβ,
\vs{44}καὶ Κεϊλὰμ, καὶ Ἀκιεζὶ, καὶ Κεζὶβ, καὶ Βαθησὰρ, καὶ Αἰλώμ· πόλεις δέκα, καὶ αἱ κῶμαι αὐτῶν·
\vs{45}Ἀκκαρὼν, καὶ αἱ κῶμαι αὐτῆς, καὶ αἱ ἐπαύλεις αὐτῶν,
\vs{46}ἀπὸ Ἀκκαρὼν Γεμνά· καὶ πᾶσαι ὅσαι εἰσὶ πλησίον Ἀσηδώθ· καὶ αἱ κῶμαι αὐτῶν,
\vs{47}Ἀσιεδὼθ, καὶ αἱ κῶμαι αὐτῆς, καὶ αἱ ἐπαύλεις αὐτῆς· Γάζα, καὶ αἱ κῶμαι αὐτῆς, καὶ ἐπαύλεις αὐτῆς ἕως τοῦ χειμάῤῥου Αἰγύπτου, καὶ ἡ θάλασσα ἡ μεγάλη διορίζει.

\vs{48}Καὶ ἐν τῇ ὀρεινῇ Σαμὶρ, καὶ Ἰεθὲρ, καὶ Σωχὰ,
\vs{49}καὶ Ῥεννὰ, καὶ πόλις Γραμμάτων, αὕτη Δαβὶρ,
\vs{50}καὶ Ἀνὼν, καὶ Ἒς, καὶ Μὰν, καὶ Αἰσὰμ,
\vs{51}καὶ Γοσὸμ, καὶ Χαλοὺ, καὶ Χαννὰ, καὶ Γηλώμ· πόλεις ἕνδεκα, καὶ αἱ κῶμαι αὐτῶν·
\vs{52}Αἰρὲμ, καὶ Ῥεμνὰ, καὶ Σομὰ,
\vs{53}καὶ Ἰεμαῒν, καὶ Βαιθαχοὺ, καὶ Φακουὰ,
\vs{54}καὶ Εὐμὰ, καὶ πόλις Ἀρβὸκ, αὕτη ἐστὶ Χεβρὼν, καὶ Σωραίθ· πόλεις ἐννέα, καὶ αἱ ἐπαύλεις αὐτῶν·
\vs{55}Μαὼρ, καὶ Χερμὲλ, καὶ Ὀζὶβ, καὶ Ἰτὰν,
\vs{56}καὶ Ἰαριὴλ, καὶ Ἀρικὰμ, καὶ Ζακαναῒμ,
\vs{57}καὶ Γαβαὰ, καὶ Θαμναθά· πόλεις ἐννέα, καὶ αἱ κῶμαι αὐτῶν·
\vs{58}Αἰλουὰ, καὶ Βηθσοὺρ, καὶ Γεδδὼν,
\vs{59}καὶ Μαγαρὼθ, καὶ Βαιθανὰμ, καὶ Θεκούμ· πόλεις ἓξ, καὶ αἱ κῶμαι αὐτῶν·
\vs{59a}Θεκὼ, καὶ Ἐφραθὰ, αὕτη ἐστὶ Βαιθλεὲμ, καὶ Φαγὼρ, καὶ Αἰτὰν, καὶ Κουλὸν, καὶ Τατὰμ, καὶ Θωβὴς, καὶ Καρὲμ, καὶ Γαλὲμ καὶ Θεθὴρ, καὶ Μανοχώ· πόλεις ἕνδεκα, καὶ αἱ κῶμαι αὐτῶν·
\vs{60}Καριαθβαὰλ, αὕτη ἡ πόλις Ἰαρὶμ, καὶ Σωθηβᾶ· πόλεις δύο, καὶ αἱ ἐπαύλεις αὐτῶν·
\vs{61}καὶ Βαδδαργεὶς, καὶ Θαραβαὰμ, καὶ Αἰνὼν,
\vs{62}καὶ Αἰοχιοζὰ, καὶ Ναφλαζὼν, καὶ αἱ πόλεις Σαδῶν, καὶ Ἀρκάδης· πόλεις ἑπτὰ, καὶ αἱ κῶμαι αὐτῶν.

\vs{63}Καὶ ὁ Ἰεβουσαῖος κατῴκει ἐν Ἱερουσαλὴμ, καὶ οὐκ ἠδυνήθησαν οἱ υἱοὶ Ἰούδα ἀπολέσαι αὐτούς· καὶ κατῴκησαν οἱ Ἰεβουσαῖοι ἐν Ἱερουσαλὴμ ἕως τῆς ἡμέρας ταύτης.

\ch{16}
Καὶ ἐγένετο τὰ ὅρια υἱῶν Ἰωσὴφ ἀπὸ τοῦ Ἰορδάνου τοῦ κατὰ Ἱεριχὼ ἀπὸ ἀνατολῶν· καὶ ἀναβήσεται ἀπὸ Ἱεριχὼ εἰς τὴν ὀρεινὴν, τὴν ἔρημον, εἰς Βαιθὴλ Λουζά.
\vs{2}Καὶ ἐξελεύσεται εἰς Βαιθὴλ, καὶ παρελεύσεται ἐπὶ τὰ ὅρια τοῦ Ἀχαταρωθί.
\vs{3}Καὶ διελεύσεται ἐπὶ τὴν θάλασσαν ἐπὶ τὰ ὅρια Ἀπταλὶμ ἕως τῶν ὁρίων Βαιθωρὼν τὴν κάτω, καὶ ἔσται ἡ διέξοδος αὐτῶν ἐπὶ τὴν θάλασσαν.
\vs{4}Καὶ ἐκληρονόμησαν οἱ υἱοὶ Ἱωσὴφ, Ἐφραῒμ καὶ Μανασσή.

\vs{5}Καὶ ἐγενήθη ὅρια υἱῶν Ἐφραῒμ κατὰ δήμους αὐτῶν· καὶ ἐγενήθη τὰ ὅρια τῆς κληρονομίας αὐτῶν ἀπʼ ἀνατολῶν Ἀταρὼθ, καὶ Ἐρὼκ ἕως Βαιθωρὼν τὴν ἄνω, καὶ Γαζαρά.
\vs{6}Καὶ ἐλεύσεται τὰ ὅρια ἐπὶ τὴν θάλασσαν εἰς Ἰκασμὼν ἀπὸ Βοῤῥᾶ Θερμᾶ· περιελεύσεται ἐπʼ ἀνατολὰς εἰς Θηνασὰ, καὶ Σέλλης, καὶ παρελεύσεται ἀπʼ ἀνατολῶν εἰς Ἰανῶκὰ,
\vs{7}καὶ εἰς Μαχὼ, καὶ Ἀταρὼθ, καὶ αἱ κῶμαι αὐτῶν· καὶ ἐλεύσεται ἐπὶ Ἱεριχὼ, καὶ διεκβαλεῖ ἐπὶ τὸν Ἰορδάνην.
\vs{8}Καὶ ἀπὸ Τάφου πορεύσεται τὰ ὅρια ἐπὶ θάλασσαν ἐπὶ Χελκανα· καὶ ἔσται ἡ διέξοδος αὐτῶν ἐπὶ θάλασσαν· αὕτη ἡ κληρονομία φυλῆς Ἐφραῒμ κατὰ δήμους αὐτῶν.

\vs{9}Καὶ αἱ πόλεις αἱ ἀφορισθεῖσαι τοῖς υἱοῖς Ἐφραῒμ ἀναμέσον τῆς κληρονομίας νἱῶν Μανασσή, πᾶσαι αἱ πόλεις καὶ αἱ κῶμαι αὐτῶν.
\vs{10}Καὶ οὐκ ἀπώλεσεν Ἐφραῒμ τὸν Χαναναῖον τὸν κατοικοῦντα ἐν Γάζέρ· καὶ κατῴκει ὁ Χαναναῖος ἐν τῷ Ἐφραῒμ ἕως τῆς ἡμέρας ταυτης, ἕως ἀνέβη Φαραὼ βασιλεὺς Αἰγύπτου, καὶ ἔλαβεν αὐτὴν, καὶ ἐνέπρησεν αὐτὴν ἐν πυρί· καὶ τοὺς Χαναναίους, καὶ τοὺς Φερεζαίους, καὶ τοὺς κατοικοῦντας ἐν Γαζὲρ, ἐξεκέντησαν· καὶ ἔδωκεν αὐτὴν Φαραὼ ἐν φερνῇ τῇ θυγατρὶ αὐτοῦ.

\ch{17}
Καὶ ἐγένετο τὰ ὅρια φυλῆς υἱῶν Μανασσῆ, ὅτι οὗτος πρωτότοκος τῷ Ἰωσὴφ, τῷ Μαχὶρ πρωτοτόκῳ Μανασσῆ πατρὶ Γαλαὰδ, ἀνὴρ γὰρ πολεμιστὴς ἦν, ἐν τῇ Γαλααδίτιδι καὶ ἐν τῇ Βασανίτιδι.
\vs{2}Καὶ ἐγενήθη τοῖς υἱοῖς Μανασσῆ τοῖς λοιποῖς κατὰ δήμους αὐτῶν· τοῖς υἱοῖς Ἰεζὶ, καὶ τοῖς υἱοῖς Κελὲζ, καὶ τοῖς υἱοῖς Ἰεζιὴλ, καὶ τοῖς υἱοῖς Συχὲμ, καὶ τοῖς υἱοῖς Συμαρὶμ, καὶ τοῖς υἱοῖς Ὀφέρ· οὗτοι ἄρσενες κατὰ δήμους αὐτῶν.

\vs{3}Καὶ τῷ Σαλπαὰδ υἱῷ Ὀφὲρ οὐκ ἦσαν αὐτῷ υἱοὶ ἀλλʼ ἢ θυγατέρες· καὶ ταῦτα τὰ ὀνόματα τῶν θυγατέρων Σαλπαάδ· Μααλὰ, καὶ Νουὰ, καὶ Ἐγλὰ, καὶ Μελχὰ, καὶ Θερσά.
\vs{4}Καὶ ἔστησαν ἐναντίον Ἐλεάζαρ τοῦ ἱερέως, καὶ ἐναντίον Ἰησοῦ, καὶ ἐναντίον τῶν ἀρχόντων, λέγουσαι, ὁ Θεὸς ἐνετείλατο διὰ χειρὸς Μωυσῆ δοῦναι ἡμῖν κληρονομίαν ἐν μέσῳ τῶν ἀδελφῶν ἡμῶν· καὶ ἐδόθη αὐταῖς διὰ προστάγματος Κυρίου κλῆρος ἐν τοῖς ἀδελφοῖς τοῦ πατρὸς αὐτῶν.
\vs{5}Καὶ ἔπεσεν ὁ σχοινισμὸς αὐτῶν ἀπὸ Ἀνάσσα, καὶ πεδίον Λαβὲκ ἐκ τῆς γῆς Γαλαὰδ, ἥ ἐστι πέραν τοῦ Ἰορδάνου·
\vs{6}Ὅτι θυγατέρες υἱῶν Μανασσὴ ἐκληρονόμησαν κλῆρον ἐν μέσῳ τῶν ἀδελφῶν αὐτῶν· ἡ δὲ γῆ Γαλαὰδ ἐγενήθη τοῖς υἱοῖς Μανασσὴ τοῖς καταλελειμμένοις.

\vs{7}Καὶ ἐγενήθη ὅρια υἱῶν Μανασσῆ Δηλανὰθ, ἥ ἐστι κατὰ πρόσωπον υἱῶν Ἀνὰθ, καὶ πορεύεται ἐπὶ τὰ ὅρια ἐπὶ Ἰαμὶν καὶ Ἰασσὶβ ἐπὶ πηγὴν Θαφθώθ.
\vs{8}Τῷ Μανασσῇ ἔσται· καὶ Θαφὲθ ἐπὶ τῶν ὁρίων Μανασσῆ, τοῖς υἱοῖς Ἐφραΐμ.
\vs{9}Καὶ καταβήσεται τὰ ὅρια ἐπὶ φάραγγα Καρανὰ ἐπὶ Λίβα κατὰ φάραγγα Ἰαριὴλ· τερέμινθος τῷ Ἐφραῒμ ἀναμέσον πόλεως Μανασσῆ· καὶ ὅρια Μανασσῆ ἐπὶ τὸν βοῤῥᾶν εἰς τὸν χειμάῤῥουν· καὶ ἔσται αὐτοῦ ἡ διέξοδος θάλασσα
\vs{10}ἀπὸ Λιβὸς τῷ Ἐφραῒμ, καὶ ἐπὶ Βοῤῥᾶν Μανασσῇ· καὶ ἔσται ἡ θάλασσα ὅρια αὐτοῖς· καὶ ἐπὶ Ἀσὴβ συνάψουσιν ἐπὶ Βοῤῥᾶν, καὶ τῷ Ἰσσάχαρ ἀπὸ ἀνατολῶν.
\vs{11}Καὶ ἔσται Μανασσῇ ἐν Ἰσσάχαρ καὶ ἐν Ἀσὴρ Βαιθοὰν καὶ αἱ κῶμαι αὐτῶν, καὶ τοὺς κατοικοῦντας Δὼρ, καὶ τὰς κώμας αὐτῆς, καὶ τοὺς κατοικοῦντας Μαγεδδὼ, καὶ τὰς κώμας αὐτῆς, καὶ τὸ τρίτον τῆς Μαφετὰ, καὶ τὰς κώμας αὐτῆς.

\vs{12}Καὶ οὐκ ἠδυνάσθησαν οἱ υἱοὶ Μανασσῆ ἐξολοθρεῦσαι τὰς πόλεις ταύτας· καὶ ἤρχετο ὁ Χαναναῖος κατοικεῖν ἐν τῇ γῇ ταύτῃ.
\vs{13}Καὶ ἐγενήθη καὶ ἐπεὶ κατίσχυσαν οἱ υἱοὶ Ἰσραὴλ, καὶ ἐποίησαν τοὺς Χαναναίους ὑπηκόους, ἐξολοθρεῦσαι δὲ αὐτοὺς οὐκ ἐξωλόθρευσαν.

\vs{14}Ἀντεῖπαν δὲ οἱ νἱοὶ Ἰωσὴφ τῷ Ἰησοῦ, λέγοντες, διατί ἐκληρονόμησας ἡμᾶς κλῆρον ἕνα, καὶ σχοίνισμα ἕν; ἐγὼ δὲ λαὸς πολύς εἰμι, καὶ ὁ Θεὸς εὐλόγησέ με.
\vs{15}Καὶ εἶπεν αὐτοῖς Ἰησοῦς, εἰ λαὸς πολὺς εἶ, ἀνάβηθι εἰς τὸν δρυμὸν, καὶ ἐκκάθαρον σεαυτῷ εἰ στενοχωρεῖ σε τὸ ὄρος τὸ Ἐφραΐμ.
\vs{16}Καὶ εἶπαν, οὐκ ἀρέσκει ἡμῖν τὸ ὄρος τὸ Ἐφραΐμ· καὶ ἵππος ἐπίλεκτος, καὶ σίδηρος τῷ Χαναναίῳ τῷ κατοικοῦντι ἐν αὐτῷ ἐν Βαιθσὰν, καὶ ἐν ταῖς κώμαις αὐτῆς, ἐν τῇ κοιλάδι Ἰεζραέλ.
\vs{17}Καὶ εἶπεν Ἰησοῦς τοῖς υἱοῖς Ἰωσὴφ, εἰ λαὸς πολὺς εἶ καὶ ἰσχὺν μεγάλην ἔχεις, οὐκ ἔσται σοι κλῆρος εἷς·
\vs{18}Ὁ γὰρ δρυμὸς ἔσται σοι, ὅτι δρυμός ἐστι καὶ ἐκκαθαριεῖς αὐτὸν, καὶ ἔσται σοι· καὶ ὅταν ἐξολοθρεύ· σῃς τὸν Χαναναῖον, ὅτι ἵππος ἐπίλεκτος αὐτῷ ἐστι· σὺ γὰρ ὑπερισχύεις αὐτοῦ.

\ch{18}
Καὶ ἐξεκκλησιάσθη πᾶσα συναγωγὴ υἱῶν Ἰσραὴλ εἰς Σηλὼ, καὶ ἔπηξαν ἐκεῖ τὴν σκηνὴν τοῦ μαρτυρίου· καὶ ἡ γῆ ἐκρατήθη ὑπʼ αὐτῶν.

\vs{2}Καὶ κατελείφθησαν οἱ υἱοὶ Ἰσραὴλ, οἳ οὐκ ἐκληρονόμησαν, ἑπτὰ φυλαί.
\vs{3}Καὶ εἶπεν Ἰησοῦς τοῖς υἱοῖς Ἰσραὴλ, ἕως τίνος ἐκλυθήσεσθε κληρονομῆσαι τὴν γῆν, ἣν ἔδωκε Κύριος ὁ Θεὸς ἡμῶν;
\vs{4}Δότε ἐξ ὑμῶν ἄνδρας τρεῖς ἐκ φυλῆς, καὶ ἀναστάντες διελθέτωσαν τὴν γῆν, καὶ διαγραψάτωσαν αὐτὴν ἐναντίον μου, καθὰ δεήσει διελεῖν αὐτήν. Καὶ διέλθοσαν πρὸς αὐτόν·
\vs{5}καὶ διεῖλεν αὐτοῖς ἑπτὰ μερίδας· Ἰούδας στήσεται αὐτοῖς ὅριον ἀπὸ Λιβὸς, καὶ οἱ υἱοὶ Ἰωσὴφ στήσονται αὐτοῖς ἀπὸ Βοῤῥᾶ.
\vs{6}Ὑμεῖς δὲ μερίσατε τὴν γῆν ἑπτὰ μερίδας, καὶ ἐνέγκατε ὧδε πρὸς μὲ, καὶ ἐξοίσω ὑμῖν κλῆρον ἔναντι Κυρίου τοῦ Θεοῦ ἡμῶν.
\vs{7}Οὐ γάρ ἐστι μερὶς τοῖς υἱοῖς Λευὶ ἐν ὑμῖν· ἱερατεία γὰρ Κυρίου μερὶς αὐτοῦ· καὶ Γὰδ καὶ Ῥουβὴν καὶ τὸ ἥμισυ φυλῆς Μανασσῆ ἐλάβοσαν τὴν κληρονομίαν αὐτῶν πέραν τοῦ Ἰορδάνου ἐπʼ ἀνατολῆς, ἣν ἔδωκεν αὐτοῖς Μωυσῆς ὁ παῖς Κυρίου.

\vs{8}Καὶ ἀναστάντες οἱ ἄνδρες ἐπορεύθησαν· καὶ ἐνετείλατο Ἰησοῦς τοῖς ἀνδράσι τοῖς πορευομένοις χωροβατῆσαι τὴν γῆν, λέγων, πορεύεσθε καὶ χωροβατήσατε τὴν γῆν, καὶ παραγενήθητε πρὸς μὲ, καὶ ὧδε ἐξοίσω ὑμῖν κλῆρον ἔναντι Κυρίου ἐν Σηλώ.
\vs{9}Καὶ ἐπορεύθησαν, καὶ ἐχωροβάτησαν τὴν γῆν· καὶ εἴδοσαν αὐτὴν, καὶ ἔγραψαν αὐτὴν κατὰ πόλεις, ἑπτα μερίδας εἰς βιβλίον, καὶ ἤνεγκαν πρὸς Ἰησοῦν.
\vs{10}Καὶ ἐνέβαλεν αὐτοῖς Ἰησοῦς κλῆρον ἐν Σηλὼ ἔναντι Κυρίου.

\vs{11}Καὶ ἐξῆλθεν ὁ κλῆρος φυλῆς Βενιαμὶν πρῶτος κατὰ δήμους αὐτῶν· καὶ ἐξῆλθεν ὅρια τοῦ κλήρου αὐτῶν ἀναμέσον υἱῶν Ἰούδα καὶ ἀναμέσον τῶν υἱῶν Ἰωσήφ.

\vs{12}Καὶ ἐγενήθη αὐτῶν τὰ ὅρια ἀπὸ Βοῤῥᾶ· ἀπὸ τοῦ Ἱορδάνου προσαναβήσεται τὰ ὅρια κατὰ νὼτου Ἱεριχὼ ἀπὸ Βοῤῥᾶ, καὶ ἀναβήσεται ἐπὶ τὸ ὄρος ἐπὶ τὴν θάλασσαν, καὶ ἔσται αὐτοῦ ἡ διέξοδος ἡ Μαβδαρίτις Βαιθών.
\vs{13}Καὶ διελεύσεται ἐκεῖθεν τὰ ὅρια Λοῦζὰ ἐπὶ νώτου Λουζὰ ἀπὸ Λιβὸς αὐτῆς· αὕτη ἐστὶ Βαιθήλ· καὶ καταβήσεται τὰ ὅρια Μααταρὼβ Ὀρὲχ ἐπὶ τὴν ὀρεινὴν, ἥ ἐστι πρὸς Λίβα Βαιθωρὼν ἡ κάτω.

\vs{14}Καὶ διελεύσεται τὰ ὅρια καὶ παρελεύσεται ἐπὶ τὸ μέρος τὸ βλέπον παρὰ θάλασσαν ἀπὸ Λιβὸς ἀπὸ τοῦ ὄρους ἐπὶ πρόσωπον Βαιθωρὼν Λίβα· καὶ ἔσται αὐτοῦ ἡ διέξοδος εἰς Καριὰθ Βαάλ· αὕτη ἐστὶ Καριαθιαρὶν, πόλις υἱῶν Ἰούδα· τοῦτό ἐστι τὸ μέρος τὸ πρὸς θάλασσαν.

\vs{15}Καὶ μέρος τὸ πρὸς Λίβα ἀπὸ μέρους Καριὰθ Βαάλ· καὶ διελεύσεται ὅρια εἰς Γασὶν, ἐπὶ πηγὴν ὕδατος Ναφθώ.
\vs{16}Καὶ καταβήσεται τὰ ὅρια ἐπὶ μέρους, τοῦτό ἐστι κατὰ πρόσωπον νάπης Σοννὰμ, ὅ ἐστιν ἐκ μέρους Ἐμὲκ Ῥαφαῒν ἀπὸ Βοῤῥᾶ, καὶ καταβήσεται Γαίεννα ἐπὶ νῶτον Ἰεβουσαὶ ἀπὸ Λιβός· καταβήσεται ἐπὶ πηγὴν Ῥωγήλ·
\vs{17}Καὶ διελεύσεται ἐπὶ πηγὴν Βαιθσαμύς· καὶ παρελεύσεται ἐπὶ Γαλιλὼθ, ἥ ἐστιν ἀπέναντι πρὸς ἀνάβασις Αἰθαμίν· καὶ καταβήσεται ἐπὶ λίθον Βαιὼν υἱῶν Ῥουβήν·
\vs{18}καὶ διελεύσεται κατὰ νώτου Βαιθάραβα ἀπὸ Βοῤῥᾶ, καὶ καταβήσεται ἐπὶ τὰ ὅρια ἐπὶ νῶτον θάλασσαν ἀπὸ Βοῤῥᾶ.
\vs{19}Καὶ ἔσται ἡ διέξοδος τῶν ὁρίων ἐπὶ λοφιὰν τῆς θαλάσσης τῶν ἁλῶν ἐπὶ Βοῤῥᾶν εἰς μέρος τοῦ Ἰορδάνου ἀπὸ Λιβός· ταῦτα τὰ ὅριά ἐστιν ἀπὸ Λιβός.

\vs{20}Καὶ ὁ Ἰορδάνης ὁριεῖ ἀπὸ μέρους ἀνατολῶν· αὕτη ἡ κληρονομία υἱῶν Βενιαμὶν, τὰ ὅρια αὐτῆς κύκλῳ κατὰ δήμους.

\vs{21}Καὶ ἐγενήθησαν αἱ πόλεις τῶν υιῶν Βενιαμὶν κατὰ δήμους αὐτῶν Ἱερειχὼ, καὶ Βεθεγαιὼ, καὶ Ἀμεκασὶς,
\vs{22}καὶ Βαιθαβαρὰ, καὶ Σαρὰ, καὶ Βησανὰ,
\vs{23}καὶ Αἰεὶν, καὶ Φαρὰ, καὶ Ἐφραθὰ,
\vs{24}καὶ Καραφὰ, καὶ Κεφιρὰ, καὶ Μονὶ, καὶ Γαβαὰ, πόλεις δώδεκα· καὶ αἱ κῶμαι αὐτῶν,
\vs{25}Γαβαὼν, καὶ Ῥαμὰ, καὶ Βεηρωθὰ,
\vs{26}καὶ Μασσημὰ, καὶ Μιρὼν, καὶ Ἀμωκὴ,
\vs{27}καὶ Φιρὰ, καὶ Καφὰν, καὶ Νακὰν, καὶ Ζεληκὰν, καὶ Θαρεηλὰ,
\vs{28}καὶ Ἰηβοῦς· αὕτη ἐστὶν Ἱερουσαλήμ· καὶ Γαβαὼθ, Ἰαρὶμ, πόλεις δεκατρεῖς, καὶ αἱ κῶμαι αὐτῶν· αὕτη ἡ κληρονομία υἱῶν Βενιαμὶν κατὰ δήμους αὐτῶν.

\ch{19}
Καὶ ἐξῆλθεν ὁ δεύτερος κλῆρος τῶν υἱῶν Συμεών· καὶ ἐγενήθη ἡ κληρονομία αὐτῶν ἀναμέσον κλήρων υἱῶν Ἰούδα.
\vs{2}Καὶ ἐγενήθη ὁ κλῆρος αὐτῶν Βηρσαβεὲ, καὶ Σαμαὰ, καὶ Καλαδὰμ,
\vs{3}καὶ Ἀρσωλὰ, καὶ Βωλὰ, καὶ Ἰασὸν,
\vs{4}καὶ Ἐρθουλὰ, καὶ Βουλὰ, καὶ Ἑρμὰ,
\vs{5}καὶ Σικελὰκ, καὶ Βαιθμαχερὲβ, καὶ Σαρσουσὶν,
\vs{6}καὶ Βαθαρὼθ, καὶ οἱ ἀγροὶ αὐτῶν· πόλεις δεκατρεῖς, καὶ αἱ κῶμαι αὐτῶν.
\vs{7}Ἐρεμμὼν, καὶ Θαλχὰ, καὶ Ἰεθὲρ, καὶ Ἀσάν· πόλεις τέσσαρες καὶ αἱ κῶμαι αὐτῶν,
\vs{8}κύκλῳ τῶν πόλεων αὐτῶν ἕως Βαλὲκ πορευομένων Βαμὲθ κατὰ Λίβα· αὕτη ἡ κληρονομία φυλῆς υἱῶν Συμεὼν κατὰ δήμους αὐτῶν.
\vs{9}Ἀπὸ τοῦ κλήρου τοῦ Ἰούδα ἡ κληρονομία φυλῆς υἱῶν Συμεὼν, ὅτι ἐγενήθη ἡ μερὶς υἱῶν Ἰούδα μείζων τῆς αὐτῶν· καὶ ἐκληρονόμησαν οἱ υἱοὶ Συμεὼν ἐν μέσῳ τοῦ κλήρου αὐτῶν.

\vs{10}Καὶ ἐξῆλθεν ὁ κλῆρος ὁ τρίτος τῷ Ζαβουλὼν κατὰ δήμους αὐτῶν· ἔσται τὰ ὅρια τῆς κληρονομίας αὐτῶν, Ἐσεδεκγωλὰ ὅρια αὐτῶν,
\vs{11}ἡ θάλασσα καὶ Μαγελδὰ, καὶ συνάψει ἐπὶ Βαιθάραβὰ εἰς τὴν φάραγγα, ἥ ἐστι κατὰ πρόσωπον Ἰεκμάν.
\vs{12}Καὶ ἀνέστρεψεν ἀπὸ Σεδδοὺκ ἐξ ἐναντίας ἀπὸ ἀνατολῶν Βαιθσαμὺς ἐπὶ τὰ ὅρια Χασελωθαὶθ, καὶ διελεύσεται ἐπὶ Δαβιρὼθ, καὶ προσαναβήσεται ἐπὶ Φαγγαί.
\vs{13}Καὶ ἐκεῖθεν περιελεύσεται ἐξ ἐναντίας ἐπʼ ἀνατολὰς ἐπὶ Γεβερὲ ἐπὶ πόλιν Κατασὲμ, καὶ διελεύσεται ἐπὶ Ῥεμμωναὰ Μαθαραοζά.
\vs{14}Καὶ περιελεύσεται ὅρια ἐπὶ Βοῤῥᾶν ἐπὶ Ἀμὼθ, καὶ ἔσται ἡ διέξοδος αὐτῶν ἐπὶ Γαιφαὴλ,
\vs{15}καὶ Κατανὰθ, καὶ Ναβαὰλ, καὶ Συμοὼν, καὶ Ἱεριχὼ, καὶ Βαιθμάν.
\vs{16}Αὕτη ἡ κληρονομία τῆς φυλῆς υἱῶν Ζαβουλὼν κατὰ δήμους αὐτῶν, πόλεις καὶ αἱ κῶμαι αὐτῶν.

\vs{17}Καὶ τῷ Ἰσσὰχαρ ἐξῆλθεν ὁ κλῆρος ὁ τέταρτος.
\vs{18}Καὶ ἐγενήθη τὰ ὅρια αὐτῶν Ἰαζὴλ, καὶ Χασαλὼθ, καὶ Σουνὰμ,
\vs{19}καὶ Ἀγὶν, καὶ Σιωνὰ, καὶ Ῥεηρὼθ, καὶ Ἀναχερὲθ,
\vs{20}καὶ Δαβιρὼν, καὶ Κισὼν, καὶ Ῥεβές,
\vs{21}καὶ Ῥεμμὰς, καὶ Ἰεὼν, καὶ Τομμὰν, καὶ Αἰμαρὲκ, καὶ Βηρσαφής.
\vs{22}Καὶ συνάψει τὰ ὅρια ἐπὶ Γαιθβὼρ, καὶ ἐπὶ Σαλὶμ κατὰ θάλασσαν, καὶ Βαιθσαμύς· καὶ ἔσται αὐτοῦ ἡ διέξοδος τῶν ὁρίων ὁ Ἰορδάνης.
\vs{23}Αὕτη ἡ κληρονομία φυλῆς υἱῶν Ἰσσάχαρ κατὰ δήμους αὐτῶν, αἱ πόλεις καὶ αἱ κῶμαι αὐτῶν.

\vs{24}Καὶ ἐξῆλθεν ὁ κλῆρος ὁ πέμπτος Ἀσὴρ κατὰ δήμους αὐτῶν.
\vs{25}Καὶ ἐγενήθη τὰ ὅρια αὐτῶν Ἐξελεκὲθ, καὶ Ἀλὲφ, καὶ Βαιθὸκ, καὶ Κεὰφ,
\vs{26}καὶ Ἐλιμελὲχ, καὶ Ἀμιὴλ, καὶ Μαασά· καὶ συνάψει τῷ Καρμήλῳ κατὰ θάλασσαν, καὶ τῷ Σιὼν, καὶ Λαβανάθ.
\vs{27}Καὶ ἐπιστρέψει ἀπὸ ἀνατολῶν ἡλίου καὶ Βαιθεγενὲθ, καὶ συνάψει τῷ Ζαβουλὼν καὶ Ἐκγαῖ, καὶ Φθαιὴλ κατὰ Βοῤῥᾶν, καὶ εἰσελεύσεται ὅρια Σαφθαιβαιθμὲ, καὶ Ἰναὴλ, καὶ διελεύσεται εἰς Χωβαμασομὲλ,
\vs{28}καὶ Ἐλβὼν, καὶ Ῥαὰβ, καὶ Ἐμεμαὼν, καὶ Κανθὰν ἕως Σιδῶνος τῆς μεγάλης.
\vs{29}Καὶ ἀναστρέψει τὰ ὅρια εἰς Ῥαμὰ, καὶ ἕως πηγῆς Μασφασσὰτ, καὶ τῶν Τυρίων· καὶ ἀναστρέψει τὰ ὅρια ἐπὶ Ἰασὶφ, καὶ ἔσται ἡ διέξοδος αὐτοῦ ἡ θάλασσα, καὶ Ἀπολὲβ, καὶ Ἐχοζὸβ,
\vs{30}καὶ Ἀρχὸβ, καὶ Ἀφὲκ, καὶ Ῥααύ.
\vs{31}Αὕτη ἡ κληρονομία φυλῆς υἱῶν Ἀσὴρ κατὰ δήμους αὐτῶν, πόλεις καὶ αἱ κῶμαι αὐτῶν.

\vs{32}Καὶ τῷ Νεφθαλὶ ἐξῆλθεν ὁ κλῆρος ὁ ἕκτος.
\vs{33}Καὶ ἐγενήθη τὰ ὅρια αὐτῶν Μοολὰμ, καὶ Μωλὰ, καὶ Βεσεμιῒν, καὶ Ἀρμὲ, καὶ Ναβὸκ, καὶ Ἰεφθαμαὶ, ἕως Δωδάμ· καὶ ἐγενήθησαν αἱ διέξοδοι αὐτοῦ Ἰορδάνης.
\vs{34}Καὶ ἐπιστρέψει τὰ ὅρια ἐπὶ θάλασσαν ἐν Ἀθθαβὼρ, καὶ διελεύσεται ἐκεῖθεν Ἰακανὰ, καὶ συνάψει τῷ Ζαβουλὼν ἀπὸ Νότου, καὶ Ἀσὴρ συνάψει κατὰ θάλασσαν, καὶ ὁ Ἰορδάνης ἀπὸ ἀνατολῶν ἡλίου.

\vs{35}Καὶ αἱ πόλεις τειχήρεις τῶν Τυρίων, Τυρὸς, καὶ Ὠμαθαδακὲθ, καὶ Κενερέθ,
\vs{36}καὶ Ἀρμαὶθ, καὶ Ἀραὴλ, καὶ Ἀσὼρ,
\vs{37}καὶ Κάδες, καὶ Ἀσσαρὶ, καὶ πηγὴ Ἀσὸρ,
\vs{38}καὶ Κερωὲ, καὶ Μεγαλααρὶμ, καὶ Βαιθθαμὲ, καὶ Θεσσαμύς.
\vs{39}Αὕτη ἡ κληρονομία φυλῆς υἱῶν Νεφθαλί.

\vs{40}Καὶ τῷ Δὰν ἐξῆλθεν ὁ κλῆρος ὁ ἕβδομος·
\vs{41}Καὶ ἐγενήθη τὰ ὅρια αὐτῶν Σαρὰθ, καὶ Ἀσὰ, καὶ πόλεις Σαμμαὺς,
\vs{42}καὶ Σαλαμὶν, καὶ Ἀμμὸν, καὶ Σιλαθὰ,
\vs{43}καὶ Ἐλὼν, καὶ Θαμναθὰ, καὶ Ἀκκαρὼν,
\vs{44}καὶ Ἀλκαθὰ, καὶ Βεγεθὼν, καὶ Γεβεελὰν,
\vs{45}καὶ Ἀζὼρ, καὶ Βαναιβακὰτ, καὶ Γεθρεμμὼν,
\vs{46}καὶ ἀπὸ θαλάσσης Ἱερὰκων ὅριον πλησίον Ἰόππης.
\vs{47}Αὕτη ἡ κληρονομία φυλῆς υἱῶν Δὰν κατὰ δήμους αὐτῶν, αἱ πόλεις αὐτῶν καὶ αἱ κῶμαι αὐτῶν·
\vs{47a}καὶ οὐκ ἐξέθλιψαν οἱ υἱοὶ Δὰν τὸν Ἀμοῤῥαῖον τὸν θλίβοντα αὐτοὺς ἐν τῷ ὄρει· καὶ οὐκ εἴων αὐτοὺς οἱ Ἀμοῥῥαῖοι καταβῆναι εἰς τὴν κοιλάδα, καὶ ἔθλιψαν ἀπʼ αὐτῶν τὸ ὅριον τῆς μερίδος αὐτῶν.

\vs{48}Καὶ ἐπορεύθησαν οἱ υἱοὶ Δὰν καὶ ἐπολέμησαν τὴν Λαχὶς, καὶ κατελάβοντο αὐτὴν, καὶ ἐπάταξαν αὐτὴν ἐν στόματι μαχαίρας· καὶ κατῴκησαν αὐτὴν καὶ ἐκάλεσαν τὸ ὄνομα αὐτῆς Λασενδάν·
\vs{48a}καὶ ὁ Ἀμοῤῥαῖος ὑπέμεινε τοῦ κατοικεῖν ἐν Ἐλὼμ καὶ ἐν Σαλαμίν· καὶ ἐβαρύνθη ἡ χεὶρ τοῦ Ἐφραὶμ ἐπʼ αὐτοὺς, καὶ ἐγένοντο αὐτοῖς εἰς φόρον.

\vs{49}Καὶ ἐπορεύθησαν ἐμβατεύσαι τὴν γῆν κατὰ τὸ ὅριον αὐτῶν· καὶ ἔδωκαν οἱ υἱοὶ Ἰσραὴλ κλῆρον τῷ Ἰησοῖ τῷ υἱῷ Ναυῆ ἐν αὐτοῖς
\vs{50}διὰ προστάγματος τοῦ Θεοῦ, καὶ ἔδωκαν αὐτῷ τὴν πόλιν, ἣν ᾐτήσατο, Θαμνασαρὰχ, ἥ ἐστιν ἐν τῷ ὄρει Ἐφραίμ· καὶ ᾠκοδόμησε τὴν πόλιν, καὶ κατῴκει ἐν αὐτῇ.

\vs{51}Αὗται αἱ διαιρέσεις ἃς κατεκληρονόμησεν Ἐλεάζαρ ὁ ἱερεὺς, καὶ Ἰησοῦς ὁ τοῦ Ναυῆ, καὶ οἱ ἄρχοντες τῶν πατριῶν ἐν ταῖς φυλαῖς Ἰσραὴλ κατὰ κλήρους ἐν Σηλὼ ἔναντι Κυρίου, παρὰ τὰς θύρας τῆς σκηνῆς τοῦ μαρτυρίου· καὶ ἐπορεύθησαν ἐμβατεῦσαι τὴν γῆν.

\ch{20}
Καὶ ἐλάλησε Κύριος τῷ Ἰησοῖ, λέγων,
\vs{2}λάλησον τοῖς υἱοῖς Ἰσραὴλ, λέγων, δότε τὰς πόλεις τῶν φυγαδευτηρίων, ἃς εἶπα πρὸς ὑμᾶς διὰ Μωυσῆ.
\vs{3}Φυγαδευτήριον τῷ φονευτῇ τῷ πατάξαντι ψυχὴν ἀκουσίως· καὶ ἔσονται ὑμῖν αἱ πόλεις φυγαδευτήριον, καὶ οὐκ ἀποθανεῖται ὁ φονευτὴς ὑπὸ τοῦ ἀγχιστεύοντος τὸ αἷμα, ἕως ἂν καταστῇ· ἐναντίον τῆς συναγωγῆς εἰς κρίσιν.

\vs{7}Καὶ διέστειλε τὴν Κάδης ἐν τῇ Γαλιλαίᾳ ἐν τῷ ὄρει τῷ Νεφθαλὶ, καὶ Συχὲμ ἐν τῷ ὄρει τῷ Ἐφραὶμ, καὶ τὴν πόλιν Ἀρβὸκ, αὕτη ἐστὶ Χεβρὼν, ἐν τῷ ὄρει τῷ Ἰούδα.

\vs{8}Καὶ ἐν τῷ πέραν τοῦ Ἰορδάνου ἔδωκε Βοσὸρ ἐν τῇ ἐρήμῳ ἐν τῷ πεδίῳ ἀπὸ τῆς φυλῆς Ῥουβὴν, καὶ Ἀρημὼθ ἐν τῇ Γαλαὰδ ἐκ τῆς φυλῆς Γὰδ, καὶ τὴν Γαυλὼν ἐν τῇ Βασανίτιδι ἐκ τῆς φυλῆς Μανασσῆ.

\vs{9}Αὗται αἱ πόλεις αἱ ἐπίκλητοι τοῖς υἱοῖς Ἰσραὴλ καὶ τῷ προσηλύτῳ τῷ προσκειμένῳ ἐν αὐτοῖς, καταφυγεῖν ἐκεῖ παντὶ παίοντι ψυχὴν ἀκουσίως, ἵνα μὴ ἀποθάνῃ ἐν χειρὶ τοῦ ἀγχιστεύοντος τὸ αἷμα, ἕως ἄν καταστῇ ἔναντι τῆς συναγωγῆς εἰς κρίσιν.

\ch{21}
Καὶ προσήλθοσαν οἱ ἀρχιπατριῶται τῶν υἱῶν Λευὶ πρὸς Ἐλεάζαρ τὸν ἱερέα, καὶ Ἰησοῦν τὸν τοῦ Ναυῆ, καὶ πρὸς τοὺς ἀρχιφύλους πατριῶν ἐκ τῶν φυλῶν Ἰσραήλ·
\vs{2}Καὶ εἶπον πρὸς αὐτοὺς ἐν Σηλὼ ἐν γῇ Χαναὰν, λέγοντες, ἐνετείλατο Κύριος ἐν χειρὶ Μωυσῆ δοῦναι ἡμῖν πόλεις κατοικεῖν, καὶ τὰ περισπόρια τοῖς κτήνεσιν ἡμῶν.
\vs{3}Καὶ ἔδωκαν οἱ υἱοὶ Ἰσραὴλ τοῖς Λευίταις ἐν τῷ κατακληρονομεῖν διὰ προστάγματος Κυρίου τὰς πόλεις καὶ τὰ περισπόρια αὐτῶν.

\vs{4}Καὶ ἐξῆλθεν ὁ κλῆρος τῷ δήμῳ Καάθ· καὶ ἐγένετο τοῖς υἱοῖς Ἀαρὼν τοῖς ἱερεῦσι τοῖς Λευίταις ἀπὸ φυλῆς Ἰούδα καὶ ἀπὸ φυλῆς Συμεὼν καὶ ἀπὸ φυλῆς Βενιαμὶν κληρωτὶ, πόλεις δεκατρεῖς.

\vs{5}Καὶ τοῖς υἱοῖς Καὰθ τοῖς καταλελειμμένοις ἐκ τῆς φυλῆς Ἐφραὶμ καὶ ἐκ τῆς φυλῆς Δὰν καὶ ἀπὸ τοῦ ἡμίσους φυλῆς Μανασσῆ κληρωτὶ, πόλεις δέκα.

\vs{6}Καὶ τοῖς υἱοῖς Γεδσὼν ἀπὸ τῆς φυλῆς Ἰσσάχαρ καὶ ἀπὸ τῆς φυλῆς Ἀσὴρ καὶ ἀπὸ τῆς φυλῆς Νεφθαλὶ καὶ ἀπὸ τοῦ ἡμίσους φυλῆς Μανασσῆ ἐν τῇ Βασὰν, πόλεις δεκατρεῖς.

\vs{7}Καὶ τοῖς υἱοῖς Μεραρὶ κατὰ δήμους αὐτῶν ἀπὸ φυλῆς Ῥουβὴν καὶ ἀπὸ φυλῆς Γὰδ καὶ ἀπὸ φυλῆς Ζαβουλὼν κληρωτὶ, πόλεις δώδεκα.

\vs{8}Καὶ ἔδωκαν οἱ υἱοὶ Ἰσραὴλ τοῖς Λευίταις τὰς πόλεις καὶ τὰ περισπόρια αὐτῶν, ὃν τρόπον ἐνετείλατο Κύριος τῷ Μωυσῇ, κληρωτί.

\vs{9}Καὶ ἔδωκεν ἡ φυλὴ υἱῶν Ἰούδα καὶ ἡ φυλὴ υἱῶν Συμεὼν καὶ ἀπὸ τῆς φυλῆς υἱῶν Βενιαμὶν τὰς πόλεις ταύτας· καὶ ἐπεκλήθησαν
\vs{10}τοῖς υἱοῖς Ἀαρὼν ἀπὸ τοῦ δήμου τοῦ Καὰθ τῶν υἱῶν Λευὶ, ὅτι τούτοις ἐγενήθη ὁ κλῆρος.
\vs{11}Καὶ ἔδωκεν αὐτοῖς τὴν Καριαθαρβὸκ μητρόπολιν τῶν Ἐνάκ· αὕτη ἐστὶ Χεβρὼν ἐν τῷ ὄρει Ἰούδα· τὰ δὲ περισπόρια κύκλῳ αὐτῆς,
\vs{12}καὶ τοὺς ἀγροὺς τῆς πόλεως, καὶ τὰς κώμας αὐτῆς ἔδωκεν Ἰησοῦς τοῖς υἱοῖς Χάλεβ υἱοῦ Ἰεφοννὴ ἐν κατασχέσει.

\vs{13}Καὶ τοῖς υἱοῖς Ἀαρὼν ἔδωκε τὴν πόλιν φυγαδευτήριον τῷ φονεύσαντι, τὴν Χεβρὼν, καὶ τὰ ἀφωρισμένα τὰ σὺν αὐτῇ· καὶ τὴν Λεμνὰ, καὶ τὰ ἀφωρισμένα τὰ πρὸς αὐτῇ·
\vs{14}Καὶ τὴν Αἰλὼμ, καὶ τὰ ἀφωρισμένα αὐτῇ· καὶ τὴν Τεμὰ, καὶ τὰ ἀφωρισμένα αὐτῇ·
\vs{15}Καὶ τὴν Γελλὰ, καὶ τὰ ἀφωρισμένα αὐτῇ· καὶ τὴν Δαβὶρ, καὶ τὰ ἀφωρισμένα αὐτῇ·
\vs{16}Καὶ Ἀσὰ, καὶ τὰ ἀφωρισμένα αὐτῇ· καὶ Τανὺ, καὶ τὰ ἀφωρισμένα αὐτῇ· καὶ Βαιθσαμὺς, καὶ τὰ ἀφωρισμένα αὐτῇ· πόλεις ἐννέα παρὰ τῶν δύο φυλῶν τούτων.
\vs{17}Καὶ παρὰ τῆς φυλῆς Βενιαμὶν, τὴν Γαβαὼν, καὶ τὰ ἀφωρισμένα αὐτῇ· καὶ Γαθὲθ, καὶ τὰ ἀφωρισμένα αὐτῇ·
\vs{18}Καὶ Ἀναθὼθ, καὶ τὰ ἀφωρισμένα αὐτῇ· καὶ Γάμαλα, καὶ τὰ ἀφωρισμένα αὐτῇ· πόλεις τέσσαρες.
\vs{19}Πᾶσαι αἱ πόλεις υἱῶν Ἀαρὼν τῶν ἱερέων, δεκατρεῖς.

\vs{20}Καὶ τοῖς δήμοις υἱοῖς Καὰθ τοῖς Λευίταις τοῖς καταλελειμμένοις ἀπὸ τῶν υἱῶν Καὰθ, καὶ ἐγενήθη ἡ πόλις τῶν ἱερέων αὐτῶν ἀπὸ φυλῆς Ἐφραίμ·
\vs{21}καὶ ἔδωκαν αὐτοῖς τὴν πόλιν τοῦ φυγαδευτηρίου τὴν τοῦ φονεύσαντος, τὴν Συχὲμ, καὶ τὰ ἀφωρισμένα αὐτῇ· καὶ Γάζαρα καὶ τὰ πρὸς αὐτὴν, καὶ τὰ ἀφωρισμένα αὐτῇ·
\vs{22}Καὶ Βαιθωρὼν, καὶ τὰ ἀφωρισμένα τὰ αὐτῇ· πόλεις τέσσαρες.
\vs{23}Καὶ ἐκ τῆς φυλῆς Δὰν, τὴν Ἑλκωθαὶμ, καὶ τὰ ἀφωρισμένα αὐτῇ· καὶ τὴν Γεθεδὰν, καὶ τὰ ἀφωρισμένα αὐτῇ·
\vs{24}Καὶ Αἰλὼν, καὶ τὰ ἀφωρισμένα αὐτῇ· καὶ τὴν Γεθερεμμὼν, καὶ τὰ ἀφωρισμένα αὐτῇ· πόλεις τέσσαρες.
\vs{25}Καὶ ἀπὸ τοῦ ἡμίσους φυλῆς Μανασσῆ, τὴν Τανὰχ, καὶ τὰ ἀφωρισμένα αὐτῇ· καὶ τὴν Ἰεβαθὰ, καὶ τὰ ἀφωρισμένα αὐτῇ· πόλεις δύο.
\vs{26}Πᾶσαι πόλεις δέκα, καὶ τὰ ἀφωρισμένα αὐτῇ τὰ πρὸς αὐταῖς, τοῖς δήμοις υἱῶν Καὰθ τοῖς ὑπολελειμμένοις.

\vs{27}Καὶ τοῖς υἱοῖς Γεδσὼν τοῖς Λευίταις ἐκ τοῦ ἡμίσους φυλῆς Μανασσῆ τὰς πόλεις τὰς ἀφωρισμένας τοῖς φονεύσασι, τὴν Γαυλὼν ἐν τῇ Βασανίτιδι, καὶ τὰ ἀφωρισμένα αὐτῇ· καὶ τὴν Βοσορὰν, καὶ τὰ ἀφωρισμένα αὐτῇ· πόλεις δύο.
\vs{28}Καὶ ἐκ τῆς φυλῆς Ἰσσάχαρ, τὴν Κισὼν, καὶ τὰ ἀφωρισμένα αὐτῇ· καὶ τὴν Δεββὰ, καὶ τὰ ἀφωρισμένα αὐτῇ·
\vs{29}Καὶ τὴν Ῥεμμὰθ, καὶ τὰ ἀφωρισμένα αὐτῇ· καὶ Πηγὴν γραμμάτων, καὶ τὰ ἀφωρισμένα αὐτῇ· πόλεις τέσσαρες.
\vs{30}Καὶ ἐκ τῆς φυλῆς Ἀσὴρ τὴν Βασελλὰν, καὶ τὰ ἀφωρισμένα αὐτῇ· καὶ τὴν Δαββὼν, καὶ τὰ ἀφωρισμένα αὐτῇ·
\vs{31}Καὶ Χελκὰτ, καὶ τὰ ἀφωρισμένα αὐτῇ· καὶ τὴν Ῥαὰβ, καὶ τὰ ἀφωρισμένα αὐτῇ· πόλεις τέσσαρες.
\vs{32}Καὶ ἐκ τῆς φυλῆς Νεφθαλὶ, τὴν πόλιν τὴν ἀφωρισμένην τῷ φονεύσαντι, τὴν Κάδης ἐν τῇ Γαλιλαίᾳ, καὶ τὰ ἀφωρισμένα αὐτῇ· καὶ τὴν Νεμμὰθ, καὶ τὰ ἀφωρισμένα αὐτῇ· καὶ Θεμμὼν, καὶ τὰ ἀφωρισμένα αὐτῇ· πόλεις τρεῖς.
\vs{33}Πᾶσαι αἱ πόλεις τοῦ Γεδσὼν κατὰ δήμους αὐτῶν, πόλεις δεκατρεῖς.

\vs{34}Καὶ τῷ δήμῳ υἱῶν Μεραρὶ τοῖς Λευίταις τοῖς λοιποῖς ἐκ τῆς φυλῆς Ζαβουλὼν, τὴν Μαὰν, καὶ τὰ περισπόρια αὐτῆς· καὶ τὴν Κάδης, καὶ τὰ περισπόρια αὐτῆς·
\vs{35}Καὶ Σελλὰ, καὶ τὰ περισπόρια αὐτῆς· πόλεις τρεῖς.
\vs{36}Καὶ πέραν τοῦ Ἰορδάνου τοῦ κατὰ Ἰεριχὼ ἐκ τῆς φυλῆς Ῥουβὴν, τὴν πόλιν τὸ φυγαδευτήριον τοῦ φονεύσαντος, τὴν Βοσὸρ ἐν τῇ ἐρήμῳ· τὴν Μισὼ, καὶ τὰ περισπόρια αὐτῆς· καὶ τὴν Ἰαζὴρ, καὶ τὰ περισπόρια αὐτῆς·
\vs{37}καὶ τὴν Δεκμὼν, καὶ τὰ περισπόρια αὐτῆς· καὶ τὴν Μαφὰ, καὶ τὰ περισπόρια αὐτῆς· πόλεις τέσσαρες.
\vs{38}Καὶ ἀπὸ τῆς φυλῆς Γὰδ, τὴν πόλιν τὸ φυγαδευτήριον τοῦ φονεύσαντος, καὶ τὴν Ῥαμὼθ ἐν τῇ Γαλαὰδ, καὶ τὰ περισπόρια αὐτῆς· τὴν Καμὶν, καὶ τὰ περισπόρια αὐτῆς·
\vs{39}καὶ τὴν Ἐσβὼν, καὶ τὰ περισπόρια αὐτῆς· καὶ τὴν Ἰαζὴρ, καὶ τὰ περισπόρια αὐτῆς· πᾶσαι αἱ πόλεις τέσσαρες.
\vs{40}Πᾶσαι αἱ πόλεις τοῖς υἱοῖς Μεραρὶ κατὰ δήμους αὐτῶν τῶν καταλελειμμένων ἀπὸ τῆς φυλῆς τῆς Λευί· καὶ ἐγενήθη τὰ ὅρια αἱ πόλεις δεκαδύο.

\vs{41}Πᾶσαι πόλεις τῶν Λευιτῶν ἐν μέσῳ κατασχέσεως υἱῶν Ἰσραὴλ, τεσσαρακονταοκτὼ πόλεις, καὶ τὰ περισπόρια αὐτῶν
\vs{42}κύκλῳ τῶν πόλεων τούτων· πόλις καὶ τὰ περισπόρια κύκλῳ τῆς πόλεως πάσαις ταῖς πόλεσι ταύταις·
\vs{42a}καὶ συνετέλεσεν Ἰησοῦς διαμερίσας τὴν γῆν ἐν τοῖς ὁρίοις αὐτῶν·
\vs{42b}καὶ ἔδωκαν οἱ υἱοὶ Ἰσραὴλ μερίδα τῷ Ἰησοῖ διὰ πρόσταγμα Κυρίου· ἔδωκαν αὐτῷ τὴν πόλιν, ἣν ᾐτήσατο· τὴν Θαμνασαχὰρ ἔδωκαν αὐτῷ ἐν τῷ ὄρει Ἐφραίμ·
\vs{42c}καὶ ᾠκοδόμησεν Ἰησοῦς τὴν πόλιν, καὶ ᾤκησεν ἐν αὐτῇ·
\vs{42d}καὶ ἔλαβεν Ἰησοῦς τὰς μαχαίρας τὰς πετρίνας, ἐν αἷς περιέτεμε τοὺς υἱοὺς Ἰσραὴλ τοὺς γενομένους ἐν τῇ ὁδῷ ἐν τῇ ἐρήμῳ, καὶ ἔθηκεν αὐτὰς ἐν Θαμνασαχάρ.

\vs{43}Καὶ ἔδωκε Κύριος τῷ Ἰσραὴλ πᾶσαν τὴν γῆν, ἣν ὤμοσε δοῦναι τοῖς πατράσιν αὐτῶν· καὶ κατεκληρονόμησαν αὐτὴν, καὶ κατῴκησαν ἐν αὐτῇ.
\vs{44}Καὶ κατέπαυσεν αὐτοὺς Κύριος κυκλόθεν, καθότι ὤμοσε τοῖς πατράσιν αὐτῶν· οὐκ ἀνέστη οὐθεὶς κατενώπιον αὐτῶν ἀπὸ πάντων τῶν ἐχθρῶν αὐτῶν· πάντας τοὺς ἐχθροὺς αὐτῶν παρέδωκε Κύριος εἰς τὰς χεῖρας αὐτῶν.
\vs{45}Οὐ διέπεσεν ἀπὸ πάντων τῶν ῥημάτων τῶν καλῶν, ὧν ἐλάλησε Κύριος τοῖς υἱοῖς Ἰσραὴλ· πάντα παρεγένετο.

\ch{22}
Τότε συνεκάλεσεν Ἰησοῦς τοὺς υἱοὺς Ῥουβὴν, καὶ τοὺς υἱοὺς Γὰδ, καὶ τὸ ἥμισυ φυλῆς Μανασσῆ,
\vs{2}καὶ εἶπεν αὐτοῖς, ὑμεῖς ἀκηκόατε πάντα ὅσα ἐνετείλατο ὑμῖν Μωυσῆς ὁ παῖς Κυρίου, καὶ ὑπηκούσατε τῆς φωνῆς μου κατὰ πάντα ὅσα ἐνετείλατο ὑμῖν.
\vs{3}Οὐκ ἐγκαταλελοίπατε τοὺς ἀδελφοὺς ὑμῶν ταύτας τὰς ἡμέρας πλείους· ἕως τῆς σήμερον ἡυμέρας ἐφυλάξατε τὴν ἐντολὴν Κυρίου τοῦ Θεοῦ ὑμῶν.
\vs{4}Νῦν δὲ κατέπαυσε Κύριος ὁ Θεὸς ἡμῶν τοὺς ἀδελφοὺς ἡμῶν, ὃν τρόπον εἶπεν αὐτοῖς· νῦν οὖν ἀποστραφέντες, ἀπέλθατε εἰς τοὺς οἴκους ὑμῶν, καὶ εἰς τὴν γῆν τῆς κατασχέσεως ὑμῶν, ἣν ἔδωκεν ὑμῖν Μωυσῆς ἐν τῷ πέραν τοῦ Ἰορδάνου.
\vs{5}Ἀλλὰ φυλάξασθε σφόδρα ποιεῖν τὰς ἐντολὰς καὶ τὸν νόμον, ὃν ἐνετείλατο ἡμῖν ποιεῖν Μωυσῆς ὁ παῖς Κυρίου· ἀγαπᾷν Κύριον τὸν Θεὸν ἡμῶν, πορεύεσθαι πάσαις ταῖς ὁδοῖς αὐτοῦ, φυλάξασθαι τὰς ἐντολὰς αὐτοῦ, καὶ προσκεῖσθαι αὐτῷ, καὶ λατρεύειν αὐτῷ ἐξ ὅλης τῆς διανοίας ὑμῶν, καὶ ἐξ ὅλης τῆς ψυχῆς ὑμῶν.
\vs{6}Καὶ εὐλόγησεν αὐτοὺς Ἰησοῦς, καὶ ἐξαπέστειλεν αὐτούς· καὶ ἐπορεύθησαν εἰς τοὺς οἴκους αὐτῶν.

\vs{7}Καὶ τῷ ἡμίσει φυλῆς Μανασσῆ ἔδωκε Μωυσῆς ἐν τῇ Βασανίτιδι, καὶ τῷ ἡμίσει ἔδωκεν Ἰησοῦς μετὰ τῶν ἀδελφῶν αὐτοῦ ἐν τῷ πέραν τοῦ Ἰορδάνου παρὰ θάλασσαν· καὶ ἡνίκα ἐξαπέστειλεν αὐτοὺς Ἰησοῦς εἰς τοὺς οἴκους αὐτῶν καὶ εὐλόγησεν αὐτούς.
\vs{8}Καὶ ἐν χρήμασι πολλοῖς ἀπήλθοσαν εἰς τοὺς οἴκους αὐτῶν· καὶ κτήνη πολλὰ σφόδρα, καὶ ἀργύριον, καὶ χρυσίον, καὶ σίδηρον, καὶ ἱματισμὸν πολὺν, διείλαντο τὴν προνομὴν τῶν ἐχθρῶν μετὰ τῶν ἀδελφῶν αὐτῶν.

\vs{9}Καὶ ἐπορεύθησαν οἱ υἱοὶ Ῥουβὴν, καὶ οἱ υἱοὶ Γὰδ, καὶ τὸ ἥμισυ φυλῆς Μανασσῆ ἀπὸ τῶν υἱῶν Ἰσραὴλ ἐν Σηλὼ ἐν γῇ Χαναὰν ἀπελθεῖν εἰς τὴν Γαλαὰδ εἰς γῆν κατασχέσεως αὐτῶν, ἣν ἐκληρονόμησαν αὐτὴν διὰ προστάγματος Κυρίου ἐν χειρὶ Μωυσῆ.

\vs{10}Καὶ ἦλθον εἰς Γαλαὰδ τοῦ Ἰορδάνου, ἥ ἐστιν ἐν γῇ Χαναάν· καὶ ᾠκοδόμησαν οἱ υἱοὶ Ῥουβὴν, καὶ οἱ υἱοὶ Γὰδ, καὶ τὸ ἥμισυ φυλῆς Μανασσῆ ἐκεῖ βωμὸν ἐπὶ τοῦ Ἰορδάνου, βωμὸν μέγαν τοῦ ἰδεῖν.

\vs{11}Καὶ ἤκουσαν οἱ υἱοὶ Ἰσραὴλ λεγόντων, Ἰδοῦ ᾠκοδομήκασιν οἱ υἱοὶ Ῥουβὴν, καὶ οἱ υἱοὶ Γὰδ, καὶ τὸ ἥμισυ φυλῆς Μανασσῆ βωμὸν ἐφʼ ὁρίων γῆς Χαναὰν ἐπὶ τοῦ Γαλαὰδ τοῦ Ἰορδάνου ἐν τῷ πέραν υἱῶν Ἰσραήλ.
\vs{12}Καὶ συνηθροίσθησαν πάντες οἱ υἱοὶ Ἰσραὴλ εἰς Σηλὼ, ὥστε ἀναβάντες ἐκπολεμῆσαι αὐτούς.

\vs{13}Καὶ ἀπέστειλαν οἱ υἱοὶ Ἰσραὴλ πρὸς τοὺς υἱοὺς Ῥουβὴν καὶ πρὸς τοὺς υἱοὺς Γὰδ καὶ πρὸς τοὺς υἱοὺς ἥμισυ φυλῆς Μανασσῆ εἰς γῆν Γαλαὰδ, τόν τε Φινεὲς υἱὸν Ἐλεάζαρ υἱοῦ Ἀαρὼν τοῦ ἀρχιερέως,
\vs{14}καὶ δέκα τῶν ἀρχόντων μετʼ αὐτοῦ· ἄρχων εἷς ἀπὸ οἴκου πατριᾶς ἀπὸ πασῶν φυλῶν Ἰσραήλ· ἄρχοντες οἴκων πατριῶν εἰσι χιλίαρχοι Ἰσραήλ.
\vs{15}Καὶ παρεγένοντο πρὸς τοὺς υἱοὺς Ῥουβὴν, καὶ πρὸς τοὺς υἱοὺς Γὰδ, καὶ πρὸς τοὺς ἡμίσεις φυλῆς Μανασσῆ εἰς γῆν Γαλαάδ· καὶ ἐλάλησαν πρὸς αὐτοὺς, λέγοντες,
\vs{16}τάδε λέγει πᾶσα ἡ συναγωγὴ Κυρίου, τίς ἡ πλημμέλεια αὕτη, ἣν ἐπλημμελήσατε ἐναντίον τοῦ Θεοῦ Ἰσραὴλ, ἀποστραφῆναι σήμερον ἀπὸ Κυρίου, οἰκοδομήσαντες ὑμῖν ἑαυτοῖς βωμὸν, ἀποστάτας ὑμᾶς γενέσθαι ἀπὸ τοῦ Κυρίου;
\vs{17}Μὴ μικρὸν ἡμῖν τὸ ἁμάρτημα Φογὼρ, ὅτι οὐκ ἐκαθαρίσθημεν ἀπʼ αὐτοῦ ἕως τῆς ἡμέρας ταύτης; καὶ ἐγενήθη πληγὴ ἐν τῇ συναγωγῇ Κυρίου.
\vs{18}Καὶ ὑμεῖς ἀπεστράφητε σήμερον ἀπὸ Κυρίου· καὶ ἔσται ἐὰν ἀποστῆτε σήμερον ἀπὸ Κυρίου, καὶ αὔριον ἐπὶ πάντα Ἰσραὴλ ἔσται ἡ ὀργή.
\vs{19}Καὶ νῦν εἰ μικρὰ ἡ γῆ ὑμῶν τῆς κατασχέσεως ὑμῶν, διάβητε εἰς τὴν γῆν τῆς Κυρίου κατασχέσεως, οὗ κατασκηνοῖ ἐκεῖ ἡ σκηνὴ Κυρίου, καὶ κατακληρονομήσετε ἐν ἡμῖν· καὶ μὴ ἀπὸ Θεοῦ ἀποστάται γενήθητε, καὶ ὑμεῖς μηδʼ ἀπόστητε ἀπὸ Κυρίου, διὰ τὸ οἰκοδομῆσαι ὑμᾶς βωμὸν ἔξω τοῦ θυσιαστηρίου Κυρίου τοῦ Θεοῦ ἡμῶν.
\vs{20}Οὐκ ἰδοὺ Ἄχαρ ὁ τοῦ Ζαρᾶ πλημμελείᾳ ἐπλημμέλησεν ἀπὸ τοῦ ἀναθέματος, καὶ ἐπὶ πᾶσαν συναγωγὴν Ἰσραὴλ ἐγενήθη ὀργή; καὶ οὗτος εἷς αὐτὸς ἀπέθανε τῇ ἑαυτοῦ ἁμαρτίᾳ.

\vs{21}Καὶ ἀπεκρίθησαν οἱ υἱοὶ Ῥουβὴν, καὶ οἱ υἱοὶ Γὰδ, καὶ τὸ ἥμισυ φυλῆς Μανασσῆ, καὶ ἐλάλησαν τοῖς χιλιάρχοις Ἰσραὴλ, λέγοντες,
\vs{22}ὁ Θεὸς Θεὸς Κύριός ἐστι, καὶ ὁ Θεὸς Θεὸς αὐτὸς οἶδε, καὶ Ἰσραὴλ αὐτὸς γνώσεται· εἰ ἐν ἀποστασίᾳ ἐπλημμελήσαμεν ἔναντι τοῦ Κυρίου, μὴ ῥύσαιτο ἡμᾶς ἐν τῇ ἡμέρᾳ ταύτῃ.
\vs{23}Καὶ εἰ ᾠκοδομήσαμεν ἑαυτοῖς βωμὸν, ὥστε ἀποστῆναι ἀπὸ Κυρίου τοῦ Θεοῦ ἡμῶν, ὥστε ἀναβιβάσαι ἐπʼ αὐτὸν θυσίαν ὁλοκαυτωμάτων, ὥστε ποιῆσαι ἐπʼ αὐτοῦ θυσίαν σωτηρίου, Κύριος ἐκζητήσει.

\vs{24}Ἀλλʼ ἕνεκεν εὐλαβείας ῥήματος ἐποιήσαμεν τοῦτο, λέγοντες, ἵνα μὴ εἴπωσιν αὔριον τὰ τέκνα ὑμῶν τοῖς τέκνοις ἡμῶν, τί ὑμῖν Κυρίῳ τῷ Θεῷ Ἰσραήλ;
\vs{25}Καὶ ὅρια ἔθηκε Κύριος ἀναμέσον ἡμῶν καὶ ὑμῶν τὸν Ἰορδάνην, καὶ οὐκ ἔστιν ὑμῖν μερὶς Κυρίου· καὶ ἀπαλλοτριώσουσιν οἱ υἱοὶ ὑμῶν τοὺς υἱῶν ἡμῶν, ἵνα μὴ σέβωνται Κύριον.
\vs{26}Καὶ εἴπαμεν ποιῆσαι οὕτω, τοῦ οἰκοδομῆσαι τὸν βωμὸν τοῦτον οὐχ ἕνεκεν καρπωμάτων οὐδὲ ἕνεκεν θυσιῶν,
\vs{27}ἀλλʼ ἵνα ᾖ τοῦτο μαρτύριον ἀναμέσον ἡμῶν καὶ ὑμῶν, καὶ ἀναμέσον τῶν γενεῶν ἡμῶν μεθʼ ἡμᾶς, τοῦ λατρεύειν λατρείαν Κυρίου ἐναντίον αὐτοῦ, ἐν τοῖς καρπώμασιν ἡμῶν καὶ ἐν ταῖς θυσίαις ἡμῶν καὶ ἐν ταῖς θυσίαις τῶν σωτηρίων ἡμῶν· καὶ οὐκ ἐροῦσι τὰ τέκνα ὑμῶν τοῖς τέκνοις ἡμῶν αὔριον, οὐκ ἔστιν ὑμῖν μερὶς Κυρίου.
\vs{28}Καὶ εἴπαμεν, ἐὰν γένηταί ποτε καὶ λαλήσωσι πρὸς ἡμᾶς, ἢ ταῖς γενεαῖς ἡμῶν αὔριον, καὶ ἐροῦσιν, ἴδετε ὁμοίωμα τοῦ θυσιαστηρίου Κυρίου, ὃ ἐποίησαν οἱ πατέρες ἡμῶν οὐχ ἕνεκεν καρπωμάτων οὐδὲ ἕνεκεν θυσιῶν, ἀλλὰ μαρτύριόν ἐστιν ἀναμέσον ὑμῶν καὶ ἀναμέσον ἡμῶν, καὶ ἀναμέσον τῶν υἱῶν ἡμῶν.
\vs{29}Μὴ γένοιτο οὖν ἡμᾶς ἀποστραφῆναι ἀπὸ Κυρίου ἐν τῇ σήμερον ἡμέρᾳ ἀποστῆναι ἀπὸ Κυρίου, ὥστε οἰκοδομῆσαι ἡμᾶς θυσιαστήριον τοῖς καρπώμασι, καὶ ταῖς θυσίαις Σαλαμὶν, καὶ τῇ θυσίᾳ τοῦ σωτηρίου, πλὴν τοῦ θυσιαστηρίου Κυρίου ὅ ἐστιν ἐναντίον τῆς σκηνῆς αὐτοῦ.

\vs{30}Καὶ ἀκούσας Φινεὲς ὁ ἱερεὺς καὶ πάντες οἱ ἄρχοντες τῆς συναγωγῆς Ἰσραὴλ οἳ ἦσαν μετʼ αὐτοῦ, τοὺς λόγους οὓς ἐλάλησαν οἱ υἱοὶ Ῥουβὴν καὶ οἱ υἱοὶ Γὰδ καὶ τὸ ἥμισυ φυλῆς Μανασσῆ, καὶ ἤρεσεν αὐτοῖς.
\vs{31}Καὶ εἶπε Φινεὲς ὁ ἱερεὺς τοῖς υἱοῖς Ῥουβὴν καὶ τοῖς υἱοῖς Γὰδ καὶ τῷ ἡμίσει φυλῆς Μανασσῆ, σήμερον ἐγνώκαμεν ὅτι μεθʼ ἡμῶν Κύριος, διότι οὐκ ἐπλημμελήσατε ἐναντίον Κυρίου πλημμέλειαν, καὶ ὅτι ἐῤῥύσασθε τοὺς υἱοὺς Ἰσραὴλ ἐκ χειρὸς Κυρίου.
\vs{32}Καὶ ἀπέστρεψε Φινεὲς ὁ ἱερεὺς καὶ οἱ ἄρχοντες ἀπὸ τῶν υἱῶν Ῥουβὴν καὶ ἀπὸ τῶν υἱῶν Γὰδ καὶ ἀπὸ τοῦ ἡμίσους φυλῆς Μανασσῆ ἐκ τῆς Γαλαὰδ εἰς γῆν Χαναὰν πρὸς τοὺς υἱοὺς Ἰσραήλ· καὶ ἀπεκρίθησαν αὐτοῖς τοὺς λόγους.
\vs{33}Καὶ ἤρεσε τοῖς υἱοῖς Ἰσραήλ· καὶ ἐλάλησαν πρὸς τοὺς υἱοὺς Ἰσραὴλ, καὶ εὐλόγησαν τὸν Θεὸν υἱῶν Ἰσραήλ· καὶ εἶπαν μηκέτι ἀναβῆναι πρὸς αὐτοὺς εἰς πόλεμον ἐξολοθρεῦσαι τὴν γῆν τῶν υἱῶν Ῥουβὴν καὶ τῶν υἱων Γὰδ καὶ τοῦ ἡμίσους φυλῆς Μανασσῆ· καὶ κατῴκησαν ἐπʼ αὐτῆς.

\vs{34}Καὶ ἐπωνόμασεν Ἰησοῦς τὸν βωμὸν τῶν Ῥουβὴν καὶ τῶν Γὰδ καὶ τοῦ ἡμίσους φυλῆς Μανασσῆ, καὶ εἶπεν, ὅτι μαρτύριόν ἐστιν ἀναμέσον αὐτῶν, ὅτι Κύριος ὁ Θεὸς αὐτῶν ἐστι.

\ch{23}
Καὶ ἐγένετο μεθʼ ἡμέρας πλείους μετὰ τὸ καταπαῦσαι Κύριον τὸν Ἰσραὴλ ἀπὸ πάντων τῶν ἐχθρῶν αὐτοῦ κυκλόθεν, καὶ Ἰησοῦς πρεσβύτερος προβεβηκὼς ταῖς ἡμέραις.
\vs{2}Καὶ συνεκάλεσεν Ἰησοῦς πάντας τοὺς υἱοὺς Ἰσραὴλ καὶ τὴν γερουσίαν αὐτῶν καὶ τοὺς ἄρχοντας αὐτῶν καὶ τοὺς δικαστὰς αὐτῶν καὶ τοὺς γραμματεῖς αὐτῶν, καὶ εἶπε πρὸς αὐτοὺς, ἐγὼ γεγήρακα καὶ προβέβηκα ταῖς ἡμέραις·
\vs{3}Ὑμεῖς δὲ ἑωράκατε ὅσα ἐποίησε Κύριος ὁ Θεὸς ἡμῶν πᾶσι τοῖς ἔθνεσι τούτοις ἀπὸ προσώπου ἡμῶν, ὅτι Κύριος ὁ Θεὸς ὑμῶν ὁ ἐκπολεμήσας ὑμῖν.
\vs{4}Ἴδετε ὅτι ἐπέῤῥιφα ὑμῖν τὰ ἔθνη τὰ καταλελειμμένα ὑμῖν ταῦτα ἐν τοῖς κλήροις εἰς τὰς φυλὰς ὑμῶν, ἀπὸ τοῦ Ἰορδάνου πάντα τὰ ἔθνη, καὶ ἐξωλόθρευσα, καὶ ἀπὸ τῆς θαλάσσης τῆς μεγάλης ὁριεῖ ἐπὶ δυσμὰς ἡλίου.

\vs{5}Κύριος δὲ ὁ Θεὸς ἡμῶν οὗτος ἐξολοθρεύσει αὐτοὺς ἀπὸ προσώπου ἡμῶν, ἕως ἂν ἀπόλωνται· καὶ ἀποστελεῖ αὐτοῖς τὰ θηρία τὰ ἄγρια, ἕως ἂν ἐξολοθρεύσῃ αὐτοὺς καὶ τοὺς βασιλεῖς αὐτῶν ἀπὸ προσώπου ὑμῶν, καὶ κατακληρονομήσετε τὴν γῆν αὐτῶν, καθὰ ἐλάλησε Κύριος ὁ Θεὸς ἡμῶν ὑμῖν.
\vs{6}Κατισχύσατε οὖν σφόδρα φυλάσσειν καὶ ποιεῖν πάντα τὰ γεγραμμένα ἐν τῷ βιβλίῳ τοῦ νόμου Μωυσῆ, ἵνα μὴ ἐκκλίνητε εἰς δεξιὰ ἢ εὐώνυμα,
\vs{7}ὅπως μὴ εἰσέλθητε εἰς τὰ ἔθνη τὰ καταλελειμμένα ταῦτα· καὶ τὰ ὀνόματα τῶν θεῶν αὐτῶν οὐκ ὀνομασθήσεται ἐν ὑμῖν, οὐδὲ μὴ λατρεύσητε, οὐδὲ μὴ προσκυνήσητε αὐτοῖς,
\vs{8}ἀλλὰ Κυρίῳ τῷ Θεῷ ἡμῶν προσκολληθήσεσθε, καθάπερ ἐποιήσατε ἕως τῆς ἡμέρας ταύτης.
\vs{9}Καὶ ἐξολοθρεύσει αὐτοὺς Κύριος ἀπὸ προσώπου ὑμῶν ἔθνη μεγάλα καὶ ἰσχυρά· καὶ οὐδεὶς ἀντέστη κατενώπιον ἡμῶν ἕως τῆς ἡμέρας ταύτης.
\vs{10}Εἷς ὑμῶν ἐδίωξε χιλίους, ὅτι Κύριος ὁ Θεὸς ἡμῶν οὗτος ἐξεπολέμει ὑμῖν, καθάπερ εἶπεν ἡμῖν.

\vs{11}Καὶ φυλάξασθε σφόδρα τοῦ ἀγαπᾷν Κύριον τὸν Θεὸν ἡμῶν.
\vs{12}Ἐὰν γὰρ ἀποστραφῆτε καὶ προσθῆσθε τοῖς ὑπολειφθεῖσιν ἔθνεσι τούτοις τοῖς μεθʼ ὑμῶν, καὶ ἐπιγαμίας ποιήσητε πρὸς αὐτοὺς, καὶ συγκαταμιγῆτε αὐτοῖς καὶ αὐτοὶ ὑμῖν,
\vs{13}γινώσκετε ὅτι οὐ μὴ προσθῇ Κύριος τοῦ ἐξολοθρεῦσαι τὰ ἔθνη ταῦτα ἀπὸ προσώπου ὑμῶν· καὶ ἔσονται ὑμῖν εἰς παγίδας, καὶ εἰς σκάνδαλα, καὶ εἰς ἥλους ἐν ταῖς πτέρναις ὑμῶν, καὶ εἰς βολίδας ἐν τοῖς ὀφθαλμοῖς ὑμῶν, ἕως ἂν ἀπόλησθε ἀπὸ τῆς γῆς τῆς ἀγαθῆς ταύτης, ἣν ἔδωκεν ὑμῖν Κύριος ὁ Θεὸς ὑμῶν.

\vs{14}Ἐγὼ δὲ ἀποτρέχω τὴν ὁδὸν, καθὰ καὶ πάντες οἱ ἐπὶ τῆς γῆς· καὶ γνώσεσθε τῇ καρδίᾳ ὑμῶν καὶ τῇ ψυχῇ ὑμῶν, διότι οὐκ ἔπεσεν εἷς λόγος ἀπὸ πάντων τῶν λόγων, ὧν εἶπε Κύριος ὁ Θεὸς ἡμῶν πρὸς πάντα τὰ ἀνήκοντα ἡμῖν, οὐ διεφώνησεν ἐξ αὐτῶν.
\vs{15}Καὶ ἔσται ὃν τρόπον ἥκει πρὸς ἡμᾶς πάντα τὰ ῥήματα τὰ καλὰ, ἃ ἐλάλησε Κύριος ἐφʼ ὑμᾶς· οὕτως ἐπάξει Κύριος ὁ Θεὸς ἐφʼ ὑμᾶς πάντα τὰ ῥήματα τὰ πονηρὰ ἕως ἂν ἐξολοθρεύσῃ ὑμᾶς ἀπὸ τῆς γῆς τῆς ἀγαθῆς ταύτης, ἧς ἔδωκε Κύριος ὑμῖν,
\vs{16}ἐν τῷ παραβῆναι ὑμᾶς τὴν διαθήκην Κυρίου τοῦ Θεοῦ ἡμῶν, ἣν ἐνετείλατο ἡμῖν, καὶ πορευθέντες λατρεύσητε θεοῖς ἑτέροις καὶ προσκυνήσητε αὐτοῖς.

\ch{24}
Καὶ συνήγαγεν Ἰησοῦς πάσας φυλὰς Ἰσραὴλ εἰς Σηλὼ, καὶ συνεκάλεσε τοὺς πρεσβυτέρους αὐτῶν καὶ τοὺς γραμματεῖς αὐτῶν καὶ τοὺς δικαστὰς αὐτῶν, καὶ ἔστησεν αὐτοὺς ἀπέναντι τοῦ Θεοῦ.

\vs{2}Καὶ εἶπεν Ἰησοῦς πρὸς πάντα τὸν λαὸν, τάδε λέγει Κύριος ὁ Θεὸς Ἰσραὴλ, πέραν τοῦ ποταμοῦ παρῴκησαν οἱ πατέρες ὑμῶν τὸ ἀπαρχῆς, Θάρα ὁ πατὴρ Ἁβραὰμ, καὶ ὁ πατὴρ Ναχὼρ, καὶ ἐλάτρευσαν θεοῖς ἑτέροις.
\vs{3}Καὶ ἔλαβον τὸν πατέρα ὑμῶν τὸν Ἁβραὰμ ἐκ τοῦ πέραν τοῦ ποταμοῦ, καὶ ὡδήγησα αὐτὸν ἐν πάσῃ τῇ γῇ· καὶ ἐπλήθυνα αὐτοῦ σπέρμα, καὶ ἔδωκα αὐτῷ τὸν Ἰσαὰκ,
\vs{4}καὶ τῷ Ἰσαὰκ τὸν Ἰακὼβ καὶ τὸν Ἡσαῦ· καὶ ἔδωκα τῷ Ἡσαῦ τὸ ὄρος τὸ Σηεὶρ κληρονομῆσαι αὐτῷ· καὶ Ἰακὼβ καὶ οἱ υἱοὶ αὐτοῦ κατέβησαν εἰς Αἴγυπτον, καὶ ἐγένοντο ἐκεῖ εἰς ἔθνος μέγα καὶ πολὺ καὶ κραταιόν· καὶ ἐκάκωσαν αὐτοὺς οἱ Αἰγύπτιοι.
\vs{5}Καὶ ἐπάταξα τὴν Αἴγυπτον ἐν σημείοις οἷς ἐποίησα ἐν αὐτοῖς.
\vs{6}Καὶ μετὰ ταῦτα ἐξήγαγε τοὺς πατέρας ἡμῶν ἐξ Αἰγύπτου, καὶ εἰσήλθατε εἰς τὴν θάλασσαν τὴν ἐρυθράν· καὶ κατεδίωξαν οἱ Αἰγύπτιοι ὀπίσω τῶν πατέρων ἡμῶν ἐν ἅρμασι καὶ ἐν ἵπποις εἰς τὴν θάλασσαν τὴν ἐρυθράν.
\vs{7}Καὶ ἀνεβοήσαμεν πρὸς Κύριον· καὶ ἔδωκε νεφέλην καὶ γνόφον ἀναμέσον ἡμῶν καὶ ἀναμέσον τῶν Αἰγυπτίων, καὶ ἐπήγαγεν ἐπʼ αὐτοὺς τὴν θάλασσαν, καὶ ἐκάλυψεν αὐτούς· καὶ εἴδοσαν οἱ ὀφθαλμοὶ ὑμῶν ὅσα ἐποίησε Κύριος ἐν γῇ· Αἰγύπτῳ· καὶ ἦτε ἐν τῇ ἐρήμῳ ἡμέρας πλείους.

\vs{8}Καὶ ἤγαγεν ἡμᾶς εἰς γῆν Ἀμοῤῥαίων τῶν κατοικούντων πέραν τοῦ Ἰορδάνου, καὶ παρέδωκεν αὐτοὺς Κύριος εἰς τὰς χεῖρας ἡμῶν· καὶ κατεκληρονομήσατε τὴν γῆν αὐτῶν, καὶ ἐξωλοθρεύσατε αὐτοὺς ἀπὸ προσώπου ὑμῶν.

\vs{9}Καὶ ἀνέστη Βαλὰκ ὁ τοῦ Σεπφὼρ βασιλεὺς Μωὰβ, καὶ παρετάξατο τῷ Ἰσραὴλ, καὶ ἀποστείλας ἐκάλεσε τὸν Βαλαὰμ ἀράσασθαι ἡμῖν.
\vs{10}Καὶ οὐκ ἠθέλησε Κύριος ὁ Θεός σου ἀπολέσαι σε· καὶ εὐλογίαις εὐλόγησεν ἡμᾶς, καὶ ἐξείλατο ἡμᾶς ἐκ χειρῶν αὐτῶν, καὶ παρέδωκεν αὐτούς.
\vs{11}Καὶ διέβητε τὸν Ἰορδάνην, καὶ παρεγενήθητε εἰς Ἱεριχώ· καὶ ἐπολέμησαν πρὸς ἡμᾶς οἱ κατοικοῦντες Ἱεριχὼ ὁ Ἀμοῤῥαῖος, καὶ ὁ Χαναναῖος, καὶ ὁ Φερεζαῖος, καὶ ὁ Εὑαῖος, καὶ ὁ Ἰεβουσαῖος, καὶ ὁ Χετταῖος, καὶ ὁ Γεργεσαῖος· καὶ παρέδωκεν αὐτοὺς Κύριος εἰς τὰς χεῖρας ἡμῶν.
\vs{12}Καὶ ἐξαπέστειλε προτέραν ὑμῶν τὴν σφηκίαν· καὶ ἐξαπέστειλεν αὐτοὺς ἀπὸ προσώπου ἡμῶν δώδεκα βασιλεῖς τῶν Ἀμοῤῥαίων, οὐκ ἐν τῇ ῥομφαίᾳ σου οὐδὲ ἐν τῷ τόξῳ σου.

\vs{13}Καὶ ἔδωκεν ὑμῖν γῆν ἐφʼ ἣν οὐκ ἐκοπιάσατε ἐπʼ αὐτῆς, καὶ πόλεις ἃς οὐκ ᾠκοδομήκατε, καὶ κατῳκίσθητε ἐν αὐταῖς, καὶ ἀμπελῶνας καὶ ἐλαιῶνας οὓς οὐκ ἐφυτεύσατε ὑμεῖς, ἔδεσθε.

\vs{14}Καὶ νῦν φοβήθητε Κύριον, καὶ λατρεύσατε αὐτῷ ἐν εὐθύτητι καὶ ἐν δικαιοσύνῃ, καὶ περιέλεσθε τοὺς θεοὺς τοὺς ἀλλοτρίους, οἷς ἐλάτρευσαν οἱ πατέρες ἡμῶν ἐν τῷ πέραν τοῦ ποταμοῦ καὶ ἐν Αἰγύπτῳ, καὶ λατρεύσατε Κυρίῳ.
\vs{15}Εἰ δὲ μὴ ἀρέσκει ὑμῖν λατρεύειν Κυρίῳ, ἐκλέξασθε ὑμῖν αὐτοῖς σήμερον τίνι λατρεύσητε, εἴτε τοῖς θεοῖς τῶν πατέρων ὑμῶν, τοῖς ἐν τῷ πέραν τοῦ ποταμοῦ, εἴτε τοῖς θεοῖς τῶν Ἀμοῤραίων, ἐν οἷς ὑμεῖς κατοικεῖτε ἐπὶ τῆς γῆς αὐτῶν· ἐγὼ δὲ καὶ ἡ οἰκία μου λατρεύσομεν Κυρίῳ, ὅτι ἅγιός ἐστι.

\vs{16}Καὶ ἀποκριθεὶς ὁ λαὸς εἶπε, μὴ γένοιτο ἡμῖν καταλιπεῖν Κύριον, ὥστε λατρεύειν θεοῖς ἑτέροις.
\vs{17}Κύριος ὁ Θεὸς ἡμῶν, αὐτὸς Θεός ἐστιν· αὐτὸς ἀνήγαγεν ἡμᾶς καὶ τοὺς πατέρας ἡμῶν ἐξ Αἰγύπτου, καἰ διεφύλαξεν ἡμᾶς ἐν πάσῃ τῇ ὁδῷ ᾗ ἐπορεύθημεν ἐν αὐτῇ, καὶ ἐν πᾶσι τοῖς ἔθνεσιν οὓς παρήλθομεν διʼ αὐτῶν.
\vs{18}Καὶ ἐξέβαλε Κύριος τὸν Ἀμοῤῥαῖον καὶ πάντα τὰ ἔθνη τὰ κατοικοῦντα τὴν γῆν ἀπὸ προσώπου ἡμῶν· ἀλλὰ καὶ ἡμεῖς λατρεύσομεν Κυρίῳ, οὗτος γὰρ Θεὸς ἡμῶν ἐστι.

\vs{19}Καὶ εἶπεν Ἰησοῦς πρὸς τὸν λαὸν, οὐ μὴ δύνησθε λατρεύειν Κυρίῳ, ὅτι Θεὸς ἅγιός ἐστι· καὶ ζηλώσας οὗτος οὐκ ἀνήσει τὰ ἁμαρτήματα ὑμῶν καὶ τὰ ἀνομήματα ὑμῶν,
\vs{20}ἡνίκα ἂν ἐγκαταλίπητε Κύριον καὶ λατρεύσητε θεοῖς ἑτέροις· καὶ ἐπελθὼν κακώσει ὑμᾶς καὶ ἐξαναλώσει ὑμᾶς ἀνθʼ ὧν εὖ ἐποίησεν ὑμᾶς.
\vs{21}Καὶ εἶπεν ὁ λαὸς πρὸς Ἰησοῦν, οὐχὶ, ἀλλὰ Κυρίῳ λατρεύσομεν.

\vs{22}Καὶ εἶπεν Ἰησοῦς πρὸς τὸν λαὸν, μάρτυρες ὑμεῖς καθʼ ὑμῶν, ὅτι ὑμεῖς ἐξελέξασθε Κυρίῳ λατρεύειν αὐτῷ.
\vs{23}Καὶ νῦν περιέλεσθε τοὺς θεοὺς τοὺς ἀλλοτρίους τοὺς ἐν ὑμῖν, καὶ εὐθύνατε τὴν καρδίαν ὑμῶν πρὸς Κύριον Θεὸν Ἰσραήλ.
\vs{24}Καὶ εἶπεν ὁ λαὸς πρὸς Ἰησοῦν, Κυρίῳ λατρεύσομεν καὶ τῆς φωνῆς αὐτοῦ ἀκουσόμεθα.

\vs{25}Καὶ διέθετο Ἰησοῦς διαθήκην πρὸς τὸν λαὸν ἐν τῇ ἡμέρᾳ ἐκείνῃ, καὶ ἔδωκεν αὐτῷ νόμον καὶ κρίσιν ἐν Σηλὼ ἐνώπιον τῆς σκηνῆς τοῦ Θεοῦ Ἰσραήλ.
\vs{26}Καὶ ἔγραψε τὰ ῥήματα ταῦτα εἰς βιβλίον νόμων τοῦ Θεοῦ· καὶ ἔλαβε λίθον μέγαν, καὶ ἔστησεν αὐτὸν Ἰησοῦς ὑπὸ τὴν τέρμινθον ἀπέναντι Κυρίου.
\vs{27}Καὶ εἶπεν Ἰησοῦς πρὸς τὸν λαὸν, ἰδοὺ ὁ λίθος οὗτος ἔσται ἐν ὑμῖν εἰς μαρτύριον, ὅτι αὐτὸς ἀκήκοε πάντα τὰ λεχθέντα αὐτῷ ὑπὸ Κυρίου· ὅτι ἐλάλησε πρὸς ὑμᾶς σήμερον, καὶ οὗτος ἔσται ἐν ὑμῖν εἰς μαρτύριον ἐπʼ ἐσχάτων τῶν ἡμερῶν, ἡνίκα ἂν ψεύσησθε Κυρίῳ τῷ Θεῷ μου.
\vs{28}Καὶ ἀπέστειλεν Ἰησοῦς τὸν λαὸν, καὶ ἐπορεύθησαν ἕκαστος εἰς τὸν τόπον αὐτοῦ.
\vs{29}Καὶ ἐλάτρευσεν Ἰσραὴλ τῷ Κυρίῳ πάσας τὰς ἡμέρας Ἰησοῦ, καὶ πάσας τὰς ἡμέρας τῶν πρεσβυτέρων ὅσοι ἐφείλκυσαν τὸν χρόνον μετὰ Ἰησοῦ, καὶ ὅσοι εἴδοσαν πάντα τὰ ἔργα Κυρίου ὅσα ἐποίησε τῷ Ἰσραήλ.

\vs{30}Καὶ ἐγένετο μετʼ ἐκεῖνα καὶ ἀπέθανεν Ἰησοῦς υἱὸς Ναυῆ δοῦλος Κυρίου ἐκατὸν δέκα ἐτῶν.
\vs{31}Καὶ ἔθαψαν αὐτὸν πρὸς τοῖς ὁρίοις τοῦ κλήρου αὐτοῦ ἐν Θαμνασαρὰχ ἐν τῷ ὄρει τῷ Ἐφραὶμ ἀπὸ βοῤῥᾶ τοῦ ὄρους τοῦ Γαλαάδ·
\vs{31a}ἐκεῖ ἔθηκαν μετʼ αὐτοῦ εἰς τὸ μνῆμα εἰς ὃ ἔθαψαν αὐτὸν ἐκεῖ τὰς μαχαίρας τὰς πετρίνας, ἐν αἷς περιέτεμε τοὺς υἱοὺς Ἰσραὴλ ἐν Γαλγάλοις, ὅτε ἐξήγαγεν αὐτοὺς ἐξ Αἰγύπτου, καθὰ συνέταξεν αὐτοῖς Κύριος· καὶ ἐκεῖ εἰσιν ἕως τῆς σήμερον ἡμέρας.

\vs{32}Καὶ τὰ ὀστᾶ Ἰωσὴφ ἀνήγαγον οἱ υἱοὶ Ἰσραὴλ ἐξ Αἰγύπτου, καὶ κατώρυξαν ἐν Σικίμοις, ἐν τῇ μερίδι τοῦ ἀγροῦ οὗ ἐκτήσατο Ἰακὼβ παρὰ τῶν Ἀμοῤῥαίων τῶν κατοικούντων ἐν Σικίμοις ἀμνάδων ἑκατὸν, καὶ ἔδωκεν αὐτὴν Ἰωσὴφ ἐν μερίδι.

\vs{33}Καὶ ἐγένετο μετὰ ταῦτα καὶ Ἐλεάζαρ υἱὸς Ἀαρὼν ὁ ἀρχιερεὺς ἐτελεύτησε, καὶ ἐτάφη ἐν Γαβαὰρ Φινεὲς τοῦ υἱοῦ αὐτοῦ, ἣν ἔδωκεν αὐτῷ ἐν τῷ ὄρει τῷ Ἐφραίμ.

\vs{33a}Ἐν ἐκείνῃ τῇ ἡμέρᾳ λαβόντες οἱ υἱοὶ Ἰσραὴλ τὴν κιβωτὸν τοῦ Θεοῦ, περιεφέροσαν ἐν ἑαυτοῖς· καὶ Φινεὲς ἱεράτευσεν ἀντὶ Ἐλεάζαρ τοῦ πατρὸς αὐτοῦ ἕως ἀπέθανε, καὶ κατωρύγη ἐν Γαβαὰρ τῇ ἑαυτοῦ·
\vs{33b}οἱ δὲ υἱοὶ Ἰσραὴλ ἀπήλθοσαν ἕκαστος εἰς τὸν τόπον αὐτῶν, καὶ εἰς τὴν ἑαυτῶν πόλιν· καὶ ἐσέβοντο οἱ υἱοὶ Ἰσραὴλ τὴν Ἀστάρτην, καὶ Ἀσταρὼθ, καὶ τοὺς θεοὺς τῶν ἐθνῶν τῶν κύκλῳ αὐτῶν· καὶ παρέδωκεν αὐτοὺς Κύριος εἰς χεῖρας Ἐγλὼμ τῷ βασιλεῖ Μωὰβ, καὶ ἐκυρίευσεν αὐτῶν ἔτη δεκαοκτώ.


\def\book{ΚΡΙΤΑΙ}
\biblebook{ΚΡΙΤΑΙ}


\lettrine[lines=2, loversize=0.2, nindent=0em, findent=.25em]{\textcolor{bookheadingcolor}{Κ}}{ΑΙ} ἐγένετο μετὰ τὴν τελευτὴν Ἰησοῦ, καὶ ἐπηρώτων οἱ υἱοὶ Ἰσραὴλ διὰ τοῦ Κυρίου, λέγοντες, τίς ἀναβήσεται ἡμῖν πρὸς τοὺς Χαναναίους ἀφηγούμενος τοῦ πολεμῆσαι πρὸς αὐτούς;
\vs{2}Καὶ εἶπε Κύριος, Ἰούδας ἀναβήσεται· ἰδοὺ δέδωκα τὴν γῆν ἐν χειρὶ αὐτοῦ.
\vs{3}Καὶ εἶπεν Ἰούδας τῷ Συμεὼν ἀδελφῷ αὐτοῦ, ἀνάβηθι μετʼ ἐμοῦ ἐν τῷ κλήρῳ μου, καὶ παραταξώμεθα πρὸς τοὺς Χαναναίους, καὶ πορεύσομαι κᾀγὼ μετὰ σοῦ ἐν τῷ κλήρῳ σου· καὶ ἐπορεύθη μετʼ αὐτοῦ Συμεών.
\vs{4}Καὶ ἀνέβη Ἰούδας· καὶ παρέδωκε Κύριος τὸν Χαναναῖον καὶ τὸν Φερεζαῖον εἰς τὰς χεῖρας αὐτῶν· καὶ ἔκοψαν αὐτοὺς ἐν Βεζὲκ εἰς δέκα χιλιάδας ἀνδρῶν·
\vs{5}Καὶ κατέλαβον τὸν Ἀδωνιβεζὲκ ἐν τῇ Βεζὲκ, καὶ παρετάξαντο πρὸς αὐτόν· καὶ ἔκοψαν τὸν Χαναναῖον καὶ Φερεζαῖον.
\vs{6}Καὶ ἔφυγεν Ἀδωνιβεζέκ· καὶ κατέδραμον ὀπίσω αὐτοῦ, καὶ ἐλάβοσαν αὐτὸν, καὶ ἀπέκοψαν τὰ ἄκρα τῶν χειρῶν αὐτοῦ καὶ τὰ ἄκρα τῶν ποδῶν αὐτοῦ.
\vs{7}Καὶ εἶπεν Ἀδωνιβεζὲκ, ἑβδομήκοντα βασιλεῖς, τὰ ἄκρα τῶν χειρῶν αὐτῶν καὶ τὰ ἄκρα τῶν ποδῶν αὐτῶν ἀποκεκομμένοι, ἦσαν συλλέγοντες τὰ ὑποκάτω τῆς τραπέζης μου· καθὼς οὖν ἐποίησα, οὕτως ἀνταπέδωκέ μοι ὁ Θεός. καὶ ἄγουσιν αὐτὸν εἰς Ἱερουσαλήμ, καὶ ἀπέθανεν ἐκεῖ.

\vs{8}Καὶ ἐπολέμουν υἱοὶ Ἰούδα τὴν Ἱερουσαλὴμ, καὶ κατελάβοντο αὐτὴν, καὶ ἐπάταξαν αὐτὴν ἐν στόματι ῥομφαίας, καὶ τὴν πόλιν ἐνέπρησαν ἐν πυρί.
\vs{9}Καὶ μετὰ ταῦτα κατέβησαν οἱ υἱοὶ Ἰούδα πολεμῆσαι πρὸς τὸν Χαναναῖον τὸν κατοικοῦντα τὴν ὀρεινὴν καὶ τὸν Νότον καὶ τὴν πεδινήν.
\vs{10}Καὶ ἐπορεύθη Ἰούδας πρὸς τὸν Χαναναῖον τὸν κατοικοῦντα ἐν Χεβρών· καὶ ἐξῆλθε Χεβρὼν ἐξ ἐναντίας· καὶ τὸ ὄνομα ἦν Χεβρὼν τὸ πρότερον Καριαθαρβοκσεφέρ· καὶ ἐπάτάξαν τὸν Σεσσὶ καὶ Ἀχιμὰν καὶ Θολμὶ γεννήματα τοῦ Ἐνάκ.
\vs{11}Καὶ ἀνέβησαν ἐκεῖθεν πρὸς τοὺς κατοικοῦντας Δαβίρ· τὸ δὲ ὄνομα τῆς Δαβὶρ ἦν ἔμπροσθεν Καριαθσεφὲρ, πόλις Γραμμάτων.

\vs{12}Καὶ εἶπε Χάλεβ, ὃς ἂν πατάξῃ τὴν πόλιν τῶν Γραμμάτων καὶ προκαταλάβηται αὐτὴν, δώσω αὐτῷ τὴν Ἀσχὰ θυγατέρα μου εἰς γυναῖκα.
\vs{13}Καὶ προκατελάβετο αὐτὴν Γοθονιὴλ υἱὸς Κενὲζ ἀδελφοῦ Χάλεβ ὁ νεώτερος· καὶ ἔδωκεν αὐτῷ Χάλεβ τὴν Ἀσχὰ θυγατέρα αὐτοῦ εἰς γυναῖκα.
\vs{14}Καὶ ἐγένετο ἐν τῇ εἰσόδῳ αὐτῆς, καὶ ἐπέσεισεν αὐτὴν Γοθονιὴλ τοῦ αἰτῆσαι παρὰ τοῦ πατρὸς αὐτῆς ἀγρόν· καὶ ἐγόγγυζε καὶ ἔκραζεν ἀπὸ τοῦ ὑποζυγίου, εἰς γῆν Νότου ἐκδέδοσαί με· καὶ εἶπεν αὐτῇ Χάλεβ, τί ἐστί σοι;
\vs{15}Καὶ εἶπεν αὐτῷ Ἀσχὰ, δὸς δή μοι εὐλογίαν, ὅτι εἰς γῆν Νότου ἐκδέδοσαί με, καὶ δώσεις μοι λύτρωσιν ὕδατος· καὶ ἔδωκεν αὐτῇ Χάλεβ κατὰ τὴν καρδίαν αὐτῆς λύτρωσιν μετεώρων καὶ λύτρωσιν ταπεινῶν.

\vs{16}Καὶ οἱ υἱοὶ Ἰοθὸρ τοῦ Κιναίου τοῦ γαμβροῦ Μωυσῆ ἀνέβησαν ἐκ πόλεως τῶν φοινίκων μετὰ τῶν υἱῶν Ἰούδα εἰς τὴν ἔρημον τὴν οὖσαν ἐν τῷ Νότῳ Ἰούδα, ἥ ἐστιν ἐπὶ καταβάσεως Ἀρὰδ, καὶ κατῴκησαν μετὰ τοῦ λαοῦ.

\vs{17}Καὶ ἐπορεύθη Ἰούδας μετὰ Συμεὼν τοῦ ἀδελφοῦ αὐτοῦ, καὶ ἔκοψε τὸν Χαναναῖον τὸν κατοικοῦντα Σεφὲθ, καὶ ἐξωλόθρευσαν αὐτους· καὶ ἐκάλεσε τὸ ὄνομα τῆς πόλεως, Ἀνάθεμα.
\vs{18}Καὶ οὐκ ἐκληρονόμησεν Ἰούδας τὴν Γάζαν οὐδὲ τὰ ὅρια αὐτῆς, οὐδὲ τὴν Ἀσκάλωνα οὐδὲ τὰ ὅρια αὐτῆς, καὶ τὴν Ἀκκαρὼν οὐδὲ τὰ ὅρια αὐτῆς, τὴν Ἄζωτον οὐδὲ τὰ περισπόρια αὐτῆς.
\vs{19}Καὶ ἦν Κύριος μετὰ Ἰούδα· καὶ ἐκληρονόμησε τὸ ὄρος, ὅτι οὐκ ἠδυνάσθησαν ἐξολοθρεῦσαι τοὺς κατοικοῦντας τὴν κοιλάδα, ὅτι Ῥηχὰβ διεστείλατο αὐτοῖς.
\vs{20}Καὶ ἔδωκαν τῷ Χάλεβ τὴν Χεβρὼν, καθὼς ἐλάλησε Μωυσῆς· καὶ ἐκληρονόμησεν ἐκεῖθεν τὰς τρεῖς πόλεις τῶν υἱῶν Ἐνάκ.

\vs{21}Καὶ τὸν Ἰεβουσαῖον τὸν κατοικοῦντα ἐν Ἱερουσαλὴμ οὐκ ἐκληρονόμησαν οἱ υἱοὶ Βενιαμίν· καὶ κατῴκησεν ὁ Ἰεβουσαῖος μετὰ τῶν υἱῶν Βενιαμὶν ἐν Ἱερουσαλὴμ ἕως τῆς ἡμέρας ταύτης.

\vs{22}Καὶ ἀνέβησαν οἱ υἱοὶ Ἰωσὴφ καί γε αὐτοὶ εἰς Βαιθήλ· καὶ Κύριος ἦν μετʼ αὐτῶν.
\vs{23}Καὶ παρενέβαλον, καὶ κατεσκέψαντο Βαιθήλ· τὸ δὲ ὄνομα τῆς πόλεως ἦν ἔμπροσθεν Λουζά.

\vs{24}Καὶ εἶδον οἱ φυλάσσοντες, καὶ ἰδοὺ ἀνὴρ ἐξεπορεύετο ἐκ τῆς πόλεως, καὶ ἔλαβον αὐτόν· καὶ εἶπον αὐτῷ, δεῖξον ἡμῖν τῆς πόλεως τὴν εἴσοδον, καὶ ποιήσομεν μετὰ σου ἔλεος.
\vs{25}Καὶ ἔδειξεν αὐτοῖς τὴν εἴσοδον τῆς πόλεως· καὶ ἐπάταξαν τὴν πόλιν ἐν στόματι ῥομφαίας· τὸν δὲ ἄνδρα καὶ τὴν συγγένειαν αὐτοῦ ἐξαπέστειλαν.
\vs{26}Καὶ ἐπορεύθη ὁ ἀνὴρ εἰς γῆν Χεττίν· καὶ ᾠκοδόμησεν ἐκεῖ πόλιν, καὶ ἐκάλεσε τὸ ὄνομα αὐτῆς Λουζά· τοῦτο ὄνομα αὐτῆς ἕως τῆς ἡμέρας ταύτης.

\vs{27}Καὶ οὐκ ἐξῇρε Μανασσῆ τὴν Βαιθσὰν, ἥ ἐστι Σκυθῶν πόλις, οὐδὲ τὰς θυγατέρας αὐτῆς οὐδὲ τὰ περίοικα αὐτῆς, οὐδὲ τὴν Θανὰκ οὐδὲ τὰς θυγατέρας αὐτῆς, οὐδὲ τοὺς κατοικοῦντας Δὼρ οὐδὲ τὰς θυγατέρας αὐτῆς, οὐδὲ τὸν κατοικοῦντα Βαλὰκ οὐδὲ τὰ περίοικα αὐτῆς οὐδὲ τὰς θυγατέρας αὐτῆς, οὐδὲ τοὺς κατοικοῦντας Μαγεδὼ οὐδὲ τὰ περίοικα αὐτῆς καὶ τὰς θυγατέρας αὐτῆς, οὐδὲ τοὺς κατοικοῦντας Ἰεβλαὰμ οὐδὲ τὰ περίοικα αὐτῆς, οὐδὲ τὰς θυγατέρας αὐτῆς· καὶ ἤρξατο ὁ Χαναναῖος κατοικεῖν ἐν τῇ γῇ ταύτῃ.
\vs{28}Καὶ ἐγένετο ὅτε ἐνίσχυσεν Ἰσραὴλ, καὶ ἐποίησε τὸν Χαναναῖον εἰς φόρον, καὶ ἐξαίρων οὐκ ἐξῇρεν αὐτόν.
\vs{29}Καὶ Ἐφραὶμ οὐκ ἐξῇρε τὸν Χαναναῖον τὸν κατοικοῦντα ἐν Γαζέρ· καὶ κατῴκησεν ὁ Χαναναῖος ἐν μέσῳ αὐτοῦ ἐν Γαζὲρ, καὶ ἐγένετο εἰς φόρον.
\vs{30}Καὶ Ζαβουλὼν οὐκ ἐξῇρε τοὺς κατοικοῦντας Κέδρων, οὐδὲ τοὺς κατοικοῦντας Δωμανά· καὶ κατῴκησεν ὁ Χαναναῖος ἐν μέσῳ αὐτῶν, καὶ ἐγένετο αὐτῷ εἰς φόρον. Καὶ Ἀσὴρ οὐκ ἐξῇρε τοὺς κατοικοῦντας Δωμανά· καὶ κατῴκησεν ὁ Χαναναῖος ἐν μέσῳ αὐτῶν, καὶ ἐγένετο αὐτῷ εἰς φόρον.
\vs{31}Καὶ Ἀσὴρ οὐκ ἐξῇρε τοὺς κατοικοῦντας Ἀκχὼ, καὶ ἐγένετο αὐτῷ εἰς φόρον, καὶ τοὺς κατοικοῦντας Δὼρ, καὶ τοὺς κατοικοῦντας Σιδῶνα, καὶ τοὺς κατοικοῦντας Δαλὰφ, τὸν Ἀσχαζὶ, καὶ τὸν Χεβδὰ, καὶ τὸν Ναῒ, καὶ τὸν Ἐρεώ.
\vs{32}Καὶ κατῴκησεν ὁ Ἀσὴρ ἐν μέσῳ τοῦ Χαναναίου τοῦ κατοικοῦντος τὴν γῆν, ὅτι οὐκ ἠδυνήθη ἐξάραι αὐτόν.
\vs{33}Καὶ Νεφθαλὶ οὐκ ἐξῇρε τοὺς κατοικοῦντας Βαιθσαμῦς, καὶ τοὺς κατοικοῦντας Βαιθανάχ· καὶ κατῴκησε Νεφθαλὶ ἐν μέσῳ τοῦ Χαναναίου τοῦ κατοικοῦντος τὴν γῆν· οἱ δὲ κατοικοῦντες Βαιθσαμὺς καὶ τὴν Βαιθενὲθ, ἐγένοντο αὐτοῖς εἰς φόρον.

\vs{34}Καὶ ἐξέθλιψεν ὁ Ἀμοῤῥαῖος τοὺς υἱοὺς Δὰν εἰς τὸ ὄρος, ὅτι οὐκ ἀφῆκαν αὐτὸν καταβῆναι εἰς τὴν κοιλάδα.
\vs{35}Καὶ ἤρξατο ὁ Ἀμοῤῥαῖος κατοικεῖν ἐν τῷ ὄρει τῷ ὀστρακώδει, ἐν ᾧ αἱ ἅρκτοι καὶ ἐν ᾧ αἱ ἀλώπεκες, ἐν τῷ Μυρσινῶνι, καὶ ἐν Θαλαβίν, καὶ ἐβαρύνθη ἡ χεὶρ οἴκου Ἰωσὴφ ἐπὶ τὸν Ἀμοῤῥαῖον, καὶ ἐγενήθη αὐτοῖς εἰς φόρον.
\vs{36}Καὶ τὸ ὅριον τοῦ Ἀμοῤῥαίου ἀπὸ τῆς ἀναβάσεως Ἀκραβὶν ἀπὸ τῆς πέτρας καὶ ἐπάνω.

\ch{2}
Καὶ ἀνέβη ἄγγελος Κυρίου ἀπὸ Γαλγὰλ ἐπὶ τὸν κλαυθμῶνα καὶ ἐπὶ Βαιθὴλ καὶ ἐπὶ τὸν οἶκον Ἰσραὴλ, καὶ εἶπε πρὸς αὐτοὺς, τάδε λέγει Κύριος, ἀνεβίβασα ὑμᾶς ἐξ Αἰγύπτου, καὶ εἰσήγαγον ὑμᾶς εἰς τὴν γῆν ἣν ὤμοσα τοῖς πατράσιν ὑμῶν· καὶ εἶπα, οὐ διασκεδάσω τὴν διαθήκην μου τὴν μεθʼ ὑμῶν εἰς τὸν αἰῶνα.
\vs{2}Καὶ ὑμεῖς οὐ διαθήσεσθε διαθήκην τοῖς ἐγκαθημένοις εἰς τὴν γῆν ταύτην, οὐδὲ τοῖς θεοῖς αὐτῶν προσκυνήσετε, ἀλλὰ τὰ γλυπτὰ αὐτῶν συντρίψετε, τὰ θυσιαστήρια αὐτῶν καθελεῖτε· καὶ οὐκ εἰσηκούσατε τῆς φωνῆς μου, ὅτι ταῦτα ἐποιήσατε.
\vs{3}Κᾀγὼ εἶπον, οὐ μὴ ἐξάρω αὐτοὺς ἐκ προσώπου ὑμῶν, καὶ ἔσονται ὑμῖν εἰς συνοχὰς, καὶ οἱ θεοὶ αὐτῶν ἔσονται ὑμῖν εἰς σκάνδαλον.
\vs{4}Καὶ ἐγένετο ὡς ἐλάλησεν ὁ ἄγγελος Κυρίου τοὺς λόγους τούτους πρὸς πάντας υἱοὺς Ἰσραὴλ, καὶ ἐπῇραν ὁ λαὸς τὴν φωνὴν αὐτῶν καὶ ἔκλαυσαν.
\vs{5}Καὶ ἐπωνόμασαν τὸ ὄνομα τοῦ τόπου ἐκείνου, Κλαυθμῶνες· καὶ ἐθυσίασαν ἐκεῖ τῷ Κυρίῳ.

\vs{6}Καὶ ἐξαπέστειλεν Ἰησοῦς τὸν λαὸν, καὶ ἦλθεν ἀνὴρ εἰς τὴν κληρονομίαν αὐτοῦ κατακληρονομῆσαι τὴν γῆν.
\vs{7}Καὶ ἐδούλευσεν ὁ λαὸς τῷ Κυρίῳ πάσας τὰς ἡμέρας Ἰησοῦ καὶ πάσας τὰς ἡμέρας τῶν πρεσβυτέρων, ὅσοι ἐμακροημέρευσαν μετὰ Ἰησοῦ, ὅσοι ἔγνωσαν πᾶν τὸ ἔργον Κυρίου τὸ μέγα ὅσα ἐποίησεν ἐν τῷ Ἰσραήλ.

\vs{8}Καὶ ἐτελεύτησεν Ἰησοῦς υἱὸς Ναυῆ δοῦλος Κυρίου, υἱὸς ἑκατὸν δέκα ἐτῶν.
\vs{9}Καὶ ἔθαψαν αὐτὸν ἐν ὁρίῳ τῆς κληρονομίας αὐτοῦ ἐν Θαμναθαρὲς, ἐν ὄρει Ἐφραὶμ ἀπὸ Βοῤῥᾶ τοῦ ὄρους Γαάς.
\vs{10}Καὶ πᾶσα ἡ γενεὰ ἐκείνη προσετέθησαν πρὸς τοὺς πατέρας αὐτῶν· καὶ ἀνέστη γενεὰ ἑτέρα μετʼ αὐτοὺς, οἳ οὐκ ἔγνωσαν τὸν Κύριον, καί γε τὸ ἔργον ὃ ἐποίησεν ἐν τῷ Ἰσραήλ.
\vs{11}Καὶ ἐποίησαν οἱ υἱοὶ Ἰσραὴλ τὸ πονηρὸν ἐνώπιον Κυρίου, καὶ ἐλάτρευσαν τοῖς Βααλίμ.
\vs{12}Καὶ ἐγκατέλιπον τὸν Κύριον τὸν Θεὸν τῶν πατέρων αὐτῶν, τὸν ἐξαγαγόντα αὐτοὺς ἐκ γῆς Αἰγύπτου, καὶ ἐπορεύθησαν ὀπίσω θεῶν ἑτέρων ἀπὸ τῶν θεῶν τῶν ἐθνῶν τῶν περικύκλῳ αὐτῶν, καὶ προσεκύνησαν αὐτοῖς· καὶ παρώργισαν τὸν Κύριον,
\vs{13}καὶ ἐγκατέλιπον αὐτὸν, καὶ ἐλάτρευσαν τῷ Βάαλ καὶ ταῖς Ἀστάρταις.

\vs{14}Καὶ ὠργίσθη θυμῷ Κύριος ἐν τῷ Ἰσραήλ· καὶ παρέδωκεν αὐτοὺς εἰς χεῖρας προνομευόντων, καὶ κατεπρονόμευσαν αὐτούς· καὶ ἀπέδοτο αὐτοὺς ἐν χερσὶ τῶν ἐχθρῶν αὐτῶν κυκλόθεν, καὶ οὐκ ἠδυνήθησαν ἔτι ἀντιστῆναι κατὰ πρόσωπον τῶν ἐχθρῶν αὐτῶν ἐν πᾶσιν οἷς ἐπορεύοντο·
\vs{15}καὶ χεὶρ Κυρίου ἦν ἐπʼ αὐτοὺς εἰς κακὰ, καθὼς ἐλάλησε Κύριος καὶ καθὼς ὤμοσε Κύριος αὐτοῖς, καὶ ἐξέθλιψεν αὐτοὺς σφόδρα.

\vs{16}Καὶ ἤγειρε Κύριος κριτὰς, καὶ ἔσωσεν αὐτοὺς Κύριος ἐκ χειρὸς τῶν προνομευόντων αὐτούς· καί γε τῶν κριτῶν οὐχ ὑπήκουσαν,
\vs{17}ὅτι ἐξεπόρνευσαν ὀπίσω θεῶν ἑτέρων, καὶ προσεκύνησαν αὐτοῖς· καὶ ἐξέκλιναν ταχὺ ἐκ τῆς ὁδοῦ, ἧς ἐπορεύθησαν οἱ πατέρες αὐτῶν τοῦ εἰσακούειν τῶν λόγων Κυρίου· οὐκ ἐποίησαν οὕτω.
\vs{18}Καὶ ὅτι ἤγειρε Κύριος αὐτοῖς κριτὰς, καὶ ἦν Κύριος μετὰ τοῦ κριτοῦ, καὶ ἔσωσεν αὐτοὺς ἐκ χειρὸς ἐχθρῶν αὐτῶν πάσας τὰς ἡμέρας τοῦ κριτοῦ, ὅτι παρεκλήθη Κύριος ἀπὸ τοῦ στεναγμοῦ αὐτῶν ἀπὸ προσώπου τῶν πολιορκούντων αὐτοὺς καὶ ἐκθλιβόντων αὐτούς.
\vs{19}Καὶ ἐγένετο ὡς ἀπέθνησκεν ὁ κριτὴς, καὶ ἀπέστρεψαν καὶ πάλιν διέφθειραν ὑπὲρ τοὺς πατέρας αὐτῶν πορεύεσθαι ὀπίσω θεῶν ἑτέρων, λατρεύειν αὐτοῖς καὶ προσκυνεῖν αὐτοῖς· οὐκ ἀπέῤῥιψαν τὰ ἐπιτηδεύματα αὐτῶν, καὶ τὰς ὁδοὺς αὐτῶν τὰς σκληράς.

\vs{20}Καὶ ὠργίσθη θυμῷ Κύριος ἐν τῷ Ἰσραήλ· καὶ εἶπεν, ἀνθʼ ὧν ὅσα ἐγκατέλιπον τὸ ἔθνος τοῦτο τὴν διαθήκην μου ἣν ἐνετειλάμην τοῖς πατράσιν αὐτῶν, καὶ οὐκ εἰσήκουσαν τῆς φωνῆς μου,
\vs{21}καὶ ἐγὼ οὐ προσθήσω τοῦ ἐξᾶραι ἄνδρα ἐκ προσώπου αὐτῶν ἀπὸ τῶν ἐθνῶν, ὧν κατέλιπεν Ἰησοῦς υἱὸς Ναυῆ ἐν τῇ γῇ· καὶ ἀφῆκε
\vs{22}τοῦ πειρᾶσαι ἐν αὐτοῖς τὸν Ἰσραὴλ, εἰ φυλάσσονται τὴν ὁδὸν Κυρίου πορεύεσθαι ἐν αὐτῇ, ὃν τρόπον ἐφύλαξαν οἱ πατέρες αὐτῶν, ἢ οὔ.
\vs{23}Καὶ ἀφήσει Κύριος τὰ ἔθνη ταῦτα τοῦ μὴ ἐξᾶραι αὐτὰ τὸ τάχος, καὶ οὐ παρέδωκεν αὐτὰ ἐν χειρὶ Ἰησοῦ.

\ch{3}
Καὶ ταῦτα τὰ ἔθνη ἃ ἀφῆκε Κύριος αὐτὰ ὥστε πειρᾶσαι ἐν αὐτοῖς τὸν Ἰσραὴλ, πάντας τοὺς μὴ ἐγνωκότας τοὺς πολέμους Χαναάν.
\vs{2}Πλὴν διὰ τὰς γενεὰς υἱῶν Ἰσραὴλ τοῦ διδάξαι αὐτοὺς πόλεμον, πλὴν οἱ ἔμπροσθεν αὐτῶν οὐκ ἔγνωσαν αὐτά.
\vs{3}Τὰς πέντε σατραπείας τῶν ἀλλοφύλων, καὶ πάντα τὸν Χαναναῖον, καὶ τὸν Σιδώνιον, καὶ τὸν Εὐαῖον τὸν κατοικοῦντα τὸν Λίβανον ἀπὸ τοῦ ὄρους τοῦ Ἀερμὼν ἕως Λαβωεμάθ.
\vs{4}Καὶ ἐγένετο ὥστε πειρᾶσαι ἐν αὐτοῖς τὸν Ἰσραὴλ, γνῶναι εἰ ἀκούσονται τὰς ἐντολὰς Κυρίου, ἃς ἐνετείλατο τοῖς πατράσιν αὐτῶν ἐν χειρὶ Μωυσῆ.

\vs{5}Καὶ οἱ υἱοὶ Ἰσραὴλ κατῴκησαν ἐν μέσῳ τοῦ Χαναναίου, καὶ τοῦ Χετταίου, καὶ τοῦ Ἀμοῤῥαίου, καὶ τοῦ Φερεζαίου, καὶ τοῦ Εὐαίου, καὶ τοῦ Ἰεβουσαίου.
\vs{6}Καὶ ἔλαβον τὰς θυγατέρας αὐτῶν ἑαυτοῖς εἰς γυναῖκας, καὶ τὰς θυγατέρας αὐτῶν ἔδωκαν τοῖς υἱοῖς αὐτῶν, καὶ ἐλάτρευσαν τοῖς θεοῖς αὐτῶν.
\vs{7}Καὶ ἐποίησαν οἱ υἱοὶ Ἰσραὴλ τὸ πονηρὸν ἐναντίον Κυρίου· καὶ ἐπελάθοντο Κυρίου τοῦ Θεοῦ αὐτῶν, καὶ ἐλάτρευσαν τοῖς Βααλὶμ καὶ τοῖς ἄλσεσι.
\vs{8}Καὶ ὠργίσθη θυμῷ Κύριος ἐν τῷ Ἰσραὴλ, καὶ ἀπέδοτο αὐτοὺς ἐν χειρὶ Χουσαρσαθαὶμ βασιλέως Συρίας ποταμῶν· καὶ ἐδούλευσαν οἱ υἱοὶ Ἰσραὴλ τῷ Χουσαρσαθαὶμ ἔτη ὀκτώ.

\vs{9}Καὶ ἐκέκραξαν οἱ υἱοὶ Ἰσραὴλ πρὸς Κύριον· καὶ ἤγειρε Κύριος σωτῆρα τῷ Ἰσραὴλ, καὶ ἔσωσεν αὐτοὺς, τὸν Γοθονιὴλ υἱὸν Κενὲζ ἀδελφοῦ Χάλεβ τὸν νεώτερον ὑπὲρ αὐτόν.
\vs{10}Καὶ ἐγένετο ἐπʼ αὐτὸν πνεῦμα Κυρίου, καὶ ἔκρινε τὸν Ἰσραήλ· καὶ ἐξῆλθεν εἰς πόλεμον πρὸς Χουσαρσαθαίμ· καὶ παρέδωκε Κύριος ἐν χειρὶ αὐτοῦ τὸν Χουσαρσαθαὶμ βασιλέα Συρίας ποταμῶν· καὶ ἐκραταιώθη χεὶρ αὐτοῦ ἐπὶ τὸν Χουσαρσαθαίμ.
\vs{11}Καὶ ἡσύχασεν ἡ γῆ ἔτη τεσσαράκοντα· καὶ ἀπέθανε Γοθονιὴλ υἱὸς Κενέζ.

\vs{12}Καὶ προσέθεντο οἱ υἱοὶ Ἰσραὴλ ποιῆσαι τὸ πονηρὸν ἐνώπιον Κυρίου· καὶ ἐνίσχυσε Κύριος τὸν Ἐγλὼμ βασιλέα Μωὰβ ἐπὶ τὸν Ἰσραὴλ, διὰ τὸ πεποιηκέναι αὐτοὺς τὸ πονηρὸν ἔναντι Κυρίου.
\vs{13}Καὶ συνήγαγε πρὸς ἑαυτὸν πάντας τοὺς υἱοὺς Ἀμμὼν καὶ Ἀμαλὴκ, καὶ ἐπορεύθη καὶ ἐπάταξε τὸν Ἰσραὴλ, καὶ ἐκληρονόμησε τὴν πόλιν τῶν φοινίκων.
\vs{14}Καὶ ἐδούλευσαν οἱ υἱοὶ Ἰσραὴλ τῷ Ἐγλὼμ βασιλεῖ Μωὰβ ἔτη δεκαοκτώ.

\vs{15}Καὶ ἐκέκραξαν οἱ υἱοὶ Ἰσραὴλ πρὸς Κύριον· καὶ ἤγειρεν αὐτοῖς σωτῆρα, τὸν Ἀὼδ υἱὸν Γηρὰ υἱὸν τοῦ Ἰεμενὶ ἄνδρα ἀμφοτεροδέξιον· καὶ ἐξαπέστειλαν οἱ υἱοὶ Ἰσραὴλ δῶρα ἐν χειρὶ αὐτοῦ τῷ Ἐγλὼμ βασιλεῖ Μωάβ.
\vs{16}Καὶ ἐποίησεν ἑαυτῷ Ἀὼδ μάχαιραν δίστομον, σπιθαμῆς τὸ μῆκος αὐτῆς· καὶ περιεζώσατο αὐτὴν ὑπὸ τὸν μανδύαν ἐπὶ τὸν μηρὸν αὐτοῦ τὸν δεξιόν.
\vs{17}Καὶ ἐπορεύθη, καὶ προσήνεγκε τὰ δῶρα τῷ Ἐγλὼμ βασιλεῖ Μωάβ· καὶ Ἐγλὼμ ἀνὴρ ἀστεῖος σφόδρα.

\vs{18}Καὶ ἐγένετο ἡνίκα συνετέλεσεν Ἀὼδ προσφέρων τὰ δῶρα, καὶ ἐξαπέστειλε τοὺς φέροντας τὰ δῶρα,
\vs{19}καὶ αὐτὸς ὑπέστρεψεν ἀπὸ τῶν γλυπτῶν τῶν μετὰ τῆς Γαλγάλ· καὶ εἶπεν Ἀὼδ, λόγος μοι κρύφιος πρὸς σὲ, βασιλεῦ· καὶ εἶπεν Ἐγλὼμ πρὸς αὐτὸν, σιώπα· καὶ ἐξαπέστειλεν ἀφʼ ἑαυτοῦ πάντας τοὺς ἐφεστῶτας ἐπʼ αὐτόν,
\vs{20}καὶ Ἀὼδ εἰσῆλθε πρὸς αὐτόν· καὶ αὐτὸς ἐκάθητο ἐν τῷ ὑπερῴῳ τῷ θερινῷ τῷ ἑαυτοῦ μονώτατος· καὶ εἶπεν Ἀὼδ, λόγος Θεοῦ μοι πρὸς σὲ, βασιλεῦ· καὶ ἐξανέστη ἀπὸ τοῦ θρόνου Ἐγλὼμ ἐγγὺς αὐτοῦ.
\vs{21}Καὶ ἐγένετο ἅμα τῷ ἀναστῆναι αὐτὸν, καὶ ἐξέτεινεν Ἀὼδ τὴν χεῖρα τὴν ἀριστερὰν αὐτοῦ, καὶ ἔλαβε τὴν μάχαιραν ἐπάνωθεν τοῦ μηροῦ αὐτοῦ τοῦ δεξιοῦ, καὶ ἐνέπηξεν αὐτὴν ἐν τῇ κοιλίᾳ αὐτοῦ,
\vs{22}καὶ ἐπεισήνεγκε καί γε τὴν λαβὴν ὀπίσω τῆς φλογὸς, καὶ ἀπέκλεισε τὸ στέαρ κατὰ τῆς φλογὸς, ὅτι οὐκ ἐξέσπασε τὴν μάχαιραν ἐκ τῆς κοιλίας αὐτοῦ.

\vs{23}Καὶ ἐξῆλθεν Ἀὼδ τὴν προστάδα· καὶ ἐξῆλθε τοὺς διατεταγμένους, καὶ ἀπέκλεισε τὰς θύρας τοῦ ὑπερῴου κατʼ αὐτοῦ, καὶ ἐσφήνωσε.
\vs{24}Καὶ αὐτὸς ἐξῆλθε· καὶ οἱ παῖδες αὐτοῦ ἐπῆλθον καὶ εἶδον, καὶ ἰδοὺ αἱ θύραι τοῦ ὑπερῴου ἐσφηνωμέναι· καὶ εἶπαν, μήποτε ἀποκενοῖ τοὺς πόδας αὐτοῦ ἐν τῷ ταμείῳ τῷ θερινῷ;
\vs{25}Καὶ ὑπέμειναν ἕως ᾑσχύνοντο· καὶ ἰδοὺ οὐκ ἔστιν ὁ ἀνοίγων τὰς θύρας τοῦ ὑπερῴου· καὶ ἔλαβον τὴν κλεῖδα, καὶ ἤνοιξαν· καὶ ἰδοὺ ὁ κύριος αὐτῶν πεπτωκὼς ἐπὶ τὴν γῆν τεθνηκώς.

\vs{26}Καὶ Ἀὼδ διεσώθη ἕως ἐθορυβοῦντο, καὶ οὐκ ἦν ὁ προσνοῶν αὐτῷ· καὶ αὐτὸς παρῆλθε τὰ γλυπτὰ, καὶ διεσώθη εἰς Σετειρωθά.

\vs{27}Καὶ ἐγένετο ἡνίκα ἦλθεν Ἀὼδ εἰς γῆν Ἰσραὴλ, καὶ ἐσάλπισεν ἐν κερατίνῃ ἐν τῷ ὄρει Ἐφραὶμ, καὶ κατέβησαν σὺν αὐτῷ οἱ υἱοὶ Ἰσραὴλ ἀπὸ τοῦ ὄρους, καὶ αὐτὸς ἔμπροσθεν αὐτῶν.
\vs{28}Καὶ εἶπε πρὸς αὐτοὺς, κατάβητε ὀπίσω μου, ὅτι παρέδωκε Κύριος ὁ Θεὸς τοὺς ἐχθροὺς ἡμῶν τὴν Μωὰβ ἐν χειρὶ ἡμῶν· καὶ κατέβησαν ὀπίσω αὐτοῦ, καὶ προκατελάβοντο τὰς διαβάσεις τοῦ Ἰορδάνου τῆς Μωὰβ, καὶ οὐκ ἀφῆκεν ἄνδρα διαβῆναι.
\vs{29}Καὶ ἐπάταξαν τὴν Μωὰβ τῇ ἡμέρᾳ ἐκείνῃ ὡσεὶ δέκα χιλιάδας ἀνδρῶν, πᾶν λιπαρὸν καὶ πάντα ἄνδρα δυνάμεως, καὶ οὐ διεσώθη ὁ ἀνήρ.
\vs{30}Καὶ ἐνετράπη Μωὰβ ἐν τῇ ἡμέρᾳ ἐκείνῃ ὑπὸ χεῖρα Ἰσραὴλ, καὶ ἡσύχασεν ἡ γῆ ὀγδοήκοντα ἔτη· καὶ ἔκρινεν αὐτοὺς Ἀὼδ ἕως οὗ ἀπέθανε.

\vs{31}Καὶ μετʼ αὐτὸν ἀνέστη Σαμεγὰρ υἱὸς Δινὰχ, καὶ ἐπάταξε τοὺς ἀλλοφύλους εἰς ἑξακοσίους ἄνδρας ἐν τῷ ἀροτρόποδι τῶν βοῶν· καὶ ἔσωσε καί γε αὐτὸς τὸν Ἰσραήλ.

\ch{4}
Καὶ προσέθεντο οἱ υἱοὶ Ἰσραὴλ ποιῆσαι τὸ πονηρὸν ἐνώπιον Κυρίου· καὶ Ἀὼδ ἀπέθανε.
\vs{2}Καὶ ἀπέδοτο τοὺς υἱοὺς Ἰσραὴλ Κύριος ἐν χειρὶ Ἰαβὶν βασιλέως Χαναὰν, ὃς ἐβασίλευσεν ἐν Ἀσώρ· καὶ ὁ ἂρχων τῆς δυνάμεως αὐτοῦ Σισάρα, καὶ αὐτὸς κατῴκει ἐν Ἀρισὼθ τῶν ἐθνῶν.
\vs{3}Καὶ ἐκέκραξαν οἱ υἱοὶ Ἰσραὴλ πρὸς Κύριον, ὅτι ἐννακόσια ἅρματα σιδηρᾶ ἦν αὐτῷ· καὶ αὐτὸς ἔθλιψε τὸν Ἰσραὴλ κατακράτος εἴκοσι ἔτη.

\vs{4}Καὶ Δεββῶρα γυνὴ προφῆτις, γυνὴ Λαφιδὼθ, αὕτη ἔκρινε τὸν Ἰσραὴλ ἐν τῷ καιρῷ ἐκείνῳ.

\vs{5}Καὶ αὐτὴ ἐκάθητο ὑπὸ φοίνικα Δεββῶρα ἀναμέσον τῆς Ῥαμὰ καὶ ἀναμέσον τῆς Βαιθὴλ ἐν τῷ ὄρει Ἐφραίμ· καὶ ἀνέβαινον πρὸς αὐτὴν οἱ υἱοὶ Ἰσραὴλ εἰς κρίσιν.

\vs{6}Καὶ ἀπέστειλε Δεββῶρα καὶ ἐκάλεσε τὸν Βαρὰκ υἱὸν Ἀβινεὲμ ἐκ Κάδης Νεφθαλὶ, καὶ εἶπε πρὸς αὐτὸν, οὐχὶ ἐνετείλατο Κύριος ὁ Θεὸς Ἰσραήλ σοι; καὶ ἀπελεύσῃ εἰς ὄρος Θαβὼρ, καὶ λήψῃ μετὰ σεαυτοῦ δέκα χιλιάδας ἀνδρῶν ἐκ τῶν υἱῶν Νεφθαλὶ, καὶ ἐκ τῶν υἱῶν Ζαβουλὼν.

\vs{7}Καὶ ἐπάξω πρὸς σὲ εἰς τὸν χειμάῤῥουν Κισῶν ἐπὶ τὸν Σισάρα ἄρχοντα τῆς δυνάμεως Ἰαβὶν, καὶ τὰ ἅρματα αὐτοῦ καὶ τὸ πλῆθος αὐτοῦ, καὶ παραδώσω αὐτὸν εἰς χεῖράς σου.

\vs{8}Καὶ εἶπε πρὸς αὐτὴν Βαρὰκ, ἐὰν πορευθῇς μετʼ ἐμοῦ, πορεύσομαι, καὶ ἐὰν μὴ πορευθῇς, οὐ πορεύσομαι· ὅτι οὐκ οἶδα τὴν ἡμέραν ἐν ᾗ εὐοδοῖ Κύριος τὸν ἄγγελον μετʼ ἐμοῦ.
\vs{9}Καὶ εἶπε, πορευομένη πορεύσομαι μετὰ σοῦ· πλὴν γίνωσκε ὅτι οὐκ ἔσται τὸ προτέρημά σου ἐπὶ τὴν ὁδὸν ἣν σὺ πορεύῃ, ὅτι ἐν χειρὶ γυναικὸς ἀποδώσεται Κύριος τὸν Σισάρα· καὶ ἀνέστη Δεββῶρα, καὶ ἐπορεύθη μετὰ τοῦ Βαρὰκ ἐκ Κάδης.
\vs{10}Καὶ ἐβόησε Βαρὰκ τὸν Ζαβουλὼν καὶ τὸν Νεφθαλὶ ἐκ Κάδης, καὶ ἀνέβησαν κατὰ πόδας αὐτοῦ δέκα χιλιάδες ἀνδρῶν, καὶ ἀνέβη Δεββῶρα μετʼ αὐτοῦ.

\vs{11}Καὶ Χαβὲρ ὁ Κιναῖος ἐχωρίσθη ἀπὸ Καινᾶ ἀπὸ τῶν υἱῶν Ἰωβὰβ γαμβροῦ Μωυσῆ· καὶ ἔπηξε τὴν σκηνὴν αὐτοῦ ἕως δρυὸς πλεονεκτούντων, ἥ ἐστιν ἐχόμενα Κεδές.

\vs{12}Καὶ ἀνηγγέλη Σισάρᾳ, ὅτι ἀνέβη Βαρὰκ υἱὸς Ἀβινεὲμ εἰς ὄρος Θαβώρ.
\vs{13}Καὶ ἐκάλεσε Σισάρα πάντα τὰ ἅρματα αὐτοὺ ἐννακόσια ἅρματα σιδηρᾶ, καὶ πάντα τὸν λαὸν τὸν μετʼ αὐτοῦ ἀπὸ Ἀρισὼθ τῶν ἐθνῶν εἰς τὸν χειμάῤῥουν Κισῶν.

\vs{14}Καὶ εἶπε Δεββῶρα πρὸς Βαρὰκ, ἀνάστηθι, ὅτι αὕτη ἡ ἡμέρα ἐν ᾗ παρέδωκε Κύριος τὸν Σισάρα ἐν τῇ χειρί σου, ὅτι Κύριος ἐξελεύσεται ἔμπροσθέν σου· καὶ κατέβη Βαρὰκ κατὰ τοῦ ὄρους Θαβὼρ, καὶ δέκα χιλιάδες ἀνδρῶν ὀπίσω αὐτοῦ.
\vs{15}Καὶ ἐξέστησε Κύριος τὸν Σισάρα καὶ πάντα τὰ ἅρματα αὐτοῦ καὶ πᾶσαν τὴν παρεμβολὴν αὐτοῦ ἐν στόματι ῥομφαίας ἐνώπιον Βαράκ· καὶ κατέβη Σισάρα ἐπάνωθεν τοὺ ἄρματος αὐτοῦ, καὶ ἔφυγε τοῖς ποσὶν αὐτοῦ.
\vs{16}Καὶ Βαρὰκ διώκων ὀπίσω τῶν ἁρμάτων καὶ ὀπίσω τῆς παρεμβολῆς ἕως Ἀρισὼθ τῶν ἐθνῶν· καὶ ἔπεσε πᾶσα παρεμβολὴ Σισάρα ἐν στόματι ῥομφαίας· οὐ κατελείφθη ἕως ἑνός.
\vs{17}Καὶ Σισάρα ἔφυγε τοῖς ποσὶν αὐτοῦ εἰς σκηνὴν Ἰαὴλ γυναικὸς Χαβὲρ ἑταίρου τοῦ Κιναίου· ὅτι εἰρήνη ἦν ἀναμέσον Ἰαβὶν βασιλέως Ἀσὼρ καὶ ἀναμέσον τοῦ οἴκου Χαβὲρ τοῦ Κιναίου.
\vs{18}Καὶ ἐξῆλθεν Ἰαὴλ εἰς συνάντησιν Σισάρα, καὶ εἶπεν αὐτῷ, ἔκκλινον, κύριέ μου, ἔκκλινον πρὸς μὲ, μὴ φοβοῦ· καὶ ἐξέκλινε πρὸς αὐτῆν εἰς τὴν σκηνήν· καὶ περιέβαλεν αὐτὸν ἐπιβολαίῳ.

\vs{19}Καὶ εἶπε Σισάρα πρὸς αὐτὴν, πότισόν με δὴ μικρὸν ὑδωρ, ὅτι ἐδίψησα· καὶ ἤνοιξε τὸν ἀσκὸν τοῦ γάλακτος, καὶ ἐπότισεν αὐτὸν, καὶ περιέβαλεν αὐτόν.
\vs{20}Καὶ εἶπε πρὸς αὐτὴν Σισάρα, στῆθι δὴ ἐπὶ τὴν θύραν τῆς σκηνῆς, καὶ ἔσται ἐὰν ἀνὴρ ἔλθῃ πρὸς σὲ, καὶ ἐρωτήσῃ σε, καὶ εἴπῃ, εἰ ἔστιν ὧδε ἀνήρ; καὶ ἐρεῖς, οὐκ ἔστι.
\vs{21}Καὶ ἔλαβεν Ἰαὴλ γυνὴ Χαβὲρ τὸν πάσσαλον τῆς σκηνῆς, καὶ ἔθηκε τὴν σφύραν ἐν τῇ χειρὶ αὐτῆς, καὶ εἰσῆλθε πρὸς αὐτὸν ἐν κρύφῇ, καὶ ἔπηξε τὸν πάσσαλον ἐν τῷ κροτάφῳ αὐτοῦ, καὶ διεξῆλθεν ἐν τῇ γῇ· καὶ αὐτὸς ἐξεστῶς ἐσκοτώθη, καὶ ἀπέθανε.
\vs{22}Καὶ ἰδοὺ Βαρὰκ διώκων τὸν Σισάρα· καὶ ἐξῆλθεν Ἰαὴλ εἰς συνάντησιν αὐτῷ, καὶ εἶπεν αὐτῷ, δεῦρο καὶ δείξω σοι τὸν ἄνδρα ὃν σὺ ζητεῖς· καὶ εἰσῆλθε πρὸς αὐτήν· καὶ ἰδοὺ Σισάρα ἐῤῥιμμένος νεκρὸς, καὶ ὁ πάσσαλος ἐν τῷ κροτάφῳ αὐτοῦ.
\vs{23}Καὶ ἐτρόπωσεν ὁ Θεὸς τὸν Ἰαβὶν βασιλέα Χαναὰν ἐν τῇ ἡμέρᾳ ἐκείνῃ ἔμπροσθεν υἱῶν Ἰσραήλ.

\vs{24}Καὶ ἐπορεύετο χεὶρ τῶν υἱῶν Ἰσραὴλ πορευομένη καὶ σκληρυνομένη ἐπὶ Ἰαβὶν βασιλέα Χαναὰν, ἕως οὗ ἐξωλόθρευσαν τὸν Ἰαβὶν βασιλέα Χαναάν.

\ch{5}
Καὶ ἦσαν Δεββῶρα καὶ Βαρὰκ υἱὸς Ἀβινεὲμ ἐν τῇ ἡμέρᾳ ἐκείνῃ, λέγοντες,

\vs{2}Ἀπεκαλύφθη ἀποκάλυμμα ἐν Ἰσραὴλ ἐν τῷ ἑκουσιασθῆναι λαὸν, εὐλογεῖτε Κύριον.
\vs{3}Ἀκούσατε βασιλεῖς, καὶ ἐνωτίσασθε σατράπαι· ᾄσομαι ἐγώ εἰμι τῷ Κυρίῳ ἐγώ εἰμι, ψαλῶ τῷ Κυρίῳ τῷ Θεῷ Ἰσραήλ.
\vs{4}Κύριε, ἐν τῇ ἐξόδῳ σου ἐν Σηεὶρ, ἐν τῷ ἀπαίρειν σε ἐξ ἀγροῦ Ἐδὼμ, γῆ ἐσείσθη, καὶ ὁ οὐρανὸς ἔσταξε δρόσους, καὶ αἱ νεφέλαι ἔσταξαν ὕδωρ.
\vs{5}Ὄρη ἐσαλεύθησαν ἀπὸ προσώπου Κυρίου Ἐλωῒ, τοῦτο Σινὰ ἀπὸ προσώπου Κυρίου Θεοῦ Ἰσραήλ.
\vs{6}Ἐν ἡμέραις Σαμεγὰρ υἱοῦ Ἀνὰθ, ἐν ἡμέραις Ἰαὴλ, ἐξέλιπον ὁδοὺς, καὶ ἐπορεύθησαν ἀτραποὺς, ἐπορεύθησαν ὁδοὺς διεστραμμένας.
\vs{7}Ἐξέλιπον δυνατοὶ ἐν Ἰσραὴλ, ἐξέλιπον ἕως οὗ ἀνέστη Δεββῶρα, ἕως οὗ ἀνέστη μήτηρ ἐν Ἰσραήλ.
\vs{8}Ἐξελέξαντο θεοὺς καινοὺς, τότε ἐπολέμησαν πόλεις ἀρχόντων· θυρεὸς ἐὰν ὀφθῇ καὶ λόγχη ἐν τεσσαράκοντα χιλιάσιν ἐν Ἰσραήλ.

\vs{9}Ἡ καρδία μου εἰς τὰ διατεταγμένα τῷ Ἰσραήλ· οἱ ἑκούσιαζόμενοι ἐν λαῷ εὐλογεῖτε Κύριον.
\vs{10}Ἐπιβεβηκότες ἐπὶ ὄνου θηλείας μεσημβρίας, καθήμενοι ἐπὶ κριτηρίου, καὶ πορευόμενοι ἐπὶ ὁδοὺς συνέδρων ἐφʼ ὁδῷ,
\vs{11}διηγεῖσθε, ἀπὸ φωνῆς ἀνακρουομένων ἀναμέσον ὑδρευομένων· ἐκεῖ δώσουσι δικαιοσύνας· Κύριε δικαιοσύνας αὔξησον ἐν Ἰσραήλ· τότε κατέβη εἰς τὰς πόλεις λαὸς Κυρίου.
\vs{12}Ἐξεγείρου, ἐξεγείρου, Δεββῶρα· ἐξεγείρου, ἐξεγείρου, λάλησον ᾠδήν· ἀνάστα Βαρὰκ, καὶ αἰχμαλώτισον αἰχμάλωσίαν σου υἱὸς Ἀβινεέμ.
\vs{13}Τότε κατέβη κατάλειμμα τοῖς ἰσχυροῖς· λαὸς Κυρίου κατέβη αὐτῷ ἐν τοῖς κραταιοῖς ἐξ ἐμοῦ.

\vs{14}Ἐφραῒμ ἐξεῤῥίζωσεν αὐτοὺς ἐν τῷ Ἀμαλὴκ, ὀπίσω σου Βενιαμὶν ἐν τοῖς λαοῖς σου· ἐν ἐμοὶ Μαχὶρ κατέβησαν ἐξερευνῶντες· καὶ ἀπὸ Ζαβουλὼν ἕλκοντες ἐν ῥάβδῳ διηγήσεως γραμματέως.
\vs{15}Καὶ ἀρχηγοι ἐν Ἰσσάχαρ μετὰ Δεββώρας καὶ Βαράκ· οὕτω Βαρὰκ ἐν κοιλάσιν ἀπέστειλεν ἐν ποσὶν αὐτοῦ, εἰς τὰς μερίδας Ῥουβὴν, μεγάλοι ἐξικνούμενοι καρδίαν.
\vs{16}Εἰς τί ἐκάθισαν ἀναμέσον τῆς διγομίας τοῦ ἀκοῦσαι συρισμοῦ ἀγελῶν εἰς διαιρέσεις Ῥουβήν; μεγάλοι ἐξετασμοὶ καρδίας.
\vs{17}Γαλαάδ ἐν τῷ πέραν τοῦ Ἰορδάνου οὗ ἐσκήνωσε· καὶ Δὰν εἰς τί παροικεῖ πλοίοις; Ἀσὴρ ἐκάθισε παραλίαν θαλασσῶν, καὶ ἐπὶ διεξόδοις αὐτοῦ σκηνώσει.
\vs{18}Ζαβουλὼν λαὸς ὠνείδισε ψυχὴν αὐτοῦ εἰς θάνατον, καὶ Νεφθαλὶ ἐπὶ ὕψη ἀγροῦ

\vs{19}ἦλθον αὐτῶν. Βασιλεῖς παρετάξαντο, τότε ἐπολέμησαν βασιλεῖς Χαναὰν ἐν Θαναὰχ ἐπὶ ὕδατι Μαγεδδὼ, δῶρον ἀργυρίου οὐκ ἔλαβον.
\vs{20}Ἐξ οὐρανοῦ παρετάξαντο οἱ ἀστέρες, ἐκ τρίβων αὐτῶν παρετάξαντο μετὰ Σισάρα.
\vs{21}Χειμάῤῥους Κισῶν ἐξέσυρεν αὐτοὺς, χειμάῤῥους ἀρχαίων, χειμάῤῥους Κισῶν· καταπατήσει αὐτὸν ψυχή μου δυνατή.
\vs{22}Ὅτε ἐνεποδίσθησαν πτέρναι ἵππου, σπουδῇ ἔσπευσαν ἰσχυροὶ αὐτοῦ
\vs{23}καταρᾶσθαι Μηρὼζ, εἶπεν ἄγγελος Κυρίου, καταρᾶσθε· ἐπικατάρατος πᾶς ὁ κατοικῶν αὐτὴν, ὅτι οὐκ ἤλθοσαν εἰς βοήθειαν Κυρίου, εἰς βοήθειαν ἐν δυνατοῖς.

\vs{24}Εὐλογηθείη ἐν γυναιξὶν Ἰαὴλ γυνὴ Χαβὲρ τοῦ Κιναίου, ἀπὸ γυναικῶν ἐν σκηναῖς εὐλογηθείη.
\vs{25}Ὕδωρ ᾔτησε, γάλα ἔδωκεν ἐν λεκάνῃ· ὑπερεχόντων προσήνεγκε βούτυρον.
\vs{26}Χεῖρα αὐτῆς ἀριστερὰν εἰς πάσσαλον ἐξέτεινε, καὶ δεξιὰν αὐτῆς εἰς σφύραν κοπίωντων, καὶ ἐσφυροκόπησε Σισάρα, διήλωσε κεφαλὴν αὐτοῦ καὶ ἐπάταξε, διήλωσε κρόταφον αὐτοῦ.
\vs{27}Ἀναμέσον τῶν ποδῶν αὐτῆς κατεκυλίσθη· ἔπεσε καὶ ἐκοιμήθη ἀναμέσον τῶν ποδῶν αὐτῆς, κατακλιθεὶς ἔπεσε· καθὼς κατεκλίθη ἐκεῖ ἔπεσεν ἐξοδευθείς.

\vs{28}Διὰ τῆς θυρίδος παρέκυψε μήτηρ Σισάρα ἐκτὸς τοῦ τοξικοῦ, διότι ᾐσχύνθη ἅρμα αὐτοῦ; διότι ἐχρόνισαν πόδες ἁρμάτων αὐτοῦ;
\vs{29}Αἱ σοφαὶ ἄρχουσαι αὐτῆς ἀπεκρίθησαν πρὸς αὐτὴν, καὶ αὐτὴ ἀπέστρεψε λόγους αὐτῆς ἑαυτῇ,
\vs{30}οὐχ εὑρήσουσιν αὐτὸν διαμερίζοντα σκῦλα; οἰκτίρμων οἰκτειρήσει εἰς κεφαλὴν ἀνδρός· σκῦλα βαμμάτων τῷ Σισάρᾳ, σκῦλα βαμμάτων ποικιλίας, βάμματα ποικιλτῶν αὐτὰ τῷ τραχήλῳ αὐτοῦ σκῦλα.
\vs{31}Οὕτως ἀπόλοιντο πάντες οἱ ἐχθροί σου, Κύριε· καὶ οἱ ἀγαπῶντες αὐτὸν, ὡς ἔξοδος ἡλίου ἐν δυνάμει αὐτοῦ.

Καὶ ἡσύχασεν ἡ γῆ τεσσαράκοντα ἔτη.

\ch{6}
Καὶ ἐποίησαν οἱ υἱοὶ Ἰσραὴλ τὸ πονηρὸν ἐνώπιον Κυρίου, καὶ ἔδωκεν αὐτοὺς Κύριος ἐν χειρὶ Μαδιὰμ ἑπτὰ ἔτη.
\vs{2}Καὶ ἴσχυσε χεὶρ Μαδιὰμ ἐπὶ Ἰσραήλ· καὶ ἐποίησαν ἑαυτοῖς οἱ υἱοὶ Ἰσραὴλ ἀπὸ προσώπου Μαδιὰμ τὰς τρυμαλιὰς τὰς ἐν τοῖς ὄρεσι, καὶ τὰ σπήλαια, καὶ τὰ κρεμαστά.
\vs{3}Καὶ ἐγένετο ἐὰν ἔσπειραν οἱ υἱοὶ Ἰσραὴλ, καὶ ἀνέβαινον Μαδιὰμ καὶ Ἀμαλὴκ, καὶ οἱ υἱοὶ ἀνατολῶν συνανέβαινον αὐτοῖς, καὶ παρενέβαλον εἰς αὐτοὺς,
\vs{4}καὶ διέφθειρον τοὺς καρποὺς αὐτῶν ἕως ἐλθεῖν εἰς Γάζαν· καὶ οὐ κατελείποντο ὑπόστασιν ζωῆς ἐν τῇ γῇ Ἰσραὴλ, οὐδὲ ἐν τοῖς ποιμνίοις ταῦρον καὶ ὄνον.
\vs{5}Ὅτι αὐτοὶ καὶ αἱ κτήσεις αὐτῶν ἀνέβαινον, καὶ αἱ σκηναὶ αὐτῶν παρεγίνοντο, καθὼς ἀκρὶς εἰς πλῆθος, καὶ αὐτοῖς καὶ ταῖς καμήλοις αὐτῶν οὐκ ἦν ἀριθμός· καὶ ἤρχοντο εἰς τὴν γῆν Ἰσραὴλ, καὶ διέφθειρον αὐτήν.
\vs{6}Καὶ ἐπτώχευσεν Ἰσραὴλ σφόδρα ἀπὸ προσώπου Μαδιάμ.
\vs{7}Καὶ ἐβόησαν οἱ υἱοὶ Ἰσραὴλ πρὸς Κύριον ἀπὸ προσώπου Μαδιάμ.

\vs{8}Καὶ ἐξαπέστειλε Κύριος ἄνδρα προφήτην πρὸς τοὺς νἱοὺς Ἰσραήλ· καὶ εἶπεν αὐτοῖς, τάδε λέγει Κύριος ὁ Θεὸς Ἰσραὴλ, ἐγώ εἰμι ὃς ἀνήγαγον ὑμᾶς ἐκ γῆς Αἰγύπτου, καὶ ἐξήγαγον ὑμᾶς ἐξ οἴκου δουλείας ὑμῶν·
\vs{9}καὶ ἐῤῥυσάμην ὑμᾶς ἐκ χειρὸς Αἰγύπτου καὶ ἐκ χειρὸς πάντων τῶν θλιβόντων ὑμᾶς, καὶ ἐξέβαλον αὐτοὺς ἐκ προσώπου ὑμῶν· καὶ ἔδωκα ὑμῖν τὴν γῆν αὐτῶν.
\vs{10}Καὶ εἶπα ὑμῖν, ἐγὼ Κύριος ὁ Θεὸς ὑμῶν· οὐ φοβηθήσεσθε τοὺς θεοὺς τοῦ Ἀμοῤῥαίου, ἐν οἷς ὑμεῖς κάθησθε ἐν τῇ γῇ αὐτῶν· καὶ οὐκ εἰσηκούσατε τῆς φωνῆς μου.

\vs{11}Καὶ ἦλθεν ἄγγελος Κυρίου, καὶ ἐκάθισεν ὑπὸ τὴν τερέμινθον τὴν ἐν Ἐφραθὰ ἐν γῇ Ἰωὰς πατρὸς τοῦ Ἐσδρί· καὶ Γεδεὼν ὁ υἱὸς αὐτοῦ ῥαβδίζων σῖτον ἐν ληνῷ εἰς ἐκφυγεῖν ἀπὸ προσώπου τοῦ Μαδίαμ.
\vs{12}Καὶ ὤφθη αὐτῷ ὁ ἄγγελος Κυρίου, καὶ εἶπε πρὸς αὐτὸν, Κύριος μετὰ σοῦ, ἰσχυρὸς τῶν δυνάμεων.
\vs{13}Καὶ εἶπε πρὸς αὐτὸν Γεδεὼν, ἐν ἐμοὶ, Κύριέ μου· καὶ εἰ ἔστι Κύριος μεθʼ ἡμῶν, εἰς τί εὗρεν ἡμᾶς τὰ κακὰ ταῦτα; καὶ ποῦ ἐστι πάντα τὰ θαυμάσια αὐτοῦ, ἃ διηγήσαντο ἡμῖν οἱ πατέρες ἡμῶν, λέγοντες, μὴ οὐχὶ ἐξ Αἰγύπτου ἀνήγαγεν ἡμᾶς Κύριος; καὶ νῦν ἐξέῤῥιψεν ἡμᾶς καὶ ἔδωκεν ἡμᾶς ἐν χειρὶ Μαδίαμ.
\vs{14}Καὶ ἐπέστρεψε πρὸς αὐτὸν ὁ ἄγγελος Κυρίου, καὶ εἶπε, πορεύου ἐν τῇ ἰσχύϊ σου ταύτῃ, καὶ σώσεις τὸν Ἰσραὴλ ἐκ χειρὸς Μαδιάμ· ἰδοὺ ἐξαπέστειλά σε.
\vs{15}Καὶ εἶπε πρὸς αὐτὸν Γεδεὼν, ἐν ἐμοὶ, Κύριέ μου, ἐν τίνι σώσω τὸν Ἰσραήλ; ἰδοὺ ἡ χιλιάς μου ἠσθένησεν ἐν Μανασσῇ, καὶ ἐγώ εἰμι μικρότερος ἐν οἴκῳ τοῦ πατρός μου.
\vs{16}Καὶ εἶπε πρὸς αὐτὸν ὁ ἄγγελος Κυρίου, Κύριος ἔσται μετὰ σοῦ, καὶ πατάξεις τὴν Μαδιὰμ ὡσεὶ ἄνδρα ἕνα.
\vs{17}Καὶ εἶπε πρὸς αὐτὸν Γεδεὼν, εἰ δὴ εὗρον ἔλεος ἐν ὀφθαλμοῖς σου, καὶ ποιήσεις μοι σήμερον πᾶν ὅτι, ἐλάλησας μετʼ ἐμοῦ,
\vs{18}μὴ χωρισθῇς ἐντεῦθεν ἕως τοῦ ἐλθεῖν με πρὸς σὲ, καὶ ἐξοίσω τὴν θυσίαν καὶ θύσω ἐνώπιόν σου· καὶ εἶπεν, ἐγώ εἰμι καθίσομαι ἕως τοῦ ἐπιστρέψαι σε.

\vs{19}Καὶ Γεδεὼν εἰσῆλθε, καὶ ἐποίησεν ἔριφον αἰγῶν καὶ οἰφὶ ἀλεύρου ἄζυμα, καὶ τὰ κρέα ἔθηκεν ἐν τῷ κοφίνῳ, καὶ τὸν ζωμὸν ἔβαλεν ἐν τῇ χύτρᾳ, καὶ ἐξήνεγκεν αὐτὰ πρὸς αὐτὸν ὑπὸ τὴν τερέμινθον, καὶ προσήγγισε.
\vs{20}Καὶ εἶπε πρὸς αὐτὸν ὁ ἄγγελος τοῦ Θεοῦ, λάβε τὰ κρέα καὶ τὰ ἄζυμα, καὶ θὲς πρὸς τὴν πέτραν ἐκείνην, καὶ τὸν ζωμὸν ἐχόμενα ἔκχεε· καὶ ἐποιήσεν οὕτως.
\vs{21}Καὶ ἐξέτεινεν ὁ ἄγγελος Κυρίου τὸ ἄκρον τῆς ῥάβδου τῆς ἐν τῇ χειρὶ αὐτοῦ, καὶ ἥψατο τῶν κρεῶν καὶ τὼν ἀζύμων· καὶ ἀνέβη πῦρ ἐκ τῆς πέτρας, καὶ κατέφαγε τὰ κρέα καὶ τοὺς ἀζύμους· καὶ ὁ ἄγγελος Κυρίου ἐπορεύθη ἀπʼ ὀφθαλμῶν αὐτοῦ.

\vs{22}Καὶ εἶδε Γεδεὼν, ὅτι ἄγγελος Κυρίου οὗτός ἐστι· καὶ εἶπε Γεδεὼν, ἆ ἆ, Κύριέ μου Κύριε, ὅτι εἶδον τὸν ἄγγελον Κυρίου πρόσωπον πρὸς πρόσωπον.
\vs{23}Καὶ εἶπεν αὐτῷ Κύριος, εἰρήνη σοι, μὴ φοβοῦ, οὐ μὴ ἀποθάνῃς.

\vs{24}Καὶ ᾠκοδόμησεν ἐκεῖ Γεδεὼν θυσιαστήριον τῷ Κυρίῳ, καὶ ἐπεκάλεσεν αὐτῷ, εἰρήνη Κυρίου, ἕως τῆς ἡμέρας ταύτης, ἔτι αὐτοῦ ὄντος ἐν Ἐφραθὰ πατρὸς τοῦ Ἐσδρί.
\vs{25}Καὶ ἐγένετο ἐν τῇ νυκτὶ ἐκείνῃ, καὶ εἶπεν αὐτῷ Κύριος, λάβε τὸν μόσχον τὸν ταῦρον ὅς ἐστι τῷ πατρί σου, καὶ μόσχον δεύτερον ἑπταετῇ, καὶ καθελεῖς τὸ θυσιαστήριον τοῦ Βάαλ ὅ ἐστι τῷ πατρί σου, καὶ τὸ ἄλσος τὸ ἐπʼ αὐτὸ ὀλοθρεύσεις.
\vs{26}Καὶ οἰκοδομήσεις θυσιαστήριον τῷ Κυρίῳ τῷ Θεῷ σου ἐπὶ κορυφὴν Μαωζὶ τούτου ἐν τῇ παρατάξει καὶ λήψῃ τὸν μόσχον τὸν δεύτερον, καὶ ἀνοίσεις ὁλοκαυτώματα ἐν τοῖς ξύλοις τοῦ ἄλσοῦς, οὗ ἐξολοθρεύσεις.
\vs{27}Καὶ ἔλαβε Γεδεὼν δέκα ἄνδρας ἀπὸ τῶν δούλων ἑαυτοῦ, καὶ ἐποίησεν ὃν τρόπον ἐλάλησε πρὸς αὐτὸν Κύριος· καὶ ἐγενήθη ὡς ἐφοβήθη τὸν οἶκον τοῦ πατρὸς αὐτοῦ καὶ τοὺς ἄνδρας τῆς πόλεως τοῦ ποιῆσαι ἡμέρας, καὶ ἐποίησε νυκτός.

\vs{28}Καὶ ὤρθρισαν οἱ ἄνδρες τῆς πόλεως τοπρωΐ· καὶ ἰδοὺ καθῄρητο τὸ θυσιαστήριον τοῦ Βάαλ, καὶ τὸ ἄλσος τὸ ἐπʼ αὐτῷ ὠλόθρευτο· καὶ εἶδον τὸν μόσχον τὸν δεύτερον, ὃν ἀνήνεγκεν ἐπὶ τὸ θυσιαστήριον τὸ ᾠκοδομημένον.
\vs{29}Καὶ εἶπεν ἀνὴρ πρὸς τὸν πλησίον αὐτοῦ, τίς ἐποίησε τὸ ῥῆμα τοῦτο; καὶ ἐπεζήτησαν καὶ ἠρεύνησαν, καὶ ἔγνωσαν ὅτι Γεδεὼν υἱὸς Ἰωὰς ἐποίησε τὸ ῥῆμα τοῦτο.
\vs{30}Καὶ εἶπαν οἱ ἄνδρες τῆς πόλεως πρὸς Ἰωὰς, ἐξένεγκε τὸν υἱόν σου, καὶ ἀποθανέτω, ὅτι καθεῖλε τὸ θυσιαστήριον τοῦ Βάαλ, καὶ ὅτι ὠλόθρευσε τὸ ἄλσος τὸ ἐπʼ αὐτῷ.
\vs{31}Καὶ εἶπε Γεδεὼν υἱὸς Ἰωὰς τοῖς ἀνδράσι πᾶσιν, οἳ ἐπανέστησαν αὐτῷ, μὴ ὑμεῖς νῦν δικάζεσθε ὑπὲρ τοῦ Βάαλ; ἢ ὑμεῖς σώσετε αὐτόν; ὃς ἐὰν δικάσηται αὐτῷ, θανατωθήτω ἕως πρωΐ· εἰ θεός ἐστι, δικαζέσθω αὐτῷ, ὅτι καθεῖλε τὸ θυσιαστήριον αὐτοῦ.
\vs{32}Καὶ ἐκάλεσεν αὐτὸ ἐν τῇ ἡμέρᾳ ἐκείνῃ Ἱεροβάαλ, λέγων, δικαζέσθω ἐν αὐτῷ ὁ Βάαλ, ὅτι καθῃρέθη τὸ θυσιαστήριον αὐτοῦ.

\vs{33}Καὶ πᾶσα Μαδιὰμ, καὶ Ἀμαλὴκ, καὶ υἱοὶ ἀνατολῶν συνήχθησαν ἐπὶ τοαυτὸ, καὶ παρενέβαλον ἐν τῇ κοιλάδι Ἰεζραέλ.
\vs{34}Καὶ πνεῦμα Κυρίου ἐνέδυσε τὸν Γεδεὼν, καὶ ἐσάλπισεν ἐν κερατίνῃ, καὶ ἐβόησεν Ἀβιέζερ ὀπίσω αὐτοῦ.
\vs{35}Καὶ ἀγγέλους ἐξαπέστειλεν εἰς πάντα Μανασσῆ, καὶ ἐν Ἀσὴρ, καὶ ἐν Ζαβουλὼν, καὶ ἐν Νεφθαλί· καὶ ἀνέβη εἰς συνάντησιν αὐτῶν.

\vs{36}Καὶ εἶπε Γεδεὼν πρὸς τὸν Θεὸν, εἰ σὺ σώζεις ἐν χειρί μου τὸν Ἰσραὴλ, καθὼς ἐλάλησας,
\vs{37}ἰδοὺ ἐγὼ τίθημι τὸν πόκον τοῦ ἐρίου ἐν τῇ ἅλωνι· ἐὰν δρόσος γένηται ἐπὶ τὸν πόκον μόνον, καὶ ἐπὶ πᾶσαν τὴν γῆν ξηρασία, γνώσομαι ὅτι σώσεις ἐν χειρί μου τὸν Ἰσραὴλ, καθὼς ἐλάλησας.
\vs{38}Καὶ ἐγένετο οὕτως· καὶ ὤρθρισε τῇ ἐπαύριον, καὶ ἐξεπίασε τὸν πόκον, καὶ ἔσταξε δρόσος ἀπὸ τοῦ πόκου πλήρης λεκάνη ὕδατος.
\vs{39}Καὶ εἶπε Γεδεὼν πρὸς τὸν Θεὸν, μὴ δὴ ὀργισθήτω ὁ θυμός σου ἐν ἐμοὶ, καὶ λαλήσω ἔτι ἅπαξ· πειράσω δὴ καί γε ἔτι ἅπαξ ἐν τῷ πόκῳ· καὶ γενέσθω ἡ ξηρασία ἐπὶ τὸν πόκον μόνον, καὶ ἐπὶ πᾶσαν τὴν γῆν γενηθήτω δρόσος.
\vs{40}Καὶ ἐποίησεν ὁ Θεὸς οὕτως ἐν τῇ νυκτὶ ἐκείνῃ· καὶ ἐγένετο ξηρασία ἐπὶ τὸν πόκον μόνον, καὶ ἐπὶ πᾶσαν τὴν γῆν ἐγενήθη δρόσος.

\ch{7}
Καὶ ὤρθρισεν Ἱεροβάαλ, αὐτός ἐστι Γεδεὼν, καὶ πᾶς ὁ λαὸς μετʼ αὐτοῦ, καὶ παρενέβαλον ἐπὶ πηγὴν Ἀράδ· καὶ παρεμβολὴ Μαδιὰμ ἦν αὐτῷ ἀπὸ Βοῤῥᾶ ἀπὸ Γαβααθαμωραὶ ἐν κοιλάδι.

\vs{2}Καὶ εἶπε Κύριος πρὸς Γεδεὼν, πολὺς ὁ λαὸς ὁ μετὰ σοῦ, ὥστε μὴ παραδοῦναί με τὴν Μαδιὰμ ἐν χειρὶ αὐτῶν, μὴ ποτε καυχήσηται Ἰσραὴλ ἐπʼ ἐμὲ, λέγων, ἡ χείρ μου ἔσωσέ με.
\vs{3}Καὶ νῦν λάλησον δὴ ἐν ὠσὶ τοῦ λαοῦ, λέγων, τίς ὁ φοβούμενος καὶ δειλός; ἐπιστραφέτω καὶ ἐκχωρείτω ἀπὸ ὄρους Γαλαάδ· καὶ ἐπέστρεψεν ἀπὸ τοῦ λαοῦ εἴκοσι καὶ δύο χιλιάδες, καὶ δέκα χιλιάδες ὑπελείφθησαν.
\vs{4}Καὶ εἶπε Κύριος πρὸς Γεδεὼν, ἔτι ὁ λαὸς πολύς ἐστι· κατένεγκον αὐτοὺς πρὸς τὸ ὕδωρ, καὶ ἐκκαθαρῶ σοι αὐτὸν ἐκεῖ· καὶ ἔσται ὃν ἐὰν εἴπω πρὸς σὲ, οὗτος πορεύσεται σὺν σοὶ, αὐτὸς πορεύσεται σὺν σοί· καὶ πᾶς ὃν ἂν εἴπω πρὸς σὲ, οὗτος οὐ πορεύσεται μετὰ σοῦ, αὐτὸς οὐ πορεύσεται μετὰ σοῦ.
\vs{5}Καὶ κατήνεγκε τὸν λαὸν πρὸς τὸ ὕδωρ· καὶ εἶπε Κύριος πρὸς Γεδεὼν, πᾶς ὃς ἂν λάψῃ τῇ γλώσσῃ αὐτοῦ ἀπὸ τοῦ ὕδατος ὡς ἐὰν λάψῃ ὁ κύων, στήσεις αὐτὸν κατὰ μόνας, καὶ πᾶς ὃς ἐὰν κλίνῃ ἐπὶ τὰ γόνατα αὐτοῦ πιεῖν.
\vs{6}Καὶ ἐγένετο ὁ ἀριθμὸς τῶν λαψάντων ἐν χειρὶ αὐτῶν πρὸς τὸ στόμα αὐτῶν, τριακόσιοι ἄνδρες· καὶ πᾶν τὸ κατάλοιπον τοῦ λαοῦ ἔκλιναν ἐπὶ τὰ γόνατα αὐτῶν πιεῖν ὕδωρ.
\vs{7}Καὶ εἶπε Κύριος πρὸς Γεδεὼν, ἐν τοῖς τριακοσίοις ἀνδράσι τοῖς λάψασι σώσω ὑμᾶς, καὶ δώσω τὴν Μαδιὰμ ἐν χειρί σου, καὶ πᾶς ὁ λαὸς πορεύσονται ἀνὴρ εἰς τὸν τόπον αὐτοῦ.
\vs{8}Καὶ ἔλαβον τὸν ἐπισιτισμὸν τοῦ λαοῦ ἐν χειρὶ αὐτῶν, καὶ τὰς κερατίνας αὐτῶν· καὶ τὸν πάντα ἄνδρα Ἰσραὴλ ἐξαπέστειλεν ἄνδρα εἰς σκηνὴν αὐτοῦ· καὶ τοὺς τριακοσίους ἄνδρας κατίσχυσε· καὶ ἡ παρεμβολὴ Μαδιὰμ ἦσαν αὐτοῦ ὑποκάτω ἐν τῇ κοιλάδι.

\vs{9}Καὶ ἐγενήθη ἐν τῇ νυκτὶ ἐκείνῃ, καὶ εἶπε πρὸς αὐτὸν Κύριος, ἀνάστα, κατάβηθι ἐν τῇ παρεμβολῇ, ὅτι παρέδωκα αὐτὴν ἐν τῇ χειρί σου.
\vs{10}Καὶ εἰ φοβῇ σὺ καταβῆναι, κατάβηθι σὺ καὶ Φαρὰ τὸ παιδάριόν σου εἰς τὴν παρεμβολὴν,
\vs{11}καὶ ἀκούσῃ τί λαλήσουσι, καὶ μετὰ τοῦτο ἰσχύσουσιν αἱ χεῖρές σου καὶ καταβήσῃ ἐν τῇ παρεμβολῇ· καὶ κατέβη αὐτὸς καὶ Φαρὰ τὸ παιδάριον αὐτοῦ πρὸς ἀρχὴν τῶν πεντήκοντα, οἳ ἦσαν ἐν τῇ παρεμβολῇ.
\vs{12}Καὶ Μαδιὰμ καὶ Ἀμαλὴκ καὶ πάντες οἱ υἱοὶ ἀνατολῶν βεβλημένοι ἐν τῇ κοιλάδι ὡς ἀκρὶς εἰς πλῆθος, καὶ ταῖς καμήλοις αὐτῶν οὐκ ἦν ἀριθμὸς, ἀλλʼ ἦσαν ὡς ἡ ἄμμος ἡ ἐπὶ χείλους τῆς θαλάσσης εἰς πλῆθος.

\vs{13}Καὶ ἦλθε Γεδεὼν, καὶ ἰδοὺ ἀνὴρ ἐξηγούμενος τῷ πλησίον αὐτοῦ ἐνύπνιον, καὶ εἶπεν, ἰδοὺ ἐνυπνιασάμην ἐνύπνιον, καὶ ἰδοὺ μαγὶς ἄρτου κριθίνου στρεφομένη ἐν τῇ παρεμβολῇ Μαδιὰμ, καὶ ἦλθεν ἕως τῆς σκηνῆς, καὶ ἐπάταξεν αὐτὴν, καὶ ἔπεσε, καὶ ἀνέστρεψεν αὐτὴν ἄνω, καὶ ἔπεσεν ἡ σκηνή.
\vs{14}Καὶ ἀπεκρίθη ὁ πλησίον αὐτοῦ, καὶ εἶπεν, οὐκ ἔστιν αὕτη εἰ μὴ ῥομφαία Γεδεὼν υἱοῦ Ἰωὰς ἀνδρὸς Ἰσραήλ· παρέδωκεν ὁ Θεὸς ἐν χειρὶ αὐτοῦ τὴν Μαδιὰμ καὶ πᾶσαν τὴν παρεμβολήν.

\vs{15}Καὶ ἐγένετο ὡς ἤκουσε Γεδεὼν τὴν ἐξήγησιν τοῦ ἐνυπνίου καὶ τὴν σύγκρισιν αὐτοῦ, καὶ προσεκύνησε Κυρίῳ, καὶ ὑπέστρεψεν εἰς τὴν παρεμβολὴν Ἰσραὴλ, καὶ εἶπεν, ἀνάστητε, ὅτι παρέδωκε Κύριος ἐν χειρὶ ἡμῶν τὴν παρεμβολὴν Μαδιάμ.
\vs{16}Καὶ διεῖλε τοὺς τριακοσίους ἄνδρας εἰς τρεῖς ἀρχὰς, καὶ ἔδωκε κερατίνας ἐν χειρὶ πάντων, καὶ ὑδρίας κενὰς, καὶ λαμπάδας ἐν ταῖς ὑδρίαις,
\vs{17}καὶ εἶπε πρὸς αὐτοὺς, ἀπʼ ἐμοῦ ὄψεσθε, καὶ οὕτω ποιήσετε· καὶ ἰδοὺ ἐγὼ εἰσπορεύομαι ἐν ἀρχῇ τῆς παρεμβολῆς, καὶ ἔσται καθὼς ἂν ποιήσω, οὕτω ποιήσετε.
\vs{18}Καὶ σαλπιῶ ἐν τῇ κερατίνῃ ἐγὼ, καὶ πάντες μετʼ ἐμοῦ σαλπιεῖτε ἐν ταῖς κερατίναις κύκλῳ ὅλης τῆς παρεμβολῆς, καὶ ἐρεῖτε, τῷ Κυρίῳ καὶ τῷ Γεδεών.

\vs{19}Καὶ εἰσῆλθε Γεδεὼν καὶ οἱ ἐκατὸν ἄνδρες οἱ μετʼ αὐτοῦ ἐν ἀρχῇ τῆς παρεμβολῆς ἐν ἀρχῇ τῆς φυλακῆς μέσης· καὶ ἐγείροντες ἤγειραν τοὺς φυλάσσοντας, καὶ ἐσάλπισαν ἐν ταῖς κερατίναις, καὶ ἐξετίναξαν τὰς ὑδρίας τὰς ἐν ταῖς χερσὶν αὐτῶν.
\vs{20}Καὶ ἐσάλπισαν αἱ τρεῖς ἀρχαὶ ἐν ταῖς κερατίναις, καὶ συνέτριψαν τὰς ὑδρίας, καὶ ἐκράτησαν ἐν χερσὶν ἀριστεραῖς αὐτῶν τὰς λαμπάδας, καὶ ἐν χερσὶ δεξιαῖς αὐτῶν τὰς κερατίνας τοῦ σαλπίζειν· καὶ ἀνέκραξαν, ῥομφαία τῷ Κυρίῳ καὶ τῷ Γεδεών.
\vs{21}Καὶ ἔστησεν ἀνὴρ ἐφʼ ἑαυτῷ κύκλῳ τῆς παρεμβολῆς· καὶ ἔδραμε πᾶσα ἡ παρεμβολὴ, καὶ ἐσήμαναν, καὶ ἔφυγον.
\vs{22}Καὶ ἐσάλπισαν ἐν ταῖς τριακοσίαις κερατίναις· καὶ ἔθηκε Κύριος τὴν ῥομφαίαν ἀνδρὸς ἐν τῷ πλησίον αὐτοῦ ἐν πάσῃ τῇ παρεμβολῇ.
\vs{23}Καὶ ἔφυγεν ἡ παρεμβολὴ ἕως Βηθσεὲδ Ταγαραγαθὰ Ἀβελμεουλὰ ἐπὶ Ταβάθ· καὶ ἐβόησαν ἀνὴρ Ἰσραὴλ ἀπὸ Νεφθαλὶ καὶ ἀπὸ Ἀσὴρ, καὶ ἀπὸ παντὸς Μανασσῆ, καὶ ἐδίωξαν ὀπίσω Μαδιάμ.

\vs{24}Καὶ ἀγγέλους ἀπέστειλε Γεδεὼν ἐν παντὶ ὄρει Ἐφραὶμ, λέγων, κατάβητε εἰς συνάντησιν Μαδιὰμ, καὶ καταλάβετε ἑαυτοῖς τὸ ὕδωρ ἕως Βαιθηρὰ καὶ τὸν Ἰορδάνην· καὶ ἐβόησε πᾶς ἀνὴρ Ἐφραὶμ, καὶ προκατελάβοντο τὸ ὕδωρ ἕως Βαιθηρὰ καὶ τὸν Ἰορδάνην.
\vs{25}Καὶ συνελάβοντο τοὺς ἄρχοντας Μαδιὰμ, καὶ τὸν Ὠρὴβ καὶ τὸν Ζήβ· καὶ ἀπέκτειναν τὸν Ὠρὴβ ἐν Σοὺρ Ὠρὴβ, καὶ τὸν Ζὴβ ἀπέκτειναν ἐν Ἱακεφζήφ· καὶ κατεδίωξαν τὴν Μαδιάμ· καὶ τὴν κεφαλὴν Ὠρὴβ καὶ Ζὴβ ἤνεγκαν πρὸς Γεδεὼν ἀπὸ πέραν τοῦ Ἰορδάνου.

\ch{8}
Καὶ εἶπαν πρὸς Γεδεὼν ἀνὴρ Ἐφραὶμ, τί τὸ ῥῆμα τοῦτο ἐποίησας ἡμῖν, τοῦ μὴ καλέσαι ἡμᾶς ὅτε ἐπορεύθης παρατάξασθαι ἐν Μαδιάμ; καὶ διελέξαντο πρὸς αὐτὸν ἰσχυρῶς.
\vs{2}Καὶ εἶπε πρὸς αὐτοὺς, τί ἐποίησα νῦν καθὼς ὑμεῖς; ἢ οὐχὶ κρείττῶν ἐπιφυλλὶς Ἐφραὶμ ἢ τρυγητὸς Ἀβιέζερ;
\vs{3}Ἐν χειρὶ ὑμῶν παρέδωκε Κύριος τοὺς ἄρχοντας Μαδιὰμ, τὸν Ὠρὴβ καὶ τὸν Ζήβ· καὶ τί ἠδυνήθην ποιῆσαι ὡς ὑμεῖς; τότε ἀνέθη τὸ πνεῦμα αὐτῶν ἀπʼ αὐτοῦ ἐν τῷ λαλῆσαι αὐτὸν τὸν λόγον τοῦτον.

\vs{4}Καὶ ἦλθε Γεδεὼν ἐπὶ τὸν Ἰορδανὴν, καὶ διέβη αὐτὸς καὶ οἱ τριακόσιοι ἄνδρες οἱ μετʼ αὐτοῦ πεινῶντες καὶ διώκοντες.
\vs{5}Καὶ εἶπε τοῖς ἀνδράσι Σοκχὼθ, δότε δὴ ἄρτους εἰς τροφὴν τῷ λαῷ τούτῳ τῷ ἐν ποσί μου, ὅτι ἐκλείπουσι, καὶ ἰδοὺ ἐγώ εἰμι διώκων ὀπίσω τοῦ Ζεβεὲ καὶ Σαλμανὰ βασιλέων Μαδιάμ.
\vs{6}Καὶ εἶπον οἱ ἄρχοντες Σοκχὼθ, μὴ χεὶρ Ζεβεὲ καὶ Σαλμανὰ νῦν ἐν χειρί σου, ὅτι δώσομεν τῇ δυνάμει σου ἄρτους;
\vs{7}Καὶ εἶπε Γεδεὼν, διὰ τοῦτο ἐν τῷ δοῦναι Κύριον τὸν Ζεβεὲ καὶ τὸν Σαλμανὰ ἐν χειρί μου, καὶ ἐγὼ ἀλοήσω τὰς σάρκας ὑμῶν ἐν ταῖς ἀκάνθαις τῆς ἐρήμου, καὶ ἐν ταῖς Βαρκηνίμ.
\vs{8}Καὶ ἀνέβη ἐκεῖθεν εἰς Φανουὴλ, καὶ ἐλάλησε πρὸς αὐτοὺς ὡσαύτως· καὶ ἀπεκρίθησαν αὐτῷ οἱ ἄνδρες Φανουὴλ ὃν τρόπον ἀπεκρίθησαν ἄνδρες Σοκχώθ.
\vs{9}Καὶ εἶπε Γεδεὼν πρὸς ἄνδρας Φανουὴλ, ἐν ἐπιστροφῇ μου μετʼ εἰρήνης, κατασκάψω τὸν πύργον τοῦτον.

\vs{10}Καὶ Ζεβεὲ καὶ Σαλμανὰ ἐν Καρκὰρ, καὶ ἡ παρεμβολὴ αὐτῶν μετʼ αὐτῶν ὡσεὶ δεκαπέντε χιλιάδες, πάντες οἱ καταλελειμμένοι ἀπὸ πάσης παρεμβολῆς ἀλλοφύλων· καὶ οἱ πεπτωκότες, ἑκατὸν εἴκοσι χιλιάδες ἀνδρῶν σπωμένων ῥομφαίαν.
\vs{11}Καὶ ἀνέβη Γεδεὼν ὁδὸν τῶν σκηνούντων ἐν σκηναῖς ἀπὸ ἀνατολῶν τῆς Ναβαὶ καὶ Ἰεγεβάλ· καὶ ἐπάταξε τὴν παρεμβολὴν, καὶ ἡ παρεμβολὴ ἦν πεποιθυῖα.
\vs{12}Καὶ ἔφυγον Ζεβεὲ καὶ Σαλμανά· καὶ ἐδίωξεν ὀπίσω αὐτῶν, καὶ ἐκράτησε τοὺς δύο βασιλεῖς Μαδιὰμ τὸν Ζεβεὲ καὶ τὸν Σαλμανὰ, καὶ πᾶσαν τὴν παρεμβολὴν ἐξέστησε.

\vs{13}Καὶ ἐπέστρεψε Γεδεὼν υἱὸς Ἰωὰς ἀπὸ τῆς παρατάξεως ἀπὸ ἐπάνωθεν τῆς παρατάξεως Ἀρές.
\vs{14}Καὶ συνέλαβε παιδάριον ἀπὸ τῶν ἀνδρῶν Σοκχὼθ, καὶ ἐπηρώτησεν αὐτόν· καὶ ἔγραψε πρὸς αὐτὸν ὀνόματα τῶν ἀρχόντων Σοκχὼθ καὶ τῶν πρεσβυτέρων αὐτῶν, ἑβδομήκοντα καὶ ἑπτὰ ἄνδρας.
\vs{15}Καὶ παρεγένετο Γεδεὼν πρὸς τοὺς ἄρχοντας Σοκχὼθ, καὶ εἶπεν, ἰδοὺ Ζεβεὲ καὶ Σαλμανὰ, ἐν οἷς ὠνειδίσατέ με, λέγοντες, μὴ χεὶρ Ζεβεὲ καὶ Σαλμανὰ νῦν ἐν χειρί σου, ὅτι δώσομεν τοῖς ἀνδράσι σου τοῖς ἐκλείπουσιν ἄρτους;
\vs{16}Καὶ ἔλαβε τοὺς πρεσβυτέρους τῆς πόλεως ἐν ταῖς ἀκάνθαις τῆς ἐρήμου καὶ ταῖς Βαρκηνὶμ, καὶ ἠλόησεν ἐν αὐτοῖς τοὺς ἄνδρας τῆς πόλεως·
\vs{17}Καὶ τὸν πύργον Φανουὴλ κατέστρεψε, καὶ ἀπέκτεινε τοὺς ἄνδρας τῆς πόλεως.

\vs{18}Καὶ εἶπε πρὸς Ζεβεὲ καὶ Σαλμανὰ, ποῦ οἱ ἄνδρες, οὓς ἀπεκτείνατε ἐν Θαβώρ; καὶ εἶπαν, ὡς σὺ, ὡς αὐτοὶ, εἰς ὁμοίωμα υἱοῦ βασιλέως.
\vs{19}Καὶ εἶπε Γεδεὼν, ἀδελφοί μου καὶ υἱοὶ τῆς μητρός μου ἦσαν· ζῇ Κύριος· εἰ ἐζωογονήκειτε αὐτοὺς, οὐκ ἂν ἀπέκτεινα ὑμᾶς.
\vs{20}Καὶ εἶπεν Ἰεθὲρ τῷ πρωτοτόκῳ αὐτοῦ, ἀναστὰς, ἀπόκτεινον αὐτούς· καὶ οὐκ ἔσπασε τὸ παιδάριον τὴν ῥομφαίαν αὐτοῦ, ὅτι ἐφοβήθη, ὅτι ἔτι νεώτερος ἦν.
\vs{21}Καὶ εἶπε Ζεβεὲ καὶ Σαλμανὰ, ἀνάστα σὺ, καὶ συνάντησον ἡμῖν, ὅτι ὡς ἀνδρὸς ἡ δύναμίς σου· καὶ ἀνέστη Γεδεὼν, καὶ ἀπέκτεινε τὸν Ζεβεὲ καὶ τὸν Σαλμανά· καὶ ἔλαβε τοὺς μηνίσκους τοὺς ἐν τοῖς τραχήλοις τῶν καμήλων αὐτῶν.

\vs{22}Καὶ εἶπον ἀνὴρ Ἰσραὴλ πρὸς Γεδεὼν, κύριε, ἄρξον ἡμῶν καὶ σὺ, καὶ ὁ υἱός σου, καὶ ὁ υἱὸς τοῦ υἱοῦ σου, ὅτι σὺ ἔσωσας ἡμᾶς ἐκ χειρὸς Μαδιάμ.
\vs{23}Καὶ εἶπε πρὸς αὐτοὺς Γεδεὼν, οὐκ ἄρξω ἐγὼ, καὶ οὐκ ἄρξει ὁ υἱός μου ἐν ὑμῖν· Κύριος ἄρξει ὑμῶν.
\vs{24}Καὶ εἶπε πρὸς αὐτοὺς Γεδεὼν, αἰτήσομαι παρʼ ὑμῶν αἴτημα, καὶ δότε μοι ἀνὴρ ἐνώτιον ἐκ σκύλων αὐτοῦ· ὅτι ἐνώτια χρυσᾶ αὐτοῖς, ὅτι ἦσαν Ἰσμαηλῖται.
\vs{25}Καὶ εἶπαν, διδόντες δώσομεν· καὶ ἀνέπτυξε τὸ ἱμάτιον αὐτοῦ, καὶ ἔβαλεν ἐκεῖ ἀνὴρ ἐνώτιον σκύλων αὐτοῦ.
\vs{26}Καὶ ἐγένετο ὁ σταθμὸς τῶν ἐνωτίων τῶν χρυσῶν ὧν ᾔτησε, χίλιοι καὶ ἑπτακόσιοι χρυσοὶ, πάρεξ τῶν μηνίσκων καὶ τῶν στραγγαλίδων καὶ τῶν ἱματίων καὶ πορφυρίδων τῶν ἐπὶ βασιλεῦσι Μαδιὰμ, καὶ ἐκτὸς τῶν περιθεμάτων ἃ ἦν ἐν τοῖς τραχήλοις τῶν καμήλων αὐτῶν.
\vs{27}Καὶ ἐποίησεν αὐτὸ Γεδεὼν εἰς ἐφὼδ, καὶ ἔστησεν αὐτὸ ἐν πόλει αὐτοῦ ἐν Ἐφραθά καὶ ἐξεπόρνευσε πᾶς Ἰσραὴλ ὀπίσω αὐτοῦ ἐκεῖ· καὶ ἐγένετο τῷ Γεδεὼν καὶ τῷ οἴκῳ αὐτοῦ εἰς σκῶλον.

\vs{28}Καὶ συνεστάλη Μαδιὰμ ἐνώπιον υἱῶν Ἰσραὴλ, καὶ οὐ προσέθηκαν ἆραι κεφαλὴν αὐτῶν· καὶ ἡσύχασεν ἡ γῆ τεσσαράκοντα ἔτη ἐν ἡμέραις Γεδεών.
\vs{29}Καὶ ἐπορεύθη Ἱεροβάαλ υἱὸς Ἰωὰς, καὶ ἐκάθισεν ἐν οἴκῳ αὐτοῦ.
\vs{30}Καὶ τῷ Γεδεὼν ἦσαν υἱοὶ ἑβδομήκοντα ἐκπορευόμενοι ἐκ μηρῶν αὐτοῦ, ὅτι γυναῖκες πολλαὶ ἦσαν αὐτῷ.
\vs{31}Καὶ παλλακὴ αὐτοῦ ἦν ἐν Συχὲμ, καὶ ἔτεκεν αὐτῷ καί γε αὐτὴ υἱὸν, καὶ ἔθηκε τὸ ὄνομα Ἀβιμέλεχ.
\vs{32}Καὶ ἀπέθανε Γεδεὼν υἱὸς Ἰωὰς ἐν πόλει αὐτοῦ, καὶ ἐτάφη ἐν τῷ τάφῳ Ἰωὰς τοῦ πατρὸς αὐτοῦ ἐν Ἐφραθὰ Ἀβὶ Ἐσδρί.

\vs{33}Καὶ ἐγενήθη ὡς ἀπέθανε Γεδεὼν, καὶ ἐπέστρεψαν οἱ υἱοὶ Ἰσραὴλ, καὶ ἐξεπόρνευσαν ὀπίσω τῶν Βααλὶμ, καὶ ἔθηκαν ἑαυτοῖς τῷ Βάαλ διαθήκην τοῦ εἶναι αὐτοῖς αὐτὸν εἰς θεόν.
\vs{34}Καὶ οὐκ ἐμνήσθησαν οἱ υἱοὶ Ἰσραὴλ Κυρίου τοῦ Θεοῦ τοῦ ῥυσαμένου αὐτοὺς ἐκ χειρὸς πάντων τῶν θλιβόντων αὐτοὺς κυκλόθεν.
\vs{35}Καὶ οὐκ ἐποίησαν ἔλεος μετὰ τοῦ οἴκου Ἱεροβάαλ, αὐτός ἐστι Γεδεὼν, κατὰ πάντα τὰ ἀγαθὰ ἃ ἐποίησε μετὰ Ἰσραήλ.

\ch{9}
Καὶ ἐπορεύθη Ἀβιμέλεχ υἱὸς Ἱεροβάαλ εἰς Συχὲμ πρὸς ἀδελφοὺς μητρὸς αὐτοῦ· καὶ ἐλάλησε πρὸς αὐτοὺς καὶ πρὸς πᾶσαν συγγένειαν οἴκου πατρὸς μητρὸς αὐτοῦ, λέγων,
\vs{2}λαλήσατε δὴ ἐν τοῖς ὠσὶ πάντων τῶν ἀνδρῶν Συχὲμ, τί τὸ ἀγαθὸν ὑμῖν, κυριεῦσαι ὑμῶν ἑβδομήκοντα ἄνδρας πάντας υἱοὺς Ἱεροβάαλ, ἢ κυριεύειν ὑμῶν ἄνδρα ἕνα; καὶ μνήσθητε ὅτι ὀστοῦν ὑμῶν καὶ σὰρξ ὑμῶν εἰμι.
\vs{3}Καὶ ἐλάλησαν περὶ αὐτοῦ οἱ ἀδελφοὶ τῆς μητρὸς αὐτοῦ ἐν τοῖς ὠσὶ πάντων τῶν ἀνδρῶν Συχὲμ πάντας τοὺς λόγους τούτους· καὶ ἔκλινεν ἡ καρδία αὐτῶν ὀπίσω Ἀβιμέλεχ, ὅτι εἶπαν, ἀδελφός ἡμῶν ἐστι.
\vs{4}Καὶ ἔδωκαν αὐτῷ ἑβδομήκοντα ἀργυρίου ἐξ οἴκου Βααλβερίθ· καὶ ἐμισθώσατο ἑαυτῷ Ἀβιμέλεχ ἄνδρας κενοὺς καὶ δειλοὺς, καὶ ἐπορεύθησαν ὀπίσω αὐτοῦ.
\vs{5}Καὶ εἰσῆλθεν εἰς τὸν οἶκον τοῦ πατρὸς αὐτοῦ εἰς Ἐφραθὰ, καὶ ἀπέκτεινε τοὺς ἀδελφοὺς αὐτοῦ υἱοὺς Ἱεροβάαλ, ἑβδομήκοντα ἄνδρας ἐπὶ λίθον ἕνα· καὶ κατελείφθη Ἰωάθαμ υἱὸς Ἱεροβάαλ ὁ νεώτερος, ὅτι ἐκρύβη.

\vs{6}Καὶ συνήχθησαν πάντες ἄνδρες Σικίμων, καὶ πᾶς οἶκος Βηθμααλὼ, καὶ ἐπορεύθησαν, καὶ ἐβασίλευσαν τὸν Ἀβιμέλεχ πρὸς τῇ βαλάνῳ τῇ εὑρετῇ τῆς στάσεως τῆς ἐν Σικίμοις.

\vs{7}Καὶ ἀνηγγέλη τῷ Ἰωάθαμ, καὶ ἐπορεύθη, καὶ ἔστη ἐπὶ κορυφὴν ὄρους Γαριζὶν, καὶ ἐπῇρε τὴν φωνὴν αὐτοῦ, καὶ ἔκλαυσε, καὶ εἶπεν αὐτοῖς, ἀκούσατέ μου ἄνδρες Σικίμων, καὶ ἀκούσεται ὑμῶν ὁ Θεός.

\vs{8}Πορευόμενα ἐπορεύθη τὰ ξύλα τοῦ χρίσαι ἐφʼ ἑαυτὰ βασιλέα, καὶ εἶπον τῇ ἐλαίᾳ, βασίλευσον ἐφʼ ἡμῶν.
\vs{9}Καὶ εἶπεν αὐτοῖς ἡ ἐλαία, μὴ ἀπολείψασα τὴν πιότητά μου, ἐν ᾗ δοξάσουσι τὸν Θεὸν ἄνδρες, πορεύσομαι κινεῖσθαι ἐπὶ τῶν ξύλων;
\vs{10}Καὶ εἶπον τὰ ξύλα τῇ συκῇ, δεῦρο, βασίλευσον ἐφʼ ἡμῶν.
\vs{11}Καὶ εἶπεν αὐτοῖς ἡ συκῆ, μὴ ἀπολείψασα ἐγὼ τὴν γλυκύτητά μου καὶ τὰ γεννήματά μου τὰ ἀγαθὰ, πορεύσομαι κινεῖσθαι ἐπὶ τῶν ξύλων;
\vs{12}Καὶ εἶπαν τὰ ξύλα πρὸς τὴν ἄμπελον, δεῦρο, βασίλευσον ἐφʼ ἡμῶν.
\vs{13}Καὶ εἶπεν αὐτοῖς ἡ ἄμπελος, μὴ ἀπολείψασα τὸν οἶνόν μου τὸν εὐφραίνοντα Θεὸν καὶ ἀνθρώπους, πορεύσομαι κινεῖσθαι ἐπὶ τῶν ξύλων;
\vs{14}Καὶ εἶπαν πάντα τὰ ξύλα τῇ ῥάμνῳ, δεῦρο σὺ, βασίλευσον ἐφʼ ἡμῶν.
\vs{15}Καὶ εἶπεν ἡ ῥάμνος πρὸς τὰ ξύλα, εἰ ἐν ἀληθείᾳ χρίετέ με ὑμεῖς τοῦ βασιλεύειν ἐφʼ ὑμᾶς, δεῦτε, ὑπόστητε ἐν τῇ σκιᾷ μου· καὶ εἰ μὴ, ἐξέλθοι πῦρ ἀπʼ ἐμοῦ καὶ καταφάγοι τὰς κέδρους τοῦ Λιβάνου.

\vs{16}Καὶ νῦν εἰ ἐν ἀληθείᾳ καὶ τελειότητι ἐποιήσατε, καὶ ἐβασιλεύσατε τὸν Ἀβιμέλεχ, καὶ εἰ ἀγαθωσύνην ἐποιήσατε μετὰ Ἱεροβάαλ, καὶ μετὰ τοῦ οἴκου αὐτοῦ, καὶ εἰ ὡς ἀνταπόδοσις χειρὸς αὐτοῦ ἐποιήσατε αὐτῷ,
\vs{17}ὡς παρετάξατο ὁ πατήρ μου ὑπὲρ ὑμῶν, καὶ ἐξέῤῥιψε τὴν ψυχὴν αὐτοῦ ἐξεναντίας, καὶ ἐῤῥύσατο ὑμᾶς ἐκ χειρὸς Μαδιὰμ,
\vs{18}καὶ ὑμεῖς ἐπανέστητε ἐπὶ τὸν οἶκον τοῦ πατρός μου σήμερον, καὶ ἀπεκτείνατε τοὺς υἱοὺς αὐτοῦ ἑβδομήκοντα ἄνδρας ἐπὶ λίθον ἕνα, καὶ ἐβασίλευσατε τὸν Ἀβιμέλεχ υἱὸν παιδίσκης αὐτοῦ ἐπὶ τοὺς ἄνδρας Σικίμων, ὅτι ἀδελφὸς ὑμῶν ἐστι·
\vs{19}Καὶ εἰ ἐν ἀληθείᾳ καὶ τελειότητι ἐποιήσατε μετὰ Ἱεροβάαλ, καὶ μετὰ τοῦ οἴκου αὐτοῦ ἐν τῇ ἡμέρᾳ ταύτῃ, εὐφρανθείητε ἐν Ἀβιμέλεχ, καὶ εὐφρανθείη καί γε αὐτὸς ἐφʼ ὑμῖν·
\vs{20}Εἰ δὲ οὐ, ἐξέλθοι πῦρ ἀπὸ Ἀβιμέλεχ, καὶ καταφάγοι τοὺς ἄνδρας Σικίμων καὶ τὸν οἴκον Βηθμααλώ· καὶ ἐξέλθοι πῦρ ἀπὸ ἀνδρῶν Σικίμων, καὶ ἐκ τοῦ οἴκου Βηθμααλὼ, καὶ καταφάγοι τὸν Ἀβιμέλεχ.

\vs{21}Καὶ ἔφυγεν Ἰωάθαμ καὶ ἀπέδρα, καὶ ἐπορεύθη ἕως Βαιὴρ, καὶ ᾤκησεν ἐκεῖ ἀπὸ προσώπου Ἀβιμέλεχ ἀδελφοῦ αὐτοῦ.

\vs{22}Καὶ ἦρξεν Ἀβιμέλεχ ἐπὶ Ἰσραὴλ τρία ἔτη.
\vs{23}Καὶ ἐξαπέστειλεν ὁ Θεὸς πνεῦμα πονηρὸν ἀναμέσον Ἀβιμέλεχ καὶ ἀναμέσον τῶν ἀνδρῶν Σικίμων· καὶ ἠθέτισαν ἄνδρες Σικίμων ἐν τῷ οἴκῳ Ἀβιμέλεχ,
\vs{24}τοῦ ἐπαγαγεῖν τὴν ἀδικίαν τῶν ἑβδομήκοντα υἱῶν Ἱεροβάαλ, καὶ τὰ αἵματα αὐτῶν τοῦ θεῖναι ἐπὶ Ἀβιμέλεχ τὸν ἀδελφὸν αὐτῶν, ὃς ἀπέκτεινεν αὐτοὺς, καὶ ἐπὶ ἄνδρας Σικίμων, ὅτι ἐνίσχυσαν τὰς χεῖρας αὐτοῦ ἀποκτεῖναι τοὺς ἀδελφοὺς αὐτοῦ.
\vs{25}Καὶ ἔθηκαν αὐτῷ οἱ ἄνδρες Σικίμων ἐνεδρεύοντας ἐπὶ τὰς κεφαλὰς τῶν ὀρέων, καὶ διήρπαζον πάντα ὃς παρεπορεύετο ἐπʼ αὐτοὺς ἐν τῇ ὁδῷ· καὶ ἀπηγγέλη τῷ βασιλεῖ Ἀβιμέλεχ.

\vs{26}Καὶ ἦλθε Γαὰλ υἱὸς Ἰωβὴλ, καὶ οἱ ἀδελφοὶ αὐτοῦ, καὶ παρῆλθον ἐν Σικίμοις, καὶ ἤλπισαν ἐν αὐτῷ οἱ ἄνδρες Σικίμων.
\vs{27}Καὶ ἐξῆλθον εἰς ἀγρὸν, καὶ ἐτρύγησαν τοὺς ἀμπελῶνας αὐτῶν, καὶ ἐπάτησαν, καὶ ἐποίησαν Ἐλλουλίμ· καὶ εἰσήνεγκαν εἰς οἶκον θεοῦ αὐτῶν, καὶ ἔφαγον καὶ ἔπιον, καὶ κατηράσαντο τὸν Ἀβιμέλεχ.
\vs{28}Καὶ εἶπε Γαὰλ υἱὸς Ἰωβὴλ, τίς ἐστιν Ἀβιμέλεχ, καὶ τίς ἐστιν υἱὸς Συχὲμ, ὅτι δουλεύσομεν αὐτῷ; οὐχ υἱὸς Ἱεροβάαλ, καὶ Ζεβοὺλ ἐπίσκοπος αὐτοῦ, δοῦλος αὐτοῦ σὺν τοῖς ἀνδράσιν Ἐμμὼρ πατρὸς Συχέμ; καὶ τί ὅτι δουλεύσομεν αὐτῷ ἡμεῖς;
\vs{29}Καὶ τίς δῴη τὸν λαὸν τοῦτον ἐν χειρί μου; καὶ μεταστήσω τὸν Ἀβιμέλεχ, καὶ ἐρῶ πρὸς αὐτὸν, πλήθυνον τὴν δύναμίν σου καὶ ἔξελθε.

\vs{30}Καὶ ἤκουσε Ζεβοὺλ ἄρχων τῆς πόλεως τοὺς λόγους Γαὰλ υἱοῦ Ἰωβὴλ, καὶ ὠργίσθη θυμῷ αὐτός.
\vs{31}Καὶ ἀπέστειλεν ἀγγέλους πρὸς Ἀβιμέλεχ ἐν κρυφῇ, λέγων, ἰδοὺ Γαὰλ υἱὸς Ἰωβὴλ καὶ οἱ ἀδελφοὶ αὐτοῦ ἔρχονται εἰς Συχὲμ, καὶ ἰδοὺ αὐτοὶ περικάθηνται τὴν πόλιν ἐπὶ σέ.
\vs{32}Καὶ νῦν ἀνάστηθι νυκτὸς, καὶ ὁ λαὸς ὁ μετὰ σου, καὶ ἐνέδρευσον ἐν τῷ ἀγρῷ.
\vs{33}Καὶ ἔσται τοπρωῒ ἅμα τῷ ἀνατεῖλαι τὸν ἥλιον, ὀρθριεῖς καὶ ἐκτενεῖς ἐπὶ τὴν πόλιν· καὶ ἰδοὺ αὐτὸς καὶ ὁ λαὸς ὁ μετʼ αὐτοῦ ἐκπορεύονται πρὸς σὲ, καὶ ποιήσεις αὐτῷ ὅσα ἂν εὕρῃ ἡ χείρ σου.

\vs{34}Καὶ ἀνέστη Ἀβιμέλεχ καὶ πᾶς ὁ λαὸς μετʼ αὐτοῦ νυκτὸς, καὶ ἐνήδρευσαν ἐπὶ Συχὲμ τέτρασιν ἀρχαῖς.
\vs{35}Καὶ ἐξῆλθε Γαὰλ υἱὸς Ἰωβὴλ, καὶ ἔστη πρὸς τῇ θύρᾳ τῆς πύλης τῆς πόλεως· καὶ ἀνέστη Ἀβιμέλεχ καὶ ὁ λαὸς ὁ μετʼ αὐτοῦ ἀπὸ τοῦ ἐνέδρου.
\vs{36}Καὶ εἶδε Γαὰλ υἱὸς Ἰωβὴλ τὸν λαὸν, καὶ εἶπε πρὸς Ζεβοὺλ, ἰδοὺ λαὸς καταβαίνει ἀπὸ τῶν κεφαλῶν τῶν ὀρέων· καὶ εἶπε πρὸς αὐτὸν Ζεβοὺλ, τὴν σκιὰν τῶν ὀρέων σὺ βλέπεις ὡς ἄνδρας.
\vs{37}Καὶ προσέθετο ἔτι Γαὰλ τοῦ λαλῆσαι, καὶ εἶπεν, ἰδοὺ λαὸς καταβαίνων κατὰ θάλασσαν ἀπὸ τοῦ ἐχόμενα ὀμφαλοῦ τῆς γῆς, καὶ ἀρχὴ ἑτέρα ἔρχεται διʼ ὁδοῦ Ἥλων Μαωνενίμ.
\vs{38}Καὶ εἶπε πρὸς αὐτὸν Ζεβοὺλ, καὶ ποῦ ἐστι τὸ στόμα σου ὡς ἐλάλησας, τίς ἐστιν Ἀβιμέλεχ, ὅτι δουλεύσομεν αὐτῷ; μὴ οὐχὶ οὗτος ὁ λαὸς ὃν ἐξουδένωσας; ἔξελθε δὴ νῦν καὶ παράταξαι αὐτῷ.
\vs{39}Καὶ ἐξῆλθε Γαὰλ ἐνώπιον ἀνδρῶν Συχὲμ, καὶ παρετάξατο πρὸς Ἀβιμέλεχ.
\vs{40}Καὶ ἐδίωξεν αὐτὸν Ἀβιμέλεχ, καὶ ἔφυγεν ἀπὸ προσώπου αὐτοῦ· καὶ ἔπεσον τραυματίαι πολλοὶ ἕως τῆς θύρας τῆς πύλης.

\vs{41}Καὶ εἰσῆλθεν Ἀβιμέλεχ ἐν Ἀρημά· καὶ ἐξέβαλε Ζεβοὺλ τὸν Γαὰλ καὶ τοὺς ἀδελφοὺς αὐτοῦ, μὴ οἰκεῖν ἐν Συχέμ.

\vs{42}Καὶ ἐγένετο τῇ ἐπαύριον καὶ ἐξῆλθεν ὁ λαὸς εἰς τὸν ἀγρὸν, καὶ ἀνήγγειλε τῷ Ἀβιμέλεχ.
\vs{43}Καὶ ἔλαβε τὸν λαὸν, καὶ διεῖλεν αὐτοὺς εἰς τρεῖς ἀρχὰς, καὶ ἐνήδρευσεν ἐν ἀγρῷ· καὶ εἶδε, καὶ ἰδοὺ λαὸς ἐξῆλθεν ἐκ τῆς πόλεως, καὶ ἀνέστη ἐπʼ αὐτοὺς, καὶ ἐπάταξεν αὐτούς.
\vs{44}Καὶ Ἀβιμέλεχ καὶ οἱ ἀρχηγοὶ οἱ μετʼ αὐτοῦ ἐξέτειναν, καὶ ἔστησαν παρὰ τὴν θύραν τῆς πύλης τῆς πόλεως· καὶ αἱ δύο ἀρχαὶ ἐξέτειναν ἐπὶ πάντας τοὺς ἐν τῷ ἀγρῷ, καὶ ἐπάταξαν αὐτούς.
\vs{45}Καὶ Ἀβιμέλεχ παρετάσσετο ἐν τῇ πόλει ὅλην τὴν ἡμέραν ἐκείνην, καὶ κατελάβετο τὴν πόλιν, καὶ τὸν λαὸν τὸν ἐν αὐτῇ ἀπέκτεινε, καὶ τὴν πόλιν καθεῖλε, καὶ ἔσπειρεν αὐτὴν ἅλας.

\vs{46}Καὶ ἤκουσαν πάντες οἱ ἄνδρες πύργου Συχὲμ, καὶ ἦλθον εἰς συνέλευσιν Βαιθηλβερίθ.
\vs{47}Καὶ ἀνηγγέλη τῷ Ἀβιμέλεχ, ὅτι συνήχθησαν πάντες οἱ ἄνδρες πύργου Συχέμ.
\vs{48}Καὶ ἀνέβη Ἀβιμέλεχ εἰς ὄρος Σελμὼν, καὶ πᾶς ὁ λαὸς ὁ μετʼ αὐτοῦ· καὶ ἔλαβεν Ἀβιμέλεχ τὰς ἀξίνας ἐν τῇ χειρὶ αὐτοῦ, καὶ ἔκοψε κλάδον ξύλου, καὶ ᾖρε, καὶ ἔθηκεν ἐπὶ ὤμων αὐτοῦ· καὶ εἶπε τῷ λαῷ τῷ μετʼ αὐτοῦ, ὃ εἴδετέ με ποιοῦντα, ταχέως ποιήσατε ὡς ἐγώ.
\vs{49}Καὶ ἔκοψαν καί γε ἀνὴρ κλάδον πᾶς ἀνὴρ, καὶ ἐπορεύθησαν ὀπίσω Ἀβιμέλεχ, καὶ ἐπέθηκαν ἐπὶ τὴν συνέλευσιν, καὶ ἐνεπύρισαν ἐπʼ αὐτοὺς τὴν συνέλευσιν ἐν πυρί· καὶ ἀπέθανον καί γε πάντες οἱ ἄνδρες πύργου Σικίμων, ὡσεὶ χίλιοι ἄνδρες καὶ γυναῖκες.

\vs{50}Καὶ ἐπορεύθη Ἀβιμέλεχ ἐκ Βαιθηγβερὶθ, καὶ παρενέβαλεν ἐν Θήβης, καὶ κατέλαβεν αὐτήν.
\vs{51}Καὶ πύργος ἰσχυρὸς ἦν ἐν μέσῳ τῆς πόλεως· καὶ ἔφυγον ἐκεῖ πάντες οἱ ἄνδρες καὶ αἱ γυναῖκες τῆς πόλεως, καὶ ἔκλεισαν ἔξωθεν αὐτῶν, καὶ ἀνέβησαν ἐπὶ τὸ δῶμα τοῦ πύργου.
\vs{52}Καὶ ἦλθεν Ἀβιμέλεχ ἕως τοῦ πύργου, καὶ παρετάξαντο αὐτῷ· καὶ ἤγγισεν Ἀβιμέλεχ ἕως τῆς θύρας τοῦ πύργου τοῦ ἐμπρῆσαι αὐτὸν ἐν πυρί.
\vs{53}Καὶ ἔῤῥιψε γυνὴ μία κλάσμα ἐπιμύλιον ἐπὶ κεφαλὴν Ἀβιμέλεχ, καὶ ἔκλασε τὸ κρανίον αὐτοῦ.
\vs{54}Καὶ ἐβόησε ταχὺ πρὸς τὸ παιδάριον τὸ αἶρον τὰ σκεύη αὐτοῦ, καὶ εἶπεν αὐτῷ, σπάσον τὴν ῥομφαίαν μου καὶ θανάτωσόν με, μή ποτε εἴπωσι, γυνὴ ἀπέκτεινεν αὐτόν· καὶ ἐξεκέντησεν αὐτὸν τὸ παιδάριον αὐτοῦ, καὶ ἀπέθανε.
\vs{55}Καὶ εἶδεν ἀνὴρ Ἰσραὴλ ὅτι ἀπέθανεν Ἀβιμέλεχ· καὶ ἐπορεύθησαν ἀνὴρ εἰς τὸν τόπον αὐτοῦ.

\vs{56}Καὶ ἐπέστρεψεν ὁ Θεὸς τὴν πονηρίαν Ἀβιμέλεχ, ἣν ἐποίησε τῷ πατρὶ αὐτοῦ, ἀποκτεῖναι τοὺς ἑβδομήκοντα ἀδελφοὺς αὐτοῦ.
\vs{57}Καὶ τὴν πᾶσαν πονηρίαν ἀνδρῶν Συχὲμ ἐπέστρεψεν ὁ Θεὸς εἰς κεφαλὴν αὐτῶν· καὶ ἐπῆλθεν ἐπʼ αὐτοὺς ἡ κατάρα Ἰωάθαμ υἱοῦ Ἰεροβάαλ.

\ch{10}
Καὶ ἀνέστη μετὰ Ἀβιμέλεχ τοῦ σῶσαι τὸν Ἰσραὴλ Θωλὰ υἱὸς Φουὰ, υἱὸς πατραδέλφου αὐτοῦ, ἀνὴρ Ἰσσάχαρ· καὶ αὐτὸς ᾤκει ἐν Σαμὶρ ἐν ὄρει Ἐφραίμ.
\vs{2}Καὶ ἔκρινε τὸν Ἰσραὴλ εἴκοσι τρία ἔτη, καὶ ἀπέθανε, καὶ ἐτάφη ἐν Σαμίρ.

\vs{3}Καὶ ἀνέστη μετʼ αὐτὸν Ἰαῒρ ὁ Γαλαὰδ, καὶ ἔκρινε τὸν Ἰσραὴλ εἴκοσι δύο ἔτη.
\vs{4}Καὶ ἦσαν αὐτῷ τριάκοντα καὶ δύο υἱοὶ ἐπιβαίνοντες ἐπὶ τριάκοντα δύο πώλους· καὶ τριάκοντα δύο πόλεις αὐτοῖς· καὶ ἐκάλουν αὐτὰς ἐπαύλεις Ἰαῒρ ἕως τῆς ἡμέρας ταύτης ἐν γῇ Γαλαάδ.
\vs{5}Καὶ ἀπέθανεν Ἰαῒρ, καὶ ἐτάφη ἐν Ῥαμνών.

\vs{6}Καὶ προσέθεντο οἱ υἱοὶ Ἰσραὴλ τοῦ ποιῆσαι τὸ πονηρὸν ἐνώπιον Κυρίου, καὶ ἐδούλευσαν τοῖς Βααλὶμ, καὶ ταῖς Ἀσταρὼθ, καὶ τοῖς θεοῖς Ἀρὰμ, καὶ τοῖς θεοῖς Σιδῶνος, καὶ τοῖς θεοῖς Μωὰβ, καὶ τοῖς θεοῖς υἱῶν Ἀμμὼν, καὶ τοῖς θεοῖς Φυλιστιῒμ, καὶ ἐγκατέλιπον τὸν Κύριον, καὶ οὐκ ἐδούλευσαν αὐτῷ.
\vs{7}Καὶ ὠργίσθη θυμῷ Κύριος ἐν Ἰσραὴλ, καὶ ἀπέδοτο αὐτοὺς ἐν χειρὶ Φυλιστιῒμ, καὶ ἐν χειρὶ υἱῶν Ἀμμών.
\vs{8}Καὶ ἔθλιψαν καὶ ἔθλασαν τοὺς υἱοὺς Ἰσραὴλ ἐν τῷ καιρῷ ἐκείνῳ ὀκτωκαίδεκα ἔτη, τοὺς πάντας υἱοὺς Ἰσραὴλ τοὺς ἐν τῷ πέραν τοῦ Ἰορδάνου ἐν γῇ τοῦ Ἀμοῤῥὶ τοῦ ἐν Γαλαάδ.
\vs{9}Καὶ διέβησαν οἱ υἱοὶ Ἀμμὼν τὸν Ἰορδάνην παρατάξασθαι πρὸς Ἰούδαν, καὶ Βενιαμὶν, καὶ πρὸς Ἐφραίμ· καὶ ἐθλίβησαν οἱ υἱοὶ Ἰσραὴλ σφόδρα.

\vs{10}Καὶ ἐβόησαν οἱ υἱοὶ Ἰσραὴλ πρὸς Κύριον, λέγοντες, ἡμάρτομέν σοι, ὅτι ἐγκατελίπομεν τὸν Θεὸν, καὶ ἐδουλεύσαμεν τῷ Βααλίμ.
\vs{11}Καὶ εἶπε Κύριος πρὸς τοὺς υἱοὺς Ἰσραὴλ, μὴ οὐχὶ ἐξ Αἰγύπτου, καὶ ἀπὸ τοῦ Ἀμοῤῥαίου, καὶ ἀπὸ υἱῶν Ἀμμὼν, καὶ ἀπὸ Φυλιστιῒμ,
\vs{12}καὶ Σιδωνίων, καὶ Ἀμαλὲκ, καὶ Μαδιὰμ, οἳ ἔθλιψαν ὑμᾶς; καὶ ἐβοήσατε πρὸς μὲ, καὶ ἔσωσα ὑμᾶς ἐκ χειρὸς αὐτῶν;
\vs{13}Καὶ ὑμεῖς ἐγκατελίπετέ με, καὶ ἐδουλεύσατε θεοῖς ἑτέροις· διὰ τοῦτο οὐ προσθήσω τοῦ σῶσαι ὑμᾶς.
\vs{14}Πορεύεσθε, καὶ βοήσατε πρὸς τοὺς θεοὺς οὓς ἐξελέξασθε ἑαυτοῖς, καὶ αὐτοὶ σωσάτωσαν ὑμᾶς ἐν καιρῷ θλίψεως ὑμῶν.
\vs{15}Καὶ εἶπαν οἱ υἱοὶ Ἰσραὴλ πρὸς Κύριον, ἡμάρτομεν, ποίησον σὺ ἡμῖν κατὰ πὰν τὸ ἀγαθὸν ἐν ὀφθαλμοῖς σου, πλὴν ἐξελοῦ ἡμᾶς ἐν τῇ ἡμέρᾳ ταύτῃ.
\vs{16}Καὶ ἐξέκλιναν τοὺς θεοὺς τοὺ ἀλλοτρίους ἐκ μέσου αὐτῶν, καὶ ἐδούλευσαν τῷ Κυρίῳ μόνῳ· καὶ ὠλιγώθη ἡ ψυχὴ αὐτοῦ ἐν κόπῳ Ἰσραήλ.

\vs{17}Καὶ ἀνέβησαν οἱ υἱοὶ Ἀμμὼν, καὶ παρενέβαλον ἐν Γαλαάδ· καὶ συνήχθησαν οἱ υἱοὶ Ἰσραὴλ, καὶ παρενέβαλον ἐν τῇ σκοπίᾳ.
\vs{18}Καὶ εἶπον ὁ λαὸς οἱ ἄρχοντες Γαλαὰδ, ἀνὴρ πρὸς τὸν πλησίον αὐτοῦ, τίς ὁ ἀνὴρ ὅστις ἂν ἄρξεται παρατάξασθαι πρὸς υἱοῖς Ἀμμὼν, καὶ ἔσται εἰς ἄρχοντα πᾶσι τοῖς κατοικοῦσι Γαλαάδ;

\ch{11}
Καὶ Ἰεφθάε ὁ Γαλααδίτης ἐπῃρμένος δυνάμει, καὶ αὐτὸς υἱοῖς γυναικὸς πόρνης, ἣ ἐγέννησε τῷ Γαλαὰδ τὸν Ἰεφθάε.
\vs{2}Καὶ ἔτεκεν ἡ γυνὴ Γαλαὰδ αὐτῷ υἱούς· καὶ ἡδρύνθησαν οἱ υἱοὶ τῆς γυναικὸς, καὶ ἐξέβαλον τὸν Ἰεφθάε, καὶ εἶπαν αὐτῷ, οὐ κληρονομήσεις ἐν τῷ οἴκῳ τοῦ πατρὸς ἡμῶν, ὅτι υἱὸς γυναικὸς ἑταίρας σύ.

\vs{3}Καὶ ἔφυγεν Ἰεφθάε ἀπὸ προσώπου τῶν ἀδελφῶν αὐτοῦ, καὶ ᾤκησεν ἐν γῇ Τώβ· καὶ συνεστράφησαν πρὸς Ἰεφθάε ἄνδρες κενοὶ, καὶ ἐξῆλθον μετʼ αὐτοῦ.

\vs{4}Καὶ ἐγένετο ἡνίκα παρετάξαντο οἱ υἱοὶ Ἀμμῶν μετὰ Ἰσραὴλ,
\vs{5}καὶ ἐπορεύθησαν οἱ πρεσβύτεροι Γαλαὰδ λαβεῖν τὸν Ἰεφθάε ἀπὸ τῆς γῆς Τὼβ,
\vs{6}καὶ εἶπαν τῷ Ἰεφθάε, δεῦρο καὶ ἔσῃ ἡμῖν εἰς ἀρχηγὸν, καὶ παραταξόμεθα πρὸς υἱοὺς Ἀμμών.
\vs{7}Καὶ εἶπεν Ἰεφθάε τοῖς πρεσβυτέροις Γαλαὰδ, οὐχὶ ὑμεῖς ἐμισήσατέ με, καὶ ἐξεβάλετέ με ἐκ τοῦ οἴκου τοῦ πατρός μου, καὶ ἐξαπεστείλατέ με ἀφʼ ὑμῶν; καὶ διατί ἤλθατε πρὸς μὲ νῦν ἡνίκα χρῄζετε;
\vs{8}Καὶ εἶπαν οἱ προσβύτεροι Γαλαὰδ πρὸς Ἰεφθάε, διὰ τοῦτο νῦν ἐπεστρέψαμεν πρὸς σὲ, καὶ πορεύσῃ μεθʼ ἡμῶν, καὶ παρατάξῃ πρὸς υἱοὺς Ἀμμὼν, καὶ ἔσῃ ἡμῖν εἰς ἄρχοντα πᾶσι τοῖς κατοικοῦσι Γαλαάδ.
\vs{9}Καὶ εἶπεν Ἰεφθάε πρὸς τοὺς πρεσβυτέρους Γαλαὰδ, εἰ ἐπιστρέφετέ με ὑμεῖς παρατάξασθαι ἐν υἱοῖς Ἀμμὼν, καὶ παραδῷ αὐτοὺς Κύριος ἐνώπιον ἐμοῦ, καὶ ἐγὼ ὑμῖν ἔσομαι εἰς ἄρχοντα.
\vs{10}Καὶ εἶπαν οἱ πρεσβύτεροι Γαλαὰδ πρὸς Ἰεφθάε, Κύριος ἔστω ἀκούων ἀναμέσον ἡμῶν, εἰ μὴ κατὰ τὸ ῥῆμά σου οὕτω ποιήσομεν.

\vs{11}Καὶ ἐπορεύθη Ἰεφθάε μετὰ τῶν πρεσβυτέρων Γαλαὰδ, καὶ ἔθηκαν αὐτὸν ὁ λαὸς ἐπʼ αὐτοὺς εἰς κεφαλὴν καὶ εἰς ἀρχηγόν· καὶ ἐλάλησεν Ἰεφθαέ πάντας τοὺς λόγους αὐτοῦ ἐνώπιον Κυρίου ἐν Μασσηφά.

\vs{12}Καὶ ἀπέστειλεν Ἰεφθάε ἀγγέλους πρὸς βασιλέα υἱῶν Ἀμμὼν, λέγων, τί ἐμοὶ καὶ σοὶ, ὅτι ἦλθες πρὸς μὲ τοῦ παρατάξασθαι ἐν τῇ γῇ μου;
\vs{13}Καὶ εἶπε βασιλεὺς υἱῶν Ἀμμὼν πρὸς τοὺς ἀγγέλους Ἰεφθάε, ὅτι ἔλαβεν Ἰσραὴλ τὴν γῆν μου ἐν τῷ ἀναβαίνειν αὐτὸν ἐξ Αἰγύπτου ἀπὸ Ἀρνὼν ἕως Ἰαβὸκ καὶ ἕως τοῦ Ἰορδάνου· καὶ νῦν ἐπίστρεψον αὐτὰς ἐν εἰρήνῃ, καὶ πορεύσομαι.

\vs{14}Καὶ προσέθηκεν ἔτι Ἰεφθάε, καὶ ἀπέστειλεν ἀγγέλους πρὸς βασιλέα υἱῶν Ἀμμών.
\vs{15}Καὶ εἶπεν αὐτῷ, οὕτω λέγει Ἰεφθάε, οὐκ ἔλαβεν Ἰσραὴλ τὴν γῆν Μωὰβ, καὶ τὴν γῆν υἱῶν Ἀμμών,
\vs{16}ὅτι ἐν τῷ ἀναβαίνειν αὐτοὺς ἐξ Αἰγύπτου ἐπορεύθη Ἰσραὴλ ἐν τῇ ἐρήμῳ ἕως θαλάσσης Σὶφ, καὶ ἦλθεν εἰς Κάδης.
\vs{17}Καὶ ἀπέστειλεν Ἰσραὴλ ἀγγέλους πρὸς βασιλέα Ἐδὼμ, λέγων, παρελεύσομαι δὴ ἐν τῇ γῇ σου· καὶ οὐκ ἤκουσε βασιλεὺς Ἐδώμ· καί γε πρὸς βασιλέα Μωὰβ ἀπέστειλε, καὶ οὐκ εὐδόκησε· καὶ ἐκάθισεν Ἰσραὴλ ἐν Κάδης,
\vs{18}καὶ ἐπορεύθη ἐν τῇ ἐρήμῳ, καὶ ἐκύκλωσε τὴν γῆν Ἐδὼμ καὶ τὴν γῆν Μωάβ· καὶ ἦλθεν ἀπὸ ἀνατολῶν ἡλίου τῇ γῇ Μωὰβ, καὶ παρενέβαλεν ἐν πέραν Ἀρνὼν, καὶ οὐκ εἰσῆλθεν ἐν ὁρίοις Μωὰβ, ὅτι Ἀρνὼν ὅριον Μωάβ.
\vs{19}Καὶ ἀπέστειλεν Ἰσραὴλ ἀγγέλους πρὸς Σηὼν βασιλέα τοῦ Ἀμοῤῥαίου βασιλέα Ἐσεβῶν, καὶ εἶπεν αὐτῷ Ἰσραήλ, παρέλθωμεν δὴ ἐν τῇ γῇ σου ἕως τοῦ τόπου ἡμῶν.
\vs{20}Καὶ οὐκ ἐνεπίστευσε Σηὼν τῷ Ἰσραὴλ παρελθεῖν ἐν τῷ ὁρίῳ αὐτοῦ· καὶ συνῆξε Σηὼν πάντα τὸν λαὸν αὐτοῦ, καὶ παρενέβαλον εἰς Ἰασὰ, καὶ παρετάξατο πρὸς Ἰσραήλ.
\vs{21}Καὶ παρέδωκε Κύριος ὁ Θεὸς Ἰσραὴλ τὸν Σηὼν καὶ πάντα τὸν λαὸν αὐτοῦ ἐν χειρὶ Ἰσραὴλ, καὶ ἐπάταξεν αὐτόν· καὶ ἐκληρονόμησεν Ἰσραὴλ πᾶσαν τὴν γῆν τοῦ Ἀμοῤῥαίου τοῦ κατοικοῦντος τὴν γῆν ἐκείνην
\vs{22}ἀπὸ Ἀρνὼν καὶ ἕως τοῦ Ἰαβὸκ, καὶ ἀπὸ τοῦ ἐρήμου ἕως τοῦ Ἰορδάνου.
\vs{23}Καὶ νῦν Κύριος ὁ Θεὸς Ἰσραὴλ ἐξῇρε τὸν Ἀμοῤῥαῖον ἀπὸ προσώπου λαοῦ αὐτοῦ Ἰσραήλ, καὶ σὺ κληρονομήσεις αὐτόν;
\vs{24}Οὐχὶ ἃ ἐὰν κληρονομήσει σε Χαμὼς ὁ θεὸς σοῦ, αὐτὰ κληρονομήσεις, καὶ τοὺς πάντας οὓς ἐξῇρε Κύριος ὁ Θεὸς ἡμῶν ἀπὸ προσώπου ἡμῶν, αὐτοὺς κληρονομήσομεν;
\vs{25}Καὶ νῦν μὴ ἐν ἀγαθῷ ἀγαθώτερος σὺ ὑπὲρ Βαλὰκ υἱὸν Σεπφὼρ βασιλέως Μωάβ; μὴ μαχόμενος ἐμαχέσατο μετὰ Ἰσραὴλ, ἢ πολεμῶν ἐπολέμησεν αὐτὸν,
\vs{26}ἐν τῷ οἰκῆσαι ἐν Ἐσεβὼν καὶ ἐν τοῖς ὁρίοις αὐτῆς, καὶ ἐν γῇ Ἀροὴρ καὶ ἐν τοῖς ὁρίοις αὐτῆς, καὶ ἐν πάσαις ταῖς πόλεσι ταῖς παρὰ τὸν Ἰορδάνην, τριακόσια ἔτη; καὶ διατί οὐκ ἐῤῥύσω αὐτοὺς ἐν τῷ καιρῷ ἐκείνῳ;
\vs{27}Καὶ νῦν ἐγώ εἰμι οὐχ ἥμαρτόν σοι, καὶ σὺ ποιεῖς μετʼ ἐμοῦ πονηρίαν τοῦ παρατάξασθαι ἐν ἐμοί· κρίναι Κύριος ὁ κρίνων σήμερον ἀναμέσον υἱῶν Ἰσραὴλ καὶ ἀναμέσον υἱῶν Ἀμμών.

\vs{28}Καὶ οὐκ ἤκουσε βασιλεὺς υἱῶν Ἀμμὼν τῶν λόγων Ἰεφθάε, ὧν ἀπέστειλε πρὸς αὐτόν.
\vs{29}Καὶ ἐγένετο ἐπὶ Ἰεφθάε πνεῦμα Κυρίου, καὶ παρῆλθε τὸν Γαλαὰδ, καὶ τὸν Μανασσῆ, καὶ παρῆλθε τὴν σκοπιὰν Γαλαὰδ εἰς τὸ πέραν υἱῶν Ἀμμών.

\vs{30}Καὶ ηὔξατο Ἰεφθάε εὐχὴν τῷ Κυρίῳ, καὶ εἶπεν, ἐὰν διδοὺς δῷς μοι τοὺς υἱοὺς Ἁμμὼν ἐν τῆ χειρί μου,
\vs{31}καὶ ἔσται ὁ ἐκπορευόμενος ὃς ἂν ἐξέλθῃ ἀπὸ τῆς θύρας τοῦ οἴκου μου εἰς συνάντησίν μου ἐν τῷ ἐπιστρέφειν με ἐν εἰρήνῃ ἀπὸ υἱῶν Ἀμμὼν, καὶ ἔσται τῷ Κυρίῳ, ἀνοίσω αὐτὸν ὁλοκαύτωμα.

\vs{32}Καὶ παρῆλθεν Ἰεφθάε πρὸς υἱοὺς Ἀμμὼν παρατάξασθαι πρὸς αὐτούς· καὶ παρέδωκεν αὐτοὺς Κύριος ἐν χειρὶ αὐτοῦ.
\vs{33}Καὶ ἐπάταξεν αὐτοὺς ἀπὸ Ἀροὴρ ἕως ἐλθεῖν ἄχρις Ἀρνὼν ἐν ἀριθμῷ εἴκοσι πόλεις, καὶ ἕως Ἐβελχαρμὶμ, πληγὴν μεγάλην σφόδρα· καὶ συνεστάλησαν οἱ υἱοὶ Ἀμμὼν ἀπὸ προσώπου υἱῶν Ἰσραήλ.

\vs{34}Καὶ ἦλθεν Ἰεφθάε εἰς Μασσηφὰ εἰς τὸν οἶκον αὐτοῦ· καὶ ἰδοὺ ἡ θυγάτηρ αὐτοῦ ἐξεπορεύετο εἰς ὑπάντησιν ἐν τυμπάνοις καὶ χοροῖς· καὶ αὕτη ἦν μονογενὴς αὐτῷ· οὐκ ἦν αὐτῷ ἕτερος υἱὸς ἢ θυγάτηρ.
\vs{35}Καὶ ἐγένετο ὡς εἶδεν αὐτὴν αὐτός, διέῤῥηξε τὰ ἱμάτια αὐτοῦ, καὶ εἶπεν, ἆ ἆ, θυγάτηρ μου, ταραχῇ ἐτάραξάς με, καὶ σὺ ἦς ἐν τῷ ταράχῳ μου, καὶ ἐγώ εἰμι ἤνοιξα κατὰ σοῦ τὸ στόμα μου πρὸς Κύριον, καὶ οὐ δυνήσομαι ἀποστρέψαι.
\vs{36}Ἡ δὲ εἶπε πρὸς αὐτὸν, πάτερ, ἤνοιξας τὸ στόμα σου πρὸς Κύριον; ποίησόν μοι ὃν τρόπον ἐξῆλθεν ἐκ στόματός σου, ἐν τῷ ποιῆσαί σοι Κύριον ἐκδίκησιν τῶν ἐχθρῶν σου ἀπὸ τῶν υἱῶν Ἀμμών.
\vs{37}Καὶ ἥδε εἶπε πρὸς τὸν πατέρα αὐτῆς, ποιησάτω δὴ ὁ πατήρ μου τὸν λόγον τοῦτον· ἔασόν με δύο μῆνας, καὶ πορεύσομαι καὶ καταβήσομαι ἐπὶ τὰ ὄρη, καὶ κλαύσομαι ἐπὶ τὰ παρθένιά μου ἐγώ εἰμι καὶ αἱ συνεταιρίδες μου.
\vs{38}Καὶ εἶπε, πορεύου· καὶ ἀπέστειλεν αὐτὴν δύο μῆνας· καὶ ἐπορεύθη αὐτὴ καὶ αἱ συνεταιρίδες αὐτῆς, καὶ ἔκλαυσεν ἐπὶ τὰ παρθένια αὐτῆς ἐπὶ τὰ ὄρη.

\vs{39}Καὶ ἐγένετο ἐν τέλει τῶν δύο μηνῶν, καὶ ἐπέστρεψε πρὸς τὸν πατέρα αὐτῆς· καὶ ἐποίησεν ἐν αὐτῇ τὴν εὐχὴν αὐτοῦ ἣν ηὔξατο· καὶ αὕτη οὐκ ἔγνω ἄνδρα· καὶ ἐγένετο εἰς πρόσταγμα ἐν Ἰσραήλ·
\vs{40}Ἀπὸ ἡμερῶν εἰς ἡμέρας ἐπορεύοντο θυγατέρες Ἰσραὴλ θρηνεῖν τὴν θυγατέρα Ἰεφθάε τοῦ Γαλααδίτου ἐπὶ τέσσαρας ἡμέρας ἐν τῷ ἐνιαυτῷ.

\ch{12}
Καὶ ἐβόησεν ἀνὴρ Ἐφραὶμ, καὶ παρῆλθαν εἰς Βοῤῥᾶν, καὶ εἶπαν πρὸς Ἰεφθάε, διατί παρῆλθες παρατάξασθαι ἐν υἱοῖς Ἀμμὼν, καὶ ἡμᾶς οὐ κέκληκας πορευθῆναι μετὰ σοῦ; τὸν οἶκόν σου ἐμπρήσομεν ἐπὶ σὲ ἐν πυρί.
\vs{2}Καὶ εἶπε πρὸς αὐτοὺς Ἰεφθὰε, ἀνὴρ μαχητὴς ἤμην ἐγὼ καὶ ὁ λαός μου, καὶ οἱ υἱοὶ Ἀμμὼν σφόδρα· καὶ ἐβόησα ὑμᾶς, καὶ οὐκ ἐσώσατέ με ἐκ χειρὸς αὐτῶν.
\vs{3}Καὶ εἶδον ὅτι οὐκ εἶ σωτὴρ, καὶ ἔθηκα τὴν ψυχήν μου ἐν χειρί μου, καὶ παρῆλθον πρὸς υἱοὺς Ἀμμὼν, καὶ ἔδωκεν αὐτοὺς Κύριος ἐν χειρί μου· καὶ εἰς τί ἀνέβητε ἐπʼ ἐμὲ ἐν τῇ ἡμέρᾳ ταύτῃ παρατάξασθαι ἐν ἐμοί;

\vs{4}Καὶ συνέστρεψεν Ἰεφθάε πάντας τοὺς ἄνδρας Γαλαὰδ, καὶ παρετάξατο τῷ Ἐφραὶμ, καὶ ἐπάταξαν ἄνδρες Γαλαὰδ τὸν Ἐφραὶμ, ὅτι εἶπαν οἱ διασωζόμενοι τοῦ Ἐφραὶμ, ὑμεῖς Γαλαὰδ ἐν μέσῳ τοῦ Ἐφραὶμ καὶ ἐν μέσῳ τοῦ Μανασσῆ.
\vs{5}Καὶ προκατελάβετο Γαλαὰδ τὰς διαβάσεις τοῦ Ἰορδάνου τοῦ Ἐφραίμ· καὶ εἶπαν αὐτοῖς οἱ διασωζόμενοι Ἐφραὶμ, διαβῶμεν· καὶ εἶπαν αὐτοῖς οἱ ἄνδρες Γαλαὰδ, μὴ Ἐφραθίτης εἶ; καὶ εἶπεν, οὔ.
\vs{6}Καὶ εἶπαν αὐτῷ, εἶπον δὴ στάχυς· καὶ οὐ κατεύθυνε τοῦ λαλῆσαι οὕτως· καὶ ἐπελάβοντο αὐτοῦ, καὶ ἔθυσαν αὐτὸν πρὸς τὰς διαβάσεις τοῦ Ἰορδάνου· καὶ ἔπεσαν ἐν τῷ καιρῷ ἐκείνῳ ἀπὸ Ἐφραὶμ δύο καὶ τεσσαράκοντα χιλιάδες.

\vs{7}Καὶ ἔκρινεν Ἰεφθάε τὸν Ἰσραὴλ ἓξ ἔτη· καὶ ἀπέθανεν Ἰεφθάε ὁ Γαλααδίτης, καὶ ἐτάφη ἐν πόλει αὐτοῦ Γαλαάδ.

\vs{8}Καὶ ἔκρινε μετʼ αὐτὸν τὸν Ἰσραὴλ Ἀβαισσὰν ἀπὸ Βηθλεέμ.
\vs{9}Καὶ ἦσαν αὐτῷ τριάκοντα υἱοὶ, καὶ τριάκοντα θυγατέρες, ἃς ἐξαπέστειλεν ἔξω, καὶ τριάκοντα θυγατέρας εἰσήνεγκε τοῖς υἱοῖς αὐτοῦ ἔξωθεν· καὶ ἔκρινε τὸν Ἰσραὴλ ἑπτὰ ἔτη.
\vs{10}Καὶ ἀπέθανεν Ἀβαισσὰν, καὶ ἐτάφη ἐν Βηθλεέμ.

\vs{11}Καὶ ἔκρινε μετʼ αὐτὸν τὸν Ἰσραὴλ Αἰλὼμ ὁ Ζαβουλωνίτης δέκα ἔτη.
\vs{12}Καὶ ἀπέθανεν Αἰλὼμ ὁ Ζαβουλωνίτης, καὶ ἐτάφη ἐν Αἰλὼμ ἐν γῇ Ζαβουλών.

\vs{13}Καὶ ἔκρινε μετʼ αὐτὸν τὸν Ἰσραὴλ Ἀβδὼν υἱὸς Ἑλλὴλ ὁ Φαραθωνίτης.
\vs{14}Καὶ ἦσαν αὐτῷ τεσσαράκοντα υἱοὶ, καὶ τριάκοντα υἱῶν υἱοὶ ἐπιβαίνοντες ἐπὶ ἑβδομήκοντα πώλους· καὶ ἔκρινε τὸν Ἰσραὴλ ὀκτὼ ἔτη.
\vs{15}Καὶ ἀπέθανεν Ἀβδὼν υἱὸς Ἐλλὴλ ὁ Φαραθωνίτης, καὶ ἐτάφη ἐν Φαραθὼν ἐν γῇ Ἐφραὶμ ἐν ὄρει τοῦ Ἀμαλήκ.

\ch{13}
Καὶ προσέθηκαν ἔτι οἱ υἱοὶ Ἰσραὴλ ποιῆσαι τὸ πονηρὸν ἐνώπιον Κυρίου· καὶ παρέδωκεν αὐτοὺς Κύριος ἐν χειρὶ Φυλιστιῒμ τεσσαράκοντα ἔτη.

\vs{2}Καὶ ἦν ἀνὴρ εἷς ἀπὸ Σαραὰ ἀπὸ δήμου συγγενείας τοῦ Δανὶ, καὶ ὄνομα αὐτῷ Μανωὲ, καὶ γυνὴ αὐτοῦ στεῖρα καὶ οὐκ ἔτεκε.
\vs{3}Καὶ ὤφθη ἄγγελος Κυρίου πρὸς τὴν γυναῖκα, καὶ εἶπε πρὸς αὐτὴν, ἰδοὺ σὺ στεῖρα καὶ οὐ τέτοκας, καὶ συλλήψῃ υἱόν.
\vs{4}Καὶ νῦν φύλαξαι δὴ, καὶ μὴ πίῃς οἶνον καὶ μέθυσμα, καὶ μὴ φάγῃς πᾶν ἀκάθαρτον,
\vs{5}ὅτι ἰδοὺ σὺ ἐν γαστρὶ ἔχεις καὶ τέξῃ υἱόν· καὶ σίδηρος ἐπὶ τὴν κεφαλὴν αὐτοῦ οὐκ ἀναβήσεται, ὅτι Ναζὶρ Θεοῦ ἔσται τὸ παιδάριον ἀπὸ τῆς κοιλίας· καὶ αὐτὸς ᾄρξεται σῶσαι τὸν Ἰσραὴλ ἐκ χειρὸς Φυλιστιΐμ.

\vs{6}Καὶ εἰσῆλθεν ἡ γυνὴ, καὶ εἶπε τῷ ἀνδρὶ αὐτῆς, λέγουσα, ἄνθρωπος Θεοῦ ἦλθε πρὸς μὲ, καὶ εἶδος αὐτοῦ ὡς εἶδος ἀγγέλου Θεοῦ, φοβερὸν σφόδρα· καὶ οὐκ ἠρώτησα αὐτὸν πόθεν ἐστὶ, καὶ τὸ ὄνομα αὐτοῦ οὐκ ἀπήγγειλέ μοι.
\vs{7}Καὶ εἶπέ μοι, ἰδοὺ σὺ ἐν γαστρὶ ἔχεις καὶ τέξῃ υἱόν· καὶ νῦν μὴ πίῃς οἶνον καὶ μέθυσμα, καὶ μὴ φάγῃς πᾶν ἀκάθαρτον, ὅτι Θεοῦ ἅγιον ἔσται τὸ παιδάριον ἀπὸ γαστρὸς ἕως ἡμέρας θανάτου αὐτοῦ.

\vs{8}Καὶ προσηύξατο Μανωὲ πρὸς Κύριον, καὶ εἶπεν, ἐν ἐμοὶ Κύριε ἀδωναϊὲ τὸν ἄνθρωπον τοῦ Θεοῦ ὃν ἀπέστειλας· ἐλθέτω δὴ ἔτι πρὸς ἡμᾶς, καὶ συμβιβασάτω ἡμᾶς τί ποιήσωμεν τῷ παιδίῳ τῷ τικτομένῳ.

\vs{9}Καὶ εἰσήκουσεν ὁ Θεὸς τῆς φωνῆς Μανωὲ, καὶ ἦλθεν ὁ ἄγγελος τοῦ Θεοῦ ἔτι πρὸς τὴν γυναῖκα· καὶ αὕτη ἐκάθητο ἐν ἀγρῷ, καὶ Μανωὲ ὁ ἀνὴρ αὐτῆς οὐκ ἦν μετʼ αὐτῆς.
\vs{10}Καὶ ἐτάχυνεν ἡ γυνὴ καὶ ἔδραμε καὶ ἀνήγγειλε τῷ ἀνδρὶ αὐτῆς, καὶ εἶπε πρὸς αὐτὸν, ἰδοὺ ὦπται πρὸς μὲ ὁ ἀνὴρ ὃς ἦλθεν ἐν ἡμέρᾳ πρὸς μέ.

\vs{11}Καὶ ἀνέστη καὶ ἐπορεύθη Μανωὲ ὀπίσω τῆς γυναικὸς αὐτοῦ, καὶ ἦλθε πρὸς τὸν ἄνδρα, καὶ εἶπεν αὐτῷ, εἰ σὺ εἶ ὁ ἀνὴρ, ὁ λαλήσας πρὸς τὴν γυναῖκα; καὶ εἶπεν ὁ ἄγγελος, ἐγώ.
\vs{12}Καὶ εἶπε Μανωὲ, νῦν ἐλεύσεται ὁ λόγος· τίς ἔσται κρίσις τοῦ παιδίου καὶ τὰ ποιήματα αὐτοῦ;
\vs{13}Καὶ εἶπεν ὁ ἄγγελος Κυρίου πρὸς Μανωὲ, ἀπὸ πάντων ὧν εἴρηκα πρὸς τὴν γυναῖκα, φυλάξεται·
\vs{14}Ἀπὸ παντὸς ὃ ἐκπορεύεται ἐξ ἀμπέλου τοῦ οἴνου, οὐ φάγεται, καὶ οἶνον καὶ μέθυσμα μὴ πιέτω, καὶ πᾶν ἀκάθαρτον μὴ φαγέτω· πάντα ὅσα ἐνετειλάμην αὐτῇ, φυλάξεται.

\vs{15}Καὶ εἶπε Μανωὲ πρὸς τὸν ἄγγελον Κυρίου, κατάσχωμεν ὧδέ σε, καὶ ποιήσωμεν ἐνώπιόν σου ἔριφον αἰγῶν.
\vs{16}Καὶ εἶπεν ὁ ἄγγελος Κυρίου πρὸς Μανωὲ, ἐὰν κατάσχῃς, οὐ φάγομαι ἀπὸ τῶν ἄρτων σου· καὶ ἐὰν ποιήσῃς ὁλοκαύτωμα, τῷ Κυρίῳ ἀνοίσεις αὐτό· ὅτι οὐκ ἔγνω Μανωὲ, ὅτι ἄγγελος Κυρίου αὐτός.
\vs{17}Καὶ εἶπε Μανωὲ πρὸς τὸν ἄγγελον Κυρίου, τί τὸ ὄνομά σοι, ὅτι ἔλθοι τὸ ῥῆμά σου, καὶ δοξάσομέν σε;
\vs{18}Καὶ εἶπεν αὐτῷ ὁ ἄγγελος Κυρίου, εἰς τί τοῦτο ἐρωτᾷς τὸ ὄνομά μου; καὶ αὐτό ἐστι θαυμαστόν.
\vs{19}Καὶ ἔλαβε Μανωὲ τὸν ἔριφον τῶν αἰγῶν καὶ τὴν θυσίαν, καὶ ἀνήνεγκεν ἐπὶ τὴν πέτραν τῷ Κυρίῳ· καὶ διεχώρισε ποιῆσαι, καὶ Μανωὲ καὶ ἡ γυνὴ αὐτοῦ βλέποντες.
\vs{20}Καὶ ἐγένετο ἐν τῷ ἀναβῆναι τὴν φλόγα ἐπάνω τοῦ θυσιαστηρίου ἕως τοῦ οὐρανοῦ, καὶ ἀνέβη ὁ ἄγγελος Κυρίου ἐν τῇ φλογί. καὶ Μανωὲ καὶ ἡ γυνὴ αὐτοῦ βλέποντες, καὶ ἔπεσον ἐπὶ πρόσωπον αὐτῶν ἐπὶ τὴν γῆν.
\vs{21}Καὶ οὐ προσέθηκεν ἔτι ὁ ἄγγελος Κυρίου ὀφθῆναι πρὸς Μανωὲ καὶ πρὸς τὴν γυναῖκα αὐτοῦ· τότε ἔγνω Μανωὲ, ὅτι ἄγγελος Κυρίου οὗτος.
\vs{22}Καὶ εἶπε Μανωὲ πρὸς τὴν γυναῖκα αὐτοῦ, θανάτῳ ἀποθανούμεθα ὅτι Θεὸν εἴδομεν.
\vs{23}Καὶ εἶπεν αὐτῷ ἡ γυνὴ αὐτοῦ, εἰ ἤθελεν ὁ Κύριος θανατῶσαι ἡμᾶς, οὐκ ἂν ἔλαβεν ἐκ χειρὸς ἡμῶν ὁλοκαύτωμα καὶ θυσίαν, καὶ οὐκ ἂν ἔδειξεν ἡμῖν ταῦτα πάντα, καὶ καθὼς καιρός οὐκ ἂν ἠκούτισεν ἡμᾶς ταῦτα.

\vs{24}Καὶ ἔτεκεν ἡ γυνὴ υἱὸν, καὶ ἐκάλεσε τὸ ὄνομα αὐτοῦ, Σαμψών· καὶ ἡδρύνθη τὸ παιδάριον, καὶ εὐλόγησεν αὐτὸ Κύριος.
\vs{25}Καὶ ἤρξατο πνεῦμα Κυρίου συνεκπορεύεσθαι αὐτῷ ἐν παρεμβολῇ Δὰν, καὶ ἀναμέσον Σαραὰ καὶ ἀναμέσον Ἐσθαόλ.

\ch{14}
Καὶ κατέβη Σαμψὼν εἰς Θαμναθὰ, καὶ εἶδε γυναῖκα ἐν Θαμναθὰ ἀπὸ τῶν θυγατέρων τῶν ἀλλοφύλων.
\vs{2}Καὶ ἀνέβη καὶ ἀπήγγειλε τῷ πατρὶ αὐτοῦ καὶ τῇ μητρὶ αὐτοῦ, καὶ εἶπε, γυναῖκα ἑώρακα ἐν Θαμναθὰ ἀπὸ τῶν θυγατέρων Φυλιστιῒμ, καὶ νῦν λάβετε αὐτήν μοι εἰς γυναῖκα.
\vs{3}Καὶ εἶπεν αὐτῷ ὁ πατὴρ αὐτοῦ, καὶ ἡ μήτηρ αὐτοῦ, μὴ οὐκ εἰσὶ θυγατέρες τῶν ἀδελφῶν σου, καὶ ἐκ παντὸς τοῦ λαοῦ μου γυνὴ, ὅτι συ πορεύῃ λαβεῖν γυναῖκα ἀπὸ τῶν ἀλλοφύλων τῶν ἀπεριτμήτων;

Καὶ εἶπε Σαμψὼν πρὸς τὸν πατέρα αὐτοῦ, ταύτην λάβε μοι, ὅτι αὕτη εὐθεῖα ἐν ὀφθαλμοῖς μου.
\vs{4}Καὶ ὁ πατὴρ αὐτοῦ καὶ ἡ μήτηρ αὐτοῦ οὐκ ἔγνωσαν ὅτι παρὰ Κυρίου ἐστὶν, ὅτι ἐκδίκησιν αὐτὸς ζητεῖ ἐκ τῶν ἀλλοφύλων· καὶ ἐν τῷ καιρῷ ἐκείνῳ οἱ ἀλλόφυλοι κυριεύοντες ἐν Ἰσραήλ.
\vs{5}Καὶ κατέβη Σαμψὼν καὶ ὁ πατὴρ αὐτοῦ καὶ ἡ μήτηρ αὐτοῦ εἰς Θαμναθά· καὶ ἦλθεν ἕως τοῦ ἀμπελῶνος Θαμναθὰ, καὶ ἰδοὺ σκύμνος λέοντος ὠρυόμενος εἰς συνάντησιν αὐτοῦ.
\vs{6}Καὶ ἥλατο ἐπʼ αὐτὸν πνεῦμα Κυρίου, καὶ συνέτριψεν αὐτὸν ὡσεὶ συντρίψει ἔριφον αἰγῶν, καὶ οὐδὲν ἦν ἐν ταῖς χερσὶν αὐτοῦ· καὶ οὐκ ἀπήγγειλε τῷ πατρὶ αὐτοῦ καὶ τῇ μητρὶ αὐτοῦ ὃ ἐποίησε.
\vs{7}Καὶ κατέβησαν καὶ ἐλάλησαν τῇ γυναικὶ, καὶ ηὐθύνθη ἐν ὀφθαλμοῖς Σαμψών.

\vs{8}Καὶ ὑπέστρεψε μεθʼ ἡμέρας λαβεῖν αὐτὴν, καὶ ἐξέκλινεν ἰδεῖν τὸ πτῶμα τοῦ λέοντος, καὶ ἰδοὺ συναγωγὴ μελισσῶν ἐν τῷ στόματι τοῦ λέοντος καὶ μέλι.
\vs{9}Καὶ ἐξεῖλεν αὐτὸ εἰς χεῖρας αὐτοῦ, καὶ ἐπορεύετο πορευόμενος καὶ ἐσθίων· καὶ ἐπορεύθη πρὸς τὸν πατέρα αὐτοῦ καὶ πρὸς τὴν μητέρα αὐτοῦ, καὶ ἔδωκεν αὐτοῖς καὶ ἔφαγον, καὶ οὐκ ἀπήγγειλεν αὐτοῖς ὅτι ἀπὸ στόματος τοῦ λέοντος ἐξεῖλε τὸ μέλι.

\vs{10}Καὶ κατέβη ὁ πατὴρ αὐτοῦ πρὸς τὴν γυναῖκα, καὶ ἐποίησεν ἐκεῖ Σαμψὼν πότον ἡμέρας ἑπτὰ, ὅτι οὕτως ποιοῦσιν οἱ νεανίσκοι.
\vs{11}Καὶ ἐγένετο ὅτε εἶδον αὐτὸν, καὶ ἔλαβον τριάκοντα κλητοὺς, καὶ ἦσαν μετʼ αὐτοῦ.

\vs{12}Καὶ εἶπεν αὐτοῖς Σαμψὼν, πρόβλημα ὑμῖν προβάλλομαι, ἐὰν ἀπαγγέλλοντες ἀπαγγείλητε αὐτὸ ἐν ταῖς ἑπτὰ ἡμέραις τοῦ πότου καὶ εὕρητε, δώσω ὑμῖν τριάκοντα σινδόνας καὶ τριάκοντα στολὰς ἱματίων.
\vs{13}Καὶ ἐὰν μὴ δύνησθε ἀπαγγεῖλαί μοι, δώσετε ὑμεῖς ἐμοὶ τριάκοντα ὀθόνια καὶ τριάκοντα ἀλλασσομένας στολὰς ἱματίων· καὶ εἶπαν αὐτῷ, προβάλου τὸ πρόβλημά σου, καὶ ἀκουσόμεθα αὐτό.
\vs{14}Καὶ εἶπεν αὐτοῖς, τὶ βρωτὸν ἐξῆλθεν ἐκ βιβρώσκοντος, καὶ ἀπὸ ἰσχυροῦ γλυκύ· καὶ οὐκ ἠδύναντο ἀπαγγεῖλαι τὸ πρόβλημα ἐπὶ τρεῖς ἡμέρας.

\vs{15}Καὶ ἐγένετο ἐν τῇ ἡμέρᾳ τῇ τετάρτῃ, καὶ εἶπαν τῇ γυναικὶ Σαμψὼν, ἀπάτησον δὴ τὸν ἄνδρα σου, καὶ ἀπαγγειλάτω σοι τὸ πρόβλημα, μή ποτε κατακαύσωμέν σε καὶ τὸν οἶκον τοῦ πατρός σου ἐν πυρί· ἢ ἐκβιᾶσαι ἡμᾶς κεκλήκατε;
\vs{16}Καὶ ἔκλαυσεν ἡ γυνὴ Σαμψὼν πρὸς αὐτὸν, καὶ εἶπε, πλὴν μεμίσηκάς με καὶ οὐκ ἠγάπησάς με, ὅτι τὸ πρόβλημα ὃ προεβάλου τοῖς υἱοῖς τοῦ λαοῦ μου, οὐκ ἀπήγγειλάς μοι αὐτό· καὶ εἶπεν αὐτῇ Σαμψὼν, εἰ τῷ πατρί μου καὶ τῇ μητρί μου οὐκ ἀπήγγελκα, σοὶ ἀπαγγείλω;
\vs{17}Καὶ ἔκλαυσε πρὸς αὐτὸν ἐπὶ τὰς ἑπτὰ ἡμέρας, ἃς ἦν αὐτοῖς ὁ πότος· καὶ ἐγένετο ἐν τῇ ἡμέρᾳ τῇ ἑβδόμῃ, καὶ ἀπήγγειλεν αὐτῇ, ὅτι παρηνόχλησεν αὐτῷ· καὶ αὐτὴ ἀπήγγειλε τοῖς υἱοῖς τοῦ λαοῦ αὐτῆς.
\vs{18}Καὶ εἶπαν αὐτῷ οἱ ἄνδρες τῆς πόλεως ἐν τῇ ἡμέρᾳ τῇ ἑβδόμῃ πρὸ τοῦ ἀνατεῖλαι τὸν ἥλιον, τί γλυκύτερον μέλιτος, καὶ τί ἰσχυρότερον λέοντος; καὶ εἶπεν αὐτοῖς Σαμψὼν, εἰ μὴ ἠροτριάσατε ἐν τῇ δαμάλει μου, οὐκ ἂν ἔγνωτε τὸ πρόβλημά μου.
\vs{19}Καὶ ἥλατο ἐπʼ αὐτὸν πνεῦμα Κυρίου, καὶ κατέβη εἰς Ἀσκάλωνα, καὶ ἐπάταξεν ἐξ αὐτῶν τριάκοντα ἄνδρας, καὶ ἔλαβε τὰ ἱμάτια αὐτῶν, καὶ ἔδωκε τὰς στολὰς τοῖς ἀπαγγείλασι τὸ πρόβλημα· Καὶ ὠργίσθη θυμῷ Σαμψὼν, καὶ ἀνέβη εἰς τὸν οἶκον τοῦ πατρὸς αὐτοῦ.
\vs{20}Καὶ ἐγένετο ἡ γυνὴ Σαμψὼν ἑνὶ τῶν φίλων αὐτοῦ, ὧν ἐφιλίασε.

\ch{15}
Καὶ ἐγένετο μεθʼ ἡμέρας ἐν ἡμέραις θερισμοῦ πυρῶν, καὶ ἐπεσκέψατο Σαμψὼν τὴν γυναῖκα αὐτοῦ ἐν ἐρίφῳ αἰγῶν, καὶ εἶπεν, εἰσελεύσομαι πρὸς τὴν γυναῖκά μου καὶ εἰς τὸ ταμεῖον· καὶ οὐκ ἔδωκεν αὐτὸν ὁ πατὴρ αὐτῆς εἰσελθεῖν.
\vs{2}Καὶ εἶπεν ὁ πατὴρ αὐτὴς, λέγων, εἶπα ὅτι μισῶν ἐμίσησας αὐτὴν, καὶ ἔδωκα αὐτὴν ἑνὶ τῶν ἐκ τῶν φίλων σου· μὴ οὐχὶ ἡ ἀδελφὴ αὐτῆς ἡ νεωτέρα ἀγαθωτέρα ὑπὲρ αὐτήν; ἔστω δή σοι ἀντὶ αὐτῆς.

\vs{3}Καὶ εἶπεν αὐτοῖς Σαμψὼν, ἠθὼωμαι καὶ τὸ ἅπαξ ἀπὸ ἀλλοφύλων, ὅτι ποιῶ ἐγὼ μετʼ αὐτῶν πονηρίαν·
\vs{4}Καὶ ἐπορεύθη Σαμψὼν, καὶ συνέλαβε τριακοσίας ἀλώπεκας, καὶ ἔλαβε λαμπάδας, καὶ ἐπέστρεψε κέρκον πρὸς κέρκον, καὶ ἔθηκε λαμπάδα μίαν ἀναμέσον τῶν δύο κέρκων καὶ ἔδησε,
\vs{5}καὶ ἐξέκαυσε πῦρ ἐν ταῖς λαμπάσι, καὶ ἐξαπέστειλεν ἐν τοῖς στάχυσι τῶν ἀλλοφύλων· καὶ ἐκάησαν ἀπὸ ἅλωνος καὶ ἕως σταχύων ὀρθῶν, καὶ ἕως ἀμπελῶνος καὶ ἐλαίας.
\vs{6}Καὶ εἶπαν οἱ ἀλλόφυλοι, τίς ἐποίησε ταῦτα; καὶ εἶπαν, Σαμψὼν ὁ νυμφίος τοῦ Θαμνὶ, ὅτι ἔλαβε τὴν γυναῖκα αὐτοῦ, καὶ ἔδωκεν αὐτὴν τῷ ἐκ τῶν φίλων αὐτοῦ· καὶ ἀνέβησαν οἱ ἀλλόφυλοι, καὶ ἐνέπρησαν αὐτὴν καὶ τὸν οἶκον τοῦ πατρὸς αὐτῆς ἐν πυρί.

\vs{7}Καὶ εἶπεν αὐτοῖς Σαμψὼν, ἐὰν ποιήσητε οὕτως ταύτην, ὅτι ἦ μὴν ἐκδικήσω ἐν ὑμῖν, καὶ ἔσχατον κοπάσω.
\vs{8}Καὶ ἐπάταξεν αὐτοὺς κνήμην ἐπὶ μηρὸν πληγὴν μεγάλην· καὶ κατέβη καὶ ἐκάθισεν ἐν τρυμαλιᾷ τῆς πέτρας Ἠτάμ.

\vs{9}Καὶ ἀνέβησαν οἱ ἀλλόφυλοι, καὶ παρενέβαλον ἐν Ἰούδα, καὶ ἐξεῤῥίφησαν ἐν Λεχί.
\vs{10}Καὶ εἶπαν ἀνὴρ Ἰούδα, εἰς τί ἀνέβητε ἐφʼ ἡμᾶς; καὶ εἶπον οἱ ἀλλόφυλοι, δῆσαι τὸν Σαμψὼν ἀνέβημεν, καὶ ποιῆσαι αὐτῷ ὃν τρόπον ἐποίησεν ἡμῖν.
\vs{11}Καὶ κατέβησαν τρισχίλιοι ἀπὸ Ἰούδα ἄνδρες εἰς τρυμαλιὰν πέτρας Ἠτὰμ, καὶ εἶπαν πρὸς Σαμψὼν, οὐκ οἶδας ὅτι κυριεύουσιν οἱ ἀλλόφυλοι ἡμῶν; καὶ τί τοῦτο ἐποίησας ἡμῖν; καὶ εἶπεν αὐτοῖς Σαμψὼν, ὃν τρόπον ἐποίησάν μοι, οὕτως ἐποίησα αὐτοῖς·
\vs{12}Καὶ εἶπαν αὐτῷ, δῆσαί σε κατέβημεν τοῦ δοῦναί σε ἐν χειρὶ ἀλλοφύλων· καὶ εἶπεν αὐτοῖς Σαμψὼν, ὀμόσατέ μοι μή ποτε συναντήσητε ἐν ἐμοὶ ὑμεῖς.
\vs{13}Καὶ εἶπον αὐτῷ, λέγοντες, οὐχὶ, ὅτι ἀλλʼ ἢ δεσμῷ δήσομέν σε, καὶ παραδώσωμέν σε ἐν χειρὶ αὐτῶν, καὶ θανάτῳ οὐ θανατώσωμέν σε· καὶ ἔδησαν αὐτὸν ἐν δυσὶ καλωδίοις καινοῖς, καὶ ἀνήνεγκαν αὐτὸν ἀπὸ τῆς πέτρας ἐκείνης.

\vs{14}Καὶ ἦλθον ἕως σιαγόνος· καὶ οἱ ἀλλόφυλοι ἠλάλαξαν, καὶ ἔδραμον εἰς συνάντησιν αὐτοῦ· καὶ ἥλατο ἐπʼ αὐτὸν πνεῦμα Κυρίου· καὶ ἐγενήθη τὰ καλώδια τὰ ἐπὶ βραχίοσιν αὐτοῦ ὡσεὶ στυππίον ὃ ἐξεκαύθη ἐν πυρί· καὶ ἐτάκησαν δεσμοὶ αὐτοῦ ἀπὸ χειρῶν αὐτοῦ.
\vs{15}Καὶ εὗρε σιαγόνα ὄνου ἐξεῤῥιμμένην, καὶ ἐξέτεινε τὴν χεῖρα αὐτοῦ καὶ ἔλαβεν αὐτὴν, καὶ ἐπάταξεν ἐν αὐτῇ χιλίους ἄνδρας.
\vs{16}Καὶ εἶπε Σαμψὼν, ἐν σιαγόνι ὄνου ἐξαλείφων ἐξήλειψα αὐτοὺς, ὅτι ἐν τῇ σιαγόνι τοῦ ὄνου ἐπάταξα χιλίους ἄνδρας.
\vs{17}Καὶ ἐγένετο ὡς ἐπαύσατο λαλῶν, καὶ ἔῤῥιψε τὴν σιαγόνα ἐκ τῆς χειρὸς αὐτοῦ· καὶ ἐκάλεσε τὸν τόπον ἐκεῖνον, ἀναίρεσις σιαγόνος.

\vs{18}Καὶ ἐδίψησε σφόδρα, καὶ ἔκλαυσε πρὸς Κύριον, καὶ εἶπε, σὺ εὐδόκησας ἐν χειρὶ δούλου σου τὴν σωτηρίαν τὴν μεγάλην ταύτην, καὶ νῦν ἀποθανοῦμαι τῷ δίψει, καὶ ἐμπεσοῦμαι ἐν χειρὶ τῶν ἀπεριτμήτων;
\vs{19}Καὶ ἔῤῥηξεν ὁ Θεὸς τὸν λάκκον τὸν ἐν τῇ σιαγόνι, καὶ ἐξῆλθεν ἐκ αὐτοῦ ὕδωρ, καὶ ἔπιε· καὶ ἐπέστρεψε τὸ πνεῦμα αὐτοῦ καὶ ἔζησε. διὰ τοῦτο ἐκλήθη τὸ ὄνομα αὐτῆς, Πηγὴ τοῦ ἐπικαλουμένου, ἥ ἐστιν ἐν σιαγόνι, ἕως τῆς ἡμέρας ταύτης.

\vs{20}Καὶ ἔκρινε τὸν Ἰσραὴλ ἐν ἡμέραις ἀλλοφύλων εἴκοσι ἔτη.

\ch{16}
Καὶ ἐπορεύθη Σαμψὼν εἰς Γάζαν, καὶ εἶδεν ἐκεῖ γυναῖκα πόρνην, καὶ εἰσῆλθε πρὸς αὐτήν.
\vs{2}Καὶ ἀνηγγέλη τοῖς Γαζαίοις, λέγοντες, ἥκει Σαμψὼν ὧδε. καὶ ἐκύκλωσαν, καὶ ἐνήδρευσαν ἐπʼ αὐτὸν ὅλην τὴν νύκτα ἐν τῇ πύλῃ τῆς πόλεως· καὶ ἐκώφευσαν ὅλην τὴν νύκτα, λέγοντες, ἕως διαφαύσῃ ὁ ὄρθρος, καὶ φονεύσωμεν αὐτόν.
\vs{3}Καὶ ἐκοιμήθη Σαμψὼν ἕως μεσονυκτίου, καὶ ἀνέστη ἐν ἡμίσει τῆς νυκτὸς, καὶ ἐπελάβετο τῶν θυρῶν τῆς πύλης τῆς πόλεως σὺν τοῖς δυσὶ σταθμοῖς, καὶ ἀνεβάσταζεν αὐτὰς σὺν τῷ μοχλῷ, καὶ ἔθηκεν ἐπὶ ὤμων αὐτοῦ· καὶ ἀνέβη ἐπὶ τὴν κορυφὴν τοῦ ὄρους τοῦ ἐπὶ προσώπου τοῦ Χεβρῶν, καὶ ἔθηκεν αὐτὰ ἐκεῖ.

\vs{4}Καὶ ἐγένετο μετὰ τοῦτο, καὶ ἠγάπησε γυναῖκα ἐν Ἀλσωρήχ· καὶ ὄνομα αὐτῇ Δαλιδά.
\vs{5}Καὶ ἀνέβησαν πρὸς αὐτὴν οἱ ἄρχοντες τῶν ἀλλοφύλων, καὶ εἶπαν αὐτῇ, ἀπάτησον αὐτὸν, καὶ ἴδε ἐν τίνι ἡ ἰσχὺς αὐτοῦ ἡ μεγάλη, καὶ ἐν τίνι δυνησόμεθα αὐτῷ, καὶ δήσομεν αὐτὸν τοῦ ταπεινῶσαι αὐτόν· καὶ ἡμεῖς δώσομέν σοι ἀνὴρ χιλίους καὶ ἑκατὸν ἀργυρίου.

\vs{6}Καὶ εἶπε Δαλιδὰ πρὸς Σαμψὼν, ἀπάγγειλον δή μοι ἐν τίνι ἡ ἰσχύς σου ἡ μεγάλη, καὶ ἐν τίνι δεθήσῃ τοῦ ταπεινωθῆναί σε.
\vs{7}Καὶ εἶπε πρὸς αὐτὴν Σαμψὼν, ἐὰν δήσωσί με ἐν ἑπτὰ νευραῖς ὑγραῖς μὴ διεφθαρμέναις, καὶ ἀσθενήσω καὶ ἔσομαι ὡς εἷς τῶν ἀνθρώπων.
\vs{8}Καὶ ἀνήνεγκαν αὐτῇ οἱ ἄρχοντες τῶν ἀλλοφύλων ἑπτὰ νευρὰς ὑγρὰς μὴ διεφθαρμένας, καὶ ἔδησεν αὐτὸν ἐν αὐταῖς.
\vs{9}Καὶ τὸ ἔνεδρον αὐτῇ ἐκάθητο ἐν τῷ ταμείῳ· καὶ εἶπεν αὐτῷ, ἀλλόφυλοι ἐπὶ σὲ Σαμψών· καὶ διέσπασε τὰς νευρὰς ὡς εἴ τις ἀποσπάσοι στρέμμα στυππίου ἐν τῷ ὀσφρανθῆναι αὐτὸ πυρὸς, καὶ οὐκ ἐγνώσθη ἡ ἰσχὺς αὐτοῦ.

\vs{10}Καὶ εἶπε Δαλιδὰ πρὸς Σαμψὼν, ἰδοὺ ἐπλάνησάς με, καὶ ἐλάλησας πρὸς μὲ ψευδῆ· νῦν οὖν ἀνάγγειλόν μοι ἐν τίνι δεθήσῃ.
\vs{11}Καὶ εἶπε πρὸς αὐτὴν, ἐὰν δεσμεύοντες δήσωσί με ἐν καλωδίοις καινοῖς οἷς οὐκ ἐγένετο ἐν αὐτοῖς ἔργον, καὶ ἀσθενήσω καὶ ἔσομαι ὡς εἶς τῶν ἀνθρώπων.
\vs{12}Καὶ ἔλαβε Δαλιδὰ καλώδια καινὰ, καὶ ἔδησεν αὐτὸν ἐν αὐτοῖς, καὶ τὰ ἔνεδρα ἐξῆλθεν ἐκ τοῦ ταμείου· καὶ εἶπεν, ἀλλόφυλοι ἐπὶ σὲ Σαμψών· καὶ διέσπασεν αὐτὰ ἀπὸ βραχιόνων αὐτοῦ ὡσεὶ σπαρτίον.

\vs{13}Καὶ εἶπε Δαλιδὰ πρὸς Σαμψὼν, ἰδοὺ ἐπλάνησάς με, καὶ ἐλάλησας πρὸς μὲ ψευδῆ· ἀνάγγειλον δή μοι ἐν τίνι δεθήσῃ· καὶ εἶπε πρὸς αὐτὴν, ἐὰν ὑφάνῃς τὰς ἑπτὰ σειρὰς τῆς κεφαλῆς μου σὺν τῷ διάσματι, καὶ ἐγκρούσῃς τῷ πασσάλῳ εἰς τὸν τοῖχον, καὶ ἔσομαι ὡς εἷς τῶν ἀνθρώπων ἀσθενής.
\vs{14}Καὶ ἐγένετο ἐν τῷ κοιμᾶσθαι αὐτὸν, καὶ ἔλαβε Δαλιδὰ τὰς ἑπτὰ σειρὰς τῆς κεφαλῆς αὐτοῦ, καὶ ὕφανεν ἐν τῷ διάσματι, καὶ ἔπηξεν τῷ πασσάῳ εἰς τὸν τοῖχον, καὶ εἶπεν, ἀλλόφυλοι ἐπὶ σὲ Σαμψών· καὶ ἐξυπνίσθη ἀπὸ τοῦ ὕπνου αὐτοῦ, καὶ ἐξῇρε τὸν πάσσαλον τοῦ ὑφάσματος ἐκ τοῦ τοίχου.

\vs{15}Καὶ εἶπε πρὸς Σαμψὼν Δαλιδὰ, πῶς λέγεις, ἠγάπηκά σε, καὶ ἡ καρδία σου οὐκ ἔστι μετʼ ἐμοῦ; τοῦτο τρίτον ἐπλάνησάς με καὶ οὐκ ἀπήγγειλάς μοι ἐν τίνι ἡ ἰσχύς σου ἡ μεγάλη.
\vs{16}Καὶ ἐγένετο ὅτε ἐξέθλιψεν αὐτὸν ἐν λόγοις αὐτῆς πάσας τὰς ἡμέρας, καὶ ἐστενοχώρησεν αὐτὸν, καὶ ὠλιγοψύχησεν ἕως τοῦ ἀποθανεῖν.
\vs{17}Καὶ ἀνήγγειλεν αὐτῇ πᾶσαν τὴν καρδίαν αὐτοῦ, καὶ εἶπεν αὐτῇ, σίδηρος οὐκ ἀνέβη ἐπὶ τὴν κεφαλήν μου, ὅτι ἅγιος Θεοῦ ἐγώ εἰμι ἀπὸ κοιλίας μητρός μου· ἐὰν οὖν ξυρήσωμαι, ἀποστήσεται ἀπʼ ἐμοῦ ἡ ἰσχύς μου καὶ ἀσθενήσω, καὶ ἔσομαι ὡς πάντες οἱ ἄνθρωποι.

\vs{18}Καὶ εἶδε Δαλιδὰ, ὅτι ἀπήγγειλεν αὐτῇ πᾶσαν τὴν καρδίαν αὐτοῦ· καὶ ἀπέστειλε καὶ ἐκάλεσε τοὺς ἄρχοντας τῶν ἀλλοφύλων, λέγουσα, ἀνάβητε ἔτι τὸ ἅπαξ τοῦτο, ὅτι ἀπήγγειλέ μοι πᾶσαν τὴν καρδίαν αὐτοῦ· καὶ ἀνέβησαν πρὸς αὐτὴν οἱ ἄρχοντες τῶν ἀλλοφύλων, καὶ ἀνήνεγκαν τὸ ἀργύριον ἐν χερσὶν αὐτῶν.
\vs{19}Καὶ ἐκοίμισε Δαλιδὰ τὸν Σαμψὼν ἐπὶ τὰ γόνατα αὐτῆς· καὶ ἐκάλεσεν ἄνδρα, καὶ ἐξύρησε τὰς ἑπτὰ σειρὰς τῆς κεφαλῆς αὐτοῦ, καὶ ἤρξατο ταπεινῶσαι αὐτὸν, καὶ ἀπέστη ἡ ἰσχὺς αὐτοῦ ἀπʼ αὐτοῦ.
\vs{20}Καὶ εἶπε Δαλιδά, ἀλλόφυλοι ἐπὶ σὲ Σαμψών. καὶ ἐξυπνίσθη ἐκ τοῦ ὕπνου αὐτοῦ, καὶ εἶπεν, ἐξελεύσομαι ὡς ἅπαξ καὶ ἅπαξ, καὶ ἐκτιναχθήσομαι· καὶ αὐτὸς οὐκ ἔγνω ὅτι ὁ Κύριος ἀπέστη ἀπάνωθεν αὐτοῦ.
\vs{21}Καὶ ἐκράτησαν αὐτὸν οἱ ἀλλόφυλοι, καὶ ἐξέκοψαν τοὺς ὀφθαλμοὺς αὐτοῦ, καὶ κατήνεγκαν αὐτὸν εἰς Γάζαν, καὶ ἐπέδησαν αὐτὸν ἐν πέδαις χαλκείαις· καὶ ἦν ἀλήθων ἐν οἴκῳ τοῦ δεσμωτηρίου.
\vs{22}Καὶ ἤρξατο θρὶξ τῆς κεφαλῆς αὐτοῦ βλαστάνειν καθὼς ἐξυρήσατο.

\vs{23}Καὶ οἱ ἄρχοντες τῶν ἀλλοφύλων συνήχθησαν θυσαιάσαι θυσίασμα μέγα τῷ Δαγὼν θεῷ αὐτῶν, καὶ εὐφρανθῆναι, καὶ εἶπαν, ἔδωκεν ὁ θεὸς ἐν χειρὶ ἡμῶν τὸν Σαμψὼν τὸν ἐχθρὸν ἡμῶν.
\vs{24}Καὶ εἶδον αὐτὸν ὁ λαὸς, καὶ ὕμνησαν τὸν θεὸν αὐτῶν, ὅτι παρέδωκεν ὁ θεὸς ἡμῶν τὸν ἐχθρὸν ἡμῶν ἐν χειρὶ ἡμῶν, τὸν ἐρημοῦντα τὴν γῆν ἡμῶν, καὶ ὃς ἐπλήθυνε τοὺς τραυματίας ἡμῶν.
\vs{25}Καὶ ὅτε ἠγαθύνθη ἡ καρδία αὐτῶν, καὶ εἶπαν, καλέσατε τὸν Σαμψὼν ἐξ οἴκου φυλακῆς, καὶ παιξάτω ἐνώπιον ἡμῶν· καὶ ἐκάλεσαν τὸν Σαμψὼν ἐξ οἴκου δεσμωτηρίου, καὶ ἔπαιζεν ἐνώπιον αὐτῶν· καὶ ἐῤῥάπιζον αὐτὸν, καὶ ἔστησαν αὐτὸν ἀναμέσον τῶν κιόνων.
\vs{26}Καὶ εἶπε Σαμψὼν πρὸς τὸν νεανίαν τὸν κρατοῦντα τὴν χεῖρα αὐτοῦ, ἄφες με, καὶ ψηλαφήσω τοὺς κίονας ἐφʼ οἷς ὁ οἶκος ἐπʼ αὐτοὺς, καὶ ἐπιστηριχθήσομαι ἐπʼ αὐτούς.
\vs{27}Καὶ ὁ οἶκος πλήρης τῶν ἀνδρῶν καὶ τῶν γυναικῶν, καὶ ἐκεῖ πάντες οἱ ἄρχοντες τῶν ἀλλοφύλων, καὶ ἐπὶ τὸ δῶμα ὡσεὶ τρισχίλιοι ἄνδρες καὶ γυναῖκες οἱ θεωροῦντες ἐν παιγνίαις Σαμψών.

\vs{28}Καὶ ἔκλαυσε Σαμψὼν πρὸς Κύριον, καὶ εἶπεν, ἀδωναϊὲ Κύριε μνήσθητι δή μου, καὶ ἐνίσχυσόν με ἔτι τὸ ἅπαξ τοῦτο Θεὲ, καὶ ἀνταποδώσω ἀνταπόδοσιν μίαν περὶ τῶν δύο ὀφθαλμῶν μου τοῖς ἀλλοφύλοις.
\vs{29}Καὶ περιέλαβε Σαμψὼν τοὺς δύο κίονας τοῦ οἴκου ἐφʼ οὓς ὁ οἶκος εἱστήκει, καὶ ἐπεστηρίχθη ἐπʼ αὐτοὺς, καὶ ἐκράτησεν ἕνα τῇ δεξιᾷ αὐτοῦ, καὶ ἕνα τῇ ἀριστερᾷ αὐτοῦ.
\vs{30}Καὶ εἶπε Σαμψὼν, ἀποθανέτω ψυχή μου μετὰ τῶν ἀλλοφύλων· καὶ ἐβάσταξεν ἐν ἰσχύϊ· καὶ ἔπεσεν ὁ οἶκος ἐπὶ τοὺς ἄρχοντας, καὶ ἐπὶ πάντα τὸν λαὸν τὸν ἐν αὐτῷ· καὶ ἦσαν οἱ τεθνηκότες οὓς ἐθανάτωσε Σαμψὼν ἐν τῷ θανάτῳ αὐτοῦ, πλείους ἢ οὓς ἐθανάτωσεν ἐν τῇ ζωῇ αὐτοῦ.

\vs{31}Καὶ κατέβησαν οἱ ἀδελφοὶ αὐτοῦ, καὶ ὁ οἶκος τοῦ πατρὸς αὐτοῦ, καὶ ἔλαβον αὐτόν· καὶ ἀνέβησαν καὶ ἔθαψαν αὐτὸν ἀναμέσον Σαραὰ καὶ ἀναμέσον Ἐσθαὸλ ἐν τῷ τάφῳ Μανωὲ τοῦ πατρὸς αὐτοῦ· καὶ αὐτὸς ἔκρινε τὸν Ἰσραὴλ εἴκοσι ἔτη.

\ch{17}
Καὶ ἐγένετο ἀνὴρ ἀπὸ ὄρους Ἐφραὶμ, καὶ ὄνομα αὐτῷ Μιχαίας.
\vs{2}Καὶ εἶπε τῇ μητρὶ αὐτοῦ, οἱ χίλιοι καὶ ἑκατὸν οὗς ἔλαβες ἀργυρίου σεαυτῇ, καί με ἠράσω, καὶ προσεῖπας ἐν ὠσί μου, ἰδοὺ τὸ ἀργύριον παρʼ ἐμοὶ, ἐγὼ ἔλαβον αὐτό. καὶ εἶπεν ἡ μήτηρ αὐτοῦ, εὐλογητὸς ὁ υἱός μου τῷ Κυρίῳ.
\vs{3}Καὶ ἀπέδωκε τοὺς χιλίους καὶ ἑκατὸν τοῦ ἀργυρίου τῇ μητρὶ αὐτοῦ· καὶ εἶπεν ἡ μήτηρ αὐτοῦ, ἁγιάζουσα ἡγίασα τὸ ἀργύριον τῷ Κυρίῳ ἐκ τῆς χειρός μου τῷ υἱῷ μου τοῦ ποιῆσαι γλυπτὸν καὶ χωνευτὸν, καὶ νῦν ἀποδώσω αὐτό σοι.
\vs{4}Καὶ ἀπέδωκε τὸ ἀργύριον τῇ μητρὶ αὐτοῦ· καὶ ἔλαβεν ἡ μήτηρ αὐτοῦ διακοσίους ἀργυρίου, καὶ ἔδωκεν αὐτὸ ἀργυροκόπῳ, καὶ ἐποίησεν αὐτὸ γλυπτὸν καὶ χωνευτόν· καὶ ἐγενήθη ἐν οἴκῳ Μιχαία.
\vs{5}Καὶ ὁ οἶκος Μιχαία αὐτῷ οἶκος Θεοῦ· καὶ ἐποίησεν ἐφὼδ καὶ θαραφίν· καὶ ἐπλήρωσε τὴν χεῖρα ἀπὸ ἑνὸς υἱῶν αὐτοῦ, καὶ ἐγένετο αὐτῷ εἰς ἱερέα.

\vs{6}Ἐν δὲ ταῖς ἡμέραις ἐκείναις οὐκ ἦν βασιλεὺς ἐν Ἰσραήλ· ἀνὴρ τὸ εὐθὲς ἐν ὀφθαλμοῖς αὐτοῦ ἐποίει.

\vs{7}Καὶ ἐγενήθη νεανίας ἐκ Βηθλεέμ δήμου Ἰούδα, καὶ αὐτὸς Λευίτης, καὶ οὗτος παρῴκει ἐκεῖ.
\vs{8}Καὶ ἐπορεύθη ὁ ἀνὴρ ἀπὸ Βηθλεὲμ τῆς πόλεως Ἰούδα παροικῆσαι ἐν ᾧ ἐὰν εὕρῃ τόπῳ· καὶ ἦλθεν ἕως ὄρους Ἐφραὶμ, καὶ ἕως οἴκου Μιχαία τοῦ ποιῆσαι ὁδὸν αὐτοῦ.
\vs{9}Καὶ εἶπεν αὐτῷ Μιχαίας, πόθεν ἔρχῃ; καὶ εἶπε πρὸς αὐτὸν, Λευίτης εἰμὶ ἐκ Βηθλεὲμ Ἰούδα, καὶ ἐγὼ πορεύομαι παροικῆσαι ἐν ᾧ ἐὰν εὕρω τόπῳ.
\vs{10}Καὶ εἶπεν αὐτῷ Μιχαίας, κάθου μετʼ ἐμοῦ, καὶ γίνου μοι εἰς πατέρα καὶ εἰς ἱερέα, καὶ ἐγὼ δώσω σοι δέκα ἀργυρίου εἰς ἡμέραν, καὶ στολὴν ἱματίων, καὶ τὰ πρὸς ζωήν σου.
\vs{11}Καὶ ἐπορεύθη ὁ Λευίτης, καὶ ἤρξατο παροικεῖν παρὰ τῷ ἀνδρί· καὶ ἐγενήθη ὁ νεανίας αὐτῷ ὡς εἷς ἀπὸ υἱῶν αὐτοῦ.
\vs{12}Καὶ ἐπλήρωσε Μιχαίας τὴν χεῖρα τοῦ Λευίτου, καὶ ἐγένετο αὐτῷ εἰς ἱερέα, καὶ ἐγένετο ἐν τῷ οἴκῳ Μιχαία.
\vs{13}Καὶ εἶπε Μιχαίας, νῦν ἔγνων ὅτι ἀγαθυνεῖ μοι Κύριος, ὅτι ἐγένετό μοι ὁ Λευίτης εἰς ἱερέα.

\ch{18}
Ἐν ταῖς ἡμέραις ἐκείναις οὐκ ἦν βασιλεὺς ἐν Ἰσραήλ· καὶ ἐν ταῖς ἡμέραις ἐκείναις ἡ φυλὴ Δὰν ἐζήτει ἑαυτῇ κληρονομίαν κατοικῆσαι, ὅτι οὐκ ἐνέπεσεν αὐτῇ ἕως τῆς ἡμέρας ἐκείνης ἐν μέσῳ φυλῶν υἱῶν Ἰσραὴλ κληρονομία.
\vs{2}Καὶ ἀπέστειλαν οἱ υἱοὶ Δὰν ἀπὸ δήμων αὐτῶν πέντε ἄνδρας υἱοὺς δυνάμεως, ἀπὸ Σαραὰ καὶ ἀπὸ Ἐσθαὸλ τοῦ κατασκέψασθαι τὴν γῆν καὶ ἐξιχνιάσαι αὐτήν· καὶ εἶπαν πρὸς αὐτοὺς, πορεύεσθε καὶ ἐξιχνιάσατε τὴν γῆν· καὶ ἦλθον ἕως ὄρους Ἐφραὶμ ἕως οἴκου Μιχαία· καὶ ηὐλίσθησαν αὐτοὶ ἐκεῖ
\vs{3}ἐν οἴκῳ Μιχαία, καὶ αὐτοὶ ἐπέγνωσαν τὴν φωνὴν τοῦ νεανίσκου τοῦ Λευίτου, καὶ ἐξέκλιναν ἐκεῖ· καὶ εἶπαν αὐτῷ, τίς ἤνεγκέ σε ὧδε; καὶ σὺ τί ποιεῖς ἐν τῷ τόπῳ τούτῳ; καὶ τί σοι ὧδε;
\vs{4}Καὶ εἶπε πρὸς αὐτοὺς, οὕτω καὶ οὕτως ἐποίησέ μοι Μιχαίας, καὶ ἐμισθώσατό με, καὶ ἐγενόμην αὐτῷ εἰς ἱερέα.
\vs{5}Καὶ εἶπαν αὐτῷ, ἐπερώτησον δὴ ἐν τῷ Θεῷ, καὶ γνωσόμεθα εἰ εὐοδωθήσεται ἡ ὁδὸς ἡμῶν, ἐν ᾗ ἡμεῖς πορευόμεθα ἐν αὐτῇ.
\vs{6}Καὶ εἶπεν αὐτοῖς ὁ ἱερεὺς, πορεύεσθε ἐν εἰρήνῃ· ἐνώπιον Κυρίου ἡ ὁδὸς ὑμῶν, ἐν ᾗ πορεύεσθε ἐν αὐτῇ.

\vs{7}Καὶ ἐπορεύθησαν οἱ πέντε ἄνδρες, καὶ ἦλθον εἰς Λαισά· καὶ εἶδαν τὸν λαὸν τὸν ἐν μέσῳ αὐτῆς καθήμενον ἐπʼ ἐλπίδι, ὡς κρίσις Σιδωνίων ἡσυχάζουσα, καὶ οὐκ ἔστι διατρέπων ἢ καταισχύνων λόγον ἐν τῇ γῇ, κληρονόμος ἐκπιέζων θησαυροὺς, καὶ μακράν εἰσι Σιδωνίων, καὶ λόγον οὐκ ἔχουσι πρὸς ἄνθρωπον.
\vs{8}Καὶ ἦλθον οἱ πέντε ἄνδρες πρὸς τοὺς ἀδελφοὺς αὐτῶν εἰς Σαραὰ καὶ Ἐσθαὸλ, καὶ εἶπον τοῖς ἀδελφοῖς αὐτῶν, τί ὑμεῖς κάθησθε;
\vs{9}Καὶ εἶπαν, ἀνάστητε, καὶ ἀναβῶμεν ἐπʼ αὐτοὺς, ὅτι εἴδομεν τὴν γῆν, καὶ ἰδοὺ ἀγαθὴ σφόδρα, καὶ ὑμεῖς ἡσυχάζετε, μὴ ὀκνήσητε τοῦ πορευθῆναι, καὶ εἰσελθεῖν τοῦ κληρονομῆσαι τὴν γῆν.
\vs{10}Καὶ ἡνίκα ἐὰν ἔλθητε, εἰσελεύσεσθε πρὸς λαὸν ἐπʼ ἐλπίδι, καὶ ἡ γῆ πλατεῖα, ὅτι ἔδωκεν αὐτὴν ὁ Θεὸς ἐν χειρὶ ὑμῶν· τόπος ὅπου οὐκ ἔστιν ἐκεῖ ὑστέρημα παντὸς ῥήματος τῶν ἐν τῇ γῇ.

\vs{11}Καὶ ἀπῇραν ἐκεῖθεν ἀπὸ δήμων τοῦ Δὰν ἀπὸ Σαραὰ καὶ ἀπὸ Ἐσθαὸλ, ἑξακόσιοι ἄνδρες ἐζωσμένοι σκεύη παρατάξεως.
\vs{12}Καὶ ἀνέβησαν καὶ παρενέβαλον ἐν Καριαθιαρὶμ ἐν Ἰούδα. διὰ τοῦτο ἐκλήθη ἐν ἐκείνῳ τῷ τόπῳ, παρεμβολὴ Δὰν, ἕως τῆς ἡμέρας ταύτης· ἰδοὺ ὀπίσω Καριαθιαρίμ.

\vs{13}Καὶ παρῆλθον ἐκεῖθεν ὄρος Ἐφραὶμ, καὶ ἦλθον ἕως οἴκου Μιχαία.
\vs{14}Καὶ ἀπεκρίθησαν οἱ πέντε ἄνδρες οἱ πορευόμενοι κατασκέψασθαι τὴν γῆν Λαισὰ, καὶ εἶπαν πρὸς τοὺς ἀδελφοὺς, ἔγνωτε ὅτι ἐστὶν ἐν τῷ οἴκῳ τούτῳ ἐφὼδ καὶ θεραφὶν καὶ γλυπτὸν καὶ χωνευτόν· καὶ νῦν γνῶτε ὅ, τι ποιήσετε.
\vs{15}Καὶ ἐξέκλιναν ἐκεῖ, καὶ εἰσῆλθον εἰς τὸν οἶκον τοῦ νεανίσκου τοῦ Λευίτου, εἰς τὸν οἶκον Μιχαία, καὶ ἠρώτησαν αὐτὸν εἰς εἰρήνην.
\vs{16}Καὶ οἱ ἑξακόσιοι ἄνδρες οἱ ἀνεζωσμένοι τὰ σκεύη τῆς παρατάξεως αὐτῶν ἑστῶτες παρὰ θύρας τῆς πύλης, οἱ ἐκ τῶν υἱῶν Δάν.
\vs{17}Καὶ ἀνέβησαν οἱ πέντε ἄνδρες οἱ πορευθέντες κατασκέψασθαι τὴν γῆν,
\vs{18}καὶ εἰσῆλθον ἐκεῖ εἰς οἶκον Μιχαία, καὶ ὁ ἱερεὺς ἑστώς. Καὶ ἔλαβον τὸ γλυπτὸν καὶ τὸ ἐφὼδ καὶ τὸ θεραφὶν καὶ τὸ χωνευτόν· καὶ εἶπε πρὸς αὐτοὺς ὁ ἱερεὺς, τί ὑμεῖς ποιεῖτε;
\vs{19}Καὶ εἶπαν αὐτῷ, κώφευσον, ἐπίθες τὴν χεῖρά σου ἐπὶ τὸ στόμα σου, καὶ δεῦρο μεθʼ ἡμῶν, καὶ γένου ἡμῖν εἰς πατέρα καὶ εἰς ἱερέα· μὴ ἀγαθὸν εἶναί σε ἱερέα οἴκου ἀνδρὸς ἑνὸς, ἢ γενέσθαι σε ἱερέα φυλῆς καὶ οἴκου εἰς δῆμον Ἰσραήλ;
\vs{20}Καὶ ἠγαθύνθη ἡ καρδία τοῦ ἱερέως, καὶ ἔλαβε τὸ ἐφὼδ καὶ τὸ θεραφὶν καὶ τὸ γλυπτὸν καὶ τὸ χωνευτὸν, καὶ ἦλθεν ἐν μέσῳ τοῦ λαοῦ.

\vs{21}Καὶ ἐπέστρεψαν καὶ ἀπῆλθον, καὶ ἔθηκαν τὰ τέκνα καὶ τὴν κτῆσιν καὶ τὸ βάρος ἔμπροσθεν αὐτῶν.

\vs{22}Αὐτοὶ ἐμάκρυναν ἀπὸ οἴκου Μιχαία, καὶ ἰδοὺ Μιχαίας καὶ οἱ ἄνδρες οἱ ἐν ταῖς οἰκίαις ταῖς μετὰ οἴκου Μιχαία ἐβόησαν, καὶ κατελάβοντο τοὺς υἱοὺς Δάν.
\vs{23}Καὶ ἐπέστρεψαν οἱ υἱοὶ Δὰν τὸ πρόσωπον αὐτῶν, καὶ εἶπαν τῷ Μιχαίᾳ, τί ἐστί σοι, ὅτι ἐβόησας;
\vs{24}Καὶ εἶπε Μιχαίας, ὅτι τὸ γλυπτόν μου, ὃ ἐποίησα, ἐλάβετε, καὶ τὸν ἱερέα, καὶ ἐπορεύθητε· καὶ τί μοι ἔτι; καὶ τί τοῦτο λέγετε πρὸς μὲ, τί κράζεις;
\vs{25}Καὶ εἶπαν πρὸς αὐτὸν οἱ υἱοὶ Δὰν, μὴ ἀκουσθήτω δὴ φωνή σου μεθʼ ἡμῶν, μή ποτε συναντήσωσιν ὑμῖν ἄνδρες πικροὶ ψυχῇ, καὶ προσθήσουσι ψυχήν σου, καὶ τὴν ψυχὴν τοῦ οἴκου σου.
\vs{26}Καὶ ἐπορεύθησαν οἱ υἱοὶ Δὰν εἰς ὁδὸν αὐτῶν· καὶ εἶδε Μιχαίας, ὅτι δυνατώτεροί εἰσιν ὑπὲρ αὐτόν· καὶ ἐπέστρεψεν εἰς τὸν οἶκον αὐτοῦ.

\vs{27}Καὶ οἱ υἱοὶ Δὰν ἔλαβον ὃ ἐποιησε Μιχαίας, καὶ τὸν ἱερέα ὃς ἦν αὐτῷ, καὶ ἦλθον ἐπὶ Λαισὰ, ἐπὶ λαὸν ἡσυχάζοντα καὶ πεποιθότα ἐπʼ ἐλπίδι· καὶ ἐπάταξαν αὐτοὺς ἐν στόματι ῥομφαίας, καὶ τὴν πόλιν ἐνέπρησαν ἐν πυρί.
\vs{28}Καὶ οὐκ ἦν ὁ ῥυόμενος, ὅτι μακράν ἐστιν ἀπὸ Σιδωνίων, καὶ λόγος οὐκ ἔστιν αὐτοῖς μετὰ ἀνθρώπου· καὶ αὐτὴ ἐν τῇ κοιλάδι τοῦ οἴκου Ῥαάβ· καὶ ᾠκοδόμησαν τὴν πόλιν, καὶ κατεσκήνωσαν ἐν αὐτῇ,
\vs{29}καὶ ἐκάλεσαν τὸ ὄνομα τῆς πόλεως Δὰν, ἐν ὀνόματι Δὰν πατρὸς αὐτῶν, ὃς ἐτέχθη τῷ Ἰσραήλ· καὶ ἦν Οὐλαμαῒς ὄνομα τῆς πόλεως τοπρότερον.

\vs{30}Καὶ ἔστησαν ἑαυτοῖς οἱ υἱοὶ Δὰν τὸ γλυπτόν· καὶ Ἰωνάθαν υἱὸς Γηρσὼν υἱὸς Μανασσῆ αὐτὸς καὶ οἱ υἱοὶ αὐτοῦ ἦσαν ἱερεῖς τῇ φυλῇ Δὰν ἕως ἡμέρας τῆς ἀποικίας τῆς γῆς.
\vs{31}Καὶ ἔθηκαν ἑαυτοῖς τὸ γλυπτὸν ὁ ἐποίησε Μιχαίας, πάσας τὰς ἡμέρας ἃς ἦν ὁ οἶκος τοῦ Θεοῦ ἐν Σηλώμ· καὶ ἐγένετο ἐν ταῖς ἡμέραις ἐκείναις οὐκ ἦν βασιλεὺς ἐν Ἰσραήλ.

\ch{19}
Καὶ ἐγένετο ἀνὴρ Λευίτης παροικῶν ἐν μηροῖς ὄρους Ἐφραὶμ, καὶ ἔλαβεν αὐτῷ γυναῖκα παλλακὴν ἀπὸ Βηθλεὲμ Ἰούδα.
\vs{2}Καὶ ἐπορεύθη ἀπʼ αὐτοῦ ἡ παλλακὴ αὐτοῦ, καὶ ἀπῆλθε παρʼ αὐτοῦ εἰς οἶκον πατρὸς αὐτῆς εἰς Βηθλεὲμ Ἰούδα, καὶ ἦν ἐκεῖ ἡμέρας μηνῶν τεσσάρων.

\vs{3}Καὶ ἀνέστη ὁ ἀνὴρ αὐτῆς, καὶ ἐπορεύθη ὀπίσω αὐτῆς τοῦ λαλῆσαι ἐπὶ καρδίαν αὐτῆς, τοῦ ἐπιστρέψαι αὐτὴν αὐτῷ· καὶ νεανίας αὐτοῦ μετʼ αὐτοῦ, καὶ ζεῦγος ὄνων· ἡ δὲ εἰσήνεγκεν αὐτὸν εἰς οἶκον πατρὸς αὐτῆς· καὶ εἶδεν αὐτὸν ὁ πατὴρ τῆς νεάνιδος, καὶ ηὐφράνθη εἰς συνάντησιν αὐτοῦ.
\vs{4}Καὶ κατέσχεν αὐτὸν ὁ γαμβρὸς αὐτοῦ ὁ πατὴρ τῆς νεάνιδος, καὶ ἐκάθισε μετʼ αὐτοῦ ἐπὶ τρεῖς ἡμέρας, καὶ ἔφαγον καὶ ἔπιον, καὶ ηὐλίσθησαν ἐκεῖ.
\vs{5}Καὶ ἐγένετο τῇ ἡμέρᾳ τῇ τετάρτῃ, καὶ ὤρθρισαν τοπρωῒ καὶ ἀνέστη τοῦ πορευθῆναι, καὶ εἶπεν ὁ πατὴρ τῆς νεάνιδος πρὸς τὸν νυμφίον αὐτοῦ, στήρισον τὴν καρδιάν σου ψωμῷ ἄρτου, καὶ μετὰ τοῦτο πορεύσεσθε.
\vs{6}Καὶ ἐκάθισαν καὶ ἔφαγον οἱ δύο ἐπὶ τὸ αὐτὸ καὶ ἔπιον· καὶ εἶπεν ὁ πατὴρ τῆς νεάνιδος πρὸς τὸν ἄνδρα, ἄγε δὴ αὐλίσθητι, καὶ ἀγαθυνθήσεται ἡ καρδία σου.
\vs{7}Καὶ ἀνέστη ὁ ἀνὴρ τοῦ πορεύεσθαι αὐτός· καὶ ἐβιάσατο αὐτὸν ὁ γαμβρὸς αὐτοῦ, καὶ ἐκάθισε καὶ ηὐλίσθη ἐκεῖ.

\vs{8}Καὶ ὤρθρισε τοπρωῒ τῇ ἡμέρᾳ τῇ πέμπτῃ τοῦ πορευθῆναι· καὶ εἶπεν ὁ πατὴρ τῆς νεάνιδος, στήρισον δὴ τὴν καρδίαν σου, καὶ στράτευσον ἕως κλῖναι τὴν ἡμέραν· καὶ ἔφαγον οἱ δύο.
\vs{9}Καὶ ἀνέστη ὁ ἀνὴρ τοῦ πορευθῆναι αὐτὸς, καὶ ἡ παλλακὴ αὐτοῦ, καὶ ὁ νεανίας αὐτοῦ· καὶ εἶπεν αὐτῷ ὁ γάμβρὸς αὐτοῦ ὁ πατὴρ τῆς νεάνιδος, ἰδοὺ δὴ ἠσθένησεν ἡμέρα εἰς τὴν ἑσπέραν· αὐλίσθητι ὧδε, καὶ ἀγαθυνθήσεται ἡ καρδία σου, καὶ ὀρθριεῖτε αὔριον εἰς ὁδὸν ὑμῶν, καὶ πορεύσῃ εἰς τὸ σκήνωμά σου.
\vs{10}Καὶ οὐκ εὐδόκησεν ὁ ἀνὴρ αὐλισθῆναι, καὶ ἀνέστη καὶ ἀπῆλθε, καὶ ἦλθεν ἕως ἀπέναντι Ἰεβοὺς, αὕτη ἐστὶν Ἱερουσαλὴμ, καὶ μετʼ αὐτοῦ ζεῦγος ὄνων ἐπισεσαγμένων, καὶ ἡ παλλακὴ αὐτοῦ μετʼ αὐτοῦ.

\vs{11}Καὶ ἦλθοσαν ἕως Ἰεβούς· καὶ ἡ ἡμέρα προβεβήκει σφόδρα, καὶ εἴπεν ὁ νεανίας πρὸς τὸν κύριον αὐτοῦ, δεῦρο δὴ καὶ ἐκκλίνωμεν εἰς πόλιν τοῦ Ἰεβουσὶ ταύτην, καὶ αὐλισθῶμεν ἐν αὐτῇ.
\vs{12}Καὶ εἶπε πρὸς αὐτὸν ὁ κύριος αὐτοῦ, οὐκ ἐκκλινοῦμεν εἰς πόλιν ἀλλοτρίαν, ἐν ᾗ οὐκ ἔστιν ἀπὸ υἱῶν Ἰσραὴλ ὧδε, καὶ παρελευσόμεθα ἕως Γαβαά.
\vs{13}Καὶ εἶπε τῷ νεανίᾳ αὐτοῦ, δεῦρο καὶ ἐγγίσωμεν ἑνὶ τῶν τόπων, καὶ αὐλισθησόμεθα ἐν Γαβαᾷ ἢ ἐν Ῥαμᾷ.
\vs{14}Καὶ παρῆλθον καὶ ἐπορεύθησαν, καὶ ἔδυ αὐτοῖς ὁ ἥλιος ἐχόμενα τῆς Γαβαὰ, ἥ ἐστιν ἐν τῷ Βενιαμίν.
\vs{15}Καὶ ἐξέκλιναν ἐκεῖ τοῦ εἰσελθεὶν αὐλισθῆναι ἐν Γαβαᾷ· καὶ εἰσῆλθον, καὶ ἐκάθισαν ἐν τῇ πλατείᾳ τῆς πόλεως, καὶ οὐκ ἦν ἀνὴρ συνάγων αὐτοὺς εἰς οἰκίαν αὐλισθῆναι.

\vs{16}Καὶ ἰδοὺ ἀνὴρ πρεσβύτης ἤρχετο ἐξ ἔργων αὐτοῦ ἐξ ἀγροῦ ἐν ἑσπέρᾳ, καὶ ὁ ἀνὴρ ἦν ἐξ ὄρους Ἐφραὶμ, καὶ αὐτὸς παρῴκει ἐν Γαβαᾷ, καὶ οἱ ἄνδρες τοῦ τόπου υἱοὶ Βενιαμίν.
\vs{17}Καὶ ᾖρε τοὺς ὀφθαλμοὺς αὐτοῦ, καὶ εἶδε τὸν ὁδοιπόρον ἄνδρα ἐν τῇ πλατείᾳ τῆς πόλεως· καὶ εἶπεν ὁ ἀνὴρ ὁ πρεσβύτης, ποῦ πορεύῃ, καὶ πόθεν ἔρχῃ;
\vs{18}Καὶ εἶπε πρὸς αὐτὸν, παραπορευόμεθα ἡμεῖς ἀπὸ Βηθλεὲμ Ἰούδα ἕως μηρῶν ὄρους Ἐφραίμ· ἐκεῖθεν ἐγώ εἰμι, καὶ ἐπορεύθην ἕως Βηθλεὲμ Ἰούδα, καὶ εἰς τὸν οἶκόν μου ἐγὼ πορεύομαι, καὶ οὐκ ἔστιν ἀνὴρ συνάγων με εἰς τὴν οἰκίαν.
\vs{19}Καί γε ἄχυρα καὶ χορτάσματά ἐστι τοῖς ὄνοις ἡμῶν, καὶ ἄρτος καὶ οἶνός ἐστιν ἐμοὶ καὶ τῇ παιδίσκῃ καὶ τῷ νεανίσκῳ μετὰ τῶν παίδων σου· οὐκ ἔστιν ὑστέρημα παντὸς πράγματος.
\vs{20}Καὶ εἶπεν ὁ ἀνὴρ πρεσβύτης, εἰρήνη σοι· πλὴν πᾶν τὸ ὑστέρημά σου ἐπʼ ἐμὲ, πλὴν ἐν τῇ πλατείᾳ οὐ μὴ αὐλισθήσῃ.
\vs{21}Καὶ εἰσήνεγκεν αὐτὸν εἰς τὸν οἶκον αὐτοῦ, καὶ τόπον ἐποίησε τοῖς ὄνοις, καὶ αὐτοὶ ἐνίψαντο τοὺς πόδας αὐτῶν, καὶ ἔφαγον καὶ ἔπιον.

\vs{22}Αὐτοὶ δὲ ἀγαθύνοντες καρδίαν αὐτῶν· καὶ ἰδοὺ ἄνδρες τῆς πόλεως υἱοὶ παρανόμων ἐκύκλωσαν τὴν οἰκίαν, κρούοντες ἐπὶ τὴν θύραν· καὶ εἶπον πρὸς τὸν ἄνδρα τὸν κύριον τοῦ οἴκου τὸν πρεσβύτην, λέγοντες, ἐξένεγκε τὸν ἄνδρα ὅς εἰσῆλθεν εἰς τὴν οἰκίαν σου, ἵνα γνῶμεν αὐτόν.
\vs{23}Καὶ ἐξῆλθε πρὸς αὐτοὺς ὁ ἀνὴρ ὁ κύριος τοῦ οἴκου, καὶ εἶπε, μὴ ἀδελφοί, μὴ κακοποιήσητε δὴ μετὰ τὸ εἰσελθεῖν τὸν ἄνδρα τοῦτον εἰς τὴν οἰκίαν μου, μὴ ποιήσητε τὴν ἀφροσύνην ταύτην.
\vs{24}Ἴδε ἡ θυγάτηρ μου ἡ παρθένος, καὶ ἡ παλλακὴ αὐτοῦ· ἐξάξω αὐτὰς, καὶ ταπεινώσατε αὐτὰς, καὶ ποιήσατε αὐταῖς τὸ ἀγαθὸν ἐν ὀφθαλμοῖς ὑμῶν, καὶ τῷ ἀνδρὶ τούτῳ μή ποιήσητε τὸ ῥῆμα τῆς ἀφροσύνης ταύτης.
\vs{25}Καὶ οὐκ εὐδόκησαν οἱ ἄνδρες τοῦ εἰσακοῦσαι αὐτοῦ· καὶ ἐπελάβετο ὁ ἀνηρ τῆς παλλακῆς αὐτοῦ, καὶ ἐξήγαγεν αὐτὴν πρὸς αὐτοὺς ἔξω· καὶ ἔγνωσαν αὐτὴν, καὶ ἐνέπαιζον ἐν αὐτῇ ὅλην τὴν νύκτα ἕως τοπρωῒ, καὶ ἐξαπέστειλαν αὐτὴν ὡς ἀνέβη τοπρωΐ.

\vs{26}Καὶ ἦλθεν ἡ γυνὴ πρὸς τὸν ὄρθρον, καὶ ἔπεσε παρὰ τὴν θύραν τοῦ οἴκου οὗ ἦν αὐτῆς ἐκεῖ ὁ ἀνὴρ, ἕως οὗ διέφαυσε.
\vs{27}Καὶ ἀνέστη ὁ ἀνὴρ αὐτῆς τοπρωῒ, καὶ ἤνοιξε τὰς θύρας τοῦ οἴκου, καὶ ἐξῆλθε τοῦ πορευθῆναι τὴν ὁδὸν αὐτοῦ· καὶ ἰδοὺ ἡ γυνὴ αὐτοῦ ἡ παλλακὴ πεπτωκυῖα παρὰ τὰς θύρας τοῦ οἴκου, καὶ αἱ χεῖρες αὐτῆς ἐπὶ τὸ πρόθυρον.
\vs{28}Καὶ εἶπε πρὸς αὐτὴν, ἀνάστα καὶ ἀπέλθωμεν· καὶ οὐκ ἀπεκρίθη, ὅτι ἦν νεκρά· καὶ ἔλαβεν αὐτὴν ἐπὶ τὸν ὄνον, καὶ ἐπορεύθη εἰς τὸν τόπον αὐτοῦ.

\vs{29}Καὶ ἔλαβε τὴν ῥομφαίαν, καὶ ἐκράτησε τὴν παλλακὴν αὐτοῦ· καὶ ἐμέλισεν αὐτὴν εἰς δώδεκα μέλη, καὶ ἀπέστειλεν αὐτὰ ἐν παντὶ ὁρίῳ Ἰσραήλ.
\vs{30}Καὶ ἐγένετο πᾶς ὁ βλέπων ἔλεγεν, οὐκ ἐγένετο καὶ οὐχ ἑὼραται ἀπὸ ἡμέρας ἀναβάσεως υἱῶν Ἰσραὴλ ἐκ γῆς Αἰγύπτου ἕως τῆς ἡμέρας ταύτης ὠς αὐτή· θέσθε ὑμῖν αὐτοῖς βουλὴν ἐπʼ αὐτὴν, καὶ λαλήσατε.

\ch{20}
Καὶ ἐξῆλθον πάντες οἱ υἱοὶ Ἰσραὴλ, καὶ ἐξεκκλησιάσθη ἡ συναγωγὴ ὡς ἀνὴρ εἷς ἀπὸ Δὰν καὶ ἕως Βηρσαβεὲ, καὶ γῇ τοῦ Γαλαὰδ, πρὸς Κύριον εἰς Μασσηφά.
\vs{2}Καὶ ἐστάθησαν κατὰ πρόσωπον Κυρίου πᾶσαι αἱ φυλαὶ τοῦ Ἰσραὴλ ἐν ἐκκλησίᾳ τοῦ λαοῦ τοῦ Θεοῦ, τετρακόσιαι χιλιάδες ἀνδρῶν πεζῶν ἕλκοντες ῥομφαίαν.
\vs{3}Καὶ ἤκουσαν οἱ υἱοὶ Βενιαμεὶν, ὅτι ἀνέβησαν οἱ υἱοὶ Ἰσραὴλ εἰς Μασσηφά· Καὶ ἐλθόντες εἶπαν οἱ υἱοὶ Ἰσραὴλ, λαλήσατε, ποῦ ἐγένετο ἡ πονηρία αὕτη;
\vs{4}Καὶ ἀπεκριθη ὁ ἀνὴρ ὁ Λευίτης, ὁ ἀνὴρ τῆς γυναικὸς τῆς φονευθείσης, καὶ εἶπεν, εἰς Γαβαὰ τῆς Βενιαμὶν ἦλθον ἐγὼ καὶ ἡ παλλακή μου τοῦ αὐλισθῆναι,
\vs{5}καὶ ἀνέστησαν ἐπʼ ἐμὲ οἱ ἄνδρες τῆς Γαβαὰ, καὶ ἐκύκλωσαν ἐπʼ ἐμὲ ἐπὶ τὴν οἰκίαν νυκτός· ἐμὲ ἠθέλησαν φονεῦσαι, καὶ τὴν παλλακήν μου ἐταπείνωσαν, καὶ ἀπέθανε.
\vs{6}Καὶ ἐκράτησα τὴν παλλακήν μου, καὶ ἐμέλισα αὐτὴν, καὶ ἐξαπέστειλα ἐν παντὶ ὁρίῳ κληρονομίας υἱῶν Ἰσραήλ· ὅτι ἐποίησαν ζέμα καὶ ἀπόπτωμα ἐν Ἰσραήλ.
\vs{7}Ἰδοὺ πάντες ὑμεῖς υἱοὶ Ἰσραὴλ, δότε ἑαυτοῖς λόγον καὶ βουλὴν ἐκεῖ.

\vs{8}Καὶ ἀνέστη πᾶς ὁ λαὸς ὡς ἀνὴρ εἷς, λέγοντες, οὐκ ἀπελευσόμεθα ἀνὴρ εἰς σκήνωμα αὐτοῦ, καὶ οὐκ ἐπιστρέψομεν ἀνὴρ εἰς τὸν οἶκον αὐτοῦ.
\vs{9}Καὶ νῦν τοῦτο τὸ ῥῆμα, ὃ ποιηθήσεται τῇ Γαβαᾷ· ἀναβησόμεθα ἐπʼ αὐτὴν ἐν κλήρῳ.
\vs{10}Πλὴν ληψόμεθα δέκα ἄνδρας τοῖς ἑκατὸν εἰς πάσας φυλὰς Ἰσραὴλ, καὶ ἑκατὸν τοῖς χιλίοις, καὶ χιλίους τοῖς μυρίοις, λαβεῖν ἐπισιτισμὸν τοῦ ποιῆσαι ἐλθεῖν αὐτοὺς εἰς Γαβαὰ Βενιαμὶν, ποιῆσαι αὐτῇ κατὰ πᾶν τὸ ἀπόπτωμα, ὃ ἐποίησεν ἐν Ἰσραήλ.
\vs{11}Καὶ συνήχθη πᾶς ἀνὴρ Ἰσραὴλ εἰς τὴν πόλιν ὡς ἀνὴρ εἷς.

\vs{12}Καὶ ἀπέστειλαν αἱ φυλαὶ Ἰσραὴλ ἄνδρας ἐν πάσῃ φυλῇ Βενιαμεὶν, λέγοντες, τίς ἡ πονηρία αὕτη ἡ γενομένη ἐν ὑμῖν;
\vs{13}Καὶ νῦν δότε τοὺς ἄνδρας υἱοὺς παρανόμων τοὺς ἐν Γαβαὰ, καὶ θανατώσομεν αὐτοὺς, καὶ ἐκκαθαριοῦμεν πονηρίαν ἀπὸ Ἰσραήλ· καὶ οὐκ εὐδόκησαν οἱ υἱοὶ Βενιαμὶν ἀκοῦσαι τῆς φωνῆς τῶν ἀδελφῶν αὐτῶν υἱῶν Ἰσραήλ.
\vs{14}Καὶ συνήχθησαν οἱ υἱοὶ Βενιαμὶν ἀπὸ τῶν πόλεων αὐτῶν εἰς Γαβαὰ ἐξελθεῖν εἰς παράταξιν πρὸς υἱοὺς Ἰσραήλ.
\vs{15}Καὶ ἐπεσκέπησαν οἱ υἱοὶ Βενιαμὶν ἐν τῇ ἡμέρᾳ ἐκείνῃ ἀπὸ τῶν πόλεων εἰκοσιτρεῖς χιλιάδες ἀνὴρ ἕλκων ῥομφαίαν, ἐκτὸς τῶν οἰκούντων τὴν Γαβαά, οἱ ἐπεσκέπησαν ἑπτακόσιοι ἄνδρες
\vs{16}ἐκλεκτοὶ ἐκ παντὸς λαοῦ ἀμφοτεροδέξιοι· πάντες οὗτοι σφενδονῆται ἐν λίθοις πρὸς τρίχα, καὶ οὐκ ἐξαμαρτάνοντες.
\vs{17}Καὶ ἀνὴρ Ἰσραὴλ ἐπεσκέπησαν ἐκτὸς τοῦ Βενιαμὶν τετρακόσιαι χιλιάδες ἀνδρῶν ἑλκόντων ῥομφαίαν· πάντες οὗτοι ἄνδρες παρατάξεως.

\vs{18}Καὶ ἀνέστησαν καὶ ἀνέβησαν εἰς Βαιθὴλ, καὶ ἠρώτησαν ἐν τῷ Θεῷ· καὶ εἶπαν οἱ υἱοὶ Ἰσραὴλ, τίς ἀναβήσεται ἡμῖν ἐν ἀρχῇ εἰς παράταξιν πρὸς υἱοὺς Βενιαμίν; καὶ εἶπε Κύριος, Ἰούδας ἐν ἀρχῇ ἀναβήσεται ἀφηγούμενος.
\vs{19}Καὶ ἀνέστησαν οἱ υἱοὶ Ἰσραὴλ τοπρωῒ, καὶ παρενέβαλον ἐπὶ Γαβαά.

\vs{20}Καὶ ἐξῆλθον πᾶς ἀνὴρ Ἰσραὴλ εἰς παράταξιν πρὸς Βενιαμὶν, καὶ συνῆψαν αὐτοῖς ἐπὶ Γαβαά.
\vs{21}Καὶ ἐξῆλθον οἱ υἱοὶ Βενιαμὶν ἀπὸ τῆς Γαβαὰ, καὶ διέφθειραν ἐν Ἰσραὴλ ἐν τῇ ἡμέρᾳ ἐκείνῃ δύο καὶ εἴκοσι χιλιάδας ἀνδρῶν ἐπὶ τὴν γῆν.

\vs{22}Καὶ ἐνίσχυσαν ἀνῆρ Ἰσραὴλ, καὶ προσέθηκαν συνάψαι παράταξιν ἐν τῷ τόπῳ ὅπου συνῆψαν ἐν τῇ ἡμέρᾳ τῇ πρώτῃ.
\vs{23}Καὶ ἀνέβησαν οἱ υἱοὶ Ἰσραὴλ, καὶ ἔκλαυσαν ἐνώπιον Κυρίου ἕως ἑσπέρας, καὶ ἠρώτησαν ἐν Κυρίῳ, λέγοντες, εἰ προσθῶμεν ἐγγίσαι εἰς παράταξιν πρὸς υἱοὺς Βενιαμὶν ἀδελφοὺς ἡμῶν; καὶ εἶπε Κύριος, ἀνάβητε πρὸς αὐτούς.
\vs{24}Καὶ προσῆλθον οἱ υἱοὶ Ἰσραὴλ πρὸς υἱοὺς Βενιαμὶν ἐν τῇ ἡμέρᾳ τῇ δευτέρᾳ.
\vs{25}Καὶ ἐξῆλθον οἱ υἱοὶ Βενιαμὶν εἰς συνάντησιν αὐτοὶς ἀπὸ τῆς Γαβαὰ ἐν τῇ ἡμέρᾳ τῇ δευτέρᾳ, καὶ διέφθειραν ἀπὸ υἱῶν Ἰσραὴλ ἔτι ὀκτωκαίδεκα χιλιάδας ἀνδρῶν ἐπὶ τὴν γῆν· πάντες οὗτοι ἕλκοντες ῥομφαίαν.

\vs{26}Καὶ ἀνέβησαν πάντες οἱ υἱοὶ Ἰσραὴλ καὶ πᾶς ὁ λαὸς, καὶ ἦλθον εἰς Βαιθήλ· καὶ ἔκλαυσαν, καὶ ἐκάθισαν ἐκεῖ ἐνώπιον Κυρίου· καὶ ἐνήστευσαν ἐν τῇ ἡμέρᾳ ἐκείνῃ ἕως ἑσπέρας, καὶ ἀνήνεγκαν ὁλοκαυτώσεις καὶ τελείας ἐνώπιον Κυρίου,
\vs{27}ὅτι ἐκεῖ κιβωτὸς διαθήκης Κυρίου τοῦ Θεοῦ ἐν ταῖς ἡμέραις ἐκείναις,
\vs{28}καὶ Φινεὲς υἱὸς Ἐλεάζαρ υἱοῦ Ἀαρὼν παρεστηκὼς ἐνώπιον αὐτῆς ἐν ταῖς ἡμέραις ἐκείναις· καὶ ἐπηρώτησαν οἱ υἱοὶ Ἰσραὴλ ἐν Κυρίῳ, λέγοντες, εἰ προσθῶμεν ἔτι ἐξελθεῖν εἰς παράταξιν πρὸς υἱοὺς Βενιαμὶν ἀδελφοὺς ἡμῶν; καὶ εἶπε Κύριος, ἀνάβητε, αὔριον δώσω αὐτοὺς εἰς χεῖρας ὑμῶν.
\vs{29}Καὶ ἔθηκαν οἱ υἱοῖ Ἰσραὴλ ἔνεδρα τῇ Γαβαὰ κύκλῳ.

\vs{30}Καὶ ἀνέβησαν οἱ υἱοὶ Ἰσραὴλ πρὸς υἱοὺς Βενιαμὶν ἐν τῇ ἡμέρᾳ τῇ τρίτῃ, καὶ συνῆψαν πρὸς τὴν Γαβαὰ ὡς ἅπαξ καὶ ἅπαξ.
\vs{31}Καὶ ἐξῆλθον οἱ υἱοὶ Βενιαμὶν εἰς συνάντησιν τοῦ λαοῦ, καὶ ἐξεκενώθησαν ἐκ τῆς πόλεως, καὶ ἤρξαντο πατάσσειν ἀπὸ τοῦ λαοῦ τραυματίας ὡς ἅπαξ καὶ ἅπαξ ἐν ταῖς ὁδοῖς, ἥ ἐστι μία ἀναβαίνουσα εἰς Βαιθὴλ, καὶ μία εἰς Γαβαὰ ἐν ἀγρῷ, ὡς τριάκοντα ἄνδρας ἐν Ἰσραήλ.
\vs{32}Καὶ εἶπαν οἱ υἱοὶ Βενιαμὶν, πίπτουσιν ἐνώπιον ἡμῶν ὡς τὸ πρῶτον. καὶ οἱ υἱοὶ Ἰσραὴλ εἶπαν, Φύγωμεν, καὶ ἐκκενώσωμεν αὐτοὺς ἀπὸ τῆς πόλεως εἰς τὰς ὁδούς· καὶ ἐποίησαν οὕτω.

\vs{33}Καὶ πᾶς ἀνὴρ ἀνέστη ἐκ τοῦ τόπου αὐτῶν, καὶ συνῆψαν ἐν Βάαλ Θαμάρ· καὶ τὸ ἔνεδρον Ἰσραὴλ ἐπήρχετο ἐκ τοῦ τόπου αὐτοῦ ἀπὸ Μαρααγαβέ.
\vs{34}Καὶ ἦλθον ἐξεναντίας Γαβαὰ δέκα χιλιάδες ἀνδρῶν ἐκλεκτῶν ἐκ παντὸς Ἰσραήλ· καὶ παράταξις βαρεῖα· καὶ αὐτοὶ οὐκ ἔγνωσαν, ὅτι φθάνει ἀπʼ αὐτοὺς ἡ κακία.
\vs{35}Καὶ ἐπάταξε Κύριος τὸν Βενιαμὶν ἐνώπιον υἱῶν Ἰσραήλ· καὶ διέφθειραν οἱ υἱοὶ Ἰσραὴλ ἐκ τοῦ Βενιαμὶν ἐν τῇ ἡμέρᾳ ἐκείνῃ εἴκοσι καὶ πέντε χιλιάδας καὶ ἑκατὸν ἄνδρας· πάντες οὗτοι εἷλκον ῥομφαίαν.
\vs{36}Καὶ εἶδον οἱ υἱοὶ Βενιαμὶν ὅτι ἐπλήγησαν· καὶ ἔδωκεν ἀνὴρ Ἰσραὴλ τῷ Βενιαμὶν τόπον, ὅτι ἤλπισαν πρὸς τὸ ἔνεδρον ὃ ἔθηκαν ἐπὶ τῇ Γαβαᾷ.

\vs{37}Καὶ ἐν τῷ αὐτοὺς ὑποχωρῆσαι, καὶ τὸ ἔνεδρον ἐκινήθη· καὶ ἐξέτειναν ἐπὶ τὴν Γαβαὰ, καὶ ἐξεχύθη τὸ ἔνεδρον, καὶ ἐπάταξαν τὴν πόλιν ἐν στόματι ῥομφαίας.

\vs{38}Καὶ σημεῖον ἦν τοῖς υἱοῖς Ἰσραὴλ μετὰ τοῦ ἐνέδρου τῆς μάχης ἀνενέγκαι αὐτοὺς σύσσημον καπνοῦ ἀπὸ τῆς πόλεως.
\vs{39}Καὶ εἶδον οἱ υἱοὶ Ἰσραὴλ, ὅτι προκατελάβετο τὸ ἔνεδρον τὴν Γαβαὰ, καὶ ἔστησαν ἐν τῇ παρατάξει· καὶ Βενιαμὶν ἤρξατο πατάσσειν τραυματίας ἐν ἀνδράσιν Ἰσραὴλ ὡς τριάκοντα ἄνδρας· ὅτι εἶπαν, πάλιν πτώσει πίπτουσιν ἐνώπιον ἡμῶν ὡς ἡ παράταξις ἡ πρώτη.

\vs{40}Καὶ τὸ σύσσημον ἀνέβη ἐπὶ πλεῖον ἐπὶ τῆς πόλεως ὡς στύλος καπνοῦ· καὶ ἐπέβλεψε Βενιαμὶν ὀπίσω αὐτοῦ, καὶ ἰδοὺ ἀνέβη ἡ συντέλεια τῆς πόλεως ἕως οὐρανοῦ.

\vs{41}Καὶ ἀνὴρ Ἰσραὴλ ἐπέστρεψε· καὶ ἔσπευσαν ἄνδρες Βενιαμὶν, ὅτι εἶδον ὅτι συνήντησεν ἐπʼ αὐτοὺς ἡ πονηρία.
\vs{42}Καὶ ἐπέβλεψαν ἐνώπιον υἱῶν Ἰσραὴλ εἰς τὴς ὁδὸν τῆς ἐρήμου, καὶ ἔφυγον· καὶ ἡ παράταξις ἔφθασεν ἐπʼ αὐτοὺς, καὶ οἱ ἀπὸ τῶν πόλεων διέφθειρον αὐτοὺς ἐν μέσῳ αὐτῶν.

\vs{43}Καὶ κατέκοπτον τὸν Βενιαμὶν, καὶ ἐδίωξαν αὐτὸν ἀπὸ Νουὰ κατὰ πόδα αὐτοῦ ἕως ἀπέναντι Γαβαὰ πρὸς ἀνατολὰς ἡλίου.
\vs{44}Καὶ ἔπεσον ἀπὸ Βενιαμὶν ὀκτωκαίδεκα χιλιάδες ἀνδρῶν· οἱ πάντες οὗτοι ἄνδρες δυνάμεως.

\vs{45}Καὶ ἐπέβλεψαν οἱ λοιποὶ, καὶ ἔφευγον εἰς τὴν ἔρημον πρὸς τὴν πέτραν τοῦ Ῥέμμών· καὶ ἐκαλαμήσαντο ἐξ αὐτῶν οἱ υἱοὶ Ἰσραὴλ πεντακισχιλίους ἄνδρας· καὶ κατέβησαν ὀπίσω αὐτῶν οἱ υἱοὶ Ἰσραὴλ ἕως Γεδᾶν, καὶ ἐπάταξαν ἐξ αὐτῶν δισχιλίους ἄνδρας.
\vs{46}Καὶ ἐγένοντο πάντες οἱ πεπτωκότες ἀπὸ Βενιαμὶν, εἰκοσιπέντε χιλιάδες ἀνδρῶν ἑλκόντων ῥομφαίαν ἐν τῇ ἡμέρᾳ ἐκείνῃ· οἱ πάντες οὗτοι ἄνδρες δυνάμεως.
\vs{47}Καὶ ἐπέβλεψαν οἱ λοιποὶ, καὶ ἔφυγον εἰς τὴν ἔρημον πρὸς τὴν πέτραν τοῦ Ῥεμμὼν ἑξακόσιοι ἄνδρες, καὶ ἐκάθισαν ἐν πέτρᾳ Ῥεμμὼν τέσσαρας μῆνας.

\vs{48}Καὶ οἱ υἱοὶ Ἰσραὴλ ἐπέστρεψαν πρὸς υἱοὺς Βενιαμὶν, καὶ ἐπάταξαν αὐτοὺς ἐν στόματι ῥομφαίας ἀπὸ πόλεως Μεθλὰ καὶ ἕως κτήνους, καὶ ἕως παντὸς τοῦ εὑρισκομένου εἰς πάσας τὰς πόλεις· καὶ τὰς πόλεις τὰς εὑρεθείσας ἐνέπρησαν ἐν πυρί.

\ch{21}
Καὶ οἱ υἱοὶ Ἰσραὴλ ὤμοσαν ἐν Μασσηφὰθ, λέγοντες, ἀνὴρ ἐξ ἡμῶν οὐ δώσει θυγατέρα αὐτοῦ τῷ Βενιαμὶν εἰς γυναῖκα.
\vs{2}Καὶ ἦλθεν ὁ λαὸς εἰς Βαιθὴλ, καὶ ἐκάθισαν ἐκεῖ ἕως ἑσπέρας ἐνώπιον τοῦ Θεοῦ· καὶ ᾖραν φωνὴν αὐτῶν, καὶ ἔκλαυσαν κλαυθμὸν μέγαν,
\vs{3}καὶ εἶπαν, εἰς τί Κύριε Θεὲ Ισραὴλ ἐγενήθη αὕτη, τοῦ ἐπισκεπῆναι σήμερον ἀπὸ Ἰσραὴλ φυλὴν μίαν;
\vs{4}Καὶ ἐγένετο τῇ ἐπαύριον, καὶ ὤρθρισεν ὁ λαὸς, καὶ ᾠκοδόμησαν ἐκεῖ θυσιαστήριον, καὶ ἀνήνεγκαν ὁλοκαυτώσεις καὶ τελείας.

\vs{5}Καὶ εἶπαν οἱ υἱοὶ Ἰσραὴλ, τίς οὐκ ἀνέβη ἐν τῇ ἐκκλησίᾳ ἀπὸ πασῶν φυλῶν Ἰσραὴλ πρὸς Κύριον; ὅτι ὁ ὅρκος μέγας ἦν τοῖς οὐκ ἀναβεβηκόσι πρὸς Κύριον εἰς Μασσηφὰθ, λέγοντες, θανάτῳ θανατωθήσεται.

\vs{6}Καὶ παρεκλήθησαν οἱ υἱοὶ Ἰσραὴλ πρὸς Βενιαμὶν ἀδελφὸν αὐτῶν, καὶ εἶπαν, ἐξεκόπη σήμερον φυλὴ μία ἀπὸ Ἰσραήλ.
\vs{7}Τί ποιήσωμεν αὐτοῖς τοῖς περισσοῖς τοῖς ὑπολειφθεῖσιν εἰς γυναῖκας; καὶ ἡμεῖς ὠμόσαμεν ἐν Κυρίῳ τοῦ μὴ δοῦναι αὐτοῖς ἀπὸ τῶν θυγατέρων ἡμῶν εἰς γυναῖκας.
\vs{8}Καὶ εἶπαν, τίς εἷς ἀπὸ φυλῶν Ἰσραὴλ, ὃς οὐκ ἀνέβη πρὸς Κύριον εἰς Μασσηφάθ; Καὶ ἰδοὺ οὐκ ἦλθεν ἀνὴρ εἰς τὴν παρεμβολὴν ἀπὸ Ἰαβεῖς Γαλαὰδ εἰς τὴν ἐκκλησίαν.
\vs{9}Καὶ ἐπεσκέπη ὁ λαὸς, καὶ οὐκ ἦν ἐκεῖ ἀνὴρ ἀπὸ οἰκούντων Ἰαβὶς Γαλαάδ.

\vs{10}Καὶ ἀπέστειλεν ἐκεῖ ἡ συναγωγὴ δώδεκα χιλιάδας ἀνδρῶν ἀπὸ υἱῶν τῆς δυνάμεως, καὶ ἐνετείλαντο αὐτοῖς λέγοντες, πορεύεσθε καὶ πατάξατε τοὺς οἰκοῦντας Ἰαβεῖς Γαλαὰδ ἐν στόματι ῥομφαίας.
\vs{11}Καὶ τοῦτο ποιήσετε· πᾶν ἄρσεν καὶ πᾶσαν γυναῖκα εἰδυῖαν κοίτην ἄρσενος, ἀναθεματιεῖτε· τὰς δὲ παρθένους, περιποιήσεσθε· καὶ ἐποίησαν οὕτως.

\vs{12}Καὶ εὗρον ἀπὸ οἰκούντων Ἰαβεῖς Γαλαὰδ, τετρακοσίας νεάνιδας παρθένους, αἵτινες οὐκ ἔγνωσαν ἄνδρα εἰς κοίτην ἄρσενος, καὶ ἤνεγκαν αὐτὰς εἰς τὴν παρεμβολὴν εἰς Σηλὼμ τὴν ἐν γῇ Χαναάν.

\vs{13}Καὶ ἀπέστειλαν πᾶσα ἡ συναγωγὴ, καὶ ἐλάλησαν πρὸς τοὺς υἱοὺς Βενιαμεὶν ἐν τῇ πέτρᾳ Ῥεμμὼν, καὶ ἐκάλεσαν αὐτοὺς εἰς εἰρήνην.
\vs{14}Καὶ ἐπέστρεψε Βενιαμὶν πρὸς τοὺς υἱοὺς Ἰσραὴλ ἐν τῷ καιρῷ ἐκείνῳ, καὶ ἔδωκαν αὐτοῖς οἱ υἱοὶ Ἰσραὴλ τὰς γυναῖκας ἃς ἐζωοποίησαν ἀπὸ τῶν θυγατέρων Ἰαβὶς Γαλαάδ· καὶ ἤρεσεν αὐτοῖς οὕτω.

\vs{15}Καὶ ὁ λαὸς παρεκλήθη ἐπὶ τῷ Βενιαμὶν, ὅτι ἐποίησε Κύριος διακοπὴν ἐν ταῖς φυλαῖς Ἰσραήλ.

\vs{16}Καὶ εἶπον οἱ πρεσβύτεροι τῆς συναγωγῆς, τί ποιήσωμεν τοῖς περισσοῖς εἰς γυναῖκας; ὅτι ἠφανίσθη ἀπὸ Βενιαμὶν γυνή.
\vs{17}Καὶ εἶπαν, κληρονομία διασωζομένων τῶν Βενιαμίν· καὶ οὐκ ἐξαλειφθήσεται φυλὴ ἀπὸ Ἰσραὴλ,
\vs{18}ὅτι ἡμεῖς οὐ δυνησόμεθα δοῦναι αὐτοῖς γυναῖκας ἀπὸ τῶν θυγατέρων ἡμῶν, ὅτι ὠμόσαμεν ἐν υἱοῖς Ἰσραὴλ, λέγοντες, ἐπικατάρατος ὁ διδοὺς γυναῖκα τῷ Βενιαμίν.

\vs{19}Καὶ εἶπαν, ἰδοὺ δὴ ἑορτὴ Κυρίου ἐν Σηλὼμ ἀφʼ ἡμερῶν εἰς ἡμέρας, ἥ ἐστιν ἀπὸ Βοῤῥᾶ τῆς Βαιθὴλ, κατʼ ἀνατολὰς ἡλίου ἐπὶ τῆς ὁδοῦ τῆς ἀναβαινούσης ἀπὸ Βαιθὴλ εἰς Συχὲμ, καὶ ἀπὸ Νότου τῆς Λεβωνᾶ.
\vs{20}Καὶ ἐνετείλαντο τοῖς υἱοῖς Βενιαμεὶν, λέγοντες, πορεύεσθε καὶ ἐνεδρεύσατε ἐν τοῖς ἀμπελῶσι,
\vs{21}καὶ ὄψεσθε, καὶ ἰδοὺ, ἐὰν ἐξέλθωσιν αἱ θυγατέρες τῶν οἰκούντων Σηλὼ χορεύειν ἐν τοῖς χοροῖς, καὶ ἐξελεύσεσθε ἐκ τῶν ἀμπελώνων, καὶ ἁρπάσατε αὑτοῖς ἀνὴρ γυναῖκα ἀπὸ τῶν θυγατέρων Σηλὼμ, καὶ πορεύεσθε εἰς γῆν Βενιαμίν.
\vs{22}Καὶ ἔσται ὅταν ἔλθωσιν οἱ πατέρες αὐτῶν ἢ οἱ ἀδελφοὶ αὐτῶν κρίνεσθαι πρὸς ἡμᾶς, καὶ ἐροῦμεν αὐτοῖς, ἔλεος ποιήσατε ἡμῖν αὐτὰς, ὅτι οὐκ ἐλάβομεν ἀνὴρ γυναῖκα αὐτοῦ ἐν τῇ παρατάξει, ὅτι οὐχ ὑμεῖς ἐδώκατε αὐτοῖς, ὡς κλῆρος πλημμελήσατε.

\vs{23}Καὶ ἐποίησαν οὕτως οἱ υἱοὶ Βενιαμίν· καὶ ἔλαβον γυναῖκας εἰς ἀριθμὸν αὐτῶν ἀπὸ τῶν χορευουσῶν ὧν ἥρπασαν· καὶ ἐπορεύθησαν, καὶ ὑπέστρεψαν εἰς τὴν κληρονομίαν αὐτῶν· καὶ ᾠκοδόμησαν τὰς πόλεις, καὶ ἐκάθισαν ἐν αὐταῖς.
\vs{24}Καὶ περιεπάτησαν ἐκεῖθεν οἱ υἱοὶ Ἰσραὴλ ἐν τῷ καιρῷ ἐκείνῳ ἀνὴρ εἰς φυλὴν αὐτοῦ καὶ εἰς συγγένειαν αὐτοῦ· καὶ ἐξῆλθον ἐκεῖθεν ἀνὴρ εἰς τὴν κληρονομίαν αὐτοῦ.
\vs{25}Ἐν δὲ ταῖς ἡμέραις ἐκείναις οὐκ ἦν βασιλεὺς ἐν Ἰσραήλ· ἀνὴρ τὸ εὐθὲς ἐνώπιον αὐτοῦ ἐποίει.


\def\book{ΡΟΥΘ}
\biblebook{ΡΟΥΘ}


\lettrine[lines=2, loversize=0.2, nindent=0em, findent=.25em]{\textcolor{bookheadingcolor}{Κ}}{ΑΙ} ἐγένετο ἐν τῷ κρίνειν τοὺς κριτὰς, καὶ ἐγένετο λιμὸς ἐν τῇ γῇ· καὶ ἐπορεύθη ἀνὴρ ἀπὸ Βηθλεὲμ Ἰούδα τοῦ παροικῆσαι ἐν ἀγρῷ Μωὰβ, αὐτὸς καὶ ἡ γυνὴ αὐτοῦ, καὶ οἱ δύο υἱοὶ αὐτοῦ.
\vs{2}Καὶ ὄνομα τῷ ἀνδρὶ Ἐλιμέλεχ, καὶ ὄνομα τῇ γυναικὶ αὐτοῦ Νωεμὶν, καὶ ὄνομα τοῖς δυσὶν υἱοῖς αὐτοῦ Μααλὼν, καὶ Χελαιὼν, Ἐφραθαῖοι ἐκ Βηθλεὲμ τῆς Ἰούδα· καὶ ἤλθοσαν εἰς ἀγρὸν Μωὰβ, καὶ ἦσαν ἐκεῖ.

\vs{3}Καὶ ἀπέθανεν Ἐλιμέλεχ ὁ ἀνὴρ τῆς Νωεμὶν, καὶ κατελείφθη αὕτη καὶ οἱ δύο υἱοὶ αὐτῆς.
\vs{4}Καὶ ἐλάβοσαν ἑαυτοῖς γυναῖκας Μωαβίτιδας· ὄνομα τῇ μιᾷ, Ὀρφά· καὶ ὄνομα τῇ δευτέρᾳ, Ῥούθ· καὶ κατῴκησαν ἐκεῖ ὡς δέκα ἔτη.
\vs{5}Καὶ ἀπέθανον καί γε ἀμφότεροι Μααλὼν καὶ Χελαιών· καὶ κατελείφθη ἡ γυνὴ ἀπὸ τοῦ ἀνδρὸς αὐτῆς, καὶ ἀπὸ τῶν δύο υἱῶν αὐτῆς.

\vs{6}Καὶ ἀνεστη αὕτη καὶ αἱ δύο νύμφαι αὐτῆς, καὶ ἀπέστρεψαν ἐξ ἀγροῦ Μωὰβ, ὅτι ἤκουσεν ἐν ἀγρῷ Μωὰβ ὅτι ἐπέσκεπται Κύριος τὸν λαὸν αὐτοῦ, δοῦναι αὐτοῖς ἄρτους.
\vs{7}Καὶ ἐξῆλθεν ἐκ τοῦ τόπου οὗ ἦν ἐκεῖ, καὶ αἱ δύο νύμφαι αὐτῆς μετʼ αὐτῆς· καὶ ἐπορεύοντο ἐν τῇ ὁδῷ τοῦ ἐπιστρέψαι εἰς τὴν γῆν Ἰούδα.

\vs{8}Καὶ εἶπε Νωεμὶν, ταῖς δυσὶ νύμφαις αὐτῆς, πορεύεσθε δὴ, ἀποστράφητε ἑκάστη εἰς οἶκον μητρὸς αὐτῆς· ποιήσαι Κύριος μεθʼ ὑμῶν ἔλεος, καθὼς ἐποιήσατε μετὰ τῶν τεθνηκότων καὶ μετʼ ἐμοῦ·
\vs{9}Δῴη Κύριος ὑμῖν καὶ εὕρητε ἀνάπαυσιν ἑκάστη ἐν οἴκῳ ἀνδρὸς αὐτῆς· καὶ κατεφίλησεν αὐτάς· καὶ ἐπῇραν τὴν φωνὴν αὐτῶν, καὶ ἔλαυσαν.
\vs{10}Καὶ εἶπαν αὐτῇ, μετὰ σου ἐπιστρέφομεν εἰς τὸν λαόν σου.

\vs{11}Καὶ εἶπε Νωεμὶν, ἐπιστράφητε δὴ θυγατέρες μου· καὶ ἱνατί πορεύεσθε μετʼ ἐμοῦ; μὴ ἔτι μοι υἱοὶ ἐν τῇ κοιλίᾳ μου, καὶ ἔσονται ὑμῖν εἰς ἄνδρας;
\vs{12}Ἐπιστράφητε δὴ θυγατέρες μου, διότι γεγήρακα τοῦ μὴ εἶναι ἀνδρί· ὅτι εἶπα, ὅτι ἐστί μοι ὑπόστασις τοῦ γενηθῆναί με ἀνδρὶ, καὶ τέξομαι υἱούς·
\vs{13}Μὴ αὐτοὺς προσδέξεσθε ἕως οὗ ἁδρυνθώσιν; ἢ αὐτοῖς κατασχεθήσεσθε τοῦ μὴ γενέσθαι ἀνδρί; μὴ δὴ θυγατέρες μου, ὅτι ἐπικράνθη μοι ὑπὲρ ὑμᾶς, ὅτι ἐξῆλθεν ἐν ἐμοὶ χεὶρ Κυρίου.

\vs{14}Καὶ ἐπῇραν τὴν φωνὴν αὐτῶν, καὶ ἔκλαυσαν ἔτι· καὶ κατεφίλησεν Ὀρφὰ τὴν πενθερὰν αὐτῆς, καὶ ἐπέστρεψεν εἰς τὸν λαὸν αὐτῆς· Ῥοὺθ δὲ ἠκολούθησεν αὐτῇ.

\vs{15}Καὶ εἶπε Νωεμῖν πρὸς Ῥοὺθ, ἰδοὺ ἀνέστρεψε σύννυμφός σου πρὸς λαὸν αὐτῆς καὶ πρὸς τοὺς θεοὺς αὐτῆς· ἐπιστράφηθι δὴ καὶ σὺ ὀπίσω τῆς συννύμφου σου.
\vs{16}Εἶπε δὴ Ῥοῦθ, μὴ ἀπάντησαί μοι τοῦ καταλιπεῖν σε, ἢ ἀποστρέψαι ὄπισθέν σου, ὅτι σὺ ὅπου ἐὰν πορευθῇς, πορεύσομαι, καὶ οὗ ἐὰν αὐλισθῇς, αὐλισθήσομαι· ὁ λαός σου, λαός μου, καὶ ὁ Θεός σου, Θεός μου·
\vs{17}Καὶ οὗ ἐὰν ἀποθάνῃς, ἀποθανοῦμαι, κᾀκεῖ ταφήσομαι· τάδε ποιήσαι μοι Κύριος, καὶ τάδε προσθείη, ὅτι θάνατος διαστελεῖ ἀναμέσον ἐμοῦ καὶ σοῦ.
\vs{18}Ἰδοῦσα δὲ Νωεμὶν ὅτι κραταιοῦται αὐτὴ τοῦ πορεύεσθαι μετʼ αὐτῆς, ἐκόπασε τοῦ λαλῆσαι πρὸς αὐτὴν ἔτι.

\vs{19}Ἐπορεύθησαν δὲ ἀμφότεραι, ἕως τοῦ παραγενέσθαι αὐτὰς εἰς Βηθλεέμ· καὶ ἐγένετο ἐν τῷ ἐλθεῖν αὐτὰς εἰς Βηθλεὲμ, καὶ ἤχησε πᾶσα ἡ πόλις ἐπʼ αὐταῖς, καὶ εἶπον, εἰ αὕτη ἐστὶ Νωεμίν;
\vs{20}Καὶ εἶπε πρὸς αὐτὰς, μὴ δὴ καλεῖτέ με Νωεμίν· καλέσατέ με πικρὰν, ὅτι ἐπικράνθη ἐν ἐμοὶ ὁ ἱκανὸς σφόδρα.
\vs{21}Ἐγὼ πλήρης ἐπορεύθην, καὶ κενὴν ἀπέστρεψέ με ὁ Κύριος· καὶ ἱνατί καλεῖτέ με Νωεμὶν, καὶ Κύριος ἐταπείνωσέ με, καὶ ὁ ἱκανὸς ἐκάκωσέ με;

\vs{22}Καὶ ἐπέστρεψε Νωεμὶν καὶ Ῥοὺθ ἡ Μωαβίτις ἡ νύμφη αὐτῆς ἐπιστρέφουσαι ἐξ ἀγροῦ Μωάβ· αὗται δὲ παρεγενήθησαν εἰς Βηθλεὲμ ἐν ἀρχῇ θερισμοῦ κριθῶν.

\ch{2}
Καὶ τῇ Νωεμὶν ἀνὴρ γνώριμος τῷ ἀνδρὶ αὐτῆς, ὁ δὲ ἀνὴρ δυνατὸς ἰσχύϊ ἐκ τῆς συγγενείας Ἐλιμέλεχ, καὶ ὄνομα αὐτῷ Βοόζ.
\vs{2}Καὶ εἶπε Ῥοὺθ ἡ Μωαβίτις πρὸς Νωεμὶν, πορεῦθω δὴ εἰς ἀγρὸν, καὶ συνάξω ἐν τοῖς στάχυσι κατόπισθεν οὗ ἐὰν εὕρω χάριν ἐν ὀφθαλμοῖς αὐτοῦ. εἶπε δὲ αὐτῇ, πορεύου, θύγατερ.
\vs{3}Καὶ ἐπορεύθη· καὶ ἐλθοῦσα συνέλεξεν ἐν τῷ ἀγρῷ κατόπισθε τῶν θεριζόντων· καὶ περιέπεσε περιπτώματι τῇ μερίδι τοῦ ἀγροῦ Βοὸζ, τοῦ ἐκ τῆς συγγενείας Ἐλιμέλεχ.

\vs{4}Καὶ ἰδοὺ Βοὸζ ἦλθεν ἐκ Βηθλεὲμ, καὶ εἶπε τοῖς θερίζουσι, Κύριος μεθʼ ὑμῶν· καὶ εἶπον αὐτῷ, εὐλογήσαι σε Κύριος.
\vs{5}Καὶ εἶπε Βοὸζ τῷ παιδαρίῳ αὐτοῦ τῷ ἐφεστῶτι ἐπὶ τοὺς θερίζοντας, τίνος ἡ νεᾶνις αὕτη;
\vs{6}Καὶ ἀπεκρίθη τὸ παιδάριον τὸ ἐφεστὸς ἐπὶ τοὺς θερίζοντας, καὶ εἶπεν, ἡ παῖς ἡ Μωαβίτις ἐστὶν ἡ ἀποστραφεῖσα μετὰ Νωεμὶν ἐξ ἀγροῦ Μωάβ·
\vs{7}Καὶ εἶπε, συλλέξω δὴ καὶ συνάξω ἐν τοῖς δράγμασιν ὄπισθεν τῶν θεριζόντων· καὶ ἦλθε καὶ ἔστη ἀπὸ πρωΐθεν καὶ ἕως ἑσπέρας, οὐ κατέπαυσεν ἐν τῷ ἀγρῷ μικρόν.

\vs{8}Καὶ εἶπε Βοὸζ πρὸς Ῥοὺθ, οὐκ ἤκουσας θύγατερ; μὴ πορευθῇς ἐν ἀγρῷ συλλέξαι ἑτέρῳ· καὶ σὺ οὐ πορεύσῃ ἐντεῦθεν, ὧδε κολλήθητι μετὰ τῶν κορασίων μου.
\vs{9}Οἱ ὀφθαλμοί σου εἰς τὸν ἀγρὸν οὗ ἐὰν θερίζωσι, καὶ πορεύσῃ κατόπισθεν αὐτῶν· ἰδοὺ ἐνετειλάμην τοῖς παιδαρίοις τοῦ μὴ ἅψασθαί σου· καὶ ὅτε διψήσεις καὶ πορευθήσῃ εἰς τὰ σκεύη, καὶ πίεσαι ὅθεν ἐὰν ὑδρεύωνται τὰ παιδάρια.
\vs{10}Καὶ ἔπεσεν ἐπὶ πρόσωπον αὐτῆς, καὶ προσεκύνησεν ἐπὶ τὴν γῆν, καὶ εἶπε πρὸς αὐτὸν, τί ὅτι εὗρον χάριν ἐν ὀφθαλμοῖς σου τοῦ ἐπιγνῶναί με, καὶ ἐγώ εἰμι ξένη;

\vs{11}Καὶ ἀπεκρίθη Βοὸζ, καὶ εἶπεν αὐτῇ, ἀπαγγελίᾳ ἀπηγγέλη μοι ὅσα πεποίηκας μετὰ τῆς πενθερᾶς σου μετὰ τὸ ἀποθανεῖν τὸν ἄνδρα σου· καὶ πῶς κατέλιπες τὸν πατέρα σου καὶ τὴν μητέρα σου, καὶ τὴν γῆν γενέσεώς σου, καὶ ἐπορεύθης πρὸς λαὸν ὃν οὐκ ᾔδεις ἐχθὲς καὶ τρίτης.
\vs{12}Ἀποτίσαι Κύριος τὴν ἐργασίαν σου· γένοιτο ὁ μισθός σου πλήρης παρὰ Κυρίου Θεοῦ Ἰσραὴλ, πρὸς ὃν ἦλθες πεποιθέναι ὑπὸ τὰς πτέρυγας αὐτοῦ.
\vs{13}Ἡ δὲ εἶπεν, εὕροιμι χάριν ἐν ὀφθαλμοῖς σου κύριε, ὅτι παρεκάλεσάς με, καὶ ὅτι ἐλάλησας ἐπὶ καρδίαν τῆς δούλης σου, καὶ ἰδοὺ ἐγὼ ἔσομαι ὡς μία τῶν παιδισκῶν σου.

\vs{14}Καὶ εἶπεν αὐτῇ Βοὸζ, ἤδη ὥρα τοῦ φαγεῖν, πρόσελθε ὧδε καὶ φάγεσαι τῶν ἄρτων, καὶ βάψεις τὸν ψωμόν σου ἐν τῷ ὄξει· καὶ ἐκάθισε Ῥοὺθ ἐκ πλαγίων τῶν θεριζόντων· καὶ ἐβούνισεν αὐτῇ Βοὸζ ἄλφιτον, καὶ ἔφαγε καὶ ἐνεπλήσθη καὶ κατέλιπε,

\vs{15}Καὶ ἀνέστη τοῦ συλλέγειν· καὶ ἐνετείλατο Βοὸζ τοῖς παιδαρίοις αὐτοῦ, λέγων, καί γε ἀναμέσον τῶν δραγμάτων συλλεγέτω, καὶ μὴ καταισχύνητε αὐτήν·
\vs{16}Καὶ βαστάζοντες βαστάσατε αὐτῇ, καί γε παραβάλλοντες παραβαλεῖτε αὐτῇ ἐκ τῶν βεβουνισμένων, καὶ φάγεται, καὶ συλλέξει, καὶ οὐκ ἐπιτιμήσετε αὐτῇ.
\vs{17}Καὶ συνέλεξεν ἐν τῷ ἀγρῷ ἕως ἑσπέρας, καὶ ἐῤῥάβδισεν ἃ συνέλεξε, καὶ ἐγενήθη ὡς οἰφὶ κριθῶν.

\vs{18}Καὶ ᾖρε καὶ εἰσῆλθεν εἰς τὴν πόλιν· καὶ εἶδεν ἡ πενθερὰ αὐτῆς ἃ συνέλεξε· καὶ ἐξενέγκασα Ῥοὺθ ἔδωκεν αὐτῇ ἃ κατέλιπεν ἐξ ὧν ἐνεπλήσθη.
\vs{19}Καὶ εἶπεν αὐτῇ ἡ πενθερὰ αὐτῆς, ποῦ συνέλεξας σήμερον καὶ ποῦ ἐποίησας; εἴη ὁ ἐπιγνούς σε εὐλογημένος· καὶ ἀνήγγειλε Ῥοὺθ τῇ πενθερᾷ αὐτῆς ποῦ ἐποίησε, καὶ εἶπε, τὸ ὄνομα τοῦ ἀνδρὸς μεθʼ οὗ ἐποίησα σήμερον Βοόζ.
\vs{20}Εἶπε δὲ Νωεμὶν τῇ νύμφῃ αὐτῆς, εὐλογητός ἐστι τῷ Κυρίῳ, ὅτι οὐκ ἐγκατέλιπε τὸ ἔλεος αὐτοῦ μετὰ τῶν ζώντων καὶ μετὰ τῶν τεθνηκότων· καὶ εἶπεν αὐτῇ Νωεμὶν, ἐγγίζει ἡμῖν ὁ ἀνὴρ, ἐκ τῶν ἀγχιστευόντων ἡμῖν ἐστι.
\vs{21}Καὶ εἶπε Ῥοὺθ πρὸς τὴν πενθερὰν αὐτῆς, καί γε ὅτι εἶπε πρὸς μέ, μετὰ τῶν κορασίων τῶν ἐμῶν προσκολλήθητι, ἕως ἂν τελέσωσιν ὅλον τὸν ἀμητὸν ὃς ὑπάρχει μοι.

\vs{22}Καὶ εἶπε Νωεμὶν πρὸς Ῥοὺθ τὴν νύμφην αὐτῆς, ἀγαθὸν θύγατερ, ὅτι ἐξῆλθες μετὰ τῶν κορασίων αὐτου, καὶ οὐκ ἀπαντήσονταί σοι ἐν ἀγρῷ ἑτέρῳ.
\vs{23}Καὶ προσεκολλήθη Ῥοὺθ τοῖς κορασίοις τοῦ Βοὸζ τοῦ συλλέγειν, ἕως τοῦ συντελέσαι τὸν θερισμὸν τῶν κριθῶν καὶ τῶν πυρῶν.

\ch{3}
Καὶ ἐκάθισε μετὰ τῆς πενθερᾶς αὐτῆς· εἶπε δὲ αὐτῇ Νωεμὶν ἡ πενθερὰ αὐτῆς, θύγατερ, οὐ μὴ ζητήσω σοι ἀνάπαυσιν, ἵνα εὖ γένηταί σοι;
\vs{2}Καὶ νῦν οὐχὶ Βοὸζ γνώριμος ἡμῶν, οὗ ἦς μετὰ τῶν κορασίων αὐτοῦ; ἰδοὺ αὐτὸς λικμᾷ τὸν ἅλωνα τῶν κριθῶν ταύτῃ τῇ νυκτί.

\vs{3}Σὺ δὲ λούσῃ, καὶ ἀλείψῃ, καὶ περιθήσεις τὸν ἱματισμόν σου ἐπὶ σὲ, καὶ ἀναβήσῃ ἐπὶ τὸν ἅλω· μὴ γνωρισθῇς τῷ ἀνδρὶ ἕως τοῦ συντελέσαι αὐτὸν τοῦ φαγεῖν καὶ πιεῖν.
\vs{4}Καὶ ἔσται ἐν τῷ κοιμηθῆναι αὐτὸν, καὶ γνώσῃ τὸν τόπον ὅπου κοιμᾶται ἐκεῖ, καὶ ἐλεύσῃ καὶ ἀποκαλύψεις τὰ πρὸς ποδῶν αὐτοῦ, καὶ κοιμηθήσῃ, καὶ αὐτὸς ἀπαγγελεῖ σοι ἃ ποιήσεις.
\vs{5}Εἶπε δὲ Ῥοὺθ πρὸς αὐτὴν, πάντα ὅσα ἂν εἴπῃς, ποιήσω.

\vs{6}Καὶ κατέβη εἰς τὸν ἅλω, καὶ ἐποίησε κατὰ πάντα, ὅσα ἐνετείλατο αὐτῇ ἡ πενθερὰ αὐτῆς.
\vs{7}Καὶ ἔφαγε Βοὸζ καὶ ἔπιε, καὶ ἠγαθύνθη ἡ καρδία αὐτοῦ, καὶ ἦλθε κοιμηθῆναι ἐν μερίδι τῆς στοιβῆς· ἡ δὲ ἦλθεν ἐν κρυφῇ, καὶ ἀπεκάλυψε τὰ πρὸς ποδῶν αὐτοῦ.
\vs{8}Ἐγένετο δὲ ἐν τῷ μεσονυκτίῳ, καὶ ἐξέστη ὁ ἀνὴρ, καὶ ἐταράχθη, καὶ ἰδοὺ γυνὴ κοιμᾶται πρὸς ποδῶν αὐτοῦ.
\vs{9}Εἶπε δὲ, τίς εἶ σύ; ἡ δὲ εἶπεν, ἐγὼ εἰμι Ῥοὺθ ἡ δούλη σου, καὶ περιβαλεῖς τὸ πτερύγιόν σου ἐπὶ τὴν δούλην σου, ὅτι ἀγχιστεὺς εἶ σύ.
\vs{10}Καὶ εἶπε Βοὸζ, εὐλογημένη σὺ τῷ Κυρίῳ Θεῷ, θύγατερ, ὅτι ἠγάθυνας τὸ ἔλεός σου τὸ ἔσχατον ὑπὲρ τὸ πρῶτον, μὴ πορευθῆναί σε ὀπίσω νεανιῶν, εἴτοι πτωχὸς εἴτοι πλούσιος.
\vs{11}Καὶ νῦν θύγατερ μὴ φοβοῦ, πάντα ὅσα ἐὰν εἴπῃς ποιήσω σοι· οἶδε γὰρ πᾶσα φυλὴ λαοῦ μου ὅτι γυνὴ δυνάμεως εἶ σύ.
\vs{12}Καὶ νῦν ὁ ἀληθῶς ἀγχιστεὺς ἐγώ εἰμι· καί γε ἐστὶν ἀγχιστεὺς ἐγγίων ὑπὲρ ἐμέ.
\vs{13}Αὐλίσθητι τὴν νύκτα, καὶ ἔσται τοπρωῒ ἐὰν ἀγχιστεύσῃ σε, ἀγαθόν· ἀγχιστευέτω· ἐὰν δὲ μὴ βούληται ἀγχιστεῦσαί σε, ἀγχιστεύσω σε ἐγώ· ζῇ Κύριος· κοιμήθητι ἕως τοπρωΐ.

\vs{14}Καὶ ἐκοιμήθη πρὸς ποδῶν αὐτοῦ ἕως πρωΐ· ἡ δὲ ἀνέστη πρὸ τοῦ ἐπιγνῶναι ἄνδρα τὸν πλησίον αὐτοῦ· καὶ εἶπε Βοὸζ, μὴ γνωσθήτω, ὅτι ἦλθε γυνὴ εἰς τὸν ἅλω.

\vs{15}Καὶ εἶπεν αὐτῇ, Φέρε τὸ περίζωμα τὸ ἐπάνω σου· καὶ ἐκράτησεν αὐτὸ, καὶ ἐμέτρησεν ἓξ κριθῶν, καὶ ἐπέθηκεν ἐπʼ αὐτὴν, καὶ εἰσῆλθεν εἰς τὴν πόλιν.

\vs{16}Καὶ Ῥοὺθ εἰσῆλθε πρὸς τὴν πενθερὰ αὐτῆς· ἡ δὲ εἶπεν αὐτῇ, θύγατερ· καὶ εἶπεν αὐτῇ πάντα ὅσα ἐποίησεν αὐτῇ ὁ ἀνήρ.
\vs{17}Καὶ εἶπεν αὐτῇ, τὰ ἕξ τῶν κριθῶν ταῦτα ἔδωκέ μοι, ὅτι εἶπε πρὸς μὲ, μὴ εἰσέλθῃς κενὴ πρὸς τὴν πενθεράν σου.
\vs{18}Ἡ δὲ εἶπε, κάθου θύγατερ, ἕως τοῦ ἐπιγνῶναί σε πῶς οὐ πεσεῖται ῥῆμα· οὐ γὰρ μὴ ἡσυχάσῃ ὁ ἀνὴρ ἕως ἂν τελεσθῇ τὸ ῥῆμα σήμερον.

\ch{4}
Καὶ Βοὸζ ἀνέβη ἐπὶ τὴν πύλην, και ἐκαθισεν ἐκεῖ, καὶ ἰδοὺ ὁ ἀγχιστεὺς παρεπορεύετο, ὃν ἐλάλησε Βοόζ· καὶ εἶπε πρὸς αὐτὸν Βοὸζ, ἐκκλίνας κάθισον ὧδε κρύφιε· καὶ ἐξέκλινε καὶ ἐκάθισε.
\vs{2}Καὶ ἔλαβε Βοὸζ δέκα ἄνδρας ἀπὸ τῶν πρεσβυτέρων τῆς πόλεως, καὶ εἶπε, καθίσατε ὧδε· καὶ ἐκάθισαν.

\vs{3}Καὶ εἶπε Βοὸζ τῷ ἀγχιστεῖ, τὴν μερίδα τοῦ ἀγροῦ ἥ ἐστι τοῦ ἀδελφοῦ ἡμῶν τοῦ Ἐλιμέλεχ, ἣ δέδοται Νωεμὶν τῇ ἐπιστρεφούσῃ ἐξ ἀγροῦ Μωὰβ,
\vs{4}κᾀγὼ εἶπα, ἀποκαλύψω τὸ οὖς σου λέγων, κτῆσαι ἐναντίον τῶν καθημένων, καὶ ἐναντίον τῶν πρεσβυτέρων τοῦ λαοῦ μου. εἰ ἀγχιστεύεις, ἀγχίστευε· εἰ δὲ μὴ ἀγχιστεύεις, ἀνάγγειλόν μοι, καὶ γνώσομαι, ὅτι οὐκ ἔστι πάρεξ σοῦ τοῦ ἀγχιστεῦσαι, κᾀγώ εἰμι μετὰ σέ· ὁ δὲ εἶπεν, ἐγώ εἰμι, ἀγχιστεῦσαι, κἀγώ εἰμι, ἀγχιστεύσω.
\vs{5}Καὶ εἶπε Βοὸζ, ἐν ἡμέρᾳ τοῦ κτήσασθαί σε τὸν ἀγρὸν ἐκ χειρὸς Νωεμὶν καὶ παρὰ Ῥοὺθ τῆς Μωαβίτιδος γυναικὸς τοῦ τεθνηκότος, καὶ αὐτὴν κτήσασθαί σε δεῖ, ὥστε ἀναστῆσαι τὸ ὄνομα τοῦ τεθνηκότος ἐπὶ τῆς κληρονομίας αὐτοῦ.
\vs{6}Καὶ εἶπεν ὁ ἀγχιστεὺς, οὐ δυνήσομαι ἀγχιστεῦσαι ἐμαυτῷ, μή ποτε διαφθείρω τὴν κληρονομίαν μου· ἀγχίστευσον σεαυτῷ τὴν ἀγχιστείαν μου, ὅτι οὐ δυνήσομαι ἀγχιστεῦσαι.

\vs{7}Καὶ τοῦτο τὸ δικαίωμα ἔμπροσθεν ἐν τῷ Ἰσραὴλ ἐπὶ τὴν ἀγχιστείαν, καὶ ἐπὶ τὸ ἀντάλλαγμα τοῦ στῆσαι πάντα λόγον· καὶ ὑπελύετο ἀνὴρ τὸ ὑπόδημα αὐτοῦ, καὶ ἐδίδου τῷ πλησίον αὐτοῦ τῷ ἀγχιστεύοντι τὴν ἀγχιστείαν αὐτοῦ· καὶ τοῦτο ἦν μαρτύριον ἐν Ἰσραήλ.
\vs{8}Καὶ εἶπεν ὁ ἀγχιστεὺς τῷ Βοὸζ, κτῆσαι σεαυτῷ τὴν ἀγχιστείαν μου· καὶ ὑπελύσατο τὸ ὑπόδημα αὐτοῦ, καὶ ἔδωκεν αὐτῷ.

\vs{9}Καὶ εἶπε Βοὸζ τοῖς πρεσβυτέροις καὶ παντὶ τῷ λαῷ, μάρτυρες ὑμεῖς σήμερον, ὅτι κέκτημαι πάντα τὰ τοῦ Ἐλιμέλεχ, καὶ πάντα ὅσα ὑπάρχει τῷ Χελαιὼν καὶ τῷ Μααλὼν ἐκ χειρὸς Νωεμίν.
\vs{10}Καί γε Ῥοὺθ τὴν Μωαβίτιν τὴν γυναῖκα Μααλὼν κέκτημαι ἐμαυτῷ εἰς γυναῖκα, τοῦ ἀναστῆσαι τὸ ὄνομα τοῦ τεθνηκότος ἐπὶ τῆς κληρονομίας αὐτοῦ, καὶ οὐκ ἐξολοθρευθήσεται τὸ ὄνομα τοῦ τεθνηκότος ἐκ τῶν ἀδελφῶν αὐτοῦ, καὶ ἐκ τῆς φυλῆς λαοῦ αὐτοῦ· μάρτυρες ὑμεῖς σήμερον.

\vs{11}Καὶ εἴποσαν πᾶς ὁ λαὸς οἱ ἐν τῇ πύλῃ, μάρτυρες· καὶ οἱ πρεσβύτεροι εἴποσαν, δῴη Κύριος τὴν γυναῖκά σου, τὴν εἰσπορευομένην εἰς τὸν οἶκόν σου, ὡς Ῥαχὴλ καὶ ὡς Λίαν, αἳ ᾠκοδόμησαν ἀμφότεραι τὸν οἶκον τοῦ Ἰσραὴλ, καὶ ἐποίησαν δύναμιν ἐν Ἐφραθᾷ, καὶ ἔσται ὄνομα ἐν Βηθλεέμ.
\vs{12}Καὶ γένοιτο οἶκός σου, ὡς οἶκος Φαρὲς, ὃν ἔτεκε Θάμαρ τῷ Ἰούδᾳ, ἐκ τοῦ σπέρματος οὗ δώσει Κύριός σοι ἐκ τῆς παιδίσκης ταύτης.

\vs{13}Καὶ ἔλαβε Βοὸζ τὴν Ῥοὺθ, καὶ ἐγενήθη αὐτῷ εἰς γυναῖκα, καὶ εἰσῆλθε πρὸς αὐτήν. καὶ ἔδωκεν αὐτῇ Κύριος κύησιν, καὶ ἔτεκεν υἱόν.
\vs{14}Καὶ εἶπαν αἱ γυναῖκες πρὸς Νωεμὶν, εὐλογητὸς Κύριος, ὃς οὐ κατέλυσέ σοι σήμερον τὸν ἀγχιστέα, καὶ καλέσαι τὸ ὄνομά σου ἑν Ἰσραήλ.
\vs{15}Καὶ ἔσται σοι εἰς ἐπιστρέφοντα ψυχὴν, καὶ τοῦ διαθρέψαι τὴν πολιάν σου, ὅτι ἡ νύμφη ἡ ἀγαπήσασά σε, ἔτεκεν αὐτὸν, ἥ ἐστιν ἀγαθή σοι ὑπὲρ ἑπτὰ υἱούς.
\vs{16}Καὶ ἔλαβε Νωεμὶν τὸ παιδίον, καὶ ἔθηκεν εἰς τὸν κόλπον αὐτῆς, καὶ ἐγενήθη αὐτῷ εἰς τιθηνόν.

\vs{17}Καὶ ἐκάλεσαν αὐτοῦ αἱ γείτονες ὄνομα, λέγουσαι, ἐτέχθη υἱὸς τῇ Νωεμίν. καὶ ἐκάλεσαν τὸ ὄνομα αὐτοῦ, Ὠβήδ· οὗτος πατὴρ Ἰεσσαὶ πατρὸς Δαυίδ.
\vs{18}Καὶ αὗται αἱ γενέσεις Φαρές. Φαρές ἐγέννησε τὸν Ἐσρώμ·
\vs{19}Ἐσρὼμ ἐγέννησε τὸν Ἀράμ· καὶ Ἀρὰμ ἐγέννησε τὸν Ἀμιναδάβ·
\vs{20}Καὶ Ἀμιναδὰβ ἐγέννησε τὸν Ναασσών· καὶ Ναασσὼν ἐγέννησε τὸν Σαλμών·
\vs{21}Καὶ Σαλμὼν ἐγέννησε τὸν Βοόζ· καὶ Βοὸζ ἐγέννησε τὸν Ὠβήδ·
\vs{22}Καὶ Ὡβὴδ ἐγέννησε τὸν Ἰεσσαί· καὶ Ἰεσσαὶ ἐγέννησε τὸν Δαυίδ.


\def\book{ΒΑΣΙΛΕΙΩΝ Α}
\biblebook{ΒΑΣΙΛΕΙΩΝ Α}


\lettrine[lines=2, loversize=0.2, nindent=0em, findent=.25em]{\textcolor{bookheadingcolor}{Ἄ}}{ΝΘΡΩΠΟΣ} ἦν ἐξ Ἀρμαθαὶμ Σιφὰ, ἐξ ὄρους Ἐφραὶμ, καὶ ὄνομα αὐτῷ Ἑλκανὰ υἱὸς Ἰερεμεὴλ υἱοῦ Ἡλιοὺ υἱοῦ Θοκὲ ἐν
\vs{2}Νασὶβ Ἐφραίμ. Καὶ τοῦτῳ δύο γυναῖκες· ὄνομα τῇ μιᾷ, Ἄννα· καὶ ὄνομα τῇ δευτέρᾳ, Φεννάνα. Καὶ ἦν τῇ Φεννάνᾳ παιδία· καὶ τῇ Ἄννᾳ οὐκ ἦν παιδίον.

\vs{3}Καὶ ἀνέβαινεν ὁ ἄνθρωπος ἐξ ἡμερῶν εἰς ἡμέρας ἐκ πόλεως αὐτοῦ ἐξ Ἀρμαθαὶμ προσκυνεῖν καὶ θύειν Κυρίῳ τῷ Θεῷ σαβαὼθ εἰς Σηλώμ· καὶ ἐκεῖ Ἡλὶ καὶ οἱ δύο υἱοὶ αὐτοῦ Ὀφνὶ καὶ Φινεὲς ἱερεῖς τοῦ Κυρίου.

\vs{4}Καὶ ἐγενήθη ἡμέρα, καὶ ἔθυσεν Ἑλκανὰ, καὶ ἔδωκε τῇ Φεννάνᾳ γυναικὶ αὐτοῦ, καὶ τοῖς υἱοῖς αὐτῆς μερίδας.
\vs{5}Καὶ τῇ Ἄννᾳ ἔδωκε μερίδα μίαν, ὅτι οὐκ ἦν αὐτῇ παιδίον, πλὴν ὅτι τὴν Ἄνναν ἠγάπα Ἑλκανὰ ὑπὲρ ταύτην· καὶ Κύριος ἀπέκλεισε τὰ περὶ τὴν μήτραν αὐτῆς,
\vs{6}ὅτι οὐκ ἔδωκεν αὐτῇ Κύριος παιδίον κατὰ τὴν θλίψιν αὐτῆς, καὶ κατὰ τὴν ἀθυμίαν τῆς θλίψεως αὐτῆς· καὶ ἠθύμει διὰ τοῦτο, ὅτι συνέκλεισε Κύριος τὰ περὶ τὴν μήτραν αὐτῆς τοῦ μὴ δοῦναι αὐτῇ παιδίον.
\vs{7}Οὕτως ἐποίει ἐνιαυτὸν κατʼ ἐνιαυτὸν, ἐν τῷ ἀναβαίνειν αὐτὴν εἰς οἶκον Κυρίου· καὶ ἠθύμει, καὶ ἔκλαιε, καὶ οὐκ ἤσθιε.

\vs{8}Καὶ εἶπεν αὐτῇ Ἑλκανὰ ὁ ἀνὴρ αὐτῆς, Ἄννα· καὶ εἶπεν αὐτῷ, ἰδοὺ ἐγώ, κύριε· καὶ εἶπεν αὐτῇ, τί ἐστί σοι ὅτι κλαίεις; καὶ ἱνατί οὐκ ἐσθίεις; καὶ ἱνατί τύπτει σε ἡ καρδία σου; οὐκ ἀγαθὸς ἐγώ σοι ὑπὲρ δέκα τέκνα;

\vs{9}Καὶ ἀνέστη Ἄννα μετὰ τὸ φαγεῖν αὐτοὺς ἐν Σηλὼμ, καὶ κατέστη ἐνώπιον Κυρίου· καὶ Ἡλὶ ὁ ἱερεὺς, ἐπὶ τοῦ δίφρου ἐπὶ τῶν φλιῶν ναοῦ Κυρίου.

\vs{10}Καὶ αὐτὴ κατώδυνος ψυχῇ, καὶ προσηύξατο πρὸς Κύριον, καί κλαίουσα ἔκλαυσε.
\vs{11}Καὶ ηὔξατο εὐχὴν Κυρίῳ, λέγουσα, Ἀδωναῒ Κύριε ἐλωὲ σαβαὼθ, ἐὰν ἐπιβλέπων ἐπιβλέψῃς ἐπὶ τὴν ταπείνωσιν τῆς δούλης σου, καὶ μνησθῇς μου, καὶ δῷς τῇ δούλῃ σου σπέρμα ἀνδρῶν, καὶ δώσω αὐτὸν ἐνώπιόν σου δοτὸν ἕως ἡμέρας θανάτου αὐτοῦ, καὶ οἶνον καὶ μέθυσμα οὐ πίεται, καὶ σίδηρος οὐκ ἀναβήσεται ἐπὶ τὴν κεφαλὴν αὐτοῦ.

\vs{12}Καὶ ἐγενήθη ὅτε ἐπλήθυνε προσευχομένη ἐνώπιον Κύριου, καὶ Ἡλὶ ὁ ἱερεὺς ἐφύλαξε τὸ στόμα αὐτῆς.
\vs{13}Καὶ αὕτη ἐλάλει ἐν τῇ καρδίᾳ αὐτῆς, καὶ τὰ χείλη αὐτῆς ἐκινεῖτο, καὶ φωνὴ αὐτῆς οὐκ ἠκούετο· καὶ ἐλογίσατο αὐτὴν Ἡλὶ εἰς μεθύουσαν.
\vs{14}Καὶ εἶπεν αὐτῇ τὸ παιδάριον Ἡλὶ, ἕως πότε μεθυσθήσῃ; περιελοῦ τὸν οἶνόν σου, καὶ πορεύου ἐκ προσώπου Κυρίου.
\vs{15}Καὶ ἀπεκρίθη Ἄννα, καὶ εἶπεν, οὐχὶ κύριε· γυνὴ ἡ σκληρὰ ἡμέρα ἐγώ εἰμι, καὶ οἶνον καὶ μέθυσμα οὐ πέπωκα, καὶ ἐκχέω τὴν ψυχήν μου ἐνώπιον Κυρίου.
\vs{16}Μὴ δῷς τὴν δούλην σου εἰς θυγατέρα λοιμὴν, ὅτι ἐκ πλήθους ἀδολεσχίας μου ἐκτέτακα ἕως νῦν.
\vs{17}Καὶ ἀπεκρίθη Ἡλὶ, καὶ εἶπεν αὐτῇ, πορεύου εἰς εἰρήνην· ὁ Θεὸς Ἰσραὴλ δῴη σοι πᾶν αἴτημά σου, ὃ ᾐτήσω παρʼ αὐτοῦ.
\vs{18}Καὶ εἶπεν, εὗρεν ἡ δούλη σου χάριν ἐν ὀφθαλμοῖς σου· καὶ ἐπορεύθη ἡ γυνὴ εἰς τὴν ὁδὸν αὐτῆς· καὶ εἰσῆλθεν εἰς τὸ κατάλυμα αὐτῆς, καὶ ἔφαγε μετὰ τοῦ ἀνδρὸς αὐτῆς καὶ ἔπιε, καὶ τὸ πρόσωπον αὐτῆς οὐ συνέπεσεν ἔτι.

\vs{19}Καὶ ὀρθρίζουσι τοπρωῒ καὶ προσκυνοῦσι τῷ Κυρίῳ, καὶ πορεύονται τὴν ὁδὸν αὐτῶν· καὶ εἰσῆλθεν Ἑλκανὰ εἰς τὸν οἴκον αὐτοῦ Ἀρμαθαὶμ, καὶ ἔγνω τὴν Ἄνναν γυναῖκα αὐτοῦ· καὶ ἐμνήσθη αὐτῆς Κύριος, καὶ συνέλαβε.
\vs{20}Καὶ ἐγενήθη τῷ καιρῷ τῶν ἡμερῶν, καὶ ἔτεκεν υἱὸν, καὶ ἐκάλεσε τὸ ὄνομα αὐτοῦ Σαμουὴλ, καὶ εἶπεν, ὅτι παρὰ Κυρίου Θεοῦ σαβαὼθ ᾐτησάμην αὐτόν.

\vs{21}Καὶ ἀνέβη ὁ ἄνθρωπος Ἑλκανὰ καὶ πᾶς ὁ οἶκος αὐτοῦ θῦσαι ἐν Σηλὼμ τὴν θυσίαν τῶν ἡμερῶν, καὶ τὰς εὐχὰς αὐτοῦ, καὶ πάσας τὰς δεκάτας τῆς γῆς αὐτοῦ.
\vs{22}Καὶ Ἄννα οὐκ ἀνέβη μετʼ αὐτοῦ, ὅτι εἶπε τῷ ἀνδρὶ αὐτῆς, ἕως τοῦ ἀναβῆναι τὸ παιδάριον, ἐὰν ἀπογαλακτίσω αὐτὸ, καὶ ὀφθήσεται τῷ προσώπῳ Κυρίου, καὶ καθήσεται ἕως αἰῶνος ἐκεῖ.
\vs{23}Καὶ εἶπεν αὐτῇ Ἑλκανὰ ὁ ἀνὴρ αὐτῆς, ποίει τὸ ἀγαθὸν ἐν ὀφθαλμοῖς σου, κάθου ἕως ἂν ἀπογαλακτίσῃς αὐτό· ἀλλὰ στήσαι Κύριος τὸ ἐξελθὸν ἐκ τοῦ στόματός σου· καὶ ἐκάθισεν ἡ γυνὴ καὶ ἐθήλασε τὸν υἱὸν αὐτῆς, ἕως ἂν ἀπογαλακτίσῃ αὐτόν.

\vs{24}Καὶ ἀνέβη μετʼ αὐτοῦ εἰς Σηλὼμ ἐν μόσχῳ τριετίζοντι, καὶ ἄρτοις, καὶ οἰφὶ σεμιδάλεως, καὶ νέβελ οἴνου· καὶ εἰσῆλθεν εἰς οἶκον Κυρίου ἐν Σηλὼμ, καὶ τὸ παιδάριον μετʼ αὐτῶν.
\vs{25}Καὶ προσήγαγον ἐνώπιον Κυρίου· καὶ ἔσφαξεν ὁ πατὴρ αὐτοῦ τὴν θυσίαν, ἣν ἐποίει ἐξ ἡμερῶν εἰς ἡμέρας τῷ Κυρίῳ· καὶ προσήγαγε τὸ παιδάριον, καὶ ἔσφαξε τὸν μόσχον· καὶ προσήγαγεν Ἄννα ἡ μήτηρ τοῦ παιδαρίου πρὸς Ἡλὶ,
\vs{26}καὶ εἶπεν, ἐν ἐμοί κύριε ζῇ ἡ ψυχή σου, ἐγὼ ἡ γυνὴ ἡ καταστᾶσα ἐνώπιόν σου μετὰ σοῦ ἐν τῷ προσεύξασθαι πρὸς Κύριον.
\vs{27}Ὑπὲρ τοῦ παιδαρίου τούτου προσηυξάμην· καὶ ἐδωκέ μοι Κύριος τὸ αἴτημά μου ὃ ᾐτησάμην παρʼ αὐτοῦ.
\vs{28}Κᾀγὼ κιχρῶ αὐτὸν τῷ Κυρίῳ πάσας τὰς ἡμέρας ἃς ζῇ αὐτὸς, χρῆσιν τῷ Κυρίῳ, καὶ εἶπεν,

\ch{2}
Ἐστερέωθη ἡ καρδία μου ἐν Κυρίῳ, ὑψώθη κέρας μου ἐν Θεῷ μου, ἐπλατύνθη ἐπʼ ἐχθρούς μου τὸ στόμα μου, εὐφράνθην ἐν σωτηρίᾳ σου.
\vs{2}Ὅτι οὐκ ἔστιν ἅγιος ὡς Κύριος, καὶ οὐκ ἔστι δίκαιος ὡς ὁ Θεὸς ἡμῶν, οὐκ ἔστιν ἅγιος πλήν σου.
\vs{3}Μὴ καυχᾶσθε, καὶ μὴ λαλεῖτε ὑψηλά· μὴ ἐξελθέτω μεγαλοῤῥημοσύνη ἐκ τοῦ στόματος ὑμῶν, ὅτι Θεὸς γνώσεων Κύριος, καὶ Θεὸς ἑτοιμάζων ἐπιτηδεύματα αὐτοῦ.
\vs{4}Τόξον δυνατῶν ἠσθένησε, καὶ ἀσθενοῦντες περιεζώσαντο δυνάμιν.
\vs{5}Πλήρεις ἄρτων ἠλαττώθησαν, καὶ οἱ πεινῶντες παρῆκαν γῆν· ὅτι στεῖρα ἔτεκεν ἑπτὰ, καὶ ἡ πολλὴ ἐν τέκνοις ἠσθένησε.
\vs{6}Κύριος θανατοῖ καὶ ζωογονεῖ, κατάγει εἰς ᾅδου καὶ ἀνάγει.
\vs{7}Κύριος πτωχίζει καὶ πλουτίζει, ταπεινοῖ καὶ ἀνυψοῖ.
\vs{8}Ἀνιστᾷ ἀπὸ γῆς πένητα, καὶ ἀπὸ κοπρίας ἐγείρει πτωχὸν, καθίσαι μετὰ δυναστῶν λαοῦ, καὶ θρόνον δόξης κατακληρονομῶν αὐτοῖς,
\vs{9}διδοὺς εὐχὴν τῷ εὐχομένῳ· καὶ εὐλόγησεν ἔτη δικαίου, ὅτι οὐκ ἐν ἰσχύϊ δυνατὸς ἀνήρ.
\vs{10}Κύριος ἀσθενῆ ποιήσει ἀντίδικον αὐτοῦ, Κύριος ἅγιος· μὴ καυχάσθω ὁ φρόνιμος ἐν τῇ φρονήσει αὐτοῦ, καὶ μὴ καυχάσθω ὁ δυνατὸς ἐν τῇ δυνάμει αὐτοῦ, καὶ μὴ καυχάσθω ὁ πλούσιος ἐν τῷ πλούτῳ αὐτοῦ· ἀλλʼ ἐν τούτῳ καυχάσθω ὁ καυχώμενος, συνιεῖν καὶ γινώσκειν τὸν Κύριον, καὶ ποιεῖν κρίμα καὶ δικαιοσύνην ἐν μέσῳ τῆς γῆς. Κύριος ἀνέβη εἰς οὐρανοὺς, καὶ ἐβρόντησεν· αὐτὸς κρινεῖ ἄκρα γῆς, καὶ δίδωσιν ἰσχὺν τοῖς βασιλεῦσιν ἡμῶν, καὶ ὑψώσει κέρας χριστοῦ αὐτοῦ.

\vs{11}Καὶ κατέλιπεν αὐτὸν ἐκεῖ ἐνώπιον Κυρίου, καὶ ἀπῆλθεν εἰς Ἀρμαθαίμ· καὶ τὸ παιδάριον ἦν λειτουργῶν τῷ προσώπῳ Κυρίου ἐνώπιον Ἡλὶ τοῦ ἱερέως.
\vs{12}Καὶ οἱ υἱοὶ Ἡλὶ τοῦ ἱερέως υἱοὶ λοιμοί, οὐκ εἰδότες τὸν Κύριον.
\vs{13}Καὶ τὸ δικαίωμα τοῦ ἱερέως παρὰ τοῦ λαοῦ παντὸς τοῦ θύοντος· καὶ ἤρχετο τὸ παιδάριον τοῦ ἱερέως ὡς ἂν ἡψήθη τὸ κρέας, καὶ κρεάγρα τριόδους ἐν τῇ χειρὶ αὐτοῦ,
\vs{14}καὶ ἐπάταξεν αὐτὴν εἰς τὸν λέβητα τὸν μέγαν ἢ εἰς τὸ χαλκεῖον ἢ εἰς τὴν χύτραν, καὶ πᾶν ὃ ἐὰν ἀνέβη ἐν τῇ κρεάγρᾳ, ἐλάμβανεν ἑαυτῷ ὁ ἱερεύς· κατὰ τάδε ἐποίουν παντὶ Ἰσραὴλ τοῖς ἐρχομένοις θῦσαι Κυρίῳ ἐν Σηλώμ.
\vs{15}Καὶ πρὶν θυμιαθῆναι τὸ στέαρ, ἤρχετο τὸ παιδάριον τοῦ ἱερέως, καὶ ἔλεγε τῷ ἀνδρὶ τῷ θύοντι, δὸς κρέας ὀπτῆσαι τῷ ἱερεῖ, καὶ οὐ μὴ λάβω παρὰ σοῦ κρέας ἑφθὸν ἐκ τοῦ λέβητος.
\vs{16}Καὶ ἔλεγεν ὁ ἀνὴρ ὁ θύων, θυμιαθήτω πρῶτον ὡς καθήκει τὸ στέαρ, καὶ λάβε σεαυτῷ ἐκ πάντων ὧν ἐπιθυμεῖ ἡ ψυχή σου· καὶ εἶπεν, οὐχί· ὅτι νῦν δώσεις· καὶ ἐὰν μὴ, λήψομαι κραταιῶς.
\vs{17}Καὶ ἦν ἡ ἁμαρτία ἐνώπιον Κυρίου τῶν παιδαρίων μεγάλη σφόδρα, ὅτι ἠθέτουν τὴν θυσίαν Κυρίου.

\vs{18}Καὶ Σαμουὴλ ἦν λειτουργῶν ἐνώπιον Κυρίου, παιδάριον περιεζωσμένον ἐφοὺδ βάδ·
\vs{19}Καὶ διπλοΐδα μικρὰν ἐποίησεν αὐτῷ ἡ μήτηρ αὐτοῦ, καὶ ἀνέφερεν αὐτῷ ἐξ ἡμερῶν εἰς ἡμέρας ἐν τῷ ἀναβαίνειν αὐτὴν μετὰ τοῦ ἀνδρὸς αὐτῆς θῦσαι τὴν θυσίαν τῶν ἡμερῶν.
\vs{20}Καὶ εὐλόγησεν Ἡλὶ τὸν Ἑλκανὰ καὶ τὴν γυναῖκα αὐτοῦ, λέγων, ἀποτίσαι σοι Κύριος σπέρμα ἐκ τῆς γυναικὸς ταύτης, ἀντὶ τοῦ χρέους οὗ ἔχρησας τῷ Κυρίῳ· καὶ ἀπῆλθεν ὁ ἄνθρωπος εἰς τὸν τὸπον αὐτοῦ.

\vs{21}Καὶ ἐπεσκέψατο Κύριος τὴν Ἄνναν, καὶ ἔτεκεν ἔτι τρεῖς υἱοὺς, καὶ δύο θυγατέρας· καὶ ἐμεγαλύνθη τὸ παιδάριον Σαμουὴλ ἐνώπιον Κυρίου.

\vs{22}Καὶ Ἡλὶ πρεσβύτης σφόδρα· καὶ ἤκουσεν ἃ ἐποίουν οἱ υἱοὶ αὐτοῦ τοῖς υἱοῖς Ἰσραήλ·
\vs{23}Καὶ εἶπεν αὐτοῖς, ἱνατί ποιεῖτε κατὰ τὸ ῥῆμα τοῦτο, ὃ ἐγὼ ἀκούω ἐκ στόματος παντὸς τοῦ λαοῦ Κυρίου;
\vs{24}Μὴ τέκνα, ὅτι οὐκ ἀγαθὴ ἡ ἀκοὴ ἣν ἐγὼ ἀκούω· μὴ ποιεῖτε οὕτως, ὅτι οὐκ ἀγαθαὶ αἱ ἀκοαὶ ἃς ἐγὼ ἀκούω τοῦ μὴ δουλεύειν λαὸν Θεῷ.
\vs{25}Ἐὰν ἁμαρτάνων ἁμάρτῃ ἀνὴρ εἰς ἄνδρα, καὶ προσεύξονται ὑπὲρ αὐτοῦ πρὸς Κύριον· καὶ ἐὰν τῷ Κυρίῳ ἁμάρτῃ, τίς προσεύξεται ὑπὲρ αὐτοῦ; Καὶ οὐκ ἤκουον τῆς φωνῆς τοῦ πατρὸς αὐτῶν, ὅτι βουλόμενος ἐβούλετο Κύριος διαφθεῖραι αὐτούς.
\vs{26}Καὶ τὸ παιδάριον Σαμουὴλ ἐπορεύετο, καὶ ἦν ἀγαθὸν μετὰ Κυρίου καὶ μετὰ ανθρώπων.

\vs{27}Καὶ ἦλθεν ὁ ἄνθρωπος Θεοῦ πρὸς Ἡλὶ, καὶ εἶπε, τάδε λέγει Κύριος, ἀποκαλυφθεὶς ἀπεκαλύφθην πρὸς οἶκον τοῦ πατρός σου, ὄντων αὐτῶν ἐν γῇ Αἰγύπτῳ δούλων τῷ οἴκῳ Φαραὼ.
\vs{28}Καὶ ἐξελεξάμην τὸν οἶκον τοῦ πατρός σου ἐκ πάντων τῶν σκήπτρων Ἰσραὴλ ἐμοὶ ἱερατεύειν, τοῦ ἀναβαίνειν ἐπὶ θυσιαστήριόν μου, καὶ θυμιᾷν θυμίαμα, καὶ αἴρειν ἐφούδ· καὶ ἔδωκα τῷ οἴκῳ τοῦ πατρός σου τὰ πάντα τοῦ πυρὸς υἱῶν Ἰσραὴλ εἰς βρῶσιν.
\vs{29}Καὶ ἱνατί ἐπέβλεψας ἐπὶ τὸ θυμίαμά μου καὶ εἰς τὴν θυσίαν μου ἀναιδεῖ ὀφθαλμῷ; καὶ ἐδόξασας τοὺς υἱούς σου ὑπὲρ ἐμὲ ἐνευλογεῖσθαι ἀπαρχῆς πάσης θυσίας τοῦ Ἰσραὴλ ἔμπροσθέν μου;
\vs{30}Διὰ τοῦτο τάδε λέγει Κύριος ὁ Θεὸς Ἰσραὴλ, εἶπα, ὁ οἶκός σου καὶ ὁ οἶκος τοῦ πατρός σου διελεύσεται ἐνώπιόν μου ἕως αἰῶνος· καὶ νῦν φησι Κὺριος, μηδαμῶς ἐμοὶ, ὅτι ἀλλʼ ἢ τοὺς δοξάζοντάς με δοξάσω, καὶ ὁ ἐξουθενῶν με ἀτιμωθήσεται.

\vs{31}Ἰδοὺ ἔρχονται ἡμέραι, καὶ ἐξολοθρεύσω τὸ σπέρμα σου καὶ τὸ σπέρμα οἴκου πατρός σου·
\vs{32}Καὶ οὐκ ἔσται σοι πρεσβύτης ἐν οἴκῳ μου πάσας τὰς ἡμέρας.
\vs{33}Καὶ ἄνδρα οὐκ ἐξολοθρεύσω σοι ἀπὸ τοῦ θυσιαστηρίου μου, ἐκλείπειν τοὺς ὀφθαλμοὺς αὐτοῦ, καὶ καταῤῥεῖν τὴν ψυχὴν αὐτοῦ· καὶ πᾶς περισσεύων οἴκου σου πεσοῦνται ἐν ῥομφαίᾳ ἀνδρῶν.
\vs{34}Καὶ τοῦτό σοι τὸ σημεῖον ὃ ἥξει ἐπὶ τοὺς δύο υἱούς σου, Ὀφνὶ καὶ Φινεὲς, ἐν μιᾷ ἡμέρᾳ ἀποθανοῦνται ἀμφότεροι.
\vs{35}Καὶ ἀναστήσω ἐμαυτῷ ἱερέα πιστὸν, ὃς πάντα τὰ ἐν τῇ καρδίᾳ μου καὶ τὰ ἐν τῇ ψυχῇ μου ποιήσει καὶ οἰκοδομήσω αὐτῷ οἶκον πιστὸν, καὶ διελεύσεται ἐνώπιον χριστοῦ μου πάσας τὰς ἡμέρας.
\vs{36}Καὶ ἔσται ὁ περισσεύων ἐν οἴκῳ σον, ἥξει προσκυνεῖν αὐτῷ ὀβολοῦ ἀργυρίου, λέγων, παράῤῥιψόν με ἐπὶ μίαν τῶν ἱερατειῶν σου φαγεῖν ἄρτον.

\ch{3}
Καὶ τὸ παιδάριον Σαμουὴλ ἦν λειτουργῶν τῷ Κυρίῳ ἐνώπιον Ἡλὶ τοῦ ἱερέως· καὶ ῥῆμα Κυρίου ἦν τίμιον ἐν ταῖς ἡμέραις ἐκείναις, οὐκ ἦν ὅρασις διαστέλλουσα.

\vs{2}Καὶ ἐγένετο ἐν τῇ ἡμέρα ἐκείνῃ, καὶ Ἡλὶ ἐκάθευδεν ἐν τῷ τόπῳ αὐτοῦ, καὶ οἱ ὀφθαλμοὶ αὐτοῦ ἤρξαντο βαρύνεσθαι, καὶ οὐκ ἠδύναντο βλέπειν.
\vs{3}Καὶ ὁ λύχνος τοῦ Θεοῦ πρὶν ἐπισκευασθῆναι, καὶ Σαμουὴλ ἐκάθευδεν ἐν τῷ ναῷ, οὗ ἡ κιβωτὸς τοῦ Θεοῦ.
\vs{4}Καὶ ἐκάλεσε Κύριος, Σαμουὴλ Σαμουήλ· καὶ εἶπεν, ἰδοὺ ἐγώ.
\vs{5}Καὶ ἔδραμε πρὸς Ἡλὶ, καὶ εἶπεν, ἰδοὺ ἐγὼ, ὅτι κέκληκάς με· καὶ εἶπεν, οὐ κέκληκά σε, ἀνάστρεφε, κάθευδε· καὶ ἀνέστρεψε καὶ ἐκάθευδε.
\vs{6}Καὶ προσέθετο Κύριος, καὶ ἐκάλεσε, Σαμουὴλ Σαμουήλ· καὶ ἐπορεύθη πρὸς Ἡλὶ τὸ δεύτερον, καὶ εἶπεν, ἰδοὺ ἐγὼ, ὅτι κέκληκάς με· καὶ εἶπεν, οὐ κέκληκά σε, ἀνάστρεφε, κάθευδε.
\vs{7}Καὶ Σαμουὴλ πρὶν γνῶναι Θεὸν, καὶ ἀποκαλυφθῆναι αὐτῷ ῥῆμα Κυρίου.
\vs{8}Καὶ προσέθετο Κύριος καλέσαι Σαμουὴλ ἐν τρίτῳ· καὶ ἀνέστη καὶ ἐπορεύθη πρὸς Ἡλὶ, καὶ εἶπεν, ἰδοὺ ἐγὼ, ὅτι κέκληκάς με· καὶ ἐσοφίσατο Ἡλὶ ὅτι Κύριος κέκληκε τὸ παιδάριον.
\vs{9}Καὶ εἶπεν, ἀνάστρεφε, κάθευδε τέκνον· καὶ ἔσται ἐὰν καλέσῃ σε, καὶ ἐρεῖς, λάλει, ὅτι ἀκούει ὁ δοῦλός σου· καὶ ἐπορεύθη Σαμουὴλ, καὶ ἐκοιμήθη ἐν τῷ τόπῳ αὐτοῦ.
\vs{10}Καὶ ἦλθε Κύριος καὶ κατέστη, καὶ ἐκάλεσεν αὐτὸν ὡς ἅπαξ καὶ ἅπαξ· καὶ εἶπε Σαμουὴλ, λάλει, ὅτι ἀκούει ὁ δοῦλός σου.

\vs{11}Καὶ εἶπε Κύριος πρὸς Σαμουὴλ, ἰδοὺ ἐγὼ ποιῶ τὰ ῥήματά μου ἐν Ἰσραὴλ, παντὸς ἀκούοντος αὐτὰ, ἠχήσει ἀμφότερα τὰ ὦτα αὐτοῦ.
\vs{12}Ἐν τῇ ἡμέρᾳ ἐκείνῃ ἐπεγερῶ ἐπὶ Ἡλὶ πάντα ὅσα ἐλάλησα εἰς τὸν οἶκον αὐτοῦ, ἄρξομαι καὶ ἐπιτελέσω.
\vs{13}Καὶ ἀνήγγελκα αὐτῷ ὅτι ἐκδικῶ ἐγὼ τὸν οἶκον αὐτοῦ ἕως αἰῶνος ἐν ἀδικίαις υἱῶν αὐτοῦ, ὅτι κακολογοῦντες Θεὸν οἱ υἱοὶ αὐτοῦ, καὶ οὐκ ἐνουθέτει αὐτούς.
\vs{14}Καὶ οὐδʼ οὕτως· ὤμοσα τῷ οἴκῳ Ἡλὶ, εἰ ἐξιλασθήσεται ἀδικία οἴκου Ἡλὶ, ἐν θυμιάματι καὶ ἐν θυσίαις ἕως αἰῶνος.

\vs{15}Καὶ κοιμᾶται Σαμουὴλ ἕως πρωῒ, καὶ ὤρθρισε τοπρωῒ καὶ ἤνοιξε τὰς θύρας οἴκου Κυρίου· καὶ Σαμουὴλ ἐφοβήθη ἀπαγγεῖλαι τὴν ὅρασιν.
\vs{16}Καὶ εἶπεν Ἡλὶ πρὸς Σαμουὴλ, Σαμουὴλ τέκνον· καὶ εἶπεν, ἰδοὺ ἐγώ.
\vs{17}Καὶ εἶπε, τί τὸ ῥήμα τὸ λαληθὲν πρὸς σέ; μὴ δὴ κρήψῃς ἀπʼ ἐμοῦ· τάδε ποιήσαι σοι ὁ Θεὸς καὶ τάδε προσθείη, ἐὰν κρύψῃς ἀπʼ ἐμοῦ ῥῆμα ἐκ πάντων τῶν λόγων τῶν λαληθέντων σοι ἐν τοῖς ὠσί σου.
\vs{18}Καὶ ἀπήγγειλε Σαμουὴλ πάντας τοὺς λόγους, καὶ οὐκ ἔκρυψεν ἀπʼ αὐτοῦ· καὶ εἶπεν Ἡλὶ, Κύριος αὐτὸς, τὸ ἀγαθὸν ἐνώπιον αὐτοῦ ποιήσει.

\vs{19}Καὶ ἐμεγαλύνθη Σαμουὴλ, καὶ ἦν Κύριος μετʼ αὐτοῦ, καὶ οὐκ ἔπεσεν ἀπὸ πάντων τῶν λόγων αὐτοῦ ἐπὶ τὴν γῆν.
\vs{20}Καὶ ἔγνωσαν πᾶς Ἰσραὴλ ἀπὸ Δὰν καὶ ἕως Βηρσαβεὲ, ὅτι πιστὸς Σαμονὴλ εἰς προφήτην τῷ Κυρίῳ.
\vs{21}Καὶ προσέθετο Κύριος δηλωθῆναι ἐν Σηλὼμε, ὅτι ἀπεκαλύφθη Κύριος πρὸς Σαμουήλ· καὶ ἐπιστεύθη Σαμουὴλ τοῦ προφήτης γενέσθαι τῷ Κυρίῳ εἰς πάντα Ἰσραὴλ ἀπʼ ἄκρων τῆς γῆς καὶ ἕως ἄκρων· καὶ Ἡλὶ πρεσβύτης σφόδρα, καὶ οἱ υἱοὶ αὐτοῦ πορεύομενοι ἐπορεύοντο, καὶ πονηρὰ ἡ ὁδὸς αὐτῶν ἐνώπιον Κυρίου.

\ch{4}
Καὶ ἐγενήθη ἐν ταῖς ἡμέραις ἐκείναις, καὶ συναθροίζονται ἀλλόφυλοι ἐπὶ Ἰσραὴλ εἰς πόλεμον· καὶ ἐξῆλθεν Ἰσραὴλ εἰς ἀπάντησιν αὐτοῖς εἰς πόλεμον, καὶ παρεμβάλλουσιν ἐπὶ Ἀβενέζερ· καὶ οἱ ἀλλόφυλοι παρεμβάλλουσιν ἐν Ἀφέκ.
\vs{2}Καὶ παρατάσσονται ἀλλόφυλοι εἰς πόλεμον ἐπὶ Ἰσραήλ· καὶ ἔκλινεν ὁ πόλεμος, καὶ ἔπταισεν ἀνὴρ Ἰσραὴλ ἐνώπιον ἀλλοφύλων, καὶ ἐπλήγησαν ἐν τῇ παρατάξει ἐν ἀγρῷ τέσσαρες χιλιάδες ἀνδρῶν.

\vs{3}Καὶ ἦλθεν ὁ λαὸς εἰς τὴν παρεμβολὴν, καὶ εἶπαν οἱ πρεσβύτεροι Ἰσραὴλ, κατὰ τί ἔπταισεν ἡμᾶς Κύριος σήμερον ἐνώπιον ἀλλοφύλων; λάβωμεν τὴν κιβωτὸν τοῦ Θεοῦ ἡμῶν ἐκ Σηλώμ, καὶ ἐξελθέτω ἐκ μέσου ἡμῶν, καὶ σώσει ἡμᾶς ἐκ χειρὸς ἐχθρῶν ἡμῶν.

\vs{4}Καὶ ἀπέστειλεν ὁ λαὸς εἰς Σηλὼμ, καὶ αἴρουσιν ἐκεῖθεν τὴν κιβωτὸν Κυρίου καθημένου χερουβίμ· καὶ ἀμφότεροι οἱ υἱοὶ Ἡλὶ μετὰ τῆς κιβωτοῦ, Ὀφνὶ καὶ Φινεές.
\vs{5}Καὶ ἐγενήθη ὡς ἦλθεν ἡ κιβωτὸς Κυρίου εἰς τὴν παρεμβολὴν, καὶ ἀνέκραξεν πᾶς Ἰσραὴλ φωνῇ μεγάλῃ, καὶ ἤχησεν ἡ γῆ.
\vs{6}Καὶ ἤκουσαν οἱ ἀλλόφυλοι τῆς κραυγῆς, καὶ εἶπον οἱ ἀλλόφυλοι, τίς ἡ κραυγὴ ἡ μεγάλη αὕτη ἐν τῇ παρεμβολῇ τῶν Ἑβραίων; καὶ ἔγνωσαν ὅτι κιβωτὸς Κυρίου ἥκει εἰς τὴν παρεμβολήν.
\vs{7}Καὶ ἐφοβήθησαν οἱ ἀλλόφυλοι, καὶ εἶπον, οὗτοι οἱ θεοὶ ἥκασι πρὸς αὐτοὺς εἰς τὴν παρεμβολήν. Οὐαὶ ἡμῖν, ἐξελοῦ ἡμᾶς Κύριε σήμερον, ὅτι οὐ γέγονε τοιαύτη ἐχθὲς καὶ τρίτην·
\vs{8}οὐαὶ ἡμῖν, τίς ἐξελεῖται ἡμᾶς ἐκ χειρὸς τῶν θεῶν τῶν στερεῶν τούτων; οὗτοι οἱ θεοὶ οἱ πατάξαντες τὴν Αἴγυπτον ἐν πάσῃ πληγῇ, καὶ ἐν τῇ ἐρήμῳ.
\vs{9}Κραταιοῦσθε καὶ γίνεσθε εἰς ἄνδρας ἀλλόφυλοι, ὅπως μὴ δουλεύσητε τοῖς Ἑβραίοις, καθὼς ἐδούλευσαν ἡυῖν, καὶ ἔσεσθε εἰς ἄνδρας, καὶ πολεμήσατε αὐτούς.

\vs{10}Καὶ ἐπολέμησαν αὐτούς· καὶ πταίει ἀνὴρ Ἰσραὴλ, καὶ ἔφυγεν ἕκαστος εἰς σκήνωμα αὐτοῦ· καὶ ἐγένετο πληγὴ μεγάλη σφόδρα· καὶ ἔπεσον ἐξ Ἰσραὴλ τριάκοντα χιλιάδες ταγμάτων.
\vs{11}Καὶ κιβωτὸς τοῦ Θεοῦ ἐλήφθη, καὶ ἀμφότεροι οἱ υἱοὶ Ἡλὶ ἀπέθανον, Ὀφνὶ καὶ Φινεές.

\vs{12}Καὶ ἔδραμεν ἀνὴρ Ἰεμιναῖος ἐκ τῆς παρατάξεως, καὶ ἦλθεν εἰς Σηλὼμ ἐν τῇ ἡμέρᾳ ἐκείνῃ, καὶ τὰ ἱμάτια αὐτοῦ διεῤῥωγότα, καὶ γῆ ἐπὶ τῆς κεφαλῆς αὐτοῦ.
\vs{13}Καὶ ἦλθε, καὶ ἰδοὺ Ἡλὶ ἐπὶ τοῦ δίφρου παρὰ τὴν πύλην σκοπεύων τὴν ὁδὸν, ὅτι ἦν καρδία αὐτοῦ ἐξεστηκυῖα περὶ τῆς κιβωτοῦ τοῦ Θεοῦ· καὶ ὁ ἄνθρωπος εἰσῆλθεν εἰς τὴν πόλιν ἀπαγγεῖλαι· καὶ ἀνεβόησεν ἡ πόλις.
\vs{14}Καὶ ἤκουσεν Ἡλὶ τὴν φωνὴν τῆς βοῆς, καὶ εἶπε, τίς ἡ φωνῆ τῆς βοῆς ταύτης; καὶ ὁ ἄνθρωπος σπεύσας εἰσῆλθε, καὶ ἀπήγγειλε τῷ Ἡλί.
\vs{15}Καὶ Ἡλὶ υἱὸς ἐνενήκοντα ἐτῶν, καὶ οἱ ὀφθαλμοὶ αὐτοῦ ἐπανέστησαν, καὶ οὐκ ἐπέβλεπε. Καὶ εἶπεν Ἡλὶ τοῖς ἀνδράσιν τοῖς περιεστηκόσιν αὐτῷ, τίς ἡ φωνὴ τοῦ ἤχου τούτου;
\vs{16}καὶ ὁ ἀνὴρ σπεύσας προσῆλθε πρὸς Ἡλὶ, καὶ εἶπεν αὐτῷ, ἐγώ εἰμι ὁ ἥκων ἐκ τῆς παρεμβολῆς, κᾀγὼ πέφευγα ἐκ τῆς παρατάξεως σήμερον· καὶ εἶπεν Ἡλὶ, τί τὸ γεγονὸς ῥῆμα, τέκνον;
\vs{17}Καὶ ἀπεκρίθη τὸ παιδάριον, καὶ εἶπε, πέφευγεν ἀνὴρ Ἰσραὴλ ἐκ προσώπου ἀλλοφύλων, καὶ ἐγένετο πληγὴ μεγάλη ἐν τῷ λαῷ, καὶ ἀμφότεροι οἱ υἱοί σου τεθνήκασι, καὶ ἡ κιβωτὸς τοῦ Θεοῦ ἐλήφθη.
\vs{18}Καὶ ἐγένετο ὡς ἐμνήσθη τῆς κιβωτοῦ τοῦ Θεοῦ, καὶ ἔπεσεν ἀπὸ τοῦ δίφρου ὀπισθίως ἐχόμενος τῆς πύλης, καὶ συνετρίβη ὁ νῶτος αὐτοῦ, καὶ ἀπέθανεν, ὅτι πρεσβύτης ὁ ἄνθρωπος καὶ βαρύς· καὶ αὐτὸς ἔκρινε τὸν Ἰσραὴλ εἴκοσι ἔτη.

\vs{19}Καὶ νύμφη αὐτοῦ γυνὴ Φινεὲς συνειληφυῖα τοῦ τεκεῖν, καὶ ἤκουσε τὴν ἀγγελίαν, ὅτι ἐλήφθη ἡ κιβωτὸς τοῦ Θεοῦ, καὶ ὅτι τέθνηκεν ὁ πενθερὸς αὐτῆς καὶ ὁ ἀνὴρ αὐτῆς, καὶ ἔκλαυσε καὶ ἔτεκεν, ὅτι ἐπεστράφησαν ἐπʼ αὐτὴν ὠδῖνες αὐτῆς·
\vs{20}Καὶ ἐν τῷ καιρῷ αὐτῆς ἀποθνήσκει· καὶ εἶπον αὐτῇ αἱ γυναῖκες αἱ παρεστηκυῖαι αὐτῇ, μὴ φοβοῦ, ὅτι υἱὸν τέτοκασ· καὶ οὐκ ἀπεκρίθη, καὶ οὐκ ἐνόησαν ἡ καρδία αὐτῆς.
\vs{21}Καὶ ἐκάλεσε τὸ παιδάριον Οὐαιβαρχαβώθ ὑπὲρ τῆς κιβωτοῦ τοῦ Θεοῦ, καὶ ὑπὲρ τοῦ πενθεροῦ αὐτῆς, καὶ ὑπὲρ τοῦ ἀνδρὸς αὐτῆς.
\vs{22}Καὶ εἶπαν, ἀπῴκισται δόξα Ἰσραὴλ ἐν τῷ ληφθῆναι τὴν κιβωτὸν Κυρίου.

\ch{5}
Καὶ ἀλλόφυλοι ἔλαβον τὴν κιβωτὸν τοῦ Θεοῦ, καὶ εἰσήνεγκαν αὐτὴν ἐξ Ἀβενεζὲρ εἰς Ἄζωτον.
\vs{2}Καὶ ἔλαβον ἀλλόφυλοι τὴν κιβωτὸν Κυρίου, καὶ εἰσήνεγκαν αὐτὴν εἰς οἴκον Δαγὼν, καὶ παρέστησαν αὐτὴν παρὰ Δαγών.
\vs{3}Καὶ ὤρθρισαν οἱ Ἀζώτιοι, καὶ εἰσῆλθον εἰς οἶκον Δαγών· καὶ εἶδον, καὶ ἰδοὺ Δαγὼν πεπτωκὼς ἐπὶ πρόσωπον αὐτοῦ ἐνώπιον κιβωτοῦ τοῦ Θεοῦ· καὶ ἤγειραν τὸν Δαγὼν, καὶ κατέστησαν εἰς τὸν τόπον αὐτοῦ· καὶ ἐβαρύνθη χεὶρ Κυρίου ἐπὶ τοὺς Ἀζωτίους, καὶ ἐβασάνισεν αὐτούς· καὶ ἐπάταξεν αὐτοὺς εἰς τὰς ἕδρας αὐτῶν, τὴν Ἄζωτον καὶ τὰ ὅρια αὐτῆς.
\vs{4}Καὶ ἐγένετο ὅτε ὤρθρισαν τοπρωῒ, καὶ ἰδοὺ Δαγὼν, πεπτωκὼς ἐπὶ πρόσωπον αὐτοῦ ἐνώπιον κιβωτοῦ διαθήκης Κυρίου· καὶ κεφαλὴ Δαγὼν καὶ ἀμφότερα τὰ ἴχνη χειρῶν αὐτοῦ ἀφῃρημένα ἐπὶ τὰ ἐμπρόσθια ἁμαφὲθ ἕκαστοι, καὶ ἀμφότεροι οἱ καρποὶ τῶν χειρῶν αὐτοῦ πεπτωκότες ἐπὶ τὸ πρόθυρον, πλὴν ἡ ῥάχις Δαγὼν ὑπελείφθη.
\vs{5}Διὰ τοῦτο οὐκ ἐπιβαίνουσιν οἱ ἱερεῖς Δαγὼν, καὶ πᾶς ὁ εἰσπορευόμενος εἰς οἶκον Δαγὼν, ἐπὶ βαθμὸν οἴκου Δαγὼν ἐν Ἀζώτῳ ἕως τῆς ἡμέρας ταύτης, ὅτι ὑπερβαίνοντες ὑπερβαίνουσι.

\vs{6}Καὶ ἐβαρύνθη ἡ χεὶρ Κυρίου ἐπὶ Ἄζωτον, καὶ ἐπήγαγεν αὐτοῖς, καὶ ἐξέζεσεν αὐτοῖς εἰς τὰς ναῦς, καὶ μέσον τῆς χώρας αὐτῆς ἀνεφύησαν μύες· καὶ ἐγένετο σύγχυσις θανάτου μεγάλη ἐν τῇ πόλει.
\vs{7}Καὶ εἶδον οἱ ἄνδρες Ἀζώτου ὅτι οὕτως, καὶ λέγουσιν, ὅτι οὐ καθήσεται κιβωτὸς τοῦ Θεοῦ Ἰσραὴλ μεθʼ ἡμῶν, ὅτι σκληρὰ χεὶρ αὐτοῦ ἐφʼ ἡμᾶς καὶ ἐπὶ Δαγὼν θεὸν ἡμῶν.
\vs{8}Καὶ ἀποστέλλουσι καὶ συνάγουσι τοὺς σατράπας τῶν ἀλλοφύλων πρὸς αὐτοὺς, καὶ λέγουσι, τί ποιήσωμεν τῇ κιβωτῷ Θεοῦ Ἰσραήλ; καὶ λέγουσιν οἱ Γεθαῖοι, μετελθέτω κιβωτὸς τοῦ Θεοῦ πρὸς ἡμᾶς· καὶ μετῆλθε κιβωτὸς τοῦ Θεοῦ Ἰσραὴλ εἰς Γέθ.

\vs{9}Καὶ ἐγενήθη μετὰ τὸ μετελθεῖν αὐτὴν, καὶ γίνεται χεὶρ Κυρίου τῇ πόλει, τάραχος μέγας σφόδρα· καὶ ἐπάταξε τοὺς ἄνδρας τῆς πόλεως ἀπὸ μικροῦ ἕως μεγάλου, καὶ ἐπάταξεν αὐτοὺς εἰς τὰς ἕδρας αὐτῶν· καὶ ἐποίησαν οἱ Γεθαῖοι ἑαυτοῖς ἕδρας.

\vs{10}Καὶ ἐξαποστέλλουσι τὴν κιβωτὸν τοῦ Θεοῦ εἰς Ἀσκάλωνα· καὶ ἐγενήθη ὡς εἰσῆλθε κιβωτὸς Θεοῦ εἰς Ἀσκάλωνα, καὶ ἐβόησαν οἱ Ἀσκαλωνῖται, λέγοντες, τί ἀπεστρέψατε τὴν κιβωτὸν τοῦ Θεοῦ Ἰσραὴλ πρὸς ἡμᾶς θανατῶσαι ἡμᾶς καὶ τὸν λαὸν ἡμῶν;
\vs{11}Καὶ ἐξαποστέλλουσι καὶ συνάγουσι τοὺς σατράπας τῶν ἀλλοφύλων, καὶ εἶπον, ἐξαποστείλατε τὴν κιβωτὸν τοῦ Θεοῦ Ἰσραήλ, καὶ καθισάτω εἰς τὸν τόπον αὐτῆς, καὶ οὐ μὴ θανατώσῃ ἡμᾶς καὶ τὸν λαὸν ἡμῶν· Ὅτι ἐγενήθη σύγχυσις ἐν ὅλῃ τῇ πόλει βαρεῖα σφόδρα, ὡς εἰσῆλθεν κιβωτὸς Θεοῦ Ἰσραὴλ ἐκεῖ.
\vs{12}καὶ οἱ ζῶντες καὶ οὐκ ἀποθανόντες ἐπλήγησαν εἰς τὰς ἕδρας· καὶ ἀνέβη ἡ κραυγὴ τῆς πόλεως εἰς τὸν οὐρανόν.

\ch{6}
Καὶ ἦν ἡ κιβωτὸς ἐν ἀγρῷ τῶν ἀλλοφύλων ἑπτὰ μῆνας, καὶ ἐξέζεσεν ἡ γῆ αὐτῶν μύας.
\vs{2}Καὶ καλοῦσιν ἀλλόφυλοι τοὺς ἱερεῖς καὶ τοὺς μάντεις καὶ τοὺς ἐπαοιδοὺς αὐτῶν, λέγοντες, τί ποιήσωμεν τῇ κιβωτῷ Κυρίου; γνωρίσατε ἡμῖν ἐν τίνι ἀποστελοῦμεν αὐτὴν εἰς τὸν τόπον αὐτῆς.
\vs{3}Καὶ εἶπαν, εἰ ἐξαποστέλλετε ὑμεῖς τὴν κιβωτὸν διαθήκης Κυρίου Θεοῦ Ἰσραήλ, μὴ δὴ ἐξαποστείλητε αὐτὴν κενήν, ἀλλὰ ἀποδιδόντες ἀπόδοτε αὐτῇ τῆς βασάνου, καὶ τότε ἰαθήσεσθε, καὶ ἐξιλασθήσεται ὑμῖν· μὴ οὐκ ἀποστῇ ἡ χεὶρ αὐτοῦ ἀφʼ ὑμῶν;
\vs{4}Καὶ λέγουσι, τί τὸ τῆς βασάνου ἀποδώσομεν αὐτῇ; καὶ εἶπαν, κατα ἀριθμὸν τῶν σατραπῶν τῶν ἀλλοφύλων πέντε ἕδρας χρυσᾶς, ὅτι πταῖσμα ἐν ὑμῖν καὶ τοῖς ἄρχουσιν ὑμῶν καὶ τῷ λαῷ,
\vs{5}καὶ μῦς χρυσοῦς ὁμοίωμα τῶν μυῶν ὑμῶν τῶν διαφθειρόντων τὴν γῆν· καὶ δώσετε τῷ Κυρίῳ δόξαν, ὅπως κουφίσῃ τὴν χεῖρα αὐτοῦ ἀφʼ ὑμῶν, καὶ ἀπὸ τῶν θεῶν ὑμῶν, καὶ ἀπὸ τῆς γῆς ὑμῶν.
\vs{6}Καὶ ἱνατί βαρύνετε τὰς καρδίας ὑμῶν, ὡς ἐβάρυνεν Αἴγυπτος καὶ Φαραὼ τὴν καρδίαν αὐτῶν; οὐχὶ ὅτε ἐνέπαιξεν αὐτοῖς, ἐξαπέστειλαν αὐτοὺς καὶ ἀπῆλθον;

\vs{7}Καὶ νῦν λάβετε καὶ ποιήσατε ἅμαξαν καινὴν, καὶ δύο βόας πρωτοτοκούσας ἄνευ τῶν τέκνων· καὶ ζεύξατε τὰς βόας ἐν τῇ ἁμάξῃ, καὶ ἀπαγάγετε τὰ τὲκνα ἀπὸ ὄπισθεν αὐτῶν εἰς οἴκον.
\vs{8}Καὶ λήψεσθε τὴν κιβωτὸν, καὶ θήσετε αὐτὴν ἐπὶ τὴν ἅμαξαν, καὶ τὰ σκεύη τὰ χρυσᾶ ἀποδώσετε αὐτῇ τῆς βασάνου, καὶ θήσετε ἐν θέματι βερσεχθὰν ἐκ μέρους αὐτῆς· καὶ ἐξαποστελεῖτε αὐτὴν, καὶ ἀπελάσατε αὐτήν, καὶ ἀπελεύσεσθε.
\vs{9}Καὶ ὄψεσθε, εἰ ὁδὸν ὁρίων αὐτῆς πορεύσεται μετὰ Βαιθσαμὺς, αὐτὸς πεποίηκεν ἡμῖν τὴν κακίαν τὴν μεγάλην ταύτην· καὶ ἐὰν μὴ, καὶ γνωσόμεθα ὅτι οὐ χεὶρ αὐτοῦ ἧπται ἡμῶν, ἀλλὰ σύμπτωμα τοῦτο γέγονεν ἡμῖν.

\vs{10}Καὶ ἐποίησαν οἱ ἀλλόφυλοι οὕτω· καὶ ἔλαβον δύο βόας πρωτοτοκούσας, καὶ ἔζευξαν αὐτὰς ἐν τῇ ἁμάξῃ, καὶ τὰ τέκνα αὐτῶν ἀπεκώλυσαν εἰς οἶκον·
\vs{11}Καὶ ἔθεντο τὴν κιβωτὸν Κυρίου ἐπὶ τὴν ἅμαξαν, καὶ τὸ θέμα ἐργὰβ καὶ τοὺς μῦς τοὺς χρυσοῦς.
\vs{12}Καὶ κατεύθυναν αἱ βόες ἐν τῇ ὁδῷ εἰς ὁδὸν Βαιθσαμὺς, ἐν τρίβῳ ἐνὶ ἐπορεύοντο καὶ ἐκοπίων, καὶ οὐ μεθίσταντο δεξιὰ οὐδὲ ἀριστερά· καὶ οἱ σατράπαι τῶν ἀλλοφύλων ἐπορεύοντο ὀπίσω αὐτῆς ἕως ὁρίων Βαιθσαμύς.
\vs{13}Καὶ οἱ ἐν Βαιθσαμὺς ἐθέριζον θερισμὸν πυρῶν ἐν κοιλάδι· καὶ ᾖραν ὀφθαλμοὺς αὐτῶν, καὶ εἶδον κιβωτὸν Κυρίου, καὶ ηὐφράνθησαν εἰς ἀπάντησιν αὐτῆς.
\vs{14}Καὶ ἡ ἅμαξα εἰσῆλθεν εἰς ἀγρὸν Ὠσῆὲ τὸν ἐν Βαιθσαμὺς, καὶ ἔστησαν ἐκεῖ παρʼ αὐτῇ λίθον μέγαν· καὶ σχίζουσι τὰ ξύλα τῆς ἁμάξης, καὶ τὰς βόας ἀνήνεγκαν εἰς ὁλοκαύτωσιν τῷ κυρίῳ.
\vs{15}Καὶ οἱ Λευῖται ἀνήνεγκαν τὴν κιβωτὸν τοῦ Κυρίου, καὶ τὸ θέμα ἐργὰβ μετʼ αὐτῆς, καὶ τὰ ἐπʼ αὐτῆς σκεύη τὰ χρυσᾶ, καὶ ἔθεντο ἐπὶ τοῦ λίθου τοῦ μεγάλου· καὶ οἱ ἄνδρες Βαιθσαμὺς ἀνήνεγκαν ὁλοκαυτώσεις καὶ θυσίας ἐν τῇ ἡμέρᾳ ἐκείνῃ τῷ Κυρίῳ.
\vs{16}Καὶ οἱ πέντε σατράπαι τῶν ἀλλοφύλων ἑώρων, καὶ ἀνέστρεψαν εἰς Ἀσκάλωνα τῇ ἡμέρᾳ ἐκείνῃ.

\vs{17}Καὶ αὗται αἱ ἕδραι αἱ χρυσαῖ, ἃς ἀπέδωκαν οἱ ἀλλόφυλοι τῆς βασάνου τῷ Κυρίῳ· τῆς Ἀζώτου μίαν, τῆς Γάζης μίαν, τῆς Ἀσκάλωνος μίαν, τῆς Γὲθ μίαν, τῆς Ἀκκαρὼν μίαν.
\vs{18}Καὶ μῦς οἱ χρυσοῖ κατʼ ἀριθμὸν πασῶν πόλεων τῶν ἀλλοφύλων τῶν πέντε σατραπῶν ἐκ πόλεως ἐστερεωμένης καὶ ἕως κώμης τοῦ Φερεζαίου, καὶ ἕως λίθου τοῦ μεγάλου, οὗ ἐπέθηκαν ἐπʼ αὐτοῦ τὴν κιβωτὸν διαθήκης Κυρίου, τοῦ ἐν ἀγρῷ Ὠσηὲ τοῦ Βαιθσαμυσίτου.

\vs{19}Καὶ οὐκ ἠσμένισαν οἱ υἱοὶ Ἰεχονίου ἐν τοῖς ἀνδράσι Βαιθσαμὺς, ὅτι εἶδαν κιβωτὸν Κυρίου· καὶ ἐπάταξεν ἐν αὐτοῖς ἑβδομήκοντα ἄνδρας, καὶ πεντήκοντα χιλιάδας ἀνδρῶν· καὶ ἐπένθησεν ὁ λαός, ὅτι ἐπάταξε Κύριος ἐν τῷ λαῷ πληγὴν μεγάλην σφόδρα.
\vs{20}Καὶ εἶπαν οἱ ἄνδρες οἱ ἐκ Βαιθσαμὺς, τίς δυνήσεται διελθεῖν ἐνώπιον Κυρίου τοῦ Θεοῦ τοῦ ἁγίου τούτου; καὶ πρὸς τίνα ἀναβήσεται κιβωτὸς Κυρίου ἀφʼ ἡμῶν;

\vs{21}Καὶ ἀποστέλλουσιν ἀγγέλους πρὸς τοὺς κατοικοῦντας Καριαθιαρεὶμ, λέγοντες, ἀπεστρόφασιν ἀλλόφυλοι τὴν κιβωτὸν Κυρίου, κατάβητε καὶ ἀναγάγετε αὐτὴν πρὸς ἑαυτούς.

\ch{7}
Καὶ ἔρχονται οἱ ἄνδρες Καριαθιαρὶμ, καὶ ἀνάγουσι τὴν κιβωτὸν διαθήκης Κυρίου· καὶ εἰσάγουσιν αὐτὴν εἰς οἶκον ʼΑμιναδὰβ τὸν ἐν τῷ βουνῷ· καὶ τὸν Ἐλεάζαρ τὸν υἱὸν αὐτοῦ ἡγίασαν φυλάσσειν τὴν κιβωτὸν διαθήκης Κυρίου.

\vs{2}Καὶ ἐγενήθη ἀφʼ ἧς ἡμέρας ἦν ἡ κιβωτὸς ἐν Καριαθιαρὶμ, ἐπλήθυναν αἱ ἡμέραι, καὶ ἐγένετο εἴκοσι ἔτη· καὶ ἐπέβλεψε πᾶς οἶκος Ἰσραὴλ ὀπίσω Κυρίου.
\vs{3}Καὶ εἶπε Σαμουὴλ πρὸς πάντα οἶκον Ἰσραὴλ, λέγων, εἰ ἐν ὅλῃ καρδίᾳ ὑμῶν ὑμεῖς ἐπιστρέφετε πρὸς Κύριον, περιέλετε θεοὺς ἀλλοτρίους ἐκ μέσου ὑμῶν, καὶ τὰ ἄλση, καὶ ἑτοιμάσατε τὰς καρδίας ὑμῶν πρὸς Κύριον, καὶ δουλεύσατε αὐτῷ μόνῳ, καὶ ἐξελεῖται ὑμᾶς ἐκ χειρὸς ἀλλοφύλων.
\vs{4}Καὶ περιεῖλον οἱ υἱοὶ Ἰσραὴλ τὰς Βααλὶμ καὶ τὰ ἄλση ʼΑσταρὼθ, καὶ ἐδούλευσαν Κυρίῳ μόνῳ.

\vs{5}Καὶ εἶπε Σαμουήλ, ἀθροίσατε πάντα Ἰσραὴλ εἰς Μασσηφὰθ, καὶ προσεύξομαι περὶ ὑμῶν πρὸς Κύριον.
\vs{6}Καὶ συνήχθησαν εἰς Μασσηφὰθ, καὶ ὑδρεύονται ὕδωρ, καὶ ἐξέχεαν ἐνώπιον Κυρίου ἐπὶ τὴν γῆν· καὶ ἐνήστευσαν ἐν τῇ ἡμέρᾳ ἐκείνῃ, καὶ εἶπαν, ἡμαρτήκαμεν ἐνώπιον Κυρίου· καὶ ἐδίκαζε Σαμουὴλ τοὺς υἱοὺς Ἰσραὴλ εἰς Μασσηφάθ.

\vs{7}Καὶ ἤκουσαν οἱ ἀλλόφυλοι ὅτι συνηθροίσθησαν πάντες οἱ υἱοὶ Ἰσραὴλ εἰς Μασσηφάθ· καὶ ἀνέβησαν σατράπαι ἀλλοφύλων ἐπὶ Ἰσραήλ· καὶ ἀκούουσιν οἱ υἱοὶ Ἰσραήλ, καὶ ἐφοβήθησαν ἀπὸ προσώπου ἀλλοφύλων.
\vs{8}Καὶ εἶπαν οἱ υἱοὶ Ἰσραὴλ πρὸς Σαμουήλ, μὴ παρασιωπήσῃς ἀφʼ ἡμῶν τοῦ μὴ βοᾷν πρὸς Κύριον Θεόν σου, καὶ σώσει ἡμᾶς ἐκ χειρὸς ἀλλοφύλων.
\vs{9}Καὶ ἔλαβε Σαμουὴλ ἄρνα γαλαθηνὸν ἕνα, καὶ ἀνήνεγκεν αὐτὸν ὁλοκαύτωσιν σὺν παντὶ τῷ λαῷ τῷ Κυρίῳ· καὶ ἐβόησε Σαμουὴλ πρὸς Κύριον περὶ Ἰσραὴλ, καὶ ἐπήκουσεν αὐτοῦ Κύριος.
\vs{10}Καὶ ἦν Σαμουὴλ ἀναφέρων τὴν ὁλοκαύτωσιν· καὶ ἀλλόφυλοι προσῆγον εἰς πόλεμον ἐπὶ Ἰσραήλ· καὶ ἐβρόντησε Κύριος ἐν φωνῇ μεγάλῃ ἐν τῇ ἡμέρᾳ ἐκείνῃ ἐπὶ τοὺς ἀλλοφύλους, καὶ συνεχύθησαν καὶ ἔπταισαν ἐνώπιον Ἰσραήλ.
\vs{11}Καὶ ἐξῆλθον ἄνδρες Ἰσραὴλ ἐκ Μασσηφὰθ, καὶ κατεδίωξαν τοὺς ἀλλοφύλους, καὶ ἐπάταξαν αὐτοὺς ἕως ὑποκάτω τοῦ Βαιθχόρ.

\vs{12}Καὶ ἔλαβεν Σαμουὴλ λίθον ἕνα, καὶ ἔστησεν αὐτὸν ἀναμέσον Μασσηφὰθ καὶ ἀναμέσον τῆς παλαιᾶς· καὶ ἐκάλεσε τὸ ὄνομα αὐτοῦ ʼΑβενέζερ, λίθος τοῦ βοηθοῦ· καὶ εἶπεν, ἕως ἐνταῦθα ἐβοήθησεν ἡμῖν Κύριος.

\vs{13}Καὶ ἐταπείνωσε Κύριος τοὺς ἀλλοφύλους, καὶ οὐ προσέθεντο ἔτι προσελθεῖν εἰς ὅριον Ἰσραήλ· καὶ ἐγενήθη χεὶρ Κυρίου ἐπὶ τοὺς ἀλλοφύλους πάσας τὰς ἡμέρας τοῦ Σαμουήλ.
\vs{14}Καὶ ἀπεδόθησαν αἱ πόλεις ἃς ἔλαβον οἱ ἀλλόφυλοι παρὰ τῶν υἱῶν Ἰσραὴλ, καὶ ἀπέδωκαν αὐτὰς τῷ Ἰσραὴλ ἀπὸ ʼΑσκάλωνος ἕως ʼΑζόβ· καὶ τὸ ὅριον Ἰσραὴλ ἀφείλοντο ἐκ χειρὸς ἀλλοφύλων· καὶ ἦν εἰρήνη ἀναμέσον Ἰσραὴλ καὶ ἀναμέσον τοῦ ʼΑμοῤῥαίου.

\vs{15}Καὶ ἐδίκαζεν Σαμουὴλ τὸν Ἰσραὴλ πάσας τὰς ἡμέρας τῆς ζεῆς αὐτοῦ.
\vs{16}Καὶ ἐπορεύετο κατʼ ἐνιαυτὸν ἐνιαυτὸν, καὶ ἐκύκλου Βαιθὴλ καὶ τὴν Γαλγαλὰ καὶ τὴν Μασσηφάθ· καὶ ἐδίκαζε τὸν Ἰσραὴλ ἐν πᾶσι τοῖς ἡγιασμένοις τούτοις.
\vs{17}Ἡ δὲ ἀποστροφὴ αὐτοῦ εἰς ʼΑρμαθάιμ, ὅτι ἐκεῖ ἦν ὁ οἶκος αὐτοῦ· καὶ ἐδίκαζεν ἐκεῖ τὸν Ἰσραήλ, καὶ ᾠκοδόμησεν ἐκεῖ θυσιαστήριον τῷ Κυρίῳ.

\ch{8}
Καὶ ἐγένετο ὡς ἐγήρασε Σαμουὴλ, καὶ κατέστησε τούς υἱοὺς αὐτοῦ δικαστὰς τῷ Ἰσραήλ.
\vs{2}Καὶ ταῦτα τὰ ὀνόματα τῶν υἱῶν αὐτοῦ· πρωτότοκος Ἰωήλ, καὶ ὄνομα τοῦ δευτέρου ʼΑβιά, δικασταὶ ἐν Βηρσαβεέ.
\vs{3}Καὶ οὐκ ἐπορεύθησαν οἱ υἱοὶ αὐτοῦ ἐν ὁδῷ αὐτοῦ· καὶ ἐξέκλιναν ὀπίσω τῆς συντελείας, καὶ ἐλάμβανον δῶρα, καὶ ἐξέκλινον δικαιώματα.

\vs{4}Καὶ συναθροίζονται ἄνδρες Ἰσραὴλ, καὶ παραγίνονται εἰς ʼΑρμαθαὶμ πρὸς Σαμουὴλ,
\vs{5}καὶ εἶπαν αὐτῷ, ἰδοὺ, σὺ γεγήρακας, καὶ οἱ υἱοί σου οὐ πορεύονται ἐν τῇ ὁδῷ σου· καὶ νῦν κατάστησον ἐφʼ ἡμᾶς βασιλέα δικάζειν ἡμᾶς, καθὰ καὶ τὰ λοιπὰ ἔθνη.

\vs{6}Καὶ πονηρὸν τὸ ῥῆμα ἐν ὀφθαλμοῖς Σαμουήλ, ὡς εἶπαν, δὸς ἡμῖν βασιλέα δικάζειν ἡμᾶς· καὶ προσηύξατο Σαμουὴλ πρὸς Κύριον.
\vs{7}Καὶ εἶπε Κύριος πρὸς Σαμουήλ, ἄκουε τῆς φωνῆς τοῦ λαοῦ, καθὰ ἂν λαλῶσί σοι, ὅτι οὐ σὲ ἐξουθενήκασιν, ἀλλʼ ἢ ἐμὲ ἐξουθενήκασι τοῦ μὴ βασιλεύειν ἐπʼ αὐτῶν.
\vs{8}Κατὰ πάντα τὰ ποιήματα, ἃ ἐποίησάν μοι ἀφʼ ἧς ἡμέρας ἀνήγαγον αὐτοὺς ἐξ Αἰγύπτου ἕως τῆς ἡμέρας ταύτης, καὶ ἐγκατέλιπόν με, καὶ ἐδούλευον θεοῖς ἑτέροις, οὕτως αὐτοὶ ποιοῦσι καὶ σοί.
\vs{9}Καὶ νῦν ἄκουε τῆς φωνῆς αὐτῶν· πλὴν ὅτι διαμαρτυρόμενος διαμαρτύρῃ αὐτοῖς, καὶ ἀπαγγελεῖς αὐτοῖς τὸ δικαίωμα τοῦ βασιλέως ὃς βασιλεύσει ἐπʼ αὐτούς.

\vs{10}Καὶ εἶπε Σαμουὴλ πᾶν τὸ ῥῆμα τοῦ κυρίου πρὸς τὸν λαὸν τοὺς αἰτοῦντας παρʼ αὐτοῦ βασιλέα.
\vs{11}Καὶ εἶπε, τοῦτο ἔσται τὸ δικαίωμα τοῦ βασιλέως ὃς βασιλεύσει ἐφʼ ὑμᾶς· τοὺς υἱοὺς ὑμῶν λήψεται, καὶ θήσεται αὐτοὺς ἐν ἅρμασιν αὐτοῦ, καὶ ἐν ἱππεῦσιν αὐτοῦ, καὶ προτρέχοντας τῶν ἁρμάτων αὐτοῦ,
\vs{12}καὶ θέσθαι αὐτοὺς ἑαυτῷ ἑκατονάρχους καὶ χιλιάρχους, καὶ θερίζειν θερισμὸν αὐτοῦ, καὶ τρυγᾷν τρυγητὸν αὐτοῦ, καὶ ποιεῖν σκεύη πολεμικὰ αὐτοῦ, καὶ σκεύη ἁρμάτων αὐτοῦ.
\vs{13}Καὶ τὰς θυγατέρας ὑμῶν λήψεται εἰς μυρεψοὺς, καὶ εἰς μαγειρίσσας, καὶ εἰς πεσσούσας.
\vs{14}Καὶ τοὺς ἀγροὺς ὑμῶν, καὶ τοὺς ἀμπελῶνας ὑμῶν, καὶ τοὺς ἐλαιῶνας ὑμῶν τοὺς ἀγαθοὺς λήψεται, καὶ δώσει τοῖς δούλοις ἑαυτοῦ.
\vs{15}Καὶ τὰ σπέρματα ὑμῶν καὶ τοὺς ἀμπελῶνας ὑμῶν ἀποδεκατώσει, καὶ δώσει τοῖς εὐνούχοις αὐτοῦ, καὶ τοῖς δούλοις αὐτοῦ.
\vs{16}Καὶ τοὺς δούλους ὑμῶν, καὶ τὰς δούλας ὑμῶν, καὶ τὰ βουκόλια ὑμῶν τὰ ἀγαθὰ, καὶ τοὺς ὄνους ὑμῶν λήψεται, καὶ ἀποδεκατώσει εἰς τὰ ἔργα αὐτοῦ·
\vs{17}Καὶ τὰ ποίμνια ὑμῶν ἀποδεκατώσει, καὶ ὑμεῖς ἔσεσθε αὐτῷ δοῦλοι.
\vs{18}Καὶ βοήσεσθε ἐν τῇ ἡμέρᾳ ἐκείνῃ ἐκ προσώπου βασιλέως ὑμῶν οὗ ἐξελέξασθε ἑαυτοῖς, καὶ οὐκ ἐπακούσεται Κύριος ὑμῶν ἐν ταῖς ἡμέραις ἐκείναις, ὅτι ὑμεῖς ἐξελέξασθε ἑαυτοῖς βασιλέα.

\vs{19}Καὶ οὐκ ἐβούλετο ὁ λαὸς ἀκοῦσαι τοῦ Σαμουὴλ, καὶ εἶπαν αὐτῷ, οὐχὶ, ἀλλʼ ἢ βασιλεὺς ἔσται ἐφʼ ἡμᾶς.
\vs{20}Καὶ ἐσόμεθα καὶ ἡμεῖς καθὰ πάντα τὰ ἔθνη· καὶ δικάσει ἡμᾶς βασιλεὺς ἡμῶν, καὶ ἐξελεύσεται ἔμπροσθεν ἡμῶν, καὶ πολεμήσει τὸν πόλεμον ἡμῶν.
\vs{21}Καὶ ἤκουσε Σαμουὴλ πάντας τοὺς λόγους τοῦ λαοῦ, καὶ ἐλάλησεν αὐτοὺς εἰς τὰ ὦτα Κυρίου.
\vs{22}Καὶ εἶπε Κύριος πρὸς Σαμουήλ, ἄκουε τῆς φωνῆς αὐτῶν, καὶ βασίλευσον αὐτοῖς βασιλέα· καὶ εἶπε Σαμουὴλ πρὸς ἄνδρας Ἰσραήλ, ἀποτρεχέτω ἕκαστος εἰς τὴν πόλιν αὐτοῦ.

\ch{9}
Καὶ ἀνὴρ ἐξ υἱῶν Βενιαμίν, καὶ ὄνομα αὐτῷ Κίς, υἱὸς ʼΑβιὴλ, υἱοῦ Ἰαρὲδ, υἱοῦ Βαχὶρ, υἱοῦ ʼΑφὲκ, υἱοῦ ἀνδρὸς Ἰεμιναίου, ἀνὴρ δυνατός.
\vs{2}Καὶ τούτῳ υἱός, καὶ ὄνομα αὐτῷ Σαούλ, εὐμεγέθης, ἀνὴρ ἀγαθός, καὶ οὐκ ἦν ἐν υἱοῖς Ἰσραὴλ ἀγαθὸς ὑπὲρ αὐτόν, ὑπερωμίαν καὶ ἐπάνω ὑψηλὸς ὑπὲρ πᾶσαν τὴν γῆν.

\vs{3}Καὶ ἀπώλοντο αἱ ὄνοι Κὶς πατρὸς Σαούλ· καὶ εἶπε Κὶς πρὸς Σαοὺλ τὸν υἱὸν αὐτοῦ, λάβε μετὰ σεαυτοῦ ἓν τῶν παιδαρίων, καὶ ἀνάστητε καὶ πορεύθητε καὶ ζητήσατε τὰς ὄνους.

\vs{4}Καὶ διῆλθον διʼ ὄρους Ἐφράιμ, καὶ διῆλθον διὰ τῆς γῆς Σελχὰ, καὶ οὐχ εὗρον· καὶ διῆλθον διὰ τῆς γῆς Σεγαλὶμ, καὶ οὐκ ἦν· καὶ διῆλθον διὰ τῆς γῆς Ἰαμὶν, καὶ οὐχ εὗρον.
\vs{5}Αὐτῶν δὲ ἐλθόντων εἰς τὴν Σὶφ, καὶ Σαοὺλ εἶπε τῷ παιδαρίῳ αὐτοῦ τῷ μετʼ αὐτοῦ, δεῦρο καὶ ἀποστρέψωμεν, μὴ ἀνεὶς ὁ πατήρ μου τὰς ὄνους, φροντίζῃ τὰ περὶ ἡμῶν.
\vs{6}Καὶ εἶπεν αὐτῷ τὸ παιδάριον, ἰδοὺ δὴ ἄνθρωπος τοῦ Θεοῦ ἐν τῇ πόλει ταύτῃ, καὶ ὁ ἄνθρωπος ἔνδοξος, πᾶν ὃ ἐὰν λαλήσῃ παραγινόμενον παρέσται· καὶ νῦν πορευθῶμεν, ὅπως ἀπαγγείλῃ ἡμῖν τὴν ὁδὸν ἡμῶν ἐφʼ ἣν ἐπορευθημεν ἐπʼ αὐτήν.
\vs{7}Καὶ εἶπε Σαοὺλ τῷ παιδαρίῳ αὐτοῦ τῷ μετʼ αὐτοῦ, καὶ ἰδοὺ πορευσόμεθα· καὶ τί οἴσομεν τῷ ἀνθρώπῳ τοῦ Θεοῦ; ὅτι οἱ ἄρτοι ἐκλελοίπασιν ἐκ τῶν ἀγγείων ἡμῶν, καὶ πλεῖον οὐκ ἔστι μεθʼ ἡμῶν εἰσενεγκεῖν τῷ ἀνθρώπῳ τοῦ Θεοῦ τὸ ὑπάρχον ἡμῖν.
\vs{8}Καὶ προσέθετο τὸ παιδάριον ἀποκριθῆναι τῷ Σαοὺλ, καὶ εἶπεν, ἰδοὺ εὕρηται ἐν τῇ χειρί μου τέταρτον σίκλου ἀργυρίου, καὶ δώσεις τῷ ἀνθρώπῳ τοῦ Θεοῦ, καὶ ἀπαγγελεῖ ἡμῖν τὴν ὁδὸν ἡμῶν.
\vs{9}Καὶ ἔμπροσθεν ἐν Ἰσραὴλ τάδε ἔλεγεν ἕκαστος ἐν τῷ πορεύεσθαι ἐπερωτᾷν τὸν Θεὸν, δεῦρο καὶ πορευθῶμεν πρὸς τὸν βλέποντα· ὅτι τὸν προφήτην ἐκάλει ὁ λαὸς ἔμπροσθεν, ὁ βλέπων.
\vs{10}Καὶ εἶπε Σαοὺλ πρὸς τὸ παιδάριον αὐτοῦ, ἀγαθὸν τὸ ῥῆμα, δεῦρο καὶ πορευθῶμεν· καὶ ἐπορεύθησαν εἰς τὴν πόλιν οὗ ἦν ἐκεῖ ὁ ἄνθρωπος ὁ τοῦ Θεοῦ.

\vs{11}Αὐτῶν ἀναβαινόντων τὴν ἀνάβασιν τῆς πόλεως, καὶ αὐτοὶ εὑρίσκουσι τὰ κοράσια ἐξεληλυθότα ὑδρεύεσθαι ὕδωρ, καὶ λέγουσιν αὐταῖς, εἰ ἔστιν ἐνταῦθα ὁ βλέπων;
\vs{12}Καὶ ἀπεκρίθη τὰ κοράσια αὐτοῖς, καὶ λέγουσιν αὐτοῖς, ἔστιν· ἰδοὺ κατὰ πρόσωπον ὑμῶν· νῦν διὰ τὴν ἡμέραν ἥκει εἰς τὴν πόλιν, ὅτι θυσία σήμερον τῷ λαῷ ἐν Βαμᾷ.
\vs{13}Ὡς ἂν εἰσέλθητε εἰς τὴν πόλιν, οὕτως εὑρήσετε αὐτὸν ἐν τῇ πόλει, πρὶν ἀναβῆναι αὐτὸν εἰς Βαμᾶ τοῦ φαγεῖν· ὅτι οὐ μὴ φάγῃ ὁ λαὸς ἕως τοῦ εἰσελθεῖν αὐτόν, ὅτι οὗτος εὐλογεῖ τὴν θυσίαν, καὶ μετὰ ταῦτα ἐσθίουσιν οἱ ξένοι· καὶ νῦν ἀνάβητε, ὅτι διὰ τὴν ἡμέραν εὑρήσετε αὐτόν.
\vs{14}Καὶ ἀναβαίνουσι τὴν πόλιν· αὐτῶν εἰσπορευομένων εἰς μέσον τῆς πόλεως, καὶ ἰδοὺ Σαμουὴλ ἐξῆλθεν εἰς τὴν ἀπάντησιν αὐτῶν, τοῦ ἀναβῆναι εἰς Βαμᾶ.

\vs{15}Καὶ Κύριος ἀπεκάλυψε τὸ ὠτίον Σαμουὴλ ἡμέρᾳ μιᾷ ἔμπροσθεν τοῦ ἐλθεῖν πρὸς αὐτὸν Σαοὺλ, λέγων,
\vs{16}ὡς ὁ καιρὸς, αὔριον ἀποστελῶ πρὸς σὲ ἄνδρα ἐκ γῆς Βενιαμίν, καὶ χρίσεις αὐτὸν εἰς ἄρχοντα ἐπὶ τὸν λαόν μου Ἰσραήλ, καὶ σώσει τὸν λαόν μου ἐκ χειρὸς ἀλλοφύλων, ὅτι ἐπέβλεψα ἐπὶ τὴν ταπείνωσιν τοῦ λαοῦ μου, ὅτι ἦλθε βοὴ αὐτῶν πρὸς μέ.
\vs{17}Καὶ Σαμουὴλ εἶδε τὸν Σαούλ, καὶ Κύριος ἀπεκρίθη αὐτῷ, ἰδοὺ ὁ ἄνθρωπος ὃν εἶπά σοι, οὗτος ἄρξει ἐν τῷ λαῷ μου.

\vs{18}Καὶ προσήγαγε Σαοὺλ πρὸς Σαμουὴλ εἰς μέσον τῆς πόλεως, καὶ εἶπεν, ἀπάγγειλον δὴ ποῖος ὁ οἶκος τοῦ βλέποντος.
\vs{19}Καὶ ἀπεκρίθη Σαμουὴλ τῷ Σαοὺλ, καὶ εἶπεν, ἐγώ εἰμι αὐτός· ἀνάβηθι ἔμπροσθέν μου εἰς Βαμὰ, καὶ φάγε μετʼ ἐμοῦ σήμερον, καὶ ἐξαποστελῶ σε πρωῒ, καὶ πάντα τὰ ἐν τῇ καρδίᾳ σου ἀπαγγελῶ σοι.
\vs{20}Καὶ περὶ τῶν ὄνων σου τῶν ἀπολωλυιῶν σήμερον τριταίων, μὴ θῇς τὴν καρδίαν σου αὐταῖς, ὅτι εὕρηνται· καὶ τίνι τὰ ὡραῖα τοῦ Ἰσραήλ; οὐ σοὶ, καὶ τῷ οἴκῳ τοῦ πατρός σου;
\vs{21}Καὶ ἀπεκρίθη Σαοὺλ, καὶ εἶπεν, οὐχὶ ἀνδρὸς υἱὸς Ἰεμιναίου ἐγώ εἰμι τοῦ μικροῦ σκήπτρου φυλῆς Ἰσραήλ; καὶ τῆς φυλῆς τῆς ἐλαχίστης ἐξ ὅλου σκήπτρου Βενιαμίν; καὶ ἱνατί ἐλάλησας πρὸς ἐμὲ κατὰ τὸ ῥῆμα τοῦτο;

\vs{22}Καὶ ἔλαβε Σαμουὴλ τὸν Σαοὺλ καὶ τὸ παιδάριον αὐτοῦ, καὶ εἰσήγαγεν αὐτοὺς εἰς τὸ κατάλυμα, καὶ ἔθετο αὐτοῖς ἐκεῖ τόπον ἐν πρώτοις τῶν κεκλημένων ὡσεὶ ἑβδομήκοντα ἀνδρῶν.
\vs{23}Καὶ εἶπε Σαμουὴλ τῷ μαγείρῳ, δός μοι τὴν μερίδα ἣν ἔδωκά σοι, ἣν εἶπά σοι θεῖναι αὐτὴν παρὰ σοί.
\vs{24}Καὶ ἥψησεν ὁ μάγειρος τὴν κωλέαν, καὶ παρέθηκεν αὐτὴν ἐνώπιον Σαούλ· καὶ εἶπε Σαμουὴλ τῷ Σαούλ, ἰδοὺ ὑπόλειμμα, παράθες αὐτὸ ἐνώπιόν σου καὶ φάγε, ὅτι εἰς μαρτύριον τέθειταί σοι παρὰ τοὺς ἄλλους, ἀπόκνιζε· καὶ ἔφαγεν Σαοὺλ μετὰ Σαμουὴλ ἐν τῇ ἡμέρᾳ ἐκείνῃ.

\vs{25}Καὶ κατέβη ἐκ τῆς Βαμᾶ εἰς τὴν πόλιν· καὶ διέστρωσαν τῷ Σαοὺλ ἐπὶ τῷ δώματι, καὶ ἐκοιμήθη.

\vs{26}Καὶ ἐγένετο ὡς ἀνέβαινεν ὁ ὄρθρος, καὶ ἐκάλεσε Σαμουὴλ τὸν Σαοὺλ ἐπὶ τῷ δώματι, λέγων, ἀνάστα, καὶ ἐξαποστελῶ σε· καὶ ἀνέστη Σαοὺλ, καὶ ἐξῆλθεν αὐτὸς καὶ Σαμουὴλ ἕως ἔξω.
\vs{27}Αὐτῶν καταβαινόντων εἰς μέρος τῆς πόλεως, καὶ Σαμουὴλ εἶπε τῷ Σαοὺλ, εἶπον τῷ νεανίσκῳ, καὶ διελθέτω ἔμπροσθεν ἡμῶν· καὶ σὺ στῆθι ὡς σήμερον, καὶ ἄκουσον ῥῆμα Θεοῦ.

\ch{10}
Καὶ ἔλαβε Σαμουὴλ τὸν φακὸν τοῦ ἐλαίου, καὶ ἐπέχεεν ἐπὶ τὴν κεφαλὴν αὐτοῦ, καὶ ἐφίλησεν αὐτὸν, καὶ εἶπεν αὐτῷ, οὐχὶ κέχρικέ σε Κύριος εἰς ἄρχοντα ἐπὶ τὸν λαὸν αὐτοῦ ἐπὶ Ἰσραήλ; καὶ σὺ ἄρξεις ἐν λαῷ Κυρίου, καὶ σὺ σώσεις αὐτὸν ἐκ χειρὸς ἐχθρῶν αὐτοῦ· καὶ τοῦτό σοι τὸ σημεῖον, ὅτι ἔχρισέ σε Κύριος ἐπὶ κληρονομίαν αὐτοῦ εἰς ἄρχοντα·
\vs{2}Ὡς ἂν ἀπέλθῃς σήμερον ἀπʼ ἐμοῦ, καὶ εὑρήσεις δύο ἄνδρας πρὸς τοῖς τάφοις Ῥαχὴλ ἐν τῷ ὄρει Βενιαμὶν ἁλλομένους μεγάλα· καὶ ἐροῦσί σοι, εὕρηνται αἱ ὄνοι ἃς ἐπορεύθητε ζητεῖν· καὶ ἰδοὺ ὁ πατήρ σου ἀποτετίνακται τὸ ῥῆμα τῶν ὄνων, καὶ ἐδαψιλεύσατο διʼ ὑμᾶς, λέγων, τί ποιήσω ὑπὲρ τοῦ υἱοῦ μου;
\vs{3}Καὶ ἀπελεύσῃ ἐκεῖθε καὶ ἐπέκεινα ἥξεις ἕως τῆς δρυὸς Θαβώρ, καὶ εὑρήσεις ἐκεῖ τρεῖς ἄνδρας ἀναβαίνοντας πρὸς τὸν Θεὸν εἰς Βαιθὴλ, ἕνα αἴροντα τρία αἰγίδια, καὶ ἕνα αἴροντα τρία ἀγγεῖα ἄρτων, καὶ ἕνα αἴροντα ἀσκὸν οἴνου·
\vs{4}Καὶ ἐρωτήσουσί σε τὰ εἰς εἰρήνην, καὶ δώσουσί σοι δύο ἀπαρχὰς ἄρτων, καὶ λὴψῃ ἐκ τῆς χειρὸς αὐτῶν.
\vs{5}Καὶ μετὰ ταῦτα εἰσελεύσῃ εἰς τὸν βουνὸν τοῦ Θεοῦ, οὗ ἐστιν ἐκεῖ τὸ ἀνάστημα τῶν ἀλλοφύλων· ἐκεῖ Νασὶβ ὁ ἀλλόφυλος· καὶ ἔσται ὡς ἂν εἰσέλθητε ἐκεῖ εἰς τὴν πόλιν, καὶ ἀπαντήσεις χορῷ προφητῶν καταβαινόντων ἐκ τῆς Βαμᾶ, καὶ ἔμπροσθεν αὐτῶν νάβλα, καὶ τύμπανον, καὶ αὐλὸς, καὶ κινύρα, καὶ αὐτοὶ προφητεύοντες·
\vs{6}Καὶ ἐφαλεῖται ἐπὶ σὲ πνεῦμα Κυρίου, καὶ προφητεύσεις μετʼ αὐτῶν, καὶ στραφήσῃ εἰς ἄνδρα ἄλλον.
\vs{7}Καὶ ἔσται ὅταν ἥξει τὰ σημεῖα ταῦτα ἐπὶ σὲ, ποίει πάντα ὅσα ἐὰν εὕρῃ ἡ χείρ σου, ὅτι Θεὸς μετὰ σοῦ.
\vs{8}Καὶ καταβήσῃ ἔμπροσθεν τῆς Γαλγὰλ, καὶ ἰδοὺ καταβαίνω πρὸς σὲ ἀνενεγκεῖν ὁλοκαύτωσιν καὶ θυσίας εἰρηνικάς· ἑπτὰ ἡμέρας διαλείψεις ἕως τοῦ ἐλθεῖν με πρὸς σὲ, καὶ γνωρίσω σοι ἃ ποιήσεις.

\vs{9}Καὶ ἐγενήθη ὥστε ἐπιστραφῆναι τῷ ὤμῳ αὐτοῦ ἀπελθεῖν ἀπὸ Σαμουὴλ, μετέστρεψεν αὐτῷ ὁ Θεὸς καρδίαν ἄλλην· καὶ ἦλθε πάντα τὰ σημεῖα ἐν τῇ ἡμέρᾳ ἐκείνῃ.
\vs{10}Καὶ ἔρχεται ἐκεῖθεν εἰς τὸν βουνόν, καὶ ἰδοὺ χορὸς προφητῶν ἐξεναντίας αὐτοῦ· καὶ ἥλατο ἐπʼ αὐτὸν πνεῦμα Θεοῦ, καὶ προεφήτευσεν ἐν μέσῳ αὐτῶν.
\vs{11}Καὶ ἐγενήθησαν πάντες οἱ εἰδότες αὐτὸν ἐχθὲς καὶ τρίτης, καὶ εἶδον, καὶ ἰδοὺ αὐτὸς ἐν μέσῳ τῶν προφητῶν· καὶ εἶπεν ὁ λαὸς ἕκαστος πρὸς τὸν πλησίον αὐτοῦ, τί τοῦτο τὸ γεγονὸς τῷ υἱῷ Κίς; ἢ καὶ Σαοὺλ ἐν προφήταις;
\vs{12}Καὶ ἀπεκρίθη τίς αὐτῶν, καὶ εἶπε, καὶ τίς πατὴρ αὐτοῦ; καὶ διὰ τοῦτο ἐγενήθη εἰς παραβολὴν, ἢ καὶ Σαοὺλ ἐν προφήταις;
\vs{13}Καὶ συνετέλεσε προφητεύων, καὶ ἔρχεται εἰς τὸν βουνόν.

\vs{14}Καὶ εἶπεν ὁ οἰκεῖος αὐτοῦ πρὸς αὐτὸν καὶ πρὸς τὸ παιδάριον αὐτοῦ, ποῦ ἐπορεύθητε; καὶ εἶπαν, ζητεῖν τὰς ὄνους, καὶ εἴδαμεν ὅτι οὐκ εἰσί, καὶ εἰσήλθομεν πρὸς Σαμουήλ.
\vs{15}Καὶ εἶπεν ὁ οἰκεῖος πρὸς Σαοὺλ, ἀπάγγειλον δή μοι, τί εἶπέ σοι Σαμουήλ;
\vs{16}Καὶ εἶπε Σαοὺλ πρὸς τὸν οἰκεῖον αὐτοῦ, ἀπήγγειλεν ἀπαγγέλλων μοι, ὅτι εὕρηνται αἱ ὄνοι· τὸ δὲ ῥῆμα τῆς βασιλείας οὐκ ἀπήγγειλεν αὐτῷ.

\vs{17}Καὶ παρήγγειλε Σαμουὴλ παντὶ τῷ λαῷ πρὸς Κύριον εἰς Μασσηφάθ.
\vs{18}Καὶ εἶπε πρὸς υἱοὺς Ἰσραὴλ, τάδε εἶπε Κύριος ὁ Θεὸς Ἰσραὴλ, λέγων, ἐγὼ ἀνήγαγον τοὺς υἱοὺς Ἰσραὴλ ἐξ Αἰγύπτου, καὶ ἐξειλάμην ὑμᾶς ἐκ χειρὸς Φαραὼ βασιλέως Αἰγύπτου, καὶ ἐκ πασῶν τῶν βασιλειῶν τῶν θλιβουσῶν ὑμᾶς.
\vs{19}Καὶ ὑμεῖς σήμερον ἐξουδενήκατε τὸν Θεόν, ὃς αὐτός ἐστιν ὑμῶν σωτὴρ ἐκ πάντων τῶν κακῶν ὑμῶν καὶ θλίψεων ὑμῶν· καὶ εἴπατε, οὐχί, ἀλλʼ ἢ ὅτι βασιλέα καταστήσεις ἐφʼ ἡμῶν· καὶ νῦν κατάστητε ἐνώπιον Κυρίου κατὰ τὰ σκῆπτρα ὑμῶν καὶ κατὰ τὰς φυλὰς ὑμῶν.

\vs{20}Καὶ προσήγαγε Σαμουὴλ πάντα τὰ σκῆπτρα Ἰσραὴλ, καὶ κατακληροῦται σκῆπτρον Βενιαμίν.
\vs{21}Καὶ προσάγει σκῆπτρον Βενιαμὶν εἰς φυλάς, καὶ κατακληροῦται φυλὴ Ματταρί· καὶ προσάγουσι τὴν φυλὴν Ματταρὶ εἰς ἄνδρας, καὶ κατακληροῦται Σαοὺλ υἱὸς Κίς· καὶ ἐζήτει αὐτόν, καὶ οὐχ εὑρίσκετο.

\vs{22}Καὶ ἐπηρώτησε Σαμουὴλ ἔτι ἐν Κυρίῳ, εἰ ἔρχεται ὁ ἀνὴρ ἐνταῦθα; καὶ εἶπε Κύριος, ἰδοὺ αὐτὸς κέκρυπται ἐν τοῖς σκεύεσι.
\vs{23}Καὶ ἔδραμε καὶ λαμβάνει αὐτὸν ἐκεῖθεν, καὶ κατέστησεν ἐν μέσῳ τοῦ λαοῦ· καὶ ὑψώθη ὑπὲρ πάντα τὸν λαὸν ὑπερωμίαν καὶ ἐπάνω.

\vs{24}Καὶ εἶπε Σαμουὴλ πρὸς πάντα τὸν λαὸν, εἰ ἑωράκατε ὃν ἐκλέλεκται ἑαυτῷ Κύριος, ὅτι οὐκ ἔστιν ὅμοιος αὐτῷ ἐν πᾶσιν ὑμῖν; καὶ ἔγνωσαν πᾶς ὁ λαὸς, καὶ εἶπαν, ζήτω ὁ βασιλεύς.
\vs{25}Καὶ εἶπε Σαμουὴλ πρὸς τὸν λαὸν τὸ δικαίωμα τοῦ βασιλέως, καὶ ἔγραψεν ἐν βιβλίῳ, καὶ ἔθηκεν ἐνώπιον Κυρίου· καὶ ἐξαπέστειλε Σαμουὴλ πάντα τὸν λαόν, καὶ ἀπῆλθεν ἕκαστος εἰς τὸν τόπον αὐτοῦ.

\vs{26}Καὶ Σαοὺλ ἀπῆλθεν εἰς τὸν οἶκον αὐτοῦ εἰς Γαβαά· καὶ ἐπορεύθησαν υἱοὶ δυνάμεων, ὧν ἥψατο Κύριος καρδίας αὐτῶν μετὰ Σαούλ.
\vs{27}Καὶ υἱοὶ λοιμοὶ εἶπαν, τίς σώσει ἡμᾶς οὗτος; καὶ ἠτίμασαν αὐτόν, καὶ οὐκ ἤνεγκαν αὐτῷ δῶρα.

\ch{11}
Καὶ ἐγενήθη ὡς μετὰ μῆνα, καὶ ἀνέβη Νάας ὁ Ἀμμανίτης, καὶ παρεμβάλλει ἐπὶ Ἰαβὶς Γαλαάδ· καὶ εἶπαν πάντες οἱ ἄνδρες Ἰαβὶς πρὸς Νάας τὸν Ἀμμανίτην, διάθου ἡμῖν διαθήκην, καὶ δουλεύσομέν σοι.
\vs{2}Καὶ εἶπε πρὸς αὐτοὺς Νάας ὁ ʼΑμμανίτης, ἐν ταύτῃ διαθήσομαι διαθήκην ὑμῖν, ἐν τῷ ἐξορύξαι ὑμῶν πάντα ὀφθαλμὸν δεξιόν, καὶ θήσομαι ὄνειδος ἐπὶ Ἰσραήλ.
\vs{3}Καὶ λέγουσιν αὐτῷ οἱ ἄνδρες Ἰαβίς, ἄνες ἡμῖν ἑπτὰ ἡμέρας καὶ ἀποστελοῦμεν ἀγγέλους εἰς πᾶν ὅριον Ἰσραήλ· ἐὰν μὴ ᾖ ὁ σώζων ἡμᾶς, ἐξελευσόμεθα πρὸς ὑμᾶς.

\vs{4}Καὶ ἔρχονται οἱ ἄγγελοι εἰς Γαβαὰ πρὸς Σαούλ, καὶ λαλοῦσι τοὺς λόγους εἰς τὰ ὦτα τοῦ λαοῦ· καὶ ᾖραν πᾶς ὁ λαὸς τὴν φωνὴν αὐτῶν, καὶ ἔκλαυσαν.
\vs{5}Καὶ ἰδοὺ Σαοὺλ ἤρχετο μετὰ τοπρωῒ ἐξ ἀγροῦ· καὶ εἶπε Σαοὺλ, Τί ὅτι κλαίει ὁ λαός; καὶ διηγοῦνται αὐτῷ τὰ ῥήματα τῶν ἀνδρῶν Ἰαβίς.
\vs{6}Καὶ ἐφήλατο πνεῦμα Κυρίου ἐπὶ Σαοὺλ ὡς ἤκουσε τὰ ῥήματα ταῦτα, καὶ ἐθυμώθη ἐπʼ αὐτοὺς ὀργῇ αὐτοῦ σφόδρα.
\vs{7}Καὶ ἔλαβε δύο βόας, καὶ ἐμέλισεν αὐτὰς, καὶ ἀπέστειλεν εἰς πᾶν ὅριον Ἰσραὴλ ἐν χειρὶ ἀγγέλων, λέγων, ὃς οὐκ ἔστιν ἐκπορευόμενος ὀπίσω Σαοὺλ καὶ ὀπίσω Σαμουὴλ, κατὰ τάδε ποιήσουσι τοῖς βουσὶν αὐτοῦ· καὶ ἐπῆλθεν ἔκστασις Κυρίου ἐπὶ τὸν λαὸν Ἰσραὴλ, καὶ ἐβόησαν ὡς ἀνὴρ εἷς.
\vs{8}Καὶ ἐπισκέπτεται αὐτοὺς ἐν βεζὲκ ἐν Βαμᾷ πάντα ἄνδρα Ἰσραὴλ ἑξακοσίας χιλιάδας, καὶ ἄνδρας Ἰούδα ἑβδομήκοντα χιλιάδας.

\vs{9}Καὶ εἶπε τοῖς ἀγγέλοις τοῖς ἐρχομένοις, τάδε ἐρεῖτε τοῖς ἀνδράσιν Ἰαβὶς, αὔριον ὑμῖν ἡ σωτηρία διαθερμάναντος τοῦ ἡλίου· καὶ ἦλθον οἱ ἄγγελοι εἰς τὴν πόλιν, καὶ ἀπαγγέλλουσι τοῖς ἀνδράσιν Ἰαβίς, καὶ εὐφράνθησαν.
\vs{10}Καὶ εἶπον οἱ ἄνδρες Ἰαβὶς πρὸς Νάας τὸν ʼΑμμανίτην, αὔριον ἐξελευσόμεθα πρὸς ὑμᾶς, καὶ ποιήσετε ἡμῖν τὸ ἀγαθὸν ἐνώπιον ὑμῶν.

\vs{11}Καὶ ἐγενήθη μετὰ τὴν αὔριον, καὶ ἔθετο Σαοὺλ τὸν λαὸν εἰς τρεῖς ἀρχὰς, καὶ εἰσπορεύονται μέσον τῆς παρεμβολῆς ἐν φυλακῇ τῇ ἑωθινῇ, καὶ ἔτυπτον τοὺς υἱοὺς ʼΑμμὼν ἕως διεθερμάνθη ἡ ἡμέρα· καὶ ἐγενήθη καὶ ὑπολελειμμένοι διεσπάρησαν, καὶ οὐκ ὑπελείφθησαν ἐν αὐτοῖς δύο κατὰ τὸ αὐτό.

\vs{12}Καὶ εἶπεν ὁ λαὸς πρὸς Σαμουὴλ τίς ὁ εἴπας ὅτι Σαοὺλ οὐ βασιλεύσει ἡμῶν; παράδος τοὺς ἄνδρας, καὶ θανατώσομεν αὐτούς.
\vs{13}Καὶ εἶπε Σαοὺλ, οὐκ ἀποθανεῖται οὐδεὶς ἐν τῇ ἡμέρᾳ ταύτῃ, ὅτι σήμερον ἐποίησε Κύριος σωτηρίαν ἐν Ἰσραήλ.

\vs{14}Καὶ εἶπε Σαμουὴλ πρὸς τὸν λαὸν, λέγων, πορευθῶμεν εἰς Γάλγαλα, καὶ ἐγκαινίσωμεν ἐκεῖ τὴν βασιλείαν.
\vs{15}Καὶ ἐπορεύθη πᾶς ὁ λαὸς εἰς Γάλγαλα, καὶ ἔχρισε Σαμουὴλ ἐκεῖ τὸν Σαοὺλ εἰς βασιλέα ἐνώπιον Κυρίου ἐν Γαλγάλοις· καὶ ἔθυσεν ἐκεῖ θυσίας καὶ εἰρηνικὰς ἐνώπιον Κυρίου· καὶ εὐφράνθη Σαμουὴλ καὶ πᾶς Ἰσραὴλ ὥστε λίαν.

\ch{12}
Καὶ εἶπε Σαμουὴλ πρὸς πάντα Ἰσραὴλ, ἰδοὺ ἤκουσα φωνῆς ὑμῶν εἰς πάντα ὅσα εἴπατέ μοι, καὶ ἐβασίλευσα ἐφʼ ὑμᾶς βασιλέα,
\vs{2}καὶ νῦν ἰδοὺ ὁ βασιλεὺς διαπορεύεται ἐνώπιον ὑμῶν· κἀγὼ γεγήρακα καὶ καθήσομαι, καὶ οἱ υἱοί μου ἰδοὺ ἐν ὑμῖν· κᾀγὼ ἰδοὺ διελήλυθα ἐνώπιον ὑμῶν ἐκ νεότητος καὶ ἕως τῆς ἡμέρας ταύτης.
\vs{3}Ἰδοὺ ἐγώ, ἀποκρίθητε κατʼ ἐμοῦ ἐνώπιον Κυρίου καὶ ἐνώπιον χριστοῦ αὐτοῦ· μόσχον τίνος εἴληφα, ἢ ὄνον τίνος εἴληφα, ἢ τίνα κατεδυνάστευσα ὑμῶν, ἢ τίνα ἐξεπίεσα, ἢ ἐκ χειρὸς τίνος εἴληφα ἐξίλασμα καὶ ὑπόδημα; ἀποκρίθητε κατʼ ἐμοῦ, καὶ ἀποδώσω ὑμῖν.
\vs{4}Καὶ εἶπαν πρὸς Σαμουὴλ οὐκ ἠδίκησας ἡμᾶς, καὶ οὐ κατεδυνάστευσας ἡμᾶς, καὶ οὐκ ἔθλασας ἡμᾶς, καὶ οὐκ εἴληφας ἐκ χειρὸς οὐδενὸς οὐδέν.

\vs{5}Καὶ εἶπε Σαμουὴλ πρὸς τὸν λαὸν, μάρτυς Κύριος ἐν ὑμῖν, καὶ μάρτυς χριστὸς αὐτοῦ σήμερον ἐν ταύτῃ τῇ ἡμέρᾳ, ὅτι οὐχ εὑρήκατε ἐν χειρί μου οὐδέν· καὶ εἶπαν, μάρτυς.

\vs{6}Καὶ εἶπε Σαμουὴλ πρὸς τὸν λαὸν, λέγων, μάρτυς Κύριος ὁ ποιήσας τὸν Μωυσῆν καὶ τὸν ʼΑαρὼν, ὁ ἀναγαγὼν τοὺς πατέρας ἡμῶν ἐξ Αἰγύπτου.
\vs{7}Καὶ νῦν κατάστητε, καὶ δικάσω ὑμᾶς ἐνώπιον Κυρίου, καὶ ἀπαγγελῶ ὑμῖν τὴν πᾶσαν δικαιοσύνην Κυρίου, ἃ ἐποίησεν ἐν ὑμῖν καὶ ἐν τοῖς πατράσιν ὑμῶν·
\vs{8}Ὡς εἰσῆλθεν Ἰακὼβ καὶ οἱ υἱοὶ αὐτοῦ εἰς Αἴγυπτον, καὶ ἐταπείνωσεν αὐτοὺς Αἴγυπτος· καὶ ἐβόησαν οἱ πατέρες ἡμῶν πρὸς Κύριον, καὶ ἀπέστειλε Κύριος τὸν Μωυσῆν καὶ τὸν ʼΑαρών, καὶ ἐξήγαγον τοὺς πατέρας ἡμῶν ἐξ Αἰγύπτου, καὶ κατῴκισεν αὐτοὺς ἐν τῷ τόπῳ τούτῳ.
\vs{9}Καὶ ἐπελάθοντο Κυρίου τοῦ Θεοῦ αὐτῶν, καὶ ἀπέδοτο αὐτοὺς εἰς χεῖρας Σισαρά ἀρχιστρατήγῳ Ἰαβὶς βασιλέως ʼΑσὼρ, καὶ εἰς χεῖρας ἀλλοφύλων, καὶ εἰς χεῖρας βασιλέως Μωὰβ, καὶ ἐπολέμησεν ἐν αὐτοῖς.
\vs{10}Καὶ ἐβόησαν πρὸς Κύριον, καὶ ἔλεγον, ἡμάρτομεν, ὅτι ἐγκατελίπομεν τὸν Κύριον, καὶ ἐδουλεύσαμεν τοῖς Βααλὶμ καὶ τοῖς ἄλσεσι· καὶ νῦν ἐξελοῦ ἡμᾶς ἐκ χειρὸς ἐχθρῶν ἡμῶν, καὶ δουλεύσομέν σοι.
\vs{11}Καὶ ἀπέστειλε τὸν Ἰεροβάαλ, καὶ τὸν Βαρὰκ, καὶ τὸν Ἰεφθάε, καὶ τὸν Σαμουὴλ, καὶ ἐξείλατο ἡμᾶς ἐκ χειρὸς ἐχθρῶν ἡμῶν τῶν κυκλόθεν, καὶ κατῳκεῖτε πεποιθότες.
\vs{12}Καὶ ἴδετε ὅτι Νάας βασιλεὺς υἱῶν ʼΑμμὼν ἦλθεν ἐφʼ ὑμᾶς, καὶ εἴπατε, οὐχί, ἀλλʼ ἢ ὅτι βασιλεὺς βασιλεύσει ἐφʼ ἡμῶν· καὶ Κύριος ὁ Θεὸς ἡμῶν βασιλεὺς ἡμῶν.

\vs{13}Καὶ νῦν ἰδοὺ ὁ βασιλεὺς ὃν ἐξελέξασθε· καὶ ἰδοὺ δέδωκε Κύριος ἐφʼ ὑμᾶς βασιλέα.
\vs{14}Ἐὰν φοβηθῆτε τὸν Κύριον, καὶ δουλεύσητε αὐτῷ, καὶ ἀκούσητε τῆς φωνῆς αὐτοῦ, καὶ μὴ ἐρίσητε τῷ στόματι Κυρίου, καὶ ἦτε καὶ ὑμεῖς καὶ ὁ βασιλεὺς ὁ βασιλεύων ἐφʼ ὑμῶν ὀπίσω Κυρίου πορευόμενοι·
\vs{15}Ἐὰν δὲ μὴ ἀκούσητε τῆς φωνῆς Κυρίου, καὶ ἐρίσητε τῷ στόματι Κυρίου, καὶ ἔσται χεὶρ Κυρίου ἐφʼ ὑμᾶς καὶ ἐπὶ τὸν βασιλέα ὑμῶν.

\vs{16}Καὶ νῦν κατάστητε, καὶ ἴδετε τὸ ῥῆμα τὸ μέγα τοῦτο ὃ ὁ κύριος ποιήσει ἐν ὀφθαλμοῖς ὑμῶν.
\vs{17}Οὐχὶ θερισμὸς πυρῶν σήμερον; ἐπικαλέσομαι Κύριον, καὶ δώσει φωνὰς καὶ ὑετόν· καὶ γνῶτε καὶ ἴδετε, ὅτι ἡ κακία ὑμῶν μεγάλη, ἣν ἐποίησατε ἐνώπιον Κυρίου, αἰτήσαντες ἑαυτοῖς βασιλέα.

\vs{18}Καὶ ἐπεκαλέσατο Σαμουὴλ τὸν Κύριον· καὶ ἔδωκε Κύριος φωνὰς καὶ ὑετὸν ἐν τῇ ἡμέρᾳ ἐκείνῃ· καὶ ἐφοβήθησαν πᾶς ὁ λαὸς τὸν Κύριον σφόδρα καὶ τὸν Σαμουήλ.
\vs{19}Καὶ εἶπαν πᾶς ὁ λαὸς πρὸς Σαμουὴλ, πρόσευξαι ὑπὲρ τῶν δούλων σου πρὸς Κύριον Θεόν σου, καὶ οὐ μὴ ἀποθάνωμεν, ὅτι προστεθείκαμεν πρὸς πάσας τὰς ἁμαρτίας ἡμῶν κακίαν, αἰτήσαντες ἑαυτοῖς βασιλέα.

\vs{20}Καὶ εἶπε Σαμουὴλ πρὸς τὸν λαὸν, μὴ φοβεῖσθε· ὑμεῖς πεποιήκατε τὴν πᾶσαν κακίαν ταύτην· πλὴν μὴ ἐκκλίνητε ἀπὸ ὄπισθεν Κυρίου, καὶ δουλεύσατε τῷ κυρίῳ ἐν ὅλῃ καρδίᾳ ὑμῶν·
\vs{21}Καὶ μὴ παραβῆτε ὀπίσω τῶν μηθὲν ὄντων, οἳ οὐ περανοῦσιν οὐθὲν, καὶ οἳ οὐκ ἐξελοῦνται, ὅτι οὐθέν εἰσιν.
\vs{22}Ὅτι οὐκ ἀπώσεται Κύριος τὸν λαὸν αὐτοῦ διὰ τὸ ὄνομα αὐτοῦ τὸ μέγα, ὅτι ἐπιεικῶς Κύριος προσελάβετο ὑμᾶς ἑαυτῷ εἰς λαόν.
\vs{23}Καὶ ἐμοὶ μηδαμῶς τοῦ ἁμαρτεῖν τῷ Κυρίῳ ἀνιέναι τοῦ προσεύχεσθαι περὶ ὑμῶν· καὶ δουλεύσω τῷ Κυρίῳ, καὶ δείξω ὑμῖν τὴν ὁδὸν τὴν ἀγαθὴν καὶ τὴν εὐθεῖαν·
\vs{24}Πλὴν φοβεῖσθε τὸν κύριον, καὶ δουλεύσατε αὐτῷ ἐν ἀληθείᾳ καὶ ἐν ὅλῃ καρδίᾳ ὑμῶν, ὅτι ἴδετε ἃ ἐμεγάλυνε μεθʼ ὑμῶν.
\vs{25}Καὶ ἐὰν κακίᾳ κακοποιήσητε, καὶ ὑμεῖς καὶ ὁ βασιλεὺς ὑμῶν προστεθήσεσθε.

\ch{13}
\vs{2}Καὶ ἐκλέγεται ἐαυτῷ Σαοὺλ τρεῖς χιλιάδας ἀνδρῶν ἐκ τῶν ἀνδρῶν Ἰσραήλ· καὶ ἦσαν μετὰ Σαοὺλ δισχίλιοι οἱ ἐν Μαχμὰς, καὶ ἐν τῷ ὄρει Βαιθήλ, καὶ χίλιοι ἦσαν μετὰ Ἰωνάθαν ἐν Γαβαὰ τοῦ Βενιαμίν· καὶ τὸ κατάλοιπον τοῦ λαοῦ ἐξαπέστειλεν ἕκαστον εἰς τὸ σκήνωμα αὐτοῦ.

\vs{3}Καὶ ἐπάταξεν Ἰωνάθαν τὸν Νασὶβ τὸν ἀλλόφυλον τὸν ἐν τῷ βουνῷ· καὶ ἀκούουσιν οἱ ἀλλόφυλοι· καὶ Σαοὺλ σάλπιγγι σαλπίζει εἰς πᾶσαν τὴν γῆν, λέγων, ἠθετήκασιν οἱ δοῦλοι.
\vs{4}Καὶ πᾶς Ἰσραὴλ ἤκουσε λεγόντων, πέπαικε Σαοὺλ τὸν Νασὶβ τὸν ἀλλόφυλον, καὶ ᾐσχύνθησαν Ἰσραὴλ ἐν τοῖς ἀλλοφύλοις· καὶ ἀνέβησαν οἱ υἱοὶ Ἰσραὴλ ὀπίσω Σαοὺλ ἐν Γαλγάλοις.
\vs{5}Καὶ οἱ ἀλλόφυλοι συνάγονται εἰς πόλεμον ἐπὶ Ἰσραήλ· καὶ ἀναβαίνουσιν ἐπὶ Ἰσραὴλ τριάκοντα χιλιάδες ἁρμάτων, καὶ ἓξ χιλιάδες ἱππέων, καὶ λαὸς ὡς ἡ ἄμμος ἡ παρὰ τὴν θάλασσαν τῷ πλήθει· καὶ ἀναβαίνουσι καὶ παρεμβάλλουσιν ἐν Μαχμὰς ἐξ ἐναντίας Βαιθωρὼν κατὰ Νότου.

\vs{6}Καὶ ἀνὴρ Ἰσραὴλ εἶδεν ὅτι στενῶς αὐτῷ μὴ προσάγειν αὐτόν, καὶ ἐκρύβη ὁ λαὸς ἐν τοῖς σπηλαίοις, καὶ ἐν ταῖς μάνδραις, καὶ ἐν ταῖς πέτραις, καὶ ἐν τοῖς βόθροις, καὶ ἐν τοῖς λάκκοις.
\vs{7}Καὶ οἱ διαβαίνοντες διέβησαν τὸν Ἰορδάνην εἰς γῆν Γὰδ καὶ Γαλαάδ· καὶ Σαοὺλ ἔτι ἦν ἐν Γαλγάλοις, καὶ πᾶς ὁ λαὸς ἐξέστη ὀπίσω αὐτοῦ.
\vs{8}Καὶ διέλιπεν ἑπτὰ ἡμέρας τῷ μαρτυρίῳ, ὡς εἶπε Σαμουήλ, καὶ οὐ παρεγένετο Σαμουὴλ εἰς Γάλγαλα, καὶ διεσπάρη ὁ λαὸς αὐτοῦ ἀπʼ αὐτοῦ.
\vs{9}Καὶ εἶπε Σαούλ, προσαγάγετε ὅπως ποιήσω ὁλοκαύτωσιν καὶ εἰρηνικάς· καὶ ἀνήνεγκε τὴν ὁλοκαύτωσιν.

\vs{10}Καὶ ἐγένετο ὡς συνετέλεσεν ἀναφέρων τὴν ὁλοκαύτωσιν, καὶ Σαμουὴλ παραγίνεται· καὶ ἐξῆλθε Σαοὺλ εἰς ἀπάντησιν αὐτῷ εὐλογῆσαι αὐτόν.
\vs{11}Καὶ εἶπε Σαμουήλ, τί πεποίηκας; καὶ εἶπε Σαούλ, ὅτι εἶδον ὡς διεσπάρη ὁ λαὸς ἀπʼ ἐμοῦ, καὶ σὺ οὐ παρεγένου ὡς διετάξω ἐν τῷ μαρτυρίῳ τῶν ἡμερῶν, καὶ οἱ ἀλλόφυλοι συνήχθησαν εἰς Μαχμάς,
\vs{12}καὶ εἶπα, νῦν καταβήσονται οἱ ἀλλόφυλοι πρὸς μὲ εἰς Γάλγαλα, καὶ τοῦ προσώπου τοῦ κυρίου οὐκ ἐδεήθην· καὶ ἐνεκρατευσάμην, καὶ ἀνήνεγκα τὴν ὁλοκαύτωσιν.
\vs{13}Καὶ εἶπε Σαμουὴλ πρὸς Σαοὺλ, μεματαίωταί σοι, ὅτι οὐκ ἐφύλαξας τὴν ἐντολήν μου, ἣν ἐνετείλατό σοι Κύριος, ὡς νῦν ἡτοίμασε Κύριος τὴν βασιλείαν σου ἐπὶ Ἰσραὴλ ἕως αἰῶνος.
\vs{14}Καὶ νῦν ἡ βασιλεία σου οὐ στήσεταί σοι, καὶ ζητήσει Κύριος ἑαυτῷ ἄνθρωπον κατὰ τὴν καρδίαν αὐτοῦ· καὶ ἐντελεῖται Κύριος αὐτῷ εἰς ἄρχοντα ἐπὶ τὸν λαὸν αὐτοῦ, ὅτι οὐκ ἐφύλαξας ὅσα ἐνετείλατό σοι Κύριος.

\vs{15}Καὶ ἀνέστη Σαμουὴλ, καὶ ἀπῆλθεν ἐκ Γαλγάλων· καὶ τὸ κατάλειμμα τοῦ λαοῦ ἀνέβη ὀπίσω Σαοὺλ εἰς ἀπάντησιν ὀπίσω τοῦ λαοῦ τοῦ πολεμιστοῦ· αὐτῶν παραγενομένων ἐκ Γαλγάλων εἰς Γαβαὰ Βενιαμίν. Καὶ ἐπεσκέψατο Σαοὺλ τὸν λαὸν τὸν εὑρεθέντα μετʼ αὐτοῦ ὡς ἑξακοσίους ἄνδρας.
\vs{16}Καὶ Σαοὺλ καὶ Ἰωνάθαν υἱὸς αὐτοῦ καὶ ὁ λαὸς οἱ εὑρεθέντες μετʼ αὐτῶν ἐκάθισαν ἐν Γαβαὰ Βενιαμὶν, καὶ ἔκλαιον· Καὶ οἱ ἀλλόφυλοι παρεμβεβλήκεισαν ἐν Μαχμάς.
\vs{17}Καὶ ἐξῆλθε διαφθείρων ἐξ ἀγροῦ ἀλλοφύλων τρισὶν ἀρχαῖς· ἡ ἀρχὴ ἡ μία ἐπιβλέπουσα ὁδὸν Γοφερὰ ἐπὶ γὴν Σωγὰλ,
\vs{18}καὶ ἡ ἀρχὴ ἡ μία ἐπιβλέπουσα ὁδὸν Βαιθωρὼν, καὶ ἡ ἀρχὴ ἡ μία ἐπιβλέπουσα ὁδὸν Γαβαὲ τὴν εἰσκύπτουσαν ἐπὶ Γαὶ τὴν Σαβίμ.

\vs{19}Καὶ τέκτων σιδήρου οὐχ εὑρίσκετο ἐν πάσῃ γῇ Ἰσραὴλ, ὅτι εἶπον οἱ ἀλλόφυλοι, μὴ ποιήσωσιν οἱ Ἑβραῖοι ῥομφαίαν καὶ δόρυ.
\vs{20}Καὶ κατέβαινον πᾶς Ἰσραὴλ εἰς γῆν ἀλλοφύλων χαλκεύειν ἕκαστος τὸ θέριστρον αὐτοῦ καὶ τὸ σκεῦος αὐτοῦ, καὶ ἕκαστος τὴν ἀξίνην αὐτοῦ καὶ τὸ δρέπανον αὐτοῦ.
\vs{21}Καὶ ἦν ὁ τρυγητὸς ἕτοιμος τοῦ θερίζειν· τὰ δὲ σκεύη ἦν τρεῖς σίκλοι εἰς τὸν ὀδόντα, καὶ τῇ ἀξίνῃ καὶ τῷ δρεπάνῳ ὑπόστασις ἦν ἡ αὐτή.
\vs{22}Καὶ ἐγενήθη ἐν ταῖς ἡμέραις τοῦ πολέμου Μαχμὰς, καὶ οὐχ εὑρέθη ῥομφαία καὶ δόρυ ἐν χειρὶ παντὸς τοῦ λαοῦ τοῦ μετὰ Σαοὺλ καὶ μετὰ Ἰωνάθαν· καὶ εὑρέθη τῷ Σαοὺλ καὶ τῷ Ἰωνὰθαν υἱῷ αὐτοῦ.

\vs{23}Καὶ ἐξῆλθεν ἐξ ὑποστάσεως τῶν ἀλλοφύλων τὴν ἐν τῷ πέραν Μαχμάς.

\ch{14}
Καὶ γίνεται ἡ ἡμέρα, καὶ εἶπεν Ἰωνάθαν υἱὸς Σαοὺλ τῷ παιδαρίῳ τῷ αἴροντι τὰ σκεύη αὐτοῦ, δεῦρο, καὶ διαβῶμεν εἰς Μεσσὰβ τῶν ἀλλοφύλων τὴν ἐν τῷ πέραν ἐκείνῳ· καὶ τῷ πατρὶ αὐτοῦ οὐκ ἀπήγγειλε.
\vs{2}Καὶ Σαοὺλ ἐκάθητο ἐπʼ ἄκρου τοῦ βουνοῦ ὑπὸ τὴν ῥοὰν τὴν ἐν Μαγδὼν, καὶ ἦσαν μετʼ αὐτοῦ ὡς ἑξακόσιοι ἄνδρες.
\vs{3}Καὶ ʼΑχιὰ υἱὸς ʼΑχιτὼβ ἀδελφοῦ Ἰωχαβὴλ υἱοῦ Φινεὲς υἱοῦ Ἡλὶ ἱερεὺς τοῦ Θεοῦ ἐν Σηλὼμ αἴρων ἐφούδ· καὶ ὁ λαὸς οὐκ ᾔδει ὅτι πεπόρευται Ἰωνάθαν.
\vs{4}Καὶ ἀναμέσον τῆς διαβάσεως οὗ ἐζήτει Ἰωνάθαν διαβῆναι εἰς τὴν ὑπόστασιν τῶν ἀλλοφύλων, καὶ ὀδοὺς πέτρας ἐκ τούτου, καὶ ὀδοὺς πέτρας ἐκ τούτου· ὄνομα τῷ ἑνὶ Βασὲς, καὶ ὄνομα τῷ ἄλλῳ Σεννά.
\vs{5}Ἡ ὁδὸς ἡ μία ἀπὸ Βοῤῥᾶ ἐρχομένῳ Μαχμάς, καὶ ἡ ὁδὸς ἡ ἄλλη ἀπὸ Νότου ἐρχομένῳ Γαβαέ.

\vs{6}Καὶ εἶπεν Ἰωνάθαν πρὸς τὸ παιδάριον τὸ αἶρον τὰ σκεύη αὐτοῦ, δεῦρο, διαβῶμεν εἰς Μεσσὰβ τῶν ἀπεριτμήτων τούτων, εἴ τι ποιήσαι Κύριος ἡμῖν, ὅτι οὐκ ἔστι τῷ Κυρίῳ συνεχόμενον σώζειν ἐν πολλοῖς ἢ ἐν ὀλίγοις.
\vs{7}Καὶ εἶπεν αὐτῷ ὁ αἴρων τὰ σκεύη αὐτοῦ, ποίει πᾶν ὃ ἐὰν ἡ καρδία σου ἐκκλίνῃ· ἰδοὺ ἐγὼ μετὰ σοῦ, ὡς ἡ καρδία σον καρδία μοῦ.
\vs{8}Καὶ εἶπεν Ἰωνάθαν, ἰδοὺ ἡμεὶς διαβαίνομεν πρὸς τοὺς ἄνδρας, καὶ κατακυλισθησόμεθα πρὸς αὐτούς·
\vs{9}Ἐὰν τάδε εἴπωσι πρὸς ἡμᾶς, ἀπόστητε ἐκεῖ ἕως ἂν ἀπαγγείλωμεν ὑμῖν· καἰ στησόμεθα ἐφʼ ἑαυτοῖς, καὶ οὐ μὴ ἀναβῶμεν ἐπʼ αὐτούς.
\vs{10}Ἐὰν τάδε εἴπωσι πρὸς ἡμᾶς, ἀνάβητε πρὸς ἡμᾶς· καὶ ἀναβησόμεθα, ὅτι παραδέδωκεν αὐτοὺς Κύριος εἰς χεῖρας ἡμῶν· τοῦτο ἡμῖν τὸ σημεῖον.

\vs{11}Καὶ εἰσῆλθον ἀμφότεροι εἰς Μεσσὰβ τῶν ἀλλοφύλων· καὶ λέγουσιν οἱ ἀλλόφυλοι, ἰδοὺ Ἐβραῖοι ἐκπορεύονται ἐκ τῶν τρωγλῶν αὐτῶν, οὗ ἐκρύβησαν ἐκεῖ.
\vs{12}Καὶ ἀπεκρίθησαν οἱ ἄνδρες Μεσσὰβ πρὸς Ἰωνάθαν καὶ πρὸς τὸν αἴροντα τὰ σκεύη αὐτοῦ, καὶ λέγουσιν, ἀνάβητε πρὸς ἡμᾶς, καὶ γνωριοῦμεν ὑμῖν ῥῆμα· καὶ εἶπεν Ἰωνάθαν πρὸς τὸν αἴροντα τὰ σκεύη αὐτοῦ, ἀνάβηθι ὀπίσω μου, ὅτι παρέδωκεν αὐτοὺς Κύριος εἰς χεῖρας Ἰσραήλ.
\vs{13}Καὶ ἀνέβη Ἰωνάθὰν ἐπὶ τὰς χεῖρας αὐτοῦ καὶ ἐπὶ τοὺς πόδας αὐτοῦ, καὶ ὁ αἴρων τὰ σκεύη αὐτοῦ μετʼ αὐτοῦ· καὶ ἐπέβλεψαν κατὰ πρόσωπον Ἰωνάθαν, καὶ ἐπάταξεν αὐτούς, καὶ ὁ αἴρων τὰ σκεύη αὐτοῦ ἐπεδίδου ὀπίσω αὐτοῦ.
\vs{14}Καὶ ἐγενήθη ἡ πληγὴ ἡ πρώτη, ἣν ἐπάταξεν Ἰωνάθαν καὶ ὁ αἴρων τὰ σκεύη αὐτοῦ, ὡς εἴκοσι ἄνδρες ἐν βολίσι καὶ ἐν πετροβόλοις καὶ ἐν κόχλαξι τοῦ πεδίου.

\vs{15}Καὶ ἐγενήθη ἔκστασις ἐν τῇ παρεμβολῇ, καὶ ἐν ἀγρῷ· καὶ πᾶς ὁ λαός οἱ ἐν Μεσσὰβ, καὶ οἱ διαφθείροντες ἐξέστησαν, καὶ αὐτοὶ οὐκ ἤθελον ποιεῖν· καὶ ἐθάμβησεν ἡ γῆ, καὶ ἐγενήθη ἔκστασις παρὰ Κυρίου.

\vs{16}Καὶ εἶδον οἱ σκοποὶ τοῦ Σαοὺλ ἐν Γαβαὰ Βενιαμὶν, καὶ ἰδοὺ ἡ παρεμβολὴ τεταραγμένη ἔνθεν καὶ ἔνθεν.
\vs{17}Καὶ εἶπε Σαοὺλ τῷ λαῷ τῷ μετʼ αὐτοῦ, ἐπισκέψασθε δὴ, καὶ ἴδετε τίς πεπόρευται ἐξ ὑμῶν· καὶ ἐπεσκέψαντο, καὶ ἰδοὺ οὐχ εὑρίσκετο Ἰωνάθαν καὶ ὁ αἴρων τὰ σκεύη αὐτοῦ.
\vs{18}Καὶ εἶπε Σαοὺλ τῷ ʼΑχιᾷ, προσάγαγε τὸ ἐφούδ· ὅτι αὐτὸς ᾗρε τὸ ἐφοὺδ ἐν τῇ ἡμέρᾳ ἐκείνῃ ἐνώπιον Ἰσραήλ.
\vs{19}Καὶ ἐγενήθη ὡς λαλεῖ Σαοὺλ πρὸς τὸν ἱερέα, καὶ ὁ ἦχος ἐν τῇ παρεμβολῇ τῶν ἀλλοφύλων ἐπορεύετο πορευόμενος καὶ ἐπλήθυνε· καὶ εἶπε Σαοὺλ πρὸς τὸν ἱερέα, συνάγαγε τὰς χεῖράς σου.

\vs{20}Καὶ ἀνέβη Σαοὺλ καὶ πᾶς ὁ λαὸς ὁ μετʼ αὐτοῦ, καὶ ἔρχονται ἕως τοῦ πολέμου· καὶ ἰδοὺ ἐγένετο ῥομφαία ἀνδρὸς ἐπὶ τὸν πλησίον αὐτοῦ, σύγχυσις μεγάλη σφόδρα.
\vs{21}Καὶ οἱ δοῦλοι οἱ ὄντες ἐχθὲς καὶ τρίτην ἡμέραν μετὰ τῶν ἀλλοφύλων οἱ ἀναβάντες εἰς τὴν παρεμβολὴν, ἐπεστράφησαν καὶ αὐτοὶ εἶναι μετὰ Ἰσραὴλ τῶν μετὰ Σαοὺλ καὶ Ἰωνάθαν.
\vs{22}Καὶ πᾶς Ἰσραὴλ οἱ κρυπτόμενοι ἐν τῷ ὄρει Ἐφραὶμ, καὶ ἤκουσαν ὅτι πεφεύγασιν οἱ ἀλλόφυλοι· καὶ συνάπτουσι καὶ αὐτοὶ ὀπίσω αὐτῶν εἰς πόλεμον·
\vs{23}καὶ ἔσωσε Κύριος ἐν τῇ ἡμέρᾳ ἐκείνῃ τὸν Ἰσραήλ· καὶ ὁ πόλεμος διῆλθε τὴν Βαμώθ· καὶ πᾶς ὁ λαὸς ἦν μετὰ Σαοὺλ ὡς δέκα χιλιάδες ἀνδρῶν. Καὶ ἦν ὁ πόλεμος διεσπαρμένος εἰς ὅλην πόλιν ἐν τῷ ὄρει Ἐφραίμ.

\vs{24}Καὶ Σαοὺλ ἠγνόησεν ἄγνοιαν μεγάλην ἐν τῇ ἡμέρᾳ ἐκείνῃ, καὶ ἀρᾶται τῷ λαῷ, λέγων, ἐπικατάρατος ὁ ἄνθρωπος ὃς φάγεται ἄρτον ἕως ἑσπέρας· καὶ ἐκδικήσω τὸν ἐχθρόν μου· καὶ οὐκ ἐγεύσατο πᾶς ὁ λαὸς ἄρτου·
\vs{25}καὶ πᾶσα ἡ γῆ ἠρίστα. Καὶ Ἰάαλ δρυμὸς ἦν μελισσῶνος κατὰ πρόσωπον τοῦ ἀγροῦ.
\vs{26}Καὶ εἰσῆλθεν ὁ λαὸς εἰς τὸν μελισσῶνα, καὶ ἰδοὺ ἐπορεύετο λαλῶν· καὶ ἰδοὺ οὐκ ἦν ἐπιστρέφων τὴν χεῖρα αὐτοῦ εἰς τὸ στόμα αὐτοῦ, ὅτι ἐφοβήθη ὁ λαὸς τὸν ὅρκον Κυρίου.
\vs{27}Καὶ Ἰωνάθαν οὐκ ἀκηκόει ἐν τῷ ὁρκίζειν τὸν πατέρα αὐτοῦ τὸν λαόν· καὶ ἐξέτεινε τὸ ἄκρον τοῦ σκήπτρου αὐτοῦ τοῦ ἐν τῇ χειρὶ αὐτοῦ, καὶ ἔβαψεν αὐτὸ εἰς τὸ κηρίον τοῦ μέλιτος, καὶ ἐπέστρεψε τὴν χεῖρα αὐτοῦ εἰς τὸ στόμα αὐτοῦ, καὶ ἀνέβλεψαν οἱ ὀφθαλμοὶ αὐτοῦ.
\vs{28}Καὶ ἀπεκρίθη εἷς ἐκ τοῦ λαοῦ, καὶ εἶπεν, ὁρκίσας ὥρκισε τὸν λαὸν ὁ πατήρ σου, λέγων ἐπικατάρατος ὁ ἄνθρωπος ὃς φάγεται ἄρτον σήμερον· καὶ ἐξελύθη ὁ λαός.
\vs{29}Καὶ ἔγνω Ἰωνάθαν, καὶ εἶπεν, ἀπήλλαχεν ὁ πατήρ μου τὴν γῆν· ἴδε, διότι εἶδον οἱ ὀφθαλμοί μου ὅτι ἐγευσάμην βραχύ τι τοῦ μέλιτος τούτου·
\vs{30}Ἀλλʼ ὅτι ἔφαγεν ἔσθων σήμερον ὁ λαὸς τῶν σκύλων τῶν ἐχθρῶν αὐτῶν ὧν εὗρεν, ὅτι νῦν ἂν μείζων ἦν ἡ πληγὴ ἡ ἐν τοῖς ἀλλοφύλοις.

\vs{31}Καὶ ἐπάταξεν ἐν τῇ ἡμέρᾳ ἐκείνῃ ἐκ τῶν ἀλλοφύλων ἐν Μαχμάς· καὶ ἐκοπίασεν ὁ λαὸς σφόδρα.
\vs{32}Καὶ ἐκλίθη ὁ λαὸς εἰς τὰ σκῦλα· καὶ ἔλαβεν ὁ λαὸς ποίμνια, καὶ βουκόλια, καὶ τέκνα βοῶν, καὶ ἔσφαξεν ἐπὶ τὴν γῆν, καὶ ἤσθιεν ὁ λαὸς σὺν τῷ αἵματι.
\vs{33}Καὶ ἀπηγγέλη Σαοὺλ, λέγοντες, ἡμάρτηκεν ὁ λαὸς τῷ Κυρίῳ, φαγὼν σὺν τῷ αἵματι· καὶ εἶπε Σαοὺλ, ἐκ Γεθθαὶμ κυλίσατέ μοι λίθον ἐνταῦθα μέγαν.
\vs{34}Καὶ εἶπε Σαοὺλ, διασπάρητε ἐν τῷ λαῷ, καὶ εἴπατε αὐτοῖς προσαγαγεῖν ἐνταῦθα ἕκαστος τὸν μόσχον αὑτοῦ, καὶ ἕκαστος τὸ πρόβατον αὑτοῦ· καὶ σφαζέτω ἐπὶ τούτου, καὶ οὐ μὴ ἁμάρτητε τῷ Κυρίῳ τοῦ ἐσθίειν σὺν τῷ αἵματι· καὶ προσῆγεν ὁ λαὸς ἕκαστος τὸ ἐν τῇ χειρὶ αὐτοῦ, καὶ ἔσφαζον ἐκεῖ.
\vs{35}Καὶ ᾠκοδόμησεν ἐκεῖ Σαοὺλ θυσιαστήριον τῷ Κυρίῳ· τοῦτο ἤρξατο Σαοὺλ οἰκοδομῆσαι θυσιαστήριον τῷ Κυρίῳ.

\vs{36}Καὶ εἶπε Σαούλ, καταβῶμεν ὀπίσω τῶν ἀλλοφύλων τὴν νύκτα, καὶ διαρπάσωμεν ἐν αὐτοῖς ἕως διαφαύσῃ ἡμέρα, καὶ μὴ ὑπολείπωμεν ἐν αὐτοῖς ἄνδρα· καὶ εἶπαν, πᾶν τὸ ἀγαθὸν ἐνώπιόν σου ποίει· καὶ εἶπεν ὁ ἱερεὺς, προσελθωμεν ἐνταῦθα πρὸς τὸν Θεόν.

\vs{37}Καὶ ἐπηρώτησε Σαοὺλ τὸν Θεὸν, εἰ καταβῶ ὀπίσω τῶν ἀλλοφύλων, εἰ παραδώσεις αὐτοὺς εἰς χεῖρας Ἰσραήλ; καὶ οὐκ ἀπεκρίθη αὐτῷ ἐν τῇ ἡμέρᾳ ἐκείνῃ.

\vs{38}Καὶ εἶπε Σαοὺλ, προσαγάγετε ἐνταῦθα πάσας τὰς γωνίας τοῦ Ἰσραὴλ, καὶ γνῶτε καὶ ἴδετε ἐν τίνι γέγονεν ἡ ἁμαρτία αὕτη σήμερον.
\vs{39}Ὅτι ζῇ Κύριος ὁ σώσας τὸν Ἰσραὴλ, ὅτι ἐὰν ἀποκριθῇ κατὰ Ἰωνάθαν τοῦ υἱοῦ μου, θανάτῳ ἀποθανεῖται· καὶ οὐκ ἦν ὁ ἀποκρινόμενος ἐκ παντὸς τοῦ λαοῦ.
\vs{40}Καὶ εἶπε παντὶ ἀνδρὶ Ἰσραὴλ, ὑμεῖς ἔσεστε εἰς δουλείαν, καὶ ἐγὼ καὶ Ἰωνάθαν ὁ υἱός μου ἐσόμεθα εἰς δουλείαν· καὶ εἶπεν ὁ λαὸς πρὸς Σαούλ, τὸ ἀγαθὸν ἐνώπιόν σου ποίει.
\vs{41}Καὶ εἶπε Σαοὺλ, Κύριε ὁ Θεὸς Ἰσραὴλ, τί ὅτι οὐκ ἀπεκρίθης τῷ δούλῳ σου σήμερον; εἰ ἐν ἐμοὶ ἢ ἐν Ἰωνάθαν τῷ υἱῷ μου ἡ ἀδικία; Κύριε ὁ Θεὸς Ἰσραὴλ, δὸς δήλους· καὶ ἐὰν τάδε εἴπῃ, δὸς δὴ τῷ λαῷ σου Ἰσραὴλ, δὸς δὴ ὁσιότητα· καὶ κληροῦται Ἰωνάθαν καὶ Σαούλ, καὶ ὁ λαὸς ἐξῆλθε.

\vs{42}Καὶ εἶπε Σαοὺλ, βάλλετε ἀναμέσον ἐμοῦ καὶ ἀναμέσον Ἰωνάθαν τοῦ υἱοῦ μου· ὃν ἂν κατακληρώσηται Κύριος, ἀποθανέτω· καὶ εἶπεν ὁ λαὸς πρὸς Σαοὺλ, οὐκ ἔστι τὸ ῥῆμα τοῦτο· καὶ κατεκράτησε Σαοὺλ τοῦ λαοῦ, καὶ βάλλουσιν ἀναμέσον αὐτοῦ καὶ ἀναμέσον Ἰωνάθαν τοῦ υἱοῦ αὐτοῦ, καὶ κατακληροῦται Ἰωνάθαν.
\vs{43}Καὶ εἶπε Σαοὺλ πρὸς Ἰωνάθαν, ἀπάγγειλόν μοι τί πεποίηκας· καὶ ἀπήγγειλεν αὐτῷ Ἰωνάθαν, καὶ εἶπε, γευόμενος ἐγευσάμην ἐν ἄκρῳ τῷ σκήπτρῳ τῷ ἐν τῇ χειρί μου βραχὺ μέλι, καὶ ἰδοὺ ἐγὼ ἀποθνήσκω.
\vs{44}Καὶ εἶπεν αὐτῷ Σαοὺλ, τάδε ποιήσαι μοι ὁ Θεὸς καὶ τάδε προσθείη, ὅτι θανάτῳ ἀποθανῇ σήμερον.
\vs{45}Καὶ εἶπεν ὁ λαὸς πρὸς Σαοὺλ, εἰ σήμερον θανατωθήσεται ὁ ποιήσας τὴν σωτηρίαν τὴν μεγάλην ταύτην ἐν Ἰσραήλ; ζῇ Κύριος, εἰ πεσεῖται τριχὸς τῆς κεφαλῆς αὐτοῦ ἐπὶ τὴν γῆν, ὅτι ὁ λαὸς τοῦ Θεοῦ ἐποίησε τὴν ἡμέραν ταύτην· καὶ προσηύξατο ὁ λαὸς περὶ Ἰωνάθαν ἐν τῇ ἡμέρᾳ ἐκείνῃ, καὶ οὐκ ἀπέθανε.
\vs{46}Καὶ ἀνέβη Σαοὺλ ἀπὸ ὄπισθε τῶν ἀλλοφύλων· καὶ οἱ ἀλλόφυλοι ἀπῆλθον εἰς τὸν τόπον αὐτῶν.

\vs{47}Καὶ Σαοὺλ ἔλαχε τοῦ βασιλεύειν, κατακληροῦται ἔργον ἐπὶ Ἰσραήλ· καὶ ἐπολέμει κύκλῳ πάντας τοὺς ἐχθροὺς αὐτοῦ εἰς τὸν Μωὰβ, καὶ εἰς τοὺς υἱοὺς ʼΑμμὼν, καὶ εἰς τοὺς υἱοὺς Ἐδὼμ, καὶ εἰς τὸν Βαιθαιὼρ, καὶ εἰς βασιλέα Σουβὰ, καὶ εἰς τοὺς ἀλλοφύλους· οὗ ἂν ἐστράφη, ἐσώζετο.
\vs{48}Καὶ ἐποίησε δύναμιν, καὶ ἐπάταξε τὸν ʼΑμαλήκ· καὶ ἐξείλετο τὸν Ἰσραὴλ ἐκ χειρὸς τῶν καταπατούντων αὐτόν.

\vs{49}Καὶ ἦσαν οἱ υἱοὶ Σαοὺλ Ἰωνάθαν, καὶ Ἰεσσιοὺ, καὶ Μελχισά· καὶ ὀνόματα τῶν δύο θυγατέρων αὐτοῦ, ὄνομα τῇ πρωτοτόκῳ Μερόβ, καὶ ὄνομα τῇ δευτέρᾳ Μελχόλ.
\vs{50}Καὶ ὄνομα τῇ γυναικὶ αὐτοῦ ʼΑχινοὸμ, θυγάτερ ʼΑχιμάας· καὶ ὄνομα τῷ ἀρχιστρατήγῳ αὐτοῦ ʼΑβεννηρ, υἱὸς Νὴρ, υἱοῦ οἰκείου Σαούλ.
\vs{51}Καὶ Κὶς πατὴρ Σαοὺλ, καὶ Νὴρ πατὴρ ʼΑβεννὴρ υἱὸς Ἰαμὶν, υἱοῦ ʼΑβιήλ.

\vs{52}Καὶ ἦν ὁ πόλεμος κραταιὸς ἐπὶ τοὺς ἀλλοφύλους πάσας τὰς ἡμέρας Σαούλ· καὶ ἰδὼν Σαοὺλ πάντα ἄνδρα δυνατὸν, καὶ πάντα ἄνδρα υἱὸν δυνάμεως, καὶ συνήγαγεν αὐτοὺς πρὸς αὑτόν.

\ch{15}
Καὶ εἶπε Σαμουὴλ πρὸς Σαοὺλ, ἐμὲ ἀπέστειλε Κύριος χρίσαί σε εἰς βασιλεα ἐπὶ Ἰσραήλ· καὶ νῦν ἄκουε τῆς φωνῆς Κυρίου.
\vs{2}Τάδε εἶπε Κύριος σαβαὼθ, νῦν ἐκδικήσω ἃ ἐποίησεν ʼΑμαλὴκ τῷ Ἰσραὴλ, ὡς ἀπήντησεν αὐτῷ ἐν τῇ ὁδῷ ἀναβαίνοντος αὐτοῦ ἐξ Αἰγύπτου.
\vs{3}Καὶ νῦν πορεύου, καὶ πατάξεις τὸν ʼΑμαλὴκ καὶ Ἱερὶμ καὶ πάντα τὰ αὐτοῦ, καὶ οὐ περιποιήσῃ ἐξ αὐτοῦ, καὶ ἐξολοθρεύσεις αὐτόν· καὶ ἀναθεματιεῖς αὐτὸν καὶ πάντα τὰ αὐτοῦ, καὶ οὐ φείσῃ ἀπʼ αὐτοῦ· καὶ ἀποκτενεῖς ἀπὸ ἀνδρὸς καὶ ἕως γυναικὸς, καὶ ἀπὸ νηπίου ἕως θηλάζοντος, καὶ ἀπὸ μόσχου ἕως προβάτου, καὶ ἀπὸ καμήλου ἕως ὄνου.

\vs{4}Καὶ παρήγγειλε Σαοὺλ τῷ λαῷ, καὶ ἐπισκέπτεται αὐτοὺς ἐν Γαλγάλοις τετρακοσίας χιλιάδας ταγμάτων, καὶ τὸν Ἰούδαν τριάκοντα χιλιάδας ταγμάτων.
\vs{5}Καὶ ἦλθε Σαοὺλ ἕως τῶν πόλεων ʼΑμαλὴκ, καὶ ἐνήδρευσεν ἐν τῷ χειμάῤῥῳ.
\vs{6}Καὶ εἶπε Σαοὺλ πρὸς τὸν Κιναῖον, ἄπελθε καὶ ἔκκλινον ἐκ μέσου τοῦ ʼΑμαληκίτου, μὴ προσθῶ σε μετʼ αὐτοῦ, καὶ σὺ ἐποίησας ἔλεος μετὰ τῶν υἱῶν Ἰσραὴλ, ἐν τῷ ἀναβαίνειν αὐτοὺς ἐξ Αἰγύπτου· καὶ ἐξέκλινεν ὁ Κιναῖος ἐκ μέσου ʼΑμαλήκ.
\vs{7}Καὶ ἐπάταξε Σαοὺλ τὸν Ἀμαλὴκ ἀπὸ Εὐειλὰτ ἕως Σοὺρ ἐπὶ προσώπου Αἰγύπτου.
\vs{8}Καὶ συνέλαβε τὸν ʼΑγὰγ βασιλέα ʼΑμαλὴκ ζῶντα, καὶ πάντα τὸν λαὸν καὶ Ἱερὶμ ἀπέκτεινεν ἐν στόματι ῥομφαίας.
\vs{9}Καὶ περιεποιήσατο Σαοὺλ καὶ πᾶς ὁ λαὸς τὸν ʼΑγὰγ ζῶντα, καὶ τὰ ἀγαθὰ τῶν ποιμνίων, καὶ τῶν βουκολίων, καὶ τῶν ἐδεσμάτων καὶ τῶν ἀμπελώνων, καὶ πάντων τῶν ἀγαθῶν, καὶ οὐκ ἐβούλοντο ἐξολοθρεῦσαι αὐτά· καὶ πᾶν ἔργον ἠτιμωμένον καὶ ἐξουδενωμένον ἐξωλόθρευσαν.

\vs{10}Καὶ ἐγενήθη ῥῆμα Κυρίου πρὸς Σαμουὴλ, λέγων,
\vs{11}παρακέκλημαι ὅτι ἐβασίλευσα τὸν Σαοὺλ εἰς βασιλέα, ὅτι ἀπέστρεψεν ἀπὸ ὄπισθέ μου, καὶ τοὺς λόγους μου οὐκ ἐτήρησε· καὶ ἠθύμησε Σαμουὴλ, καὶ ἐβόησε πρὸς Κύριον ὅλην τὴν νύκτα.
\vs{12}Καὶ ὤρθρισε Σαμουὴλ, καὶ ἐπορεύθη εἰς ἀπάντησιν Ἰσραὴλ το πρωΐ· καὶ ἀπηγγέλη τῷ Σαοὺλ, λέγοντες ἥκει Σαμουὴλ εἰς Κάρμηλον, καὶ ἀνέστακεν αὐτῷ χεῖρα· καὶ ἐπέστρεψε τὸ ἅρμα, καὶ κατέβη εἰς Γάλγαλα πρὸς Σαούλ, καὶ ἰδοὺ αὐτὸς ἀνέφερεν ὁλοκαύτωσιν τῷ Κυρίῳ, τὰ πρῶτα τῶν σκύλων ἤνεγκεν ἐξ ʼΑμαλήκ.

\vs{13}Καὶ παρεγένετο Σαμουὴλ πρὸς Σαούλ· καὶ εἶπεν αὐτῷ Σαοὺλ, εὐλογητὸς σὺ τῷ Κυρίῳ· ἔστησα πάντα ὅσα ἐλάλησε Κύριος.
\vs{14}Καὶ εἶπε Σαμουὴλ, καὶ τίς ἡ φωνὴ τοῦ ποιμνίου τούτου ἐν τοῖς ὠσί μου, καὶ φωνὴ τῶν βοῶν ἣν ἐγὼ ἀκούω;
\vs{15}Καὶ εἶπε Σαούλ, ἐξ ʼΑμαλὴκ ἤνεγκα αὐτά, ἃ περιεποιήσατο ὁ λαός, τὰ κράτιστα τοῦ ποιμνίου, καὶ τῶν βοῶν, ὅπως τυθῇ Κυρίῳ τῷ Θεῷ σου, καὶ τὰ λοιπὰ ἐξωλόθρευσα.
\vs{16}Καὶ εἶπε Σαμουὴλ πρὸς Σαοὺλ, ἄνες, καὶ ἀπαγγελῶ σοι ἃ ἐλάλησε Κύριος πρὸς μὲ τὴν νύκτα· καὶ εἶπεν αὐτῷ, λάλησον.

\vs{17}Καὶ εἶπε Σαμουὴλ πρὸς Σαούλ, οὐχὶ μικρὸς εἶ σὺ ἐνώπιον αὐτοῦ, ἡγούμενος σκήπτρου φυλῆς Ἰσραήλ; καὶ ἔχρισέ σε Κύριος εἰς βασιλέα ἐπὶ Ἰσραήλ.
\vs{18}Καὶ ἀπέστειλέ σε Κύριος ἐν ὁδῷ, καὶ εἶπέ σοι, πορεύθητι, καὶ ἐξολόθρευσον· ἀνελεῖς τοὺς ἁμαρτάνοντας εἰς ἐμέ, τὸν Ἀμαλὴκ, καὶ πολεμήσεις αὐτοὺς ἕως συντελέσῃς αὐτούς.
\vs{19}Καὶ ἱνατί οὐκ ἤκουσας φωνῆς Κυρίου, ἀλλʼ ὥρμησας τοῦ θέσθαι ἐπὶ τὰ σκῦλα, καὶ ἐποίησας τὸ πονηρὸν ἐνώπιον Κυρίου;
\vs{20}Καὶ εἶπε Σαοὺλ πρὸς Σαμουὴλ, διὰ τὸ ἀκοῦσαί με τῆς φωνῆς τοῦ λαοῦ, καὶ ἐπορεύθην τῇ ὁδῷ ᾗ ἀπέστειλέ με Κύριος, καὶ ἤγαγον τὸν ʼΑγὰγ βασιλέα ʼΑμαλὴκ, καὶ τὸν ʼΑμαλὴκ ἐξωλόθρευσα.
\vs{21}Καὶ ἔλαβεν ὁ λαὸς τῶν σκύλων ποίμνια καὶ βουκόλια τὰ πρῶτα τοῦ ἐξολοθρεύματος, θύσαι ἐνώπιον Κυρίου Θεοῦ ἡμῶν ἐν Γαλγάλοις.
\vs{22}Καὶ εἶπε Σαμουήλ, εἰ θελητὸν τῷ Κυρίῳ ὁλοκαυτώματα καὶ θυσίας, ὡς τὸ ἀκοῦσαι φωνῆς Κυρίου; ἰδοὺ ἀκοὴ ὑπὲρ θυσίαν ἀγαθὴν, καὶ ἡ ἐπακρόασις ὑπὲρ στέαρ κριῶν.
\vs{23}Ὅτι ἁμαρτία οἰώνισμά ἐστιν, ὀδύνην καὶ πόνονς θεραφὶν ἐπάγουσιν· ὅτι ἐξουδένωσας τὸ ῥῆμα Κυρίου, καὶ ἐξουδενώσει σε Κύριος μὴ εἶναι βασιλέα ἐπὶ Ἰσραήλ.

\vs{24}Καὶ εἶπε Σαοὺλ πρὸς Σαμουὴλ, ἡμάρτηκα, ὅτι παρέβην τὸν λόγον Κυρίου καὶ τὸ ῥῆμά σου, ὅτι ἐφοβήθην τὸν λαὸν, καὶ ἤκουσα τῆς φωνῆς αὐτῶν.
\vs{25}Καὶ νῦν ἆρον δὴ τὸ ἁμάρτημά μου, καὶ ἀνάστρεψον μετʼ ἐμοῦ, καὶ προσκυνήσω Κυρίῳ τῷ Θεῷ σου.
\vs{26}Καὶ εἶπε Σαμουὴλ πρὸς Σαοὺλ, οὐκ ἀναστρέφω μετὰ σοῦ; ὅτι ἐξουδένωσας τὸ ῥῆμα Κυρίου, καὶ ἐξουδενώσει σε Κύριος τοῦ μὴ εἶναι βασιλέα ἐπὶ τὸν Ἰσραήλ.

\vs{27}Καὶ ἐπέστρεψε Σαμουὴλ τὸ πρόσωπον αὐτοῦ τοῦ ἀπελθεῖν, καὶ ἐκράτησε Σαοὺλ τοῦ πτερυγίου τῆς διπλοΐδος αὐτοῦ, καὶ διέῤῥηξεν αὐτό.
\vs{28}Καὶ εἶπε πρὸς αὐτὸν Σαμουὴλ, διέῤῥηξε Κύριος τὴν βασιλείαν σου ἀπὸ Ἰσραὴλ ἐκ χείρος σου σήμερον, καὶ δώσει αὐτὴν τῷ πλησίον σου τῷ ἀγαθῷ ὑπὲρ σέ·
\vs{29}Καὶ διαιρεθήσεται Ἰσραὴλ εἰς δύο, καὶ οὐκ ἀποστρέψει οὐδὲ μετανοήσει, ὅτι οὐχ ὡς ἄνθρωπός ἐστι τοῦ μετανοῆσαι αὐτός.
\vs{30}Καὶ εἶπε Σαοὺλ, ἡμάρτηκα, ἀλλὰ δόξασόν με δὴ ἐνώπιον πρεσβυτέρων Ἰσραὴλ καὶ ἐνώπιον λαοῦ μου, καὶ ἀνάστρεψον μετʼ ἐμοῦ, καὶ προσκυνήσω Κυρίῳ τῷ Θεῷ σου.
\vs{31}Καὶ ἀνέστρεψε Σαμουὴλ ὀπίσω Σαοὺλ, καὶ προσεκύνησε τῷ Κυρίῳ.

\vs{32}Καὶ εἶπε Σαμουὴλ, προσαγάγετέ μοι τὸν ʼΑγὰγ βασιλέα ʼΑμαλήκ· καὶ προσῆλθε πρὸς αὐτὸν ʼΑγὰγ τρέμων· καὶ εἶπεν ʼΑγὰγ, εἰ οὕτω πικρὸς ὁ θάνατος.
\vs{33}Καὶ εἶπε Σαμουὴλ πρὸς ʼΑγὰγ, καθότι ἠτέκνωσε γυναῖκας ἡ ῥομφαία σου, οὕτως ἀτεκνωθήσεται ἐκ γυναικῶν ἡ μήτηρ σου· καὶ ἔσφαξε Σαμουὴλ τὸν ʼΑγὰγ ἐνώπιον Κυρίου ἐν Γαλγάλ.

\vs{34}Καὶ ἀπῆλθε Σαμουὴλ εἰς ʼΑρμαθαίμ· καὶ Σαοὺλ ἀνέβη εἰς τὸν οἶκον αὐτοῦ εἰς Γαβαά.
\vs{35}Καὶ οὐ προσέθετο ἔτι Σαμουὴλ ἰδεῖν τὸν Σαοὺλ ἕως ἡμέρας θανάτου αὐτοῦ, ὅτι ἐπένθει Σαμουὴλ ἐπὶ Σαοὺλ, καὶ Κύριος μετεμελήθη ὅτι ἐβασίλευσε τὸν Σαοὺλ ἐπὶ Ἰσραήλ.

\ch{16}
Καὶ εἶπε Κύριος πρὸς Σαμουήλ, ἕως πότε σὺ πενθεῖς ἐπὶ Σαοὺλ, κᾀγω ἐξουδένωκα αὐτὸν, μὴ βασιλεύειν ἐπὶ Ἰσραήλ; πλῆσον τὸ κέρας σου ἐλαίου, καὶ δεῦρο ἀποστείλω σε πρὸς Ἰεσσαὶ ἕως Βηθλεὲμ, ὅτι ἑώρακα ἐν τοῖς υἱοῖς αὐτοῦ ἐμοὶ βασιλέα.
\vs{2}Καὶ εἶπε Σαμουὴλ πῶς πορευθῶ; καὶ ἀκούσεται Σαοὺλ, καὶ ἀποκτενεῖ με· καὶ εἶπε Κύριος, δάμαλιν βοῶν λάβε ἐν τῇ χειρί σου, καὶ ἐρεῖς, θῦσαι τῷ Κυρίῳ ἥκω.
\vs{3}Καὶ καλέσεις τὸν Ἰεσσαὶ εἰς τὴν θυσίαν, καὶ γνωριῶ σοι ἃ ποιήσεις· καὶ χρίσεις ὃν ἂν εἴπω πρὸς σέ.

\vs{4}Καὶ ἐποίησε Σαμουὴλ πάντα ἃ ἐλάλησεν αὐτῷ Κύριος· καὶ ἦλθεν εἰς Βηθλεὲμ, καὶ ἐξέστησαν οἱ πρεσβύτεροι τῆς πόλεως τῇ ἀπαντήσει αὐτοῦ, καὶ εἶπαν, ἢ εἰρήνη ἡ εἴσοδός σου ὁ βλέπων;
\vs{5}Καὶ εἶπεν, εἰρήνη· θῦσαι τῷ Κυρίῳ ἥκω· ἁγιάσθητε καὶ εὐφράνθητε μετʼ ἐμοῦ σήμερον· καὶ ἡγίασε τὸν Ἰεσσαὶ καὶ τοὺς υἱοὺς αὐτοῦ, καὶ ἐκάλεσεν αὐτοὺς εἰς τὴν θυσίαν.
\vs{6}Καὶ ἐγενήθη ἐν τῷ εἰσιέναι αὐτοὺς, καὶ εἶδε τὸν Ἑλιάβ, καὶ εἶπεν, ἀλλὰ καὶ ἐνώπιον Κυρίου χριστὸς αὐτοῦ.
\vs{7}Καὶ εἶπε Κύριος πρὸς Σαμουὴλ, μὴ ἐπιβλέψῃς ἐπὶ τὴν ὄψιν αὐτοῦ, μηδὲ εἰς τὴν ἕξιν μεγέθους αὐτοῦ, ὅτι ἐξουδένωκα αὐτόν· ὅτι οὐχ ὡς ἐμβλέψεται ἄνθρωπος, ὄψεται ὁ Θεός· ὅτι ἄνθρωπος ὅψεται εἰς πρόσωπον, ὁ δὲ Θεὸς ὄψεται εἰς καρδίαν.
\vs{8}Καὶ ἐκάλεσεν Ἰεσσαὶ τὸν ʼΑμιναδὰβ, καὶ παρῆλθε κατὰ πρόσωπον Σαμουήλ. καὶ εἶπεν οὐδὲ τοῦτον ἐξελέξατο ὁ Θεός.
\vs{9}Καὶ παρήγαγεν Ἰεσσαὶ τὸν Σαμά· καὶ εἶπε, καὶ ἐν τούτῳ οὐκ ἐξελέξατο Κύριος.
\vs{10}Καὶ παρήγαγεν Ἰεσσαὶ τοὺς ἑπτὰ υἱοὺς αὐτοῦ ἐνώπιον Σαμουήλ· καὶ εἶπε Σαμουὴλ, οὐκ ἐξελέξατο Κύριος ἐν τούτοις.

\vs{11}Καὶ εἶπε Σαμουὴλ πρὸς Ἰεσσαὶ, ἐκλελοίπασι τὰ παιδάρια; καὶ εἶπεν, ἔτι ὁ μικρὸς ἰδοὺ ποιμαίνει ἐν τῷ ποιμνίῳ· καὶ εἶπε Σαμουὴλ πρὸς Ἰεσσαὶ, ἀπόστειλον καὶ λάβε αὐτὸν, ὅτι οὐ μὴ κατακλιθῶμεν ἕως τοῦ ἐλθεῖν αὐτόν.
\vs{12}Καὶ ἀπέστειλε καὶ εἰσήγαγεν αὐτόν· καὶ αὐτὸς πυῤῥάκης μετὰ κάλλους ὀφθαλμῶν, καὶ ἀγαθὸς ὁράσει Κυρίῳ· καὶ εἶπε Κύριος πρὸς Σαμουὴλ, ἀνάστα καὶ χρίσον τὸν Δαυείδ, ὅτι οὗτός ἐστιν ἀγαθός.
\vs{13}Καὶ ἔλαβε Σαμουὴλ τὸ κέρας τοῦ ἐλαίου, καὶ ἔχρισεν αὐτὸν ἐν μέσῳ τῶν ἀδελφῶν αὐτοῦ· καὶ ἐφήλατο πνεῦμα Κυρίου ἐπὶ Δαυὶδ ἀπὸ τῆς ἡμέρας ἐκείνης καὶ ἐπάνω· καὶ ἀνέστη Σαμουὴλ, καὶ ἀπῆλθεν εἰς ʼΑρμαθαίμ.

\vs{14}Καὶ πνεῦμα Κυρίου ἀπέστη ἀπὸ Σαοὺλ, καὶ ἔπνιγεν αὐτὸν πνεῦμα πονηρὸν παρὰ Κυρίου.
\vs{15}Καὶ εἶπαν οἱ παῖδες Σαοὺλ πρὸς αὐτόν, ἰδοὺ δὴ, πνεῦμα Κυρίου πονηρὸν πνίγει σε·
\vs{16}Εἰπάτωσαν δὴ οἱ δοῦλοί σου ἐνώπιόν σου, καὶ ζητησάτωσαν τῷ κυρίῳ ἡμῶν ἄνδρα εἰδότα ψάλλειν ἐν κινύρᾳ· καὶ ἔσται ἐν τῷ εἶναι πνεῦμα πονηρὸν ἐπὶ σοὶ, καὶ ψαλῇ ἐν τῇ κινύρᾳ αὐτοῦ, καὶ ἀγαθόν σοι ἔσται καὶ ἀναπαύσει σε.
\vs{17}Καὶ εἶπε Σαοὺλ πρὸς τοὺς παῖδας αὐτοῦ, ἴδετε δή μοι ἄνδρα ὀρθῶς ψάλλοντα, καὶ εἰσαγάγετε αὐτὸν πρὸς μέ.
\vs{18}Καὶ ἀπεκρίθη εἷς τῶν παιδαρίων αὐτοῦ, καὶ εἶπεν, ἰδοὺ ἑώρακα υἱὸν τῷ Ἰεσσαὶ Βηθλεεμίτην, καὶ αὐτὸν εἰδότα ψαλμὸν, καὶ ὁ ἀνὴρ συνετὸς, καὶ πολεμιστὴς, καὶ σοφὸς λόγῳ, καὶ ὁ ἀνὴρ ἀγαθὸς τῷ εἴδει, καὶ Κύριος μετʼ αὐτοῦ.
\vs{19}Καὶ ἀπέστειλε Σαοὺλ ἀγγέλους πρὸς Ἰεσσαὶ, λέγων, ἀπόστειλον πρὸς μὲ τὸν υἱόν σου Δαυὶδ τὸν ἐν τῷ ποιμνίῳ σου.
\vs{20}Καὶ ἔλαβεν Ἰεσσαὶ γόμορ ἄρτων, καὶ ἀσκὸν οἴνου, καὶ ἔριφον αἰγῶν ἕνα, καὶ ἐξαπέστειλεν ἐι· χειρὶ Δαυὶδ τοῦ υἱοῦ αὐτοῦ πρὸς Σαούλ.

\vs{21}Καὶ εἰσῆλθε Δαυὶδ πρὸς Σαοὺλ, καὶ παρειστήκει ἐνώπιον αὐτοῦ, καὶ ἠγάπησεν αὐτὸν σφόδρα· καὶ ἐγενήθη αὐτῷ αἴρων τὰ σκεύη αὐτοῦ.
\vs{22}Καὶ ἀπέστειλε Σαοὺλ πρὸς Ἰεσσαὶ, λέγων, παριστάσθω δὴ Δαυὶδ ἐνώπιον ἐμοῦ, ὅτι εὗρε χάριν ἐν ὀφθαλμοῖς μου.
\vs{23}Καὶ ἐγενήθη ἐν τῷ εἶναι πνεῦμα πονηρὸν ἐπὶ Σαοὺλ, καὶ ἐλάμβανε Δαυὶδ τὴν κινύραν, καὶ ἔψαλλεν ἐν χειρὶ αὐτοῦ, καὶ ἀνέψυχε Σαοὺλ, καὶ ἀγαθὸν αὐτῷ, καὶ ἀφίστατο ἀπʼ αὐτοῦ τὸ πνεῦμα τὸ πονηρόν.

\ch{17}
Καὶ συνάγουσιν ἀλλόφυλοι τὰς παρεμβολὰς αὐτῶν εἰς πόλεμον, καὶ συνάγονται εἰς Σοκχὼθ τῆς Ἰουδαίας, καὶ παρεμβάλλουσιν ἀναμέσον Σοκχὼθ, καὶ ἀναμέσον ʼΑζηκὰ Ἐφερμέν.
\vs{2}Καὶ Σαοὺλ καὶ οἱ ἄνδρες Ἰσραὴλ συνάγονται, καὶ παρεμβάλλουσιν ἐν τῇ κοιλάδι αὐτοὶ, καὶ παρατάσσονται εἰς πόλεμον ἐξεναντίας τῶν ἀλλοφύλων.
\vs{3}Καὶ οἱ ἀλλόφυλοι ἵστανται ἐπὶ τοῦ ὄρους ἐνταῦθα, καὶ Ἰσραὴλ ἵσταται ἐπὶ τοῦ ὄρους ἐνταῦθα, καὶ ὁ αὐλὼν ἀναμέσον αὐτῶν.

\vs{4}Καὶ ἐξῆλθεν ἀνὴρ δυνατὸς ἐκ τῆς παρατάξεως τῶν ἀλλοφύλων, Γολιὰθ ὄνομα αὐτῷ ἐκ Γὲθ, ὕψος αὐτοῦ τεσσάρων πήχεων καὶ σπιθαμῆς.
\vs{5}Καὶ περικεφαλαῖα ἐπὶ τῆς κεφαλῆς αὐτοῦ, καὶ θώρακα ἁλυσιδωτὸν αὐτὸς ἐνδεδυκώς· καὶ ὁ σταθμὸς τοῦ θώρακος αὐτοῦ, πέντε χιλιάδες σίκλων χαλκοῦ καὶ σιδήρου.
\vs{6}Καὶ κνημῖδες χαλκαῖ ἐπὶ τῶν σκελῶν αὐτοῦ, καὶ ἀσπὶς χαλκῆ ἀναμέσον τῶν ὤμων αὐτοῦ.
\vs{7}Καὶ ὁ κοντὸς τοῦ δόρατος αὐτοῦ ὡσεὶ μέσακλον ὑφαινόντων, καὶ ἡ λόγχη αὐτοῦ ἑξακοσίων σίκλων σιδήρου· καὶ ὁ αἴρων τὰ ὅπλα αὐτοῦ προεπορεύετο αὐτοῦ.
\vs{8}Καὶ ἔστη καὶ ἀνεβόησεν εἰς τὴν παράταξιν Ἰσραὴλ, καὶ εἶπεν αὐτοῖς, τί ἐκπορεύεσθε παρατάξασθαι πολέμῳ ἐξεναντίας ἡμῶν; οὐκ ἐγώ εἰμι ἀλλόφυλος, καὶ ὑμεῖς Ἐβραῖοι τοῦ Σαούλ; ἐκλέξασθε ἐαυτοῖς ἄνδρα, καὶ καταβήτω πρὸς με.
\vs{9}Καὶ ἐὰν δυνηθῇ πολεμῆσαι πρός με, καὶ ἐαν πατάξῃ με, καὶ ἐσόμεθα ὑμῖν εἰς δούλους· ἐὰν δὲ ἐγὼ δυνηθῶ καὶ πατάξω αὐτὸν, ἔσεσθε ἡμῖν εἰς δούλους, καὶ δουλεύσετε ἡμῖν.
\vs{10}Καὶ εἶπεν ὁ ἀλλόφυλος, ἰδοὺ ἐγὼ ὠνείδισα τὴν παράταξιν Ἰσραὴλ σήμερον ἐν τῇ ἡμέρᾳ ταύτῃ· δότε μοι ἄνδρα, καὶ μονομαχήσομεν ἀμφότεροι.

\vs{11}Καὶ ἤκουσε Σαοὺλ καὶ πᾶς Ἰσραὴλ τὰ ῥήματα τοῦ ἀλλοφύλου ταῦτα, καὶ ἐξέστησαν καὶ ἐφοβήθησαν σφόδρα.

\vs{32}Καὶ εἶπε Δαυὶδ πρὸς Σαοὺλ, μὴ δὴ συμπεσέτω καρδία τοῦ κυρίου μου ἐπʼ αὐτόν· ὁ δοῦλός σου πορεύσεται καὶ πολεμήσει μετὰ τοῦ ἀλλοφύλου τούτου.
\vs{33}Καὶ εἶπε Σαοὺλ πρὸς τὸν Δαυὶδ, οὐ μὴ δυνήσῃ πορευθῆναι πρὸς τὸν ἀλλόφυλον τοῦ πολεμεῖν μετʼ αὐτοῦ, ὅτι παιδάριον εἶ σὺ, καὶ αὐτὸς ἀνὴρ πολεμιστὴς ἐκ νεότητος αὐτοῦ.

\vs{34}Καὶ εἶπε Δαυὶδ πρὸς Σαοὺλ, ποιμαίνων ἦν ὁ δοῦλός σου τῷ πατρὶ αὐτοῦ ἐν τῷ ποιμνίῳ· καὶ ὅταν ἤρχετο ὁ λέων, καὶ ἡ ἄρκος, καὶ ἐλάμβανε πρόβατον ἐκ τῆς ἀγέλης,
\vs{35}καὶ ἐξεπορευόμην ὀπίσω αὐτοῦ καὶ ἐπάταξα αὐτόν, καὶ ἐξέσπασα ἐκ τοῦ στόματος αὐτοῦ· καὶ εἰ ἐπανίστατο ἐπʼ ἐμὲ, καὶ ἐκράτησα τοῦ φάρυγγος αὐτοῦ, καὶ ἐπάταξα, καὶ ἐθανάτωσα αὐτόν·
\vs{36}Καὶ τὸν λέοντα καὶ τὴν ἄρκον ἔτυπτεν ὁ δοῦλός σου, καὶ ἔσται ὁ ἀλλόφυλος ὁ ἀπερίτμητος ὡς ἓν τούτων· οὐχὶ πορεύσομαι καὶ πατάξω αὐτὸν, καὶ ἀφελῶ σήμερον ὄνειδος ἐξ Ἰσραήλ; διότι τίς ὁ ἀπερίτμητος οὗτος, ὃς ὠνείδισε παράταξιν Θεοῦ ζῶντος;
\vs{37}Κύριος ὃς ἐξείλατό με ἐκ χειρὸς τοῦ λέοντος καὶ ἐκ χειρὸς τῆς ἄρκτου, αὐτὸς ἐξελεῖταί με ἐκ χειρὸς τοῦ ἀλλοφύλου τοῦ ἀπεριτμήτου τούτου· καὶ εἶπε Σαοὺλ πρὸς Δαυὶδ, πορεύου, καὶ ἔσται Κύριος μετὰ σοῦ.

\vs{38}Καὶ ἐνέδυσε Σαοὺλ τὸν Δαυὶδ μανδύαν, καὶ τὴν περικεφαλαίαν χαλκῆν περὶ τὴν κεφαλὴν αὐτοῦ,
\vs{39}καὶ ἔζωσε τὸν Δαυὶδ τὴν ῥομφαίαν αὐτοῦ ἐπάνω τοῦ μανδύου αὐτοῦ· καὶ ἐκοπίασε περιπατήσας ἅπαξ καὶ δίς· καὶ εἶπε Δαυὶδ πρὸς Σαοὺλ, οὐ μὴ δύνωμαι πορευθῆναι ἐν τούτοις, ὅτι οὐ πεπείραμαι· καὶ ἀφαιροῦσιν αὐτὰ ἀπʼ αὐτοῦ.
\vs{40}Καὶ ἔλαβε τὴν βακτηρίαν αὐτοῦ ἐν τῇ χειρὶ αὐτοῦ, καὶ ἐξελέξατο ἑαυτῷ πέντε λίθους λείους ἐκ τοῦ χειμάῤῥου, καὶ ἔθετο αὐτοὺς ἐν τῷ καδίῳ τῷ ποιμενικῷ τῷ ὄντι αὐτῷ εἰς συλλογὴν, καὶ σφενδόνη αὐτοῦ ἐν τῇ χειρὶ αὐτοῦ· καὶ προσῆλθε πρὸς τὸν ἄνδρα τὸν ἀλλόφυλον.

\vs{42}Καὶ εἶδε Γολιὰθ τὸν Δαυὶδ, καὶ ἐξητίμασεν αὐτόν· ὅτι αὐτὸς ἦν παιδάριον, καὶ αὐτὸς πυῤῥάκης μετὰ κάλλους ὀφθαλμὼν.
\vs{43}Καὶ εἶπεν ὁ ἀλλόφυλος πρὸς Δαυὶδ, ὡσεὶ κύων ἐγώ εἰμι, ὅτι σὺ ἔρχῃ ἐπʼ ἐμὲ ἐν ῥάβδῳ καὶ λίθοις; καὶ εἶπε Δαυὶδ, οὐχί, ἀλλʼ ἢ χείρωυ κυνός· καὶ κατηράσατο ὁ ἀλλόφυλος τὸν Δαυὶδ ἐν τοῖς θεοῖς αὐτοῦ.
\vs{44}Καὶ εἶπεν ὁ ἀλλόφυλος πρὸς Δαυὶδ, δεῦρο πρὸς μὲ, καὶ δώσω τὰς σάρκας σου τοῖς πετεινοῖς τοῦ οὐρανοῦ καὶ τοῖς κτήνεσι τῆς γῆς.

\vs{45}Καὶ εἶπε Δαυὶδ πρὸς τὸν ἀλλόφυλον, σὺ ἔρχῃ πρὸς μὲ ἐν ῥομφαίᾳ καὶ ἐν δόρατι καὶ ἐν ἀσπίδι, κᾀγὼ πορεύομαι πρὸς σὲ ἐν ὀνόματι Κυρίου Θεοῦ σαβαὼθ παρατάξεως Ἰσραὴλ, ἣν ὠνείδισας σήμερον,
\vs{46}καὶ ἀποκλείσει σε Κύριος σήμερον εἰς τὴν χεῖρά μου, καὶ ἀποκτενῶ σε, καὶ ἀφελῶ τὴν κεφαλήν σου ἀπὸ σοῦ, καὶ δώσω τὰ κῶλά σου καὶ τὰ κῶλα παρεμβολῆς ἀλλοφύλων ἐν ταύτῃ τῇ ἡμέρᾳ τοῖς πετεινοῖς τοῦ οὐρανοῦ καὶ τοῖς θηρίοις τῆς γῆς· καὶ γνώσεται πᾶσα ἡ γῆ, ὅτι ἔστι Θεὸς ἐν Ἰσραήλ.
\vs{47}Καὶ γνώσεται πᾶσα ἡ ἐκκλησία αὕτη, ὅτι οὐκ ἐν ῥομφαίᾳ καὶ δόρατι σώζει Κύριος, ὅτι τοῦ Κυρίου ὁ πόλεμος, καὶ παραδώσει Κύριος ὑμᾶς εἰς χεῖρας ἡμῶν.

\vs{48}Καὶ ἀνέστη ὁ ἀλλόφυλος καὶ ἐπορεύθη εἰς συνάντησιν Δαυίδ.
\vs{49}Καὶ ἐξέτεινε Δαυὶδ τὴν χεῖρα αὐτοῦ εἰς τὸ κάδιον, καὶ ἔλαβεν ἐκεῖθεν λίθον ἕνα, καὶ ἐσφενδόνησε, καὶ ἐπάταξε τὸν ἀλλόφυλον εἰς τὸ μέτωπον αὐτοῦ, καὶ διέδυ ὁ λίθος διὰ τῆς περικεφαλαίας εἰς τὸ μέτωπον αὐτοῦ, καὶ ἔπεσεν ἐπὶ πρόσωπον αὐτοῦ ἐπὶ τὴν γῆν.
\vs{51}Καὶ ἔδραμε Δαυὶδ, καὶ ἐπέστη ἐπʼ αὐτὸν, καὶ ἔλαβε τὴν ῥομφαίαν αὐτοῦ, καὶ ἐθανάτωσεν αὐτὸν, καὶ ἀφεῖλε τὴν κεφαλὴν αὐτοῦ· καὶ εἶδον οἱ ἀλλόφυλοι, ὅτι τέθνηκεν ὁ δυνατὸς αὐτῶν, καὶ ἔφυγον.

\vs{52}Καὶ ἀνίστανται ἄνδρες Ἰσραὴλ καὶ Ἰούδα, καὶ ἠλάλαξαν, καὶ κατεδίωξαν ὀπίσω αὐτῶν ἕως εἰσόδου Γὲθ, καὶ ἕως τῆς πύλης ʼΑσκάλωνος· καὶ ἔπεσον τραυματίαι τῶν ἀλλοφύλων ἐν τῇ ὁδῷ τῶν πυλῶν καὶ ἕως Γὲθ, καὶ ἕως ʼΑκκαρών.
\vs{53}Καὶ ἀνέστρεψαν ἄνδρες Ἰσραὴλ ἐκκλίνοντες ὀπίσω τῶν ἀλλοφύλων, καὶ κατεπάτουν τὰς παρεμβολὰς αὐτῶν.
\vs{54}Καὶ ἔλαβε Δαυὶδ τὴν κεφαλὴν τοῦ ἀλλοφύλου, καὶ ἤνεγκεν αὐτὴν εἰς Ἰερουσαλὴμ, καὶ τὰ σκεύη αὐτοῦ ἔθηκεν ἐν τῷ σκηνώματι αὐτοῦ.

\ch{18}
\vs{6}Καὶ ἐξῆλθον αἱ χορεύουσαι εἰς συνάντησιν Δαυὶδ ἐκ πασῶν πόλεων Ἰσραὴλ ἐν τυμπάνοις, καὶ ἐν χαρμοσύνῃ, καὶ ἐν κυμβάλοις.
\vs{7}Καὶ ἐξῆρχον αἱ γυναῖκες, καὶ ἔλεγον, ἐπάταξε Σαοὺλ ἐν χιλιάσιν αὐτοῦ, καὶ Δαυὶδ ἐν μυριάσιν αὐτοῦ.
\vs{8}Καὶ πονηρὸν ἐφάνη τὸ ῥῆμα ἐν ὀφθαλμοῖς Σαοὺλ περὶ τοῦ λόγου τούτου, καὶ εἶπε, τῷ Δαυὶδ ἔδωκαν τὰς μυριάδας, καὶ ἐμοὶ ἔδωκαν τὰς χιλιάδας.

\vs{12}Καὶ ἐφοβήθη Σαοὺλ ἀπὸ προσώπου Δαυὶδ,
\vs{13}καὶ ἀπέστησεν αὐτὸν ἀπʼ αὐτοῦ, καὶ κατέστησεν αὐτὸν ἑαυτῷ χιλίαρχον· καὶ ἐξεπορεύετο καὶ εἰσεπορεύετο ἔμπροσθεν τοῦ λαοῦ.
\vs{14}Καὶ ἦν Δαυὶδ ἐν πάσαις ταῖς ὁδοῖς αὐτοῦ συνιῶν, καὶ Κύριος ἦν μετʼ αὐτοῦ.
\vs{15}Καὶ εἶδε Σαοὺλ ὡς αὐτὸς συνιεῖ σφόδρα, καὶ εὐλαβεῖτο ἀπὸ προσώπου αὐτοῦ.
\vs{16}Καὶ πᾶς Ἰσραὴλ καὶ Ἰούδας ἠγάπα τὸν Δαυὶδ, ὅτι αὐτὸς εἰσεπορεύετο καὶ ἐξεπορεύετο πρὸ προσώπου τοῦ λαοῦ.

\vs{20}Καὶ ἠγάπησε Μελχὸλ ἡ θυγάτηρ Σαοὺλ τὸν Δαυίδ· καὶ ἀπηγγέλη τῷ Σαοὺλ, καὶ ηὐθύνθη ἐν τοῖς ὀφθαλμοῖς αὐτοῦ.
\vs{21}Καὶ εἶπε Σαοὺλ, δώσω αὐτὴν αὐτῷ, καὶ ἔσται αὐτῷ εἰς σκάνδαλον· καὶ ἦν ἐπὶ Σαοὺλ χεὶρ ἀλλοφύλων.
\vs{22}Καὶ ἐνετείλατο Σαοὺλ τοῖς παισὶν αὐτοῦ, λέγων, λαλήσατε ὑμεῖς λάθρα τῷ Δαυὶδ, λέγοντες, ἰδοὺ θέλει ἐν σοὶ ὁ βασιλεὺς, καὶ πάντες οἱ παῖδες αὐτοῦ ἀγαπῶσί σε, καὶ σὺ ἐπιγάμβρευσον τῷ βασιλεῖ.
\vs{23}Καὶ ἐλάλησαν οἱ παῖδες Σαοὺλ εἰς τὰ ὦτα Δαυὶδ τὰ ῥήματα ταῦτα· καὶ εἶπε Δαυὶδ, εἰ κούφον ἐν ὀφθαλμοῖς ὑμῶν ἐπιγαμβρεῦσαι βασιλεῖ; κᾀγὼ ἀνὴρ ταπεινὸς, καὶ οὐχὶ ἔνδοξος.
\vs{24}Καὶ ἀπήγγειλαν οἱ παῖδες Σαοὺλ αὐτῷ κατὰ τὰ ῥήματα ταῦτα, ἃ ἐλάλησε Δαυίδ.
\vs{25}Καὶ εἶπε Σαοὺλ, τάδε ἐρεῖτε τῷ Δαυὶδ, οὐ βούλεται ὁ βασιλεὺς ἐν δόματι ἀλλʼ ἢ ἐν ἑκατὸν ἀκροβυστίαις ἀλλοφύλων ἐκδικῆσαι ἐχθροὺς τοῦ βασιλέως· καὶ Σαοὺλ ἐλογίσατο ἐμβαλεῖν αὐτὸν εἰς χεῖρας τῶν ἀλλοφύλων.
\vs{26}Καὶ ἀπαγγέλλουσιν οἱ παῖδες Σαοὺλ τῷ Δαυὶδ τὰ ῥήματα ταῦτα, καὶ ηὐθύνθη ὁ λόγος ἐν ὀφθαλμοῖς Δαυὶδ ἐπιγαμβρεῦσαι τῷ βασιλεῖ.

\vs{27}Καὶ ἀνέστη Δαυὶδ, καὶ ἐπορεύθη αὐτὸς καὶ οἱ ἄνδρες αὐτοῦ, καὶ ἐπάταξεν ἐν τοῖς ἀλλοφύλοις ἑκατὸν ἄνδρας· καὶ ἀνήνεγκε τὰς ἀκροβυστίας αὐτῶν, καὶ ἐπιγαμβρεύεται τῷ βασιλεῖ, καὶ δίδωσιν αὐτῷ τὴν Μελχὸλ θυγατέρα αὐτοῦ αὐτῷ εἰς γυναῖκα.
\vs{28}Καὶ εἶδε Σαοὺλ ὅτι Κύριος μετὰ Δαυὶδ, καὶ πᾶς Ἰσραὴλ ἠγάπα αὐτόν.
\vs{29}Καὶ προσέθετο εὐλαβεῖσθαι ἀπὸ Δαυὶδ ἔτι.

\ch{19}
Καὶ ἐλάλησε Σαοὺλ πρὸς Ἰωνάθαν τὸν υἱὸν αὐτοῦ, καὶ πρὸς πὰντας τοὺς παῖδας αὐτοῦ θανατῶσαι τὸν Δαυίδ.
\vs{2}Καὶ Ἰωνάθαν ὁ υἱὸς Σαοὺλ ἡρεῖτο τὸν Δαυὶδ σφόδρα· καὶ ἀπήγγειλεν Ἰωνάθαν τῷ Δαυὶδ, λέγων, Σαοὺλ ζητεῖ θανατῶσαί σε· φύλαξαι οὖν αὔριον πρωῒ, καὶ κρύβηθι καὶ κάθισον κρυδῇ·
\vs{3}Καὶ ἐγὼ ἐξελεύσομαι καὶ στήσομαι ἐχόμενος τοῦ πατρός μου ἐν ἀγρῷ οὗ ἐὰν ᾖς ἐκεῖ· καὶ ἐγὼ λαλήσω περὶ σοῦ πρὸς τὸν πατέρα μου, καὶ ὄψομαι ὅ, τι ἐὰν ᾖ, καὶ ἀπαγγελῶ σοι.

\vs{4}Καὶ ἐλάλησεν Ἰωνάθαν περὶ Δαυὶδ ἀγαθὰ πρὸς Σαοὺλ τὸν πατέρα αὐτοῦ, καὶ εἶπε πρὸς αὐτὸν, μὴ ἁμαρτησάτω ὁ βασιλεὺς εἰς τὸν δοῦλόν σου Δαυὶδ, ὅτι οὐχ ἡμάρτηκεν εἰς σὲ, καὶ τὰ ποιήματα αὐτοῦ ἀγαθὰ σφόδρα·
\vs{5}Καὶ ἔθετο τὴν ψυχὴν αὐτοῦ ἐν τῇ χειρὶ αὐτοῦ, καὶ ἐπάταξε τὸν ἀλλόφυλον, καὶ ἐποίησε Κύριος σωτηρίαν μεγάλην, καὶ πᾶς Ἰσραὴλ εἶδον, καὶ ἐχάρησαν· καὶ ἱνατί ἁμαρτάνεις εἰς αἷμα ἀθῶον θανατῶσαι τὸν Δαυὶδ δωρεάν;
\vs{6}Καὶ ἤκουσε Σαοὺλ τῆς φωνῆς Ἰωνάθαν· καὶ ὤμοσε Σαοὺλ, λέγων, ζῇ Κύριος, εἰ ἀποθανεῖται.
\vs{7}Καὶ ἐκάλεσεν Ἰωνάθαν τὸν Δαυὶδ, καὶ ἀπήγγειλεν αὐτῷ πάντα τὰ ῥήματα ταῦτα· καὶ εἰσήγαγεν Ἰωνάθαν τὸν Δαυὶδ πρὸς Σαοὺλ, καὶ ἦν ἐνώπιον αὐτοῦ, ὡς ἐχθὲς καὶ τρίτην ἡμέραν.

\vs{8}Καὶ προσέθετο ὁ πόλεμος γενέσθαι πρὸς Σαοὺλ, καὶ κατίσχυσε Δαυὶδ, καὶ ἐπολέμησε τοὺς ἀλλοφύλους, καὶ ἐπάταξεν ἐν αὐτοῖς πληγὴν μεγάλην σφόδρα, καὶ ἔφυγον ἐκ προσώπου αὐτοῦ.

\vs{9}Καὶ ἐγένετο πνεῦμα Θεοῦ πονηρὸν ἐπὶ Σαοὺλ, καὶ αὐτὸς ἐν οἴκῳ καθεύδων, καὶ δόρυ ἐν τῇ χειρὶ αὐτοῦ, καὶ Δαυὶδ ἔψαλλε ταῖς χερσὶν αὐτοῦ.
\vs{10}Καὶ ἐζήτει Σαοὺλ πατάξαι τὸ δόρυ εἰς Δαυίδ· καὶ ἀπέστη Δαυὶδ ἐκ προσώπου Σαούλ· καὶ ἐπάταξε τὸ δόρυ εἰς τὸν τοῖχον· καὶ Δαυὶδ ἀνεχώρησε καὶ διεσώθη.
\vs{11}Καὶ ἐγενήθη ἐν τῇ νυκτὶ ἐκείνῃ, καὶ ἀπέστειλε Σαοὺλ ἀγγελους εἰς οἶκον Δαυὶδ φυλάξαι αὐτὸν, τοῦ θανατῶσαι αὐτὸν πρωΐ· καὶ ἀπήγγειλε τῷ Δαυὶδ Μελχὸλ ἡ γυνὴ αὐτοῦ, λέγουσα, ἐὰν μὴ σὺ σώσῃς τὴν ψυχὴν σαυτοῦ τὴν νύκτα ταύτην, αὔριον θανατωθήσῃ.
\vs{12}Καὶ κατάγει ἡ Μελχὸλ τὸν Δαυὶδ διὰ τῆς θυρίδος, καὶ ἀπῆλθε καὶ ἔφυγε καὶ σώζεται.
\vs{13}Καὶ ἔλαβεν ἡ Μελχὸλ τὰ κενοτάφια, καὶ ἔθετο ἐπὶ τὴν κλίνην, καὶ ἧπαρ τῶν αἰγῶν ἔθετο πρὸς κεφαλῆς αὐτοῦ, καὶ ἐκάλυψεν αὐτὰ ἱματίῳ.

\vs{14}Καὶ ἀπέστειλε Σαοὺλ ἀγγέλους λαβεῖν τὸν Δαυὶδ, καὶ λέγουσιν ἐνοχλεῖσθαι αὐτόν.
\vs{15}Καὶ ἀποστέλλει ἐπὶ τὸν Δαυὶδ, λέγων, ἀγάγετε αὐτὸν ἐπὶ τῆς κλίνης πρὸς μὲ τοῦ θανατῶσαι αὐτόν.
\vs{16}Καὶ ἔρχονται οἱ ἄγγελοι, καὶ ἰδοὺ τὰ κενοτάφια ἐπὶ τῆς κλίνης, καὶ ἧπαρ τῶν αἰγῶν πρὸς κεφαλῆς αὐτοῦ.
\vs{17}Καὶ εἶπε Σαοὺλ τῇ Μελχὸλ, ἱνατί οὕτως παρελογίσω με, καὶ ἐξαπέστειλας τὸν ἐχθρόν μου, καὶ διεσώθη; καὶ εἶπε Μελχὸλ τῷ Σαοὺλ, αὐτὸς εἶπεν, ἐξαπόστειλόν με, εἰ δὲ μὴ, θανατώσω σε.

\vs{18}Καὶ Δαυὶδ ἔφυγε καὶ διεσώθη, καὶ παραγίνεται πρὸς Σαμουὴλ εἰς Ἀρμαθαὶμ, καὶ ἀπαγγέλλει αὐτῷ πάντα ὅσα ἐποίησεν αὐτῷ Σαούλ· καὶ ἐπορεύθη Σαμουὴλ καὶ Δαυὶδ, καὶ ἐκάθισαν ἐν Ναυὰθ ἐν Ῥαμᾶ.

\vs{19}Καὶ ἀπηγγέλη τῷ Σαοὺλ, λέγοντες, ἰδοὺ Δαυὶδ ἐν Ναυὰθ ἐν Ῥαμᾴ.
\vs{20}Καὶ ἀπέστειλε Σαοὺλ ἀγγέλους λαβεῖν τὸν Δαυὶδ, καὶ εἶδον τὴν ἐκκλησίαν τῶν προφητῶν, καὶ Σαμουὴλ εἱστήκει καθεστηκὼς ἐπʼ αὐτῶν· καὶ ἐγενήθη ἐπὶ τούς ἀγγέλους τοῦ Σαοὺλ πνεῦμα Θεοῦ, καὶ προφητεύουσι.
\vs{21}Καὶ ἀπηγγέλη τῷ Σαοὺλ, καὶ ἀπέστειλεν ἀγγέλους ἑτέρους, καὶ ἐπροφήτευσαν καὶ αὐτοί· καὶ προσέθετο Σαοὺλ ἀποστεῖλαι ἀγγέλους τρίτους, καὶ ἐπροφήτευσαν καὶ αὐτοί.
\vs{22}Καὶ ἐθυμώθη ὀργῇ Σαοὺλ, καὶ ἐπορεύθη καὶ αὐτὸς εἰς Ἁρμαθαὶμ, καὶ ἔρχεται ἕως τοῦ φρέατος τοῦ ἅλω τοῦ ἐν τῷ Σεφί, καὶ ἠρώτησε καὶ εἶπε, ποῦ Σαμουὴλ καὶ Δαυίδ; καὶ εἶπαν, ἰδοὺ ἐν Ναυὰθ ἐν Ῥαμᾷ.
\vs{23}Καὶ ἐπορεύθη ἐκεῖθεν εἰς Ναυὰθ ἐν Ῥαμᾷ· καὶ ἐγενήθη καὶ ἐπʼ αὐτῷ πνεῦμα Θεοῦ, καὶ ἐπορεύετο προφητεύων ἕως τοῦ ἐλθεῖν αὐτὸν εἰς Ναυὰθ ἐν Ῥαμᾷ.
\vs{24}Καὶ ἐξεδύσατο τὰ ἱμάτια αὐτοῦ, καὶ ἐπροφήτευσεν ἐνώπιον αὐτῶν· καὶ ἔπεσε γυμνὸς ὅλην τὴν ἡμέραν ἐκείνην καὶ ὅλην τὴν νύκτα· διὰ τοῦτο ἔλεγον, εἰ καὶ Σαοὺλ ἐν προφήταις;

\ch{20}
Καὶ ἀπέδρα Δαυὶδ ἐκ Ναυὰθ ἐν Ῥαμᾷ, καὶ ἔρχεται ἐνώπιον Ἰωνάθαν, καὶ εἶπε, τί πεποίηκα, καὶ τί τὸ ἀδίκημά μου, καὶ τί ἡμάρτηκα ἐνώπιον τοῦ πατρός σου, ὅτι ἐπιζητεῖ τὴν ψυχήν μου;
\vs{2}Καὶ εἶπεν αὐτῷ Ἰωνάθαν, μηδαμῶς σοι, οὐ μὴ ἀποθάνῃς· ἰδοὺ οὐ μὴ ποιήσῃ ὁ πατήρ μου ῥῆμα μέγα ἢ μικρὸν, καὶ οὐκ ἀποκαλύψει τὸ ὠτίον μου· καὶ τί ὅτι κρύψει ὁ πατήρ μου ἀπʼ ἐμοῦ τὸ ῥῆμα τοῦτο; οὐκ ἔστι τοῦτο.
\vs{3}Καὶ ἀπεκρίθη Δαυὶδ τῷ Ἰωνάθαν, καὶ εἶπε, γινώσκων οἶδεν ὁ πατήρ σου, ὅτι, εὕρηκα χάριν ἐν ὀφθαλμοῖς σου, καὶ εἶπε, μὴ γνῶναι τοῦτο Ἰωνάθαν, μὴ οὐ βούληται· ἀλλὰ ζῇ Κύριος καὶ ζῇ ἡ ψυχή σου, ὅτι καθὼς εἶπον, ἐμπέπλησται ἀναμέσον ἐμοῦ καὶ τοῦ θανάτου.
\vs{4}Καὶ εἶπεν Ἰωνάθαν πρὸς Δαυὶδ, τί ἐπιθυμεῖ ἡ ψυχή σου, καὶ τί ποιήσω σοι;

\vs{5}Καὶ εἶπε Δαυὶδ πρὸς Ἰωνάθαν, ἰδοὺ δὴ νεομηνία αὔριον, καὶ ἐγὼ καθίσας οὐ καθήσομαι φαγεῖν, καὶ ἐξαποστελεῖς με, καὶ κρυβήσομαι ἐν τῷ πεδίῳ ἕως δείλης.
\vs{6}Καὶ ἐὰν ἐπισκεπτόμενος ἐπισκέψηταί με ὁ πατήρ σου, καὶ ἐρεῖς, παραιτούμενος παρῃτήσατο ἀπʼ ἐμοῦ Δαυὶδ δραμεῖν ἕως εἰς Βηθλεὲμ τὴν πόλιν αὐτοῦ, ὅτι θυσία τῶν ἡμερῶν ἐκεῖ ὅλῃ τῇ φυλῇ.
\vs{7}Ἐὰν τάδε εἴπῃ, ἀγαθῶς, εἰρήνη τῷ δούλῳ σου· καὶ ἐὰν σκληρῶς ἀποκριθῇ σοι, γνῶθι ὅτι συντετέλεσται ἡ κακία παρʼ αὐτοῦ.
\vs{8}Καὶ ποιήσεις ἔλεος μετὰ δούλου σου, ὅτι εἰσήγαγες εἰς διαθήκην Κυρίου τὸν δοῦλόν σου μετὰ σεαυτοῦ· καὶ εἰ ἔστιν ἀδικία ἐν τῷ δούλῳ σου, θανάτωσόν με σὺ, καὶ ἕως τοῦ πατρός σου ἱνατί οὕτως εἰσάγεις με;

\vs{9}Καὶ εἶπεν Ἰωνάθαν, μηδαμῶς σοι· ὅτι ἐὰν γινῶσκων γνῶ ὅτι συντετέλεσται ἡ κακία παρὰ τοῦ πατρός μου τοῦ ἐλθεῖν ἐπὶ σὲ, καὶ ἐὰν μὴ ᾖ εἰς τὰς πόλεις σου, ἐγὼ ἀπαγγελῶ σοι.
\vs{10}Καὶ εἶπε Δαυὶδ πρὸς Ἰωνάθαν, τίς ἀπαγγείλῃ μοι, ἐὰν ἀποκριθῇ ὁ πατήρ σου σκληρῶς;
\vs{11}Καὶ εἶπεν Ἰωνάθαν πρὸς Δαυὶδ, πορεύου, καὶ μὲνε εἰς ἀγρόν· καὶ ἐκπορεύονται ἀμφότεροι εἰς ἀγρόν.

\vs{12}Καὶ εἶπεν Ἰωνάθαν πρὸς Δαυὶδ, Κύριος ὁ Θεὸς Ἰσραὴλ οἶδεν, ὅτι ἀνακρινῶ τὸν πάτερα μου ὡς ἂν ὁ καιρὸς, τρισσῶς, καὶ ἰδοὺ ἀγαθὸν ᾖ περὶ Δαυὶδ, καὶ οὐ μὴ ἀποστείλω πρὸς σὲ εἰς ἀγρὸν,
\vs{13}τάδε ποιήσαι ὁ Θεὸς τῷ Ἰωνάθαν καὶ τάδε προσθείη· ὅτι ἀνοίσω τὰ κακὰ ἐπὶ σὲ, καὶ ἀποκαλύψω τὸ ὠτίον σου, καὶ ἐξαποστελῶ σε καὶ ἀπελεύσῃ εἰς εἰρήνην, καὶ ἔσται Κύριος μετὰ σοῦ καθὼς ἦν μετὰ τοῦ πατρός μου.
\vs{14}Καὶ ἐὰν μὲν ἔτι μου ζῶντος, καὶ ποιήσεις ἔλεος μετʼ ἐμοῦ· καὶ ἐὰν θανάτῳ ἀποθάνω,
\vs{15}οὐκ ἐξαρεῖς ἔλεός σου ἀπὸ τοῦ οἴκου μου ἕως τοῦ αἰῶνος· καὶ εἰ μὴ, ἐν τῷ ἐξαίρειν Κύριον τοὺς ἐχθροὺς Δαυὶδ ἔκαστον ἀπὸ τοῦ προσώπου τῆς γῆς,
\vs{16}εὑρεθῆναι τὸ ὄνομα τοῦ Ἰωνάθαν ἀπὸ τοῦ οἴκου Δαυὶδ, καὶ ἐκζητῆσαι Κύριος ἐχθροὺς τοῦ Δαυίδ.
\vs{17}Καὶ προσέθετο ἔτι Ἰωνάθαν ὀμόσαι τῷ Δαυίδ, ὅτι ἠγάπησε ψυχὴν ἀγαπῶντος αὐτόν.

\vs{18}Καὶ εἶπεν Ἰωνάθαν, αὔριον νεομηνία, καὶ ἐπισκεπήσῃ, ὅτι ἐπισκεπήσεται καθέδρα σου.
\vs{19}Καὶ τρισσεύσεις καὶ ἐπισκεψῃ καὶ ἥξεις εἰς τὸν τόπον σου οὗ κρυβῇς ἐν τῇ ἡμέρᾳ τῇ ἐργασίμῃ, καὶ καθήσῃ παρὰ τὸ ἐργὰβ ἐκεῖνο.
\vs{20}Καὶ ἐγὼ τρισσεύσω ταῖς σχίζαις ἀκοντίζων, ἐκπέμπων εἰς τὴν ʼΑματταρί.
\vs{21}Καὶ ἰδοὺ ἀποστέλλω τὸ παιδάριον, λέγων, δεῦρο, εὗρέ μοι τὴν σχίζαν. Ἐὰν εἴπω λέγων τῷ παιδαρίῳ, ὧδε ἡ σχίζα ἀπὸ σοῦ καὶ ὧδε, λάβε αὐτήν· παραγίνου, ὅτι εἰρήνη σοι, καὶ οὐκ ἔστι λόγος, ζῇ Κύριος·
\vs{22}ἐὰν τάδε εἴπω τῷ νεανίσκῳ, ὧδε ἡ σχίζα ἀπὸ σοῦ καὶ ἐπέκεινα· πορεύου, ὅτι ἐξαπέσταλκέ σε Κύριος.
\vs{23}Καὶ τὸ ῥῆμα ὃ ἐλαλήσαμεν ἐγὼ καὶ σὺ, ἰδοὺ Κύριος μάρτυς ἀναμέσον ἐμοῦ καὶ σοῦ ἕως αἰῶνος.

\vs{24}Καὶ κρύπτεται Δαυὶδ ἐν ἀγρῷ, καὶ παραγίνεται ὁ μὴν, καὶ ἔρχεται ὁ βασιλεὺς ἐπὶ τὴν τράπεζαν τοῦ φαγεῖν.
\vs{25}Καὶ ἐκάθισεν ἐπὶ τὴν καθέδραν αὐτοῦ ὡς ἅπαξ καὶ ἅπαξ ἐπὶ τῆς καθέδρας παρὰ τοῖχον, καὶ προέφθασε τὸν Ἰωνάθαν, καὶ ἐκάθισεν Ἀβεννὴρ ἐκ πλαγίων Σαοὺλ, καὶ ἐπεσκέπη ὁ τόπος Δαυίδ.
\vs{26}Καὶ οὐκ ἐλάλησε Σαοὺλ ἐν τῇ ἡμέρᾳ ἐκείνῃ, ὅτι εἴρηκε, σύμπτωμα φαίνεται, μὴ καθαρὸς εἶναι, ὅτι οὐ κεκαθάρισται.

\vs{27}Καὶ ἐγενήθη τῇ ἐπαύριον τοῦ μηνὸς τῇ ἡμέρᾳ τῇ δευτέρᾳ, καὶ ἐπεσκέπη ὁ τόπος τοῦ Δαυίδ· καὶ εἶπε Σαοὺλ πρὸς Ἰωνάθαν τὸν υἱὸν αὐτοῦ, τί ὅτι οὐ παραγέγονεν ὁ υἱὸς Ἰεσσαὶ καὶ ἐχθὲς καὶ σήμερον ἐπὶ τὴν τράπεζαν;
\vs{28}Καὶ ἀπεκρίθη Ἰωνάθαν τῷ Σαοὺλ, καὶ εἶπεν αὐτῷ, παρῄτηται παρʼ ἐμοῦ Δαυὶδ ἕως εἰς Βηθλεὲμ τὴν πόλιν αὐτοῦ πορευθῆναι.
\vs{29}Καὶ εἶπεν, ἐξαπόστειλον δή με, ὅτι θυσία τῆς φυλῆς ἡμῖν ἐν τῇ πόλει, καὶ ἐνετείλαντο πρὸς μὲ οἱ ἀδελφοί μου· καὶ νῦν εἰ εὕρηκα χάριν ἐν ὀφθαλμοῖς σου, διαβήσομαι δὴ καὶ ὄψομαι τοὺς ἀδελφούς μου· διὰ τοῦτο οὐ παραγέγονεν ἐπὶ τὴν τράπεζαν τοῦ βασιλέως.

\vs{30}Καὶ ἐθυμώθη ὀργῇ Σαοὺλ ἐπὶ Ἰωνάθαν σφόδρα, καὶ εἶπεν αὐτῷ, υἱὲ κορασίων αὐτομολούντων, οὐ γὰρ οἶδα ὅτι μέτοχος εἶ σὺ τῷ υἱῷ Ἰεσσαὶ εἰς αἰσχύνην σου, καὶ εἰς αἰσχύνην ἀποκαλύψεως μητρός σου;
\vs{31}ὅτι πάσας τὰς ἡμέρας ἃς ὁ υἱὸς Ἰεσσαὶ ζῇ ἐπὶ τῆς γῆς, οὐχ ἑτοιμασθήσεται ἡ βασιλεία σου· νῦν οὖν ἀποστείλας λάβε τὸν νεανίαν, ὅτι υἱὸς θανάτου οὗτος.
\vs{32}Καὶ ἀπεκρίθη Ἰωνάθαν τῷ Σαοὺλ, ἵνατί ἀποθνήσκει; τί πεποίηκε;
\vs{33}Καὶ ἐπῇρε Σαοὺλ τὸ δόρυ ἐπὶ Ἰωνάθαν τοῦ θανατῶσαι αὐτόν· καὶ ἔγνω Ἰωνάθαν ὅτι συντετέλεσται ἡ κακία αὕτη παρὰ τοῦ πατρὸς αὐτοῦ θανατῶσαι τὸν Δαυίδ.
\vs{34}Καὶ ἀνεπήδησεν Ἰωνάθαν ἀπὸ τῆς τραπέζης ἐν ὀργῇ θυμοῦ, καὶ οὐκ ἔφαγεν ἐν τῇ δευτέρᾳ τοῦ μηνὸς ἄρτον, ὅτι ἐθρανσθη ἐπὶ τὸν Δαυὶδ, ὅτι συνετέλεσεν ἐπʼ αὐτὸν ὁ πατὴρ αὐτοῦ.

\vs{35}Καὶ ἐγενήθη πρωὶ, καὶ ἐξῆλθεν Ἰωνάθαν εἰς ἀγρόν, καθὼς ἐτάξατο εἰς τὸ μαρτύριον Δαυίδ, καὶ παιδάριον μικρὸν μετʼ αὐτοῦ.
\vs{36}Καὶ εἶπε τῷ παιδαρίω, δράυε, εὗρέ μοι τὰς σχίζας ἐν αἷς ἐγὼ ἀκοντίζω· και τὸ παιδάμε ἔδραυε, και αὐτὸς ἠκόντισε τῇ σχίζῃ, καὶ παρήγαγεν αὐτήν.
\vs{37}Καὶ ἦλθεν τὸ παιδάριον ἕως τοῦ τόπου τῆς σχίζης οὗ ἠκόντιζεν Ἰωνάθαν· καὶ ἀνεβόησεν Ἰωνάθαν ὀπίσω τοῦ νεανίου, καὶ εἶπεν, ἐκεῖ ἡ σχίζα ἀπὸ σοῦ καὶ ἐπέκεινα·
\vs{38}Καὶ ἀνεβόησεν Ἰωνάθαν ὀπίσω τοῦ παιδαρίου αὐτοῦ, λέγων, ταχύνας σπεῦσον, καὶ μὴ στῇς· καὶ ἀνέλεξε τὸ παιδάριον Ἰωνάθαν τὰς σχίζας, καὶ ἤνεγκε τὰς σχίζας πρὸς τὸν κύριον αὐτοῦ.
\vs{39}καὶ τὸ παιδάριον οὐκ ἔγνω οὐθὲν, πάρεξ Ἰωνάθαν καὶ Δαυίδ.
\vs{40}καὶ Ἰωνάθαν ἔδωκε τὰ σκεύη αὐτοῦ ἐπὶ τὸ παιδάριον αὐτοῦ, καὶ εἶπε τῷ παιδαρίῳ αὐτοῦ, πορεύου, εἴσελθε εἰς τὴν πόλιν.

\vs{41}Καὶ ὡς εἰσῆλθε τὸ παιδάριον, καὶ Δαυὶδ ἀνέστη ἀπὸ τοῦ ἀργὰβ, καὶ ἔπεσεν ἐπὶ πρόσωπον αὐτοῦ, καὶ προσεκύνησεν αὐτῷ τρίς, καὶ κατεφίλησεν ἕκαστος τὸν πλησίον αὐτοῦ, καὶ ἔκλαυσεν ἕκαστος τῷ πλησίον αὐτοῦ, ἕως συντελείας μεγάλης.
\vs{42}Καὶ εἶπεν Ἰωνάθαν τῷ Δαυὶδ, πορεύου εἰς εἰρήνην, καὶ ὡς ὀμωμόκαμεν ἡμεῖς ἀμφότεροι ἐν ὀνόματι Κυρίου, λέγοντες, Κύριος ἔσται μάρτυς ἀνὰμέσον ἐμοῦ καὶ σοῦ, καὶ ἀνὰμέσον τοῦ σηέρματός μου καὶ ἀναμέσον τοῦ σπέρματός σον ἕως αἰῶνος·

\ch{21}καὶ ἀνέστη Δαυὶδ καὶ ἀπῆλθε, καὶ Ἰωνάθαν εἰσῆλθεν εἰς τὴν πόλιν.

\vs{2}Καὶ ἔρχεται Δαυὶδ εἰς Νομβᾶ πρὸς Ἀβιμέλεχ τὸν ἱερέα· καὶ ἐξέστη Ἀβιμέλεχ τῇ ἀπαντήσει αὐτοῦ, καὶ εἶπεν αὐτῷ, τί ὅτι σὺ μόνος, καὶ οὐθεὶς μετὰ σοῦ;
\vs{3}Καὶ εἶπε Δαυὶδ τῷ ἱερεῖ, ὁ βασιλεὺς ἐντέταλταί μοι ῥῆμα σήμερον, καὶ εἶπέ μοι, μηδεὶς γνώτω τὸ ῥῆμα περὶ οὗ ἐγὼ ἀποστέλλω σε, καὶ ὑπὲρ οὗ ἐγὼ ἐντέταλμαί σοι· καὶ τοῖς παιδαρίοις διαμεμαρτύρημαι ἐν τῷ τόπῳ τῷ λεγομένῳ, Θεοῦ πίστις φελλανὶ μαεμωνί.
\vs{4}Καὶ νῦν εἰ εἰσὶν ὑπὸ τὴν χεῖρά σου πέντε ἄρτοι, δὸς εἰς χεῖρά μου τὸ εὑρεθέν.
\vs{5}Καὶ ἀπεκρίθη ὁ ἱερεὺς τῷ Δαυὶδ, καὶ εἶπεν, οὔκ εἰσιν ἄρτοι βέβηλοι ὑπὸ τὴν χεῖρά μου, ὅτι ἀλλʼ ἢ ἄρτοι ἅγιοί εἰσιν· εἰ πεφυλαγμένα τὰ παιδάριά ἐστι πλὴν ἀπὸ γυναικὸς, καὶ φάγεται.
\vs{6}Καὶ ἀπεκρίθη Δαυὶδ τῷ ἱερεῖ, καὶ εἶπεν αὐτῷ, ἀλλὰ ἀπὸ γυναικὸς ἀπεσχήμεθα ἐχθὲς καὶ τρίτην ἡμέραν· ἐν τῷ ἐξελθεῖν με εἰς ὁδὸν γέγονε πάντα τὰ παιδία ἡγνισμένα, καὶ αὐτὴ ἡ ὁδὸς βέβηλος, διότι ἁγιασθήσεται σήμερον διὰ τὰ σκεύη μου.

\vs{7}Καὶ ἔδωκεν αὐτῷ Ἀβιμέλεχ ὁ ἱερεὺς τοὺς ἄρτους τῆς προθέσεως, ὅτι ἐκεῖ οὐκ ἦν ἄρτοι, ἀλλʼ ἢ ἄρτοι τοῦ προσώπου οἱ ἀφῃρημένοι ἐκ προσώπου Κυρίου, τοῦ παρατεθῆναι ἄρτον θερμὸν ᾗ ἡμέρᾳ ἔλαβεν αὐτούς.

\vs{8}Καὶ ἐκεῖ ἦν ἓν τῶν παιδαρίων τοῦ Σαοὺλ ἐν τῇ ἡμέρᾳ ἐκείνῃ συνεχόμενος Νεεσσαρὰν ἐνώπιον Κυρίου, καὶ ὄνομα αὐτῷ Δωὴκ Σύρος, νέμων τὰς ἡμιόνους Σαούλ.

\vs{9}Καὶ εἶπε Δαυὶδ πρὸς Ἀβιμέλεχ, ἴδε εἰ ἔστιν ἐνταῦθα ὑπὸ τὴν χεῖρά σου δόρυ ἢ ῥομφαία, ὅτι τὴν ῥομφαίαν μου καὶ τὰ σκεύη οὐκ εἴληφα ἐν τῇ χειρί μου, ὅτι ἦν τὸ ῥῆμα τοῦ βασιλέως κατὰ σπουδήν.
\vs{10}Καὶ εἶπεν ὁ ἱερεὺς, ἰδοὺ ἡ ῥομφαία Γολιὰθ τὸυ ἀλλοφύλου, ὃν ἐπάταξας ἐν τῇ κοιλάδι Ἠλᾶ· καὶ αὕτη ἐνειλημένη ἦν ἐν ἱματίῳ· εἰ ταύτην λήψῃ, σεαυτῷ λάβε, ὅτι οὐκ ἔστιν ἑτέρα πάρεξ ταύτης ἐνταῦθα· καὶ εἶπε Δαυὶδ, ἰδοὺ οὐκ ἔστιν ὥσπερ αὐτή· δός μοι αὐτήν.

\vs{11}Καὶ ἔδωκεν αὐτὴν αὐτῷ· καὶ ἀνέστη Δαυὶδ, καὶ ἔφυγεν ἐν τῇ ἡμέρᾳ ἐκείνῃ ἐκ προσώπου Σαούλ· καὶ ἦλθε Δαυὶδ πρὸς Ἀγχοῦς βασιλέα Γέθ.
\vs{12}Καὶ εἶπον οἱ παῖδες Ἀγχοῦς πρὸς αὐτὸν, οὐχὶ οὗτος Δαυὶδ ὁ βασιλεὺς τῆς γῆς; οὐχὶ τούτῳ ἐξῆρχον αἱ χορεύουσαι, λέγουσαι, ἐπάταξε Σαοὺλ ἐν χιλιάσιν αὐτοῦ, καὶ Δαυὶδ ἐν μυριάσιν αὐτοῦ;
\vs{13}Καὶ ἔθετο Δαυὶδ τὰ ῥήματα ἐν τῇ καρδίᾳ αὐτοῦ, καὶ ἐφοβήθη σφόδρα ἀπὸ προσώπου Ἀγχοῦς βασιλέως Γέθ.
\vs{14}Καὶ ἠλλοίωσε τὸ πρόσωπον αὐτοῦ ἐνώπιον αὐτοῦ, καὶ προσεποιήσατο ἐν τῇ ἡμέρᾳ ἐκείνῃ, καὶ ἐτυμπάνιζεν ἐπὶ ταῖς θύραις τῆς πόλεως, καὶ παρεφέρετο ἐν ταῖς χερσὶν αὐτοῦ, καὶ ἔπιπτεν ἐπὶ τὰς θύρας τῆς πύλης, καὶ τὰ σίελα αὐτοῦ κατέῤῥει ἐπὶ τὸν πώγωνα αὐτοῦ.
\vs{15}Καὶ εἶπεν Ἀγχοῦς πρὸς τοὺς παῖδας αὐτοῦ, ἰδοὺ ἴδετε ἄνδρα ἐπίληπτον, ἱνατί εἰσηγάγετε αὐτὸν πρὸς μέ;
\vs{16}Μὴ ἐλαττοῦμαι ἐπιλήπτων ἐγὼ, ὅτι εἰσαγηόχατε αὐτὸν ἐπιληπτεύεσθαι πρὸς μέ; οὗτος οὐκ εἰσελεύσεται εἰς οἰκίαν.

\ch{22}
Καὶ ἀπῆλθεν ἐκεῖθεν Δαυὶδ, καὶ διεσώθη, καὶ ἔρχεται εἰς τὸ σπήλαιον τὸ Ὀδολλάμ· καὶ ἀκούουσιν οἱ ἀδελφοὶ αὐτοῦ, καὶ ὁ οἶκος τοῦ πατρὸς αὐτοῦ, καὶ καταβαίνουσι πρὸς αὐτὸν ἐκεῖ.
\vs{2}Καὶ συνήγοντο πρὸς αὐτὸν πᾶς ἐν ἀνάγκῇ, καὶ πᾶς ὑπόχρεως, καὶ πᾶς κατώδυνος ψυχῇ, καὶ ἦν ἐπʼ αὐτῶν ἡγούμενος, καὶ ἦσαν μετʼ αὐτοῦ ὡς τετρακόσιοι ἄνδρες.

\vs{3}Καὶ ἀπῆλθε Δαυὶδ ἐκεῖθεν εἰς Μασσηφὰθ τῆς Μωὰβ, καὶ εἶπε πρὸς βασιλέα Μωὰβ, γινέσθωσαν δὴ ὁ πατήρ μου καὶ ἡ μήτηρ μου παρὰ σοί, ἕως ὅτου γνῶ τί ποιήσει μοι ὁ Θεός.
\vs{4}Καὶ παρεκάλεσε τὸ πρόσωπον τοῦ βασιλέως Μωὰβ, καὶ κατῴκουν μετʼ αὐτοῦ πάσας τὰς ἡμέρας, ὄντος τοῦ Δαυὶδ ἐν τῇ περιοχῇ.
\vs{5}Καὶ εἶπε Γὰδ ὁ προφήτης πρὸς Δαυὶδ, μὴ κάθου ἐν τῇ περιοχῇ· πορεύου, καὶ ἥξεις εἰς γῆν Ἰούδα· καὶ ἐπορεύθη Δαυὶδ, καὶ ἦλθε καὶ ἐκάθισεν ἐν πόλει Σαρίκ.

\vs{6}Καὶ ἤκουσε Σαοὺλ, ὅτι ἔγνωσται Δαυὶδ, καὶ οἱ ἄνδρες οἱ μετʼ αὐτοῦ· καὶ Σαοὺλ ἐκάθητο ἐν τῷ βουνῷ ὑπὸ τὴν ἄρουραν τὴν ἐν Ῥαμᾷ, καὶ τὸ δόρυ ἐν τῇ χειρὶ αὐτοῦ, καὶ πάντες οἱ παῖδες αὐτοῦ παρειστήκεισαν αὐτῷ.
\vs{7}Καὶ εἶπε Σαοὺλ πρὸς τοὺς παῖδας αὐτοῦ τοὺς παρεστηκότας αὐτῷ, ἀκούσατε δὴ υἱοὶ Βενιαμὶν, εἰ ἀληθῶς πᾶσιν ὑμῖν δώσει ὁ υἱὸς Ἰεσσαὶ ἀγροὺς καὶ ἀμπελῶνας, καὶ πάντας ὑμᾶς τάξει ἑκατοντάρχους καὶ χιλιάρχους,
\vs{8}ὅτι σύγκεισθε πάντες ὑμεῖς ἐπʼ ἐμὲ, καὶ οὐκ ἔστιν ὁ ἀποκαλύπτων τὸ ὠτίον μου, ἐν τῷ διαθέσθαι τὸν υἱόν μου διαθήκην μετὰ τοῦ υἱοῦ Ἰεσσαὶ, καὶ οὐκ ἔστι πονῶν περὶ ἐμοῦ ἐξ ὑμῶν, καὶ ἀποκαλύπτων τὸ ὠτίον μου, ὅτι ἐπήγειρεν ὁ υἱός μου τὸν δοῦλόν μου ἐπʼ ἐμὲ εἰς ἐχθρὸν, ὡς ἡ ἡμέρα αὕτη;

\vs{9}Καὶ ἀποκρίνεται Δωὴκ ὁ Σύρος ὁ καθεστηκὼς ἐπὶ τάς ἡμιόνους Σαοὺλ, καὶ εἶπεν, ἑωρακα τὸν υἱὸν Ἰεσσαὶ παραγινόμενον εἰς Νομβᾶ πρὸς Ἀβιμέλεχ υἱὸν Ἀχιτὼβ τὸν ἱερέα.
\vs{10}Καὶ ἠρώτα αὐτῷ διὰ τοῦ Θεοῦ, καὶ ἐπισιτισμὸν ἔδωκεν αὐτῷ, καὶ τὴν ῥομφαίαν Γολιὰθ τοῦ ἀλλοφύλου ἔδωκεν αὐτῷ.

\vs{11}Καὶ ἀπέστειλεν ὁ βασιλεὺς καλέσαι τὸν Ἀβιμέλεχ υἱὸν Ἀχιτὼβ καὶ πάντας τοὺς υἱοὺς τοῦ πατρὸς αὐτοῦ τοὺς ἱερεῖς τοὺς ἐν Νομβᾶ· καὶ παρεγένοντο πάντες πρὸς τὸν βασιλέα.
\vs{12}Καὶ εἶπε Σαοὺλ, ἄκουε δὴ υἱὲ Ἀχιτώβ· καὶ εἶπεν, ἰδοὺ ἐγὼ, λάλει κύριε.
\vs{13}Καὶ εἶπεν αὐτῷ Σαοὺλ, ἱνατί συνέθου κατʼ ἐμοῦ σὺ καὶ ὁ υἱὸς Ἰεσσαὶ, δοῦναί σε αὐτῷ ἄρτον καὶ ῥομφαίαν, καὶ ἐρωτᾷν αὐτῷ διὰ τοῦ Θεοῦ, θέσθαι αὐτὸν ἐπʼ ἐμὲ εἰς ἐχθρὸν, ὡς ἡ ἡμέρα αὕτη;
\vs{14}Καὶ ἀπεκρίθη τῷ βασιλεῖ, καὶ εἶπε, καὶ τίς ἐν πᾶσι τοῖς δούλοις σου ὡς Δαυὶδ πιστὸς, καὶ γαμβρὸς τοῦ βασιλέως, καὶ ἄρχων παντὸς παραγγέλματός σου, καὶ ἔνδοξος ἐν τῷ οἰκῳ σου;
\vs{15}Ἢ σήμερον ἦργμαι ἐρωτᾷν αὐτῷ διὰ τοῦ Θεοῦ; μηδαμῶς· μὴ δότω ὁ βασιλεὺς κατὰ τοῦ δούλου αὐτοῦ λόγον, καὶ ἐφʼ ὅλον τὸν οἶκον τοῦ πατρός μου, ὅτι οὐκ ᾔδει ὁ δοῦλός σου ἐν πᾶσι τούτοις ῥῆμα μικρὸν, ἢ μέγα.

\vs{16}Καὶ εἶπεν ὁ βασιλεὺς Σαοὺλ, θανάτῳ ἀποθανῇ Ἀβιμέλεχ σὺ, καὶ πᾶς ὁ οἶκος τοῦ πατρός σου.
\vs{17}Καὶ εἶπεν ὁ βασιλεὺς τοῖς παρατρέχουσι τοῖς ἐφεστηκόσι πρὸς αὐτὸν, προσαγάγετε καὶ θανατοῦτε τοὺς ἱερεῖς τοῦ Κυρίου, ὅτι ἡ χεὶρ αὐτῶν μετὰ Δαυὶδ, καὶ ὅτι ἔγνωσαν ὅτι φεύγει αὐτὸς, καὶ οὐκ ἀπεκάλυψαν τὸ ὠτίον μου· καὶ οὐκ ἐβουλήθησαν οἱ παῖδες τοῦ βασιλέως ἐπενεγκεῖν τὰς χεῖρας αὐτῶν ἁπαντῆσαι εἰς τοὺς ἱερεῖς Κυρίου.
\vs{18}Καὶ εἶπεν ὁ βασιλεὺς τῷ Δωὴκ, ἐπιστρέφου σὺ, καὶ ἀπάντα εἰς τοὺς ἱερεῖς· καὶ ἐπεστράφη Δωὴκ ὁ Σύρος, καὶ ἐθανάτωσε τοὺς ἱερεῖς τοῦ Κυρίου ἐν τῇ ἡμέρᾳ ἐκείνῃ, τριακοσίους καὶ πέντε ἄνδρας, πάντας αἴροντας ἐφούδ.
\vs{19}Καὶ τὴν Νομβᾶ τὴν πόλιν τῶν ἱερέων ἐπάταξεν ἐν στόματι ῥομφαίας ἀπὸ ἀνδρὸς ἕως γυναικὸς, ἀπὸ νηπίου ἕως θηλάζοντος, καὶ μόσχου, καὶ ὄνου, καὶ προβάτου.

\vs{20}Καὶ διασώζεται υἱὸς εἷς τῷ Ἀβιμέλεχ υἱῷ Ἀχιτὼβ, καὶ ὄνομα αὐτῷ Ἀβιάθαρ, καὶ ἔφυγεν ὀπίσω Δαυίδ.
\vs{21}Καὶ ἀπήγγειλεν Ἀβιάθαρ τῷ Δαυὶδ, ὅτι ἐθανάτωσε Σαοὺλ πάντας τοὺς ἱερεῖς τοῦ Κυρίου.
\vs{22}Καὶ εἶπε Δαυὶδ τῷ Ἀβιάθαρ, ἤδειν ὅτι ἐν τῇ ἡμέρᾳ ἐκείνῃ, ὅτι Δωὴκ ὁ Σύρος ὅτι ἀπαγγέλλων ἀπαγγελεῖ τῷ Σαούλ· ἐγώ εἰμι αἴτιος τῶν ψυχῶν οἴκου τοῦ πατρός σου.
\vs{23}Κάθου μετʼ ἐμοῦ· μὴ φοβοῦ, ὅτι οὗ ἐὰν ζητῶ τῇ ψυχῇ μου τόπον, ζητήσω καὶ τῇ ψυχῇ σου, ὅτι πεφύλαξαι σὺ παρʼ ἐμοί.

\ch{23}
Καὶ ἀπηγγέλη τῷ Δαυὶδ, λέγοντες, ἰδοὺ οἱ ἀλλόφυλοι πολεμοῦσιν ἐν τῇ Κεϊλὰ, καὶ αὐτοὶ διαρπάζουσι, καταπατοῦσι τοὺς ἅλω.
\vs{2}Καὶ ἐπηρώτησε Δαυὶδ διὰ τοῦ Κυρίου, λέγων, εἰ πορευθῶ, καὶ πατάξω τοὺς ἀλλοφύλους τούτους; καὶ εἶπε Κύριος, πορεύου, καὶ πατάξεις ἐν τοῖς ἀλλοφύλοις τούτοις, καὶ σώσεις τὴν Κεϊλά.
\vs{3}Καὶ εἶπον οἱ ἄνδρες τοῦ Δαυὶδ πρὸς αὐτὸν, ἰδοὺ ἡμεῖς ἐνταῦθα ἐν τῇ Ἰουδαίᾳ φοβούμεθα, καὶ πῶς ἔσται ἐὰν πορευθῶμεν εἰς Κεϊλὰ, εἰς τὰ σκῦλα τῶν ἀλλοφύλων εἰσπορευσόμεθα;
\vs{4}Καὶ προσέθετο Δαυὶδ ἔτι ἐπερωτῆσαι διὰ τοῦ Κυρίου· καὶ ἀπεκρίθη αὐτῷ Κύριος, καὶ εἶπεν αὐτῷ, ἀνάστηθι καὶ κατάβηθι εἰς Κεϊλὰ, ὅτι ἐγὼ παραδίδωμι τοὺς ἀλλοφύλους εἰς χεῖράς σου.
\vs{5}Καὶ ἐπορεύθη Δαυὶδ καὶ οἱ ἄνδρες οἱ μετʼ αὐτοῦ εἰς Κεϊλὰ, καὶ ἐπολέμησε τοῖς ἀλλοφύλοις· καὶ ἔφυγον ἐκ προσώπου αὐτοῦ, καὶ ἀπήγαγε τὰ κτήνη αὐτῶν, καὶ ἐπάταξεν ἐν αὐτοῖς πληγὴν μεγάλην, καὶ ἔσωσε Δαυὶδ τοὺς κατοικοῦντας Κεϊλά.
\vs{6}Καὶ ἐγένετο ἐν τῷ φεύγειν Ἀβιάθαρ υἱὸν Ἀχιμέλεχ πρὸς Δαυὶδ, καὶ αὐτὸς μετὰ Δαυὶδ εἰς Κεϊλὰ κατέβη ἔχων ἐφοὺδ ἐν τῇ χειρὶ αὐτοῦ.

\vs{7}Καὶ ἀπηγγέλη τῷ Σαοὺλ, ὅτι ἥκει ὁ Δαυὶδ εἰς Κεϊλά· καὶ εἶπε Σαοὺλ, πέπρακεν αὐτὸν ὁ Θεὸς εἰς τὰς χεῖράς μου, ὅτι ἀποκέκλεισται εἰσελθὼν εἰς πόλιν θυρῶν καὶ μοχλῶν.
\vs{8}Καὶ παρήγγειλε Σαοὺλ παντὶ τῷ λαῷ καταβαίνειν εἰς πόλεμον εἰς Κεϊλὰ, συνέχειν τὸν Δαυὶδ καὶ τοὺς ἄνδρας αὐτοῦ.
\vs{9}Καὶ ἔγνω Δαυὶδ, ὅτι οὐ παρασιωπᾷ Σαοὺλ περὶ αὐτοῦ τὴν κακίαν· καὶ εἶπε Δαυὶδ πρὸς Ἀβιάθαρ τὸν ἱερέα, προσάγαγε τὸ ἐφοὺδ Κυρίου.
\vs{10}Καὶ εἶπε Δαυίδ, Κύριε ὁ Θεὸς Ἰσραήλ, ἀκούων ἀκήκοεν ὁ δοῦλός σου, ὅτι ζητεῖ Σαοὺλ ἐλθεῖν ἐπὶ Κεϊλὰ διαφθεῖραι τὴν πόλιν διʼ ἐμέ.
\vs{11}Εἰ ἀποκλεισθήσεται; καὶ νῦν εἰ καταβήσεται Σαοὺλ, καθὼς ἤκουσεν ὁ δοῦλός σου; Κύριε ὁ Θεὸς Ἰσραήλ ἀπάγγειλον τῷ δούλῳ σου· καὶ εἶπε Κύριος, ἀποκλεισθήσεται.

\vs{13}Καὶ ἀνέστη Δαυὶδ καὶ οἱ ἄνδρες οἱ μετʼ αὐτοῦ ὡς τετρακόσιοι, καὶ ἐξῆλθον ἐκ Κεϊλὰ, καὶ ἐπορεύοντο οὗ ἐὰν ἐπορεύοντο· καὶ τῷ Σαοὺλ ἀπηγγέλη, ὅτι διασέσωσται Δαυὶδ ἐκ Κεϊλὰ, καὶ ἀνῆκε τοῦ ἐλθεῖν.
\vs{14}Καὶ ἐκάθισεν ἐν Μασερὲμ ἐν τῇ ἐρήμῳ ἐν τοῖς στενοῖς, καὶ ἐκάθετο ἐν τῇ ἐρήμῳ ἐν τῷ ὄρει Ζὶφ, ἐν τῇ γῇ τῇ αὐχμώδει· καὶ ἐζήτει αὐτὸν Σαοὺλ πάσας τὰς ἡμέρας, καὶ οὐ παρέδωκεν αὐτὸν Κύριος εἰς τὰς χεῖρας αὐτοῦ.
\vs{15}Καὶ εἶδε Δαυὶδ, ὅτι ἐξέρχεται Σαοὺλ τοῦ ζητεῖν τὸν Δαυίδ· καὶ Δαυὶδ ἦν ἐν τῷ ὄρει τῷ αὐχμώδει ἐν τῇ Καινῇ Ζίφ.

\vs{16}Καὶ ἀνέστη Ἰωνάθαν υἱὸς Σαοὺλ καὶ ἐπορεύθη πρὸς Δαυὶδ εἰς Καινὴν, καὶ ἐκραταίωσε τὰς χεῖρας αὐτοῦ ἐν Κυρίῳ,
\vs{17}καὶ εἶπε πρὸς αὐτὸν, μὴ φοβοῦ, ὅτι οὐ μὴ εὕρῃ σε ἡ χεὶρ Σαοὺλ τοῦ πατρός μου, καὶ σὺ βασιλεύσεις ἐπὶ Ἰσραὴλ, καὶ ἐγὼ ἔσομαί σοι εἰς δεύτερον, καὶ Σαοὺλ ὁ πατήρ μου οἶδεν οὕτως.
\vs{18}Καὶ διέθεντο ἀμφότεροι διαθήκην ἐνώπιον Κυρίου· καὶ ἐκάθητο Δαυὶδ ἐν Καινῇ, καὶ Ἰωνάθαν ἀπῆλθεν εἰς οἶκον αὐτοῦ.

\vs{19}Καὶ ἀνέβησαν οἱ Ζιφαῖοι ἐκ τῆς αὐχμώδους πρὸς Σαοὺλ ἐπὶ τὸν βουνὸν, λέγοντες, οὐκ ἰδοὺ Δαυὶδ κέκρυπται παρʼ ἡμῖν ἐν Μεσσαρὰ ἐν τοῖς στενοῖς ἐν τῇ Καινῇ ἐν τῷ βουνῷ τοῦ Ἑχελᾶ τοῦ ἐκ δεξιῶν τοῦ Ἰεσσαιμοῦ;
\vs{20}Καὶ νῦν πᾶν τὸ πρὸς ψυχὴν τοῦ βασιλέως εἰς κατάβασιν, καταβαινέτω πρὸς ἡμᾶς· κεκλείκασιν αὐτὸν εἰς τὰς χεῖρας τοῦ βασιλέως.
\vs{21}Καὶ εἶπεν αὐτοῖς Σαοὺλ, εὐλογημένοι ὑμεῖς τῷ Κυρίῳ, ὅτι ἐπονέσατε περὶ ἐμοῦ·
\vs{22}Πορεύθητε δὴ καὶ ἑτοιμάσατε ἔτι, καὶ γνῶτε τὸν τόπον αὐτοῦ οὗ ἔσται ὁ ποὺς αὐτοῦ ἐν τάχει ἐκεῖ οὗ εἴπατε, μή ποτε πανουργεύσηται.
\vs{23}Καὶ ἴδετε καὶ γνῶτε, καὶ πορεύσομαι μεθʼ ὑμῶν· καὶ ἔσται εἰ ἔστιν ἐπὶ τῆς γῆς, καὶ ἐξερευνήσω αὐτὸν ἐν πάσαις χιλιάσιν Ἰούδα.

\vs{24}Καὶ ἀνέστησαν οἱ Ζιφαῖοι, καὶ ἐπορεύθησαν ἔμπροσθεν Σαούλ· καὶ Δαυὶδ καὶ οἱ ἄνδρες αὐτοῦ ἐν τῇ ἐρήμῳ τῇ Μαὼν καθʼ ἑσπέραν ἐκ δεξιῶν τοῦ Ἰεσσαιμοῦ.

\vs{25}Καὶ ἐπορεύθη Σαοὺλ καὶ οἱ ἄνδρες αὐτοῦ ζητεῖν αὐτόν· καὶ ἀπήγγειλαν τῷ Δαυὶδ, καὶ κατέβη εἰς τὴν πέτραν τὴν ἐν τῇ ἐρήμῳ Μαών· καὶ ἤκουσε Σαοὺλ, καὶ κατεδίωξεν ὀπίσω Δαυὶδ εἰς τὴν ἔρημον Μαών.
\vs{26}Καὶ πορεύονται Σαοὺλ καὶ οἱ ἄνδρες αὐτοῦ ἐκ μέρους τοῦ ὄρους ἐκ τούτου, καὶ ἦν Δαυὶδ καὶ οἱ ἄνδρες αὐτοῦ ἐκ μέρους τοῦ ὄρους ἐκ τούτου· καὶ ἦν Δαυὶδ σκεπαζόμενος πορεύεσθαι ἀπὸ προσώπου Σαούλ· καὶ Σαοὺλ καὶ οἱ ἄνδρες αὐτοῦ παρενέβαλον ἐπὶ Δαυὶδ καὶ τοὺς ἄνδρας αὐτοῦ, συλλαβεῖν αὐτούς.

\vs{27}Καὶ πρὸς Σαοὺλ ἦλθεν ἄγγελος, λέγων, σπεῦδε καὶ δεῦρο, ὅτι ἀλλόφυλοι ἐπέθεντο ἐπὶ τὴν γῆν.
\vs{28}Καὶ ἀνέστρεψε Σαοὺλ μὴ καταδιώκειν ὀπίσω Δαυίδ, καὶ ἐπορεύθη εἰς συνάντησιν τῶν ἀλλοφύλων· διὰ τοῦτο ἐπεκλήθη ὁ τόπος ἐκεῖνος, πέτρα ἡ μερισθεῖσα.

\ch{24}
Καὶ ἀνέστη Δαυὶδ ἐκεῖθεν, καὶ ἐκάθισεν ἐν τοῖς στενοῖς Ἐνγαδδί.
\vs{2}Καὶ ἐγενήθη ὡς ἀνέστρεψε Σαοὺλ ἀπὸ ὄπισθεν τῶν ἀλλοφύλων, καὶ ἀπηγγέλη αὐτῷ, λεγόντων, ὅτι Δαυὶδ ἐν τῇ ἐρήμῳ Ἐνγαδδί.
\vs{3}Καὶ ἔλαβε μεθʼ ἑαυτοῦ τρεῖς χιλιάδας ἀνδρῶν ἐκλεκτοὺς ἐκ παντὸς Ἰσραὴλ, καὶ ἐπορεύθη ζητεῖν τὸν Δαυὶδ καὶ τοὺς ἄνδρας αὐτοῦ ἐπὶ πρόσωπον Σαδδαιέμ.
\vs{4}Καὶ ἦλθεν εἰς τὰς ἀγέλας τῶν ποιμνίων τὰς ἐπὶ τῆς ὁδοῦ, καὶ ἦν ἐκεῖ σπήλαιον· καὶ Σαοὺλ εἰσῆλθε παρασκευάσασθαι, καὶ Δαυὶδ καὶ οἱ ἄνδρες αὐτοῦ ἐσώτερον τοῦ σπηλαίου ἐκάθηντο.
\vs{5}Καὶ εἶπον οἱ ἄνδρες Δαυὶδ πρὸς αὐτὸν, ἰδοὺ ἡ ἡμέρα αὕτη, ἣν εἶπε Κύριος πρὸς σὲ παραδοῦναι τὸν ἐχθρόν σου εἰς τὰς χεῖράς σου, καὶ ποιήσεις αὐτῷ ὡς ἀγαθὸν ἐν ὀφθαλμοῖς σου· καὶ ἀνέστη Δαυὶδ, καὶ ἀφεῖλε τὸ πτερύγιον τῆς διπλοΐδος τοῦ Σαοὺλ λαθραίως.

\vs{6}Καὶ ἐγενήθη μετὰ ταῦτα, καὶ ἐπάταξε καρδία Δαυὶδ αὐτὸν, ὅτι ἀφεῖλε τὸ πτερύγιον τῆς διπλοΐδος αὐτοῦ.
\vs{7}Καὶ εἶπε Δαυὶδ πρὸς τοὺς ἄνδρας αὐτοῦ, μηδαμῶς μοι παρὰ Κυρίου, εἰ ποιήσω τὸ ῥῆμα τοῦτο τῷ κυρίῳ μου τῷ χριστῷ Κυρίου, ἐπενέγκαι χεῖρά μου ἐπʼ αὐτὸν, ὅτι χριστὸς Κυρίου ἐστὶν οὗτος.
\vs{8}Καὶ ἔπεισε Δαυὶδ τοὺς ἄνδρας αὐτοῦ ἐν λόγοις, καὶ οὐκ ἔδωκεν αὐτοῖς ἀναστάντας θῦσαι τὸν Σαούλ· καὶ ἀνέστη Σαοὺλ καὶ κατέβη τὴν ὁδόν.

\vs{9}Καὶ ἀνέστη Δαυὶδ ὀπίσω αὐτοῦ ἐκ τοῦ σπηλαίου· καὶ ἐβόησε Δαυὶδ ὀπίσω Σαοὺλ, λέγων, κύριε βασιλευ· καὶ ἐπέβλεψε Σαοὺλ εἰς τᾶ ὀπίσω αὐτοῦ, καὶ ἔκυψε Δαυὶδ ἐπὶ πρόσωπον αὐτοῦ ἐπὶ τὴν γῆν, καὶ προσεκύνησεν αὐτῷ.

\vs{10}Καὶ εἶπε Δαυὶδ πρὸς Σαοὺλ, ἱνατί ἀκούεις τῶν λόγων τοῦ λαοῦ, λεγόντων, ἰδοὺ Δαυὶδ ζητεῖ τὴν ψυχήν σου;
\vs{11}Ἰδοὺ ἐν τῇ ἡμέρᾳ ταύτῃ ἑωράκασιν οἱ ὀφθαλμοί σου ὡς παρέδωκέ σε Κύριος σήμερον εἰς χεῖράς μου ἐν τῷ σπηλαίῳ, καὶ οὐκ ἠβουλήθην ἀποκτεῖναί σε, καὶ ἐφεισάμην σου, καὶ εἶπα, οὐκ ἐποίσω χεῖρά μου ἐπὶ κύριόν μου, ὅτι χριστὸς Κυρίου οὗτός ἐστι.
\vs{12}καὶ ἰδοὺ τὸ πτερύγιον τῆς διπλοΐδος σου ἐν τῇ χειρί μου, ἐγὼ ἀφῄρηκα τὸ πτερύγιον, καὶ οὐκ ἀπέκτανκά σε· καὶ γνῶθι καὶ ἴδε σήμερον, ὅτι οὐκ ἔστι κακία ἐν τῇ χειρί μου οὐδὲ ἀσέβεια καὶ ἀθέτησις, καὶ οὐχ ἡμάρτηκα εἰς σὲ, καὶ σὺ δεσμεύεις τὴν ψυχήν μου λαβεῖν αὐτήν.
\vs{13}Δικάσαι Κύριος ἀναμέσον ἐμοῦ καὶ σοῦ, καὶ ἐκδικήσαι σοι Κύριος ἐκ σου· καὶ ἡ χείρ μου οὐκ ἔσται ἐπὶ σὲ,
\vs{14}καθὼς λέγεται ἡ παραβολὴ ἡ ἀρχαῖα, ἐξ ἀνόμων ἐξελεύσεται πλημμέλεια· καὶ ἡ χείρ μου οὐκ ἔσται ἐπὶ σέ.
\vs{15}Καὶ νῦν ὀπίσω τίνος σὺ ἐκπορεύῃ βασιλεῦ Ἰσραήλ; ὀπίσω τίνος καταδιώκεις σύ; ὀπίσω κυνὸς τεθνηκότος, καὶ ὀπίσω ψύλλου ἑνός;
\vs{16}Γένοιτο Κύριος εἰς κριτὴν καὶ δικαστὴν ἀναμέσον ἐμοῦ καὶ ἀναμέσον σοῦ, ἴδοι Κύριος καὶ κρίναι τὴν κρίσιν μου, καὶ δικάσαι μοι ἐκ χειρός σου.

\vs{17}Καὶ ἐγένετο, ὡς συνετέλεσε Δαυὶδ τὰ ῥήματα ταῦτα λαλῶν πρὸς Σαοὺλ, καὶ εἶπε Σαοὺλ, ἡ φωνή σου αὕτη, τέκνον Δαυίδ; καὶ ᾖρε Σαοὺλ τὴν φωνὴν αὐτοῦ, καὶ ἔκλαυσε.
\vs{18}Καὶ εἶπε Σαοὺλ πρὸς Δαυὶδ, δίκαιος σὺ ὑπὲρ ἐμὲ, ὅτι σὺ ἀνταπέδωκάς μοι ἀγαθά, ἐγὼ δὲ ἀνταπέδωκά σοι κακά.
\vs{19}Καὶ σὺ ἀπήγγειλάς μοι σήμερον ἃ ἐποίησάς μοι ἀγαθὰ, ὡς ἀπέκλεισέ με Κύριος εἰς χεῖράς σου σήμερον, καὶ οὐκ ἀπέκτεινάς με.
\vs{20}Καὶ ὅτι εἰ εὕροι τις τὸν ἐχθρὸν αὐτοῦ ἐν θλίψει, καὶ ἐκπέμψοι αὐτὸν ἐν ὁδῷ ἀγαθῇ, καὶ Κύριος ἀποτίσει αὐτῷ ἀγαθὰ, καθὼς πεποίηκας σήμερον.
\vs{21}Καὶ νῦν ἰδοὺ ἐγὼ γινώσκω, ὅτι βασιλεύων βασιλεύσεις, καὶ στήσεται ἐν χειρί σου ἡ βασιλεία Ἰσραήλ.
\vs{22}Καὶ νῦν ὄμοσόν μοι ἐν Κυρίῳ, ὅτι οὐκ ἐξολοθρεύσεις τὸ σπέρμα μου ὀπίσω μοῦ, οὐκ ἀφανιεῖς τὸ ὄνομά μου ἐκ τοῦ οἴκου τοῦ πατρός μου.
\vs{23}Καὶ ὤμοσε Δαυὶδ τῷ Σαούλ· καὶ ἀπῆλθε Σαοὺλ εἰς τὸν τόπον αὐτοῦ, καὶ Δαυὶδ καὶ οἱ ἄνδρες αὐτοῦ ἀνέβησαν εἰς τὴν Μεσσερὰ στενήν.

\ch{25}
Καὶ ἀπέθανε Σαμουὴλ, καὶ συναθροίζονται πᾶς Ἰσραὴλ, καὶ κόπτονται αὐτὸν, καὶ θάπτουσιν αὐτὸν ἐν οἴκῳ αὐτοῦ ἐν Ἁρμαθαίμ· καὶ ἀνέστη Δαυὶδ, καὶ κατέβη εἰς τὴν ἔρημον Μαών.

\vs{2}Καὶ ἦν ἄνθρωπος ἐν τῇ Μαὼν, καὶ τὰ ποίμνια αὐτοῦ ἐν τῷ Καρμήλῳ, καὶ ὁ ἄνθρωπος μέγας σφόδρα· καὶ τούτῳ ποίμνια τρισχίλια, καὶ αἶγες χίλιαι· καὶ ἐγενήθη ἐν τῷ κείρειν τὸ ποίμνιον αὐτοῦ ἐν τῷ Καρμήλῳ.
\vs{3}Καὶ ὄνομα τῷ ἀνθρώπῳ Νάβαλ, καὶ ὄνομα τῇ γυναικὶ αὐτοῦ Ἀβιγαία· καὶ ἡ γυνὴ αὐτοῦ ἀγαθὴ συνέσει καὶ ἀγαθὴ τῷ εἴδει σφόδρα· καὶ ὁ ἄνθρωπος σκληρὸς καὶ πονηρὸς ἐν ἐπιτηδεύμασι, καὶ ὁ ἄνθρωπος κυνικός.
\vs{4}καὶ ἤκουσε Δαυὶδ ἐν τῇ ἐρήμῳ, ὅτι κείρει Νάβαλ ὁ Καρμήλιος τὸ ποίμνιον αὐτοῦ.
\vs{5}Καὶ ἀπέστειλε Δαυὶδ δέκα παιδάρια, καὶ εἶπε τοῖς παιδαρίοις, ἀνάβητε εἰς Κάρμηλον καὶ ἀπέλθατε πρὸς Νάβαλ, καὶ ἐρωτήσατε αὐτὸν ἐπὶ τῷ ὀνόματί μου εἰς εἰρήνην,
\vs{6}καὶ ἐρεῖτε τάδε· εἰς ὥρας καὶ σὺ ὑγιαίνων καὶ ὁ οἶκός σου, καὶ πάντα τὰ σὰ ὑγιαίνοντα.

\vs{7}Καὶ νῦν ἰδοὺ ἀκήκοα ὅτι κείρουσί σοι νῦν οἱ ποιμένες σου οἳ ἦσαν μεθʼ ἡμῶν ἐν τῇ ἐρήμῳ, καὶ οὐκ ἀπεκωλύσαμεν αὐτοὺς, καὶ οὐκ ἐνετειλάμεθα αὐτοῖς οὐθὲν πάσας τὰς ἡμέρας ὄντων αὐτῶν ἐν Καρμήλῳ.
\vs{8}Ἐρώτησον τὰ παιδάριά σου, καὶ ἀπαγγελοῦσί σοι· καὶ εὑρέτωσαν τὰ παιδάριά σου χάριν ἐν ὀφθαλμοῖς σου, ὅτι ἐφʼ ἡμέραν ἀγαθὴν ἥκομεν· δὸς δὴ ὃ ἐὰν εὕρῃ ἡ χείρ σου τῷ υἱῷ σου τῷ Δαυίδ.

\vs{9}Καὶ ἔρχονται τὰ παιδάρια, καὶ λαλοῦσι τοὺς λόγους τούτους πρὸς Νάβαλ κατὰ πάντα τὰ ῥήματα ταῦτα ἐν τῷ ὀνόματι Δαυίδ· καὶ ἀνεπήδησε
\vs{10}καὶ ἀπεκρίθη Νάβαλ τοῖς παισὶ Δαυὶδ, καὶ εἶπε, τίς ὁ Δαυὶδ, καὶ τίς ὁ υἱὸς Ἰεσσαί; σήμερον πεπληθυμένοι εἰσὶν οἱ δοῦλοι ἀναχωροῦντες ἕκαστος ἐκ προσώπου τοῦ κυρίου αὐτοῦ.
\vs{11}Καὶ λήψομαι τοὺς ἄρτους μου καὶ τὸν οἶνόν μου καὶ τὰ θύματά μου ἃ τέθυκα τοῖς κείρουσί μου τὰ πρόβατα, καὶ δώσω αὐτὰ ἀνδράσιν, οἷς οὐκ οἶδα πόθεν εἰσί;
\vs{12}Καὶ ἀπεστράφησαν τὰ παιδάρια Δαυὶδ εἰς ὁδὸν αὐτῶν, καὶ ἀνέστρεψαν καὶ ἦλθον, καὶ ἀνήγγειλαν τῷ Δαυὶδ κατὰ τὰ ῥήματα ταῦτα.
\vs{13}Καὶ εἶπε Δαυὶδ τοῖς ἀνδράσιν αὐτοῦ, ζώσασθε ἕκαστος τὴν ῥομφαίαν αὐτοῦ· καὶ ἀνέβησαν ὀπίσω Δαυὶδ, ὡς τετρακόσιοι ἄνδρες· καὶ οἱ διακόσιοι ἐκάθισαν μετὰ τῶν σκευῶν.

\vs{14}Καὶ τῇ Ἀβιγαίᾳ γυναικὶ Νάβαλ ἀπήγγειλεν ἓν τῶν παιδαρίων, λέγων, ἰδοὺ Δαυὶδ ἀπέστειλεν ἀγγέλους ἐκ τῆς ἐρήμου εὐλογῆσαι τὸν κύριον ἡμῶν, καὶ ἐξέκλινεν ἀπʼ αὐτῶν.
\vs{15}Καὶ οἱ ἄνδρες ἀγαθοὶ ἡμῖν σφόδρα, οὐκ ἀπεκώλυσαν ἡμᾶς, οὐδὲ ἐνετείλαντο ἡμῖν οὐδὲν πάσας τὰς ἡμέρας ἃς ἦμεν παρʼ αὐτοῖς.
\vs{16}Καὶ ἐν τῷ εἶναι ἡμᾶς ἐν ἀγρῷ, ὡς τεῖχος ἦσαν περὶ ἡμᾶς καὶ τὴν νύκτα καὶ τὴν ἡμέραν, πάσας τὰς ἡμέρας ἃς ἦμεν παρʼ αὐτοῖς ποιμαίνοντες τὸ ποίμνιον.
\vs{17}Καὶ νῦν γνῶθι καὶ ἴδε σὺ τί ποιήσεις, ὅτι συντετέλεσται ἡ κακία εἰς τὸν κύριον ἡμῶν καὶ εἰς τὸν οἶκον αὐτοῦ· καὶ οὗτος υἱὸς λοιμὸς, καὶ οὐκ ἔστι λαλῆσαι πρὸς αὐτόν.

\vs{18}Καὶ ἔσπευσεν Ἀβιγαία καὶ ἔλαβε διακοσίους ἄρτους, καὶ δύο ἀγγεῖα οἴνου, καὶ πέντε πρόβατα πεποιημένα, καὶ πέντε οἰφὶ ἀλφίτου, καὶ γόμορ ἓν σταφίδος, καὶ διακοσίας παλάθας, καὶ ἔθετο ἐπὶ τοὺς ὄνους.
\vs{19}Καὶ εἶπε τοῖς παιδαρίοις αὐτῆς, προπορεύεσθε ἔμπροσθέν μου, καὶ ἰδοὺ, ἐγὼ ὀπίσω ὑμῶν παραγίνομαι· καὶ τῷ ἀνδρὶ αὐτῆς οὐκ ἀπήγγειλε.
\vs{20}Καὶ ἐγενήθη, αὐτῆς ἐπιβεβηκυίης ἐπὶ τὴν ὄνον καὶ καταβαινούσης ἐν σκέπῃ τοῦ ὄρους, καὶ ἰδοὺ Δαυὶδ καὶ οἱ ἄνδρες αὐτοῦ κατέβαινον εἰς συνάντησιν αὐτῆς, καὶ ἀπήντησεν αὐτοῖς.
\vs{21}Καὶ Δαυὶδ εἶπεν, ἴσως εἰς ἄδικον πεφύλακα πάντα τὰ αὐτοῦ ἐν τῇ ἐρήμῳ, καὶ οὐκ ἐνετειλάμεθα λαβεῖν ἐκ πάντων τῶν αὐτοῦ οὐθὲν, καὶ ἀνταπέδωκέ μοι πονηρὰ ἀντὶ ἀγαθῶν.
\vs{22}Τάδε ποιήσαι ὁ Θεὸς τῷ Δαυὶδ καὶ τάδε προσθείη, εἰ ὑπολείψομαι ἐκ πάντων τῶν τοῦ Νάβαλ ἕως πρωΐ οὐροῦντα πρὸς τοῖχον.

\vs{23}Καὶ εἶδεν Ἀβιγαία τὸν Δαυίδ, καὶ ἔσπευσε καὶ κατεπήδησεν ἀπὸ τῆς ὄνου, καὶ ἔπεσεν ἐνώπιον Δαυὶδ ἐπὶ πρόσωπον αὐτῆς, καὶ προσεκύνησεν αὐτῷ ἐπὶ τὴν γῆν
\vs{24}ἐπὶ τοὺς πόδας αὐτοῦ, καὶ εἶπεν, ἐν ἐμοὶ κύριέ μου ἡ ἀδικία μου, λαλησάτω δὴ ἡ δούλη σου εἰς τὰ ὦτά σου, καὶ ἄκουσον λόγων τῆς δούλης σου.
\vs{25}Μὴ δὴ θέσθω ὁ κύριός μου καρδίαν αὐτοῦ ἐπὶ τὸν ἄνθρωπον τὸν λοιμὸν τοῦτον, ὅτι κατὰ τὸ ὄνομα αὐτοῦ οὗτός ἐστι· Νάβαλ ὄνομα αὐτῷ, καὶ ἀφροσύνη μετʼ αὐτοῦ· καὶ ἐγὼ ἡ δούλη σου οὐκ εἶδον τὰ παιδάρια τοῦ κυρίου μου ἃ ἀπέστειλας.

\vs{26}Καὶ νῦν κύριέ μου, ζῇ Κύριος καὶ ζῇ ἡ ψυχή σου, καθὼς ἐκώλυσέ σε Κύριος τοῦ μὴ ἐλθεῖν εἰς αἷμα ἀθῶον, καὶ σώζειν τὴν χεῖρά σου σοί· καὶ νῦν γένοιντο ὡς Νάβαλ οἱ ἐχθροί σου καὶ οἱ ζητοῦντες τῷ κυρίῳ μου κακά.
\vs{27}Καὶ νῦν λάβε τὴν εὐλογίαν ταύτην, ἣν ἐνήνοχεν ἡ δούλη σου τῷ κυρίῳ μου, καὶ δώσεις τοῖς παιδαρίοις τοῖς παρεστηκόσι τῷ κυρίῳ μου.
\vs{28}Ἆρον δὴ τὸ ἀνόμημα τῆς δούλης σου, ὅτι ποιῶν ποιήσει Κύριος τῷ κυρίῳ μου οἶκον πιστὸν, ὅτι πόλεμον κυρίου μου ὁ Κύριος πολεμεῖ, καὶ κακία οὐχ εὑρεθήσεται ἐν σοὶ πώποτε.
\vs{29}Καὶ ἀναστήσεται ἄνθρωπος καταδιώκων σε καὶ ζητῶν τὴν ψυχήν σου, καὶ ἔσται ψυχὴ κυρίου μου ἐνδεδεμένη ἐν δεσμῷ τῆς ζωῆς παρὰ Κυρίῳ τῷ Θεῷ, καὶ ψυχὴν ἐχθρῶν σου σφενδονήσεις ἐν μέσῳ τῆς σφενδόνης.
\vs{30}Καὶ ἔσται ὅτε ποιήσῃ Κύριος τῷ κυρίῳ μου πάντα ὅσα ἐλάλησεν ἀγαθὰ ἐπὶ σὲ, καὶ ἐντελεῖταί σοι εἰς ἡγούμενον ἐπὶ Ἰσραὴλ,
\vs{31}καὶ οὐκ ἔσται σοι τοῦτο βδελυγμὸς καὶ σκάνδαλον τῷ κυρίῳ μου, ἐκχέαι αἷμα ἀθῶον δωρεὰν, καὶ σῶσαι χεῖρα κυρίῳ μου αὐτῷ· καὶ ἀγαθώσαι Κύριος τῷ κυρίῳ μοῦ, καὶ μνησθήσῃ τῆς δούλης σου ἀγαθῶσαι αὐτῇ.

\vs{32}Καὶ εἶπε Δαυὶδ τῇ Ἀβιγαίᾳ, εὐλογητὸς Κύριος ὁ Θεὸς Ἰσραὴλ, ὃς ἀπέστειλέ σε σήμερον ἐν ταύτῃ εἰς ἀπάντησίν μοι,
\vs{33}καὶ εὐλογητὸς ὁ τρόπος σου, καὶ εὐλογημένη σὺ ἡ ἀποκωλύσασά με σήμερον ἐν ταύτῃ μὴ ἐλθεῖν εἰς αἵματα, καὶ σῶσαι χεῖρά μου ἐμοί.
\vs{34}Πλὴν ὅτι ζῇ Κύριος ὁ Θεὸς Ἰσραὴλ, ὃς ἀπεκώλυσέ σε σήμερον τοῦ κακοποιῆσαί σε, ὅτι εἰ μὴ ἔσπευσας καὶ παρεγένου εἰς ἀπάντησίν μοι, τότε εἶπα, εἰ ὑπολειφθήσεται τῷ Νάβαλ ἕως φωτὸς τοῦ πρωῒ οὐρῶν πρὸς τοῖχον.
\vs{35}Καὶ ἔλαβε Δαυὶδ ἐκ χειρὸς αὐτῆς πάντα ἃ ἔφερεν αὐτῷ, καὶ εἶπεν αὐτῇ, ἀνάβηθι εἰς εἰρήνην εἰς οἶκόν σου· βλέπε, ἤκουσα τῆς φωνῆς σου, καὶ ἠρέτισα τὸ πρόσωπόν σου.

\vs{36}Καὶ παρεγενήθη Ἀβιγαία πρὸς Νάβαλ· καὶ ἰδοὺ αὐτῷ πότος ἐν οἴκῳ αὐτοῦ, ὡς πότος βασιλέως· καὶ ἡ καρδία Νάβαλ ἀγαθὴ ἐπʼ αὐτόν· καὶ αὐτὸς μεθύων ἕως σφόδρα· καὶ οὐκ ἀπήγγειλεν αὐτῷ ῥῆμα μικρὸν ἢ μέγα ἕως φωτὸς τοῦ πρωΐ.
\vs{37}Καὶ ἐγένετο πρωῒ, ὡς ἐξένηψεν ἀπὸ τοῦ οἴνου Νάβαλ, ἀπήγγειλεν ἡ γυνὴ αὐτοῦ τὰ ῥήματα ταῦτα· καὶ ἐναπέθανεν ἡ καρδία αὐτοῦ ἐν αὐτῷ, καὶ αὐτὸς γίνεται ὡς λίθος.

\vs{38}Καὶ ἐγένετο ὡσεὶ δέκα ἡμέραι, καὶ ἐπάταξε Κύριος τὸν Νάβαλ, καὶ ἀπέθανε.
\vs{39}Καὶ ἤκουσε Δαυὶδ, καὶ εἶπεν, εὐλογητὸς Κύριος, ὃς ἔκρινε τὴν κρίσιν τοῦ ὀνειδισμοῦ μου ἐκ χειρὸς Νάβαλ, καὶ τὸν δοῦλον αὐτοῦ περιεποιήσατο ἐκ χειρὸς κακῶν, καὶ τὴν κακίαν Νάβαλ ἀπέστρεψε Κύριος εἰς κεφαλὴν αὐτοῦ.

Καὶ ἀπέστειλε Δαυὶδ καὶ ἐλάλησε περὶ Ἀβιγαίας, λαβεῖν αὐτὴν ἑαυτῷ εἰς γυναῖκα.
\vs{40}Καὶ ἦλθον οἱ παῖδες Δαυὶδ πρὸς Ἀβιγαίαν εἰς Κάρμηλον, καὶ ἐλάλησαν αὐτῇ, λέγοντες, Δαυὶδ ἀπέστειλεν ἡμᾶς πρὸς σὲ, λαβεῖν σε αὐτῷ εἰς γυναῖκα.
\vs{41}Καὶ ἀνέστη καὶ προσεκύνησεν ἐπὶ τὴν γῆν ἐπὶ πρόσωπον, καὶ εἶπεν, ἰδοὺ ἡ δούλη σου εἰς παιδίσκην νίψαι πόδας τῶν παίδων σου.
\vs{42}Καὶ ἀνέστη Ἀβιγαία, καὶ ἐπέβη ἐπὶ τὴν ὄνον, καὶ πέντε κοράσια ἠκολούθουν αὐτῇ, καὶ ἐπορεύθη ὀπίσω τῶν παίδων Δαυὶδ, καὶ γίνεται αὐτῷ εἰς γυναῖκα.
\vs{43}Καὶ τὴν Ἀχινάαμ ἔλαβε Δαυὶδ ἐξ Ἰεζραὲλ, καὶ ἀμφότεραι ἦσαν αὐτῷ γυναῖκες.
\vs{44}Καὶ Σαοὺλ ἔδωκε Μελχὸλ τὴν θυγατέρα αὐτοῦ τὴν γυναῖκα Δαυὶδ τῷ Φαλτὶ υἱῷ Ἀμὶς τῷ ἐκ Ῥομμᾶ.

\ch{26}
Καὶ ἔρχονται οἱ Ζιφαῖοι ἐκ τῆς αὐχμώδους πρὸς τὸν Σαοὺλ εἰς τὸν βουνὸν, λέγοντες, ἰδοὺ Δαυὶδ σκεπάζεται μεθʼ ἡμῶν ἐν τῷ βουνῷ τῷ ʼΕχελὰ, κατὰ πρόσωπον τοῦ Ἰεσσεμοῦ.
\vs{2}Καὶ ἀνέστη Σαοὺλ, καὶ κατέβη εἰς τὴν ἔρημον Ζὶφ, καὶ μετʼ αὐτοὺ τρεῖς χιλιάδες ἀνδρῶν ἐκλεκτοὶ ἐξ Ἰσραὴλ, ζητεῖν τὸν Δαυὶδ ἐν τῇ ἐρήμῳ Ζίφ.
\vs{3}Καὶ παρενέβαλε Σαοὺλ ἐν τῷ βουνῷ τῷ Ἐχελὰ τῷ ἐπὶ προσώπου τοῦ Ἰεσσεμοῦ ἐπὶ τῆς ὁδοῦ, καὶ Δαυὶδ ἐκάθισεν ἐν τῇ ἐρήμῳ· καὶ εἶδε Δαυὶδ, ὅτι ἥκει Σαοὺλ ὀπίσω αὐτοῦ εἰς τὴν ἔρημον.
\vs{4}Καὶ ἀπέστειλε Δαυὶδ κατασκόπους, καὶ ἔγνω ὅτι ἥκει Σαοὺλ ἕτοιμος ἐκ Κεϊλά.

\vs{5}Καὶ ἀνέστη Δαυὶδ λάθρα, καὶ εἰσπορεύεται εἰς τὸν τόπον οὗ ἐκάθευδεν ἐκεῖ Σαοὺλ, καὶ ἐκεῖ Ἀβεννὴρ υἱὸς Νὴρ ἀρχιστράτηγος αὐτοῦ· καὶ Σαοὺλ ἐκάθευδεν ἐν λαμπήνῃ, καὶ ὁ λαὸς παρεμβεβληκὼς κύκλῳ αὐτοῦ.
\vs{6}Καὶ ἀπεκρίθη Δαυὶδ, καὶ εἶπε πρὸς Ἀβιμέλεχ τὸν Χετταῖον, καὶ πρὸς Ἀβεσσὰ υἱὸν Σαρουίας ἀδελφὸν Ἰωὰβ, λέγων, τίς εἰσελεύσεται μετʼ ἐμοῦ πρὸς Σαοὺλ εἰς τὴν παρεμβολήν; καὶ εἶπεν Ἀβεσσὰ, ἐγὼ εἰσελεύσομαι μετὰ σοῦ.

\vs{7}Καὶ εἰσπορεύεται Δαυὶδ καὶ Ἀβεσσὰ εἰς τὸν λαὸν τὴν νύκτα· καὶ ἰδοὺ Σαοὺλ καθεύδων ὕπνῳ ἐν λαμπήνῃ, καὶ τὸ δόρυ αὐτοῦ ἐμπεπηγὸς εἰς τὴν γῆν πρὸς κεφαλῆς αὐτοῦ, καὶ Ἀβεννὴρ καὶ ὁ λαὸς αὐτοῦ ἐκάθευδε κύκλῳ αὐτοῦ.
\vs{8}Καὶ εἶπεν Ἀβεσσὰ πρὸς Δαυὶδ, ἀπέκλεισε Κύριος σήμερον τὸν ἐχθρόν σου εἰς χεῖράς σου· καὶ νῦν πατάξω αὐτὸν τῷ δόρατι εἰς τὴν γῆν ἅπαξ, καὶ οὐ δευτερώσω αὐτῷ.
\vs{9}Καὶ εἶπε Δαυὶδ πρὸς Ἀβεσσὰ, μὴ ταπεινώσῃς αὐτὸν, ὅτι τίς ἐποίσει χεῖρα αὐτοῦ ἐπὶ χριστὸν Κυρίου καὶ ἀθωωθήσεται;
\vs{10}Καὶ εἶπε Δαυὶδ, ζῇ Κύριος, ἐὰν μὴ Κύριος παίσῃ αὐτὸν, ἢ ἡμέρα αὐτοῦ ἔλθῃ καὶ ἀποθάνῃ, ἢ εἰς πόλεμον καταβῇ καὶ προστεθῇ.
\vs{11}Μηδαμῶς μοι παρὰ Κυρίου ἐπενεγκεῖν χεῖρά μου ἐπὶ χριστὸν Κυρίου· καὶ νῦν λάβε δὴ τὸ δόρυ ἀπὸ προσκεφαλῆς αὐτοῦ, καὶ τὸν φακὸν τοῦ ὕδατος, καὶ ἀπέλθωμεν ἡμεῖς καθʼ ἑαυτούς.
\vs{12}Καὶ ἔλαβε Δαυὶδ τὸ δόρυ, καὶ τὸν φακὸν τοῦ ὕδατος ἀπὸ προσκεφαλῆς αὐτοῦ, καὶ ἀπῆλθον καθʼ ἑαυτούς· καὶ οὐκ ἦν ὁ βλέπων, καὶ οὐκ ἦν ὁ γινώσκων, καὶ οὐκ ἦν ὁ ἐξεγειρόμενος, πάντες ὑπνοῦντες, ὅτι θάμβος Κυρίου ἐπέπεσεν ἐπʼ αὐτούς.

\vs{13}Καὶ διέβη Δαυὶδ εἰς τὸ πέραν, καὶ ἔστη ἐπὶ τὴν κορυφὴν τοῦ ὄρους μακρόθεν, καὶ πολλὴ ἡ ὁδὸς ἀναμέσον αὐτῶν.
\vs{14}Καὶ προσεκαλέσατο Δαυὶδ τὸν λαὸν, καὶ τῷ Ἀβεννὴρ ἐλάλησε, λέγεν, οὐκ ἀποκριθήσῃ Ἀβεννήρ; καὶ ἀπεκρίθη ʼΑβεννὴρ, καὶ εἶπε, τίς εἶ σὺ ὁ καλῶν;
\vs{15}Καὶ εἶπε Δαυὶδ πρὸς Ἀβεννὴρ, οὐκ ἀνὴρ σύ; καὶ τίς, ὡς σὺ, ἐν Ἰσραήλ; καὶ διατί οὐ φυλάσσεις τὸν κύριόν σου τὸν βασιλέα; ὅτι εἰσῆλθεν εἷς ἐκ τοῦ λαοῦ διαφθεῖραι τὸν κύριόν σου τὸν βασιλέα.
\vs{16}Καὶ οὐκ ἀγαθὸν τὸ ῥῆμα τοῦτο, ὃ πεποίηκας· ζῇ Κύριος, ὅτι υἱοὶ θανατώσεως ὑμεῖς, οἱ φυλάσσοντες τὸν βασιλέα τὸν κύριον ὑμῶν τὸν χριστὸν Κυρίου· καὶ νῦν ἴδε δὴ τὸ δόρυ τοῦ βασιλέως, καὶ ὁ φακὸς τοῦ ὕδατος, ποῦ ἐστι τὰ πρὸς κεφαλῆς αὐτοῦ;

\vs{17}Καὶ ἐπέγνω Σαοὺλ τὴν φωνὴν Δαυὶδ, καὶ εἶπεν, ἡ φωνή σου αὕτη, τέκνον Δαυίδ; καὶ εἶπε Δαυίδ, δοῦλός σου κύριε βασιλεῦ.
\vs{18}Καὶ εἶπεν, ἱνατί τοῦτο καταδιώκει ὁ κύριος ὀπίσω τοῦ δούλου αὐτοῦ; ὅτι τί ἡμάρτηκα; καὶ τί εὑρέθη ἐν ἐμοὶ ἀδίκημα;
\vs{19}Καὶ νῦν ἀκουσάτω ὁ κύριός μου ὁ βασιλεὺς τὸ ῥῆμα τοῦ δούλου αὐτοῦ· εἰ ὁ Θεὸς ἐπισείει σε ἐπʼ ἐμὲ, ὀσφρανθείη θυσία σου· καὶ εἰ υἱοὶ ἀνθρώπων, ἐπικατάρατοι οὗτοι ἐνώπιον Κυρίου, ὅτι ἐξέβαλόν με σήμερον μὴ ἐστηρίθαι ἐν κληρονομίᾳ Κυρίου, λέγοντες, πορεύου, δούλευε θεοῖς ἑτέροις.
\vs{20}Καὶ νῦν μὴ πέσοι τὸ αἷμά μου ἐπὶ τὴν γῆν ἐξεναντίας προσώπου Κυρίου, ὅτι ἐξελήλυθεν ὁ βασιλεὺς Ἰσραὴλ ζητεῖν ψυχήν μου, καθὼς καταδιώκει ὁ νυκτικόραξ ἐν τοῖς ὄρεσι.

\vs{21}Καὶ εἶπε Σαοὺλ, ἡμάρτηκα· ἐπίστρεφε, τέκνον Δαυὶδ, ὅτι οὐ κακοποιήσω σε, ἀνθʼ ὧν ἔντιμος ψυχή μου ἐν ὀφθαλμοῖς σου, καὶ ἐν τῇ σήμερον μεματαίωμαι καὶ ἠγνόηκα πολλὰ σφόδρα.
\vs{22}Καὶ ἀπεκρίθη Δαυὶδ, καὶ εἶπεν, ἰδοὺ τὸ δόρυ τοῦ βασιλέως· διελθέτω εἷς τῶν παιδαρίων καὶ λαβέτω αὐτό·
\vs{23}Καὶ Κύριος ἐπιστρέψει ἑκάστῳ κατὰ τὰς δικαιοσύνας αὐτοῦ καὶ τὴν πίστιν αὐτοῦ· ὡς παρέδωκέ σε Κύριος σήμερον εἰς χεῖράς μου, καὶ οὐκ ἠθέλησα ἐπενεγκεῖν χεῖρά μου ἐπὶ χριστὸν Κυρίου.
\vs{24}Καὶ ἰδοὺ καθὼς ἐμεγαλύνθη ἡ ψυχή σου σήμερον ἐν ταύτῃ ἐν ὀφθαλμοῖς μου, οὕτως μεγαλυνθείη ἡ ψυχή μου ἐνώπιον Κυρίου, καὶ σκεπάσαι με, καὶ ἐξελεῖταί με ἐκ πάσης θλίψεως.
\vs{25}Καὶ εἶπε Σαοὺλ πρὸς Δαυὶδ, εὐλογημένος σὺ, τέκνον· καὶ ποιῶν ποιήσεις, καὶ δυνάμενος δυνήσῃ· καὶ ἀπῆλθε Δαυὶδ εἰς τὴν ὁδὸν αὐτοῦ, καὶ Σαοὺλ ἀπέστρεψεν εἰς τὸν τόπον αὐτοῦ.

\ch{27}
Καὶ εἶπε Δαυὶδ ἐν τῇ καρδίᾳ αὐτοῦ, λέγων, νῦν προστεθήσομαι ἐν ἡμέρᾳ μιᾷ εἰς χεῖρας Σαούλ· καὶ οὐκ ἔστι μοι ἀγαθὸν ἐὰν μὴ σωθῶ εἰς γῆν ἀλλοφύλων, καὶ ἀνῇ ἀπʼ ἐμοῦ Σαοὺλ τοῦ ζητεῖν με εἰς πᾶν ὅριον Ἰσραὴλ, καὶ σωθήσομαι ἐκ χειρὸς αὐτοῦ.
\vs{2}Καὶ ἀνέστη Δαυὶδ καὶ οἱ ἐξακόσιοι ἄνδρες οἱ μετʼ αὐτοῦ, καὶ ἐπορεύθη πρὸς Ἀγχοῦς υἱὸν Ἀμμὰχ βασιλέα Γέθ.
\vs{3}Καὶ ἐκάθισε Δαυὶδ μετὰ Ἀγχοῦς, αὐτὸς καὶ οἱ ἄνδρες αὐτοῦ ἕκαστος καὶ ὁ οἶκος αὐτοῦ, καὶ Δαυὶδ καὶ ἀμφότεραι αἱ γυναῖκες αὐτοῦ, Ἀχινάαμ Ἰεζραηλίτις καὶ Ἀβιγαία ἡ γυνὴ Νάβαλ τοῦ Καρμηλίου.
\vs{4}Καὶ ἀνηγγέλη τῷ Σαοὺλ ὅτι πέφευγε Δαυὶδ εἰς Γὲθ, καὶ οὐ προσέθετο ἔτι ζητεῖν αὐτόν.

\vs{5}Καὶ εἶπε Δαυὶδ πρὸς Ἀγχοῦς, εἰ δὴ εὕρηκεν ὁ δοῦλός σου χάριν ἐν ὀφθαλμοῖς σου, δότωσαν δή μοι τόπον ἐν μιᾷ τῶν πόλεων τῶν κατʼ ἀγρὸν, καὶ καθήσομαι ἐκεῖ· καὶ ἱνατί κάθηται ὁ δοῦλός σου ἐν πόλει βασιλευομένῃ μετὰ σοῦ;
\vs{6}Καὶ ἔδωκεν αὐτῷ ἐν τῇ ἡμέρᾳ ἐκείνῃ τὴν Σεκελάκ· διὰ τοῦτο ἐγενήθη Σεκελὰκ τῷ βασιλεῖ τῆς Ἰουδαίας ἕως τῆς ἡμέρας ταύτης.
\vs{7}Καὶ ἐγενήθη ὁ ἀριθμὸς τῶν ἡμερῶν ὧν ἐκάθισε Δαυὶδ ἐν ἀγρῷ τῶν ἀλλοφύλων, τέσσαρας μῆνας.

\vs{8}Καὶ ἀνέβαινε Δαυὶδ καὶ οἱ ἄνδρες αὐτοῦ, καὶ ἐπετίθεντο ἐπὶ πάντα τὸν Γεσιρὶ καὶ ἐπὶ τὸν Ἀμαληκίτην· καὶ ἰδοὺ ἡ γῆ κατῳκεῖτο ἀπὸ ἀνηκόντων ἡ ἀπὸ Γελαμψοὺρ τετειχισμένων καὶ ἕως γῆς Αἰγύπτου.
\vs{9}Καὶ ἔτυπτε τὴν γῆν, καὶ οὐκ ἐζωογόνει ἄνδρα ἢ γυναῖκα· καὶ ἐλάμβανον ποίμνια καὶ βουκόλια καὶ ὄνους καὶ καμήλους καὶ ἱματισμὸν, καὶ ἀνέστρεψαν καὶ ἤρχοντο πρὸς Ἀγχοῦς.
\vs{10}Καὶ εἶπεν Ἀγχοῦς πρὸς Δαυὶδ ἐπὶ τίνα ἐπέθεσθε σήμερον; καὶ εἶπε Δαυὶδ πρὸς Ἀγχοῦς, κατὰ Νότον τῆς Ἰουδαίας καὶ κατὰ Νότον Ἰεσμεγὰ καὶ κατὰ Νότον τοῦ Κενεζί.
\vs{11}Καὶ ἄνδρα καὶ γυναῖκα οὐκ ἐζωογόνησα τοῦ εἰσαγαγεῖν εἰς Γὲθ, λέγων, μὴ ἀναγγειλωσιν εἰς Γὲθ καθʼ ἡμῶν, λέγοντες, τάδε Δαυὶδ ποιεῖ· καὶ τόδε τὸ δικαίωμα αὐτοῦ πάσας τὰς ἡμέρας ἃς ἐκάθητο Δαυὶδ ἐν ἀγρῷ τῶν ἀλλοφύλων.
\vs{12}Καὶ ἐπιστεύθη Δαυὶδ ἐν τῷ Ἀγχοῦς σφόδρα, λέγων, ἤσχυνται αἰσχυνόμενος ἐν τῷ λαῷ αὐτοῦ ἐν Ἰσραὴλ, καὶ ἔσται μοι δοῦλος εἰς τὸν αἰῶνα.

\ch{28}
Καὶ ἐγενήθη ἐν ταῖς ἡμέραις ἐκείναις, καὶ συναθροίζονται ἀλλόφυλοι ἐν ταῖς παρεμβολαῖς αὐτῶν ἐξελθεῖν πολεμεῖν μετὰ Ἰσραήλ· καὶ εἶπεν Ἀγχοῦς πρὸς Δαυὶδ, γινώσκων γνώσῃ, ὅτι μετʼ ἐμοῦ ἐξελεύσῃ εἰς πόλεμον σὺ, καὶ οἱ ἄνδρες σου.
\vs{2}Καὶ εἶπε Δαυὶδ πρὸς Ἀγχοῦς, οὕτω νῦν γνώσῃ ἃ ποιήσει ὁ δοῦλός σου· καὶ εἶπεν Ἀγχοῦς πρὸς Δαυὶδ, οὕτως ἀρχισωματοφύλακα θήσομαί σε πάσας τὰς ἡμέρας

\vs{3}Καὶ Σαμουὴλ ἀπέθανε, καὶ ἐκόψαντο αὐτὸν πᾶς Ἰσραὴλ, καὶ θάπτουσιν αὐτὸν ἐν Ἀρμαθαὶμ ἐν πόλει αὐτοῦ. καὶ Σαοὺλ περιεῖλε τοὺς ἐγγαστριμύθους καὶ τοὺς γνώστας ἀπὸ τῆς γῆς.
\vs{4}Καὶ συναθροίζονται οἱ ἀλλόφυλοι, καὶ ἔρχονται καὶ παρεμβάλλουσιν εἰς Σωνάμ· καὶ συναθροίζει Σαοὺλ πάντα ἄνδρα Ἰσραὴλ, καὶ παρεμβάλλουσιν εἰς Γελβουέ.
\vs{5}Καὶ εἶδε Σαοὺλ τὴν παρεμβολὴν τῶν ἀλλοφύλων, καὶ ἐφοβήθη, καὶ ἐξέστη ἡ καρδία αὐτοῦ σφόδρα.
\vs{6}Καὶ ἐπηρώτησε Σαοὺλ διὰ Κυρίου, καὶ οὐκ ἀπεκρίθη αὐτῷ Κύριος ἐν τοῖς ἐνυπνίοις καὶ ἐν τοῖς δήλοις καὶ ἐν τοῖς προφήταις.

\vs{7}Καὶ εἶπε Σαοὺλ τοῖς παισὶν αὐτοῦ, ζητήσατέ μοι γυναῖκα ἐγγαστρίμυθον, καὶ πορεύσομαι πρὸς αὐτὴν, καὶ ζητήσω ἐν αὐτῇ· καὶ εἶπαν οἱ παἷδες αὐτοῦ πρὸς αὐτὸν ἰδοὺ γυνὴ ἐγγαστρίμυθος ἐν Ἀενδώρ.

\vs{8}Καὶ συνεκαλύψατο Σαοὺλ, καὶ περιεβάλετο ἱμάτια ἕτερα, καὶ πορεύεται αὐτὸς καὶ δύο ἄνδρες μετʼ αὐτοῦ, καὶ ἔρχονται πρὸς τὴν γυναῖκα νυκτὸς, καὶ εἶπεν αὐτῇ, μάντευσαι δή μοι ἐν τῷ ἐγγαστριμύθῳ, καὶ ἀνάγαγέ μοι ὃν ἐὰν εἴπω σοι.
\vs{9}Καὶ εἶπεν αὐτῷ ἡ γυνὴ, ἰδοὺ δὴ σὺ οἶδας ὅσα ἐποίησε Σαοὺλ, ὡς ἐξωλόθρευσε τοὺς ἐγγαστριμύθους καὶ τοὺς γνώστας ἀπὸ τῆς γῆς, καὶ ἱνατί σὺ παγιδεύεις τὴν ψυχήν μοι θανατῶσαι αὐτήν;
\vs{10}Καὶ ὤμοσεν αὐτῇ Σαοὺλ, λέγων, ζῇ Κύριος, εἰ ἀπαντήσεταί σοι ἀδικία ἐν τῷ λόγῳ τούτῳ.
\vs{11}Καὶ εἶπεν ἡ γυνὴ, τίνα ἀναγάγω σοι; καὶ εἶπε, τὸν Σαμουὴλ ἀνάγαγέ μοι.

\vs{12}Καὶ εἶδεν ἡ γυνὴ τὸν Σαμουὴλ, καὶ ἀνεβόησε φωνῇ μεγάλῃ· καὶ εἶπεν ἡ γυνὴ πρὸς Σαοὺλ, ἱνατί παρελογίσω με; καὶ σὺ εἶ Σαούλ.
\vs{13}Καὶ εἶπεν αὐτῇ ὁ βασιλεὺς, μὴ φοβοῦ, εἶπον τίνα ἑώρακας· καὶ εἶπεν αὐτῷ ἡ γυνὴ, θεοὺς ἑώρακα ἀναβαίνοντας ἐκ τῆς γῆς.
\vs{14}Καὶ εἶπεν αὐτῇ, τί ἔγνως; καὶ εἶπεν αὐτῷ, ἄνδρα ὄρθιον ἀναβαίνοντα ἐκ τῆς γῆς, καὶ οὗτος διπλοΐδα ἀναβεβλημένος· καὶ ἔγνω Σαοὺλ, ὅτι οὗτος Σαμουὴλ, καὶ ἔκυψεν ἐπὶ πρόσωπον αὐτοῦ ἐπὶ τὴν γῆν, καὶ προσεκύνησεν αὐτῷ.

\vs{15}Καὶ εἶπε Σαμουὴλ, ἱνατί παρηνώχλησάς μοι ἀναβῆναί με; καὶ εἶπε Σαοὺλ, θλίβομαι σφόδρα, καὶ οἱ ἀλλόφυλοι πολεμοῦσιν ἐν ἐμοὶ, καὶ ὁ Θεὸς ἀφέστηκεν ἀπʼ ἐμοῦ, καὶ οὐκ ἐπακήκοέ μοι ἔτι καὶ ἐν χειρὶ τῶν προφητῶν καὶ ἐν τοῖς ἐνυπνίοις· καὶ νῦν κέκληκά σε γνωρίσαι μοι τί ποιήσω.
\vs{16}Καὶ εἶπε Σαμουὴλ, ἱνατί ἐπερωτᾷς με, καὶ Κύριος ἀφέστηκεν ἀπὸ σοῦ, καὶ γέγονε μετὰ τοῦ πλησίον σου;
\vs{17}Καὶ πεποίηκε Κύριος σοι, καθὼς ἐλάλησε Κύριος ἐν χειρί μου, καὶ διαῤῥήξει Κύριος τὴν βασιλείαν σου ἐκ χειρός σου, καὶ δώσει αὐτὴν τῷ πλησίον σου τῷ Δαυὶδ,
\vs{18}διότι οὐκ ἤκουσας φωνῆς Κυρίου, καὶ οὐκ ἐποίησας θυμὸν ὀργῆς αὐτοῦ ἐν Ἀμαλὴκ, διὰ τοῦτο τὸ ῥῆμα ἐποίησε Κύριός σοι ἐν τῇ ἡμέρᾳ ταύτῃ.
\vs{19}Καὶ παραδώσει Κύριος τὸν Ἰσραὴλ μετὰ σοῦ εἰς χεῖρας ἀλλοφύλων, καὶ αὔριον σὺ καὶ οἱ υἱοί σου μετὰ σοῦ πεσοῦνται, καὶ τὴν παρεμβολὴν Ἰσραὴλ δώσει Κύριος εἰς χεῖρας ἀλλοφύλων.

\vs{20}Καὶ ἔσπευσε Σαοὺλ καὶ ἔπεσεν ἑστηκὼς ἐπὶ τὴν γῆν, καὶ ἐφοβήθη σφόδρα ἀπὸ τῶν λόγων Σαμουὴλ, καὶ ἐν αὐτῷ οὐκ ἦν ἰσχὺς ἔτι, οὐ γὰρ ἔφαγεν ἄρτον ὅλην τὴν ἡμέραν καὶ ὅλην τὴν νύκτα ἐκείνην.
\vs{21}Καὶ εἰσῆλθεν ἡ γυνὴ πρὸς Σαοὺλ, καὶ εἶδεν ὅτι ἔσπευσε σφόδρα, καὶ εἶπε πρὸς αὐτὸν, ἰδοὺ δὴ ἤκουσεν ἡ δούλη σου τῆς φωνῆς σου, καὶ ἐθέμην τὴν ψυχήν μου ἐν τῇ χειρί μου, καὶ ἤκουσα τοὺς λόγους οὓς ἐλάλησάς μοι.
\vs{22}Καὶ νῦν ἄκουσον δὴ φωνῆς τῆς δούλης σου, καὶ παραθήσω ἐνώπιόν σου ψωμὸν ἄρτου, καὶ φάγε, καὶ ἔσται σοὶ ἰσχὺς, ὅτι πορεύῃ ἐν ὁδῷ.
\vs{23}Καὶ οὐκ ἐβουλήθη φαγεῖν· καὶ παρεβιάζοντο αὐτὸν οἱ παῖδες αὐτοῦ καὶ ἡ γυνὴ, καὶ ἤκουσε τῆς φωνῆς αὐτῶν, καὶ ἀνέστη ἀπὸ τῆς γῆς, καὶ ἐκάθισεν ἐπὶ τὸν δίφρον.
\vs{24}Καὶ τῇ γυναικὶ ἦν δάμαλις νομὰς ἐν τῇ οἰκίᾳ· καὶ ἔσπευσε καὶ ἔθυσεν αὐτήν· καὶ ἔλαβεν ἄλευρα καὶ ἐφύρασε, καὶ ἔπεψεν ἄζυμα,
\vs{25}καὶ προσήγαγεν ἐνώπιον Σαοὺλ, καὶ ἐνώπιον τῶν παίδων αὐτοῦ· καὶ ἔφαγον, καὶ ἀνέστησαν καὶ ἀπῆλθον τὴν νύκτα ἐκείνην.

\ch{29}
Καὶ συναθροίζουσιν ἀλλόφυλοι πάσας τὰς παρεμβολὰς αὐτῶν εἰς Ἀφὲκ, καὶ Ἰσραὴλ παρενέβαλεν ἐν Ἀεδὼρ τὴν ἐν Ἰεζραέλ.
\vs{2}Καὶ οἱ σατράπαι τῶν ἀλλοφύλων παρεπορεύοντο εἰς ἑκατοντάδας καὶ χιλιάδας, καὶ Δαυὶδ καὶ οἱ ἄνδρες αὐτοῦ παρεπορεύοντο ἐπʼ ἐσχάτων μετὰ Ἀγχοῦς.
\vs{3}Καὶ εἶπον οἱ σατράπαι τῶν ἀλλοφύλων, τίνες οἱ διαπορευόμενοι οὗτοι; καὶ εἶπεν Ἀγχοῦς πρὸς τοὺς στρατηγοὺς τῶν ἀλλοφύλων, οὐκ οὗτος Δαυὶδ ὁ δοῦλος Σαοὺλ βασιλέως Ἰσραήλ; γέγονε μεθʼ ἡμῶν ἡμέρας τοῦτο δεύτερον ἔτος, καὶ οὐχ εὕρηκα ἐν αὐτῷ οὐθὲν ἀφʼ ἧς ἡμέρας ἐνέπεσε πρὸς μὲ καὶ ἕως τῆς ἡμέρας ταύτης.
\vs{4}Καὶ ἐλυπήθησαν ἐπʼ αὐτῷ οἱ στρατηγοὶ τῶν ἀλλοφύλων, καὶ λέγουσιν αὐτῷ, ἀπόστρεψον τὸν ἄνδρα, καὶ ἀποστραφήτω εἰς τὸν τόπον αὐτοῦ, οὗ κατέστησας αὐτὸν ἐκεῖ, καὶ μὴ ἐρχέσθω μεθʼ ἡμῶν εἰς τὸν πόλεμον, καὶ μὴ γινέσθω ἐπίβουλος τῆς παρεμβολῆς· καὶ ἐν τίνι διαλλαγήσεται οὗτος τῷ κυρίῳ αὐτοῦ; οὐχὶ ἐν ταῖς κεφαλαῖς τῶν ἀνδρῶν ἐκείνων;
\vs{5}οὐχ οὗτος Δαυὶδ, ᾧ ἐξῆρχον ἐν χοροῖς, λέγοντες, ἐπάταξε Σαοὺλ ἐν χιλιάσιν αὐτοῦ, καὶ Δαυὶδ ἐν μυριάσιν αὐτοῦ;

\vs{6}Καὶ ἐκάλεσεν Ἀγχοῦς τὸν Δαυὶδ, καὶ εἶπεν αὐτῷ, ζῇ Κύριος, ὅτι εὐθὺς σὺ καὶ ἀγαθὸς ἐν ὀφθαλμοῖς μου, καὶ ἡ ἔξοδός σου καὶ ἡ εἴσοδός σου μετʼ ἐμοῦ ἐν τῇ παρεμβολῇ, καὶ ὅτι οὐχ εὕρηκα κατὰ σοῦ κακίαν ἀφʼ ἧς ἡμέρας ἥκεις πρὸς μὲ ἕως τῆς σήμερον ἡμέρας· καὶ ἐν ὀφθαλμοῖς τῶν σατραπῶν οὐκ ἀγαθὸς σύ.
\vs{7}Καὶ νῦν ἀνάστρεφε καὶ πορεύου εἰς εἰρήνην, καὶ οὐ μὴ ποιήσῃς κακίαν ἐν ὀφθαλμοῖς τῶν σατραπῶν τῶν ἀλλοφύλων.

\vs{8}Καὶ εἶπε Δαυὶδ πρὸς Ἀγχοῦς, τί πεποίηκά σοι καὶ τί εὗρες ἐν τῷ δούλῳ σου ἀφʼ ἧς ἡμέρας ἤμην ἐνώπιόν σου καὶ ἕως τῆς ἡμέρας ταύτης, ὅτι οὐ μὴ ἔλθω πολεμήσας τοὺς ἐχθροὺς τοῦ κυρίου μου τοῦ βασιλέως;

\vs{9}Καὶ ἀπεκρίθη Ἀγχοῦς πρὸς Δαυὶδ, οἶδα ὅτι ἀγαθὸς σὺ ἐν ὀφθαλμοῖς μου, ἀλλʼ οἱ σατράπαι τῶν ἀλλοφύλων λέγουσιν, οὐχ ἥξει μεθʼ ἡμῶν εἰς πόλεμον.
\vs{10}Καὶ νῦν ὄρθρισον τοπρωῒ σὺ καὶ οἱ παῖδες τοῦ κυρίου σου οἱ ἥκοντες μετὰ σοῦ, καὶ πορεύεσθε εἰς τὸν τόπον οὗ κατέστησα ὑμᾶς ἐκεῖ, καὶ λόγον λοιμὸν μὴ θῇς ἐν καρδίᾳ σου, ὅτι ἀγαθὸς σὺ ἐνώπιόν μου· καὶ ὀρθρίσατε ἐν τῇ ὁδῷ καὶ φωτισάτω ὑμῖν, καὶ πορεύθητε.

\vs{11}Καὶ ὤρθρισε Δαυὶδ αὐτὸς καὶ οἱ ἄνδρες αὐτοῦ ἀπελθεῖν καὶ φυλάσσειν τὴν γῆν τῶν ἀλλοφύλων, καὶ οἱ ἀλλόφυλοι ἀνέβησαν πολεμεῖν ἐπὶ Ἰεζραήλ.

\ch{30}
Καὶ ἐγενήθη εἰσελθόντος Δαυὶδ καὶ τῶν ἀνδρῶν αὐτοῦ τὴν Σεκελὰκ τῇ ἡμέρᾳ τῇ τρίτῃ, καὶ Ἀμαλὴκ ἐπέθετο ἐπὶ τὸν Νότον καὶ ἐπὶ τὴν Σεκελὰκ, καὶ ἐπάταξε τὴν Σεκελὰκ καὶ ἐνεπύρισαν αὐτὴν ἐν πυρί.
\vs{2}Καὶ τὰς γυναῖκας καὶ πάντα τὰ ἐν αὐτῇ ἀπὸ μικροῦ ἕως μεγάλου οὐκ ἐθανάτωσαν ἄνδρα καὶ γυναῖκα, ἀλλʼ ᾐχμαλὼτευσαν, καὶ ἀπῆλθον εἰς τὴν ὁδὸν αὐτῶν.

\vs{3}Καὶ ἦλθε Δαυὶδ καὶ οἱ ἄνδρες αὐτοῦ εἰς τὴν πόλιν, καὶ ἰδοὺ ἐμπεπύρισται ἐν πυρὶ, αἱ δὲ γυναῖκες αὐτῶν καὶ οἱ υἱοὶ αὐτῶν καὶ αἱ θυγατέρες αὐτῶν ᾐχμαλωτευμένοι.
\vs{4}Καὶ ᾖρε Δαυὶδ καὶ οἱ ἄνδρες αὐτοῦ τὴν φωνὴν αὐτῶν, καὶ ἔκλαυσαν ἕως ὅτου οὐκ ἦν ἐν αὐτοῖς ἰσχὺς ἔτι τοῦ κλαίειν.
\vs{5}Καὶ ἀμφότεραι αἱ γυναῖκες Δαυὶδ ᾐχμαλωτεύθησαν, Ἀχιναὰμ ἡ Ἰεζραηλίτις, καὶ Ἀβιγαία ἡ γυνὴ Νάβαλ τοῦ Καρμηλίου.
\vs{6}Καὶ ἐθλίβη Δαυὶδ σφόδρα, ὅτι εἶπεν ὁ λαὸς λιθοβολῆσαι αὐτὸν, ὅτι κατώδυνος ψυχὴ παντὸς τοῦ λαοῦ ἑκάστου ἐπὶ τοὺς υἱοὺς αὐτοῦ καὶ ἐπὶ τὰς θυγατέρας αὐτοῦ· καὶ ἐκραταιώθη Δαυὶδ ἐν Κυρίῳ Θεῷ αὐτοῦ.

\vs{7}Καὶ εἶπε Δαυὶδ πρὸς Ἀβιάθαρ τὸν ἱερέα υἱὸν Ἀχιμέλεχ, προσάγαγε τὸ ἐφούδ.
\vs{8}Καὶ ἐπηρώτησε Δαυὶδ διὰ τοῦ Κυρίου, λέγων, εἰ καταδιώξω ὀπίσω τοῦ γεδδοὺρ τούτου; εἰ καταλήμψομαι αὐτούς; καὶ εἶπεν αὐτῷ, καταδίωκε, ὅτι καταλαμβάνων καταλήμψῃ αὐτοὺς, καὶ ἐξαιρούμενος ἐξελῇ.
\vs{9}Καὶ ἐπορεύθη Δαυὶδ αὐτὸς καὶ οἱ ἑξακόσιοι ἄνδρες μετʼ αὐτοῦ, καὶ ἔρχονται ἕως τοῦ χειμάῤῥου Βοσὸρ, καὶ οἱ περισσοὶ ἔστησαν.
\vs{10}Καὶ κατεδίωξεν ἐν τετρακοσίοις ἀνδράσιν· ὑπέστησαν δὲ διακόσιοι ἄνδρες οἵτινες ἐκάθισαν πέραν τοῦ χειμάῤῥου τοῦ Βοσόρ.

\vs{11}Καὶ εὑρίσκουσιν ἄνδρα Αἰγύπτιον ἐν ἀγρῷ, καὶ λαμβάνουσιν αὐτὸν, καὶ ἄγουσιν αὐτὸν πρὸς Δαυίδ· καὶ διδόασιν αὐτῷ ἄρτον καὶ ἔφαγε, καὶ ἐπότισαν αὐτὸν ὕδωρ·
\vs{12}καὶ διδόασιν αὐτῷ κλάσμα παλάθης καὶ ἔφαγε, καὶ κατέστη τὸ πνεῦμα αὐτοῦ ἐν αὐτῷ· ὅτι οὐ βεβρώκει ἄρτον, καὶ οὐ πεπώκει ὕδωρ τρεῖς ἡμέρας καὶ τρεῖς νύκτας.
\vs{13}Καὶ εἶπεν αὐτῷ Δαυὶδ, τίνος σὺ εἶ, καὶ πόθεν εἶ; καὶ εἶπε τὸ παιδάριον τὸ Αἰγύπτιον, ἐγώ εἰμι δοῦλος ἀνδρὸς Ἀμαληκίτου, καὶ κατέλιπέ με ὁ κύριός μου, ὅτι ἠνωχλήθην ἐγὼ σήμερον τριταῖος.
\vs{14}Καὶ ἡμεῖς ἐπεθέμεθα ἐπὶ τὸν Νότον τοῦ Χελεθὶ, καὶ ἐπὶ τὰ τῆς Ἰουδαίας μέρη, καὶ ἐπὶ Νότον Χελοὺβ, καὶ τὴν Σεκελὰκ ἐνεπυρίσαμεν ἐν πυρί.
\vs{15}Καὶ εἶπεν, αὐτῷ Δαυὶδ, εἰ κατάξεις με ἐπὶ τὸ γεδδοὺρ τοῦτο; καὶ εἶπεν, ὄμοσον δή μοι κατὰ τοῦ Θεοῦ μὴ θανατώσειν με, καὶ μὴ παραδοῦναί με εἰς χεῖρας τοῦ κυρίου μου, καὶ κατάξω σε ἐπὶ τὸ γεδδοὺρ τοῦτο.

\vs{16}Καὶ κατήγαγεν αὐτὸν ἐκεῖ, καὶ ἰδοὺ οὗτοι διακεχυμένοι ἐπὶ πρόσωπον πάσης τῆς γῆς, ἐσθίοντες καὶ πίνοντες καὶ ἑορτάζοντες ἐν πᾶσι τοῖς σκύλοις τοῖς μεγάλοις, οἷς ἔλαβον ἐκ γῆς ἀλλοφύλων καὶ ἐκ γῆς Ἰούδα.
\vs{17}Καὶ ἦλθεν ἐπʼ αὐτοὺς Δαυὶδ, καὶ ἐπάταξεν αὐτοὺς ἀπὸ ἑωσφόρου ἕως δείλης καὶ τῇ ἐπαύριον· καὶ οὐκ ἐσώθη ἐξ αὐτῶν ἀνὴρ, ὅτι ἀλλʼ ἢ τετρακόσια παιδάρια, ἃ ἦν ἐπιβεβηκότα ἐπὶ τὰς καμήλους, καὶ ἔφυγον.
\vs{18}Καὶ ἀφείλατο Δαυὶδ πάντα ἃ ἔλαβον οἱ Ἀμαληκῖται, καὶ ἀμφότερας τὰς γυναῖκας αὐτοῦ ἐξείλατο.
\vs{19}Καὶ οὐ διεφώνησεν αὐτοῖς ἀπὸ μικροῦ ἕως μεγάλου, καὶ ἀπὸ τῶν σκύλων, καὶ ἕως υἱῶν καὶ θυγατέρων, καὶ ἕως πάντων ὧν ἔλαβον αὐτῶν, καὶ πάντα ἐπέστρεψε Δαυίδ.
\vs{20}Καὶ ἔλαβε πάντα τὰ ποίμνια, καὶ τὰ βουκόλια, καὶ ἀπήγαγεν ἔμπροσθεν τῶν σκύλων· καὶ τοῖς σκύλοις ἐκείνοις ἐλέγετο, ταῦτα τὰ σκῦλα Δαυίδ.

\vs{21}Καὶ παραγίνεται Δαυὶδ πρὸς τοὺς διακοσὶους ἄνδρας τούς ὑπολειφθέντας τοῦ πορεύεσθαι ὀπίσω Δαυὶδ, καὶ ἐκάθισεν αὐτοὺς ἐν τῷ χειμάῤῥῳ τοῦ Βοσὸρ, καὶ ἐξῆλθον εἰς ἀπάντησιν Δαυὶδ καὶ εἰς ἀπάντησιν τοῦ λαοῦ τοῦ μετʼ αὐτοῦ· καὶ προσήγαγε Δαυὶδ ἕως τοῦ λαοῦ, καὶ ἠρώτησαν αὐτὸν τὰ εἰς εἰρήνην.

\vs{22}Καὶ ἀπεκρίθη πᾶς ἀνὴρ λοιμὸς καὶ πονηρὸς τῶν ἀνδρῶν τῶν πολεμιστῶν τῶν πορευθέντων μετὰ Δαυὶδ, καὶ εἶπον, ὅτι οὐ κατεδίωξαν μεθʼ ἡμῶν, οὐ δώσομεν αὐτοῖς ἐκ τῶν σκύλων ὧν ἐξειλόμεθα, ὅτι ἀλλʼ ἢ ἕκαστος τὴς γυναῖκα αὐτοῦ, καὶ τὰ τέκνα αὐτοῦ ἀπαγέσθωσαν, καὶ ἀποστρεφέτωσαν.
\vs{23}Καὶ εἶπε Δαυὶδ, οὐ ποιήσετε οὕτως, μετὰ τὸ παραδοῦναι τὸν Κύριον ἡμῖν, καὶ φυλάξαι ἡμᾶς, καὶ παρέδωκε Κύριος τὸν γεδδοὺρ τὸν ἐπερχόμενον ἐφʼ ἡμᾶς, εἰς χεῖρας ἡμῶν.
\vs{24}Καὶ τίς ἐπακούσεται ὑμῶν τῶν λόγων τούτων; ὅτιδοὐχ ἧττον ἡμῶν εἰσι, διότι κατὰ τὴν μερίδα τοῦ καταβαίνοντος εἰς τὸν πόλεμον, οὕτως ἔσται ἡ μερὶς τοῦ καθημένου ἐπὶ τὰ σκεύη, κατὰ τὸ αὐτὸ μεριοῦνται.
\vs{25}Καὶ ἐγενήθη ἀπὸ τῆς ἡμέρας ἐκείνης καὶ ἐπάνω, καὶ ἐγένετο εἰς προστάγμα καὶ εἱς δικαίωμα τῷ Ἰσραὴλ ἕως τῆς σήμερον.

\vs{26}Καὶ ἦλθε Δαυὶδ εἰς Σεκελὰκ, καὶ ἀπέστειλε τοῖς πρεσβυτέροις τῶν σκύλων Ἰούδα καὶ τοῖς πλησίον αὐτοῦ, λέγων, ἰδοὺ ἀπὸ τῶν σκύλων τῶν ἐχθρῶν Κυρίου,
\vs{27}τοῖς ἐν Βαιθσοὺρ, καὶ τοῖς Ῥαμᾷ Νότου, καὶ τοῖς ἐν Γεθὸρ,
\vs{28}καὶ τοῖς ἐν Ἀροὴρ, καὶ τοῖς ἐν Ἀμμαδὶ, καὶ τοῖς ἐν Σαφὶ, καὶ τοῖς ἐν Ἐσθιὲ,
\vs{28a}καὶ τοὶς ἐν Γὲθ, καὶ τοῖς ἐν Κιμὰθ, καὶ τοῖς ἐν Σαφὲκ, καὶ τοῖς ἐν Θημὰθ,
\vs{29}καὶ τοῖς ἐν Καρμήλῳ, καὶ τοῖς ἐν ταῖς πόλεσι τοῦ Ἱερεμεὴλ, καὶ τοῖς ἐν ταῖς πόλεσι τοῦ Κενεζὶ,
\vs{30}καὶ τοῖς ἐν Ἱεριμοὺθ, καὶ τοῖς ἐν Βηρσαβεὶ, καὶ τοῖς ἐν Νομβὲ,
\vs{31}καὶ τοῖς ἐν Χεβρὼν, καὶ πάντας τοὺς τόπους οὓς διῆλθε Δαυὶδ ἐκεῖ αὐτὸς καὶ οἱ ἄνδρες αὐτοῦ.

\ch{31}
Καὶ οἱ ἀλλόφύλοι ἐπολέμουν ἐπὶ Ἰσραήλ· καὶ ἔφυγον οἱ ἄνδρες Ἰσραὴλ ἐκ προσώπου τῶν ἀλλοφύλων, καὶ πίπτουσι τραυματίαι ἐν τῷ ὄρει τῷ Γελβουέ.
\vs{2}Καὶ συνάπτουσιν οἱ ἀλλόφυλοι τῷ Σαοὺλ καὶ τοῖς υἱοῖς αὐτοῦ, καὶ τύπτουσιν ἀλλόφυλοι τὸν Ἰωναθάν, καὶ τὸν Ἀμιναδὰβ, καὶ τὸν Μελχισᾶ υἱὸν Σαούλ.
\vs{3}Καὶ βαρύνεται ὁ πόλεμος ἐπὶ Σαοὺλ, καὶ εὑρίσκουσιν αὐτὸν οἱ ἀκοντισταὶ ἄνδρες τόξοται, καὶ ἐτραυματίσθη εἰς τὰ ὑποχόνδρια.
\vs{4}Καὶ εἶπε Σαοὺλ πρὸς τὸν αἴροντα τὰ σκεύη αὐτοῦ, σπάσαι τὴν ῥομφαίαν σου καὶ ἀποκέντησόν με ἐν αὐτῇ, μὴ ἔλθωσιν οἱ ἀπερίτμητοι οὗτοι καὶ ἀποκεντήσωσί με καὶ ἐμπαίξωσιν ἐμοί· καὶ οὐκ ἐβούλετο ὁ αἴρων τὰ σκεύη αὐτοῦ, ὅτι ἐφοβήθη σφόδρα· καὶ ἔλαβε Σαοὺλ τὴν ῥομφαίαν, καὶ ἐπέπεσεν ἐπʼ αὐτήν.
\vs{5}Καὶ εἶδεν ὁ αἴρων τὰ σκεύη αὐτοῦ ὅτι τέθνηκε Σαοὺλ, καὶ ἐπέπεσε καὶ αὐτὸς ἐπὶ τὴν ῥομφαίαν αὐτοῦ, καὶ ἀπέθανε μετʼ αὐτοῦ.
\vs{6}Καὶ ἀπέθανε Σαοὺλ, καὶ οἱ τρεῖς υἱοὶ αὐτοῦ, καὶ ὁ αἴρων τὰ σκεύη αὐτοῦ, ἐν τῇ ἡμέρᾳ ἐκείνῃ κατὰ τὸ αὐτό.

\vs{7}Καὶ εἶδον οἱ ἄνδρες Ἰσραὴλ οἱ ἐν τῷ πέραν τῆς κοιλάδος καὶ οἱ ἐν τῷ πέραν τοῦ Ἰορδάνου, ὅτι ἔφυγον οἱ ἄνδρες Ἰσραὴλ, καὶ ὅτι τέθνηκε Σαοὺλ καὶ οἱ υἱοὶ αὐτοῦ, καὶ καταλείπουσι τὰς πόλεις αὐτῶν καὶ φεύγουσι· καὶ ἔρχονται οἱ ἀλλόφυλοι καὶ κατοικοῦσιν ἐν αὐταῖς.

\vs{8}Καὶ ἐγενήθη τῇ ἐπαύριον ἔρχονται οἱ ἀλλόφυλοι ἐκδιδύσκειν τοὺς νεκροὺς, καὶ εὑρίσκουσι τὸν Σαοὺλ καὶ τοὺς τρεῖς υἱοὺς αὐτοῦ πεπτωκότας ἐπὶ τὰ ὄρη Γελβούε.

\vs{9}Καὶ ἀποστρέφουσιν αὐτὸν, καὶ ἐξέδυσαν τὰ σκεύη αὐτοῦ, καὶ ἀποστέλλουσιν αὐτὰ εἰς γῆν ἀλλοφύλων, κύκλῳ εὐαγγελίζοντες τοῖς εἰδώλοις αὐτῶν καὶ τῷ λαῷ.
\vs{10}Καὶ ἀνέθηκαν τὰ σκεύη αὐτοῦ εἰς τὸ Ἀσταρτεῖον, καὶ τὸ σῶμα αὐτοῦ κατέπηξαν ἐν τῷ τείχει Βαιθσάμ.

\vs{11}Καὶ ἀκούουσιν οἱ κατοικοῦντες Ἰαβὶς τῆς Γαλααδίτιδος ἃ ἐποίησαν οἱ ἀλλόφυλοι τῷ Σαούλ.
\vs{12}Καὶ ἀνέστησαν πᾶς ἀνὴρ δυνάμεως, καὶ ἐπορεύθησαν ὅλην τὴν νύκτα, καὶ ἔλαβον τὸ σῶμα Σαοὺλ καὶ τὸ σῶμα Ἰωνάθαν τοῦ υἱοῦ αὐτοῦ ἀπὸ τοῦ τείχους Βαιθσὰμ, καὶ φέρουσιν αὐτοὺς εἰς Ἰαβὶς, καὶ κατακαίουσιν αὐτοὺς ἐκεῖ.
\vs{13}Καὶ λαμβάνουσιν τὰ ὀστᾶ αὐτῶν, καὶ θάπτουσιν ὑπὸ τὴν ἄρουραν τὴν ἐν Ἰαβὶς, καὶ νηστεύουσιν ἑπτὰ ἡμέρας.


\def\book{ΒΑΣΙΛΕΙΩΝ Β}
\biblebook{ΒΑΣΙΛΕΙΩΝ Β}


\lettrine[lines=2, loversize=0.2, nindent=0em, findent=.25em]{\textcolor{bookheadingcolor}{Κ}}{ΑΙ} ἐγένετο μετὰ τὸ ἀποθανεῖν Σαοὺλ, καὶ Δαυὶδ ἀνέστρεψε τύπτων τὸν Ἀμαλὴκ, καὶ ἐκάθισε Δαυὶδ ἐν Σεκελὰκ ἡμέρας δύο.
\vs{2}Καὶ ἐγενήθη τῇ ἡμέρᾳ τῇ τρίτῃ, καὶ ἰδοὺ ἀνὴρ ἦλθεν ἐκ τῆς παρεμβολῆς ἐκ τοῦ λαοῦ Σαοὺλ, καὶ τὰ ἱμάτια αὐτοῦ διεῤῥωγότα, καὶ γῆ ἐπὶ τῆς κεφαλῆς αὐτοῦ· καὶ ἐγένετο ἐν τῷ εἰσελθεῖν αὐτὸν πρὸς Δαυὶδ, καὶ ἔπεσεν ἐπὶ τὴν γῆν καὶ προσεκύνησεν αὐτῷ.

\vs{3}Καὶ εἶπεν αὐτῷ Δαυὶδ, πόθεν σὺ παραγίνῃ; καὶ εἶπεν πρὸς αὐτὸν, ἐκ τῆς παρεμβολῆς Ἰσραὴλ ἐγὼ διασέσωσμαι.
\vs{4}Καὶ εἶπεν αὐτῷ Δαυὶδ, τίς ὁ λόγος οὗτος; ἀπάγγειλόν μοι· καὶ εἶπεν, ὅτι ἔφυγεν ὁ λαὸς ἐκ τοῦ πολέμου, καὶ πεπτώκασι πολλοὶ ἐκ τοῦ λαοῦ καὶ ἀπέθανον, καὶ Σαοὺλ καὶ Ἰωνάθαν ὁ υἱὸς αὐτοῦ ἀπέθανε.

\vs{5}Καὶ εἶπε Δαυὶδ τῷ παιδαρίῳ τῷ ἀπαγγελλοντι αὐτῷ, πῶς οἶδας ὅτι τέθνηκε Σαοὺλ καὶ Ἰωνάθαν ὁ υἱὸς αὐτοῦ;
\vs{6}Καὶ εἶπε τὸ παιδάριον τὸ ἀπαγγέλλον αὐτῷ, περιπτώματι περιέπεσον ἐν τῷ ὄρει τῷ Γελβουὲ, καὶ ἰδοὺ Σαοὺλ ἐπεστήρικτο ἐπὶ τὸ δόρυ αὐτοῦ, καὶ ἰδοὺ τὰ ἅρματα καὶ οἱ ἱππάρχαι συνῆψαν αὐτῷ.
\vs{7}Καὶ ἐπέβλεψεν ἐπὶ τὰ ὀπίσω αὐτοῦ, καὶ εἶδέ με, καὶ ἐκάλεσέ με· καὶ εἶπα, ἰδοὺ ἐγώ.
\vs{8}Καὶ εἶπέ μοι, τίς εἶ σύ; καὶ εἶπα, Ἀμαληκίτης ἐγώ εἰμι.
\vs{9}Καὶ εἶπε πρὸς μὲ, στῆθι δὴ ἐπάνω μου καὶ θανάτωσόν με, ὅτι κατέσχε με σκότος δεινὸν, ὅτι πᾶσα ἡ ψυχή μου ἐν ἐμοί.
\vs{10}Καὶ ἐπέστην ἐπʼ αὐτὸν καὶ ἐθανάτωσα αὐτὸν, ὅτι ᾔδειν ὅτι οὐ ζήσεται μετὰ τὸ πεσεῖν αὐτόν· καὶ ἔλαβον τὸ βασίλειον τὸ ἐπὶ τὴν κεφαλὴν αὐτοῦ, καὶ τὸν χλιδόνα τὸν ἐπὶ τοῦ βραχίονος αὐτοῦ, καὶ ἐνήνοχα αὐτὰ τῷ κυρίῳ μου ὧδε.

\vs{11}Καὶ ἐκράτησε Δαυὶδ τῶν ἱματίων αὐτοῦ, καὶ διέῤῥηξεν αὐτά· καὶ πάντες οἱ ἄνδρες οἱ μετʼ αὐτοῦ διέῤῥηξαν τὰ ἱμάτια αὐτῶν.
\vs{12}Καὶ ἐκόψαντο καὶ ἔκλαυσαν καὶ ἐνήστευσαν ἕως δείλης ἐπὶ Σαοὺλ καὶ ἐπὶ Ἰωνάθαν τὸν υἱὸν αὐτοῦ καὶ ἐπὶ τὸν λαὸν Ἰούδα καὶ ἐπὶ τὸν οἶκον Ἰσραὴλ, ὅτι ἐπλήγησαν ἐν ῥομφαίᾳ.

\vs{13}Καὶ εἶπε Δαυὶδ τῷ παιδαρίῳ τῷ ἀπαγγέλλοντι αὐτῷ, πόθεν εἶ σύ; καὶ εἶπεν, υἱὸς ανδρὸς παροίκου Ἀμαληκίτου ἐγώ εἰμι.

\vs{14}Καὶ εἶπεν αὐτῷ Δαυὶδ, πῶς οὐκ ἐφοβήθης ἐπενεγκεῖν χεῖρά σου διαφθεῖραι τὸν χριστὸν Κυρίου;
\vs{15}Καὶ ἐκάλεσε Δαυὶδ ἓν τῶν παιδαρίων αὐτοῦ, καὶ εἶπε, προσελθὼν ἀπάντησον αὐτῷ· καὶ ἐπάταξεν αὐτὸν, καὶ ἀπέθανε.
\vs{16}Καὶ εἶπε πρὸς αὐτὸν Δαυὶδ, τὸ αἷμά σου ἐπὶ τὴν κεφαλήν σου, ὅτι τὸ στόμα σου ἀπεκρίθη κατὰ σοῦ, λέγον, ὅτι ἐγὼ ἐθανάτωσα τὸν χριστὸν Κυρίου.

\vs{17}Καὶ ἐθρήνησε Δαυὶδ τὸν θρῆνον τοῦτον ἐπὶ Σαοὺλ καὶ ἐπὶ Ἰωνάθαν τὸν υἱὸν αὐτοῦ.
\vs{18}Καὶ εἶπε τοῦ διδάξαι τοὺς υἱοὺς Ἰούδα· ἰδοὺ γέγραπται ἐπὶ βιβλίου τοῦ εὐθοῦς.

\vs{19}Στήλωσον Ἰσραὴλ ὑπὲρ τῶν τεθνηκότων ἐπὶ τὰ ὕψη σου τραυματιῶν· πῶς ἔπεσαν δυνατοί;
\vs{20}Μὴ ἀναγγείλητε ἐν Γὲθ, καὶ μὴ εὐαγγελίσησθε ἐν ταῖς ἐξόδοις Ἀσκάλωνος, μή ποτε εὐφρανθῶσι θυγατέρες ἀλλοφύλων, μή ποτε ἀγαλλιάσωνται θυγατέρες τῶν ἀπεριτμήτων.
\vs{21}Ὄρη τὰ ἐν Γελβουὲ μὴ καταβάτω δρόσος καὶ μὴ ὑετὸς ἐφʼ ὑμᾶς, καὶ ἀγροὶ ἀπαρχῶν, ὅτι ἐκεῖ προσωχθίσθη θυρεὸς δυνατῶν· θυρεὸς Σαοὺλ οὐκ ἐχρίσθη ἐν ἐλαίῳ.
\vs{22}Ἀφʼ αἵματος τραυματιῶν καὶ ἀπὸ στέατος δυνατῶν τόξον Ἰωνάθαν οὐκ ἀπεστράφη κενὸν εἰς τὰ ὀπίσω, καὶ ῥομφαία Σαοὺλ οὐκ ἀνέκαμψε κενή.
\vs{23}Σαοὺλ καὶ Ἰωνάθαν οἱ ἠγαπημένοι καὶ ὡραῖοι οὐ διακεχωρισμένοι, εὐπρεπεῖς ἐν τῇ ζωῇ αὐτῶν, καὶ ἐν τῷ θανάτῳ αὐτῶν οὐ διεχωρίσθησαν· ὑπὲρ ἀετοὺς κοῦφοι, καὶ ὑπὲρ λέοντας ἐκραταιώθησαν.
\vs{24}Θυγατέρες Ἰσραὴλ ἐπὶ Σαοὺλ κλαύσατε, τὸν ἐνδιδύσκοντα ὑμᾶς κόκκινα μετὰ κόσμου ὑμῶν, τὸν ἀναφέροντα κόσμον χρυσοῦν ἐπὶ τὰ ἐνδύματα ὑμῶν.
\vs{25}Πῶς ἔπεσαν δυνατοὶ ἐν μέσῳ τοῦ πολέμου, Ἰωνάθαν ἐπὶ τὰ ὕψη σου τραυματίαι;
\vs{26}Ἀλγῶ ἐπὶ σοὶ, ἀδελφέ μου Ἰωνάθαν, ὡραιώθης μοι σφόδρα, ἐθαυμαστώθη ἡ ἀγάπησίς σου ἐμοὶ ὑπὲρ ἀγάπησιν γυναικῶν.
\vs{27}Πῶς ἔπεσαν δυνατοὶ, καὶ ἀπώλοντο σκεύη πολεμικά;

\ch{2}
Καὶ ἐγένετο μετὰ ταῦτα καὶ ἐπηρώτησε Δαυὶδ ἐν Κυρίῳ, λέγων, εἰ ἀναβῶ εἰς μίαν τῶν πόλεων Ἰούδα; καὶ εἶπε Κύριος πρὸς αὐτὸν, ἀνάβηθι· καὶ εἶπε Δαυὶδ, ποῦ ἀναβῶ; καὶ εἶπεν, εἰς Χεβρών.
\vs{2}Καὶ ἀνέβη ἐκεῖ Δαυὶδ εἰς Χεβρὼν, καὶ ἀμφότεραι αἱ γυναῖκες αὐτοῦ, Ἀχινάαμ ἡ Ἰεζραηλίτις, καὶ Ἀβιγαία ἡ γυνὴ Νάβαλ τοῦ Καρμηλίου,
\vs{3}καὶ οἱ ἄνδρες οἱ μετʼ αὐτοῦ ἕκαστος, καὶ ὁ οἶκος αὐτοῦ, καὶ κατῴκουν ἐν ταῖς πόλεσι Χεβρών.

\vs{4}Καὶ ἔρχονται ἄνδρες τῆς Ἰουδαίας, καὶ χρίουσι τὸν Δαυὶδ ἐκεῖ τοῦ βασιλεύειν ἐπὶ τὸν οἶκον Ἰούδα· Καὶ ἀπήγγειλαν τῷ Δαυὶδ, λέγοντες, ὅτι οἱ ἄνδρες Ἰαβὶς τῆς Γαλααδίτιδος ἔθαψαν τὸν Σαούλ.
\vs{5}Καὶ ἀπέστειλε Δαυὶδ ἀγγέλους πρὸς τοὺς ἡγουμένους Ἰαβὶς τῆς Γαλααδίτιδος, καὶ εἶπε πρὸς αὐτοὺς Δαυὶδ, εὐλογημένοι ὑμεῖς τῷ Κυρίῳ, ὅτι ἐποιήσατε τὸ ἔλεος τοῦτο ἐπὶ τὸν κύριον ὑμῶν, ἐπὶ Σαοὺλ τὸν χριστὸν Κυρίου, καὶ ἐθάψατε αὐτὸν καὶ Ἰωνάθαν τὸν υἱὸν αὐτοῦ.
\vs{6}Καὶ νῦν ποιῆσαι Κύριος μεθʼ ὑμῶν ἔλεος καὶ ἀλήθειαν· καί γε ἐγὼ ποιήσω μεθʼ ὑμῶν τὸ ἀγαθὸν τοῦτο, ὅτι ἐποιήσατε τὸ ῥῆμα τοῦτο.
\vs{7}Καὶ νῦν κραταιούσθωσαν αἱ χεῖρες ὑμῶν, καὶ γίνεσθε εἰς υἱοὺς δυνατοὺς, ὅτι τέθνηκεν ὁ κύριος ὑμῶν Σαοὺλ, καί γε ἐμὲ κέχρικεν ὁ οἶκος Ἰούδα ἐφʼ ἑαυτὸν εἰς βασιλέα.

\vs{8}Καὶ Ἀβεννὴρ υἱὸς Νὴρ ἀρχιστράτηγος τοῦ Σαοὺλ ἔλαβε τὸν Ἰεβοσθὲ υἱὸν Σαοὺλ, καὶ ἀνεβίβασεν αὐτὸν ἐκ τῆς παρεμβολῆς εἰς Μαναὲμ,
\vs{9}καὶ ἐβασίλευσεν αὐτὸν ἐπὶ τὴν Γαλααδίτιν, καὶ ἐπὶ τὸν Θασιρὶ, καὶ ἐπὶ τὸν Ἰεζραὴλ, καὶ ἐπὶ τὸν Ἐφραὶμ, καὶ ἐπὶ τὸν Βενιαμὶν, καὶ ἐπὶ πάντα Ἰσραήλ.
\vs{10}Τεσσαράκοντα ἐτῶν Ἰεβοσθὲ υἱὸς Σαοὺλ, ὅτε ἐβασίλευσεν ἐπὶ τὸν Ἰσραὴλ· καὶ δύο ἔτη ἐβασίλευσε, πλὴν τοῦ οἴκου Ἰούδα, οἳ ἦσαν ὀπίσω Δαυίδ.

\vs{11}Καὶ ἐγένοντο αἱ ἡμέραι ἃς Δαυὶδ ἐβασίλευσεν ἐν Χεβρὼν ἐπὶ τὸν οἶκον Ἰούδα, ἑπτὰ ἔτη καὶ μῆνας ἕξ.

\vs{12}Καὶ ἐξῆλθεν Ἀβεννὴρ υἱὸς Νὴρ, καὶ οἱ παῖδες Ἰεβοσθὲ υἱοῦ Σαοὺλ ἐκ Μαναὲμ εἰς Γαβαών·
\vs{13}Καὶ Ἰωὰβ υἱὸς Σαρουία, καὶ οἱ παῖδες Δαυὶδ ἐξῆλθον ἐκ Χεβρὼν, καὶ συναντῶσιν αὐτοῖς ἐπὶ τὴν κρήνην τὴν Γαβαὼν ἐπὶ τὸ αὐτὸ, καὶ ἐκάθισαν οὗτοι ἐπὶ τὴν κρήνην ἐντεῦθεν, καὶ οὗτοι ἐπὶ τὴν κρήνην ἐντεῦθεν.
\vs{14}Καὶ εἶπεν Ἀβεννὴρ πρὸς Ἰωὰβ, ἀναστήτωσαν δὴ τὰ παιδάρια, καὶ παιξάτωσαν ἐνώπιον ὑμῶν· καὶ εἶπεν Ἰωὰβ, ἀναστήτωσαν.
\vs{15}Καὶ ἀνέστησαν καὶ παρῆλθον ἐν ἀριθμῷ τῶν παίδων Βενιαμὶν δώδεκα τῶν Ἰεβοσθὲ υἱοῦ Σαοὺλ, καὶ δώδεκα ἐκ τῶν παίδων Δαυίδ
\vs{16}Καὶ ἐκράτησαν ἕκαστος τῇ χειρὶ τὴν κεφαλὴν τοῦ πλησίον αὐτοῦ, καὶ μάχαιρα αὐτοῦ εἰς πλευρὰν τοῦ πλησίον αὐτοῦ, καὶ πίπτουσι κατὰ τὸ αὐτό· καὶ ἐκλήθη τὸ ὄνομα τοῦ τόπου ἐκείνου, μερὶς τῶν ἐπιβούλων, ἥ ἐστιν ἐν Γαβαών.
\vs{17}Καὶ ἐγένετο ὁ πόλεμος σκληρὸς ὥστε λίαν ἐν τῇ ἡμέρᾳ ἐκείνῃ· καὶ ἔπταισεν Ἀβεννὴρ καὶ ἄνδρες Ἰσραὴλ ἐνώπιον παίδων Δαυίδ.
\vs{18}Καὶ ἐγένοντο ἐκεῖ τρεῖς υἱοὶ Σαρουία, Ἰωὰβ, καὶ Ἀβεσσὰ, καὶ Ἀσαήλ· καὶ Ἀσαὴλ κοῦφος τοῖς ποσὶν αὐτοῦ· ὡσεὶ μία δορκὰς ἐν ἀγρῷ.

\vs{19}Καὶ κατεδίωξεν Ἀσαὴλ ὀπίσω Ἀβεννὴρ, καὶ οὐκ ἐξέκλινε τοῦ πορεύεσθαι εἰς δεξιὰ οὐδὲ εἰς ἀριστερὰ κατόπισθεν Ἀβεννήρ.
\vs{20}Καὶ ἐπέβλεψεν Ἀβεννὴρ εἰς τὰ ὀπίσω αὐτοῦ, καὶ εἶπεν, εἰ σὺ εἶ αὐτὸς Ἀσαήλ; καὶ εἶπεν, ἐγώ εἰμι.
\vs{21}Καὶ εἶπεν αὐτῷ Ἀβεννὴρ, ἔκκλινον σὺ εἰς τὰ δεξιὰ ἢ εἰς τὰ ἀριστερὰ, καὶ κατάσχε σεαυτῷ ἓν τῶν παιδαρίων, καὶ λάβε σεαυτῷ τὴν πανοπλίαν αὐτοῦ· καὶ οὐκ ἠθέλησεν Ἀσαὴλ ἐκκλῖναι ἐκ τῶν ὄπισθεν αὐτοῦ.
\vs{22}Καὶ προσέθετο ἔτι Ἀβεννὴρ λέγων τῷ Ἀσαὴλ, ἀπόστηθι ἀπʼ ἐμοῦ, ἵνα μὴ πατάξω σε εἰς τὴν γῆν· καὶ πῶς ἀρῶ τὸ πρόσωπόν μου πρὸς Ἰωάβ;
\vs{23}Καὶ ποῦ ἐστι ταῦτα; ἐπίστρεφε πρὸς Ἰωὰβ τὸν ἀδελφόν σου. Καὶ οὐκ ἐβούλετο τοῦ ἀποστῆναι· καὶ τύπτει αὐτὸν Ἀβεννὴρ ἐν τῷ ὀπίσω τοῦ δόρατος ἐπὶ τὴν ψόαν, καὶ διεξῆλθε τὸ δόρυ ἐκ τῶν ὀπίσω αὐτοῦ, καὶ πίπτει ἐκεῖ καὶ ἀποθνήσκει ὑποκάτω αὐτοῦ· καὶ ἐγένετο πᾶς ὁ ἐρχόμενος ἕως τοῦ τόπου οὗ ἔπεσεν ἐκεῖ Ἀσαὴλ καὶ ἀπέθανε, καὶ ὑφίστατο.
\vs{24}Καὶ κατεδίωξεν Ἰωὰβ καὶ Ἀβεσσὰ ὀπίσω Ἀβεννὴρ, καὶ ὁ ἥλιος ἔδυνε· καὶ αὐτοὶ εἰσῆλθον ἕως τοῦ βουνοῦ Ἀμμὰν, ὅ ἐστιν ἐπὶ προσώπου Γαῒ, ὁδὸν ἔρημον Γαβαών.

\vs{25}Καὶ συναθροίζονται οἱ υἱοὶ Βενιαμὶν οἱ ὀπίσω Ἀβεννὴρ, καὶ ἐγενήθησαν εἰς συνάντησιν μίαν, καὶ ἔστησαν ἐπὶ κεφαλὴν βουνοῦ ἑνός.
\vs{26}Καὶ ἐκάλεσεν Ἀβεννὴρ Ἰωὰβ, καὶ εἶπε, μὴ εἰς νῖκος καταφάγεται ἡ ῥομφαία; ἢ οὐκ οἴδας ὅτι πικρὰ ἔσται εἰς τὰ ἔσχατα; καὶ ἕως πότε οὐ μὴ εἴπῃς τῷ λαῷ ἀποστρέφειν ἀπὸ ὄπισθε τῶν ἀδελφῶν ἡμῶν;
\vs{27}Καὶ εἶπεν Ἰωὰβ, ζῇ Κύριος, ὅτι εἰ μὴ ἐλάλησας, διότι τότε ἐκ πρωϊόθεν ἀνέβη ἂν ὁ λαὸς ἕκαστος κατόπισθε τοῦ ἀδελφοῦ αὐτοῦ.
\vs{28}Καὶ ἐσάλπισεν Ἰωὰβ τῇ σάλπιγγι, καὶ ἀπέστησαν πᾶς ὁ λαὸς, καὶ οὐ κατεδίωξαν ὀπίσω τοῦ Ἰσραὴλ, καὶ οὐ προσέθεντο ἔτι τοῦ πολεμεῖν.

\vs{29}Καὶ Ἀβεννὴρ καὶ οἱ ἄνδρες αὐτοῦ ἀπῆλθον εἰς δυσμὰς ὅλην τὴν νύκτα ἐκείνην, καὶ διέβαινον τὸν Ἰορδάνην, καὶ ἐπορεύθησαν ὅλην τὴν παρατείνουσαν, καὶ ἔρχονται εἰς τὴν παρεμβολήν.
\vs{30}Καὶ Ἰωὰβ ἀνέστρεψεν ὄπισθεν ἀπὸ τοῦ Ἀβεννὴρ, καὶ συνήθροισε πάντα τὸν λαὸν, καὶ ἐπεσκέπησαν τῶν παίδων Δαυὶδ ἐννεακαίδεκα ἄνδρες,
\vs{31}καὶ Ἀσαήλ. Καὶ οἱ παῖδες Δαυὶδ ἐπάταξαν τῶν υἱῶν Βενιαμὶν τῶν ἀνδρῶν Ἀβεννὴρ τριακοσίους ἑξήκοντα ἄνδρας παρʼ αὐτοῦ.

\vs{32}Καὶ αἴρουσι τὸν Ἀσαὴλ, καὶ θάπτουσιν αὐτὸν ἐν τῷ τάφῳ τοῦ πατρὸς αὐτοῦ ἐν Βηθλεέμ· καὶ ἐπορεύθη Ἰωὰβ καὶ οἱ ἄνδρες οἱ μετʼ αὐτοῦ ὅλην τὴν νύκτα, καὶ διέφαυσεν αὐτοῖς ἐν Χεβρών.

\ch{3}
Καὶ ἐγένετο ὁ πόλεμος ἐπὶ πολὺ ἀναμέσον τοῦ οἴκου Σαοὺλ καὶ ἀναμέσον τοῦ οἴκου Δαυίδ· καὶ ὁ οἶκος Δαυὶδ ἐπορεύετο καὶ ἐκραταιοῦτο, καὶ ὁ οἶκος Σαοὺλ ἐπορεύετο καὶ ἠσθένει.
\vs{2}Καὶ ἐτέχθησαν τῷ Δαυὶδ υἱοὶ ἐν Χεβρών· καὶ ἦν ὁ πρωτότοκος αὐτοῦ Ἀμνὼν τῆς Ἀχινόομ τῆς Ἰεζραηλίτιδος.
\vs{3}Καὶ ὁ δεύτερος αὐτοῦ Δαλουΐα τῆς Ἀβιγαίας τῆς Καρμηλίας, καὶ ὁ τρίος, Ἀβεσσαλὼμ υἱὸς Μααχὰ θυγατρὸς Θολμὶ βασιλέως Γεσσὶρ,
\vs{4}καὶ ὁ τέταρτος Ὀρνία υἱὸς Ἀγγὶθ, καὶ ὁ πέμπτος Σαφατία τῆς Ἀβιτὰλ,
\vs{5}καὶ ὁ ἕκτος Ἰεθεραὰμ τῆς Αἰγὰλ γυναικὸς Δαυίδ· οὗτοι ἐτέχθησαν τῷ Δαυὶδ ἐν Χεβρών.

\vs{6}Καὶ ἐγένετο ἐν τῷ εἶναι τὸν πόλεμον ἀναμέσον τοῦ οἴκου Σαοὺλ καὶ ἀναμέσον τοῦ οἴκου Δαυὶδ, καὶ Ἀβεννὴρ ἦν κρατῶν τοῦ οἴκου Σαούλ.
\vs{7}Καὶ τῷ Σαοὺλ παλλακὴ Ῥεσφὰ θυγάτηρ Ἰώλ· καὶ εἶπεν Ἰεβοσθὲ υἱὸς Σαοὺλ πρὸς Ἀβεννὴρ, τί ὅτι εἰσῆλθες πρὸς τὴν παλλακὴν τοῦ πατρός μου;
\vs{8}Καὶ ἐθυμώθη σφόδρα Ἀβεννὴρ περὶ τοῦ λόγου τούτου τῷ Ἰεβοσθέ. καὶ εἶπεν Ἀβεννὴρ πρὸς αὐτὸν, μὴ κεφαλὴ κυνὸς ἐγώ εἰμι; ἐποίησα σήμερον ἔλεος μετὰ τοῦ οἴκου Σαοὺλ τοῦ πατρός σου, καὶ περὶ ἀδελφῶν καὶ περὶ γνωρίμων, καὶ οὐκ ηὐτομόλησα εἰς τὸν οἶκον Δαυὶδ, καὶ ἐπιζητεῖς ἐπʼ ἐμὲ σὺ ὑπὲρ ἀδικίας γυναικὸς σήμερον;
\vs{9}Τάδε ποιήσαι ὁ Θεὸς τῷ Ἀβεννὴρ καὶ τάδε προσθείη αὐτῷ, ὅτι καθὼς ὤμοσε Κύριος τῷ Δαυὶδ, ὅτι οὕτως ποιήσω αὐτῷ ἐν τῇ ἡμέρᾳ ταύτῃ,
\vs{10}περιελεῖν τὴν βασιλείαν ἀπὸ τοῦ οἴκου Σαοὺλ, καὶ τοῦ ἀναστῆσαι τὸν θρόνον Δαυὶδ ἐπὶ Ἰσραὴλ καὶ ἐπὶ τὸν Ἰούδαν ἀπὸ Δὰν ἕως Βηρσαβεέ.
\vs{11}Καὶ οὐκ ἠδυνάσθη ἔτι Ἰεβοσθὲ ἀποκριθῆναι τῷ Ἀβεννὴρ ῥῆμα, ἀπὸ τοῦ φοβεῖσθαι αὐτόν.

\vs{12}Καὶ ἀπέστειλεν Ἀβεννὴρ ἀγγέλους πρὸς Δαυὶδ εἰς Θαιλὰμ οὗ ἦν, παραχρῆμα, λέγων, διάθου διαθήκην σου μετʼ ἐμοῦ, καὶ ἰδοὺ ἡ χείρ μου μετὰ σοῦ ἐπιστρέψαι πρὸς σὲ πάντα τὸν οἶκον Ἰσραήλ.
\vs{13}Καὶ εἶπε Δαυὶδ, καλῶς ἐγὼ διαθήσομαι πρὸς σὲ διαθήκην· πλὴν λόγον ἕνα ἐγὼ αἰτοῦμαι παρὰ σοῦ, λέγων, οὐκ ὄψει τὸ πρόσωπόν μου, ἐὰν μὴ ἀγάγῃς τὴν Μελχὸλ θυγατέρα Σαοὺλ παραγινομένου σου ἰδεῖν τὸ πρόσωπόν μου.
\vs{14}Καὶ ἐξαπέστειλε Δαυὶδ πρὸς Ἰεβοσθὲ υἱον Σαοὺλ ἀγγέλους, λέγων, ἀπόδος μοι τῆν γυναῖκά μου τὴν Μελχὸλ, ἣν ἔλαβον ἐν ἑκατὸν ἀκροβυστίαις ἀλλοφύλων.
\vs{15}Καὶ ἀπέστειλεν Ἰεβοσθὲ, καὶ ἔλαβεν αὐτὴν παρὰ τοῦ ἀνδρὸς αὐτῆς παρὰ Φαλτιὴλ υἱοῦ Σελλῆς.
\vs{16}Καὶ ἐπορεύετο ὁ ἀνὴρ αὐτῆς μετʼ αὐτῆς κλαίων ὀπίσω αὐτῆς ἕως Βαρακίμ· καὶ εἶπε πρὸς αὐτὸν Ἀβεννὴρ, πορεύου, ἀνάστρεφε· καὶ ἀνέστρεψε.

\vs{17}Καὶ εἶπεν Ἀβεννὴρ πρὸς τοὺς πρεσβυτέρους Ἰσραὴλ, λέγων, χθὲς καὶ τρίτην ἐζητεῖτε τὸν Δαυὶδ βασιλεύειν ἐφʼ ὑμᾶς.
\vs{18}Καὶ νῦν ποιήσατε, ὅτι Κύριος ἐλάλησε περὶ Δαυὶδ, λέγων, ἐν χειρὶ τοῦ δούλου μου Δαυὶδ σώσω τὸν Ἰσραὴλ ἐκ χειρὸς ἀλλοφύλων, καὶ ἐκ χειρὸς πάντων τῶν ἐχθρῶν αὐτῶν.
\vs{19}Καὶ ἐλάλησεν Ἀβεννὴρ ἐν τοῖς ὠσὶ Βενιαμίν· καὶ ἐπορεύθη Ἀβεννὴρ τοῦ λαλῆσαι εἰς τὰ ὦτα τοῦ Δαυὶδ εἰς Χεβρὼν πάντα ὅσα ἤρεσεν ἐν ὀφθαλμοῖς Ἰσραὴλ καὶ ἐν ὀφθαλμοῖς οἴκου Βενιαμίν.
\vs{20}Καὶ ἦλθεν Ἀβεννὴρ πρὸς Δαυὶδ εἰς Χεβρὼν, καὶ μετʼ αὐτοῦ εἴκοσι ἄνδρες· καὶ ἐποίησε Δαυὶδ τῷ Ἀβεννὴρ καὶ τοῖς ἀνδράσι τοῖς μετʼ αὐτοῦ πότον.
\vs{21}Καὶ εἶπεν Ἀβεννὴρ πρὸς Δαυὶδ, ἀναστήσομαι δὴ καὶ πορεύσομαι καὶ συναθροίσω πρὸς κύριόν μου τὸν βασιλέα πάντα Ἰσραὴλ· καὶ διαθήσομαι μετʼ αὐτοῦ διαθήκην, καὶ βασιλεύσεις ἐπὶ πᾶσιν οἷς ἐπιθυμεῖ ἡ ψυχή σου. Καὶ ἀπέστειλε Δαυὶδ τὸν Ἀβεννὴρ, καὶ ἐπορεύθη ἐν εἰρήνῃ.

\vs{22}Καὶ ἰδοὺ οἱ παῖδες Δαυὶδ καὶ Ἰωὰβ παρεγένοντο ἐκ τῆς ἐξοδίας, καὶ σκῦλα πολλὰ ἔφερον μεθʼ ἑαυτῶν· καὶ Ἀβεννὴρ οὐκ ἦν μετὰ Δαυὶδ εἰς Χεβρὼν, ὅτι ἀπεστάλκει αὐτὸν, καὶ ἀπεληλύθει ἐν εἰρήνῃ.
\vs{23}Καὶ Ἰωὰβ καὶ πᾶσα ἡ στρατιὰ αὐτοῦ ἤλθοσαν, καὶ ἀπηγγέλη τῷ Ἰωὰβ, λέγοντες, ἥκει Ἀβεννὴρ υἱὸς Νὴρ πρὸς Δαυὶδ, καὶ ἀπέσταλκεν αὐτὸν, καὶ ἀπῆλθεν ἐν εἰρήνῃ.
\vs{24}Καὶ εἰσῆλθεν Ἰωὰβ πρὸς τὸν βασιλέα, καὶ εἶπε, τί τοῦτο ἐποίησας; ἰδοὺ ἦλθεν Ἀβεννὴρ πρὸς σὲ, καὶ ἱνατί ἐξαπέσταλκας αὐτὸν, καὶ ἀπελήλυθεν ἐν εἰρήνῃ;
\vs{25}Ἢ οὐκ οἶδας τὴν κακίαν Ἀβεννὴρ υἱοῦ Νὴρ, ὅτι ἀπατῆσαί σε παρεγένετο, καὶ γνῶναι τὴν ἔξοδόν σου καὶ τὴν εἴσοδόν σου, καὶ γνῶναι ἅπαντα ὅσα σὺ ποιεῖς;

\vs{26}Καὶ ἀνέστρεψεν Ἰωὰβ ἀπὸ τοῦ Δαυὶδ, καὶ ἀπέστειλεν ἀγγέλους πρὸς Ἀβεννὴρ ὀπίσω, καὶ ἐπιστρέφουσιν αὐτὸν ἀπὸ τοῦ φρέατος τοῦ Σεειρὰμ· καὶ Δαυὶδ οὐκ ᾔδει.
\vs{27}Καὶ ἐπέστρεψε τὸν Ἀβεννὴρ εἰς Χεβρὼν, καὶ ἐξέκλινεν αὐτὸν Ἰωὰβ ἐκ πλαγίων τῆς πύλης λαλῆσαι πρὸς αὐτὸν, ἐνεδρεύων· καὶ ἐπάταξεν αὐτὸν ἐκεῖ εἰς τὴν ψόαν, καὶ ἀπέθανεν ἐν τῷ αἵματι Ἀσαὴλ τοῦ ἀδελφοῦ Ἰωάβ.

\vs{28}Καὶ ἤκουσε Δαυὶδ μετὰ ταῦτα, καὶ εἶπεν, ἀθῶός εἰμι ἐγὼ καὶ ἡ βασιλεία μου ἀπὸ Κυρίου καὶ ἕως αἰῶνος ἀπὸ τῶν αἱμάτων Ἀβεννὴρ υἱοῦ Νήρ.
\vs{29}Καταντησάτωσαν ἐπὶ κεφαλὴν Ἰωὰβ καὶ ἐπὶ πάντα τὸν οἶκον τοῦ πατρὸς αὐτοῦ, καὶ μὴ ἐκλείποι ἐκ τοῦ οἴκου Ἰωὰβ γονοῤῥυὴς, καὶ λεπρὸς, καὶ κρατῶν σκυτάλης, καὶ πίπτων ἐν ῥομφαίᾳ, καὶ ἐλασσούμενος ἄρτοις.
\vs{30}Ἰωὰβ δὲ καὶ Ἀβεσσὰ ὁ ἀδελφὸς αὐτοῦ διαπαρετηροῦντο τὸν Ἀβεννὴρ, ἀνθʼ ὧν ἐθανάτωσε τὸν Ἀσαὴλ τὸν ἀδελφὸν αὐτῶν ἐν Γαβαὼν, ἐν τῷ πολέμῳ.

\vs{31}Καὶ εἶπε Δαυὶδ πρὸς Ἰωὰβ καὶ πρὸς πάντα τὸν λαὸν τὸν μετʼ αὐτοῦ, διαῤῥήξατε τὰ ἱμάτια ὑμῶν, καὶ περιζώσασθε σάκκους, καὶ κόπτεσθε ἐνώπιον Ἀβεννήρ· καὶ ὁ βασιλεὺς Δαυὶδ ἐπορεύετο ὀπίσω τῆς κλίνης.
\vs{32}Καὶ θάπτουσι τὸν Ἀβεννὴρ ἐν Χεβρών· καὶ ᾖρεν ὁ βασιλεὺς τὴν φωνὴν αὐτοῦ καὶ ἔκλαυσεν ἐπὶ τοῦ τάφου αὐτοῦ, καὶ ἔκλαυσε πᾶς ὁ λαὸς ἐπὶ Ἀβεννήρ.

\vs{33}Καὶ ἐθρήνησεν ὁ βασιλεὺς ἐπὶ Ἀβεννὴρ, καὶ εἶπεν, εἰ κατὰ τὸν θάνατον Νάβαλ ἀποθανεῖται Ἀβεννήρ;
\vs{34}Αἱ χεῖρές σου οὐκ ἐδέθησαν, οἱ πόδες σου οὐκ ἐν πέδαις· οὐ προσήγαγεν ὡς Νάβαλ, ἐνώπιον υἱῶν ἀδικίας ἔπεσας· καὶ συνήχθη πᾶς ὁ λαὸς τοῦ κλαῦσαι αὐτόν.
\vs{35}Καὶ ἦλθε πᾶς ὁ λαὸς περιδειπνῆσαι τὸν Δαυὶδ ἄρτοις ἔτι οὔσης ἡμέρας· καὶ ὤμοσε Δαυὶδ, λέγων, τάδε ποιήσαι μοι ὁ Θεὸς καὶ τάδε προσθείη, ὅτι ἐὰν μὴ δύῃ ὁ ἥλιος, οὐ μὴ γεύσομαι ἄρτου ἢ ἀπὸ παντὸς τινός.
\vs{36}Καὶ ἔγνω πᾶς ὁ λαὸς, καὶ ἤρεσεν ἐνώπιον αὐτῶν πάντα ὅσα ἐποίησεν ὁ βασιλεὺς ἐνώπιον τοῦ λαοῦ.
\vs{37}Καὶ ἔγνω πᾶς ὁ λαὸς καὶ πᾶς Ἰσραὴλ ἐν τῇ ἡμέρᾳ ἐκείνῃ, ὅτι οὐκ ἐγένετο παρὰ τοῦ βασιλέως θανατῶσαι τὸν Ἀβεννὴρ υἱὸν Νήρ.

\vs{38}Καὶ εἶπεν ὁ βασιλεὺς πρὸς τοὺς παῖδας αὐτοῦ, οὐκ οἴδατε, ὅτι ἡγούμενος μέγας πέπτωκεν ἐν τῇ ἡμέρᾳ ταύτῃ ἐν τῷ Ἰσραήλ;
\vs{39}Καὶ ὅτι ἐγώ εἰμι συγγενὴς σήμερον, καὶ καθεσταμένος ὑπὸ βασιλέως; οἱ δὲ ἄνδρες οὗτοι υἱοὶ Σαρουίας σκληρότεροί μου εἰσίν· ἀποδῷ Κύριος τῷ ποιοῦντι τὰ πονηρὰ κατὰ τὴν κακίαν αὐτοῦ.

\ch{4}
Καὶ ἤκουσεν Ἰεβοσθὲ υἱὸς Σαοὺλ, ὅτι τέθνηκεν Ἀβεννὴρ υἱὸς Νὴρ ἐν Χεβρὼν, καὶ ἐξελύθησαν αἱ χεῖρες αὐτοῦ, καὶ πάντες οἱ ἄνδρες Ἰσραὴλ παρείθησαν.
\vs{2}Καὶ δύο ἄνδρες ἡγούμενοι συστρεμμάτων τῷ Ἰεβοσθὲ υἱῷ Σαούλ· ὄνομα τῷ ἑνὶ Βαανὰ, καὶ ὄνομα τῷ δευτέρῳ Ῥηχὰβ, υἱοὶ Ῥεμμὼν τοῦ Βηρωθαίου ἐκ τῶν υἱῶν Βενιαμίν· ὅτι Βηρὼθ ἐλογίζετο τοῖς υἱοῖς Βενιαμίν·
\vs{3}Καὶ ἀπέδρασαν οἱ Βηρωθαῖοι εἰς Γεθαὶμ, καὶ ἦσαν ἐκεῖ παροικοῦντες ἕως τῆς ἡμέρας ταύτης.

\vs{4}Καὶ τῷ Ἰωνάθαν υἱῷ Σαοὺλ υἱὸς πεπληγὼς τοὺς πόδας υἱὸς ἐτῶν πέντε, καὶ οὗτος ἐν τῷ ἐλθεῖν τὴν ἀγγελίαν Σαοὺλ καὶ Ἰωνάθαν τοῦ υἱοῦ αὐτοῦ ἐξ Ἰεζραὴλ, καὶ ᾖρεν αὐτὸν ἡ τιθηνὸς αὐτοῦ, καὶ ἔφυγε· καὶ ἐγένετο ἐν τῷ σπεύδειν αὐτὸν καὶ ἀναχωρεῖν, καὶ ἔπεσε καὶ ἐχωλάνθη, καὶ ὄνομα αὐτῷ Μεμφιβοσθέ.

\vs{5}Καὶ ἐπορεύθησαν υἱοὶ Ῥεμμὼν τοῦ Βηρωθαίου Ῥηχὰβ καὶ Βαανὰ, καὶ εἰσῆλθον εν τῷ καύματι τῆς ἡμέρας εἰς οἶκον Ἰεβοσθὲ, καὶ αὐτὸς ἐκάθευδεν ἐν τῇ κοίτῃ τῆς μεσημβρίας.
\vs{6}Καὶ ἰδοὺ ἡ θυρωρὸς τοῦ οἴκου ἐκάθαιρε πυροὺς, καὶ ἐνύσταξε καὶ ἐκάθευδε· καὶ Ῥηχὰβ καὶ Βαανὰ οἱ ἀδελφοὶ διέλαθον,
\vs{7}καὶ εἰσῆλθον εἰς τὸν οἶκον· καὶ Ἰεβοσθὲ ἐκάθευδεν ἐπὶ τῆς κλίνης αὐτοῦ ἐν τῷ κοιτῶνι αὐτοῦ· καὶ τύπτουσιν αὐτὸν, καὶ θανατοῦσιν αὐτὸν, καὶ ἀφαιροῦσι τὴν κεφαλὴν αὐτοῦ· καὶ ἔλαβον τὴν κεφαλὴν αὐτοῦ, καὶ ἀπῆλθον ὁδὸν τὴν κατὰ δυσμὰς ὅλην τὴν νύκτα.

\vs{8}Καὶ ἤνεγκαν τὴν κεφαλὴν Ἰεβοσθὲ τῷ Δαυὶδ εἰς Χεβρὼν, καὶ εἶπον πρὸς τὸν βασιλέα, ἰδοὺ ἡ κεφαλὴ Ἰεβοσθὲ υἱοῦ Σαοὺλ τοῦ ἐχθροῦ σου, ὃς ἐζήτει τὴν ψυχήν σου, καὶ ἔδωκε Κύριος τῷ κυρίῳ βασιλεῖ ἐκδίκησιν τῶν ἐχθρῶν αὐτοῦ, ὡς ἡ ἡμέρα αὕτη, ἐκ Σαοὺλ τοῦ ἐχθροῦ σου καὶ ἐκ τοῦ σπέρματος αὐτοῦ.

\vs{9}Καὶ ἀπεκρίθη Δαυὶδ τῷ Ῥηχὰβ καὶ τῷ Βαανὰ ἀδελφῷ αὐτοῦ υἱοῖς Ῥεμμὼν τοῦ Βηρωθαίου, καὶ εἶπεν αὐτοῖς, ζῇ Κύριος, ὃς ἐλυτρώσατο τὴν ψυχήν μου ἐκ πάσης θλίψεως, ὅτι ὁ ἀπαγγείλας μοι ὅτι τέθνηκε Σαοὺλ,
\vs{10}καὶ αὐτὸς ἦν ὡς εὐαγγελιζόμενος ἐνώπιόν μου, καὶ κατέσχον αὐτὸν καὶ ἀπέκτεινα αὐτὸν ἐν Σεκελὰκ, ᾧ ἔδει με δοῦναι εὐαγγέλια.
\vs{11}Καὶ νῦν ἄνδρες πονηροὶ ἀπεκτάγκασιν ἄνδρα δίκαιον ἐν τῷ οἴκῳ αὐτοῦ ἐπὶ τῆς κοίτης αὐτοῦ· καὶ νῦν ἐκζητήσω τὸ αἷμα αὐτοῦ ἐκ χειρὸς ὑμῶν, καὶ ἐξολοθρεύσω ὑμᾶς ἐκ τῆς γῆς.
\vs{12}Καὶ ἐνετείλατο Δαυὶδ τοῖς παιδαρίοις αὐτοῦ, καὶ ἀποκτείνουσιν αὐτούς, καὶ κολοβοῦσι τὰς χεῖρας αὐτῶν καὶ τοὺς πόδας αὐτῶν, καὶ ἐκρέμασαν αὐτοὺς ἐπι τῆς κρήνης ἐν Χεβρὼν, καὶ τὴν κεφαλὴν Ἰεβοσθὲ ἔθαψαν ἐν τῷ τάφῳ Ἀβεννὴρ υἱοῦ Νήρ.

\ch{5}
Καὶ παραγίνονται πᾶσαι αἱ φυλαὶ Ἰσραὴλ πρὸς Δαυὶδ εἰς Χεβρὼν, καὶ εἶπαν αὐτῷ, ἰδοὺ ὀστᾶ σου, καὶ σάρκες σου ἡμεῖς.
\vs{2}Καὶ ἐχθὲς καὶ τρίτην ὄντος Σαοὺλ βασιλέως ἐφʼ ἡμῖν, σὺ ἦσθα ὁ ἐξάγων καὶ εἰσάγων τὸν Ἰσραήλ· καὶ εἶπε Κύριος πρὸς σὲ, σὺ ποιμανεῖς τὸν λαόν μου τὸν Ἰσραὴλ, καὶ σὺ ἔσῃ εἰς ἡγούμενον ἐπὶ τὸν λαόν μου Ἰσραήλ.
\vs{3}Καὶ ἔρχονται πάντες οἱ πρεσβύτεροι Ἰσραὴλ πρὸς τὸν βασιλέα εἰς Χεβρὼν, καὶ διέθετο αὐτοῖς ὁ βασιλεὺς Δαυὶδ διαθήκην ἐν Χεβρὼν ἐνώπιον Κυρίου· καὶ χρίουσι τὸν Δαυὶδ εἰς βασιλέα ἐπὶ πάντα Ἰσραήλ.
\vs{4}Υἱὸς τριάκοντα ἐτῶν Δαυὶδ ἐν τῷ βασιλεῦσαι αὐτόν, καὶ τεσσαράκοντα ἔτη ἐβασίλευσεν.
\vs{5}Ἑπτὰ ἔτη καὶ μῆνας ἓξ ἐβασίλευσεν ἐν Χεβρὼν ἐπὶ τὸν Ἰούδαν, καὶ τριάκοντα τρία ἔτη ἐβασίλευσεν ἐπὶ πάντα Ἰσραὴλ καὶ Ἰούδαν ἐν Ἱερουσαλήμ.

\vs{6}Καὶ ἀπῆλθε Δαυὶδ καὶ οἱ ἄνδρες αὐτοῦ εἰς Ἰερουσαλὴμ πρὸς τὸν Ἱεβουσαῖον τὸν κατοικοῦντα τὴν γῆν· καὶ ἐῤῥέθη τῷ Δαυὶδ, οὐκ εἰσελεύσῃ ὧδε, ὅτι ἀντέστησαν οἱ τυφλοὶ καὶ οἱ χωλοὶ, λέγοντες, ὅτι οὐκ εἰσελεύσεται Δαυὶδ ὧδε.
\vs{7}Καὶ προκατελάβετο Δαυὶδ τὴν περιοχὴν Σιών· αὕτη ἡ πόλις τοῦ Δαυίδ.
\vs{8}Καὶ εἶπε Δαυὶδ τῇ ἡμέρᾳ ἐκείνῃ, πᾶς τύπτων Ἰεβουσαῖον, ἁπτέσθω ἐν παραξιφίδι καὶ τοὺς χωλοὺς καὶ τοὺς τυφλοὺς καὶ τοὺς μισοῦντας τὴν ψυχὴν Δαυίδ. Διὰ τοῦτο ἐροῦσι, τυφλοὶ καὶ χωλοὶ οὐκ εἰσελεύσοντάι εἰς οἶκον Κυρίου.
\vs{9}Καὶ ἐκάθισε Δαυὶδ ἐν τῇ περιοχῇ, καὶ ἐκλήθη αὕτη ἡ πόλις Δαυίδ· καὶ ᾠκοδόμησεν αὐτὴν πόλιν κύκλῳ ἀπὸ τῆς ἄκρας, καὶ τὸν οἶκον αὐτοῦ.
\vs{10}Καὶ διεπορεύετο Δαυὶδ πορευόμενος καὶ μεγαλυνόμενος, καὶ Κύριος παντοκράτωρ μετʼ αὐτοῦ.

\vs{11}Καὶ ἀπέστειλε Χειρὰμ βασιλεὺς Τύρου ἀγγέλους πρὸς Δαυὶδ, καὶ ξύλα κέδρινα, καὶ τέκτονας ξύλων, καὶ τέκτονας λίθων, καὶ ᾠκοδόμησαν οἶκον τῷ Δαυίδ.
\vs{12}Καὶ ἔγνω Δαυὶδ, ὅτι ἡτοίμασεν αὐτὸν Κύριος εἰς βασιλέα ἐπὶ Ἰσραὴλ, καὶ ὅτι ἐπῄρθη ἡ βασιλεία αὐτοῦ διὰ τὸν λαὸν αὐτοῦ Ἰσραήλ.

\vs{13}Καὶ ἔλαβε Δαυὶδ ἔτι γυναῖκας καὶ παλλακὰς ἐξ Ἰερουσαλὴμ, μετὰ τὸ ἐλθεῖν αὐτὸν ἐκ Χεβρών· καὶ ἐγένοντο τῷ Δαυὶδ ἔτι υἱοὶ καὶ θυγατέρες.
\vs{14}Καὶ ταῦτα τὰ ὀνόματα τῶν γεννηθέντων αὐτῷ ἐν Ἱερουσαλήμ· Σαμμοὺς, καὶ Σωβὰβ, καὶ Νάθαν, καὶ Σαλωμὼν,
\vs{15}καὶ Ἐβεὰρ, καὶ Ἐλισουὲ, καὶ Ναφὲκ, καὶ Ἰεφιὲς,
\vs{16}καὶ Ἐλισαμὰ, καὶ Ἐλιδαὲ, καὶ Ἐλιφαλὰθ,
\vs{16a}Σαμαὲ, Ἰεσσιβὰθ, Νάθαν, Γαλαμαὰν, Ἰεβαὰρ, Θεησοῦς, Ἐλιφαλὰτ, Ναγὲδ, Ναφὲκ, Ἰανάθαν, Λεασαμὺς, Βααλιμὰθ, Ἐλιφαάθ.

\vs{17}Καὶ ἤκουσαν οἱ ἀλλόφυλοι ὅτι κέχρισται Δαυὶδ βασιλεὺς ἐπὶ Ἰσραὴλ, καὶ ἀνέβησαν πάντες οἱ ἀλλόφυλοι ζητεῖν τὸν Δαυίδ· καὶ ἤκουσε Δαυὶδ, καὶ κατέβη εἰς τὴν περιοχήν.
\vs{18}Καὶ οἱ ἀλλόφυλοι παραγίνονται, καὶ συνέπεσαν εἰς τὴν κοιλάδα τῶν Τιτάνων.

\vs{19}Καὶ ἠρώτησε Δαυὶδ διὰ Κυρίου, λέγων, εἰ ἀναβῶ πρὸς τοὺς ἀλλοφύλους, καὶ παραδώσεις αὐτοὺς εἰς τὰς χεῖράς μου; καὶ εἶπε Κύριος πρὸς Δαυὶδ, ἀνάβαινε, ὅτι παραδιδοὺς παραδώσω τοὺς ἀλλοφύλους εἰς τὰς χεῖράς σου.
\vs{20}Καὶ ἦλθε Δαυὶδ ἐκ τῶν ἐπάνω διακοπῶν, καὶ ἔκοψε τοὺς ἀλλοφύλους ἐκεῖ· καὶ εἶπε Δαυὶδ, διέκοψε Κύριος τοὺς ἐχθροὺς ἀλλοφύλους ἐνώπιον ἐμοῦ, ὡς διακόπτεται ὕδατα· διὰ τοῦτο ἐκλήθη τὸ ὄνομα τοῦ τόπου ἐκείνου, Ἐπάνω διακοπῶν.
\vs{21}Καὶ καταλιμπάνουσιν ἐκεῖ τοὺς θεοὺς αὐτῶν, καὶ ἐλάβοσαν αὐτοὺς Δαυὶδ καὶ οἱ ἄνδρες οἱ μετʼ αὐτοῦ.

\vs{22}Καὶ προσέθεντο ἔτι ἀλλόφυλοι τοῦ ἀναβῆναι, καὶ συνέπεσαν ἐν τῇ κοιλάδι τῶν Τιτάνων.
\vs{23}Καὶ ἐπηρώτησε Δαυὶδ διὰ Κυρίου· καὶ εἶπε Κύριος, οὐκ ἀναβήσῃ εἰς συνάντησιν αὐτῶν, ἀποστρέφου ἀπʼ αὐτῶν, καὶ παρέσῃ αὐτοῖς πλησίον τοῦ κλαυθμῶνος.
\vs{24}Καὶ ἔσται ἐν τῷ ἀκοῦσαί σε τὴν φωνὴν τοῦ συγκλεισμοῦ ἀπὸ τοῦ ἄλσους τοῦ κλαυθμῶνος, τότε καταβήσῃ πρὸς αὐτούς, ὅτι τότε ἐξελεύσεται Κύριος ἔμπροσθέν σου κόπτειν ἐν τῷ πολέμῳ τῶν ἀλλοφύλων.
\vs{25}Καὶ ἐποιήσε Δαυὶδ καθὼς ἐνετείλατο αὐτῷ Κύριος, καὶ ἐπάταξε τοὺς ἀλλοφύλους ἀπὸ Γαβαὼν ἕως τῆς γῆς Γαζηρά.

\ch{6}
Καὶ συνήγαγεν ἔτι Δαυὶδ πάντα νεανίαν ἐξ Ἰσραὴλ, ὡς ἑβδομήκοντα χιλιάδας.
\vs{2}Καὶ ἀνέστη καὶ ἐπορεύθη Δαυὶδ καὶ πᾶς ὁ λαὸς ὁ μετʼ αὐτοῦ καὶ ἀπὸ τῶν ἀρχόντων Ἰούδα ἐν ἀναβάσει τοῦ ἀναγαγεῖν ἐκεῖθεν τὴν κιβωτὸν τοῦ Θεοῦ, ἐφʼ ἣν ἐπεκλήθη τὸ ὄνομα τοῦ Κυρίου τῶν δυνάμεων καθημένου ἐπὶ τῶν χερουβὶν ἐπʼ αὐτῆς.

\vs{3}Καὶ ἐπεβίβασεν τὴν κιβωτὸν Κυρίου ἐφʼ ἅμαξαν καινὴν, καὶ ᾖραν αὐτὴν ἐξ οἴκου Ἀμιναδὰβ τοῦ ἐν τῷ βουνῷ· καὶ Ὀζὰ καὶ οἱ ἀδελφοὶ αὐτοῦ υἱοὶ Ἀμιναδὰβ ἦγαν τὴν ἅμαξαν σὺν τῇ κιβωτῷ.
\vs{4}Καὶ οἱ ἀδελφοὶ αὐτοῦ ἐπορεύοντο ἔμπροσθεν τῆς κιβωτοῦ.
\vs{5}Καὶ Δαυὶδ καὶ υἱοὶ Ἰσραὴλ παίζοντες ἐνώπιον Κυρίου ἐν ὀργάνοις ἡρμοσμένοις ἐν ἰσχύϊ, καὶ ἐν ᾠδαῖς, καὶ ἐν κινύραις, καὶ ἐν νάβλαις, καὶ ἐν τυμπάνοις, καὶ ἐν κυμβάλοις, καὶ ἐν αὐλοῖς.

\vs{6}Καὶ παραγίνονται ἕως ἅλω Ναχώρ· καὶ ἐξέτεινεν Ὀζὰ τὴν χεῖρα αὐτοῦ ἐπὶ τὴν κιβωτὸν τοῦ Θεοῦ κατασχεῖν αὐτὴν, καὶ ἐκράτησεν αὐτὴν, ὅτι περιέσπασεν αὐτὴν ὁ μόσχος.
\vs{7}Καὶ ἐθυμώθη ὀργῇ Κύριος τῷ Ὀζᾷ, καὶ ἔπαισεν αὐτὸν ἐκεῖ ὁ Θεός, καὶ ἀπέθανεν ἐκεῖ παρὰ τὴν κιβωτὸν τοῦ Κυρίου ἐνώπιον τοῦ Θεοῦ.
\vs{8}Καὶ ἠθύμησε Δαυὶδ ὑπὲρ οὗ διέκοψε Κύριος διακοπὴν ἐν τῷ Ὀζᾷ, καὶ ἐκλήθη ὁ τόπος ἐκεῖνος, διακοπὴ Ὀζᾶ, ἕως τῆς ἡμέρας ταύτης.
\vs{9}Καὶ ἐφοβήθη Δαυὶδ τὸν Κύριον ἐν τῇ ἡμέρᾳ ἐκείνῃ, λέγων, πῶς εἰσελεύσεται πρὸς μὲ ἡ κιβωτὸς Κυρίου;
\vs{10}Καὶ οὐκ ἐβούλετο Δαυὶδ τοῦ ἐκκλῖναι πρὸς αὐτὸν τὴν κιβωτὸν διαθήκης Κυρίου εἰς τὴν πόλιν Δαυίδ· καὶ ἀπέκλινεν αὐτὴν Δαυὶδ εἰς οἶκον Ἀβεδδαρὰ τοῦ Γεθαίου.
\vs{11}Καὶ ἐκάθισεν ἡ κιβωτὸς τοῦ Κυρίου εἰς οἶκον Ἀβεδδαρὰ τοῦ Γεθαίου μῆνας τρεῖς· καὶ εὐλόγησε Κύριος ὅλον τὸν οἶκον Ἀβεδδαρὰ, καὶ πάντα τὰ αὐτοῦ.

\vs{12}Καὶ ἀπηγγέλη τῷ βασιλεῖ Δαυὶδ, λέγοντες, εὐλόγησε Κύριος τὸν οἶκον Ἀβεδδαρὰ, καὶ πάντα τὰ αὐτοῦ, ἕνεκα τῆς κιβωτοῦ τοῦ Θεοῦ· καὶ ἐπορεύθη Δαυὶδ, καὶ ἀνήγαγε τὴν κιβωτὸν τοῦ Κυρίου ἐκ τοῦ οἴκου Ἀβεδδαρὰ εἰς τὴν πόλιν Δαυὶδ ἐν εὐφροσύνῃ.
\vs{13}Καὶ ἦσαν μετʼ αὐτοῦ αἴροντες τὴν κιβωτὸν ἑπτὰ χοροὶ, καὶ θύμα μόσχος καὶ ἄρνες.
\vs{14}Καὶ Δαυὶδ ἀνεκρούετο ἐν ὀργάνοις ἡρμοσμένοις ἐνώπιον Κυρίου, καὶ ὁ Δαυὶδ ἐνδεδυκὼς στολὴν ἔξαλλον.
\vs{15}Καὶ Δαυὶδ καὶ πᾶς ὁ οἶκος Ἰσραὴλ ἀνήγαγον τὴν κιβωτὸν Κυρίου μετὰ κραυγῆς καὶ μετὰ φωνῆς σάλπιγγος.

\vs{16}Καὶ ἐγένετο τῆς κιβωτοῦ παραγινομένης ἕως πόλεως Δαυίδ, καὶ Μελχὸλ ἡ θυγάτηρ Σαοὺλ διέκυπτε διὰ τῆς θυρίδος, καὶ εἶδε τὸν βασιλέα Δαυὶδ ὀρχούμενον, καὶ ἀνακρουόμενον ἐνώπιον Κυρίου, καὶ ἐξουδένωσεν αὐτὸν ἐν τῇ καρδίᾳ αὐτῆς.

\vs{17}Καὶ φέρουσι τὴν κιβωτὸν τοῦ Κυρίου, καὶ ἀνέθηκαν αὐτὴν εἰς τὸν τόπον αὐτῆς εἰς μέσον τῆς σκηνῆς, ἧς ἔπηξεν αὐτῇ Δαυὶδ· καὶ ἀνήνεγκε Δαυιδ ὁλοκαυτώματα ἐνώπιον Κυρίου, εἰρηνικάς.
\vs{18}Καὶ συνετέλεσε Δαυὶδ συναναφέρων τὰς ὁλοκαυτώσεις καὶ τὰς εἰρηνικὰς, καὶ εὐλόγησε τὸν λαὸν ἐν ὀνόματι Κυρίου τῶν δυνάμεων.
\vs{19}Καὶ διεμέρισε παντὶ τῷ λαῷ εἰς πᾶσαν τὴν δύναμιν τοῦ Ἰσραὴλ ἀπὸ Δὰν ἕως Βηρσαβεὲ, καὶ ἀπὸ ἀνδρὸς ἕως γυναικὸς, ἑκάστῳ κολλυρίδα ἄρτου, καὶ ἐσχαρίτην, καὶ λάγανον ἀπὸ τηγάνου· καὶ ἀπῆλθε πᾶς ὁ λαὸς ἕκαστος εἰς τὸν οἶκον αὐτοῦ.

\vs{20}Καὶ ἐπέστρεψε Δαυὶδ εὐλογῆσαι τὸν οἶκον αὐτοῦ· καὶ ἐξῆλθε Μελχὸλ ἡ θυγάτηρ Σαοὺλ εἰς ἀπάντησιν Δαυὶδ, καὶ εὐλόγησεν αὐτὸν, καὶ εἶπε, τί δεδόξασται σήμερον ὁ βασιλεὺς Ἰσραὴλ, ὃς ἀπεκαλύφθη σήμερον ἐν ὀφθαλμοῖς παιδισκῶν τῶν δούλων ἑαυτοῦ, καθὼς ἀποκαλύπτεται ἀποκαλυφθεὶς εἷς τῶν ὀρχουμένων;
\vs{21}Καὶ εἶπε Δαυὶδ πρὸς Μελχόλ, ἐνώπιον Κυρίου ὀρχήσομαι· εὐλογητὸς Κύριος ὃς ἐξελέξατό με ὑπὲρ τὸν πατέρα σου καὶ ὑπὲρ πάντα τὸν οἶκον αὐτοῦ, του καταστῆσαί με εἰς ἡγούμενον ἐπὶ τὸν λαὸν αὐτοῦ ἐπὶ τὸν Ἰσραὴλ·
\vs{22}καὶ παίξομαι καὶ ὀρχήσομαι ἐνώπιον Κυρίου, καὶ ἀποκαλυφθήσομαι ἔτι οὕτως, καὶ ἔσομαι ἀχρεῖος ἐν ὀφθαλμοῖς σου, καὶ μετὰ τῶν παιδισκῶν, ὧν εἶπάς με μὴ δοξασθῆναι.
\vs{23}Καὶ τῇ Μελχὸλ θυγατρὶ Σαοὺλ οὐκ ἐγένετο παιδίον ἕως τῆς ἡμέρας τοῦ ἀποθανεῖν αὐτήν.

\ch{7}
Καὶ ἐγένετο ὅτε ἐκάθισεν ὁ βασιλεὺς ἐν τῷ οἴκῳ αὐτοῦ, καὶ Κύριος κατεκληρονόμησεν αὐτὸν κύκλῳ ἀπὸ πάντων τῶν ἐχθρῶν αὐτοῦ τῶν κύκλῳ,
\vs{2}καὶ εἶπεν ὁ βασιλεὺς πρὸς Νάθαν τὸν προφήτην, ἰδοὺ δὴ ἐγὼ κατοικῶ ἐν οἴκῳ κεδρίνῳ, καὶ ἡ κιβωτὸς τοῦ Θεοῦ κάθηται ἐν μέσῳ τῆς σκηνῆς.
\vs{3}Καὶ εἶπε Νάθαν πρὸς τὸν βασιλέα, πάντα ὅσα ἂν ἐν τῇ καρδίᾳ σου, βάδιζε καὶ ποίει, ὅτι Κύριος μετὰ σοῦ.

\vs{4}Καὶ ἐγένετο τῇ νυκτὶ ἐκείνῃ, καὶ ἐγένετο ῥῆμα Κυρίου πρὸς Νάθαν, λέγων, πορεύου,
\vs{5}καὶ εἶπον πρὸς τὸν δοῦλόν μου Δαυὶδ, τάδε λέγει Κύριος, οὐ σὺ οἰκοδομήσεις μοι οἶκον τοῦ κατοικῆσαί με.
\vs{6}Ὅτι οὐ κατῴκηκα ἐν οἴκῳ ἀφʼ ἧς ἡμέρας ἀνήγαγον τοὺς υἱοὺς Ἰσραὴλ ἐξ Αἰγύπτου ἕως τῆς ἡμέρας ταύτης, καὶ ἤμην ἐυπεριπατῶν ἐν καταλύματι καὶ ἐν σκηνῇ,
\vs{7}ἐν πᾶσιν οἷς διῆλθον ἐν παντὶ Ἰσραήλ· εἰ λαλῶν ἐλάλησα πρὸς μίαν φυλὴν τοῦ Ἰσραὴλ, ᾧ ἐνετειλάμην ποιμαίνειν τὸν λαόν μου Ἰσραὴλ, λέγων, ἱνατί οὐκ ᾠκοδομήκατέ μοι οἶκον κέδρινον;

\vs{8}Καὶ νῦν τάδε ἐρεῖς τῷ δούλῳ μου Δαυὶδ, τάδε λέγει Κύριος παντοκράτωρ, ἔλαβόν σε ἐκ τῆς μάνδρας τῶν προβάτων, τοῦ εἶναί σε εἰς ἡγούμενον ἐπὶ τὸν λαόν μου ἐπὶ τὸν Ἰσραὴλ,
\vs{9}καὶ ἤμην μετὰ σοῦ ἐν πᾶσιν οἷς ἐπορεύου, καὶ ἐξωλόθρευσα πάντας τοὺς ἐχθρούς σου ἀπὸ προσώπου σου, καὶ ἐποίησά σε ὀνομαστὸν κατὰ τὸ ὄνομα τῶν μεγάλων τῶν ἐπὶ τῆς γῆς.
\vs{10}Καὶ θήσομαι τόπον τῷ λαῷ μου τῷ Ἰσραὴλ, καὶ καταφυτεύσω αὐτὸν, καὶ κατασκηνώσει καθʼ ἑαυτὸν, καὶ οὐ μεριμνήσει οὐκέτι· καὶ οὐ προσθήσει υἱὸς ἀδικίας τοῦ ταπεινῶσαι αὐτὸν, καθὼς ἀπʼ ἀρχῆς,
\vs{11}ἀπὸ τῶν ἡμερῶν ὧν ἔταξα κριτὰς ἐπὶ τὸν λαόν μου Ἰσραήλ· καὶ ἀναπαύσω σε ἀπὸ πάντων τῶν ἐχθρῶν σου· καὶ ἀπαγγελεῖ σοι Κύριος, ὅτι οἶκον οἰκοδομήσεις αὐτῷ.
\vs{12}Καὶ ἔσται ἐὰν πληρωθῶσιν αἱ ἡμέραι σου, καὶ κοιμηθήσῃ μετὰ τῶν πατέρων σου, καὶ ἀναστήσω τὸ σπέρμα σου μετὰ σὲ, ὃς ἔσται ἐκ τῆς κοιλίας σου, καὶ ἑτοιμάσω τὴν βασιλείαν αὐτοῦ.
\vs{13}Αὐτὸς οἰκοδομήσει μοι οἶκον τῷ ὀνόματί μου, καὶ ἀνορθώσω τὸν θρόνον αὐτοῦ ἕως εἰς τὸν αἰῶνα.
\vs{14}Ἐγὼ ἔσομαι αὐτῷ εἰς πατέρα, καὶ αὐτὸς ἔσται μοι εἰς υἱόν· καὶ ἐὰν ἔλθῃ ἡ ἀδικία αὐτοῦ, καὶ ἐλέγξω αὐτὸν ἐν ῥάβδῳ ἀνδρῶν, καὶ ἐν ἁφαῖς υἱῶν ἀνθρώπων·
\vs{15}Τὸ δὲ ἔλεός μου οὐκ ἀποστήσω ἀπʼ αὐτοῦ, καθὼς ἀπέστησα ἀφʼ ὧν ἀπέστησα ἐκ προσώπου μου.
\vs{16}Καὶ πιστωθήσεται ὁ οἶκος αὐτοῦ, καὶ ἡ βασιλεία αὐτοῦ ἕως αἰῶνος ἐνώπιόν μου· καὶ ὁ θρόνος αὐτοῦ ἔσται ἀνωρθωμένος εἰς τον αἰῶνα.

\vs{17}Κατὰ πάντας τοὺς λόγους τούτους, καὶ κατὰ πᾶσαν τὴν ὅρασιν ταύτην, οὕτως ἐλάλησε Νάθαν πρὸς Δαυίδ.

\vs{18}Καὶ εἰσῆλθεν ὁ βασιλεὺς Δαυὶδ, καὶ ἐνάθισεν ἐνώπιον Κυρίου, καὶ εἶπε, τίς εἰμι ἐγὼ, Κύριέ μου Κύριε, καὶ τίς ὁ οἶκός μου, ὅτι ἠγάπησάς με ἕως τούτων;
\vs{19}Καὶ κατεσμικρύνθην μικρὸν ἐνώπίον σου, Κύριέ μου Κύριε, καὶ ἐλάλησας ὑπὲρ τοῦ οἴκου τοῦ δούλου σου εἰς μακράν· οὗτος δὲ ὁ νόμος τοῦ ἀνθρώπου, Κύριέ μου Κύριε;
\vs{20}Καὶ τί προσθήσει Δαυὶδ ἔτι τοῦ λαλῆσαι πρὸς σέ; καὶ νῦν σὺ οἶδας τὸν δοῦλόν σου, Κύριέ μου Κύριε,
\vs{21}καὶ διὰ τὸν δοῦλόν σου πεποίηκας, καὶ κατὰ τὴν καρδίαν σου ἐποίησας πᾶσαν τὴν μεγαλωσύνην ταύτην, γνωρίσαι τῷ δούλῳ σου, ἕνεκεν τοῦ μεγαλύναι σε
\vs{22}Κύριέ μου· ὅτι οὐκ ἔστιν ὡς σὺ, καὶ οὐκ ἔστι Θεὸς πλὴν σοῦ ἐν πᾶσιν οἷς ἠκούσαμεν ἐν τοῖς ὠσὶν ἡμῶν.
\vs{23}Καὶ τίς ὡς ὁ λαός σου Ἰσραὴλ ἔθνος ἄλλο ἐν τῇ γῇ; ὡς ὡδήγησεν αὐτὸν ὁ Θεὸς τοῦ λυτρώσασθαι αὐτῷ λαὸν, τοῦ θέσθαι σε ὄνομα, τοῦ ποιῆσαι μεγαλωσύνην καὶ ἐπιφάνειαν, τοῦ ἐκβαλεῖν σε ἐκ προσώπου τοῦ λαοῦ σου, οὓς ἐλυτρώσω σεαυτῷ ἐξ Αἰγύπτου, ἔθνη καὶ σκηνώματα;
\vs{24}Καὶ ἡτοίμασας σεαυτῷ τὸν λαόν σου Ἰσραὴλ εἰς λαὸν ἕως αἰῶνος, καὶ σὺ Κύριε ἐγένου αὐτοῖς εἰς Θεόν.
\vs{25}Καὶ νῦν, κύριέ μου, ῥῆμα ὃ ἐλάλησας περὶ τοῦ δούλου σου καὶ τοῦ οἴκου αὐτοῦ, πίστωσον ἕως τοῦ αἰῶνος, Κύριε παντοκράτωρ θεὲ τοῦ Ἰσραήλ· καὶ νῦν καθὼς ἐλάλησας,
\vs{26}Μεγαλυνθείη τὸ ὄνομά σου ἕως αἰῶνος.
\vs{27}Κύριε παντοκράτωρ Θεὸς Ἰσραὴλ, ἀπεκάλυψας τὸ ὠτίον τοῦ δούλου σου, λέγων, οἶκον οἰκοδομήσω σοι· διὰ τοῦτο εὗρεν ὁ δοῦλός σου τὴν καρδίαν ἑαυτοῦ τοῦ προσεύξασθαι πρὸς σὲ τὴν προσευχὴν ταύτην.
\vs{28}Καὶ νῦν, Κύριέ μου Κύριε, σὺ εἶ Θεὸς, καὶ οἱ λόγοι σου ἔσονται ἀληθιναὶ, καὶ ἐλάλησας ὑπὲρ τοῦ δούλου σου τὰ ἀγαθὰ ταῦτα.
\vs{29}Καὶ νῦν ἄρξαι καὶ εὐλόγησον τὸν οἶκον τοῦ δούλου σου, τοῦ εἶναι εἰς τὸν αἰῶνα ἐνώπιόν σου· ὅτι σὺ Κύριέ μου Κύριε ἐλάλησας, καὶ ἀπὸ τῆς εὐλογίας σου εὐλογηθήσεται ὁ οἶκος τοῦ δούλου σου τοῦ εἶναι εἰς τὸν αἰῶνα.

\ch{8}
Καὶ ἐγένετο μετὰ ταῦτα, καὶ ἐπάταξε Δαυὶδ τοὺς ἀλλοφύλους, καὶ ἐτροπώσατο αὐτούς· καὶ ἔλαβε Δαυὶδ τὴν ἀφωρισμένην ἐκ χειρὸς τῶν ἀλλοφύλων.

\vs{2}Καὶ ἐπάταξε Δαυὶδ τὴν Μωὰβ, καὶ διεμέτρησεν αὐτοὺς ἐν σχοινίοις, κοιμίσας αὐτοὺς ἐπὶ τὴν γῆν· καὶ ἐγένετο τὰ δύο σχοινίσματα τοῦ θανατῶσαι, καὶ τὰ δύο σχοινίσματα ἐζώγρησε· καὶ ἐγένετο Μωὰβ τῷ Δαυὶδ εἰς δούλους φέροντας ξένια.

\vs{3}Καὶ ἐπάταξε Δαυὶδ τὸν Ἁδρααζὰρ υἱὸν Ῥαὰβ, βασιλέα Σουβὰ, πορευομένου αὐτοῦ ἐπιστῆσαι τὴν χεῖρα αὐτοῦ ἐπὶ τὸν ποταμὸν Εὐφράτην.
\vs{4}Καὶ προκατελάβετο Δαυὶδ τῶν αὐτοῦ χίλια ἅρματα, καὶ ἑπτὰ χιλιάδας ἱππέων, καὶ εἴκοσι χιλιάδας ἀνδρῶν πεζῶν· καὶ παρέλυσε Δαυὶδ πάντα τὰ ἅρματα, καὶ ὑπελείπετο ἑαυτῷ ἑκατὸν ἅρματα.
\vs{5}Καὶ παραγίνεται Συρία Δαμασκοῦ βοηθῆσαι τῷ Ἀδρααζὰρ βασιλεῖ Σουβὰ, καὶ ἐπάταξε Δαυὶδ ἐν τῷ Σύρῳ εἴκοσι δύο χιλιάδας ἀνδρῶν.
\vs{6}Καὶ ἔθετο Δαυὶδ φρουρὰν ἐν Συρίᾳ τῇ κατὰ Δαμασκὸν, καὶ ἐγένετο ὁ Σύρος τῷ Δαυὶδ εἰς δούλους φέροντας ξένια· καὶ ἔσωσε Κύριος τὸν Δαυὶδ ἐν πᾶσιν οἷς ἐπορεύετο.
\vs{7}Καὶ ἔλαβε Δαυὶδ τοὺς χλιδῶνας τοὺς χρυσοῦς, οἳ ἦσαν ἐπὶ τῶν παίδων τῶν Ἀδρααζὰρ βασιλέως Σουβὰ, καὶ ἤνεγκεν αὐτὰ εἰς Ἱερουσαλήμ· καὶ ἔλαβεν αὐτὰ Σουσακὶμ βασιλεὺς Αἰγύπτου, ἐν τῷ ἀναβῆναι αὐτὸν εἰς Ἱερουσαλὴμ ἐν ἡμέραις Ῥοβοὰμ υἱοῦ Σαλωμῶντος.
\vs{8}Καὶ ἐκ τῆς Μετεβὰκ καὶ ἐκ τῶν ἐκλεκτῶν πόλεων τοῦ Ἀδρααζὰρ ἔλαβεν ὁ βασιλεὺς Δαυὶδ χαλκὸν πολὺν σφόδρα· ἐν αὐτῷ ἐποίησε Σαλωμὼν τὴν θάλασσαν τὴν χαλκῆν, καὶ τοὺς στύλους, καὶ τοὺς λουτῆρας, καὶ πάντα τὰ σκεύη.

\vs{9}Καὶ ἤκουσε Θοοὺ ὁ βασιλεὺς Ἡμὰθ, ὅτι ἐπάταξε Δαυὶδ πᾶσαν τὴν δύναμιν Ἀδρααζὰρ,
\vs{10}καὶ ἀπέστειλε Θοοὺ Ἰεδδουρὰμ τὸν υἱὸν αὐτοῦ πρὸς βασιλέα Δαυὶδ ἐρωτῆσαι αὐτὸν τὰ εἰς εἰρήνην, καὶ εὐλογῆσαι αὐτὸν ὑπὲρ οὗ ἐπολέμησε τὸν Ἀδρααζὰρ, καὶ ἐπάταξεν αὐτὸν, ὅτι ἀντικείμενος ἦν τῷ Ἀδρααζάρ· καὶ ἐν ταῖς χερσὶν αὐτοῦ ἦσαν σκεύη ἀργυρᾶ, καὶ σκεύη χρυσᾶ, καὶ σκεύη χαλκᾶ.
\vs{11}Καὶ ταῦτα ἡγίασεν ὁ βασιλεὺς Δαυὶδ τῷ Κυρίῳ, μετὰ τοῦ ἀργυρίου καὶ μετὰ τοῦ χρυσίου οὗ ἡγίασεν ἐκ πασῶν τῶν πόλεων ὧν κατεδυνάστευσεν,
\vs{12}ἐκ τῆς Ἰδουμαίας, καὶ ἐκ τῆς Μωὰβ, καὶ ἐκ τῶν υἱῶν Ἀμμὼν, καὶ ἐκ τῶν ἀλλοφύλων, καὶ ἐξ Ἀμαλὴκ, καὶ ἐκ τῶν σκύλων Ἀδρααζὰρ υἱοῦ Ῥαὰβ βασιλέως Σουβά.

\vs{13}Καὶ ἐποίησε Δαυὶδ ὄνομα· καὶ ἐν τῷ ἀνακάμπτειν αὐτὸν ἐπάταξε τὴν Ἰδουμαίαν ἐν Γεβελὲμ εἰς ὀκτωκαίδεκα χιλιάδας.
\vs{14}Καὶ ἔθετο ἐν τῇ Ἰδουμαίᾳ φρουρὰν, ἐν πάσῃ τῇ Ἰδουμαίᾳ· καὶ ἐγένοντο πάντες οἱ Ἰδουμαῖοι δοῦλοι τῷ βασιλεῖ· καὶ ἔσωσε Κύριος τὸν Δαυὶδ ἐν πᾶσιν οἷς ἐπορεύετο.

\vs{15}Καὶ ἐβασίλευσε Δαυὶδ ἐπὶ πάντα Ἰσραήλ· καὶ ἦν Δαυὶδ ποιῶν κρίμα καὶ δικαιοσύνην ἐπὶ πάντα τὸν λαὸν αὐτοῦ.
\vs{16}Καὶ Ἰωὰβ υἱὸς Σαρουίας ἐπὶ τῆς στρατιᾶς· καὶ Ἰωσαφὰτ υἱὸς Ἀχιλοὺδ ἐπὶ τῶν ὑπομνημάτων·
\vs{17}καὶ Σαδὼκ υἱὸς Ἀχιτὼβ καὶ Ἀχιμέλεχ υἱὸς Ἀβιάθαρ ἱερεῖς· καὶ Σασὰ ὁ γραμματεύς·
\vs{18}καὶ Βαναίας υἱὸς Ἰωδαὲ σύμβουλος· καὶ ὁ Χελεθὶ, καὶ ὁ Φελετὶ, καὶ οἱ υἱοὶ Δαυὶδ αὐλάρχαι ἦσαν.

\ch{9}
Καὶ εἶπε Δαυὶδ, εἰ ἔστιν ἔτι ὑπολελειμμένος ἐν τῷ οἴκῳ Σαοὺλ, καὶ ποιήσω μετʼ αὐτοῦ ἔλεος ἕνεκεν Ἰωνάθαν;
\vs{2}Καὶ ἐκ τοῦ οἴκου Σαοὺλ ἦν παῖς, καὶ ὄνομα αὐτῷ Σιβά· καὶ καλοῦσιν αὐτὸν πρὸς Δαυίδ· καὶ εἶπε πρὸς αὐτὸν ὁ βασιλεὺς, σὺ εἶ Σιβά; καὶ εἶπεν, ἐγὼ δοῦλος σός.
\vs{3}Καὶ εἶπεν ὁ βασιλεὺς, εἰ ὑπολέλειπται ἐκ τοῦ οἴκου Σαοὺλ ἔτι ἀνὴρ, καὶ ποιήσω μετʼ αὐτοῦ ἔλεος Θεοῦ; καὶ εἶπε Σιβὰ πρὸς τὸν βασιλέα, ἔτι ἐστὶν υἱὸς τῷ Ἰωνάθαν πεπληγὼς τοὺς πόδας.
\vs{4}Καὶ εἶπεν ὁ βασιλεὺς, ποῦ οὗτος; καὶ εἶπε Σιβὰ πρὸς τὸν βασιλέα, ἰδοὺ ἐν οἴκῳ Μαχεὶρ υἱοῦ Ἀμιὴλ ἐκ τῆς Λοδάβαρ.
\vs{5}Καὶ ἀπέστειλεν ὁ βασιλεὺς Δαυὶδ, καὶ ἔλαβεν αὐτὸν ἐκ τοῦ οἴκου Μαχὶρ υἱοῦ Ἀμιὴλ ἐκ τῆς Λαδάβαρ.

\vs{6}Καὶ παραγίνεται Μεμφιβοσθὲ υἱὸς Ἰωνάθαν υἱοῦ Σαοὺλ πρὸς τὸν βασιλέα Δαυίδ, καὶ ἔπεσεν ἐπὶ πρόσωπον αὐτοῦ, καὶ προσεκύνησεν αὐτῷ· καὶ εἶπεν αὐτῷ Δαυίδ, Μεμφιβοσθέ; καὶ εἶπεν, ἰδοὺ ὁ δοῦλός σου.
\vs{7}Καὶ εἶπεν αὐτῷ Δαυίδ, μὴ φοβοῦ, ὅτι ποιῶν ποιήσω μετὰ σοῦ ἔλεος διὰ Ἰωνάθαν τὸν πατέρα σου, καὶ ἀποκαταστήσω σοι πάντα ἀγρὸν Σαοὺλ πατρὸς τοῦ πατρός σου, καὶ σὺ φαγῃ ἄρτον ἐπὶ τῆς τραπέζης μου διαπαντός.
\vs{8}Καὶ προσεκύνησε Μεμφιβοσθὲ, καὶ εἶπε, τίς εἰμι ὁ δοῦλός σου, ὅτι ἐπέβλεψας ἐπὶ τὸν κύνα τὸν τεθνηκότα τὸν ὅμοιον ἐμοί;

\vs{9}Καὶ ἐκάλεσεν ὁ βασιλεὺς Σιβὰ τὸ παιδάριον Σαοὺλ, καὶ εἶπε πρὸς αὐτὸν, πάντα ὅσα ἐστὶ τῷ Σαοὺλ καὶ ὅλῳ τῷ οἴκῳ αὐτοῦ δέδωκα τῷ υἱῷ τοῦ κυρίου σου.
\vs{10}Καὶ ἐργᾷ αὐτῷ τὴν γῆν σὺ, καὶ οἱ υἱοί σου, καὶ οἱ δοῦλοί σου, καὶ εἰσοίσεις τῷ υἱῷ τοῦ κυρίου σου ἄρτους, καὶ ἔδεται ἄρτους· καὶ Μεμφιβοσθὲ υἱὸς τοῦ κυρίου σου φάγεται διαπαντὸς ἄρτον ἐπὶ τῆς τραπέζης μου· καὶ τῷ Σιβᾷ ἦσαν πεντεκαίδεκα υἱοὶ, καὶ εἴκοσι δοῦλοι.
\vs{11}Καὶ εἶπε Σιβὰ πρὸς τὸν βασιλέα, κατὰ πάντα ὅσα ἐντέταλται ὁ κύριός μου ὁ βασιλεὺς τῷ δούλῳ αὐτοῦ, οὕτως ποιήσει ὁ δοῦλός σου· καὶ Μεμφιβοσθὲ ἤσθιεν ἐπὶ τῆς τραπέζης Δαυὶδ καθὼς εἷς τῶν υἱῶν αὐτοῦ τοῦ βασιλέως.
\vs{12}Καὶ τῷ Μεμφιβοσθὲ υἱὸς μικρὸς ἦν, καὶ ὄνομα αὐτῷ Μιχά· καὶ πᾶσα ἡ κατοίκησις τοῦ οἴκου Σιβὰ δοῦλοι τοῦ Μεμφιβοσθέ.
\vs{13}Καὶ Μεμφιβοσθὲ κατῴκει ἐν Ἱερουσαλὴμ, ὅτι ἐπὶ τῆς τραπέζης τοῦ βασιλέως αὐτὸς διαπαντὸς ἤσθιε, καὶ αὐτὸς ἦν χωλὸς ἀμφοτέροις τοῖς ποσὶν αὐτοῦ.

\ch{10}
Καὶ ἐγένετο μετὰ ταῦτα καὶ ἀπέθανε βασιλεὺς υἱῶν Ἀμμὼν, καὶ ἐβασίλευσεν Ἀννὼν υἱὸς αὐτοῦ ἀντʼ αὐτοῦ.
\vs{2}Καὶ εἶπε Δαυὶδ, ποιήσω ἔλεος μετὰ Ἀννὼν υἱοῦ Ναὰς, ὃν τρόπον ἐποίησεν ὁ πατὴρ αὐτοῦ μετʼ ἐμοῦ ἔλεος. Καὶ ἀπέστειλεν Δαυὶδ παρακαλέσαι αὐτὸν ἐν χειρὶ τῶν δούλων αὐτοῦ περὶ τοῦ πατρὸς αὐτοῦ· καὶ παρεγένοντο οἱ παῖδες Δαυὶδ εἰς τὴν γῆν υἱῶν Ἀμμών.
\vs{3}Καὶ εἶπον οἱ ἄρχοντες υἱῶν Ἀμμὼν πρὸς Ἀννῶν τὸν κύριον αὐτῶν, μὴ παρὰ τὸ δοξάζειν Δαυὶδ τὸν πατέρα σου ἐνώπιόν σου, ὅτι ἀπέστειλέ σοι παρακαλοῦντας; ἀλλʼ ὅπως οὐχὶ ἐρευνήσωσι τὴν πόλιν καὶ κατασκοπήσωσιν αὐτὴν καὶ τοῦ κατασκέψασθαι αὐτὴν ἀπέστειλε Δαυὶδ τοὺς παῖδας αὐτοῦ πρὸς σέ;
\vs{4}Καὶ ἔλαβεν Ἀννὼν τοὺς παῖδας Δαυὶδ, καὶ ἐξύρησε τοὺς πώγωνας αὐτῶν, καὶ ἀπέκοψε τοὺς μανδύας αὐτῶν ἐν τῷ ἡμίσει ἕως τῶν ἰσχίων αὐτῶν, καὶ ἐξαπέστειλεν αὐτούς.

\vs{5}Καὶ ἀπήγγειλαν τῷ Δαυὶδ ὑπὲρ τῶν ἀνδρῶν, καὶ ἀπέστειλεν εἰς ἀπαντὴν αὐτῶν, ὅτι ἦσαν οἱ ἄνδρες ἠτιμασμένοι σφόδρα· καὶ εἶπεν ὁ βασιλεὺς, καθίσατε ἐν Ἱεριχὼ ἕως τοῦ ἀνατεῖλαι τοὺς πώγωνας ὑμῶν, καὶ ἐπιστραφήσεσθε.

\vs{6}Καὶ εἶδον οἱ υἱοὶ Ἀμμὼν ὅτι κατῃσχύνθησαν ὁ λαὸς Δαυίδ· καὶ ἀπέστειλαν οἱ υἱοὶ Ἀμμὼν, καὶ ἐμισθώσαντο τὴν Συρίαν Βαιθραὰμ, καὶ τὴν Συρίαν Σουβὰ, καὶ Ῥοὼβ, εἴκοσι χιλιάδας πεζῶν, καὶ τὸν βασιλέα Ἀμαλὴκ χιλίους ἄνδρας, καὶ Ἰστὼβ δώδεκα χιλιάδας ἀνδρῶν.

\vs{7}Καὶ ἤκουσε Δαυὶδ, καὶ ἀπέστειλε τὸν Ἰωὰβ καὶ πᾶσαν τὴν δύναμιν τοὺς δυνατούς.
\vs{8}Καὶ ἐξῆλθον οἱ υἱοὶ Ἀμμὼν καὶ παρετάξαντο πόλεμον παρὰ τῇ θύρᾳ τῇς πύλης, Συρίας Σουβὰ καὶ Ῥοὼβ καὶ Ἰστὼβ καὶ Ἀμαλὴκ μόνοι ἐν ἀγρῷ.
\vs{9}Καὶ εἶδεν Ἰωὰβ ὅτι ἐγενήθη πρὸς αὐτὸν ἀντιπρόσωπον τοῦ πολέμου ἐκ τοῦ κατὰ πρόσωπον ἐξεναντίας καὶ ἐκ τοῦ ὄπισθεν, καὶ ἐπελέξατο ἐκ πάντων τῶν νεανιῶν Ἰσραὴλ, καὶ παρετάξαντο ἐξ ἐναντίας Συρίας.
\vs{10}Καὶ τὸ κατάλοιπον τοῦ λαοῦ ἔδωκεν ἐν χειρὶ Ἀβεσσὰ τοῦ ἀδελφοῦ αὐτοῦ, καὶ παρετάξαντο ἐξεναντίας υἱῶν Ἀμμών.
\vs{11}Καὶ εἶπεν, εὰν κραταιωθῇ Συρία ὑπὲρ ἐμὲ, καὶ ἔσεσθέ μοι εἰς σωτηρίαν· καὶ ἐὰν κραταιωθῶσιν υἱοὶ Ἀμμὼν ὑπὲρ σὲ, καὶ ἐσόμεθα τοῦ σῶσαί σε.
\vs{12}Ἀνδρίζου καὶ κραταιωθῶμεν ὑπὲρ τοῦ λαοῦ ἡμῶν καὶ περὶ τῶν πόλεων τοῦ Θεοῦ ἡμῶν, καὶ Κύριος ποιήσει τὸ ἀγαθὸν ἐν ὀφθαλμοῖς αὐτοῦ.

\vs{13}Καὶ προσῆλθεν Ἰωὰβ καὶ ὁ λαὸς αὐτοῦ μετʼ αὐτοῦ εἰς πόλεμον πρὸς Συρίαν, καὶ ἔφυγαν ἀπὸ προσώπου αὐτοῦ.
\vs{14}Καὶ οἱ υἱοὶ Ἀμμὼν εἶδαν ὅτι ἔφυγε Συρία, καὶ ἔφυγαν ἀπὸ προσώπου Ἀβεσσὰ, καὶ εἰσῆλθον εἰς τὴν πόλιν· καὶ ἀνέστρεψεν Ἰωὰβ ἀπὸ τῶν υἱῶν Ἀμμὼν, καὶ παρεγένετο εἰς Ἱερουσαλήμ.

\vs{15}Καὶ εἶδε Συρία ὅτι ἔπτάισεν ἔμπροσθεν Ἰσραὴλ, καὶ συνήχθησαν ἐπὶ τὸ αὐτό.
\vs{16}Καὶ ἀπέστειλεν Ἀδρααζὰρ, καὶ συνήγαγε τὴν Συρίαν τὴν ἐκ τοῦ πέραν τοῦ ποταμοῦ Χαλαμὰκ, καὶ παρεγένοντο Αἱλάμ· καὶ Σωβὰκ ἄρχων τῆς δυνάμεως Ἀδρααζὰρ ἔμπροσθεν αὐτῶν.

\vs{17}Καὶ ἀπηγγέλη τῷ Δαυὶδ, καὶ συνήγαγε τὸν πάντα Ἰσραὴλ, καὶ διέβη τὸν Ἰορδάνην, καὶ παρεγένετο εἰς Αἰλάμ· καὶ παρετάξατο Συρία ἀπέναντι Δαυὶδ, καὶ ἐπολέμησαν μετʼ αὐτοῦ.
\vs{18}Καὶ ἔφυγε Συρία ἀπὸ πρόσωπου Ἰσραήλ· καὶ ἀνεῖλε Δαυὶδ ἐκ τῆς Συρίας ἑπτακόσια ἅρματα, καὶ τεσσαράκοντα χιλιάδας ἱππέων, καὶ τὸν Σωβὰκ τὸν ἄρχοντα τῆς δυνάμεως αὐτοῦ ἐπάταξε, καὶ ἀπέθανεν ἐκεῖ.
\vs{19}Καὶ εἶδαν πάντες οἱ βασιλεῖς οἱ δοῦλοι Ἀδρααζὰρ ὅτι ἔπταισαν ἔμπροσθεν Ἰσραὴλ, καὶ ηὐτομόλησαν μετὰ Ἰσραὴλ, καὶ ἐδούλευσαν αὐτοῖς· καὶ ἐφοβήθη Συρία τοῦ σῶσαι ἔτι τοὺς υἱοὺς Ἀμμών.

\ch{11}
Καὶ ἐγένετο, ἐπιστρέψαντος τοῦ ἐνιαυτοῦ εἰς τὸν καιρὸν τῆς ἐξοδίας τῶν βασιλέων, καὶ ἀπέστειλε Δαυὶδ τὸν Ἰωὰβ, καὶ τοὺς παῖδας αὐτοῦ μετʼ αὐτοῦ, καὶ τὸν πάντα Ἰσραὴλ, καὶ διέφθειραν τοὺς υἱοὺς Ἀμμών· καὶ διεκάθισαν ἐπὶ Ῥαββάθ· καὶ Δαυὶδ ἐκάθισεν ἐν Ἱερουσαλήμ.

\vs{2}Καὶ ἐγένετο πρὸς ἑσπέραν, καὶ ἀνέστη Δαυὶδ ἀπὸ τῆς κοίτης αὐτοῦ, καὶ περιεπάτει ἐπὶ τοῦ δώματος τοῦ οἴκου τοῦ βασιλέως, καὶ εἶδε γυναῖκα λουομένην ἀπὸ τοῦ δώματος, καὶ ἡ γυνὴ καλὴ τῷ εἴδει σφόδρα.
\vs{3}Καὶ ἀπέστειλε Δαυὶδ, καὶ ἐζήτησε τὴν γυναῖκα, καὶ εἶπεν, οὐχὶ αὕτη Βηρσαβεὲ θυγάτηρ Ἐλιὰβ γυνὴ Οὐρίου τοῦ Χετταίου;

\vs{4}Καὶ ἀπέστειλεν Δαυὶδ ἀγγέλους, καὶ ἔλαβεν αὐτὴν, καὶ εἰσῆλθε πρὸς αὐτὴν, καὶ ἐκοιμήθη μετʼ αὐτῆς· καὶ αὕτη ἁγιαζομένη ἀπὸ ἀκαθαρσίας αὐτῆς, καὶ ἀπέστρεψεν εἰς τὸν οἶκον αὐτῆς.
\vs{5}Καὶ ἐν γαστρὶ ἔλαβεν ἡ γυνή· καὶ ἀποστείλασα ἀπήγγειλε τῷ Δαυὶδ, καὶ εἶπεν, ἐγώ εἰμι ἐν γαστρὶ ἔχω.
\vs{6}Καὶ ἀπέστειλε Δαυὶδ πρὸς Ἰωὰβ, λέγων, ἀπόστειλον πρὸς μὲ τὸν Οὐρίαν τὸν Χετταῖον· καὶ ἀπέστειλεν Ἰωὰβ τὸν Οὐρίαν πρὸς Δαυίδ.

\vs{7}Καὶ παραγίνεται Οὐρίας καὶ εἰσῆλθε πρὸς αὐτὸν, καὶ ἐπηρώτησε Δαυὶδ εἰς εἰρήνην Ἰωὰβ, καὶ εἰς εἰρήνην τοῦ λαοῦ, καὶ εἰς εἰρήνην τοῦ πολέμου.
\vs{8}Καὶ εἶπε Δαυὶδ τῷ Οὐρίᾳ, Κατάβηθι εἰς τὸν οἶκόν σου, καὶ νίψαι τοὺς πόδας σου· καὶ ἐξῆλθεν Οὐρίας ἐξ οἴκου τοῦ βασιλέως, καὶ ἐξῆλθεν ὀπίσω αὐτοῦ ἄρσις τοῦ βασιλέως.
\vs{9}Καὶ ἐκοιμήθη Οὐρίας παρὰ τῇ θύρᾳ τοῦ βασιλέως μετὰ τῶν δούλων τοῦ κυρίου αὐτοῦ, καὶ οὐ κατέβη εἰς τὸν οἶκον αὐτοῦ.
\vs{10}Καὶ ἀνήγγειλαν τῷ Δαυὶδ, λέγοντες, ὅτι Οὐ κατέβη Οὐρίας εἰς τὸν οἶκον αὐτοῦ· καὶ εἶπε Δαυὶδ πρὸς Οὐρίαν, οὐχὶ ἐξ ὁδοῦ σὺ ἔρχῃ; τί ὅτι οὐ κατέβης εἰς τὸν οἶκόν σου;
\vs{11}Καὶ εἶπεν Οὐρίας πρὸς Δαυὶδ, ἡ κιβωτὸς, καὶ Ἰσραὴλ, καὶ Ἰούδας κατοικοῦσιν ἐν σκηναῖς, καὶ ὁ κύριός μου Ἰωὰβ, καὶ οἱ δοῦλοι τοῦ κυρίου μου ἐπὶ πρόσωπον τοῦ ἀγροῦ παρεμβάλλουσι, καὶ ἐγὼ εἰσελεύσομαι εἰς τὸν οἶκόν μου τοῦ φαγεῖν, καὶ πιεῖν, καὶ κοιμηθῆναι μετὰ τῆς γυναικός μου; πῶς; ζῇ ἡ ψυχή σου, εἰ ποιήσω τὸ ῥῆμα τοῦτο.
\vs{12}Καὶ εἶπε Δαυὶδ πρὸς Οὐρίαν, κάθισον ἐνταῦθα καί γε σήμερον, καὶ αὔριον ἐξαποστελῶ σε· καὶ ἐκάθισεν Οὐρίας ἐν Ἱερουσαλὴμ ἐν τῇ ἡμέρᾳ ἐκείνῃ καὶ τῇ ἐπαύριον.

\vs{13}Καὶ ἐκάλεσεν αὐτὸν Δαυὶδ, καὶ ἔφαγεν ἐνώπιον αὐτοῦ, καὶ ἔπιε, καὶ ἐμέθυσεν αὐτὸν, καὶ ἐξῆλθεν ἑσπέρας τοῦ κοιμηθῆναι ἐπὶ τῆς κοίτης αὐτοῦ μετὰ τῶν δούλων τοῦ κυρίου αὐτοῦ, καὶ εἰς τὸν οἶκον αὐτοῦ οὐ κατέβη.

\vs{14}Καὶ ἐγένετο πρωῒ, καὶ ἔγραψε Δαυὶδ βιβλίον πρὸς Ἰωὰβ, καὶ ἀπέστειλεν ἐν χειρὶ Οὐρίου.
\vs{15}Καὶ ἔγραψεν ἐν βιβλίῳ, λέγων, εἰσάγαγε τὸν Οὐρίαν ἐξεναντίας τοῦ πολέμου τοῦ κραταιοῦ, καὶ ἀποστραφήσεσθε ἀπὸ ὄπισθεν αὐτοῦ, καὶ πληγήσεται καὶ ἀποθανεῖται.

\vs{16}Καὶ ἐγενήθη ἐν τῷ φυλάσσειν Ἰωὰβ ἐπὶ τὴν πόλιν, καὶ ἔθηκε τὸν Οὐρίαν εἰς τὸν τόπον οὗ ᾔδει ὅτι ἄνδρες δυνάμεως ἐκεῖ.
\vs{17}Καὶ ἐξῆλθον οἱ ἄνδρες τῆς πόλεως, καὶ ἐπολέμουν μετὰ Ἰωάβ· καὶ ἔπεσαν ἐκ τοῦ λαοῦ ἐκ τῶν δούλων Δαυὶδ, καὶ ἀπέθανε καί γε Οὐρίας ὁ Χετταῖος.

\vs{18}Καὶ ἀπέστειλεν Ἰωὰβ, καὶ ἀπήγγειλε τῷ Δαυὶδ πάντας τοὺς λόγους τοῦ πολέμου λαλῆσαι πρὸς τὸν βασιλέα·
\vs{19}καὶ ἐνετείλατο τῷ ἀγγέλῳ, λέγων, ἐν τῷ συντελέσαι πάντας τοὺς λόγους τοῦ πολέμου λαλῆσαι πρὸς τὸν βασιλέα,
\vs{20}καὶ ἔσται ἐὰν ἀναβῇ ὁ θυμὸς τοῦ βασιλέως, καὶ εἴπῃ σοι, τί ὅτι ἠγγίσατε πρὸς τὴν πόλιν πολεμῆσαι; οὐκ ᾔδειτε ὅτι τοξεύσουσιν ἀπάνωθεν τοῦ τείχους;
\vs{21}Τίς ἐπάταξε τὸν Ἀβιμέλεχ υἱὸν Ἱεροβάαλ υἱοῦ Νήρ; οὐχὶ γυνὴ ἔῤῥιψε κλάσμα μύλου ἐπʼ αὐτὸν ἀπὸ ἄνωθεν τοῦ τείχους, καὶ ἀπέθανεν ἐν Θαμασί; ἱνατί προσηγάγετε πρὸς τὸ τεῖχος; καὶ ἐρεῖς, καί γε ὁ δοῦλός σου Οὐρίας ὁ Χετταῖος ἀπέθανε.

\vs{22}Καὶ ἐπορεύθη ὁ ἄγγελος Ἰωὰβ πρὸς τὸν βασιλέα εἰς Ἱερουσαλὴμ, καὶ παρεγένετο καὶ ἀπήγγειλε τῷ Δαυὶδ πάντα ὅσα ἀπήγγειλεν αὐτῷ Ἰωὰβ, πάντα τὰ ῥήματα τοῦ πολέμου· καὶ ἐθυμώθη Δαυὶδ πρὸς Ἰωὰβ, καὶ εἶπε πρὸς τὸν ἄγγελον, ἱνατί προσηγάγετε πρὸς τὴν πόλιν τοῦ πολεμῆσαι; οὐκ ᾔδειτε ὅτι πληγήσεσθε ἀπὸ τοῦ τείχους; τίς ἐπάταξε τὸν Ἀβιμέλεχ υἱὸν Ἱεροβάαλ; οὐχὶ γυνὴ ἔῤῥιψεν ἐπʼ αὐτὸν κλάσμα μύλου ἀπὸ τοῦ τείχους, καὶ ἀπέθανεν ἐν Θαμασί; ἱνατί προσηγάγετε πρὸς τὸ τεῖχος;
\vs{23}Καὶ εἶπεν ὁ ἄγγελος πρὸς Δαυὶδ, ὅτι ἐκραταίωσαν ἐφʼ ἡμᾶς οἱ ἄνδρες, καὶ ἐξῆλθον ἐφʼ ἡμᾶς εἰς τὸν ἀγρὸν, καὶ ἐγενήθημεν ἐπʼ αὐτοὺς ἕως τῆς θύρας τῆς πύλης.
\vs{24}Καὶ ἐτόξευσαν οἱ τοξεύοντες πρὸς τοὺς παῖδάς σου ἀπάνωθεν τοῦ τείχους, καὶ ἀπέθανον τῶν παίδων τοῦ βασιλέως, καί γε ὁ δοῦλός σου Οὐρίας ὁ Χετταῖος ἀπέθανε.
\vs{25}Καὶ εἶπε Δαυὶδ πρὸς τὸν ἄγγελον, τάδε ἐρεῖς πρὸς Ἰωάβ, μὴ πονηρὸν ἔστω ἐν ὀφθαλμοῖς σου τὸ ῥῆμα τοῦτο, ὅτι ποτὲ μὲν οὕτως καὶ ποτὲ οὕτως φάγεται ἡ μάχαιρα· κραταίωσον τὸν πόλεμόν σου εἰν τὴν πόλιν, καὶ κατάσπασον αὐτὴν, καὶ κραταίωσον αὐτήν.

\vs{26}Καὶ ἤκουσεν ἡ γυνὴ Οὐρίου ὅτι ἀπέθανεν Οὐρίας ὁ ἀνὴρ αὐτῆς, καὶ ἐκόψατο τὸν ἄνδρα αὐτῆς.
\vs{27}Καὶ διῆλθε τὸ πένθος, καὶ ἀπέστειλεν Δαυὶδ, καὶ συνήγαγεν αὐτὴν εἰς τὸν οἶκον αὐτοῦ, καὶ ἐγενήθη αὐτῷ εἰς γυναῖκα, καὶ ἔτεκεν αὐτῷ υἱόν· καὶ πονηρὸν ἐφάνη τὸ ῥῆμα ὃ ἐποίησε Δαυὶδ ἐν ὀφθαλμοῖς Κυρίου.

\ch{12}
Καὶ ἀπέστειλε Κύριος τὸν Νάθαν τὸν προφήτην πρὸς Δαυίδ· καὶ εἰσῆλθε πρὸς αὐτὸν, καὶ εἶπεν αὐτῷ, δύο ἦσαν ἄνδρες ἐν πόλει μιᾷ, εἷς πλούσιος, καὶ εἷς πένης.
\vs{2}Καὶ τῷ πλουσίῳ ἦν ποίμνια καὶ βουκόλια πολλὰ σφόδρα.
\vs{3}Καὶ τῷ πένητι οὐδὲν ἀλλʼ ἢ ἀμνὰς μία μικρὰ, ἣν ἐκτήσατο καὶ περιεποιήσατο, καὶ ἐξέθρεψεν αὐτὴν, καὶ ἡδρύνθη μετʼ αὐτοῦ καὶ μετὰ τῶν υἱῶν αὐτοῦ ἐπὶ τὸ αὐτὸ, ἐκ τοῦ ἄρτου αὐτοῦ ἤσθιε, καὶ ἐκ τοῦ ποτηρίου αὐτοῦ ἔπινε, καὶ ἐν τῷ κόλπῷ αὐτοῦ ἐκάθευδε, καὶ ἦν αὐτῷ ὡς θυγάτηρ.
\vs{4}Καὶ ἦλθε πάροδος τῷ ἀνδρὶ τῷ πλουσίῳ, καὶ ἐφείσατο λαβεῖν ἐκ τῶν ποιμνίων αὐτοῦ, καὶ ἐκ τῶν βουκολίων αὐτοῦ, τοῦ ποιῆσαι τῷ ξένῳ ὁδοιπόρῳ τῷ ἐλθόντι πρὸς αὐτὸν, καὶ ἔλαβε τὴν ἀμνάδα τοῦ πένητος. καὶ ἐποίησεν αὐτὴν τῷ ἀνδρὶ τῷ ἐλθόντι πρὸς αὐτόν.
\vs{5}Καὶ ἐθυμώθη ὀργῇ Δαυὶδ σφόδρα τῷ ἀνδρὶ, καὶ εἶπε Δαυὶδ πρὸς Νάθαν, ζῇ Κύριος, ὅτι υἱὸς θανάτου ὁ ἀνὴρ ὁ ποιήσας τοῦτο·
\vs{6}Καὶ τὴν ἀμνάδα ἀποτίσει ἑπταπλασίονα, ἀνθʼ ὧν ὅτι ἐποίησε τὸ ῥῆμα τοῦτο, καὶ περὶ οὗ οὐκ ἐφείσατο.

\vs{7}Καὶ εἶπε Νάθαν πρὸς Δαυὶδ, σὺ εἶ ὁ ἀνὴρ ὁ ποιήσας τοῦτο· τάδε λέγει Κύριος ὁ Θεὸς Ἰσραὴλ, ἐγώ εἰμι ὁ χρίσας σε εἰς βασιλέα ἐπὶ Ἰσραὴλ, καὶ ἐγώ εἰμι ἐῤῥυσάμην σε ἐκ χειρὸς Σαοὺλ,
\vs{8}καὶ ἔδωκά σοι τὸν οἶκον τοῦ κυρίου σου, καὶ τὰς γυναῖκας τοῦ κυρίου σου ἐν τῷ κόλπῳ σου, καὶ ἔδωκά σοι τὸν οἶκον Ἰσραὴλ καὶ Ἰούδα· καὶ εἰ μικρόν ἐστι, προσθήσω σοι κατὰ ταῦτα.
\vs{9}Τί ὅτι ἐφαύλισας τὸν λόγον Κυρίου, τοῦ ποιῆσαι τὸ πονηρὸν ἐν ὀφθαλμοῖς αὐτοῦ; τὸν Οὐρίαν τὸν Χετταῖον ἐπάταξας ἐν ῥομφαίᾳ, καὶ τὴν γυναῖκα αὐτοῦ ἔλαβες σεαυτῷ εἰς γυναῖκα, καὶ αὐτὸν ἀπέκτεινας ἐν ῥομφαίᾳ υἱῶν Ἀμμών.
\vs{10}Καὶ νὺν οὐκ ἀποστήσεται ῥομφαία ἐκ τοῦ οἴκου σου ἕως αἰῶνος, ἀνθʼ ὧν ὅτι ἐξουδένωσάς με, καὶ ἔλαβες τὴν γυναῖκα τοῦ Οὐρίου τοῦ Χετταίου, τοῦ εἶναί σοι εἰς γυναῖκα.
\vs{11}Τάδε λέγει Κύριος, ἰδοὺ ἐγὼ ἐξεγείρω ἐπὶ σὲ κακὰ ἐκ τοῦ οἴκου σου, καὶ λήψομαι τὰς γυναῖκάς σου κατʼ ὀφθαλμούς σου, καὶ δώσω τῷ πλησίον σου, καὶ κοιμηθήσεται μετὰ τῶν γυναικῶν σου ἐναντίον τοῦ ἡλίου τούτου.
\vs{12}Ὅτι σὺ ἐποίησας κρυβῇ, κᾀγὼ ποιήσω τὸ ῥῆμα τοῦτο ἐναντίον παντὸς Ἰσραὴλ, καὶ ἀπέναντι τοῦ ἡλίου τούτου.

\vs{13}Καὶ εἶπε Δανὶδ τῷ Νάθαν, ἡμάρτηκα τῷ Κυρίῳ· καὶ εἶπε Νάθαν πρὸς Δαυὶδ, καὶ Κύριος παρεβίβασε τὸ ἁμάρτημά σου· οὐ μὴ ἀποθάνῃς.
\vs{14}Πλὴν ὅτι παροργίζων παρώργισας τοὺς ἐχθροὺς Κυρίου ἐν τῷ ῥήματι τούτῳ, καί γε ὁ υἱός σου ὁ τεχθείς σοι θανάτῳ ἀποθανεῖται.

\vs{15}Καὶ ἀπῆλθε Νάθαν εἰς τὸν οἶκον αὐτοῦ· καὶ ἔθραυσε Κύριος τὸ παιδίον ὃ ἔτεκεν ἡ γυνὴ Οὐρίου τοῦ Χετταίου τῷ Δαυὶδ, καὶ ἠῤῥώστησε.
\vs{16}Καὶ ἐζήτησε Δαυὶδ τὸν Θεὸν περὶ τοῦ παιδαρίου, καὶ ἐνήστευσε Δαυὶδ νηστείαν, καὶ εἰσῆλθε καὶ ηὐλίσθη ἐπὶ τῆς γῆς.
\vs{17}Καὶ ἀνέστησαν ἐπʼ αὐτὸν οἱ πρεσβύτεροι τοῦ οἴκου αὐτοῦ ἐγεῖραι αὐτὸν ἀπὸ τῆς γῆς, καὶ οὐκ ἠθέλησε, καὶ οὐ συνέφαγεν αὐτοῖς ἄρτον.

\vs{18}Καὶ ἐγένετο ἐν τῇ ἡμέρᾳ τῇ ἑβδόμῃ, καὶ ἀπέθανε τὸ παιδάριον· καὶ ἐφοβήθησαν οἱ δοῦλοι Δαυὶδ ἀναγγεῖλαι αὐτῷ, ὅτι τέθνηκε τὸ παιδάριον, ὅτι εἶπον, ἰδοὺ ἐν τῷ τὸ παιδάριον ἔτι ζῇν ἐλαλήσαμεν πρὸς αὐτὸν, καὶ οὐκ εἰσήκουσε τῆς φωνῆς ἡμῶν· καὶ πῶς εἴπωμεν πρὸς αὐτὸν ὅτι τέθνηκε τὸ παιδάριον, καὶ ποιήσει κακά;
\vs{19}Καὶ συνῆκε Δαυὶδ, ὅτι οἱ παῖδες αὐτοῦ ψιθυρίζουσι, καὶ ἐνόησε Δαυὶδ ὅτι τέθνηκε τὸ παιδάριον· καὶ εἶπε Δαυὶδ πρὸς τοὺς παῖδας αὐτοῦ, εἰ τέθνηκε τὸ παιδάριον; καὶ εἶπαν, τέθνηκε.
\vs{20}Καὶ ἀνέστη Δαυὶδ ἐκ τῆς γῆς, καὶ ἐλούσατο, καὶ ἠλείψατο, καὶ ἤλλαξε τὰ ἱμάτια αὐτοῦ, καὶ εἰσῆλθεν εἰς τὸν οἶκον τοῦ Θεοῦ, καὶ προσεκύνησεν αὐτῷ, καὶ εἰσῆλθεν εἰς τὸν οἶκον αὐτοῦ, καὶ ᾔτησεν ἄρτον φαγεῖν, καὶ παρέθηκαν αὐτῷ ἄρτον, καὶ ἔφαγε.
\vs{21}Καὶ εἶπαν οἱ παῖδες αὐτοῦ πρὸς αὐτόν, τί τὸ ῥῆμα τοῦτο ὃ ἐποίησας ἕνεκα τοῦ παιδαρίου; ἔτι ζῶντος ἐνήστευες καὶ ἔκλαιες καὶ ἠγρύπνεις, καὶ ἡνίκα ἀπέθανε τὸ παιδάριον, ἀνέστης, καὶ ἔφαγες ἄρτον, καὶ πέπωκας;
\vs{22}Καὶ εἶπε Δαυὶδ, ἐν τῷ τὸ παιδάριον ἔτι ζῇν ἐνήστευσα καὶ ἔκλαυσα, ὅτι εἶπα, τίς οἶδεν εἰ ἐλεήσει με Κύριος, καὶ ζήσεται τὸ παιδάριον;
\vs{23}Καὶ νῦν τέθνηκεν, ἱνατί τοῦτο ἐγὼ νηστεύω; μὴ δυνήσομαι ἐπιστρέψαι αὐτὸν ἔτι; ἐγὼ πορεύσομαι πρὸς αὐτόν, καὶ αὐτὸς οὐκ ἀναστρέψει πρὸς μέ.

\vs{24}Καὶ παρεκάλεσε Δαυὶδ Βηρσαβεὲ τὴν γυναῖκα αὐτοῦ, καὶ εἰσῆλθε πρὸς αὐτὴν, καὶ ἐκοιμήθη μετʼ αὐτῆς, καὶ συνέλαβε καὶ ἔτεκεν υἱὸν, καὶ ἐκάλεσε τὸ ὄνομα αὐτοῦ Σαλωμὼν, καὶ Κύριος ἠγάπησεν αὐτόν.
\vs{25}Καὶ ἀπέστειλεν ἐν χειρὶ Νάθαν τοῦ προφήτου, καὶ ἐκάλεσε τὸ ὄνομα αὐτοῦ Ἰεδδεδὶ, ἕνεκεν Κυρίου.

\vs{26}Καὶ ἐπολέμησεν Ἰωὰβ ἐν Ῥαββὰθ υἱῶν Ἀμμὼν, καὶ κατέλαβε τὴν πόλιν τῆς βασιλείας.
\vs{27}Καὶ ἀπέστειλεν Ἰωὰβ ἀγγέλους πρὸς Δαυὶδ, καὶ εἶπεν, ἐπολέμησα ἐν Ῥαββὰθ, καὶ κατελαβόμην τὴν πόλιν τῶν ὑδάτων.
\vs{28}Καὶ νῦν συνάγαγε τὸ κατάλοιπον τοῦ λαοῦ, καὶ παρέμβαλε ἐπὶ τὴν πόλιν, καὶ προκαταλαβοῦ αὐτὴν, ἵνα μὴ προκαταλάβωμαι ἐγὼ τὴν πόλιν, καὶ κληθῇ τὸ ὄνομά μου ἐπʼ αὐτήν.

\vs{29}Καὶ συνήγαγε Δαυὶδ πάντα τὸν λαὸν, καὶ ἐπορεύθη εἰς Ῥαββὰθ, καὶ ἐπολέμησεν ἐν αὐτῇ, καὶ κατελάβετο αὐτήν.
\vs{30}Καὶ ἔλαβε τὸν στέφανον Μολχὸμ τοῦ βασιλέως αὐτῶν ἀπὸ τῆς κεφαλῆς αὐτοῦ, καὶ ὁ σταθμὸς αὐτοῦ τάλαντον χρυσίου, καὶ λίθου τιμίου, καὶ ἦν ἐπὶ τῆς κεφαλῆς Δαυὶδ, καὶ σκῦλα τῆς πόλεως ἐξήνεγκε πολλὰ σφόδρα.
\vs{31}Καὶ τὸν λαὸν τὸν ὄντα ἐν αὐτῇ ἐξήγαγε, καὶ ἔθηκεν ἐν τῷ πρίονι, καὶ ἐν τοῖς τριβόλοις τοῖς σιδηροῖς, καὶ ὑποτομεῦσι σιδηροῖς, καὶ διήγαγεν αὐτοὺς διὰ τοῦ πλινθίου· καὶ οὕτως ἐποίησε πάσαις ταῖς πόλεσιν υἱῶν Ἀμμών· καὶ ἐπέστρεψε Δαυὶδ καὶ πᾶς ὁ λαὸς εἰς Ἰερουσαλήμ.

\ch{13}
Καὶ ἐγενήθη μετὰ ταῦτα καὶ τῷ Ἀβεσσαλὼμ υἱῷ Δαυὶδ ἀδελφὴ καλὴ τῷ εἴδει σφόδρα, καὶ ὄνομα αὐτῇ Θημὰρ, καὶ ἠγάπησεν αὐτὴν Ἀμνὼν υἱὸς Δαυίδ.
\vs{2}Καὶ ἐθλίβετο Ἀμνὼν ὥστε ἀῤῥωστεῖν διὰ Θημὰρ τὴν ἀδελφὴν αὐτοῦ, ὅτι παρθένος ἦν αὕτη, καὶ ὑπέρογκον ἐν ὀφθαλμοῖς Ἀμνὼν τοῦ ποιῆσαί τι αὐτῇ.
\vs{3}Καὶ ἦν τῷ Ἀμνὼν ἑταῖρος, καὶ ὄνομα αὐτῷ Ἰωναδὰβ, υἱὸς Σαμαὰ τοῦ ἀδελφοῦ Δαυίδ· καὶ Ἰωναδὰβ ἀνὴρ σοφὸς σφόδρα,
\vs{4}καὶ εἶπεν αὐτῷ, τί σοι ὅτι σὺ οὕτως ἀσθενὴς, υἱὲ τοῦ βασιλέως, τὸ πρωῒ πρωΐ; οὐκ ἀπαγγέλλεις μοι; καὶ εἶπεν αὐτῷ Ἀμνὼν, Θημὰρ τὴν ἀδελφὴν Ἀβεσσαλὼμ τοῦ ἀδελφοῦ μου ἐγὼ ἀγαπῶ.
\vs{5}Καὶ εἶπεν αὐτῷ Ἰωναδὰβ, κοιμήθητι ἐπὶ τῆς κοίτης σου καὶ μαλακίσθητι, καὶ εἰσελεύσεται ὁ πατήρ σου τοῦ ἰδεῖν σε, καὶ ἐρεῖς πρὸς αὐτὸν, ἐλθέτω δὴ Θημὰρ ἡ ἀδελφή μου, καὶ ψωμισάτω με, καὶ ποιησάτω κατʼ ὀφθαλμούς μου βρῶμα, ὅπως ἴδω καὶ φάγω ἐκ τῶν χειρῶν αὐτῆς.
\vs{6}Καὶ ἐκοιμήθη Ἀμνὼν καὶ ἠῤῥώστησε· καὶ εἰσῆλθεν ὁ βασιλεὺς ἰδεῖν αὐτόν· καὶ εἶπεν Ἀμνὼν πρὸς τὸν βασιλέα, ἐλθέτω δὴ Θημὰρ ἡ ἀδελφή μου πρὸς μὲ, καὶ κολλυρισάτω ἐν ὀφθαλμοῖς μου δύο κολλυρίδας, καὶ φάγομαι ἐκ τῆς χειρὸς αὐτῆς.

\vs{7}Καὶ ἀπέστειλε Δαυὶδ πρὸς Θημὰρ εἰς τὸν οἶκον, λέγων, πορεύθητι δὴ εἰς τὸν οἶκον τοῦ ἀδελφοῦ σου, καὶ ποίησον αὐτῷ βρῶμα.
\vs{8}Καὶ ἐπορεύθη Θημὰρ εἰς τὸν οἶκον Ἀμνὼν ἀδελφοῦ αὐτῆς, καὶ αὐτὸς κοιμώμενος· καὶ ἔλαβε τὸ σταῖς καὶ ἐφύρασε, καὶ ἐκολλύρισε κατʼ ὀφθαλμοὺς αὐτοῦ, καὶ ἥψησε τὰς κολλυρίδας.
\vs{9}Καὶ ἔλαβε τὸ τήγανον καὶ κατεκένωσεν ἐνώπιον αὐτοῦ, καὶ οὐκ ἠθέλησε φαγεῖν· καὶ εἶπεν Ἀμνὼν, ἐξαγάγετε πάντα ἄνδρα ἐπάνωθέν μου· καὶ ἐξήγαγον πάντα ἄνδρα ἀπὸ ἐπάνωθεν αὐτοῦ.
\vs{10}Καὶ εἶπεν Ἀμνὼν πρὸς Θημὰρ, εἰσένεγκε τὸ βρῶμα εἰς τὸ ταμιεῖον, καὶ φάγομαι ἐκ τῆς χειρός σου· καὶ ἔλαβε Θημὰρ τὰς κολλυρίδας ἃς ἐποίησε, καὶ εἰσήνεγκε τῷ Ἀμνὼν ἀδελφῷ αὐτῆς εἰς τὸν κοιτῶνα.
\vs{11}Καὶ προσήγαγεν αὐτῷ τοῦ φαγεῖν, καὶ ἐπελάβετο αὐτῆς, καὶ εἶπεν αὐτῇ, δεῦρο, κοιμήθητι μετʼ ἐμοῦ, ἀδελφή μου.
\vs{12}Καὶ εἶπεν αὐτῷ, μὴ ἀδελφέ μου, μὴ ταπεινώσῃς με, διότι οὐ ποιηθήσεται οὕτως ἐν Ἰσραὴλ, μὴ ποιήσῃς τὴν ἀφροσύνην ταύτην.
\vs{13}Καὶ ἐγὼ ποῦ ἀποίσω τὸ ὄνειδός μου; καὶ σὺ ἔσῃ ὡς εἷς τῶν ἀφρόνων ἐν Ἰσραήλ· καὶ νῦν λάλησον δὴ πρὸς τὸν βασιλέα, ὅτι οὐ μὴ κωλύσῃ με ἀπὸ σοῦ.
\vs{14}Καὶ οὐκ ἠθέλησεν Ἀμνὼν τοῦ ἀκοῦσαι τῆς φωνῆς αὐτῆς· καὶ ἐκραταίωσεν ὑπὲρ αὐτὴν, καὶ ἐταπείνωσεν αὐτήν, καὶ ἐκοιμήθη μετʼ αὐτῆς.

\vs{15}Καὶ ἐμίσησεν αὐτὴν Ἀμνὼν μῖσος μέγα σφόδρα, ὅτι μέγα τὸ μῖσος ὃ ἐμίσησεν αὐτὴν μπὲρ τὴν ἀγάπην ἣν ἠγάπησεν αὐτὴν, ὅτι μείζων ἡ κακία ἡ ἐσχάτη ἢ ἡ πρώτη· καὶ εἶπεν αὐτῇ Ἀμνὼν, ἀνάστηθι, καὶ πορεύου.
\vs{16}Καὶ εἶπεν αὐτῷ Θημὰρ περὶ τῆς κακίας τῆς μεγάλης ταύτης ὑπὲρ ἑτέραν ἣν ἐποίησας μετʼ ἐμοῦ, τοῦ ἐξαποστεῖλαί με· καὶ οὐκ ἐθέλησεν Ἀμνὼν ἀκοῦσαι τῆς φωνῆς αὐτῆς.
\vs{17}Καὶ ἐκάλεσε τὸ παιδάριον αὐτοῦ τὸν προεστηκότα τοῦ οἴκου, καὶ εἶπεν αὐτῷ, ἐξαποστείλατε δὴ ταύτην ἀπʼ ἐμοῦ ἔξω, καὶ ἀπόκλεισον τὴν θύραν ὀπίσω αὐτῆς.
\vs{18}Καὶ ἐπʼ αὐτῆς ἦν χιτὼν καρπωτὸς, ὅτι οὕτως ἐνεδιδύσκοντο αἱ θυγατέρες τοῦ βασιλέως αἱ παρθένοι τοὺς ἐπενδύτας αὐτῶν· καὶ ἐξήγαγεν αὐτὴν ὁ λειτουργὸς αὐτοῦ ἔξω, καὶ ἀπέκλεισε τὴν θύραν ὀπίσω αὐτῆς.

\vs{19}Καὶ ἔλαβε Θημὰρ σποδὸν, καὶ ἐπέθηκεν ἐπὶ τὴν κεφαλὴν αὐτῆς· καὶ τὸν χιτῶνα τὸν καρπωτὸν τὸν ἐπʼ αὐτῆς διέῤῥηξε· καὶ ἐπέθηκε τὰς χεῖρας αὐτῆς ἐπὶ τὴν κεφαλὴν αὐτῆς, καὶ ἐπορεύθη πορευομένη καὶ κράζουσα.
\vs{20}Καὶ εἶπε πρὸς αὐτὴν Ἀβεσσαλὼμ ὁ ἀδελφὸς αὐτῆς, μὴ Ἀμνὼν ὁ ἀδελφός σου ἐγένετο μετὰ σοῦ; καὶ νῦν ἀδελφή μου κώφευσον, ὅτι ἀδελφός σου ἐστί· μὴ θῇς τὴν καρδίαν σου τοῦ λαλῆσαι τὸ ῥῆμα τοῦτο· καὶ ἐκάθισε Θημὰρ χηρεύουσα ἐν τῷ οἴκῳ Ἀβεσσαλὼμ τοῦ ἀδελφοῦ αὐτῆς.

\vs{21}Καὶ ἤκουσεν ὁ βασιλεὺς Δαυὶδ πάντας τοὺς λόγους τούτους, καὶ ἐθυμώθη σφόδρα, καὶ οὐκ ἐλύπησε τὸ πνεῦμα Ἀμνὼν τοῦ υἱοῦ αὐτοῦ, ὅτι ἠγάπα αὐτὸν, ὅτι πρωτότοκος αὐτοῦ ἦν.
\vs{22}Καὶ οὐκ ἐλάλησεν Ἀβεσσαλὼμ μετὰ Ἀμνὼν ἀπὸ πονηροῦ ἕως ἀγαθοῦ, ὅτι ἐμίσει Ἀβεσσαλὼμ τὸν Ἀμνὼν ἐπὶ λόγου οὗ ἐταπείνωσε Θημὰρ τὴν ἀδελφὴν αὐτοῦ.
\vs{23}Καὶ ἐγένετο εἰς διετηρίδα ἡμερῶν, καὶ ἦσαν κείροντες τῷ Ἀβεσσαλὼμ ἐν Βελασὼρ τῇ ἐχόμενα Ἐφραὶμ, καὶ ἐκάλεσεν Ἀβεσσαλὼμ πάντας τοὺς υἱοὺς τοῦ βασιλέως.
\vs{24}Καὶ ἦλθεν Ἀβεσσαλὼμ πρὸς τὸν βασιλέα, καὶ εἶπεν, ἰδοὺ δὴ κείρουσι τῷ δούλῳ σου, πορευθήτω δὴ ὁ βασιλεὺς καὶ οἱ παῖδες αὐτοῦ μετὰ τοῦ δούλου σου.
\vs{25}Καὶ εἶπεν ὁ βασιλεὺς πρὸς Ἀβεσσαλὼμ, μὴ δὴ υἱέ μου, μὴ πορεύθῶμεν πάντες ἡμεῖς, καὶ οὐ μὴ καταβαρυνθῶμεν ἐπὶ σέ· καὶ ἐβιάσατο αὐτόν, καὶ οὐκ ἠθέλησε τοῦ πορευθῆναι, καὶ εὐλόγησεν αὐτόν.
\vs{26}Καὶ εἶπεν Ἀβεσσαλὼμ πρὸς αὐτὸν, καὶ εἰ μή, πορευθήτω δὴ μεθʼ ἡμῶν Ἀμνὼν ὁ ἀδελφός μου· καὶ εἶπεν αὐτῷ ὁ βασιλεὺς, ἱνατί πορευθῇ μετὰ σοῦ;
\vs{27}Καὶ ἐβιάσατο αὐτὸν Ἀβεσσαλώμ, καὶ ἀπέστειλε μετʼ αὐτοῦ τὸν Ἀμνὼν καὶ πάντας τοὺς υἱοὺς τοῦ βασιλέως· καὶ ἐποίησεν Ἀβεσσαλὼμ πότον κατὰ τὸν πότον τοῦ βασιλέως.

\vs{28}Καὶ ἐνετείλατο Ἀβεσσαλὼμ τοῖς παιδαρίοις αὐτοῦ, λέγων, ἴδετε ὡς ἂν ἀγαθυνθῇ ἡ καρδία Ἀμνὼν ἐν τῷ οἴνῳ, καὶ εἴπω πρὸς ὑμᾶς, πατάξατε τὸν Ἀμνὼν, καὶ θανατώσατε αὐτόν· μὴ φοβηθῆτε, ὅτι οὐχὶ ἐγώ εἰμι ὁ ἐντελλόμενος ὑμῖν; ἀνδρίζεσθε καὶ γίνεσθε εἰς υἱοὺς δυνάμεως.
\vs{29}Καὶ ἐποίησαν τὰ παιδάρια Ἀβεσσαλὼμ τῷ Ἀμνὼν, καθὰ ἐνετείλατο αὐτοῖς Ἀβεσσαλώμ· καὶ ἀνέστησαν πάντες οἱ υἱοὶ τοῦ βασιλέως, καὶ ἐπεκάθισαν ἀνὴρ ἐπὶ τὴν ἡμιονον αὐτοῦ, καὶ ἔφυγον.

\vs{30}Καὶ ἐγένετο, αὐτῶν ὄντων ἐν τῇ ὁδῷ, καὶ ἡ ἀκοὴ ἦλθε πρὸς Δαυὶδ, λέγων, ἐπάταξεν Ἀβεσσαλὼμ πάντας τοὺς υἱοὺς τοῦ βασιλέως, καὶ οὐ κατελείφθη ἐξ αὐτῶν οὐδὲ εἷς.
\vs{31}Καὶ ἀνέστη ὁ βασιλεὺς καὶ διέῤῥηξε τὰ ἱμάτια αὐτοῦ καὶ ἐκοιμήθη ἐπὶ τὴν γῆν, καὶ πάντες οἱ παῖδες αὐτοῦ οἱ περιεστῶτες αὐτῷ διέῤῥηξαν τὰ ἱμάτια αὐτῶν.
\vs{32}Καὶ ἀπεκρίθη Ἰωναδὰβ υἱὸς Σαμαὰ ἀδελφοῦ Δαυὶδ, καὶ εἶπε, μὴ εἰπάτω ὁ κύριός μου ὁ βασιλεὺς ὅτι πάντα τὰ παιδάρια τοὺς υἱοὺς τοῦ βασιλέως ἐθανάτωσεν, ὅτι Ἀμνὼν μονώτατος ἀπέθανεν, ὅτι ἐπὶ στόματος Ἀβεσσαλὼμ ἦν κείμενος ἀπὸ τῆς ἡμέρας ἧς ἐταπείνωσε Θημὰρ τὴν ἀδελφὴν αὐτοῦ.
\vs{33}Καὶ νῦν μὴ θέσθω ὁ κύριός μου ὁ βασιλεὺς ἐπὶ τὴν καρδίαν αὐτοῦ ῥῆμα, λέγων, πάντες οἱ υἱοὶ τοῦ βασιλέως ἀπέθανον· ὅτι ἀλλʼ ἢ Ἀμνὼν μονώτατος ἀπέθανε.
\vs{34}Καὶ ἀπέδρα Ἀβεσσαλώμ.

Καὶ ᾖρε τὸ παιδάριον ὁ σκοπὸς τοὺς ὀφθαλμοὺς αὐτοῦ, καὶ εἶδε· καὶ ἰδοὺ λαὸς πολὺς πορευόμενος ἐν τῇ ὁδῷ ὄπισθεν αὐτοῦ ἐκ πλευρᾶς τοῦ ὄρους ἐν τῇ καταβάσει· καὶ παρεγένετο ὁ σκοπὸς καὶ ἀπήγγειλε τῷ βασιλεῖ, καὶ εἶπεν, ἄνδρας ἑώρακα ἐκ τῆς ὁδοῦ τὴς Ὠρωνῆν ἐκ μέρους τοῦ ὄρους.
\vs{35}Καὶ εἶπεν Ἰωναδὰβ πρὸς τὸν βασιλέα, ἰδοὺ οἱ υἱοὶ τοῦ βασιλέως πάρεισι· κατὰ τὸν λόγον τοῦ δούλου σου, οὕτως ἐγένετο.
\vs{36}Καὶ ἐγένετο ἡνίκα συνετέλεσε λαλῶν, καὶ ἰδοὺ οἱ υἱοὶ τοῦ βασιλέως ἦλθον, καὶ ἐπῇραν τὴν φωνὴν αὐτῶν καὶ ἔκλαυσαν· καί γε ὁ βασιλεὺς καὶ πάντες οἱ παῖδες αὐτοῦ ἔκλαυσαν κλαυθμὸν μέγαν σφόδρα.

\vs{37}Καὶ Ἀβεσσαλὼμ ἔφυγε, καὶ ἐπορεύθη, πρὸς Θολμὶ υἱὸν Ἐμιοὺδ βασιλέα Γεδσοὺρ εἰς γῆν Χαμααχάδ· καὶ ἐπένθησεν ὁ βασιλεὺς Δαυὶδ ἐπὶ τὸν υἱὸν αὐτοῦ πάσας τὰς ἡμέρας.
\vs{38}Καὶ Ἀβεσσαλὼμ ἀπέδρα, καὶ ἐπορεύθη εἰς Γεδσοὺρ, καὶ ἦν ἐκεῖ ἔτη τρία.
\vs{39}Καὶ ἐκόπασεν ὁ βασιλεὺς Δαυὶδ τοῦ ἐξελθεῖν πρὸς Ἀβεσσαλὼμ, ὅτι παρεκλήθη ἐπὶ Ἀμνὼν, ὅτι ἀπέθανε.

\ch{14}
Καὶ ἔγνω Ἰωὰβ υἱὸς Σαρουίας ὅτι ἡ καρδία τοῦ βασιλέως ἐπὶ Ἀβεσσαλώμ.
\vs{2}Καὶ ἀπέστειλεν Ἰωὰβ εἰς Θεκωὲ, καὶ ἔλαβεν ἐκεῖθεν γυναῖκα σοφὴν, καὶ εἶπε πρὸς αὐτὴν, πένθησον δὴ, καὶ ἔνδυσαι ἱμάτια πενθικὰ, καὶ μὴ ἀλείψῃ ἔλαιον, καὶ ἔσῃ ὡς γυνὴ πενθοῦσα ἐπὶ τεθνηκότι τοῦτο ἡμέρας πολλὰς,
\vs{3}καὶ ἐλεύσῃ πρὸς τὸν βασιλέα, καὶ λαλήσεις πρὸς αὐτὸν κατὰ τὸ ῥῆμα τοῦτο· καὶ ἔθηκεν Ἰωὰβ τοὺς λόγους ἐν τῷ στόματι αὐτῆς.

\vs{4}Καὶ εἰσῆλθεν ἡ γυνὴ ἡ Θεκωῖτις πρὸς τὸν βασιλέα, καὶ ἔπεσεν ἐπὶ πρόσωπον αὐτῆς εἰς τὴν γῆν, καὶ προσεκύνησεν αὐτῷ, καὶ εἶπε, σῶσον βασιλεῦ, σῶσον.
\vs{5}Καὶ εἶπε πρὸς αὐτὴν ὁ βασιλεὺς, τί ἐστι σοι;

\vs{6}Ἡ δὲ εἶπε, καὶ μάλα γυνὴ χήρα ἐγώ εἰμι, καὶ ἀπέθανεν ὁ ἀνήρ μου. Καί γε τῇ δούλῃ σου δύο υἱοὶ, καὶ ἐμαχέσαντο ἀμφότεροι ἐν τῷ ἀγρῷ, καὶ οὐκ ἦν ὁ ἐξαιρούμενος ἀναμέσον αὐτῶν· καὶ ἔπαισεν ὁ εἷς τὸν ἕνα ἀδελφὸν αὐτοῦ, καὶ ἐθανάτωσεν αὐτόν.
\vs{7}Καὶ ἰδοὺ ἐπανέστη ὅλη ἡ πατριὰ πρὸς τὴν δούλην σου, καὶ εἶπαν, δὸς τὸν παίσαντα τὸν ἀδελφὸν αὐτοῦ, καὶ θανατώσομεν αὐτὸν ἀντὶ τῆς ψυχῆς τοῦ ἀδελφοῦ αὐτοῦ οὗ ἀπέκτεινε, καὶ ἐξαροῦμεν καί γε τὸν κληρονόμον ὑμῶν· καὶ σβέσουσι τὸν ἄνθρακά μου τὸν καταλειφθέντα ὥστε μὴ θέσθαι τῷ ἀνδρί μου κατάλειμμα καὶ ὄνομα ἐπὶ προσώπου τῆς γῆς.

\vs{8}Καὶ εἶπεν ὁ βασιλεὺς πρὸς τὴν γυναῖκα, ὑγιαίνουσα βάδιζε εἰς τὸν οἶκόν σου, κᾀγὼ ἐντελοῦμαι περὶ σοῦ.
\vs{9}Καὶ εἶπεν ἡ γυνὴ ἡ Θεκωΐτις πρὸς τὸν βασιλέα, ἐπʼ ἐμὲ, κύριέ μου βασιλεῦ, ἡ ἀνομία καὶ ἐπὶ τὸν οἶκον τοῦ πατρός μου, καὶ ὁ βασιλεὺς καὶ ὁ θρόνος αὐτοῦ ἀθῶος.
\vs{10}Καὶ εἶπεν ὁ βασιλεύς, τίς ὁ λαλῶν πρὸς σὲ, καὶ ἄξεις αὐτὸν πρὸς ἐμὲ, καὶ οὐ προσθήσει ἔτι ἅψασθαι αὐτοῦ;
\vs{11}Καὶ εἶπε, μνημονευσάτω δὴ ὁ βασιλεὺς τὸν Κύριον Θεὸν αὐτοῦ πληθυνθῆναι ἀγχιστέα τοῦ αἵματος τοῦ διαφθεῖραι, καὶ οὐ μὴ ἐξάρωσι τὸν υἱόν μου· καὶ εἶπε, ζῇ Κύριος, εἶ πεσεῖται ἀπὸ τῆς τριχὸς τοῦ υἱοῦ σου ἐπὶ τὴν γῆν.

\vs{12}Καὶ εἶπεν ἡ γυνὴ, λαλησάτω δὴ ἡ δούλη σου πρὸς τὸν κύριόν μου βασιλέα ῥῆμα· καὶ εἶπε, λάλησον.
\vs{13}Καὶ εἶπεν ἡ γυνὴ, ἱνατί ἐλογίσω τοιοῦτο ἐπὶ λαὸν Θεοῦ; ἢ ἐκ στόματος τοῦ βασιλέως ὁ λόγος οὗτος ὡς πλημμέλεια, τοῦ μὴ ἐπιστρέψαι τὸν βασιλέα τὸν ἐξωσμένον αὐτοῦ;
\vs{14}Ὅτι θανάτῳ ἀποθανούμεθα, καὶ ὥσπερ τὸ ὕδωρ τὸ καταφερόμενον ἐπὶ τῆς γῆς ὃ οὐ συναχθήσεται, καὶ λήψεται ὁ Θεὸς ψυχὴν, καὶ λογιζόμενος τοῦ ἐξῶσαι ἀπʼ αὐτοῦ ἐξεωσμένον.
\vs{15}Καὶ νῦν ὃ ἦλθον λαλῆσαι πρὸς τὸν βασιλέα τὸν κύριόν μου τὸ ῥῆμα τοῦτο, ὅτι ὄψεταί με ὁ λαὸς, καὶ ἐρεῖ ἡ δούλη σου, λαλησάτω δὴ πρὸς τὸν κύριόν μον τὸν βασιλέα, εἴπως ποιήσει ὁ βασιλεὺς τὸ ῥῆμα τῆς δούλης αὐτοῦ
\vs{16}ὅτι ἀκούσει ὁ βασιλεύς· ῥυσάσθω τὴν δούλην αὐτοῦ ἐκ χειρὸς τοῦ ἀνδρὸς τοῦ ζητοῦντος ἐξᾶραί με καὶ τὸν υἱόν μου ἀπὸ κληρονομίας Θεοῦ.
\vs{17}Καὶ εἶπεν ἡ γυνή, εἰ ἤδη ὁ λόγος τοῦ κυρίου μου τοῦ βασιλέως εἰς θυσίας· ὅτι καθὼς ἄγγελος Θεοῦ, οὕτως ὁ κύριός μου ὁ βασιλεὺς, τοῦ ἀκούειν τὸ ἀγαθὸν καὶ τὸ πονηρόν· καὶ Κύριος ὁ Θεός σου ἔσται μετὰ σοῦ.

\vs{18}Καὶ ἀπεκρίθη ὁ βασιλεὺς, καὶ εἶπε πρὸς τὴν γυναῖκα, μὴ δὴ κρύψῃς ἀπʼ ἐμοῦ ῥῆμα, ὃ ἐγὼ ἐπερωτῶ σε· καὶ εἶπεν ἡ γυνή, λαλησάτω δὴ ὁ κύριός μου ὁ βασιλεύς.
\vs{19}Καὶ εἶπεν ὁ βασιλεὺς, μὴ ἡ χεὶρ Ἰωὰβ ἐν παντὶ τούτῳ μετὰ σοῦ; καὶ εἶπεν ἡ γυνὴ τῷ βασιλεῖ, ζῇ ἡ ψυχή σου, κύριέ μου βασιλεῦ, εἰ ἔστιν εἰς τὰ δεξιὰ ἢ εἰς τὰ ἀριστερὰ ἐκ πάντων ὧν ἐλάλησεν ὁ κύριός μου ὁ βασιλεὺς, ὅτι ὁ δοῦλός σου Ἰωὰβ αὐτὸς ἐνετείλατό μοι, καὶ αὐτὸς ἔθετο ἐν τῷ στόματι τῆς δούλης σου πάντας τοὺς λόγους τούτους.
\vs{20}Ἕνεκεν τοῦ περιελθεῖν τὸ πρόσωπον τοῦ ῥήματος τούτου, ὃ ἐποίησεν ὁ δοῦλός σου Ἰωὰβ τὸν λόγον τοῦτον· καὶ ὁ κύριός μου σοφὸς καθὼς σοφία ἀγγέλου τοῦ Θεοῦ, τοῦ γνῶναι πάντα τὰ ἐν τῇ γῇ.

\vs{21}Καὶ εἶπεν ὁ βασιλεὺς πρὸς Ἰωὰβ, ἰδοὺ δὴ ἐποίησά σοι κατὰ τὸν λόγον σου τοῦτον· πορεύου, ἐπιστρέψον τὸ παιδάριον τὸν Ἀβεσσαλώμ.
\vs{22}Καὶ ἔπεσεν Ἰωὰβ ἐπὶ πρόσωπον αὐτοῦ ἐπὶ τὴν γῆν, καὶ προσεκύνησε, καὶ εὐλόγησε τὸν βασιλέα· καὶ εἶπεν Ἰωὰβ, σήμερον ἔγνω ὁ δοῦλός σου ὅτι εὗρον χάριν ἐν ὀφθαλμοῖς σου, κύριέ μου βασιλεῦ, ὅτι ἐποίησεν ὁ κύριός μου ὁ βασιλεὺς τὸν λόγον τοῦ δούλου αὐτοῦ.
\vs{23}Καὶ ἀνέστη Ἰωὰβ, καὶ ἐπορεύθη εἰς Γεδσούρ, καὶ ἤγαγε τὸν Ἀβεσσαλὼμ εἰς Ἰερουσαλήμ.
\vs{24}Καὶ εἶπεν ὁ βασιλεὺς, ἀποστραφήτω εἰς τὸν οἶκον αὐτοῦ, καὶ τὸ πρόσωπόν μου μὴ βλεπέτω· καὶ ἀπέστρεψεν Ἀβεσσαλὼμ εἰς τὸν οἶκον αὐτοῦ, καὶ τὸ πρόσωπον τοῦ βασιλέως οὐκ εἶδε.

\vs{25}Καὶ ὡς Ἀβεσσαλὼμ οὐκ ἦν ἀνὴρ ἐν παντὶ Ἰσραὴλ αἰνετὸς σφόδρα, ἀπὸ ἴχνους ποδὸς αὐτοῦ καὶ ἕως κορυφῆς αὐτοῦ οὐκ ἦν ἐν αὐτῷ μῶμος.
\vs{26}Καὶ ἐν τῷ κείρεσθαι αὐτὸν τὴν κεφαλὴν αὐτοῦ, καὶ ἐγένετο ἀπʼ ἀρχῆς ἡμερῶν εἰς ἡμέρας ὡς ἂν ἐκείρετο, ὅτι κατεβαρύνετο ἐπʼ αὐτὸν, καὶ κειρόμενος αὐτὴν ἔστησε τὴν τρίχα τῆς κεφαλῆς αὐτοῦ διακοσίους σίκλους ἐν τῷ σίκλῳ τῷ βασιλικῷ.
\vs{27}Καὶ ἐτέχθησαν τῷ Ἀβεσσαλὼμ τρεῖς υἱοὶ, καὶ θυγάτηρ μία, καὶ ὄνομα αὐτῇ Θημάρ· αὕτη ἦν γυνὴ καλὴ σφόδρα· καὶ γίνεται γυνὴ Ῥοβοὰμ υἱῷ Σαλωμὼν, καὶ τίκτει αὐτῷ τὸν Ἀβία.

\vs{28}Καὶ ἐκάθισεν Ἀβεσσαλὼμ ἐν Ἱερουσαλὴμ δύο ἔτη ἡμερῶν, καὶ τὸ πρόσωπον τοῦ βασιλέως οὐκ εἶδε.
\vs{29}Καὶ ἀπέστειλεν Ἀβεσσαλὼμ πρὸς Ἰωὰβ ἀποστεῖλαι αὐτὸν πρὸς τὸν βασιλέα, καὶ οὐκ ἠθέλησεν ἐλθεῖν πρὸς αὐτόν· καὶ ἀπέστειλεν ἐκ δευτέρου πρὸς αὐτὸν, καὶ οὐκ ἠθέλησε παραγενέσθαι.
\vs{30}Καὶ εἶπεν Ἀβεσσαλὼμ πρὸς τοὺς παῖδας αὐτοῦ, ἴδετε, ἡ μερὶς ἐν ἀγρῷ τοῦ Ἰωὰβ ἐχόμενά μου, καὶ αὐτῷ ἐκεῖ κριθαὶ, πορεύεσθε καὶ ἐμπρήσατε αὐτὴν ἐν πυρί· καὶ ἐνέπρησαν οἱ παῖδες Ἀβεσσαλὼμ τὴν μερίδα· καὶ παραγίνονται οἱ δοῦλοι Ἰωὰβ πρὸς αὐτὸν διεῤῥηχότες τὰ ἱμάτια αὐτῶν, καὶ εἶπον, ἐνεπύρισαν οἱ δοῦλοι Ἀβεσσαλὼμ τὴν μερίδα ἐν πυρί.
\vs{31}Καὶ ἀνέστη Ἰωὰβ, καὶ ἦλθε πρὸς Ἀβεσσαλὼμ εἰς τὸν οἶκον, καὶ εἶπε πρὸς αὐτὸν, ἱνατί ἐνεπύρισαν οἱ παῖδές σου τὴν μερίδα τὴν ἐμὴν ἐν πυρί;
\vs{32}Καὶ εἶπεν Ἀβεσσαλὼμ πρὸς Ἰωὰβ, ἰδοὺ ἀπέστειλα πρὸς σὲ, λέγων, ἧκε ὧδε, καὶ ἀποστελῶ σε πρὸς τὸν βασιλέα, λέγων, ἱνατί ἦλθον ἐκ Γεδσούρ; ἀγαθόν μοι ἦν εἶναι ἐκεῖ· καὶ νῦν ἰδοὺ τὸ πρόσωπον τοῦ βασιλέως οὐκ εἶδον· εἰ δέ ἐστιν ἐν ἐμοὶ ἀδικία, καὶ θανάτωσόν με.

\vs{33}Καὶ εἰσῆλθεν Ἰωὰβ πρὸς τὸν βασιλέα, καὶ ἀπήγγειλεν αὐτῷ· καὶ ἐκάλεσε τὸν Ἀβεσσαλώμ· καὶ εἰσῆλθε πρὸς τὸν βασιλέα, καὶ προσεκύνησεν αὐτῷ, καὶ ἔπεσεν ἐπὶ πρόσωπον αὐτοῦ ἐπὶ τὴν γῆν, καὶ κατὰ πρόσωπον τοῦ βασιλέως· καὶ κατεφίλησεν ὁ βασιλεὺς τὸν Ἀβεσσαλώμ.

\ch{15}
Καὶ ἐγένετο μετὰ ταῦτα καὶ ἐποίησεν ἑαυτῷ Ἀβεσσαλὼμ ἅρματα, καὶ ἵππους, καὶ πεντήκοντα ἄνδρας παρατρέχειν ἔμπροσθεν αὐτοῦ.
\vs{2}Καὶ ὤρθρισεν Ἀβεσσαλὼμ, καὶ ἔστη ἀνὰ χεῖρα τῆς ὁδοῦ τῆς πύλης· καὶ ἐγένετο πᾶς ἀνὴρ ᾧ ἐγένετο κρίσις, ἦλθε πρὸς τὸν βασιλέα εἰς κρίσιν, καὶ ἐβόησε πρὸς αὐτὸν Ἀβεσσαλὼμ, καὶ ἔλεγεν αὐτῷ, ἐκ ποίας πόλεως σὺ εἶ; καὶ εἶπεν, ἐκ μιᾶς φυλῶν Ἰσραὴλ ὁ δοῦλός σου.
\vs{3}Καὶ εἶπε πρὸς αὐτὸν ὁ Ἀβεσσαλὼμ, ἰδοὺ οἱ λόγοι σου ἀγαθοὶ καὶ εὔκολοι, καὶ ὁ ἀκούων οὐκ ἔστι σοι παρὰ τοῦ βασιλέως.
\vs{4}Καὶ εἶπεν Ἀβεσσαλὼμ, τίς με καταστήσει κριτὴν ἐν τῇ γῇ, καὶ ἐπʼ ἐμὲ ἐλεύσεται πᾶς ἀνὴρ ᾧ ἐὰν ᾖ ἀντιλογία καὶ κρίσις, καὶ δικαιώσω αὐτόν;
\vs{5}Καὶ ἐγένετο ἐν τῷ ἐγγίζειν ἄνδρα τοῦ προσκυνῆσαι αὐτῷ, καὶ ἐξέτεινε τὴν χεῖρα αὐτοῦ, καὶ ἐπελαμβάνετο αὐτοῦ, καὶ κατεφίλησεν αὐτόν.
\vs{6}Καὶ ἐποίησεν Ἀβεσσαλὼμ κατὰ τὸ ῥῆμα τοῦτο παντὶ Ἰσραὴλ τοῖς παραγινομένοις εἰς κρίσιν πρὸς τὸν βασιλέα, καὶ ἰδιοποιεῖτο Ἀβεσσαλὼμ τὴν καρδίαν ἀνδρῶν Ἰσραήλ.

\vs{7}Καὶ ἐγένετο ἀπὸ τέλους τεσσαράκοντα ἐτῶν, καὶ εἶπεν Ἀβεσσαλὼμ πρὸς τὸν πατέρα αὐτοῦ, πορεύσομαι δὴ, καὶ ἀποτίσω τὰς εὐχάς μου, ἃς ηὐξάμην τῷ Κυρίῳ ἐν Χεβρών·
\vs{8}Ὅτι εὐχὴν ηὔξατο ὁ δοῦλός σου ἐν τῷ οἰκεῖν με ἐν Γεδσοὺρ ἐν Συρίᾳ, λέγων, ἐὰν ἐπιστρέφων ἐπιστρέψῃ με Κύριος εἰς Ἰερουσαλὴμ, καὶ λατρεύσω τῷ Κυρίῳ.
\vs{9}Καὶ εἶπεν αὐτῷ ὁ βασιλεὺς, Βάδιζε εἰς εἰρήνην· καὶ ἀναστὰς ἐπορεύθη εἰς Χεβρών.

\vs{10}Καὶ ἀπέστειλεν Ἀβεσσαλὼμ κατασκόπους ἐν πάσαις φυλαῖς Ἰσραὴλ, λέγων, ἐν τῷ ἀκοῦσαι ὑμᾶς τὴν φωνὴν τῆς κερατίνης, καὶ ἐρεῖτε, βεβασίλευκε βασιλεὺς Ἀβεσσαλὼμ ἐν Χεβρών.
\vs{11}Καὶ μετὰ Ἀβεσσαλὼμ ἐπορεύθησαν διακόσιοι ἄνδρες ἐξ Ἰερουσαλὴμ κλητοί· καὶ πορευόμενοι τῇ ἁπλότητι αὐτῶν, καὶ οὐκ ἔγνωσαν πᾶν ῥῆμα.
\vs{12}Καὶ ἀπέστειλεν Ἀβεσσαλὼμ τῷ Ἀχειτόφελ τῷ Θεκωνεὶ, σύμβουλον Δαυὶδ, ἐκ τῆς πόλεως αὐτοῦ εἰς Γωλὰ, ἐν τῷ θυσιάζειν αὐτόν· καὶ ἐγένετο σύντριμμα ἰσχυρόν· καὶ ὁ λαὸς ὁ πορευόμενος καὶ πολὺς μετὰ Ἀβεσσαλώμ.

\vs{13}Καὶ παρεγένετο ἀπαγγέλων πρὸς Δαυὶδ, λέγων, ἐγενήθη ἡ καρδία ἀνδρῶν Ἰσραὴλ ὀπίσω Ἀβεσσαλώμ.
\vs{14}Καὶ εἶπε Δαυὶδ πᾶσι τοῖς παισὶν αὐτοῦ τοῖς μετʼ αὐτοῦ τοῖς ἐν Ἰερουσαλὴμ, ἀνάστητε καὶ φύγωμεν, ὅτι οὐκ ἔστιν ἡμῖν σωτηρία ἀπὸ προσώπου Ἀβεσσαλώμ· ταχύνατε τοῦ πορευθῆναι, ἵνα μὴ ταχύνῃ καὶ καταλάβῃ ἡμᾶς, καὶ ἐξώσῃ ἐφʼ ἡμᾶς τὴν κακίαν, καὶ πατάξῃ τὴν πόλιν ἐν στόματι μαχαίρης.
\vs{15}Καὶ εἶπον οἱ παῖδες τοῦ βασιλέως πρὸς τὸν βασιλέα, κατὰ πάντα ὅσα αἱρεῖται ὁ κύριος ἡμῶν ὁ βασιλεὺς, ἰδοὺ οἱ παῖδές σου.

\vs{16}Καὶ ἐξῆλθεν ὁ βασιλεὺς καὶ πᾶς ὁ οἶκος αὐτοῦ τοῖς ποσὶν αὐτῶν· καὶ ἀφῆκεν ὁ βασιλεὺς δέκα γυναῖκας τῶν παλλακῶν αὐτοῦ φυλάσσειν τὸν οἶκον.
\vs{17}Καὶ ἐξῆλθεν ὁ βασιλεὺς καὶ πάντες οἱ παῖδες αὐτοῦ πεζῇ, καὶ ἔστησαν ἐν οἴκῳ τῷ μακράν.
\vs{18}Καὶ πάντες οἱ παῖδες αὐτοῦ ἀνὰ χεῖρα αὐτοῦ παρῆγον, καὶ πᾶς Χελεθὶ καὶ πᾶς ὁ Φελεθὶ, καὶ ἔστησαν ἐπὶ τῆς ἐλαίας ἐν τῇ ἐρήμῳ· καὶ πᾶς ὁ λαὸς παρεπορεύετο ἐχόμενος αὐτοῦ, καὶ πάντες οἱ περὶ αὐτὸν, καὶ πάντες οἱ ἁδροὶ, καὶ πάντες οἱ μαχηταὶ ἑξακόσιοι ἄνδρες· καὶ παρῆσαν ἐπὶ χεῖρα αὐτοῦ· καὶ πᾶς ὁ Χελεθι, καὶ πᾶς ὁ Φελεθὶ, καὶ πάντες οἱ Γεθαῖοι οἱ ἑξακόσιοι ἄνδρες οἱ ἐλθόντες τοῖς ποσὶν αὐτῶν ἐκ Γὲθ, καὶ πορευόμενοι ἐπὶ πρόσωπον τοῦ βασιλέως.

\vs{19}Καὶ εἶπεν ὁ βασιλεὺς πρὸς Ἐθὶ τὸν Γεθαῖον, ἱνατί πορεύῃ καὶ σὺ μεθʼ ἡμῶν; ἐπίστρεφε, καὶ οἴκει μετὰ τοῦ βασιλέως, ὅτι ξένος εἶ σὺ, καὶ ὅτι μετῴκηκας σὺ ἐκ τοῦ τόπου σου.
\vs{20}Εἰ ἐχθὲς παραγέγονας, καὶ σήμερον κινήσω σε μεθʼ ἡμῶν; καί γε μεταναστήσεις τὸν τόπον σου· χθὲς ἡ ἐξέλευσίς σου, καὶ σήμερον μετακινήσω σε μεθʼ ἡμῶν τοῦ πορευθῆναι; καὶ ἐγὼ πορεύσομαι οὗ ἐὰν ἐγὼ πορευθῶ· ἐπιστρέφου καὶ ἐπίστρεψον τοὺς ἀδελφούς σου μετὰ σοῦ, καὶ Κύριος ποιήσει μετὰ σοῦ ἔλεος καὶ ἀλήθειαν.
\vs{21}Καὶ ἀπεκρίθη Ἐθὶ τῷ βασιλεῖ, καὶ εἶπε, ζῇ Κύριος καὶ ζῇ ὁ κύριός μου ὁ βασιλεὺς, ὅτι εἰς τὸν τόπον οὗ ἐὰν ᾖ ὁ κύριός μου, καὶ ἐὰν εἰς θάνατον καὶ ἐὰν εἰς ζωὴν, ὅτι ἐκεῖ ἔσται ὁ δοῦλός σου.
\vs{22}Καὶ εἶπεν ὁ βασιλεὺς πρὸς Ἐθὶ, δεῦρο, καὶ διάβαινε μετʼ ἐμοῦ· καὶ παρῆλθεν Ἐθὶ ὁ Γεθαῖος καὶ ὁ βασιλεὺς, καὶ πάντες οἱ παῖδες αὐτοῦ καὶ πᾶς ὁ ὄχλος ὁ μετʼ αὐτοῦ.

\vs{23}Καὶ πᾶσα ἡ γῆ ἔκλαιε φωνῇ μεγάλῃ· καὶ πᾶς ὁ λαὸς παρεπορεύοντο ἐν τῷ χειμάῤῥῳ τῶν Κέδρων· καὶ ὁ βασιλεὺς διέβη τὸν χειμάῤῥουν Κέδρων· καὶ πᾶς ὁ λαὸς καὶ ὁ βασιλεὺς παρεπορεύοντο ἐπὶ πρόσωπον ὁδοῦ τὴν ἔρημον.

\vs{24}Καὶ ἰδοὺ καί γε Σαδὼκ καὶ πάντες οἱ Λευῖται μετʼ αὐτοῦ, αἴροντες τὴν κιβωτὸν διαθήκης Κυρίου ἀπὸ Βαιθάρ· καὶ ἔστησαν τὴν κιβωτὸν τοῦ Θεοῦ· καὶ ἀνέβη Ἀβιάθαρ ἕως ἐπαύσατο πᾶς ὁ λαὸς παρελθεῖν ἐκ τῆς πόλεως.
\vs{25}Καὶ εἶπεν ὁ βασιλεὺς πρὸς τὸν Σαδὼκ, ἀπόστρεψον τὴν κιβωτὸν τοῦ Θεοῦ εἰς τὴν πόλιν· ἐὰν εὕρω χάριν ἐν ὀφθαλμοῖς Κυρίου, καὶ ἐπιστρέψει με, καὶ δείξει μοι αὐτὴν καὶ τὴν εὐπρέπειαν αὐτῆς.
\vs{26}Καὶ ἐὰν εἴπῃ οὕτως, οὐκ ἠθέληκα ἐν σοί· ἰδοὺ ἐγώ εἰμι, ποιείτω μοι κατὰ τὸ ἀγαθὸν ἐν ὀφθαλμοῖς αὐτοῦ.

\vs{27}Καὶ εἶπεν ὁ βασιλεὺς τῷ Σαδὼκ τῷ ἱερεῖ, ἴδετε, σὺ ἐπιστρέφεις εἰς τὴν πόλιν ἐν εἰρήνῃ, καὶ Ἀχιμάας ὁ υἱός σου, καὶ Ἰωνάθαν ὁ υἱὸς Ἀβιάθαρ οἱ δύο υἱοὶ ὑμῶν μεθʼ ὑμῶν.
\vs{28}Ἴδετε, ἐγώ εἰμι στρατεύομαι ἐν Ἀραβὼθ τῆς ἐρήμου, ἕως τοῦ ἐλθεῖν ῥῆμα παρʼ ὑμῶν τοῦ ἀπαγγεῖλαί μοι.
\vs{29}Καὶ ἀπέστρεψε Σαδὼκ καὶ Ἀβιάθαρ τὴν κιβωτὸν τοῦ Θεοῦ εἰς Ἰερουσαλὴμ, καὶ ἐκάθισεν ἐκεῖ.

\vs{30}Καὶ Δαυὶδ ἀνέβαινεν ἐν τῇ ἀναβάσει τῶν ἐλαιῶν ἀναβαίνων καὶ κλαίων, καὶ τὴν κεφαλὴν ἐπικεκαλυμμένος, καὶ αὐτὸς ἐπορεύετο ἀνυπόδετος· καὶ πὰς ὁ λαὸς ὁ μετʼ αὐτοῦ ἐπεκάλυψεν ἀνὴρ τὴν κεφαλὴν αὐτοῦ· καὶ ἀνέβαινον ἀναβαίνοντες καὶ κλαίοντες.
\vs{31}Καὶ ἀνηγγέλη Δαυὶδ, λέγοντες, καὶ Ἀχιτόφελ ἐν τοῖς συστρεφομένοις μετὰ Ἀβεσσαλώμ· καὶ εἶπε Δαυὶδ, διασκέδασον δὴ τὴν βουλὴν Ἀχιτόφελ, Κύριε ὁ Θεός μου.

\vs{32}Καὶ ἦν Δαυὶδ ἐρχόμενος ἕως τοῦ Ῥὼς, οὗ προσεκύνησεν ἐκεῖ τῷ Θεῷ· καὶ ἰδοὺ εἰς ἀπαντὴν αὐτῷ Χουσὶ ὁ ἀρχιεταῖρος Δαυίδ διεῤῥηχὼς τὸν χιτῶνα αὐτοῦ, καὶ γῆ ἐπὶ τῆς κεφαλῆς αὐτοῦ.
\vs{33}Καὶ εἶπεν αὐτῷ Δαυὶδ ἐὰν μὲν διαβῇς μετʼ ἐμοῦ, καὶ ἔσῃ ἐπʼ ἐμὲ εἰς βάσταγμα·
\vs{34}καὶ ἐὰν ἐπιστρέψῃς ἐπὶ τὴν πόλιν, καὶ ἐρεῖς τῷ Ἀβεσσαλὼμ, διεληλύθασιν οἱ ἀδελφοί σου, καὶ ὁ βασιλεὺς κατόπισθέ μου διελήλυθεν ὁ πατήρ σου· καὶ νῦν παῖς σού εἰμὶ, βασιλεῦ, ἔασόν με ζῆσαι· παῖς τοῦ πατρὸς σου ἤμην τότε καὶ ἀρτίως, καὶ νῦν ἐγὼ δοῦλος σός· καὶ διασκεδάσεις μοι τὴν βουλὴς Ἀχιτόφελ.
\vs{35}Καὶ ἰδοὺ ἐκεῖ μετὰ σοῦ Σαδὼκ καὶ Ἀβιάθαρ οἱ ἱερεῖς· καὶ ἔσται πὰν ῥῆμα ὁ ἐὰν ἀκούσῃς ἐξ οἴκου τοῦ βασιλέως, καὶ ἀπαγγελεῖς τῷ Σαδὼκ καὶ τῷ Ἀβιάθαρ τοῖς ἱερεῦσιν.
\vs{36}Ἰδοὺ ἐκεῖ μετʼ αὐτῶν δύο υἱοὶ αὐτῶν, Ἀχιμάας υἱὸς τῷ Σαδὼκ, καὶ Ἰωνάθαν υἱὸς τῷ Ἀβιάθαρ· καὶ ἀποστελεῖτε ἐν χειρὶ αὐτῶν πρὸς μὲ πὰν ῥῆμα ὃ ἐὰν ἀκούσητε.
\vs{37}Καὶ εἰσῆλθε Χουσὶ ὁ ἐταῖρος Δαυὶδ εἰς τὴν πόλιν, καὶ Ἀβεσσαλὼμ ἄρτι εἰσεπορεύετο εἰς Ἰερουσαλήμ.

\ch{16}
Καὶ Δαυὶδ παρῆλθε βραχύ τι ἀπὸ τῆς Ῥὼς, καὶ ἰδοὺ Σιβὰ τὸ παιδάριον Μεμφιβοσθὲ εἰς ἀπαντὴν αὐτοῦ· καὶ ζεῦγος ὄνων ἐπισεσαγμένων, καὶ ἐπʼ αὐτοῖς διακόσιοι ἄρτοι, καὶ ἑκατὸν σταφίδες, καὶ ἑκατὸν φοίνικες καὶ νέβελ οἴνου.
\vs{2}Καὶ εἶπεν ὁ βασιλεὺς πρὸς Σιβὰ, τί ταῦτά σοι; καὶ εἶπε Σιβὰ, τὰ ὑποζύγια τῇ οἰκίᾳ τοῦ βασιλέως τοῦ ἐπικαθῆσθαι, καὶ οἱ ἄρτοι καὶ οἱ φοίνικες εἰς βρῶσιν τοῖς παιδαρίοις, καὶ ὁ οἶνος πιεῖν τοῖς ἐκλελυμένοις ἐν τῇ ἐρήμῳ.
\vs{3}Καὶ εἶπεν ὁ βασιλεὺς, καὶ ποῦ ὁ υἱὸς τοῦ κυρίου σου; καὶ εἶπε Σιβὰ πρὸς τὸν βασιλέα, ἰδοὺ κάθηται ἐν Ἱερουσαλὴμ, ὅτι εἶπε, σήμερον ἐπιστρέψουσί μοι οἶκος Ἰσραὴλ τὴν βασιλείαν τοῦ πατρός μου.
\vs{4}Καὶ εἶπεν ὁ βασιλεὺς τῷ Σιβᾷ, ἰδού σοι πάντα ὅσα ἐστὶ Μεμφιβοσθέ· καὶ εἶπε Σιβὰ προσκυνήσας, εὕροιμι χάριν ἐν ὀφθαλμοῖς σου κύριέ μου βασιλεῦ.

\vs{5}Καὶ ἦλθεν ὁ βασιλεὺς Δαυὶδ ἕως Βαουρίμ· καὶ ἰδοὺ ἐκεῖθεν ἀνὴρ ἐξεπορεύετο ἐκ συγγενείας οἴκου Σαοὺλ, καὶ ὄνομα αὐτῷ Σεμεῒ υἱὸς Γηρά· ἐξῆλθε ἐκπορευόμενος καὶ καταρώμενος,
\vs{6}καὶ λιθάζων ἐν λίθοις τὸν Δαυὶδ, καὶ πάντας τοὺς παῖδας τοῦ βασιλέως Δαυίδ· καὶ πᾶς ὁ λαὸς ἦν, καὶ πάντες οἱ δυνατοὶ, ἐκ δεξιῶν καὶ ἐξ εὐωνύμων τοῦ βασιλέως.
\vs{7}Καὶ οὕτως ἔλεγε Σεμεῒ ἐν τῷ καταρᾶσθαι αὐτὸν, ἔξελθε ἔξελθε ἀνὴρ αἱμάτων καὶ ἀνὴρ ὁ παράνομος.
\vs{8}Ἐπέστρεψεν ἐπὶ σὲ Κύριος πάντα τὰ αἵματα τοῦ οἴκου Σαοὺλ, ὅτι ἐβασίλευσας ἀντʼ αὐτοῦ· καὶ ἔδωκε Κύριος τὴν βασιλείαν ἐν χειρὶ Ἀβεσσαλὼμ τοῦ υἱοῦ σου· καὶ ἰδοὺ σὺ ἐν τῇ κακίᾳ σου, ὅτι ἀνὴρ αἱμάτων σύ.

\vs{9}Καὶ εἶπεν Ἀβεσσὰ υἱὸς Σαρουίας πρὸς τὸν βασιλέα, ἱνατί καταρᾶται ὁ κύων ὁ τεθνηκὼς οὗτος τὸν κύριόν μου τὸν βασιλέα; διαβήσομαι δὴ καὶ ἀφελῶ τὴν κεφαλὴν αὐτοῦ.
\vs{10}Καὶ εἶπεν ὁ βασιλεὺς, τί ἐμοὶ καὶ ὑμῖν, υἱοὶ Σαρουίας; καὶ ἄφετε αὐτὸν, καὶ οὕτως καταράσθω, ὅτι Κύριος εἶπεν αὐτῷ καταρᾶσθαι τὸν Δαυίδ· καὶ τίς ἐρεῖ, ὡς τί ἐποίησας οὕτως;
\vs{11}Καὶ εἶπε Δαυὶδ πρὸς Ἀβεσσὰ καὶ πρὸς πάντας τοὺς παῖδας αὐτοῦ, ἰδοὺ ὁ υἱός μου ὁ ἐξελθὼν ἐκ τῆς κοιλίας μου ζητεῖ τὴν ψυχήν μου, καὶ προσέτι νῦν ὁ υἱὸς τοῦ Ἰεμινί· ἄφετε αὐτὸν καταρᾶσθαι, ὅτι εἶπεν αὐτῷ Κύριος.
\vs{12}Εἴπως ἴδοι Κύριος ἐν τῇ ταπεινώσει μου, καὶ ἐπιστρέψει μοι ἀγαθὰ ἀντὶ τῆς κατάρας αὐτοῦ τῇ ἡμέρᾳ ταύτῃ.

\vs{13}Καὶ ἐπορεύθη Δαυὶδ καὶ πάντες οἱ ἄνδρες αὐτοῦ ἐν τῇ ὁδῷ· καὶ Σεμεῒ ἐπορεύετο ἐκ πλευρᾶς τοῦ ὄρους ἐχόμενα αὐτοῦ πορευόμενος καὶ καταρώμενος καὶ λιθάζων ἐν λίθοις ἐκ πλαγίων αὐτοῦ καὶ τῷ χοῒ πάσσων.
\vs{14}Καὶ ἦλθεν ὁ βασιλεὺς καὶ πᾶς ὁ λαὸς μετʼ αὐτοῦ ἐκλελυμένοι, καὶ ἀνέψυξαν ἐκεῖ.

\vs{15}Καὶ Ἀβεσσαλὼμ καὶ πᾶς ἀνὴρ Ἰσραὴλ εἰσῆλθον εἰς Ἰερουσαλὴμ, καὶ Ἀχιτόφελ μετʼ αὐτοῦ.
\vs{16}Καὶ ἐγενήθη ἡνίκα ἦλθε Χουσὶ ὁ ἀρχιεταῖρος Δαυὶδ πρὸς Ἀβεσσαλὼμ, καὶ εἶπε Χουσὶ πρὸς Ἀβεσσαλὼμ, ζήτω ὁ βασιλεύς.
\vs{17}Καὶ εἶπεν Ἀβεσσαλὼμ πρὸς Χουσὶ, τοῦτο τὸ ἔλεός σου μετὰ τοῦ ἑταίρου σου; ἱνατί οὐκ ἀπῆλθες μετὰ τοῦ ἑταίρου σου;
\vs{18}Καὶ εἶπε Χουσὶ πρὸς Ἀβεσσαλὼμ, οὐχὶ, ἀλλὰ κατόπισθεν οὗ ἐξελέξατο Κύριος καὶ ὁ λαὸς οὗτος καὶ πᾶς ἀνὴρ Ἰσραὴλ, αὐτῷ ἔσομαι, καὶ μετὰ αὐτοῦ καθήσομαι.
\vs{19}Καὶ τὸ δεύτερον, τίνι ἐγὼ δουλεύσω; οὐχὶ ἐνώπιον τοῦ υἱοῦ αὐτοῦ; καθάπερ ἐδούλευσα ἐνώπιον τοῦ πατρός σου, οὕτως ἔσομαι ἐνώπιόν σου.

\vs{20}Καὶ εἶπεν Ἀβεσσαλὼμ πρὸς Ἀχιτόφελ, φέρετε ἑαυτοῖς βουλὴν τί ποιήσωμεν.
\vs{21}Καὶ εἶπεν Ἀχιτόφελ πρὸς Ἀβεσσαλὼμ, εἴσελθε πρὸς τὰς παλλακὰς τοῦ πατρός σου, ἃς κατέλιπε φυλάσσειν τὸν οἶκον αὐτοῦ, καὶ ἀκούσεται πᾶς Ἰσραὴλ, ὅτι κατῄσχυνα, τὸν πατέρα σοῦ, καὶ ἐνισχύσουσιν αἱ χεῖρες πάντων τῶν μετὰ σοῦ.
\vs{22}Καὶ ἔπηξαν τὴν σκηνὴν τῷ Ἀβεσσαλὼμ ἐπὶ τὸ δῶμα, καὶ εἰσῆλθεν Ἀβεσσαλὼμ πρὸς τὰς παλλακὰς τοῦ πατρὸς αὐτοῦ κατʼ ὀφθαλμοὺς παντὸς Ἰσραήλ.
\vs{23}Καὶ ἡ βουλὴ Ἀχιτόφελ, ἣν ἐβουλεύσατο ἐν ταῖς ἡμέραις ταῖς πρώταις, ὃν τρόπον ἐπερωτήσῃ τις ἐν λόγῳ τοῦ Θεοῦ· οὕτως πᾶσα ἡ βουλὴ τοῦ Ἀχιτόφελ καί γε τῷ Δαυὶδ καί γε τῷ Ἀβεσσαλώμ.

\ch{17}
Καὶ εἶπεν Ἀχιτόφελ πρὸς Ἀβεσσαλὼμ, ἐπιλέξω δὴ ἐμαυτῷ δώδεκα χιλιάδας ἀνδρῶν, καὶ ἀναστήσομαι καὶ καταδιώξω ὀπίσω Δαυὶδ τὴν νύκτα.
\vs{2}Καὶ ἐπελεύσομαι ἐπʼ αὐτὸν, καὶ αὐτὸς κοπιῶν καὶ ἐκλελυμένος χερσὶ, καὶ ἐκστήσω αὐτὸν, καὶ φεύξεται πᾶς ὁ λαὸς ὁ μετʼ αὐτοῦ, καὶ πατάξω τὸν βασιλέα μονώτατον·
\vs{3}Καὶ ἐπιστρέψω πάντα τὸν λαὸν πρὸς σὲ, ὃν τρόπον ἐπιστρέφει ἡ νύμφη πρὸς τὸν ἄνδρα αὐτῆς· πλὴν ψυχὴν ἀνδρὸς ἑνὸς σὺ ζητεῖς, καὶ παντὶ τῷ λαῷ ἔσται εἰρήνη.
\vs{4}Καὶ εὐθὺς ὁ λόγος ἐν ὀφθαλμοῖς Ἀβεσσαλὼμ, καὶ ἐν ὀφθαλμοῖς πάντων τῶν πρεσβυτέρων Ἰσραήλ.

\vs{5}Καὶ εἶπεν Ἀβεσσαλὼμ, καλέσατε δὴ καί γε τὸν Χουσὶ τὸν Ἀραχὶ, καὶ ἀκούσωμεν τί ἐν τῷ στόματι αὐτοῦ, καί γε αὐτοῦ.
\vs{6}Καὶ εἰσῆλθε Χουσὶ πρὸς Ἀβεσσαλὼμ, καὶ εἶπεν Ἀβεσσαλὼμ πρὸς αὐτὸν, λέγων, κατὰ τὸ ῥῆμα τοῦτο ἐλάλησεν Ἀχιτόφελ· ποιήσομεν κατὰ τὸν λόγον αὐτοῦ; εἰ δὲ μὴ, σὺ λάλησον.

\vs{7}Καὶ εἶπε Χουσὶ πρὸς Ἀβεσσαλὼμ, οὐκ ἀγαθὴ αὕτη ἡ βουλὴ ἣν ἐβουλεύσατο Ἀχιτόφελ τὸ ἅπαξ τοῦτο.
\vs{8}Καὶ εἶπε Χουσὶ, σὺ οἶδας τὸν πατέρα σου καὶ τοὺς ἄνδρας αὐτοῦ, ὅτι δυνατοί εἰσι σφόδρα καὶ κατάπικροι τῇ ψυχῇ αὐτῶν, ὡς ἄρκος ἠτεκνωμένη ἐν ἀγρῷ, καὶ ὡς ὗς τραχεῖα ἐν τῇ πεδίῳ· ὁ πατήρ σου ἀνὴρ πολεμιστὴς, καὶ οὐ μὴ καταλύσῃ τὸν λαόν.
\vs{9}Ἰδοὺ γὰρ αὐτὸς νῦν κέκρυπται ἐν ἑνὶ τῶν βουνῶν ἢ ἐν ἑνὶ τῶν τόπων· καὶ ἔσται ἐν τῷ ἐπιπεσεῖν αὐτοῖς ἐν ἀρχῇ, καὶ ἀκούσῃ ἀκούων, καὶ εἴπῃ, ἐγενήθη θραῦσις ἐν τῷ λαῷ τῷ ὀπίσω Ἀβεσσαλώμ·
\vs{10}Καί γε αὐτὸς υἱὸς δυνάμεως, οὗ ἡ καρδία καθὼς ἡ καρδία τοῦ λέοντος, τηκομένη τακήσεται· ὅτι· οἶδε πᾶς Ἰσραὴλ, ὅτι δυνατὸς ὁ πατήρ σου, καὶ υἱοὶ δυνάμεως οἱ μετʼ αὐτοῦ.
\vs{11}Ὅτι οὕτως συμβουλεύων ἐγὼ συνεβούλευσα, καὶ συναγόμενος συναχθήσεται ἐπὶ σὲ πᾶς Ἰσραὴλ ἀπὸ Δὰν καὶ ἕως Βηρσαβεὲ, ὡς ἡ ἄμμος ἡ ἐπὶ τῆς θαλάσσης εἰς πλῆθος· καὶ τὸ πρόσωπόν σου πορευόμενον ἐν μέσῳ αὐτῶν.
\vs{12}Καὶ ἥξομεν πρὸς αὐτὸν εἰς ἕνα τῶν τόπων οὗ ἐὰν εὕρωμεν αὐτὸν ἐκεῖ, καὶ παρεμβαλοῦμεν ἐπʼ αὐτὸν, ὡς πίπτει δρόσος ἐπὶ τὴν γῆν, καὶ οὐχ ὑπολειψόμεθα ἐν αὐτῷ καὶ τοῖς ἀνδράσιν αὐτοῦ τοῖς μετʼ αὐτοῦ, καί γε ἕνα.
\vs{13}Καὶ ἐὰν εἰς τὴν πόλιν συναχθῇ, καὶ λήψεται πᾶς Ἰσραὴλ πρὸς τὴν πόλιν ἐκείνην σχοινία, καὶ συροῦμεν αὐτὴν ἕως εἰς τὸν χειμάῤῥουν, ὅπως μὴ καταλειφθῇ ἐκεῖ μηδὲ λίθος.

\vs{14}Καὶ εἶπεν Ἀβεσσαλὼμ, καὶ πᾶς ἀνὴρ Ἰσραὴλ, ἀγαθὴ ἡ βουλὴ Χουσὶ τοῦ Ἀραχὶ ὑπὲρ τὴν βουλὴν Ἀχιτόφελ· καὶ Κύριος ἐνετείλατο διασκεδάσαι τὴν βουλὴν τοῦ Ἀχιτόφελ τὴν ἀγαθὴν, ὅπως ἂν ἐπαγάγῃ Κύριος ἐπὶ Ἀβεσσαλὼμ τὰ κακὰ πάντα.

\vs{15}Καὶ εἶπε Χουσὶ ὁ τοῦ Ἀραχὶ πρὸς Σαδὼκ καὶ Ἀβιάθαρ τοὺς ἱερεῖς, οὕτως καὶ οὕτως συνεβούλευσεν Ἀχιτόφελ τῷ Ἀβεσσαλὼμ καὶ τοῖς πρεσβυτέροις Ἰσραήλ· καὶ οὕτως καὶ οὕτως συνεβούλευσα ἐγώ.
\vs{16}Καὶ νῦν ἀποστείλατε ταχὺ καὶ ἀναγγείλατε τῷ Δαυὶδ, λέγοντες, μὴ αὐλισθῇς τὴν νύκτα ἐν Ἀραβὼθ τῆς ἐρήμου, καί γε διαβαίνων σπεῦσον, μήποτε καταπείσῃ τὸν βασιλέα, καὶ πάντα τὸν λαὸν τὸν μετʼ αὐτοῦ.

\vs{17}Καὶ Ἰωνάθαν καὶ Ἀχιμάας εἱστήκεισαν ἐν τῇ πηγῇ Ῥωγὴλ, καὶ ἐπορεύθη ἡ παιδίσκη, καὶ ἀνήγγειλεν αὐτοῖς, καὶ αὐτοὶ πορεύονται καὶ ἀναγγέλλουσι τῷ βασιλεῖ Δαυίδ· ὅτι οὐκ ἠδύναντο ὀφθῆναι τοῦ εἰσελθεῖν τἰς τὴν πόλιν.
\vs{18}Καὶ εἶδεν αὐτοὺς παιδάριον, καὶ ἀνήγγειλε τῷ Ἀβεσσαλώμ· καὶ ἐπορεύθησαν οἱ δύο ταχέως, καὶ εἰσῆλθαν εἰς οἰκίαν ἀνδρὸς ἐν Βαουρίμ· καὶ αὐτῷ λάκκος ἐν τῇ αὐλῇ, καὶ κατέβησαν ἐκεῖ.
\vs{19}Καὶ ἔλαβεν ἡ γυνὴ, καὶ διεπέτασε τὸ ἐπικάλυμμα ἐπὶ πρόσωπον τοῦ λάκκου, καὶ ἔψυξεν ἐπʼ αὐτῷ ἀραφὼθ, καὶ οὐκ ἐγνώσθη ῥῆμα.
\vs{20}Καὶ ἦλθον οἱ παῖδες Ἀβεσσαλὼμ πρὸς τὴν γυναῖκα εἰς τὴν οἰκίαν, καὶ εἶπαν, ποῦ Ἀχιμάας καὶ Ἰωνάθαν; καὶ εἶπεν αὐτοῖς ἡ γυνὴ, παρῆλθαν μικρὸν το ὕδατος· καὶ ἐζήτησαν, καὶ οὐχ εὗραν, καὶ ἀνέστρεψαν εἰς Ἰερουσαλήμ.
\vs{21}Ἐγένετο δὲ μετὰ τὸ ἀπελθεῖν αὐτοὺς, καὶ ἀνέβησαν ἐκ τοῦ λάκκου, καὶ ἐπορεύθησαν· καὶ ἀπήγγειλαν τῷ βασιλεῖ Δαυὶδ, καὶ εἶπαν πρὸς Δαυὶδ, ἀνάστητε καὶ διάβητε ταχέως τὸ ὕδωρ, ὅτι οὕτως ἐβουλεύσατο περὶ ὑμῶν Ἀχιτόφελ.

\vs{22}Καὶ ἀνέστη Δαυὶδ καὶ πᾶς ὁ λαὸς ὁ μετʼ αὐτοῦ, καὶ διέβησαν τὸν Ἰορδάνην ἕως τοῦ φωτὸς τοῦ πρωῒ, ἕως ἑνὸς οὐκ ἔλαθεν ὃς οὐ διῆλθε τὸν Ἰορδάνην.

\vs{23}Καὶ Ἀχιτόφελ εἶδεν ὅτι οὐκ ἐγενήθη ἡ βουλὴ αὐτοῦ, καὶ ἐπέσαξε τὴν ὄνον αὐτοῦ, καὶ ἀνέστη καὶ ἀπῆλθεν εἰς τὸν οἶκον αὐτοῦ εἰς τὴν πόλιν αὐτοῦ· καὶ ἐνετείλατο τῷ οἴκῳ αὐτοῦ, καὶ ἀπήγξατο καὶ ἀπέθανεν καὶ ἐτάφη ἐν τῷ τάφῳ τοῦ πατρὸς αὐτοῦ.

\vs{24}Καὶ Δαυεὶδ διῆλθεν εἰς Μαναΐμ καὶ Ἀβεσσαλὼμ διέβη τὸν Ἰορδάνην αὐτὸς καὶ πᾶς ἀνὴρ Ἰσραὴλ μετʼ αὐτοῦ.
\vs{25}Καὶ τὸν Ἀμεσσαῒ κατέστησεν Ἀβεσσαλὼμ ἀντὶ Ἰωὰβ ἐπὶ τῆς δυνάμεως. Καὶ Ἀμεσσαῒ υἱὸς ἀνδρὸς, καὶ ὄνομα αὐτῷ Ἰεθὲρ ὁ Ἰεζραηλίτης· οὗτος εἰσῆλθε πρὸς Ἀβιγαίαν θυγατέρα Νάας ἀδελφὴν Σαρουίας μητρὸς Ἰωάβ.
\vs{26}Καὶ παρενέβαλε πᾶς Ἰσραὴλ καὶ Ἀβεσσαλὼμ εἰς τὴν γῆν Γαλαάδ.

\vs{27}Καὶ ἐγένετο ἡνίκα ἦλθε Δαυὶδ εἰς Μαναῒμ, καὶ Οὐεσβὶ υἱὸς Νάας ἐκ Ῥαββὰθ υἱῶν Ἀμμὼν, καὶ Μαχὶρ υἱὸς Ἀμιὴλ ἐκ Λωδαβὰρ, καὶ Βερζελλὶ ὁ Γαλααδίτης ἐκ Ῥωγελλὶμ,
\vs{28}ἤνεγκαν δέκα κοίτας ἀμφιτάπους, καὶ λέβητας δέκα, καὶ σκεύη κεράμου, καὶ πυροὺς, καὶ κριθὰς, καὶ ἄλευρον, καὶ ἄλφιτον, καὶ κύαμον, καὶ φακὸν,
\vs{29}καὶ μέλι, καὶ βούτυρον καὶ πρόβατα καὶ σαφὼθ βοῶν· καὶ προσήνεγκαν τῷ Δαυὶδ, καὶ τῷ λαῷ τῷ μετʼ αὐτοῦ φαγεῖν· ὅτι εἶπεν, ὁ λαὸς πεινῶν καὶ ἐκλελυμένος καὶ διψῶν ἐν τῇ ἐρήμῳ.

\ch{18}
Καὶ ἐπεσκέψατο Δαυὶδ τὸν λαὸν τὸν μετʼ αὐτοῦ, καὶ κατέστησεν ἐπʼ αὐτῶν χιλιάρχους καὶ ἑκατοντάρχους.
\vs{2}Καὶ ἀπέστειλε Δαυὶδ τὸν λαὸν τὸ τρίτον ἐν χειρὶ Ἰωὰβ, καὶ τὸ τρίτον ἐν χειρὶ Ἀβεσσὰ υἱοῦ Σαρουίας ἀδελφοῦ Ἰωὰβ, καὶ τὸ τρίτον ἐν χειρὶ Ἐθὶ τοῦ Γεθαίου· καὶ εἶπε Δαυὶδ πρὸς τὸν λαὸν, ἐξελθὼν ἐξελεύσομαι καί γε ἐγὼ μεθʼ ὑμῶν.
\vs{3}Καὶ εἶπαν, οὐκ ἐξελεύσῃ, ὅτι ἐὰν φυγῇ φύγωμεν, οὐ θήσουσιν ἐφʼ ἡμᾶς καρδίαν· καὶ ἐὰν ἀποθάνωμεν τὸ ἥμισυ ἡμῶν, οὐ θήσουσιν ἐφʼ ἡμᾶς καρδίαν, ὅτι σὺ ὡς ἡμεῖς δέκα χιλιάδες· καὶ νῦν ἀγαθὸν, ὅτι ἔσῃ ἡμῖν ἐν τῇ πόλει βοήθεια τοῦ βοηθεῖν.
\vs{4}Καὶ εἶπε πρὸς αὐτοὺς ὁ βασιλεὺς, ὃ ἐὰν ἀρέσῃ ἐν ὀφθαλμοῖς ὑμῶν, ποιήσω· καὶ ἔστη ὁ βασιλεὺς ἀνὰ χεῖρα τῆς πύλης· καὶ πᾶς ὁ λαὸς ἐξεπορεύετο εἰς ἑκατοντάδας καὶ εἰς χιλιάδας.

\vs{5}Καὶ ἐνετείλατο ὁ βασιλεὺς τῷ Ἰωὰβ καὶ τῷ Ἀβεσσὰ καὶ τῷ Ἐθὶ, λέγων, φείσασθέ μου τοῦ παιδαρίου τοῦ Ἀβεσσαλώμ· καὶ πᾶς ὁ λαὸς ἤκουσεν ἐντελλομένου τοῦ βασιλέως πᾶσι τοῖς ἄρχουσιν ὑπὲρ Ἀβεσσαλώμ.

\vs{6}Καὶ ἐξῆλθε πᾶς ὁ λαὸς εἰς τὸν δρυμὸν ἐξεναντίας Ἰσραήλ· καὶ ἐγένετο ὁ πόλεμος ἐν τῷ δρυμῷ Ἐφραίμ.
\vs{7}Καὶ ἔπταισεν ἐκεῖ ὁ λαὸς Ἰσραὴλ ἐνώπιον τῶν παίδων Δαυὶδ, καὶ ἐγένετο ἡ θραῦσις μεγάλη ἐν τῇ ἡμέρᾳ ἐκείνῃ, εἴκοσι χιλιάδες ἀνδρῶν.
\vs{8}Καὶ ἐγένετο ἐκεῖ ὁ πόλεμος διεσπαρμένος ἐπὶ πρόσωπον πάσης τῆς γῆς· καὶ ἐπλεόνασεν ὁ δρυμὸς τοῦ καταφαγεῖν ἐκ τοῦ λαοῦ, ὑπὲρ οὓς κατέφαγεν ἐν τῷ λαῷ ἡ μάχαιρα τῇ ἡμέρᾳ ἐκαίνῃ.
\vs{9}Καὶ συνήντησεν Ἀβεσσαλὼμ ἐνώπιον τῶν παίδων Δαυίδ· καὶ Ἀβεσσαλὼμ ἦν ἐπιβεβηκὼς ἐπὶ τοῦ ἡμιόνου αὐτοῦ, καὶ εἰσῆλθεν ὁ ἡμίονος ὑπὸ τὸ δάσος τῆς δρυὸς τῆς μεγάλης, καὶ περιεπλάκη ἡ κεφαλὴ αὐτοῦ ἐν τῇ δρυῒ, καὶ ἐκρεμάσθη ἀναμέσον τοῦ οὐρανοῦ καὶ ἀναμέσον τῆς γῆς, καὶ ὁ ἡμίονος ὑποκάτω αὐτοῦ παρῆλθε.

\vs{10}Καὶ εἶδεν ἀνὴρ εἷς, καὶ ἀνήγγειλε τῷ Ἰωὰβ, καὶ εἶπεν, ἰδοὺ ἑώρακα τὸν Ἀβεσσαλὼμ κρεμάμενον ἐν τῇ δρυΐ.
\vs{11}Καὶ εἶπεν Ἰωὰβ τῷ ἀνδρὶ τῇ ἀναγγέλλοντι αὐτῷ, καὶ ἰδοὺ ἑώρακας· τί ὅτι οὐκ ἐπάταξας αὐτὸν ἐκεῖ εἰς τὴν γῆν; καὶ ἐγὼ ἂν δεδώκειν σοι δέκα ἀργυρίου καὶ παραζώνην μίαν.
\vs{12}Εἶπε δὲ ὁ ἀνὴρ πρὸς Ἰωὰβ, καὶ ἐγώ εἰμι ἵστημι ἐπὶ τὰς χεῖράς μου χιλίους σίκλους ἀργυρίου, οὐ μὴ ἐπιβάλω τὴν χεῖρά μου ἐπὶ τὸν υἱὸν τοῦ βασιλέως· ὅτι ἐν τοῖς ὠσὶν ἡμῶν ἐνετείλατο ὁ βασιλεύς σοι καὶ τῷ Ἀβεσσὰ καὶ τῷ Ἐθὶ, λέγων, φυλάξατέ μοι τὸ παιδάριον τὸν Ἀβεσσαλὼμ,
\vs{13}μὴ ποιῆσαι ἐν τῇ ψυχῇ αὐτοῦ ἄδικον· καὶ πᾶς ὁ λόγος οὐ λήσεται ἀπὸ τοῦ βασιλέως, καὶ σὺ στήσῃ ἐξεναντίας.
\vs{14}Καὶ εἶπεν Ἰωὰβ, τοῦτο ἐγὼ ἄρξομαι, οὐχ οὕτως μενῶ ἐνώπιόν σου· καὶ ἔλαβεν Ἰωὰβ τρία βέλη ἐν τῇ χειρὶ αὐτοῦ, καὶ ἐνέπηξεν αὐτὰ ἐν τῇ καρδίᾳ Ἀβεσσαλὼμ, ἔτι αὐτοῦ ζῶντος ἐν τῇ καρδίᾳ τῆς δρυός.
\vs{15}Καὶ ἐκύκλωσαν δέκα παιδάρια αἴροντα τὰ σκεύη Ἰωὰβ, καὶ ἐπάταξαν τὸν Ἀβεσσαλὼμ, καὶ ἐθανάτωσαν αὐτόν.

\vs{16}Καὶ ἐσάλπισεν Ἰωὰβ ἐν κερατίνῃ, καὶ ἀπέστρεψεν ὁ λαὸς τοῦ μὴ διώκειν ὀπίσω Ἰσραὴλ, ὅτι ἐφείδετο Ἰωὰβ τοῦ λαοῦ.
\vs{17}Καὶ ἔλαβε τὸν Ἀβεσσαλὼμ, καὶ ἔῤῥιψεν αὐτὸν εἰς χάσμα μέγα ἐν τῷ δρυμῷ εἰς τὸν βόθυνον τὸν μέγαν, καὶ ἐστήλωσεν ἐπʼ αὐτὸν σωρὸν λίθων μέγαν σφόδρα· καὶ πᾶς Ἰσραὴλ ἔφυγεν ἀνὴρ εἰς τὸ σκήνωμα αὐτοῦ.
\vs{18}Καὶ Ἀβεσσαλὼμ ἔτι ζῶν ἔλαβε καὶ ἔστησεν ἑαυτῷ τὴν στήλην ἐν ᾗ ἐλήφθη, καὶ ἐστήλωσεν αὐτὴν λαβεῖν τὴν στήλην τὴν ἐν τῇ βασιλέως, ὅτι εἶπεν, ὅτι οὐκ ἔστιν αὐτῷ υἱὸς ἕνεκν τοῦ ἀναμνῆσαι τὸ ὄνομα αὐτοῦ· καὶ ἐκάλεσε τὴν στήλην, Χεὶρ Ἀβεσσαλὼμ, ἕως τῆς ἡμέρας ταύτης.

\vs{19}Καὶ Ἀχιμάας υἱὸς Σαδὼκ εἶπε, δράμω δὴ καὶ εὐαγγελιῶ τῷ βασιλεῖ, ὅτι ἔκρινε Κύριος ἐκ χειρὸς τῶν ἐχθρῶν αὐτοῦ.
\vs{20}Καὶ εἶπεν αὐτῷ Ἰωὰβ, οὐκ ἀνὴρ εὐαγγελίας σὺ ἐν τῇ ἡμέρᾳ ταύτῃ, καὶ εὐαγγελιῇ ἐν ἡμέρᾳ ἄλλῃ· ἐν δὲ τῇ ἡμέρᾳ ταύτῃ οὐκ εὐαγγελιῇ, οὗ εἵνεκεν ὁ υἱὸς τοῦ βασιλέως ἀπέθανε.
\vs{21}Καὶ εἶπεν Ἰωὰβ τῷ Χουσὶ, Βαδίσας ἀνάγγειλον τῷ βασιλεῖ ὅσα εἶδες· καὶ προσεκύνησε Χουσεὶ τῷ Ἰωὰβ, καὶ ἐξῆλθε.
\vs{22}Καὶ προσέθετο ἕτι Ἀχιμάας υἱὸς Σαδὼκ, καὶ εἶπε πρὸς Ἰωὰβ, καὶ ἔστω, ὅτι δράμω καί γε ἐγὼ ὀπίσω τοῦ Χουσί· καὶ εἶπεν Ἰωὰβ, ἱνατί σὺ τοῦτο τρέχεις, υἱέ μου; δεῦρο, οὐκ ἔστι σοι εὐαγγέλια εἰς ὠφέλειαν πορευομένῳ.
\vs{23}Καὶ εἶπε, τί γὰρ ἐὰν δράμω; καὶ εἶπεν αὐτῷ Ἰωὰβ, δράμε· καὶ ἔδραμεν Ἀχιμάας τὴν ὁδὸν τὴν τοῦ Κεχὰρ, καὶ ὑπερέβη τὸν Χουσί.

\vs{24}Καὶ Δαυὶδ ἐκάθητο ἀναμέσον τῶν δύο πυλῶν· καὶ ἐπορεύθη ὁ σκοπὸς εἰς τὸ δῶμα τῆς πύλης πρὸς τὸ τεῖχος, καὶ ἐπῇρε τοὺς ὀφθαλμοὺς αὐτοῦ, καὶ εἶδε. καὶ ἰδοὺ ἀνὴρ τρέχων μόνος ἐνώπιον αὐτοῦ.
\vs{25}Καὶ ἀνεβόησεν ὁ σκοπὸς καὶ ἀπήγγειλε τῷ βασιλεῖ· καὶ εἶπεν ὁ βασιλεὺς, εἰ μόνος ἐστὶν, εὐαγγέλια ἐν τῷ στόματι αὐτοῦ· καὶ ἐπορεύετο πορευόμενος καὶ ἐγγίζων.
\vs{26}Καὶ εἶδεν ὁ σκοπὸς ἄνδρα ἕτερον τρέχοντα· καὶ ἐβόησεν ὁ σκοπὸς πρὸς τῇ πύλῃ, καὶ εἶπε, καὶ ἰδοὺ ἀνὴρ ἕτερος τρέχων μόνος· καὶ εἶπεν ὁ βασιλεὺς, καί γε οὗτος εὐαγγελιζόμενος.
\vs{27}Καὶ εἶπεν ὁ σκοπός, ἐγὼ ὁρῶ τὸν δρόμον τοῦ πρώτου ὡς δρόμον Ἀχιμάας υἱοῦ Σαδώκ· καὶ εἶπεν ὁ βασιλεὺς, ἀνὴρ ἀγαθὸς οὗτος, καί γε εἰς εὐαγγελίαν ἀγαθὴν ἐλεύσεται.

\vs{28}Καὶ ἐβόησεν Ἀχιμάας, καὶ εἶπε πρὸς τὸν βασιλέα, εἰρήνη· καὶ προσεκύνησε τῷ βασιλεῖ ἐπὶ πρόσωπον αὐτοῦ ἐπὶ τὴν γῆν, καὶ εἶπεν, εὐλογητὸς Κύριος ὁ Θεός σου, ὃς ἀπέκλεισε τοὺς ἄνδρας τοὺς ἐπαραμένους τὴν χεῖρα αὐτῶν ἐν τῷ κυρίῳ μου τῷ βασιλεῖ.
\vs{29}Καὶ εἶπεν ὁ βασιλεὺς, εἰρήνη τῷ παιδαρίῳ τῷ Ἀβεσσαλώμ; καὶ εἶπεν Ἀχιμάας, εἶδον τὸ πλῆθος τὸ μέγα τοῦ ἀποστεῖλαι τὸν δοῦλον τοῦ βασιλέως Ἰωὰβ καὶ τὸν δοῦλόν σου, καὶ οὐκ ἔγνων τί ἐκεῖ.
\vs{30}Καὶ εἶπεν ὁ βασιλεὺς, ἐπίστρεψον, στηλώθητι ὧδε· καὶ ἐπεστράφη, καὶ ἔστη.

\vs{31}Καὶ ἰδοὺ ὁ Χουσὶ παρεγένετο, καὶ εἶπε τῷ βασιλεῖ, εὐαγγελισθήτω ὁ κύριός μου ὁ βασιλεὺς, ὅτι ἔκρινέ σοι Κύριος σήμερον ἐκ χειρὸς πάντων τῶν ἐπεγειρομένων ἐπὶ σέ.
\vs{32}Καὶ εἶπεν ὁ βασιλεὺς πρὸς τὸν Χουσὶ, εἰ εἰρήνη τῷ παιδαρίῳ τῷ Ἀβεσσαλώμ; καὶ εἶπεν ὁ Χουσὶ, γένοιντο ὡς τὸ παιδάριον οἱ ἐχθροὶ τοῦ κυρίου μου τοῦ βασιλέως, καὶ πάντες ὅσοι ἐπανέστησαν ἐπʼ αὐτὸν εἰς κακά.

\ch{19}Καὶ ἐταράχθη ὁ βασιλεὺς, καὶ ἀνέβη εἰς τὸ ὑπερῷον τῆς πύλης, καὶ ἔκλαυσε· καὶ οὕτως εἶπεν ἐν τῷ πορεύεσθαι αὐτὸν, υἱέ μου Ἀβεσσαλὼμ, υἱέ μου, υἱέ μου Ἀβεσαλώμ· τίς δῴη τὸν θάνατόν μου ἀντὶ σοῦ; ἐγὼ ἀντὶ σοῦ Ἀβεσσαλὼμ, υἱέ μου, υἱέ μου.

\vs{2}Καὶ ἀνηγγέλη τῷ Ἰωὰβ, λέγοντες, ἰδοῦ ὁ βασιλεὺς κλαίει καὶ πενθεῖ ἐπὶ Ἀβεσσαλώμ.
\vs{3}Καὶ ἐγένετο ἡ σωτηρία ἐν τῇ ἡμέρᾳ ἐκείνῃ εἰς πένθος παντὶ τῷ λαῷ, ὅτι ἤκουσεν ὁ λαὸς ἐν τῇ ἡμέρᾳ ἐκείνῃ, λέγων, ὅτι λυπεῖται ὁ βασιλεὺς ἐπὶ τῷ υἱῷ αὐτοῦ.
\vs{4}Καὶ διεκλέπτετο ὁ λαὸς ἐν τῇ ἡμέρᾳ ἐκείνῃ τοῦ εἰσελθεῖν εἰς τὴν πόλιν, καθὼς διακλέπτεται ὁ λαὸς οἱ αἰσχυνόμενοι ἐν τῷ αὐτοὺς φεύγειν ἐν τῷ πολέμῳ.
\vs{5}Καὶ ὁ βασιλεὺς ἔκρυψε τὸ πρόσωπον αὐτοῦ· καὶ ἔκραξεν ὁ βασιλεὺς φωνῇ μεγάλῃ, λέγων, υἱέ μου Ἀβεσσαλὼμ, Ἀβεσσαλὼμ υἱὲ μου.

\vs{6}Καὶ εἰσῆλθεν Ἰωὰβ πρὸς τὸν βασιλέα εἰς τὸν οἶκον, καὶ εἶπε, κατῄσχυνας σήμερον τὰ πρόσωπα πάντων τῶν δούλων σου τῶν ἐξαιρουμένων σε σήμερον, καὶ τὴν ψυχὴν τῶν υἱῶν σου, καὶ τῶν θυγατέρων σου, καὶ τὴν ψυχὴν τῶν γυναικῶν σου, καὶ τῶν παλλακῶν σου,
\vs{7}τοῦ ἀγαπᾷν τοὺς μισοῦντάς σε, καὶ μισεῖν τοὺς ἀγαπῶντάς σε· καὶ ἀνήγγειλας σήμερον, ὅτι οὐκ εἰσὶν οἱ ἄρχοντές σου, οὐδὲ παῖδες· ὅτι ἔγνωκα σήμερον, ὅτι εἰ Ἀβεσσαλὼμ ἔζη, πάντες ἡμεῖς σήμερον νεκροί, ὅτι τότε τὸ εὐθὲς ἦν ἐν ὀφθαλμοῖς σου.
\vs{8}Καὶ νῦν ἀναστὰς ἔξελθε καὶ λάλησον εἰς τὴν καρδίαν τῶν δούλων σου, ὅτι ἐν Κυρίῳ ὤμοσα, ὅτι εἰ μὴ ἐκπορεύσῃ σήμερον, εἰ αὐλισθήσεται ἀνὴρ μετὰ σοῦ τὴν νύκτα ταύτην· καὶ ἐπίγνωθι σεαυτῷ, καὶ κακόν σοι τοῦτο ὑπὲρ πᾶν τὸ κακὸν τὸ ἐπελθόν σοι ἐκ νεότητός σου ἕως τοῦ νῦν.
\vs{9}Καὶ ἀνέστη ὁ βασιλεὺς καὶ ἐκάθισεν ἐν τῇ πύλῃ· καὶ πᾶς ὁ λαὸς ἀνήγγειλαν, λέγοντες, ἰδοὺ ὁ βασιλεὺς κάθηται ἐν τῇ πύλῃ· καὶ εἰσῆλθε πᾶς ὁ λαὸς κατὰ πρόσωπον τοῦ βασιλέως ἐπὶ τὴν πύλην· καὶ Ἰσραὴλ ἔφυγεν ἀνὴρ εἰς τὰ σκηνώματα αὐτοῦ.

\vs{10}Καὶ ἦν πᾶς ὁ λαὸς κρινόμενος ἐν πάσαις φυλαῖς Ἰσραὴλ, λέγοντες, ὁ βασιλεὺς Δαυὶδ ἐῤῥύσατο ἡμᾶς ἀπὸ πάντων τῶν ἐχθρῶν ἡμῶν, καὶ αὐτὸς ἐξείλετο ἡμᾶς ἐκ χειρὸς ἀλλοφύλων· καὶ νῦν πέφευγεν ἀπὸ τῆς γῆς, καὶ ἀπὸ τῆς βασιλείας αὐτοῦ, καὶ ἀπὸ Ἀβεσσαλώμ.
\vs{11}Καὶ Ἀβεσσαλὼμ, ὃν ἐχρίσαμεν ἐφʼ ἡμῶν, ἀπέθανεν ἐν τῷ πολέμῳ· καὶ νῦν ἱνατί ὑμεῖς κωφεύετε τοῦ ἐπιστρέψαι τὸν βασιλέα; καὶ τὸ ῥῆμα παντὸς Ἰσραὴλ ἦλθε πρὸς τὸν βασιλέα.

\vs{12}Καὶ ὁ βασιλεὺς Δαυὶδ ἀπέστειλε πρὸς Σαδὼκ καὶ πρὸς Ἀβιάθαρ τοῦς ἱερεῖς, λέγων, λαλήσατε πρὸς τοὺς πρεσβυτέρους Ἰούδα, λέγοντες, ἱνατί γίνεσθε ἔσχατοι τοῦ ἐπιστρέψαι τὸν βασιλέα εἰς τὸν οἶκον αὐτοῦ; καὶ λόγος παντὸς Ἰσραὴλ ἦλθε πρὸς τὸν βασιλέα εἰς τὸν οἶκον αὐτοῦ.
\vs{13}Ἀδελφοί μου ὑμεῖς, ὀστᾶ μου καὶ σάρκες μου ὑμεῖς· ἱνατί γίνεσθε ἔσχατοι τοῦ ἐπιστρέψαι τὸν βασιλέα εἰς τὸν οἶκον αὐτοῦ;
\vs{14}καὶ τῷ Ἀμεσσαῒ; ἐρεῖτε, οὐχὶ ὀστοῦν μου καὶ σάρξ μου σύ; καὶ νῦν τάδε ποιήσαι μοι ὁ Θεὸς, καὶ τάδε προσθείη, εἰ μὴ ἄρχων δυνάμεως ἔσῃ ἐνώπιον ἐμοῦ πάσας τὰς ἡμέρας ἀντὶ Ἰωάβ.
\vs{15}Καὶ ἔκλινε τὴν καρδίαν παντὸς ἀνδρὸς Ἰούδα ὡς ἀνδρὸς ἑνός· καὶ ἀπέστειλαν πρὸς τὸν βασιλεα, λέγοντες, ἐπιστράφηθι σὺ καὶ πάντες οἱ δοῦλοί σου.
\vs{16}Καὶ ἐπέστρεψεν ὁ βασιλεὺς, καὶ ἦλθεν ἕως τοῦ Ἰορδάνου· καὶ ἄνδρες Ἰούδα ἦλθαν εἰς Γάλγαλα τοῦ πορεύεσθαι εἰς ἀπαντὴν τοῦ βασιλέως, διαβιβάσαι τὸν βασιλέα τὸν Ἰορδάνην.

\vs{17}Καὶ ἐτάχυνε Σεμεῒ υἱὸς Γηρὰ υἱοῦ τοῦ Ἰεμινὶ ἐκ Βαουρὶμ, καὶ κατέβη μετὰ ἀνδρὸς Ἰούδα εἰς ἀπαντὴν τοῦ βασιλέως Δαυίδ,
\vs{18}καὶ χίλιοι ἄνδρες μετʼ αὐτοῦ ἐκ τοῦ Βενιαμὶν, καὶ Σιβὰ τὸ παιδάριον τοῦ οἴκου Σαοὺλ, καὶ πεντεκαίδεκα υἱοὶ αὐτοῦ μετʼ αὐτοῦ, καὶ εἴκοσι δοῦλοι αὐτοῦ μετʼ αὐτοῦ· καὶ κατεύθυναν τὸν Ἰορδάνην ἔμπροσθεν τοῦ βασιλέως,
\vs{19}καὶ ἐλειτούργησαν τὴν λειτουργίαν τοῦ διαβιβάσαι τὸν βασιλέα· καὶ διέβη ἡ διάβασις τοῦ ἐξεγεῖραι τὸν οἶκον τοῦ βασιλέως, καὶ τοῦ ποιῆσαι τὸ εὐθὲς ἐν ὀφθαλμοῖς αὐτοῦ. Καὶ Σεμεῒ υἱὸς Γηρὰ ἔπεσεν ἐπὶ πρόσωπον αὐτοῦ ἐνώπιον τοῦ βασιλέως, διαβαίνοντος αὐτοῦ τὸν Ἰορδάνην,
\vs{20}καὶ εἶπε πρὸς τὸν βασιλέα, μὴ δὴ λογισάσθω ὁ κύριός μου ἀνομίαν, καὶ μὴ μνησθῇς ὅσα ἠδίκησεν ὁ παῖς σου ἐν τῇ ἡμέρᾳ ᾗ ὁ κύριός μου ἐξεπορεύετο ἐξ Ἰερουσαλὴμ, τοῦ θέσθαι τὸν βασιλέα εἰς τὴν καρδίαν αὐτοῦ.
\vs{21}Ὅτι ἔγνω ὁ δοῦλός σου ὅτι ἐγὼ ἥμαρτον, καὶ ἰδοὺ ἐγὼ ἦλθον σήμερον πρότερος παντὸς Ἰσραὴλ καὶ οἴκου Ἰωσὴφ, τοῦ καταβῆναί με εἰς ἀπαντὴν τοῦ κυρίου μου τοῦ βασιλέως.

\vs{22}Καὶ ἀπεκρίθη Ἀβεσσαὲ υἱὸς Σαρουίας, καὶ εἶπε, μὴ ἀντὶ τούτου οὐ θανατωθήσεται Σεμεῒ, ὅτι κατηράσατο τὸν χριστὸν Κυρίου;
\vs{23}Καὶ εἶπε Δαυίδ, τί ἐμοὶ καὶ ὑμῖν, υἱοὶ Σαρουίας, ὅτι γίνεσθέ μοι σήμερον εἰς ἐπίβουλον; σήμερον οὐ θανατωθήσεταί τις ἀνὴρ ἐξ Ἰσραήλ· ὅτι οὐκ οἴδα εἰ σήμερον βασιλεύω ἐγὼ ἐπὶ τὸν Ἰσραήλ.
\vs{24}Καὶ εἶπεν ὁ βασιλεὺς πρὸς Σεμεῒ, οὐ μὴ ἀποθάνῃς· καὶ ὤμοσεν αὐτῷ ὁ βασιλεύς.

\vs{25}Καὶ Μεμφιβοσθὲ υἱὸς υἱοῦ Σαοὺλ κατέβη εἰς ἀπαντὴν τοῦ βασιλέως, καὶ οὐκ ἐθεράπευσε τοὺς πόδας αὐτοῦ, οὐδὲ ὠνυχίσατο, οὐδὲ ἐποίησε τὸν μύστακα αὐτοῦ, καὶ τὰ ἱμάτια αὐτοῦ οὐκ ἀπέπλυνεν, ἀπὸ τῆς ἡμέρας ἧς ἀπῆλθεν ὁ βασιλεὺς, ἕως τῆς ἡμέρας ἧς αὐτὸς παρεγένετο ἐν εἰρήνῃ.

\vs{26}Καὶ ἐγένετο ὅτε εἰσῆλθεν εἰς Ἱερουσαλὴμ εἰς ἀπάντησιν τοῦ βασιλέως, καὶ εἶπεν αὐτῷ ὁ βασιλεὺς, τί ὅτι οὐκ ἐπορεύθης μετʼ ἐμοῦ, Μεμφιβοσθέ;
\vs{27}Καὶ εἶπε πρὸς αὐτὸν Μεμφιβοσθὲ, κύριέ μου βασιλεῦ, ὁ δοῦλός μου παρελογίσατό με, ὅτι εἶπεν ὁ παῖς σου αὐτῷ, ἐπίσαξόν μοι τὴν ὄνον, καὶ ἐπιβῶ ἐπʼ αὐτὴν, καὶ πορεύσομαι μετὰ τοῦ βασιλέως, ὅτι χωλὸς ὁ δοῦλός σου.
\vs{28}Καὶ μεθώδευσεν ἐν τῷ δούλῳ σου πρὸς τὸν κύριόν μου τὸν βασιλέα· καὶ ὁ κύριός μου ὁ βασιλεὺς ὡς ἄγγελος τοῦ Θεοῦ, καὶ ποίησον τὸ ἀγαθὸν ἐν ὀφθαλμοῖς σου.
\vs{29}Ὅτι οὐκ ἦν πᾶς ὁ οἶκος τοῦ πατρός μου, ἀλλʼ ἢ ὅτι ἄνδρες θανάτου τῷ κυρίῳ μου τῷ βασιλεῖ, καὶ ἔθηκας τὸν δοῦλόν σου ἐν τοῖς ἐσθίουσι τὴν τράπεζάν σου· καὶ τί ἐστι μοὶ ἔτι δικαίωμα, καὶ τοῦ κεκραγέναι με ἔτι πρὸς τὸν βασιλέα;

\vs{30}Καὶ εἶπεν αὐτῷ ὁ βασιλεὺς, ἱνατί λαλεῖς ἔτι τοὺς λόγους σου; εἶπον, σὺ καὶ Σιβὰ διελεῖσθε τὸν ἀγρόν.
\vs{31}Καὶ εἶπε Μεμφιβοσθὲ πρὸς τὸν βασιλέα, καί γε τὰ πάντα λαβέτω, μετὰ τὸ παραγενέσθαι τὸν κύριόν μου τὸν βασιλέα ἐν εἰρήνῃ εἰς τὸν οἶκον αὐτοῦ.

\vs{32}Καὶ Βερζελλὶ ὁ Γαλααδίτης κατέβη ἐκ Ῥωγελλὶμ, καὶ διέβη μετὰ τοῦ βασιλέως τὸν Ἰορδάνην ἐκπέμψαι αὐτὸν τὸν Ἰορδάνην.
\vs{33}Καὶ Βερζελλὶ ἀνὴρ πρεσβύτερος σφόδρα, υἱὸς ὀγδοήκοντα ἐτῶν, καὶ αὐτὸς διέθρεψε τὸν βασιλέα ἐν τῷ οἰκεῖν αὐτὸν ἐν Μαναῒμ, ὅτι ἀνὴρ μέγας ἦν σφόδρα.
\vs{34}Καὶ εἶπεν ὁ βασιλεὺς πρὸς Βερζελλὶ, σὺ διαβήσῃ μετʼ ἐμοῦ, καὶ διαθρέψω τὸ γῆράς σου μετʼ ἐμοῦ ἐν Ἱερουσαλήμ.
\vs{35}Καὶ εἶπε Βερζελλὶ πρὸς τὸν βασιλέα, πόσαι ἡμέραι ἐτῶν ζωῆς μου, ὅτι ἀναβήσομαι μετὰ τοῦ βασιλέως εἰς Ἱερουσαλήμ;
\vs{36}Υἱὸς ὀγδοήκοντα ἐτῶν ἐγώ εἰμι σήμερον· εἰ μὴν γνώσομαι ἀναμέσον ἀγαθοῦ καὶ κακοῦ; εἰ γεύσεται ὁ δοῦλός σου ἔτι ὃ φάγομαι ἢ πίομαι; ἢ ἀκούσομαι ἔτι φωνὴν ᾀδόντων καὶ ᾀδουσῶν; καὶ ἱνατί ἔσται ἔτι ὁ δοῦλός σου εἰς φορτίον ἐπὶ τὸν κύριόν μου τὸν βασιλέα;
\vs{37}Ὡς βραχὺ διαβήσεται ὁ δοῦλός σου τὸν Ἰορδάνην μετὰ τοῦ βασιλέως· καὶ ἱνατί ἀνταποδίδωσί μοι ὁ βασιλεὺς τὴν ἀνταπόδοσιν ταύτην;
\vs{38}Καθισάτω δὴ ὁ δοῦλός σου, καὶ ἀποθανοῦμαι ἐν τῇ πόλει μου παρὰ τῷ τάφῳ τοῦ πατρός μου καὶ τῆς μητρός μου· καὶ ἰδοὺ ὁ δοῦλός σου Χαμαὰμ διαβήσεται μετὰ τοῦ κυρίου μου τοῦ βασιλέως· καὶ ποίησον αὐτῷ τὸ ἀγαθὸν ἐν ὀφθαλμοῖς σου.
\vs{39}Καὶ εἶπεν ὁ βασιλεὺς, μετʼ ἐμοῦ διαβήτω Χαμαὰμ, κᾀγὼ ποιήσω αὐτῷ τὸ ἀγαθὸν ἐν ὀφθαλμοῖς μου, καὶ πάντα ὅσα ἂν ἐκλέξῃ ἐπʼ ἐμοὶ, ποιήσω σοι.

\vs{40}Καὶ διέβη πᾶς ὁ λαὸς τὸν Ἰορδάνην, καὶ ὁ βασιλεὺς διέβη, καὶ κατεφίλησεν ὁ βασιλεὺς τὸν Βερζελλὶ, καὶ εὐλόγησεν αὐτὸν, καὶ ἐπέστρεψεν εἰς τὸν τόπον αὐτοῦ.
\vs{41}Καὶ διέβη ὁ βασιλεὺς εἰς Γάλγαλα, καὶ Χαμαὰμ διέβη μετʼ αὐτοῦ· καὶ πᾶς ὁ λαὸς Ἰούδα διαβαίνοντες μετὰ τοῦ βασιλέως, καί γε τὸ ἥμισυ τοῦ λαοῦ Ἰσραήλ.

\vs{42}Καὶ ἰδοὺ πᾶς ἀνὴρ Ἰσραὴλ παρεγένοντο πρὸς τὸν βασιλέα, καὶ εἶπε πρὸς τὸν βασιλέα, τί ὅτι ἔκλεψάν σε οἱ ἀδελφοὶ ἡμῶν ἀνὴρ Ἰούδα, καὶ διεβίβασαν τὸν βασιλέα καὶ τὸν οἶκον αὐτοῦ τὸν Ἰορδάνην, καὶ πάντες ἄνδρες Δαυὶδ μετʼ αὐτοῦ;
\vs{43}Καὶ ἀπεκρίθη πᾶς ἀνὴρ Ἰούδα πρὸς ἄνδρα Ἰσραὴλ, καὶ εἶπαν, διότι ἐγγίζει πρὸς μὲ ὁ βασιλεύς· καὶ ἱνατί οὕτως ἐθυμώθης περὶ τοῦ λόγου τούτου; μὴ βρώσει ἐφάγαμεν ἐκ τοῦ βασιλέως, ἢ δόμα ἔδωκεν, ἢ ἄρσιν ᾖρεν ἡμῖν;
\vs{44}Καὶ ἀπεκρίθη ἀνὴρ Ἰσραὴλ τῷ ἀνδρὶ Ἰούδα, καὶ εἶπε, δέκα χεῖρές μοι ἐν τῷ βασιλεῖ, καὶ πρωτότοκος ἐγὼ ἢ σὺ, καί γε ἐν τῷ Δαυίδ εἰμι ὑπὲρ σέ· καὶ ἱνατί τοῦτο ὕβρισάς με, καὶ οὐκ ἐλογίσθη ὁ λόγος μου πρῶτός μοι τοῦ Ἰούδα ἐπιστρέψαι τὸν βασιλέα ἐμοί; καὶ ἐσκληρύνθη ὁ λόγος ἀνδρὸς Ἰούδα ὑπὲρ τὸν λόγον ἀνδρὸς Ἰσραήλ.

\ch{20}
Καὶ ἐκεῖ ἐπικαλούμενος υἱὸς παράνομος, καὶ ὄνομα αὐτῷ Σαβεὲ, υἱὸς Βοχορί ἀνὴρ ὁ Ἰεμινὶ, καὶ ἐσάλπισε τῇ κερατίνῃ, καὶ εἶπεν, οὐκ ἔστιν ἡμῖν μερὶς ἐν Δαυὶδ, οὐδὲ κληρονομία ἡμῖν ἐν τῷ υἱῷ Ἰεσσαί· ἀνὴρ εἰς τὰ σκηνώματά σου, Ἰσραήλ.
\vs{2}Καὶ ἀνέβη πᾶς ἀνὴρ Ἰσραὴλ ἀπὸ ὄπισθεν Δαυὶδ ὀπίσω Σαβεὲ υἱοῦ Βοχορί· καὶ ἀνὴρ Ἰούδα ἐκολλήθη τῷ βασιλεῖ αὐτῶν, ἀπὸ τοῦ Ἰορδάνου καὶ ἕως Ἰερουσαλήμ.

\vs{3}Καὶ εἰσῆλθε Δαυὶδ εἰς τὸν οἶκον αὐτοῦ εἰς Ἱερουσαλήμ· καὶ ἔλαβεν ὁ βασιλεὺς τὰς δέκα γυναῖκας τὰς παλλακὰς αὐτοῦ, ἃς ἀφῆκε φυλάσσειν τὸν οἶκον, καὶ ἔδωκεν αὐτὰς ἐν οἴκῳ φυλακῆς, καὶ διέθρεψεν αὐτὰς, καὶ πρὸς αὐτὰς οὐκ εἰσῆλθε· καὶ ἦσαν συνεχόμεναι ἕως ἡμέρας θανάτου αὐτῶν χῆραι ζῶσαι.

\vs{4}Καὶ εἶπεν ὁ βασιλεὺς πρὸς Ἀμεσσαῒ, βόησόν μοι τὸν ἄνδρα Ἰούδα τρεῖς ἡμέρας, σὺ δὲ αὐτοῦ στῆθι.
\vs{5}Καὶ ἐπορεύθη Ἀμεσσαῒ τοῦ βοῆσαι τὸν Ἰούδαν, καὶ ἐχρόνισεν ἀπὸ τοῦ καιροῦ οὗ ἐτάξατο αὐτῷ Δαυίδ.
\vs{6}Καὶ εἶπε Δαυὶδ πρὸς Ἀμεσσαῒ, νῦν κακοποιήσει ἡμᾶς Σαβεὲ υἱὸς Βοχορὶ ὑπὲρ Ἀβεσσαλώμ· καὶ νῦν σὺ λάβε μετὰ σεαυτοῦ τοὺς παῖδας τοῦ κυρίου σου, καὶ καταδίωξον ὀπίσω αὐτοῦ, μή ποτε ἑαυτῷ εὕρῃ πόλεις ὀχυρὰς, καὶ σκιάσει τοὺς ὀφθαλμοὺς ἡμῶν.
\vs{7}Καὶ ἐξῆλθον ὀπίσω αὐτοὺ Ἀβεσσαῒ καὶ οἱ ἄνδρες Ἰωὰβ, καὶ ὁ Χερεθὶ, καὶ ὁ Φελεθὶ, καὶ πάντες οἱ δυνατοὶ, καὶ ἐξῆλθον ἐξ Ἱερουσαλὴμ διῶξαι ὀπίσω Σαβεὲ υἱοῦ Βοχορί.

\vs{8}Καὶ αὐτοὶ παρὰ τῷ λίθῳ τῷ μεγάλῳ τῷ ἐν Γαβαών· καὶ Ἀμεσσαῒ εἰσῆλθεν ἔμπροσθεν αὐτῶν· καὶ Ἰωὰβ περιεζωσμένος μανδύαν τὸ ἔνδυμα αὐτοῦ, καὶ ἐπʼ αὐτῷ ἐζωσμένος μάχαιραν ἐζευγμένην ἐπὶ τῆς ὀσφύος αὐτοῦ ἐν κολεῷ αὐτῆς· καὶ ἡ μάχαιρα ἐξῆλθε· καὶ αὐτὴ ἐξῆλθε καὶ ἔπεσε.

\vs{9}Καὶ εἶπεν Ἰωὰβ τῷ Ἀμεσσαῒ, εἰ ὑγιαίνεις σὺ, ἀδελφέ; καὶ ἐκράτησεν ἡ χεὶρ ἡ δεξιὰ Ἰωὰβ τοῦ πώγωνος Ἀμεσσαῒ τοῦ καταφιλῆσαι αὐτόν.
\vs{10}Καὶ Ἀμεσσαῒ οὐκ ἐφυλάξατο τὴν μάχαιραν τὴν ἐν τῇ χειρὶ Ἰωάβ· καὶ ἔπαισεν αὐτὸν ἐν αὐτῇ Ἰωὰβ εἰς τὴν ψόαν, καὶ ἐξεχύθη ἡ κοιλία αὐτοῦ εἰς τὴν γῆν, καὶ οὐκ ἐδευτέρωσεν αὐτῷ, καὶ ἀπέθανε· καὶ Ἰωὰβ καὶ Ἀβεσσαῒ ὁ ἀδελφὸς αὐτοῦ ἐδίωξεν ὀπίσω Σαβεὲ υἱοῦ Βοχορί.
\vs{11}Καὶ ἀνὴρ ἔστη ἐπʼ αὐτὸν τῶν παιδαρίων Ἰωὰβ, καὶ εἶπε, τίς ὁ βουλόμενος Ἰωὰβ, καὶ τίς τοῦ Δαυὶδ, ὀπίσω Ἰωάβ;
\vs{12}Καὶ Ἀμεσσαῒ πεφυρμένος ἐν τῷ αἵματι ἐν μέσὼ τῆς τρίβου· καὶ εἶδεν ἀνὴρ, ὅτι εἱστήκει πᾶς ὁ λαὸς, καὶ ἀπέστρεψε τὸν Ἀμεσσαῒ ἐκ τῆς τρίβου εἰς ἀγρόν· καὶ ἐπέῤῥιψεν ἐπʼ αὐτὸν ἱμάτιον, καθʼ ὅτι εἶδε πάντα τὸν ἐρχόμενον ἐπʼ αὐτὸν ἑστηκότα.
\vs{13}Ἡνίκα δὲ ἔφθασεν ἐκ τῆς τρίβου, παρῆλθε πᾶς ἀνὴρ Ἰσραὴλ ὀπίσω Ἰωὰβ τοῦ διῶξαι ὀπίσω Σαβεὲ υἱοῦ Βοχορί.

\vs{14}Καὶ διῆλθεν ἐν πάσαις φυλαῖς Ἰσραὴλ εἰς Ἀβὲλ καὶ εἰς Βαθμαχά· καὶ πάντες ἐν Χαῤῥὶ καὶ ἐξεκκλησιάσθησαν, καὶ ἦλθον κατόπισθεν αὐτοῦ.
\vs{15}Καὶ παρεγενήθησαν καὶ ἐπολιόρκουν ἐπʼ αὐτὸν ἐν Ἀβὲλ καὶ Φερμαχά· καὶ ἐξέχεαν πρόσχωμα πρὸς τὴν πόλιν, καὶ ἔστη ἐν τῷ προτειχίσματι· καὶ πᾶς ὁ λαὸς ὁ μετὰ Ἰωὰβ ἐνοοῦσαν καταβαλεῖν τὸ τεῖχος.
\vs{16}Καὶ ἐβόησε γυνὴ σοφὴ ἐκ τοῦ τείχους, καὶ εἶπεν, ἀκούσατε ἀκούσατε, εἴπατε δὴ πρὸς Ἰωὰβ, ἔγγισον ἕως ὧδε, καὶ λαλήσω πρὸς αὐτόν.

\vs{17}Καὶ προσήγγισε πρὸς αὐτήν· καὶ εἶπεν ἡ γυνὴ, εἰ σὺ εἶ Ἰωάβ; ὁ δὲ εἶπεν, ἐγώ. εἶπε δὲ αὐτῷ, ἄκουσον τοὺς λόγους τῆς δούλης σου· καὶ εἶπεν Ἰωὰβ, ἀκούω ἐγώ εἰμι.
\vs{18}Καὶ εἶπε λέγουσα, λόγον ἐλάλησαν ἐν πρώτοις, λέγοντες, ἠρωτημένος ἠρωτήθη ἐν τῇ Ἀβὲλ καὶ ἐν Δὰν εἰ ἐξέλιπον ἃ ἔθεντο οἱ πιστοὶ τοῦ Ἰσραήλ· ἐρωτῶντες ἐπερωτήσουσιν ἐν Ἀβὲλ, καὶ οὕτως εἰ ἐξέλιπον.
\vs{19}Ἐγώ εἰμι εἰρηνικὰ τῶν στηριγμάτων Ἰσραήλ· σὺ δὲ ζητεῖς θανατῶσαι πόλιν καὶ μητρόπολιν ἐν Ἰσραήλ· ἱνατί καταποντίζεις κληρονομίαν Κυρίου;
\vs{20}Καὶ ἀπεκρίθη Ἰωὰβ, καὶ εἶπεν, ἵλεώς μοι ἵλεώς μοι, εἰ καταποντιῶ καὶ εἰ φθερῶ.
\vs{21}Οὐχ οὕτως ὁ λόγος, ὅτι ἀνὴρ ἐξ ὄρους Ἐφραὶμ, Σαβεὲ υἱὸς Βοχορὶ ὄνομα αὐτοῦ, καὶ ἐπῇρε τὴν χεῖρα αὐτοῦ ἐπὶ τὸν βασιλέα Δαυίδ. δότε αὐτόν μοι μόνον, καὶ ἀπελεύσομαι ἀπάνωθεν τῆς πόλεως. Καὶ εἶπεν ἡ γυνὴ πρὸς Ἰωὰβ, ἰδοὺ ἡ κεφαλὴ αὐτοῦ ῥιφήσεται πρὸς σὲ διὰ τοῦ τείχους.

\vs{22}Καὶ εἰσῆλθεν ἡ γυνὴ πρὸς πάντα τὸν λαὸν, καὶ ἐλάλησε πρὸς πᾶσαν τὴν πόλιν ἐν τῇ σοφίᾳ αὐτῆς· καὶ ἀφεῖλε τὴν κεφαλὴν Σαβεὲ υἱοῦ Βοχορί· καὶ ἀφεῖλε καὶ ἔβαλε πρὸς Ἰωάβ· καὶ ἐσάλπισεν ἐν κερατίνῃ, καὶ διεσπάρησαν ἀπὸ τῆς πόλεως ἀπʼ αὐτοῦ ἀνὴρ εἰς τὰ σκηνώματα αὐτοῦ· καὶ Ἰωὰβ ἀπέστρεψεν εἰς Ἰερουσαλὴμ πρὸς τὸν βασιλέα.

\vs{23}Καὶ ὁ Ἰωὰβ πρὸς πάσῃ τῇ δυνάμει Ἰσραήλ· καὶ Βαναίας υἱὸς Ἰωδαὲ ἐπὶ τοῦ Χερεθὶ, καὶ ἐπὶ τοῦ Φελεθί·
\vs{24}Καὶ Ἀδωνιρὰμ ἐπὶ τοῦ φόρου· καὶ Ἰωσαφὰθ υἱὸς Ἀχιλοὺθ ἀναμιμνήσκων.
\vs{25}Καὶ Σουσὰ γραμματεύς· καὶ Σαδὼκ καὶ Ἀβιάθαρ ἱερεῖς·
\vs{26}Καί γε Ἰρὰς ὁ Ἰαρὶν ἦν ἱερεὺς τῷ Δαυίδ.

\ch{21}
Καὶ ἐγένετο λιμὸς ἐν ταῖς ἡμέραις Δαυὶδ τρία ἔτη, ἐνιαυτὸς ὁ ἐχόμενος ἐνιαυτοῦ· καὶ ἐζήτησε Δαυὶδ τὸ πρόσωπον Κυρίου· καὶ εἶπε Κύριος, ἐπὶ Σαοὺλ καὶ ἐπὶ τὸν οἶκον αὐτοῦ ἀδικία ἐν θανάτῳ αἱμάτων αὐτοῦ, περὶ οὗ ἐθανάτωσε τοὺς Γαβαωνίτας.
\vs{2}Καὶ ἐκάλεσεν ὁ βασιλεὺς Δαυὶδ τοὺς Γαβαωνίτας, καὶ εἶπε πρὸς αὐτούς· καὶ οἱ Γαβαωνίται οὐχ υἱοὶ Ἰσραήλ εἰσιν, ὅτι ἀλλʼ ἢ ἐκ τοῦ ἐλλείμματος τοῦ Ἀμοῤῥαίου, καὶ οἱ υἱοὶ Ἰσραὴλ ὤμοσαν αὐτοῖς· καὶ ἐζήτησε Σαοὺλ πατάξαι αὐτοὺς ἐν τῷ ζηλῶσαι αὐτὸν τοὺς υἱοὺς Ἰσραὴλ καὶ Ἰούδα.

\vs{3}Καὶ εἶπε Δαυὶδ πρὸς τοὺς Γαβαωνίτας, τί ποιήσω ὑμῖν, καὶ ἐν τίνι ἐξιλάσωμαι, καὶ εὐλογήσετε τὴν κληρονομίαν Κυρίου;
\vs{4}Καὶ εἶπαν αὐτῷ οἱ Γαβαωνίται, οὐκ ἔστιν ἡμῖν ἀργύριον ἢ χρυσίον μετὰ Σαοὺλ καὶ μετὰ τοῦ οἴκου αὐτοῦ, καὶ οὐκ ἔστιν ἡμῖν ἀνὴρ θανατῶσαι ἐν Ἰσραήλ. Καὶ εἶπε, τί ὑμεῖς λέγετε, καὶ ποιήσω ὑμῖν;
\vs{5}καὶ εἶπαν πρὸς τὸν βασιλέα, ὁ ἀνὴρ ὃς συνετέλεσεν ἐφʼ ἡμᾶς καὶ ἐδίωξεν ἡμᾶς, ὃς παρελογίσατο ἐξολοθρεῦσαι ἡμᾶς, ἀφανίσωμεν αὐτὸν, τοῦ μὴ ἑστάναι αὐτὸν ἐν παντὶ ὁρίῳ Ἰσραήλ.
\vs{6}Δότω ἡμῖν ἑπτὰ ἄνδρας ἐκ τῶν υἱῶν αὐτοῦ, καὶ ἐξηλιάσωμεν αὐτοὺς τῷ Κυρίῳ ἐν τῷ Γαβαὼν Σαοὺλ ἐκλεκτοὺς Κυρίου· καὶ εἶπεν ὁ βασιλεὺς, ἐγὼ δώσω.

\vs{7}Καὶ ἐφείσατο ὁ βασιλεὺς ἐπὶ Μεμφιβοσθὲ υἱὸν Ἰωνάθαν υἱοῦ Σαοὺλ διὰ τὸν ὅρκον Κυρίου τὸν ἀναμέσον αὐτῶν, καὶ ἀναμέσον Δαυὶδ, καὶ ἀναμέσον Ἰωνάθαν υἱοῦ Σαούλ.

\vs{8}Καὶ ἔλαβεν ὁ βασιλεὺς τοὺς δύο υἱοὺς Ῥεσφὰ θυγατρὸς Ἁϊᾶ, οὓς ἔτεκε τῷ Σαοὺλ, τὸν Ἐρμωνοῒ καὶ τὸν Μεμφιβοσθὲ, καὶ τοὺς πέντε υἱοὺς τῆς Μιχὸλ θυγατρὸς Σαοὺλ, οὓς ἔτεκε τῷ Ἐσδριὴλ υἱῷ Βερζελλὶ τῷ Μωουλαθί·
\vs{9}Καὶ ἔδωκεν αὐτοὺς ἐν χειρὶ τῶν Γαβαωνιτῶν, καὶ ἐξηλίασαν αὐτοὺς ἐν τῷ ὄρει ἔναντι Κυρίου· καὶ ἔπεσαν οἱ ἑπτὰ αὐτοὶ ἐπὶ τὸ αὐτό· καὶ αὐτοὶ δὲ ἐθανατώθησαν ἐν ἡμέραις θερισμοῦ ἐν πρώτοις, ἐν ἀρχῇ θερισμοῦ κριθῶν.
\vs{10}Καὶ ἔλαβε Ῥεσφὰ θυγάτηρ Ἁϊᾶ τὸν σάκκον, καὶ ἔπηξεν αὐτῇ πρὸς τὴν πέτραν ἐν ἀρχῇ θερισμοῦ κριθῶν, ἕως ἔσταξεν ἐπʼ αὐτοὺς ὕδωρ ἐκ τοῦ οὐρανοῦ· καὶ οὐκ ἔδωκε τὰ πετεινὰ τοῦ οὐρανοῦ καταπαῦσαι ἐπʼ αὐτοὺς ἡμέρας, καὶ τὰ θηρία τοῦ ἀγροῦ νυκτός.

\vs{11}Καὶ ἀπηγγέλη τῷ Δαυὶδ ὅσα ἐποίησε Ῥεσφὰ θυγάτηρ Ἁϊᾶ παλλακὴ Σαούλ· καὶ ἐξελύθησαν, καὶ κατέλαβεν αὐτοὺς Δὰν υἱὸς Ἰωὰ ἐκ τῶν ἀπογόνων τῶν γιγάντων.
\vs{12}Καὶ ἐπορεύθη Δαυὶδ καὶ ἔλαβε τὰ ὀστᾶ Σαοὺλ, καὶ τὰ ὀστᾶ Ἰωνάθαν τοῦ υἱοῦ αὐτοῦ, παρὰ τῶν ἀνδρῶν υἱῶν Ἰαβὶς Γαλαὰδ, οἳ ἔκλεψαν αὐτοὺς ἐκ τῆς πλατείας Βαιθσὰν, ὅτι ἔστησαν αὐτοὺς ἐκεῖ οἱ ἀλλόφυλοι ἐν τῇ ἡμέρᾳ ᾗ ἐπάταξαν οἱ ἀλλόφυλοι τὸν Σαοὺλ ἐν Γελβουέ.
\vs{13}Καὶ ἁνένεγκεν ἐκεῖθεν τὰ ὀστᾶ Σαοὺλ καὶ τὰ ὀστᾶ Ἰωνάθαν τοῦ υἱοῦ αὐτοῦ, καὶ συνήγαγε τὰ ὀστᾶ τῶν ἐξηλιασμένων.
\vs{14}Καὶ ἔθαψαν τὰ ὀστᾶ Σαοὺλ καὶ τὰ ὀστᾶ Ἰωνάθαν τοῦ υἱοῦ αὐτοῦ καὶ τὰ ὀστᾶ τῶν ἡλιασθένων ἐν γῇ Βενιαμὶν ἐν τῇ πλευρᾷ ἐν τῷ τάφῳ Κὶς τοῦ πατρὸς αὐτοῦ· καὶ ἐποίησαν πάντα ὅσα ἐνετείλατο ὁ βασιλεύς· καὶ ἐπήκουσεν ὁ Θεὸς τῇ γῇ μετὰ ταῦτα.

\vs{15}Καὶ ἐγενήθη ἔτι πόλεμος τοῖς ἀλλοφύλοις μετὰ Ἰσραήλ· καὶ κατέβη Δαυὶδ καὶ οἱ παῖδες αὐτοῦ μετʼ αὐτοῦ, καὶ ἐπολέμησαν μετὰ τῶν ἀλλοφύλων· καὶ ἐπορεύθη Δαυίδ.
\vs{16}Καὶ Ἰεσβὶ, ὃς ἦν ἐν τοῖς ἐκγόνοις τοῦ Ῥαφά, καὶ ὁ σταθμὸς τοῦ δόρατος αὐτοῦ, τριακοσίων σίκλων ὁλκῇ χαλκοῦ, καὶ αὐτὸς περιεζωσμένος κορύνην, καὶ διενοεῖτο τοῦ πατάξαι τὸν Δαυίδ.
\vs{17}Καὶ ἐβοήθησεν αὐτῷ Ἀβεσσὰ υἱὸς Σαρουίας, καὶ ἐπάταξε τὸν ἀλλόφυλον καὶ ἐθανάτωσεν αὐτόν· τότε ὤμοσαν οἱ ἄνδρες Δαυὶδ, λέγοντες, οὐκ ἐξελεύσῃ ἔτι μεθʼ ἡμῶν εἰς πόλεμον, καὶ οὐ μὴ σβέσῃς τὸν λύχνον Ἰσραήλ.

\vs{18}Καὶ ἐγενήθη μετὰ ταῦτα ἔτι πόλεμος ἐν Γὲθ μετὰ τῶν ἀλλοφύλων· τότε ἐπάταξε Σεβοχὰ ὁ Ἀστατωθὶ τὸν Σὲφ ἐν τοῖς ἐγγόνοις τοῦ Ῥαφά.

\vs{19}Καὶ ἐγένετο ὁ πόλεμος ἐν Ῥὸμ μετὰ τῶν ἀλλοφύλων· καὶ ἐπάταξεν Ἐλεανὰν υἱὸς Ἀριωργὶμ ὁ Βαιθλεεμίτης τὸν Γολιὰθ τὸν Γεθαῖον· καὶ τὸ ξύλον τοῦ δόρατος αὐτοῦ ὡς ἀντίον ὑφαινόντων.
\vs{20}Καὶ ἐγένετο ἔτι πόλεμος ἐν Γέθ· καὶ ἦν ἀνὴρ μαδὼν, καὶ οἱ δάκτυλοι τῶν χειρῶν αὐτοῦ, καὶ οἱ δάκτυλοι τῶν ποδῶν αὐτοῦ ἓξ καὶ ἓξ, εἰκοσιτέσσαρες ἀριθμῷ· καί γε αὐτὸς ἐτέχθη τῷ Ῥαφᾴ.
\vs{21}Καὶ ὠνείδισε τὸν Ἰσραὴλ, καὶ ἐπάταξεν αὐτὸν Ἰωνάθαν υἱὸς Σεμεῒ ἀδελφοῦ Δαυίδ.

\vs{22}Οἱ τέσσαρες οὗτοι ἐτέχθησαν ἀπόγονοι τῶν γιγάντων ἐν Γὲθ τῷ Ῥαφὰ οἶκος, καὶ ἔπεσαν ἐν χειρὶ Δαυὶδ, καὶ ἐν χειρὶ τῶν δούλων αὐτοῦ.

\ch{22}
Καὶ ἐλάλησε Δαυὶδ τῷ Κυρίῳ τοὺς λόγους τῆς ᾠδῆς ταύτης, ἐν ᾗ ἡμέρᾳ ἐξείλετο αὐτὸν Κύριος ἐκ χειρὸς πάντων τῶν ἐχθρῶν αὐτοῦ, καὶ ἐκ χειρὸς Σαούλ.
\vs{2}Καὶ εἶπεν ᾠδή·

Κύριε πέτρα μου, καὶ ὀχύρωμά μου, καὶ ἐξαιρούμενός με ἐμοί,
\vs{3}ὁ Θεός μου, φύλαξ μου ἔσται μοι, πεποιθὼς ἔσομαι ἐπʼ αὐτῷ· ὑπερασπιστής μου, καὶ κέρας σωτηρίας μου, ἀντιλήπτωρ μου, καὶ καταφυγή μου σωτηρίας μου, ἐξ ἀδίκου σώσεις με.

\vs{4}Αἰνετὸν ἐπικαλέσομαι Κύριον, καὶ ἐκ τῶν ἐχθρῶν μου σωθήσομαι.
\vs{5}Ὅτι περιέσχον με συντριμμοὶ θανάτου, χείμαῤῥοι ἀνομίας ἐθάμβησάν με.
\vs{6}Ὠδῖνες θανάτου ἐκύκλωσάν με, προέφθασάν με σκληρότητες θανάτου.
\vs{7}Ἐν τῷ θλίβεσθαί με ἐπικαλέσομαι τὸν Κύριον, καὶ πρὸς τὸν Θεόν μου βοήσομαι, καὶ ἐπακούσεται ἐκ ναοῦ αὐτοῦ φωνῆς μου, καὶ ἡ κραυγή μου ἐν τοῖς ὠσὶν αὐτοῦ.

\vs{8}Καὶ ἐταράχθη καὶ ἐσείσθη ἡ γῆ, καὶ τὰ θεμέλια τοῦ οὐρανοῦ συνεταράχθησαν καὶ ἐσπαράχθησαν, ὃτι ἐθυμώθη Κύριος αὐτοῖς.
\vs{9}Ἀνέβη καπνὸς ἐν τῇ ὀργῇ αὐτοῦ, καὶ πῦρ ἐκ στόματος αὐτοῦ κατέδεται· ἄνθρακες ἐξεκαύθησαν ἀπʼ αὐτοῦ.
\vs{10}Καὶ ἔκλινεν οὐρανοὺς καὶ κατέβη, καὶ γνόφος ὑποκάτω τῶν ποδῶν αὐτοῦ.
\vs{11}Καὶ ἐπεκάθισεν ἐπὶ τῷ χερουβὶν καὶ ἐπετάσθη, καὶ ὤφθη ἐπὶ πτερύγων ἀνέμου.
\vs{12}Καὶ ἔθετο σκότος ἀποκρυφὴν αὐτοῦ· κύκλῳ αὐτοῦ ἡ σκηνὴ αὐτοῦ σκότος ὑδάτων, ἐπάχυνεν ἐν νεφέλαις ἀέρος.
\vs{13}Ἀπὸ τοῦ φέγγους ἐναντίον αὐτοῦ ἐξεκαύθησαν ἄνθρακες πυρός.
\vs{14}Ἐβρόντησεν ἐξ οὐρανοῦ Κύριος, καὶ ὁ ὕψιστος ἔδωκε φωνὴν αὐτοῦ.
\vs{15}Καὶ ἀπέστειλε βέλη, καὶ ἐσκόρπισεν αὐτούς· καὶ ἢστραψεν ἀστραπὴν, καὶ ἐξέστησεν αὐτούς.
\vs{16}Καὶ ὤφθησαν ἀφέσεις θαλάσσης, καὶ ἀπεκαλύφθη θεμέλια τῆς οἰκουμένης ἐν τῇ ἐπιτιμήσει Κυρίου, ἀπὸ πνοῆς πνεύματος θυμοῦ αὐτοῦ.
\vs{17}Ἀπέστειλεν ἐξ ὕξ ὕψους καὶ ἔλαβέ με, εἵλκυσέ με ἐξ ὑδάτων πολλῶν.
\vs{18}Ἐῤῥύσατό με ἐξ ἐχθρῶν μου ἰσχύος, ἐκ τῶν μισούντων με, ὅτι ἐκραταιώθησαν ὑπὲρ ἐμέ.

\vs{19}Προέφθασάν με ἡμέραι θλίψεώς μου· καὶ ἐγένετο Κύριος ἐπιστήριγμά μου;
\vs{20}καὶ ἐξήγαγέ με εἰς πλατυσμὸν, καὶ ἐξείλετό με, ὅτι ηὐδόκησες ἐν ἐμοί.
\vs{21}Καὶ ἀνταπέδωκέ μοι Κύριος κατὰ τὴν δικαιοσύνην μου, καὶ κατὰ τὴν καθαριότητα τῶν χειρῶν μου ἀνταπέδωκέ μοι.
\vs{22}Ὅτι ἐφύλαξα ὁδοὺς Κυρίου, καὶ οὐκ ἠσέβησα ἀπὸ τοῦ Θεοῦ μου.
\vs{23}Ὅτι πάντα τὰ κρίματα αὐτοῦ κατεναντίον μου καὶ τὰ δικαιώματα αὐτοῦ, οὐκ ἀπέστην ἀπʼ αὐτῶν.
\vs{24}Καὶ ἔσομαι ἄμωμος αὐτῷ, καὶ προφυλάξομαι ἀπὸ τῆς ἀνομίας μου.
\vs{25}Καὶ ἀποδώσει μοι Κύριος κατὰ τὴν δικαιοσύνην μου, καὶ κατὰ τὴν καθαριότητα τῶν χειρῶν μου ἐνώπιον τῶν ὀφθαλμῶν αὐτοῦ.

\vs{26}Μετὰ ὁσίου ὁσιωθήσῃ, καὶ μετὰ ἀνδρὸς τελείου τελειωθήσῃ·
\vs{27}Καὶ μετὰ ἐκλεκτοῦ ἐκλεκτὸς ἔσῃ, καὶ μετὰ στρεβλοῦ στρεβλωθήσῃ.
\vs{28}Καὶ τὸν λαὸν τὸν πτωχὸν σώσεις, καὶ ὀφθαλμοὺς ἐπὶ μετεώρων ταπεινώσεις.
\vs{29}Ὅτι σὺ ὁ λύχνος μου Κύριε, καὶ Κύριος ἐκλάμψει μοι τὸ σκότος μου.
\vs{30}Ὅτι ἐν σοὶ δραμοῦμαι μονόζωνος, καὶ ἐν τῷ Θεῷ μου ὑπερβήσομαι τεῖχος.

\vs{31}Ὁ ἰσχυρὸς ἄμωμος ἡ ὁδὸς αὐτοῦ· τὸ ῥῆμα Κυρίου κραταιὸν πεπυρωμένον· ὑπερασπιστής ἐστι πᾶσι τοῖς πεποιθόσιν ἐπʼ αὐτόν.
\vs{32}Τίς ἰσχυρὸς πλὴν Κυρίου; καὶ τίς κτίστης ἔσται πλὴν τοῦ Θεοῦ ἡμῶν;
\vs{33}Ὁ ἰσχυρὸς ὁ κραταιῶν με δυνάμει, καὶ ἐξετίναξεν ἄμωμον τὴν ὁδόν μου.
\vs{34}Τιθεὶς τοὺς πόδας μου ὡς ἐλάφων, καὶ ἐπὶ τὰ ὕψη ἱστῶν με.
\vs{35}Διδάσκων χεῖράς μου εἰς πόλεμον, καὶ κατάξας τόξον χαλκοῦν ἐν βραχίονί μου.
\vs{36}Καὶ ἔδωκάς μοι ὑπερασπισμὸν σωτηρίας μου, καὶ ἡ ὑπακοή σου ἐπλήθυνέ με εἰς πλατυσμὸν εἰς τὰ διαβήματά μου ὑποκάτω μου,
\vs{37}καὶ οὐκ ἐσαλεύθησαν τὰ σκέλη μου.

\vs{38}Διώξω ἐχθρούς μου, καὶ ἀφανιῶ αὐτοὺς, καὶ οὐκ ἀναστρέψω ἕως ἂν συντελέσω αὐτούς.
\vs{39}Καὶ θλάσω αὐτοὺς καὶ οὐκ ἀναστήσονται, καὶ πεσοῦνται ὑπὸ τοὺς πόδας μου.
\vs{40}Καὶ ἐνισχύσεις με δυνάμει εἰς πόλεμον, κάμψεις τοὺς ἐπιστανομένους μοι ὑποκάτω μου.
\vs{41}Καὶ τοὺς ἐχθρούς μου ἔδωκάς μοι νῶτον, τοὺς μισοῦντάς με, καὶ ἐθανάτωσας αὐτούς.
\vs{42}Βοήσονται, καὶ οὐκ ἔστι βοηθὸς, πρὸς Κύριον, καὶ οὐκ ἐπήκουσεν αὐτῶν.
\vs{43}Καὶ ἐλέανα αὐτοὺς ὡς χοῦν γῆς, ὡς πηλὸν ἐξόδων ἐλέπτυνα αὐτούς.
\vs{44}Καὶ ῥύσῃ με ἐκ μάχης λαῶν, φυλάξεις με εἰς κεφαλὴν ἐθνῶν· λαὸς ὃν οὐκ ἔγνω ἐδούλευσάν μοι.
\vs{45}Υἱοὶ ἀλλότριοι ἐψεύσαντό μοι, εἰς ἀκοὴν ὠτίου ἤκουσάν μου.
\vs{46}Υἱοὶ ἀλλότριοι ἀποῤῥιφήσονται, καὶ σφαλοῦσιν ἐκ τῶν συγκλεισμῶν αὐτῶν.

\vs{47}Ζῇ Κύριος, καὶ εὐλογητὸς ὁ φύλαξ μου, καὶ ὑψωθήσεται ὁ Θεός μου ὁ φύλαξ τῆς σωτηρίας μου.
\vs{48}Ἰσχυρὸς Κύριος ὁ διδοὺς ἐκδικήσεις ἐμοὶ, παιδεύων λαοὺς ὑποκάτω μου,
\vs{49}καὶ ἐξάγων με ἐξ ἐχθρῶν μου· καὶ ἐκ τῶν ἐπεγειρομένων μοι ὑψώσεις με, ἐξ ἀνδρὸς ἀδικημάτων ῥύσῃ με.
\vs{50}Διὰ τοῦτο ἐξομολογήσομαί σοι Κύριε ἐν τοῖς ἔθνεσι, καὶ ἐν τῷ ὀνόματί σου ψαλῶ.
\vs{51}Μεγαλύνων τὰς σωτηρίας βασιλέως αὐτοῦ, καὶ ποιῶν ἔλεος τῷ χριστῷ αὐτοῦ τῷ Δαυὶδ, καὶ τῷ σπέρματι αὐτοῦ ἕως αἰῶνος.

Καὶ οὗτοι οἱ λόγοι Δαυὶδ οἱ ἔσχατοι·

\ch{23}
Πιστὸς Δαυὶδ υἱὸς Ἰεσσαί, καὶ πιστὸς ἀνὴρ ὃν ἀνέστησε Κύριος ἐπὶ χριστὸν Θεοῦ Ἰακὼβ, καὶ εὐπρεπεῖς ψαλμοὶ Ἰσραήλ.

\vs{2}Πνεῦμα Κυρίου ἐλάλησεν ἐν ἐμοὶ, καὶ ὁ λόγος αὐτοῦ ἐπὶ γλώσσης μου.
\vs{3}Λέγει ὁ Θεὸς Ἰσραὴλ, ἐμοὶ ἐλάλησε φύλαξ ἐξ Ἰσραὴλ παραβολήν· εἶπὸν ἐν ἀνθρώπῳ, πῶς κραταιώσητε φόβον χριστοῦ;
\vs{4}Καὶ ἐν φωτὶ Θεοῦ πρωΐας, ἀνατείλαι ἥλιος τοπρωῒ, οὗ Κύριος παρῆλθεν ἐκ φέγγους, καὶ ὡς ἐξ ὑετοῦ χλόης ἀπὸ γῆς.
\vs{5}Οὐ γὰρ οὕτως ὁ οἶκός μου μετὰ ἰσχυροῦ, διαθήκην γὰρ αἰώνιον ἔθετό μοι ἑτοίμην ἐν παντὶ καιρῷ πεφυλαγμένην· ὅτι πᾶσα σωτηρία μου καὶ πᾶν θέλημα, ὅτι οὐ μὴ βλαστήσῃ ὁ παράνομος.
\vs{6}Ὥσπερ ἄκανθα ἐξωσμένη πάντες οὗτοι, ὅτι οὐ χειρὶ ληφθήσονται,
\vs{7}καὶ ἀνὴρ οὐ κοπιάσει ἐν αὐτοῖς· καὶ πλῆρες σιδήρου, καὶ ξύλον δόρατος, καὶ ἐν πυρὶ καύσει, καὶ καυθήσονται αἰσχύνην αὐτῶν.

\vs{8}Ταῦτα τὰ ὀνόματα τῶν δυνατῶν Δαυίδ· Ἰεβοσθὲ ὁ Χαναναῖος ἄρχων τοῦ τρίτου ἐστίν· Ἀδινὼν ὁ Ἀσωναῖος, οὗτος ἐσπάσατο τὴν ῥομφαίαν αὐτοῦ ἐπὶ ὀκτακοσίους στρατιώτας εἰσάπαξ.
\vs{9}Καὶ μετʼ αὐτὸν Ἐλεανὰν υἱὸς πατραδέλφου αὐτοῦ υἱὸς Δουδὶ τοῦ ἐν τοῖς τρισὶ δυνατοῖς μετὰ Δαυίδ· καὶ ἐν τῷ ὀνειδίσαι αὐτὸν ἐν τοῖς ἀλλοφύλοις, συνήχθησαν ἐκεῖ εἰς πόλεμον, καὶ ἀνέβησεν ἀνὴρ Ἰσραήλ.
\vs{10}Αὐτὸς ἀνέστη καὶ ἐπάταξεν ἐν τοῖς ἀλλοφύλοις, ἕως οὗ ἐκοπίασεν ἡ χεὶρ αὐτοῦ, καὶ προσεκολλήθη ἡ χεὶρ αὐτοῦ πρὸς τὴν μάχαιραν· καὶ ἐποίησε Κύριος σωτηρίαν μεγάλην ἐν τῇ ἡμέρᾳ ἐκείνῃ· καὶ ὁ λαὸς ἐκάθητο ὀπίσω αὐτοῦ πλὴν ἐκδιδύσκειν.

\vs{11}Καὶ μετʼ αὐτὸν Σαμαΐα υἱὸς Ἄσα ὁ Ἀρουχαῖος· καὶ συνήχθησαν οἱ ἀλλόφυλοι εἰς Θηρία· καὶ ἦν ἐκεῖ μερὶς τοῦ ἀγροῦ πλήρης φακοῦ· καὶ ὁ λαὸς ἔφυγεν ἐκ προσώπου ἀλλοφύλων.
\vs{12}Καὶ ἐστηλώθη ἐν μέσῳ τῆς μερίδος, καὶ ἐξείλατο αὐτὴν, καὶ ἐπάταξε τοὺς ἀλλοφύλους· καὶ ἐποίησε Κύριος σωτηρίαν μεγάλην.

\vs{13}Καὶ κατέβησαν τρεῖς ἀπὸ τῶν τριάκοντα, καὶ κατέβησαν εἰς Κασὼν πρὸς Δαυὶδ, εἰς τὸ σπήλαιον Ὀδολλάμ· καὶ τάγμα τῶν ἀλλοφύλων, καὶ παρενέβαλον ἐν τῇ κοιλάδι Ῥαφαΐν.
\vs{14}Καὶ Δαυὶδ τότε ἐν τῇ περιοχῇ, καὶ τὸ ὑπόστεμα τῶν ἀλλοφύλων τότε ἐν Βηθλεέμ.
\vs{15}Καὶ ἐπεθύμησε Δαυὶδ καὶ εἶπε, τίς ποτιεῖ με ὕδωρ ἐκ τοῦ λάκκου τοῦ ἐν Βηθλεὲμ τοῦ ἐν τῇ πύλῃ; τὸ δὲ σύστεμα τῶν ἀλλοφύλων τότε ἐν Βηθλεέμ.
\vs{16}Καὶ διέῤῥηξαν οἱ τρεῖς δυνατοὶ ἐν τῇ παρεμβολῇ τῶν ἀλλοφύλων, καὶ ὑδρεύσαντο ὕδωρ ἐκ τοῦ λάκκου τοῦ ἐν Βηθλεὲμ τοῦ ἐν τῇ πύλῃ· καὶ ἔλαβαν, καὶ παρεγένοντο πρὸς Δαυίδ, καὶ οὐκ ἡθέλησε πιεῖν αὐτό· καὶ ἔσπεισεν αὐτὸ τῷ Κυρίῳ.
\vs{17}Καὶ εἶπεν, ἵλεώς μοι Κύριε τοῦ ποιῆσαι τοῦτο, εἰ αἷμα τῶν ἀνδρῶν τῶν πορευθέντων ἐν ταῖς ψυχαῖς αὐτῶν πίομαι· καὶ οὐκ ἠθέλησε πιεῖν αὐτό. Ταῦτα ἐποίησαν οἱ τρεῖς δυνατοί.

\vs{18}Καὶ Ἀβεσσὰ ὁ ἀδελφὸς Ἰωάβ υἱὸς Σαρουίας αὐτὸς ἄρχων ἐν τοῖς τρισὶ, καὶ αὐτὸς ἐξήγειρε τὸ δόρυ αὐτοῦ ἐπὶ τριακοσίους τραυματίας· καὶ αὐτῷ ὄνομα ἐν τοῖς τρισὶν.
\vs{19}Ἐκ τῶν τριῶν ἐκείνων ἔνδοξος, καὶ ἐγένετο αὐτοῖς εἰς ἄρχοντα, καὶ ἕως τῶν τριῶν οὐκ ἦλθε.

\vs{20}Καὶ Βαναίας υἱὸς Ἰωδαὲ ἀνὴρ αὐτὸς πολλοστὸς ἔργοις, ἀπὸ Καβεσεήλ, καὶ αὐτὸς ἐπάταξε τοὺς δύο υἱοὺς Ἀριὴλ τοῦ Μωάβ· καὶ αὐτός κατέβη καὶ ἐπάταξε τὸν λέοντα ἐν μέσῳ τοῦ λάκκου ἐν τῇ ἡμέρᾳ τῆς χιόνος.
\vs{21}Αὐτὸς ἐπάταξε τὸν ἄνδρα τὸν Αἰγύπτιον, ἄνδρα ὁρατὸν, ἐν δὲ τῇ χειρὶ τοῦ Αἰγυπτίου δόρυ ὡς ξύλον διαβάθρας· καὶ κατέβη πρὸς αὐτὸν ἐν ῥἀβδῳ, καὶ ἥρπασε τὸ δόρυ ἐκ τῆς χειρὸς τοῦ Αἰγυπτίου, καὶ ἀπέκτεινεν αὐτὸν ἐν τῷ δόρατι αὐτοῦ.
\vs{22}Ταῦτα ἐποίησε Βαναίας υἱὸν Ἰωδαὲ, καὶ αὐτῷ ὄνομα ἐν τοῖς τρισὶ τοῖς δυνατοῖς,
\vs{23}ἐκ τῶν τριῶν ἔνδοξος, καὶ πρὸς τοὺς τρεῖς οὐκ ἦλθε· καὶ ἔταξεν αὐτὸν Δαυὶδ πρὸς τὰς ἀκοὰς αὐτοῦ.

Καὶ ταῦτα τὰ ὀνόματα τῶν δυνατῶν Δαυὶδ τοῦ βασιλέως.
\vs{24}Ἀσαὴλ ἀδελφὸς Ἰωάβ· οὗτος ἐν τοῖς τριάκοντα· Ἐλεανὰν υἱὸς Δουδὶ πατραδέλφου αὐτοῦ ἐν Βηθλεέμ·
\vs{25}Σαιμὰ ὁ Ῥουδαῖος·
\vs{26}Σελλὴς ὁ Κελωθί· Ἴρας υἱὸς Ἴσκα ὁ Θεκωΐτης·
\vs{27}Ἀβιέζερ ὁ Ἀνωθίτης, ἐκ τῶν υἱῶν τοῦ Ἀνωθίτου·
\vs{28}Ἐλλὼν ὁ Ἀωΐτης· Νοερὲ ὁ Νετωφατίτης·
\vs{29}Ἐσθαὶ υἱὸς Ῥιβὰ ἐκ Γαβαὲθ υἱὸς Βενιαμὶν τοῦ Ἐφραθαίου· Ἀσμὼθ ὁ Βαρδιαμίτης·
\vs{32}Ἐμασοὺ ὁ Σαλαβωνίτης· υἱοὶ Ἀσάν, Ἰωνάθαν·
\vs{33}Σαμνὰν ὁ Ἁρωδίτης· Ἀμνὰν υἱὸς Ἀραῒ Σαραουρίτης·
\vs{34}Ἀλιφαλὲθ υἱὸς τοῦ Ἀσβίτου, υἱὸς τοῦ Μαχαχαχί· Ἐλιὰβ υἱὸς Ἀχιτόφελ τοῦ Γελωνίτου·
\vs{35}Ἀσαραῒ ὁ Καρμήλιος τοῦ Οὐραιοερχί·
\vs{36}Γάαλ υἱὸς Ναθανά· πολυδυνάμεως υἱὸς Γαλααδδί·
\vs{37}Ἐλιὲ ὁ Ἀμμανίτης·
\vs{37a}Ἁδροὶ ἀπὸ χειμάῤῥων·
\vs{37b}Γαδαβιὴλ υἱὸς τοῦ Ἀραβωθαίου·
\vs{37c}Γελωρὲ ὁ Βηθωραῖος αἴρων τὰ σκεύη· Ἰωὰβ υἱὸς Σαρουίας·
\vs{38}Ἴρας ὁ Ἐθιραῖος· Γηρὰβ ὁ Ἐθεναῖος·
\vs{39}Οὐρίας ὁ Χετταῖος· οἱ πάντες τριάκοντα καὶ ἑπτά.

\ch{24}
Καὶ προσέθετο ὀργὴν Κύριος ἐκκαῆναι ἐν Ἰσραήλ, καὶ ἐπέσεισε τὸν Δαυὶδ ἐν αὐτοῖς, λέγων, βάδιζε, ἀρίθμησον τὸν Ἰσραὴλ καὶ τὸν Ἰούδαν.
\vs{2}Καὶ εἶπεν ὁ βασιλεὺς πρὸς Ἰωὰβ ἄρχοντα τῆς ἰσχύος τὸν μετʼ αὐτοῦ, διέλθε δὴ πάσας φυλὰς Ἰσραὴλ καὶ Ἰούδα, ἀπὸ Δὰν καὶ ἕως Βηρσαβεέ, καὶ ἐπίσκεψαι τὸν λαὸν, καὶ γνώσομαι τὸν ἀριθμὸν τοῦ λαοῦ.
\vs{3}Καὶ εἶπεν Ἰωὰβ πρὸς τὸν βασιλέα, καὶ προσθείη Κύριος ὁ Θεὸς πρὸς τὸν λαὸν ὥσπερ αὐτοὺς καὶ ὥσπερ αὐτοὺς ἑκατονταπλασίονα, καὶ ὀφθαλμοὶ τοῦ κυρίου μου τοῦ βασιλέως ὁρῶντες· καὶ ὁ κύριός μου ὁ βασιλεὺς ἱνατί βούλεται ἐν τῷ λόγῳ τούτῳ;
\vs{4}Καὶ ὑπερίσχυσεν ὁ λόγος τοῦ βασιλέως πρὸς Ἰωὰβ καὶ εἰς τοὺς ἄρχοντας τῆς δυνάμεως·

Καὶ ἐξῆλθεν Ἰωὰβ καὶ οἱ ἄρχοντες τῆς ἰσχύος ἐνώπιον τοῦ βασιλέως ἐπισκέψασθαι τὸν λαὸν τὸν Ἰσραήλ.
\vs{5}Καὶ διέβησαν τὸν Ἰορδάνην, καὶ παρενέβαλον ἐν Ἀροὴρ ἐκ δεξιῶν τῆς πόλεως τῆς ἐν μέσῳ τῆς φάραγγος Γὰδ καὶ Ἐλιέζερ.
\vs{6}Καὶ ἦλθον εἰς Γαλαὰδ καὶ εἰς γῆν Θαβασὼν, ἥ ἐστιν Ἀδασαὶ, καὶ παρεγένοντο εἰς Δανιδὰν καὶ Οὐδὰν, καὶ ἐκύκλωσαν Σιδῶνα.
\vs{7}Καὶ ἦλθον εἰς Μάψαρ Τύρου, καὶ εἰς πάσας τὰς πόλεις τοῦ Εὐαίου καὶ τοῦ Χαναναίου· καὶ ἦλθαν κατὰ Νότον Ἰούδα εἰς Βηρσαβεὲ,
\vs{8}καὶ περιώδευσαν ἐν πάσῃ τῇ γῇ· καὶ παρεγένοντο ἀπὸ τέλους ἐννέα μηνῶν καὶ εἴκοσι ἡμερῶν εἰς Ἰερουσαλήμ.
\vs{9}Καὶ ἔδωκεν Ἰωὰβ τὸν ἀριθμὸν τῆς ἐπισκέψεως τοῦ λαοῦ πρὸς τὸν βασιλέα· καὶ ἐγένετο Ἰσραὴλ, ὀκτακόσιαι χιλιάδες ἀνδρῶν δυνάμεως σπωμένων ῥομφαίαν· καὶ ἀνὴρ Ἰούδα, πεντακόσιαι χιλιάδες ἀνδρῶν μαχητῶν.

\vs{10}Καὶ ἐπάταξε καρδία Δαυὶδ αὐτὸν μετὰ τὸ ἀριθμῆσαι τὸν λαόν· καὶ εἶπε Δαυὶδ πρὸς Κύριον, ἥμαρτον σφόδρα ὃ ἐποίησα νῦν Κύριε· παραβίβασον δὴ τὴν ἀνομίαν τοῦ δούλου σου, ὅτι ἐμωράνθην σφόδρα.

\vs{11}Καὶ ἀνέστη Δαυὶδ τοπρωΐ· καὶ λόγος Κυρίου ἐγένετο πρὸς Γὰδ τὸν προφήτην τὸν ὁρῶντα, λέγων,
\vs{12}πορεύθητι, καὶ λάλησον πρὸς Δαυὶδ, λέγων, τάδε λέγει Κύριος, τρία ἐγώ εἰμι αἴρω ἐπὶ σὲ, καὶ ἔκλεξαι σεαυτῷ ἓν ἐξ αὐτῶν, καὶ ποιήσω σοι.
\vs{13}Καὶ εἰσῆλθε Γὰδ πρὸς Δαυὶδ, καὶ ἀνήγγειλε, καὶ εἶπεν αὐτῷ, ἔκλεξαι σεαυτῷ γενέσθαι, εἰ ἔλθῃ σοι τρία ἔτη λιμὸς ἐν τῇ γῇ σου, ἢ τρεῖς μῆνας φεύγειν σε ἔμπροσθεν τῶν ἐχθρῶν σου, καὶ ἔσονται διώκοντές σε, ἢ γενέσθαι τρεῖς ἡμέρας θάνατον ἐν τῇ γῇ σου· νῦν οὖν γνῶθι καὶ ἴδε τί ἀποκριθῶ τῷ ἀποστείλαντί με ῥῆμα.
\vs{14}Καὶ εἶπε Δαυὶδ πρὸς Γὰδ, στενά μοι πάντοθεν σφόδρα ἐστίν· ἐμπεσοῦμαι δὴ εἰς χεῖρας Κυρίου, ὅτι πολλοὶ οἱ οἰκτιρμοὶ αὐτοῦ σφόδρα· εἰς δὲ χεῖρας ἀνθρώπου οὐ μὴ ἐμπέσω.

\vs{15}Καὶ ἐξελέξατο ἑαυτῷ Δαυὶδ τὸν θάνατον· καὶ ἡμέραι θερισμοῦ πυρῶν· καὶ ἔδωκε Κύριος θάνατον ἐν Ἰσραὴλ ἀπὸ πρωΐθεν ἕως ὥρας ἀρίστου, καὶ ἤρξατο ἡ θραῦσις ἐν τῷ λαῷ· καὶ ἀπέθανεν ἐκ τοῦ λαοῦ ἀπὸ Δὰν καὶ ἕως Βηρσαβεὲ, ἑβδομήκοντα χιλιάδες ἀνδρῶν.
\vs{16}Καὶ ἐξέτεινεν ὁ ἄγγελος τοῦ Θεοῦ τὴν χεῖρα αὐτοῦ εἰς Ἱερουσαλὴμ τοῦ διαφθεῖραι αὐτὴν, καὶ παρεκλήθη Κύριος ἐπὶ τῇ κακίᾳ, καὶ εἶπε τῷ ἀγγέλῳ τῷ διαφθείροντι ἐν τῷ λαῷ, πολὺ νῦν, ἄνες τὴν χεῖρά σου· καὶ ὁ ἄγγελος Κυρίου ἦν παρὰ τῇ ἅλῳ Ὀρνὰ τοῦ Ἰεβουσαίου.
\vs{17}Καὶ εἶπε Δαυὶδ πρὸς Κύριον, ἐν τῷ ἰδεῖν αὐτὸν τὸν ἄγγελον τὸν τύπτοντα ἐν τῷ λαῷ, καὶ εἶπεν, ἰδοὺ, ἐγώ εἰμι ἠδίκησα· καὶ οὗτοι τὰ πρόβατα τί ἐποίησαν; γενέσθω δὴ ἡ χείρ σου ἐν ἐμοὶ, καὶ ἐν τῷ οἴκῳ τοῦ πατρός μου.

\vs{18}Καὶ ἦλθε Γὰδ πρὸς Δαυὶδ ἐν τῇ ἡμέρᾳ ἐκείνῃ, καὶ εἶπεν αὐτῷ, ἀνάβηθι, καὶ στῆσον τῷ Κυρίῳ θυσιαστήριον ἐν τῷ ἅλωνι Ὀρνὰ τοῦ Ἰεβουσαίου.
\vs{19}Καὶ ἀνέβη Δαυὶδ κατὰ τὸν λόγον Γὰδ, καθʼ ὃν τρόπον ἐνετείλατο αὐτῷ Κύριος.
\vs{20}Καὶ διέκυψεν Ὀρνὰ, καὶ εἶδε τὸν βασιλέα καὶ τοὺς παῖδας αὐτοῦ παραπορευομένους ἐπάνω αὐτοῦ· καὶ ἐξῆλθεν Ὀρνὰ, καὶ προσεκύνησε τῷ βασιλεῖ ἐπὶ πρόσωπον αὐτοῦ ἐπὶ τὴν γῆν.
\vs{21}Καὶ εἶπεν Ὀρνὰ, τί ὅτι ἦλθεν ὁ κύριός μου ὁ βασιλεὺς πρὸς τὸν δοῦλον αὐτοῦ; καὶ εἶπε Δαυὶδ, κτήσασθαι παρὰ σοῦ τὸν ἅλωνα τοῦ οἰκοδομῆσαι θυσιαστήριον τῷ Κυρίῳ, καὶ συσχεθῇ ἡ θραῦσις ἐπάνω τοῦ λαοῦ.
\vs{22}Καὶ εἶπεν Ὀρνὰ πρὸς Δαυὶδ, λαβέτω καὶ ἀνενεγκάτω ὁ κύριός μου ὁ βασιλεὺς τῷ Κυρίῳ τὸ ἀγαθὸν ἐν ὀφθαλμοῖς αὐτοῦ· ἰδοὺ οἱ βόες εἰς ὁλοκαύτωμα, καὶ οἱ τροχοὶ καὶ τὰ σκεύη τῶν βοῶν εἰς ξύλα.
\vs{23}Τὰ πάντα ἔδωκεν Ὀρνὰ τῷ βασιλεῖ· καὶ εἶπεν Ὀρνὰ πρὸς τὸν βασιλέα, Κύριος ὁ Θεός σου εὐλογήσαι σε.
\vs{24}Καὶ εἶπεν ὁ βασιλεὺς πρὸς Ὀρνὰ, οὐχὶ, ὅτι ἀλλὰ κτώμενος κτήσομαι παρὰ σοῦ ἐν ἀναλλάγματι, καὶ οὐκ ἀνοίσω τῷ Κυρίῳ μου Θεῷ ὁλοκαύτωμα δωρεάν. Καὶ ἐκτήσατο Δαυὶδ τὸν ἅλωνα καὶ τοὺς βόας ἐν ἀργυρίῳ σίκλων πεντήκοντα.
\vs{25}Καὶ ᾠκοδόμησεν ἐκεῖ Δαυὶδ θυσιαστήριον Κυρίῳ, καὶ ἀνήνεγκεν ὁλοκαυτώσεις καὶ εἰρηνικάς· καὶ προσέθηκε Σαλωμὼν ἐπὶ τὸ θυσιαστήριον ἐπʼ ἐσχάτῳ, ὅτι μικρὸν ἦν ἐν πρώτοις· καὶ ἐπήκουσε Κύριος τῇ γῇ, καὶ συνεσχέθη ἡ θραῦσις ἐπάνωθεν Ἰσραήλ.


\def\book{ΒΑΣΙΛΕΙΩΝ Γ}
\biblebook{ΒΑΣΙΛΕΙΩΝ Γ}


\lettrine[lines=2, loversize=0.2, nindent=0em, findent=.25em]{\textcolor{bookheadingcolor}{Κ}}{ΑΙ} ὁ βασιλεὺς Δαυὶδ πρεσβύτερος προβεβηκὼς ἡμέραις, καὶ περιέβαλλον αὐτὸν ἱματίοις, καὶ οὐκ ἐθερμαίνετο.
\vs{2}Καὶ εἶπον οἱ παῖδες αὐτοῦ, ζητησάτωσαν τῷ βασιλεῖ παρθένον νεάνιδα, καὶ παραστήσεται τῷ βασιλεῖ, καὶ ἔσται αὐτὸν θάλπουσα, καὶ κοιμηθήσεται μετʼ αὐτοῦ, καὶ θερμανθήσεται ὁ κύριός μου ὁ βασιλεύς.
\vs{3}Καὶ ἐζήτησαν νεάνιδα καλὴν ἐκ παντὸς ὁρίου Ἰσραήλ· καὶ εὗρον τὴν Ἀβισὰγ τὴν Σωμανίτιν, καὶ ἤνεγκαν αὐτὴν πρὸς τὸν βασιλέα.
\vs{4}Καὶ ἡ νεᾶνις καλὴ ἕως σφόδρα· καὶ ἦν θάλπουσα τὸν βασιλέα, καὶ ἐλειτούργει αὐτῷ· καὶ ὁ βασιλεὺς οὐκ ἔγνω αὐτήν.

\vs{5}Καὶ Ἀδωνίας υἱὸς Ἀγγὶθ ἐπῄρετο, λέγων, ἐγὼ βασιλεύσω· καὶ ἐποίησεν ἑαυτῷ ἅρματα καὶ ἱππεῖς, καὶ πεντήκοντα ἄνδρας παρατρέχειν ἔμπροσθεν αὐτοῦ.
\vs{6}Καὶ οὐκ ἀπεκώλυσεν αὐτὸν ὁ πατὴρ αὐτοῦ οὐδέποτε, λέγων, διατί σὺ ἐποίησας; καὶ γε αὐτὸς ὡραῖος τῇ ὄψει σφόδρα, καὶ αὐτὸν ἔτεκεν ὀπίσω Ἀβεσσαλώμ.
\vs{7}Καὶ ἐγένοντο οἱ λόγοι αὐτοῦ μετὰ Ἰωὰβ τοῦ υἱοῦ Σαρουίας, καὶ μετὰ Ἀβιάθαρ τοῦ ἱερέως, καὶ ἐβοήθουν ὀπίσω Ἀδωνίου.
\vs{8}Καὶ Σαδὼκ ὁ ἱερεὺς, καὶ Βαναίας υἱὸς Ἰωδαὲ, καὶ Νάθαν ὁ προφήτης, καὶ Σεμεῒ, καὶ Ῥησὶ, καὶ υἱοὶ δυνατοὶ τοῦ Δαυὶδ, οὐκ ἦσαν ὀπίσω Ἀδωνίου.
\vs{9}Καὶ ἐθυσίασεν Ἀδωνίας πρόβατα καὶ μόσχους καὶ ἄρνας μετὰ αἰθῆ τοῦ Ζωελεθὶ, ὃς ἦν ἐχόμενα τῆς Ῥωγήλ· καὶ ἐκάλεσε πάντας τοὺς ἀδελφοὺς αὐτοῦ, καὶ πάντας τοὺς ἁδροὺς Ἰούδα παῖδας τοῦ βασιλέως.
\vs{10}Καὶ Νάθαν τὸν προφήτην, καὶ Βαναίαν, καὶ τοὺς δυνατοὺς, καὶ τὸν Σαλωμὼν ἀδελφὸν αὐτοῦ, οὐκ ἐκάλεσε.

\vs{11}Καὶ εἶπε Νάθαν πρὸς Βηρσαβεὲ μητέρα Σαλωμὼν, λέγων, οὐκ ἤκουσας ὅτι ἐβασίλευσεν Ἀδωνίας υἱὸς Ἀγγὶθ, καὶ ὁ κύριος ἡμῶν Δαυὶδ οὐκ ἔγνω;
\vs{12}Καὶ νῦν δεῦρο, συμβουλεύσω σοι δὴ συμβουλίαν, καὶ ἐξελοῦ τὴν ψυχήν σου, καὶ τὴν ψυχὴν τοῦ υἱοῦ σου Σαλωμών.
\vs{13}Δεῦρο εἴσελθε πρὸς τὸν βασιλέα Δαυὶδ, καὶ ἐρεῖς πρὸς αὐτὸν, λέγουσα, οὐχὶ σὺ κύριέ μου βασιλεῦ ὤμοσας τῇ δούλῃ σου, λέγων, ὅτι ὁ υἱός σου Σαλωμὼν βασιλεύσει μετʼ ἐμέ, καὶ αὐτὸς καθιεῖται ἐπὶ τοῦ θρόνου μου; καὶ τί ὅτι ἐβασίλευσεν Ἀδωνίας;
\vs{14}Καὶ ἰδοὺ ἔτι λαλούσης σου ἐκεῖ μετα τοῦ βασιλέως, καὶ ἐγὼ εἰσελεύσομαι ὀπίσω σου, καὶ πληρώσω τοὺς λόγους σου.

\vs{15}Καὶ εἴσῆλθε Βηρσαβεὲ πρὸς τὸν βασιλέα εἰς τὸ ταμεῖον· καὶ ὁ βασιλεὺς πρεσβύτης σφόδρα· καὶ Ἀβισὰλ ἡ Σωμανῖτις ἦν λειτουργοῦσα τῷ βασιλεῖ.
\vs{16}Καὶ ἔκυψε Βηρσαβεὲ, καὶ προσεκύνησε τῷ βασιλεῖ· καὶ εἶπεν ὁ βασιλεύς, τί ἔστι σοί;
\vs{17}Ἡ δὲ εἶπε, κύριε, σὺ ὤμοσας ἐν Κυρίῳ τῷ Θεῷ σου τῇ δούλῃ σου, λέγων, ὅτι ὁ υἱός σου Σαλωμὼν βασιλεύσει μετʼ ἐμέ, καὶ καθήσεται ἐπὶ τοῦ θρόνου μου.
\vs{18}Καὶ νῦν ἰδοὺ Ἀδωνίας ἐβασίλευσε, καὶ σὺ κύριέ μου βασιλεῦ οὐκ ἔγνως.
\vs{19}Καὶ ἐθυσίασε μόσχους καὶ ἄρνας καὶ πρόβατα εἰς πλῆθος, καὶ ἐκάλεσε πάντας τοὺς υἱοὺς τοῦ βασιλέως, καὶ Ἀβιάθαρ τὸν ἱερέα, καὶ Ἰωὰβ τὸν ἄρχοντα τῆς δυνάμεως· καὶ τὸν Σαλωμὼν τὸν δοῦλόν σου οὐκ ἐκάλεσε.
\vs{20}Καὶ σὺ κύριέ μου βασιλεῦ, οἱ ὀφθαλμοὶ παντὸς Ἰσραὴλ πρὸς σέ, ἀπάγγειλαι αὐτοῖς τίς καθήσεται ἐπὶ τοῦ θρόνου τοῦ κυρίου μου τοῦ βασιλέως μετʼ αὐτόν.
\vs{21}Καὶ ἔσται ὡς ἂν κοιμηθῇ ὁ κύριός μου ὁ βασιλεὺς μετὰ τῶν πατέρων αὐτοῦ, καὶ ἔσομαι ἐγὼ καὶ Σαλωμὼν ὁ υἱός μου ἁμαρτωλοί.

\vs{22}Καὶ ἰδοὺ ἔτι αὐτῆς λαλούσης μετὰ τοῦ βασιλέως, καὶ Νάθαν ὁ προφήτης ἦλθε.
\vs{23}Καὶ ἀνηγγέλη τῷ βασιλεῖ, ἰδοὺ Νάθαν ὁ προφήτης· καὶ εἰσῆλθε κατὰ πρόσωπον τοῦ βασιλέως, καὶ προσεκύνησε τῷ βασιλεῖ κατὰ πρόσωπον αὐτοῦ ἐπὶ τὴν γῆν.
\vs{24}Καὶ εἶπε Νὰθαν, κύριέ μου βασιλεῦ, σὺ εἶπας, Ἀδωνίας βασιλεύσει ὀπίσω μου, καὶ αὐτὸς καθήσεται ἐπὶ τοῦ θρόνου μου;
\vs{25}Ὅτι κατέβη σήμερον, καὶ ἐθυσίασε μόσχους καὶ ἄρνας καὶ πρόβατα εἰς πλῆθος, καὶ ἐκάλεσε πάντας τοὺς υἱοὺς τοῦ βασιλέως, καὶ τοὺς ἄρχοντας τῆς δυνάμεως, καὶ Ἀβιάθαρ τὸν ἱερέα· καὶ ἰδού εἰσιν ἐσθίοντες καὶ πίνοντες ἐνώπιον αὐτοῦ, καὶ εἶπαν, ζήτω ὁ βασιλεὺς Ἀδωνίας.
\vs{26}Καὶ ἐμὲ αὐτὸν τὸν δοῦλόν σου, καὶ Σαδὼκ τὸν ἱερέα, καὶ Βαναίαν υἱὸν Ἰωδαὲ, καὶ Σαλωμὼν τὸν δοῦλόν σου, οὐκ ἐκάλεσεν.
\vs{27}Εἰ διὰ τοῦ κυρίου μου τοῦ βασιλέως γέγονε τὸ ῥῆμα τοῦτο, καὶ οὐκ ἐγνώρισας τῷ δούλῳ σου τίς καθήσεται ἐπὶ τὸν θρόνον τοῦ κυρίου μου βασιλέως μετʼ αὐτόν;

\vs{28}Καὶ ἀπεκρίθη ὁ βασιλεὺς Δαυὶδ, καὶ εἶπε, καλέσατέ μοι τὴν Βηρσαβεέ· καὶ εἰσῆλθεν ἐνώπιον τοῦ βασιλέως, καὶ ἔστη ἐνώπιον αὐτοῦ.
\vs{29}Καὶ ὤμοσεν ὁ βασιλεὺς, καὶ εἶπε, ζῇ Κύριος ὃς ἐλυτρώσατο τὴν ψυχήν μου ἐκ πάσης θλίψεως,
\vs{30}ὅτι καθὼς ὤμοσά σοι ἐν Κυρίῳ Θεῷ Ἰσραὴλ, λέγων, ὅτι Σαλωμὼν ὁ υἱός σου βασιλεύσει μετʼ ἐμὲ, καὶ αὐτὸς καθήσεται ἐπὶ τοῦ θρόνου μου ἀντʼ ἐμοῦ, ὅτι οὕτω ποιήσω τῇ ἡμέρᾳ ταύτῃ.
\vs{31}Καὶ ἔκυψε Βηρσαβεὲ ἐπὶ πρόσωπον ἐπὶ τὴν γῆν, καὶ προσεκύνησε τῷ βασιλεῖ, καὶ εἶπε, ζήτω ὁ κύριός μου ὁ βασιλεὺς Δαυὶδ εἰς τὸν αἰῶνα.

\vs{32}Καὶ εἶπεν ὁ βασιλεὺς Δαυίδ, Καλέσατέ μοι Σαδὼκ τὸν ἱερέα, καὶ Νάθαν τὸν προφήτην, καὶ Βαναίαν υἱὸν Ἰωδαέ καὶ εἰσῆλθον ἐνώπιον τοῦ βασιλέως.
\vs{33}Καὶ εἶπεν ὁ βασιλεὺς αὐτοῖς, λάβετε τοὺς δούλους τοῦ κυρίου ὑμῶν μεθʼ ὑμῶν, καὶ ἐπιβιβάσατε τὸν υἱόν μου Σαλωμὼν ἐπὶ τὴν ἡμίονον τὴν ἐμὴν, καὶ καταγάγετε αὐτὸν εἰς τὴν Γιὼν,
\vs{34}καὶ χρισάτω αὐτὸν ἐκεῖ Σαδὼκ ὁ ἱερεὺς καὶ Νάθαν ὁ προφήτης εἰς βασιλέα ἐπὶ Ἰσραὴλ, καὶ σαλπίσατε κερατίνῃ, καὶ ἐρεῖτε, ζήτω ὁ βασιλεὺς Σαλωμών.
\vs{35}Καὶ καθήσεται ἐπὶ τοῦ θρόνου μου, καὶ βασιλεύσει ἀντʼ ἐμοῦ· καὶ ἐγὼ ἐνετειλάμην τοῦ εἶναι εἰς ἡγούμενον ἐπὶ Ἰσραὴλ καὶ Ἰούδαν.
\vs{36}Καὶ ἀπεκρίθη Βαναίας υἱὸς Ἰωδαὲ τῷ βασιλεῖ, καὶ εἶπε, γένοιτο οὕτως· πιστώσαι Κύριος ὁ Θεὸς τοῦ κυρίου μου τοῦ βασιλέως·
\vs{37}καθὼς ἦν Κύριος μετὰ τοῦ κυρίου μου τοῦ βασιλέως, οὕτως εἴη μετὰ Σαλωμὼν, καὶ μεγαλύναι τὸν θρόνον αὐτοῦ ὑπὲρ τὸν θρόνον τοῦ κυρίου μου τοῦ βασιλέως Δαυίδ.

\vs{38}Καὶ κατέβη Σαδὼκ ὁ ἱερεὺς, καὶ Νάθαν ὁ προφήτης, καὶ Βαναίας υἱὸς Ἰωδαὲ, καὶ ὁ Χερεθὶ, καὶ ὁ Φελεθὶ, καὶ ἐπεκάθισαν τὸν Σαλωμὼν ἐπὶ τὴν ἡμίονον τοῦ βασιλέως Δαυὶδ, καὶ ἀπήγαγον αὐτὸν εἰς τὴν Γιών.
\vs{39}Καὶ ἔλαβε Σαδὼκ ὁ ἱερεὺς τὸ κέρας τοῦ ἐλαίου ἐκ τῆς σκηνῆς, καὶ ἔχρισε τὸν Σαλωμὼν, καὶ ἐσάλπισε τῇ κερατίνῃ· καὶ εἶπε πᾶς ὁ λαός, ζήτω ὁ βασιλεὺς Σαλωμών.
\vs{40}Καὶ ἀνέβη πᾶς ὁ λαὸς ὀπίσω αὐτοῦ, καὶ ἐχόρευον ἐν χοροῖς καὶ εὐφραινόμενοι εὐφροσύνην μεγάλην, καὶ ἐῤῥάγη ἡ γῆ ἐν τῇ φωνῇ αὐτῶν.

\vs{41}Καὶ ἤκουσεν Ἀδωνίας καὶ πάντες οἱ κλητοὶ αὐτοῦ, καὶ αὐτοὶ συνετέλεσαν φαγεῖν· καὶ ἤκουσεν Ἰωὰβ τὴν φωνὴν τῆς κερατίνης, καὶ εἶπε, τίς ἡ φωνὴ τῆς πόλεως ἠχούσης;
\vs{42}Ἔτι αὐτοῦ λαλοῦντος, καὶ ἰδοὺ Ἰωνάθαν υἱὸς Ἀβιάθαρ τοῦ ἱερέως εἰσῆλθε· καὶ εἶπεν Ἀδωνίας, εἴσελθε, ὅτι ἀνὴρ δυνάμεως εἶ σὺ, καὶ ἀγαθὰ εὐαγγέλισαι.
\vs{43}Καὶ ἀπεκρίθη Ἰωνάθαν, καὶ εἶπε, καὶ μάλα ὁ κύριος ἡμῶν ὁ βασιλεὺς Δαυὶδ ἐβασίλευσε τὸν Σαλωμὼν,
\vs{44}καὶ ἀπέστειλε μετʼ αὐτοῦ ὁ βασιλεὺς τὸν Σαδὼκ τὸν ἱερέα, καὶ Νάθαν τὸν προφήτην, καὶ Βαναίαν τὸν υἱὸν Ἰωδαὲ, καὶ τὸν Χερεθὶ, καὶ τὸν Φελεθὶ, καὶ ἐπεκάθισαν αὐτὸν ἐπὶ τὴν ἡμίονον τοῦ βασιλέως·
\vs{45}Καὶ ἔχρισαν αὐτὸν Σαδὼκ ὁ ἱερεὺς καὶ Νάθαν ὁ προφήτης ἐν τῇ Γιὼν, καὶ ἀνέβησαν ἐκεῖθεν εὐφραινόμενοι, καὶ ἤχησεν ἡ πόλις· αὕτη ἡ φωνὴ ἣν ἠκούσατε.
\vs{46}Καὶ ἐκάθισε Σαλωμὼν ἐπὶ θρόνον βασιλείας.
\vs{47}Καὶ εἰσῆλθον οἱ δοῦλοι τοῦ βασιλέως εὐλογῆσαι τὸν κύριον ἡμῶν τὸν βασιλέα Δαυὶδ, λέγοντες, ἀγαθύναι ὁ Θεὸς τὸ ὄνομα Σαλωμὼν ὑπὲρ τὸ ὄνομά σου, καὶ μεγαλύναι τὸν θρόνον αὐτοῦ ὑπὲρ τὸν θρόνον σου· καὶ προσεκύνησεν ὁ βασιλεὺς ἐπὶ τὴν κοίτην.
\vs{48}Καί γε οὕτως εἶπεν ὁ βασιλεύς, εὐλογητὸς Κύριος ὁ Θεὸς Ἰσραὴλ, ὃς ἔδωκε σήμερον ἐκ τοῦ σπέρματός μου καθήμενον ἐπὶ τοῦ θρόνου μου, καὶ οἱ ὀφθαλμοί μου βλέπουσι.

\vs{49}Καὶ ἐξέστησαν πάντες οἱ κλητοὶ τοῦ Ἀδωνίου, καὶ ἦλθον ἀνὴρ εἰς τὴν ὁδὸν αὐτοῦ.
\vs{50}Καὶ Ἀδωνίας ἐφοβήθη ἀπὸ προσώπου Σαλωμὼν, καὶ ἀνέστη καὶ ἀπῆλθε καὶ ἐπελάβετο τῶν κεράτων τοῦ θυσιαστηρίου.
\vs{51}Καὶ ἀνηγγέλη τῷ Σαλωμὼν, λέγοντες, ἰδοὺ Ἀδωνίας ἐφοβήθη τὸν βασιλέα Σαλωμὼν, καὶ κατέχει τῶν κεράτων τοῦ θυσιαστηρίου, λέγων, ὀμοσάτω μοι σήμερον Σαλωμὼν, εἰ οὐ θανατώσει τὸν δοῦλον αὐτοῦ ἐν ῥομφαίᾳ.
\vs{52}Καὶ εἶπε Σαλωμών, ἐὰν γένηται εἰς υἱὸν δυνάμεως, εἰ πεσεῖται τῶν τριχῶν αὐτοῦ ἐπὶ τὴν γῆν· καὶ ἐὰν κακία εὑρεθῇ ἐν αὐτῷ, θανατωθήσεται.
\vs{53}Καὶ ἀπέστειλεν ὁ βασιλεὺς Σαλωμὼν, καὶ κατήνεγκαν αὐτὸν ἀπάνωθεν τοῦ θυσιαστηρίου· καὶ εἰσῆλθε, καὶ προσεκύνησε τῷ βασιλεῖ Σαλωμών· καὶ εἶπεν αὐτῷ Σαλωμὼν, δεῦρο εἰς τὸν οἶκόν σου.

\ch{2}
Καὶ ἤγγισαν αἱ ἡμέραι Δαυὶδ ἀποθανεῖν αὐτὸν, καὶ ἀπεκρίνατο Σαλωμὼν υἱῷ αὐτοῦ, λέγων,
\vs{2}ἐγώ εἰμι πορεύομαι ἐν ὁδῷ πάσης τῆς γῆς· καὶ ἰσχύσεις, καὶ ἔσῃ εἰς ἄνδρα,
\vs{3}καὶ φυλάξεις φυλακὴν Κυρίου Θεοῦ σου τοῦ πορεύεσθαι ἐν ταῖς ὁδοῖς αὐτοῦ, φυλάσσειν τὰς ἐντολὰς αὐτοῦ καὶ τὰ δικαιώματα καὶ τὰ κρίματα τὰ γεγραμμένα ἐν τῷ νόμῳ Μωυσέως· ἵνα συνήσῃς ἃ ποιήσεις κατὰ πάντα ὅσα ἂν ἐντείλωμαί σοι·
\vs{4}Ἵνα στήσῃ Κύριος τὸν λόγον αὐτοῦ ὃν ἐλάλησε, λέγων, ἐὰν φυλάξωσιν οἱ υἱοί σου τὴν ὁδὸν αὐτῶν πορεύεσθαι ἐνώπιόν μου ἐν ἀληθείᾳ, ἐν ὅλῃ καρδίᾳ αὐτῶν, λέγων, οὐκ ἐξολοθρευθήσεταί σοι ἀνὴρ ἐπάνωθεν θρόνου Ἰσραήλ.
\vs{5}Καί γε σὺ ἔγνως ὅσα ἐποίησέ μοι Ἰωὰβ υἱὸς Σαρουίας, ὅσα ἐποίησε τοῖς δυσὶν ἄρχουσι τῶν δυνάμεων Ἰσραὴλ, τῷ Ἀβεννὴρ υἱῷ Νὴρ, καὶ τῷ Ἀμεσσαῒ υἱῷ Ἰεθὲρ, καὶ ἀπέκτεινεν αὐτοὺς, καὶ ἔταξε τὰ αἵματα πολέμου ἐν εἰρήνῃ, καὶ ἔδωκεν αἷμα ἀθῶον ἐν τῇ ζώνῃ αὐτοῦ τῇ ἐν τῇ ὀσφύϊ αὐτοῦ, καὶ ἐν τῷ ὑποδήματι αὐτοῦ τῷ ἐν τῷ ποδὶ αὐτοῦ.
\vs{6}Καὶ ποιήσεις κατὰ τὴν σοφίαν σου, καὶ σὺ κατάξεις τὴν πολιὰν αὐτοῦ ἐν εἰρήνῃ εἰς ᾅδου.
\vs{7}Καὶ τοῖς υἱοῖς Βερζελλὶ τοῦ Γαλααδίτου ποιήσεις ἔλεος, καὶ ἔσονται ἐν τοῖς ἐσθίουσιν τὴν τράπεζάν σου· ὅτι οὕτως ἤγγισάν μοι ἐν τῷ με ἀποδιδράσκειν ἀπὸ προσώπου Ἀβεσσαλὼμ τοῦ ἀδελφοῦ σου.
\vs{8}Καὶ ἰδοὺ μετὰ σοῦ Σεμεῒ υἱὸς Γηρὰ υἱὸς τοῦ Ἰεμινὶ ἐκ Βαουρὶμ, καὶ αὐτὸς κατηράσατό με κατάραν ὀδυνηρὰν τῇ ἡμέρᾳ ᾗ ἐπορευόμην εἰς παρεμβολάς· καὶ αὐτὸς κατέβη εἰς ἀπαντήν μου εἰς τὸν Ἰορδάνην, καὶ ὤμοσα αὐτῷ ἐν Κυρίῳ, λέγων, εἰ θανατώσω σε ἐν ῥομφαίᾳ.
\vs{9}Καὶ οὐ μὴ ἀθωώσῃς αὐτὸν, ὅτι ἀνὴρ σοφὸς εἶ σύ, καὶ γνώσῃ ἃ ποιήσεις αὐτῷ, καὶ κατάξεις τὴν πολιὰν αὐτοῦ ἐν αἵματι εἰς ᾅδου.

\vs{10}Καὶ ἐκοιμήθη Δαυὶδ μετὰ τῶν πατέρων αὐτοῦ, καὶ ἐτάφη ἐν πόλει Δαυίδ.
\vs{11}Καὶ αἱ ἡμέραι ἃς ἐβασίλευσε Δαυὶδ ἐπὶ τὸν Ἰσραὴλ, τεσσαράκοντα ἔτη· ἐν Χεβρὼν ἐβασίλευσεν ἑπτά ἔτη, καὶ ἐν Ἱερουσαλὴμ τριάκοντα τρία ἔτη.

\vs{12}Καὶ Σαλωμὼν ἐκάθισεν ἐπὶ θρόνου Δαυὶδ τοῦ πατρὸς αὐτοῦ, καὶ ἡτοιμάσθη ἡ βασιλεία αὐτοῦ σφόδρα.
\vs{13}Καὶ εἰσῆλθεν Ἀδωνίας υἱὸς Ἀγγὶθ πρὸς Βηρσαβεὲ μητέρα Σαλωμὼν, καὶ προσεκύνησεν αὐτῇ· ἡ δὲ εἶπεν, εἰρήνη ἡ εἴσοδός σου; καὶ εἶπεν, εἰρήνη·
\vs{14}λόγος μοι πρὸς σέ. Καὶ εἶπεν αὐτῷ, λάλησον.
\vs{15}Καὶ εἶπεν αὐτῇ, σὺ οἶδας, ὅτι ἐμοὶ ἦν βασιλεία, καὶ ἐπʼ ἐμὲ ἔθετο πᾶς Ἰσραὴλ τὸ πρόσωπον αὐτοῦ εἰς βασιλέα· καὶ ἐστράφη ἡ βασιλεία, καὶ ἐγένετο τῷ ἀδελφῷ μου, ὅτι παρὰ Κυρίου ἐγενήθη αὐτῷ.
\vs{16}Καὶ νῦν αἴτησιν μίαν ἐγὼ αἰτοῦμαι παρὰ σοῦ, μὴ ἀποστρέψῃς τὸ πρόσωπόν σου· καὶ εἶπεν αὐτῷ Βηρσαβεὲ, λάλει.
\vs{17}Καὶ εἶπεν αὐτῇ, εἰπον δὴ πρὸς Σαλωμὼν τὸν βασιλέα, ὅτι οὐκ ἀποστρέψει τὸ πρόσωπον αὐτοῦ ἀπὸ σοῦ, καὶ δώσει μοι τὴν Ἀβισὰγ τὴν Σωμανίτιν εἰς γυναῖκα.
\vs{18}Καὶ εἶπε Βηρσαβεὲ, καλῶς· ἐγὼ λαλήσω περὶ σοῦ τῷ βασιλεῖ.

\vs{19}Καὶ εἰσῆλθε Βηρσαβεὲ πρὸς τὸν βασιλέα Σαλωμὼν αλῆσαι αὐτῷ περὶ Ἀδωνίου· καὶ ἐξανέστη ὁ βασιλεὺς εἰς ἀπαντὴν αὐτῇ, καὶ κατεφίλησεν αὐτὴν, καὶ ἐκάθισεν ἐπὶ τοῦ θρόνου· καὶ ἐτέθη θρόνος τῇ μητρὶ τοῦ βασιλέως, καὶ ἐκάθισεν ἐκ δεξιῶν αὐτοῦ.
\vs{20}Καὶ εἶπεν αὐτῷ, αἴτησιν μίαν μικρὰν ἐγὼ αἰτοῦμαι παρὰ σοῦ, μὴ ἀποστρέψῃς τὸ πρόσωπόν μου· καὶ εἶπεν αὐτῇ ὁ βασιλεύς, αἴτησαι, μήτερ ἐμὴ, καὶ οὐκ ἀποστρέψω σε.
\vs{21}Καὶ εἶπε, δοθήτω δὴ Ἀβισὰ ἡ Σωμανίτις τῷ Αδωνίᾳ τῷ ἀδελφῷ σου εἰς γυναῖκα.
\vs{22}Καὶ ἀπεκρίθη ὁ βασιλεὺς Σαλωμὼν, καὶ εἶπε τῇ μητρὶ αὐτοῦ, καὶ ἱνατί σὺ ᾔτησαι τὴν Ἀβισὰλ τῷ Ἀδωνίᾳ; καὶ αἴτησαι αὐτῷ τὴν βασιλείαν, ὅτι οὗτος ἀδελφός μου ὁ μέγας ὑπὲρ ἐμέ, καὶ αὐτῷ Ἀβιάθαρ ὁ ἱερεὺς, καὶ αὐτῷ Ἰωὰβ ὁ υἱὸς Σαρουίας ἀρχιστράτηγος ἑταῖρος.
\vs{23}Καὶ ὤμοσεν ὁ βασιλεὺς Σαλωμὼν κατὰ τοῦ Κυρίου, λέγων, τάδε ποιήσαι μοι ὁ Θεὸς καὶ τάδε παροσθείη, ὅτι κατὰ τῆς ψυχῆς αὐτοῦ ἐλάλησεν Ἀδωνίας τὸν λόγον τοῦτον.
\vs{24}Καὶ νῦν ζῇ Κύριος ὃς ἡτοίμασε με καὶ ἔθετό με ἐπὶ τὸν θρόνον Δαυὶδ τοῦ πατρός μου, καὶ αὐτὸς ἐποίησέ μοι οἶκον καθὼς ἐλάλησε Κύριος, ὅτι σήμερον θανατωθήσεται Ἀδωνίας.
\vs{25}Καὶ ἐξαπέστειλεν ὁ βασιλεὺς Σαλωμὼν ἐν χειρὶ Βαναίου υἱοῦ Ἰωδαὲ, καὶ ἀνεῖλεν αὐτὸν, καὶ ἀπέθανεν Ἀδωνίας ἐν τῇ ἡμέρᾳ ἐκείνῃ.

\vs{26}Καὶ τῷ Ἀβιάθαρ τῷ ἱερεῖ εἶπεν ὁ βασιλεύς, ἀπότρεχε σὺ εἰς Ἀναθὼθ εἰς ἀγρόν σου, ὅτι ἀνὴρ θανάτου εἶ σὺ ἐν τῇ ἡμέρᾳ ταύτῃ· καὶ οὐ θανατώσω σε, ὅτι ᾖρας τὴν κιβωτὸν τῆς διαθήκης Κυρίου ἐνώπιον τοῦ πατρός μου, καὶ ὅτι ἐκακουχήθης ἐν πᾶσιν οἷς ἐκακουχήθη ὁ πατήρ μου.
\vs{27}Καὶ ἐξέβαλε Σαλωμὼν τὸν Ἀβιάθαρ τοῦ μὴ εἶναι ἱερέα τοῦ Κυρίου, πληρωθῆναι τὸ ῥῆμα Κυρίου, ὃ ἐλάλησεν ἐπὶ τὸν οἶκον Ἡλὶ ἐν Σηλώμ.

\vs{28}Καὶ ἡ ἀκοὴ ἦλθεν ἕως Ἰωὰβ υἱοῦ Σαρουίας, ὅτι Ἰωὰβ ἦν κεκλικὼς ὀπίσω Ἀδωνίου, καὶ ὀπίσω Σαλωμὼν οὐκ ἔκλινε· καὶ ἔφυγεν Ἰωὰβ εἰς τὸ σκήνωμα τοῦ Κυρίου, καὶ κατέσχε τῶν κεράτων τοῦ θυσιαστηρίου.
\vs{29}Καὶ ἀπηγγέλη τῷ Σαλωμὼν, λέγοντες, ὅτι πέφευγεν Ἰωὰβ εἰς τὴν σκηνὴν τοῦ Κυρίου, καὶ ἰδοὺ κατέχει τῶν κεράτων τοῦ θυσιαστηρίου· καὶ ἀπέστειλε Σαλωμὼν ὁ βασιλεὺς πρὸς Ἰωὰβ, λέγων, τί γέγονέ σοι, ὅτι πέφευγας εἰς τὸ θυσιαστήριον; καὶ εἶπεν Ἰωὰβ, ὅτι ἐφοβήθην ἀπὸ προσώπου σου, καὶ ἔφυγον πρὸς Κύριον· καὶ ἀπέστειλε Σαλωμὼν τὸν Βαναίου υἱὸν Ἰωδαὲ, λέγων, πορεύου καὶ ἄνελε αὐτὸν, καὶ θάψον αὐτόν.

\vs{30}Καὶ ἦλθε Βαναίας υἱὸς Ἰωδαὲ πρὸς Ἰωὰβ εἰς τὴν σκηνὴν τοῦ Κυρίου, καὶ εἶπεν αὐτῷ τάδε λέγει ὁ βασιλεὺς, ἔξελθε· καὶ εἶπεν Ἰωὰβ, οὐκ ἐκπορεύωμαι, ὅτι ὧδε ἀποθανοῦμαι· καὶ ἐπέστρεψε Βαναίας υἱὸς Ἰωδαὲ, καὶ εἶπε τῷ βασιλεῖ, λέγων, τάδε λελάληκεν Ἰωὰβ, καὶ τάδε ἀποκέκριταί μοι.
\vs{31}Καὶ εἶπεν αὐτῷ ὁ βασιλεύς, πορεύου, καὶ ποίησον αὐτῷ καθὼς εἴρηκε, καὶ ἄνελε αὐτὸν· καὶ θάψεις αὐτὸν, καὶ ἐξαρεῖς σήμερον τὸ αἷμα ὃ δωρεὰν ἐξέχεεν, ἀπʼ ἐμοῦ καὶ ἀπὸ τοῦ οἴκου τοῦ πατρός μου.
\vs{32}Καὶ ἐπέστρεψε Κύριος τὸ αἷμα τῆς ἀδικίας αὐτοῦ εἰς κεφαλὴν αὐτοῦ, ὡς ἀπήντησε τοῖς δυσὶν ἀνθρώποις τοῖς δικαίοις καὶ ἀγαθοῖς ὑπὲρ αὐτὸν, καὶ ἀπέκτεινεν αὐτοὺς ἐν ῥομφαίᾳ, καὶ ὁ πατήρ μου Δαυὶδ οὐκ ἔγνω τὸ αἷμα αὐτῶν, τὸν Ἀβεννὴρ υἱὸν Νὴρ ἀρχιστράτηγον Ἰσραὴλ, καὶ τὸν Ἀμεσσὰ υἱὸν Ἰεθὲρ ἀρχιστράτηγον Ἰούδα.
\vs{33}Καὶ ἐπεστράφη τὰ αἵματα αὐτῶν εἰς κεφαλὴν αὐτοῦ, καὶ εἰς κεφαλὴν τοῦ σπέρματος αὐτοῦ εἰς τὸν αἰῶνα· καὶ τῷ Δαυὶδ καὶ τῷ σπέρματι αὐτοῦ καὶ τῷ οἴκῳ αὐτοῦ καὶ τῷ θρόνῳ αὐτοῦ γένοιτο εἰρήνη ἕως αἰῶνος παρὰ Κυρίου.
\vs{34}Καὶ ἀνέβη Βαναίας υἱὸς Ἰωδαὲ, καὶ ἀπήντησεν αὐτῷ, καὶ ἐθανάτωσεν αὐτὸν, καὶ ἔθαψεν αὐτὸν ἐν τῷ οἴκῳ αὐτοῦ ἐν τῇ ἐρήμῳ.

\vs{35}Καὶ ἔδωκεν ὁ βασιλεὺς τὸν Βαναίου υἱὸν Ἰωδαὲ ἀντʼ αὐτοῦ ἐπὶ τὴν στρατηγίαν· καὶ ἡ βασιλεία κατωρθοῦτο ἐν Ἰερουσαλήμ· καὶ Σαδὼκ τὸν ἱερέα ἔδωκεν αὐτὸν ὁ βασιλεὺς εἰς ἱερέα πρῶτον ἀντὶ Ἀβιάθαρ. Καὶ Σαλωμὼς υἱὸς Δαυὶδ ἐβασίλευσεν ἐπὶ Ἰσραὴλ καὶ Ἰούδα ἐν Ἱερουσαλήμ·
\vs{35a}καὶ ἔδωκε Κύριος φρόνησιν τῷ Σαλωμὼν, καὶ σοφίαν πολλὴν σφόδρα, καὶ πλάτος καρδίας, ὡς ἡ ἄμμος ἡ παρὰ τὴν θάλασσαν.

\vs{35b}Καὶ ἐπληθύνθη ἡ φρόνησις Σαλωμὼν σφόδρα ὑπὲρ τὴν φρόνησιν πάντων υἱῶν ἀρχαίων, καὶ ὑπὲρ πάντας φρονίμους Αἰγύπτου·
\vs{35c}Καὶ ἔλαβε τὴν θυγατέρα Φαραὼ, καὶ εἰσήγαγεν αὐτὴν εἰς πόλιν Δαυὶδ ἕως συντελέσαι αὐτὸν οἰκοδομῆσαι τὸν οἶκον αὐτοῦ, καὶ τὸν οἶκον Κυρίου ἐν πρώτοις, καὶ τὸ τεῖχος Ἱερουσαλὴμ κυκλόθεν· ἐν ἑπτὰ ἔτεσιν ἐποίησε καὶ συνετέλεσε.

\vs{35d}Καὶ ἦν τῷ Σαλωμὼν ἐβδομήκοντα χιλιάδες αἴροντες ἄρσιν, καὶ ὀγδοήκοντα χιλιάδες λατόμων ἐν τῷ ὄρει·
\vs{35e}καὶ ἐποίησε Σαλωμὼν τὴν θάλασσαν, καὶ τὰ ὑποστηρίγματα, καὶ τοὺς λουτῆρας τοὺς μεγάλους, καὶ τοὺς στύλους, καὶ τἠν κρήνην τῆς αὐλῆς, καὶ τὴν θάλασσαν τὴν χαλκῆν· καὶ ᾠκοδόμησε τὴν ἄκραν ἔπαλξιν ἐπʼ αὐτῆς, διεκόψεν τὴν πόλιν Δαυίδ.
\vs{35f}οὕτως θυγάτηρ Φαραὼ ἀνέβαινεν ἐκ τῆς πόλεως Δαυὶδ εἰς τὸν οἶκον αὐτῆς, ὃν ᾠκοδόμησεν αὐτῇ· τότε ᾠκοδόμησε τὴν ἄκραν·
\vs{35g}καὶ Σαλωμὼν ἀνέφερε τρεῖς ἐν τῷ ἐνιαυτῷ ὁλοκαυτώσεις καὶ εἰρηνικὰς ἐπὶ τὸ θυσιαστήριον ὃ ᾠκοδόμησε τῷ Κυρίῳ, καὶ ἐθυμία ἐνώπιον Κυρίου, καὶ συνετέλεσε τὸν οἶκον.
\vs{35h}Καὶ οὗτοι οἱ ἄρχοντες οἱ καθεσταμένοι ἐπὶ τὰ ἔργα τοῦ Σαλωμὼν, τρεῖς χιλιάδες καὶ ἑξακόσιοι ἐπιστάται τοῦ λαοῦ τῶν ποιούντων τὰ ἔργα·
\vs{35i}καὶ ᾠκοδόμησε τὴν Ἀσσοὺρ, καὶ τὴν Μαγδὼ, καὶ τὴν Γαζὲρ, καὶ τὴν Βαιθωρὼν ἐπάνω, καὶ τὰ Βαλλάθ·
\vs{35k}πλὴν μετὰ τὸ οἰκοδομῆσαι αὐτὸν τὸν οἶκον τοῦ Κυρίου, καὶ τὸ τεῖχος Ἱερουσαλὴμ κύκλῳ, μετὰ ταῦτα ᾠκοδόμησε τὰς πόλεις ταύτας.

\vs{35l}Καὶ ἐν τῷ ἔτι Δαυὶδ ζῇν, ἐνετείλατο τῷ Σαλωμὼν, λέγων, ἰδοὺ μετὰ σοῦ Σεμεῒ υἱὸς Γηρὰ υἱὸς τοῦ σπέρματος τοῦ Ἰεμινὶ ἐκ Χεβρών·
\vs{35m}οὗτος κατηράσατό με κατάραν ὀδυνηρὰν ἐν ᾗ ἡμέρᾳ ἐπορευόμην εἰς παρεμβολάς·
\vs{35n}καὶ αὐτὸς κατέβαινεν εἰς ἀπαντήν μοι ἐπὶ τὸν Ἰορδάνην, καὶ ὤμοσα αὐτῷ κατὰ τοῦ Κυρίου, λέγων, εἰ θανατωθήσεται ἐν ῥομφαίᾳ·
\vs{35o}καὶ νῦν μὴ ἀθωώσῃς αὐτὸν, ὅτι ἀνὴρ φρόνιμος σύ· καὶ γνώσῃ ἃ ποιήσεις αὐτῷ, καὶ κατάξεις τὴν πολιὰν αὐτοῦ ἐν αἵματι εἰς ᾅδου.

\vs{36}Καὶ ἐκάλεσεν ὁ βασιλεὺς τὸν Σεμεῒ, καὶ εἶπεν αὐτῷ, ὠκοδόμησον σεαυτῷ οἶκον ἐν Ἱερουσαλὴμ καὶ κάθου ἐκεῖ, καὶ οὐκ ἐξελεύσῃ ἐκεῖθεν οὐδαμοῦ.
\vs{37}Καὶ ἔσται ἐν τῇ ἡμέρᾳ τῆς ἐξόδου σου καὶ διαβήσῃ τὸν χείμαῤῥον Κέδρων, γινώσκων γνώσῃ ὅτι θανάτῳ ἀποθανῇ· τὸ αἷμά σου ἔσται ἐπὶ τὴν κεφαλήν σου· καὶ ὥρκισεν αὐτὸν ὁ βασιλεὺς ἐν τῇ ἡμέρᾳ ἐκείνῃ.
\vs{38}Καὶ εἶπε Σεμεῒ πρὸς τὸν βασιλέα, ἀγαθὸν τὸ ῥῆμα ὃ ἐλάλησας, κύριέ μου βασιλεῦ· οὕτω ποιήσει ὁ δοῦλός σου· καὶ ἐκάθισε Σεμεῒ ἐν Ἱερουσαλὴμ τρία ἔτη.

\vs{39}Καὶ ἐγενήθη μετὰ τὰ τρία ἔτη, καὶ ἀπέδρασαν δύο δοῦλοι τοῦ Σεμεῒ πρὸς Ἀγχοὺς υἱὸν Μααχὰ βασιλέα Γέθ· καὶ ἀπηγγέλη τῷ Σεμεῒ, λέγοντες, ἰδοὺ οἱ δοῦλοί σου ἐν Γέθ.
\vs{40}Καὶ ἀνέστη Σεμεῒ καὶ ἐπέσαξε τὴν ὄνον αὐτοῦ, καὶ ἐπορεύθη εἰς Γὲθ πρὸς Ἀγχοὺς τοῦ ἐκζητῆσαι τοὺς δούλους αὐτοῦ· καὶ ἐπορεύθη Σεμεῒ, καὶ ἤγαγε τοὺς δούλους αὐτοῦ ἐκ Γέθ.
\vs{41}Καὶ ἀπηγγέλη τῷ Σαλωμὼν, λέγοντες, ὅτι ἐπορεύθη Σεμεῒ ἐξ Ἱερουσαλὴμ εἰς Γὲθ, καὶ ἀνέστρεψε τοὺς δούλους αὐτοῦ.
\vs{42}Καὶ ἀπέστειλεν ὁ βασιλεὺς καὶ ἐκάλεσε τὸν Σεμεῒ, καὶ εἶπε πρὸς αὐτόν, οὐχὶ ὥρκισά σε κατὰ τοῦ Κυρίου, καὶ ἐπεμαρτυράμην σοι, λέγων, ἐν ᾗ ἂν ἡμέρᾳ ἐξέλθῃς ἐξ Ἱερουσαλὴμ, καὶ πορευθῇς εἰς δεξιὰ ἢ εἰς ἀριστερά, γινώσκων γνώσῃ ὅτι θανάτῳ ἀποθανῇ;
\vs{43}Καὶ τί ὅτι οὐκ ἐφύλαξας τὸν ὅρκον Κυρίου, καὶ τὴν ἐντολὴν ἣν ἐνετειλάμην κατὰ σοῦ;

\vs{44}Καὶ εἶπεν ὁ βασιλεὺς πρὸς Σεμεῒ, σὺ οἶδας πᾶσαν τὴν κακίαν σου ἣν οἶδεν ἡ καρδία σου, ἃ ἐποίησας Δαυὶδ τῷ πατρί μου, καὶ ἀνταπέδωκε Κύριος τὴν κακίαν σου εἰς κεφαλήν σου.
\vs{45}Καὶ ὁ βασιλεὺς Σαλωμὼν εὐλογημένος, καὶ ὁ θρόνος Δαυὶδ ἔσται ἕτοιμος ἐνώπιον Κυρίου εἰς τὸν αἰῶνα.
\vs{46}Καὶ ἐνετείλατο ὁ βασιλεὺς Σαλωμὼν τῷ Βαναίᾳ υἱῷ Ἰωδαὲ, καὶ ἐξῆλθε καὶ ἀνεῖλεν αὐτόν.

\vs{46a}Καὶ ἦν ὁ βασιλεὺς Σαλωμὼν φρόνιμος σφόδρα καὶ σοφός· καὶ Ἰούδα καὶ Ἰσραὴλ πολλοὶ σφόδρα, ὡς ἡ ἄμμος ἡ ἐπὶ τῆς θαλάσσης εἰς πλῆθος, ἐσθίοντες καὶ πίνοντες καὶ χαίροντες·
\vs{46b}καὶ Σαλωμὼν ἦν ἄρχων ἐν πάσαις ταῖς βασιλείαις· καὶ ἦσαν προσφέροντες δῶρα, καὶ ἐδούλευον τῷ Σαλωμὼν πάσας τὰς ἡμέρας τῆς ζωῆς αὐτοῦ·
\vs{46c}καὶ Σαλωμὼν ἤρξατο ἀνοίγειν τὰ δυναστεύματα τοῦ Λιβάνου·
\vs{46d}καὶ αὐτὸς ᾠκοδόμησε τὴν Θερμαὶ ἐν τῇ ἐρήμῳ·
\vs{46e}καὶ τοῦτο τὸ ἄριστον τῷ Σαλωμών· τριάκοντα κόροι σεμιδάλεως, καὶ ἑξήκοντα κόροι ἀλεύρου κεκοπανισμένου, δέκα μόσχοι ἐκλεκτοὶ, καὶ εἴκοσι βόες νομάδες, καὶ ἑκατὸν πρόβατα, ἐκτὸς ἐλάφων καὶ δορκάδων καὶ ὀρνίθων ἐκλεκτῶν νομάδων·
\vs{46f}ὅτι ἦν ἄρχων ἐν παντὶ πέραν τοῦ ποταμοῦ ἀπὸ Ῥαφὶ ἕως Γάζης ἐν πᾶσι τοῖς βασιλεῦσι πέραν τοῦ ποταμοῦ·;
\vs{46g}καὶ ἦν αὐτῷ εἰρήνη ἐκ πάντων τῶν μερῶν αὐτοῦ κυκλόθεν· καὶ κατῴκει Ἰούδα καὶ Ἰσραὴλ πεποιθότες, ἕκαστος ὑπὸ τὴν ἄμπελον αὐτοῦ, καὶ ὑπὸ τὴν συκῆν αὐτοῦ, ἐσθίοντες καὶ πίνοντες καὶ ἑορτάζοντες ἀπὸ Δὰν καὶ ἕως Βηρσαβεὲ πάσας τὰς ἡμέρας Σαλωμών.

\vs{46h}Καὶ οὗτοι οἱ ἄρχοντες τοῦ Σαλωμών. Ἀζαρίου υἱὸς Σαδὼκ τοῦ ἱερέως, καὶ Ὀρνίου υἱὸς Νάθαν ἄρχων τῶν ἐφεστηκότων· καὶ ἔδραμεν ἐπὶ τὸν οἶκον αὐτοῦ· καὶ Σουβὰ γραμματεὺς, καὶ Βασὰ υἱὸς Ἀχιθαλὰμ ἀναμιμνήσκων, καὶ Ἀβὶ υἱὸς Ἰωὰβ ἀρχιστράτηγος, καὶ Ἀχιρὲ υἱὸς Ἐδραῒ ἐπὶ τὰς ἄρσεις, καὶ Βαναίας υἱὸς Ἰωδαὲ ἐπὶ τῆς αὐλαρχίας καὶ ἐπὶ τοῦ πλινθίου, καὶ Καχοὺρ υἱὸς Νάθαν ὁ σύμβουλος.

\vs{46i}Καὶ ἦσαν τῷ Σαλωμὼν τεσσαράκοντα χιλιάδες τοκάδες ἵπποι εἰς ἅρματα, καὶ δώδεκα χιλιάδες ἵππων·
\vs{46k}καὶ ἦν ἄρχων ἐν πᾶσι τοῖς βασιλεῦσιν ἀπὸ τοῦ ποταμοῦ καὶ ἕως γῆς ἀλλοφύλων καὶ ἕως ὁρίων Αἰγύπτου·
\vs{46l}καὶ Σαλωμὼν υἱὸς Δαυὶδ ἐβασίλευσεν ἐπὶ Ἰσραὴλ καὶ Ἰούδα ἐν Ἱερουσαλήμ.

\ch{3}
\vs{2}Πλὴν ὁ λαὸς ἦσαν θυμιῶντες ἐπὶ τοῖς ὑψηλοῖς, ὅτι οὐκ ᾠκοδομήθη οἶκος τῷ Κυρίῳ ἕως τοῦ νῦν.
\vs{3}Καὶ ἠγάπησε Σαλωμὼν τὸν Κύριον πορεύεσθαι ἐν τοῖς προστάγμασι Δαυὶδ τοῦ πατρὸς αὐτοῦ, πλὴν ἐν τοῖς ὑψηλοῖς ἔθυε καὶ ἐθυμία.
\vs{4}Καὶ ἀνέστη καὶ ἐπορεύθη εἰς Γαβαὼν θῦσαι ἐκεῖ, ὅτι αὕτη ὑψηλοτάτη, καὶ μεγάλη· χιλίαν ὁλοκαύτωσιν ἀνήνεγκε Σαλωμὼν ἐπὶ τὸ θυσιαστήριον ἐν Γαβαών.

\vs{5}Καὶ ὤφθη Κύριος τῷ Σαλωμὼν ἐν ὕπνῳ τὴν νύκτα, καὶ εἶπε Κύριος πρὸς Σαλωμὼν, αἴτησαί τι αἴτημα σεαυτῷ.
\vs{6}Καὶ εἶπε Σαλωμὼν, σὺ ἐποιήσας μετὰ τοῦ δούλου σου Δαυὶδ τοῦ πατρός μου ἔλεος μέγα, καθὼς διῆλθεν ἐνώπιόν σου ἐν ἀληθείᾳ καὶ ἐν δικαιοσύνῃ, καὶ ἐν εὐθύτητι καρδίας μετὰ σοῦ, καὶ ἐφύλαξας αὐτῷ τὸ ἔλεος τὸ μέγα τοῦτο, δοῦναι τὸν υἱὸν αὐτοῦ ἐπὶ τοῦ θρόνου αὐτοῦ, ὡς ἡ ἡμέρα αὕτη.
\vs{7}Καὶ νῦν, Κύριε ὁ Θεός μου, σὺ ἔδωκας τὸν δοῦλόν σου ἀντὶ Δαυὶδ τοῦ πατρός μου· καὶ ἐγώ εἰμι παιδάριον μικρὸν, καὶ οὐκ οἶδα τὴν ἔξοδόν μου καὶ τὴν εἴσοδόν μου.
\vs{8}Ὁ δὲ δοῦλός σου ἐν μέσῳ τοῦ λαοῦ σου, ὃν ἐξελέξω, λαὸν πολὺν, ὃς οὐκ ἀριθμηθήσεται.
\vs{9}Καὶ δώσεις τῷ δούλῳ σου καρδίαν ἀκούειν καὶ διακρίνειν τὸν λαόν σου ἐν δικαιοσύνῃ, καὶ τοῦ συνιεῖν ἀναμέσον ἀγαθοῦ καὶ κακοῦ· ὅτι τίς δυνηθήσεται κρίνειν τὸν λαόν σου τὸν βαρὺν τοῦτον;

\vs{10}Καὶ ἤρεσεν ἐνώπιον Κυρίου, ὅτι ᾐτῄσατο Σαλωμὼν τὸ ῥῆμα τοῦτο.
\vs{11}Καὶ εἶπε Κύριος πρὸς αὐτὸν, ἀνθʼ ὧν ᾐτήσω παρʼ ἐμοῦ τὸ ῥῆμα τοῦτο, καὶ οὐκ ᾐτήσω σεαυτῷ ἡμέρας πολλὰς, καὶ οὐκ ᾐτήσω πλοῦτον, οὐδὲ ᾐτήσω ψυχὰς ἐχθρῶν σου, ἀλλʼ ᾐτήσω σεαυτῷ τοῦ συνιεῖν τοῦ εἰσακούειν κρίμα,
\vs{12}ἰδοὺ πεποίηκα κατὰ τὸ ῥῆμά σου· ἰδοὺ δέδωκά σοι καρδίαν φρονίμην καὶ σοφήν· ὡς σὺ οὐ γέγονεν ἔμπροσθέν σου, καὶ μετὰ σὲ οὐκ ἀναστήσεται ὅμοιός σοι.
\vs{13}Καὶ ἃ οὐκ ᾐτήσω δέδωκά σοι, καὶ πλοῦτον καὶ δόξαν, ὡς οὐ γέγονεν ἀνὴρ ὅμοιός σοι ἐν βασιλεῦσι.
\vs{14}Καὶ ἐὰν πορευθῇς ἐν τῇ ὁδῷ μου φυλάσσειν τὰς ἐντολάς μου καὶ τὰ προστάγματά μου, ὡς ἐπορεύθη Δαυὶδ ὁ πατήρ σου, καὶ πληθυνῶ τὰς ἡμέρας σου.
\vs{15}Καὶ ἐξυπνίσθη Σαλωμὼν, καὶ ἰδοὺ ἐνύπνιον· καὶ ἀνέστη καὶ παραγίνεται εἰς Ἱερουσαλήμ, καὶ ἔστη κατὰ πρόσωπον τοῦ θυσιαστηρίου τοῦ κατὰ πρόσωπον κιβωτοῦ διαθήκης Κυρίου ἐν Σιὼν, καὶ ἀνήγαγεν ὁλοκαυτώσεις, καὶ ἐποίησεν εἰρηνικὰς, καὶ ἐποίησε πότον μέγαν ἑαυτῷ καὶ πᾶσι τοῖς παισὶν αὐτοῦ.

\vs{16}Τότε ὤφθησαν δύο γυναῖκες πόρναι τῷ βασιλεῖ, καὶ ἔστησαν ἐνώπιον αὐτοῦ.
\vs{17}Καὶ εἶπεν ἡ γυνὴ ἡ μία, ἐν ἐμοὶ κύριε, ἐγὼ καὶ ἡ γυνὴ αὕτη ᾠκοῦμεν ἐν οἴκῳ ἑνί, καὶ ἐτέκομεν ἐν τῷ οἴκῳ.
\vs{18}Καὶ ἐγενήθη ἐν τῇ ἡμέρᾳ τῇ τρίτῃ τεκούσης μου, ἔτεκε καὶ ἡ γυνὴ αὕτη· καὶ ἡμεῖς κατὰ τὸ αὐτό· καὶ οὐκ ἔστιν οὐθεὶς μεθʼ ἡμῶν πάρεξ ἀμφοτέρων ἡμῶν ἐν τῷ οἴκῳ.
\vs{19}Καὶ ἀπέθανεν ὁ υἱὸς τῆς γυναικὸς ταύτης τὴν νύκτα, ὡς ἐπεκοιμήθη ἐπʼ αὐτόν.
\vs{20}Καὶ ἀνέστη μέσης τῆς νυκτὸς, καὶ ἔλαβε τὸν υἱόν μου ἐκ τῶν ἀγκαλῶν μου, καὶ ἐκοίμισεν αὐτὸν ἐν τῷ κόλπῳ αὐτῆς, καὶ τὸν υἱὸν αὐτῆς τὸν τεθνηκότα ἐκοίμισεν ἐν τῷ κόλπῳ μου.
\vs{21}Καὶ ἀνέστην τοπρωῒ θηλάσαι τὸν υἱόν μου, καὶ ἐκεῖνος ἦν τεθνηκώς· καὶ ἰδοὺ κατενόησα αὐτὸν πρωῒ, καὶ ἰδοὺ οὐκ ἦν ὁ υἱός μου ὃν ἔτεκον.
\vs{22}Καὶ εἶπεν ἡ γυνὴ ἡ ἑτέρα, οὐχί, ἀλλὰ ὁ υἱός μου ὁ ζῶν, ὁ δὲ υἱός σου ὁ τεθνηκώς· καὶ ἐλάλησαν ἐνώπιον τοῦ βασιλέως.

\vs{23}Καὶ εἶπεν ὁ βασιλεὺς αὐταῖς, σὺ λέγεις, οὗτος ὁ υἱός μου ὁ ζῶν, καὶ ὁ υἱὸς ταύτης ὁ τεθνηκώς· καὶ σὺ λέγεις, οὐχί, ἀλλὰ ὁ υἱός μου ὁ ζῶν, καὶ ὁ υἱός σου ὁ τεθνηκώς.
\vs{24}Καὶ εἶπεν ὁ βασιλεύς, λάβετε μάχαιραν· καὶ προσήνεγκαν τὴν μάχαιραν ἐνώπιον τοῦ βασιλέως.
\vs{25}Καὶ εἶπεν ὁ βασιλεύς, διέλετε τὸ παιδίον τὸ ζῶν τὸ θηλάζον εἰς δύο, καὶ δότε τὸ ἥμισυ αὐτοῦ ταύτῃ, καὶ τὸ ἥμισυ αὐτοῦ ταύτῃ.
\vs{26}Καὶ ἀπεκρίθη ἡ γυνὴ ἧς ἦν ὁ υἱὸς ὁ ζῶν, καὶ εἶπε πρὸς τὸν βασιλέα, ὅτι ἐταράχθη ἡ μήτρα αὐτῆς ἐπὶ τῷ υἱῷ αὐτῆς, καὶ εἶπεν, ἐν ἐμοὶ κύριε, δότε αὐτῇ τὸ παιδίον, καὶ θανάτῳ μὴ θανατώσητε αὐτό· καὶ αὕτη εἶπε, μήτε ἐμοὶ, μήτε αὐτῇ ἔστω, διέλετε.
\vs{27}Καὶ ἀπεκρίθη ὁ βασιλεὺς, καὶ εἶπε, δότε τὸ παιδίον τῇ εἰπούσῃ, δότε αὐτῇ αὐτὸ, καὶ θανάτῳ μὴ θανατώσητε αὐτὸ, αὕτη ἡ μήτηρ αὐτοῦ.
\vs{28}Καὶ ἤκουσαν πᾶς Ἰσραὴλ τὸ κρίμα τοῦτο ὃ ἔκρινεν ὁ βασιλεὺς, καὶ ἐφοβήθησαν ἀπὸ προσώπου τοῦ βασιλέως, ὅτι εἶδον ὅτι φρόνησις Θεοῦ ἐν αὐτῷ τοῦ ποιεῖν δικαίωμα.

\ch{4}
Καὶ ἦν ὁ βασιλεὺς Σαλωμὼν βασιλεύων ἐπὶ Ἰσραήλ.
\vs{2}Καὶ οὗτοι ἄρχοντες οἳ ἦσαν αὐτῷ· Ἀζαρίας υἱὸς Σαδώκ·
\vs{3}Ἐλιὰφ καὶ Ἀχιὰ υἱὸς Σηβὰ γραμματεῖς· καὶ Ἰωσαφὰτ υἱὸς Ἀχιλοὺδ ἀναμιμνήσκων·
\vs{4}Καὶ Βαναίας υἱὸν Ἰωδαὲ ἐπὶ τῆς δυνάμεως· καὶ Σαδὼκ καὶ Ἀβιάθαρ ἱερεῖς·
\vs{5}Καὶ Ὀρνία υἱὸς Νάθαν ἐπὶ τῶν καθεσταμένων· καὶ Ζαβοὺθ υἱὸς Νάθαν ἑταῖρος τοῦ βασιλέως·
\vs{6}Καὶ Ἀχισὰρ ἦν οἰκονόμος· καὶ Ἐλιὰκ ὁ οἰκονόμος· καὶ Ἐλιὰβ υἱὸς Σὰφ ἐπὶ τῆς πατριᾶς· καὶ Ἀδωνιρὰμ υἱὸς Αὐδῶν ἐπὶ τῶν φόρων.

\vs{7}Καὶ τῷ Σαλωμὼν δώδεκα καθεστάμενοι ἐπὶ πάντα Ἰσραήλ, χορηγεῖν τῷ βασιλεῖ καὶ τῷ οἴκῳ αὐτοῦ· μῆνα ἐν τῷ ἐνιαυτῷ ἐγίνετο ἐπὶ τὸν ἕνα χορηγεῖν.
\vs{8}Καὶ ταῦτα τὰ ὀνόματα αὐτῶν. Βεὲν υἱὸς Ὢρ ἐν ὄρει Ἐφραὶμ εἷς.
\vs{9}Υἱὸς Δακὰρ ἐν Μακὲς, καὶ ἐν Σαλαβὶν, καὶ Βαιθσαμὺς, καὶ Ἐλὼν ἕως Βηθανὰν εἷς.
\vs{10}Υἱὸς Ἐσδὶ, ἐν Ἀραβὼθ, αὐτοῦ Σωχὼ καὶ πᾶσα ἡ γῆ Ὀφέρ.
\vs{11}Υἱοῦ Ἀμιναδὰβ πᾶσα Νεφθαδὼρ, Τεφὰθ θυγάτηρ Σαλωμὼν ἦν αὐτῷ εἰς γυναῖκα, εἷς.
\vs{12}Βανὰ υἱὸς Ἀχιλοὺθ τὴν Ἰθαανὰχ, καὶ Μαγεδδὼ, καὶ πᾶς ὁ οἶκος Σὰν ὁ παρὰ Σεσαθὰν ὑποκάτω τοῦ Ἐσραὲ, καὶ ἐκ Βηθσὰν ἕως Σαβελμαουλᾶ, ἕως Μαεβὲρ Λουκάμ, εἷς.
\vs{13}Υἱὸς Ναβὲρ ἐν Ῥαβὼθ Γαλαὰδ, τούτῳ σχοίνισμα Ἐργὰβ ἐν τῇ Βασὰν, ἑξήκοντα πόλεις μεγάλαι τειχήρεις καὶ μοχλοὶ χαλκοῖ, εἷς.
\vs{14}Ἀχιναδὰβ υἱὸς Σαδδὼ Μααναΐμ.
\vs{15}Ἀχιμαὰς ἐν Νεφθαλίμ, καὶ οὗτος ἔλαβε τὴν Βασεμμὰθ θυγατέρα Σαλωμὼν εἰς γυναῖκα, εἷς.
\vs{16}Βαανὰ υἱὸς Χουσὶ ἐν Ἀσὴρ καὶ ἐν Βααλὼθ, εἷς.
\vs{17}Σεμεῒ υἱὸς Ἠλὰ ἐν τῷ Βενιαμίν.
\vs{18}Γαβὲρ υἱὸς Ἀδαῒ ἐν τῇ γῇ Γὰδ Σηὼν βασιλέως τοῦ Ἐσεβὼν καὶ Ὢγ βασιλέως τοῦ Βασὰν, καὶ νασὲφ εἷς ἐν γῇ Ἰούδα.
\vs{19}Ἰωσαφὰτ υἱὸς Φουασοὺδ ἐν Ἱσσάχαρ.

\ch{5}Καὶ ἐχορήγουν οἱ καθεστάμενοι οὕτως τῷ βσαιλεῖ Σαλωμών· καὶ πάντα τὰ διαγγέλματα ἐπὶ τὴν τράπεζαν τοῦ βασιλέως ἕκαστος μῆνα αὐτοῦ, οὐ παραλλάσσουσι λόγον. Καὶ τὰς κριθὰς καὶ τὸ ἄχυρον τοῖς ἵπποις καὶ τοῖς ἅρμασιν ᾖρον εἰς τὸν τόπον οὗ ἂν ᾖ ὁ βασιλεὺς, ἕκαστος κατὰ τὴν σύνταξιν αὐτοῦ.

\vs{2}Καὶ ταῦτα τὰ δέοντα τῷ Σαλωμών· ἐν ἡμέρᾳ μιᾷ τριάκοντα κόροι σεμιδάλεως, καὶ ἑξήκοντα κόροι ἀλεύρου κεκοπανισμένου,
\vs{3}καὶ δέκα μόσχοι ἐκλεκτοὶ, καὶ εἴκοσι βόες νομάδες, καὶ ἑκατὸν πρόβατα, ἐκτὸς ἐλάφων, καὶ δορκάδων ἐκλεκτῶν σιτευτά.
\vs{4}Ὅτι ἦν ἄρχων πέραν τοῦ ποταμοῦ, καὶ ἦν αὐτῷ εἰρήνη ἐκ πάντων τῶν μερῶν κυκλόθεν.

\vs{9}Καὶ ἔδωκε Κύριος φρόνησιν τῷ Σαλωμὼν καὶ σοφίαν πολλὴν σφόδρα καὶ χύμα καρδίας, ὡς ἡ ἄμμος ἡ παρὰ τὴν θάλασσαν.
\vs{10}Καὶ ἐπληθύνθη Σαλωμὼν σφόδρα ὑπὲρ τὴν φρόνησιν πάντων ἀρχαίων ἀνθρώπων, καὶ ὑπὲρ πάντας φρονίμους Αἰγύπτου.
\vs{11}Καὶ ἐσοφίσατο ὑπὲρ πάντας τοὺς ἀνθρώπους· καὶ ἐσοφίσατο ὑπὲρ Γαιθὰν τὸν Ζαρίτην, καὶ τὸν Αἰνὰν, καὶ τὸν Χαλκὰδ καὶ Δαράλα υἱοὺς Μάλ.
\vs{12}Καὶ ἐλάλησε Σαλωμὼν τρισχιλίας παραβολὰς, καὶ ἦσαν ᾠδαι αὐτοῦ πεντακισχίλιαι.
\vs{13}Καὶ ἐλάλησεν ὑπὲρ τῶν ξύλων ἀπὸ τῆς κέδρου τῆς ἐν τῷ Λιβάνῳ, καὶ ἕως τῆς ὑσσώπου τῆς ἐκπορευομένης διὰ τοῦ τοίχου· καὶ ἐλάλησε περὶ τῶν κτηνῶν καὶ περὶ τῶν πετεινῶν καὶ περὶ τῶν ἑρπετῶν καὶ περὶ τῶν ἰχθύων.
\vs{14}Καὶ παρεγίνοντο πάντες οἱ λαοὶ ἀκοῦσαι τῆς σοφίας Σαλωμών· καὶ παρὰ πάντων τῶν βασιλέων τῆς γῆς, ὅσοι ἤκουον τῆς σοφίας αὐτοῦ·

\vs{14a}Καὶ ἔλαβε Σαλωμὼν τὴν θυγατέρα Φαραὼ αὐτῷ εἰς γυναῖκα, καὶ εἰσήγαγεν αὐτὴν εἰς τὴν πόλιν Δαυὶδ ἕως συντελέσαι αὐτὸν τὸν οἶκον Κυρίου, καὶ τὸν οἶκον ἑαυτοῦ, καὶ τὸ τεῖχος Ἱερουσαλήμ·
\vs{14b}τότε ἀνέβη Φαραὼ βασιλεὺς Αἰγύπτου, καὶ προκατελάβετο τὴν Γαζὲρ, καὶ ἐνεπύρισεν αὐτὴν, καὶ τὸν Χανανίτην τὸν κατοικοῦντα ἐν Μεργάβ· καὶ ἔδωκεν αὐτὰς Φαραὼ ἀποστολὰς θυγατρὶ αὐτοῦ γυναικὶ Σαλωμών· καὶ Σαλωμὼν ᾠκοδόμησε τὴν Γαζέρ.

\vs{15}Καὶ ἀπέστειλε Χιρὰμ βασιλεὺς Τύρου τοὺς παῖδας αὐτοῦ χρίσαι τὸν Σαλωμὼν ἀντὶ Δαυὶδ τοῦ πατρὸς αὐτοῦ, ὅτι ἀγαπῶν ἦν Χιρὰμ τὸν Δαυὶδ πάσας τὰς ἡμέρας.
\vs{16}Καὶ ἀπέστειλε Σαλωμὼν πρὸς Χιρὰμ, λέγων,
\vs{17}Σὺ οἶδας τὸν πατέρα μου Δαυὶδ, ὅτι οὐκ ἠδύνατο οἰκοδομῆσαι οἶκον τῷ ὀνόματι Κυρίου Θεοῦ μου ἀπὸ προσώπου τῶν πολέμων τῶν κυκλωσάντων αὐτὸν, ἕως τοῦ δοῦναι Κύριον αὐτοὺς ὑπὸ τὰ ἴχνη τῶν ποδῶν αὐτοῦ.
\vs{18}Καὶ νῦν ἀνέπαυσε Κύριος ὁ Θεός μου ἐμοὶ κυκλόθεν, οὐκ ἔστιν ἐπίβουλος καὶ οὐκ ἔστιν ἁμάρτημα πονηρόν.
\vs{19}Καὶ ἰδοὺ ἐγὼ λέγω οἰκοδομῆσαι οἶκον τῷ ὀνόματι Κυρίου Θεοῦ μου, καθὼς ἐλάλησε Κύριος ὁ Θεὸς πρὸς Δαυὶδ τὸν πατέρα μου, λέγων, ὁ υἱός σου ὃν δώσω ἀντὶ σοῦ ἐπὶ τὸν θρόνον σου, οὗτος οἰκοδομήσει τὸν οἶκον τῷ ὀνόματί μου.
\vs{20}Καὶ νῦν ἔντειλαι, καὶ κοψάτωσάν μοι ξύλα ἐκ τοῦ Λιβάνου· καὶ ἰδοὺ οἱ δοῦλοί μου μετὰ τῶν δούλων σου, καὶ τὸν μισθὸν δουλείας σου δώσω σοι κατὰ πάντα ὅσα ἂν εἴπῃς, ὅτι σὺ οἶδας, ὅτι οὐκ ἔστιν ἡμῖν εἰδὼς ξύλα κόπτειν καθὼς οἱ Σιδώνιοι.

\vs{21}Καὶ ἐγενήθη καθὼς ἤκουσε Χιρὰμ τῶν λόγων Σαλωμὼν, ἐχάρη σφόδρα, καὶ εἶπεν, εὐλογητὸς ὁ Θεὸς σήμερον, ὃς ἔδωκε τῷ Δαυὶδ υἱὸν φρόνιμον ἐπὶ τὸν λαὸν τὸν πολὺν τοῦτον.
\vs{22}Καὶ ἀπέστειλε πρὸς Σαλωμὼν, λέγων, ἀκήκοα περὶ πάντων ὧν ἀπέσταλκας πρὸς μέ· ἐγὼ ποιήσω πᾶν θέλημά σου· ξύλα κέδρινα καὶ πεύκινα
\vs{23}οἱ δοῦλοί μου κατάξουσιν αὐτὰ ἐκ τοῦ Λιβάνου εἰς τὴν θάλασσαν, ἐγὼ θήσομαι αὐτὰ σχεδίας, ἕως τοῦ τόπου οὗ ἐὰν ἀποστείλῃς πρὸς μὲ, καὶ ἐκτινάξω αὐτὰ ἐκεῖ, καὶ σὺ ἀρεῖς· καὶ ποιήσεις τὸ θέλημά μου, τοῦ δοῦναι ἄρτους τῷ οἴκῳ μου.

\vs{24}Καὶ ἦν Χιρὰμ διδοὺς τῷ Σαλωμὼν κέδρους καὶ πεύκας καὶ πᾶν θέλημα αὐτοῦ.
\vs{25}Καὶ Σαλωμὼν ἔδωκε τῷ Χιρὰμ εἴκοσι χιλιάδας κόρους πυροῦ καὶ μαχεὶρ τῷ οἴκῳ αὐτοῦ, καὶ εἴκος χιλιάδας βαὶθ ἐλαίου κεκομμένου· κατὰ τοῦτο ἐδίδου Σαλωμὼν τῷ Χιρὰμ κατʼ ἐνιαυτόν.
\vs{26}Καὶ Κύριος ἔδωκε σοφίαν τῷ Σαλωμὼν καθὼς ἐλάλησεν αὐτῷ· καὶ ἦν εἰρήνη ἀναμέσον Χιρὰμ καὶ ἀναμέσον Σαλωμὼν, καὶ διέθεντο διαθήκην ἀναμέσον αὐτῶν.

\vs{27}Καὶ ἀνήνεγκεν ὁ βασιλεὺς φόρον ἐκ παντὸς Ἰσραὴλ, καὶ ἦν ὁ φόρος τριάκοντα χιλιάδες ἀνδρῶν.
\vs{28}Καὶ ἀπέστειλεν αὐτοὺς εἰς τὸν Λίβανον, δέκα χιλιάδες ἐν τῷ μηνὶ ἀλλασσόμενοι· μῆνα ἦσαν ἐν τῷ Λιβάνῳ, καὶ δύο μῆνας ἐν οἴκῳ αὐτῶν· καὶ Ἀδωνιρὰμ ἐπὶ τοῦ φόρου.
\vs{29}Καὶ ἦν τῷ Σαλωμὼν ἑβδομήκοντα χιλιάδες αἴροντες ἄρσιν, καὶ ὀγδοήκοντα χιλιάδες λατόμων ἐν τῷ ὄρει,
\vs{30}χωρὶς τῶν ἀρχόντων τῶν καθεσταμένων ἐπὶ τῶν ἔργων τῷ Σαλωμὼν, τρεῖς χιλιάδες καὶ ἑξακόσιοι ἐπιστάται οἱ ποιοῦντες τὰ ἔργα.
\vs{32}Καὶ ἡτοίμασαν τοὺς λίθους καὶ τὰ ξύλα τρία ἔτη.

\ch{6}
Καὶ ἐγενήθη ἐν τῷ τεσσαρακοστῷ καὶ τετρακοσιοστῷ ἔτει τῆς ἐξόδου υἱῶν Ἰσραὴλ ἐξ Αἰγύπτου, τῷ ἔτει τῷ τετάρτῳ ἐν μηνὶ τῷ δευτέρῳ βασιλεύοντος τοῦ βασιλέως Σαλωμὼν ἐπὶ τὸν Ἰσραὴλ,
\vs{1a}καὶ ἐνετείλατο ὁ βασιλεὺς ἵνα ἄρωσι λίθους μεγάλους τιμίους εἰς τὸν θεμέλιον τοῦ οἴκου, καὶ λίθους ἀπελεκήτους.
\vs{1b}Καὶ ἐπελέκησαν οἱ υἱοὶ Σαλωμὼν, καὶ οἱ υἱοὶ Χιρὰμ, καὶ ἔβαλαν αὐτούς.

\vs{1c}Ἐν τῷ ἔτει τῷ τετάρτῷ ἐθεμελίωσε τὸν οἶκον Κυρίου ἐν μηνὶ Ζιοῦ, καὶ τῷ δευτέρῳ μηνί.
\vs{1d}Ἐν ἑνδεκάτῳ ἐνιαυτῷ, ἐν μηνὶ Βαὰλ, οὗτος ὁ μὴν ὁ ὄγδοος, συνετελέσθη ὁ οἶκος εἰς πάντα λόγον αὐτοῦ, καὶ εἰς πᾶσαν διάταξιν αὐτοῦ.
\vs{2}Καὶ ὁ οἶκος ὃν ᾠκοδόμησεν ὁ βασιλεὺς τῷ Κυρίῳ, τεσσαράκοντα ἐν πήχει μῆκος αὐτοῦ, καὶ εἴκοσι ἐν πήχει πλάτος αὐτοῦ, καὶ πέντε καὶ εἴκοσι ἐν πήχει τὸ ὕψος αὐτοῦ·
\vs{3}Καὶ τὸ αἰλὰμ κατὰ πρόσωπον τοῦ ναοῦ, εἴκοσι ἐν πήχει μῆκος αὐτοῦ εἰς τὸ πλάτος τοῦ οἴκου, κατὰ πρόσωπον τοῦ οἴκου, καὶ ᾠκοδόμησε τὸν οἶκον, καὶ συνετέλεσεν αὐτόν.
\vs{4}Καὶ ἐποίησε τῷ οἴκῳ θυρίδας παρακυπτομένας κρυπτάς.

\vs{5}Καὶ ἔδωκεν ἐπὶ τὸν τοῖχον τοῦ οἴκου μέλαθρα κυκλόθεν τῷ ναῷ καὶ τῷ δαβίρ.
\vs{6}Ἡ πλευρὰ ἡ ὑποκάτω πέντε πήχεων ἐν πήχει τὸ πλάτος αὐτῆς, καὶ τὸ μέσον ἓξ, καὶ ἡ τρίτη ἑπτὰ ἐν πήχει τὸ πλάτος αὐτῆς· ὅτι διάστημα ἔδωκε τῷ οἴκῳ κυκλόθεν ἔξωθεν τοῦ οἴκου, ὅπως μὴ ἐπιλαμβάνωνται τῶν τοίχων τοῦ οἴκου.
\vs{7}Καὶ ὁ οἴκος ἐν τῷ οἰκοδομεῖσθαι αὐτὸν λίθοις ἀκροτομοις ἀργοῖς ᾠκοδομήθη· καὶ σφύρα καὶ πέλεκυς καὶ πᾶν σκεῦος σιδηροῦν οὐκ ἠκούσθη ἐν τῷ οἴκῳ ἐν τῷ οἰκοδομεῖσθαι αὐτόν.
\vs{8}Καὶ ὁ πυλὼν τῆς πλευρᾶς τῆς ὑποκάτωθεν ὑπὸ τὴν ὠμίαν τοῦ οἴκου τὴν δεξιὰν, καὶ ἑλικτὴ ἀνάβασις εἰς τὸ μέσον, καὶ ἐκ τῆς μέσης ἐπὶ τὸ τριόροφα.
\vs{9}Καὶ ᾠκοδόμησε τὸν οἶκον καὶ συνετέλεσεν αὐτόν· καὶ ἐκοιλοστάθμησε τὸν οἶκον κέδροις.
\vs{10}Καὶ ᾠκοδόμησε τοὺς ἐνδέσμους, διʼ ὅλου τοῦ οἴκου πέντε ἐν πήχει τὸ ὕψος αὐτοῦ, καὶ συνέσχε τὸν σύνδεσμον ἐν ξύλοις κεδρίνοις.

\vs{15}Καὶ ᾠκοδόμησε τοὺς τοίχους τοῦ οἴκου ἔσωθεν διὰ ξύλων κεδρίνων ἀπὸ τοῦ ἐδάφους τοῦ οἴκου καὶ ἕως τῶν τοίχων καὶ ἕως τῶν δοκῶν· ἐκοιλοστάθμησε συνεχόμενα ξύλοις ἔσωθεν· καὶ περιέσχε τὸ ἔσω τοῦ οἴκου ἐν πλευραῖς πευκίναις.
\vs{16}Καὶ ᾠκοδόμησε τοὺς εἴκοσι πήχεις ἀπʼ ἄκρου τοῦ τοίχου τὸ πλευρὸν τὸ ἓν ἀπὸ τοῦ ἐδάφους ἕως τῶν δοκῶν· καὶ ἐποίησεν ἐκ τοῦ δαβὶρ εἰς τὸ ἁγιον τῶν ἁγίων.
\vs{17}Καὶ τεσσαράκοντα πήχεων ἦν ὁ ναὸς κατὰ πρόσωπον τοῦ δαβὶρ ἐν μέσῳ τοῦ οἴκου ἔσωθεν,
\vs{19}δοῦναι ἐκεῖ τὴν κιβωτὸν διαθήκης Κυρίου.
\vs{20}Εἴκοσι πήχεις μῆκος, καὶ εἴκοσι πήχεις πλάτος, καὶ εἴκοσι πήχεις τὸ ὕψος αὐτοῦ· Καὶ περιέσχεν αὐτὸ χρυσίῳ συγκεκλεισμένῳ· καὶ ἐποίησε θυσιαστήριον
\vs{21}κατὰ πρόσωπον τοῦ δαβὶρ, καὶ περιέσχεν αὐτὸ χρυσίῳ.
\vs{22}Καὶ ὅλον τὸν οἶκον περιέσχε χρυσίῳ, ἕως συντελείας παντὸς τοῦ οἴκου.

\vs{23}Καὶ ἐποίησεν ἐν τῷ δαβὶρ δύο χερουβὶμ δέκα πήχεων μέγεθος ἐσταθμωμένον·
\vs{24}Καὶ πέντε πήχεων πτερύγιον τοῦ χερουβὶμ τοῦ ἑνὸς, καὶ πέντε πήχεων πτερύγιον αὐτοῦ τὸ δεύτερον, ἐν πήχει δέκα ἀπὸ μέρους πτερυγίου αὐτοῦ εἰς μέρος πτερυγίου αὐτοῦ.
\vs{25}Οὕτως τῷ χερουβὶμ τῷ δευτέρῳ, ἐν μέτρῳ ἑνὶ συντέλεια μία ἀμφοτέροις.
\vs{26}Καὶ τὸ ὕψος τοῦ χερουβὶμ τοῦ ἑνὸς δέκα ἐν πήχει· καὶ οὕτω τῷ χερουβὶμ τῷ δευτέρῳ.
\vs{27}Καὶ ἀμφότερα χερουβὶμ ἐν μέσῳ τοῦ οἴκου τοῦ ἐσωτάτου· καὶ διεπέτασε τὰς πτέρυγας αὐτῶν, καὶ ἥπτετο πτέρυξ μία τοῦ τοίχου, καὶ πτέρυξ χερουβὶμ τοῦ δευτέρου ἥπτετο τοῦ τοίχου τοῦ δευτέρου· καὶ αἱ πτέρυγες αὐτῶν ἐν μέσῳ τοῦ οἴκου ἥπτοντο πτέρυξ πτέρυγος.
\vs{28}Καὶ περιέσχε τὰ χερουβὶμ χρυσίῳ.

\vs{29}Πάντας τοὺς τοίχους τοῦ οἴκου κύκλῳ ἐγκολαπτὰ ἔγραψε γραφίδι χερουβὶμ, καὶ φοίνικας τῷ ἐσωτέρῳ καὶ τῷ ἐξωτέρῳ.
\vs{30}Καὶ τὸ ἔδαφος τοῦ οἴκου περιέσχε χρυσίῳ τοῦ ἐσωτάτου καὶ τοῦ ἐξωτάτου.

\vs{31}Καὶ τῷ θυρώματι τοῦ δαβὶρ ἐποίησε θύρας ξύλων ἀρκευθίνων, στοαὶ τετραπλῶς,
\vs{34}ἐν ἀμφοτέραις ταῖς θύραις ξύλα πεύκινα· δύο πτυχαὶ ἡ θύρα ἡ μία καὶ στροφεῖς αὐτῶν, καὶ δύο πτυχαὶ ἡ θύρα ἡ δευτέρα στρεφόμενα ἐγκεκολαμμένα χερουβὶμ,
\vs{35}καὶ φοίνικες, καὶ διαπεπετασμένα πέταλα, καὶ περιεχόμενα χρυσίῳ καταγομένῳ ἐπὶ τὴν ἐκτύπωσιν.
\vs{36}Καὶ ᾠκοδόμησε τὴν αὐλὴν τὴν ἐσωτάτην· τρεῖς στίχους ἀπελεκήτων, καὶ στίχος κατειργασμένης κέδρου κυκλόθεν·
\vs{36a}καὶ ᾠκοδόμησε τὸ καταπέτασμα τῆς αὐλῆς τοῦ αἰλὰμ τοῦ οἴκου τοῦ κατὰ πρόσωπον τοῦ ναοῦ.

\ch{7}
Καὶ ἀπέστειλεν ὁ βασιλεὺς Σαλωμὼν, καὶ ἔλαβε τὸν Χιρὰμ ἐκ Τύρου,
\vs{2}υἱὸν γυναικὸς χήρας, καὶ οὗτος ἀπὸ τῆς φυλῆς τῆς Νεφθαλίμ, καὶ ὁ πατὴρ αὐτοῦ ἀνὴρ Τύριος· τέκτων χαλκοῦ, καὶ πεπληρωμένος τῆς τέχνης καὶ συνέσεως καὶ ἐπιγνώσεως τοῦ ποιεῖν πᾶν ἔργον ἐν χαλκῷ· καὶ εἰσηνέχθη πρὸς τὸν βασιλέα Σαλωμών· καὶ ἐποίησε πάντα τὰ ἔργα.

\vs{3}Καὶ ἐχώνευσε τοὺς δύο στύλους τῷ αἰλὰμ τοῦ οἴκου· ὀκτωκαίδεκα πήχεις ὕψος τοῦ στύλου· καὶ περίμετρον τεσσαρεσκαίδεκα πήχεις ἐκύκλου αὐτὸν τὸ πάχος τοῦ στύλου· τεσσάρων δσκτύλων τὰ κοιλώματα· καὶ οὕτως ὁ στύλος ὁ δεύτερος·
\vs{4}Καὶ δύο ἐπιθέματα ἐποίησε δοῦναι ἐπὶ τὰς κεφαλὰς τῶν στύλων χωνευτά· πέντε πήχεις τὸ ὕψος τοῦ ἐπιθέματος τοῦ ἑνὸς, καὶ πέντε πήχεις τὸ ὕψος τοῦ ἐπιθέματος τοῦ δευτέρου
\vs{5}Καὶ ἐποίησε δύο δίκτυα περικαλύψαι τὸ ἐπίθεμα τῶν στύλων· καὶ δίκτυον τῷ ἐπιθέματι τῷ ἑνὶ, καὶ δίκτυον τῷ ἐπιθέματι τῷ δευτέρῳ.
\vs{6}Καὶ ἔργον κρεμαστὸν, δύο στίχοι ῥοῶν χαλκῶν, δεδικτυωμένοι, ἔργον κρεμαστὸν, στίχος ἐπὶ στίχον· καὶ οὕτως ἐποίησε τῷ ἐπιθέματι τῷ δευτέρῳ.
\vs{7}Καὶ ἔστησε τοὺς στύλους τοῦ αἰλὰμ τοῦ ναοῦ· καὶ ἔστησε τὸν στύλον τὸν ἕνα, καὶ ἐπεκάλεσε τὸ ὄνομα αὐτοῦ Ἰαχούμ· καὶ ἔστησε τὸν στύλον τὸν δεύτερον, καὶ ἐπεκάλεσε τὸ ὄνομα αὐτοῦ Βολώζ.
\vs{8}Καὶ ἐπὶ τῶν κεφαλῶν τῶν στύλων ἔργον κρίνου κατὰ τὸ αἰλὰμ τεσσάρων πηχῶν·
\vs{9}καὶ μέλαθρον ἐπʼ ἀμφοτέρων τῶν στύλων· καὶ ἐπάνωθεν τῶν πλευρῶν ἐπίθεμα τὸ μέλαθρον τῷ πάχει.

\vs{10}Καὶ ἐποίησε τὴν θάλασσαν δέκα ἐν πήχει ἀπὸ τοῦ χείλους αὐτῆς ἕως τοῦ χείλους αὐτῆς, στρογγυλόν κύκλῳ τὸ αὐτό· πέντε ἐν πήχει τὸ ὕψος αὐτῆς· καὶ συνηγμένη τρεῖς καὶ τριάκοντα ἐν πήχει.
\vs{11}Καὶ ὑποστηρίγματα ὑποκάτωθεν τοῦ χείλους αὐτῆς κυκλόθεν ἐκύκλουν αὐτὴν δέκα ἐν πήχει κυκλόθεν·
\vs{12}καὶ τὸ χεῖλος αὐτῆς ὡς ἔργον χείλους ποτηρίου βλαστὸς κρίνου· καὶ τὸ πάχος αὐτοῦ παλαιστής.
\vs{13}Καὶ δώδεκα βόες ὑποκάτω τῆς θαλάσσης, οἱ τρεῖς ἐπιβλέποντες βοῤῥὰν, καὶ οἱ τρεῖς ἐπιβλέποντες θάλασσαν, καὶ οἱ τρεῖς ἐπιβλέποντες Νότον, καὶ οἱ τρεῖς ἐπιβλέποντες ἀνατολήν· καὶ πάντα τὰ ὀπίσθια εἰς τὸν οἶκον, καὶ ἡ θάλασσα ἐπʼ αὐτῶν ἐπάνωθεν.

\vs{14}Καὶ ἐποίησε δέκα μεχωνὼθ χαλκᾶς· πέντε πήχεις μῆκος τῆς μεχωνὼθ τῆς μιᾶς, καὶ τέσσαρες πήχεις τὸ πλάτος αὐτῆς, καὶ ἓξ ἐν πήχει ὕψος αὐτῆς.
\vs{15}Καὶ τοῦτο τὸ ἔργον τῶν μεχωνὼθ συγκλειστὸν αὐτοῖς, καὶ συγκλειστὸν ἀναμέσον τῶν ἐξεχομένων.
\vs{16}Καὶ ἐπὶ τὰ συγκλείσματα αὐτῶν ἀναμέσον ἐξεχομένων λέοντες καὶ βόες καὶ χερουβὶμ, καὶ ἐπὶ τῶν ἐξεχομένων, οὕτως καὶ ἐπάνωθεν, καὶ ὑποκάτωθεν τῶν λεόντων καὶ τῶν βοῶν χῶραι, ἔργον καταβάσεως.
\vs{17}Καὶ τέσσαρες τροχοὶ χαλκοῖ τῇ μεχωνὼθ τῇ μιᾷ, καὶ τὰ προσέχοντα χαλκᾶ καὶ τέσσαρα μέρη αὐτῶν, ὠμίαι ὑποκάτω τῶν λουτήρων.
\vs{18}Καὶ χεῖρες ἐν τοῖς τροχοῖς ἐν τῇ μεχωνώθ. Καὶ τὸ ὕψος τοῦ τροχοῦ τοῦ ἑνὸς πήχεος καὶ ἡμίσους.
\vs{19}Καὶ τὸ ἔργον τῶν τροχῶν ἔργον τροχῶν ἅρματος· αἱ χεῖρες αὐτῶν καὶ οἱ νῶτοι αὐτῶν καὶ ἡ πραγματεία αὐτῶν πάντα χωνευτά.
\vs{20}Αἱ τέσσαρες ὠμίαι ἐπὶ τῶν τεσσάρων γωνιῶν τῆς μεχωνὼθ τῆς μιᾶς, ἐκ τῆς μεχωνὼθ οἱ ὦμοι αὐτῆς.
\vs{21}Καὶ ἐπὶ τῆς κεφαλῆς τῆς μεχωνὼθ ἥμισυ τοῦ πήχεος μέγεθος αὐτῆς στρογγύλον κύκλῳ ἐπὶ τῆς κεφαλῆς τῆς μεχωνώθ· καὶ ἀρχὴ χειρῶν αὐτῆς καὶ τὰ συγκλείσματα αὐτῆς· καὶ ἠνοίγετο ἐπὶ τὰς ἀρχὰς τῶν χειρῶν αὐτῆς.
\vs{22}Καὶ τὰ συγκλείσματα αὐτῆς χερουβὶμ καὶ λέοντες καὶ φοίνικες ἑστῶτα, ἐχόμενον ἕκαστον κατὰ πρόσωπον ἔσω καὶ τὰ κυκλόθεν.
\vs{23}Κατʼ αὐτὴν ἐποίησε πάσας τὰς δέκα μεχωνὼθ, τάξιν μίαν καὶ μέτρον ἓν πάσαις.
\vs{24}Καὶ ἐποίησε δέκα χυτροκαύλους χαλκοῦς, τεσσαράκοντα χοεῖς χωροῦντα τὸν ἕνα χυτρόκαυλον μετρήσει τεσσάρων πήχων· χυτρόκαυλος ὁ εἷς ἐπὶ τῇ μεχωνὼθ τῇ μιᾷ ταῖς δέκα μεχωνώθ.
\vs{25}Καὶ ἔθετο τὰς πέντε μεχωνὼθ ἀπὸ τῆς ὠμίας τοῦ οἴκου ἐκ δεξιῶν, καὶ πέντε ἀπὸ τῆς ὠμίας τοῦ οἴκου ἐξ ἀριστερῶν· καὶ ἡ θάλασσα ἀπὸ τῆς ὡμίας τοῦ οἴκου ἐκ δεξιῶν κατʼ ἀνατολὰς ἀπὸ τοῦ κλίτους τοῦ Νότου.

\vs{26}Καὶ ἐποίησε Χιρὰμ τοὺς λέβητας καὶ τὰς θερμαστρεῖς καὶ τὰς φιάλας· καὶ συνετέλεσε Χιρὰμ ποιῶν πάντα τὰ ἔργα ἃ ἐποίησε τῷ βασιλεῖ Σαλωμὼν ἐν οἴκῳ Κυρίου·
\vs{27}Στύλους δύο, καὶ τὰ στρεπτὰ τῶν στύλων ἐπὶ τῶν κεφαλῶν τῶν στύλων δύο· καὶ τὰ δίκτυα δύο τοῦ καλύπτειν ἀμφότερα τὰ στρεπτὰ τῶν γλυφῶν τὰ ὄντα ἐπὶ τῶν στύλων.
\vs{28}Τὰς ῥοὰς τετρακοσίας ἀμφοτέροις τοῖς δικτύοις, δύο στίχοι ῥοῶν τῷ δικτύῳ τῷ ἑνὶ, περικαλύπτειν ἀμφότερα τὰ ὄντα τὰ στρεπτὰ τῆς μεχωνὼθ ἐπʼ ἀμφοτέροις τοῖς στύλοις·
\vs{29}Καὶ τὰ μεχωνὼθ δέκα, καὶ τοὺς χυτροκαύλους δέκα ἐπὶ τῶν μεχωνώθ·
\vs{30}Καὶ τὴν θάλασσαν μίαν, καὶ τοὺς βόας δώδεκα ὑποκάτω τῆς θαλάσσης·
\vs{31}Καὶ τοὺς λέβητας καὶ τὰς θερμαστρεῖς καὶ τὰς φιάλας καὶ πάντα τὰ σκεύη, ἃ ἐποίησε Χιρὰμ τῷ βασιλεῖ Σαλωμὼν τῷ οἴκῳ Κυρίου· καὶ οἱ στύλοι τεσσαράκοντα καὶ ὀκτὼ τοῦ οἴκου τοῦ βασιλέως καὶ τοῦ οἴκου Κυρίου· πάντα τὰ ἔργα τοῦ βασιλέως ἃ ἐποίησε Χιρὰμ χαλκᾶ ἄρδην.
\vs{32}Οὐκ ἦν σταθμὸς τοῦ χαλκοῦ οὗ ἐποίησε πάντα τὰ ἔργα ταῦτα ἐκ πλήθους σφόδρα· οὐκ ἦν τέρμα τῶν σταθμῶν τοῦ χαλκοῦ.
\vs{33}Ἐν τῷ περιοίκῳ τοῦ Ἰορδάνου ἐχώνευσεν αὐτὰ ἐν τῷ πάχει τῆς γῆς ἀναμέσον Σοκχὼθ καὶ ἀναμέσον Σειρά.

\vs{34}Καὶ ἔλαβεν ὁ βασιλεὺς Σαλωμὼν τὰ σκεύη ἃ ἐποίησεν ἐν οἴκῳ Κυρίου, τὸ θυσιαστήριον τὸ χρυσοῦν, καὶ τὴν τράπεζαν ἐφʼ ἧς οἱ ἄρτοι τῆς προσφορᾶς, χρυσῆν,
\vs{35}καὶ τὰς λυχνίας πέντε ἐξ ἀριστερῶν, καὶ πέντε ἐκ δεξιῶν κατὰ πρόσωπον τοῦ δαβὶρ χρυσᾶς συγκλειομένας, καὶ τὰ λαμπάδια, καὶ τοὺς λύχνους, καὶ τὰς ἐπαρύστρις χρυσᾶς.
\vs{36}Καὶ τὰ πρόθυρα, καὶ οἱ ἧλοι, καὶ αἱ φιάλαι, καὶ τὰ τρυβλία, καὶ αἱ θυΐσκαι χρυσαῖ, συγκλειστά· καὶ τὰ θυρώματα τῶν θυρῶν τοῦ οἴκου τοῦ ἐσωτάτου ἁγίου τῶν ἁγίων, καὶ τὰς θύρας τοῦ ναοῦ χρυσᾶς.

\vs{37}Καὶ ἀνεπληρώθη τὸ ἔργον ὃ ἐποίησε Σαλωμὼν οἴκου Κυρίου· καὶ εἰσήνεγκε Σαλωμὼν τὰ ἅγια Δαυὶδ τοῦ πατρὸς αὐτοῦ, καὶ πάντα τὰ ἅγια Σαλωμὼν, τὸ ἀργύριον καὶ τὸ χρυσίον καὶ τὰ σκεύη ἔδωκεν εἰς τοὺς θησαυροὺς οἴκου Κυρίου.

\vs{38}Καὶ τὸν οἶκον ἑαυτῷ ᾠκοδόμησε Σαλωμὼν τρισκαίδεκα ἔτεσι·
\vs{39}Καὶ ᾠκοδόμησε τὸν οἶκον δρυμῷ τοῦ Λιβάνου· ἑκατὸν πήχεις μῆκος αὐτοῦ, καὶ πεντήκοντα πήχεις πλάτος αὐτοῦ, καὶ τριάκοντα πηχῶν ὕψος αὐτοῦ· καὶ τριῶν στύχων στύλων κεδρίνων, καὶ ὠμίαι κέδριναι τοῖς στύλοις.
\vs{40}Καὶ ἐφάτνωσε τὸν οἶκον ἄνωθεν ἐπὶ τῶν πλευρῶν τῶν στύλων· καὶ ἀριθμὸς τῶν στύλων τεσσαράκοντα καὶ πέντε ὁ στίχος,
\vs{41}καὶ μέλαθρα τρία, καὶ χῶρα ἐπὶ χώραν τρισσῶς.
\vs{42}Καὶ πάντα τὰ θυρώματα, καὶ αἱ χῶραι τετράγωνοι μεμελαθρωμέναι· καὶ ἀπὸ τοῦ θυρώματος ἐπὶ θύραν τρισσῶς.
\vs{43}Καὶ τὸ αἰλὰμ τῶν στύλων, πεντήκοντα μῆκος, καὶ πεντήκοντα ἐν πλάτει ἐζυγωμένα αἰλὰμ ἐπὶ πρόσωπον αὐτῶν· καὶ στύλοι καὶ πάχος ἐπὶ πρόσωπον αὐτῆς τοῖς αἰλαμίν.
\vs{44}Καὶ τὸ αἰλὰμ τῶν θρόνων οὗ κρινεῖ ἐκεῖ, αἰλὰμ τοῦ κριτηρίου.

\vs{45}Καὶ ὁ οἶκος αὐτῶν ἐν ᾧ καθήσεται ἐκεῖ, αὐλὴ μία ἐξελισσομένη τούτοις κατὰ τὸ ἔργον τοῦτο· Καὶ οἶκον τῇ θυγατρὶ Φαραὼ ἣν ἔλαβε Σαλωμὼν, κατὰ τὸ αἰλὰμ τοῦτο.

\vs{46}Πάντα ταῦτα ἐκ λίθων τιμίων κεκολαμμένα ἐκ διαστήματος ἔσωθεν καὶ ἐκ τοῦ θεμελίου ἕως τῶν γεισῶν· καὶ ἔξωθεν εἰς τὴν αὐλὴν τὴν μεγάλην,
\vs{47}τὴν τεθεμελιωμένην ἐν τιμίοις λίθοις μεγάλοις, λίθοις δεκαπήχεσι καὶ τοῖς ὀκταπήχεσι·
\vs{48}Καὶ ἐπάνωθεν τιμίοις κατὰ τὸ μέτρον ἀπελεκήτων, καὶ κέδροις.
\vs{49}Τῆς αὐλῆς τῆς μεγάλης κύκλῳ τρεῖς στίχοι ἀπελεκήτων, καὶ στίχος κεκολαμμένης κέδρου·
\vs{50}καὶ συνετέλεσε Σαλωμὼν ὅλον τὸν οἶκον αὐτοῦ.

\ch{8}
Καὶ ἐγένετο ὡς συνετέλεσε Σαλωμὼν τοῦ οἰκοδομῆσαι τὸν οἶκον Κυρίου καὶ τὸν οἶκον αὐτοῦ μετὰ εἴκοσι ἔτη, τότε ἐξεκκλησίασεν ὁ βασιλεὺς Σαλωμὼν πάντας τοὺς πρεσβυτέρους Ἰσραὴλ ἐν Σιὼν, τοῦ ἐνεγκεῖν τὴν κιβωτὸν διαθήκης Κυρίου ἐκ πόλεως Δαυὶδ, αὕτη ἐστὶ Σιὼν,
\vs{2}ἐν μηνὶ Ἀθανίν.

\vs{3}Καὶ ᾖραν οἱ ἱερεῖς τὴν κιβωτὸν
\vs{4}καὶ τὸ σκήνωμα τοῦ μαρτυρίου καὶ τὰ σκεύη τὰ ἅγια τὰ ἐν τῷ σκηνώματι τοῦ μαρτυρίου.
\vs{5}Καὶ ὁ βασιλεὺς καὶ πὰς Ἰσραὴλ ἔμπροσθεν τῆς κιβωτοῦ θύοντες πρόβατα, βόας, ἀναρίθμητα·
\vs{6}Καὶ εἰσφέρουσιν οἱ ἱερεῖς τὴν κιβωτὸν εἰς τὸν τόπον αὐτῆς, εἰς τὸ δαβὶρ τοῦ οἴκου, εἰς τὰ ἅγια τῶν ἁγίων, ὑπὸ τὰς πτέρυγας τῶν χερουβίν.
\vs{7}Ὅτι τὰ χερουβὶμ διαπεπετασμένα ταῖς πτέρυξιν ἐπὶ τὸν τόπον τῆς κιβωτοῦ· καὶ περιεκάλυπτον τὰ χερουβὶμ ἐπὶ τὴν κιβωτὸν καὶ ἐπὶ τὰ ἅγια αὐτῆς ἐπάνωθεν.
\vs{8}Καὶ ὑπερεῖχον τὰ ἡγιασμένα· καὶ ἐνεβλέποντο αἱ κεφαλαὶ τῶν ἡγιασμένων ἐκ τῶν ἁγίων εἰς πρόσωπον τοῦ δαβὶρ, καὶ οὐκ ὠπτάνοντο ἔξω.
\vs{9}Οὐκ ἦν ἐν τῇ κιβωτῷ πλὴν δύο πλάκες λίθιναι, πλάκες τῆς διαθήκης ἃς ἔθηκε Μωυσῆς ἐν Χωρὴβ, ἃς διέθετο Κύριος μετὰ τῶν υἱῶν Ἰσραὴλ ἐν τῷ ἐκπορεύεσθαι αὐτοὺς ἐκ γῆς Αἰγύπτου.

\vs{10}Καὶ ἐγένετο ὡς ἐξῆλθον οἱ ἱερεῖς ἐκ τοῦ ἁγίου, καὶ ἡ νεφέλη ἔπλησε τὸν οἶκον.
\vs{11}Καὶ οὐκ ἠδύναντο οἱ ἱερεῖς στήκειν λειτουργεῖν ἀπὸ προσώπου τῆς νεφέλης, ὅτι ἔπλησε δόξα Κυρίου τὸν οἶκον.

\vs{14}Καὶ ἀπέστρεψεν ὁ βασιλεὺς τὸ πρόσωπον αὐτοῦ, καὶ εὐλόγησεν ὁ βασιλεὺς πάντα Ἰσραήλ· καὶ πᾶσα ἐκκλησία Ἰσραὴλ εἱστήκει·
\vs{15}Καὶ εἶπεν, εὐλογητὸς Κύριος ὁ Θεὸς Ἰσραὴλ σημερον, ὃς ἐλάλησεν ἐν τῷ στόματι αὐτοῦ περὶ Δαυὶδ τοῦ πατρός μου καὶ ἐν ταῖς χερσὶν αὐτοῦ ἐπλήρωσε, λέγον,
\vs{16}ἀφʼ ἧς ἡμέρας ἐξήγαγον τὸν λαόν μου τὸν Ἰσραὴλ ἐξ Αἰγύπτου, οὐκ ἐξελεξάμην ἐν πόλει ἐν ἑνὶ σκήπτρῳ Ἰσραὴλ τοῦ οἰκοδομῆσαι οἶκον τοῦ εἶναι τὸ ὄνομά μου ἐκεῖ· καὶ ἐξελεξάμην ἐν Ἱερουσαλὴμ εἶναι τὸ ὄνομά μου ἐκεῖ· καὶ ἐξελεξάμην τὸν Δαυὶδ τοῦ εἶναι ἐπὶ τὸν λαόν μου τὸν Ἰσραήλ.
\vs{17}Καὶ ἐγένετο ἐπὶ τῆς καρδίας τοῦ πατρός μου οἰκοδομῆσαι οἶκον τῷ ὀνόματι Κυρίου Θεοῦ Ἰσραήλ.
\vs{18}Καὶ εἶπε Κύριος πρὸς Δαυὶδ τὸν πατέρα μου, ἀνθʼ ὧν ἦλθεν ἐπὶ τὴν καρδίαν σου τοῦ οἰκοδομῆσαι οἶκον τῷ ὀνόματί μου, καλῶς ἐποίησας ὅτι ἐγενήθη ἐπὶ τὴν καρδίαν σου.
\vs{19}Πλὴν σὺ οὐκ οἰκοδομήσεις τὸν οἶκον, ἀλλʼ ἢ ὁ υἱός σου ὁ ἐξελθὼν ἐκ τῶν πλευρῶν σου, οὗτος οἰκοδομήσει τὸν οἶκον τῷ ὀνόματί μου.
\vs{20}Καὶ ἀνέστησε Κύριος τὸ ῥῆμα αὐτοῦ ὃ ἐλάλησε· καὶ ἀνέστην ἀντὶ Δαυὶδ τοῦ πατρός μου, καὶ ἐκάθισα ἐπὶ τοῦ θρόνου Ἰσραὴλ, καθὼς ἐλάλησε Κύριος, καὶ ᾠκοδόμησα τὸν οἶκον τῷ ὀνόματι Κυρίου Θεοῦ Ἰσραήλ.
\vs{21}Καὶ ἐθέμην ἐκεῖ τόπον τῇ κιβωτῷ, ἐν ᾗ ἐστιν ἐκεῖ διαθήκη Κυρίου ἣν διέθετο Κύριος μετὰ τῶν πατέρων ἡμῶν ἐν τῷ ἐξαγαγεῖν αὐτὸν αὐτοὺς ἐκ γῆς Αἰγύπτου.

\vs{22}Καὶ ἀνέστη Σαλωμὼν κατὰ πρόσωπον τοῦ θυσιαστηρίου Κυρίου ἐνώπιον πάσης ἐκκλησίας Ἰσραήλ· καὶ διεπέτασε τὰς χεῖρας αὐτοῦ εἰς τὸν οὐρανὸν,
\vs{23}καὶ εἶπε, Κύριε ὁ Θεὸς Ἰσραὴλ, οὐκ ἔστιν ὡς σὺ Θεὸς ἐν τῷ οὐρανῷ ἄνω καὶ ἐπὶ τῆς γῆς κάτω, φυλάσσων διαθήκην καὶ ἔλεος τῷ δούλῳ σου τῷ πορευομένῳ ἐνώπιόν σου ἐν ὅλῃ τῇ καρδίᾳ αὐτοῦ,
\vs{24}ἃ ἐφύλαξας τῷ δούλῳ σου Δαυὶδ τῷ πατρί μου· καὶ γὰρ ἐλάλησας ἐν τῷ στόματί σου, καὶ ἐν χερσί σου ἐπλήρωσας, ὡς ἡ ἡμέρα αὕτη.
\vs{25}Καὶ νῦν Κύριε ὁ Θεὸς Ἰσραὴλ, φύλαξον τῷ δούλῳ σου Δαυὶδ τῷ πατρί μου ἃ ἐλάλησας αὐτῷ, λέγων, οὐκ ἐξαρθήσεταί σου ἀνὴρ ἐκ προσώπου μου καθήμενος ἐπὶ θρόνου Ἰσραὴλ, πλὴν ἐὰν φυλάξωνται τὰ τέκνα σου τὰς ὁδοὺς αὐτῶν τοῦ πορεύεσθαι ἐνώπιόν μου καθὼς ἐπορεύθης ἐνώπιον ἐμοῦ.
\vs{26}Καὶ νῦν, Κύριε ὁ Θεὸς Ἰσραὴλ, πιστωθήτω δὴ τὸ ῥῆμά σου τῷ Δαυὶδ τῷ πατρί μου.

\vs{27}Ὅτι εἰ ἀληθῶς κατοικήσει ὁ Θεὸς μετὰ ἀνθρώπων ἐπὶ τῆς γῆς; εἰ ὁ οὐρανὸς καὶ ὁ οὐρανὸς τοῦ οὐρανοῦ οὐκ ἀρκέσουσί σοι, πλὴν καὶ ὁ οἶκος οὗτος ὃν ᾠκοδόμησα τῷ ὀνόματί σου;
\vs{28}Καὶ ἐπιβλέψῃ ἐπὶ τὴν δέησίν μου Κύριε ὁ Θεὸς Ἰσραὴλ, ἀκούειν τῆς προσευχῆς ἧς ὁ δοῦλός σου προσεύχεται ἐνώπιόν σου πρὸς σὲ σήμερον,
\vs{29}τοῦ εἶναι τοὺς ὀφθαλμούς σου ἠνεῳγμένους εἰς τὸν οἶκον τοῦτον ἡμέρας καὶ νυκτὸς, εἰς τὸν τὸπον ὃν εἶπας, ἔσται τὸ ὄνομά μου ἐκεῖ, τοῦ εἰσακούειν τῆς προσευχῆς ἧς προσεύχεται ὁ δοῦλός σου εἰς τὸν τόπον τοῦτον ἡμέρας καὶ νυκτός.
\vs{30}Καὶ εἰσακούσῃ τῆς δεήσεως τοῦ δούλου σου καὶ τοῦ· λαοῦ σου Ἰσραὴλ ἃ ἂν προσεύξωνται εἰς τὸν τόπον τοῦτον· καὶ σὺ εἰσακούσῃ ἐν τῷ τόπῳ τῆς κατοικήσεώς σου ἐν οὐρανῷ· καὶ ποιήσεις καὶ ἵλεως ἔσῃ.

\vs{31}Ὅσα ἂν ἁμάρτῃ ἕκαστος τῷ πλησίον αὐτοῦ, καὶ ἐὰν λάβῃ ἐπʼ αὐτὸν ἀρὰν τοῦ ἀράσασθαι αὐτὸν, καὶ ἔλθῃ καὶ ἐξαγορεύσῃ κατὰ πρόσωπον τοῦ θυσιαστηρίου σου ἐν τῷ οἴκῳ τούτῳ,
\vs{32}καὶ σὺ εἰσακούσῃ ἐκ τοῦ οὐρανοῦ καὶ ποιήσεις· καὶ κρινεῖς τὸν λαόν σου Ἰσραὴλ, ἀνομηθῆναι ἄνομον, δοῦναι τὴν ὁδὸν αὐτοῦ εἰς κεφαλὴν αὐτοῦ, καὶ τοῦ δικαιῶσαι δίκαιον, δοῦναι αὐτῷ κατὰ τὴν δικαιοσύνην αὐτοῦ.

\vs{33}Ἐν τῷ πταῖσαι τὸν λαόν σου Ἰσραὴλ ἐνώπιον ἐχθρῶν, ὅτι ἁμαρτήσονταί σοι, καὶ ἐπιστρέψουσι καὶ ἐξομολογήσονται τῷ ὀνόματί σου, καὶ προσεύξονται καὶ δεηθήσονται ἐν τῷ οἴκῳ τούτῳ,
\vs{34}καὶ σὺ εἰσακούσῃ ἐκ τοῦ οὐρανοῦ, καὶ ἵλεως ἔσῃ ταῖς ἁμαρτίαις τοῦ λαοῦ σου Ἰσραὴλ, καὶ ἐπιστρέψεις αὐτοὺς εἰς τὴν γῆν ἣν ἔδωκας τοῖς πατράσιν αὐτῶν.

\vs{35}Ἐν τῷ συσχεθῆναι τὸν οὐρανὸν καὶ μὴ γενέσθαι ὑετὸν, ὅτι ἁμαρτήσονταί σοι, καὶ προσεύξονται εἰς τὸν τόπον τοῦτον, καὶ ἐξομολογήσονται τῷ ὀνόματί σου, καὶ ἀπὸ τῶν ἁμαρτιῶν αὐτῶν ἀποστρέψουσιν ὅταν ταπεινώσῃς αὐτοὺς,
\vs{36}καὶ εἰσακούσῃ ἐκ τοῦ οὐρανοῦ, καὶ ἵλεως ἔσῃ ταῖς ἁμαρτίαις τοῦ δούλου σου καὶ τοῦ λαοῦ σου Ἰσραήλ· ὅτι δηλώσεις αὐτοῖς τὴν ὁδὸν τὴν ἀγαθὴν πορεύεσθαι ἐν αὐτῇ, καὶ δώσεις ὑετὸν ἐπὶ τὴν γῆν ἣν ἔδωκας τῷ λαῳ σου ἐν κληρονομίᾳ.

\vs{37}Λιμὸς ἐὰν γένηται, θάνατος ἐὰν γένηται, ὅτι ἔσται ἐμπυρισμὸς, βροῦχος, ἐρυσίβη ἐὰν γένηται, καὶ ἐὰν θλίψῃ αὐτὸν ὁ ἐχθρὸς αὐτοῦ ἐν μιᾷ τῶν πόλεων αὐτοῦ, πᾶν συνάντημα, πᾶν πόνον,
\vs{38}πᾶσαν προσευχὴν, πᾶσαν δέησιν ἐὰν γένηται παντὶ ἀνθρώπῳ, ὡς ἂν γνῶσιν ἕκαστος ἁφὴν καρδίας αὐτοῦ, καὶ διαπετάσῃ τὰς χεῖρας αὐτοῦ εἰς τὸν οἶκον τοῦτον,
\vs{39}καὶ σὺ εἰσακούσῃ ἐκ τοῦ οὐρανοῦ ἐξ ἑτοίμου κατοικητηρίου σου, καὶ ἵλεως ἔσῃ, καὶ ποιήσεις καὶ δώσεις ἀνδρὶ κατὰ τὰς ὁδοὺς αὐτοῦ, καθὼς ἂν γνῷς τὴν καρδίαν αὐτοῦ, ὅτι σὺ μονώτατος οἶδας τὴν καρδίαν πάντων υἱῶν ἀνθρώπων,
\vs{40}ὅπως φοβῶνταί σε πάσας τὰς ἡμέρας ὅσας αὐτοὶ ζῶσιν ἐπὶ τῆς γῆς, ἧς ἔδωκας τοῖς πατράσιν ἡμῶν.

\vs{41}Καὶ τῷ ἀλλοτρίῳ ὃς οὐκ ἔστιν ἀπὸ λαοῦ σοῦ οὗτος,
\vs{42}καὶ ἥξουσι καὶ προσεύξονται εἰς τὸν τόπον τοῦτον,
\vs{43}καὶ σὺ εἰσακούσῃ ἐκ τοῦ οὐρανοῦ ἐξ ἑτοίμου κατοικητηρίου σου, καὶ ποιήσεις κατὰ πάντα ὅσα ἂν ἐπικαλέσηταί σε ὁ ἀλλότριος, ὅπως γνῶσι πάντες οἱ λαοὶ τὸ ὄνομά σου, καὶ φοβῶνταί σε, καθὼς ὁ λαός σου Ἰσραὴλ, καὶ γνῶσιν ὅτι τὸ ὄνομά σου ἐπικέκληται ἐπὶ τὸν οἶκον τοῦτον ὃν ᾠκοδόμησα.

\vs{44}Ὅτι ἐξελεύσεται ὁ λαός σου εἰς πόλεμον ἐπὶ τοὺς ἐχθροὺς αὐτοῦ ἐν ὁδῷ ᾗ ἐπιστρέψεις αὐτοὺς, καὶ προσεύξονται ἐν ὀνόματι Κυρίου ὁδὸν τῆς πόλεως ἧς ἑξελέξω ἐν αὐτῇ, καὶ τοῦ οἴκου οὗ ᾠκοδόμησα τῷ ὀνόματί σου,
\vs{45}καὶ σὺ εἰσακούσῃ ἐκ τοῦ οὐρανοῦ τῆς δεήσεως αὐτῶν, καὶ τῆς προσευχῆς αὐτῶν, καὶ ποιήσεις τὸ δικαίωμα αὐτοῖς.

\vs{46}Ὅτι ἁμαρτήσονταί σοι, ὅτι οὐκ ἔστιν ἄνθρωπος ὃς οὐχ ἁμαρτήσεται, καὶ ἐπάξεις αὐτοὺς καὶ παραδώσεις αὐτοὺς ἐνώπιον ἐχθρῶν, καὶ αἰχμαλωτιοῦσιν οἱ αἰχμαλωτίζοντες εἰς γῆν μακρὰν ἢ ἐγγὺς,
\vs{47}καὶ ἐπιστρέψουσι καρδίας αὐτῶν ἐν τῇ γῇ οὗ μετήχθησαν ἐκεῖ, καὶ ἐπιστρέψωσιν ἐν γῇ μετοικίας αὐτῶν, καὶ δεηθῶσί σου, λέγοντες, ἡμάρτομεν, ἠδικήσαμεν, ἠνομήσαμεν,
\vs{48}καὶ ἐπιστρέψωσι πρὸς σὲ ἐν ὅλῃ καρδίᾳ αὐτῶν καὶ ἐν ὅλῃ ψυχῇ αὐτῶν ἐν τῇ γῇ ἐχθρῶν αὐτῶν οὗ μετήγαγες αὐτοὺς, καὶ προσεύξονται πρὸς σὲ ὁδὸν γῆς αὐτῶν ἧς ἔδωκας τοῖς πατράσιν αὐτῶν, καὶ τῆς πόλεως ἧς ἐξελέξω, καὶ τοῦ οἴκου οὗ ᾠκοδόμηκα τῷ ὀνόματί σου,
\vs{49}καὶ εἰσακούσῃ ἐκ τοῦ οὐρανοῦ ἐξ ἑτοίμου κατοικητηρίου σου,
\vs{50}καὶ ἵλεως ἔσῃ ταῖς ἀδικίας αὐτῶν αἷς ἥμαρτόν σοι, καὶ κατὰ πάντα τὰ ἀθετήματα αὐτῶν ἃ ἠθέτησάν σοι, καὶ δώσεις αὐτοὺς εἰς οἰκτιρμοὺς ἐνώπιον αἰχμαλωτευόντων αὐτοὺς, καὶ οἰκτειρήσουσιν εἰς αὐτοὺς,
\vs{51}ὅτι λαός σου καὶ κληρονομία σου, οὓς ἐξήγαγες ἐκ γῆς Αἰγύπτου ἐκ μέσου χωνευτηρίου σιδήρου.
\vs{52}Καὶ ἔστωσαν οἱ ὀφθαλμοί σου καὶ τὰ ὦτά σου ἠνεῳγμένα εἰς τὴν δέησιν τοῦ δούλου σου, καὶ εἰς τὴν δέησιν τοῦ λαοῦ σου Ἰσραὴλ, εἰσακούειν αὐτῶν ἐν πᾶσιν οἷς ἂν ἐπικαλέσωνταί σε.
\vs{53}Ὅτι σὺ διέστειλας αὐτοὺς σεαυτῷ εἰς κληρονομίαν ἐκ πάντων τῶν λαῶν τῆς γῆς, καθὼς ἐλάλησας ἐν χειρὶ δούλου σου Μωυσῆ, ἐν τῷ ἐξαγαγεῖν σε τοὺς πατέρας ἡμῶν ἐκ γῆς Αἰγύπτου, Κύριε Κύριε.

\vs{53a}Τότε ἐλάλησε Σαλωμὼν ὑπὲρ τοῦ οἴκου, ὡς συνετέλεσε τοῦ οἰκοδομῆσαι αὐτὸν, Ἥλιον ἐγνώρισεν ἐν οὐρανῷ· Κύριος εἶπε τοῦ κατοικεῖν ἐν γνόφῳ· οἰκοδόμησον οἶκόν μου, οἶκον εὐπρεπῆ σεαυτῷ τοῦ κατοικεῖν ἐπὶ καινότητος· οὐκ ἰδοὺ αὕτη γέγραπται ἐν βιβλίῳ τῆς ᾠδῆς;

\vs{54}Καὶ ἐγένετο ὡς συνετέλεσε Σαλωμὼν προσευχόμενος πρὸς Κύριον ὅλην τὴν προσευχὴν καὶ τὴν δέησιν ταύτην, καὶ ἀνέστη ἀπὸ προσώπου τοῦ θυσιαστηρίου Κυρίου ὀκλακὼς ἐπὶ τὰ γόνατα αὐτοῦ, καὶ αἱ χεῖρες αὐτοῦ διαπεπετασμέναι εἰς τὸν οὐρανόν.

\vs{55}Καὶ ἔστη, καὶ εὐλόγησε πᾶσαν ἐκκλησίαν Ἰσραὴλ φωνῇ μεγάλῃ, λέγων,
\vs{56}εὐλογητὸς Κύριος σήμερον ὃς ἔδωκε κατάπαυσιν τῷ λαῷ αὐτοῦ Ἰσραὴλ, κατὰ πάντα ὅσα ἐλάλησεν· οὐ διεφώνησε λόγος εἷς ἐν πᾶσι τοῖς λόγοις αὐτοῦ τοῖς ἀγαθοῖς οἷς ἐλάλησεν ἐν χειρὶ δούλου αὐτοῦ Μωυσῆ.
\vs{57}Γένοιτο Κύριος ὁ Θεὸς ἡμῶν μεθʼ ἡμῶν, καθὼς ἦν μετὰ τῶν πατέρων ἡμῶν· μὴ ἐγκαταλοίποιτο ἡμᾶς μηδὲ ἀποστρέψοιτο ἡμᾶς,
\vs{58}Γένοιτο Κύριος ὁ Θεὸς ἡμῶν μεθʼ ἡμῶν, καθὼς ἦν μετὰ τῶν πατέρων ἡμῶν· μὴ ἐγκαταλοίποιτο ἡμᾶς μηδὲ ἀποστρέψοιτο ἡμᾶς, ἐπικλῖναι καρδίας ἡμῶν ἐπʼ αὐτὸν τοῦ πορεύεσθαι ἐν πάσαις ὁδοῖς αὐτοῦ, καὶ φυλάσσειν πάσας ἐντολὰς αὐτοῦ, καὶ τὰ προστάγματα αὐτοῦ, ἃ ἐνετείλατο τοῖς πατράσιν ἡμῶν.
\vs{59}Καὶ ἔστωσαν οἱ λόγοι οὗτοι ὡς δεδέημαι ἐνώπιον Κυρίου Θεοῦ ἡμῶν, ἐγγίζοντες πρὸς Κύριον Θεὸν ἡμῶν ἡμέρας καὶ νυκτὸς, τοῦ ποιεῖν τὸ δικαίωμα τοῦ δοῦλου σου, καὶ τὸ δικαίωμα λαοῦ Ἰσραὴλ ῥῆμα ἡμέρας ἐν ἡμέρᾳ ἐνιαυτοῦ·
\vs{60}ὅπως γνῶσι πάντες οἱ λαοὶ τῆς γῆς, ὅτι Κύριος ὁ Θεὸς, αὐτὸς Θεὸς, καὶ οὐκ ἔστιν ἔτι.
\vs{61}Καὶ ἔστωσαν αἱ καρδίαι ἡμῶν τέλειαι πρὸς Κύριον Θεὸν ἡμῶν, καὶ ὁσίως πορεύεσθαι ἐν τοῖς προστάγμασιν αὐτοῦ, καὶ φυλάσσειν ἐντολὰς αὐτοῦ, ὡς ἡ ἡμέρα αὕτη.

\vs{62}Καὶ ὁ βασιλεὺς καὶ πάντες οἱ υἱοὶ Ἰσραὴλ ἔθυσαν θυσίαν ἐνώπιον Κυρίου.
\vs{63}Καὶ ἔθυσεν ὁ βασιλεὺς Σαλωμὼν τὰς θυσίας τῶν εἰρηνικῶν ἃς ἔθυσε τῷ Κυρίῳ, βοῶν δύο καὶ εἴκοσι χιλιάδας, προβάτων ἑκατὸν καὶ εἴκοσι χιλιάδας· καὶ ἐνεκαίνισε τὸν οἶκον Κυρίου ὁ βασιλεὺς καὶ πάντες οἱ υἱοὶ Ἰσραήλ.
\vs{64}Τῇ ἡμέρᾳ ἐκείνῃ ἡγίασεν ὁ βασιλεὺς τὸ μέσον τῆς αὐλῆς τὸ κατὰ πρόσωπον τοῦ οἴκου Κυρίου· ὅτι ἐποίησεν ἐκεῖ τὴν ὁλοκαύτωσιν καὶ τὰς θυσίας καὶ τὰ στέατα τῶν εἰρηνικῶν, ὅτι τὸ θυσιαστήριον τὸ χαλκοῦν τὸ ἐνώπιον Κυρίου μικρὸν τοῦ μὴ δύνασθαι τὴν ὁλοκαύτωσιν καὶ τὰς θυσίας τῶν εἰρηνικῶν ὑπενεγκεῖν.

\vs{65}Καὶ ἑποίησε Σαλωμὼν τὴν ἑορτὴν ἐν τῇ ἡμέρᾳ ἐκείνῃ, καὶ πᾶς Ἰσραὴλ μετʼ αὐτοῦ, ἐκκλησία μεγάλη ἀπὸ τῆς εἰσόδου Ἡμὰθ ἕως ποταμοῦ Αἰγύπτου, ἐνώπιον Κυρίου Θεοῦ ἡμῶν ἐν τῷ οἴκῳ ᾧ ᾠκοδόμησεν, ἐσθίων καὶ πίνων καὶ εὐφραινόμενος ἐνώπιον Κυρίου Θεοῦ ἡμῶν ἑπτὰ ἡμέρας.
\vs{66}Καὶ ἐν τῇ ἡμέρᾳ τῇ ὀγδόῃ ἐξαπέστειλε τὸν λαόν· καὶ εὐλόγησαν τὸν βασιλέα, καὶ ἀπῆλθεν ἕκαστος εἰς τὰ σκηνώματα αὐτοῦ χαίροντες· καὶ ἀγαθὴ ἡ καρδία ἐπὶ τοῖς ἀγαθοῖς οἷς ἐποίησε Κύριος τῷ Δαυὶδ δούλῳ αὐτοῦ, καὶ τῷ Ἰσραὴλ λαῷ αὐτοῦ.

\ch{9}
Καὶ ἐγενήθη ὡς συνετέλεσε Σαλωμὼν οἰκοδομεῖν τὸν οἶκον Κυρίου, καὶ τὸν οἶκον τοῦ βασιλέως, καὶ πᾶσαν τὴν πραγματείαν Σαλωμὼν, ὅσα ἠθέλησε ποιῆσαι,
\vs{2}καὶ ὤφθη Κύριος τῷ Σαλωμὼν δεύτερον, καθὼς ὤφθη ἐν Γαβαών.

\vs{3}Καὶ εἶπε πρὸς αὐτὸν Κύριος, ἤκουσα τῆς φωνῆς τῆς προσευχῆς σου, καὶ τῆς δεήσεώς σου ἧς ἐδεήθης ἐνώπιόν μου· πεποίηκά σοι κατὰ πᾶσαν τὴν προσευχήν σου· ἡγίακα τὸν οἶκον τοῦτον ὃν ᾠκοδόμησας τοῦ θέσθαι τὸ ὄνομά μου ἐκεῖ εἰς τὸν αἰῶνα, καὶ ἔσονται οἱ ὀφθαλμοί μου ἐκεῖ καὶ ἡ καρδία μου πάσας τὰς ἡμέρας.
\vs{4}Καὶ σὺ ἐὰν πορευθῇς ἐνώπιον ἐμοῦ, καθὼς ἐπορεύθη Δαυὶδ ὁ πατήρ σου, ἐν ὁσιότητι καρδίας καὶ ἐν εὐθύτητι, καὶ τοῦ ποιεῖν κατὰ πάντα ἃ ἐνετειλάμην αὐτῷ, καὶ τὰ προστάγματά μου καὶ τὰς ἐντολάς μου φυλάξῃς,
\vs{5}καὶ ἀναστήσω τὸν θρόνον τῆς βασιλείας σου ἐν Ἰσραὴλ εἰς τὸν αἰῶνα, καθὼς ἐλάλησα Δαυὶδ πατρί σου, λέγων, οὐκ ἐξαρθήσεταί σοι ἀνὴρ ἡγούμενος ἐν Ἰσραήλ.
\vs{6}Ἐὰν δὲ ἀποστραφέντες ἀποστραφῆτε ὑμεῖς καὶ τὰ τέκνα ὑμῶν ἀπʼ ἐμοῦ, καὶ μὴ φυλάξητε τὰς ἐντολάς μου καὶ τὰ προστάγματά μου ἃ ἔδωκε Μωυσῆς ἐνώπιον ὑμῶν, καὶ πορευθῆτε καὶ δουλεύσητε θεοῖς ἑτέροις καὶ προσκυνήσητε αὐτοῖς,
\vs{7}καὶ ἐξαρῶ τὸν Ἰσραὴλ ἀπὸ τῆς γῆς ἧς ἔδωκα αὐτοῖς, καὶ τὸν οἶκον τοῦτον ὃν ἡγίασα τῷ ὀνόματί μου ἀποῤῥίψω ἐκ προσώπου μου· καὶ ἔσται Ἰσραὴλ εἰς ἀφανισμὸν καὶ εἰς λάλημα εἰς πάντας τοὺς λαούς.
\vs{8}Καὶ ὁ οἶκος οὗτος ἔσται ὁ ὑψηλὸς, πᾶς ὁ διαπορευόμενος διʼ αὐτοῦ ἐκστήσεται καὶ συριεῖ, καὶ ἐροῦσιν, ἕνεκεν τίνος ἐποίησε Κύριος οὕτως τῇ γῇ ταύτῃ καὶ τῷ οἴκῳ τούτῳ;
\vs{9}Καὶ ἐροῦσιν, ἀνθʼ ὧν ἐγκατέλιπον Κύριον Θεὸν αὐτῶν, ὃς ἐξήγαγε τοὺς πατέρας αὐτῶν ἐξ Αἰγύπτου, ἐξ οἴκου δουλείας, καὶ ἀντελάβοντο θεῶν ἀλλοτρίων καὶ προσεκύνησαν αὐτοῖς καὶ ἐδούλευσαν αὐτοῖς, διὰ τοῦτο ἐπήγαγε Κύριος ἐπʼ αὐτοὺς τὴν κακίαν ταύτην.

\vs{9a}Τότε ἀνήγαγε Σαλωμὼν τὴν θυγατέρα Φαραὼ ἐκ πόλεως Δαυὶδ εἰς οἶκον αὐτοῦ, ὃν ᾠκοδόμησεν ἐαυτῷ ἐν ταῖς ἡμέραις ἐκείναις.

\vs{10}Εἴκοσι ἔτη ἐν οἷς ᾠκοδόμησε Σαλωμὼν τοὺς δύο οἴκους, τὸν οἶκον Κυρίου καὶ τὸν οἶκον τοῦ βασιλέως,
\vs{11}Χιρὰμ βασιλεὺς Τύρου ἀντελάβετο τοῦ Σαλωμὼν ἐν ξύλοις κεδρίνοις, καὶ ἐν ξύλοις πευκίνοις, καὶ ἐν χρυσίῳ, καὶ ἐν χρυσίῳ, καὶ ἐν παντὶ θελήματι αὐτοῦ· τότε ἔδωκεν ὁ βασιλεὺς τῷ Χιρὰμ εἴκοσι πόλεις ἐν τῇ γῇ τῇ Γαλιλαίᾳ.
\vs{12}Καὶ ἐξήλθε Χιρὰμ ἐκ Τύρου, καὶ ἐπορεύθη εἰς τὴν Γαλιλαίαν τοῦ ἰδεῖν τὰς πόλεις ἃς ἔδωκεν αὐτῷ Σαλωμών· καὶ οὐκ ἤρεσαν αὐτῷ.
\vs{13}Καὶ εἶπε, τί αἱ πόλεις αὗται ἃς ἔδωκάς μοι ἀδελφέ; καὶ ἐκάλεσεν αὐτὰς Ὅριον ἕως τῆς ἡμέρας ταύτης.
\vs{14}Καὶ ἤνεγκε Χιρὰμ τῷ Σαλωμὼν ἑκατὸν καὶ εἴκοσι τάλαντα χρυσίου.
\vs{26}Καὶ ναῦν ὑπὲρ οὗ ἐποίησεν ὁ βασιλεὺς Σαλωμὼν ἐν Γασίων Γαβὲρ τὴν οὖσαν ἐχομένην Αἰλὰθ ἐπὶ τοῦ χείλους τῆς ἐσχάτης θαλάσσης ἐν γῇ Ἐδώμ.
\vs{27}Καὶ ἀπέστειλε Χιρὰμ ἐν τῇ νηῒ τῶν παίδων αὐτοῦ ἄνδρας ναυτικοὺς ἐλαύνειν εἰδότας θάλασσαν μετὰ τῶν παίδων Σαλωμών.
\vs{28}Καὶ ἦλθον εἰς Σωφιρὰ, καὶ ἔλαβον ἐκεῖθεν χρυσίου ἑκατὸν καὶ εἴκοσι τάλαντα, καὶ ἤνεγκαν τῷ βασιλεῖ Σαλωμών.

\ch{10}
Καὶ βασίλισσα Σαβὰ ἤκουσε τὸ ὄνομα Σαλωμὼν καὶ τὸ ὄνομα Κυρίου, καὶ ἦλθε πειράσαι αὐτὸν ἐν αἰνίγμασι.
\vs{2}Καὶ ἦλθεν εἰς Ἰερουσαλὴμ ἐν δυνάμει βαρείᾳ σφόδρα· καὶ κάμηλοι αἴρουσαι ἡδύσματα καὶ χρυσὸν πολὺν σφόδρα καὶ λίθον τίμιον· καὶ εἰσῆλθε πρὸς Σαλωμὼν, καὶ ἐλάλησεν αὐτῷ πάντα ὅσα ἦν ἐν τῇ καρδίᾳ αὐτῆς.
\vs{3}Καὶ ἀπήγγειλεν αὐτῇ Σαλωμὼν πάντας τοὺς λόγους αὐτῆς· οὐκ ἦν λόγος παρεωραμένος παρὰ τοῦ βασιλέως, ὃν οὐκ ἀπήγγειλεν αὐτῇ.
\vs{4}Καὶ εἶδε βασίλισσα Σαβὰ πᾶσαν τὴν φρόνησιν Σαλωμὼν, καὶ τὸν οἶκον ὃν ᾠκοδόμησε,
\vs{5}καὶ τὰ βρώματα Σαλωμὼν, καὶ τὴν καθέδραν παίδων αὐτοῦ, καὶ τὴν στάσιν λειτουργῶν αὐτοῦ, καὶ τὸν ἱματισμὸν αὐτοῦ, καὶ τοὺς οἰνοχόους αὐτοῦ, καὶ τὴν ὁλοκαύτωσιν αὐτοῦ ἣν ἀνέφερεν ἐν οἴκῳ Κυρίου, καὶ ἐξ ἑαυτῆς ἐγένετο·
\vs{6}Καὶ εἶπε πρὸς τὸν βασιλέα Σαλωμὼν, ἀληθινὸς ὁ λόγος ὃν ἤκουσα ἐν τῇ γῇ μου περὶ τοῦ λόγου σου καὶ περὶ τῆς φρονήσεώς σου.
\vs{7}Καὶ οὐκ ἐπίστευσα τοῖς λαλοῦσί μοι, ἕως ὅτου παρεγενόμην καὶ ἑωράκασιν οἱ ὀφθαλμοί μου· καὶ ἰδοὺ οὐκ εἰσὶ τὸ ἥμισυ καθὼς ἀπήγγειλάν μοι· προστέθεικας ἀγαθὰ πρὸς αὐτὰ ἐπὶ πᾶσαν τὴν ἀκοὴν ἣν ἤκουσα ἐν τῇ γῇ μου.
\vs{8}Μακάριαι αἱ γυναῖκές σου, μακάριοι οἱ παῖδές σου οὗτοι οἱ παρεστηκότες ἐνώπιόν σου διόλου, οἱ ἀκούοντες πᾶσαν τὴν φρόνησίν σου.
\vs{9}Γένοιτο Κύριος ὁ Θεός σου εὐλογημένος, ὃς ἠθέλησεν ἐν σοὶ δοῦναί σε ἐπὶ θρόνου Ἰσραὴλ, διὰ τὸ ἀγαπᾷν Κύριον τὸν Ἰσραὴλ στῆσαι εἰς τὸν αἰῶνα· καὶ ἔθετό σε βασιλέα ἐπʼ αὐτοὺς, τοῦ ποιεῖν κρίμα ἐν δικαιοσύνῃ καὶ ἐν κρίμασιν αὐτῶν.

\vs{10}Καὶ ἔδωκε τῷ Σαλωμὼν ἑκατὸν εἴκοσι τάλαντα χρυσίου, καὶ ἡδύσματα πολλὰ σφόδρα, καὶ λίθον τίμιον· οὐκ ἐληλύθει κατὰ τὰ ἡδύσματα ἐκεῖνα ἔτι εἰς πλῆθος, ἃ ἔδωκε βασίλισσα Σαβὰ τῷ βασιλεῖ Σαλωμών.

\vs{11}Καὶ ἡ ναῦς Χιρὰμ ἡ αἴρουσα τὸ χρυσίον ἐκ Σουφὶρ, ἤνεγκε ξύλα πελεκητὰ πολλὰ σφόδρα καὶ λίθον τίμιον.
\vs{12}Καὶ ἐποίησεν ὁ βασιλεὺς τὰ ξύλα τὰ πελεκητὰ ὑποστηρίγματα τοῦ οἴκου Κυρίου καὶ τοῦ οἴκου τοῦ βασιλέως, καὶ νάβλας καὶ κινύρας τοῖς ᾠδοῖς· οὐκ ἐληλύθει τοιαῦτα ξύλα ἀπελέκητα ἐπὶ τῆς γῆς, οὐδὲ ὤφθησάν που ἕως τῆς ἡμέρας ταύτης.
\vs{13}Καὶ ὁ βασιλεὺς Σαλωμὼν ἔδωκε τῇ βασιλίσσῃ Σαβὰ πάντα ὅσα ἠθέλησεν, ὅσα ᾐτήσατο, ἐκτὸς πάντων ὧν ἐδεδώκει αὐτῇ διὰ χειρὸς τοῦ βασιλέως Σαλωμών· καὶ ἀπεστράφη, καὶ ἦλθεν εἰς τὴν γῆν αὐτῆς αὐτὴ, καὶ πάντες οἱ παῖδες αὐτῆς.

\vs{14}Καὶ ἦν ὁ σταθμὸς τοῦ χρυσίου τοῦ ἐληλυθότος τῷ Σαλωμὼν ἐν ἐνιαυτῷ ἑνὶ, ἑξακόσια καὶ ἑξηκονταὲξ τάλαντα χρυσίου,
\vs{15}χωρὶς τῶν φόρων τῶν ὑποτεταγμένων καὶ τῶν ἐμπόρων καὶ πάντων τῶν βασιλέων τοῦ πέραν καὶ τῶν σατραπῶν τῆς γῆς.

\vs{16}Καὶ ἐποίησε Σαλωμὼν τριακόσια δόρατα χρυσᾶ ἐλατά· τριακόσιοι χρυσοῖ ἐπῆσαν ἐπὶ τὸ δόρυ τὸ ἕν·
\vs{17}Καὶ τριακόσια ὅπλα χρυσᾶ ἐλατά· καὶ τρεῖς μναῖ ἐνῆσαν χρυσοῦ εἰς τὸ ὅπλον τὸ ἕν· καὶ ἔδωκεν αὐτὰ ὁ βασιλεὺς εἰς οἶκον δρυμοῦ τοῦ Λιβάνου.

\vs{18}Καὶ ἐποίησεν ὁ βασιλεὺς θρόνον ἐλεφάντινον μέγαν, καὶ περιεχρύσωσεν αὐτὸν χρυσίῳ δοκίμῳ.
\vs{19}Ἓξ ἀναβαθμοὶ τῷ θρόνῳ, καὶ προτομαὶ μόσχων τῷ θρόνῳ ἐκ τῶν ὀπίσω αὐτοῦ, καὶ χεῖρες ἔνθεν καὶ ἔνθεν ἐπὶ τοῦ τόπου τῆς καθέδρας, καὶ δύο λέοντες ἑστηκότες παρὰ τὰς χεῖρας,
\vs{20}καὶ δώδεκα λέοντες ἑστῶτες ἐκεῖ ἐπὶ τῶν ἓξ ἀναβαθμῶν ἔνθεν καὶ ἔνθεν· οὐ γέγονεν οὕτως πάσῃ βασιλείᾳ.
\vs{21}Καὶ πάντα τὰ σκεύη τὰ ὑπὸ τοῦ Σαλωμὼν γεγονότα χρυσᾶ, καὶ λουτῆρες χρυσοῖ, καὶ πάντα τὰ σκεύη οἴκου δρυμοῦ τοῦ Λιβάνου χρυσίῳ συγκεκλεισμένα. οὐκ ἦν ἀργύριον, ὅτι οὐκ ἦν λογιζόμενον ἐν ταῖς ἡμέραις Σαλωμών·
\vs{22}Ὅτι ναῦς Θαρσὶς τῷ βασιλεῖ Σαλωμὼν ἐν τῇ θαλάσσῃ μετὰ τῶν νηῶν Χιράμ· μία διὰ τριῶν ἐτῶν ἤρχετο τῷ βασιλεῖ ναῦς ἐκ Θαρσὶς χρυσίου καὶ ἀργυρίου καὶ λίθων τορευτῶν καὶ πελεκητῶν.

\vs{22a}Αὕτη ἦν ἡ πραγματεία τῆς προνομῆς ἧς ἀνήνεγκαν ὁ βασιλεὺς Σαλωμὼν οἰκοδομῆσαι τὸν οἶκον Κυρίου, καὶ τὸν οἶκον τοῦ βασιλέως, καὶ τὸ τεῖχος Ἱερουσαλὴμ, καὶ τὴν ἄκραν, τοῦ περιφράξαι τὸν φραγμὸν τῆς πόλεως Δαυὶδ, καὶ τὴν Ἀσσοὺρ, καὶ τὴν Μαγδὰν, καὶ τὴν Γαζὲρ, καὶ τὴν Βαιθωρὼν τὴν ἀνωτέρω, καὶ τὴν Ἰεθερμὰθ, καὶ πάσας τὰς πόλεις τῶν ἁρμάτων, καὶ πάσας τὰς πόλεις τῶν ἱππέων, καὶ τὴν πραγματείαν Σαλωμὼν, ἣν ἐπραγματεύσατο οἰκοδομῆσαι ἐν Ἱερουσαλὴμ καὶ ἐν πάσῃ τῇ γῇ, τοῦ μὴ κατάρξαι αὐτοῦ
\vs{22b}πάντα τὸν λαὸν τὸν ὑπολελειμμένον ὑπὸ τοῦ Χετταίου καὶ τοῦ Ἀμοῤῥαίου καὶ τοῦ Φερεζαίου καὶ τοῦ Χαναναίου καὶ τοῦ Εὐαίου καὶ τοῦ Ἰεβουσαίου καὶ τοῦ Γεργεσαίου, τῶν μὴ ἐκ τῶν υἱῶν Ἰσραὴλ ὄντων, τὰ τέκνα αὐτῶν τὰ ὑπολελειμμένα μετʼ αὐτοῦ ἐν τῇ γῇ, οὓς οὐκ ἐδύναντο οἱ υἱοὶ Ἰσραὴλ ἐξολοθρεῦσαι αὐτοὺς, καὶ ἀνήγαγεν αὐτοὺς Σαλωμὼν εἰς φόρον ἕως τῆς ἡμέρας ταύτης·
\vs{22c}καὶ ἐκ τῶν υἱῶν Ἰσραὴλ οὐκ ἔδωκε Σαλωμὼν πρᾶγμα, ὅτι αὐτοὶ ἦσαν ἄνδρες οἱ πολεμισταὶ, καὶ παῖδες αὐτοῦ καὶ ἄρχοντες καὶ τρισσοὶ αὐτοῦ, καὶ ἄρχοντες τῶν ἁρμάτων αὐτοῦ, καὶ ἱππεῖς αὐτοῦ.

\vs{23}Καὶ ἐμεγαλύνθη Σαλωμὼν ὑπὲρ πάντας τοὺς βασιλεῖς τῆς γῆς πλούτῳ καὶ φρονήσει.
\vs{24}Καὶ πάντες βασιλεῖς τῆς γῆς ἐζήτουν τὸ πρόσωπον Σαλωμὼν, τοῦ ἀκοῦσαι τῆς φρονήσεως αὐτοῦ ἧς ἔδωκε Κύριος τῇ καρδίᾳ αὐτοῦ.
\vs{25}Καὶ αὐτοὶ ἔφερον ἕκαστος τὰ δῶρα, σκεύη χρυσᾶ, καὶ ἱματισμὸν, στακτὴν, καὶ ἡδύσματα, καὶ ἵππους, καὶ ἡμιόνους τὸ κατʼ ἐνιαυτὸν ἐνιαυτῷ.
\vs{26}Καὶ ἦσαν τῷ Σαλωμὼν τέσσαρες χιλιάδες θήλειαι ἵπποι εἰς ἅρματα, καὶ δώδεκα χιλιάδες ἱππέων· καὶ ἔθετο αὐτὰς ἐν ταῖς πόλεσι τῶν ἁρμάτων καὶ μετὰ τοῦ βασιλέως ἐν Ἱερουσαλήμ·
\vs{26a}καὶ ἦν ἡγούμενος πάντων τῶν βασιλέων ἀπὸ τοῦ ποταμοῦ καὶ ἕως γῆς ἀλλοφύλων καὶ ἕως ὁρίων Αἰγύπτου.

\vs{27}Καὶ ἔδωκεν ὁ βασιλεὺς τὸ χρυσίον καὶ τὸ ἀργύριον ἐν Ἱερουσαλὴμ ὡς λίθους, καὶ τὰς κέδρους ἔδωκεν ὡς συκαμίνους τὰς ἐν τῇ πεδινῇ εἰς πλῆθος.
\vs{28}Καὶ ἡ ἔξοδος Σαλωμῶν τῶν ἱππέων καὶ ἐξ Αἰγύπτου, καὶ ἐκ Θεκουὲ ἔμποροι τοῦ βασιλέως· καὶ ἐλάμβανον ἐκ Θεκουὲ ἐν ἀλλάγματι.
\vs{29}Καὶ ἀνέβαινεν ἡ ἔξοδος ἐξ Αἰγύπτου ἅρμα ἀντὶ ἑκατὸν ἀργυρίου, καὶ ἵππος ἀντὶ πεντήκοντα ἀργυρίου· καὶ οὕτως πᾶσι τοῖς βασιλεῦσι Χεττιῒν, καὶ βασιλεῦσι Συρίας κατὰ θάλασσαν ἐξεπορεύοντο.

\ch{11}
Καὶ ὁ βασιλεὺς Σαλωμὼν ἦν φιλογύνης. Καὶ ἦσαν αὐτῷ γυναῖκες ἄρχουσαι ἑπτακόσιαι, καὶ παλλακαὶ τριακόσιαι. Καὶ ἔλαβε γυναῖκας ἀλλοτρίας, καὶ τὴν θυγατέρα Φαραὼ, Μωαβίτιδας, Ἀμμανίτιδας, Σύρας, καὶ Ἰδουμαίας, Χετταίας, καὶ Ἀμοῤῥαίας,
\vs{2}ἐκ τῶν ἐθνῶν ὧν ἀπεῖπε Κύριος τοῖς υἱοῖς Ἰσραὴλ, οὐκ εἰσελεύσεσθε εἰς αὐτοὺς, καὶ αὐτοὶ οὐκ εἰσελεύσονται εἰς ὑμᾶς, μὴ ἐκκλίνωσι τὰς καρδίας ὑμῶν ὀπίσω εἰδώλων αὐτῶν· εἰς αὐτοὺς ἐκολλήθη Σαλωμὼν τοῦ ἀγαπῆσαι.
\vs{4}Καὶ ἐγενήθη ἐν καιρῷ γήρους Σαλωμὼν, καὶ οὖκ ἦν ἡ καρδία αὐτοῦ τελεία μετὰ Κυρίου Θεοῦ αὐτοῦ, καθὼς ἡ καρδία Δαυὶδ τοῦ πατρὸς αὐτοῦ. Καὶ ἐξέκλιναν γυναῖκες αἱ ἀλλότριαι τὴν καρδίαν αὐτοῦ ὀπίσω θεῶν αὐτῶν.
\vs{5}Τότε ᾠκοδόμησε Σαλωμὼν ὑψηλὸν τῷ Χαμὼς εἰδώλῳ Μωὰβ, καὶ τῷ βασιλεῖ αὐτῶν εἰδώλῳ υἱῶν Ἀμμὼν,
\vs{6}καὶ τῇ Ἀστάρτῃ βδελύγματι Σιδωνίων.
\vs{7}Καὶ οὕτως ἐποίησε πάσαις ταῖς γυναιξὶν αὐτοῦ ταῖς ἀλλοτρίαις, αἳ ἐθυμίων καὶ ἔθυον τοῖς εἰδώλοις αὐτῶν,
\vs{8}καὶ ἐποίησε Σαλωμὼν τὸ πονηρὸν ἐνώπιον Κυρίου· οὐκ ἐπορεύθη ὀπίσω Κυρίου, ὡς Δαυὶδ ὁ πατὴρ αὐτοῦ.

\vs{9}Καὶ ὠργίσθη Κύριος ἐπὶ Σαλωμὼν, ὅτι ἐξέκλινε καρδίαν αὐτοῦ ἀπὸ Κυρίου Θεοῦ Ἰσραὴλ, τοῦ ὀφθέντος αὐτῷ δὶς,
\vs{10}καὶ ἐντειλαμένου αὐτῷ ὑπὲρ τοῦ λόγου τούτου, τὸ παράπαν μὴ πορευθῆναι ὀπίσω θεῶν ἑτέρων, καὶ φυλάξασθαι ποιῆσαι ἃ ἐνετείλατο αὐτῷ Κύριος ὁ Θεός· οὐδʼ ἦν ἡ καρδία αὐτοῦ τελεία μετὰ Κυρίου, κατὰ τὴν καρδίαν Δαυὶδ τοῦ πατρὸς αὐτοῦ.
\vs{11}Καὶ εἶπε Κύριος πρὸς Σαλωμὼν, ἀνθʼ ὧν ἐγένετο ταῦτα μετὰ σοῦ, καὶ οὐκ ἐφύλαξας τὰς ἐντολάς μου καὶ τὰ προστάγματά μου ἃ ἐνετειλάμην σοι, διαῤῥήσσων διαῤῥήξω τὴν βασιλείαν σου ἐκ χειρός σου, καὶ δώσω αὐτὴν τῷ δούλῳ σου.
\vs{12}Πλὴν ἐν ταῖς ἡμέραις σου οὐ ποιήσω αὐτὰ διὰ Δαυὶδ τὸν πατέρα σου· ἐκ χειρὸς υἱοῦ σου λήψομαι αὐτήν.
\vs{13}Πλὴν ὅλην τὴν βασιλείαν οὐ μὴ λάβω· σκῆπτρον ἓν δώσω τῷ υἱῷ σου διὰ Δαυὶδ τὸν δοῦλόν μου, καὶ διὰ Ἱερουσαλὴμ τὴν πόλιν ἣν ἐξελεξάμην.

\vs{14}Καὶ ἤγειρε Κύριος σατὰν τῷ Σαλωμὼν τὸν Ἄδερ τὸν Ἰδουμαῖον, καὶ τὸν Ἐσρὼμ υἱὸν Ἐλιαδαὲ τὸν ἐν Ρααμὰ, Ἀδαδέζερ βασιλέα Σουβὰ κύριον αὐτοῦ· καὶ συνηθροίσθησαν ἐπʼ αὐτὸν ἄνδρες, καὶ ἦν ἄρχων συστρέμματος, καὶ προκατελάβετο τὴν Δαμασέκ· καὶ ἦσαν σατὰν τῷ Ἰσραὴλ πάσας τὰς ἡμέρας Σαλωμών· καὶ Ἄδερ ὁ Ἰδουμαῖος ἐκ τοῦ σπέρματος τῆς βασιλείας ἐν Ἰδουμαίᾳ.
\vs{15}Καὶ ἐγένετο ἐν τῷ ἐξολοθρεῦσαι Δαυὶδ τὸν Ἐδὼμ ἐν τῷ πορευθῆναι Ἰωὰβ ἄρχοντα τῆς στρατιᾶς θάπτειν τοὺς τραυματίας, καὶ ἔκοψαν πᾶν ἀρσενικὸν ἐν τῇ Ἰδουμαίᾳ·
\vs{16}ὅτι ἓξ μῆνας ἐνεκάθητο ἐκεῖ Ἰωὰβ καὶ πᾶς Ἰσραὴλ ἐν τῇ Ἰδουμαίᾳ, ἕως ὅτου ἐξωλόθρευσε πᾶν ἀρσενικὸν ἐν τῇ Ἰδουμαίᾳ·
\vs{17}Καὶ ἀπέδρα Ἄδερ αὐτὸς καὶ πάντες ἄνδρες Ἰδουμαῖοι τῶν παίδων τοῦ πατρὸς αὐτοῦ μετʼ αὐτοῦ, καὶ εἰσῆλθον εἰς Αἴγυπτον· καὶ Ἄδερ παιδάριον μικρόν.
\vs{18}Καὶ ἀνίστανται ἄνδρες ἐκ τῆς πόλεως Μαδιὰμ, καὶ ἄρχονται εἰς Φαρὰν, καὶ λαμβάνουσιν ἄνδρας μεθʼ αὑτῶν, καὶ ἔρχονται πρὸς Φαραὼ βασιλέα Αἰγύπτου· καὶ εἰσῆλθεν Ἄδερ πρὸς Φαραὼ, καὶ ἔδωκεν αὐτῷ οἶκον, καὶ ἄρτους διέταξεν αὐτῷ.
\vs{19}Καὶ εὗρεν Ἄδερ χάριν ἐναντίον Φαραὼ σφόδρα, καὶ ἔδωκεν αὐτῷ γυναῖκα ἀδελφὴν τῆς γυναικὸς αὐτοῦ, ἀδελφὴν Θεκεμίνας μείζω. Καὶ ἔτεκεν αὐτῷ ἡ ἀδελφὴ Θεκεμίνας τῷ Ἄδερ τὸν Γανηβὰθ υἱὸν αὐτῆς·
\vs{20}καὶ ἐξέθρεψεν αὐτὸν Θεκεμίνα ἐν μέσῳ υἱῶν Φαραώ· καὶ ἦν Γανηβὰθ ἐν μέσῳ υἱῶν Φαραώ.

\vs{21}Καὶ Ἄδερ ἤκουσεν ἐν Αἰγύπτῳ ὅτι κεκοίμηται Δαυὶδ μετὰ τῶν πατέρων αὐτοῦ, καὶ ὅτι τέθνηκεν Ἰωὰβ ὁ ἄρχων τῆς στρατιᾶς, καὶ εἶπεν Ἄδερ πρὸς Φαραὼ, ἐξαπόστειλόν με, καὶ ἀποστρέψω εἰς τὴν γῆν μου.
\vs{22}Καὶ εἶπε Φαραὼ τῷ Ἄδερ, τίνι σὺ ἐλαττονῇ μετʼ ἐμοῦ; καὶ ἰδοὺ σὺ ζητεῖς ἀπελθεῖν εἰς τὴν γῆν σου; καὶ εἶπεν αὐτῷ Ἄδερ, ὅτι ἐξαποστέλλων ἐξαποστελεῖς με· καὶ ἀνέστρεψεν Ἄδερ εἰς τὴν γῆν αὐτοῦ·
\vs{25}αὕτη ἡ κακία ἣν ἐποίησεν Ἄδερ· καὶ ἐβαρυθύμησεν Ἰσραὴλ, καὶ ἐβασίλευσεν ἐν τῇ Ἐδώμ.

\vs{26}Καὶ Ἱεροβοὰμ υἱὸς Ναβὰτ ὁ Ἐφραθὶ ἐκ τῆς Σαριρὰ, υἱὸς γυναικὸς χήρας, δοῦλος Σαλωμών.
\vs{27}Καὶ τοῦτο τὸ πρᾶγμα ὡς ἐπῇρατο χεῖρας ἐπὶ βασιλέα Σαλωμών· καὶ ὁ βασιλεὺς Σαλωμὼν ᾠκοδόμησε τὴν ἄκραν, συνέκλεισε τὸν φραγμὸν τῆς πόλεως Δαυὶδ τοῦ πατρὸς αὐτοῦ.
\vs{28}Καὶ ὁ ἄνθρωπος Ἱεροβοὰμ ἰσχυρὸς δυνάμει· καὶ εἶδε Σαλωμὼν τὸ παιδάριον ὅτι ἀνὴρ ἔργων ἐστὶ, καὶ κατέστησεν αὐτὸν ἐπὶ τὰς ἄρσεις οἴκου Ἰωσήφ.

\vs{29}Καὶ ἐγενήθη ἐν τῷ καιρῷ ἐκείνῳ, καὶ Ἱεροβοὰμ ἐξῆλθεν ἐξ Ἱερουσαλὴμ, καὶ εὗρεν αὐτὸν Ἀχιὰ ὁ Σηλωνίτης ὁ προφήτης ἐν τῇ ὁδῷ, καὶ ἀπέστησεν αὐτὸν ἐκ τῆς ὁδοῦ· καὶ Ἀχιὰ περιβεβλημένος ἱματίῳ καινῷ, καὶ ἀμφότεροι μόνοι ἐν τῷ πεδίῳ.
\vs{30}Καὶ ἐπελάβετο Ἀχιὰ τοῦ ἱματίου αὐτοῦ τοῦ καινοῦ τοῦ ἐπʼ αὐτῷ, καὶ διέῤῥηξεν αὐτὸ δώδεκα ῥήγματα,
\vs{31}καὶ εἶπε τῷ Ἱεροβοὰμ, λάβε σεαυτῷ δέκα ῥήγματα, ὅτι τάδε λέγει Κύριος ὁ Θεὸς Ἰσραὴλ, ἰδοὺ ἐγὼ ῥήσσω τὴν βασιλείαν ἐκ χειρὸς Σαλωμὼν, καὶ δώσω σοι δέκα σκῆπτρα.
\vs{32}Καὶ δύο σκῆπτρα ἔσονται αὐτῷ διὰ τὸν δοῦλόν μου Δαυὶδ, καὶ διὰ Ἱερουσαλὴμ τὴν πόλιν ἣν ἐξελεξάμην ἐν αὐτῇ ἐκ πασῶν φυλῶν Ἰσραήλ.
\vs{33}Ἀνθʼ ὧν ἐγκατέλιπέ με, καὶ ἐποίησε τῇ Ἀστάρτῃ βδελύγματι Σιδωνίων, καὶ τῷ Χαμὼς, καὶ τοῖς εἰδώλοις Μωὰβ, καὶ τῷ βασιλεῖ αὐτῶν προσοχθίσματι υἱῶν Ἀμμὼν, καὶ οὐκ ἐπορεύθη ἐν ταῖς ὁδοῖς μου τοῦ ποιῆσαι τὸ εὐθὲς ἐνώπιον ἐμοῦ, ὡς Δαυὶδ ὁ πατὴρ αὐτοῦ.
\vs{34}Καὶ οὐ μὴ λάβω τὴν βασιλείαν ὅλην ἐκ χειρὸς αὐτοῦ, διότι ἀντιτασσόμενος ἀντιτάξομαι αὐτῳ πάσας τὰς ἡμέρας τῆς ζωῆς αὐτοῦ, διὰ τὸν Δαυὶδ τὸν δοῦλόν μου ὃν ἐξελεξάμην αὐτόν.
\vs{35}Καὶ λήψομαι τὴν βασιλείαν ἐκ χειρὸς τοῦ υἱοῦ αὐτοῦ, καὶ δώσω σοι τὰ δέκα σκῆπτρα.
\vs{36}Τῷ δὲ υἱῷ αὐτοῦ δώσω τὰ δύο σκῆπτρα, ὅπως ᾖ θέσις τῷ δούλῳ μου Δαυὶδ πάσας τὰς ἡμέρας ἐνώπιον ἐμοῦ ἐν Ἱερουσαλὴμ τῇ πόλει, ἣν ἐξελεξάμην ἐμαυτῷ τοῦ θέσθαι τὸ ὄνομά μου ἐκεῖ.
\vs{37}Καὶ σὲ λήψομαι, καὶ βασιλεύσεις ἐν οἷς ἐπιθυμεῖ ἡ ψυχή σου, καὶ σὺ ἔσῃ βασιλεὺς ἐπὶ τὸν Ἰσραήλ.
\vs{38}Καὶ ἔσται ἐὰν φυλάξῃς πάντα ὅσα ἂν ἐντείλωμαί σοι, καὶ πορευθῇς ἐν ταῖς ὁδοῖς μου, καὶ ποιήσῃς τὸ εὐθὲς ἐνώπιον ἐμοῦ, τοῦ φυλάξασθαι τὰ προστάγματά μου καὶ τὰς ἐντολάς μου, καθὼς ἐποίησε Δαυὶδ ὁ δοῦλός μου, καὶ ἔσομαι μετὰ σοῦ καὶ οἰκοδομήσω σοι οἶκον πιστὸν, καθὼς ᾠκοδόμησα τῷ Δαυίδ.

\vs{40}Καὶ ἐζήτησε Σαλωμὼν θανατῶσαι τὸν Ἱεροβοάμ· καὶ ἀνέστη καὶ ἀπέδρα εἰς Αἴγυπτον πρὸς Σουσακὶμ βασιλέα Αἰγύπτου, καὶ ἦν ἐν Αἰγύπτῳ ἕως οὗ ἀπέθανε Σαλωμών.

\vs{41}Καὶ τὰ λοιπὰ τῶν λόγων Σαλωμὼν, καὶ πάντα ὅσα ἐποίησε, καὶ πᾶσαν τὴν φρόνησιν αὐτοῦ, οὐκ ἰδοὺ ταῦτα γέγραπται ἐν βιβλίῳ ῥημάτων Σαλωμών;
\vs{42}Καὶ αἱ ἡμέραι ἃς ἐβασίλευε Σαλωμὼν ἐν Ἱερουσαλὴμ ἐπὶ πάντα Ἰσραὴλ τεσσαράκοντα ἔτη.
\vs{43}Καὶ ἐκοιμήθη Σαλωμὼν μετὰ τῶν πατέρων αὐτοῦ, καὶ ἔθαψαν αὐτὸν ἐν πόλει Δαυὶδ τοῦ πατρὸς αὐτοῦ· καὶ ἐγενήθη ὡς ἤκουσεν Ἱεροβοὰμ υἱὸς Ναβὰτ, καὶ αὐτοῦ ἔτι ὄντος ἐν Αἰγύπτῳ ὡς ἔφυγεν ἐκ προσώπου Σαλωμὼν καὶ ἐκάθητο ἐν Αἰγύπτῳ, κατευθύνει καὶ ἔρχεται εἰς τὴν πόλιν αὐτοῦ εἰς τὴν γῆν Σαριρὰ τὴν ἐν ὄρει Ἐφραίμ. Καὶ ὁ βασιλεὺς Σαλωμὼν ἐκοιμήθη μετὰ τῶν πατέρων αὐτοῦ, καὶ ἐβασίλευσε Ῥοβοὰμ ὁ υἱὸς αὐτοῦ ἀντʼ αὐτοῦ.

\ch{12}
Καὶ πορεύεται βασιλεὺς Ῥοβοὰμ εἰς Σίκιμα, ὅτι εἰς Σίκιμα ἤρχοντο πᾶς Ἰσραὴλ βασιλεῦσαι αὐτόν.
\vs{3}Καὶ ἐλάλησεν ὁ λαὸς πρὸς τὸν βασιλέα Ῥοβοὰμ, λέγοντες,
\vs{4}Ὁ πατήρ σου ἐβάρυνε τὸν κλοιὸν ἡμῶν, καὶ σὺ νῦν κούφισον ἀπὸ τῆς δουλείας τοῦ πατρός σου τῆς σκληρᾶς, καὶ ἀπὸ τοῦ κλοιοῦ αὐτοῦ τοῦ βαρέως, οὗ ἔδωκεν ἐφʼ ἡμᾶς, καὶ δουλεύσομέν σοι.
\vs{5}Καὶ εἶπε πρὸς αὐτοὺς, ἀπέλθετε ἕως ἡμερῶν τριῶν, καὶ ἀναστρέψατε πρὸς μέ· καὶ ἀπῆλθον.

\vs{6}Καὶ ἀπήγγειλεν ὁ βασιλεὺς τοῖς πρεσβυτέροις, οἳ ἦσαν παρεστῶτες ἐνώπιον Σαλωμὼν τοῦ πατρὸς αὐτοῦ ἔτι ζῶντος αὐτοῦ, λέγων, πῶς ὑμεῖς βουλεύεσθε καὶ ἀποκριθῶ τῷ λαῷ τούτῳ λόγον;
\vs{7}Καὶ ἐλάλησαν πρὸς αὐτὸν, λέγοντες, εἰ ἐν τῇ ἡμέρᾳ ταύτῃ ἔσῃ δοῦλος τῷ λαῷ τούτῳ, καὶ δουλεύσεις αὐτοῖς, καὶ λαλήσεις πρὸς αὐτοὺς λόγους ἀγαθοὺς, καὶ ἔσονταί σοι δοῦλοι πάσας τὰς ἡμέρας.

\vs{8}Καὶ ἐγκατέλιπε τὴν βουλὴν τῶν πρεσβυτέρων ἃ συνεβουλεύσαντο αὐτῷ, καὶ συνεβουλεύσατο μετὰ τῶν παιδαρίων τῶν ἐκτραφέντων μετʼ αὐτοῦ τῶν παρεστηκότων πρὸ προσώπου αὐτοῦ.
\vs{9}Καὶ εἶπεν αὐτοῖς, τί ὑμεῖς συμβουλεύετε; καὶ τί ἀποκριθῶ τῷ λαῷ τούτῳ τοῖς λέγουσι πρὸς μὲ, λεγόντων, κούφισον ἀπὸ τοῦ κλοιοῦ οὗ ἔδωκεν ὁ πατήρ σου ἐφʼ ἡμᾶς;

\vs{10}Καὶ ἐλάλησαν πρὸς αὐτὸν τὰ παιδάρια τὰ ἐκτραφέντα μετʼ αὐτοῦ οἱ παρεστηκότες πρὸ προσώπου αὐτοῦ, λέγοντες, τάδε λαλήσεις τῷ λαῷ τούτῳ τοῖς λαλήσασι πρὸς σὲ, λέγοντες, ὁ πατήρ σου ἐβάρυνε τὸν κλοιὸν ἡμῶν, καὶ σὺ νῦν κούφισον ἀφʼ ἡμῶν· τάδε λαλήσεις πρὸς αὐτοὺς, ἡ μικρότης μου παχυτέρα τῆς ὀσφύος τοῦ πατρός μου.
\vs{11}Καὶ νῦν ὁ πατήρ μου ἐπεσάσσετο ὑμᾶς κλοιῷ βαρεῖ, κᾀλὼ προσθήσω ἐπὶ τὸν κλοιὸν ὑμῶν· ὁ πατήρ μου ἐπαίδευσεν ὑμᾶς ἐν μάστιξιν, ἐγὼ δὲ παιδεύσω ὑμᾶς ἐν σκορπίοις.

\vs{12}Καὶ παρεγένοντο πᾶς Ἰσραὴλ πρὸς τὸν βασιλέα Ῥοβοὰμ ἐν τῇ ἡμέρᾳ τῇ τρίτῃ, καθότι ἐλάλησεν αὐτοῖς ὁ βασιλεὺς, λέγων, ἀναστράφητε πρὸς μὲ τῇ ἡμέρᾳ τῇ τρίτῃ.
\vs{13}Καὶ ἀπεκρίθη ὁ βασιλεὺς πρὸς τὸν λαὸν σκληρά· καὶ ἐγκατέλιπε Ῥοβοὰμ τὴν βουλὴν τῶν πρεσβυτέρων ἃ συνεβουλεύσαντο αὐτῷ,
\vs{14}καὶ ἐλάλησε πρὸς αὐτοὺς κατὰ τὴν βουλὴν τῶν παιδαρίων, λέγων, ὁ πατήρ μου ἐβάρυνε τὸν κλοιὸν ὑμῶν, κᾀγὼ προσθήσω ἐπὶ τὸν κλοιὸν ὑμῶν· ὁ πατήρ μου ἐπαίδευσεν ὑμᾶς ἐν μάστιξι, κᾀγὼ παιδεύσω ὑμᾶς ἐν σκορπίοις.

\vs{15}Καὶ οὐκ ἤκουσεν ὁ βασιλεὺς τοῦ λαοῦ, ὅτι ἦν μεταστροφὴ παρὰ Κυρίου, ὅπως στήσῃ τὸ ῥῆμα αὐτοῦ ὃ ἐλάλησεν ἐν χειρὶ Ἀχιὰ τοῦ Σηλωνίτου περὶ Ἱεροβοὰμ υἱοῦ Ναβάτ.
\vs{16}Καὶ εἶδον πᾶς Ἰσραὴλ, ὅτι οὐκ ἤκουσεν ὁ βασιλεὺς αὐτῶν· καὶ ἀπεκρίθη ὁ λαὸς τῷ βασιλεῖ, λέγων, τίς ἡμῖν μερὶς ἐν Δαυίδ; καὶ οὐκ ἔστιν ἡμῖν κληρονομία ἐν υἱῷ Ἰεσσαί· ἀπότρεχε Ἰσραὴλ, εἰς τὰ σκηνώματά σου· νύν βόσκε τὸν οἶκόν σου, Δαυίδ· καὶ ἀπῆλθεν Ἰσραὴλ εἰς τὰ σκηνώματα αὐτοῦ.

\vs{18}Καὶ ἀπέστειλεν ὁ βασιλεὺς τὸν Ἀδωνιρὰμ τὸν ἐπὶ τοῦ φόρου, καὶ ἐλιθοβόλησαν αὐτὸν ἐν λίθοις καὶ ἀπέθανε· καὶ ὁ βασιλεὺς Ῥοβοὰμ ἔφθασεν ἀναβῆναι τοῦ φυγεῖν εἰς Ἱερουσαλήμ.

\vs{19}Καὶ ἠθέτησεν Ἰσραὴλ εἰς τὸν οἶκον Δαυὶδ ἕως τῆς ἡμέρας ταύτης.
\vs{20}Καὶ ἐγένετο ὡς ἤκουσε πᾶς Ἰσραὴλ ὅτι ἀνέκαμψεν Ἱεροβοὰμ ἐξ Αἰγύπτου, καὶ ἀπέστειλαν καὶ ἐκάλεσαν αὐτὸν εἰς τὴν συναγωγὴν, καὶ ἐβασίλευσαν αὐτὸν ἐπὶ Ἰσραήλ· καὶ οὐκ ἦν ὀπίσω οἴκου Δαυὶδ πάρεξ σκήπτρου Ἰούδα καὶ Βενιαμὶν μόνοι.

\vs{21}Καὶ Ῥοβοὰμ εἰσῆλθεν εἰς Ἱερουσαλὴμ, καὶ ἐξεκκλησίασε τὴν συναγωγὴν Ἰούδα καὶ σκῆπτρον Βενιαμὶν ἑκατὸν καὶ εἴκοσι χιλιάδας νεανιῶν ποιούντων πόλεμον, τοῦ πολεμεῖν πρὸς οἶκον Ἰσραὴλ, ἐπιστρέψαι τὴν βασιλείαν Ῥοβοὰμ υἱῷ Σαλωμών.
\vs{22}Καὶ ἐγένετο λόγος Κυρίου πρὸς Σαμαίαν ἄνθρωπον τοῦ Θεοῦ, λέγων,
\vs{23}εἶπον τῷ Ῥοβοὰμ υἱῷ Σαλωμὼν βασιλεῖ Ἰούδα, καὶ πρὸς πάντα οἶκον Ἰούδα καὶ Βενιαμὶν, καὶ τῷ καταλοίπῳ τοῦ λαοῦ, λέγων,
\vs{24}τάδε λέγει Κύριος, οὐκ ἀναβήσεσθε οὐδὲ πολεμήσετε μετὰ τῶν ἀδελφῶν ὑμῶν υἱῶν Ἰσραήλ· ἀποστρεφέτω ἕκαστος εἰς τὸν οἶκον ἑαυτοῦ, ὅτι παρʼ ἐμοῦ γέγονε τὸ ῥῆμα τοῦτο· καὶ ἤκουσαν τοῦ λόγου Κυρίου, καὶ κατέπαυσαν τοῦ πορευθῆναι κατὰ τὸ ῥῆμα Κυρίου.

\vs{24a}Καὶ ὁ βασιλεὺς Σαλωμὼν κοιμᾶται μετὰ τῶν πατέρων αὐτοῦ, καὶ θάπτεται μετὰ τῶν πατέρων αὐτοῦ, ἐν πόλει Δαυίδ· καὶ ἐβασίλευσε Ῥοβοὰμ υἱὸς αὐτοῦ ἀντʼ αὐτοῦ ἐν Ἱερουσαλὴμ, υἱὸς ὢν ἑκκαίδεκα ἐτῶν ἐν τῷ βασιλεύειν αὐτὸν, καὶ δώδεκα ἔτη ἐβασίλευσεν ἐν Ἱερουσαλήμ· καὶ ὄνομα τῆς μητρὸς αὐτοῦ Ναανὰν, θυγάτηρ Ἄνα υἱοῦ Ναὰς βασιλέως υἱῶν Ἀμμών· καὶ ἐποίησε τὸ πονηρὸν ἐνώπιον Κυρίου, καὶ οὐκ ἐπορεύθη ἐν ὁδῷ Δαυὶδ τοῦ πατρὸς αὐτοῦ.

\vs{24b}Καὶ ἦν ἄνθρωπος ἐξ ὄρους Ἐφραὶμ δοῦλος τῷ Σαλωμὼν, καὶ ὄνομα αὐτῷ Ἱεροβοὰμ, καὶ ὄνομα τῆς μητρὸς αὐτοῦ Σαριρὰ, γυνὴ πόρνη· καὶ ἔδωκεν αὐτὸν Σαλωμὼν εἰς ἄρχοντα σκυτάλης ἐπὶ ἄρσεις οἴκου Ἰωσήφ· καὶ ᾠκοδόμησε τῷ Σαλωμὼν τὴν Σαριρὰ τὴν ἐν ὄρει Ἐφραίμ· καὶ ἦσαν αὐτῷ τριακόσια ἅρματα ἵππων· οὗτος ᾠκοδόμησε τὴν ἄκραν ἐν ταῖς ἄρσεσιν οἴκου Ἐφραὶμ, οὗτος συνέκλεισε τὴν πόλιν Δαυὶδ, καὶ ἦν ἐπαιρόμενος ἐπὶ τὴν βασιλείαν·
\vs{24c}καὶ ἐζήτει Σαλωμὼν θανατῶσαι αὐτόν· καὶ ἐφοβήθη, καὶ ἀπέδρα αὐτὸς πρὸς Σουσακὶμ βασιλέα Αἰγύπτου, καὶ ἦν μετʼ αὐτοῦ ἕως ἀπέθανε Σαλωμών·

\vs{24d}Καὶ ἤκουσεν Ἱεροβοὰμ ἐν Αἰγύπτῳ ὅτι τέθνηκε Σαλωμὼν, καὶ ἐλάλησεν εἰς τὰ ὦτα Σουσακὶμ βασιλέως Αἰγύπτου, λέγων, ἐξαπόστειλόν με, καὶ ἀπελεύσομαι ἐγὼ εἰς τὴν γῆν μου· καὶ εἶπεν αὐτῷ Σουσακὶμ, αἴτησαί τι αἴτημα, καὶ δώσω σοι·
\vs{24e}καὶ Σουσακὶμ ἔδωκε τῷ Ἱεροβοὰμ τὴν Ἀνὼ ἀδελφὴν Θεκεμίνας τὴν πρεσβυτέραν τῆς γυναικὸς αὐτοῦ αὐτῷ εἰς γυναῖκα· αὕτη ἦν μεγάλη ἐν μέσῳ τῶν θυγατέρων τοῦ βασιλέως, καὶ ἔτεκε τῷ Ἱεροβοὰμ τὸν Ἀβιὰ υἱὸν αὐτοῦ·
\vs{24f}καὶ εἶπεν Ἱεροβοὰμ πρὸς Σουσακὶμ, ὄντως ἐξαπόστειλόν με, καὶ ἀπελεύσομαι·

Καὶ ἐξῆλθεν Ἱεροβοὰμ ἐξ Αἰγύπτου, καὶ ἦλθεν εἰς γῆν Σαριρὰ τὴν ἐν ὄρει Ἐφραίμ· καὶ συνάγεται ἐκεῖ πᾶν σκῆπτρον Ἐφραίμ· καὶ ᾠκοδόμησεν ἐκεῖ Ἱεροβοὰμ χάρακα·

\vs{24g}Καὶ ἠῤῥώστησε τὸ παιδάριον αὐτοῦ ἀῤῥωστίᾳ κραταιᾷ σφόδρα· καὶ ἐπορεύθη Ἱεροβοὰμ ἐρωτῆσαι περὶ τοῦ παιδαρίου· καὶ εἶπε πρὸς Ἀνὼ τὴν γυναῖκα αὐτοῦ, ἀνάστηθι, πορεύου, ἐπερώτησον τὸν Θεὸν περὶ τοῦ παιδαρίου, εἰ ζήσεται ἐκ τῆς ἀῤῥωστίας αὐτοῦ·
\vs{24h}καὶ ἄνθρωπος ἦν ἐν Σηλὼμ, καὶ ὄνομα αὐτῷ Ἀχιὰ, καὶ οὗτος ἦν υἱὸς ἑξήκοντα ἐτῶν, καὶ ῥῆμα Κυρίου μετʼ αὐτοῦ· καὶ εἶπεν Ἱεροβοὰμ πρὸς τὴν γυναῖκα αὐτοῦ, ἀνάστηθι, καὶ λάβε εἰς τὴν χεῖρά σου τῷ ἀνθρώπῳ τοῦ Θεοῦ ἄρτους, καὶ κολλύρια τοῖς τέκνοις αὐτοῦ, καὶ σταφυλὴν, καὶ στάμνον μέλιτος·
\vs{24i}καὶ ἀνέστη ἡ γυνὴ, καὶ ἔλαβεν εἰς τὴν χεῖρα αὐτῆς ἄρτους, καὶ δύο κολλύρια, καὶ σταφυλὴν, καὶ στάμνον μέλιτος τῷ Ἀχιά· καὶ ὁ ἄνθρωπος πρεσβύτερος, καὶ οἱ ὀφθαλμοὶ αὐτοῦ ἠμβλυώπουν τοῦ ἰδεῖν·
\vs{24k}καὶ ἀνέστη ἐκ Σαριρὰ καὶ πορεύεται· καὶ ἐγένετο ἐλθούσης αὐτὴς εἰς τὴν πόλιν πρὸς Ἀχιὰ τὸν Σηλωνίτην, καὶ εἶπεν Ἀχιὰ τῷ παιδαρίῳ αὐτοῦ, ἔξελθε δὴ εἰς ἀπαντὴν Ἀνὼ τῇ γυναικὶ Ἱεροβοὰμ, καὶ ἐρεῖς αὐτῇ, εἴσελθε, καὶ μὴ στῇς, ὅτι τάδε λέγει Κύριος, σκληρὰ ἐγὼ ἐπαποστέλλω ἐπὶ σέ·
\vs{24l}καὶ εἰσῆλθεν Ἀνὼ πρὸς τὸν ἄνθρωπον τοῦ Θεοῦ, καὶ εἴπεν αὐτῇ Ἀχιὰ, ἱνατί ἐνήνοχάς μοι ἄρτους, καὶ σταφυλὴν, καὶ κολλύρια, καὶ στάμνον μέλιτος; τάδε λέγει Κύριος, ἰδοὺ σὺ ἀπελεύσῃ ἀπʼ ἐμοῦ, καὶ ἔσται εἰσελθούσης σου τὴν πόλιν εἰς Σαριρὰ, καὶ τὰ κοράσιά σου ἐξλεύσονταί σοι εἰς συνάντησιν, καὶ ἐροῦσί σοι, τὸ παιδάριον τέθνηκεν·
\vs{24m}ὅτι τάδε λέγει Κύριος, ἰδοὺ ἐγὼ ἐξολοθρεύσω τοῦ Ἱεροβοὰμ οὐροῦντα πρὸς τοῖχον, καὶ ἔσονται οἱ τεθνηκότες τοῦ Ἱεροβοὰμ ἐν τῇ πόλει, καταφάγονται οἱ κύνες, καὶ τὸν τεθνηκότα ἐν τῷ ἀγρῷ καταφάγεται τὰ πετεινὰ τοῦ οὐρανοῦ, καὶ τὸ παιδάριον κόψεται, οὐαὶ κύριε, ὅτι εὑρέθη ἐν αὐτῷ ῥῆμα καλὸν περὶ τοῦ Κυρίου·

\vs{24n}Καὶ ἀπῆλθεν ἡ γυνὴ, ὡς ἤκουσε· καὶ ἐγένετο ὡς εἰσῆλθεν εἰς τὴν Σαριρὰ, καὶ τὸ παιδάριον ἀπέθανε· καὶ ἐξῆλθεν ἡ κραυγὴ εἰς ἀπαντήν· καὶ ἐπορεύθη Ἱεροβοὰμ εἰς Σίκιμα τὴν ἐν ὄρει Ἐφραὶμ, καὶ συνήθροισεν ἐκεῖ τὰς φυλὰς τοῦ Ἰσραὴλ, καὶ ἀνέβη ἐκεῖ Ῥοβοὰμ υἱὸς Σαλωμών·
\vs{24o}καὶ λόγος Κυρίου ἐγένετο πρὸς Σαμαίαν τὸν Ἐνλαμὶ, λέγων, λάβε σεαυτῷ ἱμάτιον καινὸν τὸ οὐκ εἰσεληλυθὸς εἰς ὕδωρ, καὶ ῥῆξον αὐτὸ δώδεκα ῥήγματα, καὶ δώσεις τῷ Ἱεροβοὰμ, καὶ ἐρεῖς αὐτῷ, τάδε λέγει Κύριος, λάβε σεαυτῷ δέκα ῥήγματα τοῦ περιβαλέσθαι σε· καὶ ἔλαβεν Ἱεροβοάμ· καὶ εἶπε Σαμαίας, τάδε λέγει Κύριος ἐπὶ τὰς δέκα φυλὰς τοῦ Ἰσραήλ.

\vs{24p}Καὶ εἶπεν ὁ λαὸς πρὸς Ῥοβοὰμ υἱὸν Σαλωμὼν, ὁ πατήρ σου ἐβάρυνε τὸν κλοιὸν αὐτοῦ ἐφʼ ἡμᾶς, καὶ ἐβάρυνε τὰ βρώματα τῆς τραπέζης αὐτοῦ· καὶ νῦν κουφιεῖς ἐφʼ ἡμᾶς, καὶ δουλεύσομέν σοι·
\vs{24q}καὶ εἶπε Ῥοβοὰμ πρὸς τὸν λαὸν, ἔτι τριῶν ἡμερῶν, καὶ ἀποκριθήσομαι ὑμῖν ῥῆμα· καὶ εἶπε Ῥοβοὰμ, εἰσαγάγετέ μοι τοὺς πρεσβυτέρους, καὶ συμβουλεύσομαι μετʼ αὐτῶν τί ἀποκριθῶ τῷ λαῷ ῥῆμα ἐν τῇ ἡμέρᾳ τῇ τρίτῃ· Καὶ ἐλάλησε Ῥοβοὰμ εἰς τὰ ὦτα αὐτῶν, καθὼς ἀπέστειλεν ὁ λαὸς πρὸς αὐτόν· καὶ εἶπον οἱ πρεσβύτεροι τοῦ λαοῦ, οὕτως ἐλάλησε πρὸς σὲ ὁ λαός·

\vs{24r}Καὶ διεσκέδασε Ῥοβοὰμ τὴν βουλὴν αὐτῶν, καὶ οὐκ ἤρεσεν ἐνώπιον αὐτοῦ· καὶ ἀπέστειλε, καὶ εἰσήγαγε τοὺς συντρόφους αὐτοῦ, καὶ ἐλάλησεν αὐτοῖς, ταῦτα καὶ ταῦτα ἀπέσταλκεν ὁ λαὸς πρὸς μὲ, λέγων· καὶ εἶπαν οἱ σύντροφοι αὐτοῦ, οὕτως λαλήσεις πρὸς τὸν λαὸν, λέγων, ἡ μικρότης μου παχυτέρα ὑπὲρ τὴν ὀσφῦν τοῦ πατρός μου· ὁ πατήρ μου ἐμαστίγου ὑμᾶς μάστιξιν, ἐγὼ δὲ κατάρξω ὑμᾶς ἐν σκορπίοις.

\vs{24s}Καὶ ἤρεσε τὸ ῥῆμα ἐνώπιον Ῥοβοάμ· καὶ ἀπεκρίθη τῷ λαῷ, καθὼς συνεβούλευσαν αὐτῷ οἱ σύντροφοι αὐτοῦ τὰ παιδάρια·
\vs{24t}καὶ εἶπε πᾶς ὁ λαὸς ὡς ἀνὴρ εἷς ἕκαστος τῷ πλησίον αὐτοῦ, καὶ ἀνέκραξαν ἅπαντες, λέγοντες, οὐ μερὶς ἡμῖν ἐν Δαυὶδ, οὐδὲ κληρονομία ἐν υἱῷ Ἰεσσαί· ἕκαστος εἰς τὰ σκηνώματά σου Ἰσραὴλ, ὅτι ὁ ἄνθρωπος οὗτος οὐκ εἰς ἄρχοντα οὐδὲ εἰς ἡγούμενον·
\vs{24u}καὶ διεσπάρη πᾶς ὁ λαὸς ἐκ Σικίμων, καὶ ἀπῆλθον ἕκαστος εἰς τὸ σκήνωμα αὐτοῦ· Καὶ κατεκράτησε Ῥοβοὰμ, καὶ ἀπῆλθε, καὶ ἀνέβη ἐπὶ τὸ ἅρμα αὐτοῦ, καὶ εἰσῆλθεν εἰς Ἱερουσαλήμ· καὶ πορεύονται ὀπίσω αὐτοῦ πᾶν σκῆπτρον Ἰούδα, καὶ πᾶν σκῆπτρον Βενιαμίν.
\vs{24x}Καὶ ἐγένετο ἐνισταμένου τοῦ ἐνιαυτοῦ, καὶ συνήθροισε Ῥοβοὰμ πάντα ἄνδρα Ἰούδα καὶ Βενιαμὶν, καὶ ἀνέβη τοῦ πολεμεῖν πρὸς Ἱεροβοὰμ εἰς Σίκιμα·
\vs{24y}καὶ ἐγένετο ῥῆμα Κυρίου πρὸς Σαμαίαν ἄνθρωπον τοῦ Θεοῦ, λέγων, εἶπον τῷ Ῥοβοὰμ βασιλεῖ Ἰούδα, καὶ πρὸς πάντα οἶκον Ἰούδα καὶ Βενιαμὶν, καὶ πρὸς τὸ κατάλειμμα τοῦ λαοῦ, λέγων, τάδε λέγει Κύριος, οὐκ ἀναβήσεσθε οὐδὲ πολεμήσετε πρὸς τοὺς ἀδελφοὺς ὑμῶν υἱοὺς Ἰσραὴλ, ἀναστρέφετε ἕκαστος εἰς τὸν οἶκον αὐτοῦ, ὅτι παρʼ ἐμοῦ γέγονε τὸ ῥῆμα τοῦτο·
\vs{24z}καὶ ἤκουσαν τοῦ λόγου Κυρίου, καὶ ἀνέσχον μὴ πορευθῆναι κατὰ τὸ ῥῆμα Κυρίου.

\vs{25}Καὶ ᾠκοδόμησεν Ἱεροβοὰμ τὴν Σίκιμα τὴν ἐν ὄρει Ἐφραὶμ, καὶ κατῴκει ἐν αὐτῇ, καὶ ἐξῆλθεν ἐκεῖθεν καὶ ᾠκοδόμησε τὴν Φανουήλ.
\vs{26}Καὶ εἶπεν Ἱεροβοὰμ ἐν τῇ καρδίᾳ αὐτοῦ, ἰδοὺ νῦν ἐπιστρέψει ἡ βασιλεία εἰς οἴκον Δαυίδ·
\vs{27}Ἐὰν ἀναβῇ ὁ λαὸς οὗτος ἀναφέρειν θυσίαν ἐν οἴκῳ Κυρίου εἰς Ἱερουσαλὴμ, καὶ ἐπιστραφήσεται καρδία τοῦ λαοῦ πρὸς Κύριον καὶ κύριον αὐτῶν, πρὸς Ῥοβοὰμ βασιλέα Ἰούδα, καὶ ἀποκτενοῦσί με.
\vs{28}Καὶ ἐβουλεύσατο ὁ βασιλεὺς, καὶ ἐπορεύθη, καὶ ἐποίησε δύο δαμάλεις χρυσᾶς, καὶ εἶπε πρὸς τὸν λαὸν, ἱκανούσθω ὑμῖν ἀναβαίνειν εἰς Ἱερουσαλήμ· ἰδοὺ θεοί σου Ἰσραὴλ οἱ ἀναγαγόντες σε ἐκ γῆς Αἰγύπτου.
\vs{29}Καὶ ἔθετο τὴν μίαν ἐν Βαιθὴλ, καὶ τὴν μίαν ἔδωκεν ἐν Δάν.
\vs{30}Καὶ ἐγένετο ὁ λόγος οὗτος εἰς ἁμαρτίαν· καὶ ἐπορεύετο ὁ λαὸς πρὸ προσώπου τῆς μιᾶς ἕως Δὰν, καὶ εἴασαν τὸν οἶκον Κυρίου
\vs{31}Καὶ ἐποίησεν οἴκους ἐφʼ ὑψηλῶν, καὶ ἐποίησεν ἱερεῖς μέρος τι ἐκ τοῦ λαοῦ, οἳ οὐκ ἦσαν ἐκ τῶν υἱῶν Λευί.

\vs{32}Καὶ ἐποίησεν Ἱεροβοὰμ ἑορτὴν ἐν τῷ μηνὶ τῷ ὀγδόῳ ἐν τῇ πεντεκαιδεκάτῃ ἡμέρᾳ τοῦ μηνὸς κατὰ τὴν ἑορτὴν τὴν ἐν γῇ Ἰούδα, καὶ ἀνέβη ἐπὶ τὸ θυσιαστήριον ὃ ἐποίησεν ἐν Βαιθὴλ τοῦ θύειν ταῖς δαμάλεσιν αἷς ἐποίησε, καὶ παρέστησεν ἐν Βαιθὴλ τοὺς ἱερεῖς τῶν ὑψηλῶν ὧν ἐποίησε.
\vs{33}Καὶ ἀνέβη ἐπὶ τὸ θυσιαστήριον ὃ ἐποίησε, τῇ πεντεκαιδεκάτῃ ἡμέρᾳ ἐν τῷ μηνὶ τῷ ὀγδόῳ ἐν τῇ ἑορτῇ ᾗ ἐπλάσατο ἀπὸ καρδίας αὐτοῦ· καὶ ἐποίησεν ἑορτὴν τοῖς υἱοῖς Ἰσραὴλ, καὶ ἀνέβη ἐπὶ τὸ θυσιαστήριον τοῦ ἐπιθῦσαι.

\ch{13}
Καὶ ἰδοὺ ἄνθρωπος τοῦ Θεοῦ ἐξ Ἰούδα παρεγένετο ἐν λόγῳ Κυρίου εἰς Βαιθὴλ, καὶ Ἱεροβοὰμ εἱστήκει ἐπὶ τὸ θυσιαστήριον ἐπιθῦσαι.
\vs{2}Καὶ ἐπεκάλεσε πρὸς τὸ θυσιαστήριον ἐν λόγῳ Κυρίου, καὶ εἶπε, θυσιαστήριον, θυσιαστήριον, τάδε λέγει Κύριος, ἰδοὺ υἱὸς τίκτεται τῷ οἴκῳ Δαυὶδ, Ἰωσίας ὄνομα αὐτῷ, καὶ θύσει ἐπὶ σὲ τοὺς ἱερεῖς τῶν ὑψηλῶν τῶν ἐπιθυόντων ἐπὶ σὲ, καὶ ὀστᾶ ἀνθρώπων καύσει ἐπὶ σέ.
\vs{3}Καὶ δώσει ἐν τῇ ἡμέρᾳ ἐκείνῃ τέρας, λέγων, τοῦτο τὸ ῥῆμα ὃ ἐλάλησε Κύριος, λέγων, ἰδοὺ τὸ θυσιαστήριον ῥήγνυται, καὶ ἐκχυθήσεται ἡ πιότης ἡ ἐπʼ αὐτῷ.

\vs{4}Καὶ ἐγένετο ὡς ἤκουσεν ὁ βασιλεὺς Ἱεροβοὰμ τῶν λόγων τοῦ ἀνθρώπου τοῦ Θεοῦ τοῦ ἐπικαλεσαμένου ἐπὶ τὸ θυσιαστήριον τὸ ἐν Βαιθὴλ, καὶ ἐξέτεινεν ὁ βασιλεὺς τὴν χεῖρα αὐτοῦ ἀπὸ τοῦ θυσιαστηρίου, λέγων, συλλάβετε αὐτόν· καὶ ἰδοὺ ἐξηράνθη ἡ χεὶρ αὐτοῦ, ἣν ἐξέτεινεν ἐπʼ αὐτὸν, καὶ οὐκ ἐδυνήθη ἐπιστρέψαι αὐτὴν πρὸς αὐτόν.
\vs{5}Καὶ τὸ θυσιαστήριον ἐῤῥάγη, καὶ ἐξεχύθη ἡ πιότης ἀπὸ τοῦ θυσιαστηρίου, κατὰ τὸ τέρας ὃ ἔδωκεν ὁ ἄνθρωπος τοῦ Θεοῦ ἐν λόγῳ Κυρίου.
\vs{6}Καὶ εἶπεν ὁ βασιλεὺς Ἱεροβοὰμ τῷ ἀνθρώπῳ τοῦ Θεοῦ, δεήθητι τοῦ πρόσώπου Κυρίου τοῦ Θεοῦ σου, καὶ ἐπιστρεψάτω ἡ χείρ μου πρὸς ἐμέ· καὶ ἐδεήθη ὁ ἄνθρωπος τοῦ Θεοῦ τοῦ προσώπου Κυρίου, καὶ ἐπέστρεψε τὴν χεῖρα τοῦ βασιλέως πρὸς αὐτὸν, καὶ ἐγένετο καθὼς τὸ πρότερον.

\vs{7}Καὶ ἐλάλησεν ὁ βασιλεὺς πρὸς τὸν ἄνθρωπον τοῦ Θεοῦ, εἴσελθε μετʼ ἐμοῦ εἰς οἶκον, καὶ ἀρίστησον, καὶ δώσω σοι δόμα.
\vs{8}Καὶ εἶπεν ὁ ἄνθρωπος τοῦ Θεοῦ πρὸς τὸν βασιλέα, ἐὰν δῷς μοι τὸ ἥμισυ τοῦ οἴκου σου, οὐκ εἰσελεύσομαι μετὰ σοῦ, οὐδὲ μὴ φάγω ἄρτον, οὐδὲ μὴ πίω ὕδωρ ἐν τῷ τόπῳ τούτῳ·
\vs{9}ὅτι οὕτως ἐνετείλατό μοι Κύριος ἐν λόγῳ, λέγων, μὴ φάγῃς ἄρτον καὶ μὴ πίῃς ὕδωρ καὶ μὴ ἐπιστρέψῃς ἐν τῇ ὁδῷ ᾗ ἐπορεύθης ἐν αὐτῇ.
\vs{10}Καὶ ἀπῆλθεν ἐν ὁδῷ ἄλλῃ, καὶ οὐκ ἀνέστρεψεν ἐν τῇ ὁδῷ ᾗ ἦλθεν ἐν αὐτῇ εἰς Βαιθήλ.

\vs{11}Καὶ προφήτης εἷς πρεσβύτης κατῴκη ἐν Βαιθὴλ, καὶ ἔρχονται οἱ υἱοὶ αὐτοῦ καὶ διηγήσαντο αὐτῷ πάντα τὰ ἔργα ἃ ἐποίησεν ὁ ἄνθρωπος τοῦ Θεοῦ ἐν τῇ ἡμέρᾳ ἐκείνῃ ἐν Βαιθὴλ, καὶ τοὺς λόγους οὓς ἐλάλησε τῷ βασιλεῖ, καὶ ἐπέστρεψαν τὸ πρόσωπον τοῦ πατρὸς αὐτῶν.
\vs{12}Καὶ ἐλάλησε πρὸς αὐτοὺς ὁ πατὴρ αὐτῶν, λέγων, ποίᾳ ὁδῷ πεπόρευται; καὶ δεικνύουσιν αὐτῷ οἱ υἱοὶ αὐτοῦ τὴν ὁδὸν ἐν ᾗ ἀνῆλθεν ὁ ἄνθρωπος τοῦ Θεοῦ ὁ ἐλθὼν ἐξ Ἰούδα.
\vs{13}Καὶ εἶπε τοῖς υἱοῖς αὐτοῦ, ἐπισάξατέ μοι τὸν ὄνον· καὶ ἐπέσαξαν αὐτῷ τὸν ὄνον, καὶ ἐπέβη ἐπʼ αὐτὸν,
\vs{14}καὶ ἐπορεύθη κατόπισθεν τοῦ ἀνθρώπου τοῦ Θεοῦ, καὶ εὗρεν αὐτὸν καθήμενον ὑπὸ δρῦν, καὶ εἶπεν αὐτῷ, εἰ σὺ εἶ ὁ ἄνθρωπος τοῦ Θεοῦ ὁ ἐληλυθὼς ἐξ Ἰούδα; καὶ εἶπεν αὐτῷ, ἐγώ.
\vs{15}Καὶ εἶπεν αὐτῷ, δεῦρο μετʼ ἐμοῦ, καὶ φάγε ἄρτον.
\vs{16}Καὶ εἶπεν, οὐ μὴ δύνωμαι τοῦ ἐπιστρέψαι μετὰ σοῦ, οὐδὲ μὴ φάγομαι ἄρτον, οὐδὲ πίομαι ὕδωρ ἐν τῷ τόπῳ τούτῳ·
\vs{17}Ὅτι οὕτως ἐντέταλταί μοι ἐν λόγῳ Κύριος, λέγων, μὴ φάγῃς ἄρτον ἐκεῖ καὶ μὴ πίῃς ὕδωρ καὶ μὴ ἐπιστρέψῃς ἐκεῖ ἐν τῇ ὁδῷ ᾗ ἐπορεύθης ἐν αὐτῇ.

\vs{18}Καὶ εἶπε πρὸς αὐτὸν, κᾀγὼ προφήτης εἰμὶ καθὼς σὺ, καὶ ἄγγελος λελάληκε πρὸς μὲ ἐν ῥήματι Κυρίου, λέγων, ἐπίστρεψον αὐτὸν πρὸς σεαυτὸν εἰς τὸν οἶκόν σου, καὶ φαγέτω ἄρτον, καὶ πιέτω ὕδωρ· καὶ ἐψεύσατο αὐτῷ·
\vs{19}Καὶ ἐπέστρεψεν αὐτὸν, καὶ ἔφαγεν ἄρτον καὶ ἔπιεν ὕδωρ ἐν τῴ οἴκῳ αὐτοῦ.

\vs{20}Καὶ ἐγένετο αὐτῶν καθημένων ἐπὶ τῆς τραπέζης, καὶ ἐγένετο λόγος Κυρίου πρὸς τὸν προφήτην τὸν ἐπιστρέψαντα αὐτόν·
\vs{21}καὶ εἶπε πρὸς τὸν ἄνθρωπον τοῦ Θεοῦ τὸν ἥκοντα ἐξ Ἰούδα, λέγων, τάδε λέγει Κύριος, ἀνθʼ ὧν παρεπίκρανας τὸ ῥῆμα Κυρίου, καὶ οὐκ ἐφύλαξας τὴν ἐντολὴν ἣν ἐνετείλατό σοι Κύριος ὁ Θεός σου,
\vs{22}καὶ ἐπέστρεψας, καὶ ἔφαγες ἄρτον καὶ ἔπιες ὕδωρ ἐν τῷ τόπῳ τούτῳ ᾧ ἐλάλησε πρὸς σὲ, λέγων, οὐ μὴ φάγῃς ἄρτον καὶ μὴ πίῃς ὕδωρ, οὐ μὴ εἰσέλθῃ τὸ σῶμά σου εἰς τὸν τάφον τῶν πατέρων σου.

\vs{23}Καὶ ἐγένετο μετὰ τὸ φαγεῖν αὐτὸν ἄρτον καὶ πιεῖν ὕδωρ, καὶ ἐπέσαξεν αὐτῷ τὸν ὄνον, καὶ ἐπέστρεψε, καὶ ἀπῆλθε.
\vs{24}Καὶ εὗρεν αὐτὸν λέων ἐν τῇ ὁδῷ, καὶ ἐθανάτωσεν αὐτόν· καὶ ἦν τὸ σῶμα αὐτοῦ ἐῥῥιμμένον ἐν τῇ ὁδῷ, καὶ ὁ ὄνος εἱστήκει παρʼ αὐτὸ, καὶ ὁ λέων εἱστήκει παρὰ τὸ σῶμα.
\vs{25}Καὶ ἰδοὺ ἄνδρες παραπορευόμενοι καὶ εἶδον τὸ θνησιμαῖον ἐῤῥιμμένον ἐν τῇ ὁδῷ, καὶ ὁ λέων εἱστήκει ἐχόμενα τοῦ θνησιμαίου· καὶ εἰσῆλθον, καὶ ἐλάλησαν ἐν τῇ πόλει οὗ ὁ προφήτης ὁ πρεσβύτης κατῴκει ἐν αὐτῄ.
\vs{26}Καὶ ἤκουσεν ὁ ἐπιστρέψας αὐτὸν ἐκ τῆς ὁδοῦ, καὶ εἶπεν, ὁ ἄνθρωπος τοῦ Θεοῦ οὗτός ἐστιν ὃς παρεπίκρανε τὸ ῥῆμα Κυρίου·
\vs{28}Καὶ ἐπορεύθη καὶ εὗρε τὸ σῶμα αὐτοῦ ἐῤῥιμμένον ἐν τῇ ὁδῷ, καὶ ὁ ὄνος, καὶ ὁ λέων εἱστήκεισαν παρὰ τὸ σῶμα· καὶ οὐκ ἔφαγεν ὁ λέων τὸ σῶμα τοῦ ἀνθρώπου τοῦ Θεοῦ, καὶ οὐ συνέτριψε τὸν ὄνον.

\vs{29}Καὶ ᾖρεν ὁ προφήτης τὸ σῶμα τοῦ ἀνθρώπου τοῦ Θεοῦ, καὶ ἐπέθηκεν αὐτὸ ἐπὶ τὸν ὄνον, καὶ ἐπέστρεψεν αὐτὸν εἰς τὴν πόλιν
\vs{30}ὁ προφήτης, τοῦ θάψαι αὐτὸν ἐν τῷ τάφῳ ἑαυτοῦ, καὶ ἐκόψαντο αὐτὸν, οὐαὶ ἀδελφέ.
\vs{31}Καὶ ἐγένετο μετὰ τὸ κόψασθαι αὐτὸν, καὶ εἶπε τοῖς υἱοῖς αὐτοῦ, λέγων, ἐὰν ἀποθάνω, θάψατέ με ἐν τῷ τάφῳ τούτῳ οὗ ὁ ἄνθρωπος τοῦ Θεοῦ τέθαπται ἐν αὐτῷ, παρὰ τὰ ὀστᾶ αὐτοῦ θέτε με, ἵνα σωθῶσι τὰ ὀστᾶ μου μετὰ τῶν ὀστῶν αὐτοῦ.
\vs{32}Ὅτι γινόμενον ἔσται τὸ ῥῆμα ὃ ἐλάλησεν ἐν λόγῳ Κυρίου ἐπὶ τὸ θυσιαστήριον ἐν Βαιθὴλ καὶ ἐπὶ τοὺς οἴκους τοὺς ὑψηλοὺς τοὺς ἐν Σαμαρείᾳ.

\vs{33}Καὶ μετὰ τὸ ῥῆμα τοῦτο οὐκ ἐπέστρεψεν Ἱεροβοὰμ ἀπὸ τῆς κακίας αὐτοῦ, καὶ ἐπέστρεψε καὶ ἐποίησεν ἐκ μέρους τοῦ λαοῦ ἱερεῖς ὑψηλῶν· ὁ βουλόμενος ἐπλήρου τὴν χεῖρα αὐτοῦ, καὶ ἐγένετο ἱερεὺς εἰς τὰ ὑψηλά.
\vs{34}Καὶ ἐγένετο τὸ ῥῆμα τοῦτο εἰς ἁμαρτίαν τῷ οἴκῳ Ἱεροβοὰμ, καὶ εἰς ὄλεθρον, καὶ εἰς ἀφανισμὸν ἀπὸ προσώπου τῆς γῆς.

\ch{14}
\vs{21}Καὶ Ῥοβοὰμ υἱὸς Σαλωμὼν ἐβασίλευσεν ἐπὶ Ἰούδαν· υἱὸς τεσσαράκοντα καὶ ἑνὸς ἐνιαυτῶν Ῥοβοὰμ ἐν τῷ βασιλεύειν αὐτόν· καὶ ἑπτακαίδεκα ἔτη ἐβασίλευσεν ἐν Ἱερουσαλὴμ τῇ πόλει, ἣν ἐξελέξατο Κύριος θέσθαι τὸ ὄνομα αὐτοῦ ἐκεῖ ἐκ πασῶν φυλῶν τοῦ Ἰσραήλ· καὶ τὸ ὄνομα τῆς μητρὸς αὐτοῦ Νααμὰ ἡ Ἀμμωνίτις.
\vs{22}Καὶ ἐποίησε Ῥοβοὰμ τὸ πονηρὸν ἐνώπιον Κυρίου· καὶ παρεζήλωσεν αὐτὸν ἐν πᾶσιν οἷς ἐποίησαν οἱ πατέρες αὐτῶν ἐν ταῖς ἁμαρτίαις αὐτῶν αἷς ἥμαρτον.
\vs{23}Καὶ ᾠκοδόμησαν ἑαυτοῖς ὑψηλὰ καὶ στήλας καὶ ἄλση ἐπὶ πάντα βουνὸν ὑψηλὸν, καὶ ὑποκάτω παντὸς ξύλου συσκίου.
\vs{24}Καὶ σύνδεσμος ἐγενήθη ἐν τῇ γῇ, καὶ ἐποίησαν ἀπὸ πάντων τῶν βδελυγμάτων τῶν ἐθνῶν ὧν ἐξῇρε Κύριος ἀπὸ προσώπου υἱῶν Ἰσραήλ.

\vs{25}Καὶ ἐγένετο ἐν τῷ ἐνιαυτῷ τῷ πέμπτῳ βασιλεύοντος Ῥοβοὰμ, ἀνέβη Σουσακεὶμ βασιλεὺς Αἰγύπτου ἐπὶ Ἱερουσαλὴμ,
\vs{26}καὶ ἔλαβε πάντας τοὺς θησαυροὺς οἴκου Κυρίου καὶ τοὺς θησαυροὺς οἴκου τοῦ βασιλέως, καὶ τὰ δόρατα τὰ χρυσᾶ ἃ ἔλαβε Δαυὶδ ἐκ χειρὸς τῶν παίδων Ἀδραζαὰρ βασιλέως Σουβὰ, καὶ εἰσήνεγκεν αὐτὰ εἰς Ἱερουσαλὴμ τὰ πάντα ἃ ἔλαβεν, ὅπλα τὰ χρυσᾶ ὅσα ἐποίησε Σαλωμὼν, καὶ ἀπήνεγκεν αὐτὰ εἰς Αἴγυπτον.
\vs{27}Καὶ ἐποίησε Ῥοβοὰμ ὁ βασιλεὺς ὅπλα χαλκᾶ ἀντʼ αὐτῶν· καὶ ἐπέθεντο ἐπʼ αὐτὸν οἱ ἡγούμενοι τῶν παρατρεχόντων οἱ φυλάσσοντες τὸν πυλῶνα οἴκου βασιλέως.
\vs{28}Καὶ ἐγένετο ὅτε εἰσεπορεύετο ὁ βασιλεὺς εἰς οἶκον Κυρίου, καὶ ᾖρον αὐτὰ οἱ παρατρέχοντες καὶ ἀπηρείδοντο αὐτὰ εἰς τὸ θεὲ τῶν παρατρεχόντων.

\vs{29}Καὶ τὰ λοιπὰ τῶν λόγων Ῥοβοὰμ καὶ πάντα ἃ ἐποίησεν, οὐκ ἰδοὺ ταῦτα γεγραμμένα ἐν βιβλίῳ λόγων τῶν ἡμερῶν τοῖς βασιλεῦσιν Ἰούδα;
\vs{30}Καὶ πόλεμος ἦν ἀναμέσον Ῥοβοὰμ καὶ ἀναμέσον Ἱεροβοὰμ πάσας τὰς ἡμέρας.
\vs{31}Καὶ ἐκοιμήθη Ῥοβοὰμ μετὰ τῶν πατέρων αὐτοῦ, καὶ θάπτεται μετὰ τῶν πατέρων αὐτοῦ, ἐν πόλει Δαυίδ· καὶ ἐβασίλευσεν Ἀβιοὺ ὁ υἱὸς αὐτοῦ ἀντʼ αὐτοῦ.

\ch{15}
Καὶ ἐν τῷ ὀκτωκαιδεκάτῳ ἔτει βασιλεύοντος Ἱεροβοὰμ υἱοῦ Ναβὰτ, βασιλεύει Ἀβιοὺ υἱὸς Ῥοβοὰμ ἐπὶ Ἰούδαν·
\vs{2}Καὶ τρία ἔτη ἐβασίλευσεν ἐπὶ Ἱερουσαλήμ· καὶ ὄνομα τῆς μητρὸς αὐτοῦ Μααχὰ, θυγάτηρ Ἀβεσσαλώμ.
\vs{3}Καὶ ἐπορεύθη ἐν ταῖς ἁμαρτίαις τοῦ πατρὸς αὐτοῦ αἷς ἐποίησεν ἐνώπιον αὐτοῦ, καὶ οὐκ ἦν ἡ καρδία αὐτοῦ τελεία μετὰ Κυρίου Θεοῦ αὐτοῦ, ὡς ἡ καρδία τοῦ πατρὸς αὐτοῦ.
\vs{4}Ὅτι διὰ Δαυὶδ ἔδωκεν αὐτῷ Κύριος κατάλειμμα, ἵνα στήσῃ τὰ τέκνα αὐτοῦ μετʼ αὐτὸν, καὶ στήσῃ τὴν Ἱερουσαλὴμ·
\vs{5}ὡς ἐποίησε Δαυὶδ τὸ εὐθὲς ἐνώπιον Κυρίου, οὐκ ἐξέκλινεν ἀπὸ πάντων ὧν ἐνετείλατο αὐτῷ πάσας τὰς ἡμέρας τῆς ζωῆς αὐτοῦ.

\vs{7}Καὶ τὰ λοιπὰ τῶν λόγων Ἀβιοὺ καὶ πάντα ἃ ἐποίησεν, οὐκ ἰδοὺ ταῦτα γεγραμμένα ἐπὶ βιβλίῳ λόγων τῶν ἡμερῶν τοῖς βασιλεῦσιν Ἰούδα; καὶ πόλεμος ἦν ἀναμέσον Ἀβιοὺ καὶ ἀναμέσον Ἱεροβοάμ.
\vs{8}Καὶ ἐκοιμήθη Ἀβιοὺ μετὰ τῶν πατέρων αὐτοῦ ἐν τῷ εἰκοστῷ καὶ τετάρτῳ ἔτει τοῦ Ἱεροβοὰμ, καὶ θάπτεται μετὰ τῶν πατέρων αὐτοῦ ἐν πόλει Δαυίδ· καὶ βασιλεύει Ἀσὰ υἱὸς αὐτοῦ ἀντʼ αὐτοῦ.

\vs{9}Ἐν τῷ ἐνιαυτῷ τετάρτῳ καὶ εἰκοστῷ τοῦ Ἱεροβοὰμ βασιλέως Ἰσραὴλ, βασιλεύει Ἀσὰ ἐπὶ Ἰούδαν·
\vs{10}Καὶ τεσσαράκοντα καὶ ἓν ἔτος ἐβασίλευσεν ἐν Ἱερουσαλήμ· καὶ ὄνομα τῆς μητρὸς αὐτοῦ Ἀνὰ, θυγάτηρ Ἀβεσσαλώμ.
\vs{11}Καὶ ἐποίησεν Ἀσὰ τὸ εὐθὲς ἐνώπιον Κυρίου, ὡς Δαυὶδ ὁ πατὴρ αὐτοῦ.
\vs{12}Καὶ ἀφεῖλε τὰς τελετὰς ἀπὸ τῆς γῆς, καὶ ἐξαπέστειλε πάντα τὰ ἐπιτηδεύματα ἃ ἐποίησαν οἱ πατέρες αὐτοῦ.
\vs{13}Καὶ τὴν Ἀνὰ τὴν μητέρα ἑαυτοῦ μετέστησε τοῦ μὴ εἶναι ἡγουμένην, καθὼς ἐποίησε σύνοδον ἐν τῷ ἄλσει αὐτῆς· καὶ ἐξέκοψεν Ἀσὰ τὰς καταδύσεις αὐτῆς, καὶ ἐνέπρησε πυρὶ ἐν τῷ χειμάῤῥῳ τῶν Κέδρων.
\vs{14}Τὰ δὲ ὑψηλὰ οὐκ ἐξῇρε· πλὴν ἡ καρδία Ἀσὰ ἦν τελεία μετὰ Κυρίου πάσας τὰς ἡμέρας αὐτοῦ.
\vs{15}Καὶ εἰσήνεγκε τοὺς κίονας τοῦ πατρὸς αὐτοῦ, καὶ τοὺς κίονας αὐτοῦ εἰσήνεγκεν εἰς τὸν οἶκον Κυρίου ἀργυροῦς καὶ χρυσοῦς, καὶ σκεύη.

\vs{16}Καὶ πόλεμος ἦν ἀναμέσον Ἀσὰ καὶ ἀναμέσον Βαασὰ βασιλέως Ἰσραὴλ πάσας τὰς ἡμέρας αὐτῶν.
\vs{17}Καὶ ἀνέβη Βαασὰ βασιλεὺς Ἰσραὴλ ἐπὶ Ἰούδαν, καὶ ᾠκοδόμησε τὴν Ῥαμὰ, τοῦ μὴ εἶναι ἐκπορευόμενον καὶ εἰσπορευόμενον τῷ Ἀσὰ βασιλεῖ Ἰούδα.

\vs{18}Καὶ ἔλαβεν Ἀσὰ σύμπαν τὸ ἀργύριον καὶ τὸ χρυσίον τὸ εὑρεθὲν ἐν τοῖς θησαυροῖς οἴκου Κυρίου καὶ ἐν τοῖς θησαυροῖς τοῦ οἴκου τοῦ βασιλέως, καὶ ἔδωκεν αὐτὰ εἰς χεῖρας παίδων αὐτοῦ· καὶ ἐξαπέστειλεν αὐτοὺς ὁ βασιλεὺς Ἀσὰ πρὸς υἱὸν Ἄδερ υἱὸν Ταβερεμὰ υἱοῦ Ἀζὶν βασιλέως Συρίας τοῦ κατοικοῦντος ἐν Δαμασκῷ, λέγων,
\vs{19}διάθου διαθήκην ἀναμέσον ἐμοῦ καὶ ἀναμέσον σοῦ, καὶ ἀναμέσον τοῦ πατρός μου καὶ τοῦ πατρός σου· ἰδοὺ ἐξαπέσταλκά σοι δῶρα ἀργύριον καὶ χρυσίον· δεῦρο, διασκέδασον τὴν διαθήκην σου τὴν πρὸς Βαασὰ βασιλέα Ἰσραὴλ, καὶ ἀναβήσεται ἀπʼ ἐμοῦ.
\vs{20}Καὶ ἤκουσεν υἱὸς Ἄδερ τοῦ βασιλέως Ἀσά, καὶ ἀπέστειλε τοὺς ἄρχοντας τῶν δυνάμεων αὐτοῦ ταῖς πόλεσι τοῦ Ἰσραὴλ, καὶ ἐπάταξαν τὴν Ἀῒν, τὴν Δὰν, καὶ τὴν Ἀβὲλ οἴκου Μααχὰ, καὶ πᾶσαν τὴν Χεννερὲθ, ἕως πάσης τῆς γῆς Νεφθαλί.
\vs{21}Καὶ ἐγένετο ὡς ἤκουσε Βαασὰ, καὶ διέλιπε τοῦ οἰκοδομεῖν τὴν Ῥαμὰ, καὶ ἀνέστρεψεν εἰς Θερσά.

\vs{22}Καὶ ὁ βασιλεὺς Ἀσὰ παρήγγειλε παντὶ Ἰούδα εἰς ἐνακὶμ, καὶ αἴρουσι τοὺς λίθους τῆς Ῥαμὰ, καὶ τὰ ξύλα αὐτῆς ἃ ᾠκοδόμησε Βαασά· καὶ ᾠκοδόμησεν ἐν αὐτοῖς ὁ βασιλεὺς Ἀσὰ πᾶν βουνὸν Βενιαμὶν καὶ τὴν σκοπιάν.

\vs{23}Καὶ τὰ λοιπὰ τῶν λόγων Ἀσὰ, καὶ πᾶσα ἡ δυναστεία αὐτοῦ ἣν ἐποίησε, καὶ τὰς πόλεις ἃς ᾠκοδόμησεν, οὐκ ἰδοὺ ταῦτα γεγραμμένα ἐστὶν ἐπὶ βιβλίῳ λόγων τῶν ἡμερῶν τοῖς βασιλεῦσιν Ἰούδα; πλὴν ἐν τῷ καιρῷ τοῦ γήρως αὐτοῦ ἐπόνεσε τοὺς πόδας αὐτοῦ.
\vs{24}Καὶ ἐκοιμήθη Ἀσὰ μετὰ τῶν πατέρων αὐτοῦ, καὶ θάπτεται μετὰ τῶν πατέρων αὐτοῦ ἐν πόλει Δαυίδ πατρὸς αὐτοῦ· καὶ βασιλεύει Ἰωσαφὰτ υἱὸς αὐτοῦ ἀντʼ αὐτοῦ.

\vs{25}Καὶ Ναβὰτ υἱὸς Ἱεροβοὰμ βασιλεύει ἐπὶ Ἰσραὴλ ἐν ἔτει δευτέρῳ τοῦ Ἀσὰ βασιλέως Ἰούδα, καὶ ἐβασίλευσεν ἐν Ἰσραὴλ ἔτη δύο.
\vs{26}Καὶ ἐποίησε τὸ πονηρὸν ἐνώπιον Κυρίου, καὶ ἐπορεύθη ἐν ὁδῷ τοῦ πατρὸς αὐτοῦ καὶ ἐν ταῖς ἁμαρτίαις αὐτοῦ αἷς ἐξήμαρτε τὸν Ἰσραήλ.

\vs{27}Καὶ περιεκάθισεν αὐτὸν Βαασὰ υἱὸς Ἀχιὰ ἐπὶ τὸν οἶκον Βελαὰν υἱοῦ Ἀχιὰ, καὶ ἐχάραξεν αὐτὸν ἐν Γαβαθὼν τῇ τῶν ἀλλοφύλων· καὶ Ναβὰτ καὶ πᾶς Ἰσραὴλ περιεκάθητο ἐπὶ Γαβαθών.
\vs{28}Καὶ ἐθανάτωσεν αὐτὸν Βαασὰ ἐν ἔτει τρίτῳ τοῦ Ἀσὰ υἱοῦ Ἀσὰ βασιλέως Ἰούδα, καὶ ἐβασίλευσεν.
\vs{29}Καὶ ἐγένετο ὡς ἐβασίλευσε, καὶ ἐπάταξεν ὅλον τὸν οἶκον Ἱεροβοὰμ, καὶ οὐχ ὑπελείπετο πᾶσαν πνοὴν τοῦ Ἱεροβοὰμ ἕως τοῦ ἐξολοθρεῦσαι αὐτὸν, κατὰ τὸ ῥῆμα Κυρίου ὃ ἐλάλησεν ἐν χειρὶ δούλου αὐτοῦ Ἀχιὰ τοῦ Σηλωνίτου·
\vs{30}περὶ τῶν ἁμαρτιῶν Ἱεροβοὰμ, ὅς ἐξήμαρτε τὸν Ἰσραήλ, καὶ ἐν τῷ παροργισμῷ αὐτοῦ ᾧ παρώργισε τὸν Κύριον Θεὸν τοῦ Ἰσραήλ.
\vs{31}Καὶ τὰ λοιπὰ τῶν λόγων Ναβὰτ καὶ πάντα ἃ ἐποίησεν, οὐκ ἰδοὺ ταῦτα γεγραμμένα ἐστὶν ἐν βιβλίῳ λόγων τῶν ἡμερῶν τοῖς βασιλεῦσιν Ἰσραήλ;

\vs{33}Καὶ ἐν τῷ ἔτει τῷ τρίτῳ τοῦ Ἀσὰ βασιλέως Ἰούδα βασιλεύει Βαασὰ υἱὸς Ἀχιὰ ἐπὶ Ἰσραὴλ ἐν Θερσᾷ εἴκοσι καὶ τέσσαρα ἔτη.
\vs{34}Καὶ ἐποίησε τὸ πονηρὸν ἐνώπιον Κυρίου, καὶ ἐπορεύθη ἐν ὁδῷ Ἱεροβοὰμ υἱοῦ Ναβὰτ καὶ ἐν ταῖς ἁμαρτίαις αὐτοῦ ὡς ἐξήμαρτε τὸν Ἰσραήλ.

\ch{16}
Καὶ ἐγένετο λόγος Κυρίου ἐν χειρὶ Ἰοὺ υἱοῦ Ἀνανὶ πρὸς Βαασά.
\vs{2}Ἀνθʼ ὧν ὕψωσά σε ἀπὸ τῆς γῆς, καὶ ἔδωκά σε ἡγούμενον ἐπὶ τὸν λαόν μου Ἰσραήλ, καὶ ἐπορεύθης ἐν τῇ ὁδῷ Ἱεροβοὰμ, καὶ ἐξήμαρτες τὸν λαόν μου τὸν Ἰσραὴλ, τοῦ παροργίσαι με ἐν τοῖς ματαίοις αὐτῶν,
\vs{3}ἰδοὺ ἐγὼ ἐξεγείρω ὀπίσω Βαασὰ, καὶ ὄπισθεν τοῦ οἴκου αὐτοῦ, καὶ δώσω τὸν οἶκόν σου ὡς τὸν οἶκον Ἱεροβοὰμ υἱοῦ Ναβάτ.
\vs{4}Τὸν τεθνηκότα τοῦ Βαασὰ ἐν τῇ πόλει καταφάγονται αὐτὸν οἱ κύνες, καὶ τὸν τεθνηκότα αὐτοῦ ἐν τῷ πεδίῳ καταφάγονται αὐτὸν τὰ πετεινὰ τοῦ οὐρανοῦ.

\vs{5}Καὶ τὰ λοιπὰ τῶν λόγων Βαασὰ καὶ πάντα ἃ ἐποίησε, καὶ αἱ δυναστεῖαι αὐτοῦ, οὐκ ἰδοὺ ταῦτα γεγραμμένα ἐν βιβλίῳ λόγων τῶν ἡμερῶν τῶν βασιλέων Ἰσραήλ;
\vs{6}Καὶ ἐκοιμήθη Βαασὰ μετὰ τῶν πατέρων αὐτοῦ, καὶ θάπτεται ἐν Θερσᾷ, καὶ βασιλεύει Ἠλὰ υἱὸς αὐτοῦ ἀντʼ αὐτοῦ.

\vs{7}Καὶ ἐν χειρὶ Ἰοὺ υἱοῦ Ἀνανὶ ἐλάλησε Κύριος ἐπὶ Βαασὰ καὶ ἐπὶ τὸν οἶκον αὐτοῦ, πᾶσαν τὴν κακίαν ἣν ἐποίησεν ἐνώπιον Κυρίου τοῦ παροργίσαι αὐτὸν ἐν τοῖς ἔργοις τῶν χειρῶν αὐτου, τοῦ εἶναι κατὰ τὸν οἶκον Ἱεροβοὰμ, καὶ ὑπὲρ τοῦ πατάξαι αὐτόν.

\vs{8}Καὶ Ἠλὰ υἱὸς Βαασὰ ἐβασίλευσεν ἐπὶ Ἰσραὴλ δύο ἔτη ἐν Θερσᾷ.
\vs{9}Καὶ συνέστρεψεν ἐπʼ αὐτὸν Ζαμβρὶ ὁ ἄρχων τῆς ἡμίσους τῆς ἵππου, καὶ αὐτὸς ἦν ἐν Θερσᾷ πίνων μεθύων ἐν τῷ οἴκῳ Ὠσᾶ τοῦ οἰκονόμου ἐν Θερσᾷ.
\vs{10}Καὶ εἰσῆλθε Σαμβρὶ καὶ ἐπάταξεν αὐτὸν καὶ ἐθανάτωσεν αὐτὸν, καὶ ἐβασίλευσεν ἀντʼ αὐτοῦ.
\vs{11}Καὶ ἐγενήθη ἐν τῷ βασιλεῦσαι αὐτὸν ἐν τῷ καθίσαι αὐτὸν ἐπὶ τοῦ θρόνου αὐτοῦ, καὶ ἐπάταξεν ὅλον τὸν οἶκον Βαασὰ,
\vs{12}κατὰ τὸ ῥῆμα ὃ ἐλάλησε Κύριος ἐπὶ τὸν οἶκον Βαασὰ, καὶ πρὸς Ἰοὺ τὸν προφήτην
\vs{13}περὶ πασῶν τῶν ἁμαρτιῶν Βαασὰ καὶ Ἠλὰ τοῦ υἱοῦ αὐτοῦ, ὡς ἐξήμαρτε τὸν Ἰσραὴλ, τοῦ παροργίσαι Κύριον τὸν Θεὸν Ἰσραὴλ ἐν τοῖς ματαίοις αὐτῶν.
\vs{14}Καὶ τὰ λοιπὰ τῶν λόγων Ἠλὰ ἃ ἐποίησεν, οὐκ ἰδοὺ ταῦτα γεγραμμένα ἐν βιβλίῳ λόγων τῶν ἡμερῶν τῶν βασιλέων Ἰσραήλ;

\vs{15}Καὶ Ζαμβρὶ ἐβασίλευσεν ἐν Θερσᾷ ἡμέρας ἑπτά· καὶ ἡ παρεμβολὴ Ἰσραὴλ ἐπὶ Γαβαθὼν τὴν τῶν ἀλλοφύλων.
\vs{16}Καὶ ἤκουσεν ὁ λαὸς ἐν τῇ παρεμβολῇ, λεγόντων, συνεστράφη Ζαμβρὶ καὶ ἔπαισν τὸν βασιλέα· καὶ ἐβασίλευσαν ἐν Ἰσραὴλ τὸν Ἀμβρὶ τὸν ἡγούμενον τῆς στρατιᾶς ἐπὶ Ἰσραὴλ ἐν τῇ ἡμέρᾳ ἐκείνῃ ἐν τῇ παρεμβολῇ.
\vs{17}Καὶ ἀνέβη Ἀμβρὶ καὶ πᾶς Ἰσραὴλ μετʼ αὐτοῦ ἐκ Γαβαθὼν, καὶ περιεκάθισαν ἐπὶ Θερσᾷ.
\vs{18}Καὶ ἐγενήθη ὡς εἶδε Ζαμβρὶ ὅτι προκατείληπται αὐτοῦ ἡ πόλις, καὶ πορεύεται εἰς ἄντρον τοῦ οἴκου τοῦ βασιλέως, καὶ ἐνεπύρισεν ἐπʼ αὐτὸν τὸν οἶκον τοῦ βασιλέως, καὶ ἀπέθανεν
\vs{19}ὑπὲρ τῶν ἁμαρτιῶν αὐτοῦ ὧν ἐποίησε, τοῦ ποιῆσαι τὸ πονηρὸν ἐνώπιον Κυρίου πορευθῆναι ἐν ὁδῷ Ἱεροβοὰμ υἱοῦ Ναβὰτ, καὶ ἐν ταῖς ἁμαρτίαις αὐτοῦ ὡς ἐξήμαρτε τὸν Ἰσραήλ.
\vs{20}Καὶ τὰ λοιπὰ τῶν λόγων Ζαμβρὶ καὶ τὰς συνάψεις αὐτοῦ ἃς συνῆψεν, οὐκ ἰδοὺ ταῦτα γεγραμμένα ἐν βιβλίῳ λόγων τῶν ἡμερῶν τῶν βασιλέων Ἰσραήλ;

\vs{21}Τότε μερίζεται ὁ λαὸς Ἰσραήλ· ἥμισυ τοῦ λαοῦ γίνεται ὀπίσω Θαμνὶ υἱοῦ Γωνὰθ τοῦ βασιλεῦσαι αὐτὸν, καὶ τὸ ἥμισυ τοῦ λαοῦ γίνεται ὀπίσω Ἀμβρί.
\vs{22}Ὁ λαὸς ὁ ὢν ὀπίσω Ἀμβρὶ ὑπερεκράτησε τὸν λαὸν τὸν ὀπίσω Θαμνὶ υἱοῦ Γωνάθ· καὶ ἀπέθανε Θαμνὶ καὶ Ἰωρὰμ ὁ ἀδελφὸς αὐτοῦ ἐν τῷ καιρῷ ἐκείνῳ, καὶ ἐβασίλευσεν Ἀμβρὶ μετὰ Θαμνί.

\vs{23}Ἐν τῷ ἔτει τῷ τριακοστῷ καὶ πρώτῳ τοῦ βασιλέως Ἀσὰ βασιλεύει Ἀμβρὶ ἐπὶ Ἰσραὴλ δώδεκα ἔτη· ἐν Θερσᾷ βασιλεύει ἓξ ἔτη.
\vs{24}Καὶ ἐκτήσατο Ἀμβρὶ τὸ ὄρος τὸ Σεμερὼν παρὰ Σεμὴρ τοῦ κυρίου τοῦ ὄρους ἐν δύο ταλάντων ἀργυρίου· καὶ ᾠκοδόμησε τὸ ὄρος, καὶ ἐπεκάλεσαν τὸ ὄνομα τοῦ ὄρους οὗ ᾠκοδόμησαν ἐπὶε τῷ ὀνόματι Σεμὴρ τοῦ κυρίου τοῦ ὄρους, Σεμηρών.
\vs{25}Καὶ ἐποίησεν Ἀμβρὶ τὸ πονηρὸν ἐνώπιον Κυρίου, καὶ ἐπονηρεύσατο ὑπὲρ πάντας τοὺς γενομένους ἔμπροσθεν αὐτοῦ.
\vs{26}Καὶ ἐπορεύθη ἐν πάσῃ ὁδῷ Ἱεροβοὰμ υἱοῦ Ναβὰτ, καὶ ἐν ταῖς ἁμαρτίαις αὐτοῦ αἷς ἐξήμαρτε τὸν Ἰσραὴλ, τοῦ παροργίσαι τὸν Κύριον Θεὸν Ἰσραὴλ ἐν τοῖς ματαίοις αὐτῶν.
\vs{27}Καὶ τὰ λοιπὰ τῶν λόγων Ἀμβρὶ καὶ πάντα ἃ ἐποίησε, καὶ πᾶσα ἡ δυναστεία αὐτοῦ, οὐκ ἰδοὺ ταῦτα γεγραμμένα ἐν βιβλίῳ λόγων τῶν ἡμερῶν τῶν βασιλέων Ἰσραήλ;

\vs{28}Καὶ ἐκοιμήθη Ἀμβρὶ μετὰ τῶν πατέρων αὐτοῦ, καὶ θάπτεται ἐν Σαμαρείᾳ, καὶ βασιλεύει Ἀχαὰβ ὁ υἱὸς αὐτοῦ ἀντʼ αὐτοῦ.

\vs{28a}Καὶ ἐν τῷ ἐνιαυτῷ τῷ ἑνδεκάτῳ ἔτει τοῦ Ἀμβρὶ βασιλεύει Ἰωσαφὰτ υἱὸς Ἀσὰ ἐτῶν τριάκοντα καὶ πέντε ἐν τῇ βασιλείᾳ αὐτοῦ, καὶ εἰκοσιπέντε ἔτη ἐβασίλευσεν ἐν Ἱερουσαλήμ· καὶ ὄνομα τῆς μητρὸς αὐτοὐ Γαζουβὰ, θυγάτηρ Σελί·
\vs{28b}καὶ ἐπορεύθη ἐν τῇ ὁδῷ Ἀσὰ τοῦ πατρὸς αὐτοῦ, καὶ οὐκ ἐξέκλινεν ἀπʼ αὐτῆς τοῦ ποιεῖν τὸ εὐθὲς ἐνώπιον Κυρίου· πλὴν τῶν ὑψηλῶν οὐκ ἐξῇραν· ἔθυον ἐν τοῖς ὑψηλοῖς, καὶ ἐθυμίων·
\vs{28c}καὶ ἃ συνέθετο Ἰωσαφὰτ μετὰ βασιλέως Ἰσραὴλ, καὶ πᾶσα ἡ δυναστεία ἣν ἐποίησε, καὶ οὓς ἐπολέμησεν, οὐκ ἰδοὺ ταῦτα γεγραμμένα ἐν βιβλίῳ λόγων τῶν ἡμερῶν τῶν βασιλέων Ἰούδα;
\vs{28d}καὶ τὰ λοιπὰ τῶν συμπλοκῶν ἃς ἐπέθεντο ἐν ταῖς ἡμέραις Ἀσὰ τοῦ πατρὸς αὐτοῦ ἐξῇρεν ἀπὸ τῆς γῆς·
\vs{28e}καὶ βασιλεὺς οὐκ ἦν ἐν Συρίᾳ· Νασίβ.

\vs{28f}Καὶ ὁ βασιλεὺς Ἰωσαφὰτ ἐποίησε ναῦν εἰς Θαρσὶς πορεύεσθαι εἰς Σωφὶρ ἐπὶ τὸ χρυσίον· καὶ οὐκ ἐπορεύθη, ὅτι συνετρίβη ἡ ναῦς ἐν Γασιὼν Γαβέρ·
\vs{28g}τότε εἶπεν βασιλεὺς Ἰσραὴλ πρὸς Ἰωσαφὰτ, ἐξαποστελῶ τοὺς παῖδάς σου καὶ τὰ παιδάριά μου ἐν τῇ νηΐ· καὶ οὐκ ἐβούλετο Ἰωσαφάτ·
\vs{28h}καὶ ἐκοιμήθη Ἰωσαφὰτ μετὰ τῶν πατέρων αὐτοῦ, καὶ θάπτεται μετὰ τῶν πατέρων αὐτοῦ ἐν πόλει Δαυίδ· καὶ ἐβασίλευσεν Ἰωρὰμ υἱὸς αὐτοῦ ἀντʼ αὐτοῦ.

\vs{29}Ἐν ἔτει δευτέρῳ τοῦ Ἰωσαφὰτ βασιλέως Ἰούδα, Ἀχαὰβ υἱὸς Ἀμβρὶ ἐβασίλευσεν ἐπὶ Ἰσραὴλ ἐν Σαμαρείᾳ εἴκοσι καὶ δύο ἔτη.
\vs{30}Καὶ ἐποίησεν Ἀχαὰβ τὸ πονηρὸν ἐνώπιον Κυρίου, καὶ ἐπονηρεύσατο ὑπὲρ πάντας τοὺς ἔμπροσθεν αὐτοῦ.
\vs{31}Καὶ οὐκ ἦν αὐτῷ ἱκανὸν τοῦ πορεύεσθαι ἐν ταῖς ἁμαρτίαις Ἱεροβοὰμ υἱοῦ Ναβὰτ, καὶ ἔλαβε γυναῖκα τὴν Ἰεζάβελ θυγατέρα Ἰεθεβαὰλ βασιλέως Σιδωνίων· καὶ ἐπορεύθη καὶ ἐδούλευσε τῷ Βάαλ, καὶ προσεκύνησεν αὐτῷ.
\vs{32}Καὶ ἔστησε θυσιαστήριον τῷ Βάαλ ἐν οἴκῳ τῶν προσοχθισμάτων αὐτοῦ, ὃν ᾠκοδόμησεν ἐν Σαμαρείᾳ.
\vs{33}Καὶ ἐποίησεν Ἀχαὰβ ἄλσος· καὶ προσέθηκεν Ἀχαὰβ τοῦ ποιῆσαι παροργίσματα, τοῦ παροργίσαι τὸν Κύριον Θεὸν τοῦ Ἰσραὴλ, καὶ τὴν ψυχὴν αὐτοῦ ἐξολοθρευθῆναι, ἐκακοποίησεν ὑπὲρ πάντας τοὺς βασιλεῖς Ἰσραὴλ τοὺς γενομένους ἔμπροσθεν αὐτοῦ.

\vs{34}Καὶ ἐν ταῖς ἡμέραις αὐτοῦ ᾠκοδόμησεν Ἀχιὴλ ὁ Βαιθηλίτης τὴν Ἱεριχώ· ἐν τῷ Ἀβιρὼν πρωτοτόκῳ αὐτοῦ ἐθεμελίωσεν αὐτὴν, καὶ τῷ Σεγοὺβ τῷ νεωτέρῳ αὐτοῦ ἐπέστησε θύρας αὐτῆς, κατὰ τὸ ῥῆμα Κυρίου, ὃ ἐλάλησεν ἐν χειρὶ Ἰησοῦ υἱοῦ Ναυῆ.

\ch{17}
Καὶ εἶπεν Ἠλιοὺ ὁ προφήτης Θεσβίτης ὁ ἐκ Θεσβῶν τῆς Γαλαὰδ πρὸς Ἀχαὰβ, ζῇ Κύριος ὁ Θεὸς τῶν δυνάμεων ὁ Θεὸς Ἰσραὴλ, ᾧ παρέστην ἐνώπιον αὐτοῦ, εἰ ἔσται τὰ ἔτη ταῦτα δρόσος καὶ ὑετὸς, ὅτι εἰ μὴ διὰ στόματος λόγου μου.

\vs{2}Καὶ ἐγένετο ῥῆμα Κυρίου πρὸς Ἠλιοὺ,
\vs{3}πορεύου ἐντεῦθεν κατὰ ἀνατολὰς, καὶ κρύβηθι ἐν τῷ χειμάῤῥῳ Χοῤῥαθ τοῦ ἐπὶ προσώπου τοῦ Ἰορδάνου.
\vs{4}Καὶ ἔσται ἐκ τοῦ χειμάῤῥου πίεσαι ὕδωρ, καὶ τοῖς κόραξιν ἐντελοῦμαι διατρέφειν σε ἐκεῖ.
\vs{5}Καὶ ἐποίησεν Ἠλιοὺ κατὰ τὸ ῥῆμα Κυρίου, καὶ ἐκάθισεν ἐν τῷ χειμάῤῥῳ Χοῤῥὰθ ἐπὶ προσώπου τοῦ Ἰορδάνου.
\vs{6}Καὶ οἱ κόρακες ἔφερον αὐτῷ ἄρτους τοπρωῒ, καὶ κρέα τοδείλης, καὶ ἐκ τοῦ χειμάῤῥου ἔπινεν ὕδωρ.
\vs{7}Καὶ ἐγένετο μεθʼ ἡμέρας, καὶ ἐξηράνθη ὁ χειμάῤῥους, ὅτι οὐκ ἐγένετο ὑετὸς ἐπὶ τῆς γῆς.

\vs{8}Καὶ ἐγένετο ῥῆμα Κυρίου πρὸς Ἠλιοὺ,
\vs{9}ἀνάστηθι, καὶ πορεύου εἰς Σαρεπτὰ τῆς Σιδωνίας· ἰδοὺ ἐντέταλμαι ἐκεῖ γυναικὶ χήρᾳ τοῦ διατρέφειν σε.
\vs{10}Καὶ ἀνέστη καὶ ἐπορεύθη εἰς Σαρεπτὰ, καὶ ἦλθεν εἰς τὸν πυλῶνα τῆς πόλεως· καὶ ἰδοὺ ἐκεῖ γυνὴ χήρα συνέλεγε ξύλα· καὶ ἐβόησεν ὀπίσω αὐτῆς Ἠλιοὺ, καὶ εἶπεν αὐτῇ, λάβε δή μοι ὀλίγον ὕδωρ εἰς ἄγγος, καὶ πίομαι.
\vs{11}Καὶ ἐπορεύθη λαβεῖν, καὶ ἐβόησεν ὀπίσω αὐτῆς Ἠλιοὺ, καὶ εἶπε, λήψῃ δή μοι ψωμὸν ἄρτου τοῦ ἐν τῇ χειρί σου.
\vs{12}Καὶ εἶπεν ἡ γυνὴ, ζῇ Κύριος ὁ Θεός σου, εἰ ἔστι μοι ἐγκρυφίας, ἀλλʼ ἢ ὅσον δρὰξ ἀλεύρου ἐν τῇ ὑδρίᾳ, καὶ ὀλίγον ἔλαιον ἐν τῷ καψάκῃ· καὶ ἰδοὺ ἐγὼ συλλέξω δύο ξυλάρια, καὶ εἰσελεύσομαι καὶ ποιήσω αὐτὸ ἐμαυτῇ καὶ τοῖς τέκνοις μου, καὶ φαγόμεθα, καὶ ἀποθανούμεθα.

\vs{13}Καὶ εἶπε πρὸς αὐτὴν Ἠλιοὺ, θάρσει, εἴσελθε καὶ ποίησον κατὰ τὸ ῥῆμά σου· ἀλλὰ ποίησόν μοι ἐκεῖθεν ἐγκρυφίαν μικρὸν, καὶ ἐξοίσεις μοι ἐν πρώτοις, σαυτῇ δὲ καὶ τοῖς τέκνοις σου ποιήσεις ἐπʼ ἐσχάτῳ·
\vs{14}Ὅτι τάδε λέγει Κύριος, ἡ ὑδρία τοῦ ἀλεύρου οὐκ ἐκλείψει, καὶ ὁ καψάκης τοῦ ἐλαίου οὐκ ἐλαττονήσει, ἕως ἡμέρας τοῦ δοῦναι Κύριον τὸν ὑετὸν ἐπὶ τῆς γῆς.
\vs{15}Καὶ ἐπορεύθη ἡ γυνὴ, καὶ ἐποίησε, καὶ ἤσθιεν αὐτὴ καὶ αὐτὸς καὶ τὰ τέκνα αὐτῆς.
\vs{16}Καὶ ἡ ὑδρία τοῦ ἀλεύρου οὐκ ἐξέλιπε, καὶ ὁ καψάκης τοῦ ἐλαίου οὐκ ἠλαττονήθη, κατὰ τὸ ῥῆμα Κυρίου ὃ ἐλάλησεν ἐν χειρὶ Ἠλιού.

\vs{17}Καὶ ἐγένετο μετὰ ταῦτα, καὶ ἠῤῥώστησεν ὁ υἱὸς τῆς γυναικὸς τῆς κυρίας τοῦ οἴκου· καὶ ἦν ἡ ἀῤῥωστία αὐτοῦ κραταιὰ σφόδρα ἕως οὐχ ὑπελείφθη ἐν αὐτῷ πνεῦμα.
\vs{18}Καὶ εἶπε πρὸς Ἠλιοὺ, τί ἐμοὶ καὶ σοὶ ἄνθρωπε τοῦ Θεοῦ; εἰσῆλθες πρὸς μὲ τοῦ ἀναμνῆσαι ἀδικίας μου, καὶ θανατῶσαι τὸν υἱόν μου;

\vs{19}Καὶ εἶπεν Ἠλιοὺ πρὸς τὴν γυναῖκα, δός μοι τὸν υἱόν σου· καὶ ἔλαβεν αὐτὸν ἐκ τοῦ κόλπου αὐτῆς, καὶ ἀνήνεγκεν αὐτὸν εἰς τὸ ὑπερῷον ἐν ᾧ αὐτὸς ἐκάθητο ἐκεῖ, καὶ ἐκοίμισεν αὐτὸν ἐπὶ τῆς κλίνης.
\vs{20}Καὶ ἀνεβόησεν Ἠλιοὺ, καὶ εἶπεν, οἵ μοι Κύριε, ὁ μάρτυς τῆς χήρας μεθʼ ἧς ἐγὼ κατοικῶ μετʼ αὐτῆς, σὺ κεκάκωκας τοῦ θανατῶσαι τὸν υἱὸν αὐτῆς.
\vs{21}Καὶ ἐνεφύσησε τῷ παιδαρίῳ τρὶς, καὶ ἐπεκαλέσατο τὸν Κύριον, καὶ εἶπε, Κύριε ὁ Θεός μου, ἐπιστραφήτω δὴ ἡ ψυχὴ τοῦ παιδαρίου τούτου εἰς αὐτόν.
\vs{22}Καὶ ἐγένετο οὕτως· καὶ ἀνεβόησε τὸ παιδάριον.
\vs{23}Καὶ κατήγαγεν αυτὸ ἀπὸ τοῦ ὑπερῴου εἰς τὸν οἶκον, καὶ ἔδωκεν αὐτὸ τῇ μητρὶ αὐτοῦ· καὶ εἶπεν Ἠλιοὺ, βλέπε, ζῇ ὁ υἱός σου.
\vs{24}Καὶ εἶπεν ἡ γυνὴ πρὸς Ἠλιοὺ, ἰδοὺ ἔγνωκα ὅτι σὺ ἄνθρωπος Θεοῦ, καὶ ῥῆμα Κυρίου ἐν τῷ στόματί σου ἀληθινόν.

\ch{18}
Καὶ ἐγένετο μεθʼ ἡμέρας πολλὰς, καὶ ῥῆμα Κυρίου ἐγένετο πρὸς Ἠλιοὺ ἐν τῷ ἐνιαυτῷ τῷ τρίτῳ, λέγων, πορεύθητι, καὶ ὄφθητι τῷ Ἀχαὰβ, καὶ δώσω ὑετὸν ἐπὶ πρόσωπον τῆς γῆς.
\vs{2}Καὶ ἐπορεύθη Ἠλιοὺ τοῦ ὀφθῆναι τῷ Ἀχαὰβ, καὶ ἡ λιμὸς κραταιὰ ἐν Σαμαρείᾳ.

\vs{3}Καὶ ἐκάλεσεν Ἀχαὰβ τὸν Ἀβδιοὺ τὸν οἰκονόμον· καὶ Ἀβδιοὺ ἦν φοβούμενος τὸν Κύριον σφόδρα.
\vs{4}Καὶ ἐγένετο ἐν τῷ τύπτειν τὴν Ἰεζάβελ τοὺς προφήτας Κυρίου, καὶ ἔλαβεν Ἀβδιοὺ ἑκατὸν ἄνδρας προφήτας καὶ κατέκρυψεν αὐτοὺς κατὰ πεντήκοντα ἐν σπηλαίῳ, καὶ διέτρεφεν αὐτοὺς ἐν ἄρτῳ καὶ ὕδατι.
\vs{5}Καὶ εἶπεν Ἀχαὰβ πρὸς Ἀβδιοὺ, δεῦρο, καὶ διέλθωμεν ἐπὶ τὴν γῆν καὶ ἐπὶ πηγὰς τῶν ὑδάτων καὶ ἐπὶ χειμάῤῥους, ἐὰν πῶς εὕρωμεν βοτάνην, καὶ περιποιησώμεθα ἵππους καὶ ἡμιόνους, καὶ οὐκ ἐξολοθρευθήσονται ἀπὸ τῶν σκηνῶν.
\vs{6}Καὶ ἐμέρισαν ἑαυτοῖς τὴν ὁδὸν τοῦ διελθεῖν αὐτήν· Ἀχαὰβ ἐπορεύθη ἐν ὁδῷ μιᾷ, καὶ Ἀβδιοὺ ἐπορεύθη ἐν ὁδῷ ἄλλῃ μόνος.
\vs{7}Καὶ ἦν Ἀβδιοὺ ἐν τῇ ὁδῷ μόνος· καὶ ἦλθεν Ἠλιοὺ εἰς συνάντησιν αὐτοῦ μόνος· καὶ Ἀβδιοὺ ἔσπευσε καὶ ἔπεσεν ἐπὶ πρόσωπον αὐτοῦ, καὶ εἶπεν, εἰ σὺ εἶ αὐτὸς, κύριέ μου Ἠλιού;
\vs{8}Καὶ εἶπεν Ἠλιοὺ αὐτῷ, ἑγώ· πορεύου, λέγε τῷ κυρίῳ σου, ἰδοὺ Ἠλιού.
\vs{9}Καὶ εἶπεν Ἀβδιοὺ, τί ἡμάρτηκα, ὅτι δίδως τὸν δοῦλόν σου εἰς χεῖρα Ἀχαὰβ τοῦ θανατῶσαί με;
\vs{10}Ζῇ Κύριος ὁ Θεός σου, εἰ ἔστιν ἔθνος ἢ βασιλεία, οὗ οὐκ ἀπέστειλεν ὁ κύριός μου ζητεῖν σε· καὶ εἰ εἶπον, οὐκ ἔστι, καὶ ἐνέπρησε τὴν βασιλείαν καὶ τὰς χώρας αὐτῆς, ὅτι οὐχ εὕρηκέ σε.
\vs{11}Καὶ νῦν σὺ λέγεις, πορεύου, ἀνάγγελε τῷ κυρίῳ σου, ἰδοὺ Ἠλιού.
\vs{12}Καὶ ἔσται ἐὰν ἐγὼ ἀπέλθω ἀπὸ σοῦ, καὶ πνεῦμα Κυρίου ἀρεῖ σε εἰς τὴν γῆν ἣν οὐκ οἶδα, καὶ εἰσελεύσομαι ἀπαγγεῖλαι τῷ Ἀχαὰβ, καὶ οὐχ εὑρήσει σε, καὶ ἀποκτενεῖ με· καὶ ὁ δοῦλός σου ἐστὶ φοβούμενος τὸν Κύριον ἐκ νεότητος αὐτοῦ.
\vs{13}Ἢ οὐκ ἀπηγγέλη σοι τῷ κυρίῳ μου, οἷα πεποίηκα ἐν τῷ ἀποκτείνειν τὴν Ἰεζάβελ τοὺς προφήτας Κυρίου, καὶ ἔκρυψα ἀπὸ τῶν προφητῶν Κυρίου ἑκατὸν ἄνδρας, ἀνὰ πεντήκοντα ἐν σπηλαίῳ, καὶ ἔθρεψα ἐν ἄρτοις καὶ ὕδατι;
\vs{14}Καὶ νῦν σὺ λέγεις μοι, πορεύου, λέγε τῷ κυρίῳ σου, ἰδοὺ Ἠλιού· καὶ ἀποκτενεῖ με.
\vs{15}Καὶ εἶπεν Ἠλιοὺ, ζῇ Κύριος τῶν δυνάμεων ᾧ παρέστην ἐνώπιον αὐτοῦ, ὅτι σήμερον ὀφθήσομαι αὐτῷ.

\vs{16}Καὶ ἐπορεύθη Ἀβδιοὺ εἰς συναντὴν τῷ Ἀχαὰβ, καὶ ἀπήγγειλεν αὐτῷ· καὶ ἐξέδραμεν Ἀχαὰβ, καὶ ἐπορεύθη εἰς συνάντησιν Ἠλιού.
\vs{17}Καὶ ἐγένετο ὡς εἶδεν Ἀχαὰβ τὸν Ἠλιού, καὶ εἶπεν Ἀχαὰβ πρὸς Ἠλιοὺ, εἰ σὺ εἶ αὐτὸς ὁ διαστρέφων τὸν Ἰσραήλ;
\vs{18}Καὶ εἶπεν Ἠλιοὺ, οὐ διαστρέφω τὸν Ἰσραὴλ, ὅτι ἀλλʼ ἢ σὺ καὶ οἶκος τοῦ πατρός σου, ἐν τῷ καταλιμπάνειν ὑμᾶς τὸν Κύριον Θεὸν ὑμῶν, καὶ ἐπορεύθης ὀπίσω τῶν Βααλίμ.
\vs{19}Καὶ νῦν ἀπόστειλον, συνάθροισον πρὸς μὲ πάντα Ἰσραὴλ εἰς ὄρος τὸ Καρμήλιον, καὶ τοὺς προφήτας τῆς αἰσχύνης τετρακοσίους καὶ πεντήκοντα, καὶ τοὺς προφήτας τῶν ἀλσῶν τετρακοσίους ἐσθίοντας τράπεζαν Ἰεζάβελ.

\vs{20}Καὶ ἀπέστειλεν Ἀχαὰβ εἰς πάντα Ἰσραὴλ, καὶ ἐπισυνήγαγεπάντας τοὺς προφήτας εἰς ὄρος τὸ Καρμήλιον.

\vs{21}Καὶ προσήγαγεν Ἠλιοὺ πρὸς πάντας· καὶ εἶπεν αὐτοῖς Ἠλιοὺ, ἕως πότε ὑμεῖς χωλανεῖτε ἐπʼ ἀμφοτέραις ταῖς ἰγνύαις· εἰ ἔστι Κύριος ὁ Θεὸς, πορεύεσθε ὀπίσω αὐτοῦ· εἰ δὲ Βάαλ, πορεύεσθε ὀπίσω αὐτοῦ· καὶ οὐκ ἀπεκρίθη ὁ λαὸς λόγον.
\vs{22}Καὶ εἶπεν Ἠλιοὺ πρὸς τὸν λαόν, ἐγὼ ὑπολέλειμμαι προφήτης τοῦ Κυρίου μονώτατος· καὶ οἱ προφῆται τοῦ Βάαλ τετρακόσιοι καὶ πεντήκοντα ἄνδρες, καὶ οἱ προφῆται τοῦ ἄλσους τετρακόσιοι·
\vs{23}Δότωσαν ἡμῖν δύο βόας, καὶ ἐκλεξάσθωσαν ἑαυτοῖς τὸν ἕνα, καὶ μελισάτωσαν, καὶ ἐπιθέτωσαν ἐπὶ τῶν ξύλων, καὶ πῦρ μὴ ἐπιθέτωσαν· καὶ ἐγὼ ποιήσω τὸν βοῦν τὸν ἄλλον, καὶ πῦρ οὐ μὴ ἐπιθῶ.
\vs{24}Καὶ βοᾶτε ἐν ὀνόματι θεῶν ὑμῶν, καὶ ἐγὼ ἐπικαλέσομαι ἐν τῷ ὀνόματι Κυρίου τοῦ Θεοῦ μου· καὶ ἔσται ὁ θεὸς ὃς ἂν ἐπακούσῃ ἐν πυρὶ, οὗτος Θεός· καὶ ἀπεκρίθησαν πᾶς ὁ λαὸς, καὶ εἶπον, καλὸν τὸ ῥῆμα ὃ ἐλάλησας.

\vs{25}Καὶ εἶπεν Ἠλιοὺ τοῖς προφήταις τῆς αἰσχύνης, ἐκλέξασθε ἑαυτοῖς τὸν μόσχον τὸν ἕνα, καὶ ποιήσατε πρῶτοι, ὅτι πολλοὶ ὑμεῖς· καὶ ἐπικαλέσασθε ἐν ὀνόματι θεοῦ ὑμῶν, καὶ πῦρ μὴ ἐπιθῆτε.
\vs{26}Καὶ ἔλαβον τὸν μόσχον καὶ ἐποίησαν, καὶ ἐπεκαλοῦντο ἐν ὀνόματι τοῦ Βάαλ ἐκ πρωΐθεν ἕως μεσημβρίας, καὶ εἶπον, ἐπάκουσον ἡμῶν ὁ Βάαλ, ἐπάκουσον ἡμῶν· καὶ οὐκ ἦν φωνὴ, καὶ οὐκ ἦν ἀκρόασις· καὶ διέτρεχον ἐπὶ τοῦ θυσιαστηρίου οὗ ἐποίησαν.
\vs{27}Καὶ ἐγένετο μεσημβρία, καὶ ἐμυκτήρισεν αὐτοὺς Ἠλιοὺ ὁ Θεσβίτης, καὶ εἶπεν, ἐπικαλεῖσθε ἐν φωνῇ μεγάλῃ, ὅτι θεός ἐστιν· ὅτι ἀδολεσχία αὐτῷ ἐστι, καὶ ἅμα μή ποτε χρηματίζει αὐτὸς, ἢ μή ποτε καθεύδει αὐτὸς, καὶ ἐξαναστήσεται.
\vs{28}Καὶ ἐπεκαλοῦντο ἐν φωνῇ μεγάλῃ, καὶ κατετέμνοντο κατὰ τὸν ἐθισμὸν αὐτῶν ἐν μαχαίραις καὶ σειρομάσταις ἕως ἐκχύσεως αἵματος ἐπʼ αὐτοὺς,
\vs{29}καὶ προεφήτευον ἕως οὗ παρῆλθε τὸ δειλινόν· καὶ ἐγένετο ὡς ὁ καιρὸς τοὺ ἀναβῆναι τὴν θυσίαν, καὶ ἐλάλησεν Ἠλιοὺ ὁ Θεσβίτης πρὸς τοὺς προφήτας τῶν προσοχθισμάτων, λέγων, μετάστητε ἀπὸ τοῦ νῦν, καὶ ἐγὼ ποιήσω τὸ ὁλοκαύτωμά μου· καὶ μετέστησαν, καὶ ἀπῆλθον.

\vs{30}Καὶ εἶπεν Ἠλιοὺ πρὸς τὸν λαὸν, προσαγάγετε πρὸς μέ· καὶ προσήγαγε πᾶς ὁ λαὸς πρὸς αὐτόν.
\vs{31}Καὶ ἔλαβεν Ἡλιοὺ δώδεκα λίθους κατὰ ἀριθμὸν φυλῶν τοῦ Ἰσραὴλ, ὡς ἐλάλησε Κύριος πρὸς αὐτὸν, λέγων, Ἰσραὴλ ἔσται τὸ ὄνομά σου.
\vs{32}Καὶ ᾠκοδόμησε τοὺς λίθους ἐν ὀνόματι Κυρίου, καὶ ἰάσατο τὸ θυσιαστήριον τὸ κατεσκαμμένον· καὶ ἐποίησε θάλασσαν χωροῦσαν δύο μετρητὰς σπέρματος κυκλόθεν τοῦ θυσιαστηρίου·
\vs{33}Καὶ ἐστοίβασε τὰς σχίδακας ἐπὶ τὸ θυσιαστήριον ὃ ἐποίησε, καὶ ἐμέλισε τὸ ὁλοκαύτωμα καὶ ἐπέθηκεν ἐπὶ τὰς σχίδακας, καὶ ἐστοίβασεν ἐπὶ τὸ θυσιαστήριον,
\vs{34}καὶ εἶπε, λάβετέ μοι τέσσαρας ὑδρίας ὕδατος, καὶ ἐπιχέετε ἐπὶ τὸ ὁλοκαύτωμα καὶ ἐπὶ τὰς σχίδακας· καὶ ἐποίησαν οὕτως. Καὶ εἶπε, δευτερώσατε· καὶ ἐδευτέρωσαν· καὶ εἶπε, τρισσώσατε· καὶ ἐτρίσσευσαν.
\vs{35}Καὶ διεπορεύετο τὸ ὕδωρ κύκλῳ τοῦ θυσιαστηρίου, καὶ τὴν θάλασσαν ἔπλησαν ὕδατος.

\vs{36}Καὶ ἀνεβόησεν Ἠλιοὺ εἰς τὸν οὐρανὸν, καὶ εἶπε, Κύριε ὁ Θεὸς Ἁβραὰμ καὶ Ἰσαὰκ καὶ Ἰσραήλ, ἐπάκουσόν μου Κύριε, ἐπάκουσόν μου σήμερον ἐν πυρί, καὶ γνώτωσαν πᾶς ὁ λαὸς οὗτος, ὅτι σὺ εἶ Κύριος ὁ Θεὸς Ἰσραὴλ, καὶ ἐγὼ δοῦλός σου, καὶ διὰ σὲ πεποίηκα τὰ ἔργα ταῦτα.
\vs{37}Ἐπάκουσόν μου Κύριε, ἐπάκουσόν μου, καὶ γνώτω ὁ λαὸς οὗτος, ὅτι σὺ εἶ Κύριος ὁ Θεὸς, καὶ σὺ ἔστρεψας τὴν καρδίαν τοῦ λαοῦ τούτου ὀπίσω.
\vs{38}Καὶ ἔπεσε πῦρ παρὰ Κυρίου ἐκ τοῦ οὐρανοῦ, καὶ κατέφαγε τὰ ὁλοκαυτώματα καὶ τὰς σχίδακας καὶ τὸ ὕδωρ τὸ ἐν τῇ θαλάσσῃ, καὶ τοὺς λίθους καὶ τὸν χοῦν ἐξέλειξε τὸ πῦρ.

\vs{39}Καὶ ἔπεσε πᾶς ὁ λαὸς ἐπὶ πρόσωπον αὐτῶν, καὶ εἶπον, ἀληθῶς Κύριος ὁ Θεὸς αὐτὸς ὁ Θεός.
\vs{40}Καὶ εἶπεν Ἠλιοὺ πρὸς τὸν λαὸν, συλλάβετε τοὺς προφήτας τοῦ Βάαλ, μηδεὶς σωθήτω ἐξ αὐτῶν· καὶ συνελαβον αὐτοὺς, καὶ κατάγει αὐτοὺς Ἠλιοὺ εἰς τὸν χειμάῤῥουν Κισσῶν, καὶ ἔσφαξεν αὐτοὺς ἐκεῖ.

\vs{41}Καὶ εἶπεν Ἠλιοὺ τῷ Ἀχαὰβ, ἀνάβηθι, καὶ φάγε καὶ πίε, ὅτι φωνὴ τῶν ποδῶν τοῦ ὑετοῦ.
\vs{42}Καὶ ἀνέβη Ἀχαὰβ τοῦ φαγεῖν καὶ πιεῖν· καὶ Ἠλιοὺ ἀνέβη ἐπὶ τὴν Κάρμηλον καὶ ἔκυψεν ἐπὶ τὴν γῆν, καὶ ἔθηκε τὸ πρόσωπον αὐτοῦ ἀναμέσον τῶν γονάτων αὐτοῦ,
\vs{43}καὶ εἶπε τῷ παιδαρίῳ αὐτοῦ, ἀνάβηθι, καὶ ἐπίβλεψον ὁδὸν τῆς θαλάσσης· καὶ ἐπέβλεψε τὸ παιδάριον, καὶ εἶπεν, οὐκ ἔστιν οὐθέν· καὶ εἶπεν Ἠλιοὺ, καὶ σὺ ἐπίστρεψον ἑπτάκις.
\vs{44}Καὶ ἐπέστρεψε τὸ παιδάριον ἑπτάκις· καὶ ἐγένετο ἐν τῷ ἑβδόμῳ, καὶ ἰδοὺ νεφέλη μικρὰ ὡς ἴχνος ἀνδρὸς ἀνάγουσα ὕδωρ· καὶ εἶπεν, ἀνάβηθι, καὶ εἶπον Ἀχαὰβ, ζεῦξον τὸ ἅρμα σου καὶ κατάβηθι, μὴ καταλάβῃ σε ὁ ὑετός.
\vs{45}Καὶ ἐγένετο ἕως ὧδε καὶ ὧδε, καὶ ὁ οὐρανὸς συνεσκότασε νεφέλαις καὶ πνεύματι, καὶ ἐγένετο ὑετὸς μέγας· καὶ ἔκλαιε καὶ ἐπορεύετο Ἀχαὰβ ἕως Ἰεζράελ.
\vs{46}Καὶ χεὶρ Κυρίου ἐπὶ τὸν Ἠλιοὺ, καὶ συνέσφιγξε τὴν ὀσφὺν αὐτοῦ, καὶ ἔτρεχεν ἔμπροσθεν Ἀχαὰβ εἰς Ἰεζράελ.

\ch{19}
Καὶ ἀνήγγειλεν Ἀχαὰβ τῇ Ἰεζάβελ γυναικὶ αὐτοῦ πάντα ἃ ἐποίησεν Ἠλιοὺ, καὶ ὡς ἀπέκτεινε τοὺς προφήτας ἐν ῥομφαίᾳ.
\vs{2}Καὶ ἀπέστειλεν Ἰεζάβελ πρὸς Ἠλιοὺ, καὶ εἶπεν, εἰ σὺ εἶ Ἠλιοὺ καὶ ἐγὼ Ἰεζάβελ, τάδε ποιήσαι μοι ὁ θεὸς καὶ τάδε προσθείη, ὅτι ταύτην τὴν ὥραν αὔριον θήσομαι τὴν ψυχήν σου καθὼς ψυχὴν ἑνὸς ἐξ αὐτῶν.
\vs{3}Καὶ ἐφοβήθη Ἠλιοὺ, καὶ ἀνέστη καὶ ἀπῆλθε κατὰ τὴν ψυχὴν αὐτοῦ, καὶ ἔρχεται εἰς Βηρσαβεὲ γῆν Ἰούδα, καὶ ἀφῆκε τὸ παιδάριον αὐτοῦ ἐκεῖ.

\vs{4}Καὶ αὐτὸς ἐπορεύθη ἐν τῇ ἐρήμῳ ὁδὸν ἡμέρας, καὶ ἦλθε καὶ ἐκάθισεν ὑποκάτω ῥαθμὲν, καὶ ᾐτήσατο τὴν ψυχὴν αὐτοῦ ἀποθανεῖν· καὶ εἶπεν, ἱκανούσθω νῦν, λάβε δὴ τὴν ψυχήν μου ἀπʼ ἐμοῦ Κύριε, ὅτι οὐ κρείσσων ἐγώ εἰμι ὑπὲρ τοὺς πατέρας μου.
\vs{5}Καὶ ἐκοιμήθη καὶ ὕπνωσεν ἐκεῖ ὑπὸ φυτόν· καὶ ἰδού τις ἥψατο αὐτοῦ, καὶ εἶπεν αὐτῷ, ἀνάστηθι καὶ φάγε.
\vs{6}Καὶ ἐπέβλεψεν Ἠλιού· καὶ ἰδοὺ πρὸς κεφαλῆς αὐτοῦ ἐγκρυφίας ὀλυρίτης καὶ καψάκης ὕδατος· καὶ ἀνέστη καὶ ἔφαγε καὶ ἔπιε, καὶ ἐπιστρέψας ἐκοιμήθη.
\vs{7}Καὶ ἐπέστρεψεν ὁ ἄγγελος Κυρίου ἐκ δευτέρου, καὶ ἥψατο αὐτοῦ, καὶ εἶπεν αὐτῷ, ἀνάστα, φάγε, ὅτι πολλὴ ἀπὸ σοῦ ἡ ὁδός.
\vs{8}Καὶ ἀνέστη, καὶ ἔφαγε, καὶ ἔπιε· καὶ ἐπορεύθη ἐν ἰσχύϊ τῆς ρώσεως ἐκείνης τεσσαράκοντα ἡμέρας καὶ τεσσαράκοντα νύκτας ἕως ὄρους Χωρήβ.

\vs{9}Καὶ εἰσῆλθεν ἐκεῖ εἰς τὸ σπήλαιον, καὶ κατέλυσεν ἐκεῖ· καὶ ἰδοὺ ῥῆμα Κυρίου πρὸς αὐτὸν, καὶ εἶπε, τί σὺ ἐνταῦθα Ἠλιού;
\vs{10}Καὶ εἶπεν Ἠλιοὺ, ζηλῶν ἐζήλωκα τῷ Κυρίῳ παντοκράτορι, ὅτι ἐγκατέλιπόν σε οἱ υἱοὶ Ἰσραήλ· τὰ θυσιαστήριά σου κατέσκαψαν, καὶ τοὺς προφήτας σου ἀπέκτειναν ἐν ῥομφαίᾳ καὶ ὑπολέλειμμαι ἐγὼ μονώτατος, καὶ ζητοῦσι τὴν ψυχήν μου λαβεῖν αὐτήν.
\vs{11}Καὶ εἶπεν, ἐξελεύσῃ αὔριον, καὶ στήσῃ ἐνώπιον Κυρίου ἐν τῷ ὄρει· ἰδοὺ παρελεύσεται Κύριος. Καὶλἰδοὺ πνεῦμα μέγα κραταιὸν διαλύον ὄρη καὶ συντρίβον πέτρας ἐνώπιον Κυρίου, οὐκ ἐν τῷ πνεύματι Κύριος· καὶ μετὰ τό πνεῦμα συσσεισμὸς, οὐκ ἐν τῷ συσσεισμῷ Κύριος·
\vs{12}Καὶ μετὰ τὸν συσσεισμὸν πῦρ, οὐκ ἐν τῷ πυρὶ Κύριος· καὶ μετὰ τὸ πῦρ φωνὴ αὔρας λεπτῆς.

\vs{13}Καὶ ἐγένετο ὡς ἤκουσεν Ἠλιοὺ, καὶ ἐπεκάλυψε τὸ πρόσωπον αὐτοῦ ἐν τῇ μηλωτῇ αὐτοῦ, καὶ ἐξῆλθε καὶ ἔστη ὑπὸ σπήλαιον· καὶ ἰδοὺ πρὸς αὐτὸν φωνὴ, καὶ εἶπε, τί σὺ ἐνταῦθα Ἠλιού;
\vs{14}Καὶ εἶπεν Ἠλιού, ζηλῶν ἐζήλωκα τῷ Κυρίῳ παντοκράτορι, ὅτι ἐγκατέλιπον τὴν διαθήκην σου οἱ υἱοὶ Ἰσραήλ· καὶ τὰ θυσιαστήριά σου καθεῖλαν, καὶ τοὺς προφήτας σου ἀπέκτειναν ἐν ῥομφαίᾳ, καὶ ὑπολέλειμμαι ἐγὼ μονώτατος, καὶ ζητοῦσι τὴν ψυχήν μου λαβεῖν αὐτήν.
\vs{15}Καὶ εἶπε Κύριος πρὸς αὐτὸν, πορεύου, ἀνάστρεφε εἰς τὴν ὁδόν σου, καὶ ἥξεις εἰς τὴν ὁδὸν ἐρήμου Δαμασκοῦ· καὶ ἥξεις καὶ χρίσεις τὸν Ἀζαὴλ εἰς βασιλέα τῆς Συρίας·
\vs{16}Καὶ τὸν Ἰοὺ υἱὸν Ναμεσσὶ χρίσεις εἰς βασιλέα ἐπὶ Ἰσραήλ· καὶ τὸν Ἑλισαιὲ υἱὸν Σαφὰτ χρίσεις εἰς προφήτην ἀντὶ σοῦ.
\vs{17}Καὶ ἔσται τὸν σωζόμενον ἐκ ῥομφαίας Ἀζαὴλ, θανατώσει Ἰού· καὶ τὸν σωζόμενον ἐκ ῥομφαίας Ἰοὺ, θανατώσει Ἑλισαιέ.
\vs{18}Καὶ καταλείψεις ἐν Ἰσραὴλ ἑπτὰ χιλιάδας ἀνδρῶν, πάντα γόνατα ἃ οὐκ ὤκλασαν γόνυ τῷ Βάαλ, καὶ πᾶν στόμα ὃ οὐ προσεκύνησεν αὐτῷ.

\vs{19}Καὶ ἀπῆλθεν ἐκεῖθεν καὶ εὑρίσκει τὸν Ἑλισαιὲ υἱὸν Σαφὰτ, καὶ αὐτὸς ἠροτρία ἐν βουσί· δώδεκα ζεύγη ἐνώπιον αὐτοῦ, καὶ αὐτὸς ἐν τοῖς δώδεκα· καὶ ἀπῆλθεν ἐπʼ αὐτὸν, καὶ ἐπέῥῥιψε τὴν μηλωτὴν αὐτοῦ ἐπʼ αὐτόν.
\vs{20}Καὶ κατέλιπεν Ἑλισαιὲ τὰς βόας, καὶ κατέδραμεν ὀπίσω Ἠλιοὺ, καὶ εἶπε, καταφιλήσω τὸν πατέρα μου, καὶ ἀκολουθήσω ὀπίσω σοῦ· καὶ εἶπεν Ἠλιοὺ, ἀνάστρεφε, ὅτι πεποίηκά σοι.
\vs{21}Καὶ ἀνέστρεψεν ἐξ ὄπισθεν αὐτοῦ· καὶ ἔλαβε τὰ ζεύγη τῶν βοῶν, καὶ ἔθυσε καὶ ἥψησεν αὐτὰ ἐν τοῖς σκεύεσι τῶν βοῶν, καὶ ἔδωκε τῷ λαῷ, καὶ ἔφαγον· καὶ ἀνέστη καὶ ἐπορεύθη ὀπίσω Ἠλιοὺ, καὶ ἐλειτούργει αὐτῷ.

\ch{20}
Καὶ ἀμπελὼν εἷς ἦν τῷ Ναβουθαὶ τῷ Ἰεζραηλίτῃ παρὰ τῇ ἅλῳ Ἀχαὰβ βασιλέως Σαμαρείας.
\vs{2}Καὶ ἐλάλησεν Ἀχαὰβ πρὸς Ναβουθαὶ, λέγων, δός μοι τὸν ἀμπελῶνά σου, καὶ ἔσται μοι εἰς κῆπον λαχάνων, ὅτι ἐγγίζων οὗτος τῷ οἴκῳ μου, καὶ δώσω σοι ἀμπελῶνα ἄλλον ἀγαθὸν ὑπὲρ αὐτόν· εἰ δὲ ἀρέσκει ἐνώπιόν σου, δώσω σοι ἀργύριον ἄλλαγμα ἀμπελῶνός σου τούτου, καὶ ἔσται μοι εἰς κῆπον λαχάνων.
\vs{3}Καὶ εἶπε Ναβουθαὶ πρὸς Ἀχαάβ, μὴ γένοιτό μοι παρὰ Θεοῦ μου δοῦναι κληρονομίαν πατέρων μου σοί.

\vs{4}Καὶ ἐγένετο τὸ πνεῦμα Ἀχαὰβ τεταραγμένον, καὶ ἐκοιμήθη ἐπὶ τῆς κλίυης αὐτοῦ, καὶ συνεκάλυψε τὸ πρόσωπον αὐτοῦ, καὶ οὐκ ἔφαγεν ἄρτον.
\vs{5}Καὶ εἰσῆλθεν Ἰεζάβελ ἡ γυνὴ αὐτοῦ πρὸς αὐτὸν, καὶ ἐλάλησε πρὸς αὐτὸν, τί τὸ πνεῦμά σου τεταραγμένον, καὶ οὐκ εἶ σὺ ἐσθίων ἄρτον;
\vs{6}Καὶ εἶπε πρὸς αὐτὴν, ὅτι ἐλάλησα πρὸς Ναβουθαὶ τὸν Ἰεζραηλίτην, λέγων, δός μοι τὸν ἀμπελῶνά σου ἀργυρίου· εἰ δὲ βούλῃ, δώσω σοι ἀμπελῶνα ἄλλον ἀντʼ αὐτοῦ· καὶ εἶπεν οὐ δώσω σοι κληρονομίαν πατέρων μου.
\vs{7}Καὶ εἶπε πρὸς αὐτὸν Ἰεζάβελ ἡ γυνὴ αὐτοῦ, σὺ νῦν οὕτω ποιεῖς βασιλέα ἐπὶ Ἰσραήλ; ἀνάστηθι καὶ φάγε ἄρτον καὶ σαυτοῦ γενοῦ, ἐγὼ δὲ δώσω σοι τὸν ἀμπελῶνα Ναβουθαὶ τοῦ Ἰεζραηλίτου.

\vs{8}Καὶ ἔγραψε βιβλίον ἐπὶ τῷ ὀνόματι Ἀχαὰβ, καὶ ἐσφραγίσατο τῇ σφραγίδι αὐτοῦ· καὶ ἀπέστειλε τὸ βιβλίον πρὸς τοὺς πρεσβυτέρους καὶ τοὺς ἐλευθέρους τοὺς κατοικοῦντας μετὰ Ναβουθαί.
\vs{9}Καὶ ἐγέγραπτο ἐν τοῖς βιβλίοις, λέγων, νηστεύσατε νηστείαν, καὶ καθίσατε τὸν Ναβουθαὶ ἐν ἀρχῇ τοῦ λαοῦ·
\vs{10}Καὶ ἐγκαθίσατε δύο ἄνδρας υἱοὺς παρανόμων ἐξεναντίας αὐτοῦ, καὶ καταμαρτυρησάτωσαν αὐτοῦ, λέγοντες, εὐλόγησε Θεὸν καὶ βασιλέα· καὶ ἐξαγαγέτωσαν αὐτὸν, καὶ λιθοβολησάτωσαν αὐτὸν, καὶ ἀποθανέτω.

\vs{11}Καὶ ἐποίησαν οἱ ἄνδρες τῆς πόλεως αὐτοῦ οἱ πρεσβύτεροι καὶ οἱ ἐλεύθεροι οἱ κατοικοῦντες ἐν τῇ πόλει αὐτοῦ, καθὼς ἀπέστειλε πρὸς αὐτοὺς Ἰεζάβελ, καὶ καθὰ ἐγέγραπτο ἐν τοῖς βιβλίοις οἷς ἀπέστειλε πρὸς αὐτούς.
\vs{12}Καὶ ἐκάλεσαν νηστείαν, καὶ ἐκάθισαν τὸν Ναβουθαὶ ἐν ἀρχῇ τοῦ λαοῦ.
\vs{13}Καὶ εἰσῆλθον δύο ἄνδρες υἱοὶ παρανόμων, καὶ ἐκάθισαν ἐξεναντίας αὐτοῦ, καὶ κατεμαρτύρησαν αὐτοῦ, λέγοντες, εὐλόγηκας Θεὸν καὶ βασιλέα· καὶ ἐξήγαγον αὐτὸν ἔξω τῆς πόλεως, καὶ ἐλιθοβόλησαν αὐτὸν ἐν λίθοις, καὶ ἀπέθανε.
\vs{14}Καὶ ἀπέστειλαν πρὸς Ἰεζάβελ, λέγοντες, λελιθοβόληται Ναβουθαὶ, καὶ τέθνηκε.

\vs{15}Καὶ ἐγένετο ὡς ἤκουσεν Ἰεζάβελ, καὶ εἶπε πρὸς Ἀχαὰβ, ἀνάστα, κληρονόμει τὸν ἀμπελῶνα Ναβουθαὶ τοῦ Ἰεζραηλίτου, ὃς οὐκ ἔδωκέ σοι ἀργυρίου, ὅτι οὐκ ἔστι Ναβουθαὶ ζῶν, ὅτι τέθνηκε.
\vs{16}Καὶ ἐγένετο ὡς ἤκουσεν Ἀχαὰβ ὅτι τέθνηκε Ναβουθαὶ ὁ Ἰεζραηλίτης, καὶ διέῤῥηξε τὰ ἱμάτια αὐτοῦ, καὶ περιεβάλετο σάκκον· καὶ ἐγένετο μετὰ ταῦτα, καὶ ἀνέστη καὶ κατέβη Ἀχαὰβ εἰς τὸν ἀμπελῶνα Ναβουθαὶ τοῦ Ἰεζραηλίτου κληρονομῆσαι αὐτόν.

\vs{17}Καὶ εἶπε Κύριος πρὸς Ἠλιοὺ τὸν Θεσβίτην, λέγων,
\vs{18}ἀνάστηθι καὶ κατάβηθι εἰς ἀπαντὴν Ἀχαὰβ βασιλέως Ἰσραὴλ τοῦ ἐν Σαμαρείᾳ, ὅτι οὗτος ἐν ἀμπελῶνι Ναβουθαὶ, ὅτι καταβέβηκεν ἐκεῖ κληρονομῆσαι αὐτόν.
\vs{19}Καὶ λαλήσεις πρὸς αὐτὸν, λέγων, τάδε λέγει Κύριος, ὡς σὺ ἐφόνευσας καὶ ἐκληρονόμησας, διὰ τοῦτο τάδε λέγει Κύριος, ἐν παντὶ τόπῳ ᾧ ἔλειξαν αἱ ὗες καὶ οἱ κύνες τὸ αἷμα Ναβουθαὶ, ἐκεῖ λείξουσιν οἱ κύνες τὸ αἷμά σου, καὶ αἱ πόρναι λούσονται ἐν τῷ αἵματί σου.
\vs{20}Καὶ εἶπεν Ἀχαὰβ πρὸς Ἠλιοὺ, εἰ εὕρηκάς με ὁ ἐχθρός μου; καὶ εἶπεν, εὕρηκα· διότι μάτην πέπρασαι ποιῆσαι τὸ πονηρὸν ἐνώπιον Κυρίου, παροργίσαι αὐτόν.
\vs{21}Ἰδοὺ ἐγὼ ἐπάγω ἐπὶ σὲ κακά· καὶ ἐκκαύσω ὀπίσω σου, καὶ ἐξολοθρεύσω τοῦ Ἀχαὰβ οὐροῦντα πρὸς τοῖχον, καὶ συνεχόμενον καὶ ἐγκαταλελειμμένον ἐν Ἰσραήλ.
\vs{22}Καὶ δώσω τὸν οἶκόν σου ὡς τὸν οἶκον Ἱεροβοὰμ υἱοῦ Ναβὰτ, καὶ ὡς τὸν οἶκον Βαασὰ υἱοῦ Ἀχιὰ, περὶ τῶν παροργισμάτων ὧν παρώργισας καὶ ἐξήμαρτες τὸν Ἰσραὴλ.
\vs{23}Καὶ τῇ Ἰεζάβελ ἐλάλησε Κύριος, λέγων, οἱ κύνες καταφάγονται αὐτὴν ἐν τῷ προτειχίσματι τοῦ Ἰεζράελ·
\vs{24}Τὸν τεθνηκότα τοῦ Ἀχαὰβ ἐν τῇ πόλει φάγονται οἱ κύνες, καὶ τὸν τεθνηκότα αὐτοῦ ἐν τῷ πεδίῳ φάγονται τὰ πετεινὰ τοῦ οὐρανοῦ.

\vs{25}Πλὴν ματαίως Ἀχαὰβ, ὃς ἐπράθη ποιῆσαι τὸ πονηρὸν ἐνώπιον Κυρίου, ὡς μετέθηκεν αὐτὸν Ἰεζάβελ ἡ γυνὴ αὐτοῦ.
\vs{26}Καὶ ἐβδελύχθη σφόδρα πορεύεσθαι ὀπίσω τῶν βδελυγμάτων, κατὰ πάντα ἃ ἐποίησεν ὁ Ἀμοῤῥαῖος, ὃν ἐξωλόθρευσε Κύριος ἀπὸ προσώπου υἱῶν Ἰσραήλ.

\vs{27}Καὶ ὑπὲρ τοῦ λόγου ὡς κατενύγη Ἀχαὰβ ἀπὸ προσώπου τοῦ Κυρίου, καὶ ἐπορεύετο κλαίων, καὶ διέῤῥηξε τὸν χιτῶνα αὐτοῦ, καὶ ἐζώσατο σάκκον ἐπὶ τὸ σῶμα αὐτοῦ, καὶ ἐνήστευσε· καὶ περιεβάλετο σάκκον ἐν τῇ ἡμέρᾳ ᾗ ἐπάταξε Ναβουθαὶ τὸν Ἰεζραηλίτην, καὶ ἐπορεύθη.
\vs{28}Καὶ ἐγένετο ῥῆμα Κυρίου ἐν χειρὶ δούλου αὐτοῦ Ἠλιοὺ περὶ Ἀχαὰβ, καὶ εἶπε Κύριος,
\vs{29}ἑώρακας ὡς κατενύγη Ἀχαὰβ ἀπὸ προσώπου μου; οὐκ ἐπάξω τὴν κακίαν ἐν ταῖς ἡμέραις αὐτοῦ, ἀλλʼ ἐν ταῖς ἡμέραις τοῦ υἱοῦ αὐτοῦ ἐπάξω τὴν κακίαν.

\ch{21}
Καὶ συνήθροισεν υἱὸς Ἄδερ πᾶσαν τὴν δύναμιν αὐτοῦ, καὶ ἀνέβη καὶ περιεκάθισεν ἐπὶ Σαμάρειαν, καὶ τριακονταδύο βασιλεῖς μετʼ αὐτοῦ, καὶ πᾶς ἵππος καὶ ἅρμα· καὶ ἀνέβησαν καὶ περιεκάθισαν ἐπὶ Σαμάρειαν, καὶ ἐπολέμησαν ἐπʼ αὐτήν.
\vs{2}Καὶ ἀπέστειλε πρὸς Ἀχαὰβ βασιλέα Ἰσραὴλ εἰς τὴν πόλιν, καὶ εἶπε πρὸς αὐτὸν, τάδε λέγει υἱὸς Ἄδερ,
\vs{3}τὸ ἀργύριόν σου καὶ τὸ χρυσίον σου ἐμόν ἐστι, καὶ αἱ γυναῖκές σου καὶ τὰ τέκνα σου ἐμά ἐστι.
\vs{4}Καὶ ἀπεκρίθη βασιλεὺς Ἰσραὴλ, καὶ εἶπε, καθὼς ἐλάλησας κύριέ μου βασιλεῦ, σὸς ἐγώ εἰμι καὶ πάντα τὰ ἐμά.

\vs{5}Καὶ ἀνέστρεψαν οἱ ἄγγελοι, καὶ εἶπαν, τάδε λέγει ὁ υἱὸς Ἄδερ, ἐγὼ ἀπέστειλα πρὸς σὲ, λέγων, τὸ ἀργύριόν σου καὶ τὸ χρυσίον σου καὶ τὰς γυναῖκας καὶ τὰ τέκνα σου δώσεις ἐμοὶ,
\vs{6}ὅτι ταύτην τὴν ὥραν αὔριον ἀποστελῶ τοὺς παῖδάς μου πρὸς σὲ, καὶ ἐρευνήσουσι τὸν οἶκόν σου καὶ τοὺς οἴκους τῶν παίδων σου, καὶ ἔσται πάντα τὰ ἐπιθυμήματα τῶν ὀφθαλμῶν αὐτῶν ἐφʼ ἃ ἂν ἐπιβάλωσι τὰς χεῖρας αὐτῶν, καὶ λήψονται.
\vs{7}Καὶ ἐκάλεσεν ὁ βασιλεὺς Ἰσραὴλ πάντας τοὺς πρεσβυτέρους τῆς γῆς, καὶ εἶπε, γνῶτε δὴ καὶ ἴδετε ὅτι κακίαν οὗτος ζητεῖ, ὅτι ἀπέσταλκε πρὸς μὲ περὶ τῶν γυναικῶν μου, καὶ περὶ τῶν υἱῶν μου, καὶ περὶ τῶν θυγατέρων μου· τὸ ἀργύριόν μου καὶ τὸ χρυσίον μου οὐκ ἀπεκώλυσα ἀπʼ αὐτοῦ.
\vs{8}Καὶ εἶπαν αὐτῷ οἱ πρεσβύτεροι καὶ πᾶς ὁ λαὸς, μὴ ἀκούσῃς, καὶ μὴ θελήσῃς.
\vs{9}Καὶ εἶπε τοῖς ἀγγέλοις υἱοῦ Ἄδερ, λέγετε τῷ κυρίῳ ὑμῶν, πάντα ὅσα ἀπέσταλκας πρὸς τὸν δοῦλόν σου ἐν πρώτοις ποιήσω, τὸ δὲ ῥῆμα τοῦτο οὐ δυνήσομαι ποιῆσαι· καὶ ἀπῇραν οἱ ἄνδρες, καὶ ἐπέστρεψαν αὐτῷ λόγον.

\vs{10}Καὶ ἀπέστειλε πρὸς αὐτὸν υἱὸς Ἄδερ, λέγων, τάδε ποιήσαι μοι ὁ Θεὸς καὶ τάδε προσθείη, εἰ ἐκποιήσει ὁ χοῦς Σαμαρείας ταῖς ἀλώπεξι παντὶ τῷ λαῷ τοῖς πεζοῖς μου.
\vs{11}Καὶ ἀπεκρίθη ὁ βασιλεὺς Ἰσραὴλ, καὶ εἶπεν, ἱκανούσθω· μὴ καυχάσθω ὁ κυρτὸς, ὡς ὁ ὀρθός.
\vs{12}Καὶ ἐγένετο ὅτε ἀπεκρίθη αὐτῷ τὸν λόγον τοῦτον, πίνων ἦν αὐτὸς καὶ πάντες οἱ βασιλεῖς οἱ μετʼ αὐτοῦ ἐν σκηναῖς· καὶ εἶπε τοῖς παισὶν αὐτοῦ, οἰκοδομήσατε χάρακα· καὶ ἔθεντο χάρακα ἐπὶ τὴν πόλιν.

\vs{13}Καὶ ἰδοὺ προφήτης εἷς προσῆλθε τῷ Ἀχαὰβ βασιλεῖ Ἰσραὴλ, καὶ εἶπε, τάδε λέγει Κύριος, εἰ ἑώρακας τὸν ὄχλον τὸν μέγαν τοῦτον; ἰδοὺ ἐγὼ δίδωμι αὐτὸν σήμερον εἰς χεῖρας σὰς, καὶ γνώσῃ ὅτι ἐγὼ Κύριος.
\vs{14}Καὶ εἶπεν Ἀχαὰβ, ἐν τίνι; καὶ εἶπε, τάδε λέγει Κύριος, ἐν τοῖς παιδαρίοις τῶν ἀρχόντων τῶν χωρῶν· καὶ εἶπεν Ἀχαὰβ, τίς συνάψει τὸν πόλεμον; καὶ εἶπε, σύ.

\vs{15}Καὶ ἐπεσκέψατο Ἀχαὰβ τοὺς ἄρχοντας τὰ παιδάρια τῶν χωρῶν, καὶ ἐγένοντο διακόσια τριάκοντα· καὶ μετὰ ταῦτα ἐπεσκέψατο τὸν λαὸν πάντα υἱὸν δυνάμεως, ἑπτὰ χιλιάδας.
\vs{16}Καὶ ἐξῆλθε μεσημβρίας, καὶ υἱὸς Ἄδερ πίνων μεθύων ἐν Σοκχὼθ αὐτὸς καὶ οἱ βασιλεῖς, τριάκοντα καὶ δύο βασιλεῖς συμβοηθοὶ αὐτοῦ.
\vs{17}Καὶ ἐξῆλθον ἄρχοντες παιδάρια τῶν χωρῶν ἐν πρώτοις· καὶ ἀποστέλλουσι καὶ ἀπαγγέλλουσι τῷ βασιλεῖ Συρίας, λέγοντες, ἄνδρες ἐξεληλύθασιν ἐκ Σαμαρείας.
\vs{18}Καὶ εἶπεν αὐτοῖς, εἰ εἰς εἰρήνην ἐκπορεύονται, συλλαβεῖν αὐτοὺς ζῶντας· καὶ εἰ εἰς πόλεμον, ζῶντας συλλαβεῖν αὐτούς·
\vs{19}καὶ μὴ ἐξελθάτωσαν ἐκ τῆς πόλεως ἄρχοντα τὰ παιδάρια τῶν χωρῶν. Καὶ ἡ δύναμις ὀπίσω αὐτῶν
\vs{20}ἐπάταξεν ἕκαστος τὸν παρʼ αὐτοῦ· καὶ ἐδευτέρωσιν ἕκαστος τὸν παρʼ αὐτοῦ· καὶ ἔφυγε Συρία· καὶ κατεδίωξεν αὐτοὺς Ἰσραήλ· καὶ σώζεται υἱὸς Ἄδερ βασιλεὺς Συρίας ἐφʼ ἵππου ἱππέως.
\vs{21}Καὶ ἐξῆλθεν βασιλεὺς Ἰσραὴλ, καὶ ἔλαβε πάντας τοὺς ἵππους καὶ τὰ ἅρματα, καὶ ἐπάταξε πληγὴν μεγάλην ἐν Συρίᾳ.
\vs{22}Καὶ προσῆλθεν ὁ προφήτης πρὸς βασιλέα Ἰσραὴλ, καὶ εἶπε, κραταιοῦ καὶ γνῶθι καὶ ἴδε τί ποιήσεις, ὅτι ἐπιστρέφοντος τοῦ ἐνιαυτοῦ υἱὸς Ἄδερ βασιλεὺς Συρίας ἀναβαίνει ἐπὶ σὲ.

\vs{23}Καὶ οἱ παῖδες βασιλέως Συρίας· καὶ εἶπον, θεὸς ὀρέων Θεὸς Ἰσραὴλ καὶ οὐ θεὸς κοιλάδων, διὰ τοῦτο ἐκραταίωσεν ὑπὲρ ἡμᾶς· ἐὰν δὲ πολεμήσωμεν αὐτοὺς κατʼ εὐθὺ, εἰ μὴν κραταιώσωμεν ὑπὲρ αὐτούς.
\vs{24}Καὶ τὸ ῥῆμα τοῦτο ποίησον· ἀπόστησον τοὺς βασιλεῖς ἕκαστον εἰς τὸν τόπον αὐτῶν, καὶ θοῦ ἀντʼ αὐτῶν σατράπας,
\vs{25}καὶ ἀλλάξομέν σοι δύναμιν κατὰ τὴν δύναμιν τὴν πεσοῦσαν, καὶ ἵππον κατὰ τὴν ἵππον, καὶ ἅρματα κατὰ τὰ ἅρματα, καὶ πολεμήσομεν πρὸς αὐτοὺς κατʼ εὐθὺ, καὶ κραταιώσομεν ὑπὲρ αὐτούς· καὶ ἤκουσε τῆς φωνῆς αὐτοῦ, καὶ ἐποίησεν οὕτως.

\vs{26}Καὶ ἐγένετο ἐπιστρέψαντος τοῦ ἐνιαυτοῦ, καὶ ἐπεσκέψατο υἱὸς Ἄδερ τὴν Συρίαν, καὶ ἀνέβη εἰς Ἀφεκὰ εἰς πόλεμον ἐπὶ Ἰσραήλ.
\vs{27}Καὶ οἱ υἱοὶ Ἰσραὴλ ἐπεσκέπησαν, καὶ παρεγένοντο εἰς ἀπαντὴν αὐτῶν· καὶ παρενέβαλεν Ἰσραὴλ ἐξεναντίας αὐτῶν ὡσεὶ δύο ποίμνια αἰγῶν· καὶ Συρία ἔπλησε τὴν γῆν.

\vs{28}Καὶ προσῆλθεν ὁ ἄνθρωπος τοῦ Θεοῦ, καὶ εἶπε τῷ βασιλεῖ Ἰσραὴλ, τάδε λέγει Κύριος, ἀνθʼ ὧν εἶπε Συρία, θεὸς ὀρέων Κύριος ὁ Θεὸς Ἰσραὴλ καὶ οὐ θεὸς κοιλάδων αὐτὸς, καὶ δώσω τὴν δύναμιν τὴν μεγάλην ταύτην εἰς χεῖρα σὴν, καὶ γνώσῃ ὅτι ἐγὼ Κύριος.
\vs{29}Καὶ παρεμβάλλουσιν οὗτοι ἀπέναντι τούτων ἑπτὰ ἡμέρας· καὶ ἐγένετο ἐν τῇ ἡμέρᾳ τῇ ἑβδόμῃ, καὶ προσήγαγεν ὁ πόλεμος, καὶ ἐπάταξεν Ἰσραὴλ τὴν Συρίαν ἑκατὸν χιλιάδας πεζῶν μιᾷ ἡμέρᾳ.
\vs{30}Καὶ ἔφυγον οἱ κατάλοιποι εἰς Ἀφεκὰ εἰς τὴν πόλιν, καὶ ἔπεσε τὸ τεῖχος ἐπὶ εἴκοσι καὶ ἑπτὰ χιλιάδας ἀνδρῶν τῶν καταλοίπων· καὶ υἱὸς Ἄδερ ἔφυγε καὶ εἰσῆλθεν εἰς τὸν οἶκον τοῦ κοιτῶνος, εἰς τὸ ταμιεῖον.

\vs{31}Καὶ εἶπε τοῖς παισὶν αὐτοῦ, οἶδα ὅτι βασιλεῖς Ἰσραὴλ βασιλεῖς ἐλέους εἰσίν· ἐπιθώμεθα δὴ σάκκους ἐπὶ τὰς ὀσφύας ἡμῶν, καὶ σχοινία ἐπὶ τὰς κεφαλὰς ἡμῶν, καὶ ἐξέλθωμεν πρὸς βασιλεα Ἰσραὴλ, εἴπως ζωογονήσει τὰς ψυχὰς ἡμῶν.
\vs{32}Καὶ περιεζώσαντο σάκκους ἐπὶ τὰς ὀσφύας αὐτῶν, καὶ ἔθεσαν σχοινία ἐπὶ τὰς κεφαλὰς αὐτῶν, καὶ εἶπον τῷ βασιλεῖ Ισραὴλ, δοῦλός σου υἱὸς Ἄδερ λέγει, ζησάτω δὴ ἡ ψυχὴ ἡμῶν· καὶ εἶπεν, εἰ ἔτι ζῇ, ἀδελφός μου ἐστί.
\vs{33}Καὶ οἱ ἄνδρες οἰωνίσαντο, καὶ ἐσπείσαντο· καὶ ἀνελέξαντο τὸν λόγον ἐκ τοῦ στόματος αὐτοῦ, καὶ εἶπον, ἀδελφός σου υἱὸς Ἄδερ· καὶ εἶπεν, εἰσέλθατε καὶ λάβετε αὐτόν· καὶ ἐξῆλθε πρὸς αὐτὸν υἱὸς Ἄδερ, καὶ ἀναβιβάζουσιν αὐτὸν πρὸς αὐτὸν ἐπὶ τὸ ἅρμα.
\vs{34}Καὶ εἶπε πρὸς αὐτὸν, τὰς πόλεις ἃς ἔλαβεν ὁ πατήρ μου παρὰ τοῦ πατρός σου ἀποδώσω σοι· καὶ ἐξόδους θήσεις σεαυτῷ ἐν Δαμασκῷ, καθὼς ἔθετο ὁ πατήρ μου ἐν Σαμαρείᾳ· καὶ ἐγὼ ἐν διαθήκῃ ἐξαποστελῶ σε. Καὶ διέθετο αὐτῷ διαθήκην, καὶ ἐξαπέστειλεν αὐτόν.

\vs{35}Καὶ ἄνθρωπος εἷς ἐκ τῶν υἱῶν τῶν προφητῶν εἶπε πρὸς τὸν πλησίον αὐτοῦ ἐν λόγῳ Κυρίου, πάταξον δή με· καὶ οὐκ ἠθέλησεν ὁ ἄνθρωπος πατάξαι αὐτόν.
\vs{36}Καὶ εἶπε πρὸς αὐτὸν, ἀνθʼ ὧν οὐκ ἤκουσας τῆς φωνῆς Κυρίου, καὶ ἰδοὺ σὺ ἀποτρέχεις ἀπʼ ἐμοῦ, καὶ πατάξει σε λέων· καὶ ἀπῆλθεν ἀπʼ αὐτοῦ, καὶ εὑρίσκει αὐτὸν λέων, καὶ ἐπάταξεν αὐτόν.
\vs{37}Καὶ εὑρίσκει ἄνθρωπον ἄλλον, καὶ εἶπε, πάταξόν με δή· καὶ ἐπάταξεν αὐτὸν ὁ ἄνθρωπος, πατάξας καὶ συνέτριψε.

\vs{38}Καὶ ἐπορεύθη ὁ προφήτης καὶ ἔστη τῷ βασιλεῖ Ἰσραὴλ ἐπὶ τῆς ὁδοῦ, καὶ κατεδήσατο ἐν τελαμῶνι τοὺς ὀφθαλμοὺς αὐτοῦ.
\vs{39}Καὶ ἐγένετο ὡς παρεπορεύετο ὁ βασιλεὺς, καὶ οὗτος ἐβόα πρὸς τὸν βασιλέα, καὶ εἶπεν, ὁ δοῦλός σου ἐξῆλθεν ἐπὶ τὴν στρατιὰν τοῦ πολέμου, καὶ ἰδοὺ ἀνὴρ εἰσήγαγε πρὸς μὲ ἄνδρα, καὶ εἶπε πρὸς μὲ, φύλαξον τοῦτον τὸν ἄνδρα· ἐὰν δὲ ἐκπηδῶν ἐκπηδήσῃ, καὶ ἔσται ἡ ψυχή σου ἀντὶ τῆς ψυχῆς αὐτοῦ, ἢ τάλαντον ἀργυρίου στήσεις.
\vs{40}Καὶ ἐγενήθη, περιεβλέψατο ὁ δοῦλός σου ὧδε καὶ ὧδε, καὶ οὗτος οὐκ ἦν· καὶ εἶπε πρὸς αὐτὸν ὁ βασιλεὺς Ἰσραὴλ, ἰδοὺ καὶ τὰ ἔνεδρα παρʼ ἐμοὶ ἐφόνευσας·
\vs{41}Καὶ ἔσπευσε καὶ ἀφεῖλε τὸν τελαμῶνα ἀπὸ τῶν ὀφθαλμῶν αὐτοῦ· καὶ ἐπέγνω αὐτὸν ὁ βασιλεὺς Ἰσραὴλ, ὅτι ἐκ τῶν προφητῶν οὗτος.
\vs{42}Καὶ εἶπε πρὸς αὐτὸν, τάδε λέγει Κύριος, διότι ἐξήνεγκας σὺ ἄνδρα ὀλέθριον ἐκ τῆς χειρός σου, καὶ ἔσται ἡ ψυχή σου ἀντὶ τῆς ψυχῆς αὐτοῦ, καὶ ὁ λαός σου ἀντὶ τοῦ λαοῦ αὐτοῦ.
\vs{43}Καὶ ἀπῆλθεν ὁ βασιλεὺς Ἰσραὴλ συγκεχυμένος καὶ ἐκλελυμένος, καὶ ἔρχεται εἰς Σαμάρειαν.

\ch{22}
Καὶ ἐκάθισε τὰ τρία ἔτη, καὶ οὐκ ἦν πόλεμος ἀναμέσον Συρίας καὶ ἀναμέσον Ἰσραήλ.
\vs{2}Καὶ ἐγενήθη ἐν τῷ ἐνιαυτῷ τῷ τρίτῳ, καὶ κατέβη Ἰωσαφὰτ βασιλεὺς Ἰούδα πρὸς βασιλέα Ἰσραήλ.
\vs{3}Καὶ εἶπε βασιλεὺς Ἰσραὴλ πρὸς τοὺς παῖδας αὐτοῦ, εἰ οἴδατε ὅτι ἡμῖν Ῥεμμὰθ Γαλαὰδ, καὶ ἡμεῖς σιωπῶμεν λαβεῖν αὐτὴν ἐκ χειρὸς βασιλέως Συρίας;
\vs{4}Καὶ εἶπε βασιλεὺς Ἰσραὴλ πρὸς Ἰωσαφὰτ, ἀναβήσῃ μεθʼ ἡμῶν εἰς Ῥεμμὰθ Γαλαὰδ εἰς πόλεμον;
\vs{5}Καὶ εἶπεν Ἰωσαφὰτ, καθὼς ἐγὼ, καὶ σὺ οὕτως· καθὼς ὁ λαός μου, ὁ λαός σου· καθὼς οἱ ἵπποι μου, οἱ ἵπποι σου.

Καὶ εἶπεν Ἰωσαφὰτ βασιλεὺς Ἰούδα πρὸς βασιλέα Ἰσραὴλ, ἐπερωτήσατε δὴ σήμερον τὸν Κύριον.
\vs{6}Καὶ συνήθροισεν ὁ βασιλεὺς Ἰσραὴλ πάντας τοὺς προφήτας ὡς τετρακοσίους ἄνδρας, καὶ εἶπεν αὐτοῖς ὁ βασιλεὺς, εἰ πορευθῶ εἰς Ῥεμμὰθ Γαλαὰδ εἰς πόλεμον ἢ ἐπισχῶ; καὶ εἶπον, ἀνάβαινε, καὶ διδοὺς δώσει Κύριος εἰς χεῖρας τοῦ βασιλέως.

\vs{7}Καὶ εἶπεν Ἰωσαφὰτ πρὸς βασιλέα Ἰσραὴλ, οὐκ ἔστιν ὧδε προφήτης τοῦ κυρίου, καὶ ἐπερωτήσομεν τὸν κύριον διʼ αὐτοῦ;
\vs{8}Καὶ εἶπεν ὁ βασιλεὺς Ἰσραὴλ πρὸς Ἰωσαφὰτ, εἷς ἐστιν ἀνὴρ εἰς τὸ ἐπερωτῆσαι δὀ αὐτοῦ τὸν Κύριον, καὶ ἐγὼ μεμίσηκα αὐτὸν, ὄτι οὐ λαλεῖ περὶ ἐμοῦ καλὰ, ἀλλʼ ἢ κακὰ, Μιχαίας υἱὸς Ἰεμβλαά· καὶ εἶπεν Ἰωσαφὰτ βασιλεὺς Ἰούδα, μὴ λεγέτω ὁ βασιλεὺς οὕτως.

\vs{9}Καὶ ἐκάλεσεν ὁ βασιλεὺς Ἰσραὴλ εὐνοῦχον ἕνα, καὶ εἶπε, τοτάχος Μιχαίαν υἱὸν Ἰεμβλαά.
\vs{10}Καὶ ὁ βασιλεὺς Ἰσραὴλ καὶ Ἰωσαφὰτ βασιλεὺς Ἰούδα ἐκάθηντο ἀνὴρ ἐπὶ τοῦ θρόνου αὐτοῦ ἔνοπλοι ἐν ταῖς πύλαις Σαμαρείας· καὶ πάντες οἱ προφῆται ἐπροφήτευον ἐνώπιον αὐτῶν.
\vs{11}Καὶ ἐποίησεν ἑαυτῷ Σεδεκίας υἱὸς Χαναὰν κέρατα σιδηρᾶ, καὶ εἶπε, τάδε λέγει Κύριος, ἐν τούτοις κερατιεῖς τὴν Συρίαν ἕως συντελεσθῇ.
\vs{12}Καὶ πάντες οἱ προφῆται ἐπροφήτευον οὕτως, λέγοντες, ἀνάβαινε εἰς Ῥεμμὰθ Γαλαὰδ, καὶ εὐοδώσει, καὶ δώσει Κύριος εἰς χεῖράς σου καὶ τὸν βασιλέα Συρίας.

\vs{13}Καὶ ὁ ἄγγελος ὁ πορευθεὶς καλέσαι τὸν Μιχαίαν, ἐλάλησεν αὐτῷ, λέγων, ἰδοὺ δὴ λαλοῦσι πάντες οἱ προφῆται ἐν στόματι ἐνὶ καλὰ περὶ τοῦ βασιλέως, γίνου δὴ καὶ σὺ εἰς τοὺς λόγους σου κατὰ τοὺς λόγους ἑνὸς τούτων, καὶ λάλησον καλά.
\vs{14}Καὶ εἶπε Μιχαίας, ζῇ Κύριος, ὅτι ἃ ἐὰν εἴπῃ Κύριος πρὸς μὲ, ταῦτα λαλήσω.

\vs{15}Καὶ ἦλθε πρὸς τὸν βασιλέα· καὶ εἶπεν αὐτῷ ὁ βασιλεὺς, Μιχαία, εἰ ἀναβῶ εἰς Ῥεμμὰθ Γαλαὰδ εἰς πόλεμον, ἢ ἐπισχῶ; καὶ εἶπεν, ἀνάβαινε, καὶ εὐοδώσει Κύριος εἰς χεῖρα τοῦ βασιλέως.
\vs{16}Καὶ εἶπεν αὐτῷ ὁ βασιλεὺς, ποσάκις ἐγὼ ὁρκίζω σε, ὃπως λαλήσῃς πρὸς μὲ ἀλήθειαν ἐν ὀνόματι Κυρίου;
\vs{17}Καὶ εἶπεν, οὐχ οὕτως· ἑώρακα πάντα τὸν Ἰσραὴλ διεσπαρμένον ἐν τοῖς ὄρεσιν ὡς ποίμνιον ᾧ οὐκ ἔστι ποιμήν· καὶ εἶπε Κύριος, οὐ κύριος τούτοις Θεός; ἕκαστος εἰς τὸν οἶκον αὐτοῦ ἐν εἰρήνῃ ἀναστρεφέτω.

\vs{18}Καὶ εἶπε βασιλεὺς Ἰσραὴλ πρὸς Ἰωσαφὰτ βασιλέα Ἰούδα, οὐκ εἶπα πρὸς σε, ὅτι οὐ προφητεύει οὗτός μοι καλὰ, διότι ἀλλʼ ἢ κακά;
\vs{19}Καὶ εἶπε Μιχαίας, οὐχ οὕτως· οὐκ ἐγώ· ἄκουε ῥῆμα Κυρίου· οὐχ οὕτως. Εἶδον Θεὸν Ἰσραὴλ καθήμενον ἐπὶ θρόνου αὐτοῦ, καὶ πᾶσα ἡ στρατιὰ τοῦ οὐρανοῦ εἱστήκει περὶ αὐτὸν ἐκ δεξιῶν αὐτοῦ καὶ ἐξ εὐωνύμων αὐτοῦ.
\vs{20}Καὶ εἶπε Κύριος, τίς ἀπατήσει τὸν Ἀχαὰβ βασιλέα Ἰσραὴλ, καὶ ἀναβήσεται, καὶ πεσεῖται ἐν Ῥεμμὰθ Γαλαάδ; καὶ εἶπεν οὗτος οὕτως, καὶ οὗτος οὕτως.
\vs{21}Καὶ ἐξῆλθε πνεῦμα καὶ ἔστη ἐνώπιον Κυρίου, καὶ εἶπεν, ἐγὼ ἀπατήσω αὐτόν.
\vs{22}Καὶ εἶπε πρὸς αὐτὸν Κύριος, ἐν τίνι; καὶ εἶπεν, ἐξελεύσομαι, καὶ ἔσομαι πνεῦμα ψευδὲς εἰς τὸ στόμα πάντων τῶν προφητῶν αὐτοῦ· καὶ εἶπεν, ἀπατήσεις, καί γε δυνήσῃ· ἔξελθε καὶ ποίησον οὕτως.
\vs{23}Καὶ νῦν ἰδοὺ ἔδωκε Κύριος πνεῦμα ψευδὲς ἐν στόματι πάντων τῶν προφητῶν σου τούτων, καὶ Κύριος ἐλάλησεν ἐπὶ σὲ κακά.

\vs{24}Καὶ προσῆλθε Σεδεκίας υἱὸς Χαναὰν, καὶ ἐπάταξε τὸν Μιχαίαν ἐπὶ τὴν σιαγόνα, καὶ εἶπε, ποῖον πνεῦμα Κύριου τὸ λαλῆσαν ἐν σοί;
\vs{25}Καὶ εἶπε Μιχαίας, ἰδοὺ σὺ ὄψῃ τῇ ἡμέρᾳ ἐκείνῃ, ὅταν εἰσέλθῃς ταμεῖον τοῦ ταμείου τοῦ κρυβῆναι ἐκεῖ.
\vs{26}Καὶ εἶπεν ὁ βασιλεὺς Ἰσραὴλ, λάβετε τὸν Μιχαίαν, καὶ ἀποστρέψατε αὐτὸν πρὸς Σεμὴρ τὸν βασιλέα τῆς πόλεως· καὶ τῷ Ἰωὰς υἱῷ τοῦ βασιλέως
\vs{27}εἶπον θέσθαι τοῦτον ἐν φυλακῇ, καὶ ἐσθίειν αὐτὸν ἄρτον θλίψεως καὶ ὕδωρ θλίψεως ἕως τοῦ ἐπιστρέψαι με ἐν εἰρήνῃ.
\vs{28}Καὶ εἶπε Μιχαίας, ἐὰν ἐπιστρέφων ἐπιστρέψῃς ἐν εἰρήνῃ, οὐ λελάληκε Κύριος ἐν ἐμοί.

\vs{29}Καὶ ἀνέβη βασιλεὺς Ἰσραὴλ καὶ Ἰωσαφὰτ βασιλεὺς Ἰούδα μετʼ αὐτοῦ εἰς Ῥεμμὰθ Γαλαάδ.
\vs{30}Καὶ εἶπε βασιλεὺς Ἰσραὴλ πρὸς Ἰωσαφὰτ βασιλέα Ἰούδα, συγκαλύψομαι καὶ εἰσελεύσομαι εἰς τὸν πόλεμον, καὶ σὺ ἔνδυσαι τὸν ἱματισμόν μου· καὶ συνεκαλύψατο βασιλεὺς Ἰσραὴλ, καὶ εἰσῆλθεν εἰς τὸν πόλεμον.
\vs{31}Καὶ βασιλεὺς Συρίας ἐνετείλατο τοῖς ἄρχουσι τῶν ἁρμάτων αὐτοῦ τριάκοντα καὶ δυσὶ, λέγων, μὴ πολεμεῖτε μικρὸν καὶ μέγαν, ἀλλʼ ἢ τὸν βασιλέα Ἰσραὴλ μονώτατον.
\vs{32}Καὶ ἐγένετο ὡς εἶδον οἱ ἄρχοντες τῶν ἁρμάτων τὸν Ἰωσαφὰτ βασιλέα Ἰούδα, καὶ αὐτοὶ εἶπαν, φαίνεται βασιλεὺς Ἰσραήλ οὗτος, καὶ ἐκύκλωσαν αὐτὸν πολεμῆσαι· καὶ ἀνέκραξεν Ἰωσαφάτ.
\vs{33}Καὶ ἐγένετο ὡς εἶδον οἱ ἄρχοντες τῶν ἁρμάτων ὅτι οὐκ ἔστι βασιλεὺς Ἰσραὴλ οὗτος, καὶ ἀπέστρεψαν ἀπʼ αὐτοῦ.

\vs{34}Καὶ ἐπέτεινεν εἷς τὸ τόξον εὐστόχως, καὶ ἐπάταξε τὸν βασιλέα Ἰσραὴλ ἀναμέσον τοῦ πνεύμονος καὶ ἀναμέσον τοῦ θώρακος· καὶ εἶπε τῷ ἡνιόχῳ αὐτοῦ, ἐπίστρεψον τὰς χεῖράς σου καὶ ἐξάγαγέ με ἐκ τοῦ πολέμου, ὅτι τέτρωμαι.
\vs{35}Καὶ ἐτροπώθη ὁ πόλεμος ἐν τῇ ἡμέρᾳ ἐκείνῃ, καὶ ὁ βασιλεὺς ἦν ἑστηκὼς ἐπὶ τοῦ ἅρματος ἐξεναντίας Συρίας ἀπὸ πρωῒ ἕως ἑσπέρας, καὶ ἀπέχυνε τὸ αἷμα ἀπὸ τῆς πληγῆς εἰς τὸν κόλπον τοῦ ἅρματος, καὶ ἀπέθανεν ἑσπέρας, καὶ ἐξεπορεύετο τὸ αἷμα τῆς τροπῆς ἕως τοῦ κόλπου τοῦ ἅρματος.
\vs{36}Καὶ ἔστη ὁ στρατοκῆρυξ δύνοντος τοῦ ἡλίου, λέγων, ἕκαστος εἰς τὴν ἑαυτοῦ πόλιν καὶ εἰς τὴν ἑαυτοῦ γῆν,
\vs{37}ὅτι τέθνηκεν ὁ βασιλεύς· καὶ ἦλθον εἰς Σαμάρειαν, καὶ ἔθαψαν τὸν βασιλέα ἐν Σαμαρείᾳ.
\vs{38}Καὶ ἀπένιψαν τὸ ἅρμα ἐπὶ τὴν κρήνην Σαμαρείας· καὶ ἐξέλιξαν αἱ ὗες καὶ οἱ κύνες τὸ αἷμα, καὶ αἱ πόρναι ἐλούσαντο ἐν τῷ αἵματι, κατὰ τὸ ῥῆμα Κυρίου ὃ ἐλάλησε.

\vs{39}Καὶ τὰ λοιπὰ τῶν λόγων Ἀχαὰβ καὶ πάντα ἃ ἐποίησε, καὶ οἴκον ἐλεφάντινον ὃν ᾠκοδόμησε, καὶ πάσας τὰς πόλεις ἃς ἐποίησεν, οὐκ ἰδοὺ ταῦτα γέγραπται ἐν βιβλίῳ λόγων τῶν ἡμερῶν τῶν βασιλέων Ἰσραήλ;
\vs{40}Καὶ ἐκοιμήθη Ἀχαὰβ μετὰ τῶν πατέρων αὐτοῦ, καὶ ἐβασίλευσεν Ὀχοζίας υἱὸς αὐτοῦ ἀντʼ αὐτοῦ.

\vs{41}Καὶ Ἰωσαφὰτ υἱὸς Ἀσὰ ἐβασίλευσεν ἐπὶ Ἰούδαν· ἐν ἔτει τετάρτῷ τοῦ Ἀχαὰβ βασιλέως Ἰσραὴλ ἐβασίλευσεν
\vs{42}Ἰωσαφὰτ, Υἱὸς τρίακοντα καὶ πέντε ἐτῶν ἐν τῷ βασιλεύειν αὐτόν, καὶ εἴκοσι καὶ πέντε ἔτη ἐβασίλευσεν ἐν Ἰερουσαλὴμ· καὶ ὄνομα τῇ μητρὶ αὐτοῦ Ἀζουβὰ θυγάτηρ Σαλαΐ.
\vs{43}Καὶ ἐπορεύθη ἐν πάση ὁδῷ Ἀσὰ τοῦ πατρὸς αὐτοῦ, οὐκ ἐξέκλινεν ἀπʼ αὐτῆς τοῦ ποιῆσαι τὸ εὐθὲς ἐν ὀφθαλμοῖς Κυρίου.
\vs{44}Πλὴν τῶν ὑψηλῶν οὐκ ἐξῇρεν· ἔτι ὁ λαὸς ἐθυσίαζε καὶ ἐθυμίων ἐν τοῖς ὑψηλοῖς.
\vs{45}Καὶ εἰρήνευσεν Ἰωσαφὰτ μετὰ βασιλέως Ἰσραήλ.

\vs{46}Καὶ τὰ λοιπὰ τῶν λόγων Ἰωσαφὰτ, καὶ αἱ δυναστεῖαι αὐτοῦ ὅσα ἐποίησεν, οὐκ ἰδοὺ ταῦτα γεγραμμένα ἐν βιβλίῳ λόγων τῶν ἡμερῶν βασιλέων Ἰούδα;
\vs{51}Καὶ ἐκοιμήθη Ἰωσαφὰτ μετὰ τῶν πατέρων αὐτοῦ, καὶ ἐτάφη παρὰ τοῖς πατράσιν αὐτοῦ ἐν πόλει Δαυὶδ τοῦ πατρὸς αὐτοῦ, καὶ ἐβασίλευσεν Ἰωρὰμ υἱὸς αὐτοῦ ἀντʼ αὐτοῦ.

\vs{52}Καὶ Ὀχοζίας υἱὸς Ἀχαὰβ ἐβασίλευσεν ἐπὶ Ἰσραὴλ ἐν Σαμαρείᾳ· ἐν ἔτει ἑπτακαιδεκάτῳ Ἰωσαφὰτ βασιλέως Ἰούδα, Ὀχοζίας υἱὸς Ἀχαὰβ ἐβασίλευσεν ἐν Ἰσραὴλ ἐν Σαμαρείᾳ δύο ἔτη.
\vs{53}Καὶ ἐποίησε τὸ πονηρὸν ἐναντίον Κυρίου, καὶ ἐπορεύθη ἐν ὁδῶ Ἀχαὰβ τοῦ πατρὸς αὐτοῦ καὶ ἐν ὁδῷ Ἰεζάβελ τῆς μητρὸς αὐτοῦ, καὶ ἐν ταῖς ἁμαρτίαις οἴκου Ἱεροβοὰμ υἱοῦ Ναβὰτ ὃς ἐξήμαρτε τὸν Ἰσραήλ·
\vs{54}Καὶ ἐδούλευσε τοῖς Βααλὶμ καὶ προσεκύνησεν αὐτοῖς, καὶ παρώργισε τὸν Κύριον Θεὸν Ἰσραὴλ, κατὰ πάντα τὰ γενόμενα ἔμπροσθεν αὐτοῦ.


\def\book{ΒΑΣΙΛΕΙΩΝ Δ}
\biblebook{ΒΑΣΙΛΕΙΩΝ Δ}


\lettrine[lines=2, loversize=0.2, nindent=0em, findent=.25em]{\textcolor{bookheadingcolor}{Κ}}{ΑΙ} ἠθέτησε Μωὰβ ἐν Ἰσραὴλ μετὰ τὸ ἀποθανεῖν Ἀχαάβ.

\postdropcapindent\vs{2}Καὶ ἔπεσεν Ὀχοζίας διὰ τοῦ δικτυωτοῦ τοῦ ἐν τῷ ὑπερῴῳ αὐτοῦ τῷ ἐν Σαμαρείᾳ, καὶ ἠῤῥώστησε· καὶ ἀπέστειλεν ἀγγέλους, καὶ εἶπε πρὸς αὐτοὺς, δεῦτε καὶ ἐπιζητήσατε ἐν τῷ Βάαλ μυΐαν θεὸν Ἀκκαρὼν, εἰ ζήσομαι ἐκ τῆς ἀῤῥωστίας μου ταύτης· καὶ ἐπορεύθησαν ἐπερωτῆσαι διʼ αὐτοῦ.
\vs{3}Καὶ ἄγγελος Κυρίου ἐκάλεσεν Ἠλιοὺ τὸν Θεσβίτην, λέγων, ἀναστὰς δεῦρο εἰς συνάντησιν τῶν ἀγγέλων Ὀχοζίου βασιλέως Σαμαρείας, καὶ λαλήσεις πρὸς αὐτοὺς, εἰ παρὰ τὸ μὴ εἶναι Θεὸν ἐν Ἰσραὴλ, ὑμεῖς πορεύεσθε ἐπιζητῆσαι ἐν τῷ Βάαλ μυΐαν θεὸν Ἀκκαρών; καὶ οὐχ οὕτως·
\vs{4}Ὅτι τάδε λέγει Κύριος, ἡ κλίνη ἐφʼ ἧς ἀνέβης ἐκεῖ, οὐ καταβήσῃ ἀπʼ αὐτῆς, ὅτι θανάτῳ ἀποθανῇ· καὶ ἐπορεύθη Ἠλιοὺ, καὶ εἶπε πρὸς αὐτούς.

\vs{5}Καὶ ἐπεστράφησαν οἱ ἄγγελοι πρὸς αὐτόν· καὶ εἶπε πρὸς αὐτοὺς, τί ὅτι ἐπεστρέψατε;
\vs{6}Καὶ εἶπαν πρὸς αὐτὸν, ἀνὴρ ἀνέβη εἰς συνάντησιν ἡμῶν, καὶ εἶπε πρὸς ἡμᾶς, δεῦτε, ἐπιστράφητε πρὸς τὸν βασιλέα τὸν ἀποστείλαντα ὑμᾶς, καὶ λαλήσατε πρὸς αὐτὸν, τάδε λέγει Κύριος, εἰ παρὰ τὸ μὴ εἶναιι Θεὸν ἐν Ἰσραὴλ, σὺ πορεύῃ ἐπιζητῆσαι ἐν τῷ Βάαλ μυΐαν θεὸν Ἀκκαρών; οὐχ οὕτως· ἡ κλίνη ἐφʼ ἧς ἀνέβης ἐκεῖ, οὐ καταβήσῃ ἀπʼ αὐτῆς, ὅτι θανάτῳ ἀποθανῇ. Καὶ ἐπιστρέψαντες ἀπήγγειλαν τῷ βασιλεῖ καθὰ ἐλάλησεν Ἠλιού·
\vs{7}καὶ ἐλάλησε πρὸς αὐτοὺς, τίς ἡ κρίσις τοῦ ἀνδρὸς τοῦ ἀναβάντος εἰς συνάντησιν ὑμῖν καὶ λαλήσαντος πρὸς ὑμᾶς τοὺς λόγους τούτους;
\vs{8}Καὶ εἶπον πρὸς αὐτὸν, ἀνὴρ δασὺς, καὶ ζώνην δερματίνην περιεζωσμένος τὴν ὀσφῦν αὐτοῦ· καὶ εἶπεν, Ἠλιοὺ ὁ Θεσβίτης οὗτός ἐστι.

\vs{9}Καὶ ἀπέστειλε πρὸς αὐτὸν πεντηκόνταρχον καὶ τοὺς πεντήκοντα αὐτοῦ, καὶ ἀνέβη πρὸς αὐτόν· καὶ ἰδοὺ Ἠλιοὺ ἐκάθητο ἐπὶ τῆς κορυφῆς τοῦ ὄρους· καὶ ἐλάλησεν ὁ πεντηκόνταρχος πρὸς αὐτὸν, καὶ εἶπεν, ἄνθρωπε τοῦ Θεοῦ, ὁ βασιλεὺς ἐκάλεσέ σε, κατάβηθι.
\vs{10}Καὶ ἀπεκρίθη Ἠλιοὺ, καὶ εἶπε πρὸς τὸν πεντηκόνταρχον, καὶ εἰ ἄνθρωπος Θεοῦ ἐγὼ, καταβήσεται πῦρ ἐκ τοῦ οὐρανοῦ, καὶ καταφάγεταί σε καὶ τοὺς πεντήκοντά σου· καὶ κατέβη πῦρ ἐκ τοῦ οὐρανοῦ, καὶ κατέφαγεν αὐτὸν καὶ τοὺς πεντήκοντα αὐτοῦ.
\vs{11}Καὶ προσέθετο ὁ βασιλεὺς, καὶ ἀπέστειλε πρὸς αὐτὸν ἄλλον πεντηκόνταρχον καὶ τοὺς πεντήκοντα αὐτοῦ· καὶ ἐλάλησεν ὁ πεντηκόνταρχος πρὸς αὐτὸν, καὶ εἶπεν, ἄνθρωπε τοῦ Θεοῦ, τάδε λέγει ὁ βασιλεὺς, ταχέως κατάβηθι.
\vs{12}Καὶ ἀπεκρίθη Ἠλιοὺ καὶ ἐλάλησε πρὸς αὐτὸν, καὶ εἶπεν, εἰ ἄνθρωπος Θεοῦ ἐγὼ, καταβήσεται πῦρ ἐκ τοῦ οὐρανοῦ, καὶ καταφάγεταί σε καὶ τοὺς πεντήκοντά σου· καὶ κατέβη πῦρ ἐκ τοῦ οὐρανοῦ, καὶ κατέφαγεν αὐτὸν καὶ τοὺς πεντήκοντα αὐτοῦ.
\vs{13}Καὶ προσέθετο ὁ βασιλεὺς ἔτι ἀποστεῖλαι ἡγούμενον καὶ τοὺς πεντήκοντα αὐτοῦ· καὶ ἦλθεν ὁ πεντηκόνταρχος ὁ τρίτος, καὶ ἔκαμψεν ἐπὶ τὰ γόνατα αὐτοῦ κατέναντι Ἠλιοὺ, καὶ ἐδεήθη αὐτοῦ, καὶ ἐλάλησε πρὸς αὐτὸν, καὶ εἶπεν, ἄνθρωπε τοῦ Θεοῦ, ἐντιμωθήτω ἡ ψυχή μου, καὶ ἡ ψυχὴ τῶν δούλων σου τούτων τῶν πεντήκοντα ἐν ὀφθαλμοῖς σου.
\vs{14}Ἰδοὺ κατέβη πῦρ ἐκ τοῦ οὐρανοῦ, καὶ κατέφαγε τοὺς δύο πεντηκοντάρχους τοὺς πρώτους· καὶ νῦν ἐντιμωθήτω δὴ ἡ ψυχή μου ἐν ὀφθαλμοῖς σου.
\vs{15}Καὶ ἐλάλησεν ἄγγελος Κυρίου πρὸς Ἠλιοὺ, καὶ εἶπε, κατάβηθι μετʼ αὐτοῦ, μὴ φοβηθῇς ἀπὸ προσώπου αὐτῶν· καὶ ἀνέστη Ἠλιοὺ καὶ κατέβη μετʼ αὐτοῦ πρὸς τὸν βασιλέα.
\vs{16}Καὶ ἐλάλησε πρὸς αὐτὸν, καὶ εἶπεν Ἠλιού, τάδε λέγει Κύριος, τί ὅτι ἀπέστειλας ἀγγέλους ἐκζητῆσαι ἐν τῷ Βάαλ μυΐαν θεὸν Ἀκκαρών; οὐχ οὕτως· ἡ κλίνη ἐφʼ ἧς ἀνέβης ἐκεῖ, οὐ καταβήσῃ ἀπʼ αὐτῆς, ὅτι θανάτῳ ἀποθανῇ.

\vs{17}Καὶ ἀπέθανε κατὰ τὸ ῥῆμα Κυρίου ὃ ἐλάλησεν Ἠλιού.
\vs{18}Καὶ τὰ λοιπὰ τῶν λόγων Ὀχοζίου ἃ ἐποίησεν, οὐκ ἰδοὺ ταῦτα γεγραμμένα ἐν βιβλίῳ λόγων τῶν ἡμερῶν τοῖς βασιλεῦσιν Ἰσραήλ;
\vs{18a}καὶ Ἰωρὰμ υἱὸς Ἀχαὰβ βασιλεύει ἐπὶ Ἰσραὴλ ἐν Σαμαρείᾳ ἔτη δεκαδύο, ἐν ἔτει ὀκτωκαιδεκάτῳ Ἰωσαφὰτ βασιλέως Ἰούδα·
\vs{18b}καὶ ἐποίησε τὸ πονηρὸν ἐνώπιον Κυρίου· πλὴν οὐχ ὡς οἱ ἀδελφοὶ αὐτοῦ, οὐδὲ ὡς ἡ μήτηρ αὐτοῦ·
\vs{18c}καὶ ἀπέστησε τὰς στήλας τοῦ Βάαλ ἃς ἐποίησεν ὁ πατὴρ αὐτοῦ, καὶ συνέτριψεν αὐτάς· πλὴν ἐν ταῖς ἁμαρτίαις οἴκου Ἰεροβοὰμ, ὃς ἐξήμαρτε τὸν Ἰσραὴλ, ἐκολλήθη, οὐκ ἀπέστη ἀπʼ αὐτῶν·
\vs{18d}καὶ ἐθυμώθη ὀργῇ Κύριος εἰς τὸν οἶκον Ἀχαάβ.

\ch{2}
Καὶ ἐγένετο ἐν τῷ ἀνάγειν Κύριον ἐν συσσεισμῷ τὸν Ἠλιοὺ ὡς εἰς τὸν οὐρανὸν, καὶ ἐπορεύθη Ἠλιοὺ καὶ Ἑλισαιὲ ἐκ Γαλγάλων.
\vs{2}Καὶ εἶπεν Ἠλιοὺ πρὸς Ἑλισαιὲ, κάθου δὴ ἐνταῦθα, ὅτι ὁ Θεὸς ἀπέσταλκέ με ἕως Βαιθήλ· καὶ εἶπεν Ἑλισαιὲ, ζῇ Κύριος καὶ ζῇ ἡ ψυχή σου, εἰ ἐγκαταλείψω σε· καὶ ἦλθον εἰς Βαιθήλ.
\vs{3}Καὶ ἦλθον οἱ υἱοὶ τῶν προφητῶν οἱ ἐν Βαιθὴλ πρὸς Ἑλισαιὲ, καὶ εἶπον πρὸς αὐτὸν, εἰ ἔγνως, ὅτι Κύριος σήμερον λαμβάνει τὸν κύριόν σου ἐπάνωθεν τῆς κεφαλῆς σου; καὶ εἶπε, κἀγὼ ἔγνωκα, σιωπᾶτε.
\vs{4}Καὶ εἶπεν Ἠλιοὺ πρὸς Ἑλισαιὲ, κάθου δὴ ἐνταῦθα, ὅτι Κύριος ἀπέσταλκέ με εἰς Ἱεριχώ· καὶ εἶπε, ζῇ Κύριος καὶ ζῇ ἡ ψυχή σου, εἰ ἐγκαταλείψω σε· καὶ ἦλθον εἰς Ἱεριχώ.

\vs{5}Καὶ ἤγγισαν οἱ υἱοῖ τῶν προφητῶν οἱ ἐν Ἱεριχὼ πρὸς Ἑλισαιὲ, καὶ εἶπον πρὸς αὐτὸν, εἰ ἔγνως, ὅτι σήμερον λαμβάνει Κύριος τὸν κύρίον σου ἐπάνωθεν τῆς κεφαλῆς σου; καὶ εἶπε, καί γε ἐγὼ ἔγνων, σιωπᾶτε.
\vs{6}Καὶ εἶπεν αὐτῷ Ἠλιοὺ, κάθου δὴ ὧδε, ὅτι Κύριος ἀπέσταλκέ με ἕως εἰς τὸν Ἰορδάνην· καὶ εἶπεν Ἑλισαιὲ, ζῇ Κύριος καὶ ζῇ ἡ ψυχή σου, εἰ ἐγκαταλείψω σε· καὶ ἐπορεύθησαν ἀμφότεροι,
\vs{7}καὶ πεντήκοντα ἄνδρες υἱοὶ τῶν προφητῶν, καὶ ἔστησαν ἐξεναντίας μακρόθεν· καὶ ἀμφότεροι ἔστησαν ἐπὶ τοῦ Ἰορδάνου.
\vs{8}Καὶ ἔλαβεν Ἠλιοὺ τὴν μηλωτὴν αὐτοῦ καὶ εἴλησε καὶ ἐπάταξε τὸ ὕδωρ, καὶ διῃρέθη τὸ ὕδωρ ἔνθα καὶ ἔνθα· καὶ διέβησαν ἀμφότεροι ἐν ἐρήμῳ.

\vs{9}Καὶ ἐγένετο ἐν τῷ διαβῆναι αὐτοὺς, καὶ Ἠλιοὺ εἶπε πρὸς Ἑλισαιὲ, αἴτησαι τί ποιήσω σοι πρὶν ἢ ἀναληφθῆναί με ἀπὸ σοῦ· καὶ εἶπεν Ἑλισαιὲ, γενηθήτω δὴ διπλᾶ ἐν πνεύματί σου ἐπʼ ἐμέ.
\vs{10}Καὶ εἶπεν Ἠλιοὺ, ἐσκλήρυνας τοῦ αἰτήσασθαι· ἐὰν ἴδῃς με ἀναλαμβανόμενον ἀπὸ σοῦ, καὶ ἔσται σοι οὕτως· καὶ ἐὰν μὴ, οὐ μὴ γένηται.

\vs{11}Καὶ ἐγένετο αὐτῶν πορευομένων, ἐπορεύοντο καὶ ἐλάλουν· καὶ ἰδοὺ ἅρμα πυρὸς καὶ ἵπποι πυρὸς, καὶ διέστειλεν ἀναμέσον ἀμφοτέρων· καὶ ἀνελήμφθη Ἠλιοὺ ἐν συσσεισμῷ ὡς εἰς τὸν οὐρανόν.
\vs{12}Καὶ Ἑλισαιὲ ἑώρα, καὶ ἐβόα, πάτερ, πάτερ, ἅρμα Ἰσραὴλ καὶ ἱππεὺς αὐτοῦ· καὶ οὐκ εἶδεν αὐτὸν ἔτι· καὶ ἐπελάβετο τῶν ἱματίων αὐτοῦ, καὶ διέῤῥηξεν αὐτὰ εἰς δύο ῥήγματα.
\vs{13}Καὶ ὕψωσε τὴν μηλωτὴν Ἠλιοὺ, ἣ ἔπεσεν ἐπάνωθεν Ἑλισαιέ· καὶ ἐπέστρεψεν Ἑλισαιὲ, καὶ ἔστη ἐπὶ τοῦ χείλους τοῦ Ἰορδάνου,
\vs{14}καὶ ἔλαβε τὴν μηλωτὴν Ἠλιοὺ, ἣ ἔπεσεν ἐπάνωθεν αὐτοῦ, καὶ ἐπάταξε τὸ ὕδωρ, καὶ εἶπε, ποῦ ὁ Θεὸς Ἠλιοὺ ἀφφώ; καὶ ἐπάταξε τὰ ὕδατα, καὶ διεῤῥάγησαν ἔνθα καὶ ἔνθα· καὶ διέβη Ἑλισαιέ.

\vs{15}Καὶ εἶδον αὐτὸν οἱ υἱοὶ τῶν προφητῶν οἱ ἐν Ἱεριχὼ ἐξεναντίας, καὶ εἶπον, ἐπαναπέπαυται τὸ πνεῦμα Ἠλιοὺ ἐπὶ Ἑλισαιέ· καὶ ἦλθον εἰς συναντὴν αὐτοῦ, καὶ προσεκύνησαν αὐτῷ ἐπὶ τὴν γῆν,
\vs{16}καὶ εἶπον πρὸς αὐτὸν, ἰδοὺ δὴ μετὰ τῶν παίδων σου πεντήκοντα ἄνδρες υἱοὶ δυνάμεως· πορευθέντες δὴ ζητησάτωσαν τὸν κύριόν σου, μή ποτε ᾖρεν αὐτὸν πνεῦμα Κυρίου, καὶ ἔῤῥιψεν αὐτὸν ἐν τῷ Ἰορδάνῃ ἢ ἐφʼ ἓν τῶν ὀρέων ἢ ἐφʼ ἕνα τῶν βουνῶν· καὶ εἶπεν Ἑλισαιὲ, οὐκ ἀποστελεῖτε.
\vs{17}Καὶ παρεβιάσαντο αὐτὸν, ἕως οὗ ᾐσχύνετο· καὶ εἶπεν, ἀποστείλατε· καὶ ἀπέστειλαν πεντήκοντα ἄνδρας, καὶ ἐζήτησαν τρεῖς ἡμέρας, καὶ οὐχ εὗρον αὐτόν.
\vs{18}Καὶ ἀνέστρεψαν πρὸς αὐτόν, καὶ αὐτὸς ἐκάθητο ἐν Ἱεριχώ· καὶ εἶπεν Ἑλισαῖε, οὐκ εἶπον πρὸς ὑμᾶς, μὴ πορευθῆτε;

\vs{19}Καὶ εἶπον οἱ ἄνδρες τῆς πόλεως πρὸς Ἑλισαιὲ, ἰδοὺ ἡ κατοίκησις τῆς πόλεως ἀγαθὴ, καθὼς ὁ κύριος βλέπει, καὶ τὰ ὕδατα πονηρὰ, καὶ ἡ γῆ ἀτεκνουμένη.
\vs{20}Καὶ εἶπεν Ἑλισαιὲ λάβετέ μοι ὑδρίσκην καινὴν, καὶ θέτε ἐκεῖ ἅλα· καὶ ἔλαβον, καὶ ἤνεγκαν πρὸς αὐτόν.
\vs{21}Καὶ ἐξῆλθεν Ἐλισαιὲ εἰς τὴν διέξοδον τῶν ὑδάτων, καὶ ἔῤῥιψεν ἐκεῖ ἅλα, καὶ εἶπε, τάδε λέγει Κύριος, ἴαμαι τὰ ὕδατα ταῦτα, οὐκ ἔσται ἔτι ἐκεῖθεν θάνατος καὶ ἀτεκνουμένη.
\vs{22}Καὶ ἰάθησαν τὰ ὕδατα ἕως τῆς ἡμέρας ταύτης, κατὰ τὸ ῥῆμα Ἑλισαιὲ ὃ ἐλάλησε.

\vs{23}Καὶ ἀνέβη ἐκεῖθεν εἰς Βαιθήλ· καὶ ἀναβαίνοντος αὐτοῦ ἐν τῇ ὁδῷ καὶ παιδάρια μικρὰ ἐξῆλθον ἐκ τῆς πόλεως καὶ κατέπαιζον αὐτοῦ, καὶ εἶπον αὐτῷ, ἀνάβαινε φαλακρὲ, ἀνάβαινε.
\vs{24}Καὶ ἐξένευσεν ὀπίσω αὐτῶν, καὶ εἶδεν αὐτὰ, καὶ κατηράσατο αὐτοῖς ἐν ὀνόματι Κυρίου· καὶ ἰδοὺ ἐξῆλθον δύο ἄρκοι ἐκ τοῦ δρυμοῦ, καὶ ἀνέῤῥηξαν ἀπʼ αὐτῶν τεσσαράκοντα καὶ δύο παῖδας.
\vs{25}Καὶ ἐπορεύθη ἐκεῖθεν εἰς τὸ ὄρος τὸ Καρμήλιον, κᾀκεῖθεν ἐπέστρεψεν εἰς Σαμάρειαν.

\ch{3}
Καὶ Ἰωρὰμ υἱὸς Ἀχαὰμ ἐβασίλευσεν ἐν Ἰσραὴλ ἐν ἔτει ὀκτωκαιδεκάτῳ Ἰωσαφὰτ βασιλέως Ἰούδα, καὶ ἐβασίλευσε δώδεκα ἔτη.
\vs{2}Καὶ ἐποίησε τὸ πονηρὸν ἐν ὀφθαλμοῖς Κυρίου· πλὴν οὐχ ὡς ὁ πατὴρ αὐτοῦ, καὶ οὐχ ὡς ἡ μήτηρ αὐτοῦ· καὶ μετέστησε τὰς στήλας τοῦ Βάαλ, ἃς ἐποίησεν ὁ πατὴρ αὐτοῦ.
\vs{3}Πλὴν ἐν τῇ ἁμαρτίᾳ Ἱεροβοὰμ υἱοῦ Ναβὰτ ὃς ἐξήμαρτε τὸν Ἰσραὴλ, ἐκολλήθη, οὐκ ἀπέστη ἀπʼ αὐτῆς.

\vs{4}Καὶ Μωσὰ βασιλεὺς Μωὰβ ἦν νωκὴδ, καὶ ἐπέστρεφε τῷ βασιλεῖ Ἰσραὴλ ἐν τῇ ἐπαναστάσει ἑκατὸν χιλιάδας ἀρνῶν, καὶ ἑκατὸν χιλιάδας κριῶν ἐπὶ πόκων.
\vs{5}Καὶ ἐγένετο μετὰ τὸ ἀποθανεῖν Ἀχαὰβ, καὶ ἠθέτησε βασιλεὺς Μωὰβ ἐν βασιλεῖ Ἰσραήλ.

\vs{6}Καὶ ἐξῆλθεν ὁ βασιλεὺς Ἰωρὰμ ἐν τῇ ἡμέρᾳ ἐκείνῃ ἐκ Σαμαρείας, καὶ ἐπεσκέψατο τὸν Ἰσραήλ.
\vs{7}Καὶ ἐπορεύθη καὶ ἐξαπέστειλε πρὸς Ἰωσαφὰτ βασιλέα Ἰούδα, λέγων, βασιλεὺς Μωὰβ ἠθέτησεν ἐν ἐμοί· εἰ πορεύσῃ μετʼ ἐμοῦ εἰς Μωὰβ εἰς πόλεμον; καὶ εἶπεν, ἀναβήσομαι· ὅμοιός μοι, ὅμοιός σοι· ὡς ὁ λαός μου, ὁ λαός σου· ὡς οἱ ἵπποι μου, οἱ ἵπποι σου.
\vs{8}καὶ εἶπε, ποίᾳ ὁδῷ ἀναβῶ; καὶ εἶπεν, ὁδὸν ἔρημον Ἐδώμ.
\vs{9}Καὶ ἐπορεύθη ὁ βασιλεὺς Ἰσραὴλ καὶ ὁ βασιλεὺς Ἰούδα καὶ ὁ βασιλεὺς Ἐδὼμ, καὶ ἐκύκλωσαν ὁδὸν ἑπτὰ ἡμερῶν· καὶ οὐκ ἦν ὕδωρ τῇ παρεμβολῇ καὶ τοῖς κτήνεσι τοῖς ἐν τοῖς ποσὶν αὐτῶν.

\vs{10}Καὶ εἶπεν ὁ βασιλεὺς Ἰσραὴλ, ὢ, ὅτι κέκληκε Κύριος τοὺς τρεῖς βασιλεῖς παρερχομένους δοῦναι αὐτοὺς ἐν χειρὶ Μωάβ.
\vs{11}Καὶ εἶπεν Ἰωσαφὰτ, οὐκ ἔστιν ὧδε προφήτης τοῦ Κυρίου, καὶ ἐπιζητήσωμεν τὸν Κύριον παρʼ αὐτοῦ; καὶ ἀπεκρίθη εἷς τῶν παίδων τοῦ βασιλέως Ἰσραὴλ, καὶ εἶπεν, ὧδε Ἑλισαιὲ υἱὸς Σαφὰτ, ὃς ἐπέχεεν ὕδωρ ἐπὶ χεῖρας Ἠλιού.
\vs{12}Καὶ εἶπεν Ἰωσαφὰτ, ἔστιν αὐτῷ ῥῆμα Κυρίου· καὶ κατέβη πρὸς αὐτὸν βασιλεὺς Ἰσραὴλ, καὶ Ἰωσαφὰτ βασιλεὺς Ἰούδα, καὶ βασιλεὺς Ἐδώμ.

\vs{13}Καὶ εἶπεν Ἑλισαιὲ πρὸς βασιλέα Ἰσραὴλ, τί ἐμοὶ καὶ σοί; δεῦρο πρὸς τοὺς προφήτας τοῦ πατρός σου καὶ τοὺς προφήτας τῆς μητπός σου· καὶ εἶπεν αὐτῷ ὁ βασιλεὺς Ἰσραήλ, μή ὅτι κέκληκε Κύριος τοὺς τρεῖς βασιλεῖς τοῦ παραδοῦναι αὐτοὺς εἰς χεῖρας Μωάβ;
\vs{14}Καὶ εἶπεν Ἑλισαιὲ, ζῇ Κύριος τῶν δυνάμεων ᾧ παρέστην ἐνώπιον αὐτοῦ, ὅτι εἰ μὴ πρόσωπον Ἰωσαφὰτ βασιλέως Ἰούδα ἐγὼ λαμβάνω, εἰ ἐπέβλεψα πρὸς σὲ, καὶ εἶδόν σε.
\vs{15}Καὶ νῦν λάβε μοι ψάλλοντα· καὶ ἐγένετο ὡς ἔψαλλεν ὁ ψάλλων, καὶ ἐγένετο ἐπʼ αὐτὸν χεὶρ Κυρίου,
\vs{16}καὶ εἶπε, τάδε λέγει Κύριος, ποιήσατε τὸν χειμάῤῥουν τοῦτον βοθύνους βοθύνους,
\vs{17}ὅτι τάδε λέγει Κύριος, οὐκ ὄψεσθε πνεῦμα, καὶ οὐκ ὄψεσθε ὑετὸν, καὶ ὁ χειμάῤῥους οὗτος πλησθήσεται ὕδατος, καὶ πίεσθε ὑμεῖς καὶ αἱ κτήσεις ὑμῶν καὶ τὰ κτήνη ὑμῶν.
\vs{18}Καὶ κούφη αὑτὴ ἐν ὀφθαλμοῖς Κυρίου· καὶ παραδώσω τὴν Μωὰβ ἐν χειρὶ ὑμῶν.
\vs{19}Καὶ πατάξετε πᾶσαν πόλιν ὀχυρὰν, καὶ πᾶν ξύλον ἀγαθὸν καταβαλεῖτε, καὶ πάσας πηγὰς ὕδατος ἐμφράξεσθε, καὶ πᾶσαν μερίδα ἀγαθὴν ἀχρειώσετε ἐν λίθοις.

\vs{20}Καὶ ἐγένετο πρωῒ ἀναβαινούσης τῆς θυσίας, καὶ ἰδοὺ ὕδατα ἤρχοντο ἐξ ὁδοῦ Ἐδώμ, καὶ ἐπλήσθη ἡ γῆ ὕδατος.

\vs{21}Καὶ πᾶσα Μωὰβ ἤκουσαν ὅτι ἀνέβησαν οἱ τρεῖς βασιλεῖς πολεμεῖν αὐτούς· καὶ ἀνεβόησαν ἐκ παντὸς περιεζωσμένοι ζώνην· καὶ εἶπον, ὤ· καὶ ἔστησαν ἐπὶ τοῦ ὁρίου.
\vs{22}Καὶ ὤρθρισαν τοπρωῒ, καὶ ὁ ἥλιος ἀνέτειλεν ἐπὶ τὰ ὕδατα· καὶ εἶδε Μωὰβ ἐξεναντίας τὰ ὕδατα πυῤῥὰ ὡς αἷμα,
\vs{23}καὶ εἶπαν, αἷμα τοῦτο τῆς ῥομφαίας· καὶ ἐμαχέσαντο οἱ βασιλεῖς, καὶ ἐπάταξεν ἀνὴρ τὸν πλησίον αὐτοῦ· καὶ νῦν ἐπὶ τὰ σκῦλα Μωάβ.
\vs{24}Καὶ εἰσῆλθον εἰς τὴν παρεμβολὴν Ἰσραήλ· καὶ Ἰσραὴλ ἀνέστησαν καὶ ἐπάταξαν τὴν Μωὰβ, καὶ ἔφυγον ἀπὸ προσώπου αὐτῶν· καὶ εἰσῆλθον εἰσπορευόμενοι καὶ τύπτοντες τὴν Μωάβ,
\vs{25}καὶ τὰς πόλεις καθεῖλον, καὶ πᾶσαν μερίδα ἀγαθὴν ἔῤῥιψαν ἀνὴρ τὸν λίθον καὶ ἐνέπλησαν αὐτὴν, καὶ πᾶσαν πηγὴν ἐνέφραξαν, καὶ πᾶν ξύλον ἀγαθὸν κατέβαλον ἕως τοῦ καταλιπεῖν τοὺς λίθους τοῦ τοίχου καθῃρημένους· καὶ ἐκύκλευσαν οἱ σφενδονῆται, καὶ ἐπάταξαν αὐτήν.
\vs{26}Καὶ εἶδεν ὁ βασιλεὺς Μωὰβ ὅτι ἐκραταίωσεν ὑπὲρ αὐτὸν ὁ πόλεμος· καὶ ἔλαβε μεθʼ ἑαυτοῦ ἑπτακοσίους ἄνδρας ἐσπασμένους ῥομφαίαν διακόψαι πρὸς βασιλέα Ἐδὼμ, καὶ οὐκ ἠδυνήθησαν.
\vs{27}Καὶ ἔλαβε τὸν υἱὸν αὐτοῦ τὸν πρωτότοκον ὃν ἐβασίλευσεν ἀντʼ αὐτοῦ, καὶ ἀνήνεγκεν αὐτὸν ὁλοκαύτωμα ἐπὶ τοῦ τείχους, καὶ ἐγένετο μετάμελος μέγας ἐπὶ Ἰσραήλ· καὶ ἀπῆραν ἀπʼ αὐτοῦ, καὶ ἐπέστρεψαν εἰς τὴν γῆν.

\ch{4}
Καὶ γυνὴ μία ἀπὸ τῶν υἱῶν τῶν προφητῶν ἐβόα πρὸς τὸν Ἑλισαῖε, λέγουσα, ὁ δοῦλός σου ἀνήρ μου ἀπέθανε, καὶ σὺ ἔγνως, ὅτι δοῦλὸς σου ἦν φοβούμενος τὸν Κύριον· καὶ ὁ δανειστὴς ἦλθε λαβεῖν τοὺς δύο υἱούς μου ἑαυτῷ εἰς δούλους.
\vs{2}Καὶ εἶπεν Ἑλισαιὲ, τί ποιήσω σοι; ἀνάγγειλόν μοι τί ἐστι σοι ἐν τῷ οἴκῳ; ἡ δὲ εἶπεν, οὐκ ἔστι τῇ δούλῃ σου οὐδὲν ἐν τῷ οἴκῳ, ὅτι ἀλλʼ ἢ ὃ ἀλείψομαι ἔλαιον.
\vs{3}Καὶ εἶπε πρὸς αὐτήν, δεῦρο, αἴτησαι σεαυτῇ σκεύη ἔξωθεν παρὰ πάντων τῶν γειτόνων σκεύη κενὰ, μὴ ὀλιγώσῃς.
\vs{4}Καὶ εἰσελεύσῃ καὶ ἀποκλείσεις τὴν θύραν κατὰ σοῦ καὶ κατὰ τῶν υἱῶν σου, καὶ ἀποχεεῖς εἰς τὰ σκεύη ταῦτα, καὶ τὸ πληρωθὲν ἀρεῖς.
\vs{5}Καὶ ἀπῆλθε παρʼ αὐτοῦ, καὶ ἀπέκλεισε τὴν θύραν καθʼ ἑαυτῆς καὶ κατὰ τῶν υἱῶν αὐτῆς· αὐτοὶ προσήγγιζον πρὸς αὐτὴν, καὶ αὐτὴ ἐπέχεεν ἕως ἐπλήσθησαν τὰ σκεύη.
\vs{6}Καὶ εἶπε πρὸς τοὺς υἱοὺς αὐτῆς, ἐγγίσατε ἔτι πρὸς μὲ τὸ σκεῦος· καὶ εἶπον αὐτῇ, οὐκ ἔστιν ἔτι σκεῦος· καὶ ἔστη τὸ ἔλαιον.
\vs{7}Καὶ ἦλθε, καὶ ἀπήγγειλε τῷ ἀνθρώπῳ τοῦ Θεοῦ· καὶ εἶπεν Ἑλισαιὲ, δεῦρο καὶ ἀπόδου τὸ ἔλαιον, καὶ ἀποτίσεις τοὺς τόκους σου, καὶ σὺ καὶ οἱ υἱοί σου ζήσεσθε ἐν τῷ ἐπιλοίπῳ ἐλαίῳ.

\vs{8}Καὶ ἐγένετο ἡμέρα, καὶ διέβη Ἐλισαιὲ εἰς Σωμὰν, καὶ ἐκεῖ γυνὴ μεγάλη, καὶ ἐκράτησεν αὐτὸν φαγεῖν ἄρτον· καὶ ἐγένετο ἀφʼ ἱκανοῦ τοῦ εἰσπορεύεσθαι αὐτὸν, ἐξέκλινε τοῦ ἐκεῖ φαγεῖν.
\vs{9}Καὶ εἶπεν ἡ γυνὴ πρὸς τὸν ἄνδρα αὐτῆς, ἰδοὺ δὴ ἔγνων ὅτι ἄνθρωπος τοῦ Θεοῦ ἅγιος οὗτος διαπορεύεται ἐφ ἡμᾶς διαπαντός.
\vs{10}Ποιήσωμεν δὴ αὐτῷ ὑπερῷον τόπον μικρόν, καὶ θῶμεν αὐτῷ ἐκεῖ κλίνην, καὶ τράπεζαν, καὶ δίφρον, καὶ λυχνίαν· καὶ ἔσται ἐν τῷ εἰσπορεύεσθαι πρὸς ἡμᾶς, καὶ ἐκκλινεῖ ἐκεῖ.

\vs{11}Καὶ ἐγένετο ἡμέρα, καὶ εἰσῆλθεν ἐκεῖ, καὶ ἐξέκλινεν εἰς τὸ ὑπερῷον, καὶ ἐκοιμήθη ἐκεῖ.
\vs{12}Καὶ εἶπε πρὸς Γιεζὶ τὸ παιδάριον αὐτοῦ, κάλεσόν μοι τὴν Σωμανίτιν ταύτην· καὶ ἐκάλεσεν αὐτήν, καὶ ἔστη ἐνώπιον αὐτοῦ.
\vs{13}Καὶ εἶπεν αὐτῷ, εἶπον δὴ πρὸς αὐτὴν, ἰδοὺ, ἐξέστησας ἡμῖν πᾶσαν τὴν ἔκστασιν ταύτην, τί δεῖ ποιῆσαί σοι; εἰ ἔστι λόγος σοι πρὸς τὸν βασιλέα, ἢ πρὸς τὸν ἄρχοντα τῆς δυνάμεως; ἡ δὲ εἶπεν, ἐν μέσῳ τοῦ λαοῦ ἐγώ εἰμι οἰκῶ.
\vs{14}Καὶ εἶπε πρὸς Γιεζὶ, τί δεῖ ποιῆσαι αὐτῇ; καὶ εἶπε Γιεζὶ τὸ παιδάριον αὐτοῦ, καὶ μάλα υἱὸς οὐκ ἔστιν αὐτῇ· καὶ ὁ ἀνὴρ αὐτῆς πρεσβύτης.

\vs{15}Καὶ ἐκάλεσεν αὐτήν, καὶ ἔστη παρὰ τὴν θύραν.
\vs{16}Καὶ εἶπεν Ἐλισαιὲ πρὸς αὐτήν, εἰς τὸν καιρὸν τοῦτον, ὡς ἡ ὥρα, ζῶσα σὺ, περιειληφυῖα υἱόν· ἡ δὲ εἶπε, μή κύριε, μὴ διαψεύσῃ τὴν δούλην σου.
\vs{17}Καὶ ἐν γαστρὶ ἔλαβεν ἡ γυνή, καὶ ἔτεκεν υἱὸν εἰς τὸν καιρὸν τοῦτον, ὡς ἡ ὥρα, ζῶσα, ὡς ἐλάλησε πρὸς αὐτὴν Ἑλισαιὲ.

\vs{18}Καὶ ἡδρύνθη τὸ παιδάριον· καὶ ἐγένετο ἡνίκα ἐξῆλθε πρὸς τὸν πατέρα αὐτοῦ πρὸς τοὺς θερίζοντας,
\vs{19}καὶ εἶπε πρὸς τὸν πατέρα αὐτοῦ, τὴν κεφαλήν μου, τὴν κεφαλήν μου· καὶ εἶπε τῷ παιδαρίῳ, ἆρον αὐτὸν πρὸς τὴν μητέρα αὐτοῦ.
\vs{20}Καὶ ᾖρεν αὐτὸν πρὸς τὴν μητέρα αὐτοῦ, καὶ ἐκοιμήθη ἐπὶ τῶν γονάτων αὐτῆς ἕως μεσημβρίας, καὶ ἀπέθανε.
\vs{21}Καὶ ἀνήνεγκεν αὐτὸν, καὶ ἐκοίμισεν αὐτὸν ἐπὶ τὴν κλίνην τοῦ ἀνθρώπου τοῦ Θεοῦ· καὶ ἀπέκλεισε κατʼ αὐτοῦ, καὶ ἐξῆλθε,
\vs{22}καὶ ἐκάλεσε τὸν ἄνδρα αὐτῆς, και εἶπεν, ἀπόστειλον δή μοι ἓν τῶν παιδαρίων καὶ μίαν τῶν ὄνων, καὶ δραμοῦμαι ἕως τοῦ ἀνθρώπου τοῦ Θεοῦ, καὶ ἐπιστρέψω.
\vs{23}Καὶ εἶπε, τι ὅτι σὺ πορεύῃ πρὸς αὐτὸν σήμερον; οὐ νεομηνία, οὐδὲ σάββατον· ἡ δὲ εἶπεν, εἰρήνη.

\vs{24}Καὶ ἐπέσαξε τὴν ὄνον, καὶ εἶπε πρὸς τὸ παιδάριον αὐτῆς, ἄγε, πορεύου, μὴ ἐπίσχῃς μοι τοῦ ἐπιβῆναι ὅτι ἐὰν εἴπω σοι· δεῦρο καὶ πορεύσῃ, καὶ ἐλεύσῃ πρὸς τὸν ἄνθρωπον τοῦ θεοῦ εἰς ὄρος τὸ Καρμήλιον.
\vs{25}Καὶ ἐπορεύθη καὶ ἦλθεν ἕως τοῦ ἀνθρώπου τοῦ Θεοῦ εἰς τὸ ὄρος· καὶ ἐγένετο ὡς εἶδεν Ἑλισαιὲ ἐρχομένην αὐτὴν, καὶ εἶπε πρὸς Γιεζὶ τὸ παιδάριον αὐτοῦ, ἰδοὺ δὴ ἡ Σωμανίτις ἐκείνη.
\vs{26}Νῦν δράμε εἰς ἀπαντὴν αὐτῆς, καὶ ἐρεῖς, ἡ εἰρήνη σοι; ἡ εἰρήνη τῷ ἀνδρί σου; ἡ εἰρήνη τῷ παιδαρίῳ; ἡ δὲ εἶπεν, εἰρήνη.
\vs{27}Καὶ ἦλθε πρὸς Ἑλισαιὲ εἰς τὸ ὄρος, καὶ ἐπελάβετο τῶν ποδῶν αὐτοῦ· καὶ ἤγγιζε Γιεζὶ ἀπώσασθαι αὐτήν. καὶ εἶπεν Ἑλισαιὲ, ἄφες αὐτὴν, ὅτι ἡ ψυχὴ αὐτῆς κατώδυνος αὐτῇ, καὶ Κύριος ἀπέκρυψεν ἀπʼ ἐμοῦ καὶ σοῦ καὶ οὐκ ἀνήγγειλέ μοι.
\vs{28}Ἡ δὲ εἶπε, μὴ ᾐτησάμην υἱὸν παρὰ τοῦ κυρίου μου; ὅτι οὐκ εἶπα, οὐ πλανήσεις μετʼ ἐμοῦ;

\vs{29}Καὶ εἶπεν Ἐλισαιὲ τῷ Γιεζί, ζῶσαι τὴν ὀσθύν σου, καὶ λάβε τὴν βακτηρίαν μου ἐν τῇ χειρί σου, καὶ δεῦρο, ὅτι ἐὰν εὕρῃς ἄνδρα οὐκ εὐλογήσεις αὐτὸν, καὶ ἐὰν εὐλογήσῃ σε ἀνὴρ οὐκ ἀποκριθήσῃ αὐτῷ· καὶ ἐπιθήσεις τὴν βακτηρίαν μου ἐπὶ πρόσωπον τοῦ παιδαρίου.
\vs{30}Καὶ εἶπεν ἡ μήτηρ τοῦ παιδαρίου, ζῇ Κύριος καὶ ζῇ ἡ ψυχή σου, εἰ ἐνκαταλείψω σε· καὶ ἀνέστη Ἑλισαιὲ, καὶ ἐπορεύθη ὀπίσω αὐτῆς.
\vs{31}Καὶ Γιεζὶ διῆλθεν ἔμπροσθεν αὐτῆς, καὶ ἀπέθηκε τὴν βακτηρίαν ἐπὶ πρόσωπον τοῦ παιδαρίου· καὶ οὐκ ἦν φωνὴ καὶ οὖκ ἦν ἀκρόασις· καὶ ἐπέστρεψεν εἰς ἀπαντὴν αὐτοῦ, καὶ ἀπήγγειλεν αὐτῷ, λέγων, οὐκ ἠγέρθη τὸ παιδάριον.

\vs{32}Καὶ εἰσῆλθεν Ἑλισαιὲ εἰς τὸν οἶκον, καὶ ἰδοὺ τὸ παιδάριον τεθνηκός κεκοιμισμένον ἐπὶ τὴν κλίνην αὐτοῦ·
\vs{33}Καὶ εἰσῆλθεν Ἑλισαιὲ εἰς τὸν οἶκον, καὶ ἀπέκλεισε τὴν θύραν κατὰ τῶν δύο ἑαυτῶν, καὶ προσηύξατο πρὸς Κύριον.
\vs{34}Καὶ ἀνέβη καὶ ἐκοιμήθη ἐπὶ τὸ παιδάριον· καὶ ἔθηκε τὸ στόμα αὐτοῦ ἐπὶ τὸ στόμα αὐτοῦ, καὶ τοὺς ὀφθαλμοὺς αὐτοῦ ἐπὶ τοὺς ὀφθαλμοὺς αὐτοῦ, καὶ τὰς χεῖρας αὐτοῦ ἐπὶ τὰς χεῖρας αὐτοῦ· καὶ διέκαμψεν ἐπʼ αὐτὸν, καὶ διεθερμάνθη ἡ σὰρξ τοῦ παιδαρίου.
\vs{35}Καὶ ἐπέστρεψε, καὶ ἐπορεύθη ἐν τῇ οἰκίᾳ ἔνθεν καὶ ἔνθεν· καὶ ἀνέβη καὶ συνέκαμψεν ἐπὶ τὸ παιδάριον ἕως ἑπτάκις· καὶ ἤνοιξε τὸ παιδάριον τοὺς ὀφθαλμοὺς αὐτοῦ.
\vs{36}Καὶ ἐξεβόησεν Ἑλισαιὲ πρὸς Γιεζὶ, καὶ εἶπε, κάλεσον τὴν Σωμανίτην ταύτην· καὶ ἐκάλεσε, καὶ εἰσῆλθε πρὸς αὐτόν· καὶ εἶπεν Ἑλισαιὲ, λάβε τὸν υἱόν σου.
\vs{37}Καὶ εἰσῆλθεν ἡ γυνὴ, καὶ ἔπεσεν ἐπὶ τοὺς πόδας αὐτοῦ, καὶ προσεκύνησεν ἐπὶ τὴν γῆν· καὶ ἔλαβε τὸν υἱὸν αὐτῆς, καὶ ἐξῆλθε.

\vs{38}Καὶ Ἑλισαιὲ ἐπέστρεψεν εἰς Γάλγαλα· καὶ ὁ λιμὸς ἐν τῇ γῇ, καὶ υἱοὶ τῶν προφητῶν ἐκάθηντο ἐνώπιον αὐτοῦ· καὶ εἶπεν Ἑλισαιὲ τῷ παιδαρίῳ αὐτοῦ, ἐπίστησον τὸν λέβητα τὸν μέγαν, καὶ ἕψε ἕψεμα τοῖς υἱοῖς τῶν προφητῶν.
\vs{39}Καὶ ἐξῆλθεν εἰς τὸν ἀγρὸν συλλέξαι ἀριώθ· καὶ εὗρεν ἄμπελον ἐν τῷ ἀγρῷ, καὶ συνέλεξεν ἀπʼ αὐτῆς τολύπην ἀγρίαν πλῆρες τὸ ἱμάτιον αὐτοῦ, καὶ ἐνέβαλεν εἰς τὸν λέβητα τοῦ ἑψέματος, ὅτι οὐκ ἔγνωσαν, καὶ ἐνέχει τοῖς ἀνδράσι φαγεῖν·
\vs{40}καὶ ἐγένετο ἐν τῷ ἐσθίειν αὐτοὺς ἐκ τοῦ ἑψέματος, καὶ ἰδοὺ ἀνεβόησαν, καὶ εἶπαν, θάνατος ἐν τῷ λέβητι, ἄνθρωπε τοῦ Θεοῦ· καὶ οὐκ ἠδύναντο φαγεῖν.
\vs{41}Καὶ εἶπε, λάβετε ἅλευρον, καὶ ἐμβάλετε εἰς τὸν λέβητα· καὶ εἶπεν Ἑλισαιὲ πρὸς Γιεζὶ τὸ παιδάριον, ἔγχει τῷ λαῷ καὶ ἐσθιέτωσαν· καὶ οὐκ ἐγενήθη ἐκεῖ ἔτι ῥῆμα πονηρὸν ἐν τῷ λέβητι.

\vs{42}Καὶ ἀνὴρ διῆλθεν ἐκ Βαιθαρισὰ, καὶ ἤνεγκε πρὸς τὸν ἄνθρωπον τοῦ Θεοῦ πρωτογεννημάτων εἴκοσι ἄρτους κριθίνους καὶ παλάθας· καὶ εἶπε, δότε τῷ λαῷ καὶ ἐσθιέτωσαν.
\vs{43}Καὶ εἶπεν ὁ λειτουργὸς αὐτοῦ, τί δῶ τοῦτο ἐνώπιον ἑκατὸν ἀνδρῶν; καὶ εἶπε, δὸς τῷ λαῷ καὶ ἐσθιέτωσαν, ὅτι τάδε λέγει Κύριος, φάγονται καὶ καταλείψουσι.
\vs{44}Καὶ ἔφαγον καὶ κατέλιπον, κατὰ τὸ ῥῆμα Κυρίου.

\ch{5}
Καὶ Ναιμὰν ὁ ἄρχων τῆς δυνάμεως Συρίας ἦν ἀνὴρ μέγας ἐνώπιον τοῦ κυρίου αὐτοῦ, καὶ τεθαυμασμένος προσώπῳ, ὅτι ἐν ὑτῶ ἔδωκε Κύριος σωτηρίαν Συρίᾳ· καὶ ὁ ἀνὴρ ἦν δυνατὸς ἰσχύϊ, λελεπρωμένος.
\vs{2}Καὶ Συρία ἐξῆλθον μονόζωνοι, καὶ ᾐχμαλώτευσαν ἐκ γῆς Ἰσραὴλ νεάνιδα μικρὰν, καὶ ἦν ἐνώπιον τῆς γυναικὸς Ναιμάν.
\vs{3}Ἡ δὲ εἶπε τῇ κυρίᾳ αὐτῆς, ὄφελον ὁ κύριός μου ἐνώπιον τοῦ προφήτου τοῦ Θεοῦ τοῦ ἐν Σαμαρείᾳ, τότε ἀποσυνάξει αὐτὸν ἀπὸ τῆς λέπρας αὐτοῦ.
\vs{4}Καὶ εἰσῆλθε καὶ ἀπήγγειλε τῷ κυρίῳ ἑαυτῆς, καὶ εἶπεν, οὕτως καὶ οὕτως ἐλάλησεν ἡ νεᾶνις ἡ ἐκ γῆς Ἰσραήλ.

\vs{5}Καὶ εἶπε βασιλεὺς Συρίας πρὸς Ναιμὰν, δεῦρο, εἴσελθε καὶ ἐξαποστελῶ βιβλίον πρὸς βασιλέα Ἰσραήλ· καὶ ἐπορεύθη, καὶ ἔλαβεν ἐν τῇ χειρὶ αὐτοῦ δέκα τάλαντα ἀργυρίου, καὶ ἑξακισχιλίους χρυσοῦς, καὶ δέκα ἀλλασσομένας στολὰς.
\vs{6}Καὶ ἤνεγκε τὸ βιβλίον πρὸς τὸν βασιλέα Ἰσραὴλ, λέγων, καὶ νῦν ὡς ἂν ἔλθῃ τὸ βιβλίον τοῦτο πρὸς σὲ, ἰδοὺ ἀπέστειλα πρὸς σὲ Ναιμὰν τὸν δοῦλόν μου, καὶ ἀποσυνάξεις αὐτὸν ἀπὸ τῆς λέπρας αὐτοῦ.
\vs{7}Καὶ ἐγένετο ὡς ἀνέγνω βασιλεὺς Ἰσραὴλ τὸ βιβλίον, διέῤῥηξε τὰ ἱμάτια αὐτοῦ, καὶ εἶπεν, ὁ Θεὸς ἐγὼ τοῦ θανατῶσαι καὶ ζωοποιῆσαι, ὅτι οὗτος ἀποστέλλει πρὸς μὲ ἀποσυνάξαι ἄνδρα ἀπὸ τῆς λέπρας αὐτοῦ; ὅτι πλὴν γνῶτε δὴ καὶ ἴδετε ὅτι προφασίζεται οὗτός μοι.

\vs{8}Καὶ ἐγένετο ὡς ἤκουσεν Ἐλισαιὲ, ὅτι διέῤῥηξεν ὁ βασιλεὺς Ἰσραὴλ τὰ ἱμάτια αὐτοῦ, καὶ ἀπέστειλε πρὸς τὸν βασιλέα Ἰσραὴλ, λέγων, ἱνατί διέῤῥηξας τὰ ἱμάτιά σου; ἐλθέτω δὴ πρὸς μὲ Ναιμὰν, καὶ γνώτω ὅτι ἐστι προφήτης ἐν Ἰσραήλ.

\vs{9}Καὶ ἦλθε Ναιμὰν ἐν ἵππῳ καὶ ἅρματι, καὶ ἔστη ἐπὶ θύρας οἴκου Ἐλισαιέ.
\vs{10}Καὶ ἀπέστειλεν Ἐλισαιὲ ἄγγελον πρὸς αὐτὸν, λέγων, πορευθεὶς λοῦσαι ἑπτάκις ἐν τῷ Ἰορδάνῃ, καὶ ἐπιστρέψει ἡ σάρξ σου σοὶ καὶ καθαρισθήσῃ.
\vs{11}Καὶ ἐθυμώθη Ναιμάν καὶ ἀπῆλθε, καὶ εἶπεν, ἰδοὺ εἶπον, πρὸς μὲ πάντως ἐξελεύσεται καὶ στήσεται, καὶ ἐπικαλέσεται ἐν ὀνόματι Θεοῦ αὐτοῦ, καὶ ἐπιθήσει τὴν χεῖρα αὐτοῦ ἐπὶ τὸν τόπον, καὶ ἀποσυνάξει τὸ λεπρόν.
\vs{12}Οὐχὶ ἀγαθὸς Ἀβανὰ καὶ Φαρφὰρ ποταμοὶ Δαμασκοῦ ὑπὲρ πάντα τὰ ὕδατα Ἰσραὴλ; οὐχὶ πορευθεὶς λούσομαι ἐν αὐτοῖς, καὶ καθαρισθήσομαι; καὶ ἐξέκλινε καὶ ἀπῆλθεν ἐν θυμῷ.
\vs{13}Καὶ ἤγγισαν οἱ παῖδες αὐτοῦ, καὶ ἐλάλησαν πρὸς αὐτὸν, μέγαν λόγον ἐλάλησεν ὁ προφήτης πρὸς σέ· οὐχὶ ποιήσεις; καὶ ὅτι εἶπε πρὸς σέ, λοῦσαι καὶ καθαρίσθητι.
\vs{14}Καὶ κατέβη Ναιμὰν καὶ ἐβαπτίσατο ἐν τῷ Ἰορδάνῃ ἑπτάκις κατὰ τὸ ῥῆμα Ἐλισαῖε· καὶ ἐπέστρεψεν ἡ σὰρξ αὐτοῦ ὡς σὰρξ παιδαρίου μικροῦ, καὶ ἐκαθαρίσθη.

\vs{15}Καὶ ἐπέστρεψε πρὸς Ἐλισαιὲ αὐτὸς καὶ πᾶσα ἡ παρεμβολὴ αὐτοῦ, καὶ ἦλθε καὶ ἔστη ἐνώπιον αὐτοῦ, καὶ εἶπεν, ἰδοὺ ἔγνωκα ὅτι οὐκ ἔστι Θεὸς ἐν πάσῃ τῇ γῇ, ὅτι ἀλλʼ ἢ ἐν τῷ Ἰσραήλ· καὶ νῦν λάβε τὴν εὐλογίαν παρὰ τοῦ δούλου σου.
\vs{16}Καὶ εἶπεν Ἐλισαιὲ ζῇ Κύριος ᾧ παρέστην ἐνώπιον αὐτοῦ, εἰ λήψομαι· καὶ παρεβιάσατο αὐτὸν λαβεῖν, καὶ ἠπείθησε·
\vs{17}Καὶ εἶπε Ναιμὰν, καὶ εἰ μὴ, δοθήτω δὴ τῷ δούλῳ σου γόμος ζεῦγος ἡμιόνων, καὶ σύ μοι δώσεις ἐκ τῆς γῆς τῆς πυῤῥᾶς, ὅτι οὐ ποιήσει ἔτι ὁ δοῦλός σου ὁλοκαύτωμα καὶ θυσίασμα θεοῖς ἑτέροις ἀλλʼ ἢ τῷ Κυρίῳ τῷ ῥήματι τούτῳ.
\vs{18}Καὶ ἱλάσεται Κύριος τῷ δούλῳ σου ἐν τῷ εἰσπορεύεσθαι τὸν κύριόν μου εἰς οἶκον Ῥεμμὰν προσκυνῆσαι ἐκεῖ· καὶ αὐτὸς ἐπαναπαύσεται ἐπὶ τῆς χειρός μου, καὶ προσκυνήσω ἐν οἴκῳ Ῥεμμὰν ἐν τῷ προσκυνεῖν αὐτὸν ἐν οἴκῳ Ῥεμμάν· καὶ ἱλάσεται δὴ Κύριος τῷ δούλῳ σου ἐν τῷ λόγῳ τούτῳ.
\vs{19}Καὶ εἶπεν Ἐλισαιὲ πρὸς Ναιμὰν, δεῦρο εἰς εἰρήνην· καὶ ἀπῆλθεν ἀπʼ αὐτοῦ εἰς δεβραθὰ τῆς γῆς.

\vs{20}Καὶ εἶπε Γιεζὶ τὸ παιδάριον Ἐλισαιὲ, ἰδοὺ ἐφείσατο ὁ κύριός μου τοῦ Ναιμὰν τοῦ Σύρου τούτου, τοῦ μὴ λαβεῖν ἐκ χειρὸς αὐτοῦ ἃ ἐνήνοχε· ζῇ Κύριος, ὅτι εἰ μὴ δραμοῦμαι ὀπίσω αὐτοῦ, καὶ λήψομαι ἀπʼ αὐτοῦ τι.
\vs{21}Καὶ ἐδίωξε Γιεζὶ ὀπίσω τοῦ Ναιμάν· καὶ εἶδεν αὐτὸν Ναιμὰν τρέχοντα ὀπίσω αὐτοῦ, καὶ ἐπέστρεψεν ἀπὸ τοῦ ἅρματος εἰς ἀπαντὴν αὐτοῦ.
\vs{22}Καὶ εἶπεν, εἰρήνη· ὁ κύριός μου ἀπεστειλέ με, λέγων, ἰδοὺ νῦν ἦλθον πρὸς μὲ δύο παιδάρια ἐξ ὄρους Ἐφραὶμ ἀπὸ τῶν υἱῶν τῶν προφητῶν· δὸς δὴ αὐτοῖς τάλαντον ἀργυρίου, καὶ δύο ἀλλασσομένας στολάς.
\vs{23}Καὶ εἶπε, λάβε διτάλαντον ἀργυρίου· καὶ ἔλαβε δύο τάλαντα ἀρλυρίον ἐν δυσὶ θυλάκοις, καὶ δύο ἀλλασσομένας στολὰς, καὶ ἔδωκεν ἐπὶ δύο παιδάρια αὐτοῦ, καὶ ᾖραν ἔμπροσθεν αὐτοῦ.
\vs{24}Καὶ ἦλθεν εἰς τὸ σκοτεινὸν, καὶ ἔλαβεν ἐκ τῶν χειρῶν αὐτῶν, καὶ παρέθετο ἐν οἴκῳ, καὶ ἐξαπέστειλε τοὺς ἄνδρας.

\vs{25}Καὶ αὐτὸς εἰσῆλθε, καὶ παρειστήκει πρὸς τὸν κύριον αὐτοῦ· καὶ εἶπε πρὸς αὐτὸν Ἐλισαιὲ, πόθεν Γιεζί; καὶ εἶπε Γιεζί, Οὐ πεπόρευται ὁ δοῦλός σου ἔνθα καὶ ἔνθα.
\vs{26}Καὶ εἶπε πρὸς αὐτὸν Ἐλισαιὲ, οὐχὶ ἡ καρδία μου ἐπορεύθη μετὰ σοῦ ὅτε ἐπέστρεψεν ὁ ἀνὴρ ἀπὸ τοῦ ἅρματος εἰς συναντήν σοι; καὶ νῦν ἔλαβες τὸ ἀργύριον, καὶ νῦν ἔλαβες τὰ ἱμάτια, καὶ ἐλαιῶνας καὶ ἀμπελῶνας καὶ πρόβατα καὶ βόας καὶ παῖδας καὶ παιδίσκας.
\vs{27}Καὶ ἡ λέπρα Ναιμὰν κολληθήσεται ἐν σοὶ καὶ ἐν τῷ σπέρματί σου εἰς τὸν αἰῶνα· καὶ ἐξῆλθεν ἐκ προσώπου αὐτοῦ λελεπρωμένος ὡσεὶ χιών.

\ch{6}
Καὶ εἶπον υἱοὶ τῶν προφητῶν πρὸς Ἐλισαιὲ, ἰδοὺ δὴ ὁ τόπος ἐν ᾧ ἡμεῖς οἰκοῦμεν ἐνώπιόν σου στενὸς ἀφʼ ἡμῶν.
\vs{2}Πορευθῶμεν δὴ ἕως τοῦ Ἰορδάνου, καὶ λάβωμεν ἐκεῖθεν ἀνὴρ εἷς δοκὸν μίαν, καὶ ποιήσωμεν ἑαυτοῖς ἐκεῖ τοῦ οἰκεῖν ἐκεῖ· καὶ εἶπε, δεῦτε.
\vs{3}Καὶ εἶπεν ὁ εἷς ἐπιεικῶς, δεῦρο μετὰ τῶν δούλων σου· καὶ εἶπεν, ἐγὼ πορεύσομαι.
\vs{4}Καὶ ἐπορεύθη μετʼ αὐτῶν, καὶ ἦλθον εἰς τὸν Ἰορδάνην, καὶ ἔτεμνον τὰ ξύλα.
\vs{5}Καὶ ἰδοὺ ὁ εἷς καταβάλλων τὴν δοκὸν, καὶ τὸ σιδήριον ἐξέπεσεν εἰς τὸ ὕδωρ, καὶ ἐβόησεν, ὤ κύριε, καὶ αὐτὸ κεκρυμμένον.
\vs{6}Καὶ εἶπεν ὁ ἄνθρωπος τοῦ Θεοῦ, ποῦ ἔπεσε; καὶ ἔδειξεν αὐτῷ τὸν τόπον· καὶ ἀπέκνισε ξύλον καὶ ἔῤῥιψεν ἐκεῖ, καὶ ἐπεπόλασεν τὸ σιδήριον.
\vs{7}Καὶ εἴρηκεν, ὕψωσον σεαυτῷ· καὶ ἐξέτεινε τὴν χεῖρα, καὶ ἔλαβεν αὐτό.

\vs{8}Καὶ ὁ βασιλεὺς Συρίας ἦν πολεμῶν ἐν Ἰσραήλ· καὶ ἐβουλεύσατο πρὸς τοὺς παῖδας αὐτοῦ, λέγων, εἰς τὸν τόπον τόνδε τινὰ ἐλμωνὶ παρεμβαλῶ.
\vs{9}Καὶ ἀπέστειλεν Ἐλισαιὲ πρὸς τὸν βασιλέα Ἰσραὴλ, λέγων, φύλαξαι μὴ παρελθεῖν ἐν τῷ τόπῳ τούτῳ, ὅτι ἐκεῖ Συρία κέκρυπται.
\vs{10}Καὶ ἀπέστειλεν ὁ βασιλεὺς Ἰσραὴλ εἰς τὸν τόπον ὃν εἶπεν αὐτῷ Ἐλισαιὲ, καὶ ἐφυλάξατο ἐκεῖθεν οὐ μίαν οὐδὲ δύο.

\vs{11}Καὶ ἐξεκινήθη ἡ ψυχὴ βασιλέως Συρίας περὶ τοῦ λόγου τούτου· καὶ ἐκάλεσε τοὺς παῖδας αὐτοῦ, καὶ εἶπε πρὸς αὐτοὺς, οὐκ ἀναγγελεῖτέ μοι τίς προδίδωσί με βασιλεῖ Ἰσραήλ;
\vs{12}Καὶ εἶπεν εἷς τῶν παίδων αὐτοῦ, οὐχί κύριέ μου βασιλεῦ, ὅτι Ἐλισαιὲ ὁ προφήτης ὁ ἐν Ἰσραὴλ ἀναγγέλλει τῷ βασιλεῖ Ἰσραὴλ πάντας τοὺς λόγους, οὓς ἐὰν λαλήσῃς ἐν τῷ ταμείῳ τοῦ κοιτῶνός σου.
\vs{13}Καὶ εἶπε, δεῦτε ἴδετε ποῦ οὗτος, καὶ ἀποστείλας λήψομαι αὐτόν· καὶ ἀπήγγειλαν αὐτῷ, λέγοντες, ἰδοὺ ἐν Δωθαΐμ.

\vs{14}Καὶ ἀπέστειλεν ἐκεῖ ἵππον καὶ ἅρμα καὶ δύναμιν βαρεῖαν, καὶ ἦλθον νυκτὸς καὶ περιεκύκλωσαν τὴν πόλιν.
\vs{15}Καὶ ὤρθρισεν ὁ λειτουργὸς Ἐλισαιὲ ἀναστῆναι, καὶ ἐξῆλθε· καὶ ἰδοὺ δύναμις κυκλοῦσα τὴν πόλιν, καὶ ἵππος καὶ ἅρμα· καὶ εἶπε τὸ παιδάριον πρὸς αὐτὸν, ὦ κύριε, πῶς ποιήσομεν;
\vs{16}Καὶ εἶπεν Ἐλισαιὲ, μὴ φοβοῦ, ὅτι πλείους οἱ μεθʼ ἡμῶν ὑπὲρ τοὺς μετʼ αὐτῶν.
\vs{17}Καὶ προσηύξατο Ἐλισαιὲ, καὶ εἶπε, Κύριε, διάνοιξον δὴ τοὺς ὀφθαλμοὺς τοῦ παιδαρίου καὶ ἰδέτω· καὶ διήνοιξε Κύριος τοὺς ὀφθαλμοὺς αὐτοῦ καὶ εἶδε· καὶ ἰδοὺ τὸ ὄρος πλῆρες ἵππων, καὶ ἅρμα πυρὸς περικύκλῳ Ἐλισαιέ.
\vs{18}Καὶ κατέβησαν πρὸς αὐτόν· καὶ προσηύξατο πρὸς Κύριον, καὶ εἶπε, πάταξον δὴ τὸ ἔθνος τοῦτο ἀορασίᾳ· καὶ ἐπάταξεν αὐτοὺς ἀορασίᾳ, κατὰ τὸ ῥῆμα Ἑλισαιὲ.
\vs{19}Καὶ εἶπε πρὸς αὐτοὺς Ἑλισαιὲ, οὐχὶ αὕτη ἡ πόλις καὶ αὕτη ἡ ὁδός· δεῦτε ὀπίσω μου, καὶ ἄξω ὑμᾶς πρὸς τὸν ἄνδρα ὃν ζητεῖτε· καὶ ἀπήγαγεν αὐτοὺς πρὸς Σαμάρειαν.
\vs{20}Καὶ ἐγένετο ὡς εἰσῆλθον εἰς Σαμάρειαν, καὶ εἶπεν Ἑλισαιὲ, ἄνοιξον δή Κύριε τοὺς ὀφθαλμοὺς αὐτῶν καὶ ἰδέτωσαν· καὶ διήνοιξε Κύριος τοὺς ὀφθαλμοὺς αὐτῶν, καὶ εἶδου· καὶ ἰδοὺ ἦσαν ἐν μέσῳ Σαμαρείας.

\vs{21}Καὶ εἶπεν ὁ βασιλεὺς Ἰσραήλ πρὸς Ἑλισαιὲ, ὡς εἶδεν αὐτοὺς, εἰ πατάξας πατάξω, πάτερ;
\vs{22}Καὶ εἶπεν, Οὐ πατάξεις, εἰ μὴ οὓς ᾐχμαλώτευσας ἐν ῥομφαίᾳ σου καὶ τόξῳ σου σὺ τύπτεις· παράθες ἄρτους καὶ ὕδωρ ἐνώπιον αὐτῶν, καὶ φαγέτωσαν καὶ πιέτωσαν, καὶ ἀπελθέτωσαν πρὸς τὸν κύριον αὐτῶν.
\vs{23}Καὶ παρέθηκεν αὐτοῖς παράθεσιν μεγάλην, καὶ ἔφαγον καὶ ἔπιον· καὶ ἀπέστειλεν αὐτοὺς, καὶ ἀπῆλθον πρὸς τὸν κύριον αὐτῶν· καὶ οὐ προσέθεντο ἔτι μονόζωνοι Συρίας τοῦ ἐλθεῖν εἰς γῆν Ἰσραήλ.

\vs{24}Καὶ ἐγένετο μετὰ ταῦτα, καὶ ἤθροισεν υἱὸς Ἅδερ βασιλεὺς Συρίας πᾶσαν τὴν παρεμβολὴν αὐτοῦ, καὶ ἀνέβη, καὶ περιεκάθισαν ἐπὶ Σαμάρειαν.
\vs{25}Καὶ ἐγένετο λιμὸς μέγας ἐν Σαμαρείᾳ· καὶ ἰδοὺ περιεκάθηντο ἐπʼ αὐτὴν ἕως οὗ ἐγενήθη κεφαλὴ ὄνου πεντήκοντα ἀργυρίου, καὶ τέταρτον τοῦ κάβου κόπρου περιστερῶν πέντε ἀργυρίου.

\vs{26}Καὶ ἦν ὁ βασιλεὺς Ἰσραὴλ διαπορευόμενος ἐπὶ τοῦ τείχους· καὶ γυνὴ ἐβόησε πρὸς αὐτὸν, λέγουσα, Σῶσον κύριε βασιλεῦ.
\vs{27}Καὶ εἶπεν αὐτῇ, μὴ σὲ σώσαι Κύριος, πόθεν σώσω σε; μὴ ἀπὸ ἅλωνος ἢ ἀπὸ ληνοῦ;
\vs{28}Καὶ εἶπεν αὐτῇ ὁ βασιλεύς, τί ἐστι σοι; καὶ εἶπεν ἡ γυνή, αὕτη εἶπε πρὸς μέ, δὸς τὸν υἱόν σου καὶ φαγόμεθα αὐτὸν σήμερον, καὶ τὸν υἱόν μου φαγόμεθα αὐτὸν αὔριον.
\vs{29}Καὶ ἡψήσαμεν τὸν υἱόν μου καὶ ἐφάγομεν αὐτὸν, καὶ εἶπον πρὸς αὐτὴν τῇ ἡμέρᾳ τῇ δευτέρᾳ, δὸς τὸν υἱόν σου καὶ φάγωμεν αὐτόν· καὶ ἔκρυψε τὸν υἱὸν αὐτῆς.
\vs{30}Καὶ ἐγένετο ὡς ἤκουσεν ὁ βασιλεὺς Ἰσραὴλ τοὺς λόγους τῆς γυναικὸς, διέῤῥηξε τὰ ἱμάτια αὐτοῦ, καὶ αὐτὸς διεπορεύετο ἐπὶ τοῦ τείχους, καὶ εἶδεν ὁ λαὸς τὸν σάκκον ἐπὶ τῆς σαρκὸς αὑτοῦ ἔσωθεν.
\vs{31}Καὶ εἶπε, τάδε ποιήσαι μοι ὁ Θεὸς καὶ τάδε προσθείη, εἰ στήσεται ἡ κεφαλὴ Ἑλισαιὲ ἐπʼ αὐτῷ σήμερον.

\vs{32}Καὶ Ἑλισαιὲ ἐκάθητο ἐν τῷ οἴκῳ αὐτοῦ, καὶ οἱ πρεσβύτεροι ἐκάθηντο μετʼ αὐτοῦ· καὶ ἀπέστειλεν ἄνδρα πρὸ προσώπου αὐτοῦ· πρὶν ἐλθεῖν τὸν ἄγγελον πρὸς αὐτόν, καὶ αὐτὸς εἶπε πρὸς τοὺς πρεσβυτέρους, εἰ εἴδετε ὅτι ἀπέστειλεν ὁ υἱὸς τοῦ φονευτοῦ οὗτος ἀφελεῖν τὴν κεφαλήν μου; ἴδετε ὡς ἂν ἔλθῃ ὁ ἄγγελος, ἀποκλείσατε τὴν θύραν, καὶ παραθλίψατε αὐτὸν ἐν τῇ θύρᾳ· οὐχὶ φωνὴ τῶν ποδῶν τοῦ κυρίου αὐτοῦ κατόπισθεν αὐτοῦ;
\vs{33}Ἔτι αὐτοῦ λαλοῦντος μετʼ αὐτῶν, καὶ ἰδοὺ ἄγγελος κατέβη πρὸς αὐτὸν, καὶ εἶπεν, ἰδοὺ αὕτη ἡ κακία παρὰ Κυρίου· τί ὑπομείνω τῷ Κυρίῳ ἔτι;

\ch{7}
Καὶ εἶπεν Ἑλισαιὲ ἄκουσον λόγον Κυρίου· τάδε λέλει Κύριος, ὡς ἡ ὥρα αὕτη, αὔριον μέτρον σεμιδάλεως σίκλου, καὶ δίμετρον κριθῶν σίκλου, ἐν ταῖς πύλαις Σαμαρείας.
\vs{2}Καὶ ἀπεκρίθη ὁ τριστάτης ἐφʼ ὃν ὁ βασιλεὺς ἐπανεπαύετο ἐπὶ τὴν χεῖρα αὐτοῦ τῷ Ἑλισαιὲ, καὶ εἶπεν, ἰδοὺ ποιήσει Κύριος καταράκτας ἐν οὐρανῷ, μὴ ἔσται τὸ ῥῆμα τοῦτο; καὶ Ἑλισαιὲ εἶπεν, ἰδοὺ σὺ ὄψει τοῖς ὀφθαλμοῖς σου, καὶ ἐκεῖθεν οὐ φάγῃ.

\vs{3}Καὶ τέσσαρες ἄνδρες ἦσαν λεπροὶ παρὰ τὴν θύραν τῆς πόλεως, καὶ εἶπεν ἀνὴρ πρὸς τὸν πλησίον αὐτοῦ, τί ἡμεῖς καθήμεθα ὧδε ἕως ἀποθάνωμεν;
\vs{4}Ἐὰν εἴπωμεν, εἰσέλθωμεν εἰς τὴν πόλιν, καὶ ὁ λιμὸς ἐν τῇ πόλει, καὶ ἀποθανούμεθα ἐκεῖ· καὶ ἐὰν καθίσωμεν ὧδε, καὶ ἀποθανούμεθα· καὶ νῦν δεῦτε, καὶ ἐμπεσωμεν εἰς τὴν παρεμβολὴν Συρίας· ἐὰν ζωογονήσωσιν ἡμᾶς, καὶ ζησόμεθα· καὶ ἐὰν θανατώσωσιν ἡμᾶς, καὶ ἀποθανούμεθα.
\vs{5}Καὶ ἀνέστησαν ἐν τῷ σκότει εἰσελθεῖν εἰς τὴν παρεμβολὴν Συρίας· καὶ ἦλθον εἰς μέρος παρεμβολῆς Συρίας, καὶ ἰδοὺ οὐκ ἔστιν ἀνὴρ ἐκεῖ.
\vs{6}Καὶ Κύριος ἀκουστὴν ἐποίησε παρεμβολὴν τὴν Συρίας φωνὴν ἅρματος καὶ φωνὴν ἵππου, φωνὴν δυνάμεως μεγάλης· καὶ εἶπεν ἀνὴρ πρὸς τὸν ἀδελφὸν αὐτοῦ, νῦν ἐμισθώσατο ἐφʼ ἡμᾶς ὁ βασιλεὺς Ἰσραὴλ τοὺς βασιλέας τῶν Χετταίων καὶ τοὺς βασιλέας Αἰγύπτου τοῦ ἐλθεῖν ἐφʼ ἡμᾶς.
\vs{7}Καὶ ἀνέστησαν καὶ ἀπέδρασαν ἐν τῷ σκότει καὶ ἐγκατέλιπον τὰς σκηνὰς αὐτῶν, καὶ τοὺς ἵππους αὐτῶν, καὶ τοὺς ὄνους αὐτῶν ἐν τῇ παρεμβολῇ ὡς ἔστι, καὶ ἔφυγον πρὸς τὴν ψυχὴν ἑαυτῶν.

\vs{8}Καὶ εἰσῆλθον οἱ λεπροὶ οὗτοι ἕως μέρους τῆς παρεμβολῆς, καὶ εἰσῆλθον εἰς σκηνὴν μίαν, καὶ ἔφαγον, καὶ ἔπιον, καὶ ἦραν ἐκεῖθεν ἀργύριον, καὶ χρυσίον, καὶ ἱματισμόν· καὶ ἐπορεύθησαν, καὶ ἐπέστρεψαν ἐκεῖθεν, καὶ εἰσῆλθον εἰς σκηνὴν ἄλλην, καὶ ἔλαβον ἐκεῖθεν καὶ ἐπορεύθησαν, καὶ κατέκρυψαν.
\vs{9}Καὶ εἶπεν ἀνὴρ πρὸς τὸν πλησιον αὐτοῦ, οὐχ οὕτως ἡμεῖς ποιοῦμεν· ἡ ἡμέρα αὕτη, ἡμέρα εὐαγγελίας ἐστὶ, καὶ ἡμεῖς σιωπῶμεν, καὶ μένομεν ἕως φωτὸς τοῦ πρωῒ, καὶ εὑρήσομεν ἀνομίαν· καὶ νῦν δεῦρο, καὶ εἰσέλθωμεν καὶ ἀναγγείλωμεν εἰς τὸν οἶκον τοῦ βασιλέως.

\vs{10}Καὶ εἰσῆλθον καὶ ἐβόησαν πρὸς τὴν πύλην τῆς πόλεως, καὶ ἀνήγγειλαν αὐτοῖς, λέγοντες, εἰσήλθομεν εἰς τὴν παρεμβολὴν Συρίας, καὶ ἰδοὺ οὐκ ἔστιν ἐκεῖ ἀνὴρ καὶ φωνὴ ἀνθρώπου, ὅτι εἰ μὴ ἵππος δεδεμένος καὶ ὄνος, καὶ αἱ σκηναὶ αὐτῶν ὡς εἰσί.
\vs{11}Καὶ ἐβόησαν οἱ θυρωροὶ, καὶ ἀνήγγειλαν εἰς τὸν οἶκον τοῦ βασιλέως ἔσω.

\vs{12}Καὶ ἀνέστη ὁ βασιλεὺς νυκτὸς, καὶ εἶπε πρὸς τοὺς παῖδας αὐτοῦ, ἀναγγελῶ δὴ ὑμῖν ἃ ἐποίησεν ἡμῖν Συρία· ἔγνωσαν ὅτι πεινῶμεν ἡμεῖς, καὶ ἐξῆλθαν ἐκ τῆς παρεμβολῆς καὶ ἐκρύβησαν ἐν τῷ ἀγρῷ, λέγοντες, ὅτι ἐξελεύσονται ἐκ τῆς πόλεως, καὶ συλληψόμεθα αὐτοὺς ζῶντας, καὶ εἰς τὴν πόλιν εἰσελευσόμεθα.
\vs{13}Καὶ ἀπεκρίθη εἷς τῶν παίδων αὐτοῦ καὶ εἶπε, λαβέτωσαν δὴ πέντε τῶν ἵππων τῶν ὑπολελειμμένων οἳ κατελείφθησεν ὧδε, ἰδού εἰσι πρὸς πᾶν τὸ πλῆθος Ἰσραὴλ τὸ ἐκλεῖπον, καὶ ἀποστελοῦμεν ἐκεῖ καὶ ὀψόμεθα.
\vs{14}Καὶ ἔλαβον δύο ἐπιβάτας ἵππων· καὶ ἀπέστειλεν ὁ βασιλεὺς Ἰσραὴλ ὀπίσω τοῦ βασιλέως Συρίας, λέγων, δεῦτε, καὶ ἴδετε.
\vs{15}Καὶ ἐπορεύθησαν ὀπίσω αὐτῶν ἕως τοῦ Ἰορδάνου. καὶ ἰδοὺ πᾶσα ἡ ὁδὸς πλήρης ἱματίων καὶ σκευῶν ὧν ἔῤῥιψε Συρία ἐν τῷ θαμβεῖσθαι αὐτούς· καὶ ἐπέστρεψαν οἱ ἄγγελοι καὶ ἀνήγγειλαν τῷ βασιλεῖ.

\vs{16}Καὶ ἐξῆλθεν ὁ λαὸς καὶ διήρπασεν τὴν παρεμβολὴν Συρίας· καὶ ἐγένετο μέτρον σεμισάλεως σίκλου, κατὰ τὸ ῥῆμα Κυρίου, καὶ δίμετρον κριθῶν σίκλου.
\vs{17}Καὶ ὁ βασιλεὺς κατέστησε τὸν τριστάτην ἐφʼ ὃν ὁ βασιλεὺς ἐπανεπαύετο τῇ χειρὶ αὐτοῦ ἐπὶ τῆς πυλῆς· καὶ συνεπάτησεν αὐτὸν ὁ λαὸς ἐν τῇ πύλῃ, καὶ ἀπέθανε καθὰ ἐλάλησεν ὁ ἄνθρωπος τοῦ Θεοῦ, ὃς ἐλάλησεν ἐν τῷ καταβῆναι τὸν ἄγγελον πρὸς αὐτόν.
\vs{18}Καὶ ἐγένετο καθὰ ἐλάλησεν Ἑλισαιὲ πρὸς τὸν βασιλέα, λέγων, δίμετρον κριθῆς σίκλου καὶ μέτρον σεμιδάλεως σίκλου· καὶ ἔσται ὡς ἡ ὥρα αὔριον ἐν τῇ πύλῃ Σαμαρείας.
\vs{19}Καὶ ἀπεκρίθη ὁ τριστάτης τῷ Ἑλισαιὲ, καὶ εἶπεν, ἰδοὺ Κύριος ποιεῖ καταράκτας ἐν τῷ οὐρανῷ, μὴ ἔσται τὸ ῥῆμα τοῦτο; καὶ εἶπεν Ἑλισαιὲ, ἰδοὺ ὄψει τοῖς ὀφθαλμοῖς σου, καὶ ἐκεῖθεν οὐ μὴ φάγῃ.
\vs{20}Καὶ ἐγένετο οὕτως, καὶ συνεπάτησαν αὐτὸν ὁ λαὸς ἐν τῇ πύλῃ, καὶ ἀπέθανε.

\ch{8}
Καὶ Ἑλισαιὲ ἐλάλησε πρὸς τὴν γυναῖκα, ἧς ἐζωπύρησε τὸν υἱὸν, λέγων, ἀνάστηθι καὶ δεῦρο σὺ καὶ ὁ οἶκός σου, καὶ παροίκει οὗ ἐὰν παροικήσῃς, ὅτι κέκληκε Κύριος λιμὸν ἐπὶ τὴν γῆν, καί γε ἦλθεν ἐπὶ τὴν γῆν ἑπτὰ ἔτη.
\vs{2}Καὶ ἀνέστη ἡ γυνὴ, καὶ ἐποίησε κατὰ τὸ ῥῆμα Ἑλισαιὲ καὶ αὐτὴ καὶ ὁ οἶκος αὐτῆς, καὶ παρῴκει ἐν γῇ ἀλλοφύλων ἑπτὰ ἔτη.

\vs{3}Καὶ ἐγένετο μετὰ τὸ τέλος τῶν ἑπτὰ ἐτῶν, καὶ ἐπέστρεψεν ἡ γυνὴ ἐκ γῆς ἀλλοφύλων εἰς τὴν πόλιν, καὶ ἦλθε βοῆσαι πρὸς τὸν βασιλέα περὶ τοῦ οἴκου ἑαυτῆς καὶ περὶ τῶν ἀγρῶν αὐτῆς.
\vs{4}Καὶ ὁ βασιλεὺς ἐλάλει πρὸς Γιεζὶ τὸ παιδάριον Ἑλισαιὲ τοῦ ἀνθρώπου τοῦ Θεοῦ, λέγων, διήγησαι δὴ ἐμοὶ πάντα τὰ μεγάλα ἃ ἐποίησεν Ἑλισαιέ.
\vs{5}Καὶ ἐγένετο αὐτοῦ ἐξηγουμένου τῷ βασιλεῖ, ὡς ἑζωπύρησεν υἱὸν τεθνηκότα, καὶ ἰδοὺ ἡ γυνὴ ἧς ἐζωπύρησε τὸν υἱὸν αὐτῆς Ἑλισαιὲ, βοῶσα πρὸς τὸν βασιλέα περὶ τοῦ οἴκου ἑαυτῆς καὶ περὶ τῶν ἀγρῶν ἑαυτῆς· καὶ εἶπε Γιεζὶ, Κύριε βασιλεῦ, αὕτη ἡ γυνὴ, καὶ οὗτος ὁ υἱὸς αὐτῆς, ὃν ἐζωπύρησεν Ἑλισαιέ.
\vs{6}Καὶ ἐπηρώτησεν ὁ βασιλεὺς τὴν γυναῖκα· καὶ διηγήσατο αὐτῷ· καὶ ἔδωκεν αὐτῇ ὁ βασιλεὺς εὐνοῦχον ἕνα, λέγων, ἐπίστρεψον πάντα τὰ αὐτῆς, καὶ πάντα τὰ γεννήματα τοῦ ἀγροῦ ἀπὸ τῆς ἡμέρας ἧς κατέλιπε τὴν γῆν ἕως τοῦ νῦν.

\vs{7}Καὶ ἦλθεν Ἑλισαιὲ εἰς Δαμασκόν· καὶ υἱὸς Ἄδερ βασιλεὺς Συρίας ἠῤῥώστησε, καὶ ἀνήγγειλαν αὐτῷ, λέγοντες, ἥκει ὁ ἄνθρωπος τοῦ Θεοῦ ἕως ᾧδε.
\vs{8}Καὶ εἶπεν ὁ βασιλεὺς πρὸς Ἀζαὴλ, λάβε ἐν τῇ χειρί σου μαναὰ, καὶ δεῦρο εἰς ἀπαντὴν τοῦ ἀνθρώπου τοῦ Θεοῦ, καὶ ἐπιζήτησον τὸν Κύριον παρʼ αὐτοῦ, λέγων, εἰ ζήσομαι ἐκ τῆς ἀῤῥωστίας μου ταύτης;
\vs{9}Καὶ ἐπορεύθη Ἀζαὴλ εἰς ἀπαντὴν αὐτοῦ, καὶ ἔλαβε μαναὰ ἐν τῇ χειρὶ αὐτοῦ, καὶ πάντα τὰ ἀγαθὰ Δαμασκοῦ, ἄρσιν τεσσαράκοντα καμήλων, καὶ ἦλθε καὶ ἔστη ἐνώπιον αὐτοῦ, καὶ εἶπε πρὸς Ἑλισαιὲ, υἱός σου υἱὸς Ἅδερ βασιλεὺς Συρίας ἀπέστειλέ με πρὸς σὲ ἐπερωτῆσαι, λέγων, εἰ ζήσομαι ἐκ τῆς ἀῤῥωστίας μου ταύτης;
\vs{10}Καὶ εἶπεν Ἑλισαιὲ, δεῦρο, εἶπον, ζωῇ ζήσῃ, καὶ ἔδειξέ μοι Κύριος ὅτι θανάτῳ ἀποθανῇ.
\vs{11}Καὶ παρέστη τῷ προσώπῳ αὐτοῦ, καὶ ἔθηκεν ἕως αἰσχύνης· καὶ ἔκλαυσεν ὁ ἄνθρωπος τοῦ Θεοῦ.
\vs{12}Καὶ εἶπεν Ἀζαὴλ, τί ὅτι ὁ κύριός μου κλαίει; καὶ εἶπεν, ὅτι οἶδα ὅσα ποιήσεις τοῖς υἱοῖς Ἰσραὴλ κακά· τὰ ὀχυρώματα αὐτῶν ἐξαποστελεῖς ἐν πυρὶ, καὶ τοὺς ἐκλεκτοὺς αὐτῶν ἐν ῥομφαίᾳ ἀποκτενεῖς, καὶ τὰ νήπια αὐτῶν ἐνσείσεις, καὶ τὰς ἐν γαστρὶ ἐχούσας αὐτῶν ἀναῥῥήξεις.
\vs{13}Καὶ εἶπεν Ἀζαὴλ, τίς ἐστιν ὁ δοῦλός σου, ὁ κύων ὁ τεθνηκὼς, ὅτι ποιήσει τὸ ῥῆμα τοῦτο; καὶ εἶπεν Ἑλισαιὲ, ἔδειξέ μοι Κύριός σε βασιλεύοντα ἐπὶ Συρίαν.
\vs{14}Καὶ ἀπῆλθεν ἀπὸ Ἑλισαιὲ, καὶ εἰσῆλθε πρὸς τὸν κύριον αὐτοῦ, καὶ εἶπεν αὐτῷ, τί εἶπέ σοι Ἑλισαιέ; καὶ εἶπεν, εἶπέ μοι, ζωῇ ζήσῃ.
\vs{15}Καὶ ἐγένετο τῇ ἐπαύριον, καὶ ἔλαβε τὸ μαχβὰρ καὶ ἔβαψεν ἐν τῷ ὕδατι, καὶ περιέβαλεν ἐπὶ τὸ πρόσωπον αὐτοῦ, καὶ ἀπέθανε· καὶ ἐβασίλευσεν Ἀζαὴλ ἀντʼ αὐτοῦ.

\vs{16}Ἐν ἔτει πέμπτῳ τῷ Ἰωρὰμ υἱῷ Ἀχαὰβ βασιλεῖ Ἰσραὴλ, καὶ Ἰωσαφὰτ βασιλεῖ Ἰούδα, ἐβασίλευσεν Ἰωρὰμ υἱὸς Ἰωσαφὰτ βασιλεὺς Ἰούδα.
\vs{17}Υἱὸς τριάκοντα καὶ δύο ἐτῶν ἦν ἐν τῷ βασιλεύειν αὐτὸν, καὶ ὀκτὼ ἔτη ἐβασίλευσεν ἐν Ἱερουσαλήμ.
\vs{18}Καὶ ἐπορεύθη ἐν ὁδῷ βασιλέων Ἰσραὴλ καθὼς ἐποίησεν οἶκος Ἀχαὰβ, ὅτι θυγάτηρ Ἀχαὰβ ἦν αὐτῷ εἰς γυναῖκα, καὶ ἐποίησε τὸ πονηρὸν ἐνώπιον Κυρίου.
\vs{19}Καὶ οὐκ ἠθέλησε Κύριος διαφθεῖραι τὸν Ἰούδαν διὰ Δαυὶδ τὸν δοῦλον αὐτοῦ, καθὼς εἶπε δοῦναι αὐτῷ λύχνον καὶ τοῖς υἱοῖς αὐτοῦ πάσας τὰς ἡμέρας.

\vs{20}Ἐν ταῖς ἡμέραις αὐτοῦ ἠθέτησεν Ἐδὼμ ὑποκάτωθεν χειρὸς Ἰούδα, καὶ ἐβασίλευσαν ἐφʼ ἑαυτὸν βασιλέα.
\vs{21}Καὶ ἀνέβη Ἰωρὰμ εἰς Σιὼρ, καὶ πάντα τὰ ἅρματα τὰ μετʼ αὐτοῦ· καὶ ἐγένετο αὐτοῦ ἀναστάντος, καὶ ἐπάταξε τὸν Ἐδὼμ τὸν κυκλώσαντα ἐπʼ αὐτὸν, καὶ τοὺς ἄρχοντας τῶν ἁρμάτων, καὶ ἔφυγεν ὁ λαὸς εἰς τὰ σκηνώματα αὐτῶν.
\vs{22}Καὶ ἠθέτησεν Ἐδὼμ ὑποκάτω τῆς χειρὸς Ἰούδα ἕως τῆς ἡμέρας ταύτης· τότε ἠθέτησε Λοβνὰ ἐν τῷ καιρῷ ἐκείνῳ.

\vs{23}Καὶ τὰ λοιπὰ τῶν λόγων Ἰωρὰμ καὶ πάντα ὅσα ἐποίησεν, οὐκ ἰδοὺ ταῦτα γέγραπται ἐπὶ βιβλίῳ λόγων τῶν ἡμερῶν τοῖς βασιλεῦσιν Ἰούδα;
\vs{24}Καὶ ἐκοιμήθη Ἰωρὰμ μετὰ τῶν πατέρων αὐτοῦ, καὶ ἐτάφη μετὰ τῶν πατέρων αὐτοῦ ἐν πόλει Δαυὶδ τοῦ πατρὸς αὐτοῦ· καὶ ἐβασίλευσεν Ὀχοζίας υἱὸς αὐτοῦ ἀντʼ αὐτοῦ.

\vs{25}Ἐν ἔτει δωδεκάτῳ τῷ Ἰωρὰμ υἱῷ Ἀχαὰβ βασιλεῖ Ἰσραὴλ ἐβασίλευσεν Ὀχοζίας υἱὸς Ἰωράμ.
\vs{26}Υἱὸς εἴκοσι καὶ δύο ἐτῶν Ὀχοζίας ἐν τῷ βασιλεύειν αὐτόν, καὶ ἐνιαυτὸν ἕνα ἐβασίλευσεν ἐν Ἱερουσαλὴμ, καὶ ὄνομα τῆς μητρὸς αὐτοῦ Γοθολία θυγάτηρ Ἀμβρὶ βασιλέως Ἰσραήλ.
\vs{27}Καὶ ἐπορεύθη ἐν ὁδῷ οἴκου Ἀχαὰβ, καὶ ἐποίησε τὸ πονηρὸν ἐνώπιον Κυρίου καθὼς ὁ οἶκος Ἀχαάβ.
\vs{28}Καὶ ἐπορεύθη μετὰ Ἰωρὰμ υἱοῦ Ἀχαὰβ εἰς πόλεμον μετὰ Ἁζαὴλ βασιλέως ἀλλοφύλων ἐν Ῥεμμὼθ Γαλαὰδ, καὶ ἐπάταξαν οἱ Σύροι τὸν Ἰωράμ.
\vs{29}Καὶ ἐπέστρεψεν ὁ βασιλεὺς Ἰωρὰμ τοῦ ἰατρευθῆναι ἐν Ἰεζράελ ἀπὸ τῶν πληγῶν ὧν ἐπάταξαν αὐτὸν ἐν Ῥεμμὼθ, ἐν τῷ πολεμεῖν αὐτὸν μετὰ Ἁζαὴλ βασιλέως Συρίας· καὶ Ὀχοζίας υἱὸς Ἰωρὰμ κατέβη τοῦ ἰδεῖν τὸν Ἰωρὰμ υἱὸν Ἀχαὰβ ἐν Ἰεζράελ, ὅτι ἠῤῥώστει αὐτός.

\ch{9}
Καὶ Ἐλισαιὲ ὁ προφήτης ἐκάλεσεν ἕνα τῶν υἱῶν τῶν προφητῶν, καὶ εἶπεν αὐτῷ, ζῶσαι τὴν ὀσφύν σου, καὶ λάβε τὸν φακὸν τοῦ ἐλαίου τούτου ἐν τῇ χειρί σου, καὶ δεῦρο εἰς Ῥεμμὼθ Γαλαάδ.
\vs{2}Καὶ εἰσελεύσῃ ἐκεῖ, καὶ ὄψει ἐκεῖ Ἰοὺ υἱὸν Ἰωσαφὰτ υἱοῦ Ναμεσσὶ, καὶ εἰσελεύσῃ καὶ ἀναστήσεις αὐτὸν ἐκ μέσου τῶν ἀδελφῶν αὐτοῦ, καὶ εἰσάξεις αὐτὸν εἰς τὸ ταμεῖον ἐν ταμείῳ.
\vs{3}Καὶ λήψῃ τὸν φακὸν τοῦ ἐλαίου, καὶ ἐπιχεεῖς ἐπὶ τὴν κεφαλὴν αὐτοῦ, καὶ εἰπον, τάδε λέγει Κύριος, κέχρικά σε εἰς βασιλέα ἐπὶ Ἰσραήλ· καὶ ἀνοίξεις τὴν θύραν, καὶ φεύξῃ καὶ οὐ μενεῖς.
\vs{4}Καὶ ἐπορεύθη τὸ προδάριον ὁ παιφήτης εἰς Ῥεμμὼθ Γαλαάδ.

\vs{5}Καὶ εἰσῆλθε· καὶ ἰδοὺ οἱ ἄρχοντες τῆς δυνάμεως ἐκάθηντο, καὶ εἶπε, λόγος μοι πρὸς σὲ ὁ ἄρχων· καὶ εἶπεν Ἰοὺ, πρὸς τίνα ἐκ πάντων ἡμῶν; καὶ εἶπε, πρὸς σέ ὁ ἄρχων.
\vs{6}Καὶ ἀνέστη καὶ εἰσῆλθεν εἰς τὸν οἶκον, καὶ ἐπέχεε τὸ ἔλαιον ἐπὶ τὴν κεφαλὴν αὐτοῦ, καὶ εἶπεν αὐτῷ, τάδε λέγει Κύριος ὁ Θεὸς Ἰσραήλ, κέχρικά σε εἰς βασιλέα ἐπὶ λαὸν Κυρίου ἐπὶ τὸν Ἰσραὴλ.
\vs{7}Καὶ ἐξολοθρεύσεις τὸν οἶκον Ἀχαὰβ τοῦ κυρίου σου ἐκ προσώπου σου, καὶ ἐκδικὴσεις τὰ αἵματα τῶν δούλων μου τῶν προφητῶν, καὶ τὰ αἵματα πάντων τῶν δούλων Κυρίου ἐκ χειρὸς Ἰεζάβελ,
\vs{8}καὶ ἐκ χειρὸς ὅλου τοῦ οἴκου Ἀχαὰβ, καὶ ἐξολθρεύσεις τῷ οἴκῳ Ἀχαὰβ οὐροῦντα πρὸς τοῖχον, καὶ συνεχόμενον καὶ ἐγκαταλελειμμένον ἐν Ἰσραήλ.
\vs{9}Καὶ δώσω τὸν οἶκον Ἀχαὰβ ὡς τὸν οἶκον Ἰεροβοὰμ υἱοῦ Ναβὰτ, καὶ ὡς τὸν οἶκον Βαασὰ υἱοῦ Ἀχιά.
\vs{10}Καὶ τὴν Ἰεζάβελ καταφάγονται οἱ κύνες ἐν τῇ μερίδι Ἰσζράελ, καὶ οὐκ ἔστιν ὁ θάπτων· καὶ ἤνοιξε τὴν θύραν καὶ ἔφυγε.

\vs{11}Καὶ Ἰοὺ ἐξῆλθε πρὸς τοὺς παῖδας τοῦ κυρίου αὐτοῦ, καὶ εἶπαν αὐτῷ, εἰρήνη; τί ὅτι εἰσῆλθεν ὁ ἐπίληπτος οὗτος πρὸς σέ; καὶ εἶπεν αὐτοῖς, ὑμεῖς οἴδατε τὸν ἄνδρα καὶ τὴν ἀδολεσχίαν αὐτοῦ.
\vs{12}Καὶ εἶπον, ἄδικον, ἀπάγγειλον δὴ ἡμῖν. Καὶ εἶπεν Ἰοὺ πρὸς αὐτοὺς, οὕτω καὶ οὕτω ἐλάλησε πρὸς μὲ, λέγων, καὶ εἶπε, τάδε λέγει Κύριος, κέχρικά σε εἰς βασιλέα ἐπὶ Ἰσραήλ.
\vs{13}καὶ ἀκούσαντες ἔσπευσαν, καὶ ἔλαβεν ἕκαστος τὸ ἱμάτιον αὐτοῦ, καὶ ἔθηκαν ὑποκάτω αὐτοῦ ἐπὶ τὸ γαρὲμ τῶν ἀναβάθμων· καὶ ἐσάλπισαν ἐν κερατίνῃ, καὶ εἶπαν, ἐβασίλευσεν Ἰού.

\vs{14}Καὶ συνεστράφη Ἰοὺ υἱὸς Ἰωσαφὰτ υἱοῦ Ναμεσσὶ πρὸς Ἰωράμ· καὶ Ἰωρὰμ αὐτὸς ἐφύλασσεν ἐν Ῥεμμὼθ Γαλαὰδ, καὶ πᾶς Ἰσραὴλ ἀπὸ προσώπου Ἀζαὴλ βασιλέως Συρίας.
\vs{15}Καὶ ἀπέστρεψεν Ἰωρὰμ ὁ βασιλεὺς ἰατρευθῆναι ἐν Ἰσζράελ ἀπὸ τῶν πληγῶν ὧν ἔπαισαν αὐτὸν οἱ Σύροι ἐν τῷ πολεμεῖν αὐτὸς μετὰ Ἀζαὴλ βασιλέως Συρίας.

Καὶ εἶπεν Ἰοὺ, εἰ ἔστι ψυχὴ ὑμῶν μετʼ ἐμοῦ, μὴ ἐξελθέτω ἐκ τῆς πόλεως διαπεφευγὼς τοῦ πορευθῆναι καὶ ἀπαγγεῖλαι ἐν Ἰεζράελ.
\vs{16}Καὶ ἵππευσε καὶ ἐπορεύθη Ἰοὺ, καὶ κατέβη εἰς Ἰεζράελ, ὅτι Ἰωρὰμ βασιλεὺς Ἰσραὴλ ἐθεραπεύετο ἐν τῷ Ἰεζράελ ἀπὸ τῶν τοξευμάτων, ὧν κατετόξευσαν αὐτὸν οἱ Ἀραμὶν ἐν τῇ Ῥαμμὰθ ἐν τῷ πολέμῳ μετὰ Ἁζαὴλ βασιλέως Συρίας, ὅτι αὐτὸς δυνατὸς καὶ ἀνὴρ δυνάμεως· καὶ Ὀχοζίας βασιλεὺς Ἰούδα κατέβη ἰδεῖν τὸν Ἰωράμ.
\vs{17}Καὶ ὁ σκοπὸς ἀνέβη ἐπὶ τὸν πύργον Ἰεζράελ, καὶ εἶδε τὸν κονιορτὸν Ἰοὺ ἐν τῷ παραγίνεσθαι αὐτὸν, καὶ εἶπε, Κονιορτὸν ἐγὼ βλέπω· καὶ εἶπεν Ἰωρὰμ, λάβε ἐπιβάτην, καὶ ἀπόστειλον ἔμπροσθεν αὐτῶν, καὶ εἰπάτω ἡ εἰρήνη.
\vs{18}Καὶ ἐπορεύθη ἐπιβάτης ἵππου εἰς ἀπαντὴν αὐτῶν, καὶ εἶπε, τάδε λέγει ὁ βασιλεὺς, ἡ εἰρήνη· καὶ εἶπεν Ἰού, τί σοι καὶ εἰρήνῃ; ἐπίστρεφε εἰς τὰ ὀπίσω μου· καὶ ἀπήγγειλεν ὁ σκοπὸς, λέγων, ἦλθεν ὁ ἄγγελος ἕως αὐτῶν, καὶ οὐκ ἀνέστρεψε.
\vs{19}Καὶ ἀπέστειλεν ἐπιβάτην ἵππου δεύτερον, καὶ ἦλθε πρὸς αὐτὸν καὶ εἶπε, τάδε λέγει ὁ βασιλεύς, ἡ εἰρήνη· καὶ εἶπεν Ἰοὺ, τί σοι καὶ εἰρήνῃ; ἐπιστρέφου εἰς τὰ ὀπίσω μου.
\vs{20}Καὶ ἀπήγγειλεν ὁ σκοπὸς, λέγων, ἦλθεν ἕως αὐτῶν καὶ οὐκ ἀνέστρεψε, καὶ ὁ ἄγων ἦγε τὸν Ἰοὺ υἱὸν Ναμεσσὶ, ὅτι ἐν παραλλαγῇ ἐγένετο.
\vs{21}Καὶ εἶπεν Ἰωρὰμ, ζεῦξον· καὶ ἔζευξεν ἅρμα· καὶ ἐξῆλθεν Ἰωρὰμ βασιλεὺς Ἰσραὴλ, καὶ Ὀχοζίας βασιλεὺς Ἰούδα, ἀνὴρ ἐν τῷ ἅρματι αὐτοῦ, καὶ ἐξῆλθον εἰς ἀπαντὴν Ἰοὺ, καὶ εὗρον αὐτὸν ἐν τῇ μερίδι Ναβουθαὶ τοῦ Ἰεζραηλίτου.

\vs{22}Καὶ ἐγένετο ὡς εἶδεν Ἰωρὰμ τὸν Ἰοὺ, καὶ εἶπεν, ἡ εἰρήνη Ἰοὺ; καὶ εἶπεν Ἰού, τί εἰρήνη; ἔτι αἱ πορνεῖαι Ἰεζάβελ τῆς μητρός σου καὶ τὰ φάρμακα αὐτῆς τὰ πολλά.
\vs{23}Καὶ ἐπέστρεψεν Ἰωρὰμ τὰς χεῖρας αὐτοῦ, καὶ ἔφυγε· καὶ εἶπε πρὸς Ὀχοζίαν, δόλος Ὀχοζία.
\vs{24}Καὶ ἔπλησεν Ἰοὺ τὴν χεῖρα αὐτοῦ ἐν τῷ τόξῳ, καὶ ἐπάταξε τὸν Ἰωρὰμ ἀνὰμέσον τῶν βραχιόνων αὐτοῦ, καὶ ἐξῆλθε τὸ βέλος αὐτοῦ διὰ τῆς καρδίας αὐτοῦ, καὶ ἔκαμψεν ἐπὶ τὰ γόνατα αὐτοῦ.
\vs{25}Καὶ εἶπε πρὸς Βαδεκὰρ τὸν τριστάτην αὐτοῦ, ῥίψον αὐτὸν ἐν τῇ μερίδι ἀγροῦ Ναβουθαὶ τοῦ Ἰεζραηλίτου, ὅτι μνημονεύω ἐγὼ καὶ σὺ ἐπιβεβηκότες ἐπὶ ζεύγη ὀπίσω Ἀχαὰβ τοῦ πατρὸς αὐτοῦ, καὶ Κύριος ἔλαβεν ἐπʼ αὐτὸν τὸ λῆμμα τοῦτο·
\vs{26}Εἰ μὴ τὰ αἵματα Ναβουθαὶ καὶ τὰ αἵματα τῶν υἱῶν αὐτοῦ εἶδον ἐχθὲς, φησὶ Κύριος, καὶ ἀνταποδώσω αὐτῷ ἐν τῇ μερίδι ταύτῃ, φησὶ Κύριος· καὶ νῦν ἄρας δὴ ῥίψον αὐτὸν ἐν τῇ μερίδι κατὰ τὸ ῥῆμα Κυρίου.

\vs{27}Καὶ Ὀχοζίας βασιλεὺς Ἰούδα εἶδε καὶ ἔφυγεν ὁδὸν Βαιθγάν· καὶ ἐδίωξεν ὀπίσω αὐτοῦ Ἰοὺ, καὶ εἶπε, καί γε αὐτόν· καὶ ἐπάταξεν αὐτὸν πρὸς τῷ ἅρματι ἐν τῷ ἀναβαίνειν Γαῒ, ἥ ἐστιν Ἰεβλαάμ· καὶ ἔφυγεν εἰς Μαγεδδὼ, καὶ ἀπέθανεν ἐκεῖ.
\vs{28}Καὶ ἐπεβίβασαν αὐτὸν οἱ παῖδες αὐτοῦ ἐπὶ τὸ ἅρμα, καὶ ἤγαγον αὐτὸν εἰς Ἱερουσαλὴμ, καὶ ἔθαψαν αὐτὸν ἐν τῷ τάφῳ αὐτοῦ ἐν πόλει Δαυίδ.

\vs{29}Καὶ ἐν ἔτει ἑνδεκάτῳ Ἰωρὰμ βασιλέως Ἰσραὴλ ἐβασίλευσεν Ὀχοζίας ἐπὶ Ἰούδαν.

\vs{30}Καὶ ἦλθεν Ἰοὺ ἐπὶ Ἰεζράελ· καὶ Ἰεζάβελ ἤκουσε, καὶ ἐστιμμίσατο τοὺς ὀφθαλμοὺς αὐτῆς, καὶ ἠγάθυνε τὴν κεφαλὴν αὐτῆς, καὶ διέκυψε διὰ τῆς θυρίδος.
\vs{31}Καὶ Ἰοὺ εἰσεπορεύετο ἐν τῇ πόλει, καὶ εἶπε, ἡ εἰρήνη Ζαμβρὶ ὁ φονευτὴς τοῦ κυρίου αὐτοῦ;
\vs{32}Καὶ ἐπῇρε τὸ πρόσωπον αὐτοῦ εἰς τὴν θυρίδα, καὶ εἶδεν αὐτὴν, καὶ εἶπε, τίς εἶ σύ; κατάβηθι μετʼ ἐμοῦ· καὶ κατέκυψαν πρὸς αὐτὸν δύο εὐνοῦχοι.
\vs{33}Καὶ εἶπε, κυλίσατε αὐτήν· καὶ ἐκύλισαν αὐτὴν, καὶ ἐῤῥαντίσθη τοῦ αἵματος αὐτῆς πρὸς τὸν τοῖχον καὶ πρὸς τοὺς ἵππους, καὶ συνεπάτησαν αὐτήν.
\vs{34}Καὶ εἰσῆλθε καὶ ἔφαγε καὶ ἔπιεν, καὶ εἶπε, ἐπισκέψασθε δὴ τὴν κατηραμένην ταύτην, καὶ θάψατε αὐτὴν, ὅτι θυλάτηρ βασιλέως ἐστί.
\vs{35}Καὶ ἐπορεύθησαν θάψαι αὐτὴν, καὶ οὐχ εὗρον ἐν αὐτῇ ἄλλο τι ἢ τὸ κρανίον καὶ οἱ πόδες καὶ τὰ ἴχνη τῶν χειρῶν.
\vs{36}Καὶ ἐπέστρεψαν καὶ ἀνήγγειλαν αὐτῷ· καὶ εἶπε, λόγος Κυρίου ὃν ἐλάλησεν ἐν χειρὶ Ἠλιοὺ τοῦ Θεσβίτου, λέγων, ἐν τῇ μερίδι Ἰεζρὰελ καταφάγονται οἱ κύνες τὰς σάρκας Ἰεζάβελ.
\vs{37}Καὶ ἔσται τὸ θνησιμαῖον Ἰεζάβελ ὡς κοπρία ἐπὶ προσώπου τοῦ ἀγροῦ ἐν τῇ μερίδι Ἰεζραλ, ὥστε μὴ εἰπεῖν αὐτοὺς, Ἰεζάβελ.

\ch{10}
Καὶ τῷ Ἀχαὰβ ἑβδομήκοντα υἱοὶ ἐν Σαμαρείᾳ· καὶ ἔγραψεν Ἰοὺ βιβλίον, καὶ ἀπέστειλεν ἐν Σαμαρείᾳ πρὸς τοὺς ἄρχοντας Σαμαρείας, καὶ πρὸς τοὺς πρεσβυτέρους, καὶ πρὸς τοὺς τιθηνοὺς Αχαὰβ, λέγων,
\vs{2}Καὶ νῦν ὡς ἂν ἔλθῃ τὸ βιβλίον τοῦτο πρὸς ὑμᾶς, καὶ μεθʼ ὑμῶν οἱ υἱοὶ τοῦ κυρίου ὑμῶν, καὶ μεθʼ ὑμῶν τὸ ἅρμα καὶ οἱ ἵπποι καὶ πόλεις ὀχυραὶ καὶ τὰ ὅπλα,
\vs{3}καὶ ὄψεσθε τὸν ἀγαθὸν καὶ τὸν εὐθῆ ἐν τοῖς υἱοῖς τοῦ κυρίου ὑμῶν, καὶ καταστήσετε αὐτὸν ἐπὶ τὸν θρόνον τοῦ πατρὸς αὐτοῦ, καὶ πολεμεῖτε ὑπὲρ τοῦ οἴκου τοῦ κυρίου ὑμῶν.
\vs{4}Καὶ ἐφοβήθησαν σφόδρα, καὶ εἶπον, ἰδοὺ οἱ δύο βασιλεῖς οὐκ ἔστησαν κατὰ πρόσωπον αὐτοῦ, καὶ πῶς στησόμεθα ἡμεῖς;
\vs{5}Καὶ ἀπέστειλαν οἱ ἐπὶ τοῦ οἴκου καὶ οἱ ἐπὶ τῆς πόλεως καὶ οἱ πρεσβύτεροι καὶ οἱ τιθηνοὶ πρὸς Ἰοὺ, λέγοντες, παῖδές σου καὶ ἡμεῖς, καὶ ὅσα ἐὰν εἴπῃς πρὸς ἡμᾶς ποιήσομεν· οὐ βασιλεύσομεν ἄνδρα, τὸ ἀγαθὸν ἐν ὀφθαλμοῖς σου ποιήσομεν.

\vs{6}Καὶ ἔγραψε πρὸς αὐτοὺς Ἰοὺ βιβλίον δεύτερον, λέγων, εἰ ἐμοὶ ὑμεῖς, καὶ τῆς φωνῆς μου ὑμεῖς εἰσακούετε, λάβετε τὴν κεφαλὴν ἀνδρῶν τῶν υἱῶν τοῦ κυρίου ὑμῶν, καὶ ἐνέγκατε πρὸς μὲ, ὡς ἡ ὥρα αὔριον ἐν Ἰεζράελ· καὶ οἱ υἱοὶ τοῦ βασιλέως ἦσαν ἑβδομήκοντα ἄνδρες, οὗτοι ἁδροὶ τῆς πόλεως ἐξέτρεφον αὐτούς.
\vs{7}Καὶ ἐγένετο ὡς ἦλθε τὸ βιβλίον πρὸς αὐτοὺς, καὶ ἔλαβον τοὺς υἱοὺς τοῦ βασιλέως, καὶ ἔσφαξαν αὐτοὺς ἑβδομήκοντα ἄνδρας· καὶ ἔθηκαν τὰς κεφαλὰς αὐτῶν ἐν καρτάλλοις, καὶ ἀπέστειλαν αὐτὰς πρὸς αὐτὸν εἰς Ἰεζράελ.
\vs{8}Καὶ ἦλθεν ὁ ἄγγελος καὶ ἀπήγγειλε, λέγων, ἤνεγκαν τὰς κεφαλὰς τῶν υἱῶν τοῦ βασιλέως· καὶ εἶπε, θέτε αὐτὰς βουνοὺς δύο παρὰ τὴν θύραν τῆς πύλης εἰς πρωΐ.
\vs{9}Καὶ ἐγένετο πρωΐ καὶ ἐξῆλθε καὶ ἔστη, καὶ εἶπε πρὸς πάντα τὸν λαόν, δίκαιοι ὑμεῖς· ἰδοὺ ἐγώ εἰμι συνεστράφην ἐπὶ τὸν κύριόν μου, καὶ ἀπέκτεινα αὐτόν· καὶ τίς ἐπάταξε πάντας τούτους;
\vs{10}Ἴδετε ἀφφὼ, ὅτι οὐ πεσεῖται ἀπὸ τοῦ ῥήματος Κυρίου εἰς τὴν γῆν οὗ ἐλάλησε Κύριος ἐπὶ τὸν οἶκον Ἀχαάβ· καὶ Κύριος ἐποίησεν ὅσα ἐλάλησεν ἐν χειρὶ δούλου αὐτοῦ Ἠλιού.
\vs{11}Καὶ ἐπάταξεν Ἰοὺ πάντας τοὺς καταλειφθέντας ἐν τῷ οἴκῳ Ἀχαὰβ ἐν Ἰεζράελ, καὶ πάντας τοὺς ἁδροὺς αὐτοῦ, καὶ τοὺς γνωστοὺς αὐτοῦ, καὶ τοὺς ἱερεῖς αὐτοῦ, ὥστε μὴ καταλιπεῖν αὐτοὺ κατάλειμμα.

\vs{12}Καὶ ἀνέστη καὶ ἐπορεύθῃ εἰς Σαμάρειαν, αὐτὸς ἐν βαιθακὰθ τῶν ποιμένων ἐν τῇ ὁδῷ.
\vs{13}Καὶ Ἰοὺ εὗρε τοὺς ἀδελφοὺς Ὀχοζίου βασιλέως Ἰούδα, καὶ εἶπε, τίνες ὑμεῖς; καὶ εἶπον, ἀδελφοὶ Ὀχοζίου ἡμεῖς, καὶ κατέβημεν εἰς εἰρήνην τῶν υἱῶν τοῦ βασιλέως, καὶ τῶν υἱῶν τῆς δυναστευούσης.
\vs{14}Καὶ εἶπε, συλλάβετε αὐτοὺς ζῶντας· καὶ ἔσφαξαν αὐτοὺς εἰς βαιθακὰθ τεσσαράκοντα καὶ δύο ἄνδρας· οὐ κατέλιπεν ἄνδρα ἐξ αὐτῶν.

\vs{15}Καὶ ἐπορεύθη ἐκεῖθεν καὶ εὗρε τὸν Ἰωναδὰβ υἱὸν Ῥηχὰβ εἰς ἀπαντὴν αὐτοῦ, καὶ εὐλόγησεν αὐτόν· καὶ εἶπε πρὸς αὐτὸν Ιοὺ, εἰ ἔστι καρδία σου μετὰ καρδίας μου εὐθεῖα καθὼς ἡ καρδία μου μετὰ τῆς καρδίας σου; καὶ εἶπεν Ἰωναδάβ ἔστι· καὶ εἶπεν Ἰοὺ, καὶ εἰ ἔστι, δὸς τὴν χεῖρά σου· καὶ ἔδωκε τὴν χεῖρα αὐτοῦ· καὶ ἀνεβίβασεν αὐτὸν πρὸς αὐτὸν ἐπὶ τὸ ἅρμα,
\vs{16}καὶ εἶπε πρὸς αὐτὸν, δεῦρο μετʼ ἐμοῦ, καὶ ἴδε ἐν τῷ ζηλῶσαί με τῷ Κυρίῳ· καὶ ἐπεκάθισε αὐτὸν ἐν τῷ ἅρματι αὐτοῦ.

\vs{17}Καὶ εἰσῆλθεν εἰς Σαμάρειαν· καὶ ἐπάταξε πάντας τοὺς καταλειφθέντας τοῦ Ἀχαὰβ ἐν Σαμαρείᾳ ἕως τοῦ ἀφανίσαι αὐτὸν κατὰ τὸ ῥῆμα Κυρίου, ὃ ἐλάλησε πρὸς Ἠλιού.
\vs{18}Καὶ συνήθροισεν Ἰοὺ πάντα τὸν λαὸν, καὶ εἶπε πρὸς αὐτοὺς, Ἀχαὰβ ἐδούλευσε τῷ Βάαλ ὀλίγα, Ἰοὺ δουλεύσει αὐτῷ πολλά.
\vs{19}Καὶ νῦν πάντες οἱ προφῆται τοῦ Βάαλ πάντας τοὺς δούλους αὐτοῦ καὶ τοὺς ἱερεῖς αὐτοῦ καλέσατε πρὸς μὲ, ἀνὴρ μὴ ἐπισκεπήτω, ὅτι θυσία μεγάλη μοι τῷ Βάαλ· πᾶς ὃς ἐὰν ἐπισκεπῇ, οὐ ζήσεται· καὶ Ἰοὺ ἐποίησεν ἐν πτερνισμῷ, ἵνʼ ἀπολέσῃ τοὺς δούλους τοῦ Βάαλ.

\vs{20}Καὶ εἶπεν Ἰοὺ, ἁγιάσατε ἱερείαν τῷ Βάαλ· καὶ ἐκήρυξν.
\vs{21}Καὶ ἀπέστειλεν Ἰοὺ ἐν παντὶ Ἰσραὴλ, λέγων, καὶ νῦν πάντες οἱ δοῦλοι, καὶ πάντες οἱ ἱερεῖς αὐτοῦ, καὶ πάντες οἱ προθῆται αὐτοῦ, μηδεὶς ἀπολιπέσθω, ὅτι θυσίαν μεγάλην ποιῶ· ὃς ἂν ἀπολειφθῇ, οὐ ζήσεται· καὶ ἦλθον πάντες οἱ δοῦλοι τοῦ Βάαλ, καὶ πάντες οἱ ἱερεῖς αὐτοῦ, καὶ πάντες οἱ προφῆται αὐτοῦ· οὐ κατελείφθη ἀνὴρ ὃς οὐ παρεγένετο· καὶ εἰσῆλθον εἰς τὸν οἶκον τοῦ Βάαλ· καὶ ἐπλήσθη ὁ οἶκος τοῦ Βάαλ στόμα εἰς στόμα.
\vs{22}Καὶ εἶπε τῷ ἐπὶ τοῦ οἴκου μεσθάαλ, ἐξάγαγε ἔνδυμα πᾶσι τοῖς δούλοις τοῦ Βάαλ· καὶ ἐξήνεγκεν αὐτοῖς ὁ στολιστής.
\vs{23}Καὶ εἰσῆλθεν Ἰοὺ καὶ Ἰωναδὰβ υἱὸς Ῥηχὰβ εἰς οἶκον τοῦ Βάαλ, καὶ εἶπε τοῖς δούλοις τοῦ Βάαλ, ἐρευνήσατε καὶ ἴδετε, εἰ ἔστι μεθʼ ὑμῶν τῶν δούλων Κυρίου, ὅτι ἀλλʼ ἢ οἱ δοῦλοι τοῦ Βάαλ μονώτατοι.
\vs{24}Καὶ εἰσῆλθε τοῦ ποιῆσαι τὰ θύματα καὶ τὰ ὁλοκαυτώματα· καὶ Ἰοὺ ἔταξεν ἑαυτῷ ἔξω ὀγδοήκοντα ἄνδρας, καὶ εἶπεν, ἀνὴρ ὃς ἐὰν διασωθῇ ἀπὸ τῶν ἀνδρῶν ὧν ἐγὼ ἀνάγω ἐπὶ χεῖρα ὑμῶν, ἡ ψυχὴ αὐτοῦ ἀντὶ τῆς ψυχῆς αὐτοῦ.

\vs{25}Καὶ ἐγένετο ὡς συνετέλεσε ποιῶν τὴν ὁλοκαύτωσιν, καὶ εἶπεν Ἰοὺ τοῖς παρατρέχουσιν καὶ τοῖς τριστάταις, εἰσελθόντες πατάξατε αὐτοὺς, μὴ ἐξελθάτω ἐξ αὐτῶν ἀνήρ· καὶ ἐπάταξαν αὐτοὺς ἐν στόματι ῥομφαίας, καὶ ἔῤῥιψαν οἱ παρατρέχοντες καὶ οἱ τριστάται, καὶ ἐπορεύθησαν ἕως πόλεως οἴκου τοῦ Βάαλ.
\vs{26}Καὶ ἐξήνεγκαν τὴν στήλην τοῦ Βάαλ, καὶ ἐνέπρησαν αὐτήν.
\vs{27}Καὶ κατέσπασαν τὰς στήλας τοῦ Βάαλ, καὶ ἔπαξαν αὐτὸν εἰς λυτρῶνας ἕως τῆς ἡμέρας ταύτης.
\vs{28}Καὶ ἠφάνισεν Ἰοὺ τὸν Βάαλ ἐξ Ἰσραήλ.

\vs{29}Πλὴν ἁμαρτιῶν Ἰεροβοὰμ υἱοῦ Ναβὰτ ὃς ἐξήμαρτε τὸν Ἰσραὴλ, οὐκ ἀπέστη Ἰοὺ ἀπὸ ὄπισθεν αὐτῶν· αἱ δαμάλεις αἱ χρυσαῖ ἐν Βαιθὴλ, καὶ ἐν Δάν.

\vs{30}Καὶ εἶπε Κύριος πρὸς Ἰοὺ, ἀνθʼ ὧν ὅσα ἠγάθυνας ποιῆσαι τὸ εὐθὲς ἐν ὀφθαλμοῖς μου κατὰ πάντα ὅσα ἐν τῇ καρδίᾳ μου ἐποίησας τῷ οἴκῳ Ἀχαὰβ, υἱοὶ τέταρτοι καθήσονταί σοι ἐπὶ θρόνου Ἰσραήλ.
\vs{31}Καὶ Ἰοὺ οὐκ ἐφύλαξε πορεύεσθαι ἐν νόμῳ Κυρίου θεοῦ Ἰσραὴλ ἐν ὅλῃ καρδίᾳ αὐτοῦ· οὐκ ἀπέστη ἑπάνωθεν ἁμαρτιῶν Ἰεροβοὰμ ὃς ἐξήμαρτε τὸν Ἰσραήλ.
\vs{32}Ἐν ταῖς ἡμέραις ἐκείναις ἤρξατο Κύριος συγκόπτειν ἐν τῷ Ἰσραήλ· καὶ ἐπάταξεν αὐτοὺς Ἀζαὴλ ἐν παντὶ ὁρίῳ Ἰσραὴλ,
\vs{33}ἀπὸ τοῦ Ἰορδάνου κατʼ ἀνατολὰς ἡλίου πᾶσαν τὴν Γαλαάδ τοῦ Γαδδὶ, καὶ τοῦ Ῥουβὴν, καὶ τοῦ Μανασσῆ, ἀπὸ Ἀροὴρ, ἥ ἐστιν ἐπὶ τοῦ χείλους χειμάῤῥου Ἀρνὼν, καὶ τὴν Γαλαὰδ καὶ τὴν Βασάν.

\vs{34}Καὶ τὰ λοιπὰ τῶν λόγων Ἰοὺ καὶ πάντα ὅσα ἐποίησε, καὶ πᾶσα ἡ δυναστεία αὐτοῦ, καὶ τὰς συνάψεις ἃς συνῆψεν, οὐχὶ ταῦτα γεγραμμένα ἐπὶ βιβλίου λόγων τῶν ἡμερῶν τοῖς βασιλεῦσιν Ἰσραήλ;
\vs{35}Καὶ ἐκοιμήθη Ἰοὺ μετὰ τῶν πατέρων αὐτοῦ, καὶ ἔθαψαν αὐτὸν ἐν Σαμαρείᾳ· καὶ ἐβασίλευσεν Ἰωάχαζ υἱὸς αὐτοῦ ἀντʼ αὐτοῦ.
\vs{36}Καὶ αἱ ἡμέραι ἃς ἐβασίλευσεν Ἰοὺ ἐπὶ Ἰσραὴλ εἰκοσιοκτὼ ἔτη ἐν Σαμαρείᾳ.

\ch{11}
Καὶ Γοθολία ἡ μήτηρ Ὀχοζίου εἶδεν ὅτι ἀπέθανεν ὁ υἱὸς αὐτῆς, καὶ ἀπώλεσε πᾶν τὸ σπέρμα τῆς βασιλείας.
\vs{2}Καὶ ἔλαβεν Ἰωσαβεὲ θυγάτηρ τοῦ βασιλέως Ἰωρὰμ ἀδελφὴ Ὀχοζίου τὸν Ἰωὰς υἱὸν ἀδελφοῦ αὐτῆς, καὶ ἔκλεψεν αὐτὸν ἐκ μέσου τῶν υἱῶν τοῦ βασιλέως τῶν θανατουμένων, αὐτὸν καὶ τὴν τροφὸν αὐτοῦ ἐν τῷ ταμείῳ τῶν κλινῶν, καὶ ἔκρυψεν αὐτὸν ἀπὸ προσώπου Γοθολίας, καὶ οὐκ ἐθανατώθη.
\vs{3}Καὶ ἦν μετʼ αὐτῆς κρυβόμενος ἐν οἴκῳ Κυρίου ἓξ ἔτη· καὶ Γοθολία βασιλεύουσα ἐπὶ τῆς γῆς.

\vs{4}Καὶ ἐν τῷ ἔτει τῷ ἑβδόμῳ ἀπέστειλεν Ἰωδαὲ, καὶ ἔλαβε τοὺς ἑκατοντάρχους τῶν Χοῤῥὶ καὶ τῶν Ῥασίμ, καὶ ἀπήγαγεν αὐτοὺς πρὸς αὐτὸν εἰς οἶκον Κυρίου, καὶ διέθετο αὐτοῖς διαθήκην Κυρίου, καὶ ὥρκωσε· καὶ ἔδειξεν αὐτοῖς Ἰωδαὲ τὸν υἱὸν τοῦ βασιλέως,
\vs{5}καὶ ἐνετείλατο αὐτοῖς, λέγων, οὗτος ὁ λόγος ὃν ποιήσετε· Τὸ τρίτον ἐξ ὑμῶν εἰσελθέτω τὸ σάββατον,
\vs{6}καὶ φυλάξατε φυλακὴν οἴκου τοῦ βασιλέως ἐν τῷ πυλῶνι, καὶ τὸ τρίτον ἐν τῇ πύλῃ τῶν ὁδῶν, καὶ τὸ τρίτον τῆς πύλης ὀπίσω τῶν παρατρεχόντων, καὶ φυλάξτε τὴν φυλακὴν τοῦ οἴκου.
\vs{7}Καὶ δύο χεῖρες ἐν ὑμῖν, πᾶς ὁ ἐκπορευόμενος τὸ σάββατον, καὶ φυλάξουσι τὴν φυλακὴν οἴκου Κυρίου πρὸς τὸν βασιλέα.
\vs{8}Καὶ κυκλώσατε ἐπὶ τὸν βασιλέα κύκλῳ, ἀνὴρ καὶ τὸ σκεῦος αὐτοῦ ἐν χειρὶ αὐτοῦ, καὶ ὁ εἰσπορευόμενος εἰς τὰς σαδηρὼθ, ἀποθανεῖται· καὶ ἔσονται μετὰ τοῦ βασιλέως ἐν τῷ ἐκπορεύεσθαι αὐτὸν καὶ ἐν τῷ εἰσπορεύεσθαι αὐτόν.

\vs{9}Καὶ ἐποίησαν οἱ ἑκατόνταρχοι πάντα ὅσα ἐνετείλατο Ἰωδαὲ ὁ συνετός· καὶ ἔλαβεν ἀνὴρ τοὺς ἄνδρας αὐτοῦ καὶ τοὺς εἰσπορευομένους τὸ σάββατον μετὰ τῶν ἐκπορευομένων τὸ σάββατον, καὶ εἰσῆλθον πρὸς Ἰωδαὲ τὸν ἱερέα.
\vs{10}Καὶ ἔδωκεν ὁ ἱερεὺς τοῖς ἑκατοντάρχοις τοὺς σειρομάστας καὶ τοὺς τρισσοὺς τοῦ βασιλέως Δαυὶδ τοὺς ἐν οἴκῳ Κυρίου.
\vs{11}Καὶ ἔστησαν οἱ παρατρέχοντες ἀνὴρ καὶ τὸ σκεῦος αὐτοῦ, ἐν τῇ χειρὶ αὐτοῦ ἀπὸ τῆς ὠμίας τοῦ οἴκου τῆς δεξιᾶς ἕως τῆς ὠμίας τοῦ οἴκου τῆς εὐωνύμου τοῦ θυσιαστηρίου καὶ τοῦ οἴκου, ἐπὶ τὸν βασιλέα κύκλῳ.
\vs{12}Καὶ ἐξαπέστειλε τὸν υἱὸν τοῦ βασιλέως, καὶ ἔδωκεν ἐπʼ αὐτὸν νεζὲρ καὶ τὸ μαρτύριον, καὶ ἐβασίλευσεν αὐτὸν καὶ ἔχρισεν αὐτόν· καὶ ἐκρότησαν τῇ χειρὶ, καὶ εἶπαν, ζήτω ὁ βασιλεύς.

\vs{13}Καὶ ἤκουσε Γοθολία τὴν φωνὴν τῶν τρεχόντων τοῦ λαοῦ, καὶ εἰσῆλθε πρὸς τὸν λαὸν εἰς οἶκον Κυρίου,
\vs{14}καὶ εἶδε, καὶ ἰδοὺ ὁ βασιλεὺς εἱστήκει ἐπὶ τοῦ στύλου κατὰ τὸ κρίμα· καὶ οἱ ᾠδοὶ καὶ αἱ σάλπιγγες πρὸς τὸν βασιλέα, καὶ πᾶς ὁ λαὸς τῆς γῆς χαίρων καὶ σαλπίζων ἐν σάλπιγξι· καὶ διέῤῥηξε Γοθολία τὰ ἱμάτια ἑαυτῆς, καὶ ἐβόησε σύνδεσμος, σύνδεσμος.
\vs{15}Καὶ ἐνετείλατο Ἰωδαὲ ὁ ἱερεὺς τοῖς ἑκατοντάρχοις τοῖς ἐπισκόποις τῆς δυνάμεως, καὶ εἶπε πρὸς αὐτούς, ἐξαγάγετε αὐτὴν ἔσωθεν τῶν σαδηρὼθ, ὁ εἰσπορευόμενος ὀπίσω αὐτῆς θανάτῳ θανατωθήσεται ἐν ῥομφαίᾳ· ὅτι εἶπεν ὁ ἱερεύς, καὶ μὴ ἀποθάνῃ ἐν οἴκῳ Κυρίου.
\vs{16}Καὶ ἐπέθηκαν αὑτῇ χεῖρας, καὶ εἰσῆλθον ὁδὸν εἰσόδου τῶν ἵππων οἴκου τοῦ βασιλέως, καὶ ἀπέθανεν ἐκεῖ.

\vs{17}Καὶ διέθετο Ἰωδαὲ διαθήκην ἀναμέσον Κυρίου καὶ ἀναμέσον τοῦ βασιλέως καὶ ἀναμέσον τοῦ λαοῦ, τοῦ εἶναι εἰς λαὸν τῷ Κυρίῳ· καὶ ἀναμέσον τοῦ βασιλέως καὶ ἀναμέσον τοῦ λαοῦ.
\vs{18}Καὶ εἰσῆλθε πᾶς ὁ λαὸς τῆς γῆς εἰς οἶκον τοῦ Βάαλ, καὶ κατέσπασαν αὐτὸν, καὶ τὰ θυσιαστήρια αὐτοῦ καὶ τὰς εἰκόνας αὐτοῦ συνέτριψαν ἀγαθῶς· καὶ τὸν Μαθὰν τὸν ἱερέα τοῦ Βάαλ ἀπέκτειναν κατὰ πρόσωπον τῶν θυσιαστηρίων· καὶ ἔθηκεν ὁ ἱερεὺς ἐπισκόπους εἰς τὸν οἶκον Κυρίου.
\vs{19}Καὶ ἔλαβε τοὺς ἑκατοντάρχους, καὶ τὸν Χοῤῥεὶ, καὶ τὸν Ῥασὶμ, καὶ πάντα τὸν λαὸν τῆς γῆς, καὶ κατήγαγον τὸν βασιλέα ἐξ οἴκου Κυρίου· καὶ εἰσῆλθον ὁδὸν πύλης τῶν παρατρεχόντων οἴκου τοῦ βασιλέως, καὶ ἐκάθισαν αὐτὸν ἐπὶ θρόνου τῶν βασιλέων.
\vs{20}Καὶ ἐχάρη πᾶς ὁ λαὸς τῆς γῆς, καὶ ἡ πόλις ἡσύχασε· καὶ τὴν Γοθολίαν ἐθανάτωσαν ἐν ῥομφαίᾳ ἐν οἴκῳ τοῦ βασιλέως.

\ch{12}
Υἱὸς ἐπτὰ ἐτῶν Ἰωὰς ἐν τῷ βασιλεύειν αὐτόν.

\vs{2}Ἐν ἔτει ἑβδόμῳ τῷ Ἰοὺ ἐβασίλευσεν Ἰωὰς, καὶ τεσσαράκοντα ἔτη ἐβασίλευσεν ἐν Ἱερουσαλὴμ, καὶ ὄνομα τῆς μητρὸς αὐτοῦ Σαβιὰ ἐκ τῆς Βηρσαβεέ.
\vs{3}Καὶ ἐποίησεν Ἰωὰς τὸ εὐθὲς ἐνώπιον Κυρίου πάσας τὰς ἡμέρας ἃς ἐφώτισεν αὐτὸν Ἰωδαὲ ὁ ἱερεύς.
\vs{4}Πλὴν τῶν ὑψηλῶν οὐ μετεστάθησαν, καὶ ἐκεῖ ἔτι ὁ λαὸς ἐθυσίαζε, καὶ ἐθυμίων ἐν τοῖς ὑψηλοῖς.

\vs{5}Καὶ εἶπεν Ἰωὰς πρὸς τοὺς ἱερεῖς, πᾶν τὸ ἀργύριον τῶν ἁγίων τὸ εἰσοδιαζόμενον ἐν τῷ οἴκῳ Κυρίου, ἀργύριον συντιμήσεως, ἀνὴρ ἀργύριον λαβὼν συντιμήσεως, πᾶν ἀργύριον ὃ ἐὰν ἀναβῇ ἐπὶ καρδίαν ἀνδρὸς ἐνεγκεῖν ἐν οἴκῳ Κυρίου,
\vs{6}λαβέτωσαν ἑαυτοῖς οἱ ἱερεῖς, ἀνὴρ ἀπὸ τῆς πράσεως αὐτοῦ, καὶ αὐτοὶ κρατήσουσι τὸ βεδὲκ τοῦ οἴκου εἰς πάντα οὗ ἐὰν εὑρεθῇ ἐκεῖ βεδέκ.

\vs{7}Καὶ ἐγενήθη ἐν τῷ εἰκοστῷ καὶ τρίτῳ ἔτει τῷ βασιλεῖ Ἰωὰς οὐκ ἐκραταίωσαν οἱ ἱερεῖς τὸ βεδέκ τοῦ οἴκου.
\vs{8}Καὶ ἐκάλεσεν Ἰωὰς ὁ βασιλεὺς Ἰωδαὲ τὸν ἱερέα καὶ τοὺς ἱερεῖς, καὶ εἶπε πρὸς αὐτούς, τί ὅτι οὐκ ἐκραταιοῦτε τὸ βεδὲκ τοῦ οἴκου; καὶ νῦν μὴ λάβητε ἀργύριον ἀπὸ τῶν πράσεων ὑμῶν, ὅτι εἰς τὸ βεδὲκ τοῦ οἴκου δώσετε αὐτό.
\vs{9}Καὶ συνεφώνησαν οἱ ἱερεῖς τοῦ μὴ λαβεῖν ἀργύριον παρὰ τοῦ λαοῦ, καὶ τοῦ μὴ ἐνισχῦσαι τὸ βεδὲκ τοῦ οἴκου.
\vs{10}Καὶ ἔλαβεν Ἰωδαὲ ὁ ἱερεὺς κιβωτὸν μίαν, καὶ ἔτρησε τρώγλην ἐπὶ τῇς σανίδος αὐτῆς, καὶ ἔδωκεν αὐτὴν παρὰ ἁμμαζειβὶ ἐν τῷ οἴκῳ ἀδρος οἴκου Κυρίου· και ἔδωκας οἱ ἱερεῖς οἱ ἱεπεῖς φυλάσσοντες τὸν σταθμὸν πᾶν τὸ ἀργύριον τὸ εὑρεθὲν ἐν οἴκῳ Κυρίου·

\vs{11}Καὶ ἐγένετο ὡς εἶδον ὅτι πολὺ τὸ ἀργύριον ἐν τῇ κιβωτῷ, καὶ ἀνέβη ὁ γραμματεὺς τοῦ βασιλέως καὶ ὁ ἱερεὺς ὁ μέγας, καὶ ἔσθιγξαν καὶ ἠρίθμησαν τὸ ἀργύριον τὸ εὑρεθὲν ἐν οἴκῳ Κυρίου.
\vs{12}Καὶ ἔδωκαν τὸ ἀργύριον τὸ ἑτοιμασθὲν ἐπὶ χεῖρας ποιούντων τὰ ἔργα τῶν ἐπισκοπῶν οἴκου· Κυρίου, καὶ ἐξέδοσαν τοῖς τέκτοσι τῶν ξύλων, καὶ τοῖς οἱκοδόμοις τοῖς ποιοῦσιν ἐν οἴκῳ Κυρίου,
\vs{13}καὶ τοῖς τειχισταῖς, καὶ τοῖς λατόμοις τῶν λίθων τοῦ κτήσασθαι ξύλα καὶ λίθους λατομητοὺς τοῦ κατασχεῖν τὸ βεδὲκ οἴκου Κυρίου, εἰς πάντα ὅσα ἐξωδιάσθη ἐπὶ τὸν οἶκον τοῦ κραταιῶσαι.
\vs{14}Πλὴν οὐ ποιηθήσονται οἴκῳ Κυρίου θύραι ἀργυραῖ, ἧλοι, φιάλαι, καὶ σάλπιγγες, πᾶν σκεῦος χρυσοῦν, καὶ σκεῦος ἀργυροῦν, ἐκ τοῦ ἀργυρίου τοῦ εἰσενεχθέντος ἐν οἴκῳ Κυρίου,
\vs{15}ὅτι τοῖς ποιοῦσι τὰ ἔργα δώσουσιν αὐτό· καὶ ἐκραταίωσαν ἐν αὐτῷ τὸν οἶκον Κυρίου.
\vs{16}Καὶ οὐκ ἐξελογίζοντο τοὺς ἄνδρας οἷς ἐδίδουν τὸ ἀργύριον ἐπὶ χεῖρας αὐτῶν δοῦναι τοῖς ποιοῦσι τὰ ἔργα, ὅτι ἐν πίστει αὐτῶν ποιοῦσιν.
\vs{17}Ἀργύριον περὶ ἁμαρτίας, καὶ ἀργύριον περὶ πλημμελείας, ὅ, τι εἰσηνέχθη ἐν οἴκῳ Κυρίου, τοῖς ἱερεῦσιν ἐγένετο.

\vs{18}Τότε ἀνέβη Ἀζαὴλ βασιλεὺς Συρίας, καὶ ἐπολέμησεν ἐπὶ Γὲθ, καὶ προκατελάβετο αὐτήν· καὶ ἔταξεν Ἀζαὴλ τὸ πρόσωπον αὐτοῦ ἀναβῆναι ἐπὶ Ἱερουσαλήμ.
\vs{19}Καὶ ἔλαβεν Ἰωὰς βασιλεὺς Ἰούδα πάντα τὰ ἅγια ὅσα ἡγίασεν Ἰωσαφὰτ καὶ Ἰωρὰμ καὶ Ὀχοζίας οἱ πατέρες αὐτοῦ καὶ βασιλεῖς Ἰούδα, καὶ τὰ ἅγια αὐτοῦ, καὶ πᾶν τὸ χρυσίον τὸ εὑρεθὲν ἐν θησαυροῖς οἴκου Κυρίου καὶ οἴκου τοῦ βασιλέως, καὶ ἀπέστειλεν τῷ Ἀζαὴλ βασιλεῖ Συρίας, καὶ ἀνέβη ἀπὸ Ἱερουσαλήμ.

\vs{20}Καὶ τὰ λοιπὰ τῶν λόγων Ἰωὰς καὶ πάντα ὅσα ἐποίησεν, οὐκ ἰδοὺ ταῦτα γεγραμμένα ἐπὶ βιβλίῳ λόγων τῶν ἡμερῶν τοῖς βασιλεῦσιν Ἰούδα;
\vs{21}Καὶ ἀνέστησαν οἱ δοῦλοι αὐτοῦ καὶ ἔδησαν πάντα σύνδεσμον, καὶ ἐπάταξαν τὸν Ἰωὰς ἐν οἴκῳ Μαλλὼ τῷ ἐν Σελά.
\vs{22}Καὶ Ἰεζιρχὰρ υἱὸς Ἱεμουὰθ, καὶ Ἰεζεβοὺθ ὁ υἱὸς αὐτοῦ Σωμὴρ, οἱ δοῦλοι αὐτοῦ, ἐπάταξαν αὐτὸν καὶ ἀπέθανε, καὶ ἔθαψαν αὐτὸν μετὰ τῶν πατέρων αὐτοῦ ἑν πόλει Δαυίδ· καὶ ἐβασίλευσεν Ἀμεσσείας υἱὸς αὐτοῦ ἀντʼ αὐτοῦ.

\ch{13}
Ἐν ἔτει εἰκοστῷ καὶ τρίτῳ ἔτει τῷ Ἰωὰς υἱῷ Ὀχοζίου βασιλεῖ Ἰούδα ἐβασίλευσεν Ἰωάχαζ υἱὸς Ἰοὺ ἐν Σαμαρείᾳ ἑπτακαίδεκα ἔτη.
\vs{2}Καὶ ἐποίησε τὸ πονηρὸν ἐν ὀφθαλμοῖς Κυρίου, καὶ ἐπορεύθη ὀπίσω ἁμαρτιῶν Ἱεροβοὰμ υἱοῦ Ναβὰτ, ὃς ἐξήμαρτε τὸν Ἰσραὴλ, οὐκ ἀπέστη ἀπʼ αὐτῆς.

\vs{3}Καὶ ὠργίσθη θυμῷ Κύριος ἐν τῷ Ἰσραὴλ, καὶ ἔδωκεν αὐτοὺς ἐν χειρὶ Ἀζαὴλ βασιλέως Συρίας, καὶ ἐν χειρὶ υἱοῦ Ἄδερ υἱοῦ Ἀζαὴλ πάσας τὰς ἡμέρας.
\vs{4}Καὶ ἐδεήθη Ἰωάχαζ τοῦ προσώπου Κυρίου, καὶ ἐπήκουσεν αὐτοῦ Κύριος, ὅτι εἶδε τὴν θλίψιν Ἰσραὴλ, ὅτι ἔθλιψεν αὐτοὺς βασιλεὺς Συρίας.
\vs{5}Καὶ ἔδωκε Κύριος σωτηρίαν τῷ Ἰσραὴλ, καὶ ἐξῆλθεν ὑποκάτωθεν χειρὸς Συρίας· καὶ ἐκάθισαν οἱ υἱοὶ Ἰσραὴλ ἐν τοῖς σκηνώμασιν αὐτῶν καθὼς ἐχθὲς καὶ τρίτης.
\vs{6}Πλὴν οὐκ ἀπέστησαν ἀπὸ ἁμαρτιῶν οἴκου Ἱεροβοὰμ ὃς ἐξήμαρτε τὸν Ἰσραὴλ, ἐν αὐτῇ ἐπορεύθη· καὶ γε τὸ ἄλσος ἐστάθη ἐν Σαμαρείᾳ.
\vs{7}Ὅτι οὐχ ὑπελείφθη τῷ Ἰωάχαζ λαὸς, ἀλλʼ ἢ πεντήκοντα ἱππεῖς καὶ δέκα ἅρματα καὶ δέκα χιλιάδες πεζῶν, ὅτι ἀπώλεσεν αὐτοὺς βασιλεὺ Συρίας, καὶ ἔθεντο αὐτοὺς ὡς χοῦν εἰς καταπάτησιν.

\vs{8}Καὶ τὰ λοιπὰ τῶν λόγων Ἰωάχαζ καὶ πάντα ὅσα ἐποίησε καὶ αἱ δυναστεῖαι αὐτοῦ, οὐχὶ ταῦτα γεγραμμένα ἐπὶ βιβλίῳ λόγων τῶν ἡμερῶν τοῖς βασιλεῦσιν Ἰσραήλ;
\vs{9}Καὶ ἐκοιμήθη Ἰωάχαζ μετὰ τῶν πατέρων αὐτοῦ, καὶ ἔθαψαν αὐτὸν ἐν Σαμαρείᾳ· καὶ ἐβασίλευσεν Ἰωὰς υἱὸς αὐτοῦ ἀντʼ αὐτοῦ.

\vs{10}Ἐν ἔτει τριακοστῷ καὶ ἑβδόμῳ ἔτει τῷ Ἰωὰς βασιλεῖ Ἰούδα ἐβασίλευσεν Ἰωὰς υἱὸς Ἰωάχαζ ἐπὶ Ἰσραὴλ ἐν Σαμαρείᾳ ἑκκαίδεκα ἔτη.
\vs{11}Καὶ ἐποίησε τὸ πονηρὸν ἐν ὀφθαλμοῖς Κυρίου· οὐκ ἀπέστη ἀπὸ πάσης Ἰεροβοὰμ υἱοῦ Ναβὰτ ἁμαρτίας, ὃς ἐξήμαρτε τὸν Ἰσραὴλ· ἐν αὐτῇ ἐπορεύθη.
\vs{12}Καὶ τὰ λοιπὰ τῶν λόγων Ἰωὰς καὶ πάντα ὅσα ἐποίησε, καὶ αἱ δυναστεῖαι αὐτοῦ ἃς ἐποίησε μετὰ Ἀμεσσίου βασιλέως Ἰούδα, οὐχὶ ταῦτα γεγραμμένα ἐπὶ βιβλίῳ λόγων τῶν ἡμερῶν τοῖς βασιλεῦσιν Ἰσραήλ;
\vs{13}Καὶ ἐκοιμήθη Ἰωὰς μετὰ τῶν πατέρων αὐτοῦ, καὶ Ἰεροβοὰμ ἐκάθισεν ἐτὶ τοῦ θρόνου αὐτοῦ, καὶ ἐτάφη ἐν Σαμαρείᾳ μετὰ τῶν βασιλέων Ἰσραήλ.

\vs{14}Καὶ Ἑλισαιὲ ἠῤῥώστησε τὴν ἀῤῥωστίαν αὐτοῦ, διʼ ἣν ἀπέθανε· καὶ κατέβη πρὸς αὐτὸν Ἰωὰς βασιλεὺς Ἰσραὴλ, καὶ ἔκλαυσεν ἐπὶ πρόσωπον αὐτοῦ, καὶ εἶπε, πάτερ πάτερ, ἅρμα Ἰσραὴλ καὶ ἱππεὺς αὐτοῦ.
\vs{15}Καὶ εἶπεν αὐτῷ Ἐλισαιὲ, λάβε τόξον καὶ βέλη· καὶ ἔλαβε πρὸς ἑαυτὸν τόξον καὶ βέλη.
\vs{16}Καὶ εἶπε τῷ βασιλεῖ, ἐπιβίβασον τὴν χεῖρά σου ἐπὶ τὸ τόξον· καὶ ἐπεβίβασεν Ἰωὰς τὴν χεῖρα αὐτοῦ· καὶ ἐπέθηκεν Ἑλισαιὲ τὰς χεῖρας αὐτοῦ ἐπὶ τὰς χεῖρας τοῦ βασιλέως,
\vs{17}καὶ εἶπεν, ἄνοιξον τὴν θυρίδα κατʼ ἀνατολάς· καὶ ἤνοιξε· καὶ εἶπεν Ἑλισαιὲ, τόζευσον· καὶ ἐτόζευσε· καὶ εἶπε, βέλος σωτηρίας τῷ Κυρίῳ, καὶ βέλος σωτηρίας ἐν Συρίᾳ, καὶ πατάξεις τὴν Συρίαν ἐν Ἀφὲκ ἕως συντελείας.
\vs{18}Καὶ εἶπεν αὐτῷ Ἑλισαιὲ, λάβε τόξα· καὶ ἔλαβε· καὶ εἶπε τῷ βασιλεῖ Ἰσραὴλ, πάταξον εἰς τὴν γῆν· καὶ ἐπάταξεν ὁ βασιλεὺς τρὶς, καὶ ἔστη·
\vs{19}Καὶ ἐλυπήθη ἐπʼ αὐτῷ ὁ ἄνθρωπος τοῦ Θεοῦ, καὶ εἶπεν, εἰ ἐπάταξας πεντάκις ἢ ἑξάκις, τότε ἂν ἐπάταξας τὴν Συρίαν ἕως συντελείας, καὶ νῦν τρὶς πατάξεις τὴν Συρίαν.

\vs{20}Καὶ ἀπέθανεν Ἑλισαιὲ, καὶ ἔθαψαν αὐτόν· καὶ μονόζωνοι Μωὰβ ἦλθον ἐν τῇ γῇ ἐλθόντος τοῦ ἐνιαυτοῦ.
\vs{21}Καὶ ἐγένετο αὐτῶν θαπτόντων τὸν ἄνδρα, καὶ ἰδοὺ εἶδον τὸν μονόζωνον, καὶ ἔῤῥιψαν τὸν ἄνδρα ἐν τῷ τάφῳ Ἑλισαιέ· καὶ ἐπορεύθη καὶ ἥψατο τῶν ὀστέων Ἑλισαιὲ, καὶ ἔζησε καὶ ἀνέστη ἐπὶ τοὺς πόδας αὐτοῦ.

\vs{22}Καὶ Ἀζαὴλ ἐξέθλιψε τὸν Ἰσραὴλ πάσας τὰς ἡμέρας Ἰωάχαζ.
\vs{23}Καὶ ἠλέησε Κύριος αὐτοὺς καὶ ᾠκτείρησεν αὐτούς, καὶ ἐπέβλεψεν ἐπʼ αὐτοὺς διὰ τὴν διαθήκην αὐτοῦ τὴν μετὰ Ἁβραὰμ καὶ Ἰσαὰκ καὶ Ἰακὼβ, καὶ οὐκ ἠθέλησε Κύριος διαφθεῖραι αὐτοὺς, καὶ οὐκ ἀπέῤῥιψεν αὐτοὺς ἀπὸ τοῦ προσώπου αὐτοῦ.
\vs{24}Καὶ ἀπέθανεν Ἀζαὴλ βασιλεὺς Συρίας, καὶ ἐβασίλευσεν υἱὸς Ἄδερ υἱὸς αὐτοῦ ἀντʼ αὐτοῦ.
\vs{25}Καὶ ἐπέστρεψεν Ἰωὰς υἱὸς Ἰωάχαζ, καὶ ἔλαβε τὰς πόλεις ἐκ χειρὸς υἱοῦ Ἄδερ υἱοῦ Ἄζαὴλ, ἃς ἔλαβεν ἐκ χειρὸς Ἰωάχας τοῦ πατρὸς αὐτοῦ ἐν τῷ πολέμῳ· τρὶς ἐπάταξεν αὐτὸν Ἰωὰς, καὶ ἐπέστρεψε τὰς πόλεις Ἰσραήλ.

\ch{14}
Ἐν ἔτει δευτέρῳ τῷ Ἰωὰς υἱῷ Ἰωάχαζ βασιλεῖ Ἰσραὴλ, καὶ ἐβασιλευσεν Ἀμεσσίας υἱὸς Ἰωὰς βασιλεὺς Ἰούδα.
\vs{2}Υἱὸς εἴκοσι καὶ πέντε ἐτῶν ἦν ἐν τῷ βασιλεύειν αὐτὸν, καὶ εἴκοσι καὶ ἐννέα ἔτη ἐβασίλευσεν ἐν Ἰερουσαλὴμ, καὶ ὄνομα τῆς μητρὸς αὐτοῦ Ἰωαδὶμ ἐξ Ἱερουσαλήμ.
\vs{3}Καὶ ἐποίησε τὸ εὐθὲς ἐν ὀφθαλμοῖς Κυρίου, πλὴν οὐχ ὡς Δαυὶδ ὁ πατὴρ αὐτοῦ· κατὰ πάντα ὅσα ἐποίησεν Ἰωὰς ὁ πατὴρ αὐτοῦ ἐποίησε.
\vs{4}Πλὴν τὰ ὑψηλὰ οὐκ ἐξῇρεν· ἔτι ὁ λαὸς ἐθυσίαζε καὶ ἐθυμίων ἐν τοῖς ὑψηλοῖς.
\vs{5}Καὶ ἐγένετο ὅτε κατίσχυσεν ἡ βασιλεία ἐν χειρὶ αὐτοῦ, καὶ ἐπάταξε τοὺς δούλους αὐτοῦ τοὺς πατάξαντας τὸν πατέρα αὐτοῦ.
\vs{6}Καὶ τοὺς υἱοὺς τῶν παταξάντων οὐκ ἐθανάτωσε, καθὼς γέγραπται ἐν βιβλίῳ νόμων Μωυσῆ, ὡς ἐνετείλατο Κύριος, λέγων, οὐκ ἀποθανοῦνται πατέρες ὑπὲρ υἱῶν, καὶ υἱοὶ οὐκ ἀποθανοῦνται ὑπὲρ πατέρων, ὅτι ἀλλʼ ἢ ἕκαστος ἐν ταῖς ἁμαρτίαις αὐτοῦ ἀποθανεῖται.
\vs{7}Αὐτὸς ἐπάταξε τὴν Ἐδὼμ ἐν γεμελὲδ δέκα χιλιάδας, καὶ συνέλαβε τὴν πέτραν ἐν τῷ πολέμῳ, καὶ ἐκάλεσε τὸ ὄνομα αὐτῆς Ἰεθοὴλ ἕως τῆς ἡμέρας ταύτης.

\vs{8}Τότε ἀπέστειλεν Ἀμεσσίας ἀγγέλους πρὸς Ἰωὰς υἱὸν Ἰωάχαζ υἱοῦ Ἰοὺ βασιλέως Ἰσραὴλ, λέγων, δεῦρο ὀφθῶμεν προσώποις.
\vs{9}Καὶ ἀπέστειλεν Ἰωὰς βασιλεὺς Ἰσραὴλ πρὸς Ἀμεσσίαν βασιλέα Ἰούδα, λέγων, ὁ ἄκαν ὁ ἐν τῷ Λιβάνῳ ἀπέστειλε πρὸς τὴν κέδρον τὴν ἐν τῷ Λιβάνῳ, λέγων δὸς τὴν θυγατέρα σου τῷ υἱῷ μου εἰς γυναῖκα· καὶ διῆλθον τὰ θηρία τοῦ ἀγροῦ τὰ ἐν τῷ Λιβάνῳ, καὶ συνεπάτησαν τὴν ἄκανα.
\vs{10}Τύπτων ἐπάταξας τὴν Ἰδουμαίαν, καὶ ἐπῇρέ σε καρδία σου· ἐνδοξάσθητι καθήμενος ἐν τῷ οἴκῳ σου, καὶ ἱνατί ἐρίζεις ἐν κακίᾳ σου; καὶ πεσῇ σὺ καὶ Ἰούδας μετὰ σοῦ.

\vs{11}Καὶ οὐκ ἤκουσεν Ἀμεσσίας· καὶ ἀνέβη Ἰωὰς βασιλεὺς Ἰσραὴλ, καὶ ὤφθησαν προσώποις αὐτὸς καὶ Ἀμεσσίας βασιλεὺς Ἰούδα ἐν Βαιθσαμὺς τῇ τοῦ Ἰούδα·
\vs{12}Καὶ ἔπταισεν Ἰούδας ἀπὸ προσώπου Ἰσραὴλ, καὶ ἔφυγεν ἀνὴρ εἰς τὸ σκήνωμα αὐτοῦ.
\vs{13}Καὶ τὸν Ἀμεσσίαν υἱὸν Ἰωὰς υἱοῦ Ὀχοζίου συνέλαβεν Ἰωὰς βασιλεὺς Ἰσραὴλ ἐν Βαιθσαμύς· καὶ ἦλθεν εἰς Ἱερουσαλήμ, καὶ καθεῖλεν ἐν τῷ τείχει Ἱερουσαλὴμ ἐν τῇ πύλῃ Ἐφραὶμ ἕως πύλης τῆς γωνίας τετρακοσίους πήχεις.
\vs{14}Καὶ ἔλαβε τὸ χρυσίον, καὶ τὸ ἀργύριον, καὶ πάντα τὰ σκεύη τὰ εὑρεθέντα ἐν οἴκῳ Κυρίου, καὶ ἐν θησαυροῖς οἴκου τοῦ βασιλέως, καὶ τοὺς υἱοὺς τῶν συμμίξεων, καὶ ἀπέστρεψεν εἰς Σαμάρειαν.

\vs{15}Καὶ τὰ λοιπὰ τῶν λόγων Ἰωὰς ὅσα ἐποίησεν ἐν δυναστείᾳ αὐτοῦ, ἃ ἐπολέμησε μετὰ Ἀμεσσείου βασιλέως Ἰούδα, οὐχὶ ταῦτα γεγραμμένα ἐπὶ βιβλίῳ λόγων τῶν ἡμερῶν τοῖς βασιλεῦσιν Ἰσραήλ;
\vs{16}Καὶ ἐκοιμήθη Ἰωὰς μετὰ τῶν πατέρων αὐτοῦ, καὶ ἐτάφη ἐν Σαμαρείᾳ μετὰ τῶν βασιλέων Ἰσραήλ· καὶ ἐβασίλευσεν Ἰεροβοὰμ υἱὸς αὐτοῦ ἀντʼ αὐτοῦ.

\vs{17}Καὶ ἔζησεν Ἀμεσσίας υἱὸς Ἰωὰς βασιλεὺς Ἰούδα μετὰ τὸ ἀποθανεῖν Ἰωὰς υἱὸν Ἰωάχαζ βασιλέα Ἰσραὴλ, πεντεκαίδεκα ἔτη.
\vs{18}Καὶ τὰ λοιπὰ τῶν λόγων Ἀμεσσίου καὶ πάντα ἅσα ἐποίησεν, οὐχὶ ταῦτα γεγραμμένα ἐπὶ βιβλίῳ λόγων τῶν ἡμερῶν τοῖς βασιλεῦσιν Ἰούδα;
\vs{19}Καὶ συνεστράφησαν ἐπʼ αὐτὸν σύστρεμμα ἐν Ἱερουσαλὴμ, καὶ ἔφυγεν εἰς Λαχίς· καὶ ἀπέστειλαν ὀπίσω αὐτοῦ εἰς Λαχὶς, καὶ ἐθανάτωσαν αὐτὸν ἐκεῖ.
\vs{20}Καὶ ῃραν αὐτὸν ἐφʼ ἵππων, καὶ ἐτάφη ἐν Ἱερουσαλὴμ μετὰ τῶν πατέρων αὐτοῦ ἐν πόλει Δαυίδ.

\vs{21}Καὶ ἔλαβε πᾶς ὁ λαὸς Ἰούδα τὸν Ἀζαρίαν, καὶ αὐτὸς υἱὸς ἑκκαίδεκα ἐτῶν, καὶ ἐβασίλευσαν αὐτὸς ἀντὶ τοῦ πατρὸς αὐτοῦ Ἀμεσσίου.
\vs{22}Αὐτὸς ᾠκοδόμησε τὴν Αἰλὼθ, καὶ ἐπέστρεψεν αὐτὴν τῷ Ἰούδα μετὰ τὸ κοιμηθῆναι τὸν βασιλέα μετὰ τῶν πατέρων αὐτοῦ.

\vs{23}Ἐν ἔτει πεντεκαιδεκάτῳ τοῦ Ἀμεσσίου υἱῷ Ἰωὰς βασιλεῖ Ἰούδα ἐβασίλευσεν Ἱεροβοὰμ υἱὸς Ἰωὰς ἐπὶ Ἰσραὴλ ἐν Σαμαρείᾳ τεσσαράκοντα καὶ ἓν ἔτος.
\vs{24}Καὶ ἐποίησε τὸ πονηρὸν ἐνώπιον Κυρίου· οὐκ ἀπέστη ἀπὸ πασῶν ἁμαρτιῶν Ἰεροβοὰμ υἱοῦ Ναβὰτ, ὃς ἐξήμαρτε τὸν Ἰσραήλ.
\vs{25}Αὐτὸς ἀπέστησε τὸ ὅριον Ἰσραὴλ ἀπὸ εἰσόδου Αἰμὰθ ἕως τῆς θαλάσσης τῆς Ἄραβα, κατὰ τὸ ῥῆμα Κυρίου Θεοῦ Ἰσραὴλ ὃ ἐλάλησεν ἐν χειρὶ δούλου αὐτοῦ Ἰῶνα υἱοῦ Ἀμαθὶ τοῦ προφήτου τοῦ ἐκ Γεθχοφέρ.
\vs{26}Ὅτι εἶδε Κύριος τὴν ταπείνωσιν Ἰσραὴλ πικρὰν σφόδρα, καὶ ὀλιγοστοὺς συνεχομένους, καὶ ἐσπανισμένους, καὶ ἐγκαταλελειμμένους, καὶ οὐκ ἦν ὁ βοηθῶν τῷ Ἰσραήλ.
\vs{27}Καὶ οὐκ ἐλάλησε Κύριος ἐξαλεῖψαι τὸ σπέρμα Ἰσραὴλ ὑποκάτωθεν τοῦ οὐρανοῦ· καὶ ἔσωσεν αὐτοὺς διὰ χειρὸς Ἱεροβοὰμ υἱοῦ Ἰωάς.

\vs{28}Καὶ τὰ λοιπὰ τῶν λόγων Ἱεροβοὰμ καὶ πάντα ὅσα ἐποίησε, καὶ αἱ δυναστεῖαι αὐτοῦ, ὅσα ἐπολέμησε, καὶ ὅσα ἐπέστρεψε τὴν Δαμασκὸν, καὶ τὴν Αἰμὰθ τῷ Ἰούδα ἐν Ἰσραὴλ, οὐχὶ ταῦτα γεγραμμένα ἐπὶ βιβλίῳ λόγων τῶν ἡμερῶν τοῖς βασιλεῦσιν Ἰσραήλ;
\vs{29}Καὶ ἐκοιμήθη Ἱεροβοὰμ μετὰ τῶν πατέρων αὐτοῦ μετὰ βασιλέων Ἰσραὴλ, καὶ ἐβασίλευσε Ζαχαρίας υἱὸς αὐτοῦ ἀντʼ αὐτοῦ.

\ch{15}
Ἐν ἔτει εἰκοστῷ καὶ ἑβδόμῳ τῷ Ἱεροβοὰμ βασιλεῖ Ἰσραὴλ ἐβασίλευσεν Ἀζαρίας υἱὸς Ἀμεσσίου βασιλέως Ἰούδα.
\vs{2}Υἱὸς ἑκκαίδεκα ἐτῶν ἦν ἐν τῷ βασιλεύειν αὐτὸν, καὶ πεντηκονταδύο ἔτη ἐβασίλευσεν ἐν Ἱερουσαλὴμ, καὶ ὄνομα τῇ μητρὶ αὐτοῦ Ἱεχελία ἐξ Ἱερουσαλήμ.
\vs{3}Καὶ ἐποίησε τὸ εὐθὲς ἐν ὀφθαλμοῖς Κυρίου, κατὰ πάντα ὅσα ἐποίησεν Ἀμεσσίας ὁ πατὴρ αὐτοῦ.
\vs{4}Πλὴν τῶν ὑψηλῶν οὐκ ἐξῇρεν, ἔτι ὁ λαὸς ἐθυσίαζε καὶ ἐθυμίων ἐν τοῖς ὑψηλοῖς.

\vs{5}Καὶ ἥψατο Κύριος τὸν βασιλέα, καὶ ἦν λελεπρωμένος ἕως ἡμέρας θανάτου αὐτοῦ· καὶ ἐβασίλευσεν ἐν οἴκῳ ἀφφουσώθ· καὶ Ἰωάθαμ υἱὸς τοῦ βασιλέως ἐπὶ τῷ οἴκῳ κρίνων τὸν λαὸν τῆς γῆς.

\vs{6}Καὶ τὰ λοιπὰ τῶν λόγων Ἀζαρίου καὶ πάντα ὅσα ἐποίησεν, οὐχὶ ταῦτα γεγραμμένα ἐπὶ βιβλίῳ λόγων τῶν ἡμερῶν τοῖς βασιλεῦσιν Ἰούδα;
\vs{7}Καὶ ἐκοιμήθη Ἀζαρίας μετὰ τῶν πατέρων αὐτοῦ, καὶ ἔθαψαν αὐτὸν μετὰ τῶν πατέρων αὐτοῦ ἐν πόλει Δαυὶδ, καὶ ἐβασίλευσεν Ἰωάθαν υἱὸς αὐτοῦ ἀντʼ αὐτοῦ.

\vs{8}Ἐν ἔτει τριακοστῷ καὶ ὀγδόῳ τῷ Ἀζαρίου βασιλεῖ Ἰούδα ἐβασίλευσε Ζαχαρίας υἱὸς Ἰεροβοὰμ ἐπὶ Ἰσραὴλ ἐν Σαμαρείᾳ ἑξάμηνον.
\vs{9}Καὶ ἐποίησε τὸ πονηρὸν ἐν ὀφθαλμοῖς Κυρίου καθὰ ἐποίησαν οἱ πατέρες αὐτοῦ· οὐκ ἀπέστη ἀπὸ πασῶν τῶν ἁμαρτιῶν Ἰεροβοὰμ υἱοῦ Ναβὰτ, ὃς ἐξήμαρτε τὸν Ἰσραήλ.
\vs{10}Καὶ συνεστράφησαν ἐπʼ αὐτὸν Σελλοὺμ υἱὸς Ἰαβίς· καὶ ἐπάταξαν αὐτὸν Κεβλαὰμ καὶ ἐθανάτωσαν αὐτὸν, καὶ ἐβασίλευσεν ἀντʼ αὐτοῦ.
\vs{11}Καὶ τὰ λοιπὰ τῶν λόγων Ζαχαρίου, ἰδού εἰσι γεγραμμένα ἐπὶ βιβλίῳ λόγων τῶν ἡμερῶν τοῖς βασιλεῦσιν Ἰσραήλ.
\vs{12}Ὁ λόγος Κυρίου ὃν ἐλάλησε πρὸς Ἰοὺ, λέγων, υἱοὶ τέταρτοι καθήσονταί σοι ἐπὶ θρόνου Ἰσραήλ· καὶ ἐγένετο οὕτως.

\vs{13}Καὶ Σελλοὺμ υἱὸς Ἰαβὶς ἐβασίλευσε· καὶ ἐν ἔτει τριακοστῷ καὶ ἐννάτῳ Ἀζαρίᾳ βασιλεῖ Ἰούδα ἐβασίλευσε Σελλοὺμ μῆνα ἡμερῶν ἐν Σαμαρείᾳ.
\vs{14}Καὶ ἀνέβη Μαναὴμ υἱὸς Γαδδὶ ἐκ Θαρσιλὰ, καὶ ἦλθεν εἰς Σαμάρειαν, καὶ ἐπάταξε τὸν Σελλοὺμ υἱὸν Ἰαβὶς ἐν Σαμαρείᾳ, καὶ ἐθανάτωσεν αὐτόν.
\vs{15}Καὶ τὰ λοιπὰ τῶν λόγων Σελλοὺμ, καὶ ἡ συστροφὴ αὐτοῦ ᾗ συνεστράφη, ἰδού εἰσι γεγραμμένα ἐπὶ βιβλίῳ λόγων τῶν ἡμερῶν τοῖς βασιλεῦσιν Ἰσραήλ.

\vs{16}Τότε ἐπάταξε Μαναὴμ καὶ τὴν Θερσὰ καὶ πάντα τὰ ἐν αὐτῇ, καὶ τὰ ὅρια αὐτῆς ἀπὸ Θερσὰ, ὅτι οὐκ ἤνοιξαν αὐτῷ, καὶ ἐπάταξεν αὐτὴν, καὶ τὰς ἐν γαστρὶ ἐχούσας ἀνέῤῥηξεν.

\vs{17}Ἐν ἔτει τριακοστῷ καὶ ἐννάτῳ τῷ Ἀζαρίᾳ βασιλεῖ Ἰούδα ἐβασίλευσε Μαναὴμ υἱὸς Γαδδὶ ἐπὶ Ἰσραὴλ ἐν Σαμαρείᾳ δέκα ἔτη.
\vs{18}Καὶ ἐποίησε τὸ πονηρὸν ἐν ὀφθαλμοῖς Κυρίου, οὐκ ἀπέστη ἀπὸ πασῶν ἁμαρτιῶν Ἱεροβοὰμ υἱοῦ Ναβὰτ ὃς ἐξήμαρτε τὸν Ἰσραήλ.
\vs{19}Ἐν ταῖς ἡμέραις αὐτοῦ ἀνέβη Φονὰ βασιλεὺς Ἀσσυρίων ἐπὶ τὴν γῆν, καὶ Μαναὴμ ἔδωκε τῷ Φουὰ χίλια τάλαντα ἀργυρίου εἶναι τὴν χεῖρα αὐτοῦ μετʼ αὐτοῦ.
\vs{20}Καὶ ἐξήνεγκε Μαναὴμ τὸ ἀργύριον ἐπὶ τὸν Ἰσραὴλ ἐπὶ πᾶν δυνατὸν ἰσχύϊ, δοῦναι τῷ βασιλεῖ τῶν Ἀσσυρίων, πεντήκοντα σίκλους τῷ ἀνδρὶ τῷ ἑνί· καὶ ἀπέστρεψε βασιλεὺς Ἀσσυρίων, καὶ οὐκ ἔστη ἐκεῖ ἐν τῇ γῇ.
\vs{21}Καὶ τὰ λοιπὰ τῶν λόγων Μαναὴμ καὶ πάντα ὅσα ἐποίησεν, οὐκ ἰδοὺ ταῦτα γεγραμμένα ἐπὶ βιβλίῳ λόγων τῶν ἡμερῶν τοῖς βασιλεῦσιν Ἰσραήλ;
\vs{22}Καὶ ἐκοιμήθη Μαναὴμ μετὰ τῶν πατέρων αὐτοῦ, καὶ ἐβασίλευσε Φακεσίας υἱὸς αὐτοῦ ἀντʼ αὐτοῦ.

\vs{23}Ἐν ἔτει πεντηκοστῷ τοῦ Ἀζαρίου βασιλεῖ Ἰούδα ἐβασίλευσε Φακεσίας υἱὸς Μαναὴμ ἐπὶ Ἰσραὴλ ἐν Σαμαρείᾳ δύο ἔτη.
\vs{24}Καὶ ἐποίησε τὸ πονηρὸν ἐν ὀφθαλμοῖς Κυρίου, οὐκ ἀπέστη ἀπὸ ἁμαρτιῶν Ἰεροβοὰμ υἱοῦ Ναβὰτ ὃς ἐξήμαρτε τὸν Ἰσραήλ.
\vs{25}Καὶ συνεστράφη ἐπʼ αὐτὸν Φακεὲ υἱὸς Ῥομελίου ὁ τριστάτης αὐτοῦ, καὶ ἐπάταξεν αὐτὸν ἐν Σαμαρείᾳ ἐναντίον οἴκου τοῦ βασιλέως μετὰ τοῦ Ἀργὸβ καὶ μετὰ τοῦ Ἀρία, καὶ μετʼ αὐτοῦ πεντήκοντα ἄνδρες ἀπὸ τῶν τετρακοσίων· καὶ ἐθανάτωσε αὐτὸν, καὶ ἐβασίλευσεν ἀντʼ αὐτοῦ.
\vs{26}Καὶ τὰ λοιπὰ τῶν λόγων Φακεσίου καὶ πάντα ὅσα ἐποίησεν, ἰδού εἰσι γεγραμμένα ἐπὶ βιβλίῳ λόγων τῶν ἡμερῶν τοῖς βασιλεῦσιν Ἰσραήλ.

\vs{27}Ἐν ἔτει πεντηκοστῷ καὶ δευτέρῳ τοῦ Ἀζαρίου βασιλεῖ Ἰούδα ἐβασίλευσε Φακεὲ υἱὸς Ῥομελίου ἐπὶ Ἰσραὴλ ἐν Σαμαρείᾳ εἴκοσι ἔτη.
\vs{28}Καὶ ἐποίησε τὸ πονηρὸν ἐν ὀφθαλμοῖς Κυρίου, οὐκ ἀπέστη ἀπὸ πασῶν ἁμαρτιῶν Ἰεροβοὰμ υἱοῦ Ναβὰτ ὃς ἐξήμαρτε τὸν Ἰσραήλ.
\vs{29}Ἐν ταῖς ἡμέραις Φακεὲ βασιλέως Ἰσραὴλ ἦλθε Θαλγαθφελλασὰρ βασιλεὺς Ἀσσυρίων, καὶ ἔλαβε τὴν Αἲν καὶ τὴν Ἀβὲλ καὶ τὴν Θαμααχὰ καὶ τὴν Ἀνιὼχ καὶ τὴν Κενὲζ καὶ τὴν Ἀσὼρ καὶ τὴν Γαλαὰν καὶ τὴν Γαλιλαίαν, πᾶσαν γῆν Νεφθαλὶ, καὶ ἀπῳκισεν αὐτοὺς εἰς Ἀσσυρίους.
\vs{30}Καὶ συνέστρεψε σύστρεμμα Ὠσηὲ υἱὸς Ἠλὰ ἐπὶ Φακεὲ υἱὸν Ῥομελίου, καὶ ἐπάταξεν αὐτὸν, καὶ ἐθανάτωσε, καὶ ἐβασίλευσεν ἀντʼ αὐτοῦ, ἐν ἔτει εἰκοστῷ Ἰωάθαμ υἱοῦ Ἀζαρίου.
\vs{31}Καὶ τὰ λοιπὰ τῶν λόγων Φακεὲ καὶ πάντα ὅσα ἐποίησεν, ἰδοὺ ταῦτα γεγραμμένα ἐπὶ βιβλίῳ λόγων τῶν ἡμερῶν τοῖς βασιλεῦσιν Ἰσραήλ.

\vs{32}Ἐν ἔτει δευτέρῳ Φακεὲ υἱοῦ Ῥομελίου βασιλεῖ Ἰσραὴλ ἐβασίλευσεν Ἰωάθαμ υἱὸς Ἀζαρίου βασιλέως Ἰούδα.
\vs{33}Υἱὸς εἴκοσι καὶ πέντε ἐτῶν ἦν ἐν τῷ βασιλεύειν αὐτὸν, καὶ ἐκκαίδεκα ἔτη ἐβασίλευσεν ἐν Ἱερουσαλήμ, καὶ ὄνομα τῆς μητρὸς αὐτοῦ Ἰερουσὰ θυγάτηρ Σαδώκ.
\vs{34}Καὶ ἐποίησε τὸ εὐθὲς ἐν ὀφθαλμοῖς Κυρίου, κατὰ πάντα ὅσα ἐποίησεν Ἀζαρίας ὁ πατὴρ αὐτοῦ.
\vs{35}Πλὴν τὰ ὑψηλὰ οὐκ ἐξῇρεν, ἔτι ὁ λαὸς ἐθυσίαζε καὶ ἐθυμία ἐν τοῖς ὑψηλοῖς· αὐτὸς ᾠκοδόμησε τὴν πύλην οἴκου Κυρίου τὴν ἐπάνω.
\vs{36}Καὶ τὰ λοιπὰ τῶν λόγων Ἰωάθαμ καὶ πάντα ὅσα ἐποίησεν, οὐχὶ ταῦτα γεγραμμένα ἐπὶ βιβλίῳ λόγων τῶν ἡμερῶν τοῖς βασιλεῦσιν Ἰούδα;

\vs{37}Ἐν ταῖς ἡμέραις ἐκείναις ἤρξατο Κύριος ἐξαποστέλλειν ἐν Ἰούδᾳ τὸν Ῥαασσὼν βασιλέα Συρίας, καὶ τὸν Φακεὲ υἱὸν Ῥομελίου.
\vs{38}Καὶ ἐκοιμήθη Ἰωάθαμ μετὰ τῶν πατέρων αὐτοῦ, καὶ ἐτάφη μετὰ τῶν πατέρων αὐτοῦ ἐν πόλει Δαυὶδ τοῦ πατρὸς αὐτοῦ· καὶ ἐβασίλευσεν Ἄχαζ υἱὸς αὐτοῦ ἀντʼ αὐτοῦ.

\ch{16}
Ἐν ἔτει ἑπτακαιδεκάτῳ Φακεὲ υἱοῦ Ῥομελίου ἐβασίλευσεν Ἄχαζ υἱὸς Ἰωάθαμ βασιλέως Ἰούδα.
\vs{2}Υἱὸς εἴκοσι ἐτῶν ἦν Ἄχαζ ἐν τῷ βασιλεύειν αὐτὸν, καὶ ἑκκαίδεκα ἔτη ἐβασίλευσεν ἐν Ἱερουσαλήμ· καὶ οὐκ ἐποίησε τὸ εὐθὲς ἐν ὀφθαλμοῖς Κυρίου Θεοῦ αὐτοῦ πιστῶς, ὡς Δαυὶδ ὁ πατὴρ αὐτοῦ.
\vs{3}Καὶ ἐπορεύθη ἐν ὁδῷ βασιλέων Ἰσραὴλ, καί γε τὸν υἱὸν αὐτοῦ διῆγεν ἐν πυρὶ, κατὰ τὰ βδελύγματα τῶν ἐθνῶν ὧν ἐξῇρε Κύριος ἀπὸ προσώπου τῶν υἱῶν Ἰσραήλ.
\vs{4}Καὶ ἐθυσίαζε καὶ ἐθυμία ἐν τοῖς ὑψηλοῖς, καὶ ἐπὶ τῶν βουνῶν, καὶ ὑποκάτω παντὸς ξύλου ἀλσώδους.

\vs{5}Τότε ἀνέβη Ῥαασσὼν βασιλεὺς Συρίας καὶ Φακεὲ υἱὸς Ῥομελίου βασιλεὺς Ἰσραὴλ εἰς Ἱερουσαλὴμ εἰς πόλεμον, καὶ ἐπολιόρκουν ἐπὶ Ἄχαζ, καὶ οὐκ ἠδύναντο πολεμεῖν.
\vs{6}Ἐν τῷ καιρῷ ἐκείνῳ ἐπέστρεψε Ῥαασσὼν βασιλεὺς Συρίας τὴν Αἰλὰθ τῇ Συρίᾳ, καὶ ἐξέβαλε τοὺς Ἰουδαίους ἐξ Αἰλὰθ, καὶ Ἰδουμαῖοι ἦλθον εἰς Αἰλὰθ, καὶ κατῴκησαν ἐκεῖ ἕως τῆς ἡμέρας ταύτης.
\vs{7}Καὶ ἀπέστειλεν Ἄχαζ ἀγγέλους πρὸς Θαλγαθφελλασὰρ βασιλέα Ἀσσυρίων, λέγων, δοῦλός σου καὶ υἱός σου ἐγὼ, ἀνάβηθι, σῶσόν με ἐκ χειρὸς βασιλέως Συρίας, καὶ ἐκ χειρὸς βασιλέως Ἰσραὴλ, τῶν ἐπανισταμένων ἐπʼ ἐμέ.
\vs{8}Καὶ ἔλαβεν Ἄχαζ ἀργύριον καὶ χρυσίον τὸ εὑρεθὲν ἐν θησαυροῖς οἴκου Κυρίου καὶ οἴκου τοῦ βασιλέως, καὶ ἀπέστειλε τῷ βασιλεῖ δῶρα.
\vs{9}Καὶ ἤκουσεν αὐτοῦ βασιλεὺς Ἀσσυρίων· καὶ ἀνέβη βασιλεὺς Ἀσσυρίων εἰς Δαμασκὸν, καὶ συνέλαβεν αὐτὴν, καὶ ἀπῴκισεν αὐτὴν, καὶ τὸν Ῥαασσὼν βασιλέα ἐθανάτωσε.

\vs{10}Καὶ ἐπορεύθη βασιλεὺς Ἄχαζ εἰς Δαμασκὸν εἰς ἀπαντὴν Θαλγαθφελλασὰρ βασιλεῖ Ἀσσυρίων εἰς Δαμασκόν· καὶ εἶδε τὸ θυσιαστήριον ἐν Δαμασκῷ· καὶ ἀπέστειλεν ὁ βασιλεὺς Ἄχαζ πρὸς Οὐρίαν τὸν ἱερέα τὸ ὁμοίωμα τοῦ θυσιαστηρίου καὶ τὸν ῥυθμὸν αὐτοῦ καὶ πᾶσαν ποίησιν αὐτοῦ.
\vs{11}Καὶ ᾠκοδόμησεν Οὐρίας ὁ ἱερεὺς τὸ θυσιαστήριον, κατὰ πάντα ὅσα ἀπέστειλεν ὁ βασιλεὺς Ἄχαζ ἐκ Δαμασκοῦ.

\vs{12}Καὶ εἶδεν ὁ βασιλεὺς τὸ θυσιαστήριον, καὶ ἀνέβη ἐπʼ αὐτὸ,
\vs{13}καὶ ἐθυμίασε τὴν ὁλοκαύτωσιν αὐτοῦ, καὶ τὴν θυσίαν αὐτοῦ, καὶ τὴν σπονδὴν αὐτοῦ, καὶ προσέχεε τὸ αἷμα τῶν εἰρηνικῶν τῶν αὐτοῦ ἐπὶ τὸ θυσιαστήριον τὸ χαλκοῦν
\vs{14}τὸ ἀπέναντι Κυρίου· καὶ προσήγαγε τὸ πρόσωπον τοῦ οἴκου Κυρίου, ἀπὸ τοῦ ἀναμέσον τοῦ θυσιαστηρίου καὶ ἀπὸ τοῦ ἀναμέσον τοῦ οἴκου Κυρίου· καὶ ἔδειξεν αὐτὸ ἐπὶ μηρὸν τοῦ θυσιαστηρίου κατὰ βοῤῥᾶν.
\vs{15}Καὶ ἐνετείλατο ὁ βασιλεὺς Ἄχαζ τῷ Οὐρίᾳ τῷ ἱερεῖ, λέγων, ἐπὶ τὸ θυσιαστήριον τὸ μέγα πρόσφερε τὴν ὁλοκαύτωσιν τὴν πρωϊνὴν καὶ τὴν θυσίαν τὴν ἑσπερινήν, καὶ τὴν ὁλοκαύτωσιν τοῦ βασιλέως καὶ τὴν θυσίαν αὐτοῦ, καὶ τὴν ὁλοκαύτωσιν παντὸς τοῦ λαοῦ, καὶ τὴν θυσίαν αὐτῶν καὶ τὴν σπονδὴν αὐτῶν, καὶ πᾶν αἷμα ὁλοκαυτώσεως, καὶ πᾶν αἷμα θυσίας ἐπʼ αὐτῶ ἐκχεεῖς· καὶ τὸ θυσιαστήριον τὸ χαλκοῦν ἔσται μοι εἰς τοπρωΐ.
\vs{16}Καὶ ἐποίησεν Οὐρίας ὁ ἱερεὺς κατὰ πάντα ὅσα ἐνετείλατο αὐτῷ ὁ βασιλεὺς Ἄχαζ.
\vs{17}Καὶ συνέκοψεν ὁ βασιλεὺς Ἄχαζ τὰ συνκλείσματα τῶν μεχωνὼθ, καὶ μετῇρεν ἀπʼ αὐτῶν τὸν λουτῆρα, καὶ τὴν θάλασσαν καθεῖλεν ἀπὸ τῶν βοῶν τῶν χαλκῶν τῶν ὑποκάτω αὐτῆς, καὶ ἔδωκεν αὐτὴν ἐπὶ βάσιν λιθίνην.
\vs{18}Καὶ τὸν θεμέλιον τῆς καθέδρας ᾠκοδόμησεν ἐν οἴκῳ Κυρίου, καὶ τὴν εἴσοδον τοῦ βασιλέως τὴν ἔξω ἐπέστρεψεν ἐν οἴκῳ Κυρίου ἀπὸ προσώπου βασιλέως Ἀσσυρίων.

\vs{19}Καὶ τὰ λοιπὰ τῶν λόγων Ἄχαζ ὅσα ἐποίησεν, οὐχὶ ταῦτα γεγραμμένα ἐπὶ βιβλίῳ λόγων τῶν ἡμερῶν τοῖς βασιλεῦσιν Ἰούδα;
\vs{20}Καὶ ἐκοιμήθη Ἄχαζ μετὰ τῶν πατέρων αὐτοῦ, καὶ ἐτάφη ἐν πόλει Δαυὶδ, καὶ ἐβασίλευσεν Ἐζεκίας υἱὸς αὐτοῦ ἀντʼ αὐτοῦ.

\ch{17}
Ἐν ἔτει δωδεκάτῳ τοῦ Ἄχαζ βασιλέως Ἰούδα ἐβασίλευσεν Ὠσηὲ υἱὸς Ἠλὰ ἐν Σαμαρείᾳ ἐπὶ Ἰσραὴλ ἐννέα ἔτη.
\vs{2}Καὶ ἐποίησε τὸ πονηρὸν ἐν ὀφθαλμοῖς Κυρίου, πλὴν οὐχ ὡς οἱ βασιλεῖς Ἰσραὴλ οἳ ἦσαν ἔμπροσθεν αὐτοῦ.

\vs{3}Ἐπʼ αὐτὸν ἀνέβη Σαλαμανασσὰρ βασιλεὺς Ἀσσυρίων· καὶ ἐγενήθη αὐτῷ Ὠσηὲ δοῦλος, καὶ ἐπέστρεψεν αὐτῷ μαναά.
\vs{4}Καὶ εὗρε βασιλεὺς Ἀσσυρίων ἐν τῷ Ὠσηὲ ἀδικίαν, ὅτι ἀπέστειλεν ἀγγέλους πρὸς Σηγὼρ βασιλέα Αἰγύπτου, καὶ οὐκ ἤνεγκε μαναὰ τῷ βασιλεῖ Ἀσσυρίων ἐν τῷ ἐνιαυτῷ ἐκείνῳ· καὶ ἐπολιόρκησεν αὐτὸν ὁ βασιλεὺς Ἀσσυρίων, καὶ ἔδησεν αὐτὸν ἐν οἴκῳ φυλακῆς.
\vs{5}Καὶ ἀνέβη ὁ βασιλεὺς Ἀσσυρίων ἐν πάσῃ τῇ γῇ, καὶ ἀνέβη εἰς Σαμάρειαν, καὶ ἐπολιόρκησεν ἐπʼ αὐτὴν τρία ἔτη.

\vs{6}Ἐν ἔτει ἐννάτῳ Ὠσῆε συνέλαβε βασιλεὺς Ἀσσυρίων τὴν Σαμάρειαν, καὶ ἀπῴκισεν Ἰσραὴλ εἰς Ἀσσυρίους, καὶ κατῴκισεν αὐτοὺς ἐν Ἀλαὲ καὶ ἐν Ἀβὼρ ποταμοῖς Γωζὰν, καὶ ὄρη Μήδων.
\vs{7}Καὶ ἐγένετο ὅτι ἥμαρτον οἱ υἱοὶ Ἰσραὴλ τῷ Κυρίῳ Θεῷ αὐτῶν τῷ ἀναγαγόντι αὐτοὺς ἐκ γῆς Αἰγύπτου ὑποκάτωθεν χειρὸς Φαραὼ βασιλέως Αἰγύπτου, καὶ ἐφοβήθησαν θεοὺς ἐτέρους,
\vs{8}καὶ ἐπορεύθησαν τοῖς δικαιώμασι τῶν ἐθνῶν ὧν ἐξῇρε Κύριος ἐκ προσώπου υἱῶν Ἰσραὴλ, καὶ οἱ βασιλεῖς Ἰσραὴλ ὅσοι ἐποίησαν,
\vs{9}καὶ ὅσοι ἠμφιέσαντο οἱ υἱοὶ Ἰσραὴλ λόγους, οὐχ οὕτως κατὰ Κυρίου Θεου αὐτῶν· καὶ ᾠκοδόμησαν ἑαυτοῖς ὑψηλὰ ἐν πάσαις ταῖς πόλεσιν αὐτῶν ἀπὸ πύργου φυλασσόντων ἕως πόλεως ὀχυρᾶς,
\vs{10}καὶ ἐστήλωσαν ἑαυτοῖς στήλας καὶ ἄλση ἐπὶ παντὶ βουνῷ ὑψηλῷ καὶ ὑποκάτω παντὸς ξύλου ἀλσώδους,
\vs{11}καὶ ἐθυμίασαν ἐκεῖ ἐν πᾶσιν ὑψηλοῖς, καθὼς τὰ ἔθνη ἃ ἀπῴκησε Κύριος ἐκ προσώπου αὐτῶν, καὶ ἐποίησαν κοινωνοὺς, καὶ ἐχάραξαν τοῦ παροργίσαι τὸν Κύριον,
\vs{12}καὶ ἐλάτρευσαν τοῖς εἰδώλοις οἷς εἶπε Κύριος αὐτοῖς, οὐ ποιήσετε τὸ ῥῆμα τοῦτο τῷ Κυρίῳ.

\vs{13}Καὶ διεμαρτύρατο Κύριος ἐν τῷ Ἰσραὴλ καὶ ἐν τῷ Ἰούδα καὶ ἐν χειρὶ πάντων τῶν προφητῶν αὐτοῦ παντὸς ὁρῶντος, λέγων, ἀποστράφητε ἀπὸ τῶν ὁδῶν ὑμῶν τῶν πονηρῶν, καὶ φυλάξατε τὰς ἐντολάς μου, καὶ τὰ δικαιώματά μου, καὶ πάντα τὸν νόμον ὃν ἐνετειλάμην τοῖς πατράσιν ὑμῶν, ὅσα ἀπέστειλα αὐτοῖς ἐν χειρὶ τῶν δούλων μου τῶν προφητῶν.
\vs{14}Καὶ οὐκ ἤκουσαν, καὶ ἐσκλήρυναν τὸν νῶτον αὐτῶν ὑπὲρ τὸν νώτον τῶν πατέρων αὐτῶν.
\vs{15}Καὶ τὰ μαρτύρια αὐτοῦ ὅσα διεμαρτύρατο αὐτοῖς οὐκ ἐφύλαξαν, καὶ ἐπορεύθησαν ὀπίσω τῶν ματαίων, καὶ ἐματαιώθησαν, καὶ ὀπίσω τῶν ἐθνῶν τῶν περικύκλῳ αὐτῶν, ὧν ἐνετείλατο Κύριος αὐτοῖς μὴ ποιῆσαι κατὰ ταῦτα.
\vs{16}Ἐγκατέλιπον τὰς ἐντολὰς Κυρίου Θεοῦ αὐτῶν, καὶ ἐποίησαν ἑαυτοῖς χώνευμα δύο δαμάλεις, καὶ ἐποίησαν ἄλση, καὶ προσεκύνησαν πάσῃ τῇ δυνάμει τοῦ οὐρανοῦ, καὶ ἐλάτρευσαν τῷ Βάαλ.
\vs{17}Καὶ διῆγον τοὺς υἱοὺς αὐτῶν καὶ τὰς θυγατέρας αὐτῶν ἐν πυρὶ, καὶ ἐμαντεύοντο μαντείας, καὶ οἰωνίζοντο· καὶ ἐπράθησαν τοῦ ποιῆσαι τὸ πονηρὸν ἐν ὀφθαλμοῖς Κυρίου παροργίσαι αὐτόν.

\vs{18}Καὶ ἐθυμώθη Κύριος σφόδρα ἐν τῷ Ἰσραὴλ, καὶ ἀπέστησεν αὐτοὺς ἀπὸ τοῦ προσώπου αὐτοῦ, καὶ οὐχ ὑμελείφθη πλὴν φυλὴ Ἰούδα μονωτάτη.
\vs{19}Καί γε Ἰούδας οὐκ ἐφύλαξε τὰς ἐντολὰς Κυρίου τοῦ Θεοῦ αὐτῶν· καὶ ἐπορεύθησαν ἐν τοῖς δικαιώμασιν Ἰσραὴλ οἷς ἐποίησαν, καὶ ἀπεώσαντο τὸν Κύριος.

\vs{20}Καὶ ἐθυμώθη Κύριος παντὶ σπέρματι Ἰσραὴλ, καὶ ἐσάλευσεν αὐτοὺς, καὶ ἔδωκεν αὐτοὺς ἐν χειρὶ διαρπαζόντων αὐτοὺς, ἕως οὗ ἀπέῤῥιψεν αὐτοὺς ἀπὸ προσώπου αὐτοῦ.
\vs{21}Ὅτι πλὴν Ἰσραὴλ ἐπάνωθεν οἴκου Δαυὶδ, καὶ ἐβασίλευσαν τὸν Ἱεροβοὰμ υἱὸν Ναβάτ· καὶ ἐξέωσεν Ἱεροβοὰμ τὸν Ἰσραὴλ ἐξόπισθε Κυρίου, καὶ ἐξήμαρτεν αὐτοὺς ἁμαρτίαν μεγάλην.
\vs{22}Καὶ ἐπορεύθησαν οἱ υἱοὶ Ἰσραὴλ ἐν πάσῃ ἁμαρτία Ἱεροβοὰμ ἧς ἐποιήσεν, οὐκ ἀπέστησαν ἀπʼ αὐτῆς
\vs{23}ἕως οὗ μετέστησε Κύριος τὸν Ἰσραὴλ ἀπὸ προσώπου αὐτοῦ, καθὼς ἐλάλησε Κύριος ἐν χειρὶ πάντων τῶν δούλων αὐτοῦ τῶν προφητῶν· καὶ ἀπῳκίσθη Ἰσραὴλ ἐπάνωθεν τῆς γῆς αὐτοῦ εἰς Ἀσσυρίους ἕως τῆς ἡμέρας ταύτης.

\vs{24}Καὶ ἤγαγε βασιλεὺς Ἀσσυρίων ἐκ Βαβυλῶνος τὸν ἐκ Χουνθὰ, ἀπὸ Ἀϊὰ, καὶ ἀπὸ Αἱμὰθ, καὶ Σεπφαρουαῒν, καὶ κατῳκίσθησαν ἐν πόλεσι Σαμαρείας ἀντὶ τῶν υἱῶν Ἰσραὴλ, καὶ ἐκληρονόμησαν τὴν Σαμάρειαν, καὶ κατῴκίσθησαν ἐν ταῖς πόλεσιν αὐτῆς.
\vs{25}Καὶ ἐγένετο ἐν ἀρχῇ τῆς καθέδρας αὐτῶν οὐκ ἐφοβήθησαν τὸν Κύριον, καὶ ἀπέστειλε Κύριος ἐν αὐτοῖς τοὺς λέοντας, καὶ ἦσαν ἀποκτενίοντες ἐν αὐτοῖς.
\vs{26}Καὶ εἶπαν τῷ βασιλεῖ Ἀσσυρίων λέγοντες, τὰ ἔθνη ἃ ἀπῴκισας καὶ ἀντεκάθισας ἐν πόλεσι Σαμαρείας οὐκ ἔγνωσαν τὸ κρίμα τοῦ Θεοῦ τῆς γῆς, καὶ ἀπέστειλεν εἰς αὐτοὺς τοὺς λέοντας, καὶ ἰδού εἰσι θανατοῦντες αὐτοὺς, καθότι οὐκ οἴδασι τὸ κρίμα τοῦ Θεοῦ τῆς γῆς.
\vs{27}Καὶ ἐνετείλατο ὁ βασιλεὺς Ἀσσυρίων, λέγων, ἀπαγάγετε ἐκεῖθεν, καὶ πορευέσθωσαν, καὶ κατοικήτωσαν ἐκεῖ, καὶ φωτιοῦσιν αὐτοὺς τὸ κρίμα τοῦ Θεοῦ τῆς γῆς.
\vs{28}Καὶ ἤγαγον ἕνα τῶν ἱερέων ὧν ἀπῴκισαν ἀπὸ Σαμαρείας, καὶ ἐκάθισεν ἐν Βαιθὴλ, καὶ ἦν φωτίζων αὐτοὺς πῶς φοβηθῶσι τὸν Κύριον.

\vs{29}Καὶ ἦσαν ποιοῦντες ἔθνη ἔθνη θεοὺς αὐτῶν· καὶ ἔθηκαν ἐν οἴκῳ τῶν ὑψηλῶν ὧν ἐποίησαν οἱ Σαμαρεῖται, ἔθνη ἔθνη ἐν ταῖς πόλεσιν αὐτῶν ἐν αἷς κατῴκουν.
\vs{30}Καὶ οἱ ἄνδρες Βαβυλῶνος ἐποίησαν τὴν Σωκχὼθ Βενὶθ, καὶ οἱ ἄνδρες Χοὺθ ἐποίησαν τὴν Ἐργέλ, καὶ οἱ ἄνδρες Αἱμὰθ ἐποίησαν τὴν Ἀσιμὰθ,
\vs{31}καὶ οἱ Εὐαῖοι ἐποίησαν τὴν Ἐβλαζὲρ καὶ τὴν Θαρθὰκ, καὶ ὁ Σεπφαρουαῒμ ἡνίκα κατέκαιον τοὺς υἱοὺς αὐτῶν ἐν πυρὶ τῷ Ἀδραμέλεχ καὶ Ἀνημελὲχ θεοῖς Σεπφαρουαΐμ.
\vs{32}Καὶ ἦσαν φοβούμενοι τὸν Κύριον· καὶ κατῴκισαν τὰ βδελύγματα αὐτῶν ἐν τοῖς οἴκοις τῶν ὑψηλῶν ἃ ἐποίησαν ἐν Σαμαρείᾳ, ἔθνος ἔθνος ἐν πόλει ἐν ᾗ κατῴκουν ἐν αὐτῇ· καὶ ἦσαν φοβούμενοι τὸν Κύριον· καὶ ἐποίησαν ἑαυτοῖς ἱερεῖς τῶν ὑψηλῶν, καὶ ἐποίησαν ἑαυτοῖς ἐν οἴκῳ τῶν ὑψηλῶν.
\vs{33}Καὶ τὸν Κύριον ἐφοβοῦντο, καὶ τοῖς θεοῖς αὐτῶν ἐλάτρευον κατὰ τὸ κρίμα τῶν ἐθνῶν, ὅθεν ἀπῴκισαν αὐτοὺς ἐκεῖθεν.

\vs{34}Ἕως τῆς ἡμέρας ταύτης αὐτοὶ ἐποίουν κατὰ τὸ κρίμα αὐτῶν· αὐτοὶ φοβοῦνται, καὶ αὐτοὶ ποιοῦσι κατὰ τὰ δικαιώματα αὐτῶν, καὶ κατὰ τὴν κρίσιν αὐτῶν, καὶ κατὰ τὸν νόμον, καὶ κατὰ τὴν ἐντολὴν ἣν ἐνετείλατο Κύριος τοῖς υἱοῖς Ἰακὼβ, οὗ ἔθηκε τὸ ὄνομα αὐτοῦ Ἰσραήλ.
\vs{35}Καὶ διέθετο Κύριος μετʼ αὐτῶν διαθήκην, καὶ ἐνετείλατο αὐτοῖς, λέγων, οὐ φοβηθήσεσθε θεοὺς ἑτέρους, καὶ οὐ προσκυνήσετε αὐτοῖς, καὶ οὐ λατρεύσετε αὐτοῖς, καὶ οὐ θυσιάσετε αὐτοῖς·
\vs{36}Ὅτι ἀλλʼ ἢ τῷ Κυρίῳ ὃς ἀνήγαγεν ὑμᾶς ἐκ γῆς Αἰγύπτου ἐν ἰσχύϊ μεγάλῃ καὶ ἐν βραχίονι ὑψηλῷ, αὐτὸν φοβηθήσεσθε, καὶ αὐτῷ προσκυνήσετε, αὐτῷ θύσετε.
\vs{37}Τὰ δικαιώματα καὶ τὰ κρίματα καὶ τὸν νόμον καὶ τὰς ἐντολὰς, ἃς ἔγραψεν ὑμῖν ποιεῖν, φυλάσσεσθε πάσας τὰς ἡμέρας, καὶ οὐ φοβηθήσεσθε θεοὺς ἑτέρους.
\vs{38}Καὶ τὴν διαθήκην ἣν διέθετο μεθʼ ὑμῶν οὐκ ἐπιλήσεσθε· καὶ οὐ φοβηθήσεσθε θεοὺς ἑτέρους·
\vs{39}Ἀλλʼ ἢ τὸν Κύριον Θεὸν ὑμῶν φοβηθήσεσθε, καὶ αὐτὸς ἐξελεῖται ὑμᾶς ἐκ πάντων τῶν ἐχθρῶν ὑμῶν.

\vs{40}Καὶ οὐκ ἀκούσεσθε ἐτὶ τῷ κρίματι αὐτῶν, ὃ αὐτοὶ ποιοῦσι.
\vs{41}Καὶ ἦσαν τὰ ἔθνη ταῦτα φοβούμενοι τὸν Κύριον, καὶ τοῖς γλυπτοῖς αὐτῶν ἦσαν δουλεύοντες· καί γε οἱ υἱοὶ καὶ υἱοὶ τῶν υἱῶν αὐτῶν, καθὰ ἐποίησαν οἱ πατέρες αὐτῶν, ποιοῦσιν ἕως τῆς ἡμέρας ταύτης.

\ch{18}
Καὶ ἐγένετο ἐν ἔτει τρίτῳ τῷ Ὠσηὲ υἱῷ Ἠλὰ βασιλεῖ Ἰσραὴλ ἐβασίλευσεν Ἐζεκίας υἱὸς Ἄχαζ βασιλέως Ἰούδα.
\vs{2}Υἱὸς εἴκοσι καὶ πέντε ἐτῶν ἐν τῷ βασιλεύειν αὐτὸν, καὶ εἴκοσι καὶ ἐννέα ἔτη ἐβασίλευσεν ἐν Ἱερουσαλὴμ, καὶ ὄνομα τῇ μητρὶ αὐτοῦ Ἄβού, θυγάτηρ Ζαχαρίου.
\vs{3}Καὶ ἐποίησε τὸ εὐθὲς ἐν ὀφθαλμοῖς Κυρίου κατὰ πάντα ὅσα ἐποίησε Δαυὶδ ὁ πατὴρ αὐτοῦ.
\vs{4}Αὐτὸς ἐξῆρε τὰ ὑψηλὰ, καὶ συνέτριψε τὰς στήλας, καὶ ἐξωλέθρευσε τὰ ἄλση, καὶ τὸν ὄφιν τὸν χαλκοῦν ὃν ἐποίησε Μωυσῆς, ὅτι ἕως τῶν ἡμερῶν ἐκείνων ἦσαν οἱ υἱοὶ Ἰσραὴλ θυμιῶντες αὐτῷ· καὶ ἐκάλεσεν αὐτὸν Νεεσθάν.
\vs{5}Ἐν Κυρίῳ Θεῷ Ἰσραὴλ ἤλπισεν, καὶ μετʼ αὐτὸν οὐκ ἐγενήθη ὅμοιος αὐτῷ ἐν βασιλεῦσιν Ἰούδα, καὶ ἐν τοῖς γενομένοις ἔμπροσθεν αὐτοῦ.
\vs{6}Καὶ ἐκολλήθη τῷ Κυρίῳ, οὐκ ἀπέστη ὄπισθεν αὐτοῦ, καὶ ἐφύλαξε τὰς ἐντολὰς αὐτοῦ ὅσας ἐνετείλατο Μωυσῇ.

\vs{7}Καὶ ἦν Κύριος μετʼ αὐτοῦ, καὶ ἐν πᾶσιν οἷς ἐποίει, συνῆκε· καὶ ἠθέτησεν ἐν τῷ βασιλεῖ Ἀσσυρίων, καὶ οὐκ ἐδούλευσεν αὐτῷ.
\vs{8}Αὐτὸς ἐπάταξεν τοὺς ἀλλοφύλους ἕως Γάζης, καὶ ἕως ὁρίου αὐτῆς ἀπὸ πύργου φυλασσόντων καὶ ἕως πόλεως ὀχυρᾶς.

\vs{9}Καὶ ἐγένετο ἐν τῷ ἔτει τῷ τετάρτῳ βασιλεῖ Ἐζεκίᾳ, αὐτὸς ἐνιαυτὸς ὁ ἕβδομος τῷ Ὠσηὲ υἱῷ Ἠλὰ βασιλεῖ Ἰσραήλ, ἀνέβη Σαλαμανασσὰρ βασιλεὺς Ἀσσυρίων ἐπὶ Σαμάρειαν, καὶ ἐπολιόρκει ἐπʼ αὐτὴν,
\vs{10}καὶ κατελάβετο αὐτὴν ἀπὸ τέλους τριῶν ἐτῶν ἐν ἔτει ἕκτῳ τῷ Ἑζεκίᾳ, αὐτὸς ἐνιαυτὸς ἔννατος τῷ Ὠσηὲ βασιλεῖ Ἰσραὴλ, καὶ συνελήμφθη Σαμάρεια.
\vs{11}Καὶ ἀπῴκισε βασιλεὺς Ἀσσυρίων τὴν Σαμάρειαν εἰς Ἀσσυρίους, καὶ ἔθηκεν αὐτοὺς ἐν Ἁλαὲ καὶ ἐν Ἀβὼρ ποταμῷ Γωζὰν καὶ ὄρὴ Μήδων·
\vs{12}ἀνθʼ ὧν ὅτι οὐκ ἤκουσαν τῆς φωνῆς Κυρίου Θεοῦ αὐτῶν, καὶ παρέβησαν τὴν διαθήκην αὐτοῦ πάντα ὅσα ἐνετείλατο Μωυσῆς ὁ δοῦλος Κυρίου, καὶ οὐκ ἤκουσαν καὶ οὐκ ἐποίησαν.

\vs{13}Καὶ τῷ τεσσαρεσκαιδεκάτῳ ἔτει ἔτει τοῦ βασιλέως βασιλέως Ἑζεκίου ἀνέβη Σενναχηρὶμ βασιλεὺς Ἀσσυρίων ἐπὶ τὰς πόλεις Ἰούδα τὰς ὀχυρὰς, καὶ συνέλαβεν αὐτάς.
\vs{14}Καὶ ἀπέστειλεν Ἐζεκίας βασιλεὺς Ἰούδα ἀγγέλους πρὸς βασιλέα Ἀσσυρίων εἰς Λαχὶς, λέγων, ἡμάρτηκα, ἀποστράφηθι ἀπʼ ἐμοῦ· ὃ ἐὰν ἐπιθῇς ἐπʼ ἐμὲ, βαστάσω· καὶ ἐπέθηκεν ὁ βασιλεὺς Ἀσσυρίων ἐπὶ Ἑζεκίαν βασιλέα Ἰούδα τριακόσια τάλαντα ἀργυρίου, καὶ τριάκοντα τάλαντα χρυσίου.
\vs{15}Καὶ ἔδωκεν Ἑζεκίας πᾶν τὸ ἀργύριον τὸ εὑρεθὲν ἐν οἴκῳ Κυρίου, καὶ ἐν θησαυροῖς οἴκου τοῦ βασιλέως.
\vs{16}Ἐν τῷ καιρῷ ἐκείνῳ συνέκοψεν Ἑζεκίας τὰς θύρας ναοῦ, καὶ τὰ ἐστηριγμένα ἃ ἐχρύσωσεν Ἑζεκίας ὁ βασιλεὺς Ἰούδα, καὶ ἔδωκεν αὐτὰ βασιλεῖ Ἀσσυρίων.

\vs{17}Καὶ ἀπέστειλε βασιλεὺς Ἀσσυρίων τὸν Θανθὰν καὶ τὸν Ῥαφὶς καὶ τὸν Ῥαψάκην ἐκ Λαχὶς πρὸς τὸν βασιλέα Ἑζεκίαν ἐν δυνάμει βαρείᾳ ἐπὶ Ἱερουσαλήμ· καὶ ἀνέβησαν καὶ ἦλθον εἰς Ἱερουσαλὴμ, καὶ ἔστησαν ἐν τῷ ὑδραγωγῷ τῆς κολυμβήθρας τῆς ἄνω, ἥ ἐστιν ἐν τῇ ὁδῷ τοῦ ἀγροῦ τοῦ γναφέως.
\vs{18}Καὶ ἐβόησαν πρὸς Ἐζεκίαν· καὶ ἦλθον πρὸς αὐτὸν Ἑλιακὶμ υἱὸς Χελκίου ὁ οἰκονόμος, καὶ Σωμνὰς ὁ γραμματεὺς, καὶ Ἰωὰς ὁ νἱὸς Σαφὰτ ὁ ἀναμιμνήσκων.

\vs{19}Καὶ εἶπε πρὸς αὐτοὺς Ῥαψάκης, εἴπατε δὴ πρὸς Ἐζεκίαν, τάδε λέγει ὁ βασιλεὺς, ὁ μέγας βασιλεὺς Ἀσσυρίων, τί ἡ πεποίθησις αὕτη ἣν πέποιθας;
\vs{20}Εἶπας, πλὴν λόγοι χειλέων, Βουλὴ καὶ δύναμις εἰς πόλεμον· νῦν οὖν τίνι πεποιθὼς ἠθέτησας ἐν ἐμοί;
\vs{21}Νῦν ἰδοὺ πέποιθας σαυτῷ ἐπὶ τὴν ῥἁβδον τὴν καλαμίνην τὴς τεθλασμένην ταύτην, ἐπʼ Αἴγυπτον; ὃς ἂν στηριχθῇ ἀνὴρ ἐπʼ αὐτὴν, καὶ εἰσελεύσεται εἰς τὴν χεῖρα αὐτοῦ, καὶ τρήσει αὐτήν· οὕτως Φαραὼ βασιλεὺς Αἰγύπτου πᾶσι τοῖς πεποιθόσιν ἐπʼ αὐτόν.
\vs{22}Καὶ ὅτι εἶπας πρὸς μὲ, ἐπὶ Κύριον Θεὸν πεποίθαμεν· οὐχὶ αὐτὸς οὗτος ἀπέστησεν Ἐζεκίας τὰ ὑψηλὰ αὐτοῦ καὶ τὰ θυσιαστήρια αὐτοῦ, καὶ εἶπε τῷ Ἰούδᾳ καὶ τῇ Ἱερουσαλὴμ, ἐνώπιον τοῦ θυσιαστηρίου τούτου προσκυνήσετε ἐν Ἰερουσαλήμ;
\vs{23}Καὶ νῦν μίχθητε δὴ τῷ κυρίῳ μου βασιλεῖ Ἀσσυρίων, καὶ δώσω σοι δισχιλίους ἵππους, εἰ δυνήσῃ δοῦναι σεαυτῷ ἐπιβάτας ἐπʼ αὐτούς.
\vs{24}Καὶ πῶς ἀποστρέψεις τὸ πρόσωπον τοπάρχου ἑνὸς τῶν δούλων τοῦ κυρίου μου τῶν ἐλαχίστων; καὶ ἤλπισας σαυτῷ ἐπʼ Αἴγυπτον εἰς ἅρματα καὶ ἱππεῖς.
\vs{25}Καὶ νῦν μὴ ἄνευ Κυρίου ἀνέβημεν ἐπὶ τὸν τόπον τοῦτον τοῦ διαφθεῖραι αὐτόν; Κύριος εἶπε πρὸς μὲ, ἀνάβηθι ἐπὶ τὴν γῆν ταύτην καὶ διάφθειρον αὐτήν.

\vs{26}Καὶ εἶπεν Ἑλιακὶμ υἱὸς Χελκίου καὶ Σωμνὰς καὶ Ἰωὰς πρὸς Ῥαψάκην, λάλησον δὴ πρὸς τοὺς παῖδάς σου Συριστὶ, ὅτι ἀκούομεν ἡμεῖς· καὶ οὐ λαλήσεις μεθʼ ἡμῶν Ἰουδαϊστὶ· καὶ ἱνατί λαλεῖς ἐν τοῖς ὠσὶ τοῦ λαοῦ τοῦ ἐπὶ τοῦ τείχους;
\vs{27}Καὶ εἶπε πρὸς αὐτοὺς Ῥαψάκης, μὴ ἐπὶ τὸν κύριόν σου καὶ πρὸς σὲ ἀπέστειλέ με ὁ κύριός μου λαλῆσαι τοὺς λόγους τούτους; οὐχὶ ἐπὶ τοὺς ἄνδρας τοὺς καθημένους ἐπὶ τοῦ τείχους, τοῦ φαγεῖν τὴν κόπρον αὐτῶν, καὶ πιεῖν τὸ οὖρον αὐτῶν μεθʼ ὑμῶν ἅμα;

\vs{28}Καὶ ἔστη Ῥαψάκης καὶ ἐβόησε φωνῇ μεγάλῃ Ἰουδαϊστὶ, καὶ ἐλάλησε καὶ εἶπεν, ἀκούσατε τοὺς λόγους τοῦ μεγάλου βασιλέως Ἀσσυρίων.
\vs{29}Τάδε λέγει ὁ βασιλεὺς, μὴ ἐπαιρέτω ὑμᾶς Ἐζεκίας λόγοις, ὅτι οὐ μὴ δύνηται ὑμᾶς ἐξελέσθαι ἐκ χειρὸς αὐτοῦ.
\vs{30}Καὶ μὴ ἐπελπιζέτω ὑμᾶς Ἐζεκίας πρὸς Κὺριον, λέγων, ἐξαιρούμενος ἐξελεῖται Κύριος, οὐ μὴ παραδοθῇ ἡ πόλις αὕτη ἐν χειρὶ βασιλέως Ἀσσυρίων.
\vs{31}Μὴ ἀκούετε Ἐζεκίου· ὅτι τάδε λέγει ὁ βασιλεὺς Ἀσσυρίων, ποιήσατε μετʼ ἐμοῦ εὐλογίαν, καὶ ἐξέλθατε πρὸς μὲ, καὶ πίεται ἀνὴρ τὴν ἄμπελον αὐτοῦ, καὶ ἀνὴρ τὴν συκῆν αὐτοῦ φάγεται, καὶ πίεται ὕδωρ τοῦ λάκκου αὐτοῦ,
\vs{32}ἕως ἔλθω καὶ λάβω ὑμᾶς εἰς γῆν ὡς γῆ ὑμῶν, γῆ σίτου καὶ οἴνου καὶ ἄρτου καὶ ἀμπελώνων, γῆ ἐλαίας ἐλαίου καὶ μέλιτος, καὶ ζήσετε καὶ οὐ μὴ ἀποθάνητε· καὶ μὴ ἀκούετε Ἐζεκίου, ὅτι ἀπατᾷ ὑμᾶς, λέγων, Κύριος ῥύσεται ὑμᾶς.
\vs{33}Μὴ ῥυόμενοι ἐῤῥύσαντο οἱ θεοὶ τῶν ἐθνῶν ἕκαστος τὴν ἑαυτοῦ χώραν ἐκ χειρὸς βασιλέως Ἀσσυρίων;
\vs{34}Ποῦ ἐστιν ὁ θεὸς Αἱμὰθ, καὶ Ἀρφάδ; ποῦ ἐστιν ὁ θεὸς Σεπφαρουαῒμ, Ἀνὰ, καὶ Ἀβὰ, ὅτι ἐξείλαντο Σαμάρειαν ἐκ χειρός μου;
\vs{35}Τίς ἐν πᾶσι τοῖς θεοῖς τῶν γαιῶν, οἳ ἐξείλαντο τὰς γᾶς αὐτῶν ἐκ χειρός μου, ὅτι ἐξελεῖται Κύριος τὴν Ἱερουσαλὴμ ἐκ χειρός μου;

\vs{36}Καὶ ἐκώφευσαν καὶ οὐκ ἀπεκρίθησαν αὐτῷ λόγον, ὅτι ἐντολὴ τοῦ βασιλέως, λέγων, οὐκ ἀποκριθήσεσθε αὐτῷ.
\vs{37}Καὶ εἰσῆλθεν Ἐλιακὶμ υἱὸς Χελκίου ὁ οἰκονόμος, καὶ Σωμνὰς ὁ γραμματεὺς, καὶ Ἰωὰς υἱὸς Σαφὰτ ὁ ἀναμιμνήσκων πρὸς Ἐζεκίαν, διεῤῥηχότες τὰ ἱμάτια, καὶ ἀνήγγειλαν αὐτῷ τοὺς λόγους Ῥαψάκου.

\ch{19}
Καὶ ἐγένετο ὡς ἤκουσεν ὁ βασιλεὺς Ἐζεκίας, καὶ διέῤῥηξε τὰ ἱμάτια αὐτοῦ, καὶ περιεβάλετο σάκκον, καὶ εἰσῆλθεν εἰς οἶκον Κυρίου.
\vs{2}Καὶ ἀπέστειλεν Ἑλιακὶμ τὸν οἰκονόμον, καὶ Σωμνὰν τὸν γραμματέα, καὶ τοὺς πρεσβυτέρους τῶν ἱερέων περιβεβλημένους σάκκους, πρὸς Ἡσαΐαν τὸν προφήτην υἱὸν Ἀμώς.
\vs{3}Καὶ εἶπον πρὸς αὐτὸν, τάδε λέγει Ἐζεκίας, ἡμέρα θλίψεως καὶ ἐλεγμοῦ καὶ παροργισμοῦ ἡ ἡμέρα αὕτη· ὅτι ἦλθον υἱοὶ ἕως ὠδίνων, καὶ ἰσχὺς οὐκ ἔστι τῇ τικτούσῃ.
\vs{4}Εἴ πως εἰσακούσεται Κύριος ὁ Θεός σου πάντας τοὺς λόγους Ῥαψάκου, ὃν ἀπέστειλεν αὐτὸν βασιλεὺς Ἀσσυρίων ὁ κύριος αὐτοῦ ὀνειδίζειν Θεὸν ζῶντα, καὶ βλασφημεῖν ἐν λόγοις οἷς ἤκουσε Κύριος ὁ Θεός σου, καὶ λήψῃ προσευχὴν περὶ τοῦ λείμματος τοῦ εὑρισκομένου.

\vs{5}Καὶ ἦλθον οἱ παῖδες τοῦ βασιλέως Ἐζεκίου πρὸς Ἡσαΐαν.
\vs{6}Καὶ εἶπεν αὐτοῖς Ἡσαΐας, Τάδε ἐρεῖτε πρὸς τὸν κύριον ὑμῶν, τάδε λέγει Κύριος, μὴ φοβηθῇς ἀπὸ τῶν λόγων ὧν ἤκουσας, ὧν ἐβλασφήμησαν τὰ παιδάρια βασιλέως Ἀσσυρίων.
\vs{7}Ἰδοὺ ἐγὼ δίδωμι ἐν αὐτῷ πνεῦμα, καὶ ἀκούσεται ἀγγελίαν, καὶ ἀποστραφήσεται εἰς τὴν γῆν αὐτοῦ· καὶ καταβαλῶ αὐτὸν ἐν ῥομφαίᾳ ἐν τῇ γῇ αὐτοῦ.

\vs{8}Καὶ ἐπέστρεψε Ῥαψάκης, καὶ εὗρε τὸν βασιλέα Ἀσσυρίων πολεμοῦντα ἐπὶ Λοβνὰ, ὅτι ἤκουσεν ὅτι ἀπῇρεν ἐκ Λαχίς.
\vs{9}Καὶ ἤκουσε περὶ Θαρακὰ βασιλέως Αἰθιόπων, λέγων, ἰδοὺ ἐξῆλθε πολεμεῖν μετὰ σοῦ· καὶ ἐπέστρεψε, καὶ ἀπέστειλεν ἀγγέλους πρὸς Ἐζεκίαν, λέγων,
\vs{10}μὴ ἐπαιρέτω σε ὁ Θεός σου ἐφʼ ᾧ σὺ πέποιθας ἐν αὐτῷ, λέγων, οὐ μὴ παραδοθῇ Ἱερουσαλὴμ εἰς χεῖρας βασιλέως Ἀσσυρίων.
\vs{11}Ἰδοὺ σὺ ἤκουσας πάντα ὅσα ἐποίησαν βασιλεῖς Ἀσσυρίων πάσαις ταῖς γαίαις τοῦ ἀναθεματίσαι αὐτάς· καὶ σὺ ῥυσθήσῃ;
\vs{12}Μὴ ἐξαιρούμενοι ἐξείλαντο αὐτοὺς οἱ θεοὶ τῶν ἐθνῶν, οὓς διέφθειραν οἱ πατρές μου, τήν τε Γωζὰν, καὶ τὴν Χαῤῥὰν, καὶ τὴν Ῥαφὶς, καὶ υἱοὺς Ἐδὲμ τοὺς ἐν Θαεσθέν;
\vs{13}Ποῦ ἐστιν ὁ βασιλεὺς Αἱμὰθ, καὶ ὁ βασιλεὺς Ἀρφάδ; καὶ ποῦ ἐστιν ὁ βασιλεὺς τῆς πόλεως Σεπφαρουαῒν, Ἀνὰ, καὶ Ἀβά;

\vs{14}Καὶ ἔλαβεν Ἐζεκίας τὰ βιβλία ἐκ χειρὸς τῶν ἀγγέλων, καὶ ἀνέγνω αὐτά· καὶ ἀνέβη εἰς οἶκον Κυρίου, καὶ ἀνέπτυξεν αὐτὰ Ἐζεκίας ἐναντίον Κυρίου,
\vs{15}καὶ εἶπε, Κύριε ὁ Θεὸς Ἰσραὴλ ὁ καθήμενος ἐπὶ τῶν χερουβὶμ, σὺ εἶ ὁ Θεὸς μόνος ἐν πάσαις ταῖς βασιλείαις τῆς γῆς, σὺ ἐποίησας τὸν οὐρανὸν καὶ τὴν γῆν.
\vs{16}Κλῖνον Κύριε τὸ οὖς σου καὶ ἄκουσον, ἄνοιξον Κύριε τοὺς ὀφθαλμούς σου καὶ ἴδε, καὶ ἄκουσον τοὺς λόγους Σενναχηρὶμ οὓς ἀπέστειλεν ὀνειδίζειν Θεὸν ζῶντα.
\vs{17}Ὅτι ἀληθείᾳ Κύριε ἠρήμωσαν βασιλεῖς Ἀσσυρίων τὰ ἔθνη,
\vs{18}καὶ ἔδωκαν τοὺς θεοὺς αὐτῶν εἰς τὸ πῦρ, ὅτι οὐ θεοί εἰσιν, ἀλλʼ ἢ ἔργα χειρῶν ἀνθρώπων ξύλα καὶ λίθος· καὶ ἀπώλεσαν αὐτούς.
\vs{19}Καὶ νῦν, Κύριε ὁ Θεὸς ἡμῶν, σῶσον ἡμᾶς ἐκ χειρὸς αὐτοῦ, καὶ γνώσονται πᾶσαι αἱ βασιλεῖαι τῆς γῆς, ὅτι σὺ Κύριος ὁ Θεὸς μόνος.

\vs{20}Καὶ ἀπέστειλεν Ἡσαΐας υἱὸς Ἀμὼς πρὸς Ἐζεκίαν, λέγων, τάδε λέγει Κύριος ὁ Θεὸς τῶν δυνάμεων Θεὸς Ἰσραὴλ, ἃ προσηύξω πρὸς μὲ περὶ Σενναχηρὶμ βασιλέως Ἀσσυρίων, ἤκουσα.
\vs{21}Οὗτος ὁ λόγος ὃν ἐλάλησε Κύριος ἐπʼ αὐτὸν, ἐξουδένησέ σε καὶ ἐμυκτήρισέ σε παρθένος θυγάτηρ Σιὼν, ἐπὶ σοὶ κεφαλὴν αὐτῆς ἐκίνησε θυγάτηρ Ἱερουσαλήμ.
\vs{22}Τίνα ὠνείδισας, καὶ τίνα ἐβλασφήμησας; καὶ ἐπὶ τίνα ὕψωσας φωνὴν, καὶ ἦρας εἰς ὕψος τοὺς ὀφθαλμούς σου; εἰς τὸν ἅγιον τοῦ Ἰσραήλ;

\vs{23}Ἐν χειρὶ ἀγγέλων σου ὠνείδισας Κύριον, καὶ εἶπας, ἐν τῷ πλήθει τῶν ἁρμάτων μου ἐγὼ ἀναβήσομαι εἰς ὕψος ὀρέων μηροὺς τοῦ Λιβάνου, καὶ ἔκοψα τὸ μέγεθος τῆς κέδρου αὐτοῦ, τὰ ἐκλεκτὰ κυπαρίσσων αὐτοῦ, καὶ ἦλθον εἰς μέσον δρυμοῦ καὶ Καρμήλου.
\vs{24}Ἐγὼ ἔψυξα καὶ ἔπιον ὕδατα ἀλλότρια, καὶ ἐξερήμωσα τῷ ἴχνει τοῦ ποδός μου πάντας ποταμοὺς περιοχῆς.
\vs{25}Ἔπλασα αὐτὴν, συνήγαγον αὐτήν· καὶ ἐγενήθη εἰς ἐπάρσεις ἀποικεσιῶν μαχίμων πόλεις ὀχυράς.
\vs{26}Καὶ οἱ ἐνοικοῦντες ἐν αὐταῖς ἠσθένησαν τῇ χειρὶ, ἔπτηξαν καὶ κατῃσχύνθησαν· ἐγένοντο χόρτος ἀγροῦ, ἢ χλωρὰ βοτάνη, χλόη δωμάτων, καὶ πάτημα ἀπέναντι ἑστηκότος.
\vs{27}Καὶ τὴν καθέδραν σου καὶ τὴν ἔξοδόν σου ἔγνων, καὶ τὸν θυμόν σου ἐπʼ ἐμὲ,
\vs{28}διὰ τὸ ὀργισθῆναί σε ἐπʼ ἐμὲ, καὶ τὸ στρῆνός σου ἀνέβη ἐν τοῖς ὠσί μου· καὶ θήσω τὰ ἄγκιστρά μου ἐν τοῖς μυκτῆρσί σου, καὶ χαλινὸν ἐν τοῖς χείλεσί σου, καὶ ἀποστρέψω σε ἐν τῇ ὁδῷ ᾗ ἦλθες ἐν αὐτῇ.

\vs{29}Καὶ τοῦτό σοι τὸ σημεῖον· φάγε τοῦτον τὸν ἐνιαυτὸν αὐτόματα, καὶ τῷ ἔτει τῷ δευτέρῳ τὰ ἀνατέλλοντα, καὶ ἔτει τρίτῳ σπορὰ καὶ ἀμητὸς καὶ φυτεία ἀμπελώνων, καὶ φάγεσθε τὸν καρπὸν αὐτῶν.
\vs{30}Καὶ προσθήσει τὸν διασεσωσμένον οἴκου Ἰούδα τὸ ὑπολειφθὲν ῥίζαν κάτω, καὶ ποιήσει καρπὸν ἄνω.
\vs{31}Ὅτι ἐξ Ἱερουσαλὴμ ἐξελεύσεται κατάλειμμα, καὶ ἀνασωζόμενος ἐξ ὄρους Σιών· ὁ ζῆλος Κυρίου τῶν δυνάμεων ποιήσει τοῦτο.
\vs{32}Οὐχ οὕτως;

Τάδε λέγει Κύριος πρὸς βασιλέα Ἀσσυρίων, οὐκ εἰσελεύσεται εἰς τὴν πόλιν ταύτην, καὶ οὐ τοξεύσει ἐκεῖ βέλος, καὶ οὐ προφθάσει ἐπʼ αὐτὴν θυρεὸς, καὶ οὐ μὴ ἐκχέῃ πρὸς αὐτὴν πρόσχωμα.

\vs{33}Τῇ ὁδῷ ᾗ ἦλθεν, ἐν αὐτῇ ἀποστραφήσεται, καὶ εἰς τὴν πόλιν ταύτην οὐκ εἰσελεύσεται, λέγει Κύριος.
\vs{34}Καὶ ὑπερασπιῶ ὑπὲρ τῆς πόλεως ταύτης διʼ ἐμὲ καὶ διὰ Δαυὶδ τὸν δοῦλόν μου.

\vs{35}Καὶ ἐγένετο νυκτὸς, καὶ ἐξῆλθεν ἄγγελος Κυρίου καὶ ἐπάταξεν ἐν τῇ παρεμβολῇ τῶν Ἀσσυρίων ἑκατὸν ὀγδοηκονταπέντε χιλιάδας· καὶ ὤρθρισαν τοπρωῒ, καὶ ἰδοὺ πάντες σώματα νεκρά.
\vs{36}Καὶ ἀπῇρε καὶ ἐπορεύθη καὶ ἀπέστρεψε Σενναχηρὶμ βασιλεὺς Ἀσσυρίων, καὶ ᾤκησεν ἐν Νινευῆ.
\vs{37}Καὶ ἐγένετο αὐτοῦ προσκυνοῦντος ἐν οἴκῳ Μεσερὰχ τοῦ θεοῦ αὐτοῦ, καὶ Ἀδραμέλεχ καὶ Σαρασὰρ οἱ υἱοὶ αὐτοῦ ἐπάταξαν αὐτὸν ἐν μαχαίρᾳ· καὶ αὐτοὶ ἐσώθησαν εἰς γῆν Ἀραράθ· καὶ ἐβασίλευσεν Ἀσορδὰν ὁ υἱὸς αὐτοῦ ἀντʼ αὐτοῦ.

\ch{20}
Ἐν ταῖς ἡμέραις ἐκείναις ἠῤῥώστησεν Ἐζεκίας εἰς θάνατον· καὶ εἰσῆλθε πρὸς αὐτὸν Ἡσαΐας υἱὸς Ἀμὼς ὁ προφήτης, καὶ εἶπε πρὸς αὐτὸν, τάδε λέγει Κύριος, ἔντειλαι τῷ οἴκῳ σου, ἀποθνήσκεις σὺ καὶ οὐ ζήσῃ.
\vs{2}Καὶ ἐπέστρεψεν Ἐζεκίας πρὸς τὸν τοῖχον, καὶ ηὔξατο πρὸς Κύριον λέγων
\vs{3}Κύριε, μνήσθητι δὴ ὅσα περιεπάτησα ἐνώπιόν σου ἐν ἀληθείᾳ καὶ καρδίᾳ πλήρει, καὶ τὸ ἀγαθὸν ἐν ὀφθαλμοῖς σου ἐποίησα· καὶ ἔκλαυσεν Ἐζεκίας κλαυθμῷ μεγάλῳ.

\vs{4}Καὶ ἦν Ἡσαΐας ἐν τῇ αὐλῇ τῇ μέσῃ, καὶ ῥῆμα Κυρίου ἐγένετο πρὸς αὐτὸν, λέγων,
\vs{5}ἐπίστρεψον, καὶ ἐρεῖς πρὸς Ἐζεκίαν τὸν ἡγούμενον τοῦ λαοῦ μου, τάδε λέγει Κύριος ὁ Θεὸς Δαυὶδ τοῦ πατρός σου, ἤκουσα τῆς προσευχῆς σου, εἶδον τὰ δάκρυά σου· ἰδοὺ ἐγὼ ἰάσομαί σε· τῇ ἡμέρᾳ τῇ τρίτῃ ἀναβήσῃ εἰς οἶκον Κυρίου.
\vs{6}Καὶ προσθήσω ἐπὶ τὰς ἡμέρας σου πεντεκαίδεκα ἔτη· καὶ ἐκ χειρὸς· βασιλέως Ἀσσυρίων σώσω σε καὶ τὴν πόλιν ταύτην, καὶ ὑπερασπιῶ ὑπὲρ τῆς πόλεως ταύτης διʼ ἐμὲ καὶ διὰ Δαυὶδ τὸν δοῦλόν μου.
\vs{7}Καὶ εἶπε, λαβέτωσαν παλάθην σύκων, καὶ ἐπιθέτωσαν ἐπὶ τὸ ἕλκος, καὶ ὑγιάσει.
\vs{8}Καὶ εἶπεν Ἐζεκίας πρὸς Ἡσαΐαν, τί τὸ σημεῖον ὅτι ἰάσεταί με Κύριος, καὶ ἀναβήσομαι εἰς οἶκον Κυρίου τῇ ἡμέρᾳ τῇ τρίτῃ;
\vs{9}Καὶ εἶπεν Ἡσαΐας, τοῦτο τὸ σημεῖον παρὰ Κυρίου, ὅτι ποιήσει Κύριος τὸν λόγον ὃν ἐλάλησε· πορεύσεται ἡ σκιὰ δέκα βαθμοὺς, ἐὰν ἐπιστρέψῃ δέκα βαθμούς.
\vs{10}Καὶ εἶπεν Ἐζεκίας, κοῦφον τὴν σκιὰν κλῖναι δέκα βαθμούς· οὐχὶ, ἀλλʼ ἐπιστραφήτω ἡ σκιὰ ἐν τοῖς ἀναβαθμοῖς δέκα βαθμοὺς εἰς τὰ ὀπίσω.
\vs{11}Καὶ ἐβόησεν Ἡσαΐας ὁ προφήτης πρὸς Κύριον, καὶ ἐπέστρεψεν ἡ σκιὰ ἐν τοῖς ἀναβαθμοῖς εἰς τὰ ὀπίσω δέκα βαθμούς.

\vs{12}Ἐν τῷ καιρῷ ἐκείνῳ ἀπέστειλε Μαρωδὰχ Βαλαδὰν υἱὸς Βαλαδὰν βασιλεὺς Βαβυλῶνος βιβλία καὶ μαναὰ πρὸς Ἐζεκίαν, ὅτι ἤκουσεν ὅτι ἠῤῥώστησεν Ἐζεκίας.
\vs{13}Καὶ ἐχάρη ἐπʼ αὐτοῖς Ἐζεκίας, καὶ ἔδειξεν αὐτοῖς ὅλον τὸν οἶκον τοῦ νεχωθὰ, τὸ ἀργύριον καὶ τὸ χρυσίον, τὰ ἀρώματα καὶ τὸ ἔλαιον τὸ ἀγαθὸν, καὶ τὸν οἶκον τῶν σκευῶν, καὶ ὅσα εὑρέθη ἐν τοῖς θησαυροῖς αὐτοῦ· οὐκ ἦν λόγος ὃν οὐκ ἔδειξεν αὐτοῖς Ἐζεκίας ἐν τῷ οἴκῳ αὐτοῦ καὶ ἐν πάσῃ τῇ ἐξουσίᾳ αὐτοῦ.

\vs{14}Καὶ εἰσῆλθεν Ἡσαΐας ὁ προφήτης πρὸς τὸν βασιλέα Ἑζεκίαν, καὶ εἶπε πρὸς αὐτόν, τί ἐλάλησαν οἱ ἄνδρες οὖτοι, καὶ πόθεν ἥκασι πρὸς σέ; καὶ εἶπεν Ἐζεκίας, ἐκ γῆς πόῤῥωθεν ἥκασιν πρὸς μὲ, ἐκ Βαβυλῶνος.
\vs{15}Καὶ εἶπε, τί εἶδον ἐν τῷ οἴκῳ σου; καὶ εἶπε πάντα ὅσα ἐν τῷ οἴκῳ μου εἶδον· οὐκ ἦν ἐν τῷ οἴκῳ μου ὃ οὐκ ἔδειξα αὐτοῖς, ἀλλὰ καὶ τὰ ἐν τοῖς θησαυροῖς μου.
\vs{16}Καὶ εἶπεν Ἡσαΐας πρὸς Ἐζεκίαν, ἄκουσον λόγον Κυρίου,
\vs{17}ἰδοὺ ἡμέραι ἔρχονται, καὶ ληφθήσεται πάντα τὰ ἐν τῷ οἴκῳ σου, καὶ ὅσα ἐθησαύρισαν οἱ πατέρες σου ἕως τῆς ἡμέρας ταύτης, εἰς Βαβυλῶνα· καὶ οὐχ ὑπολειφθήσεται ῥῆμα ὃ εἶπε Κύριος.
\vs{18}Καὶ οἱ υἱοί σου οἳ ἐξελεύσονται ἐκ σοῦ οὓς γεννήσεις, λήψεται, καὶ ἔσονται εὐνοῦχοι ἐν τῷ οἴκῳ τοῦ βασιλέως Βαβυλῶνος.
\vs{19}Καὶ εἶπεν Ἐζεκίας πρὸς Ἡσαΐαν, ἀγαθὸς ὁ λόγος Κυρίου ὃν ἐλάλησεν· ἔστω εἰρήνη ἐν ταῖς ἡμέραις μου.

\vs{20}Καὶ τὰ λοιπὰ τῶν λόγων Ἐζεκίου καὶ πᾶσα ἡ δυναστεία αὐτοῦ καὶ ὅσα ἐποίησε, τὴν κρήνην καὶ τὸν ὑδραγωγόν, καὶ εἰσήγαγε τὸ ὕδωρ εἰς τὴν πόλιν, οὐχὶ ταῦτα γεγραμμένα ἐπὶ βιβλίῳ λόγων τῶν ἡμερῶν τοῖς βασιλεῦσιν Ἰούδα;
\vs{21}Καὶ ἐκοιμήθη Ἐζεκίας μετὰ τῶν πατέρων αὐτοῦ, καὶ ἐβασίλευσε Μανασσῆς υἱὸς αὐτοῦ ἀντʼ αὐτοῦ.

\ch{21}
Υἱὸς δώδεκα ἐτῶν Μανασσὴς ἐν τῷ βασιλεύειν αὐτόν, καὶ πεντήκοντα καὶ πέντε ἔτη ἐβασίλευσεν ἐν Ἰερουσαλήμ, καὶ ὄνομα τῇ μητρὶ αὐτοῦ Ἁψιβά.
\vs{2}Καὶ ἐποίησε τὸ πονηρὸν ἐν ὀφθαλμοῖς Κυρίου, κατὰ τὰ βδελύγματα τῶν ἐθνῶν ὧν ἐξῆρε Κύριος ἀπὸ προσώπον τῶν υἱῶν Ἰσραήλ.
\vs{3}Καὶ ἐπέστρεψε καὶ ᾠκοδόμησε τὰ ὑψηλὰ ἃ κατέσπασεν Ἑζεκίας ὁ πατὴρ αὐτοῦ, καὶ ἀνέστησε θυσιαστήριον τῇ Βάαλ, καὶ ἐποίησε τὰ ἄλση καθὼς ἐποίησεν Ἀχαὰβ βασιλεὺς Ἰσραήλ, καὶ προσεκύνησε πάσῃ τῇ δυνάμει τοῦ οὐρανοῦ, καὶ ἐδούλευσεν αὐτοῖς.
\vs{4}Καὶ ᾠκοδόμησε θυσιαστήριον ἐν ὄκῳ Κυρίου, ὡς εἶπεν, ἐν Ἰερουσάλὴμ θήσω τὸ ὄνομά μου.
\vs{5}Καὶ ᾠκοδόμησε θυσιαστήριον πάσῃ τῇ δυνάμει τοῦ οὐρανοῦ ἐν ταῖς δυσὶν αὐλαῖς οἴκου Κυρίου.
\vs{6}Καὶ διῆγε τοὺς υἱοὺς αὐτοῦ ἐν πυρί, καὶ ἐκληδονίζετο καὶ οἰωνίζετο, καὶ ἐποίησε τεμένη, καὶ γνώστας ἐπλήθυνε τοῦ ποιεῖν τὸ πονηρὸν ἐν ὀφθαλμοῖς Κυρίου παροργίσαι αὐτόν.
\vs{7}Καὶ ἔθηκε τὸ γλυπτὸν τοῦ ἄλσους ἐν τῷ οἴκῳ ᾧ εἶπε Κύριος πρὸς Δαυὶδ καὶ πρὸς Σαλωμὼν τὸν υἱὸν αὐτοῦ, ἐν τῷ οἴκῳ τούτῳ καὶ ἐν Ἰερουσαλὴμ ᾗ ἐξελεξάμην ἐκ πασῶν φυλῶν τοῦ Ἰσραὴλ καὶ θήσω τὸ ὄνομά μου εἰς τὸν αἰῶνα,
\vs{8}καὶ οὐ προσθήσω τοῦ σαλεῦσαι τὸν πόδα Ἰσραὴλ ἀπὸ τῆς γῆς ἧς ἔδωκα τοῖς πατράσιν αὐτῶν, οἵ τινες φυλάξουσι πάντα ὅσα ἐνετειλάμην κατὰ πᾶσαν τὴν ἐντολὴν ἣν ἐνετείλατο αὐτοῖς ὁ δοῦλός μου Μωυσῆς.
\vs{9}Καὶ οὐκ ἤκουσαν, καὶ ἐπλάνησεν αὐτοὺς Μανασσῆς τοῦ ποιῆσαι τὸ πονηρὸν ἐν ὀφθαλμοῖς Κυρίου ὑπὲρ τὰ ἔθνη, ἃ ἠθάνισε Κύριος ἐκ προσώπου υἱῶν Ἰσραήλ.

\vs{10}Καὶ ἐλάλησε Κύριος ἐν χειρὶ δούλων αὐτοῦ τῶν προφητῶν, λέγων,
\vs{11}ἀνθʼ ὧν ὅσα ἐποίησε Μανασσῆς ὁ βασιλεὺς Ἰούδα τὰ βδελύγματα ταῦτα τὰ πονηρὰ ἀπὸ πάντων ὧν ἐποίησεν ὁ Ἀμοῤῥαῖος ὁ ἔμπροσθεν, καὶ ἐξήμαρτε καί γε Ἰούδαν ἐν τοῖς εἰδώλοις αὐτῶν,
\vs{12}οὐχ οὕτως· τάδε λέγει Κύριος ὁ Θεὸς Ἰσραήλ, ἰδοὺ ἐγὼ φέρω κακὰ ἐπὶ Ἰερουσαλὴμ καὶ Ἰούδαν, ὥστε παντὸς ἀκούοντος ἠχήσει ἀμφότερα τὰ ὦτα αὐτοῦ·
\vs{13}Καὶ ἐκτενῶ ἐπὶ Ἱερουσαλὴμ τὸ μέτρον Σαμαρείας καὶ τὸ στάθμιον οἴκου Ἀχαάβ· καὶ ἀπαλείψω τὴν Ἱερουσαλὴμ, καθὼς ἀπαλείφεται ὁ ἀλάβαστρος ἀπαλειφόμενος καὶ καταστρέφεται ἐπὶ πρόσωπον αὐτοῦ.
\vs{14}Καὶ ἀπεώσομαι τὸ ὑπόλειμμα τῆς κληρονομίας μου, καὶ παραδώσω αὐτοὺς εἰς χεῖρας ἐχθρῶν αὐτῶν, καὶ ἔσονται εἰς διαρπαγὴν καὶ εἰς προνομὴν πᾶσι τοῖς ἐχθροῖς αὐτῶν,
\vs{15}ἀνθʼ ὧν ὅσα ἐποίησαν τὸ πονηρὸν ἐν ὀφθαλμοῖς μου, καὶ ἦσαν παροργίζοντές με ἀπὸ τῆς ἡμέρας ἧς ἐξήγαγον τοὺς πατέρας αὐτῶν ἐξ Αἰγύπτου καὶ ἕως τῆς ἡμέρας ταύτης.
\vs{16}Καί γε αἷμα ἀθῶον ἐξέχεε Μανασσῆς πολὺ σφόδρα ἕως οὗ ἔπλησε τὴν Ἱερουσαλὴμ στόμα εἰς στόμα, πλὴν ἀπὸ τῶν ἁμαρτιῶν αὐτοῦ ὧν ἐξήμαρτε τὸν Ἰούδαν τοῦ ποιῆσαι τὸ πονηρὸν ἐν ὀφθαλμοῖς Κυρίου.

\vs{17}Καὶ τὰ λοιπὰ τῶν λόγων Μανασσῆ καὶ πάντα ὅσα ἐποίησε, καὶ ἡ ἁμαρτία αὐτοῦ ἣν ἥμαρτεν, οὐχὶ ταῦτα γεγραμμένα ἐπὶ βιβλίῳ λόγων τῶν ἡμερῶν τοῖς βασιλεῦσιν Ἰούδα;
\vs{18}Καὶ ἐκοιμήθη Μανασσὴς μετὰ τῶν πατέρων αὐτοῦ, καὶ ἐτάφη ἐν τῷ κήπῳ τοῦ οἴκου αὐτοῦ ἐν κήπῳ Ὀζά· καὶ ἐβασίλευσεν Ἀμὼς υἱὸς αὐτοῦ ἀντʼ αὐτοῦ.

\vs{19}Υἱὸς εἴκοσι καὶ δύο ἐτῶν Ἀμὼς ἐν τῷ βασιλεύειν αὐτόν, καὶ δύο ἔτη ἐβασίλευσεν ἐν Ἱερουσαλήμ, καὶ ὄνομα τῇ μητρὶ αὐτοῦ Μεσολλάμ, θυγάτηρ Ἀροὺς ἐξ Ἰετέβα.
\vs{20}Καὶ ἐποίησε τὸ πονηρὸν ἐν ὀφθαλμοῖς Κυρίου καθὼς ἐποίησε Μανασσὴς, ὁ πατὴρ αὐτοῦ.
\vs{21}Καὶ ἐπορεύθη ἐν πάσῃ ὁδῷ ᾗ ἐπορεύθη ὁ πατὴρ αὐτοῦ, καὶ ἐλάτρευσε τοῖς εἰδώλοις οἷς ἐλάτρευσεν ὁ πατὴρ αὐτοῦ, καὶ προσεκύνησεν αὐτοῖς.
\vs{22}Καὶ ἐγκατέλιπε τὸν Κύριον Θεὸν τῶν πατέρων αὐτοῦ, καὶ οὐκ ἐπορεύθη ἐν ὁδῷ Κυρίου.
\vs{23}Καὶ συνεστράφησαν οἱ παῖδες Ἀμὼς πρὸς αὐτὸν, καὶ ἐθανάτωσαν τὸν βασιλέα ἐν τῷ οἴκῳ αὐτοῦ.
\vs{24}Καὶ ἐπάταξεν ὁ λαὸς τῆς γῆς πάντας τοὺς συστραφέντας ἐπὶ τὸν βασιλέα Ἀμώς, καὶ ἐβασίλευσεν ὁ λαὸς τῆς γῆς τὸν Ἰωσίαν υἱὸν αὐτοῦ ἀντʼ αὐτοῦ.

\vs{25}Καὶ τὰ λοιπὰ τῶν λόγων Ἀμὼς ὅσα ἐποίησεν, οὐκ ἰδοὺ ταῦτα γεγραμμένα ἐπὶ βιβλίῳ λόγων τῶν ἡμερῶν τοῖς βασιλεῦσιν Ἰούδα;
\vs{26}Καὶ ἔθαψαν αὐτὸν ἑν τῷ τάφῳ αὐτοῦ ἑν τῷ κήτῳ Ὀζά, καὶ ἐβασίλευσεν Ἰωσίας υἱὸς αὐτοῦ ἀντʼ αὐτοῦ.

\ch{22}
Υἱὸς ὀκτὼ ἐτῶν Ἰωσείας ἐν τῷ βασιλεύειν αὐτὸν, καὶ τριάκοντα καὶ ἓν ἔτος ἐβασίλευσεν ἐν Ἰερουσαλὴμ, καὶ ὄνομα τῇ μητρὶ αὐτοῦ Ἰεδία, θυλάτηρ Ἐδεϊὰ ἐκ Βασουρώθ.
\vs{2}Καὶ ἐποίησε τὸ εὐθὲς ἐν ὀφθαλμοῖς Κυρίου, καὶ ἐπορεύθη ἐν πάσῃ ὁδῷ Δαυὶδ τοῦ πατρὸς αὐτοῦ, οὐκ ἀπέστη δεξιὰ καὶ ἀριστερά.

\vs{3}Καὶ ἐγενήθη ἐν τῷ ὀκτωκαιδεκάτῳ ἔτει τῷ βασιλεῖ Ἰωσείᾳ, ἐν τῷ μηνὶ τῷ ὀγδόῳ ἀπέστειλεν ὁ βασιλεὺς τὸν Σαπφὰν υἱὸν Ἐζελίου υἱοῦ Μεσολλὰμ τὸν γραμματέα οἴκου Κυρίου, λέγων,
\vs{4}ἀνάβηθι πρὸς Χελκίαν τὸν ἱερέα τὸν μέγαν, καὶ σφράγισον τὸ ἀργύριον τὸ εἰσενεχθὲν ἐν οἴκῳ Κυρίου, ὃ συνήγαγον οἱ φυλάσσοντες τὸν σταθμὸν παρὰ τοῦ λαοῦ.
\vs{5}Καὶ δότωσαν αὐτὸ ἐπὶ χεῖρα ποιούντων τὰ ἔργα τῶν καθεσταμένων ἐν οἴκῳ Κυρίου· καὶ ἔδωκεν αὐτὸ τοῖς ποιοῦσι τὰ ἔργα τοῖς ἐν οἴκῳ Κυρίου τοῦ κατισχῦσαι τὸ βεδὲκ τοῦ οἴκου,
\vs{6}τοῖς τέκτοσι καὶ τοῖς οἰκοδόμοις καὶ τοῖς τειχισταῖς, καὶ τοῦ κτήσασθαι ξύλα καὶ λίθους λατομητοὺς, τοῦ κραταιῶσαι τὸ βεδὲκ τοῦ οἴκου.
\vs{7}Πλὴν οὐκ ἐξελογίζοντο αὐτοὺς τὸ ἀργύριον τὸ διδόμενον αὐτοῖς, ὅτι ἐν πίστει αὐτοὶ ποιοῦσι.

\vs{8}Καὶ εἶπε Χελκίας ὁ ἱερεὺς ὁ μέγας πρὸς Σαπφὰν τὸν γραμματέα, βιβλίον τοῦ νόμου εὗρον ἐν οἴκῳ Κυρίου· καὶ ἔδωκε Χελκίας τὸ βιβλίον πρὸς Σαπφὰν, καὶ ἀνέγνω αὐτό.
\vs{9}Καὶ εἰσῆλθεν ἐν οἴκῳ Κυρίου πρὸς τὸν βασιλέα, καὶ ἀπέστρεψε τῷ βασιλεῖ ῥῆμα, καὶ εἶπεν, ἐχώνευσαν οἱ δοῦλοί σου τὸ ἀργύριον τὸ εὑρεθὲν ἐν οἴκῳ Κυρίου, καὶ ἔδωκαν αὐτὸ ἐπὶ χεῖρα ποιούντων τὰ ἔργα καθεσταμένων ἐν οἴκῳ Κυρίου.
\vs{10}Καὶ εἶπε Σαπφὰν ὁ γραμματεὺς πρὸς τὸν βασιλέα, λέγων, βιβλίον ἔδωκέ μοι Χελκίας ὁ ἱερεύς· καὶ ἀνέγνω αὐτὸ Σαπφὰν ἐνώπιον τοῦ βασιλέως.
\vs{11}Καὶ ἐγένετο ὡς ἤκουσεν ὁ βασιλεὺς τοὺς λόγους βιβλίου τοῦ νόμου, καὶ διέῤῥηξε τὰ ἱμάτια ἑαυτοῦ.
\vs{12}Καὶ ἐνετείλατο ὁ βασιλεὺς τῷ Χελκίᾳ τῷ ἱερεῖ, καὶ τῷ Ἀχικὰμ υἱῷ Σαπφὰν, καὶ τῷ Ἀχοβὼρ υἱῷ Μιχαίου, καὶ τῷ Σαπφὰν τῷ γραμματεῖ, καὶ τῷ Ἀσαΐᾳ δούλῳ τοῦ βασιλέως, λέγων,
\vs{13}δεῦτε, ἐκζητήσατε τὸν Κύριον περὶ ἐμοῦ, καὶ περὶ παντὸς τοῦ λαοῦ, καὶ περὶ παντὸς τοῦ Ἰούδα, καὶ περὶ τῶν λόγων τοῦ βιβλίου τοῦ εὑρεθέντος τούτου, ὅτι μεγάλη ἡ ὀργὴ Κυρίου ἐκκεκαυμένη ἐν ἡμῖν, ὑπὲρ οὗ οὐκ ἤκουσαν οἱ πατέρες ἡμῶν τῶν λόγων τοῦ βιβλίου τούτου τοῦ ποιεῖν κατὰ πάντα τὰ γεγραμμένα καθʼ ἡμῶν.

\vs{14}Καὶ ἐπορεύθη Χελκίας ὁ ἱερεὺς, καὶ Ἀχικὰμ, καὶ Ἀχοβὼρ, καὶ Σαπφὰν, καὶ Ἀσαΐας πρὸς Ὄλδαν τὴν προφῆτιν μητέρα Σελλὴμ υἱοῦ Θεκουαὺ υἱοῦ Ἀρὰς τοῦ ἱματιοφύλακος· καὶ αὕτη κατῴκει ἐν Ἰερουσαλὴμ ἐν τῇ Μασενᾷ· καὶ ἐλάλησαν πρὸς αὐτήν.

\vs{15}Καὶ εἶπεν αὐτοῖς, τάδε λέγει Κύριος ὁ θεὸς Ἰσραὴλ, εἴπατε τῷ ἀνδρὶ τῷ ἀποστείλαντι ὑμᾶς πρὸς μέ,
\vs{16}τάδε λέγει Κύριος, ἰδοὺ ἐγὼ ἐπάγω κακὰ ἐπὶ τὸν τόπον τοῦτον, καὶ ἐπὶ τοὺς ἐνοικοῦντας αὐτὸν πάντας τοὺς λόγους τοῦ βιβλίου οὓς ἀνέγνω βασιλεὺς Ἰούδα,
\vs{17}ἀνθʼ ὧν ἐγκατέλιπόν με, καὶ ἐθυμίων θεοῖς ἑτέροις, ὅπως παροργίσωσί με ἐν τοῖς ἔργοις τῶν χειρῶν αὐτῶν, καὶ ἐκκαυθήσεται θυμός μου ἐν τῷ τόπῳ τούτῳ καὶ οὐ σβεσθήσεται.
\vs{18}Καὶ πρὸς βασιλέα Ἰούδα τὸν ἀποστείλαντα ὑμᾶς ἐπιζητῆσαι τὸν Κύριον, τάδε ἐρεῖτε πρὸς αὐτὸν, τάδε λέγει Κύριος ὁ Θεὸς Ἰσραὴλ, οἱ λόγοι οὓς ἤκουσας,
\vs{19}ἀνθʼ ὧν ὅτι ἡπαλύνθη ἡ καρδία σου, καὶ ἐνετράπης ἀπὸ προσώπου, ὡς ἤκουσας ὅσα ἐλάλησα ἐπὶ τὸν τόπον τοῦτον καὶ ἐπὶ τοὺς ἐνοικοῦντας αὐτὸν, τοῦ εἶναι εἰς ἀφανισμὸν καὶ εἰς κατάραν, καὶ διέῤῥηξας τὰ ἱμάτιά σου καὶ ἔκλαυσας ἐνώπιόν μου, καί γε ἐγὼ ἤκουσα, λέγει Κύριος.
\vs{20}Οὐχ οὕτως· ἰδοὺ προστίθημί σε πρὸς τοὺς πατέρας σου, καὶ συναχθήσῃ εἰς τὸν τάφον σου ἐν εἰρήνῃ, καὶ οὐκ ὀφθήσεται ἐν τοῖς ὀφθαλμοῖς σου ἐν πᾶσι τοῖς κακοῖς οἷς ἐγώ εἰμι ἐπάγω ἐπὶ τὸν τόπον τοῦτον.

\ch{23}
Καὶ ἐπέστρεψαν τῷ βασιλεῖ τὸ ῥῆμα· καὶ ἀπέστειλεν ὁ βασιλεὺς, καὶ συνήγαγε πρὸς ἑαυτὸν πάντας τοὺς πρεσβυτέρους Ἰούδα καὶ Ἰερουσαλήμ.
\vs{2}Καὶ ἀνέβη ὁ βασιλεὺς εἰς οἶκον Κυρίου, καὶ πᾶς ἀνὴρ Ἰούδα καὶ πάντες οἱ κατοικοῦντες ἐν Ἰερουσαλὴμ μετʼ αὐτοῦ, καὶ οἱ ἱερεῖς, καὶ οἱ προφῆται, καὶ πᾶς ὁ λαὸς ἀπὸ μικροῦ καὶ ἕως μεγάλου, καὶ ἀνέγνω ἐν ὠσὶν αὐτῶν πάντας τοὺς λόγους τοῦ βιβλίου τῆς διαθήκης τοῦ εὑρεθέντος ἐν οἴκῳ Κυρίου.
\vs{3}Καὶ ἔστη ὁ βασιλεὺς πρὸς τὸν στύλον, καὶ διέθετο διαθήκην ἐνώπιον Κυρίου, τοῦ πορεύεσθαι ὀπίσω Κυρίου, τοῦ φυλάσσειν τὰς ἐντολὰς αὐτοῦ, καὶ τὰ μαρτύρια αὐτοῦ, καὶ τὰ δικαιώματα αὐτοῦ ἐν πάσῃ καρδίᾳ καὶ ἐν πάσῃ ψυχῇ, τοῦ ἀναστῆσαι τοὺς λόγους τῆς διαθήκης ταύτης, τὰ γεγραμμένα ἐπὶ τὸ βιβλίον τοῦτο· καὶ ἔστη πᾶς ὁ λαὸς ἐν τῇ διαθήκῃ.

\vs{4}Καὶ ἐνετείλατο ὁ βασιλεὺς τῷ Χελκίᾳ τῷ ἱερεῖ τῷ μεγάλῳ καὶ τοῖς ἱερεῦσι τῆς δευτερώσεως καὶ τοῖς φυλάσσουσι τὸν σταθμὸν, ἐξαγαγεῖν ἐκ τοῦ ναοῦ Κυρίου πάντα τὰ σκεύη τὰ πεποιημένα τῷ Βάαλ καὶ τῷ ἄλσει καὶ πάσῃ τῇ δυνάμει τοῦ οὐρανοῦ· καὶ κατέκαυσεν αὐτὰ ἔξω Ἰερουσαλὴμ ἐν σαλημὼθ Κέδρών, καὶ ἔλαβεν τὸν χοῦν αὐτῶν εἰς Βαιθήλ.
\vs{5}Καὶ κατέκαυσε τοὺς χωμαρὶμ οὓς ἔδωκαν βασιλεῖς Ἰούδα, καὶ ἐθυμίων ἐν τοῖς ὑψηλοῖς καὶ ἐν ταῖς πόλεσιν Ἰούδα καὶ τοῖς περικύκλῳ Ἰερουσαλήμ, καὶ τοὺς θυμιῶντας τῷ Βάαλ καὶ τῷ ἡλίῳ, καὶ τῇ σελήνῃ, καὶ τοῖς μαζουρὼθ, καὶ πάσῃ τῇ δυνάμει τοῦ οὐρανοῦ.

\vs{6}Καὶ ἐξήνεγκε τὸ ἄλσος ἐξ οἴκου Κυρίου ἔξωθεν Ἰερουσαλὴμ εἰς τὸν χειμάῤῥουν Κέδρών, καὶ κατέκαυσεν αὐτὸν ἐν τῷ χειμάῤῥῳ Κεδρων, καὶ ἐλέπτυνεν εἰς χοῦν· καὶ ἔῤῥιψε τὸν χοῦν αὐτοῦ εἰς τὸν τάφον τῶν υἱῶν τοῦ λαοῦ.
\vs{7}καὶ καθεῖλε τὸν οἶκον τῶν καδησὶμ τῶν ἐν τῷ οἴκῳ Κυρίου, οὗ αἱ γυναῖκες ὕφαινον ἐκεῖ χεττιῒμ τῷ ἄλσει.
\vs{8}Καὶ ἀνήγαγεν πάντας τοὺς ἱερεῖς ἐκ πόλεων Ἰούδα, καὶ ἐμίανε τὰ ὑψηλὰ οὗ ἐθυμίασαν ἐκεῖ οἱ ἱερεῖς ἀπὸ Γαιβὰλ καὶ ἕως Βηρσάβεέ· καὶ καθεῖλε τὸν οἶκον τῶν πυλῶν τὸν παρὰ τὴν θύραν τῆς πύλης Ἰησοῦ ἄρχοντος τῆς πόλεως, τῶν ἐξ ἀριστερῶν ἀνδρὸς ἐν τῇ πύλῃ τῆς πόλεως.
\vs{9}Πλὴν οὐκ ἀνέβησαν οἱ ἱερεῖς τῶν ὑψηλῶν πρὸς τὸ θυσιαστήριον Κυρίου ἐν Ἰερουσαλήμ, ὅτι εἰ μὴ ἔφαγον ἄζυμα ἐν μέσῳ τῶν ἀδελφῶν αὐτῶν.
\vs{10}Καὶ ἐμίανε τὸν Ταφὲθ τὸν ἐν φάραγγι υἱοῦ Ἑννὸμ, τοῦ διαγαγεῖν ἄνδρα τὸν υἱὸν αὐτοῦ καὶ ἄνδρα τὴν θυγατέρα αὐτοῦ τῷ Μόλοχ ἐν πυρί.

\vs{11}Καὶ κατέκαυσε τοὺς ἵππους οὓς ἔδωκαν βασιλεῖς Ἰούδα τῷ ἡλίῳ ἐν τῇ εἰσόδῳ οἴκου Κυρίου εἰς τὸ γαζοφυλάκιον Ναθὰν βασιλέως τοῦ εὐνούχου ἐν φαρουρείμ· καὶ τὸ ἅρμα τοῦ ἡλίου κατέκαυσεν πυρί,
\vs{12}καὶ τὰ θυσιαστήρια τὰ ἐπὶ τοῦ δώματος τοῦ ὑπερῴου Ἄχαζ, ἃ ἐποίησαν βασιλεὶς Ἰούδα· καὶ τὰ θυσιαστήρια ἃ ἐποίησε Μανασσῆς ἐν ταῖς δυσὶν αὐλαῖς οἴκου Κυρίου καθεῖλεν ὁ βασιλεὺς καὶ κατέσπασεν ἐκεῖθεν, καὶ ἔῤῥιψε τὸν χοῦν αὐτῶν εἰς τὸν χειμάῤῥουν Κέδρων.
\vs{13}Καὶ τὸν οἶκον τὸν ἐπὶ πρόσωπον Ἱερουσαλὴμ τὸν ἐκ δεξιῶν τοῦ ὄρους τοῦ Μοσθὰθ, ὃν ᾠκοδόμησε Σαλωμὼν βασιλεὺς Ἰσραὴλ τῇ Ἀστάρτῃ προσοχθίσματι Σιδωνίων, καὶ τῷ Χαμὼς προσοχθίσματι Μωὰβ, καὶ τῷ Μολὸχ βδελύγματι υἱῶν Ἀμμὼν, ἐμίανεν ὁ βασιλεύς.
\vs{14}Καὶ συνέτριψε τὰς στήλας, καὶ ἐξωλόθρευσε τὰ ἄλση, καὶ ἔπλησε τοὺς τόπους αὐτῶν ὀστέων ἀνθρώπων.

\vs{15}Καί γε τὸ θυσιαστήριον τὸ ἐν Βαιθὴλ τὸ ὑψηλὸν ὃ ἐποίησεν Ἱεροβοὰμ υἱὸς Ναβὰτ, ὃς ἐξήμαρτε τὸν Ἰσραὴλ, καί γε τὸ θυσιαστήριον ἐκεῖνο τὸ ὑψηλὸν κατέσπασε, καὶ συνέτριψε τοὺς λίθους αὐτοῦ καὶ ἐλέπτυνεν εἰς χοῦν, καὶ κατέκαυσε τὸ ἄλσος.
\vs{16}Καὶ ἐξένευσεν Ἰωσίας καὶ εἶδε τοὺς τάφους τοὺς ἐκεῖ ἐν τῇ πόλει, καὶ ἀπέστειλε, καὶ ἔλαβε τὰ ὀστᾶ ἐκ τῶν τάφων, καὶ κατέκαυσεν ἐπὶ τὸ θυσιαστήριον, καὶ ἐμίανεν αὐτὸ κατὰ τὸ ῥῆμα Κυρίου ὃ ἐλάλησεν ὁ ἄνθρωπος τοῦ Θεοῦ ἐν τῷ ἑστᾶναι Ἱεροβοὰμ ἐν τῇ ἑορτῇ ἐπὶ τὸ θυσιαστήριον· καὶ ἐπιστρέψας ᾖρε τοὺς ὀφθαλμοὺς αὐτοῦ ἐπὶ τὸν τάφον τοῦ ἀνθρώπου τοῦ Θεοῦ τοῦ λαλήσαντος τοὺς λόγους τούτους.
\vs{17}Καὶ εἶπεν, τί τὸ σκόπελον ἐκεῖνο ὃ ἐγὼ ὁρῶ; καὶ εἶπον αὐτῷ οἱ ἄνδρες τῆς πόλεως, ὁ ἄνθρωπος τοῦ Θεοῦ ὁ ἐξεληλυθὼς ἐξ Ἰούδα, καὶ ἐπικαλεσάμενος τοὺς λόγους τούτους οὓς ἐπεκαλέσατο ἐπὶ τὸ θυσιαστήριον Βαιθήλ.
\vs{18}Καὶ εἶπεν, ἄφετε αὐτὸν, ἀνὴρ μὴ κινησάτωσαν τὰ ὀστᾶ αὐτοῦ· καὶ ἐῤῥύσθησαν τὰ ὀστᾶ αὐτοῦ μετὰ τῶν ὀστῶν τοῦ προφήτου τοῦ ἥκοντος ἐκ Σαμαρείας.

\vs{19}Καί γε πάντας τοὺς οἴκους τῶν ὑψηλῶν τοὺς ἐν ταῖς πόλεσιν Σαμαρείας, οὓς ἐποίησαν βασιλεῖς Ἰσραὴλ παροργίζειν Κύριον, ἀπέστησεν Ἰωσίας, καὶ ἐποίησεν ἐν αὐτοῖς πάντα τὰ ἔργα ἃ ἐποίησεν ἐν Βαιθήλ.
\vs{20}Καὶ ἐθυσίασε πάντας τοὺς ἱερεῖς τῶν ὑψηλῶν τοὺς ὄντας ἐκεῖ ἐπὶ τῶν θυσιαστηρίων, καὶ κατέκαυσε τὰ ὀστᾶ τῶν ἀνθρώπων ἐπʼ αὐτὰ, καὶ ἐπεστράφη εἰς Ἱερουσαλήμ.

\vs{21}Καὶ ἐνετείλατο ὁ βασιλεὺς παντὶ τῷ λαῷ, λέγων, ποιήσατε πάσχα τῷ Κυρίῳ Θεῷ ἡμῶν, καθὼς γέγραπται ἐπὶ βιβλίου τῆς διαθήκης ταύτης.
\vs{22}Ὅτι οὐκ ἐγενήθη τὸ πάσχα τοῦτο ἀφʼ ἡμερῶν τῶν κριτῶν οἳ ἔκρινον τὸν Ἰσραὴλ, καὶ πάσας τὰς ἡμέρας βασιλέων Ἰσραὴλ καὶ βασιλέων Ἰούδα.
\vs{23}ὅτι ἀλλʼ ἢ τῷ ὀκτωκαιδεκάτῳ ἔτει τοῦ βασιλέως Ἰωσίον ἐγενήθη τὸ πάσχα τῷ Κυρίῳ ἐν Ἰερουσαλήμ.

\vs{24}Καί γε τοὺς θελητὰς, καὶ τοὺς γνωριστὰς καὶ τὰ θεραφὶν, καὶ τὰ εἴδωλα, καὶ πάντα τὰ προσοχθίσματα τὰ γεγονότα ἐν τῇ γῇ Ἰούδα καὶ ἐν Ἱερουσαλὴμ ἐξῇρεν Ἰωσίας, ἵνα στήσῃ τοὺς λόγους τοῦ νόμου τοὺς γεγραμμένους ἐπὶ τοῦ βιβλίον, οὗ εὗρε Χελκίας ὁ ἱερεὺς ἐν οἴκῳ Κυρίου.
\vs{25}Ὅμοιος αὐτῷ οὐκ ἐγενήθη ἔμπροσθεν αὐτοῦ βασιλεὺς, ὃς ἐπέστρεψε πρὸς Κύριον ἐν ὅλῃ καρδίᾳ αὐτοῦ, καὶ ἐν ὅλῃ ψυχῇ αὐτοῦ, καὶ ἐν ὅλῃ ἰσχύϊ αὐτοῦ κατὰ πάντα τὸν νόμον Μωυσῆ, καὶ μετʼ αὐτὸν οὐκ ἀνέστη ὅμοιος αὐτῷ.
\vs{26}Πλὴν οὐκ ἀπεστράφη Κύριος ἀπὸ θυμοῦ τῆς ὀργῆς αὐτοῦ τῆς μεγάλης οὗ ἐθυμώθη ὀργῇ αὐτοῦ ἐν τῷ Ἰούδᾳ ἐπὶ τοὺς παροργισμοὺς, οὓς παρώργισεν αὐτὸν Μανασσῆς.
\vs{27}Καὶ εἶπε Κύριος, καί γε τὸν Ἰούδα ἀποστήσω ἀπὸ τοῦ προσώπου μου, καθὼς ἀπέστησα τὸν Ἰσραὴλ, καὶ ἀπεώσομαι τὴν πόλιν ταύτην ἣν ἐξελεξάμην, τὴν Ἱερουσαλὴμ, καὶ τὸν οἶκον οὗ εἶπον, ἔσται τὸ ὄνομά μου ἐκεῖ.
\vs{28}Καὶ τὰ λοιπὰ τῶν λόγων Ἰωσίου καὶ πάντα ὅσα ἐποίησεν, οὐχὶ ταῦτα γεγραμμένα ἐπὶ βιβλίῳ λόγων ἡμερῶν τοῖς βασιλεῦσιν Ἰούδα;

\vs{29}Ἐν δὲ ταῖς ἡμέραις αὐτοῦ ἀνέβη Φαραὼ Νεχαὼ βασιλεὺς Αἰγύπτου ἐπὶ βασιλέα Ἀσσυρίων ἐπὶ ποταμὸν Εὐφράτην· καὶ ἐπορεύθη Ἰωσίας εἰς ἀπαντὴν αὐτοῦ, καὶ ἐθανάτωσεν αὐτὸν Νεχαὼ ἐν Μαγεδδὼ ἐν τῷ ἰδεῖν αὐτόν.
\vs{30}Καὶ ἐπεβίβασαν αὐτὸν οἱ παῖδες αὐτοῦ νεκρὸν ἐκ Μαγεδδὼ, καὶ ἤγαγον αὐτὸν εἰς Ἱερουσαλὴμ, καὶ ἔθαψαν αὐτὸν ἐν τῷ τάφῳ αὐτοῦ· καὶ ἔλαβεν ὁ λαὸς τῆς γῆς τὸν Ἰωάχαζ υἱὸν Ἰωσίου, καὶ ἔχρισαν αὐτὸν, καὶ ἐβασίλευσαν αὐτὸν ἀντὶ τοῦ πατρὸς αὐτοῦ.

\vs{31}Υἱὸς εἴκοσι καὶ τριῶν ἐτῶν ἦν Ἰωάχαζ ἐν τῷ βασιλεύειν αὐτὸν, καὶ τρίμηνον ἐβασίλευσεν ἐν Ἱερουσαλὴμ, καὶ ὄνομα τῇ μητρὶ αὐτοῦ Ἀμιτὰλ, θυγάτηρ Ἱερεμίου ἐκ Λοβνά.
\vs{32}Καὶ ἐποίησε τὸ πονηρὸν ἐν ὀφθαλμοῖς Κυρίου, κατὰ πάντα ὅσα ἐποίησαν οἱ πατέρες αὐτοῦ.
\vs{33}Καὶ μετέστησεν αὐτὸν Φαραὼ Νεχαὼ ἐν Ῥαβλαὰμ ἐν γῇ Ἐμὰθ τοῦ μὴ βασιλεύειν ἐν Ἱερουσαλὴμ, καὶ ἔδωκε ζημίαν ἐπὶ τὴν γῆν ἑκατὸν τάλαντα ἀργυρίου καὶ ἑκατὸν τάλαντα χρυσίου.
\vs{34}Καὶ ἐβασίλευσε Φαραὼ Νεχαὼ ἐπʼ αὐτοὺς τὸν Ἐλιακὶμ υἱὸν Ἰωσίου βασιλέως Ἰούδα ἀντὶ Ἰωσίου τοῦ πατρὸς αὐτοῦ· καὶ ἐπέστρεψε τὸ ὄνομα αὐτοῦ Ἰωακίμ· καὶ τὸν Ἰωάχαζ ἔλαβε καὶ εἰσήνεγκεν εἰς Αἴγυπτον, καὶ ἀπέθανεν ἐκεῖ.
\vs{35}Καὶ τὸ ἀργύριον καὶ τὸ χρυσίον ἔδωκεν Ἰωακὶμ τῷ Φαραῷ, πλὴν ἐτιμογράφησε τὴν γῆν τοῦ δοῦναι τὸ ἀργύριον ἐπὶ στόματος Φαραώ ἀνὴρ κατὰ τὴν συντίμησιν αὐτοῦ ἔδωκαν τὸ ἀργύριον καὶ τὸ χρυσίον μετὰ τοῦ λαοῦ τῆς γῆς τοῦ δοῦναι τῷ Φαραῷ Νεχαώ.

\vs{36}Υἱὸς εἴκοσι καὶ πέντε ἐτῶν Ἰωακὶμ ἐν τῷ βασιλεύειν αὐτὸν, καὶ ἕνδεκα ἔτη ἐβασίλευσεν ἐν Ἱερουσαλὴμ, καὶ ὄνομα τῇ μητρὶ αὐτοῦ Ἰελδὰφ θυγάτηρ Φαδαῒλ ἐκ Ῥουμά.
\vs{37}Καὶ ἐποίησε τὸ πονηρὸν ἐν ὀφθαλμοῖς Κυρίου, κατὰ πάντα ὅσα ἐποίησαν οἱ πατέρες αὐτοῦ.

\ch{24}
Ἐν ταῖς ἡμέραις αὐτοῦ ἀνέβη Ναβουχοδονόσορ βασιλεὺς Βαβυλῶνος, καὶ ἐγενήθη αὐτῷ Ἰωακὶμ δοῦλος τρία ἔτη· καὶ ἐπέστρεψε καὶ ἠθέτησεν ἐν αὐτῷ.
\vs{2}Καὶ ἀπέστειλε Κύριος αὐτῷ τοὺς μονοζώνους τῶν Χαλδαίων, καὶ τοὺς μονοζώνους Συρίας, καὶ τοὺς μονοζώνους Μωὰβ, καὶ τοὺς μονοζώνους υἱῶν Ἀμμὼν, καὶ ἐξαπέστειλεν αὐτοὺς ἐν τῇ γῇ Ἰούδα τοῦ κατισχῦσαι κατὰ τὸν λόγον Κυρίου, ὃν ἐλάλησεν ἐν χειρὶ τῶν δούλων αὐτοῦ τῶν προφητῶν.
\vs{3}Πλὴν ἐπὶ τὸν θυμὸν Κυρίου ἦν ἐν τῷ Ἰούδα, ἀποστῆσαι αὐτὸν ἀπὸ τοῦ προσώπου αὐτὸν ἐν ἁμαρτίαις Μανασσὴ κατὰ πάντα ὅσα ἐποίησε.
\vs{4}Καί γε τὸ αἷμα ἀθῶον, ἐξέχεε, καὶ ἔπλησε τὴν Ἰερουσαλὴμ αἵματος ἀθώου, καὶ οὐκ ἠθέλησε Κύριος ἱλασθῆναι.
\vs{5}Καὶ τὰ λοιπὰ τῶν λόγων Ἰωακὶμ καὶ πάντα ὅσα ἐποίησεν, οὐκ ἰδοὺ ταῦτα γεγραμμένα ἐπὶ βιβλίῳ λόγων τῶν ἡμερῶν τοῖς βασιλεῦσιν Ἰούδα;

\vs{6}καὶ ἐκοιμήθη Ἰωακεὶμ μετὰ τῶν πατέρων αὐτοῦ, καὶ ἐβασίλευσεν Ἰωακεὶμ υἱὸς αὐτοῦ ἀντʼ αὐτοῦ.
\vs{7}Καὶ οὐ προσέθετο ἔτι βασιλεὺς Αἰγύπτου ἐξελθεῖν ἐκ τῆς γῆς αὐτοῦ, ὅτι ἔλαβε βασιλεὺς Βαβυλῶνος ἀπὸ τοῦ χειμάῤῥου Αἰγύπτου ἕως τοῦ ποταμοῦ Εὐφράτου πάντα ὅσα ἦν τοῦ βασιλέως Αἰγύπτου.

\vs{8}Υἱὸς ὀκτωκαίδεκα ἐτῶν Ἰωαχὶμ ἐν τῷ βασιλεύειν αὐτὸν, καὶ τρίμηνον ἐβασίλευσεν ἐν Ἱερουσαλὴμ, καὶ ὄνομα τῇ μητρὶ αὐτοῦ Νέσθα, θυγάτηρ Ἐλλαναθὰμ, ἐξ Ἱερουσαλήμ.
\vs{9}Καὶ ἐποίησε τὸ πονηρὸν ἐν ὀφθαλμοῖς Κυρίου, κατὰ πάντα ὅσα ἐποίησεν ὁ πατὴρ αὐτοῦ.

\vs{10}Ἐν τῷ καιρῷ ἐκείνῳ ἀνέβη Ναβουχοδονόσορ βασιλεὺς Βαβυλῶνος εἰς Ἱερουσαλὴμ, καὶ ἦλθεν ἡ πόλις ἐν περιοχῇ.
\vs{11}Καὶ εἰσῆλθε Ναβουχοδονόσορ βασιλεὺς Βαβυλῶνος εἰς πόλιν, καὶ οἱ παῖδες αὐτοῦ ἐπολιόρκουν ἐπʼ αὐτήν.
\vs{12}Καὶ ἐξῆλθεν Ἰωαχὶμ βασιλεὺς Ἰούδα ἐπὶ βασιλέα Βαβυλῶνος, αὐτὸς καὶ οἱ παῖδες αὐτοῦ, καὶ ἡ μήτηρ αὐτοῦ, καὶ οἱ ἄρχοντες αὐτοῦ, καὶ οἱ εὐνοῦχοι αὐτοῦ· καὶ ἔλαβεν αὐτὸν βασιλεὺς Βαβυλῶνος ἐν τῷ ὀγδόῳ ἔτει τῆς βασιλείας αὐτοῦ.
\vs{13}Καὶ ἐξήνεγκεν ἐκεῖθεν πάντας τοὺς θησαυροὺς οἴκου Κυρίου, καὶ τοὺς θησαυροὺς οἴκου τοῦ βασιλέως, καὶ συνέκοψε πάντα τὰ σκεύη τὰ χρυσᾶ ἃ ἐποίησε Σαλωμὼν ὁ βασιλεὺς Ἰσραὴλ ἐν τῷ ναῷ Κυρίου κατὰ τὸ ῥῆμα Κυρίου.
\vs{14}Καὶ ἀπῴκισε τὴν Ἱερουσαλὴμ καὶ πάντας τοὺς ἄρχοντας καὶ τοὺς δυνατοὺς ἰσχύϊ αἰχμαλωσίας δέκα χιλιάδας αἰχμαλωτίσας, καὶ πᾶν τέκτονα καὶ τὸν συγκλείοντα, καὶ οὐχ ὑπελείφθη πλὴν οἱ πτωχοὶ τῆς γῆς.
\vs{15}Καὶ ἀπῴκισε τὸν Ἰωαχὶμ εἰς Βαβυλῶνα, καὶ τὴν μητέρα τοῦ βασιλέως, καὶ τὰς γυναῖκας τοῦ βασιλέως, καὶ τοὺς εὐνούχους αὐτοῦ· καὶ τοὺς ἰσχυροὺς τῆς γῆς ἀπήγαγεν εἰς ἀποικεσίαν ἐξ Ἱερουσαλὴμ εἰς Βαβυλῶνα·
\vs{16}Καὶ πάντας τοὺς ἄνδρας τῆς δυνάμεως ἑπτακισχιλίους, καὶ τὸν τέκτονα καὶ τὸν συγκλείοντα χιλίους· πάντες δυνατοὶ ποιοῦντες πόλεμον· καὶ ἤγαγεν αὐτοὺς βασιλεὺς Βαβυλῶνος μετοικεσίαν εἰς Βαβυλῶνα.
\vs{17}Καὶ ἐβασίλευσε βασιλεὺς Βαβυλῶνος τὸν Βατθανίαν υἱὸν αὐτοῦ ἀντʼ αὐτοῦ, καὶ ἐπέθηκε τὸ ὄνομα αὐτοῦ, Σεδεκία.

\vs{18}Υἱὸς εἴκοσι καὶ ἑνὸς ἑνιαυτῶν Σεδεκίας ἐν τῷ βασιλεύειν αὐτὸν, καὶ ἔνδεκα ἔτη ἐβασίλευσεν ἐν Ἱερουσαλὴμ, καὶ ὄνομα τῇ μητρὶ αὐτοῦ Ἀμιτὰλ, θυγὰτηρ Ἱερεμίου.
\vs{19}Καὶ ἐποίησε τὸ πονηρὸν ἐνώπιον Κυρίου, κατὰ πάντα ὅσα ἐποίησεν Ἰωακίμ.
\vs{20}Ὅτι ἐπὶ τὸν θυμὸν Κυρίου ἦν ἐπὶ Ἱερουσαλὴμ καὶ ἐν τῷ Ἰούδα, ἕως ἀπέῤῥιψεν αὐτοὺς ἀπὸ προσώπου αὐτοῦ· καὶ ἠθέτησε Σεδεκίας ἐν τῷ βασιλεῖ Βαβυλῶνος.

\ch{25}
Καὶ ἐγενήθη ἐν τῷ ἔτει τῷ ἐννάτῳ τῆς βασιλείας αὐτοῦ ἐν τῷ μηνὶ τῷ δεκάτῳ, ἦλθε Ναβουχοδονόσορ ὁ βασιλεὺς Βαβυλῶνος, καὶ πᾶσα ἡ δύναμις αὐτοῦ ἐπὶ Ἱερουσαλήμ· καὶ παρενέβαλεν ἐπʼ αὐτὴν, καὶ ᾠκοδόμησεν ἐπʼ αὐτὴν περίτειχος κύκλῳ.
\vs{2}Καὶ ἦλθεν ἡ πόλις ἐν περιοχῇ ἕως τοῦ ἑνδεκάτου ἔτους τοῦ βασιλέως Σεδεκίου ἐννάτῃ τοῦ μηνός.
\vs{3}Καὶ ἐνίσχυσεν ὁ λιμὸς ἐν τῇ πόλει, καὶ οὐκ ἦσαν ἄρτοι τῷ λαῷ τῆς γῆς.
\vs{4}Καὶ ἐῤῥάγη ἡ πόλις, καὶ πάντες οἱ ἄνδρες τοῦ πολέμου ἐξῆλθον νυκτὸς ὁδὸν πύλης τῆς ἀναμέσον τῶν τειχῶν, αὕτη ἐστὶ τοῦ κήπου τοῦ βασιλέως, καὶ οἱ Χαλδαῖοι ἐπὶ τὴν πόλιν κύκλῳ· καὶ ἐπορεύθη ὁδὸν τὴν Ἄραβα·
\vs{5}Καὶ ἐδίωξεν ἡ δύναμις τῶν Χαλδαίων ὀπίσω τοῦ βασιλέως, καὶ κατέλαβον αὐτὸν ἐν Ἀραβὼθ Ἰεριχὼ, καὶ πᾶσα ἡ δύναμις αὐτοῦ διεσπάρη ἐπάνωθεν αὐτοῦ.
\vs{6}Καὶ συνέλαβον τὸν βασιλέα, καὶ ἤγαγον αὐτὸν πρὸς βασιλέα Βαβυλῶνος εἰς Ῥεβλαθά· καὶ ἐλάλησε μετʼ αὐτοῦ κρίσιν.
\vs{7}Καὶ τοὺς υἱοὺς Σεδεκίου ἔσφαξε κατʼ ὀφθαλμοὺς αὐτοῦ, καὶ τοὺς ὀφθαλμοὺς Σεδεκίου ἐξετύφλωσε, καὶ ἔδησεν αὐτὸν ἐν πέδαις, καὶ ἤγαγεν εἰς Βαβυλῶνα.

\vs{8}Καὶ ἐν τῷ μηνὶ τῷ πέμπτῳ ἑβδόμῃ τοῦ μηνὸς, αὐτὸς ἐνιαυτὸς ἐννεακαιδέκατος τῷ Ναβουχοδονόσορ βασιλεῖ Βαβυλῶνος, ἦλθε Ναβουζαρδὰν ὁ ἀρχιμάγειρος ἑστὼς ἐνώπιον βασιλέως Βαβυλῶνος εἰς Ἱερουσαλήμ·
\vs{9}Καὶ ἐνέπρησε τὸν οἶκον Κυρίου, καὶ τὸν οἶκον τοῦ βασιλέως, καὶ πάντας τοὺς οἴκους Ἱερουσαλὴμ, καὶ πᾶν οἶκον ἐνέπρησεν ὁ ἀρχιμάγειρος.
\vs{10}Καὶ τὸ τεῖχος Ἱερουσαλὴμ κυκλόθεν κατέσπασεν ἡ δύναμις τῶν χαλδαίων.
\vs{11}Καὶ τὸ περισσὸν τοῦ λαοῦ τὸ καταλειφθὲν ἐν τῇ πόλει, καὶ τοὺς ἐμπεπτωκότας οἳ ἐνέπεσον πρὸς τὸν βασιλέα Βαβυλῶνος, καὶ τὸ λοιπὸν τοῦ στηρίγματος μετῇρε Ναβουζαρδὰν ὁ ἀρχιμάγειρος.
\vs{12}Καὶ ἀπὸ τῶν πτωχῶν τῆς γῆς ὑπέλιπεν ὁ ἀρχιμάγειρος εἰς ἀμπελουργοὺς καὶ εἰς γαβίν.

\vs{13}Καὶ τοὺς στύλους τοὺς χαλκοῦς τοὺς ἐν οἴκῳ Κυρίου, καὶ τὰς μεχωνὼθ, καὶ τὴν θάλασσαν τὴν χαλκῆν τὴν ἐν οἴκῳ Κυρίου συνέτριψαν οἱ Χαλδαῖοι, καὶ ᾖραν τὸν χαλκὸν αὐτῶν εἰς Βαβυλῶνα.
\vs{14}Καὶ τοὺς λέβητας, καὶ τὰ ἰαμὶν, καὶ τὰς φιάλας, καὶ τὰς θυΐσκας, καὶ πάντα τὰ σκεύη τὰ χαλκᾶ ἐν οἷς λειτουργοῦσιν ἐν αὐτοῖς, ἔλαβε.
\vs{15}Καὶ τὰ πυρεῖα, καὶ τὰς φιάλας τὰς χρυσᾶς καὶ τὰς ἀργυρᾶς ἔλαβεν ὁ ἀρχιμάγειρος,
\vs{16}στύλους δύο, καὶ τὴν θάλασσαν μίαν, καὶ τὰς μεχωνὼθ ἃς ἐποίησε Σαλωμὼν τῷ οἴκῳ Κυρίου· οὐκ ἦν σταθμὸς τοῦ χαλκοῦ πάντων τῶν σκευῶν.
\vs{17}Ὀκτωκαίδεκα πήχεων ὕψος τοῦ στύλου τοῦ ἑνὸς, καὶ τὸ χωθὰρ ἐπʼ αὐτοῦ τὸ χαλκοῦν· καὶ τὸ ὕψος τοῦ χωθὰρ τριῶν πήχεων· σαβαχὰ, καὶ ῥοαὶ ἐπὶ τῷ χωθὰρ κύκλῳ τὰ πάντα χαλκᾶ, καὶ κατὰ ταῦτα τῷ στύλῳ τῷ δευτέρῳ ἐπὶ τῷ σαβαχά.

\vs{18}Καὶ ἔλαβεν ὁ ἀρχιμάγειρος τὸν Σαραίαν ἱερέα τὸν πρῶτον, καὶ τὸν Σοφονίαν υἱὸν τῆς δευτερώσεως, καὶ τοὺς τρεῖς τοὺς φυλάσσοντας τὸν σταθμόν.
\vs{19}Καὶ ἐκ τῆς πόλεως ἔλαβον εὐνοῦχον ἕνα, ὅς ἦν ἐπιστάτης τῶν ἀνδρῶν τῶν πολεμιστῶν, καὶ πέντε ἄνδρας τῶν ὁρώντων τὸ πρόσωπον τοῦ βασιλέως τοὺς εὑρεθέντας ἐν τῇ πόλει, καὶ τὸν γραμματέα τοῦ ἄρχοντος τῆς δυνάμεως τὸν ἐκτάσσοντα τὸν λαὸν τῆς γῆς, καὶ ἑξήκοντα ἄνδρας τοῦ λαοῦ τῆς γῆς τοὺς εὑρεθέντας ἐν τῇ πόλει.
\vs{20}Καὶ ἔλαβεν αὐτοὺς Ναβουζαρδὰν ὁ ἀρχιμάγειρος, καὶ ἤγαγεν αὐτοὺς πρὸς τὸν βασιλέα Βαβυλῶνος εἰς Ῥεβλαθά.
\vs{21}Καὶ ἔπαισεν αὐτοὺς ὁ βασιλεὺς Βαβυλῶνος, καὶ ἐθανάτωσεν αὐτοὺς εἰς Ῥεβλαθὰ ἐν γῇ Αἰμάθ· καὶ ἀπῳκίσθη Ἰούδας ἐπάνωθεν τῆς γῆς αὐτοῦ.

\vs{22}Καὶ ὁ λαὸς ὁ καταλειφθεὶς ἐν τῇ γῇ Ἰούδα οὓς κατέλιπε Ναβουχοδονόσορ βασιλεὺς Βαβυλῶνος, καὶ κατέστησεν ἐπʼ αὐτῶν τὸν Γοδολίαν υἱὸν Ἀχικὰμ υἱὸν Σαφάν.
\vs{23}Καὶ ἤκουσαν πάντες οἱ ἄρχοντες τῆς δυνάμεως αυτοὶ καὶ οἱ ἄνδρες αὐτῶν, ὅτι κατέστησε βασιλεὺς Βαβυλῶνος τὸν Γοδολίαν, καὶ ἦλθον πρὸς Γοδολίαν εἰς Μασσηφὰθ, καὶ Ἰσμαὴλ υἱὸς Ναθανίου, καὶ Ἰωνὰ υἱὸς Καρὴθ, καὶ Σαραίας υἱὸς Θαναμὰθ ὁ Νετωφαθίτης, καὶ Ἰεζονίας υἱὸς τοῦ Μαχαθὶ, αὐτοὶ καὶ οἱ ἄνδρες αὐτῶν·
\vs{24}Καὶ ὤμοσεν Γοδολίας αὐτοῖς, καὶ τοῖς ἀνδράσιν αὐτῶν, καὶ εἶπεν αὐτοῖς, μὴ φοβεῖσθε πάροδον τῶν Χαλδαίων, καθίσατε ἐν τῇ γῇ, καὶ δουλεύσατε τῷ βασιλεῖ Βαβυλῶνος, καὶ καλῶς ἔσται ὑμῖν.

\vs{25}Καὶ ἐγενήθη ἐν τῷ ἑβδόμῳ μηνὶ ἦλθεν Ἰσμαὴλ υἱὸς Ναθανίου υἱοῦ Ἐλισαμὰ ἐκ τοῦ σπέρματος τῶν βασιλέων, καὶ δέκα ἄνδρες μετʼ αὐτοῦ, καὶ ἐπάταξε τὸν Γοδολίαν καὶ ἀπέθανε, καὶ τοὺς Ἰουδαίους καὶ τοὺς Χαλδαίους, οἳ ἦσαν μετʼ αὐτοῦ ἐν Μασσηφά.
\vs{26}Καὶ ἀνέστη πᾶς ὁ λαὸς ἀπὸ μικροῦ ἕως μεγάλου καὶ οἱ ἄρχοντες τῶν δυνάμεων, καὶ εἰσῆλθον εἰς Αἴγυπτον, ὅτι ἐφοβήθησαν ἀπὸ προσώπου τῶν Χαλδαίων.

\vs{27}Καὶ ἐγενήθη ἐν τῷ τριακοστῷ καὶ ἑβδόμῳ ἔτει τῆς ἀποικίας τοῦ Ἰωαχὶμ βασιλέως Ἰούδα, ἐν τῷ δωδεκάτῳ μηνὶ, ἑβδόμῃ καὶ εἰκάδι τοῦ μηνὸς, ὕψωσεν Εὐιαλμαρωδὲκ βασιλεὺς Βαβυλῶνος ἐν τῷ ἐνιαυτῷ τῆς βασιλείας αὐτοῦ τὴν κεφαλὴν Ἰωαχὶμ τοῦ βασιλέως Ἰούδα, καὶ ἐξήγαγεν αὐτὸν ἐξ οἴκου φυλακῆς αὐτοῦ.
\vs{28}Καὶ ἐλάλησε μετʼ αὐτοῦ ἀγαθὰ, καὶ ἔδωκε τὸν θρόνον αὐτοῦ ἐπάνωθεν τῶν θρόνων τῶν βασιλέων τῶν μετʼ αὐτοῦ ἐν Βαβυλῶνι.
\vs{29}Καὶ ἠλλοίωσε τὰ ἱμάτια τῆς φυλακῆς αὐτοῦ, καὶ ἤσθιεν ἄρτον διαπαντὸς ἐνώπιον αὐτοῦ πάσας τὰς ἡμέρας τῆς ζωῆς αὐτοῦ.
\vs{30}Καὶ ἡ ἑστιατορία αὐτοῦ ἑστιατορία διαπαντὸς ἐδόθη αὐτῷ ἐξ οἴκου τοῦ βασιλέως, λόγον ἡμέρας ἐν τῇ ἡμέρᾳ αὐτοῦ, πάσας τὰς ἡμέρας τῆς ζωῆς αὐτοῦ.


\def\book{ΠΑΡΑΛΕΙΠΟΜΕΝΩΝ Α}
\biblebook{ΠΑΡΑΛΕΙΠΟΜΕΝΩΝ Α}


\lettrine[lines=2, loversize=0.2, nindent=0em, findent=.25em]{\textcolor{bookheadingcolor}{Ἀ}}{ΔΑΜ,} Σὴθ, Ἐνὼς,
\vs{2}καὶ Καϊνᾶν, Μαλελεὴλ, Ἰάρεδ,
\vs{3}Ἐνὼχ, Μαθουσάλα, Λάμεχ,
\vs{4}Νῶε· υἱοὶ Νῶε, Σὴμ, Χὰμ, Ἰάφεθ.

\vs{5}Υἱοὶ Ἰάφεθ, Γαμὲρ, Μαγὼγ, Μαδαῒμ, Ἰωϋὰν, Ἑλισὰ, Θοβὲλ, Μοσὸχ, καὶ Θίρας.
\vs{6}Καὶ οἱ υἱοὶ Γαμὲρ, Ἀσχανὰζ, καὶ Ῥιφὰθ, καὶ Θοργαμά.
\vs{7}Καὶ οἱ υἱοὶ Ἰωϋὰν, Ἑλισὰ, καὶ Θαρσὶς, Κίτιοι, καὶ Ῥόδιοι.

\vs{8}Καὶ υἱοὶ Χὰμ, Χοὺς, καὶ Μεσραῒμ, Φοὺδ, καὶ Χαναάν.
\vs{9}Καὶ υἱοὶ Χοὺς, Σαβὰ, καὶ Εὐιλὰ, καὶ Σαβαθὰ, καὶ Ῥεγμὰ, καὶ Σεβεθαχά· καὶ υἱοὶ Ῥεγμὰ, Σαβὰ, καὶ Δαδάν.
\vs{10}Καὶ Χοὺς ἐγέννησε τὸν Νεβρώδ· οὗτος ἤρξατο εἶναι γίγας κυνηγὸς ἐπὶ τῆς γῆς.

\vs{17}Υἱοὶ Σὴμ, Αἰλὰμ, καὶ Ἀσσοὺρ,
\vs{24}καὶ Ἀρφαξὰδ, Σάλα,
\vs{25}Ἔβερ, Φαλὲγ, Ῥαγὰν,
\vs{26}Σεροὺχ, Ναχὼρ, Θάῤῥα,
\vs{27}Ἁβραάμ·

\vs{28}Υἱοὶ δὲ Ἁβραὰμ, Ἰσαὰκ, καὶ Ἰσμαήλ.
\vs{29}Αὗται δὲ αἱ γενέσεις αὐτῶν· πρωτότοκος Ἰσμαὴλ, Ναβαιὼθ, καὶ Κηδὰρ, Ναβδεὴλ, Μασσὰμ,
\vs{30}Μασμὰ, Ἰδουμὰ, Μασσὴ, Χονδὰν, Θαιμὰν,
\vs{31}Ἰεττοὺρ, Ναφὲς, Κεδμά· οὗτοι υἱοὶ Ἰσμαήλ.

\vs{32}Καὶ υἱοὶ Χεττούρας παλλακῆς Ἁβραάμ· καὶ ἔτεκεν αὐτῷ τὸν Ζεμβρὰμ, Ἰεξὰν, Μαδιὰμ, Μαδὰμ, Σοβὰκ, Σωέ· καὶ υἱοὶ Ἰεξὰν, Δαιδὰν, καὶ Σαβαΐ.
\vs{33}Καὶ υἱοὶ Μαδιὰμ, Γεφὰρ, καὶ Ὀφὲρ, καὶ Ἐνὼχ, καὶ Ἀβιδὰ, καὶ Ἐλδαδά· πάντες οὗτοι υἱοὶ Χεττούρας.

\vs{34}Καὶ ἐγέννησεν Ἁβραὰμ τὸν Ἰσαάκ· καὶ υἱοὶ Ἰσαὰκ, Ἰακὼβ, καὶ Ἡσαῦ.
\vs{35}Υἱοὶ Ἡσαῦ, Ἐλιφὰζ, καὶ Ῥαγουὴλ, καὶ Ἰεοὺλ, καὶ Ἰεγλὸμ, καὶ Κορέ.
\vs{36}Υἱοὶ Ἐλιφὰζ, Θαιμὰν, καὶ Ὠμὰρ, Σωφὰρ, καὶ Γοωθὰμ, καὶ Κενὲζ, καὶ Θαμνὰ, καὶ Ἀμαλήκ.
\vs{37}Καὶ υἱοὶ Ῥαγουὴλ, Ναχὲς, Ζαρὲ, Σομὲ, καὶ Μοζέ.
\vs{38}Υἱοὶ Σηῒρ, Λωτὰν, Σωβὰλ, Σεβεγὼν, Ἀνὰ, Δησὼν, Ὠσὰρ, καὶ Δισάν.
\vs{39}Καὶ υἱοὶ Λωτὰν, Χοῤῥὶ, καὶ Αἰμὰν· ἀδελφὴ δὲ Λωτὰν Θαμνά. Υἱοὶ Σωβὰλ, Ἀλὼν, Μαχανὰθ, Ταιβὴλ, Σωφὶ, καὶ Ὠνὰν·
\vs{40}υἱοὶ δὲ Σεβεγὼν, Ἀῒθ, καὶ Σωνάν.
\vs{41}Υἱοὶ Σωνὰν. Δαισών· υἱοὶ δὲ Δαισὼν, Ἐμερὼν, καὶ Ἀσεβὼν, καὶ Ἰεθρὰμ, καὶ Χαῤῥάν.
\vs{42}Καὶ υἱοὶ Ὡσὰρ, Βαλαὰμ, καὶ Ζουκὰμ, καὶ Ἀκάν· υἱοὶ Δισὰν, Ὢς καὶ Ἀράν.

\vs{43}Καὶ οὗτοι οἱ βασιλεῖς αὐτῶν· Βαλὰκ υἱὸς Βεὼρ, καὶ ὄνομα τῇ πόλει αὐτοῦ Δενναβά.
\vs{44}Καὶ ἀπέθανε Βαλὰκ, καὶ ἐβασίλευσεν ἀντʼ αὐτοῦ Ἰωβὰβ υἱὸς Ζαρὰ ἐκ Βοσόῤῥας.
\vs{45}Καὶ ἀπέθανεν Ἰωβὰβ καὶ ἐβασίλευσεν ἀντʼ αὐτοῦ Ἁσὸμ ἐκ γῆς Θαιμανών.
\vs{46}Καὶ ἀπέθανεν Ἁσόμ, καὶ ἐβασίλευσεν ἀντʼ αὐτοῦ Ἀδὰδ υἱὸς Βαρὰδ, ὁ πατάξας Μαδιὰμ ἐν τῷ πεδίῳ Μωὰβ· καὶ ὄνομα τῇ πόλει αὐτοῦ Γεθαίμ.
\vs{47}Καὶ ἀπέθανεν Ἁδὰδ, καὶ ἐβασίλευσεν ἀντʼ αὐτοῦ Σεβλὰ ἐκ Μασεκκάς.
\vs{48}Καὶ ἀπέθανε Σεβλὰ, καὶ ἐβασίλευσεν ἀντʼ αὐτοῦ Σαοὺλ ἐκ Ῥωβὼθ τῆς παρὰ ποταμόν.
\vs{49}Καὶ ἀπέθανε Σαοὺλ, καὶ ἐβασίλευσεν ἀντʼ αὐτοῦ Βαλαεννὼρ υἱὸς Ἀχωβώρ.
\vs{50}Καὶ ἀπέθανε Βαλαεννὼρ, καὶ ἐβασίλευσεν ἀντʼ αὐτοῦ Ἀδὰδ υἱὸς Βαρὰδ, καὶ ὄνομα τῇ πόλει αὐτοῦ, Φογώρ.

\vs{51}Ἡγεμόνες Ἐδώμ· ἡγεμὼν Θαμνὰ, ἡγεμὼν Γωλαδὰ, ἡγεμὼν Ἰεθὲρ,
\vs{52}ἡγεμὼν Ἐλιβαμὰς, ἡγεμὼν Ἠλὰς, ἡγεμὼν Φινὼν,
\vs{53}ἡγεμὼν Κενέζ, ἡγεμὼν Θαιμὰν, ἡγεμὼν Βαβσὰρ,
\vs{54}ἡγεμὼν Μαγεδιὴλ, ἡγεμὼν Ζαφωΐν· οὗτοι ἡγεμόνες Ἐδώμ.

\ch{2}
Ταῦτα τὰ ὀνόματα τῶν υἱῶν Ἰσραήλ· Ῥουβὴν, Συμεὼν, Λευὶ, Ἰούδα, Ἰσσάχαρ, Ζαβουλὼν,
\vs{2}Δὰν, Ἰωσὴφ, Βενιαμὶν, Νεφθαλὶ, Γὰδ, Ἀσήρ.

\vs{3}Υἱοὶ Ἰούδα, Ἢρ, Αὐνὰν, Σηλώμ· τρεῖς ἐγεννήθησαν αὐτῷ ἐκ τῆς θυγατρὸς Σαύας τῆς Χανανίτιδος· καὶ ἦν Ἢρ ὁ πρωτότοκος Ἰούδα πονηρὸς ἐναντίον Κυρίου, καὶ ἀπέκτεινεν αὐτόν·
\vs{4}Καὶ Θάμαρ ἡ νύμφη αὐτοῦ ἔτεκεν αὐτῷ τὸν Φαρὲς, καὶ τὸν Ζαρα· πάντες υἱοὶ Ἰούδα πέντε.

\vs{5}Υἱοὶ Φαρὲς, Ἐσρὼμ, καὶ Ἰεμουήλ.
\vs{6}Καὶ υἱοὶ Ζαρὰ, Ζαμβρὶ, καὶ Αἰθὰμ, καὶ Αἰμουὰν, καὶ Καλχὰλ, καὶ Δαρὰδ, πάντες πέντε.

\vs{7}Καὶ υἱοὶ Χαρμὶ, Ἀχὰρ ὁ ἐμποδοστάτης Ἰσραὴλ, ὃς ἠθέτησεν εἰς τὸ ἀνάθεμα.
\vs{8}Καὶ υἱοὶ Αἰθὰμ, Ἀζαρίας.
\vs{9}Καὶ υἱοὶ Ἐσρὼμ οἳ ἐτέχθησαν αὐτῷ, ὁ Ἱεραμεὴλ, καὶ ὁ Ἀρὰμ, καὶ ὁ Χαλέβ.

\vs{10}Καὶ Ἀρὰμ ἐγέννησε τὸν Ἀμιναδὰβ, καὶ Ἀμιναδὰβ ἐγέννησε τὸν Ναασσὼν ἄρχοντα οἴκου Ἰούδα,
\vs{11}καὶ Ναασσὼν ἐγέννησε τὸν Σαλμὼν, καὶ Σαλμὼν ἐγέννησε τὸν Βοὸζ,
\vs{12}καὶ Βοὸζ ἐγέννησε τὸν Ὠβὴδ, καὶ Ὠβὴδ ἐγέννησε τὸν Ἰεσσαὶ,
\vs{13}καὶ Ἰεσσαὶ ἐγέννησε τὸν πρωτότοκον αὐτοῦ τὸν Ἐλιὰβ, Ἀμιναδὰβ ὁ δεύτερος, Σαμαὰ ὁ τρίτος,
\vs{14}Ναθαναὴλ ὁ τέταρτος, Ζαβδαῒ ὁ πέμπτος,
\vs{15}Ἀσὰμ ὁ ἕκτος, Δαυὶδ ὁ ἕβδομος.
\vs{16}Καὶ ἡ ἀδελφὴ αὐτῶν Σαρουία, καὶ Ἀβιγαία· καὶ υἱοὶ Σαρουία, Ἀβισὰ, καὶ Ἰωὰβ, καὶ Ἀσαὴλ, τρεῖς.
\vs{17}Καὶ Ἀβειγαία ἐγέννησε τὸν Ἀμεσσάβ· καὶ πατὴρ Ἀμεσσὰβ Ἰοθὸρ ὁ Ἰσμαηλίτης.

\vs{18}Καὶ Χαλὲβ υἱὸς Ἐσρὼμ ἔλαβε τὴν Γαζουβὰ γυναῖκα, καὶ τὴν Ἰεριώθ· καὶ οὗτοι υἱοὶ αὐτῆς, Ἰασὰρ, καὶ Σουβὰβ, καὶ Ἀρδών.
\vs{19}Καὶ ἀπέθανε Γαζουβὰ, καὶ ἔλαβεν ἑαυτῷ Χαλὲβ τὴν Ἐφρὰθ, καὶ ἔτεκεν αὐτῷ τὸν Ὥρ.
\vs{20}Καὶ Ὣρ ἐγέννησε τὸν Οὐρί· καὶ Οὐρὶ ἐγέννησε τὸν Βεσελεήλ.
\vs{21}Καὶ μετὰ ταῦτα εἰσῆλθεν Ἐσρὼν πρὸς τὴν θυγατέρα Μαχὶρ πατρὸς Γαλαὰδ, καὶ αὐτὸς ἔλαβεν αὐτὴν, καὶ αὐτὸς ἑξηκονταπέντε ἐτῶν ἦν· καὶ ἔτεκεν αὐτῷ τὸν Σερούχ.
\vs{22}Καὶ Σεροὺχ ἐγέννησε τὸν Ἰαΐρ. καὶ ἦσαν αὐτῷ εἴκοσι καὶ τρεῖς πόλεις ἐν τῇ Γαλαάδ.
\vs{23}Καὶ ἔλαβε Γεδσοὺρ καὶ Ἀρὰμ τὰς κώμας Ἰαῒρ ἐξ αὐτῶν, τὴν Κανὰθ καὶ τὰς κώμας αὐτῆς, ἑξήκοντα πόλεις· πᾶσαι αὗται υἱῶν Μαχὶρ πατρὸς Γαλαάδ.
\vs{24}Καὶ μετὰ τὸ ἀποθανεῖν Ἐσρὼν, ἦλθε Χαλὲβ εἰς Ἐφραθά· καὶ ἡ γυνὴ Ἐσρὼν Ἀβιά· καὶ ἔτεκεν αὐτῷ τὸν Ἀσχὼ πατέρα Θεκωέ.

\vs{25}καὶ ἦσαν οἱ υἱοὶ Ἱεραμεὴλ πρωτοτόκου Ἐσρὼν, ὁ πρωτότοκος Ῥὰμ, καὶ Βαναὰ, καὶ Ἀρὰμ, καὶ Ἀσὰν ἀδελφὸς αὐτοῦ.
\vs{26}Καὶ ἦν γυνὴ ἑτέρα τῷ Ἱεραμεὴλ, καὶ ὄνομα αὐτῇ Ἀτάρα· αὕτη ἐστὶ μήτηρ Ὀζόμ.
\vs{27}Καὶ ἦσαν υἱοὶ Ῥὰμ πρωτοτόκου Ἱεραμεὴλ, Μαὰς, καὶ Ἰαμὶν, καὶ Ἀκόρ.
\vs{28}Καὶ ἦσαν υἱοὶ Ὀζὸμ, Σαμαῒ, καὶ Ἰαδαέ· καὶ υἱοὶ Σαμαῒ, Ναδὰβ καὶ Ἀβισούρ.
\vs{29}Καὶ ὄνομα τῆς γυναικὸς Ἀβισοὺρ, Ἀβιχαία· καὶ ἔτεκεν αὐτῷ τὸν Ἀχαβὰρ, καὶ τὸν Μωήλ.
\vs{30}Καὶ υἱοὶ Ναδὰβ, Σαλὰδ, καὶ Ἀπφαίν· καὶ ἀπέθανε Σαλὰδ οὐκ ἔχων τέκνα.
\vs{31}Καὶ υἱοὶ Ἀπφαὶν, Ἰσεμιήλ· καὶ υἱοὶ Ἰσεμιὴλ, Σωσάν· καὶ υἱοὶ Σωσὰν, Δαδαί.
\vs{32}Καὶ υἱοὶ Δαδαὶ, Ἀχισαμὰς, Ἰεθὲρ, Ἰωνάθαν· καὶ ἀπέθανεν Ἰεθὲρ οὐκ ἔχων τέκνα.
\vs{33}Καὶ υἱοὶ Ἰωνάθαν, Φαλὲθ, καὶ Ὁζάμ· οὗτοι ἦσαν υἱοὶ Ἱεραμεήλ.

\vs{34}Καὶ οὐκ ἦσαν τῷ Σωσὰν υἱοὶ, ἀλλʼ ἢ θυγατέρες· καὶ τῷ Σωσὰν παῖς Αἰγύπτιος, καὶ ὄνομα αὐτῷ Ἰωχήλ.
\vs{35}Καὶ ἔδωκε Σωσὰν τὴν θυγατέρα αὐτοῦ τῷ Ἰωχὴλ παιδὶ αὐτοῦ εἰς γυναῖκα, καὶ ἔτεκεν αὐτῷ τὸν Ἐθὶ,
\vs{36}καὶ Ἐθὶ ἐγέννησε τὸν Ναθὰν, καὶ Ναθὰν ἐγέννησε τὸν Ζαβὲδ,
\vs{37}καὶ Ζαβὲδ ἐγέννησε τὸν Ἀφαμὴλ, καὶ Ἀφαμὴλ ἐγέννησε τὸν Ὠβὴδ,
\vs{38}καὶ Ὠβὴδ ἐγέννησε τὸν Ἰηοὺ, καὶ Ἰηοὺ ἐγέννησε τὸν Ἀζαρίαν,
\vs{39}καὶ Ἀζαρίας ἐγέννησε τὸν Χελλὴς, καὶ Χελλὴς ἐγέννησε τὸν Ἐλεασὰ,
\vs{40}καὶ Ἐλεασὰ ἐγέννησε τὸν Σοσομαῒ, καὶ Σοσομαῒ ἐγέννησε τὸν Σαλοὺμ,
\vs{41}καὶ Σαλοὺμ ἐγέννησε τὸν Ἰεχεμίαν, καὶ Ἰεχεμίας ἐγέννησε τὸν Ἐλισαμὰ, καὶ Ἐλισαμὰ ἐγέννησε τὸν Ἰσμαήλ.

\vs{42}Καὶ υἱοὶ Χαλὲβ ἀδελφοῦ Ἱεραμεὴλ, Μαρισὰ ὁ πρωτότοκος αὐτοῦ· οὗτος πατὴρ Ζίφ· καὶ υἱοὶ Μαρισὰ πατρὸς Χεβρών.
\vs{43}Καὶ υἱοὶ Χεβρὼν, Κορὲ, καὶ Θαπφοὺς, καὶ Ῥεκὸμ, καὶ Σαμαά.
\vs{44}Καὶ Σαμαὰ ἐγέννησε τὸν Ῥαὲμ πατέρα Ἰεκλὰν, καὶ Ἰεκλὰν ἐγέννησε τὸν Σαμαΐ.
\vs{45}Καὶ υἱὸς αὐτοῦ Μαών· καὶ Μαὼν πατὴρ Βαιθσούρ.
\vs{46}Καὶ Γαιφὰ ἡ παλλακὴ Χαλὲβ ἐγέννησε τὸν Ἀρὰμ, καὶ τὸν Μοσὰ, καὶ τὸν Γεζουέ.
\vs{47}Καὶ υἱοὶ Ἀδδαῒ, Ῥαγὲμ, καὶ Ἰωάθαμ, καὶ Σωγὰρ, καὶ Φαλὲκ, καὶ Γαιφὰ, καὶ Σαγαέ.
\vs{48}Καὶ ἡ παλλακὴ Χαλὲβ Μωχὰ ἐγέννησε τὸν Σαβὲρ, καὶ τὸν Θαράμ.
\vs{49}Καὶ ἐγέννησε Σαγαὲ πατέρα Μαδμηνὰ, καὶ τὸν Σαοὺ πατέρα Μαχαβηνὰ, καὶ πατέρα Γαιβάλ· καὶ θυγάτηρ Χαλὲβ, Ἀσχά.

\vs{50}Οὗτοι ἦσαν υἱοὶ Χαλέβ· υἱοὶ Ὢρ πρωτοτόκου Ἐφραθά· Σωβὰλ πατὴρ Καριαθιαρὶμ,
\vs{51}Σαλωμὼν πατὴρ Βαιθὰ, Λαμμὼν πατὴρ Βαιθαλαὲμ, καὶ Ἀρὶμ πατὴρ Βεθγεδώρ.
\vs{52}Καὶ ἦσαν υἱοὶ τῷ Σωβὰλ πατρὶ Καριαθιαρὶμ Ἀραὰ, καὶ Αἰσὶ, καὶ Ἀμμανὶθ,
\vs{53}καὶ Οὐμασφαὲ, πόλεις Ἰαῒρ, Αἰθαλὶμ, καὶ Μιφιθὶμ, καὶ Ἡσαμαθὶμ, καὶ Ἡμασαραΐμ· ἐκ τούτων ἐξήλθοσαν οἱ Σαραθαῖοι, καὶ υἱοὶ Ἐσθαάμ.
\vs{54}Υἱοὶ Σαλωμὼν Βαιθαλαὲμ, ὁ Νετωφατὶ, Ἀταρὼθ οἴκου Ἰωὰβ, καὶ ἥμισυ τῆς Μαλαθὶ, Ἠσαρὶ·
\vs{55}Πατριαὶ γραμματέων κατοικοῦντες ἐν Ἰάβις Θαργαθιῒμ, καὶ Σαμαθιῒμ, καὶ Σωχαθίμ· οὗτοι οἱ Κιναῖοι οἱ ἐλθόντες ἐξ Αἱμὰθ πατρὸς οἴκου Ῥηχάβ.

\ch{3}
Καὶ οὗτοι ἦσαν υἱοὶ Δαυὶδ οἱ τεχθέντες αὐτῷ ἐν Χεβρών· ὁ πρωτότοκος Ἀμνὼν τῇ Ἀχιναὰμ τῇ Ἰεζραηλίτιδι· ὁ δεύτερος Δαμνιὴλ τῇ Ἀβιγαίᾳ τῇ Καρμηλίᾳ·
\vs{2}Ὁ τρίτος Ἀβεσσαλὼμ, υἱὸς Μωχὰ θυγατρὸς Θολμαῒ βασιλέως Γεδσούρ· ὁ τέταρτος Ἀδωνία υἱὸς Ἀγγίθ·
\vs{3}Ὁ πέμπτος Σαφατία τῆς Ἀβιτὰλ· ὁ ἕκτος Ἰεθραὰμ τῇ Ἀγλᾷ γυναικὶ αὐτοῦ.
\vs{4}Ἓξ ἐγεννήθησαν αὐτῷ ἐν Χεβρών· καὶ ἐβασίλευσεν ἐκεῖ ἑπτὰ ἔτη, καὶ ἑξάμηνον· καὶ τριάκοντα καὶ τρία ἔτη ἐβασίλευσεν ἐν Ἱερουσαλήμ.
\vs{5}Καὶ οὗτοι ἐτέχθησαν αὐτῷ ἐν Ἱερουσαλήμ· Σαμαὰ, Σωβὰβ, Νάθαν, καὶ Σαλωμών· τέσσαρες τῇ Βηρσαβεὲ θυγατρὶ Ἀμιήλ·
\vs{6}Καὶ Ἐβαὰρ, καὶ Ἐλισὰ, καὶ Ἐλιφαλὴθ,
\vs{7}καὶ Ναγαὶ, καὶ Ναφὲκ, καὶ Ἰαφιὲ,
\vs{8}καὶ Ἑλισαμὰ, καὶ Ἐλιαδὰ, καὶ Ἐλιφαλὰ, ἐννέα.
\vs{9}Πάντες υἱοὶ Δαυὶδ, πλὴν τῶν υἱῶν τῶν παλλακῶν, καὶ Θήμαρ ἀδελφὴ αὐτῶν.

\vs{10}Υἱοὶ Σαλωμὼν, Ῥοβοάμ, Ἀβιὰ υἱὸς αὐτοῦ, Ἀσὰ υἱὸς αὐτοῦ, Ἰωσαφὰτ υἱὸς αὐτοῦ,
\vs{11}Ἰωρὰμ υἱὸς αὐτοῦ, Ὀχοζίας υἱὸς αὐτοῦ, Ἰωὰς υἱὸς αὐτοῦ,
\vs{12}Ἀμασίας υἱὸς αὐτοῦ, Ἀζαρίας υἱὸς αὐτοῦ, Ἰωάθαν υἱὸς αὐτοῦ,
\vs{13}Ἄχαζ υἱὸς αὐτοῦ, Ἐζεκίας υἱὸς αὐτοῦ, Μανασσῆς υἱὸς αὐτοῦ,
\vs{14}Ἀμὼν υἱὸς αὐτοῦ, Ἰωσία υἱὸς αὐτοῦ.
\vs{15}Καὶ υἱοὶ Ἰωσία, πρωτότοκος Ἰωανὰν, ὁ δεύτερος Ἰωακεὶμ, ὁ τρίτος Σεδεκίας, ὁ τέταρτος Σαλούμ.
\vs{16}Καὶ υἱοὶ Ἰωακείμ, Ἰεχονίας υἱὸς αὐτοῦ, Σεδεκίας υἱὸς αὐτοῦ.
\vs{17}Καὶ υἱοὶ Ἰεχονία, Ἀσὶρ, Σαλαθιὴλ υἱὸς αὐτοῦ,
\vs{18}Μελχειρὰμ, καὶ Φαδαΐας, καὶ Σανεσὰρ, καὶ Ἰεκιμία, καὶ Ὡσαμὰθ, καὶ Ναβαδίας.

\vs{19}Καὶ υἱοὶ Φαδαΐας, Ζοροβάβελ, καὶ Σεμεΐ· καὶ υἱοὶ Ζοροβάβελ, Μοσολλὰμ, καὶ Ἀνανία, καὶ Σαλωμεθὶ ἀδελφὴ αὐτῶν,
\vs{20}καὶ Ἀσουβὲ, καὶ Ὀὸλ, καὶ Βαραχὶα, καὶ Ἀσαδία, καὶ Ἀσοβὲδ, πέντε.

\vs{21}Καὶ υἱοὶ Ἀνανία, Φαλεττία, καὶ Ἰεσίας υἱὸς αὐτοῦ, Ῥαφὰλ υἱὸς αὐτοῦ, Ὀρνὰ υἱὸς αὐτοῦ, Ἀβδία υἱὸς αὐτοῦ, Σεχενίας υἱὸς αὐτοῦ.
\vs{22}Καὶ υἱὸς Σεχενία, Σαμαΐα· καὶ υἱοὶ Σαμαΐα, Χαττοὺς, καὶ Ἰωὴλ, καὶ Βεῤῥὶ, καὶ Νωαδία, καὶ Σαφὰθ, ἕξ.

\vs{23}Καὶ υἱοὶ Νωαδία, Ἐλιθενὰν, καὶ Ἐζεκία, καὶ Ἐζρικὰμ, τρεῖς.

\vs{24}Καὶ υἱοὶ Ἐλιθενὰν, Ὀδολία, καὶ Ἑλιασεβὼν, καὶ Φαδαΐα, καὶ Ἀκοὺβ, καὶ Ἰωανὰν, καὶ Δαλααΐα, καὶ Ἀνὰν, ἑπτά.

\ch{4}
Καὶ υἱοὶ Ἰούδα, Φαρὲς, Ἐσρὼμ, καὶ Χαρμὶ, καὶ Ὢρ, Σουβὰλ,
\vs{2}καὶ Ῥάδα υἱὸς αὐτοῦ· καὶ Σουβὰλ ἐγέννησε τὸν Ἰέθ· καὶ Ἰὲθ ἐγέννησε τὸν Ἀχιμαῒ, καὶ τὸν Λαάδ· αὗται αἱ γενέσεις τοῦ Ἀραθί.
\vs{3}Καὶ οὗτοι υἱοὶ Αἰτὰμ, Ἰεζραὴλ, καὶ Ἰεσμὰν, καὶ Ἰεβδάς· καὶ ὄνομα ἀδελφῆς αὐτῶν Ἐσηλεββών.
\vs{4}Καὶ Φανουὴλ πατὴρ Γεδὼρ, καὶ Ἰαζὴρ πατὴρ Ὠσάν· οὗτοι υἱοὶ Ὢρ τοῦ πρωτοτόκου Ἐφραθὰ πατρὸς Βαιθαλαέν.

\vs{5}Καὶ τῷ Ἀσοὺρ πατρὶ Θεκωὲ ἦσαν δύο γυναῖκες, Ἀωδὰ, καὶ Θοαδά.
\vs{6}Καὶ ἔτεκεν αὐτῷ Ἀωδὰ τὸν Ὠχαία, καὶ τὸν Ἠφὰλ, καὶ τὸν Θαιμὰν, καὶ τὸν Ἀασθήρ· πάντες οὗτοι υἱοὶ Ἀωδᾶς.
\vs{7}Καὶ υἱοὶ Θοαδᾶς, Σερὲθ, καὶ Σαὰρ, καὶ Ἐσθανάμ.
\vs{8}Καὶ Κωὲ ἐγέννησε τὸν Ἐνὼβ, καὶ τὸν Σαβαθά· καὶ γεννήσεις ἀδελφοῦ Ῥηχὰβ, υἱοῦ Ἰαρίν.
\vs{9}Καὶ ἦν Ἰγαβὴς ἔνδοξος ὑπὲρ τοὺς ἀδελφοὺς αὐτοῦ· καὶ ἡ μήτηρ ἐκάλεσε τὸ ὄνομα αὐτοῦ Ἰγαβὴς, λέγουσα, ἔτεκον ὡς γαβής.
\vs{10}Καὶ ἐπεκαλέσατο Ἰγαβὴς τὸν Θεὸν Ἰσραὴλ, λέγων, ἐὰν εὐλογῶν εὐλογήσῃς με, καὶ πληθύνῃς τὰ ὅριά μου, καὶ ᾖ ἡ χείρ σου μετʼ ἐμοῦ, καὶ ποιήσῃς γνῶσιν τοῦ μὴ ταπεινῶσαί με· καὶ ἐπήγαγεν ὁ Θεὸς πάντα ὅσα ᾐτήσατο.

\vs{11}Καὶ Χαλὲβ πατὴρ Ἀσχὰ ἐγέννησε τὸν Μαχίρ· οὗτος πατὴρ Ἀσσαθών.
\vs{12}Ἐγέννησε τὸν Βαθραίαν, καὶ τὸν Βεσσηὲ, καὶ τὸν Θαιμὰν πατέρα πόλεως Ναᾶς ἀδελφοῦ Ἐσελὼμ τοῦ Κενεζί· οὗτοι ἄνδρες Ῥηχάβ.
\vs{13}Καὶ υἱοὶ Κενὲζ, Γοθονιὴλ, καὶ Σαραΐα· καὶ υἱοὶ Γοθονιὴλ, Ἀθάθ.
\vs{14}Καὶ Μαναθὶ ἐγέννησε τὸν Γοφερά. καὶ Σαραΐα ἐγέννησε τὸν Ἰωβὰβ, πατέρα Ἀγεαδδαῒρ, ὅτι τέκτονες ἦσαν.
\vs{15}Καὶ υἱοὶ Χαλὲβ υἱοῦ Ἰεφοννὴ, Ἢρ, Ἀδὰ, καὶ Νοόμ· καὶ υἱοὶ Ἀδὰ, Κενέζ.
\vs{16}Καὶ υἱοὶ Ἀλεὴλ, Ζὶβ, καὶ Ζεφὰ, καὶ Θιριὰ, καὶ Ἐσερήλ.
\vs{17}Καὶ υἱοὶ Ἐσρὶ, Ἰεθὲρ, Μωρὰδ, καὶ Ἄφερ, καὶ Ἰαμών· καὶ ἐγέννησεν Ἰεθὲρ τὸν Μαρὼν, καὶ τὸν Σεμεῒ, καὶ τὸν Ἰεσβὰ πατέρα Ἐσθαίμων·
\vs{18}Καὶ ἡ γυνὴ αὐτοῦ αὕτη Ἀδία, ἔτεκε τὸν Ἰάρεδ πατέρα Γεδὼρ, καὶ τὸν Ἀβὲρ πατέρα Σωχὼν, καὶ τὸν Χετιὴλ πατέρα Ζαμών· καὶ οὗτοι υἱοὶ Βετθία θυγατρὸς Φαραὼ, ἣν ἔλαβε Μωρήδ.
\vs{19}Καὶ υἱοὶ γυναικὸς τῆς Ἰδουίας ἀδελφῆς Ναχαῒμ πατρὸς Κεϊλὰ, Γαρμὶ, καὶ Ἐσθαιμὼν Νωχαθί.
\vs{20}Καὶ υἱοὶ Σεμὼν, Αμνὼν, καὶ Ἀνὰ υἱὸς Φανὰ, καὶ Ἰνών· καὶ υἱοὶ Σεῒ, Ζωὰν, καὶ υἱοὶ Ζωάβ.

\vs{21}Υἱοὶ Σηλὼμ υἱοῦ Ἰούδα, Ἢρ πατὴρ Ληχὰβ, καὶ Λααδὰ πατὴρ Μαρισά· καὶ γενέσεις οἰκείων Ἐφραθαβὰκ τῷ οἴκῳ Ἐσοβὰ,
\vs{22}καὶ Ἰωακὶμ, καὶ ἄνδρες Χωζηβὰ, καὶ Ἰωὰς, καὶ Σαρὰφ, οἳ κατῴκησαν ἐν Μωάβ· καὶ ἀπέστρεψεν αὐτος ἀβεδηρὶν, ἀθουκιΐμ·
\vs{23}Οὗτοι κεραμεῖς οἱ κατοικοῦντες ἐν Ἀταῒμ καὶ Γαδιρὰ μετὰ τοῦ βασιλέως, ἐν τῇ βασιλείᾳ αὐτοῦ ἐνίσχυσαν, καὶ κατῴκησαν ἐκεῖ.

\vs{24}Υἱοὶ Σεμεὼν, Ναμουὴλ, καὶ Ἰαμὶν, Ἰαρὶβ, Ζαρὲς, Σαοὺλ,
\vs{25}Σαλὲμ υἱὸς αὐτοῦ, Μαβασὰμ υἱὸς αὐτοῦ, Μασμὰ υἱὸς αὐτοῦ,
\vs{26}Σεμεῒ υἱὸς αὐτοῦ·
\vs{27}Τῷ Σεμεῒ υἱοὶ ἑκκαίδεκα, καὶ θυγατέρες ἕξ· καὶ τοῖς ἀδελφοῖς αὐτῶν οὐκ ἦσαν υἱοὶ πολλοί· καὶ πᾶσαι αἱ πατριαὶ αὐτῶν οὐκ ἐπλεόνασαν ὡς υἱοὶ Ἰούδα.
\vs{28}Καὶ κατῴκησαν ἐν Βηρσαβεὲ, καὶ Μωλαδὰ, καὶ ἐν Ἑσερσουὰλ,
\vs{29}καὶ ἐν Βαλαὰ, καὶ ἐν Αἰσὲμ, καὶ ἐν Θωλὰδ,
\vs{30}καὶ ἐν Ἑρμὰ, καὶ ἐν Σικελὰγ,
\vs{31}καὶ ἐν Βαιθμαριμὼθ, καὶ Ἡμισουσεωσὶν, καὶ οἴκου Βαρουσεωρίμ· αὗται αἱ πόλεις αὐτῶν ἕως βασιλέως Δαυίδ.
\vs{32}Καὶ ἐπαύλεις αὐτῶν Αἰτὰν, καὶ Ἢν, Ῥεμνὼν, καὶ Θοκκὰ, καὶ Αἰσὰρ, πόλεις πέντε.
\vs{33}Καὶ πᾶσαι ἐπαύλεις αὐτῶν κύκλῳ τῶν πόλεων τούτων ἕως Βάαλ· αὕτη κατάσχεσις αὐτῶν, καὶ ὁ καταλοχισμὸς αὐτῶν.
\vs{34}Καὶ Μοσωβὰβ, καὶ Ἰεμολὸχ, καὶ Ἰωσία υἱὸς Ἀμασία,
\vs{35}καὶ Ἰωὴλ, καὶ Ἰηοὺ υἱὸς Ἀσαβία, υἱὸς Σαραῦ, υἱὸς Ἀσιὴλ,
\vs{36}καὶ Ἐλιωναῒ, καὶ Ἰωκαβὰ, καὶ Ἰασουία, καὶ Ἀσαΐα, καὶ Ἰεδιὴλ, καὶ Ἰσμαὴλ, καὶ Βαναίας, καὶ Ζουζὰ υἱὸς Σαφαῒ,
\vs{37}υἱοῦ Ἀλὼν, υἱοῦ Ἰεδιὰ, υἱοῦ Σεμρὶ, υἱοῦ Σαμαίου.
\vs{38}Οὗτοι οἱ διελθόντες ἐν ὀνόμασιν ἀρχόντων ἐν ταῖς γενέσεσιν αὐτῶν, καὶ ἐν οἴκοις πατριῶν αὐτῶν ἐπληθύνθησαν εἰς πλῆθος.

\vs{39}Καὶ ἐπορεύθησαν ἕως τοῦ ἐλθεῖν Γέραρα ἕως τῶν ἀνατολῶν τῆς Γαὶ, τοῦ ζητῆσαι νομὰς τοῖς κτήνεσιν αὐτῶν.
\vs{40}Καὶ εὗρον νομὰς πλεὶονας καὶ ἀγαθάς· καὶ ἡ γῆ πλατεῖα ἐναντίον αὐτῶν, καὶ εἰρήνη καὶ ἡσυχία, ὅτι ἐκ τῶν υἱῶν Χὰμ τὼν κατοικούντων ἐκεῖ ἔμπροσθεν.
\vs{41}Καὶ ἤλθοσαν οὗτοι οἱ γεγραμμένοι ἐπʼ ὀνόματος ἐν ἡμέραις Ἐζεκίου βασιλέως Ἰούδα, καὶ ἐπάταξαν τοὺς οἴκους αὐτῶν καὶ τοὺς Μιναίους οὕς εὕροσαν ἐκεῖ, καὶ ἀνεθεμάτισαν αὐτοὺς ἕως τῆς ἡμέρας ταύτης· καὶ ᾤκησαν ἀντʼ αὐτῶν, ὅτι νομαὶ τοῖς κτήνεσιν αὐτῶν ἐκεῖ.
\vs{42}Καὶ ἐξ αὐτῶν ἀπὸ τῶν υἱῶν Συμεὼν ἐπορεύθησαν εἰς ὄρος Σηὶρ ἄνδρες πεντακόσιοι, καὶ Φαλαεττία, καὶ Νωαδία, καὶ Ῥαφαΐα, καὶ Ὀζιὴλ υἱοὶ Ἰεσὶ ἄρχοντες αὐτῶν.
\vs{43}Καὶ ἐπάταξαν τοὺς καταλοίπους τοὺς καταλειφθέντας τοῦ Ἀμαλὴκ ἕως ἡμέρας ταύτης.

\ch{5}
Καὶ υἱοὶ Ῥουβὴν πρωτοτόκου Ἰσραήλ· ὅτι οὗτος ὁ πρωτότοκος, καὶ ἐν τῷ ἀναβῆναι ἐπὶ τὴν κοίτην τοῦ πατρὸς αὐτοῦ ἔδωκεν εὐλογίαν αὐτοῦ τῷ υἱῷ αὐτοῦ Ἰωσὴφ υἱῷ Ἰσραήλ, καὶ οὐκ ἐγενεαλογήθη εἰς πρωτοτοκεῖα,
\vs{2}ὅτι Ἰούδας δυνατὸς ἰσχύι καὶ ἐν τοῖς ἀδελφοῖς αὐτοῦ, καὶ εἰς ἡγούμενον ἐξ αὐτοῦ, καὶ ἡ εὐλογία τοῦ Ἰωσήφ·
\vs{3}Υἱοὶ Ῥουβὴν πρωτοτόκου Ἰσραὴλ, Ἐνὼχ, καὶ Φαλλοὺς, Ἀσρὼμ, καὶ Χαρμί.
\vs{4}Υἱοὶ Ἰωὴλ, Σεμεῒ, καὶ Βαναία υἱὸς αὐτοῦ· καὶ υἱοὶ Γοὺγ υἱοῦ Σεμεῒ,
\vs{5}υἱὸς αὐτοῦ Μιχὰ, υἱὸς αὐτοῦ Ῥηχά, υἱὸς αὐτοῦ Ἰωὴλ,
\vs{6}υἱὸς αὐτοῦ Βεὴλ, ὃν μετῴκισε Θαγλαφαλλασὰρ βασιλεὺς Ἀσσούρ· οὗτος ἄρχων τῶν Ῥουβήν.

\vs{7}Καὶ ἀδελφοὶ αὐτοῦ τῇ πατρίδι αὐτοῦ ἐν τοῖς καταλοχισμοῖς αὐτῶν κατὰ γενέσεις αὐτῶν, ὁ ἄρχων Ἰωὴλ, καὶ Ζαχαρία,
\vs{8}καὶ Βαλὲκ υἱὸς Ἀζοὺζ, υἱὸς Σαμὰ, υἱὸς Ἰωήλ· οὗτος κατῴκησεν ἐν Ἀροὴρ, καὶ ἐπὶ Ναβαῦ, καὶ Βεελμασσών.
\vs{9}Καὶ πρὸς ἀνατολὰς κατῴκησεν ἕως ἐρχομένων τῆς ἐρήμου, ἀπὸ τοῦ ποταμοῦ Εὐφράτου, ὅτι κτήνη αὐτῶν πολλὰ ἐν γῇ Γαλαάδ.
\vs{10}Καὶ ἐν ἡμέραις Σαοὺλ ἐποίησαν πόλεμον πρὸς τοὺς παροίκους, καὶ ἔπεσον ἐν χερσὶν αὐτῶν κατοικοῦντες ἐν σκηναῖς αὐτῶν πάντες κατʼ ἀνατολὰς τῆς Γαλαάδ.

\vs{11}Υἱοὶ Γὰδ κατέναντι αὐτῶν κατῴκησαν ἐν γῇ Βασὰν ἕως Σελά·
\vs{12}Ἰωὴλ πρωτότοκος, καὶ Σαφὰμ ὁ δεύτερος, καὶ Ἰανὶν ὁ γραμματεὺς ἐν Βασάν.
\vs{13}Καὶ οἱ ἀδελφοὶ αὐτῶν κατʼ οἴκους πατριῶν αὐτῶν, Μιχαὴλ, Μοσολλὰμ, καὶ Σεβεὲ, καὶ Ἰωρεὲ, καὶ Ἰωαχὰν, καὶ Ζουὲ, καὶ Ὠβὴδ, ἑπτά.
\vs{14}Οὗτοι υἱοὶ Ἀβιχαία υἱοῦ Οὐρὶ, υἱοῦ Ἰδαῒ, υἱοῦ Γαλαὰδ, υἱοῦ Μιχαὴλ, υἱοῦ Ἰεσαῒ, υἱοῦ Ἰεδδαῒ, υἱοῦ Βοὺζ ἀδελφοῦ
\vs{15}υἱοῦ Ἀβδιὴλ, υἱοῦ Γουνὶ, ἄρχων οἴκου πατριῶν.
\vs{16}Κατῴκουν ἐν Γαλαὰδ, ἐν Βασὰν, καὶ ἐν ταῖς κώμαις αὐτῶν, καὶ πάντα τὰ περίχωρα Σαρὼν ἕως ἐξόδου.
\vs{17}Πάντων ὁ καταλοχισμὸς ἐν ἡμέραις Ἰωάθαμ βασιλέως Ἰούδα, καὶ ἐν ἡμέραις Ἱεροβοὰμ βασιλέως Ἰσραήλ.

\vs{18}Υἱοὶ Ῥουβὴν καὶ Γὰδ καὶ ἥμισυ φυλῆς Μανασσῆ ἐξ υἱῶν δυνάμεως, ἄνδρες αἴροντες ἀσπίδας καὶ μάχαιραν, καὶ τείνοντες τόξον, καὶ δεδιδαγμένοι πόλεμον, τεσσαράκοντα καὶ τέσσαρες χιλιάδες καὶ ἑπτακόσιοι καὶ ἑξήκοντα ἐκπορευόμενοι εἰς παράταξιν.
\vs{19}Καὶ ἐποίουν πόλεμον μετὰ τῶν Ἀγαρηνῶν, καὶ Ἰτουραίων, καὶ Ναφισαίων, καὶ Ναδαβαίων,
\vs{20}καὶ κατίσχυσαν ἐπʼ αὐτῶν· καὶ ἐδόθησαν εἰς χεῖρας αὐτῶν Ἀγαραῖοι, καὶ πάντα τὰ σκηνώματα αὐτῶν, ὅτι πρὸς τὸν Θεὸν ἐβόησαν ἐν τῷ πολέμῳ, καὶ ἐπήκουσεν αὐτοῖς, ὅτι ἤλπισαν ἐπʼ αὐτόν.
\vs{21}Καὶ ᾐχμαλώτευσαν τὴν ἀποσκευὴν αὐτῶν, καμήλους πεντακισχιλίας, καὶ προβάτων διακοσίας πεντήκοντα χιλιάδας, ὄνους δισχιλίους, καὶ ψυχὰς ἀνδρῶν ἑκατὸν χιλιάδας.
\vs{22}Ὅτι τραυματίαι πολλοὶ ἔπεσον, ὅτι παρὰ τοῦ Θεοῦ ὁ πόλεμος· καὶ κατῴκησαν ἀντʼ αὐτῶν ἕως μετοικεσίας.

\vs{23}Καὶ οἱ ἡμίσεις φυλῆς Μανασσῆ κατῴκησαν ἀπὸ Βασὰν ἕως Βαὰλ, Ἐρμὼν, καὶ Σανὶρ, καὶ ὄρος Ἀερμών· καὶ ἐν τῷ Λιβάνῳ αὐτοὶ ἐπλεονάσθησαν.
\vs{24}Καὶ οὗτοι ἀρχηγοὶ οἴκου πατριῶν αὐτῶν· Ὀφὲρ, καὶ Σεῒ, καὶ Ἐλιὴλ, καὶ Ἱερεμία, καὶ Ὠδουΐα, καὶ Ἰεδιήλ· ἄνδρες ἰσχυροὶ δυνάμει, ἄνδρες ὀνομαστοὶ, ἄρχοντες τῶν οἴκων πατριῶν αὐτῶν.

\vs{25}Καὶ ἠθέτησαν ἐν Θεῷ πατέρων αὐτῶν, καὶ ἐπόρνευσαν ὀπίσω θεῶν τῶν λαῶν τῆς γῆς, οὓς ἐξῇρεν ὁ Θεὸς ἀπὸ προσώπου αὐτῶν.
\vs{26}Καὶ ἐπήγειρεν ὁ Θεὸς Ἰσραὴλ τὸ πνεῦμα Φαλὼχ βασιλέως Ἀσσοὺρ, καὶ τὸ πνεῦμα Θαγλαφαλλασὰρ βασιλέως Ἀσσοὺρ, καὶ μετῴκισε τὸν Ῥουβὴν, καὶ τὸν Γαδδὶ, καὶ τὸ ἥμισυ φυλῆς Μανασσῆ, καὶ ἤγαγεν αὐτοὺς εἰς Χαὰχ, καὶ Χαβὼρ, καὶ ἐπὶ ποταμὸν Γωζὰν ἕως τῆς ἡμέρας ταύτης.

\vs{27}Υἱοὶ Λευὶ, Γεδσὼν, Καὰθ, καὶ Μεραρί.
\vs{28}Καὶ υἱοὶ Καὰθ, Ἄμβραμ, καὶ Ἰσσαὰρ, Χεβρὼν, καὶ Ὀζιήλ.
\vs{29}Καὶ υἱοὶ Ἄμβραμ, Ἀαρὼν, καὶ Μωυσῆς, καὶ Μαριάμ· καὶ υἱοὶ Ἀαρών, Ναδὰβ, καὶ Ἀβιοὺδ, Ἐλεάζαρ, καὶ Ἰθάμαρ.
\vs{30}Ἐλεάζαρ ἐγέννησε τὸν Φινεὲς, Φινεὲς ἐγέννησε τὸν Αβισοὺ,
\vs{31}Αβισοὺ ἐγέννησεν τὸν Βοκκὶ, καὶ Βοκκὶ ἐγέννησε τὸν Ὀζὶ,
\vs{32}Ὀζὶ ἐγέννησε τὸν Ζαραία, Ζαραία ἐγέννησε τὸν Μαριὴλ,
\vs{33}καὶ Μαριὴλ ἐγέννησε τὸν Ἀμαρία, καὶ Ἀμαρία ἐγέννησε τὸν Ἀχιτὼβ,
\vs{34}καὶ Ἀχιτὼβ ἐγέννησε τὸν Σαδὼκ, καὶ Σαδὼκ ἐγέννησε τὸν Ἀχιμάας,
\vs{35}καὶ Ἀχιμάας ἐγέννησε τὸν Ἀζαρίαν, καὶ Ἀζαρίας ἐγέννησε τὸν Ἰωανὰν,
\vs{36}καὶ Ἰωανὰν ἐγέννησε τὸν Ἀζαρίαν, οὗτος ἱεράτευσεν ἐν τῷ οἴκῳ ᾧ ᾠκοδόμησε Σαλωμὼν ἐν Ἱερουσαλήμ.
\vs{37}Καὶ ἐγέννησεν Ἀζαρίας τὸν Ἀμαρία, καὶ Ἀμαρία ἐγέννησε τὸν Ἀχιτὼβ,
\vs{38}καὶ Ἀχιτὼβ ἐγέννησε τὸν Σαδὼκ, καὶ Σαδὼκ ἐγέννησε τὸν Σαλὼμ,
\vs{39}καὶ Σαλὼμ ἐγέννησε τὸν Χελκίαν, καὶ Χελκίας ἐγέννησε τὸν Ἀζαρίαν,
\vs{40}καὶ Ἀζαρίας ἐγέννησε τὸν Σαραία, καὶ Σαραίας ἐγέννησε τὸν Ἰωσαδάκ.
\vs{41}Καὶ Ἰωσαδὰκ ἐπορεύθη ἐν τῇ μετοικίᾳ μετὰ Ἰούδα καὶ Ἱερουσαλὴμ ἐν χειρὶ Ναβουχοδονόσορ.

\ch{6}
Υἱοὶ Λευὶ, Γεδσὼν, Καὰθ, καὶ Μεραρί.
\vs{2}Καὶ ταῦτα τὰ ὀνόματα τῶν υἱῶν Γεδσὼν, Λοβενὶ, καὶ Σεμεΐ.
\vs{3}Υἱοὶ Καὰθ, Ἄμβραμ, καὶ Ἰσσαὰρ, Χεβρὼν, καὶ Ὀζιήλ.
\vs{4}Υἱοὶ Μεραρὶ, Μοολὶ, καὶ ὁ Μουσί· καὶ αὗται αἱ πατριαὶ τοῦ Λευὶ κατὰ πατριὰς αὐτῶν.
\vs{5}Τῷ Γεδσὼν, τῷ Λοβενὶ υἱῷ αὐτοῦ, Ἰὲθ υἱὸς αὐτοῦ, Ζαμμὰθ υἱὸς αὐτοῦ,
\vs{6}Ἰωὰβ υἱὸς αὐτοῦ, Ἀδδὶ υἱὸς αὐτοῦ, Ζαρὰ υἱὸς αὐτοῦ, Ἰεθρὶ υἱὸς αὐτοῦ.
\vs{7}Υἱοὶ Καὰθ, Ἀμιναδὰβ υἱὸς αὐτοῦ, Κορὲ υἱὸς αὐτοῦ, Ἀσὴρ υἱὸς αὐτοῦ,
\vs{8}Ἑλκανὰ υἱὸς αὐτοῦ, Ἀβισὰφ υἱὸς αὐτοῦ, Ἀσὴρ υἱὸς αὐτοῦ,
\vs{9}Θαὰθ υἱὸς αὐτοῦ, Οὐριὴλ υἱὸς αὐτοῦ, Ὀζία υἱὸς αὐτοῦ, Σαοὺλ υἱὸς αὐτοῦ.
\vs{10}Καὶ υἱοὶ Ἑλκανὰ, Ἀμεσσὶ, καὶ Ἀχιμὼθ,
\vs{11}Ἑλκανὰ υἱὸς αὐτοῦ, Σουφὶ υἱὸς αὐτοῦ, Καιναὰθ υἱὸς αὐτοῦ,
\vs{12}Ἐλιὰβ υἱὸς αὐτοῦ, Ιεροβοὰμ υἱὸς αὐτοῦ, Ἑλκανὰ υἱὸς αὐτοῦ.
\vs{13}Υἱοὶ Σαμουὴλ, ὁ πρωτότοκος Σανὶ, καὶ Ἀβιά.
\vs{14}Υἱοὶ Μεραρὶ, Μοολὶ, Λοβενὶ υἱὸς αὐτοῦ, Σεμεῒ υἱὸς αὐτοῦ, Ὀζὰ υἱὸς αὐτοῦ,
\vs{15}Σαμαὰ υἱὸς αὐτοῦ, Ἀγγία υἱὸς αὐτοῦ, Ἀσαῒας υἱὸς αὐτοῦ.

\vs{16}Καὶ οὗτοι οὓς κατέστησε Δαυὶδ ἐπὶ χεῖρας ᾀδόντων ἐν οἴκῳ Κυρίου ἐν τῇ καταπαύσει τῆς κιβωτοῦ.
\vs{17}Καὶ ἦσαν λειτουργοῦντες ἐναντίον τῆς σκηνῆς τοῦ μαρτυρίου ἐν ὀργάνοις, ἕως οὗ ᾠκοδόμησε Σαλωμὼν τὸν οἶκον Κυρίου ἐν Ἱερουσαλήμ· καὶ ἔστησαν κατὰ τὴν κρίσιν αὐτῶν ἐπὶ τὰς λειτουργίας αὐτῶν.

\vs{18}Καὶ οὗτοι οἱ ἑστηκότες, καὶ υἱοὶ αὐτῶν ἐκ τῶν υἱῶν τοῦ Καὰθ, Αἰμὰν ὁ ψαλτῳδὸς υἱὸς Ἰωὴλ, υἱοῦ Σαμουὴλ,
\vs{19}υἱοῦ Ἑλκανὰ, υἱοῦ Ἱεροβοὰμ, υἱοῦ Ἐλιὴλ, υἱοῦ Θοοὺ,
\vs{20}υἱοῦ Σοὺφ, υἱοῦ Ἑλκανὰ, υἱοῦ Μαὰθ, υἱοῦ Ἀμαθὶ,
\vs{21}υἱοῦ Ἑλκανὰ, υἱοῦ Ἰωὴλ, υἱοῦ Ἀζαρία, υἱοῦ Σαφανία,
\vs{22}υἱοῦ Θαὰθ, υἱοῦ Ἀσὴρ, υἱοῦ Ἀβιασὰφ, υἱοῦ Κορὲ,
\vs{23}υἱοῦ Ἰσαὰρ, υἱοῦ Καὰθ, υἱοῦ Λευὶ, υἱοῦ Ἰσραήλ.
\vs{24}Καὶ ὁ ἀδελφὸς αὐτοῦ Ἀσὰρ ὁ ἑστηκὼς ἐν δεξιᾷ αὐτοῦ· Ἀσὰρ υἱὸς Βαραχία, υἱοῦ Σαμαὰ,
\vs{25}υἱοῦ Μιχαὴλ, υἱοῦ Βαασία, υἱοῦ Μελχία,
\vs{26}υἱοῦ Ἀθανὶ, υἱοῦ Ζααραῒ, υἱοῦ Ἀδαῒ,
\vs{27}υἱοῦ Αἰθὰμ, υἱοῦ Ζαμμὰμ, υἱοῦ Σεμεῒ,
\vs{28}υἱοῦ Ἰεὲθ, υἱοῦ Γεδσὼν, υἱοῦ Λευί.
\vs{29}Καὶ υἱοὶ Μεραρὶ οἱ ἀδελφοὶ αὐτῶν ἐξ ἀριστερῶν· Αἰθὰμ υἱὸς Κισὰ, υἱοῦ Ἀβαῒ, υἱοῦ Μαλῶχ,
\vs{30}υἱοῦ Ἀσεβὶ, υἱοῦ Ἀμεσσία,
\vs{31}υἱοῦ Βανὶ, υἱοῦ Σεμὴρ,
\vs{32}υἱοῦ Μοολὶ, υἱοῦ Μουσὶ, υἱοῦ Μεραρὶ, υἱοῦ Λευί.
\vs{33}Καὶ οἱ ἀδελφοὶ αὐτῶν κατʼ οἴκους πατριῶν αὐτῶν, οἱ Λευῖται οἱ δεδομένοι εἰς πᾶσαν ἐργασίαν λειτουργίας σκηνῆς οἴκου τοῦ Θεοῦ.

\vs{34}Καὶ Ἀαρὼν καὶ υἱοὶ αὐτοῦ θυμιῶντες ἐπὶ τὸ θυσιαστήριον τῶν ὁλοκαυτωμάτων, καὶ ἐπὶ τὸ θυσιαστήριον τῶν θυμιαμάτων εἰς πᾶσαν ἐργασίαν ἅγια τῶν ἁγίων, καὶ ἐξιλάσκεσθαι περὶ Ἰσραὴλ, κατὰ πάντα ὅσα ἐνετείλατο Μωυσῆς παῖς τοῦ Θεοῦ.
\vs{35}Καὶ οὗτοι υἱοὶ Ἀαρών· Ἐλεάζαρ υἱὸς αὐτοῦ, Φινεὲς υἱὸς αὐτοῦ, Ἀβισοὺ υἱὸς αὐτοῦ,
\vs{36}Βοκκὶ υἱὸς αὐτοῦ, Ὀζὶ υἱὸς αὐτοῦ, Σαραῒα υἱὸς αὐτοῦ,
\vs{37}Μαριὴλ υἱὸς αὐτοῦ, Ἀμαρία υἱὸς αὐτοῦ, Ἀχιτὼβ υἱὸς αὐτοῦ,
\vs{38}Σαδὼκ υἱὸς αὐτοῦ, Ἀχιμάας υἱὸς αὐτοῦ.

\vs{39}Καὶ αὗται αἱ κατοικίαι αὐτῶν ἐν ταῖς κώμαις αὐτῶν, ἐν τοῖς ὁρίοις αὐτῶν, τοῖς υἱοῖς Ἀαρὼν τῇ πατριᾷ αὐτῶν τοῖς Κααθὶ, ὅτι αὐτοῖς ἐγένετο ὁ κλῆρος.
\vs{40}Καὶ ἔδωκαν αὐτοῖς τὴν Χεβρὼν ἐν γῇ Ἰούδα, καὶ τὰ περισπόρια αὐτῆς κύκλῳ αὐτῆς.
\vs{41}Καὶ τὰ πεδία τῆς πόλεως, καὶ τὰς κώμας αὐτῆς ἔδωκαν τῷ Χαλὲβ υἱῷ Ἰεφοννή.
\vs{42}Καὶ τοῖς υἱοῖς Ἀαρὼν ἔδωκαν τὰς πόλεις τῶν φυγαδευτηρίων, τὴν Χεβρὼν, καὶ τὴν Λοβνὰ καὶ τὰ περισπόρια αὐτῆς, καὶ τὴν Σελνὰ καὶ τὰ περισπόρια αὐτῆς, καὶ τὴν Ἐσθαμὼ καὶ τὰ περισπόρια αὐτῆς,
\vs{43}καὶ τὴν Ἰεθὰρ καὶ τὰ περισπόρια αὐτῆς, καὶ τὴν Δαβὶρ καὶ τὰ περισπόρια αὐτῆς,
\vs{44}καὶ τὴν Ἀσὰν καὶ τὰ περισπόρια αὐτῆς, καὶ τὴν Βαιθσαμὺς καὶ τὰ περισπόρια αὐτῆς·
\vs{45}Καὶ ἐκ φυλῆς Βενιαμὶν τὴν Γαβαῒ καὶ τὰ περισπόρια αὐτῆς, καὶ τὴν Γαλεμὰθ καὶ τὰ περισπόρια αὐτῆς, καὶ τὴν Ἀναθὼθ καὶ τὰ περισπόρια αὐτῆς· πᾶσαι αἱ πόλεις αὐτῶν τρισκαίδεκα πόλεις κατὰ πατριὰς αὐτῶν.

\vs{46}Καὶ τοῖς υἱοῖς Καὰθ τοῖς καταλοίποις ἐκ τῶν πατριῶν ἐκ τῆς φυλῆς ἐκ τοῦ ἡμίσους φυλῆς Μανασσῆ, κλήρῳ πόλεις δέκα.
\vs{47}Καὶ τοῖς υἱοῖς Γεδσὼν κατὰ πατριὰς αὐτῶν ἐκ φυλῆς Ἰσσάχαρ, ἐκ φυλῆς Ἀσὴρ, ἀπὸ φυλῆς Νεφθαλὶ, ἐκ φυλῆς Μανασσῆ ἐν τῇ Βασὰν, πόλεις τρισκαίδεκα.
\vs{48}Καὶ τοῖς υἱοῖς Μεραρὶ κατὰ πατριὰς αὐτῶν ἐκ φυλῆς Ῥουβὴν, ἐκ φυλῆς Γὰδ, ἐκ φυλῆς Ζαβουλὼν, κλήρῳ πόλεις δεκαδύο.
\vs{49}Καὶ ἔδωκαν οἱ υἱοὶ Ἰσραὴλ τοῖς Λευίταις τὰς πόλεις καὶ τὰ περισπόρια αὐτῶν.
\vs{50}Καὶ ἔδωκαν ἐν κλήρῳ ἐκ φυλῆς υἱῶν Ἰούδα, καὶ ἐκ φυλῆς υἱῶν Συμεὼν, καὶ ἐκ φυλῆς υἱῶν Βενιαμὶν τὰς πόλεις ταύτας ἃς ἐκάλεσαν αὐτὰς ἐπʼ ὀνόματος.

\vs{51}Καὶ ἀπὸ τῶν πατριῶν υἱῶν Καὰθ, καὶ ἐγένοντο πόλεις τῶν ὁρίων αὐτῶν ἐκ φυλῆς Ἐφραίμ.
\vs{52}Καὶ ἔδωκαν αὐτοῖς τὰς πόλεις τῶν φυγαδευτηρίων, τὴν Συχὲμ καὶ τὰ περισπόρια αὐτῆς ἐν ὄρει Ἐφραὶμ, καὶ τὴν Γαζὲρ καὶ τὰ περισπόρια αὐτῆς,
\vs{53}καὶ τὴν Ἰεκμαὰν καὶ τὰ περισπόρια αὐτῆς, καὶ τὴν Βαιθωρὼν καὶ τὰ περισπόρια αὐτῆς,
\vs{54}καὶ τὴν Αἰλὼν καὶ τὰ περισπόρια αὐτῆς, καὶ τὴν Γεθρεμμὼν καὶ τὰ περισπόρια αὐτῆς·
\vs{55}Καὶ ἀπὸ τοῦ ἡμίσους φυλῆς Μανασσῆ τὴν Ἀνὰρ καὶ τὰ περισπόρια αὐτῆς, καὶ τὴν Ἰεμβλάαν καὶ τὰ περισπόρια αὐτῆς, κατὰ πατριὰν τοῖς υἱοῖς Καὰθ τοῖς καταλοίποις.

\vs{56}Τοῖς υἱοῖς Γεδσὼν ἀπὸ πατριῶν ἡμίσους φυλῆς Μανασσῆ τὴν Γωλὰν ἐκ τῆς Βασὰν καὶ τὰ περιπόλια αὐτῆς, καὶ τὴν Ἀσηρὼθ καὶ τὰ περιπόλια αὐτῆς·
\vs{57}Καὶ ἐκ φυλῆς Ἰσσάχαρ τὴν Κέδες καὶ τὰ περισπόρια αὐτῆς, καὶ τὴν Δεβερὶ καὶ τὰ περισπόρια αὐτῆς, καὶ τὴν Δαβὼρ καὶ τὰ περισπόρια αὐτῆς,
\vs{58}καὶ τὴν Ῥαμὼθ, καὶ τὴν Αἰνὰν καὶ τὰ περισπόρια αὐτῆς.
\vs{59}Καὶ ἐκ φυλῆς Ἀσὴρ τὴν Μαασὰλ καὶ τὰ περισπόρια αὐτῆς, καὶ τὴν Ἀβδὼν καὶ τὰ περισπόρια αὐτῆς,
\vs{60}καὶ τὴν Ἀκὰκ καὶ τὰ περισπόρια αὐτῆς, καὶ τὴν Ῥοὼβ καὶ τὰ περισπόρια αὐτῆς·
\vs{61}Καὶ ἀπὸ φυλῆς Νεφθαλὶ τὴν Κέδες ἐν τῇ Γαλιλαίᾳ καὶ τὰ περισπόρια αὐτῆς, καὶ τὴν Χαμὼθ καὶ τὰ περισπόρια αὐτῆς, καὶ τὴν Καριαθαῒμ καὶ τὰ περισπόρια αὐτῆς.

\vs{62}Τοῖς υἱοῖς Μεραρὶ τοῖς καταλοίποις ἐκ φυλῆς Ζαβουλὼν τὴν Ῥεμμὼν καὶ τὰ περισπόρια αὐτῆς, καὶ τὴν Θαβὼρ καὶ τὰ περισπόρια αὐτῆς,
\vs{63}ἐκ τοῦ πέραν τοῦ Ἰορδάνου τὴν Ἱεριχὼ κατὰ δυσμὰς τοῦ Ἰορδάνου· ἐκ φυλῆς Ῥουβὴν τὴν Βοσὸρ ἐν τῇ ἐρήμῳ καὶ τὰ περισπόρια αὐτῆς, καὶ τὴν Ἰασὰ καὶ τὰ περισπόρια αὐτῆς,
\vs{64}καὶ τὴν Καδμὼθ καὶ τὰ περισπόρια αὐτῆς, καὶ τὴν Μαεφλὰ καὶ τὰ περισπόρια αὐτῆς·
\vs{65}Ἐκ φυλῆς Γὰδ τὴν Ῥαμμὼθ Γαλαὰδ καὶ τὰ περισπόρια αὐτῆς, καὶ τὴν Μααναῒμ καὶ τὰ περισπόρια αὐτῆς,
\vs{66}καὶ τὴν Ἐσεβὼν καὶ τὰ περισπόρια αὐτῆς, καὶ τὴν Ἰαζὴρ καὶ τὰ περισπόρια αὐτῆς.

\ch{7}
Καὶ τοῖς υἱοῖς Ἰσσάχαρ, Θωλὰ, καὶ Φουὰ, καὶ Ἰασοὺβ, καὶ Σεμερὼν, τέσσαρες.
\vs{2}Καὶ υἱοὶ Θωλὰ, Ὀζὶ, Ῥαφαΐα, καὶ ʼΙεριὴλ, καὶ ʼΙαμαῒ, καὶ ʼΙεμασὰν, καὶ Σαμουὴλ, ἄρχοντες οἴκων πατριῶν αὐτῶν τῷ Θωλὰ, ἰσχυροὶ δυνάμει κατὰ γενέσεις αὐτῶν, ὁ ἀριθμὸς αὐτῶν ἐν ἡμέραις Δαυὶδ, εἴκοσι καὶ δύο χιλιάδες καὶ ἑξακόσιοι.
\vs{3}Καὶ υἱοὶ ʼΟζὶ, ʼΙεζραῒα· καὶ υἱοὶ ʼΙεζραΐ, Μιχαὴλ, ʼΑβδιοὺ, καὶ ʼΙωὴλ, καὶ ʼΙεσία, πέντε, ἄρχοντες πάντες.

\vs{4}Καὶ ἐπʼ αὐτῶν, κατὰ γενέσεις αὐτῶν, κατʼ οἴκους πατριῶν αὐτῶν, ἰσχυροὶ παρατάξασθαι εἰς πόλεμον, τριάκοντα καὶ ἓξ χιλιάδες, ὅτι ἐπλήθυναν γυναῖκας καὶ υἱούς.
\vs{5}Καὶ ἀδελφοὶ αὐτῶν εἰς πάσας πατριὰς Ἰσσάχαρ, καὶ ἰσχυροὶ δυνάμει, ὀγδοήκοντα καὶ ἑπτὰ χιλιάδες, ὁ ἀριθμὸς αὐτῶν τῶν πάντων.

\vs{6}Υἱοὶ Βενιαμὶν, καὶ Βαλὲ, καὶ Βαχὶρ, καὶ ʼΙεδιὴλ, τρεῖς.
\vs{7}Καὶ υἱοὶ Βαλὲ, ʼΕσεβὼν, καὶ Ὀζὶ, καὶ Ὀζιὴλ, καὶ Ἰεριμοὺθ, καὶ Οὐρὶ, πέντε, ἄρχοντες οἴκων πατρικῶν ἰσχυροὶ δυνάμει· καὶ ὁ ἀριθμὸς αὐτῶν, εἴκοσι καὶ δύο χιλιάδες καὶ τριακοντατέσσαρες.
\vs{8}Καὶ υἱοὶ Βαχὶρ, Ζεμιρὰ, καὶ Ἰωὰς, καὶ Ἐλιέζερ, καὶ Ἐλιθενὰν, καὶ Ἀμαρία, καὶ ʼΙεριμοὺθ, καὶ Ἀβιοὺδ, καὶ Ἀναθὼθ, καὶ ʼΕληεμέθ· πάντες οὗτοι υἱοὶ Βαχίρ.
\vs{9}Καὶ ὁ ἀριθμὸς αὐτῶν κατὰ γενέσεις αὐτῶν, ἄρχοντες οἴκων πατριῶν αὐτῶν ἰσχυροὶ δυνάμει, εἴκοσι χιλιάδες καὶ διακόσιοι.
\vs{10}Καὶ υἱοὶ ʼΙεδιὴλ, Βαλαάν· καὶ υἱοὶ Βαλαὰν, Ἰαοὺς, καὶ Βενιαμὶν, καὶ Ἀὼθ, καὶ Χανανὰ, καὶ Ζαιθὰν, καὶ Θαρσὶ, καὶ Ἀχισαάρ.
\vs{11}Πάντες οὗτοι υἱοὶ ʼΙεδιὴλ, ἄρχοντες τῶν πατριῶν ἰσχυροὶ δυνάμει, ἑπτα καίδεκα χιλιάδες καὶ διακόσιοι, ἐκπορευόμενοι δυνάμει πολεμεῖν.
\vs{12}Καὶ Σαπφὶν, καὶ Ἀπφὶν, καὶ υἱοὶ Ὢρ, ʼΑσὼμ, υἱὸς αὐτοῦ Ἀόρ.

\vs{13}Υἱοὶ Νεφθαλὶ, Ἰασιὴλ, Γωνὶ, καὶ Ἀσὴρ, καὶ Σελλοὺμ, υἱοὶ αὐτοῦ, Βαλὰμ υἱὸς αὐτοῦ.

\vs{14}Υἱοὶ Μανασσῆ, Ἐσριὴλ, ὃν ἔτεκεν ἡ παλλακὴ αὐτοῦ ἡ Σύρα, ἔτεκε δὲ αὐτῷ καὶ Μαχὶρ πατέρα Γαλαάδ.
\vs{15}Καὶ Μαχὶρ ἔλαβε γυναῖκα τῷ Ἀπφὶν καὶ Σαπφίν· καὶ ὄνομα ἀδελφῆς αὐτοῦ Μοωχὰ, καὶ ὄνομα τῷ δευτέρῳ Σαπφαάδ· ἐγεννήθησαν δὲ τῷ Σαπφαὰδ θυγατέρες.
\vs{16}Καὶ ἔτεκε Μοωχὰ γυνὴ Μαχὶρ υἱὸν, καὶ ἐκάλεσε τὸ ὄνομα αὐτοῦ Φαρές· καὶ ὄνομα ἀδελφοῦ αὐτοῦ Σοῦρος· υἱοὶ αὐτοῦ Οὐλὰμ, καὶ ʼΡοκόμ.
\vs{17}Καὶ υἱοὶ Οὐλὰμ, Βαδάμ· οὗτοι υἱοὶ Γαλαὰδ, υἱοῦ Μαχὶρ, υἱοῦ Μανασσῆ.
\vs{18}Καὶ ἡ ἀδελφὴ αὐτοῦ ἡ Μαλεχὲθ ἔτεκε τὸν Ἰσοὺδ, καὶ τὸν Ἀβιέζερ, καὶ τὸν Μαελά.
\vs{19}Καὶ ἦσαν υἱοὶ Σεμιρὰ, ʼΑῒμ, καὶ Συχὲμ, καὶ Λακὶμ, καὶ Ἀνιάν.

\vs{20}Καὶ υἱοὶ Ἐφραὶμ, Σωθαλὰθ, καὶ Βαρὰδ υἱὸς αὐτοῦ, καὶ Θαὰθ υἱὸς αὐτοῦ, ʼΕλαδὰ υἱὸς αὐτοῦ, Σαὰθ υἱὸς αὐτοῦ,
\vs{21}καὶ Ζαβὰδ υἱὸς αὐτοῦ, Σωθελὲ υἱὸς αὐτοῦ, καὶ ʼΑζὲρ, καὶ ʼΕλεάδ· καὶ ἀπέκτειναν αὐτοὺς οἱ ἄνδρες Γὲθ οἱ τεχθέντες ἐν τῇ γῇ, ὅτι κατέβησαν τοῦ λαβεῖν τὰ κτήνη αὐτῶν.
\vs{22}Καὶ ἐπένθησεν Ἐφραὶμ ὁ πατὴρ αὐτῶν ἡμέρας πολλάς· καὶ ἦλθον ἀδελφοὶ αὐτοῦ τοῦ παρακαλέσαι αὐτόν.
\vs{23}Καὶ εἰσῆλθε πρὸς τὴν γυναῖκα αὐτοῦ, καὶ ἔλαβεν ἐν γαστρὶ, καὶ ἔτεκεν υἱόν· καὶ ἐκάλεσε τὸ ὄνομα αὐτοῦ Βεριὰ, ὅτι ἐν κακοῖς ἐγένετο ἐν οἴκῳ μου.
\vs{24}Καὶ ἡ θυγάτηρ αὐτοῦ Σαραά· καὶ ἐν ἐκείνοις τοῖς καταλοίποις· καὶ ᾠκοδόμησε τὴν Βαιθωρὼν τὴν κάτω καὶ τὴν ἄνω· καὶ υἱοὶ Ὀζὰν Σεηρὰ,
\vs{25}καὶ Ῥαφὴ υἱὸς αὐτοῦ, Σαρὰφ καὶ Θαλεὲς υἱοὶ αὐτοῦ, Θαὲν υἱὸς αὐτοῦ.
\vs{26}Τῷ Λααδὰν υἱῷ αὐτοῦ υἱὸς Ἀμιοὺδ, υἱὸς Ἐλισαμαῒ,
\vs{27}υἱὸς Νοὺν, υἱὸς Ἰησουὲ, υἱοὶ αὐτοῦ.

\vs{28}Καὶ κατάσχεσις αὐτῶν καὶ κατοικία αὐτῶν Βαιθὴλ καὶ αἱ κῶμαι αὐτῆς, κατʼ ἀνατολὰς Νοαρὰν, πρὸς δυσμαῖς Γάζερ καὶ αἱ κῶμαι αὐτῆς, καὶ Συχὲμ καὶ αἱ κῶμαι αὐτῆς ἕως Γάζης, καὶ αἱ κῶμαι αὐτῆς,
\vs{29}καὶ ἕως ὁρίων υἱῶν Μανασσῆ, Βαιθσαὰν καὶ αἱ κῶμαι αὐτῆς, Θανὰχ καὶ αἱ κῶμαι αὐτῆς, Μαγεδδὼ καὶ αἱ κῶμαι αὐτῆς, Δὼρ καὶ αἱ κῶμαι αὐτῆς· ἐν ταύτῃ κατῴκησαν υἱοὶ Ἰωσὴφ υἱοῦ Ἰσραήλ.

\vs{30}Υἱοὶ Ἀσὴρ, Ἰεμνὰ, καὶ Σουΐα, καὶ Ἰσουῒ, καὶ Βεριὰ, καὶ Σορὲ ἀδελφὴ αὐτῶν.
\vs{31}Καὶ υἱοὶ Βεριὰ, Χάβερ, καὶ Μελχιήλ· οὗτος πατὴρ Βερθαΐθ.
\vs{32}Καὶ Χάβερ ἐγέννησε τὸν Ἰαφλὴτ, καὶ τὸν Σαμὴρ, καὶ τὸν Χωθὰν, καὶ τὴν Σωλὰ ἀδελφὴν αὐτῶν.
\vs{33}Καὶ υἱοὶ Ἰαφλὴτ, Φασὲκ, καὶ Βαμαὴλ, καὶ Ἀσίθ· οὗτοι υἱοὶ Ἰαφλήτ.
\vs{34}Καὶ υἱοὶ Σεμμὴρ, Ἀχὶρ, καὶ Ῥοογὰ, καὶ ʼΙαβὰ, καὶ Αρὰμ,
\vs{35}καὶ Βανὴ ʼΕλὰμ ἀδελφοῦ αὐτοῦ Σωφὰ, καὶ ʼΙμανὰ, καὶ Σελλὴς, καὶ ʼΑμάλ.
\vs{36}Υἱοὶ Σωφὰς, Σουὲ, καὶ Ἁρναφὰρ, καὶ Σουδὰ, καὶ Βαρὶν, καὶ ʼΙμρὰν,
\vs{37}καὶ Βασὰν, καὶ ʼΩὰ, καὶ Σαμὰ, καὶ Σαλισὰ, καὶ ʼΙεθρὰ, καὶ Βεηρά.
\vs{38}Καὶ υἱοὶ Ἰεθὴρ, Ἰεφινὰ, καὶ Φασφὰ, καὶ Ἀρά.
\vs{39}καὶ υἱοὶ ʼΟλὰ, Ὀρὲχ, Ἀνιὴλ, καὶ Ῥασιά.

\vs{40}Πάντες οὗτοι υἱοὶ Ἀσὴρ, πάντες ἄρχοντες πατριῶν, ἐκλεκτοὶ ἰσχυροὶ δυνάμει, ἄρχοντες ἡγούμενοι· ὁ ἀριθμὸς αὐτῶν εἰς παράταξιν τοῦ πολεμεῖν, ἀριθμὸς αὐτῶν ἄνδρες εἰκοσιὲξ χιλιάδες.

\ch{8}
Καὶ Βενιαμὶν ἐγέννησε Βαλὲ πρωτότοκον αὐτοῦ, καὶ ʼΑσβὴλ τὸν δεύτερον, Ἀαρὰ τὸν τρίτον,
\vs{2}Νωὰ τὸν τέταρτον, καὶ Ῥαφὰ τὸν πέμπτον.
\vs{3}Καὶ ἦσαν υἱοὶ τῷ Βαλὲ, Ἀδὶρ, καὶ Γηρὰ, καὶ Ἀβιοὺδ,
\vs{4}καὶ Ἀβεσσουὲ, καὶ Νοαμὰ, καὶ Ἀχιὰ,
\vs{5}καὶ Γερὰ, καὶ Σεφουφὰμ, καὶ Οὐράμ.
\vs{6}Οὗτοι υἱοὶ Ἀὼδ, οὗτοί εἰσιν ἄρχοντες πατριῶν τοῖς κατοικοῦσι Γαβεέ· καὶ μετῴκισαν αὐτοὺς εἰς Μαχαναθὶ,
\vs{7}καὶ Νοομὰ, καὶ Ἀχιὰ, καὶ Γηρά· οὗτος ἰεγλαὰμ, καὶ ἐγέννησε τὸν ʼΑζὰ, καὶ τὸν Ἰαχιχώ.

\vs{8}Καὶ Σααρὶν ἐγέννησεν ἐν τῷ πεδίῳ Μωὰβ μετὰ τὸ ἀποστεῖλαι αὐτὸν ʼΩσὶν καὶ τὴν Βααδὰ γυναῖκα αὐτοῦ.

\vs{9}Καὶ ἐγέννησεν ἐκ τῆς ʼΑδὰ γυναικὸς αὐτοῦ τὸν Ἰωλὰβ, καὶ τὸν Σεβιὰ, καὶ τὸν Μισὰ, καὶ τὸν Μελχὰς,
\vs{10}καὶ τὸν Ἰεβοὺς, καὶ τὸν Ζαβιὰ καὶ τὸν Μαρμά· οὗτοι ἄρχοντες πατριῶν.
\vs{11}Καὶ ἐκ τῆς Ὠσὶν ἐγέννησε τὸν Ἀβιτὼλ, καὶ τὸν Ἀλφαάλ.
\vs{12}Καὶ υἱοὶ Ἀλφαὰλ, Ὠβὴδ, Μισαὰλ, Σεμμήρ· οὗτος ᾠκοδόμησε τὴν Ὠνὰν, καὶ τὴν Λὼδ καὶ τὰς κώμας αὐτῆς·
\vs{13}Καὶ Βεριὰ, καὶ Σαμά· οὗτοι ἄρχοντες τῶν πατριῶν τοῖς κατοικοῦσιν Αἰλὰμ, καὶ οὗτοι ἐξεδίωξαν τοὺς κατοικοῦντας Γέθ.
\vs{14}καὶ ἀδελφὸς αὐτοῦ Σωσὴκ, καὶ ʼΑριμὼθ,
\vs{15}καὶ Ζαβαδία, καὶ Ὠρὴδ, καὶ Ἔδερ,
\vs{16}καὶ Μιχαὴλ, καὶ ʼΙεσφὰ, καὶ Ἰωδὰ, υἱοὶ Βεριά.
\vs{17}Καὶ Ζαβαδία, καὶ Μοσολλὰμ, καὶ Ἀζακὶ, καὶ Ἀβὰρ,
\vs{18}καὶ Ἰσαμαρὶ, καὶ Ἰεξλίας, καὶ Ἰωβὰβ, υἱοὶ Ἐλφαάλ.
\vs{19}Καὶ Ἰακὶμ, καὶ Ζαχρὶ, καὶ Ζαβδὶ,
\vs{20}καὶ Ἐλιωναῒ, καὶ Σαλαθὶ, καὶ Ἐλιηλὶ,
\vs{21}καὶ Ἀδαΐα, καὶ Βαραΐα, καὶ Σαμαρὰθ, υἱοὶ Σαμαΐθ.
\vs{22}καὶ Ἰεσφὰν, καὶ Ὠβὴδ, καὶ Ἐλεὴλ,
\vs{23}καὶ Ἀβδὼν, καὶ Ζεχρὶ, καὶ Ἀνὰν,
\vs{24}καὶ Ἀνανία, καὶ Ἀμβρὶ, καὶ Αἰλὰμ, καὶ Ἀναθὼθ,
\vs{25}καὶ Ἰαθὶν, καὶ Ἰεφαδίας, καὶ Φανουὴλ, υἱοὶ Σωσήκ.
\vs{26}Καὶ Σαμσαρὶ, καὶ Σααρίας, καὶ Γοθολία,
\vs{27}καὶ Ἰαρασία, καὶ Ἐριὰ, καὶ Ζεχρὶ υἱὸς Ἰροάμ.
\vs{28}Οὗτοι ἄρχοντες πατριῶν κατὰ γενέσεις αὐτῶν ἄρχοντες· οὗτοι κατῴκησαν ἐν Ἱερουσαλήμ.

\vs{29}Καὶ ἐν Γαβαὼν κατῴκησε πατὴρ Γαβαών· καὶ ὄνομα γυναικὶ αὐτοῦ Μοαχά.
\vs{30}Καὶ ὁ υἱὸς αὐτῆς ὁ πρωτότοκος Ἀβδὼν, καὶ Σοὺρ, καὶ Κὶς, καὶ Βαὰλ, καὶ Ναδὰβ,
\vs{31}καὶ Νὴρ, καὶ Γεδοὺρ καὶ ἀδελφὸς αὐτοῦ, καὶ Ζακχοὺρ, καὶ Μακελώθ.
\vs{32}Καὶ Μακελὼθ ἐγέννησε τὸν Σαμαά· καὶ γὰρ οὗτοι κατέναντι τῶν ἀδελφῶν αὐτῶν κατῴκησαν ἐν Ἱερουσαλὴμ μετὰ τῶν ἀδελφῶν αὐτῶν.
\vs{33}Καὶ Νὴρ ἐγέννησε τὸν Κὶς, καὶ Κὶς ἐγέννησε τὸν Σαοὺλ, καὶ Σαοὺλ ἐγέννησε τὸν Ἰωνάθαν, καὶ τὸν Μελχισουὲ, καὶ τὸν Ἀμιναδὰβ, καὶ τὸν Ἀσαβάλ.
\vs{34}καὶ υἱὸς Ἰωνάθαν Μεριβαάλ· καὶ Μεριβαὰλ ἐγέννησε τὸν Μιχά.
\vs{35}καὶ υἱοὶ Μιχὰ, Φιθὼν, καὶ Μελὰχ, καὶ Θαρὰχ, καὶ Ἀχάζ.
\vs{36}καὶ Ἀχὰζ ἐγέννησε τὸν Ἰαδά· καὶ Ἰαδὰ ἐγέννησε τὸν Σαλαιμὰθ, καὶ τὸν Ἀσμὼθ, καὶ τὸν Ζαμβρί· καὶ Ζαμβρὶ ἐγέννησε τὸν Μαισά.
\vs{37}Καὶ Μαισὰ ἐγέννησε τὸν Βαανά· Ῥαφαία υἱὸς αὐτοῦ, Ἐλασὰ υἱὸς αὐτοῦ, Ἐσὴλ υἱὸς αὐτοῦ.

\vs{38}Καὶ τῷ Ἐσὴλ ἓξ υἱοί· καὶ ταῦτα τὰ ὀνόματα αὐτῶν· Ἐζρικὰμ πρωτότοκος αὐτοῦ, καὶ Ἰσμαὴλ, καὶ Σαραΐα, καὶ Ἀβδία, καὶ Ἀνὰν, καὶ Ἀσά· πάντες οὗτοι υἱοὶ Ἐσήλ.
\vs{39}Καὶ υἱοὶ Ἀσὴλ ἀδελφοῦ αὐτοῦ, Αἰλὰμ πρωτότοκος αὐτοῦ, καὶ Ἰὰς ὁ δεύτερος, καὶ Ἐλιφαλὲτ ὁ τρίτος.
\vs{40}Καὶ ἦσαν υἱοὶ Αἰλὰμ ἰσχυροὶ ἄνδρες δυνάμει, τείνοντες τόξον, καὶ πληθύνοντες υἱοὺς καὶ υἱοὺς τῶν υἱῶν, ἑκατὸν πεντήκοντα· πάντες οὗτοι ἐξ υἱῶν Βενιαμίν.

\ch{9}
Καὶ πᾶς Ἰσραὴλ ὁ συλλοχισμὸς αὐτῶν· καὶ οὗτοι καταγεγραμμένοι ἐν βιβλίῳ τῶν βασιλέων Ἰσραὴλ καὶ Ἰούδα, μετὰ τῶν ἀποικισθέντων εἰς Βαβυλῶνα ἐν ταῖς ἀνομίαις αὐτῶν,
\vs{2}καὶ οἱ κατοικοῦντες πρότερον ἐν ταῖς κατασχέσεσιν αὐτῶν ἐν ταῖς πόλεσιν Ἰσραὴλ, οἱ ἱερεῖς, οἱ Λευῖται, καὶ οἱ δεδομένοι.

\vs{3}Καὶ ἐν Ἱερουσαλὴμ κατῴκησαν ἀπὸ τῶν υἱῶν Ἰούδα, καὶ ἀπὸ τῶν υἱῶν Βενιαμὶν, καὶ ἀπὸ τῶν υἱῶν Ἐφραῒμ, καὶ Μανασσῆ.
\vs{4}Καὶ Γνωθὶ, καὶ υἱὸς Σαμιοὺδ, υἱοῦ Ἀμρὶ, υἱοῦ Ἀμβραῒμ, υἱοῦ Βουνὶ, υἱοῦ υἱῶν Φαρὲς, υἱοῦ Ἰούδα.
\vs{5}Καὶ ἐκ τῶν Σηλωνὶ, Ἀσαΐα πρωτότοκος αὐτοῦ, καὶ ὁ ὑοὶ αὐτοῦ.
\vs{6}Ἐκ τῶν υἱῶν Ζαρὰ, Ἰεὴλ, καὶ ἀδελφοὶ αὐτῶν ἑξακόσιοι καὶ ἐννενήκοντα.

\vs{7}Καὶ ἐκ τῶν υἱῶν Βενιαμὶν Σαλὼμ υἱὸς Μοσλλὰμ, υἱοῦ Ὠδουΐα, υἱοῦ Ἀσινοῦ,
\vs{8}καὶ Ἰεμναὰ υἱὸς Ἱεροβοὰμ, καὶ Ἠλώ· οὗτοι υἱοὶ Ὀζὶ υἱοῦ Μαχίρ· καὶ Μοσολλὰμ υἱὸς Σαφατία, υἱοῦ Ῥαγουὴλ, υἱοῦ Ἰεμναῒ,
\vs{9}καὶ ἀδελφοὶ αὐτῶν κατὰ γενέσεις αὐτῶν ἐννακόσιοι πεντηκονταὲξ, πάντες οἱ ἄνδρες ἄρχοντες πατριῶν κατʼ οἴκους πατριῶν αὐτῶν.

\vs{10}Καὶ ἀπὸ τῶν ἱερέων, Ἰωδαὲ, καὶ Ἰωαρὶμ, καὶ Ἰαχὶν,
\vs{11}καὶ Ἀζαρία υἱὸς Χελκία υἱοῦ Μοσολλὰμ, υἱοῦ Σαδὼκ, υἱοῦ Μαραϊὼθ, υἱοῦ Ἀχιτὼβ ἡγουμένου οἴκου τοῦ Θεοῦ,
\vs{12}καὶ Ἀδαΐα υἱὸς Ἰραὰμ, υἱοῦ Φασχὼρ, υἱοῦ Μελχία, καὶ Μαασαία υἱὸς Ἀδιὴλ, υἱοῦ Ἐζιρὰ, υἱοῦ Μοσολλὰμ, υἱοῦ Μασελμὼθ, υἱοῦ Ἐμμὴρ,
\vs{13}καὶ ἀδελφοὶ αὐτῶν ἄρχοντες οἴκων πατριῶν αὐτῶν, χίλιοι καὶ ἑπτακόσιοι καὶ ἑξήκοντα, ἰσχυροὶ δυνάμει εἰς ἐργασίαν λειτουργίας οἴκου τοῦ Θεοῦ.

\vs{14}Καὶ ἐκ τῶν Λευιτῶν, Σαμαΐα υἱὸς Ἀσὼβ, υἱοῦ Ἐζρικὰμ, υἱοῦ Ἀσαβία, ἐκ τῶν υἱῶν Μεραρί.
\vs{15}Καὶ Βακβακὰρ, καὶ Ἀρὴς, καὶ Γαλαὰλ, καὶ Ματθανίας υἱὸς Μιχὰ, υἱοῦ Ζεχρὶ, υἱοῦ Ἀσάφ.
\vs{16}Καὶ Ἀβδία υἱὸς Σαμία, υἱοῦ Γαλαὰλ, υἱοῦ Ἰδιθοὺν, καὶ Βαραχία υἱὸς Ὀσσὰ, υἱοῦ Ἑλκανὰ, ὁ κατοικῶν ἐν ταῖς κώμαις Νωτεφατί.
\vs{17}Οἱ πυλωροὶ, Σαλὼμ, Ἀκοὺμ, Τελμὼν, καὶ Διμὰν, καὶ ἀδελφοὶ αὐτῶν· Σαλὼμ ὁ ἄρχων,
\vs{18}καὶ ἕως ταύτης ἐν τῇ πύλῃ τοῦ βασιλέως κατʼ ἀνατολάς· αὗται αἱ πύλαι τῶν παρεμβολῶν υἱῶν Λευί.
\vs{19}Καὶ Σελλοὺμ υἱὸς Κορὲ, υἱοῦ Ἀβιασὰφ, υἱοῦ Κορέ· καὶ οἱ ἀδελφοὶ αὐτοῦ εἰς οἶκον πατρὸς αὐτοῦ, οἱ Κορῖται ἐπὶ τῶν ἔργων τῆς λειτουργίας φυλάσσοντες τὰς φυλακὰς τῆς σκηνῆς· καὶ πατέρες αὐτῶν ἐπὶ τῆς παρεμβολῆς Κυρίου φυλάσσοντες τὴν εἴσοδον.

\vs{20}Καὶ Φινεὲς υἱὸς Ἐλεάζαρ ἡγούμενος ἦν ἐπʼ αὐτῶν ἔμπροσθεν Κυρίου, καὶ οὗτοι μετʼ αὐτοῦ.
\vs{21}Ζαχαρίας υἱὸς Μοσολλαμὶ πυλωρὸς τῆς θύρας τῆς σκηνῆς τοῦ μαρτυρίου.
\vs{22}Πάντες οἱ ἐκλεκτοὶ ἐπὶ τῆς πύλης ἐν ταῖς πύλαις διακόσιοι καὶ δεκαδύο· οὗτοι ἐν ταῖς αὐλαῖς αὐτῶν, ὁ καταλοχισμὸς αὐτῶν· τούτους ἔστησε Δαυὶδ καὶ Σαμουὴλ ὁ βλέπων τῇ πίστει αὐτῶν.
\vs{23}Καὶ οὗτοι καὶ οἱ υἱοὶ αὐτῶν ἐπὶ τῶν πύλων ἐν οἴκῳ Κυρίου, καὶ ἐν οἴκῳ τῆς σκηνῆς τοῦ φυλάσσειν.
\vs{24}Κατὰ τοὺς τέσσαρας ἀνέμους ἦσαν αἱ πύλαι, κατὰ ἀνατολὰς, θάλασσαν, Βοῤῥᾶν, Νότον.
\vs{25}Καὶ ἀδελφοὶ αὐτῶν ἐν ταῖς αὐλαῖς αὐτῶν τοῦ εἰσπορεύεσθαι κατὰ ἑπτὰ ἡμέρας ἀπὸ καιροῦ εἰς καιρὸν μετὰ τούτων·
\vs{26}Ὅτι ἐν πίστει εἰσὶ τέσσαρες δυνατοὶ τῶν πυλῶν· καὶ οἱ Λευῖται ἦσαν ἐπὶ τῶν παστοφορίων, καὶ ἐπὶ τῶν θησαυρῶν οἴκου τοῦ Θεοῦ παρεμβάλλουσιν,
\vs{27}ὅτι ἐπʼ αὐτοὺς ἡ φυλακή· καὶ οὗτοι ἐπὶ τῶν κλειδῶν τοπρωῒ πρωῒ ἀνοίγειν τὰς θύρας τοῦ ἱεροῦ.

\vs{28}Καὶ ἐξ αὐτῶν ἐπὶ τὰ σκεύη τῆς λειτουργίας, ὅτι ἐν ἀριθμῷ εἰσοίσουσι, καὶ ἐν ἀριθμῷ ἐξόσουσι.
\vs{29}Καὶ ἐξ αὐτῶν καθεσταμένοι ἐπὶ τὰ σκεύη, καὶ ἐπὶ πάντα σκεύη τὰ ἅγια, καὶ ἐπὶ τῆς σεμιδάλεως, τοῦ οἴνου, τοῦ ἐλαίου, τοῦ λιβανωτοῦ, καὶ τῶν ἀρωμάτων.
\vs{30}Καὶ ἀπὸ τῶν υἱῶν τῶν ἱερέων ἦσαν μυρεψοὶ τοῦ μύρου, καὶ εἰς τὰ ἀρώματα.
\vs{31}Καὶ Ματθαθίας ἐκ τῶν Λευιτῶν, οὗτος ὁ πρωτότοκος τῷ Σαλὼμ τῷ Κορείτῃ, ἐν τῇ πίστει ἐπὶ τὰ ἔργα τῆς θυσίας τοῦ τηγάνου τοῦ μεγάλου ἱερέως.
\vs{32}Καὶ Βαναΐας ὁ Κααθίτης ἐκ τῶν ἀδελφῶν αὐτῶν ἐπὶ τῶν ἄρτων τῆς προθέσεως, τοῦ ἑτοιμάσαι σάββατον κατὰ σάββατον.
\vs{33}Καὶ οὗτοι ψαλτῳδοὶ ἄρχοντες τῶν πατριῶν τῶν Λευιτῶν διατεταγμέναι ἐφημερίαι, ὅτι ἡμέρα καὶ νὺξ ἐπʼ αὐτοῖς ἐν τοῖς ἔργοις.
\vs{34}Οὗτοι ἄρχοντες τῶν πατριῶν τῶν Λευιτῶν κατὰ γενέσεις αὐτῶν, ἄρχοντες οὗτοι κατῴκησαν ἐν Ἱερουσαλήμ.

\vs{35}Καὶ ἐν Γαβαὼν κατῴκησε πατὴρ Γαβαὼν Ἰεήλ· καὶ ὄνομα γυναικὸς αὐτοῦ Μοωχά.
\vs{36}Καὶ υἱὸς αὐτοῦ ὁ πρωτότοκος Ἀβδὼν, καὶ Σοὺρ, καὶ Κὶς, καὶ Βαὰλ, καὶ Νὴρ, καὶ Ναδὰβ,
\vs{37}καὶ Γεδοὺρ καὶ ἀδελφὸς, καὶ Ζακχοὺρ, καὶ Μακελώθ.
\vs{38}Καὶ Μακελὼθ ἐγέννησε τὸν Σαμαά· καὶ οὗτοι ἐν μέσῳ τῶν ἀδελφῶν αὐτῶν κατῴκησαν ἐν Ἱερουσαλὴμ, ἐν μέσῳ τῶν ἀδελφῶν αὐτῶν.

\vs{39}Καὶ Νὴρ ἐγέννησε τὸν Κὶς, καὶ Κὶς ἐγέννησε τὸν Σαοὺλ, καὶ Σαοὺλ ἐγέννησε τὸν Ἰωνάθαν, καὶ τὸν Μελχισουὲ, καὶ τὸν Ἀμιναδὰβ, καὶ τὸν Ἀσαβάλ.
\vs{40}Καὶ υἱὸς Ἰωνάθαν Μεριβάλ· καὶ Μεριβαὰλ ἐγέννησε τὸν Μιχά.
\vs{41}Καὶ υἱοὶ Μιχὰ Φιθὼν, καὶ Μαλὰχ, καὶ Θαράχ.
\vs{42}Καὶ Ἀχὰζ ἐγέννησε τὸν Ἰαδά· καὶ Ἰαδὰ ἐγέννησε τὸν Γαλεμὲθ, καὶ τὸν Γαζμὼθ, καὶ τὸν Ζαμβρί· και Ζαμβρὶ ἐγέννησε τὸν Μασσά.
\vs{43}Καὶ Μασσὰ ἐγέννησε τὸν Βαανὰ, καὶ Ῥαφαΐα υἱὸς αὐτοῦ, Ἐλασὰ υἱὸς αὐτοῦ, Ἐσὴλ υἱὸς αὐτοῦ.
\vs{44}Καὶ τῷ Ἐσὴλ ἓξ υἱοί· καὶ ταῦτα τὰ ὀνόματα αὐτῶν· Ἐζρικὰμ πρωτότοκος αὐτοῦ, καὶ Ἰσμαὴλ, καὶ Σαραΐα, καὶ Ἀβδία, καὶ Ἀνὰν, καὶ Ἀσά· οὗτοι υἱοὶ Ἐσήλ.

\ch{10}
Καὶ ἀλλόφυλοι ἐπολέμησαν πρὸς τὸν Ἰσραὴλ, καὶ ἔφυγον ἀπὸ προσώπου ἀλλοφύλων, καὶ ἔπεσον τραυματίαι ἐν ὄρει Γελβουέ.
\vs{2}Καὶ κατεδίωξαν οἱ ἀλλόφυλοι ὀπίσω Σαοὺλ καὶ ὀπίσω τῶν υἱῶν αὐτοῦ· καὶ ἐπάταξαν ἀλλόφυλοι τὸν Ἰωνάθαν, καὶ τὸν Ἀμιναδὰβ, καὶ τὸν Μελχισουὲ, υἱοὺς Σαούλ.
\vs{3}Καὶ ἐβαρύνθη ὁ πόλεμος ἐπὶ Σαούλ· καὶ εὗρον αὐτὸν οἱ τοξόται ἐν τόξοις καὶ πόνοις, καὶ ἐπονέσαν ἀπὸ τῶν τόξων.
\vs{4}Καὶ εἶπε Σαοὺλ τῷ αἴροντι τὰ σκεύη αὐτοῦ, σπάσαι τὴν ῥομφαίαν σου, καὶ ἐκκέντησόν με ἐν αὐτῇ, μὴ ἔλθωσιν οἱ ἀπερίτμητοι οὗτοι καὶ ἐμπαίξωσί μοι· καὶ οὐκ ἐβούλετο ὁ αἴρων τὰ σκεύη αὐτοῦ, ὅτι ἐφοβεῖτο σφόδρα· καὶ ἔλαβε Σαοὺλ τὴν ῥομφαίαν, καὶ ἐπέπεσεν ἐπʼ αὐτήν.
\vs{5}Καὶ εἶδεν ὁ αἴρων τὰ σκεύη αὐτοῦ, ὅτι ἀπέθανε Σαοὺλ, καὶ ἔπεσε καὶ αὐτὸς ἐπὶ τὴν ῥομφαίαν αὐτοῦ.
\vs{6}Καὶ ἀπέθανε Σαοὺλ, καὶ τρεῖς υἱοὶ αὐτοῦ ἐν τῇ ἡμέρᾳ ἐκείνῃ· καὶ πᾶς ὁ οἶκος αὐτοῦ ἐπὶ τὸ αὐτὸ ἀπέθανε.
\vs{7}Καὶ εἶδε πᾶς ἀνὴρ Ἰσραὴλ ὁ ἐν τῷ αὐλῶνι ὅτι ἔφυγεν Ἰσραὴλ, καὶ ὅτι ἀπέθανε Σαοὺλ καὶ οἱ υἱοὶ αὐτοῦ, καὶ κατέλιπον τὰς πόλεις αὐτῶν καὶ ἔφυγον· καὶ ἦλθον ἀλλόφυλοι, καὶ κατῴκησαν ἐν αὐταῖς.

\vs{8}Καὶ ἐγένετο τῇ ἐχομένῃ, καὶ ἦλθον ἀλλόφυλοι τοῦ σκυλεύειν τοὺς τραυματίας, καὶ εὗρον τὸν Σαοὺλ καὶ τοὺς υἱοὺς αὐτοῦ πεπτωκότας ἐν τῷ ὄρει Γελβουέ.
\vs{9}Καὶ ἐξέδυσαν αὐτὸν, καὶ ἔλαβον τὴν κεφαλὴν αὐτοῦ καὶ τὰ σκεύη αὐτοῦ, καὶ ἀπέστειλαν εἰς γῆν ἀλλοφύλων κύκλῳ τοῦ εὐαγγελίσασθαι τοῖς εἰδώλοις αὐτῶν, καὶ τῷ λαῷ·
\vs{10}Καὶ ἔθηκαν τὰ σκεύη αὐτῶν ἐν οἴκῳ θεοῦ αὐτῶν· καὶ τὴν κεφαλὴν αὐτοῦ ἔθηκαν ἐν οἴκῳ Δαγών.

\vs{11}Καὶ ἤκουσαν πάντες οἱ κατοικοῦντες Γαλαὰδ ἅπαντα ἃ ἐποίησαν οἱ ἀλλόφυλοι τῷ Σαοὺλ καὶ τῷ Ἰσραήλ.
\vs{12}Καὶ ἠγέρθησαν ἐκ Γαλαὰδ πᾶς ἀνὴρ δυνατὸς, καὶ ἔλαβον τὸ σῶμα Σαοὺλ καὶ τὸ σῶμα τῶν υἱῶν αὐτοῦ, καὶ ἤνεγκαν αὐτὰ εἰς Ἰαβὶς, καὶ ἔθαψαν τὰ ὀστᾶ αὐτῶν ὑπὸ τὴν δρὺν ἐν Ἰαβίς· καὶ ἐνήστευσαν ἑπτὰ ἡμέρας.
\vs{13}Καὶ ἀπέθανε Σαοὺλ ἐν ταῖς ἀνομίαις αὐτοῦ, αἷς ἠνόμησε τῷ Θεῷ κατὰ τὸν λόγον Κυρίου, διότι οὐκ ἐφύλαξεν, ὅτι ἐπηρώτησε Σαοὺλ ἐν τῷ ἐγγαστριμύθῳ τοῦ ζητῆσαι, καὶ ἀπεκρίνατο αὐτῷ Σαμουὴλ ὁ προφήτης,
\vs{14}καὶ οὐκ ἐζήτησε Κύριον· καὶ ἀπέκτεινεν αὐτὸν, καὶ ἐπέστρεψε τὴν βασιλείαν τῷ Δαυὶδ υἱῷ Ἰεσσαί.

\ch{11}
Καὶ ἦλθε πᾶς Ἰσραὴλ πρὸς Δαυὶδ ἐν Χεβρὼν, λέγοντες, ἰδοὺ ὀστᾶ σου καὶ σάρκες σου ἡμεῖς,
\vs{2}καὶ ἐχθὲς καὶ τρίτην ὄντος Σαοὺλ βασιλέως, σὺ ἦσθα ὁ ἐξάγων καὶ εἰσάγων τὸν Ἰσραήλ· καὶ εἶπεν Ἰσραὴλ Κύριός σοι, σὺ ποιμανεῖς τὸν λαόν μου τὸν Ἰσραὴλ, καὶ σὺ ἔσῃ εἰς ἡγούμενον ἐπὶ Ἰσραήλ.
\vs{3}Καὶ ἦλθον πάντες πρεσβύτεροι Ἰσραὴλ πρὸς τὸν βασιλέα εἰς Χεβρών· καὶ διέθετο αὐτοῖς ὁ βασιλεὺς Δαυὶδ διαθήκην ἐν Χεβρὼν ἔναντι Κυρίου· καὶ ἔχρισαν τὸν Δαυὶδ εἰς βασιλέα ἐπὶ Ἰσραὴλ κατὰ τὸν λόγον Κυρίου διὰ χειρὸς Σαμουήλ.

\vs{4}Καὶ ἐπορεύθη ὁ βασιλεὺς καὶ ἄνδρες αὐτοῦ εἰς Ἱερουσαλὴμ, αὕτη Ἰεβοὺς, καὶ ἐκεῖ οἱ κατοικοῦντες τὴν γῆν εἶπον τῷ Δανιδ, οὐκ εἰσελεύση ὧδε·
\vs{5}καὶ προκατελάβετο τὴν περιοχὴν Σιών· αὕτη ἡ πόλις Δαυίδ.
\vs{6}Καὶ εἶπε Δαυὶδ, πᾶς τύπτων Ἰεβουσαῖον ἐν πρώτοις, καὶ ἔσται εἰς ἄρχοντα καὶ εἰς στρατηγόν· καὶ ἀνέβη ἐπʼ αὐτὴν ἐν πρώτοις Ἰωὰβ υἱὸς Σαρουία, καὶ ἐγένετο εἰς ἄρχοντα.
\vs{7}Καὶ ἐκάθισε Δαυὶδ ἐν τῇ περιοχῇ· διὰ τοῦτο ἐκάλεσεν αὐτὴν πόλιν Δαυίδ.
\vs{8}Καὶ ᾠκοδόμησε τὴν πόλιν κύκλῳ.
\vs{9}Καὶ ἐπορεύετο Δαυὶδ πορευόμενος καὶ μεγαλυνόμενος, καὶ Κύριος παντοκράτωρ μετʼ αὐτοῦ.
\vs{10}καὶ οὗτοι οἱ ἄρχοντες τῶν δυνατῶν, οἳ ἦσαν τῷ Δαυὶδ, οἱ κατισχύοντες μετʼ αὐτοῦ ἐν τῇ βασιλείᾳ αὐτοῦ μετὰ παντὸς Ἰσραὴλ, τοῦ βασιλεῦσαι αὐτὸν κατὰ τὸν λόγον Κυρίου ἐπὶ Ἰσραήλ.

\vs{11}Καὶ οὗτος ὁ ἀριθμὸς τῶν δυνατῶν τοῦ Δαυίδ· Ἰεσεβαδὰ υἱὸς ʼΑχαμὰν πρῶτος τῶν τριάκοντα· οὗτος ἐσπάσατο τὴν ῥομφαίαν αὐτοῦ ἅπαξ ἐπὶ τριακοσίους τραυματίας ἐν καιρῷ ἑνί.
\vs{12}Καὶ μετʼ αὐτὸν Ἐλεάζαρ υἱὸς Δωδαῒ ὁ Ἀχωχί· οὗτος ἦν ἐν τοῖς τρισὶ δυνατοῖς.
\vs{13}Οὗτος ἦν μετὰ Δαυὶδ ἐν Φασσδαμὶν, καὶ οἱ ἀλλόφυλοι συνήχθησαν ἐκεῖ εἰς πόλεμον, καὶ ἦν μερὶς τοῦ ἀγροῦ πλήρης κριθῶν, καὶ ὁ λαὸς ἔφυγεν ἀπὸ προσώπου ἀλλοφύλων.
\vs{14}Καὶ ἔστη ἐν μέσῳ τῆς μερίδος, καὶ ἔσωσεν αὐτὴν, καὶ ἐπάταξε τοὺς ἀλλοφύλους, καὶ ἐποίησε Κύριος σωτηρίαν μεγάλην.

\vs{15}Καὶ κατέβησαν τρεῖς ἐκ τῶν τριάκοντα ἀρχόντων εἰς τὴν πέτραν πρὸς Δαυὶδ εἰς τὸ σπήλαιον Ὀδολλὰμ, καὶ παρεμβολὴ τῶν ἀλλοφύλων ἐν τῇ κοιλάδι τῶν γιγάντων.
\vs{16}Καὶ Δαυὶδ τότε ἐν τῇ περιοχῇ, καὶ τὸ σύστημα τῶν ἀλλοφύλων τότε ἐν Βηθλεέμ.
\vs{17}Καὶ ἐπεθύμησε Δαυὶδ, καὶ εἶπε, τίς ποτιεῖ με ὕδωρ ἐκ τοῦ λάκκου Βηθλεὲμ τοῦ ἐν τῇ πύλῃ;
\vs{18}Καὶ διέῤῥηξαν οἱ τρεῖς τὴν παρεμβολὴν τῶν ἀλλοφύλων· καὶ ὑδρεύσαντο ὕδωρ ἐκ τοῦ λάκκου τοῦ ἐν Βηθλεὲμ, ὃς ἦν ἐν τῇ πύλῃ, καὶ ἔλαβον καὶ ἤλθον πρὸς Δαυίδ· καὶ οὐκ ἠθέλησε Δαυὶδ τοῦ πιεῖν αὐτὸ, καὶ ἔσπεισεν αὐτὸ τῷ Κυρίῳ,
\vs{19}καὶ εἶπεν, ἵλεώς μοι ὁ θεὸς τοῦ ποιῆσαι τὸ ῥῆμα τοῦτο· εἰ αἷμα ἀνδρῶν τούτων πίομαι ἐν ψυχαῖς αὐτῶν; ὅτι ἐν ψυχαῖς αὐτῶν ἤνεγκαν· καὶ οὐκ ἐβούλετο πιεῖν αὐτό· ταῦτα ἐποίησαν οἱ τρεῖς δυνατοί.

\vs{20}Καὶ Ἀβισὰ ἀδελφὸς Ἰωὰβ, οὗτος ἦν ἄρχων τῶν τριῶν· οὗτος ἐσπάσατο τὴν ῥομφαίαν αὐτοῦ ἐπὶ τριακοσίους τραυματίας ἐν καιρῷ ἑνὶ, καὶ οὗτος ἦν ὀνομαστὸς ἐν τοῖς τρισίν.
\vs{21}Ἀπὸ τῶν τριῶν ὑπὲρ τοὺς δύο ἔνδοξος, καὶ ἦν αὐτοῖς εἰς ἄρχοντα, καὶ ἕως τῶν τριῶν οὐκ ἤρχετο.

\vs{22}Καὶ Βαναία υἱὸς Ἰωδαὲ υἱὸς ἄνδρος δυνατοῦ, πολλὰ ἔργα αὐτοῦ ὑπὲρ Καβασαήλ· οὗτος ἐπάταξε τοὺς δύο ἀριὴλ Μωὰβ, καὶ οὗτος κατέβη καὶ ἐπάταξε τὸν λέοντα ἐν τῷ λάκκῳ ἐν ἡμέρᾳ χιόνος.
\vs{23}Καὶ οὗτος ἐπάταξε τὸν ἄνδρα τὸν Αἰγύπτιον, ἄνδρα ὁρατὸν πεντάπηχυν, καὶ ἐν χειρὶ τοῦ Αἰγυπτίου δόρυ ὡς ἀντίον ὑφαινόντων· καὶ κατέβη ἐπʼ αὐτὸν Βαναία ἐν ῥάβδῳ, καὶ ἀφείλατο ἐκ τῆς χειρὸς τοῦ Αἰγυπτίου τὸ δόρυ, καὶ ἀπέκτεινεν αὐτὸν ἐν τῷ δόρατι αὐτοῦ.
\vs{24}Ταῦτα ἐποίησε Βαναία υἱὸς Ἰωδαὲ, καὶ τούτῳ ὄνομα ἐν τοῖς τρισὶ τοῖς δυνατοῖς. Ὑπὲρ τοὺς τριάκοντα ἦν ἔνδοξος οὗτος, καὶ πρὸς τοὺς τρεῖς οὐκ ἤρχετο· καὶ κατέστησεν αὐτὸν Δαυὶδ ἐπὶ τὴν πατριὰν αὐτοῦ.
\vs{25}Ὑπὲρ τοὺς τριάκοντα ἦν ἔνδοξος οὗτος, καὶ πρὸς τοὺς τρεῖς οὐκ ἤρχετο· καὶ κατέστησεν αὐτὸν Δαυεὶδ ἐπὶ τὴν πατριὰν αὐτοῦ.

\vs{26}Καὶ δυνατοὶ τῶν δυνάμεων, Ἀσαὴλ ἀδελφὸς Ἰωὰβ, Ἐλεανὰν υἱὸς Δωδωὲ ἐκ Βηθλεὲμ,
\vs{27}Σαμαὼθ ὁ ʼΑρωρὶ, Χελλὴς ὁ Φελωνὶ,
\vs{28}Ὠρὰ υἱὸς Ἐκκὶς ὁ Θεκωί, Ἀβιέζερ ὁ Ἀναθωθὶ,
\vs{29}Σοβοχαὶ ὁ Οὐσαθὶ, Ἠλὶ ὁ Ἀχωνὶ,
\vs{30}Μαραῒ ὁ Νετωφαθὶ, Χθαὸδ υἱὸς Νοοζὰ ὁ Νετωφαθὶ,
\vs{31}Αἰρεὶ υἱὸς Ῥεβιὲ ἐκ βουνοῦ Βενιαμὶν, Βαναίας ὁ Φαραθωνὶ,
\vs{32}Οὐρὶ ἐκ Ναχαλὶ Γάας, Ἀβιὴλ ὁ Γαραβαιθὶ,
\vs{33}Ἀζβὼν ὁ Βαρμὶ, Ἐλιαβὰ ὁ Σαλαβωνὶ,
\vs{34}υἱὸς Ἀσὰμ τοῦ Γιζωνίτου, Ἰωνάθαν υἱὸς Σωλὰ ὁ Ἀραρὶ,
\vs{35}Ἀχὶμ υἱὸς Ἀχὰρ ὁ ʼΑραρὶ, Ἐλφὰτ υἱὸς Θυροφὰρ
\vs{36}ὁ Μεχωραθρὶ, Ἀχία ὁ Φελλωνὶ,
\vs{37}ʼΗσερὲ ὁ Χαρμαδαῒ, Νααραὶ υἱὸς Ἀζοβαὶ,
\vs{38}Ἰωὴλ υἱὸς Νάθαν, Μεβαὰλ υἱὸς ʼΑγαρὶ,
\vs{39}Σελὴ ὁ Ἀμμωνὶ, Ναχὼρ ὁ Βηρωθὶ, αἴρων σκεύη υἱῷ Σαρουία,
\vs{40}Ἰρὰ ὁ Ἰεθρὶ, Γαβὴρ ὁ Ἰεθρὶ,
\vs{41}Οὐρία ὁ Χεττὶ, Ζαβὲτ υἱὸς Ἀχαϊὰ,
\vs{42}Ἀδινὰ υἱὸς Σαιζὰ τοῦ Ῥουβὴν ἄρχων, καὶ ἐπʼ αὐτῷ τριάκοντα,
\vs{43}ʼΑνὰν υἱὸς Μοωχὰ, καὶ Ἰωσαφὰτ ὁ Ματθανί,
\vs{44}Ὀζία ὁ Ἀσταρωθὶ, Σαμαθὰ καὶ Ἰεϊὴλ υἱοὶ Χωθὰν τοῦ Ἀραρὶ,
\vs{45}Ἰεδιὴλ υἱὸς Σαμερὶ, καὶ Ἰωζαὲ ὁ ἀδελφὸς αὐτοῦ ὁ Θωσαῒ,
\vs{46}Ἐλιὴλ ὁ Μαωὶ, καὶ Ἰαριβὶ, καὶ Ἰωσία υἱὸς αὐτοῦ, Ἐλλαὰμ, καὶ Ἰεθαμὰ ὁ Μωαβίτης,
\vs{47}Δαλιὴλ, καὶ Ὠβὴθ, καὶ Ἰεσσειὴλ ὁ Μεσωβία.

\ch{12}
Καὶ οὗτοι οἱ ἐλθόντες πρὸς Δαυὶδ εἰς Σικελὰγ, ἔτι συνεχομένου ἀπὸ προσώπου Σαοὺλ υἱοῦ Κίς· καὶ οὗτοι ἐν τοῖς δυνατοῖς βοηθοῦντες ἐν πολέμῳ,
\vs{2}καὶ τόξῳ ἐκ δεξιῶν καὶ ἐκ ἀριστερῶν, καὶ σφενδονῆται ἐν λίθοις καὶ τόξοις· ἐκ τῶν ἀδελφῶν Σαοὺλ ἐκ Βενιαμὶν
\vs{3}ὁ ἄρχων Ἀχιέζερ, καὶ Ἰωὰς υἱὸς Ἀσμὰ τοῦ Γαβαθίτου, καὶ Ἰωὴλ, καὶ Ἰωφαλὴτ υἱοὶ Ἀσμὼθ, καὶ Βερχία, καὶ Ἰηοὺλ ὁ Ἀναθωθὶ,
\vs{4}καὶ Σαμαΐας ὁ Γαβαωνείτης δυνατὸς ἐν τοῖς τριάκοντα, καὶ ἐπὶ τῶν τριάκοντα
\vs{5}Ἱερεμία, καὶ Ἰεζιὴλ, καὶ Ἰωανὰν, καὶ Ἰωαζαβὰθ ὁ Γαδαραθιῒμ,
\vs{6}Ἀζαῒ, καὶ Ἀριμοὺθ, καὶ Βααλιὰ, καὶ Σαμαραΐα, καὶ Σαφατίας ὁ Χαραιφιὴλ,
\vs{7}Ἑλκανὰ, καὶ Ἰησουνὶ, καὶ Ὀζριὴλ, καὶ Ἰωζαρὰ, καὶ Σοβοκὰμ, καὶ οἱ Κορῖται,
\vs{8}καὶ Ἰελία καὶ Ζαβαδία υἱοὶ ʼΙροὰμ, καὶ οἱ τοῦ Γεδώρ.

\vs{9}Καὶ ἀπὸ τοῦ Γαδδὶ ἐχωρίσθησαν πρὸς Δαυὶδ ἀπὸ τῆς ἐρήμου ἰσχυροὶ δυνατοὶ ἄνδρες παρατάξεως πολέμου, αἴροντες θυρεοὺς καὶ δόρατα, καὶ πρόσωπον λέοντος τὰ πρόσωπα αὐτῶν, καὶ κοῦφοι ὡς δορκάδες ἐπὶ τῶν ὀρέων τῷ τάχει·
\vs{10}ʼΑζὰ ὁ ἄρχων, Ἀβδία ὁ δεύτερος, Ἐλιὰβ ὁ τρίτος,
\vs{11}Μασμανὰ ὁ τέταρτος, Ἱερεμίας ὁ πέμπτος,
\vs{12}Ἰεθὶ ὁ ἕκτος, Ἐλιὰβ ὁ ἕβδομος,
\vs{13}Ἰωανὰν ὁ ὄγδοος, Ἐλιαζὲρ ὁ ἔννατος,
\vs{14}Ἱερεμία ὁ δέκατος, Μελχαβαναὶ ὁ ἑνδέκατος.
\vs{15}Οὗτοι ἐκ τῶν υἱῶν Γὰδ ἄρχοντες τῆς στρατιᾶς, εἷς τοῖς ἑκατὸν μικρὸς, καὶ μέγας τοῖς χιλίοις.
\vs{16}Οὗτοι οἱ διαβάντες τὸν Ἰορδάνην ἐν τῷ μηνὶ τῷ πρώτῳ· καὶ οὗτος πεπληρωκὼς ἐπὶ πᾶσαν κρηπίδα αὐτοῦ· καὶ ἐξεδίωξαν πάντας τοὺς κατοικοῦντας αὐλῶνας ἀπὸ ἀνατολῶν ἕως δυσμῶν.

\vs{17}Καὶ ἦλθον ἀπὸ τῶν υἱῶν Βενιαμὶν καὶ Ἰούδα εἰς βοήθειαν τοῦ Δαυίδ.
\vs{18}Καὶ Δαυὶδ ἐξῆλθεν εἰς ἀπάντησιν αὐτῶν, καὶ εἶπεν αὐτοῖς, εἰ εἰς εἰρήνην ἥκατε πρὸς μέ, εἴη μοι καρδία καθʼ ἑαυτὴν ἐφʼ ὑμᾶς· καὶ εἰ τοῦ παραδοῦναί με τοῖς ἐχθροῖς μου οὐκ ἐν ἀληθείᾳ χειρὸς, ἴδοι ὁ Θεὸς τῶν πατέρων ὑμῶν καὶ ἐλέγξαιτο.
\vs{19}Καὶ πνεῦμα ἐνέδυσε τὸν Ἀμασαὶ ἄρχοντα τῶν τριάκοντα, καὶ εἶπε, πορεύου καὶ ὁ λαός σου Δαυὶδ υἱὸς Ἰεσσαὶ, εἰρήνη εἰρήνη σοι, καὶ εἰρήνη τοῖς βοηθοῖς σου, ὅτι ἐβοήθησέ σοι ὁ Θεός σου· καὶ προσεδέξατο αὐτοὺς Δαυὶδ, καὶ κατέστησεν αὐτοὺς ἄρχοντας τῶν δυνάμεων.

\vs{20}Καὶ ἀπὸ Μανασσῆ προσεχώρησαν πρὸς Δαυὶδ ἐν τῷ ἐλθεῖν τοὺς ἀλλοφύλους ἐπὶ Σαοὺλ εἰς πόλεμον· καὶ οὐκ ἐβοήθησεν αὐτοῖς, ὅτι ἐν βουλῇ ἐγένετο παρὰ τῶν στρατηγῶν τῶν ἀλλοφύλων λεγόντων, ἐν ταῖς κεφαλαῖς τῶν ἀνδρῶν ἐκείνων ἐπιστρέψει πρὸς τὸν κυρίον αὐτοῦ Σαούλ.
\vs{21}Ἐν τῷ πορευθῆναι τὸν Δαυίδ εἰς Σικελὰγ προσεχώρησαν αὐτῷ ἀπὸ Μανασσῆ, Ἐδνὰ, καὶ Ἰωζαβὰθ, καὶ Ῥωδιὴλ, καὶ Μιχαὴλ, καὶ Ἰωσαβαὶθ, καὶ Ἐλιμοὺθ, καὶ Σεμαθὶ· ἀρχηγοὶ χιλιάδων εἰσὶ τοῦ Μανασσῆ.
\vs{22}Καὶ αὐτοὶ συνεμάχησαν τῷ Δαυὶδ ἐπὶ τὸν γεδδοὺρ, ὅτι δυνατοὶ ἰσχύος πάντες· καὶ ἦσαν ἡγούμενοι ἐν τῇ στρατιᾷ ἐν τῇ δυνάμει,
\vs{23}ὅτι ἡμέραν ἐξ ἡμέρας ἤρχοντο πρὸς Δαυὶδ εἰς δύναμιν μεγάλην ὡς δύναμις τοῦ Θεοῦ.

\vs{24}Καὶ ταῦτα τὰ ὀνόματα τῶν ἀρχόντων τῆς στρατιᾶς, οἱ ἐλθόντες πρὸς Δαυὶδ εἰς Χεβρὼν τοῦ ἀποστρέψαι τὴν βασιλείαν Σαοὺλ πρὸς αὐτὸν κατὰ τὸν λόγον Κυρίου.
\vs{25}Υἱοὶ Ἰούδα θυρεοφόροι καὶ δορατοφόροι, ἓξ χιλιάδες καὶ ὀκτακόσιοι δυνατοὶ παρατάξεως.
\vs{26}Τῶν υἱῶν Συμεὼν, δυνατοὶ ἰσχύος εἰς παράταξιν, ἑπτὰ χιλιάδες καὶ ἑκατόν.
\vs{27}Τῶν υἱῶν Λευὶ, τετρακισχίλιοι καὶ ἑξακόσιοι.
\vs{28}Καὶ Ἰωαδὰς ὁ ἡγούμενος τῷ Ἀαρὼν, καὶ μετʼ αὐτοῦ τρεῖς χιλιάδες καὶ ἑπτακόσιοι.
\vs{29}Καὶ Σαδὼκ νέος δυνατὸς ἰσχύϊ, καὶ τῆς πατρικῆς οἰκίας αὐτοῦ ἄρχοντες εἰκοσιδύο.
\vs{30}Καὶ τῶν υἱῶν Βενιαμὶν τῶν ἀδελφῶν Σαοὺλ, τρεῖς χιλιάδες· καὶ ἔτι τὸ πλεῖστον αὐτῶν ἀπεσκόπει τὴν φυλακὴν οἴκου Σαούλ.
\vs{31}Καὶ ἀπὸ υἱῶν Ἐφραὶμ, εἴκοσι χιλιάδες καὶ ὀκτακόσιοι, δυνατοὶ ἰσχύϊ ἄνδρες ὀνομαστοὶ κατʼ οἴκους πατριῶν αὐτῶν.
\vs{32}Καὶ ἀπὸ τοῦ ἡμίσους φυλῆς Μανασσῆ, δεκαοκτὼ χιλιάδες, καὶ οἳ ὠνομάσθησαν ἐν ὀνόματι τοῦ βασιλεῦσαι τὸν Δαυίδ.
\vs{33}Καὶ ἀπὸ τῶν υἱῶν Ἰσσάχαρ γινώσκοντες σύνεσιν εἰς τοὺς καιροὺς, γινώσκοντες τί ποιῆσαι Ἰσραὴλ, διακόσιοι, καὶ πάντες ἀδελφοὶ αὐτῶν μετʼ αὐτῶν.

\vs{34}Καὶ ἀπὸ Ζαβουλὼν ἐκπορευόμενοι εἰς παράταξιν πολέμου ἐν πᾶσι σκεύεσι πολεμικοῖς πεντήκοντα χιλιάδες βοηθῆσαι τῷ Δαυὶδ οὐ χεροκένως.
\vs{35}Καὶ ἀπὸ Νεφθαλὶ ἄρχοντες χίλιοι, καὶ μετʼ αὐτῶν ἐν θυρεοῖς καὶ δόρασι, τριακονταεπτὰ χιλιάδες.
\vs{36}Καὶ ἀπὸ τῶν Δανιτῶν παρατασσόμενοι εἰς πόλεμον, εἰκοσιοκτὼ χιλιάδες καὶ ὀκτακόσιοι.
\vs{37}Καὶ ἀπὸ τοῦ Ἀσὴρ ἐκπορευόμενοι βοηθῆσαι εἰς πόλεμον, τεσσαράκοντα χιλιάδες.
\vs{38}Καὶ ἐκ πέραν τοῦ Ἰορδάνου ἀπὸ Ῥουβὴν, καὶ Γαδδὶ, καὶ ἀπὸ τοῦ ἡμίσους φυλῆς Μανασσῆ ἐν πᾶσι σκεύεσι πολεμικοῖς, ἑκατὸν εἴκοσι χιλιάδες.

\vs{39}Πάντες οὗτοι ἄνδρες πολεμισταὶ παρατασσόμενοι παράταξιν ἐν ψυχῇ εἰρηνικῇ· καὶ ἦλθον εἰς Χεβρὼν τοῦ βασιλεῦσαι τὸν Δαυὶδ ἐπὶ πάντα Ἰσραήλ· καὶ ὁ κατάλοιπος Ἰσραὴλ ψυχὴ μία τοῦ βασιλεῦσαι τὸν Δαυίδ.
\vs{40}Καὶ ἦσαν ἐκεῖ ἡμέρας τρεῖς ἐσθίοντες καὶ πίνοντες, ὅτι ἡτοίμασαν οἱ ἀδελφοὶ αὐτῶν.
\vs{41}Καὶ οἱ ὁμοροῦντες αὐτοῖς ἕως Ἰσσάχαρ καὶ Ζαβουλὼν καὶ Νεφθαλὶ, ἔφερον αὐτοῖς ἐπὶ τῶν καμήλων καὶ τῶν ὄνων καὶ τῶν ἡμιόνων καὶ ἐπὶ τῶν μόσχων βρώματα, ἄλευρα, παλάθας, σταφίδας, οἶνον, καὶ ἔλαιον, μόσχους καὶ πρόβατα εἰς πλῆθος, ὅτι εὐφροσύνη ἐν Ἰσραήλ.

\ch{13}
Καὶ ἐβουλεύσατο Δαυὶδ μετὰ τῶν χιλιάρχων καὶ τῶν ἑκατοντάρχων, παντὶ ἡγουμένῳ.
\vs{2}Καὶ εἶπε Δαυὶδ πάσῃ ἐκκλησίᾳ Ἰσραὴλ, εἰ ἐφʼ ὑμῖν ἀγαθὸν, καὶ παρὰ Κυρίου τοῦ Θεοῦ ἡμῶν εὐοδωθῇ, ἀποστείλωμεν πρὸς τοὺς ἀδελφοὺς ἡμῶν τοὺς ὑπολελειμμένους ἐν πάσῃ γῇ Ἰσραὴλ, καὶ μετʼ αὐτῶν οἱ ἱερεῖς οἱ Λευῖται ἐν πόλεσι κατασχέσεως αὐτῶν, καὶ συναχθήσονται πρὸς ἡμᾶς,
\vs{3}καὶ μετενέγκωμεν τὴν κιβωτὸν τοῦ Θεοῦ ἡμῶν πρὸς ἡμᾶς, ὅτι οὐκ ἐζήτησαν αὐτὴν ἀφʼ ἡμερῶν Σαούλ.
\vs{4}Καὶ εἶπε πᾶσα ἡ ἐκκλησία τοῦ ποιῆσαι οὕτως, ὅτι εὐθὺς ὁ λόγος ἐν ὀφθαλμοῖς παντὸς τοῦ λαοῦ.

\vs{5}Καὶ ἐξεκκλησίασε Δαυὶδ τὸν πάντα Ἰσραὴλ ἀπὸ ὁρίων Αἰγύπτου καὶ ἕως εἰσόδου Ἡμὰθ τοῦ εἰσενέγκαι τὴν κιβωτὸν τοῦ Θεοῦ ἐκ πόλεως Ἰαρίμ.
\vs{6}Καὶ ἀνήγαγεν αὐτὴν Δαυίδ· καὶ πᾶς Ἰσραὴλ ἀνέβη εἰς πόλιν Δαυὶδ, ἣ ἦν τοῦ Ἰούδα, τοῦ ἀναγαγεῖν ἐκεῖθεν τὴν κιβωτὸν τοῦ Θεοῦ Κυρίου καθημένου ἐπὶ χερουβὶμ, οὗ ἐπεκλήθη ὄνομα αὐτοῦ.
\vs{7}Καὶ ἐπέθηκαν τὴν κιβωτὸν τοῦ Θεοῦ ἐφʼ ἅμαξαν καινὴν ἐξ οἴκου Ἀμιναδάβ· καὶ Ὀζὰ καὶ οἱ ἀδελφοὶ αὐτοῦ ἦγον τὴν ἅμαξαν.

\vs{8}Καὶ Δαυὶδ καὶ πᾶς Ἰσραὴλ παίζοντες ἐναντίον τοῦ Θεοῦ ἐν πάσῃ δυνάμει, καὶ ἐν ψαλτῳδοῖς, καὶ ἐν κινύραις, καὶ ἐν νάβλαις, ἐν τυμπάνοις, καὶ ἐν κυμβάλοις, καὶ ἐν σάλπιγξι.
\vs{9}Καὶ ἤλθοσαν ἕως τῆς ἅλωνος· καὶ ἐξέτεινεν Ὀζὰ τὴν χεῖρα αὐτοῦ τοῦ κατασχεῖν τὴν κιβωτὸν, ὅτι ἐξέκλινεν αὐτὴν ὁ μόσχος.
\vs{10}Καὶ ἐθυμώθη Κύριος ὀργῇ ἐπὶ Ὀζά, καὶ ἐπάταξεν αὐτὸν ἐκεῖ διὰ τὸ ἐκτεῖναι τὴν χεῖρα αὑτοῦ ἐπὶ τὴν κιβωτὸν, καὶ ἀπέθανεν ἐκεῖ ἀπέναντι τοῦ Θεοῦ·
\vs{11}Καὶ ἠθύμησε Δαυὶδ, ὅτι διέκοψε Κύριος διακοπὴν ἐν Ὀζὰ, καὶ ἐκάλεσε τὸν τόπον ἐκεῖνον, Διακοπὴ Ὀζὰ, ἕως τῆς ἡμέρας ταύτης.
\vs{12}Καὶ ἐφοβήθη Δαυὶδ τὸν Θεὸν ἐν τῇ ἡμέρᾳ ἐκείνῃ, λέγων, πῶς εἰσοίσω τὴν κιβωτὸν τοῦ Θεοῦ πρὸς ἐμαυτόν;
\vs{13}Καὶ οὐκ ἀπέστρεψε Δαυὶδ τὴν κιβωτὸν πρὸς ἑαυτὸν εἰς πόλιν Δαυὶδ, καὶ ἐξέκλινεν αὐτὴν εἰς οἶκον Ἀβεδδαρὰ τοῦ Γεθαίου.

\vs{14}Καὶ ἐκάθισεν ἡ κιβωτὸς τοῦ Θεοῦ ἐν οἴκῳ Ἀβεδδαρὰ τρεῖς μῆνας· καὶ εὐλόγησεν ὁ Θεὸς Ἀβεδδαρὰ, καὶ πάντα τὰ αὐτοῦ.

\ch{14}
Καὶ ἀπέστειλε Χειρὰμ βασιλεὺς Τύρου ἀγγέλους πρὸς Δαυὶδ, καὶ ξύλα κέδρινα, καὶ οἰκοδόμους, καὶ τέκτονας ξύλων, τοῦ οἰκοδομῆσαι αὐτῷ οἶκον.
\vs{2}Καὶ ἔγνω Δαυὶδ ὅτι ἡτοίμασεν αὐτὸν Κύριος εἰς βασιλέα ἐπὶ Ἰσραὴλ, ὅτι ηὐξήθη εἰς ὕψος ἡ βασιλεία αὐτοῦ διὰ τὸν λαὸν αὐτοῦ Ἰσραήλ.

\vs{3}Καὶ ἔλαβε Δαυὶδ ἔτι γυναῖκας ἐν Ἱερουσαλήμ· καὶ ἐτέχθησαν Δαυὶδ ἔτι υἱοὶ καὶ θυγατέρες.
\vs{4}Καὶ ταῦτα τὰ ὀνόματα αὐτῶν τῶν τεχθέντων, οἳ ἦσαν αὐτῷ ἐν Ἱερουσαλήμ· Σαμαὰ, Σωβὰβ, Νάθαν, καὶ Σαλωμὼν,
\vs{5}καὶ Βαὰρ, καὶ Ἐλισὰ, καὶ Ἐλιφαλὴθ,
\vs{6}καὶ Ναγὲθ, καὶ Ναφὰθ, καὶ Ἰαφιὲ,
\vs{7}καὶ Ἐλισαμαὲ, καὶ Ἐλιαδὲ, καὶ Ἐλιφαλά.

\vs{8}Καὶ ἤκουσαν ἀλλόφυλοι ὅτι ἐχρίσθη Δαυὶδ βασιλεὺς ἐπὶ πάντα Ἰσραήλ· καὶ ἀνέβησαν πάντες οἱ ἀλλόφυλοι ζητῆσαι τὸν Δαυίδ· καὶ ἤκουσε Δαυὶδ, καὶ ἐξῆλθεν εἰς ἀπάντησιν αὐτοῖς·
\vs{9}Καὶ ἀλλόφυλοι ἦλθον, καὶ συνέπεσον ἐν τῇ κοιλάδι τῶν γιγάντων.
\vs{10}Καὶ ἐπηρώτησε Δαυὶδ διὰ τοῦ Θεοῦ, λέγων, εἰ ἀναβῶ ἐπὶ τοὺς ἀλλοφύλους, καὶ δώσεις αὐτοὺς εἰς τὰς χεῖράς μου; καὶ εἶπεν αὐτῷ Κύριος, ἀνάβηθι, καὶ δώσω αὐτοὺς εἰς τὰς χεῖράς σου.
\vs{11}Καὶ ἀνέβη εἰς Βαὰλ Φαρασὶν, καὶ ἐπάταξεν αὐτοὺς ἐκεῖ Δανίδ· καὶ εἶπε Δαυὶδ, διέκοψεν ὁ Θεὸς τοὺς ἐχθρούς μου ἐν χειρί μου, ὡς διακοπὴν ὕδατος· διὰ τοῦτο ἐκάλεσε τὸ ὄνομα τοῦ τόπου ἐκείνου, Διακοπὴ Φαρσίν.
\vs{12}Καὶ ἐγκατέλιπον ἐκεῖ τοὺς θεοὺς αὐτῶν οἱ ἀλλόφυλοι· καὶ εἶπε Δαυὶδ κατακαῦσαι ἐν πυρί.

\vs{13}Καὶ προσέθεντο ἔτι ἀλλόφυλοι, καὶ συνέπεσαν ἔτι ἐν τῇ κοιλάδι τῶν γιγάντων.
\vs{14}Καὶ ἠρώτησε Δαυὶδ ἔτι ἐν Θεῷ· καὶ εἶπεν αὐτῷ ὁ Θεὸς, οὐ πορεύσῃ ὀπίσω αὐτῶν· ἀποστρέφου ἀπʼ αὐτῶν, καὶ παρέσῃ αὐτοῖς πλησίον τῶν ἀπίων.
\vs{15}Καὶ ἔσται ἐν τῷ ἀκοῦσαί σε τὴν φωνὴν τοῦ συσσεισμοῦ αὐτῶν ἄκρων τῶν ἀπίων, τότε εἰσελεύσῃ εἰς τὸν πόλεμον, ὅτι ἐξῆλθεν ὁ Θεὸς ἔμπροσθέν σου τοῦ πατάξαι τὴν παρεμβολὴν τῶν ἀλλοφύλων.
\vs{16}Καὶ ἐποίησε καθὼς ἐνετείλατο αὐτῷ ὁ Θεός· καὶ ἐπάταξε τὴν παρεμβολὴν τῶν ἀλλοφύλων ἀπὸ Γαβαὼν ἕως Γαζηρά.
\vs{17}Καὶ ἐγένετο ὄνομα Δαυὶδ ἐν πάσῃ τῇ γῇ, καὶ Κύριος ἔδωκε τὸν φόβον αὐτοῦ ἐπὶ πάντα τὰ ἔθνη.

\ch{15}
Καὶ ἐποίησεν αὐτῷ οἰκίας ἐν πόλει Δαυὶδ, καὶ ἡτοίμασε τὸν τόπον τῇ κιβωτῷ τοῦ Θεοῦ, καὶ ἐποίησεν αὐτῇ σκηνήν.
\vs{2}Τότε εἶπε Δαυὶδ, οὐκ ἔστιν ἆραι τὴν κιβωτὸν τοῦ Θεοῦ, ἀλλʼ ἢ τοὺς Λευίτας, ὅτι αὐτοὺς ἐξελέξατο Κύριος αἴρειν τὴν κιβωτὸν Κυρίου, καὶ λειτουργεῖν αὐτῷ ἕως αἰῶνος.

\vs{3}Καὶ ἐξεκκλησίασε Δαυὶδ τὸν πάντα Ἰσραὴλ ἐν Ἱερουσαλὴμ, τοῦ ἀνενέγκαι τὴν κιβωτὸν Κυρίου εἰς τὸν τόπον ὃν ἡτοίμασεν αὐτῇ.
\vs{4}Καὶ συνήγαγε Δαυὶδ τοὺς υἱοὺς Ἀαρὼν τοὺς Λευίτας.
\vs{5}Τῶν υἱῶν Καὰθ, Οὐριὴλ ὁ ἄρχων καὶ οἱ ἀδελφοὶ αὐτοῦ, ἑκατὸν εἴκοσι.
\vs{6}Τῶν υἱῶν Μεραρὶ, Ἀσαΐα ὁ ἄρχων καὶ οἱ ἀδελφοὶ αὐτοῦ, διακόσιοι εἴκοσι.
\vs{7}Τῶν υἱῶν Γεδσὼν, Ἰωὴλ ὁ ἄρχων καὶ οἱ ἀδελφοὶ αὐτοῦ, ἑκατὸν τριάκοντα·
\vs{8}Τῶν υἱῶν Ἐλισαφὰτ, Σεμεῒ ὁ ἄρχων καὶ οἱ ἀδελφοὶ αὐτοῦ, διακόσιοι.
\vs{9}Τῶν υἱῶν Χεβρὼμ, Ἐλιὴλ ὁ ἄρχων καὶ οἱ ἀδελφοὶ αὐτοῦ, ὀγδοήκοντα.
\vs{10}Τῶν υἱῶν Ὀζιὴλ, Ἀμιναδὰβ ὁ ἄρχων καὶ οἱ ἀδελφοὶ αὐτοῦ, ἑκατὸν δεκαδύο.

\vs{11}Καὶ ἐκάλεσε Δαυὶδ τὸν Σαδὼκ καὶ Ἀβιάθαρ τοὺς ἱερεῖς, καὶ τοὺς Λευίτας, τὸν Οὐριὴλ, Ἀσαΐαν, καὶ Ἰωὴλ, καὶ Σεμαίαν, καὶ Ἐλιὴλ, καὶ Ἀμιναδὰβ,
\vs{12}καὶ εἶπεν αὐτοῖς, ὑμεῖς ἄρχοντες πατριῶν τῶν Λευιτῶν, ἁγνίσθητε ὑμεῖς καὶ οἱ ἀδελφοὶ ὑμῶν, καὶ ἀνοίσετε τὴν κιβωτὸν τοῦ Θεοῦ Ἰσραὴλ, οὗ ἡτοίμασα αὐτῇ·
\vs{13}Ὅτι οὐκ ἐν τῷ πρότερον ὑμᾶς εἶναι, διέκοψεν ὁ Θεὸς ἡμῶν ἐν ἡμῖν, ὅτι οὐκ ἐξεζητήσαμεν ἐν κρίματι.
\vs{14}Καὶ ἡγνίσθησαν οἱ ἱερεῖς καὶ οἱ Λευῖται, τοῦ ἀνενέγκαι τὴν κιβωτὸν Θεοῦ Ἰσραήλ.
\vs{15}Καὶ ἔλαβον οἱ υἱοὶ τῶν Λευιτῶν τὴν κιβωτὸν τοῦ Θεοῦ, ὡς ἐνετείλατο Μωυσῆς ἐν λόγῳ Θεοῦ κατὰ τὴν γραφὴν, ἐν ἀναφορεῦσιν ἐπʼ αὐτούς.

\vs{16}Καὶ εἶπε Δανὶδ τοῖς ἄρχουσι τῶν Λευιτῶν, στήσατε τοὺς ἀδελφοὺς αὐτῶν τοὺς ψαλτῳδοὺς ἐν ὀργάνοις, νάβλαις, κινύραις, καὶ κυμβάλοις τοῦ φωνῆσαι εἰς ὕψος ἐν φωνῇ εὐφροσύνης.
\vs{17}Καὶ ἔστησαν οἱ Λευῖται τὸν Αἰμὰν υἱὸν Ἰωήλ· ἐκ τῶν ἀδελφῶν αὐτοῦ Ἀσὰφ υἱὸς Βαραχία· καὶ ἐκ τῶν νἱῶν Μεραρὶ ἀδελφῶν αὐτοῦ Αἰθὰν υἱὸς Κισαίου·
\vs{18}Καὶ μετʼ αὐτῶν οἱ ἀδελφοὶ αὐτῶν οἱ δεύτεροι, Ζαχαρίας, καὶ Ὀζιὴλ, καὶ Σεμιραμὼθ, καὶ Ἰεϊὴλ, καὶ Ἐλιωὴλ, καὶ Ἐλιὰβ, καὶ Βαναία, καὶ Μαασαΐα, καὶ Ματθαθία, καὶ Ἐλιφενὰ, καὶ Μακελλία, καὶ Ἀβδεδὸμ, καὶ Ἰεϊὴλ, καὶ Ὀζίας, οἱ πυλωροί.·
\vs{19}Καὶ οἱ ψαλτῳδοὶ, Αἰμὰν, Ἀσὰφ, καὶ Αἰθὰν ἐν κυμβάλοις χαλκοῖς τοῦ ἀκουσθῆναι ποιῆσαι.
\vs{20}Ζαχαρίας, καὶ Ὀζιὴλ, Σεμιραμὼθ, Ἰεϊὴλ, Ὠνι, Ἐλιὰβ, Μαασαίας, Βαναίας ἐν νάβλαις ἐπὶ ἀλαιμώθ.
\vs{21}Καὶ Ματταθίας, καὶ Ἐλιφαλοὺ, καὶ Μακενία, καὶ Ἀβδεδὸμ, καὶ Ἰεϊὴλ, καὶ Ὀζίας ἐν κινύραις ἁμασενὶθ τοῦ ἐνισχῦσαι.

\vs{22}Καὶ Χωνενία ἄρχων τῶν Λευιτῶν ἄρχων τῶν ᾠδῶν, ὅτι συνετὸς ἦν.
\vs{23}Καὶ Βαραχία καὶ Ἐλκανὰ πυλωροὶ τῆς κιβωτοῦ.
\vs{24}Καὶ Σομνία, καὶ Ἰωσαφὰτ, καὶ Ναθαναὴλ, καὶ Ἀμασαῒ, καὶ Ζαχαρία, καὶ Βαναΐα, καὶ Ἐλιέζερ οἱ ἱερεῖς σαλπίζοντες ταῖς σάλπιγξιν ἔμπροσθεν τῆς κιβωτοῦ τοῦ Θεοῦ· καὶ Ἀβδεδὸμ καὶ Ἰεΐα πυλωροὶ τῆς κιβωτοῦ τοῦ Θεοῦ.

\vs{25}Καὶ ἦν Δαυὶδ καὶ οἱ πρεσβύτεροι Ἰσραὴλ καὶ οἱ χιλίαρχοι οἱ πορευόμενοι τοῦ ἀναγαγεῖν τὴν κιβωτὸν τῆς διαθήκης ἐξ οἴκου Ἀβδεδὸμ ἐν εὐφροσύνῃ.
\vs{26}Καὶ ἐγένετο ἐν τῷ κατισχῦσαι τὸν Θεὸν τοὺς Λευίτας αἴροντας τὴν κιβωτὸν τῆς διαθήκης Κυρίου, καὶ ἔθυσαν ἀνʼ ἑπτὰ μόσχους, καὶ ἀνʼ ἑπτὰ κριούς.
\vs{27}Καὶ Δαυὶδ περιεζωσμένος ἐν στολῇ βυσσίνῃ, καὶ πάντες οἱ Λευῖται αἴροντες τὴν κιβωτὸν διαθήκης Κυρίου, καὶ οἱ ψαλτῳδοὶ, καὶ Χωνενίας ὁ ἄρχων τῶν ᾠδῶν τῶν ᾀδόντων, καὶ ἐπὶ Δαυὶδ στολὴ βυσσίνη.
\vs{28}Καὶ πᾶς Ἰσραὴλ ἀνάγοντες τὴν κιβωτὸν διαθήκης Κυρίου ἐν σημασίᾳ, καὶ ἐν φωνῇ σωφὲρ, καὶ ἐν σάλπιγξι, καὶ ἐν κυμβάλοις, ἀναφωνοῦντες ἐν νάβλαις καὶ ἐν κινύραις.
\vs{29}Καὶ ἐγένετο ἡ κιβωτὸς διαθήκης Κυρίου, καὶ ἦλθεν ἕως πόλεως Δαυίδ· καὶ Μελχὸλ ἡ θυγάτηρ Σαοὺλ παρέκυψε διὰ τῆς θυρίδος, καὶ εἶδε τὸν βασιλέα Δαυὶδ ὀρχούμενον καὶ παίζοντα, καὶ ἐξουδένωσεν αὐτὸν ἐν τῇ ψυχῇ αὐτῆς.

\ch{16}
Καὶ εἰσήνεγκαν τὴν κιβωτὸν τοῦ Θεοῦ, καὶ ἀπηρείσαντο αὐτὴν ἐν μέσῳ τῆς σκηνῆς ἧς ἔπηξεν αὐτῇ Δαυὶδ, καὶ προσήνεγκαν ὁλοκαυτώματα καὶ σωτηρίου ἐναντίον τοῦ Θεοῦ.
\vs{2}Καὶ συνετέλεσε Δαυὶδ ἀναφέρων ὁλοκαυτώματα καὶ σωτηρίου, καὶ εὐλόγησε τὸν λαὸν ἐν ὀνόματι Κυρίου.
\vs{3}Καὶ διεμέρισε παντὶ ἀνδρὶ Ἰσραὴλ ἀπὸ ἀνδρὸς καὶ ἕως γυναικὸς, τῷ ἀνδρὶ ἄρτον ἕνα ἀρτοκοπικὸν, καὶ ἀμορείτην.
\vs{4}Καὶ ἔταξε κατὰ πρόσωπον τῆς κιβωτοῦ διαθήκης Κυρίου ἐκ τῶν Λευιτῶν λειτουργοῦντας ἀναφωνοῦντας, καὶ ἐξομολογεῖσθαι καὶ αἰνεῖν Κύριον τὸν Θεὸν Ἰσραήλ·
\vs{5}Ἀσὰφ ὁ ἡγούμενος, καὶ δευτερεύων αὐτῷ Ζαχαρίας, Ἰεϊὴλ, Σεμιραμὼθ, καὶ Ἰεϊὴλ, Ματταθίας, Ἐλιὰβ, καὶ Βαναΐας, καὶ Ἀβδεδομ· καὶ Ἰεϊὴλ ἐν ὀργάνοις, νάβλαις, κινύραις, καὶ Ἀσὰφ ἐν κυμβάλοις ἀναφωνῶν·
\vs{6}Καὶ Βαναίας καὶ Ὀζιὴλ οἱ ἱερεῖς ἐν ταῖς σάλπιγξι διαπαντὸς ἐναντίον τῆς κιβωτοῦ τῆς διαθήκης τοῦ Θεοῦ.

\vs{7}Ἐν τῇ ἡμέρᾳ ἐκείνῃ. Τότε ἔταξε Δαυὶδ ἐν ἀρχῇ τοῦ αἰνεῖν τὸν Κύριον ἐν χειρὶ Ἀσὰφ καὶ τῶν ἀδελφῶν αὐτοῦ.

\vs{8}Ὠδη. Ἐξομολογεῖσθε τῷ Κυρίῳ, ἐπικαλεῖσθε αὐτὸν ἐν ὀνόματι αὐτοῦ, γνωρίσατε ἐν λαοῖς τὰ ἐπιτηδεύματα αὐτοῦ.
\vs{9}Ἄσατε αὐτῷ καὶ ὑμνήσατε αὐτῷ, διηγήσασθε πᾶσι τὰ θαυμάσια αὐτοῦ, ἃ ἐποίησε Κύριος.
\vs{10}Αἰνεῖτε ἐν ὀνόματι ἁγίῳ αὐτοῦ, εὐφρανθήσεται καρδία ζητοῦσα τὴν εὐδοκίαν αὐτοῦ.
\vs{11}Ζητήσατε τὸν Κύριον καὶ ἰσχύσατε, ζητήσατε τὸ πρόσωπον αὐτοῦ διαπαντός.
\vs{12}Μνημονεύετε τὰ θαυμάσια αὐτοῦ ἃ ἐποίησε, τέρατα καὶ κρίματα τοῦ στόματος αὐτοῦ.
\vs{13}σπέρμα Ἰσραὴλ παῖδες αὐτοῦ, υἱοὶ Ἰακὼβ ἐκλεκτοὶ αὐτοῦ.
\vs{14}Αὐτὸς Κύριος ὁ Θεὸς ἡμῶν, ἐν πάσῃ τῇ γῇ τὰ κρίματα αὐτοῦ.
\vs{15}Μνημονεύωμεν εἰς αἰῶνα διαθήκης αὐτοῦ, λόγον αὐτοῦ ὃν ἐνετείλατο εἰς χιλίας γενεὰς,
\vs{16}ὃν διέθετο τῷ Ἁβραὰμ, καὶ τὸν ὅρκον αὐτοῦ τῷ Ἰσαάκ.
\vs{17}Ἔστησεν αὐτὸν τῷ Ἰακὼβ εἰς πρόσταγμα, τῷ Ἰσραὴλ διαθήκην αἰώνιον,
\vs{18}λέγων, σοὶ δώσω τὴν γῆν Χαναὰν σχοίνισμα κληρονομίας ὑμῶν.
\vs{19}Ἐν τῷ γενέσθαι αὐτοὺς ὀλιγοστοὺς ἀριθμῷ, ὡς ἐσμικρύνθησαν, καὶ παρῴκησαν ἐν αὐτῇ,
\vs{20}καὶ ἐπορεύθησαν ἀπὸ ἔθνους εἰς ἔθνος, καὶ ἀπὸ βασιλείας εἰς λαὸν ἕτερον,
\vs{21}οὐκ ἀφῆκεν ἄνδρα τοῦ δυναστεῦσαι αὐτοὺς, καὶ ἤλεγξε περὶ αὐτῶν βασιλεῖς.
\vs{22}Μὴ ἅψησθε τῶν χριστῶν μου, καὶ ἐν τοῖς πρόφήταις μου μὴ πονηρεύεσθε.

\vs{23}Ἄσατε τῷ Κυρίῳ πᾶσα ἡ γῆ, ἀναγγείλατε ἐξ ἡμέρας εἰς ἡμέραν σωτηρίαν αὐτοῦ.
\vs{23a}Ἐξηγεῖσθε ἐν τοῖς ἔθνεσι τὴν δόξαν αὐτοῦ, ἐν πᾶσι τοῖς λαοῖς τὰ θαυμάσια αὐτοῦ.
\vs{25}Ὅτι μέγας Κύριος καὶ αἰνετὸς σφόδρα, φοβερός ἐστιν ἐπὶ πάντας τοὺς θεούς.
\vs{26}Ὅτι πάντες οἱ θεοὶ τῶν ἐθνῶν εἴδωλα, καὶ ὁ Θεὸς ἡμῶν οὐρανοὺς ἐποίησε.
\vs{27}Δόξα καὶ ἔπαινος κατὰ πρόσωπον αὐτοῦ, ἰσχὺς καὶ καύχημα ἐν τόπῳ αὐτοῦ.
\vs{28}Δότε τῷ Κυρίῳ αἱ πατριαὶ τῶν ἐθνῶν, δότε τῷ Κυρίῳ δόξαν καὶ ἰσχὺν,
\vs{29}δότε τῷ Κυρίῳ δόξαν ὀνόματι αὐτοῦ· λάβετε δῶρα καὶ ἐνέγκατε κατὰ πρόσωπον αὐτοῦ, καὶ προσκυνήσατε Κυρίῳ ἐν αὐλαῖς ἁγίαις αὐτοῦ.
\vs{30}Φοβηθήτω ἀπὸ προσώπου αὐτοῦ πᾶσα ἡ γῆ, κατορθωθήτω ἡ γῆ, καὶ μὴ σαλευθήτω.
\vs{31}Εὐφρανθήτω ὁ οὐρανὸς, καὶ ἀγαλλιάσθω ἡ γῆ, καὶ εἰπάτωσαν ἐν τοῖς ἔθνεσι, Κύριος βασιλεύων.
\vs{32}Βομβήσει ἡ θάλασσα σὺν τῷ πληρώματι, καὶ ξὺλον ἀγροῦ καὶ πάντα τὰ ἐν αὐτῷ
\vs{33}Τότε εὐφρανθήσεται τὰ ξύλατοῦ δρυμοῦ ἀπὸ προσώπου Κυρίου, ὅτι ἦλθε κρίναι τὴν γῆν.
\vs{34}Ἐξομολογεῖσθε τῷ Κυρίῳ, ὅτι ἀγαθὸν, ὅτι εἰς τὸν αἰῶνα τὸ ἔλεος αὐτοῦ.
\vs{35}Καὶ εἴπατε, σῶσον ἡμᾶς, ὁ Θεὸς τῆς σωτηρίας ἡμῶν καὶ ἄθροισον ἡμᾶς. καὶ ἐξελοῦ ἡμᾶς ἐκ τῶν ἐθνῶν, τοῦ αἴνειν τὸ ὄνομα τὸ ἅγιόν σου, καὶ καυχᾶσθαι ἐν ταῖς αἰνέσεσί σου.
\vs{36}Εὐλογημένος Κύριος ὁ Θεὸς Ἰσραὴλ ἀπὸ τοῦ αἰῶνος καὶ ἕως τοῦ αἰῶνος·

Καὶ ἐρεῖ πᾶς ὁ λαὸς, ἀμήν· καὶ ᾔνεσαν τῷ Κυρίῳ.

\vs{37}Καὶ κατέλιπον ἐκεῖ ἔναντι τῆς κιβωτοῦ διαθήκης Κυρίου τὸν Ἀσὰφ καὶ τοὺς ἀδελφοὺς αὐτοῦ, τοῦ λειτουργεῖν ἐναντίον τῆς κιβωτοῦ διαπαντὸς τὸ τῆς ἡμέρας εἰς ἡμέραν.
\vs{38}Καὶ Ἀβδεδὸμ καὶ οἱ ἀδελφοὶ αὐτοῦ, ἑξήκοντα καὶ ὀκτώ· καὶ Ἀβδεδὸμ υἱὸς Ἰδιθοὺν, καὶ Ὀσὰ, εἰς τοὺς πυλωρούς.
\vs{39}Καὶ τὸν Σαδὼκ τὸν ἱερέα καὶ τοὺς ἀδελφοὺς αὐτοῦ τοὺς ἱερεῖς ἐναντίον τῆς σκηνῆς Κυρίου ἐν βαμὰ τῇ ἐν Γαβαὼν,
\vs{40}τοῦ ἀναφέρειν ὁλοκαυτώματα τῷ Κυρίῳ ἐπὶ τοῦ θυσιαστηρίου τῶν ὁλοκαυτωμάτων διαπαντὸς τοπρωῒ καὶ τοεσπέρας, καὶ κατὰ πάντα τὰ γεγραμμένα ἐν νόμῳ Κυρίου ὅσα ἐνετείλατο ἐφʼ υἱοῖς Ἰσραὴλ ἐν χειρὶ Μωυσῆ τοῦ θεράποντος τοῦ Θεοῦ.
\vs{41}Καὶ μετʼ αὐτοῦ Αἰμὰν καὶ Ἰδιθοὺν, καὶ οἱ λοιποὶ ἐκλεγέντες ἐπʼ ὀνόματος τοῦ αἰνεῖν τὸν Κύριον, ὅτι εἰς τὸν αἰῶνα τὸ ἔλεος αὐτοῦ.
\vs{42}Καὶ μετʼ αὐτῶν σάλπιγγες καὶ κὺμβαλα τοῦ ἀναφωνεῖν καὶ ὄργανα τῶν ῷδῶν τοῖ Θεοῦ, οἱ δὲ υἱοὶ Ἰδιθοὺν εἰς τὴν πύλην.

\vs{43}Καὶ ἐπορεύθη πᾶς ὁ λαὸς ἕκαστος εἰς τὸν οἶκον αὐτοῦ, καὶ ἐπέστρεψε Δαυὶδ τοῦ εὐλογῆσαι τὸν οἶκον αὐτοῦ.

\ch{17}
Καὶ ἐγένετο ὡς κατῴκησε Δαυὶδ ἐν οἴκῳ αὐτοῦ, καὶ εἶπε Δαυὶδ πρὸς Νάθαν τὸν προφήτην, ἰδοὺ ἐγὼ κατοικῶ ἐν οἴκῳ κεδρίνῳ, καὶ ἡ κιβωτὸς διαθήκης Κυρίου ὑποκάτω δέῤῥεων.

\vs{2}Καὶ εἶπε Νάθαν πρὸς Δαυὶδ, πᾶν τὸ ἐν τῇ ψυχῇ σου ποίει, ὅτι Θεὸς μετὰ σοῦ.

\vs{3}Καὶ ἐγένετο ἐν τῇ νυκτὶ ἐκείνῃ, καὶ ἐγένετο λόγος Κυρίου πρὸς Νάθαν·
\vs{4}Πορεύου καὶ εἶπον πρὸς Δαυὶδ τὸν δοῦλόν μου, οὕτως εἶπε Κύριος, οὐ σὺ οἰκοδομήσεις μοι οἶκον τοῦ κατοικῆσαί με ἐν αὐτῷ.
\vs{5}Ὅτι οὐ κατῴκησα ἐν οἴκῳ ἀπὸ τῆς ἡμέρας ἧς ἀνήγαγον τὸν Ἰσραὴλ ἕως τῆς ἡμέρας ταύτης,
\vs{6}καὶ ἤμην ἐν σκηνῇ καὶ ἐν καλὺμματι ἐν πᾶσιν οἷς διῆλθον ἐν παντὶ Ἰσραήλ· εἰ λαλῶν ἐλάλησα πρὸς μίαν φυλὴν τοῦ Ἰσραὴλ οἷς ἐνετειλάμην τοῦ ποιμαίνειν τὸν λαόν μου, λέγων, ὅτι οὐκ ᾠκοδομήσατέ μοι οἶκον κέδρινον;
\vs{7}Καὶ νῦν οὕτως ἐρεῖς τῷ δούλῳ μου Δαυὶδ, τάδε λέγει Κύριος παντοκράτωρ, ἐγὼ ἔλαβον σε ἐκ τῆς μάνδρας ἐξόπισθεν τῶν ποιμνίων τοῦ εἶναι εἰς ἡγούμενον ἐπὶ τὸν λαόν μου Ἰσραήλ·
\vs{8}Καὶ ἤμην μετὰ σοῦ ἐν πᾶσιν οἷς ἐπορεύθης, καὶ ἐξωλόθρευσα πάντας τοὺς ἐχθρούς σου ἀπὸ προσώπου σου, καὶ ἐποίησά σοι ὄνομα κατὰ τὸ ὄνομα τῶν μεγάλων τῶν ἐπὶ τῆς γῆς.
\vs{9}Καὶ θήσομαι τόπον τῷ λαῷ μου Ἰσραὴλ, καὶ καταφυτεύσω αὐτὸν, καὶ κατασκηνώσει καθʼ ἑαυτὸν, καὶ οὐ μεριμνήσει ἔτι, καὶ οὐ προσθήσει υἱὸς ἀδικίας τοῦ ταπεινῶσαι αὐτὸν καθὼς ἀρχῆς, καὶ ἀφʼ ἡμερῶν ὧν ἔταξα κριτὰς ἐπὶ τὸν λαόν μου Ἰσραήλ·
\vs{10}καὶ ἐταπείνωσα πάντας τοὺς ἐχθρούς σου, καὶ αὐξήσω σε, καὶ οἶκον οἰκοδομήσει σοι Κύριος.
\vs{11}Καὶ ἔσται ὅταν πληρωθῶσιν ἡμέραι σου καὶ κοιμηθήσῃ μετὰ τῶν πατέρων σου, καὶ ἀναστήσω τὸ σπέρμα σου μετὰ σὲ ὅς ἔσται ἐκ τῆς κοιλίας σου, καὶ ἑτοιμάσω τὴν βασιλείαν αὐτοῦ.
\vs{12}Αὐτὸς οἰκοδομήσει μοι οἶκον, καὶ ἀνορθώσω τὸν θρόνον αὐτοῦ ἕως αἰῶνος.
\vs{13}Ἐγὼ ἔσομαι αὐτῷ εἰς πατέρα, καὶ αὐτὸς ἔσται μοι εἰς υἱόν· καὶ τὸ ἔλεός μου οὐκ ἀποστήσω ἀπʼ αὐτοῦ, ὡς ἀπέστησα ἀπὸ τῶν ὄντων ἔμπροσθέν σου.
\vs{14}Καὶ πιστώσω αὐτὸν ἐν οἴκῳ μου καὶ ἐν βασιλείᾳ αὐτοῦ ἕως αἰῶνος, καὶ ὁ θρόνος αὐτοῦ ἔσται ἀνωρθωμένος ἕως αἰῶνος.

\vs{15}Κατὰ πάντας τοὺς λόγους τούτους, καὶ κατὰ πᾶσαν τὴν ὅρασιν ταύτην, οὕτως ἐλάλησε Νάθαν πρὸς Δαυίδ.

\vs{16}Καὶ ἦλθεν ὁ βασιλεὺς Δαυὶδ καὶ ἐκὰθισεν ἀπέναντι Κυρίου, καὶ εἶπε, τίς εἰμι ἐγὼ Κύριε ὁ Θεός; καὶ τίς ὁ οἶκός μου, ὅτι ἠγάπησάς με ἕως αἰῶνος;
\vs{17}Καὶ ἐσμικρύνθη ταῦτα ἐνώπιόν σου ὁ Θεὸς, καὶ ἐλάλησας ἐπὶ τὸν οἶκον τοῦ παιδός σου ἐκ μακρῶν, καὶ ἐπεῖδές με ὡς ὅρασις ἀνθρώπου, καὶ ὕψωσάς με Κύριε ὁ Θεός.
\vs{18}Τί προσθήσει ἔτι Δαυὶδ πρὸς σέ τοῦ δοξάσαι; καὶ σὺ τὸν δοῦλόν σου οἶδας,
\vs{19}καὶ κατὰ τὴν καρδίαν σου ἐποίησας τὴν πᾶσαν μεγαλωσύνην.
\vs{20}Κύριε, οὐκ ἔστιν ὅμοιός σοι, καὶ οὐκ ἔστι Θεὸς πλὴν σοῦ, κατὰ πάντα ὅσα ἠκούσαμεν ἐν ὠσὶν ἡμῶν.
\vs{21}Καὶ οὐκ ἔστιν ὡς ὁ λαός σου Ἰσραὴλ ἔθνος ἔτι ἐτὶ τῆς γῆς, ὡς ὡδήγησεν αὐτὸν ὁ Θεὸς τοῦ λυτρώσασθαι λαὸν ἑαυτῷ, τοῦ θέσθαι ἑαυτῷ ὄνομα μέγα καὶ ἐπιφανὲς, τοῦ ἐκβαλεῖν ἀπὸ προσώπου λαοῦ σου οὕς ἐλυτρώσω ἐξ Αἰγύπτου ἔθνη.
\vs{22}Καὶ ἔδωκας τὸν λαόν σου Ἰσραὴλ, σεαυτῷ λαὸν ἕως αἰῶνος, καὶ σὺ Κύριος ἐγενήθης αὐτοῖς εἰς Θεόν.
\vs{23}Καὶ νῦν, Κύριε, ὁ λόγος σου ὅν ἐλάλησας πρὸς τὸν παῖδά σου καὶ ἐπὶ τὸν οἶκον αὐτοῦ, πιστωθήτω ἕως αἰῶνος· καὶ ποιήσον καθὼς ἐλάλησας,
\vs{24}καὶ πιστωθήτω καὶ μεγαλυνθήτω τὸ ὄνομά σου ἕως αἰῶνος, λεγόντων, Κύριε Κύριε παντοκράτωρ Θεὸς Ἰσραὴλ, καὶ ὁ οἶκος Δαυὶδ παιδός σου ἀνωρθωμένος ἐναντίον σου.
\vs{25}Ὅτι σὺ Κύριος ὁ Θεός μου ἤνοιξας τὸ οὖς τοῦ παιδός σου τοῦ οἰκοδομῆσαι αὐτῷ οἶκον, διὰ τοῦτο εὗρεν ὁ παῖς σου τοῦ προσεύξασθαι κατὰ πρόσωπόν σου.
\vs{26}Καὶ νῦν, Κύριε, σὺ εἶ αὐτὸς Θεὸς, καὶ ἐλάλησας ἐπὶ τὸν δοῦλόν σου τὰ ἀγαθὰ ταῦτα.
\vs{27}Καὶ νῦν, ἦρξαι τοῦ εὐλογῆσαι τὸν οἶκον τοῦ παιδός σου, τοῦ εἶναι εἰς τὸν αἰῶνα ἐναντίον σου· ὅτι σὺ Κύριε εὑλόγησας, καὶ εὐλόγησον εἰς τὸν αἰῶνα.

\ch{18}
Καὶ ἐγένετο μετὰ ταῦτα, καὶ ἐπάταξε Δαυὶδ τοὺς ἀλλοφύλους καὶ ἐτροπώσατο αὐτοὺς, καὶ ἔλαβε τὴν Γὲθ καὶ τὰς κώμας αὐτῆς ἐκ χειρὸς ἀλλοφύλων.

\vs{2}Καὶ ἐπάταξε τὴν Μωὰβ, καὶ ἦσαν Μωὰβ παῖδες τῷ Δαυὶδ φέροντες δῶρα·

\vs{3}Καὶ ἐπάταξε Δαυὶδ τὸν Ἀδρααζὰρ βασιλέα Σουβὰ Ἠμάθ, πορευομένου αὐτοῦ ἐπιστῆσαι χεῖρα αὐτοῦ ἐπὶ ποταμὸν Εὐφράτην·
\vs{4}Καὶ προκατελάβετο Δαυὶδ αὐτῶν χίλια ἅρματα καὶ ἑπτὰ χιλιάδας ἵππων καὶ εἴκοσι χιλιάδας ἀνδρῶν πεζῶν· καὶ παρέλυσεν Δαυεὶδ πάντα τὰ ἅρματα, καὶ ὑπελίπετο ἐξ αὐτῶν ἑκατὸν ἅρματα.
\vs{5}Καὶ ἦλθεν Σύρος ἐκ Δαμασκοῦ βοηθῆσαι Ἀδρααζὰρ βασιλεῖ Σουβά, καὶ ἐπάταξε Δαυὶδ ἐν τῷ Σύρῳ εἴκοσι καὶ δύο χιλιάδας ἀνδρῶν.
\vs{6}Καὶ ἔθετο Δαυὶδ φρουρὰν ἐν Συρίᾳ τῇ κατὰ Δαμασκὸν, καὶ ἦσαν τῷ Δαυεὶδ εἰς παῖδας φέροντας δῶρα· καὶ ἔσωσε Κύριος Δαυίδ ἐν πᾶσιν οἷς ἐπορεύετο.
\vs{7}Καὶ ἔλαβεν Δαυεὶδ τοὺς κλοιοὺς τοὺς χρυσοῦς οἳ ἦσαν ἐπὶ τοὺς παῖδας Ἁδραάζαρ, καὶ ἤνεγκεν αὐτοὺς εἰς Ἰερουσαλήμ.
\vs{8}Καὶ ἐκ τῆς Μαγαβὲθ καὶ ἐκ τῶν ἐκλεκτῶν πόλεων τῶν Ἀδρααζὰρ ἔλαβε Δαυὶδ χαλκὸν πολὺν σφόδρα· ἐξ αὐτοῦ ἐποίησε Σαλωμὼν τὴν θάλασσαν τὴν χαλκῆν, καὶ τοὺς στύλους καὶ τὰ σκεύη τὰ χαλκᾶ.

\vs{9}καὶ ἤκουσε Θωὰ βασιλεὺς Ἠμὰθ, ὅτι ἐπάταξε Δαυὶδ τὴν πᾶσαν δύναμιν Ἀδρααζὰρ βασιλέως Σουβά·
\vs{10}Καὶ ἀπέστειλεν τὸν Ἀδουρὰμ υἱὸν αὐτοῦ πρὸς τὸν βασιλέα Δαυὶδ τοῦ ἐρωτῆσαι αὐτὸν τὰ εἰς εἰρήνην, καὶ του εὐλογῆσαι αὐτὸν ὑπὲρ οὗ ἐπολέμησε τὸν Ἁδρααζαρ, καὶ ἐπάταξεν αὐτόν, ὅτι ἀνὴρ πολέμιος Θῶα ἦν τῷ Ἁδραάζαρ· καὶ πάντα τὰ σκεύη τὰ χρυσᾶ, καὶ τὰ ἀργυρᾶ,
\vs{11}Καὶ τὰ χαλκᾶ, καὶ ταῦτα ἡγίασεν ὁ βασιλεὺς Δαυὶδ τῷ κυρίῳ, μετὰ τοῦ ἀργυρίου καὶ τοῦ χρυσίου οὗ ἔλαβεν ἐκ πάντων τῶν ἐθνῶν, ἐξ Ἰδουμαίας, καὶ Μωὰβ, καὶ ἐξ υἱῶν Ἀμμὼν, καὶ ἐκ τῶν ἀλλοφύλων, καὶ ἐξ Ἀμαλήκ.

\vs{12}Καὶ Ἀβεσὰ υἱὸς Σαρουίας ἐπάταξε τὴν Ἰδουμαίαν ἐν κοιλάδι τῶν ἁλῶν, ὀκτὼκαὶδεκα χιλιάδας.
\vs{13}Καὶ ἔθετο ἐν τῇ κοιλάδι φρουρὰς, καὶ ἦσαν πάντες οἱ Ἰδουμαῖοι παῖδες Δαυίδ· καὶ ἔσωζε Κύριος τὸν Δανὶδ ἐν πᾶσιν οἷς ἐπορεύετο.

\vs{14}Καὶ ἐβασίλευσε Δαυὶδ ἐπὶ πάντα Ἰσραὴλ, καὶ ἦν ποιῶν κρίμα καὶ δικαιοσύνην τῷ παντὶ λαῷ αὐτοῦ.
\vs{15}Καὶ Ἰωὰβ υἱὸς Σαρουίας ἐπὶ τῆς στρατιᾶς, καὶ Ἰωσαφὰτ υἱὸς Ἀχιλοὺδ ὁ ὑπομνηματογράφος,
\vs{16}καὶ Σαδὼκ υἱὸς Ἀχειτὼβ καὶ Αχειμέλεχ υἱὸς Ἀβιάθαρ οἱ ἱερεῖς, καὶ Σουσὰ γραμματεύς,
\vs{17}καὶ Βαναίας υἱὸς Ἰωδαὲ ἐπὶ τοῦ Χερεθὶ καὶ ἐπὶ τοῦ Φελεθί. καὶ υἱοὶ Δαυὶδ οἱ πρῶτοι διάδοχοι τοῦ βασιλέως.

\ch{19}
Καὶ ἐγένετο μετὰ ταῦτα ἀπέθανε Ναὰς βασιλεῦς υἱῶν Ἀμμών, καὶ ἐβασίλευσεν Ἀνὰν υἱὸς αὐτοῦ ἀντʼ αὐτοῦ.
\vs{2}Καὶ εἶπε Δαυίδ, Ποιήσω ἔλεος μετὰ Ἀνὰν υἱοῦ Ναὰς, ὡς ἐποίησεν ὁ πατὴρ αὐτοῦ μετʼ ἐμοῦ ἔλεος· καὶ ἀγγέλοους ἀγγέλους Δαυὶδ τοῦ παρακαλέσαι αὐτὸν περὶ τοῦ πατρὸς αὐτοῦ· καὶ ἦλθον παῖδες Δαυὶδ εἰς γῆν νἱῶν Αμμὼν πρὸς Ἀνὰν τοῦ παρακαλέσαι αὐτόν.
\vs{3}Καὶ εἶπον ἄρχοντες υἱῶν Ἀμμὼν πρὸς Ἀνάν, Μὴ δοξάζων Δαυὶδ τὸν πατέρα σου ἐναντίον σου ἀπέστειλέ σοι παρακαλοῦντας; οὐχ ὅπως ἐξερενήσωσι τὴν πόλιν, καὶ τοῦ κατασκοπῆσαι τὴν γῆν, ἦλθον παῖδες αὐτοῦ πρὸς σέ;
\vs{4}Καὶ ἔλαβεν Ἁνὰν τοὺς παῖδας Ααυὶδ, καὶ ἐξύρησεν αὐτοὺς, καὶ ἀφεῖλε τῶν μανδυῶν αὐτῶν τὸ ἥμισυ ἕως τῆς ἀναβολῆς, καὶ ἀπέστειλεν αὐτούς.
\vs{5}Καὶ ἦλθον ἀπαγγεῖλαι τῷ Δαυὶδ περὶ τῶν ἀνδρῶν· καὶ ἀπέστειλεν εἰς ἀπάντησιν αὐτοῖς, ὅτι ἦσαν ἠτιμωμένοι σφόδρα· καὶ εἶπεν ὁ βασιλεύς, Καθίσατε ἐν Ἱερειχὼ ἕως τοῦ ἀνατεῖλαι τοὺς πώγωνας ὑμῶν, καὶ ἀνακάμψατε.

\vs{6}Καὶ εἶδον οἱ υἱοὶ Ἀμμὼν ὅτι ἠπχὺνθη λαὸς Δαυίδ, καὶ ἀπέστειλεν Ἁνὰν καὶ οἱ υἱοὶ Ἀμμὼν χίλια τάλαντα ἀργυρίου τοῦ μισθώσασθαι ἑαυτοῖς ἐκ Συρίας Μεσοποταμίας καὶ ἐκ Συρίας Μααχὰ καὶ παρὰ Σωβὰλ ἅρματα καὶ ἱππεῖς.
\vs{7}καὶ ἐμισθώσαντο ἑαυτοῖς δύο καὶ τριάκοντα χιλιάδας ἁρμάτων, καὶ τὸν βασιλέα Μααχὰ καὶ τὸν λαὸν αὐτοῦ· καὶ ἦλθον καὶ παρενέβαλον κατέναντι Μηδαβά· καὶ οἱ υἱοὶ Ἀμμὼν συνήχθησαν ἐκ τῶν πόλεων αὐτῶν, καὶ ἦλθον εἰς τὸ πολεμῆσαι.

\vs{8}Καὶ ἤκουσεν Δαυίδ, καὶ απέστειλε τὸν Ἰωὰβ καὶ πᾶσαν τὴν στρατίαν τῶν δυνατῶν.
\vs{9}καὶ ἐξῆλθον οἱ υἱοὶ Ἀμμὼν, καὶ παρατάσσονται εἰς πόλεμον παρὰ τὸν πυλῶνα τῆς πόλεως· καὶ οἱ βασιλεῖς οἱ ἐλθόντες παρενέβαλον καθʼ ἑαυτοὺς ἐν τῷ πεδίῳ·
\vs{10}καὶ εἶδεν Ἰωὰβ ὅτι γεγόνασιν ἀντιπρόσωποι τοῦ πολεμεῖν πρὸς αὐτὸν κατὰ πρόσωπον καὶ ἐξόπισθε, καὶ ἐξελέξατο ἐκ παντὸς νεανίου ἐξ Ἰσραὴλ, καὶ παρετάξαντο ἐναντίον τοῦ Σύρου.
\vs{11}καὶ τὸ κατάλοιπον τοῦ λαοῦ ἔδωκεν ἐν χειρὶ Ἀβεσσὰ ἀδελφοῦ αὐτοῦ, καὶ παρετάξαντο ἐξεναντίας υἱῶν Ἀμμών.
\vs{12}καὶ εἶπεν, Ἐὰν κρατήσῃ ὑπὲρ ἐμὲ Σύρος, καὶ ἔσῃ μοι εἰς σωτηρίαν· καὶ ἐὰν υἱοὶ Ἀμμὼν κρατήσωσιν ὑπὲρ σέ, καὶ σώσω σε.
\vs{13}καὶ Κύριος τὸ ἀγαθὸν ποιήσει.

\vs{14}Καὶ παρετάξατο Ἰωὰβ καὶ ὁ λαὸς ὁ μετʼ αὐτοῦ κατέναντι Σύρων εἰς πόλεμον, καὶ ἔφυγον ἀπʼ αὐτοῦ.
\vs{15}καὶ οἱ υἱοὶ Ἀμμὼν εἶδον ὅτι ἔφυγον Σύροι, καὶ ἔφυγον καὶ αὐτοὶ ἀπὸ προσώπου Ἰωὰβ καὶ ἀπὸ προσώπου ἀδελφοῦ αὐτοῦ, καὶ ἦλθον εἰς τὴν πόλιν· καὶ ἦλθεν Ἰωὰβ εἰς Ἰερουσαλήμ.

\vs{16}Καὶ εἶδεν Σύρος ὅτι ἐτροπώσατο αὐτὸν Ἰσραήλ, καὶ ἀπέστειλαν ἀγγέλους· καὶ ἐξήγαγον τὸν Σύρον ἐκ τοῦ πέραν τοῦ ποταμοῦ, καὶ Σωφὰρ ἀρχιστράτηγος δυνάμεως Ἁδραάζαρ ἔμπροσθεν αὐτῶν.
\vs{17}καὶ ἀπηγγέλη τῷ Δαυείδ, καὶ συνήγαγεν τὸν πάντα Ἰσραήλ, καὶ διέβη τὸν Ἰορδάνην καὶ ἦλθεν ἐπʼ αὐτοὺς καὶ παρετάξατο ἐπʼ αὐτούς. καὶ παρατάσσεται Σύρος ἐξ ἐναντίας Δαυεὶδ καὶ ἐπολέμησαν αὐτόν.
\vs{18}καὶ ἔφυγεν Σύρος ἀπὸ προσώπου Δαυείδ, καὶ ἀπέκτεινεν Δαυεὶδ ἀπὸ τοῦ Σύρου ἐπτὰ χιλιάδας ἁρμάτων καὶ τεσσεράκοντα χιλιάδας πεζῶν, καὶ τὸν Σαφὰθ ἀρχιστράτηγον δυνάμεως ἀπέκτεινεν.
\vs{19}καὶ εἶδον παῖδες Ἁδραάζαρ ὅτι ἐπταίκασιν ἀπὸ προσώπου Ἰσραήλ, καὶ διέθεντο μετὰ Δαυεὶδ καὶ ἐδούλευσαν αὐτῷ· καὶ οὐκ ἠθέλησεν Σύρος τοῦ βοηθῆσαι Ἀμμὼν ἔτι.

\ch{20}
Καὶ ἐγένετο ἐν τῷ ἐπιόντι ἔτει ἐν τῇ ἐξόδῳ τῶν βασιλέων καὶ ἤγαγεν Ἰωὰβ πᾶσαν τὴν δὺναμιν τῆς στρατείας, καὶ ἔφθειραν τὴν χώραν υἱῶν Ἀμμών· καὶ ἦλθεν καὶ περιεκάθισεν τὴν Ῥάββαν. καὶ Δαυεὶδ ἐκάθητο ἐν Ἰερουσαλήμ· καὶ ἐπάταξεν τὴν Ῥαββὰ καὶ κατέσκαψεν αὐτήν.
\vs{2}καὶ ἔλαβεν Δαυεὶδ τὸν στέφανον Μολχὸλ βασιλέως αὐτῶν ἀπὸ τῆς κεφαλῆς αὐτοῦ, καὶ εὑρέθη ὁ σταθμὸς αὐτοῦ τάλαντον χρυσίου, καὶ ἐν αὐτῷ λίθος τίμιος, καὶ ἦν ἐπὶ τὴν κεφαλὴν Δαυείδ· καὶ σκῦλα τῆς πόλεως ἐξήνεγκεν πολλὰ σφόδρα.
\vs{3}καὶ τὸν λαὸν τὸν ἐν αὐτῇ ἐξήγαγεν καὶ διέπρισεν πρίοσιν καὶ ἐν σκεπάρνοις σιδηροῖς, καὶ οὕτως ἐποίησεν Δαυεὶδ τοῖς παισὶν υἱοῖς Ἀμμών· καὶ ἀνέστρεψεν Δαυεὶδ καὶ πᾶς ὁ λαὸς αὐτοῦ εἰς Ἰερουσαλήμ.

\vs{4}Καὶ ἐγένετο μετὰ ταῦτα καὶ ἐγένετο ἔτι πόλεμος ἐν Γάζερ μετὰ τῶν ἀλλοφύλων· τότε ἐπάταξεν Σοβοχαὶ Θωσαθεὶ τὸν Σαφοὺτ ἀπὸ τῶν υἱῶν τῶν γιγάντων καὶ ἐταπείνωσεν αὐτόν.

\vs{5}Καὶ ἐγένετο ἔτι πόλεμος μετὰ τῶν ἀλλοφύλων, καὶ ἐπάταξεν Ἐλλὰν υἱὸς Ἰαεὶρ τὸν Ἐλεμεὲ ἀδελφὸν Γολιὰθ τοῦ Γεθθαίου, καὶ ξύλον δόρατος αὐτοῦ ὡς ἀντίον ὑφαινόντων.

\vs{6}Καὶ ἐγένετο ἔτι πόλεμος ἐν Γέθ, καὶ ἦν ἀνὴρ ὑπερμεγέθης, καὶ δάκτυλοι αὐτοῦ ἓξ καὶ ἓξ, εἴκοσι τέσσαρες· καὶ οὗτος ἦν ἀπόγονος γιγάντων.
\vs{7}Καὶ ὠνείδισεν τὸν Ἰσραήλ, καὶ ἐπάταξεν αὐτὸν Ἰωναθὰν υἱὸς Σαμαά, υἱὸς ἀδελφοῦ Δαυείδ.
\vs{8}Οὗτος ἐγένετο Ῥαφὰ ἐν Γέθ· πάντες ἦσαν τέσσαρες γίγαντες, καὶ ἔπεσον ἐν χειρὶ Δαυεὶδ καὶ ἐν χειρὶ παίδων αὐτοῦ.

\ch{21}
Καὶ ἔστη διάβολος ἐν τῷ Ἰσραήλ, καὶ ἐπέσεισεν τὸν Δαυὶδ τοῦ ἀριθμῆσαι τὸν ἰσραήλ.
\vs{2}Καὶ εἶπεν ὁ βασιλεὺς δαυὶδ πρὸς Ἰωὰβ καὶ τοὺς ἄρχοντας τῆς δυνάμεως, πορεύθητε, ἀριθμήσατε τὸν Ἰσραὴλ ἀρὸ Βηρσαβεὲ καὶ ἕως Δᾶν, καὶ ἐνέγκατε πρὸς μὲ, καὶ γνώσομαι τὸν ἀριθμὸν αὐτῶν.
\vs{3}Καὶ εἶπεν Ἰωάβ, Προσθείη Κύριος ἐπὶ τὸν λαὸν αὐτοῦ, ὡς αὐτοὶ ἑκατονταπλασίως, καὶ οἱ ὀφθαλμοὶ τοῦ κυρίου μυυ τοῦ βασιλέως βλέποντες· πάντες τῷ κυρίῳ μου παῖδεσ· ἱνατί ζητεῖ κύριός μου τοῦτο; ἳνα μὴ γένηται εἰς ἁμαρτίαν τῷ Ἰσραήλ.
\vs{4}Τὸ δὲ ρῆμα τοῦ βασιλέως ἴσχυσεν ἐπὶ Ιωὰβ, καὶ ἐξῆλθεν Ιωὰβ, καὶ διῆθεν ἐν παντὶ Ἰσραὴλ, καὶ ἦλθεν εἰς Ἰερουσλήμ.
\vs{5}Καὶ ἔδωκεν Ἰωὰβ τὸν ἀριθμὸν τῆς ἐπισκέψεως τοῦ λαοῦ τῷ Δαυίδ· καὶ ἦν πᾶς Ἰσραὴλ χίλιαι χιλιάδες καὶ ἑκατὸν χιλιάδες ἀνδρῶν ἐσπασμένων μάχαιραν·
\vs{6}Καὶ τὸν Λευὶ καὶ τὸν Βενιαμεὶν οὐκ ἠρίθμησεν ἐν μέσῳ αὐτῶν, ὅτι κατίσχυσν λόγος τοῦ βασιλέως τὸν Ἰωάβ.

\vs{7}καὶ πονηρὸν ἐναντίον τοῦ Θεοῦ περὶ τοῦ πράγματος τούτου, καὶ ἐπάταξεν τὸν Ἰσραήλ.
\vs{8}Καὶ εἶπε Δαυὶδ πρὸς τὸν Θεόν, ἡμάρτηκα σφόδρα, ὅτι ἐποίησα τὸ πρᾶγμα τοῦτο, καὶ νῦν περίελε δὴ τὴν κακίαν παιδός σου, ὅτι ἐματαιώθην σφόδρα.

\vs{9}Καὶ ἐλάλησε Κύριος πρὸς Γὰδ τὸν ὁρῶντα,
\vs{10}Πορεύου καὶ λάλησον πρὸς Δαυὶδ, λέγων, Οὕτως λέγει Κύριος, τρία αἴρω ἐγὼ ἐπι σέ, ἔκλεξαι σεαυτῷ ἓν ἐξ αὐτῶν, καὶ ποιήσω σοι.
\vs{11}καὶ ἦλθεν Γὰδ πρὸς Δαυὶδ, καὶ εἶπεν αὐτῷ, Οὕτως λέγει Κύριος, Ἔκλεξαι σεαυτῷ
\vs{12}ἢ τρία ἔτη λιμοῦ, ἢ τρεῖς μῆνας φεύγειν σε ἐκ προσώπου ἐχθρῶν σου, καὶ μάχαιρα ἐξ ἐχθρῶν σου τοῦ ἐξολεθρεῦσαι, ἢ τρεῖς ἡμέρας ῥομφαίαν Κυρίου καὶ θάνατον ἐν τῇ γῇ, καὶ ἄγγελος Κυρίου ἐξολεθρεύων ἐν πάσῃ κληρονομίᾳ Ἰσραήλ· καὶ νῦν ἴδε τί ἀποκριθῶ τῷ ἀποστείλαντι λόγον.

\vs{13}Καὶ εἶπεν Δαυεὶδ πρὸς Γάδ Στενά μοι καὶ τὰ τρία σφόδρα· ἐμπεσοῦμαι δὴ εἰς χεῖρας Κυρίου, ὅτι πολλοὶ οἱ οἰκτειρμοὶ αὐτοῦ σφόδρα, καὶ εἰς χεῖρας ἀνθρώπων οὐ μὴ ἐμπέσω.

\vs{14}Καὶ ἔδωκεν Κύριος θάνατον ἐν Ἰσραήλ, καὶ ἔπεσον ἐξ Ἰσραὴλ ἑβδομήκοντα χιλιάδες ἀνδρῶν.
\vs{15}καὶ ἀπέστειλεν Θεὸς ἄγγελο εἰς Ἰερουσαλὴμ τοῦ ἐξολοθρεῦσαι αὐτήν. καὶ ὡς ἐξολόθρευεν, εἶδε Κύριος, καὶ μετεμελήθη ἐπὶ τῇ κακίᾳ. καὶ εἶπε τῷ ἀγγέλῳ τῷ ἐξολοθρεύονντι, Ἱκανούσθω σοι, ἄνες τὴν χεῖρά σου· καὶ ὁ ἄγγελος Κυρίου ἑστὼς ἐν τῷ ἅλῳ Ὀρνὰ τοῦ Ἰεβουσαίου.
\vs{16}Καὶ ἐπῆρε Δαυὶδ τοὺς ὀφθαλμοὺς αὐτοῦ, καὶ εἶδε τὸν ἄγγελον Κυρίου ἑστῶτα ἀναμέσον τῆς γῆς καὶ τοῦ οὐρανοῦ, καὶ ἡ ῥομφαία αὐτοῦ ἐσπασμένη ἐν τῇ χειρὶ αὐτοῦ ἐκτεταμένη ἐπὶ Ἰερουσαλὴμ, καὶ ἔπεσε Δαυὶδ καὶ οἱ πρεσβύτεροι περιβεβλημένοι ἐν σάκκοις ἐπὶ πρόσωπον αὐτῶν.
\vs{17}Καὶ εἶπεν Δαὶδ πρὸς τὸν Θθεόν, Οὐκ ἐγὼ εἶπα τοῦ ἀριθμῆσαι τῷ λαῷ; καὶ ἐγώ εἰμι ὁ ἁμαρτών, κακοποιῶν ἐκακοποίησα, καὶ ταῦτα τὰ πρόβατα τί ἐποίησαν; Κύριε ὁ Θθεός, γενηθήτω ἡ χείρ σου ἐν ἐμοὶ καὶ ἐν τῷ οἴκῳ τοῦ πατρός μου, καὶ μὴ ἐν τῷ λαῷ σου εἰς ἀπώλειαν, Κύριε.

\vs{18}Καὶ ἄγγελος Κυρίου εἶπε τῷ Γὰδ τοῦ εἰπεῖν πρὸς Δαυὶδ ἵνα ἀναβῇ τοῦ στῆσαι θυσιαστήριον κυρίῳ ἐν ἅλῳ Ὀρνὰ τοῦ Ἰεβουσαίου.
\vs{19}Καὶ ἀνέβη Δαυὶδ κατὰ τὸν λόγον Γὰδ, ὃν ἐλάλησεν ἐν ὀνόματι Κυρίου.
\vs{20}Καὶ ἐπέστρεψεν Ὀρνά, καὶ εἶδε τὸν βασιλέα καὶ τέσσαρας υἱοὺς αὐτοῦ μετʼ αὐτοῦ μεθʼ ἁχαβίν· καὶ Ὀρνὰ ἦν ἀλοῶν πυρούς.
\vs{21}Καὶ ἦλθεν Δαυὶδ πρὸς Ὀρνᾶν, καὶ Ὀρνὰ ἐξῆλθεν ἐκ τῆς ἅλω καὶ προσεκύνησεν τῷ Δαυὶδ τῷ προσώπῳ ἐπὶ τὴν γῆν.
\vs{22}Καὶ εἶπε Δαυὶδ πρὸς Ὀρνὰ, δός μοι τὸν τόπον σου τῆς ἅλω, καὶ οἰκοδομήσω ἐπʼ αὐτῷ θυσιαστήριον τῷ κυρίῳ· ἐν ἀργυρίῳ ἀξίῳ δός μοι αὐτόν, καὶ παύσεται ἡ πληγὴ ἐκ τοῦ λαοῦ.
\vs{23}Καὶ εἶπεν Ὀρνὰ πρὸς Δαυίδ, Λάβε σεαυτῷ, καὶ ποιησάτω ὁ κύριός μου ὁ βασιλεὺς τὸ ἀγαθὸν ἐναντίον εαὐτοῦ· ἴδε δέδωκα τοὺς μόσχους εἰς ὁλοκαύτωσιν, καὶ τὸ ἄροτρον εἰς ξύλα, καὶ τὸν σῖτον εἰς θυαίαν, τὰ πάντα δέδωκα.
\vs{24}Καὶ εἶπεν ὁ βασιλεὺς Δαυὶδ τῷ Ὀρνὰ, Οὐχὶ, ὅτι ἀγοράζων ἀγοράσω ἐν ἀργυρίῳ ἀξίῳ, ὅτι οὐ μὴ λάβω ἅ ἐστί σοι Κυρίῳ, τοῦ ἀνενέγκαι ὁλοκαύτωσιν δωρεὰν Κυρίῳ.
\vs{25}καὶ ἔδωκεν Δαυὶδ τῷ Ὀρνὰ ἐν τῷ τόπῳ αὐτοῦ σίκλους χρυσίου ὁλκῆς ἑξακοσίους.
\vs{26}Καὶ ᾠκοδόμησεν ἐκεῖ Δανδ θυσιαστήριον Κυρίῳ, καὶ ἀνήνεγκεν ὁλοκαυτώματα καὶ σωτηρίου· καὶ ἐβόησε πρὸς Κύριον, καὶ ἐπήκουσεν αὐτῷ ἐν πυρὶ ἐκ τοῦ οὐρανοῦ ἐπὶ τὸ θυσιαστήριον τῆς ὁλοκαυτώσεως καὶ κατανάλωσεν τὴν ὁλοκαύτωσιν.
\vs{27}Καὶ εἶπε Κύριος πρὸς τὸν ἄγγελον· καὶ κατέθηκε τὴν ῥομφαίαν εἰς τὸν κολεὸν αὐτῆς.

\vs{28}Ἐν τῷ καιρῷ ἐκείνῳ ἐν τῷ ἰδεῖν τὸν Δαυὶδ ὅτι ἐπήκουσεν αὐτῷ Κύριος ἐν ἅλῳ Ὀρνὰ τοῦ Ἰεβουσαίου, καὶ ἐθυσίασεν ἐκεῖ.
\vs{29}Καὶ σκηνὴ Κυρίου ἣν ἐποίησε Μωυσῆς ἐν τῇ ἐρήμῳ, καὶ θυσιαστήριον τῶν ὁλοκαυτωμάτων ἐν τῷ καιρῷ ἐκείνῳ ἐν Βαμὰ ἐν Γαβαών.
\vs{30}Καὶ οὐκ ἐδύνατο Δαυὶδ τοῦ πορευθῆναι ἔμπροσθεν αὐτοῦ τοῦ ζητῆσαι τὸν Θεόν, ὅτι οὐ κατέσπευσεν ἀπὸ προσώπου τῆς ῥομφαίας ἀγγέλου Κυρίου.

\ch{22}
Καὶ εἶπε Δαυὶδ, Οὗτός ἐστιν ὁ οἶκος Κυρίου τοῦ Θεοῦ, καὶ τοῦτο τὸ θυσιαστήριον εἰς ὁλοκαύτωσιν τῷ Ἰσραήλ.

\vs{2}Καὶ εἶπε Δαυὶδ συναγαγεῖν πάντας τοὺς προσηλύτους τοὺς ἐν γῇ Ἰσραὴλ, καὶ κατέστησε λατόμους λατομῆσαι λίθους ξυστοὺς τοῦ οἰκοδομῆσαι οἶκον τῷ Θεῷ.
\vs{3}Καὶ σίδηρον πολὺν εἰς τοὺς ἥλους τῶν θυρωμάτων καὶ τῶν πυλῶν, καὶ τοὺς στροφεῖς ἡτοίμασε Δαυὶδ καὶ χαλκὸν εἰς πλῆθος, οὐκ ἦν σταθμός.
\vs{4}Καὶ ξύλα κέδρινα, οὐκ ἦν ἀριθμός· ὅτι ἐφέροσαν οἱ Σιδώνιοι καὶ οἱ Τύριοι ξύλα κέδρινα εἰς πλῆθος τῷ Δαυίδ.
\vs{5}Καὶ εἶπε Δαυὶδ, Σαλωμὼν ὁ υἱός μου παιδάριον ἁπαλὸν, καὶ ὁ οἶκος τοῦ οἰκοδομῆσαι τῷ Κυρίῳ εἰς μεγαλωσύνην ἄνω, εἰς ὄνομα καὶ εἰς δόξαν εἰς πᾶσαν τὴν γῆν· ἑτοιμάσω αὐτῷ· καὶ ἡτοίμασε Δαυεὶδ εἰς πλῆθος ἔμπροσθεν τῆς τελευτῆς αὐτοῦ.

\vs{6}Καὶ ἐκάλεσε Σαλωμὼν τὸν υἱὸν αὐτοῦ, καὶ ἐνετείλατο αὐτῷ τοῦ οἰκοδομῆσαι τὸν οἶκον τῷ Κυρίῳ Θεῷ Ἰσραήλ.
\vs{7}Καὶ εἶπε Δαυὶδ Σαλωμών, τέκνον, ἐμοὶ ἐγένετο ἐπὶ ψυχῇ τοῦ οἰκοδομῆσαι οἶκον τῷ ὀνόματι Κυρίου Θεοῦ.
\vs{8}Καὶ ἐγένετό μοι λόγος Κυρίου, λέγων, αἷμα εἰς πλῆθος ἐξέχεας, καὶ πολέμους μεγάλους ἐποίησας· οὐκ οἰκοδομήσεις οἶκον τῷ ὀνόματί μου, ὅτι αἵματα πολλὰ ἐξέχεας ἐπὶ τὴν γῆν ἐναντίον ἐμου.
\vs{9}Ἰδοὺ υἱὸς τίκτεταί σοι, οὗτος ἔσται ἀνὴρ ἀναπαύσεως, καὶ ἀναπαύσω αὐτὸν ἀπὸ πάντων τῶν ἐχθρῶν αὐτοῦ κυκλόθεν, ὅτι Σαλωμὼν ὄνομα αὐτῷ, καὶ εἰρήνην καὶ ἡσυχίαν δώσω ἐπὶ Ἰσραὴλ ἐν ταῖς ἡμέραις αὐτοῦ.
\vs{10}Οὗτος οἰκοδομήσει οἶκον τῷ ὀνόματί μου, καὶ οὗτος ἔσται μοι εἰς υἱὸν, κἀγὼ αὐτῷ εἰς πατέρα, καὶ ἀνορθώσω θρόνον βασιλείας αὐτοῦ ἐν Ἰσραὴλ ἕως αἰῶνος.
\vs{11}Καὶ νῦν, υἱέμου, ἔσται μετὰ σοῦ Κύριος, καὶ εὐοδώσει, καὶ οἰκοδομήσεις οἶκον τῷ Κυρίῳ Θεῷ σου. ὡς ἐλάλησε περὶ σοῦ.
\vs{12}Ἀλλʼ ἢ δῴη σοι σοφίαν καὶ σύνεσιν Κύριος καὶ κατισχύσαι σε ἐπὶ Ἰσραὴλ, καὶ τοῦ φυλάσσεσθαι καὶ τοῦ ποιεῖν τὸν νόμον Κυρίου τοῦ Θεοῦ σου.
\vs{13}Τότε εὐοδώσει ἐὰν φυλάξῃς τοῦ ποιεῖν τὰ προστάγματα καὶ τὰ κρίματα ἃ ἐνετείλατο Κύριος τῷ Μωυσῇ ἐπὶ Ἰσραήλ· ἀνδρίζου καὶ ἴσχυε, μὴ φοβοῦ μηδὲ πτοηθῇς.

\vs{14}Καὶ ἰδοὺ ἐγὼ κατὰ τὴν πτωχείαν μου ἡτοίμασα εἰς οἶκον Κυρίου χρυσίου ταλάντων ἑκατὸν χιλιάδας, καὶ ἀργυρίου ταλάντων χιλίας χιλιάδας, καὶ χαλκὸν καὶ σίδηρον οὗ οὐκ ἔστι σταθμὸς, ὅτιι εἰς πλῆθός ἐστι· καὶ ξύλα καὶ λίθους ἡτοίμασα, καὶ πρὸς ταῦτα πρόσθες.
\vs{15}Καὶ μετὰ σοῦ πρόσθες εἰς πλῆθος ποιούντων ἔργα, τεχνῖται καὶ οἰκοδόμοι λίθων, καὶ τέκτονες ξύλων, καὶ πᾶς σοφὸς ἐν παντὶ ἔργῳ,
\vs{16}ἐν παντὶ ἔργῳ, ἐν χρυσίῳ καὶ ἀργυρίῳ, χαλκῷ καὶ ἐν σιδήρῳ, οὐκ ἔστιν ἀριθμός· ἀνάστηθι καὶ ποίει, καὶ Κύριος μετὰ σοῦ.

\vs{17}Καὶ ἐνετείλατο Δαυὶδ τοῖς πᾶσιν ἄρχουσιν Ἰσραὴλ ἀντιλαβέσθαι τῷ Σαλωμὼν υἱῷ αὐτοῦ.
\vs{18}Οὐχὶ Κύριος μεθʼ ὑμῶν; καὶ ἀνέπαυσεν ὑμᾶς κυκλόθεν, ὅτι ἔδωκεν ἐν ὑμῶν τοὺς κατοικοῦντας τὴν γῆν, καὶ ὑπετάγη ἡ γῆ ἐναντίον Κυρίου καὶ ἐναντίον λαοῦ αὐτοῦ.
\vs{19}Νῦν δότε καρδίας ὑμῶν καὶ ψυχὰς ὑμῶν τοῦ ζητῆσαι τῷ κυρίῳ θεῷ ὑμῶν, καὶ ἐγέρθητε καὶ οἰκοδομήσατε ἁγίασμα τῷ Θεῷ ὑμῶν, τοῦ εἰσενέγκαι τὴν κιβωτὸν διαθήκης Κυρίου, καὶ σκεύη τὰ ἅγια τοῦ Θεοῦ εἰς οἶκον τὸν οἰκοδομούμενον τῷ ὀνόματι Κυρίου.

\ch{23}
Καὶ Δαυὶδ πρεσβύτης καὶ πλήρης ἡμερῶν, καὶ ἐβασίλευσε Σαλωμὼν τὸν υἱὸν αὐτοῦ ἀνθʼ αὐτοῦ ἐπὶ Ἰσραήλ.
\vs{2}καὶ συνήγαγε τοὺς πάντας ἄρχοντας Ἰσραὴλ καὶ τοὺς ἱερεῖς καὶ τοὺς Λευείτας.
\vs{3}Καὶ ἠρίθμησαν οἱ Λευῖται ἀπὸ τριακονταετοῦς καὶ ἐπάνω, καὶ ἐγένετο ὁ ἀριθμὸς αὐτῶν κατὰ κεφαλὴν αὐτῶν εἰς ἄνδρας τριάκοντα καὶ ὀκτὼ χιλιάδας.
\vs{4}Ἀπὸ τῶν ἐργοδιωκτῶν ἐπὶ τὰ ἔργα οἴκου Κυρίου εἴκοσι τέσσαρες χιλιάδες, καὶ γραμματεῖς καὶ κριταὶ ἑξακισχίλιοι,
\vs{5}καὶ τέσσαρες χιλιάδες πυλωροὶ καὶ τώσσαρες χιλιάδες αἰνοῦντες τῷ κυρίῳ ἐν ὀργάνοιν οἷς ἐποίησε τοῦ αἰνεῖν τῷ κυρίῳ.

\vs{6}καὶ διεῖλεν αὐτοὺς Δαυὶδ ἐφημερίας τοῖς υἱοῖς Λευῖ, τῷ Γεδσών, Καάθ, καὶ Μαραρί·
\vs{7}καὶ τῷ Γεδσὼν, Ἐδὰν, καὶ Σεμεΐ.
\vs{8}γἱοὶ τῷ Ἐδάν, ἄρχων Ἰεϊὴλ, καὶ Ζηθὸμ, καὶ Ἰωήλ, καὶ τρεῖς.
\vs{9}γἱοὶ Σεμεῒ, Σαλωμὶθ, Ἰεϊὴλ, καὶ Δὰν, τρεῖς· οὗτοι ἄρχοντες πατριῶν τῶν Ἐδάν·
\vs{10}καὶ τοῖς υἱοῖς Σεμεΐ· Ἰὲθ, καὶ Ζιζὰ, καὶ Ἰῶας, καὶ Βεριά· οὗτοι υἱοὶ Σεμεΐ, τέσσαρες.
\vs{11}καὶ ἦν Ἰὲθ ὁ ἄρχων· καὶ Ζιζὰ ὁ δεύτερος· καὶ Ἰωὰς καὶ Βεριὰ οὐκ ἐπλήθυναν υἱοὺς, καὶ ἐγένετο εἰς οἶκον πατριᾶς εἰς ἐπίσκεψιν μίαν.

\vs{12}γἱοὶ Καάθ, Ἀμβρὰμ, Ἰσαάρ, Χεβρών, Ὀζιὴλ, τέσσαρες.
\vs{13}γἱοὶ Ἀμβρὰμ, Ἀαρὼν καὶ Μωυσῆς· καὶ διεστάλη Ἀαρὼν τοῦ ἁγιασθῆναι ἅγια ἁγίων, αὐτὸς καὶ οἱ υἱοὶ αὐτοῦ ἕως αἰῶνος, τοῦ θυμιᾷν ἐναντίον τοῦ κυρίου, λειτουργεῖν καὶ ἐπεύχεσθαι ἐπὶ τῷ ὀνόματι αὐτοῦ ἕως αἰῶνος.
\vs{14}Καὶ Μωυσῆς ἄνθρωπος τοῦ Θεοῦ, υἱοὶ αὐτοῦ ἐκλήθησαν εἰς φυλὴν τοῦ Λευεί.
\vs{15}γἱοὶ Μωυσῆ, Γηρσὰμ, καὶ Ἐλιέζερ.
\vs{16}γἱοὶ Γηρσὰμ, Σουβαὴλ ὁ ἄρχων.
\vs{17}Καὶ ἦσαν υἱοὶ τῷ Ἐλιέζερ, Ῥαβιὰ ὁ ἄρχων· καὶ οὐκ ἦσαν τῷ Ἐλιέζερ υἱοὶ ἕτεροι· καὶ υἱοὶ Ῥαβιὰ ηὐξήθησαν εἰς ὕψος.
\vs{18}γἱοὶ Ἰσαάρ, Σαλωμὼθ ὁ ἄρχων.
\vs{19}γἱοὶ Χεβρὼν, Ἰεριὰ ὁ ἄρχων, Ἀμαριὰ ὁ δεύτερος, Ἱεζιὴλ ὁ τρίτος, Ἰκεμίας ὁ τέταρτος.
\vs{20}γἱοὶ Ὀζιὴλ, Μιχὰ ὁ ἄρχων, καὶ Ἰσιὰ ὁ δεύτερος.

\vs{21}γἱοὶ Μεραρί, Μοολὶ καὶ Ὁ Μουσί· υἱοὶ Μοολὶ, Ἐλεαζαρ, καὶ Κείς. καὶ Κίς.
\vs{22}Καὶ ἀπέθανεν Ἐλεάζαρ· καὶ οὐκ ἦσαν αὐτῷ υἱοὶ, ἀλλʼ ἢ θυγατέρες· καὶ ἔλαβον αὐτὰς υἱοὶ Κὶς ἀδελφοὶ αὐτῶν.
\vs{23}Υἱοὶ Μουσὶ, Μοολὶ, καὶ Ἐδὲρ, καὶ Ἰαριμὼθ, τρεῖς.

\vs{24}Οὗτοι υἱοὶ Λευὶ κατʼ οἴκους πατριῶν αὐτῶν, ἄρχοντες τῶν πατριῶν αὐτῶν κατὰ τὴν ἐπίσκεψιν αὐτῶν, κατὰ τὸν ἀριθμὸν ὀνομάτων αὐτῶν, κατὰ κεφαλὴν αὐτῶν, ποιοῦντες τὰ ἔργα λειτουργείας οἴκου Κυρίου ἀπὸ εἰκοσαετοῦς καὶ ἐπάνω.
\vs{25}Ὅτι εἶπε Δαυὶδ, κατέπαυσε Κύριος ὁ Θεὸς Ἰσραὴλ τῷ λαῷ αὐτοῦ, καὶ κατεσκήνωσεν ἐν Ἱερουσαλὴμ ἕως αἰῶνος.
\vs{26}Καὶ οἱ Λευῖται οὐκ ἦσαν αἴροντες τὴν σκηνὴν καὶ τὰ πάντα σκεύη αὐτῆς εἰς τὴν λειτουργείαν αὐτῆς·
\vs{27}Ὅτι ἐν τοῖς λόγοις Δαυὶδ τοῖς ἐσχάτοις ἐστὶν ὁ ἀριθμὸς υἱῶν Λευὶ ἀπὸ εἰκοσαετοῦς καὶ ἐπάνω·
\vs{28}Ὅτι ἔστησεν αὐτοὺς ἐπὶ χειρὶ Ἀαρὼν, τοῦ λειτουργεῖν ἐν οἴκῳ Κυρίου ἐπὶ τὰς αὐλὰς, καὶ ἐπὶ τὰ παστοφόρια, καὶ ἐπὶ τὸν καθαρισμὸν τῶν πάντων ἁγίων, καὶ ἐπὶ τὰ ἔργα λειτουργείας οἴκου τοῦ Θεοῦ,
\vs{29}καὶ εἰς τοὺς ἄρτους τῆς προθέσεως, καὶ εἰς τὴν σεμίδαλιν τῆς θυσίας, καὶ εἰς τὰ λάγανα τὰ ἄζυμα, καὶ εἰς τήγανον, καὶ εἰς τὴν πεφυραμένην, καὶ εἰς πᾶν μέτρον,
\vs{30}καὶ τοῦ στῆναι πρωῒ τοῦ αἰνεῖν καὶ ἐξομολογεῖσθαι τῷ Κυρίῳ, καὶ οὕτω τοεσπέρας·
\vs{31}Καὶ ἐπὶ πάντων τῶν ἀναφερομένων ὁλοκαυτωμάτων τῷ Κυρίῳ ἐν τοῖς σαββάτοις καὶ ἐν ταῖς νεομηνίαις καὶ ἐν ταῖς ἑορταῖς, κατὰ ἀριθμὸν, κατὰ τὴν κρίσιν ἐπʼ αὐτοῖς διαπαντὸς τῷ Κυρίῳ.
\vs{32}Καὶ φυλάξουσι τὰς φυλακὰς σκηνῆς τοῦ μαρτυρίου, καὶ τὴν φυλακὴν τοῦ ἁγίου, καὶ τὰς φυλακὰς υἱῶν Ἀαρὼν ἀδελφῶν αὐτῶν, τοῦ λειτουργεῖν ἐν οἴκῳ Κυρίου.

\ch{24}
Καὶ τοὺς υἱοὺς Ἀαρὼν διαιρέσει Ναδὰβ, καὶ Ἀβιοὺδ, καὶ Ἐλεάζαρ, καὶ Ἰθάμαρ.
\vs{2}Καὶ ἀπέθανε Ναδὰβ καὶ Ἀβιοὺδ ἐναντίον τοῦ πατρὸς αὐτῶν, καὶ υἱοὶ οὐκ ἦσαν αὐτοῖς· καὶ ἱεράτευσεν Ἐλεάζαρ καὶ Ἰθάμαρ υἱοὶ Ἀαρών·
\vs{3}Καὶ διεῖλεν αὐτοὺς Δαυὶδ, καὶ Σαδὼκ ἐκ τῶν υἱῶν Ἐλεάζαρ, καὶ Ἀχιμέλεχ ἐκ τῶν υἱῶν Ἰθάμαρ, κατὰ τὴν ἐπίσκεψιν αὐτῶν, κατὰ τὴν λειτουργείαν αὐτῶν, κατʼ οἴκους πατριῶν αὐτῶν.

\vs{4}Καὶ εὑρέθησαν οἱ υἱοὶ Ἐλεάζαρ πλείους εἰς ἄρχοντας τῶν δυνατῶν παρὰ τοὺς υἱοὺς Ἰθάμαρ· καὶ διεῖλεν αὐτοὺς, τοῖς υἱοῖς Ἐλεάζαρ ἄρχοντας εἰς οἴκους πατριῶν ἑκκαίδεκα, τοῖς υἱοῖς Ἰθάμαρ κατʼ οἴκους πατριῶν ὀκτώ.
\vs{5}Καὶ διεῖλεν αὐτοὺς κατὰ κλήρους τούτους πρὸς τούτους, ὅτι ἦσαν ἄρχοντες τῶν ἁγίων, καὶ ἄρχοντες Κυρίου ἐν τοῖς υἱοῖς Ἐλεάζαρ καὶ ἐν τοῖς υἱοῖς Ἰθάμαρ.

\vs{6}Καὶ ἔγραψεν αὐτοὺς Σαμαΐας υἱὸς Ναθαναὴλ ὁ γραμματεὺς ἐκ τοῦ Λευὶ κατέναντι τοῦ βασιλέως καὶ τῶν ἀρχόντων, καὶ Σαδὼκ ὁ ἱερεὺς, καὶ Ἀχιμέλεχ υἱὸς Ἀβιάθαρ, καὶ ἄρχοντες τῶν πατριῶν τῶν ἱερέων καὶ τῶν Λευιτῶν οἴκου πατριᾶς, εἷς εἷς τῷ Ἐλεάζαρ, καὶ εἷς εἷς τῷ Ἰθάμαρ.

\vs{7}Καὶ ἐξῆλθεν ὁ κλῆρος ὁ πρῶτος τῷ Ἰωαρὶμ, τῷ Ἰεδίᾳ ὁ δεύτερος,
\vs{8}τῷ Χαρὶβ ὁ τρίτος, τῷ Σεωρὶμ ὁ τέταρτος,
\vs{9}τῷ Μελχίᾳ ὁ πέμπτος, τῷ Μεϊαμὶν ὁ ἕκτος,
\vs{10}τῷ Κὼς ὁ ἕβδομος, τῷ Ἀβίᾳ ὁ ὄγδοος,
\vs{11}τῷ Ἰησοῦ ὁ ἔννατος, τῷ Σεχενίᾳ ὁ δέκατος,
\vs{12}τῷ Ἐλιαβὶ ὁ ἑνδέκατος, τῷ Ἰακὶμ ὁ δωδέκατος,
\vs{13}τῷ Ὀπφᾷ ὁ τρισκαιδέκατος, τῷ Ἰεσβαὰλ ὁ τεσσαρεσκαιδέκατος, τῷ Βελγὰ ὁ πεντεκαιδέκατος,
\vs{14}τῷ Ἐμμὴρ ὁ ἑκκαιδέκατος, τῷ Χηζὶν ὁ ἑπτακαιδέκατος,
\vs{15}τῷ Ἀφεσὴ ὁ ὀκτωκαιδέκατος, τῷ Φεταίᾳ ὁ ἐννεακαιδέκατος,
\vs{16}τῷ Ἐζεκὴλ ὁ εἰκοστὸς, τῷ Ἀχὶμ ὁ εἷς καὶ εἰκοστὸς,
\vs{17}τῷ Γαμοὺλ ὁ δεύτερος καὶ εἰκοστὸς, τῷ Ἀδαλλαὶ ὁ τρίτος καὶ εἰκοστὸς,
\vs{18}τῷ Μαασαὶ ὁ τέταρτος καὶ εἰκοστός.

\vs{19}Αὕτη ἡ ἐπίσκεψις αὐτῶν κατὰ τὴν λειτουργίαν αὐτῶν τοῦ εἰσπορεύεσθαι εἰς οἶκον Κυρίου κατὰ τὴν κρίσιν αὐτῶν διὰ χειρὸς Ἀαρὼν πατρὸς αὐτῶν, ὡς ἐνετείλατο Κύριος ὁ Θεὸς Ἰσραήλ.

\vs{20}Καὶ τοῖς υἱοῖς Λευὶ τοῖς καταλοίποις, τοῖς υἱοῖς Ἀμβρὰμ, Σωβαήλ· τοῖς υἱοῖς Σωβαὴλ Ἰεδία.
\vs{21}Τῷ Ῥααβία ὁ ἄρχων.
\vs{22}Καὶ τῷ Ἰσααρὶ, Σαλωμώθ· τοῖς υἱοῖς Σαλωμὼθ, Ἰάθ.
\vs{23}υἱοὶ Ἐκδιοῦ, Ἀμαδία ὁ δεύτερος, Ἰαζιὴλ ὁ τρίτος, Ἰεκμοὰμ ὁ τέταρτος.
\vs{24}Τοῖς υἱοῖς Ὀζιὴλ, Μιχά· υἱοὶ Μιχὰ, Σαμήρ·
\vs{25}Ἀδελφὸς Μιχὰ, Ἰσία· υἱὸς Ἰσία, Ζαχαρία.
\vs{26}Υἱοὶ Μεραρὶ, Μοολὶ καὶ ὁ Μουσί·
\vs{27}υἱοὶ Ὀζία τοῦ Μεραρὶ τῷ Ὀζίᾳ· υἱοὶ αὐτοῦ Ἰσοὰμ, καὶ Σακχοὺρ, καὶ Ἀβαΐ.
\vs{28}Τῷ Μοολὶ Ἐλεάζαρ, καὶ Ἰθάμαρ· καὶ ἀπέθανεν Ἐλεάζαρ καὶ οὐκ ἦσαν αὐτῷ υἱοί.
\vs{29}Τῷ Κὶς, υἱοὶ τοῦ Κὶς Ἱεραμεήλ.
\vs{30}Καὶ υἱοὶ τοῦ Μουσὶ, Μοολὶ, καὶ Ἐδὲρ, καὶ Ἰεριμώθ· οὗτοι υἱοὶ τῶν Λευιτῶν κατʼ οἴκους πατριῶν αὐτῶν.
\vs{31}Καὶ ἔλαβον καὶ αὐτοὶ κλήρους καθὼς οἱ ἀδελφοὶ αὐτῶν υἱοὶ Ἀαρὼν ἐναντίον τοῦ βασιλέως, καὶ Σαδὼκ, καὶ Ἀχιμέλεχ, καὶ οἱ ἄρχοντες τῶν πατριῶν τῶν ἱερέων καὶ τῶν Λευειτῶν πατριάρχαι Ἀραὰβ, καθὼς οἱ ἀδελφοὶ αὐτοῦ οἱ νεώτεροι.

\ch{25}
Καὶ ἔστησε Δαυεὶδ ὁ βασιλεὺς καὶ οἱ ἄρχοντες τῆς δυνάμεως εἰς τὰ ἔργα τοὺς υἱοὺς Ἀσὰφ, καὶ Αἰμὰν, καὶ Ἰδιθοὺν, τοὺς ἀποφεγγομένους ἐν κινύραις, καὶ ἐν νάβλαις, καὶ ἐν κυμβάλοις· καὶ ἐγένετο ὁ ἀριθμὸς αὐτῶν κατὰ κεφαλὴν αὐτῶν ἐργαζομένων ἐν τοῖς ἔργοις αὐτῶν.

\vs{2}Υἱοὶ Ἀσάφ, Σακχοὺρ, Ἰωσὴφ, καὶ Ναθαυίας, καὶ Ἐραήλ· υἱοὶ Ἀσὰφ ἐχόμενοι τοῦ βασιλέως.

\vs{3}Τῷ Ἰδιθοὺν, υἱοὶ Ἰδιθοὺν, Τοδολίας, καὶ Σουρὶ, καὶ Ἰσέας, καὶ Ἀσαβίας, καὶ Ματθαθίας, ἕξ μετὰ τὸν πατέρα αὐτῶν Ἰδιθὺν, ἐν κινύρᾳ ἀνακρουόμενοι ἐξομολόγησιν καὶ αἴνεσιν τῷ κυρίῳ.

\vs{4}Τῷ Αἰμὰν, υἱοὶ Αἰμὰν, Βουκίας, καὶ Ματθανίας, καὶ Ὀζιὴλ, καὶ Σουβαὴλ, καὶ Ἰεριμὼθ, καὶ Ἀνανίας, καὶ Ἀνὰν, καὶ Ἑλιαθὰ, καὶ Γοδολλαθὶ, καὶ Ῥωμετθιέζερ, καὶ Ἰεσβασακὰ, καὶ Μαλλιθὶ, καὶ Ὠθηρὶ, καὶ Μεαζώθ.
\vs{5}Πάντες οὗτοι υἱοὶ τῷ Αἰμὰν τῷ ἀνακρουομένῳ τῷ βασιλεῖ ἐν λόγοις Θεοῦ, ὑψῶσαι κέρας· καὶ ἔδωκεν ὁ Θεὸς τῷ Αἰμὰν υἱοὺς τεσσαρεσκαίδεκα, καὶ θυγατέρας τρεῖς.
\vs{6}Πάντες οὗτοι μετὰ τοῦ πατρὸς αὐτῶν ὑμνῳδοῦντες ἐν οἴκῳ Θεοῦ, ἐν κυμβάλοις, καὶ ἐν νάβλαις, καὶ ἐν κινύραις εἰς τὴν δουλείαν οἴκου τοῦ Θεοῦ, ἐχόμενα τοῦ βασιλέως, καὶ Ἀσὰφ, καὶ Ἰδιθοὺν, καὶ Αἱμάν.

\vs{7}Καὶ ἐγένετο ὁ ἀριθμὸς αὐτῶν μετὰ τοὺς ἀδελφοὺς αὐτῶν δεδιδαγμένοι ᾄδειν Κυρίῳ πᾶς συνιὼν, διακόσιοι ὀγδοήκοντα καὶ ὀκτώ.

\vs{8}Καὶ ἔβαλον καὶ αὐτοὶ κλήρους ἐφημεριῶν κατὰ τὸν μικρὸν καὶ κατὰ τὸν μέγαν τελείων καὶ μανθανόντων.
\vs{9}Καὶ ἐξῆλθεν ὁ κλῆρος ὁ πρῶτος υἱῶν αὑτοῦ καὶ ἀδελφῶν αὐτοῦ τῷ Ἀσὰφ τοῦ Ἰωσὴφ, Γοδολίας· ὁ δεύτερος Ἡνεία, υἱοὶ αὐτοῦ καὶ ἀδελφοὶ αὐτοῦ δεκαδύο·
\vs{10}Ὁ τρίτος Ζακχοὺρ, υἱοὶ αὐτοῦ καὶ ἀδελφοὶ αὐτοῦ δεκαδύο·
\vs{11}Ὁ τέταρτος Ἰεσρὶ, υἱοὶ αὐτοῦ καὶ ἀδελφοὶ αὐτοῦ δεκαδύο·
\vs{12}Ὁ πέμπτος Νάθαν, υἱοὶ αὐτοῦ καὶ ἀδελφοὶ αὐτοῦ δεκαδύο·
\vs{13}Ὁ ἕκτος Βουκίας, υἱοὶ αὐτοῦ καὶ ἀδελφοὶ αὐτοῦ δεκαδύο·
\vs{14}Ὁ ἕβδομος Ἰσεριὴλ, υἱοὶ αὐτοῦ καὶ ἀδελφοὶ αὐτοῦ δεκαδύο·
\vs{15}Ὁ ὄγδοος Ἰωσία, υἱοὶ αὐτοῦ καὶ ἀδελφοὶ αὐτοῦ δεκαδύο·
\vs{16}Ὁ ἔννατος Ματθανίας, υἱοὶ αὐτοῦ καὶ ἀδελφοὶ αὐτοῦ δεκαδύο·
\vs{17}Ὁ δέκατος Σεμεΐα, υἱοὶ αὐτοῦ καὶ ἀδελφοὶ αὐτοῦ δεκαδύο·
\vs{18}Ὁ ἑνδέκατος Ἀσριὴλ, υἱοὶ αὐτοῦ καὶ ἀδελφοὶ αὐτοῦ δεκαδύο·
\vs{19}Ὁ δωδέκατος Ἀσαβία, υἱοὶ αὐτοῦ καὶ ἀδελφοὶ αὐτοῦ δεκαδύο·
\vs{20}Ὁ τρισκαιδέκατος Σουβαὴλ, υἱοὶ αὐτοῦ καὶ ἀδελφοὶ αὐτοῦ δεκαδύο·
\vs{21}Ὁ τεσσαρεσκαιδέκατος Ματθαθίας, υἱοὶ αὐτοῦ καὶ ἀδελφοὶ αὐτοῦ δεκαδύο·
\vs{22}Ὁ πεντεκαιδέκατος Ἰεριμὼθ, υἱοὶ αὐτοῦ καὶ ἀδελφοὶ αὐτοῦ δεκαδύο·
\vs{23}Ὁ ἑκκαιδέκατος Ἁνανίας, υἱοὶ αὐτοῦ καὶ ἀδελφοὶ αὐτοῦ δεκαδύο·
\vs{24}Ὁ ἑπτακαιδέκατος Ἰεσβασακὰ, υἱοὶ αὐτοῦ καὶ ἀδελφοὶ αὐτοῦ δεκαδύο·
\vs{25}Ὁ ὀκτωκαιδέκατος Ἀνανίας, υἱοὶ αὐτοῦ καὶ ἀδελφοὶ αὐτοῦ δεκαδύο·
\vs{26}Ὁ ἐννεακαιδέκατος Μαλλιθὶ, υἱοὶ αὐτοῦ καὶ ἀδελφοὶ αὐτοῦ δεκαδύο·
\vs{27}Ὁ εἰκοστὸς Ἑλιαθὰ, υἱοὶ αὐτοῦ καὶ ἀδελφοὶ αὐτοῦ δεκαδύο·
\vs{28}Ὁ εἰκοστὸς πρῶτος Ὠθηρὶ, υἱοὶ αὐτοῦ καὶ ἀδελφοὶ αὐτοῦ δεκαδύο·
\vs{29}Ὁ εἰκοστὸς δεύτερος Γοδολλαθὶ, υἱοὶ αὐτοῦ καὶ ἀδελφοὶ αὐτοῦ δεκαδύο·
\vs{30}Ὁ εἰκοστὸς τρίτος Μεαζὼθ, υἱοὶ αὐτοῦ καὶ ἀδελφοὶ αὐτοῦ δεκαδύο·
\vs{31}Ὁ εἰκοστὸς τέταρτος Ῥωμετθιέζερ, υἱοὶ αὐτοῦ καὶ ἀδελφοὶ αὐτοῦ δεκαδύο·

\ch{26}
Καὶ εἰς διαιρέσεις τῶν πυλῶν, υἱοὶ Κορεῒμ Μοσελλεμία ἐκ τῶν υἱῶν Ἀσάφ.
\vs{2}Καὶ τῷ Μοσελλαμίᾳ υἱὸς Ζαχαρίας ὁ πρωτότοκος, Ἰαδιὴλ ὁ δεύτερος, Ζαβαδία ὁ τρίτος, Ἰενουὴλ ὁ τέταρτος,
\vs{3}Ἰωλὰμ ὁ πέμπτος, Ἰωνάθαν ὁ ἕκτος, Ἐλιωναῒ ὁ ἕβδομος. Ἀβδεδὸμ ὁ ὄγδοος.
\vs{4}Καὶ τῷ Ἀβδεδὸμ υἱοὶ, Σαμαίας ὁ πρωτότοκος, Ἰωζαβὰθ ὁ δεύτερος, Ἰωὰθ ὁ τρίος, Σαχὰρ ὁ τέταρτος, Ναθαναὴλ ὁ πέμπτος,
\vs{5}Ἀμιὴλ ὁ ἕκτος, Ἰσσάχαρ ὁ ἕβδομος, Φελαθὶ ὁ ὄγδοος, ὅτι εὐλόγησεν αὐτὸν ὁ Θεός.
\vs{6}Καὶ τῷ Σαμαίᾳ υἱῶ αὑτοῦ ἐτέχθησαν υἱοὶ τοῦ πρωτοτόκου Ῥωσαὶ εἰς τὸν οἶκον τὸν πατρικὸν αὐτοῦ, ὅτι δυνατοὶ ἦσαν.
\vs{7}Υἱοὶ Σαμαῒ, Ὀθνὶ καὶ Ῥαφαὴλ, καὶ Ὠβὴδ, καὶ Ἐλζαβὰθ, καὶ Ἀχιοὺδ, υἱοὶ δυνατοὶ, Ἑλιοῦ, καὶ Σαβαχία, καὶ Ἰσβακώμ.
\vs{8}Πάντες ἀπὸ τῶν υἱῶν Ἀβδεδὸμ, αὐτοὶ καὶ οἱ υἱοὶ αὐτῶν καὶ οἱ ἀδελφοὶ αὐτῶν ποιοῦντες δυνατῶς ἐν τῇ ἐργασίᾳ, οἱ πάντες ἑξηκονταδύο τῷ Ἀβδεδόμ.

\vs{9}Καὶ τῷ Μοσελλεμίᾳ υἱοὶ καὶ ἀδελφοὶ δεκακαιοκτὼ δυνατοί.
\vs{10}Καὶ τῷ Ὀσᾷ τῶν υἱῶν Μεραρὶ υἱοὶ φυλάσσοντες τὴν ἀρχὴν, ὅτι οὐκ ἦν πρωτότοκος· καὶ ἐποίησεν αὐτὸν ὁ πατὴρ αὐτοῦ ἄρχοντα τῆς διαιρέσεως τῆς δευτέρας.
\vs{11}Χελκίας ὁ δεύτερος, Ταβλαὶ ὁ τρίτος, Ζαχαρίας ὁ τέταρτος· πάντες οὗτοι υἱοὶ καὶ ἀδελφοὶ τῷ Ὀσᾷ τρισκαίδεκα.

\vs{12}Τούτοις αἱ διαιρέσεις τῶν πυλῶν τοῖς ἄρχουσι τῶν δυνατῶν ἐφημερίαι, καθὼς οἱ ἀδελφοὶ αὐτῶν λειτουργεῖν ἐν οἴκῳ Κυρίου.
\vs{13}Καὶ ἔβαλον κλήρους κατὰ τὸν μικρὸν καὶ κατὰ τὸν μέγαν κατʼ οἴκους πατριῶν αὐτῶν εἰς πυλῶνα καὶ πυλῶνα.
\vs{14}Καὶ ἔπεσεν ὁ κλῆρος τῶν πρὸς ἀνατολὰς τῷ Σελεμίᾳ, καὶ Ζαχαρίᾳ· υἱοὶ Σωὰζ τῷ Μελχίᾳ ἔβαλον κλήρους, καὶ ἐξῆλθεν ὁ κλῆρος Βοῥῥᾶ.
\vs{15}Τῷ Ἀβδεδὸμ Νότον κατέναντι οἴκου Ἐσεφίμ.
\vs{16}Εἰς δεύτερον τῷ Ὀσᾷ πρὸς δυσμαῖς μετὰ τὴν πύλην παστοφορίου τῆς ἀναβάσεως· φυλακὴ κατέναντι φυλακῆς.
\vs{17}Πρὸς ἀνατολὰς ἓξ τὴν ἡμέραν· Βοῤῥᾶ τῆς ἡμέρας τέσσαρες· Νότον τῆς ἡμέρας τέσσαρες· καὶ εἰς τὸν Ἐσεφὶμ δύο
\vs{18}εἰς διαδεχομένους· καὶ τῷ Ὀσᾷ πρὸς δυσμαῖς μετὰ τὴν πύλην τοῦ παστοφορίου τρεῖς· φυλακὴ κατέναντι φυλακῆς τῆς ἀναβάσεως πρὸς ἀνατολὰς τῆς ἡμέρας ἓξ, καὶ τῷ Βοῤῥᾷ τέσσαρες, καὶ τῷ Νότῳ τέσσαρες, καὶ Ἐσεφὶμ δύο εἰς διαδεχομένους, καὶ πρὸς δυσμαῖς τέσσαρες, καὶ εἰς τὸν τρίβον δύο διαδεχομένους.
\vs{19}Αὗται αἱ διαιρέσεις τῶν πυλωρῶν τοῖς υἱοῖς τοῦ Κορὲ, καὶ τοῖς υἱοῖς Μεραρί.

\vs{20}Καὶ οἱ Λευῖται ἀδελφοὶ αὐτῶν ἐπὶ τῶν θησαυρῶν οἴκου Κυρίου, καὶ ἐπὶ τῶν θησαυρῶν τῶν καθηγιασμένων.
\vs{21}Υἱοὶ Λαδὰν οὗτοι, υἱοὶ τῷ Γηρσωνί· τῷ Λαδὰν ἄρχοντες πατριῶν, τῷ Λαδὰν, τῷ Γηρσωνὶ Ἰεϊήλ.
\vs{22}Υἱοὶ Ἰεϊὴλ Ζεθὸμ καὶ Ἰωὴλ, οἱ ἀδελφοὶ ἐπὶ τῶν θησαυρῶν οἴκου Κυρίου.
\vs{23}Τῷ Ἀμβρὰμ καὶ Ἰσσαὰρ, Χεβρὼν, καὶ Ὀζιήλ.
\vs{24}Καὶ Σουβαὴλ ὁ τοῦ Γηρσὰμ τοῦ Μωυσῆ ἐπὶ τῶν θησαυρῶν.
\vs{25}Καὶ τῷ ἀδελφῷ αὐτοῦ Ἐλιέζερ Ῥαβίας υἱὸς, καὶ Ἰωσίας, καὶ Ἰωρὰμ, καὶ Ζεχρὶ, καὶ Σαλωμώθ.
\vs{26}Αὐτὸς Σαλωμὼθ καὶ οἱ ἀδελφοὶ αὐτοῦ ἐπὶ πάντων τῶν θησαυρῶν τῶν ἁγίων, οὓς ἡγίασε Δαυὶδ ὁ βασιλεὺς καὶ οἱ ἄρχοντες τῶν πατριῶν, χιλίαρχοι καὶ ἑκατόνταρχοι καὶ ἀρχηγοὶ τῆς δυνάμεως,
\vs{27}ἃ ἔλαβεν ἐκ πόλεων καὶ ἐκ τῶν λαφύρων, καὶ ἡγίασεν ἀπʼ αὐτῶν τοῦ μὴ καθυστερῆσαι τὴν οἰκοδομὴν τοῦ οἴκου τοῦ Θεοῦ·
\vs{28}καὶ ἐπὶ πάντων τῶν ἁγίων τοῦ Θεοῦ Σαμουὴλ τοῦ προφήτου, καὶ Σαοὺλ τοῦ Κὶς, καὶ Ἀβεννὴρ τοῦ Νὴρ, καὶ Ἰωὰβ τοῦ Σαρουία, πᾶν ὃ ἡγίασαν διὰ χειρὸς Σαλωμὼθ καὶ τῶν ἀδελφῶν αὐτοῦ.

\vs{29}Τῷ Ἰσσααρὶ Χωνενία, καὶ υἱοὶ τῆς ἐργασίας τῆς ἔξω ἐπὶ τὸν Ἰσραὴλ τοῦ γραμματεύειν καὶ διακρίνειν.
\vs{30}Τῷ Χεβρωνὶ Ἀσαβίας καὶ οἱ ἀδελφοὶ αὐτοῦ υἱοὶ δυνατοὶ χίλιοι καὶ ἐπτακόσιοι ἐπὶ τῆς ἐπισκέψεως τοῦ Ἰσραὴλ πέραν τοῦ Ἰορδάνου πρὸς δυσμαῖς, εἰς πᾶσαν λειτουργίαν Κυρίου καὶ ἐργασίαν τοῦ βασιλέως.
\vs{31}Τοῦ Χεβρωνὶ Οὐρίας ὁ ἄρχων τῶν Χεβρωνὶ κατὰ γενέσεις αὐτῶν, κατὰ πατριὰς, ἐν τῷ τεσσαρακοστῷ ἔτει τῆς βασιλείας αὐτοῦ ἐπεσκέπησαν, καὶ εὑρέθη ἀνὴρ δυνατὸς ἐν αὐτοῖς ἐν Ἰαζὴρ τῆς Γαλααδίτιδος·
\vs{32}Καὶ οἱ ἀδελφοὶ αὐτοῦ υἱοὶ δυνατοὶ δισχίλιοι ἑπτακόσιοι οἱ ἄρχοντες τῶν πατριῶν, καὶ κατέστησεν αὐτοὺς Δαυὶδ ὁ βασιλεὺς ἐπὶ τοῦ Ῥουβηνὶ, καὶ Γαδδὶ, καὶ ἡμίσους φυλῆς Μανασσῆ εἰς πᾶν πρόσταγμα Κυρίου καὶ λόγον βασιλέως.

\ch{27}
Καὶ υἱοὶ Ἰσραὴλ κατὰ ἀριθμὸν αὐτῶν ἄρχοντες τῶν πατριῶν, χιλίαρχοι καὶ ἑκατόνταρχοι, καὶ γραμματεῖς οἱ λειτουργοῦντες τῷ βασιλεῖ καὶ εἰς πᾶν λόγον τοῦ βασιλέως κατὰ διαιρέσεις, πᾶν λόγον τοῦ εἰσπορευομένου καὶ ἐκπορευομένου μῆνα ἐκ μηνὸς, εἰς πάντας τοὺς μῆνας τοῦ ἐνιαυτοῦ, διαίρεσις μία εἴκοσι καὶ τέσσαρες χιλιάδες.

\vs{2}Καὶ ἐπὶ τῆς διαιρέσεως τῆς πρώτης τοῦ μηνὸς τοῦ πρώτου, Ἰσβοὰζ ὁ τοῦ Ζαβδιὴλ, ἐπὶ τῆς διαιρέσεως αὐτοῦ εἴκοσι καὶ τέσσαρες χιλιάδες·
\vs{3}Ἀπὸ τῶν υἱῶν Φαρὲς, ἄρχων πάντων τῶν ἀρχόντων τῆς δυνάμεως τοῦ μηνὸς τοῦ πρώτου.
\vs{4}Καὶ ἐπὶ τῆς διαιρέσεως τοῦ μηνὸς τοῦ δευτέρου Δωδία ὁ Ἐκχὼκ, καὶ ἐπὶ τῆς διαιρέσεως αὐτοῦ, καὶ Μακελλὼθ ὁ ἡγούμενος, καὶ ἐπὶ τῆς διαιρέσεως αὐτοῦ εἴκοσι καὶ τέσσαρες χιλιάδες ἄρχοντες δυνάμεως.
\vs{5}Ὁ τρίτος τὸν μῆνα τὸν τρίτον Βαναίας ὁ τοῦ Ἰωδαὲ ὁ ἱερεὺς ὁ ἄρχων, καὶ ἐπὶ τῆς διαιρέσεως αὐτοῦ εἴκοσι καὶ τέσσαρες χιλιάδες.
\vs{6}Αὐτὸς Βαναίας ὁ δυνατώτερος τῶν τριάκοντα καὶ ἐπὶ τῶν τριάκοντα· καὶ ἐπὶ τῆς διαιρέσεως αὐτοῦ Ζαβὰδ ὁ υἱὸς αὐτοῦ.
\vs{7}Ὁ τέταρτος εἰς τὸν μῆνα τὸν τέταρτον Ἀσαὴλ ὁ ἀδελφὸς Ἰωὰβ, καὶ Ζαβαδίας υἱὸς αὐτοῦ, καὶ οἱ ἀδελφοὶ, καὶ ἐπὶ τῆς διαιρέσεως αὐτοῦ εἴκοσι καὶ τέσσαρες χιλιάδες.
\vs{8}Ὁ πέμπτος τῷ μηνὶ τῷ πέμπτῳ ὁ ἡγούμενος Σαμαὼθ ὁ Ἰεσραὲ, καὶ ἐπὶ τῆς διαιρέσεως αὐτοῦ εἴκοσι καὶ τέσσαρες χιλιάδες.
\vs{9}Ὁ ἕκτος τῷ μηνὶ τῷ ἕκτῳ Ὁδουίας ὁ τοῦ Ἐκκῆς ὁ Θεκωΐτης, καὶ ἐπὶ τῆς διαιρέσεως αὐτοῦ εἴκοσι καὶ τέσσαρες χιλιάδες.
\vs{10}Ὁ ἕβδομος τῷ μηνὶ τῷ ἐβδόμῳ Χελλὴς ὁ ἐκ Φαλλοῦς ἀπὸ τῶν υἱῶν Ἐφραὶμ, καὶ ἑπὶ τῆς διαιρέσεως αὐτοῦ εἴκοσι καὶ τέσσαρες χιλιάδες.
\vs{11}Ὁ ὄγδοος τῷ μηνὶ τῷ ὀγδόῳ Σοβοχαῒ ὁ Οὐσαθὶ τῷ Ζαραῒ, καὶ ἐπὶ τῆς διαιρέσεως αὐτοῦ εἴκοσι καὶ τέσσαρες χιλιάδες.
\vs{12}Ὁ ἔννατος τῷ μηνὶ τῷ ἐννάτῳ Ἀβιέζερ ὁ ἐξ Ἀναθὼθ ὁ ἐκ γῆς Βενιαμὶν, καὶ ἐπὶ τῆς διαιρέσεως αὐτοῦ τέσσαρες καὶ εἴκοσι χιλιάδες.
\vs{13}Ὁ δέκατος τῷ μηνὶ τῷ δεκάτῳ Μεηρὰ ὁ ἐκ Νετωφαθὶ τῷ Ζαραῒ, καὶ ἐπὶ τῆς διαιρέσεως αὐτοῦ εἴκοσι καὶ τέσσαρες χιλιάδες.
\vs{14}Ὁ ἑνδέκατος τῷ μηνὶ τῷ ἑνδεκάτῳ Βαναίας ὁ ἐκ Φαραθὼν ἐκ τῶν υἱῶν Ἐφραὶμ, καὶ ἐπὶ τῆς διαιρέσεως αὐτοῦ εἴκοσι καὶ τέσσαρες χιλιάδες.
\vs{15}Ὁ δωδέκατος εἰς τὸν μῆνα τὸν δωδέκατον Χολδία ὁ ἐκ Νετωφαθὶ τῷ Γοθονιὴλ, καὶ ἐπὶ τῆς διαιρέσεως αὐτοῦ εἴκοσι καὶ τέσσαρες χιλιάδες.

\vs{16}Καὶ ἐπὶ τῶν φυλῶν Ἰσραὴλ, τῷ Ῥουβὴν ἡγούμενος Ἐλιέζερ ὁ τοῦ Ζεχρί· τῷ Συμεὼν, Σαφατίας ὁ τοῦ Μααχά.
\vs{17}Τῷ Λευὶ, Ἀσαβίας ὁ τοῦ Καμουήλ· τῷ Ἀαρὼν, Σαδώκ.
\vs{18}Τῷ Ἰούδα, Ἐλιὰβ τῶν ἀδελφῶν Δαυίδ· τῷ Ἰσσάχαρ, Ἀμβρὶ ὁ τοῦ Μιχαήλ.
\vs{19}Τῷ Ζαβουλὼν, Σαμαΐας ὁ τοῦ Ἀβδίου· τῷ Νεφθαλὶ, Ἰεριμὼθ ὁ τοῦ Ὀζιήλ.
\vs{20}Τῷ Ἐφραὶμ, Ὠσὴ ὁ τοῦ Ὀζίου· τῷ ἡμίσει φυλῆς Μανασσῆ, Ἰωὴλ υἱὸς Φαδαΐα.
\vs{21}Τῷ ἡμίσει φυλῆς Μανασσῆ τῷ ἐν γῇ Γαλαὰδ, Ἰαδαῒ ὁ τοῦ Ζαδαίου· τοῖς υἱοῖς Βενιαμὶν, Ἰασιὴλ ὁ τοῦ Ἀβεννήρ.
\vs{22}Τῷ Δὰν, Ἀζαριὴλ ὁ τοῦ Ἰρωάβ· οὗτοι πατριάρχαι τῶν φυλῶν Ἰσραήλ.

\vs{23}Καὶ οὐκ ἔλαβε Δαυὶδ τὸν ἀριθμὸν αὐτῶν ἀπὸ εἰκοσαετοῦς καὶ κάτω, ὅτι εἶπε Κύριος πληθῦναι τὸν Ἰσραὴλ ὡς τοὺς ἀστέρας τοῦ οὐρανοῦ.
\vs{24}Καὶ Ἰωὰβ ὁ τοῦ Σαρουία ἤρξατο ἀριθμεῖν ἐν τῷ λαῷ, καὶ οὐ συνετέλεσε· καὶ ἐγένετο ἐν τούτοις ὀργὴ ἐπὶ Ἰσραήλ· καὶ οὐ κατεχωρίσθη ὁ ἀριθμὸς ἐν βιβλίῳ λόγων τῶν ἡμερῶν τοῦ βασιλέως Δαυίδ.

\vs{25}Καὶ ἐπὶ τῶν θησαυρῶν τοῦ βασιλέως, Ἀσμὼθ ὁ τοῦ Ὀδιὴλ, καὶ ἐπὶ τῶν θησαυρῶν τῶν ἐν ἀγρῷ καὶ ἐν ταῖς κώμαις καὶ ἐν τοῖς ἐποικίοις καὶ ἐν τοῖς πύργοις, Ἰωνάθαν ὁ τοῦ Ὀζίου.
\vs{26}Καὶ ἐπὶ τῶν γεωργούντων τὴν γῆν τῶν ἐργαζομένων, Ἐσδρὶ ὁ τοῦ Χελούβ.
\vs{27}Καὶ ἐπὶ τῶν χωρίων, Σεμεῒ ὁ ἐκ Ῥαὴλ, καὶ ἐπὶ τῶν θησαυρῶν τῶν ἐν τοῖς χωρίοις τοῦ οἴνου, Ζαβδὶ ὁ τοῦ Σεφνί.
\vs{28}Καὶ ἐπὶ τῶν ἐλαιώνων, καὶ ἐπὶ τῶν συκαμίνων τῶν ἐν τῇ πεδινῇ, Βαλλανὰν ὁ Γεδωρίτης· ἐπὶ δὲ τῶν θησαυρῶν τοῦ ἐλαίου, Ἰωάς.
\vs{29}Καὶ ἐπὶ τῶν βοῶν τῶν νομάδων τῶν ἐν τῷ Σαρὼν, Σατραῒ ὁ Σαρωνίτης· καὶ ἐπὶ τῶν βοῶν τῶν ἐν τοῖς αὐλῶσι, Σωφὰτ, ὁ τοῦ Ἀδλί·
\vs{30}Ἐπὶ δὲ τῶν καμήλων, Ἀβίας ὁ Ἰσμαηλίτης· ἐπὶ δὲ τῶν ὄνων, Ἰαδίας ὁ ἐκ Μεραθών.
\vs{31}Καὶ ἐπὶ τῶν προβάτων, Ἰαζὶζ ὁ Ἀγαρίτης· πάντες οὗτοι προστάται ὑπαρχόντων Δαυὶδ τοῦ βασιλέως.

\vs{32}Καὶ Ἰωνάθαν ὁ πατράδελφος Δαυὶδ σύμβουλος, ἄνθρωπος συνετός· καὶ Ἰεὴλ ὁ τοῦ Ἀχαμὶ μετὰ τῶν υἱῶν τοῦ βασιλέως.
\vs{33}Ἀχιτόφελ σύμβουλος τοῦ βασιλέως, καὶ Χουσὶ ὁ πρῶτος φίλος τοῦ βασιλέως·
\vs{34}Καὶ μετὰ τοῦτον Ἀχιτόφελ ἐχόμενος Ἰωδαὲ ὁ τοῦ Βαναίου, καὶ Ἀβιάθαρ· καὶ Ἰωὰβ ἀρχιστράτηγος τοῦ βασιλέως.

\ch{28}
Καὶ ἐξεκκλησίασε Δαυὶδ πάντας τοὺς ἄρχοντας Ἰσραὴλ, ἄρχοντας τῶν κριτῶν, καὶ πάντας τοὺς ἄρχοντας τῶν ἐφημεριῶν τῶν περὶ τὸ σῶμα τοῦ βασιλέως, καὶ ἄρχοντας τῶν χιλιάδων καὶ τῶν ἑκατοντάδων, καὶ τοὺς γαζοφύλακας, καὶ τοὺς ἐπὶ τῶν ὑπαρχόντων αὐτοῦ, καὶ πάσης τῆς κτήσεως τοῦ βασιλέως, καὶ τῶν υἱῶν αὐτοῦ, σὺν τοῖς εὐνούχοις, καὶ τοὺς δυνάστας, καὶ τοὺς μαχητὰς τῆς στρατιᾶς ἐν Ἱερουσαλήμ.

\vs{2}Καὶ ἔστη Δαυὶδ ἐν μέσῳ τῆς ἐκκλησίας, καὶ εἶπεν, ἀκούσατέ μου ἀδελφοί μου, καὶ λαός μου· ἐμοὶ ἐγένετο ἐπὶ καρδίαν οἰκοδομῆσαι οἶκον ἀναπαύσεως τῆς κιβωτοῦ διαθήκης Κυρίου, καὶ στάσιν ποδῶν Κυρίου ἡμῶν, καὶ ἡτοίμασα τὰ εἰς τὴν κατασκήνωσιν ἐπιτήδεια.
\vs{3}Καὶ ὁ Θεὸς εἶπεν, οὐκ οἰκοδομήσεις ἐμοὶ οἶκον τοῦ ἐπονομάσαι τὸ ὄνομά μου ἐπʼ αὐτῷ, ὅτι ἄνθρωπος πολεμιστὴς εἶ σὺ, καὶ αἷμα ἐξέχεας.
\vs{4}Καὶ ἐξελέξατο Κύριος ὁ Θεὸς Ἰσραὴλ ἐν ἐμοὶ ἀπὸ παντὸς οἴκου πατρός μου εἶναι βασιλέα ἐπὶ Ἰσραὴλ εἰς τὸν αἰῶνα, καὶ ἐν Ἰούδα ᾑρέτικε τὸ βασίλειον, καὶ ἐξ οἴκου Ἰούδα τὸν οἶκον τοῦ πατρός μου· καὶ ἐν τοῖς υἱοῖς τοῦ πατρός μου, ἐν ἐμοὶ ἠθέλησε τοῦ γενέσθαι με εἰς βασιλέα ἐπὶ παντὶ Ἰσραήλ.
\vs{5}Καὶ ἀπὸ πάντων τῶν υἱῶν μου, ὅτι πολλοὺς υἱοὺς ἔδωκέ μοι Κύριος, ἐξελέξατο ἐν Σαλωμὼν τῷ υἱῷ μου καθίσαι αὐτὸν ἐπὶ θρόνου βασιλείας Κυρίου ἐπὶ τὸν Ἰσραήλ.
\vs{6}Καὶ εἶπέ μοι ὁ Θεὸς, Σαλωμὼν ὁ υἱός σου οἰκοδομήσει τὸν οἶκόν μου καὶ τὴν αὐλήν μου, ὅτι ᾑρέτικα ἐν αὐτῷ εἶναί μου υἱὸν, κᾀγὼ ἔσομαι αὐτῷ εἰς πατέρα.
\vs{7}Καὶ κατορθώσω τὴν βασιλείαν αὐτοῦ ἕως αἰῶνος, ἐὰν ἰσχύσῃ τοῦ φυλάξασθαι τὰς ἐντολάς μου, καὶ τὰ κρίματά μου, ὡς ἡ ἡμέρα αὕτη.
\vs{8}Καὶ νῦν κατὰ πρόσωπον πάσης ἐκκλησίας Κυρίου, καὶ ἐν ὠσὶ Θεοῦ ἡμῶν, φυλάξασθε καὶ ζητήσατε πάσας τὰς ἐντολὰς Κυρίου τοῦ Θεοῦ ἡμῶν, ἵνα κληρονομήσητε τὴν γῆν τὴν ἀγαθὴν, καὶ κατακληρονομήσητε τοῖς υἱοῖς ὑμῶν μεθʼ ὑμᾶς ἕως αἰῶνος.

\vs{9}Καὶ νῦν Σαλωμὼν υἱὲ, γνῶθι τὸν Θεὸν τῶν πατέρων σου, καὶ δούλευε αὐτῷ ἐν καρδίᾳ τελείᾳ· καὶ ψυχῇ θελούσῃ, ὅτι πάσας καρδίας ἐτάζει Κύριος, καὶ πᾶν ἐνθύμημα γινώσκει· ἐὰν ζητήσῃς αὐτὸν, εὑρεθήσεταί σοι, καὶ ἐὰν καταλείψῃς αὐτὸν, καταλείψει σε εἰς τέλος.
\vs{10}Ἴδε νῦν, ὅτι Κύριος ᾑρέτικέ σε οἰκοδομῆσαι αὐτῷ οἶκον εἰς ἁγίασμα, ἴσχυε καὶ ποίει.

\vs{11}Καὶ ἔδωκε Δαυὶδ Σαλωμὼν τῷ υἱῷ αὐτοῦ τὸ παράδειγμα τοῦ ναοῦ καὶ τῶν οἴκων αὐτοῦ, καὶ τῶν ζακχῶν αὐτοῦ, καὶ τῶν ὑπερῴων, καὶ τῶν ἀποθηκῶν τῶν ἐσωτέρων, καὶ τοῦ οἴκου τοῦ ἐξιλασμοῦ,
\vs{12}καὶ τὸ παράδειγμα ὃ εἶχεν ἐν πνεύματι αὐτοῦ, τῶν αὐλῶν οἴκου Κυρίου, καὶ πάντων τῶν παστοφορίων τῶν κύκλῳ τῶν εἰς τὰς ἀποθήκας οἴκου Κυρίου, καὶ τῶν ἀποθηκῶν τῶν ἁγίων,
\vs{13}καὶ τῶν καταλυμάτων, καὶ τῶν ἐφημεριῶν τῶν ἱερέων καὶ τῶν Λευιτῶν εἰς πᾶσαν ἐργασίαν λειτουργίας οἴκου Κυρίου, καὶ τῶν ἀποθηκῶν τῶν λειτουργησίμων σκευῶν τῆς λατρείας οἴκου Κυρίου.
\vs{14}Καὶ τὸν σταθμὸν τῆς ὁλκῆς αὐτῶν τῶν τε χρυσῶν καὶ ἀργυρῶν λυχνιῶν τὴν ὁλκὴν ἔδωκεν αὐτῷ,
\vs{15}καὶ τῶν λύχνων.
\vs{16}Ἔδωκεν αὐτῷ ὁμοίως τὸν σταθμὸν τῶν τραπεζῶν τῆς προθέσεως, ἑκάστης τραπέζης χρυσῆς, καὶ ὡσαύτως τῶν ἀργυρῶν,
\vs{17}καὶ τῶν κρεαγρῶν καὶ σπονδείων καὶ τῶν φιαλῶν τῶν χρυσῶν· καὶ τὸν σταθμὸν τῶν χρυσῶν καὶ τῶν ἀργυρῶν, καὶ θυΐσκων κεφουρὲ, ἑκάστου σταθμοῦ.
\vs{18}Καὶ τῶν τοῦ θυσιαστηρίου τῶν θυμιαμάτων ἐκ χρυσίου δοκίμου σταθμὸν ὑπέδειξεν αὐτῷ, καὶ τὸ παράδειγμα τοῦ ἅρματος τῶν χερουβὶμ τῶν διαπεπετασμένων ταῖς πτέρυξι, καὶ σκιαζόντων ἐπὶ τῆς κιβωτοῦ διαθήκης Κυρίου·
\vs{19}πάντα ἐν γραφῇ χειρὸς Κυρίου ἔδωκε Δαυὶδ Σαλωμὼν, κατὰ τὴν περιγενηθεῖσαν αὐτῷ σύνεσιν τῆς κατεργασίας τοῦ παραδείγματος.

\vs{20}Καὶ εἶπε Δαυὶδ Σαλωμὼν τῷ υἱῷ αὐτοῦ, ἴσχυε καὶ ἀνδρίζου καὶ ποίει, μὴ φοβοῦ μηδὲ πτοηθῇς, ὅτι Κύριος ὁ Θεός μου μετὰ σοῦ, οὐκ ἀνήσει σε, καὶ οὐ μὴ ἐγκαταλίπῃ ἕως τοῦ συντελέσαι σε πᾶσαν ἐργασίαν λειτουργίας οἴκου Κυρίου· καὶ ἰδοὺ τὸ παράδειγμα τοῦ ναοῦ καὶ τοῦ οἴκου αὐτοῦ, καὶ ζακχῶ αὐτοῦ, καὶ τὰ ὑπερῷα καὶ τὰς ἀποθήκας τὰς ἐσωτέρας, καὶ τὸν οἶκον τοῦ ἱλασμοῦ, καὶ τὸ παράδειγμα οἴκου Κυρίου.
\vs{21}Καὶ ἰδοὺ αἱ ἐφημερίαι τῶν ἱερέων καὶ τῶν Λευιτῶν εἰς πᾶσαν λειτουργίαν οἴκου Κυρίου, καὶ μετὰ σοῦ ἐν πάσῃ πραγματείᾳ, καὶ πᾶς πρόθυμος ἐν σοφίᾳ κατὰ πᾶσαν τέχνην, καὶ οἱ ἄρχοντες καὶ πᾶς ὁ λαὸς εἰς πάντας τοὺς λόγους σου.

\ch{29}
Καὶ εἶπε Δαυὶδ ὁ βασιλεὺς πάσῃ τῇ ἐκκλησίᾳ, Σαλωμὼν ὁ υἱός μου, εἰς ὃν ᾑρέτικεν ἐν αὐτῷ Κύριος, νέος καὶ ἁπαλὸς, καὶ τὸ ἔργον μέγα, ὅτι οὐκ ἀνθρώπῳ, ἀλλʼ ἢ Κυρίῳ Θεῷ.
\vs{2}Κατὰ πᾶσαν τὴν δύναμιν ἡτοίμακα εἰς οἶκον Θεοῦ μου χρυσὶον, ἀργύριον, χαλκόν, σίδηρον, ξύλα, λίθους σοὰμ, καὶ πληρώσεως λίθους πολυτελεῖς καὶ ποικίλους, καὶ πάντα λίθον τίμιον, καὶ Πάριον πολύν.
\vs{3}Καὶ ἔτι ἐν τῷ εὐδοκῆσαί με ἐν οἴκῳ Θεοῦ μου, ἔστι μοι ὃ περιπεποίημαι χρυσίον καὶ ἀργύριον, καὶ ἰδοὺ δέδωκα εἰς οἶκον Θεοῦ μου εἰς ὕψος, ἐκτὸς ὧν ἡτοίμακα εἰς τὸν οἶκον τῶν ἁγίων,
\vs{4}τρισχίλια τάλαντα χρυσίου τοῦ ἐκ Σουφὶρ, καὶ ἑπτακισχίλια τάλαντα ἀργυρίου δοκίμου, ἐξαλειφῆναι ἐν αὐτοῖς τοὺς τοίχους τοῦ ἱεροῦ,
\vs{5}εἰς τὸ χρυσίον τῷ χρυσίῳ, καὶ εἰς τὸ ἀργύριον τῷ ἀργυρίῳ, καὶ εἰς πᾶν ἔργον διὰ χειρὸς τῶν τεχνιτῶν· καὶ τίς ὁ προθυμούμενος πληρῶσαι τὰς χεῖρας αὐτοῦ σήμερον Κυρίῳ;

\vs{6}Καὶ προεθυμήθησαν ἄρχοντες πατριῶν, καὶ οἱ ἄρχοντες τῶν υἱῶν Ἰσραὴλ, καὶ οἱ χιλίαρχοι καὶ οἱ ἑκατόνταρχοι, καὶ οἱ προστάται τῶν ἔργων, καὶ οἱ οἰκοδόμοι τοῦ βασιλέως.
\vs{7}Καὶ ἔδωκαν εἰς τὰ ἔργα τοῦ οἴκου Κυρίου χρυσίου τάλαντα πεντακισχίλια, καὶ χρυσοῦς μυρίους, καὶ ἀργυρίου ταλάντων δέκα χιλιάδας, καὶ χαλκοῦ τάλαντα μύρια ὀκτακισχίλια, καὶ σιδήρου ταλάντων χιλιάδας ἑκατόν.
\vs{8}Καὶ οἷς εὑρέθη παρʼ αὐτοῖς λίθος, ἔδωκαν εἰς τὰς ἀποθήκας οἴκου Κυρίου διὰ χειρὸς Ἰεϊὴλ τοῦ Γεδσωνί.
\vs{9}Καὶ εὐφράνθη ὁ λαὸς ὑπὲρ τοῦ προθυμηθῆναι, ὅτι ἐν καρδίᾳ πλήρει προεθυμήθησαν τῷ Κυρίῳ· καὶ Δαυὶδ ὁ βασιλεὺς εὐφράνθη μεγάλως.

\vs{10}Καὶ εὐλόγησεν ὁ βασιλεὺς Δαυὶδ τὸν Κύριον ἐνώπιον τῆς ἐκκλησίας, λέγων,

\vs{11}Εὐλογητὸς εἶ Κύριε ὁ Θεὸς Ἰσραὴλ, ὁ πατὴρ ἡμῶν, ἀπὸ τοῦ αἰῶνος καὶ ἕως τοῦ αἰῶνος. Σοὶ Κύριε ἡ μεγαλωσύνη, καὶ ἡ δύναμις, καὶ τὸ καύχημα, καὶ ἡ νίκη, καὶ ἡ ἰσχὺς, ὅτι σὺ πάντων τῶν ἐν τῷ οὐρανῷ καὶ ἐπὶ τῆς γῆς δεσπόζεις· ἀπὸ προσώπου σου ταράσσεται πᾶς βασιλεὺς, καὶ ἔθνος.
\vs{12}Παρὰ σοῦ ὁ πλοῦτος καὶ ἡ δόξα, σὺ πάντων ἄρχεις Κύριε, ὁ ἄρχων πάσης ἀρχῆς, καὶ ἐν χειρί σου ἰσχὺς καὶ δυναστεία, καὶ ἐν χειρί σου παντοκράτωρ μεγαλῦναι καὶ κατισχύσαι τὰ πάντα.
\vs{13}Καὶ νῦν, Κύριε ἐξομολογούμεθά σοι, καὶ αἰνοῦμεν τὸ ὄνομα τῆς καυχήσεώς σου.
\vs{14}Καὶ τίς εἰμι ἐγὼ, καὶ τίς ὁ λαός μου ὅτι ἰσχύσαμεν προθυμηθῆναί σοι κατὰ ταῦτα; ὅτι σὰ τὰ πάντα, καὶ ἐκ τῶν σῶν δεδώκαμέν σοι.
\vs{15}Ὅτι πάροικοι ἐσμὲν ἐναντίον σου, καὶ παροικοῦντες ὡς πάντες οἱ πατέρες ἡμων· ὡς σκιὰ αἱ ἡμέραι ἡμῶν ἐπὶ γῆς, καὶ οὐκ ἔστιν ὑπομονή.
\vs{16}Κύριε ὁ Θεὸς ἡμῶν, πρὸς πᾶν τὸ πλῆθος τοῦτο ὃ ἡτοίμακα οἰκοδομηθῆναι οἶκον τῷ ὀνόματι τῷ ἁγίῳ σου, ἐκ χειρός σου ἐστὶ, καὶ σοὶ τὰ πάντα.
\vs{17}Καὶ ἔγνων, Κύριε, ὅτι σὺ εἶ ὁ ἐτάζων καρδίας, καὶ δικαιοσύνην ἀγαπᾷς· ἐν ἁπλότητι καρδίας προεθυμήθην ταῦτα πάντα, καὶ νῦν τὸν λαόν σου τὸν εὑρεθέντα ὧδε εἶδον ἐν εὐφροσύνῃ προθυμηθέντα σοι.
\vs{18}Κύριε ὁ Θεὸς Ἁβραὰμ, καὶ Ἰσαὰκ, καὶ Ἰσραὴλ, τῶν πατέρων ἡμῶν, φύλαξον ταῦτα ἐν διανοίᾳ καρδίας λαοῦ σου εἰς τὸν αἰῶνα, καὶ κατεύθυνον τὰς καρδίας αὐτῶν πρὸς σέ.
\vs{19}Καὶ Σαλωμὼν τῷ υἱῷ μου δὸς καρδίαν ἀγαθὴν ποιεῖν τὰς ἐντολάς σου, καὶ τὰ μαρτύριά σου, καὶ τὰ προστάγματά σου, καὶ τοῦ ἐπὶ τέλος ἀγαγεῖν τὴν κατασκευὴν τοῦ οἴκου σου.

\vs{20}Καὶ εἶπε Δαυὶδ πάσῃ τῇ ἐκκλησίᾳ, εὐλογήσατε Κύριον τὸν Θεὸν ἡμῶν· καὶ εὐλόγησε πᾶσα ἡ ἐκκλησία Κύριον τὸν Θεὸν τῶν πατέρων αὐτῶν· καὶ κάμψαντες τὰ γόνατα προσεκύνησαν Κυρίῳ, καὶ τῷ βασιλεῖ.
\vs{21}Καὶ ἔθυσε Δαυὶδ τῷ Κυρίῳ θυσίας, καὶ ἀνήνεγκεν ὁλοκαυτώματα τῷ Θεῷ τῇ ἐπαύριον τῆς πρώτης ἡμέρας μόσχους χιλίους, κριοὺς χιλίους, ἄρνας χιλίους, καὶ τὰς σπονδὰς αὐτῶν, καὶ θυσίας εἰς πλῆθος παντὶ τῷ Ἰσραήλ.
\vs{22}Καὶ ἔφαγον καὶ ἔπιον ἐναντίον τοῦ Κυρίου ἐν ἐκείνῃ τῇ ἡμέρᾳ μετὰ χαρᾶς· καὶ ἐβασίλευσαν ἐκ δευτέρου τὸν Σαλωμὼν υἱὸν Δαυὶδ, καὶ ἔχρισαν αὐτὸν τῷ Κυρίῳ εἰς βασιλέα, καὶ Σαδὼκ εἰς ἱερωσύνην.

\vs{23}Καὶ ἐκάθισε Σαλωμὼν ἐπὶ θρόνου Δαυὶδ τοῦ πατρὸς αὐτοῦ, καὶ εὐδοκήθη, καὶ ὑπήκουσαν αὐτοῦ πᾶς Ἰσραήλ.
\vs{24}Οἱ ἄρχοντες, καὶ οἱ δυνάσται, καὶ πάντες υἱοὶ Δαυὶδ τοῦ βασιλέως τοῦ πατρὸς αὐτοῦ, ὑπετάγησαν αὐτῷ.
\vs{25}Καὶ ἐμεγάλυνε Κύριος τὸν Σαλωμὼν ἐπάνωθεν παντὸς Ἰσραὴλ, καὶ ἔδωκεν αὐτῷ δόξαν βασιλέως, ὃ οὐκ ἐγένετο ἐπὶ παντὸς βασιλέως ἔμπροσθεν αὐτοῦ.

\vs{26}Καὶ Δαυὶδ υἱὸς Ἰεσσαὶ ἐβασίλευσεν ἐπὶ Ἰσραὴλ
\vs{27}ἔτη τεσσράκοντα, ἐν Χεβρὼν ἔτη ἑπτὰ, καὶ ἐν Ἱερουσαλὴμ ἔτη τριακοντατρία.
\vs{28}Καὶ ἐτελεύτησεν ἐν γήρᾳ καλῷ, πλήρης ἡμερῶν, πλούτῳ καὶ δόξῃ, καὶ ἐβασίλευσε Σαλωμὼν υἱὸς αὐτοῦ ἀντʼ αὐτοῦ.
\vs{29}Οἱ δὲ λοιποὶ λόγοι τοῦ βασιλέως Δαυὶδ οἱ πρότεροι καὶ οἱ ὕστεροι γεγραμμένοι εἰσὶν ἐν λόγοις Σαμουὴλ τοῦ βλέποντος, καὶ ἐπὶ λόγων Νάθαν τοῦ προφήτου, καὶ ἐπὶ λόγων Γὰδ τοῦ βλέποντος,
\vs{30}περὶ πάσης τῆς βασιλείας αὐτοῦ, καὶ τῆς δυναστείας αὐτοῦ, καὶ οἱ καιροὶ οἳ ἐγένοντο ἐπʼ αὐτῷ, καὶ ἐπὶ τὸν Ἰσραὴλ, καὶ ἐπὶ πάσας βασιλείας τῆς γῆς.


\def\book{ΠΑΡΑΛΕΙΠΟΜΕΝΩΝ Β}
\biblebook{ΠΑΡΑΛΕΙΠΟΜΕΝΩΝ Β}


\lettrine[lines=2, loversize=0.2, nindent=0em, findent=.25em]{\textcolor{bookheadingcolor}{Κ}}{ΑΙ} ἐνίσχυσε Σαλωμὼν υἱὸς Δαυὶδ ἐπὶ τὴν βασιλείαν αὐτοῦ, καὶ Κύριος ὁ Θεὸς αὐτοῦ μετʼ αὐτοῦ, καὶ ἐμεγάλυνεν αὐτὸν εἰς ὕψος.
\vs{2}Καὶ εἶπε Σαλωμὼν πρὸς πάντα Ἰσραὴλ, τοῖς χιλιάρχοις, καὶ τοῖς ἑκατοντάρχοις, καὶ τοῖς κριταῖς, καὶ πᾶσι τοῖς ἄρχουσιν ἐναντίον Ἰσραὴλ τοῖς ἄρχουσι τῶν πατριῶν.
\vs{3}καὶ ἐπορεύθη Σαλωμὼν καὶ πᾶσα ἡ ἐκκλησία εἰς τὴν ὑψηλὴν τὴν ἐν Γαβαὼν, οὗ ἐκεῖ ἦν ἡ σκηνὴ τοῦ μαρτυρίου τοῦ Θεοῦ, ἣν ἐποίησε Μωυσῆς παῖς Κυρίου ἐν τῇ ἐρήμῳ.
\vs{4}Ἀλλὰ κιβωτὸν τοῦ Θεοῦ ἀνήνεγκε Δαυὶδ ἐκ πόλεως Καριαθιαρὶμ, ὅτι ἡτοίμασεν αὐτῇ Δαυὶδ, ὅτι ἔπηξεν αὐτῇ σκηνὴν ἐν Ἱερουσαλήμ.
\vs{5}Καὶ τὸ θυσιαστήριον τὸ χαλκοῦν ὃ ἐποίησε Βεσελεὴλ υἱὸς Οὐρίου υἱοῦ Ὢρ, ἐκεῖ ἦν ἔναντι τῆς σκηνῆς Κυρίου· καὶ ἐξεζήτησεν αὐτὸ Σαλωμὼν καὶ ἡ ἐκκλησία,
\vs{6}καὶ ἤνεγκε Σαλωμὼν ἐκεῖ ἐπὶ τὸ θυσιαστήριον τὸ χαλκοῦν ἐνώπιον Κυρίου τὸ ἐν τῇ σκηνῇ, καὶ ἤνεγκεν ἐπʼ αὐτῷ ὁλοκαύτωσιν χιλίαν.

\vs{7}Ἐν τῇ νυκτὶ ἐκείνῃ ὤφθη Θεὸς τῷ Σαλωμὼν, καὶ εἶπεν αὐτῷ, αἴτησαι τί σοι δῶ.
\vs{8}Καὶ εἶπε Σαλωμὼν πρὸς τὸν Θεὸν, σὺ ἐποίησας μετὰ Δαυὶδ τοῦ πατρός μου ἔλεος μέγα, καὶ ἐβασίλευσάς με ἀντʼ αὐτοῦ.
\vs{9}Καὶ νῦν, Κύριε ὁ Θεὸς, πιστωθήτω δὴ τὸ ὄνομά σου ἐπὶ Δαυὶδ τὸν πατέρα μου, ὅτι σὺ ἐβασίλευσάς με ἐπὶ λαὸν πολὺν, ὡς ὁ χοῖς τῆς γῆς.
\vs{10}Νῦν σοφίαν καὶ σύνεσιν δός μοι, καὶ ἐξελεύσομαι ἐνώπιον τοῦ λαοῦ τούτου καὶ εἰσελεύσομαι, ὅτι τίς κρινεῖ τὸν λαόν σου τὸν μέγαν τοῦτον;

\vs{11}Καὶ εἶπεν ὁ Θεὸς πρὸς Σαλωμὼν, ἀνθʼ ὧν ἐγένετο τοῦτο ἐν τῇ καρδίᾳ σου, καὶ οὐκ ᾐτήσω πλοῦτον χρημάτων, οὐδὲ δόξαν, οὐδὲ τὴν ψυχὴν τῶν ὑπεναντίων, καὶ ἡμέρας πολλὰς οὐκ ᾐτήσω, καὶ ᾔτησας σεαυτῷ σοφίαν καὶ σύνεσιν, ὅπως κρίνῃς τὸν λαόν μου, ἐφʼ ὃν ἐβασίλευσά σε ἐπʼ αὐτὸν,
\vs{12}τὴν σοφίαν καὶ τὴν σύνεσιν δίδωμί σοι, καὶ πλοῦτον καὶ χρήματα καὶ δόξαν δώσω σοι, ὡς οὐκ ἐγενήθη ὅμοιός σοι ἐν τοῖς βασιλεῦσι τοῖς ἔμπροσθέν σου, καὶ μετὰ σὲ οὐκ ἔσται οὕτως.

\vs{13}Καὶ ἦλθε Σαλωμὼν ἐκ βαμὰ τῆς ἐν Γαβαὼν εἰς Ἱερουσαλὴμ πρὸ προσώπου τῆς σκηνῆς τοῦ μαρτυρίου, καὶ ἐβασίλευσεν ἐπὶ Ἰσραήλ.

\vs{14}Καὶ συνήγαγε Σαλωμὼν ἅρματα καὶ ἱππεῖς, καὶ ἐγένοντο αὐτῷ χιλία καὶ τετρακόσια ἅρματα, καὶ δώδεκα χιλιάδες ἱππέων· καὶ κατέλιπεν αὐτὰ ἐν πόλεσι τῶν ἁρμάτων, καὶ ὁ λαὸς μετὰ τοῦ βασιλέως ἐν Ἱερουσαλήμ.
\vs{15}Καὶ ἔθηκεν ὁ βασιλεὺς τὸ ἀργύριον καὶ τὸ χρυσίον ἐν Ἱερουσαλὴμ ὡς λίθους, καὶ τὰς κέδρους ἐν τῇ Ἰουδαίᾳ ὡς συκαμίνους τὰς ἐν τῇ πεδινῇ εἰς πλῆθος.
\vs{16}Καὶ ἡ ἔξοδος τῶν ἵππων Σαλωμὼν ἐξ Αἰγύπτου, καὶ ἡ τιμὴ τῶν ἐμπόρων τοῦ βασιλέως πορεύεσθαι, καὶ ἠγόραζον,
\vs{17}καὶ ἐνέβαινον καὶ ἐξῆγον ἐξ Αἰγύπτου ἅρμα ἓν ἑξακοσίων ἀργυρίου, καὶ ἵππον πεντήκοντα καὶ ἑκατὸν ἀργυρίου· καὶ οὕτω πᾶσι τοῖς βασιλεῦσι τῶν Χετταίων καὶ τοῖς βασιλεῦσι Συρίας ἐν χερσὶν αὐτῶν ἔφερον.

\vs{18}Καὶ εἶπε Σαλωμὼν τοῦ οἰκοδομῆσαι οἶκον τῷ ὀνόματι Κυρίου, καὶ οἶκον τῷ βασιλείᾳ αὐτοῦ.

\ch{2}Καὶ συνήγαγε Σαλωμὼν ἑβδομήκοντα χιλιάδας ἀνδρῶν νωτοφόρων, καὶ ὀγδοήκοντα χιλιάδας λατόμων ἐν τῷ ὄρει, καὶ οἱ ἐπιστάται ἐπʼ αὐτῶν τρισχίλιοι ἑξακόσιοι.

\vs{2}Καὶ ἀπέστειλε Σαλωμὼν πρὸς Χιρὰμ βασιλέα Τύρου, λέγων, ὡς ἐποίησας μετὰ Δαυὶδ τοῦ πατρός μου, καὶ ἀπέστειλας αὐτῷ κέδρους τοῦ οἰκοδομῆσαι ἑαυτῷ οἶκον κατοικῆσαι ἐν αὐτῷ,
\vs{3}καὶ ἰδοὺ ἐγὼ ὁ υἱὸς αὐτοῦ οἰκοδομῶ οἶκον τῷ ὀνόματι Κυρίου Θεοῦ μου, ἁγιάσαι αὐτὸν αὐτῷ τοῦ θυμιᾷν ἀπέναντι αὐτοῦ θυμίαμα καὶ πρόθεσιν διαπαντὸς, καὶ τοῦ ἀναφέρειν ὁλοκαυτώματα διαπαντὸς τοπρωῒ καὶ τοδείλης, καὶ ἐν τοῖς σαββάτοις, καὶ ἐν ταῖς νουμηνίαις, καὶ ἐν ταῖς ἑορταῖς τοῦ Κυρίου Θεοῦ ἡμῶν· εἰς τὸν αἰῶνα τοῦτο ἐπὶ τὸν Ἰσραήλ.
\vs{4}Καὶ ὁ οἶκος ὃν ἐγὼ οἰκοδομῶ, μέγας, ὅτι μέγας Κύριος ὁ Θεὸς ἡμῶν παρὰ πάντας τοὺς θεούς.
\vs{5}Καὶ τίς ἰσχύσει οἰκοδομῆσαι αὐτῷ οἶκον; ὅτι ὁ οὐρανὸς, καὶ ὁ οὐρανὸς τοῦ οὐρανοῦ οὐ φέρουσι τὴς δόξαν αὐτοῦ· καὶ τίς ἐγὼ οἰκοδομῶν αὐτῷ οἶκον; ὅτι ἀλλʼ ἢ τοῦ θυμιᾷν κατέναντι αὐτοῦ.
\vs{6}Καὶ νῦν ἀπόστειλόν μοι ἄνδρα σοφὸν καὶ εἰδότα τοῦ ποιῆσαι ἐν τῷ χρυσίῳ, καὶ ἐν τῷ ἀργυρίῳ, καὶ ἐν τῷ χαλκῷ, καὶ ἐν τῷ σιδήρῳ, καὶ ἐν τῇ πορφύρᾳ, καὶ ἐν τῷ κοκκίνῳ, καὶ ἐν τῇ ὑακίνθῳ, καὶ ἐπιστάμενον γλύψαι γλυφὴν μετὰ τῶν σοφῶν τῶν μετʼ ἐμοῦ ἐν Ἰούδα καὶ ἐν Ἱερουσαλὴμ, ἃ ἡτοίμασε Δαυὶδ ὁ πατήρ μου.
\vs{7}Καὶ ἀπόστειλόν μοι ξύλα κέδρινα καὶ ἀρκεύθινα καὶ πεύκινα ἐκ τοῦ Λιβάνου, ὅτι ἐγὼ οἶδα ὡς οἱ δοῦλοί σου οἴδασι κόπτειν ξύλα ἐκ τοῦ Λιβάνου· καὶ ἰδοὺ οἱ παῖδές σου μετὰ τῶν παίδων μου
\vs{8}πορεύσονται ἑτοιμάσαι μοι ξύλα εἰς πλῆθος, ὅτι ὁ οἶκος ὃν ἐγὼ οἰκοδομῶ μέγας καὶ ἔνδοξος.
\vs{9}Καὶ ἰδοὺ τοῖς ἐργαζομένοις τοῖς κόπτουσι ξύλα, εἰς βρώματα δέδωκα σῖτον εἰς δόματα τοῖς παισί σου κόρων πυροῦ εἴκοσι χιλιάδας, καὶ κριθῶν κόρων εἴκοσι χιλιάδας, καὶ οἴνου μέτρων εἴκοσι χιλιάδας, καὶ ἐλαίου μέτρων εἴκοσι χιλιάδας.

\vs{10}Καὶ εἶπε Χιρὰμ βασιλεὺς Τύρου ἐν γραφῇ, καὶ ἀνέστειλε πρὸς Σαλωμὼν, λέγων, ἐν τῷ ἀγαπῆσαι Κύριον τὸν λαὸν αὐτοῦ, ἔδωκέ σε ἐπʼ αὐτοὺς βασιλέα.
\vs{11}Καὶ εἶπε Χιρὰμ, εὐλογητὸς Κύριος ὁ Θεὸς Ἰσραὴλ, ὃς ἐποίησε τὸν οὐρανὸν καὶ τὴν γῆν, ὃς ἔδωκε τῷ Δαυὶδ τῷ βασιλεῖ υἱὸν σοφὸν, καὶ ἐπιστάμενον ἐπιστήμην καὶ σύνεσιν, ὃς οἰκοδομήσει οἶκον τῷ Κυρίῳ, καὶ οἶκον τῇ βασιλείᾳ αὐτοῦ.
\vs{12}Καὶ νῦν ἀπέστειλά σοι ἄνδρα σοφὸν καὶ εἰδότα σύνεσιν Χιρὰμ τὸν πατέρα μου,
\vs{13}ἡ μήτηρ αὐτοῦ ἀπὸ θυγατέρων Δὰν, καὶ ὁ πατὴρ αὐτοῦ ἀνὴρ Τύριος, εἰδότα ποιῆσαι ἐν χρυσίῳ, καὶ ἐν ἀργυρίῳ, καὶ ἐν χαλκῷ, καὶ ἐν σιδήρῳ, καὶ ἐν λίθοις καὶ ξύλοις, καὶ ὑφαίνειν ἐν τῇ πορφύρᾳ, καὶ ἐν τῇ ὑακίνθῳ, καὶ ἐν τῇ βύσσῳ, καὶ ἐν τῷ κοκκίνῳ, καὶ γλύψαι γλυφὰς, καὶ διανοεῖσθαι πᾶσαν διανόησιν, ὅσα ἂν δῷς αὐτῷ μετὰ τῶν σοφῶν σου, καὶ σοφῶν Δαυὶδ κυρίου μου πατρός σου.
\vs{14}Καὶ νῦν τὸν σῖτον καὶ τὴν κριθὴν, καὶ τὸ ἔλαιον, καὶ τὸν οἶνον ἃ εἶπεν ὁ κύριός μου, ἀποστειλάτω τοῖς παισὶν αὐτοῦ.
\vs{15}Καὶ ἡμεῖς κόψομεν ξύλα ἐκ τοῦ Λιβάνου κατὰ πᾶσαν τὴν χρείαν σου, καὶ ἄξομεν αὐτὰ σχεδίαις ἐπὶ θάλασσαν Ἰόππης, καὶ σὺ ἄξεις αὐτὰ εἰς Ἱερουσαλήμ.

\vs{16}Καὶ συνήγαγε Σαλωμὼν πάντας τοὺς ἄνδρας τοὺς προσηλύτους τοὺς ἐν γῇ Ἰσραὴλ μετὰ τὸν ἀριθμὸν ὃν ἠρίθμησεν αὐτοὺς Δαυὶδ ὁ πατὴρ αὐτοῦ, καὶ εὑρέθησαν ἑκατὸν πεντήκοντα χιλιάδες καὶ τρισχίλιοι ἑξακόσιοι.
\vs{17}Καὶ ἐποίησεν ἐξ αὐτῶν ἑβδομήκοντα χιλιάδας νωτοφόρων, καὶ ὀγδοήκοντα χιλιάδας λατόμων, καὶ τρισχιλίους ἑξακοσίους ἐργοδιώκτας ἐπὶ τὸν λαόν.

\ch{3}
Καὶ ἤρξατο Σαλωμὼν τοῦ οἰκοδομεῖν τὸν οἶκον Κυρίου ἐν Ἱερουσαλὴμ ἐν ὄρει τοῦ Ἀμωρία, οὗ ὤφθη Κύριος τῷ Δαυὶδ πατρὶ αὐτοῦ, ἐν τῷ τόπῳ ᾧ ἡτοίμασε Δαυὶδ ἐν ἅλῳ Ὀρνὰ τοῦ Ἰεβουσαίου.
\vs{2}Καὶ ἤρξατο οἰκοδομῆσαι ἐν τῷ μηνὶ τῷ δευτέρῳ ἐν τῷ ἔτει τῷ τετάρτῳ τῆς βασιλείας αὐτοῦ.

\vs{3}Καὶ ταῦτα ἤρξατο Σαλωμὼν τοῦ οἰκοδομῆσαι τὸν οἶκον τοῦ Θεοῦ· μῆκος πήχεων ἡ διαμέτρησις ἡ πρώτη πήχεων ἑξήκοντα, καὶ εὖρος πήχεων εἴκοσι.
\vs{4}Καὶ αἰλὰμ κατὰ πρόσωπον τοῦ οἴκου, μῆκος ἐπὶ πρόσωπον πλάτους τοῦ οἴκου πήχεων εἴκοσι, καὶ ὕψος πήχεων ἑκατὸν εἴκοσι, καὶ κατεχρύσωσεν αὐτὸν ἔσωθεν χρυσίῳ καθαρῷ.
\vs{5}Καὶ τὸν οἶκον τὸν μέγαν ἐξύλωσε ξύλοις κεδρίνοις, καὶ κατεχρύσωσε χρυσίῳ καθαρῷ, καὶ ἔγλυψεν ἐπʼ αὐτοῦ φοίνικας καὶ χαλαστά.
\vs{6}Καὶ ἐκόσμησε τὸν οἶκον λίθοις τιμίοις εἰς δόξαν, καὶ ἐχρύσωσε χρυσίῳ χρυσίου τοῦ ἐκ Φαρουΐμ.
\vs{7}Καὶ ἐχρύσωσε τὸν οἶκον, καὶ τοὺς τοίχους αὐτοῦ, καὶ τοὺς πυλῶνας, καὶ τὰ ὀροφώματα, καὶ τὰ θυρώματα χρυσίῳ, καὶ ἔγλυψε χερουβὶμ ἐπὶ τῶν τοίχων.

\vs{8}Καὶ ἐποίησε τὸν οἶκον τοῦ ἁγίου τῶν ἁγίων, μῆκος αὐτοῦ ἐπὶ πρόσωπον, πλάτος τοῦ οἴκου πήχεων εἴκοσι, καὶ τὸ μῆκος πήχεων εἴκοσι, καὶ ἐχρύσωσεν αὐτὸν χρυσίῳ καθαρῷ εἰς χερουβὶμ εἰς τάλαντα ἑξακόσια.
\vs{9}Καὶ ὁλκὴ τῶν ἥλων, ὁλκὴ τοῦ ἑνὸς πεντήκοντα σίκλοι χρυσίου, καὶ τὸ ὑπερῷον ἐχρύσωσε χρυσίῳ.

\vs{10}Καὶ ἐποίησεν ἐν τῷ οἴκῳ τῷ ἁγίῳ τῶν ἁγίων χερουβὶμ δύο, ἔργον ἐκ ξύλων· καὶ ἐχρύσωσεν αὐτὰ χρυσίῳ.
\vs{11}Καὶ αἱ πτέρυγες τῶν χερουβὶμ τὸ μῆκος πήχεων εἴκοσι, καὶ ἡ πτέρυξ ἡ μία πήχεων πέντε ἁπτομένη τοῦ τοίχου τοῦ οἴκου, καὶ ἡ πτέρυξ ἡ ἑτέρα πήχεων πέντε ἁπτομένη τῆς πτέρυγος τοῦ χερουβὶμ τοῦ ἑτέρου.
\vs{13}Καὶ αἱ πτέρυγες τῶν χερουβὶμ τούτων διαπεπετασμέναι πήχεων εἴκοσι, καὶ αὐτὰ ἑστηκότα ἐπὶ τοὺς πόδας αὐτῶν, καὶ τὰ πρόσωπα αὐτῶν εἰς τὸν οἶκον.
\vs{14}Καὶ ἐποίησε τὸ καταπέτασμα ὑακίνθου, καὶ πορφύρας, καὶ κοκκίνου, καὶ βύσσου, καὶ ὕφανεν ἐν αὐτῷ χερουβίμ.

\vs{15}Καὶ ἐποίησεν ἔμπροσθεν τοῦ οἴκου στύλους δύο, πήχεων τριακονταπέντε τὸ ὕψος, καὶ τὰς κεφαλὰς αὐτῶν πήχεων πέντε.
\vs{16}Καὶ ἐποίησε σερσερὼθ ἐν τῷ δαβὶρ, καὶ ἔδωκεν ἐπὶ τῶν κεφαλῶν τῶν στύλων· καὶ ἐποίησε ῥοΐσκους ἑκατὸν, καὶ ἔθηκεν ἐπὶ τῶν χαλαστῶν.
\vs{17}Καὶ ἔστησε τοὺς στύλους κατὰ πρόσωπον τοῦ ναοῦ, ἕνα ἐκ δεξιῶν, καὶ τὸν ἕνα ἐξ εὐωνύμων· καὶ ἐκάλεσε τὸ ὄνομα τοῦ ἐκ δεξιῶν, Κατόρθωσις, καὶ τὸ ὄνομα τοῦ ἐξ ἀριστερῶν, Ἰσχύς.

\ch{4}
Καὶ ἐποίησε θυσιαστήριον χαλκοῦν, εἴκοσι πήχεων τὸ μῆκος, καὶ εἴκοσι πήχεων τὸ εὖρος, καὶ δέκα πήχεων τὸ ὕψος.
\vs{2}Καὶ ἐποίησε τὴν θάλασσαν χυτὴν, δέκα πήχεων τὴν διαμέτρησιν, στρογγύλην κυκλόθεν, καὶ πέντε πήχεων τὸ ὕψος, καὶ τὸ κύκλωμα τριάκοντα πήχεων.
\vs{3}Καὶ ὁμοίωμα μόσχων ὑποκάτω αὐτῆς, κύκλῳ κυκλοῦσιν αὐτήν· δέκα πήχεις περιέχουσι τὸν λουτῆρα κυκλόθεν· δύο γένη ἐχώνευσαν τοὺς μόσχους ἐν τῇ χωνεύσει αὐτῶν
\vs{4}ᾗ ἐποίησαν αὐτοὺς δώδεκα μόσχους, οἱ τρεῖς βλέποντες Βοῤῥᾶν, καὶ οἱ τρεῖς δυσμὰς, καὶ οἱ τρεῖς Νότον, καὶ οἱ τρεῖς κατʼ ἀνατολὰς, καὶ ἡ θάλασσα ἐπʼ αὐτῶν ἄνω, ἦσαν τὰ ὀπίσθια αὐτῶν ἔσω.
\vs{5}Καὶ τὸ πάχος αὐτῆς παλαιστὴς, καὶ τὸ χεῖλος αὐτῆς ὡς χεῖλος ποτηρίου, διαγεγλυμμένα βλαστοὺς κρίνου, χωροῦσαν μετρητὰς τρισχιλίους· καὶ ἐξετέλεσε.

\vs{6}Καὶ ἐποίησε λουτῆρας δέκα, καὶ ἔθηκε τοὺς πέντε ἐκ δεξιῶν καὶ τοὺς πέντε ἐξ ἀριστερῶν, τοῦ πλύνειν ἐν αὐτοῖς τὰ ἔργα τῶν ὁλοκαυτωμάτων, καὶ ἀποκλύζειν ἐν αὐτοῖς, καὶ ἡ θάλασσα εἰς τὸ νίπτεσθαι τοὺς ἱερεῖς ἐν αὐτῇ.

\vs{7}Καὶ ἐποίησε τὰς λυχνίας τὰς χρυσᾶς δέκα κατὰ τὸ κρίμα αὐτῶν, καὶ ἔθηκεν ἐν τῷ ναῷ, πέντε ἐκ δεξιῶν καὶ πέντε ἐξ ἀριστερῶν.

\vs{8}Καὶ ἐποίησε τραπέζας δέκα, καὶ ἔθηκεν ἐν τῷ ναῷ, πέντε ἐκ δεξιῶν καὶ πέντε ἐξ εὐωνύμων, καὶ ἐποίησε φιάλας χρυσᾶς ἑκατὸν,
\vs{9}καὶ ἐποίησε τὴν αὐλὴν τῶν ἱερέων, καὶ τὴν αὐλὴν τὴν μεγάλην, καὶ θύρας τῇ αὐλῇ, καὶ θυρώματα αὐτῶν κατακεχαλκωμένα χαλκῷ.
\vs{10}Καὶ τὴν θάλασσαν ἔθηκεν ἀπὸ γωνίας τοῦ οἴκου ἐκ δεξιῶν ὡς πρὸς ἀνατολὰς κατέναντι.

\vs{11}Καὶ ἐποίησε Χιρὰμ τὰς κρεάγρας, καὶ τὰ πυρεῖα, καὶ τὴν ἐσχάραν τοῦ θυσιαστηρίου, καὶ πάντα τὰ σκεύη αὐτοῦ· καὶ συνετέλεσε Χιρὰμ ποιῆσαι πᾶσαν τὴν ἐργασίαν ἣν ἐποίησε Σαλωμὼν τῷ βασιλεῖ ἐν οἴκῳ τοῦ Θεοῦ,
\vs{12}στύλους δύο, καὶ ἐπʼ αὐτῶν γελὰθ τῇ χωθαρὲθ ἐπὶ τῶν κεφαλῶν τῶν στύλων δύο, καὶ δίκτυα δύο συγκαλύψαι τὰς κεφαλὰς τῶν χωθαρὲθ ἅ ἐστιν ἐπὶ τῶν κεφαλῶν τῶν στύλων,
\vs{13}καὶ κώδωνας χρυσοῦς τετρακοσίους εἰς τὰ δύο δίκτυα, καὶ δύο γένη ῥοΐσκων ἐν τῷ δικτύῳ τῷ ἑνὶ τοῦ συγκαλύψαι τὰς δύο γωλὰθ τῶν χωθαρὲθ ἅ ἐστιν ἐπάνω τῶν στύλων·
\vs{14}Καὶ τοὺς μεχωνὼθ ἐποίησε δέκα, καὶ τοὺς λουτῆρας ἐποίησεν ἐπὶ τοὺς μεχωνὼθ,
\vs{15}καὶ τὴν θάλασσαν μίαν, καὶ τοὺς μόσχους τοὺς δώδεκα ὑποκάτω αὐτῆς,
\vs{16}καὶ τοὺς ποδιστῆρας, καὶ τοὺς ἀναλημπτῆρας, καὶ τοὺς λέβητας, καὶ τὰς κρεάγρας, καὶ πάντα τὰ σκεύη αὐτῶν ἃ ἐποίησε Χιρὰμ, καὶ ἀνήνεγκε τῷ βασιλεῖ Σαλωμὼν ἐν οἴκῳ Κυρίου, χαλκοῦ καθαροῦ.
\vs{17}Ἐν τῷ περιχώρῳ τοῦ Ἰορδάνου ἐχώνευσεν αὐτὰ ὁ βασιλεὺς ἐν τῷ πάχει τῆς γῆς ἐν οἴκῳ Σοκχὼθ καὶ ἀναμέσον Σαρηδαθά.

\vs{18}Καὶ ἐποίησε Σαλωμὼν πάντα τὰ σκεύη ταῦτα εἰς πλῆθος σφόδρα, ὅτι οὐκ ἐξέλιπεν ὁλκὴ τοῦ χαλκοῦ.
\vs{19}Καὶ ἐποίησε Σαλωμὼν πάντα τὰ σκεύη οἴκου Κυρίου, καὶ τὸ θυσιαστήριον τὸ χρυσοῦν, καὶ τὰς τραπέζας, καὶ ἐπʼ αὐτῶν ἄρτοι προθέσεως,
\vs{20}καὶ τὰς λυχνίας, καὶ τοὺς λύχνους τοῦ φωτὸς κατὰ τὸ κρίμα καὶ κατὰ πρόσωπον τοῦ δαβὶρ χρυσίου καθαροῦ,
\vs{21}καὶ λαβίδες αὐτῶν, καὶ οἱ λύχνοι αὐτῶν, καὶ τὰς φιάλας, καὶ τὰς θυΐσκας, καὶ τὰ πυρεῖα χρυσίου καθαροῦ,
\vs{22}καὶ ἡ θύρα τοῦ οἴκου ἡ ἐσωτέρα εἰς τὰ ἅγια τῶν ἁγίων, καὶ τὰς θύρας τοῦ οἴκου τοῦ ναοῦ χρυσᾶς· καὶ συνετελέσθη πᾶσα ἡ ἐργασία ἣν ἐποίησε Σαλωμὼν ἐν οἴκῳ Κυρίου.

\ch{5}
Καὶ εἰσήνεγκε Σαλωμὼν τὰ ἅγια Δαυὶδ τοῦ πατρὸς αὐτοῦ, τὸ ἀργύριον, καὶ τὸ χρυσίον, καὶ τὰ σκεύη, καὶ ἔδωκεν εἰς θησαυρὸν οἴκου Κυρίου.

\vs{2}Τότε ἐξεκκλησίασε Σαλωμὼν πάντας τοὺς πρεσβυτέρους Ἰσραὴλ, καὶ πάντας τοὺς ἄρχοντας τῶν φυλῶν τοὺς ἡγουμένους πατριῶν υἱῶν Ἰσραὴλ εἰς Ἱερουσαλὴμ, τοῦ ἀνενέγκαι κιβωτὸν διαθήκης Κυρίου ἐκ πόλεως Δαυὶδ, αὕτη Σιών.
\vs{3}Καὶ ἐξεκκλησιάσθησαν πρὸς τὸν βασιλέα πᾶς Ἰσραὴλ ἐν τῇ ἑορτῇ, οὗτος ὁ μὴν ἕβδομος.
\vs{4}Καὶ ἦλθον πάντες οἱ πρεσβύτεροι Ἰσραὴλ, καὶ ἔλαβον πάντες οἱ Λευῖται τὴν κιβωτὸν,
\vs{5}καὶ τὴν σκηνὴν τοῦ μαρτυρίου, καὶ πάντα τὰ σκεύη τὰ ἅγια τὰ ἐν τῇ σκηνῇ, καὶ ἀνήνεγκαν αὐτὴν οἱ ἱερεῖς καὶ οἱ Λευῖται.
\vs{6}Καὶ ὁ βασιλεὺς Σαλωμὼν, καὶ πᾶσα συναγωγὴ Ἰσραὴλ, καὶ οἱ φοβούμενοι, καὶ οἱ ἐπισυνηγμένοι αὐτῶν ἔμπροσθεν τῆς κιβωτοῦ θύοντες μόσχους καὶ πρόβατα, οἳ οὐκ ἀριθμηθήσονται καὶ οἳ οὐ λογισθήσονται ἀπὸ τοῦ πλήθους.
\vs{7}Καὶ εἰσήνεγκαν οἱ ἱερεῖς τὴν κιβωτὸν διαθήκης Κυρίου εἰς τὸν τόπον αὐτῆς, εἰς τὸ δαβὶρ τοῦ οἴκου εἰς τὰ ἅγια τῶν ἁγίων, ὑποκάτω τῶν πτερύγων τῶν χερουβίμ.
\vs{8}Καὶ ἦν τὰ χερουβὶμ διαπεπετακότα τὰς πτέρυγας αὐτῶν ἐπὶ τὸν τόπον τῆς κιβωτοῦ, καὶ συνεκάλυπτε τὰ χερουβὶμ ἐπὶ τὴν κιβωτὸν, καὶ ἐπὶ τοὺς ἀναφορεῖς αὐτῆς ἐπάνωθεν,
\vs{9}καὶ ὑπερεῖχον οἱ ἀναφορεῖς, καὶ ἐβλέποντο αἱ κεφαλαὶ τῶν ἀναφορέων ἐκ τῶν ἁγίων εἰς πρόσωπον τοῦ δαβὶρ, οὐκ ἐβλέποντο ἔξω, καὶ ἦσαν ἐκεῖ ἕως τῆς ἡμέρας ταύτης.
\vs{10}Οὐκ ἦν ἐν τῇ κιβωτῷ πλὴν δύο πλάκες ἃς ἔθηκε Μωυσῆς ἐν Χωρὴβ, ἃ διέθετο Κύριος μετὰ τῶν υἱῶν Ἰσραὴλ, ἐν τῷ ἐξελθεῖν αὐτοὺς ἐκ γῆς Αἰγύπτου.

\vs{11}Καὶ ἐγένετο ἐν τῷ ἐξελθεῖν τοὺς ἱερεῖς ἐκ τῶν ἁγίων, ὅτι πάντες οἱ ἱερεῖς οἱ εὑρεθέντες ἡγιάσθησαν, οὐκ ἦσαν διατεταγμένοι κατʼ ἐφημερίαν.
\vs{12}Καὶ οἱ Λευῖται οἱ ψαλτῳδοὶ πάντες τοῖς υἱοῖς Ἀσὰφ τῷ Αἰμὰν τῷ Ἰδιθοὺν καὶ τοῖς υἱοῖς αὐτοῦ, καὶ τοῖς ἀδελφοῖς αὐτοῦ τῶν ἐνδεδυμένων στολὰς βασσίνας ἐν κυμβάλοις καὶ ἐν νάβλαις καὶ ἐν κινύραις, ἑστηκότες κατέναντι τοῦ θυσιαστηρίου, καὶ μετʼ αὐτῶν ἱερεῖς ἑκατὸν εἴκοσι σαλπίζοντες ταῖς σάλπιγξι.
\vs{13}Καὶ ἐγένετο μία φωνὴ ἐν τῷ σαλπίζειν καὶ ἐν τῷ ψαλτῳδεῖν καὶ ἐν τῷ ἀναφωνεῖν φωνῇ μιᾷ τοῦ ἐξομολογεῖσθαι καὶ αἰνεῖν τῷ Κυρίῳ· καὶ ὡς ὕψωσαν φωνὴν ἐν σάλπιγξι, καὶ ἐν κυμβάλοις, καὶ ἐν ὀργάνοις τῶν ᾠδῶν, καὶ ἔλεγον, ἐξομολογεῖσθε τῷ Κυρίῳ, ὅτι ἀγαθὸν, ὅτι εἰς τὸν αἰῶνα τὸ ἔλεος αὐτοῦ· καὶ ὁ οἶκος ἐνεπλήσθη νεφέλης δόξης Κυρίου.
\vs{14}Καὶ οὐκ ἠδύναντο οἱ ἱερεῖς τοῦ στῆναι λειτουργεῖν ἀπὸ προσώπου τῆς νεφέλης, ὅτι ἐνέπλησε δόξα Κυρίου τὸν οἶκον τοῦ Θεοῦ.

\ch{6}
Τότε εἰπε Σαλωμὼν, Κύριος εἶπε τοῦ κατασκηνῶσαι ἐν γνόφῳ,
\vs{2}καὶ ἐγὼ ᾠκοδόμηκα οἶκον τῷ ὀνόματί σου ἅγιόν σοι καὶ ἕτοιμον τοῦ κατασκηνῶσαι εἰς τοὺς αἰῶνας.

\vs{3}Καὶ ἐπέστρεψεν ὁ βασιλεὺς τὸ πρόσωπον αὐτοῦ, καὶ εὐλόγησε τὴν πᾶσαν ἐκκλησίαν Ἰσραὴλ, καὶ πᾶσα ἡ ἐκκλησία Ἰσραὴλ παρειστήκει,
\vs{4}καὶ εἶπεν, εὐλογητὸς Κύριος ὁ Θεὸς Ἰσραὴλ, ὡς ἐλάλησεν ἐν στόματι αὐτοῦ πρὸς Δαυὶδ τὸν πατέρα μου, καὶ ἐν χερσὶν αὐτοῦ ἐπλήρωσε, λέγων,
\vs{5}ἀπὸ τῆς ἡμέρας ἧς ἀνήγαγον τὸν λαόν μου ἐκ γῆς Αἰγύπτου, οὐκ ἐξελεξάμην ἐν πόλει ἀπὸ πασῶν φυλῶν Ἰσραὴλ, τοῦ οἰκοδομῆσαι οἶκον τοῦ εἶναι τὸ ὄνομά μου ἐκεῖ, καὶ οὐκ ἐξελεξάμην ἐν ἀνδρὶ τοῦ εἶναι εἰς ἡγούμενον ἐπὶ τὸν λαόν μου Ἰσραήλ·
\vs{6}Καὶ ἐξελεξάμην τὴν Ἱερουσαλὴμ γενέσθαι τὸ ὄνομά μου ἐκεῖ, καὶ ἐξελεξάμην ἐν Δαυὶδ τοῦ εἶναι ἐπὶ τὸν λαόν μου Ἰσραήλ.
\vs{7}Καὶ ἐγένετο ἐπὶ καρδίαν Δαυὶδ τοῦ πατρός μου, τοῦ οἰκοδομῆσαι οἶκον τῷ ὀνόματι Κυρίου Θεοῦ Ἰσραήλ.
\vs{8}Καὶ εἶπε Κύριος πρὸς Δαυὶδ πατέρα μου, διότι ἐγένετο ἐπὶ καρδίαν σου τοῦ οἰκοδομῆσαι οἶκον τῷ ὀνόματί μου, καλῶς ἐποίησας, ὅτι ἐγένετο ἐπὶ τὴν καρδίαν σου.
\vs{9}Πλὴν σὺ οὐκ οἰκοδομήσεις τὸν οἶκον, ὅτι ὁ υἱός σου ὃς ἐξελεύσεται ἐκ τῆς ὀσφύος σου, οὗτος οἰκοδομήσει τὸν οἶκον τῷ ὀνόματί μου.
\vs{10}Καὶ ἀνέστησε Κύριος τὸν λόγον τοῦτον, ὃν ἐλάλησε, καὶ ἐγενήθην ἀντὶ Δαυὶδ πατρός μου, καὶ ἐκάθισα ἐπὶ τὸν θρόνον Ἰσραὴλ, καθὼς ἐλάλησε Κύριος· καὶ ᾠκοδόμησα τὸν οἶκον τῷ ὀνόματι Κυρίου Θεοῦ Ἰσραὴλ,
\vs{11}καὶ ἔθηκα ἐκεῖ τὴν κιβωτὸν ἐν ᾗ ἐκεῖ διαθήκη Κυρίου, ἣν διέθετο τῷ Ἰσραήλ.

\vs{12}Καὶ ἔστη κατέναντι τοῦ θυσιαστηρίου Κυρίου ἔναντι πάσης ἐκκλησίας Ἰσραὴλ, καὶ διεπέτασε τὰς χεῖρας αὐτοῦ·
\vs{13}Ὅτι ἐποίησε Σαλωμὼν βάσιν χαλκῆν, καὶ ἔθηκεν αὐτὴν ἐν μέσῳ τῆς αὐλῆς τοῦ ἱεροῦ, πέντε πήχεων τὸ μῆκος αὐτῆς, καὶ πέντε πήχεων τὸ εὖρος αὐτῆς, καὶ τριῶν πήχεων τὸ ὕψος αὐτῆς· καὶ ἔστη ἐπʼ αὐτῆς, καὶ ἔπεσεν ἐπὶ τὰ γόνατα ἔναντι πάσης ἐκκλησίας Ἰσραὴλ, καὶ διεπέτασε τὰς χεῖρας αὐτοῦ εἰς τὸν οὐρανὸν,
\vs{14}καὶ εἶπε,

Κύριε ὁ Θεὸς Ἰσραὴλ, οὐκ ἔστιν ὅμοιός σοι Θεὸς ἐν οὐρανῷ καὶ ἐπὶ τῆς γῆς, φυλάσσων τὴν διαθήκην καὶ τὸ ἔλεος τοῖς παισὶ σου τοῖς πορευομένοις ἐναντίον σου ἐν ὅλῃ καρδίᾳ·
\vs{15}Ἃ ἐφύλαξας τῷ παιδί σου Δαυὶδ τῷ πατρί μου, ἃ ἐλάλησας αὐτῷ, λέγων· καὶ ἐλάλησας ἐν στόματί σου, καὶ ἐν χερσί σου ἐπλήρωσας, ὡς ἡ ἡμέρα αὕτη.
\vs{16}Καὶ νῦν, Κύριε ὁ Θεὸς Ἰσραὴλ, φύλαξον τῷ παιδί σου τῷ Δαυὶδ τῷ πατρί μου ἃ ἐλάλησας αὐτῷ, λέγων, οὐκ ἐκλείψει σοι ἀνὴρ ἀπὸ προσώπου μου καθήμενος ἐπὶ θρόνου Ἰσραὴλ, πλὴν ἐὰν φυλάξωσιν οἱ υἱοί σου τὴν ὁδὸν αὐτῶν τοῦ πορεύεσθαι ἐν τῷ νόμῳ μου, ὡς ἐπορεύθης ἐναντίον μου.
\vs{17}Καὶ νῦν, Κύριε ὁ Θεὸς Ἰσραὴλ, πιστωθήτω δὴ τὸ ῥῆμά σου ὃ ἐλάλησας τῷ παιδί σου τῷ Δαυίδ.

\vs{18}Ὅτι εἰ ἀληθῶς κατοικήσει Θεὸς μετὰ ἀνθρώπων ἐπὶ τῆς γῆς; εἰ ὁ οὐρανὸς καὶ ὁ οὐρανὸς τοῦ οὐρανοῦ οὐκ ἀρκέσουσί σοι, καὶ τίς ὁ οἶκος οὗτος ὃν ᾠκοδόμησα;
\vs{19}Καὶ ἐπιβλέψῃ ἐπὶ τὴν προσευχὴν παιδός σου καὶ ἐπὶ τὴν δέησίν μου, Κύριε ὁ Θεὸς, τοῦ ἐπακοῦσαι τῆς δεήσεως καὶ τῆς προσευχῆς ἧς ὁ παῖς σου προσεύχεται ἐναντίον σου σήμερον,
\vs{20}τοῦ εἶναι ὀφθαλμούς σου ἀνεῳγμένους ἐπὶ τὸν οἶκον τοῦτον ἡμέρας καὶ νυκτὸς εἰς τὸν τόπον τοῦτον, ὃν εἶπας ἐπικληθῆναι τὸ ὄνομά σου ἐκεῖ, τοῦ ἀκοῦσαι τῆς προσευχῆς ἧς προσεύχεται ὁ παῖς σου εἰς τὸν τόπον τοῦτον·
\vs{21}Καὶ ἀκούσῃ τῆς δεήσεως τοῦ παιδός σου, καὶ λαοῦ σου Ἰσραὴλ, ἃ ἂν προσεύξωνται εἰς τὸν τόπον τοῦτον· καὶ σὺ εἰσακούσῃ ἐν τῷ τόπῳ τῆς κατοικήσεώς σου ἐκ τοῦ οὐρανοῦ, καὶ ἀκούσῃ καὶ ἵλεως ἔσῃ.

\vs{22}Ἐὰν ἁμάρτῃ ἀνὴρ τῷ πλησίον αὐτοῦ καὶ λάβῃ ἐπʼ αὐτὸν ἀρὰν τοῦ ἀρᾶσθαι αὐτὸν, καὶ ἔλθῃ καὶ ἀράσηται κατέναντι τοῦ θυσιαστηρίου ἐν τῷ οἴκῳ τούτῳ,
\vs{23}καὶ σὺ εἰσακούσῃ ἐκ τοῦ οὐρανοῦ, καὶ ποιήσεις, καὶ κρινεῖς τοὺς δούλους σου, τοῦ ἀποδοῦναι τῷ ἀνόμῳ, καὶ ἀποδοῦναι ὁδοὺς αὐτοῦ εἰς κεφαλὴν αὐτοῦ, καὶ τοῦ δικαιῶσαι δίκαιον, τοῦ ἀποδοῦναι αὐτῷ κατὰ τὴν δικαιοσύνην αὐτοῦ.

\vs{24}Καὶ ἐὰν θραυσθῇ ὁ λαός σου Ἰσραὴλ κατέναντι τοῦ ἐχθροῦ, ἐὰν ἁμάρτωσί σοι, καὶ ἐπιστρέψωσι καὶ ἐξομολογήσωνται τῷ ὀνόματί σου, καὶ προσεύξωνται καὶ δεηθῶσιν ἐναντίον σου ἐν τῷ οἴκῳ τούτῳ,
\vs{25}καὶ σὺ εἰσακούσῃ ἐκ τοῦ οὐρανοῦ, καὶ ἵλεως ἔσῃ ταῖς ἁμαρτίαις λαοῦ σου Ἰσραὴλ, καὶ ἀποστρέψεις αὐτοὺς εἰς τὴν γῆν ἣν ἔδωκας αὐτοῖς καὶ τοῖς πατράσιν αὐτῶν.

\vs{26}Ἐν τῷ συσχεθῆναι τὸν οὐρανὸν καὶ μὴ γενέσθαι ὑετὸν, ὅτι ἁμαρτήσονταί σοι, καὶ προσεύξονται εἰς τὸν τόπον τοῦτον, καὶ αἰνέσουσι τὸ ὄνομά σου, καὶ ἀπὸ τῶν ἁμαρτιῶν αὐτῶν ἐπιστρέψουσιν, ὅτι ταπεινώσεις αὐτοὺς,
\vs{27}καὶ σὺ εἰσακούσῃ ἐκ τοῦ οὐρανοῦ, καὶ ἵλεως ἔσῃ ταῖς ἁμαρτίαις τῶν παίδων καὶ τοῦ λαοῦ σου Ἰσραὴλ, ὅτι δηλώσεις αὐτοῖς τὴν ὁδὸν τὴν ἀγαθὴν, ἐν ᾗ πορεύσονται ἐν αὐτῇ, καὶ δώσεις ὑετὸν ἐπὶ τὴν γῆν σου, ἣν ἔδωκας τῷ λαῷ σου εἰς κληρονομίαν.

\vs{28}Λιμὸς ἐὰν γένηται ἐπὶ τῆς γῆς, θάνατος ἐὰν γένηται, ἀνεμοφθορία καὶ ἴκτερος, ἀκρὶς καὶ βροῦχος ἐὰν γένηται, καὶ ἐὰν θλίψῃ αὐτὸν ὁ ἐχθρὸς κατέναντι τῶν πόλεων αὐτῶν, κατὰ πᾶσαν πληγὴν καὶ πάντα πόνον,
\vs{29}καὶ πᾶσα προσευχὴ, καὶ πᾶσα δέησις ἣ ἐὰν γένηται παντὶ ἀνθρώπῳ καὶ παντὶ λαῷ σου Ἰσραὴλ, ἐὰν γνῷ ἄνθρωπος τὴν ἁφὴν αὐτοῦ καὶ τὴν μαλακίαν αὐτοῦ, καὶ διαπετάσῃ τὰς χεῖρας αὐτοῦ εἰς τὸν οἶκον τοῦτον,
\vs{30}καὶ σὺ εἰσακούσῃ ἐκ τοῦ οὐρανοῦ ἐξ ἑτοίμου κατοικητηρίου σου, καὶ ἱλάσῃ, καὶ δώσεις ἀνδρὶ κατὰ τὰς ὁδοὺς αὐτοῦ, ὡς ἂν γνῷς τὴν καρδίαν αὐτοῦ, ὅτι μόνος γινώσκεις τὴν καρδίαν υἱῶν ἀνθρώπων,
\vs{31}ὅπως φοβῶνται πάσας ὁδούς σου πάσας τὰς ἡμέρας ἃς αὐτοὶ ζῶσιν ἐπὶ πρόσωπον τῆς γῆς, ἧς ἔδωκας τοῖς πατράσιν ἡμῶν.

\vs{32}Καὶ πᾶς ἀλλότριος ὃς οὐκ ἐκ τοῦ λαοῦ σου Ἰσραήλ ἐστιν αὐτὸς, καὶ ἔλθῃ ἐκ γῆς μακρόθεν διὰ τὸ ὄνομά σου τὸ μέγα, καὶ τὴν χεῖρά σου τὴν κραταιὰν, καὶ τὸν βραχίονά σου τὸν ὑψηλὸν, καὶ ἔλθωσι καὶ προσεύξωνται εἰς τὸν τόπον τοῦτον,
\vs{33}καὶ σὺ εἰσακούσῃ ἐκ τοῦ οὐρανοῦ ἐξ ἑτοίμου κατοικητηρίου σου, καὶ ποιήσεις κατὰ πάντα ὅσα ἂν ἐπικαλέσηταί σε ὁ ἀλλότριος, ὅπως γνῶσι πάντες οἱ λαοὶ τῆς γῆς τὸ ὄνομά σου, καὶ τοῦ φοβεῖσθαί σε, ὡς ὁ λαός σου Ἰσραὴλ, καὶ τοῦ γνῶναι ὅτι τὸ ὄνομά σου ἐπικέκληται ἐπὶ τὸν οἶκον τοῦτον, ὃν ᾠκοδόμησα.

\vs{34}Ἐὰν δὲ ἐξέλθῃ ὁ λαός σου εἰς πόλεμον ἐπὶ τοὺς ἐχθροὺς αὐτοῦ ἐν ὁδῷ ᾗ ἀποστελεῖς αὐτούς, καὶ προσεύξωνται πρὸς σὲ κατὰ τὴν ὁδὸν τῆς πόλεως ταύτης ἣν ἐξελέξω ἐν αὐτῇ, καὶ οἴκου οὗ ᾠκοδόμηκα τῷ ὀνόματί σου,
\vs{35}καὶ ἀκούσῃ ἐκ τοῦ οὐρανοῦ τῆς προσευχῆς αὐτῶν καὶ τῆς δεήσεως αὐτῶν, καὶ ποιήσεις τὸ δικαίωμα αὐτῶν.

\vs{36}Ὅτι ἁμαρτήσονταί σοι, ὅτι οὐκ ἔσται ἄνθρωπος ὃς οὐχ ἁμαρτήσεται, καὶ πατάξεις αὐτοὺς καὶ παραδώσεις αὐτοὺς κατὰ πρόσωπον ἐχθρῶν καὶ αἰχμαλωτεύσουσιν αὐτοὺς οἱ αἰχμαλωτεύοντες αὐτοὺς εἰς γῆν ἐχθρῶν εἰς γῆν μακρὰν ἢ ἐγγὺς,
\vs{37}καὶ ἐπιστρέψωσι καρδίαν αὐτῶν ἐν τῇ γῇ αὐτῶν οὗ μετήχθησαν ἐκεῖ, καί γε ἐπιστρέψωσι καὶ δεηθῶσί σου ἐν τῇ αἰχμαλωσίᾳ αὐτῶν, λέγοντες, ἡμάρτομεν, ἠνομήσαμεν, ἠδικήσαμεν,
\vs{38}καὶ ἐπιστρέψωσι πρὸς σὲ ἐν ὅλῃ καρδίᾳ καὶ ἐν ὅλῃ ψυχῇ αὐτῶν ἐν γῇ αἰχμαλωτευσάντων αὐτοὺς, ὅπου ᾐχμαλώτευσαν αὐτοὺς, καὶ προσεύξωνται ὁδὸν γῆς αὐτῶν ἧς ἔδωκας τοῖς πατράσιν αὐτῶν, καὶ τῆς πόλεως ἧς ἐξελέξω, καὶ τοῦ οἴκου οὗ ᾠκοδόμησα τῷ ὀνόματί σου,
\vs{39}καὶ ἀκούσῃ ἐκ τοῦ οὐρανοῦ ἐξ ἑτοίμου κατοικητηρίου σου τῆς προσευχῆς αὐτῶν καὶ τῆς δεήσεως αὐτῶν, καὶ ποιήσεις κρίματα, καὶ ἵλεως ἔσῃ τῷ λαῷ τῷ ἁμαρτῶντί σοι.

\vs{40}Καὶ νῦν Κύριε ἔστωσαν δὴ οἱ ὀφθαλμοί σου ἀνεῳγμένοι, καὶ τὰ ὦτά σου ἐπήκοα εἰς τὴν δέησιν τοῦ τόπου τούτου.
\vs{41}Καὶ νῦν ἀνάστηθι Κύριε ὁ Θεὸς εἰς τὴν κατάπαυσίν σου, σὺ καὶ ἡ κιβωτὸς τῆς ἰσχύος σου· ἱερεῖς σου Κύριε ὁ Θεὸς ἐνδύσαιντο σωτηρίαν, καὶ οἱ υἱοί σου εὐφρανθήτωσαν ἐν ἀγαθοῖς.
\vs{42}Κύριε ὁ Θεὸς, μὴ ἀποστρέψῃς τὸ πρόσωπον τοῦ χριστοῦ σου, μνήσθητι τὰ ἐλέη Δαυὶδ τοῦ δούλου σου.

\ch{7}
Καὶ ὡς συνετέλεσε Σαλωμὼν προσευχόμενος, καὶ τὸ πῦρ κατέβη ἐκ τοῦ οὐρανοῦ, καὶ κατέφαγε τὰ ὁλοκαυτώματα καὶ τὰς θυσίας, καὶ δόξα Κυρίου ἔπλησε τὸν οἶκον.
\vs{2}Καὶ οὐκ ἠδύναντο οἱ ἱερεῖς εἰσελθεῖν εἰς τὸν οἶκον Κυρίου ἐν τῷ καιρῷ ἐκείνῳ ὅτι ἔπλησε δόξα Κυρίου τὸν οἶκον.
\vs{3}Καὶ πάντες οἱ υἱοὶ Ισραὴλ ἑώρων καταβαῖνον τὸ πῦρ, καὶ ἡ δόξα Κυρίου ἐπὶ τὸν οἶκον· καὶ ἔπεσον ἐπὶ πρόσωπον ἐπὶ τὴν γῆν ἐπὶ τὸ λιθόστρωτον, καὶ προσεκύνησαν καὶ ᾔνουν τῷ Κυρίῳ, ὅτι ἀγαθὸν, ὅτι εἰς τὸν αἰῶνα τὸ ἔλεος αὐτοῦ.

\vs{4}Καὶ ὁ βασιλεὺς καὶ πᾶς ὁ λαὸς θύοντες θύματα ἔναντι Κυρίου.
\vs{5}Καὶ ἐθυσίασεν ὁ βασιλεὺς Σαλωμὼν τὴν θυσίαν μόσχων εἴκοσι καὶ δύο χιλιάδας, βοσκημάτων ἑκατὸν καὶ εἴκοσι χιλιάδας, καὶ ἐνεκαίνισε τὸν οἶκον τοὔ Θεοῦ ὁ βασιλεὺς καὶ πᾶς ὁ λαός.
\vs{6}Καὶ οἱ ἱερεῖς ἐπὶ τὰς φυλακὰς αὐτῶν ἑστηκότως, καὶ οἱ Λευῖται ἐν ὀργάνοις ᾠδῶν Κυρίου τοῦ Δαυὶδ τοῦ βασιλέως, τοῦ ἐξομολογεῖσθαι ἔναντι Κυρίου, ὅτι εἰς τὸν αἰῶνα τὸ ἔλεος αὐτοῦ, ἐν ὕμνοις Δαυὶδ διὰ χειρὸς αὐτῶν· καὶ οἱ ἱερεῖς σαλπίζοντες ταῖς σάλπιγξιν ἐναντίον αὐτῶν, καὶ πᾶς Ἰσραὴλ ἑστηκώς.
\vs{7}Καὶ ἡγίασε Σαλωμὼν τὸ μέσον τῆς αὐλῆς τῆς ἐν οἴκῳ Κυρίου, ὅτι ἐποίησεν ἐκεῖ τὰ ὁλοκαυτώματα καὶ τὰ στέατα τῶν σωτηρίων, ὅτι τὸ θυσιαστήριον τὸ χαλκοῦν ὃ ἐποίησε Σαλωμὼν, οὐκ ἐξεποίει δέξασθαι τὰ ὁλοκαυτώματα καὶ τὰ μαναὰ καὶ τὰ στέατα.

\vs{8}Καὶ ἐποίησε Σαλωμὼν τὴν ἑορτὴν ἐν τῷ καιρῷ ἐκείνῳ ἑπτὰ ἡμέρας, καὶ πᾶς Ἰσραὴλ μετʼ αὐτοῦ, ἐκκλησία μεγάλη σφόδρα ἀπὸ εἰσόδου Αἰμὰθ καὶ ἕως χειμάῤῥου Αἰγύπτου.
\vs{9}Καὶ ἐποίησεν ἐν τῇ ἡμέρᾳ τῇ ὀγδόῃ ἐξόδιον, ὅτι ἐγκαινισμὸν τοῦ θυσιαστηρίου ἐποίησεν ἑπτὰ ἡμέρας ἑορτήν.
\vs{10}Καὶ ἐν τῇ τρίτῃ καὶ εἰκοστῇ τοῦ μηνὸς τοῦ ἑβδόμου ἀπέστειλε τὸν λαὸν εἰς τὰ σκηνώματα αὐτῶν εὐφραινομένους, καὶ ἀγαθῇ καρδίᾳ ἐπὶ τοῖς ἀγαθοῖς οἷς ἐποίησε Κύριος τῷ Δαυὶδ, καὶ τῷ Σαλωμῶντι, καὶ τῷ Ἰσραὴλ λαῷ αὐτοῦ.

\vs{11}Καὶ συνετέλεσε Σαλωμὼν τὸν οἶκον Κυρίου, καὶ τὸν οἶκον τοῦ βασιλέως· καὶ πάντα ὅσα ἠθέλησεν ἐν τῇ ψυχῇ Σαλωμὼν τοῦ ποιῆσαι ἐν οἴκῳ Κυρίου καὶ ἐν οἴκῳ αὐτοῦ, εὐωδώθη.

\vs{12}Καὶ ὤφθη Κύριος τῷ Σαλωμὼν τὴν νύκτα, καὶ εἶπεν αὐτῷ, ἤκουσα τῆν προσευχῆς σου, καὶ ἐξελεξάμην ἐν τῷ τόπῳ τούτῳ ἐμαυτῷ εἰς οἶκον θυσίας.
\vs{13}Ἐὰν συσχῶ τὸν οὐρανὸν καὶ μὴ γένηται ὑετὸς, καὶ ἐὰν ἐντείλωμαι τῇ ἀκρίδι καταφαγεῖν τὸ ξύλον, καὶ ἐὰν ἀποστείλω θάνατον ἐν τῷ λαῷ μου,
\vs{14}καὶ ἐὰν ἐντραπῇ ὁ λαός μου ἐφʼ οὓς ἐπικέκληται τὸ ὄνομά μου ἐπʼ αὐτοὺς, καὶ προσεύξωνται καὶ ζητήσωσι τὸ πρόσωπόν μου, καὶ ἀποστρέψωσιν ἀπὸ τῶν ὁδῶν αὐτῶν τῶν πονηρῶν, καὶ ἐγὼ εἰσακούσομαι ἐκ τοῦ οὐρανοῦ, καὶ ἵλεως ἔσομαι ταῖς ἁμαρτίαις αὐτῶν, καὶ ἰάσομαι τὴν γῆν αὐτῶν.
\vs{15}Καὶ νῦν οἱ ὀφθαλμοί μου ἔσονται ἀνεῳγμένοι, καὶ τὰ ὦτά μου ἐπήκοα τῇ προσευχῇ τοῦ τόπου τούτου.
\vs{16}Καὶ νῦν ἐξελεξάμην καὶ ἡγίακα τὸν οἶκον τοῦτον, τοῦ εἶναι ὄνομά μου ἐκεῖ ἕως αἰῶνος, καὶ ἔσονται οἱ ὀφθαλμοί μου καὶ ἡ καρδία μου ἐκεῖ πάσας τὰς ἡμέρας.

\vs{17}Καὶ σὺ ἐὰν πορευθῇς ἐναντίον μου ὡς Δαυὶδ ὁ πατήρ σου, καὶ ποιήσῃς κατὰ πάντα ἃ ἐνετειλάμην σοι, καὶ τὰ προστάγματά μου καὶ τὰ κρίματά μου φυλάξῃ,
\vs{18}καὶ ἀναστήσω τὸν θρόνον τῆς βασιλείας σου ὡς διεθέμην Δαυὶδ τῷ πατρί σου, λέγων, οὐκ ἐξαρθήσεταί σοι ἡγούμενος ἀνὴρ ἐν Ἰσραήλ.

\vs{19}Καὶ ἐὰν ἀποστρέψητε ὑμεῖς, καὶ ἐγκαταλείπητε τὰ προστάγματά μου καὶ τὰς ἐντολάς μου ἃς ἔδωκα ἐναντίον ὑμῶν, καὶ πορευθῆτε καὶ λατρεύσητε θεοῖς ἑτέροις καὶ προσκυνήσητε αὐτοῖς,
\vs{20}καὶ ἐξαρῶ ὑμᾶς ἀπὸ τῆς γῆς ἧς ἔδωκα αὐτοῖς· καὶ τὸν οἶκον τοῦτον ὃν ἡγίασα τῷ ὀνόματί μου ἀποστρέψω ἐκ προσώπου μου, καὶ δώσω αὐτὸν εἰς παραβολὴν καὶ εἰς διήγημα ἐν πᾶσι τοῖς ἔθνεσι.
\vs{21}Καὶ ὁ οἶκος οὗτος ὁ ὑψηλὸς πᾶς ὁ διαπορευόμενος αὐτὸν ἐκστήσεται, καὶ ἐρεῖ, Χάριν τίνος ἐποίησε Κύριος τῇ γῇ ταύτῃ καὶ τῷ οἴκῳ τούτῳ;
\vs{22}Καὶ ἐροῦσι, διότι ἐγκατέλιπον Κυρίον τὸν Θεὸν τῶν πατέρων αὐτῶν, τὸν ἐξαγαγόντα αὐτοὺς ἐκ γῆς Αἰγύπτου, καὶ ἀντελάβοντο θεῶν ἑτέρων, καὶ προσεκύνησαν αὐτοῖς, καὶ ἐδούλευσαν αὐτοῖς, καὶ διὰ τοῦτο ἐπήγαγεν ἐπʼ αὐτοὺς πᾶσαν τὴν κακίαν ταύτην.

\ch{8}
Καὶ ἐγένετο μετὰ εἴκοσι ἔτη ἐν οἷς ᾠκοδόμησε Σαλωμὼν τὸν οἶκον Κυρίου, καὶ τὸν οἶκον αὑτοῦ,
\vs{2}καὶ τὰς πόλεις ἃς ἔδωκε Χιρὰμ τῷ Σαλωμὼν, ᾠκοδόμησεν αὐτὰς Σαλωμὼν, καὶ κατῴκισεν ἐκεῖ τοὺς υἱοὺς Ἰσραήλ.

\vs{3}Καὶ ἦλθε Σαλωμὼν εἰς Βαισωβὰ, καὶ κατίσχυσεν αὐτήν.
\vs{4}Καὶ ᾠκοδόμησε τὴν Θοεδμὸρ ἐν τῇ ἐρήμῳ, καὶ πάσας τὰς πόλεις τὰς ὀχυρὰς ἃς ᾠκοδόμησεν ἐν Ἠμάθ.
\vs{5}Καὶ ᾠκοδόμησε τὴν Βαιθωρὼν τὴν ἄνω καὶ τὴν Βαιθωρὼν τὴν κάτω, πόλεις ὀχυράς· τείχη, πύλαι, καὶ μοχλοί·
\vs{6}Καὶ τὴν Βαλαὰθ, καὶ πάσας τὰς πόλεις τὰς ὀχυρὰς αἳ ἦσαν τῷ Σαλωμὼν, καὶ πάσας τὰς πόλεις τῶν ἁρμάτων, καὶ τὰς πόλεις τῶν ἱππέων· καὶ ὅσα ἐπεθύμησε Σαλωμὼν κατὰ τὴν ἐπιθυμίαν τοῦ οἰκοδομῆσαι ἐν Ἱερουσαλὴμ, καὶ ἐν τῷ Λιβάνῳ, καὶ ἐν πάσῃ τῇ βασιλείᾳ αὐτοῦ.

\vs{7}Πᾶς ὁ λαὸς ὁ καταλειφθεὶς ἀπὸ τοῦ Χετταίου, καὶ τοῦ Ἀμοῥῥαίου, καὶ τοῦ Φερεζαίου, καὶ τοῦ Εὐαίου, καὶ τοῦ Ἰεβουσαίου, οἳ οὐκ εἰσὶν ἐκ τοῦ Ἰσραὴλ,
\vs{8}ἀλλʼ ἦσαν ἐκ τῶν υἱῶν αὐτῶν τῶν καταλειφθέντων μετʼ αὐτοὺς ἐν τῇ γῇ, οὓς οὐκ ἐξωλόθρευσαν οἱ υἱοὶ Ἰσραὴλ, καὶ ἀνήγαγεν αὐτοὺς Σαλωμὼν εἰς φόρον ἕως τῆς ἡμέρας ταύτης.
\vs{9}Καὶ ἐκ τῶν υἱῶν Ἰσραὴλ οὐκ ἔδωκε Σαλωμὼν εἰς παῖδας τῇ βασιλείᾳ αὐτοῦ· ὅτι ἰδοὺ ἄνδρες πολεμισταὶ καὶ ἄρχοντες, καὶ οἱ δυνατοὶ καὶ ἄρχοντες ἁρμάτων καὶ ἱππέων.
\vs{10}Καὶ οὗτοι ἄρχοντες τῶν προστατῶν βασιλέως Σαλωμὼν, πεντήκοντα καὶ διακόσιοι ἐργοδιωκτοῦντες ἐν τῷ λαῷ.

\vs{11}Καὶ τὴν θυγατέρα Φαραὼ ἀνήγαγε Σαλωμὼν ἐκ πόλεως Δαυὶδ εἰς τὸν οἶκον ὃν ᾠκοδόμησεν αὐτῇ, ὅτι εἶπεν, οὐ κατοικήσει ἡ γυνή μου ἐν πόλει Δαυὶδ τοῦ βασιλέως Ἰσραὴλ, ὅτι ἅγιός ἐστιν οὗ εἰσῆλθεν ἐκεῖ κιβωτὸς Κυρίου.

\vs{12}Τότε ἀνήνεγκε Σαλωμὼν ὁλοκαυτώματα τῷ Κυρίῳ ἐπὶ τὸ θυσιαστήριον, ὃ ᾠκοδόμησε Κυρίῳ ἀπέναντι τοῦ ναοῦ
\vs{13}κατὰ τὸν λόγον ἡμέρας ἐν ἡμέρᾳ, τοῦ ἀναφέρειν κατὰ τὰς ἐντολὰς Μωυσῆ ἐν τοῖς σαββάτοις, καὶ ἐν τοῖς μησὶ, καὶ ἐν ταῖς ἑορταῖς, τρεῖς καιροὺς τοῦ ἐνιαυτοῦ, ἐν τῇ ἑορτῇ τῶν ἀζύμων, καὶ ἐν τῇ ἐορτῇ τῶν ἐβδομάδων, καὶ ἐν τῇ ἑορτῇ τῶν σκηνῶν.
\vs{14}Καὶ ἔστησε κατὰ τὴν κρίσιν Δαυὶδ τοῦ πατοὸς αὐτοῦ τὰς διαιρέσεις τῶν ἱερέων, καὶ κατὰ τὰς λειτουργίας αὐτῶν· καὶ οἱ Λευῖται ἐπὶ τὰς φυλακὰς αὐτῶν, τοῦ αἰνεῖν καὶ λειτουργεῖν κατέναντι τῶν ἱερέων κατὰ τὸν λόγον ἡμέρας ἐν τῇ ἡμερᾳ· καὶ οἱ πυλωροὶ κατὰ τὰς διαιρέσεις αὐτῶν εἰς πύλην καὶ πύλην, ὅτι οὕτως ἐντολαὶ Δαυὶδ ἀνθρώπου τοῦ Θεοῦ·
\vs{15}Οὐ παρῆλθον τὰς ἐντολὰς τοῦ βασιλέως περὶ τῶν ἱερέων καὶ τῶν Λευιτῶν εἰς πάντα λόγον, καὶ εἰς τοὺς θησαυρούς.
\vs{16}Καὶ ἡτοιμάσθη πᾶσα ἡ ἐργασία ἀφʼ ἧς ἡμέρας ἐθεμελιώθη ἕως οὗ ἐτελείωσε Σαλωμὼν τὸν οἶκον Κυρίου.

\vs{17}Τότε ᾤχετο Σαλωμὼν εἰς Γασιὼν Γαβὲρ, καὶ εἰς τὴν Αἰλὰθ τὴν παραθαλασσίαν ἐν γῇ Ἰδουμαίᾳ.
\vs{18}Καὶ ἀπέστειλε Χιρὰμ ἐν χειρὶ παίδων αὐτοῦ πλοῖα καὶ παῖδας εἰδότας θάλασσαν, καὶ ᾤχοντο μετὰ τῶν παίδων Σαλωμὼν εἰς Σωφιρὰ, καὶ ἔλαβον ἐκεῖθεν τὰ τετρακόσια καὶ πεντήκοντα τάλαντα χρυσίου, καὶ ἦλθον πρὸς τὸν βασιλέα Σαλωμών.

\ch{9}
Καὶ βασίλισσα Σαβὰ ἤκουσε τὸ ὄνομα Σαλωμὼν, καὶ ἦλθε τοῦ πειρᾶσαι Σαλωμὼν ἐν αἰνίγμασιν εἰς Ἱερουσαλὴμ ἐν δυνάμει βαρείᾳ σφόδρα, καὶ κάμηλοι αἴρουσαι ἀρώματα εἰς πλῆθος, καὶ χρυσίον, καὶ λίθον τίμιον· καὶ ἦλθε πρὸς Σαλωμὼν, καὶ ἐλάλησε πρὸς αὐτὸν πάντα ὅσα ἦν ἐν τῇ ψυχῇ αὐτῆς.
\vs{2}Καὶ ἀνήγγειλεν αὐτῇ Σαλωμὼν πάντας τοὺς λόγους αὐτῆς, καὶ οὐ παρῆλθε λόγος ἀπὸ Σαλωμὼν ὃν οὐκ ἀπήγγειλεν αὐτῇ.

\vs{3}Καὶ εἶδε βασίλισσα Σαβὰ τὴν σοφίαν Σαλωμὼν καὶ τὸν οἶκον ὃν ᾠκοδόμησε,
\vs{4}καὶ τὰ βρώματα τῶν τραπεζῶν, καὶ καθέδραν παίδων αὐτοῦ, καὶ στάσιν λειτουργῶν αὐτοῦ, καὶ ἱματισμὸν αὐτῶν, καὶ οἰνοχόους αὐτοῦ, καὶ στολισμὸν αὐτῶν, καὶ τὰ ὁλοκαυτώματα ἃ ἀνέφερεν ἐν οἴκῳ Κυρίου, καὶ ἐξ ἑαυτῆς ἐγένετο.
\vs{5}Καὶ εἶπε πρὸς τὸν βασιλέα, ἀληθινὸς ὁ λόγος ὃν ἤκουσα ἐν τῇ γῇ μου περὶ τῶν λόγων σου, καὶ περὶ τῆς σοφίας σου.
\vs{6}Καὶ οὐκ ἐπίστευσα τοῖς λόγοις ἕως οὗ ἦλθον καὶ εἶδον οἱ ὀφθαλμοί μου, καὶ ἰδοὺ οὐκ ἀπηγγέλη μοι ἥμισυ τοῦ πλήθους τῆς σοφίας σου· προσέθηκας ἐπὶ τὴν ἀκοὴν ἣν ἤκουσα.
\vs{7}Μακάριοι οἱ ἄνδρες σου, μακάριοι οἱ παῖδες οὗτοι οἱ παρεστηκότες σοι διαπαντὸς καὶ ἀκούοντες τὴς σοφίαν σου.
\vs{8}Ἔστω Κύριος ὁ Θεός σου εὐλογημένος, ὃς ἠθέλησεν ἐν σοὶ τοῦ δοῦναί σε ἐπὶ θρόνον αὐτοῦ εἰς βασιλέα Κυρίῳ Θεῷ σου· ἐν τῷ ἀγαπῆσαι Κύριον τὸν Θεόν σου τὸν Ἰσραὴλ τοῦ στῆσαι αὐτὸν εἰς αἰῶνα, καὶ ἔδωκέ σε ἐπʼ αὐτοὺς εἰς βασιλέα τοῦ ποιῆσαι κρίμα καὶ δικαιοσύνην.
\vs{9}Καὶ ἔδωκε τῷ βασιλεῖ ἑκατὸν εἴκοσι τάλαντα χρυσίου, καὶ ἀρώματα εἰς πλῆθος πολὺ, καὶ λίθον τίμιον· καὶ οὐκ ἦν κατὰ τὰ ἀρώματα ἐκεῖνα ἃ ἔδωκε βασίλισσα Σαβὰ τῷ βασιλεῖ Σαλωμών.

\vs{10}Καὶ οἱ παῖδες Σαλωμὼν καὶ οἱ παῖδες Χιρὰμ ἔφερον χρυσίον τῷ Σαλωμὼν ἐκ Σουφὶρ, καὶ ξύλα πεύκινα, καὶ λίθον τίμιον.
\vs{11}Καὶ ἐποίησεν ὁ βασιλεὺς τὰ ξύλα τὰ πεύκινα ἀναβάσεις τῷ οἴκῳ Κυρίου, καὶ τῷ οἴκῳ τοῦ βασιλέως, καὶ κιθάρας καὶ νάβλας τοῖς ᾠδοῖς· καὶ οὐκ ὤφθησαν τοιαῦτα ἔμπροσθεν ἐν γῇ Ἰούδα.
\vs{12}Καὶ ὁ βασιλεὺς Σαλωμὼν ἔδωκε τῇ βασιλίσσῃ Σαβὰ πάντα τὰ θελήματα αὐτῆς ἃ ᾔτησεν, ἐκτὸς πάντων ὧν ἤνεγκε τῷ βασιλεῖ Σαλωμὼν· καὶ ἀπέστρεψεν εἰς τὴν γῆν αὐτῆς.

\vs{13}Καὶ ἦν ὁ σταθμὸς τοῦ χρυσίου τοῦ ἐνεχθέντος τῷ Σαλωμὼν ἐν ἐνιαυτῷ ἑνὶ, ἑξακόσια ἑξηκονταὲξ τάλαντα χρυσίου,
\vs{14}πλὴν τῶν ἀνδρῶν τῶν ὑποτεταγμένων καὶ τῶν ἐμπορευομένων ὧν ἔφερον· καὶ πάντων τῶν βασιλέων τῆς Ἀραβίας καὶ σατραπῶν τῆς γῆς, πάντες ἔφερον χρυσίον καὶ ἀργύριον τῷ βασιλεῖ Σαλωμών.
\vs{15}Καὶ ἐποίησεν ὁ βασιλεὺς Σαλωμὼν διακοσίους θυρεοὺς χρυσοῦς ἐλατοὺς, ἑξακόσιοι χρυσοῖ καθαροὶ ἐπῆσαν ἐπὶ τὸν ἕνα θυρεόν.
\vs{16}Καὶ τριακοσίας ἀσπίδας ἐλατὰς χρυσᾶς, τριακοσίων χρυσῶν ἀνεφέρετο ἐπὶ τὴν ἀσπίδα ἑκάστην, καὶ ἔδωκεν αὐτὰς ὁ βασιλεὺς ἐν οἴκῳ δρυμοῦ τοῦ Λιβάνου.

\vs{17}Καὶ ἐποίησεν ὁ βασιλεὺς θρόνον ἐλεφαντίνων ὀδόντων μέγαν, καὶ κατεχρύσωσεν αὐτὸν χρυσίῳ δοκίμῳ.
\vs{18}Καὶ ἓξ ἀναβαθμοὶ τῷ θρόνῳ ἐνδεδεμένοι χρυσίῳ, καὶ ἀγκῶνες ἔνθεν καὶ ἔνθεν ἐπὶ τοῦ θρόνου τῆς καθέδρας, καὶ δύο λέοντες ἑστηκότες παρὰ τοὺς ἀγκῶνας,
\vs{19}καὶ δώδεκα λέοντες ἑστηκότες ἐκεῖ ἐπὶ τῶν ἓξ ἀναβαθμῶν ἔνθεν καὶ ἔνθεν· οὐκ ἐγενήθη οὕτως ἐν πάσῃ τῇ βασιλείᾳ.

\vs{20}Καὶ πάντα τὰ σκεύη τοῦ βασιλέως Σαλωμὼν χρυσίου, καὶ πάντα τὰ σκεύη οἴκου δρυμοῦ τοῦ Λιβάνου χρυσίῳ κατειλημμένα· οὐκ ἦν ἀργύριον λογιζόμεμον ἐν ἡμέραις Σαλωμὼν εἰς οὐθέν.
\vs{21}Ὅτι ναῦς τῷ βασιλεῖ ἐπορεύετο εἰς Θαρσεῖς μετὰ τῶν παίδων Χιρὰμ, ἅπαξ διὰ τριῶν ἐτῶν ἤρχετο πλοῖα ἐκ Θαρσεὶς τῷ βασιλεῖ γέμοντα χρυσίου καὶ ἀργυρίου, καὶ ὀδόντων ἐλεφαντίνων, καὶ πιθήκων.

\vs{22}Καὶ ἐμεγαλύνθη Σαλωμὼν ὑπὲρ πάντας τοὺς βασιλεῖς καὶ πλούτῳ καὶ σοφίᾳ.
\vs{23}Καὶ πάντες οἱ βασιλεῖς τῆς γῆς ἐζήτουν τὸ πρόσωπον Σαλωμὼν ἀκοῦσαι τῆς σοφίας αὐτοῦ, ἧς ἔδωκεν ὁ Θεὸς ἐν καρδίᾳ αὐτοῦ.
\vs{24}Καὶ αὐτοὶ ἔφερον ἕκαστος τὰ δῶρα αὐτοῦ, σκεύη ἀργυρᾶ καὶ σκεύη χρυσᾶ, καὶ ἱματισμὸν, στακτὴν καὶ ἡδύσματα, ἵππους καὶ ἡμιόνους τὸ κατʼ ἐνιαυτὸν ἐνιαυτόν.

\vs{25}Καὶ ἦσαν τῷ Σαλωμὼν τέσσαρες χιλιάδες θήλειαι ἵπποι εἰς ἅρματα, καὶ δώδεκα χιλιάδες ἱππέων, καὶ ἔθετο αὐτοὺς ἐν πόλεσι τῶν ἁρμάτων, καὶ μετὰ τοῦ βασιλέως ἐν Ἱερουσαλήμ.
\vs{26}Καὶ ἦν ἡγούμενος πάντων τῶν βασιλέων ἀπὸ τοῦ ποταμοῦ καὶ ἕως γῆς ἀλλοφύλων, καὶ ἕως ὁρίων Αἰγύπτου.
\vs{27}Καὶ ἔδωκεν ὁ βασιλεὺς τὸ χρυσίον καὶ τὸ ἀργύριον ἐν Ἱερουσαλὴμ ὡς λίθους, καὶ τὰς κέρδους ὡς συκαμίνους τὰς ἐν τῇ πεδινῇ εἰς πλῆθος.
\vs{28}Καὶ ἡ ἔξοδος τῶν ἵππων ἐξ Αἰγύπτου τῷ Σαλωμὼν καὶ ἐκ πάσης τῆς γῆς.

\vs{29}Καὶ οἱ κατάλοιποι λόγοι Σαλωμὼν οἱ πρῶτοι καὶ οἱ ἔσχατοι, ἰδοὺ οὗτοι γεγραμμένοι ἐπὶ τῶν λόγων Νάθαν τοῦ προφήτου, καὶ ἐπὶ τῶν λόγων Ἀχία τοῦ Σηλωνίτου, καὶ ἐν ταῖς ὁράσεσιν Ἰωὴλ τοῦ ὁρῶντος περὶ Ἱεροβοὰμ υἱοῦ Ναβάτ.
\vs{30}Καὶ ἐβασίλευσε Σαλωμὼν ἐπὶ πάντα Ἰσραὴλ τεσσαράκοντα ἔτη.
\vs{31}Καὶ ἐκοιμήθη Σαλωμὼν, καὶ ἔθαψαν αὐτὸν ἐν πόλει Δαυὶδ τοῦ πατρὸς αὐτοῦ, καὶ ἐβασιλευσε Ῥοβοὰμ υἱὸς αὐτοῦ ἀντʼ αὐτοῦ.

\ch{10}
Καὶ ἦλθε Ῥοβοὰμ εἰς Συχὲμ, ὅτι εἰς Συχὲμ ἤρχετο πᾶς Ἰσραὴλ βασιλεῦσαι αὐτόν.

\vs{2}Καὶ ἐγένετο ὡς ἤκουσεν Ἱεροβοὰμ υἱὸς Ναβὰτ, καὶ αὐτὸς ἐν Αἰγύπτῳ, ὡς ἔφυγεν ἀπὸ προσώπου Σαλωμὼν τοῦ βασιλέως, καὶ κατῴκησεν Ἱεροβοὰμ ἐν Αἰγύπτῳ, καὶ ἀπέστρεψεν Ἱεροβοὰμ ἐξ Αἰγύπτου.
\vs{3}Καὶ ἀπέστειλαν καὶ ἐκάλεσαν αὐτόν· καὶ ἦλθεν Ἱεροβοὰμ καὶ πᾶσα ἡ ἐκκλησία πρὸς Ῥοβοὰμ, λέγοντες,
\vs{4}ὁ πατήρ σου ἐσκλήρυνε τὸν ζυγὸν ἡμῶν, καὶ νῦν ἄφες ἀπὸ τῆς δουλείας τοῦ πατρός σου τῆς σκληρᾶς, καὶ ἀπὸ τοῦ ζυγοῦ αὐτοῦ τοῦ βαρέος, οὗ ἔδωκεν ἐφʼ ἡμᾶς, καὶ δουλεύσομέν σοι.
\vs{5}Καὶ εἶπεν αὐτοῖς, πορεύεσθε ἕως τριῶν ἡμερῶν, καὶ ἔρχεσθε πρὸς μέ· καὶ ἀπῆλθεν ὁ λαός.

\vs{6}Καὶ συνήγαγεν ὁ βασιλεὺς Ῥοβοὰμ τοὺς πρεσβυτέρους τοὺς ἑστηκότας ἐναντίον τοῦ Σαλωμὼν τοῦ πατρὸς αὐτοῦ ἐν τῷ ζῇν αὐτὸν, λέγων, πῶς ὑμεῖς βουλεύεσθε τοῦ ἀποκριθῆναι τῷ λαῷ τούτῳ λόγον;
\vs{7}Καὶ ἐλάλησαν αὐτῷ, λέγοντες, ἐὰν ἐν τῇ σήμερον γένῃ εἰς ἀγαθὸν τῷ λαῷ τούτῳ, καὶ εὐδοκήσῃς, καὶ λαλήσῃς αὐτοῖς λόγους ἀγαθοὺς, καὶ ἔσονταί σοι παῖδες πάσας τὰς ἡμέρας.
\vs{8}Καὶ κατέλιπε τὴν βουλὴν τῶν πρεσβυτέρων, οἳ συνεβουλεύσαντο αὐτῷ· καὶ συνεβουλεύσατο μετὰ τῶν παιδαρίων τῶν συνεκτραφέντων μετʼ αὐτοῦ τῶν ἑστηκότων ἐναντίον αὐτοῦ.
\vs{9}Καὶ εἶπεν αὐτοῖς, τί ὑμεῖς βουλεύεσθε, καὶ ἀποκριθήσομαι λόγον τῷ λαῷ τούτῳ, οἳ ἐλάλησαν πρὸς μὲ, λέγοντες, ἄνες ἀπὸ τοῦ ζυγοῦ οὗ ἔδωκεν ὁ πατήρ σου ἐφʼ ἡμᾶς;
\vs{10}Καὶ ἐλάλησαν αὐτῷ τὰ παιδάρια τὰ ἐκτραφέντα μετʼ αὐτοῦ, λέγοντες, οὕτως λαλήσεις τῷ λαῷ τῷ λαλήσαντι πρὸς σὲ, λέγων, ὁ πατήρ σου ἐβάρυνε τὸν ζυγὸν ἡμῶν, καὶ σὺ ἄφες ἀφʼ ἡμῶν· οὕτως ἐρεῖς, ὁ μικρὸς δάκτυλός μου παχύτερος τῆς ὀσφύος τοῦ πατρός μου.
\vs{11}Καὶ νῦν ὁ πατήρ μου ἐπαίδευσεν ὑμᾶς ζυγῷ βαρεῖ, κᾀγὼ προσθήσω ἐπὶ τὸν ζυγὸν ὑμῶν· ὁ πατήρ μου ἐπαίδευσεν ὑμᾶς ἐν μάστιξι, κᾀγὼ παιδεύσω ὑμᾶς ἐν σκορπίοις.

\vs{12}Καὶ ἦλθεν Ἱεροβοὰμ καὶ πᾶς ὁ λαὸς πρὸς Ῥοβοὰμ τῇ ἡμέρᾳ τῇ τρίτῃ, ὡς ἐλάλησεν ὁ βασιλεὺς, λέγων, ἐπιστρέψατε πρὸς μὲ ἐν τῇ ἡμέρᾳ τῇ τρίτῃ.
\vs{13}Καὶ ἀπεκρίθη ὁ βασιλεὺς σκληρὰ, καὶ ἐγκατέλιπεν ὁ βασιλεὺς Ῥοβοὰμ τὴν βουλὴν τῶν πρεσβυτέρων,
\vs{14}καὶ ἐλάλησε πρὸς αὐτοὺς κατὰ τὴν βουλὴν τῶν νεωτέρων, λέγων, ὁ πατήρ μου ἐβάρυνε τὸν ζυγὸν ὑμῶν, καὶ ἐγὼ προσθήσω ἐπʼ αὐτόν· ὁ πατήρ μου ἐπαίδευσεν ὑμᾶς ἐν μάστιξι, καὶ ἐγὼ παιδεύσω ὑμᾶς ἐν σκορπίοις.

\vs{15}Καὶ οὐκ ἤκουσεν ὁ βασιλεὺς τοῦ λαοῦ, ὅτι ἦν μεταστροφὴ παρὰ τοῦ Θεοῦ, λέγων, ἀνέστησε Κύριος τὸν λόγον αὐτοῦ, ὃν ἐλάλησεν ἐν χειρὶ Ἀχία τοῦ Σηλωνίτου περὶ Ἱεροβοὰμ υἱοῦ Ναβὰτ
\vs{16}καὶ παντὸς Ἰσραὴλ, ὅτι οὐκ ἤκουσεν ὁ βασιλεὺς αὐτῶν· καὶ ἀπεκρίθη ὁ λαὸς πρὸς τὸν βασιλέα, λέγων, τίς ἡμῶν ἡ μερὶς ἐν Δαυὶδ καὶ κληρονομία ἐν υἱῷ Ἰεσσαί; εἰς τὰ σκηνώματά σου, Ἰσραήλ· νῦν βλέπε τὸν οἶκόν σου, Δαυίδ. καὶ ἐπορεύθη πᾶς Ἰσραὴλ εἰς τὰ σκηνώματα αὐτοῦ.
\vs{17}Καὶ ἄνδρες Ἰσραὴλ καὶ οἱ κατοικοῦντες ἐν πόλεσιν Ἰούδα, καὶ ἐβασίλευσαν ἐπʼ αὐτῶν Ῥοβοάμ.

\vs{18}Καὶ ἀπέστειλεν ἐπʼ αὐτοὺς Ῥοβοὰμ ὁ βασιλεὺς τὸν Ἀδωνιρὰμ τὸν ἐπὶ τοῦ φόρου, καὶ ἐλιθοβόλησαν αὐτὸν οἱ υἱοὶ Ἰσραὴλ λίθοις, καὶ ἀπέθανε· καὶ ὁ βασιλεὺς Ῥοβοὰμ ἔσπευσε τοῦ ἀναβῆναι εἰς τὸ ἅρμα, τοῦ φυγεῖν εἰς Ἱερουσαλήμ.
\vs{19}Καὶ ἠθέτησεν Ἰσραὴλ ἐν τῷ οἴκῳ Δαυὶδ ἕως τῆς ἡμέρας ταύτης.

\ch{11}
Καὶ ἦλθε Ῥοβοὰμ εἰς Ἱερουσαλὴμ, καὶ ἐξεκκλησίασε τὸν Ἰούδαν καὶ Βενιαμὶν ἑκατὸν ὀγδοήκοντα χιλιάδας νεανίσκων ποιούντων πόλεμον, καὶ ἐπολέμει πρὸς Ἰσραὴλ τοῦ ἐπιστρέψαι τὴν βασιλείαν τῷ Ῥοβοάμ.
\vs{2}Καὶ ἐγένετο λόγος Κυρίου πρὸς Σαμαίαν ἄνθρωπον τοῦ Θεοῦ, λέγων,
\vs{3}εἶπον πρὸς Ῥοβοὰμ τὸν τοῦ Σαλωμὼν καὶ πάντα Ἰούδαν καὶ Βενιαμὶν, λέγων,
\vs{4}τάδε λέγει Κύριος, οὐκ ἀναβήσεσθε, καὶ οὐ πολεμήσεσθε πρὸς τοὺς ἀδελφοὺς ὑμῶν· ἀποστρέφετε ἕκαστος εἰς τὸν οἶκον αὐτοῦ, ὅτι παρʼ ἐμοῦ ἐγένετο τὸ ῥῆμα τοῦτο· καὶ ἐπήκουσαν τοῦ λόγου Κυρίου, καὶ ἀπεστράφησαν τοῦ μὴ πορευθῆναι ἐπὶ Ἱεροβοάμ.

\vs{5}Καὶ κατῴκησε Ῥοβοὰμ εἰς Ἱερουσαλὴμ, καὶ ᾠκοδόμησε πόλεις τειχήρεις ἐν τῇ Ἰουδαίᾳ.
\vs{6}Καὶ ᾠκοδόμησε τὴν Βηθλεὲμ, καὶ Αἰτὰν, καὶ Θεκωὲ,
\vs{7}καὶ Βαιθσουρὰ, καὶ τὴν Σοχὼθ, καὶ τὴν Ὀδολλὰμ,
\vs{8}καὶ τὴν Γὲθ, καὶ τὴν Μαρισὰν, καὶ τὴν Ζὶφ,
\vs{9}καὶ τὴν Ἀδωραὶ, καὶ Λαχὶς, καὶ τὴν Ἀζηκά,
\vs{10}καὶ τὴν Σαραὰ, καὶ τὴν Αἰλὼμ, καὶ τὴν Χεβρὼν, ἥ ἐστι τοῦ Ἰούδα καὶ Βενιαμὶν, πόλεις τειχήρεις.
\vs{11}Καὶ ὠχύρωσεν αὐτὰς τειχήρεις, καὶ ἔδωκεν ἐν αὐταῖς ἡγουμένους, καὶ παραθέσεις βρωμάτων, ἔλαιον καὶ οἶνον,
\vs{12}κατὰ πόλιν καὶ κατὰ πόλιν θυρεοὺς καὶ δόρατα, καὶ κατίσχυσεν αὐτὰς εἰς πλῆθος σφόδρα· καὶ ἦσαν αὐτῷ Ἰούδα καὶ Βενιαμίν.

\vs{13}Καὶ οἱ ἱερεῖς καὶ οἱ Λευῖται οἳ ἦσαν ἐν παντὶ Ἰσραὴλ συνήχθησαν πρὸς αὐτὸν ἐκ πάντων τῶν ὁρίων·
\vs{14}Ὅτι ἐγκατέλιπον οἱ Λευῖται τὰ σκηνώματα τῆς κατασχέσεως αὐτῶν, καὶ ἐπορεύθησαν πρὸς Ἰούδα εἰς Ἱερουσαλὴμ, ὅτι ἐξέβαλεν αὐτοὺς Ἱεροβοὰμ καὶ οἱ υἱοὶ αὐτοῦ μὴ λειτουργεῖν Κυρίῳ,
\vs{15}καὶ κατέστησεν ἑαυτῷ ἱερεῖς τῶν ὑψηλῶν καὶ τοῖς εἰδώλοις καὶ τοῖς ματαίοις καὶ τοῖς μόσχοις, ἃ ἐποίησεν Ἱεροβοάμ.
\vs{16}Καὶ ἐξέβαλεν αὐτοὺς ἀπὸ φυλῶν Ἰσραὴλ, οἳ ἔδωκαν καρδίαν αὐτῶν τοῦ ζητῆσαι Κύριον Θεὸν Ἰσραήλ· καὶ ἦλθον εἰς Ἱερουσαλὴμ θῦσαι Κυρίῳ Θεῷ τῶν πατέρων αὐτῶν.
\vs{17}Καὶ κατίσχυσαν τὴν βασιλείαν Ἰούδα· καὶ κατίσχυσε Ῥοβοὰμ τὸν τοῦ Σαλωμὼν εἰς ἔτη τρία, ὅτι ἐπορεύθη ἐν ταῖς ὁδοῖς Δαυὶδ καὶ Σαλωμὼν ἔτη τρία.

\vs{18}Καὶ ἔλαβεν ἑαυτῷ Ῥοβοὰμ γυναῖκα τὴν Μοολὰθ θυγατέρα Ἱεριμοὺθ υἱοῦ Δαυὶδ, καὶ Ἀβιγαίαν θυγατέρα Ἑλιὰβ τοῦ Ἰεσσαί.
\vs{19}Καὶ ἔτεκεν αὐτῷ υἱοὺς, τὸν Ἰεοὺς, καὶ τὸν Σαμαρία, καὶ τὸν Ζαάμ.
\vs{20}Καὶ μετὰ ταῦτα ἔλαβεν ἑαυτῷ τὴν Μααχὰ θυγατέρα Ἀβεσσαλὼμ, καὶ ἔτεκεν αὐτῷ τὸν Ἀβιὰ, καὶ τὸν Ἰετθὶ, καὶ τὸν Ζηζὰ, καὶ τὸν Σαλημώθ.
\vs{21}Καὶ ἠγάπησε Ῥοβοὰμ τὴν Μααχά θυγατέρα Ἀβεσσαλὼμ ὑπὲρ πάσας τὰς γυναῖκας αὐτοῦ καὶ τὰς παλλακὰς αὐτοῦ, ὅτι γυναῖκας δεκαοκτὼ εἶχε καὶ παλλακὰς ἑξήκοντα· καὶ ἐγέννησεν υἱοὺς εἴκοσι καὶ ὀκτὼ, καὶ θυγατέρας ἑξήκοντα.
\vs{22}Καὶ κατέστησεν εἰς ἄρχοντα Ἀβιὰ τὸν τῆς Μααχὰ εἰς ἡγούμενον ἐν τοῖς ἀδελφοῖς αὐτοῦ, ὅτι βασιλεῦσαι διενοεῖτο αὐτόν.
\vs{23}Καὶ ηὐξήθη παρὰ πάντας τοὺς υἱοὺς αὐτοῦ ἐν πᾶσι τοῖς ὁρίοις Ἰούδα καὶ Βενιαμὶν, καὶ ἐν ταῖς πόλεσι ταῖς ὀχυραῖς, καὶ ἔδωκεν αὐταῖς τροφὰς πλῆθος πολὺ, καὶ ᾐτήσατο πλῆθος γυναικῶν.

\ch{12}
Καὶ ἐγένετο ὡς ἡτοιμάσθη ἡ βασιλεία Ῥοβοὰμ, καὶ ὡς κατεκρατήθη, ἐγκατέλιπε τὰς ἐντολὰς Κυρίου, καὶ πᾶς Ἰσραὴλ μετʼ αὐτοῦ.

\vs{2}Καὶ ἐγένετο ἐν τῷ ἔτει τῷ πέμπτῳ τῆς βασιλείας Ῥοβοὰμ, ἀνέβη Σουσακὶμ βασιλεὺς Αἰγύπτου ἐπὶ Ἱερουσαλὴμ, ὅτι ἥμαρτον ἐναντίον Κυρίου,
\vs{3}ἐν χιλίοις καὶ διακοσίοις ἅρμασι καὶ ἑξήκοντα χιλιάσιν ἵππων, καὶ οὐκ ἦν ἀριθμὸς τοῦ πλήθους τοῦ ἐλθόντος μετʼ αὐτοῦ ἐξ Αἰγύπτου, Λίβυες, Τρωγοδύται, καὶ Αἰθίοπες.
\vs{4}Καὶ κατεκράτησαν τῶν πόλεων τῶν ὀχυρῶν, αἳ ἦσαν ἐν Ἰούδα, καὶ ἦλθον εἰς Ἱερουσαλήμ.

\vs{5}Καὶ Σαμαίας ὁ προφήτης ἦλθε πρὸς Ῥοβοὰμ, καὶ πρὸς τοὺς ἄρχοντας Ἰούδα τοὺς συναχθέντας εἰς Ἱερουσαλὴμ ἀπὸ προσώπου Σουσακὶμ, καὶ εἶπεν αὐτοῖς, οὕτως εἶπε Κύριος, ὑμεῖς ἐγκατελίπετέ με, καὶ ἐγὼ ἐγκαταλείψω ὑμᾶς ἐν χειρὶ Σουσακίμ.
\vs{6}Καὶ ᾐσχύνθησαν οἱ ἄρχοντες Ἰσραὴλ καὶ ὁ βασιλεὺς, καὶ εἶπαν, δίκαιος ὁ κύριος.
\vs{7}Καὶ ἐν τῷ ἰδεῖν Κύριον ὅτι ἐνετράπησαν, καὶ ἐγένετο λόγος Κυρίου πρὸς Σαμαίαν, λέγων, ἐνετράπησαν, οὐ καταφθερῶ αὐτοὺς, καὶ δώσω αὐτοὺς ὡς μικρὸν εἰς σωτηρίαν, καὶ οὐ μὴ στάξῃ ὁ θυμός μου ἐν Ἱερουσαλὴμ·
\vs{8}ὅτι ἔσονται εἰς παῖδας, καὶ γνώσονται τὴν δουλείαν μου, καὶ τὴν δουλείαν τῆς βασιλείας τῆς γῆς.

\vs{9}Καὶ ἀνέβη Σουσακὶμ βασιλεὺς Αἰγύπτου ἐπὶ Ἱερουσαλὴμ, καὶ ἔλαβε τοὺς θησαυροὺς τοὺς ἐν οἴκῳ Κυρίου, καὶ τοὺς θησαυροὺς τοὺς ἐν οἴκῳ τοῦ βασιλέως, τὰ πάντα ἔλαβε· καὶ ἔλαβε τοὺς θυρεοὺς τοὺς χρυσοῦς οὓς ἐποίησε Σαλωμών.
\vs{10}Καὶ ἐποίησεν ὁ βασιλεὺς Ῥοβοὰμ θυρεοὺς χαλκοῦς ἀντʼ αὐτῶν· καὶ κατέστησεν ἐπʼ αὐτὸν Σουσακὶμ ἄρχοντας παρατρεχόντων, τοὺς φυλάσσοντας τὸν πυλῶνα τοῦ βασιλέως.
\vs{11}Καὶ ἐγένετο ἐν τῷ εἰσελθεῖν τὸν βασιλέα εἰς οἶκον Κυρίου, εἰσεπορεύοντο οἱ φυλάσσοντες, καὶ οἱ παρατρέχοντες, καὶ οἱ ἐπιστρέφοντες εἰς ἀπάντησιν τῶν παρατρεχόντων.
\vs{12}Καὶ ἐν τῷ ἐντραπῆναι αὐτὸν, ἀπεστράφη ἀπʼ αὐτοῦ ὀργὴ Κυρίου, καὶ οὐκ εἰς καταφθορὰν εἰς τέλος· καὶ γὰρ ἐν Ἰούδα ἦσαν λόγοι ἀγαθοί.

\vs{13}Καὶ κατίσχυσεν ὁ βασιλεὺς Ῥοβοὰμ ἐν Ἱερουσαλὴμ, καὶ ἐβασίλευσε· καὶ τεσσαράκοντα καὶ ἑνὸς ἐτῶν Ῥοβοὰμ ἐν τῷ βασιλεῦσαι αὐτὸν, καὶ ἑπτακαίδεκα ἔτη ἐβασίλευσεν ἐν Ἱερουσαλὴμ, ἐν τῇ πόλει ᾗ ἐξελέξατο Κύριος ἐπονομάσαι τὸ ὄνομα αὐτοῦ ἐκεῖ ἐκ πασῶν φυλῶν υἱῶν Ἰσραὴλ, καὶ τὸ ὄνομα τῆς μητρὸς αὐτοῦ Νοομμὰ ἡ Ἀμμανίτις.
\vs{14}Καὶ ἐποίησε τὸ πονηρὸν, ὅτι οὐ κατεύθυνε τὴν καρδίαν αὐτοῦ ἐκζητῆσαι τὸν κύριον.

\vs{15}Καὶ λόγοι Ῥοβοὰμ οἱ πρῶτοι καὶ ἔσχατοι οὐκ ἰδοὺ γεγραμμένοι ἐν τοῖς λόγοις Σαμαία τοῦ προφήτου, καὶ Ἀδδὼ τοῦ ὁρῶντος, καὶ πράξεις αὐτοῦ; καὶ ἐπολέμησε Ῥοβοὰμ τὸν Ἱεροβοὰμ πάσας τὰς ἡμέρας.
\vs{16}Καὶ ἀπέθανε Ῥοβοὰμ μετὰ τῶν πατέρων αὐτοῦ, καὶ ἐτάφη ἐν πόλει Δαυὶδ, καὶ ἐβασίλευσεν Ἀβιὰ υἱὸς αὐτοῦ ἀντʼ αὐτοῦ.

\ch{13}
Ἐν τῷ ὀκτωκαιδεκάτῳ ἔτει τῆς βασιλείας Ἱεροβοὰμ ἐβασίλευσεν Ἀβιὰ ἐπὶ Ἰούδαν.
\vs{2}Τρία ἔτη ἐβασίλευσεν ἐν Ἱερουσαλὴμ, καὶ ὄνομα τῇ μητρὶ αὐτοῦ Μααχὰ, θυγάτηρ Οὐριὴλ ἀπὸ Γαβαών.

Καὶ πολεμος ἦν ἀναμέσον Ἀβιὰ καὶ ἀναμέσον Ἱεροβοάμ.
\vs{3}Καὶ παρετάξατο Ἀβιὰ ἐν δυνάμει πολεμισταῖς δυνάμεως τετρακοσίαις χιλιάσιν ἀνδρῶν δυνατῶν· καὶ Ἱεροβοὰμ παρετάξατο πρὸς αὐτὸν πόλεμον ἐν ὀκτακοσίαις χιλιάσι, δυνατοὶ πολεμισταὶ δυνάμεως.

\vs{4}Καὶ ἀνέστη Ἀβιὰ ἀπὸ τοῦ ὄρους Σομόρων, ὅ ἐστιν ἐν τῷ ὄρει Ἐφραὶμ, καὶ εἶπεν, ἀκούσατε Ἱεροβοὰμ καὶ πᾶς Ἰσραήλ·
\vs{5}Οὐχ ὑμῖν γνῶναι ὅτι Κύριος ὁ Θεὸς Ἰσραὴλ ἔδωκε βασιλέα ἐπὶ τὸν Ἰσραὴλ εἰς τὸν αἰῶνα τῷ Δαυὶδ καὶ τοῖς υἱοῖς αὐτοῦ διαθήκῃ ἁλός;
\vs{6}Καὶ ἀνέστη Ἱεροβοὰμ ὁ τοῦ Ναβὰτ ὁ παῖς Σαλωμὼν τοῦ Δαυὶδ, καὶ ἀπέστη ἀπὸ τοῦ κυρίου αὐτοῦ·
\vs{7}Καὶ συνήχθησαν πρὸς αὐτὸν ἄνδρες λοιμοὶ υἱοὶ παράνομοι, καὶ ἀνέστη πρὸς Ῥοβοὰμ τὸν τοῦ Σαλωμὼν, καὶ Ῥοβοὰμ ἦν νεώτερος καὶ δειλὸς τῇ καρδίᾳ, καὶ οὐκ ἀντέστη κατὰ πρόσωπον αὐτοῦ.
\vs{8}Καὶ νῦν ὑμεῖς λέγετε ἀντιστῆναι κατὰ πρόσωπον βασιλείας Κυρίου διὰ χειρὸς υἱῶν Δαυίδ· καὶ ὑμεῖς πλῆθος πολὺ, καὶ μεθʼ ὑμῶν μόσχοι χρυσοῖ οὓς ἐποίησεν ὑμῖν Ἱεροβοὰμ εἰς θεούς.
\vs{9}Ἢ οὐκ ἐξεβάλετε τοὺς ἱερεῖς Κυρίου τοὺς υἱοὺς Ἀαρὼν καὶ τοὺς Λευίτας, καὶ ἐποιήσατε ἑαυτοῖς ἱερεῖς ἐκ τοῦ λαοῦ τῆς γῆς πάσης; ὁ προσπορευόμενος πληρῶσαι τὰς χεῖρας ἐν μόσχῳ ἐκ βοῶν καὶ κριοῖς ἑπτὰ, καὶ ἐγίνετο εἰς ἱερέα τῷ μὴ ὄντι θεῷ.
\vs{10}Καὶ ἡμεῖς Κύριον τὸν Θεὸν ἡμῶν οὐκ ἐγκατελίπομεν, καὶ οἱ ἱερεῖς αὐτοῦ λειτουργοῦσι τῷ Κυρίῳ οἱ υἱοὶ Ἀαρὼν καὶ οἱ Λευῖται, καὶ ἐν ταῖς ἐφημερίαις αὐτῶν
\vs{11}θυμιῶσι τῷ Κυρίῳ ὁλοκαύτωμα πρωῒ καὶ δείλης, καὶ θυμίαμα συνθέσεως, καὶ προθέσεις ἄρτων ἐπὶ τῆς τραπέζης τῆς καθαρᾶς, καὶ ἡ λυχνία ἡ χρυσῆ καὶ οἱ λυχνοὶ τῆς καύσεως ἀνάψαι δείλης· ὅτι φυλάσσομεν ἡμεῖς τὰς φυλακὰς Κυρίου τοῦ Θεοῦ τῶν πατέρων ἡμῶν, καὶ ὑμεῖς ἐγκατελίπετε αὐτόν.
\vs{12}Καὶ ἰδοὺ μεθʼ ἡμῶν ἐν ἀρχῇ Κύριος καὶ οἱ ἱερεῖς αὐτοῦ, καὶ αἱ σάλπιγγες τῆς σημασίας τοῦ σημαίνειν ἐφʼ ἡμᾶς. Οἱ υἱοὶ τοῦ Ἰσραὴλ, μὴ πολεμήσητε πρὸς Κύριον Θεὸν τῶν πατέρων ἡμῶν, ὅτι οὐκ εὐοδώσεται ὑμῖν.

\vs{13}Καὶ Ἱεροβοὰμ ἀπέστρεψε τὸ ἔνεδρον ἐλθεῖν αὐτῷ ἐκ τῶν ὄπισθεν, καὶ ἐγένετο ἔμπροσθεν Ἰούδα, καὶ τὸ ἔνεδρον ἐκ τῶν ὄπισθεν.
\vs{14}Καὶ ἀπέστρεψεν Ἰούδας, καὶ ἰδοὺ αὐτοῖς ὁ πόλεμος ἐκ τῶν ἔμπροσθεν καὶ ἐκ τῶν ὄπισθεν, καὶ ἐβόησαν πρὸς Κύριον, καὶ οἱ ἱερεῖς ἐσάλπισαν ταῖς σάλπιγξι.
\vs{15}Καὶ ἐβόησαν ἄνδρες Ἰούδα· καὶ ἐγένετο ἐν τῷ βοᾷν ἄνδρας Ἰούδα, καὶ Κύριος ἐπάταξε τὸν Ἱεροβοὰμ καὶ τὸν Ἰσραὴλ ἐναντίον Ἀβιὰ καὶ Ἰούδα.
\vs{16}Καὶ ἔφυγον οἱ υἱοὶ Ἰσραὴλ ἀπὸ προσώπου Ἰούδα, καὶ παρέδωκεν αὐτοὺς Κύριος εἰς τὰς χεῖρας αὐτῶν.
\vs{17}Καὶ ἐπάταξεν ἐν αὐτοῖς Ἀβιὰ καὶ ὁ λαὸς αὐτοῦ πληγὴν μεγάλην· καὶ ἔπεσον τραυματίαι ἀπὸ Ἰσραὴλ πεντακόσιαι χιλιάδες ἄνδρες δυνατοί.
\vs{18}Καὶ ἐταπεινώθησαν οἱ υἱοὶ Ἰσραὴλ ἐν τῇ ἡμέρᾳ ἐκείνῃ, καὶ κατίσχυσαν οἱ υἱοὶ Ἰούδα, ὅτι ἤλπισαν ἐπὶ Κύριον Θεὸν τῶν πατέρων αὐτῶν.
\vs{19}Καὶ κατεδίωξεν Ἀβιὰ ὀπίσω Ἱεροβοὰμ, καὶ προκατελάβετο παρʼ αὐτοῦ τὰς πόλεις, τὴν Βαιθὴλ καὶ τὰς κώμας αὐτῆς, καὶ τὴν Ἰεσυνὰ καὶ τὰς κώμας αὐτῆς, καὶ τὴν Ἐφρὼν καὶ τὰς κώμας αὐτῆς.
\vs{20}Καὶ οὐκ ἔσχεν ἰσχὺν Ἱεροβοὰμ ἔτι πάσας τὰς ἡμέρας Ἀβιὰ, καὶ ἐπάταξεν αὐτὸν Κύριος, καὶ ἐτελεύτησε.

\vs{21}Καὶ κατίσχυσεν Ἀβιὰ, καὶ ἔλαβεν ἑαυτῷ γυναῖκας δεκατέσσαρας, καὶ ἐγέννησεν υἱοὺς εἰκοσιδύο καὶ ἐκκαίδεκα θυγατέρας.

\vs{22}Καὶ οἱ λοιποὶ λόγοι Ἀβιὰ καὶ αἱ πράξεις αὐτοῦ καὶ οἱ λόγοι αὐτοῦ γεγραμμένοι ἐπὶ βιβλίῳ τοῦ προφήτου Ἀδδώ.

\vs{23}Καὶ ἀπέθανεν Ἀβιὰ μετὰ τῶν πατέρων αὐτοῦ, καὶ ἔθαψαν αὐτὸν ἐν πόλει Δαυὶδ, καὶ ἐβασίλευσεν Ἀσὰ υἱὸς αὐτοῦ ἀντʼ αὐτοῦ. ἐν ταῖς ἡμέραις Ἀσὰ ἡσύχασεν ἡ γῆ Ἰούδα δέκα ἔτη.

\ch{14}
Καὶ ἐποίησε τὸ καλὸν καὶ τὸ εὐθὲς ἐνώπιον Κυρίου τοῦ Θεοῦ αὐτοῦ.
\vs{2}Καὶ ἀπέστησε τὰ θυσιαστήρια τῶν ἀλλοτρίων καὶ τὰ ὑψηλὰ, καὶ συνέτριψε τὰς στήλας, καὶ ἐξέκοψε τὰ ἄλση,
\vs{3}καὶ εἶπε τῷ Ἰούδα ἐκζητῆσαι τὸν Κύριον Θεὸν τῶν πατέρων αὐτῶν, καὶ ποιῆσαι τὸν νόμον καὶ τὰς ἐντολάς.
\vs{4}Καὶ ἀπέστησεν ἀπὸ πασῶν πόλεων Ἰούδα τὰ θυσιαστήρια καὶ τὰ εἴδωλα, καὶ εἰρήνευσε πόλεις τειχήρεις ἐν γῇ Ἰούδα,
\vs{5}ὅτι εἰρήνευσεν ἡ γῆ, καὶ οὐκ ἦν αὐτῷ πόλεμος ἐν τοῖς ἔτεσι τούτοις, ὅτι κατέπαυσε Κύριος αὐτῷ.
\vs{6}Καὶ εἶπε τῷ Ἰούδα, οἰκοδομήσωμεν τὰς πόλεις ταύτας, καὶ ποιήσωμεν τείχη καὶ πύργους καὶ πύλας καὶ μοχλοὺς, ἐνώπιον τῆς γῆς κυριεύσομεν· ὅτι καθὼς ἐξεζητήσαμεν Κύριον τὸν Θεὸν ἡμῶν, ἐξεζήτησεν ἡμᾶς, καὶ κατέπαυσεν ἡμᾶς κυκλόθεν, καὶ εὐώδωσεν ἡμῖν.
\vs{7}Καὶ ἐγένετο δύναμις τῷ Ἀσὰ ὁπλοφόρων αἰρόντων θυρεοὺς καὶ δόρατα ἐν γῇ Ἰούδα τριακόσιαι χιλιάδες, καὶ ἐν γῇ Βενιαμὶν πελτασταὶ καὶ τοξόται διακόσιαι καὶ ὀγδοήκοντα χιλιάδες, πάντες οὗτοι πολεμισταὶ δυνάμεως.

\vs{8}Καὶ ἐξῆλθεν ἐπʼ αὐτοὺς Ζαρὲ ὁ Αἰθίοψ ἐν δυνάμει ἐν χιλίαις χιλιάσι καὶ ἅρμασι τριακοσίοις, καὶ ἦλθεν ἕως Μαρησά.
\vs{9}Καὶ ἐξῆλθεν Ἀσὰ εἰς συνάντησιν αὐτῷ, καὶ παρετάξατο πόλεμον ἐν τῇ φάραγγι κατὰ Βοῤῥᾶν Μαρησά.
\vs{10}Καὶ ἐβόησεν Ἀσὰ πρὸς Κύριον Θεὸν αὐτοῦ, καὶ εἶπε, Κύριε, οὐκ ἀδυνατεῖ παρὰ σοὶ σώζειν ἐν πολλοῖς καὶ ἐν ὀλίγοις· κατίσχυσον ἡμᾶς Κύριε ὁ Θεὸς ἡμῶν, ὅτι ἐπὶ σοὶ πεποιθαμεν, καὶ ἐπὶ τῷ ὀνόματί σου ἤλθομεν ἐπὶ τὸ πλῆθος τὸ πολὺ τοῦτο· Κύριε ὁ Θεὸς ἡμῶν, μὴ κατισχυσάτω πρὸς σὲ ἄνθρωπος.
\vs{11}Καὶ ἐπάταξε Κύριος τοὺς Αἰθίοπας ἐναντίον Ἰούδα, καὶ ἔφυγον Αἰθίοπες,
\vs{12}καὶ κατεδίωξεν αὐτοὺς Ἀσὰ καὶ ὁ λαὸς αὐτοῦ ἕως Γεδώρ· καὶ ἔπεσον Αἰθίοπες ὥστε μὴ εἶναι ἐν αὐτοῖς περιποίησιν, ὅτι συνετρίβησαν ἐνώπιον Κυρίου καὶ ἐναντίον τῆς δυνάμεως αὐτοῦ, καὶ ἐσκύλευσαν σκῦλα πολλά.
\vs{13}Καὶ ἐξέκοψαν τὰς κώμας αὐτῶν κύκλῳ Γεδὼρ, ὅτι ἐγενήθη ἔκστασις Κυρίου ἐπʼ αὐτοὺς, καὶ ἐσκύλευσαν πάσας τὰς πόλεις αὐτῶν, ὅτι πολλὰ σκῦλα ἐγενήθη αὐτοῖς.
\vs{14}Καί γε σκηνὰς κτήσεων, καὶ τοὺς Ἀλιμαζονεῖς ἐξέκοψαν, καὶ ἔλαβον πρόβατα πολλὰ καὶ καμήλους, καὶ ἐπέστρεψαν εἰς Ἱερουσαλήμ.

\ch{15}
Καὶ Ἀζαρίας υἱὸς Ὡδὴδ, ἐγένετο ἐπʼ αὐτὸν πνεῦμα Κυρίου.
\vs{2}Καὶ ἐξῆλθεν εἰς ἀπάντησιν Ἀσὰ καὶ παντὶ Ἰούδα καὶ Βενιαμὶν, καὶ εἶπεν, ἀκούσατέ μου Ἀσὰ, καὶ πᾶς Ἰούδα καὶ Βενιαμίν· Κύριος μεθʼ ὑμῶν ἐν τῷ εἶναι ὑμᾶς μετʼ αὐτοῦ, καὶ ἐὰν ἐκζητήσητε αὐτὸν εὑρεθήσεται ὑμῖν, καὶ ἐὰν ἐγκαταλείπητε αὐτὸν, ἐγκαταλείψει ὑμᾶς.
\vs{3}Καὶ ἡμέραι πολλαὶ τῷ Ἰσραὴλ ἐν οὐ Θεῷ ἀληθινῷ, καὶ οὐχ ἱερέως ὑποδεικνῦντος, καὶ ἐν οὐ νόμῳ.
\vs{4}Καὶ ἐπιστρέψει αὐτοὺς ἐπὶ Κύριον Θεὸν Ἰσραὴλ, καὶ εὑρεθήσεται αὐτοῖς.
\vs{5}Καὶ ἐν ἐκείνῳ τῷ καιρῷ οὐκ ἔστιν εἰρήνη τῷ ἐκπορευομένῳ καὶ τῷ εἰσπορευομένῳ, ὅτι ἔκστασις Κυρίου ἐπὶ πάντας τοὺς κατοικοῦντας τὰς χώρας.
\vs{6}Καὶ πολεμήσει ἔθνος πρὸς ἔθνος καὶ πόλις πρὸς πόλιν, ὅτι ὁ Θεὸς ἐξέστησεν αὐτοὺς ἐν πάσῃ θλίψει.
\vs{7}Καὶ ὑμεῖς ἰσχύσατε, καὶ μὴ ἐκλυέσθωσαν αἱ χεῖρες ὑμῶν, ὅτι ἔστι μισθὸς τῇ ἐργασίᾳ ὑμῶν.

\vs{8}Καὶ ἐν τῷ ἀκοῦσαι τοὺς λόγους τούτους καὶ τὴν προφητείαν Ἀδὰδ τοῦ προφήτου, καὶ κατίσχυσε καὶ ἐξέβαλε τὰ βδελύγματα ἀπὸ πάσης τῆς γῆς Ἰούδα καὶ Βενιαμὶν καὶ ἀπὸ τῶν πόλεων ὧν κατέσχεν Ἱεροβοὰμ ἐν ὄρει Ἐφραὶμ, καὶ ἐνεκαίνισε τὸ θυσιαστήριον Κυρίου, ὃ ἦν ἔμπροσθεν τοῦ ναοῦ Κυρίου.
\vs{9}Καὶ ἐξεκκλησίασε τὸν Ἰούδαν καὶ Βενιαμὶν καὶ τοὺς προσηλύτους τοὺς παροικοῦντας μετʼ αὐτοῦ ἀπὸ Ἐφραὶμ καὶ ἀπὸ Μανασσῆ καὶ ἀπὸ Συμεὼν, ὅτι προσετέθησαν πρὸς αὐτὸν πολλοὶ τοῦ Ἰσραὴλ ἐν τῷ ἰδεῖν αὐτοὺς, ὅτι Κύριος ὁ Θεὸς αὐτοῦ μετʼ αὐτοῦ.
\vs{10}Καὶ συνήχθησαν εἰς Ἱερουσαλὴμ ἐν τῷ μηνὶ τῷ τρίτῳ ἐν τῷ ἔτει τῷ πεντεκαιδεκάτῳ τῆς βασιλείας Ἀσά.
\vs{11}Καὶ ἔθυσε τῷ Κυρίῳ ἐν τῇ ἡμέρᾳ ἐκείνῃ ἀπὸ τῶν σκύλων ὧν ἤνεγκαν, μόσχους ἑπτακοσίους καὶ πρόβατα ἑπτακισχίλια.

\vs{12}Καὶ διῆλθεν ἐν διαθήκῃ ζητῆσαι Κύριον Θεὸν τῶν πατέρων αὐτῶν ἐξ ὅλης τῆς καρδίας καὶ ἐξ ὅλης τῆς ψυχῆς αὐτῶν.
\vs{13}Καὶ πᾶς ὃς ἐὰν μὴ ἐκζητήσῃ τὸν Κύριον Θεὸν τοῦ Ἰσραὴλ, ἀποθανεῖται ἀπὸ νεωτέρου ἕως πρεσβυτέρου, ἀπὸ ἀνδρὸς ἕως γυναικός.
\vs{14}Καὶ ὤμοσαν ἐν Κυρίῳ ἐν φωνῇ μεγάλῃ καὶ ἐν σάλπιγξι καὶ ἐν κερατίναις.
\vs{15}Καὶ εὐφράνθησαν πᾶς Ἰούδα περὶ τοῦ ὅρκου, ὅτι ἐξ ὅλης τῆς ψυχῆς ὤμοσαν, καὶ ἐν πάσῃ θελήσει ἐζήτησαν αὐτὸν· καὶ εὑρέθη αὐτοῖς, καὶ κατέπαυσε Κύριος αὐτοῖς κυκλόθεν.

\vs{16}Καὶ τὴν Μααχὰ τὴν μητέρα αὐτοῦ μετέστησε τοῦ μὴ εἶναι τῇ Ἀστάρτῃ λειτουργοῦσαν, καὶ κατέκοψε τὸ εἴδωλον, καὶ κατέκαυσεν ἐν χειμάῤῥῳ Κέδρων.
\vs{17}Πλὴν τὰ ὑψηλὰ οὐκ ἀπέστησαν, ἔτι ὑπῆρχεν ἐν τῷ Ἰσραήλ· ἀλλʼ ἡ καρδία Ἀσὰ ἐγένετο πλήρης πάσας τὰς ἡμέρας αὐτοῦ.
\vs{18}Καὶ εἰσήνεγκε τὰ ἅγια Δαυὶδ τοῦ πατρὸς αὐτοῦ, καὶ τὰ ἅγια οἴκου τοῦ Θεοῦ, ἀργύριον καὶ χρυσίον καὶ σκεύη.
\vs{19}Καὶ πόλεμος οὖκ ἦν μετʼ αὐτοῦ ἕως τοῦ πέμπτου καὶ τριακοστοῦ ἔτους τῆς βασιλείας Ἀσά.

\ch{16}
Καὶ ἐν τῷ ὀγδόῳ καὶ τριακοστῷ ἔτει τῆς βασιλείας Ἀσὰ, ἀνέβη βασιλεὺς Ἰσραὴλ ἐπὶ Ἰούδαν, καὶ ᾠκοδόμησε τὴν Ῥαμὰ τοῦ μὴ δοῦναι ἔξοδον καὶ εἴσοδον τῷ Ἀσὰ βασιλεῖ Ἰούδα.

\vs{2}Καὶ ἔλαβεν Ἀσὰ ἀργύριον καὶ χρυσίον ἐκ θησαυρῶν οἴκου Κυρίου καὶ οἴκου τοῦ βασιλέως, καὶ ἀπέστειλε πρὸς τὸν υἱὸν τοῦ Ἄδερ βασιλέως Συρίας τὸν κατοικοῦντα ἐν Δαμασκῷ, λέγων,
\vs{3}διάθου διαθήκην ἀναμέσον ἐμοῦ καὶ σοῦ, καὶ ἀναμέσον τοῦ πατρός μου καὶ ἀναμέσον τοῦ πατρός σου· ἰδοὺ ἀπέσταλκά σοι χρυσίον καὶ ἀργύριον· δεῦρο καὶ διασκέδασον ἀπʼ ἐμοῦ τὸν Βαασὰ βασιλέα Ἰσραὴλ, καὶ ἀπελθέτω ἀπʼ ἐμοῦ.

\vs{4}Καὶ ἤκουσεν υἱὸς Ἄδερ τοῦ βασιλέως Ἀσὰ, καὶ ἀπέστειλε τοὺς ἄρχοντας τῆς δυνάμεως αὐτοῦ ἐπὶ τὰς πόλεις Ἰσραὴλ, καὶ ἐπάταξε τὴν Ἀϊὼν, καὶ τὴν Δὰν, καὶ τὴν Ἀβελμαῒν, καὶ πάσας τὰς περιχώρους Νεφθαλί.

\vs{5}Καὶ ἐγένετο ἐν τῷ ἀκοῦσαι Βαασὰ, ἀπέλιπε τοῦ μηκέτι οἰκοδομεῖν τὴν Ῥαμὰ, καὶ κατέπαυσε τὸ ἔργον αὐτοῦ.
\vs{6}Καὶ Ἀσὰ βασιλεὺς ἔλαβε πάντα τὸν Ἰούδαν, καὶ ἔλαβε τοὺς λίθους τῆς Ῥαμὰ καὶ τὰ ξύλα αὐτῆς, ἃ ᾠκοδόμησε Βαασὰ, καὶ ᾠκοδόμησεν ἐν αὐτοῖς τὴν Γαβαὲ καὶ τὴν Μασφά.

\vs{7}Καὶ ἐν τῷ καιρῷ ἐκείνῳ ἦλθεν Ἀνανὶ ὁ προφήτης πρὸς Ἀσὰ βασιλέα Ἰούδα, καὶ εἶπεν αὐτῷ, ἐν τῷ πεποιθέναι σε ἐπὶ βασιλέα Συρίας, καὶ μὴ πεποιθέναι σε ἐπὶ Κύριον Θεόν σου, διὰ τοῦτο ἐσώθη ἡ δύναμις Συρίας ἀπὸ τῆς χειρός σου.
\vs{8}Οὐχ οἱ Αἰθίοπες καὶ Λίβυες ἦσαν εἰς δύναμιν πολλὴν, εἰς θάρσος, εἰς ἱππεῖς, εἰς πλῆθος σφόδρα; καὶ ἐν τῷ πεποιθέναι σε ἐπὶ Κύριον παρέδωκεν εἰς τὰς χεῖράς σου;
\vs{9}Ὅτι οἱ ὀφθαλμοὶ Κυρίου ἐπιβλέπουσιν ἐν πάσῃ τῇ γῇ, κατισχύσαι ἐν πάσῃ καρδίᾳ πλήρει πρὸς αὐτόν· ἠγνόηκας ἐπὶ τούτῳ, ἀπὸ τοῦ νῦν ἔσται μετὰ σοῦ πόλεμος.
\vs{10}Καὶ ἐθυμώθη Ἀσὰ τῷ προφήτῃ, καὶ παρέθετο αὐτὸν εἰς φυλακὴν, ὅτι ὠργίσθη ἐπὶ τούτῳ, καὶ ἐλυμῄνατο Ἀσὰ ἐν τῷ λαῷ ἐν τῷ καιρῷ ἐκείνῳ.

\vs{11}Καὶ ἰδοὺ οἱ λόγοι Ἀσὰ, οἱ πρῶτοι καὶ οἱ ἔσχατοι, γεγραμμένοι ἐν βιβλίῳ βασιλέων Ἰούδα καὶ Ἰσραήλ.

\vs{12}Καὶ ἐμαλακίσθη Ἀσὰ ἐν τῷ ἔτει τῷ ἐννάτῳ καὶ τριακοστῷ τῆς βασιλείας αὐτοῦ τοὺς πόδας, ἕως σφόδρα ἐμαλακίσθη· καὶ ἐν τῇ μαλακίᾳ αὐτοῦ οὐκ ἐζήτησε τὸν Κύριον, ἀλλὰ τοὺς ἰατρούς.
\vs{13}Καὶ ἐκοιμήθη Ἀσὰ μετὰ τῶν πατέρων αὐτοῦ, καὶ ἐτελεύτησεν ἐν τῷ τεσσαρακοστῷ ἔτει τῆς βασιλείας αὐτοῦ.
\vs{14}Καὶ ἔθαψαν αὐτὸν ἐν τῷ μνήματι, ᾧ ὤρυξεν ἑαυτῷ ἐν πόλει Δαυὶδ, καὶ ἐκοίμισαν αὐτὸν ἐπὶ τῆς κλίνης, καὶ ἔπλησαν ἀρωμάτων καὶ γένη μύρων μυρεψῶν, καὶ ἐποίησαν αὐτῷ ἐκφορὰν μεγάλην ἕως σφόδρα.

\ch{17}
Καὶ ἐβασίλευσεν Ἰωσαφὰτ υἱὸς αὐτοῦ ἀντʼ αὐτοῦ, καὶ κατίσχυσεν Ἰωσαφὰτ ἐπὶ τὸν Ἰσραήλ.
\vs{2}Καὶ ἔδωκε δύναμιν ἐν πάσαις ταῖς πόλεσιν Ἰούδα ταῖς ὀχυραῖς, καὶ κατέστησεν ἡγουμένους ἐν πάσαις ταῖς πόλεσιν Ἰούδα, καὶ ἐν πόλεσιν Ἐφραὶμ, ἃς προκατελάβετο Ἀσὰ ὁ πατὴρ αὐτοῦ.
\vs{3}Καὶ ἐγένετο Κύριος μετὰ Ἰωσαφὰτ, ὅτι ἐπορεύθη ἐν ὁδοῖς τοῦ πατρὸς αὐτοῦ ταῖς πρώταις, καὶ οὐκ ἐξεζήτησε τὰ εἴδωλα,
\vs{4}ἀλλὰ Κύριον τὸν Θεὸν τοῦ πατρὸς αὐτοῦ ἐξεζήτησε, καὶ ἐν ταῖς ἐντολαῖς τοῦ πατρὸς αὐτοῦ ἐπορεύθη, καὶ οὐχ ὡς τὰ ἔργα τοῦ Ἰσραήλ.
\vs{5}Καὶ κατεύθυνε Κύριος τὴν βασιλείαν ἐν χειρὶ αὐτοῦ, καὶ ἔδωκε πᾶς Ἰούδα δῶρα τῷ Ἰωσαφὰτ, καὶ ἐγένετο αὐτῷ πλοῦτος καὶ δόξα πολλή.
\vs{6}Καὶ ὑψώθη ἡ καρδία αὐτοῦ ἐν ὁδῷ Κυρίου, καὶ ἐξῇρε τὰ ὑψηλὰ καὶ τὰ ἄλση ἀπὸ τῆς γῆς Ἰούδα.

\vs{7}Καὶ ἐν τῷ ἔτει τῷ τρίτῳ ἔτει τῆς βασιλείας αὐτοῦ ἀπέστειλε τοὺς ἡγουμένους αὐτοῦ καὶ τοὺς υἱοὺς τῶν δυνατῶν, τὸν Ἀβδιὰν, καὶ Ζαχαρίαν, καὶ Ναθαναὴλ, καὶ Μιχαίαν, τοῦ διδάσκειν ἐν πόλεσιν Ἰούδα.
\vs{8}Καὶ μετʼ αὐτῶν οἱ Λευῖται, Σαμαίας, καὶ Ναθανίας, καὶ Ζαβδίας, καὶ Ἀσιὴλ, καὶ Σεμιραμὼθ, καὶ Ἰωνάθαν, καὶ Ἀδωνίας, καὶ Τωβίας, καὶ Τωβαδωνίας, Λευῖται, καὶ οἱ μετʼ αὐτῶν Ἐλισαμὰ, καὶ Ἰωρὰμ, οἱ ἱερεῖς.
\vs{9}Καὶ ἐδίδασκον ἐν Ἰούδα, καὶ μετʼ αὐτῶν βίβλος νόμου Κυρίου, καὶ διῆλθον ἐν ταῖς πόλεσιν Ἰούδα, καὶ ἐδίδασκον τὸν λαόν.

\vs{10}Καὶ ἐγένετο ἔκστασις Κυρίου ἐπὶ πάσαις ταῖς βασιλείαις τῆς γῆς κύκλῳ Ἰούδα, καὶ οὐκ ἐπολέμουν πρὸς Ἰωσαφάτ.
\vs{11}Καὶ ἀπὸ τῶν ἀλλοφύλων ἔφερον τῷ Ἰωσαφὰτ δῶρα καὶ ἀργύριον καὶ δόματα· καὶ οἱ Ἀραβες ἔφερον αὐτῷ κριοὺς προβάτων ἑπτακισχιλίους ἑπτακοσίους.
\vs{12}Καὶ ἦν Ἰωσαφὰτ πορευόμενος μείζων ἕως εἰς ὕψος, καὶ ᾠκοδόμησεν ἐν τῇ Ἰουδαίᾳ οἰκήσεις καὶ πόλεις ὀχυράς.
\vs{13}Καὶ ἔργα πολλὰ ἐγένετο αὐτῷ ἐν τῇ Ἰουδαίᾳ· καὶ ἄνδρες πολεμισταὶ δυνατοὶ ἰσχύοντες ἐν Ἱερουσαλήμ.

\vs{14}Καὶ οὗτος ὁ ἀριθμὸς αὐτῶν κατʼ οἴκους πατριῶν αὐτῶν· καὶ τῷ Ἰούδα χιλίαρχοι, Ἔδνας ὁ ἄρχων, καὶ μετʼ αὐτοῦ υἱοὶ δυνατοὶ δυνάμεως τριακόσιαι χιλιάδες·
\vs{15}Καὶ μετʼ αὐτὸν, Ἰωανὰν ὁ ἡγούμενος, καὶ μετʼ αὐτοῦ διακόσιαι ὀγδοήκοντα χιλιάδες·
\vs{16}Καὶ μετʼ αὐτὸν Ἀμασίας ὁ τοῦ Ζαρὶ, ὁ προθυμούμενος τῷ κυρίῳ, καὶ μετʼ αὐτοῦ διακόσιαι χιλιάδες δυνατοὶ δυνάμεως.
\vs{17}Καὶ ἐκ τοῦ Βενιαμὶν δυνατὸς δυνάμεως καὶ Ἐλιαδὰ, καὶ μετʼ αὐτοῦ τοξόται καὶ πελτασταὶ διακόσιαι χιλιάδες·
\vs{18}Καὶ μετʼ αὐτὸν Ἰωζαβὰδ, καὶ μετʼ αὐτοῦ ἑκατὸν ὀγδοήκοντα χιλιάδες δυνατοὶ πολέμου.
\vs{19}Οὗτοι οἱ λειτουργοῦντες τῷ βασιλεῖ, ἐκτὸς ὧν ἔδωκεν ὁ βασιλεὺς ἐν ταῖς πόλεσι ταῖς ὀχυραῖς ἐν πάσῃ τῇ Ἰουδαίᾳ.

\ch{18}
Καὶ ἐγενήθη τῷ Ἰωσαφὰτ ἔτι πλοῦτος καὶ δόξα πολλὴ, καὶ ἐπεγαμβρεύσατο ἐν οἴκῳ Ἀχαάβ.
\vs{2}Καὶ κατέβη διὰ τέλους ἐτῶν πρὸς Ἀχαὰβ εἰς Σαμάρειαν· καὶ ἔθυσεν αὐτῷ Ἀχαὰβ πρόβατα καὶ μόσχους πολλοὺς, καὶ τῷ λαῷ τῷ μετʼ αὐτοῦ, καὶ ἠγάπα αὐτὸν τοῦ συναναβῆναι μετʼ αὐτοῦ εἰς Ῥαμὼθ τῆς Γαλααδίτιδος.
\vs{3}Καὶ εἶπεν Ἀχαὰβ βασιλεὺς Ἰσραὴλ πρὸς Ἰωσαφὰτ βασιλέα Ἰούδα, εἰ πορεύσῃ μετʼ ἐμοῦ εἰς Ῥαμὼθ τῆς Γαλααδίτιδος; καὶ εἶπεν αὐτῷ, ὡς ἐγὼ, οὕτω καὶ σύ· ὡς ὁ λαός σου, καὶ ὁ λαός μου μετὰ σοῦ εἰς πόλεμον.

\vs{4}Καὶ εἶπεν Ἰωσαφὰτ πρὸς βασιλέα Ἰσραὴλ, ζήτησον δὴ σήμερον τὸν Κύριον.
\vs{5}Καὶ συνήγαγεν ὁ βασιλεὺς Ἰσραὴλ τοὺς προφήτας τετρακοσίους ἄνδρας, καὶ εἶπεν αὐτοῖς, εἰ πορευθῶ εἰς Ῥαμὼθ Γαλαὰδ εἰς πόλεμον, ἢ ἐπίοχω; καὶ εἶπαν, ἀνάβαινε, καὶ δώσει ὁ Θεὸς εἰς τὰς χεῖρας τοῦ βασιλέως.
\vs{6}Καὶ εἶπεν Ἰωσαφὰτ, οὐκ ἔστιν ὧδε προφήτης τοῦ κυρίου ἔτι, καὶ ἐπιζητήσομεν παρʼ αὐτοῦ;
\vs{7}Καὶ εἶπε βασιλεὺς Ἰσραὴλ πρὸς Ἰωσαφὰτ, ἔτι ἀνὴρ εἷς τοῦ ζητῆσαι τὸν Κύριον διʼ αὐτοῦ, καὶ ἐγὼ ἐμίσησα αὐτὸν, ὅτι οὐκ ἔστι προφητεύων περὶ ἐμοῦ εἰς ἀγαθὰ, ὅτι πᾶσαι αἱ ἡμέραι αὐτοῦ εἰς κακὰ, οὗτος Μιχαίας υἱὸς Ἰεμβλά· καὶ εἶπεν Ἰωσαφὰτ, μὴ λαλείτω ὁ βασιλεὺς οὕτως.

\vs{8}Καὶ ἐκάλεσεν ὁ βασιλεὺς εὐνοῦχον ἕνα, καὶ εἶπε, τάχος Μιχαίαν υἱὸν Ἰεμβλά.
\vs{9}Καὶ βασιλεὺς Ἰσραὴλ καὶ Ἰωσαφὰτ βασιλεὺς Ἰούδα καθήμενοι ἕκαστος ἐπὶ θρόνου αὐτοῦ, καὶ ἐνδεδυμένοι στολὰς, καθήμενοι ἐν τῷ εὐρυχώρῳ θύρας πύλης Σαμαρείας, καὶ πάντες οἱ προφῆται προεφήτευον ἐναντίον αὐτῶν.
\vs{10}Καὶ ἐποίησεν ἑαυτῷ Σεδεκίας υἱὸς Χαναὰν κέρατα σιδηρᾶ, καὶ εἶπε, τάδε λέγει Κύριος, ἐν τούτοις κερατιεῖς τὴν Συρίαν ἕως ἂν συντελεσθῇ.
\vs{11}Καὶ πάντες οἱ προφῆται προεφήτευον οὕτω, λέγοντες, ἀνάβαινε εἰς Ῥαμὼθ Γαλαὰδ, καὶ εὐοδωθήσῃ, καὶ δώσει Κύριος εἰς χεῖρας τοῦ βασιλέως.

\vs{12}Καὶ ὁ ἄγγελος ὁ πορευθεὶς τοῦ καλέσαι τὸν Μιχαίαν, ἐλάλησεν αὐτῷ, λέγων, ἰδοὺ ἐλάλησαν οἱ προφῆται ἐν στόματι ἑνὶ ἀγαθὰ περὶ τοῦ βασιλέως, καὶ ἔστωσαν δὴ οἱ λόγοι σου ὡς ἑνὸς αὐτῶν, καὶ λαλήσεις ἀγαθά.
\vs{13}Καὶ εἶπε Μιχαίας, ζῇ Κύριος, ὅτι ὃ ἐὰν εἴπῃ ὁ Θεὸς πρὸς μὲ, αὐτὸ λαλήσω.

\vs{14}Καὶ ἦλθε πρὸς τὸν βασιλέα, καὶ εἶπεν αὐτῷ ὁ βασιλεὺς, Μιχαία, εἰ πορευθῶ εἰς Ῥαμὼθ Γαλαὰδ εἰς πόλεμον, ἢ ἐπίσχω; καὶ εἶπεν, ἀνάβαινε, καὶ εὐοδώσεις, καὶ δοθήσονται εἰς χεῖρας ὑμῶν.
\vs{15}Καὶ εἶπεν αὐτῷ ὁ βασιλεὺς, ποσάκις ὁρκίζω σε ἵνα μὴ λαλήσῃς πρὸς μὲ πλὴν τὴν ἀλήθειαν ἐν ὀνόματι Κυρίου;
\vs{16}Καὶ εἶπεν, εἶδον τὸν Ἰσραὴλ διεσπαρμένους ἐν τοῖς ὄρεσιν, ὡς πρόβατα οἷς οὐκ ἔστι ποιμήν· καὶ εἶπε Κύριος, οὐκ ἔχουσιν ἡγούμενον οὗτοι, ἀναστρεφέτωσαν ἕκαστος εἰς τὸν οἶκον αὐτοῦ ἐν εἰρήνῃ.

\vs{17}Καὶ εἶπεν ὁ βασιλεὺς Ἰσραὴλ πρὸς Ἰωσαφὰτ, οὐκ εἶπόν σοι, ὅτι οὐ προφητεύει περὶ ἐμοῦ ἀγαθὰ ἀλλʼ ἢ κακά;
\vs{18}Καὶ εἶπεν, οὐχ οὕτως· ἀκούσατε λόγον Κυρίου. εἶδον τὸν Κύριον καθήμενον ἐπὶ θρόνου αὐτοῦ, καὶ πᾶσα δύναμις τοῦ οὐρανοῦ παρειστήκει ἐκ δεξιῶν αὐτοῦ καὶ ἐξ ἀριστερῶν αὐτοῦ.
\vs{19}Καὶ εἶπε Κύριος, τίς ἀπατήσει τὸν Ἀχαὰβ βασιλέα Ἰσραὴλ, καὶ ἀναβήσεται, καὶ πεσεῖται ἐν Ῥαμὼθ Γαλαάδ; καὶ οὗτος εἶπεν οὕτως, καὶ οὗτος εἶπεν οὕτως.
\vs{20}Καὶ ἐξῆλθε τὸ πνεῦμα καὶ ἔστη ἐνώπιον Κυρίου, καὶ εἶπεν, ἐγὼ ἀπατήσω αὐτόν· καὶ εἶπε Κύριος, ἐν τίνι;
\vs{21}Καὶ εἶπεν, ἐξελεύσομαι καὶ ἔσομαι πνεῦμα ψευδὲς ἐν στόματι πάντων τῶν προφητῶν αὐτοῦ. καὶ εἶπεν, ἀπατήσεις καὶ δυνήσῃ, ἔξελθε καὶ ποίησον οὕτω.
\vs{22}Καὶ νῦν ἰδοὺ ἔδωκε Κύριος πνεῦμα ψευδὲς ἐν στόματι τῶν προφητῶν σου τούτων, καὶ Κύριος ἐλάλησεν ἐπὶ σὲ κακά.

\vs{23}Καὶ ἤγγισε Σεδεκίας υἱὸς Χαναὰν, καὶ ἐπάταξε τὸν Μιχαίαν ἐπὶ τὴν σιαγόνα, καὶ εἶπεν αὐτῷ, ποίᾳ τῇ ὁδῷ παρῆλθε πνεῦμα Κυρίου παρʼ ἐμοῦ τοῦ λαλῆσαι πρὸς σέ;
\vs{24}Καὶ εἶπε Μιχαίας, ἰδοὺ ὄψῃ ἐν τῇ ἡμέρᾳ ἐκείνῃ, ἐν ᾗ εἰσελεύσῃ ταμεῖον ἐκ ταμείου τοῦ κατακρυβῆναι.

\vs{25}Καὶ εἶπε βασιλεὺς Ἰσραὴλ, λάβετε τὸν Μιχαίαν καὶ ἀποστρέψατε πρὸς Ἐμὴρ ἄρχοντα τῆς πόλεως, καὶ πρὸς Ἰωὰς ἄρχοντα υἱὸν τοῦ βασιλέως,
\vs{26}καὶ ἐρεῖτε, οὕτως εἶπεν ὁ βασιλεὺς, ἀπόθεσθε τοῦτον εἰς οἶκον φυλακῆς, καὶ ἐσθιέτω ἄρτον θλίψεως καὶ ὕδωρ θλίψεως, ἕως τοῦ ἐπιστρέψαι με ἐν εἰρήνῃ.
\vs{27}Καὶ εἶπε Μιχαίας, ἐὰν ἐπιστρέφων ἐπιστρέψῃς ἐν εἰρήνῃ, οὐκ ἐλάλησε Κύριος ἐν ἐμοί· καὶ εἶπεν, ἀκούσατε λαοὶ πάντες.

\vs{28}Καὶ ἀνέβη βασιλεὺς Ἰσραὴλ, καὶ Ἰωσαφὰτ βασιλεὺς Ἰούδα, εἰς Ῥαμὼθ Γαλαάδ.
\vs{29}Καὶ εἶπε βασιλεὺς Ἰσραὴλ πρὸς Ἰωσαφὰτ, κατακάλυψόν με, καὶ εἰσελεύσομαι εἰς τὸν πόλεμον, καὶ σὺ ἔνδυσαι τὸν ἱματισμόν μου· καὶ συνεκαλύψατο βασιλεὺς Ἰσραὴλ, καὶ εἰσῆλθεν εἰς τὸν πόλεμον.
\vs{30}Καὶ βασιλεὺς Συρίας ἐνετείλατο τοῖς ἄρχουσι τῶν ἁρμάτων τοῖς μετʼ αὐτοῦ, λέγων, μὴ πολεμεῖτε τὸν μικρὸν καὶ τὸν μέγαν, ἀλλʼ ἢ τὸν βασιλέα Ἰσραὴλ μόνον.
\vs{31}Καὶ ἐγένετο ὡς εἶδον οἱ ἄρχοντες τῶν ἁρμάτων τὸν Ἰωσαφὰτ, καὶ αὐτοὶ εἶπαν, βασιλεὺς Ἰσραήλ ἐστι, καὶ ἐκύκλωσαν αὐτὸν τοῦ πολεμεῖν· καὶ ἐβόησεν Ἰωσαφὰτ, καὶ Κύριος ἔσωσεν αὐτὸν, καὶ ἀπέστρεψεν αὐτοὺς ὁ Θεὸς ἀπʼ αὐτοῦ.
\vs{32}Καὶ ἐγένετο ὡς εἶδον οἱ ἄρχοντες τῶν ἁρμάτων ὅτι οὐκ ἦν βασιλεὺς Ἰσραὴλ, καὶ ἀπέστρεψαν ἀπʼ αὐτοῦ.
\vs{33}Καὶ ἀνὴρ ἔτεινε τόξον εὐστόχως, καὶ ἐπάταξε τὸν βασιλέα Ἰσραὴλ ἀναμέσον τοῦ πνεύμονος καὶ ἀναμέσον τοῦ θώρακος· καὶ εἶπε τῷ ἡνίοχῳ, ἐπίστρεφε τὴν χεῖρά σου, ἐξάγαγέ με ἐκ τοῦ πολέμου, ὅτι ἐπόνεσα.
\vs{34}Καὶ ἐτροπώθη ὁ πόλεμος ἐν τῇ ἡμέρᾳ ἐκείνῃ, καὶ ὁ βασιλεὺς Ἰσραὴλ ἦν ἑστηκὼς ἐπὶ τοῦ ἅρματος ἐξεναντίας Συρίας ἕως ἑσπέρας, καὶ ἀπέθανε δύνοντος τοῦ ἡλίου.

\ch{19}
Καὶ ἐπέστρεψεν Ἰωσαφὰτ βασιλεὺς Ἰούδα εἰς τὸν οἶκον αὐτοῦ εἰς Ἱερουσαλήμ.
\vs{2}Καὶ ἐξῆλθεν εἰς ἀπάντησιν αὐτοῦ Ἰηοὺ ὁ τοῦ Ἀνανὶ ὁ προφήτης, καὶ εἶπεν αὐτῷ, βασιλεὺς Ἰωσαφὰτ, εἰ ἁμαρτωλῷ σὺ βοηθεῖς, ἢ μισουμένῳ ὑπὸ Κυρίου φιλιάζεις; διὰ τοῦτο ἐγένετο ἐπὶ σὲ ὀργὴ παρὰ Κυρίου.
\vs{3}Ὅτι ἀλλʼ ἢ λόγοι ἀγαθοὶ ηὑρέθησαν ἐν σοὶ, ὅτι ἐξῇρας τὰ ἄλση ἀπὸ τῆς γῆς Ἰούδα, καὶ κατηύθυνας τὴν καρδίαν σου ἐκζητῆσαι τὸν Κύριον.

\vs{4}Καὶ κατῴκησεν Ἰωσαφὰτ εἰς Ἱερουσαλήμ· καὶ πάλιν ἐξῆλθεν εἰς τὸν λαὸν ἀπὸ Βηρσαβεὲ ἕως ὄρους Ἐφραὶμ, καὶ ἐπέστρεψεν αὐτοὺς ἐπὶ Κύριον Θεὸν τῶν πάτερων αὐτῶν.
\vs{5}Καὶ κατέστησε τοὺς κριτὰς ἐν πάσαις ταῖς πόλεσιν Ἰούδα ταῖς ὀχυραῖς, ἐν πόλει καὶ πόλει.
\vs{6}Καὶ εἶπε τοῖς κριταῖς, ἴδετε τί ὑμεῖς ποιεῖτε, ὅτι οὐκ ἀνθρώπῳ ὑμεῖς κρίνετε, ἀλλʼ ἢ τῷ Κυρίῳ, καὶ μεθʼ ὑμῶν λόγοι τῆς κρίσεως.
\vs{7}Καὶ νῦν γενέσθω φόβος Κυρίου ἐφʼ ὑμᾶς, καὶ φυλάσσετε καὶ ποιήσατε, ὅτι οὐκ ἔστι μετὰ Κυρίου Θεοῦ ἡμῶν ἀδικία, οὐδὲ θαυμάσαι πρόσωπον, οὐδὲ λαβεῖν δῶρα.

\vs{8}Καί γε ἐν Ἱερουσαλὴμ κατέστησεν Ἰωσαφὰτ τῶν ἱερέων καὶ τῶν Λευιτῶν καὶ τῶν πατριαρχῶν Ἰσραὴλ εἰς κρίσιν Κυρίου, καὶ κρίνειν τοὺς κατοικοῦντας ἐν Ἱερουσαλήμ.
\vs{9}Καὶ ἐνετείλατο πρὸς αὐτοὺς, λέγων, οὕτω ποιήσετε ἐν φόβῳ Κυρίου, ἐν ἀληθείᾳ, καὶ ἐν πλήρει καρδίᾳ·
\vs{10}Πᾶς ἀνὴρ κρίσιν τὴν ἐλθοῦσαν ἐφʼ ὑμᾶς τῶν ἀδελφῶν ὑμῶν τῶν κατοικούντων ἐν ταῖς πόλεσιν αὐτῶν ἀναμέσον αἷμα αἵματος, καὶ ἀναμέσον τοῦ προστάγματος καὶ ἐντολῆς, καὶ δικαιώματα καὶ κρίματα, καὶ διαστελεῖσθε αὐτοῖς, καὶ οὐχ ἁμαρτήσονται τῷ κυρίῳ, καὶ οὐκ ἔσται ὀργὴ ἐφʼ ὑμᾶς, καὶ ἐπὶ τοὺς ἀδελφοὺς ὑμῶν· οὕτως ποιήσετε, καὶ οὐχ ἁμαρτήσεσθε.
\vs{11}Καὶ ἰδοὺ Ἀμαρίας ὁ ἱερεὺς ἡγούμενος ἐφʼ ὑμᾶς εἰς πάντα λόγον Κυρίου, καὶ Ζαβδίας υἱὸς Ἰσμαὴλ ὁ ἡγούμενος εἰς οἶκον Ἰούδα πρὸς πάντα λόγον βασιλέως, καὶ οἱ γραμματεῖς καὶ οἱ Λευῖται πρὸ προσώπου ὑμῶν· ἰσχύσατε καὶ ποιήσατε, καὶ ἔσται Κύριος μετὰ τοῦ ἀγαθοῦ.

\ch{20}
Καὶ μετὰ ταῦτα ἦλθον οἱ υἱοὶ Μωὰβ, καὶ υἱοὶ Ἀμμὼν, καὶ μετʼ αὐτῶν ἐκ τῶν Μιναίων πρὸς Ἰωσαφὰτ εἰς πόλεμον.
\vs{2}Καὶ ἦλθον καὶ ὑπέδειξαν τῷ Ἰωσαφὰτ, λέγοντες, ἥκει ἐπὶ σὲ πλῆθος πολὺ ἐκ πέραν τῆς θαλάσσης ἀπὸ Συρίας, καὶ ἰδού εἰσιν ἐν Ἀσασὰν Θαμὰρ, αὕτη ἐστὶν Ἐγγαδί.
\vs{3}Καὶ ἐφοβήθη, καὶ ἔδωκε Ἰωσαφὰτ πρόσωπον αὐτοῦ ἐκζητῆσαι τὸν Κύριον, καὶ ἐκήρυξε νηστείαν ἐν παντὶ Ἰούδα.
\vs{4}Καὶ συνήχθη Ἰούδα ἐκζητῆσαι τὸν Κύριον, καὶ ἀπὸ πασῶν τῶν πόλεων Ἰούδα ἦλθον ζητῆσαι τὸν Κύριον.

\vs{5}Καὶ ἀνέστη Ἰωσαφὰτ ἐν ἐκκλησίᾳ Ἰούδα ἐν Ἱερουσαλὴμ ἐν οἴκῳ Κυρίου κατὰ πρόσωπον τῆς αὐλῆς τῆς καινῆς.
\vs{6}Καὶ εἶπε, Κύριε ὁ Θεὸς τῶν πατέρων μου, οὐχὶ σὺ εἶ Θεὸς ἐν οὐρανῷ ἄνω, καὶ σὺ κυριεύεις πασῶν τῶν βασιλειῶν τῶν ἐθνῶν; καὶ ἐν τῇ χειρί σου ἰσχὺς δυναστείας, καὶ οὐκ ἔστι πρὸς σὲ ἀντιστῆναι;
\vs{7}Οὐχὶ σὺ ὁ Κύριος ὁ ἐξολοθρεύσας τοὺς κατοικοῦντας τὴν γῆν ταύτην ἀπὸ προσώπου τοῦ λαοῦ σου Ἰσραὴλ, καὶ ἔδωκας αὐτὴν σπέρματι Ἁβραὰμ τῷ ἠγαπημένῳ σου εἰς τὸν αἰῶνα;
\vs{8}Καὶ κατῴκησαν ἐν αὐτῇ, καὶ ᾠκοδόμησαν ἐν αὐτῇ ἁγίασμα τῷ ὀνόματί σου, λέγοντες,
\vs{9}ἐὰν ἐπέλθῃ ἐφʼ ἡμᾶς κακὰ, ῥομφαία, κρίσις, θάνατος, λιμὸς, στησόμεθα ἐναντίον τοῦ οἴκου τούτου καὶ ἐναντίον σου, ὅτι τὸ ὄνομά σου ἐπὶ τῷ οἴκῳ τούτῳ, καὶ βοησόμεθα πρὸς σὲ ἀπὸ τῆς θλίψεως, καὶ ἀκούσῃ καὶ σώσεις.
\vs{10}Καὶ νῦν ἰδοὺ οἱ υἱοὶ Ἀμμὼν, καὶ Μωὰβ, καὶ ὄρος Σηείρ· εἰς οὓς οὐκ ἔδωκας τῷ Ἰσραὴλ διελθεῖν διʼ αὐτῶν, ἐξελθόντων αὐτῶν ἐκ γῆς Αἰγύπτου, ὅτι ἐξέκλιναν ἀπʼ αὐτῶν, καὶ οὐκ ἐξωλόθρευσαν αὐτούς·
\vs{11}Καὶ νῦν ἰσοὺ αὐτοὶ ἐπιχειροῦσιν ἐφʼ ἡμᾶς ἐξελθεῖν ἐκβαλεῖν ἡμᾶς ἀπὸ τῆς κληρονομίας ἡμῶν, ἧς ἔδωκας ἡμῖν.
\vs{12}Κύριε ὁ Θεὸς ἡμῶν, οὐ κρινεῖς ἐν αὐτοῖς; ὅτι οὐκ ἔστιν ἡμῖν ἰσχὺς τοῦ ἀντιστῆναι πρὸς τὸ πλῆθος τὸ πολὺ τοῦτο τὸ ἐλθὸν ἐφʼ ἡμᾶς, καὶ οὐκ οἴδαμεν τί ποιήσωμεν αὐτοῖς, ἀλλʼ ἢ ἐπὶ σοὶ οἱ ὀφθαλμοὶ ἡμῶν.

\vs{13}Καὶ πᾶς Ἰούδα ἑστηκὼς ἔναντι Κυρίου, καὶ τὰ παιδία αὐτῶν, καὶ αἱ γυναῖκες αὐτῶν.
\vs{14}Καὶ τῷ Ὀζιὴλ τῷ τοῦ Ζαχαρίου, τῶν υἱῶν Βαναίου, τῶν υἱῶν Ἐλεϊὴλ τοῦ Ματθανίου τοῦ Λευίτου ἀπὸ τῶν υἱῶν Ἀσὰφ, ἐγένετο ἐπʼ αὐτὸν πνεῦμα Κυρίου ἐν τῇ ἐκκλησίᾳ.
\vs{15}Καὶ εἶπεν, ἀκούσατε πᾶς Ἰούδα καὶ οἱ κατοικοῦντες ἐν Ἱερουσαλὴμ καὶ ὁ βασιλεὺς Ἰωσαφάτ· τάδε λέγει Κύριος ὑμῖν αὐτοῖς, μὴ φοβεῖσθε, μηδὲ πτοηθῆτε ἀπὸ προσώπου τοῦ ὄχλου τοῦ πολλοῦ τούτου, ὅτι οὐχ ὑμῖν ἐστιν ἡ παράταξις, ἀλλʼ ἢ τῷ Θεῷ.
\vs{16}Αὔριον κατάβητε ἐπʼ αὐτούς· ἰδοὺ ἀναβαίνουσι κατὰ τὴν ἀνάβασιν Ἀσσεῖς, καὶ εὑρήσετε αὐτοὺς ἐπʼ ἄκρου ποταμοῦ τῆς ἐρήμου Ἰεριήλ.
\vs{17}Οὐχ ὑμῖν ἐστι πολεμῆσαι· ταῦτα σύνετε, καὶ ἴδετε τὴν σωτηρίαν Κυρίου μεθʼ ὑμῶν Ἰούδα καὶ Ἱερουσαλήμ· μὴ φοβηθῆτε, μηδὲ πτοηθῆτε αὔριον ἐξελθεῖν εἰς ἀπάντησιν αὐτοῖς, καὶ Κύριος μεθʼ ὑμῶν.
\vs{18}Καὶ κύψας Ἰωσαφὰτ ἐπὶ πρόσωπον αὐτοῦ καὶ πᾶς Ἰούδα καὶ οἱ κατοικοῦντες Ἱερουσαλὴμ, ἔπεσον ἔναντι Κυρίου προσκυνῆσαι Κυρίῳ.
\vs{19}Καὶ ἀνέστησαν οἱ Λευῖται ἀπὸ τῶν υἱῶν Καὰθ, καὶ ἀπὸ τῶν υἱῶν Κορὲ, αἰνεῖν Κυρίῳ Θεῷ Ἰσραὴλ ἐν φωνῇ μεγάλῃ εἰς ὕψος.

\vs{20}Καὶ ὤρθρισαν πρωῒ, καὶ ἐξῆλθον εἰς τὴν ἔρημον Θεκωέ· καὶ ἐν τῷ ἐξελθεῖν αὐτοὺς, ἔστη Ἰωσαφὰτ καὶ ἐβόησε, καὶ εἶπεν, ἀκούσατέ μου Ἰούδα καὶ οἱ κατοικοῦντες ἐν Ἱερουσαλήμ· ἐμπιστεύσατε ἐν Κυρίῳ Θεῷ ἡμῶν, καὶ ἐμπιστευθήσεσθε· ἐμπιστεύσατε ἐν προφήτῃ αὐτοῦ, καὶ εὐοδωθήσεσθε.
\vs{21}Καὶ ἐβουλεύσατο μετὰ τοῦ λαοῦ, καὶ ἔστησε ψαλτῳδοὺς καὶ αἰνοῦντας, ἐξομολογεῖσθαι καὶ αἰνεῖν τὰ ἅγια ἐν τῷ ἐξελθεῖν ἔμπροσθεν τῆς δυνάμεως, καὶ ἔλεγον, ἐξομολογεῖσθε τῷ Κυρίῳ, ὅτι εἰς τὸν αἰῶνα τὸ ἔλεος αὐτοῦ.

\vs{22}Καὶ ἐν τῷ ἄρξασθαι αὐτοὺς τῆς αἰνέσεως καὶ τῆς ἐξομολογήσεως, ἔδωκε Κύριος πολεμεῖν τοὺς υἱοὺς Ἀμμὼν ἐπὶ Μωὰβ καὶ ὄρος Σηεὶρ τοὺς ἐξελθόντας ἐπὶ Ἰούδαν, καὶ ἐτροπώθησαν.
\vs{23}Καὶ ἀνέστησαν οἱ υἱοὶ Ἀμμὼν καὶ Μωὰβ ἐπὶ τοὺς κατοικοῦντας ὄρος Σηεὶρ, ἐξολοθρεῦσαι καὶ ἐκτρίψαι αὐτούς· καὶ ὡς συνετέλεσαν τοὺς κατοικοῦντας Σηεὶρ, ἀνέστησαν εἰς ἀλλήλους τοῦ ἐξολοθρευθῆναι.

\vs{24}Καὶ Ἰούδας ἦλθεν ἐπὶ τὴν σκοπιὰν τῆς ἐρήμου, καὶ ἐπέβλεψε, καὶ εἶδε τὸ πλῆθος· καὶ ἰδοὺ πάντες νεκροὶ πεπτωκότες ἐπὶ τῆς γῆς, οὐκ ἦν σωζόμενος.
\vs{25}Καὶ ἐξῆλθεν Ἰωσαφὰτ καὶ ὁ λαὸς αὐτοῦ σκυλεῦσαι τὰ σκῦλα αὐτῶν, καὶ εὗρον κτήνη πολλὰ, καὶ ἀποσκευὴν, καὶ σκῦλα, καὶ σκεύη ἐπιθυμητὰ, καὶ ἐσκύλευσαν ἐν αὐτοῖς· καὶ ἐγένοντο ἡμέραι τρεῖς σκυλευόντων αὐτῶν τὰ σκῦλα, ὅτι πολλὰ ἦν.
\vs{26}Καὶ ἐγένετο τῇ ἡμέρᾳ τῇ τετάρτῃ ἐπισυνήχθησαν εἰς τὸν αὐλῶνα τῆς εὐλογίας, ἐκεῖ γὰρ ηὐλόγησαν τὸν Κύριον· διὰ τοῦτο ἐκάλεσαν τὸ ὄνομα τοῦ τόπου ἐκείνου Κοιλὰς Εὐλογίας, ἕως τῆς ἡμέρας ταύτης.

\vs{27}Καὶ ἐπέστρεψε πᾶς ἀνὴρ Ἰούδα εἰς Ἱερουσαλὴμ, καὶ Ἰωσαφὰτ ἡγούμενος αὐτῶν ἐν εὐφροσύνῃ μεγάλῃ, ὅτι εὔφραινεν αὐτοὺς Κύριος ἀπὸ τῶν ἐχθρῶν αὐτῶν.
\vs{28}Καὶ εἰσῆλθον εἰς Ἱερουσαλὴμ ἐν νάβλαις καὶ κινύραις καὶ ἐν σάλπιγξιν εἰς οἶκον Κυρίου.
\vs{29}Καὶ ἐγένετο ἔκστασις Κυρίου ἐπὶ πάσας τὰς βασιλείας τῆς γῆς, ἐν τῷ ἀκοῦσαι αὐτοὺς ὅτι Κύριος ἐπολέμησε πρὸς τοὺς ὑπεναντίους Ἰσραήλ.
\vs{30}Καὶ εἰρήνευσεν ἡ βασιλεία Ἰωσαφὰτ, καὶ κατέπαυσεν αὐτῷ ὁ Θεὸς αὐτοῦ κυκλόθεν.

\vs{31}Καὶ ἐβασίλευσεν Ἰωσαφὰτ ἐπὶ τὸν Ἰούδαν, ὢν ἐτῶν τριακονταπέντε ἐν τῷ βασιλεῦσαι αὐτὸν, καὶ εἴκοσι καὶ πέντε ἔτη ἐβασίλευσεν ἐν Ἱερουσαλὴμ, καὶ ὄνομα τῇ μητρὶ αὐτοῦ Ἀζουβὰ, θυγάτηρ Σαλί.
\vs{32}Καὶ ἐπορεύθη ἐν ταῖς ὁδοῖς τοῦ πατρὸς αὐτοῦ Ἀσὰ, καὶ οὐκ ἐξέκλινε τοῦ ποιῆσαι τὸ εὐθὲς ἐνώπιον Κυρίου.
\vs{33}Ἀλλὰ καὶ τὰ ὑψηλὰ ἔτι ὑπῆρχε, καὶ ἔτι ὁ λαὸς οὐ κατεύθυνε τὴν καρδίαν αὐτῶν πρὸς Κύριον τὸν Θεὸν τῶν πατέρων αὐτῶν.

\vs{34}Καὶ οἱ λοιποὶ λόγοι Ἰωσαφὰτ οἱ πρῶτοι καὶ οἱ ἔσχατοι, ἰδοὺ γεγραμμένοι ἐν λόγοις Ἰηοὺ τοῦ Ἀνανὶ, ὃς κατέγραψε βιβλίον βασιλέων Ἰσραήλ.

\vs{35}Καὶ μετὰ ταῦτα ἐκοινώνησεν Ἰωσαφὰτ βασιλεὺς Ἰούδα πρὸς Ὀχοζίαν βασιλέα Ἰσραὴλ, καὶ οὗτος ἠνόμησεν·
\vs{36}ἐν τῷ ποιῆσαι καὶ πορευθῆναι πρὸς αὐτὸν, τοῦ ποιῆσαι πλοῖα τοῦ πορευθῆναι εἰς Θαρσεῖς· καὶ ἐποίησε πλοῖα ἐν Γασίων Γαβέρ.
\vs{37}Καὶ προεφήτευσεν Ἐλιέζερ ὁ τοῦ Δωδία ἀπὸ Μαρισὴς ἐπὶ Ἰωσαφὰτ, λέγων, ὡς ἐφιλίασας τῷ Ὀχοζίᾳ, ἔθραυσε Κύριος τὸ ἔργον σου, καὶ συνετρίβη τὰ πλοῖά σου· καὶ οὐκ ἐδυνάσθη πορευθῆναι εἰς Θαρσεῖς.

\ch{21}
Καὶ ἐκοιμήθη Ἰωσαφὰτ μετὰ τῶν πατέρων αὐτοῦ, καὶ ἐτάφη ἐν πόλει Δαυίδ. Καὶ ἐβασίλευσεν Ἰωρὰν υἱὸς αὐτοῦ ἀντʼ αὐτοῦ.
\vs{2}Καὶ αὐτῷ ἀδελφοὶ υἱοὶ Ἰωσαφὰτ ἓξ, Ἀζαρίας, καὶ Ἰεϊὴλ, καὶ Ζαχαρίας, καὶ Ἀζαρίας, καὶ Μιχαὴλ, καὶ Ζαφατίας· πάντες οὗτοι υἱοὶ Ἰωσαφὰτ βασιλέως Ἰούδα.
\vs{3}Καὶ ἔδωκεν αὐτοῖς ὁ πατὴρ αὐτῶν δόματα πολλὰ, ἀργύριον καὶ χρυσίον, καὶ ὅπλα μετὰ τῶν πόλεων τετειχισμένων ἐν Ἰούδα, καὶ τὴν βασιλείαν ἔδωκε τῷ Ἰωρὰμ, ὅτι οὗτος ὁ πρωτότοκος.
\vs{4}Καὶ ἀνέστη Ἰωρὰμ ἐπὶ τὴν βασιλείαν αὐτοῦ, καὶ ἐκραταιώθη, καὶ ἀπέκτεινε πάντας τοὺς ἀδελφοὺς αὐτοῦ ἐν ῥομφαίᾳ, καὶ ἀπὸ τῶν ἀρχόντων Ἰσραήλ.

\vs{5}Ὄντος αὐτοῦ τριάκοντα καὶ δύο ἐτῶν, κατέστη Ἰωρὰμ ἐπὶ τὴν βασιλείαν αὐτοῦ, καὶ ὀκτὼ ἔτη ἐβασίλευσεν ἐν Ἱερουσαλήμ.
\vs{6}Καὶ ἐπορεύθη ἐν ὁδῷ βασιλέων Ἰσραὴλ, ὡς ἐποίησεν οἶκος Ἀχαὰβ, ὅτι θυγάτηρ Ἀχαὰβ ἦν αὐτοῦ γυνὴ, καὶ ἐποίησε τὸ πονηρὸν ἐναντίον Κυρίου.
\vs{7}Καὶ οὐκ ἐβούλετο Κύριος ἐξολοθρεῦσαι τὸν οἶκον Δαυὶδ, διὰ τὴν διαθήκην ἣν διέθετο τῷ Δαυὶδ, καὶ ὡς εἶπεν αὐτῷ δοῦναι αὐτῷ λύχνον καὶ τοῖς υἱοῖς αὐτοῦ πάσας τὰς ἡμέρας.

\vs{8}Ἐν ταῖς ἡμέραις ἐκείναις ἀπέστη Ἐδὼμ ἀπὸ τοῦ Ἰούδα, καὶ ἐβασίλευσαν ἐφʼ ἑαυτοὺς βασιλέα.
\vs{9}Καὶ ᾤχετο Ἰωρὰμ μετὰ τῶν ἀρχόντων, καὶ πᾶσα ἡ ἵππος μετʼ αὐτοῦ· καὶ ἐγένετο καὶ ἠγέρθη νυκτὸς, καὶ ἐπάταξεν Ἐδὼμ τὸν κυκλοῦντα αὐτὸν, καὶ τοὺς ἄρχοντας τῶν ἁρμάτων, καὶ ἔφυγεν ὁ λαὸς εἰς τὰ σκηνώματα αὐτῶν.
\vs{10}Καὶ ἀπέστη ἀπὸ Ἰούδα Ἐδὼμ ἕως τῆς ἡμέρας ταύτης· τότε ἀπέστη Λομνὰ ἐν τῷ καιρῷ ἐκείνῳ ἀπὸ χειρὸς αὐτοῦ, ὅτι ἐγκατέλιπε Κύριον τὸν Θεὸν τῶν πατέρων αὐτοῦ.
\vs{11}Καὶ γὰρ αὐτὸς ἐποίησεν ὑψηλὰ ἐν ταῖς πόλεσιν Ἰούδα, καὶ ἐξεπόρνευσε τοὺς κατοικοῦντας ἐν Ἱερουσαλὴμ, καὶ ἀπεπλάνησε τὸν Ἰούδαν.

\vs{12}Καὶ ἦλθεν αὐτῷ ἐν γραφῇ παρὰ Ἠλιοὺ τοῦ προφήτου, λέγων, τάδε λέγει Κύριος Θεὸς Δαυὶδ τοῦ πατρός σου, ἀνθʼ ὧν οὐκ ἐπορεύθης ἐν ὁδῷ Ἰωσαφὰτ τοῦ πατρός σου, καὶ ἐν ὁδοῖς Ἀσὰ βασιλέως Ἰούδα,
\vs{13}καὶ ἐπορεύθης ἐν ὁδοῖς βασιλέων Ἰσραὴλ, καὶ ἐξεπόρνευσας τὸν Ἰούδαν καὶ τοὺς κατοικοῦντας ἐν Ἱερουσαλὴμ, ὡς ἐξεπόρνευσεν οἶκος Ἀχαὰβ, καὶ τοὺς ἀδελφούς σου υἱοὺς τοῦ πατρός σου τοὺς ἀγαθοὺς ὑπὲρ σὲ ἀπέκτεινας,
\vs{14}ἰδοὺ Κύριος πατάξει σε πληγὴν μεγάλην ἐν τῷ λαῷ σου, καὶ ἐν τοῖς υἱοῖς σου, καὶ ἐν γυναιξί σου, καὶ ἐν πάσῃ τῇ ἀποσκευῇ σου.
\vs{15}Καὶ σὺ ἐν μαλακίᾳ πονηρᾷ, ἐν νόσῳ κοιλίας, ἕως οὗ ἐξέλθῃ ἡ κοιλία σου μετὰ τῆς μαλακίας ἐξ ἡμερῶν εἰς ἡμέρας.

\vs{16}Καὶ ἐπήγειρε Κύριος ἐπὶ Ἰωρὰμ τοὺς ἀλλοφύλους, καὶ τοὺς Ἄραβας, καὶ τοὺς ὁμόρους τῶν Αἰθιόπων.
\vs{17}Καὶ ἀνέβησαν ἐπὶ Ἰούδαν, καὶ κατεδυνάστευον, καὶ ἀπέστρεψαν πᾶσαν τὴν ἀποσκευὴν ἣν εὗρον ἐν οἴκῳ τοῦ βασιλέως, καὶ τοὺς υἱοὺς αὐτοῦ, καὶ τὰς θυγατέρας αὐτοῦ, καὶ οὐ κατελείφθη αὐτῷ υἱὸς, ἀλλʼ ἢ Ὀχοζίας ὁ μικρότατος τῶν υἱῶν αὐτοῦ.
\vs{18}Καὶ μετὰ ταῦτα πάντα ἐπάταξεν αὐτὸν Κύριος εἰς τὴν κοιλίαν μαλακίαν ᾗ οὐκ ἔστιν ἰατρεία.
\vs{19}Καὶ ἐγένετο ἐξ ἡμερῶν εἰς ἡμέρας· καὶ ὡς ἦλθε καιρὸς τῶν ἡμερῶν ἡμέρας δύο, ἐξῆλθεν ἡ κοιλία αὐτοῦ μετὰ τῆς νόσου, καὶ ἀπέθανεν ἐν μαλακίᾳ πονηρᾷ· καὶ οὐκ ἐποίησεν ὁ λαὸς αὐτοῦ ἐκφορὰν, καθὼς ἐκφορὰν πατέρων αὐτοῦ.
\vs{20}Ἦν τριάκοντα καὶ δύο ἐτῶν ὅτε ἐβασίλευσε, καὶ ὀκτὼ ἔτη ἐβασίλευσεν ἐν Ἱερουσαλήμ· καὶ ἐπορεύθη οὐκ ἐν ἐπαίνῳ, καὶ ἐτάφη ἐν πόλει Δαυὶδ, καὶ οὐκ ἐν τάφοις τῶν βασιλέων.

\ch{22}
Καὶ ἐβασίλευσαν οἱ κατοικοῦντες ἐν Ἱερουσαλὴμ τὸν Ὀχοζίαν υἱὸν αὐτοῦ τὸν μικρὸν ἀντʼ αὐτοῦ, ὅτι πάντας τοὺς πρεσβυτέρους ἀπέκτεινε τὸ ἐπελθὸν ἐπʼ αὐτοὺς λῃστήριον, οἱ Ἄραβες καὶ οἱ Ἀλιμαζονεῖς· καὶ ἐβασίλευσεν Ὀχοζίας υἱὸς Ἰωρὰμ βασιλέως Ἰούδα.

\vs{2}Ὢν ἐτῶν εἴκοσι Ὀχοζίας ἐβασίλευσε, καὶ ἐνιαυτὸν ἕνα ἐβασίλευσεν ἐν Ἱερουσαλὴμ, καὶ ὄνομα τῇ μητρὶ αὐτοῦ Γοθολία, θυγάτηρ Ἀμβρί.
\vs{3}Καὶ οὗτος ἐπορεύθη ἐν ὁδῷ οἴκου Ἀχαὰβ, ὅτι μήτηρ αὐτοῦ ἦν σύμβουλος τοῦ ἁμαρτάνειν.
\vs{4}Καὶ ἐποίησε τὸ πονηρὸν ἐναντίον Κυρίου ὡς οἶκος Ἀχαὰβ, ὅτι αὐτοὶ ἦσαν αὐτῷ σύμβουλοι μετὰ τὸ ἀποθανεῖν τὸν πατέρα αὐτοῦ, τοῦ ἐξολοθρεῦσαι αὐτὸν,
\vs{5}καὶ ἐν ταῖς βουλαῖς αὐτῶν ἐπορεύθη· καὶ ἐπορεύθη μετὰ Ἰωρὰμ υἱοῦ Ἀχαὰβ βασιλέως Ἰσραὴλ εἰς πόλεμον ἐπὶ Ἀζαὴλ βασιλέα Συρίας εἰς Ῥαμὼθ Γαλαάδ· καὶ ἐπάταξαν οἱ τοξόται τὸν Ἰωράμ.
\vs{6}Καὶ ἐπέστρεψεν Ἰωρὰμ τοῦ ἰατρευθῆναι εἰς Ἰεζράελ ἀπὸ τῶν πληγῶν ὧν ἐπάταξαν αὐτὸν οἱ Σύροι ἐν Ῥαμὼθ ἐν τῷ πολεμεῖν αὐτὸν πρὸς Ἀζαὴλ βασιλέα Συρίας.

Καὶ Ὀχοζίας υἱὸς Ἰωρὰμ βασιλεὺς Ἰούδα κατέβη θεάσασθαι τὸν Ἰωρὰμ υἱὸν Ἀχαὰβ εἰς Ἰεζράελ, ὅτι ἠῥῥώστει.
\vs{7}Καὶ παρὰ τοῦ Θεοῦ ἐγένετο καταστροφὴ Ὀχοζία ἐλθεῖν πρὸς Ἰωράμ· καὶ ἐν τῷ ἐλθεῖν αὐτὸν, ἐξῆλθε μετʼ αὐτοῦ Ἰωρὰμ πρὸς Ἰηοὺ υἱὸν Ναμεσσεῒ χριστὸν Κυρίου εἰς τὸν οἶκον Ἀχαάβ.

\vs{8}Καὶ ἐγένετο ὡς ἐξεδίκησεν Ἰηοὺ τὸν οἶκον Ἀχαὰβ, καὶ εὗρε τοὺς ἄρχοντας Ἰούδα καὶ τοὺς ἀδελφοὺς Ὀχοζίου λειτουργοῦντας τῷ Ὀχοζίᾳ, καὶ ἀπέκτεινεν αὐτούς.
\vs{9}Καὶ εἶπε τοῦ ζητῆσαι τὸν Ὀχοζίαν· καὶ κατέλαβον αὐτὸν ἰατρευόμενον ἐν Σαμαρείᾳ, καὶ ἤγαγον αὐτὸν πρὸς Ἰηοὺ, καὶ ἀπέκτεινεν αὐτὸν· καὶ ἔθαψαν αὐτὸν, ὅτι εἶπαν, υἱὸς Ἰωσαφάτ ἐστιν, ὃς ἐζήτησε τὸν Κύριον ἐν ὅλῃ τῇ καρδίᾳ αὐτοῦ.

Καὶ οὐκ ἦν ἐν οἴκῳ Ὀχοζίᾳ κατισχῦσαι δύναμιν περὶ τῆς βασιλείας.
\vs{10}Καὶ Γοθολία ἡ μήτηρ Ὀχοζίου εἶδεν ὅτι τέθνηκεν ὁ υἱὸς αὐτῆς, καὶ ἠγέρθη καὶ ἀπώλεσε πᾶν τὸ σπέρμα τῆς βασιλείας ἐν οἴκῳ Ἰούδα.
\vs{11}Καὶ ἔλαβεν Ἰωσαβεὲθ θυγάτηρ τοῦ βασιλέως τὸν Ἰωὰς υἱὸν Ὀχοζίου καὶ ἔκλεψεν αὐτὸν ἐκ μέσου υἱῶν τοῦ βασιλέως τῶν θανατουμένων, καὶ ἔδωκεν αὐτὸν καὶ τὴν τροφὸν αὐτοῦ εἰς ταμεῖον τῶν κλινῶν, καὶ ἔκρυψεν αὐτὸν Ἰωσαβεὲθ θυγάτηρ τοῦ βασιλέως Ἰωρὰμ ἀδελφὴ Ὀχοζίου γυνὴ Ἰωδαὲ τοῦ ἱερέως, καὶ ἔκρυψεν αὐτὸν ἀπὸ προσώπου τῆς Γοθολίας, καὶ οὐκ ἀπέκτεινεν αὐτόν.
\vs{12}Καὶ ἦν μετʼ αὐτοῦ ἐν οἴκῳ τοῦ Θεοῦ κατακεκρυμμένος ἓξ ἔτη, καὶ Γοθολία ἐβασίλευσεν ἐπὶ τῆς γῆς.

\ch{23}
Καὶ ἐν τῷ ἔτει τῷ ὀγδόῳ ἐκραταίωσεν Ἰωδαὲ, καὶ ἔλαβε τοὺς ἑκατοντάρχους, τὸν Ἀζαρίαν υἱὸν Ἰωρὰμ, καὶ τὸν Ἰσμαὴλ υἱὸν Ἰωανὰν, καὶ τὸν Ἀζαρίαν υἱὸν Ὠβὴδ, καὶ τὸν Μαασαίαν υἱὸν Ἀδία, καὶ τὸν Ἐλισαφὰν υἱὸν Ζαχαρίου, μεθʼ ἑαυτοῦ εἰς οἶκον Κυρίου.
\vs{2}Καὶ ἐκύκλωσαν τὸν Ἰούδαν, καὶ συνήγαγον τοὺς Λευίτας ἐκ πασῶν τῶν πόλεων Ἰούδα, καὶ ἄρχοντας πατριῶν τοῦ Ἰσραὴλ, καὶ ἦλθον εἰς Ἱερουσαλήμ.
\vs{3}Καὶ διέθεντο πᾶσα ἡ ἐκκλησία Ἰούδα διαθήκην ἐν οἴκῳ τοῦ Θεοῦ μετὰ τοῦ βασιλέως· καὶ ἔδειξεν αὐτοῖς τὸν υἱὸν τοῦ βασιλέως, καὶ εἶπεν αὐτοῖς, ἰδοὺ ὁ υἱὸς τοῦ βασιλέως βασιλευσάτω, καθὼς ἐλάλησε Κύριος ἐπὶ τὸν οἶκον Δαυίδ.
\vs{4}Νῦν ὁ λόγος οὗτος, ὃν ποιήσετε· τὸ τρίτον ἐξ ὑμῶν εἰσπορευέσθωσαν τὸ σάββατον τῶν ἱερέων καὶ τῶν Λευιτῶν καὶ εἰς τὰς πύλας τῶν εἰσόδων,
\vs{5}καὶ τὸ τρίτον ἐν οἴκῳ τοῦ βασιλέως, καὶ τὸ τρίτον ἐν τῇ πύλῃ τῇ μέσῃ, καὶ πᾶς ὁ λαὸς ἐν αὐλαῖς οἴκου Κυρίου.
\vs{6}Καὶ μὴ εἰσελθέτω εἰς οἶκον Κυρίου, ἐὰν μὴ οἱ ἱερεῖς καὶ οἱ Λευῖται καὶ οἱ λειτουργοῦντες τῶν Λευιτῶν· αὐτοὶ εἰσελεύσονται ὅτι ἅγιοί εἰσι, καὶ πᾶς ὁ λαὸς φυλασσέτω φυλακὰς Κυρίου.
\vs{7}Καὶ κυκλώσουσιν οἱ Λευῖται τὸν βασιλέα κύκλῳ ἀνδρὸς σκεῦος σκεῦος ἐν χειρὶ αὐτοῦ, καὶ ὁ εἰσπορευόμενος εἰς τὸν οἶκον ἀποθανεῖται, καὶ ἔσονται μετὰ τοῦ βασιλέως ἐκπορευομένου καὶ εἰσπορευομένου αὐτοῦ.

\vs{8}Καὶ ἐποίησαν οἱ Λευῖται καὶ πᾶς Ἰούδα κατὰ πάντα ὅσα ἐνετείλατο αὐτοῖς Ἰωδαὲ ὁ ἱερεύς· καὶ ἔλαβον ἕκαστος τοὺς ἄνδρας αὐτοῦ ἀπʼ ἀρχῆς τοῦ σαββάτου ἕως ἐξόδου τοῦ σαββάτου, ὅτι οὐ κατέλυσεν Ἰωδαὲ ὁ ἱερεὺς τὰς ἐφημερίας.
\vs{9}Καὶ ἔδωκεν Ἰωδαὲ τὰς μαχαίρας καὶ τοὺς θυρεοὺς καὶ τὰ ὅπλα ἃ ἦν τοῦ βασιλέως Δαυὶδ ἐν οἴκῳ τοῦ θεοῦ.
\vs{10}Καὶ ἔστησε τὸν λαὸν πάντα ἕκαστον ἐν τοῖς ὅπλοις αὐτοῦ, ἀπὸ τῆς ὠμίας τοῦ οἴκου τῆς δεξιᾶς ἕως τῆς ὠμίας τῆς ἀριστερᾶς τοῦ θυσιαστηρίου καὶ τοῦ οἴκου, ἐπὶ τὸν βασιλέα κύκλῳ.
\vs{11}Καὶ ἐξήγαγε τὸν υἱὸν τοῦ βασιλέως, καὶ ἔδωκεν ἐπʼ αὐτὸν τὸ βασίλειον καὶ τὰ μαρτύρια, καὶ ἐβασίλευσαν καὶ ἔχρισαν αὐτὸν Ἰωδαὲ ὁ ἱερεὺς καὶ οἱ υἱοὶ αὐτοῦ, καὶ εἶπαν, ζήτω ὁ βασιλεύς.

\vs{12}Καὶ ἤκουσε Γοθολία τὴν φωνὴν τοῦ λαοῦ τρεχόντων, καὶ ἐξομολογουμένων, καὶ αἰνούντων τὸν βασιλέα· καὶ εἰσῆλθε πρὸς τὸν βασιλέα εἰς οἶκον Κυρίου.
\vs{13}Καὶ εἶδε, καὶ ἰδοὺ ὁ βασιλεὺς ἐπὶ τῆς στάσεως αὐτοῦ, καὶ ἐπὶ τῆς εἰσόδου οἱ ἄρχοντες καὶ οἱ σάλπιγγες· καὶ οἱ ἄρχοντες περὶ τὸν βασιλέα, καὶ πᾶς ὁ λαὸς τῆς γῆς ηὐφράνθη, καὶ ἐσάλπισαν ταῖς σάλπιγξι, καὶ οἱ ᾄδοντες ἐν τοῖς ὀργάνοις ᾠδοὶ, καὶ ὑμνοῦντες αἶνον· καὶ διέῥῥηξε Γοθολία τὴν στολὴν αὐτῆς, καὶ ἐβόησεν, ἐπιτιθέμενοι ἐπιτίθεσθε.
\vs{14}Καὶ ἐξῆλθεν Ἰωδαὲ ὁ ἱερεὺς, καὶ ἐνετείλατο Ἰωδαὲ ὁ ἱερεὺς τοῖς ἑκατοντάρχοις, καὶ τοῖς ἀρχηγοῖς τῆς δυνάμεως, καὶ εἶπεν αὐτοῖς, ἐκβάλετε αὐτὴν ἐκτὸς τοῦ οἴκου, καὶ εἰσέλθατε ὀπίσω αὐτῆς, καὶ ἀποθανέτω μαχαίρᾳ, ὅτι εἶπεν ὁ ἱερεὺς, μὴ ἀποθανέτω ἐν οἴκῳ Κυρίου.
\vs{15}Καὶ ἔδωκαν αὐτῇ ἄνεσιν, καὶ διῆλθε διὰ τῆς πύλης τῶν ἱππέων τοῦ οἴκου τοῦ βασιλέως, καὶ ἐθανάτωσαν αὐτὴν ἐκεῖ.

\vs{16}Καὶ διέθετο Ἰωδαὲ διαθήκην ἀναμέσον αὐτοῦ καὶ τοῦ λαοῦ καὶ τοὺ βασιλέως, εἶναι λαὸν τῷ Κυρίῳ.
\vs{17}Καὶ εἰσῆλθε πᾶς ὁ λαὸς τῆς γῆς εἰς οἶκον Βάαλ, καὶ κατέσπασαν αὐτὸν καὶ τὰ θυσιαστήρια αὐτοῦ, καὶ τὰ εἴδωλα αὐτοῦ ἐλέπτυναν, καὶ τὸν Ματθὰν ἱερέα Βάαλ ἐθανάτωσαν ἐναντίον τῶν θυσιαστηρίων αὐτοῦ.
\vs{18}Καὶ ἐνεχείρισεν Ἰωδαὲ ὁ ἱερεὺς τὰ ἔργα οἴκου Κυρίου διὰ χειρὸς ἱερέων καὶ Λευιτῶν, καὶ ἀνέστησε τὰς ἐφημερίας τῶν ἱερέων καὶ τῶν Λευιτῶν, ἃς διέστειλε Δαυὶδ ἐπὶ τὸν οἶκον Κυρίου, καὶ ἀνενέγκαι ὁλοκαυτώματα Κυρίῳ, καθὼς γέγραπται ἐν νόμῳ Μωυσῆ, ἐν εὐφροσύνῃ καὶ ἐν ᾠδαῖς διὰ χειρὸς Δαυίδ.
\vs{19}Καὶ ἔστησαν οἱ πυλωροὶ ἐπὶ τὰς πύλας οἴκου Κυρίου, καὶ οὐκ εἰσελεύσεται ἀκάθαρτος εἰς πᾶν πρᾶγμα.
\vs{20}Καὶ ἔλαβε τοὺς πατριάρχας, καὶ τοὺς δυνατοὺς, καὶ τοὺς ἄρχοντας τοῦ λαοῦ, καὶ πάντα τὸν λαὸν τῆς γῆς, καὶ ἐπεβίβασαν τὸν βασιλέα εἰς οἶκον Κυρίου, καὶ εἰσῆλθε διὰ τῆς πύλης τῆς ἐσωτέρας εἰς τὸν οἶκον τοῦ βασιλέως, καὶ ἐκάθισαν τὸν βασιλέα ἐπὶ τοῦ θρόνου τῆς βασιλείας.
\vs{21}Καὶ ηὐφράνθη πᾶς ὁ λαὸς τῆς γῆς, καὶ ἡ πόλις ἡσύχασε, καὶ τὴν Γοθολίαν ἐθανάτωσαν.

\ch{24}
Ὢν ἐτῶν ἑπτὰ Ἰωὰς ἐν τῷ βασιλεύειν αὐτὸν, καὶ τεσσαράκοντα ἔτη ἐβασίλευσεν ἐν Ἱερουσαλὴμ. καὶ ὄνομα τῇ μητρὶ αὐτοῦ Σαβιὰ ἐκ Βηρσαβεέ.
\vs{2}Καὶ ἐποίησεν Ἰωὰς τὸ εὐθὲς ἐνώπιον Κυρίου πάσας τὰς ἡμέρας Ἰωδαὲ τοῦ ἱερέως.
\vs{3}Καὶ ἔλαβεν Ἰωδαὲ δύο γυναῖκας ἑαυτῷ, καὶ ἐγέννησαν υἱοὺς καὶ θυγατέρας.

\vs{4}Καὶ ἐγένετο μετὰ ταῦτα, καὶ ἐγένετο ἐπὶ καρδίαν Ἰωὰς ἐπισκευάσαι τὸν οἶκον Κυρίου.
\vs{5}Καὶ συνήγαγε τοὺς ἱερεῖς καὶ τοὺς Λευίτας, καὶ εἶπεν αὐτοῖς, ἐξέλθατε εἰς τὰς πόλεις Ἰούδα, καὶ συναγάγετε ἀπὸ παντὸς Ἰσραὴλ ἀργύριον κατισχῦσαι τὸν οἶκον Κυρίου ἐνιαυτὸν κατʼ ἐνιαυτὸν, καὶ σπεύσατε λαλῆσαι· καὶ οὐκ ἔσπευσαν οἱ Λευῖται.

\vs{6}Καὶ ἐκάλεσεν ὁ βασιλεὺς Ἰωὰς τὸν Ἰωδαὲ τὸν ἄρχοντα, καὶ εἶπεν αὐτῷ, διατί οὐκ ἐπεσκέψω περὶ τῶν Λευιτῶν τοῦ εἰσενέγκαι ἀπὸ Ἰούδα καὶ Ἱερουσαλὴμ τὸ κεκριμένον ὑπὸ Μωυσῆ ἀνθρώπου τοῦ Θεοῦ, ὅτι ἐξεκκλησίασε τὸν Ἰσραὴλ εἰς τὴν σκηνὴν τοῦ μαρτυρίου;
\vs{7}Ὅτι Γοθολία ἦν ἡ ἄνομος, καὶ οἱ υἱοὶ αὐτῆς κατέσπασαν τὸν οἶκον τοῦ Θεοῦ· καὶ γὰρ τὰ ἅγια οἴκου Κυρίου ἐποίησαν ταῖς Βααλίμ.

\vs{8}Καὶ εἶπεν ὁ βασιλεὺς, γενηθήτω γλωσσόκομον, καὶ τεθήτω ἐν πύλῃ οἴκου Κυρίου ἔξω.
\vs{9}Καὶ κηρυξάτωσαν ἐν Ἰούδα καὶ ἐν Ἱερουσαλὴμ, εἰσενέγκαι Κυρίῳ καθὼς εἶπε Μωυσῆς παῖς τοῦ Θεοῦ ἐπὶ τὸν Ἰσραὴλ ἐν τῇ ἐρήμῳ.
\vs{10}Καὶ ἔδωκαν πάντες ἄρχοντες καὶ πᾶς ὁ λαὸς, καὶ εἰσέφερον καὶ ἐνέβαλον εἰς τὸ γλωσσόκομον ἕως οὗ ἐπληρώθη.
\vs{11}Καὶ ἐγένετο ὡς εἰσέφερον τὸ γλωσσόκομον πρὸς τοὺς προστάτας τοῦ βασιλέως διὰ χειρὸς τῶν Λευιτῶν, καὶ ὡς εἶδον ὅτι ἐπλεόνασε τὸ ἀργύριον, καὶ ἦλθεν ὁ γραμματεὺς τοῦ βασιλέως καὶ ὁ προστάτης τοῦ ἱερέως τοῦ μεγάλου, καὶ ἐξεκένωσεν τὸ γλωσσόκομον, καὶ κατέστησαν εἰς τὸν τόπον αὐτοῦ· οὕτως ἐποίουν ἡμέραν ἐξ ἡμέρας, καὶ συνήγαγον ἀργύριον πολύ.
\vs{12}Καὶ ἔδωκεν αὐτὸ ὁ βασιλεὺς καὶ Ἰωδαὲ ὁ ἱερεὺς τοῖς ποιοῦσι τὰ ἔργα εἰς ἐργασίαν οἴκου Κυρίου· καὶ ἐμισθοῦντο λατόμους καὶ τέκτονας ἐπισκευάσαι τὸν οἶκον Κυρίου, καὶ χαλκεῖς σιδήρου καὶ χαλκοῦ ἐπισκευάσαι τὸν οἶκον Κυρίου.
\vs{13}Καὶ ἐποίουν οἱ ποιοῦντες τὰ ἔργα, καὶ ἀνέβη μῆκος τῶν ἔργων ἐν χερσὶν αὐτῶν, καὶ ἀνέστησαν τὸν οἶκον Κυρίου ἐπὶ τὴν στάσιν αὐτοῦ, καὶ ἐνίσχυσαν.
\vs{14}Καὶ ὡς συνετέλεσαν, ἤνεγκαν πρὸς τὸν βασιλέα καὶ πρὸς Ἰωδαὲ τὸ κατάλοιπον τοῦ ἀργυρίου, καὶ ἐποίησαν σκεύη εἰς οἶκον Κυρίου, σκεύη λειτουργικὰ ὁλοκαυτωμάτων, καὶ θυΐσκας χρυσᾶς καὶ ἀργυρᾶς, καὶ ἀνήνεγκαν ὁλοκαυτώσεις ἐν οἴκῳ Κυρίου διαπαντὸς πάσας τὰς ἡμέρας Ἰωδαέ.

\vs{15}Καὶ ἐγήρασεν Ἰωδαὲ πλήρης ἡμερῶν, καὶ ἐτελεύτησεν ὢν ἑκατὸν καὶ τριάκοντα ἐτῶν ἐν τῷ τελευτᾷν αὐτόν.
\vs{16}Καὶ ἔθαψαν αὐτὸν ἐν πόλει Δαυὶδ μετὰ τῶν βασιλέων, ὅτι ἐποίησεν ἀγαθωσύνην μετὰ Ἰσραὴλ καὶ μετὰ τοῦ Θεοῦ καὶ τοῦ οἴκου αὐτοῦ.

\vs{17}Καὶ ἐγένετο μετὰ τὴν τελευτὴν Ἰωδαὲ εἰσῆλθον οἱ ἄρχοντες Ἰούδα, καὶ προσεκύνησαν τὸν βασιλέα· τότε ἐπήκουσεν αὐτοῖς ὁ βασιλεύς.
\vs{18}Καὶ ἐγκατέλιπον τὸν οἶκον Κυρίον Θεοῦ τῶν πατέρων αὐτῶν, καὶ ἐδούλευον ταῖς Ἀστάρταις καὶ τοῖς εἰδώλοις, καὶ ἐγένετο ὀργὴ ἐπὶ Ἰούδαν καὶ ἑπὶ Ἱερουσαλὴμ ἐν τῇ ἡμέρᾳ ταύτῃ.
\vs{19}Καὶ ἀπέστειλε πρὸς αὐτοὺς προφήτας ἐπιστρέψαι πρὸς Κύριον, καὶ οὐκ ἤκουσαν· καὶ διεμαρτύρατο αὐτοῖς, καὶ οὐχ ὑπήκουσαν.
\vs{20}Καὶ πνεῦμα Θεοῦ ἐνέδυσε τὸν Ἀζαρίαν τὸν τοῦ Ἰωδαὲ τὸν ἱερέα, καὶ ἀνέστη ἐπάνω τοῦ λαοῦ, καὶ εἶπε, τάδε λέγει Κύριος, τί παραπορεύεσθε τὰς ἐντολὰς Κυρίου; καὶ οὐκ εὐοδωθήσεσθε· ὅτι ἐγκατελίπετε τὸν Κύριον, καὶ ἐγκαταλείψει ὑμᾶς.
\vs{21}Καὶ ἐπέθεντο αὐτῷ, καὶ ἐλιθοβόλησαν αὐτὸν διʼ ἐντολῆς Ἰωὰς τοῦ βασιλέως ἐν αὐλῇ οἴκου Κυρίου.
\vs{22}Καὶ οὐκ ἐμνήσθη Ἰωὰς τοῦ ἐλέους οὗ ἐποίησεν Ἰωδαὲ ὁ πατὴρ αὐτοῦ μετʼ αὐτοῦ, καὶ ἐθανάτωσε τὸν υἱὸν αὐτοῦ· καὶ ὡς ἀπέθνησκεν, εἶπεν, ἴδοι Κύριος καὶ κρινάτω.

\vs{23}Καὶ ἐγένετο μετὰ τὴν συντέλειαν τοῦ ἐνιαυτοῦ ἀνέβη ἐπʼ αὐτὸν δύναμις Συρίας, καὶ ἦλθεν ἐπὶ Ἰούδαν καὶ ἐπὶ Ἱερουσαλὴμ, καὶ κατέφθειραν πάντας τοὺς ἄρχοντας τοῦ λαοῦ ἐν τῷ λαῷ, καὶ πάντα τὰ σκῦλα αὐτῶν ἀπέστειλαν τῷ βασιλεῖ Δαμασκοῦ.
\vs{24}Ὅτι ἐν ὀλίγοις ἀνδράσιν παρεγένετο δύναμις Συρίας, καὶ ὁ Θεὸς παρέδωκεν εἰς τὰς χεῖρας αὐτῶν δύναμιν πολλὴν σφόδρα, ὅτι ἐγκατέλιπον Κύριον τὸν Θεὸν τῶν πατέρων αὐτῶν· καὶ μετὰ Ἰωὰς ἐποίησε κρίματα.

\vs{25}Καὶ μετὰ τὸ ἀπελθεῖν αὐτοὺς ἀπʼ αὐτοῦ, ἐν τῷ ἐγκαταλιπεῖν αὐτὸν ἐν μαλακίαις μεγάλαις, καὶ ἐπέθεντο αὐτῷ οἱ παῖδες αὐτοῦ ἐν αἵμασιν υἱοῦ Ἰωδαὲ τοῦ ἱερέως, καὶ ἐθανάτωσαν αὐτὸν ἐπὶ τῆς κλίνης αὐτοῦ, καὶ ἀπέθανε· καὶ ἔθαψαν αὐτὸν ἐν πόλει Δαυὶδ, καὶ οὐκ ἔθαψαν αὐτὸν ἐν τῷ τάφῳ τῶν βασιλέων.
\vs{26}Καὶ οἱ ἐπιθέμενοι ἐπʼ αὐτὸν Ζαβὲδ ὁ τοῦ Σαμαὰθ ὁ Ἀμμανίτης, καὶ Ἰωζαβὲδ ὁ τοῦ Σαμαρὴθ ὁ Μωαβίτης,
\vs{27}καὶ οἱ υἱοὶ αὐτοῦ πάντες, καὶ προσῆλθον αὐτῷ οἱ πέντε· καὶ τὰ λοιπὰ ἰδοὺ γεγραμμένα ἐπὶ τὴν γραφὴν τῶν βασιλέων· καὶ ἐβασίλευσεν Ἀμασίας υἱὸς αὐτοῦ ἀντʼ αὐτοῦ.

\ch{25}
Ὢν εἴκοσι καὶ πέντε ἐτῶν ἐβασίλευσεν Ἀμασίας, καὶ εἰκοσιεννέα ἔτη ἐβασίλευσεν ἐν Ἱερουσαλὴμ, καὶ ὄνομα τῇ μητρὶ αὐτοῦ Ἰωαδαὲν ἀπὸ Ἱερουσαλήμ.
\vs{2}Καὶ ἐποίησε τὸ εὐθὲς ἐνώπιον Κυρίου, ἀλλʼ οὐκ ἐν καρδίᾳ πλήρει.
\vs{3}Καὶ ἐγένετο ὡς κατέστη ἡ βασιλεία ἐν χειρὶ αὐτοῦ, καὶ ἐθανάτωσε τοὺς παῖδας αὐτοῦ τοὺς φονεύσαντας τὸυ βασιλέα πατέρα αὐτοῦ·
\vs{4}Καὶ τοὺς υἱοὺς αὐτῶν οὐκ ἀπέκτεινε, κατὰ τὴν διαθήκην τοῦ νόμου Κύριου, καθὼς γέγραπται, ὡς ἐνετείλατο Κύριος, λέγων, οὐκ ἀποθανοῦνται πατέρες ὑπὲρ τέκνων, καὶ υἱοὶ οὐκ ἀποθανοῦνται ὑπὲρ πατέρων, ἀλλʼ ἢ ἕκαστος τῇ ἑαυτοῦ ἁμαρτίᾳ ἀποθανοῦνται.

\vs{5}Καὶ συνήγαγεν Ἀμασίας τὸν οἶκον Ἰούδα, καὶ ἀνέστησεν αὐτοὺς κατʼ οἴκους πατριῶν αὐτῶν εἰς χιλιάρχους καὶ ἑκατοντάρχους ἐν παντὶ Ἰούδᾳ καὶ Ἱερουσαλήμ· καὶ ἠρίθμησεν αὐτοὺς ἀπὸ εἰκοσαετοῦς καὶ ἐπάνω, καὶ εὗρεν αὐτοὺς τριακοσίας χιλιάδας ἐξελθεῖν εἰς πόλεμον δυνατοὺς, κρατοῦντας δόρυ καὶ θυρεόν.
\vs{6}Καὶ ἐμισθώσατο ἀπὸ Ἰσραὴλ ἑκατὸν χιλιάδας δυνατοὺς ἰσχύϊ ἑκατὸν ταλάντων ἀργυρίου.

\vs{7}Καὶ ἄνθρωπος τοῦ Θεοῦ ἦλθε πρὸς αὐτὸν, λέγων, βασιλεῦ, οὐ πορεύσεται μετὰ σοῦ δύναμις Ἰσραὴλ, ὅτι οὐκ ἔστι Κύριος μετὰ Ἰσραὴλ πάντων τῶν υἱῶν Ἐφραίμ.
\vs{8}Ὅτι ἐὰν ὑπολάβῃς κατισχῦσαι ἐν τούτοις, καὶ τροπώσεταί σε Κύριος ἐναντίον τῶν ἐχθρῶν, ὅτι ἐστὶ παρὰ Κυρίου καὶ ἰσχῦσαι καὶ τροπώσασθαι.
\vs{9}Καὶ εἶπεν Ἀμασίας τῷ ἀνθρώπῳ τοῦ Θεοῦ, καὶ τί ποιήσω τὰ ἑκατὸν τάλαντα ἃ ἔδωκα τῇ δυνάμει Ἰσραήλ; καὶ εἶπεν ὁ ἄνθρωπος τοῦ Θεοῦ, ἔστι τῷ Κυρίῳ δοῦναί σοι πλεῖστα τούτων.

\vs{10}Καὶ διεχώρισεν Ἀμασίας τῇ δυνάμει τῇ ἐλθούσῃ πρὸς αὐτὸν ἀπὸ Ἐφραὶμ, ἀπέλθεῖν εἰς τὸν τόπον αὐτῶν· καὶ ἐθυμώθησαν σφόδρα ἐπὶ Ἰούδαν, καὶ ἐπέστρεψαν εἰς τὸν τόπον αὐτῶν ἐν ὀργῇ θυμοῦ.
\vs{11}Καὶ Ἀμασίας κατίσχυσε καὶ παρέλαβε τὸν λαὸν αὐτοῦ, καὶ ἐπορεύθη εἰς τὴν κοιλάδα τῶν ἁλῶν, καὶ ἐπάταξεν ἐκεῖ τοὺς υἱοὺς Σηεὶρ, δέκα χιλιάδας.
\vs{12}Καὶ δέκα χιλιάδας ἐζώγρησαν οἱ υἱοὶ Ἰούδα, καὶ ἔφερον αὐτοὺς ἐπὶ τὸ ἄκρον τοῦ κρημνοῦ, καὶ κατεκρήμνιζον αὐτοὺς ἀπὸ τοῦ ἄκρου τοῦ κρημνοῦ, καὶ πάντες διεῤῥήγνυντο.
\vs{13}Καὶ υἱοὶ τῆς δυνάμεως οὓς ἀπέστρεψεν Ἀμασίας τοῦ μὴ πορευθῆναι μετʼ αὐτοῦ εἰς πόλεμον, καὶ ἐπέθεντο ἐπὶ τὰς πόλεις Ἰούδα ἀπὸ Σαμαρείας ἕως Βαιθωρών· καὶ ἐπάταξαν ἐν αὐτοῖς τρεῖς χιλιάδας, καὶ ἐσκύλευσαν σκῦλα πολλά.

\vs{14}Καὶ ἐγένετο μετὰ τὸ ἐλθεῖν Ἀμασίαν πατάξαντος τὴν Ἰδουμαίαν, καὶ ἤνεγκε πρὸς αὐτὸν τοὺς θεοὺς υἱῶν Σηεὶρ, καὶ ἔστησεν αὐτοὺς αὐτῷ εἰς θεοὺς, καὶ ἐναντίον αὐτῶν προσεκύνει, καὶ αὐτὸς αὐτοῖς ἔθυε.
\vs{15}Καὶ ἐγένετο ὀργὴ Κυρίου ἐπὶ Ἀμασίαν, καὶ ἀπέστειλεν αὐτῷ προφήτην, καὶ εἶπεν αὐτῷ, τί ἐζήτησας τοὺς θεοὺς τοῦ λαοῦ, οἳ οὐκ ἐξείλοντο τὸν λαὸν ἑαυτῶν ἐκ χειρός σου;
\vs{16}Καὶ ἐγένετο ἐν τῷ λαλῆσαι αὐτῷ πρὸς αὐτὸν, καὶ εἶπεν αὐτῷ, μὴ σύμβουλον τοῦ βασιλέως δέδωκά σε; πρόσεχε ἵνα μὴ μαστιγωθῇς· καὶ ἐσιώπησεν ὁ προφήτης, καὶ εἶπεν, ὅτι γινώσκω, ὅτι ἐβούλετο ἐπὶ σοὶ τοῦ καταφθεῖραί σε, ὅτι ἐποίησας τοῦτο, καὶ οὐκ ἐπήκουσας τῆς συμβουλίας μου.

\vs{17}Καὶ ἐβουλεύσατο Ἀμασίας ὁ βασιλεὺς Ἰούδα, καὶ ἀπέστειλε πρὸς Ἰωὰς υἱὸν Ἰωάχαζ υἱοῦ Ἰηοὺ βασιλέα Ἰσραὴλ, λέγων, δεῦρο, καὶ ὀφθῶμεν προσώποις.
\vs{18}Καὶ ἀπέστειλεν Ἰωὰς βασιλεὺς Ἰσραὴλ πρὸς Ἀμασίαν βασιλέα Ἰούδα, λέγων, ὁ ἀκχοὺχ ὁ ἐν τῷ Λιβάνῳ ἀπέστειλε πρὸς τὴν κέδρον τὴν ἐν τῷ Λιβάνῳ, λέγων, δὸς τὴν θυγατέρα σου τῷ υἱῷ μου εἰς γυναῖκα, καὶ ἰδοὺ ἐλεύσεται τὰ θηρία τοῦ ἀγροῦ τὰ ἐν τῷ Λιβάνῳ· καὶ ἦλθον τὰ θηρία, καὶ κατεπάτησαν τὸν ἀκχούχ.
\vs{19}Εἶπας, ἰδοὺ ἐπάταξα τὴν Ἰδουμαίαν, καὶ ἐπαίρει σε ἡ καρδία σου ἡ βαρεῖα· νῦν κάθισον ἐν οἴκῳ σου, καὶ ἱνατί συμβάλλεις ἐν κακίᾳ, καὶ πεσῇ σὺ καὶ Ἰούδας μετὰ σοῦ;

\vs{20}Καὶ οὐκ ἤκουσεν Ἀμασίας, ὅτι παρὰ Κυρίου ἐγένετο τοῦ παραδοῦναι αὐτὸν εἰς χεῖρας, ὅτι ἐξεζήτησε τοὺς θεοὺς τῶν Ἰδουμαίων.
\vs{21}Καὶ ἀνέβη Ἰωὰς βασιλεὺς Ἰσραὴλ, καὶ ὤφθησαν ἀλλήλοις αὐτὸς καὶ Ἀμασίας βασιλεὺς Ἰούδα ἐν Βαιθσαμὺς, ἥ ἐστι τοῦ Ἰούδα.
\vs{22}Καὶ ἐτροπώθη Ἰούδας κατὰ πρόσωπον Ἰσραὴλ, καὶ ἔφυγεν ἕκαστος εἰς τὸ σκήνωμα αὐτοῦ.
\vs{23}Καὶ τὸν Ἀμασίαν βασιλέα Ἰούδα τὸν τοῦ Ἰωὰς υἱοῦ Ἰωάχαζ κατέλαβεν Ἰωὰς βασιλεὺς Ἰσραὴλ ἐν Βαιθσαμὺς, καὶ εἰσήγαγεν αὐτὸν εἰς Ἱερουσαλὴμ· καὶ κατέσπασεν ἀπὸ τοῦ τείχους Ἱερουσαλὴμ ἀπὸ πύλης Ἐφραὶμ ἕως πύλης γωνίας τετρακοσίους πήχεις.
\vs{24}Καὶ πᾶν τὸ χρυσίον καὶ τὸ ἀργύριον, καὶ πάντα τὰ σκεύη τὰ εὑρεθέντα ἐν οἴκῳ Κυρίου καὶ παρὰ τῷ Ἀβδεδὸμ, καὶ τοὺς θησαυροὺς οἴκου τοῦ βασιλέως, καὶ τοὺς υἱοὺς τῶν συμμίξεων, καὶ ἐπέστρεψεν εἰς Σαμάρειαν.

\vs{25}Καὶ ἔζησεν Ἀμασίας ὁ τοῦ Ἰωὰς βασιλεὺς Ἰούδα μετὰ τὸ ἀποθανεῖν Ἰωὰς τὸν τοῦ Ἰωάχαζ βασιλέα Ἰσραὴλ ἔτη δεκαπέντε.
\vs{26}Καὶ οἱ λοιποὶ λόγοι Ἀμασίου οἱ πρῶτοι καὶ οἱ ἔσχατοι οὐκ ἰδοὺ γεγραμμένοι ἐπὶ βιβλίου βασιλέων Ἰούδα καὶ Ἰσραήλ;
\vs{27}Καὶ ἐν τῷ καιρῷ ᾧ ἀπέστη Ἀμασίας ἀπὸ Κυρίου, καὶ ἐπέθεντο αὐτῷ ἐπίθεσιν, καὶ ἔφυγεν ἀπὸ Ἱερουσαλὴμ εἰς Λαχίς· καὶ ἀπέστειλαν κατόπισθεν αὐτοῦ εἰς Λαχίς, καὶ ἐθανάτωσαν αὐτὸν ἐκεῖ.
\vs{28}Καὶ ἀνέλαβον αὐτὸν ἐπὶ τῶν ἵππων, καὶ ἔθαψαν αὐτὸν μετὰ τῶν πατέρων αὐτοῦ ἐν πόλει Δαυίδ.

\ch{26}
Καὶ ἔλαβε πᾶς ὁ λαὸς τῆς γῆς τὸν Ὀζίαν, καὶ αὐτὸς υἱὸς ἑκκαίδεκα ἐτῶν, καὶ ἐβασίλευσαν αὐτὸν ἀντὶ τοῦ πατρὸς αὐτοῦ Ἀμασίου.
\vs{2}Αὐτὸς ᾠκοδόμησε τὴν Αἰλάθ, αὐτὸς ἐπέστρεψεν αὐτὴν τῷ Ἰούδα, μετὰ τὸ κοιμηθῆναι τὸν βασιλέα μετὰ τῶν πατέρων αὐτοῦ.

\vs{3}Υἱὸς ἑκκαίδεκα ἐτῶν ἐβασίλευσεν Ὀζίας, καὶ πεντήκοντα καὶ δύο ἔτη ἐβασίλευσεν ἐν Ἱερουσαλὴμ, καὶ ὄνομα τῇ μητρὶ αὐτοῦ Ἰεχελία ἀπὸ Ἱερουσαλήμ.
\vs{4}Καὶ ἐποίησε τὸ εὐθὲς ἐνώπιον Κυρίου κατὰ πάντα ὅσα ἐποίησεν Ἀμασίας ὁ πατὴρ αὐτοῦ.
\vs{5}Καὶ ἦν ἐκζητῶν τὸν Κύριον ἐν ταῖς ἡμέραις Ζαχαρίου τοῦ συνιόντος ἐν φόβῳ Κυρίου, καὶ ἐν ταῖς ἡμέραις αὐτοῦ ἐζήτησε τὸν Κύριον, καὶ εὐώδωσεν αὐτῷ Κύριος.

\vs{6}Καὶ ἐξῆλθε καὶ ἐπολέμησε πρὸς τοὺς ἀλλοφύλους, καὶ κατέσπασε τὰ τείχη Γὲθ, καὶ τὰ τείχη Ἰαβνὴρ, καὶ τὰ τείχη Ἀζώτου, καὶ ᾠκοδόμησε πόλεις Ἀζώτου, καὶ ἐν τοῖς ἀλλοφύλοις.
\vs{7}Καὶ κατίσχυσεν αὐτὸν Κύριος ἐπὶ τοὺς ἀλλοφύλους, καὶ ἐπὶ τοὺς Ἄραβας τοὺς κατοικοῦντας ἐπὶ τῆς πέτρας, καὶ ἐπὶ τοὺς Μιναίους.
\vs{8}Καὶ ἔδωκαν οἱ Μιναῖοι δῶρα τῷ Ὀζίᾳ, καὶ ἦν τὸ ὄνομα αὐτοῦ ἕως εἰσόδου Αἰγύπτου, ὅτι κατίσχυσεν ἕως ἄνω.

\vs{9}Καὶ ᾠκοδόμησεν Ὀζίας πύργους ἐν Ἱερουσαλὴμ, καὶ ἐπὶ τὴν πύλην τῆς γωνίας καὶ ἐπὶ τὴν πύλην τῆς φάραγγος, καὶ ἐπὶ τῶν γενιῶν, καὶ κατίσχυσε.
\vs{10}Καὶ ᾠκοδόμησε πύργους ἐν τῇ ἐρήμῳ, καὶ ἐλατόμησε λάκκους πολλοὺς, ὅτι κτήνη πολλὰ ὑπῆρχεν αὐτῷ ἐν σεφηλᾷ· καὶ ἐν τῇ πεδινῇ, καὶ ἀμπελουργοὶ ἐν τῇ ὀρεινῇ καὶ ἐν τῷ Καρμήλῳ, ὅτι γεωργὸς ἦν.
\vs{11}Καὶ ἐγένετο τῷ Ὀζίᾳ δύναμις ποιοῦσα πόλεμον, καὶ ἐκπορευομένη εἰς παράταξιν εἰς πόλεμον καὶ εἱσπορευομένη εἰς παράταξιν εἰς ἀριθμόν· καὶ ἦν ὁ ἀριθμὸς αὐτῶν διὰ χειρὸς Ἰεϊὴλ τοῦ γραμματέως, καὶ Μαασίου τοῦ κριτοῦ, διὰ χειρὸς Ἀνανίου τοῦ διαδόχου τοῦ βασιλέως.
\vs{12}Πᾶς ὁ ἀριθμὸς τῶν πατριαρχῶν τῶν δυνατῶν εἰς πόλεμον δισχίλιοι ἑξακόσιοι,
\vs{13}καὶ μετʼ αὐτῶν δύναμις πολεμικὴ, τριακόσιαι χιλιάδες καὶ ἑπτακισχίλιοι καὶ πεντακόσιοι· οὗτοι οἱ ποιοῦντες πόλεμον ἐν δυνάμει ἰσχύος βοηθῆσαι τῷ βασιλεῖ ἐπὶ τοὺς ὑπεναντίους.
\vs{14}Καὶ ἡτοίμασεν αὐτοῖς Ὀζίας πάσῃ τῇ δυναμει θυρεοὺς καὶ δόρατα καὶ περικεφαλαίας καὶ θώρακας καὶ τόξα καὶ εἰς λίθους σφενδόνας·
\vs{15}Καὶ ἐποίησεν ἐν Ἱερουσαλὴμ μηχανὰς μεμηχανευμένας λογιστοῦ, τοῦ εἶναι ἐπὶ τῶν πύργων καὶ ἐπὶ τῶν γωνιῶν, βάλλειν βέλεσι καὶ λίθοις μεγάλοις· καὶ ἠκούσθη ἡ κατασκευὴ αὐτῶν ἕως πόῤῥω· ὅτι ἐθαυμαστώθη τοῦ βοηθῆναι ἕως οὗ κατίσχυσε.

\vs{16}Καὶ ὡς κατίσχυσεν, ὑψώθη ἡ καρδία αὐτοῦ τοῦ καταφθεῖραι· καὶ ἠδίκησεν ἐν Κυρίῳ Θεῷ αὐτοῦ, καὶ εἰσῆλθεν εἰς τὸν ναὸν Κυρίου θυμιάσαι ἐπὶ τὸ θυσιαστήριον τῶν θυμιαμάτων.
\vs{17}Καὶ εἰσῆλθεν ὀπίσω αὐτοῦ Ἀζαρίας ὁ ἱερεύς, καὶ μετʼ αὐτοῦ ἱερεῖς τοῦ Κυρίου ὀγδοήκοντα υἱοὶ δυνατοί.
\vs{18}Καὶ ἔστησαν ἐπὶ Ὀζίαν τὸν βασιλέα, καὶ εἶπαν αὐτῷ, οὐ σοὶ, Ὀζία, θυμιάσαι τῷ Κυρίῳ, ἀλλʼ ἢ τοῖς ἱερεῦσιν υἱοῖς Ἀαρὼν τοῖς ἡγιασμένοις θῦσαι· ἔξελθε ἐκ τοῦ ἁγιάσματος, ὅτι ἀπέστης ἀπὸ Κυρίου· καὶ οὐκ ἔσται σοι τοῦτο εἰς δόξαν παρὰ Κυρίου Θεοῦ.

\vs{19}Καὶ ἐθυμώθη Ὀζίας, καὶ ἐν τῇ χειρὶ αὐτοὺ τὸ θυμιατήριον τοῦ θυμιάσαι ἐν τῷ ναῷ· καὶ ἐν τῷ θυμωθῆναι αὐτὸν πρὸς τοὺς ἱερεῖς, καὶ ἡ λέπρα ἀνέτειλεν ἐν τῷ μετώπῳ αὐτοῦ ἐναντίον τῶν ἱερέων ἐν οἴκῳ Κυρίου ἐπάνω τοῦ θυσιαστηρίου τῶν θυμιαμάτων.
\vs{20}Καὶ ἐπέστρεψε πρὸς αὐτὸν Ἀζαρίας ὁ ἱερεὺς ὁ πρῶτος, καὶ οἱ ἱερεῖς, καὶ ἰδοὺ αὐτὸς λεπρὸς ἐν τῷ μετώπῳ, καὶ κατέσπευσαν αὐτὸν ἐκεῖθεν, καὶ γὰρ αὐτὸς ἔσπευσεν ἐξελθεῖν, ὅτι ἤλεγξεν αὐτὸν Κύριος.
\vs{21}Καὶ Ὀζίας ὁ βασιλεὺς ἦν λεπρὸς ἕως ἡμέρας τῆς τελευτῆς αὐτοῦ, καὶ ἐν οἴκῳ ἀπφουσὼθ ἐκάθητο λεπρὸς, ὅτι ἀπεσχίσθη ἀπὸ οἴκου Κυρίου· καὶ Ἰωάθαν ὁ υἱὸς αὐτοῦ ἐπὶ τῆς βασιλείας αὐτοῦ κρίνων τὸν λαὸν τῆς γῆς.

\vs{22}Καὶ οἱ λοιποὶ λόγοι Ὀζίου οἱ πρῶτοι καὶ οἱ ἔσχατοι, γεγραμμένοι ὑπὸ Ἰεσσίου τοῦ προφήτου.
\vs{23}Καὶ ἐκοιμήθη Ὀζίας μετὰ τῶν πατέρων αὐτοῦ, καὶ ἔθαψαν αὐτὸν μετὰ τῶν πατέρων αὐτοῦ ἐν τῷ πεδίῳ τῆς ταφῆς τῶν βασιλέων, ὅτι εἶπαν ὅτι λεπρός ἐστι· καὶ ἐβασίλευσεν Ἰωάθαμ υἱὸς αὐτοῦ ἀντʼ αὐτοῦ.

\ch{27}
Υἱὸς εἴκοσι καὶ πέντε ἐτῶν Ἰωάθαμ ἐν τῷ βασιλεῦσαι αὐτὸν, καὶ ἑκκαίδεκα ἔτη ἐβασίλευσεν ἐν Ἱερουσαλὴμ, καὶ ὄνομα τῆς μητρὸς αὐτοῦ Ἱερουσὰ θυγάτηρ Σαδώκ.
\vs{2}Καὶ ἐποίησε τὸ εὐθὲς ἐνώπιον Κυρίου, κατὰ πάντα ἃ ἐποίησεν Ὀζίας ὁ πατὴρ αὐτοῦ, ἀλλʼ οὐκ εἰσῆλθεν εἰς τὸν ναὸν Κυρίου. Καὶ ἔτι ὁ λαὸς κατεφθείρετο.
\vs{3}Αὐτὸς ᾠκοδόμησε τὴν πύλην οἴκου Κυρίου τὴν ὑψηλὴν, καὶ ἐν τείχει Ὀπέλ ᾠκοδόμησε πολλὰ,
\vs{4}ἐν ὄρει Ἰούδα, καὶ ἐν τοῖς δρυμοῖς καὶ οἰκήσεις καὶ πύργους.
\vs{5}Αὐτὸς ἐμαχέσατο πρὸς βασιλέα υἱῶν Ἀμμὼν, καὶ κατίσχυσεν ἐπʼ αὐτόν· καὶ ἐδίδουν αὐτῷ οἱ υἱοὶ Ἀμμὼν καὶ κατʼ ἐνιαυτὸν ἑκατὸν τάλαντα ἀργυρίου, καὶ δέκα χιλιάδας κόρων πυροῦ, καὶ κριθῶν δέκα χιλιάδας· ταῦτα ἔφερεν αὐτῷ βασιλεὺς υἱῶν Ἀμμὼν κατʼ ἐνιαυτὸν ἐν τῷ πρώτῳ ἔτει καὶ ἐν τῷ δευτέρῳ καὶ τῷ τρίτῳ.
\vs{6}Κατίσχυσεν Ιωάθαμ, ὅτι ἡτοίμασεν τὰς ὁδοὺς αὐτοῦ ἐναντίον Κυρίου Θεοῦ αὐτοῦ.

\vs{7}Καὶ οἱ λοιποὶ λόγοι Ἰωάθαμ καὶ ὁ πόλεμος καὶ αἱ πράξεις αὐτοῦ, ἰδοὺ γεγραμμέναι ἐπὶ βιβλίῳ βασιλέων Ἰούδα καὶ Ἰσραήλ.
\vs{9}Καὶ ἐκοιμήθη Ἰωάθαμ μετὰ τῶν πατέρων αὐτοῦ, καὶ ἐτάφη ἐν πόλει Δαυὶδ, καὶ ἐβασίλευσεν Ἄχαζ υἱὸς αὐτοῦ ἀντʼ αὐτοῦ.

\ch{28}
Υἱὸς εἴκοσι καὶ πέντε ἐτῶν ἦν Ἄχαζ ἐν τῷ βασιλεύειν αὐτὸν, καὶ ἑκκαίδεκα ἔτη ἐβασίλευσεν ἐν Ἱερουσαλήμ· καὶ οὐκ ἐποίησε τὸ εὐθὲς ἐνώπιον Κυρίου, ὡς Δαυὶδ ὁ πατὴρ αὐτοῦ.
\vs{2}Καὶ ἐπορεύθη κατὰ τὰς ὁδοὺς βασιλέων Ἰσραήλ· καὶ γὰρ γλυπτὰ ἐποίησε, καὶ τοῖς εἰδώλοις αὐτῶν
\vs{3}ἐν γὲ Βενεννόμ· καὶ διῆγε τὰ τέκνα αὐτοῦ διὰ πυρὸς κατὰ τὰ βδελύγματα τῶν ἐθνῶν, ὧν ἐξωλόθρευσε Κύριος ἀπὸ προσώπου υἱῶν Ἰσραήλ.
\vs{4}Καὶ ἐθυμία ἐπὶ τῶν ὑψηλῶν, καὶ ἐπὶ τῶν δωμάτων, καὶ ὑποκάτω παντὸς ξύλου ἀλσώδους.

\vs{5}Καὶ παρέδωκεν αὐτὸν Κύριος ὁ Θεὸς αὐτοῦ διὰ χειρὸς βασιλέως Συρίας, καὶ ἐπάταξεν ἐν αὐτῷ, καὶ ᾐχμαλώτευσεν ἐξ αὐτῶν αἰχμαλωσίαν πολλὴν, καὶ ἤγαγεν εἰς Δαμασκόν· καὶ εἰς χεῖρας βασιλέως Ἰσραὴλ παρέδωκεν αὐτὸν, καὶ ἐπάταξεν ἐν αὐτῷ πληγὴν μεγάλην.
\vs{6}Καὶ ἀπέκτεινε Φακεὲ ὁ τοῦ Ῥομελία βασιλεὺς Ἰσραὴλ ἐν Ἰούδᾳ ἐν μιᾷ ἡμέρᾳ ἑκατὸν εἴκοσι χιλιάδας ἀνδρῶν δυνατῶν ἰσχύϊ, ἐν τῷ καταλιπεῖν αὐτοὺς Κύριον τὸν Θεὸν τῶν πατέρων αὐτῶν.
\vs{7}Καὶ ἀπέκτεινε Ζεχρὶ ὁ δυνατὸς τοῦ Ἐφραὶμ τὸν Μαασίαν τὸν υἱὸν τοῦ βασιλέως, καὶ τὸν Ἐζρικὰν ἡγούμενον τοῦ οἴκου αὐτοῦ, καὶ τὸν Ἐλκανὰ τὸν διάδοχον τοῦ βασιλέως.
\vs{8}Καὶ ᾐχμαλώτισαν οἱ υἱοὶ Ἰσραὴλ ἀπὸ τῶν ἀδελφῶν αὐτῶν τριακοσίας χιλιάδας, γυναῖκας καὶ υἱοὺς καὶ θυγατέρας· καὶ σκῦλα πολλὰ ἐσκύλευσαν ἐξ αὐτῶν, καὶ ἤνεγκαν τὰ σκῦλα εἰς Σαμάρειαν.

\vs{9}Καὶ ἐκεῖ ἦν ὁ προφήτης τοῦ Κυρίου, Ὠδὴδ ὄνομα αὐτῷ· καὶ ἐξῆλθεν εἰς ἀπάντησιν τῆς δυνάμεως τῶν ἐρχομένων εἰς Σαμάρειαν, καὶ εἶπεν αὐτοῖς, ἰδοὺ ὀργὴ Κυρίου Θεοῦ τῶν πατέρων ὑμῶν ἐπὶ Ἰούδαν, καὶ παρέδωκεν αὐτοὺς εἰς τὰς χεῖρας ὑμῶν, καὶ ἀπεκτείνατε ἐν αὐτοῖς ἐν ὀργῇ, καὶ ἕως τῶν οὐρανῶν ἔφθακε.
\vs{10}Καὶ νῦν υἱοὺς Ἰούδα καὶ Ἱερουσαλὴμ ὑμεῖς λέγετε κατακτήσασθαι εἰς δούλους καὶ δούλας· οὐκ ἰδού εἰμι μεθʼ ὑμῶν μαρτυρῆσαι Κυρίῳ Θεῷ ὑμῶν;
\vs{11}Καὶ νῦν ἀκούσατέ μου, καὶ ἀποστρέψατε τὴν αἰχμαλωσίαν ἣν ᾐχμαλωτεύσατε τῶν ἀδελφῶν ὑμῶν, ὅτι ὀργὴ θυμοῦ Κυρίου ἐφʼ ὑμῖν.

\vs{12}Καὶ ἀνέστησαν ἄρχοντες ἀπὸ τῶν υἱῶν Ἐφραὶμ, Οὐδείας ὁ τοῦ Ἰωανοῦ, καὶ Βαραχίας ὁ τοῦ Μοσολαμὼθ, καὶ Ἐζεκίας ὁ τοῦ Σελλὴμ, καὶ Ἀμασίας ὁ τοῦ Ἐλδαῒ ἐπὶ τοὺς ἐρχομένους ἀπὸ τοῦ πολέμου,
\vs{13}καὶ εἶπαν αὐτοῖς, οὐ μὴ εἰσαγάγητε τὴν αἰχμαλωσίαν ὧδε πρὸς ἡμᾶς, ὅτι εἰς τὸ ἁμαρτάνειν τῷ Κυρίῳ ἐφʼ ἡμᾶς, ὑμεῖς λέγετε προσθεῖναι ἐπὶ ταῖς ἁμαρτίαις ἡμῶν, καὶ ἐπὶ τὴν ἄγνοιαν ἡμῶν, ὅτι πολλὴ ἡ ἁμαρτία ἡμῶν, καὶ ὀργὴ θυμοῦ Κυρίου ἐπὶ τὸν Ἰσραήλ.
\vs{14}Καὶ ἀφῆκαν οἱ πολεμισταὶ τὴν αἰχμαλωσίαν καὶ τὰ σκῦλα ἐναντίον τῶν ἀρχόντων καὶ πάσης τῆς ἐκκλησίας.
\vs{15}Καὶ ἀνέστησαν ἄνδρες οἳ ἐπεκλήθησαν ἐν ὀνόματι, καὶ ἀντελάβοντο τῆς αἰχμαλωσίας, καὶ πάντας τοὺς γυμνοὺς περιέβαλον ἀπὸ τῶν σκύλων, καὶ ἐνέδυσαν αὐτοὺς καὶ ὑπέδησαν αὐτοὺς, καὶ ἔδωκαν φαγεῖν καὶ ἀλείψασθαι, καὶ ἀντελάβοντο καὶ ἐν ὑποζυγίοις παντὸς ἀσθενοῦντος, καὶ κατέστησαν αὐτοὺς εἰς Ἱεριχὼ πόλιν φοινίκων πρὸς τοὺς ἀδελφοὺς αὐτῶν, καὶ ἐπέστρεψαν εἰς Σαμάρειαν.

\vs{16}Ἐν τῷ καιρῷ ἐκείνῳ ἀπέστειλεν ὁ βασιλεὺς Ἄχας πρὸς βασιλέα Ἀσσοὺρ βοηθῆσαι αὐτῷ
\vs{17}καὶ ἐν τούτῳ, ὅτι οἱ Ἰδουμαῖοι ἐπέθεντο, καὶ ἐπάταξαν ἐν Ἰούδα, καὶ ᾐχμαλώτισαν αἰχμαλωσίαν.
\vs{18}Καὶ οἱ ἀλλόφυλοι ἐπέθεντο ἐπὶ τὰς πόλεις τῆς πεδινῆς, καὶ ἀπὸ Λιβὸς τοῦ Ἰούδα, καὶ ἔλαβον τὴν Βαιθσαμὺς, καὶ τὰ ἐν οἴκῳ Κυρίου, καὶ τὰ ἐν οἴκῳ τοῦ βασιλέως καὶ τῶν ἀρχόντων, καὶ ἔδωκαν τῷ βασιλεῖ τὴν Ἀϊλὼν, καὶ τὴν Γαληρὼ, καὶ τὴν Σωχὼ καὶ τὰς κώμας αὐτῆς, καὶ τὴν Θαμνὰ καὶ τὰς κώμας αὐτῆς, καὶ τὴν Γαμζὼ καὶ τὰς κώμας αὐτῆς· καὶ κατῴκησαν ἐκεῖ.
\vs{19}Ὅτι ἐταπείνωσε Κύριος τὸν Ἰούδαν διὰ Ἄχαζ βασιλέα Ἰούδα, ὅτι ἀπέστη ἀποστάσει ἀπὸ Κυρίου.
\vs{20}Καὶ ἦλθεν ἐπʼ αὐτὸν Θαλγαφελλασὰρ βασιλεὺς Ἀσσοὺρ, καὶ ἔθλιψεν αὐτόν.
\vs{21}Καὶ ἔλαβεν Ἄχας τὰ ἐν οἴκῳ Κυρίου, καὶ τὰ ἐν οἴκῳ τοῦ βασιλέως καὶ τῶν ἀρχόντων, καὶ ἔδωκε τῷ βασιλεῖ Ἀσσούρ· καὶ οὐκ εἰς βοήθειαν αὐτῷ ἦν,
\vs{22}ἀλλʼ ἢ τῷ θλιβῆναι αὐτόν· καὶ προσέθηκε τοῦ ἀποστῆναι ἀπὸ Κυρίου, καὶ εἶπεν ὁ βασιλεὺς Ἄχαζ,
\vs{23}ἐκζητήσω τοὺς θεοὺς Δαμασκοῦ τοὺς τύπτοντάς με· καὶ εἶπεν, ὅτι θεοὶ βασιλέως Συρίας αὐτοὶ κατισχύσουσιν αὐτοὺς, αὐτοῖς τοίνυν θύσω, καὶ ἀντιλήψονταί μου· καὶ αὐτοὶ ἐγένοντο αὐτῷ εἰς σκῶλον καὶ παντὶ Ἰσραήλ.

\vs{24}Καὶ ἀπέστησεν Ἄχαζ τὰ σκεύη οἴκου Κυρίου, καὶ κατέκοψεν αὐτὰ, καὶ ἔκλεισε τὰς θύρας οἴκου Κυρίου, καὶ ἐποίησεν ἑαυτῷ θυσιαστήρια ἐν πάσῃ γωνίᾳ ἐν Ἱερουσαλὴμ,
\vs{25}καὶ ἐν πάσῃ πόλει καὶ πόλει ἐν Ἰούδᾳ ἐποίησεν ὑψηλὰ θυμιᾷν θεοῖς ἀλλοτρίοις, καὶ παρώργισαν Κύριον τὸν Θεὸν τῶν πατέρων αὐτῶν.
\vs{26}Καὶ οἱ λοιποὶ λόγοι αὐτοῦ καὶ αἱ πράξεις αὐτοῦ αἱ πρῶται καὶ ἔσχαται, ἰδοὺ γεγραμμέναι ἐπὶ βιβλίῳ βασιλέων Ἰούδα καὶ Ἰσραήλ.
\vs{27}Καὶ ἐκοιμήθη Ἄχαζ μετὰ τῶν πατέρων αὐτοῦ, καὶ ἐτάφη ἐν πόλει Δαυὶδ, ὅτι οὐκ εἰσήνεγκαν αὐτὸν εἰς τοὺς τάφους τῶν βασιλέων Ἰσραὴλ, καὶ ἐβασίλευσεν Ἐζεκίας υἱὸς αὐτοῦ ἀντʼ αὐτοῦ.

\ch{29}
Καὶ Ἐζεκίας ἐβασίλευσεν ὢν εἴκοσι καὶ πέντε ἐτῶν, καὶ εἴκοσι ἐννέα ἔτη ἐβασίλευσεν ἐν Ἱερουσαλὴμ, καὶ ὄνομα τῇ μητρὶ αὐτοῦ Ἀβιὰ, θυγάτηρ Ζαχαρίου.
\vs{2}Καὶ ἐποίησε τὸ εὐθὲς ἐνώπιον Κυρίου κατὰ πάντα ὅσα ἐποίησε Δαυὶδ ὁ πατὴρ αὐτοῦ.

\vs{3}Καὶ ἐγένετο ὡς ἔστη ἐπὶ τῆς βασιλείας αὐτοῦ, ἐν τῷ μηνὶ τῷ πρώτῳ ἀνέῳξε τὰς θύρας οἴκου Κυρίου καὶ ἐπεσκεύασεν αὐτάς.
\vs{4}Καὶ εἰσήγαγε τοὺς ἱερεῖς καὶ τοὺς Λευίτας, καὶ κατέστησεν αὐτοὺς εἰς τὸ κλίτος τὸ πρὸς ἀνατολὰς,
\vs{5}καὶ εἶπεν αὐτοῖς, ἀκούσατε οἱ Λευῖται· νῦν ἁγνίσθητε, καὶ ἁγνίσατε τὸν οἶκον Κυρίου Θεοῦ τῶν πατέρων ὑμῶν, καὶ ἐκβάλετε τὴν ἀκαθαρσίαν ἐκ τῶν ἁγίων·
\vs{6}Ὅτι ἀπέστησαν οἱ πατέρες ἡμῶν, καὶ ἐποίησαν τὸ πονηρὸν ἐναντίον Κυρίου Θεοῦ ἡμῶν, καὶ ἐγκατέλιπαν αὐτὸν, καὶ ἀπέστρεψαν τὸ πρόσωπον αὐτῶν ἀπὸ τῆς σκηνῆς Κυρίου, καὶ ἔδωκαν αὐχένα,
\vs{7}καὶ ἀπέκλεισαν τὰς θύρας τοῦ ναοῦ, καὶ ἔσβεσαν τοὺς λύχνους, καὶ θυμίαμα οὐκ ἐθυμίασαν, καὶ ὁλοκαυτώματα οὐ προσήνεγκαν ἐν τῷ ἁγίῳ Θεῷ Ἰσραήλ.
\vs{8}Καὶ ὠργίσθη ὀργῇ Κύριος ἐπὶ τὸν Ἰούδαν καὶ τὴν Ἱερουσαλὴμ, καὶ ἔδωκεν αὐτοὺς εἰς ἔκστασιν καὶ εἰς ἀφανισμὸν καὶ εἰς συρισμὸν ὡς ὑμεῖς ὁρᾶτε τοῖς ὀφθαλμοῖς ὑμῶν.
\vs{9}Καὶ ἰδοὺ πεπλήγασιν οἱ πατέρες ὑμῶν ἐν μαχαίρᾳ, καὶ οἱ υἱοὶ ὑμῶν καὶ αἱ θυγατέρες ὑμῶν καὶ αἱ γυναῖκες ὑμῶν ἐν αἰχμαλωσίᾳ ἐν γῇ οὐκ αὐτῶν, ὃ καὶ νῦν ἐστιν.
\vs{10}Ἐπὶ τούτοις νῦν ἐστιν ἐπὶ καρδίας διαθέσθαι διαθήκην μου, διαθήκην Κυρίου Θεοῦ Ἰσραὴλ, καὶ ἀποστρέψει τὴν ὀργὴν τοῦ θυμοῦ αὐτοῦ ἀφʼ ἡμῶν.
\vs{11}Καὶ νῦν μὴ διαλίπητε, ὅτι ἐν ὑμῖν ᾑρέτικε Κύριος στῆναι ἐναντίον αὐτοῦ λειτουργεῖν, καὶ εἶναι αὐτῷ λειτουργοῦντας καὶ θυμιῶντας.

\vs{12}Καὶ ἀνέστησαν οἱ Λευῖται, Μαὰθ ὁ τοῦ Ἀμασὶ, καὶ Ἰωὴλ ὁ τοῦ Ἀζαρίου ἐκ τῶν υἱῶν Καάθ· καὶ ἐκ τῶν υἱῶν Μεραρὶ, Κὶς ὁ τοῦ Ἀβδὶ, καὶ Ἀζαρίας ὁ τοῦ Ἰλαελήλ· καὶ ἀπὸ τῶν υἱῶν Γεδσωνὶ, Ἰωδαὰδ ὁ τοῦ Ζεμμὰθ, καὶ Ἰωαδάμ· οὗτοι υἱοὶ Ἰωαχά.
\vs{13}Καὶ τῶν υἱῶν Ἐλισαφὰν, Ζαμβρὶ, καὶ Ἰεϊήλ· καὶ τῶν υἱῶν Ἀσὰφ, Ζαχαρίας καὶ Ματθανίας·
\vs{14}Καὶ τῶν υἱῶν Αἰμὰν, Ἰεϊὴλ καὶ Σεμεΐ· καὶ τῶν υἱῶν Ἰδιθοὺν, Σαμαίας, καὶ Ὀζιήλ.
\vs{15}Καὶ συνήγαγον τοὺς ἀδελφοὺς αὐτῶν, καὶ ἡγνίσθησαν κατὰ τὴν ἐντολὴν τοῦ βασιλέως διὰ προστάγματος Κυρίου, καθαρίσαι τὸν οἶκον Κυρίου.
\vs{16}Καὶ εἰσῆλθον οἱ ἱερεῖς ἔσω εἰς τὸν οἶκον Κυρίου ἁγνίσαι, καὶ ἐξέβαλον πᾶσαν τὴν ἀκαθαρσίαν τὴν εὑρεθεῖσαν ἐν τῷ οἴκῳ Κυρίου, καὶ εἰς τὴν αὐλὴν οἴκου Κυρίου· καὶ ἐδέξαντο οἱ Λευῖται ἐκβαλεῖν εἰς τὸν χειμάῤῥουν Κέδρων ἔξω.

\vs{17}Καὶ ἤρξατο τῇ ἡμέρᾳ τῇ πρώτῃ νουμηνίᾳ τοῦ πρώτου μηνὸς ἁγνίσαι, καὶ τῇ ἡμέρᾳ τῇ ὀγδόῃ τοῦ μηνὸς εἰσῆλθαν εἰς τὸν ναὸν Κυρίου, καὶ ἥγνισαν τὸν οἶκον Κυρίου ἐν ἡμέραις ὀκτὼ, καὶ τῇ ἡμέρᾳ τῇ τρισκαιδεκάτῃ τοῦ μηνὸς τοῦ πρώτου συνετέλεσαν.

\vs{18}Καὶ εἰσῆλθαν ἔσω πρὸς Ἐξεκίαν τὸν βασιλέα, καὶ εἶπαν, ἡγνίσαμεν πάντα τὰ ἐν οἴκῳ Κυρίου, τὸ θυσιαστήριον τῆς ὁλοκαυτώσεως καὶ τὰ σκεύη αὐτοῦ, καὶ τὴν τράπεζαν τῆς προθέσεως καὶ τὰ σκεύη αὐτῆς,
\vs{19}καὶ πάντα τὰ σκεύη ἃ ἐμίανεν ὁ βασιλεὺς Ἄχαζ ἐν τῇ βασιλείᾳ αὐτοῦ ἐν τῇ ἀποστασίᾳ αὐτοῦ, ἡτοιμάκαμεν καὶ ἡγνίσαμεν· ἰδού ἐστιν ἐναντίον τοῦ θυσιαστήριου Κυρίου.

\vs{20}Καὶ ὤρθρισεν Ἐζεκίας ὁ βασιλεὺς, καὶ συνήγαγε τοὺς ἄρχοντας τῆς πόλεως, καὶ ἀνέβη εἰς οἶκον Κυρίου.
\vs{21}Καὶ ἀνήνεγκε μόσχους ἑπτὰ, κριοὺς ἑπτὰ, ἀμνοὺς ἑπτὰ, χιμάρους αἰγῶν ἑπτὰ περὶ ἁμαρτίας, περὶ τῆς βασιλείας, καὶ περὶ τῶν ἁγίων, καὶ περὶ Ἰσραήλ· καὶ εἶπε τοῖς υἱοῖς Ἀαρὼν τοῖς ἱερεῦσιν ἀναβαίνειν ἐπὶ τὸ θυσιαστήριον Κυρίου.
\vs{22}Καὶ ἔθυσαν τοὺς μόσχους, καὶ ἐδέξαντο οἱ ἱερεῖς τὸ αἷμα, καὶ προσέχεαν ἐπὶ τὸ θυσιαστήριον· καὶ ἔθυσαν τοὺς κριοὺς, καὶ προσέχεαν τὸ αἷμα ἐπὶ τὸ θυσιαστήριον· καὶ ἔθυσαν τοὺς ἀμνοὺς, καὶ περιέχεον τὸ αἷμα τῷ θυσιαστηρίῳ.
\vs{23}Καὶ προσήγαγον τοὺς χιμάρους τοὺς περὶ ἁμαρτίας ἐναντίον τοῦ βασιλέως καὶ τῆς ἐκκλησίας, καὶ ἐπέθηκαν τὰς χεῖρας αὐτῶν ἐπʼ αὐτούς.
\vs{24}Καὶ ἔθυσαν αὐτοὺς οἱ ἱερεῖς, καὶ ἐξιλάσαντο τὸ αἷμα αὐτῶν πρὸς τὸ θυσιαστήριον, καὶ ἐξιλάσαντο περὶ παντὸς Ἰσραὴλ, ὅτι εἶπεν ὁ βασιλεὺς, περὶ παντὸς Ἰσραὴλ ἡ ὁλοκαύτωσις, καὶ τὰ περὶ ἁμαρτίας.

\vs{25}Καὶ ἔστησε τοὺς Λευίτας ἐν οἴκῳ Κυρίου ἐν κυμβάλοις, καὶ ἐν νάβλαις, καὶ ἐν κινύραις κατὰ τὴν ἐντολὴν Δαυὶδ τοῦ βασιλέως, καὶ Γὰδ τοῦ ὁρῶντος τῷ βασιλεῖ, καὶ Νάθαν τοῦ προφήτου, ὅτι διὰ ἐντολῆς Κυρίου τὸ πρόσταγμα ἐν χειρὶ τῶν προφητῶν.
\vs{26}Καὶ ἔστησαν οἱ Λευῖται ἐν ὀργάνοις Δαυὶδ, καὶ οἱ ἱερεῖς ταῖς σάλπιγξι.
\vs{27}Καὶ εἶπεν Ἐζεκίας ἀνένεγκαι τὴν ὁλοκαύτωσιν ἐπὶ τὸ θυσιαστήριον· καὶ ἐν τῷ ἄρξασθαι ἀναφέρειν τὴν ὁλοκαύτωσιν, ἤρξαντο ᾄδειν Κυρίῳ, καὶ σάλπιγγες πρὸς τὰ ὄργανα Δαυὶδ βασιλέως Ἰσραήλ.
\vs{28}Καὶ πᾶσα ἡ ἐκκλησία προσεκύνει, καὶ οἱ ψαλτῳδοὶ ᾄδοντες, καὶ σάλπιγγες σαλπίζουσαι ἕως οὗ συνετελέσθη ἡ ὁλοκαύτωσις.
\vs{29}Καὶ ὡς συνετέλεσαν ἀναφέροντες, ἔκαμψεν ὁ βασιλεὺς καὶ πάντες οἱ εὑρεθέντες, καὶ προσεκύνησαν.

\vs{30}Καὶ εἶπεν Ἐζεκίας ὁ βασιλεὺς καὶ οἱ ἄρχοντες τοῖς Λευίταις, ὑμνεῖν τὸν Κύριον ἐν λόγοις Δαυὶδ καὶ Ἀσὰφ τοῦ προφήτου· καὶ ὕμνουν ἐν εὐφροσύνῃ, καὶ ἔπεσον καὶ προσεκύνησαν.

\vs{31}Καὶ ἀπεκρίθη Ἐζεκίας, καὶ εἶπε, νῦν ἐπληρώσατε τὰς χεῖρας ὑμῶν Κυρίῳ, προσαγάγετε καὶ φέρετε θυσίας αἰνέσεως εἰς οἶκον Κυρίου· καὶ ἀνήνεγκεν ἡ ἐκκλησία θυσίας καὶ αἰνέσεις εἰς οἶκον Κυρίου, καὶ πᾶς πρόθυμος τῇ καρδίᾳ ὁλοκαυτώσεις.
\vs{32}Καὶ ἐγένετο ὁ ἀριθμὸς τῆς ὁλοκαυτώσεως ἧς ἀνήνεγκεν ἡ ἐκκλησία, μόσχοι ἑβδομήκοντα, κριοὶ ἑκατὸν, ἀμνοὶ διακόσιοι· εἰς ὁλοκαύτωσιν Κυρίῳ πάντα ταῦτα.
\vs{33}Καὶ οἱ ἠγιασμένοι μόσχοι ἑξακόσιοι, πρόβατα τρισχίλια.
\vs{34}Ἀλλʼ ἢ οἱ ἱερεῖς ἦσαν ὀλίγοι, καὶ οὐκ ἠδύναντο ἐκδεῖραι τὴν ὁλοκαύτωσιν, καὶ ἀντελάβοντο αὐτῶν οἱ ἀδελφοὶ αὐτῶν οἱ Λευῖται ἕως οὗ συνετελέσθη τὸ ἔργον, καὶ ἕως οὗ ἡγνίσθησαν οἱ ἱερεῖς· ὅτι οἱ Λευῖται προθύμως ἥγνισαν παρὰ τοὺς ἱερεῖς.
\vs{35}Καὶ ἡ ὁλοκαύτωσις πολλὴ ἐν τοῖς στέασι τῆς τελειώσεως τοῦ σωτηρίου καὶ τῶν σπονδῶν τῆς ὁλοκαυτώσεως· καὶ κατωρθώθη τὸ ἔργον ἐν οἴκῳ Κυρίου.

\vs{36}Καὶ ηὐφράνθη Ἐζεκίας καὶ πᾶς ὁ λαὸς, διὰ τὸ ἡτοιμακέναι τὸν Θεὸν τῷ λαῷ, ὅτι ἐξάπινα ἐγένετο ὁ λόγος.

\ch{30}
Καὶ ἀπέστειλεν Ἐζεκίας ἐπὶ πάντα Ἰσραὴλ καὶ Ἰούδα, καὶ ἐπιστολὰς ἔγραψεν ἐπὶ τὸν Ἐφραὶμ καὶ Μανασσῆ, ἐλθεῖν εἰς οἶκον Κυρίου, εἰς Ἱερουσαλὴμ, ποιῆσαι τὸ φασὲκ τῷ Κυρίῳ Θεῷ Ἰσραήλ.
\vs{2}Καὶ ἐβουλεύσατο ὁ βασιλεὺς καὶ οἱ ἄρχοντες καὶ πᾶσα ἡ ἐκκλησία ἐν Ἱερουσαλὴμ ποιῆσαι τὸ φασὲκ τῷ μηνὶ τῷ δευτέρῳ.
\vs{3}Οὐ γὰρ ἠδυνάσθησαν ποιῆσαι αὐτὸ ἐν τῷ καιρῷ ἐκείνῳ, ὅτι οἱ ἱερεῖς οὐχ ἠγνίσθησαν ἱκανοὶ, καὶ ὁ λαὸς οὐ συνήχθη εἰς Ἱερουσαλήμ.
\vs{4}Καὶ ἤρεσεν ὁ λόγος ἐναντίον τοῦ βασιλέως καὶ ἐναντίον τῆς ἐκκλησίας.
\vs{5}Καὶ ἔστησαν λόγον διελθεῖν κήρυγμα ἐν παντὶ Ἰσραὴλ ἀπὸ Βηρσαβεὲ ἕως Δὰν, ἐλθόντας ποιῆσαι τὸ φασὲκ Κυρίῳ Θεῷ Ἰσραὴλ εἰς Ἱερουσαλὴμ, ὅτι πλῆθος οὐκ ἐποίησεν κατὰ τὴν γραφήν.

\vs{6}Καὶ ἐπορεύθησαν οἱ τρέχοντες σὺν ταῖς ἐπιστολαῖς παρὰ τοῦ βασιλέως καὶ τῶν ἀρχόντων εἰς πάντα Ἰσραὴλ καὶ Ἰούδαν κατὰ τὸ πρόσταγμα τοῦ βασιλέως, λέγοντες, οἱ υἱοὶ Ἰσραὴλ ἐπιστρέψατε πρὸς Κύριον Θεὸν Ἁβραὰμ καὶ Ἰσαὰκ καὶ Ἰσραὴλ, καὶ ἐπιστρέψατε τοὺς ἀνασεσωσμένους τοὺς καταλειφθέντας ἀπὸ χειρὸς βασιλέως Ἀσσούρ.
\vs{7}Καὶ μὴ γίνεσθε καθὼς οἱ πατέρες ὑμῶν καὶ οἱ ἀδελφοὶ ὑμῶν, οἳ ἀπέστησαν ἀπὸ Κυρίου Θεοῦ πατέρων αὐτῶν, καὶ παρέδωκεν αὐτοὺς εἰς ἐρήμωσιν καθὼς ὑμεῖς ὁρᾶτε.
\vs{8}Καὶ νῦν μὴ σκληρύνητε τὰς καρδίας ὑμῶν ὡς οἱ πατέρες ὑμῶν, δότε δόξαν Κυρίῳ τῷ Θεῷ, καὶ εἰσέλθετε εἰς τὸ ἁγίασμα αὐτοῦ ὃ ἡγίασεν εἰς τὸν αἰῶνα, καὶ δουλεύσατε τῷ Κυρίῳ Θεῷ ὑμῶν, καὶ ἀποστρέψει ἀφʼ ὑμῶν θυμὸν ὀργῆς.
\vs{9}Ὅτι ἐν τῷ ἐπιστρέφειν ὑμᾶς πρὸς Κύριον, οἱ ἀδελφοὶ ὑμῶν καὶ τὰ τέκνα ὑμῶν ἔσονται ἐν οἰκτιρμοῖς ἔναντι πάντων τῶν αἰχμαλωτισάντων αὐτοὺς, καὶ ἀποστρέψει εἰς τὴν γῆν ταύτην· ὅτι ἐλεήμων καὶ οἰκτίρμων Κύριος ὁ Θεὸς ἡμῶν, καὶ οὐκ ἀποστρέψει τὸ πρόσωπον αὐτοῦ ἀφʼ ὑμῶν, ἐὰν ἐπιστρέψωμεν πρὸς αὐτόν.

\vs{10}Καὶ ἦσαν οἱ τρέχοντες διαπορευόμενοι πόλιν ἐκ πόλεως ἐν τῷ ὄρει Ἐφραὶμ, καὶ Μανασσῆ, καὶ ἕως Ζαβουλών· καὶ ἐγένοντο ὡς καταγελῶντες αὐτῶν, καὶ καταμωκώμενοι.
\vs{11}Ἀλλὰ ἄνθρωποι Ἀσὴρ καὶ ἀπὸ Μανασσῆ καὶ ἀπὸ Ζαβουλὼν ἐνετράπησαν, καὶ ἦλθον εἰς Ἱερουσαλὴμ καὶ εἰς Ἰούδα.
\vs{12}Καὶ ἐγένετο χεὶρ Κυρίου δοῦναι αὐτοῖς καρδίαν μίαν ἐλθεῖν, τοῦ ποιῆσαι κατὰ τὰ προστάγματα τοῦ βασιλέως καὶ τῶν ἀρχόντων ἐν λόγῳ Κυρίου.

\vs{13}Καὶ συνήχθησαν εἰς Ἱερουσαλὴμ λαὸς πολὺς τοῦ ποιῆσαι τὴν ἑορτὴν τῶν ἀζύμων ἐν τῷ μηνὶ τῷ δευτέρῳ, ἐκκλησία πολλὴ σφόδρα.
\vs{14}Καὶ ἀνέστησαν, καὶ καθεῖλαν τὰ θυσιαστήρια τὰ ἐν Ἱερουσαλὴμ, καὶ πάντα ἐν οἷς ἐθυμίων τοῖς ψευδέσι, κατέσπασαν καὶ ἔῤῥιψαν εἰς τὸν χειμάῤῥουν Κέδρων.
\vs{15}Καὶ ἔθυσαν τὸ φασὲκ τῇ τεσσαρεσκαιδεκάτῃ τοῦ μηνὸς τοῦ δευτέρου· καὶ οἱ ἱερεῖς καὶ οἱ Λευῖται ἐνετράπησαν καὶ ἥγνισαν, καὶ εἰσήνεγκαν ὁλοκαυτώματα ἐν οἴκῳ Κυρίου.

\vs{16}Καὶ ἔστησαν ἐπὶ τὴν στάσιν αὐτῶν, κατὰ τὸ κρίμα αὐτῶν, κατὰ τὴν ἐντολὴν Μωυσῆ ἀνθρώπου τοῦ Θεοῦ· καὶ οἱ ἱερεῖς ἐδέχοντο τὰ αἵματα ἐκ χειρὸς τῶν Λευιτῶν.
\vs{17}Ὅτι πλῆθος τῆς ἐκκλησίας οὐχ ἡγνίσθη, καὶ οἱ Λευῖται ἦσαν τοῦ θύειν τὸ φασὲκ παντὶ τῷ μὴ δυναμένῳ ἁγνισθῆναι τῷ Κυρίῳ.
\vs{18}Ὅτι πλεῖστον τοῦ λαοῦ ἀπὸ Ἐφραὶμ, καὶ Μανασσῆ, καὶ Ἰσσάχαρ, καὶ Ζαβουλὼν, οὐκ ἥγνισαν, ἀλλʼ ἔφαγον τὸ φασὲκ παρὰ τὴν γραφήν· τοῦτο καὶ προσηύξατο Ἐζεκίας περὶ αὐτῶν, λέγων, Κύριος ἀγαθὸς ἐξιλάσθω ὑπὲρ
\vs{19}πάσης καρδίας κατευθυνούσης ἐκζητῆσαι Κύριον τὸν Θεὸν τῶν πατέρων αὐτῶν, καὶ οὐ κατὰ τὴν ἁγνείαν τῶν ἁγίων.
\vs{20}Καὶ ἐπήκουσε Κύριος τῷ Ἐζεκίᾳ, καὶ ἰάσατο τὸν λαόν.

\vs{21}Καὶ ἐποίησαν οἱ υἱοὶ Ἰσραὴλ οἱ εὑρεθέντες ἐν Ἱερουσαλὴμ, τὴν ἑορτὴν τῶν ἀζύμων ἑπτὰ ἡμέρας ἐν εὐφροσύνῃ μεγάλῃ, καὶ καθυμνοῦντες τῷ Κυρίῳ ἡμέραν καθʼ ἡμέραν, καὶ οἱ ἱερεῖς καὶ οἱ Λευῖται ἐν ὀργάνοις τῷ Κυρίῳ.
\vs{22}Καὶ ἐλάλησεν Ἐζεκίας ἐπὶ πᾶσαν καρδίαν τῶν Λευιτῶν καὶ τῶν συνιόντων σύνεσιν ἀγαθὴν τῷ Κυρίῳ· καὶ συνετέλεσαν τὴν ἑορτὴν τῶν ἀζύμων ἑπτὰ ἡμέρας, θύοντες θυσίαν σωτηρίου, καὶ ἐξομολογούμενοι τῷ Κυρίῳ Θεῷ τῶν πατέρων αὐτῶν.

\vs{23}Καὶ ἐβουλεύσατο ἡ ἐκκλησία ἅμα ποιῆσαι ἑπτὰ ἡμέρας ἄλλας· καὶ ἐποίησαν ἑπτὰ ἡμέρας ἐν εὐφροσύνῃ.
\vs{24}Ὅτι Ἐζεκίας ἀπήρξατο τῷ Ἰούδα τῇ ἐκκλησίᾳ χιλίους μόσχους καὶ ἑπτακισχίλια πρόβατα, καὶ οἱ ἄρχοντες ἀπήρξαντο τῷ λαῷ μόσχους χιλίους καὶ πρόβατα δέκα χιλιάδας, καὶ τὰ ἅγια τῶν ἱερέων εἰς πλῆθος.
\vs{25}Καὶ ηὐφράνθη πᾶσα ἡ ἐκκλήσια οἱ ἱερεῖς καὶ οἱ Λευῖται, καὶ πᾶσα ἡ ἐκκλησία Ἰούδα, καὶ οἱ εὑρεθέντες ἐξ Ἱερουσαλὴμ, καὶ οἱ προσήλυτοι οἱ ἐλθόντες ἀπὸ γῆς Ἰσραὴλ, καὶ οἱ κατοικοῦντες Ἰούδα.
\vs{26}Καὶ ἐγένετο εὐφροσύνη μεγάλη ἐν Ἱερουσαλήμ· ἀπὸ ἡμερῶν Σαλωμὼν υἱοῦ Δαυὶδ βασιλέως Ἰσραὴλ οὐκ ἐγένετο τοιαύτη ἑορτὴ ἐν Ἱερουσαλήμ.
\vs{27}Καὶ ἀνέστησαν οἱ ἱερεῖς οἱ Λευῖται καὶ εὐλόγησαν τὸν λαὸν· καὶ ἐπηκούσθη ἡ φωνὴ αὐτῶν, καὶ ἦλθεν ἡ προσευχὴ αὐτῶν εἰς τὸ κατοικητήριον τὸ ἅγιον αὐτοῦ εἰς τὸν οὐρανόν.

\ch{31}
Καὶ ὡς συνετελέσθη πάντα ταῦτα, ἐξῆλθε πᾶς Ἰσραὴλ οἱ εὑρεθέντες ἐν πόλεσιν Ἰούδα, καὶ συνέτριψαν τὰς στήλας, καὶ ἔκοψαν τὰ ἄλση, καὶ κατέσπασαν τὰ ὑψηλὰ καὶ τοὺς βωμοὺς ἀπὸ πάσης τῆς Ἰουδαίας καὶ Βενιαμὶν, καὶ ἐξ Ἐφραὶμ, καὶ ἀπὸ Μανασσῆ ἕως εἰς τέλος· καὶ ἐπέστρεψαν πᾶς Ἰσραὴλ ἕκαστος εἰς τὴν κληρονομίαν αὐτοῦ, καὶ εἰς τὰς πόλεις αὐτῶν.

\vs{2}Καὶ ἔταξεν Ἐζεκίας τὰς ἐφημερίας τῶν ἱερέων καὶ τῶν Λευιτῶν, καὶ τὰς ἐφημερίας ἑκάστου κατὰ τὴν ἑαυτοῦ λειτουργίαν, τοῖς ἱερεῦσι καὶ τοῖς Λευίταις, εἰς τὴν ὁλοκαύτωσιν, καὶ εἰς τὴν θυσίαν τοῦ σωτηρίου, καὶ αἰνεῖν, καὶ ἐξομολογεῖσθαι, καὶ λειτουργεῖν ἐν ταῖς πύλαις ἐν ταῖς αὐλαῖς οἴκου Κυρίου.
\vs{3}Καὶ μερὶς τοῦ βασιλέως ἐκ τῶν ὑπαρχόντων αὐτοῦ εἰς τὰς ὁλοκαυτώσεις τὴν πρωϊνὴν καὶ τὴν δειλινὴν, καὶ ὁλοκαυτώσεις εἰς τὰ σάββατα, καὶ εἰς τὰς νουμηνίας, καὶ εἰς τὰς ἑορτὰς τὰς γεγραμμένας ἐν τῷ νόμῳ Κυρίου.

\vs{4}Καὶ εἶπαν τῷ λαῷ τοῖς κατοικοῦσιν ἐν Ἱερουσαλὴμ, δοῦναι τὴν μεριδα τῶν ἱερέων καὶ τῶν Λευιτῶν, ὅπως κατισχύσωσιν ἐν τῇ λειτουργίᾳ οἴκου Κυρίου.
\vs{5}Καὶ ὡς προσέταξεν τὸν λόγον, ἐπλεόνασεν Ἰσραὴλ ἀπαρχὴν σίτου, καὶ οἴνου, καὶ ἐλαίου, καὶ μέλιτος, καὶ πᾶν γέννημα ἀγροῦ, καὶ ἐπιδέκατα πάντα εἰς πλῆθος ἤνεγκαν οἱ υἱοὶ Ἰσραὴλ καὶ Ἰούδα.
\vs{6}Καὶ οἱ κατοικοῦντες ἐν ταῖς πόλεσιν Ἰούδα καὶ αὐτοὶ ἤνεγκαν ἐπιδέκατα μόσχων καὶ προβάτων, καὶ ἐπιδέκατα αἰγῶν, καὶ ἡγίασαν τῷ Κυρίῳ Θεῷ αὐτῶν, καὶ εἰσήνεγκαν καὶ ἔθηκαν σωροὺς σωρούς.
\vs{7}Ἐν τῷ μηνὶ τῷ τρίτῳ ἤρξαντο οἱ σωροὶ θεμελιοῦσθαι, καὶ ἐν τῷ μηνὶ τῷ ἑβδόμῳ συνετελέσθησαν.
\vs{8}Καὶ ἦλθεν Ἐζεκίας καὶ οἱ ἄρχοντες, καὶ εἶδον τοὺς σωροὺς, καὶ ηὐλόγησαν τὸν Κύριον καὶ τὸν λαὸν αὐτοῦ Ἰσραήλ·
\vs{9}Καὶ ἐπυνθάνετο Ἐζεκίας τῶν ἱερέων καὶ τῶν Λευιτῶν ὑπὲρ τῶν σωρῶν.
\vs{10}Καὶ εἶπε πρὸς αὐτὸν Ἀζαρίας ὁ ἱερεὺς ὁ ἄρχων εἰς οἶκον Σαδὼκ, καὶ εἶπεν, ἐξ οὗ ἦρκται ἡ ἀπαρχὴ φέρεσθαι εἰς οἶκον Κυρίου, ἐφάγομεν καὶ ἐπίομεν καὶ κατελίπομεν ἕως εἰς πλῆθος, ὅτι Κύριος ηὐλόγησε τὸν λαὸν αὐτοῦ, καὶ κατελίπομεν ἐπὶ τὸ πλῆθος τοῦτο.

\vs{11}Καὶ εἶπεν Ἐζεκίας ἔτι ἑτοιμάσαι παστοφόρια εἰς οἶκον Κυρίου· καὶ ἡτοίμασαν,
\vs{12}καὶ ἤνεγκαν ἐκεῖ τὰς ἀπαρχὰς καὶ τὰ ἐπιδέκατα ἐν πίστει· καὶ ἐπʼ αὐτῶν ἐπιστάτης Χωνενίας ὁ Λευίτης, καὶ Σεμεῒ ὁ ἀδελφὸς αὐτοῦ διαδεχόμενος·
\vs{13}Καὶ Ἰεϊὴλ, καὶ Ὀζίας, καὶ Ναὲθ, καὶ Ἀσαὴλ, καὶ Ἰεριμὼθ, καὶ Ἰωζαβὰδ, καὶ Ἐλιὴλ, καὶ ὁ Σαμαχία, καὶ Μαὰθ, καὶ Βαναΐας, καὶ οἱ υἱοὶ αὐτοῦ καθεσταμένοι διὰ Χωνενίου καὶ Σεμεῒ τοῦ ἀδελφοῦ αὐτοῦ, καθὼς προσέταξεν Ἐζεκίας ὁ βασιλεὺς καὶ Ἀζαρίας ὁ ἡγούμενος οἴκου Κυρίου.

\vs{14}Καὶ Κορὴ ὁ τοῦ Ἰεμνὰ ὁ Λευίτης ὁ πυλωρὸς κατὰ ἀνατολὰς ἐπὶ τῶν δομάτων, δοῦναι τὰς ἀπαρχὰς Κυρίου, καὶ τὰ ἅγια τῶν ἁγίων,
\vs{15}διὰ χειρὸς Ὀδὸμ, καὶ Βενιαμὶν, καὶ Ἰησοὺς, καὶ Σεμεῒ, καὶ Ἀμαρίας, καὶ Σεχονίας, διὰ χειρὸς τῶν ἱερέων ἐν πίστει, δοῦναι τοῖς ἀδελφοῖς αὐτῶν κατὰ τὰς ἐφημερίας, κατὰ τὸν μέγαν καὶ τὸν μικρὸν,
\vs{16}ἐκτὸς τὴς ἐπιγονῆς τῶν ἀρσενικῶν ἀπὸ τριετοῦς καὶ ἐπάνω, παντὶ τῷ εἰσπορευομένῳ εἰς οἶκον Κυρίου, εἰς λόγον ἡμερῶν εἰς ἡμέραν, εἰς λειτουργείαν ἐφημερίαις διατάξεως αὐτῶν·
\vs{17}Οὗτος ὁ καταλοχισμὸς τῶν ἱερέων κατʼ οἴκους πατριῶν· καὶ οἱ Λευῖται ἐν ταῖς ἐφημερίαις αὐτῶν ἀπὸ εἰκοσαετοῦς καὶ ἐπάνω ἐν διατάξει,
\vs{18}ἐγκαταλοχίσαι ἐν πάσῃ ἐπιγονῇ υἱῶν αὐτῶν καὶ θυγατέρων αὐτῶν εἰς πᾶν πλῆθος, ὅτι ἐν πίστει ἥγνισαν τὸ ἅγιον·
\vs{19}Τοῖς υἱοῖς Ἀαρὼν τοῖς ἱερατεύουσι, καὶ οἱ ἀπὸ τῶν πόλεων αὐτῶν ἐν πάσῃ πόλει καὶ πόλει ἄνδρες οἳ ὠνομάσθησαν ἐν ὀνόματι, δοῦναι μερίδα παντὶ ἀρσενικῷ ἐν τοῖς ἱερεῦσι, καὶ παντὶ καταριθμουμένῳ ἐν τοῖς Λευίταις.

\vs{20}Καὶ ἐποίησεν οὕτως Ἐζεκίας ἐν παντὶ Ἰούδα, καὶ ἐποίησε τὸ καλὸν καὶ τὸ εὐθὲς ἐναντίον τοῦ Κυρίου Θεοῦ αὐτοῦ.
\vs{21}Καὶ ἐν παντὶ ἔργῳ ᾧ ἤρξατο ἐν ἐργασίᾳ ἐν οἴκῳ Κυρίου, καὶ ἐν τῷ νόμῳ καὶ ἐν τοῖς προστάγμασιν, ἐξεζήτησε τὸν Θεὸν αὐτοῦ ἐξ ὅλης ψυχῆς αὐτοῦ, καὶ ἐποίησε καὶ εὐοδώθη.

\ch{32}
Καὶ μετὰ τοὺς λόγους τούτους καὶ τὴν ἀλήθειαν ταύτην ἦλθε Σενναχηρὶμ βασιλεὺς Ἀσσυρίων, καὶ ἦλθεν ἐπὶ Ἰούδαν, καὶ παρενέβαλεν ἐπὶ τὰς πόλεις τὰς τειχήρεις, καὶ εἶπε προκαταλαβέσθαι αὐτάς.

\vs{2}Καὶ εἶδεν Ἐζεκίας ὅτι ἥκει Σενναχηρὶμ, καὶ τὸ πρόσωπον αὐτοῦ τοῦ πολεμῆσαι ἐπὶ Ἱερουσαλήμ.
\vs{3}Καὶ ἐβουλεύσατο μετὰ τῶν πρεσβυτέρων αὐτοῦ καὶ τῶν δυνατῶν, ἐμφράξαι τὰ ὕδατα τῶν πηγῶν ἃ ἦν ἔξω τῆς πόλεως, καὶ συνεπίσχυσαν αὐτῷ.
\vs{4}Καὶ συνήγαγε λαὸν πολὺν, καὶ ἐνέφραξε τὰ ὕδατα τῶν πηγῶν, καὶ τὸν ποταμὸν τὸν διορίζοντα διὰ τῆς πόλεως, λέγων, μὴ ἔλθῃ βασιλεὺς Ἀσσοὺρ, καὶ εὕρῃ ὕδωρ πολὺ, καὶ κατισχύσῃ.

\vs{5}Καὶ κατίσχυσεν Ἐζεκίας, καὶ ᾠκοδόμησε πᾶν τὸ τεῖχος τὸ κατεσκαμμένον, καὶ πύργους, καὶ ἔξω προτείχισμα ἄλλο, καὶ κατίσχυσε τὸ ἀνάλημμα τῆς πόλεως Δαυὶδ, καὶ κατεσκεύασεν ὅπλα πολλά.
\vs{6}Καὶ ἔθετο ἄρχοντας τοῦ πολέμου ἐπὶ τὸν λαὸν, καὶ συνήχθησαν πρὸς αὐτὸν ἐπὶ τὴν πλατεῖαν τῆς πύλης τῆς φάραγγος, καὶ ἐλάλησεν ἐπὶ καρδίαν αὐτῶν, λέγων,
\vs{7}ἰσχύσατε καὶ ἀνδρίζεσθε, καὶ μὴ φοβηθῆτε, μηδὲ πτοηθῆτε ἀπὸ προσώπου βασιλέως Ἀσσοὺρ, καὶ ἀπὸ προσώπου παντὸς τοῦ ἔθνους τοῦ μετʼ αὐτοῦ, ὅτι μεθʼ ἡμῶν πλείονες ἢ μετʼ αὐτοῦ.
\vs{8}Μετὰ αὐτοῦ βραχίονες σάρκινοι, μεθʼ ἡμῶν δὲ Κύριος ὁ Θεὸς ἡμῶν τοῦ σώζειν καὶ τοῦ πολεμεῖν τὸν πόλεμον ἡμῶν· καὶ κατεθάρσησεν ὁ λαὸς ἐπὶ τοῖς λόγοις Ἐζεκίου βασιλέως Ἰούδα.

\vs{9}Καὶ μετὰ ταῦτα ἀπέστειλε Σενναχηρὶμ βασιλεὺς Ἀσσυρίων τοὺς παῖδας ἑαυτοῦ ἐπὶ Ἱερουσαλὴμ, καὶ αὐτὸς ἐπὶ Λαχὶς, καὶ πᾶσα ἡ στρατιὰ μετʼ αὐτοῦ, καὶ ἀπέστειλε πρὸς Ἐζεκίαν βασιλέα Ἰούδα, καὶ πρὸς πάντα Ἰούδα τὸν ἐν Ἱερουσαλὴμ, λέγων,
\vs{10}οὕτως λέγει Σενναχηρὶμ βασιλεὺς Ἀσσυρίων, ἐπὶ τί ὑμεῖς πεποίθατε, καὶ καθήσεσθε ἐν τῇ περιοχῇ ἐν Ἱερουσαλήμ;
\vs{11}οὐχὶ Ἐζεκίας ἀπατᾷ ὑμᾶς τοῦ παραδοῦναι ὑμᾶς εἰς θάνατον καὶ εἰς λιμὸν καὶ εἰς δίψαν, λέγων, Κύριος ὁ Θεὸς ἡμῶν σώσει ἡμᾶς ἐκ χειρὸς βασιλέως Ἀσσούρ;
\vs{12}Οὐχ οὗτός ἐστιν Ἐζεκίας ὃς περιεῖλε τὰ θυσιαστήρια αὐτοῦ, καὶ τὰ ὑψηλὰ αὐτοῦ, καὶ εἶπε τῷ Ἰούδα καὶ τοῖς κατοικοῦσιν ἐν Ἱερουσαλὴμ, λέγων, κατέναντι τοῦ θυσιαστηρίου τούτου προσκυνήσετε, καὶ ἐπʼ αὐτῷ θυμιάσατε;
\vs{13}Οὐ γνώσεσθε ὅ, τι ἐποίησα ἐγὼ καὶ οἱ πατέρες μου πᾶσι τοῖς λαοῖς τῶν χωρῶν; μὴ δυνάμενοι ἠδύναντο θεοὶ τῶν ἐθνῶν πάσης τῆς γῆς σῶσαι τὸν λαὸν αὐτῶν ἐκ χειρός μου;
\vs{14}Τίς ἐν πᾶσι τοῖς θεοῖς τῶν ἐθνῶν τούτων οὓς ἐξωλόθρευσαν οἱ πατέρες μου; μὴ ἐδύναντο σῶσαι τὸν λαὸν αὐτῶν ἐκ χειρός μου, ὅτι δυνήσεται ὁ Θεὸς ὑμῶν σῶσαι ὑμᾶς ἐκ χειρός μου;
\vs{15}Νῦν οὖν μὴ ἀπατάτω ὑμᾶς Ἐζεκίας, καὶ μὴ πεποιθέναι ὑμᾶς ποιείτω κατὰ ταῦτα, καὶ μὴ πιστεύετε αὐτῷ, ὅτι οὐ μὴ δύνηται ὁ θεὸς παντὸς ἔθνους καὶ βασιλείας τοῦ σῶσαι τὸν λαὸν αὐτοῦ ἐκ χειρός μου καὶ ἐκ χειρὸς πατέρων μου, ὅτι ὁ Θεὸς ὑμῶν οὐ μὴ σώσει ὑμᾶς ἐκ χειρός μου.
\vs{16}Καὶ ἔτι ἐλάλησαν οἱ παῖδες αὐτοῦ ἐπὶ τὸν Κύριον Θεὸν, καὶ ἐπὶ Ἐζεκίαν παῖδα αὐτοῦ.

\vs{17}Καὶ βιβλίον ἔγραψεν ὀνειδίζειν τὸν Κύριον Θεὸν Ἰσραὴλ, καὶ εἶπε περὶ αὐτοῦ, λέγων, ὡς οἱ θεοὶ τῶν ἐθνῶν τῆς γῆς οὐκ ἐξείλαντο λαοὺς αὐτῶν ἐκ χειρός μου, οὕτως οὐ μὴ ἐξέληται ὁ Θεὸς Ἐζεκίου λαὸν αὐτοῦ ἐκ χειρός μου.
\vs{18}Καὶ ἐβόησε φωνῇ μεγάλῃ Ἰουδαϊστὶ ἐπὶ τὸν λαὸν Ἱερουσαλὴμ τὸν ἐπὶ τοῦ τείχους, τοῦ βοηθῆσαι αὐτοῖς, καὶ κατασπάσαι, ὅπως προκαταλάβωνται τὴν πόλιν.
\vs{19}Καὶ ἐλάλησεν ἐπὶ Θεὸν Ἱερουσαλὴμ, ὡς καὶ ἐπὶ θεοὺς λαῶν τῆς γῆς, ἔργα χειρῶν ἀνθρώπων.

\vs{20}Καὶ προσηύξατο Ἐζεκίας ὁ βασιλεὺς, καὶ Ἡσαΐας υἱὸς Ἀμὼς ὁ προφήτης περὶ τούτων, καὶ ἐβόησαν εἰς τὸν οὐρανόν.
\vs{21}Καὶ ἀπέστειλε Κύριος ἄγγελον, καὶ ἐξέτριψε πάντα δυνατὸν καὶ πολεμιστὴν καὶ ἄρχοντα καὶ στρατηγὸν ἐν τῇ παρεμβολῇ βασιλέως Ἀσσούρ· καὶ ἀπέστρεψε μετὰ αἰσχύνης προσώπου εἰς τὴν γῆν ἑαυτοῦ, καὶ ἦλθεν εἰς οἶκον θεοῦ αὐτοῦ· καὶ τῶν ἐξελθόντων ἐκ κοιλίας αὐτοῦ κατέβαλον αὐτὸν ἐν ῥομφαίᾳ.
\vs{22}Καὶ ἔσωσε Κύριος τὸν Ἐζεκίαν καὶ τοὺς κατοικοῦντας ἐν Ἱερουσαλὴμ ἐκ χειρὸς Σενναχηρὶμ βασιλέως Ἀσσοὺρ, καὶ ἐκ χειρὸς πάντων, καὶ κατέπαυσεν αὐτοὺς κυκλόθεν.
\vs{23}Καὶ πολλοὶ ἔφερον δῶρα τῷ Κυρίῳ εἰς Ἱερουσαλὴμ, καὶ δόματα τῷ Ἐζεκίᾳ βασιλεῖ Ἰούδα, καὶ ὑπερῄρθη κατʼ ὀφθαλμοὺς πάντων τῶν ἐθνῶν μετὰ ταῦτα.

\vs{24}Ἐν ταῖς ἡμέραις ἐκείναις ἠῤῥώστησεν Ἐζεκίας ἕως θανάτου, καὶ προσηύξατο πρὸς Κύριον· καὶ ἐπήκουσεν αὐτῷ, καὶ σημεῖον ἔδωκεν αὐτῷ.
\vs{25}Καὶ οὐ κατὰ τὸ ἀνταπόδομα ὃ ἔδωκεν αὐτῷ ἀνταπέδωκεν Ἐζεκίας, ἀλλὰ ὑψώθη ἡ καρδία αὐτοῦ, καὶ ἐγένετο ἐπʼ αὐτὸν ὀργὴ καὶ ἐπὶ Ἰούδαν καὶ Ἱερουσαλήμ.
\vs{26}Καὶ ἐταπεινώθη Ἐζεκίας ἀπὸ τοῦ ὕψους τῆς καρδίας αὐτοῦ, αὐτὸς καὶ οἱ κατοικοῦντες Ἱερουσαλὴμ, καὶ οὐκ ἐπῆλθεν ἐπʼ αὐτοὺς ὀργὴ Κυρίου ἐν ταῖς ἡμέραις Ἐζεκίου.
\vs{27}Καὶ ἐγένετο τῷ Ἐζεκίᾳ πλοῦτος καὶ δόξα πολλὴ σφόδρα· καὶ θησαυροὺς ἐποίησεν αὐτῷ ἀργυρίου καὶ χρυσίου καὶ τοῦ λίθου τοῦ τιμίου, καὶ εἰς τὰ ἀρώματα, καὶ ὁπλοθήκας, καὶ εἰς σκεύη ἐπιθυμητὰ,
\vs{28}καὶ πόλεις εἰς τὰ γεννήματα τοῦ σίτου καὶ οἴνου καὶ ἐλαίου, καὶ κώμας καὶ φάτνας παντὸς κτήνους, καὶ μάνδρας εἰς τὰ ποίμνια,
\vs{29}καὶ πόλεις ἃς ᾠκοδόμησεν αὐτῷ, καὶ ἀποσκευὴν προβάτων καὶ βοῶν εἰς πλῆθος, ὅτι ἔδωκεν αὐτῷ Κύριος ἀποσκευὴν πολλὴν σφόδρα.

\vs{30}Αὐτὸς Ἐζεκίας ἐνέφραξεν τὴν ἔξοδον τοῦ ὕδατος Γειῶν τὸ ἄνω, καὶ κατηύθυνεν αὐτὰ κάτω πρὸς Λίβα τῆς πόλεως Δαυίδ· καὶ εὐοδώθη Ἐζεκίας ἐν πᾶσι τοῖς ἔργοις αὐτοῦ.
\vs{31}Καὶ οὕτως τοῖς πρεσβευταῖς τῶν ἀρχόντων ἀπὸ Βαβυλῶνος, τοῖς ἀποσταλεῖσι πρὸς αὐτὸν πυθέσθαι παρʼ αὐτοῦ τὸ τέρας ὃ ἐγένετο ἐπὶ τῆς γῆς, ἐγκατέλιπεν αὐτὸν Κύριος τοῦ πειράσαι αὐτὸν, εἰδέναι τὰ ἐν τῇ καρδίᾳ αὐτοῦ.

\vs{32}Καὶ τὰ λοιπὰ τῶν λόγων Ἐζεκίου, καὶ τὸ ἔλεος αὐτοῦ, ἰδοὺ γέγραπται ἐν τῇ προφητείᾳ Ἡσαΐου υἱοῦ Ἀμὼς τοῦ προφήτου, καὶ ἐπὶ βιβλίου βασιλέων Ἰούδα καὶ Ἰσραήλ.
\vs{33}Καὶ ἐκοιμήθη Ἐζεκίας μετὰ τῶν πατέρων αὐτοῦ, καὶ ἔθαψαν αὐτὸν ἐν ἀναβάσει τάφων υἱῶν Δαυίδ· καὶ δόξαν καὶ τιμὴν ἔδωκαν αὐτῷ ἐν τῷ θανάτῳ αὐτοῦ πᾶς Ἰούδα, καὶ οἱ κατοικοῦντες ἐν Ἱερουσαλήμ· καὶ ἐβασίλευσε Μανασσῆς υἱὸς αὐτοῦ ἀντʼ αὐτοῦ.

\ch{33}
Ὢν δεκαδύο ἐτῶν Μανασσῆς ἐν τῷ βασιλεῦσαι αὐτὸν, καὶ πεντηκονταπέντε ἔτη ἐβασίλευσεν ἐν Ἱερουσαλήμ.
\vs{2}Καὶ ἐποίησε τὸ πονηρὸν ἐναντίον Κυρίου ἀπὸ πάντων τῶν βδελυγμάτων τῶν ἐθνῶν, οὓς ἐξωλόθρευσε Κύριος ἀπὸ προσώπου τῶν υἱῶν Ἰσραήλ.
\vs{3}Καὶ ἐπέστρεψε καὶ ᾠκοδόμησε τὰ ὑψηλὰ, ἃ κατέσπασεν Ἐζεκίας ὁ πατὴρ αὐτοῦ, καὶ ἔστησε στήλας τοῖς Βααλὶμ, καὶ ἐποίησεν ἄλση, καὶ προσεκύνησεν πάσῃ τῇ στρατιᾷ τοῦ οὐρανοῦ, καὶ ἐδούλευσεν αὐτοῖς.
\vs{4}Καὶ ᾠκοδόμησε θυσιαστήρια ἐν οἴκῳ Κυρίου, οὗ εἶπε Κύριος, ἐν Ἱερουσαλὴμ ἔσται τὸ ὄνομά μου εἰς τὸν αἰῶνα.
\vs{5}Καὶ ᾠκοδόμησε θυσιαστήρια πάσῃ τῇ στρατιᾷ τοῦ οὐρανοῦ ἐν ταῖς δυσὶν αὐλαῖς οἴκου Κυρίου.
\vs{6}Καὶ αὐτὸς διήγαγε τὰ τέκνα αὐτοῦ ἐν πυρὶ ἐν γὲ Βενεννόμ· καὶ ἐκληδονίζετο, καὶ οἰωνίζετο, καὶ ἐφαρμακεύετο, καὶ ἐποίησεν ἐγγαστριμύθους καὶ ἐπᾳοιδοὺς, καὶ ἐπλήθυνε τοῦ ποιῆσαι τὸ πονηρὸν ἐναντίον Κυρίου τοῦ παροργίσαι αὐτόν.
\vs{7}Καὶ ἔθηκε τὸ γλυπτὸν, τὸ χωνευτὸν, εἰκόνα ἣν ἐποίησεν ἐν οἴκῳ Θεοῦ, οὗ εἶπε Θεὸς πρὸς Δαυὶδ καὶ πρὸς Σαλωμὼν υἱὸν αὐτοῦ, ἐν τῷ οἴκῳ τούτῳ καὶ Ἱερουσαλὴμ ἣν ἐξελεξάμην ἐκ πασῶν φυλῶν Ἰσραὴλ, θήσω τὸ ὄνομά μου εἰς τὸν αἰῶνα.
\vs{8}Καὶ οὐ προσθήσω σαλεῦσαι τὸν πόδα Ἰσραὴλ ἀπὸ τῆς γῆς ἧς ἔδωκα τοῖς πατράσιν αὐτῶν, πλὴν ἐὰν φυλάσσωνται τοῦ ποιῆσαι πάντα ἃ ἐνετειλάμην αὐτοῖς κατὰ πάντα τὸν νόμον καὶ τὰ προστάγματα καὶ τὰ κρίματα ἐν χειρὶ Μωυσῆ.
\vs{9}Καὶ ἐπλάνησε Μανασσῆς τὸν Ἰούδαν καὶ τοὺς κατοικοῦντας ἐν Ἱερουσαλὴμ, τοῦ ποιῆσαι τὸ πονηρὸν ὑπὲρ πάντα τὰ ἔθνη ἃ ἐξῇρε Κύριος ἀπὸ προσώπου υἱῶι Ἰσραήλ.

\vs{10}Καὶ ἐλάλησε Κύριος ἐπὶ Μανασσῆ καὶ ἐπὶ τὸν λαὸν αὐτοῦ, καὶ οὐκ ἐπήκουσαν.
\vs{11}Καὶ ἤγαγε Κύριος ἐπʼ αὐτοὺς τοὺς ἄρχοντας τῆς δυνάμεως τοῦ βασιλέως Ἀσσοὺρ, καὶ κατέλαβον τὸν Μανασσῆ ἐν δεσμοῖς, καὶ ἔδησαν αὐτὸν ἐν πέδαις, καὶ ἤγαγον εἰς Βαβυλῶνα.
\vs{12}Καὶ ὡς ἐθλίβη, ἐζήτησε τὸ πρόσωπον Θεοῦ τοῦ Κυρίου αὐτοῦ, καὶ ἐταπεινώθη σφόδρα ἀπὸ προσώπου Θεοῦ πατέρων αὐτοῦ,
\vs{13}καὶ προσηύξατο πρὸς αὐτόν· καὶ ἐπήκουσεν αὐτοῦ καὶ ἐπήκουσε τῆς βοῆς αὐτοῦ, καὶ ἐπέστρεψεν αὐτὸν εἰς Ἱερουσαλὴμ ἐπὶ τὴν βασιλείαν αὐτοῦ, καὶ ἔγνω Μανασσῆς ὅτι Κύριος αὐτός ἐστι Θεός.

\vs{14}Καὶ μετὰ ταῦτα ᾠκοδόμησε τεῖχος ἔξω τῆς πόλεως Δαυὶδ ἀπὸ Λιβὸς κατὰ Νότον ἐν τῷ χειμάῤῥῳ, καὶ κατὰ τὴν εἴσοδον τὴν διὰ τῆς πύλης τῆς ἰχθυϊκῆς ἐκπορευομένων τὴν πύλην τὴν κυκλόθεν, καὶ εἰς Ὄπελ, καὶ ὕψωσε σφόδρα, καὶ κατέστησεν ἄρχοντας τῆς δυνάμεως ἐν πάσαις ταῖς πόλεσι ταῖς τειχήρεσιν ἐν Ἰούδα.
\vs{15}Καὶ περιεῖλε τοὺς θεοὺς τοὺς ἀλλοτρίους καὶ τὸ γλυπτὸν ἐξ οἴκου Κυρίου, καὶ πάντα τὰ θυσιαστήρια, ἃ ᾠκοδόμησεν ἐν ὄρει οἴκου Κυρίου καὶ ἐν Ἱερουσαλὴμ, καὶ ἔξωθεν τῆς πόλεως·
\vs{16}Καὶ κατώρθωσε τὸ θυσιαστήριον Κυρίου, καὶ ἐθυσίασεν ἐπʼ αὐτὸ θυσίαυ σωτηρίου καὶ αἰνέσεως· καὶ εἶπε τῷ Ἰούδα, τοῦ δουλεύειν Κυρίῳ Θεῷ Ἰσραήλ.
\vs{17}Πλὴν ἔτι ὁ λαὸς ἐπὶ τῶν ὑψηλῶν ἐθυσίαζε, πλὴν Κυρίῳ Θεῷ αὐτῶν.

\vs{18}Καὶ τὰ λοιπὰ τῶν λόγων Μανασσῆ, καὶ ἡ προσευχὴ αὐτοῦ ἡ πρὸς τὸν Θεὸν, καὶ λόγοι τῶν ὁρώντων τῶν λαλούντων πρὸς αὐτὸν ἐπʼ ὀνόματι Θεοῦ Ἰσραήλ,
\vs{19}ἰδοὺ ἐπὶ λόγων προσευχῆς αὐτοῦ, καὶ ἐπήκουσεν αὐτοῦ· καὶ πᾶσαι αἱ ἁμαρτίαι αὐτοῦ καὶ ἀποστάσεις αὐτοῦ, καὶ οἱ τόποι ἐφʼ οἷς ᾠκοδόμησεν ἐν αὐτοῖς τὰ ὑψηλὰ, καὶ ἔστησεν ἐκεῖ ἄλση καὶ γλυπτὰ, πρὸ τοῦ ἐπιστρέψαι, ἰδοὺ γέγραπται ἐπὶ τῶν λόγων τῶν ὁρώντων.
\vs{20}Καὶ ἐκοιμήθη Μανασσῆς μετὰ τῶν πατέρων αὐτοῦ, καὶ ἔθαψαν αὐτὸν ἐν παραδείσῳ οἴκου αὐτοῦ· καὶ ἐβασίλευσεν ἀντʼ αὐτοῦ Ἀμὼν υἱὸς αὐτοῦ.

\vs{21}Ὢν ἐτῶν εἴκοσι καὶ δύο Ἀμὼν ἐν τῷ βασιλεύειν αὐτὸν, καὶ δύο ἔτη ἑβασιλευσεν ἐν Ἱερουσαλήμ.
\vs{22}Καὶ ἐποίησε τὸ πονηρὸν ἐνώπιον Κυρίου, ὡς ἐποίησε Μανασσῆς ὁ πατὴρ αὐτοῦ· καὶ πᾶσι τοῖς εἰδώλοις οἷς ἐποίησε Μανασσῆς ὁ πατὴρ αὐτοῦ, ἔθυεν Ἀμὼν καὶ ἐδούλευσεν αὐτοῖς.
\vs{23}Καὶ οὐκ ἐταπεινώθη ἐναντίον Κυρίου ὡς ἐταπεινώθη Μανασσῆς ὁ πατὴρ αὐτοῦ, ὅτι υἱὸς αὐτοῦ Ἀμὼν ἐπλήθυνε πλημμέλειαν.
\vs{24}Καὶ ἐπέθεντο αὐτῷ οἱ παῖδες αὐτοῦ, καὶ ἐπάταξαν αὐτὸν ἐν οἴκῳ αὐτοῦ.
\vs{25}Καὶ ἐπάταξεν ὁ λαὸς τῆς γῆς τοῦς ἐπιθεμένους ἐπὶ τὸν βασιλέα Ἀμὼν, καὶ ἐβασίλευσεν ὁ λαὸς τῆς γῆς τὸν Ἰωσίαν υἱὸν αὐτοῦ ἀντʼ αὐτοῦ.

\ch{34}
Ὢν ὀκτὼ ἐτῶν Ἰωσίας ἐν τῷ βασιλεῦσαι αὐτὸν, καὶ τριάκοντα καὶ ἓν ἔτος ἐβασίλευσεν ἐν Ἱερουσαλήμ.
\vs{2}Καὶ ἐποίησε τὸ εὐθὲς ἐναντίον Κυρίου, καὶ ἐπορεύθη ἐν ὁδοῖς Δαυὶδ τοῦ πατρὸς αὐτοῦ, καὶ οὐκ ἐξέκλινε δεξιὰ καὶ ἀριστερά.
\vs{3}Καὶ ἐν τῷ ὀγδόῳ ἔτει τῆς βασιλείας αὐτοῦ, καὶ αὐτὸς ἔτι παιδάριον ἤρξατο τοῦ ζητῆσαι Κύριον τὸν Θεὸν Δαυὶδ τοῦ πατρὸς αὐτοῦ· καὶ ἐν τῷ δωδεκάτῳ ἔτει τῆς βασιλείας αὐτοῦ ἤρξατο τοῦ καθαρίσαι τὸν Ἰούδαν καὶ τὴν Ἱερουσαλὴμ ἀπὸ τῶν ὑψηλῶν, καὶ τῶν ἄλσεων, καὶ ἀπὸ τῶν περιβωμίων, καὶ ἀπὸ τῶν χωνευτῶν.
\vs{4}Καὶ κατέσπασε τὰ κατὰ πρόσωπον αὐτοῦ θυσιαστήρια τῶν Βααλὶμ, καὶ τὰ ὑψηλὰ τὰ ἐπʼ αὐτῶν· καὶ ἔκοψε τὰ ἄλση καὶ τὰ γλυπτὰ, καὶ τὰ χωνευτὰ συνέτριψε, καὶ ἐλέπτυνε καὶ ἔῤῥιψεν ἐπὶ πρόσωπον τῶν μνημάτων τῶν θυσιαζόντων αὐτοῖς.
\vs{5}Καὶ ὀστᾶ ἱερέων κατέκαυσεν ἐπὶ τὰ θυσιαστήρια, καὶ ἐκαθάρισε τὸν Ἰούδαν καὶ τὴν Ἱερουσαλὴμ,
\vs{6}καὶ ἐν πόλεσι Μανασσῆ, καὶ Ἐφραὶμ, καὶ Συμεὼν, καὶ Νεφθαλὶ, καὶ τοῖς τόποις αὐτῶν κύκλῳ.
\vs{7}Καὶ κατέσπασε τὰ θυσιαστήρια, καὶ τὰ ἄλση, καὶ εἴδωλα κατέκοψε λεπτὰ, καὶ πάντα τὰ ὑψηλὰ ἔκοψεν ἀπὸ πάσης τῆς γῆς Ἰσραὴλ, καὶ ἀπέστρεψεν εἰς Ἱερουσαλήμ.

\vs{8}Καὶ ἐν τῷ ἔτει τῷ ὀκτωκαιδεκάτῳ τῆς βασιλείας αὐτοῦ τοῦ καθαρίσαι τὴν γῆν καὶ τὸν οἶκον, ἀπέστειλε τὸν Σαφὰν υἱὸν Ἐσελία, καὶ τὸν Μαασὰ ἄρχοντα τῇς πόλεως, καὶ τὸν Ἰουὰχ υἱὸν Ἰωάχαζ τὸν ὑπομνηματογράφον αὐτοῦ, κραταιῶσαι τὸν οἶκον Κυρίου τοῦ Θεοῦ αὐτοῦ.
\vs{9}Καὶ ἦλθον πρὸς Χελκίαν τὸν ἱερέα τὸν μέγαν, καὶ ἔδωκαν τὸ ἀργύριον τὸ εἰσενεχθὲν εἰς οἶκον Θεοῦ, ὃ συνήγαγον οἱ Λευῖται φυλάσσοντες τὴν πύλην ἐκ χειρὸς Μανασσῆ καὶ Ἐφραὶμ, καὶ τῶν ἀρχόντων, καὶ ἀπὸ παντὸς καταλοίπου ἐν Ἰσραὴλ, καὶ υἱῶν Ἰούδα καὶ Βενιαμὶν, καὶ οἰκούντων ἐν Ἱερουσαλήμ.
\vs{10}Καὶ ἔδωκαν αὐτὸ ἐπὶ χεῖρα ποιούντων τὰ ἔργα, οἱ καθεσταμένοι ἐν οἴκῳ Κυρίου, καὶ ἔδωκαν αὐτὸ ποιοῦσι τὰ ἔργα οἳ ἐποίουν ἐν οἴκῳ Κυρίου, ἐπισκευάσαι καὶ κατισχύσαι τὸν οἶκον.
\vs{11}Καὶ ἔδωκαν τοῖς τέκτοσι καὶ τοῖς οἰκοδόμοις, ἀγοράσαι λίθους τετραπέδους καὶ ξύλα εἰς δοκοὺς στεγάσαι τοὺς οἴκους, οὓς ἐξωλόθρευσαν βασιλεῖς Ἰούδα.
\vs{12}Καὶ οἱ ἄνδρες ἐν πίστει ἐπὶ τῶν ἔργων· καὶ ἐπʼ αὐτῶν ἐπίσκοποι, Ἰὲθ καὶ Ἀβδίας οἱ Λευῖται ἐξ υἱῶν Μεραρὶ, καὶ Ζαχαρίας καὶ Μοσολλὰμ ἐκ τῶν υἱῶν Καὰθ ἐπισκοπεῖν, καὶ πᾶς Λευίτης, καὶ πᾶς συνιὼν ἐν ὀργάνοις ᾠδῶν.
\vs{13}Καὶ ἐπὶ τῶν νωτοφόρων, καὶ ἐπὶ πάντων τῶν ποιούντων τὰ ἔργα, ἐργασίᾳ καὶ ἐργασίᾳ· καὶ ἀπὸ τῶν Λευιτῶν γπαμματεῖς καὶ κριταὶ καὶ πυλωροί.

\vs{14}Καὶ ἐν τῷ ἐκφέρειν αὐτοὺς τὸ ἀργύριον τὸ εἰσοδιασθὲν εἰς οἶκον Κυρίου, εὗρε Χελκίας ὁ ἱερεὺς βιβλίον νόμου Κυρίου διὰ χειρὸς Μωυσῆ.
\vs{15}Καὶ ἀπεκρίθη Χελκίας, καὶ εἶπε πρὸς Σαφὰν τὸν γραμματέα, βιβλίον νόμου εὗρον ἐν οἴκῳ Κυρίου· καὶ ἔδωκε Χελκίας τὸ βιβλίον τῷ Σαφάν.
\vs{16}Καὶ εἰσήνεγκε Σαφὰν τὸ βιβλίον πρὸς τὸν βασιλέα, καὶ ἀπέδωκεν ἔτι τῷ βασιλεῖ λόγον, πᾶν τὸ δοθὲν ἀργύριον ἐν χειρὶ τῶν παίδων σον τῶν ποιούντων.
\vs{17}Καὶ ἐχώνευσαν τὸ ἀργύριον τὸ εὑρεθὲν ἐν οἴκῳ Κυρίου, καὶ ἔδωκαν ἐπὶ χεῖρα τῶν ἐπισκόπων, καὶ ἐπὶ χεῖρα τῶν ποιούντων τὴν ἐργασίαν.

\vs{18}Καὶ ἀπήγγειλε Σαφὰν ὁ γραμματεὺς τῷ βασιλεῖ λόγον, λέγων, βιβλίον δέδωκέ μοι Χελκίας ὁ ἱερεύς· καὶ ἀνέγνω αὐτὸ Σαφὰν ἐναντίον τοῦ βασιλέως.
\vs{19}Καὶ ἐγένετο ὡς ἤκουσεν ὁ βασιλεὺς τοὺς λόηους τοῦ νόμου, καὶ διέῤῥηξε τὰ ἱμάτια αὐτοῦ.
\vs{20}Καὶ ἐνετείλατο ὁ βασιλεὺς τῷ Χελκίᾳ καὶ τῷ Ἀχικὰμ υἱῷ Σαφὰν καὶ τῷ Ἀβδὸμ υἱῷ Μιχαία καὶ τῷ Σαφὰν τῷ γραμματεῖ καὶ τῷ Ἀσαΐᾳ παιδὶ τοῦ βασιλέως, λέγων,
\vs{21}πορεύθητε, ζητήσατε τὸν Κύριον περὶ ἐμοῦ καὶ περὶ παντὸς τοῦ καταλειφθέντος ἐν Ἰσραὴλ καὶ Ἰούδα περὶ τῶν λόγων τοῦ βιβλίου τοῦ εὑρεθέντος, ὅτι μέγας ὁ θυμὸς Κυρίου ἐκκέκαυται ἐν ἡμῖν, διότι οὐκ εἰσήκουσαν οἱ πατέρες ἡμῶν τῶν λόγων Κυρίου, τοῦ ποιῆσαι κατὰ πάντα τὰ γεγραμμένα ἐν τῷ βιβλίῳ τούτῳ.

\vs{22}Καὶ ἐπορεύθη Χελκίας, καὶ οἷς εἶπεν ὁ βασιλεὺς, πρὸς Ὀλδὰν τὴν προφῆτιν, γυναῖκα Σελλὴμ υἱοῦ Θεκωὲ υἱοῦ Ἀρὰς, φυλάσσουσαν τὰς ἐντολὰς, καὶ αὕτη κατῴκει ἐν Ἱερουσαλὴμ ἐν μασαναὶ, καὶ ἐλάλησαν αὐτῇ κατὰ ταῦτα.

\vs{23}Καὶ εἶπεν αὐτοῖς, οὕτως εἶπε Κύριος ὁ Θεὸς Ἰσραὴλ, εἴπατε τῷ ἀνδρὶ τῷ ἀποστείλαντι ὑμᾶς πρὸς μὲ,
\vs{24}οὕτω λέγει Κύριος, ἰδοὺ ἐγὼ ἐπάγω ἐπὶ τὸν τόπον τοῦτον κακὰ, τοὺς πάντας λόγους τοὺς γεγραμμένους ἐν τῷ βιβλίῳ τῷ ἀνεγνωσμένῳ ἐναντίον τοῦ βασιλέως Ἰούδα,
\vs{25}ἀνθʼ ὧν ἐγκατέλιπόν με καὶ ἐθυμίασαν θεοῖς ἀλλοτρίοις, ἵνα παροργίσωσίν με ἐν πᾶσιν τοῖς ἔργοις τῶν χειρῶν αὐτῶν· καὶ ἐξεκαύθη ὁ θυμός μου ἐν τῷ τόπῳ τούτῳ, καὶ οὐ σβεσθήσεται.
\vs{26}Καὶ ἐπὶ βασιλέα Ἰούδα τὸν ἀποστείλαντα ὑμᾶς τοῦ ζητῆσαι τὸν Κύριον, οὕτως ἐρεῖτε αὐτῷ, οὕτω λέγει Κύριος ὁ Θεὸς Ἰσραὴλ, τοὺς λόγους οὓς ἤκουσας,
\vs{27}καὶ ἐνετράπη ἡ καρδία σου, καὶ ἐταπεινώθης ἀπὸ προσώπου μου ἐν τῷ ἀκοῦσαί σε τοὺς λόγους μου ἐπὶ τὸν τόπον τοῦτον καὶ ἐπὶ τοὺς κατοικοῦντας αὐτὸν, καὶ ἐταπεινώθης ἐναντίον μου, καὶ διέῤῥηξας τὰ ἱμάτιά σου, καὶ ἔκλαυσας κατεναντίον μου, καὶ ἐγὼ ἤκουσα, φησὶ Κύριος.
\vs{28}Ἰδοὺ προστίθημί σε πρὸς τοὺς πατέρας σου, καὶ προστεθήσῃ πρὸς τὰ μνήματά σου ἐν εἰρήνῃ, καὶ οὐκ ὄψονται οἱ ὀφθαλμοί σου ἐν πᾶσι τοῖς κακοῖς οἷς ἐγὼ ἐπάγω ἐπὶ τὸν τόπον τοῦτον, καὶ ἐπὶ τοὺς κατοικοῦντας αὐτόν. Καὶ ἀπέδωκαν τῷ βασιλεῖ λόγον.

\vs{29}Καὶ ἀπέστειλεν ὁ βασιλεὺς, καὶ συνήγαγε τοὺς πρεσβυτέρους Ἰούδα καὶ Ἱερουσαλήμ.
\vs{30}Καὶ ἀνέβη ὁ βασιλεὺς εἰς οἶκον Κυρίου, καὶ πᾶς Ἰούδα, καὶ οἱ κατοικοῦντες Ἱερουσαλὴμ, καὶ οἱ ἱερεῖς, καὶ οἱ Λευῖται, καὶ πᾶς ὁ λαὸς ἀπὸ μικροῦ ἕως μεγάλου, καὶ ἀνέγνω ἐν ὠσὶν αὐτῶν πάντας λόγους βιβλίου τῆς διαθήκης τοὺς εὑρεθέντας ἐν οἴκῳ Κυρίου.
\vs{31}Καὶ ἔστη ὁ βασιλεὺς ἐπὶ τὸν στύλον, καὶ διέθετο διαθήκην ἐναντίον Κυρίου, τοῦ πορευθῆναι ἐνώπιον Κυρίου, τοῦ φυλάσσειν τὰς ἐντολὰς αὐτοῦ, καὶ μαρτύρια, καὶ προστάγματα αὐτοῦ ἐν ὅλῃ καρδίᾳ, καὶ ἐν ὅλῃ ψυχῇ, ὥστε ποιεῖν τοὺς λόγους τῆς διαθήκης τοὺς γεγραμμένους ἐπὶ τῷ βιβλίῳ τούτῳ.
\vs{32}Καὶ ἔστησε πάντας τοὺς εὑρεθέντας ἐν Ἱερουσαλὴμ καὶ Βενιαμίν· καὶ ἐποίησαν οἱ κατοικοῦντες Ἱερουσαλὴμ διαθήκην ἐν οἴκῳ Κυρίου Θεοῦ πατέρων αὐτῶν.

\vs{33}Καὶ περιεῖλεν Ἰωσίας τὰ πάντα βδελύγματα ἐκ πάσης τῆς γῆς, ἣ ἦν υἱῶν Ἰσραὴλ, καὶ ἐποίησε πάντας τοὺς εὑρεθέντας ἐν Ἱερουσαλὴμ καὶ ἐν Ἰσραὴλ, τοῦ δουλεύειν Κυρίῳ Θεῷ αὐτῶν πάσας τὰς ἡμέρας αὐτοῦ· οὐκ ἐξέκλινεν ἀπὸ ὄπισθε Κυρίου Θεοῦ πατέρων αὐτοῦ.

\ch{35}
Καὶ ἐποίησεν Ἰωσίας τὸ φασὲχ τῷ Κυρίῳ Θεῷ αὐτοῦ, καὶ ἔθυσε τὸ φασὲκ τῇ τεσσαρεσκαιδεκάτῃ ἡμέρᾳ τοῦ μηνὸς τοῦ πρώτου.
\vs{2}Καὶ ἔστησε τοὺς ἱερεῖς ἐπὶ τὰς φυλακὰς αὐτῶν, καὶ κατίσχυσεν αὐτοὺς εἰς τὰ ἔργα οἴκου Κυρίου.
\vs{3}Καὶ εἶπε τοῖς Λευίταις τοῖς δυνατοῖς ἐν παντὶ Ἰσραὴλ, τοῦ ἁγιασθῆναι αὐτοὺς τῷ Κυρίῳ· καὶ ἔθηκαν τὴν κιβωτὸν τὴν ἁγίαν εἰς τὸν οἶκον ὃν ᾠκοδόμησε Σαλωμὼν υἱὸς Δαυὶδ τοῦ βασιλέως Ἰσραήλ. καὶ εἶπεν ὁ βασιλεὺς, οὐκ ἔστιν ὑμῖν ἐπʼ ὤμων ἆραι οὐδέν· νῦν οὖν λειτουργήσατε τῷ Κυρίῳ Θεῷ ὑμῶν, καὶ τῷ λαῷ αὐτοῦ Ἰσραήλ.
\vs{4}Καὶ ἑτοιμάσθητε κατʼ οἴκους πατριῶν ὑμῶν, καὶ κατὰ τὰς ἐφημερίας ὑμῶν, κατὰ τὴν γραφὴν Δαυὶδ βασιλέως Ἰσραὴλ, καὶ διὰ χειρὸς Σαλωμὼν υἱοῦ αὐτοῦ.
\vs{5}Καὶ στῆτε ἐν τῷ οἴκῳ κατὰ τὰς διαιρέσεις οἴκων πατριῶν ὑμῶν τοῖς ἀδελφοῖς ὑμῶν υἱοῖς τοῦ λαοῦ, καὶ μερὶς οἴκου πατριᾶς τοῖς Λευίταις.
\vs{6}Καὶ θύσατε τὸ φασὲκ, καὶ ἑτοιμάσατε τοῖς ἀδελφοῖς ὑμῶν, τοῦ ποιῆσαι κατὰ τὸν λόγον Κυρίου διὰ χειρὸς Μωυσῆ.

\vs{7}Καὶ ἀπήρξατο Ἰωσίας τοῖς υἱοῖς τοῦ λαοῦ πρόβατα, καὶ ἀμνοὺς, καὶ ἐρίφους ἀπὸ τῶν τέκνων τῶν αἰγῶν, πάντα εἰς τὸ φασὲκ, καὶ πάντας τοὺς εὑρεθέντας εἰς ἀριθμὸν τριάκοντα χιλιάδας, καὶ μόσχων τρεῖς χιλιάδας, ταῦτα ἀπὸ τῆς ὑπάρξεως τοῦ βασιλέως.
\vs{8}Καὶ οἱ ἄρχοντες αὐτοῦ ἀπήρξαντο τῷ λαῷ καὶ τοῖς ἱερεῦσι καὶ τοῖς Λευίταις· ἔδωκε δὲ Χελκίας καὶ Ζαχαρίας καὶ Ἰεϊὴλ οἱ ἄρχοντες τοῖς ἱερεῦσιν οἴκου Θεοῦ, καὶ ἔδωκαν εἰς τὸ φασὲκ πρόβατα καὶ ἀμνοὺς καὶ ἐρίφους διαχίλια ἑξακόσια, καὶ μόσχους τριακοσίους.
\vs{9}Καὶ Χωνενίας, καὶ Βαναίας, καὶ Σαμαίας, καὶ Ναθαναὴλ ἀδελφὸς αὐτοῦ, καὶ Ἀσαβίας, καὶ Ἰεϊὴλ, καὶ Ἰωζαβὰδ, ἄρχοντες τῶν Λευιτῶν, ἀπήρξαντο τοῖς Λευίταις εἰς τὸ φασὲχ πρόβατα πεντακισχίλια, καὶ μόσχους πεντακοσίους.

\vs{10}Καὶ κατωρθώθη ἡ λειτουργεία, καὶ ἔστησαν οἱ ἱερεῖς ἐπὶ τὴν στάσιν αὐτῶν, καὶ οἱ Λευῖται ἐπὶ τὰς διαιρέσεις αὐτῶν κατὰ τὴν ἐντολὴν τοῦ βασιλέως.
\vs{11}Καὶ ἔθυσαν τὸ φασὲκ, καὶ προσέχεαν οἱ ἱερεῖς τὸ αἷμα ἐκ χειρὸς αὐτῶν, καὶ οἱ Λευῖται ἐξέδειραν.
\vs{12}Καὶ ἡτοίμασαν τὴν ὁλοκαύτωσιν παραδοῦναι αὐτοῖς κατὰ τὴν διαίρεσιν κατʼ οἴκους πατριῶν τοῖς υἱοῖς τοῦ λαοῦ, τοῦ προσάγειν τῷ Κυρίῳ, ὡς γέγραπται ἐν βίβλῳ Μωυσῆ· καὶ οὕτως εἰς τὸ πρωΐ.
\vs{13}Καὶ ὤπτησαν τὸ φασὲκ ἐν πυρὶ κατὰ τὴν κρίσιν, καὶ τὰ ἅγια ἥψησαν ἐν τοῖς χαλκείοις καὶ ἐν τοῖς λέβησι, καὶ εὐωδώθη, καὶ ἔδραμον πρὸς πάντας τοὺς υἱοὺς τοῦ λαοῦ.

\vs{14}Καὶ μετὰ τὸ ἑτοιμάσαι αὐτοῖς καὶ τοῖς ἱερεῦσιν ὅτι οἱ ἱερεῖς ἐν τῷ ἀναφέρειν τὰ ὁλοκαυτώματα καὶ τὰ στέατα ἕως νυκτὸς, καὶ οἱ Λευῖται ἡτοίμασαν αὐτοῖς, καὶ τοῖς ἀδελφοῖς αὐτῶν υἱοῖς Ἀαρών.
\vs{15}Καὶ οἱ ψαλτῳδοὶ υἱοὶ Ἀσὰφ ἐπὶ τῆς στάσεως αὐτῶν κατὰ τὰς ἐντολὰς Δαυὶδ, καὶ Ἀσὰφ, καὶ Αἰμὰν, καὶ Ἰδιθὼμ οἱ προφῆται τοῦ βασιλέως· καὶ οἱ ἄρχοντες καὶ οἱ πυλωροὶ πύλης καὶ πύλης, οὐκ ἦν αὐτοῖς κινεῖσθαι ἀπὸ τῆς λειτουργίας τῶν ἁγίων, ὅτι οἱ ἀδελφοὶ αὐτῶν οἱ Λευῖται ἡτοίμασαν αὐτοῖς.
\vs{16}Καὶ κατωρθώθη καὶ ἡτοιμάσθη πᾶσα ἡ λειτουργία Κυρίου ἐν τῇ ἡμέρᾳ ἐκείνῃ τοῦ ποιῆσαι τὸ φασὲκ, καὶ ἐνεγκεῖν τὰ ὁλοκαυτώματα ἐπὶ τὸ θυσιαστήριον Κυρίου κατὰ τὴν ἐντολὴν τοῦ βασιλέως Ἰωσίου.
\vs{17}Καὶ ἐποίησαν οἱ υἱοὶ Ἰσραὴλ οἱ εὑρεθέντες τὸ φασὲκ ἑν τῷ καιρῷ ἐκείνῳ, καὶ τὴν ἑορτὴν τῶν ἀζύμων ἑπτὰ ἡμέρας.

\vs{18}Καὶ οὐκ ἐγένετο φασὲκ ὅμοιον αὐτῷ ἐν Ἰσραὴλ, ἀπὸ ἡμερῶν Σαμουὴλ τοῦ προφήτου καὶ παντὸς βασιλέως Ἰσραὴλ· οὐκ ἐποίησαν τὸ φασὲκ ὃ ἐποίησεν Ἰωσίας, καὶ οἱ ἱερεῖς, καὶ οἱ Λευῖται, καὶ πᾶς Ἰούδα καὶ Ἰσραὴλ ὁ εὑρεθεὶς, καὶ οἱ κατοικοῦντες ἐν Ἱερουσαλὴμ, τῷ Κυρίῳ.
\vs{19}Τῷ ὀκτωκαιδεκάτῳ ἔτει τῆς βασιλείας Ἰωσίου ἐποιήθη τὸ φασὲκ τοῦτο· μετὰ ταῦτα πάντα ἃ ἔδρασεν Ἰωσίας ἐν τῷ οἴκῳ,
\vs{19a}καὶ τοὺς ἐγγαστριμύθους καὶ τοὺς γνώστας καὶ τὰ θεραφὶν καὶ τὰ εἴδωλα καὶ τὰ καρησὶμ ἃ ἦν ἐν γῇ Ἰούδα καὶ ἐν Ἱερουσαλὴμ, ἐνεπύρισεν ὁ βασιλεὺς Ἰωσίας, ἵνα στήσῃ τοὺς λόγους τοῦ νόμου τοὺς γεγραμμένους ἐπὶ τοῦ βιβλίου οὗ εὗρε Χελκίας ὁ ἱερεὺς ἐν τῷ οἴκῳ Κυρίου·
\vs{19b}ὅμοιος αὐτῷ οὐκ ἐγενήθη ἔμπροσθεν αὐτοῦ, ὃς ἐπέστρεψε πρὸς Κύριον ἐν ὅλῃ καρδίᾳ αὐτοῦ καὶ ἐν ὅλῃ ψυχῇ αὐτοῦ καὶ ἐν ὅλῃ τῇ ἰσχύϊ αὐτοῦ κατὰ πάντα τὸν νόμον Μωυσῆ, καὶ μετʼ αὐτὸν οὐκ ἀνέστη ὅμοιος·
\vs{19c}πλὴν οὐκ ἀπεστράφη Κύριος ἀπὸ ὀργῆς θυμοῦ αὐτοῦ τοῦ μεγάλου, οὗ ὠργίσθη θυμῷ Κύριος ἐν τῷ Ἰούδα, ἐπὶ πάντα τὰ παροργίσματα αὐτοῦ ἂ παρώργισε Μανασσῆς·
\vs{19d}καὶ εἶπε Κύριος, καί γε τὸν Ἰούδαν ἀποστήσω ἀπὸ προσώπου μου, καθὼς ἀπέστησα τὸν Ἰσραὴλ, καὶ ἀπωσάμην τὴν πόλιν ἣν ἐξελεξάμην τὴν Ἱερουσαλὴμ, καὶ τὸν οἶκον ὃν εἶπα, ἔσται τὸ ὄνομά μου ἐκεῖ.

\vs{20}Καὶ ἀνέβη Φαραὼ Νεχαὼ βασιλεὺς Αἰγύπτου ἐπὶ τὸν βασιλέα Ἀσσυρίων ἐπὶ τὸν ποταμὸν Εὐφράτην, καὶ ἐπορεύθη βασιλεὺς Ἰωσίας εἰς συνάντησιν αὐτῷ.
\vs{21}Καὶ ἀπέστειλε πρὸς αὐτὸν ἀγγέλους, λέγων, τί ἐμοὶ καὶ σοί βασιλεῦ Ἰούδα; οὐκ ἐπὶ σὲ ἥκω σήμερον πόλεμον πολεμῆσαι· καὶ ὁ Θεὸς εἶπε τοῦ κατασπεῦσαι με· πρόσεχε ἀπὸ τοῦ Θεοῦ τοῦ μετʼ ἐμοῦ, μὴ καταφθείρῃ σε.
\vs{22}Καὶ οὐκ ἀπέστρεψεν Ἰωσίας τὸ πρόσωπον αὐτοῦ ἀπʼ αὐτοῦ, ἀλλʼ ἢ πολεμεῖν αὐτὸν ἐκραταιώθη, καὶ οὐκ ἤκουσε τῶν λόγων Νεχαὼ διὰ στόματος Θεοῦ, καὶ ἦλθε τοῦ πολεμῆσαι ἐν τῷ πεδίῳ Μαγεδδώ.
\vs{23}Καὶ ἐτόξευσαν οἱ τοξόται ἐπὶ βασιλέα Ἰωσίαν· καὶ εἶπεν ὁ βασιλεὺς τοῖς παισὶν αὐτοῦ, ἐξαγάγετέ με, ὅτι ἐπόνεσα σφόδρα.
\vs{24}Καὶ ἐξήγαγον αὐτὸν οἱ παῖδες αὐτοῦ ἀπὸ τοῦ ἅρματος, καὶ ἀνεβίβασαν αὐτὸν ἐπὶ τὸ ἅρμα τὸ δευτερεῦον ὃ ἦν αὐτῷ, καὶ ἤγαγον αὐτὸν εἰς Ἱερουσαλὴμ, καὶ ἀπέθανε, καὶ ἐτάφη μετὰ τῶν πατέρων αὐτοῦ· καὶ πᾶς Ἰούδα καὶ Ἱερουσαλὴμ ἐπένθησαν ἐπὶ Ἰωσίαν.
\vs{25}Καὶ ἐθρήνησεν Ἱερεμίας ἐπὶ Ἰωσίαν, καὶ εἶπαν πάντες οἱ ἄρχοντες καὶ αἱ ἄρχουσαι θρῆνον ἐπὶ Ἰωσίαν ἕως τῆς σήμερον· καὶ ἔδωκαν αὐτὸν εἰς πρόσταγμα ἐπὶ Ἰσραὴλ, καὶ ἰδοὺ γέγραπται ἐπὶ τῶν θρήνων.

\vs{26}Καὶ ἦσαν οἱ λοιποὶ λόγοι Ἰωσίου καὶ ἡ ἐλπὶς αὐτοῦ γεγραμμένα ἐν νόμῳ Κυρίου·
\vs{27}καὶ οἱ λόγοι αὐτοῦ οἱ πρῶτοι καὶ οἱ ἔσχατοι, ἰδοὺ γεγραμμένοι ἐπὶ βιβλίῳ βασιλέων Ἰσραὴλ καὶ Ἰούδα.

\ch{36}
Καὶ ἔλαβεν ὁ λαὸς τῆς γῆς τὸν Ἰωάχαζ υἱὸν Ἰωσίου, καὶ ἔχρισαν αὐτὸν, καὶ κατέστησαν αὐτὸν ἀντὶ τοῦ πατρὸς αὐτοῦ εἰς βασιλέα ἐπὶ Ἱερουσαλήμ.
\vs{2}Υἱὸς εἴκοσι καὶ τριῶν ἐτῶν Ἰωάχαζ ἐν τῷ βασιλεύειν αὐτὸν, καὶ τρίμηνον ἐβασίλευσεν ἐν Ἱερουσαλὴμ,
\vs{2a}καὶ ὄνομα τῆς μητρὸς αὐτοῦ Ἀμιτὰλ, θυγάτηρ Ἱερεμίου ἐκ Λοβνά·
\vs{2b}καὶ ἐποίησε τὸ πονηρὸν ἐνώπιον Κυρίου κατὰ πάντα ἃ ἐποίησαν οἱ πατέρες αὐτοῦ·
\vs{2c}καὶ ἔδησεν αὐτὸν Φαραὼ Νεχαὼ ἐν Δεβλαθὰ ἐν γῇ Αἰμὰθ, τοῦ μὴ βασιλεύειν αὐτὸν ἐν Ἱερουσαλήμ.
\vs{3}Καὶ μετήγαγεν αὐτὸν ὁ βασιλεὺς εἰς Αἴγυπτον, καὶ ἐπέβαλε φόρον ἐπὶ τὴν γῆν, ἑκατὸν τάλαντα ἀργυρίου καὶ τάλαντον χρυσίου.
\vs{4}Καὶ κατέστησε Φαραὼ Νεχαὼ τὸν Ἐλιακὶμ υἱὸν Ἰωσίου βασιλέα ἐπὶ Ἰούδα ἀντὶ Ἰωσίου τοῦ πατρὸς αὐτοῦ, καὶ μετέστρεψε τὸ ὄνομα αὐτοῦ Ἰωακίμ· καὶ τὸν Ἰωαχαζ ἀδελφὸν αὐτοῦ ἔλαβε Φαραὼ Νεχαὼ, καὶ εἰσήγαγεν αὐτὸν εἰς Αἴγυπτον, καὶ ἀπέθανεν ἐκεῖ·
\vs{4a}καὶ τὸ ἀργύριον καὶ τὸ χρυσίον ἔδωκε τῷ Φαραῷ· τότε ἤρξατο ἡ γῆ φορολογεῖσθαι τοῦ δοῦναι τὸ ἀργύριον ἐπὶ στόμα Φαραώ· καὶ ἕκαστος κατὰ δύναμιν ἀπῄτει τὸ ἀργύριον καὶ τὸ χρυσίον παρὰ τοῦ λαοῦ τῆς γῆς, δοῦναι Φαραῷ Νεχαῷ.

\vs{5}Ὢν εἴκοσι καὶ πέντε ἐτῶν Ἰωακὶμ ἐν τῷ βασιλεύειν αὐτὸν, καὶ ἕνδεκα ἔτη ἐβασίλευσεν ἐν Ἱερουσαλὴμ, καὶ ὄνομα τῆς μητρὸς αὐτοῦ Ζεχωρὰ, θυγάτηρ Νηρίου ἐκ Ῥαμά. καὶ ἐποίησε τὸ πονηρὸν ἐναντίον Κυρίου κατὰ πάντα ὅσα ἐποίησαν οἱ πατέρες αὐτοῦ.
\vs{5a}Ἐν ταῖς ἡμέραις αὐτοῦ ἦλθε Ναβουχοδονόσορ ὁ βασιλεὺς Βαβυλῶνος εἰς τὴν γῆν, καὶ ἦν αὐτῷ δουλεύων τρία ἔτη, καὶ ἀπέστη ἀπʼ αὐτοῦ·
\vs{5b}καὶ ἀπέστειλε Κύριος ἐπʼ αὐτοὺς τοὺς Χαλδαίους, καὶ λῃστήρια Σύρων, καὶ λῃστήρια Μωαβιτῶν, καὶ υἱῶν Ἀμμὼν, καὶ τῆς Σαμαρείας, καὶ ἀπέστησαν μετὰ τὸν λόγον τοῦτον κατὰ τὸν λόγον Κυρίου ἐν χειρὶ τῶν παίδων αὐτοῦ τῶν προφητῶν·
\vs{5c}πλὴν θυμὸς Κυρίου ἦν ἐπὶ Ἰούδαν, τοῦ ἀποστῆναι αὐτὸν ἀπὸ προσώπου αὐτοῦ διὰ τὰς ἁμαρτίας Μανασσῆ ἐν πᾶσιν οἷς ἐποίησε,
\vs{5d}καὶ ἐν αἵματι ἀθώῳ ᾧ ἐξέχεεν Ἰωακὶμ, καὶ ἔπλησε τὴν Ἱερουσαλὴμ αἵματος ἀθώου, καὶ οὐκ ἠθέλησε Κύριος ἐξολοθρεῦσαι αὐτούς.
\vs{6}Καὶ ἀνέβη ἐπʼ αὐτὸν Ναβουχοδονόσορ βασιλεὺς Βαβυλῶνος, καὶ ἔδησεν αὐτὸν ἐν χαλκαῖς πέδαις, καὶ ἀπήγαγεν αὐτὸν εἰς Βαβυλῶνα.
\vs{7}Καὶ μήρος τῶν σκευῶν σἴκου Κυρίου ἀπήνεγκεν εἰς Βαβυλῶνα, καὶ ἔθηκεν αὐτὰ ἐν τῷ ναῷ αὐτοῦ ἐν Βαβυλῶνι.

\vs{8}Καὶ τὰ λοιπὰ τῶν λόγων Ἰωακὶμ καὶ πάντα ἃ ἐποίησεν, οὐκ ἰδοὺ ταῦτα γεγραμμένα ἐν βιβλίῳ λόγων τῶν ἡμερῶν τοῖς βασιλεῦσιν Ἰούδα; καὶ ἐκοιμήθη Ἰωακὶμ μετὰ τῶν πατέρων αὐτοῦ, καὶ ἐτάφη ἐν γανοζαὴ μετὰ τὼν πατέρων αὐτοῦ, καὶ ἐβασίλευσεν Ἰεχονίας υἱὸς αὐτοῦ ἀντʼ αὐτοῦ.

\vs{9}Ὀκτὼ ἐτῶν Ἰεχονίας ἐν τῷ βασιλεύειν αὐτὸν, καὶ τρίμηνον καὶ δέκα ἡμέρας ἐβασίλευσεν ἐν Ἱερουσαλὴμ, καὶ ἐποίησε τὸ πονηρὸν ἐνώπιον Κυρίου.
\vs{10}Καὶ ἐπιστρέφοντος τοῦ ἐνιαυτοῦ, ἀπέστειλεν ὁ βασιλεὺς Ναβουχοδονόσορ, καὶ εἰσήνεγκεν αὐτὸν εἰς Βαβυλῶνα μετὰ τῶν σκευῶν τῶν ἐπιθυμητῶν οἴκου Κυρίου· καὶ ἐβασίλευσε Σεδεκίαν ἀδελφὸν τοῦ πατρὸς αὐτοῦ ἐπὶ Ἰούδαν καὶ Ἱερουσαλήμ.

\vs{11}Ἐτῶν εἴκοσι υἱὸς καὶ ἑνὸς ἔτου Σεδεκίας ἐν τῷ βασιλεύειν αὐτὸν, καὶ ἕνδεκα ἔτη ἐβασίλευσεν ἐν Ἱερουσαλήμ.
\vs{12}Καὶ ἐποίησε τὸ πονηρὸν ἐνώπιον Κορίου Θεοῦ αὐτοῦ, οὐκ ἐνετράπη ἀπὸ προσώπου Ἱερεμίου τοῦ προφήτου καὶ ἐκ στόματος Κυρίου,
\vs{13}ἐν τῷ τὰ πρὸς τὸν βασιλέα Ναβουχοδονόσορ ἀθετῆσαι, ἃ ὥρκισεν αὐτὸν κατὰ τοῦ Θεοῦ, καὶ ἐσκλήρυνε τὸν τράχηλον αὐτοῦ καὶ τὴν καρδίαν αὐτοῦ κατίσχυσε, τοῦ μὴ ἐπιστρέψαι πρὸς Κύριον Θεὸν Ἰσραήλ.
\vs{14}Καὶ πάντες οἱ ἔνδοξοι Ἰούδα, καὶ οἱ ἱερεῖς, καὶ ὁ λαὸς τῆς γῆς ἐπλήθυναν τοῦ ἀθετῆσαι ἀθετήματα βδελυγμάτων ἐθνῶν, καὶ ἐμίαναν τὸν οἶκον Κυρίου τὸν ἐν Ἱερουσαλήμ.
\vs{15}Καὶ ἐξαπέστειλε Κύριος ὁ Θεὸς τῶν πατέρων αὐτῶν ἐν χειρὶ τῶν προθητῶν αὐτοῦ, ὀρθρίζων καὶ ἀποστέλλων τοὺς ἀγγέλους αὐτοῦ, ὅτι ἦν φειδόμενος τοῦ λαοῦ αὐτοῦ, καὶ τοῦ ἁγιάσματος αὐτοῦ.
\vs{16}Καὶ ἦσαν μυκτηρίζοντες τοὺς ἀγγέλους αὐτοῦ, καὶ ἐξουθενοῦντες τοὺς λόγους αὐτοῦ, καὶ ἐμπαίζοντες ἐν τοῖς προφήταις αὐτοῦ, ἕως ἀνέβη ὁ θυμὸς Κυρίου ἐν τῷ λαῷ αὐτοῦ, ἕως οὐκ ἦν ἴαμα.

\vs{17}Καὶ ἤγαγεν ἐπʼ αὐτοὺς βασιλέα Χαλδαίων, καὶ ἀπέκτεινε τοὺς νεανίσκους αὐτῶν ἐν ῥομφαίᾳ ἐν οἴκῳ ἁγιάσματος αὐτοῦ· καὶ οὐκ ἐφείσατο τοῦ Σεδεκίου, καὶ τὰς παρθένους αὐτῶν οὐκ ἠλέησε, καὶ τοὺς πρεσβυτέρους αὐτῶν ἀπήγαγον· τὰ πάντα παρέδωκεν ἐν χερσὶν αὐτῶν.
\vs{18}Καὶ πάντα τὰ σκεύη οἴκου τοῦ Θεοῦ τὰ μεγάλα καὶ τὰ μικρὰ, καὶ τοὺς θησαυροὺς, οἴκου καὶ Κυρίου, καὶ πάντας τοὺς θησαυροὺς τοῦ βασιλέως καὶ τῶν μεγιστάνων, πάντα εἰσήνεγκεν εἰς Βαβυλῶνα.

\vs{19}Καὶ ἐνέπρησε τὸν οἶκον Κυρίου, καὶ κατέσκαψε τὸ τεῖχος Ἱερουσαλὴμ, καὶ τὰς βάρεις αὐτῆς ἐνέπρησεν ἐν πυρί, καὶ πᾶν σκεῦος ὡραῖον εἰς ἀφανισμόν.
\vs{20}Καὶ ἀπῴκισε τοὺς καταλοίπους εἰς Βαβυλῶνα, καὶ ἦσαν αὐτῷ καὶ τοῖς υἱοῖς αὐτοῦ εἰς δούλους ἕως βασιλείας Μήδων,
\vs{21}τοῦ πληρωθῆναι λόγον Κυρίου διὰ στόματος Ἱερεμίου, ἕως τοῦ προσδέξασθαι τὴν γῆν τὰ σάββατα αὐτῆς σαββατίσαι, πάσας τὰς ἡμέρας ἐρημώσεως αὐτῆς σαββατίσαι εἰς συμπλήρωσιν ἐτῶν ἑβδομήκοντα.

\vs{22}Ἔτους πρώτου Κύρου βασιλέως Περσῶν, μετὰ τὸ πληρωθῆναι ῥῆμα Κυρίου διὰ στὸματος Ἱερεμίου, ἐξήγειρε Κύριος τὸ πνεῦμα Κύρου βασιλέως Περσῶν, καὶ παρήγγειλε κηρύξαι ἐν πάσῃ τῇ βασιλείᾳ αὐτοῦ ἑν γραπτῷ, λέγων,

\vs{23}Τάδε λέγει Κῦρος βασιλεὺς Περσῶν πάσαις ταῖς βασιλείαις τῆς γῆς, ἔδωκέ μοι Κύριος ὁ Θεὸς τοῦ οὐρανοῦ, καὶ αὐτὸς ἐνετείλατό μοι οἰκοδομῆσαι οἶκον αὐτῷ ἐν Ἱερουσαλὴμ ἐν τῇ Ἰουδαίᾳ· τίς ἐξ ὑμῶν ἐκ παντὸς τοῦ λαοῦ αὐτοῦ; ἔσται Θεὸς αὐτοῦ μετʼ αὐτοῦ, καὶ ἀναβήτω.


\def\book{ΕΣΔΡΑΣ}
\biblebook{ΕΣΔΡΑΣ}


\lettrine[lines=2, loversize=0.2, nindent=0em, findent=.25em]{\textcolor{bookheadingcolor}{Κ}}{ΑΙ} ἐν τῷ πρώτῳ ἔτει Κύρου τοῦ βασιλέως Περσῶν, τοῦ τελεσθῆναι λόγον Κυρίου ἀπὸ στόματος Ἱερεμίου, ἐξήγειρε Κύριος τὸ πνεῦμα Κύρου βασιλέως Περσῶν, καὶ παρήγγειλε φωνὴν ἐν πάσῃ βασιλείᾳ αὐτοῦ, καί γε ἐν γραπτῷ, λέγων,

\vs{2}Οὕτως εἶπε Κύρος βασιλεὺς Περσῶν, πάσας τὰς βασιλείας τῆς γῆς ἔδωκέ μοι Κύριος ὁ Θεὸς τοῦ οὐρανοῦ, καὶ αὐτὸς ἐπεσκέψατο ἐπʼ ἐμὲ τοῦ οἰκοδομῆσαι οἶκον αὐτῷ ἐν Ἱερουσαλὴμ τῇ ἐν τῇ Ἰουδαίᾳ.
\vs{3}Τίς ἐν ὑμῖν ἀπὸ παντὸς τοῦ λαοῦ αὐτοῦ; καὶ ἔσται ὁ Θεὸς αὐτοῦ μετʼ αὐτοῦ, καὶ ἀναβήσεται εἰς Ἱερουσαλὴμ τὴν ἐν τῇ Ἰουδαίᾳ, καὶ οἰκοδομησάτω τὸν οἶκον Θεοῦ Ἰσραὴλ· αὐτὸς ὁ Θεὸς ὁ ἐν Ἱερουσαλήμ.
\vs{4}Καὶ πᾶς ὁ καταλιπόμενος ἀπὸ πάντων τῶν τόπων οὗ αὐτὸς παροικεῖ ἐκεῖ, καὶ λήψονται αὐτὸν ἄνδρες τοῦ τόπου αὐτοῦ ἐν ἀργυρίῳ, καὶ χρυσίῳ, καὶ ἀποσκευῇ, καὶ κτήνεσι μετὰ τοῦ ἑκουσίου εἰς οἶκον τοῦ Θεοῦ τὸν ἐν Ἱερουσαλήμ.

\vs{5}Καὶ ἀνέστησαν ἄρχοντες τῶν πατριῶν τῶν Ἰούδα καὶ Βενιαμεὶν, καὶ οἱ ἱερεῖς καὶ οἱ Λευεῖται, πάντων ὧν ἐξήγειρεν ὁ Θεὸς τὸ πνεῦμα αὐτῶν τοῦ ἀναβῆναι οἰκοδομῆσαι τὸν οἶκον Κυρίου τὸν ἐν Ἱερουσαλήμ.
\vs{6}Καὶ πάντες οἱ κυκλόθεν ἐνίσχυσαν ἐν χερσὶν αὐτῶν ἐν σκεύεσιν ἀργυρίου, ἐν χρυσῷ, ἐν ἀποσκευῇ, καὶ ἐν κτήνεσι, καὶ ἐν ξενίοις, πάρεξ τῶν ἑκουσίων.

\vs{7}Καὶ ὁ βασιλεὺς Κύρος ἐξήνεγκε τὰ σκεύη οἴκου Κυρίου, ἃ ἔλαβε Ναβουχοδονόσορ ἀπὸ Ἱερουσαλὴμ καὶ ἔδωκεν αὐτὰ ἐν οἴκῳ θεοῦ αὐτοῦ·
\vs{8}Καὶ ἐξήνεγκεν αὐτὰ Κύρος ὁ βασιλεὺς Περσῶν ἐπὶ χεῖρα Μιθραδάτου γασβαρηνοῦ, καὶ ἠρίθμησεν αὐτὰ τῷ Σασαβασὰρ τῷ ἄρχοντι τοῦ Ἰούδα.
\vs{9}Καὶ οὗτος ὁ ἀριθμὸς αὐτῶν· ψυκτῆρες χρυσοῖ τριάκοντα, καὶ ψυκτῆρες ἀργυροῖ χίλιοι, παρηλλαγμένα ἐννέα καὶ εἴκοσι, κεφουρῆς χρυσοῖ τριάκοντα, καὶ ἀργυροῖ διπλοῖ τετρακόσια δέκα,
\vs{10}καὶ σκεύη ἕτερα χίλια.
\vs{11}Πάντα τὰ σκεύη τῷ χρυσῷ καὶ τῷ ἀργυρῷ πεντακισχίλια τετρακόσια, τὰ πάντα ἀναβαίνοντα μετὰ Σασαβασὰρ ἀπὸ τῆς ἀποικίας ἐκ Βαβυλῶνος εἰς Ἱερουσαλήμ.

\ch{2}
Καὶ οὗτοι οἱ υἱοὶ τῆς χώρας οἱ ἀναβαίνοντες ἀπὸ τῆς αἰχμαλωσίας τῆς ἀποικίας, ἧς ἀπῴκισε Ναβουχοδονόσορ βασιλεὺς Βαβυλῶνος εἰς Βαβυλῶνα, καὶ ἐπέστρεψαν εἰς Ἱερουσαλὴμ καὶ Ἰούδα ἀνὴρ εἰς πόλιν αὐτοῦ·
\vs{2}Οἳ ἦλθον μετὰ Ζοροβάβελ, Ἰησοῦς, Νεεμίας, Σαραΐας, Ῥεελίας, Μαρδοχαῖος, Βαλασὰν, Μασφὰρ, Βαγουαὶ, Ῥεοὺμ, Βαανά· ἀνδρῶν ἀριθμὸς λαοῦ Ἰσραήλ.

\vs{3}Υἱοὶ Φαρὲς, δισχίλιοι ἑκατὸν ἑβδομηκονταδύο.

\vs{4}Υἱοὶ Σαφατία, τριακόσιοι ἑβδομηκονταδύο.

\vs{5}Υἱοὶ Ἄρες, ἑπτακόσιοι ἑβδομηκονταπέντε.

\vs{6}Υἱοὶ Φαὰθ Μωὰβ τοῖς υἱοῖς Ἰησουὲ Ἰωὰβ, δισχίλιοι ὀκτακόσιοι δεκαδύο.

\vs{7}Υἱοὶ Αἰλὰμ, χίλιοι διακόσιοι πεντηκοντατέσσαρες.

\vs{8}Υἱοὶ Ζατθουὰ, ἐννακόσιοι τεσσαρακονταπέντε.

\vs{9}Υἱοὶ Ζακχοὺ, ἑπτακόσιοι ἑξήκοντα.

\vs{10}Υἱοὶ Βανουὶ, ἑξακόσιοι τεσσαρακονταδύο.

\vs{11}Υἱοὶ Βαβαῒ, ἑξακόσιοι εἰκοσιτρεῖς.

\vs{12}Υἱοὶ Ἀσγὰδ, χίλιοι διακόσιοι εἰκοσιδύο.

\vs{13}Υἱοὶ Ἀδωνικὰμ, ἑξακόσιοι ἑξηκονταέξ.

\vs{14}Υἱοὶ Βαγουὲ, δισχίλιοι πεντηκονταέξ.

\vs{15}Υἱοὶ Ἀδδὶν, τετρακόσιοι πεντηκοντατέσσαρες.

\vs{16}Υἱοὶ Ἀτὴρ τῷ Ἐζεκίᾳ, ἐννενηκονταοκτώ.

\vs{17}Υἱοὶ Βασσοῦ, τριακόσιοι εἰκοσιτρεῖς.

\vs{18}Υἱοὶ Ἰωρὰ, ἑκατὸν δεκαδύο.

\vs{19}Υἱοὶ Ἀσοὺμ, διακόσιοι εἰκοσιτρεῖς.

\vs{20}Υἱοὶ Γαβὲρ, ἐννενηκονταπέντε.

\vs{21}Υἱοὶ Βεθλαὲμ, ἑκατὸν εἰκοσιτρεῖς.

\vs{22}Υἱοὶ Νετωφὰ, πεντηκονταέξ.

\vs{23}Υἱοὶ Ἀναθὼθ, ἑκατὸν εἰκοσιοκτώ.

\vs{24}Υἱοὶ Ἀζμὼθ, τεσσαρακοντατρεῖς.

\vs{25}Υἱοὶ Καριαθιαρὶμ, Χαφιρὰ, καὶ Βηρὼθ, ἑπτακόσιοι τεσσαρακοντατρεῖς.

\vs{26}Υἱοὶ τῆς Ῥαμὰ καὶ Γαβαὰ, ἑξακόσιοι εἰκοσιεῖς.

\vs{27}Ἄνδρες Μαχμὰς, ἑκατὸν εἰκοσιδύο.

\vs{28}Ἄνδρες Βαιθὴλ καὶ Ἀϊὰ, τετρακόσιοι εἰκοσιτρεῖς.

\vs{29}Υἱοὶ Ναβοῦ, πεντηκονταδύο.

\vs{30}Υἱοὶ Μαγεβὶς, ἑκατὸν πεντηκονταέξ.

\vs{31}Υἱοὶ Ἠλαμὰρ, χίλιοι διακόσιοι πεντηκοντατέσσαρες.

\vs{32}Υἱοὶ Ἠλὰμ, τριακόσιοι εἴκοσι.

\vs{33}Υἱοὶ Λοδαδὶ καὶ Ὠνὼ, ἑπτακόσιοι εἰκοσιπέντε.

\vs{34}Υἱοὶ Ἱεριχὼ, τριακόσιοι τεσσαρακονταπέντε.

\vs{35}Υἱοὶ Σεναὰ, τρισχίλιοι ἑξακόσιοι τριάκοντα.

\vs{36}Καὶ οἱ ἱερεῖς υἱοὶ Ἰεδουὰ τῷ οἴκῳ Ἰησοῖ, ἐννακόσιοι ἑβδομηκοντατρεῖς.
\vs{37}Υἱοὶ Ἐμμὴρ, χίλιοι πεντηκονταδύο.
\vs{38}Υἱοὶ Φασσοὺρ, χίλιοι διακόσιοι τεσσαρακονταεπτά.
\vs{39}Υἱοὶ Ἠρὲμ, χίλιοι ἑπτά.

\vs{40}Καὶ οἱ Λευῖται υἱοὶ Ἰησοῦ καὶ Καδμιὴλ τοῖς υἱοῖς Ὠδουΐα, ἑβδομηκοντατέσσαρες.

\vs{41}Οἱ ᾄδοντες υἱοὶ Ἀσὰφ, ἑκατὸν εἰκοσιοκτώ.

\vs{42}Υἱοὶ τῶν πυλωρῶν, υἱοὶ Σελλοὺμ, υἱοὶ Ἀτὴρ, υἱοὶ Τελμὼν, υἱοὶ Ἀκοὺβ, υἱοὶ Ἀτιτὰ, υἱοὶ Σωβαῒ, οἱ πάντες ἑκατὸν τριακονταεννέα.

\vs{43}Οἱ Ναθινὶμ, υἱοὶ Σουθία, υἱοὶ Ἀσουφὰ, υἱοὶ Ταβαὼθ,
\vs{44}υἱοὶ Κάδης, υἱοὶ Σιαὰ, υἱοὶ Φαδὼν,
\vs{45}υἱοὶ Λαβανὼ, υἱοὶ Ἀγαβὰ, υἱοὶ Ἀκοὺβ,
\vs{46}υἱοὶ Ἀγὰβ, υἱοὶ Σελαμὶ, υἱοὶ Ἀνὰν,
\vs{47}υἱοὶ Γεδδὴλ, υἱοὶ Γαὰρ, υἱοὶ Ῥαϊὰ,
\vs{48}υἱοὶ Ῥασὼν, υἱοὶ Νεκωδὰ, υἱοὶ Γαζὲμ,
\vs{49}υἱοὶ Ἀζὼ, υἱοὶ Φασὴ, υἱοὶ Βασὶ,
\vs{50}υἱοὶ Ἀσενὰ, υἱοὶ Μοουνὶμ, υἱοὶ Νεφουσὶμ,
\vs{51}υἱοὶ Βακβοὺκ, υἱοὶ Ἀκουφὰ, υἱοὶ Ἀροὺρ,
\vs{52}υἱοὶ Βασαλὼθ, υἱοὶ Μαουδὰ, υἱοὶ Ἀρσὰ,
\vs{53}υἱοὶ Βαρκὸς, υἱοὶ Σισάρα, υἱοὶ Θεμὰ,
\vs{54}υἱοὶ Νασθιὲ, υἱοὶ Ἀτουφά·
\vs{55}Υἱοὶ δούλων Σαλωμὼν, υἱοὶ Σωταῒ, υἱοὶ Σεφηρὰ, υἱοὶ Φαδουρὰ,
\vs{56}υἱοὶ Ἰεηλὰ, υἱοὶ Δαρκὼν, υἱοὶ Γεδὴλ,
\vs{57}υἱοὶ Σαφατία, υἱοὶ Ἀτὶλ, υἱοὶ Φαχερὰθ, υἱοὶ Ἀσεβωεὶμ, υἱοὶ Ἡμεΐ.
\vs{58}Πάντες οἱ Ναθανὶμ, καὶ υἱοὶ Ἀβδησελμὰ, τριακόσιοι ἐννενηκονταδύο.

\vs{59}Καὶ οὗτοι οἱ ἀναβάντες ἀπὸ Θελμελὲχ, Θελαρησὰ, Χεροὺβ, Ἡδὰν, Ἐμμὴρ· καὶ οὐκ ἐδυνάσθησαν τοῦ ἀναγγεῖλαι οἶκον πατριᾶς αὐτῶν καὶ σπέρμα αὐτῶν, εἰ ἐξ Ἰσραήλ εἰσιν·
\vs{60}Υἱοὶ Δαλαΐα, υἱοὶ Βουὰ, υἱοὶ Τωβίου, υἱοὶ Νεκωδὰ, ἑξακόσιοι πεντηκονταδύο.
\vs{61}Καὶ ἀπὸ τῶν υἱῶν τῶν ἱερέων υἱοὶ Λαβεία, υἱοὶ Ἀκκοὺς, υἱοὶ Βερζελλαῒ, ὃς ἔλαβεν ἀπὸ τῶν θυγατέρων Βερζελλαῒ τοῦ Γαλααδίτου γυναῖκα, καὶ ἐκλήθη ἐπὶ τῷ ὀνόματι αὐτῶν.
\vs{62}Οὗτοι ἐζήτησαν γραφὴν αὐτῶν οἱ μεθωεσὶμ, καὶ οὐχ εὑρέθησαν, καὶ ἠγχιστεύθησαν ἀπὸ τῆς ἱερατείας.
\vs{63}Καὶ εἶπεν ἀθερσασθὰ αὐτοῖς τοῦ μὴ φαγεῖν ἀπὸ τοῦ ἁγίου τῶν ἁγίων, ἕως ἀναστῇ ἱερεὺς τοῖς φωτίζουσι καὶ τοῖς τελείοις.

\vs{64}Πᾶσα δὲ ἡ ἐκκλησία ὁμοῦ ὡσεὶ τέσσαρες μυριάδες δισχίλιοι τριακόσιοι ἑξήκοντα,
\vs{65}χωρὶς δούλων αὐτῶν καὶ παιδισκῶν αὐτῶν, οὗτοι ἑπτακισχίλιοι τριακόσιοι τριακονταεπτά· καὶ οὗτοι ᾄδοντες καὶ ᾄδουσαι διακόσιοι.
\vs{66}Ἵπποι αὐτῶν, ἑπτακόσιοι τριακονταέξ· ἡμίονοι αὐτῶν, διακόσιοι τεσσαρακονταπέντε·
\vs{67}Κάμηλοι αὐτῶν, τετρακόσιοι τριακονταπέντε· ὄνοι αὐτῶν, ἑξακισχίλιοι ἑπτακόσιοι εἴκοσι.

\vs{68}Καὶ ἀπὸ ἀρχόντων πατριῶν ἐν τῷ εἰσελθεῖν αὐτοὺς εἰς οἶκον Κυρίου τὸν ἐν Ἱερουσαλὴμ, ἡκουσιάσαντο εἰς οἶκον τοῦ Θεοῦ, τοῦ στῆσαι αὐτὸν ἐπὶ τὴν ἑτοιμασίαν αὐτοῦ·
\vs{69}ὡς ἡ δύναμις αὐτῶν, ἔδωκαν εἰς θησαυρὸν τοῦ ἔργου χρυσίον καθαρὸν μναὶ ἓξ μυριάδες καὶ χίλιαι, καὶ ἀργυρίου μνὰς πεντακισχιλίας, καὶ κόθωνοι τῶν ἱερέων ἑκατόν.

\vs{70}Καὶ ἐκάθισαν οἱ ἱερεῖς, καὶ οἱ Λευῖται, καὶ οἱ ἀπὸ τοῦ λαοῦ, καὶ οἱ ᾄδοντες, καὶ οἱ πυλωροὶ, καὶ οἱ Ναθινὶμ ἐν πόλεσιν αὐτῶν, καὶ πᾶς Ἰσραὴλ ἐν πόλεσιν αὐτῶν.

\ch{3}
Καὶ ἔφθασεν ὁ μὴν ὁ ἕβδομος, καὶ οἱ υἱοὶ Ἰσραὴλ ἐν πόλεσιν αὐτῶν, καὶ συνήχθη ὁ λαὸς ὡς ἀνὴρ εἷς εἰς Ἱερουσαλήμ.
\vs{2}Καὶ ἀνέστη Ἰησοῦς ὁ τοῦ Ἰωσεδὲκ καὶ οἱ ἀδελφοὶ αὐτοῦ ἱερεῖς, καὶ Ζοροβάβελ ὁ τοῦ Σαλαθιὴλ καὶ οἱ ἀδελφοὶ αὐτοῦ, καὶ ᾠκοδόμησαν τὸ θυσιαστήριον Θεοῦ Ἰσραὴλ, τοῦ ἀνενέγκαι ἐπʼ αὐτὸ ὁλοκαυτώσεις, κατὰ τὰ γεγραμμένα ἐν νόμῳ Μωυσῆ ἀνθρώπου τοῦ Θεοῦ.

\vs{3}Καὶ ἡτοίμασαν τὸ θυσιαστήριον ἐπὶ τὴν ἑτοιμασίαν αὐτοῦ, ὅτι ἐν καταπλήξει ἐπʼ αὐτοὺς ἀπὸ τῶν λαῶν τῶν γαιῶν· καὶ ἀνέβη ἐπʼ αὐτὸ ὁλοκαύτωσις τῷ Κυρίῳ τοπρωῒ καὶ εἰς ἑσπέραν.
\vs{4}Καὶ ἐποίησαν τὴν ἑορτὴν τῶν σκηνῶν κατὰ τὸ γεγραμμένον, καὶ ὁλοκαυτώσεις ἡμέραν ἐν ἡμέρᾳ ἐν ἀριθμῷ ὡς ἡ κρίσις, λόγον ἡμέρας ἐν ἡμέρᾳ αὐτοῦ·
\vs{5}Καὶ μετὰ τοῦτο ὁλοκαυτώσεις ἐνδελεχισμοῦ, καὶ εἰς τὰς νουμηνίας καὶ εἰς πάσας ἑορτὰς τῷ Κυρίῳ τὰς ἡγιασμένας, καὶ παντὶ ἑκουσιαζομένῳ ἑκούσιον τῷ Κυρίῳ.
\vs{6}Ἐν ἡμέρᾳ μιᾷ τοῦ μηνὸς τοῦ ἑβδόμου ἤρξαντο ἀναφέρειν ὁλοκαυτώσεις τῷ Κυρίῳ, καὶ ὁ οἶκος τοῦ Κυρίου οὐκ ἐθεμελιώθη.
\vs{7}Καὶ ἔδωκαν ἀργύριον τοῖς λατόμοις καὶ τοῖς τέκτοσι, καὶ βρώματα καὶ ποτὰ, καὶ ἔλαιον τοῖς Σιδωνίοις καὶ τοῖς Τυρίοις, ἐνέγκαι ξύλα κέδρινα ἀπὸ τοῦ Λιβάνου πρὸς θάλασσαν Ἰόππης, κατʼ ἐπιχώρησιν Κύρου βασιλέως Περσῶν ἐπʼ αὐτούς.

\vs{8}Καὶ ἐν τῷ ἔτει τῷ δευτέρῳ τοῦ ἐλθεῖν αὐτοὺς εἰς οἶκον τοῦ Θεοῦ ἐν Ἱερουσαλὴμ, ἐν μηνὶ τῷ δευτέρῳ ἤρξατο Ζοροβάβελ ὁ τοῦ Σαλαθιὴλ καὶ Ἰησοῦς ὁ τοῦ Ἰωσεδὲκ, καὶ οἱ κατάλοιποι τῶν ἀδελφῶν αὐτῶν οἱ ἱερεῖς καὶ οἱ Λευῖται, καὶ πάντες οἱ ἐρχόμενοι ἀπὸ τῆς αἰχμαλωσίας εἰς Ἱερουσαλὴμ, καὶ ἔστησαν τοὺς Λευίτας ἀπὸ εἰκοσαετοῦς καὶ ἐπάνω ἐπὶ τοὺς ποιοῦντας τὰ ἔργα ἐν οἴκῳ Κυρίου.
\vs{9}Καὶ ἔστη Ἰησοῦς καὶ οἱ υἱοὶ αὐτοῦ καὶ οἱ ἀδελφοὶ αὐτοῦ, Καδμιὴλ καὶ οἱ υἱοὶ αὐτοῦ υἱοὶ Ἰούδα ἐπὶ τοὺς ποιοῦντας τὰ ἔργα ἐν οἴκῳ τοῦ Θεοῦ· υἱοὶ Ἠναδὰδ, υἱοὶ αὐτῶν καὶ οἱ ἀδελφοὶ αὐτῶν οἱ Λευῖται.

\vs{10}Καὶ ἐθεμελίωσαν τοῦ οἰκοδομῆσαι τὸν οἶκον Κυρίου· καὶ ἔστησαν οἱ ἱερεῖς ἐστολισμένοι ἐν σάλπιγξι, καὶ οἱ Λευῖται υἱοὶ Ἀσὰφ ἐν κυμβάλοις τοῦ αἰνεῖν τὸν Κύριον ἐπὶ χεῖρας Δαυὶδ βασιλέως Ἰσραήλ.
\vs{11}Καὶ ἀπεκρίθησαν ἐν αἴνῳ καὶ ἀνθομολογήσει τῷ Κυρίῳ, ὅτι ἀγαθὸν, ὅτι εἰς τὸν αἰῶνα τὸ ἔλεος αὐτοῦ ἐπὶ Ἰσραήλ· καὶ πᾶς ὁ λαὸς ἐσήμαινε φωνῇ μεγάλῃ αἰνεῖν τῷ Κυρίῳ ἐπὶ τῇ θεμελιώσει τοῦ οἴκου Κυρίου.
\vs{12}Καὶ πολλοὶ ἀπὸ τῶν ἱερέων καὶ τῶν Λευιτῶν καὶ ἄρχοντες τῶν πατριῶν οἱ πρεσβύτεροι οἳ εἴδοσαν τὸν οἴκον τὸν πρῶτον ἐν θεμελιώσει αὐτοῦ, καὶ τοῦτον τὸν οἶκον ἐν ὀφθαλμοῖς αὐτῶν, ἔκλαιον φωνῇ μεγάλῃ· καὶ ὁ ὄχλος ἐν σημασίᾳ μετʼ εὐφροσύνης τοῦ ὑψῶσαι ᾠδήν.
\vs{13}Καὶ οὐκ ἦν ὁ λαὸς ἐπιγινώσκων φωνὴν σημασίας τῆς εὐφροσύνης ἀπὸ τῆς φωνῆς τοῦ κλαυθμοῦ τοῦ λαοῦ, ὅτι ὁ λαὸς ἐκραύγασε φωνῇ μεγάλῃ, καὶ ἡ φωνὴ ἠκούετο ἕως ἀπὸ μακρόθεν.

\ch{4}
Καὶ ἤκουσαν οἱ θλίβοντες Ἰούδα καὶ Βενιαμὶν, ὅτι υἱοὶ τῆς ἀποικίας οἰκοδομοῦσιν οἶκον τῷ Κυρίῳ Θεῷ Ἰσραὴλ,
\vs{2}καὶ ἤγγισαν πρὸς Ζοροβάβελ καὶ πρὸς τοὺς ἄρχοντας τῶν πατριῶν, καὶ εἶπον αὐτοῖς, οἰκοδομήσομεν μεθʼ ὑμῶν, ὅτι ὡς ὑμεῖς ἐκζητοῦμεν τῷ Θεῷ ἡμῶν, καὶ αὐτῷ ἡμεῖς θυσιάζομεν ἀπὸ ἡμερῶν Ἀσαραδὰν βασιλέως Ἀσσοὺρ τοῦ ἐνέγκαντος ἡμᾶς ὧδε.

\vs{3}Καὶ εἶπε πρὸς αὐτοὺς Ζοροβάβελ καὶ Ἰησοῦς καὶ οἱ κατάλοιποι τῶν ἀρχόντων τῶν πατριῶν τοῦ Ἰσραὴλ, οὐχ ἡμῖν καὶ ὑμῖν τοῦ οἰκοδομῆσαι οἶκον τῷ Θεῷ ἡμῶν, ὅτι ἡμεῖς ἐπὶ τοαυτὸ οἰκοδομήσομεν τῷ Κυρίω Θεῷ ἡμῶν, ὡς ἐνετείλατο ἡμῖν Κύρος ὁ βασιλεὺς Περσῶν.
\vs{4}Καὶ ἦν ὁ λαὸς τῆς γῆς ἐκλύων τὰς χεῖρας τοῦ λαοῦ Ἰούδα, καὶ ἐνεπόδιζον αὐτοὺς οἰκοδομεῖν,
\vs{5}καὶ μισθούμενοι ἐπʼ αὐτοὺς βουλευόμενοι τοῦ διασκεδάσαι βουλὴν αὐτῶν πάσας τὰς ἡμέρας Κύρου βασιλέως Περσῶν, καὶ ἕως βασιλείας Δαρείου βασιλείας Περσῶν.

\vs{6}Καὶ ἐν βασιλείᾳ Ἀσσουήρου, καὶ ἐν ἀρχῇ βασιλείας αὐτοῦ ἔγραψαν ἐπιστολὴν ἐπὶ οἰκοῦντας Ἰούδα καὶ Ἱερουσαλήμ.
\vs{7}Καὶ ἐν ἡμέραις Ἀρθασασθὰ ἔγραψεν ἐν εἰρήνῃ Μιθραδάτῃ Ταβεὴλ καὶ τοῖς λοιποῖς συνδούλοις· πρὸς Ἀρθασασθὰ βασιλέα Περσῶν ἔγραψεν ὁ φορολόγος γραφὴν Συριστὶ καὶ ἡρμηνευμένην.
\vs{8}Ῥεοὺμ βαλτὰμ καὶ Σαμψὰ ὁ γραμματεὺς ἔγραψαν ἐπιστολὴν μίαν κατὰ Ἱερουσαλὴμ τῷ Ἀρθασασθὰ βασιλεῖ.
\vs{9}Τάδε ἔκρινε Ῥεοὺμ βαλτὰμ καὶ Σαμψὰ ὁ γραμματεὺς καὶ οἱ κατάλοιποι σύνδουλοι ἡμῶν, Δειναῖοι, Ἀφαρσαθαχαῖοι, Ταρφαλαῖοι, Ἀφαρσαῖοι, Ἀρχυαῖοι, Βαβυλώνιοι, Σουσαναχαῖοι, Δαυαῖοι,
\vs{10}καὶ οἱ κατάλοιποι ἐθνῶν ὧν ἀπῴκισεν Ἀσσεναφὰρ ὁ μέγας καὶ ὁ τίμιος, καὶ κατῴκισεν αὐτοὺς ἐν πόλεσι τῆς Σομόρων καὶ τὸ κατάλοιπον πέραν τοῦ ποταμοῦ.
\vs{11}Αὕτη ἡ διαταγὴ τῆς ἐπιστολῆς, ἧς ἀπέστειλαν πρὸς αὐτόν· πρὸς Ἀρθασασθὰ βασιλέα παῖδές σου ἄνδρες πέραν τοῦ ποταμοῦ.

\vs{12}Γνωστὸν ἔστω τῷ βασιλεῖ, ὅτι οἱ Ἰουδαῖοι οἱ ἀναβάντες ἀπὸ σοῦ πρὸς ἡμᾶς ἤλθοσαν εἰς Ἱερουσαλὴμ τὴν πόλιν τὴν ἀποστάτιν καὶ πονηρὰν, ἣν οἰκοδομοῦσι· καὶ τὰ τείχη αὐτῆς κατηρτισμένα εἰσὶ, καὶ θεμελίους αὐτῆς ἀνύψωσαν.
\vs{13}Νῦν οὖν γνωστὸν ἔστω τῷ βασιλεῖ, ὅτι ἐὰν ἡ πόλις ἐκείνη ἀνοικοδομηθῇ, καὶ τὰ τείχη αὐτῆς καταρτισθῶσι, φόροι οὐκ ἔσονταί σοι, οὐδὲ δώσουσι· καὶ τοῦτο βασιλεῖς κακοποιεῖ,
\vs{14}καὶ ἀσχημοσύνην βασιλέως οὐκ ἔξεστιν ἡμῖν ἰδεῖν· διὰ τοῦτο ἐπέμψαμεν καὶ ἐγνωρίσαμεν τῷ βασιλεῖ,
\vs{15}ἵνα ἐπισκέψηται ἐν βίβλῳ ὑπομνηματισμοῦ τῶν πατέρων σου, καὶ εὑρήσεις, καὶ γνώσῃ, ὅτι ἡ πόλις ἐκείνη πόλις ἀποστάτις, καὶ κακοποιοῦσα βασιλεῖς καὶ χώρας, καὶ φυγαδείαι δούλων γίνονται ἐν μέσῳ αὐτῆς ἀπὸ ἡμερῶν αἰῶνος, διὰ ταῦτα ἡ πόλις αὕτη ἠρημώθη.
\vs{16}Γνωρίζομεν οὖν ἡμεῖς τῷ βασιλεῖ, ὅτι ἂν ἡ πόλις ἐκείνη οἰκοδομηθῇ, καὶ τὰ τείχη αὐτῆς καταρτισθῇ, οὐκ ἔστι σοι εἰρήνη.

\vs{17}Καὶ ἀπέστειλεν ὁ βασιλεὺς πρὸς Ῥεοὺμ βαλτὰμ καὶ Σαμψὰ γραμματέα καὶ τοὺς καταλοίπους συνδούλους αὐτῶν τοὺς οἰκοῦντας ἐν Σαμαρείᾳ καὶ τοὺς καταλοίπους πέραν τοῦ ποταμοῦ, εἰρήνην· καὶ φησὶν,
\vs{18}ὁ φορολόγος ὃν ἀπεστείλατε πρὸς ἡμᾶς, ἐκλήθη ἔμπροσθεν ἐμοῦ·
\vs{19}Καὶ παρʼ ἐμοῦ ἐτέθη γνώμη, καὶ ἐπεσκεψάμεθα καὶ εὕραμεν, ὅτι ἡ πόλις ἐκείνη ἀφʼ ἡμερῶν αἰῶνος ἐπὶ βασιλεῖς ἐπαίρεται, καὶ ἀποστάσεις καὶ φυγαδείαι γίνονται ἐν αὐτῇ·
\vs{20}Καὶ βασιλεῖς ἰσχυροὶ ἐγένοντο ἐν Ἱερουσαλὴμ, καὶ ἐπικρατοῦντες ὅλης τῆς πέραν τοῦ ποταμοῦ, καὶ φόροι πλήρεις καὶ μέρος δίδονται αὐτοῖς.
\vs{21}Καὶ νῦν θέτε γνώμην, καταργῆσαι τοὺς ἄνδρας ἐκείνους, καὶ ἡ πόλις ἐκείνη οὐκ οἰκοδομηθήσεται ἔτι· Ὅπως ἀπὸ τῆς γνώμης
\vs{22}πεφυλαγμένοι ἦτε ἄνεσιν ποιῆσαι περὶ τούτου, μή ποτε πληθυνθῇ ἀφανισμὸς εἰς κακοποίησιν βασιλεῦσι.

\vs{23}Τότε ὁ φορολόγος τοῦ Ἀρθασασθὰ βασιλέως ἀνέγνω ἐνώπιον Ῥεοὺμ βαλτὰμ καὶ Σαμψὰ γραμματέως καὶ συνδούλων αὐτοῦ· καὶ ἐπορεύθησαν σπουδῇ εἶς Ἱερουσαλὴμ καὶ ἐν Ἰούδα, καὶ κατήργησαν αὐτοὺς ἐν ἵπποις καὶ δυνάμει.
\vs{24}Τότε ἤργησε τὸ ἔργον οἴκου τοῦ Θεοῦ τὸ ἐν Ἱερουσαλήμ· καὶ ἦν ἀργοῦν ἕως δευτέρου ἔτους τῆς βασιλείας Δαρείου βασιλέως Περσῶν.

\ch{5}
Καὶ προεφήτευσεν Ἀγγαῖος ὁ προφήτης καὶ Ζαχαρίας ὁ τοῦ Ἀδδὼ προφητείαν ἐπὶ τοὺς Ἰουδαίους τοὺς ἐν Ἰούδα καὶ Ἱερουσαλὴμ ἐν ὀνόματι Θεοῦ Ἰσραὴλ ἐπʼ αὐτούς.
\vs{2}Τότε ἀνέστησαν Ζοροβάβελ ὁ τοῦ Σαλαθιὴλ καὶ Ἰησοῦς υἱὸς Ἰωσεδὲκ, καὶ ἤρξαντο οἰκοδομῆσαι τὸν οἶκον τοῦ Θεοῦ τὸν ἐν Ἱερουσαλήμ· καὶ μετʼ αὐτῶν οἱ προφῆται τοῦ Θεοῦ βοηθοῦντες αὐτοῖς.

\vs{3}Ἐν αὐτῷ τῷ καιρῷ ἦλθεν ἐπʼ αὐτοὺς Θανθαναῒ ἔπαρχος πέραν τοῦ ποταμοῦ, καὶ Σαθαρβουζαναῒ, καὶ οἱ σύνδουλοι αὐτῶν, καὶ τοιάδε εἶπαν αὐτοῖς, τίς ἔθηκεν ὑμῖν γνώμην τοῦ οἰκοδομῆσαι τὸν οἶκον τοῦτον, καὶ τὴν χορηγίαν ταύτην καταρτίσασθαι;
\vs{4}Τότε ταῦτα εἴποσαν αὐτοῖς, τίνα ἐστὶ τὰ ὀνόματα τῶν ἀνδρῶν τῶν οἰκοδομούντων τὴν πόλιν ταύτην;
\vs{5}Καὶ οἱ ὀφθαλμοὶ τοῦ Θεοῦ ἐπὶ τὴν αἰχμαλωσίαν Ἰούδα, καὶ οὐ κατήργησαν αὐτοὺς ἕως γνώμη τῷ Δαρείῳ ἀπηνέχθη· καὶ τότε ἀπεστάλη τῷ φορολόγῳ ὑπὲρ τούτου
\vs{6}διασάφησις ἐπιστολῆς, ἧς ἀπέστειλε Θανθαναῒ, ὁ ἔπαρχος τοῦ πέραν τοῦ ποταμοῦ, καὶ Σαθαρβουζαναῒ καὶ οἱ σύνδουλοι αὐτῶν Ἀφαρσαχαῖοι οἱ ἐν τῷ πέραν τοῦ ποταμοῦ, Δαρείῳ τῷ βασιλεῖ·
\vs{7}Ῥήμασιν ἀπέστειλαν πρὸς αὐτόν· καὶ τάδε γέγραπται ἐν αὐτῷ·

\vs{8}Δαρείῳ τῷ βασιλεῖ εἰρήνη πᾶσα. Γνωστὸν ἔστω τῷ βασιλεῖ, ὅτι ἐπορεύθημεν εἰς τὴν Ἰουδαίαν χώραν εἰς οἶκον τοῦ Θεοῦ τοῦ μεγάλου, καὶ αὐτὸς οἰκοδομεῖται λίθοις ἐκλεκτοῖς, καὶ ξύλα ἐντίθεται ἐν τοῖς τοίχοις, καὶ τὸ ἔργον ἐκεῖνο ἐπιδέξιον γίνεται, καὶ εὐοδοῦται ἐν ταῖς χερσὶν αὐτῶν.
\vs{9}Τότε ἠρωτήσαμεν τοὺς πρεσβυτέρους ἐκείνους, καὶ οὕτως εἴπαμεν αὐτοῖς, τίς ἔθηκεν ὑμῖν γνώμην τὸν οἶκον τοῦτον οἰκοδομῆσαι, καὶ τὴν χορηγίαν ταύτην καταρτίσασθαι;
\vs{10}Καὶ τὰ ὀνόματα αὐτῶν ἠρωτήσαμεν αὐτοὺς γνωρίσαι σοι, ὥστε γράψαι σοι τὰ ὀνόματα τῶν ἀνδρῶν τῶν ἀρχόντων αὐτῶν.
\vs{11}Καὶ τοιοῦτο τὸ ῥῆμα ἀπεκρίθησαν ἡμῖν, λέγοντες, ἡμεῖς ἐσμὲν δοῦλοι τοῦ Θεοῦ τοῦ οὐρανοῦ καὶ τῆς γῆς, καὶ οἰκοδομοῦμεν τὸν οἶκον ὃς ἦν ᾠκοδομημένος πρὸ τούτου ἔτη πολλὰ, καὶ βασιλεὺς τοῦ Ἰσραὴλ μέγας ᾠκοδόμησεν αὐτὸν, καὶ κατηρτίσατο αὐτὸν αὐτοῖς.
\vs{12}Ἀφότε δὲ παρώργισαν οἱ πατέρες ἡμῶν τὸν Θεὸν τοῦ οὐρανοῦ, ἔδωκεν αὐτοὺς εἰς χεῖρας Ναβουχοδονόσορ βασιλέως Βαβυλῶνος τοῦ Χαλδαίου, καὶ τὸν οἴκον τοῦτον κατέλυσε, καὶ τὸν λαὸν ἀπῴκισεν εἰς Βαβυλῶνα.
\vs{13}Αλλʼ ἐν ἔτει πρώτῳ Κύρου τοῦ βασιλέως, Κύρος ὁ βασιλεὺς ἔθετο γνώμην τὸν οἶκον τοῦ Θεοῦ τοῦτον οἰκοδομηθῆναι·
\vs{14}Καὶ τὰ σκεύη τοῦ οἴκου τοῦ Θεοῦ τὰ χρυσᾶ καὶ τὰ ἀργυρᾶ, ἃ Ναβουχοδονόσορ ἐξήνεγκεν ἀπὸ τοῦ οἴκου τοῦ ἐν Ἱερουσαλὴμ, καὶ ἀπήνεγκεν αὐτὰ εἰς τὸν ναὸν του βασιλέως, ἐξήνεγκεν αὐτὰ Κύρος ὁ βασιλεὺς ἀπὸ τοῦ ναοῦ τοῦ βασιλέως, καὶ ἔδωκε τῷ Σαβανασὰρ τῷ θησαυροφύλακι, τῷ ἐπὶ τοῦ θησαυροῦ,
\vs{15}καὶ εἶπεν αὐτῷ, πάντα τὰ σκεύη λάβε καὶ πορεύου, θὲς αὐτὰ ἐν τῷ οἴκῳ τῷ ἐν Ἱερουσαλὴμ εἰς τὸν τόπον αὐτῶν.
\vs{16}Τότε Σαβανασὰρ ἐκεῖνος ἦλθε καὶ ἔδωκε θεμελίους τοῦ οἴκου τοῦ Θεοῦ ἐν Ἱερουσαλὴμ, καὶ ἀπὸ τότε ἕως τοῦ νῦν ᾠκοδομήθη, καὶ οὐκ ἐτελέσθη.
\vs{17}Καὶ νῦν εἰ ἐπὶ τὸν βασιλέα ἀγαθὸν ἐπισκεπήτω ἐν τῷ οἴκῳ τῆς γάζης τοῦ βασιλέως Βαβυλῶνος, ὅπως γνῷς ὅτι ἀπὸ βασιλέως Κύρου ἐτέθη γνώμη οἰκοδομῆσαι τὸν οἶκον τοῦ Θεοῦ ἐκεῖνον τὸν ἐν Ἱερουσαλήμ· καὶ γνοὺς ὁ βασιλεὺς περὶ τούτου, πεμψάτω πρὸς ἡμᾶς.

\ch{6}
Τότε Δαρεῖος ὁ βασιλεὺς ἔθηκε γνώμην, καὶ ἐπεσκέψατο ἐν ταῖς βιβλιοθήκαις ὅπου ἡ γάζα κεῖται ἐν Βαβυλῶνι.
\vs{2}Καὶ εὑρέθη ἐν πόλει ἐν τῇ βάρει κεφαλὶς μία, καὶ τοῦτο γεγραμμένον ἐν αὐτῇ ὑπόμνημα.

\vs{3}Ἐν ἔτει πρώτῳ Κύρου βασιλέως, Κύρος ὁ βασιλεὺς ἔθηκε γνώμην περὶ οἴκου ἱεροῦ Θεοῦ τοῦ ἐν Ἱερουσαλήμ· οἶκος οἰκοδομηθήτω, καὶ τόπος οὗ θυσιάζουσι τὰ θυσιάσματα· καὶ ἔθηκεν ἔπαρμα ὕψος πήχεις ἑξήκοντα, πλάτος αὐτοῦ πήχεων ἑξήκοντα·
\vs{4}Καὶ δόμοι λίθινοι κραταιοὶ τρεῖς, καὶ δόμος ξύλινος εἷς, καὶ ἡ δαπάνη ἐξ οἴκου τοῦ βασιλέως δοθήσεται.
\vs{5}Καὶ τὰ σκεύη οἴκου τοῦ Θεοῦ τὰ ἀργυρᾶ καὶ τὰ χρυσᾶ, ἃ Ναβουχοδονόσορ ἐξήνεγκεν ἀπὸ τοῦ οἴκου τοῦ ἐν Ἱερουσαλὴμ, καὶ ἐκόμισεν εἰς Βαβυλῶνα, καὶ δοθήτω καὶ ἀπελθέτω εἰς τὸν ναὸν τὸν ἐν Ἱερουσαλὴμ ἐπὶ τόπου οὗ ἐτέθη ἐν οἴκῳ τοῦ Θεοῦ.

\vs{6}Νῦν δώσετε ἔπαρχοι πέραν τοῦ ποταμοῦ Σαθαρβουζαναῒ, καὶ οἱ σύνδουλοι αὐτῶν Ἀφαρσαχαῖοι οἱ ἐν τῷ πέραν τοῦ ποταμοῦ μακρὰν ὄντες ἐκεῖθεν,
\vs{7}νῦν ἄφετε τὸ ἔργον οἴκου τοῦ Θεοῦ· οἱ ἀφηγούμενοι τῶν Ἰουδαίων καὶ οἱ πρεσβύτεροι τῶν Ἰουδαίων οἶκον τοῦ Θεοῦ ἐκεῖνον οἰκοδομείτωσαν ἐπὶ τοῦ τόπου αὐτοῦ.
\vs{8}Καὶ ἀπʼ ἐμοῦ γνώμη ἐτέθη, μή ποτε τὶ ποιήσητε μετὰ τῶν πρεσβυτέρων τῶν Ἰουδαίων τοῦ οἰκοδομηθῆναι οἶκον τοῦ Θεοῦ ἐκεῖνον· καὶ ἀπὸ ὑπαρχόντων βασιλέως τῶν φόρων πέραν τοῦ ποταμοῦ ἐπιμελῶς δαπάνη ἔστω διδομένη τοῖς ἀνδράσιν ἐκείνοις τὸ μὴ καταργηθῆναι.
\vs{9}Καὶ ὃ ἂν ὑστέρημα, καὶ υἱοὺς βοῶν καὶ κριῶν, καὶ ἀμνοὺς εἰς ὁλοκαυτώσεις τῷ Θεῷ τοῦ οὐρανοῦ, πυροὺς, ἅλας, οἶνον, ἔλαιον, κατὰ τὸ ῥῆμα ἱερέων τῶν ἐν Ἱερουσαλὴμ ἔστω διδόμενον αὐτοῖς, ἡμέραν ἐν ἡμέρᾳ, ὃ ἐὰν αἰτήσωσιν,
\vs{10}ἵνα ὦσιν εὐωδίας προσφέροντες τῷ Θεῷ τοῦ οὐρανοῦ, καὶ προσεύχωνται εἰς ζωὴν τοῦ βασιλέως καὶ υἱῶν αὐτοῦ.
\vs{11}Καὶ ἀπʼ ἐμοῦ ἐτέθη γνώμη, ὅτι πᾶς ἄνθρωπος ὃς ἀλλάξει τὸ ῥῆμα τοῦτο, καθαιρεθήσεται ξύλον ἐκ τῆς οἰκίας αὐτοῦ, καὶ ὠρθωμένος πληγήσεται ἐπʼ αὐτοῦ, καὶ ὁ οἶκος αὐτοῦ τὸ κατʼ ἐμὲ ποιηθήσεται.
\vs{12}Καὶ ὁ Θεὸς οὗ κατασκηνοῖ τὸ ὄνομα ἐκεῖ, καταστρέψαι πάντα βασιλέα καὶ λαὸν ὃς ἐκτενεῖ τὴν χεῖρα αὐτοῦ ἀλλάξαι ἢ ἀφανίσαι τὸν οἶκον τοῦ Θεοῦ τὸν ἐν Ἱερουσαλήμ· ἐγὼ Δαρεῖος ἔθηκα γνώμην, ἐπιμελῶς ἔσται.

\vs{13}Τότε Θανθαναῒ ὁ ἔπαρχος πέραν τοῦ ποταμοῦ, Σαθαρβουζαναῒ, καὶ οἱ σύνδουλοι αὐτοῦ, πρὸς ὃ ἀπέστειλε Δαρεῖος βασιλεὺς, οὕτως ἐποίησαν ἐπιμελῶς.
\vs{14}Καὶ οἱ πρεσβύτεροι τῶν Ἰουδαίων ᾠκοδομοῦσαν καὶ οἱ Λευῖται ἐν προφητείᾳ Ἀγγαίου τοῦ προφήτου, καὶ Ζαχαρίου υἱοῦ Ἀδδώ· καὶ ἀνῳκοδόμησαν καὶ κατηρτίσαντο ἀπὸ γνώμης Θεοῦ Ἰσραὴλ, καὶ ἀπὸ γνώμης Κύρου, καὶ Δαρείου, καὶ Ἀρθασασθὰ βασιλέων Περσῶν.

\vs{15}Καὶ ἐτέλεσαν τὸν οἶκον τοῦτον ἕως ἡμέρας τρίτης μηνὸς Ἀδὰρ, ὅ ἐστιν ἔτος ἕκτον τῆς βασιλείας Δαρείου τοῦ βασιλέως.

\vs{16}Καὶ ἐποίησαν οἱ υἱοὶ Ἰσραὴλ, οἱ ἱερεῖς καὶ οἱ Λευῖται, καὶ οἱ κατάλοιποι υἱῶν ἀποικεσίας ἐγκαίνια τοῦ οἴκου τοῦ Θεοῦ ἐν εὐφροσύνῃ.
\vs{17}Καὶ προσήνεγκαν εἰς τὰ ἐγκαίνια τοῦ οἴκου τοῦ Θεοῦ μόσχους ἑκατὸν, κριοὺς διακοσίους, ἀμνοὺς τετρακοσίους, χιμάῤῥους αἰγῶν ὑπὲρ ἁμαρτίας ὑπὲρ παντὸς Ἰσραὴλ δώδεκα εἰς ἀριθμὸν φυλῶν Ἰσραήλ.
\vs{18}Καὶ ἔστησαν τοὺς ἱερεῖς ἐν διαιρέσεσιν αὐτῶν, καὶ τοὺς Λευίτας ἐν μερισμοῖς αὐτῶν, ἐπὶ δουλείας Θεοῦ ἐν Ἱερουσαλὴμ, κατὰ τὴν γραφὴν βίβλου Μωυσῆ.

\vs{19}Καὶ ἐποίησαν οἱ υἱοὶ τῆς ἀποικεσίας τὸ πάσχα τῇ τεσσαρεσκαιδεκάτῃ τοῦ μηνὸς τοῦ πρώτου.
\vs{20}Ὅτι ἐκαθαρίσθησαν οἱ ἱερεῖς καὶ Λευῖται, ἕως εἷς πάντες καθαροί· καὶ ἔσφαξαν τὸ πάσχα τοῖς πᾶσιν υἱοῖς τῆς ἀποικεσίας καὶ τοῖς ἀδελφοῖς αὐτῶν τοῖς ἱερεῦσι καὶ ἑαυτοῖς.
\vs{21}Καὶ ἔφαγον υἱοὶ Ἰσραὴλ τὸ πάσχα, οἱ ἀπὸ τῆς ἀποικεσίας, καὶ πᾶς ὁ χωριζόμενος τῆς ἀκαθαρσίας ἐθνῶν τῆς γῆς πρὸς αὐτοῦς, τοῦ ἐκζητῆσαι Κύριον Θεὸν Ἰσραήλ.
\vs{22}Καὶ ἐποίησαν τὴν ἑορτὴν τῶν ἀζύμων ἑπτὰ ἡμέρας ἐν εὐφροσύνῃ, ὅτι εὔφρανεν αὐτοὺς Κύριος, καὶ ἐπέστρεψε καρδίαν βασιλέως Ἀσσοὺρ ἐπʼ αὐτοὺς κραταιῶσαι τὰς χεῖρας αὐτῶν ἐν ἔργοις οἴκου τοῦ Θεοῦ Ἰσραήλ.

\ch{7}
Καὶ μετὰ τὰ ῥήματα ταῦτα ἐν βασιλείᾳ Ἀρθασασθὰ βασιλέως Περσῶν, ἀνέβη Ἔσδρας υἱὸς Σαραίου, υἱοῦ Ἀζαρίου, υἱοῦ Χελκία,
\vs{2}υἱοῦ Σελοὺμ, υἱοῦ Σαδδοὺκ, υἱοῦ Ἀχιτὼβ,
\vs{3}υἱοῦ Σαμαρία, υἱοῦ Ἐσριὰ, υἱοῦ Μαρεὼθ,
\vs{4}υἱοῦ Ζαραΐα, υἱοῦ Ὀζίου, υἱοῦ Βοκκὶ,
\vs{5}υἱοῦ Ἀβισουὲ, υἱοῦ Φινεὲς, υἱοῦ Ἐλεάζαρ, υἱοῦ Ἀαρὼν τοῦ ἱερέως τοῦ πρώτου·
\vs{6}Αὐτὸς Ἔσδρας ἀνέβη ἐκ Βαβυλῶνος, καὶ αὐτὸς γραμματεὺς ταχὺς ἐν νόμῳ Μωυσῆ, ὃν ἔδωκε Κύριος ὁ Θεὸς Ἰσραήλ· καὶ ἔδωκεν αὐτῷ ὁ βασιλεὺς, ὅτι χεὶρ Κυρίου Θεοῦ αὐτοῦ ἐπʼ αὐτὸν ἐν πᾶσιν οἷς ἐζήτει αὐτός.
\vs{7}Καὶ ἀνέβησαν ἀπὸ τῶν υἱῶν Ἰσραὴλ, καὶ ἀπὸ τῶν ἱερέων, καὶ ἀπὸ τῶν Λευιτῶν. καὶ οἱ ᾄδοντες, καὶ οἱ πυλωροὶ, καὶ οἱ Ναθινὶμ, εἰς Ἱερουσαλὴμ ἐν ἔτει ἑβδόμῳ τῷ Ἀρθασασθὰ τῷ βασιλεῖ.
\vs{8}Καὶ ἤλθοσαν εἰς Ἱερουσαλὴμ τῷ μηνὶ τῷ πέμπτῳ, τοῦτο τὸ ἔτος ἕβδομον τῷ βασιλεῖ·
\vs{9}Ὅτι ἐν μιᾷ τοῦ μηνὸς τοῦ πρώτου αὐτὸς ἐθεμελίωσε τὴν ἀνάβασιν τὴν ἀπὸ Βαβυλῶνος· ἐν δὲ τῇ πρώτῃ τοῦ μηνὸς τοῦ πέμπτου ἤλθοσαν εἰς Ἱερουσαλὴμ, ὅτι χεὶρ Θεοῦ αὐτοῦ ἦν ἀγαθὴ ἐπʼ αὐτόν.
\vs{10}ὅτι Ἔσδρας ἔδωκεν ἐν καρδίᾳ αὐτοῦ ζητῆσαι τὸν νόμον, καὶ ποιεῖν καὶ διδάσκειν ἐν Ἰσραὴλ προστάγματα καὶ κρίματα.

\vs{11}Καὶ αὕτη ἡ διασάφησις τοῦ διατάγματος, οὗ ἔδωκεν Ἀρθασασθὰ τῷ Ἔσδρα τῷ ἱερεῖ τῷ γραμματεῖ βιβλίου λόγων ἐντολῶν Κυρίου καὶ προσταγμάτων αὐτοῦ ἐπὶ τὸν Ἰσραήλ.

\vs{12}Ἀρθασασθὰ βασιλεὺς βασιλέων Ἔσδρᾳ γραμματεῖ νόμου Κυρίου τοῦ Θεοῦ τοῦ οὐρανοῦ· τετελέσθω λόγος καὶ ἡ ἀπόκρισις.
\vs{13}Ἀπʼ ἐμοῦ ἐτέθη γνώμη, ὅτι πᾶς ὁ ἑχουσιαζόμενος ἐν βασιλείᾳ μου ἀπὸ λαοῦ Ἰσραὴλ καὶ ἱερέων καὶ Λευιτῶν πορευθῆναι εἰς Ἱερουσαλὴμ, μετὰ σοῦ πορευθῆναι.
\vs{14}Ἀπὸ προσώπου τοῦ βασιλέως καὶ τῶν ἑπτὰ συμβούλων ἀπεστάλη ἐπισκέψασθαι ἐπὶ τὴν Ἰουδαίαν καὶ εἰς Ἱερουσαλὴμ νόμῳ Θεοῦ αὐτῶν τῷ ἐν χειρί σου·
\vs{15}Καὶ εἰς οἶκον Κυρίου, ἀργύριον καὶ χρυσίον, ὃ ὁ βασιλεὺς καὶ οἱ σύμβουλοι ἑκουσιάσθησαν τῷ Θεῷ τοῦ Ἰσραὴλ τῷ ἐν Ἱερουσαλὴμ κατασκηνοῦντι.
\vs{16}Καὶ πᾶν ἀργύριον καὶ χρυσίον, ὅ, τι ἐὰν εὕρῃς ἐν πάσῃ χώρᾳ Βαβυλῶνος μετὰ ἑκουσιασμοῦ τοῦ λαοῦ, καὶ ἱερέων τῶν ἑκουσιαζομένων εἰς οἶκον Θεοῦ τὸν ἐν Ἱερουσαλήμ.
\vs{17}Καὶ πάντα προσπορευόμενον τοῦτον ἑτοίμως ἔνταξον ἐν βιβλίῳ τούτῳ, μόσχους, κριοὺς, ἀμνοὺς, καὶ θυσίας αὐτῶν, καὶ σπονδὰς αὐτῶν· καὶ προσοίσεις αὐτὰ ἐπὶ τοῦ θυσιαστηρίου τοῦ οἴκου τοῦ Θεοῦ ὑμῶν τοῦ ἐν Ἱερουσαλήμ.
\vs{18}Καὶ εἴ τι ἐπὶ σὲ καὶ τοὺς ἀδελφούς σου ἀγαθυνθῇ ἐν καταλοίπῳ τοῦ ἀργυρίου καὶ τοῦ χρυσίου ποιῆσαι, ὡς ἀρεστὸν τῷ Θεῷ ὑμῶν ποιήσατε.
\vs{19}Καὶ τὰ σκεύη τὰ διδόμενά σοι εἰς λειτουργίαν οἴκου Θεοῦ, παράδος ἐνώπιον τοῦ Θεοῦ ἐν Ἱερουσαλήμ.
\vs{20}Καὶ κατάλοιπον χρείας οἴκου Θεοῦ σου, ὃ ἂν φανῇ σοι δοῦναι, δώσεις ἀπὸ οἴκων γάζης βασιλέως καὶ ἀπʼ ἐμοῦ.

\vs{21}Ἐγὼ Ἀρθασασθὰ βασιλεὺς ἔθηκα γνώμην πάσαις ταῖς γάζαις ταῖς ἐν πέρα τοῦ ποταμοῦ, ὅτι πᾶν ὃ ἂν αἰτήσῃ ὑμᾶς Ἔσδρας ὁ ἱερεὺς καὶ γραμματεὺς τοῦ Θεοῦ τοῦ οὐρανοῦ, ἑτοίμως γινέσθω·
\vs{22}ἕως ἀργυρίου ταλάντων ἑκατὸν, καὶ ἕως πυροῦ κόρων ἑκατὸν, καὶ ἕως οἶνου βατῶν ἑκατὸν, καὶ ἕως ἐλαίου βατῶν ἑκατὸν, καὶ ἅλας οὗ οὐκ ἔστι γραφή.
\vs{23}Πᾶν ὅ ἐστιν ἐν γνώμῃ Θεοῦ τοῦ οὐρανοῦ, γινέσθω· προσέχετε μήτις ἐπιχειρήσῃ εἰς τὸν οἶκον Θεοῦ τοῦ οὐρανοῦ, μή ποτε γένηται ὀργὴ ἐπὶ τὴν βασιλείαν τοῦ βασιλέως καὶ τῶν υἱῶν αὐτοῦ.
\vs{24}Καὶ ὑμῖν ἐγνώρισται ἐν πᾶσι τοῖς ἱερεῦσι, καὶ τοῖς Λευίταις, ᾄδουσι, πυλωροῖς, Ναθινὶμ, καὶ λειτουργοῖς οἴκου Θεοῦ τοῦτο, φόρος μὴ ἔστω σοι, οὐκ ἐξουσιάσεις καταδουλοῦσθαι αὐτούς.
\vs{25}Καὶ σὺ Ἔσδρα, ὡς ἡ σοφία τοῦ Θεοῦ ἐν χειρί σου, κατάστησον γραμματεῖς καὶ κριτὰς, ἵνα ὦσι κρίνοντες παντὶ τῷ λαῷ τῷ ἐν πέρα τοῦ ποταμοῦ πᾶσι τοῖς εἰδόσι νόμον τοῦ Θεοῦ σου, καὶ τῷ μὴ εἰδότι γνωριεῖτε.

\vs{26}Καὶ πὰς ὃς ἂν μὴ ᾖ ποιῶν νόμον τοῦ Θεοῦ καὶ νόμον τοῦ βασιλέως ἑτοίμως, τὸ κρίμα ἔσται γινόμενον ἐξ αὐτοῦ, ἐάν τε εἰς θάνατον, ἐάν τε εἰς παιδείαν, ἐάν τε εἰς ζημίαν τοῦ βίου. ἐάν τε εἰς παράδοσιν.

\vs{27}Εὐλογητὸς Κύριος ὁ Θεὸς τῶν πατέρων ἡμῶν, ὃς ἔδωκεν ἐν καρδίᾳ τοῦ βασιλέως οὕτως, τοῦ δοξάσαι τὸν οἶκον Κυρίου τὸν ἐν Ἱερουσαλὴμ,
\vs{28}καὶ ἐπʼ ἐμὲ ἔκλινεν ἔλεος ἐν ὀφθαλμοῖς τοῦ βασιλέως καὶ τῶν συμβούλων αὐτοῦ, καὶ πάντων τῶν ἀρχόντων τοῦ βασιλέως, τῶν ἐπῃρμένων· καὶ ἐγὼ ἐκραταιώθην ὡς χεὶρ Θεοῦ ἡ ἀγαθὴ ἐπʼ ἐμέ, καὶ συνῆξα ἀπὸ Ἰσραὴλ ἄρχοντας ἀναβῆναι μετʼ ἐμοῦ.

\ch{8}
Καὶ οὗτοι οἱ ἄρχοντες πατριῶν αὐτῶν οἱ ὁδηγοὶ ἀναβαίνοντες μετʼ ἐμοῦ ἐν βασιλείᾳ Ἀρθασασθὰ τοῦ βασιλέως Βαβυλῶνος.
\vs{2}Ἀπὸ υἱῶν Φινεὲς, Γηρσών· ἀπὸ υἱῶν Ἰθάμαρ, Δανιήλ· ἀπὸ υἱῶν Δαυὶδ, Ἀττούς.
\vs{3}Ἀπὸ υἱῶν Σαχανία, καὶ ἀπὸ υἱῶν Φόρος, Ζαχαρίας, καὶ μετ αὐτοῦ τὸ συστρεμμα ἑκατὸν καὶ πεντήκοντα.
\vs{4}Ἀπὸ υἱῶν Φαὰθ Μωὰβ, Ἐλιανὰ υἱὸς Σαραΐα, καὶ μετʼ αὐτοῦ διακόσιοι τὰ ἀρσενικά.
\vs{5}Καὶ ἀπὸ υἱῶν Ζαθόης, Σεχενίας υἱὸς Ἀζιὴλ, καὶ μετʼ αὐτοῦ τριακόσια τὰ ἀρσενικά.
\vs{6}Καὶ ἀπὸ τῶν υἱῶν Ἀδὶν, Ὠβὴθ υἱὸς Ἰωνάθαν, καὶ μετʼ αὐτοῦ πεντήκοντα τὰ ἀρσενικά.
\vs{7}Καὶ ἀπὸ υἱῶν Ἠλὰμ, Ἰσαΐας υἱὸς Ἀθελία, καὶ μετʼ αὐτοῦ ἑβδομήκοντα τὰ ἀρσενικά.
\vs{8}Καὶ ἀπὸ υἱῶν Σαφατία, Ζαβαδίας υἱὸς Μιχαὴλ, καὶ μετʼ αὐτοῦ ὀγδοήκοντα τὰ ἀρσενικά.
\vs{9}Καὶ ἀπὸ υἱῶν Ἰωὰβ, Ἀβαδία υἱὸς Ἰεϊὴλ, καὶ μετʼ αὐτοῦ διακόσιοι δεκαοκτὼ τὰ ἀρσενικά.
\vs{10}Καὶ ἀπὸ τῶν υἱῶν Βαανὶ, Σελιμοὺθ υἱὸς Ἰωσεφία, καὶ μετʼ αὐτοῦ ἑκατὸν ἑξήκοντα τὰ ἀρσενικά.
\vs{11}Καὶ ἀπὸ υἱῶν Βαβὶ, Ζαχαρίας υἱὸς Βαβὶ, καὶ μετʼ αὐτοῦ εἰκοσιοκτὼ τὰ ἀρσενικά.
\vs{12}Καὶ ἀπὸ υἱῶν Ἀσγὰδ, Ἰωανὰν υἱὸς Ἀκκατὰν, καὶ μετʼ αὐτοῦ ἑκατὸν δέκα τὰ ἀρσενικά.
\vs{13}Καὶ ἀπὸ υἱῶν Ἀδωνικὰμ ἔσχατοι, καὶ ταῦτα τὰ ὀνόματα αὐτῶν, Ἐλιφαλὰτ, Ἰεὴλ, καὶ Σαμαΐα, καὶ μετʼ αὐτῶν ἑξήκοντα τὰ ἀρσενικά.
\vs{14}Καὶ ἀπὸ υἱῶν Βαγουαῒ, Οὐθαῒ καὶ Ζαβοὺδ, καὶ μετʼ αὐτοῦ ἑβδομήκοντα τὰ ἀρσενικά.

\vs{15}Καὶ συνῆξα αὐτοὺς πρὸς τὸν ποταμὸν τὸν ἐρχόμενον πρὸς τὸν Εὐὶ, καὶ παρενεβάλομεν ἐκεῖ ἡμέρας τρεῖς· καὶ συνῆκα ἐν τῷ λαῷ καὶ ἐν τοῖς ἱερεῦσι, καὶ ἀπὸ υἱῶν Λευὶ οὐχ εὗρον ἐκεῖ.
\vs{16}Καὶ ἀπέστειλα τῷ Ἐλεάζαρ, τῷ Ἀριὴλ, τῷ Σεμεΐᾳ, καὶ τῷ Ἀλωνὰμ, καὶ τῷ Ἰαρὶβ, καὶ τῷ Ἐλνάθαμ, καὶ τῷ Νάθαν, καὶ τῷ Ζαχαρίᾳ, καὶ τῷ Μεσολλὰμ, καὶ τῷ Ἰωαρὶμ, καὶ τῷ Ἐλνάθαν, συνιέντας.
\vs{17}Καὶ ἐξήνεγκα αὐτοὺς ἐπὶ ἄρχοντας ἐν ἀργυρίῳ τοῦ τόπου, καὶ ἔθηκα ἐν στόματι αὐτῶν λόγους λαλῆσαι πρὸς τοὺς ἀδελφοὺς αὐτῶν τῶν Ἀθινεὶμ ἐν ἀργυρίῳ τοῦ τόπου, τοῦ ἐνέγκαι ἡμῖν ᾄδοντας εἰς οἶκον Θεοῦ ἡμῶν.
\vs{18}Καὶ ἤλθοσαν ἡμῖν ὡς χεὶρ Θεοῦ ἡμῶν ἀγαθὴ ἐφʼ ἡμᾶς, ἀνὴρ σαχὼν ἀπὸ υἱῶν Μοολὶ, υἱοῦ Λευὶ, υἱοῦ Ἰσραήλ· καὶ ἀρχὴν ἦλθον οἱ υἱοὶ αὐτοῦ καὶ ἀδελφοὶ αὐτοῦ δεκαοκτώ·
\vs{19}Καὶ τὸν Ἀσεβία, καὶ τὸν Ἰσαΐα ἀπὸ τῶν υἱῶν Μεραρὶ, ἀδελφοὶ αὐτοῦ καὶ υἱοὶ αὐτοῦ εἴκοσι.
\vs{20}Καὶ ἀπὸ τῶν Ναθινὶμ, ὧν ἔδωκε Δαυὶδ καὶ οἱ ἄρχοντες εἰς δουλείαν τῶν Λευιτῶν, Ναθινὶμ διακόσιοι εἴκοσι, πάντες συνήχθησαν ἐν ὀνόμασι.

\vs{21}Καὶ ἐκάλεσα ἐκεῖ νηστείαν ἐπὶ τὸν ποταμὸν Ἀουὲ, τοῦ ταπεινωθῆναι ἐνώπιον τοῦ Θεοῦ ἡμῶν, ζητῆσαι παρʼ αὐτοῦ ὁδὸν εὐθείαν ἡμῖν καὶ τοῖς τέκνοις ἡμῶν καὶ πάσῃ τῇ κτήσει ἡμῶν·
\vs{22}Ὅτι ᾐσχύνθην αἰτήσασθαι παρὰ τοῦ βασιλέως δύναμιν καὶ ἱππεῖς σῶσαι ἡμᾶς ἀπὸ ἐχθροῦ ἐν τῇ ὁδῷ, ὅτι εἴπαμεν τῷ βασιλεῖ, λέγοντες, χεὶρ τοῦ Θεου ἡμῶν ἐπὶ πάντας τοὺς ζητοῦντας αὐτὸν εἰς ἀγαθόν· καὶ κράτος αὐτοῦ, καὶ θυμὸς αὐτοῦ, ἐπὶ πάντας τοὺς ἐγκαταλείποντας αὐτόν.
\vs{23}Καὶ ἐνηστεύσαμεν, καὶ ἐζητήσαμεν παρὰ τοῦ Θεοῦ ἡμῶν περὶ τούτου, καὶ ἐπήκουσεν ἡμῖν.

\vs{24}Καὶ διέστειλα ἀπὸ ἀρχόντων τῶν ἱερέων δώδεκα, τῷ Σαραΐα, τῷ Ἀσαβία, καὶ μετʼ αὐτῶν ἀπὸ ἀδελφῶν αὐτῶν δέκα.
\vs{25}Καὶ ἔστησα αὐτοῖς τὸ ἀργύριον καὶ τὸ χρυσίον καὶ τὰ σκεύη ἀπαρχῆς οἴκου Θεοῦ ἡμῶν, ἃ ὕψωσεν ὁ βασιλεὺς καὶ οἱ σύμβουλοι αὐτοῦ καὶ οἱ ἄρχοντες αὐτοῦ, καὶ πᾶς Ἰσραὴλ οἱ εὑρισκόμενοι.
\vs{26}Καὶ ἔστησα ἐπὶ χεῖρας αὐτῶν ἀργυρίου τάλαντα ἑξακόσια πεντήκοντα, καὶ σκεύη ἀργυρᾶ ἑκατὸν, καὶ τάλαντα χρυσίου ἑκατὸν,
\vs{27}καὶ χαφουρῆ χρυσοῖ εἴκοσι εἰς τὴν ὁδὸν χίλιοι, καὶ σκεύη χαλκοῦ στίλβοντος ἀγαθοῦ διάφορα ἐπιθυμητὰ ἐν χρυσίῳ.
\vs{28}Καὶ εἶπα πρὸς αὐτοὺς, ὑμεῖς ἅγιοι τῷ Κυρίῳ, καὶ τὰ σκεύη ἅγια, καὶ τὸ ἀργύριον καὶ τὸ χρυσίον ἑκούσια τῷ Κυρίῳ Θεῷ πατέρων ἡμῶν.
\vs{29}Ἀγρυπνεῖτε καὶ τηρεῖτε ἕως στῆτε ἐνώπιον ἀρχόντων τῶν ἱερέων καὶ τῶν Λευιτῶν καὶ τῶν ἀρχόντων τῶν πατριῶν ἐν Ἱερουσαλὴμ εἰς σκηνὰς οἴκου Κυρίου.
\vs{30}Καὶ ἐδέξαντο οἱ ἱερεῖς καὶ οἱ Λευῖται σταθμὸν τοῦ ἀργυρίου καὶ τοῦ χρυσίου καὶ τῶν σκευῶν, ἐνεγκεῖν εἰς Ἱερουσαλὴμ εἰς οἶκον Θεοῦ ἡμῶν.

\vs{31}Καὶ ἐξῄραμεν ἀπὸ τοῦ ποταμοῦ τοῦ Ἀουὲ ἐν τῇ δωδεκάτῃ τοῦ μηνὸς τοῦ πρώτου τοῦ ἐλθεῖν εἰς Ἱερουσαλήμ· καὶ χεὶρ Θεοῦ ἡμῶν ἦν ἐφʼ ἡμῖν, καὶ ἐῤῥύσατο ἡμᾶς ἀπὸ χειρὸς ἐχθροῦ καὶ πολεμίου ἐν τῇ ὁδῷ.
\vs{32}Καὶ ἤλθομεν εἰς Ἱερουσαλὴμ, καὶ ἐκαθίσαμεν ἐκεῖ ἡμέρας τρεῖς.
\vs{33}Καὶ ἐγενήθη τῇ ἡμέρᾳ τῇ τετάρτῃ ἐστήσαμεν τὸ ἀργύριον καὶ τὸ χρυσίον καὶ τὰ σκεύη ἐν οἴκῳ Θεοῦ ἡμῶν ἐπὶ χεῖρα Μεριμὼθ υἱοῦ Οὐρία τοῦ ἱερέως, καὶ μετʼ αὐτοῦ Ἐλεάζαρ υἱὸς Φινεὲς, καὶ μετʼ αὐτῶν Ἰωζαβὰδ υἱὸς Ἰησοῦ, καὶ Νωαδία υἱὸς Βαναΐα οἱ Λευῖται.
\vs{34}Ἐν ἀριθμῷ καὶ ἐν σταθμῷ τὰ πάντα, καὶ ἐγράφη πᾶς ὁ σταθμός.

\vs{35}Ἐν τῷ καιρῷ ἐκείνῳ οἱ ἐλθόντες ἐκ τῆς αἰχμαλωσίας υἱοὶ τῆς παροικίας, προσήνεγκαν ὁλοκαυτώσεις τῷ Θεῷ Ἰσραὴλ, μόσχους δώδεκα περὶ παντὸς Ἰσραὴλ, κριοὺς ἐννενηκανταὲξ, ἀμνοὺς ἑβδομηκονταεπτὰ, χιμάρους περὶ ἁμαρτίας δώδεκα, τὰ πάντα ὁλοκαυτώματα τῷ Κυρίῳ·
\vs{36}Καὶ ἔδωκαν τὸ νόμισμα τοῦ βασιλέως τοῖς διοικηταῖς τοῦ βασιλέως καὶ ἐπάρχοις πέραν τοῦ ποταμοῦ· καὶ ἐδόξασαν τὸν λαὸν καὶ τὸν οἶκον τοῦ Θεοῦ.

\ch{9}
Καὶ ὡς ἐτελέσθη ταῦτα, ἤγγισαν πρὸς μὲ οἱ ἄρχοντες, λέγοντες, οὐκ ἐχωρίσθη ὁ λαὸς Ἰσραὴλ καὶ οἱ ἱερεῖς καὶ οἱ Λευῖται ἀπὸ λαῶν τῶν γαιῶν ἐν μακρύμμασιν αὐτῶν, τῷ Χανανὶ ὁ Ἐθὶ, ὁ Φερεζὶ, ὁ Ἰεβουσὶ, ὁ Ἀμμωνὶ, ὁ Μωαβὶ, καὶ ὁ Μοσερὶ, καὶ ὁ Ἀμοῤῥὶ,
\vs{2}ὅτι ἐλάβοσαν ἀπὸ θυγατέρων αὐτῶν ἑαυτοῖς καὶ τοῖς υἱοῖς αὐτῶν· καὶ παρήχθη σπέρμα τὸ ἅγιον ἐν λαοῖς τῶν γαιῶν, καὶ χεὶρ τῶν ἀρχόντων ἐν τῇ ἀσυνθεσίᾳ ταύτῃ ἐν ἀρχῇ.
\vs{3}Καὶ ὡς ἤκουσα τὸν λόγον τοῦτον, διέῤῥηξα τὰ ἱμάτιά μου, καὶ ἐπαλλόμην, καὶ ἔτιλλον ἀπὸ τῶν τριχῶν τῆς κεφαλῆς μου καὶ ἀπὸ τοῦ πώγωνός μου, καὶ ἐκαθήμην ἠρεμάζων.
\vs{4}Καὶ συνήχθησαν πρὸς μὲ πᾶς ὁ διώκων λόγον Θεοῦ Ἰσραὴλ ἐπὶ ἀσυνθεσίᾳ τῆς ἀποικίας· κᾀγὼ καθήμενος ἠρεμάζων ἕως τῆς θυσίας τῆς ἑσπερινῆς.

\vs{5}Καὶ ἐν θυσίᾳ τῇ ἑσπερινῇ ἀνέστην ἀπὸ ταπεινώσεώς μου· καὶ ἐν τῷ διαῤῥήξαί με τὰ ἱμάτιά μου, καὶ ἐπαλλόμην, καὶ κλίνω ἐπὶ τὰ γόνατά μου, καὶ ἐκπετάζω τὰς χεῖράς μου πρὸς Κύριον τὸν Θεὸν,
\vs{6}καὶ εἶπα, Κύριε, ἠσχύνθην καὶ ἐνετράπην τοῦ ὑψῶσαι, Θεέ μου, τὸ πρόσωπόν μου πρὸς σὲ, ὅτι αἱ ἀνομίαι ἡμῶν ἐπληθύνθησαν ὑπὲρ κεφαλῆς ἡμῶν, καὶ αἱ πλημμέλειαι ἡμῶν ἐμεγαλύνθησαν ἕως εἰς τὸν οὐρανόν.
\vs{7}Ἀπὸ ἡμερῶν πατέρων ἡμῶν ἐσμὲν ἐν πλημμελείᾳ μεγάλῃ ἕως τῆς ἡμέρας ταύτης· καὶ ἐν ταῖς ἀνομίαις ἡμῶν παρεδόθημεν ἡμεῖς καὶ οἱ βασιλεῖς ἡμῶν καὶ οἱ υἱοὶ ἡμῶν ἐν χειρὶ βασιλέων τῶν ἐθνῶν ἐν ῥομφαίᾳ, καὶ ἐν αἰχμαλωσίᾳ, καιὶ ἐν διαρπαγῇ, καὶ ἐν αἰσχύνῃ προσώπου ἡμῶν, ὡς ἡ ἡμέρα αὕτη.
\vs{8}Καὶ νῦν ἐπιεικεύσατο ἡμῖν ὁ Θεὸς ἡμῶν τοῦ καταλιπεῖν ἡμῖν εἰς σωτηρίαν, καὶ δοῦναι ἡμῖν στήριγμα ἐν τόπῳ ἁγιάσματος αὐτοῦ, τοῦ φωτίσαι ὀφθαλμοὺς ἡμῶν, καὶ δοῦναι ζωοποίησιν μικρὰν ἐν τῇ δουλείᾳ ἡμῶν.
\vs{9}Ὅτι δοῦλοι ἐσμὲν, καὶ ἐν τῇ δουλείᾳ ἡμῶν οὐκ ἐγκατέλιπεν ἡμᾶς Κύριος ὁ Θεὸς ἡμῶν· καὶ ἔκλινεν ἐφʼ ἡμᾶς ἔλεος ἐνώπιον βασιλέων Περσῶν, δοῦναι ἡμῖν ζωοποίησιν τοῦ ὑψῶσαι αὐτοὺς τὸν οἶκον τοῦ Θεοῦ ἡμῶν, καὶ ἀναστῆσαι τὰ ἔρημα αὐτῆς, καὶ τοῦ δοῦναι ἡμῖν φραγμὸν ἐν Ἰούδα καὶ Ἱερουσαλήμ.
\vs{10}Τί εἴπωμεν ὁ Θεὸς ἡμῶν μετὰ τοῦτο; ὅτι ἐγκατελίπομεν ἐντολάς σου,
\vs{11}ἃς ἔδωκας ἡμῖν ἐν χειρὶ δούλων σου τῶν προφητῶν, λέγων, ἡ γῆ εἰς ἣν εἰσπορεύεσθε κληρονομῆσαι αὐτὴν, γῆ μετακινουμένη ἐστὶν ἐν μετακινήσει λαῶν τῶν ἐθνῶν ἐν μακρύμμασιν αὐτῶν, ὧν ἔπλησαν αὐτὴν ἀπὸ στόματος ἐπὶ στόμα ἐν ἀκαθαρσίαις αὐτῶν.

\vs{12}Καὶ νῦν τὰς θυγατέρας ὑμῶν μὴ δότε τοῖς υἱοῖς αὐτῶν, καὶ ἀπὸ τῶν θυγατέρων αὐτῶν μὴ λάβητε τοῖς υἱοῖς ὑμῶν, καὶ οὐκ ἐκζητήσετε εἰρήνην αὐτῶν καὶ ἀγαθὸν αὐτῶν ἕως αἰῶνος, ὅπως ἐνισχύσητε, καὶ φάγητε τὰ ἀγαθὰ τῆς γῆς, καὶ κληροδοτήσητε τοῖς υἱοῖς ὑμῶν ἕως αἰῶνος.
\vs{13}Καὶ μετὰ πᾶν τὸ ἐρχόμενον ἐφʼ ἡμᾶς ἐν ποιήμασιν ἡμῶν τοῖς πονηροῖς καὶ ἐν πλημμελείᾳ ἡμῶν τῇ μεγάλῃ, ὅτι οὐκ ἔστιν ὡς ὁ Θεὸς ἡμῶν, ὅτι ἐκούφισας ἡμῶν τὰς ἀνομίας, καὶ ἔδωκας ἡμῖν σωτηρίαν·
\vs{14}Ὅτι ἐπεστρέψαμεν διασκεδάσαι ἐντολάς σου, καὶ ἐπιγαμβρεῦσαι τοῖς λαοῖς τῶν γαιῶν· μὴ παροξυνθῇς ἐν ἡμῖν ἕως συντελείας, τοῦ μὴ εἶναι ἐγκατάλειμμα καὶ διασωζόμενον.
\vs{15}Κύριε ὁ Θεὸς Ἰσραὴλ, δίκαιος σὺ, ὅτι κατελείφθημεν διασωζόμενοι, ὡς ἡ ἡμέρα αὕτη· ἰδοὺ ἡμεις ἐναντίον σου ἐν πλημμελείαις ἡμῶν, ὅτι οὐκ ἔστι στῆναι ἐνώπιόν σου ἐπὶ τούτῳ.

\ch{10}
Καὶ ὡς προσηύξατο Ἔσδρας, καὶ ὡς ἐξηγόρευσε κλαίων καὶ προσευχόμενος ἐνώπιον οἴκου τοῦ Θεοῦ, συνήχθησαν πρὸς αὐτὸν ἀπὸ Ἰσραὴλ ἐκκλησία πολλὴ σφόδρα, ἄνδρες καὶ γυναῖκες καὶ νεανίσκοι, ὅτι ἔκλαυσαν ὁ λαὸς, καὶ ὕψωσε κλαίων.
\vs{2}Καὶ ἀπεκρίθη Σεχενίας υἱὸς Ἰεὴλ ἀπὸ υἱῶν Ἠλὰμ, καὶ εἶπε τῷ Ἔσδρᾳ, ἡμεῖς ἠσυνθετήσαμεν τῷ Θεῷ ἡμῶν, καὶ ἐκαθίσαμεν γυναῖκας ἀλλοτρίας ἀπὸ τῶν λαῶν τῆς γῆς· καὶ νῦν ἐστιν ὑπομονὴ τῷ Ἰσραὴλ ἐπὶ τούτῳ.
\vs{3}Καὶ νῦν διαθώμεθα διαθήκην τῷ Θεῷ ἡμῶν ἐκβαλεῖν πάσας τὰς γυναῖκας, καὶ τὰ γενόμενα ἐξ αὐτῶν, ὡς ἂν βούλῃ· ἀνάστηθι, καὶ φοβέρισον αὐτοὺς ἐν ἐντολαῖς Θεοῦ ἡμῶν, καὶ ὡς ὁ νόμος, γενηθήτω.
\vs{4}Ἀνάστα, ὅτι ἐπὶ σὲ τὸ ῥῆμα, καὶ ἡμεῖς μετὰ σοῦ· κραταιοῦ καὶ ποίησον.

\vs{5}Καὶ ἀνέστη Ἔσδρας, καὶ ὥρκισε τοὺς ἄρχοντας, τοὺς ἱερεῖς, καὶ Λευείτας, καὶ πάντα Ἰσραὴλ, τοῦ ποιῆσαι κατὰ τὸ ῥῆμα τοῦτο· καὶ ὤμοσαν.
\vs{6}Καὶ ἀνέστη Ἔσδρας ἀπὸ προσώπου οἴκου τοῦ Θεοῦ, καὶ ἐπορεύθη εἰς γαζοφυλάκιον Ἰωανὰν υἱοῦ Ἐλισοὺβ, και ἐπορεύθη ἐκεῖ· ἄρτον οὐκ ἔφαγε, καὶ ὕδωρ οὐκ ἔπιεν, ὅτι ἐπένθει ἐπὶ τῇ ἀσυνθεσίᾳ τῆς ἀποικίας.
\vs{7}Καὶ παρήνεγκαν φωνὴν ἐν Ἰούδα καὶ ἐν Ἱερουσαλὴμ πᾶσι τοῖς υἱοῖς τῆς ἀποικίας, τοῦ συναθροισθῆναι εἰς Ἱερουσαλήμ.
\vs{8}Πᾶς ὃς ἂν μὴ ἔλθῃ εἰς τρεῖς ἡμέρας, ὡς ἡ βουλὴ τῶν ἀρχόντων καὶ τῶν πρεσβυτέρων, ἀναθεματισθήσεται πᾶσα ἡ ὕπαρξις αὐτοῦ, καὶ αὐτὸς διασταλήσεται ἀπὸ ἐκκλησίας τῆς ἀποικίας.

\vs{9}Καὶ συνήχθησαν πάντες ἄνδρες Ἰούδα καὶ Βενιαμὶν εἰς Ἱερουσαλὴμ εἰς τὰς τρεῖς ἡμέρας· οὗτος ὁ μὴν ὁ ἔννατος· ἐν εἰκάδι τοῦ μηνὸς ἐκάθισε πᾶς ὁ λαὸς ἐν πλατείᾳ οἴκου τοῦ Θεοῦ ἀπὸ θορύβου αὐτῶν περὶ τοῦ ῥήματος, καὶ ἀπὸ τοῦ χειμῶνος.
\vs{10}Καὶ ἀνέστη Ἔσδρας ὁ ἱερεὺς, καὶ εἶπε πρὸς αὐτοὺς, ὑμεῖς ἠσυνθετήκατε, καὶ ἐκαθίσατε γυναῖκας ἀλλοτρίας τοῦ προσθεῖναι ἐπὶ πλημμέλειαν Ἰσραήλ.
\vs{11}Καὶ νῦν δότε αἴνεσιν Κυρίῳ Θεῷ τῶν πατέρῶν ἡμῶν, καὶ ποιήσατε τὸ ἀρεστὸν ἐνώπιον αὐτοῦ, καὶ διαστάλητε ἀπὸ λαῶν τῆς γῆς καὶ ἀπὸ τῶν γυναικῶν τῶν ἀλλοτρίων.

\vs{12}Καὶ ἀπεκρίθησαν πᾶσα ἡ ἐκκλησία, καὶ εἶπαν, μέγα τοῦτο τὸ ῥῆμά σου ἐφʼ ἡμᾶς ποιῆσαι.
\vs{13}Ἀλλὰ ὁ λαὸς πολὺς, καὶ ὁ καιρὸς χειμερινὸς, καὶ οὐκ ἔστι δύναμις στῆναι ἔξω· καὶ τὸ ἔργον οὐκ εἰς ἡμέραν μίαν καὶ οὐκ εἰς δύο, ὅτι ἐπληθύναμεν τοῦ ἀδικῆσαι ἐν τῷ ῥῆματι τούτῳ·
\vs{14}Στήτωσαν δὴ ἄρχοντες ἡμῶν, καὶ πᾶσι τοῖς ἐν πόλεσιν ἡμῶν ὃς ἐκάθισε γυναῖκας ἀλλοτρίας, ἐλθέτωσαν εἰς καιροὺς ἀπὸ συνταγῶν, καὶ μετʼ αὐτῶν πρεσβύτεροι πόλεως καὶ πόλεως, καὶ κριταὶ, τοῦ ἀποστρέψαι ὀργὴν θυμοῦ Θεοῦ ἡμῶν ἐξ ἡμῶν, περὶ τοῦ ῥήματος τούτου.
\vs{15}Πλὴν Ἰωνάθαν υἱὸς Ἀσὴλ, καὶ Ἰαζίας υἱὸς Θεκωὲ μετʼ ἐμοῦ περὶ τούτου· καὶ Μεσουλλὰμ, καὶ Σαββαθαῒ ὁ Λευίτης βοηθῶν αὐτοῖς.

\vs{16}Καὶ ἐποίησαν οὕτως υἱοὶ τῆς ἀποικίας· Καὶ διεστάλησαν Ἔσρας ὁ ἱερεὺς, καὶ ἄνδρες ἄρχοντες πατριῶν τῷ οἴκῳ, καὶ πάντες ἐν ὀνόμασιν, ὅτι ἐπέστρεψαν ἐν ἡμέρᾳ μιᾷ τοῦ μηνὸς τοῦ δεκάτου ἐκζητῆσαι τὸ ῥῆμα.
\vs{17}Καὶ ἐτέλεσαν ἐν πᾶσιν ἀνδράσιν οἳ ἐκάθισαν γυναῖκας ἀλλοτρίας, ἕως ἡμέρας μιᾶς τοῦ μηνὸς τοῦ πρώτου.

\vs{18}Καὶ εὑρέθη ἀπὸ υἱῶν τῶν ἱερέων οἳ ἐκάθισαν γυναῖκας ἀλλοτρίας, ἀπὸ υἱῶν Ἰησοῦ υἱοῦ Ἰωσεδὲκ, καὶ ἀδελφοὶ αὐτοῦ Μαασία, καὶ Ἐλιέζερ, καὶ Ἰαρὶβ, καὶ Γαδαλία.
\vs{19}Καὶ ἔδωκαν χεῖρα αὐτῶν τοῦ ἐξενέγκαι γυναῖκας ἑαυτῶν, καὶ πλημμελείας κριὸν ἐκ προβάτων περὶ πλημμελήσεως αὐτῶν.
\vs{20}Καὶ ἀπὸ υἱῶν Ἐμμὴρ, Ἀνανὶ, καὶ Ζαβδία.
\vs{21}Καὶ ἀπὸ υἱῶν Ἠρὰμ, Μασαὴλ, καὶ Ἐλία, καὶ Σαμαΐα, καὶ Ἰεὴλ, καὶ Ὀζία.
\vs{22}Καὶ ἀπὸ υἱῶν Φασοὺρ, Ἐλιωναῒ, Μαασία, καὶ Ἰσμαὴλ, καὶ Ναθαναὴλ, καὶ Ἰωζαβὰδ, καὶ Ἠλασά.
\vs{23}Καὶ ἀπὸ τῶν Λευιτῶν, Ἰωζαβὰδ, καὶ Σαμοὺ, καὶ Κωλία, αὐτὸς Κωλίτας, καὶ Φεθεΐα, καὶ Ἰούδας, καὶ Ἐλιέζερ.
\vs{24}Καὶ ἀπὸ τῶν ᾀδόντων, Ἐλισάβ· καὶ ἀπὸ τῶν πυλωρῶν, Σολμὴν, καὶ Τελμὴν, καὶ Ὠδούθ.
\vs{25}Καὶ ἀπὸ Ἰσραὴλ, ἀπὸ υἱῶν Φόρος, Ῥαμία, καὶ Ἀζία, καὶ Μελχία, καὶ Μεαμὶν, καὶ Ἐλεάζαρ, καὶ Ἀσαβία, καὶ Βαναία.
\vs{26}Καὶ ἀπὸ υἱῶν Ἡλὰμ, Ματθανία, καὶ Ζαχαρία, καὶ Ἰαϊὴλ, καὶ Ἀβδία, καὶ Ἰαριμὼθ, καὶ Ἠλία.
\vs{27}Καὶ ἀπὸ υἱῶν Ζαθούα, Ἐλιωναῒ, Ἐλισοὺβ, Ματθαναῒ, καὶ Ἀρμὼθ, καὶ Ζαβὰδ, καὶ Ὀζιζά.
\vs{28}Καὶ ἀπὸ υἱῶν Βαβεῒ, Ἰωανὰν, Ἀνανία, καὶ Ζαβοὺ, καὶ Θαλί.
\vs{29}Καὶ ἀπὸ υἱῶν Βανουῒ, Μοσολλὰμ, Μαλοὺχ, Ἀδαΐας, Ἰασοὺβ, καὶ Σαλουΐα, καὶ Ῥημώθ.
\vs{30}Καὶ ἀπὸ υἱῶν Φαὰθ Μωὰβ, Ἐδνὲ, καὶ Χαλὴλ, καὶ Βαναία, Μαασία, Ματθανία, Βεσελεὴλ, καὶ Βανουῒ, καὶ Μανασσῆ·
\vs{31}Καὶ ἀπὸ υἱῶν Ἠρὰμ, Ἐλιέζερ, Ἰεσία, Μελχία, Σαμαΐας, Σεμεὼν,
\vs{32}Βενιαμὶν, Βαλοὺχ, Σαμαρία.
\vs{33}Καὶ ἀπὸ υἱῶν Ἀσὴμ, Μετθανία, Ματθαθὰ, Ζαδὰβ, Ἐλιφαλὲτ, Ἱεραμὶ, Μανασσῆ, Σεμεΐ.
\vs{34}Καὶ ἀπὸ υἱῶν Βανὶ, Μοοδία, Ἀμρὰμ, Οὐὴλ,
\vs{35}Βαναία, Βαδαία, Χελκία,
\vs{36}Οὐουανία, Μαριμὼθ, Ἐλιασὶφ,
\vs{37}Ματθανία, Ματθαναΐ· καὶ ἐποίησαν
\vs{38}οἱ υἱοἱ Βανουὶ, καὶ οἱ υἱοὶ Σεμεῒ,
\vs{39}καὶ Σελεμία, καὶ Νάθαν, καὶ Ἀδαΐα,
\vs{40}Μαχαδναβοὺ, Σεσεῒ, Σαριοὺ,
\vs{41}Ἐζριὴλ, καὶ Σελεμία, καὶ Σαμαρία,
\vs{42}καὶ Σελλοὺμ, Ἀμαρεία, Ἰωσήφ.
\vs{43}Ἀπὸ υἱῶν Ναβοὺ, Ἰαὴλ, Ματθανίας, Ζαβὰδ, Ζεβεννὰς, Ἰαδαὶ, καὶ Ἰωὴλ, καὶ Βαναία.

\vs{44}Πάντες οὗτοι ἐλάβοσαν γυναῖκας ἀλλοτρίας, καὶ ἐγέννησαν ἐξ αὐτῶν υἱούς.


\def\book{ΝΕΕΜΙΑΣ}
\biblebook{ΝΕΕΜΙΑΣ}


\lettrine[lines=2, loversize=0.2, nindent=0em, findent=.25em]{\textcolor{bookheadingcolor}{Λ}}{ΟΓΟΙ} Νεεμία υἱοῦ Χελκία. Καὶ ἐγένετο ἐν μηνὶ Χασελεῦ ἔτους εἰκοστοῦ, καὶ ἐγὼ ἤμηνν ἐν Σουσὰν ἀβιρά·
\vs{2}καὶ ἦλθεν Ἀνανὶ εἷς ἀπὸ ἀδελφῶν μου, αὐτὸς καὶ ἄνδρες Ἰούδα· καὶ ἠρώτησα αὐτοὺς περὶ τῶν σωθέντων, οἳ κατελείθησαν ἀπὸ τῆς αἰχμαλωσίας, καὶ περὶ Ἱερουσαλήμ.
\vs{3}Καὶ εἴποσαν πρὸς μέ, Οἱ καταλειπόμενοι οἱ καταλειφθέντες ἀπὸ τῆς αἰχμαλωσίας ἐκεῖ ἐν τῇ χώρᾳ, ἐν πονηρίᾳ μεγάλῃ καὶ ἐν ὀνειδισμῷ, καὶ τείχη Ἱερουσαλὴμ καθῃρημένα, καὶ αἱ πύλαι αὐτῆς ἐνεπρήσθησαν ἐν πυρί.

\vs{4}Καὶ ἐγένετο ἐν τῷ ἀκοῦσαί με τοὺς λόγους τούτους, ἐκάθισα καὶ ἔκλαυσα καὶ ἐπένθησα ἡμέρας, καὶ ἤμην νηστεύων καὶ προσευχόμενος ἐνώπιον τοῦ Θεοῦ τοῦ οὐρανοῦ.
\vs{5}Καὶ εἶπα, μὴ δὴ Κύριε ὁ Θεὸς τοῦ οὐρανοῦ, ὁ ἰσχυρὸς, ὁ μέγας καὶ φοβερὸς, φυλάσσων τὴν διαθήκην καὶ τὸ ἔλεός σου τοῖς ἀγαπῶσιν αὐτὸν καὶ τοῖς φυλάσσουσι τὰς ἐντολὰς αὐτοῦ·
\vs{6}Ἔστω δὴ τὸ οὖς σου προσέχον, καὶ οἱ ὀφθαλμοί σου ἀνεῳγμένοι, τοῦ ἀκοῦσαι προσευχὴν τοῦ δούλου σου, ἣν ἐγὼ προσεύχομαι ἐνώπιόν σου σήμερον ἡμέραν καὶ νύκτα περὶ υἱῶν Ἰσραὴλ δούλων σου· καὶ ἐξαγορεύω ἐπὶ ἁμαρτίαις υἱῶν Ἰσραὴλ αἷς ἡμάρτομέν σοι· καὶ ἐγὼ καὶ ὁ οἶκος πατρός μου ἡμάρτομεν.
\vs{7}Διαλύσει διελύσαμεν πρὸς σὲ, καὶ οὐκ ἐφυλάξαμεν τὰς ἐντολὰς καὶ τὰ προστάγματα καὶ τὰ κρίματα, ἃ ἐνετείλω τῷ Μωυσῇ παιδί σου.
\vs{8}Μνήσθητι δὴ τὸν λόγον ὃν ἐνετείλω τῷ Μωυσῇ παιδί σου, λέγων, ὑμεῖς ἐὰν ἀσυνθετήσητε, ἐγὼ διασκορπιῶ ὑμᾶς ἐν τοῖς λαοῖς·
\vs{9}καὶ ἐὰν ἐπιστρέψητε πρὸς μὲ, καὶ φυλάξητε τὰς ἐντολάς μου, καὶ ποιήσητε αὐτὰς, ἐὰν ᾖ ἡ διασπορὰ ὑμῶν ἀπʼ ἄκρου τοῦ οὐρανοῦ, ἐκεῖθεν συνάξω αὐτοὺς, καὶ εἰσάξω αὐτοὺς εἰς τὸν τόπον, ὃν ἐξελεξάμην κατασκηνῶσαι τὸ ὄνομά μου ἐκεῖ.
\vs{10}Καὶ αὐτοὶ παῖδές σου καὶ λαός σου, οὓς ἐλυτρώσω ἐν τῇ δυνάμει σου τῇ μεγάλῃ, καὶ ἐν τῇ χειρί σου τῇ κραταιᾷ.

\vs{11}Μὴ δὴ Κύριε, ἀλλὰ ἔστω τὸ οὖς σου προσέχον εἰς τὴν προσευχὴν τοῦ δούλου σου, καὶ εἰς τὴν προσευχὴν παίδων σου τῶν θελόντων φοβεῖσθαι τὸ ὄνομά σου· καὶ εὐόδωσον δὴ τῷ παιδί σου σήμερον, καὶ δὸς αὐτὸν εἰς οἰκτιρμοὺς ἐνώπιον τοῦ ἀνδρὸς τούτου· καὶ ἐγὼ ἤμην οἰνοχόος τῷ βασιλεῖ.

\ch{2}
Καὶ ἐγένετο ἐν μηνὶ Νισὰν ἔτους εἰκοστοῦ Ἀρθασασθὰ βασιλεῖ, καὶ ἦν ὁ οἶνος ἐνώπιον ἐμοῦ· καὶ ἔλαβον τὸν οἶνον, καὶ ἔδωκα τῷ βασιλεῖ· καὶ οὐκ ἦν ἕτερος ἐνώπιον αὐτοῦ.

\vs{2}Καὶ εἶπέ μοι ὁ βασιλεὺς, διὰ τί τὸ πρόσωπόν σου πονηρὸν, καὶ οὐκ εἶ μετριάζων; καὶ οὐκ ἔστι τοῦτο, εἰ μὴ πονηρία καρδίας· καὶ ἐφοβήθην πολὺ σφόδρα,
\vs{3}καὶ εἶπα τῷ βασιλεῖ, ὁ βασιλεὺς εἰς τὸν αἰῶνα ζήτω· διὰ τί οὐ μὴ γένηται πονηρὸν τὸ πρόσωπόν μου, διότι ἡ πόλις οἶκος μνημείων πατέρων μου ἠρημώθη, καὶ αἱ πύλαι αὐτῆς κατεβρώθησαν ἐν πυρί;
\vs{4}Καὶ εἶπέ μοι ὁ βασιλεύς, περὶ τίνος τοῦτο σὺ ζητεῖς; καὶ προσηυξάμην πρὸς τὸν Θεὸν τοῦ οὐρανοῦ,
\vs{5}καὶ εἶπα τῷ βασιλεῖ, εἰ ἐπὶ τὸν βασιλέα ἀγαθὸν, καὶ εἰ ἀγαθυνθήσεται ὁ παῖς σου ἐνώπιόν σοῦ, ὥστε πέμψαι αὐτὸν ἐν Ἰούδα εἰς πόλιν μνημείων πατέρων μου, καὶ ἀνοικοδομήσω αὐτήν.

\vs{6}Καὶ εἶπέ μοι ὁ βασιλεὺς, καὶ ἡ παλλακὴ ἡ καθημένη ἐχόμενα αὐτοῦ, ἕως πότε ἔσται ἡ πορεία σου, καὶ πότε ἐπιστρέψεις; καὶ ἠγαθύνθη ἐνώπιον τοῦ βασιλέως, καὶ ἀπέστειλέ με, καὶ ἔδωκα αὐτῷ ὅρον.
\vs{7}Καὶ εἶπα τῷ βασιλεῖ, εἰ ἐπὶ τὸν βασιλέα ἀγαθὸν, δότω μοι ἐπιστολὰς πρὸς τοὺς ἐπάρχους πέραν τοῦ ποταμοῦ, ὥστε παραγαγεῖν με ἕως ἔλθω ἐπὶ Ἰούδαν,
\vs{8}καὶ ἐπιστολὴν ἐπὶ Ἀσὰφ φύλακα τοῦ παραδείσου ὅς ἐστι τῷ βασιλεῖ, ὥστε δοῦναί μοι ξύλα στεγάσαι τὰς πύλας, καὶ εἰς τὸ τεῖχος τῆς πόλεως, καὶ εἰς οἶκον ὃν εἰσελεύσομαι εἰς αὐτόν· καὶ ἔδωκέ μοι ὁ βασιλεὺς ὡς χεὶρ Θεοῦ ἡ ἀγαθή.

\vs{9}Καὶ ἦλθον πρὸς τοὺς ἐπάρχους πέραν τοῦ ποταμοῦ, καὶ ἔδωκα αὐτοῖς τὰς ἐπιστολὰς τοῦ βασιλέως· καὶ ἀπέστειλε μετʼ ἐμοῦ ὁ βασιλεὺς ἀρχηγοὺς δυνάμεως καὶ ἱππεῖς.
\vs{10}Καὶ ἤκουσε Σαναβαλλὰτ ὁ Ἀρωνὶ, καὶ Τωβία ὁ δοῦλος Ἀμμωνὶ, καὶ πονηρὸν αὐτοῖς ἐγένετο, ὅτι ἥκει ὁ ἄνθρωπος ζητῆσαι ἀγαθὸν τοῖς υἱοῖς Ἰσραήλ.

\vs{11}Καὶ ἦλθον εἰς Ἱερουσαλὴμ, καὶ ἤμην ἐκεῖ ἡμέρας τρεῖς.
\vs{12}Καὶ ἀνέστην νυκτὸς ἐγὼ καὶ ἄνδρες ὀλίγοι μετʼ ἐμοῦ, καὶ οὐκ ἀπήγγειλα ἀνθρώπῳ τί ὁ Θεὸς δίδωσιν εἰς καρδίαν μου τοῦ ποιῆσαι μετὰ τοῦ Ἰσραὴλ, καὶ κτῆνος οὐκ ἔστι μετʼ ἐμοῦ, εἰ μὴ τὸ κτῆνος ᾧ ἐγὼ ἐπιβαίνω ἐπʼ αὐτῷ.
\vs{13}Καὶ ἐξῆλθον ἐν πύλῃ τοῦ γωληλὰ, καὶ πρὸς στόμα πηγῆς τῶν συκῶν, καὶ εἰς πύλην τῆς κοπρίας· καὶ ἤμην συντρίβων ἐν τῷ τείχει Ἱερουσαλὴμ ὃ αὐτοὶ καθαιροῦσι, καὶ πύλαι αὐτῆς κατεβρώθησαν πυρί.
\vs{14}Καὶ παρῆλθον ἐπὶ πύλην τοῦ αῒν, καὶ εἰς κολυμβήθραν τοῦ βασιλέως, καὶ οὐκ ἦν τόπος τῷ κτήνει παρελθεῖν ὑποκάτω μου.
\vs{15}Καὶ ἤμην ἀναβαίνων ἐν τῷ τείχει χειμάῤῥου νυκτός, καὶ ἤμην συντρίβων ἐν τῷ τείχει, καὶ ἤμην ἐν πύλῃ τῆς φάραγγος, καὶ ἐπέστρεψα.

\vs{16}Καὶ οἱ φυλάσσοντες οὐκ ἔγνωσαν τί ἐπορεύθην, καὶ τί ἐγὼ ποιῶ· καὶ τοῖς Ἰουδαίοις, καὶ τοῖς ἱερεῦσι, καὶ τοῖς ἐντίμοις, καὶ τοῖς στρατηγοῖς, καὶ τοῖς καταλοίποις τοῖς ποιοῦσι τὰ ἔργα, ἕως τότε οὐκ ἀπήγγειλα.
\vs{17}Καὶ εἶπα πρὸς αὐτοὺς, ὑμεῖς βλέπετε τὴν πονηρίαν ταύτην, ἐν ᾗ ἐσμὲν ἐν αὐτῇ, πῶς Ἱερουσαλὴμ ἔρημος, καὶ αἱ πύλαι αὐτῆς ἐδόθησαν πυρί· δεῦτε, καὶ διοικοδομήσωμεν τὸ τεῖχος Ἱερουσαλὴμ, καὶ οὐκ ἐσόμεθα ἔτι ὄνειδος.
\vs{18}Καὶ ἀπήγγειλα αὐτοῖς τὴν χεῖρα τοῦ Θεοῦ ᾗ ἐστιν ἀγαθὴ ἐπʼ ἐμὲ, καὶ πρὸς τοὺς λόγους τοῦ βασιλέως οὓς εἶπέ μοι· καὶ εἶπα, ἀναστῶμεν, καὶ οἰκοδομήσωμεν· καὶ ἐκραταιώθησαν αἱ χεῖρες αὐτῶν εἰς τὸ ἀγαθόν.

\vs{19}Καὶ ἤκουσε Σαναβαλλὰτ ὁ Ἀρωνὶ, καὶ Τωβία ὁ δοῦλος ὁ Ἀμμωνὶ, καὶ Γησὰμ ὁ Ἀραβί, καὶ ἐξεγέλασαν ἡμᾶς, καὶ ἦλθον ἐφʼ ἡμᾶς, καὶ εἶπον, τί τὸ ῥῆμα τοῦτο ὃ ὑμεῖς ποιεῖτε; ἢ ἐπὶ τὸν βασιλέα ὑμεῖς ἀποστατεῖτε;
\vs{20}Καὶ ἐπέστρεψα αὐτοῖς λόγον, καὶ εἶπα αὐτοῖς, ὁ Θεὸς τοῦ οὐρανοῦ αὐτὸς εὐοδώσει ἡμῖν, καὶ ἡμεῖς δοῦλοι αὐτοῦ καθαροὶ, καὶ οἰκοδομήσομεν· καὶ ὑμῖν οὐκ ἔστι μερὶς καὶ δικαιοσύνη καὶ μνημόσυνον ἐν Ἱερουσαλήμ.

\ch{3}
Καὶ ἀνέστη Ἐλιασοὺβ ὁ ἱερεὺς ὁ μέγας, καὶ οἱ ἀδελφοὶ αὐτοῦ οἱ ἱερεῖς, καὶ ᾠκοδόμησαν τὴν πύλην τὴν προβατικήν· αὐτοὶ ἡγίασαν αὐτὴν, καὶ ἔστησαν θύρας αὐτῆς, καὶ ἕως πύργου τῶν ἑκατὸν ἡγίασαν ἕως πύργου Ἀναμεήλ.
\vs{2}Καὶ ἐπὶ χεῖρας ἀνδρῶν υἱῶν Ἱεριχὼ, καὶ ἐπὶ χεῖρας υἱῶν Ζακχοὺρ, υἱοῦ Ἀμαρί.

\vs{3}Καὶ τὴν πύλην τὴν ἰχθυηρὰν ᾠκοδόμησαν υἱοὶ Ἀσανά· αὐτοὶ ἐστέγασαν αὐτὴν, καὶ ἐστέγασαν θύρας αὐτῆς καὶ κλεῖθρα αὐτῆς καὶ μοχλοὺς αὐτῆς.
\vs{4}Καὶ ἐπὶ χεῖρα αὐτῶν κατέσχεν ἐπὶ Ῥαμὼθ υἱοῦ Οὐρία, υἱοῦ Ἀκκώς· καὶ ἐπὶ χεῖρα αὐτῶν κατέσχε Μοσολλὰμ υἱὸς Βαραχίου, υἱοῦ Μαζεβήλ· καὶ ἐπὶ χεῖρα αὐτῶν κατέσχε Σαδὼκ υἱὸς Βαανά.
\vs{5}Καὶ ἐπὶ χεῖρα αὐτῶν κατέσχοσαν οἱ Θεκωῒμ, καὶ ἀδωρὶμ οὐκ εἰσήνεγκαν τράχηλον αὐτῶν εἰς δουλείαν αὐτῶν.

\vs{6}Καὶ τὴν πύλην ἰασαναῒ ἐκράτησαν Ἰωϊδὰ υἱὸς Φασὲκ, καὶ Μεσουλὰμ υἱὸς Βασωδία· αὐτοὶ ἐστέγασαν αὐτὴν, καὶ ἔστησαν θύρας αὐτῆς καὶ κλεῖθρα αὐτῆς καὶ μοχλοὺς αὐτῆς.
\vs{6a}Καὶ ἐπὶ χεῖρα αὐτῶν ἐκράτησαν Μαλτίας ὁ Γαβαωνίτης, καὶ Εὐάρων ὁ Μηρωνωθίτης, ἄνδρες τῆς Γαβαὼν καὶ τῆς Μασφὰ ἓως θρόνου τοῦ ἄρχοντος τοῦ πέραν τοῦ ποταμοῦ.
\vs{8}Καὶ παρʼ αὐτὸν παρησφαλίσατο Ὀζιὴλ υἱὸς Ἀραχίου πυρωτῶν· καὶ ἐπὶ χεῖρα αὐτῶν ἐκράτησεν Ἀνανίας υἱὸς τοῦ ῥωκεῒμ, καὶ κατέλιπον Ἱερουσαλὴμ ἕως τοῦ τείχους τοῦ πλατέος.
\vs{9}Καὶ ἐπὶ χεῖρα αὐτῶν ἐκράτησε Ῥαφαΐα υἱὸς Σοὺρ, ἄρχων ἡμίσους περιχώρου Ἱερουσαλήμ.
\vs{10}Καὶ ἐπὶ χεῖρα αὐτῶν ἐκράτησεν Ἰεδαΐα υἱὸς Ἑρωμὰφ, καὶ κατέναντι οἰκίας αὐτοῦ· καὶ ἐπὶ χεῖρα αὐτοῦ ἐκράτησεν Ἀττοὺθ υἱὸς Ἀσαβανία.
\vs{11}Καὶ δεύτερος ἐκράτησε Μελχίας υἱὸς Ἡρὰμ, καὶ Ἀσοὺβ υἱὸς Φαὰτ Μωὰβ, καὶ ἕως πύργου τῶν θανουρίμ.
\vs{12}Καὶ ἐπὶ χεῖρα αὐτοῦ ἐκράτησε Σαλλοὺμ υἱὸς Ἀλλωῆς, ἄρχων ἡμίσους περιχώρου Ἱερουσαλὴμ, αὐτὸς καὶ αἱ θυγατέρες αὐτοῦ.

\vs{13}Τὴν πύλην τῆς φάραγγος ἐκράτησαν Ἀνοὺν καὶ οἱ κατοικοῦντες Ζανώ· αὐτοὶ ᾠκοδόμησαν αὐτὴν, καὶ ἔστησαν θύρας αὐτῆς καὶ κλεῖθρα αὐτῆς καὶ μοχλοὺς αὐτῆς, καὶ χιλίους πήχεις ἐν τῷ τείχει ἕως τῆς πύλης τῆς κοπρίας.

\vs{14}Καὶ τὴν πύλην τῆς κοπρίας ἐκράτησε Μελχία υἱὸς Ῥηχὰβ, ἄρχων περιχώρου Βηθακχαρὶμ, αὐτὸς καὶ υἱοὶ αὐτοῦ· καὶ ἐσκέπασαν αὐτὴν, καὶ ἔστησαν θύρας αὐτῆς καὶ κλεῖθρα αὐτῆς καὶ μοχλοὺς αὐτῆς.

\vs{15}Τὴν δὲ πύλην τῆς πηγῆς ἠσφαλίσατο Σαλωμὼν υἱὸς Χολεζὲ, ἄρχων μέρους τῆς Μασφά· αὐτὸς ἐζῳκοδόμησεν αὐτὴν καὶ ἐστέγασεν αὐτὴν, καὶ ἔστησε τὰς θύρας αὐτῆς καὶ μοχλοὺς αὐτῆς· καὶ τὸ τεῖχος κολυμβήθρας τῶν κωδίων τῇ κουρᾷ τοῦ βασιλέως, καὶ ἕως τῶν κλιμάκων τῶν καταβαινουσῶν ἀπὸ πόλεως Δαυίδ.
\vs{16}Ὀπίσω αὐτοῦ ἐκράτησε Νεεμίας υἱὸς Ἀζαβοὺχ, ἄρχων ἡμίσους περιχώρου Βηθσοὺρ, ἕως κήπου τάφου Δαυὶδ, καὶ ἕως τῆς κολυμβήθρας τῆς γεγονυίας, καὶ ἕως βηθαγγαρίμ.
\vs{17}Ὀπίσω αὐτοῦ ἐκράτησαν οἱ Λευῖται, Ῥαοὺμ υἱὸς Βανί· ἐπὶ χεῖρα αὐτοῦ ἐκράτησεν Ἀσαβία ἄρχων ἡμίσους περιχώρου Κεϊλὰ τῷ περιχώρῳ αὐτοῦ.
\vs{18}Καὶ μετʼ αὐτὸν ἐκράτησαν ἀδελφοὶ αὐτῶν Βενεῒ υἱὸς Ἠναδὰδ, ἄρχων ἡμίσους περιχώρου Κεϊλά.
\vs{19}Καὶ ἐκράτησεν ἐπὶ χεῖρα αὐτοῦ Ἀζοὺρ υἱὸς Ἰησοῦ, ἄρχων τοῦ Μασφαὶ, μέτρον δεύτερον πύργου ἀναβάσεως τῆς συναπτούσης τῆς γωνίας.
\vs{20}Μετʼ αὐτὸν ἐκράτησε Βαροὺχ υἱὸς Ζαβοῦ, μέτρον δεύτερον ἀπὸ τῆς γωνίας ἕως θύρας Βηθελιασοὺβ τοῦ ἱερέως τοῦ μεγάλου.
\vs{21}Μετʼ αὐτὸν ἐκράτησε Μεραμὼθ υἱὸς Οὐρία, υἱοῦ Ἀκκὼς, μέτρον δεύτερον ἀπὸ θύρας Βηθελιασοὺβ ἕως ἐκλείψεως Βηθελιασούβ.
\vs{22}Καὶ μετʼ αὐτὸν ἐκράτησαν οἱ ἱερεῖς ἄνδρες Ἐκχεχάρ.
\vs{23}Καὶ μετʼ αὐτὸν ἐκράτησε Βενιαμὶν καὶ Ἀσοὺβ κατέναντι οἴκου αὐτῶν· καὶ μετʼ αὐτὸν ἐκράτησεν Ἀζαρίας υἱὸς Μαασίου, υἱοῦ Ἀνανία ἐχόμενα οἴκου αὐτοῦ.
\vs{24}Μετʼ αὐτὸν ἐκράτησε Βανὶ υἱὸς Ἀδὰδ, μέτρον δεύτερον ἀπὸ
\vs{25}Βηθαζαρία ἕως τῆς γωνίας, καὶ ἕως τῆς καμπῆς Φαλὰχ υἱοῦ Εὐζαῒ ἐξεναντίας τῆς γωνίας, καὶ ὁ πύργος ὁ ἐξέχων ἐκ τοῦ οἴκου τοῦ βασιλέως ὁ ἀνώτερος ὁ τῆς αὐλῆς τῆς φυλακῆς· καὶ μετʼ αὐτὸν Φαδαΐα υἱὸς Φόρος.
\vs{26}Καὶ οἱ Ναθινὶμ ἦσαν οἰκοῦντες ἐν τῷ Ὠφὰλ, ἕως κήπου πύλης τοῦ ὕδατος εἰς ἀνατολὰς, καὶ ὁ πύργος ὁ ἐξέχων.

\vs{27}Καὶ μετʼ αὐτὸν ἐκράτησαν οἱ Θεκωῒμ, μέτρον δεύτερον ἐξεναντίας τοῦ πύργου τοῦ μεγάλου τοῦ ἐξέχοντος, καὶ ἕως τοῦ τείχους τοῦ Ὀφλά.

\vs{28}Ἀνώτερον πύλης τῶν ἵππων ἐκράτησαν οἱ ἱερεῖς, ἀνὴρ ἐξεναντίας οἴκου ἑαυτοῦ.
\vs{29}Καὶ μετʼ αὐτὸν ἐκράτησε Σαδδοὺκ υἱὸς Ἐμμὴρ ἐξεναντίας οἴκου ἑαυτοῦ· καὶ μετʼ αὐτὸν ἐκράτησε Σαμαΐα υἱὸς Σεχενία φύλαξ τῆς πύλης τῆς ἀνατολῆς.
\vs{30}Μετʼ αὐτὸν ἐκράτησεν Ἀνανία υἱὸς Σελεμία, καὶ Ἀνὼμ υἱὸς Σελὲφ ὁ ἕκτος, μέτρον δεύτερον· μετʼ αὐτὸν ἐκράτησε Μεσουλὰμ υἱὸς Βαραχία ἐξεναντίας γαζοφυλακίου αὐτοῦ.
\vs{31}Μετʼ αὐτὸν ἐκράτησε Μελχία υἱὸς τοῦ Σαρεφὶ ἕως Βηθὰν Ναθινὶμ, καὶ οἱ ῥωποπῶλαι ἀπέναντι πύλης τοῦ Μαφεκὰδ καὶ ἕως ἀναβάσεως τῆς καμπῆς.
\vs{32}Καὶ ἀναμέσον τῆς πύλης τῆς προβατικῆς ἐκράτησαν οἱ χαλκεῖς καὶ οἱ ῥωποπῶλαι.

\vs{33}Καὶ ἐγένετο ἡνίκα ἤκουσε Σαναβαλλὰτ, ὅτι ἡμεῖς οἰκοδομοῦμεν τὸ τεῖχος, καὶ πονηρὸν αὐτῷ ἐφάνη, καὶ ὠργίσθη ἐπὶ πολὺ, καὶ ἐξεγέλα ἐπὶ τοῖς Ἰουδαίοις.
\vs{34}Καὶ εἶπεν ἐνώπιον τῶν ἀδελφῶν αὐτοῦ, αὕτη ἡ δύναμις Σομόρων, ὅτι οἱ Ἰουδαῖοι οὗτοι οἰκοδομοῦσι τὴν ἑαυτῶν πόλιν; ἆρα θυσιάζουσιν; ἆρα δυνήσονται; καὶ σήμερον ἰάσονται τοὺς λίθους, μετὰ τὸ χῶμα γενέσθαι γῆς καυθέντας;
\vs{35}Καὶ Τωβίας ὁ Ἀμμανίτης ἐχόμενα αὐτοῦ ἦλθε, καὶ εἶπε πρὸς αὐτοὺς, μὴ θυσιάζουσιν ἢ φάγονται ἐπὶ τοῦ τόπου αὐτῶν; οὐχὶ ἀναβήσεται ἀλώπηξ καὶ καθελεῖ τὸ τεῖχος λιθων αὐτῶν;

\vs{36}Ἄκουσον ὁ Θεὸς ἡμῶν, ὅτι ἐγενήθημεν εἰς μυκτηρισμὸν, καὶ ἐπίστρεψον ὀνειδισμὸν αὐτῶν εἰς κεφαλὴν αὐτῶν, καὶ δὸς αὐτοὺς εἰς μυκτηρισμὸν ἐν γῇ αἰχμαλωσίας,
\vs{37}καὶ μὴ καλύψῃς ἐπὶ ἀνομίαν.

\ch{4}
Καὶ ἐγένετο ὡς ἤκουσε Σαναβαλλὰτ καὶ Τωβία καὶ οἱ Ἄραβες καὶ οἱ Ἀμμανίται, ὅτι ἀνέβη ἡ φυὴ τοῖς τείχεσιν Ἱερουσαλὴμ, ὅτι ἤρξαντο αἱ διασφαγαὶ ἀναφράσσεσθαι, καὶ πονηρὸν αὐτοῖς ἐφάνη σφόδρα.
\vs{2}Καὶ συνήχθησαν πάντες ἐπιτοαυτὸ, ἐλθεῖν παρατάξασθαι ἐν Ἱερουσαλὴμ καὶ ποιῆσαι αὐτὴν ἀφανῆ.
\vs{3}Καὶ προσηυξάμεθα πρὸς τὸν Θεὸν ἡμῶν, καὶ ἐστήσαμεν προφύλακας ἐπʼ αὐτοὺς ἡμέρας καὶ νυκτὸς ἀπὸ προσώπου αὐτῶν.
\vs{4}Καὶ εἶπεν Ἰούδας, συνετρίβη ἡ ἰσχὺς τῶν ἐχθρῶν, καὶ ὁ χοῦς πολὺς, καὶ ἡμεῖς οὐ δυνησόμεθα οἰκοδομεῖν ἐν τῷ τείχει.
\vs{5}Καὶ εἶπαν οἱ θλίβοντες ἡμᾶς, οὐ γνώσονται καὶ οὐκ ὄψονται, ἕως ὅτου ἔλθωμεν εἰς μέσον αὐτῶν, καὶ φονεύσωμεν αὐτοὺς καὶ καταπαύσωμεν τὸ ἔργον.

\vs{6}Καὶ ἐγένετο ὡς ἤλθοσαν οἱ Ἰουδαῖοι οἱ οἰκοῦντες ἐχόμενα αὐτῶν, καὶ εἴποσαν ἡμῖν, ἀναβαίνουσιν ἐκ πάντων τῶν τόπων ἐφʼ ἡμᾶς.
\vs{7}Καὶ ἔστησα εἰς τὰ κατώτατα τοῦ τόπου κατόπισθεν τοῦ τείχους ἐν τοῖς σκεπεινοῖς, καὶ ἔστησα τὸν λαὸν κατὰ δήμους μετὰ ῥομφαιῶν αὐτῶν, λόγχας αὐτῶν, καὶ τόξα αὐτῶν.
\vs{8}Καὶ εἶδον καὶ ἀνέστην, καὶ εἶπα πρὸς τοὺς ἐντίμους καὶ πρὸς τοὺς στρατηγοὺς καὶ πρὸς τοὺς καταλοίπους τοῦ λαοῦ, μὴ φοβηθῆτε ἀπὸ προσώπου αὐτῶν, μνήσθητε τοῦ Θεοῦ ἡμῶν τοῦ μεγάλου καὶ φοβεροῦ, καὶ παρατάξασθε περὶ τῶν ἀδελφῶν ὑμῶν, υἱῶν ὑμῶν, θυγατέρων ὑμῶν, γυναικῶν ὑμῶν, καὶ οἴκων ὑμῶν.

\vs{9}Καὶ ἐγένετο ἡνίκα ἤκουσαν οἱ ἐχθροὶ ἡμῶν ὅτι ἐγνώσθη ἡμῖν, καὶ διεσκέδασεν ὁ Θεὸς τὴν βουλὴν αὐτῶν· καὶ ἐπεστρέψαμεν πάντες ἡμεῖς εἰς τὸ τεῖχος, ἀνὴρ εἰς τὸ ἔργον αὐτοῦ.
\vs{10}Καὶ ἐγένετο ἀπὸ τῆς ἡμέρας ἐκείνης ἥμισυ τῶν ἐκτετιναγμένων ἐποίουν τὸ ἔργον, καὶ ἥμισυ αὐτῶν ἀντείχοντο, καὶ λόγχαι καὶ θυρεοὶ καὶ τόξα καὶ θώρακες καὶ οἱ ἄρχοντες ὀπίσω παντὸς οἴκου Ἰούδα τῶν οἰκοδομούντων ἐν τῷ τείχει,
\vs{11}καὶ οἱ αἴροντες ἐν τοῖς ἀρτῆρσιν ἐν ὅπλοις· ἐν μιᾷ χειρὶ ἐποίει αὐτοῦ τὸ ἔργον, καὶ ἐν μιᾷ ἐκράτει τὴν βολίδα.
\vs{12}Καὶ οἱ οἰκοδόμοι ἀνὴρ ῥομφαίαν αὐτοῦ ἐζωσμένος ἐπὶ τὴν ὀσφῦν αὐτοῦ, καὶ ᾠκοδομοῦσαν· καὶ ὁ σαλπίζων ἐν τῇ κερατίνῃ ἐχόμενα αὐτοῦ.
\vs{13}Καὶ εἶπα πρὸς τοὺς ἐντίμους καὶ πρὸς τοὺς ἄρχοντας καὶ πρὸς τοὺς καταλοίπους τοῦ λαοῦ, τὸ ἔργον πλατὺ καὶ πολὺ, καὶ ἡμεῖς σκορπιζόμεθα ἐπὶ τοῦ τείχους μακρὰν ἀνὴρ ἀπὸ τοῦ ἀδελφοῦ αὐτοῦ.
\vs{14}Ἐν τόπῳ οὗ ἐὰν ἀκούσητε τὴν φωνὴν τῆς κερατίνης, ἐκεῖ συναχθήσεσθε πρὸς ἡμᾶς, καὶ ὁ Θεὸς ἡμῶν πολεμήσει περὶ ἡμῶν.

\vs{15}Καὶ ἡμεῖς ποιοῦντες τὸ ἔργον, καὶ ἥμισυ αὐτῶν κρατοῦντες τὰς λόγχας ἀπὸ ἀναβάσεως τοῦ ὄρθρου ἕως ἐξόδου τῶν ἄστρων.
\vs{16}Καὶ ἐν τῷ καιρῷ ἐκείνῳ εἶπα τῷ λαῷ, ἕκαστος μετὰ τοῦ νεανίσκου αὐτοῦ αὐλίσθητε ἐν μέσῳ Ἱερουσαλήμ· καὶ ἔστω ὑμῖν ἡ νὺξ προφυλακὴ, καὶ ἡ ἡμέρα ἔργον.
\vs{17}Καὶ ἤμην ἐγὼ καὶ οἱ ἄνδρες τῆς προφυλακῆς ὀπίσω μου, καὶ οὐκ ἦν ἐξ ἡμῶν ἐκδιδυσκόμενος ἀνὴρ τὰ ἱμάτια αὐτοῦ.

\ch{5}
Καὶ ἡ κραυγὴ τοῦ λαοῦ καὶ γυναικῶν αὐτῶν μεγάλη πρὸς τοὺς ἀδελφοὺς αὐτῶν τοὺς Ἰουδαίους.
\vs{2}Καὶ ἦσάν τινες λέγοντες, ἐν υἱοῖς ἡμῶν καὶ ἐν θυγατράσιν ἡμῶν ἡμεῖς πολλοὶ, καὶ λημψόμεθα σῖτον καὶ φαγόμεθα καὶ ζησόμεθα.
\vs{3}Καὶ εἰσί τινες λέγοντες, ἀγροὶ ἡμῶν καὶ ἀμπελῶνες ἡμῶν καὶ οἰκίαι ἡμῶν, ἡμεῖς διεγγυῶμεν καὶ ληψόμεθα σῖτον καὶ φαγόμεθα.
\vs{4}Καὶ εἰσί τινες λέγοντες, ἐδανεισάμεθα ἀργύριον εἰς φόρους τοῦ βασιλέως, ἀγροὶ ἡμῶν καὶ ἀμπελῶνες ἡμῶν καὶ οἰκίαι ἡμῶν.
\vs{5}Καὶ νῦν ὡς σὰρξ ἀδελφῶν ἡμῶν, σὰρξ ἡμῶν· ὡς υἱοὶ αὐτῶν, υἱοὶ ἡμῶν· καὶ ἰδοὺ ἡμεῖς καταδυναστεύομεν τοὺς υἱοὺς ἡμῶν καὶ τὰς θυγατέρας ἡμῶν εἰς δούλους, καὶ εἰσὶν ἀπὸ θυγατέρων ἡμῶν καταδυναστευόμεναι, καὶ οὐκ ἔστι δύναμις χειρῶν ἡμῶν, καὶ ἀγροὶ ἡμῶν καὶ ἀμπελῶνες ἡμῶν τοῖς ἐντίμοις.

\vs{6}Καὶ ἐλυπήθην σφόδρα καθὼς ἤκουσα τὴν κραυγὴν αὐτῶν καὶ τοὺς λόγους τούτους.
\vs{7}Καὶ ἐβουλεύσατο καρδία μου ἐπʼ ἐμέ· καὶ ἐμαχεσάμην πρὸς τους ἐντίμους καὶ τοὺς ἄρχοντας, καὶ εἶπα αὐτοῖς, ἀπαιτήσει ἀνὴρ τὸν ἀδελφὸν αὐτοῦ, ἃ ὑμεῖς ἀπαιτεῖτε; καὶ ἔδωκα ἐπʼ αὐτοὺς ἐκκλησίαν μεγάλην,
\vs{8}καὶ εἶπα αὐτοῖς, ἡμεῖς κεκτήμεθα τοὺς ἀδελφοὺς ἡμῶν τοὺς Ἰουδαίους τοὺς πωλουμένους τοῖς ἔθνεσιν ἐν ἑκουσίῳ ἡμῶν· καὶ ὑμεῖς πωλεῖτε τοὺς ἀδελφοὺς ὑμῶν, καὶ παραδοθήσονται ἡμῖν; καὶ ἡσύχασαν, καὶ οὐχ εὕροσαν λόγον.
\vs{9}Καὶ εἶπα, οὐκ ἀγαθὸς ὁ λόγος ὃν ὑμεῖς ποιεῖτε· οὐχ οὕτως ἐν φόβῳ Θεοῦ ἡμῶν ἀπελεύσεσθε ἀπὸ ὀνειδισμοῦ τῶν ἐθνῶν τῶν ἐχθρῶν ἡμῶν.
\vs{10}Καὶ οἱ ἀδελφοί μου καὶ οἱ γνωστοί μου καὶ ἐγὼ ἐθήκαμεν αὐτοῖς ἀργύριον καὶ σῖτον· ἐγκατελίπωμεν δὴ τὴν ἀπαίτησιν ταύτην.
\vs{11}Ἐπιστρέψατε δὴ αὐτοῖς ὡς σήμερον ἀγροὺς αὐτῶν καὶ ἀμπελῶνας αὐτῶν καὶ ἐλαιῶνας αὐτῶν καὶ οἰκίας αὐτῶν, καὶ ἀπὸ τοῦ ἀργυρίου τὸν σῖτον καὶ τὸν οἶνον καὶ τὸ ἔλαιον ἐξενέγκατε ἑαυτοῖς.
\vs{12}Καὶ εἶπαν, ἀποδώσομεν, καὶ παρʼ αὐτῶν οὐ ζητήσομεν, οὕτως ποιήσομεν καθὼς σὺ λέγεις· καὶ ἐκάλεσα τοὺς ἱερεῖς καὶ ὥρκισα αὐτοὺς ποιῆσαι ὡς τὸ ῥῆμα τοῦτο.

\vs{13}Καὶ τὴν ἀναβολήν μου ἐξετίναξα, καὶ εἶπα, οὕτως ἐκτινάξαι ὁ Θεὸς πάντα ἄνδρα, ὃς οὐ στήσει τὸν λόγον τοῦτον, ἐκ τοῦ οἴκου αὐτοῦ καὶ ἐκ κόπου αὐτοῦ, καὶ ἔσται οὕτως ἐκτετιναγμένος καὶ κενός· καὶ εἶπε πᾶσα ἡ ἐκκλησία, ἀμην· καὶ ᾔνεσαν τὸν Κύριον· καὶ ἐποίησεν ὁ λαὸς τὸ ῥῆμα τοῦτο.

\vs{14}Ἀπὸ ἡμέρας ἧς ἐνετείλατό μοι εἶναι εἰς ἄρχοντα αὐτῶν ἐν γῇ Ἰούδα, ἀπὸ ἔτους εἰκοστοῦ καὶ ἕως ἔτους τριακοστοῦ καὶ δευτέρου τῷ Ἀρθασασθὰ ἔτη δώδεκα, ἐγὼ καὶ οἱ ἀδελφοί μου βίαν αὐτῶν οὐκ ἔφαγον.
\vs{15}Καὶ τὰς βίας τὰς πρώτας ἃς πρὸ ἐμοῦ ἐβάρυναν ἐπʼ αὐτοὺς, καὶ ἐλάβοσαν παρʼ αὐτῶν ἐν ἄρτοις καὶ ἐν οἴνῳ ἔσχατον ἀργύριον δίδραχμα τεσσαράκοντα· καὶ οἱ ἐκτετιναγμένοι αὐτῶν ἐξουσιάζονται ἐπὶ τὸν λαόν· κᾀγὼ οὐκ ἐποίησα οὕτως ἀπὸ προσώπου φόβου Θεοῦ.
\vs{16}Καὶ ἐν ἔργῳ τοῦ τείχους τούτων οὐκ ἐκράτησα, ἀγρὸν οὐκ ἐκτησάμην, καὶ πάντες οἱ συνηγμένοι ἐκεῖ ἐπὶ τὸ ἔργον.
\vs{17}Καὶ οἱ Ἰουδαῖοι ἑκατὸν καὶ πεντήκοντα ἄνδρες, καὶ ἐρχόμενοι πρὸς ἡμᾶς ἀπὸ τῶν ἐθνῶν τῶν κύκλῳ ἡμῶν ἐπὶ τράπεζάν μου.
\vs{18}Καὶ ἦν γινόμενον εἰς ἡμέραν μίαν μόσχος εἷς, καὶ πρόβατα ἓξ ἐκλεκτὰ καὶ χίμαρος ἐγίνοντό μοι· καὶ ἀναμέσον δέκα ἡμερῶν ἐν πᾶσιν οἶνος τῷ πλήθει· καὶ σὺν τούτοις ἄρτους τῆς βίας οὐκ ἐζήτησα, ὅτι βαρεῖα ἡ δουλεία ἐπὶ τὸν λαὸν τοῦτον.

\vs{19}Μνήσθητί μου ὁ Θεὸς εἰς ἀγαθὸν πάντα ὅσα ἐποίησα τῷ λαῷ τούτῳ.

\ch{6}
Καὶ ἐγένετο καθὼς ἠκούσθη τῷ Σαναβαλλὰτ, καὶ Τωβίᾳ, καὶ τῷ Γησὰμ τῷ Ἀραβὶ, καὶ τοῖς καταλοίποις ἐχθρῶν ἡμῶν, ὅτι ᾠκοδόμησα τὸ τεῖχος, καὶ οὐ κατελείφθη ἐν αὐτοῖς πνοή· ἕως τοῦ καιροῦ ἐκείνου θύρας οὐκ ἐπέστησα ἐν ταῖς πύλαις.
\vs{2}Καὶ ἀπέστειλε Σαναβαλλὰτ καὶ Γησὰμ πρὸς μὲ, λέγων, δεῦρο καὶ συναχθῶμεν ἐπιτοαυτὸ ἐν ταῖς κώμαις ἐν πεδίῳ Ὠνώ· καὶ αὐτοὶ λογιζόμενοι ποιῆσαί μοι πονηρίαν.
\vs{3}Καὶ ἀπέστειλα ἐπʼ αὐτοὺς ἀγγέλους, λέγων, ἔργον μέγα ἐγὼ ποιῶ, καὶ οὐ δυνήσομαι καταβῆναι, μή ποτε καταπαύσῃ τὸ ἔργον· ὡς ἂν τελειώσω αὐτὸ, καταβήσομαι πρὸς ὑμᾶς.
\vs{4}Καὶ ἀπέστειλαν πρὸς μὲ ὡς τὸ ῥῆμα τοῦτο· καὶ ἀπέστειλα αὐτοῖς κατὰ ταῦτα.

\vs{5}Καὶ ἀπέστειλε πρὸς μὲ Σαναβαλλὰτ τὸν παῖδα αὐτοῦ, καὶ ἐπιστολὴν ἀνεῳγμένην ἐν χειρὶ αὐτοῦ.
\vs{6}Καὶ ἦν γεγραμμένον ἐν αὐτῇ, ἐν ἔθνεσιν ἠκούσθη ὅτι σὺ καὶ οἱ Ἰουδαῖοι λογίζεσθε ἀποστατῆσαι, διὰ τοῦτο σὺ οἰκοδομεῖς τὸ τεῖχος, καὶ σὺ ἔσῃ αὐτοῖς εἰς βασιλέα.
\vs{7}Καὶ πρὸς τούτοις προφήτας ἔστησας σεαυτῷ, ἵνα καθίσῃς ἐν Ἰερουσαλὴμ εἰς βασιλέα ἐπὶ Ἰούδα· καὶ νῦν ἀπαγγελήσονται τῷ βασιλεῖ οἱ λόγοι οὗτοι· καὶ νῦν δεῦρο βουλευσώμεθα ἐπιτοαυτό.
\vs{8}Καὶ ἀπέστειλα πρὸς αὐτὸν, λέγων, οὐκ ἐγενήθη ὡς οἱ λόγοι οὗτοι ὧς σὺ λέγεις, ὅτι ἀπὸ καρδίας σου σὺ ψεύδῃ αὐτούς.
\vs{9}Ὅτι πάντες φοβερίζουσιν ἡμᾶς, λέγοντες, ἐκλυθήσονται χεῖρες αὐτῶν ἀπὸ τοῦ ἔργου τούτου, καὶ οὐ ποιηθήσεται· καὶ νῦν ἐκραταίωσα τὰς χεῖράς μου.

\vs{10}Κᾀγὼ εἰσῆλθον εἰς οἶκον Σεμεῒ υἱοῦ Δαλαΐα υἱοῦ Μεταβεὴλ, καὶ αὐτὸς συνεχόμενος· καὶ εἶπε, συναχθῶμεν εἰς οἶκον τοῦ Θεοῦ ἐν μέσῳ αὐτοῦ, καὶ κλείσωμεν τὰς θύρας αὐτοῦ, ὅτι ἔρχονται νυκτὸς φονεῦσαί σε.
\vs{11}Καὶ εἶπα, τίς ἐστιν ὁ ἀνὴρ ὃς εἰσελεύσεται εἰς τὸν οἶκον, καὶ ζήσεται;
\vs{12}Καὶ ἐπέγνων, καὶ ἰδοὺ ὁ Θεὸς οὐκ ἀπέστειλεν αὐτὸν, ὅτι ἡ προφητεία λόγος κατʼ ἐμοῦ·
\vs{13}καὶ Τωβίας καὶ Σαναβαλλὰτ ἐμισθώσαντο ἐπʼ ἐμὲ ὄχλον ὅπως φοβηθῶ, καὶ ποιήσω οὕτως, καὶ ἁμάρτω, καὶ γένωμαι αὐτοῖς εἰς ὄνομα πονηρὸν, ὅπως ὀνειδίσωσί με.

\vs{14}Μνήσθητι ὁ Θεὸς Τωβίᾳ καὶ Σαναβαλλάτ, ὡς τὰ ποιήματα αὐτοῦ ταῦτα, καὶ τῷ Νωαδίᾳ τῷ προφήτῃ, καὶ καταλοίποις τῶν προφητῶν οἳ ἦσαν φοβερίζοντές με.

\vs{15}Καὶ ἐτελέσθη τὸ τεῖχος πέμπτῃ καὶ εἰκάδι τοῦ Ἐλοὺλ μηνὸς εἰς πεντήκοντα καὶ δύο ἡμέρας.
\vs{16}Καὶ ἐγένετο ἡνίκα ἤκουσαν πάντες οἱ ἐχθροὶ ἡμῶν, καὶ ἐφοβήθησαν πάντα τὰ ἔθνη τὰ κύκλῳ ἡμῶν, καὶ ἐπέπεσε φόβος σφόδρα ἐν ὀφθαλμοῖς αὐτῶν, καὶ ἔγνωσαν ὅτι παρὰ τοῦ Θεοῦ ἡμῶν ἐγενήθη τελειωθῆναι τὸ ἔργον τοῦτο.

\vs{17}καὶ ἐν ταῖς ἡμέραις ἐκείναις ἀπὸ πολλῶν ἐντίμων Ἰούδα ἐπιστολαὶ ἐπορεύοντο πρὸς Τωβίαν, καὶ αἱ Τωβία ἤρχοντο πρὸς αὐτούς·
\vs{18}Ὅτι πολλοὶ ἐν Ἰούδα ἔνορκοι ἦσαν αὐτῷ, ὅτι γαμβρὸς ἦν τοῦ Σεχενία υἱοῦ Ἡραέ· καὶ Ἰωνὰν υἱὸς αὐτοῦ ἔλαβε τὴν θυγατέρα Μεσουλὰμ υἱοῦ Βαραχία εἰς γυναῖκα.
\vs{19}Καὶ τοὺς λόγους αὐτοῦ ἦσαν λέγοντες πρὸς μὲ, καὶ λόγους μου ἦσαν ἐκφέροντες αὐτῷ· καὶ ἐπιστολὰς ἀπέστειλε Τωβίας φοβερίσαι με.

\ch{7}
Καὶ ἐγένετο ἡνίκα ᾠκοδομήθη τὸ τεῖχος, καὶ ἔστησα τὰς θύρας, καὶ ἐπεσκέπησαν οἱ πυλωροὶ, καὶ οἱ ᾄδοντες, καὶ οἱ Λευῖται,
\vs{2}καὶ ἐνετειλάμην τῷ Ἀνανίᾳ ἀδελφῷ μου, καὶ τῷ Ανανίᾳ ἄρχοντι τῆς βιρὰ ἐν Ἱερουσαλὴμ, ὅτι αὐτὸς ὡς ἀνὴρ ἀληθὴς καὶ φοβούμενος τὸν Θεὸν παρὰ πολλούς.
\vs{3}Καὶ εἶπα αὐτοῖς, οὐκ ἀνοιγήσονται πύλαι Ἱερουσαλὴμ ἕως ἅμα τῷ ἡλίῳ· καὶ ἔτι αὐτῶν γρηγορούντων, κλειέσθωσαν αἱ θύραι, καὶ σφηνούσθωσαν· καὶ στῆσον προφύλακας οἰκούντων ἐν Ἱερουσαλὴμ, ἀνὴρ ἐν προφυλακῇ αὐτοῦ, καὶ ἀνὴρ ἀπέναντι οἰκίας αὐτοῦ.

\vs{4}καὶ ἡ πόλις πλατεῖα καὶ μεγάλη, καὶ ὁ λαὸς ὀλίγος ἐν αὐτῇ, καὶ οὐκ ἦσαν οἰκίαι ᾠκοδομημέναι.
\vs{5}Καὶ ἔδωκεν ὁ Θεὸς εἰς τὴν καρδίαν μου, καὶ συνῆξα τοὺς ἐντίμους καὶ τοὺς ἄρχοντας καὶ τὸν λαὸν εἰς συνοδίας· καὶ εὗρον βιβλίον τῆς συνοδίας οἳ ἀνέβησαν ἐν πρώτοις· καὶ εὗρον γεγραμμένον ἐν αὐτῷ,

\vs{6}Καὶ οὗτοι υἱοὶ τῆς χώρας οἱ ἀναβάντες ἀπὸ αἰχμαλωσίας τῆς ἀποικίας ἧς ἀπῴκισε Ναβουχοδονόσορ βασιλεὺς Βαβυλῶνος, καὶ ἐπέστρεψεν εἰς Ἱερουσαλὴμ καὶ εἰς Ἰούδα ἀνὴρ εἰς τὴν πόλιν αὐτοῦ μετὰ Ζοροβάβελ,
\vs{7}καὶ Ἰησοῦ, καὶ Νεεμία, Ἀζαρία, καὶ Ῥεελμὰ, Ναεμανὶ, Μαρδοχαῖος, Βαλσὰν, Μασφαρὰθ, Ἔσδρα, Βογουΐα, Ἰναοὺμ, Βαανὰ, Μασφὰρ, ἄνδρες λαοῦ Ἰσραήλ.

\vs{8}γἱοὶ Φόρος, δισχίλιοι ἑκατὸν ἑβδομηκονταδύο.

\vs{9}γἱοὶ Σαφατία, τριακόσιοι ἑβδομηκονταδύο.

\vs{10}γἱοὶ Ἠρὰ, ἑξακόσιοι πεντηκονταδύο.

\vs{11}γἱοὶ Φαὰθ Μωὰβ τοῖς υἱοῖς Ἰησοῦ καὶ Ἰωὰβ, δισχίλιοι ἑξακόσιοι δεκαοκτώ.

\vs{12}γἱοὶ Αἰλὰμ, χίλιοι διακόσιοι πεντηκοντατέσσαρες.

\vs{13}γἱοὶ Ζαθουΐα, ὀκτακόσιοι τεσσαρακονταπέντε.

\vs{14}γἱοὶ Ζακχού, ἑπτακόσιοι ἑξήκοντα.

\vs{15}γἱοὶ Βανουῒ, ἑξακόσιοι τεσσαρακονταοκτώ.

\vs{16}γἱοὶ Βηβὶ, ἑξακόσιοι εἰκοσιοκτώ.

\vs{17}γἱοὶ Ἀσγὰδ, δισχίλιοι τριακόσιοι εἰκοσιδύο.

\vs{18}γἱοὶ Ἀδωνικὰμ, ἑξακόσιοι ἑξηκονταεπτά.

\vs{19}γἱοὶ Βαγοῒ, δισχίλιοι ἑξηκονταεπτά.

\vs{20}γἱοὶ Ἠδὶν, ἑξακόσιοι πεντηκονταπέντε.

\vs{21}γἱοὶ Ἀτὴρ τῷ Ἐζεκίᾳ, ἐννενηκονταοκτώ.

\vs{22}γἱοὶ Ἠσὰμ, τριακόσιοι εἰκοσιοκτώ.

\vs{23}γἱοὶ Βεσεῒ, τριακόσιοι εἰκοσιτέσσαρες.

\vs{24}γἱοὶ Ἀρὶφ, ἑκατὸν δώδεκα· υἱοὶ Ἀσὲν, διακόσιοι εἰκοσιτρεῖς.

\vs{25}γἱοὶ Γαβαὼν, ἐννενηκονταπέντε.

\vs{26}γἱοὶ Βαιθαλὲμ, ἑκατὸν εἰκοσιτρεῖς· υἱοὶ Ἀτωφὰ, πεντηκονταέξ.

\vs{27}γἱοὶ Ἀναθὼθ, ἑκατὸν εἰκοσιοκτώ.

\vs{28}Ἄνδρες Βηθασμὼθ, τεσσαρακονταδύο.

\vs{29}Ἄνδρες Καριαθαρὶμ, Καφιρὰ, καὶ Βηρὼθ, ἑπτακόσιοι τεσσαρακοντατρεῖς.

\vs{30}Ἄνδρες Ἀραμὰ, καὶ Γαβαὰ, ἑξακόσιοι εἴκοσι.

\vs{31}Ἄνδρες Μαχεμὰς, ἑκατὸν εἰκοσιδύο.

\vs{32}Ἄνδρες Βαιθὴλ καὶ Ἀῒ, ἑκατὸν εἰκοσιτρεῖς.

\vs{33}Ἄνδρες Ναβία, ἐκατὸν πεντηκονταδύο.

\vs{34}Ἄνδρες Ἠλαμαὰρ, χίλιοι διακόσιοι πεντηκονταδύο.

\vs{35}γἱοὶ Ἠρὰμ, τριακόσιοι εἴκοσι.

\vs{36}γἱοὶ Ἱεριχὼ, τριακόσιοι τεσσαρακονταπέντε.

\vs{37}γἱοὶ Λοδαδὶδ καὶ Ὠνὼ, ἑπτακόσιοι εἰκοσιεῖς.

\vs{38}γἱοὶ Σανανὰ, τρισχίλιοι ἐννακόσιοι τριάκοντα.

\vs{39}Οἱ ἱερεῖς υἱοὶ Ἰωδαὲ εἰς οἶκον Ἰησοῦ, ἐννακόσιοι ἑβδομηκοντατρεῖς.

\vs{40}Υἱοὶ Ἐμμὴρ, χίλιοι πεντηκονταδύο.

\vs{41}Υἱοὶ Φασεοὺρ, χίλιοι διακόσιοι τεσσαρακονταεπτά.

\vs{42}γἱοὶ Ἠρὰμ, χίλιοι δεκαεπτά.

\vs{43}Οἱ Λευῖται, υἱοὶ Ἰησοῦ τοῦ Καδμιὴλ τοῖς υἱοῖς τοῦ Οὐδουΐα, ἑβδομηκοντατέσσαρες.

\vs{44}Οἱ ᾄδοντες, υἱοὶ Ἀσὰφ, ἑκατὸν τεσσαρακονταοκτώ.

\vs{45}Οἱ πυλωροὶ, υἱοὶ Σαλοὺμ, υἱοὶ Ἀτὴρ, υἱοὶ Τελμὼν, υἱοὶ Ἀκοὺβ, υἱοὶ Ἀτιτὰ, υἱοὶ Σαβὶ, ἑκατὸν τριακονταοκτώ.

\vs{46}Οἱ Ναθινὶμ, υἱοὶ Σηὰ, υἱοὶ Ἀσφὰ, υἱοὶ Ταβαὼθ,
\vs{47}υἱοὶ Κιρὰς, υἱοὶ Ἀσουΐα, υἱοὶ Φαδὼν,
\vs{48}υἱοὶ Λαβανὰ, υἱοὶ Ἀγαβὰ, υἱοὶ Σελμεῒ,
\vs{49}υἱοὶ Ἀνὰν, υἱοὶ Γαδὴλ, υἱοὶ Γαὰρ,
\vs{50}υἱοὶ Ῥααΐα, υἱοὶ Ῥασσὼν, υἱοὶ Νεκωδὰ,
\vs{51}υἱοὶ Γηζὰμ, υἱοὶ Ὀζὶ, υἱοὶ Φεσὴ,
\vs{52}υἱοὶ Βησὶ, υἱοὶ Μεϊνὼν, υἱοὶ Νεφωσασὶ,
\vs{53}υἱοὶ Βακβοὺκ, υἱοὶ Ἀχιφὰ, υἱοὶ Ἀροὺρ,
\vs{54}υἱοὶ Βασαλὼθ, υἱοὶ Μιδὰ, υἱοὶ Ἀδασὰν,
\vs{55}υἱοὶ Βαρκουὲ, υἱοὶ Σισαρὰθ, υἱοὶ Θημὰ,
\vs{56}υἱοὶ Νισιὰ, υἱοὶ Ἀτιφά.
\vs{57}γἱοὶ δούλων Σαλωμὼν, υἱοὶ Σουτεῒ, υἱοὶ Σαφαρὰτ, υἱοὶ Φεριδὰ,
\vs{58}υἱοὶ Ἰελὴλ, υἱοὶ Δορκὼν, υἱοὶ Γαδαὴλ,
\vs{59}υἱοὶ Σαφατία, υἱοὶ Ἐττὴλ, υἱοὶ Φακαρὰθ, υἱοὶ Σαβαῒμ, υἱοὶ Ἠμίμ.
\vs{60}Πάντες οἱ Ναθινὶμ, καὶ υἱοὶ δούλων Σαλωμὼν, τριακόσιοι ἐννενηκονταδύο.

\vs{61}Καὶ οὗτοι ἀνέβησαν ἀπὸ Θελμελὲθ, Θελαρησὰ, Χαροὺβ, Ἠρὼν, Ἰεμὴρ, καὶ οὐκ ἐδυνάσθησαν ἀπαγγεῖλαι οἴκους πατριῶν αὐτῶν καὶ σπέρμα αὐτῶν, εἰ ἀπὸ Ἰσραὴλ εἰσίν·
\vs{62}γἱοὶ Δαλαία, υἱοὶ Τωβία, υἱοὶ Νεκωδὰ, ἑξακόσιοι τεσσαρακονταδύο.

\vs{63}Καὶ ἀπὸ τῶν ἱερέων, υἱοὶ Ἐβία, υἱοὶ Ἀκὼς, υἱοὶ Βερζελλὶ, ὅτι ἔλαβον ἀπὸ θυγατέρων Βερζελλὶ τοῦ Γαλααδίτου γυναῖκας, καὶ ἐκλήθησαν ἐπʼ ὀνόματι αὐτῶν.
\vs{64}Οὗτοι ἐζήτησαν γραφὴν αὐτῶν τῆς συνοδίας, καὶ οὐχ εὑρέθη· καὶ ἠγχιστεύθησαν ἀπὸ τῆς ἱερατείας.
\vs{65}Καὶ εἶπεν ἀθερσασθὰ, ἵνα μὴ φάγωσιν ἀπὸ τοῦ ἁγίου τῶν ἁγίων, ἕως ἀναστῇ ἱερεὺς φωτίσων.

\vs{66}Καὶ ἐγένετο πᾶσα ἡ ἐκκλησία ὡσεὶ τέσσαρες μυριάδες δισχίλιοι τριακόσιοι ἐξήκοντα,
\vs{67}πάρεκ δούλων αὐτῶν καὶ παιδισκῶν αὐτῶν· οὗτοι ἑπτακισχίλιοι τριακόσιοι τριακονταεπτά· καὶ ᾄδοντες καὶ ᾄδουσαι, διακόσιοι τεσσαρακονταπέντε.
\vs{69}Ὄνοι δισχίλιοι ἑπτακόσιοι.

\vs{70}Καὶ ἀπὸ μέρους ἀρχηγῶν τῶν πατριῶν ἔδωκαν εἰς τὸ ἐργον τῷ Νεεμίᾳ εἰς θησαυρὸν χρυσοῦς χιλίους, φιάλας πεντήκοντα, καὶ χωθωνὼθ τῶν ἱερέων τριάκοντα.
\vs{71}Καὶ ἀπὸ ἀρχηγῶν τῶν πατριῶν ἔδωκαν εἰς θησαυρούς τοῦ ἔργου χρυσοῦ νομίσματος δύο μυριάδας, καὶ ἀργυρίου μνᾶς δισχιλίας τριακοσίας.
\vs{72}Καὶ ἔδωκαν οἱ κατάλοιποι τοῦ λαοῦ χρυσίου δύο μυριάδας, καὶ ἀργυρίου μνᾶς δισχιλίας διακοσίας, καὶ χωθωνὼθ τῶν ἱερέων ἑξηκονταεπτά.

\vs{73}Καὶ ἐκάθισαν οἱ ἱερεῖς, καὶ Λευῖται, καὶ οἱ πυλωροὶ, καὶ οἱ ᾄδοντες, καὶ οἱ ἀπὸ τοῦ λαοῦ, καὶ οἱ Ναθινὶμ, καὶ πᾶς Ἰσραὴλ ἐν πόλεσιν αὐτῶν.

\ch{8}
Καὶ ἔφθασεν ὁ μὴν ὁ ἕβδομος, καὶ οἱ υἱοὶ Ἰσραὴλ ἐν πόλεσιν αὐτῶν· καὶ συνήχθησαν πᾶς ὁ λαὸς ὡς ἀνὴρ εἷς εἰς τὸ πλάτος τὸ ἔμπροσθεν πύλης τοῦ ὕδατος· καὶ εἶπαν τῷ Ἔσδρᾳ τῷ γραμματεῖ, ἐνέγκαι τὸ βιβλίον νόμου Μωυσῆ, ὃν ἐνετείλατο Κύριος τῷ Ἰσραήλ.
\vs{2}Καὶ ἤνεγκεν Ἔσδρας ὁ ἱερεὺς τὸν νόμον ἐνώπιον τῆς ἐκκλησίας ἀπὸ ἀνδρὸς ἕως γυναικὸς, καὶ πᾶς ὁ συνιὼν, ἀκούειν ἐν ἡμέρᾳ μιᾷ τοῦ μηνὸς τοῦ ἑβδόμου.
\vs{3}Καὶ ἀνέγνω ἐν αὐτῷ ἀπὸ τῆς ὥρας τοῦ διαφωτίσαι τὸν ἥλιον ἕως ἡμίσους τῆς ἡμέρας, ἀπέναντι τῶν ἀνδρῶν καὶ τῶν γυναικῶν, καὶ αὐτοὶ συνιέντες· καὶ ὦτα παντὸς τοῦ λαοῦ εἰς τὸ βιβλίον τοῦ νόμου.
\vs{4}Καὶ ἔστη Ἔσδρας ὁ γραμματεὺς ἐπὶ βήματος ξυλίνου, καὶ ἔστησαν ἐχόμενα αὐτοῦ Ματθαθίας, καὶ Σαμαΐας, καὶ Ἀνανίας, καὶ Οὐρίας, καὶ Χελκία, καὶ Μαασία ἐκ δεξιῶν αὐτοῦ, καὶ ἐξ ἀριστερῶν Φαδαΐας, καὶ Μισαὴλ, καὶ Μελχίας, καὶ Ἀσὼμ, καὶ Ἀσαβαδμὰ, καὶ Ζαχαρίας, καὶ Μεσολλάμ.
\vs{5}Καὶ ἤνοιξεν Ἔσδρας τὸ βιβλίον ἐνώπιον παντὸς τοῦ λαοῦ, ὅτι αὐτὸς ἦν ἐπάνω τοῦ λαοῦ· καὶ ἐγένετο ἡνίκα ἤνοιξεν αὐτὸ, ἔστη πᾶς ὁ λαός.
\vs{6}Καὶ ηὐλόγησεν Ἔσδρας Κύριον τὸν Θεὸν τὸν μέγαν· καὶ ἀπεκρίθη πᾶς ὁ λαὸς, καὶ εἶπαν, ἀμὴν, ἐπάραντες τὰς χεῖρας αὐτῶν· καὶ ἔκυψαν καὶ προσεκύνησαν τῷ Κυρίῳ ἐπὶ πρόσωπον ἐπὶ τὴν γῆν.
\vs{7}Καὶ Ἰησοῦς, καὶ Βαναΐας, καὶ Σαραβίας ἦσαν συνετίζοντες τὸν λαὸν εἰς τὸν νόμον· καὶ ὁ λαὸς ἐν τῇ στάσει αὐτοῦ.
\vs{8}Καὶ ἀνέγνωσαν ἐν βιβλίῳ νόμου τοῦ Θεοῦ, καὶ ἐδίδασκεν Ἔσδρας, καὶ διέστελλεν ἐν ἐπιστήμῃ Κυρίου, καὶ συνῆκεν ὁ λαὸς ἐν τῇ ἀναγνώσει.

\vs{9}Καὶ εἶπε Νεεμίας καὶ Ἔσδρας ὁ ἱερεὺς καὶ γραμματεὺς, καὶ οἱ Λευῖται, καὶ οἱ συνετίζοντες τὸν λαὸν, καὶ εἶπαν παντὶ τῷ λαῷ, ἡμέρα ἁγία ἐστὶ τῷ Κυρίῳ Θεῷ ἡμῶν, μὴ πενθεῖτε μηδὲ κλαίετε· ὅτι ἔκλαιε πᾶς ὁ λαὸς ὡς ἤκουσαν τοὺς λόγους τοῦ νόμου.
\vs{10}Καὶ εἶπεν αὐτοῖς, πορεύεσθε, φάγετε λιπάσματα, καὶ πίετε γλυκάσματα, καὶ ἀποστείλατε μερίδας τοῖς μὴ ἔχουσιν, ὅτι ἁγία ἐστὶν ἡ ἡμέρα τῷ Κυρίῳ ἡμῶν· καὶ μὴ διαπέσητε, ὅτι ἐστὶ Κύριος ἰσχὺς ἡμῶν.
\vs{11}Καὶ οἱ Λευῖται κατεσιώπων πάντα τὸν λαὸν, λέγοντες, σιωπᾶτε, ὅτι ἡμέρα ἁγία, καὶ μὴ καταπίπτετε.
\vs{12}Καὶ ἀπῆλθε πᾶς ὁ λαὸς φαγεῖν καὶ πιεῖν, καὶ ἀποστέλλειν μερίδας, καὶ ποιῆσαι εὐφροσύνην μεγάλην, ὅτι συνῆκαν ἐν τοῖς λόγοις οἷς ἐγνώρισεν αὐτοῖς.

\vs{13}Καὶ ἐν τῇ ἡμέρᾳ τῇ δευτέρᾳ συνήχθησαν οἱ ἄρχοντες τῶν πατριῶν σὺν τῷ παντὶ λαῷ, οἱ ἱερεῖς καὶ οἱ Λευῖται πρὸς Ἔσδραν τὸν γραμματέα, ἐπιστῆσαι πρὸς πάντας τοὺς λόγους τοῦ νόμου.
\vs{14}Καὶ εὕροσαν γεγραμμένον ἐν τῷ νόμῳ, ᾧ ἐνετείλατο Κύριος τῷ Μωυσῇ, ὅπως κατοικήσωσιν οἱ υἱοὶ Ἰσραὴλ ἐν σκῆναις ἐν ἑορτῇ ἐν μηνὶ τῷ ἑβδόμῳ,
\vs{15}καὶ ὅπως σημάνωσι σάλπιγξιν ἐν πάσαις ταῖς πόλεσιν αὐτῶν καὶ ἐν Ἱερουσαλήμ· καὶ εἶπεν Ἔσδρας, ἐξέλθετε εἰς τὸ ὄρος, καὶ ἐνέγκατε φύλλα ἐλαίας, καὶ φύλλα ξύλων κυπαρισσίνων, καὶ φύλλα μυρσίνης, καὶ φύλλα φοινίκων, καὶ φύλλα ξύλου δασέος, ποιῆσαι σκηνὰς κατὰ τὸ γεγραμμένον.
\vs{16}Καὶ ἐξῆλθεν ὁ λαὸς, καὶ ἤνεγκαν, καὶ ἐποίησαν ἑαυτοῖς σκηνὰς ἀνὴρ ἐπὶ τοῦ δώματος αὐτοῦ, καὶ ἐν ταῖς αὐλαῖς αὐτῶν, καὶ ἐν ταῖς αὐλαῖς οἴκου τοῦ Θεοῦ, καὶ ἐν πλατείαις τῆς πόλεως, καὶ ἕως πύλης Ἐφραίμ.
\vs{17}Καὶ ἐποίησαν πᾶσα ἡ ἐκκλησία, οἱ ἐπιστρέψαντες ἀπὸ τῆς αἰχμαλωσίας, σκηνὰς, καὶ ἐκάθισαν ἐν σκηναῖς· ὅτι οὐκ ἐποίησαν ἀπὸ ἡμερῶν Ἰησοῦ υἱοῦ Ναυῆ οὕτως οἱ υἱοὶ Ἰσραὴλ ἕως τῆς ἡμέρας ἐκείνης· καὶ ἐγένετο εὐφροσύνη μεγάλη.

\vs{18}Καὶ ἀνέγνω ἐν βιβλίῳ νόμου τοῦ Θεοῦ ἡμέραν ἐν ἡμέρᾳ ἀπὸ τῆς ἡμέρας τῆς πρώτης ἕως τῆς ἡμέρας τῆς ἐσχάτης· καὶ ἐποίησαν ἑορτὴν ἑπτὰ ἡμέρας, καὶ τῇ ἡμέρᾳ τῇ ὀγδόῃ ἐξόδιον κατὰ τὸ κρίμα.

\ch{9}
Καὶ ἐν ἡμέρᾳ εἰκοστῇ καὶ τετάρτῃ τοῦ μηνὸς τούτου συνήχθησαν οἱ υἱοὶ Ἰσραὴλ ἐν νηστείᾳ καὶ ἐν σάκκοις καὶ σποδῷ ἐπὶ κεφαλῆς αὐτῶν.
\vs{2}Καὶ ἐχωρίσθησαν οἱ υἱοὶ Ἰσραὴλ ἀπὸ παντὸς υἱοῦ ἀλλοτρίου, καὶ ἔστησαν καὶ ἐξηγόρευσαν τὰς ἁμαρτίας αὐτῶν, καὶ τὰς ἀνομίας τῶν πατέρων αὐτῶν.
\vs{3}Καὶ ἔστησαν ἐπὶ τῇ στάσει αὐτῶν, καὶ ἀνέγνωσαν ἐν βιβλίῳ νόμου Κυρίου Θεοῦ αὐτῶν· καὶ ἦσαν ἐξαγορεύοντες τῷ Κυρίῳ καὶ προσκυνοῦντες τῷ Κυρίῳ Θεῷ αὐτῶν.
\vs{4}Καὶ ἔστη ἐπὶ ἀναβάσει τῶν Λευιτῶν Ἰησοῦς, καὶ οἱ υἱοὶ Καδμιὴλ, Σεχενία υἱὸς Σαραβία, υἱοὶ Χωνενί· καὶ ἐβόησαν φωνῇ μεγάλῃ πρὸς Κύριον τὸν Θεὸν αὐτῶν.
\vs{5}Καὶ εἴποσαν οἱ Λευῖται Ἰησοῦς καὶ Καδμιὴλ, ἀνάστητε, εὐλογεῖτε Κύριον τὸν Θεὸν ἡμῶν ἀπὸ τοῦ αἰῶνος καὶ ἕως τοῦ αἰῶνος· καὶ εὐλογήσουσιν ὄνομα δόξης σου, καὶ ὑψώσουσιν ἐπὶ πάσῃ εὐλογίᾳ καὶ αἰνέσει.

\vs{6}Καὶ εἶπεν Ἔσδρας, σὺ εἶ αὐτὸς Κύριος μόνος, σὺ ἐποίησας τὸν οὐρανὸν καὶ τὸν οὐρανὸν τοῦ οὐρανοῦ, καὶ πᾶσαν τὴν στάσιν αὐτῶν, τὴν γῆν καὶ πάντα ὅσα ἐστὶν ἐν αὐτῇ, τὰς θαλάσσας καὶ πάντα τὰ ἐν αὐταῖς· καὶ σὺ ζωοποιεῖς τὰ πάντα, καὶ σοὶ προσκυνοῦσιν αἱ στρατιαὶ τῶν οὐρανῶν·

\vs{7}Σὺ εἶ Κύριος ὁ Θεὸς, σὺ ἐξελέξω ἐν Ἅβραμ καὶ ἐξήγαγες αὐτὸν ἐκ τῆς χώρας τῶν Χαλδαίων, καὶ ἐπέθηκας αὐτῷ ὄνομα Ἁβραάμ·
\vs{8}Καὶ εὗρες τὴν καρδίαν αὐτοῦ πιστὴν ἐνώπιόν σου, καὶ διέθου πρὸς αὐτὸν διαθήκην δοῦναι αὐτῷ τὴν γῆν τῶν Χαναναίων, καὶ Χετταίων, καὶ Ἀμοῤῥαίων, καὶ Φερεζαίων, καὶ Ἰεβουσαίων, καὶ Γεργεσαίων, καὶ τῷ σπέρματι αὐτοῦ· καὶ ἔστησας τοὺς λόγους σου, ὅτι δίκαιος σύ.

\vs{9}Καὶ εἶδες τὴν ταπείνωσιν τῶν πατέρων ἡμῶν ἐν Αἰγύπτῳ, καὶ τὴν κραυγὴν αὐτῶν ἤκουσας ἐπὶ θάλασσαν ἐρυθράν.
\vs{10}Καὶ ἔδωκας σημεῖα καὶ τέρατα ἐν Αἰγύπτῳ ἐν Φαραῷ, καὶ ἐν πᾶσι τοῖς παισὶν αὐτοῦ, καὶ ἐν παντὶ τῷ λαῷ τῆς γῆς αὐτοῦ, ὅτς ἔγνως ὅτι ὑπερηφάνησαν ἐπʼ αὐτοὺς καὶ ἐποίησας σεαυτῷ ὄνομα ὡς ἡ ἡμέρα αὕτη.
\vs{11}Καὶ τὴν θάλασσαν ἔῥῥηξας ἐνώπιον αὐτῶν, καὶ παρήλθοσαν ἐν μέσῳ τῆς θαλάσσης ἐν ξηρασίᾳ, καὶ τοὺς καταδιώξοντας αὐτοὺς ἔῤῥιψας εἰς βυθὸν, ὡσεὶ λίθον ἐν ὕδατι σφοδρῶ.

\vs{12}Καὶ ἐν στύλῳ νεφέλης ὡδήγησας αὐτοὺς ἡμέρας, καὶ ἐν στύλῳ πυρὸς τὴν νύκτα, τοῦ φωτίσαι αὐτοῖς τὴν ὁδὸν ἐν ᾗ πορεύσονται ἐν αὐτῇ.
\vs{13}καὶ ἐπὶ ὄρος Σινὰ κατέβης, καὶ ἐλάλησας πρὸς αὐτοὺς ἐξ οὐρανοῦ, καὶ ἔδωκας αὐτοῖς κρίματα εὐθέα, καὶ νόμους ἀληθείας, προστάγματα, καὶ ἐντολὰς ἀγαθάς.
\vs{14}Καὶ τὸ σάββατόν σου τὸ ἅγιον ἐγνώρισας αὐτοῖς, ἐντολὰς καὶ προστάγματα καὶ νόμον ἐνετείλω αὐτοῖς ἐν χειρὶ Μωυσῆ δούλου σου.
\vs{15}Καὶ ἄρτον ἐξ οὐρανοῦ ἔδωκας αὐτοῖς εἰς σιτοδοτίαν αὐτῶν, καὶ ὕδωρ ἐκ πέτρας ἐξήνεγκας αὐτοῖς εἰς δίψαν αὐτῶν· καὶ εἶπας αὐτοῖς εἰσελθεῖν κληρονομῆσαι τὴν γῆν ἐφʼ ἣν ἐξέτεινας τὴν χεῖρά σου δοῦναι αὐτοῖς.

\vs{16}Καὶ αὐτοὶ καὶ οἱ πατέρες ἡμῶν ὑπερηφανεύσαντο, καὶ ἐσκλήρυναν τὸν τράχηλον αὐτῶν, καὶ οὐκ ἤκουσαν τῶν ἐντολῶν σου,
\vs{17}καὶ ἀνένευσαν τοῦ εἰσακοῦσαι, καὶ οὐκ ἐμνήσθησαν τῶν θαυμασίων σου ὧν ἐποίησας μετʼ αὐτῶν· καὶ ἐσκλήρυναν τὸν τράχηλον αὐτῶν, καὶ ἔδωκαν ἀρχὴν ἐπιστρέψαι εἰς δουλείαν αὐτῶν ἐν Αἰγύπτῳ· καὶ σὺ ὁ Θεὸς ἐλεήμων καὶ οἰκτίρμων, μακρόθυμος καὶ πολυέλεος, καὶ οὐκ ἐγκατέλιπες αὐτούς.
\vs{18}Ἔτι δὲ καὶ ἐποίησαν ἑαυτοῖς μόσχον χωνευτὸν, καὶ εἶπαν, οὗτοι οἱ θεοὶ οἱ ἐξαγαγόντες ἡμᾶς ἐξ Αἰγύπτου· καὶ ἐποίησαν παροργισμοὺς μεγάλους.

\vs{19}Καὶ σὺ ἐν οἰκτιρμοῖς σου τοῖς πολλοῖς οὐκ ἐγκατέλιπες αὐτοὺς ἐν τῇ ἐρήμῳ, τὸν στύλον τῆς νεφέλης οὐκ ἐξέκλινας ἀπʼ αὐτῶν ἡμέρας, ὁδηγῆσαι αὐτοὺς ἐν τῇ ὁδῷ, καὶ τὸν στύλον τοῦ πυρὸς τὴν νύκτα, φωτίζειν αὐτοῖς τὴν ὁδὸν ἐν ᾗ πορεύσονται ἐν αὐτῇ.
\vs{20}Καὶ τὸ πνεῦμά σου τὸ ἀγαθὸν ἔδωκας συνετίσαι αὐτούς· καὶ τὸ μάννα σου οὐκ ἀφυστέρησας ἀπὸ στόματος αὐτῶν, καὶ ὕδωρ ἔδωκας αὐτοῖς ἐν τῷ δίψει αὐτῶν.
\vs{21}Καὶ τεσσαράκοντα ἔτη διέθρεψας αὐτοὺς ἐν τῇ ἐρήμῳ, οὐχ ὑστέρησας αὐτοῖς οὐδέν· ἱμάτια αὐτῶν οὐκ ἐπαλαιώθησαν, καὶ πόδες αὐτῶν οὐ διεῤῥάγησαν.

\vs{22}Καὶ ἔδωκας αὐτοῖς βασιλείας, καὶ λαοὺς ἐμέρισας αὐτοῖς· καὶ ἐκληρονόμησαν τὴν γῆν Σηὼν βασιλέως Ἐσεβὼν, καὶ τὴν γῆν Ὢγ βασιλέως τοῦ Βασάν.
\vs{23}Καὶ τοὺς υἱοὺς αὐτῶν ἐπλήθυνας ὡς τοὺς ἀστέρας τοῦ οὐρανοῦ, καὶ εἰσήγαγες αὐτοὺς εἰς τὴν γῆν ἣν εἶπας τοῖς πατράσιν αὐτῶν, καὶ ἐκληρονόμησαν αὐτήν·
\vs{24}καὶ ἐξέτριψας ἐνώπιον αὐτῶν τοὺς κατοικοῦντας τὴν γῆν τῶν Χαναναίων, καὶ ἔδωκας αὐτοὺς εἰς τὰς χεῖρας αὐτῶν καὶ τοὺς βασιλεῖς αὐτῶν καὶ τοὺς λαοὺς τῆς γῆς, ποιῆσαι αὐτοῖς ὡς ἀρεστὸν ἐνώπιον αὐτῶν.
\vs{25}Καὶ κατελάβοσαν πόλεις ὑψηλὰς, καὶ ἐκληρονόμησαν οἰκίας πλήρεις πάντων ἀγαθῶν, λάκκους λελατομημένους, ἀμπελῶνας καὶ ἐλαιῶνας, καὶ πᾶν ξύλον βρώσιμον εἰς πλῆθος· καὶ ἐφάγοσαν καὶ ἐνεπλήσθησαν καὶ ἐλιπάνθησαν, καὶ ἐτρύφησαν ἐν ἀγαθωσύνῃ σου τῇ μεγάλῃ.

\vs{26}Καὶ ἤλλαξαν, καὶ ἀπέστησαν ἀπὸ σοῦ, καὶ ἔῤῥιψαν τὸν νόμον σου ὀπίσω σώματος αὐτῶν· καὶ τοὺς προφήτας σου ἀπέκτειναν, οἳ διεμαρτύραντο ἐν αὐτοῖς ἐπιστρέψαι αὐτοὺς πρὸς σέ· καὶ ἐποίησαν παροργισμοὺς μεγάλους.
\vs{27}Καὶ ἔδωκας αὐτοὺς ἐν χειρὶ θλιβόντων αὐτούς, καὶ ἔθλιψαν αὐτούς· καὶ ἀνεβόησαν πρὸς σὲ ἐν καιρῷ θλίψεως αὐτῶν, καὶ σὺ ἐξ οὐρανοῦ σου ἤκουσας, καὶ ἐν οἰκτιρμοῖς σου τοῖς μεγάλοις ἔδωκας αὐτοῖς σωτῆρας, καὶ ἔσωσας αὐτοὺς ἐκ χειρὸς θλιβόντων αὐτούς.

\vs{28}Καὶ ὡς ἀνεπαύσαντο, ἐπέστρεψαν ποιῆσαι τὸ πονηρὸν ἐνώπιόν σου· καὶ ἐγκατέλιπες αὐτοὺς εἰς χεῖρας ἐχθρῶν αὐτῶν, καὶ κατῆρξαν ἐν αὐτοῖς· καὶ πάλιν ἀνεβόησαν πρὸς σὲ, καὶ σὺ ἐξ οὐρανοῦ εἰσήκουσας, καὶ ἐῤῥύσω αὐτοὺς ἐν οἰκτιρμοῖς σου πολλοῖς.
\vs{29}Καὶ ἐπεμαρτύρω αὐτοῖς ἐπιστρέψαι αὐτοὺς εἰς τὸν νόμον σου· καὶ οὐκ ἤκουσαν, ἀλλʼ ἐν ταῖς ἐντολαῖς σου καὶ κρίμασί σου ἡμάρτοσαν, ἃ ποιήσας αὐτὰ ἄνθρωπος ζήσεται ἐν αὐτοῖς· καὶ ἔδωκαν νῶτον ἀπειθοῦντα, καὶ τράχηλον αὐτῶν ἐσκλήρυναν καὶ οὐκ ἤκουσαν.
\vs{30}Καὶ εἵλκυσας ἐπʼ αὐτοὺς ἔτη πολλὰ, καὶ ἐπεμαρτύρω αὐτοῖς ἐν πνεύματί σου ἐν χειρὶ προφητῶν σου, καὶ οὐκ ἐνωτίσαντο, καὶ ἔδωκας αὐτοὺς ἐν χειρὶ λαῶν τῆς γῆς.
\vs{31}Καὶ σὺ ἐν οἰκτιρμοῖς σου τοῖς πολλοῖς οὐκ ἐποίησας αὐτοὺς εἰς συντέλειαν, καὶ οὐκ ἐγκατέλιπες αὐτοὺς, ὅτι ἰσχυρὸς εἶ καὶ ἐλεήμων καὶ οἰκτίρμων.

\vs{32}Καὶ νῦν ὁ Θεὸς ἡμῶν ὁ ἰσχυρὸς ὁ μέγας ὁ κραταιὸς καὶ ὁ φοβερὸς, φυλάσσων τὴν διαθήκην σου καὶ τὸ ἔλεός σου, μὴ ὀλιγωθήτω ἐνώπιόν σου πᾶς ὁ μόχθος ὃς εὗρεν ἡμᾶς, καὶ τοὺς βασιλεῖς ἡμῶν, καὶ τοὺς ἄρχοντας ἡμῶν, καὶ τοὺς ἱερεῖς ἡμῶν, καὶ τοὺς προφήτας ἡμῶν, καὶ τοὺς πατέρας ἡμῶν, καὶ ἐν παντὶ τῷ λαῷ σου ἀπὸ ἡμερῶν βασιλέων Ἀσσοὺρ καὶ ἕως τῆς ἡμέρας ταύτης.
\vs{33}Καὶ σὺ δίκαιος ἐπὶ πᾶσι τοῖς ἐρχομένοις ἐφʼ ἡμᾶς, ὅτι ἀλήθειαν ἐποίησας· καὶ ἡμεῖς ἐξημάρτομεν,
\vs{34}καὶ οἱ βασιλεῖς ἡμῶν, καὶ οἱ ἄρχοντες ἡμῶν, καὶ οἱ ἱερεῖς ἡμῶν, καὶ οἱ πατέρες ἡμῶν οὐκ ἐποίησαν τὸν νόμον σου, καὶ οὐ προσέσχον τῶν ἐντολῶν σου, καὶ τὰ μαρτύριά σου ἃ διεμαρτύρω αὐτοῖς.
\vs{35}Καὶ αὐτοὶ ἐν βασιλείᾳ σου καὶ ἐν ἀγαθωσύνῃ σου τῇ πολλῇ ᾗ ἔδωκας αὐτοῖς, καὶ ἐν τῇ γῇ τῇ πλατείᾳ καὶ λιπαρᾷ ᾗ ἔδωκας ἐνώπιον αὐτῶν, οὐκ ἐδούλευσάν σοι, καὶ οὐκ ἀπέστρεψαν ἀπὸ ἐπιτηδευμάτων αὐτῶν τῶν πονηρῶν.
\vs{36}Ἰδοὺ σήμερον ἐσμὲν δοῦλοι, καὶ ἡ γῆ ἣν ἔδωκας τοῖς πατράσιν ἡμῶν φαγεῖν τὸν καρπὸν αὐτῆς καὶ τὰ ἀγαθὰ αὐτῆς, ἰδοὺ ἐσμὲν δοῦλοι ἐπʼ αὐτῆς, καὶ οἱ καρποὶ αὐτῆς πολλοὶ
\vs{37}τοῖς βασιλεῦσιν οἷς ἔδωκας ἐφʼ ἡμᾶς ἐν ἁμαρτίαις ἡμῶν, καὶ ἐπὶ τὰ σώματα ἡμῶν ἐξουσιάζουσι, καὶ ἐν κτήνεσιν ἡμῶν, ὡς ἀρεστὸν αὐτοῖς, καὶ ἐν θλίψει μεγάλῃ ἐσμέν.

\ch{10}
Καὶ ἐν πᾶσι τούτοις ἡμεῖς διατιθέμεθα πίστιν, καὶ γράφομεν, καὶ ἐπισφραγίζουσιν ἄρχοντες ἡμῶν, Αευῖται ἡμῶν, ἱερεῖς ἡμῶν.

\vs{2}Καὶ ἐπὶ τῶν σφραγιζόντων Νεεμίας ἀρτασασθὰ υἱὸς Ἀχαλία, καὶ Σεδεκίας
\vs{3}υἱὸς Ἀραία, καὶ Ἀζαρία, καὶ Ἱερεμία,
\vs{4}Φασοὺρ, Ἀμαρία, Μελχία,
\vs{5}Ἀττοὺς, Σεβανὶ, Μαλοὺχ,
\vs{6}Ἰρὰμ, Μεραμὼθ, Ἀβδία,
\vs{7}Δανιὴλ, Γανναθὼν, Βαροὺχ,
\vs{8}Μεσουλὰμ, Ἀβία, Μιαμὶν,
\vs{9}Μααζία, Βελγαῒ, Σαμαΐα οὗτοι ἱερεῖς.

\vs{10}Καὶ οἱ Λευῖται, Ἰησοῦς υἱὸς Ἀζανία, Βαναίου ἀπὸ υἱῶν Ἠναδὰδ, Καδμιὴλ
\vs{11}καὶ οἱ ἀδελφοὶ αὐτοῦ, Σαβανία, Ὠδουΐα, Καλιτὰν,
\vs{12}Φελία, Ανὰν, Μιχὰ, Ῥοὼβ,
\vs{13}Ἀσεβίας, Ζακχὼρ, Σαραβία, Σεβανία,
\vs{14}Ὠδούμ· υἱοὶ Βανουαῒ

\vs{15}Ἄρχοντες τοῦ λαοῦ Φόρος, Φαὰθ Μωὰβ, Ἠλὰμ, Ζαθουΐα·
\vs{16}υἱοὶ Βανὶ, Ἀσγὰδ, Βηβαῒ,
\vs{17}Ἀδανία, Βαγοῒ, Ἡδὶν,
\vs{18}Ἀτὴρ, Ἐζεκία, Ἀζοὺρ,
\vs{19}Ὠδουΐα, Ἠσὰμ, Βησὶ,
\vs{20}Ἀρὶφ, Ἀναθὼθ, Νωβαῒ,
\vs{21}Μεγαφὴς, Μεσουλλὰμ, Ἠζὶρ,
\vs{22}Μεσωζεβὴλ, Σαδοὺκ, Ἰσδδούα,
\vs{23}Φαλτία, Ἀνὰν, Ἀναΐα,
\vs{24}Ὠσηὲ, Ἀνανία, Ἀσοὺβ,
\vs{25}Ἀλωὴς, Φαλαῒ, Σωβὴκ,
\vs{26}Ῥεοὺμ, Ἐσσαβανὰ, Μαασία,
\vs{27}καὶ Ἀΐα, Αἰνὰν, Ἠνὰμ,
\vs{28}Μαλοὺχ, Ἠρὰμ, Βαανὰ,

\vs{29}Καὶ οἱ κατάλοιποι τοῦ λαοῦ, οἱ ἱερεῖς, οἱ Λευῖται, οἱ πυλωροὶ, οἱ ᾄδοντες, οἱ Ναθινὶμ, καὶ πᾶς ὁ προσπορευόμενος ἀπὸ λαῶν τῆς γῆς πρὸς νόμον τοῦ Θεοῦ, γυναῖκες αὐτῶν, υἱοὶ αὐτῶν, θυγατέρες αὐτῶν· πᾶς ὁ εἰδὼς καὶ συνιὼν,
\vs{30}ἐνίσχυον ἐπὶ τοὺς ἀδελφοὺς αὐτῶν, καὶ κατηράσαντο αὐτοὺς, καὶ εἰσήλθοσαν ἐν ἀρᾷ καὶ ἐν ὅρκῳ τοῦ πορεύεσθαι ἐν νόμῳ τοῦ Θεοῦ, ὃς ἐδόθη ἐν χειρὶ Μωυσῆ δούλου τοῦ Θεοῦ, φυλάσσεσθαι καὶ ποιεῖν πάσας τὰς ἐντολὰς Κυρίου, καὶ τὰ κρίματα αὐτοῦ, καὶ τὰ προστάγματα αὐτοῦ·
\vs{31}Καὶ τοῦ μὴ δοῦναι θυγατέρας ἡμῶν τοῖς λαοῖς τοῖς γῆς, καὶ τὰς θυγατέρας αὐτῶν οὐ ληψόμεθα τοῖς υἱοῖς ἡμῶν·
\vs{32}Καὶ λαοὶ τῆς γῆς οἱ φέροντες τοὺς ἀγορασμοὺς καὶ πᾶσαν πρᾶσιν ἐν ἡμέρᾳ τοῦ σαββάτου ἀποδόσθαι, οὐκ ἀγορῶμεν παρʼ αὐτῶν ἐν σαββάτῳ καὶ ἐν ἡμέρᾳ ἁγίᾳ· καὶ ἀνήσομεν τὸ ἔτος τὸ ἕβδομον, καὶ ἀπαίτησιν πάσης χειρός.

\vs{33}Καὶ στήσομεν ἐφʼ ἡμᾶς ἐντολάς δοῦναι ἐφʼ ἡμᾶς τρίτον τοῦ διδράχμου κατʼ ἐνιαυτὸν εἰς δουλείαν οἴκου τοῦ Θεοῦ ἡμῶν,
\vs{34}εἰς ἄρτους τοῦ προσώπου, καὶ θυσίαν τοῦ ἐνδελεχισμοῦ καὶ εἰς ὁλοκαύτωμα τοῦ ἐνδελεχισμοῦ τῶν σαββάτων, τῶν νουμηνιῶν, εἰς τὰς ἑορτὰς καὶ εἰς τὰ ἅγια, καὶ τὰ περὶ ἁμαρτίας, ἐξιλάσασθαι περὶ Ἰσραὴλ, καὶ εἰς ἔργα οἴκου τοῦ Θεοῦ ἡμῶν.

\vs{35}Καὶ κλήρους ἐβάλομεν περὶ κλήρου ξυλοφορίας, οἱ ἱερεῖς καὶ οἱ Λευῖται καὶ ὁ λαὸς, ἐνέγκαι εἰς οἶκον Θεοῦ ἡμῶν, εἰς οἶκον πατριῶν ἡμῶν, εἰς καιροὺς ἀπὸ χρόνων, ἐνιαυτὸν κατʼ ἐνιαυτὸν, ἐκκαῦσαι ἑπὶ τὸ θυσιαστήριον Κυρίου Θεοῦ ἡμῶν, ὡς γέγραπται ἐν τῷ νόμῳ·
\vs{36}Καὶ ἐνέγκαι τὰ πρωτογενήματα τῆς γῆς ἡμῶν, καὶ πρωτογεννήματα καρποῦ παντὸς ξύλου ἐνιαυτὸν κατʼ ἐνιαυτὸν εἰς οἶκον Κυρίου,
\vs{37}καὶ τὰ πρωτότοκα υἱῶν ἡμῶν καὶ κτηνῶν ἡμῶν, ὡς γέγραπται ἐν τῷ νόμῳ, καὶ τὰ πρωτότοκα τῶν βοῶν ἡμῶν καὶ ποιμνίων ἡμῶν ἐνέγκαι εἰς οἶκον Θεοῦ ἡμῶν, τοῖς ἱερεῦσι τοῖς λειτουργοῦσιν ἐν οἴκῳ Θεοῦ ἡμῶν.
\vs{38}Καὶ τὴν ἀπαρχὴν σίτων ἡμῶν, καὶ τὸν καρπὸν παντὸς ξύλου, οἴνου, καὶ ἐλαίου, οἴσομεν τοῖς ἱερεῦσιν εἰς τὸ γαζοφυλάκιον οἴκου τοῦ Θεοῦ, καὶ δεκάτην γῆς ἡμῶν τοῖς Λευίταις· καὶ αὐτοὶ οἱ Λευῖται δεκατοῦντες ἐν πάσαις πόλεσι δουλείας ἡμῶν.
\vs{39}Καὶ ἔσται ὁ ἱερεὺς υἱὸς Ἀαρὼν μετὰ τοῦ Λευίτου ἐν τῇ δεκάτῃ τοῦ Λευίτου, καὶ οἱ Λευῖται ἀνοίσουσι τὴν δεκάτην τῆς δεκάδος εἰς οἶκον Θεοῦ ἡμῶν εἰς τὰ γαζοφυλάκια εἰς οἶκον τοῦ Θεοῦ.
\vs{40}Ὅτι εἰς τοὺς θησαυροὺς εἰσοίσουσιν οἱ υἱοὶ Ἰσραὴλ καὶ οἱ υἱοὶ τοῦ Λευὶ τὰς ἀπαρχὰς τοῦ σίτου, καὶ τοῦ οἴνου, καὶ τοῦ ἐλαίου, καὶ ἐκεῖ σκεύη τὰ ἅγια, καὶ οἱ ἱερεῖς καὶ οἱ λειτουργοὶ καὶ οἱ πυλωροὶ καὶ οἱ ᾄδοντες· καὶ οὐκ ἐγκαταλείψομεν τὸν οἶκον τοῦ Θεοῦ ἡμῶν.

\ch{11}
Καὶ ἐκάθισαν οἱ ἄρχοντες τοῦ λαοῦ ἐν Ἱερουσαλήμ· καὶ οἱ κατάλοιποι τοῦ λαοῦ ἐβάλοσαν κλήρους ἐνέγκαι ἕνα ἀπὸ τῶν δέκα καθίσαι ἐν Ἱερουσαλὴμ πόλει τῇ ἁγίᾳ, καὶ ἐννέα μέρη ἐν ταῖς πόλεσι.
\vs{2}Καὶ εὐλόγησεν ὁ λαὸς τοὺς πάντας ἄνδρας τοὺς ἑκουσιαζομένους καθίσαι ἐν Ἱερουσαλήμ.

\vs{3}Καὶ οὗτοι οἱ ἄρχοντες τῆς χώρας οἳ ἐκάθισαν ἐν Ἱερουσαλὴμ καὶ ἐν πόλεσιν Ἰούδα· ἐκάθισαν ἀνὴρ ἐν κατασχέσει αὐτοῦ ἐν πόλεσιν αὐτῶν Ἰσραὴλ, οἱ ἱερεῖς, καὶ οἱ Λευῖται, καὶ οἱ Ναθιναοι, καὶ οἱ υἱοὶ δούλων Σαλωμών·

\vs{4}Καὶ ἐν Ἱερουσαλὴμ ἐκάθισαν ἀπὸ υἱῶν Ἰούδα, καὶ ἀπὸ υἱῶν Βενιαμίν· Ἀπὸ υἱῶν Ἰούδα, Ἀθαΐα υἱὸς Ἀζία, υἱὸς Ζαχαρία, υἱὸς Σαμαρία, υἱὸς Σαφατία, υἱὸς Μαλελεήλ· καὶ ἀπὸ τῶν υἱῶν Φαρὲς,
\vs{5}καὶ Μαασία υἱὸς Βαροὺχ, υἱὸς Χαλαζὰ, υἱὸς Ὀζία, υἱὸς Ἀδαΐα, υἱὸς Ἰωαρὶβ, υἱὸς Ζαχαρίου, υἱὸς τοῦ Σηλωνί.
\vs{6}Πάντες υἱοὶ Φαρὲς οἱ καθήμενοι ἐν Ἱερουσαλὴμ, τετρακόσιοι ἑξηκονταοκτὼ ἄνδρες δυνάμεως.
\vs{7}Καὶ οὗτοι υἱοὶ Βενιαμὶν, Σηλὼ υἱὸς Μεσουλὰμ, υἱὸς Ἰωὰδ, υἱὸς Φαδαΐα, υἱὸς Κωλεΐα, υἱὸς Μαασίου, υἱὸς Ἐθιὴλ, υἱὸς Ἰεσία,
\vs{8}καὶ ὀπίσω αὐτοῦ Γηβὲ, Σηλὶ, ἐννακόσιοι εἰκοσιοκτώ.
\vs{9}Καὶ Ἰωὴλ υἱὸς Ζεχρὶ ἐπίσκοπος ἐπʼ αὐτούς· καὶ Ἰούδα υἱὸς Ἀσανὰ ἀπὸ τῆς πόλεως, δεύτερος.

\vs{10}Ἀπὸ τῶν ἱερέων· καὶ Ἰαδία υἱὸς Ἰωαρὶβ, Ἰαχὶν.
\vs{11}Σαραία υἱὸς Ἐλχία, υἱὸς Μεσουλὰμ, υἱὸς Σαδδοὺκ, υἱὸς Μαριὼθ, υἱὸς Αἰτὼθ, ἀπέναντι οἴκου τοῦ Θεοῦ.
\vs{12}Καὶ οἱ ἀδελφοὶ αὐτῶν ποιοῦντες τὸ ἔργον τοῦ οἴκου, ὀκτακόσιοι εἰκοσιδύο· καὶ Ἀδαΐα υἱὸς Ἱεροὰμ, υἱοῦ Φαλαλία, υἱοῦ Ἀμασὶ, υἱὸς Ζαχαρία, υἱὸς Φασσοὺρ, υἱὸς Μελχία,
\vs{13}καὶ ἀδελφοὶ αὐτοῦ ἄρχοντες πατριῶν, διακόσιοι τεσσαρακονταδύο· καὶ Ἀμασία υἱὸς Ἐσδριὴλ, υἱοῦ Μεσαριμὶθ, υἱοῦ Ἐμμὴρ,
\vs{14}καὶ ἀδελφοὶ αὐτοῦ δυνατοὶ παρατάξεως, ἐκατὸν εἰκοσιοκτώ· καὶ ἐπίσκοπος Βαδιὴλ υἱὸς τῶν μεγάλων.

\vs{15}Καὶ ἀπὸ τῶν Λευιτῶν, Σαμαΐα υἱὸς Ἐσρικὰμ, Ματθανίας υἱὸς Μιχὰ,
\vs{17}καὶ Ἰωβὴβ υἱὸς Σαμουῒ,
\vs{18}διακόσιοι ὀγδοηκοντατέσσαρες.

\vs{19}Καὶ οἱ πυλωροὶ, Ἀκοὺβ, Τελαμὶν, καὶ οἱ ἀδελφοὶ αὐτῶν, ἑκατὸν ἑβδομηκονταδύο.

\vs{22}Καὶ ἐπίσκοπος Λευιτῶν υἱὸς Βανὶ, υἱὸς Οζὶ, υἱὸς Ἀσαβία, υἱὸς Μιχά· ἀπὸ υἱῶν Ἀσὰφ, τῶν ᾀδόντων ἀπέναντι ἔργου οἴκου τοῦ Θεοῦ·
\vs{23}ὅτι ἐντολὴ τοῦ βασιλέως εἰς αὐτούς.

\vs{24}Καὶ Φαθαΐα υἱὸς Βασηζὰ πρὸς χεῖρα τοῦ βασιλέως εἰς πᾶν χρῆμα τῷ λαῷ,
\vs{25}καὶ πρὸς τὰς ἐπαύλεις ἐν ἀγρῷ αὐτῶν· καὶ ἀπὸ υἱῶν Ἰούδα ἐκάθισαν ἐν Καριαθαρβὸκ,
\vs{26}καὶ ἐν Ἰησοὺ,
\vs{27}καὶ ἐν Βηρσαβεὲ,
\vs{30}καὶ ἐπαύλεις αὐτῶν Λαχὶς καὶ ἀγροὶ αὐτῆς· καὶ παρενεβάλοσαν ἐν Βηρσαβεέ.
\vs{31}Καὶ οἱ υἱοὶ Βενιαμὶν ἀπὸ Γαβαὰ Μαχμάς.
\vs{36}Καὶ ἀπὸ τῶν Λευιτῶν μερίδες Ἰούδα τῷ Βενιαμίν.

\ch{12}
Καὶ οὗτοι οἱ ἱερεῖς καὶ οἱ Λευῖται οἱ ἀναβάντες μετὰ Ζοροβάβελ υἱοῦ Σαλαθιὴλ καὶ Ἰησοῦ· Σαραΐα, Ἱερεμία, Ἔσδρα,
\vs{2}Ἀμαρία, Μαλοὺχ,
\vs{3}Σεχενία·
\vs{7}Οὗτοι οἱ ἄρχοντες τῶν ἱερέων, καὶ ἀδελφοὶ αὐτῶν ἐν ἡμέραις Ἰησοί·

\vs{8}Καὶ οἱ Λευῖται, Ἰησοῦ, Βανουὶ, Καδμιὴλ, Σαραβία, Ἰωδαὲ, Ματθανία, ἐπὶ τῶν χειρῶν αὐτὸς, καὶ οἱ ἀδελφοὶ αὐτῶν
\vs{9}εἰς τὰς ἐφημερίας.

\vs{10}Καὶ Ἰησοὺς ἐγέννησε τὸν Ἰωακὶμ, καὶ Ἰωακὶμ ἐγέννησε τὸν Ἐλιασὶβ, καὶ Ἐλιασὶβ τὸν Ἰωδαὲ,
\vs{11}καὶ Ἰωδαὲ ἐγέννησε τὸν Ἰωνάθαν, καὶ Ἰωνάθαν ἐγέννησε τὸν Ἰαδού.
\vs{12}Καὶ ἐν ἡμέραις Ἰωακὶμ ἀδελφοὶ αὐτοῦ οἱ ἱερεῖς καὶ οἱ ἄρχοντες τῶν πατριῶν, τῷ Σαραΐα, Ἀμαρία· τῷ Ἱερεμία, Ἀυανία·
\vs{13}Τῷ Ἔσδρᾳ Μεσουλάμ· τῷ Ἀμαρία, Ἰωανάν·
\vs{14}τῷ Ἀμαλοὺχ, Ἰωνάθαν· τῷ Σεχενία, Ἰωσήφ·
\vs{15}Τῷ Ἀρὲ, Μαννάς· τῷ Μαριὼθ, Ἐλκαΐ·
\vs{16}τῷ Ἀδαδαῒ, Ζαχαρία· τῷ Γαναθὼθ, Μεσολάμ·
\vs{17}Τῷ Ἀβιὰ, Ζεχρί· τῷ Μιαμὶν, Μααδαί· τῷ Φελετὶ, τῷ βαλγὰς, Σαμουέ· τῷ Σεμία, Ἰωνάθαν·
\vs{18}Τῷ Ἰωαρὶβ,
\vs{19}Ματθαναΐ τῷ Ἐδίῳ, Ὀζί.
\vs{20}Τῷ Σαλαῒ, καλλαΐ· τῷ Ἀμὲκ, Ἀβέδ·
\vs{21}Τῷ Ἐλκίᾳ, Ἀσαβιας· τῷ Ἰσδεϊοὺ, Ναθαναήλ

\vs{22}Οἱ Λευῖται ἐν ἡμέραις Ἐλιασὶβ, Ἰωαδὰ, καὶ Ἰωὰ, καὶ Ἰωανὰν, καὶ Ἰδούα, γεγραμμένοι ἄρχοντες τῶυ πατριῶν, καὶ οἱ ἱερεῖς ἐν βασιλείᾳ Δαρείου τοῦ Πέρσου.
\vs{23}Υἱοὶ δὲ Λυεὶ ἄρχοντες τῶν πατριῶν γεγραμμένοι ἐπὶ βιβλίῳ λόγων τῶν ἡμερῶν, καὶ ἕως ἡμερῶν Ἰωανὰν υἱοῦ Ἐλισουέ.
\vs{24}Καὶ οἱ ἄρχοντες τῶν Λευιτῶν, Ἀσαβία, καὶ Σαραβία, καὶ Ἰησοῦ· καὶ υἱοὶ Καδμιὴλ καὶ ἀδελφοὶ αὐτῶν κατεναντίον αὐτῶν εἰς ὕμνον αἰνεῖν ἐν ἐντολῇ Δαυὶδ ἀνθρώπου τοῦ Θεοῦ ἐφημερίαν πρὸς ἐφημερίαν.

\vs{25}Ἐν τῷ συναγαγεῖν με τοὺς πυλωροὺς
\vs{26}ἐν ἡμέραις Ἰωακὶμ υἱοῦ Ἰησοῦ, υἱοῦ Ἰωσεδὲκ, καὶ ἐν ἡμέραις Νεεμία, καὶ Ἔσδρας ὁ ἱερεὺς γραμματεύς.

\vs{27}Καὶ ἐν ἐγκαινίοις τείχους Ἱερουσαλὴμ ἐζήτησαν τοὺς Λευίτας ἐν τοῖς τόποις αὐτῶν τοῦ ἐνέγκαι αὐτοὺς εἰς Ἱερουσαλὴμ, ποιῆσαι ἐγκαίνια καὶ εὐφροσύνην ἐν θωδαθὰ, καὶ ἐν ᾠδαῖς κυμβαλίζοντες, καὶ ψαλτήρια, καὶ κινύραι.
\vs{28}Καὶ συνήχθησαν οἱ υἱοὶ τῶν ᾀδόντων καὶ ἀπὸ τῆς περιχώρου κυκλόθεν εἰς Ἱερουσαλὴμ, καὶ ἀπὸ ἐπαύλεων,
\vs{29}καὶ ἀπὸ ἀγρῶν, ὅτι ἐπαύλεις ᾠκοδόμησαν ἑαυτοῖς οἱ ᾄδοντες ἐν Ἰερουσαλήμ.
\vs{30}Καὶ ἐκαθαρίσθησαν οἱ ἱερεῖς καὶ οἱ Λευῖται, καὶ ἐκαθάρισαν τὸν λαὸν καὶ τοὺς πυλωροὺς καὶ τὸ τεῖχος.

\vs{31}Καὶ ἀνήνεγκαν τοὺς ἄρχοντας Ἰούδα ἐπάνω τοῦ τείχους· καὶ ἔστησαν δύο περὶ αἰνέσεως μεγάλους, καὶ διῆλθον ἐκ δεξιῶν ἐπάνω τοῦ τεὶχους τῆς κοπρίας.
\vs{32}Καὶ ἐπορεύθησαν ὀπίσω αὐτῶν Ὠσαΐα καὶ ἥμισυ ἀρχόντων Ἰούδα,
\vs{33}καὶ Ἀζαρίας, καὶ Ἔσδρας, καὶ Μεσολλὰμ,
\vs{34}καὶ Ἰούδα, καὶ Βενιαμὶν, καὶ Σαμαΐας, καὶ Ἱερεμὶα.
\vs{35}Καὶ ἀπὸ τῶν υἱῶν τῶν ἱερέων ἐν σάλπιγξι, Ζαχαρίας υἱὸς Ἰωνὰθαν, υἱὸς Σαμαΐα, υἱὸς Ματθανία, υἱὸς Μιχαία, υἱὸς Ζακχοὺρ, υἱὸς Ἀσάφ·
\vs{36}Καὶ ἀδελφοὶ αὐτοῦ, Σαμαΐα, καὶ Ὀζιὴλ, Γελὼλ, Ἰαμὰ, Ἀΐα, Ναθαναὴλ, καὶ Ἰούδα, Ανανὶ, τοῦ αἰνεῖν ἐν ᾠδαῖς Δαυὶδ ἀνθρώπου τοῦ Θεοῦ· καὶ Ἔσδρας ὁ γραμματεὺς ἔμπροσθεν αὐτῶν ἐπὶ πύλης,
\vs{37}τοῦ αἰνεῖν κατέναντι αὐτῶν· καὶ ἀνέβησαν ἐπὶ κλίμακας πόλεως Δαυὶδ ἐν ἀναβάσει τοῦ τείχους ἐπάνωθεν τοῦ οἴκου Δαυὶδ, καὶ ἕως τῆς πύλης τοῦ ὕδατος
\vs{39}Ἐφραὶμ, καὶ ἐπὶ τὴν πύλην ἰχθυρὰν, καὶ πύργῳ Ἁναμεὴλ, καὶ ἕως πύλης τῆς προβατικῆς.
\vs{42}Καὶ ἠκούσθησαν οἱ ᾄδοντες, καὶ ἐπεσκέπησαν.
\vs{43}Καὶ ἔθυσαν ἐν τῇ ἡμέρᾳ ἐκείνῃ θυσιάσματα μεγάλα, καὶ ηὐφράνθησαν, ὅτι ὁ Θεὸς ηὔφρανεν αὐτοὺς μεγάλως· καὶ αἱ γυναῖκες αὐτῶν καὶ τὰ τέκνα αὐτῶν ηὐφράνθησαν, καὶ ἠκούσθη ἡ εὐφροσύνη ἐν Ἱερουσαλὴμ ἀπὸ μακρόθεν.

\vs{44}Καὶ κατέστησαν ἐν τῇ ἡμέρᾳ ἐκείνῃ ἄνδρας ἐπὶ τῶν γαζοφυλακίων, τοῖς θησαυροῖς, ταῖς ἀπαρχαῖς, καὶ ταῖς δεκάταις, καὶ τοῖς συνηγμένοις ἐν αὐτοῖς ἄρχουσι τῶν πόλεων, μερίδας τοῖς ἱερεῦσι καὶ τοῖς Λευίταις, ὅτι εὐφροσύνη ἐν Ἰούδα ἐπὶ τοὺς ἱερεῖς, καὶ ἐπὶ τοὺς Λευίτας τοὺς ἑστῶτας.
\vs{45}Καὶ ἐφύλαξαν φυλακὰς Θεοῦ αὐτῶν, καὶ φυλακὰς τοῦ καθαρισμοῦ, καὶ τοὺς ᾄδοντας, καὶ τοὺς πυλωρούς, ὡς ἐντολαὶ Δαυὶδ καὶ Σαλωμὼν υἱοῦ αὐτοῦ.
\vs{46}Ὅτι ἐν ἡμέραις Δαυὶδ Ἀσὰφ ἀπʼ ἀρχῆς πρῶτος τῶν ᾀδόντων, καὶ ὕμνον καὶ αἴνεσιν τῷ Θεῷ,
\vs{47}καὶ πᾶς Ἰσραὴλ ἐν ἡμέραις Ζοροβάβελ καὶ ἐν ταῖς ἡμέραις Νεεμίου, διδόντες μερίδας τῶν ᾀδόντων καὶ τῶν πυλωρῶν, λόγον ἡμέρας ἐν ἡμέρᾳ αὐτοῦ, καὶ ἁγιάζοντες τοῖς Λευίταις, καὶ οἱ Λευῖται ἁγιάζοντες τοῖς υἱοῖς Ἀαρών.

\ch{13}
Ἐν τῇ ἡμέρᾳ ἐκείνῃ ἀνεγνώσθη ἐν βιβλίῳ Μωυσῆ ἐν ὠσὶ τοῦ λαοῦ· καὶ εὑρέθη γεγραμμένον ἐν αὐτῷ, ὅπως μὴ εἰσέλθωσιν Ἀμμανῖται καὶ Μωαβῖται ἐν ἐκκλησίᾳ Θεοῦ ἕως αἰῶνος,
\vs{2}ὅτι οὐ συνήντησαν τοῖς υἱοῖς Ἰσραὴλ ἐν ἄρτῳ καὶ ὕδατι, καὶ ἐμισθώσαντο ἐπʼ αὐτὸν τὸν Βαλαὰμ καταράσασθαι· καὶ ἐπέστρεψεν ὁ Θεὸς ἡμῶν τὴν κατάραν εἰς εὐλογίαν.
\vs{3}Καὶ ἐγένετο ὡς ἤκουσαν τὸν νόμον, καὶ ἐχωρίσθησαν πᾶς ἐπίμικτος ἐν Ἰσραήλ.

\vs{4}Καὶ πρὸ τούτου Ἐλιασὶβ ὁ ἱερεὺς οἰκῶν ἐν γαζοφυλακίῳ οἴκου Θεοῦ ἡμῶν, ἐγγιῶν Τωβίᾳ.
\vs{5}Καὶ ἐποίησεν ἑαυτῷ γαζοφυλάκιον μέγα καὶ ἐκεῖ ἦσαν πρότερον διδόντες τὴν μαναὰ καὶ τὸν λίβανον καὶ τὰ σκεύη, καὶ τὴν δεκάτην τοῦ σίτου καὶ τοῦ οἴνου καὶ τοῦ ἐλαίου, ἐντολὴν τῶν Λευιτῶν καὶ τῶν ᾀδόντων καὶ τῶν πυλωρῶν, καὶ ἀπαρχὰς τῶν ἱερέων.
\vs{6}Καὶ ἐν παντὶ τούτῳ οὐκ ἤμην ἐν Ἱερουσαλήμ· ὅτι ἐν ἔτει τριακοστῷ καὶ δευτέρῳ τοῦ Ἀρθασασθὰ βασιλέως Βαβυλῶνος ἦλθον πρὸς τὸν βασιλέα, καὶ μετὰ τὸ τέλος τῶν ἡμερῶν ᾐτησάμην παρὰ τοῦ βασιλέως,
\vs{7}καὶ ἦλθον εἰς Ἱερουσαλήμ· καὶ συνῆκα ἐν τῇ πονηρίᾳ ᾗ ἐποίησεν Ἐλιασὶβ τῷ Τωβίᾳ, ποιῆσαι αὐτῷ γαζοφυλάκιον ἐν αὐλῇ οἴκου τοῦ Θεοῦ.

\vs{8}Καὶ πονηρόν μοι ἐφάνη σφόδρα· καὶ ἔῤῥιψα πάντα τὰ σκεύη οἴκου Τωβία ἔξω ἀπὸ τοῦ γαζοφυλακίου.
\vs{9}Καὶ εἶπα, καὶ ἐκαθάρισαν τὰ γαζοφυλάκια· καὶ ἐπέστρεψα ἐκεῖ σκεύη οἴκου τοῦ Θεοῦ, τὴν μαναὰ καὶ τὸν λίβανον.

\vs{10}Καὶ ἔγνων ὅτι μερίδες τῶν Λευιτῶν οὐκ ἐδόθησαν· καὶ ἐφύγοσαν ἀνὴρ εἰς ἀγρὸν αὐτοῦ, οἱ Λευῖται καὶ οἱ ᾄδοντες ποιοῦντες τὸ ἔργον.
\vs{11}Καὶ ἐμαχεσάμην τοῖς στρατηγοῖς, καὶ εἶπα, διὰ τί ἐγκατελείφθη ὁ οἶκος τοῦ Θεοῦ; καὶ συνήγαγον αὐτοὺς, καὶ ἔστησα αὐτοὺς ἐπὶ τῇ στάσει αὐτῶν.
\vs{12}Καὶ πᾶς Ἰούδα ἤνεγκαν δεκάτην τοῦ πυροῦ καὶ τοῦ οἴνου καὶ τοῦ ἐλαίου εἰς τοὺς θησαυροὺς
\vs{13}ἐπὶ χεῖρα Σελεμία τοῦ ἱερέως, καὶ Σαδὼκ τοῦ γραμματέως, καὶ Φαδαΐα ἀπὸ τῶν Λευιτῶν· καὶ ἐπὶ χεῖρα αὐτῶν Ἀνὰν υἱὸς Ζακχοὺρ, υἱὸς Ματθανίου, ὅτι πιστοὶ ἐλογίσθησαν, ἐπʼ αὐτοὺς μερίζειν τοῖς ἀδελφοῖς αὐτῶν.

\vs{14}Μνήσθητί μου ὁ Θεὸς ἐν ταὐτῇ, καὶ μὴ ἐξαλειφθήτω ἔλεός μου ὃ ἐποίησα ἐν οἴκῳ Κυρίου τοῦ Θεοῦ.

\vs{15}Ἐν ταῖς ἡμέραις ἐκείναις εἶδον ἐν Ἰούδα πατοῦντας ληνοὺς ἐν τῷ σαββάτῳ, καὶ φέροντας δράγματα, καὶ ἐπιγεμίζοντας ἐπὶ τοὺς ὄνους καὶ οἶνον καὶ σταφυλὴν καὶ σῦκα καὶ πᾶν βάσταγμα, καὶ φέροντας εἰς Ἱερουσαλὴμ ἐν ἡμέρᾳ τοῦ σαββάτου· καὶ ἐπεμαρτυράμην ἐν ἡμέρᾳ πράσεως αὐτῶν.
\vs{16}Καὶ ἐκάθισαν ἐν αὐτῇ φέροντες ἰχθὺν, καὶ πᾶσαν πρᾶσιν πωλοῦντες τῷ σαββάτῳ τοῖς υἱοῖς Ἰούδα καὶ ἐν Ἱερουσαλήμ.
\vs{17}Καὶ ἐμαχεσάμην τοῖς υἱοῖς Ἰούδα τοῖς ἐλευθέροις, καὶ εἶπα αὐτοῖς, τίς ὁ λόγος οὗτος ὁ πονηρὸς, ὃν ὑμεῖς ποιεῖτε, καὶ βεβηλοῦτε τὴν ἡμέραν τοῦ σαββάτου;
\vs{18}Οὐχὶ οὕτως ἐποίησαν οἱ πατέρες ὑμῶν, καὶ ἤνεγκεν ἐπʼ αὐτοὺς ὁ Θεὸς ἡμῶν καὶ ἐφʼ ἡμᾶς πάντα τὰ κακὰ ταῦτα καὶ ἐπὶ τὴν πόλιν ταύτην; καὶ ὑμεῖς προστίθετε ὀργὴν ἐπὶ Ἰσραὴλ βεβηλῶσαι τὸ σάββατον;

\vs{19}Καὶ ἐγένετο ἡνίκα κατέστησαν πύλαι ἐν Ἱερουσαλὴμ πρὸ τοῦ σαββάτου, καὶ εἶπα, καὶ ἔκλεισαν τὰς πύλας· καὶ εἶπα, ὥστε μὴ ἀνοιγῆναι αὐτὰς ἕως ὀπίσω τοῦ σαββάτου· καὶ ἐκ τῶν παιδαρίων μου ἔστησα ἐπὶ τὰς πύλας, ὥστε μὴ αἴρειν βαστάγματα ἐν ἡμέρᾳ τοῦ σαββάτου.
\vs{20}Καὶ ηὐλίσθησαν πάντες, καὶ ἐποίησαν πρᾶσιν ἔξω Ἰερουσαλὴμ ἅπαξ καὶ δίς.
\vs{21}Καὶ ἐπεμαρτυράμην ἐν αὐτοῖς, καὶ εἶπα πρὸς αὐτοὺς, διὰ τί ὑμεῖς αὐλίζεσθε ἀπέναντι τοῦ τείχους; ἐὰν δευτερώσητε, ἐκτενῶ χεῖρά μου ἐν ὑμῖν· ἀπὸ τοῦ καιροῦ ἐκείνου οὐκ ἤλθοσαν ἐν σαββάτῳ.
\vs{22}Καὶ εἶπα τοῖς Λευίταις, οἳ ἦσαν καθαριζόμενοι, καὶ ἐρχόμενοι φυλάσσοντες τὰς πύλας, ἁγιάζειν τὴν ἡμέραν τοῦ σαββάτου.

Πρὸς ταῦτα μνήσθητί μου ὁ Θεὸς, καὶ φεῖσαί μου κατὰ τὸ πλῆθος τοῦ ἐλέους σου.

\vs{23}Καὶ ἐν ταῖς ἡμέραις ἐκείναις εἶδον τοὺς Ἰουδαίους οἳ ἐκάθισαν γυναῖκας Ἀζωτίας, Ἀμμανίτιδας, Μωαβίτιδας·
\vs{24}καὶ οἱ υἱοὶ αὐτῶν ἥμισυ λαλοῦντες Ἀζωτιστὶ, καὶ οὐκ εἰσὶν ἐπιγινώσκοντες λαλεῖν Ἰουδαϊστί.
\vs{25}Καὶ ἐμαχεσάμην μετʼ αὐτῶν, καὶ καταρασάμην αὐτούς· καὶ ἐπάταξα ἐν αὐτοῖς ἄνδρας, καὶ ἐμαδάρωσα αὐτοὺς, καὶ ὤρκισα αὐτοὺς ἐν τῷ Θεῷ, ἐὰν δῶτε τὰς θυγατέρας ὑμῶν τοῖς υἱοῖς αὐτῶν, καὶ ἐὰν λάβητε ἀπὸ τῶν θυγατέρων αὐτῶν τοῖς υἱοῖς ὑμῶν.
\vs{26}Οὐχ οὕτως ἥμαρτε Σαλωμὼν βασιλεὺς Ἰσραήλ; καὶ ἐν ἔθνεσι πολλοῖς οὐκ ἦν βασιλεὺς ὅμοιος αὐτῷ, καὶ ἀγαπώμενος τῷ Θεῷ ἦν, καὶ ἔδωκεν αὐτὸν ὁ Θεὸς εἰς βασιλέα ἐπὶ πάντα Ἰσραὴλ, καὶ τοῦτον ἐξέκλιναν αἱ γυναῖκες αἱ ἀλλότριαι.
\vs{27}Καὶ ὑμῶν μὴ ἀκουσώμεθα ποιῆσαι πᾶσαν πονηρίαν ταύτην, ἀσυνθετῆσαι ἐν τῷ Θεῷ ἡμῶν, καθίσαι γυναῖκας ἀλλοτρίας.

\vs{28}Καὶ ἀπὸ υἱῶν Ἰωαδὰ τοῦ Ἐλισοὺβ τοῦ ἱερέως τοῦ μεγάλου νυμφίου τοῦ Σαναβαλλὰτ τοῦ Οὐρανίτου, καὶ ἐξέβρασα αὐτὸν ἀπʼ ἐμοῦ.
\vs{29}Μνήσθητι αὐτοῖς ὁ Θεὸς ἐπὶ ἀλχιστείᾳ τῆς ἱερατείας, καὶ διαθήκῃ τῆς ἱερατείας, καὶ τοὺς Λευίτας.

\vs{30}Καὶ ἐκαθάρισα αὐτοὺς ἀπὸ πάσης ἀλλοτριώσεως, καὶ ἔστησα ἐφημερίας τοῖς ἱερεῦσι καὶ τοῖς Λευίταις, ἀνὴρ ὡς τὸ ἔργον αὐτοῦ.
\vs{31}Καὶ τὸ δῶρον τῶν ξυλοφόρων ἐν καιροῖς ἀπὸ χρόνων, καὶ ἐν τοῖς βακχουρίοις. Μνήσθητί μου ὁ Θεὸς ἡμῶν εἰς ἀγαθωσύνην.


\def\book{ΕΣΘΗΡ}
\biblebook{ΕΣΘΗΡ}


\lettrine[lines=2, loversize=0.2, nindent=0em, findent=.25em]{\textcolor{bookheadingcolor}{Ἔ}}{ΤΟΥΣ} δευτέρου βασιλεύοντος Ἀρταξέρξου τοῦ μεγάλου βασιλέως τῇ μιᾷ τοῦ Νισὰν, ἐνύπνιον εἶδε Μαρδοχαῖος ὁ τοῦ Ἰαΐρου, τοῦ Σεμεΐου, τοῦ Κισαίου, ἐκ φυλῆς Βενιαμὶν,
\vs{1b}ἄνθρωπος Ἰουδαῖος οἰκῶν ἐν Σούσοις τῇ πόλει, ἄνθρωπος μέγας, θεραπεύων ἐν τῇ αὐλῇ τοῦ βασιλέως.
\vs{1c}Ἦν δὲ ἐκ τῆς αἰχμαλωσίας, ἦς ᾐχμαλώτευσε Ναβουχοδονόσορ βασιλεὺς Βαβυλῶνος ἐξ Ἰερουσαλὴμ, μετὰ Ἰεχονίου τοῦ βασιλέως τῆς Ἰουδαίας.

\vs{1d}Καὶ τοῦτο αὐτοῦ τὸ ἐνύπνιον· καὶ ἰδοὺ φωναὶ καὶ θόρυβος, βρονταὶ καὶ σεισμὺς, τάραχος ἐπὶ τῆς γῆς.
\vs{1e}Καὶ ἰδοὺ δύο δράκοντες μεγάλοι, ἕτομοι προῆλθον ἀμφότεροι παλαίειν· καὶ ἐγένετο αὐτῶν φωνὴ μεγάλη,
\vs{1f}καὶ τῇ φωνῇ αὐτῶν ἡτοιμάσθη πᾶν ἔθνος εἰς πόλεμον, ὥστε πολεμῆσαι δικαίων ἔθνος.
\vs{1g}Καὶ ἰδοὺ ἡμέρα σκότους καὶ γνόφου, θλιψις καὶ στενοχωρία, κάκωσις καὶ τάραχος μέγας ἐπὶ τῆς γῆς.
\vs{1h}Καὶ ἐταράχθη πᾶν ἔθνος δίκαιον φοβούμενοι τὰ ἑαυτῶν κακὰ, καὶ ἡτοιμάσθησαν ἀπολέσθαι, καὶ ἐβόησαν πρὸς τὸν Θεόν·
\vs{1i}ἀπὸ δὲ τῆς βοῆς αὐτῶν ἐγένετο ὡσανεὶ ἀπὸ μικρᾶς πηγῆς ποταμὸς μέγας, ὕδωρ πολύ.
\vs{1k}Καὶ φῶς και ὁ ἥλιος ἀνέτειλε, καὶ οἱ ταπεινοὶ ὑψώθησαν, καὶ κατέφαγον τοὺς ἐνδόξους.

\vs{1l}Καὶ διεγερθεὶς Μαρδοχαῖος ὁ ἑωρακὼς τὸ ἐνύπνιον τοῦτο, καὶ τί ὁ Θεὸς βεβούλευται ποιῆσαι, εἶχεν αὐτὸ ἐν τῇ καρδία, καὶ ἐν παντὶ λόγῳ ἤθελεν ἐπιγνῶναι αὐτὸ ἕως τῆς νυκτός.
\vs{1m}Καὶ ἡσύχασε Μαρδοχαῖος ἐν τῇ αὐλῇ. μετὰ Γαβαθὰ καὶ Θάῤῥα τῶν δύο εὐνούχων τοῦ βασιλέως, τῶν φυλασσόντων τὴν αὐλήν.
\vs{1n}Ἤκουσέ τε αὐτῶν τοὺς λογισμοὺς, καὶ τὰς μερίμνας αὐτῶν ἐξηρεύνησεν· καὶ ἔμαθεν, ὅτι ἑτοιμάζουσι τὰς χεῖρας ἐπιβαλεῖν Ἀρταξέρξῃ τῷ βασιλεῖ· καὶ ὑπέδειξε τῷ βασιλεῖ περὶ αὐτῶ.
\vs{1o}Καὶ ἐξήτασεν ὁ βασιλεὺς τοὺς δύο εὐνούχους, καὶ ὁμολογήσαντες σαντες ἀπήχθησαν.
\vs{1p}Καὶ ἔγραψεν ὁ βασιλεὺς τοὺς λόγους τούτους εἰς μνημόσυνον, καὶ Μαρδοχαῖος ἔγραψε περὶ τῶν λόγων τούτων.
\vs{1q}Καὶ ἐπέταξεν ὁ βασιλεὺς Μαρδοχαίῳ θεραπεύειν ἐν τῇ αὐλῇ, καὶ ἔδωκεων αὐτῷ δόματα περὶ τούτων.
\vs{1r}Καὶ ἦν Ἀμὰν Ἀμαδάθου Βουγαῖος ἔνδοξος ἐνώπιον τοῦ βασιλέως, καὶ ἐζήτησε κακποιῆσαι τὸν Μαρδοχαῖον καὶ τον λαὸν αὐτοῦ, ὑπὲρ τῶν δύο εὐνούχων τοῦ βασιλέως.

\vs{1s}Καὶ ἐγένετο μετὰ τοὺς λόγους τούτους ἐν ταῖς ἡμέραις Ἀρταξέρξου· οὗτος ὁ Ἀρταξέρξης ἀπὸ τῆς Ἰνδικῆς ἑκατὸν εἰκοσιεπτὰ χωρῶν ἐκράτησεν·
\vs{2}Ἐν αὐταῖς ταῖς ἡμέραις ὅτε ἐθρονίσθη βασιλεὺς Ἀρταξέρξης ἐν Σούσοις τῇ πόλει,
\vs{3}ἐν τῷ τρίτῳ ἔτει βασιλεύοντος αὐτοῦ, δοχὴν ἐποίησε τοῖς φίλοις καὶ τοῖς λοιποῖς ἔθνεσι, καὶ τοῖς Περσῶν καὶ Μήδων ἐνδόξοις, καὶ τοῖς ἄρχουσι τῶν σατραπῶν.

\vs{4}Καὶ μετὰ ταῦτα μετὰ τὸ δεῖξαι αὐτοῖς τὸν πλοῦτον τῆς βασιλείας αὐτοῦ, καὶ τὴν δόξαν τῆς εὐφροσύνης τοῦ πλούτου αὐτοῦ ἐν ἡμέραις ἑκατὸν ὀγδοήκοντα.
\vs{5}Ὅτε δὲ ἀνεπληρώθησαν αἱ ἡμέραι τοῦ γάμου, ἐποίησεν ὁ βασιλεὺς πότον τοῖς ἔθνεσι τοῖς εὑρεθεῖσιν εἰς τὴν πόλιν ἐπὶ ἡμέρας ἓξ, ἐν αὐλῇ οἴκου τοῦ βασιλέως
\vs{6}κεκοσμημένῃ βυσσίνοις καὶ καρπασίνοις τεταμένοις ἐπὶ σχοινίοις βυσσίνοις καὶ πορφυροῖς, ἐπὶ κύβοις χρυσοῖς καὶ ἀργυροῖς, ἐπὶ στύλοις Παρίνοις καὶ λιθίνοις· κλίναι χρυσαῖ καὶ ἀργυραῖ ἐπὶ λιθοστρώτου σμαραγδίτου λίθου, καὶ πιννίνου, καὶ Παρίνου λίθου· καὶ στρωμναὶ διαφανεῖς ποικίλως διηνθισμέναι, κύκλῳ ῥόδα πεπασμένα·
\vs{7}Ποτήρια χρυσᾶ καὶ ἀργυρᾶ, καὶ ἀνθράκινον κυλίκιον προκείμενον ἀπὸ ταλάντων τρισμυρίων· οἶνος πολὺς καὶ ἡδὺς, ὃν αὐτὸς ὁ βασιλεὺς ἔπινεν.
\vs{8}Ὁ δὲ πότος οὗτος οὐ κατὰ προκείμενον νόμον ἐγένετο· οὕτως δὲ ἠθέλησεν ὁ βασιλεὺς, καὶ ἐπέταξε τοῖς οἰκονόμοις ποιῆσαι τὸ θέλημα αὐτοῦ, καὶ τῶν ἀνθρώπων.
\vs{9}Καὶ Ἀστὶν ἡ βασίλισσα ἐποίησε πότον ταῖς γυναιξὶν ἐν τοῖς βασιλείοις, ὅπου ὁ βασιλεὺς Ἀρταξέρξης.

\vs{10}Ἐν δὲ τῇ ἡμέρᾳ τῇ ἑβδόμῃ ἡδέως γενόμενος ὁ βασιλεὺς, εἶπε τῷ Ἀμὰν, καὶ Βαζὰν, καὶ Θάῤῥα, καὶ Βωραζὶ, καὶ Ζαθολθὰ, καὶ Ἀβαταζὰ, καὶ Θαραβά, τοῖς ἑπτὰ εὐνούχοις τοῖς διακόνοις τοῦ βασιλέως Ἀρταξέρξου,
\vs{11}εἰσαγαγεῖν τὴν βασίλισσαν πρὸς αὐτὸν, βασιλεύειν αὐτὴν, καὶ περιθεῖναι αὐτῇ τὸ διάδημα, καὶ δεῖξαι αὐτὴν τοῖς ἄρχουσι, καὶ τοῖς ἔθνεσι τὸ κάλλος αὐτῆς, ὅτι καλὴ ἦν.
\vs{12}Καὶ οὐκ εἰσήκουσεν αὐτοῦ Ἀστὶν ἡ βασίλισσα ἐλθεῖν μετὰ τῶν εὐνούχων· καὶ ἐλυπήθη ὁ βασιλεὺς καὶ ὠργίσθη,

\vs{13}Καὶ εἶπε τοῖς φίλοις αὐτοῦ, κατὰ ταῦτα ἐλάλησεν Ἀστίν, ποιήσατε οὖν περὶ τούτου νόμον καὶ κρίσιν.
\vs{14}Καὶ προσῆλθεν αὐτῷ Ἀρκεσαῖος, καὶ Σαρσαθαῖος, καὶ Μαλισεὰρ, οἱ ἄρχοντες Περσῶν καὶ Μήδων, οἱ ἐγγὺς τοῦ βασιλέως, οἱ πρῶτοι παρακαθήμενοι τῷ βασιλεῖ,
\vs{15}καὶ ἀπήγγειλαν αὐτῷ κατὰ τοὺς νόμους, ὡς δεῖ ποιῆσαι Ἀστὶν τῇ βασιλίσσῃ, ὅτι οὐκ ἐποίησε τὰ ὑπὸ τοῦ βασιλέως προσταχθέντα διὰ τῶν εὐνούχων.

\vs{16}καὶ εἶπεν ὁ Μουχαῖος πρὸς τὸν βασιλέα καὶ τοὺς ἄρχοντας, οὐ τὸν βασιλέα μόνον ἠδίκησεν Ἀστὶν ἡ βασίλισσα, ἀλλὰ καὶ πάντας τοὺς ἄρχοντας καὶ τοὺς ἡγουμένους τοῦ βασιλέως·
\vs{17}Καὶ γὰρ διηγήσατο αὐτοῖς τὰ ῥήματα τῆς βασιλίσσης, καὶ ὡς ἀντεῖπε τῷ βασιλεῖ· ὡς οὖν ἀντεῖπε τῷ βασιλεῖ Ἀρταξέρξῃ,
\vs{18}οὕτω σήμερον αἱ τυραννίδες αἱ λοιπαὶ τῶν ἀρχόντων Περσῶν καὶ Μήδων ἀκούσασαι τὰ τῷ βασιλεῖ λεχθέντα ὑπʼ αὐτῆς, τολμήσουσιν ὁμοίως ἀτιμάσαι τοὺς ἄνδρας αὐτῶν.
\vs{19}Εἶ οὖν δοκεῖ τῷ βασιλεῖ, προσταξάτω βασιλικὸν, καὶ γραφήτω κατὰ τοὺς νόμους Μήδων καὶ Περσῶν, καὶ μὴ ἄλλως χρησάσθω, μηδὲ εἰσελθάτω ἔτι ἡ βασίλισσα πρὸς αὐτὸν, καὶ τὴν βασιλείαν αὐτῆς δότω ὁ βασιλεὺς γυναικὶ κρείττονι αὐτῆς.
\vs{20}Καὶ ἀκουσθήτω ὁ νόμος ὁ ὑπὸ τοῦ βασιλέως, ὃν ἐὰν ποιῇ ἐν τῇ βασιλείᾳ αὐτοῦ· καὶ οὕτω πᾶσαι αἱ γυναῖκες περιθήσουσι τιμὴν τοῖς ἀνδράσιν ἑαυτῶν, ἀπὸ πτωχοῦ ἕως πλουσίου.

\vs{21}Καὶ ἤρεσεν ὁ λόγος τῷ βασιλεῖ καὶ τοῖς ἄρχουσι· καὶ ἐποίησεν ὁ βασιλεὺς καθὰ ἐλάλησεν ὁ Μουχαῖος,
\vs{22}καὶ ἀπέστειλεν εἰς πᾶσαν τὴν βασιλείαν κατὰ χώραν, κατὰ τὴν λέξιν αὐτῶν, ὥστε εἶναι φόβον αὐτοῖς ἐν ταῖς οἰκίαις αὐτῶν.

\ch{2}
Καὶ μετὰ τοὺς λόγους τούτους ἐκόπασεν ὁ βασιλεὺς τοῦ θυμοῦ, καὶ οὐκ ἔτι ἐμνήσθη τῆς Ἀστίν, μνημονεύων οἷα ἐλάλησε, καὶ ὡς κατέκρινεν αὐτήν.
\vs{2}Καὶ εἶπαν οἱ διάκονοι τοῦ βασιλέως, ζητηθήτω τῷ βασιλεῖ κοράσια ἄφθορα καλὰ τῷ εἴδει·
\vs{3}Καὶ καταστήσει ὁ βασιλεὺς κωμάρχας ἐν πάσαις ταῖς χώραις τῆς βασιλείας αὐτοῦ· καὶ ἐπιλεξάτωσαν κοράσια παρθενικὰ καλὰ τῷ εἴδει εἰς Σοῦσαν τὴν πόλιν εἰς τὸν γυναικῶνα· καὶ παραδοθήτωσαν τῷ εὐνούχῳ τοῦ βασιλέως τῷ φύλακι τῶν γυναικῶν· καὶ δοθήτω σμῆγμα, καὶ ἡ λοιπὴ ἐπιμέλεια.
\vs{4}Καὶ ἡ γυνὴ ἣ ἂν ἀρέσῃ τῷ βασιλεῖ, βασιλεύσει ἀντὶ Ἀστίν· καὶ ἤρεσε τῷ βασιλεῖ τὸ πρᾶγμα, καὶ ἐποίησεν οὕτως.

\vs{5}Καὶ ἄνθρωπος ἦν Ἰουδαῖος ἐν Σούσοις τῇ πόλει, καὶ ὄνομα αὐτοῦ Μαρδοχαῖος ὁ τοῦ Ἰαΐρου, τοῦ Σεμείου, τοῦ Κισαίου, ἐκ φυλῆς Βενιαμίν,
\vs{6}ὃς ἦν αἰχμάλωτος ἐξ Ἱερουσαλὴμ, ἣν ᾐχμαλώτευσε Ναβουχοδονόσορ βασιλεὺς Βαβυλῶνος.
\vs{7}Καὶ ἦν τούτῳ παῖς θρεπτή, θυγάτηρ Ἀμιναδὰβ ἀδελφοῦ πατρὸς αὐτοῦ, καὶ ὄνομα αὐτῇ Ἐσθήρ· ἐν δὲ τῷ μεταλλάξαι αὐτῆς τοὺς γονεῖς, ἐπαίδευσεν αὐτὴν ἑαυτῷ εἰς γυναῖκα· καὶ ἦν τὸ κοράσιον καλὸν τῷ εἴδει.

\vs{8}Καὶ ὅτε ἠκούσθη τὸ τοῦ βασιλέως πρόσταγμα, συνήχθησαν πολλὰ κοράσια εἰς Σοῦσάν τὴν πόλιν ὑπὸ χεῖρα Γαῒ, καὶ ἤχθη Ἐσθὴρ πρὸς Γαῒ τὸν φύλακα τῶν γυναικῶν.
\vs{9}Καὶ ἤρεσεν αὐτῷ τὸ κοράσιον, καὶ εὗρε χάριν ἐνώπιον αὐτοῦ· καὶ ἔσπευσε δοῦναι αὐτῇ τὸ σμῆγμα, καὶ τὴν μερίδα, καὶ τὰ ἑπτὰ κοράσια τὰ ὑποδεδειγμένα αὐτῇ ἐκ βασιλικοῦ· καὶ ἐχρήσατο αὐτῇ καλῶς καὶ ταῖς ἅβραις αὐτῆς ἐν τῷ γυναικῶνι·
\vs{10}Καὶ οὐχ ὑπέδειξεν Ἐσθὴρ τὸ γένος αὐτῆς οὐδὲ τὴν πατρίδα· ὁ γὰρ Μαρδοχαῖος ἐνετείλατο αὐτῇ μὴ ἀπαγγεῖλαι.

\vs{11}Καθʼ ἑκάστην δὲ ἡμέραν περιεπάτει ὁ Μαρδοχαῖος κατὰ τὴν αὐλὴν τὴν γυναικείαν, ἐπισκοπῶν τί Ἐσθὴρ συμβήσεται.
\vs{12}Οὗτος δὲ ἦν καιρὸς κορασίου εἰσελθεῖν πρὸς τὸν βασιλέα, ὅταν ἀναπληρώσῃ μῆνας δεκαδύο· οὕτως γὰρ ἀναπληροῦνται αἱ ἡμέραι τῆς θεραπείας, μῆνας ἓξ ἀλειφομέναις ἐν σμυρνίνῳ ἐλαίῳ, καὶ μῆνας ἓξ ἐν τοῖς ἀρώμασι καὶ ἐν τοῖς σμήγμασι τῶν γυναικῶν,
\vs{13}καὶ τότε εἰσπορεύεται πρὸς τὸν βασιλέα· καὶ ᾧ ἐὰν εἴπῃ, παραδώσει αὐτὴν συνεισέρχεσθαι αὐτῷ ἀπὸ τοῦ γυναικῶνος ἕως τῶν βασιλείων·
\vs{14}Δείλης εἰσπορεύεται, καὶ πρὸς ἡμέραν ἀποτρέχει εἰς τὸν γυναικῶνα τὸν δεύτερον, οὗ Γαῒ ὁ εὐνοῦχος τοῦ βασιλέως ὁ φύλαξ τῶν γυναικῶν, καὶ οὐκ ἔτι εἰσπορεύεται πρὸς τὸν βασιλέα, ἐὰν μὴ κληθῇ ὀνόματι.

\vs{15}Ἐν δὲ τῷ ἀναπληροῦσθαι τὸν χρόνον Ἐσθὴρ τῆς θυγατρὸς Ἀμιναδὰβ ἀδελφοῦ πατρὸς Μαρδοχαίου εἰσελθεῖν πρὸς τὸν βασιλέα, οὐδὲν ἠθέτησεν ὧν ἐνετείλατο ὁ εὐνοῦχος ὁ φύλαξ τῶν γυναικῶν· ἦν γὰρ Ἐσθὴρ εὑρίσκουσα χάριν παρὰ πάντων τῶν βλεπόντων αὐτήν.
\vs{16}καὶ εἰσῆλθεν Ἐσθὴρ πρὸς Ἀρταξέρξην τὸν βασιλέα τῷ δωδεκάτῳ μηνὶ, ὅς ἐστιν Ἀδάρ, τῷ ἑβδόμῳ ἔτει τῆς βασιλείας αὐτοῦ.
\vs{17}Καὶ ἠράσθη ὁ βασιλεὺς Ἐσθὴρ, καὶ εὗρε χάριν παρὰ πάσας τὰς παρθένους, καὶ ἐπέθηκεν αὐτῇ τὸ διάδημα τὸ γυναικεῖον.
\vs{18}Καὶ ἐποίησεν ὁ βασιλεὺς πότον πᾶσι τοῖς φίλοις αὐτοῦ καὶ ταῖς δυνάμεσιν ἐπὶ ἡμέρας ἑπτά, καὶ ὕψωσε τοὺς γάμους Ἐσθήρ, καὶ ἄφεσιν ἐποίησε τοῖς ὑπὸ τὴν βασιλείαν αὐτοῦ.
\vs{19}Ὁ δὲ Μαρδοχαῖος ἐθεράπευεν ἐν τῇ αὐλῇ.
\vs{20}Ἡ δὲ Ἐσθὴρ οὐχ ὑπέδειξε τὴν πατρίδα αὐτῆς· οὕτως γὰρ ἐνετείλατο αὐτῇ Μαρδοχαῖος, φοβεῖσθαι τὸν Θεὸν, καὶ ποιεῖν τὰ προστάγματα αὐτοῦ, καθὼς ἦν μετʼ αὐτοῦ· καὶ Ἐσθὴρ οὐ μετήλλαξε τὴν ἀγωγὴν αὐτῆς.

\vs{21}Καὶ ἐλυπήθησαν οἱ δύο εὐνοῦχοι τοῦ βασιλέως οἱ ἀρχισωματοφύλακες, ὅτι προήχθη Μαρδοχαῖος, καὶ ἐζήτουν ἀποκτεῖναι Ἀρταξέρξην τὸν βασιλέα.
\vs{22}Καὶ ἐδηλώθη Μαρδοχαίῳ ὁ λόγος, καὶ ἐσήμανεν Ἐσθήρ, καὶ αὐτὴ ἐνεφάνησε τῷ βασιλεῖ τὰ τῆς ἐπιβουλῆς.
\vs{23}Ὁ δὲ βασιλεὺς ἤτασε τοὺς δύο εὐνούχους, καὶ ἐκρέμασεν αὐτούς· καὶ προσέταξεν ὁ βασιλεὺς καταχωρίσαι εἰς μνημόσυνον ἐν τῇ βασιλικῇ βιβλιοθήκῃ ὑπὲρ τῆς εὐνοίας Μαρδοχαίου, ἐν ἐγκωμίῳ.

\ch{3}
Μετὰ δὲ ταῦτα ἐδόξασεν ὁ βασιλεὺς Ἀρταξέρξης Ἀμὰν Ἀμαδαθοῦ Βουγαῖον, καὶ ὕψωσεν αὐτὸν, καὶ ἐπρωτοβάθρει πάντων τῶν φίλων αὐτοῦ,
\vs{2}καὶ πάντες οἱ ἐν τῇ αὐλῇ προσεκύνουν αὐτῷ· οὕτως γὰρ προσέταξεν ὁ βασιλεὺς ποιῆσαι· ὁ δὲ Μαρδοχαῖος οὐ προσεκύνει αὐτῷ.
\vs{3}Καὶ ἐλάλησαν οἱ ἐν τῇ αὐλῇ τοῦ βασιλέως τῷ Μαρδοχαίῳ, Μαρδοχαῖε, τί παρακούεις τὰ ὑπὸ τοῦ βασιλέως λεγόμενα;

\vs{4}Καθʼ ἑκάστην ἡμέραν ἐλάλουν αὐτῷ, καὶ οὐχ ὑπήκουεν αὐτῶν· καὶ ὑπέδειξαν τῷ Ἀμὰν, Μαρδοχαῖον τοῖς τοῦ βασιλέως λόγοις ἀντιτασσόμενον, καὶ ὑπέδειξεν αὐτοῖς ὁ Μαρδοχαῖος ὅτι Ἰουδαῖός ἐστι.
\vs{5}Καὶ ἐπιγνοὺς Ἀμὰν ὅτι οὐ προσκυνεῖ αὐτῷ Μαρδοχαῖος, ἐθυμώθη σφόδρα,
\vs{6}Καὶ ἐβουλεύσατο ἀφανίσαι πάντας τοὺς ὑπὸ τὴν Ἀρταξέρξου βασιλείαν Ἰουδαίους.

\vs{7}Καὶ ἐποίησε ψήφισμα ἐν ἔτει δωδεκάτῳ τῆς βασιλείας Ἀρταξέρξου, καὶ ἔβαλε κλήρους ἡμέραν ἐξ ἡμέρας, καὶ μῆνα ἐκ μηνὸς, ὥστε ἀπολέσαι ἐν μιᾷ ἡμέρᾳ τὸ γένος Μαρδοχαίου· καὶ ἔπεσεν ὁ κλῆρος εἰς τὴν τεσσαρεσκαιδεκάτην τοῦ μηνὸς ὅς ἐστιν Ἀδάρ.
\vs{8}Καὶ ἐλάλησε πρὸς τὸν βασιλέα Ἀρταξέρξην, λέγων, ὑπάρχει ἔθνος διεσπαρμένον ἐν τοῖς ἔθνεσιν ἐν πάσῃ τῇ βασιλείᾳ σου, οἱ δὲ νόμοι αὐτῶν ἔξαλλοι παρὰ πάντα τὰ ἔθνη, τῶν δὲ νόμων τοῦ βασιλέως παρακούουσι, καὶ οὐ συμφέρει τῷ βασιλεῖ ἐᾶσαι αὐτούς.
\vs{9}Εἰ δοκεῖ τῷ βασιλεῖ, δογματισάτω ἀπολέσαι αὐτοὺς, κᾀγὼ διαγράψω εἰς τὸ γαζοφυλάκιον τοῦ βασιλέως ἀργυρίου τάλαντα μύρια.
\vs{10}Καὶ περιελόμενος ὁ βασιλεὺς τὸν δακτύλιον, ἔδωκεν εἰς χεῖρας τῷ Ἀμὰν, σφραγίσαι κατὰ τῶν γεγραμμένων κατὰ τῶν Ἰουδαίων.
\vs{11}Καὶ εἶπεν ὁ βασιλεὺς τῷ Ἀμὰν, τὸ μὲν ἀργύριον ἔχε, τῷ δὲ ἔθνει χρῶ ὡς βούλει.

\vs{12}Καὶ ἐκλήθησαν οἱ γραμματεῖς τοῦ βασιλέως μηνὶ πρώτῳ τῇ τρισκαιδεκάτῃ, καὶ ἔγραψαν ὡς ἐπέταξεν Ἀμὰν τοῖς στρατηγοῖς καὶ τοῖς ἄρχουσι κατὰ πᾶσαν χώραν ἀπὸ Ἰνδικῆς ἕως τῆς Αἰθιοπίας, ταῖς ἑκατὸν εἰκοσιεπτὰ χώραις, τοῖς τε ἄρχουσι τῶν ἐθνῶν κατὰ τὴν αὐτῶν λέξιν, διὰ Ἀρταξέρξου τοῦ βασιλέως.
\vs{13}Καὶ ἀπεστάλη διὰ βιβλιοφόρων εἰς τὴν Ἀρταξέρξου βασιλείαν, ἀφανίσαι τὸ γένος τῶν Ἰουδαίων ἐν ἡμέρᾳ μιᾷ μηνὸς δωδεκάτου, ὅς ἐστιν Ἀδὰρ, καὶ διαρπάσαι τὰ ὑπάρχοντα αὐτῶν.

\vs{13a}Τῆς δὲ ἐπιστολῆς ἐστι τὸ ἀντίγραφον τόδε. Βασιλεὺς μέγας Ἀρταξέρξης τοῖς ἀπὸ τῆς Ἰνδικῆς ἕως τῆς Αἰθιοπίας ἑκατὸν εἰκοσιεπτὰ χωρῶν ἄρχουσι καὶ τοπάρχαις ὑποτεταγμένοις τάδε γράφει.
\vs{13b}Πολλῶν ἐπάρξας ἐθνῶν, καὶ πάσης ἐπικρατήσας οἰκουμένης, ἐβουλήθην, μὴ τῷ θράσει τῆς ἐξουσίας ἐπαιρόμενος, ἐπιεικέστερον δὲ καὶ μετὰ ἠπιότητος ἀεὶ διεξάγων τοὺς τῶν ὑποτεταγμένων ἀκυμάντους διαπαντὸς καταστῆσαι βίους, τήν τε βασιλείαν ἥμερον καὶ πορευτὴν μέχρι περάτων παρεξόμενος, ἀνανεώσασθαί τε τὴν ποθουμένην τοῖς πᾶσιν ἀνθρώποις εἰρήνην.
\vs{13c}Πυθομένου δέ μου τῶν συμβούλων, πῶς ἂν ἀχθείη τοῦτο ἐπὶ πέρας, ὁ σωφροσύνῃ παρʼ ἡμῖν διενέγκας, καὶ ἐν τῇ εὐνοίᾳ ἀπαραλλάκτως καὶ βεβαίᾳ πίστει ἀποδεδειγμένος, καὶ δεύτερον τῶν βασιλειῶν γέρας ἀπενηνεγμένος Ἀμὰν,
\vs{13d}ἐπέδειξεν ἡμῖν ἐν πάσαις ταῖς κατὰ τὴν οἰκουμένην φυλαῖς ἀναμεμίχθαι δυσμενῆ λαόν τινα, τοῖς νόμοις ἀντίθετον πρὸς πᾶν ἔθνος, τά τε τῶν βασιλέων παραπέμποντας διηνεκῶς διατάγματα, πρὸς τὸ μὴ κατατίθεσθαι τὴν ὑφʼ ἡμῶν κατευθυνομένην ἀμέμπτως συναρχίαν.
\vs{13e}Διειληφότες οὖν τόδε τὸ ἔθνος μονώτατον ἐν ἀντιπαραγωγῇ παντὶ διαπαντὸς ἀνθρώπῳ κείμενον, διαγωγὴν νόμων ξενίζουσαν παραλλάσσον, καὶ δυσνοοῦν τοῖς ἡμετέροις πράγμασι τὰ χείριστα συντελοῦν κακὰ, καὶ πρὸς τὸ μὴ τὴν βασιλείαν εὐσταθείας τυγχάνειν·
\vs{13f}προστετάχαμεν οὖν τοὺς σημαινομένους ὑμῖν ἐν τοῖς γεγραμμένοις ὑπὸ Ἀμὰν τοῦ τεταγμένου ἐπὶ τῶν πραγμάτων, καὶ δευτέρου πατρὸς ἡμῶν, πάντας σὺν γυναιξὶ καὶ τέκνοις ἀπολέσαι ὁλοριζὶ, ταῖς τῶν ἐχθρῶν μαχαίραις, ἄνευ παντὸς οἴκτου καὶ φειδοῦς, τῇ τεσσαρεσκαιδεκάτῃ τοῦ δωδεκάτου μηνὸς Ἄδαρ, τοῦ ἐνεστῶτος ἔτους,
\vs{13g}ὅπως οἱ πάλαι καὶ νῦν δυσμενεῖς ἐν ἡμέρᾳ μιᾷ βιαίως εἰς τὸν ᾅδην κατελθόντες, εἰς τὸν μετέπειτα χρόνον εὐσταθῆ καὶ ἀτάραχα παρέχωσιν ἡμῖν διὰ τέλους τὰ πράγματα.

\vs{14}Τὰ δὲ ἀντίγραφα τῶν ἐπιστολῶν ἐξετίθετο κατὰ χώραν· καὶ προσετάγη πᾶσι τοῖς ἔθνεσιν ἑτοίμους εἶναι εἰς τὴν ἡμέραν ταύτην.
\vs{15}Ἐσπεύδετο δὲ τὸ πρᾶγμα, καὶ εἰς Σοῦσαν· ὁ δὲ βασιλεὺς καὶ Ἀμὰν ἐκωθωνίζοντο· ἐταράσσετο δὲ ἡ πόλις.

\ch{4}
Ὁ δὲ Μαρδοχαῖος ἐπιγνοὺς τὸ συντελούμενον, διέῤῥηξε τὰ ἱμάτια ἑαυτοῦ, καὶ ἐνεδύσατο σάκκον, καὶ κατεπάσατο σποδόν· καὶ ἐκπηδήσας διὰ τῆς πλατείας τῆς πόλεως, ἐβόα φωνῇ μεγάλῃ, αἴρεται ἔθνος μηδὲν ἠδικηκός.
\vs{2}Καὶ ἦλθεν ἕως τῆς πύλης τοῦ βασιλέως, καὶ ἔστη· οὐ γὰρ ἦν αὐτῷ ἐξὸν εἰσελθεῖν εἰς τὴν αὐλὴν, σάκκον ἔχοντι καὶ σποδόν.
\vs{3}Καὶ ἐν πάσῃ χώρᾳ οὗ ἐξετίθετο τὰ γράμματα, κραυγὴ καὶ κοπετὸς καὶ πένθος μέγα τοῖς Ἰουδαίοις, σάκκον καὶ σποδὸν ἔστρωσαν ἑαυτοῖς.

\vs{4}Καὶ εἰσῆλθον αἱ ἅβραι καὶ οἱ εὐνοῦχοι τῆς βασιλίσσης, καὶ ἀνήγγειλαν αὐτῇ· καὶ ἐταράχθη ἀκούσασα τὸ γεγονός· καὶ ἀπέστειλε στολίσαι τὸν Μαρδοχαῖον, καὶ ἀφελέσθαι αὐτοῦ τὸν σάκκον· ὁ δὲ οὐκ ἐπείσθη.
\vs{5}Ἡ δὲ Ἐσθὴρ προσεκαλέσατο Ἀχραθαῖον τὸν εὐνοῦχον αὐτῆς, ὃς παρειστήκει αὐτῇ, καὶ ἀπέστειλε μαθεῖν αὕτη παρὰ τοῦ Μαρδοχαίου τὸ ἀκριβές.
\vs{7}Ὁ δὲ Μαρδοχαῖος ὑπέδειξεν αὐτῷ τὸ γεγονὸς, καὶ τὴν ἐπαγγελίαν ἣν ἐπηγγείλατο Ἀμὰν τῷ βασιλεῖ εἰς τὴν γάζαν ταλάντων μυρίων, ἵνα ἀπολέσῃ τοὺς Ἰουδαίους.
\vs{8}καὶ τὸ ἀντίγραφον τὸ ἐν Σούσοις ἐκτεθὲν ὑπὲρ τοῦ ἀπολέσθαι αὐτοὺς, ἔδωκεν αὐτῷ δεῖξαι τῇ Ἐσθήρ· καὶ εἶπεν αὐτῷ, ἐντείλασθαι αὐτῇ εἰσελθούσῃ παραιτήσασθαι τὸν βασιλέα, καὶ ἀξιῶσαι αὐτὸν περὶ τοῦ λαοῦ, μνησθεῖσα ἡμερῶν ταπεινώσεώς σου, ὡς ἐτράφης ἐν χειρί μου, διότι Ἀμὰν ὁ δευτερεύων τῷ βασιλεῖ ἐλάλησεν καθʼ ἡμῶν εἰς θάνατον· ἐπικάλεσαι τὸν Κύριον, καὶ λάλησον τῷ βασιλεῖ περὶ ἡμῶν, ῥύσαι ἡμᾶς ἐκ θανάτου.

\vs{9}Εἰσελθὼν δὲ ὁ Ἀχραθαῖος ἐλάλσεν αὐτῇ πάντας τοὺς λόγους τούτους.
\vs{10}Εἶπεν δὲ Ἐσθὴρ πρὸς Ἀχραθαῖον, πορεύθητι πρὸς Μαρδοχαῖον, καὶ εἶπον,
\vs{11}ὅτι τὰ ἔθνη πάντα τῆς βασιλείας γινώσκει ὅτι πᾶς ἄνθρωπος ἢ γυνὴ ὃς εἰσελεύσεται πρὸς τὸν βασιλέα εἰς τὴν αὐλὴν τὴν ἐσωτέραν ἄκλητος, οὐκ ἔστιν αὐτῷ σωτηρία· πλὴν ᾧ ἐκτείνῃ ὁ βασιλεὺς τὴν χρυσῆν ῥάβδον, οὗτος σωθήσεται· κᾀγὼ οὐ κέκλημαι εἰσελθεῖν πρὸς τὸν βασιλέα, εἰσὶν αὗται ἡμέραι τριάκοντα.
\vs{12}Καὶ ἀπήγγειλεν Ἀχραθαῖος Μαρδοχαίῳ πάντας τοὺς λόγους Ἐσθήρ.

\vs{13}Καὶ εἶπε Μαρδοχαῖος πρὸς Ἀχραθαῖον, πορεύθητι, καὶ εἰπὸν αὐτῇ, Ἐσθήρ, μὴ εἴπῃς σεαυτῇ, ὅτι σωθήσῃ μόνη ἐν τῇ βασιλείᾳ παρὰ πάντας τοὺς Ἰουδαίους·
\vs{14}Ὡς ὅτι ἐὰν παρακούσῃς ἐν τούτῳ τῷ καιρῷ, ἄλλοθεν βοήθεια καὶ σκέπη ἔσται τοῖς Ἰουδαίοις· σὺ δὲ καὶ ὁ οἶκος τοῦ πατρός σου ἀπολεῖσθε· καὶ τίς εἶδεν, εἰ εἰς τὸν καιρὸν τοῦτον ἐβασίλευσας;
\vs{15}καὶ ἐξαπέστειλεν Ἐσθὴρ τὸν ἥκοντα πρὸς αὐτὴν, πρὸς Μαρδοχαῖον, λέγουσα,
\vs{16}βαδίσας ἐκκλησίασον τοὺς Ἰουδαίους τοὺς ἐν Σούσοις, καὶ νηστεύσατε ἐπʼ ἐμοὶ, καὶ μὴ φάγητε μηδὲ πίητε ἐπὶ ἡμέρας τρεῖς νύκτα καὶ ἡμέραν· κᾀγὼ δὲ καὶ αἱ ἅβραι μου ἀσιτήσομεν· καὶ τότε εἰσελεύσομαι πρὸς τὸν βασιλέα παρὰ τὸν νόμον, ἐὰν καὶ ἀπολέσθαι με δέῃ.
\vs{17}Καὶ βαδίσας Μαρδοχαῖος ἐποίησεν ὅσα ἐνετείλατο αὐτῷ Ἐσθήρ·

\vs{17a}Καὶ ἐδεήθη Κυρίου, μνημονεύων πάντα τὰ ἔργα Κυρίου, καὶ εἶπε,
\vs{17b}Κύριε Κύριε βασιλεῦ πάντων κρατῶν, ὅτι ἐν ἐξουσίᾳ σου τὸ πᾶν ἐστι, καὶ οὐκ ἔστιν ὁ ἀντιδοξῶν σοι ἐν τῷ θέλειν σε σῶσαι τὸν Ἰσραήλ.
\vs{17c}Ὅτι σὺ ἐποιήσας τὸν οὐρανὸν καὶ τὴν γῆν, καὶ πᾶν θαυμαζόμενον ἐν τῇ ὑπʼ οὐρανόν. Καὶ Κύριος εἶ πάντων, καὶ οὐκ ἔστιν ὃς ἀντιτάξεταί σοι τῷ Κυρίῳ.
\vs{17d}Σὺ πάντα γινώσκεις· σὺ οἶδας, Κύριε, ὅτι οὐκ ἐν ὕβρει, οὐδὲ ἐν ὑπερηφανείᾳ, οὐδὲ ἐν φιλοδοξίᾳ ἐποίησα τοῦτο, τὸ μὴ προσκυνεῖν τὸν ὑπερήφανον Ἀμάν. Ὅτι ηὐδόκουν φιλεῖν πέλματα ποδῶν αὐτοῦ πρὸς σωτηρίαν Ἰσραήλ.
\vs{17e}Ἀλλʼ ἐποίησα τοῦτο, ἵνα μὴ θῶ δόξαν ἀνθρώπου ὑπεράνω δόξης Θεοῦ· καὶ οὐ προσκυνήσω οὐδένα, πλὴν σοῦ τοῦ κυρίου μου, καὶ οὐ ποιήσω αὐτὰ ἐν ὑπερηφανείᾳ.
\vs{17f}Καὶ νῦν, Κύριε ὁ Θεὸς ὁ βασιλεὺς ὁ Θεὸς Ἀβραὰμ, φεῖσαι τοῦ λαοῦ σου, ὅτι ἐπιβλέπουσιν ἡμῖν εἰς καταφθορὰν, καὶ ἐπεθύμησαν ἀπολέσαι τὴν ἐξ ἀρχῆς κληρονομίαν σου.
\vs{17g}Μὴ ὑπερίδῃς τὴν μερίδα σου, ἣν σεαυτῷ ἐλυτρώσω ἐκ γῆς Αἰγύπτου.
\vs{17h}Ἐπάκουσον τῆς δεήσεώς μου, καὶ ἱλάσθητι τῷ κλήρῳ σου, καὶ στρέψον τὸ πένθος ἡμῶν εἰς εὐωχίαν, ἵνα ζῶντες ὑμνῶμέν σου τὸ ὄνομα Κύριε, καὶ μὴ ἀφανίσῃς στόμα αἰνούντων σε Κύριε.

\vs{17i}Καὶ πᾶς Ἰσραὴλ ἐκέκραξεν ἐξ ἰσχύος αὐτῶν, ὅτι θάνατος αὐτῶν ἐν ὀφθαλμοῖς αὐτῶν.
\vs{17k}Καὶ Ἐσθὴρ ἡ βασίλισσα κατέφυγεν ἐπὶ τὸν Κύριον ἐν ἀγῶνι θανάτου κατειλημμένη, καὶ ἀφελομένη τὰ ἱμάτια τῆς δόξης αὐτῆς, ἐνεδύσατο ἱμάτια στενοχωρίας καὶ πένθους, καὶ ἀντὶ τῶν ὑπερηφάνων ἡδυσμάτων, σποδοῦ καὶ κοπριῶν ἐνέπλησε τὴν κεφαλὴν αὐτῆς· καὶ τὸ σῶμα αὐτῆς ἐταπείνωσε σφόδρα, καὶ πάντα τόπον κόσμου ἀγαλλιάματος αὐτῆς ἔπλησε στρεπτῶν τριχῶν αὐτῆς.

Καὶ ἐδεῖτο Κυρίου Θεοῦ Ἰσραὴλ, καὶ εἶπεν,
\vs{17l}Κύριέ μου βασιλεὺς ἡμῶν σὺ εἶ μόνος, βοήθησόν μοι τῇ μόνῃ, καὶ μὴ ἐχούσῃ βοηθὸν εἰ μὴ σὲ, ὅτι κίνδυνός μου ἐν χειρί μου.
\vs{17m}Ἐγὼ ἤκουον ἐκ γενετῆς μου ἐν φυλῇ πατριᾶς μου, ὅτι σὺ Κύριε ἔλαβες τὸν Ἰσραὴλ ἐκ πάντων τῶν ἐθνῶν, καὶ τοὺς πατέρας ἡμῶν ἐκ πάντων τῶν προγόνων αὐτῶν εἰς κληρονομίαν αἰώνιον, καὶ ἐποίησας αὐτοῖς ὅσα ἐλάλησας.
\vs{17n}Καὶ νῦν ἡμάρτομεν ἐνώπιόν σου, καὶ παρέδωκας ἡμᾶς εἰς χεῖρας τῶν ἐχθρῶν ἡμῶν, ἀνθʼ ὧν ἐδοξάσαμεν τοὺς θεοὺς αὐτῶν· δίκαιος εἶ Κύριε.
\vs{17o}Καὶ νῦν οὐχ ἱκανώθησαν ἐν πικρασμῷ δουλείας ἡμῶν, ἀλλʼ ἔθηκαν τὰς χεῖρας αὐτῶν ἐπὶ τὰς χεῖρας τῶν εἰδώλων αὐτῶν, ἐξᾶραι ὁρισμὸν στόματός σου, καὶ ἀφανίσαι κληρονομίαν σου, καὶ ἐμφράξαι στόμα αἰνούντων σοι καὶ σβέσαι δόξαν οἴκου σου καὶ θυσιαστηρίον σου,
\vs{17p}καὶ ἀνοῖξαι στόμα ἐθνῶν εἰς ἀρετὰς ματαίων, καὶ θαυμασθῆναι βασιλέα σάρκινον εἰς αἰῶνα.

\vs{17q}Μὴπαραδῷς Κύριε τὸ σκῆπτρόνσουτοῖς μὴοὖσι, καὶ μὴ καταγελασάτωσαν ἐν τῇ πτώσει ἡμῶν, ἀλλὰ στρέψον τὴν βουλὴν αὐτῶν ἐπʼ αὐτούς· τὸν δὲ ἀρξάμενον ἐφʼ ἡμᾶς παραδειγμάτισον.
\vs{17r}Μνήσθητι Κύριε, γνώσθητι ἐν καιρῷ θλίψεως ἡμῶν, καὶ ἐμὲ θάρσυνον, βασιλεῦ τῶν θεῶν, καὶ πάσης ἀρχῆς ἐπικρατῶν.
\vs{17s}Δὸς λόγον εὔρυθμον εἰς τὸ στόμα μου ἐνώπιον τοῦ λέοντος, καὶ μετάθες τὴν καρδίαν αὐτοῦ εἰς μῖσος τοῦ πολεμοῦντος ἡμᾶς, εἰς συντέλειαν αὐτοῦ καὶ τῶν ὁμονοούντων αὐτῷ.
\vs{17t}Ἡμᾶς δὲ ῥύσαι ἐν χειρί σου, καὶ βοήθησόν μοι τῇ μόνῃ, καὶ μὴ ἐχούσῃ εἰ μὴ σέ Κύριε·
\vs{17u}πάντων γνῶσιν ἔχεις, καὶ οἶδας ὅτι ἐμίσησα δόξαν ἀνόμων, καὶ βδελύσσομαι κοίτην ἀπεριτμήτων, καὶ παντὸς ἀλλοτρίου.
\vs{17w}Σὺ οἶδας τὴν ἀνάγκην μου, ὅτι βδελύσσομαι τὸ σημεῖον τῆς ὑπερηφανίας μου, ὅ ἐστιν ἐπὶ τῆς κεφαλῆς μου ἐν ἡμέραις ὁπτασίας μου· βδελύσσομαι αὐτὸ ὡς ῥάκος καταμηνίων, καὶ οὐ φορῶ αὐτὸ ἐν ἡμέραις ἡσυχίας μου.
\vs{17x}Καὶ οὐκ ἔφαγεν ἡ δούλη σου τράπεζαν Ἀμὰν, καὶ οὐκ ἐδόξασα συμπόσιον βασιλέως, οὐδὲ ἔπιον οἶνον σπονδῶν.
\vs{17y}Καὶ οὐκ ηὐφράνθη ἡ δούλη σου ἀφʼ ἡμέρας μεταβολῆς μου μέχρι νῦν, πλὴν ἐπὶ σοὶ, Κύριε ὁ Θεὸς Ἁβραάμ.
\vs{17z}Ὁ Θεὸς ὁ ἰσχύων ἐπὶ πάντας, εἰσάκουσον φωνὴν ἀπηλπισμένων, καὶ ῥύσαι ἡμᾶς ἐκ χειρὸς τῶν πονηρευομένων, καὶ ῥῦσαί με ἐκ τοῦ φόβου μου.

\ch{5}
Καὶ ἐγενήθη ἐν τῇ ἡμέρᾳ τῇ τρίτῃ ὡς ἐπαύσατο προσευχομένη, ἐξεδύσατο τὰ ἱμάτια τῆς θεραπείας, καὶ περιεβάλλετο τὴν δόξαν αὐτῆς.
\vs{1a}Καὶ γενηθεῖσα ἐπιφανὴς, ἐπικαλεσαμένη τὸν πάντων ἐπόπτην Θεὸν καὶ σωτῆρα, παρέλαβε τὰς δύο ἅβρας, καὶ τῇ μὲν μιᾷ ἐπηρείδετο ὡς τρυφερευομένη, ἡ δὲ ἑτέρα ἐπηκολούθει κουφίζουσα τὴν ἔνδυσιν αὐτῆς.
\vs{1b}Καὶ αὐτὴ ἐρυθριῶσα ἀκμῇ κάλλους αὐτῆς· καὶ τὸ πρόσωπον αὐτῆς ἱλαρὸν, ὡς προσφιλές· ἡ δὲ καρδία αὐτῆς ἀπεστενωμένη ἀπὸ τοῦ φόβου.
\vs{1c}Καὶ εἰσελθοῦσα πάσας τὰς θύρας, κατέστη ἐνώπιον τοῦ βασιλέως· καὶ αὐτὸς ἐκάθητο ἐπὶ τοῦ θρόνου τῆς βασιλείας αὐτοῦ, καὶ πᾶσαν στολὴν τῆς ἐπιφανείας αὐτοῦ ἐνδεδύκει, ὅλος διὰ χρυσοῦ καὶ λίθων πολυτελῶν, καὶ ἦν φοβερὸς σφόδρα.
\vs{1d}Καὶ ἄρας τὸ πρόσωπον αὐτοῦ πεπυρωμένον δόξῃ, ἐν ἀκμῇ θυμοῦ ἔβλεψεν· καὶ ἔπεσεν ἡ βασίλισσα, καὶ μετέβαλε τὸ χρῶμα αὐτῆς ἐν ἐκλύσει· καὶ κατεπέκυψεν ἐπὶ τὴν κεφαλὴν τῆς ἅβρας τῆς προπορευομένης.
\vs{1e}Καὶ μετέβαλεν ὁ Θεὸς τὸ πνεῦμα τοῦ βασιλέως εἰς πραύτητα, καὶ ἀγωνιάσας ἀνεπήδησεν ἀπὸ τοῦ θρόνου αὐτοῦ, καὶ ἀνέλαβεν αὐτὴν ἐπὶ τὰς ἀγκάλας αὐτοῦ, μέχρις οὗ κατέστη· καὶ παρεκάλει αὐτὴν λόγοις εἰρηνικοῖς, καὶ εἶπεν αὐτῇ,
\vs{1f}τί ἐστιν, Ἐσθήρ; ἐγὼ ὁ ἀδελφός σου, θάρσει, οὐ μὴ ἀποθάνῃς· ὅτι κοινὸν τὸ πρόσταγμα ἡμῶν ἐστίν, πρόσελθε.

\vs{2}Καὶ ἄρας τὴν χρυσῆν ῥάβδον, ἐπέθηκεν ἐπὶ τὸν τράχηλον αὐτῆς, καὶ ἠσπάσατο αὐτὴν, καὶ εἶπε, λάλησόν μοι.
\vs{2a}Καὶ εἶπεν αὐτῷ, εἶδόν σε κύριε ὡς ἄγγελον Θεοῦ, καὶ ἐταράχθη ἡ καρδία μου ἀπὸ φόβου τῆς δόξης σου, ὅτι θαυμαστὸς εἶ κύριε, καὶ τὸ πρόσωπόν σου χαρίτων μεστόν.
\vs{2b}Ἐν δὲ τῷ διαλέγεσθαι αὐτὴν, ἔπεσεν ἀπὸ ἐκλύσεως. Καὶ ὁ βασιλεὺς ἐταράσσετο, καὶ πᾶσα ἡ θεραπεία αὐτοῦ παρεκάλει αὐτήν.

\vs{3}Καὶ εἶπεν ὁ βασιλεύς, τί θέλεις, Ἐσθήρ; καὶ τί σου ἐστὶ τὸ ἀξίωμα; ἕως τοῦ ἡμίσους τῆς βασιλείας μου, καὶ ἔσται σοι.
\vs{4}Εἶπε δὲ Ἐσθὴρ, ἡμέρα μου ἐπίσημος σήμερόν ἐστιν· εἰ οὖν δοκεῖ τῷ βασιλεῖ, ἐλθάτω καὶ αὐτὸς καὶ Ἀμὰν εἰς τὴν δοχὴν, ἣν ποιήσω σήμερον.
\vs{5}Καὶ εἶπεν ὁ βασιλεύς, κατασπεύσατε Ἀμὰν, ὅπως ποιήσωμεν τὸν λόγον Ἐσθήρ· καὶ παραγίνονται ἀμφότεροι εἰς τὴν δοχὴν, ἣν εἶπεν Ἐσθήρ.

\vs{6}Ἐν δὲ τῷ πότῳ εἶπεν ὁ βασιλεὺς πρὸς Ἐσθήρ, τί ἐστι βασίλισσα Ἐσθήρ; καὶ ἔσται ὅσα ἀξιοῖς.
\vs{7}Καὶ εἶπε, τὸ αἴτημά μου, καὶ τὸ ἀξίωμα·
\vs{8}Εἰ εὗρον χάριν ἐνώπιον τοῦ βασιλέως, ἐλθάτω ὁ βασιλεὺς καὶ Ἀμὰν ἔτι τὴν αὔριον εἰς τὴν δοχὴν, ἣν ποιήσω αὐτοῖς, καὶ αὔριον ποιήσω τὰ αὐτά.

\vs{9}Καὶ ἐξῆλθεν ὁ Ἀμὰν ἀπὸ τοῦ βασιλέως ὑπερχαρὴς εὐφραινόμενος· ἐν δὲ τῷ ἰδεῖν Ἀμὰν Μαρδοχαῖον τὸν Ἰουδαῖον ἐν τῇ αὐλῇ, ἐθυμώθη σφόδρα.
\vs{10}Καὶ εἰσελθὼν εἰς τὰ ἴδια, ἐκάλεσε τοὺς φίλους, καὶ Ζωσάραν τὴν γυναῖκα αὐτοῦ,
\vs{11}καὶ ὑπέδειξεν αὐτοῖς τὸν πλοῦτον αὐτοῦ, καὶ τὴν δόξαν ἣν ὁ βασιλεὺς αὐτῷ περιέθηκε, καὶ ὡς ἐποίησεν αὐτὸν πρωτεύειν καὶ ἡγεῖσθαι τῆς βασιλείας.
\vs{12}Καὶ εἶπεν Ἀμὰν, οὐ κέκληκεν ἡ βασίλισσα μετὰ τοῦ βασιλέως οὐδένα εἰς τὴν δοχὴν, ἀλλʼ ἢ ἐμὲ, καὶ εἰς τὴν αὔριον κέκλημαι.
\vs{13}Καὶ ταῦτά μοι οὐκ ἀρέσκει, ὅταν ἴδω Μαρδοχαῖον τὸν Ἰουδαῖον ἐν τῇ αὐλῇ.
\vs{14}Καὶ εἶπε πρὸς αὐτὸν Ζωσάρα ἡ γυνὴ αὐτοῦ, καὶ οἱ φίλοι, κοπήτω σοι ξύλον πηχῶν πεντήκοντα, ὄρθρου δὲ εἶπον τῷ βασιλεῖ, καὶ κρεμασθήτω Μαρδοχαῖος ἐπὶ τοῦ ξύλου· σὺ δὲ εἴσελθε εἰς τὴν δοχὴν σὺν τῷ βασιλεῖ, καὶ εὐφραίνου· καὶ ἤρεσε τὸ ῥῆμα τῷ Ἀμὰν, καὶ ἡτοιμάσθη τὸ ξύλον.

\ch{6}
Ὁ δε Κύριος ἀπέστησε τὸν ὕπνον ἀπὸ τοὺ βασιλέως τὴν νύκτα ἐκείνην· καὶ εἶπε τῷ διακόνῳ αὐτοῦ εἰσφέρειν γράμματα μνημόσυνα τῶν ἡμερῶν ἀναγινώσκειν αὐτῷ.
\vs{2}Εὗρε δὲ τὰ γράμματα τὰ γραφέντα περὶ Μαρδοχαίου, ὡς ἀπήγγειλεν τῷ βασιλεῖ περὶ τῶν δύο εὐνούχων τοῦ βασιλέως ἐν τῷ φυλάσσειν αὐτοὺς, καὶ ζητῆσαι ἐπιβαλεῖν τὰς χεῖρας Ἀρταξέρξῃ.

\vs{3}Εἶπε δὲ ὁ βασιλεὺς, τίνα δόξαν ἢ χάριν ἐποιήσαμεν τῷ Μαρδοχαίῳ; καὶ εἶπαν οἱ διάκονοι τοῦ βασιλέως, οὐκ ἐποίησας αὐτῷ οὐδέν.
\vs{4}Ἐν δὲ τῷ πυνθάνεσθαι τὸν βασιλέα περὶ τῆς εὐνοίας Μαρδοχαίου, ἰδοὺ Ἀμὰν ἐν τῇ αὐλῇ· εἶπν δὲ ὁ βασιλεύς, τίς ἐν τῇ αὐλῇ; ὁ δὲ Ἀμὰν εἰσῆλθεν εἰπεῖν τῷ βασιλεῖ, κρεμάσαι τὸν Μαρδοχαῖον ἐπὶ τῷ ξύλῳ, ᾧ ἡτοίμασε.
\vs{5}Καὶ εἶπαν οἱ διάκονοι τοῦ βασιλέως, ἰδοὺ Ἀμὰν ἕστηκεν ἐν τῇ αὐλῇ· καὶ εἶπεν ὁ βασιλεύς, καλέσατε αὐτόν.

\vs{6}Εἶπε δὲ ὁ βασιλεὺς τῷ Ἀμὰν, τί ποιήσω τῷ ἀνθρώπῳ, ὃν ἐγὼ θέλω δοξάσαι; εἶπε δὲ ἐν ἑαυτῷ Ἀμάν, τίνα θέλει ὁ βασιλεὺς δοξάσαι εἰ μὴ ἐμέ;
\vs{7}Εἶπε δὲ πρὸς τὸν βασιλέα, ἄνθρωπον ὃν ὁ βασιλεὺς θέλει δοξάσαι,
\vs{8}ἐνεγκάτωσαν οἱ παῖδες τοῦ βασιλέως στολὴν βυσσίνην ἣν ὁ βασιλεὺς περιβάλλεται, καὶ ἵππον ἐφʼ ὃν ὁ βασιλεὺς ἐπιβαίνει,
\vs{9}καὶ δότω ἑνὶ τῶν φίλων τοῦ βασιλέως τῶν ἐνδόξων, καὶ στολισάτω τὸν ἄνθρωπον, ὃν ὁ βασιλεὺς ἀγαπᾷ· καὶ ἀναβιβασάτω αὐτὸν ἐπὶ τὸν ἵππον, καὶ κηρυσσέτω διὰ τῆς πλατείας τῆς πόλεως, λέγων, οὕτως ἔσται παντὶ ἀνθρώπῳ ὃν ὁ βασιλεὺς δοξάζει.
\vs{10}Εἶπε δὲ ὁ βασιλεὺς τῷ Ἀμὰν, καλῶς ἐλάλησας· οὕτως ποίησον τῷ Μαρδοχαίῳ τῷ Ἰουδαίῳ, τῷ θεραπεύοντι ἐν τῇ αὐλῇ, καὶ μὴ παραπεσάτω σου λόγος ὧν ἐλάλησας.

\vs{11}Ἔλαβε δὲ Ἀμὰν τὴν στολὴν καὶ τὸν ἵππον, καὶ ἐστόλισε τὸν Μαρδοχαῖον, καὶ ἀνεβίβασεν αὐτὸν ἐπὶ τὸν ἵππον, καὶ διῆλθε διὰ τῆς πλατείας τῆς πόλεως, καὶ ἐκήρυσσε λέγων, οὕτως ἔσται παντὶ ἀνθρώπῳ ὃν ὁ βασιλεὺς θέλει δοξάσαι.

\vs{12}Ἐπέστρεψε δὲ ὁ Μαρδοχαῖος εἰς τὴν αὐλήν· Ἁμὰν δὲ ὑπέστρεψεν εἰς τὰ ἴδια λυπούμενος κατὰ κεφαλῆς.
\vs{13}Καὶ διηγήσατο Ἀμὰν τὰ συμβεβηκότα αὐτῷ Ζωσάρᾳ τῇ γυναικὶ αὐτοῦ, καὶ τοῖς φίλοις· καὶ εἶπαν πρὸς αὐτὸν οἱ φίλοι, καὶ ἡ γυνή εἰ ἐκ γένους Ἰουδαίων Μαρδοχαῖος, ἦρξαι ταπεινοῦσθαι ἐνώπιον αὐτοῦ, πεσὼν πεσῇ· καὶ οὐ μὴ δύνῃ αὐτὸν ἀμύνασθαι, ὅτι Θεὸς ζῶν μετʼ αὐτοῦ.
\vs{14}Ἔτι αὐτῶν λαλούντων, παραγίνονται οἱ εὐνοῦχοι, ἐπισπεύδοντες τὸν Ἀμὰν ἐπὶ τὸν πότον ὃν ἡτοίμασεν Ἐσθήρ.

\ch{7}
Εἴσῆλθε δὲ ὁ βασιλεὺς καὶ Ἀμὰν, συμπιεῖν τῇ βασιλίσσῃ.
\vs{2}Εἶπε δὲ ὁ βασιλεὺς Ἐσθὴρ τῇ δευτέρᾳ ἡμέρᾳ ἐν τῷ πότῳ, τί ἐστιν, Ἐσθὴρ βασίλισσα, καὶ τί τὸ αἴτημά σου; καὶ τί τὸ ἀξίωμά σου; καὶ ἔστω σοι ἕως ἡμίσους τῆς βασιλείας μου.
\vs{3}Καὶ ἀποκριθεῖσα, εἶπεν, εἰ εὗρον χάριν ἐνώπιον τοῦ βασιλέως, δοθήτω ἡ ψυχὴ τῷ αἰτήματί μου, καὶ ὁ λαός μου τῷ ἀξιώματί μου.
\vs{4}Ἐπράθημεν γὰρ ἐγώ τε καὶ ὁ λαός μου εἰς ἀπώλειαν καὶ διαρπαγὴν καὶ δουλείαν, ἡμεῖς καὶ τὰ τέκνα ἡμῶν εἰς παῖδας καὶ παιδίσκας, καὶ παρήκουσα· οὐ γὰρ ἄξιος ὁ διάβολος τῆς αὐλῆς τοῦ βασιλέως.
\vs{5}Εἶπε δὲ ὁ βασιλεὺς, τίς οὗτος, ὅστις ἐτόλμησε ποιῆσαι τὸ πρᾶγμα τοῦτο;
\vs{6}Εἶπε δὲ Ἐσθὴρ, ἄνθρωπος ἐχθρὸς Ἀμὰν, ὁ πονηρὸς οὗτος. Ἀμὰν δὲ ἐταράχθη ἀπὸ τοῦ βασιλέως καὶ τῆς βασιλίσσης.

\vs{7}Ὁ δὲ βασιλεὺς ἐξανέστη ἀπὸ τοῦ συμποσίου εἰς τὸν κῆπον· ὁ δὲ Ἀμὰν παρῃτεῖτο τὴν βασίλισσαν· ἑώρα γὰρ ἑαυτὸν ἐν κακοῖς ὄντα.

\vs{8}Ἐπέστρεψεν δὲ ὁ βασιλεὺς ἐκ τοῦ κήπου. Ἀμὰν δὲ ἐπιπεπτώκει ἐπὶ τὴν κλίνην, ἀξιῶν τὴν βασίλισσαν· εἶπν δὲ ὁ βασιλεὺς, ὥστε καὶ τὴν γυναῖκα βιάζῃ ἐν τῇ οἰκίᾳ μου; Ἀμὰν δὲ ἀκούσας διετράπη τῷ προσώπῳ.
\vs{9}Εἶπε δὲ Βουγαθὰν εἷς τῶν εὐνούχων πρὸς τὸν βασιλέα, ἰδοὺ καὶ ξύλον ἡτοίμασεν Ἀμὰν Μαρδοχαίῳ τῷ λαλήσαντι περὶ τοῦ βασιλέως, καὶ ὤρθωται ἐν τοῖς Ἀμὰν ξύλον πηχῶν πεντήκοντα· εἶπε δὲ ὁ βασιλεὺς, σταυρωθήτω ἐπʼ αὐτοῦ.
\vs{10}Καὶ ἐκρεμάσθη Ἀμὰν ἐπὶ τοῦ ξύλου ὃ ἡτοιμάσθη Μαρδοχαίῳ· καὶ τότε ὁ βασιλεὺς ἐκόπασε τοῦ θυμοῦ.

\ch{8}
Καὶ ἐν αὐτῇ τῇ ἡμέρᾳ ὁ βασιλεὺς Ἀρταξέρξης ἐδωρήσατο Ἐσθὴρ ὅσα ὑπῆρχεν Ἀμὰν τῷ διαβόλῳ· καὶ Μαρδοχαῖος προσεκλήθη ὑπὸ τοῦ βασιλέως· ὑπέδειξε γὰρ Ἐσθὴρ, ὅτι ἐνοικείωται αὐτῇ.
\vs{2}Ἔλαβε δὲ ὁ βασιλεὺς τὸν δακτύλιον ὃν ἀφείλατο Ἀμὰν, καὶ ἔδωκεν αὐτὸν Μαρδοχαίῳ· καὶ κατέστησεν Ἐσθὴρ Μαρδοχαῖον ἐπὶ πάντων τῶν Ἀμάν.

\vs{3}Καὶ προσθεῖσα ἐλάλησε πρὸς τὸν βασιλέα, καὶ προσέπεσε πρὸς τοὺς πόδας αὐτοῦ, καὶ ἠξίου ἀφελεῖν τὴν Ἀμὰν κακίαν, καὶ ὅσα ἐποίησε τοῖς Ἰουδαίοις.
\vs{4}Ἐξέτεινε δὲ ὁ βασιλεὺς Ἐσθὴρ τὴν ῥάβδον τὴν χρυσῆν· ἐξηγέρθη δὲ Ἐσθὴρ παρεστηκέναι τῷ βασιλεῖ,
\vs{5}καὶ εἶπεν Ἐσθὴρ, εἰ δοκεῖ σοι, καὶ εὗρον χάριν, πεμφθήτω ἀποστραφῆναι τὰ γράμματα τὰ ἀπεσταλμένα ὑπὸ Ἀμὰν, τὰ γραφέντα ἀπολέσθαι τοὺς Ἰουδαίους, οἵ εἰσιν ἐν τῇ βασιλείᾳ σου.
\vs{6}Πῶς γὰρ δυνήσομαι ἰδεῖν τὴν κάκωσιν τοῦ λαοῦ μου, καὶ πῶς δυνήσομαι σωθῆναι ἐν τῇ ἀπωλείᾳ τῆς πατρίδος μου;

\vs{7}Καὶ εἶπεν ὁ βασιλεὺς πρὸς Ἐσθὴρ, εἰ πάντα τὰ ὑπάρχοντα Ἀμὰν ἔδωκα καὶ ἐχαρισάμην σοι, καὶ αὐτὸν ἐκρέμασα ἐπὶ ξύλου, ὅτι τὰς χεῖρας ἐπήνεγκε τοῖς Ἰουδαίοις, τί ἔτι ἐπιζητεῖς;
\vs{8}Γράψατε καὶ ὑμεῖς ἐκ τοῦ ὀνόματός μου, ὡς δοκεῖ ὑμῖν, καὶ σφραγίσατε τῷ δακτυλίῳ μου· ὅσα γὰρ γράφεται τοῦ βασιλέως ἐπιτάξαντος, καὶ σφραγισθῇ τῷ δακτυλίῳ μου, οὐκ ἔστιν αὐτοῖς ἀντειπεῖν.

\vs{9}Ἐκλήθησαν δὲ οἱ γραμματεῖς ἐν τῷ πρώτῳ μηνὶ, ὅς ἐστι Νισὰν, τρίτῃ καὶ εἰκάδι τοῦ αὐτοῦ ἔτους, καὶ ἐγράφη τοῖς Ἰουδαίοις, ὅσα ἐνετείλατο τοῖς οἰκονόμοις καὶ τοῖς ἄρχουσι τῶν σατραπῶν, ἀπὸ τῆς Ἰνδικῆς ἕως τῆς Αἰθιοπίας, ἑκατὸν εἰκοσιεπτὰ σατράπαις κατὰ χώραν καὶ χώραν, κατὰ τὴν αὐτῶν λέξιν.

\vs{10}Ἐγράφη δὲ διὰ τοῦ βασιλέως, καὶ ἐσφραγίσθη τῷ δακτυλίῳ αὐτοῦ· καὶ ἐξαπέστειλαν τὰ γράμματα διὰ βιβλιοφόρων,
\vs{11}ὡς ἐπέταξεν αὐτοῖς χρῆσθαι τοῖς νόμοις αὐτῶν ἐν πάσῃ πόλει, βοηθῆσαί τε αὑτοῖς, καὶ χρῆσθαι τοῖς ἀντιδίκοις αὐτῶν καὶ τοῖς ἀντικειμένοις αὐτῶν, ὡς βούλονται,
\vs{12}ἐν ἡμέρᾳ μιᾷ ἐν πάσῃ τῇ βασιλείᾳ Ἀρταξέρξου, τῇ τρισκαιδεκάτῃ τοῦ δωδεκάτου μηνὸς, ὅς ἐστιν Ἀδάρ.

\vs{12a}Ὧν ἐστιν ἀντίγραφον τῆς ἐπιστολῆς τὰ ὑπογεγραμμένα·

\vs{12b}Βασιλεὺς μέγας Ἀρταξέρξης τοῖς ἀπὸ τῆς Ἰνδικῆς ἕως τῆς Αἰθιοπίας ἑκατὸν εἰκοσιεπτὰ σατραπείαις χωρῶν ἄρχουσι, καὶ τοῖς τὰ ἡμέτερα φρονοῦσι, χαίρειν.
\vs{12c}Πολλοὶ τῇ πλείστῃ τῶν εὐεργετούντων χρηστότητι πυκνότερον τιμώμενοι, μεῖζον ἐφρόνησαν, καὶ οὐ μόνον τοὺς ὑποτεταγμένους ἡμῖν ζητοῦσι κακοποιεῖν, τόν τε κόρον οὐ δυνάμενοι φέρειν, καὶ τοῖς ἑαυτῶν εὐεργέταις ἐπιχειροῦσι μηχανᾶσθαι·
\vs{12d}καὶ τὴν εὐχαριστίαν οὐ μόνον ἐκ τῶν ἀνθρώπων ἀνταναιροῦντες, ἀλλὰ καὶ τοῖς τῶν ἀπειραγάθων κόμποις ἐπαρθέντες, τοῦ τὰ πάντα κατοπτεύοντος ἀεὶ Θεοῦ μισοπόνηρον ὑπολαμβάνουσιν ἐκφεύξεσθαι δίκην.
\vs{12e}Πολλάκις δὲ καὶ πολλοὺς τῶν ἐπʼ ἐξουσίαις τεταγμένεν τῶν πιστευθέντων χειρίζειν φίλων τὰ πράγματα, παραμυθία μετόχους αἱμάτων ἀθώων καταστήσασα περιέβαλε συμφοραῖς ἀνηκέστοις,
\vs{12f}τῷ τῆς κακοηθείας ψευδεῖ παραλογισμῷ παραλογισαμένων τὴν τῶν ἐπικρατούντων ἀκέραιον εὐγνωμοσύνην.
\vs{12g}Σκοπεῖν δὲ ἔξεστιν, οὐ τοσοῦτον ἐκ τῶν παλαιοτέρων ὡς παρεδώκαμεν ἱστοριῶν, ὅσα ἐστὶ παρὰ πόδας ὑμᾶς ἐκζητοῦντας ἀνοσίως συντετελεσμένα τῇ τῶν ἀναξίᾳ δυναστευόντων λοιμότητι·
\vs{12h}καὶ προσέχειν εἰς τὰ μετὰ ταῦτα, εἰς τὸ τὴν βασιλείαν ἀτάραχον τοῖς πᾶσιν ἀνθρώποις μετʼ εἰρήνης παρεξόμεθα
\vs{12i}χρώμενοι ταῖς μεταβολαῖς, τὰ δὲ ὑπὸ τὴν ὄψιν ἐρχόμενα διακρίνοντες ἀεὶ μετʼ ἐπιεικεστέρας ἀπαντήσεως.

\vs{12k}Ὡς γὰρ Ἀμὰν Ἀμαδαθοῦ Μακεδὼν ταῖς ἀληθείαις ἀλλότριος τοῦ τῶν Περσῶν αἵματος, καὶ πολὺ διεστηκὼς τῆς ἡμετέρας χρηστότητος ἐπιξενωθεὶς ἡμῖν,
\vs{12l}ἔτυχεν ἧς ἔχομεν πρὸς πᾶν ἔθνος φιλανθρωπίας ἐπὶ τοσοῦτον, ὥστε ἀναγορεύεσθαι ἡμῶν πατερα, καὶ προσκυνούμενον ὑπὸ πάντων τὸ δεύτερον τοῦ βασιλικοῦ θρόνου πρόσωπον διατελεῖν.
\vs{12m}Οὐκ ἐνέγκας δὲ τὴν ὑπερηφανίαν, ἐπετήδευσε τῆς ἀρχῆς στερῆσαι ἡμᾶς, καὶ τοῦ πνεύματος,
\vs{12n}τόν τε ἡμέτερον σωτῆρα καὶ διαπαντὸς εὐεργέτην Μαρδοχαῖον, καὶ τὴν ἄμεμπτον τῆς βασιλείας κοινωνὸν Ἐσθὴρ σὺν παντὶ τῷ τούτων ἔθνει, πολυπλόκοις μεθόδων παραλογισμοῖς αἰτησάμενος εἰς ἀπώλειαν.
\vs{12o}Διὰ γὰρ τῶν τρόπων τούτων ᾠήθη λαβὼν ἡμᾶς ἐρήμους, τὴν τῶν Περσῶν ἐπικράτησιν εἰς τοὺς Μακεδόνας μετάξαι.
\vs{12p}Ἡμεῖς δὲ τοὺς ὑπὸ τοῦ τρισαλιτηρίου παραδεδομένους εἰς ἀφανισμὸν Ἰουδαίους, εὑρίσκομεν οὐ κακούργους ὄντας δικαιοτάτοις δὲ πολιτευομένους νόμοις,
\vs{12q}ὄντας δὲ υἱοὺς τοῦ ὑψίστου μεγίστου ζῶντος Θεοῦ, τοῦ κατευθύνοντος ἡμῖν τε καὶ τοῖς προγόνοις ἡμῶν τὴν βασιλείαν ἐν τῇ καλλίστῃ διαθέσει.

\vs{12r}Καλῶς οὖν ποιήσετε μὴ προσχρησάμενοι τοῖς ὑπὸ Ἀμὰν Ἀμαδαθοῦ ἀποσταλεῖσι γράμμασι, διὰ τὸ αὐτὸν τὸν ταῦτα ἐξεργασάμενον πρὸς ταῖς Σούσων πύλαις ἐσταυρῶσθαι σὺν τῇ πανοικίᾳ, τὴν καταξίαν τοῦ τὰ πάντα ἐπικρατοῦντος Θεοῦ διατάχους ἀποδόντος αὐτῷ κρίσιν.
\vs{12s}Τὸ δὲ ἀντίγραφον τῆς ἐπιστολῆς ταύτης ἐκθέντες ἐν παντὶ τόπῳ μετὰ παῤῥησίας, ἐᾷν τοὺς Ἰουδαίους χρῆσθαι τοῖς ἑαυτῶν νομίμοις, καὶ συνεπισχύειν αὐτοῖς, ὅπως τοὺς ἐν καιρῷ θλίψεως ἐπιθεμένους αὐτοῖς, ἀμύνωνται τῇ τρισκαιδεκάτῃ τοῦ δωδεκάτου μηνὸς Ἀδὰρ τῇ αὐτῇ ἡμέρᾳ·
\vs{12t}Ταύτην γὰρ ὁ τὰ πάντα δυναστεύων Θεὸς ἀντʼ ὀλεθρίας τοῦ ἐκλεκτοῦ γένους, ἐποίησεν αὐτοῖς εὐφροσύνην.

\vs{12u}Καὶ ὑμεῖς οὖν ἐν ταῖς ἐπωνύμοις ὑμῶν ἑορταῖς, ἐπίσημον ἡμέραν μετὰ πάσης εὐωχίας ἄγετε, ὅπως καὶ νῦν καὶ μετὰ ταῦτα σωτήρια ᾖ ἡμῖν, καὶ τοῖς εὐνοοῦσι Πέρσαις, τοῖς δὲ ἡμῖν ἐπιβουλεύουσι, μνημόσυνον τῆς ἀπωλείας.
\vs{12x}Πᾶσα δὲ πόλις ἢ χώρα τὸ σύνολον, ἥτις κατὰ ταῦτα μὴ ποιήσῃ, δόρατι καὶ πυρὶ καταναλωθήσεται μετʼ ὀργῆς· οὐ μόνον ἀνθρώποις ἄβατος, ἀλλὰ καὶ θηρίοις καὶ πετεινοῖς εἰς τὸν ἅπαντα χρόνον ἔχθιστος κατασταθήσεται.

\vs{13}Τὰ δὲ ἀντίγραφα ἐκτιθέσθωσαν ὀφθαλμοφανῶς ἐν πάσῃ τῇ βασιλείᾳ, ἑτοίμους τε εἶναι πάντας τοὺς Ἰουδαίους εἰς ταύτην τὴν ἡμέραν, πολεμῆσαι αὐτῶν τοὺς ὑπεναντίους.

\vs{14}Οἱ μὲν οὖν ἱππεῖς ἐξῆλθον σπεύδοντες τὰ ὑπὸ τοῦ βασιλέως λεγόμενα ἐπιτελεῖν· ἐξετέθη δὲ τὸ πρόσταγμα καὶ ἐν Σούσοις.

\vs{15}Ὁ δὲ Μαρδοχαῖος ἐξῆλθεγ ἐστολισμένος τὴν βασιλικὴν στολὴν, καὶ στέφανον ἔχων χρυσοῦν, καὶ διάδημα βύσσινον πορφυροῦν· ἰδόντες δὲ οἱ ἐν Σούσοις ἐχάρησαν.
\vs{16}Τοῖς δὲ Ἰουδαίοις ἐγένετο φῶς καὶ εὐφροσύνη
\vs{17}κατὰ πόλιν καὶ χώραν, οὗ ἂν ἐξετέθη τὸ πρόσταγμα· οὗ ἂν ἐξετέθη τὸ ἔκθεμα, χαρὰ καὶ εὐφροσύνη τοῖς Ἰουδαίοις, κώθων καὶ εὐφροσύνη· καὶ πολλοὶ τῶν ἐθνῶν περιετέμοντο, καὶ Ἰουδάϊζον διὰ τὸν φόβον τῶν Ἰουδαίων.

\ch{9}
Ἐν γὰρ τῷ δωδεκάτῳ μηνὶ τρισκαιδεκάτῃ τοῦ μηνός, ὅς ἐστιν Ἀδάρ, παρῆν τὰ γράμματα τὰ γραφέντα ὑπὸ τοῦ βασιλέως.
\vs{2}Ἐν αὐτῇ τῇ ἡμέρᾳ ἀπώλοντο οἱ ἀντικείμενοι τοῖς Ἰουδαίοις· οὐδεὶς γὰρ ἀντέστη, φοβούμενος αὐτούς.
\vs{3}Οἱ γὰρ ἄρχοντες τῶν σατραπῶν, καὶ οἱ τύραννοι, καὶ οἱ βασιλικοὶ γραμματεῖς ἐτίμων τοὺς Ἰουδαίους· ὁ γὰρ φόβος Μαρδοχαίου ἐνέκειτο αὐτοῖς.
\vs{4}Προσέπεσε γὰρ τὸ πρόσταγμα τοῦ βασιλέως ὀνομασθῆναι ἐν πάσῃ τῇ βασιλείᾳ.
\vs{6}Καὶ ἐν Σούσοις τῇ πόλει ἀπέκτειναν οἱ Ἰουδαῖοι ἄνδρας πεντακοσίους,
\vs{7}τόν τε Φαρσαννὲς, καὶ Δελτφὼν, καὶ Φασγὰ, καὶ Φαραδαθὰ,
\vs{8}καὶ Βαρεὰ, καὶ Σαρβακὰ,
\vs{9}καὶ Μαρμασιμὰ, καὶ Ῥουφαῖον, καὶ Ἀρσαῖον, καὶ Ζαβουθαῖον,
\vs{10}τοὺς δέκα υἱοὺς Ἀμὰν Ἀμαδάθοῦ Βουγαίου τοῦ ἐχθροῦ τῶν Ἰουδαίων, καὶ διήρπασαν
\vs{11}ἐν αὐτῇ τῇ ἡμέρᾳ· ἐπεδόθη τε ὁ ἀριθμὸς τῷ βασιλεῖ τῶν ἀπολωλότων ἐν Σούσοις.
\vs{12}Εἶπε δὲ ὁ βασιλεὺς πρὸς Ἐσθὴρ, ἀπώλεσαν οἱ Ἰουδαῖοι ἐν Σούσοις τῇ πόλει ἄνδρας πεντακοσίους, ἐν δὲ τῇ περιχώρῳ πῶς οἴει ἐχρήσαντο; τί οὖν ἀξιοῖς ἔτι, καὶ ἔσται σοι;

\vs{13}Καὶ εἶπεν Ἐσθὴρ τῷ βασιλεῖ, δοθήτω τοῖς Ἰουδαίοις χρῆσθαι ὡσαύτως τὴν αὔριον, ὥστε τοὺς δέκα υἱοὺς Ἀμὰν κρεμάσαι.
\vs{14}Καὶ ἐπέτρεψεν οὕτως γενέσθαι, καὶ ἐξέθηκε τοῖς Ἰουδαίοις τῆς πόλεως τὰ σώματα τῶν υἱῶν Ἀμὰν κρεμάσαι·
\vs{15}Καὶ συνήχθησαν οἱ Ἰουδαῖοι ἐν Σούσοις τῇ τεσσαρεσκαιδεκάτῃ τοῦ Ἀδὰρ, καὶ ἀπέκτειναν ἄνδρας τριακοσίους, καὶ οὐδὲν διήρπασαν.

\vs{16}Οἱ δὲ λοιποὶ τῶν Ἰουδαίων οἱ ἐν τῇ βασιλείᾳ συνήχθησαν, καὶ ἑαυτοῖς ἐβοήθουν, καὶ ἀνεπαύσαντο ἀπὸ τῶν πολεμίων· ἀπώλεσαν γὰρ αὐτῶν μυρίους πεντακισχιλίους τῇ τρισκαιδεκάτῃ τοῦ Ἀδὰρ, καὶ οὐδὲν διήρπασαν.
\vs{17}Καὶ ἀνεπαύσαντο τῇ τεσσαρεσκαιδεκάτῃ τοῦ αὐτοῦ μηνός, καὶ ἦγον αὐτὴν ἡμέραν ἀναπαύσεως μετὰ χαρᾶς καὶ εὐφροσύνης.
\vs{18}Οἱ δὲ Ἰουδαῖοι ἐν Σούσοις τῇ πόλει συνήχθησαν καὶ τῇ τεσσαρεσκαιδεκάτῃ, καὶ ἀνεπαύσαντο· ἦγον δὲ καὶ τὴν πεντεκαιδεκάτην μετὰ χαρᾶς καὶ εὐφροσύνης.
\vs{19}Διὰ τοῦτο οὖν οἱ Ἰουδαῖοι οἱ διεσπαρμένοι ἐν πάσῃ χώρᾳ τῇ ἔξω, ἄγουσι τὴν τεσσαρεσκαιδεκάτην τοῦ Ἀδὰρ ἡμέραν ἀγαθὴν μετʼ εὐφροσύνης, ἀποστέλλοντες μερίδας ἕκαστος τῷ πλησίον.

\vs{20}Ἔγραψεν δὲ Μαρδοχαῖος τοὺς λόγους τούτους εἰς βιβλίον, καὶ ἐξαπέστειλε τοῖς Ἰουδαίοις, ὅσοι ἦσαν ἐν τῇ Ἀρταξέρξου βασιλείᾳ τοῖς ἐγγὺς καὶ τοῖς μακράν,
\vs{21}στῆσαι τὰς ἡμέρας ταύτας ἀγαθάς, ἄγειν τε τὴν τεσσαρεσκαιδεκάτην καὶ τὴν πεντεκαιδεκάτην τοῦ Ἀδάρ.
\vs{22}Ἐν γὰρ ταύταις ταῖς ἡμέραις ἀνεπαύσαντο οἱ Ἰουδαῖοι ἀπὸ τῶν ἐχθρῶν αὐτῶν· καὶ τὸν μῆνα ἐν ᾧ ἐστράφη αὐτοῖς, ὃς ἦν Ἀδὰρ, ἀπὸ πένθους εἰς χαρὰν, καὶ ἀπὸ ὀδύνης εἰς ἀγαθὴν ἡμέραν, ἄγειν ὅλον ἀγαθὰς ἡμέρας γάμων καὶ εὐφροσύνης, ἐξαποστέλλοντας μερίδας τοῖς φίλοις, καὶ τοῖς πτωχοῖς.

\vs{23}Καὶ προσεδέξαντο οἱ Ἰουδαῖοι, καθὼς ἔγραψεν αὐτοῖς ὁ Μαρδοχαῖος·
\vs{24}Πῶς Ἀμὰν Ἀμαδαθοῦ ὁ Μακεδὼν ἐπολέμει αὐτοὺς, καθὼς ἔθετο ψήφισμα καὶ κλῆρον ἀφανίσαι αὐτοὺς,
\vs{25}καὶ ὡς εἰσῆλθε πρὸς τὸν βασιλέα, λέγων, κρεμάσαι τὸν Μαρδοχαῖον· ὅσα δὲ ἐπεχείρησεν ἐπάξαι ἐπὶ τοὺς Ἰουδαίους κακὰ, ἐπʼ αὐτὸν ἐγένοντο, καὶ ἐκρεμάσθη αὐτὸς, καὶ τὰ τέκνα αὐτοῦ.
\vs{26}Διὰ τοῦτο ἐπεκλήθησαν αἱ ἡμέραι αὗται Φρουραὶ διὰ τοὺς κλήρους, ὅτι τῇ διαλέκτῳ αὐτῶν καλοῦνται Φρουραί· διὰ τοὺς λόγους τῆς ἐπιστολῆς ταύτης, καὶ ὅσα πεπόνθασι διὰ ταῦτα, καὶ ὅσα αὐτοῖς ἐγένετο,
\vs{27}καὶ ἔστησε· καὶ προσεδέχοντο οἱ Ἰουδαῖοι ἐφʼ ἑαυτοῖς καὶ ἐπὶ τῷ σπέρματι αὐτῶν καὶ ἐπὶ τοῖς προστεθειμένοις ἐπʼ αὐτῶν, οὐδὲ μὴν ἄλλως χρήσονται. αἱ δὲ ἡμέραι αὗται μνημόσυνον ἐπιτελούμενον κατὰ γενεὰν καὶ γενεὰν, καὶ πόλιν, καὶ πατριὰν, καὶ χώραν.
\vs{28}Αἱ δὲ ἡμέραι αὗται τῶν Φρουραὶ ἀχθήσονται εἰς τὸν ἅπαντα χρόνον, καὶ τὸ μνημόσυνον αὐτῶν οὐ μὴ ἐκλίπῃ ἐκ τῶν γενεῶν.

\vs{29}Καὶ ἔγραψεν Ἐσθὴρ ἡ βασίλισσα θυγάτηρ Ἀμιναδὰβ, καὶ Μαρδοχαῖος ὁ Ἰουδαῖος, ὅσα ἐποίησαν, τό, τε στερέωμα τῆς ἐπιστολῆς τῶν Φρουραί.
\vs{31}Καὶ Μαρδοχαῖος καὶ Ἐσθὴρ ἡ βασίλισσα ἔστησαν ἑαυτοῖς καθʼ ἑαυτῶν, καὶ τότε στήσαντες κατὰ τῆς ὑγιείας ἑαυτῶν, καὶ τὴν βουλὴν αὐτῶν.
\vs{32}Καὶ Ἐσθὴρ λόγῳ ἔστησεν εἰς τὸν αἰῶνα, καὶ ἐγράφη εἰς μνημόσυνον.

\ch{10}
Ἔγραψε δὲ ὁ βασιλεὺς ἐπὶ τὴν βασιλείαν τῆς τε γῆς καὶ τῆς θαλάσσης.
\vs{2}Καὶ τὴν ἰσχὺν αὐτοῦ καὶ ἀνδραγαθίαν, πλοῦτόν τε καὶ δόξαν τῆς βασιλείας αὐτοῦ, ἰδοὺ γέγραπται ἐν βιβλίῳ βασιλέων Περσῶν καὶ Μήδων, εἰς μνημόσυνον.
\vs{3}Ὁ δὲ Μαρδοχαῖος διεδέχετο τὸν βασιλέα Ἀρταξέρξην, καὶ μέγας ἦν ἐν τῇ βασιλείᾳ, καὶ δεδοξασμένος ὑπὸ τῶν Ἰουδαίων· καὶ φιλούμενος διηγεῖτο τὴν ἀγωγὴν παντὶ τῷ ἔθνει αὐτοῦ.

\vs{3a}Καὶ εἶπε Μαρδοχαῖος, παρὰ τοῦ Θεοῦ ἐγένετο ταῦτα.
\vs{3b}Ἐμνήσθην γὰρ περὶ τοῦ ἐνυπνίου οὗ εἶδον περὶ τῶν λόγων τούτων· οὐδὲ γὰρ παρῆλθεν ἀπʼ αὐτῶν λόγος·
\vs{3c}Ἡ μικρὰ πηγὴ ἣ ἐγένετο ποταμός, καὶ ἦν φῶς καὶ ἥλιος καὶ ὕδωρ πολύ. Ἐσθήρ ἐστιν ὁ ποταμός, ἣν ἐγάμησεν ὁ βασιλεὺς, καὶ ἐποίησε βασίλισσαν.
\vs{3d}Οἱ δὲ δύο δράκοντες, ἐγώ εἰμι καὶ Ἀμάν.
\vs{3e}Τὰ δὲ ἔθνη, τὰ ἐπισυναχθέντα ἀπολέσαι τὸ ὄνομα τῶν Ἰουδαίων.
\vs{3f}Τὸ δὲ ἔθνος τὸ ἐμόν, οὗτός ἐστιν Ἰσραὴλ, οἱ βοήσαντες πρὸς τὸν Θεὸν, καὶ σωθέντες· καὶ ἔσωσε Κύριος τὸν λαὸν αὐτοῦ, καὶ ἐῤῥύσατο Κύριος ἡμᾶς ἐκ πάντων τῶν κακῶν τούτων· καὶ ἐποίησεν ὁ Θεὸς τὰ σημεῖα, καὶ τὰ τέρατα τὰ μεγάλα, ἃ οὐ γέγονεν ἐν τοῖς ἔθνεσι.
\vs{3g}Διὰ τοῦτο ἐποίησε κλήρους δύο, ἕνα τῷ λαῷ τοῦ Θεοῦ, καὶ ἕνα πᾶσι τοῖς ἔθνεσι.
\vs{3h}Καὶ ἦλθον οἱ δύο κλῆροι οὗτοι εἰς ὥραν καὶ καιρὸν, καὶ εἰς ἡμέραν κρίσεως, ἐνώπιον τοῦ Θεοῦ καὶ πᾶσι τοῖς ἔθνεσι.
\vs{3i}Καὶ ἐμνήσθη ὁ Θεὸς τοῦ λαοῦ αὐτοῦ, καὶ ἐδικαίωσε τὴν κληρονομίαν αὑτοῦ.
\vs{3k}Καὶ ἔσονται αὐτοῖς αἱ ἡμέραι αὗται, ἐν μηνὶ Ἀδὰρ, τῇ τεσσαρεσκαιδεκάτῃ καὶ τῇ πεντεκαιδεκάτῃ τοῦ μηνὸς, μετὰ συναγωγῆς καὶ χαρᾶς καὶ εὐφροσύνης ἐνώπιον τοῦ Θεοῦ, κατὰ γενεὰς εἰς τὸν αἰῶνα ἐν τῷ λαῷ αὐτοῦ Ἰσραήλ.

\vs{3l}Ἔτους τετάρτου βασιλεύοντος Πτολεμαίου καὶ Κλεοπάτρας, εἰσήνεγκε Δοσίθεος, ὃς ἔφη εἶναι ἱερεὺς καὶ Λευίτης, καὶ Πτολεμαῖος ὁ υἱὸς αὐτοῦ, τὴν προκειμένην ἐπιστολὴν τῶν Φρουραί, ἣν ἔφασαν εἶναι, καὶ ἡρμηνευκέναι Λυσίμαχον Πτολεμαίου, τὸν ἐν Ἱερουσαλήμ.

Τέλος τῆς Ἐσθήρ.


\def\book{ΙΩΒ}
\biblebook{ΙΩΒ}


\lettrine[lines=2, loversize=0.2, nindent=0em, findent=.25em]{\textcolor{bookheadingcolor}{Ἄ}}{ΝΘΡΩΠΟΣ} τις ἦν ἐν χώρᾳ τῇ Αὐσίτιδι, ᾧ ὄνομα Ἰὼβ, καὶ ἦν ὁ ἄνθρωπος ἐκεῖνος ἀληθινὸς, ἄμεμπτος, δίκαιος, θεοσεβὴς, ἀπεχόμενος ἀπὸ παντὸς πονηροῦ πράγματος.
\vs{2}Ἐγένοντο δὲ αὐτῷ υἱοὶ ἑπτὰ καὶ θυγατέρες τρεῖς.
\vs{3}Καὶ ἦν τὰ κτήνη αὐτοῦ πρόβατα ἑπτακισχίλια, κάμηλοι τρισχίλιαι, ζεύγη βοῶν πεντακόσια, θήλειαι ὄνοι νομάδες πεντακόσιαι, καὶ ὑπηρεσία πολλὴ σφόδρα, καὶ ἔργα μεγάλα ἦν αὐτῷ ἐπὶ τῆς γῆς, καὶ ἦν ὁ ἄνθρωπος ἐκεῖνος εὐγενὴς τῶν ἀφʼ ἡλίου ἀνατολῶν.

\vs{4}Συμπορευόμενοι δὲ οἱ υἱοὶ αὐτοῦ πρὸς ἀλλήλους, ἐποιοῦσαν πότον καθʼ ἑκάστην ἡμέραν, συμπαραλαμβάνοντες ἅμα καὶ τὰς τρεῖς ἀδελφὰς αὐτῶν, ἐσθίειν καὶ πίνειν μετʼ αὐτῶν.
\vs{5}Καὶ ὡς ἂν συνετελέσθησαν αἱ ἡμέραι τοῦ πότου, ἀπέστελλεν Ἰὼβ καὶ ἐκαθάριζεν αὐτοὺς ἀνιστάμενος τοπρωῒ, καὶ προσέφερε περὶ αὐτῶν θυσίας, κατὰ τὸν ἀριθμὸν αὐτῶν, καὶ μόσχον ἕνα περὶ ἁμαρτίας περὶ τῶν ψυχῶν αὐτῶν· ἔλεγε γὰρ Ἰὼβ, μή ποτε οἱ υἱοί μου ἐν τῇ διανοίᾳ αὐτῶν κακὰ ἐνενόησαν πρὸς Θεόν· οὕτως οὖν ἐποίει Ἰὼβ πάσας τὰς ἡμέρας.

\vs{6}Καὶ ἐγένετο ὡς ἡ ἡμέρα αὕτη, καὶ ἰδοὺ ἦλθον οἱ ἄγγελοι τοῦ Θεοῦ παραστῆναι ἐνώπιον τοῦ Κυρίου, καὶ ὁ διάβολος ἦλθε μετʼ αὐτῶν.
\vs{7}Καὶ εἶπεν ὁ Κύριος τῷ διαβόλῳ, πόθεν παραγέγονας; καὶ ἀποκριθεὶς ὁ διάβολος τῷ Κυρίῳ, εἶπε, περιελθὼν τὴν γῆν καὶ ἐμπεριπατήσας τὴν ὑπʼ οὐρανὸν πάρειμι.
\vs{8}Καὶ εἶπεν αὐτῷ ὁ Κύριος, προσέσχες τῇ διανοίᾳ σου κατὰ τοῦ παιδός μου Ἰώβ; ὅτι οὐκ ἔστι κατʼ αὐτὸν ἐπὶ τῆς γῆς; ἄνθρωπος ἄμεμπτος, ἀληθινὸς, θεοσεβὴς, ἀπεχόμενος ἀπὸ παντὸς πονηροῦ πράγματος;
\vs{9}Ἀπεκρίθη δὲ ὁ διάβολος, καὶ εἶπεν ἐναντίον τοῦ Κυρίου, μὴ δωρεὰν Ἰὼβ σέβεται τὸν Κύριον;
\vs{10}Οὐ σὺ περιέφραξας τὰ ἔξω αὐτοῦ, καὶ τὰ ἔσω τῆς οἰκίας αὐτοῦ, καὶ τὰ ἔξω πάντων τῶν ὄντων αὐτοῦ κύκλῳ; τὰ δὲ ἔργα τῶν χειρῶν αὐτοῦ εὐλόγησας; καὶ τὰ κτήνη αὐτοῦ πολλὰ ἐποίησας ἐπὶ τῆς γῆς;
\vs{11}Ἀλλὰ ἀπόστειλον τὴν χεῖρά σου, καὶ ἅψαι πάντων ὧν ἔχει· ἦ μὴν εἰς πρόσωπόν σε εὐλογήσει.
\vs{12}Τότε εἶπεν ὁ Κύριος τῷ διαβόλῳ, ἰδοὺ πάντα ὅσα ἐστὶν αὐτῷ, δίδωμι ἐν τῇ χειρί σου, ἀλλʼ αὐτοῦ μὴ ἅψῃ· καὶ ἐξῆλθεν ὁ διάβολος παρὰ τοῦ Κυρίου.

\vs{13}Καὶ ἦν ὡς ἡ ἡμέρα αὕτη, οἱ υἱοὶ Ἰὼβ καὶ αἱ θυγατέρες αὐτοῦ ἔπινον οἶνον ἐν τῇ οἰκίᾳ τοῦ ἀδελφοῦ αὐτῶν τοῦ πρεσβυτέρου.
\vs{14}Καὶ ἰδοὺ ἄγγελος ἦλθε πρὸς Ἰὼβ, καὶ εἶπεν αὐτῷ, τὰ ζεύγη τῶν βοῶν ἠροτρία, καὶ αἱ θήλειαι ὄνοι ἐβόσκοντο ἐχόμεναι αὐτῶν,
\vs{15}καὶ ἐλθόντες οἱ αἰχμαλωτεύοντες ᾐχμαλώτευσαν αὐτὰς, καὶ τοὺς παῖδας ἀπέκτειναν ἐν μαχαίραις· σωθεὶς δὲ ἐγὼ μόνος ἦλθον τοῦ ἀπαγγεῖλαί σοι.
\vs{16}Ἔτι τούτου λαλοῦντος, ἦλθεν ἕτερος ἄγγελος, καὶ εἶπε πρὸς Ἰὼβ, πῦρ ἔπεσεν ἐκ τοῦ οὐρανοῦ, καὶ κατέκαυσε τὰ πρόβατα, καὶ τοὺς ποιμένας κατέφαγεν ὁμοίως· σωθεὶς δέ ἐγὼ μόνος ἦλθον τοῦ ἀπαγγεῖλαί σοι.
\vs{17}Ἔτι τούτου λαλοῦντος, ἦλθεν ἕτερος ἄγγελος, καὶ εἶπεν πρὸς Ἰὼβ, οἱ ἱππεῖς ἐποίησαν ἡμῖν κεφαλᾶς τρεῖς, καὶ ἐκύκλωσαν τὰς καμήλους, καὶ ᾐχμαλώτευσαν αὐτὰς, καὶ τοὺς παῖδας ἀπέκτειναν ἐν μαχαίραις· ἐσώθην δὲ ἐγὼ μόνος, καὶ ἦλθον τοῦ ἀπαγγεῖλαί σοι.
\vs{18}Ἔτι τούτου λαλοῦντος, ἄλλος ἄγγελος ἔρχεται λέγων τῷ Ἰὼβ, τῶν υἱῶν σου καὶ τῶν θυγατέρων σου ἐσθίοντων καὶ πινόντων παρὰ τῷ ἀδελφῷ αὐτῶν τῷ πρεσβυτέρῳ,
\vs{19}ἐξαίφνης πνεῦμα μέγα ἐπῆλθεν ἐκ τῆς ἐρήμου, καὶ ἥψατο τῶν τεσσάρων γωνιῶν τῆς οἰκίας, καὶ ἔπεσεν ἡ οἰκία ἐπὶ τὰ παιδία σου, καὶ ἐτελεύτησαν· ἐσώθην δὲ ἐγὼ μόνος, καὶ ἦλθον τοῦ ἀπαγγεῖλαί σοι.

\vs{20}Οὕτως ἀναστὰς Ἰὼβ ἔῤῥξε τὰ ἱμάτια ἑαυτοῦ, καὶ ἐκείρατο τὴν κόμην τῆς κεφαλῆς, καὶ πεσὼν χαμαὶ, προσεκύνησε, καὶ εἶπεν,
\vs{21}αὐτὸς γυμνὸς ἐξῆλθον ἐκ κοιλίας μητρός μου, γυμνὸς καὶ ἀπελεύσομαι ἐκεῖ· ὁ Κύριος ἔδωκεν, ὁ Κύριος ἀφείλατο· ὡς τῷ Κυρίῳ ἔδοξεν, οὕτως ἐγένετο· εἴη τὸ ὄνομα Κυρίου εὐλογημένον.
\vs{22}Ἐν τούτοις πᾶσι τοῖς συμβεβηκόσιν αὐτῷ οὐδὲν ἥμαρτεν Ἰὼβ ἐναντίον τοῦ Κυρίου· καὶ οὐκ ἔδωκεν ἀφροσύνην τῷ Θεῷ.

\ch{2}
Ἐγένετο δὲ ὡς ἡ ἡμέρα αὕτη, καὶ ἦλθον οἱ ἄγγελοι τοῦ Θεοῦ παραστῆναι ἔναντι Κυρίου, καὶ ὁ διάβολος ἦλθεν ἐν μέσῳ αὐτῶν παραστῆναι ἐναντίον τοῦ Κυρίου.
\vs{2}Καὶ εἶπεν ὁ Κύριος τῷ διαβόλῳ, πόθεν σὺ ἔρχῃ; τότε εἶπεν ὁ διάβολος ἐνώπιον τοῦ Κυρίου, διαπορευθεὶς τὴν ὑπʼ οὐρανὸν, καὶ ἐμπεριπατήσας τὴν σύμπασαν, πάρειμι.
\vs{3}Εἶπε δὲ ὁ Κύριος πρὸς τὸν διάβολον, προσέσχες οὖν τῷ θεράποντί μου Ἰὼβ, ὅτι οὐκ ἔστι κατʼ αὐτὸν τῶν ἐπὶ τῆς γῆς; ἄνθρωπος ἄκακος, ἀληθινὸς, ἄμεμπτος, θεοσεβὴς, ἀπεχόμενος ἀπὸ παντὸς κακοῦ, ἔτι δὲ ἔχεται ἀκακίας· σὺ δὲ εἶπας ὑπάρχοντα αὐτοῦ διακενῆς ἀπολέσαι.
\vs{4}Ὑπολαβὼν δὲ ὁ διάβολος εἶπε τῷ Κυρίῳ, δέρμα ὑπὲρ δέρματος, ὅσα ὑπάρχει ἀνθρώπῳ ὑπέρ τῆς ψυχῆς αὐτοῦ ἐκτίσει.
\vs{5}Οὐ μὴν δὲ ἀλλὰ ἀποστείλας τὴν χεῖρά σου, ἅψαι τῶν ὀστῶν αὐτοῦ καὶ τῶν σαρκῶν αὐτοῦ· ἦ μὴν εἰς πρόσωπόν σε εὐλογήσει.
\vs{6}Εἶπε δὲ ὁ Κύριος τῷ διαβόλῳ, ἰδοὺ παραδίδωμί σοι αὐτόν· μόνον τὴν ψυχὴν αὐτοῦ διαφύλαξον.

\vs{7}Ἐξῆλθε δὲ ὁ διάβολος ἀπὸ τοῦ Κυρίου· καὶ ἔπαισε τὸν Ἰὼβ ἕλκει πονηρῷ ἀπὸ ποδῶν ἕως κεφαλῆς.
\vs{8}Καὶ ἔλαβεν ὄστρακον, ἵνα τὸν ἰχῶρα ξύῃ, καὶ ἐκάθητο ἐπὶ τῆς κοπρίας ἔξω τῆς πόλεως.

\vs{9}Χρόνου δὲ πολλοῦ προβεβηκότος, εἶπεν αὐτῷ ἡ γυνὴ αὐτοῦ, μέχρι τίνος καρτερήσεις, λέγων,
\vs{9a}ἰδοὺ ἀναμένω χρόνον ἔτι μικρὸν, προσδεχόμενος τὴν ἐλπίδα τῆς σωτηρίας μου;
\vs{9b}ἰδοὺ γὰρ ἠφάνισταί σου τὸ μνημόσυνον ἀπὸ τῆς γῆς· υἱοὶ καὶ θυγατέρες, ἐμῆς κοιλίας ὠδῖνες καὶ πόνοι, οὓς εἰς τὸ κενὸν ἐκοπίασα μετὰ μόχθων·
\vs{9c}σύ τε αὐτὸς ἐν σαπρίᾳ σκωλήκων κάθησαι διανυκτερεύων αἴθριος,
\vs{9d}κᾀγὼ πλανωμένη καὶ λάτρις τόπον ἐκ τόπου καὶ οἰκίαν ἐξ οἰκίας, προσδεχομένη τὸν ἥλιον πότε δύσεται, ἵνα ἀναπαύσωμαι τῶν μόχθων μου καὶ τῶν ὀδυνῶν αἵ με νῦν συνέχουσιν·
\vs{9e}ἀλλὰ εἶπόν τι ῥῆμα εἰς Κύριον, καὶ τελεύτα.
\vs{10}Ὁ δὲ ἐμβλέψας, εἶπεν αὐτῇ, ὥσπερ μία τῶν ἀφρόνων γυναικῶν ἐλάλησας· εἰ τὰ ἀγαθὰ ἐδεξάμεθα ἐκ χειρὸς Κυρίου, τὰ κακὰ οὐχ ὑποίσομεν;

Ἐν πᾶσι τούτοις τοῖς συμβεβηκόσιν αὐτῷ οὐδὲν ἥμαρτεν Ἰὼβ τοῖς χείλεσιν ἐναντίον τοῦ Θεοῦ.

\vs{11}Ἀκούσαντες δὲ οἱ τρεῖς φίλοι αὐτοῦ τὰ κακὰ πάντα τὰ ἐπελθόντα αὐτῷ, παρεγένοντο ἕκαστος ἐκ τῆς ἰδίας χώρας πρὸς αὐτόν· Ἐλιφὰζ ὁ Θαιμανῶν βασιλεὺς, Βαλδὰδ ὁ Σαυχέων τύραννος, Σωφὰρ Μιναίων βασιλεύς· καὶ παρεγένοντο πρὸς αὐτὸν ὁμοθυμαδὸν, τοῦ παρακαλέσαι καὶ ἐπισκέψασθαι αὐτόν.
\vs{12}Ἰδόντες δὲ αὐτὸν πόῤῥωθεν, οὐκ ἐπέγνωσαν· καὶ βοήσαντες φωνῇ μεγάλῃ ἔκλαυσαν, ῥήξαντες ἕκαστος τὴν ἑαυτοῦ στολὴν, καὶ καταπασάμενοι γῆν.
\vs{13}παρεκάθισαν αὐτῷ ἑπτὰ ἡμέρας καὶ ἑπτὰ νύκτας, καὶ οὐδεὶς αὐτῶν ἐλάλησεν· ἑώρων γὰρ τὴν πληγὴν δεινὴν οὖσαν καὶ μεγάλην σφόδρα.

\ch{3}
Μετὰ τοῦτο ἤνοιξεν Ἰὼβ τὸ στόμα αὐτοῦ,
\vs{2}καὶ κατηράσατο τὴν ἡμέραν αὐτοῦ, λέγων,

\vs{3}Ἀπόλοιτο ἡ ἡμέρα ἐν ᾗ ἐγεννήθην, καὶ ἡ νὺξ ἐκείνη ᾗ εἶπαν, Ἰδοὺ ἄρσεν.
\vs{4}Ἡ νὺξ ἐκείνη εἴη σκότος, καὶ μὴ ἀναζητήσαι αὐτὴν ὁ Κύριος ἄνωθεν, μηδὲ ἔλθοι εἰς αὐτὴν φέγγος·
\vs{5}Ἐκλάβοι δὲ αὐτὴν σκότος καὶ σκιὰ θανάτου, ἐπέλθοι ἐπʼ αὐτὴν γνόφος·
\vs{6}καταραθείη ἡ ἡμέρα καὶ ἡ νὺξ ἐκείνη, ἀπενέγκοιτο αὐτὴν σκότος· μὴ εἴη εἰς ἡμέρας ἐνιαυτοῦ, μηδὲ ἀριθμηθείη εἰς ἡμέρας μηνῶν.
\vs{7}Ἀλλὰ ἡ νὺξ ἐκείνη εἴη ὀδύνη, καὶ μὴ ἔλθοι ἐπʼ αὐτὴν εὐφροσύνη, μηδὲ χαρμονή·
\vs{8}Ἀλλὰ καταράσαιτο αὐτὴν ὁ καταρώμενος τὴν ἡμέραν ἐκείνην, ὁ μέλλων τὸ μέγα κῆτος χειρώσασθαι.
\vs{9}Σκοτωθείη τὰ ἄστρα τῆς νυκτὸς ἐκείνης· ὑπομείναι, καὶ εἰς φωτισμὸν μὴ ἔλθοι, καὶ μὴ ἴδοι Ἑωσφόρον ἀνατέλλοντα.
\vs{10}Ὅτι οὐ συνέκλεισε πύλας γαστρὸς μητρός μου, ἀπήλλαξε γὰρ ἂν πόνον ἀπὸ ὀφθαλμῶν μου.

\vs{11}Διατί γὰρ ἐν κοιλίᾳ οὐκ ἐτελεύτησα; ἐκ γαστρὸς δὲ ἐξῆλθον, καὶ οὐκ εὐθὺς ἀπωλόμην;
\vs{12}Ἱνατί δὲ συνήντησάν μοι τὰ γόνατα; ἱνατί δὲ μαστοὺς ἐθήλασα;
\vs{13}νῦν ἂν κοιμηθεὶς ἡσύχασα, ὑπνώσας δὲ ἀνεπαυσάμην
\vs{14}Νῦν ἂν κοιμηθεὶς ἡσύχασα, ὑπνώσας δὲ ἀνεπαυσάμην μετὰ βασιλέων βουλευτῶν γῆς οἳ ἐγαυριῶντο ἐπὶ ξίφεσιν,
\vs{15}ἢ μετὰ ἀρχόντων, ὧν πολὺς ὁ χρυσός, οἳ ἔπλησαν τοὺς οἴκους αὐτῶν ἀργυρίου·
\vs{16}Ἢ ὥσπερ ἔκτρωμα ἐκπορευόμενον ἐκ μήτρας μητρὸς, ἢ ὥσπερ νήπιοι, οἳ οὐκ εἶδον φῶς·
\vs{17}Ἐκεῖ ἀσεβεῖς ἐξέκαυσαν θυμὸν ὀργῆς, ἐκεῖ ἀνεπαύσαντο κατάκοποι τῷ σώματι.
\vs{18}Ὁμοθυμαδὸν δὲ οἱ αἰώνιοι οὐκ ἤκουσαν φωνὴν φορολόγου.
\vs{19}Μικρὸς καὶ μέγας ἐκεῖ ἐστι, καὶ θεράπων δεδοικὼς τὸν κύριον αὐτοῦ.

\vs{20}Ἱνατί γὰρ δέδοται τοῖς ἐν πικρίᾳ φῶς; ζωὴ δὲ ταῖς ἐν ὀδύναις ψυχαῖς,
\vs{21}οἳ ἱμείρονται τοῦ θανάτου, καὶ οὐ τυγχάνουσιν ἀνορύσσοντες ὥσπερ θησαυροὺς,
\vs{22}περιχαρεῖς δὲ ἐγένοντο ἐὰν κατατύχωσι;
\vs{23}Θάνατος ἀνδρὶ ἀνάπαυμα, συνέκλεισεν γὰρ ὁ Θεὸς κατʼ αὐτοῦ.
\vs{24}Πρὸ γὰρ τῶν σίτων μου στεναγμὸς ἥκει, δακρύω δὲ ἐγὼ συνεχόμενος φόβῳ.
\vs{25}Φόβος γὰρ ὃν ἐφρόντισα ἦλθέ μοι, καὶ ὃν ἐδεδοίκειν, συνήντησέν μοι.
\vs{26}Οὔτε εἰρήνευσα, οὔτε ἡσύχασα, οὔτε ἀνεπαυσάμην, ἦλθε δέ μοι ὀργή·

\ch{4}
Ὑπολαβὼν δὲ Ἐλιφὰς ὁ Θαιμανίτης, λέγει,

\vs{2}Μὴ πολλάκις σοι λελάληται ἐν κόπῳ; ἰσχὺν δὲ ῥημάτων σου τίς ὑποίσει;
\vs{3}Εἰ γὰρ σὺ ἐνουθέτησας πολλοὺς, καὶ χεῖρας ἀσθενοῦς παρεκάλεσας,
\vs{4}ἀσθενοῦντάς τε ἐξανέστησας ῥήμασι, γόνασί τε ἀδυνατοῦσι θάρσος περιέθηκας.
\vs{5}Νῦν δὲ ἥκει ἐπὶ σὲ πόνος καὶ ἥψατό σου, σὺ ἐσπούδασας.
\vs{6}Πότερον οὐχ ὁ φόβος σου ἐστὶν ἐν ἀφροσύνῃ, καὶ ἡ ἐλπίς σου καὶ ἡ κακία τῆς ὁδοῦ σου;
\vs{7}Μνήσθητι οὖν, τίς καθαρὸς ὢν ἀπώλετο, ἢ πότε ἀληθινοὶ ὁλόῤῥιζοι ἀπώλοντο;
\vs{8}Καθʼ ὃν τρόπον εἶδον τοὺς ἀροτριῶντας τὰ ἄτοπα, οἱ δὲ σπείροντες αὐτὰ ὀδύνας θεριοῦσιν ἑαυτοῖς.
\vs{9}Ἀπὸ προστάγματος Κυρίου ἀπολοῦνται, ἀπὸ δὲ πνεύματος ὀργῆς αὐτοῦ ἀφανισθήσονται.

\vs{10}Σθένος λέοντος, φωνὴ δὲ λεαίνης, γαυρίαμα δὲ δρακόντων ἐσβέσθη.
\vs{11}Μυρμηκολέων ὤλετο παρὰ τὸ μὴ ἔχειν βορὰν, σκύμνοι δὲ λεόντων ἔλιπον ἀλλήλους.

\vs{12}Εἰ δέ τι ῥῆμα ἀληθινὸν ἐγεγόνει ἐν λόγοις σου, οὐθὲν ἄν σοι τούτων κακὸν ἀπήντησε· πότερον οὐ δέξεταί μου τὸ οὖς ἐξαίσια παρʼ αὐτοῦ;
\vs{13}Φόβῳ δὲ καὶ ἤχῳ νυκτερινῇ ἐπιπίπτων φόβος ἐπʼ ἀνθρώπους,
\vs{14}φρίκη μοι συνήντησεν καὶ τρόμος, καὶ μεγάλως μου τὰ ὀστᾶ διέσεισε,
\vs{15}καὶ πνεῦμα ἐπὶ πρόσωπόν μου ἐπῆλθεν, ἔφριξαν δέ μου τρίχες καὶ σάρκες.
\vs{16}Ἀνέστην καὶ οὐκ ἐπέγνων, εἶδον καὶ οὐκ ἦν μορφὴ πρὸ ὀφθαλμῶν μου, ἀλλʼ ἢ αὖραν καὶ φωνὴν ἤκουον.
\vs{17}Τί γάρ; μὴ καθαρὸς ἔσται βροτὸς ἐναντίον τοῦ Κυρίου; ἢ ἀπὸ τῶν ἔργων αὐτοῦ ἄμεμπτος ἀνήρ;
\vs{18}Εἰ κατὰ παίδων αὐτοῦ οὐ πιστεύει, κατὰ δὲ ἀγγέλων αὐτοῦ σκολιόν τι ἐπενόησε.

\vs{19}Τοὺς δὲ κατοικοῦντας οἰκίας πηλίνας, ἐξ ὧν καὶ αὐτοὶ ἐκ τοῦ αὐτοῦ πηλοῦ ἐσμεν, ἔπαισεν αὐτοὺς σητὸς τρόπον,
\vs{20}καὶ ἀπὸ πρωΐθεν μέχρι ἑσπέρας οὐκ ἔτι εἰσὶ, παρὰ τὸ μὴ δύνασθαι αὐτοὺς ἑαυτοῖς βοηθῆσαι, ἀπώλοντο.
\vs{21}Ἐνεφύσησε γὰρ αὐτοῖς καὶ ἐξηράνθησαν, ἀπώλοντο παρὰ τὸ μὴ ἔχειν αὐτοὺς σοφίαν.

\ch{5}
Ἐπικαλέσαι δὲ εἴ τις σοι ὑπακούσεται, ἢ εἴ τινα ἀγγέλων ἁγίων ὄψῃ·
\vs{2}Καὶ γὰρ ἄφρονα ἀναιρεῖ ὀργὴ, πεπλανημένον δὲ θανατοῖ ζῆλος.
\vs{3}Ἐγὼ δὲ ἑώρακα ἄφρονας ῥίζαν βάλλοντας, ἀλλʼ εὐθέως ἐβρώθη αὐτῶν ἡ δίαιτα.
\vs{4}Πόῤῥω γένοιντο οἱ υἱοὶ αὐτῶν ἀπὸ σωτηρίας, κολαβρισθείησαν δὲ ἐπὶ θύραις ἡσσόνων, καὶ οὐκ ἔσται ὁ ἐξαιρούμενος.
\vs{5}Ἃ γὰρ ἐκεῖνοι συνήγαγον, δίκαιοι ἔδονται, αὐτοὶ δὲ ἐκ κακῶν οὐκ ἐξαίρετοι ἔσονται· ἐκσιφωνισθείη αὐτῶν ἡ ἰσχύς.
\vs{6}Οὐ γὰρ μὴ ἐξέλθῃ ἐκ τῆς γῆς κόπος, οὐδὲ ἐξ ὀρέων ἀναβλαστήσει πόνος.
\vs{7}Ἀλλὰ ἄνθρωπος γεννᾶται κόπῳ, νεοσσοὶ δὲ γυπὸς τὰ ὑψηλὰ πέτονται.

\vs{8}Οὐ μὴν δὲ ἀλλὰ ἐγὼ δεηθήσομαι Κυρίου, Κύριον δὲ τὸν πάντων δεσπότην ἐπικαλέσομαι,
\vs{9}τὸν ποιοῦντα μεγάλα καὶ ἀνεξιχνίαστα, ἔνδοξά τε καὶ ἐξαίσια, ὧν οὐκ ἔστιν ἀριθμὸς,
\vs{10}τὸν διδόντα ὑετὸν ἐπὶ τὴν γῆν, ἀποστέλλοντα ὕδωρ ἐπὶ τὴν ὑπʼ οὐρανὸν,
\vs{11}τὸν ποιοῦντα ταπεινοὺς εἰς ὕψος, καὶ ἀπολωλότας ἐξεγείροντα,
\vs{12}διαλλάσσοντα βουλὰς πανούργων, καὶ οὐ μὴ ποιήσουσιν αἱ χεῖρες αὐτῶν ἀληθές·
\vs{13}ὁ καταλαμβάνων σοφοὺς ἐν τῇ φρονήσει, βουλὴν δὲ πολυπλόκων ἐξέστησεν.
\vs{14}Ἡμέρας συναντήσεται αὐτοῖς σκότος, τὸ δὲ μεσημβρινὸν ψηλαφήσαισαν ἶσα νυκτὶ,
\vs{15}ἀπόλοιντο δὲ ἐν πολέμῳ· ἀδύνατος δὲ ἐξέλθοι ἐκ χειρὸς δυνάστου.
\vs{16}Εἴη δὲ ἀδυνάτῳ ἐλπὶς, ἀδίκου δὲ στόμα ἐμφραχθείη.

\vs{17}Μακάριος δὲ ἄνθρωπος ὃν ἤλεγξεν ὁ Κύριος, νουθέτημα δὲ παντοκράτορος μὴ ἀπαναίνου.
\vs{18}Αὐτὸς γὰρ ἀλγεῖν ποιεῖ, καὶ πάλιν ἀποκαθίστησιν· ἔπαισε, καὶ αἱ χεῖρες αὐτοῦ ἰάσαντο.
\vs{19}Ἑξάκις ἐξ ἀναγκῶν σε ἐξελεῖται, ἐν δὲ τῷ ἑβδόμῳ οὐ μὴ ἅψηταί σου κακόν.
\vs{20}Ἐν λιμῷ ῥύσεταί σε ἐκ θανάτου, ἐν πολέμῳ δὲ ἐκ χειρὸς σιδήρου λύσει σε.
\vs{21}Ἀπὸ μάστιγος γλώσσης σε κρύψει, καὶ οὐ μὴ φοβηθῇς ἀπὸ κακῶν ἐρχομένων.
\vs{22}Ἀδίκων καὶ ἀνόμων καταγελάσῃ· ἀπὸ δὲ θηρίων ἀγρίων οὐ μὴ φοβηθῇς·
\vs{23}Θῆρες γὰρ ἄγριοι εἰρηνεύσουσί σοι.
\vs{24}Εἶτα γνώσῃ ὅτι εἰρηνεύσει σου ὁ οἶκος· ἡ δὲ δίαιτα τῆς σκηνῆς σου οὐ μὴ ἁμάρτῃ.
\vs{25}Γνώσῃ δὲ ὅτι πολὺ τὸ σπέρμα σου, τὰ δὲ τέκνα σου ἔσται ὥσπερ τὸ παμβότανον τοῦ ἀγροῦ.
\vs{26}Ἐλεύσῃ δὲ ἐν τάφῳ ὥσπερ σῖτος ὥριμος κατὰ καιρὸν θεριζόμενος, ἢ ὥσπερ θιμωνία ἅλωνος καθʼ ὥραν συγκομισθεῖσα.

\vs{27}Ἰδοὺ ταῦτα οὕτως ἐξιχνιάσαμεν· ταῦτά ἐστιν ἃ ἀκηκόαμεν· σὺ δὲ γνῶθι σεαυτῷ, εἴ τι ἔπραξας.

\ch{6}
Ὑπολαβὼν δὲ Ἰὼβ, λέγει,

\vs{2}Εἰ γάρ τις ἱστῶν στήσαι μου τὴν ὀργὴν, τὰς δὲ ὀδύνας μου ἄραι ἐν ζυγῷ ὁμοθυμαδὸν,
\vs{3}καὶ δὴ ἄμμου παραλίας βαρυτέρα ἔσται· ἀλλʼ ὡς ἔοικε τὰ ῥήματά μου ἐστὶ φαῦλα.
\vs{4}Βέλη γὰρ Κυρίου ἐν τῷ σώματί μου ἐστὶν, ὧν ὁ θυμὸς αὐτῶν ἐκπίνει μου τὸ αἷμα· ὅταν ἄρξωμαι λαλεῖν, κεντοῦσί με.
\vs{5}Τί γάρ; μὴ διακενῆς κεκράξεται ὄνος ἄγριος, ἀλλʼ ἢ τὰ σῖτα ζητῶν; εἰ δὲ καὶ ῥήξει φωνὴν βοῦς ἐπὶ φάτνης ἔχων τὰ βρώματα;
\vs{6}Εἰ βρωθήσεται ἄρτος ἄνευ ἁλός; εἰ δὲ καὶ ἔστι γεῦμα ἐν ῥήμασι κενοῖς;
\vs{7}Οὐ δύναται γὰρ παύσασθαί μου ἡ ὀργή· βρόμον γὰρ ὁρῶ τὰ σῖτά μου ὥσπερ ὀσμὴν λέοντος.

\vs{8}Εἰ γὰρ δώῃ καὶ ἔλθοι μου ἡ αἴτησις, καὶ τὴν ἐλπίδα μου δώῃ ὁ Κύριος.
\vs{9}Ἀρξάμενος ὁ Κύριος τρωσάτω με, εἰς τέλος δὲ μή με ἀνελέτω.
\vs{10}Εἴη δέ μου πόλις τάφος, ἐφʼ ἧς ἐπὶ τειχέων ἡλλόμην, ἐπʼ αὐτῆς οὐ φείσομαι· οὐ γὰρ ἐψευσάμην ῥήματα ἅγια Θεοῦ μου.
\vs{11}Τίς γάρ μου ἡ ἰσχύς, ὅτι ὑπομένω; τίς μου ὁ χρόνος, ὅτι ἀνέχεταί μου ἡ ψυχή;
\vs{12}Μὴ ἰσχὺς λίθων ἡ ἰσχύς μου; ἢ αἱ σάρκες μου εἰσὶ χάλκεαι;
\vs{13}Ἢ οὐκ ἐπʼ αὐτῷ ἐπεποίθειν; βοήθεια δὲ ἀπʼ ἐμοῦ ἄπεστιν.

\vs{14}Ἀπείπατό με ἔλεος, ἐπισκοπὴ δὲ Κυρίου ὑπερεῖδέ με.
\vs{15}Οὐ προσεῖδόν με οἱ ἐγγύτατοί μου, ὥσπερ χειμάῤῥους ἐκλείπων, ἢ ὥσπερ κῦμα παρῆλθόν με
\vs{16}Οἵτινές με διευλαβοῦντο, νῦν ἐπιπεπτώκασί μοι ὥσπερ χιὼν ἢ κρύσταλλος πεπηγώς·
\vs{17}Καθὼς τακεῖσα θέρμης γενομένης, οὐκ ἐπεγνώσθη ὅπερ ἦν,
\vs{18}οὕτω κᾀγὼ κατελείφηθν ὑπὸ πάντων, ἀπωλόμην δὲ καὶ ἔξοικος ἐγενόμην.
\vs{19}Ἴδετε ὁδοὺς Θαιμανῶν, ἀτραποὺς, Σαβῶν οἱ διορῶντες.
\vs{20}Καὶ αἰσχύνην ὀφειλήσουσιν, οἱ ἐπὶ πόλεσι καὶ χρήμασι πεποιθότες.
\vs{21}Ἀτὰρ δὲ καὶ ὑμεῖς ἐπέβητέ μοι ἀνελεημόνως, ὥστε ἰδόντες τὸ ἐμὸν τραῦμα φοβήθητε.
\vs{22}Τί γάρ; μήτι ὑμᾶς ἢτησα, ἢ τῆς παρʼ ὑμῶν ἰσχύος ἐπιδέομαι,
\vs{23}ὥστε σῶσαί με ἑξ ἐχθρῶν, ἢ ἐκ χειρὸς δυναστῶν ῥύσασθαί με;

\vs{24}Διδάξατέ με, ἐγὼ δὲ κωφεύσω· εἴ τι πεπλάνημαι, φράσατέ μοι.
\vs{25}Ἀλλʼ ὡς ἔοικε φαῦλα ἀληθινοῦ ῥήματα, οὐ γὰρ παρʼ ὑμῶν ἰσχὺν αἰτοῦμαι.
\vs{26}Οὐδὲ ἔλεγχος ὑμῶν ῥήμασί με παύσει, οὐδὲ γὰρ ὑμῶν φθέγμα ῥήματος ἀνέξομαι.
\vs{27}Πλὴν ὅτι ἐπʼ ὀρφανῷ ἐπιπίπτετε, ἐνάλλεσθε δὲ ἐπὶ φίλῳ ὑμῶν.
\vs{28}Νυνὶ δὲ εἰσβλέψας εἰς πρόσωπα ὑμῶν, οὐ ψεύσομαι.
\vs{29}Καθίσατε δὴ καὶ μὴ εἴη ἄδικον, καὶ πάλιν τῷ δικαίῳ συνέρχεσθε.
\vs{30}Οὐ γάρ ἐστιν ἐν γλώσσῃ μουἄδικον, ἢ ὁ λάρυγξ μου οὐχὶ σύνεσιν μελετᾷ.

\ch{7}
Πότερον οὐχὶ πειρατήριόν ἐστιν ὁ βίος ἀνθρώπου ἐπὶ τῆς γῆς; καὶ ὥσπερ μισθίου αὐθημερινοῦ ἡ ζωὴ αὐτοῦ;
\vs{2}Ἢ ὥσπερ θεράπων δεδοικὼς τὸν Κύριον αὐτοῦ, καὶ τετευχὼς σκιᾶς; ἢ ὥσπερ μισθωτὸς ἀναμένων τὸν μισθὸν αὐτοῦ;
\vs{3}Οὕτως κᾀγὼ ὑπέμεινα μῆνας κενοὺς, νύκτες δὲ ὀδυνῶν δεδομέναι μοι εἰσίν.
\vs{4}Ἐὰν κοιμηθῶ, λέγω, πότε ἡμέρα; ὡς δʼ ἂν ἀναστῶ, πάλιν, πότε ἑσπέρα; πλήρης δὲ γίνομαι ὀδυνῶν ἀπὸ ἑσπέρας ἕως πρωΐ.
\vs{5}Φύρεται δέ μου τὸ σῶμα ἐν σαπρίᾳ σκωλήκων, τήκω δὲ βώλακας γῆς ἀπὸ ἰχῶρος ξύων.
\vs{6}Ὁ δὲ βίος μου ἔστιν ἐλαφρότερος λαλιᾶς, ἀπόλωλε δὲ ἐν κενῇ ἐλπίδι.
\vs{7}Μνήσθητι οὖν ὅτι πνεῦμά μου ἡ ζωὴ, καὶ οὐκ ἔτι ἐπανελεύσεται ὀφθαλμός μου ἰδεῖν ἀγαθόν.
\vs{8}Οὐ περιβλέψεταί με ὀφθαλμὸς ὁρῶντός με, οἱ ὀφθαλμοί σου ἐν ἐμοί, καὶ οὐκ ἔτι εἰμί·
\vs{9}Ὥσπερ νέφος ἀποκαθαρθὲν ἀπʼ οὐρανοῦ· ἐὰν γὰρ ἄνθρωπος καταβῇ εἰς ᾅδην, οὐκ ἔτι μὴ ἀναβῇ,
\vs{10}οὐδʼ οὐ μὴ ἐπιστρέψῃ εἰς τὸν ἴδιον οἶκον, οὐδʼ οὐ μὴ ἐπιγνῶ αὐτὸν ἔτι ὁ τόπος αὐτοῦ.
\vs{11}Ἀτὰρ οὖν οὐδὲ ἐγὼ φείσομαι τῷ στόματί μου, λαλήσω ἐν ἀνάγκῃ ὤν, ἀνοίξω πικρίαν ψυχῆς μου συνεχόμενος.

\vs{12}Πότερον θάλασσα εἰμὶ ἢ δράκων, ὅτι κατέταξας ἐπʼ ἐμὲ φυλακήν;
\vs{13}Εἴπα ὅτι παρακαλέσει με ἡ κλίνη μου, ἀνοίσω δὲ πρὸς ἐμαυτὸν ἰδίᾳ λόγον τῇ κοίτῃ μου.
\vs{14}Ἐκφοβεῖς με ἐνυπνίοις, καὶ ὁράμασί με καταπλήσσεις.
\vs{15}Ἀπαλλάξεις ἀπὸ πνεύματός μου τὴν ψυχήν μου, ἀπὸ δὲ θανάτου τὰ ὀστᾶ μου.
\vs{16}Οὐ γὰρ εἰς τὸν αἰῶνα ζήσομαι, ἵνα μακροθυμήσω· ἀπόστα ἀπʼ ἐμοῦ, κενὸς γάρ μου ὁ βίος.
\vs{17}Τί γάρ ἐστιν ἄνθρωπος, ὅτι ἐμεγάλυνας αὐτόν; ἢ ὅτι προσέχεις τὸν νοῦν εἰς αὐτόν;
\vs{18}Ἢ ἐπισκοπὴν αὐτοῦ ποιήσῃ ἕως τὸ πρωΐ; καὶ εἰς ἀνάπαυσιν αὐτὸν κρινεῖς;
\vs{19}Ἕως τίνος οὐκ ἐᾷς με, οὐδὲ προΐῃ με, ἕως ἂν καταπίω τὸν πτύελόν μου;
\vs{20}Εἰ ἐγὼ ἥμαρτον, τί δυνήσομαι πρᾶξαι, ὁ ἐπιστάμενος τὸν νοῦν τῶν ἀνθρώπων; διατί ἔθου με κατεντευκτήν σου, εἰμὶ δὲ ἐπὶ σοὶ φορτίον;
\vs{21}Διατί οὐκ ἐποιήσω τῆς ἀνομίας μου λήθην, καὶ καθαρισμὸν τῆς ἁμαρτίας μου; νυνὶ δὲ εἰς γῆν ἀπελεύσομαι, ὀρθρίζων δὲ οὐκ ἔτι εἰμί.

\ch{8}
Ὑπολαβὼν δὲ Βαλδὰδ ὁ Σαυχίτης, λέγει,

\vs{2}Μέχρι τίνος λαλήσεις ταῦτα, πνεῦμα πολυῤῥῆμον τοῦ στόματός σου;
\vs{3}Μὴ ὁ Κύριος ἀδικήσει κρίνων; ἢ ὁ τὰ πάντα ποιήσας ταράξει τὸ δίκαιον;
\vs{4}Εἰ οἱ υἱοί σου ἥμαρτον ἐναντίον αὐτοῦ, ἀπέστειλεν ἐν χειρὶ ἀνομίας αὐτῶν.

\vs{5}Σὺ δὲ ὄρθριζε πρὸς Κύριον παντοκράτορα δεόμενος.
\vs{6}Εἰ καθαρὸς εἶ καὶ ἀληθινὸς, δεήσεως ἐπακούσεταί σου, ἀποκαταστήσει δέ σοι δίαιταν δικαιοσύνης.
\vs{7}Ἔσται οὖν τὰ μὲν πρῶτά σου ὀλίγα, τὰ δὲ ἔσχατά σου ἀμύθητα.

\vs{8}Ἐπερώτησον γὰρ γενεὰν πρώτην, ἐξιχνίασον δὲ κατὰ γένος πατέρων·
\vs{9}Χθιζοὶ γάρ ἐσμεν, καὶ οὐκ οἴδαμεν· σκιὰ γάρ ἐστιν ἡμῶν ἐπὶ τῆς γῆς ὁ βίος·
\vs{10}Ἢ οὐχ οὗτοί σε διδάξουσιν καὶ ἀναγγελοῦσι, καὶ ἐκ καρδίας ἐξάξουσι ῥήματα;
\vs{11}Μὴ θάλλει πάπυρος ἄνευ ὕδατος, ἢ ὑψωθήσεται βούτομον ἄνευ πότου;
\vs{12}Ἔτι ὂν ἐπὶ ῥίζης, καὶ οὐ μὴ θερισθῇ, πρὸ τοῦ πιεῖν πᾶσα βοτάνη οὐχὶ ξηραίνεται;
\vs{13}Οὕτως τοίνυν ἔσται τὰ ἔσχατα πάντων τῶν ἐπιλανθανομένων τοῦ Κυρίου· ἐλπὶς γὰρ ἀσεβοῦς ἀπολεῖται·
\vs{14}Ἀοίκητος γὰρ αὐτοῦ ἔσται ὁ οἶκος· ἀράχνη δὲ αὐτοῦ ἀποβήσεται ἡ σκηνή.
\vs{15}Ἐὰν ὑπερείσῃ τὴν οἰκίαν αὐτοῦ, οὐ μὴ στῇ· ἐπιλαβομένου δὲ αὐτοῦ, οὐ μὴ ὑπομείνῃ·
\vs{16}Ὑγρὸς γάρ ἐστιν ὑπὸ ἡλίου· καὶ ἐκ σαπρίας αὐτοῦ ὁ ῥάδαμνος αὐτοῦ ἐξελεύσεται.
\vs{17}Ἐπὶ συναγωγὴν λίθων κοιμᾶται· ἐν δὲ μέσῳ χαλίκων ζήσεται.
\vs{18}Ἐὰν καταπίῃ, ὁ τόπος ψεύσεται αὐτόν· οὐχ ἑώρακας τοιαῦτα,
\vs{19}ὅτι καταστροφὴ ἀσεβοῦς τοιαύτη, ἐκ δὲ γῆς ἄλλον ἀναβλαστήσει.

\vs{20}Ὁ γὰρ Κύριος οὐ μὴ ἀποποιήσηται τὸν ἄκακον· πᾶν δὲ δῶρον ἀσεβοῦς οὐ δέξεται.
\vs{21}Ἀληθινῶν δὲ στόμα ἐμπλήσει γέλωτος, τὰ δὲ χείλη αὐτῶν ἐξομολογήσεως.
\vs{22}Οἱ δὲ ἐχθροὶ αὐτῶν ἐνδύσονται αἰσχύνην, δίαιτα δὲ ἀσεβοῦς οὐκ ἔσται.

\ch{9}
Ὑπολαβὼν δὲ Ἰὼβ, λέγει,

\vs{2}Ἐπʼ ἀληθείας οἶδα, ὅτι οὕτως ἐστί· πῶς γὰρ ἔσται δίκαιος βροτὸς παρὰ Κυρίῳ;
\vs{3}Ἐὰν γὰρ βούληται κριθῆναι αὐτῷ, οὐ μὴ ὑπακούσῃ αὐτῷ, ἵνα μὴ ἀντείπῃ πρὸς ἕνα λόγον αὐτοῦ ἐκ χιλίων.
\vs{4}Σοφὸς γάρ ἐστι διανοίᾳ, κραταιός τε καὶ μέγας· τίς σκληρὸς γενόμενος ἐναντίον αὐτοῦ ὑπέμεινεν;
\vs{5}Ὁ παλαιῶν ὄρη καὶ οὐκ οἴδασιν, ὁ καταστρέφων αὐτὰ ὀργῇ·
\vs{6}Ὁ σείων τὴν ὑπʼ οὐρανὸν ἐκ θεμελίων, οἱ δὲ στύλοι αὐτῆς σαλεύονται·
\vs{7}Ὁ λέγων τῷ ἡλίῳ καὶ οὐκ ἀνατέλλει, κατὰ δὲ ἄστρων κατασφραγίζει·
\vs{8}Ὁ τανύσας τὸν οὐρανὸν μόνος, καὶ περιπατῶν ὡς ἐπʼ ἐδάφους ἐπὶ θαλάσσης·
\vs{9}Ὁ ποιῶν Πλειάδα καὶ Ἕσπερον καὶ Ἀρκτοῦρον, καὶ ταμεῖα Νότου·
\vs{10}Ὁ ποιῶν μεγάλα καὶ ἀνεξιχνίαστα, ἔνδοξά τε καὶ ἐξαίσια, ὧν οὐκ ἔστιν ἀριθμός.

\vs{11}Ἐὰν ὑπερβῇ με, οὐ μὴ ἴδω· ἐὰν παρέλθῃ με, οὐδʼ ὡς ἔγνων.
\vs{12}Ἐὰν ἀπαλλάξῃ, τίς ἀποστρέψει; ἢ τίς ἐρεῖ αὐτῷ, τί ἐποίησας;
\vs{13}Αὐτὸς γὰρ ἀπέστραπται ὀργὴν, ὑπʼ αὐτοῦ ἐκάμφθησαν κήτη τὰ ὑπʼ οὐρανόν

\vs{14}Ἐὰν δέ μου ὑπακούσεται, ἤ διακρίνει τὰ ῥήματά μου.
\vs{15}Ἐὰν γὰρ ὦ δίκαιος, οὐκ εἰσακούσεταί μου, τοῦ κρίματος αὐτοῦ δεηθήσομαι·
\vs{16}Ἐάν τε καλέσω καὶ μὴ ὑπακούσῃ, οὐ πιστεύω ὅτι εἰσακήκοέ μου τῆς φωνῆς.

\vs{17}Μὴ γνόφῳ με ἐκτρίψῃ· πολλὰ δέ μου τὰ συντρίμματα πεποίηκε διακενῆς.
\vs{18}Οὐκ ἐᾷ γάρ με ἀναπνεῦσαι· ἐνέπλησε δέ με πικρίας,
\vs{19}ὅτι μὲν γὰρ ἰσχύει κράτει· τίς οὖν κρίματι αὐτοῦ ἀντιστήσεται;
\vs{20}Ἐὰν γὰρ ὦ δίκαιος, τὸ στόμα μου ἀσεβήσει· ἐάν τε ᾧ ἄμεμπτος, σκολιὸς ἀποβήσομαι.
\vs{21}Εἴτε γὰρ ἠσέβησα, οὐκ οἶδα τῇ ψυχῇ, πλὴν ἀφαιρεῖταί μου ἡ ζωή.

\vs{22}Διὸ εἶπον, μέγαν καὶ δυνάστην ἀπολλύει ὀργή,
\vs{23}ὅτι φαῦλοι ἐν θανάτῳ ἐξαισίῳ, ἀλλὰ δίκαιοι καταγελῶνται,
\vs{24}παραδέδονται γὰρ εἰς χεῖρας ἀσεβοῦς, πρόσωπα κριτῶν αὐτῆς συγκαλύπτει· εἰ δὲ μὴ αὐτός ἐστι, τίς ἐστιν;
\vs{25}Ὁ δὲ βίος μου ἐστὶν ἐλαφρότερος δρομέως· ἀπέδρασαν, καὶ οὐκ εἴδοσαν.
\vs{26}Ἢ καὶ ἐστι ναυσὶν ἴχνος ὁδοῦ, ἢ ἀετοῦ πετομένου ζητοῦντος βοράν;
\vs{27}Ἐάν τε γὰρ εἶπω, ἐπιλήσομαι λαλῶν, συγκύψας τῷ προσώπῳ στενάξω·
\vs{28}Σείομαι πᾶσι τοῖς μέλεσιν, οἴδα γὰρ ὅτι οὐκ ἀθῶόν με ἐάσεις.

\vs{29}Ἐπειδὴ δέ εἰμι ἀσεβὴς, διὰ τί οὐκ ἀπέθανον;
\vs{30}Ἐὰν γὰρ ἀπολούσωμαι χιόνι, καὶ ἀποκαθάρωμαι χερσὶ καθαραῖς,
\vs{31}ἱκανῶς ἐν ῥύπῳ με ἔβαψας, ἐβδελύξατο δέ με ἡ στολή.
\vs{32}Οὐ γὰρ εἰ ἄνθρωπος κατʼ ἐμὲ, ᾧ ἀντικρινοῦμαι, ἵνα ἔλθωμεν ὁμοθυμαδὸν εἰς κρίσιν.
\vs{33}Εἴθε ἦν ὁ μεσίτης ἡμῶν, καὶ ἐλέγχων, καὶ διακούων ἀναμέσον ἀμφοτέρων.
\vs{34}Ἀπαλλαξάτω ἀπʼ ἐμοῦ τὴν ῥάβδον, ὁ δὲ φόβος αὐτοῦ μή με στροβείτω,
\vs{35}καὶ οὐ μὴ φοβηθῶ, ἀλλὰ λαλήσω· οὐ γὰρ οὕτω συνεπίσταμαι.

\ch{10}
Καμνὼν τῇ ψυχῇ μου, στένων ἐπαφήσω ἐπʼ αὐτὸν τὰ ῥήματα μου· λαλήσω πικρίᾳ ψυχῆς μου συνεχόμενος,
\vs{2}καὶ ἐρῶ πρὸς Κύριον, μή με ἀσεβεῖν δίδασκε· καὶ διατί με οὕτως ἔκρινας;
\vs{3}Ἢ καλόν σοι ἐὰν ἀδικήσω; ὅτι ἀπείπω ἔργα χειρῶν σου, βουλῇ δὲ ἀσεβῶν προσέσχες.
\vs{4}Ἢ ὥσπερ βροτὸς ὁρᾷ, καθορᾷς; ἢ καθὼς ὁρᾷ ἄνθρωπος, βλέψῃ;
\vs{5}Ἢ ὁ βίος σου ἀνθρώπινός ἐστιν, ἢ τὰ ἔτη σου ἀνδρὸς,
\vs{6}ὅτι ἀνεζήτησας τὴν ἀνομίαν μου, καὶ τὰς ἁμαρτίας μου ἐξιχνίασας;
\vs{7}Οἶδας γὰρ ὅτι οὐκ ἠσέβησα· ἀλλὰ τίς ἐστιν ὁ ἐκ τῶν χειρῶν σου ἐξαιρούμενος;

\vs{8}Αἱ χεῖρές σου ἔπλασάν με καὶ ἐποίησάν με, μετὰ ταῦτα μεταβαλών με ἔπαισας.
\vs{9}Μνήσθητι, ὅτι πηλόν με ἔπλασας, εἰς δὲ γῆν με πάλιν ἀποστρέφεις.
\vs{10}Ἢ οὐχ ὥσπερ γάλα με ἤμελξας, ἐτύρωσας δέ με ἶσα τυρῷ;
\vs{11}Δέρμα δὲ καὶ κρέας με ἐνέδυσας, ὀστέοις δὲ καὶ νεύροις με ἔνειρας.
\vs{12}Ζωὴν δὲ καὶ ἔλεος ἔθου παρʼ ἐμοὶ, ἡ δὲ ἐπισκοπή σου ἐφυλαξέ μου τὸ πνεῦμα.
\vs{13}Ταῦτα ἔχων ἐν σεαυτῷ, οἶδα ὅτι πάντα δύνασαι· ἀδυνατεῖ δέ σοι οὐθέν.

\vs{14}Ἐάν τε γὰρ ἁμάρτω, φυλάσσεις με, ἀπὸ δὲ ἀνομίας οὐκ ἀθῶόν με πεποίηκας.
\vs{15}Ἐάν τε γὰρ ἀσεβήσω, οἴμοι· ἐὰν δὲ ὦ δίκαιος, οὐ δύναμαι ἀνακύψαι, πλήρης γὰρ ἀτιμίας εἰμί.
\vs{16}Ἀγρεύομαι γὰρ ὥσπερ λέων εἰς σφαγήν· πάλιν γὰρ μεταβαλὼν δεινῶς με ὀλέκεις,
\vs{17}ἐπανακαινίζων ἐπʼ ἐμὲ τὴν ἔτασίν μου· ὀργῇ δὲ μεγάλῃ μοι ἐχρήσω, ἐπήγαγες δὲ ἐπʼ ἐμὲ πειρατήρια.

\vs{18}Ἱνατί οὖν ἐκ κοιλίας με ἐξήγαγες, καὶ οὐκ ἀπέθανον, ὀφθαλμὸς δέ με οὐκ εἶδε,
\vs{19}καὶ ὥσπερ οὐκ ὢν ἐγενόμην; διατί γὰρ ἐκ γαστρὸς εἰς μνῆμα οὐκ ἀπηλλάγην;
\vs{20}Ἢ οὐκ ὀλίγος ἐστὶν ὁ βίος τοῦ χρόνου μου; ἔασόν με ἀναπαύσασθαι μικρὸν,
\vs{21}πρὸ τοῦ με πορευθῆναι ὅθεν οὐκ ἀναστρέψω, εἰς γῆν σκοτεινὴν καὶ γνοφερὰν,
\vs{22}εἰς γῆν σκότους αἰωνίου, οὗ οὐκ ἔστι φέγγος, οὐδὲ ὁρᾷν ζωὴν βροτῶν.

\ch{11}
Ὑπολαβὼν δὲ Σωφὰρ ὁ Μιναῖος, λέγει,

\vs{2}Ὁ τὰ πολλὰ λέγων, καὶ ἀντακούσεται· ἢ καὶ ὁ εὔλαλος οἴεται εἶναι δίκαιος;
\vs{3}εὐλογημένος γεννητὸς γυναικὸς ὀλιγόβιος. Μὴ πολὺς ἐν ῥήμασι γίνου, οὐ γάρ ἐστιν ὁ ἀντικρινόμενός σοι;
\vs{4}Μὴ γὰρ λέγε, ὅτι καθαρός εἰμι τοῖς ἔργοις καὶ ἄμεμπτος ἐναντίον αὐτοῦ.

\vs{5}Ἀλλὰ πῶς ἂν ὁ Κύριος λαλήσαι πρὸς σὲ, καὶ ἀνοίξει χείλη αὐτοῦ μετὰ σοῦ;
\vs{6}Εἶτα ἀναγγελεῖ σοι δύναμιν σοφίας· ὅτι διπλοῦς ἔσται τῶν κατὰ σέ· καὶ τότε γνώσῃ, ὅτι ἄξιά σοι ἀπέβη ἀπὸ Κυρίου ὧν ἡμάρτηκας.

\vs{7}Ἢ ἴχνος Κυρίου εὑρήσεις, ἢ εἰς τὰ ἔσχατα ἀφίκου ἃ ἐποίησεν ὁ παντοκράτωρ;
\vs{8}Ὑψηλὸς ὁ οὐρανὸς, καὶ τί ποιήσεις; βαθύτερα δὲ τῶν ἐν ᾅδου, τί οἶδας;
\vs{9}ἢ μακρότερα μέτρου γῆς, ἢ εὔρους θαλάσσης;

\vs{10}Ἐὰν δὲ καταστρέψῃ τὰ πάντα, τίς ἐρεῖ αὐτῷ, τί ἐποίησας;
\vs{11}Αὐτὸς γὰρ οἶδεν ἔργα ἀνόμων, ἰδὼν δὲ ἄτοπα οὐ παρόψεται.

\vs{12}Ἄνθρωπος δὲ ἄλλως νήχεται λόγοις· βροτὸς δὲ γεννητὸς γυναικὸς, ἶσα ὄνῳ ἐρημίτῃ.

\vs{13}Εἰ γὰρ σὺ καθαρὰν ἔθου τὴν καρδίαν σου, ὑπτιάζεις δὲ χεῖρας πρὸς αὐτὸν,
\vs{14}εἰ ἄνομόν τί ἐστιν ἐν χερσί σου, πόῤῥω ποίησον αὐτὸ ἀπὸ σοῦ, ἀδικία δὲ ἐν διαίτῃ σου μὴ αὐλισθήτω·
\vs{15}Οὕτως γὰρ ἀναλάμψει σου τὸ πρόσωπον, ὥσπερ ὕδωρ καθαρὸν, ἐκδύσῃ δὲ ῥύπον, καὶ οὐ μὴ φοβηθήσῃ·
\vs{16}Καὶ τὸν κόπον ἐπιλήσῃ, ὥσπερ κῦμα παρελθὸν, καὶ οὐ πτοηθήσῃ·
\vs{17}Ἡ δὲ εὐχή σου ὥσπερ Ἑωσφόρος, ἐκ δὲ μεσημβρίας ἀνατελεῖ σοι ζωή·
\vs{18}Πεποιθώς τε ἔσῃ, ὅτι ἐστί σοι ἐλπὶς, ἐκ δὲ μερίμνης καὶ φροντίδος ἀναφανεῖταί σοι εἰρήνη·
\vs{19}Ἡσυχάσεις γὰρ, καὶ οὐκ ἔσται ὁ πολεμῶν σε· μεταβαλόμενοι δὲ πολλοί σου δεηθήσονται.
\vs{20}Σωτηρία δὲ αὐτοὺς ἀπολείψει· ἡ γὰρ ἐλπὶς αὐτῶν ἀπώλεια, ὀφθαλμοὶ δὲ ἀσεβῶν τακήσονται.

\ch{12}
Ὑπολαβὼν δὲ Ἰὼβ, λέγει,

\vs{2}Εἶτα ὑμεῖς ἐστε ἄνθρωποι, ἢ μεθʼ ὑμῶν τελευτήσει σοφία;
\vs{3}Κᾀμοὶ μὲν καρδία καθʼ ὑμᾶς ἐστι.
\vs{4}Δίκαιος γὰρ ἀνὴρ καὶ ἄμεμπτος ἐγεννήθη εἰς χλεύασμα·
\vs{5}Εἰς χρόνον γὰρ τακτὸν ἡτοίμαστο πεσεῖν ὑπὸ ἄλλων, οἴκους τε αὐτοῦ ἐκπορθεῖσθαι ὑπὸ ἀνόμων·
\vs{6}οὐ μὴν δὲ ἀλλὰ μηδεὶς πεποιθέτω πονηρὸς ὢν ἀθῶος ἔσεσθαι, ὅσοι παροργίζουσι τὸν Κύριον, ὡς οὐχὶ καὶ ἔτασις αὐτῶν ἔσται.

\vs{7}Ἀλλὰ δὴ ἐρώτησον τετράποδα ἐάν σοι εἴπωσι, πετεινὰ δὲ οὐρανοῦ ἐάν σοι ἀπαγγείλωσιν.
\vs{8}Ἐκδιήγησαι γῇ, ἐάν σοι φράσῃ, καὶ ἐξηγήσονταί σοι οἱ ἰχθύες τῆς θαλάσσης.
\vs{9}Τίς οὖν οὐκ ἔγνω ἐν πᾶσι τούτοις, ὅτι χεὶρ Κυρίου ἐποίησε ταῦτα;
\vs{10}Εἰ μὴ ἐν χειρὶ αὐτοῦ ψυχὴ πάντων ζώντων, καὶ πνεῦμα παντὸς ἀνθρώπου.

\vs{11}Οὖς μὲν γὰρ ῥήματα διακρίνει, λάρυγξ δὲ σῖτα γεύεται.
\vs{12}Ἐν πολλῷ χρόνῳ σοφία, ἐν δὲ πολλῷ βίῳ ἐπιστήμη.
\vs{13}Παρʼ αὐτῷ σοφία καὶ δύναμις, αὐτῷ βουλὴ καὶ σύνεσις.
\vs{14}Ἐὰν καταβάλῃ, τίς οἰκοδομήσει; ἐὰν κλείσῃ κατʼ ἀνθρώπων, τίς ἀνοίξει;
\vs{15}Ἐὰν κωλύσῃ τὸ ὕδωρ, ξηρανεῖ τὴν γῆν· ἐὰν δὲ ἐπαφῇ, ἀπώλεσεν αὐτὴν καταστρέψας.
\vs{16}Παρʼ αὐτῷ κράτος καὶ ἰσχὺς, αὐτῷ ἐπιστήμη καὶ σύνεσις.
\vs{17}Διάγων βουλευτὰς αἰχμαλώτους, κριτὰς δὲ γῆς ἐξέστησε.
\vs{18}Καθιζάνων βασιλεῖς ἐπὶ θρόνους, καὶ περιεδησε ζώνῃ ὀσφύας αὐτῶν.
\vs{19}Ἐξαποστέλλων ἱερεῖς αἰχμαλώτους, δυνάστας δὲ γῆς κατέστρεψε.
\vs{20}Διαλλάσσων χείλη πιστῶν, σύνεσιν δὲ πρεσβυτέρων ἔγνω.
\vs{21}Ἐκχέων ἀτιμίαν ἐπʼ ἄρχοντας, ταπεινοὺς δὲ ἰάσατο.
\vs{22}Ἀνακαλύπτων βαθέα ἐκ σκότους, ἐξήγαγε δὲ εἰς φῶς σκιὰν θανάτου.
\vs{23}Πλανῶν ἔθνη καὶ ἀπολλύων αὐτὰ, καταστρωννύων ἔθνη καὶ καθοδηγῶν αὐτά.
\vs{24}Διαλλάσσων καρδίας ἀρχόντων γῆς· ἐπλάνησε δὲ αὐτοὺς ἐν ὁδῷ ᾗ οὐκ ᾔδεισαν.
\vs{25}Ψηλαφήσαισαν σκότος καὶ μὴ φῶς, πλανηθείησαν δὲ ὥσπερ ὁ μεθύων.

\ch{13}
Ἰδοὺ ταῦτα ἑώρακέ μου ὁ ὀφθαλμὸς, καὶ ἀκήκοέ μου τὸ οὖς.
\vs{2}Καὶ οἶδα ὅσα καὶ ὑμεῖς ἐπίστασθε, καὶ οὐκ ἀσυνετώτερός εἰμι ὑμῶν.

\vs{3}Οὐ μὴν δὲ ἀλλʼ ἐγὼ πρὸς Κύριον λαλήσω, ἐλέγξω δὲ ἐναντίον αὐτοῦ ἐὰν βούληται.
\vs{4}Ὑμεῖς δὲ ἐστὲ ἰατροὶ ἄδικοι, καὶ ἰαταὶ κακῶν πάντες.
\vs{5}Εἴη δὲ ὑμῖν κωφεῦσαι, καὶ ἀποβήσεται ὑμῖν σοφία.

\vs{6}Ἀκούσατε δὲ ἔλεγχον τοῦ στόματός μου, κρίσιν δὲ χειλέων μου προσέχετε.
\vs{7}Πότερον οὐκ ἔναντι Κυρίου λαλεῖτε, ἔναντι δὲ αὐτοῦ φθέγγεσθε δόλον;
\vs{8}Ἢ ὑποστελεῖσθε; ὑμεῖς δὲ αὐτοὶ κριταὶ γίνεσθε.
\vs{9}Καλὸν γὰρ ἐὰν ἐξιχνιάσῃ ὑμᾶς, εἰ γὰρ τὰ πάντα ποιοῦντες προστεθήσεσθε αὐτῷ,
\vs{10}οὐθὲν ἧττον ἐλέγξει ὑμᾶς· εἰ δὲ καὶ κρυφῇ πρόσωπα θαυμάσεσθε,
\vs{11}πότερον οὐχ ἡ δῖνα αὐτοῦ στροβήσει ὑμᾶς; ὁ φόβος δὲ παρʼ αὐτοῦ ἐπιπεσεῖται ὑμῖν,
\vs{12}ἀποβήσεται δὲ ὑμῶν τὸ γαυρίαμα ἶσα σποδῷ, τὸ δὲ σῶμα πήλινον.

\vs{13}Κωφεύσατε ἵνα λαλήσω, καὶ ἀναπαύσωμαι θυμοῦ,
\vs{14}ἀναλαβὼν τὰς σάρκας μου τοῖς ὀδοῦσι, ψυχὴν δέ μου θήσω ἐν χειρί.
\vs{15}Ἐάν με χειρώσηται ὁ δυνάστης, ἐπεὶ καὶ ἦρκται, ἦ μὴν λαλήσω καὶ ἐλέγξω ἐναντίον αὐτοῦ·
\vs{16}Καὶ τοῦτό μοι ἀποβήσεται εἰς σωτηρίαν, οὐ γὰρ ἐναντίον αὐτοῦ δόλος εἰσελεύσεται.
\vs{17}Ἀκούσατε ἀκούσατε τὰ ῥήματά μου, ἀναγγελῶ γὰρ ὑμῶν ἀκουόντων.
\vs{18}Ἰδοὺ ἐγὼ ἐγγύς εἰμι τοῦ κρίματός μου, οἶδα ἐγὼ ὅτι δίκαιος ἀναφανοῦμαι.
\vs{19}Τίς γάρ ἐστιν ὁ κριθησόμενός μοι, ὅτι νῦν κωφεύσω καὶ ἐκλείψω;

\vs{20}Δυεῖν δέ μοι χρήσῃ, τότε ἀπὸ τοῦ προσώπου σου οὐ κρυβήσομαι.
\vs{21}Τὴν χεῖρα ἀπʼ ἐμοῦ ἀπέχου, καὶ ὁ φόβος σου μή με καταπλησσέτω.
\vs{22}Εἶτα καλέσεις, ἐγὼ δέ σοι ὑπακούσομαι, ἢ λαλήσεις, ἐγὼ δέ σοι δώσω ἀνταπόκρισιν·
\vs{23}Πόσαι εἰσὶν αἱ ἁμαρτίαι μου καὶ ἀνομίαι μου; δίδαξόν με τίνες εἰσί.

\vs{24}Διατί ἀπʼ ἐμοῦ κρύπτῃ, ἥγησαι δέ με ὑπεναντίον σοι;
\vs{25}Ἢ ὡς φῦλλον κινούμενον ὑπὸ ἀνέμου εὐλαβηθήσῃ, ἢ ὡς χόρτῳ φερομένῳ ὑπὸ πνεύματος ἀντίκεισαί μοι;
\vs{26}Ὅτι κατέγραψας κατʼ ἐμοῦ κακὰ, περιέθηκας δέ μοι νεότητος ἁμαρτίας.
\vs{27}Ἔθου δέ μου τὸν πόδα ἐν κωλύματι, ἐφύλαξας δέ μου πάντα τὰ ἔργα, εἰς δὲ ῥίζας τῶν ποδῶν μου ἀφίκου.
\vs{28}Ὃ παλαιοῦται ἶσα ἀσκῷ, ἢ ὥσπερ ἱμάτιον σητόβρωτον.

\ch{14}
Βρότος γὰρ γεννητὸς γυναικὸς, ὀλιγόβιος, καὶ πλήρης ὀργῆς·
\vs{2}ἢ ὥσπερ ἄνθος ἀνθῆσαν ἐξέπεσεν, ἀπέδρα δὲ ὥσπερ σκιὰ, καὶ οὐ μὴ στῇ.
\vs{3}Οὐχὶ καὶ τούτου λόγον ἐποιήσω, καὶ τοῦτον ἐποίησας εἰσελθεῖν ἐν κρίματι ἐνώπιόν σου;
\vs{4}Τίς γὰρ καθαρὸς ἔσται ἀπὸ ῥύπου; ἀλλʼ οὐθεὶς,
\vs{5}ἐὰν καὶ μία ἡμέρα ὁ βίος αὐτοῦ ἐπὶ τῆς γῆς· ἀριθμητοὶ δὲ μῆνες αὐτοῦ παρʼ αὐτοῦ· εἰς χρόνον ἔθου, καὶ οὐ μὴ ὑπερβῇ.

\vs{6}Ἀπόστα ἀπʼ αὐτοῦ, ἵνα ἡσυχάσῃ, καὶ εὐδοκήσῃ τὸν βίον, ὥσπερ ὁ μισθωτός.

\vs{7}Ἔστι γὰρ δένδρῳ ἐλπὶς, ἐὰν γὰρ ἐκκοπῇ, ἔτι ἐπανθήσει, καὶ ὁ ῥάδαμνος αὐτοῦ οὐ μὴ ἐκλείπῃ.
\vs{8}Ἐὰν γὰρ γηράσῃ ἐν γῇ ἡ ῥίζα αὐτοῦ, ἐν δὲ πέτρᾳ τελευτήσῃ,
\vs{9}τὸ στέλεχος αὐτοῦ ἀπὸ ὀσμῆς ὕδατος ἀνθήσει, ποιήσει δὲ θερισμὸν, ὥσπερ νεόφυτον.
\vs{10}Ἀνὴρ δὲ τελευτήσας ᾤχετο, πεσὼν δὲ βροτὸς οὐκ ἔτι ἐστί.
\vs{11}Χρόνῳ γὰρ σπανίζεται θάλασσα, ποταμὸς δὲ ἐρημωθεὶς ἐξηράνθη.
\vs{12}Ἄνθρωπος δὲ κοιμηθεὶς οὐ μὴν ἀναστῇ ἕως ἂν ὁ οὐρανὸς οὐ μὴ συῤῥαφῇ, καὶ οὐκ ἐξυπνισθήσονται ἐξ ὕπνου αὐτῶν.

\vs{13}Εἰ γὰρ ὄφελον ἐν ᾅδῃ με ἐφύλαξας, ἔκρυψας δέ με ἕως ἂν παύσηταί σου ἡ ὀργὴ, καὶ τάξῃ μοι χρόνον ἐν ᾧ μνείαν μου ποιήσῃ.
\vs{14}Ἐὰν γὰρ ἀποθάνῃ ἄνθρωπος, ζήσεται συντελέσας ἡμέρας τοῦ βίου αὐτοῦ· ὑπομενῶ ἕως πάλιν γένωμαι;
\vs{15}Εἶτα καλέσεις, ἐγὼ δέ σοι ὑπακούσομαι, τὰ δὲ ἔργα τῶν χειρῶν σου μὴ ἀποποιοῦ.
\vs{16}Ἠρίθμησας δέ μου τὰ ἐπιτηδεύματα, καὶ οὐ μὴ παρέλθῃ σε οὐδὲν τῶν ἁμαρτιῶν μου;
\vs{17}Ἐσφράγισας δέ μου τὰς ἀνομίας ἐν βαλαντίῳ, ἐπεσημῄνω δὲ εἴτι ἄκων παρέβην.

\vs{18}Καὶ πλὴν ὄρος πίπτον διαπεσεῖται, καὶ πέτρα παλαιωθήσεται ἐκ τοῦ τόπου αὐτῆς.
\vs{19}Λίθους ἐλέαναν ὕδατα, καὶ κατέκλυσεν ὕδατα ὕπτια τοῦ χώματος τῆς γῆς· καὶ ὑπομονὴν ἀνθρώπου ἀπώλεσας.
\vs{20}Ὦσας αὐτὸν εἰς τέλος, καὶ ᾤχετο· ἐπέστησας αὐτῷ τὸ πρόσωπον, καὶ ἐξαπέστειλας,
\vs{21}πολλὼν δὲ γενομένων τῶν υἱῶν αὐτοῦ, οὐκ οἶδεν· ἐὰν δὲ ὀλίγοι γένωνται, οὐκ ἐπέσταται.
\vs{22}Ἀλλʼ ἢ αἱ σάρκες αὐτοῦ ἤλγησαν, ἡ δὲ ψυχὴ αὐτοῦ ἐπένθησεν.

\ch{15}
Ὑπολαβὼν δὲ Ἐλιφὰζ ὁ Θαιμανίτης, λέγει,

\vs{2}Πότερον σοφὸς ἀπόκρισιν δώσει συνέσεως πνεῦμα, καὶ ἐνέπλησε πόνον γαστρὸς,
\vs{3}ἐλέγχων ἐν ῥήμασιν οἷς οὐ δεῖ, καὶ ἐν λόγοις οἷς οὐδὲν ὄφελος;
\vs{4}Οὐ καὶ σὺ ἀπεποιήσω φόβον, συνετελέσω δὲ ῥήματα τοιαῦτα ἔναντι τοῦ Κυρίου;
\vs{5}Ἔνοχος εἶ ῥήμασι στόματός σου, οὐδὲ διέκρινας ῥήματα δυναστῶν.
\vs{6}Ἐλέγξαι σε τὸ σὸν στόμα καὶ μὴ ἐγὼ, τὰ δὲ χείλη σου καταμαρτυρήσουσι σου.

\vs{7}Τί γάρ; μὴ πρῶτος ἀνθρώπων ἐγεννήθης; ἢ πρὸ θινῶν ἐπάγης;
\vs{8}Ἢ σύνταγμα Κυρίου ἀκήκοας; ἢ συμβούλῳ σοι ἐχρήσατο ὁ Θεός; εἰς δέ σε ἀφίκετο σοφία;
\vs{9}Τί γὰρ οἶδας, ὃ οὐκ οἴδαμεν; ἢ τί συνίεις σὺ, ὃ οὐ καὶ ἡμεῖς;
\vs{10}Καί γε πρεσβύτης καί γε παλαιὸς ἐν ἡμῖν, βαρύτερος τοῦ πατρός σου ἡμέραις.
\vs{11}Ὀλίγα ὧν ἡμάρτηκας μεμαστίγωσαι, μεγάλως ὑπερβαλλόντως λελάληκας.

\vs{12}Τί ἐτόλμησεν ἡ καρδία σου; ἢ τί ἐπήνεγκαν οἱ ὀφθαλμοί σου,
\vs{13}ὅτι θυμὸν ἔῤῥηξας ἔναντι Κυρίου, ἐξήγαγες δὲ ἐκ στόματος ῥήματα τοιαῦτα;
\vs{14}Τίς γὰρ ὢν βροτὸς, ὅτι ἔσται ἅμεμπτος; ἢ ὡς ἐσόμενος δίκαιος γεννητὸς γυναικός;
\vs{15}Εἰ κατὰ ἁγίων οὐ πιστεύει, οὐρανὸς δὲ οὐ καθαρὸς ἐναντίον αὐτοῦ.
\vs{16}Ἔα δὲ ἐβδελυγμένος καὶ ἀκάθαρτος ἀνὴρ, πίνων ἀδικίας ἶσα ποτῷ.

\vs{17}Ἀναγγελῶ δέ σοι, ἄκουέ μου, ἃ δὴ ἑώρακα, ἀναγγελῶ σοι,
\vs{18}ἃ σοφοὶ ἐροῦσι, καὶ οὐκ ἔκρυψαν πατέρες αὐτῶν.
\vs{19}Αὐτοῖς μόνοις ἐδόθη ἡ γῆ, καὶ οὐκ ἐπῆλθεν ἀλλογενὴς ἐπʼ αὐτούς.
\vs{20}Πᾶς ὁ βίος ἀσεβοῦς ἐν φροντίδι, ἔτη δὲ ἀριθμητὰ δεδομένα δυνάστῃ.
\vs{21}Ὁ δὲ φόβος αὐτοῦ ἐν ὠσὶν αὐτοῦ· ὅταν δοκῇ ἤδη εἰρηνεύειν, ἥξει αὐτοῦ ἡ καταστροφή.
\vs{22}Μὴ πιστευέτω ἀποστραφῆναι ἀπὸ σκότους, ἐντέταλται γὰρ ἤδη εἰς χεῖρας σιδήρου,
\vs{23}κατατέτακται δὲ εἰς σῖτα γυψίν· οἶδε δὲ ἐν ἑαυτῷ ὅτι μένει εἰς πτῶμα· ἡμέρα δὲ σκοτεινὴ αὐτὸν στροβήσει,
\vs{24}ἀνάγκη δὲ καὶ θλίψις αὐτὸν καθέξει, ὥσπερ στρατηγὸς πρωτοστάτης πίπτων·
\vs{25}Ὅτι ᾖρκε χεῖρας ἐναντίον τοῦ Κυρίου, ἔναντι δὲ Κυρίου παντοκράτορος ἐτραχηλίασεν.
\vs{26}Ἔδραμε δὲ ἐναντίον αὐτοῦ ὕβρει, ἐν πάχει νώτου ἀσπίδος αὐτοῦ.
\vs{27}Ὅτι ἐκάλυψε τὸ πρόσωπον αὐτοῦ ἐν στέατι αὐτοῦ, καὶ ἐποίησε περιστόμιον ἐπὶ τῶν μηρίων.
\vs{28}Αὐλισθείη δὲ πόλεις ἐρήμους, εἰσέλθοι δὲ εἰς οἴκους ἀοικήτους· ἃ δὲ ἐκεῖνοι ἡτοίμασαν, ἄλλοι ἀποίσονται.

\vs{29}Οὔτε μὴ πλουτισθῇ, οὔτε μὴ μείνῃ αὐτοῦ τὰ ὑπάρχοντα· οὐ μὴ βάλῃ ἐπὶ τὴν γῆν σκιὰν,
\vs{30}οὐδὲ μὴ ἐκφύγῃ τὸ σκότος· τὸν βλαστὸν αὐτοῦ μαράναι ἄνεμος, ἐκπέσοι δὲ αὐτοῦ τὸ ἄνθος.
\vs{31}Μὴ πιστευέτω ὅτι ὑπομενεῖ, κενὰ γὰρ ἀποβήσεται αὐτῷ.
\vs{32}Ἡ τομὴ αὐτοῦ πρὸ ὥρας φθαρήσεται, καὶ ὁ ῥάδαμνος αὐτοῦ οὐ μὴ πυκάσῃ.
\vs{33}Τρυγηθείη δὲ ὡς ὄμφαξ πρὸ ὥρας, ἐκπέσοι δὲ ὡς ἄνθος ἐλαίας.
\vs{34}Μαρτύριον γὰρ ἀσεβοῦς θάνατος· πῦρ δὲ καύσει οἴκους δωροδεκτῶν·
\vs{35}Ἐν γαστρὶ δὲ λήψεται ὀδύνας, ἀποβήσεται δὲ αὐτῷ κενὰ, ἡ δὲ κοιλία αὐτοῦ ὑποίσει δόλον.

\ch{16}
Ὑπολαβὼν δὲ Ἰὼβ, λέγει,

\vs{2}Ἀκήκοα τοιαῦτα πολλὰ, παρακλήτορες κακῶν πάντες.
\vs{3}Τί γάρ; μὴ τάξις ἐστὶ ῥήμασι πνεύματος; ἢ τί παρενοχλήσει σοι ὅτι ἀποκρίνῃ;
\vs{4}Κᾀγὼ καθʼ ὑμᾶς λαλήσω· εἰ ὑπέκειτό γε ἡ ψυχὴ ὑμῶν ἀντὶ τῆς ἐμῆς, εἶτʼ ἐναλοῦμαι ὑμῖν ῥήμασι· κινήσω δὲ καθʼ ὑμῶν κεφαλήν.
\vs{5}Εἴη δὲ ἰσχὺς ἐν τῷ στόματί μου, κίνησιν δὲ χειλέων οὐ φείσομαι.

\vs{6}Ἐὰν γὰρ λαλήσω, οὐκ ἀλγήσω τὸ τραῦμα· ἐὰν δὲ καὶ σιωπήσω, τί ἔλαττον τρωθήσομαι;
\vs{7}Νῦν δὲ κατάκοπόν με πεποίηκε μωρὸν σεσηπότα, καὶ ἐπελάβου μου.
\vs{8}Εἰς μαρτύριον ἐγενήθη, καὶ ἀνέστη ἐν ἐμοὶ τὸ ψεῦδός μου, κατὰ πρόσωπόν μου ἀνταπεκρίθη.

\vs{9}Ὀργῇ χρησάμενος κατέβαλέ με, ἔβρυξεν ἐπʼ ἐμὲ τοὺς ὀδόντας, βέλη πειρατῶν αὐτοῦ ἐπʼ ἐμοὶ ἔπεσαν.
\vs{10}Ἀκίσιν ὀφθαλμῶν ἐνήλατο, ὀξεῖ ἔπαισέ με εἰς τὰ γόνατα, ὁμοθυμαδὸν δὲ κατέδραμον ἐπʼ ἐμοί.

\vs{11}Παρέδωκε γάρ με ὁ Κύριος εἰς χείρας ἀδίκων, ἐπὶ δὲ ἀσεβέσιν ἔῤῥιψέ με.
\vs{12}Εἰρηνεύοντα διεσκέδασέ με· λαβών με τῆς κόμης διέτιλε, κατέστησέ με ὥσπερ σκοπόν.
\vs{13}Ἐκύκλωσάν με λόγχαις βάλλοντες εἰς νεφρούς μου, οὐ φειδόμενοι ἐξέχεαν εἰς τὴν γῆν τὴν χολήν μου·
\vs{14}Κατέβαλόν με πτῶμα ἐπὶ πτώματι, ἔδραμον πρὸς μὲ δυνάμενοι.
\vs{15}Σάκκον ἔῤῥαψαν ἐπὶ βύρσης μου, τὸ δὲ σθένος μου ἐν γῇ ἐσβέσθη.
\vs{16}Ἡ γαστήρ μου συγκέκαυται ἀπὸ κλαυθμοῦ, ἐπὶ δὲ βλεφάροις μου σκιά.
\vs{17}Ἄδικον δὲ οὐδὲν ἦν ἐν χερσί μου, εὐχὴ δέ μου καθαρά.

\vs{18}Γῆ μὴ ἐπικαλύψῃς ἐφʼ αἵματι τῆς σαρκός μου, μηδὲ εἴη τόπος τῇ κραυγῇ μου.
\vs{19}Καὶ νῦν ἰδοὺ ἐν οὐρανοῖς ὁ μάρτυς μου, ὁ δὲ συνίστωρ μου ἐν ὑψίστοις.
\vs{20}Ἀφίκοιτό μου ἡ δέησις πρὸς Κύριον, ἔναντι δὲ αὐτοῦ στάζοι μου ὁ ὀφθαλμός.
\vs{21}Εἴη δὲ ἔλεγχος ἀνδρὶ ἔναντι Κυρίου, καὶ υἱῷ ἀνθρώπου τῷ πλησίον αὐτοῦ.
\vs{22}Ἔτη δὲ ἀριθμητὰ ἥκασιν, ὁδῷ δὲ ᾗ οὐκ ἐπαναστραφήσομαι, πορεύσομαι.

\ch{17}
Ὀλέκομαι πνεύματι φερόμενος, δέομαι δὲ ταφῆς, καὶ οὐ τυγχάνω.
\vs{2}Λίσσομαι κάμνων, καὶ τί ποιήσας;
\vs{3}ἔκλεψαν δέ μου τὰ ὑπάρχοντα ἀλλότριοι. Τίς ἐστιν οὕτος; τῇ χειρί μου συνδεθήτω.
\vs{4}Ὅτι καρδίαν αὐτῶν ἔκρυψας ἀπὸ φρονήσεως, διὰ τοῦτο οὐ μὴ ὑψώσῃς αὐτούς.
\vs{5}Τῇ μερίδι ἀναγγελεῖ κακίας· ὀφθαλμοὶ δὲ ἐφʼ υἱοῖς ἐτάκησαν.

\vs{6}Ἔθου δέ με θρύλλημα ἐν ἔθνεσι, γέλως δὲ αὐτοῖς ἀπέβην.
\vs{7}Πεπώρωνται γὰρ ἀπὸ ὀργῆς οἱ ὀφθαλμοί μου, πεπολιόρκημαι μεγάλως ὑπὸ πάντων.
\vs{8}Θαῦμα ἔσχεν ἀληθινοὺς ἐπὶ τούτῳ, δίκαιος δὲ ἐπὶ παρανόμῳ ἐπανασταίη.
\vs{9}Σχοίη δὲ πιστὸς τὴν ἑαυτοῦ ὁδὸν, καθαρὸς δὲ χεῖρας ἀναλάβοι θάρσος.
\vs{10}Οὐ μὴν δὲ ἀλλὰ πάντες ἐρείδετε καὶ δεῦτε δὴ, οὐ γὰρ εὑρίσκω ἐν ὑμῖν ἀληθές.

\vs{11}Αἱ ἡμέραι μου παρῆλθον ἐν βρόμῳ, ἐῤῥάγη δὲ τὰ ἄρθρα τῆς καρδίας μου.
\vs{12}Νύκτα εἰς ἡμέραν ἔθηκα, φῶς ἐγγὺς ἀπὸ προσώπου σκότους.
\vs{13}Ἐὰν γὰρ ὑπομείνω, ᾅδης μου ὁ οἶκος, ἐν δὲ γνοφῳ ἔστρωταί μου ἡ στρωμνή.
\vs{14}Θάνατον ἐπεκαλεσάμην πατέρα μου εἶναι, μητέρα δέ μου καὶ ἀδελφὴν σαπρίαν.
\vs{15}Ποῦ οὖν μου ἔτι ἐστὶν ἡ ἐλπὶς, ἢ τὰ ἀγαθά μου ὄψομαι;
\vs{16}Ἢ μετʼ ἐμοῦ εἰς ᾅδην καταβήσονται; ἢ ὁμοθυμαδὸν ἐπὶ χώματος καταβησόμεθα;

\ch{18}
Ὑπολαβὼν δὲ Βαλδὰδ ὁ Σαυχίτης, λέγει,

\vs{2}Μέχρι τίνος οὐ παύσῃ; ἐπίσχες, ἵνα καὶ αὐτοὶ λαλήσωμεν.
\vs{3}Διατί δὲ ὥσπερ τετράποδα σεσιωπήκαμεν ἐναντίον σου;
\vs{4}Κέχρηταί σοι ὀργή· τί γὰρ ἐὰν σὺ ἀποθάνῃς, ἀοίκητος ἡ ὑπʼ οὐρανόν; ἢ καταστραφήσεται ὄρη ἐκ θεμελίων;

\vs{5}Καὶ φῶς ἀσεβῶν σβεσθήσεται, καὶ οὐκ ἀποβήσεται αὐτῶν ἡ φλόξ.
\vs{6}Τὸ φῶς αὐτοῦ σκότος ἐν διαίτῃ, ὁ δὲ λύχνος ἐπʼ αὐτῷ σβεσθήσεται.
\vs{7}Θηρεύσαισαν ἐλάχιστοι τὰ ὑπάρχοντα αὐτοῦ· σφάλαι δὲ αὐτοῦ ἡ βουλή.
\vs{8}Ἐμβέβληται δὲ ὁ ποῦς αὐτοῦ ἐν παγίδι, ἐν δικτύῳ ἑλιχθείη.
\vs{9}Ἔλθοισαν δὲ ἐπʼ αὐτὸν παγίδες, κατισχύσει ἐπʼ αὐτὸν διψῶντας.
\vs{10}Κέκρυπται ἐν τῇ γῇ σχοινίον αὐτοῦ, καὶ ἡ σύλληψις αὐτοῦ ἐπὶ τρίβον.
\vs{11}Κύκλῳ ὀλέσαισαν αὐτὸν ὀδύναι· πολλοὶ δὲ περὶ πόδα αὐτοῦ ἔλθοισαν ἐν λιμῷ στενῷ·
\vs{12}πτῶμα δὲ αὐτῷ ἡτοίμασται ἐξαίσιον.
\vs{13}Βρωθείησαν αὐτοῦ κλῶνες ποδῶν, κατέδεται δὲ αὐτοῦ τὰ ὡραῖα θάνατος.
\vs{14}Ἐκραγείη δέ ἐκ διαίτης αὐτοῦ ἴασις, σχοίη δὲ αὐτὸν ἀνάγκη αἰτίᾳ βασιλικῇ.
\vs{15}Κατασκηνώσει ἐν τῇ σκηνῇ αὐτοῦ ἐν νυκτὶ αὐτοῦ, κατασπαρήσονται τὰ εὐπρεπῆ αὐτοῦ θείῳ.
\vs{16}Ὑποκάτωθεν αἱ ῥίζαι αὐτοῦ ξηρανθήσονται, καὶ ἐπάνωθεν ἐπιπεσεῖται θερισμὸς αὐτοῦ.
\vs{17}Τὸ μνημόσυνον αὐτοῦ ἀπόλοιτο ἐκ γῆς, καὶ ὑπάρξει ὄνομα αὐτῷ ἐπὶ πρόσωπον ἐξωτέρω.
\vs{18}Ἀπώσειεν αὐτὸν ἐκ φωτὸς εἰς σκότος.
\vs{19}Οὐκ ἔσται ἐπίγνωστος ἐν λαῷ αὐτοῦ, οὐδὲ σεσωσμένος ἐν τῇ ὑπʼ οὐρανὸν ὁ οἶκος αὐτοῦ.
\vs{20}Ἀλλʼ ἐν τοῖς αὐτοῦ ζήσονται ἕτεροι· ἐπʼ αὐτῷ ἐστέναξαν ἔσχατοι, πρώτους δὲ ἔσχε θαῦμα.

\vs{21}Οὗτοί εἰσιν οἱ οἶκοι ἀδίκων, οὗτος δὲ ὁ τόπος τῶν μὴ εἰδότων τὸν Κύριον.

\ch{19}
Ὑπολαβὼν δὲ Ἰὼβ, λέγει,

\vs{2}Ἕως τίνος ἔγκοπον ποιήσετε ψυχήν μου, καὶ καθαιρεῖτέ με λόγοις; γνῶτε μόνον ὅτι ὁ Κύριος ἐποίησέ με οὕτως.
\vs{3}Καταλαλεῖτέ μου, οὐκ αἰσχυνόμενοί με ἐπίκεισθέ μοι.
\vs{4}Ναὶ δὴ ἐπʼ ἀληθείας ἐγὼ ἐπλανήθην, παρʼ ἐμοὶ δὲ αὐλίζεται πλάνος·
\vs{4a}λαλῆσαι ῥῆματα ἃ οὐκ ἔδει, τὰ δὲ ῥήματά μου πλανᾶται καὶ οὐκ ἐπὶ καιροῦ.
\vs{5}Ἔα δὲ, ὅτι ἐπʼ ἐμοὶ μεγαλύνεσθε, ἐνάλλεσθε δέ μοι ὀνείδει.
\vs{6}Γνῶτε οὖν ὅτι Κύριός ἐστιν ὁ ταράξας, ὀχύρωμα δὲ αὐτοῦ ἐπʼ ἐμὲ ὕψωσεν.
\vs{7}Ἰδοὺ γελῶ ὀνείδει, οὐ λαλήσω· κεκράξομαι, καὶ οὐδαμοῦ κρίμα.
\vs{8}Κύκλῳ περιῳκοδόμημαι, καὶ οὐ μὴ διαβῶ· ἐπὶ πρόσωπόν μου σκότος ἔθετο.
\vs{9}Τὴν δὲ δόξαν ἀπʼ ἐμοῦ ἐξέδυσεν, ἀφεῖλε δὲ στέφανον ἀπὸ κεφαλῆς μου.
\vs{10}Διέσπασέ με κύκλῳ καὶ ᾠχόμην, ἐξέκοψε δὲ ὥσπερ δένδρον τὴν ἐλπίδα μου.
\vs{11}Δεινῶς δέ μοι ὀργῇ ἐχρήσατο, ἡγήσατο δέ με ὥσπερ ἐχθρόν.
\vs{12}Ὁμοθυμαδὸν δὲ ἦλθον τὰ πειρατήρια αὐτοῦ ἐπʼ ἐμοὶ, ταῖς ὁδοῖς μου ἐκύκλωσαν ἐγκάθετοι.

\vs{13}Ἀπʼ ἐμοῦ ἀδελφοί μου ἀπέστησαν, ἔγνωσαν ἀλλοτρίους ἢ ἐμέ· φίλοι δέ μου ἀνελεήμονες γεγόνασιν·
\vs{14}Οὐ προσεποιήσαντό με οἱ ἐγγύτατοί μου, καὶ οἱ εἰδότες μου τὸ ὄνομα ἐπελάθοντό μου.
\vs{15}Γείτονες οἰκίας, θεράπαιναί τε μοῦ, ἀλλογενὴς ἤμην ἐναντίον αὐτῶν.
\vs{16}Θεράποντά μου ἐκάλεσα, καὶ οὐχ ὑπήκουσε· στόμα δέ μου ἐδέετο.
\vs{17}Καὶ ἱκέτευον τὴν γυναῖκά μου, προσεκαλούμην δὲ καλακευων υἱοὺς παλλακίδων μου·
\vs{18}Οἱ δὲ εἰς τὸν αἰῶνά με ἀπεποιήσαντο, ὅταν ἀναστῶ, κατʼ ἐμοῦ λαλοῦσιν.
\vs{19}Ἐβδελύξαντό με οἱ ἰδόντες με· οὓς δὴ ἠγαπήκειν, ἐπανέστησάν μοι.
\vs{20}Ἐν δέρματί μου ἐσάπησαν αἱ σάρκες μου, τὰ δὲ ὀστᾶ μου ἐν ὀδοῦσιν ἔχεται.
\vs{21}Ἐλεήσατέ με, ἐλεήσατέ με, ὦ φίλοι, χεὶρ γὰρ Κυρίου ἡ ἁψαμένη μου ἐστί.
\vs{22}Διατί με διώκετε ὥσπερ καὶ ὁ Κύριος; ἀπὸ δὲ σαρκῶν μου οὐκ ἐμπίπλασθε;

\vs{23}Τίς γὰρ ἂν δοίη γραφῆναι τὰ ῥήματά μου, τεθῆναι δὲ αὐτὰ ἐν βιβλίῳ εἰς τὸν αἰῶνα,
\vs{24}ἐν γραφείῳ σιδηρῷ καὶ μολίβῳ, ἢ ἐν πέτραις ἐγγλυφῆναι;
\vs{25}Οἶδα γὰρ ὅτι ἀένναός ἐστιν ὁ ἐκλύειν με μέλλων,
\vs{26}ἐπὶ γῆς ἀναστῆσαι τὸ δέρμα μου τὸ ἀναντλοῦν ταῦτα· παρὰ γὰρ Κυρίου ταῦτά μοι συνετελέσθη,
\vs{27}ἃ ἐγὼ ἐμαυτῷ συνεπίσταμαι, ἃ ὁ ὀφθαλμός μου ἑώρακε, καὶ οὐκ ἄλλος, πάντα δέ μοι συντετέλεσται ἐν κόλπῳ.

\vs{28}Εἰ δὲ καὶ ἐρεῖτε, τί ἐροῦμεν ἔναντι αὐτοῦ, καὶ ῥίζαν λόγου εὑρήσομεν ἐν αὐτῷ;
\vs{29}Εὐλαβήθητε δὴ καὶ ὑμεῖς ἀπὸ ἐπικαλύμματος, θυμὸς γὰρ ἐπʼ ἀνόμους ἐπελεύσεται· καὶ τότε γνώσονται, ποῦ ἐστιν αὐτῶν ἡ ὕλη

\ch{20}
Ὑπολαβὼν δὲ Σωφὰρ ὁ Μιναῖος, λέγει,

\vs{2}Οὐχ οὕτως ὑπελάμβανον ἀντερεῖν σε ταῦτα, καὶ οὐχὶ συνίετε μᾶλλον ἢ καὶ ἐγώ.
\vs{3}Παιδείαν ἐντροπῆς μου ἀκούσομαι, καὶ πνεῦμα ἐκ τῆς συνέσεως ἀποκρίνεταί μοι.

\vs{4}Μὴ ταῦτα ἔγνως ἀπὸ τοῦ ἔτι, ἀφʼ οὗ ἐτέθη ἄνθρωπος ἐπὶ τῆς γῆς;
\vs{5}Εὐφροσύνη δὲ ἀσεβῶν πτῶμα ἐξαίσιον, χαρμονὴ δὲ παρανόμων ἀπώλεια·
\vs{6}ἐὰν ἀναβῇ εἰς οὐρανὸν αὐτοῦ τὰ δῶρα, ἡ δὲ θυσία αὐτοῦ νεφῶν ἅψηται.
\vs{7}Ὅταν γὰρ δοκῇ ἤδη κατεστηρίχθαι, τότε εἰς τέλος ἀπολεῖται· οἱ δὲ εἰδότες αὐτὸν ἐροῦσι, ποῦ ἐστιν;
\vs{8}Ὥσπερ ἐνύπνιον ἐκπετασθὲν οὐ μὴ εὑρεθῇ, ἔπτη δὲ ὥσπερ φάσμα νυκτερινόν.
\vs{9}Ὀφθαλμὸς παρέβλεψε, καὶ οὐ προσθήσει, καὶ οὐκ ἔτι προσνοήσει αὐτὸν ὁ τόπος αὐτοῦ.
\vs{10}Τοὺς υἱοὺς αὐτοῦ ὀλέσαισαν ἥττονες, αἱ δὲ χεῖρες αὐτοῦ πυρσεύσαισαν ὀδύνας.
\vs{11}Ὀστᾶ αὐτοῦ ἐνεπλήσθησαν νεότητος αὐτοῦ, καὶ μετʼ αὐτοῦ ἐπὶ χώματος κοιμηθήσεται.

\vs{12}Ἐὰν γλυκανθῇ ἐν στόματι αὐτοῦ κακία, κρύψει αὐτὴν ὑπὸ τὴν γλῶσσαν αὐτοῦ,
\vs{13}οὐ φείσεται αὐτῆς, καὶ οὐκ ἐγκαταλείψει αὐτήν· καὶ συνάξει αὐτὴν ἐν μέσῳ τοῦ λάρυγγος αὐτοῦ,
\vs{14}καὶ οὐ μὴ δυνηθῇ βοηθῆσαι ἑαυτῷ· χολὴ ἀσπίδος ἐν γαστρὶ αὐτοῦ.

\vs{15}Πλοῦτος ἀδίκως συναγόμενος ἐξεμεθήσεται, ἐξ οἰκίας αὐτοῦ ἐξελκύσει αὐτὸν ἄγγελος.
\vs{16}Θυμὸν δὲ δρακόντων θηλάσειεν, ἀνέλοι δὲ αὐτὸν γλῶσσα ὄφεως.
\vs{17}Μὴ ἴδοι ἄμελξιν νομάδων, μηδὲ νομὰς μέλιτος καὶ βουτύρου.
\vs{18}Εἰς κενὰ καὶ μάταια ἐκοπίασε, πλοῦτον ἐξ οὗ οὐ γεύσεται· ὥσπερ στρίφνος ἀμάσητος, ἀκατάποτος.
\vs{19}Πολλῶν γὰρ δυνατῶν οἴκους ἔθλασε· δίαιταν δὲ ἥρπασε, καὶ οὐκ ἔστησεν.
\vs{20}Οὐκ ἔστιν αὐτοῦ σωτηρία τοῖς ὑπάρχουσιν, ἐν ἐπιθυμίᾳ αὐτοῦ οὐ σωθήσεται.
\vs{21}Οὐκ ἔστιν ὑπόλειμμα τοῖς βρώμασιν αὐτοῦ, διὰ τοῦτο οὐκ ἀνθήσει αὐτοῦ τὰ ἀγαθά.
\vs{22}Ὅταν δὲ δοκῇ ἤδη πεπληρῶσθαι, θλιβήσεται, πᾶσα δὲ ἀνάγκη ἐπʼ αὐτὸν ἐπελεύσεται.

\vs{23}Εἴ πῶς εἶ πληρῶσαι γαστέρα αὐτοῦ, ἐπαποστείλαι ἐπʼ αὐτὸν θυμὸν ὀργῆς, νίψαι ἐπʼ αὐτὸν ὀδύνας,
\vs{24}καὶ οὐ μὴ σωθῇ ἐκ χειρὸς σιδήρου· τρώσαι αὐτὸν τόξον χάλκειον.
\vs{25}Διεξέλθοι δὲ διὰ σώματος αὐτοῦ βέλος, ἄστρα δὲ ἐν διαίταις αὐτοῦ· περιπατήσαισαν ἐπʼ αὐτῷ φόβοι,
\vs{26}πᾶν δὲ σκότος αὐτῷ ὑπομείναι· κατέδεται αὐτὸν πῦρ ἄκαυστον, κακώσαι δὲ αὐτοῦ ἐπήλυτος τὸν οἶκον.
\vs{27}Ἀνακάλυψαι δὲ αὐτοῦ ὁ οὐρανὸς τὰς ἀνομίας, γῆ δὲ ἐπανασταίη αὐτῷ.
\vs{28}Ἑλκύσαι τὸν οἶκον αὐτοῦ ἀπώλεια εἰς τέλος, ἡμέρα ὀργῆς ἐπέλθοι αὐτῷ.
\vs{29}Αὕτη ἡ μέρις ἀνθρώπου ἀσεβοῦς παρὰ Κυρίου, καὶ κτῆμα ὑπαρχόντων αὐτῷ παρὰ τοῦ ἐπισκόπου.

\ch{21}
Ὑπολαβὼν δὲ Ἰὼβ, λέγει,

\vs{2}Ἀκούσατε ἀκούσατέ μου τῶν λόγων, ἵνα μὴ ᾖ μοι παρʼ ὑμῶν αὕτη ἡ παράκλησις.
\vs{3}Ἄρατέ με, ἐγὼ δὲ λαλήσω, εἶτʼ οὐ καταγελάσετέ μου.
\vs{4}Τί γάρ; μὴ ἀνθρώπου μου ἡ ἔλεγξις; ἢ διὰ τί οὐ θυμωθήσομαι;
\vs{5}Εἰσβλέψαντες εἰς ἐμὲ θαυμάσετε, χεῖρα θέντες ἐπὶ σιαγόνι.

\vs{6}Ἐάν τε γὰρ μνησθῶ, ἐσπούδακα· ἔχουσι δέ μου τὰς σάρκας ὀδύναι.
\vs{7}Διὰ τί ἀσεβεῖς ζῶσι, πεπαλαίωνται δὲ καὶ ἐν πλούτῳ;
\vs{8}Ὁ σπόρος αὐτῶν κατὰ ψυχὴν, τὰ δὲ τέκνα αὐτῶν ἐν ὀφθαλμοῖς.
\vs{9}Οἱ οἶκοι αὐτῶν εὐθηνοῦσι, φόβος δὲ οὐδαμοῦ, μάστιξ δὲ παρὰ Κυρίου οὐκ ἔστιν ἐπʼ αὐτοῖς.
\vs{10}Ἡ βοῦς αὐτῶν οὐκ ὠμοτόκησε, διεσώθη δὲ αὐτῶν ἐν γαστρὶ ἔχουσα καὶ οὐκ ἔσφαλε.
\vs{11}Μένουσι δὲ ὡς πρόβατα αἰώνια, τὰ δὲ παιδία αὐτῶν προσπαίζουσιν,
\vs{12}ἀναλαβόντες ψαλτήριον καὶ κιθάραν, καὶ εὐφραίνονται φωνῇ ψαλμοῦ.
\vs{13}Συνετέλεσαν δὲ ἐν ἀγαθοῖς τὸν βίον αὐτῶν, ἐν δὲ ἀναπαύσει ᾅδου ἐκοιμήθησαν.
\vs{14}Λέγει δὲ Κυρίῳ, ἀπόστα ἀπʼ ἐμοῦ, ὁδούς σου εἰδέναι οὐ βούλομαι.
\vs{15}Τί ἱκανὸς, ὅτι δουλεύσομεν αὐτῷ; καὶ τίς ὠφέλεια, ὅτι ἀπαντήσομεν αὐτῷ;

\vs{16}Ἐν χερσὶ γὰρ ἦν αὐτῶν τὰ ἀγαθὰ, ἔργα δὲ ἀσεβῶν οὐκ ἐφορᾷ.
\vs{17}Οὐ μὴν δὲ ἀλλὰ καὶ ἀσεβῶν λύχνος σβεσθήσεται, ἐπελεύσεται δὲ αὐτοῖς ἡ καταστροφὴ, ὠδῖνες δὲ αὐτοὺς ἕξουσιν ἀπὸ ὀργῆς.
\vs{18}Ἔσονται δὲ ὥσπερ ἄχυρα ὑπʼ ἀνέμου, ἢ ὥσπερ κονιορτὸς ὃν ὑφείλετο λαίλαψ.
\vs{19}Ἐκλείποι υἱοὺς τὰ ὑπάρχοντα αὐτοῦ, ἀνταποδώσει πρὸς αὐτὸν καὶ γνώσεται.
\vs{20}Ἴδοισαν οἱ ὀφθαλμοὶ αὐτοῦ τὴν ἑαυτοῦ σφαγὴν, ἀπὸ δὲ Κυρίου μὴ διασωθείη.
\vs{21}Ὅτι τὸ θέλημα αὐτοῦ ἐν οἴκῳ αὐτοῦ μετʼ αὐτοῦ, καὶ ἀριθμοὶ μηνῶν αὐτοῦ διῃρέθησαν.

\vs{22}Πότερον οὐχὶ ὁ Κύριός ἐστιν ὁ διδάσκων σύνεσιν καὶ ἐπιστήμην; αὐτὸς δὲ φόνους διακρίνει;
\vs{23}Οὗτος ἀποθανεῖται ἐν κράτει ἁπλοσύνης αὐτοῦ, ὅλος δὲ εὐπαθῶν καὶ εὐθηνῶν.
\vs{24}Τὰ δὲ ἔγκατα αὐτοὺ πλήρη στέατος, μυελὸς δὲ αὐτοῦ διαχεῖται.
\vs{25}Ὁ δὲ τελευτᾷ ὑπὸ πικρίας ψυχῆς, οὐ φαγὼν οὐδὲν ἀγαθόν.
\vs{26}Ὁμοθυμαδὸν δὲ ἐπὶ γῆς κοιμῶνται, σαπρία δὲ αὐτοὺς ἐκάλυψεν.

\vs{27}Ὥστε οἶδα ὑμᾶς, ὅτι τόλμῃ ἐπίκεισθέ μοι·
\vs{28}Ὥστε ἐρεῖτε, ποῦ ἐστιν οἶκος ἄρχοντος; καὶ ποῦ ἐστιν ἡ σκέπη τῶν σκηνωμάτων τῶν ἀσεβῶν;
\vs{29}Ἐρωτήσατε παραπορευομένους ὁδὸν, καὶ τὰ σημεῖα αὐτῶν οὐκ ἀπαλλοτριώσετε.
\vs{30}Ὅτι εἰς ἡμέραν ἀπωλείας κουφίζεται ὁ πονηρὸς, εἰς ἡμέραν ὀργῆς αὐτοῦ ἀπαχθήσονται.
\vs{31}Τίς ἀπαγγελεῖ ἐπὶ προσώπου αὐτοῦ τὴν ὁδὸν αὐτοῦ, καὶ αὐτὸς ἐποίησε; τίς ἀνταποδώσει αὐτῷ;
\vs{32}Καὶ αὐτὸς εἰς τάφους ἀπηνέχθη, καὶ αὐτὸς ἐπὶ σωρῶν ἠγρύπνησεν.
\vs{33}Ἐγλυκάνθησαν αὐτῷ χάλικες χειμάῤῥου, καὶ ὀπίσω αὐτοῦ πᾶς ἄνθρωπος ἀπελεύσεται, καὶ ἔμπροσθεν αὐτοῦ ἀναρίθμητοι.
\vs{34}Πῶς δὲ παρακαλεῖτέ με κενά; τὸ δὲ ἐμὲ καταπαύσασθαι ἀφʼ ὑμῶν οὐδέν.

\ch{22}
Ὑπολαβὼν δὲ Ἐλιφὰζ ὁ Θαιμανίτης, λέγει,

\vs{2}Πότερον οὐχὶ ὁ Κύριός ἐστιν ὁ διδάσκων σύνεσιν καὶ ἐπιστήμην;
\vs{3}Τί γὰρ μέλει τῷ Κυρίῳ, ἐὰν σὺ ἦσθα τοῖς ἔργοις ἄμεμπτος; ἢ ὠφέλεια, ὅτι ἀπλώσῃς τὴν ὁδόν σου;
\vs{4}Ἢ λόγον σου ποιοῦμενος ἐλέγξεις, καὶ συνεισελεύσεταί σοι εἰς κρίσιν;

\vs{5}Πότερον οὐχ ἡ κακία σου ἐστὶ πολλὴ, ἀναρίθμητοι δέ σου εἰσὶν αἱ ἁμαρτίαι;
\vs{6}Ἠνεχύραζες δὲ τοὺς ἀδελφούς σου διακενῆς, ἀμφίασιν δὲ γυμνῶν ἀφείλου.
\vs{7}Οὐδὲ ὕδωρ διψῶντας ἐπότισας, ἀλλὰ πεινώντων ἐστέρησας ψωμόν·
\vs{8}Ἐθαύμασας δέ τινων πρόσωπον, ᾤκισας δὲ τοὺς ἐπὶ τῆς γῆς.
\vs{9}Χήρας δὲ ἐξαπέστειλας κενὰς, ὀρφανοὺς δὲ ἐκάκωσας.
\vs{10}Τοιγαροῦν ἐκύκλωσάν σε παγίδες, καὶ ἐσπούδασέ σε πόλεμος ἐξαίσιος.
\vs{11}Τὸ φῶς σοι σκότος ἀπέβη, κοιμηθέντα δὲ ὕδωρ σε ἐκάλυψε.

\vs{12}Μὴ οὐχὶ ὁ τὰ ὑψηλὰ ναίων ἐφορᾷ; τοὺς δὲ ὕβρει φερομένους ἐταπείνωσε;
\vs{13}Καὶ εἶπας, τί ἔγνω ὁ ἰσχυρός; ἢ κατὰ τοῦ γνόφου κρίνει;
\vs{14}Νεφέλη ἀποκρυφὴ αὐτοῦ καὶ οὐχ ὁραθήσεται, καὶ γῦρον οὐρανοῦ διαπορεύεται.
\vs{15}Μὴ τρίβον αἰώνιον φυλάξεις, ἣν ἐπάτησαν ἄνδρες δίκαιοι,
\vs{16}οἳ συνελήφθησαν ἄωροι; ποταμὸς ἐπιῤῥέων οἱ θεμέλιοι αὐτῶν,
\vs{17}οἱ λέγοντες, Κύριος τί ποιήσει ἡμῖν; ἢ τί ἐπάξεται ἡμῖν ὁ παντοκράτωρ;
\vs{18}Ὃς δὲ ἐνέπλησε τοὺς οἴκους αὐτῶν ἀγαθῶν, βουλὴ δὲ ἀσεβῶν πόῤῥω ἀπʼ αὐτοῦ.
\vs{19}Ἰδόντες δίκαιοι ἐγέλασαν, ἄμεμπτος δὲ ἐμυκτήρισεν.
\vs{20}Εἰ μὴ ἠφανίσθη ἡ ὑπόστασις αὐτῶν, καὶ τὸ κατάλειμμα αὐτῶν καταφάγεται πῦρ.

\vs{21}Γενοῦ δὴ σκληρὸς, ἐὰν ὑπομείνῃς, εἶτα ὁ καρπός σου ἔσται ἐν ἀγαθοῖς.
\vs{22}Ἔκλαβε δὲ ἐκ στόματος αὐτοῦ ἐξηγορίαν, καὶ ἀνάλαβε τὰ ῥήματα αὐτοῦ ἐν καρδίᾳ σου.
\vs{23}Ἐὰν δὲ ἐπιστραφῇς καὶ ταπεινώσῃς σεαυτὸν ἔναντι Κυρίου, πόῤῥω ἐποίησας ἀπὸ διαίτης σου ἄδικον.
\vs{24}Θήσῃ ἐπὶ χώματι ἐν πέτρᾳ, καὶ ὡς πέτρα χειμάῤῥου Σωφίρ.
\vs{25}Ἔσται οὖν σου ὁ παντοκράτωρ βοηθὸς ἀπὸ ἐχθρῶν, καθαρὸν δὲ ἀποδώσει σε ὥσπερ ἀργύριον πεπυρωμένον.
\vs{26}Εἶτα παῤῥησιασθήσῃ ἐναντίον Κυρίου ἀναβλέψας εἰς τὸν οὐρανὸν ἱλαρῶς.
\vs{27}Εὐξαμένου δέ σου πρὸς αὐτὸν εἰσακούσεταί σου· δώσει δέ σοι ἀποδοῦναι τὰς εὐχάς.
\vs{28}Ἀποκαταστήσει δέ σοι δίαιταν δικαιοσύνης, ἐπὶ δὲ ὁδοῖς σου ἔσται φέγγος·
\vs{29}Ὅτι ἐταπείνωσας σεαυτὸν, καὶ ἐρεῖς, ὑπερηφανεύσατο, καὶ κύφοντα ὀφθαλμοῖς σώσει.
\vs{30}Ῥύσεται ἀθῶον, καὶ διασώθητι ἐν καθαραῖς χερσί σου.

\ch{23}
Ὑπολαβὼν δὲ Ἰὼβ, λέγει,

\vs{2}Καὶ δὴ οἶδα ὅτι ἐκ χειρός μου ἡ ἔλεγξίς ἐστι, καὶ ἡ χεὶρ αὐτοῦ βαρεῖα γέγονεν ἐπʼ ἐμῷ στεναγμῷ.
\vs{3}Τίς δʼ ἄρα γνοίη, ὅτι εὕροιμι αὐτὸν, καὶ ἔλθοιμι εἰς τέλος;
\vs{4}Εἴποιμι δὲ ἐμαυτοῦ κρίμα, τὸ δὲ στόμα μου ἐμπλήσαι ἐλέγχων.
\vs{5}Γνοίην δὲ ἰάματα ἅ μοι ἐρεῖ, αἰσθοίμην δὲ τίνα μοι ἀπαγγελεῖ.
\vs{6}Καὶ ἐν πολλῇ ἰσχύϊ ἐπελεύσεταί μοι, εἶτα ἐν ἀπειλῇ μοι οὐ χρήσεται.
\vs{7}Ἀλήθεια γὰρ καὶ ἔλεγχος παρʼ αὐτοῦ, ἐξαγάγοι δὲ εἰς τέλος τὸ κρίμα μου.
\vs{8}Εἰ γὰρ πρῶτος πορεύσομαι, καὶ οὐκ ἔτι εἰμὶ, τὰ δὲ ἐπʼ ἐσχάτοις, τί οἶδα;

\vs{9}Ἀριστερὰ ποιήσαντος αὐτοῦ καὶ οὐ κατέσχον, περιβαλεῖ δεξιὰ καὶ οὐκ ὄψομαι·
\vs{10}Οἶδε γὰρ ἤδη ὁδόν μου· διέκρινε δέ με ὥσπερ τὸ χρυσίον.
\vs{11}Ἐξελεύσομαι δὲ ἐν ἐντάλμασιν αὐτοῦ, ὁδοὺς γὰρ αὐτοῦ ἐφύλαξα, καὶ οὐ μὴ ἐκκλίνω ἀπὸ ἐνταλμάτων αὐτοῦ,
\vs{12}καὶ οὐ μὴ παρέλθω, ἐν δὲ κόλπῳ μου ἔκρυψα ῥήματα αὐτοῦ.

\vs{13}Εἰ δὲ καὶ αὐτὸς ἔκρινεν οὕτως, τίς ἐστιν ὁ ἀντειπὼν αὐτῷ; ὁ γὰρ αὐτὸς ἠθέλησε, καὶ ἐποίησε.
\vs{15}Διὰ τοῦτο ἐπʼ αὐτῷ ἐσπούδακα· νουθετούμενος δὲ, ἐφρόντισα αὐτοῦ.
\vs{15a}Ἐπὶ τούτῳ ἀπὸ προσώπου αὐτοῦ κατασπουδασθῶ· κατανοήσω, καὶ πτοηθήσομαι ἐξ αὐτοῦ.

\vs{16}Κύριος δὲ ἐμαλάκυνε τὴν καρδίαν μου· ὁ δὲ παντοκράτωρ ἐσπούδασέ με.
\vs{17}Οὐ γὰρ ᾔδειν ὅτι ἐπελεύσεταί μοι σκότος, πρὸ προσώπου δέ μου ἐκάλυψε γνόφος.

\ch{24}
Διατί δὲ Κύριον ἔλαθον ὧραι,
\vs{2}ἀσεβεῖς δὲ ὅριον ὑπερέβησαν, ποίμνιον σὺν ποιμένι ἁρπάσαντες;
\vs{3}Ὑποζύγιον ὀρφανῶν ἀπήγαγον, καὶ βοῦν χήρας ἠνεχύρασαν.

\vs{4}Ἐξέκλιναν ἀδυνάτους ἐξ ὁδοῦ δικαίας, ὁμοθυμαδὸν δὲ ἐκρύβησαν πρᾳεῖς γῆς.
\vs{5}Ἀπέβησαν δὲ ὥσπερ ὄνοι ἐν ἀγρῷ, ὑπὲρ ἐμοῦ ἐξελθόντες τὴν ἑαυτῶν τάξιν· ἡδύνθη αὐτῷ ἄρτος εἰς νεωτέρους.

\vs{6}Ἀγρὸν πρὸ ὥρας οὐκ αὐτῶν ὄντα ἐθέρισαν· ἀδύνατοι ἀμπελῶνας ἀσεβῶν ἀμισθὶ καὶ ἀσιτὶ εἰργάσαντο.
\vs{7}Γυμνοὺς πολλοὺς ἐκοίμισαν ἄνευ ἱματίων, ἀμφίασιν δὲ ψυχῆς αὐτῶν ἀφείλαντο.
\vs{8}Ἀπὸ ψεκάδων ὀρέων ὑγραίνονται· παρὰ τὸ μὴ ἔχειν ἑαυτοὺς σκέπην, πέτραν περιεβάλοντο.

\vs{9}Ἥρπασαν ὀρφανὸν ἀπὸ μαστοῦ, ἐκπεπτωκότα δὲ ἐταπείνωσαν·
\vs{10}Γυμνοὺς δὲ ἐκοίμισαν ἀδίκως, πεινώντων δὲ τὸν ψωμὸν ἀφείλαντο.

\vs{11}Ἐν στενοῖς ἀδίκως ἐνήδρευσαν, ὁδὸν δὲ δικαίαν οὐκ ᾔδεισαν.
\vs{12}Οἳ ἐκ πόλεως καὶ οἴκων ἰδίων ἐξεβάλοντο, ψυχὴ δὲ νηπίων ἐστέναξε μέγα.

\vs{13}Αὐτὸς δὲ διατί τούτων ἐπισκοπὴν οὐ πεποίηται; ἐπὶ γῆς ὄντων αὐτῶν καὶ οὐκ ἐπέγνωσαν, ὁδὸν δὲ δικαιοσύνης οὐκ ᾔδεισαν, οὐδὲ ἀτραποὺς αὐτῶν ἐπορεύθησαν.
\vs{14}Γνοὺς δὲ αὐτῶν τὰ ἔργα, παρέδωκεν αὐτοὺς εἰς σκότος, καὶ νυκτὸς ἔσται ὡς κλέπτης.
\vs{15}Καὶ ὀφθαλμὸς μοιχοῦ ἐφύλαξε σκότος, λέγων, οὐ προνοήσει με ὀφθαλμὸς, καὶ ἀποκρυβὴν προσώπου ἔθετο.
\vs{16}Διώρυξεν ἐν σκότει οἰκίας, ἡμέρας ἐσφράγισαν ἑαυτοὺς, οὐκ ἐπέγνωσαν φῶς.
\vs{17}Ὅτι ὁμοθυμαδὸν αὐτοῖς τὸ πρωῒ σκιὰ θανάτου, ὅτι ἐπιγνώσεται τάραχος σκιᾶς θανάτου.
\vs{18}Ἐλαφρός ἐστιν ἐπὶ πρόσωπον ὕδατος, καταραθείη ἡ μερὶς αὐτῶν ἐπὶ γῆς,
\vs{19}ἀναφανείη δὲ τὰ φυτὰ αὐτῶν ἐπὶ γῆς ξηρά· ἀγκαλίδα γὰρ ὀρφανῶν ἥρπασαν.

\vs{20}Εἶτʼ ἀνεμνήσθη αὐτοῦ ἡ ἁμαρτία· ὥσπερ δὲ ὁμίχλη δρόσου ἀφανὴς ἐγένετο· ἀποδοθείη δὲ αὐτῷ ἃ ἔπραξε, συντριβείη δὲ πᾶς ἄδικος ἶσα ξύλῳ ἀνιάτῳ.

\vs{21}Στείραν δὲ οὐκ εὖ ἐποίησε, καὶ γύναιον οὐκ ἠλέησε.
\vs{22}Θυμῷ δὲ κατέστρεψεν ἀδυνάτους· ἀναστὰς τοιγαροῦν, οὐ μὴ πιστεύσῃ κατὰ τῆς ἑαυτοῦ ζωῆς.
\vs{23}Μαλακισθεὶς, μὴ ἐλπιζέτω ὑγιασθῆναι, ἀλλὰ πεσεῖται νόσῳ.
\vs{24}Πολλοὺς γὰρ ἐκάκωσε τὸ ὕψωμα αὐτοῦ, ἐμαράνθη δὲ ὥσπερ μολόχη ἐν καύματι, ἢ ὥσπερ στάχυς ἀπὸ καλάμης αὐτόματος ἀποπεσών.
\vs{25}Εἰ δὲ μὴ, τίς ἐστιν ὁ φάμενος ψευδῆ με λέγειν, καὶ θήσει εἰς οὐδὲν τὰ ῥήματά μου;

\ch{25}
Ὑπολαβὼν δὲ Βαλδὰδ ὁ Σαυχίτης, λέγει,

\vs{2}Τί γὰρ προοίμιον ἢ φόβος παρʼ αὐτοῦ ὁ ποιῶν τὴν σύμπασαν ἐν ὑψιστῳ;
\vs{3}Μὴ γάρ τις ὑπολάβοι ὅτι ἐστὶ παρέλκυσις πειραταῖς· ἐπὶ τίνας δὲ οὐκ ἐπελεύσεται ἔνεδρα παρʼ αὐτοῦ;
\vs{4}Πῶς γὰρ ἔσται δίκαιος βροτὸς ἔναντι Κυρίου; ἢ τίς ἂν ἀποκαθαρίσαι αὐτὸν γεννητὸς γυναικός;
\vs{5}Εἰ σελήνῃ συντάσσει, καὶ οὐκ ἐπιφαύσκει, ἄστρα δὲ οὐ καθαρὰ ἐναντίον αὐτοῦ.
\vs{6}Ἔα δὲ, ἄνθρωπος σαπρία, καὶ υἱὸς ἀνθρώπου σκώληξ.

\ch{26}
Ὑπολαβὼν δὲ Ἰὼβ, λέγει,

\vs{2}Τίνι πρόσκεισαι, ἢ τίνι μέλλεις βοηθεῖν; πότερον οὐχ ᾧ πολλὴ ἰσχὺς, καὶ ᾧ βραχίων κραταιός ἐστι;
\vs{3}Τίνι συμβεβούλευσαι; οὐχ ᾧ πᾶσα σοφία; τίνι ἐπακολουθήσεις; οὐχ ᾧ μεγίστη δύναμις;
\vs{4}Τίνι ἀνήγγειλας ῥήματα; πνοὴ δὲ τίνος ἐστὶν ἡ ἐξελθοῦσα ἐκ σοῦ;

\vs{5}Μὴ γίγαντες μαιωθήσονται ὑποκάτωθεν ὕδατος καὶ τῶν γειτόνων αὐτοῦ;
\vs{6}Γυμνὸς ὁ ᾅδης ἐνώπιον αὐτοῦ, καὶ οὐκ ἔστι περιβόλαιον τῇ ἀπωλείᾳ.
\vs{7}Ἐκτείνων βορέαν ἐπʼ οὐδὲν, κρεμάζων γῆν ἐπὶ οὐδενός.
\vs{8}Δεσμεύων ὕδωρ ἐν νεφέλαις αὐτοῦ, καὶ οὐκ ἐῤῥάγη νέφος ὑποκάτω αὐτοῦ·
\vs{9}Ὁ κρατῶν πρόσωπον θρόνου, ἐκπετάζων ἐπʼ αὐτὸν νέφος αὐτοῦ.
\vs{10}Πρόσταγμα ἐγύρωσεν ἐπὶ πρόσωπον ὕδατος, μέχρι συντελείας φωτὸς μετὰ σκότους.
\vs{11}Στύλοι οὐρανοῦ ἐπετάσθησαν, καὶ ἐξέστησαν ἀπὸ τῆς ἐπιτιμήσεως αὐτοῦ.
\vs{12}Ἰσχύϊ κατέπαυσε τὴν θάλασσαν, ἐπιστήμῃ δὲ ἔστρωται τὸ κῆτος.
\vs{13}Κλεῖθρα δὲ οὐρανοῦ δεδοίκασιν αὐτόν· προστάγματι δὲ ἐθανάτωσε δράκοντα ἀποστάτην.
\vs{14}Ἰδοῦ ταῦτα μέρη ὁδοῦ αὐτοῦ, καὶ ἐπὶ ἰκμάδα λόγου ἀκουσόμεθα ἐν αὐτῷ· σθενὸς δὲ βροντῆς αὐτοῦ, τίς οἶδεν ὁπότε ποιήσει;

\ch{27}
Ἔτι δὲ προσθεὶς Ἰὼβ εἶπεν τῷ προοιμίῳ,

\vs{2}Ζῇ ὁ Θεὸς, ὃς οὕτω με κέκρικε, καὶ ὁ παντοκράτωρ ὁ πικράνας μου τὴν ψυχὴν,
\vs{3}εἰ μὴν ἔτι τῆς πνοῆς μου ἐνούσης, πνεῦμα δὲ θεῖον τὸ περιόν μοι ἐν ῥινὶ,
\vs{4}μὴ λαλήσειν τὰ χείλη μου ὄνομα, οὐδὲ ἡ ψυχή μου μελετήσει ἄδικα.
\vs{5}Μή μοι εἴη δικαίους ὑμᾶς ἀποφῆναι ἕως ἂν ἀποθάνω, οὐ γὰρ ἀπαλλάξω μου τὴν ἀκακίαν μου.
\vs{6}Δικαιοσύνῃ δὲ προσέχων οὐ μὴν προῶμαι, οὐ γὰρ σύνοιδα ἐμαυτῷ ἄτοπα πράξας.
\vs{7}Οὐ μὴν δὲ ἀλλὰ εἴησαν οἱ ἐχθροί μου ὥσπερ ἡ καταστροφὴ τῶν ἀσεβῶν, καὶ οἱ ἐπʼ ἐμὲ ἐπανιστάμενοι, ὥσπερ ἡ ἀπώλεια τῶν παρανόμων.

\vs{8}Καὶ τίς γάρ ἐστιν ἐλπὶς ἀσεβεῖ, ὅτι ἐπέχει; πεποιθὼς ἐπὶ κύριον ἄρα σωθήσεται;
\vs{9}Ἢ τὴν δέησιν αὐτοῦ εἰσακούσεται ὁ Θεός; ἢ ἐπελθούσης αὐτῷ ἀνάγκης
\vs{10}μὴ ἔχει τινὰ παῤῥησίαν ἔναντι αὐτοῦ; ἢ ὡς ἐπικαλεσαμένου αὐτοῦ εἰσακούσεται αὐτοῦ;

\vs{11}Ἀλλὰ δὴ ἀναγγελῶ ὑμῖν τί ἐστιν ἐν χειρὶ Κυρίου, ἅ ἐστι παρὰ παντοκράτορι, οὐ ψεύσομαι.
\vs{12}Ἰδοὺ πάντες οἴδατε, ὅτι κενὰ κενοῖς ἐπιβάλλετε.
\vs{13}Αὕτη ἡ μερὶς ἀνθρώπου ἀσεβοῦς παρὰ Κυρίου, κτῆμα δὲ δυναστῶν ἐλεύσεται παρὰ παντοκράτορος ἐπʼ αὐτούς.
\vs{14}Ἐὰν δἐ πολλοὶ γένωνται οἱ υἱοὶ αὐτῶν, εἰς σφαγὴν ἔσονται· ἐὰν δὲ καὶ ἀνδρωθῶσι, προσαιτήσουσιν.
\vs{15}Οἱ δὲ περιόντες αὐτοῦ ἐν θανάτῳ τελευτήσουσι, χήρας δὲ αὐτῶν οὐδεὶς ἐλεήσει.
\vs{16}Ἐὰν συναγάγῃ ὥσπερ γῆν ἀργύριον, ἶσα δὲ πηλῷ ἑτοιμάσῃ χρυσίον,
\vs{17}ταῦτα πάντα δίκαιοι περιποιήσονται, τὰ δὲ χρήματα αὐτοῦ ἀληθινοὶ καθέξουσιν.
\vs{18}Ἀπέβη δὲ ὁ οἶκος αὐτοῦ ὥσπερ σῆτες, καὶ ὥσπερ ἀράχνη.
\vs{19}Πλούσιος κοιμηθήσεται καὶ οὐ προσθήσει, ὀφθαλμοὺς αὐτοῦ διήνοιξε καὶ οὐκ ἔστι.
\vs{20}Συνήντησαν αὐτῷ ὥσπερ ὕδωρ αἱ ὀδύναι, νυκτὶ δὲ ὑφείλατο αὐτὸν γνόφος.
\vs{21}Ἀναλήψεται δὲ αὐτὸν καύσων καὶ ἀπελεύσεται, καὶ λικμήσει αὐτὸν ἐκ τοῦ τόπου αὐτοῦ.
\vs{22}Καὶ ἐπιῤῥίψει ἐπʼ αὐτὸν καὶ οὐ φείσεται, ἐκ χειρὸς αὐτοῦ φυγῇ φεύξεται.
\vs{23}Κροτήσει ἐπʼ αὐτοὺς χεῖρας αὐτῶν, καὶ συριεῖ αὐτὸν ἐκ τοῦ τόπου αὐτοῦ.

\ch{28}
Ἐστι γὰρ ἀργυρίῳ τόπος ὅθεν γίνεται, τόπος δὲ χρυσίου ὅθεν διηθεῖται.
\vs{2}Σίδηρος μὲν γὰρ ἐκ γῆς γίνεται, χαλκὸς δὲ ἶσα λίθῳ λατομεῖται.

\vs{3}Τάξιν ἔθετο σκότει, καὶ πᾶν πέρας αὐτὸς ἐξακριβάζεται, λίθος σκοτία, καὶ σκιὰ θανάτου.
\vs{4}Διακοπὴ χειμάῤῥου ἀπὸ κονίας, οἱ δὲ ἐπιλανθανόμενοι ὁδὸν δικαίαν ἠσθένησαν, ἐκ βροτῶν ἐσαλεύθησαν.
\vs{5}Γῆ, ἐξ αὐτῆς ἐξελεύσεται ἄρτος, ὑποκάτω αὐτῆς ἐστράφη ὡσεὶ πῦρ.
\vs{6}Τόπος σαπφείρου οἱ λίθοι αὐτῆς, καὶ χῶμα χρυσίον αὐτῷ.
\vs{7}Τρίβος, οὐκ ἔγνω αὐτὴν πετεινὸν, καὶ οὐ παρέβλεψεν αὐτὴν ὀφθαλμὸς γυπός·
\vs{8}Καὶ οὐκ ἐπάτησαν αὐτὸν υἱοὶ ἀλαζόνων, οὐ παρῆλθεν ἐπʼ αὐτῆς λέων.
\vs{9}Ἐν ἀκροτόμῳ ἐξέτεινε χεῖρα αὐτοῦ, κατέστρεψε δὲ ἐκ ῥιζῶν ὄρη.
\vs{10}Δίνας δὲ ποταμῶν διέῤῥηξε, πᾶν δὲ ἔντιμον εἴδέ μου ὁ ὀφθαλμός.
\vs{11}Βάθη δὲ ποταμῶν ἀνεκάλυψεν, ἔδειξε δὲ αὐτοῦ δύναμιν εἰς φῶς.

\vs{12}Ἡ δὲ σοφία πόθεν εὑρέθη; ποῖος δὲ τόπος ἐστὶ τῆς ἐπιστήμης;
\vs{13}Οὐκ οἶδε βροτὸς ὁδὸν αὐτῆς, οὐδὲ μὴν εὑρέθη ἐν ἀνθρώποις.
\vs{14}Ἄβυσσος εἶπεν, οὐκ ἔνεστιν ἐν ἐμοί· καὶ ἡ θάλασσα εἶπεν, οὐκ ἔνεστι μετʼ ἐμοῦ.
\vs{15}Οὐ δώσει συνκλεισμὸν ἀντʼ αὐτῆς, καὶ οὐ σταθήσεται ἀργύριον ἀντάλλαγμα αὐτῆς.
\vs{16}Καὶ οὐ συνβασταχθήσεται χρυσίῳ Σωφεὶρ, ἐν ὄνυχι τιμίῳ καὶ σαπφείρῳ.
\vs{17}Οὐκ ἰσωθήσεται αὐτῇ χρυσίον καὶ ὕαλος, καὶ τὸ ἄλλαγμα αὐτῆς σκεύη χρυσᾶ.
\vs{18}Μετέωρα καὶ γαβὶς οὐ μνησθήσεται, καὶ ἕλκυσον σοφίαν ὑπὲρ τὰ ἐσώτατα.
\vs{19}Οὐκ ἰσωθήσεται αὐτῇ τοπάζιον Αἰθιοπίας, χρυσίῳ καθαρῷ οὐ συμβασταχθήσεται.

\vs{20}Ἡ δὲ σοφία πόθεν εὑρέθη; ποῖος δὲ τόπος ἐστὶ τῆς συνέσεως;
\vs{21}Λέληθε πάντα ἄνθρωπον, καὶ ἀπὸ πετεινῶν τοῦ οὐρανοῦ ἐκρύβη.
\vs{22}Ἡ ἀπώλεια καὶ ὁ θάνατος εἶπαν, ἀκηκόαμεν δὲ αὐτῆς τὸ κλέος.

\vs{23}Ὁ Θεὸς εὖ συνέστησεν αὐτῆς τὴν ὁδὸν, αὐτὸς δὲ οἶδε τὸν τόπον αὐτῆς.
\vs{24}Αὐτὸς γὰρ τὴν ὑπʼ οὐρανὸν πᾶσαν ἐφορᾷ· εἰδὼς τὰ ἐν τῇ γῇ, πάντα
\vs{25}ἃ ἐποίησεν, ἀνέμων σταθμὸν, ὕδατος μέτρα ὅτε ἐποίησεν·
\vs{26}οὕτως ἰδὼν ἠρίθμησε, καὶ ὁδὸν ἐν τινάγματι φωνάς.
\vs{27}Τότε εἶδεν αὐτὴν, καὶ ἐξηγήσατο αὐτὴν, ἑτοιμάσας ἐξιχνίασεν.
\vs{28}Εἶπε δὲ ἀνθρώπῳ, Ἰδοὺ ἡ θεοσέβειά ἐστι σοφία, τὸ δὲ ἀπέχεσθαι ἀπὸ κακῶν ἐστὶν ἐπιστήμη.

\ch{29}
Ἔτι δὲ προσθεὶς Ἰὼβ εἶπε τῷ προοιμίῳ,

\vs{2}Τίς ἄν με θείη κατὰ μῆνα ἔμπροσθεν ἡμερῶν, ὧν με ὁ Θεὸς ἐφύλαξεν;
\vs{3}Ὡς ὅτε ηὔγει ὁ λύχνος αὐτοῦ ὑπὲρ κεφαλῆς μου, ὅτε τῷ φωτὶ αὐτοῦ ἐπορευόμην ἐν σκότει·
\vs{4}Ὅτε ἤμην ἐπιβρίθων ὁδοὺς, ὅτε ὁ Θεὸς ἐπισκοπὴν ἐποιεῖτο τοῦ οἴκου μου·
\vs{5}Ὅτε ἤμην ὑλώδης λίαν, κύκλῳ δέ μου οἱ παῖδες·
\vs{6}Ὅτε ἐχέοντο αἱ ὁδοί μου βουτύρῳ, τὰ δὲ ὄρη μου ἐχέοντο γάλακτι·

\vs{7}Ὅτε ἐξεπορευόμην ὄρθριος ἐν πόλει, ἐν δὲ πλατείαις ἐτίθετό μου ὁ δίφρος.
\vs{8}Ἰδόντες με νεανίσκοι ἐκρύβησαν, πρεσβῦται δὲ πάντες ἔστησαν.
\vs{9}Ἁδροὶ δὲ ἐπαύσαντο λαλοῦντες, δάκτυλον ἐπιθέντες ἐπὶ στόματι.
\vs{10}Οἱ δὲ ἀκούσαντες ἐμακάρισάν με, καὶ γλῶσσα αὐτῶν τῷ λάρυγγι αὐτῶν ἐκολλήθη.
\vs{11}Ὅτι οὖς ἤκουσε καὶ ἐμακάρισέ με, ὀφθαλμὸς δὲ ἰδών με ἐξέκλινε.
\vs{12}Διέσωσα γὰρ πτωχὸν ἐκ χειρὸς δυνάστου, καὶ ὀρφανῷ ᾧ οὐκ ἦν βοηθὸς, ἐβοήθησα.
\vs{13}Εὐλογία ἀπολλυμένου ἐπʼ ἐμὲ ἔλθοι, στόμα δὲ χήρας με εὐλόγησε·
\vs{14}Δικαιοσύνην δὲ ἐνδεδύκειν, ἠμφιασάμην δὲ κρίμα ἶσα διπλοΐδι.
\vs{15}Ὀφθαλμὸς ἤμην τυφλῶν, ποὺς δὲ χωλῶν.
\vs{16}Ἐγὼ ἤμην πατὴρ ἀδυνάτων, δίκην δὲ ἣν οὐκ ᾔδειν ἐξιχνίασα.
\vs{17}Συνέτριψα δὲ μύλας ἀδίκων, ἐκ μέσου τῶν ὀδόντων αὐτῶν ἅρπαγμα ἐξήρπασα.
\vs{18}Εἶπα δὲ, ἡ ἡλικία μου γηράσει ὥσπερ στέλεχος φοίνικος, πολὺν χρόνον βιώσω.
\vs{19}Ἡ ῥίζα διήνοικται ἐπὶ ὕδατος, καὶ δρόσος αὐλισθήσεται ἐν τῷ θερισμῷ μου.
\vs{20}Ἡ δόξα μου κενὴ μετʼ ἐμοῦ, καὶ τὸ τόξον μου ἐν χειρὶ αὐτοῦ πορεύεται.

\vs{21}Ἐμοῦ ἀκούσαντες προσέσχον, ἐσιώπησαν δὲ ἐπὶ τῇ ἐμῇ βουλῇ.
\vs{22}Ἐπὶ τῷ ἐμῷ ῥήματι οὐ προσέθεντο, περιχαρεῖς δὲ ἐγίνοντο ὁπόταν αὐτοῖς ἐλάλουν.
\vs{23}Ὥσπερ γῆ διψῶσα προσδεχομένη τὸν ὑετὸν, οὕτως οὗτοι τὴν ἐμὴν λαλιάν.
\vs{24}Ἐὰν γελάσω πρὸς αὐτοὺς, οὐ μὴ πιστεύσωσι, καὶ φῶς τοῦ προσώπου μου οὐκ ἀπέπιπτεν.
\vs{25}Ἐξελεξάμην ὁδὸν αὐτῶν, καὶ ἐκάθισα ἄρχων, καὶ κατεσκήνουν ὡσεὶ βασιλεὺς ἐν μονοζώνοις, ὃν τρόπον παθεινοὺς παρακαλῶν.

\ch{30}
Νυνὶ δὲ κατεγέλασάν μου ἐλάχιστοι, νῦν νουθετοῦσί με ἐν μέρει, ὧν ἐξουδένουν τοὺς πατέρας αὐτῶν, οὓς οὐχ ἡγησάμην ἀξίους κυνῶν τῶν ἐμῶν νομάδων.
\vs{2}Καί γε ἰσχὺς χειρῶν αὐτῶν ἱνατί μοι; ἐπʼ αὐτοὺς ἀπώλετο συντέλεια.
\vs{3}Ἐν ἐνδείᾳ καὶ λιμῷ ἄγονος, οἱ φεύγοντες ἄνυδρον ἐχθὲς συνοχὴν καὶ ταλαιπωρίαν.
\vs{4}Οἱ περικυκλοῦντες ἄλιμα ἐπὶ ἠχοῦντι, οἵτινες ἄλιμα ἦν αὐτῶν τὰ σῖτα, ἄτιμοι δὲ καὶ πεφαυλισμένοι, ἐνδεεῖς παντὸς ἀγαθοῦ, οἳ καὶ ῥὶζας ξύλων ἐμασσῶντο ὑπὸ λιμοῦ μεγάλου.

\vs{5}Ἐπανέστησάν μοι κλέπται,
\vs{6}ὧν οἱ οἶκοι αὐτῶν ἦσαν τρῶγλαι πετρῶν·
\vs{7}Ἀναμέσον εὐήχων βοήσονται οἳ ὑπὸ φρύγανα ἄγρια διῃτῶντο.
\vs{8}Ἀφρόνων υἱοὶ καὶ ἀτίμων, ὄνομα καὶ κλέος ἐσβεσμένον ἀπὸ γῆς.
\vs{9}Νυνὶ δὲ κιθάρα ἐγώ εἰμι αὐτῶν, καὶ ἐμὲ θρύλλημα ἔχουσιν.
\vs{10}Ἐβδελύξαντο δέ με ἀποστάντες μακρὰν, ἀπὸ δὲ τοῦ προσώπου μου οὐκ ἐφείσαντο πτύελον.
\vs{11}Ἀνοίξας γὰρ φαρέτραν αὐτοῦ ἐκάκωσέ με, καὶ χαλινὸν τοῦ προσώπου μου ἐξαπέστειλεν.
\vs{12}Ἐπὶ δεξιῶν βλαστοῦ ἐπανέστησαν, πόδα αὐτῶν ἐξέτειναν, καὶ ὡδοποίησαν ἐπʼ ἐμὲ· τρίβους ἀπωλείας αὐτῶν.
\vs{13}Ἐξετρίβησαν τρίβοι μου, ἐξέδυσαν γάρ μου τὴν στολήν·
\vs{14}βέλεσιν αὐτοῦ κατηκόντισέ με. Κέκριται δέ μοι ὡς βούλεται, ἐν ὀδύναις πέφυρμαι.
\vs{15}Ἐπιστρέφονταί μου αἱ ὀδύναι, ᾤχετό μου ἡ ἐλπὶς ὥσπερ πνεῦμα, καὶ ὥσπερ νέφος ἡ σωτηρία μου.

\vs{16}Καὶ νῦν ἐπʼ ἐμὲ ἐκχυθήσεται ἡ ψυχή μου, ἔχουσιν δέ με ἡμέραι ὀδυνῶν.
\vs{17}Νυκτὶ δέ μου τὰ ὀστᾶ συγκέχυται, τὰ δὲ νεῦρά μου διαλέλυται.
\vs{18}Ἐν πολλῇ ἰσχύϊ ἐπελάβετό μου τῆς στολῆς, ὥσπερ τὸ περιστόμιον τοῦ χιτῶνός μου περιέσχε με.
\vs{19}Ἥγησαι δέ με ἴσα πηλῷ, ἐν γῇ καὶ σποδῷ μου ἡ μερίς.

\vs{20}Κέκραγα δὲ πρὸς σὲ καὶ οὐκ ἀκούεις μου, ἔστησαν δὲ καὶ κατενόησάν με.
\vs{21}Ἐπέβησαν δέ μοι ἀνελεημόνως, χειρὶ κραταιᾷ με ἐμαστίγωσας.
\vs{22}Ἔταξας δέ με ἐν ὀδύναις, καὶ ἀπέῤῥιψάς με ἀπὸ σωτηρίας.
\vs{23}Οἶδα γὰρ ὅτι θάνατός με ἐκτρίψει, οἰκία γὰρ παντὶ θνητῷ γῆ.
\vs{24}Εἰ γὰρ ὄφελον δυναίμην ἐμαυτὸν χειρώσασθαι, ἢ δεηθείς γε ἑτέρου, καὶ ποιήσει μοι τοῦτο.
\vs{25}Ἐγὼ δὲ ἐπὶ παντὶ ἀδυνάτῳ ἔκλαυσα, ἐστέναξα ἰδὼν ἄνδρα ἐν ἀνάγκαις·
\vs{26}ἐγὼ δὲ ἐπέχων ἀγαθοῖς, ἰδοὺ συνήντησάν μοι μᾶλλον ἡμέραι κακῶν.

\vs{27}Ἡ κοιλία μου ἐξέζεσε καὶ οὐ σιωπήσεται, προέφθασάν με ἡμέραι πτωχείας.
\vs{28}Στένων πεπόρευμαι ἄνευ φιμοῦ, ἕστηκα δὲ ἐν ἐκκλησίᾳ κεκραγώς.
\vs{29}Ἀδελφὸς γέγονα σειρήνων, ἑταῖρος δὲ στρουθῶν.
\vs{30}Τὸ δὲ δέρμα μου ἐσκότωται μεγάλως, τὰ δὲ ὀστᾶ μου ἀπὸ καύματος.
\vs{31}Ἀπέβη δὲ εἰς πένθος μου ἡ κιθάρα, ὁ δὲ ψαλμός μου εἰς κλαυθμὸν ἐμοί.

\ch{31}
Διαθήκην ἐθέμην τοῖς ὀφθαλμοῖς μου, καὶ οὐ συνήσω ἐπὶ παρθένον.
\vs{2}Καὶ τί ἐμέρισεν ὁ Θεὸς ἄνωθεν, καὶ κληρονομία ἱκανοῦ ἐξ ὑψίστων;
\vs{3}Οὐαί, ἀπώλεια τῷ ἀδίκῳ, καὶ ἀπαλλοτρίωσις τοῖς ποιοῦσιν ἀνομίαν.
\vs{4}Οὐχὶ αὐτὸς ὄψεται ὁδόν μου, καὶ πάντα τὰ διαβήματά μου ἐξαριθμήσεται;
\vs{5}Εἰ δὲ ἤμην πεπορευμένος μετὰ γελοιαστῶν, εἰ δὲ καὶ ἐσπούδασεν ὁ πούς μου εἰς δόλον·
\vs{6}Ἕσταμαι γὰρ ἐν ζυγῷ δικαίῳ, οἶδε δὲ ὁ Κύριος τὴν ἀκακίαν μου·
\vs{7}Εἰ ἐξέκλινεν ὁ πούς μου ἐκ τῆς ὁδοῦ, εἰ δὲ καὶ τῷ ὀφθαλμῷ ἐπηκολούθησεν ἡ καρδία μου, εἰ δὲ καὶ ταῖς χερσί μου ἡψάμην δώρων,
\vs{8}σπείραιμι ἄρα καὶ ἄλλοι φάγοισαν, ἄῤῥιζος δὲ γενοίμην ἐπὶ γῆς.
\vs{9}Εἰ ἐξηκολούθησεν ἡ καρδία μου γυναικὶ ἀνδρὸς ἑτέρου, εἰ καὶ ἐγκάθετος ἐγενόμην ἐπὶ θύραις αὐτῆς,
\vs{10}ἀρέσαι ἄρα καὶ ἡ γυνή μου ἑτέρῳ, τὰ δὲ νήπιά μου ταπεινωθείη.
\vs{11}Θυμὸς γὰρ ὀργῆς ἀκατάσχετος, τὸ μιᾶναι ἀνδρὸς γυναῖκα.
\vs{12}Πῦρ γάρ ἐστι καιόμενον ἐπὶ πάντων τῶν μερῶν, οὗδʼ ἂν ἐπέλθῃ ἐκ ῥιζῶν ἀπώλεσεν.

\vs{13}Εἰ δὲ καὶ ἐφαύλισα κρίμα θεράποντός μου ἢ θεραπαίνης, κρινομένων αὐτῶν πρὸς μὲ,
\vs{14}τί γὰρ ποιήσω ἐὰν ἔτασίν μου ποιῆται ὁ Κύριος; ἐὰν δὲ καὶ ἐπισκοπήν τινα, ἀπόκρισιν ποιήσομαι;
\vs{15}Πότερον οὐχ ὡς καὶ ἐγὼ ἐγενόμην ἐν γαστρὶ, καὶ ἐκεῖνοι γεγόνασι; γεγόναμεν δὲ ἐν τῇ αὐτῇ κοιλίᾳ.

\vs{16}Ἀδύνατοι δὲ χρείαν ἥν ποτε εἶχον οὐκ ἀπέτυχον, χήρας δὲ τὸν ὀφθαλμὸν οὐκ ἐξέτηξα·
\vs{17}Εἰ δὲ καὶ τὸν ψωμόν μου ἔφαγον μόνος, καὶ οὐχὶ ὀρφανῷ μετέδωκα·
\vs{18}Ὅτι ἐκ νεότητός μου ἐξέτρεφον ὡς πατὴρ, καὶ ἐκ γαστρὸς μητρός μου ὡδήγησα·
\vs{19}Εἰ δὲ καὶ ὑπερεῖδον γυμνὸν ἀπολλύμενον, καὶ οὐκ ἠμφίασα αὐτόν·
\vs{20}Ἀδύνατοι δὲ εἰ μὴ εὐλόγησάν με, ἀπὸ δὲ κουρᾶς ἀμνῶν μοῦ ἐθερμάνθησαν οἱ ὦμοι αὐτῶν·
\vs{21}Εἰ ἐπῆρα ὀρφανῷ χεῖρα, πεποιθὼς ὅτι πολλή μοι βοήθεια περίεστιν·
\vs{22}Ἀποσταίη ἄρα ὁ ὦμός μου ἀπὸ τῆς κλειδός, ὁ δὲ βραχίων μου ἀπὸ τοῦ ἀγκῶνος συντριβείη.
\vs{23}Φόβος γὰρ Κυρίου συνέσχε με, ἀπὸ τοῦ λήμματος αὐτοῦ οὐχ ὑποίσω.

\vs{24}Εἰ ἔταξα χρυσίον εἰς χοῦν μου, εἰ δὲ καὶ λίθῳ πολυτελεῖ ἐπεποίθησα,
\vs{25}εἰ δὲ καὶ εὐφράνθην πολλοῦ πλούτου μου γενομένου, εἰ δὲ καὶ ἐπʼ ἀναριθμήτοις ἐθέμην χεῖρά μου·
\vs{26}Ἢ οὐχ ὁρῶμεν ἥλιον τὸν ἐπιθαύσκοντα ἐκλείποντα, σελήνην δὲ φθίνουσαν; οὐ γὰρ ἐπʼ αὐτοῖς ἐστί·
\vs{27}Καὶ εἰ ἠπατήθη λάθρα ἡ καρδία μου, εἰ δὲ χεῖρά μου ἐπιθεὶς ἐπὶ στόματί μου ἐφίλησα.
\vs{28}Καὶ τοῦτό μοι ἄρα ἀνομία ἡ μεγίστη λογισθείη, ὅτι ἐψευσάμην ἐναντίον Κυρίου τοῦ ὑψίστου.
\vs{29}Εἰ δὲ καὶ ἐπιχαρὴς ἐγενόμην πτώματι ἐχθρῶν μου, καὶ εἶπεν ἡ καρδία μου, εὖγε.
\vs{30}Ἀκούσαι ἄρα τὸ οὖς μου τὴν κατάραν μου, θρυληθείην δὲ ἄρα ὑπὸ λαοῦ μου κακούμενος.

\vs{31}Εἰ δὲ καὶ πολλάκις εἶπον αἱ θεράπαιναί μου, τίς ἄν δῴη ἡμῖν τῶν σαρκῶν αὐτοῦ πλησθῆναι; λίαν μου χρηστοῦ ὄντος·
\vs{32}Ἔξω δὲ οὐκ ηὐλίζετο ξένος, ἡ δὲ θύρα μου παντὶ ἐλθόντι ἀνέῳκτο·
\vs{33}Εἰ δὲ καὶ ἁμαρτὼν ἀκουσίως ἔκρυψα τὴν ἁμαρτίαν μου·
\vs{34}Οὐ γὰρ διετράπην πολυοχλίαν πλήθους, τοῦ μὴ ἐξαγορεῦσαι ἐνώπιον αὐτῶν· εἰ δὲ καὶ εἴασα ἀδύνατον ἐξελθεῖν θύραν μου κόλπῳ κενῷ·
\vs{35}Τίς δῴη ἀκούοντά μου; χεῖρα δὲ Κυρίου εἰ μὴ ἐδεδοίκειν· συγγραφὴν δὲ ἣν εἶχον κατά τινος,
\vs{36}ἐπʼ ὤμοις ἂν περιθέμενος στέφανον ἀνεγίνωσκον,
\vs{37}καὶ εἰ μὴ ῥῆξας αὐτὴν ἀπέδωκα, οὐθὲν λαβὼν παρὰ χρεωφειλέτου·

\vs{38}Εἰ ἐπʼ ἐμοί ποτε ἡ γῆ ἐστέναξεν, εἰ δὲ καὶ οἱ αὔλακες αὐτῆς ἔκλαυσαν ὁμοθυμαδόν·
\vs{39}Εἰ δὲ καὶ τὴν ἰσχὺν αὐτῆς ἔφαγον μόνος ἄνευ τιμῆς, εἰ δὲ καὶ ψυχὴν κυρίου τῆς γῆς ἐκλαβὼν ἐλύπησα·
\vs{40}Ἀντὶ πυροῦ ἄρα ἐξέλθοι μοι κνίδη, ἀντὶ δὲ κριθῆς βάτος. Καὶ ἐπαύσατο Ἰὼβ ῥήμασιν.

\ch{32}
Ἡσύχασαν δὲ καὶ οἱ τρεῖς φίλοι αὐτοῦ ἔτι ἀντειπεῖν Ἰὼβ, ἦν γὰρ Ἰὼβ δίκαιος ἐναντίον αὐτῶν.

\vs{2}Ὠργίσθη δὲ Ἐλιοὺς ὁ τοῦ Βαραχιὴλ ὁ Βουζίτης ἐκ τῆς συγγενείας Ῥὰμ, τῆς Αὐσίτιδος χώρας· ὠργίσθη δὲ τῷ Ἰὼβ σφόδρα, διότι ἀπέφῃνεν ἑαυτὸν δίκαιον ἐναντίον Κυρίου·
\vs{3}Καὶ κατὰ τῶν τριῶν δὲ φίλων ὠργίσθη σφόδρα, διότι οὐκ ἠδυνήθησαν ἀποκριθῆναι ἀντίθετα Ἰὼβ, καὶ ἔθεντο αὐτὸν εἶναι ἀσεβῆ.
\vs{4}Ἐλιοὺς δὲ ὑπέμεινε δοῦναι ἀπόκρισιν Ἰὼβ, ὅτι πρεσβύτεροι αὐτοῦ εἰσιν ἡμέραις.
\vs{5}Καὶ εἶδεν Ἐλιοὺς, ὅτι οὐκ ἔστιν ἀπόκρισις ἐν στόματι τῶν τριῶν ἀνδρῶν, καὶ ἐθυμώθη ὀργῇ αὐτοῦ.
\vs{6}Ὑπολαβὼν δὲ Ἐλιοὺς ὁ τοῦ Βαραχιὴλ ὁ Βουζίτης, εἶπε,

Νεώτερος μέν εἰμι τῷ χρόνῳ, ὑμεῖς δέ ἐστε πρεσβύτεροι, διὸ ἡσύχασα φοβηθεὶς τοῦ ὑμῖν ἀναγγεῖλαι τὴν ἐμαυτοῦ ἐπιστήμην.
\vs{7}Εἶπα δὲ, ὅτι οὐχ ὁ χρόνος ἐστὶν ὁ λαλῶν, ἐν πολλοῖς δὲ ἔτεσιν οἴδασι σοφίαν.
\vs{8}Ἀλλὰ πνεῦμά ἐστιν ἐν βροτοῖς· πνοὴ δὲ παντοκράτορός ἐστιν ἡ διδάσκουσα.
\vs{9}Οὐχ οἱ πολυχρόνιοί εἰσι σοφοὶ, οὐδʼ οἱ γέροντες οἴδασι κρίμα.
\vs{10}Διὸ εἶπα, ἀκούσατέ μου, καὶ ἀναγγελῶ ὑμῖν ἃ οἶδα.

\vs{11}Ἐνωτίζεσθέ μου τὰ ῥήματα· ἐρῶ γὰρ ὑμῶν ἀκουόντων ἄχρις οὗ ἐτάσητε λόγουις,
\vs{12}καὶ μέχρι ὑμῶν συνήσω, καὶ ἰδοὺ οὐκ ἦν τῷ Ἰὼβ ἐλέγχων ἀνταποκρινόμενος ῥήματα αὐτοῦ ἐξ ὑμῶν·
\vs{13}ἵνα μὴ εἴπητε, εὕρομεν σοφίαν Κυρίῳ προσθέμενοι.
\vs{14}Ἀνθρώπῳ δὲ ἐπετρέψατε λαλῆσαι τοιαῦτα ῥήματα.
\vs{15}Ἐπτοήθησαν, οὐκ ἀπεκρίθησαν ἔτι, ἐπαλαίωσαν ἐξ αὐτῶν λόγους.
\vs{16}Ὑπέμεινα, οὐ γὰρ ἐλάλησα, ὅτι ἔστησαν, οὐκ ἀπεκρίθησαν.
\vs{17}Ὑπολαβὼν δὲ Ἐλιοῦς, λέγει,
\vs{18}πάλιν λαλήσω, πλήρης γάρ εἰμι ῥημάτων, ὀλέκει γάρ με τὸ πνεῦμα τῆς γαστρός.
\vs{19}Ἡ δὲ γαστήρ μου ὥσπερ ἀσκὸς γλεύκους ζέων δεδεμένος, ἢ ὥσπερ φυσητὴρ χαλκέως ἐῤῥηγώς.
\vs{20}Λαλήσω, ἵνα ἀναπαύσωμαι ἀνοίξας τὰ χείλη,
\vs{21}ἄνθρωπον γὰρ οὐ μὴ αἰσχυνθῶ, ἀλλὰ μὴν οὐδὲ βροτὸν οὐ μὴ ἐντραπῶ.
\vs{22}Οὐ γὰρ ἐπίσταμαι θαυμάσαι πρόσωπα· εἰ δὲ μὴ, καὶ ἐμὲ σῆτες ἔδονται.

\ch{33}
Οὐ μὴν δὲ ἀλλὰ ἄκουσου Ἰὼβ τὰ ῥήματά μου, καὶ λαλιὰν ἐνωτίζου μου.
\vs{2}γὰρ ἤνοιξα τὸ στόμα μου, καὶ ἐλάλησεν ἡ γλῶσσά μου.
\vs{3}Καθαρά μου ἡ καρδία ῥήμασι, σύνεσις δὲ χειλέων μου καθαρὰ νοήσει.
\vs{4}Πνεῦμα θεῖον τὸ ποιῆσάν με, πνοὴ δὲ παντοκράτορος ἡ διδάσκουσά με.
\vs{5}Ἐὰν δύνῃ, δός μοι ἀπόκρισιν, πρὸς ταῦτα ὑπόμεινον, στῆθι κατʼ ἐμὲ, καὶ ἐγὼ κατὰ σέ.
\vs{6}Ἐκ πηλοῦ διήρτισαι σὺ ὡς καὶ ἐγὼ, ἐκ τοῦ αὐτοῦ διηρτίσμεθα.
\vs{7}Οὐχ ὁ φόβος μου σὲ στροβήσει, οὐδὲ ἡ χείρ μου βαρεῖα ἔσται ἐπὶ σοί.

\vs{8}Πλὴν εἶπας ἐν ὠσί μου· φωνὴν ῥημάτων σου ἀκήκοα·
\vs{9}διότι λέγεις, καθαρός εἰμι οὐχ ἁμαρτών, ἄμεμπτός εἰμι, οὐ γὰρ ἠνόμησα·
\vs{10}Μέμψιν δὲ κατʼ ἐμοῦ εὗρεν· ἥγηται δέ με ὥσπερ ὑπεναντίον.
\vs{11}Ἔθετο δὲ ἐν ξύλῳ τὸν πόδα μου, ἐφύλαξε δέ μου πάσας τὰς ὁδούς.
\vs{12}Πῶς γὰρ λέγεις, δίκαιός εἰμι, καὶ οὐκ ἐπακήκοέ μου; αἰώνιος γάρ ἐστιν ὁ ἐπάνω βροτῶν.

\vs{13}Λέγεις δέ, διατί τῆς δίκης μου οὐκ ἐπακήκοέ μου πᾶν ῥῆμα;
\vs{14}Ἐν γὰρ τῷ ἅπαξ λαλήσαι ὁ Κύριος, ἐν δὲ τῷ δευτέρῳ.
\vs{15}ἐνύπνιον ἢ ἐν μελέτῃ νυκτερινῇ, ὡς ὅταν ἐπιπίπτῃ δεινὸς φόβος ἐπʼ ἀνθρώπους, ἐπὶ νυσταγμάτων ἐπὶ κοίτης·
\vs{16}Τότε ἀνακαλύπτει νοῦν ἀνθρώπων, ἐν εἴδεσι φόβου τοιούτοις αὐτοὺς ἐξεφόβησεν.
\vs{17}Ἀποστρέψαι ἄνθρωπον ἀπὸ ἀδικίας, τὸ δὲ σῶμα αὐτοῦ ἀπὸ πτώματος ἐῤῥύσατο.
\vs{18}Ἐφείσατο δὲ τῆς ψυχῆς αὐτοῦ ἀπὸ θανάτου, καὶ μὴ πεσεῖν αὐτὸν ἐν πολέμῳ.

\vs{19}Πάλιν δὲ ἤλεγξεν αὐτὸν ἐπὶ μαλακίᾳ ἐπὶ κοίτης, καὶ πλῆθος ὀστῶν αὐτοῦ ἐνάρκησε.
\vs{20}Πᾶν δὲ βρωτὸν σίτου οὐ μὴ δύνηται προσδέξασθαι, καὶ ἡ ψυχὴ αὐτοῦ βρῶσιν ἐπιθυμήσει·
\vs{21}ἕως ἂν σαπῶσιν αὐτοῦ αἱ σάρκες, καὶ ἀποδείξῃ τὰ ὀστᾶ αὐτοῦ κένα.
\vs{22}Ἤγγισε δὲ εἰς θάνατον ἡ ψυχὴ αὐτοῦ, ἡ δὲ ζωὴ αὐτοῦ ἐν ᾅδῃ·
\vs{23}Ἐὰν ὦσι χιλιοι αγγελοι θανατηφόροι, εἷς αὐτῶν οὐ μὴ τρώσῃ αὐτόν· ἐὰν νοήσῃ τῇ καρδίᾳ ἐπιστραφῆναι πρὸς Κύριον, ἀναγγείλῃ δὲ ἀνθρώπῳ τὴν ἑαυτοῦ μέμψιν, τὴν δὲ ἄνοιαν αὐτοῦ δείξῃ,
\vs{24}ἀνθέξεται τοῦ μὴ πεσεῖν εἰς θάνατον· ἀνανεώσει δὲ αὐτοῦ τὸ σῶμα ὥσπερ ἀλοιφὴν ἐπὶ τοίχου, τὰ δὲ ὀστᾶ αὐτοῦ ἐμπλήσει μυελοῦ·
\vs{25}Ἁπαλυνεῖ δὲ αὐτοῦ τὰς σάρκας ὥσπερ νηπίου, ἀποκαταστήσει δὲ αὐτὸν ἀνδρωθέντα ἐν ἀνθρώποις.
\vs{26}Εὐξάμενος δὲ πρὸς Κύριον καὶ δεκτὰ αὐτῷ ἔσται, εἰσελεύσεται προσώπῳ ἱλαρῷ σὺν ἐξηγορίᾳ· ἀποδώσει δὲ ἀνθρώποις δικαιοσύνην.
\vs{27}Εἶτα τότε ἀπομέμψεται ἄνθρωπος αὐτὸς ἑαυτῷ, λέγων, οἷα συνετέλουν; καὶ οὐκ ἄξια ἤτασέ με ὧν ἥμαρτον.
\vs{28}Σῶσον ψυχήν μου τοῦ μὴ ἐλθεῖν εἰς διαφθορὰν, καὶ ἡ ζωή μου φῶς ὄψεται.

\vs{29}Ἰδοὺ ταῦτα πάντα ἐργᾶται ὁ ἰσχυρὸς ὁδοὺς τρεῖς μετὰ ἀνδρός.
\vs{30}Καὶ ἐῤῥύσατο τὴν ψυχήν μου ἐκ θανάτου, ἵνα ἡ ζωή μου ἐν φωτὶ αἰνῇ αὐτόν.

\vs{31}Ἐνωτίζου Ἰὼβ καὶ ἄκουέ μου· κώφευσον, καὶ ἐγώ εἰμι λαλήσω.
\vs{32}Εἰ εἰσί σοι λόγοι, ἀποκρίθητί μοι· λάλησον, θέλω γὰρ δικαιωθῆναί σε.
\vs{33}Εἰ μὴ, σὺ ἄκουσόν μου, κώφευσον καὶ διδάξω σε.

\ch{34}
Ὑπολαβὼν δὲ Ἐλιοὺς, λέγει,

\vs{2}Ἀκούσατέ μου σοφοὶ, ἐπιστάμενοι ἐνωτίζεσθε.
\vs{3}Ὅτι οὖς λόγους δοκιμάζει, καὶ λάρυγξ γεύεται βρῶσιν.
\vs{4}Κρίσιν ἑλώμεθα ἑαυτοῖς, γνῶμεν ἀναμέσον ἑαυτῶν, ὅ, τι καλόν.
\vs{5}Ὅτι εἴρηκεν Ἰὼβ, δίκαιός εἰμι· ὁ Κύριος ἀπήλλαξέ μου τὸ κρίμα.
\vs{6}Ἐψεύσατο δὲ τῷ κρίματί μου, βίαιον τὸ βέλος μου ἄνευ ἀδικίας.

\vs{7}Τίς ἀνὴρ ὥσπερ Ἰὼβ, πίνων μυκτηρισμὸν ὥσπερ ὕδωρ;
\vs{8}Οὐχ ἁμαρτῶν οὐδὲ ἀσεβήσας, ἢ οὐδʼ οὐ κοινωνήσας μετὰ ποιούντων τὰ ἄνομα, τοῦ πορευθῆναι μετὰ ἀσεβῶν.
\vs{9}Μὴ γὰρ εἴπῃς, ὅτι οὐκ ἔσται ἐπισκοπὴ ἀνδρός, καὶ ἐπισκοπὴ αὐτῷ παρὰ Κυρίου.

\vs{10}Διὸ συνετοὶ καρδίας ἀκούσατέ μου, μή μοι εἴη ἔναντι Κυρίου ἀσεβῆσαι, καὶ ἔναντι παντοκράτορος ταράξαι τὸ δίκαιον.
\vs{11}Ἀλλὰ ἀποδιδοῖ ἀνθρώπῳ καθὰ ποιεῖ ἕκαστος αὐτῶν, καὶ ἐν τρίβῳ ἀνδρὸς εὑρήσει αὐτόν.

\vs{12}Οἴει δὲ τὸν Κύριον ἄτοπα ποιήσειν, ἢ ὁ παντοκράτωρ ταράξει κρίσιν,
\vs{13}ὃς ἐποίησε τὴν γῆν; Τίς δέ ἐστιν ὁ ποιῶν τὴν ὑπʼ οὐρανὸν, καὶ τὰ ἐνόντα πάντα;
\vs{14}Εἰ γὰρ βούλοιτο συνέχειν, καὶ τὸ σνεῦμα παρʼ αὐτῷ κατασχεῖν,
\vs{15}τελευτήσει πᾶσα σὰρξ ὁμοθυμαδὸν, πᾶς δὲ βροτὸς εἰς γῆν ἀπελεύσεται, ὅθεν καὶ ἐπλάσθη.

\vs{16}Ἴδὲ μὴ νουθετῇ, ἄκουε ταῦτα, ἐνωτίζου φωνὴν ῥημάτων.
\vs{17}Ἴδε σὺ τὸν μισοῦντα ἄνομα, καὶ τὸν ὀλλύντα τοὺς πονηροὺς, ὄντα αἰώνιον δίκαιον.

\vs{18}Ἀσεβὴς ὁ λέγων βασιλεῖ, παρανομεῖς, ἀσεβέστατε τοῖς ἄρχουσιν.
\vs{19}Ὃς οὐκ ἐπαισχυνθῇ πρόσωπον ἐντίμου, οὐδὲ οἶδε τιμὴν θέσθαι ἁδροῖς θαυμασθῆναι πρόσωπα αὐτῶν.
\vs{20}Κενὰ δὲ αὐτοῖς ἀποβήσεται, τὸ κεκραγέναι καὶ δεῖσθαι ἀνδρός· ἐχρήσαντο γὰρ παρανόμως, ἐκκλινομένων ἀδυνάτων.
\vs{21}Αὐτὸς γὰρ ὁρατής ἐστιν ἔργων ἀνθρώπων, λέληθε δὲ αὐτὸν οὐδὲν ὧν πράσσουσιν.
\vs{22}Οὐδὲ ἔσται τόπος τοῦ κρυβῆναι τοὺς ποιοῦντας τὰ ἄνομα·
\vs{23}Ὅτι οὐκ ἐπʼ ἄνδρα θήσει ἔτι.
\vs{24}Ὁ γὰρ Κύριος πάντας ἐφορᾷ· ὁ καταλαμβάνων ἀνεξιχνίαστα, ἔνδοξά τε καὶ ἐξαίσια, ὧν οὐκ ἔστιν ἀριθμός.
\vs{25}Ὁ γνωρίζων αὐτῶν τὰ ἔργα, καὶ στρέψει νύκτα καὶ ταπεινωθήσοται.
\vs{26}Ἔσβεσε δὲ ἀσεβεῖς, ὁρατοὶ δὲ ἐναντίον αὐτοῦ.
\vs{27}Ὅτι ἐξέκλιναν ἐκ νόμου Θεοῦ, δικαιώματα δὲ αὐτοῦ οὐκ ἐπέγνωσαν,
\vs{28}τοῦ ἐπαγαγεῖν ἐπʼ αὐτὸν κραυγὴν πενήτων, καὶ κραυγὴν πτωχῶν εἰσακούσεται.

\vs{29}Καὶ αὐτὸς ἡσυχίαν παρέξει, καὶ τίς καταδικάσεται; καὶ κρύψει πρόσωπον, καὶ τίς ὄψεται αὐτόν; καὶ κατὰ ἔθνους, καὶ κατὰ ἀνθρώπου ὁμοῦ.
\vs{30}Βασιλεύων ἀνθρωπον ὑποκριτὴν ἀπὸ δυσκολίας λαοῦ.

\vs{31}Ὅτι πρὸς τὸν ἰσχυρὸν ὁ λέγων, εἴληφα, οὐκ ἐνεχυράσω·
\vs{32}Ἄνευ ἐμαυτοῦ ὄψομαι· σὺ δεῖξόν μοι, εἰ ἀδικίαν εἰργασάμην, οὐ μὴ προσθήσω.
\vs{33}Μὴ ἀπὸ σοῦ ἀποτίσει αὐτὴν, ὅτι σὺ ἀπώσῃ; ὅτι σὺ ἐκλέξῃ καὶ οὐκ ἐγώ; καὶ τί ἔγνως, λάλησον.
\vs{34}Διὸ συνετοὶ καρδίας ἐροῦσι ταῦτα, ἀνὴρ δὲ σοφὸς ἀκήκοέ μου τὸ ῥῆμα.
\vs{35}Ἰὼβ δὲ οὐκ ἐν συνέσει ἐλάλησε, τὰ ῥήματα αὐτοῦ οὐκ ἐν ἐπιστήμῃ.
\vs{36}Οὐ μὴν δὲ ἀλλὰ μάθε Ἰώβ, μὴ δῷς ἔτι ἀνταπόκρισιν ὥσπερ οἱ ἄφρονες·
\vs{37}Ἵνα μὴ προσθῶμεν ἐφʼ ἁμαρτίας ἡμῶν, ἀνομία δὲ ἐφʼ ἡμῖν λογισθήσεται, πολλὰ λαλούντων ῥήματα ἐναντίον τοῦ Κυρίου.

\ch{35}
Ὑπολαβὼν δὲ Ἐλιοὺς, λέγει,

\vs{2}Τί τοῦτο ἡγήσω ἐν κρίσει; σὺ τίς εἶ, ὅτι εἶπας, δίκαιός εἰμι ἔναντι Κυρίου;
\vs{4}Ἐγώ σοι δώσω ἀπόκρισιν, καὶ τοῖς τρισὶ φίλοις σου.

\vs{5}Ἀνάβλεψον εἰς τὸν οὐρανὸν, καὶ ἴδε· κατάμαθε δὲ νέφη, ὡς ὑψηλὰ ἀπὸ σοῦ.
\vs{6}Εἰ ἥμαρτες, τί πράξεις; εἰ δὲ καὶ πολλὰ ἠνόμησας, τί δύνασαι ποιῆσαι;
\vs{7}Ἐπεὶ δὲ οὖν δίκαιος εἶ, τί δώσεις αὐτῷ; ἢ τί ἐκ χειρός σου λήψεται;
\vs{8}Ἀνδρὶ τῷ ὁμοίῳ σου ἡ ἀσέβειά σου, καὶ υἱῷ ἀνθρώπου ἡ δικαιοσύνη σου.
\vs{9}Ἀπὸ πλήθους συκοφαντούμενοι κεκράξονται, βοήσονται ἀπὸ βραχίονος πολλῶν.
\vs{10}Καὶ οὐκ εἶπε, ποῦ ἐστιν ὁ Θεὸς ὁ ποιήσας με, ὁ κατατάσσων φυλακὰς νυκτερινὰς,
\vs{11}ὁ διορίζων με ἀπὸ τετραπόδων γῆς, ἀπὸ δὲ πετεινῶν οὐρανοῦ;
\vs{12}Ἐκεῖ κεκράξονται, καὶ οὐ μὴ εἰσακούσῃ, καὶ ἀπὸ ὕβρεως πονηρῶν.

\vs{13}Ἄτοπα γὰρ οὐ βούλεται ἰδεῖν ὁ Κύριος, αὐτὸς γὰρ ὁ παντοκράτωρ.
\vs{14}Ὁρατής ἐστι τῶν συντελούντων τὰ ἄνομα, καὶ σώσει με· κρίθητι δὲ ἐναντίον αὐτοῦ, εἰ δύνασαι αὐτὸν αἰνέσαι, ὡς ἔστι καὶ νῦν.
\vs{15}Ὅτι οὐκ ἔστιν ἐπισκεπτόμενος ὀργὴν αὐτοῦ, καὶ οὐκ ἔγνω παράπτωμά τι σφόδρα·
\vs{16}καὶ Ἰὼβ ματαίως ἀνοίγει τὸ στόμα αὐτοῦ, ἐν ἀγνωσίᾳ ῥήματα βαρύνει.

\ch{36}
Προσθεὶς δὲ ἔτι Ἐλιοὺς, λέγει,

\vs{2}Μεῖνόν με μικρὸν ἔτι, ἵνα διδάξω σε· ἔτι γὰρ ἐν ἐμοί ἐστι λέξις.
\vs{3}Ἀναλαβὼν τὴν ἐπιστήμην μου μακράν, ἔργοις δέ μου δίκαια ἐρῶ ἐπʼ ἀληθείας,
\vs{4}καὶ οὐκ ἄδικα ῥήματα ἀδίκως συνιεῖς.

\vs{5}Γίνωσκε δὲ, ὅτι ὁ Κύριος οὐ μὴ ἀποποιήσηται τὸν ἄκακον, δυνατὸς ἰσχύϊ καρδίας
\vs{6}ἀσεβῆ οὐ μὴ ζωοποιήσῃ, καὶ κρίμα πτωχῶν δώσει.
\vs{7}Οὐκ ἀφελεῖ ἀπὸ δικαίου ὀφθαλμοὺς αὐτοῦ, καὶ μετὰ βασιλέων εἰς θρόνον, καὶ καθιεῖ αὐτοὺς εἰς νῖκος, καὶ ὑψωθήσονται.
\vs{8}Καὶ οἱ πεπεδημένοι ἐν χειροπέδαις, συσχεθήσονται ἐν σχοινίοις πενίας.
\vs{9}Καὶ ἀναγγελεῖ αὐτοῖς τὰ ἔργα αὐτῶν, καὶ τὰ παραπτώματα αὐτῶν, ὅτι ἰσχύσουσιν.
\vs{10}Ἀλλὰ τοῦ δικαίου εἰσακούσεται· καὶ εἶπεν ὅτι ἐπιστραφήσονται ἐξ ἀδικίας.
\vs{11}Ἐὰν ἀκούσωσι, καὶ δουλεύσωσι, συντελέσουσι τὰς ἡμέρας αὐτῶν ἐν ἀγαθοῖς, καὶ τὰ ἔτη αὐτῶν ἐν εὐπρεπείαις.
\vs{12}Ἀσεβεῖς δὲ οὐ διασώζει, παρὰ τὸ μὴ βούλεσθαι αὐτοὺς εἰδέναι τὸν Κύριον, καὶ διότι νουθετούμενοι ἀνήκοοι ἦσαν.

\vs{13}Καὶ ὑποκριταὶ καρδίᾳ τάξουσι θυμόν, οὐ βοήσονται, ὅτι ἔδησεν αὐτούς.
\vs{14}Ἀποθάνοι τοίνυν ἐν νεότητι ἡ ψυχὴ αὐτῶν, ἡ δὲ ζωὴ αὐτῶν τιτρωσκομένη ὑπὸ ἀγγέλων,
\vs{15}ἀνθʼ ὧν ἔθλιψαν ἀσθενῆ καὶ ἀδύνατον, κρίμα δὲ πραέων ἐκθήσει.
\vs{16}Καὶ προσεπιηπάτησέ σε ἐκ στόματος ἐχθροῦ, ἄβυσσος κατάχυσις ὑποκάτω αὐτῆς, καὶ κατέβη τράπεζά σου πλήρης πιότητος.
\vs{17}Οὐχ ὑστερήσει δὲ ἀπὸ δικαίων κρίμα,
\vs{18}θυμὸς δὲ ἐπʼ ἀσεβεῖς ἔσται, διʼ ἀσέβειαν δώρων ὧν ἐδέχοντο ἐπʼ ἀδικίαις.

\vs{19}Μή σε ἐκκλινάτω ἑκὼν ὁ νοῦς δεήσεως ἐν ἀνάγκῃ ὄντων ἀδυνάτων,
\vs{20}καὶ πάντας τοὺς κραταιοῦντας ἰσχὺν μὴ ἐξελκύσῃς τὴν νύκτα, τοῦ ἀναβῆναι λαοὺς ἀντʼ αὐτῶν,
\vs{21}ἀλλὰ φύλαξαι μὴ πράξῃς ἄτοκα· ἐπὶ τοὺτων γὰρ ἐξείλω ἀπὸ πτωχείας.

\vs{22}Ἰδοὺ ὁ ἰσχυρὸς κραταιώσει ἐν ἰσχύι αὐτοῦ· τίς γάρ ἐστι κατʼ αὐτὸν δυνάστης;
\vs{23}Τίς δέ ἐστιν ὁ ἐτάζων αὐτοῦ τὰ ἔργα; ἢ τίς ὁ εἰπὼν, ἔπραξεν ἄδικα;
\vs{24}Μνήσθητι, ὅτι μεγάλα ἐστὶν αὐτοῦ τὰ ἔργα, ὧν ἦρξαν ἄνδρες.
\vs{25}Πᾶς ἄνθρωπος εἶδεν ἐν ἑαυτῷ, ὅσοι τιτρωσκόμενοί εἰσι βροτοί.
\vs{26}Ἰδοὺ ὁ ἰσχυρὸς πολύς, καὶ οὐ γνωσόμεθα· ἀριθμὸς ἐτῶν αὐτοῦ καὶ ἀπέραντος.
\vs{27}Ἀριθμηταὶ δὲ αὐτῷ σταγόνες ὑετοῦ, καὶ ἐπιχυθήσονται ὑετῷ εἰς νεφέλην.
\vs{28}Ῥυήσονται παλαιώματα, ἐσκίασε δὲ νέφη ἐπὶ ἀμυθήτῳ βροτῷ·
\vs{28a}ὥραν ἔθετο κτήνεσιν, οἴδασι δὲ κοίτης τάξιν·
\vs{28b}ἐπὶ τούτοις πᾶσιν οὐκ ἐξίσταταί σου ἡ διάνοια, οὐδὲ διαλλάσσεταί σου ἡ καρδία ἀπὸ σώματος.
\vs{29}Καὶ ἐὰν συνῇ ἐπεκτάσεις νεφέλης, ἰσότητα σκηνῆς αὐτοῦ,
\vs{30}ἰδοὺ ἐκτενεῖ ἐπʼ αὐτὸν ἠδὼ, καὶ ῥιζώματα τῆς θαλάσσης ἐκάλυψεν.
\vs{31}Ἐν γὰρ αὐτοῖς κρινεῖ λαούς, δώσει τροφὴν τῷ ἰσχύοντι.
\vs{32}Ἐπὶ χειρῶν ἐκάλυψε φῶς, καὶ ἐνετείλατο περὶ αὐτῆς ἐν ἀπαντῶντι.
\vs{33}Ἀναγγελεῖ περὶ αὐτοῦ φίλον αὐτοῦ Κύριος, κτῆσι καὶ περὶ ἀδικίας.

\ch{37}
Καὶ ἀπὸ ταύτης ἐταράχθη ἡ καρδία μου, καὶ ἀπεῤῥύη ἐκ τοῦ τόπου αὐτῆς.
\vs{2}Ἄκουε ἀκοὴν ἐν ὀργῇ θυμοῦ Κυρίου, καὶ μελέτη ἐκ στόματος αὐτοῦ ἐξελεύσεται.
\vs{3}Ὑποκάτω παντὸς τοῦ οὐρανοῦ ἡ ἀρχὴ αὐτοῦ, καὶ τὸ φῶς αὐτοῦ ἐπὶ πτερύγων τῆς γῆς.
\vs{4}Ὀπίσω αὐτοῦ βοήσεται φωνῇ, βροντήσει ἐν φωνῇ ὕβρεως αὐτοὺ· καὶ οὐκ ἀνταλλάξει αὐτούς, ὅτι ἀκούσει φωνὴν αὐτοῦ.
\vs{5}Βροντήσει ὁ ἰσχυρὸς ἐν φωνῇ αὐτοῦ θαυμάσια· ἐποίησε γὰρ μεγάλα ἃ οὐκ ᾔδειμεν,
\vs{6}συντάσσων χιόνι, γίνου ἐπὶ τῆς, γῆς, καὶ χειμὼν ὑετὸς, καὶ χειμὼν ὑετῶν δυναστείας αὐτοῦ.
\vs{7}Ἐν χειρὶ παντὸς ἀνθρώπου κατασφραγίζει, ἵνα γνῷ πᾶς ἄνθρωπος τὴν ἑαυτοῦ ἀσθένειαν.
\vs{8}Εἰσῆλθε δὲ θηρία ὑπὸ τὴν σκέπην, ἡσύχασαν δὲ ἐπὶ κοίτης.
\vs{9}Ἐκ ταμιείων ἐπέρχονται ὀδύναι, ἀπὸ δὲ ἀκρωτηρίων ψῦχος·
\vs{10}Καὶ ἀπὸ πνοῆς ἰσχυροῦ δώσει πάγος· οἰακίζει δὲ τὸ ὕδωρ ὡς ἐὰν βούληται,
\vs{11}καὶ ἐκλεκτὸν καταπλάσσει νεφέλη· διασκορπιεῖ νέφος φῶς αὐτοῦ,
\vs{12}καὶ αὐτὸς κυκλώματα διαστρέψει, ἐν θεεβουλαθὼθ, εἰς ἔργα αὐτῶν· πάντα ὅσα ἂν ἐντείληται αὐτοῖς, ταῦτα συντέτακται παρʼ αὐτοῦ ἐπὶ τῆς γῆς,
\vs{13}ἐάν τε εἰς παιδείαν, ἐὰν εἰς τὴν γῆν αὐτοῦ, ἐὰν εἰς ἔλεος εὑρήσει αὐτόν.

\vs{14}Ἐνωτίζου ταῦτα, Ἰώβ· στῆθι νουθετούμενος δύναμιν Κυρίου.
\vs{15}Οἴδαμεν ὅτι ὁ Θεὸς ἔθετο ἔργα αὐτοῦ, φῶς ποιήσας ἐκ σκότους.
\vs{16}Ἐπίσταται δὲ διάκρισιν νεφῶν, ἐξαίσια δὲ πτώματα πονηρῶν.
\vs{17}Σοῦ δὲ ἡ στολὴ θερμὴ, ἡσυχάζεται δὲ ἐπὶ τῆς γῆς.
\vs{18}Στερεώσεις μετʼ αὐτοῦ εἰς παλαιώματα, ἰσχυραὶ ὡς ὅρασις ἐπιχύσεως.
\vs{19}Διατὶ δίδαξόν με, τί ἐροῦμεν αὐτῷ; καὶ παυσώμεθα πολλὰ λέγοντες.
\vs{20}Μὴ βίβλος ἢ γραμματεύς μοι παρέστηκεν, ἵνα ἄνθρωπον ἑστηκὼς κατασιωπήσω;

\vs{21}Πᾶσι δὲ οὐχ ὁρατὸν τὸ φῶς, τηλαυγές ἐστιν ἐν τοῖς παλαιώμασιν, ὥσπερ τὸ παρʼ αὐτοῦ ἐπὶ νεφῶν.
\vs{22}Ἀπὸ Βοῤῥᾶ νέφη χρυσαυγοῦντα, ἐπὶ τούτοις μεγάλη ἡ δόξα καὶ τιμὴ παντοκράτορος,
\vs{23}καὶ οὐχ εὑρίσκομεν ἄλλον ὅμοιον τῇ ἰσχύϊ αὐτοῦ· ὁ τὰ δίκαια κρίνων, οὐκ οἴει ἐπακούειν αὐτόν;
\vs{24}Διὸ φοβηθήσονται αὐτὸν οἱ ἄνθρωποι, φοβηθήσονται δὲ αὐτὸν καὶ οἱ σοφοὶ καρδίᾳ.

\ch{38}
Μετὰ δὲ τὸ παύσασθαι Ἐλιοὺν τῆς λέξεως, εἶπεν ὁ Κύριος τῷ Ἰὼβ διὰ λαίλαπος καὶ νεφῶν,

\vs{2}Τίς οὗτος ὁ κρύπτων με βουλὴν, συνέχεν δὲ ῥήματα ἐν καρδίᾳ, ἐμὲ δὲ οἴεται κρύπτειν;
\vs{3}Ζώσαι ὥσπερ ἀνὴρ τὴν ὀσφύν σου· ἐρωτήσω δέ σε, σὺ δέ μοι ἀποκρίθητι.

\vs{4}Ποῦ ἦς ἐν τῷ θεμελιοῦν με τὴν γῆν; ἀπάγγειλον δέ μοι εἰ ἐπίστῃ σύνεσιν.
\vs{5}Τίς ἔθετο τὰ μέτρα αὐτῆς, εἰ οἶδας; ἢ τίς ὁ ἐπαγαγὼν σπαρτίον ἐπʼ αὐτῆς;
\vs{6}Ἐπὶ τίνος οἱ κρίκοι αὐτῆς πεπήγασι; τίς δέ ἐστιν ὁ βαλὼν λίθον γωνιαῖον ἐπʼ αὐτῆς;
\vs{7}Ὅτε ἐγενήθησαν ἄστρα, ᾔνεσάν με φωνῇ μεγάλῃ πάντες ἄγγελοί μου.
\vs{8}Ἔφραξα δὲ θάλασσαν πύλαις, ὅτε ἐμαίμασσεν ἐκ κοιλίας μητρὸς αὐγῆς ἐκπορευομένη·
\vs{9}Ἐθέμην δὲ αὐτῇ νέφος ἀμφίασιν, ὁμίχλῃ δὲ αὐτὴν ἐσπαργάνωσα.
\vs{10}Ἐθέμην δὲ αὐτῇ ὅρια, περιθεὶς κλεῖθρα καὶ πύλας.
\vs{11}Εἶπα δὲ αὐτῇ, μέχρι τούτου ἐλεύσῃ, καὶ οὐχ ὑπερβήσῃ, ἀλλʼ ἐν σεαυτῇ συντριβήσεταί σου τὰ κύματα.

\vs{12}Ἢ ἐπὶ σοῦ συντέταχα φέγγος πρωϊνόν; Ἑωσφόρος δὲ ἶδε τὴν ἑαυτοῦ τάξιν,
\vs{13}ἐπιλαβέσθαι πτερύγων γῆς, ἐκτινάξαι ἀσεβεῖς ἐξ αὐτῆς;
\vs{14}Ἢ σὺ λαβὼν γῆν πηλὸν, ἔπλασας ζῶον, καὶ λαλητὸν αὐτὸν ἔθου ἐπὶ γῆς;
\vs{15}Ἀφεῖλες δὲ ἀπὸ ἀσεβῶν τὸ φῶς, βραχίονα δὲ ὑπερηφάνων συνέτριψας;

\vs{16}Ἦλθες δὲ ἐπὶ πηγὴν θαλάσσης, ἐν δὲ ἴχνεσιν ἀβύσσου περιεπάτησας;
\vs{17}Ἀνοίγονται δέ σοι φόβῳ πύλαι θανάτου, πυλωροὶ δὲ ᾅδου ἰδόντες σε ἔπτηξαν;
\vs{18}Νενουθέτησαι δὲ τὸ εὖρος τῆς ὑπʼ οὐρανόν; ἀνάγγειλον δή μοι, πόση τίς ἐστι;

\vs{19}Ποίᾳ δὲ γῇ αὐλίζεται τὸ φῶς; σκότους δὲ ποῖος ὁ τόπος;
\vs{20}Εἰ ἀγάγοις με εἰς ὅρια αὐτῶν, εἰ δὲ καὶ ἐπίστασαι τρίβους αὐτῶν·
\vs{21}Οἶδα ἄρα ὅτι τότε γεγέννησαι, ἀριθμὸς δὲ ἐτῶν σου πολύς;

\vs{22}Ἦλθες δὲ ἐπὶ θησαυροὺς χιόνος, θησαυροὺς δὲ χαλάζης ἑώρακας;
\vs{23}Ἀπόκειται δέ σοι εἰς ὥραν ἐχθρῶν, εἰς ἡμέραν πολέμων καὶ μάχης;
\vs{24}Πόθεν δὲ ἐκπορεύεται πάχνη, ἢ διασκεδάννυται Νότος εἰς τὴν ὑπʼ οὐρανόν;
\vs{25}Τίς δὲ ἡτοίμασεν ὑετῷ λάβρῳ ῥύσιν, ὁδὸν δὲ κυδοιμῶν,
\vs{26}τοῦ ὑετίσαι ἐπὶ γῆν οὗ οὐκ ἀνήρ, ἔρημον οὗ οὐχ ὑπάρχει ἄνθρωπος ἐν αὐτῇ,
\vs{27}τοῦ χορτάσαι ἄβατον καὶ ἀοίκητον, καὶ τοῦ ἐκβλαστῆσαι ἔξοδον χλόης;

\vs{28}Τίς ἐστιν ὑετοῦ πατήρ; τίς δέ ἐστιν ὁ τετοκὼς βώλους δρόσου;
\vs{29}Ἐκ γαστρὸς δὲ τίνος ἐκπορεύεται ὁ κρύσταλλος; πάχνην δὲ ἐν οὐρανῷ τίς τέτοκεν,
\vs{30}ἣ καταβαίνει ὥσπερ ὕδωρ ῥέον; πρόσωπον ἀσεβοῦς τίς ἔπτηξε;

\vs{31}Συνῆκας δὲ δεσμὸν Πλειάδος, καὶ φραγμὸν Ὠρίωνος ἤνοιξας;
\vs{32}Ἢ διανοίξεις μαζουρὼθ ἐν καιρῷ αὐτοῦ, καὶ ἕσπερον ἐπὶ κόμης αὐτοῦ ἄξεις αὐτά;
\vs{33}Ἐπίστασαι δὲ τροπὰς οὐρανοῦ, ἢ τὰ ὑπʼ οὐρανὸν ὁμοθυμαδὸν γινόμενα;
\vs{34}καλέσεις δὲ νέφος φωνῇ, καὶ τρόμῳ ὕδατος λάβρου ὑπακούσεταί σου;
\vs{35}Ἀποστελεῖς δὲ κεραυνοὺς καὶ πορεύσονται; ἐροῦσι δέ σοι, τί ἐστι;
\vs{36}Τίς δὲ ἔδωκε γυναιξὶν ὑφάσματος σοφίαν, ἢ ποικιλτικὴν ἐπιστήμην;
\vs{37}Τίς δὲ ὁ ἀριθμῶν νέφη σοφίᾳ, οὐρανὸν δὲ εἰς γῆν ἔκλινε;
\vs{38}Κέχυται δὲ ὥσπερ γῆ κονία, κεκόλληκα δὲ αὐτὸν ὥσπερ λίθῳ κύβον.

\vs{39}Θηρεύσεις δὲ λέουσι βορὰν, ψυχὰς δὲ δρακόντων ἐμπλήσεις;
\vs{40}Δεδοίκασι γὰρ ἐν κοίταις αὐτῶν, κάθηνται δὲ ἐν ὕλαις ἐνεδρεύοντες.
\vs{41}Τίς δὲ ἡτοίμασε κόρακι βοράν; νεοσσοὶ γὰρ αὐτοῦ πρὸς Κύριον κεκράγασι πλανώμενοι, τὰ σῖτα ζητοῦντες.

\ch{39}
Εἰ ἔγνως καιρὸν τοκετοῦ τραγελάφων πέτρας, ἐφύλαξας δὲ ὠδῖνας ἐλάφων,
\vs{2}ἠρίθμησας δὲ μῆνας αὐτῶν πλήρεις τοκετοῦ αὐτῶν, ὠδῖνας δὲ αὐτῶν ἔλυσας,
\vs{3}ἐξέθρεψας δὲ αὐτῶν τὰ παιδία ἔξω φόβου, ὠδῖνας δὲ αὐτῶν ἐξαποστελεῖς,
\vs{4}ἀποῤῥήξουσι τὰ τέκνα αὐτῶν, πληθυνθήσονται ἐν γενγήματι· ἐξελεύσονται, καὶ οὐ μὴ ἀνακάμψουσιν αὐτοῖς.

\vs{5}Τίς δέ ἐστιν ὁ ἀφεὶς ὄνον ἄγριον ἐλεύθερον; δεσμοὺς δὲ αὐτοῦ τίς ἔλυσεν;
\vs{6}Ἐθέμην δὲ τὴν δίαιταν αὐτοῦ ἔρημον, καὶ τὰ σκηνώματα αὐτοῦ ἁλμυρίδα.
\vs{7}Καταγελῶν πολυοχλίας πόλεως, μέμψιν δὲ φορολόγου οὐκ ἀκούων,
\vs{8}κατασκέψεται ὄρη νομὴν αὐτοῦ, καὶ ὀπίσω παντὸς χλωροῦ ζητεῖ.

\vs{9}Βουλήσεται δέ σοι μονόκερως δουλεῦσαι, ἢ κοιμηθῆναι ἐπὶ φάτνης σου;
\vs{10}Δήσεις δὲ ἐν ἱμᾶσι ζυγὸν αὐτοῦ, ἢ ἑλκύσει σου αὔλακας ἐν πεδίῳ;
\vs{11}Πέποιθας δὲ ἐπʼ αὐτῷ, ὅτι πολλὴ ἡ ἰσχὺς αὐτοῦ, ἐπαφήσεις δὲ αὐτῷ τὰ ἔργα σου;
\vs{12}Πιστεύσεις δὲ, ὅτι ἀποδώσει σοι τὸν σπόρον, εἰσοίσει δέ σου τὸν ἅλωνα;

\vs{13}Πτέρυξ τερπομένων νεέλασσα, ἐὰν συλλάβῃ ἁσίδα καὶ νέσσα·
\vs{14}Ὅτι ἀφήσει εἰς γῆν τὰ ὠὰ αὐτῆς, καὶ ἐπὶ χοῦν θάλψει,
\vs{15}καὶ ἐπελάθετο, ὅτι ποῦς σκορπιεῖ, καὶ θηρία ἀγροῦ καταπατήσει.
\vs{16}Ἀπεσκλήρυνε τὰ τέκνα ἑαυτῆς, ὥστε μὴ ἑαυτήν· εἰς κενὸν ἐκοπίασεν ἄνευ φόβου.
\vs{17}Ὅτι κατεσιώπησεν αὐτῇ ὁ Θεὸς σοφίαν, καὶ οὐκ ἐπεμέρισεν αὐτῇ ἐν τῇ συνέσει.
\vs{18}Κατὰ καιρὸν ἐν ὕψει ὑψώσει, καταγελάσεται ἵππου, καὶ τοῦ ἐπιβάτου αὐτοῦ.

\vs{19}Ἢ σὺ περιέθηκας ἵππῳ δύναμιν, ἐνέδυσας δὲ τραχήλῳ αὐτοῦ φόβον;
\vs{20}Περιέθηκας δὲ αὐτῷ πανοπλίαν; δόξαν δὲ στηθέων αὐτοῦ τόλμῃ.
\vs{21}Ἀνορύσσων ἐν πεδίῳ γαυριᾷ, ἐκπορεύεται δὲ εἰς πεδίον ἐν ἰσχύϊ.
\vs{22}Συναντῶν βασιλεῖ καταγελᾷ, καὶ οὐ μὴν ἀποστραφῇ ἀπὸ σιδήρου.
\vs{23}Ἐπʼ αὐτῷ γαυριᾷ τόξον καὶ μάχαιρα,
\vs{24}καὶ ὀργὴ ἀφανιεῖ τὴν γῆν· καὶ οὐ μὴ πιστεύσει, ἕως ἂν σημάνῃ σάλπιγξ.
\vs{25}Σάλπιγγος δὲ σημαινούσης, λέγει, εὖγε· πόῤῥωθεν δὲ ὀσφραίνεται πολέμου σὺν ἅλματι καὶ κραυγῇ.

\vs{26}Ἐκ δὲ τῆς σῆς ἐπιστήμης ἕστηκεν ἱέραξ, ἀναπετάσας τὰς πτέρυγας, ἀκίνητος, καθορῶν τὰ πρὸς Νότον;
\vs{27}Ἐπὶ δὲ σῷ προστάγματι ὑψοῦται ἀετὸς, γὺψ δὲ ἐπὶ νοσσιᾶς αὐτοῦ καθεσθεὶς αὐλίζεται,
\vs{28}ἐπʼ ἐξοχῇ πέτρας, καὶ ἀποκρύφῳ,
\vs{29}ἐκεῖσε ὢν ζητεῖ τὰ σῖτα, πόῤῥωθεν οἱ ὀφθαλμοὶ αὐτοῦ σκοπεύουσι.
\vs{30}Νεοσσοὶ δὲ αὐτοῦ φύρονται ἐν αἵματι, οὗ δʼ ἂν ὦσι τεθνεῶτες, παραχρῆμα εὑρίσκονται.

\ch{40}
Καὶ ἀπεκρίθη Κύριος ὁ Θεὸς τῷ Ἰὼβ, καὶ εἶπε,
\vs{2}μὴ κρίσιν μετὰ ἱκανοῦ ἐκκλίνει; ἐλέγχων δὲ Θεὸν, ἀποκριθήσεται αὐτήν.
\vs{3}Ὑπολαβὼν δὲ Ἰὼβ λέγει τῷ Κυρίῳ,
\vs{4}τί ἔτι ἐγὼ κρίνομαι, νουθετούμενος καὶ ἐλέγχων Κύριον, ἀκούων τοιαῦτα οὐθὲν ὤν; ἐγὼ δὲ τίνα ἀπόκρισιν δῶ πρὸς ταῦτα; χεῖρα θήσω ἐπὶ στόματί μου.
\vs{5}Ἅπαξ λελάληκα, ἐπὶ δὲ τῷ δευτέρῳ οὐ προσθήσω.

\vs{6}Ἔτι δὲ ὑπολαβὼν ὁ Κύριος, εἶπε τῷ Ἰὼβ ἐκ τοῦ νέφους,

\vs{7}Μὴ, ἀλλὰ ζῶσαι ὥσπερ ἀνὴρ τὴν ὀσφύν σου, ἐρωτήσω δέ σε, σὺ δέ μοι ἀπόκριναι.
\vs{8}Μὴ ἀποποιοῦ μου τὸ κρίμα· οἴει δέ με ἄλλως σοι κεχρηματικέναι, ἢ ἳνα ἀναφανῇς δίκαιος;
\vs{9}Ἢ βραχίων σοί ἐστι κατὰ τοῦ Κυρίου, ἢ φωνῇ κατʼ αὐτοῦ βροντᾷς;
\vs{10}Ἀνάλαβε δὴ ὕψος καὶ δύναμιν, δόξαν δὲ καὶ τιμὴν ἀμφίασαι.
\vs{11}Ἀπόστειλον δὲ ἀγγέλους ὀργῇ, πὰντα δὲ ὑβριστὴν ταπείνωσον.
\vs{12}Ὑπερήφανον δὲ σβέσον, σῆψον δὲ ἀσεβεῖς παραχρῆμα.
\vs{13}Κρύψον δὲ εἰς γῆν ὁμοθυμαδόν, τὰ δὲ πρόσωπα αὐτῶν ἀτιμίας ἔμπλησον.
\vs{14}Ὁμολογήσω ὅτι δύναται ἡ δεξιά σου σῶσαι.

\vs{15}Ἀλλὰ δὴ ἰδοὺ θηρία παρὰ σοὶ, χόρτον ἶσα βουσὶν ἐσθίουσιν.
\vs{16}Ἰδοὺ δὴ ἡ ἰσχὺς αὐτοῦ ἐπʼ ὀσφύϊ, ἡ δὲ δύναμις αὐτοῦ ἐπʼ ὀμφαλοῦ γαστρός·
\vs{17}Ἔστησεν οὐρὰν ὡς κυπάρισσον, τὰ δὲ νεῦρα αὐτοῦ συμπέπλεκται.
\vs{18}Αἱ πλευραὶ αὐτοῦ, πλευραὶ χάλκειαι, ἡ δὲ ῥάχις αὐτοῦ σίδηρος χυτός.
\vs{19}Τουτέστιν ἀρχὴ πλάσματος Κυρίου· πεποιημένον ἐγκαταπαίζεσθαι ὑπὸ τῶν ἀγγέλων αὐτοῦ.
\vs{20}Ἐπελθὼν δὲ ἐπʼ ὄρος ἀκρότομον, ἐποίησε χαρμονὴν τετράποσιν ἐν τῷ ταρτάρῳ.
\vs{21}Ὑπὸ παντοδαπὰ δένδρα κοιμᾶται, παρὰ πάπυρον καὶ κάλαμον καὶ βούτομον.
\vs{22}Σκιάζονται δὲ ἐν αὐτῷ δένδρα μεγάλα σὺν ῥαδάμνοις, καὶ κλῶνες ἀγροῦ.
\vs{23}Ἐὰν γένηται πλημμύρα, οὐ μὴ αἰσθηθῇ· πέποιθεν, ὅτι προσκρούσει ὁ Ἰορδάνης εἰς τὸ στόμα αὐτοῦ.
\vs{24}Ἐν τῷ ὀφθαλμῷ αὐτοῦ δέξεται αὐτόν, ἐνσκολιευόμενος τρήσει ῥῖνα.

\vs{25}Ἄξεις δὲ δράκοντα ἐν ἀγκίστρῳ, περιθήσεις δὲ φορβαίεαν περὶ ῥῖνα αὐτοῦ;
\vs{26}Ἤ δήσεις κρίκον ἐν τῷ μυκτῆρι αὐτοῦ, ψελλίῳ δὲ τρυπήσεις τὸ χεῖλος αὐτοῦ;
\vs{27}Λαλήσει δέ σοι δεήσει, ἱκετηρίᾳ μαλακῶς;
\vs{28}Θήσεται δὲ μετὰ σοῦ διαθήκην; λήψῃ δὲ αὐτὸν δοῦλον αἰώνιον;
\vs{29}Παίξῃ δὲ ἐν αὐτῷ ὥσπερ ὀρνέῳ; ἢ δήσεις αὐτὸν ὥσπερ στρουθίον παιδίῳ;
\vs{30}Ἐνσιτοῦνται δὲ ἐν αὐτῷ ἔθνη, μεριτεύονται δὲ αὐτὸν Φοινίκων ἔθνη;
\vs{31}Πᾶν δὲ πλωτὸν συνελθὸν οὐ μὴ ἐνέγκωσι βύρσαν μίαν οὐρᾶς αὐτοῦ, καὶ ἐν πλοίοις ἁλιέων κεφαλὴν αὐτοῦ.
\vs{32}Ἐπιθήσεις δὲ αὐτῷ χεῖρα, μνησθεὶς πόλεμον τὸν γινόμενον ἐν στόματι αὐτοῦ, καὶ μηκέτι γινέσθω.

\ch{41}
Οὐχ ἐώρακας αὐτόν; οὐδὲ ἐπὶ τοῖς λεγομένοις τεθαύμακας;
\vs{2}Οὐ δέδοικας, ὅτι ἡτοίμασταί μοι; τίς γάρ ἐστιν ὁ ἐμοὶ ἀντιστάς;
\vs{3}Ἢ τίς ἀντιστήσεταί μοι, καὶ ὑπομενεῖ; εἰ πᾶσα ἡ ὑπʼ οὐρανὸν ἐμή ἐστιν,

\vs{4}Οὐ σιωπήσομαι διʼ αὐτόν· καὶ λόγον δυνάμεως ἐλεήσει τὸν ἶσον αὐτῷ.
\vs{5}Τίς ἀποκαλύψει πρόσωπον ἐνδύσεως αὐτοῦ, εἰς δὲ πτύξιν θώρακος αὐτοῦ τίς ἂν εἰσέλθοι;
\vs{6}Πύλας προσώπου αὐτοῦ τίς ἀνοίξει; κύκλῳ ὀδόντων αὐτοῦ φόβος.
\vs{7}Τὰ ἔγκατα αὐτοῦ ἀσπίδες χάλκεαι. σύνδεσμος δὲ αὐτοῦ, ὥσπερ σμυρίτης λίθος.
\vs{8}Εἷς τοῦ ἑνὸς κολλῶνται, πνεῦμα δὲ οὐ μὴ διέλθῃ αὐτόν.
\vs{9}Ἀνὴρ τῷ ἀδελφῷ αὐτοῦ προσκολληθήσεται· συνέχονται καὶ οὐ μὴ ἀποσπασθῶσιν.
\vs{10}Ἐν πταρμῷ αὐτοῦ ἐπιφαύσκεται φέγγος, οἱ δὲ ὀφθαλμοὶ αὐτοῦ εἶδος Ἑωσφόρου.
\vs{11}Ἐκ στόματος αὐτοῦ ἐκπορεύονται ὡς λαμπάδες καιόμεναι, καὶ διαῤῥιπτοῦνται ἑς ἐσχάραι πυρός.
\vs{12}Ἐκ μυκτήρων αὐτοῦ ἐκπορεύεται καπνὸς καμίνου καιομένης πυρὶ ἀνθράκων.
\vs{13}Ἡ ψυχὴ αὐτοῦ ἄνθρακες, φλὸξ δὲ ἐκ στόματος αὐτοῦ ἐκπορεύεται·
\vs{14}Ἐν δὲ τραχήλῳ αὐτοῦ αὐλίζεται δύναμις, ἔμπροσθεν αὐτοῦ τρέχει ἀπώλεια.
\vs{15}Σάρκες δὲ σώματος αὐτοῦ κεκόλληνται· καταχέει ἐπʼ αὐτὸν, οὐ σαλευθήσεται.
\vs{16}Ἡ καρδία αὐτοῦ πέπηγεν ὡς λίθος, ἕστηκε δὲ ὥσπερ ἄκμων ἀνήλατος.
\vs{17}Στραφέντος δὲ αὐτοῦ, φόβος θηρίοις τετράποσιν ἐπὶ γῆς ἁλλομένοις.
\vs{18}Ἐὰν συναντήσωσιν αὐτῷ λόγχαι, οὐδὲν μὴ ποιήσωσι, δόρυ, καὶ θώρακα.
\vs{19}Ἥγηται μὲν γὰρ σίδηρον ἄχυρα, χαλκὸν δὲ ὥσπερ ξύλον σαθρόν.
\vs{20}Οὐ μὴ τρώσῃ αὐτὸν τόξον χάλκεον· ἥγηται μὲν πετροβόλον χόρτον.
\vs{21}Ὡς καλάμη ἐλογίσθησαν σφυρά, καταγελᾷ δὲ σεισμοῦ πυρφόρου.
\vs{22}Ἡ στρωμνὴ αὐτοῦ ὀβελίσκοι ὀξεῖς, πᾶς δὲ χρυσὸς θαλάσσης ὑπʼ αὐτὸν ὥσπερ πηλὸς ἀμύθητος.
\vs{23}Ἀναζεῖ τὴν ἄβυσσον ὥσπερ χαλκεῖον· ἥγηται δὲ τὴν θάλασσαν ὥσπερ ἐξάλειπτρον,
\vs{24}τὸν δὲ τάρταρον τῆς ἀβύσσου ὥσπερ αἰχμάλωτον· ἐλογίσατο ἄβυσσον εἰς περίπατον.
\vs{25}Οὐκ ἔστιν οὐδὲν ἐπὶ τῆς γῆς ὅμοιον αὐτῷ, πεποιημένον ἐκαταπαίζεσθαι ὑπὸ τῶν ἀγγέλων μου.
\vs{26}Πᾶν ὑψηλὸν ὁρᾷ· αὐτὸς δὲ βασιλεὺς πάντων τῶν ἐν τοῖς ὕδασιν.

\ch{42}
Ὑπολαβὼν δὲ Ἰὼβ, λέγει τῷ Κυρίῳ,

\vs{2}Οἶδα ὅτι πάντα δύνασαι, ἀδυνατεῖ δέ σοι οὐδέν.
\vs{3}Τίς γάρ ἐστιν ὁ κρύπτων σε βουλήν; φειδόμενος δὲ ῥημάτων, καὶ σὲ οἴεται κρύπτειν; τίς δὲ ἀναγγελεῖ μοι ἃ οὐκ ᾔδειν, μεγάλα καὶ θαυμαστὰ ἃ οὐκ ἐπιστάμην;

\vs{4}Ἄκουσον δέ μου Κύριε, ἵνα κᾀγὼ λαλήσω· ἐρωτήσω δέ σε, σὺ δέ με δίδαξον.
\vs{5}Ἀκοὴν μὲν ὠτὸς ἤκουόν σου τὸ πρότερον, νυνὶ δὲ ὁ ὀφθαλμός μου ἑώρακέ σε.
\vs{6}Διὸ ἐφαύλισα ἐμαυτὸν, καὶ ἐτάκην· ἥγημαι δὲ ἐμαυτὸν γῆν καὶ σποδόν.

\vs{7}Ἐγένετο δὲ μετὰ τὸ λαλῆσαι τὸν Κύριον πάντα τὰ ῥήματα ταῦτα τῷ Ἰώβ, εἶπεν ὁ Κύριος Ἐλὺιφὰς τῷ Θαιμανὺίτῃ, ἥμαρτες σὺ, καὶ οἱ δὺο φίλοι σου· οὐ γὰρ ἐλαλήσατε ἐνώπιόν μου ἀληθὲς οὐδὲν, ὥσπερ ὁ θεράπων μου Ἰώβ.
\vs{8}Νῦν δὲ λάβετε ἑπτὰ μόσχους, καὶ ἑπτὰ κριοὺς, καὶ πορεύθητε πρὸς τὸν θεράποντά μου Ἰὼβ, καὶ ποιήσει κάρπωσιν ὑπὲρ ὑμῶν. Ἰὼβ δὲ ὁ θεράπων μου εὔξεται περὶ ὑμῶν, ὅτι εἰ μὴ πρόσωπον αὐτοῦ λήψομαι· εἰ μὴ γὰρ διʼ αὐτὸν, ἀπώλεσα ἂν ὑμᾶς· οὐ γὰρ ἐλαλήσατε ἀληθὲς κατὰ τοῦ θεράποντός μου Ἰώβ.

\vs{9}Ἐπορεύθη δὲ Ἐλιφὰζ ὁ Θαιμανίτης, καὶ Βαλδὰδ ὁ Σαυχίτης, καὶ Σωφὰρ ὁ Μιναῖος, καὶ ἐποίησαν καθὼς συνέταξεν αὐτοῖς ὁ Κύριος· καὶ ἔλυσε τὴν ἁμαρτίαν αὐτοῖς διὰ Ἰὼβ.

\vs{10}Ὁ δὲ Κύριος ηὔξησε τὸν Ἰώβ· εὐξαμένου δὲ αὐτοῦ καὶ περὶ τῶν φίλων αὐτοῦ, ἀφῆκεν αὐτοῖς τὴν ἁμαρτίαν· ἔδωκε δὲ ὁ Κύριος διπλᾶ, ὅσα ἦν ἔμπροσθεν Ἰὼβ εἰς διπλασιασμόν.
\vs{11}Ἤκουσαν δὲ πάντες οἱ ἀδελφοὶ αὐτοῦ καὶ αἱ ἀδελφαὶ αὐτοῦ πάντα τὰ συμβεβηκότα αὐτῷ, καὶ ἦλθον πρὸς αὐτὸν, καὶ πάντες ὅσοι ᾔδεισαν αὐτὸν ἐκ πρώτου· φαγόντες δὲ καὶ πιόντες παρʼ αὐτῷ παρεκάλεσαν αὐτὸν, καὶ ἐθαύμασαν ἐπὶ πᾶσιν οἷς ἐπήγαγεν ἐπʼ αὐτῷ ὁ Κύριος· ἔδωκε δὲ αὐτῷ ἕκαστος ἀμνάδα μίαν, καὶ τετράδραχμον χρυσοῦ καὶ ἀσήμου.

\vs{12}Ὁ δὲ Κύριος εὐλόγησε τὰ ἔσχατα Ἰὼβ, ἢ τὰ ἔμπροσθεν· ἦν δὲ τὰ κτήνη αὐτοῦ, πρόβατα μύρια τετρακισχίλια, κάμηλοι ἑξακισχίλιαι, ζεύγη βοῶν χίλια, ὄνοι θήλειαι νομάδες χίλιαι.
\vs{13}Γεννῶνται δὲ αὐτῷ υἱοὶ ἑπτὰ, καὶ θυγατέρες τρεῖς.
\vs{14}Καὶ ἐκάλεσε τὴν μὲν πρώτην, Ἡμέραν· τὴν δὲ δευτέραν, Κασίαν· τὴν δὲ τρίτην, Ἀμαλθαίας κέρας.
\vs{15}Καὶ οὐχ εὑρέθησαν κατὰ τὰς θυγατέρας Ἰὼβ, βελτίους αὐτῶν ἐν τῇ ὑπʼ οὐρανόν· ἔδωκε δὲ αὐταῖς ὁ πατὴρ κληρονομίαν ἐν τοῖς ἀδελφοῖς.

\vs{16}Ἔζησε δὲ Ἰὼβ μετὰ τὴν πληγὴν ἔτη ἑκατὸν ἑβδομήκοντα· τὰ δὲ πάντα ἔτη ἔζησε, διακόσια τεσσράκοντα· καὶ Σἴδεν Ἰὼβ τοὺς υἱοὺς αὐτοῦ, καὶ τοὺς υἱοὺς τῶν υἱῶν αὐτοῦ, τετάρτην γενεάν.
\vs{17}Καὶ ἐτελεύτησεν Ἰὼβ πρεσβύτερος, καὶ πλήρης ἡμερῶν·
\vs{17a}γέγραπται δὲ, αὐτὸν πάλιν ἀναστήσεσθαι μεθʼ ὧν ὁ Κύριος ἀνίστησιν.

\vs{17b}Οὗτος ἑρμηνεύεται ἐκ τῆς Συριακῆς βίβλου, ἐν μὲν γῇ κατοικῶν τῇ Αὐσίτιδι, ἐπὶ τοῖς ὁρίοις τῆς Ἰδουμαίας καὶ Ἀραβίας· προϋπῆρχε δὲ αὐτῷ ὄνομα Ἰωβάβ·
\vs{17c}λαβὼν δὲ γυναῖκα Ἀράβισσαν, γεννᾷ υἱὸν, ᾧ ὄνομα Ἐννών· ἦν δὲ αὐτὸς πατρὸς μὲν Ζαρὲ ἐκ τῶν Ἡσαῦ υἱῶν υἱὸς, μητρὸς δὲ Βοσόῤῥας, ὥστε εἶναι αὐτὸν πέμπτον ἀπὸ Ἁβραάμ·
\vs{17d}καὶ οὗτοι οἱ βασιλεῖς οἱ βασιλεύσαντες ἐν Ἐδὼμ, ἧς καὶ αὐτὸς ἦρξε χώρας· πρῶτος Βαλὰκ ὁ τοῦ Βεὼρ, καὶ ὄνομα τῇ πόλει αὐτοῦ Δενναβά· μετὰ δὲ Βαλὰκ, Ἰωβὰβ ὁ καλούμενος Ἰώβ· μετὰ δὲ τοῦτον, Ἀσὼμ ὁ ὑπάρχων ἡγεμὼν ἐκ τῆς Θαιμανίτιδος χώρας· μετὰ δὲ τοῦτον, Ἀδὰδ υἱὸς Βαρὰδ, ὁ ἐκκόψας Μαδιὰμ ἐν τῷ πεδίῳ Μωὰβ, καὶ ὄνομα τῇ πόλει αὐτοῦ Γεθαίμ·
\vs{17e}οἱ δὲ ἐλθόντες πρὸς αὐτὸν φίλοι, Ἐλιφὰς τῶν Ἡσαῦ υἱῶν Θαιμανῶν βασιλεὺς, Βαλδὰδ ὁ Σαυχαίων τύραννος, Σωφὰρ ὁ Μιναίων βασιλεύς.


\def\book{ΨΑΛΜΟΙ}
\biblebook{ΨΑΛΜΟΙ}

\begin{psalms}
\ch{1}{1}ΜΑΚΑΡΙΟΣ ἀνὴρ, ὃς οὐκ ἐπορεύθη ἐν βουλῇ ἀσεβῶν, καὶ ἐν ὁδῷ ἁμαρτωλῶν· οὐκ ἔστη, καὶ ἐπὶ καθέδρᾳ λοιμῶν οὐκ ἐκάθισεν.
\vs{2}Ἀλλ’ ἢ ἐν τῷ νόμῳ Κυρίου τὸ θέλημα αὐτοῦ, καὶ ἐν τῷ νόμῳ αὐτοῦ μελετήσει ἡμέρας καὶ νυκτός.
\vs{3}Καὶ ἔσται ὡς τὸ ξύλον τὸ πεφυτευμένον παρὰ τὰς διεξόδους τῶν ὑδάτων, ὃ τὸν καρπὸν αὐτοῦ δώσει ἐν καιρῷ αὐτοῦ· καὶ τὸ φύλλον αὐτοῦ οὐκ ἀποῤῥυήσεται, καὶ πάντα ὅσα ἂν ποιῇ κατευοδωθήσεται.

\vs{4}Οὐχ οὕτως οἱ ἀσεβεῖς, οὐχ οὕτως, ἀλλʼ ἢ ὡς ὁ χνοῦς ὃν ἐκρίπτει ὁ ἄνεμος ἀπὸ προσώπου τῆς γῆς.
\vs{5}Διὰ τοῦτο οὐκ ἀναστήσονται οἱ ἀσεβεῖς ἐν κρίσει, οὐδὲ ἁμαρτωλοὶ ἐν βουλῇ δικαίων.
\vs{6}Ὅτι γινώσκει Κύριος ὁδὸν δικαίων, καὶ ὁδὸς ἀσεβῶν ἀπολεῖται.

\ch{2}{2}Ἱνατί ἐφρύαξαν ἔθνη, καὶ λαοὶ ἐμελέτησαν κενά;
\vs{2}Παρέστησαν οἱ βασιλεῖς τῆς γῆς, καὶ οἱ ἄρχοντες συνήχθησαν ἐπιτοαυτὸ κατὰ τοῦ Κυρίου, καὶ κατὰ τοῦ Χριστοῦ αὐτοῦ.
\vs{3}Διαῤῥήξωμεν τοὺς δεσμοὺς αὐτῶν, καὶ ἀποῤῥίψωμεν ἀφʼ ἡμῶν τὸν ζυγὸν αὐτῶν.

\vs{4}Ὁ κατοικῶν ἐν οὐρανοῖς ἐκγελάσεται αὐτοὺς, καὶ ὁ Κύριος ἐκμυκτηριεῖ αὐτούς.
\vs{5}Τότε λαλήσει πρὸς αὐτοὺς ἐν ὀργῇ αὐτοῦ, καὶ ἐν τῷ θυμῷ αὐτοῦ ταράξει αὐτούς.
\vs{6}Ἐγὼ δὲ κατεστάθην βασιλεὺς ὑπʼ αὐτοῦ ἐπὶ Σιὼν ὄρος τὸ ἅγιον αὐτοῦ,
\vs{7}διαγγέλλων τὸ πρόσταγμα Κυρίου· Κύριος εἶπε πρὸς μὲ, υἱός μου εἶ σὺ, ἐγὼ σήμερον γεγέννηκά σε.
\vs{8}Αἴτησαι παρʼ ἐμοῦ, καὶ δώσω σοι ἔθνη τὴν κληρονομίαν σου, καὶ τὴν κατάσχεσίν σου τὰ πέρατα τῆς γῆς.
\vs{9}Ποιμανεῖς αὐτοὺς ἐν ῥάβδῳ σιδηρᾷ, ὡς σκεῦος κεραμέως συντρίψεις αὐτούς.

\vs{10}Καὶ νῦν βασιλεῖς σύνετε, παιδεύθητε πάντες οἱ κρίνοντες τὴν γῆν.
\vs{11}Δουλεύσατε τῷ Κυρίῳ ἐν φόβῳ, καὶ ἀγαλλιᾶσθε αὐτῷ ἐν τρόμῳ.
\vs{12}Δράξασθε παιδείας, μή ποτε ὀργισθῇ Κύριος, καὶ ἀπολεῖσθε ἐξ ὁδοῦ δικαίας· ὅταν ἐκκαυθῇ ἐν τάχει ὁ θυμὸς αὐτοῦ, μακάριοι πάντες οἱ πεποιθότες ἐπʼ αὐτῷ.

\begin{psalmheading}{\ch{3}{3} Ψαλμὸς τῷ Δαυὶδ, ὁπότε ἀπεδίδρασκεν ἀπὸ προσώπου Ἀβεσσαλὼμ τοῦ υἱοῦ αὐτοῦ.}
\end{psalmheading}
\vs{2}Κύριε τί ἐπληθύνθησαν οἱ θλίβοντές με; πολλοὶ ἐπανίστανται ἐπʼ ἐμέ.
\vs{3}Πολλοὶ λέγουσι τῇ ψυχῇ μου, οὐκ ἔστι σωτηρία αὐτῷ ἐν τῷ Θεῷ αὐτοῦ· διάψαλμα.

\vs{4}Σὺ δὲ Κύριε, ἀντιλήπτωρ μου εἶ, δόξα μου, καὶ ὑψῶν τὴν κεφαλήν μου.
\vs{5}Φωνῇ μου πρὸς Κύριον ἐκέκραξα, καὶ ἐπήκουσέ μου ἐξ ὄρους ἁγίου αὐτοῦ· διάψαλμα.
\vs{6}Ἐγὼ ἐκοιμήθην καὶ ὕπνωσα, ἐξηγέρθην, ὅτι Κύριος ἀντιλήψεταί μου.
\vs{7}Οὐ φοβηθήσομαι ἀπὸ μυριάδων λαοῦ, τῶν κύκλῳ ἐπιτιθεμένων μοι.
\vs{8}Ἀνάστα Κύριε, σῶσόν με ὁ Θεός μου· ὅτι σὺ ἐπάταξας πάντας τοὺς ἐχθραίνοντάς μοι ματαίως, ὀδόντας ἁμαρτωλῶν συνέτριψας.
\vs{9}Τοῦ Κυρίου ἡ σωτηρία, καὶ ἐπὶ τὸν λαόν σου ἡ εὐλογία σου.

\begin{psalmheading}{\ch{4}{4} Εἰς τὸ τέλος, ἐν ψαλμοῖς ᾠδὴ τᾠ Δαυίδ.}
\end{psalmheading}
\vs{2}Ἐν τῷ ἐπικαλεῖσθαί με, εἰσήκουσέ μου ὁ Θεὸς τῆς δικαιοσύνης μου· ἐν θλίψει ἐπλάτυνάς μοι· οἰκτείρησόν με, καὶ εἰσάκουσον τῆς προσευχῆς μου.

\vs{3}Υἱοὶ ἀνθρώπων, ἕως πότε βαρυκάρδιοι; ἱνατί ἀγαπᾶτε ματαιότητα, καὶ ζητεῖτε ψεῦδος; διάψαλμα.
\vs{4}Καὶ γνῶτε ὅτι ἐθαυμάστωσε Κύριος τὸν ὅσιον αὐτοῦ, Κύριος εἰσακούσεταί μου ἐν τῷ κεκραγέναι με πρὸς αὐτόν.
\vs{5}Ὀργίζεσθε καὶ μὴ ἁμαρτάνετε· ἃ λέγετε ἐν ταῖς καρδίαις ὑμῶν, ἐπὶ ταῖς κοίταις ὑμῶν κατανύγητε· διάψαλμα.
\vs{6}Θύσατε θυσίαν δικαιοσύνης, καὶ ἐλπίσατε ἐπὶ Κύριον.

\vs{7}Πολλοὶ λέγουσι, τίς δείξει ἡμῖν τὰ ἀγαθά; ἐσημειώθη ἐφʼ ἡμᾶς τὸ φῶς τοῦ προσώπου σου, Κύριε.
\vs{8}Ἔδωκας εὐφροσύνην εἰς τὴν καρδίαν μου· ἀπὸ καρποῦ σίτου καὶ οἴνου καὶ ἐλαίου αὐτῶν ἐπληθύνθησαν.
\vs{9}Ἐν εἰρήνῃ ἐπὶ τὸ αὐτὸ κοιμηθήσομαι, καὶ ὑπνώσω· ὅτι σὺ Κύριε κατὰ μόνας ἐπʼ ἐλπίδι κατῴκισάς με.

\begin{psalmheading}{\ch{5}{5} Εἰς τὸ τέλος, ὑπὲρ τῆς κληρονομούσης, ψαλμὸς τῷ Δαυίδ.}
\end{psalmheading}
\vs{2}Τὰ ῥήματά μου ἐνώτισαι Κύριε, σύνες τῆς κραυγῆς μου,
\vs{3}πρόσχες τῇ φωνῇ τῆς δεήσεώς μου, ὁ βασιλεύς μου καὶ ὁ Θεός μου· ὅτι πρὸς σὲ προσεύξομαι Κύριε,
\vs{4}τοπρωῒ εἰσακούσῃ τῆς φωνῆς μου· τοπρωῒ παραστήσομαί σοι, καὶ ἐπόψομαι.
\vs{5}Ὅτι οὐχὶ Θεὸς θέλων ἀνομίαν σὺ εἶ· οὐδὲ παροικησει σοι πονηρευόμενος,
\vs{6}οὐδὲ διαμενοῦσι παράνομοι κατέναντι τῶν ὀφθαλμῶν σου· ἐμίσησας Κύριε πάντας τοὺς ἐργαζομένους τὴν ἀνομίαν,
\vs{7}ἀπολεῖς πάντας τοὺς λαλοῦντας τὸ ψεῦδος· ἄνδρα αἱμάτων καὶ δόλιον βδελύσσεται Κύριος.
\vs{8}Ἐγὼ δὲ ἐν τῷ πλήθει τοῦ ἐλέου σου εἰσελεύσομαι εἰς τὸν οἶκόν σου, προσκυνήσω πρὸς ναὸν ἅγιόν σου ἐν φόβῳ σου.

\vs{9}Κύριε ὁδήγησόν με ἐν τῇ δικαιοσύνῃ σου ἕνεκα τῶν ἐχθρῶν μου, κατεύθυνον ἐνώπιόν σου τὴν ὁδόν μου.
\vs{10}Ὅτι οὐκ ἔστιν ἐν τῷ στόματι αὐτῶν ἀλήθεια· ἡ καρδία αὐτῶν ματαία· τάφος ἀνεῳγμένος ὁ λάρυγξ αὐτῶν· ταῖς γλώσσαις αὐτῶν ἐδολιοῦσαν.
\vs{11}Κρίνον αὐτοὺς ὁ Θεός· ἀποπεσάτωσαν ἀπὸ τῶν διαβουλιῶν αὐτῶν· κατὰ τὸ πλῆθος τῶν ἀσεβειῶν αὐτῶν ἔξωσον αὐτοὺς, ὅτι παρεπίκρανάν σε Κύριε.

\vs{12}Καὶ εὐφρανθήτωσαν ἐπὶ σοὶ πάντες οἱ ἐλπίζοντες ἐπὶ σὲ, εἰς αἰῶνα ἀγαλλιάσονται, καὶ κατασκηνώσεις ἐν αὐτοῖς· καὶ καυχήσονται ἐπὶ σοὶ πάντες οἱ ἀγαπῶντες τὸ ὄνομά σου,
\vs{13}ὅτι σὺ εὐλογήσεις δίκαιον Κύριε, ὡς ὅπλῳ εὐδοκίας ἐστεφάνωσας ἡμᾶς.

\begin{psalmheading}{\ch{6}{6} Εἰς τὸ τέλος, ἐν ὕμνοις ὑπὲρ τῆς ὀγδόης, ψαλμὸς τῷ Δαυίδ.}
\end{psalmheading}
\vs{2}Κύριε, μὴ τῷ θυμῷ σου ἐλέγξῃς με, μηδὲ τῇ ὀργῇ σου παιδεύσῃς με.
\vs{3}Ἐλέησόν με Κύριε, ὅτι ἀσθενής εἰμι· ἴασαί με Κύριε, ὅτι ἐταράχθη τὰ ὀστᾶ μου.
\vs{4}Καὶ ἡ ψυχή μου ἐταράχθη σφόδρα· καὶ σὺ Κύριε ἕως πότε;
\vs{5}Ἐπίστρεψον Κύριε, ῥῦσαι τὴν ψυχήν μου· σῶσόν με ἕνεκεν τοῦ ἐλέους σου,
\vs{6}ὅτι οὐκ ἔστιν ἐν τῷ θανάτῳ ὁ μνημονεύων σου, ἐν δὲ τῷ ᾅδῃ τίς ἐξομολογήσεταί σοι;
\vs{7}Ἐκοπίασα ἐν στεναγμῷ μου, λούσω καθʼ ἑκάστην νύκτα τὴν κλίνην μου, ἐν δάκρυσί μου τὴν στρωμνήν μου βρέξω.
\vs{8}Ἐταράχθη ἀπὸ θυμοῦ ὁ ὀφθαλμός μου, ἐπαλαιώθην ἐν πᾶσι τοῖς ἐχθροῖς μου.

\vs{9}Ἀπόστητε ἀπʼ ἐμοῦ πάντες οἱ ἐργαζόμενοι τὴν ἀνομίαν, ὅτι εἰσήκουσε Κύριος τῆς φωνῆς τοῦ κλαυθμοῦ μου.
\vs{10}Εἰσήκουσε Κύριος τῆς δεήσεώς μου, Κύριος τὴν προσευχήν μου προσεδέξατο.
\vs{11}Αἰσχυνθείησαν καὶ ταραχθείησαν σφόδρα πάντες οἱ ἐχθροί μου, ἐπιστραφείησαν καὶ αἰσχυνθείησαν σφόδρα διὰ τάχους.

\begin{psalmheading}{\ch{7}{7} Ψαλμὸς τῷ Δαυὶδ, ὃν ᾖσε τῷ Κυρίῳ ὑπὲρ τῶν λόγων Χουσὶ υἱοῦ Ἰεμενεί.}
\end{psalmheading}
\vs{2}Κύριε ὁ Θεός μου, ἐπὶ σοὶ ἤλπισα, σῶσόν με ἐκ πάντων τῶν διωκόντων με, καὶ ῥῦσαί με,
\vs{3}μή ποτε ἁρπάσῃ ὡς λέων τὴν ψυχήν μου, μὴ ὄντος λυτρουμένου, μηδὲ σώζοντος.

\vs{4}Κύριε ὁ Θεός μου, εἰ ἐποίησα τοῦτο, εἰ ἔστιν ἀδικία ἐν χερσί μου,
\vs{5}εἰ ἀνταπέδωκα τοῖς ἀνταποδιδοῦσί μοι κακὰ, ἀποπέσοιμι ἄρα ἀπὸ τῶν ἐχθρῶν μου κενός·
\vs{6}Καταδιώξαι ὁ ἐχθρὸς τὴν ψυχήν μου καὶ καταλάβοι, καὶ καταπατήσαι εἰς γῆν τὴν ζωήν μου, καὶ τὴν δόξαν μου εἰς χοῦν κατασκηνώσαι· διάψαλμα.

\vs{7}Ἀνάστηθι Κύριε ἐν ὀργῇ σου, ὑψώθητι ἐν τοῖς πέρασι τῶν ἐχθρῶν μου· ἐξεγέρθητι Κύριε ὁ Θεός μου ἐν προστάγματι ᾧ ἐνετείλω,
\vs{8}καὶ συναγωγὴ λαῶν κυκλώσει σε· καὶ ὑπὲρ ταύτης εἰς ὕψος ἐπίστρεψον.
\vs{9}Κύριος κρινεῖ λαούς· κρίνον με Κύριε κατὰ τὴν δικαιοσύνην μου, καὶ κατὰ τὴν ἀκακίαν μου ἐπʼ ἐμοί.
\vs{10}Συντελεσθήτω δὴ πονηρία ἁμαρτωλῶν, καὶ κατευθυνεῖς δίκαιον, ἐτάζων καρδίας καὶ νεφροὺς ὁ Θεός.

\vs{11}Δικαία ἡ βοήθειά μου παρὰ τοῦ Θεοῦ τοῦ σώζοντος τοὺς εὐθεῖς τῇ καρδίᾳ.
\vs{12}Ὁ Θεὸς κριτὴς δίκαιος, καὶ ἰσχυρὸς, καὶ μακρόθυμος, μὴ ὀργὴν ἐπάγων καθʼ ἑκάστην ἡμέραν.
\vs{13}Ἐὰν μὴ ἐπιστραφῆτε, τὴν ῥομφαίαν αὐτοῦ στιλβώσει, τὸ τόξον αὐτοῦ ἐνέτεινε, καὶ ἡτοίμασεν αὐτό.
\vs{14}Καὶ ἐν αὐτῷ ἡτοίμασε σκεύη θανάτου, τὰ βέλη αὐτοῦ τοῖς καιομένοις ἐξειργάσατο.

\vs{15}Ἰδοὺ ὠδίνησεν ἀδικίαν, συνέλαβε πόνον, καὶ ἔτεκεν ἀνομίαν.
\vs{16}Λάκκον ὤρυξε καὶ ἀνέσκαψεν αὐτὸν, καὶ ἐμπεσεῖται εἰς βόθρον ὃν εἰργάσατο.
\vs{17}Ἐπιστρέψει ὁ πόνος αὐτοῦ εἰς κεφαλὴν αὐτοῦ, καὶ ἐπὶ κορυφὴν αὐτοῦ ἡ ἀδικία αὐτοῦ καταβήσεται.
\vs{18}Ἐξομολογήσομαι Κυρίῳ κατὰ τὴν δικαιοσύνην αὐτοῦ, ψαλῶ τῷ ὀνόματι Κυρίου τοῦ ὑψίστου.

\begin{psalmheading}{\ch{8}{8} Εἰς τὸ τέλος, ὑπὲρ τῶν ληνῶν, ψαλμὸς τῷ Δαυίδ.}
\end{psalmheading}
\vs{2}Κύριε ὁ Κύριος ἡμῶν, ὡς θαυμαστὸν τὸ ὄνομά σου ἐν πάσῃ τῇ γῇ; ὅτι ἐπῄρθη ἡ μεγαλοπρέπειά σου ὑπεράνω τῶν οὐρανῶν.
\vs{3}Ἐκ στόματος νηπίων καὶ θηλαζόντων κατηρτίσω αἶνον· ἕνεκα τῶν ἐχθρῶν σου, τοῦ καταλῦσαι ἐχθρὸν καὶ ἐκδικητήν.

\vs{4}Ὅτι ὄψομαι τοὺς οὐρανοὺς ἔργα τῶν δακτύλων σου, σελήνην καὶ ἀστέρας, ἃ σὺ ἐθεμελίωσας·
\vs{5}Τί ἐστιν ἄνθρωπος, ὅτι μιμνήσκῃ αὐτοῦ; ἢ υἱὸς ἀνθρώπου, ὅτι ἐπισκέπτῃ αὐτόν;
\vs{6}Ἠλάττωσας αὐτὸν βραχύ τι παρʼ ἀγγέλους, δόξῃ καὶ τιμῇ ἐστεφάνωσας αὐτὸν,
\vs{7}καὶ κατέστησας αὐτὸν ἐπὶ τὰ ἔργα τῶν χειρῶν σου· πάντα ὑπέταξας ὑποκάτω τῶν ποδῶν αὐτοῦ,
\vs{8}πρόβατα καὶ βόας πάσας, ἔτι δὲ καὶ τὰ κτήνη τοῦ πεδίου,
\vs{9}τὰ πετεινὰ τοῦ οὐρανοῦ, καὶ τοὺς ἰχθύας τῆς θαλάσσης, τὰ διαπορευόμενα τρίβους θαλασσῶν.
\vs{10}Κύριε ὁ Κύριος ἡμῶν, ὡς θαυμαστὸν ὄνομά σου ἐν πάσῃ τῇ γῇ;

\begin{psalmheading}{\ch{9}{9-10} Εἰς τὸ τέλος, ὑπὲρ τῶν κρυφίων τοῦ υἱοῦ, ψαλμὸς τῷ Δαυίδ.}
\end{psalmheading}
\vs{2}Ἐξομολογήσομαι σοι Κύριε ἐν ὅλῃ καρδίᾳ μου, διηγήσομαι πάντα τὰ θαυμάσιά σου.
\vs{3}Εὐφρανθήσομαι καὶ ἀγαλλιάσομαι ἐν σοὶ, ψαλῶ τῷ ὀνόματί σου ὕψιστε.

\vs{4}Ἐν τῷ ἀποστραφῆναι τὸν ἐχθρόν μου εἰς τὰ ὀπίσω, ἀσθενήσουσι καὶ ἀπολοῦνται ἀπὸ προσώπου σου.
\vs{5}Ὅτι ἐποίησας τὴν κρίσιν μου καὶ τὴν δίκην μου, ἐκάθισας ἐπὶ θρόνου ὁ κρίνων δικαιοσύνην.
\vs{6}Ἐπετίμησας ἔθνεσι, καὶ ἀπώλετο ὁ ἀσεβὴς· τὸ ὄνομα αὐτῶν ἐξήλειψας εἰς τὸν αἰῶνα, καὶ εἰς τὸν αἰῶνα τοῦ αἰῶνος.
\vs{7}Τοῦ ἐχθροῦ ἐξέλιπον αἱ ῥομφαῖαι εἰς τέλος, καὶ πόλεις καθεῖλες· ἀπώλετο τὸ μνημόσυνον αὐτῶν μετʼ ἤχου,
\vs{8}καὶ ὁ Κύριος εἰς τὸν αἰῶνα μένει· ἡτοίμασεν ἐν κρίσει τὸν θρόνον αὐτοῦ,
\vs{9}καὶ αὐτὸς κρινεῖ τὴν οἰκουμένην ἐν δικαιοσύνῃ, κρινεῖ λαοὺς ἐν εὐθύτητι.
\vs{10}Καὶ ἐγένετο Κύριος καταφυγὴ τῷ πένητι, βοηθὸς ἐν εὐκαιρίαις, ἐν θλίψει.
\vs{11}Καὶ ἐλπισάτωσαν ἐπὶ σὲ οἱ γινώσκοντες τὸ ὄνομά σου, ὅτι οὐκ ἐγκατέλιπες τοὺς ἐκζητοῦντάς σε Κύριε.

\vs{12}Ψάλατε τῷ Κυρίῳ τῷ κατοικοῦντι ἐν Σιὼν, ἀναγγείλατε ἐν τοῖς ἔθνεσι τὰ ἐπιτηδεύματα αὐτοῦ.
\vs{13}Ὅτι ἐκζητῶν τὰ αἵματα αὐτῶν ἐμνήσθη, οὐκ ἐπελάθετο τῆς δεήσεως τῶν πενήτων.

\vs{14}Ἐλέησόν με Κύριε, ἴδε τὴν ταπείνωσίν μου ἐκ τῶν ἐχθρῶν μου, ὁ ὑψῶν με ἐκ τῶν πυλῶν τοῦ θανάτου·
\vs{15}Ὅπως ἂν ἐξαγγείλω πάσας τὰς αἰνέσεις σου ἐν ταῖς πύλαις τῆς θυγατρὸς Σιών· ἀγαλλιάσομαι ἐπὶ τῷ σωτηρίῳ σου.

\vs{16}Ἐνεπάγησαν ἔθνη ἐν διαφθορᾷ ᾗ ἐποίησαν· ἐν παγίδι ταύτῃ ᾗ ἔκρυψαν συνελήφθη ὁ ποὺς αὐτῶν.
\vs{17}Γινώσκεται Κύριος κρίματα ποιῶν, ἐν τοῖς ἔργοις τῶν χειρῶν αὐτοῦ συνελήφθη ὁ ἁμαρτωλός· ᾠδὴ διαψάλματος.
\vs{18}Ἀποστραφήτωσαν οἱ ἁμαρτωλοὶ εἰς τὸν ᾅδην, πάντα τὰ ἔθνη τὰ ἐπιλανθανόμενα τοῦ Θεοῦ.
\vs{19}Ὅτι οὐκ εἰς τέλος ἐπιλησθήσεται ὁ πτωχὸς, ἡ ὑπομονὴ τῶν πενήτων οὐκ ἀπολεῖται εἰς τὸν αἰῶνα.
\vs{20}Ἀνάστηθι Κύριε, μὴ κραταιούσθω ἄνθρωπος, κριθήτωσαν ἔθνη ἐνώπιόν σου.
\vs{21}Κατάστησον, Κύριε, νομοθέτην ἐπʼ αὐτοὺς, γνώτωσαν ἔθνη ὅτι ἀνθρωποί εἰσι· διάψαλμα.

\vs{22}Ἱνατί, Κύριε, ἀφέστηκας μακρόθεν, ὑπερορᾷς ἐν εὐκαιρίαις, ἐν θλίψει;
\vs{23}Ἐν τῷ ὑπερηφανεύεσθαι τὸν ἀσεβῆ, ἐμπυρίζεται ὁ πτωχὸς, συλλαμβάνονται ἐν διαβουλίοις οἷς διαλογίζονται.
\vs{24}Ὅτι ἐπαινεῖται ὁ ἁμαρτωλὸς ἐν ταῖς ἐπιθυμίαις τῆς ψυχῆς αὐτοῦ, καὶ ὁ ἀδικῶν ἐνευλογεῖται.
\vs{25}Παρώξυνε τὸν Κύριον ὁ ἁμαρτωλὸς, κατὰ τὸ πλῆθος τῆς ὀργῆς αὐτοῦ οὐκ ἐκζητήσει· οὐκ ἔστιν ὁ Θεὸς ἐνώπιον αὐτοῦ.
\vs{26}Βεβηλοῦνται αἱ ὁδοὶ αὐτοῦ ἐν παντὶ καιρῷ· ἀνταναιρεῖται τὰ κρίματά σου ἀπὸ προσώπου αὐτοῦ, πάντων τῶν ἐχθρῶν αὐτοῦ κατακυριεύσει.
\vs{27}Εἶπε γὰρ ἐν καρδίᾳ αὐτοῦ, οὐ μὴ σαλευθῶ ἀπὸ γενεᾶς εἰς γενεὰν ἄνευ κακοῦ.
\vs{28}Οὗ ἀρᾶς τὸ στόμα αὐτοῦ γέμει καὶ πικρίας καὶ δόλου, ὑπὸ τὴν γλῶσσαν αὐτοῦ κόπος καὶ πόνος.
\vs{29}Ἐγκάθηται ἔνεδρα μετὰ πλουσίων ἐν ἀποκρύφοις, τοῦ ἀποκτεῖναι ἀθῶον· οἱ ὀφθαλμοὶ αὐτοῦ εἰς τὸν πένητα ἀποβλέπουσιν.
\vs{30}Ἐνεδρεύει ἐν ἀποκρύφῳ ὡς λέων ἐν τῇ μάνδρᾳ αὐτοῦ· ἐνεδρεύει τοῦ ἁρπάσαι πτωχὸν, ἁρπάσαι πτωχὸν ἐν τῷ ἑλκύσαι αὐτόν· ἐν τῇ παγίδι αὐτοῦ
\vs{31}ταπεινώσει αὐτὸν, κύψει καὶ πεσεῖται ἐν τῷ αὐτὸν κατακυριεῦσαι τῶν πενήτων.
\vs{32}Εἶπε γὰρ ἐν τῇ καρδίᾳ αὐτοῦ, ἐπιλέλησται ὁ Θεὸς, ἀπέστρεψε τὸ πρόσωπον αὐτοῦ τοῦ μὴ βλέπειν εἰς τέλος.

\vs{33}Ἀνάστηθι Κύριε ὁ Θεὸς, ὑψωθήτω ἡ χείρ σου, μὴ ἐπιλάθῃ τῶν πενήτων.
\vs{34}Ἕνεκεν τίνος παρώξυνεν ὁ ἀσεβὴς τὸν Θεόν; εἶπε γὰρ ἐν καρδίᾳ αὐτοῦ, οὐ ζητήσει.
\vs{35}Βλέπεις, ὅτι σὺ πόνον καὶ θυμὸν κατανοεῖς, τοῦ παραδοῦναι αὐτοὺς εἰς χεῖράς σου· σοὶ ἐγκαταλέλειπται ὁ πτωχὸς, ὀρφανῷ σὺ ἦσθα βοηθός.
\vs{36}Σύντριψον τὸν βραχίονα τοῦ ἁμαρτωλοῦ καὶ πονηροῦ, ζητηθήσεται ἡ ἁμαρτία αὐτοῦ καὶ οὐ μὴ εὑρεθῇ.

\vs{37}Βασιλεύσει Κύριος εἰς τὸν αἰῶνα, καὶ εἰς τὸν αἰῶνα τοῦ αἰῶνος, ἀπολεῖσθε ἔθνη ἐκ τῆς γῆς αὐτοῦ.
\vs{38}Τὴν ἐπιθυμίαν τῶν πενήτων εἰσήκουσε Κύριος, τὴν ἑτοιμασίαν τῆς καρδίας αὐτῶν προσέσχε τὸ οὖς σου·
\vs{39}Κρῖναι ὀρφανῷ καὶ ταπεινῷ, ἵνα μὴ προσθῇ ἔτι μεγαλαυχεῖν ἄνθρωπος ἐπὶ τῆς γῆς.

\begin{psalmheading}{\ch{10}{11} Εἰς τὸ τέλος, ψαλμὸς τῷ Δαυὶδ.}
\end{psalmheading}
Ἐπὶ τῷ Κυρίῳ πέποιθα· πῶς ἐρεῖτε τῇ ψυχῇ μου, μεταναστεύου ἐπὶ τὰ ὄρη ὡς στρουθίον;
\vs{2}Ὅτι ἰδοὺ οἱ ἁμαρτωλοὶ ἐνέτειναν τόξον, ἡτοίμασαν βέλη εἰς φαρέτραν, τοῦ κατατοξεῦσαι ἐν σκοτομήνῃ τοὺς εὐθεῖς τῇ καρδίᾳ.
\vs{3}Ὅτι ἃ κατηρτίσω καθεῖλον, ὁ δὲ δίκαιος τί ἐποίησε;

\vs{4}Κύριος ἐν ναῷ ἁγίῳ αὐτοῦ, Κύριος, ἐν οὐρανῷ ὁ θρόνος αὐτοῦ· οἱ ὀφθαλμοὶ αὐτοῦ εἰς τὸν πένητα ἀποβλέπουσι, τὰ βλέφαρα αὐτοῦ ἐξετάζει τοὺς υἱοὺς τῶν ἀνθρώπων·
\vs{5}Κύριος ἐξετάζει τὸν δίκαιον καὶ τὸν ἀσεβῆ, ὁ δὲ ἀγαπῶν ἀδικίαν μισεῖ τὴν ἑαυτοῦ ψυχήν.
\vs{6}Ἐπιβρέξει ἐπὶ ἁμαρτωλοὺς παγίδας, πῦρ καὶ θεῖον καὶ πνεῦμα καταιγίδος ἡ μερὶς τοῦ ποτηρίου αὐτῶν.
\vs{7}Ὅτι δίκαιος Κύριος καὶ δικαιοσύνας ἠγάπησεν, εὐθύτητα εἶδε τὸ πρόσωπον αὐτοῦ.

\begin{psalmheading}{\ch{11}{12} Εἰς τὸ τέλος, ὑπὲρ τῆς ὀγδόης, ψαλμὸς τῷ Δαυίδ.}
\end{psalmheading}
\vs{2}Σῶσον με Κύριε, ὅτι ἐκλέλοιπεν ὅσιος, ὅτι ὠλιγώθησαν αἱ ἀλήθειαι ἀπὸ τῶν υἱῶν τῶν ἀνθρώπων.
\vs{3}Μάταια ἐλάλησεν ἕκαστος πρὸς τὸν πλησίον αὐτοῦ, χείλη δόλια, ἐν καρδίᾳ καὶ ἐν καρδίᾳ ἐλάλησαν.
\vs{4}Ἐξολοθρεύσαι Κύριος πάντα τὰ χείλη τὰ δόλια, καὶ γλῶσσαν μεγαλοῤῥήμονα·
\vs{5}Τοὺς εἰπόντας, τὴν γλῶσσαν ἡμῶν μεγαλυνοῦμεν, τὰ χείλη ἡμῶν παρʼ ἡμῶν ἐστι· τίς ἡμῶν Κύριός ἐστιν;

\vs{6}Ἀπὸ τῆς ταλαιπωρίας τῶν πτωχῶν, καὶ ἀπὸ τοῦ στεναγμοῦ τῶν πενήτων, νῦν ἀναστήσομαι, λέγει Κύριος· θήσομαι ἐν σωτηρίῳ, παῤῥησιάσομαι ἐν αὐτῷ.
\vs{7}Τὰ λόγια Κυρίου, λόγια ἁγνά· ἀργύριον πεπυρωμένον, δοκίμιον τῇ γῇ, κεκαθαρισμένον ἑπταπλασίως.
\vs{8}Σὺ Κύριε φυλάξεις ἡμᾶς· καὶ διατηρήσεις ἡμᾶς ἀπὸ τῆς γενεᾶς ταύτης, καὶ εἰς τὸν αἰῶνα.
\vs{9}Κύκλῳ οἱ ἀσεβεῖς περιπατοῦσι, κατὰ τὸ ὕψος σου ἐπολυώρησας τοὺς υἱοὺς τῶν ἀνθρώπων.

\begin{psalmheading}{\ch{12}{13} Εἰς τὸ τέλος, ψαλμὸς τῷ Δαυίδ.}
\end{psalmheading}
Ἕως πότε Κύριε ἐπιλήσῃ μου, εἰς τέλος; ἕως πότε ἀποστρέψεις τὸ πρόσωπόν σου ἀπʼ ἐμοῦ;
\vs{2}Ἕως τίνος θήσομαι βουλὰς ἐν ψυχῇ μου, ὀδύνας ἐν καρδίᾳ μου ἡμέρας; ἕως πότε ὑψωθήσεται ὁ ἐχθρός μου ἐπʼ ἐμέ;
\vs{3}Ἐπίβλεψον, εἰσάκουσόν μου, Κύριε ὁ Θεός μου· φώτισον τοὺς ὀφθαλμούς μου, μή ποτε ὑπνώσω εἰς θάνατον·
\vs{4}μή ποτε εἴποι ὁ ἐχθρός μου, ἴσχυσα πρὸς αὐτόν· οἱ θλίβοντές με ἀγαλλιάσονται ἐὰν σαλευθῶ.

\vs{5}Ἐγὼ δὲ ἐπὶ τῷ ἐλέει σου ἤλπισα· ἀγαλλιάσεται ἡ καρδία μου ἐν τῷ σωτηρίῳ σου.
\vs{6}Ἄσω τῷ Κυρίῳ τῷ εὐεργετήσαντί με, καὶ ψαλῶ τῷ ὀνόματι Κυρίου τοῦ ὑψίστου.

\begin{psalmheading}{\ch{13}{14} Εἰς τὸ τέλος, ψαλμὸς τῷ Δαυίδ.}
\end{psalmheading}
Εἶπεν ἄφρων ἐν καρδίᾳ αὐτοῦ, οὐκ ἔστι Θεός· διέφθειραν καὶ ἐβδελύχθησαν ἐν ἐπιτηδεύμασιν, οὐκ ἔστι ποιῶν χρηστότητα, οὐκ ἔστιν ἕως ἑνός.
\vs{2}Κύριος ἐκ τοῦ οὐρανοῦ διέκυψεν ἐπὶ τοὺς υἱοὺς τῶν ἀνθρώπων, τοῦ ἰδεῖν εἰ ἔστι συνιὼν ἢ ἐκζητῶν τὸν Θεόν.
\vs{3}Πάντες ἐξέκλιναν, ἅμα ἠχρειώθησαν, οὐκ ἔστι ποιῶν χρηστότητα, οὐκ ἔστιν ἕως ἑνός· τάφος ἀνεῳγμένος ὁ λάρυγξ αὐτῶν, ταῖς γλώσσαις αὐτῶν ἐδολιοῦσαν, ἰὸς ἀσπίδων ὑπὸ τὰ χείλη αὐτῶν· ὧν τὸ στόμα ἀρᾶς καὶ πικρίας γέμει, ὀξεῖς οἱ πόδες αὐτῶν ἐκχέαι αἷμα· σύντριμμα καὶ ταλαιπωρία ἐν ταῖς ὁδοῖς αὐτῶν, καὶ ὁδὸν εἰρήνης οὐκ ἔγνωσαν· οὐκ ἔστι φόβος Θεοῦ ἀπέναντι τῶν ὀφθαλμῶν αὐτῶν.

\vs{4}Οὐχὶ γνώσονται πάντες οἱ ἐργαζόμενοι τὴν ἀνομίαν, οἱ κατέσθοντες τὸν λαόν μου βρώσει ἄρτου; τὸν Κύριον οὐκ ἐπεκαλέσαντο.
\vs{5}Ἐκεῖ ἐδειλίασαν φόβῳ, οὗ οὐκ ἦν φόβος, ὅτι ὁ Θεὸς ἐν γενεᾷ δικαίᾳ.
\vs{6}Βουλὴν πτωχοῦ κατῃσχύνατε, ὅτι Κύριος ἐλπὶς αὐτοῦ ἐστι.
\vs{7}Τίς δώσει ἐκ Σιὼν τὸ σωτήριον τοῦ Ἰσραήλ; ἐν τῷ ἐπιστρέψαι Κύριον τὴν αἰχμαλωσίαν τοῦ λαοῦ αὐτοῦ, ἀγαλλιάσθω Ἰακὼβ, καὶ εὐφρανθήτω Ἰσραήλ.

\begin{psalmheading}{\ch{14}{15} Ψαλμὸς τῷ Δαυίδ.}
\end{psalmheading}
Κύριε, τίς παροικήσει ἐν τῷ σκηνώματί σου; καὶ τίς κατασκηνώσει ἐν τῷ ὄρει τῷ ἁγίῳ σου;

\vs{2}Πορευόμενος ἄμωμος, καὶ ἐργαζόμενος δικαιοσύνην· λαλῶν ἀλήθειαν ἐν καρδίᾳ αὐτοῦ·
\vs{3}Ὃς οὐκ ἐδόλωσεν ἐν γλώσσῃ αὐτοῦ, οὐδὲ ἐποίησε τῷ πλησίον αὐτοῦ κακὸν, καὶ ὀνειδισμὸν οὐκ ἔλαβεν ἐπὶ τοὺς ἔγγιστα αὐτοῦ·
\vs{4}Ἐξουδένωται ἐνώπιον αὐτοῦ πονηρευόμενος, τοὺς δὲ φοβουμένους Κύριον δοξάζει· ὁ ὀμνύων τῷ πλησίον αὐτοῦ καὶ οὐκ ἐθετῶν·
\vs{5}Τὸ ἀργύριον αὐτοῦ οὐκ ἔδωκεν ἐπὶ τόκῳ, καὶ δῶρα ἐπʼ ἀθώοις οὐκ ἔλαβεν· ὁ ποιῶν ταῦτα, οὐ σαλευθήσεται εἰς τὸν αἰῶνα.

\begin{psalmheading}{\ch{15}{16} Στηλογραφία τῷ Δαυίδ.}
\end{psalmheading}
Φύλαξον με Κύριε, ὅτι ἐπὶ σοὶ ἤλπισα.
\vs{2}Εἶπα τῷ Κυρίῳ, Κύριός μου εἶ σὺ, ὅτι τῶν ἀγαθῶν μου οὐ χρείαν ἔχεις.
\vs{3}Τοῖς ἁγίοις τοῖς ἐν τῇ γῇ αὐτοῦ, ἐθαυμάστωσε πάντα τὰ θελήματα αὐτοῦ ἐν αὐτοῖς.
\vs{4}Ἐπληθύνθησαν αἱ ἀσθένειαι αὐτῶν, μετὰ ταῦτα ἐτάχυναν· οὐ μὴ συναγάγω τὰς συναγωγὰς αὐτῶν ἐξ αἱμάτων, οὐδὲ μὴ μνησθῶ τῶν ὀνομάτων αὐτῶν διὰ χειλέων μου.
\vs{5}Κύριος μερὶς τῆς κληρονομίας μου καὶ τοῦ ποτηρίου μου, σὺ εἶ ὁ ἀποκαθιστῶν τὴν κληρονομίαν μου ἐμοί.
\vs{6}Σχοινία ἐπέπεσάν μοι ἐν τοῖς κρατίστοις, καὶ γὰρ ἡ κληρονομία μου κρατίστη μοι ἐστίν.

\vs{7}Εὐλογήσω τὸν Κύριον τὸν συνετίσαντά με, ἔτι δὲ καὶ ἕως νυκτὸς ἐπαίδευσάν με οἱ νεφροί μου.
\vs{8}Προωρώμην τὸν Κύριον ἐνώπιόν μου διαπαντὸς, ὅτι ἐκ δεξιῶν μου ἐστὶν ἵνα μὴ σαλευθῶ.
\vs{9}Διὰ τοῦτο ηὐφράνθη ἡ καρδία μου, καὶ ἠγαλλιάσατο ἡ γλῶσσά μου, ἔτι δὲ καὶ ἡ σάρξ μου κατασκηνώσει ἐπʼ ἐλπίδι·
\vs{10}Ὅτι οὐκ ἐγκαταλείψεις τὴν ψυχήν μου εἰς ᾅδην, οὐδὲ δώσεις τὸν ὅσιόν σου ἰδεῖν διαφθοράν.
\vs{11}Ἐγνώρισάς μοι ὁδοὺς ζωῆς, πληρώσεις με εὐφροσύνης μετὰ τοῦ προσώπου σου, τερπνότητες ἐν τῇ δεξιᾷ σου εἰς τέλος.

\begin{psalmheading}{\ch{16}{17} Προσευχὴ τοῦ Δαυίδ.}
\end{psalmheading}
Εἰσάκουσον Κύριε τῆς δικαιοσύνης μου, πρόσχες τῇ δεήσει μου· ἐνώτισαι τὴν προσευχήν μου οὐκ ἐν χείλεσι δολίοις.
\vs{2}Ἐκ προσώπου σου τὸ κρίμα μου ἐξέλθοι, οἱ ὀφθαλμοί μου ἰδέτωσαν εὐθύτητας.
\vs{3}Ἐδοκίμασας τὴν καρδίαν μου, ἐπεσκέψω νυκτὸς, ἐπύρωσάς με, καὶ οὐχ εὑρέθη ἐν ἐμοὶ ἀδικία·
\vs{4}ὅπως ἂν μὴ λαλήσῃ τὸ στόμα μου. Τὰ ἔργα τῶν ἀνθρώπων, διὰ τοὺς λόγους τῶν χειλέων σου ἐγὼ ἐφύλαξα ὁδοὺς σκληράς.
\vs{5}Κατάρτισαι τὰ διαβήματά μου ἐν ταῖς τρίβοις σου, ἵνα μὴ σαλευθῇ τὰ διαβήματά μου.

\vs{6}Ἐγὼ ἐκεκραξα, ὅτι ἐπήκουσας μου ὁ Θεός· κλῖνον τὸ οὖς σου ἐμοὶ, καὶ εἰσάκουσον τῶν ῥημάτων μου.
\vs{7}Θαυμάστωσον τὰ ἐλέη σου, ὁ σώζων τοὺς ἐλπίζοντας ἐπὶ σέ· ἐκ τῶν ἀνθεστηκότων τῇ δεξιᾷ σου,
\vs{8}φύλαξόν με ὡς κόρην ὀφθαλμοῦ· ἐν σκέπῃ τῶν πτερύγων σου σκεπάσεις με,
\vs{9}ἀπὸ προσώπου ἀσεβῶν τῶν ταλαιπωρησάντων με· οἱ ἐχθροί μου τὴν ψυχήν μου περιέσχον.
\vs{10}Τὸ στέαρ αὐτῶν συνέκλεισαν, τὸ στόμα αὐτῶν ἐλάλησεν ὑπερηφανίαν.
\vs{11}Ἐκβαλόντες με νυνὶ περιεκύκλωσάν με, τοὺς ὀφθαλμοὺς αὐτὼν ἔθεντο ἐκκλῖναι ἐν τῇ γῇ.
\vs{12}Ὑπέλαβόν με ὡσεὶ λέων ἕτοιμος εἰς θήραν, καὶ ὡσεὶ σκύμνος οἰκῶν ἐν ἀποκρύφοις.
\vs{13}Ἀνάστηθι Κύριε, πρόφθασον αὐτοὺς, καὶ ὑποσκέλισον αὐτοὺς, ῥῦσαι τὴν ψυχήν μου ἀπὸ ἀσεβοῦς, ῥομφαίαν σου
\vs{14}ἀπὸ ἐχθρῶν τῆς χειρός σου· Κύριε ἀπολύων ἀπὸ γῆς, διαμέρισον αὐτοὺς ἐν τῇ ζωῇ αὐτῶν, καὶ τῶν κεκρυμμένων σου ἐπλήσθη ἡ γαστὴρ αὐτῶν· ἐχορτάσθησαν ὑείων, καὶ ἀφῆκαν τὰ κατάλοιπα τοῖς νηπίοις αὐτῶν.

\vs{15}Ἐγὼ δὲ ἐν δικαιοσύνῃ ὀφθήσομαι τῷ προσώπῳ σου, χορτασθήσομαι ἐν τῷ ὀφθῆναι τὴν δόξαν σου.

\begin{psalmheading}{\ch{17}{18} Εἰς τὸ τέλος τῷ παιδὶ Κυρίου τῷ Δαυὶδ, ἃ ἐλάλησε τῷ Κυρίῳ, τοὺς λόγους τῆς ᾠδῆς ταύτης, ἐν ἡμέρᾳ ᾗ ἐῤῥύσατο αὐτὸν Κύριος ἐκ χειρὸς πάντων τῶν ἐχθρῶν αὐτοῦ, καὶ ἐκ χειρὸς Σαοὺλ, καὶ εἶπεν,}
\end{psalmheading}
\vs{2}Αγαπήσω σε, Κύριε ἰσχύς μου.
\vs{3}Κύριος στερέωμά μου, καὶ καταφυγή μου, καὶ ῥύστης μου· ὁ Θεός μου βοηθός μου, ἐλπιῶ ἐπʼ αὐτόν· ὑπερασπιστής μου, καὶ κέρας σωτηρίας μου, καὶ ἀντιλήπτωρ μου.
\vs{4}Αἰνῶν ἐπικαλέσομαι Κύριον, καὶ ἐκ τῶν ἐχθρῶν μου σωθήσομαι.
\vs{5}Περιέσχον με ὠδῖνες θανάτου, καὶ χείμαῤῥοι ἀνομίας ἐξετάραξάν με.
\vs{6}Ὠδῖνες ᾅδου περιεκύκλωσάν με, προέφθασάν με παγίδες θανάτου.

\vs{7}Καὶ ἐν τῷ θλίβεσθαί με ἐπεκαλεσάμην τὸν Κύριον, καὶ πρὸς τὸν Θεόν μου ἐκέκραξα· ἤκουσεν ἐκ ναοῦ ἁγίου αὐτοῦ φωνῆς μου. καὶ ἡ κραυγή μου ἐνώπιον αὐτοῦ εἰσελεύσεται εἰς τὰ ὦτα αὐτοῦ.

\vs{8}Καὶ ἐσαλεύθη, καὶ ἔντρόμος ἐγενήθη ἡ γῆ, καὶ τὰ θεμέλια τῶν ὀρέων ἐταράχθησαν, καὶ ἐσαλεύθησαν, ὅτι ὠργίσθη αὐτοῖς ὁ Θεός.
\vs{9}Ἀνέβη καπνὸς ἐν ὀργῇ αὐτοῦ, καὶ πῦρ ἀπὸ προσώπου αὐτοῦ κατεφλόγισεν, ἄνθρακες ἀνήφθησαν ἀπʼ αὐτοῦ.
\vs{10}Καὶ ἔκλινεν οὐρανὸν καὶ κατέβη, καὶ γνόφος ὑπὸ τοὺς πόδας αὐτοῦ.
\vs{11}Καὶ ἐπέβη ἐπὶ χερουβὶμ καὶ ἐπετάσθη, ἐπετάσθη ἐπὶ πτερύγων ἀνέμων.
\vs{12}Καὶ ἔθετο σκότος ἀποκρυφὴν αὐτοῦ, κύκλῳ αὐτοῦ ἡ σκηνὴ αὐτοῦ, σκοτεινὸν ὕδωρ ἐν νεφέλαις ἀέρων.
\vs{13}Ἀπὸ τῆς τηλαυγήσεως ἐνώπιον αὐτοῦ αἱ νεφέλαι διῆλθον, χάλαζα καὶ ἄνθρακες πυρός.
\vs{14}Καὶ ἐβρόντησεν ἐξ οὐρανοῦ Κύριος, καὶ ὁ ὕψιστος ἔδωκε φωνὴν αὐτοῦ.
\vs{15}Καὶ ἐξαπέστειλε βέλη καὶ ἐσκόρπισεν αὐτοὺς, καὶ ἀστραπὰς ἐπλήθυνε καὶ συνετάραξεν αὐτούς.
\vs{16}Καὶ ὤφθησαν αἱ πηγαὶ τῶν ὑδάτων, καὶ ἀνεκαλύφθη τὰ θεμέλια τῆς οἰκουμένης· ἀπὸ ἐπιτιμήσεώς σου Κύριε, ἀπὸ ἐνπνεύσεως πνεύματος ὀργῆς σου.

\vs{17}Ἐξαπέστειλεν ἐξ ὕψους καὶ ἔλαβέ με, προσελάβετό με ἐξ ὑδάτων πολλῶν.
\vs{18}Ῥύσεταί με ἐξ ἐχθρῶν μου δυνατῶν, καὶ ἐκ τῶν μισούντων με, ὅτι ἐστερεώθησαν ὑπὲρ ἐμέ.
\vs{19}Προέφθασάν με ἐν ἡμέρᾳ κακώσεώς μου, καὶ ἐγένετο Κύριος ἀντιστήριγμά μου.
\vs{20}Καὶ ἐξήγαγέ με εἰς πλατυσμὸν, ῥύσεταί με, ὅτι ἠθέλησέ με.
\vs{21}Καὶ ἀνταποδώσει μοι Κύριος κατὰ τὴν δικαιοσύνην μου, καὶ κατὰ τὴν καθαριότητα τῶν χειρῶν μου ἀνταποδώσει μοι.
\vs{22}Ὅτι ἐφύλαξα τὰς ὁδοὺς Κυρίου, καὶ οὐκ ἠσέβησα ἀπὸ τοῦ Θεοῦ μου.
\vs{23}Ὅτι πάντα τὰ κρίματα αὐτοῦ ἐνώπιόν μου, καὶ τὰ δικαιώματα αὐτοῦ οὐκ ἀπέστησαν ἀπʼ ἐμοῦ.
\vs{24}Καὶ ἔσομαι ἄμωμος μετʼ αὐτοῦ, καὶ φυλάξομαι ἀπὸ τῆς ἀνομίας μου.
\vs{25}Καὶ ἀνταποδώσει μοι Κύριος κατὰ τὴν δικαιοσύνην μου, καὶ κατὰ τὴν καθαριότητα τῶν χειρῶν μου ἐνώπιον τῶν ὀφθαλμῶν αὐτοῦ.

\vs{26}Μετὰ ὁσίου ὁσιωθήσῃ, καὶ μετὰ ἀνδρὸς ἀθώου ἀθῶος ἔσῃ·
\vs{27}Καὶ μετὰ ἐκλεκτοῦ ἐκλεκτὸς ἔσῃ, καὶ μετὰ στρεβλοῦ διαστρέψεις.
\vs{28}Ὅτι σὺ λαὸν ταπεινὸν σώσεις, καὶ ὀφθαλμοὺς ὑπερηφάνων ταπεινώσεις.
\vs{29}Ὅτι σὺ φωτιεῖς λύχνον μου Κύριε, ὁ Θεός μου φωτιεῖς τὸ σκότος μου.
\vs{30}Ὅτι ἐν σοὶ ῥυσθήσομαι ἀπὸ πειρατηρίου, καὶ ἐν τῷ Θεῷ μου ὑπερβήσομαι τεῖχος.
\vs{31}Ὁ Θεός μου, ἄμωμος ἡ ὁδὸς αὐτοῦ, τὰ λόγια Κυρίου πεπυρωμένα, ὑπερασπιστής ἐστι πάντων τῶν ἐλπιζόντων ἐπʼ αὐτόν.
\vs{32}Ὅτι τίς Θεὸς πλὴν τοῦ Κυρίου; καὶ τίς Θεὸς πλὴν τοῦ Θεοῦ ἡμῶν;

\vs{33}Ὁ Θεὸς ὁ περιζωννύων με δύναμιν, καὶ ἔθετο ἄμωμον τὴν ὁδόν μου·
\vs{34}ὁ καταρτιζόμενος τοὺς πόδας μου ὡσεὶ ἐλάφου, καὶ ἐπὶ τὰ ὑψηλὰ ἱστῶν με·
\vs{35}Διδάσκων χεῖράς μου εἰς πόλεμον· καὶ ἔθου τόξον χαλκοῦν τοὺς βραχίονάς μου,
\vs{36}καὶ ἔδωκάς με ὑπερασπισμὸν σωτηρίας μου· καὶ ἡ δεξιά σου ἀντελάβετό μου, καὶ ἡ παιδεία σου ἀνώρθωσέ με εἰς τέλος, καὶ ἡ παιδεία σου αὐτή με διδάξει.
\vs{37}Ἐπλάτυνας τὰ διαβήματά μου ὑποκάτω μου, καὶ οὐκ ἠσθένησαν τὰ ἴχνη μου.
\vs{38}Καταδιώξω τοὺς ἐχθρούς μου, καὶ καταλήψομαι αὐτοὺς, καὶ οὐκ ἀποστραφήσομαι, ἕως ἂν ἐκλείπωσιν.
\vs{39}Ἐκθλίψω αὐτοὺς, καὶ οὐ μὴ δύνωνται στῆναι, πεσοῦνται ὑπὸ τοὺς πόδας μου.
\vs{40}Καὶ περιέζωσάς με δύναμιν εἰς πόλεμον, συνεπόδισας πάντας τοὺς ἐπανισταένους ἐπʼ ἐμὲ ὑποκάτω μου.
\vs{41}Καὶ τοὺς ἐχθρούς μου ἔδωκάς μοι νῶτον, καὶ τοὺς μισοῦντάς με ἐξωλόθρευσας.
\vs{42}Ἐκέκραξαν, καὶ οὐκ ἦν ὁ σώζων· πρὸς Κύριον, καὶ οὐκ εἰσήκουεν αὐτῶν.
\vs{43}Καὶ λεπτυνῶ αὐτοὺς ὡς χοῦν κατὰ πρόσωπον ἀνέμου, ὡς πηλὸν πλατειῶν λεανῶ αὐτούς.
\vs{44}Ῥῦσαί με ἐξ ἀντιλογιῶν λαοῦ, καταστήσεις με εἰς κεφαλὴν ἐθνῶν· λαὸς ὃν οὐκ ἔγνων, ἐδούλευσέ μοι,
\vs{45}εἰς ἀκοὴν ὠτίου ὑπήκουσέ μοι· υἱοὶ ἀλλότριοι ἐψεύσαντό μοι,
\vs{46}υἱοὶ ἀλλότριοι ἐπαλαιώθησαν, καὶ ἐχώλαναν ἀπὸ τῶν τρίβων αὐτῶν.

\vs{47}Ζῇ Κύριος, καὶ εὐλογητὸς ὁ Θεός μου, καὶ ὑψωθήτω ὁ Θεὸς τῆς σωτηρίας μου.
\vs{48}Ὁ Θεὸς ὁ διδοὺς ἐκδικήσεις ἐμοὶ, καὶ ὑποτάξας λαοὺς ὑπʼ ἐμὲ,
\vs{49}ὁ ῥύστης μου ἐξ ἐχθρῶν ὀργίλων· ἀπὸ τῶν ἐπανισταμένων ἐπʼ ἐμὲ ὑψώσεις με, ἀπὸ ἀνδρὸς ἀδίκου ῥύσῃ με.
\vs{50}Διὰ τοῦτο ἐξομολογήσομαί σοι ἐν ἔθνεσι, Κύριε, καὶ τῷ ὀνόματί σου ψαλῶ.
\vs{51}Μεγαλύνων τὰς σωτηρίας τοῦ βασιλέως αὐτοῦ, καὶ ποιῶν ἔλεος τῷ χριστῷ αὐτοῦ τῷ Δαυὶδ, καὶ τῷ σπέρματι αὐτοῦ ἕως αἰῶνος.

\begin{psalmheading}{\ch{18}{19} Εἰς τὸ τέλος, ψαλμὸς τῷ Δαυίδ.}
\end{psalmheading}
\vs{2}Οἱ οὐρανοὶ διηγοῦνται δόξαν Θεοῦ, ποίησιν δὲ χειρῶν αὐτοῦ ἀναγγέλλει τὸ στερέωμα.
\vs{3}Ἡμέρα τῇ ἡμέρᾳ ἐρεύγεται ῥῆμα, καὶ νὺξ νυκτὶ ἀναγγέλλει γνῶσιν.
\vs{4}Οὐκ εἰσὶ λαλιαὶ οὐδὲ λόγοι, ὧν οὐχὶ ἀκούονται αἱ φωναὶ αὐτῶν·
\vs{5}Εἰς πᾶσαν τὴν γῆν ἐξῆλθεν ὁ φθόγγος αὐτῶν, καὶ εἰς τὰ πέρατα τῆς οἰκουμένης τὰ ῥήματα αὐτῶν·
\vs{6}ἐν τῷ ἡλίῳ ἔθετο τὸ σκήνωμα αὐτοῦ, καὶ αὐτὸς ὡς νυμφίος ἐκπορευόμενος ἐκ παστοῦ αὐτοῦ· ἀγαλλιάσεται ὡς γίγας δραμεῖν ὁδὸν αὐτοῦ.
\vs{7}Ἀπʼ ἄκρου τοῦ οὐρανοῦ ἡ ἔξοδος αὐτοῦ· καὶ τὸ κατάντημα αὐτοῦ ἕως ἄκρου τοῦ οὐρανοῦ· καὶ οὐκ ἔστιν ὃς ἀποκρυβήσεται τὴν θέρμην αὐτοῦ.

\vs{8}Ὁ νόμος τοῦ Κυρίου ἄμωμος ἐπιστρέφων ψυχὰς, ἡ μαρτυρία Κυρίου πιστὴ σοφίζουσα νήπια.
\vs{9}Τὰ δικαιώματα Κυρίου εὐθέα εὐφραίνοντα καρδίαν, ἡ ἐντολὴ Κυρίου τηλαυγὴς φωτίζουσα ὀφθαλμούς.
\vs{10}Ὁ φόβος Κυρίου ἁγνὸς διαμένων εἰς αἰῶνα αἰῶνος, τὰ κρίματα Κυρίου ἀληθινὰ δεδικαιωμένα ἐπὶ τὸ αὐτό·
\vs{11}Ἐπιθυμητὰ ὑπὲρ χρυσίον καὶ λίθον τίμιον πολὺν, καὶ γλυκύτερα ὑπὲρ μέλι καὶ κηρίον.
\vs{12}Καὶ γὰρ ὁ δοῦλός σου φυλάσσει αὐτὰ, ἐν τῷ φυλάσσειν αὐτὰ ἀνταπόδοσις πολλή.

\vs{13}Παραπτώματα τίς συνήσει; ἐκ τῶν κρυφίων μου καθάρισόν με,
\vs{14}καὶ ἀπὸ ἀλλοτρίων φεῖσαι τοῦ δούλου σου· ἐὰν μή μου κατακυριεύσωσι, τότε ἄμωμος ἔσομαι, καὶ καθαρισθήσομαι ἀπὸ ἁμαρτίας μεγάλης.
\vs{15}Καὶ ἔσονται εἰς εὐδοκίαν τὸ λόγια τοῦ στόματός μου, καὶ ἡ μελέτη τῆς καρδίας μου ἐνώπιόν σου διαπαντός· Κύριε βοηθέ μου, καὶ λυτρωτά μου.

\begin{psalmheading}{\ch{19}{20} Εἰς τὸ τέλος, ψαλμὸς τῷ Δαυίδ.}
\end{psalmheading}
\vs{2}Ἐπακούσαι σου Κύριος ἐν ἡμέρᾳ θλίψεως, ὑπερασπίσαι σου τὸ ὄνομα τοῦ Θεοῦ Ἰακώβ.
\vs{3}Ἐξαποστείλαι σοι βοήθειαν ἐξ ἁγίου, καὶ ἐκ Σιὼν ἀντιλάβοιτό σου.
\vs{4}Μνησθείη πάσης θυσίας σου, καὶ τὸ ὁλοκαύτωμά σου πιανάτω· διάψαλμα.
\vs{5}Δῴη σοι κατὰ τὴν καρδίαν σου, καὶ πᾶσαν τὴν βουλήν σου πληρώσαι.
\vs{6}Ἀγαλλιασόμεθα ἐν τῷ σωτηρίῳ σου, καὶ ἐν ὀνόματι Θεοῦ ἡμῶν μεγαλυνθησόμεθα· πληρώσαι Κύριος πάντα τὰ αἰτήματά σου.

\vs{7}Νῦν ἔγνων ὅτι ἔσωσε Κύριος τὸν χριστὸν αὐτοῦ· ἐπακούσεται αὐτοῦ ἐξ οὐρανοῦ ἁγίου αὐτοῦ, ἐν δυναστείαις ἡ σωτηρία τῆς δεξιᾶς αὐτοῦ.
\vs{8}Οὗτοι ἐν ἅρμασι καὶ οὗτοι ἐν ἵπποις, ἡμεῖς δὲ ἐν ὀνόματι Κυρίου Θεοῦ ἡμῶν μεγαλυνθησόμεθα.
\vs{9}Αὐτοὶ συνεποδίσθησαν καὶ ἔπεσαν, ἡμεῖς δὲ ἀνέστημεν καὶ ἀνωρθώθημεν.
\vs{10}Κύριε σῶσον τὸν βασιλέα καὶ ἐπάκουσον ἡμῶν, ἐν ᾗ ἂν ἡμέρᾳ ἐπικαλεσώμεθά σε.

\begin{psalmheading}{\ch{20}{21} Εἰς τὸ τέλος, ψαλμὸς τῷ Δαυίδ.}
\end{psalmheading}
\vs{2}Κύριε, ἐν τῇ δυνάμει σου εὐφρανθήσεται ὁ βασιλεὺς, καὶ ἐπὶ τῷ σωτηρίῳ σου ἀγαλλιάσεται σφόδρα.
\vs{3}Τὴν ἐπιθυμίαν τῆς ψυχῆς αὐτοῦ ἔδωκας αὐτῷ, καὶ τὴν δέησιν τῶν χειλέων αὐτοῦ οὐκ ἐστέρησας αὐτόν· διάψαλμα.
\vs{4}Ὅτι προέφθασας αὐτὸν ἐν εὐλογίαις χρηστότητος, ἔθηκας ἐπὶ τὴν κεφαλὴν αὐτοῦ στέφανον ἐκ λίθου τιμίου.
\vs{5}Ζωὴν ᾐτήσατό σε, καὶ ἔδωκας αὐτῷ μακρότητα ἡμερῶν εἰς αἰῶνα αἰῶνος.
\vs{6}Μεγάλη ἡ δόξα αὐτοῦ ἐν τῷ σωτηρίῳ σου, δόξαν καὶ μεγαλοπρέπειαν ἐπιθήσεις ἐπʼ αὐτόν.
\vs{7}Ὅτι δώσεις αὐτῷ εὐλογίαν εἰς αἰῶνα αἰῶνος, εὐφρανεῖς αὐτὸν ἐν χαρᾷ μετὰ τοῦ προσώπου σου.
\vs{8}Ὅτι ὁ βασιλεὺς ἐλπίζει ἐπὶ Κύριον, καὶ ἐν τῷ ἐλέει τοῦ ὑψίστου οὐ μὴ σαλευθῇ.

\vs{9}Εὑρεθείη ἡ χείρ σου πᾶσι τοῖς ἐχθροῖς σου, ἡ δεξιά σου εὕροι πάντας τοὺς μισοῦντάς σε.
\vs{10}Θήσεις αὐτοὺς ὡς κλίβανον πυρὸς εἰς καιρὸν τοῦ προσώπου σου, Κύριος ἐν ὀργῇ αὐτοῦ συνταράξει αὐτοὺς, καὶ καταφάγεται αὐτοὺς πῦρ.
\vs{11}Τὸν καρπὸν αὐτῶν ἀπὸ γῆς ἀπολεῖς, καὶ τὸ σπέρμα αὐτῶν ἀπὸ υἱῶν ἀνθρώπων.
\vs{12}Ὅτι ἔκλιναν εἰς σὲ κακὰ, διελογίσαντο βουλὴν, ἣν οὐ μὴ δύνωνται στῆσαι.
\vs{13}Ὅτι θήσεις αὐτοὺς νῶτον ἐν τοῖς περιλοίποις σου, ἑτοιμάσεις τὸ πρόσωπον αὐτῶν.
\vs{14}Ὑψώθητι Κύριε ἐν τῇ δυνάμει σου· ᾄσομεν καὶ ψαλοῦμεν τὰς δυναστείας σου.

\begin{psalmheading}{\ch{21}{22} Εἰς τὸ τέλος, ὑπὲρ τῆς ἀντιλήψεως τῆς ἑωθινῆς, ψαλμὸς τῷ Δαυίδ.}
\end{psalmheading}
\vs{2}Ὁ Θεός ὁ Θεός μου, πρόσχες μοι, ἱνατί ἐγκατέλιπές με; μακρὰν ἀπὸ τῆς σωτηρίας μου οἱ λόγοι τῶν παραπτωμάτων μου.
\vs{3}Ο Θεός μου, κεκράξομαι ἡμέρας πρὸς σὲ καὶ οὐκ εἰσακούσῃ, καὶ νυκτὸς καὶ οὐκ εἰς ἄνοιαν ἐμοί.

\vs{4}Σὺ δὲ ἐν ἁγίῳ κατοικεῖς, ὁ ἔπαινος τοῦ Ἰσραήλ.
\vs{5}Ἐπὶ σοὶ ἤλπισαν οἱ πατέρες ἡμῶν, ἤλπισαν καὶ ἐῤῥύσω αὐτούς.
\vs{6}Πρὸς σὲ ἐκέκραξαν καὶ ἐσώθησαν, ἐπὶ σοὶ ἤλπισαν καὶ οὐ κατῃσχύνθησαν.
\vs{7}Ἐγὼ δέ εἰμι σκώληξ καὶ οὐκ ἄνθρωπος, ὄνειδος ἀνθρώπων καὶ ἐξουδένημα λαοῦ.
\vs{8}Πάντες οἱ θεωροῦντές με ἐξεμυκτήρισάν με, ἐλάλησαν ἐν χείλεσιν, ἐκίνησαν κεφαλὴν,
\vs{9}ἤλπισεν ἐπὶ Κύριον, ῥυσάσθω αὐτὸν, σωσάτω αὐτὸν, ὅτι θέλει αὐτόν.
\vs{10}Ὅτι σὺ εἶ ὁ ἐκσπάσας με ἐκ γαστρὸς, ἡ ἐλπίς μου ἀπὸ μαστῶν τῆς μητρός μου,
\vs{11}ἐπὶ σὲ ἐπεῤῥίφην ἐκ μήτρας· ἐκ κοιλίας μητρός μου Θεός μου εἶ σύ.

\vs{12}Μὴ ἀποστῇς ἀπʼ ἐμοῦ· ὅτι θλίψις ἐγγὺς, ὅτι οὐκ ἔστιν ὁ βοηθῶν.
\vs{13}Περιεκύκλωσάν με μόσχοι πολλοί, ταῦροι πίονες περιέσχον με.
\vs{14}Ἤνοιξαν ἐπʼ ἐμὲ τὸ στόμα αὐτῶν, ὡς λέων ὁ ἁρπάζων καὶ ὠρυόμενος.
\vs{15}Ὡσεὶ ὕδωρ ἐξεχύθην, καὶ διεσκορπίσθη πάντα τὰ ὀστᾶ μου, ἐγενήθη ἡ καρδία μου ὡσεὶ κηρὸς τηκόμενος ἐν μέσῳ τῆς κοιλίας μου.
\vs{16}Ἐξηράνθη ὡσεὶ ὄστρακον ἡ ἰσχύς μου, καὶ ἡ γλῶσσά μου κεκόλληται τῷ λάρυγγί μου, καὶ εἰς χοῦν θανάτου κατήγαγές με.
\vs{17}Ὅτι ἐκύκλωσάν με κύνες πολλοὶ, συναγωγὴ πονηρευομένων περιέσχον με· ὤρυξαν χεῖράς μου, καὶ πόδας,
\vs{18}ἐξηρίθμησαν πάντα τὰ ὀστᾶ μου· αὐτοὶ δὲ κατενόησαν καὶ ἐπεῖδόν με.
\vs{19}Διεμερίσαντο τὰ ἱμάτιά μου ἑαυτοῖς, καὶ ἐπὶ τὸν ἱματισμόν μου ἔβαλον κλῆρον.

\vs{20}Σὺ δὲ Κύριε μὴ μακρύνῃς τὴν βοήθειάν μου, εἰς τὴν ἀντίληψίν μου πρόσχες.
\vs{21}Ῥῦσαι ἀπὸ ῥομφαίας τὴν ψυχήν μου, καὶ ἐκ χειρὸς κυνὸς τὴν μονογενῆ μου.
\vs{22}Σῶσόν με ἐκ στόματος λέοντος, καὶ ἀπὸ κεράτων μονοκερώτων τὴν ταπείνωσίν μου.

\vs{23}Διηγήσομαι τὸ ὄνομά σου τοῖς ἀδελφοῖς μου, ἐν μέσῳ ἐκκλησίας ὑμνήσω σε.
\vs{24}Οἱ φοβούμενοι Κύριον αἰνέσατε αὐτὸν, ἅπαν τὸ σπέρμα Ἰακὼβ δοξάσατε αὐτὸν, φοβηθήτωσαν αὐτὸν ἅπαν τὸ σπέρμα Ἰσραήλ.
\vs{25}Οτι οὐκ ἐξουδένωσεν οὐδὲ προσώχθισε τῇ δεήσει τοῦ πτωχοῦ, οὐδὲ ἀπέστρεψε τὸ πρόσωπον αὐτοῦ ἀπʼ ἐμοῦ· καὶ ἐν τῷ κεκραγέναι με πρὸς αὐτὸν εἰσήκουσέ μου.
\vs{26}Παρὰ σοῦ ὁ ἔπαινός μου ἐν ἐκκλησίᾳ μεγάλῃ, τὰς εὐχάς μου ἀποδώσω ἐνώπιον τῶν φοβουμένων αὐτόν.

\vs{27}Φάγονται πένητες καὶ ἐμπλησθήσονται, καὶ αἰνέσουσι Κύριον οἱ ἐκζητοῦντες αὐτὸν, ζήσονται αἱ καρδίαι αὐτῶν εἰς αἰῶνα αἰῶνος.
\vs{28}Μνησθήσονται καὶ ἐπιστραφήσονται πρὸς Κύριον πάντα τὰ πέρατα τῆς γῆς, καὶ προσκυνήσουσιν ἐνώπιον αὐτοῦ πᾶσαι αἱ πατριαὶ τῶν ἐθνῶν.
\vs{29}Ὅτι τοῦ Κυρίου ἡ βασιλεία, καὶ αὐτὸς δεσπόζει τῶν ἐθνῶν.
\vs{30}Ἔφαγον καὶ προσεκύνησαν πάντες οἱ πίονες τῆς γῆς· ἐνώπιον αὐτοῦ προπεσοῦνται πάντες οἱ καταβαίνοντες εἰς τὴν γῆν· καὶ ἡ ψυχή μου αὐτῷ ζῇ,
\vs{31}καὶ τὸ σπέρμα μου δουλεύσει αὐτῷ· ἀναγγελήσεται τῷ Κυρίῳ γενεὰ ἡ ἐρχομένη·
\vs{32}Καὶ ἀναγγελοῦσι τὴν δικαιοσύνην αὐτοῦ λαῷ τῷ τεχθησομένῳ, ὃν ἐποίησεν ὁ Κύριος.

\begin{psalmheading}{\ch{22}{23} Ψαλμὸς τῷ Δαυίδ.}
\end{psalmheading}
Κύριος ποιμαίνει με, καὶ οὐδέν με ὑστερήσει.
\vs{2}Εἰς τόπον χλόης ἐκεῖ με κατεσκήνωσεν· ἐπὶ ὕδατος ἀναπαύσεως ἐξέθρεψέ με·
\vs{3}Τὴν ψυχήν μου ἐπέστρεψεν· ὡδήγησέν με ἐπὶ τρίβους δικαιοσύνης, ἕνεκεν τοῦ ὀνόματος αὐτοῦ.
\vs{4}Ἐὰν γὰρ καὶ πορευθῶ ἐν μέσῳ σκιᾶς θανάτου, οὐ φοβηθήσομαι κακὰ, ὅτι σὺ μετʼ ἐμοῦ εἶ· ἡ ῥάβδος σου καὶ ἡ βακτηρία σου, αὗταί με παρεκάλεσαν.
\vs{5}Ἡτοίμασας ἐνώπίον μου τράπεζαν, ἐξεναντίας τῶν θλιβόντων με· ἐλίπανας ἐν ἐλαίῳ τὴν κεφαλήν μου, καὶ τὸ ποτήριόν σου μεθύσκον ὡς κράτιστον.
\vs{6}Καὶ τὸ ἔλεός σου καταδιώξεταί με πάσας τὰς ἡμέρας τῆς ζωῆς μου, καὶ τὸ κατοικεῖν με ἐν οἴκῳ Κυρίου εἰς μακρότητα ἡμερῶν.

\begin{psalmheading}{\ch{23}{24} Ψαλμὸς τῷ Δαυὶδ τῆς μιᾶς σαββάτου.}
\end{psalmheading}
Τοῦ Κυρίου ἡ γῆ καὶ τὸ πλήρωμα αὐτῆς, ἡ οἰκουμένη καὶ πάντες οἱ κατοικοῦντες ἐν αὐτῇ.
\vs{2}Αὐτὸς ἐπὶ θαλασσῶν ἐθεμελιώσεν αὐτὴν, καὶ ἐπὶ ποταμῶν ἡτοίμασεν αὐτήν.

\vs{3}Τίς ἀναβήσεται εἰς τὸ ὄρος τοῦ Κυρίου, καὶ τίς στήσεται ἐν τόπῳ ἁγίῳ αὐτοῦ;
\vs{4}Ἀθῶος χερσὶ καὶ καθαρὸς τῇ καρδίᾳ, ὃς οὐκ ἔλαβεν ἐπὶ ματαίῳ τὴν ψυχὴν αὐτοῦ, καὶ οὐκ ὤμοσεν ἐπὶ δόλῳ τῷ πλησίον αὐτοῦ.
\vs{5}Οὗτος λήψεται εὐλογίαν παρὰ Κυρίου, καὶ ἐλεημοσύνην παρὰ Θεοῦ σωτῆρος αὐτοῦ.
\vs{6}Αὕτη ἡ γενεὰ ζητούντων αὐτὸν, ζητούντων τὸ πρόσωπον τοῦ Θεοῦ Ἰακώβ. διάψαλμα.

\vs{7}Ἄρατε πύλας οἱ ἄρχοντες ὑμῶν, καὶ ἐπάρθητε πύλαι αἰώνιοι, καὶ εἰσελεύσεται ὁ βασιλεὺς τῆς δόξης.
\vs{8}Τίς ἐστιν οὗτος ὁ βασιλεὺς τῆς δόξης; Κύριος κραταιὸς καὶ δυνατὸς, Κύριος δυνατὸς ἐν πολέμῳ.
\vs{9}Ἄρατε πύλας οἱ ἄρχοντες ὑμῶν, καὶ ἐπάρθητε πύλαι αἰώνιοι, καὶ εἰσελεύσεται ὁ βασιλεὺς τῆς δόξης.
\vs{10}Τίς ἐστιν οὗτος ὁ βασιλεὺς τῆς δόξης; Κύριος τῶν δυνάμεων, αὐτός ἐστιν οὗτος ὁ βασιλεὺς τῆς δόξης.

\begin{psalmheading}{\ch{24}{25} Ψαλμὸς τῷ Δαυίδ.}
\end{psalmheading}
Πρὸς σὲ, Κύριε, ᾖρα τὴν ψυχήν μου.
\vs{2}Ὁ Θεός μου ἐπὶ σοὶ πέποιθα, μὴ καταισχυνθείην· μηδὲ καταγελασάτωσάν μου οἱ ἐχθροί μου,
\vs{3}καὶ γὰρ πάντες οἱ ὑπομένοντές σε οὐ μὴ καταισχυνθῶσιν· αἰσχυνθήτωσαν οἱ ἀνομοῦντες διακενῆς.
\vs{4}Τὰς ὁδούς σου, Κύριε, γνώρισόν μοι, καὶ τὰς τρίβους σου δίδαξόν με.
\vs{5}Ὁδήγησόν με ἐπὶ τὴν ἀλήθειάν σου, καὶ δίδαξόν με, ὅτι σὺ εἶ ὁ Θεὸς ὁ σωτήρ μου, καὶ σὲ ὑπέμεινα ὅλην τὴν ἡμέραν.
\vs{6}Μνήσθητι τῶν οἰκτιρμῶν σου Κύριε, καὶ τὰ ἐλέη σου, ὅτι ἀπὸ τοῦ αἰῶνος εἰσίν.
\vs{7}Ἁμαρτίας νεότητός μου, καὶ ἀγνοίας μου μὴ μνησθῇς· κατὰ τὸ ἔλεός σου μνήσθητί μου, ἕνεκεν τῆς χρηστότητος σου, Κύριε.

\vs{8}Χρηστὸς καὶ εὐθὴς ὁ Κύριος, διὰ τοῦτο νομοθετήσει ἁμαρτάνοντας ἐν ὁδῷ.
\vs{9}Ὁδηγήσει πρᾳεῖς ἐν κρίσει, διδάξει πρᾳεῖς ὁδοὺς αὐτοῦ.
\vs{10}Πᾶσαι αἱ ὁδοὶ Κυρίου ἔλεος καὶ ἀλήθεια τοῖς ἐκζητοῦσι τὴν διαθήκην αὐτοῦ καὶ τὰ μαρτύρια αὐτοῦ.
\vs{11}Ἕνεκα τοῦ ὀνόματός σου, Κύριε, καὶ ἱλάσῃ τῇ ἁμαρτίᾳ μου, πολλὴ γάρ ἐστι.
\vs{12}Τίς ἐστιν ἄνθρωπος ὁ φοβούμενος τὸν Κύριον; νομοθετήσει αὐτῷ ἐν ὁδῷ, ᾗ ᾑρετίσατο.
\vs{13}Ἡ ψυχὴ αὐτοῦ ἐν ἀγαθοῖς αὐλισθήσεται, καὶ τὸ σπέρμα αὐτοῦ κληρονομήσει γῆν.
\vs{14}Κραταίωμα Κύριος τῶν φοβουμένων αὐτὸν, καὶ ἡ διαθήκη αὐτοῦ τοῦ δηλῶσαι αὐτοῖς.

\vs{15}Οἱ ὀφθαλμοί μου διαπαντὸς πρὸς τὸν Κύριον, ὅτι αὐτὸς ἐκσπάσει ἐκ παγίδος τοὺς πόδας μου.
\vs{16}Ἐπίβλεψον ἐπʼ ἐμὲ καὶ ἐλέησόν με, ὅτι μονογενὴς καὶ πτωχός εἰμι ἐγώ.
\vs{17}Αἱ θλίψεις τῆς καρδίας μου ἐπληθύνθησαν, ἐκ τῶν ἀναγκῶν μου ἐξάγαγέ με·
\vs{18}Ἴδε τὴν ταπείνωσίν μου καὶ τὸν κόπον μου, καὶ ἄφες πάσας τὰς ἁμαρτίας μου.
\vs{19}Ἴδε τοὺς ἐχθρούς μου, ὅτι ἐπληθύνθησαν, καὶ μῖσος ἄδικον ἐμίσησάν με.
\vs{20}Φύλαξον τὴν ψυχήν μου καὶ ῥῦσαί με· μὴ καταισχυνθείην, ὅτι ἤλπισα ἐπὶ σέ.
\vs{21}Ἄκακοι καὶ εὐθεῖς ἐκολλῶντό μοι, ὅτι ὑπέμεινά σε, Κύριε.
\vs{22}Λύτρωσαι ὁ Θεὸς τὸν Ἰσραὴλ ἐκ πασῶν τῶν θλίψεων αὐτοῦ.

\begin{psalmheading}{\ch{25}{26} Τοῦ Δαυίδ.}
\end{psalmheading}
Κρίνον με, Κύριε, ὅτι ἐγὼ ἐν ἀκακίᾳ μου ἐπορεύθην, καὶ ἐπὶ τῷ Κυρίῳ ἐλπίζων οὐ μὴ σαλευθῶ.
\vs{2}Δοκίμασόν με, Κύριε, καὶ πείρασόν με, πύρωσον τοὺς νεφρούς μου καὶ τὴν καρδίαν μου.

\vs{3}Ὅτι τὸ ἔλεός σου κατέναντι τῶν ὀφθαλμῶν μου ἐστὶ, καὶ εὐηρέστησα ἐν τῇ ἀληθείᾳ σου.
\vs{4}Οὐκ ἐκάθισα μετὰ συνεδρίου ματαιότητος, καὶ μετὰ παρανομούντων οὐ μὴ εἰσέλθω·
\vs{5}Ἐμίσησα ἐκκλησίαν πονηρευομένων, καὶ μετὰ ἀσεβῶν οὐ μὴ καθίσω.
\vs{6}Νίψομαι ἐν ἀθώοις τὰς χεῖράς μου, καὶ κυκλώσω τὸ θυσιαστήριόν σου, Κύριε·
\vs{7}Τοῦ ἀκοῦσαι φωνῆς αἰνέσεως, καὶ διηγήσασθαι πάντα τὰ θαυμάσιά σου.
\vs{8}Κύριε, ἠγάπησα εὐπρέπειαν οἴκου σου, καὶ τόπον σκηνώματος δόξης σου.
\vs{9}Μὴ συναπολέσῃς μετὰ ἀσεβῶν τὴν ψυχήν μου, καὶ μετὰ ἀνδρῶν αἱμάτων τὴν ζωήν μου,
\vs{10}ὧν ἐν χερσὶν ἀνομίαι, ἡ δεξιὰ αὐτῶν ἐπλήσθη δώρων.
\vs{11}Ἐγὼ δὲ ἐν ἀκακίᾳ μου ἐπορεύθην, λύτρωσαί με καὶ ἐλέησόν με.
\vs{12}Ὁ ποῦς μου ἔστη ἐν εὐθύτητι, ἐν ἐκκλησίαις εὐλογήσω σε Κύριε.

\begin{psalmheading}{\ch{26}{27} Τοῦ Δαυὶδ, πρὸ τοῦ χρισθῆναι.}
\end{psalmheading}
Κύριος φωτισμός μου καὶ σώτηρ μου, τίνα φοβηθήσομαι; Κύριος ὑπερασπιστὴς τῆς ζωῆς μου, ἀπὸ τίνος δειλιάσω;
\vs{2}Ἐν τῷ ἐγγίζειν ἐπʼ ἐμὲ κακοῦντας, τοῦ φαγεῖν τὰς σάρκας μου, οἱ θλίβοντες με καὶ οἱ ἐχθροί μου, αὐτοὶ ἠσθένησαν καὶ ἔπεσαν.
\vs{3}Ἐὰν παρατάξηται ἐπʼ ἐμὲ παρεμβολὴ, οὐ φοβηθήσεται ἡ καρδία μου· ἐὰν ἐπαναστῇ ἐπʼ ἐμὲ πόλεμος, ἐν ταύτῃ ἐγὼ ἐλπιζω.
\vs{4}Μίαν ᾐτησάμην παρὰ Κυρίου, ταύτην ἐκζητήσω, τοῦ κατοικεῖν με ἐν οἴκῳ Κυρίου πάσας τὰς ἡμέρας τῆς ζωῆς μου, τοῦ θεωρεῖν με τὴν τερπνότητα Κυρίου, καὶ ἐπισκέπτεσθαι τὸν ναὸν αὐτοῦ.
\vs{5}Ὅτι ἔκρυψέ με ἐν σκηνῇ αὐτοῦ ἐν ἡμέρᾳ κακῶν μου, ἐσκέπασέ με ἐν ἀποκρύφῳ τῆς σκῆνης αὐτοῦ· ἐν πέτρᾳ ὕψωσέ με,
\vs{6}καὶ νῦν ἰδοὺ ὕψωσε τὴν κεφαλήν μου ἐπʼ ἐχθρούς μου· ἐκύκλωσα καὶ ἔθυσα ἐν τῇ σκηνῇ αὐτοῦ θυσίαν ἀλαλαγμοῦ, ᾄσομαι καὶ ψαλῶ τῷ Κυρίῳ·

\vs{7}Εἰσάκουσον, Κύριε, τῆς φωνῆς μου ἧς ἐκέκραξα, ἐλέησόν με, καὶ εἰσάκουσόν μου.
\vs{8}Σοὶ εἶπεν ἡ καρδία μου, ἐξεζήτησα τὸ πρόσωπόν σου, τὸ πρόσωπόν σου Κύριε ζητήσω.
\vs{9}Μὴ ἀποστρέψῃς τὸ πρόσωπόν σου ἀπʼ ἐμοῦ, μὴ ἐκκλίνῃς ἐν ὀργῇ ἀπὸ τοῦ δούλου σου· βοηθός μου γενοῦ, μὴ ἐγκαταλίπῃς με, καὶ μὴ ὑπερίδῃς με ὁ Θεὸς ὁ σωτήρ μου.
\vs{10}Ὅτι ὁ πατήρ μου καὶ ἡ μήτηρ μου ἐγκατέλιπόν με, ὁ δὲ Κύριος προσελάβετό με.
\vs{11}Νομοθέτησόν με, Κύριε, ἐν τῇ ὁδῷ σου, καὶ ὁδήγησόν με ἐν τρίβῳ εὐθείᾳ ἕνεκα τῶν ἐχθρῶν μου.
\vs{12}Μὴ παραδῷς με εἰς ψυχὰς θλιβόντων με, ὅτι ἐπανέστησάν μοι μάρτυρες ἄδικοι, καὶ ἐψεύσατο ἡ ἀδικία ἑαυτῇ.

\vs{13}Πιστεύω τοῦ ἰδεῖν τὰ ἀγαθὰ Κυρίου ἐν γῇ ζώντων.
\vs{14}Ὑπόμεινον τὸν Κυριον, ἀνδρίζου, καὶ κραταιούσθω ἡ καρδία σου, καὶ ὑπόμεινον τὸν Κύριον.

\begin{psalmheading}{\ch{27}{28} Τοῦ Δαυίδ.}
\end{psalmheading}
Πρὸς σὲ Κύριε ἐκέκραξα, ὁ Θεός μου μὴ παρασιωπήσῃς ἐπʼ ἑμοὶ, μήποτε παρασιωπήσῃς ἐπʼ ἐμοὶ, καὶ ὁμοιωθήσομαι τοῖς καταβαίνουσιν εἰς λάκκον.
\vs{2}Εἰσάκουσον τῆς φωνῆς τῆς δεήσεώς μου, ἐν τῷ δέεσθαί με πρὸς σὲ, ἐν τῷ αἴρειν με χεῖράς μου εἰς ναὸν ἅγιόν σου.
\vs{3}Μὴ συνελκύσῃς μετὰ ἁμαρτωλῶν τὴν ψυχήν μου, καὶ μετὰ ἐργαζομένων ἀδικίαν μὴ συναπολέσῃς με, τῶν λαλούντων εἰρήνην μετὰ τῶν πλησίον αὐτῶν, κακὰ δὲ ἐν ταῖς καρδίαις αὐτῶν.
\vs{4}Δὸς αὐτοῖς κατὰ τὰ ἔργα αὐτῶν, καὶ κατὰ τὴν πονηρίαν τῶν ἐπιτηδευμάτων αὐτῶν· κατὰ τὰ ἔργα τῶν χειρῶν αὐτῶν δὸς αὐτοῖς, ἀπόδος τὸ ἀνταπόδομα αὐτῶν αὐτοῖς.
\vs{5}Ὅτι οὐ συνῆκαν εἰς τὰ ἔργα Κυρίου καὶ εἰς τὰ ἔργα τῶν χειρῶν αὐτοῦ, καθελεῖς αὐτοὺς καὶ οὐ μὴ οἰκοδομήσεις αὐτούς.

\vs{6}Εὐλογητὸς Κύριος, ὅτι εἰσήκουσε τῆς φωνῆς τῆς δεήσεώς μου.
\vs{7}Κύριος βοηθός μου καὶ ὑπερασπιστής μου· ἐπʼ αὐτῷ ἤλπισεν ἡ καρδία μου, καὶ ἐβοηθήθην, καὶ ἀνέθαλεν ἡ σάρξ μου· καὶ ἐκ θελήματός μου ἐξομολογήσομαι αὐτῷ.
\vs{8}Κύριος κραταίωμα τοῦ λαοῦ αὐτοῦ, καὶ ὑπερασπιστὴς τῶν σωτηρίων τοῦ χριστοῦ αὐτοῦ ἐστι.

\vs{9}Σῶσον τὸν λαόν σου, καὶ εὐλόγησον τὴν κληρονομίαν σου, καὶ ποίμανον αὐτοὺς, καὶ ἔπαρον αὐτοὺς ἕως τοῦ αἰῶνος.

\begin{psalmheading}{\ch{28}{29} Ψαλμὸς τῷ Δαυὶδ ἐξοδίου σκηνῆς.}
\end{psalmheading}
Ἐνέγκατε τῷ Κυρίῳ υἱοὶ Θεοῦ, ἐνέγκατε τῷ Κρίῳ υἱοὺς κριῶν· ἐνέγκατε τῷ Κυρίῳ δόξαν καὶ τιμὴν,
\vs{2}ἐνέγκατε τῷ Κυρίῳ δόξαν ὀνόματι αὐτοῦ· προσκυνήσατε τῷ Κυρίῳ ἐν αὐλῇ ἁγίᾳ αὐτοῦ.

\vs{3}Φωνὴ Κυρίου ἐπὶ τῶν ὑδάτων, ὁ Θεὸς τῆς δόξης ἐβρόντησε, Κύριος ἐπὶ ὑδάτων πολλῶν.
\vs{4}Φωνὴ Κυρίου ἐν ἰσχύϊ, φωνὴ Κυρίου ἐν μεγαλοπρεπείᾳ.
\vs{5}Φωνὴ Κυρίου συντρίβοντος κέδρους, συντρίψει Κύριος τὰς κέδρους τοῦ Λιβάνου,
\vs{6}Καὶ λεπτυνεῖ αὐτὰς ὡς τὸν μόσχον τὸν Λίβανον, καὶ ὁ ἠγαπημένος ὡς υἱὸς μονοκερώτων.
\vs{7}Φωνὴ Κυρίου διακόπτοντος φλόγα πυρός.
\vs{8}Φωνὴ Κυρίου συσσείοντος ἔρημον, συνσσείσει Κύριος τὴν ἔρημον Κάδης.
\vs{9}Φωνὴ Κυρίου καταρτιζομένου ἐλάφους, καὶ ἀποκαλύψει δρυμοὺς, καὶ ἐν τῷ ναῷ αὐτοῦ πᾶς τις λέγει δόξαν.
\vs{10}Κύριος τὸν κατακλυσμὸν κατοικιεῖ· καὶ καθιεῖται Κύριος βασιλεὺς εἰς τὸν αἰῶνα.
\vs{11}Κύριος ἰσχὺν τῷ λαῷ αὐτοῦ δώσει, Κύριος εὐλογήσει τὸν λαὸν αὐτοῦ ἐν εἰρήνῃ.

\begin{psalmheading}{\ch{29}{30} Εἰς τὸ τέλος, ψαλμὸς ᾠδῆς τοῦ ἐνκαινισμοῦ τοῦ οἴκου τοῦ Δαυίδ.}
\end{psalmheading}
\vs{2}Ὑψώσω σε, Κύριε, ὅτι ὑπέλαβές με, καὶ οὐκ εὔφρανας τοὺς ἐχθρούς μου ἐπʼ ἐμέ.
\vs{3}Κύριε ὁ Θεός μου, ἐκέκραξα πρὸς σὲ, καὶ ἰάσω με·
\vs{4}Κύριε, ἀνήγαγες ἐξ ᾅδου τὴν ψυχήν μου, ἔσωσάς με ἀπὸ τῶν καταβαινόντων εἰς λάκκον.

\vs{5}Ψάλατε τῷ Κυρίῳ οἱ ὅσιοι αὐτοῦ, καὶ ἐξομολογεῖσθε τῇ μνήμῃ τῆς ἁγιωσύνης αὐτοῦ.
\vs{6}Ὅτι ὀργὴ ἐν τῷ θυμῷ αὐτοῦ, καὶ ζωὴ ἐν τῷ θελήματι αὐτοῦ· τοεσπέρας αὐλισθήσεται κλαυθμὸς, καὶ εἰς τοπρωῒ ἀγαλλίασις.

\vs{7}Ἐγὼ δὲ εἶπα ἐν τῇ εὐθηνίᾳ μου, οὐ μὴ σαλευθῶ εἰς τὸν αἰῶνα.
\vs{8}Κύριε, ἐν τῷ θελήματί σου παρέσχου τῷ κάλλει μου δύναμιν· ἀπέστρεψας δὲ τὸ πρόσωπόν σου, καὶ ἐγενήθην τεταραγμένος.
\vs{9}Πρὸς σὲ Κύριε κεκράξομαι, καὶ πρὸς τὸν Θεόν μου δεηθήσομαι.
\vs{10}Τίς ὠφέλεια ἐν τῷ αἵματί μου, ἐν τῷ καταβῆναί με εἰς διαφθοράν; μὴ ἐξομολογήσεταί σοι χοῦς; ἢ ἀναγγελεῖ τὴν ἀλήθειάν σου;
\vs{11}Ἤκουσε Κύριος, καὶ ἠλέησέ με, Κύριος ἐγενήθη βοηθός μου.
\vs{12}Ἔστρεψας τὸν κοπετόν μου εἰς χαρὰν ἐμοὶ, διεῤῥήξας τὸν σάκκον μου, καὶ περιέζωσάς με εὐφροσύνην·
\vs{13}ὅπως ἂν ψάλῃ σοι ἡ δόξα μου, καὶ οὐ μὴ κατανυγῶ. Κύριε ὁ Θεός μου, εἰς τὸν αἰῶνα ἐξομολογήσομαί σοι.

\begin{psalmheading}{\ch{30}{31} Εἰς τὸ τέλος, ψαλμὸς τῷ Δαυὶδ, ἐκστάσεως.}
\end{psalmheading}
\vs{2}Ἐπὶ σοὶ Κύριε ἤλπισα, μὴ καταισχυνθείην εἰς τὸν αἰῶνα, ἐν τῇ δικαιοσύνῃ σου ῥῦσαί με καὶ ἐξελοῦ με.
\vs{3}Κλῖνον πρὸς μὲ τὸ οὖς σου, τὰχυνον τοῦ ἐξελέσθαι με γενοῦ μοι εἰς Θεὸν ὑπερασπιστὴν καὶ εἰς οἶκον καταφυγῆς τοῦ σῶσαί με.
\vs{4}Ὅτι κραταίωμά μου καὶ καταφυγή μου εἶ σὺ, καὶ ἕνεκεν τοῦ ὀνόματός σου ὁδηγήσεις με, καὶ διαθρέψεις με.
\vs{5}Ἐξάξεις με ἐκ παγίδος ταύτης ἧς ἔκρυψάν μοι, ὅτι σὺ εἶ ὁ ὑπερασπιστής μου, Κύριε.
\vs{6}Εἰς χεῖράς σου παραθήσομαι τὸ πνεῦμά μου, ἐλυτρώσω με Κύριε ὁ Θεὸς τῆς ἀληθέιας.
\vs{7}Ἐμίσησας τοὺς διαφυλάσσοντας ματαιότητας διακενῆς· ἐγὼ δὲ ἐπὶ τῷ Κυρίῳ ἤλπισα.
\vs{8}Ἀγαλλιάσομαι καὶ εὐφρανθήσομαι ἐπὶ τῷ ἐλέει σου· ὅτι ἐπεῖδες τὴν ταπείνωσίν μου, ἔσωσας ἐκ τῶν ἀναγκῶν τὴν ψυχήν μου.
\vs{9}Καὶ οὐ συνέκλεισάς με εἰς χεῖρας ἐχθροῦ, ἔστησας ἐν εὐρυχώρῳ τοὺς πόδας μου.

\vs{10}Ἐλέησόν με, Κύριε, ὅτι θλίβομαι· ἐταράχθη ἐν θυμῷ ὁ ὀφθαλμός μου, ἡ ψυχή μου, καὶ ἡ γαστήρ μου.
\vs{11}Ὅτι ἐξέλιπεν ἐν ὀδύνῃ ἡ ζωή μου, καὶ τὰ ἔτη μου ἐν στεναγμοῖς· ἠσθένησεν ἐν πτωχείᾳ ἡ ἰσχύς μου, καὶ τὰ ὀστᾶ μου ἐταράχθησαν.
\vs{12}Παρὰ πάντας τοὺς ἐχθρούς μου ἐγενήθην ὄνειδος, καὶ τοῖς γείτοσί μου σφόδρα, καὶ φόβος τοῖς γνωστοῖς μου· οἱ θεωροῦντές με ἔξω ἔφυγον ἀπʼ ἐμοῦ·
\vs{13}Ἐπελήσθην ὡσεὶ νεκρὸς ἀπὸ καρδίας· ἐγενήθην ὡσεὶ σκεῦος ἀπολωλὸς,
\vs{14}ὅτι ἤκουσα ψόγον πολλῶν παροικούντων κυκλόθεν· ἐν τῷ συναχθῆναι αὐτοὺς ἅμα ἐπʼ ἐμὲ, τοῦ λαβεῖν τὴν ψυχήν μου ἐβουλεύσαντο.

\vs{15}Ἐγὼ δὲ ἐπὶ σοὶ ἤλπισα, Κύριε· εἶπα, σὺ εἶ ὁ Θεός μου,
\vs{16}ἐν ταῖς χερσί σου οἱ κλῆροί μου· ῥῦσαί με ἐκ χειρὸς ἐχθρῶν μου, καὶ ἐκ τῶν καταδιωκόντων με.
\vs{17}Ἐπίφανον τὸ πρόσωπόν σου ἐπὶ τὸν δοῦλόν σου, σῶσόν με ἐν τῷ ἐλέει σου.
\vs{18}Κύριε, μὴ καταισχυνθείην, ὅτι ἐπεκαλεσάμην σε· αἰσχυνθείησαν οἱ ἀσεβεῖς, καὶ καταχθείησαν εἰς ᾅδου.
\vs{19}Ἄλαλα γενηθήτω τὰ χείλη τὰ δόλια, τὰ λαλοῦντα κατὰ τοῦ δικαίου ἀνομίαν ἐν ὑπερηφανίᾳ καὶ ἐξουδενώσει.

\vs{20}Ὡς πολὺ τὸ πλῆθος τῆς χρηστότητός σου, Κύριε, ἧς ἔκρυψας τοῖς φοβουμένοις σε; ἐξειργάσω τοῖς ἐλπίζουσιν ἐπὶ σὲ, ἐναντίον τῶν υἱῶν τῶν ἀνθρώπων.
\vs{21}Κατακρύψεις αὐτοὺς ἐν ἀποκρύφῳ τοῦ προσώπου σου ἀπὸ ταραχῆς ἀνθρώπων, σκεπάσεις αὐτοὺς ἐν σκηνῇ ἀπὸ ἀντιλογίας γλωσσῶν.
\vs{22}Εὐλογητὸς Κύριος, ὅτι ἐθαυμάστωσε τὸ ἔλεος αὐτοῦ ἐν πόλει περιοχῆς.
\vs{23}Ἐγὼ δὲ εἶπα ἐν τῇ ἐκστάσει μου, ἀπέῤῥιμμαι ἀπὸ προσώπου τῶν ὀφθαλμῶν σου· διὰ τοῦτο εἰσήκουσας, Κύριε, τῆς φωνῆς τῆς δεήσεώς μου ἐν τῷ κεκραγέναι με πρὸς σέ.

\vs{24}Ἀγαπήσατε τὸν Κύριον πάντες οἱ ὅσιοι αὐτοῦ, ὅτι ἀληθείας ἐκζητεῖ Κύριος, καὶ ἀνταποδίδωσι τοῖς περισσῶς ποιοῦσιν ὑπερηφανίαν.
\vs{25}Ἀνδρίζεσθε, καὶ κραταιούσθω ἡ καρδία ὑμῶν, πάντες οἱ ἐλπίζοντες ἐπὶ Κύριον.

\begin{psalmheading}{\ch{31}{32} Συνέσεως τῷ Δαυίδ.}
\end{psalmheading}
Μακάριοι ὧν ἀφέθησαν αἱ ἀνομίαι, καὶ ὧν ἐπεκαλύφθησαν αἱ ἁμαρτίαι.
\vs{2}Μακάριος ἀνὴρ ᾧ οὐ μὴ λογίσηται Κύριος ἁμαρτίαν, οὐδέ ἐστιν ἐν τῷ στόματι αὐτοῦ δόλος.

\vs{3}Ὅτι ἐσίγησα, ἐπαλαιώθη τὰ ὀστᾶ μου, ἀπὸ τοῦ κράζειν με ὅλην τὴν ἡμέραν.
\vs{4}Ὅτι ἡμέρας καὶ νυκτὸς ἐβαρύνθη ἐπʼ ἐμὲ ἡ χείρ σου, ἐστράφην εἰς ταλαιπωρίαν ἐν τῷ ἐμπαγῆναι ἄκανθαν· διάψαλμα.
\vs{5}Τὴν ἁμαρτίαν μου ἐγνώρισα, καὶ τὴν ἀνομίαν μου οὐκ ἐκάλυψα· εἶπα, ἐξαγορεύσω κατʼ ἐμοῦ τὴν ἀνομίαν μου τῷ Κυρίῳ, καὶ σὺ ἀφῆκας τὴν ἀσέβειαν τῆς καρδίας μου· διάψαλμα.
\vs{6}Ὑπὲρ ταύτης προσεύξεται πρὸς σὲ πᾶς ὅσιος ἐν καιρῷ εὐθέτῳ· πλὴν ἐν κατακλυσμῷ ὑδάτων πολλῶν πρὸς αὐτὸν οὐκ ἐγγιοῦσι.
\vs{7}Σύ μου εἶ καταφυγὴ ἀπὸ θλίψεως τῆς περιεχούσης με, τὸ ἀγαλλίαμά μου λύτρωσαί με ἀπὸ τῶν κυκλωσάντων με· διάψαλμα.

\vs{8}Συνετιῶ σε καὶ συμβιβῶ σε ἐν ὁδῷ ταύτῃ ᾗ πορεύσῃ, ἐπιστηριῶ ἐπὶ σὲ τοὺς ὀφθαλμούς μου.
\vs{9}Μὴ γίνεσθε ὡς ἵππος καὶ ἡμίονος, οἷς οὐκ ἔστι σύνεσις· ἐν χαλινῷ καὶ κημῷ τὰς σιαγόνας αὐτῶν ἄγξαι τῶν μὴ ἐγγιζόντων πρὸς σέ.
\vs{10}Πολλαὶ αἱ μάστιγες τοῦ ἁμαρτωλοῦ, τὸν δὲ ἐλπίζοντα ἐπὶ Κύριον ἔλεος κυκλώσει.
\vs{11}Εὐφράνθητε ἐπὶ Κύριον καὶ ἀγαλλιᾶσθε δίκαιοι, καὶ καυχᾶσθε πάντες οἱ εὐθεῖς τῇ καρδίᾳ.

\begin{psalmheading}{\ch{32}{33} Τῷ Δαυίδ.}
\end{psalmheading}
Ἀγαλλιᾶσθε δίκαιοι ἐν τῷ Κυρίῳ, τοῖς εὐθέσι πρέπει αἴνεσις.
\vs{2}Ἐξομολογεῖσθε τῷ Κυρίῳ ἐν κιθάρᾳ, ἐν ψαλτηρίῳ δεκαχόρδῳ ψάλατε αὐτῷ.
\vs{3}Ἄσατε αὐτῷ ᾆσμα καινὸν, καλῶς ψάλατε ἐν ἀλαλαγμῷ.

\vs{4}Ὅτι εὐθὴς ὁ λόγος τοῦ Κυρίου, καὶ πάντα τὰ ἔργα αὐτοῦ ἐν πίστει.
\vs{5}Ἀγαπᾷ ἐγεημοσύνην καὶ κρίσιν, τοῦ ἐλέους Κυρίου πλήρης ἡ γῆ.
\vs{6}Τῷ λόγῳ τοῦ Κυρίου οἱ οὐρανοὶ ἐστερεώθησαν, καὶ τῷ πνεύματι τοῦ στόματος αὐτοῦ πᾶσα ἡ δύναμις αὐτῶν.
\vs{7}Συνάγων ὡσεὶ ἀσκὸν ὕδατα θαλάσσης, τιθεὶς ἐν θησαυροῖς ἀβύσσους.
\vs{8}Φόβηθήτω τὸν Κύριον πᾶσα ἡ γῆ, ἀπʼ αὐτοῦ δὲ σαλευθήτωσαν πάντες οἱ κατοικοῦντες τὴν οἰκουμένην.
\vs{9}Ὅτι αὐτὸς εἶπε καὶ ἐγενήθησαν, αὐτὸς ἐνετείλατο καὶ ἐκτίσθησαν.
\vs{10}Κύριος διασκεδάζει βουλὰς ἐθνῶν, ἀθετεῖ δὲ λογισμοὺς λαῶν, καὶ ἀθετεῖ βουλὰς ἀρχόντων.
\vs{11}Ἡ δὲ βουλὴ τοῦ Κυρίου εἰς τὸν αἰῶνα μένει, λογισμοὶ τῆς καρδίας αὐτοῦ ἀπὸ γενεῶν εἰς γενεάς.
\vs{12}Μακάριον τὸ ἔθνος οὗ ἐστι Κύριος ὁ Θεὸς αὐτοῦ, λαὸς ὃν ἐξελέξατο εἰς κληρονομίαν ἑαυτῷ.
\vs{13}Ἐξ οὐρανοῦ ἐπέβλεψεν ὁ Κύριος, εἶδε πάντας τοὺς υἱοὺς τῶν ἀνθρώπων.
\vs{14}Ἐξ ἑτοίμου κατοικητηρίου αὐτοῦ ἐπέβλεψεν ἐπὶ πάντας τοὺς κατοικοῦντας τὴν γῆν.
\vs{15}Ὁ πλάσας κατὰ μόνας τὰς καρδίας αὐτῶν, ὁ συνιεὶς πάντα τὰ ἔργα αὐτῶν.
\vs{16}Οὐ σώζεται βασιλεὺς διὰ πολλὴν δύναμιν, καὶ γίγας οὐ σωθήσεται ἐν πλήθει ἰσχύος αὐτοῦ.
\vs{17}Ψευδὴς ἵππος εἰς σωτηρίαν, ἐν δὲ πλήθει δυνάμεως αὐτοῦ οὐ σωθήσεται.

\vs{18}Ἰδοὺ οἱ ὀφθαλμοὶ Κυρίου ἐπὶ τοὺς φοβουμένους αὐτὸν, τοὺς ἐλπίζοντας ἐπὶ τὸ ἔλεος αὐτοῦ,
\vs{19}ῥύσασθαι ἐκ θανάτου τὰς ψυχὰς αὐτῶν, καὶ διαθρέψαι αὐτοὺς ἐν λιμῷ.
\vs{20}Ἡ ψυχὴ ἡμῶν ὑπομένει τῷ Κυρίῳ, ὅτι βοηθὸς καὶ ὑπερασπιστὴς ἡμῶν ἐστιν.
\vs{21}Ὅτι ἐν αὐτῷ εὐφρανθήσεται ἡ καρδία ἡμῶν, καὶ ἐν τῷ ὀνόματι τῷ ἁγίῳ αὐτοῦ ἠλπίσαμεν.
\vs{22}Γένοιτο τὸ ἔλεός σου, Κύριε, ἐφʼ ἡμᾶς, καθάπερ ἠλπίσαμεν ἐπὶ σέ.

\begin{psalmheading}{\ch{33}{34} Τῷ Δαυὶδ, ὁπότε ἠλλοίωσε τὸ πρόσωπον αὐτοῦ ἐναντίον Ἀβιμέλεχ, καὶ ἀπέλυσεν αὐτὸν, καὶ ἀπῆλθεν.}
\end{psalmheading}
\vs{2}Εὐλογήσω τὸν Κύριον ἐν παντὶ καιρῷ, διαπαντὸς ἡ αἴνεσις αὐτοῦ ἐν τῷ στόματί μου.
\vs{3}Ἐν τῷ Κυρίῳ ἐπαινεθήσεται ἡ ψυχή μου· ἀκουσάτωσαν πρᾳεῖς καὶ εὐφρανθήτωσαν.
\vs{4}Μεγαλύνατε τὸν Κύριον σὺν ἐμοὶ, καὶ ὑψώσωμεν τὸ ὄνομα αὐτοῦ ἐπιτοαυτό.

\vs{5}Ἐξεζήτησα τὸν Κύριον καὶ ἐπήκουσέ μου, καὶ ἐκ πασῶν τῶν παροικιῶν μου ἐῤῥύσατό με·
\vs{6}Προσέλθατε πρὸς αὐτὸν καὶ φωτίσθητε, καὶ τὰ πρόσωπα ὑμῶν οὐ μὴ καταισχυνθῇ·
\vs{7}Οὗτος ὁ πτωχὸς ἐκέκραξε, καὶ ὁ Κύριος εἰσήκουσεν αὐτοῦ, καὶ ἐκ πασῶν τῶν θλίψεων αὐτοῦ ἔσωσεν αὐτόν.
\vs{8}Παρεμβαλεῖ ἄγγελος Κυρίου κύκλῳ τῶν φοβουμένων αὐτὸν, καὶ ῥύσεται αὐτούς.
\vs{9}Γεύσασθε καὶ ἴδετε ὅτι χρηστὸς ὁ Κύριος, μακάριος ἀνὴρ ὃς ἐλπίζει ἐπʼ αὐτόν.
\vs{10}φοβήθητε τὸν Κύριον πάντες οἱ ἅγιοι αὐτοῦ, ὅτι οὐκ ἔστιν ὑστέρημα τοῖς φοβουμένοις αὐτόν.
\vs{11}Πλούσιοι ἐπτώχευσαν καὶ ἐπείνασαν, οἱ δὲ ἐκζητοῦντες τὸν Κύριον οὐκ ἐλαττωθήσονται παντὸς ἀγαθοῦ· διάψαλμα.

\vs{12}Δεῦτε τέκνα, ἀκούσατέ μου, φόβον Κυρίου διδάξω ὑμᾶς.
\vs{13}Τίς ἐστιν ἄνθρωπος ὁ θέλων ζωὴν, ἀγαπῶν ἡμέρας ἰδεῖν ἀγαθάς;
\vs{14}Παῦσον τὴν γνῶσσάν σου ἀπὸ κακοῦ, καὶ χείλη σου τοῦ μὴ λαλῆσαι δόλον.
\vs{15}Ἔκκλινον ἀπὸ κακοῦ, καὶ ποίησον ἀγαθὸν· ζήτησον εἰρήνην, καὶ δίωξον αὐτήν.

\vs{16}Οφθαλμοὶ Κυρίου ἐπὶ δικαίους, καὶ ὦτα αὐτοῦ εἰς δέησιν αὐτῶν·
\vs{17}Πρόσωπον δὲ Κυρίου ἐπὶ ποιοῦντας κακὰ, τοῦ ἐξολοθρεῦσαι ἐκ γῆς τὸ μνημόσυνον αὐτῶν.
\vs{18}Ἐκέκραξαν οἱ δίκαιοι, καὶ ὁ Κύριος εἰσήκουσεν αὐτῶν, καὶ ἐκ πασῶν τῶν θλίψεων αὐτῶν ἐῤῥύσατο αὐτούς.
\vs{19}Ἐγγὺς Κύριος τοῖς συντετριμμένοις τὴν καρδίαν, καὶ τοὺς ταπεινοὺς τῷ πνεύματι σώσει.
\vs{20}Πολλαὶ αἱ θλίψεις τῶν δικαίων, καὶ ἐκ πασῶν αὐτῶν ῥύσεται αὐτοὺς
\vs{21}ὁ Κύριος. Φυλάσσει πάντα τὰ ὀστᾶ αὐτῶν, ἓν ἐξ αὐτῶν οὐ συντριβήσεται.
\vs{22}Θάνατος ἁμαρτωλῶν πονηρὸς, καὶ οἱ μισοῦντες τὸ δίκαιον πλημμελήσουσι.
\vs{23}Λυτρώσεται Κύριος ψυχὰς δούλων αὐτοῦ, καὶ οὐ μὴ πλημμελήσουσι πάντες οἱ ἐλπίζοντες ἐπʼ αὐτόν.

\begin{psalmheading}{\ch{34}{35} Τῷ Δαυίδ.}
\end{psalmheading}
Δικάσον, Κύριε, τοὺς ἀδικοῦντάς με, πολέμησον τοὺς πολεμοῦντάς με.
\vs{2}Ἐπιλαβοῦ ὅπλου καὶ θυρεοῦ, καὶ ἀνάστηθι εἰς βοήθειάν μοι.
\vs{3}Ἔκχεον ῥομφαίαν, καὶ σύγκλεισον ἐξεναντίας τῶν καταδιωκόντων με· εἰπὸν τῇ ψυχῇ μου, σωτηρία σου ἐγώ εἰμι.
\vs{4}Αἰσχυνθείησαν καὶ ἐντραπείησαν οἱ ζητοῦντες τὴν ψυχήν μου, ἀποστραφείησαν εἰς τὰ ὀπίσω, καὶ καταισχυνθείησαν οἱ λογιζόμενοί μοι κακά.
\vs{5}Γενηθήτωσαν ὡσεὶ χοῦς κατὰ πρόσωπον ἀνέμου, καὶ ἄγγελος Κυρίου ἐκθλίβων αὐτούς.
\vs{6}Γενηθήτω ἡ ὁδὸς αὐτῶν σκότος καὶ ὀλίσθημα, καὶ ἄγγελος Κυρίου καταδιώκων αὐτούς.
\vs{7}Ὅτι δωρεὰν ἔκρυψάν μοι διαφθορὰν παγίδος αὐτῶν, μάτην ὠνείδισαν τὴν ψυχήν μου.

\vs{8}Ἐλθέτω αὐτοῖς παγὶς ἣν οὐ γινώσκουσι, καὶ ἡ θήρα ἣν ἔκρυψαν, συλλαβέτω αὐτοὺς, καὶ ἐν τῇ παγίδι πεσοῦνται ἐν αὐτῇ.
\vs{9}Ἡ δὲ ψυχή μου ἀγαλλιάσεται ἐπὶ τῷ Κυρίῳ, τερφθήσεται ἐπὶ τῷ σωτηρίῳ αὐτοῦ.
\vs{10}Πάντα τὰ ὀστᾶ μου ἐροῦσι, Κύριε τίς ὅμοιός σοι; ῥυόμενος πτωχὸν ἐκ χειρὸς στερεωτέρων αὐτοῦ, καὶ πτωχὸν καὶ πένητα ἀπὸ τῶν διαρπαζόντων αὐτόν.

\vs{11}Ἀναστάντες μάρτυρες ἄδικοι, ἃ οὐκ ἐγίνωσκον, ἐπηρώτων με.
\vs{12}Ἀνταπεδίδοσάν μοι πονηρὰ ἀντὶ καλῶν, καὶ ἀτεκνίαν τῇ ψυχῇ μου.
\vs{13}Ἐγὼ δὲ ἐν τῷ αὐτοὺς παρενοχλεῖν μοι, ἐνεδυόμην σάκκον· καὶ ἐταπείνουν ἐν νηστείᾳ τὴν ψυχήν μου, καὶ ἡ προσευχή μου εἰς κόλπον μου ἀποστραφήσεται.
\vs{14}Ὡς πλησίον ὡς ἀδελφὸν ἡμέτερον οὕτως εὐηρέστουν, ὡς πενθῶν καὶ σκυθρωπάξων οὕτως ἐταπεινούμην.
\vs{15}Καὶ κατʼ ἐμοῦ εὐφράνθησαν, καὶ συνήχθησαν, συνήχθησαν ἐπʼ ἐμὲ μάστιγες καὶ οὐκ ἔγνων· διεσχίσθησαν καὶ οὐ κατενύγησαν.
\vs{16}Ἐπείρασάν με, ἐξεμυκτήρισάν με μυκτηρισμὸν, ἔβρυξαν ἐπʼ ἐμὲ τοὺς ὀδόντας αὐτῶν.

\vs{17}Κύριε πότε ἐπόψῃ; ἀποκατάστησον τὴν ψυχήν μου ἀπὸ τῆς κακουργίας αὐτῶν, ἀπὸ λεόντων τὴν μονογενῆ μου.
\vs{18}Ἐξομολογήσομαί σοι καὶ ἐν ἐκκλησίᾳ πολλῇ, ἐν λαῷ βαρεῖ αἰνέσω σε.
\vs{19}Μὴ ἐπιχαρείησάν μοι οἱ ἐχθραίνοντές μοι ματαίως, οἱ μισοῦντές με δωρεὰν, καὶ διανεύοντες ὀφθαλμοῖς.
\vs{20}Ὅτι ἐμοὶ μὲν εἰρηνικὰ ἐλάλουν, καὶ ἐπʼ ὀργῇ δόλους διελογίζοντο.
\vs{21}Καὶ ἐπλάτυναν ἐπʼ ἐμὲ τὸ στόμα αὐτῶν, εἶπαν, εὖγε, εὖγε, εἶδον οἱ ὀφθαλμοὶ ἡμῶν.

\vs{22}Εἶδες Κύριε, μὴ παρασιωπήσῃς, Κύριε μὴ ἀποστῇς ἀπʼ ἐμοῦ.
\vs{23}Ἐξεγέρθητι Κύριε, καὶ πρόσχες τῇ κρίσει μου, ὁ Θεός μου καὶ ὁ Κύριός μου εἰς τὴν δίκην μου.
\vs{24}Κρῖνόν με Κύριε κατὰ τὴν δικαιοσύνην σου Κύριε ὁ Θεός μου, καὶ μὴ ἐπιχαρείησάν μοι.
\vs{25}Μὴ εἴποισαν ἐν καρδίαις αὐτῶν, εὖγε, εὖγε τῇ ψυχῇ ἡμῶν· μηδὲ εἴποιεν, κατεπίομεν αὐτόν.
\vs{26}Αἰσχυνθείησαν καὶ ἐντραπείησαν ἅμα οἱ ἐπιχαίροντες τοῖς κακοῖς μου· ἐνδυσάσθωσαν αἰσχύνην καὶ ἐντροπὴν οἱ μεγαλοῤῥημονοῦντες ἐπʼ ἐμέ.
\vs{27}Ἀγαλλιάσαιντο καὶ εὐφρανθείησαν οἱ θέλοντες τὴν δικαιοσύνην μου, καὶ εἰπάτωσαν διαπαντὸς, μεγαλυνθείη ὁ Κύριος, οἱ θέλοντες τὴν εἰρήνην τοῦ δούλου αὐτοῦ.
\vs{28}Καὶ ἡ γλῶσσά μου μελετήσει τὴν δικαιοσύνην σου, ὅλην τὴν ἡμέραν τὸν ἔπαινόν σου.

\begin{psalmheading}{\ch{35}{36} Εἰς τὸ τέλος, τῷ δούλῳ Κυρίου τῷ Δαυίδ.}
\end{psalmheading}
\vs{2}Φησὶν ὁ παράνομος τοῦ ἁμαρτάνειν ἐν ἑαυτῷ, οὐκ ἔστι φόβος Θεοῦ ἀπέναντι τῶν ὀφθαλμῶν αὐτοῦ.
\vs{3}Ὅτι ἐδόλωσεν ἐνώπιον αὐτοῦ, τοῦ εὑρεῖν τὴν ἀνομίαν αὐτοῦ καὶ μισῆσαι.
\vs{4}Τὰ ῥήματα τοῦ στόματος αὐτοῦ ἀνομία καὶ δόλος, οὐκ ἠβουλήθη συνιέναι τοῦ ἀγαθῦναι.
\vs{5}Ἀνομίαν ἐλογίσατο ἐπὶ τῆς κοίτης αὐτοῦ, παρέστη πάσῃ ὁδῷ οὐκ ἀγαθῇ, τῇ δὲ κακίᾳ οὐ προσώχθισε.

\vs{6}Κύριε, ἐν τῷ οὐρανῷ τὸ ἔλεός σου, καὶ ἡ ἀλήθειά σου ἕως τῶν νεφελῶν.
\vs{7}Ἡ δικαιοσύνη σου ὡς ὄρη Θεοῦ, τὰ κρίματά σου ὡσεὶ ἄβυσσος πολλή· ἀνθρώπους καὶ κτήνη σώσεις Κύριε.
\vs{8}Ὡς ἐπλήθυνας τὸ ἔλεός σου ὁ Θεός· οἱ δὲ υἱοὶ τῶν ἀνθρώπων ἐν σκέπῃ τῶν πτερύγων σου ἐλπιοῦσι.
\vs{9}Μεθυσθήσονται ἀπὸ πιότητος οἴκου σου, καὶ τὸν χειμάῤῥουν τῆς τρυφῆς σου ποτιεῖς αὐτούς.
\vs{10}Ὅτι παρὰ σοὶ πηγὴ ζωῆς, ἐν τῷ φωτί σου ὀψόμεθα φῶς.

\vs{11}Παράτεινον τὸ ἔλεός σου τοῖς γινώσκουσί σε, καὶ τὴν δικαιοσύνην σου τοῖς εὐθέσι τῇ καρδίᾳ.
\vs{12}Μὴ ἐλθέτω μοι ποῦς ὑπερηφανίας, καὶ χεὶρ ἁμαρτωλῶν μὴ σαλεύσαι με·

\vs{13}Ἐκεῖ ἔπεσον πάντες οἱ ἐργαζόμενοι τὴν ἀνομίαν, ἐξώσθησαν καὶ οὐ μὴ δύνωνται στῆναι.

\begin{psalmheading}{\ch{36}{37} Τῷ Δαυίδ.}
\end{psalmheading}
Μὴ παραζήλου ἐν πονηρευομένοις, μηδὲ ζήλου τοὺς ποιοῦντας τὴν ἀνομίαν.
\vs{2}Ὅτι ὡσεὶ χόρτος ταχὺ ἀποξηρανθήσονται, καὶ ὡσεὶ λάχανα χλόης ταχὺ ἀποπεσοῦνται.
\vs{3}Ἔλπισον ἐπὶ Κύριον, καὶ ποίει χρηστότητα, καὶ κατασκήνου τὴν γῆν, καὶ ποιμανθήσῃ ἐπὶ τῷ πλούτῳ αὐτῆς.
\vs{4}Κατατρύφησον τοῦ Κυρίου, καὶ δώσει σοι τὰ αἰτήματα τῆς καρδίας σου.
\vs{5}Ἀποκάλυψον πρὸς Κύριον τὴν ὁδόν σου, καὶ ἔλπισον ἐπʼ αὐτὸν, καὶ αὐτὸς ποιήσει.
\vs{6}Καὶ ἐξοίσει ὡς φῶς τὴν δικαιοσύνην σου, καὶ τὸ κρίμα σου ὡς μεσημβρίαν.

\vs{7}Ὑποτάγηθι τῷ Κυρίῳ, καὶ ἱκέτευσον αὐτόν· μὴ παραζήλου ἐν τῷ κατευοδουμένῳ ἐν τῇ ὁδῷ αὐτοῦ, ἐν ἀνθρώπῳ ποιοῦντι παρανομίας.
\vs{8}Παῦσαι ἀπὸ ὀργῆς καὶ ἐνκατάλιπε θυμὸν, μὴ παραζήλου ὥστε πονηρεύεσθαι.
\vs{9}Ὅτι οἱ πονηρευόμενοι ἐξολοθρευθήσονται, οἱ δὲ ὑπομένοντες τὸν Κύριον, αὐτοὶ κληρονομήσουσι τὴν γῆν.
\vs{10}Καὶ ἔτι ὀλίγον καὶ οὐ μὴ ὑπάρξῃ ἁμαρτωλὸς, καὶ ζητήσεις τὸν τόπον αὐτοῦ, καὶ οὐ μὴ εὕρῃς.
\vs{11}Οἱ δὲ πρᾳεῖς κληρονομήσουσι γῆν, καὶ κατατρυφήσουσιν ἐπὶ πλήθει εἰρήνης.

\vs{12}Παρατηρήσεται ὁ ἁμαρτωλὸς τὸν δίκαιον, καὶ βρύξει ἐπʼ αὐτὸν τοὺς ὀδόντας αὐτοῦ.
\vs{13}Ὁ δὲ Κύριος ἐκγελάσεται αὐτὸν, ὅτι προβλέτει ὅτι ἥξει ἡ ἡμέρα αὐτοῦ.
\vs{14}Ῥομφαίαν ἐσπάσαντο οἱ ἁμαρτωλοὶ ἐνέτειναν τόξον αὐτῶν, τοῦ καταβαλεῖν πτωχὸν καὶ πένητα, τοῦ σφάξαι τοὺς εὐθεῖς τῇ καρδίᾳ.
\vs{15}Ἡ ῥομφαία αὐτῶν εἰσέλθοι εἰς τὴν καρδίαν αὐτῶν, καὶ τὰ τόξα αὐτῶν συντριβείη.

\vs{16}Κρεῖσσον ὀλίγον τῷ δικαίῳ ὑπὲρ πλοῦτον ἁμαρτωλῶν πολύν.
\vs{17}Ὅτι βραχίονες ἁμαρτωλῶν συντριβήσονται, ὑποστηρίζει δὲ τοὺς δικαίους ὁ Κύριος.

\vs{18}Γινώσκει Κύριος τὰς ὁδοὺς τῶν ἀμώμων, καὶ ἡ κληρονομία αὐτῶν εἰς τὸν αἰῶνα ἔσται.
\vs{19}Οὐ καταισχυνθήσονται ἐν καιρῷ πονηρῷ, καὶ ἐν ἡμέραις λιμοῦ χορτασθήσονται·
\vs{20}Ὅτι οἱ ἁμαρτωλοὶ ἀπολοῦνται, οἱ δὲ ἐχθροὶ τοῦ Κυρίου ἅμα τῷ δοξασθῆναι αὐτοὺς καὶ ὑψωθῆναι, ἐκλείποντες ὡσεὶ καπνὸς ἐξέλιπον.
\vs{21}Δανείζεται ὁ ἁμαρτωλὸς, καὶ οὐκ ἀποτίσει, ὁ δὲ δίκαιος οἰκτείρει καὶ διδοῖ.
\vs{22}Ὅτι οἱ εὐλογοῦντες αὐτὸν κληρονομήσουσι γῆν, οἱ δὲ καταρώμενοι αὐτὸν ἐξολοθρευθήσονται.

\vs{23}Παρὰ Κυρίου τὰ διαβήματα ἀνθρώπου κατευθύνεται, καὶ τὴν ὁδὸν αὐτοῦ θελήσει.
\vs{24}Ὅταν πέσῃ οὐ καταῤῥαχθήσεται, ὅτι Κύριος ἀντιστηρίζει χεῖρα αὐτοῦ.
\vs{25}Νεώτερος ἐγενόμην, καὶ γὰρ ἐγήρασα· καὶ οὐκ εἶδον δίκαιον ἐγκαταλελειμμένον, οὐδὲ τὸ σπέρμα αὐτοῦ ζητοῦν ἄρτους.
\vs{26}Ὅλην τὴν ἡμέραν ἐλεεῖ καὶ δανείζει, καὶ τὸ σπέρμα αὐτοῦ εἰς εὐλογίαν ἔσται.

\vs{27}Ἔκκλινον ἀπὸ κακοῦ, καὶ ποίησον ἀγαθὸν, καὶ κατασκήνου εἰς αἰῶνα αἰῶνος.
\vs{28}Ὅτι Κύριος ἀγαπᾷ κρίσιν, καὶ οὐκ ἐγκαταλείψει τοὺς ὁσίους αὐτοῦ, εἰς τὸν αἰῶνα φυλαχθήσονται· ἄμωμοι ἐκδικηθήσονται, καὶ σπέρμα ἀσεβῶν ἐξολοθρευθήσεται.
\vs{29}Δίκαιοι δὲ κληρονομήσουσι γῆν, καὶ κατασκηνώσουσιν εἰς αἰῶνα αἰῶνος ἐπʼ αὐτῆς.

\vs{30}Στόμα δικαίου μελετήσει σοφίαν, καὶ ἡ γλῶσσα αὐτοῦ λαλήσει κρίσιν.
\vs{31}Ὁ νόμος τοῦ Θεοῦ αὐτοῦ ἐν καρδίᾳ αὐτοῦ, καὶ οὐχ ὑποσκελισθήσεται τὰ διαβήματα αὐτοῦ.
\vs{32}Κατανοεῖ ὁ ἁμαρτωλὸς τὸν δίκαιον, καὶ ζητεῖ τοῦ θανατῶσαι αὐτόν.
\vs{33}Ὁ δὲ Κύριος οὐ μὴ ἐγκαταλίπῃ αὐτὸν εἰς τὰς χεῖρας αὐτοῦ, οὐδὲ μὴ καταδικάσαι αὐτὸν, ὅταν κρίνηται αὐτῷ.
\vs{34}Ὑπόμεινον τὸν Κύριον, καὶ φύλαξον τὴν ὁδὸν αὐτοῦ, καὶ ὑψώσει σε τοῦ κατακληρονομῆσαι τὴν γῆν· ἐν τῷ ἐξολοθρεύεσθαι ἁμαρτωλοὺς, ὄψει.

\vs{35}Εἶδον τὸν ἀσεβῆ ὑπερυψούμενον, καὶ ἐπαιρόμενον ὡς τὰς κέδρους τοῦ Λιβάνου·
\vs{36}Καὶ παρῆλθον, καὶ ἰδοὺ οὐκ ἦν, καὶ ἐζήτησα αὐτόν, καὶ οὐχ εὑρέθη ὁ τόπος αὐτοῦ.
\vs{37}Φύλασσε ἀκακίαν καὶ ἴδε εὐθύτητα, ὅτι ἐστὶν ἐγκατάλειμμα ἀνθρώπῳ εἰρηνικῷ.
\vs{38}Οἱ δὲ παράνομοι ἐξολοθρευθήσονται ἐπιτοαυτὸ, τὰ ἐγκαταλείμματα τῶν ἀσεβῶν ἐξολοθρευθήσονται·
\vs{39}Σωτηρία δὲ τῶν δικαίων παρὰ Κυρίου, καὶ ὑπερασπιστὴς αὐτῶν ἐστιν ἐν καιρῷ θλίψεως.
\vs{40}Καὶ βοηθήσει αὐτοῖς Κύριος, καὶ ῥύσεται αὐτοὺς, καὶ ἐξελεῖται αὐτοὺς ἐξ ἁμαρτωλῶν, καὶ σώσει αὐτοὺς; ὅτι ἤλπισαν ἐπʼ αὐτόν.

\begin{psalmheading}{\ch{37}{38} Φαλμὸς τῷ Δαυὶδ εἰς ἀνάμνησιν περὶ σαββάτου.}
\end{psalmheading}
\vs{2}Κύριε, μὴ τῷ θυμῷ σου ἐλέγξῃς με, μηδὲ τῇ ὀργῇ σου παιδεύσῃς με.
\vs{3}Ὅτι τὰ βέλη σου ἐνεπάγησάν μοι, καὶ ἐπεστήριξας ἐπʼ ἐμὲ τὴν χεῖρά σου.

\vs{4}Οὐκ ἔστιν ἴασις ἐν τῇ σαρκί μου ἀπὸ προσώπου τῆς ὀργῆς σου, οὐκ ἔστιν εἰρήνη τοῖς ὀστέοις μου ἀπὸ προσώπου τῶν ἁμαρτιῶν μου.
\vs{5}Ὅτι αἱ ἀνομίαι μου ὑπερῇραν τὴν κεφαλήν μου, ὡσεὶ φορτίον βαρὺ ἐβαρύνθησαν ἐπʼ ἐμέ.
\vs{6}Προσώζεσαν καὶ ἐσάπησαν οἱ μώλωπές μου, ἀπὸ προσώπου τῆς ἀφροσύνης μου.
\vs{7}Ἐταλαιπώρησα καὶ κατεκάμφθην ἕως τέλους, ὅλην τὴν ἡμέραν σκυθρωπάζων ἐπορευόμην.
\vs{8}Ὅτι ἡ ψυχή μου ἐπλησθη ἐμπαιγμῶν, καὶ οὐκ ἔστιν ἴασις ἐν τῇ σαρκί μου.
\vs{9}Ἐκακώθην καὶ ἐταπεινώθην ἕως σφόδρα, ὠρυόμην ἀπὸ στεναγμοῦ τῆς καρδίας μου.

\vs{10}Καὶ ἐναντίον σου πᾶσα ἡ ἐπιθυμία μου, καὶ ὁ στεναγμός μου οὐκ ἀπεκρύβη ἀπὸ σοῦ.
\vs{11}Ἡ καρδία μου ἐταράχθη, ἐγκατέλιπέ με ἡ ἰσχύς μου, καὶ τὸ φῶς τῶν ὀφθαλμῶν μου οὐκ ἔστι μετʼ ἐμοῦ.
\vs{12}Οἱ φίλοι μου καὶ οἱ πλησίον μου ἐξ ἐναντίας μου ἤγγισαν καὶ ἔστησαν, καὶ οἱ ἔγγιστά μου μακρόθεν ἔστησαν,
\vs{13}καὶ ἐξεβιάζοντο οἱ ζητοῦντες τὴν ψυχήν μου· καὶ οἱ ζητοῦντες τὰ κακά μοι ἐλάλησαν ματαιότητας, καὶ δολιότητας ὅλην τὴν ἡμέραν ἐμελέτησαν.
\vs{14}Ἐγὼ δὲ ὡσεὶ κωφὸς οὐκ ἤκουον, καὶ ὡσεὶ ἄλαλος οὐκ ἀνοίγων τὸ στόμα αὐτοῦ·
\vs{15}Καὶ ἐγενόμην ὡσεὶ ἄνθρωπος οὐκ ἀκούων, καὶ οὐκ ἔχων ἐν τῷ στόματι αὐτοῦ ἐλεγμούς.

\vs{16}Ὅτι ἐπὶ σοὶ Κύριε ἤλπισα, σὺ εἰσακούσῃ Κύριε ὁ Θεός μου.
\vs{17}Ὅτι εἶπα, μή ποτε ἐπιχαρῶσί μοι οἱ ἐχθροί μου, καὶ ἐν τῷ σαλευθῆναι πόδας μου, ἐπʼ ἐμὲ ἐμεγαλοῤῥημόνησαν.
\vs{18}Ὅτι ἐγὼ εἰς μάστιγας ἕτοιμος, καὶ ἡ ἀλγηδών μου ἐνώπιόν μου διαπαντός.
\vs{19}Ὅτι τὴν ἀνομίαν μου ἀναγγελῶ, καὶ μεριμνήσω ὑπὲρ τῆς ἁμαρτίας μου.
\vs{20}Οἱ δὲ ἐχθροί μου ζῶσι, καὶ κεκραταίωνται ὑπὲρ ἐμὲ, καὶ ἐπληθύνθησαν οἱ μισοῦντές με ἀδίκως.
\vs{21}Οἱ ἀνταποδιδόντες κακὰ ἀντὶ ἀγαθῶν, ἐνδιέβαλλόν με, ἐπεὶ κατεδίωκον δικαιοσύνην.
\vs{22}Μὴ ἐγκαταλίπῃς με Κύριε ὁ Θεός μου, μὴ ἀποστῇς ἀπʼ ἐμοῦ.
\vs{23}Πρόσχες εἰς τὴν βοήθειάν μου Κύριε τῆς σωτηρίας μου.

\begin{psalmheading}{\ch{38}{39} Εἰς τὸ τέλος, τῷ Ἰδιθοὺν ᾠδὴ τῷ Δαυίδ.}
\end{psalmheading}
\vs{2}Εἶπα, φυλάξω τὰς ὁδούς μου, τοῦ μὴ ἁμαρτάνειν ἐν γλώσσῃ μου· ἐθέμην τῷ στόματί μου φυλακὴν, ἐν τῷ συστῆναι τὸν ἁμαρτωλὸν ἐναντίον μου.
\vs{3}Ἐκωφώθην καὶ ἐταπεινώθην καὶ ἐσίγησα ἐξ ἀγαθῶν, καὶ τὸ ἄλγημά μου ἀνεκαινίσθη.
\vs{4}Ἐθερμάνθη ἡ καρδία μου ἐντός μου, καὶ ἐν τῇ μελέτῃ μου ἐκκαυθήσεται πῦρ· ἐλάλησα ἐν γλώσσῃ μου,

\vs{5}Γνώρισόν μοι Κύριε τὸ πέρας μου, καὶ τὸν ἀριθμὸν τῶν ἡμερῶν μου τίς ἐστιν, ἵνα γνῶ τί ὑστερῶ ἐγώ.
\vs{6}Ἰδοὺ παλαιὰς ἔθου τὰς ἡμέρας μου, καὶ ὑπόστασίς μου ὡσεὶ οὐθὲν ἐνώπιόν σου· πλὴν τὰ σύμπαντα ματαιότης, πᾶς ἄνθρωπος ζῶν· διάψαλμα.
\vs{7}Μέντοιγε ἐν εἰκόνι διαπορεύεται ἄνθρωπος, πλὴν μάτην ταράσσεται· θησαυρίζει, καὶ οὐ γινώσκει τίνι συνάξει αὐτά.

\vs{8}Καὶ νῦν τίς ἡ ὑπομονή μου; οὐχὶ ὁ Κύριος; καὶ ἡ ὑπόστασίς μου παρὰ σοί ἐστι· διάψαλμα.
\vs{9}Ἀπὸ πασῶν τῶν ἀνομιῶν μου ῥῦσαί με, ὄνειδος ἄφρονι ἔδωκάς με.
\vs{10}Ἐκωφώθην καὶ οὐκ ἤνοιξα τὸ στόμα μου, ὅτι σὺ εἶ ὁ ποιήσας με.
\vs{11}Ἀπόστησον ἀπʼ ἐμοῦ τὰς μάστιγάς σου· ἀπὸ τῆς ἰσχύος τῆς χειρός σου ἐγὼ ἐξέλιπον.
\vs{12}Ἐν ἐλεγμοῖς ὑπὲρ ἀνομίας ἐπαίδευσας ἄνθρωπον· καὶ ἐξέτηξας ὡς ἀράχνην τὴν ψυχὴν αὐτοῦ, πλὴν μάτην ταράσσεται πᾶς ἄνθρωπος· διάψαλμα.

\vs{13}Εἰσάκουσον τῆς προσευχῆς μου Κύριε καὶ τῆς δεήσεώς μου· ἐνώτισαι τῶν δακρύων μου· μὴ παρασιωπήσῃς, ὅτι πάροικος ἐγώ εἰμι ἐν τῇ γῇ καὶ παρεπίδημος, καθὼς πάντες οἱ πατέρες μου.
\vs{14}Ἄνες μοι ἵνα ἀναψύξω πρὸ τοῦ με ἀπελθεῖν, καὶ οὐκέτι μὴ ὑπάρξω.

\begin{psalmheading}{\ch{39}{40} Εἰς τὸ τέλος, τῷ Δαυὶδ ψαλμός.}
\end{psalmheading}
\vs{2}Ὑπομένων ὑπέμεινα τὸν Κύριον, καὶ προσέσχε μοι, καὶ εἰσήκουσε τῆς δεήσεώς μου.
\vs{3}Καὶ ἀνήγαγέ με ἐκ λάκκου ταλαιπωρίας, καὶ ἀπὸ πηλοῦ ἰλύος· καὶ ἔστησεν ἐπὶ πέτραν τοὺς πόδας μου, καὶ κατεύθυνε τὰ διαβήματά μου.
\vs{4}Καὶ ἐνέβαλεν εἰς τὸ στόμα μου ᾆσμα καινὸν, ὕμνον τῷ Θεῷ ἡμῶν· ὄψονται πολλοὶ καὶ φοβηθήσονται, καὶ ἐλπιοῦσιν ἐπὶ Κύριον.
\vs{5}Μακάριος ἀνὴρ, οὗ ἐστι τὸ ὄνομα Κυρίου ἐλπὶς αὐτοῦ, καὶ οὐκ ἐπέβλεψεν εἰς ματαιότητας καὶ μανίας ψευδεῖς.

\vs{6}Πολλὰ ἐποίησας σὺ Κύριε ὁ Θεός μου τὰ θαυμάσιά σου, καὶ τοῖς διαλογισμοῖς σου οὐκ ἔστι τίς ὁμοιωθήσεταί σοι· ἀπήγγειλα καὶ ἐλάλησα, ἐπληθύνθησαν ὑπὲρ ἀριθμόν.
\vs{7}Θυσίαν καὶ προσφορὰν οὐκ ἠθέλησας, σῶμα δὲ κατηρτίσω μοι· ὁλοκαύτωμα καὶ περὶ ἁμαρτίας οὐκ ᾔτησας.
\vs{8}Τότε εἶπον, ἰδοὺ ἥκω· ἐν κεφαλίδι βιβλίου γέγραπται περὶ ἐμοῦ,
\vs{9}τοῦ ποιῆσαι τὸ θέλημά σου ὁ Θεός μου ἠβουλήθην, καὶ τὸν νόμον σου ἐν μέσῳ τῆς καρδίας μου.
\vs{10}Εὐηγγελισάμην δικαιοσύνην ἐν ἐκκλησίᾳ μεγάλῃ, ἰδοὺ τὰ χείλη μου οὐ μὴ κωλύσω. Κύριε, σὺ ἔγνως
\vs{11}τὴν δικαιοσύνην μου, οὐκ ἔκρυψα ἐν τῇ καρδίᾳ μου τὴν ἀλήθειάν σου, καὶ τὸ σωτήριόν σου εἶπα· οὐκ ἔκρυψα τὸ ἔλεός σου καὶ τὴν ἀλήθειάν σου ἀπὸ συναγωγῆς πολλῆς.

\vs{12}Σὺ δέ Κύριε μὴ μακρύνῃς τοὺς οἰκτιρμούς σου ἀπʼ ἐμοῦ, τὸ ἔλεός σου καὶ ἡ ἀλήθειά σου διαπαντὸς ἀντελάβοντό μου.
\vs{13}Ὅτι περιέσχον με κακὰ, ὧν οὐκ ἔστιν ἀριθμὸς, κατέλαβόν με αἱ ἀνομίαι μου, καὶ οὐκ ἠδυνάσθην τοῦ βλέπειν· ἐπληθύνθησαν ὑπὲρ τὰς τρίχας τῆς κεφαλῆς μου, καὶ ἡ καρδία μου ἐγκατέλιπέ με.
\vs{14}Εὐδόκησον Κύριε τοῦ ῥύσασθαί με, Κύριε εἰς τὸ βοηθῆσαί μοι πρόσχες.
\vs{15}Καταισχυνθείησαν καὶ ἐντραπείησαν ἅμα οἱ ζητοῦντες τὴν ψυχήν μου, τοῦ ἐξάραι αὐτήν· ἀποστραφείησαν εἰς τὰ ὀπίσω, καὶ ἐντραπείησαν οἱ θέλοντές μοι κακά.
\vs{16}Κομισάσθωσαν παραχρῆμα αἰσχύνην αὐτῶν, οἱ λέγοντές μοι, εὖγε, εὖγε.
\vs{17}Ἀγαλλιάσαιντο καὶ εὐφρανθείησαν ἐπὶ σοὶ πάντες οἱ ζητοῦντές σε Κύριε· καὶ εἰπάτωσαν διαπαντὸς, μεγαλυνθήτω ὁ Κύριος, οἱ ἀγαπῶντες τὸ σωτήριόν σου.
\vs{18}Ἐγὼ δὲ πτωχὸς καὶ πένης εἰμὶ, Κύριος φροντιεῖ μου· βοηθός μου καὶ ὑπερασπιστής μου εἶ σὺ ὁ Θεός μου, μὴ χρονίσῃς.

\begin{psalmheading}{\ch{40}{41} Εἰς τὸ τέλος, ψαλμὸς τῷ Δαυίδ.}
\end{psalmheading}
\vs{2}Μακαρίος ὁ συνιῶν ἐπὶ πτωχὸν καὶ πένητα, ἐν ἡμέρᾳ πονηρᾷ ῥύσεται αὐτὸν ὁ Κύριος.
\vs{3}Κύριος φυλάξαι αὐτὸν καὶ ζήσαι αὐτὸν, καὶ μακαρίσαι αὐτὸν ἐν τῇ γῇ, καὶ μὴ παραδοῖ αὐτὸν εἰς χεῖρας ἐχθροῦ αὐτοῦ.
\vs{4}Κύριος βοηθήσαι αὐτῷ ἐπὶ κλίνης ὀδύνης αὐτοῦ, ὅλην τὴν κοίτην αὐτοῦ ἔστρεψας ἐν τῇ ἀῤῥωστίᾳ αὐτοῦ.

\vs{5}Ἐγὼ εἶπα, Κύριε ἐλέησόν με, ἴασαι τὴν ψυχήν μου, ὅτι ἥμαρτόν σοι.
\vs{6}Οἱ ἐχθροί μου εἶπαν κακά μοι, πότε ἀποθανεῖται καὶ ἀπολεῖται τὸ ὄνομα αὐτοῦ;
\vs{7}Καὶ εἰ εἰσεπορεύετο τοῦ ἰδεῖν, μάτην ἐλάλει ἡ καρδία αὐτοῦ, συνήγαγεν ἀνομίαν ἑαυτῷ, ἐξεπορεύετο ἔξω, καὶ ἐλάλει
\vs{8}ἐπὶ τὸ αὐτό. Κατʼ ἐμοῦ ἐψιθύριζον πάντες οἱ ἐχθροί μου, κατʼ ἐμοῦ ἐλογίζοντο κακά μοι.
\vs{9}Λόγον παράνομον κατέθεντο κατʼ ἐμοῦ, μὴ ὁ κοιμώμενος οὐχὶ προσθήσει τοῦ ἀναστῆναι;
\vs{10}Καὶ γὰρ ὁ ἄνθρωπος τῆς εἰρήνης μου ἐφʼ ὃν ἤλπισα, ὁ ἐσθίων ἄρτους μου ἐμεγάλυνεν ἐπʼ ἐμὲ πτερνισμόν.

\vs{11}Σὺ δὲ Κύριε ἐλέησόν με, καὶ ἀνάστησόν με, καὶ ἀνταποδώσω αὐτοῖς.
\vs{12}Ἐν τούτῳ ἔγνων, ὅτι τεθέληκάς με ὅτι οὐ μὴ ἐπιχαρῇ ὁ ἐχθρός μου ἐπʼ ἐμέ.
\vs{13}Ἐμοῦ δὲ διὰ τὴν ἀκακίαν ἀντελάβου, καὶ ἐβεβαίωσάς με ἐνώπιόν σου εἰς τὸν αἰῶνα.
\vs{14}Εὐλογητὸς Κύριος ὁ Θεὸς Ἰσραὴλ ἀπὸ τοῦ αἰῶνος καὶ εἰς τὸν αἰῶνα· γένοιτο, γένοιτο.

\begin{psalmheading}{\ch{41}{42} Εἰς τὸ τέλος, εἰς σύνεσιν τοῖς υἱοῖς Κορέ.}
\end{psalmheading}
\vs{2}Ὃν τρόπον ἐπιποθεῖ ἡ ἔλαφος ἐπὶ τὰς πηγὰς τῶν ὑδάτων, οὕτως ἐπιποθεῖ ἡ ψυχή μου πρὸς σὲ ὁ Θεός.
\vs{3}Ἐδίψησεν ἡ ψυχή μου πρὸς τὸν Θεὸν τὸν ζῶντα· πότε ἥξω καὶ ὀφθήσομαι τῷ προσώπῳ τοῦ Θεοῦ;
\vs{4}Ἐγενήθη τὰ δάκρυά μου ἐμοὶ ἄρτος ἡμέρας καὶ νυκτὸς, ἐν τῷ λέγεσθαί μοι καθʼ ἑκάστην ἡμέραν, ποῦ ἐστιν ὁ Θεός σου;
\vs{5}Ταῦτα ἐμνήσθην, καὶ ἐξέχεα ἐπʼ ἐμὲ τὴν ψυχήν μου, ὅτι διελεύσομαι ἐν τόπῳ σκηνῆς θαυμαστῆς ἕως τοῦ οἴκου τοῦ Θεοῦ· ἐν φωνῇ ἀγαλλιάσεως καὶ ἐξομολογήσεως ἤχου ἑορταζόντων.

\vs{6}Ἱνατί περίλυπος εἶ ἡ ψυχή μου, καὶ ἱνατί συνταράσσεις με; ἔλπισον ἐπὶ τὸν Θεὸν, ὅτι ἐξομολογήσομαι αὐτῷ, σωτήριον τοῦ προσώπου μου,

\vs{7}Ὁ Θεός μου· πρὸς ἐμαυτὸν ἡ ψυχή μου ἐταράχθη, διὰ τοῦτο μνησθήσομαί σου ἐκ γῆς Ἰορδάνου, καὶ Ἐρμωνιεὶμ ἀπὸ ὄρους μικροῦ.
\vs{8}Ἄβυσσος ἄβυσσον ἐπικαλεῖται εἰς φωνὴν τῶν καταῤῥακτῶν σου· πάντες οἱ μετεωρισμοί σου, καὶ τὰ κύματά σου ἐπʼ ἐμὲ διῆλθον.
\vs{9}Ἡμέρας ἐντελεῖται Κύριος τὸ ἔλεος αὐτοῦ, καὶ νυκτὸς δηλώσει· παρʼ ἐμοὶ προσευχὴ τῷ Θεῷ τῆς ζωῆς μου,
\vs{10}ἐρῶ τῷ Θεῷ, ἀντιλήπτωρ μου εἶ, διατί μου ἐπελάθου; ἱνατί σκυθρωπάζων πορεύομαι ἐν τῷ ἐκθλίβειν τὸν ἐχθρόν μου;
\vs{11}Ἐν τῷ καταθλᾶσθαι τὰ ὀστᾶ μου, ὠνείδισάν με οἱ θλίβοντές με· ἐν τῷ λέγειν αὐτούς μοι καθʼ ἑκάστην ἡμέραν, ποῦ ἐστιν ὁ Θεός σου;

\vs{12}Ἱνατί περίλυπος εἶ ἡ ψυχή μου, καὶ ἱνατί συνταράσσεις με; ἔλπισον ἐπὶ τὸν Θεὸν, ὅτι ἐξομολογήσομαι αὐτῷ, ἡ σωτηρία τοῦ προσώπου μου, καὶ ὁ Θεός μου.

\begin{psalmheading}{\ch{42}{43} Ψαλμὸς τῷ Δαυίδ.}
\end{psalmheading}
Κρίνον με ὁ Θεὸς, καὶ δίκασον τὴν δίκην μου, ἐξ ἔθνους οὐχ ὁσίου, ἀπὸ ἀνθρώπου ἀδίκου καὶ δολίου ῥῦσαί με.
\vs{2}Ὅτι σὺ εἶ ὁ Θεὸς κραταίωμά μου, ἱνατί ἀπώσω με; καὶ ἱνατί σκυθρωπάζων πορεύομαι ἐν τῷ ἐκθλίβειν τὸν ἐχθρόν μου;
\vs{3}Ἐξαπόστειλον τὸ φῶς σου καὶ τὴν ἀλήθειάν σου, αὐτά με ὡδήγησαν καὶ ἤγαγόν με εἰς ὄρος ἅγιόν σου, καὶ εἰς τὰ σκηνώματά σου.
\vs{4}Καὶ εἰσελεύσομαι πρὸς τὸ θυσιαστήριον τοῦ Θεοῦ, πρὸς τὸν Θεὸν τὸν εὐφραίνοντα τὴν νεότητά μου, ἐξομολογήσομαί σοι ἐν κιθάρᾳ ὁ Θεὸς ὁ Θεός μου.

\vs{5}Ἱνατί περίλυπος εἶ ἡ ψυχή μου, καὶ ἱνατί συνταράσσεις με; ἔλπισον ἐπὶ τὸν Θεὸν, ὅτι ἐξομολογήσομαι αὐτῷ, σωτήριον τοῦ προσώπου μου, ὁ Θεός μου.

\begin{psalmheading}{\ch{43}{44} Εἰς τὸ τέλος, τοῖς υἱοῖς Κορὲ εἰς σύνεσιν ψαλμός.}
\end{psalmheading}
\vs{2}Ὁ Θεὸς ἐν τοῖς ὠσὶν ἡμῶν ἠκούσαμεν, οἱ πατέρες ἡμῶν ἀνήγγειλαν ἡμῖν, ἔργον ὃ εἰργάσω ἐν ταῖς ἡμέραις αὐτῶν, ἐν ἡμεραίς ἀρχαίαις.
\vs{3}Ἡ χείρ σου ἔθνη ἐξωλόθρευσε, καὶ κατεφύτευσας αὐτοὺς, ἐκάκωσας λαοὺς καὶ ἐξέβαλες αὐτούς.
\vs{4}Οὐ γὰρ ἐν τῇ ῥομφαίᾳ αὐτῶν ἐκληρονόμησαν γῆν, καὶ ὁ βραχίων αὐτῶν οὐκ ἔσωσεν αὐτοὺς, ἀλλὰ ἡ δεξιά σου καὶ ὁ βραχίων σου, καὶ ὁ φωτισμὸς τοῦ προσώπου σου, ὅτι εὐδόκησας ἐν αὐτοῖς.

\vs{5}Σὺ εἶ αὐτὸς ὁ βασιλεύς μου καὶ ὁ Θεός μου, ὁ ἐντελλόμενος τὰς σωτηρίας Ἰακώβ.
\vs{6}Ἐν σοὶ τοὺς ἐχθροὺς ἡμῶν κερατιοῦμεν, καὶ ἐν τῷ ὀνόματί σου ἐξουθενώσομεν τοὺς ἐπανισταμένους ἡμῖν.
\vs{7}Οὐ γὰρ ἐπὶ τῷ τόξῳ μου ἐλπιῶ, καὶ ἡ ῥομφαία μου οὐ σώσει με.
\vs{8}Ἔσωσας γὰρ ἡμᾶς ἐκ τῶν θλιβόντων ἡμᾶς, καὶ τοὺς μισοῦντας ἡμᾶς κατῄσχυνας.
\vs{9}Ἐν τῷ Θεῷ ἐπαινεθησόμεθα ὅλην τὴν ἡμέραν, καὶ ἐν τῷ ὀνόματί σου ἐξομολογησόμεθα εἰς τὸν αἰῶνα. διάψαλμα.

\vs{10}Νυνὶ δὲ ἀπώσω καὶ κατῄσχυνας ἡμᾶς, καὶ οὐκ ἐξελεύσῃ ἐν ταῖς δυνάμεσιν ἡμῶν.
\vs{11}Ἀπέστρεψας ἡμᾶς εἰς τὰ ὀπίσω παρὰ τοὺς ἐχθροὺς ἡμῶν, καὶ οἱ μισοῦντες ἡμᾶς διήρπαζον ἑαυτοῖς.
\vs{12}Ἔδωκας ἡμᾶς ὡς πρόβατα βρώσεως, καὶ ἐν τοῖς ἔθνεσι διέσπειρας ἡμᾶς·
\vs{13}Ἀπέδου τὸν λαόν σου ἄνευ τιμῆς, καὶ οὐκ ἦν πλῆθος ἐν τοῖς ἀλαλάγμασιν αὐτῶν.
\vs{14}Ἔθου ἡμᾶς ὄνειδος τοῖς γείτοσιν ἡμῶν, μυκτηρισμὸν καὶ καταγέλωτα τοῖς κύκλῳ ἡμῶν.
\vs{15}Ἔθου ἡμᾶς εἰς παραβολὴν ἐν τοῖς ἔθνεσι, κίνησιν κεφαλῆς ἐν τοῖς λαοῖς.
\vs{16}Ὅλην τὴν ἡμέραν ἡ ἐντροπή μου κατεναντίον μου ἐστὶ, καὶ ἡ αἰσχύνη τοῦ προσώπου μου ἐκάλυψέ με,
\vs{17}ἀπὸ φωνῆς ὀνειδίζοντος καὶ παραλαλοῦντος, ἀπὸ προσώπου ἐχθροῦ καὶ ἐκδιώκοντος.

\vs{18}Ταῦτα πάντα ἦλθεν ἐφʼ ἡμᾶς, καὶ οὐκ ἐπελαθόμεθά σου, καὶ οὐκ ἠδικήσαμεν ἐν διαθήκῃ σου.
\vs{19}Καὶ οὐκ ἀπέστη εἰς τὰ ὀπίσω ἡ καρδία ἡμῶν, καὶ ἐξέκλινας τὰς τρίβους ἡμῶν ἀπὸ τῆς ὁδοῦ σου.
\vs{20}Ὅτι ἐταπείνωσας ἡμᾶς ἐν τόπῳ κακώσεως, καὶ ἐπεκάλυψεν ἡμᾶς σκιὰ θανάτου.
\vs{21}Εἰ ἐπελαθόμεθα τοῦ ὀνόματος τοῦ Θεοῦ ἡμῶν, καὶ εἰ διεπετάσαμεν χεῖρας ἡμῶν πρὸς θεὸν ἀλλότριον,
\vs{22}οὐχὶ ὁ Θεὸς ἐκζητήσει ταῦτα; αὐτὸς γὰρ γινώσκει τὰ κρύφια τῆς καρδίας.
\vs{23}Ὅτι ἕνεκά σου θανατούμεθα ὅλην τὴν ἡμέραν, ἐλογίσθημεν ὡς πρόβατα σφαγῆς.

\vs{24}Ἐξεγέρθητι, ἱνατί ὑπνοῖς Κύριε; ἀνάστηθι, καὶ μὴ ἀπώσῃ εἰς τέλος.
\vs{25}Ἱνατί τὸ πρόσωπόν σου ἀποστρέφεις; ἐπιλανθάνῃ τῆς πτωχείας ἡμῶν καὶ τῆς θλίψεως ἡμῶν;
\vs{26}Ὅτι ἐταπεινώθη εἰς χοῦν ἡ ψυχὴ ἡμῶν, ἐκολλήθη εἰς γῆν ἡ γαστὴρ ἡμῶν.
\vs{27}Ἀνάστα Κύριε, βοήθησον ἡμῖν, καὶ λύτρωσαι ἡμᾶς ἕνεκεν τοῦ ὀνόματός σου.

\begin{psalmheading}{\ch{44}{45} Εἰς τὸ τέλος, ὑπὲρ τῶν ἀλλοιωθησομένων τοῖς υἱοῖς Κορὲ εἰς σύνεσιν, ᾠδὴ ὑπὲρ τοῦ ἀγαπητοῦ.}
\end{psalmheading}
\vs{2}Ἐξηρεύξατο ἡ καρδία μου λόγον ἀγαθὸν, λέγω ἐγὼ τὰ ἔργα μου τῷ βασιλεῖ· ἡ γλῶσσά μου κάλαμος γραμματέως ὀξυγράφου.
\vs{3}Ὡραῖος κάλλει παρὰ τοὺς υἱοὺς τῶν ἀνθρώπων, ἐξεχύθη χάρις ἐν χείλεσί σου, διὰ τοῦτο εὐλόγησέ σε ὁ Θεὸς εἰς τὸν αἰῶνα.

\vs{4}Περίζωσαι τὴν ῥομφαίαν σου ἐπὶ τὸν μηρόν σου δυνατέ· τῇ ὡραιότητί σου, καὶ τῷ κάλλει σου,
\vs{5}καὶ ἔντεινον, καὶ κατευοδοῦ, καὶ βασίλευε· ἕνεκεν ἀληθείας καὶ πρᾳΰτητος καὶ δικαιοσύνης, καὶ ὁδηγήσει σε θαυμαστῶς ἡ δεξιά σου.
\vs{6}Τὰ βέλη σου ἠκονημένα δυνατὲ, λαοὶ ὑποκάτω σου πεσοῦνται, ἐν καρδίᾳ τῶν ἐχθρῶν τοῦ βασιλέως.

\vs{7}Ὁ θρόνος σου ὁ Θεὸς εἰς αἰῶνα αἰῶνος, ῥάβδος εὐθύτητος ἡ ῥάβδος τῆς βασιλείας σου.
\vs{8}Ἠγάπησας δικαιοσύνην, καὶ ἐμίσησας ἀνομίαν, διὰ τοῦτο ἔχρισέ σε ὁ Θεὸς ὁ Θεός σου ἔλαιον ἀγαλλιάσεως παρὰ τοὺς μετόχους σου.

\vs{9}Σμύρνα καὶ στακτὴ καὶ κασία ἀπὸ τῶν ἱματίων σου, ἀπὸ βάρεων ἐλεφαντίνων, ἐξ ὧν ηὔφρανάν σε
\vs{10}θυγατέρες βασιλέων ἐν τῇ τιμῇ σου· παρέστη ἡ βασίλισσα ἐκ δεξιῶν σου, ἐν ἱματισμῷ διαχρύσῳ περιβεβλημένη, πεποικιλμένη.
\vs{11}Ἄκουσον θύγατερ καὶ ἴδε καὶ κλῖνον τὸ οὖς σου, καὶ ἐπιλάθου τοῦ λαοῦ σου, καὶ τοῦ οἴκου τοῦ πατρός σου.
\vs{12}Ὅτι ἐπεθύμησεν ὁ βασιλεὺς τοῦ κάλλους σου, ὅτι αὐτός ἐστιν ὁ Κύριός σου.
\vs{13}Καὶ προσκυνήσουσιν αὐτῷ θυγατέρες Τύρου ἐν δώροις, τὸ πρόσωπόν σου λιτανεῦσουσιν οἱ πλούσιοι τοῦ λαοῦ τῆς γῆς.

\vs{14}Πᾶσα ἡ δόξα αὐτῆς θυγατρὸς τοῦ βασιλέως Ἐσεβὼν, ἐν κροσσωτοῖς χρυσοῖς περιβεβλημένη, πεποικιλμένη·
\vs{15}ἀπενεχθήσονται τῷ βασιλεῖ παρθένοι ὀπίσω αὐτῆς, αἱ πλησίον αὐτῆς ἀπενεχθήσονταί σοι.
\vs{16}Ἀπενεχθήσονται ἐν εὐφροσύνῃ καὶ ἀγαλλιάσει, ἀχθήσονται εἰς ναὸν βασιλέως.
\vs{17}Ἀντὶ τῶν πατέρων σου ἐγενήθησάν σοι υἱοὶ, καταστήσεις αὐτοὺς ἄρχοντας ἐπὶ πᾶσαν τὴν γῆν.
\vs{18}Μνησθήσονται τοῦ ὀνόματός σου ἐν πάσῃ γενεᾷ καὶ γενεᾷ, διὰ τοῦτο λαοὶ ἐξομολογήσονταί σοι εἰς τὸν αἰῶνα, καὶ εἰς τὸν αἰῶνα τοῦ αἰῶνος·

\begin{psalmheading}{\ch{45}{46} Εἰς τὸ τέλος, ὑπὲρ τῶν υἱῶν Κορὲ, ὑπὲρ τῶν κρυφίων ψαλμός.}
\end{psalmheading}
\vs{2}Ὁ Θεὸς ἡμῶν καταφυγὴ καὶ δύναμις, βοηθὸς ἐν θλίψεσι ταῖς εὑρούσαις ἡμᾶς σφόδρα.
\vs{3}Διὰ τοῦτο οὐ φοβηθησόμεθα ἐν τῷ ταράσσεσθαι τὴν γῆν, καὶ μετατίθεσθαι ὄρη ἐν καρδίαις θαλασσῶν.
\vs{4}Ἤχησαν καὶ ἐταράχθησαν τὰ ὕδατα αὐτῶν, ἐταράχθησαν τὰ ὄρη ἐν τῇ κραταιότητι αὐτοῦ. διάψαλμα.
\vs{5}Τοῦ ποταμοῦ τὰ ὁρμήματα εὐφραίνουσι τὴν πόλιν τοῦ Θεοῦ, ἡγίασε τὸ σκήνωμα αὐτοῦ ὁ ὕψιστος.
\vs{6}Ὁ Θεὸς ἐν μέσῳ αὐτῆς οὐ σαλευθήσεται, βοηθήσει αὐτῇ ὁ Θεὸς τῷ προσώπῳ.
\vs{7}Ἐταράχθησαν ἔθνη, ἔκλιναν βασιλεῖαι, ἔδωκε φωνὴν αὐτοῦ, ἐσαλεύθη ἡ γῆ.
\vs{8}Κύριος τῶν δυνάμεων μεθʼ ἡμῶν, ἀντιλήπτωρ ἡμῶν ὁ Θεὸς Ἰακώβ. διάψαλμα.

\vs{9}Δεῦτε καὶ ἴδετε τὰ ἔργα τοῦ Κυρίου, ἃ ἔθετο τέρατα ἐπὶ τῆς γῆς·
\vs{10}ἀνταναιρῶν πολέμους μέχρι τῶν περάτων τῆς γῆς, τόξον συντρίψει, καὶ συγκλάσει ὅπλον, καὶ θυρεοὺς κατακαύσει ἐν πυρί.
\vs{11}Σχολάσατε καὶ γνῶτε, ὅτι ἐγώ εἰμι ὁ Θεὸς, ὑψωθήσομαι ἐν τοῖς ἔθνεσιν, ὑψωθήσομαι ἐν τῇ γῇ·
\vs{12}Κύριος τῶν δυνάμεων μεθʼ ἡμῶν, ἀντιλήπτωρ ἡμῶν ὁ Θεὸς Ἰακώβ.

\begin{psalmheading}{\ch{46}{47} Εἰς τὸ τέλος, ὑπὲρ τῶν υἱῶν Κορὲ ψαλμός.}
\end{psalmheading}
\vs{2}Παντὰ τὰ ἔθνη κροτήσατε χεῖρας, ἀλαλάξατε τῷ Θεῷ ἐν φωνῇ ἀγαλλιάσεως.
\vs{3}Ὅτι Κύριος ὕψιστος, φοβερὸς, βασιλεὺς μέγας ἐπὶ πᾶσαν τὴν γῆν.
\vs{4}Ὑπέταξε λαοὺς ἡμῖν, καὶ ἔθνη ὑπὸ τοὺς πόδας ἡμῶν·
\vs{5}Ἐξελέξατο ἡμῖν τὴν κληρονομίαν αὐτοῦ, τὴν καλλονὴν Ἰακὼβ, ἣν ἠγάπησε· διάψαλμα.

\vs{6}Ἀνέβη ὁ Θεὸς ἐν ἀλαλαγμῷ, Κύριος ἐν φωνῇ σάλπιγγος.
\vs{7}Ψάλατε τῷ Θεῷ ἡμῶν, ψάλατε· ψάλατε τῷ βασιλεῖ ἡμῶν, ψάλατε.
\vs{8}Ὅτι βασιλεὺς πάσης τῆς γῆς ὁ Θεὸς, ψάλατε συνετῶς.
\vs{9}Ἐβασίλευσεν ὁ Θεὸς ἐπὶ τὰ ἔθνη, ὁ Θεὸς κάθηται ἐπὶ θρόνου ἁγίου αὐτοῦ.
\vs{10}Ἄρχοντες λαῶν συνήχθησαν μετὰ τοῦ Θεοῦ Ἀβραὰμ, ὅτι τοῦ Θεοῦ οἱ κραταιοὶ τῆς γῆς σφόδρα ἐπῄρθησαν.

\begin{psalmheading}{\ch{47}{48} Ψαλμὸς ᾠδῆς τοῖς υἱοῖς Κορὲ δευτέρᾳ σαββάτου.}
\end{psalmheading}
\vs{2}Μεγὰς Κύριος, καὶ αἰνετὸς σφόδρα ἐν πόλει τοῦ Θεοῦ ἡμῶν, ἐν ὄρει ἁγίῳ αὐτοῦ.
\vs{3}Εὐρίζων, ἀγαλλιάματι πάσης τῆς γῆς, ὄρη Σιὼν τὰ πλευρὰ τοῦ βοῤῥᾶ, ἡ πόλις τοῦ βασιλέως τοῦ μεγάλου.
\vs{4}Ὁ Θεὸς ἐν ταῖς βάρεσιν αὐτῆς γινώσκεται, ὅταν ἀντιλαμβάνηται αὐτῆς.

\vs{5}Ὅτι ἰδοὺ οἱ βασιλεῖς τῆς γῆς συνήχθησαν, ἤλθοσαν ἐπὶ τὸ αὐτό.
\vs{6}Αὐτοὶ ἰδόντες οὕτως ἐθαύμασαν, ἐταράχθησαν, ἐσαλεύθησαν,
\vs{7}τρόμος ἐπελάβετο αὐτῶν· ἐκεῖ ὠδῖνες ὡς τικτούσης.
\vs{8}Ἐν πνεύματι βιαίῳ συντρίψεις πλοῖα Θάρσις.
\vs{9}Καθάπερ ἠκούσαμεν, οὕτως καὶ εἴδομεν, ἐν πόλει Κυρίου τῶν δυνάμεων, ἐν πόλει τοῦ Θεοῦ ἡμῶν, ὁ Θεὸς ἐθεμελίωσεν αὐτὴν εἰς τὸν αἰῶνα· διάψαλμα.

\vs{10}Ὑπελάβομεν ὁ Θεὸς τὸ ἔλεός σου ἐν μέσῳ τοῦ λαοῦ σου.
\vs{11}Κατὰ τὸ ὄνομά σου ὁ Θεὸς, οὕτως καὶ ἡ αἴνεσίς σου ἐπὶ τὰ πέρατα τῆς γῆς, δικαιοσύνης πλήρης ἡ δεξιά σου.
\vs{12}Εὐφρανθήτω τὸ ὄρος Σιὼν, ἀγαλλιάσθωσαν αἱ θυγατέρες τῆς Ἰουδαίας ἕνεκα τῶν κριμάτων σου Κύριε.

\vs{13}Κυκλώσατε Σιὼν, καὶ περιλάβετε αὐτὴν, διηγήσασθε ἐν τοῖς πύργοις αὐτῆς.
\vs{14}Θέσθε τὰς καρδίας ὑμῶν εἰς τὴν δύναμιν αὐτῆς, καὶ καταδιέλεσθε τὰς βάρεις αὐτῆς, ὅπως ἂν διηγήσησθε εἰς γενεὰν ἑτέραν.
\vs{15}Ὅτι οὗτός ἐστιν ὁ Θεὸς ἡμῶν εἰς τὸν αἰῶνα καὶ εἰς τὸν αἰῶνα τοῦ αἰῶνος, αὐτὸς ποιμανεῖ ἡμᾶς εἰς τοὺς αἰῶνας.

\begin{psalmheading}{\ch{48}{49} Εἰς τὸ τέλος, τοῖς υἱοῖς Κορὲ ψαλμός.}
\end{psalmheading}
\vs{2}Ἀκούσατε ταῦτα πάντα τὰ ἔθνη, ἐνωτίσασθε πάντες οἱ κατοικοῦντες τὴν οἰκουμένην,
\vs{3}οἵ τε γηγενεῖς καὶ οἱ υἱοὶ τῶν ἀνθρώπων, ἐπιτοαυτὸ πλούσιος καὶ πένης.
\vs{4}Τὸ στόμα μου λαλήσει σοφίαν, καὶ ἡ μελέτη τῆς καρδίας μου σύνεσιν.
\vs{5}Κλινῶ εἰς παραβολὴν τὸ οὖς μου, ἀνοίξω ἐν ψαλτηρίῳ τὸ πρόβλημά μου.

\vs{6}Ἱνατί φοβοῦμαι ἐν ἡμέρᾳ πονηρᾷ; ἡ ἀνομία τῆς πτέρνης μου κυκλώσει με.
\vs{7}Οἱ πεποιθότες ἐπὶ τῇ δυνάμει αὐτῶν, καὶ ἐπὶ τῷ πλήθει τοῦ πλούτου αὐτῶν καυχώμενοι.
\vs{8}Ἀδελφὸς οὐ λυτροῦται, λυτρώσεται ἄνθρωπος; οὐ δώσει τῷ Θεῷ ἐξίλασμα ἑαυτοῦ,
\vs{9}καὶ τὴν τιμὴν τῆς λυτρώσεως τῆς ψυχῆς αὐτοῦ· καὶ ἐκοπίασεν εἰς τὸν αἰῶνα,
\vs{10}καὶ ζήσεται εἰς τέλος· ὅτι οὐκ ὄψεται καταφθορὰν,

\vs{11}Ὅταν ἴδῃ σοφοὺς ἀποθνήσκοντας, ἐπιτοαυτὸ ἄφρων καὶ ἄνους ἀπολοῦνται, καὶ καταλείψουσιν ἀλλοτρίοις τὸν πλοῦτον αὐτῶν.
\vs{12}Καὶ οἱ τάφοι αὐτῶν οἰκίαι αὐτῶν εἰς τὸν αἰῶνα, σκηνώματα αὐτῶν εἰς γενεὰν καὶ γενεὰν, ἐπεκαλέσαντο τὰ ὀνόματα αὐτῶν ἐπὶ τῶν γαιῶν αὐτῶν.
\vs{13}Καὶ ἄνθρωπος ἐν τιμῇ ὢν, οὐ συνῆκε, παρασυνεβλήθη τοῖς κτήνεσι τοῖς ἀνοήτοις, καὶ ὡμοιώθη αὐτοῖς.
\vs{14}Αὕτη ἡ ὁδὸς αὐτῶν σκάνδαλον αὐτοῖς, καὶ μετὰ ταῦτα ἐν τῷ στόματι αὐτῶν εὐλογήσουσι· διάψαλμα.
\vs{15}Ὡς πρόβατα ἐν ᾅδῃ ἔθεντο, θάνατος ποιμανεῖ αὐτούς· καὶ κατακυριεύσουσιν αὐτῶν οἱ εὐθεῖς τοπρωῒ, καὶ ἡ βοήθεια αὐτῶν παλαιωθήσεται ἐν τῷ ᾅδῃ ἐκ τῆς δόξης αὐτῶν.
\vs{16}Πλὴν ὁ Θεὸς λυτρώσεται τὴν ψυχήν μου ἐκ χειρὸς ᾅδου, ὅταν λαμβάνῃ με· διάψαλμα.

\vs{17}Μὴ φοβοῦ ὅταν πλουτήσῃ ἄνθρωπος, καὶ ὅταν πληθυνθῇ ἡ δόξα τοῦ οἴκου αὐτοῦ.
\vs{18}Ὅτι οὐκ ἐν τῷ ἀποθνήσκειν αὐτὸν λήψεται τὰ πάντα, οὐδὲ συγκαταβήσεται αὐτῷ ἡ δόξα αὐτοῦ.
\vs{19}Ὅτι ἡ ψυχὴ αὐτοῦ ἐν τῇ ζωῇ αὐτοῦ εὐλογηθήσεται, ἐξομολογήσεταί σοι ὅταν ἀγαθύνῃς αὐτῷ.
\vs{20}Εἰσελεύσεται ἕως γενεᾶς πατέρων αὐτοῦ, ἕως αἰῶνος οὐκ ὄψεται φῶς.
\vs{21}Ἄνθρωπος ἐν τιμῇ ὢν, οὐ συνῆκε, παρασυνεβλήθη τοῖς κτήνεσι τοῖς ἀνοήτοις, καὶ ὡμοιώθη αὐτοῖς.

\begin{psalmheading}{\ch{49}{50} Ψαλμὸς τῷ Ἀσάφ.}
\end{psalmheading}
Θεὸς θεῶν Κύριος ἐλάλησε, καὶ ἐκάλεσε τὴν γῆν ἀπὸ ἀνατολῶν ἡλίου μέχρι δυσμῶν.
\vs{2}Ἐκ Σιὼν ἡ εὐπρέπεια τῆς ὡραιότητος αὐτοῦ.
\vs{3}Ὁ Θεὸς ἐμφανῶς ἥξει, ὁ Θεὸς ἡμῶν, καὶ οὐ παρασιωπήσεται· πῦρ ἐναντίον αὐτοῦ καυθήσεται, καὶ κύκλῳ αὐτοῦ καταιγὶς σφοδρά.
\vs{4}Προσκαλέσεται τὸν οὐρανὸν ἄνω, καὶ τὴν γῆν διακρῖναι τὸν λαὸν αὐτοῦ.
\vs{5}Συναγάγετε αὐτῷ τοὺς ὁσίους αὐτοῦ, τοὺς διατιθεμένους τὴν διαθήκην αὐτοῦ ἐπὶ θυσίαις.
\vs{6}Καὶ ἀναγγελοῦσιν οἱ οὐρανοὶ τὴν δικαιοσύνην αὐτοῦ, ὅτι Θεὸς κριτής ἐστι· διάψαλμα.

\vs{7}Ἄκουσον λαός μου καὶ λαλήσω σοι, Ἰσραὴλ, καὶ διαμαρτύρομαί σοι· ὁ Θεὸς ὁ Θεός σου εἰμὶ ἐγώ.
\vs{8}Οὐκ ἐπὶ ταῖς θυσίαις σου ἐλέγξω σε, τὰ δὲ ὁλοκαυτώματά σου ἐνώπιόν μου ἐστὶ διαπαντός.
\vs{9}Οὐ δέξομαι ἐκ τοῦ οἴκου σου μόσχους, οὐδὲ ἐκ τῶν ποιμνίων σου χιυάρους·
\vs{10}Ὅτι ἐμά ἐστι πάντα τὰ θηρία τοῦ δρυμοῦ, κτήνη ἐν τοῖς ὄρεσι, καὶ βόες.
\vs{11}Ἔγνωκα πάντα τὰ πετεινὰ τοῦ οὐρανοῦ, καὶ ὡραιότης ἀγροῦ μετʼ ἐμοῦ ἐστιν.
\vs{12}Ἐὰν πεινάσω, οὐ μή σοι εἴπω, ἐμὴ γάρ ἐστιν ἡ οἰκουμένη καὶ τὸ πλήρωμα αὐτῆς.
\vs{13}Μὴ φάγομαι κρέα ταύρων, ἢ αἷμα τράγων πίομαι;
\vs{14}Θῦσον τῷ Θεῷ θυσίαν αἰνέσεως, καὶ ἀπόδος τῷ ὑψίστῳ τὰς εὐχάς σου.
\vs{15}Καὶ ἐπικάλεσαί με ἐν ἡμέρᾳ θλίψεως, καὶ ἐξελοῦμαί σε, καὶ δοξάσεις με. διάψαλμα.

\vs{16}Τῷ δὲ ἁμαρτωλῷ εἶπεν ὁ Θεὸς, ἱνατί σὺ διηγῇ τὰ δικαιώματά μου, καὶ ἀναλαμβάνεις τὴν διαθήκην μου διὰ στόματός σου;
\vs{17}Σὺ δὲ ἐμίσησας παιδείαν, καὶ ἐξέβαλες τοὺς λόγους μου εἰς τὰ ὀπίσω.
\vs{18}Εἰ ἐθεώρεις κλέπτην, συνέτρεχες αὐτῷ, καὶ μετὰ μοιχῶν τὴν μερίδα σου ἐτίθεις.
\vs{19}Τὸ στόμα σου ἐπλεόνασε κακίαν, καὶ ἡ γλῶσσά σου περιέπλεκε δολιότητα.
\vs{20}Καθήμενος κατὰ τοῦ ἀδελφοῦ σου κατελάλεις, καὶ κατὰ τοῦ υἱοῦ τῆς μητρός σου ἐτίθεις σκάνδαλον.

\vs{21}Ταῦτα ἐποίησας, καὶ ἐσίγησα, ὑπέλαβες ἀνομίαν ὅτι ἔσομαί σοι ὅμοιος· ἐλέγξω σε, καὶ παραστήσω κατὰ πρόσωπόν σου.
\vs{22}Σύνετε δὴ ταῦτα οἱ ἐπιλανθανόμενοι τοῦ Θεοῦ, μήποτε ἁρπάσῃ, καὶ μὴ ᾖ ὁ ῥυόμενος.

\vs{23}Θυσία αἰνέσεως δοξάσει με, καὶ ἐκεῖ ὁδὸς ᾗ δείξω αὐτῷ τὸ σωτήριον Θεοῦ.

\begin{psalmheading}{\ch{50}{51} Εἰς τὸ τέλος, ψαλμὸς τῷ Δαυὶδ,}\vs{2}ἐν τῷ ἐλθεῖν πρὸς αὐτὸν Νάθαν τὸν προφήτην, ἡνίκα εἰσῆλθε πρὸς Βηρσαβεέ.
\end{psalmheading}

\vs{3}Ἐλέησον με ὁ Θεὸς κατὰ τὸ μέγα ἔλεός σου, καὶ κατὰ τὸ πλῆθος τῶν οἰκτιρμῶν σου ἐξάλειψον τὸ ἀνόμημά μου.
\vs{4}Ἐπιπλεῖον πλῦνόν με ἀπὸ τῆς ἀνομίας μου, καὶ ἀπὸ τῆς ἁμαρτίας μου καθάρισόν με.

\vs{5}Ὅτι τὴν ἀνομίαν μου ἐγὼ γινώσκω, καὶ ἡ ἁμαρτία μου ἐνώπιόν μου ἐστὶ διαπαντός·
\vs{6}Σοὶ μόνῳ ἥμαρτον, καὶ τὸ πονηρὸν ἐνώπιόν σου ἐποίησα· ὅπως ἂν δικαιωθῇς ἐν τοῖς λόγοις σου, καὶ νικήσῃς ἐν τῷ κρίνεσθαί σε.
\vs{7}Ἰδοὺ γὰρ ἐν ἀνομίαις συνελήφθην, καὶ ἐν ἁμαρτίαις ἐκίσσησέ με ἡ μήτηρ μου.

\vs{8}Ἰδοὺ γὰρ ἀλήθειαν ἠγάπησας, τὰ ἄδηλα καὶ τὰ κρύφια τῆς σοφίας σου ἐδήλωσάς μοι.
\vs{9}Ῥαντιεῖς με ὑσσώπῳ καὶ καθαρισθήσομαι, πλυνεῖς με καὶ ὑπὲρ χιόνα λευκανθήσομαι.
\vs{10}Ἀκουτιεῖς με ἀγαλλίασιν καὶ εὐφροσύνην, ἀγαλλιάσονται ὀστᾶ τεταπεινωμένα.
\vs{11}Ἀπόστρεψον τὸ πρόσωπόν σου ἀπὸ τῶν ἁμαρτιῶν μου, καὶ πάσας τὰς ἀνομίας μου ἐξάλειψον.
\vs{12}Καρδίαν καθαρὰν κτίσον ἐν ἐμοὶ ὁ Θεὸς, καὶ πνεῦμα εὐθὲς ἐγκαίνισον ἐν τοῖς ἐγκάτοις μου.
\vs{13}Μὴ ἀποῤῥίψῃς με ἀπὸ τοῦ προσώπου σου, καὶ τὸ πνεῦμα τὸ ἅγιόν σου μὴ ἀντανέλῃς ἀπʼ ἐμοῦ.
\vs{14}Ἀπόδος μοι τὴν ἀγαλλίασιν τοῦ σωτηρίου σου, πνεύματι ἡγεμονικῷ στήριξόν με.

\vs{15}Διδάξω ἀνόμους τὰς ὁδούς σου, καὶ ἀσεβεῖς ἐπὶ σὲ ἐπιστρέψουσι.
\vs{16}Ῥῦσαί με ἐξ αἱμάτων ὁ Θεὸς, ὁ Θεὸς τῆς σωτηρίας μου, ἀγαλλιάσεται ἡ γλῶσσά μου τὴν δικαιοσύνην σου.
\vs{17}Κύριε, τὰ χείλη μου ἀνοίξεις, καὶ τὸ στόμα μου ἀναγγελεῖ τὴν αἴνεσίν σου.
\vs{18}Ὅτι εἰ ἠθέλησας θυσίαν, ἔδωκα ἄν· ὁλοκαυτώματα οὐκ εὐδοκήσεις.
\vs{19}Θυσία τῷ Θεῷ πνεῦμα συντετριμμένον, καρδίαν συντετριμμένην καὶ τεταπεινωμένην ὁ Θεὸς οὐκ ἐξουδενώσει.

\vs{20}Ἀγάθυνον, Κύριε, ἐν τῇ εὐδοκίᾳ σου τὴν Σιὼν, καὶ οἰκοδομηθήτω τὰ τείχη Ἱερουσαλήμ.
\vs{21}Τότε εὐδοκήσεις θυσίαν δικαιοσύνης, ἀναφορὰν, καὶ ὁλοκαυτώματα· τότε ἀνοίσουσιν ἐπὶ τὸ θυσιαστήριόν σου μόσχους.

\begin{psalmheading}{\ch{51}{52} Εἰς τὸ τέλος συνέσεως τῷ Δαυὶδ,}\vs{2}ἐν τῷ ἐλθεῖν Δωὴκ τὸν Ἰδουμαῖον, καὶ ἀναγγεῖλαι τῷ Σαοὺλ, καὶ εἰπεῖν αὐτῷ, ἦλθε Δαυὶδ εἰς τὸν οἶκον Ἀβιμέλεχ.
\end{psalmheading}

\vs{3}Τί ἐγκαυχᾷ ἐν κακίᾳ ὁ δυνατὸς ἀνομίαν; ὅλην τὴν ἡμέραν
\vs{4}ἀδικίαν ἐλογίσατο ἡ γλῶσσά σου· ὡσεὶ ξυρὸν ἠκονημένον ἐποίησας δόλον.
\vs{5}Ἠγάπησας κακίαν ὑπὲρ ἀγαθωσύνην, ἀδικίαν ὑπὲρ τὸ λαλῆσαι δικαιοσύνην· διάψαλμα.
\vs{6}Ἠγάπησας πάντα τὰ ῥήματα καταποντισμοῦ, γλῶσσαν δολίαν.

\vs{7}Διὰ τοῦτο ὁ Θεὸς καθέλοι σε εἰς τέλος, ἐκτίλαι σε καὶ μεταναστεύσαι σε ἀπὸ σκηνώματος, καὶ τὸ ῥίζωμά σου ἐκ γῆς ζώντων· διάψαλμα.
\vs{8}Καὶ ὄψονται δίκαιοι καὶ φοβηθήσονται, καὶ ἐπʼ αὐτὸν γελάσονται, καὶ ἐροῦσιν,
\vs{9}ἰδοὺ ἄνθρωπος ὃς οὐκ ἔθετο τὸν Θεὸν βοηθὸν αὐτοῦ, ἀλλʼ ἐπήλπισεν ἐπὶ τὸ πλῆθος τοῦ πλούτου αὐτοῦ, καὶ ἐνεδυναμώθη ἐπὶ τῇ ματαιότητι αὐτοῦ.

\vs{10}Ἐγὼ δὲ ὡσεὶ ἐλαία κατάκαρπος ἐν τῷ οἴκῳ τοῦ Θεοῦ, ἤλπισα ἐπὶ τὸ ἔλεος τοῦ Θεοῦ εἰς τὸν αἰῶνα καὶ εἰς τὸν αἰῶνα τοῦ αἰῶνος.
\vs{11}Ἐξομολογήσομαί σοι εἰς τὸν αἰῶνα, ὅτι ἐποίησας, καὶ ὑπομενῶ τὸ ὄνομά σου, ὅτι χρηστὸν ἐναντίον τῶν ὁσίων σου.

\begin{psalmheading}{\ch{52}{53} Εἰς τὸ τέλος, ὑπὲρ μαελὲθ συνέσεως τῷ Δαυίδ.}
\end{psalmheading}
\vs{2}Εἶπεν ἄφρων ἐν καρδίᾳ αὐτοῦ, οὐκ ἔστι Θεός· διέφθειραν, καὶ ἐβδελύχθησαν ἐν ἀνομίαις· οὐκ ἔστι ποιῶν ἀγαθόν.
\vs{3}Ὁ Θεὸς ἐκ τοῦ οὐρανοῦ διέκυψεν ἐπὶ τοὺς υἱοὺς τῶν ἀνθρώπων, τοῦ ἰδεῖν εἰ ἔστι συνιῶν, ἢ ἐκζητῶν τὸν Θεόν.
\vs{4}Πάντες ἐξέκλιναν, ἅμα ἠχρειώθησαν, οὐκ ἔστι ποιῶν ἀγαθὸν, οὐκ ἔστιν ἕως ἑνός.

\vs{5}Οὐχὶ γνώσονται πάντες οἱ ἐργαζόμενοι τὴν ἀνομίαν, οἱ κατεσθίοντες τὸν λαόν μου βρώσει ἄρτου; τὸν Θεὸν οὐκ ἐπεκαλέσαντο.
\vs{6}Ἐκεῖ ἐφοβήθησαν φόβον, οὗ οὐκ ἦν φόβος· ὅτι ὁ Θεὸς διεσκόρπισεν ὀστᾶ ἀνθρωπαρέσκων, κατῃσχύνθησαν, ὅτι ὁ Θεὸς ἐξουδένωσεν αὐτούς.
\vs{7}Τίς δώσει ἐκ Σιὼν τὸ σωτήριον τοῦ Ἰσραήλ; ἐν τῷ ἀποστρέψαι Κύριον τὴν αἰχμαλωσίαν τοῦ λαοῦ αὐτοῦ, ἀγαλλιάσεται Ἰακὼβ, καὶ εὐφρανθήσεται Ἰσραήλ.

\begin{psalmheading}{\ch{53}{54} Εἰς τὸ τέλος, ἐν ὕμνοις συνέσεως τῷ Δαυὶδ,}\vs{2}ἐν τῷ ἐλθεῖν τοὺς Ζειφαίους, καὶ εἰπεῖν τῷ Σαοὺλ, οὐκ ἰδοὺ Δαυὶδ κέκρυπται παρʼ ἡμῖν;
\end{psalmheading}

\vs{3}Ὁ Θεὸς ἐν τῷ ὀνόματί σου σῶσόν με, καὶ ἐν τῇ δυνάμει σου κρῖνόν με.
\vs{4}Ὁ Θεὸς εἰσάκουσον τῆς προσευχῆς μου, ἐνώτισαι τὰ ῥήματα τοῦ στόματός μου.
\vs{5}Ὅτι ἀλλότριοι ἐπανέστησαν ἐπʼ ἐμὲ, καὶ κραταιοὶ ἐζήτησαν τὴν ψυχήν μου, οὐ προέθεντο τὸν Θεὸν ἐνώπιον αὐτῶν· διάψαλμα.

\vs{6}Ἰδοὺ γὰρ ὁ Θεὸς βοηθεῖ μοι, καὶ ὁ Κύριος ἀντιλήπτωρ τῆς ψυχῆς μου.
\vs{7}Ἀποστρέψει τὰ κακὰ τοῖς ἐχθροῖς μου, ἐν τῇ ἀληθείᾳ σου ἐξολόθρευσον αὐτούς.
\vs{8}Ἑκουσίως θύσω σοι, ἐξομολογήσομαι τῷ ὀνόματί σου Κύριε, ὅτι ἀγαθόν.
\vs{9}Ὅτι ἐκ πάσης θλίψεως ἐῤῥύσω με, καὶ ἐν τοῖς ἐχθροῖς μου ἐπεῖδεν ὁ ὀφθαλμός μου.

\begin{psalmheading}{\ch{54}{55} Εἰς τὸ τέλος, ἐν ὕμνοις συνέσεως τῷ Δαυίδ.}
\end{psalmheading}
\vs{2}Ἐνωτίσαι ὁ Θεὸς τὴν προσευχήν μου, καὶ μὴ ὑπερίδῃς τὴν δέησίν μου·
\vs{3}Πρόσχες μοι, καὶ εἰσάκουσόν μου· ἐλυπήθην ἐν τῇ ἀδολεσχίᾳ μου, καὶ ἐταράχθην
\vs{4}ἀπὸ φωνῆς ἐχθροῦ, καὶ ἀπὸ θλίψεως ἁμαρτωλοῦ· ὅτι ἐξέκλιναν ἐπʼ ἐμὲ ἀνομίαν, καὶ ἐν ὀργῇ ἐνεκότουν μοι.

\vs{5}Ἡ καρδία μου ἐταράχθη ἐν ἐμοὶ, καὶ δειλία θανάτου ἐπέπεσεν ἐπʼ ἐμέ.
\vs{6}Φόβος καὶ τρόμος ἦλθεν ἐπʼ ἐμὲ καὶ ἐκάλυψέ με σκότος.
\vs{7}Καὶ εἶπα, τίς δώσει μοι πτέρυγας ὡσεὶ περιστερᾶς; καὶ πετασθήσομαι καὶ καταπαύσω.
\vs{8}Ἰδοὺ ἐμάκρυνα φυγαδεύων, καὶ ηὐλίσθην ἐν τῇ ἐρήμῳ· διάψαλμα.
\vs{9}Προσεδεχόμην τὸν σώζοντά με ἀπὸ ὀλιγοψυχίας καὶ καταιγίδος.

\vs{10}Καταπόντισον Κύριε καὶ καταδίελε τὰς γλώσσας αὐτῶν, ὅτι εἶδον ἀνομίαν καὶ ἀντιλογίαν ἐν τῇ πόλει.
\vs{11}Ἡμέρας καὶ νυκτὸς κυκλώσει αὐτὴν ἐπὶ τὰ τείχη αὐτῆς, ἀνομία καὶ πόνος ἐν μέσῳ αὐτῆς
\vs{12}καὶ ἀδικία, καὶ οὐκ ἐξέλιπεν ἐκ τῶν πλατειῶν αὐτῆς τόκος καὶ δόλος.

\vs{13}Ὅτι εἰ ἐχθρὸς ὠνείδισέ με, ὑπήνεγκα ἂν, καὶ εἰ ὁ μισῶν ἐπʼ ἐμὲ ἐμεγαλοῤῥημόνησεν, ἐκρύβην ἂν ἀπʼ αὐτοῦ.
\vs{14}Σὺ δὲ ἄνθρωπε ἰσόψυχε, ἡγεμών μου καὶ γνωστέ μου,
\vs{15}ὃς ἐπιτοαυτὸ ἐγλύκανας ἐδέσματα, ἐν τῷ οἴκῳ τοῦ Θεοῦ ἐπορεύθημεν ἐν ὁμονοίᾳ.
\vs{16}Ἐλθέτω θάνατος ἐπʼ αὐτοὺς, καὶ καταβήτωσαν εἰς ᾅδου ζῶντες· ὅτι πονηρία ἐν ταῖς παροικίαις αὐτῶν ἐν μέσῳ αὐτῶν.

\vs{17}Ἐγὼ πρὸς τὸν Θεὸν ἐκέκραξα, καὶ ὁ Κύριος εἰσήκουσέ μου.
\vs{18}Ἑσπέρας καὶ πρωῒ καὶ μεσημβρίας διηγήσομαι, καὶ ἀπαγγελῶ, καὶ εἰσακούσεται τῆς φωνῆς μου.
\vs{19}Λυτρώσεται ἐν εἰρήνῃ τὴν ψυχήν μου ἀπὸ τῶν ἐγγιζόντων μοι, ὅτι ἐν πολλοῖς ἦσαν σὺν ἐμοί.
\vs{20}Εἰσακούσεται ὁ Θεὸς καὶ ταπεινώσει αὐτοὺς, ὁ ὑπάρχων πρὸ τῶν αἰώνων· διάψαλμα·

\vs{21}Οὐ γάρ ἐστιν αὐτοῖς ἀντάλλαγμα, καὶ οὐκ ἐφοβήθησαν τὸν Θεόν. Ἐξέτεινε τὴν χεῖρα αὐτοῦ ἐν τῷ ἀποδιδόναι· ἐβεβήλωσαν τὴν διαθήκην αὐτοῦ.
\vs{22}Διεμερίσθησαν ἀπὸ ὀργῆς τοῦ προσώπου αὐτοῦ, καὶ ἤγγισεν ἡ καρδία αὐτοῦ· ἡπαλύνθησαν οἱ λόγοι αὐτοῦ ὑπὲρ ἔλαιον, καὶ αὐτοί εἰσι βολίδες.

\vs{23}Ἐπίῤῥιψον ἐπὶ Κύριον τὴν μέριμνάν σου, καὶ αὐτός σε διαθρέψει, οὐ δώσει εἰς τὸν αἰῶνα σάλον τῷ δικαίῳ.
\vs{24}Σὺ δὲ ὁ Θεὸς κατάξεις αὐτοὺς εἰς φρέαρ διαφθορᾶς· ἄνδρες αἱμάτων καὶ δολιότητος οὐ μὴ ἡμισεύσωσι τὰς ἡμέρας αὐτῶν· ἐγὼ δὲ ἐλπιῶ ἐπὶ σε, Κύριε.

\begin{psalmheading}{\ch{55}{56} Εἰς τὸ τέλος, ὑπὲρ τοῦ λαοῦ τοῦ ἀπὸ τῶν ἁγίων μεμακρυμμένου, τῷ Δαυὶδ εἰς στηλογραφίαν, ὁπότε ἐκράτησαν αὐτὸν οἱ ἀλλόφυλοι ἐν Γέθ.}
\end{psalmheading}
\vs{2}Ἐλεήσον με ὁ Θεὸς, ὅτι κατεπάτησέ με ἄνθρωπος, ὅλην τὴν ἡμέραν πολεμῶν ἔθλιψέ με.
\vs{3}Κατεπάτησάν με οἱ ἐχθροί μου ὅλην τὴν ἡμέραν ἀπὸ ὕψους ἡμέρας, ὅτι πολλοὶ οἱ πολεμοῦντές με.

\vs{4}Φοβηθήσονται, ἐγὼ δὲ ἐλπιῶ ἐπὶ σοί.
\vs{5}Ἐν τῷ Θεῷ ἐπαινέσω τοὺς λόγους μου, ὅλην τὴν ἡμέραν ἐν τῷ Θεῷ ἤλπισα, οὐ φοβηθήσομαι τί ποιήσει μοι σάρξ.

\vs{6}Ὅλην τὴν ἡμέραν τοὺς λόγους μου ἐβδελύσσοντο, κατʼ ἐμοῦ πάντες οἱ διαλογισμοὶ αὐτῶν εἰς κακόν.
\vs{7}Παροικήσουσι καὶ κατακρύψουσιν αὐτοὶ, τὴν πτέρναν μου φυλάξουσι· καθάπερ ὑπέμεινα τῇ ψυχῇ μου.
\vs{8}Ὑπὲρ τοῦ μηθενὸς σώσεις αὐτοὺς ἐν ὀργῇ λαοὺς κατάξεις· ὁ Θεὸς
\vs{9}τὴν ζωήν μου ἐξήγγειλά σοι, ἔθου τὰ δάκρυά μου ἐνώπιόν σου, ὡς καὶ ἐν τῇ ἐπαγγελίᾳ σου.

\vs{10}Ἐπιστρέψουσιν οἱ ἐχθροί μου εἰς τὰ ὀπίσω, ἐν ᾗ ἂν ἡμέρᾳ ἐπικαλέσωμαί σε· ἰδοὺ ἔγνων ὅτι Θεός μου εἶ σύ.
\vs{11}Ἐπὶ τῷ Θεῷ αἰνέσω ῥῆμα, ἐπὶ τῷ Κυρίῳ αἰνέσω λόγον.
\vs{12}Ἐπὶ τῷ Θεῷ ἤλπισα, οὐ φοβηθήσομαι τί ποιήσει μοι ἄνθρωπος.
\vs{13}Ἐν ἐμοὶ ὁ Θεὸς αἱ εὐχαὶ, ἃς ἀποδώσω αἰνέσεώς σου.
\vs{14}Ὅτι ἐῤῥύσω τὴν ψυχήν μου ἐκ θανάτου, καὶ τοὺς πόδας μου ἐξ ὀλισθήματος, τοῦ εὐαρεστῆσαι ἐνώπιον τοῦ Θεοῦ ἐν φωτὶ ζώντων.

\begin{psalmheading}{\ch{56}{57} Εἰς τὸ τέλος, μὴ διαφθείρῃς, τῷ Δαυὶδ εἰς στηλογραφίαν, ἐν τῷ αὐτὸν ἀποδιδράσκειν ἀπὸ προσώπου Σαοὺλ εἰς τὸ σπήλαιον.}
\end{psalmheading}
\vs{2}Ἐλέησον με ὁ Θεὸς, ἐλέησόν με, ὅτι ἐπὶ σοὶ πέποιθεν ἡ ψυχή μου, καὶ ἐν τῇ σκιᾷ τῶν πτερύγων σου ἐλπιῶ, ἕως οὗ παρέλθῃ ἡ ἀνομία.
\vs{3}Κεκράξομαι πρὸς τὸν Θεὸν τὸν ὕψιστον τὸν Θεὸν τὸν εὐεργετήσαντά με· διάψαλμα.
\vs{4}Ἐξαπέστειλεν ἐξ οὐρανοῦ καὶ ἔσωσέ με, ἔδωκεν εἰς ὄνειδος τοὺς καταπατοῦντάς με· ἐξαπέστειλεν ὁ Θεὸς τὸ ἔλεος αὐτοῦ καὶ τὴν ἀλήθειαν αὐτοῦ,
\vs{5}καὶ ἐῤῥύσατο τὴν ψυχήν μου ἐκ μέσου σκύμνων· ἐκοιμήθην τεταραγμένος· υἱοὶ ἀνθρώπων, οἱ ὀδόντες αὐτῶν, ὅπλον καὶ βέλη, καὶ ἡ γλῶσσα αὐτῶν, μάχαιρα ὀξεῖα.

\vs{6}Ὑψώθητι ἐπὶ τοὺς οὐρανοὺς ὁ Θεὸς, καὶ ἐπὶ πᾶσαν τὴν γῆν ἡ δόξα σου.
\vs{7}Παγίδας ἡτοίμασαν τοῖς ποσί μου, καὶ κατέκαμψαν τὴν ψυχήν μου· ὤρυξαν πρὸ προσώπου μου βόθρον, καὶ ἐνέπεσαν εἰς αὐτόν· διάψαλμα.
\vs{8}Ἑτοίμη ἡ καρδία μου ὁ Θεὸς, ἑτοίμη ἡ καρδία μου, ᾄσομαι καὶ ψαλῶ.
\vs{9}Ἐξεγέρθητι ἡ δόξα μου, ἐξεγέρθητι ψαλτήριον καὶ κιθάρα, ἐξεγερθήσομαι ὄρθρου.
\vs{10}Ἐξομολογήσομαί σοι ἐν λαοῖς Κύριε, ψαλῶ σοι ἐν ἔθνεσιν.
\vs{11}Ὅτι ἐμεγαλύνθη ἕως τῶν οὐρανῶν τὸ ἔλεός σου, καὶ ἕως τῶν νεφελῶν ἡ ἀλήθειά σου.
\vs{12}Ὑψώθητι ἐπὶ τοὺς οὐρανοὺς ὁ Θεὸς, καὶ ἐπὶ πᾶσαν τὴν γῆν ἡ δόξα σου.

\begin{psalmheading}{\ch{57}{58} Εἰς τὸ τέλος, μὴ διαφθείρῃς, τῷ Δαυὶδ εἰς στηλογραφίαν.}
\end{psalmheading}
\vs{2}Εἰ ἀληθῶς ἄρα δικαιοσύνην λαλεῖτε, εὐθεῖα κρίνετε οἱ υἱοὶ τῶν ἀνθρώπων.
\vs{3}Καὶ γὰρ ἐν καρδίᾳ ἀνομίας ἐργάζεσθε ἐν τῇ γῇ, ἀδικίαν αἱ χεῖρες ὑμῶν συμπλέκουσιν.
\vs{4}Ἀπηλλοτριώθησαν οἱ ἁμαρτωλοὶ ἀπὸ μήτρας, ἐπλανήθησαν ἀπὸ γαστρὸς, ἐλάλησαν ψευδῆ.
\vs{5}Θυμὸς αὐτοῖς κατὰ τὴν ὁμοίωσιν τοῦ ὄφεως, ὡσεὶ ἀσπίδος κωφῆς, καὶ βυούσης τὰ ὦτα αὐτῆς,
\vs{6}ἥτις οὐκ εἰσακούσεται φωνὴν ἐπᾳδόντων, φαρμάκου τε φαρμακευομένου παρὰ σοφοῦ.

\vs{7}Ὁ Θεὸς συνέτριψε τοὺς ὀδόντας αὐτῶν ἐν τῷ στόματι αὐτῶν, τὰς μύλας τῶν λεόντων συνέθλασεν ὁ Κύριος.
\vs{8}Ἐξουδενωθήσονται ὡς ὕδωρ διαπορευόμενον, ἐντενεῖ τὸ τόξον αὐτοῦ ἕως οὗ ἀσθενήσουσιν.
\vs{9}Ὡσεὶ κηρὸς ὁ τακεῖς ἀνταναιρεθήσονται, ἔπεσε πῦρ, καὶ οὐκ εἶδον τὸν ἥλιον.
\vs{10}Πρὸ τοῦ συνιέναι τὰς ἀκάνθας ὑμῶν τὴν ῥάμνον, ὡσεὶ ζῶντας ὡσεὶ ἐν ὀργῇ καταπίεται ὑμᾶς.

\vs{11}Εὐφρανθήσεται δίκαιος, ὅταν ἴδῃ ἐκδίκησιν ἀσεβῶν, τὰς χεῖρας αὐτοῦ νίψεται ἐν τῷ αἵματι τοῦ ἁμαρτωλοῦ.
\vs{12}Καὶ ἐρεῖ ἄνθρωπος, εἰ ἄρα ἐστὶ καρπὸς τῷ δικαίῳ, ἄρα ἐστὶν ὁ Θεὸς κρίνων αὐτοὺς ἐν τῇ γῇ.

\begin{psalmheading}{\ch{58}{59} Εἰς τὸ τέλος, μὴ διαφθείρῃς, τῷ Δαυὶδ εἰς στηλογραφίαν, ὁπότε ἀπέστειλε Σαοὺλ, καὶ ἐφύλαξε τὸν οἶκον αὐτοῦ τοῦ θανατῶσαι αὐτόν.}
\end{psalmheading}
\vs{2}Ἐξελοῦ με ἐκ τῶν ἐχθρῶν μου ὁ Θεὸς, καὶ ἐκ τῶν ἐπανισταμένων ἐπʼ ἐμὲ λύτρωσαί με.
\vs{3}Ῥῦσαί με ἐκ τῶν ἐργαζομένων τὴν ἀνομίαν, καὶ ἐξ ἀνδρῶν αἱμάτων σῶσόν με.

\vs{4}Ὅτι ἰδοὺ ἐθήρευσαν τὴν ψυχήν μου, ἐπέθεντο ἐπʼ ἐμὲ κραταιοί· οὔτε ἡ ἀνομία μου, οὔτε ἡ ἁμαρτία μου Κύριε·
\vs{5}Ανευ ἀνομίας ἔδραμον καὶ κατεύθυνα· ἐξεγέρθητι εἰς συνάντησίν μου, καὶ ἴδε.
\vs{6}Καὶ σὺ Κύριε ὁ Θεὸς τῶν δυνάμεων ὁ Θεὸς τοῦ Ἰσραὴλ, πρόσχες τοῦ ἐπισκέψασθαι πάντα τὰ ἔθνη, μὴ οἰκτειρήσῃς πάντας τοὺς ἐργαζομένους τὴν ἀνομίαν· διάψαλμα.
\vs{7}Ἐπιστρέψουσιν εἰς ἑσπέραν, καὶ λιμώξουσιν ὡς κύων, καὶ κυκλώσουσιν πόλιν.

\vs{8}Ἰδοὺ ἀποφθέγξονται ἐν τῷ στόματι αὐτῶν, καὶ ῥομφαία ἐν τοῖς χείλεσιν αὐτῶν, ὅτι τίς ἤκουσε;
\vs{9}Καὶ δὺ Κύριε ἐκγελάσῃ αὐτοὺς, ἐξουδενώσεις τάντα τὰ ἔθνη.
\vs{10}Τὸ κράτος μου πρὸς σὲ φυλάξω, ὅτι σὺ ὁ Θεὸς ἀντιλήμπτωρ μου εἶ.
\vs{11}Ὁ Θεός μου, τὸ ἔλεος αὐτοῦ προφθάσει με, ὁ Θεός μου δείξει μοι ἐν τοῖς ἐχθροῖς μου.

\vs{12}Μὴ ἀποκτείνῃς αὐτοὺς, μήποτε ἐπιλάθωνται τοῦ νόμου σου· διασκόρπισον αὐτοὺς ἐν τῇ δυνάμει σου, καὶ κατάγαγε αὐτοὺς ὁ ὑπερασπιστής μου Κύριε.
\vs{13}Ἁμαρτίαν στόματος αὐτῶν, λόγον χειλέων αὐτῶν, καὶ συλληφθήτωσαν ἐν τῇ ὑπερηφανίᾳ αὐτῶν· καὶ ἐξ ἀρᾶς καὶ ψεύδους διαγγελήσονται
\vs{14}συντέλειαι, ἐν ὀργῇ συντελείας, καὶ οὐ μὴ ὑπάρξουσι· καὶ γνώσονται ὅτι ὁ Θεὸς τοῦ Ἰακὼβ δεσπόζει τῶν περάτων τῆς γῆς· διάψαλμα.
\vs{15}Ἐπιστρέψουσιν εἰς ἑσπέραν, καὶ λιμώξουσιν ὡς κύων, καὶ κυκλώσουσι πόλιν·
\vs{16}Αὐτοὶ διασκορπισθήσονται τοῦ φαγεῖν, ἐὰν δὲ μὴ χορτασθῶσι, καὶ γογγύσουσιν.

\vs{17}Ἐγὼ δὲ ᾄσομαι τῇ δυνάμει σου, καὶ ἀγαλλιάσομαι τοπρωῒ τὸ ἔλεός σου, ὅτι ἐγενήθης ἀντιλήπτωρ μου καὶ καταφυγή μου ἐν ἡμέρᾳ θλίψεώς μου.
\vs{18}Βοηθός μου, σοὶ ψαλῶ ὁ Θεός μου, ἀντιλήμπτωρ μου εἶ ὁ Θεός μου, τὸ ἔλεός μου.

\begin{psalmheading}{\ch{59}{60} Εἰς τὸ τέλος, τοῖς ἀλλοιωθησομένοις ἔτι, εἰς στηλογραφίαν τῷ Δαυὶδ εἰς διδαχὴν, ὁπότε ἐνεπύρισε τὴν Μεσοποταμίαν Συρίας, καὶ τὴν Συρίαν Σοβὰλ, καὶ ἐπέστρεψεν Ἰωὰβ, καὶ ἐπάταξε τὴν φάραγγα τῶν ἁλῶν, δώδεκα χιλιάδας.}
\end{psalmheading}
\vs{3}Ὁ Θεὸς ἀπώσω ἡμᾶς καὶ καθεῖλες ἡμᾶς, ὠργίσθης καὶ ᾠκτείρησας ἡμᾶς.
\vs{4}Συνέσεισας τὴν γῆν καὶ συνετάραξας αὐτὴν, ἴασαι τὰ συντρίμματα αὐτῆς, ὅτι ἐσαλεύθη.
\vs{5}Ἔδειξας τῷ λαῷ σου σκληρὰ, ἐπότισας ἡμᾶς οἶνον κατανύξεως.
\vs{6}Ἔδωκας τοῖς φοβουμένοις σε σημείωσιν, τοῦ φυγεῖν ἀπὸ προσώπου τόξον· διάψαλμα.
\vs{7}Ὅπως ἂν ῥυσθῶσιν οἱ ἀγαπητοί σου, σῶσον τῇ δεξιᾷ σου καὶ ἐπάκουσόν μου.

\vs{8}Ὁ Θεὸς ἐλάλησεν ἐν τῷ ἁγίῳ αὐτοῦ, ἀγαλλιάσομαι καὶ διαμεριῶ Σίκιμα, καὶ τὴν κοιλάδα τῶν σκηνῶν διαμετρήσω.
\vs{9}Ἐμός ἐστι Γαλαὰδ, καὶ ἐμός ἐστι Μανασσῆ, καὶ Ἐφραὶμ κραταίωσις τῆς κεφαλῆς μου· Ἰούδας βασιλεύς μου,
\vs{10}Μωὰβ λέβης τῆς ἐλπίδος μου, ἐπὶ τὴν Ἰδουμαίαν ἐκτενῶ τὸ ὑπόδημά μου, ἐμοὶ ἀλλόφυλοι ὑπετάγησαν.

\vs{11}Τίς ἀπάξει με εἰς πόλιν περιοχῆς; τίς ὁδηγήσει με ἕως τῆς Ἰδουμαίας;
\vs{12}Οὐχὶ σὺ ὁ Θεὸς, ὁ ἀπωσάμενος ἡμᾶς; καὶ οὐκ ἐξελεύσῃ, ὁ Θεὸς ἐν ταῖς δυνάμεσιν ἡμῶν;
\vs{13}Δὸς ἡμῖν βοήθειαν ἐκ θλίψεως, καὶ ματαία σωτηρία ἀνθρώπου.

\vs{14}Ἐν τῷ Θεῷ ποιήσομεν δύναμιν, καὶ αὐτὸς ἐξουδενώσει τοὺς θλίβοντας ἡμᾶς.

\begin{psalmheading}{\ch{60}{61} Εἰς τὸ τέλος, ἐν ὕμνοις τῷ Δαυίδ.}
\end{psalmheading}
\vs{2}Εἰσάκουσον ὁ Θεὸς τῆς δεήσεώς μου, πρόσχες τῇ προσευχῇ μου.
\vs{3}Ἀπὸ τῶν περάτων τῆς γῆς πρὸς σὲ ἐκέκραξα, ἐν τῷ ἀκηδιάσαι τὴν καρδίαν μου, ἐν πέτρᾳ ὕψωσάς με, ὁδήγησάς με,
\vs{4}ὅτι ἐγενήθης ἐλπίς μου, πύργος ἰσχύος ἀπὸ προσώπου ἐχθροῦ.
\vs{5}Παροικήσω ἐν τῷ σκηνώματί σου εἰς τοὺς αἰῶνας, σκεπασθήσομαι ἐν σκέπῃ τῶν πτερύγων σου· διάψαλμα.

\vs{6}Ὅτι σὺ ὁ Θεὸς εἰσήκουσας τῶν προσευχῶν μου, ἔδωκας κληρονομίαν τοῖς φοβουμένοις τὸ ὄνομά σου.
\vs{7}Ἡμέρας ἐφʼ ἡμέρας βασιλέως προσθήσεις, τὰ ἔτη αὐτοῦ ἕως ἡμέρας γενεᾶς καὶ γενεᾶς.
\vs{8}Διαμενεῖ εἰς τὸν αἰῶνα ἐνώπιον τοῦ Θεοῦ, ἔλεος καὶ ἀλήθειαν αὐτοῦ τίς ἐκζητήσει αὐτῶν;
\vs{9}Οὕτως ψαγῶ τῷ ὀνόματί σου εἰς τὸν αἰῶνα τοῦ αἰῶνος, τοῦ ἀποδοῦναί με τὰς εὐχάς μου ἡμέραν ἐξ ἡμέρας.

\begin{psalmheading}{\ch{61}{62} Εἰς τὸ τέλος, ὑπὲρ Ἰδιθοὺν ψαλμὸς τῷ Δαυίδ.}
\end{psalmheading}
\vs{2}Οὐχὶ τῷ Θεῷ ὑποταγήσεται ἡ ψυχή μου; παρʼ αὐτοῦ γὰρ τὸ σωτήριόν μου.
\vs{3}Καὶ γὰρ αὐτὸς Θεός μου καὶ σωτήρ μου, ἀντιλήπτωρ μου, οὐ μὴ σαλευθῶ ἐπὶ πλεῖον.
\vs{4}Ἕως πότε ἐπιτίθεσθε ἐπʼ ἄνθρωπον; φονεύετε πάντες ὡς τοίχῳ κεκλιμένῳ καὶ φραγμῷ ὠσμένῳ.
\vs{5}Πλὴν τὴν τιμήν μου ἐβουλεύσαντο ἀπώσασθαι· ἔδραμον ἐν δίψει· τῷ στόματι αὐτῶν εὐλόγουν, καὶ τῇ καρδίᾳ αὐτῶν κατηρῶντο. διάψαλμα.

\vs{6}Πλὴν τῷ Θεῷ ὑποτάγηθι, ἡ ψυχή μου, ὅτι παρʼ αὐτοῦ ἡ ὑπομονή μου.
\vs{7}Ὅτι αὐτὸς Θεός μου καὶ σωτήρ μου, ἀντιλήπτωρ μου, οὐ μὴ μεταναστεύσω.
\vs{8}Ἐπὶ τῷ Θεῷ τὸ σωτήριόν μου, καὶ ἡ δόξα μου· ὁ Θεὸς τῆς βοηθείας μου, καὶ ἡ ἐλπίς μου ἐπὶ τῷ Θεῷ.
\vs{9}Ἐλπίσατε ἐπʼ αὐτὸν πᾶσα συναγωγὴ λαοῦ· ἐκχέετε ἐνώπιον αὐτοῦ τὰς καρδίας ὑμῶν, ὅτι ὁ Θεὸς βοηθὸς ἡμῶν· διάψαλμα.

\vs{10}Πλὴν μάταιοι οἱ νἱοὶ τῶν ἀνθρώπων, ψευδεῖς οἱ υἱοὶ τῶν ἀνθρώπων ἐν ζυγοῖς τοῦ ἀδικῆσαι, αὐτοὶ ἐκ ματαιότητος ἐπιτοαυτό.
\vs{11}Μὴ ἐλπίζετε ἐπʼ ἀδικίαν, καὶ ἐπὶ ἁρπάγματα μὴ ἐπιποθεῖτε· πλοῦτος ἐὰν ῥέῃ, μὴ προστίθεσθε καρδίαν.
\vs{12}Ἅπαξ ἐλάλησεν ὁ Θεὸς, δύο ταῦτα ἤκουσα, ὅτι τὸ κράτος τοῦ Θεοῦ·
\vs{13}καὶ σοῦ, Κύριε τὸ ἔλεος, ὅτι σὺ ἀποδώσεις ἑκάστῳ κατὰ τὰ ἔργα αὐτοῦ.

\begin{psalmheading}{\ch{62}{63} Ψαλμὸς τῷ Δαυὶδ, ἐν τῷ εἶναι αὐτὸν ἐν τῇ ἐρήμῳ τῆς Ἰδουμαίας.}
\end{psalmheading}
\vs{2}Ὁ Θεὸς ὁ Θεός μου πρὸς σὲ ὀρθρίζω, ἐδίψησέ σοι ἡ ψυχή μου, ποσαπλῶς σοι ἡ σάρξ μου, ἐν γῇ ἐρήμῳ καὶ ἀβάτῳ καὶ ἀνύδρῳ,
\vs{3}οὕτως ἐν τῷ ἁγίῳ ὤφθην σοι, τοῦ ἰδεῖν τὴν δύναμίν σου καὶ τὴν δόξαν σου.
\vs{4}Ὅτι κρεῖσσον τὸ ἔλεός σου ὑπὲρ ζωὰς, τὰ χείλη μου ἐπαινέσουσί σε.
\vs{5}Οὕτως εὐλογήσω σε ἐν τῇ ζωῇ μου, ἐν τῷ ὀνόματί σου ἀρῶ τὰς χεῖράς μου.
\vs{6}Ὡσεὶ στέατος καὶ πιότητος ἐμπλησθείη ἡ ψυχή μου, καὶ χείλη ἀγαλλιάσεως αἰνέσει τὸ ὄνομά σου.

\vs{7}Εἰ ἐμνημόνευόν σου ἐπὶ τῆς στρωμνῆς μου, ἐν τοῖς ὄρθροις ἐμελέτων εἰς σέ.
\vs{8}Ὅτι ἐγενήθης βοηθός μου, καὶ ἐν τῇ σκέπῃ τῶν πτερύγων σου ἀγαλλιάσομαι.
\vs{9}Ἐκολλήθη ἡ ψυχή μου ὀπίσω σου, ἐμοῦ ἀντελάβετο ἡ δεξιά σου.
\vs{10}Αὐτοὶ δὲ εἰς μάτην ἐζήτησαν τὴν ψυχήν μου, εἰσελεύσονται εἰς τὰ κατώτατα τῆς γῆς,
\vs{11}παραδοθήσονται εἰς χεῖρας ῥομφαίας, μερίδες ἀλωπέκων ἔσονται.
\vs{12}Ὁ δὲ βασιλεὺς εὐφρανθήσεται ἐπὶ τῷ Θεῷ, ἐπαινεθήσεται πᾶς ὁ ὀμνύων ἐν αὐτῷ, ὅτι ἐνεφράγη στόμα λαλούντων ἄδικα.

\begin{psalmheading}{\ch{63}{64} Εἰς τὸ τέλος, ψαλμὸς τῷ Δαυίδ.}
\end{psalmheading}
\vs{2}Εἰσάκουσον ὁ Θεὸς τῆς προσευχῆς μου ἐν τῷ δέεσθαί με πρὸς σὲ, ἀπὸ φόβου ἐχθροῦ ἐξελοῦ τὴν ψυχήν μου.
\vs{3}Ἐσκέπασάς με ἀπὸ συστροφῆς πονηρευομένων, ἀπὸ πλήθους ἐργαζομένων ἀδικίαν·
\vs{4}Οἵτινες ἠκόνησαν ὡς ῥομφαίαν τὰς γλώσσας αὐτῶν, ἐνέτειναν τόξον πρᾶγμα πικρὸν,
\vs{5}τοῦ κατατοξεῦσαι ἐν ἀποκρύφοις ἄμωμον, ἐξάπινα κατατοξεύσουσιν αὐτὸν, καὶ οὐ φοβηθήσονται.
\vs{6}Ἐκραταίωσαν ἑαυτοῖς λόγον πονηρὸν, διηγήσαντο τοῦ κρύψαι παγίδας· εἶπαν, τίς ὄψεται αὐτούς;
\vs{7}Ἐξηρεύνησαν ἀνομίαν, ἐξέλιπον ἐξερευνῶντες ἐξερευνήσει· προσελεύσεται ἄνθρωπος, καὶ καρδία βαθεῖα,
\vs{8}καὶ ὑψωθήσεται ὁ Θεός· βέλος νηπίων ἐγενήθησαν αἱ πληγαὶ αὐτῶν,
\vs{9}καὶ ἐξουθένησαν αὐτὸν αἱ γλῶσσαι αὐτῶν· ἐταράχθησαν πάντες οἱ θεωροῦντες αὐτοὺς,
\vs{10}καὶ ἐφοβήθη πᾶς ἄνθρωπος· καὶ ἀνήγγειλαν τὰ ἔργα τοῦ Θεοῦ, καὶ τὰ ποιήματα αὐτοῦ συνῆκαν.
\vs{11}Εὐφρανθήσεται δίκαιος ἐν τῷ Κυρίῳ, καὶ ἐλπιεῖ ἐπʼ αὐτόν· καὶ ἐπαινεθήσονται πάντες οἱ εὐθεῖς τῇ καρδίᾳ.

\begin{psalmheading}{\ch{64}{65} Εἰς τὸ τέλος, ψαλμὸς τῷ Δαυίδ, ᾠδή.}
\end{psalmheading}
\vs{2}Σοὶ πρέπει ὕμνος, ὁ Θεὸς ἐν Σιὼν, καὶ σοὶ ἀποδοθήσεται εὐχή.
\vs{3}Εἰσάκουσον προσευχῆς μου, πρὸς σὲ πᾶσα σὰρξ ἥξει.
\vs{4}Λόγοι ἀνόμων ὑπερδυνάμωσαν ἡμᾶς, καὶ τὰς ἀσεβείας ἡμῶν σὺ ἱλάσῃ.
\vs{5}Μακάριος, ὃν ἐξελέξω καὶ προσελάβου, κατασκηνώσει ἐν ταῖς αὐλαῖς σου· πλησθησόμεθα ἐν τοῖς ἀγαθοῖς τοῦ οἴκου σου, ἅγιος ὁ ναός σου,
\vs{6}θαυμαστὸς ἐν δικαιοσύνῃ. ἐπάκουσον ἡμῶν ὁ Θεὸς ὁ σωτὴρ ἡμῶν, ἡ ἐλπὶς πάντων τῶν περάτων τῆς γῆς, καὶ τῶν ἐν θαλάσσῃ μακράν·
\vs{7}ἑτοιμάζων ὄρη ἐν τῇ ἰσχύϊ σου, περιεζωσμένος ἐν δυναστείᾳ·
\vs{8}Ὁ συνταράσσων τὸ κῦτος τῆς θαλάσσης, ἤχους κυμάτων αὐτῆς. Ταραχθήσονται τὰ ἔθνη,
\vs{9}καὶ φοβηθήσονται οἱ κατοικοῦντες τὰ πέρατα ἀπὸ τῶν σημείων σου· ἐξόδους πρωΐας καὶ ἑσπέρας τέρψεις.

\vs{10}Ἐπεσκέψω τὴν γῆν καὶ ἐμέθυσας αὐτὴν, ἐπλήθυνας τοῦ πλουτίσαι αὐτήν· ὁ ποταμὸς τοῦ Θεοῦ ἐπληρώθη ὑδάτων· ἡτοίμασας τὴν τροφὴν αὐτῶν, ὅτι οὕτως ἡ ἑτοιμασία.
\vs{11}Τοὺς αὔλακας αὐτῆς μέθυσον, πλήθυνον τὰ γεννήματα αὐτῆς, ἐν ταῖς σταγόσιν αὐτῆς εὐφρανθήσεται ἀνατέλλουσα.
\vs{12}Εὐλογήσεις τὸν στέφανον τοῦ ἐνιαυτοῦ τῆς χρηστότητός σου, καὶ τὰ πεδία σου πλησθήσονται πιότητος.
\vs{13}Πιανθήσεται τὰ ὄρη τῆς ἐρήμου, καὶ ἀγαλλίασιν οἱ βουνοὶ περιζώσονται.
\vs{14}Ἐνεδύσαντο οἱ κριοὶ τῶν προβάτων, καὶ αἱ κοιλάδες πληθυνοῦσι σῖτον, κεκράξονται, καὶ γὰρ ὑμνήσουσιν.

\begin{psalmheading}{\ch{65}{66} Εἰς τὸ τέλος, ᾠδὴ ψαλμοῦ ἀναστάσεως.}
\end{psalmheading}
Ἀλαλάξατε τῷ Θεῷ, πᾶσα ἡ γῆ,
\vs{2}ψάλατε δὴ τῷ ὀνόματι αὐτοῦ, δότε δόξαν αἰνέσει αὐτοῦ.
\vs{3}Εἴπατε τῷ Θεῷ ὡς φοβερὰ τὰ ἔργα σου; ἐν τῷ πλήθει τῆς δυνάμεώς σου ψεύσονταί σε οἱ ἐχθροί σου.
\vs{4}Πᾶσα ἡ γῆ προσκυνησάτωσάν σοι, καὶ ψαλάτωσάν σοι, ψαλάτωσαν τῷ ὀνόματί σου· διάψαλμα.

\vs{5}Δεῦτε καὶ ἴδετε τὰ ἔργα τοῦ Θεοῦ, φοβερὸς ἐν βουλαῖς ὑπὲρ τοὺς υἱοὺς τῶν ἀνθρώπων.
\vs{6}Ὁ μεταστρέφων τὴν θάλασσαν εἰς ξηρὰν, ἐν ποταμῷ διελεύσονται ποδί· ἐκεῖ εὐφρανθησόμεθα ἐπʼ αὐτῷ,
\vs{7}τῷ δεσπόζοντι ἐν τῇ δυναστείᾳ αὐτοῦ τοῦ αἰῶνος· οἱ ὀφθαλμοὶ αὐτοῦ ἐπὶ τὰ ἔθνη ἐπιβλέπουσιν, οἱ παραπικραίνοντες μὴ ὑψούσθωσαν ἐν ἑαυτοῖς· διάψαλμα.

\vs{8}Εὐλογεῖτε ἔθνη τὸν Θεὸν ἡμῶν, καὶ ἀκουτίσατε τὴν φωνὴν τῆς αἰνέσεως αὐτοῦ,
\vs{9}τοῦ θεμένου τὴν ψυχήν μου εἰς ζωὴν, καὶ μὴ δόντος εἰς σάλον τοὺς πόδας μου.
\vs{10}Ὅτι ἐδοκίμασας ἡμᾶς ὁ Θεὸς, ἐπύρωσας ἡμᾶς ὡς πυροῦται τὸ ἀργύριον.
\vs{11}Εἰσήγαγες ἡμᾶς εἰς τὴν παγίδα, ἔθου θλίψεις ἐπὶ τὸν νῶτον ἡμῶν,
\vs{12}ἐπεβίβασας ἀνθρώπους ἐπὶ τὰς κεφαλὰς ἡμῶν· διήλθομεν διὰ πυρὸς καὶ ὕδατος, καὶ ἐξήγαγες ἡμᾶς εἰς ἀναψυχήν.

\vs{13}Εἰσελεύσομαι εἰς τὸν οἶκόν σου ἐν ὁλοκαυτώμασιν, ἀποδώσω σοι τὰς εὐχάς μου,
\vs{14}ἃς διέστειλε τὰ χείλη μου, καὶ ἐλάλησε τὸ στόμα μου ἐν τῇ θλίψει μου.
\vs{15}Ὁλοκαυτώματα μεμυελωμένα ἀνοίσω σοι μετὰ θυμιάματος καὶ κριῶν, ποιήσω σοι βόας μετὰ χιμάρων· διάψαλμα.

\vs{16}Δεῦτε ἀκούσατε, καὶ διηγήσομαι, πάντες οἱ φοβούμενοι τὸν Θεὸν, ὅσα ἐποίησε τῇ ψυχῇ μου.
\vs{17}Πρὸς αὐτὸν τῷ στόματί μου ἐκέκραξα, καὶ ὕψωσα ὑπὸ τὴν γλῶσσάν μου.
\vs{18}Ἀδικίαν εἰ ἐθεώρουν ἐν καρδίᾳ μου, μὴ εἰσακουσάτω Κύριος.
\vs{19}Διὰ τοῦτο εἰσήκουσέ μου ὁ Θεὸς, προσέσχε τῇ φωνῇ τῆς προσευχῆς μου.
\vs{20}Εὐλογητὸς ὁ Θεὸς, ὃς οὐκ ἀπέστησε τὴν προσευχήν μου, καὶ τὸ ἔλεος αὐτοῦ ἀπʼ ἐμοῦ.

\begin{psalmheading}{\ch{66}{67} Εἰς τὸ τέλος, ἐν ὕμνοις ψαλμὸς τῷ Δαυίδ.}
\end{psalmheading}
\vs{2}Ὁ Θεὸς οἰκτειρήσαι ἡμᾶς, καὶ εὐλογήσαι ἡμᾶς, ἐπιφάναι τὸ πρόσωπον αὐτοῦ ἐφʼ ἡμᾶς· διάψαλμα.
\vs{3}Τοῦ γνῶναι ἐν τῇ γῇ τὴν ὁδόν σου, ἐν πᾶσιν ἔθνεσι τὸ σωτήριόν σου.
\vs{4}Ἐξομολογησάσθωσάν σοι λαοὶ ὁ Θεὸς, ἐξομολογησάσθωσάν σοι λαοὶ πάντες.
\vs{5}Εὐφρανθήτωσαν καὶ ἀγαλλιάσθωσαν ἔθνη, ὅτι κρινεῖς λαοὺς ἐν εὐθύτητι, καὶ ἔθνη ἐν τῇ γῇ ὁδηγήσεις· διάψαλμα.
\vs{6}Ἐξομολογησάσθωσάν σοι λαοὶ, ὁ Θεὸς, ἐξομολογησάσθωσάν σοι λαοὶ πάντες.
\vs{7}Γῆ ἔδωκε τὸν καρπὸν αὐτῆς· εὐλογήσαι ἡμᾶς ὁ Θεὸς, ὁ Θεὸς ἡμῶν,
\vs{8}εὐλογήσαι ἡμᾶς ὁ Θεὸς, καὶ φοβηθήτωσαν αὐτὸν πάντα τὰ πέρατα τῆς γῆς.

\begin{psalmheading}{\ch{67}{68} Εἰς τὸ τέλος, τῷ Δαυὶδ ψαλμὸς ᾠδῆς.}
\end{psalmheading}
\vs{2}Αναστήτω ὁ Θεὸς, καὶ διασκορπισθήτωσαν οἱ ἐχθροὶ αὐτοῦ, καὶ φυγέτωσαν οἱ μισοῦντες αὐτὸν ἀπὸ προσώπου αὐτοῦ.
\vs{3}Ὡς ἐκλείπει καπνὸς, ἐκλιπέτωσαν· ὡς τήκεται κηρὸς ἀπὸ προσώπου πυρὸς, οὕτως ἀπόλοιντο οἱ ἁμαρτωλοὶ ἀπὸ προσώπου τοῦ Θεοῦ.
\vs{4}Καὶ οἱ δίκαιοι εὐφρανθήτωσαν· ἀγαλλιάσθωσαν ἐνώπιον τοῦ Θεοῦ, τερφθήτωσαν ἐν εὐφροσύνῃ.

\vs{5}Ἄσατε τῷ Θεῷ, ψάλατε τῷ ὀνόματι αὐτοῦ, ὁδοποιήσατε τῷ ἐπιβεβηκότι ἐπὶ δυσμῶν, Κύριος ὄνομα αὐτῷ, καὶ ἀγαλλιᾶσθε ἐνώπιον αὐτοῦ.
\vs{6}ταραχθήσονται ἀπὸ προσώπου αὐτοῦ, τοῦ πατρὸς τῶν ὀρφανῶν, καὶ κριτοῦ τῶν χηρῶν, ὁ Θεὸς ἐν τόπῳ ἁγίῳ αὐτοῦ.
\vs{7}Ὁ Θεὸς κατοικίζει μονοτρόπους ἐν οἴκῳ, ἐξάγων πεπεδημένους ἐν ἀνδρείᾳ· ὁμοίως τοὺς παραπικραίνοντας, τοὺς κατοικοῦντας ἐν τάφοις.

\vs{8}Ὁ Θεὸς, ἐν τῷ ἐκπορεύεσθαί σε ἐνώπιον τοῦ λαοῦ σου, ἐν τῷ διαβαίνειν σε τὴν ἔρημον· διάψαλμα·
\vs{9}Γῆ ἐσείσθη, καὶ γὰρ οἱ οὐρανοὶ ἔσταξαν ἀπὸ προσώπου τοῦ Θεοῦ τοῦ Σινὰ, ἀπὸ προσώπου τοῦ Θεοῦ Ἰσραήλ.
\vs{10}Βροχὴν ἑκούσιον ἀφοριεῖς ὁ Θεὸς τῇ κληρονομίᾳ σου· καὶ ἠσθένησε, σὺ δὲ κατηρτίσω αὐτήν.

\vs{11}Τὰ ζῶά σου κατοικοῦσιν ἐν αὐτῇ, ἡτοίμασας ἐν τῇ χρηστότητί σου τῷ πτωχῷ.
\vs{12}Ὁ Θεὸς Κύριος δώσει ῥῆμα τοῖς εὐαγγελιζομένοις δυνάμει πολλῇ,
\vs{13}ὁ βασιλεὺς τῶν δυνάμεων τοῦ ἀγαπητοῦ, τοῦ ἀγαπητοῦ, καὶ ὡραιότητι τοῦ οἴκου διελέσθαι σκῦλα.
\vs{14}Ἐὰν κοιμηθῆτε ἀναμέσον τῶν κλήρων, πτέρυγες περιστερᾶς περιηργυρωμέναι, καὶ τὰ μετάφρενα αὐτῆς ἐν χλωρότητι χρυσίου.
\vs{15}Ἐν τῷ διαστέλλειν τὴν ἐπουράνιον βασιλεῖς ἐπʼ αὐτῆς, χιονωθήσονται ἐν Σελμών.
\vs{16}Ὄρος τοῦ Θεοῦ ὄρος πῖον, ὄρος τετυρωμένον, ὄρος πῖον.
\vs{17}Ἱνατί ὑπολαμβάνετε ὄρη τετυρωμένα; τὸ ὄρος ὃ εὐδόκησεν ὁ Θεὸς κατοικεῖν ἐν αὐτῷ· καὶ γὰρ ὁ Κύριος κατασκηνώσει εἰς τέλος.

\vs{18}Τὸ ἅρμα τοῦ Θεοῦ μυριοπλάσιον, χιλιάδες εὐθηνούντων· Κύριος ἐν αὐτοῖς ἐν Σινὰ ἐν τῷ ἁγίῳ.
\vs{19}Ἀναβὰς εἰς ὕψος, ᾐχμαλώτευσας αἰχμαλωσίαν· ἔλαβες δόματα ἐν ἀνθρώπῳ, καὶ γὰρ ἀπειθοῦντες τοῦ κατασκηνῶσαι.

\vs{20}Κύριος ὁ Θεὸς εὐλογητὸς, εὐλογητὸς Κύριος ἡμέραν καθʼ ἡμέραν, καὶ κατευοδώσει ἡμῖν ὁ Θεὸς τῶν σωτηρίων ἡμῶν· διάψαλμα.
\vs{21}Ὁ Θεὸς ἡμῶν, ὁ Θεὸς τοῦ σώζειν, καὶ τοῦ Κυρίου αἱ διέξοδοι τοῦ θανάτου.
\vs{22}Πλὴν ὁ Θεὸς συνθλάσει κεφαλὰς ἐχθρῶν αὐτοῦ, κορυφὴν τριχὸς διαπορευομένων ἐν πλημμελείαις αὐτῶν.
\vs{23}Εἶπε Κύριος, ἐκ Βασὰν ἐπιστρέψω, ἐπιστρέψω ἐν βυθοῖς θαλάσσης.
\vs{24}Ὅπως ἂν βαφῇ ὁ ποῦς σου ἐν αἵματι, ἡ γλῶσσα τῶν κυνῶν σου ἐξ ἐχθρῶν παρʼ αὐτοῦ.

\vs{25}Ἐθεωρήθησαν αἱ πορεῖαί σου ὁ Θεὸς, αἱ πορεῖαι τοῦ Θεοῦ μου τοῦ βασιλέως τοῦ ἐν τῷ ἁγίῳ.
\vs{26}Προέφθασαν ἄρχοντες ἐχόμενοι ψαλλόντων, ἐν μέσῳ νεανίδων τυμπανιστριῶν.
\vs{27}Ἐν ἐκκλησίαις εὐλογεῖτε τὸν Θεὸν, τὸν Κύριον ἐκ πηγῶν Ἰσραήλ.
\vs{28}Ἐκεῖ Βενιαμὶν νεώτερος ἐν ἐκστάσει, ἄρχοντες Ἰούδα ἡγεμόνες αὐτῶν, ἄρχοντες Ζαβουλὼν, ἄρχοντες Νεφθαλί.

\vs{29}Ἔντειλαι ὁ Θεὸς τῇ δυνάμει σου, δυνάμωσον ὁ Θεὸς τοῦτο, ὃ κατηρτίσω ἐν ἡμῖν.
\vs{30}Ἀπὸ τοῦ ναοῦ σου ἐπὶ Ἱερουσαλὴμ, σοὶ οἴσουσι βασιλεῖς δῶρα.
\vs{31}Ἐπιτίμησον τοῖς θηρίοις τοῦ καλάμου· ἡ συναγωγὴ τῶν ταύρων ἐν ταῖς δαμάλεσι τῶν λαῶν, τοῦ μὴ ἀποκλεισθῆναι τοὺς δεδοκιμασμένους τῷ ἀργυρίῳ· διασκόρπισον ἔθνη τὰ τοὺς πολέμους θέλοντα.
\vs{32}Ἥξουσι πρέσβεις ἐξ Αἰγύπτου, Αἰθιοπία προφθάσει χεῖρα αὐτῆς τῷ Θεῷ.

\vs{33}Αἱ βασιλεῖαι τῆς γῆς ᾄσατε τῷ Θεῷ, ψάλατε τῷ Κυρίῳ. διάψαλμα.
\vs{34}Ψάλατε τῷ Θεῷ τῷ ἐπιβεβηκότι ἐπὶ τὸν οὐρανὸν τοῦ οὐρανοῦ κατὰ ἀνατολὰς, ἰδοὺ δώσει ἐν τῇ φωνῇ αὐτοῦ φωνὴν δυνάμεως.
\vs{35}Δότε δόξαν τῷ Θεῷ, ἐπὶ τὸν Ἰσραὴλ ἡ μεγαλοπρέπεια αὐτοῦ, καὶ ἡ δύναμις αὐτοῦ ἐν ταῖς νεφέλαις.
\vs{36}Θαυμαστὸς ὁ Θεὸς ἐν τοῖς ὁσίοις αὐτοῦ, ὁ Θεὸς Ἰσραὴλ, αὐτὸς δώσει δύναμιν καὶ κραταίωσιν τῷ λαῷ αὐτοῦ· εὐλογητὸς ὁ Θεός.

\begin{psalmheading}{\ch{68}{69} Εἰς τὸ τέλος, ὑπὲρ τῶν ἀλλοιωθησομένων, τῶ Δαυίδ.}
\end{psalmheading}
\vs{2}Σώσον με ὁ Θεὸς, ὅτι εἰσήλθοσαν ὕδατα ἕως ψυχῆς μου.
\vs{3}Ἐνεπάγην εἰς ἰλὺν βυθοῦ, καὶ οὐκ ἔστιν ὑπόστασις· ἦλθον εἰς τὰ βάθη τῆς θαλάσσης, καὶ καταιγὶς κατεπόντισέ με.
\vs{4}Ἐκοπίασα κράζων, ἐβραγχίασεν ὁ λάρυγξ μου, ἐξέλιπον οἱ ὀφθαλμοί μου ἀπὸ τοῦ ἐλπίζειν με ἐπὶ τὸν Θεόν μου.
\vs{5}Ἐπληθύνθησαν ὑπὲρ τὰς τρίχας τῆς κεφαλῆς μου οἱ μισοῦντές με δωρεάν· ἐκραταιώθησαν οἱ ἐχθροί μου οἱ ἐκδιώκοντές με ἀδίκως· ἃ οὐχ ἥρπασα, τότε ἀπετίννυον.

\vs{6}Ὁ Θεὸς σὺ ἔγνως τὴν ἀφροσύνην μου, καὶ αἱ πλημμελειαί μου ἀπὸ σοῦ οὐκ ἐκρύβησαν.
\vs{7}Μὴ αἰσχυνθείησαν ἐπʼ ἐμὲ οἱ ὑπομένοντές σε Κύριε τῶν δυνάμεων, μὴ ἐντραπείησαν ἐπʼ ἐμὲ οἱ ζητοῦντές σε ὁ Θεὸς τοῦ Ἰσραήλ.
\vs{8}Ὅτι ἕνεκά σου ὑπήνεγκα ὀνειδισμὸν, ἐκάλυψεν ἐντροπὴ τὸ πρόσωπόν μου.
\vs{9}Ἀπηλλοτριωμένος ἐγενήθην τοῖς ἀδελφοῖς μου, καὶ ξένος τοῖς υἱοῖς τῆς μητρός μου·
\vs{10}ὅτι ὁ ζῆλος τοῦ οἴκου σου κατάφαγέ με, καὶ οἱ ὀνειδισμοὶ τῶν ὀνειδιζόντων σε ἐπέπεσον ἐπʼ ἐμέ.
\vs{11}Καὶ συνέκαμψα ἐν νηστείᾳ τὴν ψυχήν μου, καὶ ἐγενήθη εἰς ὀνειδισμοὺς ἐμοί.
\vs{12}Καὶ ἐθέμην τὸ ἔνδυμά μου σάκκον, καὶ ἐγενόμην αὐτοῖς εἰς παραβολήν.
\vs{13}Κατʼ ἐμοῦ ἠδολέσχουν οἱ καθήμενοι ἐν πύλῃ, καὶ εἰς ἐμὲ ἔψαλλον οἱ πίνοντες τὸν οἶνον.

\vs{14}Ἐγὼ δὲ τῇ προσευχῇ μου πρὸς σὲ Κύριε, καιρὸς εὐδοκίας ὁ Θεός· ἐν τῷ πλήθει τοῦ ἐλέους σου ἐπάκουσόν μου, ἐν ἀληθείᾳ τῆς σωτηρίας σου.
\vs{15}Σῶσόν με ἀπὸ πηλοῦ, ἵνα μὴ ἐνπαγῶ· ῥυσθείην ἐκ τῶν μισούντων με, καὶ ἐκ τοῦ βάθους τῶν ὑδάτων.
\vs{16}Μή με καταποντισάτω καταιγὶς ὕδατος, μηδὲ καταπιέτω με βυθὸς, μηδὲ συνσχέτω ἐπʼ ἐμὲ φρέαρ τὸ στόμα αὐτοῦ.
\vs{17}Εἰσάκουσόν μου Κύριε, ὅτι χρηστὸν τὸ ἔλεός σου, κατὰ τὸ πλῆθος τῶν οἰκτιρμῶν σου ἐπίβλεψον ἐπʼ ἐμε.
\vs{18}Καὶ μὴ ἀποστρέψῃς τὸ πρόσωπόν σου ἀπὸ τοῦ παιδός σου· ὅτι θλίβομαι, ταχὺ ἐπάκουσόν μου.
\vs{19}Πρόσχες τῇ ψυχῇ μου, καὶ λύτρωσαι αὐτὴν, ἕνεκα τῶν ἐχθρῶν μου ῥῦσαί με.

\vs{20}Σὺ γὰρ γινώσκεις τὸν ὀνειδισμόν μου, καὶ τὴν αἰσχύνην μου, καὶ τὴν ἐντροπήν μου· ἐναντίον σου πάντες οἱ θλίβοντές με.
\vs{21}Ὀνειδισμὸν προσεδόκησεν ἡ ψυχή μου καὶ ταλαιπωρίαν· καὶ ὑπέμεινα συλλυπούμενον, καὶ οὐχ ὑπῆρξε, καὶ παρακαλοῦντα, καὶ οὐχ εὗρον.
\vs{22}Καὶ ἔδωκαν εἰς τὸ βρῶμά μου χολὴν, καὶ εἰς τὴν δίψαν μου ἐπότισάν με ὄξος.
\vs{23}Γενηθήτω ἡ τράπεζα αὐτῶν ἐνώπιον αὐτῶν εἰς παγίδα, καὶ εἰς ἀνταπόδοσιν, καὶ εἰς σκάνδαλον.
\vs{24}Σκοτισθήτωσαν οἱ ὀφθαλμοὶ αὐτῶν τοῦ μὴ βλέπειν, καὶ τὸν νῶτον αὐτῶν διαπαντὸς σύγκαμψον.
\vs{25}Ἔκχεον ἐπʼ αὐτοὺς τὴν ὀργήν σου, καὶ ὁ θυμὸς τῆς ὀργῆς σου καταλάβοι αὐτούς.
\vs{26}Γενηθήτω ἡ ἔπαυλις αὐτῶν ἠρημωμένη, καὶ ἐν τοῖς σκηνώμασιν αὐτῶν μὴ ἔστω ὁ κατοικῶν·
\vs{27}ὅτι ὃν σὺ ἐπάταξας, αὐτοὶ κατεδίωξαν, καὶ ἐπὶ τὸ ἄλγος τῶν τραυμάτων μου προσέθηκαν.
\vs{28}Πρόσθες ἀνομίαν ἐπὶ τὴν ἀνομίαν αὐτῶν, καὶ μὴ εἰσελθέτωσαν ἐν δικαιοσύνῃ σου.
\vs{29}Ἐξαλειφθήτωσαν ἐκ βίβλου ζώντων, καὶ μετὰ δικαίων μὴ γραφήτωσαν.

\vs{30}Πτωχὸς καὶ ἀλγῶν εἰμι ἐγὼ, καὶ ἡ σωτηρία τοῦ προσώπου σου ἀντελάβετό μου.
\vs{31}Αἰνέσω τὸ ὄνομα τοῦ Θεοῦ μου μετʼ ᾠδῆς, μεγαλυνῶ αὐτὸν ἐν αἰνέσει·
\vs{32}καὶ ἀρέσει τῷ Θεῷ ὑπὲρ μόσχον νέον κέρατα ἐκφέροντα καὶ ὁπλάς.
\vs{33}Ἰδέτωσαν πτωχοὶ καὶ εὐφρανθήτωσαν· ἐκζητήσατε τὸν Θεὸν, καὶ ζήσεσθε.
\vs{34}Ὅτι εἰσήκουσε τῶν πενήτων ὁ Κύριος, καὶ τοὺς πεπεδημένους αὐτοῦ οὐκ ἐξουδένωσεν.
\vs{35}Αἰνεσάτωσαν αὐτὸν οἱ οὐρανοὶ καὶ ἡ γῆ, θάλασσα καὶ πάντα τὰ ἕρποντα ἐν αὐτοῖς.
\vs{36}Ὅτι ὁ Θεὸς σώσει τὴν Σιὼν, καὶ οἰκοδομηθήσονται αἱ πόλεις τῆς Ἰουδαίας, καὶ κατοικήσουσιν ἐκεῖ, καὶ κληρονομήσουσιν αὐτήν.
\vs{37}Καὶ τὸ σπέρμα τῶν δούλων αὐτοῦ καθέξουσιν αὐτὴν, καὶ οἱ ἀγαπῶντες τὸ ὄνομα αὐτοῦ κατασκηνώσουσιν ἐν αὐτῇ.

\begin{psalmheading}{\ch{69}{70} Εἰς τὸ τέλος, τῷ Δαυὶδ εἰς ἀνάμνησιν, εἰς τὸ σῶσαί με Κύριον.}
\end{psalmheading}
\vs{2}Ὁ Θεὸς εἰς τὴν βοήθειάν μου πρόσχες.
\vs{3}Αἰσχυνθείησαν καὶ ἐντραπείησαν οἱ ζητοῦντες τὴν ψυχήν μου, ἀποστραφείησαν εἰς τὰ ὀπίσω, καὶ καταισχυνθείησαν οἱ βουλόμενοί μοι κατά·
\vs{4}Ἀποστραφείησαν παραυτίκα αἰσχυνόμενοι οἱ λέγοντές μοι, εὖγε, εὖγε.
\vs{5}Ἀγαλλιάσθωσαν καὶ εὐφρανθήτωσαν ἐπὶ σοὶ πάντες οἱ ζητοῦντές σε, καὶ λεγέτωσαν διαπαντὸς, μεγαλυνθήτω ὁ Θεὸς, οἱ ἀγαπῶντες τὸ σωτήριόν σου.
\vs{6}Ἐγὼ δὲ πτωχὸς καὶ πένης, ὁ Θεὸς βοήθησόν μοι· βοηθός μου, καὶ ῥύστης μου εἶ σὺ, Κύριε μὴ χρονίσῃς.

\begin{psalmheading}{\ch{70}{71} Τῷ Δαυὶδ υἱῶν Ἰωναδὰβ, καὶ τῶν πρώτων αἰχμαλωτισθέντων.}
\end{psalmheading}
Ἐπὶ σοὶ Κύριε ἤλπισα, μὴ καταισχυνθείην εἰς τὸν αἰῶνα.
\vs{2}Ἐν τῇ δικαιοσύνῃ σου ῥῦσαί με καὶ ἐξελοῦ με, κλῖνον πρὸς μὲ τὸ οὖς σου καὶ σῶσόν με.
\vs{3}Γενοῦ μοι εἰς Θεὸν ὑπερασπιστὴν, καὶ εἰς τόπον ὀχυρὸν τοῦ σῶσαί με, ὅτι στερέωμά μου καὶ καταφυγή μου εἶ σύ.
\vs{4}Ὁ Θεός μου ῥῦσαί με ἐκ χειρὸς ἁμαρτωλοῦ, ἐκ χειρὸς παρανομοῦντος καὶ ἀδικοῦντος.
\vs{5}Ὅτι σὺ εἶ ἡ ὑπομονή μου Κύριε, Κύριε ἡ ἐλπίς μου ἐκ νεότητός μου·
\vs{6}Ἐπὶ σὲ ἐπεστηρίχθην ἀπὸ γαστρὸς, ἐκ κοιλίας μητρός μου σύ μου εἶ σκεπαστής. ἐν σοὶ ἡ ὕμνησίς μου διαπαντός.

\vs{7}Ὡσεὶ τέρας ἐγενήθην τοῖς πολλοῖς, καὶ σὺ βοηθὸς κραταιός.
\vs{8}Πληρωθήτω τὸ στόμα μου αἰνέσεως, ὅπως ὑμνήσω τὴν δόξαν σου, ὅλην τὴν ἡμέραν τὴν μεγαλοπρέπειάν σου.
\vs{9}Μὴ ἀποῤῥίψῃς με εἰς καιρὸν γήρους, ἐν τῷ ἐκλείπειν τὴν ἰσχύν μου μὴ ἐγκαταλίπῃς με.
\vs{10}Ὅτι εἶπαν οἱ ἐχθροί μου ἐμοὶ, καὶ οἱ φυλάσσοντες τὴν ψυχήν μου ἐβουλεύσαντο ἐπιτοαυτὸ,
\vs{11}λέγοντες, ὁ Θεὸς ἐγκατέλιπεν αὐτὸν, καταδιώξατε καὶ καταλάβετε αὐτὸν, ὅτι οὐκ ἔστιν ὁ ῥυόμενος.
\vs{12}Ὁ Θεὸς μὴ μακρύνῃς ἀπʼ ἐμοῦ,
\vs{13}ὁ Θεός μου εἰς τὴν βοήθειάν μου πρόσχες. Αἰσχυνθήτωσαν καὶ ἐκλιπέτωσαν οἱ ἐνδιαβάλλοντες τὴν ψυχήν μου, περιβαλλέσθωσαν αἰσχύνην καὶ ἐντροπὴν οἱ ζητοῦντες τὰ κακά μοι.

\vs{14}Ἐγὼ δὲ διαπαντὸς ἐλπιῶ, καὶ προσθήσω ἐπὶ πᾶσαν τὴν αἴνεσίν σου.
\vs{15}Τὸ στόμα μου ἐξαγγελεῖ τὴν δικαιοσύνην σου, ὅλην τὴν ἡμέραν τὴν σωτηρίαν σου· ὅτι οὐκ ἔγνων πραγματείας.
\vs{16}Εἰσελεύσομαι ἐν δυναστείᾳ Κυρίου, Κύριε μνησθήσομαι τῆς δικαιοσύνης σου μόνου.
\vs{17}Ἐδίδαξας με ὁ Θεὸς ἐκ νεότητός μου, καὶ μέχρι νῦν ἀπαγγελῶ τὰ θαυμάσιά σου,
\vs{18}καὶ ἕως γήρους καὶ πρεσβείου· ὁ Θεὸς μὴ ἐγκαταλίπῃς με, ἕως ἂν ἀπαγγείλω τὸν βραχίονά σου πάσῃ τῇ γενεᾷ τῇ ἐρχομένῃ· Τὴν δυναστείαν σου,
\vs{19}καὶ τὴν δικαιοσύνην σου ὁ Θεὸς ἕως ὑψίστων, ἃ ἐποίησας μεγαλεῖα· ὁ Θεὸς τίς ὅμοιός σοι;

\vs{20}Ὅσας ἔδειξάς μοι θλίψεις πολλὰς καὶ κακάς; καὶ ἐπιστρέψας ἐζωοποίησάς με, καὶ ἐκ τῶν ἀβύσσων τῆς γῆς πάλιν ἀνήγαγές με.
\vs{21}Ἐπλεόνασας τὴν δικαιοσύνην σου, καὶ ἐπιστρέψας παρεκάλεσάς με, καὶ ἐκ τῶν ἀβύσσων τῆς γῆς πάλιν ἀνήγαγές με.
\vs{22}Καὶ γὰρ ἐγὼ ἐξομολογήσομαί σοι ἐν σκεύει ψαλμοῦ τὴν ἀλήθειάν σου ὁ Θεὸς, ψαλῶ σοι ἐν κιθάρᾳ ὁ ἅγιος τοῦ Ἰσραήλ.
\vs{23}Ἀγαλλιάσονται τὰ χείλη μου ὅταν ψάλω σοι, καὶ ἡ ψυχή μου ἣν ἐλυτρώσω.
\vs{24}Ἔτι δὲ καὶ ἡ γλῶσσά μου ὅλην τὴν ἡμέραν μελετήσει τὴν δικαιοσύνην σου, ὅταν αἰσχυνθῶσι καὶ ἐντραπῶσιν οἱ ζητοῦντες τὰ κακά μοι.

\begin{psalmheading}{\ch{71}{72} Εἰς Σαλωμών.}
\end{psalmheading}
Ὁ Θεὸς τὸ κρίμα σου τῷ βασιλεῖ δὸς, καὶ τὴν δικαιοσύνην σου τῷ υἱῷ τοῦ βασιλέως·
\vs{2}κρίνειν τὸν λαόν σου ἐν δικαιοσύνῃ, καὶ τοὺς πτωχούς σου ἐν κρίσει.

\vs{3}Ἀναλαβέτω τὰ ὄρη εἰρήνην τῷ λαῷ σου, καὶ οἱ βουνοί·
\vs{4}ἐν δικαιοσύνῃ κρινεῖ τοὺς πτωχοὺς τοῦ λαοῦ, καὶ σώσει τοὺς υἱοὺς τῶν πενήτων· καὶ ταπεινώσει συκοφάντην,
\vs{5}καὶ συμπαραμενεῖ τῷ ἡλίῳ, καὶ πρὸ τῆς σελήνης γενεὰς γενεῶν.
\vs{6}Καταβήσεται ὡς ὑετὸς ἐπὶ πόκον, καὶ ὡσεὶ σταγόνες στάζουσαι ἐπὶ τὴν γῆν.
\vs{7}Ἀνατελεῖ ἐν ταῖς ἡμέραις αὐτοῦ δικαιοσύνη, καὶ πλῆθος εἰρήνης ἕως οὗ ἀνταναιρεθῇ ἡ σελήνη.
\vs{8}Καὶ κατακυριεύσει ἀπὸ θαλάσσης ἕως θαλάσσης, καὶ ἀπὸ ποταμοῦ ἕως περάτων τῆς οἰκουμένης.
\vs{9}Ἐνώπιον αὐτοῦ προπεσοῦνται Αἰθίοπες, καὶ οἱ ἐχθροὶ αὐτοῦ χοῦν λείξουσι.
\vs{10}Βασιλεῖς Θαρσὶς καὶ αἱ νῆσοι δῶρα προσοίσουσι, βασιλεῖς Ἀράβων καὶ Σαβὰ δῶρα προσάξουσι.
\vs{11}Καὶ προσκυνήσουσιν αὐτῷ πάντες οἱ βασιλεῖς, πάντα τὰ ἔθνη δουλεύσουσιν αὐτῷ.
\vs{12}Ὅτι ἐῤῥύσατο πτωχὸν ἐκ δυνάστου, καὶ πένητα ᾧ οὐχ ὑπῆρχε βοηθός.
\vs{13}Φείσεται πτωχοῦ καὶ πένητος, καὶ ψυχὰς πένητων σώσει.
\vs{14}Ἐκ τόκου καὶ ἐξ ἀδικίας λυτρώσεται τὰς ψυχὰς αὐτῶν, καὶ ἔντιμον τὸ ὄνομα αὐτῶν ἐνώπιον αὐτοῦ.
\vs{15}Καὶ ζήσεται, καὶ δοθήσεται αὐτῷ ἐκ τοῦ χρυσίου τῆς Ἀραβίας, καὶ προσεύξονται περὶ αὐτοῦ διαπαντός· ὅλην τὴν ἡμέραν εὐλογήσουσιν αὐτόν.
\vs{16}Ἔσται στήριγμα ἐν τῇ γῇ ἐπʼ ἄκρων τῶν ὀρέων· ὑπεραρθήσεται ὑπὲρ τὸν Λίβανον ὁ καρπὸς αὐτοῦ, καὶ ἐξανθήσουσιν ἐκ πόλεως ὡσεὶ χόρτος τῆς γῆς.

\vs{17}Ἔστω τὸ ὄνομα αὐτοῦ εὐλογημένον εἰς τοὺς αἰῶνας, πρὸ τοῦ ἡλίου διαμενεῖ τὸ ὄνομα αὐτοῦ, καὶ εὐλογηθήσονται ἐν αὐτῷ πᾶσαι αἱ φυλαὶ τῆς γῆς· πάντα τὰ ἔθνη μακαριοῦσιν αὐτόν.

\vs{18}Εὐλογητὸς Κύριος ὁ Θεὸς τοῦ Ἰσραὴλ, ὁ ποιῶν θαυμάσια μόνος,
\vs{19}καὶ εὐλογητὸν τὸ ὄνομα τῆς δόξης αὐτοῦ εἰς τὸν αἰῶνα καὶ εἰς αἰῶνα τοῦ αἰῶνος· καὶ πληρωθήσεται τῆς δόξης αὐτοῦ πᾶσα ἡ γῆ. γένοιτο, γένοιτο.

\textit{Ἐξέλιπον οἱ ὕμνοι Δαυὶδ τοῦ υἱοῦ Ἰεσσαί.}

\begin{psalmheading}{\ch{72}{73} Ψαλμὸς τῷ Ἀσάφ.}
\end{psalmheading}
Ὡς ἀγαθὸς ὁ Θεὸς τῷ Ἰσραὴλ, τοῖς εὐθέσι καρδίᾳ.
\vs{2}Ἐμοῦ δὲ παρὰ μικρὸν ἐσαλεύθησαν οἱ πόδες, παρʼ ὀλίγον ἐξεχύθη τὰ διαβήματά μου.
\vs{3}Ὅτι ἐζήλωσα ἐπὶ τοῖς ἀνόμοις, εἰρήνην ἁμαρτωλῶν θεωρῶν.

\vs{4}Ὅτι οὐκ ἔστιν ἀνάνευσις ἐν τῷ θανάτῳ αὐτῶν, καὶ στερέωμα ἐν τῇ μάστιγι αὐτῶν.
\vs{5}Ἐν κόποις ἀνθρώπων οὐκ εἰσὶ, καὶ μετὰ ἀνθρώπων οὐ μαστιγωθήσονται.
\vs{6}Διὰ τοῦτο ἐκράτησεν αὐτοὺς ἡ ὑπερηφανία, περιεβάλοντο ἀδικίαν καὶ ἀσέβειαν αὐτῶν.
\vs{7}Ἐξελεύσεται ὡς ἐκ στέατος ἡ ἀδικία αὐτῶν· διῆλθον εἰς διάθεσιν καρδίας.
\vs{8}Διενοήθησαν, καὶ ἐλάλησαν ἐν πονηρίᾳ, ἀδικίαν εἰς τὸ ὕψος ἐλάλησαν.
\vs{9}Ἔθεντο εἰς οὐρανὸν τὸ στόμα αὐτῶν, καὶ ἡ γλῶσσα αὐτῶν διῆλθεν ἐπὶ τῆς γῆς.
\vs{10}Διὰ τοῦτο ἐπιστρέψει ὁ λαός μου ἐνταῦθα, καὶ ἡμέραι πλήρεις εὑρεθήσονται ἐν αὐτοῖς.
\vs{11}Καὶ εἶπαν, πῶς ἔγνω ὁ Θεὸς, καὶ εἰ ἔστι γνῶσις ἐν τῷ ὑψίστῳ;
\vs{12}Ἰδοὺ οὗτοι οἱ ἁμαρτωλοὶ καὶ εὐθηνοῦντες εἰς τὸν αἰῶνα, κατέσχον πλούτου.

\vs{13}Καὶ εἶπα, ἄρα ματαίως ἐδικαίωσα τὴν καρδίαν μου, καὶ ἐνιψάμην ἐν ἀθώοις τὰς χεῖράς μου·
\vs{14}Καὶ ἐγενόμην μεμαστιγωμένος ὅλην τὴν ἡμέραν, καὶ ὁ ἔλεγχός μου εἰς τὰς πρωΐας.
\vs{15}Εἰ ἔλεγον, διηγήσομαι οὕτως, ἰδοὺ τῇ γενεᾷ τῶν υἱῶν σου ἠσυνθέτηκα.
\vs{16}Καὶ ὑπέλαβον τοῦ γνῶναι, τοῦτο κόπος ἐστὶν ἐναντίον μου,
\vs{17}ἕως εἰσέλθω εἰς τὸ ἁγιαστήριον τοῦ Θεοῦ, συνῶ εἰς τὰ ἔσχατα.

\vs{18}Πλὴν διὰ τὰς δολιότητας ἔθου αὐτοῖς, κατέβαλες αὐτοὺς ἐν τῷ ἐπαρθῆναι.
\vs{19}Πῶς ἐγένοντο εἰς ἐρήμωσιν; ἐξάπινα ἐξέλιπον, ἀπώλοντο διὰ τὴν ἀνομίαν αὐτῶν.
\vs{20}Ὡσεὶ ἐνύπνιον ἐξεγειρομένου, Κύριε ἐν τῇ πόλει σου τὴν εἰκόνα αὐτῶν ἐξουδενώσεις.

\vs{21}Ὅτι ηὐφράνθη ἡ καρδία μου, καὶ οἱ νεφροί μου ἠλλοιώθησαν.
\vs{22}Κᾀγὼ ἐξουδενωμένος, καὶ οὐκ ἔγνων, κτηνώδης ἐγενόμην παρὰ σοὶ,
\vs{23}κᾀγὼ διαπαντὸς μετὰ σου· ἐκράτησας τῆς χειρὸς τῆς δεξιᾶς μου,
\vs{24}ἐν τῇ βουλῇ σου ὡδήγησάς με, καὶ μετὰ δόξης προσελάβου με.
\vs{25}Τί γάρ μοι ὑπάρχει ἐν τῷ οὐρανῷ, καὶ παρὰ σοῦ τί ἠθέλησα ἐπὶ τῆς γῆς;
\vs{26}Ἐξέλιπεν ἡ καρδία μου καὶ ἡ σάρξ μου, ὁ Θεὸς τῆς καρδίας μου, καὶ ἡ μερίς μου ὁ Θεὸς εἰς τὸν αἰῶνα.

\vs{27}Ὅτι ἰδοὺ οἱ μακρύνοντες ἑαυτοὺς ἀπὸ σοῦ, ἀπολοῦνται· ἐξωλόθρευσας πάντα τὸν πορνεύοντα ἀπὸ σοῦ.
\vs{28}Ἐμοὶ δὲ τὸ προσκολλᾶσθαι τῷ Θεῷ ἀγαθόν ἐστι, τίθεσθαι ἐν τῷ Κυρίῳ τὴν ἐλπίδα μου· τοῦ ἐξαγγεῖλαι πάσας τὰς αἰνέσεις σου ἐν ταῖς πύλαις τῆς θυγατρὸς Σιών.

\begin{psalmheading}{\ch{73}{74} Συνέσεως τῷ Ἀσάφ.}
\end{psalmheading}
Ἱνατί ἀπόσω ὁ Θεὸς εἰς τέλος; ὠργίσθη ὁ θυμός σου ἐπὶ πρόβατα νομῆς σου;
\vs{2}Μνήσθητι τῆς συναγωγῆς σου ἧς ἐκτήσω ἀπʼ ἀρχῆς· ἐλυτρώσω ῥάβδον κληρονομίας σου· ὄρος Σιὼν τοῦτο ὃ κατεσκήνωσας ἐν αὐτῷ.
\vs{3}Ἔπαρον τὰς χεῖράς σου ἐπὶ τὰς ὑπερηφανίας αὐτῶν εἰς τέλος· ὅσα ἐπονηρεύσατο ὁ ἐχθρὸς ἐν τοῖς ἁγίοις σου.

\vs{4}Καὶ ἐνεκαυχήσαντο οἱ μισοῦντές σε ἐν μέσῳ τῆς ἑορτῆς σου·
\vs{5}ἔθεντο τὰ σημεῖα αὐτῶν σημεῖα, καὶ οὐκ ἔγνωσαν, ὡς εἰς τὴν εἴσοδον ὑπεράνω· ὡς ἐν δρυμῷ ξύλων ἀξίναις ἐξέκοψαν
\vs{6}τὰς θύρας αὐτῆς ἐπιτοαυτὸ, ἐν πελέκει καὶ λαξευτηρίῳ κατέῤῥαξαν αὐτήν.
\vs{7}Ἐνεπύρισαν ἐν πυρὶ τὸ ἁγιαστήριόν σου εἰς τὴν γῆν, ἐβεβήλωσαν τὸ σκήνωμα τοῦ ὀνόματός σου.
\vs{8}Εἶπαν ἐν τῇ καρδίᾳ αὐτῶν, ἡ συγγένεια αὐτῶν ἐπιτοαυτὸ, δεῦτε, καταπαύσωμεν τὰς ἑορτὰς Κυρίου ἀπὸ τῆς γῆς.
\vs{9}Τὰ σημεῖα ἡμῶν οὐκ εἴδομεν, οὐκ ἔστιν ἔτι προφήτης, καὶ ἡμᾶς οὐ γνώσεται ἔτι.

\vs{10}Ἕως πότε ὁ Θεὸς, ὀνειδιεῖ ὁ ἐχθρὸς, παροξυνεῖ ὁ ὑπεναντίος τὸ ὄνομά σου εἰς τέλος;
\vs{11}Ἱνατί ἀποστρέφεις τὴν χεῖρά σου, καὶ τὴν δεξιάν σου ἐκ μέσου τοῦ κόλπου σου εἰς τέλος;
\vs{12}Ὁ δὲ Θεὸς βασιλεὺς ἡμῶν πρὸ αἰῶνος, εἰργάσατο σωτηρίαν ἐν μέσῳ τῆς γῆς.
\vs{13}Σὺ ἐκραταίωσας ἐν τῇ δυνάμει σου τὴν θάλασσαν, σὺ συνέτριψας τὰς κεφαλὰς τῶν δρακόντων ἐπὶ τοῦ ὕδατος. Σὺ συνέτριψας τὰς κεφαλὰς τοῦ δράκοντος,
\vs{14}ἔδωκας αὐτὸν βρῶμα λαοῖς τοῖς Αἰθίοψι.
\vs{15}Σὺ διέῤῥηξας πηγὰς καὶ χειμάῤῥους, σὺ ἐξήρανας ποταμοὺς ἠθάμ.
\vs{16}Σή ἐστιν ἡ ἡμέρα, καὶ σή ἐστιν ἡ νὺξ, σὺ κατηρτίσω ἥλιον καὶ σελήνην.
\vs{17}Σὺ ἐποίησας πάντα τὰ ὅρια τῆς γῆς, θέρος καὶ ἔαρ σὺ ἐποίησας.

\vs{18}Μνήσθητι ταύτης τῆς κτίσεώς σου· ἐχθρὸς ὠνείδισε τὸν Κύριον, καὶ λαὸς ἄφρων παρώξυνε τὸ ὄνομά σου.
\vs{19}Μὴ παραδῷς τοῖς θηρίοις ψυχὴν ἐξομολογουμένην σοι, τῶν ψυχῶν τῶν πενήτων σου μὴ ἐπιλάθῃ εἰς τέλος.
\vs{20}Ἐπίβλεψον εἰς τὴν διαθήκην σου, ὅτι ἐπληρωθήσαν οἱ ἐσκοτωμένοι τῆς γῆς οἴκων ἀνομιῶν.
\vs{21}Μὴ ἀποστραφήτω τεταπεινωμένος καὶ κατῃσχυμμένος, πτωχὸς καὶ πένης αἰνέσουσι τὸ ὄνομά σου.
\vs{22}Ἀνάστα ὁ Θεὸς, δίκασον τὴν δίκην σου, μνήσθητι τῶν ὀνειδισμῶν σου τῶν ὑπὸ ἄφρονος ὅλην τὴν ἡμέραν.
\vs{23}Μὴ ἐπιλάθῃ τῆς φωνῆς τῶν ἱκετῶν σου, ἡ ὑπερηφανία τῶν μισούντων σε ἀναβαίη διαπαντὸς πρὸς σέ.

\begin{psalmheading}{\ch{74}{75} Εἰς τὸ τέλος, μὴ διαφθείρῃς, ψαλμὸς ᾠδῆς τῷ Ἀσάφ.}
\end{psalmheading}
\vs{2}Ἐξομολογησόμεθα σοι ὁ Θεὸς, ἐξομολογησόμεθα, καὶ ἐπικαλεσόμεθα τὸ ὄνομά σου· διηγήσομαι πάντα τὰ θαυμάσιά σου.
\vs{3}Ὅταν λάβω καιρὸν, ἐγὼ εὐθύτητας κρινῶ.
\vs{4}Ἐτάκη ἡ γῆ, καὶ πάντες οἱ κατοικοῦντες αὐτὴν, ἐγὼ ἐστερέωσα τοὺς στύλους αὐτῆς· διάψαλμα.

\vs{5}Εἶπα τοῖς παρανομοῦσι, μὴ παρανομεῖν, καὶ τοῖς ἁμαρτάνουσι, μὴ ὑψοῦτε κέρας.
\vs{6}Μὴ ἐπαίρετε εἰς ὕψος τὸ κέρας ὑμῶν, μὴ λαλεῖτε κατὰ τοῦ Θεοῦ ἀδικίαν.
\vs{7}Ὅτι οὔτε ἀπὸ ἐξόδων, οὔτε ἀπὸ δυσμῶν, οὔτε ἀπὸ ἐρήμων ὀρέων,
\vs{8}ὅτι ὁ Θεὸς κριτής ἐστι· τοῦτον ταπεινοῖ, καὶ τοῦτον ὑψοῖ.
\vs{9}Ὅτι ποτήριον ἐν χειρὶ Κυρίου, οἴνου ἀκράτου πλῆρες κεράσματος· καὶ ἔκλινεν ἐκ τούτου εἰς τοῦτο, πλὴν ὁ τρυγίας αὐτοῦ οὐκ ἐξεκενώθη· πίονται πάντες οἱ ἁμαρτωλοὶ τῆς γῆς.

\vs{10}Ἐγὼ δὲ ἀγαλλιάσομαι εἰς τὸν αἰῶνα, ψαλῶ τῷ Θεῷ Ἰακώβ.
\vs{11}Καὶ πάντα τὰ κέρατα τῶν ἁμαρτωλῶν συγκλάσω, καὶ ὑψωθήσεται τὰ κέρατα τοῦ δικαίου.

\begin{psalmheading}{\ch{75}{76} Εἰς τὸ τέλος ἐν ὕμνοις, ψαλμὸς τῷ Ἀσάφ· ᾠδὴ πρὸς τὸν Ἀσσύριον.}
\end{psalmheading}
\vs{2}Γνωστὸς ἐν τῇ Ἰουδαίᾳ ὁ Θεὸς, ἐν τῷ Ἰσραὴλ μέγα τὸ ὄνομα αὐτοῦ.
\vs{3}Καὶ ἐγενήθη ἐν εἰρήνῃ ὁ τόπος αὐτοῦ, καὶ τὸ κατοικητήριον αὐτοῦ ἐν Σιών.
\vs{4}Ἐκεῖ συνέτριψε τὰ κράτη τῶν τόξων, ὅπλον καὶ ῥομφαίαν καὶ πόλεμον· διάψαλμα.

\vs{5}Φωτίζεις σὺ θαυμαστῶς ἀπὸ ὀρέων αἰωνίων,
\vs{6}ἐταράχθησαν πάντες οἱ ἀσύνετοι τῇ καρδίᾳ· ὕπνωσαν ὕπνον αὐτῶν, καὶ οὐχ εὗρον οὐδὲν πάντες οἱ ἄνδρες τοῦ πλούτου ταῖς χερσὶν αὐτῶν.
\vs{7}Ἀπὸ ἐπιτιμήσεώς σου, ὁ Θεὸς Ἰακὼβ, ἐνύσταξαν οἱ ἐπιβεβηκότες τοὺς ἵππους.
\vs{8}Σὺ φοβερὸς εἶ, καὶ τίς ἀντιστήσεταί σοι ἀπὸ τῆς ὀργῆς σου;
\vs{9}Ἐκ τοῦ οὐρανοῦ ἠκούτισας κρίσιν, γῆ ἐφοβήθη καὶ ἡσύχασεν,
\vs{10}ἐν τῷ ἀναστῆναι εἰς κρίσιν τὸν Θεὸν, τοῦ σῶσαι πάντας τοὺς πρᾳεῖς τῇ καρδίᾳ. διάψαλμα.

\vs{11}Ὅτι ἐνθύμιον ἀνθρώπου ἐξομολογήσεταί σοι, καὶ ἐγκατάλειμμα ἐνθυμίου ἑορτάσει σοι.
\vs{12}Εὔξασθε καὶ ἀπόδοτε Κυρίῳ τῷ Θεῷ ἡμῶν, πάντες οἱ κύκλῳ αὐτοῦ οἴσουσι δῶρα·
\vs{13}τῷ φοβερῷ καὶ ἀφαιρουμένῳ πνεύματα ἀρχόντων, φοβερῷ παρὰ τοῖς βασιλεῦσι τῆς γῆς.

\begin{psalmheading}{\ch{76}{77} Εἰς τὸ τέλος, ὑπὲρ Ἰδιθοὺν ψαλμὸς τῷ Ἀσάφ.}
\end{psalmheading}
\vs{2}Φωνῇ μου πρὸς Κύριον ἐκέκραξα, καὶ ἡ φωνή μου πρὸς τὸν Θεὸν, καὶ προσέσχε μοι.
\vs{3}Ἐν ἡμέρᾳ θλίψεώς μου τὸν Θεὸν ἐξεζήτησα, ταῖς χερσί μου νυκτὸς ἐναντίον αὐτοῦ, καὶ οὐκ ἠπατήθην· ἀπηνῄνατο παρακληθῆναι ἡ ψυχή μου·
\vs{4}Ἐμνήσθην τοῦ Θεοῦ, καὶ εὐφράνθην, ἠδολέσχησα, καὶ ὠλιγοψύχησε τὸ πνεῦμά μου· διάψαλμα.
\vs{5}Προκατελάβοντο φυλακὰς πάντες οἱ ἐχθροί μου, ἐταράχθην καὶ οὐκ ἐλάλησα.

\vs{6}Διελογισάμην ἡμέρας ἀρχαίας, καὶ ἔτη αἰώνια
\vs{7}ἐμνήσθην, καὶ ἐμελέτησα· νυκτὸς μετὰ τῆς καρδίας μου ἠδολέσχουν, καὶ ἔσκαλλον τὸ πνεῦμά μου.
\vs{8}Μὴ εἰς τοὺς αἰῶνας ἀπώσεται Κύριος, καὶ οὐ προσθήσει τοῦ εὐδοκῆσαι ἔτι;
\vs{9}Ἢ εἰς τέλος ἀποκόψει τὸ ἔλεος ἀπὸ γενεᾶς καὶ γενεᾶς;
\vs{10}Ἢ ἐπιλήσεται τοῦ οἰκτειρῆσαι ὁ Θεὸς, ἢ συνέξει ἐν τῇ ὀργῇ αὐτοῦ τοὺς οἰκτιρμοὺς αὐτοῦ; διάψαλμα.

\vs{11}Καὶ εἶπα, νῦν ἠρξάμην, αὕτη ἡ ἀλλοίωσις τῆς δεξιᾶς τοῦ ὑψίστου.
\vs{12}Ἐμνήσθην τῶν ἔργων Κυρίου, ὅτι μνησθήσομαι ἀπὸ τῆς ἀρχῆς τῶν θαυμασίων σου,
\vs{13}καὶ μελετήσω ἐν πᾶσι τοῖς ἔργοις σου, καὶ ἐν τοῖς ἐπιτηδεύμασί σου ἀδολεσχήσω.

\vs{14}Ὁ Θεὸς ἐν τῷ ἁγίῳ ἡ ὁδός σου, τίς θεὸς μέγας ὡς ὁ Θεὸς ἡμῶν;
\vs{15}Σὺ εἶ ὁ Θεὸς ὁ ποιῶν θαυμάσια, ἐγνώρισας ἐν τοῖς λαοῖς τὴν δύναμίν σου·
\vs{16}ἐλυτρώσω ἐν τῷ βραχίονί σου τὸν λαόν σου, τοὺς υἱοὺς Ἰακὼβ καὶ Ἰωσήφ· διάψαλμα.
\vs{17}Εἴδοσάν σε ὕδατα ὁ Θεὸς, εἴδοσάν σε ὕδατα καὶ ἐφοβήθησαν, καὶ ἐταράχθησαν ἄβυσσοι.
\vs{18}Πλῆθος ἤχους ὑδάτων, φωνὴν ἔδωκαν αἱ νεφέλαι· καὶ γὰρ τὰ βέλη σου διαπορεύονται.
\vs{19}Φωνὴ τῆς βροντῆς σου ἐν τῷ τροχῷ· ἔφαναν αἱ ἀστραπαί σου τῇ οἰκουμένῃ, ἐσαλεύθη καὶ ἔντρομος ἐγενήθη ἡ γῆ.
\vs{20}Ἐν τῇ θαλάσσῃ ἡ ὁδός σου, καὶ αἱ τρίβοι σου ἐν ὕδασι πολλοῖς, καὶ τὰ ἴχνη σου οὐ γνωσθήσονται.
\vs{21}Ὡδήγησας ὡς πρόβατα τὸν λαόν σου ἐν χειρὶ Μωυσῆ καὶ Ἀαρών.

\begin{psalmheading}{\ch{77}{78} Συνέσεως τῷ Ἀσάφ.}
\end{psalmheading}
Προσέχετε λαός μου τὸν νόμον μου, κλίνατε τὸ οὖς ὑμῶν εἰς τὰ ῥήματα τοῦ στόματός μου.
\vs{2}Ἀνοίξω ἐν παραβολαῖς τὸ στόμα μου, φθέγξομαι προβλήματα ἀπʼ ἀρχῆς.
\vs{3}Ὅσα ἠκούσαμεν καὶ ἔγνωμεν αὐτὰ, καὶ οἱ πατέρες ἡμῶν διηγήσαντο ἡμῖν.
\vs{4}Οὐκ ἐκρύβη ἀπὸ τῶν τέκνων αὐτῶν εἰς γενεὰν ἑτέραν, ἀπαγγέλλοντες τὰς αἰνέσεις Κυρίου καὶ τὰς δυναστείας αὐτοῦ, καὶ τὰ θαυμάσια αὐτοῦ ἃ ἐποίησε.

\vs{5}Καὶ ἀνέστησε μαρτύριον ἐν Ἰακὼβ, καὶ νόμον ἔθετο ἐν Ἰσραήλ· ὃν ἐνετείλατο τοῖς πατράσιν ἡμῶν, γνωρίσαι αὐτὸν τοῖς υἱοῖς αὐτῶν,
\vs{6}ὅπως ἂν γνῷ γενεὰ ἑτέρα, υἱοὶ οἱ τεχθησόμενοι, καὶ ἀναστήσονται καὶ ἀπαγγελοῦσιν αὐτὰ τοῖς υἱοῖς αὐτῶν·
\vs{7}ἵνα θῶνται ἐπὶ τὸν Θεὸν τὴν ἐλπίδα αὐτῶν, καὶ μὴ ἐπιλάθωνται τῶν ἔργων τοῦ Θεοῦ, καὶ τὰς ἐντολὰς αὐτοῦ ἐκζητήσωσιν.
\vs{8}Ἵνα μὴ γένωνται ὡς οἱ πατέρες αὐτῶν, γενεὰ σκολιὰ καὶ παραπικραίνουσα, γενεὰ ἥτις οὐ κατεύθυνεν ἐν τῇ καρδίᾳ αὐτῆς, καὶ οὐκ ἐπιστώθη μετὰ τοῦ Θεοῦ τὸ πνεῦμα αὐτῆς.

\vs{9}Υἱοὶ Ἐφραὶμ ἐντεινοντες καὶ βάλλοντες τόξον, ἐστράφησαν ἐν ἡμέρᾳ πολέμου.
\vs{10}Οὐκ ἐφύλαξαν τὴν διαθήκην τοῦ Θεοῦ, καὶ ἐν τῷ νόμῳ αὐτοῦ οὐκ ἤθελον πορεύεσθαι.
\vs{11}Καὶ ἐπελάθοντο τῶν εὐεργεσιῶν αὐτοῦ καὶ τῶν θαυμασίων αὐτοῦ, ὧν ἔδειξεν αὐτοῖς·
\vs{12}ἐναντίον τῶν πατέρων αὐτῶν ἃ ἐποίησε θαυμάσια, ἐν γῇ Αἰγύπτῳ, ἐν πεδίῳ Τάνεως.
\vs{13}Διέῤῥηξε θάλασσαν, καὶ διήγαγεν αὐτοὺς· ἔστησεν ὕδατα ὡσεὶ ἀσκόν.
\vs{14}Καὶ ὡδήγησεν αὐτοὺς ἐν νεφέλῃ ἡμέρας, καὶ ὅλην τὴν νύκτα ἐν φωτισμῷ πυρός.
\vs{15}Διέῤῥηξε πέτραν ἐν ἐρήμῳ, καὶ ἐπότισεν αὐτοὺς ὡς ἐν ἀβύσσῳ πολλῇ.
\vs{16}Καὶ ἐξήγαγεν ὕδωρ ἐκ πέτρας, καὶ κατήγαγεν ὡς ποταμοὺς ὕδατα.

\vs{17}Καὶ προσέθεντο ἔτι τοῦ ἁμαρτάνειν αὐτῷ· παρεπικραναν τὸν ὕψιστον ἐν ἀνύδρῳ,
\vs{18}καὶ ἐξεπείρασαν τὸν Θεὸν ἐν ταῖς καρδίαις αὐτῶν, τοῦ αἰτῆσαι βρώματα ταῖς ψυχαῖς αὐτῶν.
\vs{19}Καὶ κατελάλησαν τοῦ Θεοῦ, καὶ εἶπαν, μὴ δυνήσεται ὁ Θεὸς ἑτοιμάσαι τράπεζαν ἐν ἐρήμῳ;
\vs{20}Ἐπεὶ ἐπάταξε πέτραν, καὶ ἐῤῥύησαν ὕδατα, καὶ χείμαῤῥοι κατεκλύσθησαν· μὴ καὶ ἄρτον δυνήσεται δοῦναι; ἢ ἑτοιμάσαι τράπεζαν τῷ λαῷ αὐτοῦ;

\vs{21}Διὰ τοῦτο ἤκουσε Κύριος καὶ ἀνεβάλετο, καὶ πῦρ ἀνήφθη ἐν Ἰακὼβ, καὶ ὀργὴ ἀνέβη ἐπὶ τὸν Ἰσραήλ.
\vs{22}Ὅτι οὐκ ἐπίστευσαν ἐν τῷ Θεῷ, οὐδὲ ἤλπισαν ἐπὶ τὸ σωτήριον αὐτοῦ.
\vs{23}Καὶ ἐνετείλατο νεφέλαις ὑπεράνωθεν, καὶ θύρας οὐρανοῦ ἀνέῳξε·
\vs{24}καὶ ἔβρεξεν αὐτοῖς μάννα φαγεῖν, καὶ ἄρτον οὐρανοῦ ἔδωκεν αὐτοῖς.
\vs{25}Ἄρτον ἀγγέλων ἔφαγεν ἄνθρωπος, ἐπισιτισμὸν ἀπέστειλεν αὐτοῖς εἰς πλησμονήν.

\vs{26}Ἀπῇρε Νότον ἐξ οὐρανοῦ, καὶ ἐπήγαγεν ἐν τῇ δυναστείᾳ
\vs{27}αὐτοῦ Λίβα. Καὶ ἔβρεξεν ἐπ αὐτοὺς ὡσεὶ χοῦν σάρκας, καὶ ὡσεῖ ἄμμον θαλασσῶν πετεινὰ πτερωτά.
\vs{28}Καὶ ἐπέπεσον εἰς μέσον τῆς παρεμβολῆς αὐτῶν, κύκλῳ τῶν σκηνωμάτων αὐτῶν.
\vs{29}Καὶ ἐφάγοσαν καὶ ἐνεπλήσθησαν σφόδρα, καὶ τὴν ἐπιθυμίαν αὐτῶν ἤνεγκεν αὐτοῖς.

\vs{30}Οὐκ ἐστερήθησαν ἀπὸ τῆς ἐπιθυμίας αὐτῶν· ἔτι τῆς βρώσεως
\vs{31}αὐτῶν οὔσης ἐν τῷ στόματι αὐτῶν, καὶ ὀργὴ τοῦ Θεοῦ ἀνέβη ἐπʼ αὐτοὺς, καὶ ἀπέκτεινεν ἐν τοῖς πίοσιν αὐτῶν, καὶ τοὺς ἐκλεκτοὺς τοῦ Ἰσραὴλ συνεπόδισεν.

\vs{32}ʼΕν πᾶσι τούτοις ἥμαρτον ἔτι, καὶ οὐκ ἐπίστευσαν τοῖς
\vs{33}θαυμασίοις αὐτοῦ. Καὶ ἐξέλιπον ἐν ματαιότητι αἱ ἡμέραι αὐτῶν, καὶ τὰ ἔτη αὐτῶν μετὰ σπουδῆς.

\vs{34}Ὅταν ἀπέκτεινεν αὐτοὺς, ἐζήτουν αὐτὸν, καὶ ἐπέστρεφον καὶ
\vs{35}ὤρθριζον πρὸς τὸν Θεόν. Καὶ ἐμνήσθησαν ὅτι ὁ Θεὸς βοηθὸς αὐτῶν ἐστι, καὶ ὁ Θεὸς ὁ ὕψιστος λυτρωτὴς αὐτῶν ἐστι.
\vs{36}Καὶ ἠγάπησαν αὐτὸν ἐν τῷ στόματι αὐτῶν, καὶ τῇ γλώσσῃ αὐτῶν
\vs{37}ἐψεύσαντο αὐτῷ· ἡ δὲ καρδία αὐτῶν οὐκ εὐθεῖα μετʼ αὐτοῦ, οὐδὲ ἐπιστώθησαν ἐν τῇ διαθήκῃ αὐτοῦ.

\vs{38}Αὐτὸς δέ ἐστιν οἰκτίρμων, καὶ ἱλάσεται ταῖς ἁμαρτίαις αὐτῶν, καὶ οὐ διαφθερεῖ· καὶ πληθυνεῖ τοῦ ἀποστρέψαι τὸν θυμὸν αὐτοῦ, καὶ οὐχὶ ἐκκαύσει πᾶσαν τὴν ὀργὴν αὐτοῦ.
\vs{39}Καὶ ἐμνήσθη ὅτι σάρξ εἰσι, πνεῦμα πορευόμενον καὶ οὐκ ἐπιστρέφον.

\vs{40}Ποσάκις παρεπίκραναν αὐτὸν ἐν τῇ ἐρήμῳ, παρώργισαν
\vs{41}αὐτὸν ἐν γῇ ἀνύδρῳ; Καὶ ἐπέστρεψαν καὶ ἐπείρασαν τὸν Θεὸν, καὶ τὸν ἅγιον τοῦ Ἰσραὴλ παρώξυναν.
\vs{42}Οὐκ ἐμνήσθησαν τῆς χειρὸς αὐτοῦ, ἡμέρας ἧς ἐλυτρώσατο αὐτοὺς ἐκ χειρὸς θλίβοντος·
\vs{43}Ὡς ἔθετο ἐν Αἰγύπτῳ τὰ σημεῖα αὐτοῦ, καὶ τὰ
\vs{44}τέρατα αὐτοῦ ἐν πεδίῳ Τάνεως· Καὶ μετέστρεψεν εἰς αἷμα τοὺς ποταμοὺς αὐτῶν, καὶ τὰ ὀμβρήματα αὐτῶν ὅπως μὴ πίωσιν·
\vs{45}ἐξαπέστειλεν εἰς αὐτοὺς κυνόμυιαν καὶ κατέφαγεν αὐτοὺς, καὶ βάτραχον, καὶ διέφθειρεν αὐτούς·
\vs{46}Καὶ ἔδωκε τῇ ἐρυσίβῃ τὸν καρπὸν αὐτῶν, καὶ τοὺς πόνους αὐτῶν τῇ ἀκρίδι.
\vs{47}Ἀπέκτεινεν ἐν χαλάζῃ τὴν ἄμπελον αὐτῶν, καὶ τὰς συκαμίνους αὐτῶν ἐν
\vs{48}τῇ πάχνῃ. Καὶ παρέδωκεν ἐν χαλάζῃ τὰ κτήνη αὐτῶν, καὶ τὴν ὕπαρξιν αὐτῶν τῷ πυρί.
\vs{49}Ἐξαπέστειλεν εἰς αὐτοὺς ὀργὴν θυμοῦ αὐτοῦ, θυμὸν καὶ ὀργὴν καὶ θλίψιν, ἀποστολὴν διʼ
\vs{50}ἀγγέλων πονηρῶν. Ὡδοποίησε τρίβον τῇ ὀργῇ αὐτοῦ, οὐκ ἐφείσατο ἀπὸ θανάτου τῶν ψυχῶν αὐτῶν, καὶ τὰ κτήνη αὐτῶν
\vs{51}εἰς θάνατον συνέκλεισε. Καὶ ἐπάταξε πᾶν πρωτότοκον ἐν γῇ Αἰγυπτῳ, ἀπαρχὴν πόνων αὐτῶν ἐν τοῖς σκηνώμασι Χάμ.
\vs{52}Καὶ ἀπῇρεν ὡς πρόβατα τὸν λαὸν αὐτοῦ, ἤγαγεν αὐτοὺς ὡσεὶ
\vs{53}ποίμνιον ἐν ἐρήμῳ. Καὶ ὡδήγησεν αὐτοὺς ἐν ἐλπίδι, καὶ οὐκ ἐδειλίασον, καὶ τοὺς ἐχθροὺς αὐτῶν ἐκάλυψε θάλασσα.
\vs{54}Καὶ εἰσήγαγεν αὐτοὺς εἰς ὄρος ἁγιάσματος αὐτοῦ, ὄρος τοῦτο ὃ
\vs{55}ἐκτήσατο ἡ δεξιὰ αὐτοῦ. Καὶ ἐξέβαλεν ἀπὸ προσώπου αὐτῶν ἔθνη, καὶ ἐκληροδότησεν αὐτοὺς ἐν σχοινίῳ κληροδοσίας, καὶ κατεσκήνωσεν ἐν τοῖς σκηνώμασιν αὐτῶν τὰς φυλὰς τοῦ Ἰσραήλ.

\vs{56}Καὶ ἐπείρασαν καὶ παρεπίκραναν τὸν Θεὸν τὸν ὕψιστον, καὶ τὰ μαρτύρια αὐτοῦ οὐκ ἐφυλάξαντο.
\vs{57}Καὶ ἀπέστρεψαν, καὶ ἠσυνθέτησαν καθὼς καὶ οἱ πατέρες αὐτῶν, μετεστράφησαν εἰς τόξον στρεβλόν.
\vs{58}Καὶ παρώργισαν αὐτὸν ἐπὶ τοῖς βουνοῖς αὐτῶν, καὶ ἐν τοῖς γλυπτοῖς αὐτῶν παρεζήλωσαν αὐτόν.

\vs{59}Ἤκουσεν ὁ Θεὸς καὶ ὑπερεῖδε, καὶ ἐξουδένωσε σφόδρα τὸν
\vs{60}ʼΙσραήλ. Καὶ ἀπώσατο τὴν σκηνὴν Σηλὼμ, σκήνωμα αὐτοῦ
\vs{61}οὗ κατεσκήνωσεν ἐν ἀνθρώποις. Καὶ παρέδωκεν εἰς αἰχμαλωσίαν τὴν ἰσχὺν αὐτῶν, καὶ τὴν καλλονὴν αὐτῶν εἰς χεῖρα
\vs{62}ἐχθρου. Καὶ συνέκλεισεν εἰς ῥομφαίαν τὸν λαὸν αὐτοῦ, καὶ τὴν κληρονομίαν αὐτοῦ ὑπερεῖδε·
\vs{63}Τοὺς νεανίσκους αὐτῶν κατέφαγε πῦρ, καὶ αἱ παρθένοι αὐτῶν οὐκ ἐπένθησαν. Οἱ
\vs{64}ἱερεῖς αὐτῶν ἐν ῥομφαίᾳ ἔπεσον, καὶ αἱ χῆραι αὐτῶν οὐ κλαυσθήσονται.

\vs{65}Καὶ ἐξηγέρθη ὡς ὁ ὑπνῶν Κύριος, ὡς δυνατὸς κεκραιπαληκὼς
\vs{66}ἐξ οἴνου. Καὶ ἐπάταξε τοὺς ἐχθροὺς αὐτοῦ εἰς τὰ ὀπίσω, ὄνειδος αἰώνιον ἔδωκεν αὐτοῖς.

\vs{67}Καὶ ἀπώσατο τὸ σκήνωμα Ἰωσὴφ, καὶ τὴν φυλὴν Ἐφραὶμ
\vs{68}οὐκ ἐξελέξατο. Καὶ ἐξελέξατο τὴν φυλὴν Ἰούδα, τὸ ὄρος τὸ
\vs{69}Σιὼν, ὃ ἠγάπησε. Καὶ ᾠκοδόμησεν ὡς μονοκερώτων τὸ ἁγίασμα αὐτοῦ, ἐν τῇ γῇ ἐθεμελίωσεν αὐτὴν εἰς τὸν αἰῶνα.
\vs{70}Καὶ ἐξελέξατο Δαυὶδ τὸν δοῦλον αὐτοῦ, καὶ ἀνέλαβεν αὐτὸν ἐκ τῶν
\vs{71}ποιμνίων τῶν προβάτων. Ἐξόπισθεν τῶν λοχευομένων ἔλαβεν αὐτὸν, ποιμαίνειν Ἰακὼβ τὸν δοῦλον αὐτοῦ, καὶ Ἰσραὴλ τὴν
\vs{72}κληρονομίαν αὐτοῦ. Καὶ ἐποίμανεν αὐτοὺς ἐν τῇ ἀκακίᾳ τῆς καρδίας αὐτοῦ, καὶ ἐν τῇ συνέσει τῶν χειρῶν αὐτοῦ, ὡδήγησεν αὐτούς.

\begin{psalmheading}{\ch{78}{79} Ψαλμὸς τῷ Ἀσάφ.}
\end{psalmheading}
Ὁ Θεὸς, ἤλθοσαν ἔθνη εἰς τὴν κληρονομίαν σου, ἐμίαναν τὸν ναὸν τὸν ἅγιόν σου· ἔθεντο Ἱερουσαλὴμ εἰς ὀπωροφυλάκιον.
\vs{2}Ἔθεντο τὰ θνησιμαῖα τῶν δούλων σου βρώματα τοῖς πετεινοῖς τοῦ οὐρανοῦ, τὰς σάρκας τῶν ὁσίων σου τοῖς θηρίοις τῆς γῆς.
\vs{3}Ἐξέχεαν τὸ αἷμα αὐτῶν ὡς ὕδωρ, κύκλῳ Ἱερουσαλὴμ, καὶ οὐκ ἦν ὁ θάπτων.
\vs{4}Ἐγενήθημεν εἰς ὄνειδος τοῖς γείτοσιν ἡμῶν, μυκτηρισμὸς καὶ χλευασμὸς τοῖς κύκλῳ ἡμῶν.

\vs{5}Ἕως πότε, Κύριε, ὀργισθήσῃ εἰς τέλος; ἐκκαυθήσεται ὡς πῦρ ὁ ζῆλός σου;
\vs{6}Ἔκχεον τὴν ὀργήν σου ἐπὶ ἔθνη τὰ μὴ ἐπεγνωκότα σε, καὶ ἐπὶ βασιλείας αἳ τὸ ὄνομά σου οὐκ ἐπεκαλέσαντο.
\vs{7}Ὅτι κατέφαγον τὸν Ἰακὼβ, καὶ τὸν τόπον αὐτοῦ ἠρήμωσαν.

\vs{8}Μὴ μνησθῇς ἡμῶν ἀνομιῶν ἀρχαίων, ταχὺ προκαταλαβέτωσαν ἡμᾶς οἱ οἰκτιρμοί σου, ὅτι ἐπτωχεύσαμεν σφόδρα.
\vs{9}Βοήθησον ἡμῖν ὁ Θεὸς ὁ σωτὴρ ἡμῶν, ἕνεκα τῆς δόξης τοῦ ὀνόματός σου Κύριε ῥῦσαι ἡμᾶς, καὶ ἱλάσθητι ταῖς ἁμαρτίαις ἡμῶν ἕνεκα τοῦ ὀνόματός σου·
\vs{10}Μή ποτε εἴπωσιν ἐν τοῖς ἔθνεσι, ποῦ ἐστιν ὁ Θεὸς αὐτῶν; καὶ γνωσθήτω ἐν τοῖς ἔθνεσιν ἐνώπιον τῶν ὀφθαλμῶν ἡμῶν ἡ ἐκδίκησις τοῦ αἵματος τῶν δούλων σου τοῦ ἐκκεχυμένου.

\vs{11}Εἰσελθέτω ἐνώπιόν σου ὁ στεναγμὸς τῶν πεπεδημένων, κατὰ τὴν μεγαλωσύνην τοῦ βραχίονός σου περιποίησαι τοὺς υἱοὺς τῶν τεθανατωμένων.
\vs{12}Ἀπόδος τοῖς γείτοσιν ἡμῶν ἑπταπλάσια εἰς τὸν κόλπον αὐτῶν τὸν ὀνειδισμὸν αὐτῶν, ὃν ὠνείδισάν σε Κύριε.
\vs{13}Ἡμεῖς γὰρ λαός σου καὶ πρόβατα νομῆς σου, ἀνθομολογησόμεθά σοι εἰς τὸν αἰῶνα, εἰς γενεὰν καὶ γενεὰν ἐξαγγελοῦμεν τὴν αἴνεσίν σου.

\begin{psalmheading}{\ch{79}{80} Εἰς τὸ τέλος, ὑπὲρ τῶν ἀλλοιωθησομένων, μαρτύριον τῷ Ἀσὰφ, ψαλμὸς ὑπὲρ τοῦ Ἀσσυρίου.}
\end{psalmheading}
\vs{2}Ὁ Ποιμαίνων τὸν Ἰσραὴλ πρόσχες, ὁ ὁδηγῶν ὡσεὶ πρόβατα τὸν Ἰωσήφ· ὁ καθήμενος ἐπὶ τῶν χερουβὶμ ἐμφάνηθι,
\vs{3}ἐναντίον Ἐφραὶμ καὶ Βενιαμὶν καὶ Μανασσῆ· ἐξέγειρον τὴν δυναστείαν σου καὶ ἐλθὲ εἰς τὸ σῶσαι ἡμᾶς.
\vs{4}Ὁ Θεὸς ἐπίστρεψον ἡμᾶς, καὶ ἐπίφανον τὸ πρόσωπόν σου, καὶ σωθησόμεθα.

\vs{5}Κύριε ὁ Θεὸς τῶν δυνάμεων, ἕως πότε ὀργίζῃ ἐπὶ τὴν προσευχὴν τοῦ δούλου σου;
\vs{6}Ψωμιεῖς ἡμᾶς ἄρτον δακρύων, καὶ ποτιεῖς ἡμᾶς ἐν δάκρυσιν ἐν μέτρῳ.
\vs{7}Ἔθου ἡμᾶς εἰς ἀντιλογίαν τοῖς γείτοσιν ἡμῶν, καὶ οἱ ἐχθροὶ ἡμῶν ἐμυκτήρισαν ἡμᾶς.
\vs{8}Κύριε ὁ Θεὸς τῶν δυνάμεων ἐπίστρεψον ἡμᾶς, καὶ ἐπίφανον τὸ πρόσωπόν σου, καὶ σωθησόμεθα· διάψαλμα.

\vs{9}Ἄμπελον ἐξ Αἰγύπτου μετῇρας, ἐξέβαλες ἔθνη καὶ κατεφύτευσας αὐτήν.
\vs{10}Ὡδοποίησας ἔμπροσθεν αὐτῆς, καὶ κατεφύτευσας τὰς ῥίζας αὐτῆς, καὶ ἐπλήσθη ἡ γῆ.
\vs{11}Ἐκάλυψεν ὄρη ἡ σκιὰ αὐτῆς, καὶ αἱ ἀναδενδράδες αὐτῆς τὰς κέδρους τοῦ Θεοῦ.
\vs{12}Ἐξέτεινε τὰ κλήματα αὐτῆς ἕως θαλάσσης, καὶ ἕως ποταμοῦ τὰς παραφυάδας αὐτῆς.
\vs{13}Ἱνατί καθεῖλες τὸν φραγμὸν αὐτῆς, καὶ τρυγῶσιν αὐτὴν πάντες οἱ παραπορευόμενοι τὴν ὁδόν;
\vs{14}Ἐλυμῄνατο αὐτὴν σῦς ἐκ δρυμοῦ, καὶ μονιὸς ἄγριος κατενεμήσατο αὐτήν.

\vs{15}Ὁ Θεὸς τῶν δυνάμεων ἐπίστρεψον δὴ, ἐπίβλεψον ἐξ οὐρανοῦ καὶ ἴδε, καὶ ἐπίσκεψαι τὴν ἄμπελον ταύτην·
\vs{16}Καὶ κατάρτισαι αὐτὴν, ἣν ἐφύτευσεν ἡ δεξιά σου, καὶ ἐπὶ υἱὸν ἀνθρώπου ὃν ἐκραταίωσας σεαυτῷ.
\vs{17}Ἐμπεπυρισμένη πυρὶ καὶ ἀνεσκαμμένη· ἀπὸ ἐπιτιμήσεως τοῦ προσώπου σου ἀπολοῦνται.
\vs{18}Γενηθήτω ἡ χείρ σου ἐπʼ ἄνδρα δεξιᾶς σου, καὶ ἐπὶ υἱὸν ἀνθρώπου, ὃν ἐκραταίωσας σεαυτῷ.

\vs{19}Καὶ οὐ μὴ ἀποστῶμεν ἀπὸ σοῦ, ζωώσεις ἡμᾶς, καὶ τὸ ὄνομά σου ἐπικαλεσόμεθα.
\vs{20}Κύριε ὁ Θεὸς τῶν δυνάμεων ἐπίστρεψον ἡμᾶς, καὶ ἐπίφανον τὸ πρόσωπόν σου, καὶ σωθησόμεθα.

\begin{psalmheading}{\ch{80}{81} Εἰς τὸ τέλος, ὑπὲρ τῶν ληνῶν ψαλμὸς τῷ Ἀσάφ.}
\end{psalmheading}
\vs{2}Ἀγαλλιᾶσθε τῷ Θεῷ τῷ βοηθῷ ἡμῶν, ἀλαλάξατε τῷ Θεῷ Ἰακώβ.
\vs{3}Λάβετε ψαλμὸν καὶ δότε τύμπανον, ψαλτήριον τερπνὸν μετὰ κιθάρας.
\vs{4}Σαλπίσατε ἐν νεομηνίᾳ σάλπιγγι, ἐν εὐσήμῳ ἡμέρᾳ ἑορτῆς ὑμῶν.

\vs{5}Ὅτι πρόσταγμα τῷ Ἰσραήλ ἐστι, καὶ κρίμα τῷ Θεῷ Ἰακώβ.
\vs{6}Μαρτύριον ἐν τῷ Ἰωσὴφ ἔθετο αὐτὸν, ἐν τῷ ἐξελθεῖν αὐτὸν ἐκ γῆς Αἰγύπτου· γλῶσσαν ἣν οὐκ ἔγνω, ἤκουσεν.

\vs{7}Ἀπέστησεν ἀπὸ ἄρσεων τὸν νῶτον αὐτοῦ· αἱ χεῖρες αὐτοῦ ἐν τῷ κοφίνῳ ἐδούλευσαν.
\vs{8}Ἐν θλίψει ἐπεκαλέσω με καὶ ἐῤῥυσάμην σε· ἐπήκουσά σου ἐν ἀποκρύφῳ καταιγίδος, ἐδοκίμασά σε ἐπὶ ὕδατος ἀντιλογίας· διάψαλμα.
\vs{9}Ἄκουσον λαός μου καὶ λαλήσω σοι, Ἰσραὴλ, καὶ διαμαρτύρομαί σοι· ἐὰν ἀκούσῃς μου,
\vs{10}οὐκ ἔσται ἐν σοὶ θεὸς πρόσφατος, οὐδὲ προσκυνήσεις θεῷ ἀλλοτρίῳ.
\vs{11}Ἐγὼ γάρ εἰμι Κύριος ὁ Θεός σου, ὁ ἀναγαγών σε ἐκ γῆς Αἰγύπτου, πλάτυνον τὸ στόμα σου καὶ πληρώσω αὐτό.
\vs{12}Καὶ οὐκ ἤκουσεν ὁ λαός μου τῆς φωνῆς μου, καὶ Ἰσραὴλ οὐ προσέσχε μοι.
\vs{13}Καὶ ἐξαπέστειλα αὐτοὺς κατὰ τὰ ἐπιτηδεύματα τῶν καρδιῶν αὐτῶν, πορεύσονται ἐν τοῖς ἐπιτηδεύμασιν αὐτῶν.

\vs{14}Εἰ ὁ λαός μου ἤκουσέ μου, Ἰσραὴλ ταῖς ὁδοῖς μου εἰ ἐπορεύθη,
\vs{15}ἐν τῷ μηδενὶ ἂν τοὺς ἐχθροὺς αὐτῶν ἐταπείνωσα, καὶ ἐπὶ τοὺς θλίβοντας αὐτοὺς ἐπέβαλον ἂν τὴν χεῖρά μου.
\vs{16}Οἱ ἐχθροὶ Κυρίου ἐψεύσαντο αὐτῷ, καὶ ἔσται ὁ καιρὸς αὐτῶν εἰς τὸν αἰῶνα,
\vs{17}καὶ ἐψώμισεν αὐτοὺς ἐκ στέατος πυροῦ, καὶ ἐκ πέτρας μέλι ἐχόρτασεν αὐτούς.

\begin{psalmheading}{\ch{81}{82} Ψαλμὸς τῷ Ἀσάφ.}
\end{psalmheading}
\vs{2}Ὁ Θεὸς ἔστη ἐν συναγωγῇ θεῶν, ἐν μέσῳ δὲ θεοὺς διακρινεῖ.
\vs{2}Ἕως πότε κρίνετε ἀδικίαν, καὶ πρόσωπα ἁμαρτωλῶν λαμβάνετε; διάψαλμα.
\vs{3}Κρίνατε ὀρφανὸν καὶ πτωχὸν, ταπεινὸν καὶ πένητα δικαιώσατε.
\vs{4}Εξέλεσθε πένητα, καὶ πτωχὸν ἐκ χειρὸς ἁμαρτωλοῦ ῥύσασθε.

\vs{5}Οὐκ ἔγνωσαν οὐδὲ συνῆκαν, ἐν σκότει διαπορεύονται· σαλευθήσονται πάντα τὰ θεμελια τῆς γῆς.
\vs{6}Ἐγὼ εἶπα, θεοί ἐστε, καὶ υἱοὶ ὑψίστου πάντες.
\vs{7}Ὑμεῖς δὲ ὡς ἄνθρωποι ἀποθνήσκετε, καὶ ὡς εἷς τῶν ἀρχόντων πίπτετε.

\vs{8}Ἀνάστα ὁ Θεὸς, κρῖνον τὴν γῆν, ὅτι σὺ κατακληρονομήσεις ἐν πᾶσι τοῖς ἔθνεσιν.

\begin{psalmheading}{\ch{82}{83} Ὠδὴ ψαλμοῦ τῷ Ἀσάφ.}
\end{psalmheading}
\vs{2}Ὁ Θεὸς, τίς ὁμοιωθήσεταί σοι; μὴ σιγήσῃς, μηδὲ καταπραΰνῃς ὁ Θεός.

\vs{3}Ὅτι ἰδοὺ οἱ ἐχθροί σου ἤχησαν· καὶ οἱ μισοῦντές σε ᾖραν κεφαλήν.
\vs{4}Ἐπὶ τὸν λαόν σου κατεπανουργεύσαντο γνώμην, καὶ ἐβουλεύσαντο κατὰ τῶν ἁγίων σου.
\vs{5}Εἶπαν, δεῦτε καὶ ἐξολοθρεύσωμεν αὐτοὺς ἐξ ἔθνους, καὶ οὐ μὴ μνησθῇ τὸ ὄνομα Ἰσραὴλ ἔτι.
\vs{6}Ὅτι ἐβουλεύσαντο ἐν ὁμονοίᾳ ἐπιτοαυτὸ, κατὰ σοῦ διαθήκην διέθεντο·
\vs{7}Τὰ σκηνώματα τῶν Ἰδουμαίων καὶ οἱ Ἰσμαηλῖται, Μωὰβ καὶ οἱ Ἀγαρηνοὶ,
\vs{8}Γεβὰλ καὶ Ἀμμὼν καὶ Ἀμαλὴκ, καὶ ἀλλόφυλοι μετὰ τῶν κατοικούντων Τύρον.
\vs{9}Καὶ γὰρ καὶ Ἀσσοὺρ συμπαρεγένετο μετʼ αὐτῶν, ἐγενήθησαν εἰς ἀντίληψιν τοῖς υἱοῖς Λώτ· διάψαλμα.

\vs{10}Ποίησον αὐτοῖς ὡς τῇ Μαδιὰμ καὶ τῷ Σεισάρᾳ, ὡς τῷ Ἰαβεὶν ἐν τῷ χειμάῤῥῳ Κεισῶν.
\vs{11}Ἐξωλοθρεύθησαν ἐν Ἀενδὼρ, ἐγενήθησαν ὡσεὶ κόπρος τῇ γῇ.
\vs{12}Θοῦ τοὺς ἄρχοντας αὐτῶν ὡς τὸν Ὠρὴβ καὶ Ζὴβ καὶ Ζεβεὲ καὶ Σαλμανὰ, πάντας τοὺς ἄρχοντας αὐτῶν·
\vs{13}Οἵτινες εἶπαν, κληρονομήσωμεν ἑαυτοῖς τὸ θυσιαστήριον τοῦ Θεοῦ.
\vs{14}Ὁ Θεός μου θοῦ αὐτοὺς ὡς τροχὸν, ὡς καλάμην κατὰ πρόσωπον ἀνέμου.
\vs{15}Ὡσεὶ πῦρ ὃ διαφλέξει δρυμὸν, ὡσεὶ φλὸξ κατακαύσαι ὄρη·
\vs{16}Οὕτως καταδιώξεις αὐτοὺς ἐν τῇ καταιγίδι σου, καὶ ἐν τῇ ὀργῇ σου ταράξεις αὐτούς.
\vs{17}Πλήρωσον τὰ πρόσωπα αὐτῶν ἀτιμίας, καὶ ζητήσουσι τὸ ὄνομά σου Κύριε.
\vs{18}Αἰσχυνθήτωσαν καὶ ταραχθήτωσαν εἰς τὸν αἰῶνα τοῦ αἰῶνος, καὶ ἐντραπήτωσαν καὶ ἀπολέσθωσαν.
\vs{19}Καὶ γνώτωσαν ὅτι ὄνομά σοι Κύριος· σὺ μόνος ὕψιστος ἐπὶ πᾶσαν τὴν γῆν.

\begin{psalmheading}{\ch{83}{84} Εἰς τὸ τέλος, ὑπὲρ τῶν ληνῶν τοῖς υἱοῖς Κορὲ ψαλμός.}
\end{psalmheading}
\vs{2}Ὡς ἀγαπητὰ τὰ σκηνώματά σου Κύριε τῶν δυνάμεων.
\vs{3}Ἐπιποθεῖ καὶ ἐκλείπει ἡ ψυχή μου εἰς τὰς αὐλὰς τοῦ Κυρίου· ἡ καρδία μου καὶ ἡ σάρξ μου ἠγαλλιάσαντο ἐπὶ Θεὸν ζῶντα·
\vs{4}Καὶ γὰρ στρουθίον εὗρεν ἑαυτῷ οἰκίαν, καὶ τρυγὼν νοσσιὰν ἑαυτῇ, οὗ θήσει τὰ νοσσία ἑαυτῆς· τὰ θυσιαστήριά σου Κύριε τῶν δυνάμεων, ὁ βασιλεύς μου καὶ ὁ Θεός μου.

\vs{5}Μακάριοι οἱ κατοικοῦντες ἐν τῷ οἴκῳ σου, εἰς τοὺς αἰῶνας τῶν αἰώνων αἰνέσουσί σε· διάψαλμα.
\vs{6}Μακάριος ἀνὴρ οὗ ἐστιν ἡ ἀντίληψις αὐτοῦ παρὰ σοῦ, Κύριε· ἀναβάσεις ἐν τῇ καρδίᾳ αὐτοῦ διέθετο,
\vs{7}εἰς τὴν κοιλάδα τοῦ κλαυθμῶνος, εἰς τὸν τόπον ὃν ἔθετο· καὶ γὰρ εὐλογίας δώσει ὁ νομοθετῶν,
\vs{8}πορεύσονται ἐκ δυνάμεως εἰς δύναμιν, ὀφθήσεται ὁ Θεὸς τῶν θεῶν ἐν Σιών.

\vs{9}Κύριε ὁ Θεὸς τῶν δυνάμεων, εἰσάκουσον τῆς προσευχῆς μου, ἐνώτισαι ὁ Θεὸς Ἰακώβ· διάψαλμα.
\vs{10}Ὑπερασπιστὰ ἡμῶν ἴδε ὁ Θεὸς, καὶ ἐπίβλεψον ἐπὶ τὸ πρόσωπον τοῦ χριστοῦ σοῦ.
\vs{11}Ὅτι κρείσσων ἡμέρα μία ἐν ταῖς αὐλαῖς σου, ὑπὲρ χιλιάδας· ἐξελεξάμην παραῤῥιπτεῖσθαι ἐν τῷ οἴκῳ τοῦ Θεοῦ μᾶλλον ἢ οἰκεῖν με ἐπὶ σκηνώμασιν ἁμαρτωλῶν.
\vs{12}Ὅτι ἔλεον καὶ ἀλήθειαν ἀγαπᾷ Κύριος, ὁ Θεὸς χάριν καὶ δόξαν δώσει· Κύριος οὐχ ὑστερήσει τὰ ἀγαθὰ τοῖς πορευομένοις ἐν ἀκακίᾳ.
\vs{13}Κύριε τῶν δυνάμεων, μακάριος ἄνθρωπος ὁ ἐλπίζων ἐπὶ σέ.

\begin{psalmheading}{\ch{84}{85} Εἰς τὸ τέλος, τοῖς υἱοῖς Κορὲ ψαλμός.}
\end{psalmheading}
\vs{2}Εὐδόκησας Κύριε τὴν γῆν σου, ἀπέστρεψας τὴν αἰχμαλωσίαν Ἰακώβ.
\vs{3}Ἀφῆκας τὰς ἀνομίας τῷ λαῷ σου, ἐκάλυψας πάσας τὰς ἁμαρτίας αὐτῶν· διάψαλμα.
\vs{4}Κατέπαυσας πᾶσαν τὴν ὀργήν σου, ἀπέστρεψας ἀπὸ ὀργῆς θυμοῦ σου.

\vs{5}Ἐπίστρεψον ἡμᾶς ὁ Θεὸς τῶν σωτηρίων ἡμῶν, καὶ ἀπόστρεψον τὸν θυμόν σου ἀφʼ ἡμῶν.
\vs{6}Μὴ εἰς τὸν αἰῶνα ὀργισθῇς ἡμῖν; ἢ διατενεῖς τὴν ὀργήν σου ἀπὸ γενεᾶς εἰς γενεάν;
\vs{7}Ὁ Θεὸς, σὺ ἐπιστρέψας ζωώσεις ἡμᾶς, καὶ ὁ λαός σου εὐφρανθήσεται ἐπὶ σοί.
\vs{8}Δεῖξον ἡμῖν Κύριε τὸ ἔλεός σου, καὶ τὸ σωτήριόν σου δῴης ἡμῖν.

\vs{9}Ἀκούσομαι τί λαλήσει ἐν ἐμοὶ Κύριος ὁ Θεὸς, ὅτι λαλήσει εἰρήνην ἐπὶ τὸν λαὸν αὐτοῦ, καὶ ἐπὶ τοὺς ὁσίους αὐτοῦ, καὶ ἐπὶ τοὺς ἐπιστρέφοντας πρὸς αὐτὸν καρδίαν.
\vs{10}Πλὴν ἐγγὺς τῶν φοβουμένων αὐτὸν τὸ σωτήριον αὐτοῦ, τοῦ κατασκηνῶσαι δόξαν ἐν τῇ γῇ ἡμῶν.
\vs{11}Ἔλεος καὶ ἀλήθεια συνήντησαν, δικαιοσύνη καὶ εἰρήνη κατεφίλησαν.
\vs{12}Ἀλήθεια ἐκ τῆς γῆς ἀνέτειλε, καὶ δικαιοσύνη ἐκ τοῦ οὐρανοῦ διέκυψε.
\vs{13}Καὶ γὰρ ὁ Κύριος δώσει χρηστότητα, καὶ ἡ γῆ ἡμῶν δώσει τὸν καρπὸν αὐτῆς.
\vs{14}Δικαιοσύνη ἐναντίον αὐτοῦ προπορεύσεται, καὶ θήσει εἰς ὁδὸν τὰ διαβήματα αὐτοῦ.

\begin{psalmheading}{\ch{85}{86} Προσευχὴ τῷ Δαυίδ.}
\end{psalmheading}
\vs{2}Κλίνον Κύριε τὸ οὖς σου, καὶ εἰσάκουσόν μου, ὅτι πτωχὸς καὶ πένης εἰμὶ ἐγώ.
\vs{2}Φύλαξον τὴν ψυχήν μου, ὅτι ὅσιός εἰμι· σῶσον τὸν δοῦλόν σου ὁ Θεὸς, τὸν ἐλπίζοντα ἐπὶ σέ.
\vs{3}Ἐλέησόν με Κύριε, ὅτι πρὸς σὲ κεκράξομαι ὅλην τὴν ἡμέραν.
\vs{4}Εὔφρανον τὴν ψυχὴν τοῦ δούλου σου, ὅτι πρὸς σὲ Κύριε ᾖρα τὴν ψυχήν μου.
\vs{5}Ὅτι σὺ Κύριε χρηστὸς καὶ ἐπιεικὴς, καὶ πολυέλεος πᾶσι τοῖς ἐπικαλουμένοις σε.
\vs{6}Ἐνώτισαι Κύριε τὴν προσευχήν μου, καὶ πρόσχες τῇ φωνῇ τῆς δεήσεώς μου.
\vs{7}Ἐν ἡμέρᾳ θλίψεώς μου ἐκέκραξα πρὸς σὲ, ὅτι εἰσήκουσάς μου.

\vs{8}Οὐκ ἔστιν ὅμοιός σοι ἐν θεοῖς Κύριε, καὶ οὐκ ἔστι κατὰ τὰ ἔργα σου.
\vs{9}Πάντα τὰ ἔθνη ὅσα ἐποίησας ἥξουσι, καὶ προσκυνήσουσιν ἐνώπιόν σου Κύριε, καὶ δοξάσουσι τὸ ὄνομά σου,
\vs{10}ὅτι μέγας εἶ σὺ, καὶ ποιῶν θαυμάσια, σὺ εἶ ὁ Θεὸς μόνος ὁ μέγας.
\vs{11}Ὁδήγησόν με Κύριε ἐν τῇ ὁδῷ σου, καὶ πορεύσομαι ἐν τῇ ἀληθείᾳ σου· εὐφρανθήτω ἡ καρδία μου, τοῦ φοβεῖσθαι τὸ ὄνομά σου.
\vs{12}Ἐξομολογήσομαί σοι Κύριε ὁ Θεός μου ἐν ὅλῃ καρδίᾳ μου, καὶ δοξάσω τὸ ὄνομά σου εἰς τὸν αἰῶνα.
\vs{13}Ὅτι τὸ ἔλεός σου μέγα ἐπʼ ἐμὲ, καὶ ἐῤῥύσω τὴν ψυχήν μου ἐξ ᾅδου κατωτάτου.

\vs{14}Ὁ Θεὸς, παράνομοι ἐπανέστησαν ἐπʼ ἐμὲ, καὶ συναγωγὴ κραταιῶν ἐζήτησαν τὴν ψυχήν μου, καὶ οὐ προέθεντό σε ἐνώπιον αὐτῶν.
\vs{15}Καὶ σὺ Κύριε ὁ Θεὸς οἰκτίρμων καὶ ἐλεήμων, μακρόθυμος καὶ πολυέλεος καὶ ἀληθινός.
\vs{16}Ἐπίβλεψον ἐπʼ ἐμὲ, καὶ ἐλέησόν με, δὸς τὸ κράτος σου τῷ παιδί σου, καὶ σῶσον τὸν υἱὸν τῆς παιδίσκης σου.
\vs{17}Ποίησον μετʼ ἐμοῦ σημεῖον εἰς ἀγαθὸν, καὶ ἰδέτωσαν οἱ μισοῦντές με, καὶ αἰσχυνθήτωσαν· ὅτι σὺ Κύριε, ἐβοήθησάς μοι, καὶ παρεκάλεσάς με.

\begin{psalmheading}{\ch{86}{87} Τοῖς υἱοῖς Κορὲ ψαλμὸς ᾠδῆς.}
\end{psalmheading}
Οἱ θεμέλιοι αὐτοῦ ἐν τοῖς ὄρεσι τοῖς ἁγίοις.
\vs{2}Ἀγαπᾷ Κύριος τὰς πύλας Σιὼν, ὑπὲρ πάντα τὰ σκηνώματα Ἰακώβ.
\vs{3}Δεδοξασμένα ἐλαλήθη περὶ σοῦ ἡ πόλις τοῦ Θεοῦ· διάψαλμα.

\vs{4}Μνησθήσομαι Ῥαὰβ, καὶ Βαβυλῶνος, τοῖς γινώσκουσί με· καὶ ἰδοὺ ἀλλόφυλοι καὶ Τύρος καὶ λαὸς Αἰθιόπων, οὗτοι ἐγεννήθησαν ἐκεῖ.
\vs{5}Μήτηρ Σιὼν ἐρεῖ ἄνθρωπος, καὶ ἄνθρωπος ἐγενήθη ἐν αὐτῇ, καὶ αὐτὸς ἐθεμελίωσεν αὐτὴν ὁ ὕψιστος.
\vs{6}Κύριος διηγήσεται ἐν γραφῇ λαῶν, καὶ ἀρχόντων τούτων τῶν γεγενημένων ἐν αὐτῇ. διάψαλμα.
\vs{7}Ὡς εὐφραινομένων πάντων ἡ κατοικία ἐν σοί.

\begin{psalmheading}{\ch{87}{88} Ὠδὴ ψαλμοῦ τοῖς υἱοῖς Κορὲ, εἰς τὸ τέλος, ὑπὲρ μαελὲθ τοῦ ἀποκριθῆναι, συνέσεως Αἰμὰν τῷ Ἰσραηλίτῃ.}
\end{psalmheading}
\vs{2}Κύριε ὁ Θεὸς τῆς σωτηρίας μου, ἡμέρας ἐκέκραξα καὶ ἐν νυκτὶ ἐναντίον σου.
\vs{3}Εἰσελθέτω ἐνώπιόν σου ἡ προσευχή μου, κλῖνον τὸ οὖς σου εἰς τὴν δέησίν μου, Κύριε.

\vs{4}Ὅτι ἐπλήσθη κακῶν ἡ ψυχή μου, καὶ ἡ ζωή μου τῷ ᾅδῃ ἤγγισε.
\vs{5}Προσελογίσθην μετὰ τῶν καταβαινόντων εἰς λάκκον, ἐγενήθην ὡς ἄνθρωπος ἀβοήθητος,
\vs{6}ἐν νεκροῖς ἐλεύθερος, ὡσεὶ τραυματίαι ἐῤῥιμμένοι καθεύδοντες ἐν τάφῳ, ὧν οὐκ ἐμνήσθης ἔτι, καὶ αὐτοὶ ἐκ τῆς χειρός σου ἀπώσθησαν.
\vs{7}Ἔθεντό με ἐν λάκκῳ κατωτάτῳ, ἐν σκοτεινοῖς καὶ ἐν σκιᾷ θανάτου.
\vs{8}Ἐπʼ ἐμὲ ἐπεστηρίχθη ὁ θυμός σου, καὶ πάντας τοὺς μετεωρισμούς σου ἐπήγαγες ἐπʼ ἐμέ· διάψαλμα.
\vs{9}Ἐμάκρυνας τοὺς γνωστούς μου ἀπʼ ἐμοῦ, ἔθεντό με βδέλυγμα ἑαυτοῖς· παρεδόθην καὶ οὐκ ἐξεπορευόμην.
\vs{10}Οἱ ὀφθαλμοί μου ἠσθένησαυ ἀπὸ πτωχείας· καὶ ἐκέκραξα πρὸς σὲ Κύριε ὅλην τὴν ἡμέραν, διεπέτασα πρὸς σὲ τὰς χεῖράς μου.

\vs{11}Μὴ τοῖς νεκροῖς ποιήσεις θαυμάσια, ἢ ἰατροὶ ἀναστήσουσι καὶ ἐξομολογήσονταί σοι;
\vs{12}Μὴ διηγήσεταί τις ἐν τάφῳ τὸ ἔλεός σου, καὶ τὴν ἀλήθειάν σου ἐν τῇ ἀπωλείᾳ;
\vs{13}Μὴ γνωσθήσεται ἐν τῷ σκότει τὰ θαυμάσιά σου, καὶ ἡ δικαιοσύνη σου ἐν γῇ ἐπιλελησμένῃ;
\vs{14}Κᾀγὼ πρὸς σὲ Κύριε ἐκέκραξα, καὶ τοπρωῒ ἡ προσευχή μου προφθάσει σε.

\vs{15}Ἱνατί Κύριε ἀπωθεῖς τὴν προσευχήν μου, ἀποστρέφεις τὸ πρόσωπόν σου ἀπʼ ἐμοῦ;
\vs{16}Πτωχός εἰμι ἐγὼ καὶ ἐν κόποις ἐκ νεότητός μου, ὑψωθεὶς δὲ ἐταπεινώθην καὶ ἐξηπορήθην.
\vs{17}Ἐπʼ ἐμὲ διῆλθον αἱ ὀργαί σου, καὶ οἱ φοβερισμοί σου ἐξετάραξάν με.
\vs{18}Ἐκύκλωσάν με ὡς ὕδωρ, ὅλην τὴν ἡμέραν περιέσχον με ἅμα.
\vs{19}Ἐμάκρυνας ἀπʼ ἐμοῦ φίλον, καὶ τοὺς γνωστούς μου ἀπὸ ταλαιπωρίας.

\begin{psalmheading}{\ch{88}{89} Συνέσεως Αἰθὰμ τῷ Ἰσραηλίτῃ.}
\end{psalmheading}
\vs{2}Τὰ ἐλέη σου Κύριε εἰς τὸν αἰῶνα ᾄσομαι, εἰς γενεὰν καὶ γενεὰν ἀπαγγελῶ τὴν ἀλήθειάν σου ἐν τῷ στόματί μου.
\vs{3}Ὅτι εἶπας, εἰς τὸν αἰῶνα ἔλεος οἰκοδομηθήσεται, ἐν τοῖς οὐρανοῖς ἑτοιμασθήσεται ἡ ἀλήθειά σου·
\vs{4}Διεθέμην διαθήκην τοῖς ἐκλεκτοῖς μου, ὤμοσα Δαυὶδ τῷ δούλῳ μου.
\vs{5}Ἕως τοῦ αἰῶνος ἑτοιμάσω τὸ σπέρμα σου, καὶ οἰκοδομήσω εἰς γενεὰν καὶ γενεὰν τὸν θρόνον σου· διάψαλμα.

\vs{6}Ἐξομολογήσονται οἱ οὐρανοὶ τὰ θαυμάσιά σου Κύριε, καὶ τὴν ἀλήθειάν σου ἐν ἐκκλησίᾳ ἁγίων.
\vs{7}Ὅτι τίς ἐν νεφέλαις ἰσωθήσεται τῷ Κυρίῳ; καὶ τίς ὁμοιωθήσεται τῷ Κυρίῳ ἐν υἱοῖς Θεοῦ;
\vs{8}Ὁ Θεὸς ἐνδοξαζόμενος ἐν βουλῇ ἁγίων, μέγας καὶ φοβερὸς ἐπὶ πάντας τοὺς περικύκλῳ αὐτοῦ.
\vs{9}Κύριε ὁ Θεὸς τῶν δυνάμεων, τίς ὅμοιός σοι; δυνατὸς εἶ Κύριε, καὶ ἡ ἀλήθειά σου κύκλῳ σου.
\vs{10}Σὺ δεσπόζεις τοῦ κράτους τῆς θαλάσσης, τὸν δὲ σάλον τῶν κυμάτων αὐτῆς σὺ καταπραΰνεις.
\vs{11}Σὺ ἐταπείνωσας ὡς τραυματίαν ὑπερήφανον, καὶ ἐν τῷ βραχίονι τῆς δυνάμεώς σου διεσκόρπισας τοὺς ἐχθρούς σου.
\vs{12}Σοί εἰσιν οἱ οὐρανοὶ, καὶ σή ἐστιν ἡ γῆ, τὴν οἰκουμένην καὶ τὸ πλήρωμα αὐτῆς σὺ ἐθεμελίωσας.
\vs{13}Τὸν Βοῤῥᾶν καὶ θάλασσας σὺ ἔκτισας, Θαβὼρ καὶ Ἑρμὼν ἐν τῷ ὀνόματί σου ἀγαλλιάσονται.
\vs{14}Σὸς ὁ βραχίων μετὰ δυναστείας· κραταιωθήτω ἡ χείρ σου, ὑψωθήτω ἡ δεξιά σου.
\vs{15}Δικαιοσύνη καὶ κρίμα ἑτοιμασία τοῦ θρόνου σου· ἔλεος καὶ ἀλήθεια προπορεύσονται πρὸ προσώπου σου.

\vs{16}Μακάριος ὁ λαὸς ὁ γινώσκων ἀλαλαγμόν· Κύριε ἐν τῷ φωτὶ τοῦ προσώπου σου πορεύσονται,
\vs{17}καὶ ἐν τῷ ὀνόματί σου ἀγαλλιάσονται ὅλην τὴν ἡμέραν, καὶ ἐν τῇ δικαιοσύνῃ σου ὑψωθήσονται.
\vs{18}Ὅτι τὸ καύχημα τῆς δυνάμεως αὐτῶν σὺ εἶ, καὶ ἐν τῇ εὐδοκίᾳ σου ὑψωθήσεται τὸ κέρας ἡμῶν·
\vs{19}ὅτι τοῦ Κυρίου ἡ ἀντίληψις, καὶ τοῦ ἁγίου Ἰσραὴλ βασιλέως ἡμῶν.

\vs{20}Τότε ἐλάλησας ἐν ὁράσει τοῖς υἱοῖς σου, καὶ εἶπας, ἐθέμην βοήθειαν ἐπὶ δυνατὸν, ὕψωσα ἐκλεκτὸν ἐκ τοῦ λαοῦ μου.
\vs{21}Εὗρον Δαυὶδ τὸν δοῦλόν μου, ἐν ἐλέει ἁγίῳ ἔχρισα αὐτόν.
\vs{22}Ἡ γὰρ χείρ μου συναντιλήψεται αὐτῷ, καὶ ὁ βραχίων μου κατισχύσει αὐτόν.
\vs{23}Οὐκ ὠφελήσει ἐχθρὸς ἐν αὐτῷ, καὶ υἱὸς ἀνομίας οὐ προσθήσει τοῦ κακῶσαι αὐτόν.
\vs{24}Καὶ συγκόψω ἀπὸ προσώπου αὐτοῦ τοὺς ἐχθροὺς αὐτοῦ, καὶ τοὺς μισοῦντας αὐτὸν τροπώσομαι.
\vs{25}Καὶ ἡ ἀλήθειά μου καὶ τὸ ἔλεός μου μετʼ αὐτοῦ, καὶ ἐν τῷ ὀνόματί μου ὑψωθήσεται τὸ κέρας αὐτοῦ.
\vs{26}Καὶ θήσομαι ἐν θαλάσσῃ χεῖρα αὐτοῦ, καὶ ἐν ποταμοῖς δεξιὰν αὐτοῦ.
\vs{27}Αὐτὸς ἐπικαλέσεταί με, πατήρ μου εἶ σὺ, Θεός μου καὶ ἀντιλήπτωρ τῆς σωτηρίας μου.
\vs{28}Κᾀγὼ πρωτότοκον θήσομαι αὐτὸν, ὑψηλὸν παρὰ τοῖς βασιλεῦσι τῆς γῆς.
\vs{29}Εἰς τὸν αἰῶνα φυλάξω αὐτῷ τὸ ἔλεός μου, καὶ ἡ διαθήκη μου πιστὴ αὐτῷ.
\vs{30}Καὶ θήσομαι εἰς τὸν αἰῶνα τοῦ αἰῶνος τὸ σπέρμα αὐτοῦ, καὶ τὸν θρόνον αὐτοῦ ὡς τὰς ἡμέρας τοῦ οὐρανοῦ.

\vs{31}Ἐὰν ἐγκαταλίπωσιν οἱ υἱοὶ αὐτοῦ τὸν νόμον μου, καὶ τοῖς κρίμασί μου μὴ πορευθῶσιν·
\vs{32}ἐὰν τὰ δικαιώματά μου βεβηλώσωσι, καὶ τὰς ἐντολάς μου μὴ φυλάξωσιν·
\vs{33}Ἐπισκέψομαι ἐν ῥάβδῳ τὰς ἀνομίας αὐτῶν, καὶ ἐν μάστιξι τὰς ἁμαρτίας αὐτῶν.
\vs{34}Τὸ δὲ ἔλεός μου οὐ μὴ διασκεδάσω ἀπʼ αὐτοῦ, οὐδὲ μὴ ἀδικήσω ἐν τῇ ἀληθείᾳ μου,
\vs{35}οὐδὲ μὴ βεβηλώσω τὴν διαθήκην μου, καὶ τὰ ἐκπορευόμενα διὰ τῶν χειλέων μου οὐ μὴ ἀθετήσω.
\vs{36}Ἅπαξ ὤμοσα ἐν τῷ ἁγίῳ μου, εἰ τῷ Δαυὶδ ψεύσομαι.
\vs{37}Τὸ σπέρμα αὐτοῦ εἰς τὸν αἰῶνα μενεῖ, καὶ ὁ θρόνος αὐτοῦ ὡς ὁ ἥλιος ἐναντίον μου,
\vs{38}καὶ ὡς ἡ σελήνη κατηρτισμένη εἰς τὸν αἰῶνα, καὶ ὁ μάρτυς ἐν οὐρανῷ πιστός· διάψαλμα.

\vs{39}Σὺ δὲ ἀπώσω καὶ ἐξουδένωσας, ἀνεβάλου τὸν χριστόν σου.
\vs{40}Κατέστρεψας τὴν διαθήκην τοῦ δούλου σου, ἐβεβήλωσας εἰς τὴν γῆν τὸ ἁγίασμα αὐτοῦ.
\vs{41}Καθεῖλες πάντας τοὺς φραγμοῦς αὐτοῦ, ἔθου τὰ ὀχυρώματα αὐτοῦ δειλίαν.
\vs{42}Διήρπασαν αὐτὸν πάντες οἱ διοδεύοντες ὁδὸν, ἐγενήθη ὄνειδος τοῖς γείτοσιν αὐτοῦ.
\vs{43}Ὕψωσας τὴν δεξιὰν τῶν ἐχθρῶν αὐτοῦ, εὔφρανας πάντας τοὺς ἐχθροὺς αὐτοῦ.
\vs{44}Ἀπέστρεψας τὴν βοήθειαν τῆς ῥομφαίας αὐτοῦ, καὶ οὐκ ἀντελάβου αὐτοῦ ἐν τῷ πολέμῳ.
\vs{45}Κατέλυσας ἀπὸ καθαρισμοῦ αὐτὸν, τὸν θρόνον αὐτοῦ εἰς τὴν γῆν κατέῤῥαξας,
\vs{46}ἐσμίκρυνας τὰς ἡμέρας τοῦ θρόνου αὐτοῦ, κατέχεας αὐτοῦ αἰσχύνην· διάψαλμα.

\vs{47}Ἕως πότε Κύριε ἀποστρέφῃ εἰς τέλος; ἐκκαυθήσεται ὡς πῦρ ἡ ὀργή σου;
\vs{48}Μνήσθητι τίς ἡ ὑπόστασίς μου· μὴ γὰρ ματαίως ἔκτισας πάντας τοὺς υἱοὺς τῶν ἀνθρώπων;
\vs{49}Τίς ἐστιν ἄνθρωπος, ὃς ζήσεται καὶ οὐκ ὄψεται θάνατον; ῥύσεται τὴν ψυχὴν αὐτοῦ ἐκ χειρὸς ᾅδου; διάψαλμα.
\vs{50}Ποῦ ἐστι τὰ ἐλέη σου τὰ ἀρχαῖα, Κύριε, ἃ ὤμοσας τῷ Δαυὶδ ἐν τῇ ἀληθείᾳ σου;
\vs{51}Μνήσθητι Κύριε τοῦ ὀνειδισμου τῶν δούλων σου οὗ ὑπέσχον ἐν τῷ κόλπῳ μου πολλῶν ἐθνῶν·
\vs{52}οὗ ὠνείδισαν οἱ ἐχθροί σου Κύριε, οὗ ὠνείδισαν τὸ ἀντάλλαγμα τοῦ χριστοῦ σου.
\vs{53}Εὐλογητὸς Κύριος εἰς τὸν αἰῶνα· γένοιτο, γένοιτο.

\begin{psalmheading}{\ch{89}{90} Προσευχὴ τοῦ Μωυσῆ ἀνθρώπου τοῦ Θεοῦ.}
\end{psalmheading}
Κύριε, καταφυγὴ ἐγενήθης ἡμῖν ἐν γενεᾷ καὶ γενεᾷ.
\vs{2}Πρὸ τοῦ ὄρη γενηθῆναι καὶ πλασθῆναι τὴν γῆν καὶ τὴν οἰκουμένην, καὶ ἀπὸ τοῦ αἰῶνος ἕως τοῦ αἰῶνος σὺ εἶ.
\vs{3}Μὴ ἀποστρέψῃς ἄνθρωπον εἰς ταπείνωσιν, καὶ εἶπας, ἐπιστρέψατε υἱοὶ ἀνθρώπων;
\vs{4}Ὅτι χίλια ἔτη ἐν ὀφθαλμοῖς σου, ὡς ἡ ἡμέρα ἡ ἐχθὲς ἥτις διῆλθε, καὶ φυλακὴ ἐν νυκτί.
\vs{5}Τὰ ἐξουδενώματα αὐτῶν ἔτη ἔσονται, τοπρωῒ ὡσεὶ χλόη παρέλθοι·
\vs{6}Τοπρωῒ ἀνθήσαι καὶ παρέλθοι, τὸ ἑσπέρας ἀποπέσοι, σκληρυνθείη καὶ ξηρανθείη.
\vs{7}Ὅτι ἐξελίπομεν ἐν τῇ ὀργῇ σου, καὶ ἐν τῷ θυμῷ σου ἐταράχθημεν.
\vs{8}Ἔθου τὰς ἀνομίας ἡμῶν ἐνώπιόν σου, ὁ αἰὼν ἡμῶν εἰς φωτισμὸν τοῦ προσώπου σου.
\vs{9}Ὅτι πᾶσαι αἱ ἡμέραι ἡμῶν ἐξέλιπον, καὶ ἐν τῇ ὀργῇ σου ἐξελίπομεν· τὰ ἔτη ἡμῶν ὡς ἀράχνη ἐμελέτων.
\vs{10}Αἱ ἡμέραι τῶν ἐτῶν ἡμῶν ἐν αὐτοῖς ἑβδομήκοντα ἔτη, ἐὰν δὲ ἐν δυναστείαις, ὀγδοήκοντα ἔτη, καὶ τὸ πλεῖον αὐτῶν κόπος καὶ πόνος· ὅτι ἐπῆλθε πραΰτης ἐφʼ ἡμᾶς, καὶ παιδευθησόμεθα.
\vs{11}Τίς γινώσκει τὸ κράτος τῆς ὀργῆς σου, καὶ ἀπὸ τοῦ φόβου τοῦ θυμοῦ σου ἐξαριθμήσασθαι;
\vs{12}τὴν δεξιάν σου οὕτως γνώρισον, καὶ τοὺς πεπαιδευμένους τῇ καρδίᾳ ἐν σοφίᾳ.

\vs{13}Ἐπίστρεψον, Κύριε· ἕως πότε; καὶ παρακλήθητι ἐπὶ τοῖς δούλοις σου.
\vs{14}Ἐνεπλήσθημεν τοπρωῒ τοῦ ἐλέους σου, καὶ ἠγαλλιασάμεθα καὶ εὐφράνθημεν· ἐν πάσαις ταῖς ἡμέραις ἡμῶν
\vs{15}εὐφρανθείημεν, ἀνθʼ ᾧν ἡμερῶν ἐταπείνωσας ἡμᾶς, ἐτῶν ὧν εἴδομεν κακά.
\vs{16}Καὶ ἴδε ἐπὶ τοὺς δούλους σου καὶ ἐπὶ τὰ ἔργα σου, καὶ ὁδήγησον τοὺς υἱοὺς αὐτῶν.
\vs{17}Καὶ ἔστω ἡ λαμπρότης Κυρίου τοῦ Θεοῦ ἡμῶν ἐφʼ ἡμᾶς, καὶ τὰ ἔργα τῶν χειρῶν ἡμῶν κατεύθυνον ἐφʼ ἡμᾶς.

\begin{psalmheading}{\ch{90}{91} Αἶνος ᾠδῆς τῷ Δαυίδ.}
\end{psalmheading}
Ὁ κατοικῶν ἐν βοηθείᾳ τοῦ ὑψίστου, ἐν σκέπῃ τοῦ Θεοῦ τοῦ οὐρανοῦ αὐλισθήσεται.
\vs{2}Ἐρεῖ τῷ Κυρίῳ, ἀντιλήπτωρ μου εἶ καὶ καταφυγή μου, ὁ Θεός μου, ἐλπιῶ ἐπʼ αὐτόν.
\vs{3}Ὅτι αὐτὸς ῥύσεταί σε ἐκ παγίδος θηρευτῶν, ἀπὸ λόγου ταραχώδους.
\vs{4}Ἐν τοῖς μεταφρένοις αὐτοῦ ἐπισκιάσει σοι, καὶ ὑπὸ τὰς πτέρυγας αὐτοῦ ἐλπιεῖς· ὅπλῳ κυκλώσει σε ἡ ἀλήθεια αὐτοῦ.
\vs{5}Οὐ φοβηθήσῃ ἀπὸ φόβου νυκτερινοῦ, ἀπὸ βέλους πετομένου ἡμέρας,
\vs{6}ἀπὸ πράγματος διαπορευομένου ἐν σκότει, ἀπὸ συμπτώματος καὶ δαιμονίου μεσημβρινοῦ.
\vs{7}Πεσεῖται ἐκ τοῦ κλίτους σου χιλιὰς, καὶ μυριὰς ἐκ δεξιῶν σου, πρὸς σὲ δὲ οὐκ ἐγγιεῖ.
\vs{8}Πλὴν τοῖς ὀφθαλμοῖς σου κατανοήσεις, καὶ ἀνταπόδοσιν ἁμαρτωλῶν ὄψει.

\vs{9}Ὅτι σὺ Κύριε ἡ ἐλπίς μου, τὸν ὕψιστον ἔθου καταφυγήν σου.
\vs{10}Οὐ προσελεύσεται πρὸς σὲ κακὰ, καὶ μάστιξ οὐκ ἐγγιεῖ τῷ σκηνώματί σου.
\vs{11}Ὅτι τοῖς ἀγγέλοις αὐτοῦ ἐντελεῖται περὶ σοῦ, τοῦ διαφυλάξαι σε ἐν πάσαις ταῖς ὁδοῖς σου.
\vs{12}Ἐπὶ χειρῶν ἀροῦσί σε, μή ποτε προσκόψῃς πρὸς λίθον τὸν πόδα σου.
\vs{13}Ἐπʼ ἀσπίδα καὶ βασιλίσκον ἐπιβήσῃ, καὶ καταπατήσεις λέοντα καὶ δράκοντα.

\vs{14}Ὅτι ἐπʼ ἐμὲ ἤλπισε, καὶ ῥύσομαι αὐτόν· σκεπάσω αὐτὸν, ὅτι ἔγνω τὸ ὄνομά μου.
\vs{15}Ἐπικαλέσεται πρὸς μὲ, καὶ εἰσακούσομαι αὐτοῦ, μετʼ αὐτοῦ εἰμι ἐν θλίψει, καὶ ἐξελοῦμαι αὐτὸν, καὶ δοξάσω αὐτόν.
\vs{16}Μακρότητι ἡμερῶν ἐμπλήσω αὐτὸν, καὶ δείξω αὐτῷ τὸ σωτήριόν μου.

\begin{psalmheading}{\ch{91}{92} Ψαλμὸς ᾠδῆς εἰς τὴν ἡμέραν τοῦ σαββάτου.}
\end{psalmheading}
\vs{2}Ἀγαθὸν τὸ ἐξομολογεῖσθαι τῷ Κυρίῳ, καὶ ψάλλειν τῷ ὀνόματί σου, ὕψιστε·
\vs{3}τοῦ ἀναγγέλλειν τοπρωῒ τὸ ἔλεός σου, καὶ τὴν ἀλήθειάν σου κατὰ νύκτα,
\vs{4}ἐν δεκαχόρδῳ ψαλτηρίῳ, μετʼ ᾠδῆς ἐν κιθάρᾳ.
\vs{5}Ὅτι εὔφρανάς με, Κύριε, ἐν τῷ ποιήματί σου, καὶ ἐν τοῖς ἔργοις τῶν χειρῶν σου ἀγαλλιάσομαι.

\vs{6}Ὡς ἐμεγαλύνθη τὰ ἔργα σου, Κύριε; σφόδρα ἐβαθύνθησαν οἱ διαλογισμοί σου.
\vs{7}Ἀνὴρ ἄφρων οὐ γνώσεται, καὶ ἀσύνετος οὐ συνήσει ταῦτα.
\vs{8}Ἐν τῷ ἀνατεῖλαι τοὺς ἁμαρτωλοὺς ὡσεὶ χόρτον, καὶ διέκυψαν πάντες οἱ ἐργαζόμενοι τὴν ἀνομίαν, ὅπως ἂν ἐξολοθρευθῶσιν εἰς τὸν αἰῶνα τοῦ αἰῶνος.
\vs{9}Σὺ δὲ ὕψιστος εἰς τὸν αἰῶνα, Κύριε.

\vs{10}Ὅτι ἰδοὺ οἱ ἐχθροί σου ἀπολοῦνται, καὶ διασκορπισθήσονται πάντες οἱ ἐργαζόμενοι τὴν ἀνομίαν.
\vs{11}Καὶ ὑψωθήσεται ὡς μονοκέρωτος τὸ κέρας μου, καὶ τὸ γῆράς μου ἐν ἐλέῳ πίονι.
\vs{12}Καὶ ἐπεῖδεν ὁ ὀφθαλμός μου ἐν τοῖς ἐχθροῖς μου, καὶ ἐν τοῖς ἐπανισταμένοις ἐπʼ ἐμὲ πονηρευομένοις ἀκούσεται τὸ οὖς μου.

\vs{13}Δίκαιος ὡς φοῖνιξ ἀνθήσει, ὡς ἡ κέδρος ἡ ἐν τῷ Λιβάνῳ πληθυνθήσεται·
\vs{14}Πεφυτευμένοι ἐν τῷ οἴκῳ Κυρίου, ἐν ταῖς αὐλαῖς τοῦ Θεοῦ ἡμῶν ἐξανθήσουσι.
\vs{15}Τότε πληθυνθήσονται ἐν γήρει πίονι, καὶ εὐπαθοῦντες ἔσονται
\vs{16}τοῦ ἀναγγεῖλαι· ὅτι εὐθὴς Κύριος ὁ Θεός μου, καὶ οὐκ ἔστιν ἀδικία ἐν αὐτῷ.

\begin{psalmheading}{\ch{92}{93} Εἰς τὴν ἡμέραν τοῦ προσαββάτου, ὅτε κατῴκισται ἡ γῆ, αἶνος ᾠδῆς τῷ Δαυίδ.}
\end{psalmheading}
Ὁ Κύριος ἐβασίλευσεν, εὐπρέπειαν ἐνεδύσατο, ἐνεδύσατο Κύριος δύναμιν καὶ περιεζώσατο· καὶ γὰρ ἐστερέωσε τὴν οἰκουμένην, ἥτις οὐ σαλευθήσεται.
\vs{2}Ἕτοιμος ὁ θρόνος σου ἀπὸ τότε, ἀπὸ τοῦ αἰῶνος σὺ εἶ.
\vs{3}Ἐπῇραν οἱ ποταμοὶ Κύριε, ἐπῇραν οἱ ποταμοὶ φωνὰς αὐτῶν,
\vs{4}ἀπὸ φωνῶν ὑδάτων πολλῶν· θαυμαστοὶ οἱ μετεωρισμοὶ τῆς θαλάσσης· θαυμαστὸς ἐν ὑψηλοῖς ὁ Κύριος.
\vs{5}Τὰ μαρτύριά σου ἐπιστώθησαν σφόδρα· τῷ οἴκῳ σου πρέπει ἁγίασμα, Κύριε, εἰς μακρότητα ἡμερῶν.

\begin{psalmheading}{\ch{93}{94} Ψαλμὸς τῷ Δαυὶδ τετράδι σαββάτου.}
\end{psalmheading}
Θεὸς ἐκδικήσεων Κύριος, ὁ Θεὸς ἐκδικήσεων ἐπαῤῥησιάσατο.
\vs{2}Ὑψώθητι ὁ κρίνων τὴν γῆν, ἀπόδος ἀνταπόδοσιν τοῖς ὑπερηφάνοις.

\vs{3}Ἕως πότε ἁμαρτωλοὶ, Κύριε, ἕως πότε ἁμαρτωλοὶ καυχήσονται;
\vs{4}Φθέγξονται καὶ λαλήσουσιν ἀδικίαν, λαλήσουσι πάντες οἱ ἐργαζόμενοι τὴν ἀνομίαν.
\vs{5}Τὸν λαόν σου, Κύριε, ἐταπείνωσαν, καὶ τὴν κληρονομίαν σου ἐκάκωσαν.
\vs{6}Χήραν καὶ ὀρφανὸν ἀπέκτειναν, καὶ προσήλυτον ἐφόνευσαν.
\vs{7}Καὶ εἶπαν, οὐκ ὄψεται Κύριος, οὐδὲ συνήσει ὁ Θεὸς τοῦ Ἰακώβ.

\vs{8}Σύνετε δὴ ἄφρονες ἐν τῷ λαῷ, καὶ μωροὶ, ποτὲ φρονήσατε.
\vs{9}Ὁ φυτεύσας τὸ οὖς, οὐχὶ ἀκούει; ἢ ὁ πλάσας τὸν ὀφθαλμὸν, οὐχὶ κατανοεῖ;
\vs{10}Ὁ παιδεύων ἔθνη, οὐχὶ ἐλέγξει; ὁ διδάσκων ἄνθρωπον γνῶσιν;
\vs{11}Κύριος γινώσκει τοὺς διαλογισμοὺς τῶν ἀνθρώπων, ὅτι εἰσὶ μάταιοι.

\vs{12}Μακάριος ὁ ἄνθρωπος ὃν ἂν σὺ παιδεύσῃς Κύριε, καὶ ἐκ τοῦ νόμου σου διδάξῃς αὐτόν·
\vs{13}τοῦ πραΰναι αὐτῷ ἀφʼ ἡμερῶν πονηρῶν, ἕως οὗ ὀρυγῇ τῷ ἁμαρτωλῷ βόθρος.
\vs{14}Ὅτι οὐκ ἀπώσεται Κύριος τὸν λαὸν αὐτοῦ, καὶ τὴν κληρονομίαν αὐτοῦ οὐκ ἐγκαταλείψει,
\vs{15}ἕως οὗ δικαιοσύνη ἐπιστρέψῃ εἰς κρίσιν, καὶ ἐχόμενοι αὐτῆς πάντες οἱ εὐθεῖς τῇ καρδίᾳ· διάψαλμα.

\vs{16}Τίς ἀναστήσεταί μοι ἐπὶ πονηρευομένους, ἢ τίς συμπαραστήσεταί μοι ἐπὶ τοὺς ἐργαζομένους τὴν ἀνομίαν;
\vs{17}Εἰ μὴ ὅτι Κύριος ἐβοήθησέ μοι, παραβραχὺ παρῴκησε τῷ ᾅδῃ ἡ ψυχή μου.
\vs{18}Εἰ ἔλεγον, σεσάλευται ὁ πούς μου, τὸ ἔλεός σου Κύριε ἐβοήθει μοι.
\vs{19}Κύριε, κατὰ τὸ πλῆθος τῶν ὀδυνῶν μου ἐν τῇ καρδίᾳ μου, αἱ παρακλήσεις σου ἠγάπησαν τὴν ψυχήν μου.

\vs{20}Μὴ συμπροσέσται σοι θρόνος ἀνομίας, ὁ πλάσσων κόπον ἐπὶ προστάγματι.
\vs{21}Θηρεύσουσιν ἐπὶ ψυχὴν δικαίου, καὶ αἷμα ἀθῶον καταδικάσονται.
\vs{22}Καὶ ἐγένετό μοι Κύριος εἰς καταφυγὴν, καὶ ὁ Θεός μου εἰς βοηθὸν ἐλπίδος μου.
\vs{23}Καὶ ἀποδώσει αὐτοῖς τὴν ἀνομίαν αὐτῶν, καὶ τὴν πονηρίαν αὐτῶν· ἀφανιεῖ αὐτοὺς Κύριος ὁ Θεὸς ἡμῶν.

\begin{psalmheading}{\ch{94}{95} Αἶνος ᾠδῆς τῷ Δαυίδ.}
\end{psalmheading}
Δεῦτε ἀγαλλιασώμεθα τῷ Κυρίῳ, ἀλαλάξωμεν τῷ Θεῷ τῷ σωτῆρι ἡμῶν.
\vs{2}Προφθάσωμεν τὸ πρόσωπον αὐτοῦ ἐν ἐξομολογήσει, καὶ ἐν ψαλμοῖς ἀλαλάξωμεν αὐτῷ.
\vs{3}Ὅτι Θεὸς μέγας Κύριος, καὶ βασιλεὺς μέγας ἐπὶ πάντας τοὺς θεούς·
\vs{4}ὅτι οὐκ ἀπώσεται Κύριος τὸν λαὸν αὐτοῦ, ὅτι ἐν τῇ χειρὶ αὐτοῦ τὰ πέρατα τῆς γῆς, καὶ τὰ ὕψη τῶν ὀρέων αὐτοῦ ἐστιν.
\vs{5}Ὅτι αὐτοῦ ἐστιν ἡ θάλασσα καὶ αὐτὸς ἐποίησεν αὐτὴν, καὶ τὴν ξηρὰν χεῖρες αὐτοῦ ἔπλασαν.

\vs{6}Δεῦτε προσκυνήσωμεν καὶ προσπέσωμεν αὐτῷ, καὶ κλαύσωμεν ἐναντίον Κυρίου τοῦ ποίησαντος ἡμᾶς.
\vs{7}Ὅτι αὐτός ἐστιν ὁ Θεὸς ἡμῶν, καὶ ἡμεῖς λαὸς νομῆς αὐτοῦ, καὶ πρόβατα χειρὸς αὐτοῦ·
\vs{8}σήμερον ἐὰν τῆς φωνῆς αὐτοῦ ἀκούσητε, μὴ σκληρύνητε τὰς καρδίας ὑμῶν, ὡς ἐν τῷ παραπικρασμῷ, κατὰ τὴν ἡμέραν τοῦ πικρασμοῦ ἐν τῇ ἐρήμῳ,
\vs{9}οὗ ἐπείρασάν με οἱ πατέρες ὑμῶν· ἐδοκίμασαν, καὶ εἶδον τὰ ἔργα μου.
\vs{10}Τεσσαράκοντα ἔτη προσώχθισα τῇ γενεᾷ ἐκείνῃ, καὶ εἶπα, ἀεὶ πλανῶνται τῇ καρδίᾳ, καὶ αὐτοὶ οὐκ ἔγνωσαν τὰς ὁδούς μου.
\vs{11}Ὡς ὤμοσα ἐν τῇ ὀργῇ μου, εἰ εἰσελεύσονται εἰς τὴν κατάπαυσίν μου.

\begin{psalmheading}{\ch{95}{96} Ὅτε ὁ οἶκος ᾠκοδόμηται μετὰ τὴν αἰχμαλωσίαν, ᾠδὴ τῷ Δαυίδ.}
\end{psalmheading}
Ἄσατε τῷ Κυρίῳ ἆσμα καινὸν, ᾄσατε τῷ Κυρίῳ πᾶσα ἡ γῆ.
\vs{2}Ἄσατε τῷ Κυρίῳ, εὐλογήσατε τὸ ὄνομα αὐτοῦ, εὐαγγελίζεσθε ἡμέραν ἐξ ἡμέρας τὸ σωτήριον αὐτοῦ.
\vs{3}Ἀναγγείλατε ἐν τοῖς ἔθνεσι τὴν δόξαν αὐτοῦ, ἐν πᾶσι τοῖς λαοῖς τὰ θαυμάσια αὐτοῦ.

\vs{4}Ὅτι μέγας Κύριος καὶ αἰνετὸς σφόδρα, φοβερός ἐστιν ἐπὶ πάντας τοὺς θεούς.
\vs{5}Ὅτι πάντες οἱ θεοὶ τῶν ἐθνῶν δαιμόνια, ὁ δὲ Κύριος τοὺς οὐρανοὺς ἐποίησεν.
\vs{6}Ἐξομολόγησις καὶ ὡραιότης ἐνώπιον αὐτοῦ, ἁγιωσύνη καὶ μεγαλοπρέπεια ἐν τῷ ἁγιάσματι αὐτοῦ.

\vs{7}Ἐνέγκατε τῷ Κυρίῳ αἱ πατριαὶ τῶν ἐθνῶν, ἐνέγκατε τῷ Κυρίῳ δόξαν καὶ τιμὴν,
\vs{8}ἐνέγκατε τῷ Κυρίῳ δόξαν ὀνόματι αὐτοῦ, ἄρατε θυσίας καὶ εἰσπορεύεσθε εἰς τὰς αὐλὰς αὐτοῦ·
\vs{9}Προσκυνήσατε τῷ Κυρίῳ ἐν αὐλῇ ἁγίᾳ αὐτοῦ, σαλευθήτω ἀπὸ προσώπου αὐτοῦ πᾶσα ἡ γῆ.
\vs{10}Εἴπατε ἐν τοῖς ἔθνεσιν, ὁ Κύριος ἐβασίλευσε· καὶ γὰρ κατώρθωσε τὴν σἰκουμένην, ἥτις οὐ σαλευθήσεται, κρινεῖ λαοῦς ἐν εὐθύτητι.
\vs{11}Εὐφραινέσθωσαν οἱ οὐρανοὶ καὶ ἀγαλλιάσθω ἡ γῆ, σαλευθήτω ἡ θάλασσα καὶ τὸ πλήρωμα αὐτῆς.
\vs{12}Χαρήσεται τὰ πεδία, καὶ πάντα τὰ ἐν αὐτοῖς· τότε ἀγαλλιάσονται πάντα τὰ ξύλα τοῦ δρυμοῦ
\vs{13}πρὸ προσώπου τοῦ Κυρίου, ὅτι ἔρχεται, ὅτι ἔρχεται κρῖναι τὴν γῆν· κρινεῖ τὴν οἰκουμένην ἐν δικαιοσύνῃ, καὶ λαοὺς ἐν τῇ ἀληθείᾳ αὐτοῦ.

\begin{psalmheading}{\ch{96}{97} Τῷ Δαυὶδ, ὅτε ἡ γῆ αὐτοῦ καθίσταται.}
\end{psalmheading}
Ὁ Κύριος ἐβασίλευσεν, ἀγαλλιάσθω ἡ γῆ, εὐφρανθήτωσαν νῆσοι πολλαί.

\vs{2}Νεφέλη καὶ γνόφος κύκλῳ αὐτοῦ, δικαιοσύνη καὶ κρίμα κατόρθωσις τοῦ θρόνου αὐτοῦ.
\vs{3}Πῦρ ἐναντίον αὐτοῦ προπορεύσεται, καὶ φλογιεῖ κύκλῳ τοὺς ἐχθροὺς αὐτοῦ.
\vs{4}Ἔφαναν αἱ ἀστραπαὶ αὐτοῦ τῇ οἰκουμένῃ, εἶδε καὶ ἐσαλεύθη ἡ γῆ·
\vs{5}τὰ ὄρη ὡσεὶ κηρὸς ἐτάκησαν ἀπὸ προσώπου Κυρίου, ἀπὸ προσώπου Κυρίου πάσης τῆς γῆς.
\vs{6}Ἀνήγγειλαν αἱ οὐρανοὶ τὴν δικαιοσύνην αὐτοῦ, καὶ εἴδοσαν πάντες οἱ λαοὶ τὴν δόξαν αὐτοῦ.

\vs{7}Αἰσχυνθήτωσαν πάντες οἱ προσκυνοῦντες τοῖς γλυπτοῖς, οἱ ἐγκαυχώμενοι ἐν τοῖς εἰδώλοις αὐτῶν· προσκυνήσατε αὐτῷ πάντες ἄγγελοι αὐτοῦ.

\vs{8}Ἤκουσε καὶ εὐφράνθη Σιὼν, καὶ ἠγαλλιάσαντο αἱ θυγατέρες τῆς Ἰουδαίας, ἕνεκεν τῶν κριμάτων σου Κύριε.
\vs{9}Ὅτι σὺ εἶ Κύριος ὁ ὕψιστος ἐπὶ πᾶσαν τὴν γῆν, σφόδρα ὑπερυψώθης ὑπὲρ πάντας τοὺς θεούς.

\vs{10}Οἱ ἀγαπῶντες τὸν Κύριον, μισεῖτε πονηρόν· φυλάσσει Κύριος τὰς ψυχὰς τῶν ὁσίων αὐτοῦ, ἐκ χειρὸς ἁμαρτωλῶν ῥύσεται αὐτούς.
\vs{11}Φῶς ἀνέτειλε τῷ δικαίῳ, καὶ τοῖς εὐθέσι τῇ καρδίᾳ εὐφροσύνη.
\vs{12}Εὐφράνθητε δίκαιοι ἐν τῷ Κυρίῳ, καὶ ἐξομολογεῖσθε τῇ μνήμῃ τῆς ἁγιωσύνης αὐτοῦ.

\begin{psalmheading}{\ch{97}{98} Ψαλμὸς τῷ Δαυίδ.}
\end{psalmheading}
Ἄσατε τῷ Κυρίῳ ᾆσμα καινὸν, ὅτι θαυμαστὰ ἐποίησεν ὁ Κύριος· ἔσωσεν αὐτῷ ἡ δεξιὰ αὐτοῦ, καὶ ὁ βραχίων ὁ ἅγιος αὐτοῦ.

\vs{2}Ἐγνώρισε Κύριος τὸ σωτήριον αὐτοῦ, ἐναντίον τῶν ἐθνῶν ἀπεκάλυψε τὴν δικαιοσύνην αὐτοῦ.
\vs{3}Ἐμνήσθη τοῦ ἐλέους αὐτοῦ τῷ Ἰακὼβ, καὶ τῆς ἀληθείας αὐτοῦ τῷ οἴκῳ Ἰσραήλ· εἴδοσαν πάντα τὰ πέρατα τῆς γῆς τὸ σωτήριον τοῦ θεοῦ ἡμῶν.

\vs{4}Ἀλαλάξατε τῷ Θεῷ πᾶσα ἡ γῆ, ᾄσατε καὶ ἀγαλλιᾶσθε καὶ ψάλατε.
\vs{5}Ψάλατε τῷ Κυρίῳ ἐν κιθάρᾳ, ἐν κιθάρᾳ καὶ φωνῇ ψαλμοῦ.
\vs{6}Ἐν σάλπιγξιν ἐλαταῖς, καὶ φωνῇ σάλπιγγος κερατίνης· ἀλαλάξατε ἐνώπιον τοῦ βασιλέως Κυρίῳ.
\vs{7}Σαλευθήτω ἡ θάλασσα καὶ τὸ πλήρωμα αὐτῆς, ἡ οἰκουμένη καὶ οἱ κατοικοῦντες αὐτήν.
\vs{8}Ποταμοὶ κροτήσουσι χειρὶ ἐπιτοαυτὸ, τὰ ὄρη ἀγαλλιάσονται.
\vs{9}Ὅτι ἥκει κρῖναι τὴν γῆν· κρινεῖ τὴν οἰκουμένην ἐν δικαιοσύνῃ, καὶ λαοὺς ἐν εὐθύτητι.

\begin{psalmheading}{\ch{98}{99} Ψαλμὸς τῷ Δαυίδ.}
\end{psalmheading}
Ὁ Κύριος ἐβασίλευσεν, ὀργιζέσθωσαν λαοί· ὁ καθήμενος ἐπὶ τῶν χερουβὶμ, σαλευθήτω ἡ γῆ.
\vs{2}Κύριος ἐν Σιὼν μέγας, καὶ ὑψηλός ἐστιν ἐπὶ πάντας τοὺς λαούς.
\vs{3}Ἐξομολογησάσθωσαν τῷ ὀνόματί σου τῷ μεγάλῳ, ὅτι φοβερὸν καὶ ἅγιόν ἐστι,
\vs{4}καὶ τιμὴ βασιλέως κρίσιν ἀγαπᾷ· σὺ ἡτοίμασας εὐθύτητας, κρίσιν καὶ δικαιοσύνην ἐν Ἰακὼβ σὺ ἐποίησας.
\vs{5}Ὑψοῦτε Κύριον τὸν Θεὸν ἡμῶν, καὶ προσκυνεῖτε τῷ ὑποποδίῳ τῶν ποδῶν αὐτοῦ, ὅτι ἅγιός ἐστι.

\vs{6}Μωυσῆς καὶ Ἀαρὼν ἐν τοῖς ἱερεῦσιν αὐτοῦ, καὶ Σαμουὴλ ἐν τοῖς ἐπικαλουμένοις τὸ ὄνομα αὐτοῦ· ἐπεκαλοῦντο τὸν Κύριον, καὶ αὐτὸς εἰσήκουεν,
\vs{7}ἐν στύλῳ νεφέλης ἐλάλει πρὸς αὐτούς· ἐφύλασσον τὰ μαρτύρια αὐτοῦ, καὶ τὰ προστάγματα ἃ ἔδωκεν αὐτοῖς.
\vs{8}Κύριε ὁ Θεὸς ἡμῶν, σὺ ἐπήκουες αὐτῶν· ὁ Θεὸς, εὐίλατος ἐγίνου αὐτοῖς, καὶ ἐκδικῶν ἐπὶ πάντα τὰ ἐπιτηδεύματα αὐτῶν.
\vs{9}Ὑψοῦτε Κύριον τὸν Θεὸν ἡμῶν, καὶ προσκυνεῖτε εἰς ὄρος ἅγιον αὐτοῦ, ὅτι ἅγιος Κύριος ὁ Θεὸς ἡμῶν.

\begin{psalmheading}{\ch{99}{100} Ψαλμὸς εἰς ἐξομολόγησιν.}
\end{psalmheading}
Ἀλαλάξατε τῷ Κυρίῳ πᾶσα ἡ γῆ,
\vs{2}δουλεύσατε τῷ Κυρίῳ ἐν εὐφροσύνῃ· εἰσέλθατε ἐνώπιον αὐτοῦ ἐν ἀγαλλιάσει.
\vs{3}Γνῶτε ὅτι Κύριος αὐτός ἐστιν ὁ Θεός· αὐτὸς ἐποίησεν ἡμᾶς, καὶ οὐχ ἡμεῖς, λαὸς αὐτοῦ καὶ πρόβατα τῆς νομῆς αὐτοῦ.
\vs{4}Εἰσέλθατε εἰς τὰς πύλας αὐτοῦ ἐν ἐξομολογήσει, τὰς αὐλὰς αὐτοῦ ἐν ὕμνοις· ἐξομολογεῖσθε αὐτῷ, αἰνεῖτε τὸ ὄνομα αὐτοῦ.
\vs{5}Ὅτι χρηστὸς Κύριος, εἰς τὸν αἰῶνα τὸ ἔλεος αὐτοῦ, καὶ ἕως γενεᾶς καὶ γενεᾶς ἡ ἀλήθεια αὐτοῦ.

\begin{psalmheading}{\ch{100}{101} Ψαλμὸς τῷ Δαυίδ.}
\end{psalmheading}
Ἔλεος καὶ κρίσιν ᾄσομαί σοι, Κύριε·
\vs{2}ψαλῶ καὶ συνήσω ἐν ὁδῷ ἀμώμῳ· πότε ἥξεις πρὸς μέ; διεπορευόμην ἐν ἀκακίᾳ καρδίας μου, ἐν μέσῳ τοῦ οἴκου μου.
\vs{3}Οὐ προεθέμην πρὸ ὀφθαλμῶν μου πρᾶγμα παράνομον, ποιοῦντας παραβάσεις ἐμίσησα· οὐκ ἐκολλήθη μοικαρδία σκαμβὴ,
\vs{4}ἐκκλίνοντος ἀπʼ ἐμοῦ τοῦ πονηροῦ οὐκ ἐγίνωσκον.
\vs{5}Τὸν καταλαλοῦντα λάθρα τοῦ πλησίον αὐτοῦ, τοῦτον ἐξεδίωκον· ὑπερηφάνῳ ὀφθαλμῷ καὶ ἀπλήστῳ καρδίᾳ, τούτῳ οὐ συνήσθιον.
\vs{6}Οἱ ὀφθαλμοί μου ἐπὶ τοὺς πιστοὺς τῆς γῆς, τοῦ συγκαθῆσθαι αὐτοὺς μετʼ ἐμοῦ· πορευόμενος ἐν ὁδῷ ἀμώμῳ, οὗτός μοι ἐλειτούργει.
\vs{7}Οὐ κατῴκει ἐν μέσῳ τῆς οἰκίας μου ποιῶν ὑπερηφανίαν· λαλῶν ἄδικα οὐ κατεύθυνεν ἐναντίον τῶν ὀφθαλμῶν μου.
\vs{8}Εἰς τὰς πρωΐας ἀπέκτενον πάντας τοὺς ἁμαρτωλοὺς τῆς γῆς, τοῦ ἐξολοθρεῦσαι ἐκ πόλεως Κυρίου πάντας τοὺς ἐργαζομένους τὴν ἀδικίαν.

\begin{psalmheading}{\ch{101}{102} Προσευχὴ τῷ πτωχῷ, ὅταν ἀκηδιάσῃ, καὶ ἐναντίον Κυρίου ἐκχέῃ τὴν δέησιν αὐτοῦ.}
\end{psalmheading}
\vs{2}Κύριε εἰσάκουσον τῆς προσευχῆς μου, καὶ ἡ κραυγή μου πρὸς σὲ ἐλθέτω.
\vs{3}Μὴ ἀποστρέψῃς τὸ πρόσωπόν σου ἀπʼ ἐμοῦ· ἐν ᾗ ἂν ἡμέρᾳ θλίβομαι, κλῖνον πρὸς μὲ τὸ οὖς σου· ἐν ᾗ ἂν ἡμέρᾳ ἐπικαλέσωμαί σε, ταχὺ εἰσάκουσόν μου.

\vs{4}Ὅτι ἐξέλιπον ὡσεὶ καπνὸς αἱ ἡμέραι μου, καὶ τὰ ὀστᾶ μου ὡσεὶ φρύγιον συνεφρύγησαν.
\vs{5}Ἐπλήγην ὡσεὶ χόρτος, καὶ ἐξηράνθη ἡ καρδία μου, ὅτι ἐπελαθόμην τοῦ φαγεῖν τὸν ἄρτον μου.
\vs{6}Ἀπὸ φωνῆς τοῦ στεναγμοῦ μου, ἐκολλήθη τὸ ὀστοῦν μου τῇ σαρκί μου.
\vs{7}Ὡμοιώθην πελεκᾶνι ἐρημικῷ, ἐγενήθην ὡσεὶ νυκτικόραξ ἐν οἰκοπέδῳ.
\vs{8}Ἠγρύπνησα, καὶ ἐγενήθην ὡσεὶ στρουθίον μονάζον ἐπὶ δώματι.
\vs{9}Ὅλην τὴν ἡμέραν ὠνείδιζόν με οἱ ἐχθροί μου, καὶ οἱ ἐπαινοῦντές με κατʼ ἐμοῦ ὤμνυον.
\vs{10}Ὅτι σποδὸν ὡσεὶ ἄρτον ἔφαγον, καὶ τὸ πόμα μου μετὰ κλαυθμοῦ ἐκίρνων,
\vs{11}ἀπὸ προσώπου τῆς ὀργῆς σου καὶ τοῦ θυμοῦ σου, ὅτι ἐπάρας κατέῤῥαξάς με.

\vs{12}Αἱ ἡμέραι μου ὡσεὶ σκιὰ ἐκλίθησαν, κᾀγὼ ὡσεὶ χόρτος ἐξηράνθην.
\vs{13}Σὺ δέ Κύριε εἰς τὸν αἰῶνα μένεις, καὶ τὸ μνημόσυνόν σου εἰς γενεὰν καὶ γενεάν.
\vs{14}Σὺ ἀναστὰς οἰκτειρήσεις τὴν Σιὼν, ὅτι καιρὸς τοῦ οἰκτειρῆσαι αὐτὴν, ὅτι ἥκει καιρός.
\vs{15}Ὅτι εὐδόκησαν οἱ δοῦλοί σου τοὺς λίθους αὐτῆς, καὶ τὸν χοῦν αὐτῆς οἰκτειρήσουσι.
\vs{16}Καὶ φοβηθήσονται τὰ ἔθνη τὸ ὄνομά σου Κύριε, καὶ πάντες οἱ βασιλεῖς τὴν δόξαν σου.

\vs{17}Ὅτι οἰκοδομήσει Κύριος τὴν Σιὼν, καὶ ὀφθήσεται ἐν τῇ δόξῃ αὐτοῦ.
\vs{18}Ἐπέβλεψεν ἐπὶ τὴν προσευχὴν τῶν ταπεινῶν, καὶ οὐκ ἐξουδένωσε τὴν δέησιν αὐτῶν.
\vs{19}Γραφήτω αὕτη εἰς γενεὰν ἑτέραν, καὶ λαὸς ὁ κτιζόμενος αἰνέσει τὸν Κύριον.
\vs{20}Ὅτι ἐξέκυψεν ἐξ ὕψους ἁγίου αὐτοῦ, Κύριος ἐξ οὐρανοῦ ἐπὶ τὴν γῆν ἐπέβλεψε,
\vs{21}τοῦ ἀκοῦσαι τοῦ στεναγμοῦ τῶν πεπεδημένων, τοῦ λῦσαι τοὺς υἱοὺς τῶν τεθανατωμένων,
\vs{22}τοῦ ἀναγγεῖλαι ἐν Σιὼν τὸ ὄνομα Κυρίου, καὶ τὴν αἴνεσιν αὐτοῦ ἐν Ἱερουσαλήμ·
\vs{23}ἐν τῷ συναχθῆναι λαοὺς ἐπιτοαυτὸ, καὶ βασιλεῖς τοῦ δουλεύειν τῷ Κυρίῳ.

\vs{24}Ἀπεκρίθη αὐτῷ ἐν ὁδῷ ἰσχύος αὐτοῦ, τὴν ὀλιγότητα τῶν ἡμερῶν μου ἀνάγγειλόν μοι·
\vs{25}μὴ ἀναγάγῃς με ἐν ἡμίσει ἡμερῶν μου, ἐν γενεᾷ γενεῶν τὰ ἔτη σου.
\vs{26}Κατʼ ἀρχὰς τὴν γῆν σὺ Κύριε ἐθεμελίωσας, καὶ ἔργα τῶν χειρῶν σου εἰσὶν οἱ οὐρανοί.
\vs{27}Αὐτοὶ ἀπολοῦνται, σὺ δὲ διαμένεις· καὶ πάντες ὡς ἱμάτιον παλαιωθήσονται, καὶ ὡσεὶ περιβόλαιον ἑλίξεις αὐτοὺς, καὶ ἀλλαγήσονται.
\vs{28}Σὺ δὲ ὁ αὐτὸς εἶ, καὶ τὰ ἔτη σου οὐκ ἐκλείψουσιν·
\vs{29}Οἱ υἱοῖ τῶν δούλων σου κατασκηνώσουσι, καὶ τὸ σπέρμα αὐτῶν εἰς τὸν αἰῶνα κατευθυνθήσεται.

\begin{psalmheading}{\ch{102}{103} Τῷ Δαυίδ.}
\end{psalmheading}
Εὐλόγει ἡ ψυχή μου τὸν Κύριον, καὶ πάντα τὰ ἐντός μου τὸ ὄνομα τὸ ἅγιον αὐτοῦ.
\vs{2}Εὐλόγει ἡ ψυχή μου τὸν Κύριον, καὶ μὴ ἐπιλανθάνου πάσας τὰς αἰνέσεις αὐτοῦ·
\vs{3}Τὸν εὐιλατεύοντα πάσαις ταῖς ἀνομίαις σου, τὸν ἰώμενον πάσας τὰς νόσους σου,
\vs{4}τὸν λυτρούμενον ἐκ φθορᾶς τὴν ζωήν σου, τὸν στεφανοῦντά σε ἐν ἐλέει καὶ οἰκτιρμοῖς,
\vs{5}τὸν ἐμπιπλῶντα ἐν ἀγαθοῖς τὴν ἐπιθυμίαν σου· ἀνακαινισθήσεται ὡς ἀετοῦ ἡ νεότης σου.

\vs{6}Ποιῶν ἐλεημοσύνας ὁ Κύριος, καὶ κρίμα πᾶσι τοῖς ἀδικουμένοις.
\vs{7}Ἐγνώρισε τὰς ὁδοὺς αὐτοῦ τῷ Μωυσῇ, τοῖς υἱοῖς Ἰσραὴλ τὰ θελήματα αὐτοῦ.
\vs{8}Οἰκτίρμων καὶ ἐλεήμων ὁ Κύριος, μακρόθυμος καὶ πολυέλεος.
\vs{9}Οὐκ εἰς τέλος ὀργισθήσεται, οὐδὲ εἰς τὸν αἰῶνα μηνιεῖ.
\vs{10}Οὐ κατὰ τὰς ἁμαρτίας ἡμῶν ἐποίησεν ἡμῖν, οὐδὲ κατὰ τὰς ἀνομίας ἡμῶν ἀνταπέδωκεν ἡμῖν.
\vs{11}Ὅτι κατὰ τὸ ὕψος τοῦ οὐρανοῦ ἀπὸ τῆς γῆς, ἐκραταίωσε Κύριος τὸ ἔλεος αὐτοῦ ἐπὶ τοὺς φοβουμένους αὐτόν.
\vs{12}Καθόσον ἀπέχουσιν ἀνατολαὶ ἀπὸ δυσμῶν, ἐμάκρυνεν ἀφʼ ἡμῶν τὰς ἀνομίας ἡμῶν.
\vs{13}Καθὼς οἰκτείρει πατὴρ υἱοὺς, ᾠκτείρησε Κύριος τοὺς φοβουμένους αὐτόν.
\vs{14}Ὅτι αὐτὸς ἔγνω τὸ πλάσμα ἡμῶν· μνήσθητι ὅτι χοῦς ἐσμεν.

\vs{15}Ἄνθρωπος, ὡσεὶ χόρτος αἱ ἡμέραι αὐτοῦ, ὡσεὶ ἄνθος τοῦ ἀγροῦ οὕτως ἐξανθήσει.
\vs{16}Ὅτι πνεῦμα διῆλθεν ἐν αὐτῷ, καὶ οὐχ ὑπάρξει, καὶ οὐκ ἐπιγνώσεται ἔτι τὸν τόπον αὐτοῦ.
\vs{17}Τὸ δὲ ἔλεος τοῦ Κυρίου ἀπὸ τοῦ αἰῶνος καὶ ἕως τοῦ αἰῶνος ἐπὶ τοὺς φοβουμένους αὐτόν· καὶ ἡ δικαιοσύνη αὐτοῦ ἐπὶ υἱοὺς υἱῶν,
\vs{18}τοῖς φυλάσσουσι τὴν διαθήκην αὐτοῦ, καὶ μεμνημένοις τῶν ἐντολῶν αὐτοῦ τοῦ ποιῆσαι αὐτάς.

\vs{19}Κύριος ἐν τῷ οὐρανῷ ἡτοίμασε τὸν θρόνον αὐτοῦ, καὶ ἡ βασιλεία αὐτοῦ πάντων δεσπόζει.
\vs{20}Εὐλογεῖτε τὸν Κύριον πάντες ἄγγελοι αὐτοῦ, δυνατοὶ ἰσχύϊ ποιοῦντες τὸν λόγον αὐτοῦ, τοῦ ἀκοῦσαι τῆς φωνῆς τῶν λόγων αὐτοῦ.
\vs{21}Εὐλογεῖτε τὸν Κύριον πᾶσαι αἱ δυνάμεις αὐτοῦ, λειτουργοὶ αὐτοῦ ποιοῦντες τὰ θελήματα αὐτοῦ.
\vs{22}Εὐλογεῖτε τὸν Κύριον πάντα τὰ ἔργα αὐτοῦ, ἐν παντὶ τόπῳ τῆς δυναστείας αὐτοῦ· εὐλόγει ἡ ψυχή μου τὸν Κύριον.

\begin{psalmheading}{\ch{103}{104} Τῷ Δαυίδ.}
\end{psalmheading}
Εὐλόγει ἡ ψυχή μου τὸν Κύριον. Κύριε ὁ Θεός μου ἐμεγαλύνθης σφόδρα· ἐξομολόγησιν καὶ εὐπρέπειαν ἐνεδύσω,
\vs{2}ἀναβαλλόμενος φῶς ὡς ἱμάτιον, ἐκτείνων τὸν οὐρανὸν ὡσεὶ δέῤῥιν·
\vs{3}Ὁ στεγάζων ἐν ὕδασι τὰ ὑπερῷα αὐτοῦ, ὁ τιθεὶς νέφη τὴν ἐπίβασιν αὐτοῦ· ὁ περιπατῶν ἐπὶ πτερύγων ἀνέμων·
\vs{4}Ὁ ποιῶν τοὺς ἀγγέλους αὐτοῦ πνεύματα, καὶ τοὺς λειτουργοὺς αὐτοῦ πῦρ φλέγον·

\vs{5}Ὁ θεμελιῶν τὴν γῆν ἐπὶ τὴν ἀσφάλειαν αὐτῆς, οὐ κλιθήσεται εἰς τὸν αἰῶνα τοῦ αἰῶνος.
\vs{6}Ἄβυσσος ὡς ἱμάτιον τὸ περιβόλαιον αὐτοῦ, ἐπὶ τῶν ὀρέων στήσονται ὕδατα.
\vs{7}Ἀπὸ ἐπιτιμήσεώς σου φεύξονται, ἀπὸ φωνῆς βροντῆς σου δειλιάσουσιν.
\vs{8}Ἀναβαίνουσιν ὄρη, καὶ καταβαίνουσι πεδία εἰς τόπον ὃν ἐθεμελίωσας αὐτοῖς.
\vs{9}Ὅριον ἔθου ὃ οὐ παρελεύσονται, οὐδὲ ἐπιστρέψουσι καλύψαι τὴν γῆν.

\vs{10}Ὁ ἐξαποστέλλων πηγὰς ἐν φάραγξιν, ἀναμέσον τῶν ὀρέων διελεύσονται ὕδατα.
\vs{11}Ποτιοῦσι πάντα τὰ θηρία τοῦ ἀγροῦ, προσδέξονται ὄναγροι εἰς δίψαν αὐτῶν.
\vs{12}Ἐπʼ αὐτὰ τὰ πετεινὰ τοῦ οὐρανοῦ κατασκηνώσει, ἐκ μέσου τῶν πετρῶν δώσουσι φωνήν.
\vs{13}Ποτίζων ὄρη ἐκ τῶν ὑπερῴων αὐτοῦ, ἀπὸ καρποῦ τῶν ἔργων σου χορτασθήσεται ἡ γῆ.

\vs{14}Ὁ ἐξανατέλλων χόρτον τοῖς κτήνεσι, καὶ χλόην τῇ δουλείᾳ τῶν ἀνθρώπων· τοῦ ἐξαγαγεῖν ἄρτον ἐκ τῆς γῆς,
\vs{15}καὶ οἶνος εὐφραίνει καρδίαν ἀνθρώπου· τοῦ ἱλαρύναι πρόσωπον ἐν ἐλαίῳ, καὶ ἄρτος καρδίαν ἀνθρώπου στηρίζει.
\vs{16}Χορτασθήσεται τὰ ξύλα τοῦ πεδίου, αἱ κέδροι τοῦ Λιβάνου ἃς ἐφύτευσεν.
\vs{17}Ἐκεῖ στρουθία ἐννοσσεύσουσι, τοῦ ἐρωδιοῦ ἡ οἰκία ἡγεῖται αὐτῶν.
\vs{18}Ὄρη τὰ ὑψηλὰ ταῖς ἐλάφοις, πέτρα καταφυγὴ τοῖς χοιρογρυλλίοις.

\vs{19}Ἐποίησε σελήνην εἰς καιροὺς, ὁ ἥλιος ἔγνω τὴν δύσιν αὐτοῦ.
\vs{20}Ἔθου σκότος καὶ ἐγένετο νὺξ, ἐν αὐτῇ διελεύσονται πάντα τὰ θηρία τοῦ δρυμοῦ.
\vs{21}Σκυμνοι ὠρυόμενοι ἁρπάσαι, καὶ ζητῆσαι παρὰ τοῦ Θεοῦ βρῶσιν αὐτοῖς.
\vs{22}Ἀνέτειλεν ὁ ἥλιος καὶ συναχθήσονται, καὶ ἐν ταῖς μάνδραις αὐτῶν κοιτασθήσονται.
\vs{23}Ἐξελεύσεται ἄνθρωπος ἐπὶ τὸ ἔργον αὐτοῦ, καὶ ἐπὶ τὴν ἐργασίαν αὐτοῦ ἕως ἑσπέρας.

\vs{24}Ὡς ἐμεγαλύνθη τὰ ἔργα σου Κύριε, πάντα ἐν σοφίᾳ ἐποίησας· ἐπληρώθη ἡ γῆ τῆς κτίσεώς σου·
\vs{25}Αὕτη ἡ θάλασσα ἡ μεγάλη καὶ εὐρύχωρος· ἐκεῖ ἑρπετὰ ὧν οὐκ ἔστιν ἀριθμὸς, ζῶα μικρὰ μετὰ μεγάλων.
\vs{26}Ἐκεῖ πλοῖα διαπορεύονται, δράκων οὗτος ὃν ἔπλασας ἐμπαίζειν αὐτῷ.
\vs{27}Πάντα πρὸς σὲ προσδοκῶσι, δοῦναι τὴν τροφὴν αὐτοῖς εὔκαιρον.
\vs{28}Δόντος σου αὐτοῖς, συλλέξουσιν· ἀνοίξαντος δέ σου τὴν χεῖρα, τὰ σύμπαντα πλησθήσονται χρηστότητος.
\vs{29}Ἀποστρέψαντος δέ σου τὸ πρόσωπον, ταραχθήσονται· ἀντανελεῖς τὸ πνεῦμα αὐτῶν, καὶ ἐκλείψουσι, καὶ εἰς τὸν χοῦν αὐτῶν ἐπιστρέψουσιν.
\vs{30}Ἐξαποστελεῖς τὸ πνεῦμά σου καὶ κτισθήσονται, καὶ ἀνακαινιεῖς τὸ πρόσωπον τῆς γῆς.

\vs{31}Ἤτω ἡ δόξα Κυρίου εἰς τὸν αἰῶνα, εὐφρανθήσεται Κύριος ἐπὶ τοῖς ἔργοις αὐτοῦ·
\vs{32}Ὁ ἐπιβλέπων ἐπὶ τὴν γῆν καὶ ποιῶν αὐτὴν τρέμειν, ὁ ἁπτόμενος τῶν ὀρέων καὶ καπνίζονται.
\vs{33}Ἄσω τῷ Κυρίῳ ἐν τῇ ζωῇ μου, ψαλῶ τῷ Θεῷ μου ἕως ὑπάρχω.
\vs{34}Ἡδυνθείη αὐτῷ ἡ διαλογή μου, ἐγὼ δὲ εὐφρανθήσομαι ἐπὶ τῷ Κυρίῳ.
\vs{35}Ἐκλείποισαν ἁμαρτωλοὶ ἀπὸ τῆς γῆς, καὶ ἄνομοι, ὥστε μὴ ὑπάρχειν αὐτούς· εὐλόγει, ἡ ψυχή μου τὸν Κύριον.

\begin{psalmheading}{\ch{104}{105} Ἀλληλούϊα.}
\end{psalmheading}
Ἐξομολογεῖσθε τῷ Κυρίῳ, καὶ ἐπικαλεῖσθε τὸ ὄνομα αὐτοῦ· ἀπαγγείλατε ἐν τοῖς ἔθνεσι τὰ ἔργα αὐτοῦ.
\vs{2}Ασατε αὐτῷ καὶ ψάλατε αὐτῷ. διηγήσασθε πάντα τὰ θαυμάσια αὐτοῦ.
\vs{3}Ἐπαινεῖσθε ἐν τῷ ὀνόματι τῷ ἁγίῳ αὐτοῦ· εὐφρανθήτω καρδία ζητούντων τὸν Κύριον.
\vs{4}Ζητήσατε τὸν Κύριον καὶ κραταιώθητε· ζητήσατε τὸ πρόσωπον αὐτοῦ διαπαντός.
\vs{5}Μνήσθητε τῶν θαυμασίων αὐτον ὧν ἐποίησε, τὰ τέρατα αὐτοῦ, καὶ τὰ κρίματα τοῦ στόματος αὐτοῦ.
\vs{6}Σπέρμα Ἁβραὰμ δοῦλοι αὐτοῦ, υἱοὶ Ἰακὼβ ἐκλεκτοὶ αὐτοῦ.

\vs{7}Αὐτὸς Κύριος ὁ Θεὸς ἡμῶν, ἐν πάσῃ τῇ γῇ τὰ κρίματα αὐτοῦ.
\vs{8}Ἐμνήσθη εἰς τὸν αἰῶνα διαθήκης αὐτοῦ, λόγου οὗ ἐνετείλατο εἰς χιλίας γενεὰς,
\vs{9}ὃν διέθετο τῷ Ἁβραὰμ, καὶ τοῦ ὅρκου αὐτοῦ τῷ Ἰσαάκ·
\vs{10}Καὶ ἔστησεν αὐτὴν τῷ Ἰακὼβ εἰς πρόσταγμα, καὶ τῷ Ἰσραὴλ εἰς διαθήκην αἰώνιον,
\vs{11}λέγων, σοὶ δώσω τὴν γῆν Χαναὰν, σχοίνισμα κληρονομίας ὑμῶν.
\vs{12}Ἐν τῷ εἶναι αὐτοὺς ἀριθμῷ βραχεῖς, ὀλιγοστοὺς καὶ παροίκους ἐν αὐτῇ,
\vs{13}καὶ διῆλθον ἐξ ἔθνους εἰς ἔθνος, καὶ ἐκ βασιλείας εἰς λαὸν ἕτερον,
\vs{14}οὐκ ἀφῆκεν ἄνθρωπον ἀδικῆσαι αὐτοὺς, καὶ ἤλεγξεν ὑπὲρ αὐτῶν βασιλεῖς·
\vs{15}Μὴ ἅψησθε τῶν χριστῶν μου, καὶ ἐν τοῖς προφήταις μου μὴ πονηρεύεσθε.
\vs{16}Καὶ ἐκάλεσε λιμὸν ἐπὶ τὴν γῆν, πᾶν στήριγμα ἄρτου συνέτριψεν.

\vs{17}Ἀπέστειλεν ἔμπροσθεν αὐτῶν ἄνθρωπον, εἰς δοῦλον ἐπράθη Ἰωσήφ.
\vs{18}Ἐταπείνωσαν ἐν πέδαις τοὺς πόδας αὐτοῦ, σίδηρον διῆλθεν ἡ ψυχὴ αὐτοῦ·
\vs{19}Μέχρι τοῦ ἐλθεῖν τὸν λόγον αὐτοῦ· τὸ λόγιον τοῦ Κυρίου ἐπύρωσεν αὐτόν.
\vs{20}Ἀπέστειλε βασιλεὺς καὶ ἔλυσεν αὐτὸν, ἄρχων λαῶν καὶ ἀφῆκεν αὐτόν.
\vs{21}Κατέστησεν αὐτὸν κύριον τοῦ οἴκου αὐτοῦ, καὶ ἄρχοντα πάσης τῆς κτήσεως αὐτοῦ,
\vs{22}τοῦ παιδεῦσαι τοὺς ἄρχοντας αὐτοῦ ὡς ἑαυτὸν, καὶ τοὺς πρεσβυτέρους αὐτοῦ σοφίσαι.

\vs{23}Καὶ εἰσῆλθεν Ἰσραὴλ εἰς Αἴγυπτον, καὶ Ἰακὼβ παρῴκησεν ἐν γῇ Χάμ.
\vs{24}Καὶ ηὔξησε τὸν λαὸν αὐτοῦ σφόδρα, καὶ ἐκραταίωσεν αὐτὸν ὑπὲρ τοὺς ἐχθροὺς αὐτοῦ.
\vs{25}Καὶ μετέστρεψε τὴν καρδίαν αὐτῶν τοῦ μισῆσαι τὸν λαὸν αὐτοῦ, τοῦ δολιοῦσθαι ἐν τοῖς δούλοις αὐτοῦ.
\vs{26}Ἐξαπέστειλε Μωυσῆν τὸν δοῦλον αὐτοῦ, Ἀαρὼν ὃν ἐξελέξατο αὐτόν.

\vs{27}Ἔθετο ἐν αὐτοῖς τοὺς λόγους τῶν σημείων αὐτοῦ, καὶ τῶν τεράτων ἐν γῇ Χάμ.
\vs{28}Ἐξαπέστειλε σκότος καὶ ἐσκότασε, καὶ παρεπίκραναν τοὺς λόγους αὐτοῦ·
\vs{29}Μετέστρεψε τὰ ὕδατα αὐτῶν εἰς αἷμα, καὶ ἀπέκτεινε τοὺς ἰχθύας αὐτῶν.
\vs{30}Ἐξῆρψεν ἡ γῆ αὐτῶν βατράχους, ἐν τοῖς ταμείοις τῶν βασιλέων αὐτῶν.
\vs{31}Εἶπε καὶ ἦλθε κυνόμυια, καὶ σκνίπες ἐν πᾶσι τοῖς ὁρίοις αὐτῶν.
\vs{32}Ἔθετο τὰς βροχὰς αὐτῶν χάλαζαν, πῦρ καταφλέγον ἐν τῇ γῇ αὐτῶν.
\vs{33}Καὶ ἐπάταξε τὰς ἀμπέλους αὐτῶν καὶ τὰς συκὰς αὐτῶν, καὶ συνέτριψε πᾶν ξύλον ὁρίου αὐτῶν.
\vs{34}Εἶπε καὶ ἦλθεν ἀκρὶς, καὶ βροῦχος οὗ οὐκ ἦν ἀριθμὸς,
\vs{35}καὶ κατέφαγε πάντα τὸν χόρτον ἐν τῇ γῇ αὐτῶν, καὶ κατέφαγε τὸν καρπὸν τῆς γῆς αὐτῶν.
\vs{36}Καὶ ἐπάταξε πᾶν πρωτότοκον ἐκ τῆς γῆς αὐτῶν, ἀπαρχὴν παντὸς πόνου αὐτῶν.
\vs{37}Καὶ ἐξήγαγεν αὐτοὺς ἐν ἀργυρίῳ καὶ χρυσίῳ, καὶ οὐκ ἦν ἐν ταῖς φυλαῖς αὐτῶν ὁ ἀσθενῶν.
\vs{38}Εὐφράνθη Αἴγυπτος ἐν τῇ ἐξόδῳ αὐτῶν, ὅτι ἐπέπεσεν ὁ φόβος αὐτῶν ἐπʼ αὐτούς.
\vs{39}Διεπέτασε νεφέλην εἰς σκέπην αὐτοῖς, καὶ πῦρ τοῦ φωτίσαι αὐτοῖς τὴν νύκτα.
\vs{40}Ἤτησαν, καὶ ἦλθεν ὀρτυγομήτρα, καὶ ἄρτον οὐρανοῦ ἐνέπλησεν αὐτούς.
\vs{41}Διέῤῥηξε πέτραν, καὶ ἐῤῥύησαν ὕδατα, ἐπορεύθησαν ἐν ἀνύδροις ποταμοί.

\vs{42}Ὅτι ἐμνήσθη τοῦ λόγου τοῦ ἁγίου αὐτοῦ, τοῦ πρὸς Ἀβραὰμ τὸν δοῦλον αὐτοῦ·
\vs{43}Καὶ ἐξήγαγε τὸν λαὸν αὐτοῦ ἐν ἀγαλλιάσει, καὶ τοὺς ἐκλεκτοὺς αὐτοῦ ἐν εὐφροσύνῃ·
\vs{44}Καὶ ἔδωκεν αὐτοῖς χώρας ἐθνῶν, καὶ πόνους λαῶν ἐκληρονόμησαν.
\vs{45}Ὅπως ἂν φυλάξωσι τὰ δικαιώματα αὐτοῦ, καὶ τὸν νόμον αὐτοῦ ἐκζητήσωσιν.

\begin{psalmheading}{\ch{105}{106} Ἀλληλούϊα.}
\end{psalmheading}
Ἐξομολογεῖσθε τῷ Κυρίῳ, ὅτι χρηστὸς, ὅτι εἰς τὸν αἰῶνα τὸ ἔλεος αὐτοῦ.
\vs{2}Τίς λαλήσει τὰς δυναστείας τοῦ Κυρίου, ἀκουστὰς ποιήσει πάσας τὰς αἰνέσεις αὐτοῦ;
\vs{3}Μακάριοι οἱ φυλάσσοντες κρίσιν, καὶ ποιοῦντες δικαιοσύνην ἐν παντὶ καιρῷ.

\vs{4}Μνήσθητι ἡμῶν Κύριε ἐν τῇ εὐδοκίᾳ τοῦ λαοῦ σου, ἐπίσκεψαι ἡμᾶς ἐν τῷ σωτηρίῳ σου·
\vs{5}τοῦ ἰδεῖν ἐν τῇ χρηστότητι τῶν ἐκλεκτῶν σου, τοῦ εὐφρανθῆναι ἐν τῇ εὐφροσύνῃ τοῦ ἔθνους σου, τοῦ ἐπαινεῖσθαι μετὰ τῆς κληρονομίας σου.

\vs{6}Ἡμάρτομεν μετὰ τῶν πατέρων ἡμῶν, ἠνομήσαμεν, ἠδικήσαμεν.
\vs{7}Οἱ πατέρες ἡμῶν ἐν Αἰγύπτῳ οὐ συνῆκαν τὰ θαυμάσιά σου, καὶ οὐκ ἐμνήσθησαν τοῦ πλήθους τοῦ ἐλέους σου· καὶ παρεπίκραναν ἀναβαίνοντες ἐν τῇ ἐρυθρᾷ θαλάσσῃ.
\vs{8}Καὶ ἔσωσεν αὐτοὺς ἕνεκεν τοῦ ὀνόματος αὐτοῦ, τοῦ γνωρίσαι τὴν δυναστείαν αὐτοῦ.
\vs{9}Καὶ ἐπετίμησε τῇ ἐρυθρᾷ θαλάσσῃ, καὶ ἐξηράνθη· καὶ ὡδήγησεν αὐτοὺς ἐν ἀβύσσῳ ὡς ἐν ἐρήμῳ·
\vs{10}Καὶ ἔσωσεν αὐτοὺς ἐκ χειρὸς μισούντων, καὶ ἐλυτρώσατο αὐτοὺς ἐκ χειρὸς ἐχθροῦ.
\vs{11}Ἐκάλυψεν ὕδωρ τοὺς θλίβοντας αὐτοὺς, εἷς ἐξ αὐτῶν οὐχ ὑπελείφθη.
\vs{12}Καὶ ἐπίστευσαν τοῖς λόγοις αὐτοῦ, καὶ ᾔνεσαν τὴν αἴνεσιν αὐτοῦ.
\vs{13}Ἐτάχυναν, ἐπελάθοντο τῶν ἔργων αὐτοῦ, οὐχ ὑπέμειναν τὴν βουλὴν αὐτοῦ.
\vs{14}Καὶ ἐπεθύμησαν ἐπιθυμίαν ἐν τῇ ἐρήμῳ, καὶ ἐπείρασαν τὸν Θεὸν ἐν ἀνύδρῳ.
\vs{15}Καὶ ἔδωκεν αὐτοῖς τὸ αἴτημα αὐτῶν, καὶ ἐξαπέστειλε πλησμονὴν εἰς τὴν ψυχὴν αὐτῶν.

\vs{16}Καὶ παρώργισαν Μωυσῆν ἐν τῇ παρεμβολῇ, καὶ Ἀαρὼν τὸν ἅγιον Κυρίου.
\vs{17}Ἠνοίχθη ἡ γῆ καὶ κατέπιε Δαθὰν, καὶ ἐκάλυψεν ἐπὶ τὴν συναγωγὴν Ἀβειρών.
\vs{18}Καὶ ἐξεκαύθη πῦρ ἐν τῇ συναγωγῇ αὐτῶν, καὶ φλὸξ κατέφλεξεν ἁμαρτωλούς.

\vs{19}Καὶ ἐποίησαν μόσχον ἐν Χωρὴβ, καὶ προσεκύνησαν τῷ γλυπτῷ·
\vs{20}Καὶ ἠλλάξαντο τὴν δόξαν αὐτῶν ἐν ὁμοιώματι μόσχου ἔσθοντος χόρτον.
\vs{21}Ἐπελάθοντο τοῦ Θεοῦ τοῦ σώζοντος αὐτοὺς, τοῦ ποιήσαντος μεγάλα ἐν Αἰγύπτῳ,
\vs{22}θαυμαστὰ ἐν γῇ Χὰμ, καὶ φοβερὰ ἐπὶ θαλάσσης ἐρυθρᾶς.
\vs{23}Καὶ εἶπε τοῦ ἐξολοθρεῦσαι αὐτοὺς, εἰ μὴ Μωυσῆς ὁ ἐκλεκτὸς αὐτοῦ ἔστη ἐν τῇ θραύσει ἐνώπιον αὐτοῦ, τοῦ ἀποστρέψαι ἀπὸ θυμοῦ ὀργῆς αὐτοῦ, τοῦ μὴ ἐξολοθρεῦσαι.

\vs{24}Καὶ ἐξουδένωσαν γῆν ἐπιθυμητὴν, καὶ οὐκ ἐπίστευσαν τῷ λόγῳ αὐτοῦ.
\vs{25}Καὶ ἐγόγγυσαν ἐν τοῖς σκηνώμασιν αὐτῶν, οὐκ εἰσήκουσαν τῆς φωνῆς Κυρίου.
\vs{26}Καὶ ἐπῇρε τὴν χεῖρα αὐτοῦ ἐπʼ αὐτοὺς, τοῦ καταβαλεῖν αὐτοὺς ἐν τῇ ἐρήμῳ,
\vs{27}καὶ τοῦ καταβαλεῖν τὸ σπέρμα αὐτῶν ἐν τοῖς ἔθνεσι, καὶ διασκορπίσαι αὐτοὺς ἐν ταῖς χώραις.

\vs{28}Καὶ ἐτελέσθησαν τῷ Βεελφεγὼρ, καὶ ἔφαγον θυσίας νεκρῶν.
\vs{29}Καὶ παρώξυναν αὐτὸν ἐν τοῖς ἐπιτηδεύμασιν αὐτῶν, καὶ ἐπληθύνθη ἐν αὐτοῖς ἡ πτῶσις.
\vs{30}Καὶ ἔστη Φινεὲς καὶ ἐξιλάσατο, καὶ ἐκόπασεν ἡ θραῦσις.
\vs{31}Καὶ ἐλογίσθη αὐτῷ εἰς δικαιοσύνην, εἰς γενεὰν καὶ γενεὰν ἕως τοῦ αἰῶνος.

\vs{32}Καὶ παρώργισαν αὐτὸν ἐπὶ ὕδατος ἀντιλογίας, καὶ ἐκακώθη Μωυσῆς διʼ αὐτούς·
\vs{33}Ὅτι παρεπίκραναν τὸ πνεῦμα αὐτοῦ, καὶ διέστειλεν ἐν τοῖς χείλεσιν αὐτοῦ.

\vs{34}Οὐκ ἐξωλόθρευσαν τὰ ἔθνη ἃ εἶπε Κύριος αὐτοῖς.
\vs{35}Καὶ ἐμίγησαν ἐν τοῖς ἔθνεσι, καὶ ἔμαθον τὰ ἔργα αὐτῶν.
\vs{36}Καὶ ἐδούλευσαν τοῖς γλυπτοῖς αὐτῶν, καὶ ἐγενήθη αὐτοῖς εἰς σκάνδαλον.
\vs{37}Καὶ ἔθυσαν τοὺς υἱοὺς αὐτῶν καὶ τὰς θυγατέρας αὐτῶν τοῖς δαιμονίοις,
\vs{38}καὶ ἐξέχεαν αἷμα ἀθῶον, αἷμα υἱῶν αὐτῶν καὶ θυγατέρων, ὧν ἔθυσαν τοῖς γλυπτοῖς Χαναάν· καὶ ἐφονοκτονήθη ἡ γῆ ἐν τοῖς αἵμασι,
\vs{39}καὶ ἐμιάνθη ἐν τοῖς ἔργοις αὐτῶν· καὶ ἐπόρνευσαν ἐν τοῖς ἐπιτηδεύμασιν αὐτῶν.

\vs{40}Καὶ ὠργίσθη θυμῷ Κύριος ἐπὶ τὸν λαὸν αὐτοῦ, καὶ ἐβδελύξατο τὴν κληρονομίαν αὐτοῦ.
\vs{41}Καὶ παρέδωκεν αὐτοὺς εἰς χεῖρας ἐχθρῶν, καὶ ἐκυρίευσαν αὐτῶν οἱ μισοῦντες αὐτούς.
\vs{42}Καὶ ἔθλιψαν αὐτοὺς οἱ ἐχθροὶ αὐτῶν, καὶ ἐταπεινώθησαν ὑπὸ τὰς χεῖρας αὐτῶν.
\vs{43}Πλεονάκις ἐῤῥύσατο αὐτοὺς, αὐτοὶ δὲ παρεπίκραναν αὐτὸν ἐν τῇ βουλῇ αὐτῶν· καὶ ἐταπεινώθησαν ἐν ταῖς ἀνομίαις αὐτῶν.
\vs{44}Καὶ εἶδε Κύριος ἐν τῷ θλίβεσθαι αὐτοὺς, ἐν τῷ αὐτὸν εἰσακοῦσαι τῆς δεήσεως αὐτῶν.
\vs{45}Καὶ ἐμνήσθη τῆς διαθήκης αὐτοῦ, καὶ μετεμελήθη κατὰ τὸ πλῆθος τοῦ ἐλέους αὐτοῦ.
\vs{46}Καὶ ἔδωκεν αὐτοὺς εἰς οἰκτιρμοὺς ἐναντίον πάντων τῶν αἰχμαλωτευσάντων αὐτούς.

\vs{47}Σῶσον ἡμᾶς Κύριε ὁ Θεὸς ἡμῶν, καὶ ἐπισυνάγαγε ἡμᾶς ἐκ τῶν ἐθνῶν, τοῦ ἐξομολογήσασθαι τῷ ὀνόματί σου τῷ ἁγίῳ, τοῦ ἐγκαυχᾶσθαι ἐν τῇ αἰνέσει σου.
\vs{48}Εὐλογητὸς Κύριος ὁ Θεὸς Ἰσραὴλ, ἀπὸ τοῦ αἰῶνος καὶ ἕως τοῦ αἰῶνος· καὶ ἐρεῖ πᾶς ὁ λαὸς, γένοιτο, γένοιτο.

\begin{psalmheading}{\ch{106}{107} Ἀλληλούϊα.}
\end{psalmheading}
Ἐξομολογεῖσθε τῷ Κυρίῳ, ὅτι χρηστὸς, ὅτι εἰς τὸν αἰῶνα τὸ ἔλεος αὐτοῦ.
\vs{2}Εἰπάτωσαν οἱ λελυτρωμένοι ὑπὸ Κυρίου, οὓς ἐλυτρώσατο ἐκ χειρὸς ἐχθροῦ,
\vs{3}καὶ ἐκ τῶν χωρῶν συνήγαγεν αὐτούς· ἀπὸ ἀνατολῶν, καὶ δυσμῶν, καὶ βοῤῥᾶ, καὶ θαλάσσης.

\vs{4}Ἐπλανήθησαν ἐν τῇ ἐρήμῳ ἐν ἀνύδρῳ· ὁδὸν πόλεως κατοικητηρίου οὐχ εὗρον·
\vs{5}Πεινῶντες καὶ διψῶντες, ἡ ψυχὴ αὐτῶν ἐν αὐτοῖς ἐξέλιπε.
\vs{6}Καὶ ἐκέκραξαν πρὸς Κύριον ἐν τῷ θλίβεσθαι αὐτοὺς, καὶ ἐκ τῶν ἀναγκῶν αὐτῶν ἐῤῥύσατο αὐτούς·
\vs{7}Καὶ ὡδήγησεν αὐτοὺς εἰς ὁδὸν εὐθεῖαν, τοῦ πορευθῆναι εἰς πόλιν κατοικητηρίου.

\vs{8}Ἐξομολογησάσθωσαν τῷ Κυρίῳ τὰ ἐλέη αὐτοῦ, καὶ τὰ θαυμάσια αὐτοῦ τοῖς υἱοῖς τῶν ἀνθρώπων.
\vs{9}Ὅτι ἐχόρτασε ψυχὴν κενὴν, καὶ πεινῶσαν ἐνέπλησεν ἀγαθῶν.
\vs{10}Καθημένους ἐν σκότει καὶ σκιᾷ θανάτου, πεπεδημένους ἐν πτωχείᾳ καὶ σιδήρῳ·
\vs{11}Ὅτι παρεπίκραναν τὰ λόγια τοῦ Θεοῦ, καὶ τὴν βουλὴν τοῦ ὑψίστου παρώξυναν·
\vs{12}Καὶ ἐταπεινώθη ἐν κόποις ἡ καρδία αὐτῶν, ἠσθένησαν καὶ οὐκ ἦν ὁ βοηθῶν.
\vs{13}Καὶ ἐκέκραξαν πρὸς Κύριον ἐν τῷ θλίβεσθαι αὐτοὺς, καὶ ἐκ τῶν ἀναγκῶν αὐτῶν ἔσωσεν αὐτούς.
\vs{14}Καὶ ἐξήγαγεν αὐτοὺς ἐκ σκότους καὶ σκιᾶς θανάτου, καὶ τοὺς δεσμοὺς αὐτῶν διέῤῥηξεν.

\vs{15}Ἐξομολογησάσθωσαν τῷ Κυρίῳ τὰ ἐλέη αὐτοῦ, καὶ τὰ θαυμάσια αὐτοῦ τοῖς υἱοῖς τῶν ἀνθρώπων.
\vs{16}Ὅτι συνέτριψε πύλας χαλκᾶς, καὶ μοχλοὺς σιδηροὺς συνέθλασεν.

\vs{17}Ἀντελάβετο αὐτῶν ἐξ ὁδοῦ ἀνομίας αὐτῶν, διὰ γὰρ τὰς ἀνομίας αὐτῶν ἐταπεινώθησαν.
\vs{18}Πᾶν βρῶμα ἐβδελύξατο ἡ ψυχὴ αὐτῶν, καὶ ἤγγισαν ἕως τῶν πυλῶν τοῦ θανάτου.
\vs{19}Καὶ ἐκέκραξαν πρὸς Κύριον ἐν τῷ θλίβεσθαι αὐτοὺς, καὶ ἐκ τῶν ἀναγκῶν αὐτῶν ἔσωσεν αὐτούς.
\vs{20}Ἀπέστειλε τὸν λόγον αὐτοῦ, καὶ ἰάσατο αὐτοὺς, καὶ ἐῤῥύσατο αὐτοὺς ἐκ τῶν διαφθορῶν αὐτῶν.

\vs{21}Ἐξομολογησάσθωσαν τῷ Κυρίῳ τὰ ἐλέη αὐτοῦ, καὶ τὰ θαυμάσια αὐτοῦ τοῖς υἱοῖς τῶν ἀνθρώπων·
\vs{22}Καὶ θυσάτωσαν αὐτῷ θυσίαν αἰνέσεως, καὶ ἐξαγγειλάτωσαν τὰ ἔργα αὐτοῦ ἐν ἀγαλλιάσει.

\vs{23}Οἱ καταβαίνοντες εἰς θάλασσαν ἐν πλοίοις, ποιοῦντες ἐργασίαν ἐν ὕδασι πολλοῖς,
\vs{24}αὐτοὶ εἶδον τὰ ἔργα Κυρίου, καὶ τὰ θαυμάσια αὐτοῦ ἐν τῷ βυθῷ.
\vs{25}Εἶπε, καὶ ἔστη πνεῦμα καταιγίδος, καὶ ὑψώθη τὰ κύματα αὐτῆς.
\vs{26}Ἀναβαίνουσιν ἕως τῶν οὐρανῶν, καὶ καταβαίνουσιν ἕως τῶν ἀβύσσων· ἡ ψυχὴ αὐτῶν ἐν κακοῖς ἐτήκετο,
\vs{27}ἐταράχθησαν, ἐσαλεύθησαν ὡς ὁ μεθύων, καὶ πᾶσα ἡ σοφία αὐτῶν κατεπόθη.
\vs{28}Καὶ ἐκέκραξαν πρὸς Κύριον ἐν τῷ θλίβεσθαι αὐτοὺς, καὶ ἐκ τῶν ἀναγκῶν αὐτῶν ἐξήγαγεν αὐτούς.
\vs{29}Καὶ ἐπέταξε τῇ καταιγίδι, καὶ ἔστη εἰς αὖραν, καὶ ἐσίγησαν τὰ κύματα αὐτῆς.
\vs{30}Καὶ εὐφράνθησαν, ὅτι ἡσύχασαν, καὶ ὡδήγησεν αὐτοὺς ἐπὶ λιμένα θελήματος αὐτῶν.

\vs{31}Ἐξομολογησάσθωσαν τῷ Κυρίῳ τὰ ἐλέη αὐτοῦ, καὶ τὰ θαυμάσια αὐτοῦ τοῖς υἱοῖς τῶν ἀνθρώπων.
\vs{32}Ὑψωσάτωσαν αὐτὸν ἐν ἐκκλησίᾳ λαοῦ, καὶ ἐν καθέδρᾳ πρεσβυτέρων αἰνεσάτωσαν αὐτόν.

\vs{33}Ἔθετο ποταμοὺς εἰς ἐρήμον, καὶ διεξόδους ὑδάτων εἰς δίψαν·
\vs{34}Γῆν καρποφόρον εἰς ἅλμην, ἀπὸ κακίας τῶν κατοικούντων ἐν αὐτῇ.
\vs{35}Ἔθετο ἔρημον εἰς λίμνας ὑδάτων, καὶ γῆν ἄνυδρον εἰς διεξόδους ὑδάτων.
\vs{36}Καὶ κατῴκισεν ἐκεῖ πεινῶντας, καὶ συνεστήσαντο πόλεις κατοικεσίας·
\vs{37}Καὶ ἔσπειραν ἀγροὺς, καὶ ἐφύτευσαν ἀμπελῶνας, καὶ ἐποίησαν καρπὸν γεννήματος.
\vs{38}Καὶ εὐλόγησεν αὐτοὺς, καὶ ἐπληθύνθησαν σφόδρα, καὶ τὰ κτήνη αὐτῶν οὐκ ἐσμίκρυνε.
\vs{39}Καὶ ὠλιγώθησαν καὶ ἐκακώθησαν ἀπὸ θλίψεως κακῶν καὶ ὀδύνης.
\vs{40}Ἐξεχύθη ἐξουδένωσις ἐπʼ ἄρχοντας αὐτῶν, καὶ ἐπλάνησεν αὐτοὺς ἐν ἀβάτῳ καὶ οὐχ ὁδῷ.
\vs{41}Καὶ ἐβοήθησε πένητι ἐκ πτωχείας, καὶ ἔθετο ὡς πρόβατα πατριάς·
\vs{42}Ὄψονται εὐθεῖς καὶ εὐφρανθήσονται, καὶ πᾶσα ἀνομία ἐμφράξει τὸ στόμα αὐτῆς.
\vs{43}Τίς σοφὸς καὶ φυλάξει ταῦτα, καὶ συνήσει τὰ ἐλέη τοῦ Κυρίου;

\begin{psalmheading}{\ch{107}{108} Ὠδὴ ψαλμοῦ τῷ Δαυίδ.}
\end{psalmheading}
\vs{2}Ἑτοίμη ἡ καρδία μου ὁ Θεὸς, ἑτοίμη ἡ καρδία μου, ᾄσομαι καὶ ψαλῶ ἐν τῇ δόξῃ μου.
\vs{3}Ἐξεγέρθητι ψαλτήριον καὶ κιθάρα, ἐξεγερθήσομαι ὄρθρου.
\vs{4}Ἐξομολογήσομαί σοι ἐν λαοῖς Κύριε, ψαλῶ σοι ἐν ἔθνεσιν.
\vs{5}Ὅτι μέγα ἐπάνω τῶν οὐρανῶν τὸ ἔλεός σου, καὶ ἕως τῶν νεφελῶν ἡ ἀλήθειά σου.
\vs{6}Ὑψώθητι ἐπὶ τοὺς οὐρανοὺς, ὁ Θεὸς, καὶ ἐπὶ πᾶσαν τὴν γῆν ἡ δόξα σου.
\vs{7}Ὅπως ἂν ῥυσθῶσιν οἱ ἀγαπητοί σου, σῶσον τῇ δεξιᾷ σου, καὶ ἐπάκουσόν μου.
\vs{8}Ὁ Θεὸς ἐλάλησεν ἐν τῷ ἁγίῳ αὐτοῦ, ὑψωθήσομαι καὶ διαμεριῶ Σίκιμα, καὶ τὴν κοιλάδα τῶν σκηνῶν διαμετρήσω.
\vs{9}Ἐμός ἐστι Γαλαὰδ, καὶ ἐμός ἐστι Μανασσῆς, καὶ Ἐφραίμ ἀντίληψις τῆς κεφαλῆς μου· Ἰούδας βασιλεύς μου,
\vs{10}Μωὰβ λέβης τῆς ἐλπίδος μου· ἐπὶ τὴν Ἰδουμαίαν ἐπιβαλῶ τὸ ὑπόδημά μου, ἐμοὶ ἀλλόφυλοι ὑπετάγησαν.

\vs{11}Τίς ἀπάξει με εἰς πόλιν περιοχῆς; ἢ τίς ὁδηγήσει με ἕως τῆς Ἰδουμαίας;
\vs{12}Οὐχὶ σὺ ὁ Θεὸς ὁ ἀπωσάμενος ἡμᾶς; καὶ οὐκ ἐξελεύσῃ ὁ Θεὸς ἐν ταῖς δυνάμεσιν ἡμῶν;
\vs{13}Δὸς ἡμῖν βοήθειαν ἐκ θλίψεως, καὶ ματαία σωτηρία ἀνθρώπου.
\vs{14}Ἐν τῷ Θεῷ ποιήσομεν δύναμιν, καὶ αὐτὸς ἐξουδενώσει τοὺς ἐχθροὺς ἡμῶν.

\begin{psalmheading}{\ch{108}{109} Εἰς τὸ τέλος, ψαλμὸς τῷ Δαυίδ.}
\end{psalmheading}
Ὁ Θεὸς τὴν αἴνεσίν μου μὴ παρασιωπήσῃς,
\vs{2}ὅτι στόμα ἁμαρτωλοῦ καὶ στόμα δολίου ἐπʼ ἐμὲ ἠνοίχθη· ἐλάλησαν κατʼ ἐμοῦ γλώσσῃ δολίᾳ,
\vs{3}καὶ λόγοις μίσους ἐκύκλωσάν με, καὶ ἐπολέμησάν με δωρέαν.
\vs{4}Ἀντὶ τοῦ ἀγαπᾷν με, ἐνδιέβαλλόν με, ἐγὼ δὲ προσηυχόμην.
\vs{5}Καὶ ἔθεντο κατʼ ἐμοῦ κακὰ ἀντὶ ἀγαθῶν, καὶ μῖσος ἀντὶ τῆς ἀγαπήσεώς μου.

\vs{6}Κατάστησον ἐπʼ αὐτὸν ἁμαρτωλὸν, καὶ διάβολος στήτω ἐκ δεξιῶν αὐτοῦ.
\vs{7}Ἐν τῷ κρίνεσθαι αὐτὸν, ἐξέλθοι καταδεδικασμένος, καὶ ἡ προσευχὴ αὐτοῦ γενέσθω εἰς ἁμαρτίαν.
\vs{8}Γενηθήτωσαν αἱ ἡμέραι αὐτοῦ ὀλίγαι, καὶ τὴν ἐπισκοπὴν αὐτοῦ λάβοι ἕτερος.
\vs{9}Γενηθήτωσαν οἱ υἱοὶ αὐτοῦ ὀρφανοὶ, καὶ ἡ γυνὴ αὐτοῦ χήρα.
\vs{10}Σαλευόμενοι μεταναστήτωσαν οἱ υἱοὶ αὐτοῦ, καὶ ἐπαιτησάτωσαν, ἐκβληθήτωσαν ἐκ τῶν οἰκοπέδων αὐτῶν.
\vs{11}Ἐξερευνησάτω δανειστὴς πάντα ὅσα ὑπάρχει αὐτῷ, καὶ διαρπασάτωσαν ἀλλότριοι τοὺς πόνους αὐτοῦ.
\vs{12}Μὴ ὑπαρξάτω αὐτῷ ἀντιλήμπτωρ, μηδὲ γενηθήτω οἰκτίρμων τοῖς ὀρφανοῖς αὐτοῦ.
\vs{13}Γενηθήτω τὰ τέκνα αὐτοῦ εἰς ἐξολόθρευσιν, ἐν γενεᾷ μιᾷ ἐξαλειφθείη τὸ ὄνομα αὐτοῦ.
\vs{14}Ἀναμνησθείη ἡ ἀνομία τῶν πατέρων αὐτοῦ ἔναντι Κυρίου, καὶ ἡ ἁμαρτία τῆς μητρὸς αὐτοῦ μὴ ἐξαλειφθείη.
\vs{15}Γενηθήτωσαν ἐναντίον Κυρίου διαπαντὸς, καὶ ἐξολοθρευθείη ἐκ γῆς τὸ μνημόσυνον αὐτῶν·

\vs{16}Ἀνθʼ ὧν οὐκ ἐμνήσθη ποιῆσαι ἔλεος, καὶ κατεδίωξεν ἄνθρωπον πένητα καὶ πτωχὸν, καὶ κατανενυγμένον τῇ καρδίᾳ τοῦ θανατῶσαι.
\vs{17}Καὶ ἠγάπησε κατάραν, καὶ ἥξει αὐτῷ, καὶ οὐκ ἠθέλησεν εὐλογίαν, καὶ μακρυνθήσεται ἀπʼ αὐτοῦ.
\vs{18}Καὶ ἐνεδύσατο κατάραν ὡς ἱμάτιον, καὶ εἰσῆλθεν ὡσεὶ ὕδωρ εἰς τὰ ἔγκατα αὐτοῦ, καὶ ὡσεὶ ἔλαιον ἐν τοῖς ὀστέοις αὐτοῦ.
\vs{19}Γενηθήτω αὐτῷ ὡς ἱμάτιον ὃ περιβάλλεται, καὶ ὡσεὶ ζώνη ἣν διαπαντὸς περιζώννυται.
\vs{20}Τοῦτο τὸ ἔργον τῶν ἐνδιαβαλλόντων με παρὰ Κυρίου, καὶ τῶν λαλούντων πονηρὰ κατὰ τῆς ψυχῆς μου.

\vs{21}Καὶ σὺ Κύριε Κύριε ποίησον μετʼ ἐμοῦ ἕνεκεν τοῦ ὀνόματός σου, ὅτι χρηστὸν τὸ ἔλεός σου.
\vs{22}Ῥῦσαί με ὅτι πτωχὸς καὶ πένης εἰμὶ ἐγὼ, καὶ ἡ καρδία μου τετάρακται ἐντός μου.
\vs{23}Ὡσεὶ σκιὰ ἐν τῷ ἐκκλῖναι αὐτὴν ἀντανῃρέθην, ἐξετινάχθην ὡσεὶ ἀκρίδες.
\vs{24}Τὰ γόνατά μου ἠσθένησαν ἀπὸ νηστείας, καὶ ἡ σάρξ μου ἠλλοιώθη διʼ ἔλαιον.
\vs{25}Κᾀγὼ ἐγενήθην ὄνειδος αὐτοῖς· εἴδοσάν με, ἐσάλευσαν κεφαλὰς αὐτῶν.

\vs{26}Βοήθησόν μοι Κύριε ὁ Θεός μου, καὶ σῶσόν με κατὰ τὸ ἔλεός σου.
\vs{27}Καὶ γνώτωσαν ὅτι ἡ χείρ σου αὕτη, καὶ σὺ Κύριε ἐποίησας αὐτήν.
\vs{28}Καταράσονται αὐτοὶ, καὶ σὺ εὐλογήσεις· οἱ ἐπανιστάμενοί μοι αἰσχυνθήτωσαν, ὁ δὲ δοῦλός σου εὐφρανθήσεται.
\vs{29}Ἐνδυσάσθωσαν οἱ ἐνδιαβάλλοντές με ἐντροπήν· καὶ περιβαλέσθωσαν ὡς διπλοΐδα αἰσχύνην αὐτῶν.
\vs{30}Ἐξομολογήσομαι τῷ Κυρίῳ σφόδρα ἐν τῷ στόματί μου, καὶ ἐν μέσῳ πολλῶν αἰνέσω αὐτόν·
\vs{31}Ὅτι παρέστη ἐκ δεξιῶν πένητος, τοῦ σῶσαι ἐκ τῶν καταδιωκόντων τὴν ψυχήν μου.

\begin{psalmheading}{\ch{109}{110} Ψαλμὸς τῷ Δαυίδ.}
\end{psalmheading}
Εἶπεν ὁ Κύριος τῷ Κυρίῳ μου, κάθου ἐκ δεξιῶν μου, ἕως ἂν θῶ τοὺς ἐχθρούς σου ὑποπόδιον τῶν ποδῶν σου.
\vs{2}Ῥάβδον δυνάμεως ἐξαποστελεῖ σοι Κύριος ἐκ Σιὼν, κατακυρίευε ἐν μέσῳ τῶν ἐχθρῶν σου.
\vs{3}Μετὰ σοῦ ἡ ἀρχὴ ἐν ἡμέρᾳ τῆς δυνάμεώς σου, ἐν ταῖς λαμπρότησι τῶν ἁγίων σου· ἐκ γαστρὸς πρὸ Ἑωσφόρου ἐγέννησά σε.
\vs{4}Ὤμοσε Κύριος καὶ οὐ μεταμεληθήσεται, σὺ ἱερεὺς εἰς τὸν αἰῶνα, κατὰ τὴν τάξιν Μελχισεδέκ.
\vs{5}Κύριος ἐκ δεξιῶν σου συνέθλασεν ἐν ἡμέρᾳ ὀργῆς αὐτοῦ βασιλεῖς.
\vs{6}Κρινεῖ ἐν τοῖς ἔθνεσι, πληρώσει πτώματα, συνθλάσει κεφαλὰς ἐπὶ γῆν πολλῶν.
\vs{7}Ἐκ χειμάῤῥου ἐν ὁδῷ πίεται, διὰ τοῦτο ὑψώσει κεφαλήν.

\begin{psalmheading}{\ch{110}{111} Ἀλληλούϊα.}
\end{psalmheading}
Ἐξομολογήσομαι σοι Κύριε ἐν ὅλῃ καρδίᾳ μου, ἐν βουλῇ εὐθέων καὶ συναγωγῇ.
\vs{2}Μεγάλα τὰ ἔργα Κυρίου, ἐξεζητημένα εἰς πάντα τὰ θελήματα αὐτοῦ.
\vs{3}Ἐξομολόγησις καὶ μεγαλοπρέπεια τὸ ἔργον αὐτοῦ, καὶ ἡ δικαιοσύνη αὐτοῦ μένει εἰς τὸν αἰῶνα τοῦ αἰῶνος.
\vs{4}Μνείαν ἐποιήσατο τῶν θαυμασίων αὐτοῦ, ἐλεήμων καὶ οἰκτίρμων ὁ Κύριος.
\vs{5}Τροφὴν ἔδωκε τοῖς φοβουμένοις αὐτόν· μνησθήσεται εἰς τὸν αἰῶνα διαθήκης αὐτοῦ.
\vs{6}Ἰσχὺν ἔργων αὐτοῦ ἀνήγγειλε τῷ λαῷ αὐτοῦ, τοῦ δοῦναι αὐτοῖς κληρονομίαν ἐθνῶν.
\vs{7}Ἔργα χειρῶν αὐτοῦ, ἀλήθεια καὶ κρίσις· πισταὶ πᾶσαι αἱ ἐντολαὶ αὐτοῦ,
\vs{8}ἐστηριγμέναι εἰς τὸν αἰῶνα τοῦ αἰῶνος, πεποιημέναι ἐν ἀληθείᾳ καὶ εὐθύτητι.
\vs{9}Λύτρωσιν ἀπέστειλε τῷ λαῷ αὐτοῦ· ἐνετείλατο εἰς τὸν αἰῶνα διαθήκην αὐτοῦ· ἅγιον καὶ φοβερὸν τὸ ὄνομα αὐτοῦ.
\vs{10}Ἀρχὴ σοφίας φόβος Κυρίου, σύνεσις δὲ ἀγαθὴ πᾶσι τοῖς ποιοῦσιν αὐτήν· ἡ αἴνεσις αὐτοῦ μένει εἰς τὸν αἰῶνα τοῦ αἰῶνος.

\begin{psalmheading}{\ch{111}{112} Ἀλληλούϊα.}
\end{psalmheading}
Μακάριος ἀνὴρ ὁ φοβούμενος τὸν Κύριον, ἐν ταῖς ἐντολαῖς αὐτοῦ θελήσει σφόδρα.
\vs{2}Δυνατὸν ἐν τῇ γῇ ἔσται τὸ σπέρμα αὐτοῦ, γενεὰ εὐθέων εὐλογηθήσεται·
\vs{3}Δόξα καὶ πλοῦτος ἐν τῷ οἴκῳ αὐτοῦ, καὶ ἡ δικαιοσύνη αὐτοῦ μένει εἰς τὸν αἰῶνα τοῦ αἰῶνος.
\vs{4}Ἐξανέτειλεν ἐν σκότει φῶς τοῖς εὐθέσιν· ἐλεήμων καὶ οἰκτίρμων καὶ δίκαιος.
\vs{5}Χρηστὸς ἀνὴρ ὁ οἰκτείρων καὶ κιχρῶν, οἰκονομήσει τοὺς λόγους αὐτοῦ ἐν κρίσει,
\vs{6}ὅτι εἰς τὸν αἰῶνα οὐ σαλευθήσεται· εἰς μνημόσυνον αἰώνιον ἔσται δίκαιος.
\vs{7}Ἀπὸ ἀκοῆς πονηρᾶς οὐ φοβηθήσεται· ἑτοίμη ἡ καρδία αὐτοῦ ἐλπίζειν ἐπὶ Κύριον·
\vs{8}Ἐστήρικται ἡ καρδία αὐτοῦ, οὐ φοβηθῇ, ἕως οὗ ἐπίδῃ ἐπὶ τοὺς ἐχθροὺς αὐτοῦ.
\vs{9}Ἐσκόρπισεν, ἔδωκε τοῖς πένησιν, ἡ δικαιοσύνη αὐτοῦ μένει εἰς τὸν αἰῶνα τοῦ αἰῶνος· τὸ κέρας αὐτοῦ ὑψωθήσεται ἐν δόξῃ.
\vs{10}Ἁμαρτωλὸς ὄψεται καὶ ὀργισθήσεται, τοὺς ὀδόντας αὐτοῦ βρύξει καὶ τακήσεται· ἐπιθυμία ἁμαρτωλοῦ ἀπολεῖται.

\begin{psalmheading}{\ch{112}{113} Ἀλληλούϊα.}
\end{psalmheading}
Αἰνεῖτε παῖδες Κύριον, αἰνεῖτε τὸ ὄνομα Κυρίου.
\vs{2}Εἴη τὸ ὄνομα Κυρίου εὐλογημένον ἀπὸ τοῦ νῦν καὶ ἕως τοῦ αἰῶνος.
\vs{3}Ἀπὸ ἀνατολῶν ἡλίου μέχρι δυσμῶν, αἰνετὸν τὸ ὄνομα Κυρίου.
\vs{4}Ὑψηλὸς ἐπὶ πάντα τὰ ἔθνη ὁ Κύριος, ἐπὶ τοὺς οὐρανοὺς ἡ δόξα αὐτοῦ.

\vs{5}Τίς ὡς Κύριος ὁ Θεὸς ἡμῶν; ὁ ἐν ὑψηλοῖς κατοικῶν,
\vs{6}καὶ τὰ ταπεινὰ ἐφορῶν ἐν τῷ οὐρανῷ, καὶ ἐν τῇ γῇ·
\vs{7}ὁ ἐγείρων ἀπὸ γῆς πτωχὸν, καὶ ἀπὸ κοπρίας ἀνυψῶν πένητα,
\vs{8}τοῦ καθίσαι αὐτὸν μετὰ ἀρχόντων, μετὰ ἀρχόντων λαοῦ αὐτοῦ·
\vs{9}ὁ κατοικίζων στείραν ἐν οἴκῳ, μητέρα ἐπὶ τέκνοις εὐφραινομένην.

\begin{psalmheading}{\ch{113}{114-115} Ἀλληλούϊα.}
\end{psalmheading}
Ἐν ἐξόδῳ Ἰσραὴλ ἐξ Αἰγύπτου, οἴκου Ἰακὼβ ἐκ λαοῦ βαρβάρου,
\vs{2}ἐγενήθη Ἰουδαία ἁγίασμα αὐτοῦ, Ἰσραὴλ ἐξουσία αὐτοῦ.
\vs{3}Ἡ θάλασσα εἶδε καὶ ἔφυγεν, ὁ Ἰορδάνης ἐστράφη εἰς τὰ ὀπίσω.
\vs{4}Τὰ ὄρη ἐσκίρτησαν ὡσεὶ κριοὶ, καὶ οἱ βουνοὶ ὡς ἀρνία προβάτων.

\vs{5}Τί σοι ἐστὶ θάλασσα ὅτι ἔφυγες; καὶ σὺ Ἰορδάνη ὅτι ἐστράφης εἰς τὰ ὀπίσω;
\vs{6}Τὰ ὄρη ὅτι ἐσκίρτησατε ὡσεὶ κριοί; καὶ οἱ βουνοὶ ὡς ἀρνία προβάτων;
\vs{7}Ἀπὸ προσώπου Κυρίου ἐσαλεύθη ἡ γῆ, ἀπὸ προσώπου τοῦ Θεοῦ Ἰακὼβ,
\vs{8}τοῦ στρέψαντος τὴν πέτραν εἰς λίμνας ὑδάτων, καὶ τὴν ἀκρότομον εἰς πηγὰς ὑδάτων.
\vs{9}Μὴ ἡμῖν Κύριε, μὴ ἡμῖν, ἀλλʼ ἢ τῷ ὀνόματί σου δὸς δόξαν, ἐπὶ τῷ ἐλέει σου καὶ τῇ ἀληθείᾳ σου·
\vs{10}μή ποτε εἴπωσι τὰ ἔθνη, ποῦ ἐστιν ὁ Θεὸς αὐτῶν;
\vs{11}Ὁ δὲ Θεὸς ἡμῶν ἐν τῷ οὐρανῷ καὶ ἐν τῇ γῇ, πάντα ὅσα ἠθέλησεν ἐποίησε.

\vs{12}Τὰ εἴδωλα τῶν ἐθνῶν, ἀργύριον καὶ χρυσίον, ἔργα χειρῶν ἀνθρώπων.
\vs{13}Στόμα ἔχουσι καὶ οὐ λαλήσουσιν, ὀφθαλμοὺς ἔχουσι καὶ οὐκ ὄψονται·
\vs{14}Ὦτα ἔχουσι καὶ οὐκ ἀκούσονται, ῥῖνας ἔχουσι καὶ οὐκ ὀσφρανθήσονται·
\vs{15}Χεῖρας ἔχουσι καὶ οὐ ψηλαφήσουσι, πόδας ἔχουσι καὶ οὐ περιπατήσουσιν, οὐ φωνήσουσιν ἐν τῷ λάρυγγι αὐτῶν.
\vs{16}ὅμοιοι αὐτοῖς γένοιντο οἱ ποιοῦντες αὐτὰ, καὶ πάντες οἱ πεποιθότες ἐπʼ αὐτοῖς.

\vs{17}Οἶκος Ἰσραὴλ ἔλπισεν ἐπὶ Κύριον, βοηθὸς καὶ ὑπερασπιστὴς αὐτῶν ἐστιν.
\vs{18}Οἶκος Ἀαρὼν ἤλπισεν ἐπὶ Κύριον, βοηθὸς καὶ ὑπερασπιστὴς αὐτῶν ἐστιν.
\vs{19}Οἱ φοβούμενοι τὸν Κύριον ἤλπισαν ἐπὶ Κύριον, βοηθὸς καὶ ὑπερασπιστὴς αὐτῶν ἐστι.

\vs{20}Κύριος μνησθεὶς ἡμῶν εὐλόγησεν ἡμᾶς, εὐλόγησε τὸν οἶκον Ἰσραὴλ, εὐλόγησε τὸν οἶκον Ἀαρών·
\vs{21}Εὐλόγησε τοὺς φοβουμένους τὸν Κύριον, τοὺς μικροὺς μετὰ τῶν μεγάλων.
\vs{22}Προσθείη Κύριος ἐφʼ ὑμᾶς, ἐφʼ ὑμᾶς καὶ ἐπὶ τοὺς υἱοὺς ὑμῶν.
\vs{23}Εὐλογημένοι ὑμεῖς τῷ Κυρίῳ, τῷ ποιήσαντι τὸν οὐρανὸν καὶ τὴν γῆν.

\vs{24}Ὁ οὐρανὸς τοῦ οὐρανοῦ τῷ Κυρίῳ, τὴν δὲ γῆν ἔδωκε τοῖς υἱοῖς τῶν ἀνθρώπων.
\vs{25}Οὐχ οἱ νεκροὶ αἰνέσουσί σε Κύριε, οὐδὲ πάντες οἱ καταβαίνοντες εἰς ᾅδου·
\vs{26}Ἀλλʼ ἡμεῖς οἱ ζῶντες εὐλογήσομεν τὸν Κύριον, ἀπὸ τοῦ νῦν καὶ ἕως τοῦ αἰῶνος.

\begin{psalmheading}{\ch{114}{116:1-9} Ἀλληλούϊα.}
\end{psalmheading}
Ἠγαπήσα, ὅτι εἰσακούσεται Κύριος τῆς φωνῆς τῆς δεήσεώς μου·
\vs{2}Ὅτι ἔκλινε τὸ οὖς αὐτοῦ ἐμοὶ, καὶ ἐν ταῖς ἡμέραις μου ἐπικαλέσομαι.
\vs{3}Περιέσχον με ὠδῖνες θανάτου, κίνδυνοι ᾅδου εὕροσάν με· θλίψιν καὶ ὀδύνην εὗρον,
\vs{4}καὶ τὸ ὄνομα Κυρίου ἐπεκαλεσάμην, ὦ Κύριε ῥῦσαι τὴν ψυχήν μου.

\vs{5}Ἐλεήμων ὁ Κύριος καὶ δίκαιος, καὶ ὁ Θεὸς ἡμῶν ἐλεεῖ.
\vs{6}Φυλάσσων τὰ νήπια ὁ Κύριος, ἐταπεινώθην καὶ ἔσωσέ με.

\vs{7}Ἐπίστρεψον ψυχή μου εἰς τὴν ἀνάπαυσίν σου, ὅτι Κύριος εὐηργέτησέ σε.
\vs{8}Ὅτι ἐξείλετο τὴν ψυχήν μου ἐκ θανάτου, τοὺς ὀφθαλμούς μου ἀπὸ δακρύων, καὶ τοὺς πόδας μου ἀπὸ ὀλισθήματος.
\vs{9}Εὐαρεστήσω ἐνώπιον Κυρίου ἐν χώρᾳ ζώντων.

\begin{psalmheading}{\ch{115}{116:10-19} Ἀλληλούϊα.}
\end{psalmheading}
Ἐπιστεύσα, διὸ ἐλάλησα, ἐγὼ δὲ ἐταπεινώθην σφόδρα.
\vs{2}Ἐγὼ δὲ εἶπα ἐν τῇ ἐκστάσει μου, πᾶς ἄνθρωπος ψεύστης.

\vs{3}Τί ἀνταποδώσω τῷ Κυρίῳ περὶ πάντων ὧν ἀνταπέδωκέ μοι;
\vs{4}Ποτήριον σωτηρίου λήψομαι, καὶ τὸ ὄνομα Κυρίου ἐπικαλέσομαι·
\vs{4a}Τὰς εὐχάς μου τῷ Κυρίῳ ἀποδώσω, ἐναντίον παντὸς τοῦ λαοῦ αὐτοῦ.

\vs{6}Τίμιος ἐναντίον Κυρίου ὁ θάνατος τῶν ὁσίων αὐτοῦ.
\vs{7}Ὦ Κύριε ἐγὼ δοῦλος σὸς, ἐγὼ δοῦλος σὸς, καὶ υἱὸς τῆς παιδίσκης σου, διέῤῥηξας τοὺς δεσμούς μου.
\vs{8}Σοὶ θύσω θυσίαν αἰνέσεως, καὶ ἐν ὀνόματι Κυρίου ἐπικαλέσομαι·
\vs{9}Τὰς εὐχάς μου τῷ Κυρίῳ ἀποδώσω, ἐναντίον παντὸς τοῦ λαοῦ αὐτοῦ,
\vs{10}ἐν αὐλαῖς οἴκου Κυρίου, ἐν μέσῳ σου Ἱερουσαλήμ.

\begin{psalmheading}{\ch{116}{117} Ἀλληλούϊα.}
\end{psalmheading}
Αἰνεῖτε τὸν Κύριον πάντα τὰ ἔθνη, ἐπαινέσατε αὐτὸν πάντες οἱ λαοί.
\vs{2}Ὅτι ἐκραταιώθη τὸ ἔλεος αὐτοῦ ἐφʼ ἡμᾶς, καὶ ἡ ἀλήθεια τοῦ Κυρίου μένει εἰς τὸν αἰῶνα.

\begin{psalmheading}{\ch{117}{118} Ἀλληλούϊα.}
\end{psalmheading}
Ἐξομολογεῖσθε τῷ Κυρίῳ, ὅτι ἀγαθὸς, ὅτι εἰς τὸν αἰῶνα τὸ ἔλεος αὐτοῦ.
\vs{2}Εἰπάτω δὴ οἶκος Ἰσραὴλ, ὅτι ἀγαθὸς, ὅτι εἰς τὸν αἰῶνα τὸ ἔλεος αὐτοῦ.
\vs{3}Εἰπάτω δὴ οἶκος Ἀαρὼν, ὅτι ἀγαθὸς, ὅτι εἰς τὸν αἰῶνα τὸ ἔλεος αὐτοῦ.
\vs{4}Εἰπάτωσαν δὴ πάντες οἱ φοβούμενοι τὸν Κύριον, ὅτι ἀγαθὸς, ὅτι εἰς τὸν αἰῶνα τὸ ἔλεος αὐτοῦ.

\vs{5}Ἐκ θλίψεως ἐπεκαλεσάμην τὸν Κύριον, καὶ ἐπήκουσέ μου εἰς πλατυσμόν.
\vs{6}Κύριος ἐμοὶ βοηθὸς, καὶ οὐ φοβηθήσομαι τί ποιησει μοι ἄνθρωπος.
\vs{7}Κύριος ἐμοὶ βοηθὸς, κᾀγὼ ἐπόψομαι τοὺς ἐχθρούς μου.
\vs{8}Ἀγαθὸν πεποιθέναι ἐπὶ Κύριον, ἢ πεποιθέναι ἐπʼ ἄνθρωπον.
\vs{9}Ἀγαθὸν ἐλπίζειν ἐπὶ Κύριον, ἢ ἐλπίζειν ἐπʼ ἄρχουσι.

\vs{10}Πάντα τὰ ἔθνη ἐκύκλωσάν με, καὶ τῷ ὀνόματι Κυρίου ἠμυνάμην αὐτούς·
\vs{11}Κυκλώσαντες ἐκύκλωσάν με, καὶ τῷ ὀνόματι Κυρίου ἠμυνάμην αὐτούς·
\vs{12}Ἐκύκλωσάν με ὡσεὶ μέλισσαι κηρίον, καὶ ἐξεκαύθησαν ὡς πῦρ ἐν ἀκάνθαις, καὶ τῷ ὀνόματι Κυρίου ἠμυνάμην αὐτούς.
\vs{13}Ὠσθεὶς ἀνετράπην τοῦ πεσεῖν, καὶ ὁ Κύριος ἀντελάβετό μου.

\vs{14}Ἰσχύς μου καὶ ὕμνησίς μου ὁ Κύριος, καὶ ἐγένετό μοι εἰς σωτηρίαν.
\vs{15}Φωνὴ ἀγαλλιάσεως καὶ σωτηρίας ἐν σκηναῖς δικαίων· δεξιὰ Κυρίου ἐποίησε δύναμιν,
\vs{16}δεξιὰ Κυρίου ὕψωσέ με· δεξιὰ Κυρίου ἐποίησε δύναμιν.
\vs{17}Οὐκ ἀποθανοῦμαι, ἀλλὰ ζήσομαι, καὶ διηγήσομαι τὰ ἔργα Κυρίου.
\vs{18}Παιδεύων ἐπαίδευσέ με ὁ Κύριος, καὶ τῷ θανάτῳ οὐ παρέδωκέ με.

\vs{19}Ἀνοίξατέ μοι πύλας δικαιοσύνης, εἰσελθὼν ἐν αὐταῖς ἐξομολογήσομαι τῷ Κυρίῳ.
\vs{20}Αὕτη ἡ πύλη τοῦ Κυρίου, δίκαιοι εἰσελεύσονται ἐν αὐτῇ.
\vs{21}Ἐξομολογήσομαί σοι, ὅτι ἐπήκουσάς μου, καὶ ἐγένου μου εἰς σωτηρίαν.
\vs{22}Λίθον ὃν ἀπεδοκίμασαν οἱ οἰκοδομοῦντες, οὗτος ἐγενήθη εἰς κεφαλὴν γωνίας.
\vs{23}Παρὰ Κυρίου ἐγένετο αὕτη, καὶ ἔστι θαυμαστὴ ἐν ὀφθαλμοῖς ἡμῶν.
\vs{24}Αὕτη ἡ ἡμέρα ἣν ἐποίησεν ὁ Κύριος, ἀγαλλιασώμεθα καὶ εὐφρανθῶμεν ἐν αὐτῇ.
\vs{25}Ὦ Κύριε σῶσον δὴ, ὦ Κύριε εὐόδωσον δή.
\vs{26}Εὐλογημένος ὁ ἐρχόμενος ἐν ὀνόματι Κυρίου· εὐλογήκαμεν ὑμᾶς ἐξ οἴκου Κυρίου.

\vs{27}Θεὸς Κύριος, καὶ ἐπέφανεν ἡμῖν· συστήσασθε ἑορτὴν ἐν τοῖς πυκάζουσιν, ἕως τῶν κεράτων τοῦ θυσιαστηρίου.
\vs{28}Θεός μου εἶ σὺ, καὶ ἐξομολογήσομαί σοι· Θεός μου εἶ σὺ, καὶ ὑψώσω σε· ἐξομολογήσομαί σοι, ὅτι ἐπήκουσάς μου, καὶ ἐγένου μοι εἰς σωτηρίαν.
\vs{29}Ἐξομολογεῖσθε τῷ Κυρίῳ, ὅτι ἀγαθὸς, ὅτι εἰς τὸν αἰῶνα τὸ ἔλεος αὐτοῦ.

\begin{psalmheading}{\ch{118}{119} Ἀλληλούϊα.}
\end{psalmheading}
Μακάριοι ἄμωμοι ἐν ὁδῷ, οἱ πορευόμενοι ἐν νόμῳ Κυρίου.
\vs{2}Μακάριοι οἱ ἐξερευνῶντες τὰ μαρτύρια αὐτοῦ, ἐν ὅλῃ καρδίᾳ ἐκζητήσουσιν αὐτόν.
\vs{3}Οὐ γὰρ οἱ ἐργαζόμενοι τὴν ἀνομίαν ἐν ταῖς ὁδοῖς αὐτοῦ ἐπορεύθησαν.
\vs{4}Σὺ ἐνετείλω τὰς ἐντολάς σου, τοῦ φυλάξασθαι σφόδρα.
\vs{5}Ὄφελον κατευθυνθείησαν αἱ ὁδοί μου, τοῦ φυλάξασθαι τὰ δικαιώματά σου.
\vs{6}Τότε οὐ μὴ αἰσχυνθῶ, ἐν τῷ με ἐπιβλέπειν ἐπὶ πάσας τὰς ἐντολάς σου.
\vs{7}Ἐξομολογήσομαί σοι ἐν εὐθύτητι καρδίας, ἐν τῷ μεμαθηκέναι με τὰ κρίματα τῆς δικαιοσύνης σου.
\vs{8}Τὰ δικαιώματά σου φυλάξω, μή με ἐγκαταλίπῃς ἕως σφόδρα.

\vs{9}Ἐν τίνι κατορθώσει νεώτερος τὴν ὁδὸν αὐτοῦ; ἐν τῷ φυλάξασθαι τοὺς λόγους σου.
\vs{10}Ἐν ὅλῃ καρδίᾳ μου ἐξεζήτησά σε, μὴ ἀπώσῃ με ἀπὸ τῶν ἐντολῶν σου.
\vs{11}Ἐν τῇ καρδίᾳ μου ἔκρυψα τὰ λόγιά σου, ὅπως ἂν μὴ ἁμάρτω σοι.
\vs{12}Εὐλογητὸς εἶ Κύριε, δίδαξόν με τὰ δικαιώματά σου.
\vs{13}Ἐν τοῖς χείλεσί μου ἐξήγγειλα πάντα τὰ κρίματα τοῦ στόματός σου.
\vs{14}Ἐν τῇ ὁδῷ τῶν μαρτυρίων σου ἐτέρφθην, ὡς ἐπὶ παντὶ πλούτῳ.
\vs{15}Ἐν ταῖς ἐντολαῖς σου ἀδολεσχήσω, καὶ κατανοήσω τὰς ὁδούς σου.
\vs{16}Ἐν τοῖς δικαιώμασί σου μελετήσω, οὐκ ἐπιλήσομαι τῶν λόγων σου.

\vs{17}Ἀνταπόδος τῷ δούλῳ σου, ζήσομαι καὶ φυλάξω τοὺς λόγους σου.
\vs{18}Ἀποκάλυψον τοὺς ὀφθαλμούς μου, καὶ κατανοήσω τὰ θαυμάσια ἐκ τοῦ νόμου σου.
\vs{19}Πάροικος ἐγώ εἰμι ἐν τῇ γῇ, μὴ ἀποκρύψῃς ἀπʼ ἐμοῦ τὰς ἐντολάς σου.
\vs{20}Ἐπεπόθησεν ἡ ψυχή μου τοῦ ἐπιθυμῆσαι τὰ κρίματά σου ἐν παντὶ καιρῷ.
\vs{21}Ἐπετίμησας ὑπερηφάνοις, ἐπικατάρατοι οἱ ἐκκλίνοντες ἀπὸ τῶν ἐντολῶν σου.
\vs{22}Περίελε ἀπʼ ἐμοῦ ὄνειδος καὶ ἐξουδένωσιν, ὅτι τὰ μαρτύριά σου ἐξεζήτησα.
\vs{23}Καὶ γὰρ ἐκάθισαν ἄρχοντες, καὶ κατʼ ἐμοῦ κατελάλουν, ὁ δὲ δοῦλός σου ἠδολέσχει ἐν τοῖς δικαιώμασί σου.
\vs{24}Καὶ γὰρ τὰ μαρτύριά σου μελέτη μου ἐστὶ, καὶ αἱ συμβουλίαι μου τὰ δικαιώματά σου.

\vs{25}Ἐκολλήθη τῷ ἐδάφει ἡ ψυχή μου, ζῆσόν με κατὰ τὸ λόγον σου.
\vs{26}Τὰς ὁδούς σου ἐξήγγειλα καὶ ἐπήκουσάς μου, δίδαξόν με τὰ δικαιώματά σου.
\vs{27}Ὁδὸν δικαιωμάτων σου συνέτισόν με, καὶ ἀδολεσχήσω ἐν τοῖς θαυμασίοις σου.
\vs{28}Ἐνύσταξεν ἡ ψυχή μου ἀπὸ ἀκηδίας, βεβαίωσόν με ἐν τοῖς λόγοις σου.
\vs{29}Ὁδὸν ἀδικίας ἀπόστησον ἀπʼ ἐμοῦ, καὶ τῷ νόμῳ σου ἐλέησόν με.
\vs{30}Ὁδὸν ἀληθείας ᾑρετισάμην, καὶ τὰ κρίματά σου οὐκ ἐπελαθόμην.
\vs{31}Ἐκολλήθην τοῖς μαρτυρίοις σου Κύριε, μή με καταισχύνῃς.
\vs{32}Ὁδὸν ἐντολῶν σου ἔδραμον, ὅταν ἐπλάτυνας τὴν καρδίαν μου.

\vs{33}Νομοθέτησόν με Κύριε τὴν ὁδὸν τῶν δικαιωμάτων σου, καὶ ἐκζητήσω αὐτὴν διαπαντός.
\vs{34}Συνέτισόν με, καὶ ἐξερευνήσω τὸν νόμον σου, καὶ φυλάξω αὐτὸν ἐν ὅλῃ καρδίᾳ μου.
\vs{35}Ὁδήγησόν με ἐν τῇ τρίβῳ τῶν ἐντολῶν σου, ὅτι αὐτὴν ἠθέλησα.
\vs{36}Κλῖνον τὴν καρδίαν μου εἰς τὰ μαρτύριά σου, καὶ μὴ εἰς πλεονεξίαν.
\vs{37}Ἀπόστρεψον τοὺς ὀφθαλμούς μου τοῦ μὴ ἰδεῖν ματαιότητα, ἐν τῇ ὁδῷ σου ζῆσόν με.
\vs{38}Στῆσον τῷ δούλῳ σου τὸ λόγιόν σου εἰς τὸν φόβον σου.
\vs{39}Περίελε τὸν ὄνειδισμόν μου ὃν ὑπώπτευσα, ὅτι τὰ κρίματά σου χρηστά.
\vs{40}Ἰδοὺ ἐπεθύμησα τὰς ἐντολάς σου, ἐν τῇ δικαιοσύνῃ σου ζῆσόν με.

\vs{41}Καὶ ἔλθοι ἐπʼ ἐμὲ τὸ ἔλεός σου Κύριε, τὸ σωτήριόν σου κατὰ τὸν λόγον σου.
\vs{42}Καὶ ἀποκριθήσομαι τοῖς ὀνειδίζουσί μοι λόγον, ὅτι ἤλπισα ἐπὶ τοῖς λόγοις σου.
\vs{43}Καὶ μὴ περιέλῃς ἐκ τοῦ στόματός μου λόγον ἀληθείας ἕως σφόδρα, ὅτι ἐπὶ τοῖς κρίμασί σου ἐπήλπισα.
\vs{44}Καὶ φυλάξω τὸν νόμον σου διαπαντὸς, εἰς τὸν αἰῶνα καὶ εἰς τὸν αἰῶνα τοῦ αἰῶνος.
\vs{45}Καὶ ἐπορευόμην ἐν πλατυσμῷ, ὅτι τὰς ἐντολάς σου ἐξεζήτησα.
\vs{46}Καὶ ἐλάλουν ἐν τοῖς μαρτυρίοις σου ἐναντίον βασιλέων, καὶ οὐκ ᾐσχυνόμην.
\vs{47}Καὶ ἐμελέτων ἐν ταῖς ἐντολαῖς σου, αἷς ἠγάπησα σφόδρα.
\vs{48}Καὶ ᾖρα τὰς χεῖράς μου πρὸς τὰς ἐντολάς σου ἃς ἠγάπησα, καὶ ἠδολέσχουν ἐν τοῖς δικαιώμασί σου.

\vs{49}Μνήσθητι τῶν λόγων σου τῷ δούλῳ σου ὧν ἐπήλπισάς με.
\vs{50}Αὕτη με παρεκάλεσεν ἐν τῇ ταπεινώσει μου, ὅτι τὸ λόγιόν σου ἔζησέ με.
\vs{51}Ὑπερήφανοι παρηνόμουν ἕως σφόδρα, ἀπὸ δὲ τοῦ νόμου σου οὐκ ἐξέκλινα.
\vs{52}Ἐμνήσθην τῶν κριμάτων σου ἀπʼ αἰῶνος Κύριε, καὶ παρεκλήθην.
\vs{53}Ἀθυμία κατέσχε με ἀπὸ ἁμαρτωλῶν τῶν ἐγκαταλιμπανόντων τὸν νόμον σου.
\vs{54}Ψαλτὰ ἦσάν μοι τὰ δικαιώματά σου, ἐν τόπῳ παροικίας μου·
\vs{55}Ἐμνήσθην ἐν νυκτὶ τοῦ ὀνόματός σου Κύριε, καὶ ἐφύλαξα τὸν νόμον σου.
\vs{56}Αὕτη ἐγενήθη μοι, ὅτι τὰ δικαιώματά σου ἐξεζήτησα.

\vs{57}Μερίς μου εἶ Κύριε, εἶπα τοῦ φυλάξασθαι τὸν νόμον σου.
\vs{58}Ἐδεήθην τοῦ προσώπου σου ἐν ὅλῃ καρδίᾳ μου, ἐλέησόν με κατὰ τὸ λόγιόν σου.
\vs{59}Διελογισάμην τὰς ὁδούς σου, καὶ ἐπέστρεψα τοὺς πόδας μου εἰς τὰ μαρτύριά σου.
\vs{60}Ἡτοιμάσθην καὶ οὐκ ἐταράχθην, τοῦ φυλάξασθαι τὰς ἐντολάς σου.
\vs{61}Σχοινία ἁμαρτωλῶν περιεπλάκησάν μοι, καὶ τοῦ νόμου σου οὐκ ἐπελαθόμην.
\vs{62}Μεσονύκτιον ἐξεγειρόμην, τοῦ ἐξομολογεῖσθαί σοι ἐπὶ τὰ κρίματα τῆς δικαιοσύνης σου.
\vs{63}Μέτοχος ἐγώ εἰμι πάντων τῶν φοβουμένων σε, καὶ τῶν φυλασσόντων τὰς ἐντολάς σου.
\vs{64}Τοῦ ἐλέους σου Κύριε πλήρης ἡ γῆ, τὰ δικαιώματά σου δίδαξόν με.

\vs{65}Χρηστότητα ἐποίησας μετὰ τοῦ δούλου σου Κύριε κατὰ τὸν λόγον σου.
\vs{66}Χρηστότητα καὶ παιδείαν καὶ γνῶσιν δίδαξόν με, ὅτι ταῖς ἐντολαῖς σου ἐπίστευσα.
\vs{67}Πρὸ τοῦ με ταπεινωθῆναι, ἐγὼ ἐπλημμέλησα, διὰ τοῦτο τὸ λόγιόν σου ἐφύλαξα.
\vs{68}Χρηστὸς εἶ σὺ Κύριε· καὶ ἐν τῇ χρηστότητί σου δίδαξόν με τὰ δικαιωματά σου.
\vs{69}Ἐπληθύνθη ἐπʼ ἐμὲ ἀδικία ὑπερηφάνων, ἐγὼ δὲ ἐν ὅλῃ καρδίᾳ μου ἐξερευνήσω τὰς ἐντολάς σου.
\vs{70}Ἐτυρώθη ὡς γάλα ἡ καρδία αὐτῶν, ἐγὼ δὲ τὸν νόμον σου ἐμελέτησα.
\vs{71}Ἀγαθόν μοι ὅτι ἐταπείνωσάς με, ὅπως ἂν μάθω τὰ δικαιώματά σου.
\vs{72}Ἀγαθόν μοι ὁ νόμος τοῦ στόματός σου, ὑπὲρ χιλιάδας χρυσίου καὶ ἀργυρίου.

\vs{73}Αἱ χεῖρές σου ἔποίησάν με καὶ ἐπλασάν με, συνέτισόν με καὶ μαθήσομαι τὰς ἐντολάς σου.
\vs{74}Οἱ φοβούμενοί σε ὄψονταί με καὶ εὐφρανθήσονται, ὅτι εἰς τοὺς λόγους σου ἐπήλπισα.
\vs{75}Ἔγνων Κύριε ὅτι δικαιοσύνη τὰ κρίματά σου, καὶ ἀληθείᾳ ἐταπείνωσάς με.
\vs{76}Γενηθήτω δὴ τὸ ἔλεός σου τοῦ παρακαλέσαι με, κατὰ τὸ λόγιόν σου τῷ δούλῳ σου.
\vs{77}Ἐλθέτωσάν μοι οἱ οἰκτιρμοί σου, καὶ ζήσομαι, ὅτι ὁ νόμος σου μελέτη μου ἐστίν.
\vs{78}Αἰσχυνθήτωσαν ὑπερήφανοι, ὅτι ἀδίκως ἠνόμησαν εἰς ἐμὲ, ἐγὼ δὲ ἀδολεσχήσω ἐν ταῖς ἐντολαῖς σου.
\vs{79}Ἐπιστρεψάτωσάν με οἱ φοβούμενοί σε, καὶ οἱ γινώσκοντες τὰ μαρτύριά σου.
\vs{80}Γενηθήτω ἡ καρδία μου ἄμωμος ἐν τοῖς δικαιώμασί σου, ὅπως ἂν μὴ αἰσχυνθῶ.

\vs{81}Ἐκλείπει εἰς τὸ σωτήριόν σου ἡ ψυχή μου, εἰς τοὺς λόγους σου ἐπήλπισα.
\vs{82}Ἐξέλιπον οἱ ὀφθαλμοί μου εἰς τὸ λόγιόν σου, λέγοντες, πότε παρακαλέσεις με;
\vs{83}Ὅτι ἐγενήθην ὡς ἀσκὸς ἐν πάχνῃ· τὰ δικαιώματά σου οὐκ ἐπελαθόμην.
\vs{84}Πόσαι εἰσὶν αἱ ἡμέραι τοῦ δούλου σου; πότε ποιήσεις μοι ἐκ τῶν καταδιωκόντων με κρίσιν;
\vs{85}Διηγήσαντό μοι παράνομοι ἀδολεσχίας, ἀλλʼ οὐχ ὡς ὁ νόμος σου Κύριε.
\vs{86}Πᾶσαι αἱ ἐντολαί σου ἀλήθεια· ἀδίκως κατεδίωξάν με, βοήθησόν μοι.
\vs{87}Παρὰ βραχὺ συνετέλεσάν με ἐν τῇ γῇ, ἐγὼ δὲ οὐκ ἐγκατέλιπον τὰς ἐντολάς σου.
\vs{88}Κατὰ τὸ ἔλεός σου ζῆσόν με, καὶ φυλάξω τὰ μαρτύρια τοῦ στόματός σου.

\vs{89}Εἰς τὸν αἰῶνα, Κύριε, ὁ λόγος σου διαμένει ἐν τῷ οὐρανῷ,
\vs{90}εἰς γενεὰν καὶ γενεὰν ἡ ἀλήθειά σου· ἐθεμελίωσας τὴν γῆν καὶ διαμένει.
\vs{91}Τῇ διατάξει σου διαμένει ἡμέρα, ὅτι τὰ σύμπαντα δοῦλα σά.
\vs{92}Εἰ μὴ ὅτι ὁ νόμος σου μελέτη μου ἐστὶ, τότε ἂν ἀπωλόμην ἐν τῇ ταπεινώσει μου.
\vs{93}Εἰς τὸν αἰῶνα οὐ μὴ ἐπιλάθωμαι τῶν δικαιωμάτων σου, ὅτι ἐν αὐτοῖς ἔζησάς με.
\vs{94}Σός εἰμι ἐγὼ, σῶσόν με, ὅτι τὰ δικαιώματά σου ἐξεζήτησα.
\vs{95}Ἐμὲ ὑπέμειναν ἁμαρτωλοὶ τοῦ ἀπολέσαι με, τὰ μαρτύριά σου συνῆκα.
\vs{96}Πάσης συντελείας εἶδον πέρας, πλατεία ἡ ἐντολή σου σφόδρα.

\vs{97}Ὡς ἠγάπησα τὸν νόμον σου Κύριε; ὅλην τὴν ἡμέραν μελέτη μου ἐστίν.
\vs{98}Ὑπὲρ τοὺς ἐχθρούς μου ἐσόφισάς με τὴν ἐντολήν σου, ὅτι εἰς τὸν αἰῶνα ἐμή ἐστιν.
\vs{99}Ὑπὲρ πάντας τοὺς διδάσκοντάς με συνῆκα, ὅτι τὰ μαρτύριά σου μελέτη μου ἐστίν.
\vs{100}Ὑπὲρ πρεσβυτέρους συνῆκα, ὅτι τὰς ἐντολάς σου ἐξεζήτησα.
\vs{101}Ἐκ πάσης ὁδοῦ πονηρᾶς ἐκώλυσα τοὺς πόδας μου, ὅπως ἂν φυλάξω τοὺς λόγους σου.
\vs{102}Ἀπὸ τῶν κριμάτων σου οὐκ ἐξέκλινα, ὅτι σὺ ἐνομοθέτησάς με.

\vs{103}Ὡς γλυκέα τῷ λάρυγγί μου τὰ λόγιά σου, ὑπὲρ μέλι τῷ στόματί μου.
\vs{104}Ἀπὸ τῶν ἐντολῶν σου συνῆκα, διὰ τοῦτο ἐμίσησα πᾶσαν ὁδὸν ἀδικίας.

\vs{105}Λύχνος τοῖς ποσί μου ὁ νόμος σου, καὶ φῶς ταῖς τρίβοις μου.
\vs{106}Ὤμοσα καὶ ἔστησα τοῦ φυλάξασθαι τὰ κρίματα τῆς δικαιοσύνης σου.
\vs{107}Ἐταπεινώθην ἕως σφόδρα Κύριε, ζῆσόν με κατὰ τὸν λόγον σου.
\vs{108}Τὰ ἑκούσια τοῦ στόματός μου εὐλόκησον δὴ Κύριε, καὶ τὰ κρίματά σου δίδαξόν με.
\vs{109}Ἡ ψυχή μου ἐν ταῖς χερσὶ σου διαπαντὸς, καὶ τὸν νόμον σου οὐκ ἐπελαθόμην.
\vs{110}Ἔθεντο ἁμαρτωλοὶ παγίδα μοι, καὶ ἐκ τῶν ἐντολῶν σου οὐκ ἐπλανήθην.
\vs{111}Ἐκληρονόμησα τὰ μαρτύριά σου εἰς τὸν αἰῶνα, ὅτι ἀγαλλιάμα τῆς καρδίας μού εἰσιν.
\vs{112}Ἔκλινα τὴν καρδίαν μου τοῦ ποιῆσαι τὰ δικαιώματά σου εἰς τὸν αἰῶνα διʼ ἀντάμειψιν.

\vs{113}Παρανόμους ἐμίσησα, τὸν δὲ νόμον σου ἠγάπησα.
\vs{114}Βοηθός μου, καὶ ἀντιλήπτωρ μου εἶ σὺ, εἰς τοὺς λόγους σου ἐπήλπισα.
\vs{115}Ἐκκλίνατε ἀπʼ ἐμοῦ οἱ πονηρευόμενοι, καὶ ἐξεραυνήσω τὰς ἐντολὰς τοῦ Θεοῦ μου.
\vs{116}Ἀντιλαβοῦ μου κατὰ τὸ λόγιόν σου, καὶ ζῆσόν με, καὶ μὴ καταισχύνῃς με ἀπὸ τῆς προσδοκίας μου.
\vs{117}Βοήθησόν μοι, καὶ σωθήσομαι, καὶ μελετήσω ἐν τοῖς δικαιώμασί σου διαπαντός.
\vs{118}Ἐξουδένωσας πάντας τοὺς ἀποστατοῦντας ἀπὸ τῶν δικαιωμάτων σου, ὅτι ἄδικον τὸ ἐνθύμημα αὐτῶν.
\vs{119}Παραβαίνοντας ἐλογισάμην πάντας τοὺς ἁμαρτωλοὺς τῆς γῆς, διὰ τοῦτο ἠγάπησα τὰ μαρτύριά σου.
\vs{120}Καθήλωσον ἐκ τοῦ φόβου σου τὰς σάρκας μου, ἀπὸ γὰρ τῶν κριμάτων σου ἐφοβήθην.

\vs{121}Ἐποίησα κρίμα καὶ δικαιοσύνην, μὴ παραδῷς με τοῖς ἀδικοῦσί με.
\vs{122}Ἔνδεξαι τὸν δοῦλόν σου εἰς ἀγαθὸν, μὴ συκοφαντησάτωσάν με ὑπερήφανοι.
\vs{123}Οἱ ὀφθαλμοί μου ἐξέλιπον εἰς τὸ σωτήριόν σου, καὶ εἰς τὸ λόγιον τῆς δικαιοσύνης σου.
\vs{124}Ποίησον μετὰ τοῦ δούλου σου κατὰ τὸ ἔλεός σου, καὶ τὰ δικαιώματά σου δίδαξόν με.
\vs{125}Δοῦλός σου εἰμὶ ἐγὼ, συνέτισόν με καὶ γνώσομαι τὰ μαρτύριά σου.
\vs{126}Καιρὸς τοῦ ποιῆσαι τῷ Κυρίῳ, διεσκέδασαν τὸν νόμον σου.
\vs{127}Διὰ τοῦτο ἠγάπησα τὰς ἐντολάς σου ὑπὲρ χρυσίον καὶ τοπάζιον.
\vs{128}Διὰ τοῦτο πρὸς πάσας τὰς ἐντολάς σου κατωρθούμην, πᾶσαν ὁδὸν ἄδικον ἐμίσησα.

\vs{129}Θαυμαστὰ τὰ μαρτύριά σου, διὰ τοῦτο ἐξηρεύνησεν αὐτὰ ἡ ψυχή μου.
\vs{130}Ἡ δήλωσις τῶν λόγων σου φωτιεῖ καὶ συνετιεῖ νηπίους.
\vs{131}Τὸ στόμα μου ἤνοιξα, καὶ εἵλκυσα πνεῦμα, ὅτι τὰς ἐντολάς σου ἐπεπόθουν.
\vs{132}Ἐπίβλεψον ἐπʼ ἐμὲ καὶ ἐλέησόν με, κατὰ τὸ κρίμα τῶν ἀγαπώντων τὸ ὄνομά σου.
\vs{133}Τὰ διαβήματά μου κατεύθυνον κατὰ τὸ λόγιόν σου, καὶ μὴ κατακυριευσάτω μου πᾶσα ἀνομία.
\vs{134}Λύτρωσαί με ἀπὸ συκοφαντίας ἀνθρώπων, καὶ φυλάξω τὰς ἐντολάς σου.
\vs{135}Τὸ πρόσωπόν σου ἐπίφανον ἐπὶ τὸν δοῦλόν σου, καὶ δίδαξόν με τὰ δικαιώματά σου.
\vs{136}Διεξόδους ὑδάτων κατέβησαν οἱ ὀφθαλμοί μου, ἐπεὶ οὐκ ἐφύλαξα τὸν νόμον σου.

\vs{137}Δίκαιος εἶ Κύριε, καὶ εὐθεῖς αἱ κρίσεις σου.
\vs{138}Ἐνετείλω δικαιοσύνην τὰ μαρτύριά σου, καὶ ἀλήθειαν σφόδρα.
\vs{139}Ἐξέτηξέ με ὁ ζῆλός σου, ὅτι ἐπελάθοντο τῶν λόγων σου οἱ ἐχθροί μου.
\vs{140}Πεπυρωμένον τὸ λόγιόν σου σφόδρα, καὶ ὁ δοῦλός σου ἠγάπησεν αὐτό.
\vs{141}Νεώτερος ἐγὼ εἰμι καὶ ἐξουδενωμένος, τὰ δικαιώματά σου οὐκ ἐπελαθόμην.
\vs{142}Ἡ δικαιοσύνη σου δικαιοσύνη εἰς τὸν αἰῶνα, καὶ ὁ νόυος σου ἀλήθεια.
\vs{143}Θλίψεις καὶ ἀνάγκαι εὕροσάν με, ἐντολαί σου μελέτη μου.
\vs{144}Δικαιοσύνη τὰ μαρτύριά σου εἰς τὸν αἰῶνα· συνέτισόν με, καὶ ζήσομαι.

\vs{145}Ἐκέκραξα ἐν ὅλῃ καρδίᾳ μου, ἐπάκουσόν μου Κύριε, τὰ δικαιώματά σου ἐκζητήσω.
\vs{146}Ἐκέκραξά σοι, σῶσόν με, καὶ φυλάξω τὰ μαρτύριά σου.
\vs{147}Προέφθασα ἐν ἀωρίᾳ καὶ ἐκέκραξα, εἰς τοὺς λόγους σου ἐπήλπισα.
\vs{148}Προέφθασαν οἱ ὀφθαλμοί μου πρὸς ὄρθρον, τοῦ μελετᾷν τὰ λόγιά σου.
\vs{149}Τῆς φωνῆς μου ἄκουσον Κύριε κατὰ τὸ ἔλεός σου, κατὰ τὸ κρίμα σου ζῆσόν με.
\vs{150}Προσήγγισαν οἱ καταδιώκοντές με ἀνομίᾳ, ἀπὸ δὲ τοῦ νόμου σου ἐμακρύνθησαν.
\vs{151}Ἐγγὺς εἶ Κύριε, καὶ πᾶσαι αἱ ὁδοί σου ἀλήθεια.
\vs{152}Κατʼ ἀρχὰς ἔγνων ἐκ τῶν μαρτυρίων σου, ὅτι εἰς τὸν αἰῶνα ἐθεμελίωσας αὐτά.

\vs{153}Ἴδε τὴν ταπείνωσίν μου καὶ ἐξελοῦ με, ὅτι τοῦ νόμου σου οὐκ ἐπελαθόμην.
\vs{154}Κρῖνον τὴν κρίσιν μου καὶ λύτρωσαί με, διὰ τὸν λόγον σου ζῆσόν με.
\vs{155}Μακρὰν ἀπὸ ἁμαρτωλῶν σωτηρία, ὅτι τὰ δικαιώματά σου οὐκ ἐξεζήτησαν.
\vs{156}Οἱ οἰκτιρμοί σου πολλοὶ Κύριε, κατὰ τὸ κρίμα σου ζῆσόν με.
\vs{157}Πολλοὶ οἱ ἐκδιώκοντές με καὶ θλίβοντές με· ἐκ τῶν μαρτυρίων σου οὐκ ἐξέκλινα.
\vs{158}Εἶδον ἀσυνετοῦντας καὶ ἐξετηκόμην, ὅτι τὰ λόγιά σου οὐκ ἐφυλάξαντο.
\vs{159}Ἴδε ὅτι τὰς ἐντολάς σου ἠγάπησα Κύριε, ἐν τῷ ἐλέει σου ζῆσόν με.
\vs{160}Ἀρχὴ τῶν λόγων σου ἀλήθεια, καὶ εἰς τὸν αἰῶνα πάντα τὰ κρίματα τῆς δικαιοσύνης σου.

\vs{161}Ἄρχοντες κατεδίωξάν με δωρεὰν, καὶ ἀπὸ τῶν λόγων σου ἐδειλίασεν ἡ καρδία μου.
\vs{162}Ἀγαλλιάσομαι ἐγὼ ἐπὶ τὰ λόγιά σου, ὡς ὁ εὑρίσκων σκῦλα πολλά.
\vs{163}Ἀδικίαν ἐμίσησα καὶ ἐβδελυξάμην, τὸν δὲ νόμον σου ἠγάπησα.
\vs{164}Ἑπτάκις τῆς ἡμέρας ᾔνεσά σε ἐπὶ τὰ κρίματα τῆς δικαιοσύνης σου.
\vs{165}Εἰρήνη πολλὴ τοῖς ἀγαπῶσι τὸν νόμον σου, καὶ οὐκ ἔστιν αὐτοῖς σκάνδαλον.
\vs{166}Προσεδόκων τὸ σωτήριόν σου Κύριε, καὶ τὰς ἐντολάς σου ἠγάπησα.
\vs{167}Ἐφύλαξεν ἡ ψυχή μου τὰ μαρτύριά σου, καὶ ἠγάπησεν αὐτὰ σφόδρα.
\vs{168}Ἐφύλαξα τὰς ἐντολάς σου καὶ τὰ μαρτύριά σου, ὅτι πᾶσαι αἱ ὁδοί μου ἐναντίον σου Κύριε.

\vs{169}Ἐγγυσάτω ἡ δέησίς μου ἐνώπιόν σου Κύριε, κατὰ τὸ λόγιόν σου συνέτισόν με.
\vs{170}Εἰσέλθοι τὸ ἀξίωμά μου ἐνώπιόν σου Κύριε, κατὰ τὸ λόγιόν σου ῥῦσαί με.
\vs{171}Ἐξερεύξαιντο τὰ χείλη μου ὕμνον, ὅταν διδάξῃς με τὰ δικαιώματά σου.
\vs{172}Φθέγξαιτο ἡ γλῶσσά μου τὰ λόγιά σου, ὅτι πᾶσαι αἱ ἐντολαί σου δικαιοσύνη.
\vs{173}Γενέσθω ἡ χείρ σου τοῦ σῶσαί με, ὅτι τὰς ἐντολάς σου ᾑρετισάμην.
\vs{174}Ἐπεπόθησα τὸ σωτήριόν σου Κύριε, καὶ ὁ νόμος σου μελέτη μου ἐστί.
\vs{175}Ζήσεται ἡ ψυχή μου καὶ αἰνέσει σε, καὶ τὰ κρίματά σου βοηθήσει μοι.
\vs{176}Ἐπλανήθην ὡς πρόβατον ἀπολωλὸς, ζήτησον τὸν δοῦλόν σου, ὅτι τὰς ἐντολάς σου οὐκ ἐπελαθόμην.

\begin{psalmheading}{\ch{119}{120} Ὠδὴ τῶν ἀναβαθμῶν.}
\end{psalmheading}
Πρὸς Κύριον ἐν τῷ θλίβεσθαί με ἐκέκραξα, καὶ εἰσήκουσέ μου.
\vs{2}Κύριε, ῥῦσαι τὴν ψυχήν μου ἀπὸ χειλέων ἀδίκων καὶ ἀπὸ γλώσσης δολίας.

\vs{3}Τί δοθείη σοι, καὶ τί προστεθείη σοι πρὸς γλῶσσαν δολίαν;
\vs{4}Τὰ βέλη τοῦ δυνατοῦ ἠκονημένα σὺν τοῖς ἄνθραξι τοῖς ἐρημικοῖς.

\vs{5}Οἴμοι ὅτι ἡ παροικία μου ἐμακρύνθη, κατεσκήνωσα μετὰ τῶν σκηνωμάτων Κηδάρ.
\vs{6}Πολλὰ παρῴκησεν ἡ ψυχή μου· μετὰ τῶν μισούντων τὴν εἰρήνην.
\vs{7}Ἤμην εἰρηνικός· ὅταν ἐλάλουν αὐτοῖς, ἐπολέμουν με δωρεάν.

\begin{psalmheading}{\ch{120}{121} Ὠδὴ τῶν ἀναβαθμῶν.}
\end{psalmheading}
Ἦρα τοὺς ὀφθαλμούς μου εἰς τὰ ὄρη, ὅθεν ἥξει ἡ βοήθειά μου.
\vs{2}Ἡ βοήθειά μου παρὰ Κυρίου τοῦ ποιήσαντος τὸν οὐρανὸν καὶ τὴν γῆν.
\vs{3}Μὴ δῴης εἰς σάλον τὸν πόδα σου, μηδὲ νυστάξῃ ὁ φυλάσσων σε.
\vs{4}Ἰδοὺ οὐ νυστάξει οὐδὲ ὑπνώσει ὁ φυλάσσων τὸν Ἰσραήλ.
\vs{5}Κύριος φυλάξει σε, Κύριος σκέπη σου ἐπὶ χεῖρα δεξιάν σου.
\vs{6}Ἡμέρας ὁ ἥλιος οὐ συγκαύσει σε, οὐδὲ ἡ σελήνη τὴν νύκτα.
\vs{7}Κύριος φυλάξαι σε ἀπὸ παντὸς κακοῦ, φυλάξει τὴν ψυχήν σου ὁ Κύριος.
\vs{8}Κύριος φυλάξει τὴν εἴσοδόν σου, καὶ τὴν ἔξοδόν σου, ἀπὸ τοῦ νῦν καὶ ἕως τοῦ αἰῶνος.

\begin{psalmheading}{\ch{121}{122} Ὠδὴ τῶν ἀναβαθμῶν.}
\end{psalmheading}
Εὐφράνθην ἐπὶ τοῖς εἰρηκόσι μοι, εἰς οἶκον Κυρίου πορευσόμεθα.
\vs{2}Ἑστῶτες ἦσαν οἱ πόδες ἡμῶν ἐν ταῖς αὐλαῖς σου Ἱερουσαλήμ.
\vs{3}Ἱερουσαλὴμ οἰκοδομουμένη ὡς πόλις, ἧς ἡ μετοχὴ αὐτῆς ἐπιτοαυτό.
\vs{4}Ἐκεῖ γὰρ ἀνέβησαν αἱ φυλαὶ, φυλαὶ Κυρίου μαρτύριον τῷ Ἰσραὴλ, τοῦ ἐξομολογήσασθαι τῷ ὀνόματι Κυρίου.
\vs{5}Ὅτι ἐκεῖ ἐκάθισαν θρόνοι εἰς κρίσιν, θρόνοι ἐπὶ οἶκον Δαυίδ.

\vs{6}Ἐρωτήσατε δὴ τὰ εἰς εἰρήνην τὴν Ἱερουσαλὴμ, καὶ εὐθηνία τοῖς ἀγαπῶσί σε.
\vs{7}Γενέσθω δὴ εἰρήνη ἐν τῇ δυνάμει σου, καὶ εὐθηνία ἐν ταῖς πυργοβάρεσί σου.
\vs{8}Ἕνεκα τῶν ἀδελφῶν μου καὶ τῶν πλησίον μου, ἐλάλουν δὴ εἰρήνην περὶ σοῦ.
\vs{9}Ἕνεκα τοῦ οἴκου Κυρίου τοῦ Θεοῦ ἡμῶν ἐξεζήτησα ἀγαθά σοι.

\begin{psalmheading}{\ch{122}{123} Ὠδὴ τῶν ἀναβαθμῶν.}
\end{psalmheading}
Πρὸς σὲ ᾖρα τοὺς ὀφθαλμούς μου, τὸν κατοικοῦντα ἐν τῷ οὐρανῷ.
\vs{2}Ἰδοὺ ὡς ὀφθαλμοὶ δούλων εἰς χεῖρας τῶν κυρίων αὐτῶν, ὡς ὀφθαλμοὶ παιδίσκης εἰς χεῖρας τῆς κυρίας αὐτῆς, οὕτως οἱ ὀφθαλμοὶ ἡμῶν πρὸς Κύριον τὸν Θεὸν ἡμῶν, ἕως οὗ οἰκτειρῆσαι ἡμᾶς.
\vs{3}Ἐλέησον ἡμᾶς Κύριε, ἐλέησον ἡμᾶς, ὅτι ἐπὶ πολὺ ἐπλήσθημεν ἐξουδενώσεως.
\vs{4}Ἐπὶ πλεῖον ἐπλήσθη ἡ ψυχὴ ἡμῶν· τὸ ὄνειδος τοῖς εὐθηνοῦσι καὶ ἡ ἐξουδένωσις τοῖς ὑπερηφάνοις.

\begin{psalmheading}{\ch{123}{124} Ὠδὴ τῶν ἀναβαθμῶν.}
\end{psalmheading}
Εἰ μὴ ὅτι Κύριος ἦν ἐν ἡμῖν, εἰπάτω δὴ Ἰσραὴλ,
\vs{2}εἰ μὴ ὅτι Κύριος ἦν ἐν ἡμῖν, ἐν τῷ ἐπαναστῆναι ἀνθρώπους ἐφʼ ἡμᾶς,
\vs{3}ἄρα ζῶντας ἂν κατέπιον ἡμᾶς· ἐν τῷ ὀργισθῆναι τὸν θυμὸν αὐτῶν ἐφʼ ἡμᾶς,
\vs{4}ἄρα τὸ ὕδωρ ἂν κατεπόντισεν ἡμᾶς· χείμαῤῥον διῆλθεν ἡ ψυχὴ ἡμῶν.
\vs{5}Ἄρα διῆλθεν ἡ ψυχὴ ἡμῶν τὸ ὕδωρ τὸ ἀνυπόστατον.

\vs{6}Εὐλογητὸς Κύριος, ὃς οὐκ ἔδωκεν ἡμᾶς εἰς θήραν τοῖς ὀδοῦσιν αὐτῶν.
\vs{7}Ἡ ψυχὴ ἡμῶν ὡς στρουθίον ἐῤῥύσθη ἐκ τῆς παγίδος τῶν θηρευόντων· ἡ παγὶς συνετρίβη, καὶ ἡμεῖς ἐῤῥύσθημεν.
\vs{8}Ἡ βοήθεια ἡμῶν ἐν ὀνόματι Κυρίου, τοῦ ποιήσαντος τὸν οὐρανὸν καὶ τὴν γῆν.

\begin{psalmheading}{\ch{124}{125} Ὠδὴ τῶν ἀναβαθμῶν.}
\end{psalmheading}
Οἱ πεποιθότες ἐπὶ Κύριον ὡς ὄρος Σιὼν, οὐ σαλευθήσεται εἰς αἰῶνα ὁ κατοικῶν Ἱερουσαλήμ·
\vs{2}ὄρη κύκλῳ αὐτῆς, καὶ ὁ Κύριος κύκλῳ τοῦ λαοῦ αὐτοῦ, ἀπὸ τοῦ νῦν καὶ ἕως τοῦ αἰῶνος.
\vs{3}Ὅτι οὐκ ἀφήσει Κύριος τὴν ῥάβδον τῶν ἁμαρτωλῶν ἐπὶ τὸν κλῆρον τῶν δικαίων, ὅπως ἂν μὴ ἐκτείνωσιν οἱ δίκαιοι ἐν ἀνομίαις χεῖρας αὐτῶν.

\vs{4}Ἀγάθυνον, Κύριε, τοῖς ἀγαθοῖς καὶ τοῖς εὐθέσι τῇ καρδίᾳ.
\vs{5}Τοὺς δὲ ἐκκλίνοντας εἰς τὰς στραγγαλιὰς, ἀπάξει Κύριος μετὰ τῶν ἐργαζομένων τὴν ἀνομίαν. εἰρήνη ἐπὶ τὸν Ἰσραήλ.

\begin{psalmheading}{\ch{125}{126} Ὠδὴ τῶν ἀναβαθμῶν.}
\end{psalmheading}
Ἐν τῷ ἐπιστρέψαι Κύριον τὴν αἰχμαλωσίαν Σιὼν, ἐγενήθημεν ὡσεὶ παρακεκλημένοι.
\vs{2}Τότε ἐπλήσθη χαρᾶς τὸ στόμα ἡμῶν, καὶ ἡ γλῶσσα ἡμῶν ἀγαλλιάσεως· τότε ἐροῦσιν ἐν τοῖς ἔθνεσιν, ἐμεγάλυνε Κύριος τοῦ ποιῆσαι μετʼ αὐτῶν.
\vs{3}Ἐμεγάλυνε Κύριος τοῦ ποιῆσαι μεθʼ ἡμῶν, ἐγενήθημεν εὐφραινόμενοι.

\vs{4}Ἐπίστρεψον, Κύριε, τὴν αἰχμαλωσίαν ἡμῶν ὡς χειμάῤῥους ἐν τῷ Νότῳ.
\vs{5}Οἱ σπείροντες ἐν δάκρυσιν, ἐν ἀγαλλιάσει θεριοῦσι.
\vs{6}Πορευόμενοι ἐπορεύοντο, καὶ ἔκλαιον βάλλοντες τὰ σπέρματα αὐτῶν· ἐρχόμενοι δὲ ἥξουσιν ἐν ἀγαλλιάσει αἴροντες τὰ δράγματα αὐτῶν.

\begin{psalmheading}{\ch{126}{127} Ὠδὴ τῶν ἀναβαθμῶν.}
\end{psalmheading}
Ἐὰν μὴ Κύριος οἰκοδομήσῃ οἶκον, εἰς μάτην ἐκοπίασαν οἱ οἰκοδομοῦντες· ἐὰν μὴ Κύριος φυλάξῃ πόλιν, εἰς μάτην ἠγρύπνησεν ὁ φυλάσσων.
\vs{2}Εἰς μάτην ὑμῖν ἐστι τὸ ὀρθρίζειν· ἐγείρεσθε μετὰ τὸ καθῆσθαι, οἱ ἐσθίοντες ἄρτον ὀδύνης, ὅταν δῷ τοῖς ἀγαπητοῖς αὐτοῦ ὕπνον.

\vs{3}Ἰδοὺ ἡ κληρονομία Κυρίου, υἱοί, ὁ μισθὸς τοῦ καρποῦ τῆς γαστρός.
\vs{4}Ὡσεὶ βέλη ἐν χειρὶ δυνατοῦ, οὕτως οἱ υἱοὶ τῶν ἐκτετιναγμένων.
\vs{5}Μακάριος ὃς πληρώσει τὴν ἐπιθυμίαν αὐτοῦ ἐξ αὐτῶν· οὐ καταισχυνθήσονται, ὅταν λαλῶσι τοῖς ἐχθροῖς αὐτῶν ἐν πύλαις.

\begin{psalmheading}{\ch{127}{128} Ὠδὴ τῶν ἀναβαθμῶν.}
\end{psalmheading}
Μακάριοι πάντες οἱ φοβούμενοι τὸν Κύριον, οἱ πορευόμενοι ἐν ταῖς ὁδοῖς αὐτοῦ.
\vs{2}Τοὺς πόνους τῶν καρπῶν σου φάγεσαι· μακάριος εἶ καὶ καλῶς σοι ἔσται.
\vs{3}Ἡ γυνή σου ὡς ἄμπελος εὐθηνοῦσα ἐν ταῖς κλίτεσι τῆς οἰκίας σου· οἱ υἱοί σου ὡς νεόφυτα ἐλαιῶν κύκλῳ τῆς τραπέζης σου.

\vs{4}Ἰδοὺ οὕτως εὐλογηθήσεται ἄνθρωπος ὁ φοβούμενος τὸν Κύριον.
\vs{5}Εὐλογήσαι σε Κύριος ἐκ Σιὼν, καὶ ἴδοις τὰ ἀγαθὰ Ἱερουσαλὴμ πάσας τὰς ἡμέρας τῆς ζωῆς σου·
\vs{6}Καὶ ἴδοις υἱοὺς τῶν υἱῶν σου· εἰρήνη ἐπὶ τὸν Ἰσραήλ.

\begin{psalmheading}{\ch{128}{129} Ὠδὴ τῶν ἀναβαθμῶν.}
\end{psalmheading}
Πλεονάκις ἐπολέμησάν με ἐκ νεότητός μου, εἰπάτω δὴ Ἰσραήλ·
\vs{2}Πλεονάκις ἐπολέμησάν με ἐκ νεότητός μου, καὶ γὰρ οὐκ ἠδυνήθησάν μοι.
\vs{3}Ἐπὶ τὸν νῶτόν μου ἐτέκταινον οἱ ἁμαρτωλοὶ, ἐμάκρυναν τὴν ἀνομίαν αὐτῶν.
\vs{4}Κύριος δίκαιος συνέκοψεν αὐχένας ἁμαρτωλῶν.

\vs{5}Αἰσχυνθήτωσαν καὶ ἀποστραφήτωσαν εἰς τὰ ὀπίσω πάντες οἱ μισοῦντες Σιών.
\vs{6}Γενηθήτωσαν ὡσεὶ χόρτος δωμάτων, ὃς πρὸ τοῦ ἐκσπασθῆναι ἐξηράνθη·
\vs{7}Οὗ οὐκ ἐπλήρωσεν τὴν χεῖρα αὐτοῦ ὁ θερίζων, καὶ τὸν κόλπον αὐτοῦ ὁ τὰ δράγματα συλλέγων.
\vs{8}Καὶ οὐκ εἶπαν οἱ παράγοντες, εὐλογία Κυρίου ἐφʼ ὑμᾶς, εὐλογήκαμεν ὑμᾶς ἐν ὀνόματι Κυρίου.

\begin{psalmheading}{\ch{129}{130} Ὠδὴ τῶν ἀναβαθμῶν.}
\end{psalmheading}
Ἐκ βαθέων ἐκέκραξά σοι Κύριε.
\vs{2}Κύριε εἰσάκουσον τῆς φωνῆς μου· γενηθήτω τὰ ὦτά σου προσέχοντα εἰς τὴν φωνὴν τῆς δεήσεώς μου.
\vs{3}Ἐὰν ἀνομίας παρατηρήσῃς Κύριε, Κύριε, τίς ὑποστήσεται;
\vs{4}Ὅτι παρὰ σοὶ ὁ ἱλασμός ἐστιν· ἕνεκεν τοῦ ὀνόματός σου
\vs{5}ὑπέμεινά σε Κύριε, ὑπέμεινεν ἡ ψυχή μου εἰς τὸν λόγον σου,
\vs{6}ἤλπισεν ἡ ψυχή μου ἐπὶ τὸν Κύριον· ἀπὸ φυλακῆς πρωΐας μέχρι νυκτός,

\vs{7}Ἐλπισάτω Ἰσραὴλ ἐπὶ τὸν Κύριον· ὅτι παρὰ τῷ Κυρίῳ τὸ ἔλεος, καὶ πολλὴ παρ αὐτῷ λύτρωσις.
\vs{8}Καὶ αὐτὸς λυτρώσεται τὸν Ἰσραὴλ ἐκ πασῶν τῶν ἀνομιῶν αὐτοῦ.

\begin{psalmheading}{\ch{130}{131} Ὠδὴ τῶν ἀναβαθμῶν.}
\end{psalmheading}
Κύριε, οὐχ ὑψώθη ἡ καρδία μου, οὐδὲ ἐμετεωρίσθησαν οἱ ὀφθαλμοί μου· οὐδὲ ἐπορεύθην ἐν μεγάλοις, οὐδὲ ἐν θαυμασίοις ὑπὲρ ἐμέ.
\vs{2}Εἰ μὴ ἐταπεινοφρόνουν, ἀλλὰ ὕψωσα τὴν ψυχήν μου· ὡς τὸ ἀπογεγαλακτισμένον ἐπὶ τὴν μητέρα αὐτοῦ, ὡς ἀνταποδώσεις ἐπὶ τὴν ψυχήν μου.
\vs{3}Ἐλπισάτω Ἰσραὴλ ἐπὶ τὸν Κύριον ἀπὸ τοῦ νῦν καὶ ἕως τοῦ αἰῶνος.

\begin{psalmheading}{\ch{131}{132} Ὠδὴ τῶν ἀναβαθμῶν.}
\end{psalmheading}
Μνήσθητι Κύριε τοῦ Δαυὶδ, καὶ πάσης τῆς πρᾳότητος αὐτοῦ·
\vs{2}Ὡς ὤμοσε τῷ Κυρίῳ, ηὔξατο τῷ Θεῷ Ἰακώβ·
\vs{3}Εἰ εἰσελεύσομαι εἰς σκήνωμα οἴκου μου, εἰ ἀναβήσομαι ἐπὶ κλίνης στρωμνῆς μου·
\vs{4}Εἰ δώσω ὕπνον τοῖς ὀφθαλμοῖς μοῦ, καὶ τοῖς βλεφάροις μου νυσταγμὸν, καὶ ἀνάπαυσιν τοῖς κροτάφοις μου·
\vs{5}Ἕως οὗ εὕρω τόπον τῷ Κυρίῳ, σκήνωμα τῷ Θεῷ Ἰακώβ.
\vs{6}Ἰδοὺ ἠκούσαμεν αὐτὴν ἐν Ἐφραθὰ, εὕρομεν αὐτὴν ἐν τοῖς πεδίοις τοῦ δρυμοῦ.
\vs{7}Εἰσελευσώμεθα εἰς τὰ σκηνώματα αὐτοῦ· προσκυνήσωμεν εἰς τὸν τόπον οὗ ἔστησαν οἱ πόδες αὐτοῦ.

\vs{8}Ἀνάστηθι Κύριε εἰς τὴν ἀνάπαυσίν σου, σὺ καὶ ἡ κιβωτὸς τοῦ ἁγιάσματός σου.
\vs{9}Οἱ ἱερεῖς σου ἐνδύσονται δικαιοσύνην, καὶ οἱ ὅσιοί σου ἀγαλλιάσονται.
\vs{10}Ἕνεκεν Δαυὶδ τοῦ δούλου σου, μὴ ἀποστρέψῃς τὸ πρόσωπον τοῦ χριστοῦ σου.

\vs{11}Ὤμοσε Κύριος τῷ Δαυὶδ ἀλήθειαν, καὶ οὐ μὴ ἀθετήσει αὐτὴν, ἐκ καρποῦ τῆς κοιλίας σου, θήσομαι ἐπὶ τοῦ θρόνον σου.
\vs{12}Ἐὰν φυλάξωνται οἱ υἱοί σου τὴν διαθήκην μου, καὶ τὰ μαρτύριά μου ταῦτα ἃ διδάξω αὐτοὺς, καὶ οἱ υἱοὶ αὐτῶν ἕως τοῦ αἰῶνος καθιοῦνται ἐπὶ τοῦ θρόνου σου.
\vs{13}Ὅτι ἐξελέξατο Κύριος τὴν Σιὼν, ᾑρετίσατο αὐτὴν εἰς κατοικίαν ἑαυτῷ.
\vs{14}Αὕτη ἡ κατάπαυσίς μου εἰς αἰῶνα αἰῶνος, ὧδε κατοικήσω ὅτι ᾑρετισάμην αὐτήν.
\vs{15}Τὴν θήραν αὐτῆς εὐλογῶν εὐλογήσω, τοὺς πτωχοὺς αὐτῆς χορτάσω ἄρτων.
\vs{16}Τοὺς ἱερεῖς αὐτῆς ἐνδύσω σωτηρίαν, καὶ οἱ ὅσιοι αὐτῆς ἀγαλλιάσει ἀγαλλιάσονται.
\vs{17}Ἐκεῖ ἐξανατελῶ κέρας τῷ Δαυὶδ, ἡτοίμασα λύχνον τῷ χριστῷ σου.
\vs{18}Τοὺς ἐχθροὺς αὐτοῦ ἐνδύσω αἰσχύνην, ἐπὶ δὲ αὐτὸν ἐξανθήσει τὸ ἁγίασμά μου.

\begin{psalmheading}{\ch{132}{133} Ὠδὴ τῶν ἀναβαθμῶν.}
\end{psalmheading}
Ἰδοὺ δὴ τί καλὸν, ἢ τί τερπνὸν, ἀλλʼ ἢ τὸ κατοικεῖν ἀδελφοὺς ἐπιτοαυτό;
\vs{2}Ὡς μῦρον ἐπὶ κεφαλῆς τὸ καταβαῖνον ἐπὶ πώγωνα, τὸν πώγωνα τὸν Ἀαρὼν, τὸ καταβαῖνον ἐπὶ τὴν ὦαν τοῦ ἐνδύματος αὐτοῦ.
\vs{3}Ὡς δρόσος Ἀερμῶν ὁ καταβαίνουσα ἐπὶ τὰ ὄρη Σιών· ὅτι ἐκεῖ ἐνετείλατο Κύριος τὴν εὐλογίαν, ζωὴν ἕως τοῦ αἰῶνος.

\begin{psalmheading}{\ch{133}{134} Ὠδὴ τῶν ἀναβαθμῶν.}
\end{psalmheading}
Ἰδοὺ δὴ εὐλογεῖτε τὸν Κύριον πάντες οἱ δοῦλοι Κυρίου, οἱ ἑστῶτες ἐν οἴκῳ Κυρίου ἐν αὐλαῖς οἴκου Θεοῦ ἡμῶν·
\vs{2}ἐν ταῖς νυξὶν ἐπάρατε τὰς χεῖρας ὑμῶν εἰς τὰ ἅγια, καὶ εὐλογεῖτε τὸν Κύριον.
\vs{3}Εὐλογήσαι σε Κύριος ἐκ Σιὼν, ὁ ποιήσας τὸν οὐρανὸν καὶ τὴν γῆν.

\begin{psalmheading}{\ch{134}{135} Ἀλληλούϊα.}
\end{psalmheading}
Αἰνεῖτε τὸ ὄνομα Κυρίου, αἰνεῖτε δοῦλοι Κύριον,
\vs{2}Οἱ ἑστῶτες ἐν οἴκω Κυρίου, ἐν αὐλαῖς οἴκου Θεοῦ ἡμῶν.
\vs{3}Αἰνεῖτε τὸν Κύριον, ὅτι ἀγαθὸς Κύριος· ψάλατε τῷ ὀνόματι αὐτοῦ, ὅτι καλόν.

\vs{4}Ὅτι τὸν Ἰακὼβ ἐξελέξατο ἑαυτῷ ὁ Κύριος, Ἰσραὴλ εἰς περιουσιασμὸν ἑαυτῷ.
\vs{5}Ὅτι ἐγὼ ἔγνωκα, ὅτι μέγας ὁ Κύριος, καὶ ὁ Κύριος ἡμῶν παρὰ πάντας τοὺς θεούς.
\vs{6}Πάντα ὅσα ἠθέλησεν ὁ Κύριος, ἐποίησεν ἐν τῷ οὐρανῷ καὶ ἐν τῇ γῇ, ἐν ταῖς θαλάσσαις καὶ ἐν πάσαις ταῖς ἀβύσσοις.
\vs{7}Ἀνάγων νεφέλας ἐξ ἐσχάτου τῆς γῆς, ἀστραπὰς εἰς ὑετὸν ἐποίησεν· ὁ ἐξάγων ἀνέμους ἐκ θησαυρῶν αὐτοῦ.
\vs{8}Ὃς ἐπάταξε τὰ πρωτότοκα Αἰγύπτου ἀπὸ ἀνθρώπου ἕως κτήνους.
\vs{9}Ἐξαπέστειλεν σημεῖα καὶ τέρατα ἐν μέσῳ σου Αἴγυπτε, ἐν Φαραῷ καὶ ἐν πᾶσι τοῖς δούλοις αὐτοῦ.
\vs{10}Ὃς ἐπάταξεν ἔθνη πολλά, καὶ ἀπέκτεινε βασιλεῖς κραταιούς·
\vs{11}τὸν Σηὼν βασιλέα τῶν Ἀμοῤῥαίων, καὶ τὸν Ὢγ βασιλέα τῆς Βασὰν, καὶ πάσας τὰς βασιλείας Χαναάν·
\vs{12}Καὶ ἔδωκε τὴν γῆν αὐτῶν κληρονομίαν, κληρονομίαν Ἰσραὴλ λαῷ αὐτοῦ.

\vs{13}Κύριε τὸ ὄνομά σου εἰς τὸν αἰῶνα, καὶ τὸ μνημόσυνόν σου εἰς γενεὰν καὶ γενεάν.
\vs{14}Ὅτι κρινεῖ Κύριος τὸν λαὸν αὐτοῦ, καὶ ἐπὶ τοῖς δούλοις αὐτοῦ παρακληθήσεται.
\vs{15}Τὰ εἴδωλα τῶν ἐθνῶν ἀργύριον καὶ χρυσίον, ἔργα χειρῶν ἀνθρώπων.
\vs{16}Στόμα ἔχουσι καὶ οὐ λαλήσουσιν, ὀφθαλμοὺς ἔχουσι καὶ οὐκ ὄψονται·
\vs{17}Ὦτα ἔχουσι καὶ οὐκ ἐνωτισθήσονται, οὐδὲ γάρ ἐστι πνεῦμα ἐν τῷ στόματι αὐτῶν.
\vs{18}Ὅμοιοι αὐτοῖς γένοιντο οἱ ποιοῦντες αὐτὰ, καὶ πάντες οἱ πεποιθότες ἐπʼ αὐτοῖς.

\vs{19}Οἶκος Ἰσραὴλ εὐλογήσατε τὸν Κύριον, οἶκος Ἀαρὼν εὐλογήσατε τὸν Κύριον·
\vs{20}Οἶκος Λευὶ, εὐλογήσατε τὸν Κύριον, οἱ φοβούμενοι τὸν Κύριον εὐλογήσατε τὸν Κύριον.
\vs{21}Εὐλογητὸς Κύριος ἐκ Σιὼν, ὁ κατοικῶν Ἱερουσαλήμ.

\begin{psalmheading}{\ch{135}{136} Ἀλληλούϊα.}
\end{psalmheading}
Ἐξομολογεῖσθε τῷ Κυρίῳ, ὅτι ἀγαθὸς, ὅτι εἰς τὸν αἰῶνα τὸ ἔλεος αὐτοῦ.
\vs{2}Ἐξομολογεῖσθε τῷ Θεῷ τῶν θεῶν, ὅτι εἰς τὸν αἰῶνα τὸ ἔλεος αὐτοῦ.
\vs{3}Ἐξομολογεῖσθε τῷ Κυρίῳ τῶν κυρίων, ὅτι εἰς τὸν αἰῶνα τὸ ἔλεος αὐτοῦ.

\vs{4}Τῷ ποιήσαντι θαυμάσια μεγάλα μόνῳ, ὅτι εἰς τὸν αἰῶνα τὸ ἔλεος αὐτοῦ.
\vs{5}Τῷ ποιήσαντι τοὺς οὐρανοὺς ἐν συνέσει, ὅτι εἰς τὸν αἰῶνα τὸ ἔλεος αὐτοῦ.
\vs{6}Τῷ στερεώσαντι τὴν γῆν ἐπὶ τῶν ὑδάτων, ὅτι εἰς τὸν αἰῶνα τὸ ἔλεος αὐτοῦ.
\vs{7}Τῷ ποιήσαντι φῶτα μεγάλα μόνῳ, ὅτι εἰς τὸν αἰῶνα τὸ ἔλεος αὐτοῦ.
\vs{8}Τὸν ἥλιον εἰς ἐξουσίαν τῆς ἡμέρας, ὅτι εἰς τὸν αἰῶνα τὸ ἔλεος αὐτοῦ·
\vs{9}Τὴν σελήνην καὶ τοὺς ἀστέρας εἰς ἐξουσίαν τῆς νυκτός, ὅτι εἰς τὸν αἰῶνα τὸ ἔλεος αὐτοῦ.

\vs{10}Τῷ πατάξαντι Αἴγυπτον σὺν τοῖς πρωτοτόκοις αὐτῶν, ὅτι εἰς τὸν αἰῶνα τὸ ἔλεος αὐτοῦ.
\vs{11}Καὶ ἐξαγαγόντι τὸν Ἰσραὴλ ἐκ μέσου αὐτῶν, ὅτι εἰς τὸν αἰῶνα τὸ ἔλεος αὐτοῦ·
\vs{12}Ἐν χειρὶ κραταιᾷ καὶ ἐν βραχίονι ὑψηλῷ, ὅτι εἰς τὸν αἰῶνα τὸ ἔλεος αὐτοῦ.
\vs{13}Τῷ καταδιελόντι τὴν ἐρυθρὰν θάλασσαν εἰς διαιρέσεις, ὅτι εἰς τὸν αἰῶνα τὸ ἔλεος αὐτοῦ·
\vs{14}Καὶ διαγαγόντι τὸν Ἰσραὴλ διὰ μέσου αὐτῆς, ὅτι εἰς τὸν αἰῶνα τὸ ἔλεος αὐτοῦ·
\vs{15}Καὶ ἐκτινάξαντι Φαραὼ καὶ τὴν δύναμιν αὐτοῦ εἰς θάλασσαν ἐρυθράν, ὅτι εἰς τὸν αἰῶνα τὸ ἔλεος αὐτοῦ.
\vs{16}Τῷ διαγαγόντι τὸν λαὸν αὐτοῦ ἐν τῇ ἐρήμῳ, ὅτι εἰς τὸν αἰῶνα τὸ ἔλεος αὐτοῦ.

\vs{17}Τῷ πατάξαντι βασιλεῖς μεγάλους, ὅτι εἰς τὸν αἰῶνα τὸ ἔλεος αὐτοῦ.
\vs{18}Καὶ ἀποκτείναντι βασιλεῖς κραταιούς, ὅτι εἰς τὸν αἰῶνα τὸ ἔλεος αὐτοῦ.
\vs{19}Τὸν Σηὼν βασιλέα τῶν Ἀμοῤῥαίων, ὅτι εἰς τὸν αἰῶνα τὸ ἔλεος αὐτοῦ.
\vs{20}Καὶ τὸν Ὢγ βασιλέα τῆς Βασάν, ὅτι εἰς τὸν αἰῶνα τὸ ἔλεος αὐτοῦ.
\vs{21}Καὶ δόντι τὴν γῆν αὐτῶν κληρονομίαν, ὅτι εἰς τὸν αἰῶνα τὸ ἔλεος αὐτοῦ·
\vs{22}κληρονομίαν Ἰσραὴλ δούλῳ αὐτοῦ, ὅτι εἰς τὸν αἰῶνα τὸ ἔλεος αὐτοῦ.

\vs{23}Ὅτι ἐν τῇ ταπεινώσει ἡμῶν ἐμνήσθη ἡμῶν ὁ Κύριος, ὅτι εἰς τὸν αἰῶνα τὸ ἔλεος αὐτοῦ.
\vs{24}Καὶ ἐλυτρώσατο ἡμᾶς ἐκ τῶν ἐχθρῶν ἡμῶν, ὅτι εἰς τὸν αἰῶνα τὸ ἔλεος αὐτοῦ.
\vs{25}Ὁ διδοὺς τροφὴν πάσῃ σαρκὶ, ὅτι εἰς τὸν αἰῶνα τὸ ἔλεος αὐτοῦ.
\vs{26}Ἐξομολογεῖσθε τῷ Θεῷ τοῦ οὐρανοῦ, ὅτι εἰς τὸν αἰῶνα τὸ ἔλεος αὐτοῦ.

\begin{psalmheading}{\ch{136}{137} Τῷ Δαυὶδ, Ἱερεμίου.}
\end{psalmheading}
Ἐπὶ τῶν ποταμῶν Βαβυλῶνος ἐκεῖ ἐκαθίσαμεν, καὶ ἐκλαύσαμεν ἐν τῷ μνησθῆναι ἡμᾶς τῆς Σιών.
\vs{2}Ἐπὶ ταῖς ἰτέαις ἐν μέσῳ αὐτῆς ἐκρεμάσαμεν τὰ ὄργανα ἡμῶν.
\vs{3}Ὅτι ἐκεῖ ἐπηρώτησαν ἡμᾶς οἱ αἰχμαλωτεύσαντες ἡμᾶς, λόγους ᾠδῶν, καὶ οἱ ἀπαγαγόντες ἡμᾶς, ὕμνον· ᾄσατε ἡμῖν ἐκ τῶν ᾠδῶν Σιών.

\vs{4}Πῶς ᾄσωμεν τὴν ᾠδὴν Κυρίου ἐπὶ γῆς ἀλλοτρίας;
\vs{5}Ἐὰν ἐπιλάθωμαί σου Ἱερουσαλήμ, ἐπιλησθείη ἡ δεξιά μου.
\vs{6}Κολληθείη ἡ γλῶσσά μου τῷ λάρυγγί μου, ἐὰν μή σου μνησθῷ· ἐὰν μὴ προανατάξωμαι τὴν Ἱερουσαλὴμ ἐν ἀρχῇ τῆς εὐφροσύνης μου.

\vs{7}Μνήσθητι, Κύριε, τῶν υἱῶν Ἐδὼμ τὴν ἡμέραν Ἱερουσαλήμ· τῶν λεγόντων, ἐκκενοῦτε ἐκκενοῦτε, ἕως τῶν θεμελίων αὐτῇς.
\vs{8}Θυγάτηρ Βαβυλῶνος ἡ ταλαίπωρος, μακάριος ὃς ἀνταποδώσει σοι τὸ ἀνταπόδομά σου, ὃ ἀνταπέδωκας ἡμῖν.
\vs{9}Μακάριος ὃς κρατήσει καὶ ἐδαφιεῖ τὰ νήπιά σου πρὸς τὴν πέτραν.

\begin{psalmheading}{\ch{137}{138} Ψαλμὸς τῷ Δαυὶδ, Ἀγγαίου καὶ Ζαχαρίου.}
\end{psalmheading}
Ἐξομολογήσομαί σοι Κύριε ἐν ὅλῃ καρδίᾳ μου, καὶ ἐναντίον ἀγγέλων ψαλῶ σοι, ὅτι ἤκουσας πάντα τὰ ῥήματα τοῦ στόματός μου.
\vs{2}Προσκυνήσω πρὸς ναὸν ἅγιόν σου, καὶ ἐξομολογήσομαι τῷ ὀνόματί σου, ἐπὶ τῷ ἐλέει σου καὶ τῇ ἀληθείᾳ σου· ὅτι ἐμεγάλυνας ἐπὶ πᾶν τὸ ὄνομα τὸ ἅγιόν σου.
\vs{3}Ἐν ᾗ ἂν ἡμέρᾳ ἐπικαλέσωμαί σε, ταχὺ ἐπάκουσόν μου· πολυωρήσεις με ἐν ψυχῇ μου δυνάμει σου.
\vs{4}Ἐξομολογησάσθωσάν σοι Κύριε πάντες οἱ βασιλεῖς τῆς γῆς, ὅτι ἤκουσαν πάντα τὰ ῥήματα τοῦ στόματός σου.
\vs{5}Καὶ ᾀσάτωσαν ἐν ταῖς ὁδοῖς Κυρίου, ὅτι μεγάλη ἡ δόξα Κυρίου.

\vs{6}Ὅτι ὑψηλὸς Κύριος, καὶ τὰ ταπεινὰ ἐφορᾷ, καὶ τὰ ὑψηλὰ ἀπομακρόθεν γινώσκει.
\vs{7}Ἐὰν πορευθῶ ἐν μέσῳ θλίψεως, ζήσεις με· ἐπʼ ὀργὴν ἐχθρῶν μου ἐξέτεινας χεῖράς σου, καὶ ἔσωσέ με ἡ δεξιά σου.
\vs{8}Κύριε ἀνταποδώσεις ὑπὲρ ἐμοῦ· Κύριε τὸ ἔλεός σου εἰς τὸν αἰῶνα, τὰ ἔργα τῶν χειρῶν σου μὴ παρίδῃς.

\begin{psalmheading}{\ch{138}{139} Εἰς τὸ τέλος, ψαλμὸς τῷ Δαυίδ.}
\end{psalmheading}
Κύριε ἐδοκίμασάς με, καὶ ἔγνως με.
\vs{2}Σὺ ἔγνως τὴν καθέδραν μου, καὶ τὴν ἔγερσίν μου· σὺ συνῆκας τοὺς διαλογισμούς μου ἀπὸ μακρόθεν.
\vs{3}Τὴν τρίβον μου καὶ τὴν σχοῖνόν μου ἑξιχνίασας· καὶ πάσας τὰς ὁδούς μου προεῖδες,
\vs{4}ὅτι οὐκ ἔστι λόγος ἄδικος ἐν γλώσσῃ μου· ἰδοὺ Κύριε, σὺ ἔγνως πάντα
\vs{5}τὰ ἔσχατα καὶ τὰ ἀρχαῖα· σὺ ἔπλασάς με καὶ ἔθηκας ἐπʼ ἐμὲ τὴν χεῖρά σου.

\vs{6}Ἐθαυμαστώθη ἡ γνῶσίς σου ἐξ ἐμοῦ, ἐκραταιώθη, οὐ μὴ δύνωμαι πρὸς αὐτήν.
\vs{7}Ποῦ πορευθῶ ἀπὸ τοῦ πνεύματός σου, καὶ ἀπὸ τοῦ προσώπου σου ποῦ φύγω;
\vs{8}Ἐὰν ἀναβῶ εἰς τὸν οὐρανὸν, σὺ ἐκεῖ εἶ· ἐὰν καταβῶ εἰς τὸν ᾅδην, πάρει.
\vs{9}Ἐὰν ἀναλάβω τὰς πτέρυγάς μου κατʼ ὀρθὸν, καὶ κατασκηνώσω εἰς τὰ ἔσχατα τῆς θαλάσσης,
\vs{10}καὶ γὰρ ἐκεῖ ἡ χείρ σου ὁδηγήσει με, καὶ καθέξει με ἡ δεξιά σου.
\vs{11}Καὶ εἶπα, ἄρα σκότος καταπατήσει με, καὶ νὺξ φωτισμὸς ἐν τῇ τρυφῇ μου.
\vs{12}Ὅτι σκότος οὐ σκοτισθήσεται ἀπὸ σοῦ, καὶ νὺξ ὡς ἡμέρα φωτισθήσεται· ὡς τὸ σκότος αὐτῆς, οὕτως καὶ τὸ φῶς αὐτῆς.
\vs{13}Ὅτι σὺ ἐκτήσω τοὺς νεφρούς μου Κύριε, ἀντελάβου μου ἐκ γαστρὸς μητρός μου.
\vs{14}Ἐξομολογήσομαί σοι, ὅτι φοβερῶς ἐθαυμαστώθης· θαυμάσια τὰ ἔργα σου, καὶ ἡ ψυχή μου γινώσκει σφόδρα.
\vs{15}Οὐκ ἐκρύβη τὸ ὀστοῦν μου ἀπὸ σοῦ, ὃ ἐποίησας ἐν κρυφῇ, καὶ ἡ ὑπόστασίς μου ἐν τοῖς κατωτάτω τῆς γῆς.
\vs{16}Ἀκατέργαστόν μου εἶδου οἱ ἀφθαλμοί σου, καὶ ἐπὶ τὸ βιβλίον σου πάντες γραφήσονται· ἡμέρας πλασθήσονται καὶ οὐθεὶς ἐν αὐτοῖς.

\vs{17}Ἐμοὶ δὲ λίαν ἐτιμήθησαν οἱ φίλοι σου ὁ Θεὸς, λίαν ἐκραταιώθησαν αἱ ἀρχαὶ αὐτῶν.
\vs{18}Ἐξαριθμήσομαι αὐτοὺς καὶ ὑπὲρ ἄμμον πληθυνθήσονται· ἐξηγέρθην, καὶ ἔτι εἰμὶ μετὰ σοῦ.

\vs{19}Ἐὰν ἀποκτείνῃς ἁμαρτωλοὺς ὁ Θεός· ἄνδρες αἱμάτων ἐκκλίνατε ἀπʼ ἐμοῦ,
\vs{20}ὅτι ἐρεῖς εἰς διαλογισμόν· λήψονται εἰς ματαιότητα τὰς πόλεις σου.
\vs{21}Οὐχὶ τοὺς μισοῦντάς σε Κύριε ἐμίσησα, καὶ ἐπὶ τοὺς ἐχθρούς σου ἐξετηκόμην;
\vs{22}Τέλειον μῖσος ἐμίσουν αὐτοὺς, εἰς ἐχθροὺς ἐγένοντό μοι.
\vs{23}Δοκίμασόν με ὁ Θεὸς, καὶ γνῶθι τὴν καρδίαν μου· ἔτασόν με, καὶ γνῶθι τὰς τρίβους μου.
\vs{24}Καὶ ἴδε εἰ ὁδὸς ἀνομίας ἐν ἐμοὶ, καὶ ὁδήγησόν με ἐν ὁδῷ αἰωνίᾳ.

\begin{psalmheading}{\ch{139}{140} Εἰς τὸ τέλος, τῷ Δαυὶδ ψαλμός.}
\end{psalmheading}
\vs{2}Ἐξελοῦ με Κύριε ἐξ ἀνθρώπου πονηροῦ, ἀπὸ ἀνδρὸς ἀδίκου ῥῦσαί με·
\vs{3}οἵτινες ἐλογίσαντο ἀδικίας ἐν καρδίᾳ, ὅλην τὴν ἡμέραν παρετάσσοντο πολέμους.
\vs{4}Ἠκόνησαν γλῶσσαν αὐτῶν ὡσεὶ ὄφεως, ἰὸς ἀσπίδων ὑπὸ τὰ χείλη αὐτῶν· διάψαλμα.
\vs{5}Φύλαξόν με Κύριε ἐκ χειρὸς ἁμαρτωλοῦ, ἀπὸ ἀνθρώπων ἀδίκων ἐξελοῦ με· οἵτινες ἐλογίσαντο τοῦ ὑποσκελίσαι τὰ διαβήματά μου.
\vs{6}Ἔκρυψαν ὑπερήφανοι παγίδα μοι· καὶ σχοινία διέτειναν παγίδας τοῖς ποσί μου, ἐχόμενα τρίβου σκάνδαλον ἔθεντό μοι· διάψαλμα.

\vs{7}Εἶπα τῷ Κυρίῳ, Θεός μου εἶ σὺ· ἐνώτισαι, Κύριε, τὴν φωνὴν τῆς δεήσεώς μου.
\vs{8}Κύριε Κύριε, δύναμις τῆς σωτηρίας μου, ἐπεσκίασας ἐπὶ τὴν κεφαλήν μου ἐν ἡμέρᾳ πολέμου.
\vs{9}Μὴ παραδῷς με, Κύριε, ἀπὸ τῆς ἐπιθυμίας μου ἁμαρτωλῷ· διελογίσαντο κατʼ ἐμοῦ, μὴ ἐγκαταλίπῃς με, μή ποτε ὑψωθῶσι· διάψαλμα.

\vs{10}Ἡ κεφαλὴ τοῦ κυκλώματος αὐτῶν, κόπος τῶν χειλέων αὐτῶν καλύψει αὐτούς.
\vs{11}Πεσοῦνται ἐπʼ αὐτοὺς ἄνθρακες πυρὸς ἐπὶ τῆς γῆς, καὶ καταβαλεῖς αὐτοὺς ἐν ταλαιπωρίαις, οὐ μὴ ὑποστῶσιν.
\vs{12}Ἀνὴρ γλωσσώδης οὐ κατευθυνθήσεται ἐπὶ τῆς γῆς· ἄνδρα ἄδικον κακὰ θηρεύσει εἰς καταφθοράν.
\vs{13}Ἔγνων ὅτι ποιήσει Κύριος τὴν κρίσιν τοῦ πτωχοῦ, καὶ τὴν δίκην τῶν πενήτων.
\vs{14}Πλὴν δίκαιοι ἐξομολογήσονται τῷ ὀνόματί σου, κατοικήσουσιν εὐθεῖς σὺν τῷ προσώπῳ σου.

\begin{psalmheading}{\ch{140}{141} Ψαλμὸς τῷ Δαυίδ.}
\end{psalmheading}
Κύριε ἐκέκραξα πρὸς σὲ, εἰσάκουσόν μου· πρόσχες τῇ φωνῇ τῆς δεήσεώς μου, ἐν τῷ κεκραγέναι με πρὸς σέ.
\vs{2}Κατευθυνθήτω ἡ προσευχή μου ὡς θυμίαμα ἐνώπιόν σου· ἔπαρσις τῶν χειρῶν μου θυσία ἑσπερινή.
\vs{3}Θοῦ, Κύριε, φυλακὴν τῷ στόματί μου, καὶ θύραν περιοχῆς περὶ τὰ χείλη μου.
\vs{4}Μὴ ἐκκλίνῃς τὴν καρδίαν μου εἰς λόγους πονηρίας, τοῦ προφασίζεσθαι προφάσεις ἐν ἁμαρτίαις, σὺν ἀνθρώποις ἐργαζομένοις τὴν ἀνομίαν, καὶ οὐ μὴ συνδοιάσω μετὰ τῶν ἐκλεκτῶν αὐτῶν.
\vs{5}Παιδεύσει με δίκαιος ἐν ἐλέει καὶ ἐλέγξει με, ἔλαιον δὲ ἁμαρτωλοῦ μὴ λιπανάτω τὴν κεφαλήν μου, ὅτι ἔτι καὶ ἡ προσευχή μου ἐν ταῖς εὐδοκίαις αὐτῶν.

\vs{6}Κατεπόθησαν ἐχόμενα πέτρας οἱ κραταιοὶ αὐτῶν· ἀκούσονται τὰ ῥήματά μου, ὅτι ἠδύνθησαν.
\vs{7}Ὡσεὶ πάχος γῆς διεῤῥάγη ἐπὶ τῆς γῆς, διεσκορπίσθη τὰ ὀστᾶ ἡμῶν παρὰ τὸν ᾅδην.
\vs{8}Ὅτι πρὸς σὲ, Κύριε Κύριε, οἱ ὀφθαλμοί μου, ἐπὶ σοὶ ἤλπισα, μὴ ἀντανέλῃς τὴν ψυχήν μου.
\vs{9}Φύλαξόν με ἀπὸ παγίδος ἧς συνεστήσαντό μοι, καὶ ἀπὸ σκανδάλων τῶν ἐργαζομένων τὴν ἀνομίαν.
\vs{10}Πεσοῦνται ἐν ἀμφιβλήστρῳ αὐτοῦ ἁμαρτωλοὶ, καταμόνας εἰμὶ ἐγὼ ἕως οὗ ἂν παρέλθω.

\begin{psalmheading}{\ch{141}{142} Συνέσεως τῷ Δαυὶδ, ἐν τῷ εἶναι αὐτὸν ἐν τῷ σπηλαίῳ, προσευχή.}
\end{psalmheading}
\vs{2}Φωνῇ μου πρὸς Κύριον ἐκέκραξα, φωνῇ μου πρὸς Κύριον ἐδεήθην.
\vs{3}Ἐκχεῶ ἐναντίον αὐτοῦ τὴν δέησίν μου, τὴν θλίψιν μου ἐνώπιον αὐτοῦ ἀπαγγελῶ.
\vs{4}Ἐν τῷ ἐκλείπειν ἐξ ἐμοῦ τὸ πνεῦμά μου, καὶ σὺ ἔγνως τὰς τρίβους μου· ἐν ὁδῷ ταύτῃ ᾗ ἐπορευόμην, ἔκρυψαν παγίδα μοι.
\vs{5}Κατενόουν εἰς τὰ δεξιὰ καὶ ἐπέβλεπον, ὅτι οὐκ ἦν ὁ ἐπιγινώσκων με· ἀπώλετο φυγὴ ἀπʼ ἐμοῦ, καὶ οὐκ ἔστιν ὁ ἐκζητῶν τὴν ψυχήν μου.
\vs{6}Πρὸς σὲ, Κύριε, ἐκέραξα, καὶ εἶπα, σὺ εἶ ἡ ἐλπίς μου, μερίς μου ἐν γῇ ζώντων.
\vs{7}Πρόσχες πρὸς τὴν δέησίν μου, ὅτι ἐταπεινώθην σφόδρα· ῥῦσαί με ἐκ τῶν καταδιωκόντων με, ὅτι ἐκραταιώθησαν ὑπὲρ ἐμέ.
\vs{8}Ἐξάγαγε ἐκ φυλακῆς τὴν ψυχήν μου, τοῦ ἐξομολογήσασθαι τῷ ὀνόματί σου, Κύριε· ἐμὲ ὑπομενοῦσι δίκαιοι, ἕως οὗ ἀνταποδῷς μοι.

\begin{psalmheading}{\ch{142}{143} Ψαλμὸς τῷ Δαυὶδ, ὅτε αὐτὸν ὁ υἱὸς κατεδιώκει.}
\end{psalmheading}
Κύριε εἰσάκουσον τῆς προσευχῆς μου, ἐνώτισαι τὴν δέησίν μου ἐν τῇ ἀληθείᾳ σου, ἐπάκουσόν μου ἐν τῇ δικαιοσύνῃ σου.
\vs{2}Καὶ μὴ εἰσέλθῃς εἰς κρίσιν μετὰ τοῦ δούλου σου, ὅτι οὐ δικαιωθήσεται ἐνώπιόν σου πᾶς ζῶν.

\vs{3}Ὅτι κατεδίωξεν ὁ ἐχθρὸς τὴν ψυχήν μου· ἐταπείνωσεν εἰς τὴν γῆν τὴν ζωήν μου, ἐκάθισέ με ἐν σκοτεινοῖς ὡς νεκροὺς αἰῶνος,
\vs{4}καὶ ἠκηδίασεν ἐπʼ ἐμὲ τὸ πνεῦμά μου· ἐν ἐμοὶ ἐταράχθη ἡ καρδία μου.
\vs{5}Ἐμνήσθην ἡμερῶν ἀρχαίων· καὶ ἐμελέτησα ἐν πᾶσι τοῖς ἔργοις σου, ἐν ποιήμασι τῶν χειρῶν σου ἐμελέτων.
\vs{6}Διεπέτασα πρὸς σὲ τὰς χεῖράς μου, ἡ ψυχή μου ὡς γῆ ἄνυδρός σοι· διάψαλμα.

\vs{7}Ταχὺ εἰσάκουσόν μου, Κύριε, ἐξέλιπε τὸ πνεῦμά μου· μὴ ἀποστρέψῃς τὸ πρόσωπόν σου ἀπʼ ἐμοῦ, καὶ ὁμοιωθήσομαι τοῖς καταβαίνουσιν εἰς λάκκον.
\vs{8}Ἀκουστὸν ποίησόν μοι τοπρωῒ τὸ ἔλεός σου, ὅτι ἐπὶ σοὶ ἤλπισα· γνώρισόν μοι, Κύριε, ὁδὸν ἐν ᾗ πορεύσομαι, ὅτι πρὸς σὲ ᾖρα τὴν ψυχήν μου.
\vs{9}Ἐξελοῦ με ἐκ τῶν ἐχθρῶν μου Κύριε, ὅτι πρὸς σὲ κατέφυγον.
\vs{10}Δίδαξόν με τοῦ ποιεῖν τὸ θέλημά σου, ὅτι Θεός μου εἶ σὺ, τὸ πνεῦμά σου τὸ ἁγαθὸν ὁδηγήσει με ἐν τῇ εὐθείᾳ.
\vs{11}Ἕνεκα τοῦ ὀνόματός σου, Κύριε, ζήσεις με, ἐν τῇ δικαιοσύνῃ σου ἐξάξεις ἐκ θλίψεως τὴν ψυχήν μου.
\vs{12}Καὶ ἐν τῷ ἐλέει σου ἐξολοθρεύσεις τοὺς ἐχθρούς μου, καὶ ἀπολεῖς πάντας τοὺς θλίβοντας τὴν ψυχήν μου, ὅτι δοῦλός σου εἰμὶ ἐγώ.

\begin{psalmheading}{\ch{143}{144} Τῷ Δαυὶδ πρὸς τὸν Γολιάδ.}
\end{psalmheading}
Εὐλογήτος Κύριος ὁ Θεός μου, ὁ διδάσκων τὰς χεῖράς μου εἰς παράταξιν, τοὺς δακτύλους μου εἰς πόλεμον·
\vs{2}Ἔλεός μου καὶ καταφυγή μου, ἀντιλήπτωρ μου καὶ ῥύστης μου, ὑπερασπιστής μου, καὶ ἐπʼ αὐτῷ ἤλπισα, ὁ ὑποτάσσων τὸν λαόν μου ὑπʼ ἐμέ.

\vs{3}Κύριε, τί ἐστιν ἄνθρωπος, ὅτι ἐγνώσθης αὐτῷ; ἢ υἱὸς ἀνθρώπου, ὅτι λογίζῃ αὐτόν;
\vs{4}Ἄνθρωπος ματαιότητι ὡμοιώθη, αἱ ἡμέραι αὐτοῦ ὡσεὶ σκιὰ παράγουσι.

\vs{5}Κύριε, κλῖνον οὐρανούς σου καὶ κατάβηθι, ἅψαι τῶν ὀρέων καὶ καπνισθήσονται.
\vs{6}Ἄστραψον ἀστραπὴν καὶ σκορπιεῖς αὐτοὺς, ἐξαπόστειλον τὰ βέλη σου καὶ συνταράξεις αὐτούς.
\vs{7}Ἐξαπόστειλον τὴν χεῖρά σου ἐξ ὕψους, ἐξελοῦ με καὶ ῥῦσαί με ἐξ ὑδάτων πολλῶν, ἐκ χειρὸς υἱῶν ἀλλοτρίων·
\vs{8}ὧν τὸ στόμα ἐλάλησε ματαιότητα, καὶ ἡ δεξιὰ αὐτῶν δεξιὰ ἀδικίας.

\vs{9}Ὁ Θεὸς, ᾠδὴν καινὴν ᾄσομαί σοι, ἐν ψαλτηρίῳ δεκαχόρδῳ ψαλῶ σοι·
\vs{10}Τῷ διδόντι τὴν σωτηρίαν τοῖς βασιλεῦσι, τῷ λυτρουμένῳ Δαυὶδ τὸν δοῦλον αὐτοῦ ἐκ ῥομφαίας πονηρᾶς.
\vs{11}Ῥῦσαί με καὶ ἐξελοῦ με ἐκ χειρὸς υἱῶν ἀλλοτρίων, ὧν τὸ στόμα ἐλάλησε ματαιότητα, καὶ ἡ δεξιὰ αὐτῶν δεξιὰ ἀδικίας·
\vs{12}ὧν οἱ υἱοὶ ὡς νεόφυτα ἱδρυμένα ἐν τῇ νεότητι αὐτῶν· αἱ θυγατέρες αὐτῶν κεκαλλωπισμέναι, περικεκοσμημέναι ὡς ὁμοίωμα ναοῦ.
\vs{13}Τὰ ταμεῖα αὐτῶν πλήρη, ἐξερευγόμενα ἐκ τούτου εἰς τοῦτο· τὰ πρόβατα αὐτῶν πολυτόκα, πληθύνοντα ἐν ταῖς ἐξόδοις αὐτῶν·
\vs{14}Οἱ βόες αὐτῶν παχεῖς· οὐκ ἔστι κατάπτωμα φραγμοῦ, οὐδὲ διέξοδος, οὐδὲ κραυγὴ ἐν ταῖς ἐπαύλεσιν αὐτῶν.
\vs{15}Ἐμακάρισαν τὸν λαὸν ᾧ ταῦτά ἐστι· μακάριος ὁ λαὸς οὗ Κύριος ὁ Θεὸς αὐτοῦ.

\begin{psalmheading}{\ch{144}{145} Αἴνεσις τοῦ Δαυίδ.}
\end{psalmheading}
Ὑψώσω σε, ὁ Θεός μου ὁ βασιλεύς μου, καὶ εὐλογήσω τὸ ὄνομά σου εἰς τὸν αἰῶνα καὶ εἰς τὸν αἰῶνα τοῦ αἰῶνος.
\vs{2}Καθʼ ἑκάστην ἡμέραν εὐλογήσω σε, καὶ αἰνέσω τὸ ὄνομά σου εἰς τὸν αἰῶνα καὶ εἰς τὸν αἰῶνα τοῦ αἰῶνος.
\vs{3}Μέγας ὁ Κύριος καὶ αἰνετὸς σφόδρα, καὶ τῆς μεγαλωσύνης αὐτοῦ οὐκ ἔστι πέρας.
\vs{4}Γενεὰ καὶ γενεὰ ἐπαινέσει τὰ ἔργα σου, καὶ τὴν δύναμίν σου ἀπαγγελοῦσι.
\vs{5}Καὶ τὴν μεγαλοπρέπειαν τῆς δόξης τῆς ἁγιωσύνης σου λαλήσουσι, καὶ τὰ θαυμάσιά σου διηγήσονται.
\vs{6}Καὶ τὴν δύναμιν τῶν φοβερῶν σου ἐροῦσι, καὶ τὴν μεγαλωσύνην σου διηγήσονται.
\vs{7}Μνήμην τοῦ πλήθους τῆς χρηστότητός σου ἐξερεύξονται, καὶ τῇ δικαιοσύνῃ σου ἀγαλλιάσονται.

\vs{8}Οἰκτίρμων καὶ ἐλεήμων ὁ Κύριος, μακρόθυμος καὶ πολυέλεος.
\vs{9}Χρηστὸς Κύριος τοῖς ὑπομένουσι, καὶ οἱ οἰκτιρμοὶ αὐτοῦ ἐπὶ πάντα τὰ ἔργα αὐτοῦ.
\vs{10}Ἐξομολογησάσθωσάν σοι, Κύριε, πάντα τὰ ἔργα σου, καὶ οἱ ὅσιοί σου εὐλογησάτωσάν σε.
\vs{11}Δόξαν τῆς βασιλείας σου ἐροῦσι, καὶ τὴν δυναστείαν σου λαλήσουσιν·
\vs{12}Τοῦ γνωρίσαι τοῖς υἱοῖς τῶν ἀνθρώπων τὴν δυναστείαν σου, καὶ τὴν δόξαν τῆς μεγαλοπρεπείας τῆς βασιλείας σου.
\vs{13}Ἡ βασιλεία σου βασιλεία πάντων τῶν αἰώνων, καὶ ἡ δεσποτία σου ἐν πάσῃ γενεᾷ καὶ γενεᾷ·
\vs{13a}πιστὸς Κύριος ἐν τοῖς λόγοις αὐτοῦ, καὶ ὅσιος ἐν πᾶσι τοῖς ἔργοις αὐτοῦ.

\vs{14}Ὑποστηρίζει Κύριος πάντας τοὺς καταπίπτοντας, καὶ ἀνορθοῖ πάντας τοὺς κατεῤῥαγμένους.
\vs{15}Οἱ ὀφθαλμοὶ πάντων εἰς σὲ ἐλπίζουσι, καὶ σὺ δίδως τὴν τροφὴν αὐτῶν ἐν εὐκαιρίᾳ.
\vs{16}Ἀνοίγεις σὺ τὰς χεῖράς σου, καὶ ἐμπιπλᾷς πᾶν ζῶον εὐδοκίας.
\vs{17}Δίκαιος Κύριος ἐν πάσαις ταῖς ὁδοῖς αὐτοῦ, καὶ ὅσιος ἐν πᾶσι τοῖς ἔργοις αὐτοῦ.

\vs{18}Ἐγγὺς Κύριος πᾶσιν τοῖς ἐπικαλουμένοις αὐτόν, πᾶσι τοῖς ἐπικαλουμένοις αὐτὸν ἐν ἀληθείᾳ.
\vs{19}Θέλημα τῶν φοβουμένων αὐτὸν ποιήσει, καὶ τῆς δεήσεως αὐτῶν ἐπακούσεται, καὶ σώσει αὐτούς.
\vs{20}Φυλάσσει Κύριος πάντας τοὺς ἀγαπῶντας αὐτόν, καὶ πάντας τοὺς ἁμαρτωλοὺς ἐξολοθρεύσει.
\vs{21}Αἴνεσιν Κυρίου λαλήσει τὸ στόμα μου, καὶ εὐλογείτω πᾶσα σὰρξ τὸ ὄνομα τὸ ἅγιον αὐτοῦ, εἰς τὸν αἰῶνα καὶ εἰς τὸν αἰῶνα τοῦ αἰῶνος.

\begin{psalmheading}{\ch{145}{146} Ἀλληλούϊα· Ἀγγαίου καὶ Ζαχαρίου.}
\end{psalmheading}
Αἰνεῖ ἡ ψυχή μου τὸν Κύριον.
\vs{2}Αἰνέσω Κύριον ἐν ζωῇ μοῦ, ψαλῶ τῷ Θεῷ μου ἕως ὑπάρχω.
\vs{3}Μὴ πεποίθατε ἐπʼ ἄρχοντας, καὶ ἐφʼ υἱοὺς ἀνθρώπων, οἷς οὐκ ἔστι σωτηρία.
\vs{4}Ἐξελεύσεται τὸ πνεῦμα αὐτοῦ καὶ ἐπιστρέψει εἰς τὴν γῆν αὐτοῦ, ἐν ἐκείνῃ τῇ ἡμέρᾳ ἀπολοῦνται πάντες οἱ διαλογισμοὶ αὐτῶν.

\vs{5}Μακάριος οὗ ὁ Θεὸς Ἰακὼβ βοηθός αὐτοῦ, ἡ ἐλπὶς αὐτοῦ ἐπὶ Κύριον τὸν Θεὸν αὐτοῦ·
\vs{6}Τὸν ποιήσαντα τὸν οὐρανὸν καὶ τὴν γῆν, τὴν θάλασσαν καὶ πάντα τὰ ἐν αὐτοῖς· τὸν φυλάσσοντα ἀλήθειαν εἰς τὸν αἰῶνα,
\vs{7}ποιοῦντα κρίμα τοῖς ἀδικουμένοις, διδόντα τροφὴν τοῖς πεινῶσι. Κύριος λύει πεπεδημένους,
\vs{8}Κύριος σοφοῖ τυφλοὺς,

Κύριος ἀνορθοῖ κατεῤῥαγμένους, Κύριος ἀγαπᾷ δικαίους,
\vs{9}Κύριος φυλάσσει τοὺς προσηλύτους· ὀρφανὸν καὶ χήραν ἀναλήμψεται, καὶ ὁδὸν ἁμαρτωλῶν ἀφανιεῖ.
\vs{10}Βασιλεύσει Κύριος εἰς τὸν αἰῶνα, ὁ Θεός σου, Σιὼν, εἰς γενεὰν καὶ γενεάν.

\begin{psalmheading}{\ch{146}{147:1-11} Ἀλληλούϊα· Ἀγγαίου καὶ Ζαχαρίου.}
\end{psalmheading}
Αἰνεῖτε τὸν Κύριον ὅτι ἀγαθὸν ψαλμός, τῷ Θεῷ ἡμῶν ἡδυνθείη αἴνεσις.
\vs{2}Οἰκοδομῶν Ἱερουσαλὴμ ὁ Κύριος, καὶ τὰς διασπορὰς τοῦ Ἰσραὴλ ἐπισυνάξει·
\vs{3}Ὁ ἰώμενος τοὺς συντετριμμένους τὴν καρδίαν, καὶ δεσμεύων τὰ συντρίμματα αὐτῶν·
\vs{4}Ὁ ἀριθμῶν πλήθη ἄστρων, καὶ πᾶσιν αὐτοῖς ὀνόματα καλῶν.
\vs{5}Μέγας ὁ Κύριος ἡμῶν, καὶ μεγάλη ἡ ἰσχὺς αὐτοῦ, καὶ τῆς συνέσεως αὐτοῦ οὐκ ἔστιν ἀριθμός.
\vs{6}Ἀναλαμβάνων πρᾳεῖς ὁ Κύριος, ταπεινῶν δὲ ἁμαρτωλοὺς ἕως τῆς γῆς.

\vs{7}Ἐξάρξατε τῷ Κυρίῳ ἐν ἐξομολογήσει, ψάλατε τῷ Θεῷ ἡμῶν ἐν κιθάρᾳ·
\vs{8}Τῷ περιβάλλοντι τὸν οὐρανὸν ἐν νεφέλαις, τῷ ἑτοιμάζοντι τῇ γῇ ὑετόν· τῷ ἐξανατέλλοντι ἐν ὄρεσι χόρτον, καὶ χλόην τῇ δουλείᾳ τῶν ἀνθρώπων·
\vs{9}καὶ διδόντι τοῖς κτήνεσι τροφὴν αὐτῶν, καὶ τοῖς νεοσσοῖς τῶν κοράκων τοῖς ἐπικαλουμένοις αὐτόν.
\vs{10}Οὐκ ἐν τῇ δυναστείᾳ τοῦ ἵππου θελήσει, οὐδὲ ἐν ταῖς κνήμαις τοῦ ἀνδρὸς εὐδοκεῖ.
\vs{11}Εὐδοκεῖ Κύριος ἐν τοῖς φοβουμένοις αὐτὸν, καὶ ἐν πᾶσιν τοῖς ἐλπιζουσιν ἐπὶ τὸ ἔλεος αὐτοῦ.

\begin{psalmheading}{\ch{147}{147:12-20} Ἀλληλούϊα· Ἀγγαίου καὶ Ζαχαρίου.}
\end{psalmheading}
Ἐπαίνει, Ἱερουσαλήμ, τὸν Κύριον, αἴνει τὸν Θεόν σου Σιών.
\vs{2}Ὅτι ἐνίσχυσεν τοὺς μοχλοὺς τῶν πυλῶν σου, εὐλόγησεν τοὺς υἱούς σου ἐν σοί.
\vs{3}Ὁ τιθεὶς τὰ ὅριά σου εἰρήνην, καὶ στέαρ πυροῦ ἐμπιπλῶν σε·
\vs{4}Ὁ ἀποστέλλων τὸ λόγιον αὐτοῦ τῇ γῇ, ἕως τάχους δραμεῖται ὁ λόγος αὐτοῦ·
\vs{5}Τοῦ διδόντος χιόνα ὡσεὶ ἔριον, ὁμίχλην ὡσεὶ σποδὸν πάσσοντος·
\vs{6}Βάλλοντος κρύσταλλον αὐτοῦ ὠσεὶ ψωμούς· κατὰ πρόσωπον ψύχους αὐτοῦ τίς ὑποστήσεται;
\vs{7}Ἀποστελεῖ τὸν λόγον αὐτοῦ, καὶ τήξει αὐτά, πνεύσει τὸ πνεῦμα αὐτοῦ, καὶ ῥυήσεται ὕδατα.
\vs{8}Ἀπαγγέλλων τὸν λόγον αὐτοῦ τῷ Ἰακώβ, δικαιώματα καὶ κρίματα αὐτοῦ τῷ Ἰσραήλ.
\vs{9}Οὐκ ἐποίησεν οὕτως παντὶ ἔθνει, καὶ τὰ κρίματα αὐτοῦ οὐκ ἐδήλωσεν αὐτοῖς.

\begin{psalmheading}{\ch{148}{148} Ἀλληλούϊα· Ἀγγαίου καὶ Ζαχαρίου.}
\end{psalmheading}
Αἰνεῖτε τὸν Κύριον ἐκ τῶν οὐρανῶν, αἰνεῖτε αὐτὸν ἐν τοῖς ὑψίστοις.
\vs{2}Αἰνεῖτε αὐτόν πάντες οἱ ἄγγελοι αὐτοῦ, αἰνεῖτε αὐτόν πᾶσαι αἱ δυνάμεις αὐτοῦ.
\vs{3}Αἰνεῖτε αὐτόν ἥλιος καὶ σελήνη, αἰνεῖτε αὐτόν πάντα τὰ ἄστρα καὶ τὸ φῶς.
\vs{4}Αἰνεῖτε αὐτόν οἱ οὐρανοὶ τῶν οὐρανῶν, καὶ τὸ ὕδωρ τὸ ὑπεράνω τῶν οὐρανῶν.
\vs{5}Αἰνεσάτωσαν τὸ ὄνομα Κυρίου· ὅτι αὐτὸς εἶπεν καὶ ἐγενήθησαν, αὐτὸς ἐνετείλατο καὶ ἐκτίσθησαν.
\vs{6}Ἔστησεν αὐτὰ εἰς τὸν αἰῶνα, καὶ εἰς τὸν αἰῶνα τοῦ αἰῶνος· πρόσταγμα ἔθετο, καὶ οὐ παρελεύσεται.

\vs{7}Αἰνεῖτε τὸν Κύριον ἐκ τῆς γῆς, δράκοντες καὶ πᾶσαι ἄβυσσοι·
\vs{8}Πῦρ, χάλαζα, χιὼν, κρύσταλλος, πνεῦμα καταιγίδος, τὰ ποιοῦντα τὸν λόγον αὐτοῦ·
\vs{9}Τὰ ὄρη καὶ πάντες βουνοί, ξύλα καρποφόρα καὶ πᾶσαι κέδροι·
\vs{10}Τὰ θηρία καὶ πάντα τὰ κτήνη, ἑρπετὰ καὶ πετεινὰ πτερωτά·
\vs{11}Βασιλεῖς τῆς γῆς καὶ πάντες λαοί, ἄρχοντες καὶ πάντες κριταὶ γῆς·
\vs{12}Νεανίσκοι καὶ παρθένοι, πρεσβύται μετὰ νεωτέρων
\vs{13}αἰνεσάτωσαν τὸ ὄνομα Κυρίου, ὅτι ὑψώθη τὸ ὄνομα αὐτοῦ μόνου· ἡ ἐξομολόγησις αὐτοῦ ἐπὶ γῆς καὶ οὐρανοῦ,
\vs{14}καὶ ὑψώσει κέρας λαοῦ αὐτοῦ· ὕμονς πᾶσι τοῖς ὁσίοις αὐτοῦ, τοῖς υἱοῖς Ἰσραήλ, λαῷ ἐγγίζοντι αὐτῷ.

\begin{psalmheading}{\ch{149}{149} Ἀλληλούϊα.}
\end{psalmheading}
Ἄσατε τῷ Κυρίῳ ᾆσμα καινόν· ἡ αἴνεσις αὐτοῦ ἐν ἐκκλησίᾳ ὁσίων.
\vs{2}Εὐφρανθήτω Ἰσραὴλ ἐπὶ τῷ ποιήσαντι αὐτὸν, καὶ υἱοὶ Σιὼν ἀγαλλιάσθωσαν ἐπὶ τῷ βασιλεῖ αὐτῶν.
\vs{3}Αἰνεσάτωσαν τὸ ὄνομα αὐτοῦ ἐν χορῷ, ἐν τυμπάνῳ καὶ ψαλτηρίῳ ψαλάτωσαν αὐτῷ.
\vs{4}Ὅτι εὐδοκεῖ Κύριος ἐν λαῷ αὐτοῦ, καὶ ὑψώσει πρᾳεῖς ἐν σωτηρίᾳ.

\vs{5}Καυχήσονται ὅσιοι ἐν δόξῃ, καὶ ἀγαλλιάσονται ἐπὶ τῶν κοιτῶν αὐτῶν.
\vs{6}Αἱ ὑψώσεις τοῦ Θεοῦ ἐν λάρυγγι αὐτῶν, καὶ ῥομφαῖαι δίστομοι ἐν ταῖς χερσὶν αὐτῶν·
\vs{7}Τοῦ ποιῆσαι ἐκδίκησιν ἐν τοῖς ἔθνεσιν, ἐλεγμοὺς ἐν τοῖς λαοῖς·
\vs{8}Τοῦ δῆσαι τοὺς βασιλεῖς αὐτῶν ἐν πέδαις, καὶ τοὺς ἐνδόξους αὐτῶν ἐν χειροπέδαις σιδηραῖς·
\vs{9}Τοῦ ποιῆσαι ἐν αὐτοῖς κρίμα ἔγγραπτον· δόξα αὕτη ἐστὶ πᾶσι τοῖς ὁσίοις αὐτοῦ.

\begin{psalmheading}{\ch{150}{150} Ἁλληλούϊα.}
\end{psalmheading}
Αἰνεῖτε τὸν Θεὸν ἐν τοῖς ἁγίοις αὐτοῦ, αἰνεῖτε αὐτὸν ἐν στερεώματι δυνάμεως αὐτοῦ.
\vs{2}Αἰνεῖτε αὐτὸν ἐπὶ ταῖς δυναστείαις αὐτοῦ, αἰνεῖτε αὐτὸν κατὰ τὸ πλῆθος τῆς μεγαλωσύνης αὐτοῦ.
\vs{3}Αἰνεῖτε αὐτὸν ἐν ἤχῳ σάλπιγγος, αἰνεῖτε αὐτὸν ἐν ψαλτηρίῳ καὶ κιθάρᾳ.
\vs{4}Αἰνεῖτε αὐτὸν ἐν τυμπάνῳ καὶ χορῷ, αἰνεῖτε αὐτὸν ἐν χορδαῖς καὶ ὀργάνῳ.
\vs{5}Αἰνεῖτε αὐτὸν ἐν κυμβάλοις εὐήχοις, αἰνεῖτε αὐτὸν ἐν κυμβάλοις ἀλαλαγμοῦ.
\vs{6}Πᾶσα πνοὴ αἰνεσάτω τὸν Κύριον.

\begin{psalmheading}{\ch{151}{} Οὗτος ὁ ψαλμὸς ἰδιόγραφος εἰς Δαυὶδ, καὶ ἔξωθεν τοῦ ἀριθμοῦ, ὅτε ἐμονομάχησε τῷ Γολιάδ.}
\end{psalmheading}
Μικρὸς ἤμην ἐν τοῖς ἀδελφοίς μου, καὶ νεώτερος ἐν τῷ οἴκῳ τοῦ πατρός μου, ἐποίμαινον τὰ πρόβατα τοῦ πατρός μου.
\vs{2}Αἱ χεῖρές μου ἐποίησαν ὄργανον, καὶ οἱ δάκτυλοί μου ἥρμοσαν ψαλτήριον.
\vs{3}Καὶ τίς ἀναγγελεῖ τῷ Κυρίῳ μου; αὐτὸς Κύριος, αὐτὸς εἰσακούει.
\vs{4}Αὐτὸς ἐξαπέστειλεν τὸν ἄγγελον αὐτοῦ, καὶ ᾖρέ με ἐκ τῶν προβάτων τοῦ πατρός μου, καὶ ἔχρισέ με ἐν τῷ ἐλαίῳ τῆς χρίσεως αὐτοῦ.
\vs{5}Οἱ ἀδελφοί μου καλοὶ καὶ μεγάλοι, καὶ οὐκ εὐδόκησεν ἐν αὐτοῖς Κύριος.
\vs{6}Ἐξῆλθον εἰς συνάντησιν τῷ ἀλλοφύλῳ, καὶ ἐπικατηράσατό με ἐν τοῖς εἰδώλοις αὐτοῦ·
\vs{7}Ἐγὼ δὲ σπασάμενος τὴν παρʼ αὐτοῦ μάχαιραν, ἀπεκεφάλισα αὐτόν, καὶ ᾖρα ὄνειδος ἐξ υἱῶν Ἰσραήλ.
\end{psalms}


\def\book{ΠΑΡΟΙΜΙΑΙ ΣΑΛΩΜΩΝΤΟΣ}
\biblebook{ΠΑΡΟΙΜΙΑΙ ΣΑΛΩΜΩΝΤΟΣ}


\lettrine[lines=2, loversize=0.2, nindent=0em, findent=.25em]{\textcolor{bookheadingcolor}{Π}}{ΑΡΟΙΜΙΑΙ} Σαλωμῶντος υἱοῦ Δαυὶδ, ὃς ἐβασίλευσεν ἐν Ἰσραήλ·
\vs{2}γνῶναι σοφίαν καὶ παιδείαν, νοῆσαί τε λόγους φρονήσεως,
\vs{3}δέξασθαί τε στροφὰς λόγων, νοῆσαί τε δικαιοσύνην ἀληθῆ, καὶ κρίμα κατευθύνειν·
\vs{4}Ἵνα δῷ ἀκάκοις πανουργίαν, παιδὶ δὲ νέῳ αἴσθησίν τε καὶ ἔννοιαν.
\vs{5}Τῶν δὲ γὰρ ἀκούσας σοφὸς σοφώτερος ἔσται, ὁ δὲ νοήμων κυβέρνησιν κτήσεται·
\vs{6}Νοήσει τε παραβολὴν καὶ σκοτεινὸν λόγον, ῥήσεις τε σοφῶν καὶ αἰνίγματα.

\vs{7}Ἀρχὴ σοφίας φόβος Κυριου, σύνεσις δὲ ἀγαθὴ πᾶσι τοῖς ποιοῦσιν αὐτήν· εὐσέβεια δὲ εἰς Θεὸν ἀρχὴ αἰσθήσεως, σοφίαν δὲ καὶ παιδείαν ἀσεβεῖς ἐξουθενήσουσιν.
\vs{8}Ἄκουε υἱὲ παιδείαν πατρός σου, καὶ μὴ ἀπώσῃ θεσμοὺς μητρός σου.
\vs{9}Στέφανον γὰρ χαρίτων δέξῃ σῇ κορυφῇ, καὶ κλοιὸν χρύσεον περὶ σῷ τραχήλῳ.

\vs{10}Υἱὲ μή σε πλανήσωσιν ἄνδρες ἀσεβεῖς, μηδὲ βουληθῇς.
\vs{11}Ἐὰν παρακαλέσωσί σε, λέγοντες, ἐλθὲ μεθʼ ἡμῶν, κοινώνησον αἵματος, κρύψωμεν δὲ εἰς γῆν ἄνδρα δίκαιον ἀδίκως,
\vs{12}καταπίωμεν δὲ αὐτὸν ὥσπερ ᾅδης ζῶντα, καὶ ἄρωμεν αὐτοῦ τὴν μνήμην ἐκ γῆς,
\vs{13}τὴν κτῆσιν αὐτοῦ τὴν πολυτελῆ καταλαβώμεθα, πλήσωμεν δὲ οἴκους ἡμετέρους σκύλων·
\vs{14}Τὸν δὲ σὸν κλῆρον βάλε ἐν ἡμῖν, κοινὸν δὲ βαλάντιον κτησώμεθα πάντες, καὶ μαρσίππιον ἓν γενηθήτω ἡμῖν·
\vs{15}Μὴ πορευθῇς ἐν ὁδῷ μετʼ αὐτῶν, ἔκκλινον δὲ τὸν πόδα σου ἐκ τῶν τρίβων αὐτῶν.
\vs{17}Οὐ γὰρ ἀδίκως ἐκτείνεται δίκτυα πτερωτοῖς.
\vs{18}Αὐτοὶ γὰρ οἱ φόνου μετέχοντες, θησαυρίζουσιν ἑαυτοῖς κακά· ἡ δὲ καταστροφὴ ἀνδρῶν παρανόμων κακή.
\vs{19}Αὗται αἱ ὁδοί εἰσι πάντων τῶν συντελούντων τὰ ἄνομα· τῇ γὰρ ἀσεβείᾳ τὴν ἑαυτῶν ψυχὴν ἀφαιροῦνται.

\vs{20}Σοφία ἐν ἐξόδοις ὑμνεῖται, ἐν δὲ πλατείαις παῤῥησίαν ἄγει.
\vs{21}Ἐπʼ ἄκρων δὲ τειχέων κηρύσσεται, ἐπὶ δὲ πύλαις δυναστῶν παρεδρεύει, ἐπὶ δὲ πύλαις πόλεως θαῤῥοῦσα λέγει,
\vs{22}ὅσον ἂν χρόνον ἄκακοι ἔχονται τῆς δικαιοσύνης, οὐκ αἰσχυνθήσονται· οἱ δὲ ἄφρονες τῆς ὕβρεως ὄντες ἐπιθυμηταί, ἀσεβεῖς γενόμενοι ἐμίσησαν αἴσθησιν,
\vs{23}καὶ ὑπεύθυνοι ἐγένοντο ἐλέγχοις· ἰδοὺ προήσομαι ὑμῖν ἐμῆς πνοῆς ῥῆσιν· διδάξω δὲ ὑμᾶς τὸν ἐμὸν λόγον.

\vs{24}Ἐπειδὴ ἐκάλουνμ, καὶ οὐχ ὑπηκούσατε· καὶ ἐξέτεινον λόγους, καὶ οὐ προσείχετε·
\vs{25}ἀλλὰ ἀκύρους ἐποιεῖτε ἐμὰς βουλὰς, τοῖς δὲ ἐμοῖς ἐλέγχοις ἠπειθήσατε·
\vs{26}Τοιγαροῦν κᾀγὼ τῇ ὑμετέρᾳ ἀπωλείᾳ ἐπιγελάσομαι, καταχαροῦμαι δὲ ἡνίκα ἔρχηται ὑμῖν ὄλεθρος·
\vs{27}Καὶ ὡς ἂν ἀφίκηται ὑμῖν ἄφνω θόρυβος, ἡ δὲ καταστροφὴ ὁμοίως καταιγίδι παρῇ, καὶ ὅταν ἔρχηται ὑμῖν θλίψις καὶ πολιορκία, ἢ ὅταν ἔρχηται ὑμῖν ὄλεθρος.
\vs{28}Ἔσται γὰρ ὅταν ἐπικαλέσησθέ με, ἐγὼ δὲ οὐκ εἰσακούσομαι ὑμῶν· ζητήσουσί με κακοὶ, καὶ οὐχ εὑρήσουσιν.
\vs{29}Ἐμίσησαν γὰρ σοφίαν, τὸν δὲ λόγον τοῦ Κυρίου οὐ προείλαντο,
\vs{30}οὐδὲ ἤθελον ἐμαῖς προσέχειν βουλαῖς, ἐμυκτήριζον δὲ ἐμοὺς ἐλέγχους·
\vs{31}Τοιγαροῦν ἔδονται τῆς ἑαυτῶν ὁδοῦ τοὺς καρποὺς, καὶ τῆς ἐαυτῶν ἀσεβείας πλησθήσονται.
\vs{32}Ἀνθʼ ὧν γὰρ ἠδίκουν νηπίους, φονευθήσονται, καὶ ἐξετασμὸς ἀσεβεῖς ὀλεῖ.
\vs{33}Ὁ δὲ ἐμοῦ ἀκούων κατασκηνώσει ἐπʼ ἐλπίδι, καὶ ἡσυχάσει ἀφόβως ἀπὸ παντὸς κακοῦ.

\ch{2}
Υἱὲ, ἐὰν δεξάμενος ῥῆσιν ἐμῆς ἐντολῆς κρύψῃς παρὰ σεαυτῷ,
\vs{2}ὑπακούσεται σοφίας τὸ οὖς σου, καὶ παραβαλεῖς καρδίαν σου εἰς σύνεσιν, παραβαλεῖς δὲ αὐτὴν ἐπὶ νουθέτησιν τῷ υἱῷ σου·

\vs{3}Ἐὰν γὰρ τὴν σοφίαν ἐπικαλέσῃ, καὶ τῇ συνέσει δῷς φωνήν σου,
\vs{4}καὶ ἐὰν ζητήσῃς αὐτὴν ὡς ἀργύριον, καὶ ὡς θησευροὺς ἐξεραυνήσῃς αὐτήν·
\vs{5}Τότε συνήσεις φόβον Κυρίου, καὶ ἐπίγνωσιν Θεοῦ εὑρήσεις.

\vs{6}Ὅτι Κύριος δίδωσι σοφίαν, καὶ ἀπὸ προσώπου αὐτοῦ γνῶσις καὶ σύνεσις.
\vs{7}Καὶ θησαυρίζει τοῖς κατορθοῦσι σωτηρίαν, ὑπερασπιεῖ τὴν πορείαν αὐτῶν,
\vs{8}τοῦ φυλάξαι ὁδοὺς δικαιωμάτων, καὶ ὁδὸν εὐλαβουμένων αὐτὸν διαφυλάξει.
\vs{9}Τότε συνήσεις δικαιοσύνην καὶ κρίμα, καὶ κατορθώσεις πάντας ἄξονας ἀγαθούς.

\vs{10}Ἐὰν γὰρ ἔλθῃ ἡ σοφία εἰς σὴν διάνοιαν, ἡ δὲ αἴσθησις τῇ σῇ ψυχῇ καλὴ εἶναι δόξῃ,
\vs{11}βουλὴ καλὴ φυλάξει σε, ἔννοια δὲ ὁσία τηρήσει σε·
\vs{12}Ἵνα ῥύσηταί σε ἀπὸ ὁδοῦ κακῆς, καὶ ἀπὸ ἀνδρὸς λαλοῦντος μηδὲν πιστόν.

\vs{13}Ὦ οἱ ἐγκαταλείποντες ὁδοὺς εὐθείας τοῦ πορεύεσθαι ἐν ὁδοῖς σκότους·
\vs{14}Οἱ εὐφραινόμενοι ἐπὶ κακοῖς καὶ χαίροντες ἐπὶ διαστροφῇ κακῇ·
\vs{15}Ὧν αἱ τρίβοι σκολιαὶ, καὶ καμπύλαι αἱ τροχιαὶ αὐτῶν,
\vs{16}τοῦ μακράν σε ποιῆσαι ἀπὸ ὁδοῦ εὐθείας, καὶ ἀλλότριον τῆς δικαίας γνώμης· υἱὲ, μή σε καταλάβῃ κακὴ βουλή·
\vs{17}Ἡ ἀπολιποῦσα διδασκαλίαν νεότητος, καὶ διαθήκην θείαν ἐπιλελησμένη.
\vs{18}Ἔθετο γὰρ παρὰ τῷ θανάτῳ τὸν οἶκον αὐτῆς, καὶ παρὰ τῷ ᾅδῃ μετὰ τῶν γηγενῶν τοὺς ἄξονας αὐτῆς.
\vs{19}Πάντες οἱ πορευόμενοι ἐν αὐτῇ οὐκ ἀναστρέψουσιν, οὐδὲ μὴ καταλάβωσι τρίβους εὐθείας· οὐ γὰρ καταλαμβάνονται ὑπὸ ἐνιαυτῶν ζωῆς.
\vs{20}Εἰ γὰρ ἐπορεύοντο τρίβους ἀγαθὰς, εὕροσαν ἂν τρίβους δικαιοσύνης λείας.
\vs{21}Ὅτι εὐθεῖς κατασκηνώσουσι γῆν, καὶ ὅσιοι ὑπολειφθήσονται ἐν αὐτῇ.
\vs{22}Ὁδοὶ ἀσεβῶν ἐκ γῆς ὀλοῦνται, οἱ δὲ παράνομοι ἐξωσθήσονται ἀπʼ αὐτῆς.

\ch{3}
Υἱὲ, ἐμῶννομίμων μὴ ἐπιλανθάνου, τὰ δὲ ῥήματά μου τηρείτω σὴ καρδία·
\vs{2}Μῆκος γὰρ βίου, καὶ ἔτη ζωῆς, καὶ εἰρήνην προσθήσουσί σοι.
\vs{3}Ἐλεημοσύναι καὶ πίστεις μὴ ἐκλειπέτωσάν σε· ἄφαψαι δὲ αὐτὰς ἐπὶ σῷ τραχήλῳ, καὶ εὑρήσεις χάριν·
\vs{4}καὶ προνοοῦ καλὰ ἐνώπιον Κυρίου καὶ ἀνθρώπων.

\vs{5}Ἴσθι πεποιθὼς ἐν ὅλῃ τῇ καρδίᾳ ἐπὶ Θεῷ, ἐπὶ δὲ σῇ σοφίᾳ μὴ ἐπαίρου.
\vs{6}Πάσαις ὁδοῖς σου γνώριζε αὐτὴν, ἵνα ὀρθοτομῇ τὰς ὁδούς σου.
\vs{7}Μὴ ἴσθι φρόνιμος παρὰ σεαυτῷ, φοβοῦ δὲ τὸν Θεὸν, καὶ ἔκκλινε ἀπὸ παντὸς κακοῦ.
\vs{8}Τότε ἴασις ἔσται τῷ σώματί σου, καὶ ἐπιμέλεια τοῖς ὀστέοις σου.

\vs{9}Τίμα τὸν Κύριον ἀπὸ σῶν δικαίων πόνων, καὶ ἀπάρχου αὐτῷ ἀπὸ σῶν καρπῶν δικαιοσύνης·
\vs{10}Ἵνα πίμπληται τὰ ταμιεῖά σου πλησμονῆς σίτῳ, οἴνῳ δὲ αἱ ληνοί σου ἐκβλύζωσιν.

\vs{11}Υἱὲ, μὴ ὀλιγώρει παιδείας Κυρίου, μηδὲ ἐκλύου ὑπʼ αὐτοῦ ἐλεγχόμενος.
\vs{12}Ὃν γὰρ ἀγαπᾷ Κύριος, ἐλέγχει, μαστιγοῖ δὲ πάντα υἱὸν ὃν παραδέχεται.

\vs{13}Μακάριος ἄνθρωπος ὃς εὗρε σοφίαν, καὶ θνητὸς ὃς εἶδε φρόνησιν.
\vs{14}Κρεῖσσον γὰρ αὐτὴν ἐμπορεύεσθαι, ἢ χρυσίου καὶ ἀργυρίου θησαυρούς.
\vs{15}Τιμιωτέρα δέ ἐστι λίθων πολυτελῶν, οὐκ ἀντιτάξεται αὐτῇ οὐδὲν πονηρόν· εὔγνωστός ἐστι πᾶσι τοῖς ἐγγίζουσιν αὐτῇ, πᾶν δὲ τίμιον οὐκ ἄξιον αὐτῆς ἐστι.
\vs{16}Μῆκος γὰρ βίου καὶ ἔτη ζωῆς ἐν τῇ δεξιᾷ αὐτῆς, ἐν δὲ τῇ ἀριστερᾷ αὐτῆς πλοῦτος καὶ δόξα·
\vs{16a}ἐκ τοῦ στόματος αὐτῆς ἐκπορεύεται δικαιοσύνη, νόμον δὲ καὶ ἔλεον ἐπὶ γλώσσης φορεῖ.
\vs{17}Αἱ ὁδοὶ αὐτῆς ὁδοὶ καλαί, καὶ πάσαι αἱ τρίβοι αὐτῆς ἐν εἰρήνῃ.
\vs{18}Ξύλον ζωῆς ἐστι πᾶσι τοῖς ἀντεχομένοις αὐτῆς, καὶ τοῖς ἐπερειδομένοις ἐπʼ αὐτὴν ὡς ἐπὶ Κύριον ἀσφαλής.

\vs{19}Ὁ Θεὸς τῇ σοφίᾳ ἐθεμελίωσε τὴν γῆν, ἡτοίμασε δὲ οὐρανοὺς φρονήσει.
\vs{20}Ἐν αἰσθήσει ἄβυσσοι ἐῤῥάγησαν, νέφη δὲ ἐῤῥύησαν δρόσους.

\vs{21}Υἱὲ, μὴ παραῤῥυῇς, τήρησον δὲ ἐμὴν βουλὴν καὶ ἔννοιαν·
\vs{22}ἵνα ζήσῃ ἡ ψυχή σου, καὶ χάρις ᾖ περὶ σῷ τραχήλῳ·
\vs{22a}ἔσται δὲ ἴασις ταῖς σαρξί σου, καὶ ἐπιμέλεια τοῖς σοῖς ὀστέοις·
\vs{23}ἵνα πορεύῃ πεποιθὼς ἐν εἰρήνῃ πάσας τὰς ὁδούς σου, ὁ δὲ πούς σου οὐ μὴ προσκόψῃ.
\vs{24}Ἐὰν γὰρ κάθῃ, ἄφοβος ἔσῃ· ἐὰν δὲ καθεύδῃς, ἡδέως ὑπνώσεις.
\vs{25}Καὶ οὐ φοβηθήσῃ πτόησιν ἐπελθοῦσαν, οὐδὲ ὁρμὰς ἀσεβῶν ἐπερχομένας.
\vs{26}Ὁ γὰρ Κύριος ἔσται ἐπὶ πασῶν ὁδῶν σου, καὶ ἐρείσει σὸν πόδα ἵνα μὴ σαλευθῇς.

\vs{27}Μὴ ἀπόσχῃ εὖ ποιεῖν ἐνδεῆ, ἡνίκα ἂν ἔχῃ ἡ χείρ σου βοηθεῖν.
\vs{28}Μὴ εἴπῃς, ἐπανελθὼν ἐπάνηκε, αὔριον δώσω, δυνατοῦ σου ὄντος εὖ ποιεῖν· οὐ γὰρ οἶδας τί τέξεται ἡ ἐπιοῦσα.
\vs{29}Μὴ τεκτῄνῃ ἐπὶ σὸν φίλον κακὰ παροικοῦντα καὶ πεποιθότα ἐπὶ σοί.

\vs{30}Μὴ φιλεχθρήσῃς πρὸς ἄνθρωπον μάτην, μήτί σε ἐργάσηται κακόν.

\vs{31}Μὴ κτήσῃ κακῶν ἀνδρῶν ὀνείδη, μηδὲ ζηλώσῃς τὰς ὁδοὺς αὐτῶν.
\vs{32}Ἀκάθαρτος γὰρ ἔναντι Κυρίου πᾶς παράνομος, ἐν δὲ δικαίοις οὐ συνεδριάζει.
\vs{33}Κατάρα Θεοῦ ἐν οἴκοις ἀσεβῶν, ἐπαύλεις δὲ δικαίων εὐλογοῦνται.
\vs{34}Κύριος ὑπερηφάνοις ἀντιτάσσεται, ταπεινοῖς δὲ δίδωσι χάριν.
\vs{35}Δόξαν σοφοὶ κληρονομήσουσιν, οἱ δὲ ἀσεβεῖς ὕψωσαν ἀτιμίαν.

\ch{4}
Ἀκούσατε, παῖδες, παιδείαν πατρὸς, καὶ προσέχετε γνῶναι ἔννοιαν.
\vs{2}Δῶρον γὰρ ἀγαθὸν δωροῦμαι ὑμῖν, τὸν ἐμὸν νόμον μὴ ἐγκαταλίπητε.
\vs{3}Υἱὸς γὰρ ἐγενόμην κᾀγὼ πατρὶ ὑπήκοος, καὶ ἀγαπώμενος ἐν προσώπῳ μητρός.
\vs{4}Οἳ ἔλεγον καὶ ἐδίδασκόν με, ἐρειδέτω ὁ ἡμέτερος λόγος εἰς σὴν καρδίαν· φύλασσε ἐντολὰς,
\vs{5}μὴ ἐπιλάθῃ· Μηδὲ παρίδῃς ῥῆσιν ἐμοῦ στόματος,
\vs{6}μηδὲ ἐγκαταλίπῃς αὐτὴν, καὶ ἀνθέξεταί σου· ἐράσθητι αὐτῆς, καὶ τηρήσει σε.
\vs{8}Περιχαράκωσον αὐτὴν, καὶ ὑψώσει σε· τίμησον αὐτὴν, ἵνα σε περιλάβῃ·
\vs{9}Ἵνα δῷ τῇ σῇ κεφαλῇ στέφανον χαρίτων, στεφάνῳ δὲ τρυφῆς ὑπερασπίσῃ σου.

\vs{10}Ἄκουε υἱὲ καὶ δέξαι ἐμοὺς λόγους, καὶ πληθυνθήσεται ἔτη ζωῆς σου, ἵνα σοι γένωνται πολλαὶ ὁδοὶ βίου.
\vs{11}Ὁδοὺς γὰρ σοφίας διδάσκω σε, ἐμβιβάζω δέ σε τροχιαῖς ὀρθαῖς.
\vs{12}Ἐὰν γὰρ πορεύῃ, οὐ συγκλεισθήσεταί σου τὰ διαβήματα· ἐὰν δὲ τρέχῃς, οὐ κοπιάσεις.
\vs{13}Ἐπιλαβοῦ ἐμῆς παιδείας, μὴ ἀφῇς, ἀλλὰ φύλαξον αὐτὴν σεαυτῷ εἰς ζωήν σου.

\vs{14}Ὁδοὺς ἀσεβῶν μὴ ἐπέλθῃς, μηδὲ ζηλώσῃς ὁδοὺς παρανόμων.
\vs{15}Ἐν ᾧ ἂν τόπῳ στρατοπεδεύσωσι, μὴ ἐπέλθῃς ἐκεὶ, ἔκκλινον δὲ ἀπʼ αὐτῶν καὶ παράλλαξον.
\vs{16}Οὐ γὰρ μὴ ὑπνώσωσιν, ἐὰν μὴ κακοποιήσωσιν· ἀφῄρηται ὁ ὕπνος αὐτῶν, καὶ οὐ κοιμῶνται.
\vs{17}Οἵδε γὰρ σιτοῦνται σῖτα ἀσεβείας, οἴνῳ δὲ παρανόμῳ μεθύσκονται.
\vs{18}Αἱ δὲ ὁδοὶ τῶν δικαίων ὁμοίως φωτὶ λάμπουσι, προπορεύονται καὶ φωτίζουσιν, ἕως κατορθώσῃ ἡ ἡμέρα.
\vs{19}Αἱ δὲ ὁδοὶ τῶν ἀσεβῶν σκοτειναὶ, οὐκ οἴδασι πῶς προσκόπτουσιν.

\vs{20}Υἱὲ ἐμῇ ῥήσει πρόσεχε, τοῖς δὲ ἐμοῖς λόγοις παράβαλλε σὸν οὖς.
\vs{21}Ὅπως μὴ ἐκλίπωσί σε αἱ πηγαί σου, φύλασσε αὐτὰς ἐν καρδίᾳ.
\vs{22}Ζωὴ γάρ ἐστι τοῖς εὑρίσκουσιν αὐτὰς, καὶ πάσῃ σαρκὶ ἴασις.
\vs{23}Πάσῃ φυλακῇ τήρει σὴν καρδίαν, ἐκ γὰρ τούτων ἔξοδοι ζωῆς.
\vs{24}Περίελε σεαυτοῦ σκολιὸν στόμα, καὶ ἄδικα χείλη μακρὰν ἀπὸ σοῦ ἄπωσαι.
\vs{25}Οἱ ὀφθαλμοί σου ὀρθὰ βλεπέτωσαν, τὰ δὲ βλέφαρά σου νευέτω δίκαια.
\vs{26}Ὀρθὰς τροχιὰς ποίει σοῖς ποσί, καὶ τὰς ὁδούς σου κατεύθυνε.
\vs{27}Μὴ ἐκκλίνῃς εἰς τὰ δεξιὰ, μηδὲ εἰς τὰ ἀριστερά, ἀπόστρεψον δὲ σὸν πόδα ἀπὸ ὁδοῦ κακῆς·
\vs{27a}ὁδοὺς γὰρ τὰς ἐκ δεξιῶν οἶδεν ὁ Θεὸς, διεστραμμέναι δέ εἰσιν αἱ ἐξ ἀριστερῶν·
\vs{27b}αὐτὸς δὲ ὀρθὰς ποιήσει τὰς τροχιάς σου, τὰς δὲ πορείας σου ἐν εἰρήνῃ προάξει.

\ch{5}
Υἱὲ, ἐμῇ σοφίᾳ πρόσεχε, ἐμοῖς δὲ λόγοις παράβαλλε σὸν οὖς,
\vs{2}ἵνα φυλάξῃς ἔννοιαν ἀγαθήν· αἴσθησις δὲ ἐμῶν χειλέων ἐντέλλεταί σοι·

\vs{3}Μὴ πρόσεχε φαύλῃ γυναικί. Μέλι γὰρ ἀποστάζει ἀπὸ χειλέων γυναικὸς πόρνης, ἣ πρὸς καιρὸν λιπαίνει σὸν φάρυγγα,
\vs{4}ὕστερον μέντοι πικρότερον χολῆς εὑρήσεις, καὶ ἠκονημένον μᾶλλον μαχαίρας διστόμου.
\vs{5}Τῆς γὰρ ἀφροσύνης οἱ πόδες κατάγουσι τοὺς χρωμένους αὐτῇ μετὰ θανάτου εἰς τὸν ᾅδην, τὰ δὲ ἴχνη αὐτῆς οὐκ ἐρείδεται.
\vs{6}Ὁδοὺς γὰρ ζωῆς οὐκ ἐπέρχεται, σφαλεραὶ δὲ αἱ τροχιαὶ αὐτῆς, καὶ οὐκ εὔγνωστοι.

\vs{7}Νῦν οὖν υἱὲ ἄκουέ μου, καὶ μὴ ἀκύρους ποιήσεις ἐμοὺς λόγους.
\vs{8}Μακρὰν ποίησον ἀπʼ αὐτῆς σὴν ὁδόν· μὴ ἐγγίσῃς πρὸς θύραις οἴκων αὐτῆς,
\vs{9}ἵνα μὴ πρόῃ ἄλλοις ζωήν σου, καὶ σὸν βίον ἀνελεήμοσιν·
\vs{10}Ἵνα μὴ πλησθῶσιν ἀλλότριοι σῆς ἰσχύος, οἱ δὲ σοὶ πόνοι εἰς οἴκους ἀλλοτρίων ἔλθεσι·
\vs{11}Καὶ μεταμεληθήσῃ ἐπʼ ἐσχάτων, ἡνίκα ἂν κατατριβῶσι σάρκες σώματός σου,
\vs{12}καὶ ἐρεῖς, πῶς ἐμίσησα παιδείαν, καὶ ἐλέγχους ἐξέκλινεν ἡ καρδία μου;
\vs{13}Οὐκ ἤκουον φωνὴν παιδεύοντός με καὶ διδάσκοντός με, οὐδὲ παρέβαλλον τὸ οὖς μου.
\vs{14}Παρʼ ὀλίγον ἐγενόμην ἐν παντὶ κακῷ, ἐν μέσῳ ἐκκλησίας καὶ συναγωγῆς.

\vs{15}Πίνε ὕδατα ἀπὸ σῶν ἀγγείων, καὶ ἀπὸ σῶν φρεάτων πηγῆς.
\vs{16}Μὴ ὑπερεκχείσθω σοι ὕδατα ἐκ τῆς σῆς πηγῆς, εἰς δὲ σὰς πλατείας διαπορευέσθω τὰ σὰ ὕδατα.
\vs{17}Ἔστω σοι μόνῳ ὑπάρχοντα, καὶ μηδεὶς ἀλλότριος μετασχέτω σοι.
\vs{18}Ἡ πηγή σου τοῦ ὕδατος ἔστω σοι ἰδία, καὶ συνευφραίνου μετὰ γυναικὸς τῆς ἐκ νεότητός σου.
\vs{19}Ἔλαφος φιλίας καὶ πῶλος σῶν χαρίτων ὁμιλείτω σοι, ἡ δὲ ἰδία ἡγείσθω σου καὶ συνέστω σοι ἐν παντὶ καιρῷ· ἐν γὰρ τῇ ταύτης φιλίᾳ συμπεριφερόμενος, πολλοστὸς ἔσῃ.
\vs{20}Μὴ πολὺς ἴσθι πρὸς ἀλλοτρίαν, μηδὲ συνέχου ἀγκάλαις τῆς μὴ ἰδίας.
\vs{21}Ἐνώπιον γάρ εἰσι τῶν τοῦ Θεοῦ ὀφθαλμῶν ὁδοὶ ἀνδρὸς, εἰς δὲ πάσας τὰς τροχιὰς αὐτοῦ σκοπεύει.
\vs{22}Παρανομίαι ἄνδρα ἀγρεύουσι, σειραῖς δὲ τῶν ἑαυτοῦ ἁμαρτιῶν ἕκαστος σφίγγεται.
\vs{23}Οὗτος τελευτᾷ μετὰ ἀπαιδεύτων, ἐκ δὲ πλήθους τῆς ἑαυτοῦ βιότητος ἐξεῤῥίφη, καὶ ἀπώλετο διʼ ἀφροσύνην.

\ch{6}
Υἱὲ, ἐὰν ἐγγυήσῃ σὸν φίλον, παραδώσεις σὴν χεῖρα ἐχθρῷ.
\vs{2}Παγὶς γὰρ ἰσχυρὰ ἀνδρὶ τὰ ἴδια χείλη, καὶ ἁλίσκεται χείλεσιν ἰδίου στόματος.
\vs{3}Ποίει υἱὲ ἃ ἐγώ σοι ἐντέλλομαι, καὶ σώζου· ἥκεις γὰρ εἰς χεῖρας κακῶν διὰ σὸν φίλον· ἴσθι μὴ ἐκλυόμενος, παρόξυνε δὲ καὶ τὸν φίλον σου ὃν ἐνεγγυήσω.
\vs{4}Μὴ δῷς ὕπνον σοῖς ὄμμασι, μηδὲ ἐπινυστάξῃς σοῖς βλεφάροις,
\vs{5}ἵνα σώζῃ ὥσπερ δορκὰς ἐκ βρόχων, καὶ ὥσπερ ὄρνεον ἐκ παγίδος.

\vs{6}Ἴθι πρὸς τὸν μύρμηκα ὦ ὀκνηρὲ, καὶ ζήλωσον ἰδὼν τὰς ὁδοὺς αὐτοῦ, καὶ γενοῦ ἐκείνου σοφώτερος.

\vs{7}Ἐκείνῳ γὰρ γεωργίου μὴ ὑπάρχοντος, μηδὲ τὸν ἀναγκάζοντα ἔχων, μηδὲ ὑπὸ δεσπότην ὢν,
\vs{8}ἐτοιμάζεται θέρους τὴν τροφὴν, πολλήν τε ἐν τῷ ἀμητῷ ποιεῖται τὴν παράθεσιν·
\vs{8a}ἢ πορεύθητι πρὸς τὴν μέλισσαν, καὶ μάθε ὡς ἐργάτις ἐστὶ, τήν τε ἐργασίαν ὡς σεμνὴν ποιεῖται·
\vs{8b}ἧς τοὺς πόνους βασιλεῖς καὶ ἰδιῶται πρὸς ὑγίειαν προσφέρονται· ποθεινὴ δέ ἐστι πᾶσι καὶ ἐπίδοξος,
\vs{8c}καίπερ οὖσα τῇ ῥώμῃ ἀσθενὴς, τὴν σοφίαν τιμήσασα προήχθη.
\vs{9}Ἕως τίνος ὀκνηρὲ κατάκεισαι; πότε δὲ ἐξ ὕπνου ἐγερθήσῃ;
\vs{10}ὀλίγον μὲν ὑπνοῖς, ὀλίγον δὲ κάθησαι, μικρὸν δὲ νυστάζεις, ὀλίγον δὲ ἐναγκαλίζῃ χερσὶ στήθη.
\vs{11}Εἶτʼ ἐνπαραγίνεταί σοι ὥσπερ κακὸς ὁδοιπόρος ἡ πενία, καὶ ἡ ἔνδεια ὥσπερ ἀγαθὸς δρομεύς·
\vs{11a}ἐὰν δὲ ἄοκνος ᾖς, ἥξει ὥσπερ πηγὴ ὁ ἀμητός σου· ἡ δὲ ἔνδεια, ὥσπερ κακὸς δρομεὺς ἀπαυτομολήσει.

\vs{12}Ἀνὴρ ἄφρων καὶ παράνομος πορεύεται ὁδοὺς οὐκ ἀγαθάς.
\vs{13}Ὁ δʼ αὐτὸς ἐννεύει ὀφθαλμῷ, σημαίνει δὲ ποδὶ, διδάσκει δὲ ἐννεύμασι δακτύλων.
\vs{14}Διεστραμμένη καρδία τεκταίνεται κακὰ, ἐν παντὶ καιρῷ ὁ τοιοῦτος ταραχὰς συνίστησιν πόλει.
\vs{15}Διὰ τοῦτο ἐξαπίνης ἔρχεται ἡ ἀπώλεια αὐτοῦ, διακοπὴ καὶ συντριβὴ ἀνίατος.

\vs{16}Ὅτι χαίρει πᾶσιν οἷς μισεῖ ὁ Θεὸς, συντρίβεται δὲ διʼ ἀκαθαρσίαν ψυχῆς.
\vs{17}Ὀφθαλμὸς ὑβριστοῦ, γλῶσσα ἄδικος· χεῖρες ἐκχέουσαι αἷμα δικαίου,
\vs{18}καὶ καρδία τεκταινομένη λογισμοὺς κακοὺς, καὶ πόδες ἐπισπεύδοντες κακοποιεῖν.
\vs{19}Ἐκκαίει ψευδῆ μάρτυς ἄδικος, καὶ ἐπιπέμπει κρίσεις ἀναμέσον ἀδελφῶν.

\vs{20}Υἱὲ, φύλασσε νόμους πατρός σου, καὶ μὴ ἀπώσῃ θεσμοὺς μητρός σου·
\vs{21}Ἄφαψαι δὲ αὐτοὺς ἐπὶ σῇ ψυχῇ διαπαντὸς, καὶ ἐλκλοίωσαι περὶ σῷ τραχήλῳ·
\vs{22}Ἡνίκα ἂν περιπατῇς, ἐπάγου αὐτὴν καὶ μετὰ σοῦ ἔστω, ὡς δʼ ἂν καθεύδῃς φυλασσέτω σε, ἵνα ἐγειρομένῳ συλλαλῇ σοι.
\vs{23}Ὅτι λύχνος ἐντολὴ νόμου καὶ φῶς, ὁδὸς ζωῆς, καὶ ἔλεγχος καὶ παιδεία,
\vs{24}τοῦ διαφυλάσσειν σε ἀπὸ γυναικὸς ὑπάνδρου, καὶ ἀπὸ διαβολῆς γλώσσης ἀλλοτρίας.

\vs{25}Μή σε νικήσῃ κάλλους ἐπιθυμία, μηδὲ ἀγρευθῇς σοῖς ὀφθαλμοῖς, μηδὲ συναρπασθῇς ἀπὸ τῶν αὐτῆς βλεφάρων.
\vs{26}Τιμὴ γὰρ πόρνης ὅση καὶ ἑνὸς ἄρτου, γυνὴ δὲ ἀνδρῶν τιμίας ψυχὰς ἀγρεύει.
\vs{27}Ἀποδήσει τις πῦρ ἐν κόλπῳ, τὰ δὲ ἱμάτια οὐ κατακαύσει;
\vs{28}ἢ περιπατήσει τις ἐπʼ ἀνθράκων πυρὸς, τοὺς δὲ πόδας οὐ κατακαύσει;
\vs{29}Οὕτως ὁ εἰσελθὼν πρὸς γυναῖκα ὕπανδρον, οὐκ ἀθωωθήσεται, οὐδὲ πᾶς ὁ ἁπτόμενος αὐτῆς.
\vs{30}Οὐ θαυμαστὸν ἐὰν ἁλῷ τις κλέπτων, κλέπτει γὰρ ἵνα ἐμπλήσῃ τὴν ψυχὴν πεινῶν.
\vs{31}Ἐὰν δὲ ἁλῷ, ἀποτίσει ἑπταπλάσια, καὶ πάντα τὰ ὑπάρχοντα αὐτοῦ δοὺς ῥύσεται ἑσυτόν.
\vs{32}Ὁ δὲ μοιχὸς διʼ ἔνδειαν φρενῶν ἀπώλειαν τῇ ψυχῇ αὐτοῦ περιποιεῖται,
\vs{33}ὀδύνας τε καὶ ἀτιμίας ὑποφέρει, τὸ δὲ ὄνειδος αὐτοῦ οὐκ ἐξαλειφθήσεται εἰς τὸν αἰῶνα.
\vs{34}Μεστὸς γὰρ ζήλου θυμὸς ἀνδρὸς αὐτῆς, οὐ φείσεται ἐν ἡμέρᾳ κρίσεως.
\vs{35}Οὐκ ἀνταλλάξεται οὐδενὸς λύτρου τὴν ἔχθραν, οὐδὲ μὴ διαλυθῇ πολλῶν δώρων.

\ch{7}
Υἱὲ φύλασσε ἐμοὺς λόγους, τὰς δὲ ἐμὰς ἐντολὰς κρύψον παρὰ σεαυτῷ·
\vs{1a}Υἱὲ τίμα τὸν Κύριον καὶ ἰσχύσεις, πλὴν δὲ αὐτοῦ μὴ φοβοῦ ἄλλον·
\vs{2}φύλαξον ἐμὰς ἐντολὰς καὶ βιώσεις, τοὺς δὲ ἐμοὺς λόγους ὥσπερ κόρας ὀμμάτων.
\vs{3}Περίθου δὲ αὐτοὺς σοῖς δακτύλοις, ἐπίγραψον δὲ ἐπὶ τὸ πλάτος τῆς καρδίας σου.

\vs{4}Εἰπὸν τὴν σοφίαν σὴν ἀδελφὴν εἶναι, τὴν δὲ φρόνησιν γνώριμον περιποίησαι σεαυτῷ.
\vs{5}Ἵνα σε τηρήσῃ ἀπὸ γυναικὸς ἀλλοτρίας καὶ πονηρᾶς, ἐάν σε λόγοις τοῖς πρὸς χάριν ἐμβάληται.

\vs{6}Ἀπὸ γὰρ θυρίδος ἐκ τοῦ οἴκου αὐτῆς εἰς τὰς πλατείας παρακύπτουσα,
\vs{7}ὃν ἂν ἴδῃ τῶν ἀφρόνων τέκνων νεανίαν ἐνδεῆ φρενῶν,
\vs{8}παραπορευόμενον παρὰ γωνίαν ἐν διόδοις οἴκων αὐτῆς, καὶ λαλοῦντα
\vs{9}ἐν σκότει ἑσπερινῷ, ἡνίκα ἂν ἡσυχία νυκτερινὴ καὶ γνοφώδης,
\vs{10}ἡ δὲ γυνὴ συναντᾷ αὐτῷ, εἶδος ἔχουσα πορνικὸν, ἣ ποιεῖ νέων ἐξίπτασθαι καρδίας.
\vs{11}Ἀνεπτερωμένη δέ ἐστι καὶ ἄσωτος, ἐν οἴκῳ δὲ οὐχ ἡσυχάζουσιν οἱ πόδες αὐτῆς.
\vs{12}Χρόνον γάρ τινα ἔξω ῥέμβεται, χρόνον δὲ ἐν πλατείαις παρὰ πᾶσαν γωνίαν ἐνεδρεύει.
\vs{13}Εἶτα ἐπιλαβομένη ἐφίλησεν αὐτὸν, ἀναιδεῖ δὲ προσώπῳ προσεῖπεν αὐτῷ,
\vs{14}θυσία εἰρηνική μοι ἐστὶ, σήμερον ἀποδίδωμι τὰς εὐχάς μου.
\vs{15}Ἕνεκα τούτου ἐξῆλθον εἰς συνάντησίν σοι, ποθοῦσα τὸ σὸν πρόσωπον, εὕρηκά σε.
\vs{16}κειρίαις τέτακα τὴν κλίνην μου, ἀμφιτάποις δὲ ἔστρωκα τοῖς ἀπʼ Αἰγύπτου.
\vs{17}Διέῤῥαγκα τὴν κοίτην μου κροκίνῳ, τὸν δὲ οἶκόν μου κινναμώμῳ·
\vs{18}Ἐλθὲ καὶ ἀπολαύσωμεν φιλίας ἕως ὄρθρου, δεῦρο καὶ ἐλκυλισθῶμεν ἔρωτι.
\vs{19}Οὐ γὰρ πάρεστιν ὁ ἀνήρ μου ἐν οἴκω, πεπόρευται δὲ ὁδὸν μακράν,
\vs{20}ἔνδεσμον ἀργυρίου λαβὼν ἐν χειρὶ αὐτοῦ, διʼ ἡμερῶν πολλῶν ἐπανήξει εἰς τὸν οἶκον αὐτοῦ.

\vs{21}Ἀπεπλάνησε δὲ αὐτὸν πολλῇ ὁμιλίᾳ, βρόχοις τε τοῖς ἀπὸ χειλέων ἐξώκειλεν αὐτόν.
\vs{22}Ὁ δὲ ἐπηκολούθησεν αὐτῇ κεπφωθείς· ὥσπερ δὲ βοῦς ἐπὶ σφαγὴν ἄγεται, καὶ ὥσπερ κύων ἐπὶ δεσμοὺς,
\vs{23}ἢ ὡς ἔλαφος τοξεύματι πεπληγὼς εἰς τὸ ἧπαρ· σπεύδει δὲ ὥσπερ ὄρνεον εἰς παγίδα, οὐκ εἰδὼς ὅτι περὶ ψυχῆς τρέχει.

\vs{24}Νῦν οὖν υἱὲ ἄκουέ μου, καὶ πρόσεχε ῥήμασι στόματός μου.
\vs{25}Μὴ ἐκκλινάτω εἰς τὰς ὁδοὺς αὐτῆς ἡ καρδία σου,
\vs{26}πολλοὺς γὰρ τρώσασα καταβέβληκε, καὶ ἀναρίθμητοί εἰσιν οὓς πεφόνευκεν.
\vs{27}Ὁδοὶ ᾅδου ὁ οἶκος αὐτῆς, κατάγουσαι εἰς τὰ ταμιεῖα τοῦ θανάτου.

\ch{8}
Σὺ τὴν σοφίαν κηρύξεις, ἵνα φρόνησίς σοι ὑπακούσῃ.
\vs{2}Ἐπὶ γὰρ τῶν ὑψηλῶν ἄκρων ἐστὶν, ἀναμέσον δὲ τῶν τρίβων ἕστηκε.
\vs{3}Παρὰ γὰρ πύλαις δυναστῶν παρεδρεύει, ἐν δὲ εἰσόδοις ὑμνεῖται.
\vs{4}Ὑμᾶς ὦ ἄνθρωποι παρακαλῶ, καὶ προΐεμαι ἐμὴν φωνὴν υἱοῖς ἀνθρώπων.
\vs{5}Νοήσατε ἄκακοι πανουργίαν, οἱ δὲ ἀπαίδευτοι ἔνθεσθε καρδίαν.
\vs{6}Εἰσακούσατέ μου, σεμνὰ γὰρ ἐρῶ, καὶ ἀνοίσω ἀπὸ χειλέων ὀρθά.
\vs{7}Ὅτι ἀλήθειαν μελετήσει ὁ φάρυγξ μου, ἐβδελυγμένα δὲ ἐναντίον ἐμοῦ χείλη ψευδῆ.
\vs{8}Μετὰ δικαιοσύνης πάντα τὰ ῥήματα τοῦ στόματός μου, οὐδὲν ἐαυτοῖς σκολιὸν οὐδὲ στραγγαλιῶδες.
\vs{9}Πάντα ἐνώπια τοῖς συνιοῦσι, καὶ ὀρθὰ τοῖς εὑρίσκουσι γνῶσιν.
\vs{10}Λάβετε παιδείαν καὶ μὴ ἀργύριον, καὶ γνῶσιν ὑπὲρ χρυσίον δεδοκιμασμένον·
\vs{11}Κρείσσων γὰρ σοφία λίθων πολυτελῶν, πᾶν δὲ τίμιον οὐκ ἄξιον αὐτῆς ἐστιν.

\vs{12}Ἐγὼ ἡ σοφία κατεσκήνωσα βουλὴν καὶ γνῶσιν, καὶ ἔννοιαν ἐγὼ ἐπεκαλεσάμην.
\vs{13}Φόβος Κυρίου μισεῖ ἀδικίαν, ὕβριν τε καὶ ὑπερηφανίαν καὶ ὁδοὺς πονηρῶν· μεμίσηκα δὲ ἐγὼ διεστραμμένας ὁδοὺς κακῶν.
\vs{14}Ἐμὴ βουλὴ καὶ ἀσφάλεια, ἐμὴ φρόνησις, ἐμὴ δὲ ἰσχύς.
\vs{15}Διʼ ἐμοῦ βασιλεῖς βασιλεύουσι, καὶ οἱ δυνάσται γράφουσιν δικαιοσύνην.
\vs{16}Διʼ ἐμοῦ μεγιστᾶνες μεγαλύνονται, καὶ τύραννοι διʼ ἐμοῦ κρατοῦσι γῆς.
\vs{17}Ἐγὼ τοὺς ἐμὲ φιλοῦντας ἀγαπῶ, οἱ δὲ ἐμὲ ζητοῦντες εὑρήσουσιν.

\vs{18}Πλοῦτος καὶ δόξα ἐμοὶ ὑπάρχει, καὶ κτῆσις πολλῶν καὶ δικαιοσύνη.
\vs{19}Βέλτιον ἐμὲ καρπίζεσθαι ὑπὲρ χρυσίον καὶ λίθον τίμιον, τὰ δὲ ἐμὰ γεννήματα κρείσσω ἀργυρίου ἐκλεκτοῦ.
\vs{20}Ἐν ὁδοῖς δικαιοσύνης περιπατῶ, καὶ ἀναμέσον τρίβων δικαιώματος ἀναστρέφομαι·
\vs{21}ἵνα μερίσω τοῖς ἐμὲ ἀγαπῶσιν ὕπαρξιν, καὶ τοὺς θησαυροὺς αὐτῶν ἐμπλήσω ἀγαθῶν·
\vs{21a}ἐὰν ἀναγγείλω ὑμῖν τὰ καθʼ ἡμέραν γινόμενα, μνημονεύσω τὰ ἐξ αἰῶνος ἀριθμῆσαι.

\vs{22}Κύριος ἔκτισέ με ἀρχὴν ὁδῶν αὐτοῦ εἰς ἔργα αὐτοῦ,
\vs{23}πρὸ τοῦ αἰῶνος ἐθεμελίωσέ με, ἐν ἀρχῇ πρὸ τοῦ τὴν γῆν ποιῆσαι,
\vs{24}καὶ πρὸ τοῦ τὰς ἀβύσσους ποιῆσαι, πρὸ τοῦ προελθεῖν τὰς πηγὰς τῶν ὑδάτων·
\vs{25}Πρὸ τοῦ ὄρη ἑδρασθῆναι, πρὸ δὲ πάντων βουνῶν, γεννᾷ με.
\vs{26}Κύριος ἐποίησε χώρας καὶ ἀοικήτους, καὶ ἄκρα οἰκούμενα τῆς ὑπʼ οὐρανῶν.
\vs{27}Ἡνίκα ἡτοίμαζε τὸν οὐρανὸν, συμπαρήμην αὐτῷ, καὶ ὅτε ἀφώριζε τὸν ἑαυτοῦ θρόνον ἐπʼ ἀνέμων,
\vs{28}καὶ ὡς ἰσχυρὰ ἐποίει τὰ ἄνω νέφη, καὶ ὡς ἀσφαλεῖς ἐτίθει πηγὰς τῆς ὑπʼ οὐρανὸν,
\vs{29}καὶ ὡς ἰσχυρὰ ἐποίει τὰ θεμέλια τῆς γῆς,
\vs{30}ἤμην παρʼ αὐτῷ ἁρμόζουσα· ἐγὼ ἤμην ᾗ προσέχαιρε· καθʼ ἡμέραν δὲ εὐφραινόμην ἐν προσώπῳ αὐτοῦ ἐν παντὶ καιρῷ,
\vs{31}ὅτε ἐνευφραίνετο τὴν οἰκουμένην συντελέσας, καὶ ἐνευφραίνετο ἐν υἱοῖς ἀνθρώπων.

\vs{32}Νῦν οὖν υἱὲ ἄκουέ μου,
\vs{34}μακάριος ἀνὴρ ὃς εἰσακούσεταί μου, καὶ ἄνθρωπος ὃς τὰς ἐμὰς ὁδοὺς φυλάξει, ἀγρυπνῶν ἐπʼ ἐμαῖς θύραις καθʼ ἡμέραν, τηρῶν σταθμοὺς ἐμῶν εἰσόδων.
\vs{35}Αἱ γὰρ ἔξοδοί μου, ἔξοδοι ζωῆς, καὶ ἑτοιμάζεται θέλησις παρὰ Κυρίου.
\vs{36}Οἱ δὲ ἁμαρτάνοντες εἰς ἐμὲ, ἀσεβοῦσιν εἰς τὰς ἑαυτῶν ψυχὰς, καὶ οἱ μισοῦντές με ἀγαπῶσι θάνατον.

\ch{9}
Ἡ σοφία ᾠκοδόμησεν ἑαυτῇ οἶκον, καὶ ὑπήρεισε στύλους ἑπτά.
\vs{2}Ἔσφαξε τὰ ἑαυτῆς θύματα, ἐκέρασεν εἰς κρατῆρα τὸν ἑαυτῆς οἶνον, καὶ ἡτοιμάσατο τὴν ἑαυτῆς τράπεζαν.
\vs{3}Ἀπέστειλε τοὺς ἑαυτῆς δούλους, συγκαλοῦσα μετὰ ὑψηλοῦ κηρύγματος ἐπὶ κρατῆρα, λέγουσα,
\vs{4}Ὅς ἐστιν ἄφρων, ἐκκλινάτω πρὸς μέ· καὶ τοῖς ἐνδεέσι φρενῶν εἶπεν,
\vs{5}ἔλθατε, φάγετε τῶν ἐμῶν ἄρτων, καὶ πίετε οἶνον ὃν ἐκέρασα ὑμῖν.

\vs{6}Ἀπολείπετε ἀφροσύνην, ἵνα εἰς τὸν αἰῶνα βασιλεύσητε· καὶ ζητήσατε φρόνησιν, καὶ κατορθώσατε ἐν γνώσει σύνεσιν.
\vs{7}Ὁ παιδεύων κακοὺς λήψεται ἑαυτῷ ἀτιμίαν· ἐλέγχων δὲ τὸν ἀσεβῆ μωμήσεται ἑαυτόν.
\vs{8}Μὴ ἔλεγχε κακοὺς, ἵνα μὴ μισήσωσί σε· ἔλεγχε σοφὸν, καὶ ἀγαπήσει σε.
\vs{9}Δίδου σοφῷ ἀφορμὴν, καὶ σοφώτερος ἔσται· γνώριζε δικαίῳ, καὶ προσθήσει τοῦ δέχεσθαι.
\vs{10}Ἀρχὴ σοφίας φόβος Κυρίου, καὶ βουλὴ ἁγίων σύνεσις·
\vs{10a}τὸ γὰρ γνῶναι νόμον, διανοίας ἐστὶν ἀγαθῆς.
\vs{11}Τούτῳ γὰρ τῷ τρόπῳ πολὺν ζήσεις χρόνον, καὶ προστεθήσεταί σοι ἔτη ζωῆς σου.

\vs{12}Υἱὲ ἐὰν σοφὸς γένῃ σεαυτῷ, σοφὸς ἔσῃ καὶ τοῖς πλησίον· ἐὰν δὲ κακὸς ἀποβῇς, μόνος ἂν ἀντλήσεις κακά·
\vs{12a}ὃς ἐρείδεται ἐπὶ ψεύδεσιν, οὗτος ποιμαίνει ἀνέμους, ὁ δʼ αὐτὸς διώξεται ὄρνεα πετόμενα·
\vs{12b}ἀπέλιπε γὰρ ὁδοὺς τοῦ ἑαυτοῦ ἀμπελῶνος, τοὺς δὲ ἄξονας τοῦ ἰδίου γεωργίου πεπλάνηται·
\vs{12c}διαπορεύεται δὲ διʼ ἀνύδρου ἐρήμου, καὶ γῆν διατεταγμένην ἐν διψώδεσι, συνάγει δὲ χερσὶν ἀκαρπίαν.

\vs{13}Γυνὴ ἄφρων καὶ θρασεῖα ἐνδεὴς ψωμοῦ γίνεται, ἣ οὐκ ἐπίσταται αἰσχύνην.
\vs{14}Ἐκάθισεν ἐπὶ θύραις τοῦ ἑαυτῆς οἴκου, ἐπὶ δίφρου ἐμφανῶς ἐν πλατείαις,
\vs{15}προσκαλουμένη τοὺς παριόντας καὶ κατευθύνοντας ἐν ταῖς ὁδοῖς αὐτῶν·
\vs{16}Ὅς ἐστιν ὑμῶν ἀφρονέστατος, ἐκκλινάτω πρὸς μέ· καὶ τοῖς ἐνδεέσι φρονήσεως παρακελεύομαι, λέγουσα,
\vs{17}ἄρτων κρυφίων ἡδέως ἅψασθε, καὶ ὕδατος κλοπῆς γλυκεροῦ.

\vs{18}Ὁ δὲ οὐκ οἶδεν ὅτι γηγενεῖς παρʼ αὐτῇ ὄλλυνται, καὶ ἐπὶ πέταυρον ᾅδου συναντᾷ·
\vs{18a}ἀλλὰ ἀποπήδησον, μὴ χρονίσῃς ἐν τῷ τόπῳ, μηδὲ ἐπιστήσῃς τὸ σὸν ὄμμα πρὸς αὐτὴν,
\vs{18b}οὕτως γὰρ διαβήσῃ ὕδωρ ἀλλότριον·
\vs{18c}ἀπὸ δὲ ὕδατος ἀλλοτρίου ἀπόσχου, καὶ ἀπὸ πηγῆς ἀλλοτρίας μὴ πίῃς
\vs{18d}ἵνα πολὺν ζησῃς χρόνον, προστεθῇ δέ σοι ἔτη ζωῆς.

\ch{10}
Υἱὸς σοφὸς εὐφραίνει πατέρα, υἱὸς δὲ ἄφρων λύπη τῇ μητρί.
\vs{2}Οὐκ ὠφελήσουσι θησαυροὶ ἀνόμους, δικαιοσύνη δὲ ῥύσεται ἐκ θανάτου.
\vs{3}Οὐ λιμοκτονήσει Κύριος ψυχὴν δικαίαν, ζωὴν δὲ ἀσεβῶν ἀνατρέψει.

\vs{4}Πενία ἄνδρα ταπεινοῖ, χεῖρες δὲ ἀνδρείων πλουτίζουσιν·
\vs{4a}υἱὸς πεπαιδευμένος σοφὸς ἔσται, τῷ δὲ ἄφρονι διακόνῳ χρήσεται.
\vs{5}Διεσώθη ἀπὸ καύματος υἱὸς νοήμων, ἀνεμόφθορος δὲ γίνεται ἐν ἀμητῷ υἱὸς παράνομος.

\vs{6}Εὐλογία Κυρίου ἐπὶ κεφαλὴν δικαίου, στόμα δὲ ἀσεβῶν καλύψει πένθος ἄωρον.
\vs{7}Μνήμη δικαίων μετʼ ἐγκωμίων, ὄνομα δὲ ἀσεβοῦς σβέννυται.
\vs{8}Σοφὸς καρδίᾳ δέξεται ἐντολὰς, ὁ δὲ ἄστεγος χείλες σκολιάζων ὑποσκελισθήσεται.
\vs{9}Ὃς πορεύεται ἁπλῶς, πορεύεται πεποιθώς· ὁ δὲ διαστρέφων τὰς ὁδοὺς αὐτοῦ, γνωσθήσεται.
\vs{10}Ὁ ἐννεύων ὀφθαλμοῖς μετὰ δόλου, συνάγει ἀνδράσι λύπας· ὁ δὲ ἐλέγχων μετὰ παῤῥησίας, εἰρηνοποιεῖ.
\vs{11}Πηγὴ ζωῆς ἐν χειρὶ δικαίου, στόμα δὲ ἀσεβοῦς καλύψει ἀπώλεια.

\vs{12}Μῖσος ἐγείρει νεῖκος, πάντας δὲ τοὺς μὴ φιλονεικοῦντας καλύπτει φιλία.
\vs{13}Ὃ ἐκ χειλέων προφέρει σοφίαν, ῥάβδῳ τύπτει ἄνδρα ἀκάρδιον.
\vs{14}Σοφοὶ κρύψουσιν αἴσθησιν, στόμα δὲ προπετοῦς ἐγγίζει συντριβῇ.
\vs{15}Κτῆσις πλουσίων πόλις ὀχυρὰ, συντριβὴ δὲ ἀσεβῶν πενία.
\vs{16}Ἔργα δικαίων ζωὴν ποιεῖ, καρποὶ δὲ ἀσεβῶν ἁμαρτίας.
\vs{17}Ὁδοὺς δικαίας ζωῆς φυλάσσει παιδεία, παιδεία δὲ ἀνεξέλεγκτος πλανᾶται.

\vs{18}Καλύπτουσιν ἔχθραν χείλη δίκαια, οἱ δὲ ἐκφέροντες λοιδορίας ἀφρονέστατοί εἰσιν.
\vs{19}Ἐκ πολυλογίας οὐκ ἐκφεύξῃ ἁμαρτίαν, φειδόμενος δὲ χειλέων νοήμων ἔσῃ.
\vs{20}Ἄργυρος πεπυρωμένος γλῶσσα δικαίου, καρδία δὲ ἀσεβοῦς ἐκλείψει.
\vs{21}Χείλη δικαίων ἐπίσταται ὑψηλὰ, οἱ δὲ ἄφρονες ἐν ἐνδείᾳ τελευτῶσιν.
\vs{22}Εὐλογία Κυρίου ἐπὶ κεφαλὴν δικαίου, αὕτη πλουτίζει, καὶ οὐ μὴ προστεθῇ αὐτῇ λύπη ἐν καρδίᾳ.

\vs{23}Ἐν γέλωτι ἄφρων πράσσει κακὰ, ἡ δὲ σοφία ἀνδρὶ τίκτει φρόνησιν.

\vs{24}Ἐν ἀπωλείᾳ ἀσεβὴς περιφέρεται, ἐπιθυμία δὲ δικαίου δεκτή.
\vs{25}Παραπορευομένης καταιγίδος ἀφανίζεται ἀσεβὴς, δίκαιος δὲ ἐκκλίνας σώζεται εἰς τὸν αἰῶνα.
\vs{26}Ὥσπερ ὄμφαξ ὀδοῦσι βλαβερὸν, καὶ καπνὸς ὄμμασιν, οὕτως παρανομία τοῖς χρωμένοις αὐτῇ.
\vs{27}Φόβος Κυρίου προστίθησιν ἡμέρας, ἔτη δὲ ἀσεβῶν ὀλιγωθήσεται.
\vs{28}Ἐγχρονίζει δικαίοις εὐφροσύνη, ἐλπὶς δὲ ἀσεβῶν ἀπολεῖται.
\vs{29}Ὀχύρωμα ὁσίου φόβος Κυρίου, συντριβὴ δὲ τοῖς ἐργαζομένοις κακά.

\vs{30}Δίκαιος εἰς τὸν αἰῶνα οὐκ ἐνδώσει, ἀσεβεῖς δὲ οὐκ οἰκήσουσι γῆν.
\vs{31}Στόμα δικαίου ἀποστάζει σοφίαν, γλῶσσα δὲ ἀδίκου ἐξολεῖται.
\vs{32}Χείλη ἀνδρῶν δικαίων ἀποστάζει χάριτας, στόμα δὲ ἀσεβῶν ἀποστρέφεται.

\ch{11}
Ζυγοί δόλιοι βδέλυγμα ἐνώπιον Κυρίου, στάθμιον δὲ δίκαιον δεκτὸν αὐτῷ.
\vs{2}Οὗ ἐὰν εἰσέλθῃ ὕβρις, ἐκεῖ καὶ ἀτιμία· στόμα δὲ ταπεινῶν μελετᾷ σοφίαν.
\vs{3}Ἀποθανὼν δίκαιος ἔλιπε μετάμελον, πρόχειρος δὲ γίνεται καὶ ἐπίχαρτος ἀσεβῶν ἀπώλεια.
\vs{5}Δικαιοσύνη ἀμώμους ὀρθοτομεῖ ὁδοὺς, ἀσέβεια δὲ περιπίπτει ἀδικίᾳ.

\vs{6}Δικαιοσύνη ἀνδρῶν ὀρθῶν ῥύεται αὐτούς, τῇ δὲ ἀπωλείᾳ αὐτῶν ἁλίσκονται παράνομοι.
\vs{7}Τελευτήσαντος ἀνδρὸς δικαίου, οὐκ ὄλλυται ἐλπίς, τὸ δὲ καύχημα τῶν ἀσεβῶν ὄλλυται.
\vs{8}Δίκαιος ἐκ θήρας ἐκδύνει, ἀντʼ αὐτοῦ δὲ παραδίδοται ὁ ἀσεβής.
\vs{9}Ἐν στόματι ἀσεβῶν παγὶς πολίταις, αἴσθησις δὲ δικαίων εὔοδος.
\vs{10}Ἐν ἀγαθοῖς δικαίων κατώρθωσε πόλις,
\vs{11}στόμασι δὲ ἀσεβῶν κατεσκάφη.

\vs{12}Μυκτηρίζει πολίτας ἐνδεὴς φρενῶν, ἀνὴρ δὲ φρόνιμος ἡσυχίαν ἄγει.
\vs{13}Ἀνὴρ δίγλωσσος ἀποκαλύπτει βουλὰς ἐν συνεδρίῳ, πιστὸς δὲ πνοῇ κρύπτει πράγματα.
\vs{14}Οἷς μὴ ὑπάρχει κυβέρνησις, πίπτουσιν ὥσπερ φύλλα, σωτηρία δὲ ὑπάρχει ἐν πολλῇ βουλῇ.

\vs{15}Πονηρὸς κακοποιεῖ ὅταν συνμίξῃ δικαίῳ, μισεῖ δὲ ἦχον ἀσφαλείας.
\vs{16}Γυνὴ εὐχάριστος ἐγείρει ἀνδρὶ δόξαν, θρόνος δὲ ἀτιμίας γυνὴ μισοῦσα δίκαια· πλούτου ὀκνηροὶ ἐνδεεῖς γίνονται, οἱ δὲ ἀνδρεῖοι ἐρείδονται πλούτῳ.
\vs{17}Τῇ ψυχῇ αὐτοῦ ἀγαθὸν ποιεῖ ἀνὴρ ἐλεήμων, ἐξολλύει δὲ αὐτοῦ σῶμα ὁ ἀνελεήμων.

\vs{18}Ἀσεβὴς ποιεῖ ἔργα ἄδικα, σπέρμα δὲ δικαίων μισθὸς ἀληθείας.
\vs{19}Υἱὸς δίκαιος γεννᾶται εἰς ζωὴν, διωγμὸς δὲ ἀσεβοῦς εἰς θάνατον.
\vs{20}Βδέλυγμα Κυρίῳ διεστραμμέναι ὁδοὶ, προσδεκτοὶ δὲ αὐτῷ πάντες ἄμωμοι ἐν ταῖς ὁδοῖς αὐτῶν.
\vs{21}Χειρὶ χεῖρας ἐμβαλὼν ἀδίκως οὐκ ἀτιμώρητος ἔσται, ὁ δὲ σπείρων δικαιοσύνην λήψεται μισθὸν πιστόν.
\vs{22}Ὥσπερ ἐνώτιον ἐν ῥινὶ ὑός, οὕτως γυναικὶ κακόφρονι κάλλος.
\vs{23}Ἐπιθυμία δικαίων πᾶσα ἀγαθὴ, ἐλπὶς δὲ ἀσεβῶν ἀπολεῖται.

\vs{24}Εἰσὶν, οἳ τὰ ἴδια σπείροντες πλείονα ποιοῦσιν· εἰσὶδέ καὶ οἳ συνάγοντες ἐλαττονοῦνται.
\vs{25}Ψυχὴ εὐλογουμένη πᾶσα ἁπλῇ, ἀνὴρ δὲ θυμώδης οὐκ εὐσχήμων.
\vs{26}Ὁ συνέχων σῖτον ὑπολείποιτο αὐτὸν τοῖς ἔθνεσιν· εὐλογία δὲ εἰς κεφαλὴν τοῦ μεταδιδόντος.
\vs{27}Τεκταινόμενος ἀγαθὰ ζητεῖ χάριν ἀγαθὴν, ἐκζητοῦντα δὲ κακὰ καταλήψεται αὐτόν.
\vs{28}Ὁ πεποιθὼς ἐπὶ πλούτῳ οὗτος πεσεῖται, ὁ δὲ ἀντιλαμβανόμενος δικαίων οὗτος ἀνατελεῖ.
\vs{29}Ὁ μὴ συμπεριφερόμενος τῷ ἑαυτοῦ οἴκῳ, κληρονομήσει ἄνεμον, δουλεύσει δὲ ἄφρων φρονίμῳ.
\vs{30}Ἐκ καρποῦ δικαιοσύνης φύεται δένδρὅν ζωῆς, ἀφαιροῦνται δὲ ἄωροι ψυχαὶ παρανόμων.
\vs{31}Εἰ ὁ μὲν δίκαιος μόλις σώζεται, ὁ ἀσεβὴς καὶ ἁμαρτωλὸς ποῦ φανεῖται;

\ch{12}
Ὁ ἀγαπῶν παιδείαν, ἀγαπᾷ αἴσθησιν· ὁ δὲ μισῶν ἐλέγχους, ἄφρων.
\vs{2}Κρείσσων ὁ εὑρὼν χάριν παρὰ Κυρίῳ, ἀνὴρ δὲ παράνομος παρασιωπηθήσεται.
\vs{3}Οὐ κατορθώσει ἄνθρωπος ἐξ ἀνόμου, αἱ δὲ ῥίζαι τῶν δικαίων οὐκ ἐξαρθήσονται.
\vs{4}Γυνὴ ἀνδρεία στέφανος τῷ ἀνδρὶ αὐτῆς· ὥσπερ δὲ ἐν ξύλῳ σκώληξ, οὕτως ἄνδρα ἀπόλλυσι γυνὴ κακοποιός.

\vs{5}Λογισμοὶ δικαίων κρίματα, κυβερνῶσι δὲ ἀσεβεῖς δόλους.

\vs{6}Λόγοι ἀσεβῶν δόλιοι, στόμα δὲ ὀρθῶν ῥύσεται αὐτούς.
\vs{7}Οὗ ἐὰν στραφῇ ὁ ἀσεβὴς, ἀφανίζεται, οἶκοι δὲ δικαίων παραμένουσι·
\vs{8}Στόμα συνετοῦ ἐγκωμιάζεται ὑπὸ ἀνδρὸς, νωθροκάρδιος δὲ μυκτηρίζεται.
\vs{9}Κρείσσων ἀνὴρ ἐν ἀτιμίᾳ δουλεύων ἑαυτῷ, ἢ τιμὴν ἑαυτῷ περιτιθεὶς καὶ προσδεόμενος ἄρτου.

\vs{10}Δίκαιος οἰκτείρει ψυχὰς κτηνῶν αὐτοῦ, τὰ δὲ σπλάγχνα τῶν ἀσεβῶν ἀνελεήμονα.
\vs{11}Ὁ ἐργαζόμενος τὴν ἑαυτοῦ γῆν, ἐμπλησθήσεται ἄρτων, οἱ δὲ διώκοντες μάταια, ἐνδεεῖς φρενῶν·
\vs{11a}ὅς ἐστιν ἡδὺς ἐν οἴνων διατριβαῖς, ἐν τοῖς ἑαυτοῦ ὀχυρώμασι καταλείψει ἀτιμίαν.

\vs{12}Ἐπιθυμίαι ἀσεβῶν κακαὶ, αἱ δὲ ῥίζαι τῶν εὐσεβῶν ἐν ὀχυρώμασι.
\vs{13}Διʼ ἁμαρτίαν χειλέων ἐμπίπτει εἰς παγίδας ἁμαρτωλὸς, ἐκφεύγει δὲ ἐξ αὐτῶν δίκαιος·
\vs{13a}ὁ βλέπων λεῖα ἐλεηθήσεται, ὁ δὲ συναντῶν ἐν πύλαις ἐκθλίψει ψυχάς.
\vs{14}Ἀπὸ καρπῶν στόματος ψυχὴ ἀνδρὸς πλησθήσεται ἀγαθῶν, ἀνταπόδομα δὲ χειλέων αὐτοῦ δοθήσεται αὐτῷ.
\vs{15}Ὁδοὶ ἀφρόνων ὀρθαὶ ἐνώπιον αὐτῶν, εἰσακούει δὲ συμβουλίας σοφός.
\vs{16}Ἄφρων αὐθημερὸν ἐξαγγέλλει ὀργὴν αὐτοῦ, κρύπτει δὲ τὴν ἑαυτοῦ ἀτιμίαν ἀνὴρ πανοῦργος.
\vs{17}Ἐπιδεικνυμένην πίστιν ἀπαγγέλλει δίκαιος, ὁ δὲ μάρτυς τῶν ἀδίκων δόλιος.

\vs{18}Εἰσὶν οἳ λέγοντες τιτρώσκουσι, μάχαιραι· γλῶσσαι δὲ σοφῶν ἰῶνται.
\vs{19}Χείλη ἀληθινὰ κατορθοῖ μαρτυρίαν, μάρτυς δὲ ταχὺς γλῶσσαν ἔχει ἄδικον.
\vs{20}Δόλος ἐν καρδίᾳ τεκταινομένου κακὰ, οἱ δὲ βουλόμενοι εἰρήνην εὐφρανθήσονται.
\vs{21}Οὐκ ἀρέσει τῷ δικαίῳ οὐδὲν ἄδικον, οἱ δὲ ἀσεβεῖς πλησθήσονται κακῶν.
\vs{22}Βδέλυγμα Κυρίῳ χείλη ψευδῆ, ὁ δὲ ποιῶν πίστεις δεκτὸς παρʼ αὐτῷ.
\vs{23}Ἀνὴρ συνετὸς θρόνος αἰσθήσεως, καρδία δὲ ἀφρόνων συναντήσεται ἀραῖς.

\vs{24}Χεὶρ ἐκλεκτῶν κρατήσει εὐχερῶς, δόλιοι δὲ ἔσονται ἐν προνομῇ.
\vs{25}Φοβερὸς λόγος καρδίαν ταράσσει ἀνδρὸς δικαίου, ἀγγελία δὲ ἀγαθὴ εὐφραίνει αὐτόν.
\vs{26}Ἐπιγνώμων δίκαιος ἑαυτοῦ φίλος ἔσται, ἁμαρτάνοντας δὲ καταδιώξεται κακὰ, ἡ δὲ ὁδὸς τῶν ἀσεβῶν πλανήσει αὐτούς.
\vs{27}Οὐκ ἐπιτεύξεται δόλιος θήρας, κτῆμα δὲ τίμιον ἀνὴρ καθαρός.
\vs{28}Ἐν ὁδοῖς δικαιοσύνης ζωὴ, ὁδοὶ δὲ μνησικάκων εἰς θάνατον.

\ch{13}
Υἱὸς πανοῦργος ὑπήκοος πατρὶ, υἱὸς δὲ ἀνήκοος ἐν ἀπωλείᾳ.
\vs{2}Ἀπὸ καρπῶν δικαιοσύνης φάγεται ἀγαθὸς, ψυχαὶ δὲ παρανόμων ὀλοῦνται ἄωροι.
\vs{3}Ὃς φυλάσσει τὸ ἑαυτοῦ στόμα τηρεῖ τὴν ἑαυτοῦ ψυχὴν, ὁ δὲ προπετὴς χείλεσι πτοήσει ἑαυτόν.
\vs{4}Ἐν ἐπιθυμίαις ἐστὶ πᾶς ἀεργὸς, χεῖρες δὲ ἀνδρείων ἐν ἐπιμελείᾳ.
\vs{5}Λόγον ἄδικον μισεῖ δίκαιος, ἀσεβὴς δὲ αἰσχύνεται, καὶ οὐκ ἕξει παῤῥησίαν.
\vs{7}Εἰσὶν οἱ πλουτίζοντες ἑαυτοὺς μηδὲν ἔχοντες, καὶ εἰσὶν οἱ ταπεινοῦντες ἑαυτοὺς ἐν πολλῷ πλούτῳ.

\vs{8}Λύτρον ἀνδρὸς ψυχῆς ὁ ἴδιος πλοῦτος, πτωχὸς δὲ οὐχ ὑφίσταται ἀπειλήν.
\vs{9}Φῶς δικαίοις διαπαντὸς, φῶς δὲ ἀσεβῶν σβέννυται·
\vs{9a}ψυχαὶ δόλιαι πλανῶνται ἐν ἁμαρτίαις, δίκαιοι δὲ οἰκτείρουσι καὶ ἐλεοῦσι.
\vs{10}Κακὸς μεθʼ ὕβρεως πράσσει κακὰ, οἱ δʼ ἑαυτῶν ἐπιγνώμονες σοφοί.
\vs{11}Ὕπαρξις ἐπισπουδαζομένη μετὰ ἀνομίας, ἐλάσσων γίνεται, ὁ δὲ συνάγων ἑαυτῷ μετʼ εὐσεβείας πληθυνθήσεται· δίκαιος οἰκτείρει καὶ κιχρᾷ.
\vs{12}Κρείσσων ἐναρχόμενος βοηθῶν καρδίᾳ, τοῦ ἐπαγγελλομένου καὶ εἰς ἐλπίδα ἄγοντος· δένδρον γὰρ ζωῆς, ἐπιθυμία ἀγαθή.
\vs{13}Ὃς καταφρονεῖ πράγματος, καταφρονηθήσεται ὑπʼ αὐτοῦ· ὁ δὲ φοβούμενος ἐντολὴν, οὗτος ὑγιαίνει·
\vs{13a}υἱῷ δολίῳ οὐδὲν ἔσται ἀγαθὸν, οἰκέτῃ δὲ σοφῷ εὔοδοι ἔσονται πράξεις, καὶ κατευθυνθήσεται ἡ ὁδὸς αὐτοῦ.

\vs{14}Νόμος σοφοῦ πηγὴ ζωῆς, ὁ δὲ ἄνους ὑπὸ παγίδος θανεῖται.
\vs{15}Σύνεσις ἀγαθὴ δίδωσι χάριν, τὸ δὲ γνῶναι νόμον διανοίας ἐστὶν ἀγαθῆς, ὁδοὶ δὲ καταφρονούντων ἐν ἀπωλείᾳ.

\vs{16}Πᾶς πανοῦργος πράσσει μετὰ γνώσεως, ὁ δὲ ἄφρων ἐξεπέτασεν ἑαυτοῦ κακίαν.
\vs{17}Βασιλεὺς θρασὺς ἐμπεσεῖται εἰς κακὰ, ἄγγελος δὲ σοφὸς ῥύσεται αὐτόν.
\vs{18}Πενίαν καὶ ἀτιμίαν ἀφαιρεῖται παιδεία, ὁ δὲ φυλάσσων ἐλέγχους δοξασθήσεται.
\vs{19}Ἐπιθυμίαι εὐσεβῶν ἡδύνουσι ψυχὴν, ἔργα δὲ ἀσεβῶν μακρὰν ἀπὸ γνώσεως.
\vs{20}Συμπορευόμενος σοφοῖς σοφὸς ἔσῃ, ὁ δὲ συμπορευόμενος ἄφροσι γνωσθήσεται.
\vs{21}Ἁμαρτάνοντας καταδιώξεται κακὰ, τοὺς δὲ δικαίους καταλήψεται ἀγαθά.
\vs{22}Ἀγαθὸς ἀνὴρ κληρονομήσει υἱοὺς υἱῶν, θησαυρίζεται δὲ δικαίοις πλοῦτος ἀσεβῶν.
\vs{23}Δίκαιοι ποιήσουσιν ἐν πλούτῳ ἔτη πολλὰ, ἄδικοι δὲ ἀπολοῦνται συντόμως.

\vs{24}Ὃς φείδεται τῆς βακτηρίας, μισεῖ τὸν υἱὸν αὐτοῦ ὁ δὲ ἀγαπῶν, ἐπιμελῶς παιδεύει.
\vs{25}Δίκαιος ἔσθων ἐμπιπλᾷ τὴν ψυχὴν αὐτοῦ, ψυχαὶ δὲ ἀσεβῶν ἐνδεεῖς.

\ch{14}
Σοφαὶ γυναῖκες ᾠκοδόμησαν οἴκους, ἡ δὲ ἄφρων κατέσκαψε ταῖς χερσὶν αὐτῆς.
\vs{2}Ὁ πορευόμενος ὀρθῶς φοβεῖται τὸν Κύριον, ὁ δὲ σκολιάζων ταῖς ὁδοῖς αὐτοῦ ἀτιμασθήσεται.
\vs{3}Ἐκ στόματος ἀφρόνων βακτηρία ὕβρεως, χείλη δὲ σοφῶν φυλάσσει αὐτούς.
\vs{4}Οὗ μή εἰσι βόες, φάτναι καθαραί· οὗ δὲ πολλὰ γεννήματα, φανερὰ βοὸς ἰσχύς.
\vs{5}Μάρτυς πιστὸς οὐ ψεύδεται, ἐκκαίει δὲ ψευδῆ μάρτυς ἄδικος.
\vs{6}Ζητήσεις σοφίαν παρὰ κακοῖς καὶ οὐχ εὑρήσεις, αἴσθησις δὲ παρὰ φρονίμοις εὐχερής.

\vs{7}Πάντα ἐναντία ἀνδρὶ ἄφρονι, ὅπλα δὲ αἰσθήσεως χείλη σοφά.
\vs{8}Σοφία πανούργων ἐπιγνώσεται τὰς ὁδοὺς αὐτῶν, ἄνοια δὲ ἀφρόνων ἐν πλάνῃ.
\vs{9}Οἰκίαι παρανόμων ὀφειλήσουσι καθαρισμὸν, οἰκίαι δὲ δικαίων δεκταί.

\vs{10}Καρδία ἀνδρὸς αἰσθητικὴ, λυπηρὰ ψυχὴ αὐτοῦ, ὅταν δὲ εὐφραίνηται οὐκ ἐπιμίγνυται ὕβρει.
\vs{11}Οἰκίαι ἀσεβῶν ἀφανισθήσονται, σκηναὶ δὲ κατορθούντων στήσονται.
\vs{12}Ἔστιν ὁδὸς ἣ δοκεῖ παρὰ ἀνθρώποις ὀρθὴ εἶναι, τὰ δὲ τελευταῖα αὐτῆς ἔρχεται εἰς πυθμένα ᾅδου.
\vs{13}Ἐν εὐφροσύναις οὐ προσμίγνυται λύπη, τελευταῖα δὲ χαρὰ εἰς πένθος ἔρχεται.
\vs{14}Τῶν ἑαυτοῦ ὁδῶν πλησθήσεται θρασυκάρδιος, ἀπὸ δὲ τῶν διανοημάτων αὐτοῦ ἀνὴρ ἀγαθός.
\vs{15}Ἄκακος πιστεύει παντὶ λόγῳ, πανοῦργος δὲ ἔρχεται εἰς μετάνοιαν.
\vs{16}Σοφὸς φοβηθεὶς ἐξέκλινεν ἀπὸ κακοῦ, ὁ δὲ ἄφρων ἑαυτῷ πεποιθὼς μίγνυται ἀνόμῳ.
\vs{17}Ὀξύθυμος πράσσει μετὰ ἀβουλίας, ἀνὴρ δὲ φρόνιμος πολλὰ ὑποφέρει.

\vs{18}Μεριοῦνται ἄφρονες κακίαν, οἱ δὲ πανοῦργοι κρατήσουσιν αἰσθήσεως.
\vs{19}Ὀλισθήσουσι κακοὶ ἔναντι ἀγαθῶν, καὶ ἀσεβεῖς θεραπεύσουσι θύρας δικαίων.
\vs{20}Φίλοι μισήσουσι φίλους πτωχούς, φίλοι δὲ πλουσίων πολλοί.
\vs{21}Ὁ ἀτιμάζων πένητας ἁμαρτάνει, ἐλεῶν δὲ πτωχοὺς μακαριστός.
\vs{22}Πλανώμενοι τεκταίνουσι κακά, ἔλεον δὲ καὶ ἀλήθειαν τεκταίνουσιν ἀγαθοί· οὐκ ἐπίστανται ἔλεον καὶ πίστιν τέκτονες κακῶν, ἐλεημοσύναι δὲ καὶ πίστεις παρὰ τέκτοσιν ἀγαθοῖς.
\vs{23}Ἐν παντὶ μεριμνῶντι ἔνεστι περισσόν, ὁ δὲ ἡδὺς καὶ ἀνάλγητος ἐν ἐνδείᾳ ἔσται.
\vs{24}Στέφανος σοφῶν πανοῦργος, ἡ δὲ διατριβὴ ἀφρόνων κακή.

\vs{25}Ῥύσεται ἐκ κακῶν ψυχὴν μάρτυς πιστὸς, ἐκκαίει δὲ ψευδῆ δόλιος.
\vs{26}Ἐν φόβῳ Κυρίου ἐλπὶς ἰσχύος, τοῖς δὲ τέκνοις αὐτοῦ καταλείπει ἔρεισμα.
\vs{27}Πρόσταγμα Κυρίου πηγὴ ζωῆς, ποιεῖ δὲ ἐκκλίνειν ἐκ παγίδος θανάτου.

\vs{28}Ἐν πολλῷ ἔθνει δόξα βασιλέως, ἐν δὲ ἐκλείψει λαοῦ συντριβὴ δυνάστου.
\vs{29}Μακρόθυμος ἀνὴρ πολὺς ἐν φρονήσει, ὁ δὲ ὀλιγόψυχος ἰσχυρῶς ἄφρων.
\vs{30}Πρᾳΰθυμος ἀνὴρ καρδίας ἰατρὸς, σὴς δὲ ὀστέων καρδία αἰσθητική·
\vs{31}Ὁ συκοφαντῶν πένητα παροξύνει τὸν ποιήσαντα αὐτὸν, ὁ δὲ τιμῶν αὐτὸν ἐλεεῖ πτωχόν.
\vs{32}Ἐν κακίᾳ αὐτοῦ ἀπωσθήσεται ἀσεβής, ὁ δὲ πεποιθὼς τῇ ἑαυτοῦ ὁσιότητι δίκαιος.

\vs{33}Ἐν καρδίᾳ ἀγαθῇ ἀνδρὸς σοφία, ἐν δὲ καρδίᾳ ἀφρόνων οὐ διαγινώσκεται.
\vs{34}Δικαιοσύνη ὑψοῖ ἔθνος, ἐλασσονοῦσι δὲ φυλὰς ἁμαρτίαι.
\vs{35}Δεκτὸς βασιλεῖ ὑπηρέτης νοήμων, τῇ δὲ ἑαυτοῦ εὐστροφίᾳ ἀφαιρεῖται ἀτιμίαν.

\ch{15}
Ὀργὴ ἀπόλλυσι καὶ φρονίμους, ἀπόκρισις δὲ ὑποπίπτουσα ἀποστρέφει θυμὸν, λόγος δὲ λυπηρὸς ἐγείρει ὀργάς.
\vs{2}Γλῶσσα σοφῶν καλὰ ἐπίσταται, στόμα δὲ ἀφρόνων ἀναγγέλλει κακά.

\vs{3}Ἐν παντὶ τόπῳ ὀφθαλμοὶ Κυρίου σκοπεύουσι κακούς τε καὶ ἀγαθούς.
\vs{4}Ἴασις γλώσσης δένδρον ζωῆς, ὁ δὲ συντηρῶν αὐτὴν πλησθήσεται πνεύματος.
\vs{5}Ἄφρων μυκτηρίζει παιδείαν πατρὸς, ὁ δὲ φυλάσσων ἐντολὰς, πανουργότερος· ἐν πλεοναζούσῃ δικαιοσύνῃ ἰσχὺς πολλὴ, οἱ δὲ ἀσεβεῖς ὁλόῤῥιζοι ἐκ γῆς ἀπολοῦνται.

\vs{6}Οἴκοις δικαίων ἰσχὺς πολλή, καρποὶ δὲ ἀσεβῶν ἀπολοῦνται.
\vs{7}Χείλη σοφῶν δέδεται αἰσθήσει, καρδίαι δὲ ἀφρόνων οὐκ ἀσφαλεῖς.
\vs{8}Θυσίαι ἀσεβῶν βδέλυγμα Κυρίῳ, εὐχαὶ δὲ κατευθυνόντων δεκταὶ παρʼ αὐτῷ.
\vs{9}Βδέλυγμα Κυρίῳ ὁδοὶ ἀσεβοῦς, διώκοντας δὲ δικαιοσύνην ἀγαπᾷ.
\vs{10}Παιδεία ἀκάκου γνωρίζεται ὑπὸ τῶν παριόντων, οἱ δὲ μισοῦντες ἐλέγχους τελευτῶσιν αἰσχρῶς.

\vs{11}Ἅδης καὶ ἀπώλεια φανερὰ παρὰ τῷ Κυρίῳ· πῶς οὐχὶ καὶ αἱ καρδίαι τῶν ἀνθρώπων;
\vs{12}Οὐκ ἀγαπήσει ἀπαίδευτος τοὺς ἐλέγχοντας αὐτόν, μετὰ δὲ σοφῶν οὐχ ὁμιλήσει.
\vs{13}Καρδίας εὐφραινομένης πρόσωπον θάλλει, ἐν δὲ λύπαις οὔσης σκυθρωπάζει.
\vs{14}Καρδία ὀρθὴ ζητεῖ αἴσθησιν, στόμα δὲ ἀπαιδεύτων γνώσεται κακά.

\vs{15}Πάντα τὸν χρόνον οἱ ὀφθαλμοὶ τῶν κακῶν προσδέχονται κακὰ, οἱ δὲ ἀγαθοὶ ἡσυχάζουσι διαπαντός.
\vs{16}Κρεῖσσον μικρὰ μερὶς μετὰ φόβου Κυρίου, ἢ θησαυροὶ μεγάλοι μετὰ ἀφοβίας.
\vs{17}Κρείσσων ξενισμὸς μετὰ λαχάνων πρὸς φιλίαν καὶ χάριν, ἢ παράθεσις μόσχων μετὰ ἔχθρας.
\vs{18}Ἀνὴρ θυμώδης παρασκευάζει μάχας· μακρόθυμος δὲ καὶ τὴν μέλλουσαν καταπρᾳΰνει·
\vs{18a}μακρόθυμος ἀνὴρ κατασβέσει κρίσεις, ὁ δὲ ἀσεβὴς ἐγείρει μᾶλλον.
\vs{19}Ὁδοὶ ἀεργῶν ἐστρωμέναι ἀκάνθαις, αἱ δὲ τῶν ἀνδρείων τετριμμέναι.
\vs{20}Υἱὸς σοφὸς εὐφραίνει πατέρα, υἱὸς δὲ ἄφρων μυκτηρίζει μητέρα αὐτοῦ.
\vs{21}Ἀνοήτου τρίβοι ἐνδεεῖς φρενῶν, ἀνὴρ δὲ φρόνιμος κατευθύνων πορεύεται.
\vs{22}Ὑπερτίθενται λογισμοὺς οἱ μὴ τιμῶντες συνέδρια, ἐν δὲ καρδίαις βουλευομένων μένει βουλή.

\vs{23}Οὐ μὴ ὑπακούσει ὁ κακὸς αὐτῇ, οὐδὲ μὴ εἴπῃ καίριόν τι καὶ καλὸν τῷ κοινῷ.

\vs{24}Ὁδοὶ ζωῆς διανοήματα συνετοῦ, ἵνα ἐκκλίνας ἐκ τοῦ ᾅδου σωθῇ.
\vs{25}Οἴκους ὑβριστῶν κατασπᾷ Κύριος, ἐστήρισε δὲ ὅριον χήρας.
\vs{26}Βδέλυγμα Κυρίῳ λογισμὸς ἄδικος, ἁγνῶν δὲ ῥήσεις σεμναί.
\vs{27}Ἐξόλλυσιν ἑαυτὸν ὁ δωρολήπτης, ὁ δὲ μισῶν δώρων λήψεις σώζεται·
\vs{27a}ἐλεημοσύναις καὶ πίστεσιν ἀποκαθαίρονται ἁμαρτίαι, τῷ δὲ φόβῳ Κυρίου ἐκκλίνει πᾶς ἀπὸ κακοῦ.

\vs{28}Καρδίαι δικαίων μελετῶσι πίστεις, στόμα δὲ ἀσεβῶν ἀποκρίνεται κακά·
\vs{28a}δεκταὶ παρὰ Κυρίῳ ὁδοὶ ἀνθρώπων δικαίων, διὰ δὲ αὐτῶν καὶ οἱ ἐχθροὶ φίλοι γίνονται.
\vs{29}Μακρὰν ἀπέχει ὁ Θεὸς ἀπὸ ἀσεβῶν, εὐχαῖς δὲ δικαίων ἐπακούει·
\vs{29a}κρείσσων ὀλίγη λῆψις μετὰ δικαιοσύνης, ἢ πολλὰ γεννήματα μετὰ ἀδικίας.

\vs{29b}Καρδία ἀνδρὸς λογιζέσθω δίκαια, ἵνα ὑπὸ τοῦ Θεοῦ διορθωθῇ τὰ διαβήματα αὐτοῦ.
\vs{30}Θεωρῶν ὀφθαλμὸς καλὰ εὐφραίνει καρδίαν, φημη δὲ ἀγαθὴ πιαίνει ὀστᾶ.
\vs{32}Ὃς ἀπωθεῖται παιδείαν, μισεῖ ἑαυτὸν· ὁ δὲ τηρῶν ἐλέγχους, ἀγαπᾷ ψυχὴν αὐτοῦ.
\vs{33}Φόβος Κυρίου παιδεία καὶ σοφία, καὶ ἀρχὴ δόξης ἀποκριθήσεται αὐτῇ.

\ch{16}\vs{2}Πάντα τὰ ἔργα τοῦ ταπεινοῦ φανερὰ παρὰ τῷ Θεῷ, οἱ δὲ ἀσεβεῖς ἐν ἡμέρᾳ κακῇ ὀλοῦνται.
\vs{5}Ἀκάθαρτος παρὰ Θεῷ πᾶς ὑψηλοκάρδιος, χειρὶ δὲ χεῖρας ἐμβαλὼν ἀδίκως οὐκ ἀθωωθήσεται·
\vs{7}ἀρχὴ ὁδοῦ ἀγαθῆς τὸ ποιεῖν τὰ δίκαια, δεκτὰ δὲ παρὰ Θεῷ μᾶλλον ἢ θύειν θυσίας·
\vs{8}ὁ ζητῶν τὸν Κύριον εὑρήσει γνῶσιν μετὰ δικαιοσύνης, οἱ δὲ ὀρθῶς ζητοῦντες αὐτὸν εὑρήσουσιν εἰρήνην.
\vs{9}Πάντα τὰ ἔργα τοῦ Κυρίου μετὰ δικαιοσύνης, φυλάσσεται δὲ ὁ ἀσεβὴς εἰς ἡμέραν κακήν.

\vs{10}Μαντεῖον ἐπὶ χείλεσι βασιλέως, ἐν δὲ κρίσει οὐ μὴ πλανηθῇ τὸ στόμα αὐτοῦ.
\vs{11}Ῥοπὴ ζυγοῦ δικαιοσύνη παρὰ Κυρίῳ, τὰ δὲ ἔργα αὐτοῦ στάθμια δίκαια.
\vs{12}Βδέλυγμα βασιλεῖ ὁ ποιῶν κακὰ, μετὰ γὰρ δικαιοσύνης ἑτοιμάζεται θρόνος ἀρχῆς.
\vs{13}Δεκτὰ βασιλεῖ χείλη δίκαια, λόγους δέ ὀρθοὺς ἀγαπᾷ.
\vs{14}Θυμὸ βασιλέως ἄγγελος θανάτου, ἀνὴρ δὲ σοφὸς ἐξιλάσεται αὐτόν.
\vs{15}Ἐν φωτὶ ζωῆς υἱὸς βασιλέως, οἱ δὲ προσδεκτοὶ αὐτῷ ὥσπερ νέφος ὄψιμον.
\vs{16}Νοσσιαὶ σοφίας αἱρετώτεραι χρυσίου, νοσσιαὶ δὲ φρονήσεως αἱρετώτεραι ὑπὲρ ἀργύριον.
\vs{17}Τρίβοι ζωῆς ἐκκλίνουσιν ἀπὸ κακῶν, μῆκος δὲ βίου ὁδοὶ δικαιοσύνης. Ὁ δεχόμενος παιδείαν ἐν ἀγαθοῖς ἔσται, ὁ δὲ φυλάσσων ἐλέγχους σοφισθήσεται· ὃς φυλάσσει τὰς ἑαυτοῦ ὁδοὺς, τηρεῖ τὴν ἑαυτοῦ ψυχήν· ἀγαπῶν δὲ ζωὴν αὐτοῦ, φείσεται στόματος αὐτοῦ.

\vs{18}Πρὸ συντριβῆς ἡγεῖται ὕβρις, πρὸ δὲ πτώματος κακοφροσύνη.
\vs{19}Κρείσσων πρᾳΰθυμος μετὰ ταπεινώσεως, ἢ ὃς διαιρεῖται σκῦλα μετὰ ὑβριστῶν.
\vs{20}Συνετὸς ἐν πράγμασιν εὑρετὴς ἀγαθῶν, πεποιθὼς δὲ ἐπὶ Θεῷ μακαριστός.
\vs{21}Τοὺς σοφοὺς καὶ συνετοὺς φαύλους καλοῦσιν, οἱ δὲ γλυκεῖς ἐν λόγῳ πλείονα ἀκούσονται.
\vs{22}Πηγὴ ζωῆς ἔννοια τοῖς κεκτημένοις, παιδεία δὲ ἀφρόνων κακή.
\vs{23}Καρδία σοφοῦ νοήσει τὰ ἀπὸ τοῦ ἰδίου στόματος, ἐπὶ δὲ χείλεσι φορέσει ἐπιγνωμοσύνην·
\vs{24}Κηρία μέλιτος λόγοι καλοί, γλύκασμα δὲ αὐτοῦ ἴασις ψυχῆς.

\vs{25}Εἰσὶν ὁδοὶ δοκοῦσαι εἶναι ὀρθαὶ ἀνδρὶ, τὰ μέντοι τελευταῖα αὐτῶν βλέπει εἰς πυθμένα ᾅδου.
\vs{26}Ἀνὴρ ἐν πόνοις πονεῖ ἑαυτῷ, καὶ ἐκβιάζεται τὴν ἀπώλειαν ἑαυτοῦ. Ὁ μέντοι σκολιὸς ἐπὶ τῷ ἑαυτοῦ στόματι φορεῖ τὴν ἀπώλειαν·
\vs{27}ἀνὴρ ἄφρων ὀρύσσει ἑαυτῷ κακὰ, ἐπὶ δὲ τῶν ἑαυτοῦ χειλέων θησαυρίζει πῦρ.
\vs{28}Ἀνὴρ σκολιὸς διαπέμπεται κακὰ, καὶ λαμπτῆρα δόλου πυρσεύσει κακοῖς, καὶ διαχωρίζει φίλους.
\vs{29}Ἀνὴρ παράνομος ἀποπειρᾶται φίλων, καὶ ἀπάγει αὐτοὺς ὁδοὺς οὐκ ἀγαθάς.

\vs{30}Στηρίζων δὲ ὀφθαλμοὺς αὐτοῦ διαλογίζεται διεστραμμένα, ὁρίζει δὲ τοῖς χείλεσιν αὐτοῦ πάντα τὰ κακά· οὗτος κάμινός ἐστι κακίας.
\vs{31}Στέφανος καυχήσεως γῆρας, ἐν δὲ ὁδοῖς δικαιοσύνης εὑρίσκεται.
\vs{32}Κρείσσων ἀνὴρ μακρόθυμος ἰσχυροῦ, ὁ δὲ κρατῶν ὀργῆς κρείσσων καταλαμβανομένου πόλιν.
\vs{33}Εἰς κόλπους ἐπέρχεται πάντα τοῖς ἀδίκοις, παρὰ δὲ Κυρίου πάντα τὰ δίκαια.

\ch{17}
Κρείσσων ψωμὸς μεθʼ ἡδονῆς ἐν εἰρήνῃ, ἢ οἶκος πολλῶν ἀγαθῶν καὶ ἀδίκων θυμάτων μετὰ μάχης.
\vs{2}Οἰκέτης νοήμων κρατήσει δεσποτῶν ἀφρόνων, ἐν δὲ ἀδελφοῖς διελεῖται μέρη.
\vs{3}Ὥσπερ δοκιμάζεται ἐν καμίνῳ ἄργυρος καὶ χρυσὸς, οὕτως ἐκλεκταὶ καρδίαι παρὰ Κυρίῳ.
\vs{4}Κακὸς ὑπακούει γλώσσης παρανόμων, δίκαιος δὲ οὐ προσέχει χείλεσι ψευδέσιν.
\vs{5}Ὁ καταγελῶν πτωχοῦ παροξύνει τὸν ποιήσαντα αὐτὸν, ὁ δὲ ἐπιχαίρων ἀπολλυμένῳ οὐκ ἀθωωθήσεται, ὁ δὲ ἐπισπλαγχνιζόμενος ἐλεηθήσεται.

\vs{6}Στέφανος γερόντων τέκνα τέκνων, καύχημα δὲ τέκνων πατέρες αὐτῶν·
\vs{6a}τοῦ πιστοῦ ὅλος ὁ κόσμος τῶν χρημάτων, τοῦ δὲ ἀπίστου οὐδὲ ὀβολός.
\vs{7}Οὐχ ἁρμόσει ἄφρονι χείλη πιστὰ, οὐδὲ δικαίῳ χείλη ψευδῆ.
\vs{8}Μισθὸς χαρίτων παιδεία τοῖς χρωμένοις, οὗ δʼ ἂν ἐπιστρέψῃ εὐοδωθήσεται.
\vs{9}Ὃς κρύπτει ἀδικήματα, ζητεῖ φιλίαν· ὃς δὲ μισεῖ κρύπτειν, διΐστησι φίλους καὶ οἰκείους.
\vs{10}Συντρίβει ἀπειλὴ καρδίαν φρονίμου, ἄφρων δὲ μαστιγωθεὶς οὐκ αἰσθάνεται.
\vs{11}Ἀντιλογίας ἐγείρει πᾶς κακὸς, ὁ δὲ Κύριος ἄγγελον ἀνελεήμονα ἐκπέμψει αὐτῷ.

\vs{12}Ἐμπεσεῖται μέριμνα ἀνδρὶ νοήμονι, οἱ δὲ ἄφρονες διαλογιοῦνται κακά.
\vs{13}Ὃς ἀποδίδωσι κακὰ ἀντὶ ἀγαθῶν, οὐ κινηθήσεται κακὰ ἐκ τοῦ οἴκου αὐτοῦ.
\vs{14}Ἐξουσίαν δίδωσι λόγοις ἀρχὴ δικαιοσύνης, προηγεῖται δὲ τῆς ἐνδείας στάσις καὶ μάχη.
\vs{15}Ὃς δίκαιον κρίνει τὸν ἄδικον, ἄδικον δὲ τὸν δίκαιον, ἀκάθαρτος καὶ βδελυκτὸς παρὰ Θεῷ.
\vs{16}Ἱνατί ὑπῆρξε χρήματα ἄφρονι; κτήσασθαι γὰρ σοφίαν ἀκάρδιος οὐ δυνήσεται·
\vs{16a}ὃς ὑψηλὸν ποιεῖ τὸν ἑαυτοῦ οἶκον, ζητεῖ συντριβήν· ὁ δὲ σκολιάζων τοῦ μαθεῖν, ἐμπεσεῖται εἰς κακά.
\vs{17}Εἰς πάντα καιρὸν φίλος ὑπαρχέτω σοι, ἀδελφοὶ δὲ ἐν ἀνάγκαις χρήσιμοι ἔστωσαν, τούτου γὰρ χάριν γεννῶνται.
\vs{18}Ἀνὴρ ἄφρων ἐπικροτεῖ καὶ ἐπιχαίρει ἑαυτῷ, ὡς καὶ ὁ ἐγγυώμενος ἐγγύῃ τῶν ἑαυτοῦ φίλων.

\vs{19}Φιλαμαρτήμων χαίρει μάχαις,
\vs{20}ὁ δὲ σκληροκάρδιος οὐ συναντᾷ ἀγαθοῖς· ἀνὴρ εὐμετάβολος γλώσσῃ ἐμπεσεῖται εἰς κακὰ,
\vs{21}καρδία δὲ ἄφρονος ὀδύνη τῷ κεκτημένῳ αὐτήν· οὐκ εὐφραίνεται πατὴρ ἐφʼ υἱῷ ἀπαιδεύτῳ, υἱὸς δὲ φρόνιμος εὐφραίνει μητέρα αὐτοῦ.
\vs{22}Καρδία εὐφραινομένη εὐεκτεῖν ποιεῖ, ἀνδρὸς δὲ λυπηροῦ ξηραίνεται τὰ ὀστᾶ.
\vs{23}Λαμβάνοντος δῶρα ἀδίκως ἐν κόλποις οὐ κατευοδοῦνται ὁδοὶ, ἀσεβὴς δὲ ἐκκλίνει ὁδοὺς δικαιοσύνης.
\vs{24}Πρόσωπον συνετὸν ἀνδρὸς σοφοῦ, οἱ δὲ ὀφθαλμοὶ τοῦ ἄφρονος ἐπʼ ἄκρα γῆς.
\vs{25}Ὀργὴ πατρὶ υἱὸς ἄφρων, καὶ ὀδύνη τῇ τεκούσῃ αὐτόν.

\vs{26}Ζημιοῦν ἄνδρα δίκαιον οὐ καλὸν, οὐδὲ ὅσιον ἐπιβουλεύειν δυνάσταις δικαίοις.
\vs{27}Ὃς φείδεται ῥῆμα προέσθαι σκληρὸν, ἐπιγνώμων· μακρόθυμος δὲ ἀνὴρ φρόνιμος.
\vs{28}Ἀνοήτῳ ἐπερωτήσαντι σοφίαν σοφία λογισθήσεται, ἐνεὸν δέ τις ἑαυτὸν ποιήσας, δόξει φρόνιμος εἶναι.

\ch{18}
Προφάσεις ζητεῖ ἀνὴρ βουλόμενος χωρίζεσθαι ἀπὸ φίλων, ἐν παντὶ δὲ καιρῷ ἐπονείδιστος ἔσται.
\vs{2}Οὐ χρείαν ἔχει σοφίας ἐνδεὴς φρενῶν, μᾶλλον γᾶρ ἄγεται ἀφροσύνῃ.
\vs{3}Ὅταν ἔλθῃ ἀσεβὴς εἰς βάθος κακῶν, καταφρονεῖ, ἐπέρχεται δὲ αὐτῷ ἀτιμία καὶ ὄνειδος.
\vs{4}Ὕδωρ βαθὺ λόγος ἐν καρδίᾳ ἀνδρὸς, ποταμὸς δὲ ἀναπηδύει καὶ πηγὴ ζωῆς.
\vs{5}Θαυμάσαι πρόσωπον ἀσεβοῦς οὐ καλὸν, οὐδὲ ὅσιον ἐκκλίνειν τὸ δίκαιον ἐν κρίσει.

\vs{6}Χείλη ἄφρονος ἄγουσιν αὐτὸν εἰς κακὰ, τὸ δὲ στόμα αὐτοῦ τὸ θρασὺ θάνατον ἐπικαλεῖται.
\vs{7}Στόμα ἄφρονος συντριβὴ αὐτῷ, τὰ δὲ χείλη αὐτοῦ παγὶς τῇ ψυχῇ αὐτοῦ.
\vs{8}Ὀκνηροὺς καταβάλλει φόβος, ψυχαὶ δὲ ἀνδρογύνων πεινάσουσιν.
\vs{9}Ὁ μὴ ἰώμενος αὐτὸν ἐν τοῖς ἔργοις αὐτοῦ, ἀδελφός ἐστι τοῦ λυμαινομένου ἑαυτόν.
\vs{10}Ἐκ μεγαλωσύνης ἰσχύος ὄνομα Κυρίου, αὐτῷ δὲ προσδραμόντες δίκαιοι ὑψοῦνται.
\vs{11}Ὕπαρξις πλουσίου ἀνδρὸς πόλις ὀχυρὰ, ἡ δὲ δόξα αὐτῆς μέγα ἐπισκιάζει.
\vs{12}Πρὸ συντριβῆς ὑψοῦται καρδία ἀνδρὸς, καὶ πρὸ δόξης ταπεινοῦνται.
\vs{13}Ὃς ἀποκρίνεται λόγον πρὶν ἀκοῦσαι, ἀφροσύνη αὐτῷ ἐστι καὶ ὄνειδος.
\vs{14}Θυμὸν ἀνδρὸς πρᾳΰνει θεράπων φρόνιμος, ὀλιγόψυχον δὲ ἄνδρα τίς ὑποίσει;
\vs{15}Καρδία φρονίμου κτᾶται αἴσθησιν, ὦτα δὲ σοφῶν ζητεῖ ἔννοιαν.
\vs{16}Δόμα ἀνθρώπου ἐμπλατύνει αὐτὸν, καὶ παρὰ δυνάσταις καθιζάνει αὐτόν.
\vs{17}Δίκαιος ἑαυτοῦ κατήγορος ἐν πρωτολογίᾳ, ὡς δʼ ἂν ἐπιβάλῃ ὁ ἀντίδικος ἐλέγχεται.

\vs{18}Ἀντιλογίας παύει σιγηρὸς, ἐν δὲ δυναστείαις ὁρίζει.
\vs{19}Ἀδελφὸς ὑπὸ ἀδελφοῦ βοηθούμενος, ὡς πόλις ὀχυρὰ καὶ ὑψηλὴ, ἰσχύει δὲ ὥσπερ τεθεμελιωμένον βασίλειον.
\vs{20}Ἀπὸ καρπῶν στόματος ἀνὴρ πίμπλησι κοιλίαν αὐτοῦ, ἀπὸ δὲ καρπῶν χειλέων αὐτοῦ ἐμπλησθήσεται.
\vs{21}Θάνατος καὶ ζωὴ ἐν χειρὶ γλώσσης, οἱ δὲ κρατοῦντες αὐτῆς ἔδονται τοὺς καρποὺς αὐτῆς.
\vs{22}Ὃς εὗρε γυναῖκα ἀγαθὴν, εὗρε χάριτας, ἔλαβε δὲ παρὰ Θεοῦ ἱλαρότητα·
\vs{22a}ὃς ἐκβάλλει γυναῖκα ἀγαθὴν, ἐκβάλλει τὰ ἀγαθὰ, ὁ δὲ κατέχων μοιχαλίδα, ἄφρων καὶ ἀσεβής.

\ch{19}
\vs{3}Ἀφροσύνη ἀνδρὸς λυμαίνεται τὰς ὁδοὺς αὐτοῦ, τὸν δὲ Θεὸν αἰτιᾶται τῇ καρδίᾳ αὐτοῦ.

\vs{4}Πλοῦτος προστίθησι φίλους πολλοὺς, ὁ δὲ πτωχὸς καὶ ἀπὸ τοῦ ὑπάρχοντος φίλου λείπεται.
\vs{5}Μάρτυς ψευδὴς οὐκ ἀτιμώρητος ἔσται, ὁ δὲ ἐγκαλῶν ἀδίκως οὐ διαφεύξεται.
\vs{6}Πολλοὶ θεραπεύουσι πρόσωπα βασιλέων, πᾶς δὲ ὁ κακὸς γίνεται ὄνειδος ἀνδρί.
\vs{7}Πᾶς ὃς ἀδελφὸν πτωχὸν μισεῖ, καὶ φιλίας μακρὰν ἔσται· ἔννοια ἀγαθὴ τοῖς εἰδόσιν αὐτὴν ἐγγιεῖ, ἀνὴρ δὲ φρόνιμος εὑρήσει αὐτήν· ὁ πολλὰ κακοποιῶν τελεσιουργεῖ κακίαν, ὃς δὲ ἐρεθίζει λόγους, οὐ σωθήσεται.

\vs{8}Ὁ κτώμενος φρόνησιν ἀγαπᾷ ἑαυτὸν, ὃς δὲ φυλάσσει φρόνησιν, εὑρήσει ἀγαθά.
\vs{9}Μάρτυς ψευδὴς οὐκ ἀτιμώρητος ἔσται, ὃς δʼ ἂν ἐκκαύσῃ κακίαν, ἀπολεῖται ὑπʼ αὐτῆς.
\vs{10}Οὐ συμφέρει ἄφρονι τρυφὴ, καὶ ἐὰν οἰκέτης ἄρξηται μεθʼ ὕβρεως δυναστεύειν.
\vs{11}Ἐλεήμων ἀνὴρ μακροθυμεῖ, τὸ δὲ καύχημα αὐτοῦ ἐπέρχεται παρανόμοις.
\vs{12}Βασιλέως ἀπειλὴ ὁμοία βρυγμῷ λέοντος· ὥσπερ δὲ δρόσος ἐπὶ χόρτῳ, οὕτως τὸ ἱλαρὸν αὐτοῦ.

\vs{13}Αἰσχύνη πατρὶ υἱὸς ἄφρων, οὐχ ἁγναὶ εὐχαὶ ἀπὸ μισθώματος ἑταίρας.
\vs{14}Οἶκον καὶ ὕπαρξιν μερίζουσι πατέρες παισὶ, παρὰ δὲ Κυρίου ἁρμόζεται γυνὴ ἀνδρί.
\vs{15}Δειλία κατέχει ἀνδρόγυνον, ψυχὴ δὲ ἀεργοῦ πεινάσει.
\vs{16}Ὃς φυλάσσει ἐντολὴν, τηρεῖ τὴν ἑαυτοῦ ψυχήν· ὁ δὲ καταφρονῶν τῶν ἑαυτοῦ ὁδῶν, ἀπολεῖται.
\vs{17}Δανείζει Θεῷ ὁ ἐλεῶν πτωχὸν, κατὰ δὲ τὸ δόμα αὐτοῦ ἀνταποδώσει αὐτῷ.
\vs{18}Παίδευε υἱόν σου, οὕτως γὰρ ἔσται εὔελπις, εἰς δὲ ὕβριν μὴ ἐπαίρου τῇ ψυχῇ σου.
\vs{19}Κακόφρων ἀνὴρ πολλὰ ζημιωθήσεται· ἐὰν δὲ λοιμεύηται, καὶ τὴν ψυχὴν αὐτοῦ προσθήσει.

\vs{20}Ἄκουε, υἱὲ, παιδείαν πατρός σου, ἵνα σοφὸς γένῃ ἐπʼ ἐσχάτων σου.
\vs{21}Πολλοὶ λογισμοὶ ἐν καρδίᾳ ἀνδρὸς, ἡ δὲ βουλὴ τοῦ Κυρίου εἰς τὸν αἰῶνα μένει.
\vs{22}Καρπὸς ἀνδρὶ ἐλεημοσύνη, κρείσσων δὲ πτωχὸς δίκαιος ἢ πλούσιος ψευδής.
\vs{23}Φόβος Κυρίου εἰς ζωὴν ἀνδρὶ· ὁ δὲ ἄφοβος αὐλισθήσεται ἐν τόποις οὗ οὐκ ἐπισκοπεῖται γνῶσις.
\vs{24}Ὁ ἐγκρύπτων εἰς τὸν κόλπον αὐτοῦ χεῖρας ἀδίκως, οὐδὲ τῷ στόματι οὐ μὴ προσενείκῃ αὐτάς.
\vs{25}Λοιμοῦ μαστιγουμένου, ἄφρων πανουργότερος γίνεται· ἐὰν δὲ ἐλέγχῃς ἄνδρα φρόνιμον, νοήσει αἴσθησιν.

\vs{26}Ὁ ἀτιμάζων πατέρα καὶ ἀπωθούμενος μητέρα αὐτοῦ, καταισχυνθήσεται καὶ ἐπονείδιστος ἔσται.
\vs{27}Υἱὸς ἀπολειπόμενος φυλάξαι παιδείαν πατρὸς, μελετήσει ῥήσεις κακάς.
\vs{28}Ὁ ἐγγυώμενος παῖδα ἄφρονα, καθυβρίσει δικαίωμα· στόμα δὲ ἀσεβῶν καταπίεται κρίσεις.
\vs{29}Ἑτοιμάζονται ἀκολάστοις μάστιγες, καὶ τιμωρίαι ὁμοίως ἄφροσιν.

\ch{20}
Ἀκολάστον οἶνος, καὶ ὑβριστικὸν μέθη, πᾶς δὲ ἄφρων τοιούτοις συμπλέκεται.
\vs{2}Οὐ διαφέρει ἀπειλὴ βασιλέως θυμοῦ λέοντος, ὁ δὲ παροξύνων αὐτὸν ἁμαρτάνει εἰς τὴν ἑαυτοῦ ψυχήν.
\vs{3}Δόξα ἀνδρὶ ἀποστρέφεσθαι λοιδορίας, πᾶς δὲ ἄφρων τοιούτοις συμπλέκεται.
\vs{4}Ὀνειδιζόμενος ὀκνηρὸς οὐκ αἰσχύνεται, ὡσαύτως καὶ ὁ δανειζόμενος σῖτον ἐν ἀμητῷ.

\vs{5}Ὕδωρ βαθὺ βουλὴ ἐν καρδίᾳ ἀνδρὸς, ἀνὴρ δὲ φρόνιμος ἐξαντλήσει αὐτήν.
\vs{6}Μέγα ἄνθρωπος, καὶ τίμιον ἀνὴρ ἐλεήμων, ἄνδρα δὲ πιστὸν ἔργον εὑρεῖν.
\vs{7}Ὃς ἀναστρέφεται ἄμωμος ἐν δικαιοσύνῃ, μακαρίους τοὺς παῖδας αὐτοῦ καταλείψει.
\vs{8}Ὅταν βασιλεὺς δίκαιος καθίσῃ ἐπὶ θρόνου, οὐκ ἐναντιοῦται ἐν ὀφθαλμοῖς αὐτοῦ πᾶν πονηρόν.
\vs{9}Τίς καυχήσεται ἁγνὴν ἔχειν τὴν καρδίαν; ἢ τίς παῤῥησιάσεται καθαρὸς εἶναι ἀπὸ ἁμαρτιῶν;
\vs{9a}Κακολογοῦντος πατέρα ἢ μητέρα σβεσθήσεται λαμπτὴρ, αἱ δὲ κόραι τῶν ὀφθαλμῶν αὐτοῦ ὄψονται σκότος.

\vs{9b}Μερὶς ἐπισπουδαζομένη ἐν πρώτοις, ἐν τοῖς τελευταίοις οὐκ εὐλογηθήσεται.
\vs{9c}Μὴ εἴπῃς, τίσομαι τὸν ἐχθρὸν, ἀλλʼ ὑπόμεινον τὸν Κύριον, ἵνα σοι βοηθήσῃ.

\vs{10}Στάθμιον μέγα καὶ μικρὸν, καὶ μέτρα δισσὰ, ἀκάθαρτα ἐνώπιον Κυρίου καὶ ἀμφότερα, καὶ ὁ ποιῶν αὐτά.
\vs{11}Ἐν τοῖς ἐπιτηδεύμασιν αὐτοῦ συμποδισθήσεται νεανίσκος μετὰ ὁσίου, καὶ εὐθεῖα ἡ ὁδὸς αὐτοῦ.
\vs{12}Οὖς ἀκούει, καὶ ὀφθαλμὸς ὁρᾷ, Κυρίου ἔργα καὶ ἀμφότερα.
\vs{13}Μὴ ἀγάπα καταλαλεῖν, ἵνα μὴ ἐξαρθῇς· διάνοιξον τοὺς ὀφθαλμούς σου, καὶ ἐμπλήσθητι ἄρτων.

\vs{23}Βδέλυγμα Κυρίῳ δισσὸν στάθμιον, καὶ ζυγὸς δόλιος οὐ καλὸν ἐνώπιον αὐτοῦ.
\vs{24}Παρὰ Κυρίου εὐθύνεται τὰ διαβήματα ἀνδρὶ, θνητὸς δὲ πῶς ἂν νοήσαι τὰς ὁδοὺς αὐτοῦ;
\vs{25}Παγὶς ἀνδρὶ ταχύ τι τῶν ἰδίων ἁγιάσαι, μετὰ γὰρ τὸ εὔξασθαι μετανοεῖν γίνεται.
\vs{26}Λικμήτωρ ἀσεβῶν βασιλεὺς σοφὸς, καὶ ἐπιβαλεῖ αὐτοῖς τροχόν.

\vs{27}Φῶς Κυρίου πνοὴ ἀνθρώπων, ὃς ἐρευνᾷ ταμιεῖα κοιλίας.
\vs{28}Ἐλεημοσύνη καὶ ἀλήθεια φυλακὴ βασιλεῖ, καὶ περικυκλώσουσιν ἐν δικαιοσύνῃ τὸν θρόνον αὐτοῦ.
\vs{29}Κόσμος νεανίαις σοφία, δόξα δὲ πρεσβυτέρων πολιαί.
\vs{30}Ὑπώπια καὶ συντρίμματα συναντᾷ κακοῖς, πληγαὶ δὲ εἰς ταμιεῖα κοιλίας.

\ch{21}
Ὥσπερ ὁρμὴ ὕδατος, οὕτως καρδία βασιλέως ἐν χειρὶ Θεοῦ, οὗ ἐὰν θέλων νεῦσαι ἐκεῖ ἔκλινεν αὐτήν.
\vs{2}Πᾶς ἀνὴρ φαίνεται ἑαυτῷ δίκαιος, κατευθύνει δὲ καρδίας Κύριος.
\vs{3}Ποιεῖν δίκαια καὶ ἀληθεύειν, ἀρεστὰ παρὰ Θεῷ μᾶλλον ἢ θυσιῶν αἷμα.
\vs{4}Μεγαλόφρων ἐν ὕβρει θρασυκάρδιος, λαμπτὴρ δὲ ἀσεβῶν ἁμαρτία.
\vs{6}Ὁ ἐνεργῶν θησαυρίσματα γλώσσῃ ψευδεῖ, μάταια διώκει ἐπὶ παγίδας θανάτου.
\vs{7}Ὄλεθρος ἀσεβέσιν ἐπιξενωθήσεται, οὐ γὰρ βούλονται πράσσειν τὰ δίκαια.
\vs{8}Πρὸς τοὺς σκολιοὺς σκολιὰς ὁδοὺς ἀποστέλλει ὁ Θεὸς, ἁγνὰ γὰρ καὶ ὀρθὰ τὰ ἔργα αὐτοῦ.
\vs{9}Κρεῖσσον οἰκεῖν ἐπὶ γωνίας ὑπαίθρου, ἢ ἐν κεκονιαμένοις μετὰ ἀδικίας καὶ ἐν οἴκῳ κοινῷ.
\vs{10}Ψυχὴ ἀσεβοῦς οὐκ ἐλεηθήσεται ὑπʼ οὐδενὸς τῶν ἀνθρώπων.
\vs{11}Ζημιουμένου ἀκολάστου πανουργότερος γίνεται ὁ ἄκακος, συνιῶν δὲ σοφὸς δέξεται γνῶσιν.
\vs{12}Συνιεῖ δίκαιος καρδίας ἀσεβῶν, καὶ φαυλίζει ἀσεβεῖς ἐν κακοῖς.

\vs{13}Ὃς φράσσει τὰ ὦτα αὐτοῦ τοῦ μὴ ἐπακοῦσαι ἀσθενοῦς, καὶ αὐτὸς ἐπικαλέσεται καὶ οὐκ ἔσται ὁ εἰσακούων.
\vs{14}Δόσις λάθριος ἀνατρέπει ὀργάς, δώρων δὲ ὁ φειδόμενος θυμὸν ἐγείρει ἰσχυρόν.
\vs{15}Εὐφροσύνη δικαίων ποιεῖν κρίμα, ὅσιος δὲ ἀκάθαρτος παρὰ κακούργοις.
\vs{16}Ἀνὴρ πλανώμενος ἐξ ὁδοῦ δικαιοσύνης, ἐν συναγωγῇ γιγάντων ἀναπαύσεται.
\vs{17}Ἀνὴρ ἐνδεὴς ἀγαπᾷ εὐφροσύνην, φιλῶν οἶνον καὶ ἔλαιον εἰς πλοῦτον·
\vs{18}Περικάθαρμα δὲ δικαίου ἄνομος.
\vs{19}Κρεῖσσον οἰκεῖν ἐν τῇ ἐρήμῳ, ἢ μετὰ γυναικὸς μαχίμου καὶ γλωσσώδους καὶ καὶ ὀργίλου.
\vs{20}Θησαυρὸς ἐπιθυμητὸς ἀναπαύσεται ἐπὶ στόματος σοφοῦ, ἄφρονες δὲ ἄνδρες καταπίονται αὐτόν.
\vs{21}Ὁδὸς δικαιοσύνης καὶ ἐλεημοσύνης εὑρήσει ζωὴν καὶ δόξαν.
\vs{22}Πόλεις ὀχυρὰς ἐπέβη σοφὸς, καὶ καθεῖλε τὸ ὀχύρωμα ἐφʼ ᾧ ἐπεποίθεισαν οἱ ἀσεβεῖς.
\vs{23}Ὃς φυλάσσει τὸ στόμα αὐτοῦ καὶ τὴν γλῶσσαν, διατηρεῖ ἐκ θλίψεως τὴν ψυχὴν αὐτοῦ.

\vs{24}Θρασὺς καὶ αὐθάδης καὶ ἀλαζὼν λοιμὸς καλεῖται, ὃς δὲ μνησικακεῖ παράνομος.
\vs{25}Ἐπιθυμίαι ὀκνηρὸν ἀποκτείνουσιν, οὐ γὰρ προαιροῦνται αἱ χεῖρες αὐτοῦ ποιεῖν τι.
\vs{26}Ἀσεβὴς ἐπιθυμεῖ ὅλην τὴν ἡμέραν ἐπιθυμίας κακὰς, ὁ δὲ δίκαιος ἐλεᾷ καὶ οἰκτείρει ἀφειδῶς.
\vs{27}Θυσίαι ἀσεβῶν βδέλυγμα Κυρίῳ, καὶ γὰρ παρανόμως προσφέρουσιν αὐτάς.
\vs{28}Μάρτυς ψευδὴς ἀπολεῖται, ἀνὴρ δὲ ὑπήκοος φυλασσόμενος λαλήσει.
\vs{29}Ἀσεβὴς ἀνὴρ ἀναιδῶς ὑφίσταται προσώπῳ, ὁ δὲ εὐθὺς αὐτὸς συνιεῖ τὰς ὁδοὺς αὐτοῦ.
\vs{30}Οὐκ ἔστι σοφία, οὐκ ἔστιν ἀνδρεία, οὐκ ἔστι βουλὴ πρὸς τὸν ἀσεβῆ.
\vs{31}Ἵππος ἑτοιμάζεται εἰς ἡμέραν πολέμου, παρὰ δὲ Κυρίου ἡ βοήθεια.

\ch{22}
Αἱρετώτερον ὄνομα καλὸν ἢ πλοῦτος πολύς, ὑπὲρ δὲ ἀργύριον καὶ χρυσίον χάρις ἀγαθή.
\vs{2}Πλούσιος καὶ πτωχὸς συνήντησαν ἀλλήλοις, ἀμφοτέρους δὲ ὁ κύριος ἐποίησε.
\vs{3}Πανοῦργος ἰδὼν πονηρὸν τιμωρούμενον κραταιῶς αὐτὸς παιδεύεται, οἱ δὲ ἄφρονες παρελθόντες ἐζημιώθησαν.
\vs{4}Γενεὰ σοφίας φόβος Κυρίου, καὶ πλοῦτος, καὶ δόξα, καὶ ζωή.
\vs{5}Τρίβολοι καὶ παγίδες ἐν ὁδοῖς σκολιαῖς, ὁ δὲ φυλάσσων τὴν ἑαυτοῦ ψυχὴν ἀφέξεται αὐτῶν.
\vs{7}Πλούσιοι πτωχῶν ἄρξουσι, καὶ οἰκέται ἰδίοις δεσπόταις δανειοῦσιν.

\vs{8}Ὁ σπείρων φαῦλα θερίσει κακὰ, πληγὴν δὲ ἔργων αὐτοῦ συντελέσει·
\vs{8a}ἄνδρα ἱλαρὸν καὶ δότην εὐλογεῖ ὁ Θεὸς, ματαιότητα δὲ ἔργων αὐτοῦ συντελέσει.
\vs{9}Ὁ ἐλεῶν πτωχὸν αὐτὸς διατραφήσεται, τῶν γὰρ ἑαυτοῦ ἄρτων ἔδωκε τῷ πτωχῷ·
\vs{9a}νίκην καὶ τιμὴν περιποιεῖται ὁ δῶρα δοὺς, τὴν μέντοι ψυχὴν ἀφαιρεῖται τῶν κεκτημένων.
\vs{10}Ἔκβαλε ἐκ συνεδρίου λοιμὸν, καὶ συνεξελεύσεται αὐτῷ νεῖκος, ὅταν γὰρ καθίσῃ ἐν συνεδρίῳ πάντας ἀτιμάζει.

\vs{11}Ἀγαπᾷ Κύριος ὁσίας καρδίας, δεκτοὶ δὲ αὐτῷ πάντες ἄμωμοι· χείλεσι ποιμαίνει βασιλεύς.
\vs{12}Οἱ δὲ ὀφθαλμοὶ Κυρίου διατηροῦσιν αἴσθησιν, φαυλίζει δὲ λόγους παράνομος.
\vs{13}Προφασίζεται, καὶ λέγει ὀκνηρὸς, λέων ἐν ταῖς ὁδοῖς, ἐν δὲ ταῖς πλατείαις φονευταί.
\vs{14}Βόθρος βαθὺς στόμα παρανόμου, ὁ δὲ μισηθεὶς ὑπὸ Κυρίου ἐμπεσεῖται εἰς αὐτόν.
\vs{14a}εἰσὶν ὁδοὶ κακαὶ ἐνώπιον ἀνδρὸς, καὶ οὐκ ἀγαπᾷ τοῦ ἀποστρέψαι ἀπʼ αὐτῶν, ἀποστρέφειν δὲ δεῖ ἀπὸ ὁδοῦ σκολιᾶς καὶ κακῆς.
\vs{15}Ἄνοια ἐξῆπται καρδίας νέου, ῥάβδος δὲ καὶ παιδεία μακρὰν ἀπʼ αὐτοῦ.

\vs{16}Ὁ συκοφαντῶν πένητα, πολλὰ ποιεῖ τὰ ἑαυτοῦ, δίδωσι δὲ πλουσίῳ ἐπʼ ἐλάσσονι.

\vs{17}Λόγοις σοφῶν παράβαλλε σὸν οὖς, καὶ ἄκουε ἐμὸν λόγον, τὴν δὲ σὴν καρδίαν ἐπίστησον, ἵνα γνῷς ὅτι καλοί εἰσι·
\vs{18}καὶ ἐὰν ἐμβάλῃς αὐτοὺς εἰς τὴν καρδίαν σου, εὐφρανοῦσί σε ἅμα ἐπὶ σοῖς χείλεσιν·
\vs{19}Ἵνα σου γένηται ἐπὶ Κύριον ἡ ἐλπὶς, καὶ γνωρίσῃ σοι τὴν ὁδόν σου.
\vs{20}Καὶ σὺ δὲ ἀπόγραψαι αὐτὰ σεαυτῷ τρισσῶς, εἰς βουλὴν καὶ γνῶσιν ἐπὶ τὸ πλάτος τῆς καρδίας σου.
\vs{21}Διδάσκω οὖν σε ἀληθῆ λόγον, καὶ γνῶσιν ἀγαθὴν ὑπακούειν, τοῦ ἀποκρίνεσθαί σε λόγους ἀληθείας τοῖς προβαλλομένοις σοι.

\vs{22}Μὴ ἀποβιάζου πένητα, πτωχὸς γὰρ ἐστι, καὶ μὴ ἀτιμάσῃς ἀσθενῆ ἐν πύλαις.
\vs{23}Ὁ γὰρ Κύριος κρινεῖ αὐτοῦ τὴν κρίσιν, καὶ ῥύσῃ σὴν ἄσυλον ψυχήν.

\vs{24}Μὴ ἴσθι ἑταῖρος ἀνδρὶ θυμώδει, φίλῳ δὲ ὀργίλῳ μὴ συναυλίζου·
\vs{25}μήποτε μάθῃς τῶν ὁδῶν αὐτοῦ, καὶ λάβῃς βρόχους τῇ σῇ ψυχῇ.

\vs{26}Μὴ δίδου σεαυτὸν εἰς ἐγγύην αἰσχυνόμενος πρόσωπον·
\vs{27}Ἐὰν γὰρ μὴ ἔχῃ πόθεν ἀποτίσῃς, λήψονται τὸ στρῶμα τὸ ὑπὸ τὰς πλευράς σου.
\vs{28}Μὴ μέταιρε ὅρια αἰώνια, ἃ ἔθεντο οἱ πατέρες σου.

\vs{29}Ὁρατικὸν ἄνδρα καὶ ὀξὺν ἐν τοῖς ἔργοις αὐτοῦ βασιλεῦσι δεῖ παρεστάναι, καὶ μὴ παρεστάναι ἀνδράσι νωθροῖς.

\ch{23}
Ἐὰν καθίσῃς δειπνεῖν ἐπὶ τραπέζης δυνάστου, νοητῶς νόει τὰ παρατιθέμενά σοι.
\vs{2}Καὶ ἐπίβαλλε τὴν χεῖρά σου, εἰδὼς ὅτι τοιαῦτά σε δεῖ παρασκευάσαι· εἰ δὲ ἀπληστότερος εἶ,
\vs{3}μὴ ἐπιθύμει τῶν ἐδεσμάτων αὐτοῦ, ταῦτα γὰρ ἔχεται ζωῆς ψευδοῦς.

\vs{4}Μὴ παρεκτείνου πένης ὢν πλουσίῳ, τῇ δὲ σῇ ἐννοίᾳ ἀπόσχου.
\vs{5}Ἐὰν ἐπιστήσῃς τὸ σὸν ὄμμα πρὸς αὐτὸν, οὐδαμοῦ φανεῖται· κατεσκεύασται γὰρ αὐτῷ πτέρυγες ὥσπερ ἀετοῦ, καὶ ὑποστρέφει εἰς τὸν οἶκον τοῦ προεστηκότος αὐτοῦ.
\vs{6}Μὴ συνδείπνει ἀνδρὶ βασκάνῳ, μηδὲ ἐπιθύμει τῶν βρωμάτων αὐτοῦ,
\vs{7}ὃν τρόπον γὰρ εἴ τις καταπίοι τρίχα, οὕτως ἐσθίει καὶ πίνει· μηδὲ πρὸς σὲ εἰσαγάγῃς αὐτὸν, καὶ φάγῃς τὸν ψωμόν σου μετʼ αὐτοῦ,
\vs{8}ἐξεμέσει γὰρ αὐτὸν, καὶ λυμανεῖται τοὺς λόγους σου τοὺς καλούς.

\vs{9}Εἰς ὦτα ἄφρονος μηδὲν λέγε, μήποτε μυκτηρίσῃ τοὺς συνετοὺς λόγους σου.
\vs{10}Μὴ μεταθῇς ὅρια αἰώνια, εἰς δὲ κτῆμα ὀρφανῶν μὴ εἰσέλθῃς·
\vs{11}Ὁ γὰρ λυτρούμενος αὐτοὺς Κύριος, κραταιός ἐστι, καὶ κρινεῖ τὴν κρίσιν αὐτῶν μετὰ σοῦ.
\vs{12}Δὸς εἰς παιδείαν τὴν καρδίαν σου, τὰ δὲ ὦτά σου ἑτοίμασον λόγοις αἰσθήσεως.

\vs{13}Μὴ ἀπόσχῃ νήπιον παιδεύειν, ὅτι ἐὰν πατάξῃς αὐτὸν ῥάβδῳ, οὐ μὴ ἀποθάνῃ.
\vs{14}Συ μὲν γὰρ πατάξεις αὐτὸν ῥάβδῳ, τὴν δὲ ψυχὴν αὐτοῦ ἐκ θανάτου ῥύσῃ.

\vs{15}Υἱὲ, ἐὰν σοφὴ γένηταί σου ἡ καρδία, εὐφρανεῖς καὶ τὴν ἐμὴν καρδίαν,
\vs{16}καὶ ἐνδιατρίψει λόγοις τὰ σὰ χείλη πρὸς τὰ ἐμὰ χείλη ἐὰν ὀρθὰ ὦσι.
\vs{17}Μὴ ζηλούτω ἡ καρδία σου ἁμαρτωλοὺς, ἀλλὰ ἐν φόβῳ Κυρίου ἴσθι ὅλην τὴν ἡμέραν.
\vs{18}Ἐὰν γὰρ τηρήσῃς αὐτὰ, ἔσται σοι ἔκγονα, ἡ δὲ ἐλπίς σου οὐκ ἀποστήσεται.

\vs{19}Ἄκουε υἱὲ, καὶ σοφὸς γίνου, καὶ κατεύθυνε ἐννοίας σῆς καρδίας.
\vs{20}Μὴ ἴσθι οἰνοπότης, μηδὲ ἐκτείνου συμβολαῖς, κρεῶν τε ἀγορασμοῖς.
\vs{21}Πᾶς γὰρ μέθυσος καὶ πορνοκόπος πτωχεύσει, καὶ ἐνδύσεται διεῤῥηγμένα καὶ ῥακώδη πᾶς ὑπνώδης.

\vs{22}Ἄκουε, υἱὲ, πατρὸς τοῦ γεννήσαντός σε, καὶ μὴ καταφρόνει ὅτι γεγήρακέ σου ἡ μήτηρ.
\vs{24}Καλῶς ἐκτρέφει πατὴρ δίκαιος, ἐπὶ δὲ υἱῷ σοφῷ εὐφραίνεται ἡ ψυχὴ αὐτοῦ.
\vs{25}Εὐφραινέσθω ὁ πατὴρ καὶ ἡ μήτηρ ἐπὶ σοὶ, καὶ χαιρέτω ἡ τεκοῦσά σε.

\vs{26}Δός μοι υἱὲ σὴν καρδίαν, οἱ δὲ σοὶ ὀφθαλμοὶ ἐμὰς ὁδοὺς τηρείτωσαν.
\vs{27}Πίθος γὰρ τετρημένος ἐστὶν ἀλλότριος οἶκος, καὶ φρέαρ στενὸν ἀλλότριον.
\vs{28}Οὗτος γὰρ συντόμως ἀπολεῖται, καὶ πᾶς παράνομος ἀναλωθήσεται.

\vs{29}Τίνι οὐαί; τίνι θόρυβος; τίνι κρίσεις; τίνι δὲ ἀηδίαι καὶ λέσχαι; τίνι συντρίμματα διακενῆς; τίνος πελιδνοὶ οἱ ὀφθαλμοί;
\vs{30}Οὐ τῶν ἐγχρονιζόντων ἐν οἴνοις; οὐ τῶν ἰχνευόντων ποῦ πότοι γίνονται; μὴ μεθύσκεσθε ἐν οἴνοις, ἀλλὰ ὁμιλεῖτε ἀνθρώποις δικαίοις καὶ ὁμιλεῖτε ἐν περιπάτοις.
\vs{31}Ἐὰν γὰρ εἰς τὰς φιάλας καὶ τὰ ποτήρια δῷς τοὺς ὀφθαλμούς σου, ὕστερον περιπατήσεις γυμνότερος ὑπέρου.
\vs{32}Τὸ δὲ ἔσχατον ὥσπερ ὑπὸ ὄφεως πεπληγὼς ἐκτείνεται, καὶ ὥσπερ ὑπὸ κεράστου διαχεῖται αὐτῷ ὁ ἰός.

\vs{33}Οἱ ὀφθαλμοί σου ὅταν ἴδωσιν ἀλλοτρίαν, τὸ στόμα σου τότε λαλήσει σκολιά.
\vs{34}Καὶ κατακείσῃ ὥσπερ ἐν καρδίᾳ θαλάσσης, καὶ ὥσπερ κυβερνήτης ἐν πολλῷ κλύδωνι.
\vs{35}Ἐρεῖς δὲ, τύπτουσί με καὶ οὐκ ἐπόνεσα, καὶ ἐνέπαιξάν μοι, ἐγὼ δὲ οὐκ ᾔδειν· πότε ὄρθρος ἔσται, ἵνα ἐλθὼν ζητήσω μεθʼ ὧν συνελεύσομαι;

\ch{24}
Υἱέ, μὴ ζηλώσῃς κακοὺς ἄνδρας, μηδὲ ἐπιθυμήσῃς εἶναι μετʼ αὐτῶν.
\vs{2}Ψευδῆ γὰρ μελετᾷ ἡ καρδία αὐτῶν, καὶ πόνους τὰ χείλη αὐτῶν λαλεῖ.
\vs{3}Μετὰ σοφίας οἰκοδομεῖται οἶκος, καὶ μετὰ συνέσεως ἀνορθοῦται.
\vs{4}Μετὰ αἰσθήσεως ἐμπίμπλανται ταμιεῖα ἐκ παντὸς πλούτου τιμίου καὶ καλοῦ.
\vs{5}Κρείσσων σοφὸς ἰσχυροῦ, καὶ ἀνὴρ φρόνησιν ἔχων γεωργίου μεγάλου.
\vs{6}Μετὰ κυβερνήσεως γίνεται πόλεμος, βοήθεια δὲ μετὰ καρδίας βουλευτικῆς.

\vs{7}Σοφία καὶ ἔννοια ἀγαθὴ ἐν πύλαις σοφῶν· σοφοὶ οὐκ ἐκκλίνουσιν ἐκ στόματος Κυρίου,
\vs{8}ἀλλὰ λογίζονται ἐν συνεδρίοις· ἀπαιδεύτοις συναντᾷ θάνατος,
\vs{9}ἀποθνήσκει δὲ ἄφρων ἐν ἁμαρτίαις· ἀκαθαρσία δὲ ἀνδρὶ λοιμῷ,
\vs{10}ἐμμολυνθήσεται ἐν ἡμέρᾳ κακῇ, καὶ ἐν ἡμέρᾳ θλίψεως ἕως ἂν ἐκλίπῃ.

\vs{11}Ῥῦσαι ἀγομένους εἰς θάνατον, καὶ ἐκπρίου κτεινομένους, μὴ φείσῃ.
\vs{12}Ἐὰν δὲ εἴπῃς, οὐκ οἶδα τοῦτον, γίνωσκε, ὅτι Κύριος καρδίας πάντων γινώσκει· καὶ ὁ πλάσας πνοὴν πᾶσιν, αὐτὸς οἶδε πάντα, ὃς ἀποδίδωσιν ἑκάστῳ κατὰ τὰ ἔργα αὐτοῦ.
\vs{13}Φάγε μέλι υἱὲ, ἀγαθὸν γὰρ κηρίον, ἵνα γλυκανθῇ σου ὁ φάρυγξ.
\vs{14}Οὕτως αἰσθητήσῃ σοφίαν τῇ σῇ ψυχῇ· ἐὰν γὰρ εὕρῃς, ἔσται καλὴ ἡ τελευτή σου, καὶ ἐλπίς σε οὐκ ἐγκαταλείψει.

\vs{15}Μὴ προσαγάγῃς ἀσεβῆ νομῇ δικαίων, μηδὲ ἀπατηθῇς χορτασίᾳ κοιλίας.
\vs{16}Ἑπτάκις γὰρ πεσεῖται δίκαιος καὶ ἀναστήσεται, οἱ δὲ ἀσεβεῖς ἀσθενήσουσιν ἐν κακοῖς.
\vs{17}Ἐὰν πέσῃ ὁ ἐχθρός σου, μὴ ἐπιχαρῇς αὐτῷ, ἐν δὲ τῷ ὑποσκελίσματι αὐτοῦ μὴ ἐπαίρου.
\vs{18}Ὅτι ὄψεται Κύριος καὶ οὐκ ἀρέσει αὐτῷ, καὶ ἀποστρέψει τὸν θυμὸν αὐτοῦ ἀπʼ αὐτοῦ.
\vs{19}Μὴ χαῖρε ἐπὶ κακοποιοῖς, μηδὲ ζήλου ἁμαρτωλούς.
\vs{20}Οὐ γὰρ μὴ γένηται ἔκγονα πονηρῷ, λαμπτὴρ δὲ ἀσεβῶν σβεσθήσεται.

\vs{21}Φοβοῦ τὸν Θεὸν υἱὲ, καὶ βασιλέα, καὶ μηθʼ ἑτέρῳ αὐτῶν ἀπειθήσῃς.
\vs{22}Ἐξαίφνης γὰρ τίσονται τοὺς ἀσεβεῖς, τὰς δὲ τιμωρίας ἀμφοτέρων τίς γνώσεται;

\vs{22a}Λόγον φυλασσόμενος υἱὸς ἀπωλείας ἐκτὸς ἔσται, [δεχόμενος δὲ ἐδέξατο αὐτόν·
\vs{22b}μηδὲν ψεῦδος ἀπὸ γλώσσης βασιλεῖ λεγέσθω, καὶ οὐδὲν ψεῦδος ἀπὸ γλώσσης αὐτοῦ οὐ μή ἐξέλθῃ·
\vs{22c}μάχαιρα γλῶσσα βασιλέως καὶ οὐ σαρκίνη, ὃς δʼ ἂν παραδοθῇ συντριβήσεται·
\vs{22d}ἐὰν γὰρ ὀξυνθῇ ὁ θυμὸς αὐτοῦ, σὺν νεύροις ἀνθρώπους ἀναλίσκει,
\vs{22e}καὶ ὀστᾶ ἀνθρώπων κατατρώγει, καὶ συγκαίει ὥσπερ φλὸξ, ὥστε ἄβρωτα εἶναι νεοσσοῖς ἀετῶν·
\vs{22f}τοὺς ἐμοὺς λόγους υἱὲ φοβήθητι, καὶ δεξάμενος αὐτοὺς μετανόει.]

Τάδε λέγει ὁ ἀνὴρ τοῖς πιστεύουσι Θεῷ, καὶ παύομαι.

\vs{22g}Ἀφρονέστατος γάρ εἰμι ἁπάντων ἀνθρώπων, καὶ φρόνησις ἀνθρώπων οὐκ ἔστιν ἐν ἐμοί.
\vs{22h}Θεὸς δεδίδαχέ με σοφίαν, καὶ γνῶσιν ἁγίων ἔγνωκα.
\vs{22i}Τίς ἀνέβη εἰς τὸν οὐρανὸν καὶ κατέβη; τίς συνήγαγεν ἀνέμους ἐν κόλπῳ; τίς συνέστρεψεν ὕδωρ ἐν ἱματίῳ; τίς ἐκράτησε πάντων τῶν ἄκρων τῆς γῆς; τί ὄνομα αὐτῷ; ἢ τί ὄνομα τοῖς τέκνοις αὐτοῦ;
\vs{22k}Πάντες γὰρ λόγοι Θεοῦ πεπυρωμένοι, ὑπερασπίζει δὲ αὐτὸς τῶν εὐλαβουμένων αὐτόν.
\vs{22l}Μὴ προσθῇς τοῖς λόγοις αὐτοῦ, ἵνα μὴ ἐλέγξῃ σε, καὶ ψευδὴς γένῃ.

\vs{22m}Δύο αἰτοῦμαι παρὰ σοῦ, μὴ ἀφέλῃς μου χάριν πρὸ τοῦ ἀποθανεῖν με.
\vs{22n}Μάταιον λόγον καὶ ψευδῆ μακράν μου ποίησον, πλοῦτον δὲ καὶ πενίαν μή μοι δῷς, σύνταξον δέ μοι τὰ δέοντα καὶ τὰ αὐτάρκη·
\vs{22o}Ἵνα μὴ πλησθεὶς ψευδὴς γένωμαι, καὶ εἴπω, τίς με ὁρᾷ; ἢ πενηθεὶς κλέψω, καὶ ὀμόσω τὸ ὄνομα τοῦ Θεοῦ.

\vs{22p}Μὴ παραδῷς οἰκέτην εἰς χεῖρας δεσπότου, μήποτε καταράσηταί σε καὶ ἀφανισθῇς.
\vs{22q}Ἔκγονον κακὸν πατέρα καταρᾶται, τὴς δὲ μητέρα οὐκ εὐλογεῖ.
\vs{22r}Ἔκγονον κακὸν δίκαιον ἑαυτὸν κρίνει, τὴν δʼ ἔξοδον αὐτοῦ οὐκ ἀπένιψεν.
\vs{22s}Ἔκγονον κακὸν ὑψηλοὺς ὀφθαλμοὺς ἔχει, τοῖς δὲ βλεφάροις αὐτοῦ ἐπαίρεται.
\vs{22t}Ἔκγονον κακὸν μαχαίρας τοὺς ὀδόντας ἔχει, καὶ τὰς μύλας, τομίδας, ὥστε ἀναλίσκειν καὶ κατεσθίειν τοὺς ταπεινοὺς ἀπὸ τῆς γῆς, καὶ τοὺς πένητας αὐτῶν ἐξ ἀνθρώπων.

\vs{23}Ταῦτα δὲ λέγω ὑμῖν τοῖς σοφοῖς ἐπιγινώσκειν· αἰδεῖσθαι πρόσωπον ἐν κρίσει οὐ καλόν.
\vs{24}Ὁ εἰπὼν τὸν ἀσεβῆ, δίκαιός ἐστιν, ἐπικατάρατος λαοῖς ἔσται καὶ μισητὸς εἰς ἔθνη.
\vs{25}Οἱ δὲ ἐλέγχοντες βελτίους φανοῦνται, ἐπʼ αὐτοὺς δὲ ἥξει εὐλογία·
\vs{26}χείλη δὲ φιλήσουσιν ἀποκρινόμενα λόγους ἀγαθούς.
\vs{27}Ἑτοίμαζε εἰς τὴν ἔξοδον τὰ ἔργα σου, καὶ παρασκευάζου εἰς τὸν ἀγρὸν, καὶ πορεύου κατόπισθέν μου, καὶ ἀνοικοδομήσεις τὸν οἶκόν σου.
\vs{28}Μὴ ἴσθι ψευδὴς μάρτυς ἐπὶ σὸν πολίτην, μηδὲ πλατύνου σοῖς χείλεσι.
\vs{29}Μὴ εἴπῃς, ὃν τρόπον ἐχρήσατό μοι, χρήσομαι αὐτῷ, τίσομαι δὲ αὐτὸν ἅ με ἠδίκησεν.
\vs{30}Ὥσπερ γεώργιον ἀνὴρ ἄφρων, καὶ ὥσπερ ἀμπελὼν ἄνθρωπος ἐνδεὴς φρενῶν.
\vs{31}Ἐὰν ἀφῇς αὐτὸν, χερσωθήσεται καὶ χορτομανήσει ὅλος, καὶ γίνεται ἐκλελειμμένος, οἱ δὲ φραγμοὶ τῶν λίθων αὐτοῦ κατασκάπτονται.
\vs{32}Ὕστερον ἐγὼ μετενόησα, ἐπέβλεψα τοῦ ἐκλέξασθαι παιδείαν.
\vs{33}Ὀλίγον νυστάζω, ὀλίγον δὲ καθυπνῶ, ὀλίγον δὲ ἐναγκαλίζομαι χερσὶ στήθη.
\vs{34}Ἐὰν δὲ τοῦτο ποιῇς, ἥξει προπορευομένη ἡ πενία σου, καὶ ἡ ἔνδειά σου ὥσπερ ἀγαθὸς δρομεύς.

\vs{35}Τῇ βδέλλῃ τρεῖς θυγατέρες ἦσαν ἀγαπήσει ἀγαπώμεναι, καὶ αἱ τρεῖς αὗται οὐκ ἐνεπίμπλασαν αὐτὴν, καὶ ἡ τετάρτη οὐκ ἠρκέσθη εἰπεῖν, ἱκανόν.
\vs{36}Ἄδης καὶ ἔρως γυναικὸς, καὶ γῆ οὐκ ἐμπιπλαμένη ὕδατος, καὶ ὕδωρ καὶ πῦρ οὐ μὴ εἴπωσιν, ἀρκεῖ.

\vs{37}Ὀφθαλμὸν καταγελῶντα πατρὸς, καὶ ἀτιμάζοντα γῆρας μητρὸς, ἐκκόψαισαν αὐτὸν κόρακες ἐκ τῶν φαράγγων, καὶ καταφάγοισαν αὐτὸν νεοσσοὶ ἀετῶν.
\vs{38}Τρία δέ ἐστιν ἀδύνατά μοι νοῆσαι, καὶ τὸ τέταρτον οὐκ ἐπιγινώσκω·
\vs{39}Ἴχνη ἀετοῦ πετομένου, καὶ ὁδοὺς ὄφεως ἐπὶ πέτρας, καὶ τρίβους νηὸς ποντοπορούσης, καὶ ὁδοὺς ἀνδρὸς ἐν νεότητι.
\vs{40}Τοιαύτη ὁδὸς γυναικὸς μοιχαλίδος, ἣ ὅτʼ ἂν πράξῃ ἀπονιψαμένη, οὐδέν φησι πεπραχέναι ἄτοπον.

\vs{41}Διὰ τριῶν σείεται ἡ γῆ, τὸ δὲ τέταρτον οὐ δύναται φέρειν·
\vs{42}Ἐὰν οἰκέτης βασιλεύσῃ, καὶ ἄφρων πλησθῇ σιτίων,
\vs{43}καὶ οἰκέτις ἐὰν ἐκβάλῃ τὴν ἑαυτῆς κυρίαν, καὶ μισητὴ γυνὴ ἐὰν τύχῃ ἀνδρὸς ἀγαθοῦ.

\vs{44}Τέσσαρα δὲ ἐλάχιστα ἐπὶ τῆς γῆς, ταῦτα δέ ἐστι σοφώτερα τῶν σοφῶν·
\vs{45}Οἱ μύρμηκες οἷς μή ἐστιν ἰσχὺς, καὶ ἑτοιμάζονται θέρους τὴν τροφήν·
\vs{46}Καὶ οἱ χοιρογρύλλιοι ἔθνος οὐκ ἰσχυρὸν, οἳ ἐποιήσαντο ἐν πέτραις τοὺς ἑαυτῶν οἴκους·
\vs{47}Ἀβασίλευτόν ἐστιν ἡ ἀκρὶς, καὶ στρατεύει ἀφʼ ἑνὸς κελεύσματος εὐτάκτως·
\vs{48}Καὶ καλαβώτης χερσὶν ἐρειδόμενος, καὶ εὐάλωτος ὢν, κατοικεῖ ἐν ὀχυρώμασι βασιλέων.

\vs{49}Τρία δέ ἐστιν ἃ εὐόδως πορεύεται, καὶ τέταρτον ὃ καλῶς διαβαίνει·
\vs{50}Σκύμνος λέοντος ἰσχυρότερος κτηνῶν, ὃς οὐκ ἀποστρέφεται, οὐδὲ καταπτήσσει κτῆνος·
\vs{51}Καὶ ἀλέκτωρ ἐμπεριπατῶν θηλείαις εὔψυχος, καὶ τράγος ἡγούμενος αἰπολίου, καὶ βασιλεὺς δημηγορῶν ἐν ἔθνει.

\vs{52}Ἐὰν πρόῃ σεαυτὸν ἐν εὐφροσύνῃ, καὶ ἐκτείνῃς τὴν χεῖρά σου μετὰ μάχης, ἀτιμασθήσῃ.
\vs{53}Ἄμελγε γάλα, καὶ ἔσται βούτυρον· ἐὰν δὲ ἐκπιέζῃς μυκτῆρας ἐξελεύσεται αἷμα, ἐὰν δὲ ἐξέλκῃς λόγους, ἐξελεύσονται κρίσεις καὶ μάχαι.

\vs{54}Οἱ ἐμοὶ λόγοι εἴρηνται ὑπὸ Θεοῦ, βασιλέως χρηματισμὸς, ὃν ἐπαίδευσεν ἡ μήτηρ αὐτοῦ.

\vs{55}Τί τέκνον τηρήσεις; τί; ῥήσεις Θεοῦ· πρωτογενὲς σοὶ λέγω υἱέ· τί τέκνον ἐμῆς κοιλίας; τί τέκνον ἐμῶν εὐχῶν;
\vs{56}Μὴ δῷς γυναιξὶ σὸν πλοῦτον, καὶ τὸν σὸν νοῦν καὶ βίον εἰς ὑστεροβουλίαν·
\vs{57}μετὰ βουλῆς πάντα ποίει, μετὰ βουλῆς οἰνοπότει. Οἱ δυνάσται θυμώδεις εἰσὶν, οἶνον δὲ μὴ πινέτωσαν,
\vs{58}ἵνα μὴ πιόντες ἐπιλάθωνται τῆς σοφίας, καὶ ὀρθὰ κρῖναι οὐ μὴ δύνωνται τοὺς ἀσθενεῖς.
\vs{59}Δίδοτε μέθην τοῖς ἐν λύπαις, καὶ οἶνον πίνειν τοῖς ἐν ὀδύναις,
\vs{60}ἵνα ἐπιλάθωνται τῆς πενίας, καὶ τῶν πόνων μὴ μνησθῶσιν ἔτι.
\vs{61}Ἄνοιγε σὸν στόμα λόγῳ Θεοῦ, καὶ κρίνε πάντας ὑγιῶς.
\vs{62}Ἄνοιγε σὸν στόμα καὶ κρίνε δικαίως, διάκρινε δὲ πένητα καὶ ἀσθενῆ.

\ch{25}
Αὗται αἱ παιδεῖαι Σαλωμῶντος αἱ ἀδιάκριτοι, ἃς ἐξεγράψαντο οἱ φίλοι Ἐζεκίου τοῦ βασιλέως τῆς Ἰουδαίας.

\vs{2}Δόξα Θεοῦ κρύπτει λόγον, δόξα δὲ βασιλέως τιμᾷ πράγματα.
\vs{3}Οὐρανὸς ὑψηλὸς, γῆ δὲ βαθεῖα, καρδία δὲ βασιλέως ἀνεξέλεγκτος.
\vs{4}Τύπτε ἀδόκιμον ἀργύριον, καὶ καθαρισθήσεται καθαρὸν ἅπαν.
\vs{5}Κτεῖνε ἀσεβεῖς ἐκ προσώπου βασιλέως, καὶ κατορθώσει ἐν δικαιοσύνῃ ὁ θρόνος αὐτοῦ.

\vs{6}Μὴ ἀλαζονεύου ἐνώπιον βασιλέως, μηδὲ ἐν τόποις δυναστῶν ὑφίστασο·
\vs{7}Κρεῖσσον γάρ σοι τὸ ῥηθῆναι, ἀνάβαινε πρὸς μὲ, ἢ ταπεινῶσαί σε ἐν προσώπῳ δυνάστου· ἃ εἶδον οἱ ὀφθαλμοί σου λέγε.

\vs{8}Μὴ πρόσπιπτε εἰς μάχην ταχέως, ἵνα μὴ μεταμεληθῇς ἐπʼ ἐσχάτων· ἡνίκα ἄν σε ὀνειδίσῃ ὁ σὸς φίλος,
\vs{9}ἀναχώρει εἰς τὰ ὀπίσω· μὴ καταφρόνει,
\vs{10}μή σε ὀνειδίσῃ μὲν ὁ φίλος, ἡ δὲ μάχη σου καὶ ἡ ἔχθρα οὐκ ἀπέσται, ἀλλὰ ἔσται σοι ἴση θανάτῳ·
\vs{10a}χάρις καὶ φιλία ἐλευθεροῖ, ἃς τήρησον σεαυτῷ, ἵνα μὴ ἐπονείδιστος γένῃ, ἀλλὰ φύλαξον τὰς ὁδούς σου εὐσυναλλάκτως.

\vs{11}Μῆλον χρυσοῦν ἐν ὁρμίσκῳ σαρδίου, οὕτως εἰπεῖν λόγον.
\vs{12}Εἰς ἐνώτιον χρυσοῦν καὶ σάρδιον πολυτελὲς δέδεται, λόγος σοφὸς εἰς εὐήκοον οὖς.
\vs{13}Ὥσπερ ἔξοδος χιόνος ἐν ἀμητῷ κατὰ καῦμα ὠφελεῖ, οὕτως ἄγγελος πιστὸς τοὺς ἀποστείλαντας αὐτόν· ψυχὰς γὰρ τῶν αὐτῷ χρωμένων ὠφελεῖ.

\vs{14}Ὥσπερ ἄνεμοι καὶ νέφη καὶ ὑετοὶ, ἐπιφανέστατα, οὕτως ὁ καυχώμενος ἐπὶ δόσει ψευδεῖ.
\vs{15}Ἐν μακροθυμίᾳ εὐοδία βασιλεῦσι, γλῶσσα δὲ μαλακὴ συντρίβει ὀστᾶ.
\vs{16}Μέλι εὑρὼν φάγε τὸ ἱκανὸν, μή ποτε πλησθεὶς ἐξεμέσῃς.
\vs{17}Σπάνιον εἴσαγε σὸν πόδα πρὸς σεαυτοῦ φίλον, μή ποτε πλησθείς σου μισήσῃ σε.
\vs{18}Ῥόπαλον καὶ μάχαιρα καὶ τόξευμα ἀκιδωτὸν, οὕτως καὶ ἀνὴρ ὁ καταμαρτυρῶν τοῦ φίλου αὐτοῦ μαρτυρίαν ψευδῆ.
\vs{19}Οδὸς κακοῦ καὶ ποὺς παρανόμου ὀλεῖται ἐν ἡμέρᾳ κακῇ.

\vs{20}Ὥσπερ ὄξος ἕλκει ἀσύμφορον, οὕτως προσπεσὸν πάθος ἐν σώματι καρδίαν λυπεῖ·
\vs{20a}ὥσπερ σὴς ἐν ἱματίῳ καὶ σκώληξ ξύλῳ, οὕτως λύπη ἀνδρὸς βλάπτει καρδίαν.

\vs{21}Ἐὰν πεινᾷ ὁ ἐχθρός σου, ψώμιζε αὐτὸν, ἐὰν διψᾷ, πότιζε αὐτόν·
\vs{22}Τοῦτο γὰρ ποιῶν ἄνθρακας πυρὸς σωρεύσεις ἐπὶ τὴν κεφαλὴν αὐτοῦ, ὁ δὲ Κύριος ἀνταποδώσει σοι ἀγαθά.
\vs{23}Ἄνεμος Βορέας ἐξεγείρει νέφη, πρόσωπον δὲ ἀναιδὲς γλῶσσαν ἐρεθίζει·
\vs{24}Κρεῖσσον οἰκεῖν ἐπὶ γωνίας δώματος, ἢ μετὰ γυναικὸς λοιδόρου ἐν οἰκίᾳ κοινῇ.
\vs{25}Ὥσπερ ὕδωρ ψυχρὸν ψυχῇ διψώσῃ προσηνὲς, οὕτως ἀγγελία ἀγαθὴ ἐκ γῆς μακρόθεν.
\vs{26}Ὥσπερ εἴ τις πηγὴν φράσσοι καὶ ὕδατος ἔξοδον λυμαίνοιτο, οὕτως ἄκοσμον δίκαιον πεπτωκέναι ἐνώπιον ἀσεβοῦς.
\vs{27}Ἐσθίειν μέλι πολὺ οὐ καλὸν, τιμᾷν δὲ χρὴ λόγους ἐνδόξους.
\vs{28}Ὥσπερ πόλις τὰ τείχη καταβεβλημένη καὶ ἀτείχιστος, οὕτως ἀνὴρ ὃς οὐ μετὰ βουλῆς τι πράσσει.

\ch{26}
Ὥσπερ δρόσος ἐν ἀμητῷ, καὶ ὥσπερ ὑετὸς ἐν θέρει, οὕτως οὐκ ἔστιν ἄφρουι τιμή.
\vs{2}Ὥσπερ ὄρνεα πέταται καὶ στρουθοί, οὕτως ἀρὰ ματαία οὐκ ἐπελεύσεται οὐδενί.
\vs{3}Ὥσπερ μάστιξ ἵππῳ καὶ κέντρον ὄνῳ, οὕτως ῥάβδος ἔθνει παρανόμῳ.
\vs{4}Μὴ ἀποκρίνου ἄφρονι πρὸς τὴν ἐκείνου ἀφροσύνην, ἵνα μὴ ὅμοιος γένῃ αὐτῷ.
\vs{5}Ἀλλὰ ἀποκρίνου ἄφρονι κατὰ τὴν ἀφροσύνην αὐτοῦ, ἵνα μὴ φαίνηται σοφὸς παρʼ ἑαυτῷ.
\vs{6}Ἐκ τῶν ὁδῶν ἑαυτοῦ ὄνειδος ποιεῖται ὁ ἀποστείλας διʼ ἀγγέλου ἄφρονος λόγον.
\vs{7}Ἄφελοῦ πορείαν σκελῶν, καὶ παρανομίαν ἐκ στόματος ἀφρόνων.
\vs{8}Ὃς ἀποδεσμεύει λίθον ἐν σφενδόνῃ, ὅμοιός ἐστι τῷ διδόντι ἄφρονι δόξαν.
\vs{9}Ἄκανθαι φύονται ἐν χειρὶ μεθύσου, δουλεία δὲ ἐν χειρὶ τῶν ἀφρόνων.
\vs{10}Πολλὰ χειμάζεται πᾶσα σὰρξ ἀφρόνων, συντρίβεται γὰρ ἡ ἔκστασις αὐτῶν.
\vs{11}Ὥσπερ κύων ὅταν ἐπέλθῃ ἐπὶ τὸν ἑαυτοῦ ἔμετον καὶ μισητὸς γένηται, οὕτως ἄφρων τῇ ἑαυτοῦ κακίᾳ ἀναστρέψας ἐπὶ τὴν ἑαυτοῦ ἁμαρτίαν·
\vs{11a}ἔστιν αἰσχύνη ἐπάγουσα ἁμαρτίαν, καὶ ἐστιν αἰσχύνη δόξα καὶ χάρις.
\vs{12}Εἶδον ἄνδρα δόξαντα παρʼ αὐτῷ σοφὸν εἶναι, ἐλπίδα μέντοι ἔσχε μᾶλλον ἄφρων αὐτοῦ.
\vs{13}Λέγει ὀκνηρὸς ἀποστελλόμενος εἰς ὁδὸν, λέων ἐν ταῖς ἐν δὲ ταῖς πλατείαις φονευταί.

\vs{14}Ὥσπερ θύρα στρέφεται ἐπὶ τοῦ στρόφιγγος, οὕτως ὀκνηρὸς ἐπὶ τῆς κλίνης αὐτοῦ.
\vs{15}Κρύψας ὀκνηρὸς τὴν χεῖρα ἐν τῷ κόλπῳ αὐτοῦ, οὐ δυνήσεται ἐπενεγκεῖν ἐπὶ στόμα.
\vs{16}Σοφώτερος ἑαυτῷ ὀκνηρὸς φαίνεται, τοῦ ἐν πλησμονῇ ἀποκομίζοντος ἀγγελίαν.

\vs{17}Ὥσπερ ὁ κρατῶν κέρκου κυνὸς, οὕτως ὁ προεστὼς ἀλλοτρίας κρίσεως.
\vs{18}Ὥσπερ οἱ ἰώμενοι προβάλλουσι λόγους εἰς ἀνθρώπους, ὁ δὲ ἀπαντήσας τῷ λόγῳ πρῶτος ὑποσκελισθήσεται·
\vs{19}Οὕτως πάντες οἱ ἐνεδρεύοντες τοὺς ἑαυτῶν φίλους, ὅταν δὲ ὁραθῶσι, λέγουσιν, ὅτι παίζων ἔπραξα.
\vs{20}Ἐν πολλοῖς ξύλοις θάλλει πῦρ, ὅπου δὲ οὐκ ἔστι δίθυμος, ἡσυχάζει μάχη.
\vs{21}Ἐσχάρα ἄνθραξι καὶ ξύλα πυρὶ, ἀνὴρ δὲ λοίδορος εἰς ταραχὴν μάχης.
\vs{22}Λόγοι κερκώπων μαλακοὶ, οὗτοι δὲ τύπτουσιν εἰς ταμιεῖα σπλάγχνων.

\vs{23}Ἀργύριον διδόμενον μετὰ δόλου, ὥσπερ ὄστρακον ἡγητέον· χείλη λεῖα καρδίαν καλύπτει λυπηράν.
\vs{24}Χείλεσι πάντα ἐπινεύει ἀποκλαιόμενος ἐχθρὸς, ἐν δὲ τῇ καρδίᾳ τεκταίνεται δόλους.
\vs{25}Ἐάν σου δέηται ὁ ἐχθρὸς μεγάλῃ τῇ φωνῇ, μὴ πεισθῇς, ἑπτὰ γάρ πονηρίαι ἐν τῇ ψυχῇ αὐτοῦ.
\vs{26}Ὁ κρύπτων ἔχθραν συνίστησι δόλον, ἐκκαλύπτει δὲ τὰς ἑαυτοῦ ἁμαρτίας εὔγνωστος ἐν συνεδρίοις.
\vs{27}Ὁ ὀρύσσων βόθρον τῷ πλησίον, ἐμπεσεῖται εἰς αὐτόν· ὁ δὲ κυλίων λίθον, ἐφʼ ἑαυτὸν κυλίει.
\vs{28}Γλῶσσα ψευδὴς μισεῖ ἀλήθειαν, στόμα δὲ ἄστεγον ποιεῖ ἀκαταστασίας.

\ch{27}
Μὴ καυχῶ τὰ εἰς αὔριον, οὐ γὰρ γινώσκεις τί τέξεται ἡ ἐπιοῦσα.
\vs{2}Ἐγκωμιαζέτω σε ὁ πέλας καὶ μὴ τὸ σὸν στόμα, ἀλλότριος καὶ μὴ τὰ σὰ χείλη.
\vs{3}Βαρὺ λίθος καὶ δυσβάστακτον ἄμμος, ὀργὴ δὲ ἄφρονος βαρυτέρα ἀμφοτέρων.
\vs{4}Ἀνελεήμων θυμὸς καὶ ὀξεῖα ὀργὴ, ἀλλʼ οὐδὲν ὑφίσταται ζῆλος.
\vs{5}Κρείσσους ἔλεγχοι ἀποκεκαλυμμένοι κρυπτομένης φιλίας.
\vs{6}Ἀξιοπιστότερά ἐστι τραύματα φίλου, ἢ ἑκούσια φιλήματα ἐχθροῦ.

\vs{7}Ψυχὴ ἐν πλησμονῇ οὖσα κηρίοις ἐμπαίζει, ψυχῇ δὲ ἐνδεεῖ καὶ τὰ πικρὰ γλυκέα φαίνεται.
\vs{8}Ὥσπερ ὅταν ὄρνεον καταπετασθῇ ἐκ τῆς ἰδίας νοσσιᾶς, οὕτως ἄνθρωπος δουλοῦται ὅταν ἀποξενωθῇ ἐκ τῶν ἰδίων τόπων.
\vs{9}Μύροις καὶ οἴνοις καὶ θυμιάμασι τέρπεται καρδία, καταῤῥήγνυται δὲ ὑπὸ συμπτωμάτων ψυχή.

\vs{10}Φίλον σὸν ἢ φίλον πατρῷον μὴ ἐγκαταλίπῃς, εἰς δὲ τὸν οἶκον τοῦ ἀδελφοῦ σου μὴ εἰσέλθῃς ἀτυχῶν· κρείσσων φίλος ἐγγὺς, ἢ ἀδελφὸς μακρὰν οἰκῶν.
\vs{11}Σοφὸς γίνου υἱὲ, ἵνα σου εὐφραίνηται ἡ καρδία, καὶ ἀπόστρεψον ἀπὸ σοῦ ἐπονειδίστους λόγους.
\vs{12}Πανοῦργος κακῶν ἐπερχομένων ἀπεκρύβη, ἄφρονες δὲ ἐπελθόντες ζημίαν τίσουσιν.
\vs{13}Ἀφελοῦ τὸ ἱμάτιον αὐτοῦ, παρῆλθε γὰρ ὑβριστὴς, ὅστις τὰ ἀλλότρια λυμαίνεται.
\vs{14}Ὃς ἂν εὐλογῇ θίλον τοπρωῒ μεγάλῃ τῇ φωνῇ, καταρωμένου οὐδὲν διαφέρειν δόξει.

\vs{15}Σταγόνες ἐκβάλλουσιν ἄνθρωπον ἐν ἡμέρᾳ χειμερινῇ ἐκ τοῦ οἴκου αὐτοῦ, ὡσαύτως καὶ γυνὴ λοίδορος ἐκ τοῦ ἰδίου οἴκου.
\vs{16}Βορέας σκληρὸς ἄνεμος, ὀνόματι δὲ ἐπιδέξιος καλεῖται.
\vs{17}Σίδηρος σίδηρον ὀξύνει, ἀνὴρ δὲ παροξύνει πρόσωπον ἑταίρου.
\vs{18}Ὃς φυτεύει συκὴν φάγεται τοὺς καρποὺς αὐτῆς, ὃς δὲ φυλάσσει τὸν ἑαυτοῦ κύριον τιμηθήσεται.
\vs{19}Ὥσπερ οὐχ ὅμοια πρόσωπα προσώποις, οὕτως οὐδὲ αἱ διάνοιαι τῶν ἀνθρώπων.
\vs{20}Ἅδης καὶ ἀπώλεια οὐκ ἐμπίμπλανται, ὡσαύτως καὶ οἱ ὀφθαλμοὶ τῶν ἀνθρώπων ἄπληστοι·
\vs{20a}βδέλυγμα Κυρίῳ στηρίζων ὀφθαλμὸν, καὶ οἱ ἀπαίδευτοι ἀκρατεῖς γλώσσῃ.
\vs{21}Δοκίμιον ἀργυρίῳ καὶ χρυσῷ πύρωσις, ἀνὴρ δὲ δοκιμάζεται διὰ στόματος ἐγκωμιαζόντων αὐτόν.
\vs{21a}καρδία ἀνόμου ἐκζητεῖ κακὰ, καρδία δὲ εὐθὴς ζητεῖ γνῶσιν.
\vs{22}Ἐὰν μαστιγοῖς ἄφρονα ἐν μέσῳ συνεδρίου ἀτιμάζων, οὐ μὴ περιέλῃς τὴν ἀφροσύνην αὐτοῦ.

\vs{23}Γνωστῶς ἐπιγνώσῃ ψυχὰς ποιμνίου σου, καὶ ἐπιστήσεις καρδίαν σου σαῖς ἀγέλαις.
\vs{24}Ὅτι οὐκ εἰς τὸν αἰῶνα ἀνδρὶ κράτος καὶ ἰσχὺς, οὐδὲ παραδίδωσιν ἐκ γενεᾶς εἰς γενεάν.
\vs{25}Ἐπιμελοῦ τῶν ἐν τῷ πεδίῳ χλωρῶν, καὶ κερεῖς πόαν, καὶ σύναγε χόρτον ὀρεινὸν,
\vs{26}ἵνα ἔχῃς πρόβατα εἰς ἱματισμόν· τίμα πεδίον, ἵνα ὠσί σοι ἄρνες.
\vs{27}Υἱὲ, παρʼ ἐμοῦ ἔχεις ῥήσεις ἰσχυρὰς εἰς τὴν ζωήν σου, καὶ εἰς τὴν ζωὴν σῶν θεραπόντων.

\ch{28}
Φεύγει ἀσεβὴς μηδενὸς διώκοντος, δίκαιος δὲ ὥσπερ λέων πέποιθε.
\vs{2}Διʼ ἁμαρτίας ἀσεβῶν κρίσεις ἐγείρονται, ἀνὴρ δὲ πανοῦργος κατασβέσει αὐτάς.
\vs{3}Ἀνδρεῖος ἐν ἀσεβείαις συκοφαντεῖ πτωχούς· ὥσπερ ὑετὸς λάβρος καὶ ἀνωφελὴς,
\vs{4}οὕτως οἱ ἐγκαταλείποντες τὸν νόμον ἐγκωμιάζουσιν ἀσέβειαν· οἱ δὲ ἀγαπῶντες τὸν νόμον, περιβάλλουσιν ἑαυτοῖς τεῖχος.
\vs{5}Ἄνδρες κακοὶ οὐ συνήσουσι κρίμα, οἱ δὲ ζητοῦντες τὸν Κύριον συνήσουσιν ἐν παντί.

\vs{6}Κρείσσων πτωχὸς πορευόμενος ἐν ἀληθείᾳ, πλουσίου ψευδοῦς.
\vs{7}Φυλάσσει νόμον υἱὸς συνετὸς, ὃς δὲ ποιμαίνει ἀσωτίαν ἀτιμάζει πατέρα.
\vs{8}Ὁ πληθύνων τὸν πλοῦτον αὐτοῦ μετὰ τόκων καὶ πλεονασμῶν, τῷ ἐλεῶντι πτωχοὺς συνάγει αὐτόν.
\vs{9}Ὁ ἐκκλίνων τὸ οὖς αὐτοῦ μὴ εἰσακοῦσαι νόμου, καὶ αὐτὸς τὴν προσευχὴν αὐτοῦ ἐβδέλυκται.

\vs{10}Ὃς πλανᾷ εὐθεῖς ἐν ὁδῷ κακῇ, εἰς διαφθορὰν αὐτὸς ἐμπεσεῖται· οἱ δὲ ἄνομοι διελεύσονται ἀγαθὰ, καὶ οὐκ εἰσελεύσονται εἰς αὐτά.
\vs{11}Σοφὸς παρʼ ἑαυτῷ ἀνὴρ πλούσιος, πένης δὲ νοήμων καταγνώσεται αὐτοῦ.
\vs{12}Διὰ βοήθειαν δικαίων πολλὴ γίνεται δόξα, ἐν δὲ τόποις ἀσεβῶν ἁλίσκονται ἄνθρωποι.

\vs{13}Ὁ ἐπικαλύπτων ἀσέβειαν ἑαυτοῦ οὐκ εὐοδωθήσεται, ὁ δὲ ἐξηγούμενος ἐλέγχους ἀγαπηθήσεται.
\vs{14}Μακάριος ἀνὴρ ὃς καταπτήσσει πάντα διʼ εὐλάβειαν, ὁ δὲ σκληρὸς τὴν καρδίαν ἐμπεσεῖται κακοῖς.
\vs{15}Λέων πεινῶν καὶ λύκος διψῶν, ὃς τυραννεῖ, πτωχὸς ὢν, ἔθνους πενιχροῦ.
\vs{16}Βασιλεὺς ἐνδεὴς προσόδων μέγας συκοφάντης, ὁ δὲ μισῶν ἀδικίαν μακρὸν χρόνον ζήσεται.

\vs{17}Ἄνδρα τὸν ἐν αἰτίᾳ φόνου ὁ ἐγγυώμενος, φυγὰς ἔσται καὶ οὐκ ἐν ἀσφαλείᾳ·
\vs{17a}παίδευε υἱὸν καὶ ἀγαπήσει σε, καὶ δώσει κόσμον τῇ σῇ ψυχῇ, οὐ μὴ ὑπακούσει ἔθνει παρανόμῳ.
\vs{18}Ὁ πορευόμενος δικαίως βεβοήθηται, ὁ δὲ σκολιαῖς ὁδοῖς πορευόμενος ἐμπλακήσεται.
\vs{19}Ὁ ἐργαζόμενος τὴν ἑαυτοῦ γῆν πλησθήσεται ἄρτων, ὁ δὲ διώκων σχολὴν πλησθήσεται πενίας.

\vs{20}Ἀνὴρ ἀξιόπιστος πολλὰ εὐλογηθήσεται, ὁ δὲ κακὸς οὐκ ἀτιμώρητος ἔσται.
\vs{21}Ὃς οὐκ αἰσχύνεται πρόσωπα δικαίων, οὐκ ἀγαθὸς, ὁ τοιοῦτος ψωμοῦ ἄρτου ἀποδώσεται ἄνδρα.
\vs{22}Σπεύδει πλουτεῖν ἀνὴρ βάσκανος, καὶ οὐκ οἶδεν ὅτι ἐλεήμων κρατήσει αὐτοῦ.

\vs{23}Ὁ ἐλέγχων ἀνθρώπου ὁδοὺς, χάριτας ἕξει μᾶλλον τοῦ γλωσσοχαριτοῦντος.
\vs{24}Ὃς ἀποβάλλεται πατέρα ἢ μητέρα, καὶ δοκεῖ μὴ ἁμαρτάνειν, οὗτος κοινωνός ἐστιν ἀνδρὸς ἀσεβοῦς.
\vs{25}Ἄπιστος ἀνὴρ κρίνει εἰκῆ, ὃς δὲ πέποιθεν ἐπὶ Κύριον ἐν ἐπιμελείᾳ ἔσται.
\vs{26}Ὃς πέποιθε θρασείᾳ καρδίᾳ, ὁ τοιοῦτος ἄφρων, ὃς δὲ πορεύεται σοφίᾳ σωθήσεται.
\vs{27}Ὃς δίδωσι πτωχοῖς οὐκ ἐνδεηθήσεται, ὃς δὲ ἀποστρέφει τὸν ὀφθαλμὸν αὐτοῦ ἐν πολλῇ ἀπορίᾳ ἔσται.
\vs{28}Ἐν τόποις ἀσεβῶν στένουσι δίκαιοι, ἐν δὲ τῇ ἐκείνων ἀπωλείᾳ πληθυνθήσονται δίκαιοι.

\ch{29}
Κρείσσων ἀνὴρ ἐλέγχων ἀνδρὸς σκληροτραχήλου, ἐξαπίνης γὰρ φλεγομένου αὐτοῦ οὐκ ἔστιν ἴασις.
\vs{2}Ἐγκωμιαζομένων δικαίων εὐφρανθήσονται λαοὶ, ἀρχόντων δὲ ἀσεβῶν στένουσιν ἄνδρες.
\vs{3}Ἀνδρὸς φιλοῦντος σοφίαν εὐφραίνεται πατὴρ αὐτοῦ, ὃς δὲ ποιμαίνει πόρνας ἀπολεῖ πλοῦτον.
\vs{4}Βασιλεὺς δίκαιος ἀνίστησι χώραν, ἀνὴρ δὲ παράνομος κατασκάπτει.
\vs{5}Ὃς παρασκευάζεται ἐπὶ πρόσωπον τοῦ ἑαυτοῦ φιλου δίκτυον, περιβάλλει αὐτὸ τοῖς ἑαυτοῦ ποσίν.
\vs{6}Ἁμαρτάνοντι ἀνδρὶ μεγάλη παγὶς, δίκαιος δὲ ἐν χαρᾷ καὶ ἐν εὐφροσύνῃ ἔσται.
\vs{7}Ἐπίσταται δίκαιος κρίνειν πενιχροῖς, ὁ δὲ ἀσεβὴς οὐ νοεῖ γνῶσιν, καὶ πτωχῷ οὐχ ὑπάρχει νοῦς ἐπιγνώμων.

\vs{8}Ἄνδρες ἄνομοι ἐξέκαυσαν πόλιν, σοφοὶ δὲ ἀπέστρεψαν ὀργήν.
\vs{9}Ἀνὴρ σοφὸς κρινεῖ ἔθνη, ἀνὴρ δὲ φαῦλος ὀργιζόμενος καταγελᾶται καὶ οὐ καταπτήσσει.
\vs{10}Ἄνδρες αἱμάτων μέτοχοι μισοῦσιν ὅσιον, οἱ δὲ εὐθεῖς ἐκζητήσουσι ψυχὴν αὐτοῦ.
\vs{11}Ὅλον τὸν θυμὸν αὐτοῦ ἐκφέρει ἄφρων, σοφὸς δὲ ταμιεύεται κατὰ μέρος.
\vs{12}Βασιλέως ὑπακούοντος λόγον ἄδικον, πάντες οἱ ὑπʼ αὐτὸν παράνομοι.
\vs{13}Δανιστοῦ καὶ χρεωφειλέτου ἀλλήλοις συνελθόντων, ἐπισκοπὴν ἀμφοτέρων ποιεῖται ὁ Κύριος.
\vs{14}Βασιλέως ἐν ἀληθείᾳ κρίνοντος πτωχοὺς, ὁ θρόνος αὐτοῦ εἰς μαρτύριον κατασταθήσεται.
\vs{15}Πληγαὶ καὶ ἔλεγχοι διδόασι σοφίαν, παῖς δὲ πλανώμενος αἰσχύνει γονεῖς αὐτοῦ.
\vs{16}Πολλῶν ὄντων ἀσεβῶν πολλαὶγίνονται ἁμαρτίαι, οἱ δὲ δίκαιοι ἐκείνων πιπτόντων κατάφοβοι γίνονται.

\vs{17}Παίδευε υἱόν σου, καὶ ἀναπαύσει σε, καὶ δώσει κόσμον τῇ ψυχῇ σου.
\vs{18}Οὐ μὴ ὑπάρξῃ ἐξηγητὴς ἔθνει παρανόμῳ, ὁ δὲ φυλάσσων τὸν νόμον μακαριστός.
\vs{19}Λόγοις οὐ παιδευθήσεται οἰκέτης σκληρός· ἐὰν γὰρ καὶ νοήσῃ, ἀλλʼ οὐχ ὑπακούσεται.
\vs{20}Ἐὰν ἴδῃς ἄνδρα ταχὺν ἐν λόγοις, γίνωσκε ὅτι ἐλπίδα ἔχει μᾶλλον ὁ ἄφρων αὐτοῦ.
\vs{21}Ὃς κατασπαταλᾷ ἐκ παιδὸς, οἰκέτης ἔσται, ἔσχατον δὲ ὀδυνηθήσεται ἐφʼ ἑαυτῷ·
\vs{22}Ἀνὴρ θυμώδης ἐγείρει νεῖκος, ἀνὴρ δὲ ὀργίλος ἐξώρυξεν ἁμαρτίαν.
\vs{23}Ὕβρις ἄνδρα ταπεινοῖ, τοὺς δὲ ταπεινόφρονας ἐρείδει δόξῃ Κύριος.

\vs{24}Ὃς μερίζεται κλέπτῃ, μισεῖ τὴν ἑαυτοῦ ψυχήν· ἐὰν δὲ ὅρκου προτεθέντος ἀκούσαντες μὴ ἀναγγείλωσι,
\vs{25}φοβηθέντες καὶ αἰσχυνθέντες ἀνθρώπους ὑπεσκελίσθησαν, ὁ δὲ πεποιθὼς ἐπὶ Κυρίῳ εὐφρανθήσεται· ἀσέβεια ἀνδρὶ δίδωσι σφάλμα, ὃς δὲ πέποιθεν ἐπὶ τῷ δεσπότῃ σωθήσεται.
\vs{26}Πολλοὶ θεραπεύουσι πρόσωπα ἡγουμένων, παρὰ δὲ Κυρίου γίνεται τὸ δίκαιον ἀνδρί.
\vs{27}Βδέλυγμα δίκαιος ἀνὴρ ἀνδρὶ ἀδίκῳ, βδέλυγμα δὲ ἀνόμῳ κατευθύνουσα ὁδός.


\ch{31}
\vs{10}Γυναῖκα ἀνδρείαν τίς εὑρήσει; τιμιωτέρα δέ ἐστι λίθων πολυτελῶν ἡ τοιαύτη.
\vs{11}Θαρσεῖ ἐπʼ αὐτῇ ἡ καρδία τοῦ ἀνδρὸς αὐτῆς· ἡ τοιαύτη καλῶν σκύλων οὐκ ἀπορήσει.
\vs{12}Ἐνεργεῖ γὰρ τῷ ἀνδρὶ εἰς ἀγαθὰ πάντα τὸν βίον.
\vs{13}Μηρυομένη ἔρια καὶ λινὸν, ἐποίησεν εὔχρηστον ταῖς χερσὶν αὐτῆς.
\vs{14}Ἐγένετο ὡσεὶ ναῦς ἐμπορευομένη μακρόθεν, συνάγει δὲ αὕτη τὸν βίον.
\vs{15}Καὶ ἀνίσταται ἐκ νυκτῶν, καὶ ἔδωκε βρώματα τῷ οἴκῳ, καὶ ἔργα ταῖς θεραπαίναις.
\vs{16}Θεωρήσασα γεώργιον ἐπρίατο, ἀπὸ δὲ καρπῶν χειρῶν αὐτῆς κατεφύτευσεν κτῆμα.
\vs{17}Ἀναζωσαμένη ἰσχυρῶς τὴν ὀσφῦν αὐτῆς ἤρεισε τοὺς βραχίονας αὐτῆς εἰς ἔργον.
\vs{18}Καὶ ἐγεύσατο ὅτι καλόν ἐστι τὸ ἐργάζεσθαι, καὶ οὐκ ἀποσβέννυται ὁ λύχνος αὐτῆς ὅλην τὴν νύκτα.
\vs{19}τοὺς πήχεις αὐτῆς ἐκτείνει ἐπὶ τὰ συμφέροντα, τὰς δὲ χεῖρας αὐτῆς ἐρείδει εἰς ἄτρακτον.
\vs{20}Χεῖρας δὲ αὐτῆς διήνοιξε πένητι, καρπὸν δὲ ἐξέτεινεν πτωχῷ.

\vs{21}Οὐ φροντίζει τῶν ἐν οἴκῳ ὁ ἀνὴρ αὐτῆς ὅταν που χρονίζῃ, πάντες γὰρ οἱ παρʼ αὐτῆς ἐνδεδυμένοι εἰαί.
\vs{22}Δισσὰς χλαίνας ἐποίησε τῷ ἀνδρὶ αὐτῆς, ἐκ δὲ βύσσου καὶ πορφύρας ἑαυτῇ ἐνδύματα.
\vs{23}Περίβλεπτος δὲ γίνεται ὁ ἀνὴρ αὐτῆς ἐν πύλαις, ἡνίκα ἂν καθίσῃ ἐν συνεδρίῳ μετὰ τῶν γερόντων κατοίκων τῆς γῆς.
\vs{24}Σινδόνας ἐποίησε καὶ ἀπέδοτο περιζώματα τοῖς Χαναναίοις. στόμα αὐτῆς διήνοιξε προσεχόντως καὶ ἐννόμως, καὶ τάξιν ἐστείλατο τῇ γλώσσῃ αὐτῆς.
\vs{25}Ἰσχὺν καὶ εὐπρέπειαν ἐνεδύσατο, καὶ εὐφράνθη ἐν ἡμέραις ἐσχάταις.
\vs{26}Στεγναὶ διατριβαὶ οἴκων αὐτῆς, σῖτα δὲ ὀκνηρὰ οὐκ ἔφαγεν.
\vs{27}Τὸ στόμα δὲ ἀνοίγει σοφῶς καὶ νομοθέσμως. Ἡ δὲ ἐλεημοσύνη αὐτῆς
\vs{28}ἀνέστησε τὰ τέκνα αὐτῆς καὶ ἐπλούτησαν, καὶ ὁ ἀνὴρ αὐτῆς ᾔνεσεν αὐτήν.
\vs{29}Πολλαὶ θυγατέρες ἐκτήσαντο πλοῦτον, πολλαὶ ἐποίησαν δύναμιν· σὺ δὲ ὑπέρκεισαι, ὑπερῇρας πάσας.
\vs{30}Ψευδεῖς ἀρέσκειαι, καὶ μάταιον κάλλος γυναικος· γυνὴ γὰρ συνετὴ εὐλογεῖται, φόβον δὲ Κυρίου αὕτη αἰνείτω.
\vs{31}Δότε αὐτῇ ἀπὸ καρπῶν χειλέων αὐτῆς, καὶ αἰνείσθω ἐν πύλαις ὁ ἀνὴρ αὐτῆς.


\def\book{ΕΚΚΛΗΣΙΑΣΤΗΣ}
\biblebook{ΕΚΚΛΗΣΙΑΣΤΗΣ}


\lettrine[lines=2, loversize=0.2, nindent=0em, findent=.25em]{\textcolor{bookheadingcolor}{Ῥ}}{ΗΜΑΤΑ} Ἐκκλησιαστοῦ υἱοῦ Δαυὶδ βασιλέως Ἰσραὴλ ἐν Ἱερουσαλήμ.
\vs{2}Ματαιότης ματαιοτήτων, εἶπεν ὁ Ἐκκλησιαστὴς, ματαιότης ματαιοτήτων, τὰ πάντα ματαιότης.

\vs{3}Τίς περίσσεια τῷ ἀνθρώπῳ ἐν παντὶ μόχθῳ αὐτοῦ ᾧ μοχθεῖ ὑπὸ τὸν ἥλιον;
\vs{4}Γενεὰ πορεύεται καὶ γενεὰ ἔρχεται, καὶ ἡ γῆ εἰς τὸν αἰῶνα ἕστηκε.
\vs{5}Καὶ ἀνατέλλει ὁ ἥλιος καὶ δύνει ὁ ἥλιος καὶ εἰς τὸν τόπον αὐτοῦ ἕλκει·
\vs{6}αὐτὸς ἀνατέλλων ἐκεῖ πορεύεται πρὸς Νότον, καὶ κυκλοῖ πρὸς Βοῤῥᾶν· κυκλοῖ κυκλῶν πορεύεται τὸ πνεῦμα, καὶ ἐπὶ κύκλους αὐτοῦ ἐπιστρέφει τὸ πνεῦμα.
\vs{7}Πάντες οἱ χείμαῤῥοι πορεύονται εἰς τὴν θάλασσαν, καὶ ἡ θάλασσα οὐκ ἔστιν ἐμπιμπλαμένη· εἰς τὸν τόπον οὗ οἱ χείμαῤῥοι πορεύονται, ἐκεῖ αὐτοὶ ἐπιστρέφουσι τοῦ πορευθῆναι.
\vs{8}Πάντες οἱ λόγοι ἔγκοποι, οὐ δυνήσεται ἀνὴρ τοῦ λαλεῖν· καὶ οὐ πλησθήσεται ὀφθαλμὸς τοῦ ὁρᾷν, καὶ οὐ πληρωθήσεται οὖς ἀπὸ ἀκροάσεως.

\vs{9}Τί τὸ γεγονός; αὐτὸ τὸ γενησόμενον· καὶ τί τὸ πεποιημένον; αὐτὸ τὸ ποιηθησόμενον· καὶ οὐκ ἔστι πᾶν πρόσφατον ὑπὸ τὸν ἥλιον.
\vs{10}Ὃς λαλήσει καὶ ἐρεῖ, ἴδε τοῦτο καινόν ἐστιν; ἤδη γέγονεν ἐν τοῖς αἰῶσι τοῖς γενομένοις ἀπὸ ἔμπροσθεν ἡμῶν.
\vs{11}Οὐκ ἔστι μνήμη τοῖς πρώτοις, καί γε τοῖς ἐσχάτοις γενομένοις οὐκ ἔσται αὐτῶν μνήμη μετὰ τῶν γενησομένων εἰς τὴν ἐσχάτην.

\vs{12}Ἐγὼ ἐκκλησιαστὴς ἐγενόμην βασιλεὺς ἐπὶ Ἰσραὴλ ἐν Ἱερουσαλήμ.
\vs{13}Καὶ ἔδωκα τὴν καρδίαν μου τοῦ ἐκζητῆσαι καὶ τοῦ κατασκέψασθαι ἐν τῇ σοφίᾳ περὶ πάντων τῶν γινομένων ὑπὸ τὸν οὐρανόν, ὅτι περισπασμὸν πονηρὸν ἔδωκεν ὁ Θεὸς τοῖς υἱοῖς τῶν ἀνθρώπων τοῦ περισπᾶσθαι ἐν αὐτῷ.

\vs{14}Εἶδον σύμπαντα τὰ ποιήματα τὰ πεποιημένα ὑπὸ τὸν ἥλιον· καὶ ἰδοὺ τὰ πάντα ματαιότης καὶ προαίρεσις πνεύματος.
\vs{15}Διεστραμμένον οὐ δυνήσεται ἐπικοσμηθῆναι, καὶ ὑστέρημα οὐ δυνήσεται ἀριθμηθῆναι.
\vs{16}Ἐλάλησα ἐγὼ ἐν καρδίᾳ μου, τῷ λέγειν, ἰδοὺ ἐγὼ ἐμεγαλύνθην, καὶ προσέθηκα σοφίαν ἐπὶ πᾶσιν οἳ ἐγένοντο ἔμπροσθέν μου ἐν Ἱερουσαλήμ· καὶ ἔδωκα καρδίαν μου τοῦ γνῶναι σοφίαν καὶ γνῶσιν.
\vs{17}Καὶ καρδία μου εἶδε πολλὰ, σοφίαν καὶ γνῶσιν, παραβολὰς καὶ ἐπιστήμην· ἔγνων ἐγώ ὅτι καί γε τοῦτό ἐστι προαίρεσις πνεύματος·
\vs{18}Ὅτι ἐν πλήθει σοφίας πλῆθος γνώσεως, καὶ ὁ προστιθεὶς γνῶσιν, προσθήσει ἄλγημα.

\ch{2}
Εἶπον ἐγὼ ἐν καρδίᾳ μου, δεῦρο δὴ πειράσω σε ἐν εὐφροσύνῃ, καὶ ἴδε ἐν ἀγαθῷ· καὶ ἰδοὺ καί γε τοῦτο ματαιότης.
\vs{2}Τῷ γέλωτι εἶπα, περιφορὰν, καὶ τῇ εὐφροσύνῃ, τί τοῦτο ποιεῖς;

\vs{3}Καὶ κατεσκεψάμην εἰ ἡ καρδία μου ἑλκύσει ὡς οἶνον τὴν σάρκα μου, καὶ καρδία μου ὡδήγησεν ἐν σοφίᾳ, καὶ τοῦ κρατῆσαι ἐπʼ εὐφροσύνην, ἕως οὗ ἴδω ποῖον τὸ ἀγαθὸν τοῖς υἱοῖς τῶν ἀνθρώπων, ὃ ποιήσουσιν ὑπὸ τὸν ἥλιον, ἀριθμὸν ἡμερῶν ζωῆς αὐτῶν.
\vs{4}Ἐμεγάλυνα ποίημἀ μου, ᾠκοδόμησά μοι οἴκους, ἐφύτευσά μοι ἀμπελῶνας,
\vs{5}ἐποίησά μοι κήπους καὶ παραδείσους, καὶ ἐφύτευσα ἐν αὐτοῖς ξύλον πᾶν καρποῦ.
\vs{6}Ἐποίησά μοι κολυμβήθρας ὑδάτων τοῦ ποτίσαι ἀπʼ αὐτῶν δρυμὸν βλαστῶντα ξύλα.
\vs{7}Ἐκτησάμην δούλους καὶ παιδίσκας, καὶ οἰκογενεῖς ἐγένοντό μοι, καί γε κτῆσις βουκολίου καὶ ποιμνίου πολλὴ ἐγένετό μοι ὑπὲρ πάντας τοὺς γενομένους ἔμπροσθέν μου ἐν Ἱερουσαλήμ.
\vs{8}Συνήγαγόν μοι καί γε ἀργύριον καί γε χρυσίον, καὶ περιουσιασμοὺς βασιλέων καὶ τῶν χωρῶν· ἐποίησά μοι ᾄδοντας καὶ ᾀδούσας, καὶ ἐντρυφήματα υἱῶν ἀνθρώπων, οἰνοχόον καὶ οἰνοχόας.

\vs{9}Καὶ ἐμεγαλύνθην καὶ προσέθηκα παρὰ πάντας τοὺς γενομένους ἀπὸ ἔμπροσθέν μου ἐν Ἱερουσαλήμ, καί γε σοφία μου ἐστάθη μοι.
\vs{10}Καὶ πᾶν ὃ ᾔτησαν οἱ ὀφθαλμοί μου, οὐκ ἀφεῖλον ἀπʼ αὐτῶν· οὐκ ἀπεκώλυσα τὴν καρδίαν μου ἀπὸ πάσης εὐφροσύνης μου, ὅτι καρδία μου εὐφράνθη ἐν παντὶ μόχθῳ μου· καὶ τοῦτο ἐγένετο μερίς μου ἀπὸ παντὸς μόχθου μου.
\vs{11}Καὶ ἐπέβλεψα ἐγὼ ἐν πᾶσι ποιήμασί μου οἷς ἐποίησαν αἱ χεῖρές μου, καὶ ἐν μόχθῳ ᾧ ἐμόχθησα τοῦ ποιεῖν, καὶ ἰδοὺ τὰ πάντα ματαιότης καὶ προαίρεσις πνεύματος, καὶ οὐκ ἔστι περίσσεια ὑπὸ τὸν ἥλιον.

\vs{12}Καὶ ἐπέβλεψα ἐγὼ τοῦ ἰδεῖν σοφίαν καὶ παραφορὰν καὶ ἀφροσύνην, ὅτι τίς ἄνθρωπος ὅς ἐπελεύσεται ὀπίσω τῆς βουλῆς; τὰ ὅσα ἐποίησεν αὐτήν.
\vs{13}Καὶ εἶδον ἐγὼ ὅτι ἐστὶ περίσσεια τῇ σοφίᾳ ὑπὲρ τὴν ἀφροσύνην, ὡς περίσσεια τοῦ φωτὸς ὑπὲρ τὸ σκότος.
\vs{14}Τοῦ σοφοῦ οἱ ὀφθαλμοὶ αὐτοῦ ἐν κεφαλῇ αὐτοῦ, καὶ ὁ ἄφρων ἐν σκότει πορεύεται· καὶ ἔγνων καί γε ἐγὼ, ὅτι συνάντημα ἓν συναντήσεται τοῖς πᾶσιν αὐτοῖς.

\vs{15}Καὶ εἶπα ἐγὼ ἐν καρδίᾳ μου, ὡς συνάντημα τοῦ ἄφρονος καί γε ἐμοὶ συναντήσεταί μοι, καὶ ἱνατί ἐσοφισάμην ἐγώ; περισσὸν ἐλάλησα ἐν καρδίᾳ μου, ὅτι καί γε τοῦτο ματαιότης, διότι ὁ ἄφρων ἐκ περισσεύματος λαλεῖ·
\vs{16}Ὅτι οὐκ ἔστιν ἡ μνήμη τοῦ σοφοῦ μετὰ τοῦ ἄφρονος εἰς τὸν αἰῶνα, καθότι ἤδη αἱ ἡμέραι ἐρχόμεναι τὰ πάντα ἐπελήσθη· καὶ πῶς ἀποθανεῖται ὁ σοφὸς μετὰ τοῦ ἄφρονος;

\vs{17}Καὶ ἐμίσησα σὺν τὴν ζωήν· ὅτι πονηρὸν ἐπʼ ἐμὲ τὸ ποίημα τὸ πεποιημένον ὑπὸ τὸν ἥλιον, ὅτι πάντα ματαιότης καὶ προαίρεσις πνεύματος.
\vs{18}Καὶ ἐμίσησα ἐγὼ σύμπαντα μόχθον μου ὃν ἐγὼ κοπιῶ ὑπὸ τὸν ἥλιον, ὅτι ἀφίω αὐτὸν τῷ ἀνθρώπῳ τῷ γινομένῳ μετʼ ἐμέ.
\vs{19}Καὶ τίς εἶδεν εἰ σοφὸς ἔσται ἢ ἄφρων; καὶ εἰ ἐξουσιάζεται ἐν παντὶ μόχθῳ μου, ᾧ ἐμόχθησα καὶ ᾧ ἐσοφισάμην ὑπὸ τὸν ἥλιον; καί γε τοῦτο ματαιότης.
\vs{20}Καὶ ἐπέστρεψα ἐγὼ τοῦ ἀποτάξασθαι τὴν καρδίαν μου ἐν παντὶ μόχθῳ μου ᾧ ἐμόχθησα ὑπὸ τὸν ἥλιον·
\vs{21}Ὅτι ἐστὶν ἄνθρωπος ὅτι μόχθος αὐτοῦ ἐν σοφίᾳ καὶ ἐν γνώσει καὶ ἐν ἀνδρίᾳ· καὶ ἄνθρωπος ᾧ οὐκ ἐμόχθησεν ἐν αὐτῷ, δώσει αὐτῷ μερίδα αὐτοῦ· καί γε τοῦτο ματαιότης καὶ πονηρία μεγάλη.
\vs{22}ὅτι γίνεται ἐν τῷ ἀνθρώπῳ ἐν παντὶ μόχθῳ αὐτοῦ καὶ ἐν προαιρέσει καρδίας αὐτοῦ ᾧ αὐτὸς μοχθεῖ ὑπὸ τὸν ἥλιον.
\vs{23}Ὅτι πᾶσαι αἱ ἡμέραι αὐτοῦ ἀλγημάτων καὶ θυμοῦ περισπασμὸς αὐτοῦ, καί γε ἐν νυκτὶ οὐ κοιμᾶται ἡ καρδία αὐτοῦ· καί γε τοῦτο ματαιότης ἐστίν.

\vs{24}Οὐκ ἔστιν ἀγαθὸν ἀνθρώπῳ, ὃ φάγεται καὶ ὃ πίεται καὶ ὃ δείξει τῇ ψυχῇ αὐτοῦ ἀγαθὸν ἐν μόχθῳ αὐτοῦ· καί γε τοῦτο εἶδον ἐγὼ ὅτι ἀπὸ χειρὸς τοῦ Θεοῦ ἐστιν·
\vs{25}Ὅτι τίς φάγεται καὶ τίς πίεται πάρεξ αὐτοῦ;
\vs{26}Ὅτι τῷ ἀνθρώπῳ τῷ ἀγαθῷ πρὸ προσώπου αὐτοῦ ἔδωκε σοφίαν καὶ γνῶσιν καὶ εὐφροσύνην, καὶ τῷ ἁμαρτάνοντι ἔδωκε περισπασμὸν τοῦ προσθεῖναι καὶ τοῦ συναγαγεῖν, τοῦ δοῦναι τῷ ἀγαθῷ πρὸ προσώπου τοῦ Θεοῦ, ὅτι καί γε τοῦτο ματαιότης καὶ προαίρεσις πνεύματος.

\ch{3}
Τοῖς πᾶσιν ὁ χρόνος, καὶ καιρὸς τῷ παντὶ πράγματι ὑπὸ τὸν οὐρανόν.
\vs{2}Καιρὸς τοῦ τεκεῖν καὶ καιρὸς τοῦ ἀποθανεῖν, καιρὸς τοῦ φυτεῦσαι καὶ καιρὸς τοῦ ἐκτίλαι τὸ πεφυτευμένον·
\vs{3}Καιρὸς τοῦ ἀποκτεῖναι καὶ καιρὸς τοῦ ἰάσασθαι, καιρὸς τοῦ καθελεῖν καὶ καιρὸς τοῦ οἰκοδομεῖν·
\vs{4}Καιρὸς τοῦ κλαῦσαι καὶ καιρὸς τοῦ γελάσαι, καιρὸς τοῦ κόψασθαι καὶ καιρὸς τοῦ ὀρχήσασθαι·
\vs{5}Καιρὸς τοῦ βαλεῖν λίθους καὶ καιρὸς τοῦ συναγαγεῖν λίθους, καιρὸς τοῦ περιλαβεῖν καὶ καιρὸς τοῦ μακρυνθῆναι ἀπὸ περιλήψεως·
\vs{6}Καιρὸς τοῦ ζητῆσαι καὶ καιρὸς τοῦ ἀπολέσαι, καιρὸς τοῦ φυλάξαι καὶ καιρὸς τοῦ ἐκβαλεῖν·
\vs{7}Καιρὸς τοῦ ῥῆξαι καὶ καιρὸς τοῦ ῥάψαι, καιρὸς τοῦ σιγᾷν καὶ καιρὸς τοῦ λαλεῖν·
\vs{8}Καιρὸς τοῦ φιλῆσαι καὶ καιρὸς τοῦ μισῆσαι, καιρὸς πολέμου καὶ καιρὸς εἰρήνης.

\vs{9}Τίς περίσσεια τοῦ ποιοῦντος ἐν οἷς αὐτὸς μοχθεῖ;

\vs{10}Εἶδον σὺν πάντα τὸν περισπασμὸν, ὃν ἔδωκεν ὁ Θεὸς τοῖς υἱοῖς τῶν ἀνθρώπων τοῦ περισπᾶσθαι ἐν αὐτῷ.
\vs{11}Τὰ σύμπαντα ἃ ἐποίησε καλὰ ἐν καιρῷ αὐτοῦ· καί γε σύμπαντα τὸν αἰῶνα ἔδωκεν ἐν καρδίᾳ αὐτῶν, ὅπως μὴ εὕρῃ ὁ ἄνθρωπος τὸ ποίημα ὁ ἐποίησεν ὁ Θεὸς ἀπʼ ἀρχῆς καὶ μέχρι τέλους.
\vs{12}Ἔγνων ὅτι οὐκ ἔστιν ἀγαθὸν ἐν αὐτοῖς, εἰ μὴ τοῦ εὐφρανθῆναι καὶ τοῦ ποιεῖν ἀγαθὸν ἐν ζωῇ αὐτοῦ·

\vs{13}Καί γε πᾶς ὁ ἄνθρωπος ὃς φάγεται καὶ πίεται, καὶ ἴδῃ ἀγαθὸν ἐν παντὶ μόχθῳ αὐτοῦ, δόμα Θεοῦ ἐστιν.
\vs{14}Ἔγνων ὅτι πάντα ὅσα ἐποίησεν ὁ Θεὸς αὐτὰ ἔσται εἰς τὸν αἰῶνα, ἐπʼ αὐτῷ οὐκ ἔστι προσθεῖναι, καὶ ἀπʼ αὐτοῦ οὐκ ἔστιν ἀφελεῖν· καὶ ὁ Θεὸς ἐποίησεν, ἵνα φοβηθῶσιν ἀπὸ προσώπου αὐτοῦ.
\vs{15}Τὸ γενόμενον ἤδη ἐστί, καὶ ὅσα τοῦ γίνεσθαι ἤδη γέγονε, καὶ ὁ Θεὸς ζητήσει τὸν διωκόμενον.

\vs{16}Καὶ ἔτι εἶδον ὑπὸ τὸν ἥλιον τόπον τῆς κρίσεως, ἐκεῖ ὁ ἀσεβής· καὶ τόπον τοῦ δικαίου, ἐκεῖ ὁ εὐσεβής.
\vs{17}Καὶ εἶπα ἐγὼ ἐν καρδίᾳ μου, σὺν τὸν δίκαιον καὶ σὺν τὸν ἀσεβῆ κρινεῖ ὁ Θεός, ὅτι καιρὸς τῷ παντὶ πράγματι καὶ ἐπὶ παντὶ τῷ ποιήματι ἐκεῖ.

\vs{18}Εἶπα ἐγὼ ἐν καρδίᾳ μου, περὶ λαλιᾶς υἱῶν τοῦ ἀνθρώπου, ὅτι διακρινεῖ αὐτοὺς ὁ Θεὸς, καὶ τοῦ δεῖξαι ὅτι αὐτοὶ κτήνη εἰσί.
\vs{19}Καί γε αὐτοῖς συνάντημα υἱῶν τοῦ ἀνθρώπου, καὶ συνάντημα τοῦ κτήνους, συνάντημα ἓν αὐτοῖς· ὡς ὁ θάνατος τούτου, οὕτως καὶ ὁ θάνατος τούτου· καὶ πνεῦμα ἓν τοῖς πᾶσι· καὶ τί ἐπερίσσευσεν ὁ ἄνθρωπος παρὰ τὸ κτῆνος; οὐδέν· ὅτι πάντα ματαιότης.
\vs{20}Τὰ πάντα εἰς τόπον ἕνα, τὰ πάντα ἐγένετο ἀπὸ τοῦ χοὸς, καὶ τὰ πάντα ἐπιστρέψει εἰς τὸν χοῦν.
\vs{21}Καὶ τίς εἶδε πνεῦμα υἱῶν τοῦ ἀνθρώπου, εἰ ἀναβαίνει αὐτὸ ἄνω; καὶ τὸ πνεῦμα τοῦ κτήνους, εἰ καταβαίνει αὐτὸ κάτω εἰς γῆν;
\vs{22}Καὶ εἶδον ὅτι οὐκ ἔστιν ἀγαθὸν εἰ μὴ ὃ εὐφρανθήσεται ὁ ἀνθρωπος ἐν ποιήμασιν αὐτοῦ, ὅτι αὐτὸ μερὶς αὐτοῦ, ὅτι τίς ἄξει αὐτὸν τοῦ ἰδεῖν ἐν ᾧ ἐὰν γένηται μετʼ αὐτόν;

\ch{4}
Καὶ ἐπέστρεψα ἐγὼ, καὶ εἶδον συμπάσας τὰς συκοφαντίας τὰς γενομένας ὑπὸ τὸν ἥλιον· καὶ ἰδοὺ δάκρυον τῶν συκοφαντουμένων, καὶ οὐκ ἔστιν αὐτοῖς παρακαλῶν, καὶ ἀπὸ χειρὸς συκοφαντούντων αὐτοῖς ἰσχὺς, καὶ οὐκ ἔστιν αὐτοῖς παρακαλῶν.
\vs{2}Καὶ ἐπῄνεσα ἐγὼ σύμπαντας τοὺς τεθνηκότας τοὺς ἤδη ἀποθανόντας ὑπὲρ τοὺς ζῶντας, ὅσοι αὐτοὶ ζῶσιν ἕως τοῦ νῦν.
\vs{3}Καὶ ἀγαθὸς ὑπὲρ τοὺς δύο τούτους ὅστις οὔπω ἐγένετο, ὃς οὐκ εἶδε σὺν πᾶν τὸ ποίημα τὸ πονηρὸν τὸ πεποιημένον ὑπὸ τὸν ἥλιον.

\vs{4}Καὶ εἶδον ἐγὼ σύμπαντα τὸν μόχθον, καὶ σύμπασαν ἀνδρίαν τοῦ ποιήματος, ὅτι αὐτὸ ζῆλος ἀνδρὸς ἀπὸ τοῦ ἑταίρου αὐτοῦ· καί γε τοῦτο ματαιότης καὶ προαίρεσις πνεύματος.
\vs{5}Ὁ ἄφρων περιέβαλε τὰς χεῖρας αὐτοῦ, καὶ ἔφαγε τὰς σάρκας αὐτοῦ.
\vs{6}Ἀγαθὸν πλήρωμα δρακὸς ἀναπαύσεως ὑπὲρ πληρώματα δύο δρακῶν μόχθου καὶ προαιρέσεως πνεύματος.

\vs{7}Καὶ ἐπέστρεψα ἐγὼ, καὶ εἶδον ματαιότητα ὑπὸ τὸν ἥλιον.
\vs{8}Ἔστιν εἷς, καὶ οὐκ ἔστι δεύτερος· καί γε υἱὸς καί γε ἀδελφὸς οὐκ ἔστιν αὐτῷ· καὶ οὐκ ἔστι περασμὸς τῷ παντὶ μόχθῳ αὐτοῦ· καί γε ὀφθαλμὸς αὐτοῦ οὐκ ἐμπίμπλαται πλούτου· καὶ τίνι ἐγὼ μοχθῶ, καὶ στερίσκω τὴν ψυχήν μου ἀπὸ ἀγαθωσύνης; καί γε τοῦτο ματαιότης καὶ περισπασμὸς πονηρός ἐστιν.
\vs{9}Ἀγαθοὶ οἱ δύο ὑπὲρ τὸν ἕνα, οἷς ἐστὶν αὐτοῖς μισθὸς ἀγαθὸς ἐν μόχθῳ αὐτῶν·
\vs{10}Ὅτι ἐὰν πέσωσιν, ὁ εἷς ἐγερεῖ τὸν μέτοχον αὐτοῦ· καὶ οὐαὶ αὐτῷ τῷ ἑνὶ, ὅταν πέσῃ καὶ μὴ ᾖ δεύτερος ἐγεῖραι αὐτόν.
\vs{11}Καί γε ἐὰν κοιμηθῶσι δύο, καὶ θέρμη αὐτοῖς, καὶ ὁ εἷς πῶς θερμανθῇ;
\vs{12}Καὶ ἐὰν ἐπικραταιωθῇ ὁ εἷς, οἱ δύο στήσονται κατέναντι αὐτοῦ, καὶ τὸ σπαρτίον τὸ ἔντριτον οὐ ταχέως ἀποῤῥαγήσεται.

\vs{13}Ἀγαθὸς παῖς πένης καὶ σοφὸς ὑπὲρ βασιλέα πρεσβύτερον καὶ ἄφρονα, ὃς οὐκ ἔγνω τοῦ προσέχειν ἔτι·
\vs{14}Ὅτι ἐξ οἴκου τῶν δεσμίων ἐξελεύσεται τοῦ βασιλεῦσαι, ὅτι καί γε ἐν βασιλείᾳ αὐτοῦ ἐγενήθη πένης.
\vs{15}Εἶδον σύμπαντας τοὺς ζῶντας τοὺς περιπατοῦντας ὑπὸ τὸν ἥλιον μετὰ τοῦ νεανίσκου τοῦ δευτέρου, ὃς στήσεται ἀντʼ αὐτοῦ.
\vs{16}Οὐκ ἔστι περασμὸς τῷ παντὶ λαῷ, τοῖς πᾶσιν οἳ ἐγένοντο ἔμπροσθεν αὐτῶν· καί γε οἱ ἔσχατοι οὐκ εὐφρανθήσονται ἐπʼ αὐτῷ· ὅτι καί γε τοῦτο ματαιότης καὶ προαίρεσις πνεύματος.

\vs{17}Φύλαξον τὸν πόδα σου, ἐν ᾧ ἐὰν πορεύῃ εἰς οἶκον τοῦ Θεοῦ· καὶ ἐγγὺς τοῦ ἀκούειν, ὑπὲρ δόμα τῶν ἀφρόνων θυσία σου, ὅτι οὐκ εἰσὶν εἰδότες τοῦ ποιῆσαι κακόν.

\ch{5}
Μὴ σπεῦδε ἐπὶ στόματί σου, καὶ καρδία σου μὴ ταχυνάτω τοῦ ἐξενέγκαι λόγον πρὸ προσώπου τοῦ Θεοῦ· ὅτι ὁ Θεὸς ἐν τῷ οὐρανῷ ἄνω, καὶ σὺ ἐπὶ τῆς γῆς· διὰ τοῦτο ἔστωσαν οἱ λόγοι σου ὀλίγοι.
\vs{2}Ὅτι παραγίνεται ἐνύπνιον ἐν πλήθει πειρασμοῦ, καὶ φωνὴ ἄφρονος ἐν πλήθει λόγων.

\vs{3}Καθὼς εὔξῃ εὐχὴν τῷ Θεῷ, μὴ χρονίσῃς τοῦ ἀποδοῦναι αὐτήν· ὅτι οὐκ ἔστι θέλημα ἐν ἄφροσι· σὺ οὖν ὅσα ἐὰν εὔξῃ, ἀπόδος.
\vs{4}Ἀγαθὸν τὸ μὴ εὔξασθαί σε, ἢ τὸ εὔξασθαί σε καὶ μὴ ἀποδοῦναι.
\vs{5}Μὴ δῷς τὸ στόμα σου τοῦ ἐξαμαρτῆσαι τὴν σάρκα σου, καὶ μὴ εἴπῃς πρὸ προσώπου τοῦ Θεοῦ, ὅτι ἄγνοιά ἐστιν· ἵνα μὴ ὀργισθῇ ὁ Θεὸς ἐπὶ φωνῇ σου, καὶ διαφθείρῃ τὰ ποιήματα χειρῶν σου.
\vs{6}Ὅτι ἐν πλήθει ἐνυπνίων καὶ ματαιοτήτων καὶ λόγων πολλῶν, ὅτι σὺ τὸν Θεὸν φοβοῦ.

\vs{7}Ἐὰν συκοφαντίαν πένητος καὶ ἁρπαγὴν κρίματος καὶ δικαιοσύνης ἴδῃς ἐν χώρᾳ, μὴ θαυμάσῃς ἐπὶ τῷ πράγματι· ὅτι ὑψηλὸς ἐπάνω ὑψηλοῦ φυλάξαι, καὶ ὑψηλοὶ ἐπʼ αὐτοῖς.
\vs{8}Καὶ περίσσεια γῆς ἐπὶ παντί ἐστι, βασιλεὺς τοῦ ἀγροῦ εἰργασμένου.

\vs{9}Ἀγαπῶν ἀργύριον οὐ πλησθήσεται ἀργυρίου· καὶ τίς ἠγάπησεν ἐν πλήθει αὐτῶν γέννημα; καί γε τοῦτο ματαιότης.
\vs{10}Ἐν πλήθει ἀγαθωσύνης ἐπληθύνθησαν ἔσθοντες αὐτήν· καὶ τί ἀνδρεία τῷ παρʼ αὐτῆς; ὅτι ἀρχὴ τοῦ ὁρᾷν ὀφθαλμοῖς αὐτοῦ.
\vs{11}Γλυκὺς ὕπνος τοῦ δούλου εἰ ὀλίγον καὶ εἰ πολὺ φάγεται, καὶ τῷ ἐμπλησθέντι τοῦ πλουτῆσαι, οὐκ ἔστιν ἀφίων αὐτὸν τοῦ ὑπνῶσαι.

\vs{12}Ἔστιν ἀῤῥωστία ἣν εἶδον ὑπὸ τὸν ἥλιον, πλοῦτον φυλασσόμενον τῷ παρʼ αὐτοῦ εἰς κακίαν αὐτῷ,
\vs{13}καὶ ἀπολεῖται ὁ πλοῦτος ἐκεῖνος ἐν περισπασμῷ πονηρῷ, καὶ ἐγέννησεν υἱὸν, καὶ οὐκ ἔστιν ἐν χειρὶ αὐτοῦ οὐδέν.
\vs{14}Καθὼς ἐξῆλθεν ἀπὸ γαστρὸς μητρὸς αὐτοῦ γυμνὸς, ἐπιστρέψει τοῦ πορευθῆναι ὡς ἥκει, καὶ οὐδὲν οὐ λήψεται ἐν μόχθῳ αὐτοῦ, ἵνα πορευθῇ ἐν χειρὶ αὐτοῦ.
\vs{15}Καί γε τοῦτο πονηρὰ ἀῤῥωστία· ὥσπερ γὰρ παρεγένετο, οὕτως καὶ ἀπελεύσεται· καὶ τίς ἡ περίσσεια αὐτοῦ ᾗ μοχθεῖ εἰς ἄνεμον;
\vs{16}Καί γε πᾶσαι αἱ ἡμέραι αὐτοῦ ἐν σκότει, καὶ ἐν πένθει, καὶ θυμῷ πολλῷ, καὶ ἀῤῥωστίᾳ, καὶ χόλῳ.

\vs{17}Ἰδοὺ, εἶδον ἐγὼ ἀγαθὸν, ὅ ἐστι καλὸν, τοῦ φαγεῖν καὶ τοῦ πιεῖν καὶ τοῦ ἰδεῖν ἀγαθωσύνην ἐν παντὶ μόχθῳ αὐτοῦ, ᾧ ἐὰν μοχθῇ ὑπὸ τὸν ἥλιον ἀριθμὸν ἡμερῶν ζωῆς αὐτοῦ ὧν ἔδωκεν αὐτῷ ὁ Θεὸς, ὅτι αὐτὸ μερὶς αὐτοῦ.
\vs{18}Καί γε πᾶς ἄνθρωπος ᾧ ἔδωκεν αὐτῷ ὁ Θεὸς πλοῦτον καὶ ὑπάρχοντα, καὶ ἐξουσίασεν αὐτῷ φαγεῖν ἀπʼ αὐτοῦ, καὶ λαβεῖν τὸ μέρος αὐτοῦ, καὶ τοῦ· εὐφρανθῆναι ἐν μόχθῳ αὐτοῦ, τοῦτο δόμα Θεοῦ ἐστιν.
\vs{19}Ὅτι οὐ πολλὰ μνησθήσεται τὰς ἡμέρας τῆς ζωῆς αὐτοῦ, ὅτι ὁ Θεὸς περισπᾷ αὐτὸν ἐν εὐφροσύνῃ καρδίας αὐτοῦ.

\ch{6}
Ἔστι πονηρία ἣν εἶδον ὑπὸ τὸν ἥλιον, καὶ πολλή ἐστιν ὑπὸ τὸν ἄνθρωπον·
\vs{2}Ἀνὴρ ᾧ δώσει αὐτῷ ὁ Θεὸς πλοῦτον καὶ ὑπαρχοντα καὶ δόξαν, καὶ οὐκ ἔστιν ὑστερῶν τῇ ψυχῇ αὐτοῦ ἀπὸ πάντων ὧν ἐπιθυμήσει, καὶ οὐκ ἐξουσιάσει αὐτῷ ὁ Θεὸς τοῦ φαγεῖν ἀπʼ αὐτοῦ, ὅτι ἀνὴρ ξένος φάγεται αὐτόν· τοῦτο ματαιότης καὶ ἀῤῥωστία πονηρά ἐστιν.

\vs{3}Ἐὰν γεννήσῃ ἀνὴρ ἑκατόν, καὶ ἔτη πολλὰ ζήσεται, καὶ πλῆθος ὅ, τι ἔσονται αἱ ἡμέραι ἐτῶν αὐτοῦ, καὶ ψυχὴ αὐτοῦ οὐ πλησθήσεται ἀπὸ τῆς ἀγαθωσύνης, καί γε ταφὴ οὐκ ἐγένετο αὐτῷ, εἶπα, ἀγαθὸν ὑπὲρ αὐτὸν τὸ ἔκτρωμα.
\vs{4}Ὅτι ἐν ματαιότητι ἦλθε, καὶ ἐν σκότει πορεύεται, καὶ ἐν σκότει ὄνομα αὐτοῦ καλυφθήσεται·
\vs{5}Καί γε ἥλιον οὐκ εἶδε, καὶ οὐκ ἔγνω ἀναπαύσεις, τούτῳ ὑπὲρ τοῦτον·
\vs{6}Καὶ ἔζησε χιλίων ἐτῶν καθόδους, καὶ ἀγαθωσύνην οὐκ εἶδε, μὴ οὐκ εἰς τόπον ἕνα πορεύεται τὰ πάντα;

\vs{7}Πᾶς μόχθος ἀνθρώπου εἰς στόμα αὐτοῦ, καί γε ἡ ψυχὴ οὐ πληρωθήσεται.
\vs{8}Ὅτι περίσσεια τῷ σοφῷ ὑπὲρ τὸν ἄφρονα, διότι ὁ πένης οἶδε πορευθῆναι κατέναντι τῆς ζωῆς.
\vs{9}Ἀγαθὸν ὅραμα ὀφθαλμῶν ὑπερπορευόμενον ψυχῇ· καί γε τοῦτο ματαιότης καὶ προαίρεσις πνεύματος.

\vs{10}Εἰ τι ἐγένετο, ἤδη κέκληται ὄνομα αὐτοῦ, καὶ ἐγνώσθη ὅ ἐστιν ἄνθρωπος, καὶ οὐ δυνήσεται κριθῆναι μετὰ τοῦ ἰσχυροτὲρου ὑπὲρ αὐτόν.
\vs{11}Ὅτι εἰσι λόγοι πολλοὶ πληθύνοντες ματαιότητα.

\vs{12}Τί περισσὸν τῷ ἀνθρώπῳ; ὅτι τίς οἶδεν ἀγαθὸν τῷ ἀνθρώπῳ ἐν τῇ ζωῇ, ἀριθμὸν ζωῆς ἡμερῶν ματαιότητος αὐτοῦ; καὶ ἐποίησεν αὐτὰ ἐν σκιᾷ· ὅτι τίς ἀπαγγελεῖ τῷ ἀνθρώπῳ, τί ἔσται ὀπίσω αὐτοῦ ὑπὸ τὸν ἥλιον;

\ch{7}
Ἀγαθὸν ὄνομα ὑπὲρ ἔλαιον ἀγαθὸν, καὶ ἡμέρα τοῦ θανάτου ὑπὲρ ἡμέραν γεννήσεως.
\vs{2}Ἀγαθὸν πορευθῆναι εἰς οἶκον πένθους ἢ ὅτι πορευθῆναι εἰς οἶκον πότου· καθότι τοῦτο τέλος παντὸς ἀνθρώπου, καὶ ὁ ζῶν δώσει ἀγαθὸν εἰς καρδίαν αὐτοῦ.
\vs{3}Ἀγαθὸν θυμὸς ὑπὲρ γέλωτα, ὅτι ἐν κακίᾳ προσώπου ἀγαθυνθήσεται καρδία.
\vs{4}Καρδία σοφῶν ἐν οἴκῳ πένθους, καὶ καρδία ἀφρόνων ἐν οἴκῳ εὐφροσύνης.

\vs{5}Ἀγαθὸν τὸ ἀκοῦσαι ἐπιτίμησιν σοφοῦ ὑπὲρ ἄνδρα ἁκούοντα ᾆσμα ἀφρόνων.
\vs{6}Ὡς φωνὴ ἀκανθῶν ὑπὸ τὸν λέβητα, οὕτως γέλως τῶν ἀφρόνων· καί γε τοῦτο ματαιότης.

\vs{7}Ὅτι ἡ συκοφαντία περιφέρει σοφὸν, καὶ ἀπόλλυσι τὴν καρδίαν εὐγενείας αὐτοῦ.
\vs{8}Ἀγαθὴ ἐσχάτη λόγων ὑπὲρ ἀρχὴν αὐτοῦ, ἀγαθὸν μακρόθυμος ὑπὲρ ὑψηλὸν πνεύματι.
\vs{9}Μὴ σπεύσῃς ἐν πνεύματί σου τοῦ θυμοῦσθαι, ὅτι θυμὸς ἐν κόλπῳ ἀφρόνων ἀναπαύσεται.
\vs{10}Μὴ εἴπῃς, τί ἐγένετο, ὅτι αἱ ἡμέραι αἱ πρότεραι ἦσαν ἀγαθαὶ ὑπὲρ ταύτας; ὅτι οὐκ ἐν σοφίᾳ ἐπηρώτησας περὶ τούτου.

\vs{11}Ἀγαθὴ σοφία μετὰ κληρονομίας, καὶ περίσσεια τοῖς θεωροῦσι τὸν ἥλιον.
\vs{12}Ὅτι ἐν σκιᾷ αὐτῆς ἡ σοφία ὡς σκιὰ ἀργυρίου, καὶ περίσσεια γνώσεως τῆς σοφίας ζωοποιήσει τὸν παρʼ αὐτῆς.

\vs{13}Ἴδε τὰ ποιήματα τοῦ Θεοῦ, ὅτι τίς δυνήσεται κοσμῆσαι ὃν ἂν ὁ Θεὸς διαστρέψῃ αὐτόν;
\vs{14}Ἐν ἡμέρᾳ ἀγαθωσύνης ζῆθι ἐν ἀγαθῷ, καὶ ἴδε ἐν ἡμέρᾳ κακίας· ἴδε, καί γε σὺν τούτῳ συμφώνως τοῦτο ἐποίησεν ὁ Θεὸς περὶ λαλιᾶς, ἵνα μὴ εὕρῃ ἄνθρωπος ὀπίσω αὐτοῦ οὐδέν.

\vs{15}Σύμπαντα εἶδον ἐν ἡμέραις ματαιότητός μου· ἐστὶ δίκαιος ἀπολλύμενος ἐν δικαίῳ αὐτοῦ, καί ἐστιν ἀσεβὴς μένων ἐν κακίᾳ αὐτοῦ.
\vs{16}Μὴ γίνου δίκαιος πολὺ, μηδὲ σοφίζου περισσὰ, μή ποτε ἐκπλαγῇς.
\vs{17}Μὴ ἀσεβήσῃς πολὺ, καὶ μὴ γίνου σκληρὸς, ἵνα μὴ ἀποθάνῃς ἐν οὐ καιρῷ σου.
\vs{18}Ἀγαθὸν τὸ ἀντέχεσθαί σε ἐν τούτῳ, καί γε ἀπὸ τούτου μὴ μιάνῃς τὴν χεῖρά σου, ὅτι φοβουμένοις τὸν Θεὸν ἐξελεύσεται τὰ πάντα.

\vs{19}Ἡ σοφία βοηθήσει τῷ σοφῷ ὑπὲρ δέκα ἐξουσιάζοντας τοὺς ὄντας ἐν τῇ πόλει.
\vs{20}Ὅτι ἄνθρωπος οὐκ ἔστι δίκαιος ἐν τῇ γῇ, ὃς ποιήσει ἀγαθὸν καὶ οὐχ ἁμαρτήσεται.
\vs{21}Καί γε εἰς πάντας λόγους οὓς λαλήσουσιν ἀσεβεῖς, μὴ θῇς καρδίαν σου, ὅπως μὴ ἀκούσῃς τοῦ δούλου σου καταρωμένου σε.
\vs{22}Ὅτι πλειστάκις πονηρεύσεταί σε, καὶ καθόδους πολλὰς κακώσει καρδίαν σου, ὅτι ὡς καί γε σὺ κατηράσω ἑτέρους.
\vs{23}Πάντα ταῦτα ἐπείρασα ἐν σοφίᾳ· εἶπα, σοφισθήσομαι· καὶ αὕτη ἐμακρύνθη ἀπʼ ἐμοῦ.
\vs{24}Μακρὰν ὑπὲρ ὃ ἦν, καὶ βαθὺ βάθος, τίς εὑρήσει αὐτό;

\vs{25}Ἐκύκλωσα ἐγὼ καὶ ἡ καρδία μου τοῦ γνῶναι καὶ τοῦ κατασκέψασθαι καὶ τοῦ ζητῆσαι σοφίαν καὶ ψῆφον, καὶ τοῦ γνῶναι ἀσεβοῦς ἀφροσύνην καὶ ὀχληρίαν καὶ περιφοράν.

\vs{26}Καὶ εὑρίσκω ἐγὼ αὐτὴν, καὶ ἐρῶ πικρότερον ὑπὲρ θάνατον· σὺν τὴν γυναῖκα ἥτις ἐστι θήρευμα, καὶ σαγῆναι καρδία αὐτῆς, δεσμὸς εἰς χεῖρας αὐτῆς· ἀγαθὸς πρὸ προσώπου τοῦ Θεοῦ ἐξαιρεθήσεται ἀπʼ αὐτῆς, καὶ ἁμαρτάνων συλληφθήσεται ἐν αὐτῇ.
\vs{27}Ἴδε τοῦτο εὗρον, εἶπεν ὁ Ἐκκλησιαστής· μία τῇ μιᾷ τοῦ εὑρεῖν λογισμὸν,
\vs{28}ὃν ἐπεζήτησεν ἡ ψυχή μου, καὶ οὐχ εὗρον· καὶ ἄνθρωπον ἕνα ἀπὸ χιλίων εὗρον, καὶ γυναῖκα ἐν πᾶσι τούτοις οὐχ εὗρον.
\vs{29}Πλὴν ἴδε τοῦτο εὗρον, ὃ ἐποίησεν ὁ Θεὸς σὺν τὸν ἄνθρωπον εὐθῆ· καὶ αὐτοὶ ἐζήτησαν λογισμοὺς πολλούς.

Τίς οἶδε σοφοὺς, καὶ τίς οἶδε λύσιν ῥήματος;

\ch{8}
Σοφία ἀνθρώπου φωτιεῖ πρόσωπον αὐτοῦ, καὶ ἀναιδὴς προσώπῳ αὐτοῦ μισηθήσεται.

\vs{2}Στόμα βασιλέως φύλαξον, καὶ περὶ λόγου ὅρκου Θεοῦ.
\vs{3}Μὴ σπουδάσῃς, ἀπὸ προσώπου αὐτοῦ πορεύσῃ· μὴ στῇς ἐν λόγῳ πονηρῷ, ὅτι πᾶν ὃ ἐὰν θελήσῃ ποιήσει,
\vs{4}καθὼς βασιλεὺς ἐξουσιάζων. καὶ τίς ἐρεῖ αὐτῷ, τί ποιεῖς;

\vs{5}Ὁ φυλάσσων ἐντολὴν, οὐ γνώσεται ῥῆμα πονηρὸν, καὶ καιρὸν κρίσεως γινώσκει καρδία σοφοῦ.
\vs{6}Ὅτι παντὶ πράγματί ἐστι καιρὸς καὶ κρίσις, ὅτι γνῶσις τοῦ ἀνθρώπου πολλὴ ἐπʼ αὐτόν.
\vs{7}Ὅτι οὐκ ἔστι γινώσκων τί τὸ ἐσόμενον, ὅτι καθὼς ἔσται, τίς ἀναγγελεῖ αὐτῷ;

\vs{8}Οὐκ ἔστιν ἄνθρωπος ἐξουσιάζων ἐν πνεύματι, τοῦ κωλύσαι σὺν τὸ πνεῦμα. καὶ οὐκ ἔστιν ἐξουσία ἐν ἡμέρᾳ θανάτου, καὶ οὐκ ἔστιν ἀποστολὴ ἐν ἡμέρᾳ πολέμου, καὶ οὐ διασώσει ἀσέβεια τὸν παρʼ αὐτῆς.

\vs{9}Καὶ σύμπαν τοῦτο εἶδον, καὶ ἔδωκα τὴν καρδίαν μου εἰς πᾶν τὸ ποίημα ὃ πεποίηται ὑπὸ τὸν ἥλιον, τὰ ὅσα ἐξουσιάσατο ὁ ἄνθρωπος ἐν ἀνθρώπῳ τοῦ κακῶσαι αὐτόν.
\vs{10}Καὶ τότε εἶδον ἀσεβεῖς εἰς τάφους εἰσαχθέντας, καὶ ἐκ τοῦ ἁγίου· καὶ ἐπορεύθησαν καὶ ἐπῃνέθησαν ἐν τῇ πόλει, ὅτι οὕτως ἐποίησαν· καί γε τοῦτο ματαιότης.

\vs{11}Ὅτι οὐκ ἔστι γινομένη ἀντίῤῥησις ἀπὸ τῶν ποιούντων τὸ πονηρὸν ταχὺ, διὰ τοῦτο ἐπληροφορήθη καρδία υἱῶν τοῦ ἀνθρώπου ἐν αὐτοῖς τοῦ ποιῆσαι τὸ πονηρόν.
\vs{12}Ὃς ἥμαρτεν ἐποίησε τὸ πονηρὸν ἀπὸ τότε καὶ ἀπὸ μακρότητος αὐτῶν· ὅτι καὶ γινώσκω ἐγὼ, ὅτι ἐστὶν ἀγαθὸν τοῖς φοβουμένοις τὸν Θεὸν, ὅπως φοβῶνται ἀπὸ προσώπου αὐτοῦ·
\vs{13}Καὶ ἀγαθὸν οὐκ ἔσται τῷ ἀσεβεῖ, καὶ οὐ μακρυνεῖ ἡμέρας ἐν σκιᾷ, ὃς οὐκ ἔστι φοβούμενος ἀπὸ προσώπου τοῦ Θεοῦ.

\vs{14}Ἔστι ματαιότης ἣ πεποίηται ἐπὶ τῆς γῆς, ὅτι εἰσὶ δίκαιοι, ὅτι φθάνει ἐπʼ αὐτοὺς ὡς ποίημα τῶν ἀσεβῶν, καί εἰσιν ἀσεβεῖς, ὅτι φθάνει πρὸς αὐτοὺς ὡς ποίημα τῶν δικαίων· εἶπα, ὅτι καί γε τοῦτο ματαιότης.
\vs{15}Καὶ ἐπῄνεσα ἐγὼ σὺν τὴν εὐφροσύνην, ὅτι οὐκ ἔστιν ἀγαθὸν τῷ ἀνθρώπῳ ὑπὸ τὸν ἥλιον, ὅτι εἰ μὴ φαγεῖν καὶ τοῦ πιεῖν καὶ τοῦ εὐφρανθῆναι· καὶ αὐτὸ συμπροσέσται αὐτῷ ἐν μόχθῳ αὐτοῦ ἡμέρας ζωῆς αὐτοῦ, ὅσας ἔδωκεν αὐτῷ ὁ Θεὸς ὑπὸ τὸν ἥλιον.

\vs{16}Ἐν οἷς ἔδωκα τὴν καρδίαν μου τοῦ γνῶναι τὴν σοφίαν, καὶ τοῦ ἰδεῖν τὸν περισπασμὸν τὸν πεποιημένον ἐπὶ τῆς γῆς, ὅτι καὶ ἐν ἡμέρᾳ καὶ ἐν νυκτὶ ὕπνον ὀφθαλμοῖς αὐτοῦ οὐκ ἔστι βλέπων.
\vs{17}Καὶ εἶδον σύμπαντα τὰ ποιήματα τοῦ Θεοῦ, ὅτι οὐ δυνήσεται ἄνθρωπος τοῦ εὑρεῖν σὺν τὸ ποίημα τὸ πεποιημένον ὑπὸ τὸν ἥλιον· ὅσα ἂν μοχθήσῃ ἄνθρωπος τοῦ ζητῆσαι, καὶ οὐχ εὑρήσει· καί γε ὅσα ἂν εἴπῃ σοφὸς τοῦ γνῶναι, οὐ δυνήσεται τοῦ εὑρεῖν· ὅτι σύμπαν τοῦτο ἔδωκα εἰς καρδίαν μου, καὶ καρδία μου σύμπαν εἶδε τοῦτο.

\ch{9}
Ὡς οἱ δίκαιοι καὶ οἱ σοφοὶ καὶ αἱ ἐργασίαι αὐτῶν ἐν χειρὶ τοῦ Θεοῦ, καί γε ἀγάπην καί γε μῖσος οὐκ ἔστιν εἰδὼς ὁ ἄνθρωπος· τὰ πάντα πρὸ προσώπου αὐτῶν.
\vs{2}Ματαιότης ἐν τοῖς πᾶσι· συνάντημα ἓν τῷ δικαίῳ καὶ τῷ ἀσεβεῖ, τῷ ἀγαθῷ καὶ τῷ κακῷ, καὶ τῷ καθαρῷ καὶ τῷ ἀκαθάρτῳ, καὶ τῷ θυσιάζοντι καὶ τῷ μὴ θυσιάζοντι· ὡς ὁ ἀγαθὸς ὡς ὁ ἁμαρτάνων, ὡς ὁ ὀμνύων καθὼς ὁ τὸν ὅρκον φοβούμενος.

\vs{3}Τοῦτο πονηρὸν ἐν παντὶ πεποιημένῳ ὑπὸ τὸν ἥλιον, ὅτι συνάντημα ἓν τοῖς πᾶσι· καί γε καρδια υἱῶν τοῦ ἀνθρώπου ἐπληρώθη πονηροῦ, καὶ περιφέρεια ἐν καρδίᾳ αὐτῶν ἐν ζωῇ αὐτῶν, καὶ ὀπίσω αὐτῶν πρὸς τοὺς νεκρούς.
\vs{4}Ὅτι τίς ὃς κοινωνεῖ πρὸς πάντας τοὺς ζῶντας; ἔστιν ἐλπὶς, ὅτι ὁ κύων ὁ ζῶν αὐτὸς ἀγαθὸς ὑπὲρ τὸν λέοντα τὸν νεκρόν·
\vs{5}Ὅτι οἱ ζῶντες γνώσονται ὅτι ἀποθανοῦνται, καὶ οἱ νεκροὶ οὐκ εἰσὶ γινώσκοντες οὐδέν· καὶ οὐκ ἔστιν αὐτοῖς ἔτι μισθὸς, ὅτι ἐπελήσθη ἡ μνήμη αὐτῶν.
\vs{6}Καί γε ἀγάπη αὐτῶν, καί γε μῖσος αὐτῶν, καί γε ζῆλος αὐτῶν ἤδη ἀπώλετο· καί γε μερὶς οὐκ ἔστιν αὐτοῖς ἔτι εἰς τὸν αἰῶνα ἐν παντὶ τῷ πεποιημένῳ ὑπὸ τὸν ἥλιον.

\vs{7}Δεῦρο φάγε ἐν εὐφροσύνῃ τὸν ἄρτον σου, καὶ πίε ἐν καρδίᾳ ἀγαθῇ οἶνόν σου, ὅτι ἤδη εὐδόκησεν ὁ Θεὸς τὰ ποιήματά σου.
\vs{8}Ἐν παντὶ καιρῷ ἔστωσαν ἱμάτιά σου λευκὰ, καὶ ἔλαιον ἐπὶ κεφαλῆς σου μὴ ὑστερησάτω.
\vs{9}Καὶ ἴδε ζωὴν μετὰ γυναικὸς ἧς ἠγάπησας πάσας τὰς ἡμέρας ζωῆς ματαιότητός σου, τὰς δοθείσας σοι ὑπὸ τὸν ἥλιον, ὅτι αὐτὸ μερίς σου ἐν τῇ ζωῇ σου, καὶ ἐν τῷ μόχθῳ σου ᾧ σὺ μοχθεῖς ὑπὸ τὸν ἥλιον.

\vs{10}Πάντα ὅσα ἂν εὕρῃ ἡ χείρ σου τοῦ ποιῆσαι, ὡς ἡ δύναμίς σου ποίησον, ὅτι οὐκ ἔστι ποίημα καὶ λογισμὸς καὶ γνῶσις καὶ σοφία ἐν ᾅδῃ, ὅπου σὺ πορεύῃ ἐκεῖ.

\vs{11}Ἐπέστρεψα καὶ εἶδον ὑπὸ τὸν ἥλιον, ὅτι οὐ τοῖς κούφοις ὁ δρόμος, καὶ οὐ τοῖς δυνατοῖς ὁ πόλεμος, καί γε οὐ τῷ σοφῷ ἄρτος, καί γε οὐ τοῖς συνετοῖς πλοῦτος, καί γε οὐ τοῖς γινώσκουσι χάρις, ὅτι καιρὸς καὶ ἀπάντημα συνατήσεται σύμπασιν αὐτοῖς.
\vs{12}Ὅτι καί γε καὶ οὐκ ἔγνω ὁ ἄνθρωπος τὸν καιρὸν αὐτοῦ, ὡς οἱ ἰχθύες οἱ θηρευόμενοι ἐν ἀμφιβλήστρῳ κακῷ, καὶ ὡς ὄρνεα τὰ θηρευόμενα ἐν παγίδι· ὡς αὐτὰ παγιδεύονται οἱ υἱοὶ τοῦ ἀνθρώπου εἰς καιρὸν πονηρὸν, ὅταν ἐπιπέσῃ ἐπʼ αὐτοὺς ἄφνω.

\vs{13}Καί γε τοῦτο εἶδον σοφίαν ὑπὸ τὸν ἥλιον, καὶ μεγάλη ἐστι πρὸς μέ·
\vs{14}Πόλις μικρὰ καὶ ἄνδρες ἐν αὐτῇ ὀλίγοι, καὶ ἔλθῃ ἐπʼ αὐτὴν βασιλεὺς μέγας καὶ κυκλώσῃ αὐτὴν, καὶ οἰκοδομήσῃ ἐπʼ αὐτὴν χάρακας μεγάλους·
\vs{15}καὶ εὕρῃ ἐν αὐτῇ ἄνδρα πένητα σοφὸν, καὶ διασώσῃ αὐτὸς τὴν πόλιν ἐν τῇ σοφίᾳ αὐτοῦ, καὶ ἄνθρωπος οὐκ ἐμνήσθη σὺν τοῦ ἀνδρὸς τοῦ πένητος ἐκείνου.
\vs{16}Καὶ εἶπα ἐγὼ, ἀγαθὴ σοφία ὑπὲρ δύναμιν· καὶ σοφία τοῦ πένητος ἐξουδενωμένη, καὶ οἱ λόγοι αὐτοῦ οὐκ εἰσακουόμενοι.

\vs{17}Λόγοι σοφῶν ἐν ἀναπαύσει ἀκούονται ὑπὲρ κραυγὴν ἐξουσιάζόντων ἐν ἀφροσύναις.

\vs{18}Ἀγαθὴ σοφία ὑπὲρ σκεύη πολέμου· καὶ ἁμαρτάνων εἷς ἀπολέσει ἀγαθωσύνην πολλήν.

\ch{10}
Μυῖαι θανατοῦσαι σαπριοῦσι σκευασίαν ἐλαίου ἡδύσματος· τίμιον ὀλίγον σοφίας ὑπὲρ δόξαν ἀφροσύνης μεγάλην.

\vs{2}Καρδία σοφοῦ εἰς δεξιὸν αὐτοῦ, καὶ καρδία ἄφρονος εἰς ἀριστερὸν αὐτοῦ.
\vs{3}Καί γε ἐν ὁδῷ ὅταν ἄφρων πορεύηται, καρδία αὐτοῦ ὑστερήσει, καὶ ἃ λογιεῖται πάντα ἀφροσύνη ἐστίν.

\vs{4}Ἐὰν πνεῦμα τοῦ ἐξουσιάζοντος ἀναβῇ ἐπὶ σὲ, τόπον σου μὴ ἀφῇς, ὅτι ἴαμα καταπαύσει ἁμαρτίας μεγάλας.
\vs{5}Ἔστι πονηρία ἣν εἶδον ὑπὸ τὸν ἥλιον, ὡς ἀκούσιον ἐξῆλθεν ἀπὸ προσώπου ἐξουσιάζοντος.
\vs{6}Ἐδόθη ὁ ἄφρων ἐν ὕψεσι μεγάλοις, καὶ πλούσιοι ἐν ταπεινῷ καθήσονται.
\vs{7}Εἶδον δούλους ἐφʼ ἵππους, καὶ ἄρχοντας πορευομένους ὡς δούλους ἐπὶ τῆς γῆς.

\vs{8}Ὁ ὀρύσσων βόθρον, εἰς αὐτὸν ἐμπεσεῖται· καὶ καθαιροῦντα φραγμὸν, δήξεται αὐτὸν ὄφις,

\vs{9}Ἐξαίρων λίθους, διαπονηθήσεται ἐν αὐτοῖς· σχίζων ξύλα, κινδυνεύσει ἐν αὐτοῖς.

\vs{10}Ἐὰν ἐκπέσῃ τὸ σιδήριον, καὶ αὐτὸς πρόσωπον ἐτάραξε· καὶ δυνάμεις δυναμώσει, καὶ περίσσεια τῷ ἀνδρὶ οὐ σοφία.

\vs{11}Ἐὰν δάκῃ ὄφις ἐν οὐ ψιθυρισμῷ, καὶ οὐκ ἔστι περίσσεια τῷ ἐπᾴδοντι.
\vs{12}Λόγοι στόματος σοφοῦ χάρις, καὶ χείλη ἄφρονος καταποντιοῦσιν αὐτόν.
\vs{13}Ἀρχὴ λόγων στόματος αὐτοῦ ἀφροσύνη, καὶ ἐσχάτη στόματος αὐτοῦ περιφέρεια πονηρὰ,
\vs{14}καὶ ὁ ἄφρων πληθύνει λόγους· οὐκ ἔγνω ἄνθρωπος τί τὸ γενόμενον, καὶ τί τὸ ἐσόμενον, ὅ, τι ὀπίσω αὐτοῦ τίς ἀναγγελεῖ αὐτῷ;
\vs{15}Μόχθος τῶν ἀφρόνων κακώσει αὐτοὺς, ὃς οὐκ ἔγνω τοῦ πορευθῆναι εἰς πόλιν.

\vs{16}Οὐαί σοι πόλις ἧς ὁ βασιλεύς σου νεώτερος, καὶ οἱ ἄρχοντές σου πρωῒ ἐσθίουσι.
\vs{17}Μακαρία σὺ γῆ, ἧς ὁ βασιλεύς σου υἱὸς ἐλευθέρων, καὶ οἱ ἄρχοντές σου πρὸς καιρὸν φάγονται ἐν δυνάμει, καὶ οὐκ αἰσχυνθήσονται.

\vs{18}Ἐν ὀκνηρίαις ταπεινωθήσεται ἡ δόκωσις, καὶ ἐν ἀργίᾳ χειρῶν στάξει ἡ οἰκία.

\vs{19}Εἰς γέλωτα ποιοῦσιν ἄρτον, καὶ οἶνον καὶ ἔλαιον τοῦ εὐφρανθῆναι ζῶντας, καὶ τοῦ ἀργυρίου ταπεινώσει ἐπακούσεται τὰ πάντα.

\vs{20}Καί γε ἐν συνειδήσει σου βασιλέα μὴ καταράσῃ, καὶ ἐν ταμιείοις κοιτώνων σου μὴ καταράσῃ πλούσιον· ὅτι πετεινὸν τοῦ οὐρανοῦ ἀποίσει τὴν φωνήν σου, καὶ ὁ ἔχων τὰς πτέρυγας ἀπαγγελεῖ λόγον σου.

\ch{11}
Ἀπόστειλον τὸν ἄρτον σου ἐπὶ πρόσωπον τοῦ ὕδατος, ὅτι ἐν πλήθει ἡμερῶν εὑρήσεις αὐτόν.
\vs{2}Δὸς μερίδα τοῖς ἑπτὰ, καί γε τοῖς ὀκτὼ, ὅτι οὐ γινώσκεις τί ἔσται πονηρὸν ἐπὶ τὴν γῆν.
\vs{3}Ἐὰν πλησθῶσι τὰ νέφη ὑετοῦ, ἐπὶ τὴν γῆν ἐκχέουσι· καὶ ἐὰν πέσῃ ξύλον ἐν τῷ Νότῳ, καὶ ἐὰν ἐν τῷ Βοῤῥᾷ, τόπῳ οὗ πεσεῖται τὸ ξυλον, ἐκεῖ ἔσται.
\vs{4}Τηρῶν ἄνεμον οὐ σπείρει, καὶ βλέπων ἐν ταῖς νεφέλαις οὐ θερίσει.
\vs{5}Ἐν οἷς οὐκ ἔστι γινώσκων τίς ἡ ὁδὸς τοῦ πνεύματος, ὡς ὀστᾶ ἐν γαστρὶ κυοφορούσης, οὕτως οὐ γνώσῃ τὰ ποιήματα τοῦ Θεοῦ, ὅσα ποιήσει τὰ σύμπαντα.
\vs{6}Ἐν τῷ πρωῒ σπεῖρον τὸ σπέρμα σου, καὶ ἐν ἑσπέρᾳ μὴ ἀφέτω ἡ χείρ σου, ὅτι οὐ γινώσκεις ποῖον στοιχήσει, ἢ τοῦτο ἢ τοῦτο, καὶ ἐὰν τὰ δύο ἐπιτοαυτὸ ἀγαθά.

\vs{7}Καὶ γλυκὺ τὸ φῶς, καὶ ἀγαθὸν τοῖς ὀφθαλμοῖς τοῦ βλέπειν σὺν τὸν ἥλιον.
\vs{8}Ὅτ καὶ ἐὰν ἔτη πολλὰ ζήσεται ὁ ἄνθρωπος, ἐν πᾶσιν αὐτοῖς εὐφρανθήσεται καὶ μνησθήσεται τὰς ἡμέρας τοῦ σκότους, ὅτι πολλαὶ ἔσονται· πᾶν τὸ ἐρχόμενον ματαιότης.

\vs{9}Εὐφραίνου νεανίσκε ἐν νεότητί σου, καὶ ἀγαθυνάτω σε ἡ καρδία σου ἐν ἡμέραις νεότητός σου, καὶ περιπάτει ἐν ὁδοῖς καρδίας σου ἄμωμος, καὶ μὴ ἐν ὁράσει ὀφθαλμῶν σου· καὶ γνῶθι ὅτι ἐπὶ πᾶσι τούτοις ἄξει σε ὁ Θεὸς ἐν κρίσει.
\vs{10}Καὶ ἀπόστησον θυμὸν ἀπὸ καρδίας σου, καὶ πάραγε πονηρίαν ἀπὸ σαρκός σου, ὅτι ἡ νεότης καὶ ἡ ἄνοια ματαιότης.

\ch{12}
Καὶ μνήσθητι τοῦ κτίσαντός σε ἐν ἡμέραις νεότητός σου, ἕως ὅτου μὴ ἔλθωσιν αἱ ἡμέραι τῆς κακίας, καὶ φθάσουσιν ἔτη ἐν οἷς ἐρεῖς, οὐκ ἔστι μοι ἐν αὐτοῖς θέλημα.
\vs{2}Ἕως οὗ μὴ σκοτισθῇ ὁ ἥλιος καὶ τὸ φῶς, καὶ ἡ σελήνη καὶ οἱ ἀστέρες, καὶ ἐπιστρέψουσι τὰ νέφη ὀπίσω τοῦ ὑετοῦ.
\vs{3}Ἐν ἡμέρᾳ ᾗ ἐὰν σαλευθῶσι φύλακες τῆς οἰκίας, καὶ διαστραφῶσιν ἄνδρες τῆς δυνάμεως, καὶ ἤργησαν αἱ ἀλήθουσαι ὅτι ὠλιγώθησαν, καὶ σκοτάσουσιν αἱ βλέπουσαι ἐν ταῖς ὀπαῖς·
\vs{4}Καὶ κλείσουσι θύρας ἐν ἀγορᾷ, ἐν ἀσθενείᾳ φωνῆς τῆς ἀληθούσης· καὶ ἀναστήσεται εἰς φωνὴν τοῦ στρουθίου, καὶ ταπεινωθήσονται πᾶσαι αἱ θυγατέρες τοῦ ᾄσματος·
\vs{5}Καὶ εἰς τὸ ὕψος ὄψονται, καὶ θάμβοι ἐν τῇ ὁδῷ, καὶ ἀνθήσῃ τὸ ἀμύγδαλον, καὶ παχυνθῇ ἡ ἀκρὶς, καὶ διασκεδασθῇ ἡ κάππαρις, ὅτι ἐπορεύθη ὁ ἄνθρωπος εἰς οἶκον αἰῶνος αὐτοῦ, καὶ ἐκύκλωσαν ἐν ἀγορᾷ οἱ κοπτόμενοι.
\vs{6}Ἕως ὅτου μὴ ἀνατραπῇ τὸ σχοινίον τοῦ ἀργυρίου, καὶ συντριβῇ τὸ ἀνθέμιον τοῦ χρυσίου, καὶ συντριβῇ ὑδρία ἐπὶ τῇ πηγῇ, καὶ συντροχάσῃ ὁ τροχὸς ἐπὶ τὸν λάκκον·
\vs{7}Καὶ ἐπιστρέψῃ ὁ χοῦς ἐπὶ τὴν γῆν ὡς ἦν, καὶ τὸ πνεῦμα ἐπιστρέψῃ πρὸς τὸν Θεὸν ὃς ἔδωκεν αὐτό.

\vs{8}Ματαιότης ματαιοτήτων, εἶπεν ὁ Ἐκκλησιαστὴς, τὰ πάντα ματαιότης.
\vs{9}Καὶ περισσὸν ὅτι ἐγένετο Ἐκκλησιαστὴς σοφὸς, ὅτι ἐδίδαξε γνῶσιν σὺν τὸν ἄνθρωπον, καὶ οὖς ἐξιχνιάσεται κόσμιον παραβολῶν.
\vs{10}Πολλὰ ἐζήτησεν Ἐκκλησιαστὴς τοῦ εὑρεῖν λόγους θελήματος, καὶ γεγραμμένον εὐθύτητος, λόγους ἀληθείας.
\vs{11}Λόγοι σοφῶν ὡς τὰ βούκεντρα, καὶ ὡς ἧλοι πεφυτευμένοι, οἳ παρὰ τῶν συνθεμάτων ἐδόθησαν ἐκ ποιμένος ἑνός.
\vs{12}Καὶ περισσὸν ἐξ αὐτῶν υἱέ μου φύλαξαι· τοῦ ποιῆσαι βιβλία πολλὰ οὐκ ἔστι περασμὸς, καὶ μελέτη πολλὴ κόπωσις σαρκός.

\vs{13}Τέλος λόγου, τὸ πᾶν ἄκουε· τὸν Θεὸν φοβοῦ, καὶ τὰς ἐντολὰς αὐτοῦ φύλασσε· ὅτι τοῦτο πᾶς ὁ ἄνθρωπος.
\vs{14}Ὅτι σύμπαν τὸ ποίημα ὁ Θεὸς ἄξει ἐν κρίσει, ἐν παντὶ παρεωραμένῳ, ἐὰν ἀγαθὸν καὶ ἐὰν πονηρόν.


\def\book{ΑΣΜΑ}
\biblebook{ΑΣΜΑ}


\lettrine[lines=2, loversize=0.2, nindent=0em, findent=.25em]{\textcolor{bookheadingcolor}{Ἄ}}{ΣΜΑ} ᾀσμάτων, ὅ ἐστι Σαλωμών.
\vs{2}Φιλησάτω με ἀπὸ φιλημάτων στόματος αὐτοῦ· ὅτι ἀγαθοὶ μαστοί σου ὑπὲρ οἶνον, καὶ ὀσμὴ μύρων σου ὑπὲρ πάντα τὰ ἀρώματα·
\vs{3}μῦρον ἐκκενωθὲν ὄνομά σου· διὰ τοῦτο νεάνιδες ἠγάπησάν σε,
\vs{4}εἵλκυσάν σε· ὀπίσω σου εἰς ὀσμὴν μύρων σου δραμοῦμεν· εἰσήνεγκέ με ὁ βασιλεὺς εἰς τὸ ταμεῖον αὐτοῦ· ἀγαλλιασώμεθα καὶ εὐφρανθῶμεν ἐν σοί· ἀγαπήσομεν μαστούς σου ὑπὲρ οἶνον· εὐθύτης ἠγάπησέ σε.

\vs{5}Μέλαινά εἰμι ἐγὼ καὶ καλὴ, θυγατέρες Ἱερουσαλὴμ, ὡς σκηνώματα Κηδὰρ, ὡς δέῤῥεις Σαλωμών.
\vs{6}Μὴ βλέψητέ με ὅτι ἐγώ εἰμι μεμελανωμένη, ὅτι παρέβλεψέ με ὁ ἥλιος· υἱοὶ μητρός μου ἐμαχέσαντο ἐν ἐμοὶ, ἔθεντό με φυλάκισσαν ἐν ἀμπελῶσιν, ἀμπελῶνα ἐμὸν οὐκ ἐφύλαξα.

\vs{7}Ἀπάγγειλόν μοι ὃν ἠγάπησεν ἡ ψυχή μου, ποῦ ποιμαίνεις, ποῦ κοιτάζεις ἐν μεσημβρίᾳ, μήποτε γένωμαι ὡς περιβαλλομένη ἐπʼ ἀγέλαις ἑταίρων σου.

\vs{8}Ἐὰν μὴ γνῷς σεαυτὴν ἡ καλὴ ἐν γυναιξὶν, ἔξελθε σὺ ἐν πτέρναις τῶν ποιμνίων, καὶ ποίμαινε τὰς ἐρίφους σου ἐπὶ σκηνώμασι τῶν ποιμένων.
\vs{9}Τῇ ἵππῳ μου ἐν ἅρμασε Φαραὼ ὡμοίωσά σε ἡ πλησίον μου.
\vs{10}Τί ὡραιώθησαν σιαγόνες σου ὡς τρυγόνος, τράχηλός σου ὡς ὁρμίσκοι;
\vs{11}Ὁμοιώματα χρυσίου ποιήσομέν σοι μετὰ στιγμάτων τοῦ ἀργυρίου.

\vs{12}Ἕως οὗ ὁ βασιλεὺς ἐν ἀνακλίσει αὐτοῦ· νάρδος μου ἔδωκεν ὀσμὴν αὐτοῦ.
\vs{13}Ἀπόδεσμος τῆς στακτῆς ἀδελφιδός μου ἐμοὶ, ἀυαμέσου τῶν μαστῶν μου αὐλισθήσεται.
\vs{14}Βότρυς τῆς κύπρου ἀδελφιδός μου ἐμοὶ, ἐν ἀμπελῶσιν Ἐνγαδδί.

\vs{15}Ἰδοὺ εἶ καλὴ ἡ πλησίον μου, ἰδοὺ εἶ καλὴ· ὀφθαλμοί σου περιστεραί.
\vs{16}Ἰδοὺ εἶ καλὸς ἀδελφιδός μου, καί γε ὡραῖος πρὸς κλίνῃ ἡμῶν σύσκιος·
\vs{17}Δοκοὶ οἴκων ἡμῶν κέδροι, φατνώματα ἡμῶν κυπάρισσοι.

\ch{2}
Ἐγὼ ἄνθος τοῦ πεδίου, κρίνον τῶν κοιλάδων.

\vs{2}Ὡς κρίνον ἐν μέσῳ ἀκανθῶν, οὕτως ἡ πλησίον μου ἀναμέσον τῶν θυγατέρων.

\vs{3}Ὡς μῆλον ἐν τοῖς ξύλοις τοῦ δρυμοῦ, οὕτως ἀδελφιδός μου ἀναμέσον τῶν υἱῶν· ἐν τῇ σκιᾷ αὐτοῦ ἐπεθύμησα, καὶ ἐκάθισα, καὶ καρπὸς αὐτοῦ γλυκὺς ἐν λάρυγγί μου.
\vs{4}Εἰσαγάγετέ με εἰς οἶκον τοῦ οἴνου, τάξατε ἐπʼ ἐμὲ ἀγάπην.
\vs{5}Στηρίσατέ με ἐν μύροις, στοιβάσατέ με ἐν μήλοις, ὅτι τετρωμένη ἀγάπης ἐγώ.
\vs{6}Εὐώνυμος αὐτοῦ ὑπὸ τὴν κεφαλήν μου, καὶ ἡ δεξιὰ αὐτοῦ περιλήψεταί με.

\vs{7}Ὥρκισα ὑμᾶς θυγατέρες Ἱερουσαλὴμ ἐν δυνάμεσι καὶ ἐν ἰσχύσεσι τοῦ ἀγροῦ· ἐὰν ἐγείρητε καὶ ἐξεγείρητε τὴν ἀγάπην ἕως οὗ θελήσῃ.

\vs{8}Φωνὴ ἀδελφιδοῦ μου, ἰδοὺ οὗτος ἥκει πηδῶν ἐπὶ τὰ ὄρη, διαλλόμενος ἐπὶ τοὺς βουνούς.

\vs{9}Ὅμοιός ἐστιν ἀδελφιδός μου τῇ δορκάδι ἢ νεβρῷ ἐλάφων ἐπὶ τὰ ὄρη Βαιθήλ· ἰδοὺ οὗτος ὀπίσω τοῦ τοίχου ἡμῶν, παρακύπτων διὰ τῶν θυρίδων, ἐκκύπτων διὰ τῶν δικτύων.
\vs{10}Ἀποκρίνεται ἀδελφιδός μου, καὶ λέγει μοι, ἀνάστα, ἐλθὲ ἡ πλησίον μου, καλή μου, περιστερά μου.
\vs{11}Ὅτι ἰδοὺ ὁ χειμὼν παρῆλθεν, ὁ ὑετὸς ἀπῆλθεν, ἐπορεύθη ἑαυτῷ.
\vs{12}Τὰ ἄνθη ὤφθη ἐν τῇ γῇ, καιρὸς τῆς τομῆς ἔφθακε, φωνὴ τῆς τρυγόνος ἠκούσθη ἐν τῇ γῇ ἡμῶν.
\vs{13}Ἡ συκὴ ἐξήνεγκεν ὀλύνθους αὐτῆς, αἱ ἄμπελοι κυπρίζουσιν, ἔδωκαν ὀσμήν· ἀνάστα, ἐλθὲ ἡ πλησίον μου, καλή μου, περιστερά μου, καὶ ἐλθὲ.

\vs{14}Σὺ περιστερά μου, ἐν σκέπῃ τῆς πέτρας, ἐχόμενα τοῦ προτειχίσματος· δεῖξόν μοι τὴν ὄψιν σου, καὶ ἀκούτισόν με τὴν φωνήν σου, ὅτι ἡ φωνή σου ἡδεῖα, καὶ ἡ ὄψις σου ὡραῖα.

\vs{15}Πιάσατε ἡμῖν ἀλώπεκας μικροὺς ἀφανίζοντας ἀμπελῶνας· καὶ αἱ ἄμπελοι ἡμῶν κυπρίζουσαι.

\vs{16}Ἀδελφιδός μου ἐμοὶ, κᾀγὼ αὐτῷ· ὁ ποιμαίνων ἐν τοῖς κρίνοις.

\vs{17}Ἕως οὗ διαπνεύσῃ ἡ ἡμέρα, καὶ κινηθῶσιν αἱ σκιαί· ἀπόστρεψον, ὁμοιώθητι σὺ ἀδελφιδε μου τῷ δόρκωνι ἢ νεβρῷ ἐλάφων ἐπὶ ὄρη κοιλωμάτων.

\ch{3}
Ἐπὶ κοίτην μου ἐν νυξὶν, ἐζήτησα ὃν ἠγάπησεν ἡ ψυχή μου· ἐζήτησα αὐτὸν, καὶ οὐχ εὗρον αὐτόν· ἐκάλεσα αὐτὸν, καὶ οὐχ ὑπήκουσέ μου.
\vs{2}Ἀναστήσομαι δὴ καὶ κυκλώσω ἐν τῇ πόλει, ἐν ταῖς ἀγοραῖς, καὶ ἐν ταῖς πλατείαις, καὶ ζητήσω ὃν ἠγάπησεν ἡ ψυχή μου· ἐζήτησα αὐτὸν, καὶ οὐχ εὗρον αὐτόν.
\vs{3}Εὕροσάν με οἱ τηροῦντες, οἱ κυκλοῦντες ἐν τῇ πόλει. Μὴ ὃν ἠγάπησεν ἡ ψυχή μου, ἴδετε;
\vs{4}Ὡς μικρὸν ὅτε παρῆλθον ἀπʼ αὐτῶν, ἕως οὗ εὗρον ὃν ἠγάπησεν ἡ ψυχή μου· ἐκράτησα αὐτὸν καὶ οὐκ ἀφῆκα αὐτὸν, ἕως οὗ εἰσήγαγον αὐτὸν εἰς οἶκον μητρός μου, καὶ εἰς ταμεῖον τῆς συλλαβούσης με.

\vs{5}Ὥρκισα ὑμᾶς θυγατέρες Ἱερουσαλὴμ ἐν ταῖς δυνάμεσι, καὶ ἐν ταῖς ἰσχύσεσι τοῦ ἀγροῦ· ἐὰν ἐγείρητε καὶ ἐξεγείρητε τὴν ἀγάπην ἕως ἂν θελήσῃ.

\vs{6}Τίς αὕτη ἡ ἀναβαίνουσα ἀπὸ τῆς ἐρήμου, ὡς στελέχη καπνοῦ τεθυμιαμένη σμύρναν καὶ λίβανον ἀπὸ πάντων κονιορτῶν μυρεψοῦ;
\vs{7}Ἰδοὺ ἡ κλίνη τοῦ Σαλωμὼν, ἑξήκοντα δυνατοὶ κύκλῳ αὐτῆς ἀπὸ δυνατῶν Ἰσραήλ·
\vs{8}Πάντες κατέχοντες ῥομφαίαν δεδιδαγμένοι πόλεμον· ἀνὴρ ῥομφαία αὐτοῦ ἐπὶ μηρὸν αὐτοῦ ἀπὸ θάμβους ἐν νυξί.

\vs{9}Φορεῖον ἐποίησεν ἑαυτῷ ὁ βασιλεὺς Σαλωμὼν ἀπὸ ξύλων τοῦ Λιβάνου·
\vs{10}Στύλους αὐτοῦ ἐποίησεν ἀργύριον, καὶ ἀνάκλιτον αὐτοῦ χρύσεον· ἐπίβασις αὐτοῦ πορφυρᾶ, ἐντὸς αὐτοῦ λιθόστρωτον, ἀγάπην ἀπὸ θυγατέρων Ἱερουσαλήμ.
\vs{11}Θυγατέρες Σιὼν ἐξέλθατε, καὶ ἴδετε ἐν τῷ βασιλεῖ Σαλωμὼν, ἐν τῷ στεφάνῳ ᾧ ἐστεφάνωσεν αὐτὸν ἡ μήτηρ αὐτοῦ, ἐν ἡμέρᾳ νυμφεύσεως αὐτοῦ, καὶ ἐν ἡμέρᾳ εὐφροσύνης καρδίας αὐτοῦ.

\ch{4}
Ἰδοὺ εἶ καλὴ ἡ πλησίον μου, ἰδοὺ εἶ καλή· ὀφθαλμοί σου περιστεραὶ, ἐκτὸς τῆς σιωπήσεώς σου· τρίχωμά σου ὡς ἀγέλαι τῶν αἰγῶν, αἳ ἀπεκαλύφθησαν ἀπὸ τοῦ Γαλαάδ.
\vs{2}Ὀδόντες σου ὡς ἀγέλαι τῶν κεκαρμένων, αἳ ἀνέβησαν ἀπὸ τοῦ λουτροῦ, αἱ πᾶσαι διδυμεύουσαι, καὶ ἀτεκνοῦσα οὐκ ἔστιν ἐν αὐταῖς.
\vs{3}Ὡς σπάρτίον τὸ κόκκικον χείλη σου, καὶ ἡ λαλιά σου ὡραῖα, ὡς λέπυρον ῥοᾶς μῆλόν σου ἐκτὸς τῆς σιωπήσεώς σου.
\vs{4}Ὡς πύργος Δαυὶδ τράχηλός σου, ὁ ᾠκοδομημένος εἰς θαλπιώθ· χίλιοι θυρεοὶ κρέμανται ἐπʼ αὐτὸν, πᾶσαι βολίδες τῶν δυνατῶν.
\vs{5}Δύο μαστοί σου ὡς δύο νεβροὶ δίδυμοι δορκάδος οἱ νεμόμενοι ἐν κρίνοις,
\vs{6}ἕως οὗ διαπνεύσῃ ἡμέρα καὶ κινηθῶσιν αἱ σκιαί· πορεύσομαι ἐμαυτῷ πρὸς τὸ ὄρος τῆς σμύρνης καὶ πρὸς τὸν βουνὸν τοῦ λιβάνου.
\vs{7}Ὅλη καλὴ εἶ πλησίον μου, καὶ μῶμος οὐκ ἔστιν ἐν σοί.

\vs{8}Δεῦρο ἀπὸ Λιβάνου νύμφη, δεῦρο ἀπὸ Λιβάνου· ἐλεύσῃ καὶ διελεύσῃ ἀπὸ ἀρχῆς Πίστεως, ἀπὸ κεφαλῆς Σανὶρ καὶ Ἑρμὼν, ἀπὸ μανδρῶν λεόντων, ἀπὸ ὀρέων παρδάλεων.
\vs{9}Ἐκαρδίωσας ἡμᾶς ἀδελφή μου νύμφη, ἐκαρδίωσας ἡμᾶς ἑνὶ ἀπὸ ὀφθαλμῶν σου, ἐν μιᾷ ἐνθέματι τραχήλων σου.
\vs{10}Τί ἐκαλλιώθησαν μαστοί σου ἀδελφή μου, νύμφη; τί ἐκαλλιώθησαν μαστοί σου ἀπὸ οἴνου, καὶ ὀσμὴ ἱματίων σου ὑπὲρ πάντα ἀρώματα;
\vs{11}Κηρίον ἀποστάζουσι χείλη σου νύμφη· μέλι καὶ γάλα ὑπὸ τὴν γλῶσσάν σου· καὶ ὀσμὴ ἱματίων σου, ὡς ὀσμὴ Λιβάνου.
\vs{12}Κῆπος κεκλεισμένος ἀδελφή μου νύμφη, κῆπος κεκλεισμένος, πηγὴ ἐσφραγισμένη·
\vs{13}Ἀποστολαί σου παράδεισος ῥοῶν μετὰ καρποῦ ἀκροδρύων, κύπροι μετὰ νάρδων·
\vs{14}Νάρδος καὶ κρόκος, κάλαμος καὶ κιννάμωμον, μετὰ πάντων ξύλων τοῦ Λιβάνου, σμύρνα, ἀλὼθ, μετὰ πάντων πρώτων μύρων,
\vs{15}πηγὴ κήπου, καὶ φρέαρ ὕδατος ζῶντος καὶ ῥοιζοῦντος ἀπὸ τοῦ Λιβάνου.

\vs{16}Ἐξεγέρθητι βοῤῥᾶ, καὶ ἔρχου Νότε, καὶ διάπνευσον κῆπόν μου, καὶ ῥευσάτωσαν ἀρώματά μου.

\ch{5}
Καταβήτω ἀδελφιδός μου εἰς κῆπον αὐτοῦ, καὶ φαγέτω καρπὸν ἀκροδρύων αὐτοῦ· εἰσῆλθον εἰς κῆπόν μου ἀδελφή μου νύμφη· ἐτρύγησα σμύρναν μου μετὰ ἀρωμάτων μου, ἔφαγον ἄρτον μου μετὰ μέλιτός μου, ἔπιον οἶνόν μου μετὰ γάλακτός μου· φάγετε πλήσιοι καὶ πίετε, καὶ μεθύσθητε ἀδελφοί.

\vs{2}Ἐγὼ καθεύδω, καὶ ἡ καρδία μου ἀγρυπνεῖ. φωνὴ ἀδελφιδοῦ μου κρούει ἐπὶ τὴν θύραν, ἀνοιξόν μοι ἡ πλησίον μου, ἄδελφή μου, περιστερά μου, τελεία μου· ὅτι ἡ κεφαλή μου ἐπλήσθη δρόσου, καὶ οἱ βόστρυχοί μου ψεκάδων νυκτός.
\vs{3}Ἐξεδυσάμην τὸν χιτῶνά μου, πῶς ἐνδύοσμαι αὐτόν; ἐνιψάμην τοὺς πόδας μου, πῶς μολυνῶ αὐτούς;
\vs{4}Ἀδελφιδός μου ἀπέστειλε χεῖρα αὐτοῦ ἀπὸ τῆς ὀπῆς, καὶ ἡ κοιλία μου ἐθροήθη ἐπʼ αὐτόν.
\vs{5}Ἀνέστην ἐγὼ ἀνοῖξαι τῷ ἀδελφιδῷ μου, χεῖρές μου ἔσταξαν σμύρναν, δάκτυλοί μου σμύρναν πλήρη ἐπὶ χεῖρας τοῦ κλείθρου.
\vs{6}Ἤνοιξα ἐγὼ τῷ ἀδελφιδῷ μου· ἀδελφιδός μου παρῆλθε· ψυχή μου ἐξῆλθεν ἐν λόγῳ αὐτοῦ· ἐζήτησα αὐτὸν καὶ οὐχ εὗρον αὐτὸν, ἐκάλεσα αὐτὸν καὶ οὐχ ὑπήκουσέ μου.
\vs{7}Εὕροσάν με οἱ φύλακες οἱ κυκλοῦντες ἐν τῇ πόλει, ἐπάταξάν με, ἐτραυμάτισάν με· ᾖραν τὸ θέριστρόν μου ἀπʼ ἐμοῦ φύλακες τῶν τειχέων.
\vs{8}Ὥρκισα ὑμᾶς θυγατέρες Ἱερουσαλὴμ ἐν ταῖς δυνάμεσι καὶ ἐν ταῖς ἰσχύσεσι τοῦ ἀγροῦ· ἐὰν εὕρητε τὸν ἀδελφιδόν μου, τί ἀπαγγείλητε αὐτῷ; ὅτι τετρωμένη ἀγάπης ἐγώ εἰμι.

\vs{9}Τί ἀδελφιδός σου ἀπὸ ἀδελφιδοῦ, ἡ καλὴ ἐν γυναιξί; τί ἀδελφιδός σου ἀπὸ ἀδελφιδοῦ, ὅτι οὕτως ὥρκισας ἡμᾶς;

\vs{10}Ἀδελφιδός μου λευκὸς καὶ πυῤῥὸς, ἐκλελοχισμένος ἀπὸ μυριάδων.
\vs{11}Κεφαλὴ αὐτοῦ χρυσίον κεφὰζ, βόστρυχοι αὐτοῦ ἐλάται, μέλανες ὡς κόραξ.
\vs{12}Ὀφθαλμοὶ αὐτοῦ ὡς περιστεραὶ ἐπὶ πληρώματα ὑδάτων, λελουσμέναι ἐν γάλακτι, καθήμεναι ἐπὶ πληρώματα.
\vs{13}Σιαγόνες αὐτοῦ ὡς φιάλαι τοῦ ἀρώματος φύουσαι μυρεψικά· χείλη αὐτοῦ κρίνα στάζοντα σμύρναν πλήρη.
\vs{14}Χεῖρες αὐτοῦ τορευταὶ χρυσαῖ πεπληρωμέναι Θαρσίς· κοιλία αὐτοῦ πυξίον ἐλεφάντινον ἐπὶ λίθου σαπφείρου.
\vs{15}Κνῆμαι αὐτοῦ στύλοι μαρμάρινοι τεθεμελιωμένοι ἐπὶ βάσεις χρυσᾶς· εἶδος αὐτοῦ ὡς Λίβανος, ἐκλεκτὸς ὡς κέδροι.
\vs{16}Φάρυγξ αὐτοῦ γλυκασμοὶ καὶ ὅλος ἐπιθυμία· οὗτος ἀδελφιδός μου καὶ οὗτος πλησίον μου, θυγατέρες Ἱερουσαλήμ.

\ch{6}
Ποῦ ἀπῆλθεν ὁ ἀδελφιδός σου ἡ καλὴ ἐν γυναιξὶ; ποῦ ἀπέβλεψεν ὁ ἀδελφιδός σου; καὶ ζητήσομεν αὐτὸν μετὰ σοῦ.

\vs{2}Ἀδελφιδός μου κατέβη εἰς κῆπον αὐτοῦ εἰς φιάλας τοῦ ἀρώματος, ποιμαίνειν ἐν κήποις, καὶ συλλέγειν κρίνα.
\vs{3}Ἐγὼ τῷ ἀδελφιδῷ μου, καὶ ἀδελφιδός μου ἐμοί, ὁ ποιμαίνων ἐν τοῖς κρίνοις.

\vs{4}Καλὴ εἶ ἡ πλησίον μου, ὡς εὐδοκία, ὡραῖα ὡς Ἱερουσαλὴμ, θάμβος ὡς τεταγμέναι.
\vs{5}Ἀπόστρεψον ὀφθαλμούς σου ἀπεναντίον μου, ὅτι αὐτοὶ ἀνεπτέρωσάν με· τρίχωμά σου ὡς ἀγέλαι τῶν αἰγῶν, αἳ ἀνεφάνησαν ἀπὸ τοῦ Γαλαάδ.
\vs{6}Ὀδόντες σου ὡς ἀγέλαι τῶν κεκαρμένων, αἳ ἀνέβησαν ἀπὸ τοῦ λουτροῦ, αἱ πᾶσαι διδυμεύουσαι, καὶ ἀτεκνοῦσα οὐκ ἔστιν ἐν αὐταῖς· ὡς σπαρτίον τὸ κόκκινον χείλη σου, καὶ ἡ λαλιά σου ὡραιᾶ.
\vs{7}Ὡς λέπυρον ῥοᾶς μῆλόν σου ἐκτὸς τῆς σιωπήσεώς σου.

\vs{8}Ἑξήκοντά εἰσι βασίλισσαι καὶ ὀγδοήκοντα παλλακαὶ, καὶ νεάνιδες ὧν οὐκ ἔστιν ἀριθμός.
\vs{9}Μία ἐστὶ περιστερά μου, τελεία μου, μία ἐστὶ τῇ μητρὶ αὐτῆς, ἐκλεκτή ἐστι τῇ τεκούσῃ αὐτήν· Εἴδοσαν αὐτὴν θυγατέρες καὶ μακαριοῦσιν αὐτὴν, βασίλισσαι καί γε παλλακαὶ, καὶ αἰνέσουσιν αὐτήν.
\vs{10}Τίς αὕτη ἡ ἐκκύπτουσα ὡσεὶ ὄρθρος, καλὴ ὡς σελήνη, ἐκλεκτὴ ὡς ὁ ἥλιος, θάμβος ὡς τεταγμέναι;

\vs{11}Εἰς κῆπον καρύας κατέβην ἰδεῖν ἐν γεννήμασι τοῦ χειμάῤῥου, ἰδεῖν εἰ ἤνθησεν ἡ ἄμπελος, ἐξήνθησαν αἱ ῥοαί·
\vs{12}Ἐκεῖ δώσω τοὺς μαστούς μου σοί· οὐκ ἔγνω ἡ ψυχή μου· ἔθετο με ἅρματα Ἀμιναδάβ.

\ch{7}
Ἐπίστρεφε ἐπίστρεφε ἡ Σουναμίτις· ἐπίστρεφε ἐπίστρεφε, καὶ ὀφόμεθα ἐν σοί.

Τί ὄψεσθε ἐν τῇ Σουναμίτιδι; ἡ ἐρχομένη ὡς χοροὶ τῶν παρεμβολῶν.

\vs{2}Ὡραιώθησαν διαβήματά σου ἐν ὑποδήμασί σου, θύγατερ ναδάβ· ῥυθμοὶ μηρῶν ὅμοιοι ὁρμίσκοις, ἔργον τεχνίτου.
\vs{3}Ὀμφαλός σου κρατὴρ τορευτὸς, μὴ ὑστερούμενος κράμα· κοιλία σου θημωνία σίτου πεφραγμένη ἐν κρίνοις.
\vs{4}Δύο μαστοί σου, ὡς δύο νεβροὶ δίδυμοι δορκάδος.
\vs{5}Ὁ τράχηλός σου ὡς πύργος ἐλεφάντινος· οἱ ὀφθαλμοί σου ὡς λίμναι ἐν Ἐσεβὼν, ἐν πύλαις θυγατρὸς πολλῶν· μυκτήρ σου, ὡς πύργος τοῦ Λιβάνου σκοπεύων πρόσωπον Δαμασκοῦ.
\vs{6}Κεφαλή σου ἐπὶ σὲ ὡς Κάρμηλος, καὶ πλόκιον κεφαλῆς σου ὡς πορφύρα· βασιλεὺς δεδεμένος ἐν παραδρομαῖς.
\vs{7}Τί ὡραιώθης, καὶ τί ἡδύνθης ἀγάπη;
\vs{8}ἐν τρυφαῖς σου τοῦτο μέγεθός σου· ὡμοιώθης τῷ φοίνικι, καὶ οἱ μαστοί σου τοῖς βότρυσιν.
\vs{9}Εἶπα, ἀναβήσομαι ἐπὶ τῷ φοίνικι, κρατήσω τῶν ὕψεων αὐτοῦ· καὶ ἔσονται δὴ μαστοί σου ὡς βότρυες τῆς ἀμπέλου, καὶ ὀσμὴ ῥινός σου ὡς μῆλα,
\vs{10}καὶ ὁ λάρυγξ σου ὡς οἶνος ὁ ἀγαθὸς, πορευόμενος τῷ ἀδελφιδῷ μου εἰς εὐθύτητα, ἱκανούμενος χείλεσί μου καὶ ὀδοῦσιν.

\vs{11}Ἐγὼ τῷ ἀδελφιδῷ μου, καὶ ἐπʼ ἐμὲ ἡ ἐπιστροφὴ αὐτοῦ.
\vs{12}Ἐλθὲ ἀδελφιδέ μου, ἐξέλθωμεν εἰς ἀγρὸν, αὐλισθῶμεν ἐν κώμαις.
\vs{13}Ὀρθρίσωμεν εἰς ἀμπελῶνας· ἴδωμεν εἰ ἤνθησεν ἡ ἄμπελος, ἤνθησεν ὁ κυπρισμὸς, ἤνθησν αἱ ῥοαί· ἐκεῖ δώσω τοὺς μαστούς μου σοί.
\vs{14}Οἱ μανδραγόραι ἔδωκαν ὀσμήν· καὶ ἐπὶ θύραις ἡμῶν πάντα ἀκρόδρυα νέα πρὸς παλαιὰ, ἀδελφιδέ μου, ἐτήρησά σοι.

\ch{8}
Τίς δῴη σε, ἀδελφιδέ μου, θηλάζοντα μαστοὺς μητρός μου; εὑροῦσά σε ἔξω φιλήσω σε, καί γε οὐκ ἐξουδενώσουσί μοι.
\vs{2}Παραλήψομαί σε, εἰσάξω σε εἰς οἶκον μητρός μου καὶ εἰς ταμεῖον τῆς συλλαβούσης με· ποτιῶ σε ἀπὸ οἴνου τοῦ μυρεψικοῦ, ἀπὸ νάματος ῥοῶν μου.

\vs{3}Εὐώνυμος αὐτοῦ ὑπὸ τὴν κεφαλήν μου, καὶ ἡ δεξιὰ αὐτοῦ περιλήψεταί με.

\vs{4}Ὥρκισα ὑμᾶς θυγατέρες Ἱερουσαλὴμ ἐν ταῖς ἰσχύσεσι τοῦ ἀγροῦ· ἐὰν ἐγείρητε καὶ ἐὰν ἐξεγείρητε τὴν ἀγάπην ἕως ἂν θελήσῃ.

\vs{5}Τίς αὕτη ἡ ἀναβαίνουσα λελευκανθισμένη, ἐπιστηριζομένη ἐπὶ τὸν ἀδελφιδὸν αὐτῆς; Ὑπὸ μῆλον ἐξήγειρά σε· ἐκεῖ ὠδίνησέ σε ἡ μήτηρ σου, ἐκεῖ ὠδίνησέ σε ἡ τεκοῦσά σε.

\vs{6}Θές με ὡς σφραγίδα ἐπὶ τὴν καρδίαν σου, ὡς σφραγίδα ἐπὶ τὸν βραχίονά σου· ὅτι κραταιὰ ὡς θάνατος ἀγάπη, σκληρὸς ὡς ᾅδης ζῆλος· περίπτερα αὐτῆς περίπτερα πυρὸς, φλόγες αὐτῆς.

\vs{7}Ὕδωρ πολὺ οὐ δυνήσεται σβέσαι τὴν ἀγάπην, καὶ ποταμοὶ οὐ συνκλύσουσιν αὐτήν· ἐὰν δῷ ἀνὴρ πάντα τὸν βίον αὐτοῦ ἐν τῇ ἀγάπῃ, ἐξουδενώσει ἐξουδενώσουσιν αὐτόν.

\vs{8}Ἀδελφή ἡμῶν μικρὰ καὶ μαστοὺς οὐκ ἔχει· τί ποιήσωμεν τῇ ἀδελφῇ ἡμῶν, ἐν ἡμέρᾳ ᾗ ἐὰν λαληθῇ ἐν αὐτῇ;
\vs{9}Εἰ τεῖχός ἐστιν, οἰκοδομήσωμεν ἐπʼ αὐτὴν ἐπάλξεις ἀργυρᾶς· καὶ εἰ θύρα ἐστὶ, διαγράψωμεν ἐπʼ αὐτὴν σανίδα κεδρίνην.
\vs{10}Ἐγὼ τεῖχος, καὶ μαστοί μου ὡς πύργοι ἐγὼ ἤμην ἐν ὀφθαλμοῖς αὐτῶν ὡς εὑρίσκουσα εἰρήνην.
\vs{11}Ἀμπελὼν ἐγενήθη τῷ Σαλωμὼν ἐν Βεεθλαμών· ἔδωκε τὸν ἀμπελῶνα αὐτοῦ τοῖς τηροῦσιν· ἀνὴρ οἴσει ἐν καρπῷ αὐτοῦ χιλίους ἀργυρίου.
\vs{12}Ἀμπελών μου ἐμὸς ἐνώπιόν μου, οἱ χίλιοι Σαλωμὼν, καὶ οἱ διακόσιοι τοῖς τηροῦσι τὸν καρπὸν αὐτοῦ.

\vs{13}Ὁ καθήμενος ἐν κήποις, ἑταῖροι προσέχοντες τῇ φωνῇ σου, ἀκούτισόν με.

\vs{14}Φύγε ἀδελφιδέ μου, καὶ ὁμοιώθητι τῇ δορκάδι, ἢ τῷ νεβρῷ τῶν ἐλάφων ἐπὶ ὄρη τῶν ἀρωμάτων.


\def\book{ΗΣΑΙΑΣ}
\biblebook{ΗΣΑΙΑΣ}


\lettrine[lines=2, loversize=0.2, nindent=0em, findent=.25em]{\textcolor{bookheadingcolor}{Ὅ}}{ΡΑΣΙΣ} ἣν εἶδεν Ἡσαΐας υἱὸς Ἀμὼς, ἣν εἶδε κατὰ τῆς Ἰουδαίας καὶ κατὰ Ἱερουσαλὴμ, ἐν βασιλείᾳ Ὀζίου, καὶ Ἰωάθαμ, καὶ Ἄχαζ, καὶ Ἐζεκίου οἳ ἐβασίλευσαν τῆς Ἰουδαίας.

\vs{2}Ἄκουε οὐρανὲ, καὶ ἐνωτίζου γῆ, ὅτι Κύριος ἐλάλησεν, υἱοὺς ἐγέννησα καὶ ὕψωσα, αὐτοὶ δέ με ἠθέτησαν.
\vs{3}Ἔγνω βοῦς τὸν κτησάμενον, καὶ ὄνος τὴν φάτνην τοῦ κυρίου αὐτοῦ· Ἰσραὴλ δέ με οὐκ ἔγνω, καὶ ὁ λαός με οὐ συνῆκεν.

\vs{4}Οὐαὶ ἔθνος ἁμαρτωλὸν, λαὸς πλήρης ἁμαρτιῶν, σπέρμα πονηρὸν, υἱοὶ ἄνομοι· ἐγκατελίπατε τὸν Κύριον, καὶ παρωργίσατε τὸν ἅγιον τοῦ Ἰσραήλ.
\vs{5}Τί ἔτι πληγῆτε προστιθέντες ἀνομίαν; πᾶσα κεφαλὴ εἰς πόνον, καὶ πᾶσα καρδία εἰς λύπην·
\vs{6}Ἀπὸ ποδῶν ἕως κεφαλῆς, οὐκ ἔστιν ἐν αὐτῷ ὁλοκληρία, οὔτε τραῦμα, οὔτε μώλωψ, οὔτε πληγὴ φλεγμαίνουσα· οὐκ ἔστι μάλαγμα ἐπιθεῖναι, οὔτε ἔλαιον, οὔτε καταδέσμους.
\vs{7}Ἡ γῆ ὑμῶν ἔρημος, αἱ πόλεις ὑμῶν πυρίκαυστοι· τὴν χώραν ὑμῶν ἐνώπιον ὑμῶν ἀλλότριοι κατεσθίουσιν αὐτὴν, καὶ ἠρήμωται κατεστραμμένη ὑπὸ λαῶν ἀλλοτρίων.
\vs{8}Ἐγκαταλειφθήσεται ἡ θυγάτηρ Σιὼν, ὡς σκηνὴ, ἐν ἀμπελῶνι, καὶ ὡς ὀπωροφυλάκιον ἐν σικυηράτῳ, ὡς πόλις πολιορκουμένη.
\vs{9}Καὶ εἰ μὴ Κύριος σαβαὼθ ἐγκατέλιπεν ἡμῖν σπέρμα, ὡς Σόδομα ἂν ἐγενήθημεν, καὶ ὡς Γόμοῤῥα ἂν ὁμοίωθημεν.

\vs{10}Ἀκούσατε λόγον Κυρίου, ἄρχοντες Σοδόμων· προσέχετε νόμον Θεοῦ, λαὸς Γομόῤῥας.
\vs{11}Τί μοι πλῆθος τῶν θυσιῶν ὑμῶν; λέγει Κύριος· πλήρης εἰμὶ ὁλοκαυτωμάτων κριῶν, καὶ στέαρ ἀρνῶν καὶ αἷμα ταύρων καὶ τράγων οὐ βούλομαι,
\vs{12}οὐδʼ ἂν ἔρχησθε ὀφθῆναί μοι· τίς γὰρ ἐξεζήτησε ταῦτα ἐκ τῶν χειρῶν ὑμῶν; πατεῖν τὴν αὐλήν μου οὐ προσθήσεθε.
\vs{13}Ἐὰν φέρητε σεμίδαλιν, μάταιον θυμίαμα, βδέλυγμά μοι ἐστι· τὰς νουμηνίας ὑμῶν, καὶ τὰ σάββατα, καὶ ἡμέραν μεγάλην οὐκ ἀνέχομαι· νηστείαν, καὶ ἀργίαν,
\vs{14}καὶ τὰς νουμηνίας ὑμῶν, καὶ τὰς ἑορτὰς ὑμῶν μισεῖ ἡ ψυχή μου· ἐγενήθητέ μοι εἰς πλησμονὴν, οὐκέτι ἀνήσω τὰς ἁμαρτίας ὑμῶν.
\vs{15}Ὅταν ἐκτείνητε τὰς χεῖρας, ἀποστρέψω τοὺς ὀφθαλμούς μου ἀφʼ ὑμῶν· καὶ ἐὰν πληθύνητε τὴν δέησιν, οὐκ εἰσακούσομαι ὑμῶν· αἱ γὰρ χεῖρες ὑμῶν αἵματος πλήρεις.

\vs{16}Λούσασθε, καθαροὶ γένεσθε, ἀφέλετε τὰς πονηρίας ἀπὸ τῶν ψυχῶν ὑμῶν, ἀπέναντι τῶν ὀφθαλμῶν μου· παύσασθε ἀπὸ τῶν πονηριῶν ὑμῶν,
\vs{17}μάθετε καλὸν ποιεῖν, ἐκζητήσατε κρίσιν, ῥύσασθε ἀδικούμενον, κρίνατε ὀρφανῷ, καὶ δικαιώσατε χήραν.

\vs{18}Καὶ δεῦτε, διελεγχθῶμεν, λέγει Κύριος· καὶ ἐὰν ὦσιν αἱ ἁμαρτίαι ὑμῶν ὡς φοινικοῦν, ὡς χιόνα λευκανῶ· ἐὰν δὲ ὦσιν ὡς κόκκινον, ὡς ἔριον λευκανῶ.
\vs{19}Καὶ ἐὰν θέλητε, καὶ εἰσακούσητέ μου, τὰ ἀγαθὰ τῆς γῆς φάγεσθε.
\vs{20}Ἐὰν δὲ μὴ θέλητε, μηδὲ εἰσακούσητέ μου, μάχαιρα ὑμᾶς κατέδεται· τὸ γὰρ στόμα Κυρίου ἐλάλησε ταῦτα.

\vs{21}Πῶς ἐγένετο πόρνη πόλις πιστὴ Σιὼν πλήρης κρίσεως; ἐν ᾗ δικαιοσύνη ἐκοιμήθη ἐν αὐτῇ, νῦν δὲ φονευταί.
\vs{22}Τὸ ἀργύριον ὑμῶν ἀδόκιμον· οἱ κάπηλοί σου μίσγουσι τὸν οἶνον ὕδατι.
\vs{23}Οἱ ἄρχοντές σου ἀπειθοῦσι, κοινωνοὶ κλεπτῶν, ἀγαπῶντες δῶρα, διώκοντες ἀνταπόδομα, ὀρφανοῖς οὐ κρίνοντες, καὶ κρίσιν χηρῶν οὐ προσέχοντες.

\vs{24}Διατοῦτο τάδε λέγει Κύριος ὁ δεσπότης σαβαὼθ, οὐαὶ οἱ ἰσχύοντες Ἰσραήλ· οὐ παύσεται γάρ μου ὁ θυμὸς ἐν τοῖς ὑπεναντίοις, καὶ κρίσιν ἐκ τῶν ἐχθρῶν μου ποιήσω.
\vs{25}Καὶ ἐπάξω τὴν χεῖρά μου ἐπὶ σὲ, καὶ πυρώσω εἰς καθαρόν, τοὺς δὲ ἀπειθοῦντας ἀπολέσω, καὶ ἀφελῶ πάντας ἀνόμους ἀπὸ σοῦ·
\vs{26}Καὶ ἐπιστήσω τοὺς κριτάς σου ὡς τὸ πρότερον, καὶ τοὺς συμβούλους σου ὡς τὸ ἀπʼ ἀρχῆς· καὶ μετὰ ταῦτα κληθήσῃ πόλις δικαιοσύνης, μητρόπολις πιστὴ Σιών·
\vs{27}μετὰ γὰρ κρίματος σωθήσεται ἡ αἰχμαλωσία αὐτῆς, καὶ μετὰ ἐλεημοσύνης.
\vs{28}Καὶ συντριβήσονται οἱ ἄνομοι καὶ οἱ ἁμαρτωλοὶ ἅμα, καὶ οἱ ἐγκαταλιπόντες τὸν Κύριον συντελεσθήσονται.
\vs{29}Διότι αἰσχυνθήσονται ἀπὸ τῶν εἰδώλων αὐτῶν ἃ αὐτοὶ ἠβούλοντο, καὶ ᾐσχύνθησαν ἐπὶ τοῖς κήποις, ἃ ἐπεθύμησαν.
\vs{30}Ἔσονται γὰρ ὡς τερέβινθος ἀποβεβληκυῖα τὰ φύλλα, καὶ ὡς παράδεισος ὕδωρ μὴ ἔχων.
\vs{31}Καὶ ἔσται ἡ ἰσχὺς αὐτῶν ὡς καλάμη στιππύου, καὶ αἱ ἐργασίαι αὐτῶν ὡς σπινθῆρες, καὶ κατακαυθήσονται οἱ ἄνομοι καὶ οἱ ἁμαρτωλοὶ ἅμα, καὶ οὐκ ἔσται ὁ σβέσων.

\ch{2}
Ὁ λόγος ὁ γενόμενος πρὸς Ἡσαΐαν υἱὸν Ἀμὼς περὶ τῆς Ἰουδαίας, καὶ περὶ Ἱερουσαλήμ.

\vs{2}Ὅτι ἔσται ἐν ταῖς ἐσχάταις ἡμέραις ἐμφανὲς τὸ ὄρος Κυρίου, καὶ ὁ οἶκος τοῦ Θεοῦ ἐπʼ ἄκρου τῶν ὀρέων, καὶ ὑψωθήσεται ὑπεράνω τῶν βουνῶν, καὶ ἥξουσιν ἐπʼ αὐτὸ πάντα τὰ ἔθνη.
\vs{3}Καὶ πορεύσονται ἔθνη πολλὰ, καὶ ἐροῦσι, δεῦτε καὶ ἀναβῶμεν εἰς τὸ ὄρος Κυρίου, καὶ εἰς τὸν οἶκον τοῦ Θεοῦ Ἰακὼβ, καὶ ἀναγγελεῖ ἡμῖν τὴν ὁδὸν αὐτοῦ, καὶ πορευσόμεθα ἐν αὐτῇ· ἐκ γὰρ Σιὼν ἐξελεύσεται νόμος, καὶ λόγος Κυρίου ἐξ Ἱερουσαλήμ.
\vs{4}Καὶ κρινεῖ ἀναμέσον τῶν ἐθνῶν, καὶ ἐξελέγξει λαὸν πολύν· καὶ συγκόψουσι τὰς μαχαίρας αὐτῶν εἰς ἄροτρα, καὶ τὰς ζιβύνας αὐτῶν εἰς δρέπανα· καὶ οὐ λήψεται ἔθνος ἐπʼ ἔθνος μάχαιραν, καὶ οὐ μὴ μάθωσιν ἔτι πολεμεῖν.

\vs{5}Καὶ νῦν ὁ οἶκος Ἰακὼβ, δεῦτε πορευθῶμεν τῷ φωτὶ Κυρίου.
\vs{6}Ἀνῆκεν γὰρ τὸν λαὸν αὐτοῦ τὸν οἶκον τοῦ Ἰσραήλ· ὅτι ἐνεπλήσθη ὡς τὸ ἀπʼ ἀρχῆς ἡ χώρα αὐτῶν κληδονισμῶν, ὡς ἡ τῶν ἀλλοφύλων, καὶ τέκνα πολλὰ ἀλλόφυλα ἐγενήθη αὐτοῖς.
\vs{7}Ἐνεπλήσθη γὰρ ἡ χώρα αὐτῶν ἀργυρίου καὶ χρυσίου, καὶ οὐκ ἦν ἀριθμὸς τῶν θησαυρῶν αὐτῶν· καὶ ἐνεπλήσθη ἡ γῆ ἵππων, καὶ οὐκ ἦν ἀριθμὸς τῶν ἁρμάτων αὐτῶν.
\vs{8}Καὶ ἐνεπλήσθη ἡ γῆ βδελυγμάτων τῶν ἔργων τῶν χειρῶν αὐτῶν, καὶ προσεκύνησαν οἷς ἐποίησαν οἱ δάκτυλοι αὐτῶν.
\vs{9}Καὶ ἔκυψεν ἄνθρωπος, καὶ ἐταπεινώθη ἀνήρ, καὶ οὐ μὴ ἀνήσω αὐτούς.

\vs{10}Καὶ νῦν εἰσέλθετε εἰς τὰς πέτρας, καὶ κρύπτεσθε εἰς τὴν γῆν ἀπὸ προσώπου τοῦ φόβου Κυρίου, καὶ ἀπὸ τῆς δόξης τῆς ἰσχύος αὐτοῦ, ὅταν ἀναστῇ θραῦσαι τὴν γῆν.
\vs{11}Οἱ γὰρ ὀφθαλμοὶ Κυρίου ὑψηλοί, ὁ δὲ ἄνθρωπος ταπεινός· καὶ ταπεινωθήσεται τὸ ὕψος τῶν ἀνθρώπων, καὶ ὑψωθήσεται Κύριος μόνος ἐν τῇ ἡμέρᾳ ἐκείνῃ.

\vs{12}Ἡμέρα γὰρ Κυρίου σαβαὼθ ἐπὶ πάντα ὑβριστὴν καὶ ὑπερήφανον καὶ ἐπὶ πάντα ὑψηλὸν καὶ μετέωρον, καὶ ταπεινωθήσονται.
\vs{13}Καὶ ἐπὶ πᾶσαν κέδρον τοῦ Λιβάνου τῶν ὑψηλῶν καὶ μετεώρων, καὶ ἐπὶ πᾶν δένδρον βαλάνου Βασάν,
\vs{14}καὶ ἐπὶ πᾶν ὑψηλὸν ὄρος, καὶ ἐπὶ πάντα βουνὸν ὑψηλόν,
\vs{15}καὶ ἐπὶ πάντα πύργον ὑψηλὸν, καὶ ἐπὶ πᾶν τεῖχος ὑψηλόν,
\vs{16}καὶ ἐπὶ πᾶν πλοῖον θαλάσσης, καὶ ἐπὶ πᾶσαν θέαν πλοίων κάλλους.
\vs{17}Καὶ ταπεινωθήσεται πᾶς ἄνθρωπος, καὶ πεσεῖται ὕβρις τῶν ἀνθρώπων, καὶ ὑψωθήσεται Κύριος μόνος ἐν τῇ ἡμέρᾳ ἐκείνῃ.
\vs{18}Καὶ τὰ χειροποίητα πάντα κατακρύψουσιν,
\vs{19}εἰσενέγκαντες εἰς τὰ σπήλαια, καὶ εἰς τὰς σχισμὰς τῶν πετρῶν, καὶ εἰς τὰς τρώγλας τῆς γῆς, ἀπὸ προσώπου τοῦ φόβου Κυρίου, καὶ ἀπὸ τῆς δόξης τῆς ἰσχύος αὐτοῦ, ὅταν ἀναστῇ θραῦσαι τὴν γῆν.
\vs{20}τῇ γὰρ ἡμέρᾳ ἐκείνῃ ἐκβαλεῖ ἄνθρωπος τὰ βδελύγματα αὐτοῦ τὰ ἀργυρᾶ καὶ τὰ χρυσᾶ, ἃ ἐποίησαν προσκυνεῖν τοῖς ματαίοις καὶ ταῖς νυκτερίσι,
\vs{21}τοῦ εἰσελθεῖν εἰς τὰς τρώγλας τῆς στερεᾶς πέτρας, καὶ εἰς τὰς σχισμὰς τῶν πετρῶν, ἀπὸ προσώπου τοῦ φόβου Κυρίου, καὶ ἀπὸ τῆς δόξης τῆς ἰσχύος αὐτοῦ, ὅταν ἀναστῇ θραῦσαι τὴν γῆν.

\ch{3}
Ἰδοὺ δὴ ὁ δεσπότης Κύριος σαβαὼθ ἀφελεῖ ἀπὸ Ἱερουσαλὴμ, καὶ ἀπὸ τῆς Ἰουδαίας, ἰσχύοντα καὶ ἰσχύουσαν, ἰσχὺν ἄρτου καὶ ἰσχὺν ὕδατος,
\vs{2}γίγαντα καὶ ἰσχύοντα, καὶ ἄνθρωπον πολεμιστὴν, καὶ δικαστὴν, καὶ προφήτην, καὶ στοχαστὴν, καὶ πρεσβύτερον,
\vs{3}καὶ πεντηκόνταρχον, καὶ θαυμαστὸν σύμβουλον, καὶ σοφὸν ἀρχιτέκτονα, καὶ συνετὸν ἀκροατήν.
\vs{4}Καὶ ἐπιστήσω νεανίσκους ἄρχοντας αὐτῶν, καὶ ἐμπαίκται κυριεύσουσιν αὐτῶν.
\vs{5}Καὶ συμπεσεῖται ὁ λαὸς, ἄνθρωπος πρὸς ἄνθρωπον, καὶ ἄνθρωπος πρὸς τὸν πλησίον αὐτοῦ· προσκόψει τὸ παιδίον πρὸς τὸν πρεσβύτην, ὁ ἄτιμος πρὸς τὸν ἔντιμον.
\vs{6}Ὅτι ἐπιλήψεται ἄνθρωπος τοῦ ἀδελφοῦ αὐτοῦ, ἢ τοῦ οἰκείου τοῦ πατρὸς αὐτοῦ, λέγων, ἱμάτιον ἔχεις, ἀρχηγὸς γενοῦ ἡμῶν, καὶ τὸ βρῶμα τὸ ἐμὸν ὑπὸ σὲ ἔστω.
\vs{7}Καὶ ἀποκριθεὶς ἐν τῇ ἡμέρᾳ ἐκείνῃ ἐρεῖ, οὐκ ἔσομαί σου ἀρχηγός· οὐ γάρ ἐστιν ἐν τῷ οἴκῳ μου ἄρτος, οὐδὲ ἱμάτιον· οὐκ ἔσομαι ἀρχηγὸς τοῦ λαοῦ τούτου.
\vs{8}Ὅτι ἀνεῖται Ἱερουσαλὴμ, καὶ ἡ Ἰουδαία συμπέπτωκε, καὶ αἱ γλῶσσαι αὐτῶν μετὰ ἀνομίας, τὰ πρὸς Κύριον ἀπειθοῦντες. Διότι νῦν ἐταπεινώθη ἡ δόξα αὐτῶν,
\vs{9}καὶ ἡ αἰσχύνη τοῦ προσώπου αὐτῶν ἀντέστη αὐτοῖς· τὴν δὲ ἁμαρτίαν αὐτῶν ὡς Σοδόμων ἀνήγγειλαν καὶ ἐνεφάνισαν· οὐαὶ τῇ ψυχῇ αὐτῶν, διότι βεβούλευνται βουλὴν πονηρὰν, καθʼ ἑαυτῶν
\vs{10}εἰπόντες, δήσωμεν τὸν δίκαιον, ὅτι δύσχρηστος ἡμῖν ἐστι· τοίνυν τὰ γεννήματα τῶν ἔργων αὐτῶν φάγονται.
\vs{11}Οὐαὶ τῷ ἀνόμῳ, πονηρὰ κατὰ τὰ ἔργα τῶν χειρῶν αὐτοῦ συμβήσεται αὐτῷ.
\vs{12}Λαός μου, οἱ πράκτορες ὑμῶν καλαμῶνται ὑμᾶς, καὶ οἱ ἀπαιτοῦντες κυριεύουσιν ὑμῶν· λαός μου, οἱ μακαρίζοντες ὑμᾶς πλανῶσιν ὑμᾶς, καὶ τὸν τρίβον τῶν ποδῶν ὑμῶν ταράσσουσιν.

\vs{13}Ἀλλὰ νῦν καταστήσεται εἰς κρίσιν Κύριος, καὶ στήσει εἰς κρίσιν τὸν λαὸν αὐτοῦ.
\vs{14}Αὐτὸς Κύριος εἰς κρίσιν ἥξει μετὰ τῶν πρεσβυτέρων τοῦ λαοῦ, καὶ μετὰ τῶν ἀρχόντων αὐτοῦ· ὑμεῖς δὲ τί ἐνεπυρίσατε τὸν ἀμπελῶνά μου, καὶ ἡ ἁρπαγὴ τοῦ πτωχοῦ ἐν τοῖς οἴκοις ὑμῶν;
\vs{15}Τί ὑμεῖς ἀδικεῖτε τὸν λαόν μου, καὶ τὸ πρόσωπον τῶν πτωχῶν καταισχύνετε;

\vs{16}Τάδε λέγει Κύριος, ἀνθʼ ὧν ὑψώθησαν αἱ θυγατέρες Σιών, καὶ ἐπορεύθησαν ὑψηλῷ τραχήλῳ καὶ ἐν νεύμασιν ὀφθαλμῶν, καὶ τῇ πορείᾳ τῶν ποδῶν ἅμα σύρουσαι τοὺς χιτῶνας, καὶ τοῖς ποσὶν ἅμα παίζουσαι·
\vs{17}Καὶ ταπεινώσει ὁ Θεὸς ἀρχούσας θυγατέρας Σιών· καὶ Κύριος ἀνακαλύψει τὸ σχῆμα αὐτῶν
\vs{18}ἐν τῇ ἡμέρᾳ ἐκείνῃ, καὶ ἀφελεῖ Κύριος τὴν δόξαν τοῦ ἱματισμοῦ αὐτῶν, τὰ ἐμπλόκια, καὶ τοὺς κοσύμβους, καὶ τοὺς μηνίσκους,
\vs{19}καὶ τὸ κάθεμα, καὶ τὸν κόσμον τοῦ προσώπου αὐτῶν,
\vs{20}καὶ τὴν σύνθεσιν τοῦ κόσμου τῆς δόξης, καὶ τοὺς χλιδῶνας, καὶ τὰ ψέλλια, καὶ τὸ ἐμπλόκιον, καὶ τοὺς δακτυλίους, καὶ τὰ περιδέξια, καὶ τὰ ἐνώτια,
\vs{21}καὶ τὰ περιπόρφυρα, καὶ τὰ μεσοπόρφυρα,
\vs{22}καὶ τὰ ἐπιβλήματα τὰ κατὰ τὴν οἰκίαν, καὶ τὰ διαφανῆ Λακωνικά,
\vs{23}καὶ τὰ βύσσινα, καὶ τὰ ὑακίνθινα, καὶ κόκκινα, καὶ τὴν βύσσον, σὺν χρυσῷ καὶ ὑακίνθῳ συγκαθυφασμένα, καὶ θέριστρα κατάκλιτα.
\vs{24}Καὶ ἔσται ἀντὶ ὀσμῆς ἡδείας, κονιορτός· καὶ ἀντὶ ζώνης, σχοινίῳ ζώσῃ· καὶ ἀντὶ τοῦ κόσμου τῆς κεφαλῆς τοῦ χρυσίου, φαλάκρωμα ἕξεις διὰ τὰ ἔργα σοῦ· καὶ ἀντὶ τοῦ χιτῶνος τοῦ μεσοπορφύρου, περιζώσῃ σάκκον.
\vs{25}Καὶ ὁ υἱός σου ὁ κάλλιστος ὃν ἀγαπᾷς, μαχαίρᾳ πεσεῖται· καὶ οἱ ἰσχύοντες ὑμῶν, μαχαίρᾳ πεσοῦνται, καὶ ταπεινωθήσονται·
\vs{26}Καὶ πενθήσουσιν αἱ θῆκαι τοῦ κόσμου ὑμῶν· καὶ καταλειφθήσῃ μόνη, καὶ εἰς τὴν γῆν ἐδαφισθήσῃ.

\ch{4}
Καὶ ἐπιλήψονται ἑπτὰ γυναῖκες ἀνθρώπου ἑνὸς, λέγουσαι, τὸν ἄρτον ἡμῶν φαγόμεθα, καὶ τὰ ἱμάτια ἡμῶν περιβαλούμεθα, πλὴν τὸ ὄνομα τὸ σὸν κεκλήσθω ἐφʼ ἡμᾶς, ἄφελε τὸν ὀνειδισμὸν ἡμῶν.

\vs{2}Τῇ δὲ ἡμέρᾳ ἐκείνῃ ἐπιλάμψει ὁ Θεὸς ἐν βουλῇ μετὰ δόξης ἐπὶ τῆς γῆς, τοῦ ὑψῶσαι καὶ δοξάσαι τὸ καταλειφθὲν τοῦ Ἰσραήλ.
\vs{3}Καὶ ἔσται τὸ ὑπολειφθὲν ἐν Σιών, καὶ τὸ καταλειφθὲν ἐν Ἱερουσαλήμ, ἅγιοι κληθήσονται πάντες οἱ γραφέντες εἰς ζωὴν ἐν Ἱερουσαλὴμ.
\vs{4}Ὅτι ἐκπλυνεῖ Κύριος τὸν ῥύπον τῶν υἱῶν καὶ τῶν θυγατέρων Σειών, καὶ τὸ αἷμα ἐκκαθαριεῖ ἐκ μέσου αὐτῶν, ἐν πνεύματι κρίσεως καὶ πνεύματι καύσεως.
\vs{5}Καὶ ἥξει, καὶ ἔσται πᾶς τόπος τοῦ ὄρους Σιὼν καὶ πάντα τὰ περικύκλῳ αὐτῆς σκιάσει νεφέλη ἡμέρας, καὶ ὡς καπνοῦ καὶ φωτὸς πυρὸς καιομένου νυκτὸς, καὶ πάσῃ τῇ δόξῃ σκεπασθήσεται.
\vs{6}Καὶ ἔσται εἰς σκιὰν ἀπὸ καύματος, καὶ ἐν σκέπῃ καὶ ἐν ἀποκρύφῳ ἀπὸ σκληρότητος καὶ ὑετοῦ.

\ch{5}
Ἄσω δὴ τῷ ἠγαπημένῳ ᾆσμα τοῦ ἀγαπητοῦ μου τῷ ἀμπελῶνί μου.

Ἀμπελὼν ἐγενήθη τῷ ἠγαπημένῳ, ἐν κέρατι, ἐν τόπῳ πίονι.
\vs{2}Καὶ φραγμὸν περιέθηκα, καὶ ἐχαράκωσα, καὶ ἐφύτευσα ἄμπελον σωρὴκ, καὶ ᾠκοδόμησα πύργον ἐν μέσῳ αὐτοῦ, καὶ προλήνιον ὤρυξα ἐν αὐτῷ, καὶ ἔμεινα τοῦ ποιῆσαι σταφυλὴν, καὶ ἐποίησεν ἀκάνθας.
\vs{3}Καὶ νῦν οἱ ἐνοικοῦντες ἐν Ἱερουσαλὴμ, καὶ ἄνθρωπος τοῦ Ἰούδα, κρίνατε ἐν ἐμοὶ καὶ ἀναμέσον τοῦ ἀμπελῶνός μου.
\vs{4}Τί ποιήσω ἔτι τῷ ἀμπελῶνί μου, καὶ οὐκ ἐποίησα αὐτῷ; διότι ἔμεινα τοῦ ποιῆσαι σταφυλὴν, ἐποίησε δὲ ἀκάνθας.
\vs{5}Νῦν δὲ ἀναγγελῶ ὑμῖν τί ἐγὼ ποιήσω τῷ ἀμπελῶνί μου· ἀφελῶ τὸν φραγμὸν αὐτοῦ, καὶ ἔσται εἰς διαρπαγήν· καὶ καθελῶ τὸν τοῖχον αὐτοῦ, καὶ ἔσται εἰς καταπάτημα·
\vs{6}Καὶ ἀνήσω τὸν ἀμπελῶνά μου, καὶ οὐ τμηθῇ, οὐδὲ μὴ σκαφῇ· καὶ ἀναβήσονται εἰς αὐτὸν, ὡς εἰς χέρσον ἄκανθαι· καὶ ταῖς νεφέλαις ἐντελοῦμαι, τοῦ μὴ βρέξαι εἰς αὐτὸν ὑετόν.
\vs{7}Ὁ γὰρ ἀμπελὼν Κυρίου σαβαὼβ, οἶκος τοῦ Ἰσραὴλ, καὶ ἄνθρωπος τοῦ Ἰούδα νεόφυτον ἠγαπημένον· ἔμεινα τοῦ ποιῆσαι κρίσιν, ἐποίησε δὲ ἀνομίαν, καὶ οὐ δικαιοσύνην, ἀλλὰ κραυγήν.

\vs{8}Οὐαὶ οἱ συνάπτοντες οἰκίαν πρὸς οἰκίαν, καὶ ἀγρὸν πρὸς ἀγρὸν ἐγγίζοντες, ἵνα τοῦ πλησίον ἀφέλωνταί τι· μὴ οἰκήσετε μόνοι ἐπὶ τῆς γῆς;
\vs{9}Ἠκούσθη γὰρ εἰς τὰ ὦτα Κυρίου σαβαὼθ ταῦτα· ἐὰν γὰρ γένωνται οἰκίαι πολλαί, εἰς ἔρημον ἔσονται μεγάλαι καὶ καλαὶ, καὶ οὐκ ἔσονται οἱ ἐνοικοῦντες ἐν αὐταῖς.
\vs{10}Οὗ γὰρ ἐργῶνται δέκα ζεύγη βοῶν, ποιήσει κεράμιον ἕν· καὶ ὁ σπείρων ἀρτάβας ἓξ, ποιήσει μέτρα τρία.

\vs{11}Οὐαὶ οἱ ἐγειρόμενοι τοπρωῒ, καὶ τὸ σίκερα διώκοντες, οἱ μένοντες τὸ ὀψέ· ὁ γὰρ οἶνος αὐτοὺς συνκαύσει.
\vs{12}Μετὰ γὰρ κιθάρας καὶ ψαλτηρίου καὶ τυμπάνων καὶ αὐλῶν τον οἶνον πίνουσι, τὰ δὲ ἔργα Κυρίου οὐκ ἐμβλέπουσι, καὶ τὰ ἔργα τῶν χειρῶν αὐτοῦ οὐ κατανοοῦσι.

\vs{13}Τοίνυν αἰχμάλωτος ὁ λαός μου ἐγενήθη, διὰ τὸ μὴ εἰδέναι αὐτοὺς τὸν Κύριον· καὶ πλῆθος ἐγενήθη νεκρῶν, διὰ λιμὸν καὶ δίψος ὕδατος.
\vs{14}Καὶ ἐπλάτυνεν ὁ ᾅδης τὴν ψυχὴν αὐτοῦ, καὶ διήνοιξε τὸ στόμα αὐτοῦ, τοῦ μὴ διαλιπεῖν· καὶ καταβήσονται οἱ ἔνδοξοι καὶ οἱ μεγάλοι καὶ οἱ πλούσιοι καὶ οἱ λοιμοὶ αὐτῆς.
\vs{15}Καὶ ταπεινωθήσεται ἄνθρωπος, καὶ ἀτιμασθήσεται ἀνήρ· καὶ οἱ ὀφθαλμοὶ οἱ μετέωροι ταπεινωθήσονται.
\vs{16}Καὶ ὑψωθήσεται Κύριος σαβαὼθ ἐν κρίματι, καὶ ὁ Θεὸς ὁ ἅγιος δοξασθήσεται ἐν δικαιοσύνῃ·
\vs{17}Καὶ βοσκηθήσονται οἱ διηρπασμένοι ὡς ταῦροι, καὶ τὰς ἐρήμους τῶν ἀπειλημμένων ἄρνες φάγονται.

\vs{18}Οὐαὶ οἱ ἐπισπώμενοι τὰς ἁμαρτίας ὡς σχοινίῳ μακρῷ, καὶ ὡς ζυγοῦ ἱμάντι δαμάλεως τὰς ἀνομίας·
\vs{19}Οἱ λέγοντες, τὸ τάχος ἐγγισάτω ἃ ποιήσει, ἵνα ἴδωμεν· καὶ ἐλθάτω ἡ βουλὴ τοῦ ἁγίου Ἰσραὴλ, ἵνα γνῶμεν.

\vs{20}Οὐαὶ οἱ λέγοντες τὸ πονηρὸν καλὸν, καὶ τὸ καλὸν πονηρὸν· οἱ τιθέντες τὸ σκότος φῶς, καὶ τὸ φῶς σκότος· οἱ τιθέντες τὸ πικρὸν γλυκὺ, καὶ τὸ γλυκὺ πικρόν.
\vs{21}Οὐαὶ οἱ συνετοὶ ἐν ἑαυτοῖς, καὶ ἐνώπιον αὐτῶν ἐπιστήμονες.
\vs{22}Οὐαὶ οἱ ἰσχύοντες ὑμῶν, οἱ πίνοντες τὸν οἶνον, καὶ οἱ δυνάσται οἱ κεραννύντες τὸ σίκερα,
\vs{23}οἱ δικαιοῦντες τὸν ἀσεβῆ ἕνεκεν δώρων, καὶ τὸ δίκαιον τοῦ δικαίου αἴροντες.

\vs{24}Διατοῦτο ὃν τρόπον καυθήσεται καλάμη ὑπὸ ἄνθρακος πυρός, καὶ συγκαυθήσεται ὑπὸ φλογὸς ἀνειμένης, ἡ ῥίζα αὐτῶν ὡς χνοῦς ἔσται, καὶ τὸ ἄνθος αὐτῶν ὡς κονιορτὸς ἀναβήσεται· οὐ γὰρ ἠθέλησαν τὸν νόμον Κυρίου σαβαὼθ, ἀλλὰ τὸ λόγιον τοῦ ἁγίου Ἰσραὴλ παρώξυναν.
\vs{25}Καὶ ἐθυμώθη ὀργῇ Κύριος σαβαὼθ ἐπὶ τὸν λαὸν αὐτοῦ, καὶ ἐπέβαλε τὴν χεῖρα ἐπʼ αὐτοὺς, καὶ ἐπάταξεν αὐτούς· καὶ παρωξύνθη τὰ ὄρη, καὶ ἐγενήθη τὰ θνησιμαῖα αὐτῶν ὡς κοπρία ἐν μέσῳ ὁδοῦ· καὶ ἐν πᾶσι τούτοις οὐκ ἀπεστράφη ὁ θυμὸς αὐτοῦ, ἀλλὰ ἔτι ἡ χεὶρ ὑψηλή.

\vs{26}Τοιγαροῦν ἀρεῖ σύσσημον ἐν τοῖς ἔθνεσι τοῖς μακρὰν, καὶ συριεῖ αὐτοὺς ἀπʼ ἄκρου τῆς γῆς· καὶ ἰδοὺ ταχὺ κούφως ἔρχονται.
\vs{27}Οὐ πεινάσουσιν, οὐδὲ κοπιάσουσιν, οὐδὲ νυστάξουσιν, οὐδὲ κοιμηθήσονται, οὐδὲ λύσουσι τὰς ζώνας αὐτῶν ἀπο τῆς ὀσφύος αὐτῶν, οὐδὲ μὴ ῥαγῶσιν οἱ ἱμάντες τῶν ὑποδημάτων αὐτῶν·
\vs{28}Ὧν τὰ βέλη ὀξέα ἐστί, καὶ τὰ τόξα αὐτῶν ἐντεταμένα· οἱ πόδες τῶν ἵππων αὐτῶν ὡς στερεὰ πέτρα ἐλογίσθησαν· οἱ τροχοὶ τῶν ἁρμάτων αὐτῶν ὡς καταιγίς.
\vs{29}Ὀργιῶσιν ὡς λέοντες, καὶ παρέστηκαν ὡς σκύμνοι λέοντος· καὶ ἐπιλήψεται, καὶ βοήσει ὡς θηρίον, καὶ ἐκβαλεῖ, καὶ οὐκ ἔσται ὁ ῥυόμενος αὐτούς.
\vs{30}Καὶ βοήσει διʼ αὐτοὺς τῇ ἡμέρᾳ ἐκείνῃ, ὡς φωνὴ θαλάσσης κυμαινούσης· καὶ ἐμβλέψονται εἰς τὴν γῆν, καὶ ἰδοὺ σκότος σκληρὸν ἐν τῇ ἀπορίᾳ αὐτῶν.

\ch{6}
Καὶ ἐγένετο τοῦ ἐνιαυτοῦ οὗ ἀπέθανεν Ὀζίας ὁ βασιλεὺς, εἶδον τὸν Κύριον καθήμενον ἐπὶ θρόνου ὑψηλοῦ καὶ ἐπῃρμένου· καὶ πλήρης ὁ οἶκος τῆς δόξης αὐτοῦ.
\vs{2}Καὶ σεραφὶμ εἱστήκεισαν κύκλῳ αὐτοῦ, ἓξ πτέρυγες τῷ ἑνὶ, καὶ ἓξ πτέρυγες τῷ ἑνί· καὶ ταῖς μὲν δυσὶ, κατεκάλυπτον τὸ πρόσωπον· ταῖς δὲ δυσὶ κατεκάλυπτον τοὺς πόδας· καὶ ταῖς δυσὶν ἐπέταντο.
\vs{3}Καὶ ἐκέκραγεν ἕτερος πρὸς τὸν ἕτερον, καὶ ἔλεγον, ἅγιος ἅγιος ἅγιος Κύριος σαβαὼθ, πλήρης πᾶσα ἡ γῆ τῆς δόξης αὐτοῦ.

\vs{4}Καὶ ἐπῄρθη τὸ ὑπέρθυρον ἀπὸ τῆς φωνῆς ἧς ἐκέκραγον, καὶ ὁ οἶκος ἐνεπλήσθη καπνοῦ.
\vs{5}Καὶ εἶπον, ὢ τάλας ἐγώ, ὅτι κατανένυγμαι, ὅτι ἄνθρωπος ὢν, καὶ ἀκάθαρτα χείλη ἔχων, ἐν μέσῳ λαοῦ ἀκάθαρτα χείλη ἔχοντος ἐγὼ οἰκῶ, καὶ τὸν βασιλέα Κύριον σαβαὼθ εἶδον τοῖς ὀφθαλμοῖς μου.
\vs{6}Καὶ ἀπεστάλη πρὸς μὲ ἓν τῶν σεραφὶμ, καὶ ἐν τῇ χειρὶ εἶχεν ἄνθρακα, ὃν τῇ λαβίδι ἔλαβεν ἀπὸ τοῦ θυσιαστηρίου,
\vs{7}καὶ ἥψατο τοῦ στόματός μου, καὶ εἶπεν, ἰδοὺ ἥψατο τοῦτο τῶν χειλέων σου, καὶ ἀφελεῖ τὰς ἀνομίας σου, καὶ τὰς ἁμαρτίας σου περικαθαριεῖ.

\vs{8}Καὶ ἤκουσα τῆς φωνῆς Κυρίου λέγοντος, τίνα ἀποστείλω, καὶ τίς πορεύσεται πρὸς τὸν λαὸν τοῦτον; καὶ εἶπα, ἰδοὺ ἐγώ εἰμι· ἀπόστειλόν με.
\vs{9}Καὶ εἶπε, πορεύθητι, καὶ εἰπὸν τῷ λαῷ τούτῳ, ἀκοῇ ἀκούσετε, καὶ οὐ μὴ συνῆτε, καὶ βλέποντες βλέψετε, καὶ οὐ μὴ ἴδητε.
\vs{10}Ἐπαχύνθη γὰρ ἡ καρδία τοῦ λαοῦ τούτου, καὶ τοῖς ὠσὶν αὐτῶν βαρέως ἤκουσαν, καὶ τοὺς ὀφθαλμοὺς ἐκάμμυσαν· μήποτε ἴδωσι τοῖς ὀφθαλμοῖς, καὶ τοῖς ὠσὶν ἀκούσωσι, καὶ τῇ καρδίᾳ συνῶσι καὶ ἐπιστρέψωσι, καὶ ἰάσομαι αὐτούς.
\vs{11}Καὶ εἶπα, ἕως πότε Κύριε; καὶ εἶπεν, ἕως ἂν ἐρημωθῶσι πόλεις παρὰ τὸ μὴ κατοικεῖσθαι, καὶ οἶκοι παρὰ τὸ μὴ εἶναι ἀνθρώπους, καὶ ἡ γῆ καταλειφθήσεται ἔρημος.
\vs{12}Καὶ μετὰ ταῦτα μακρυνεῖ ὁ Θεὸς τοὺς ἀνθρώπους, καὶ πληθυνθήσονται οἱ ἐγκαταλειφθέντες ἐπὶ τῆς γῆς·
\vs{13}καὶ ἔτι ἐπʼ αὐτῆς ἐστι τὸ ἐπιδέκατον, καὶ πάλιν ἔσται εἰς προνομὴν ὡς τερέβινθος, καὶ ὡς βάλανος ὅταν ἐκπέσῃ ἐκ τῆς θήκης αὐτῆς.

\ch{7}
Καὶ ἐγένετο ἐν ταῖς ἡμέραις Ἄχαζ τοῦ Ἰωάθαμ τοῦ υἱοῦ Ὀζίου βασιλέως Ἰούδα, ἀνέβη Ῥασὶν βασιλεὺς Ἀραμ, καὶ Φακεὲ υἱὸς Ῥομελίου βασιλεὺς Ἰσραὴλ ἐπὶ Ἱερουσαλὴμ πολεμῆσαι αὐτὴν, καὶ οὐκ ἠδυνήθησαν πολιορκῆσαι αὐτήν.
\vs{2}Καὶ ἀνηγγέλη εἰς τὸν οἶκον Δαυὶδ, λέγων, συνεφώνησεν Ἀρὰμ πρὸς τὸν Ἐφραίμ· καὶ ἐξέστη ἡ ψυχὴ αὐτοῦ, καὶ ἡ ψυχὴ τοῦ λαοῦ αὐτοῦ, ὃν τρόπον ἐν δρυμῷ ξύλον ὑπὸ πνεύματος σαλευθῇ·
\vs{3}Καὶ εἶπε Κύριος πρὸς Ἡσαΐαν, ἔξελθε εἰς συνάντησιν Ἄχαζ σὺ, καὶ ὁ καταλειφθεὶς Ἰασοὺβ ὁ υἱός σου, πρὸς τὴν κολυμβήθραν τῆς ἄνω ὁδοῦ ἀγροῦ τοῦ κναφέως.
\vs{4}Καὶ ἐρεῖς αὐτῷ, φύλαξαι τοῦ ἡσυχάσαι, καὶ μὴ φοβοῦ, μηδὲ ἡ ψυχή σου ἀσθενείτω ἀπὸ τῶν δύο ξύλων τῶν δαλῶν τῶν καπνιζομένων τούτων· ὅταν γὰρ ὀργὴ τοῦ θυμοῦ μου γένηται, πάλιν ἰάσομαι.
\vs{5}Καὶ ὁ υἱὸς τοῦ Ἀρὰμ, καὶ ὁ υἱὸς τοῦ Ῥομελίου, ὅτι ἐβουλεύσαντο βουλὴν πονηράν·
\vs{6}Ἀναβησόμεθα εἰς τὴν Ἰουδαίαν, καὶ συλλαλήσαντες αὐτοῖς, ἀποστρέψομεν αὐτοὺς πρὸς ἡμᾶς, καὶ βασιλεύσομεν αὐτῆς τὸν υἱὸν Ταβεήλ·
\vs{7}Τάδε λέγει Κύριος σαβαὼθ, οὐ μὴ μείνῃ ἡ βουλὴ αὕτη, οὐδὲ ἔσται,
\vs{8}ἀλλʼ ἡ κεφαλὴ Ἀρὰμ, Δαμασκὸς, καὶ ἡ κεφαλὴ Δαμασκοῦ, Ῥασίμ· ἀλλʼ ἔτι ἐξήκοντα καὶ πέντε ἐτῶν ἐκλείψει ἡ βασιλεία Ἐφραίμ ἀπὸ λαοῦ,
\vs{9}καὶ ἡ κεφαλὴ Ἐφραὶμ Σομόρων, καὶ ἡ κεφαλὴ Σομορων, υἱὸς τοῦ Ῥομελίου, καὶ ἐὰν μὴ πιστεύσητε, οὐδὲ μὴ συνῆτε.

\vs{10}Καὶ προσέθετο Κύριος λαλῆσαι τῷ Ἄχαζ, λέγων,
\vs{11}αἴτησαι σεαυτῷ σημεῖον παρὰ Κυρίου Θεοῦ σου εἰς βάθος, ἢ εἰς ὕψος.
\vs{12}Καὶ εἶπεν Ἄχαζ, οὐ μὴ αἰτήσω, οὐδὲ μὴ πειράσω Κύριον.
\vs{13}Καὶ εἶπεν, ἀκούσατε δὴ οἶκος Δαυίδ· μὴ μικρὸν ὑμῖν ἀγῶνα παρέχειν ἀνθρώποις, καὶ πῶς Κυρίῳ παρέχετε ἀγῶνα;
\vs{14}Διατοῦτο δώσει Κύριος αὐτὸς ὑμῖν σημεῖον· ἰδοὺ ἡ παρθένος ἐν γαστρὶ λήψεται, καὶ τέξεται υἱὸν, καὶ καλέσεις τὸ ὄνομα αὐτοῦ Ἐμμανουήλ.
\vs{15}Βούτυρον καὶ μέλι φάγεται πρινὴ γνῶναι αὐτὸν ἢ προελέσθαι πονηρὰ, ἐκλέξασθαι τὸ ἀγαθόν·
\vs{16}Διότι πρινὴ γνῶναι τὸ παιδίον ἀγαθὸν ἢ κακὸν, ἀπειθεῖ πονηρίᾳ, ἐκλέξασθαι τὸ ἀγαθόν· καὶ καταλειφθήσεται ἡ γῆ ἣν σὺ φοβῇ, ἀπὸ προσώπου τῶν δύο βασιλέων.

\vs{17}Ἀλλὰ ἐπάξει ὁ Θεὸς ἐπὶ σὲ καὶ ἐπὶ τὸν λαόν σου καὶ ἐπὶ τὸν οἶκον τοῦ πατρός σου ἡμέρας, αἳ οὔπω ἥκασιν ἀφʼ ἧς ἡμέρας ἀφεῖλεν Ἐφραὶμ ἀπὸ Ἰούδα τὸν βασιλέα τῶν Ἀσσυρίων.
\vs{18}Καὶ ἔσται ἐν τῇ ἡμέρᾳ ἐκείνῃ συριεῖ Κύριος μυίαις, ὃ κυριεύσει μέρος ποταμοῦ Αἰγύπτου, καὶ τῇ μελίσσῃ, ἥ ἐστιν ἐν χώρᾳ Ἀσσυρίων·
\vs{19}Καὶ ἐλεύσονται πάντες ἐν ταῖς φάραγξι τῆς χώρας, καὶ ἐν ταῖς τρώγλαις τῶν πετρῶν, καὶ εἰς τὰ σπήλαια, καὶ εἰς πᾶσαν ῥαγάδα.
\vs{20}Ἐν τῇ ἡμέρᾳ ἐκείνῃ ξυρήσει Κύριος ἐν τῷ ξυρῷ τῷ μεμισθωμένῳ πέραν τοῦ ποταμοῦ βασιλέως Ἀσσυρίων τὴν κεφαλὴν, καὶ τὰς τρίχας τῶν ποδῶν, καὶ τὸν πώγωνα ἀφελεῖ.
\vs{21}Καὶ ἔσται ἐν τῇ ἡμέρᾳ ἐκείνῃ θρέψει ἄνθρωπος δάμαλιν βοῶν, καὶ δύο πρόβατα·
\vs{22}Καὶ ἔσται ἀπὸ τοῦ πλεῖστον πιεῖν γάλα, βούτυρον καὶ μέλι φάγεται πᾶς ὁ καταλειφθεὶς ἐπὶ τῆς γῆς.

\vs{23}Καὶ ἔσται ἐν τῇ ἡμέρᾳ ἐκείνῃ, πᾶς τόπος οὗ ἐὰν ὦσι χίλιαι ἄμπελοι χιλίων σίκλων, εἰς χέρσον ἔσονται, καὶ εἰς ἄκανθαν.
\vs{24}Μετὰ βέλους καὶ τοξεύματος εἰσελεύσονται ἐκεῖ· ὅτι χέρσος καὶ ἄκανθα ἔσται πᾶσα ἡ γῆ,
\vs{25}καὶ πᾶν ὄρος ἠροτριωμένον ἀροτριωθήσεται· οὐ μὴ ἐπέλθῃ ἐκεῖ φόβος· ἔσται γὰρ ἀπὸ τῆς χέρσου καὶ ἀκάνθης εἰς βόσκημα προβάτου, καὶ καταπάτημα βοός.

\ch{8}
Καὶ εἶπε Κύριος πρὸς μὲ, λάβε σεαυτῷ τόμον καινοῦ μεγαλοῦ, καὶ γράψον εἰς αὐτὸν γραφίδι ἀνθρώπου, τοῦ ὀξέως προνομὴν ποιῆσαι σκύλων·
\vs{2}Πάρεστι γάρ· καὶ μάρτυράς μοι ποίησον πιστοὺς ἀνθρώπους, τὸν Οὐρίαν καὶ Ζαχαρίαν υἱὸν Βαραχίου.
\vs{3}Καὶ προσῆλθον πρὸς τὴν προφῆτιν, καὶ ἐν γαστρὶ ἔλαβε, καὶ ἔτεκεν υἱόν· καὶ εἶπε Κύριός μοι, κάλεσον τὸ ὄνομα αὐτοῦ, Ταχέως σκύλευσον, ὀξέως προνόμευσον·
\vs{4}Διότι πρινὴ γνῶναι τὸ παιδίον καλεῖν πατέρα ἢ μητέρα, λήψεται δύναμιν Δαμασκοῦ, καὶ τὰ σκῦλα Σαμαρείας ἔναντι βασιλέως Ἀσσυρίων.

\vs{5}Καὶ προσέθετο Κύριος λαλῆσαί μοι ἔτι·
\vs{6}Διὰ τὸ μὴ βούλεσθαι τὸν λαὸν τοῦτον τὸ ὕδωρ τοῦ Σιλωὰμ τὸ πορευόμενον ἡσυχῆ, ἀλλὰ βούλεσθαι ἔχειν τὸν Ῥασσὶν καὶ τὸν υἱὸν Ῥομελίου βασιλέα ἐφʼ ὑμῶν,
\vs{7}διατοῦτο ἰδοὺ Κύριος ἀνάγει ἐφʼ ὑμᾶς τὸ ὕδωρ τοῦ ποταμοῦ, τὸ ἰσχυρὸν καὶ τὸ πολὺ, τὸν βασιλέα τῶν Ἀσσυρίων, καὶ τὴν δόξαν αὐτοῦ· καὶ ἀναβήσεται ἐπὶ πᾶσαν φάραγγα ὑμῶν, καὶ περιπατήσει ἐπὶ πᾶν τεῖχος ὑμῶν,
\vs{8}καὶ ἀφελεῖ ἀπὸ τῆς Ιουδαίας ἄνθρωπον, ὃς δυνήσεται κεφαλὴν ἆραι, ἢ δυνατὸν συντελέσασθαί τι· καὶ ἔσται ἡ παρεμβολὴ αὐτοῦ ὥστε πληρῶσαι τὸ πλάτος τῆς χώρας σου, μεθʼ ἡμῶν ὁ Θεός.

\vs{9}Γνῶτε ἔθνη καὶ ἡττᾶσθε, ἐπακούσατε ἕως ἐσχάτου τῆς γῆς· ἰσχυκότες ἡττᾶσθε· ἐὰν γὰρ πάλιν ἰσχύσητε, πάλιν ἡττηθήσεσθε.
\vs{10}Καὶ ἣν ἂν βουλεύσησθε βουλὴν, διασκεδάσει Κύριος· καὶ λόγον ὃν ἐὰν λαλήσητε, οὐ μὴ ἐμμείνῃ ἐν ὑμῖν, ὅτι μεθʼ ἡμῶν ὁ Θεός.
\vs{11}Οὕτω λέγει Κύριος, Τῇ ἰσχυρᾷ χειρὶ ἀπειθοῦσι τῇ πορεὶᾳ τῆς ὁδοῦ τοῦ λαοῦ τούτου, λέγοντες,
\vs{12}μήποτε εἴπωσι, σκληρόν· πᾶν γὰρ ὃ ἐὰν εἴπῃ ὁ λαὸς οὕτος, σκληρόν ἐστι· τὸν δὲ φόβον αὐτοῦ οὐ μὴ φοβηθῆτε οὐδὲ μὴ ταραχθῆτε.
\vs{13}Κύριον αὐτὸν ἁγιάσατε, καὶ αὐτὸς ἔσται σου φόβος.
\vs{14}κᾂν ἐπʼ αὐτῷ πεποιθὼς ᾖς, ἔσται σοι εἰς ἁγίασμα, καὶ οὐχ ὡς λίθου προσκόμματι συναντήσεσθε, οὐδὲ ὡς πέτρας πτώματι· οἱ δὲ οἶκοι Ἰακὼβ ἐν παγίδι, καὶ ἐν κοιλάσματι ἐγκαθήμενοι ἐν Ἱερουσαλήμ.
\vs{15}Διατοῦτο ἀδυνατήσουσιν ἐν αὐτοῖς πολλοὶ, καὶ πεσοῦνται καὶ συντριβήσονται, καὶ ἐγγιοῦσι, καὶ ἁλώσονται ἄνθρωποι ἐν ἀσφαλείᾳ.
\vs{16}Τότε φανεροὶ ἔσονται οἱ σφραγιζόμενοι τὸν νόμον τοῦ μὴ μαθεῖν.

\vs{17}Καὶ ἐρεῖ, μενῶ τὸν Θεὸν τὸν ἀποστρέψαντα τὸ πρόσωπον αὐτοῦ ἀπὸ τοῦ οἴκου Ἰακὼβ, καὶ πεποιθὼς ἔσομαι ἐπʼ αὐτῷ.
\vs{18}Ἰδοὺ ἐγὼ καὶ τὰ παιδία ἅ μοι ἔδωκεν ὁ Θεός· καὶ ἔσται σημεῖα καὶ τέρατα ἐν τῷ οἴκῳ Ἰσραὴλ παρὰ Κυρίου σαβαὼθ, ὃς κατοικεῖ ἐν τῷ ὄρει Σιών.

\vs{19}Καὶ ἐὰν εἴπωσι πρὸς ὑμᾶς, ζητήσατε τοὺς ἐγγαστριμύθους, καὶ τοὺς ἀπὸ τῆς γῆς φωνοῦντας, τοὺς κενολογοῦντας, οἳ ἐκ τῆς κοιλίας φωνοῦσιν· οὐκ ἔθνος πρὸς Θεὸν αὐτοῦ ἐκζητήσουσι; τί ἐκζητοῦσι περὶ τῶν ζώντων τοὺς νεκρούς;
\vs{20}Νόμον γὰρ εἰς βοήθειαν ἔδωκεν, ἵνα εἴπωσιν οὐχ ὡς τὸ ῥῆμα τοῦτο, περὶ οὗ οὐκ ἔστι δῶρα δοῦναι περὶ αὐτοῦ.

\vs{21}Καὶ ἥξει ἐφʼ ὑμᾶς σκληρὰ λιμὸς, καὶ ἔσται ὡς ἂν πεινάσητε, λυπηθήσεσθε. καὶ κακῶς ἐρεῖτε τὸν ἄρχοντα καὶ τὰ πάτρια· καὶ ἀναβλέψονται εἰς τὸν οὐρανὸν ἄνω,
\vs{22}καὶ εἰς τὴν γῆν κάτω ἐμβλέψονται· καὶ ἰδοὺ ἀπορία στενὴ, καὶ σκότος, θλίψις, καὶ στενοχωρία, καὶ σκότος ὥστε μὴ βλέπειν· καὶ οὐκ ἀπορηθήσεται ὁ ἐν στενοχωρίᾳ ὢν ἕως καιροῦ.

\vs{23}Τοῦτο πρῶτον πίε· ταχὺ ποίει χώρα Ζαβυλὼν, ἡ γῆ Νεφθαλεὶμ, καὶ οἱ λοιποὶ οἱ τὴν παραλίαν, καὶ πέραν τοῦ Ἰορδάνου Γαλιλαία τῶν ἐθνῶν.

\ch{9}
Ὁ λαὸς ὁ πορευόμενος ἐν σκότει, ἴδετε φῶς μέγα· οἱ κατοικοῦντες ἐν χώρᾳ σκιᾷ θανάτου, φῶς λάμψει ἐφʼ ὑμᾶς.
\vs{2}Τὸ πλεῖστον τοῦ λαοῦ, ὃ κατήγαγες ἐν εὐφροσύνῃ σου· καὶ εὐφρανθήσονται ἐνώπιόν σου, ὡς οἱ εὐφραινόμενοι ἐν ἀμήτῳ, καὶ ὃν τρόπον οἱ διαιρούμενοι σκῦλα.
\vs{3}Διότι ἀφῄρηται ὁ ζυγὸς ὁ ἐπʼ αὐτῶν κείμενος, καὶ ἡ ῥάβδος ἡ ἐπὶ τοῦ τραχήλου αὐτῶν· τὴν γὰρ ῥάβδον τῶν ἀπαιτούντων διεσκέδασεν, ὡς τῇ ἡμέρᾳ τῇ ἐπὶ Μαδιάμ.
\vs{4}Ὅτι πᾶσαν στολὴν ἐπισυνήγμένην δόλῳ, καὶ ἱμάτιον μετὰ καταλλαγῆς ἀποτίσουσι· καὶ θελήσουσιν, εἰ ἐγένοντο πυρίκαυστοι.

\vs{5}Ὅτι παιδίον ἐγεννήθη ἡμῖν, υἱὸς καὶ ἐδόθη ἡμῖν, οὗ ἡ ἀρχὴ ἐγενήθη ἐπὶ τοῦ ὤμου αὐτοῦ, καὶ καλεῖται τὸ ὄνομα αὐτοῦ, Μεγάλης βουλῆς ἄγγελος· ἄξω γὰρ εἰρήνην ἐπὶ τοὺς ἄρχοντας, καὶ ὑγίειαν αὐτῷ.
\vs{6}Μεγάλη ἡ ἀρχὴ αὐτοῦ, καὶ τῆς εἰρήνης αὐτοῦ οὐκ ἔστιν ὅριον· ἐπὶ τὸν θρόνον Δαυεὶδ, καὶ τὴν βασιλείαν αὐτοῦ, κατορθῶσαι αὐτὴν, καὶ ἀντιλαβέσθαι ἐν κρίματι καὶ ἐν δικαιοσύνῃ, ἀπὸ τοῦ νῦν καὶ εἰς τὸν αἰῶνα· ὁ ζῆλος Κυρίου σαβαὼθ ποιήσει ταῦτα.

\vs{7}Θάνατον ἀπέστειλε Κύριος ἐπὶ Ἰακώβ, καὶ ἦλθεν ἐπὶ Ἰσραήλ.
\vs{8}Καὶ γνώσονται πᾶς ὁ λαὸς τοῦ Ἐφραὶμ, καὶ οἱ καθήμενοι ἐν Σαμαρείᾳ, ἐφʼ ὕβρει καὶ ὑψηλῇ καρδίᾳ λέγοντες,
\vs{9}πλίνθοι πεπτώκασιν, ἀλλὰ δεῦτε λαξεύσωμεν λίθους, καὶ κόψωμεν συκαμίνους καὶ κέδρους, καὶ οἰκοδομήσωμεν ἑαυτοῖς πύργον.
\vs{10}Καὶ ῥάξει ὁ Θεὸς τοὺς ἐπανισταμένους ἐπὶ ὄρος Σιὼν ἐπὶ αὐτὸν, καὶ τοὺς ἐχθροὺς διασκεδάσει·
\vs{11}Συρίαν ἀφʼ ἡλίου ἀνατολῶν, καὶ τοὺς Ἕλληνας ἀφʼ ἡλίου δυσμῶν, τοὺς κατεσθίοντας τὸν Ἰσραὴλ ὅλῳ τῷ στόματι· ἐπὶ πᾶσι τούτοις οὐκ ἀπεστράφη ὁ θυμὸς, ἀλλʼ ἔτι ἡ χεὶρ ὑψηλή.

\vs{12}Καὶ ὁ λαὸς οὐκ ἀπεστράφη, ἕως ἐπλήγη, καὶ τὸν κύριον οὐκ ἐζήτησαν.
\vs{13}Καὶ ἀφεῖλε Κύριος ἀπὸ Ἰσραὴλ κεφαλὴν καὶ οὐρὰν, μέγαν καὶ μικρὸν, ἐν μιᾷ ἡμέρᾳ, πρεσβύτην, καὶ τοὺς τὰ πρόσωπα θαυμάζοντας, αὕτη ἡ ἀρχή·
\vs{14}καὶ προφήτην διδάσκοντα ἄνομα, οὗτος ἡ οὐρά.
\vs{15}Καὶ ἔσονται οἱ μακαρίζοντες τὸν λαὸν τοῦτον πλανῶντες, καὶ πλανῶσιν, ὅπως καταπίνωσιν αὐτούς.
\vs{16}Διατοῦτο ἐπὶ τοὺς νεανίσκους αὐτῶν οὐκ εὐφρανθήσεται ὁ Κύριος, καὶ τοὺς ὀρφανοὺς αὐτῶν καὶ τὰς χήρας αὐτῶν οὐκ ἐλεήσει· ὅτι πάντε ἄνομοι καὶ πονηροί, καὶ πᾶν στόμα λαλεῖ ἄδικα· ἐπὶ πᾶσι τούτοις οὐκ ἀπεστράφη ὁ θυμός, ἀλλʼ ἔτι ἡ χεὶρ ὑψηλή.

\vs{17}Καὶ καυθήσεται ὡς πῦρ ἡ ἀνομία, καὶ ὡς ἄγρωστις ξηρὰ βρωθήσεται ὑπὸ πυρός· καὶ καυθήσεται ἐν τοῖς δάσεσι τοῦ δρυμοῦ, καὶ συγκαταφάγεται τὰ κύκλῳ τῶν βουνῶν πάντα.
\vs{18}Διὰ θυμὸν ὀργῆς Κυρίου συγκέκαυται ἡ γῆ ὅλη· καὶ ἔσται ὁ λαὸς ὡς κατακεκαυμένος ὑπὸ πυρὸς· ἄνθρωπος τὸν ἀδελφὸν αὐτοῦ οὐκ ἐλεήσει.
\vs{19}Ἀλλὰ ἐκκλινεῖ εἰς τὰ δεξιά, ὅτι πεινάσει, καὶ φάγεται ἐκ τῶν ἀριστερῶν, καὶ οὐ μὴ ἐμπλησθῇ ἄνθρωπος ἔσθων τὰς σάρκας τοῦ βραχίονος αὐτοῦ.
\vs{20}Φάγεται γὰρ Μανασσῆς τοῦ Ἐφραὶμ, καὶ Ἐφραὶμ τοῦ Μανασσῆ, ὅτι ἅμα πολιορκήσουσι τὸν Ἰούδαν. ἐπὶ τούτοις πᾶσιν οὐκ ἀπεστράφη ὁ θυμός, ἀλλʼ ἔτι ἡ χεὶρ ὑψηλή.

\ch{10}
Οὐαὶ τοῖς γράφουσι πονηρίαν, γράφοντες γὰρ, πονηρίαν γράφουσιν·
\vs{2}Ἐκκλίνοντες κρίσιν πτωχῶν, ἁρπάζοντες κρίμα πενήτων τοῦ λαοῦ μου, ὥστε εἶναι αὐτοῖς χήραν εἰς διαρπαγὴν, καὶ ὀρφανὸν εἰς προνομήν.
\vs{3}Καὶ τί ποιήσουσιν ἐν τῇ ἡμέρᾳ τῆς ἐπισκοπῆς; ἡ γὰρ θλίψις ὑμῖν πόῤῥωθεν ἥξει· καὶ πρὸς τίνα καταφεύξεσθε τοῦ βοηθηθῆναι; καὶ ποῦ καταλείψετε τὴν δόξαν ὑμῶν, τοῦ μὴ ἐμπεσεῖν εἰς ἀπαγωγήν;

\vs{4}Ἐπὶ πᾶσι τούτοις οὐκ ἀπεστράφη ἡ ὀργή, ἀλλʼ ἔτι ἡ χεὶρ ὑψηλή.

\vs{5}Οὐαὶ Ἀσσυρίοις, ἡ ῥάβδος τοῦ θυμοῦ μου, καὶ ὀργή ἐστιν ἐν ταῖς χερσὶν αὐτῶν.
\vs{6}Τὴν ὀργήν μου εἰς ἔθνος ἄνομον ἀποστελῶ, καὶ τῷ ἐμῷ λαῷ συντάξω ποιῆσαι σκῦλα καὶ προνομήν, καὶ καταπατεῖν τὰς πόλεις, καὶ θεῖναι αὐτὰς εἰς κονιορτόν.
\vs{7}Αὐτὸς δὲ οὐχ οὕτως ἐνεθυμήθη, καὶ τῇ ψυχῇ οὐχ οὕτως λελόγισται, ἀλλὰ ἀπαλλάξει ὁ νοῦς αὐτοῦ, καὶ τοῦ ἔθνη ἐξολοθρεῦσαι οὐκ ὀλίγα.
\vs{8}Καὶ ἐὰν εἴπωσιν αὐτῷ, σὺ μόνος εἶ ἄρχων·
\vs{9}Καὶ ἐρεῖ, οὐκ ἔλαβον τὴν χώραν τὴν ἐπάνω Βαβυλῶνος καὶ Χαλάνης, οὗ ὁ πύργος ᾠκοδομήθη, καὶ ἔλαβον Ἀραβίαν καὶ Δαμασκὸν καὶ Σαμάρειαν;
\vs{10}Ὃν τρόπον ταύτας ἔλαβον, καὶ πάσας τὰς ἀρχὰς λήψομαι· ὀλολύξατε τὰ γλυπτὰ ἐν Ἱερουσαλὴμ, καὶ ἐν Σαμαρείᾳ.
\vs{11}Ὃν τρόπον γὰρ ἐποίησα Σαμαρείᾳ, καὶ τοῖς χειροποιήτοις αὐτῆς, οὕτω ποιήσω καὶ Ἱερουσαλὴμ, καὶ τοῖς εἰδώλοις αὐτῆς.
\vs{12}Καὶ ἔσται ὅταν συντελέσῃ Κύριος πάντα ποιῶν ἐν τῷ ὄρει Σιὼν καὶ Ἱερουσαλήμ, ἐπισκέψομαι ἐπὶ τὸν νοῦν τὸν μέγαν ἐπὶ τὸν ἄρχοντα τῶν Ἀσσυρίων, καὶ ἐπὶ τὸ ὕψος τῆς δόξης τῶν ὀφθαλμῶν αὐτοῦ.
\vs{13}Εἶπε γὰρ, ἐν τῇ ἰσχύϊ ποιήσω, καὶ ἐν τῇ σοφίᾳ τῆς συνέσεως ἀφελῶ ὅρια ἐθνῶν, καὶ τὴν ἰσχὺν αὐτῶν προνομεύσω·
\vs{14}Καὶ σείσω πόλεις κατοικουμένας, καὶ τὴν οἰκουμένην ὅλην καταλήψομαι τῇ χειρὶ ὡς νοσσιὰν, καὶ ὡς καταλελειμμένα ὠὰ ἀρῶ· καὶ οὐκ ἔστιν ὃς διαφεύξεταί με, ἢ ἀντείπῃ μοι.
\vs{15}Μὴ δοξασθήσεται ἀξίνη ἄνευ τοῦ κόπτοντος ἐν αὐτῇ; ἢ ὑψωθήσεται πριὼν ἄνευ τοῦ ἔλκοντος αὐτόν; ὡς ἄν τις ἄρῃ ῥάβδον ἢ ξύλον· καὶ οὐχ οὕτως,
\vs{16}ἀλλὰ ἀποστελεῖ Κύριος σαβαὼθ εἰς τὴν σὴν τιμὴν ἀτιμίαν, καὶ εἰς τὴν σὴν δόξαν πῦρ καιόμενον καυθήσεται.
\vs{17}Καὶ ἔσται τὸ φῶς τοῦ Ἰσρὴλ εἰς πῦρ, καὶ ἁγιάσει αὐτὸν ἐν πυρὶ καιομένῳ, καὶ φάγεται ὡσεὶ χόρτον τὴν ὕλην·
\vs{18}τῇ ἡμέρᾳ ἐκείνῃ ἀποσβεσθήσεται τὰ ὄρη, καὶ οἱ βουνοὶ, καὶ οἱ δρυμοὶ, καὶ καταφάγεται ἀπὸ ψυχὴς ἕως σαρκῶν· καὶ ἔσται ὁ φεύγων, ὡς ὁ φεύγων ἀπὸ φλογὸς καιομένης.
\vs{19}Καὶ οἱ καταλειφθέντες ἀπʼ αὐτῶν ἀριθμὸς ἔσονται, καὶ παιδίον γράψει αὐτούς.

\vs{20}Καὶ ἔσται ἐν τῇ ἡμέρᾳ ἐκείνῃ, οὐκέτι προστεθήσεται τὸ καταλειφθὲν Ἰσραὴλ, καὶ οἱ σωθέντες τοῦ Ἰακὼβ οὐκέτι μὴ πεποιθότες ὦσιν ἐπὶ τοὺς ἀδικήσαντας αὐτοὺς, ἀλλὰ ἔσονται πεποιθότες ἐπὶ τὸν Θεὸν τὸν ἅγιον τοῦ Ἰσραὴλ τῇ ἀληθείᾳ.
\vs{21}Καὶ ἔσται τὸ καταλειφθὲν τοῦ Ἰακὼβ ἐπὶ Θεὸν ἰσχύοντα.
\vs{22}Καὶ ἐὰν γένηται ὁ λαὸς Ἰσραὴλ ὡς ἡ ἄμμος τῆς θαλάσσης, τὸ κατάλειμμα αὐτῶν σωθήσεται.
\vs{23}Λόγον συντελῶν καὶ συντέμνων ἐν δικαιοσύνῃ, ὅτι λόγον συντετμημένον Κύριος ποιήσει ἐν τῇ οἰκουμένῃ ὅλῃ.

\vs{24}Διατοῦτο τάδε λέγει Κύριος σαβαὼθ, μὴ φοβοῦ ὁ λαός μου, οἱ κατοικοῦντες ἐν Σιὼν, ἀπὸ Ἀσσυρίων, ὅτι ἐν ῥάβδῳ πατάξει σε· πληγὴν γὰρ ἐπάγω ἐπὶ σὲ, τοῦ ἰδεῖν ὁδὸν Αἰγύπτου.
\vs{25}Ἔτι γὰρ μικρὸν, καὶ παύσεται ἡ ὀργὴ, ὁ δὲ θυμός μου ἐπὶ τὴν βουλὴν αὐτῶν.
\vs{26}Καὶ ἐγερεῖ ὁ Θεὸς ἐπʼ αὐτοὺς, κατὰ τὴν πληγὴν Μαδιὰμ ἐν τόπῳ θλίψεως· καὶ ὁ θυμὸς αὐτοῦ τῇ ὁδῷ τῇ κατὰ θάλασσαν, εἰς τὴν ὁδὸν τὴν κατʼ Αἴγυπτον.
\vs{27}Καὶ ἔσται ἐν τῇ ἡμέρᾳ ἐκείνῃ, ἀφαιρεθήσεται ὁ ζυγὸς αὐτοῦ ἀπὸ τοῦ ὤμου σου, καὶ ὁ φόβος αὐτοῦ ἀπό σου, καὶ καταφθαρήσεται ὁ ζυγὸς ἀπὸ τῶν ὤμων ὑμῶν.

\vs{28}Ἥξει γὰρ εἰς τὴν πόλιν Ἀγγαί, καὶ παρελεύσεται εἰς Μαγεδὼ, καὶ ἐν Μαχμὰς θήσει τὰ σκεύη αὐτοῦ.
\vs{29}Καὶ παρελεύσεται φάραγγα, καὶ ἥξει εἰς Ἀγγαί, φόβος λήψεται Ῥαμᾶ, πόλιν Σαοὺλ, φεύξεται
\vs{30}ἡ θυγάτηρ Γαλλεὶμ, ἐπακούσεται Λαϊσα, ἐπακούσεται ἐν Ἀναθώθ.
\vs{31}Καὶ ἐξέστη Μαδεβηνὰ, καὶ οἱ κατοικοῦντες Γιββείρ.

\vs{32}Παρακαλεῖτε σήμερον ἐν ὁδῷ τοῦ μεῖναι, τῇ χειρὶ παρακαλεῖτε τὸ ὄρος τὴν θυγατέρα Σιὼν, καὶ οἱ βουνοὶ οἱ ἐν Ἱερουσαλήμ.

\vs{33}Ἰδοὺ ὁ δεσπότης Κύριος σαβαὼθ συνταράσσει τοὺς ἐνδόξους μετὰ ἰσχύος, καὶ οἱ ὑψηλοὶ τῇ ὕβρει συντριβήσονται, καὶ οἱ ὑψηλοὶ ταπεινωθήσονται,
\vs{34}καὶ πεσοῦνται ὑψηλοὶ μαχαίρᾳ, ὁ δὲ Λίβανος σὺν τοῖς ὑψηλοῖς πεσεῖται.

\ch{11}
Καὶ ἐξελεύσεται ῥάβδος ἐκ τῆς ῥίζης Ἰεσσαί, καὶ ἄνθος ἐκ τῆς ῥίζης ἀναβήσεται,
\vs{2}καὶ ἀναπαύσεται ἐπʼ αὐτὸν πνεῦμα τοῦ Θεοῦ, πνεῦμα σοφίας καὶ συνέσεως, πνεῦμα βουλῆς καὶ ἰσχύος, πνεύμα γνώσεως καὶ εὐσεβείας
\vs{3}ἐμπλήσει αὐτὸν, πνεῦμα φόβου Θεοῦ· οὐ κατὰ τὴν δόξαν κρινεῖ, οὐδὲ κατὰ τὴν λαλιὰν ἐλέγξει,
\vs{4}ἀλλὰ κρινεῖ ταπεινῷ κρίσιν, καὶ ἐλέγξει τοὺς ταπεινοὺς τῆς γῆς, καὶ πατάξει γῆν τῷ λόγῳ τοῦ στόματος αὐτοῦ, καὶ ἐν πνεύματι διὰ χειλέων ἀνελεῖ ἀσεβῆ.
\vs{5}Καὶ ἔσται δικαιοσύνῃ ἐζωσμένος τὴν ὀσφὺν αὐτοῦ, καὶ ἀληθείᾳ εἱλημένος τὰς πλευράς.

\vs{6}Καὶ συμβοσκηθήσεται λύκος μετὰ ἀρνός, καὶ πάρδαλις συναναπαύσεται ἐρίφῳ, καὶ μοσχάριον καὶ ταῦρος καὶ λέων ἅμα βοσκηθήσονται, καὶ παιδίον μικρὸν ἄξει αὐτούς.
\vs{7}Καὶ βοῦς καὶ ἄρκος ἅμα βοσκηθήσονται, καὶ ἅμα τὰ παιδία αὐτῶν ἔσονται· καὶ λέων ὡς βοῦς φάγεται ἄχυρα.
\vs{8}Καὶ παιδίον νήπιον ἐπὶ τρωγλῶν ἀσπίδων, καὶ ἐπὶ κοίτην ἐκγόνων ἀσπίδων τὴν χεῖρα ἐπιβαλεῖ.
\vs{9}καὶ οὐ μὴ κακοποιήσουσιν, οὐδὲ μὴ δύνωνται ἀπολέσαι οὐδένα ἐπὶ τὸ ὄρος τὸ ἅγιόν μου· ὅτι ἐνεπλήσθη ἡ σύμπασα τοῦ γνῶναι τὸν Κύριον, ὡς ὕδωρ πολὺ κατακαλύψαι θαλάσσας.
\vs{10}Καὶ ἔσται ἐν τῇ ἡμέρᾳ ἐκείνῃ ἡ ῥίζα τοῦ Ἰεσσαὶ, καὶ ὁ ἀνιστάμενος ἄρχειν ἐθνῶν· ἐπʼ αὐτῷ ἔθνη ἐλπιοῦσι, καὶ ἔσται ἡ ἀνάπαυσις αὐτοῦ τιμή.
\vs{11}Καὶ ἔσται τῇ ἡμέρᾳ ἐκείνῃ, προσθήσει ὁ Κύριος τοῦ δεῖξαι τὴν χεῖρα αὐτοῦ, τοῦ ζηλῶσαι τὸ καταλειφθὲν ὑπόλοιπον τοῦ λαοῦ, ὃ ἂν καταλειφθῇ ὑπὸ τῶν Ἀσσυρίων, καὶ ἀπὸ Αἰγύπτου, καὶ ἀπὸ Βαβυλωνίας, καὶ ἀπὸ Αἰθιοπίας, καὶ ἀπὸ Ἐλαμιτῶν, καὶ ἀπὸ ἡλίου ἀνατολῶν, καὶ ἐξ Ἀραβίας.
\vs{12}Καὶ ἀρεῖ σημεῖον εἰς τὰ ἔθνη, καὶ συνάξει τοὺς ἀπολομένους Ἰσραήλ, καὶ τοὺς διεσπαρμένους Ἰούδα συνάξει ἐκ τῶν τεσσάρων πτερύγων τῆς γῆς.
\vs{13}Καὶ ἀφαιρεθήσεται ὁ ζῆλος Ἐφράιμ, καὶ οἱ ἐχθροὶ Ἰούδα ἀπολοῦνται. Ἐφραὶμ οὐ ζηλώσει Ἰούδαν, καὶ Ἰούδας οὐ θλίψει Ἐφραὶμ.
\vs{14}Καὶ πετασθήσονται ἐν πλοίοις ἀλλοφύλων· θάλασσαν ἅμα προνομεύσουσι, καὶ τοὺς ἀφʼ ἡλίου ἀνατολῶν, καὶ Ἰδουμαίαν, καὶ ἐπὶ Μωὰβ πρῶτον τὰς χεῖρας ἐπιβαλοῦσιν· οἱ δὲ υἱοὶ Ἀμμὼν πρῶτοι ὑπακούσονται.

\vs{15}Καὶ ἐρημώσει Κύριος τὴν θάλασσαν Αἰγύπτου, καὶ ἐπιβαλεῖ τὴν χεῖρα αὐτοῦ ἐπὶ τὸν ποταμὸν πνεύματι βιαίῳ· καὶ πατάξει ἑπτὰ φάραγγας, ὥστε διαπορεύεσθαι αὐτὸν ἐν ὑποδήμασι.
\vs{16}Καὶ ἔσται δίοδος τῷ καταλειφθέντι μου λαῷ ἐν Αἰγύπτῳ· καὶ ἔσται τῷ Ἰσραὴλ, ὡς ἡ ἡμέρα ὅτε ἐξῆλθεν ἐκ γῆς Αἰγύπτου.

\ch{12}
Καὶ ἐρεῖς ἐν τῇ ἡμέρᾳ ἐκείνῃ, Εὐλογῶ σε, Κύριε, διότι ὠργίσθης μοι, καὶ ἀπέστρεψας τὸν θυμόν σου, καὶ ἠλέησάς με.
\vs{2}Ἰδοὺ ὁ Θεός μου σωτήρ μου, πεποιθὼς ἔσομαι ἐπʼ αὐτῷ, καὶ οὐ φοβηθήσομαι· διότι ἡ δόξα μου καὶ ἡ αἴνεσίς μου Κύριος, καὶ ἐγένετό μοι εἰς σωτηρίαν.
\vs{3}Καὶ ἀντλήσατε ὕδωρ μετʼ εὐφροσύνης ἐκ τῶν πηγῶν τοῦ σωτηρίου.
\vs{4}Καὶ ἐρεῖς ἐν τῇ ἡμέρᾳ ἐκείνῃ, ὑμνεῖτε Κύριον, βοᾶτε τὸ ὄνομα αὐτοῦ, ἀναγγείλατε ἐν τοῖς ἔθνεσι τὰ ἔνδοξα αὐτοῦ· μιμνήσκεσθε, ὅτι ὑψώθη τὸ ὄνομα αὐτοῦ.
\vs{5}Ὑμνήσατε τὸ ὄνομα Κυρίου, ὅτι ὑψηλὰ ἐποίησεν· ἀναγγείλατε ταῦτα ἐν πάσῃ τῇ γῇ.
\vs{6}Ἀγαλλιᾶσθε, καὶ εὐφραίνεσθε οἱ κατοικοῦντες Σιὼν, ὅτι ὑψώθη ὁ ἅγιος τοῦ Ἰσραὴλ ἐν μέσῳ αὐτῆς.

\ch{13}
ὍΡΑΣΙΣ ἫΝ ΕΙΔΕΝ ἩΣΑΙΑΣ ΥΙΟΣ ἈΜΩΣ ΚΑΤΑ ΒΑΒΥΛΩΝΟΣ.

\vs{2}Ἐπʼ ὄρους πεδινοῦ ἄρατε σημεῖον, ὑψώσατε τὴν φωνὴν αὐτοῖς, παρακαλεῖτε τῇ χειρί, ἀνοίξατε οἱ ἄρχοντες.
\vs{3}Ἐγὼ συντάσσω, καὶ ἐγώ αὐτούς· γίγαντες ἔρχονται πληρῶσαι τὸν θυμόν μου χαίροντες ἅμα καὶ ὑβρίζοντες.
\vs{4}Φωνὴ ἐθνῶν πολλῶν ἐπὶ τῶν ὀρέων, ὁμοία ἐθνῶν πολλῶν, φωνὴ βασιλέων καὶ ἐθνῶν συνηγμένων· Κύριος σαβαὼθ ἐντέταλται ἔθνει ὁπλομάχῳ,
\vs{5}ἔρχεσθαι ἐκ γῆς πόῤῥωθεν ἀπʼ ἄκρου θεμελίου τοῦ οὐρανοῦ, Κύριος καὶ οἱ ὁπλομάχοι αὐτοῦ, καταφθεῖραι πᾶσαν τὴν οἰκουμένην.

\vs{6}Ὀλολύζετε, ἐγγὺς γὰρ ἡμέρα Κυρίου, καὶ συντριβὴ παρὰ τοῦ Θεοῦ ἥξει.
\vs{7}Διάτοῦτο πᾶσα χεὶρ ἐκλυθήσεται, καὶ πᾶσα ψυχὴ ἀνθρώπου δειλιάσει.
\vs{8}Ταραχθήσονται οἱ πρέσβεις, καὶ ὠδῖνες αὐτοὺς ἕξουσιν, ὡς γυναικὸς τικτούσης· καὶ συμφοράσουσιν ἕτερος πρὸς τὸν ἕτερον, καὶ ἐκστήσονται, καὶ τὸ πρόσωπον αὐτῶν ὡς φλὸξ μεταβαλοῦσιν.
\vs{9}Ἰδοὺ γὰρ ἡμέρα Κυρίου ἔρχεται ἀνίατος, θυμοῦ καὶ ὀργῆς, θεῖναι τὴν οἰκουμένην ἔρημον, καὶ τοὺς ἁμαρτωλοὺς ἀπολέσαι ἐξ αὐτῆς.
\vs{10}Οἱ γὰρ ἀστέρες τοῦ οὐρανοῦ καὶ ὁ Ὠρίων καὶ πᾶς ὁ κόσμος τοῦ οὐρανοῦ, τὸ φῶς οὐ δώσουσι· καὶ σκοτισθήσεται τοῦ ἡλίου ἀνατέλλοντος, καὶ ἡ σελήνη οὐ δώσει τὸ φῶς αὐτῆς.
\vs{11}Καὶ ἐντελοῦμαι τῇ οἰκουμένῃ ὅλῃ κακὰ, καὶ τοῖς ἀσεβέσι τὰς ἁμαρτίας αὐτῶν· καὶ ἀπολῶ ὕβριν ἀνόμων, καὶ ὕβριν ὑπερηφάνων ταπεινώσω.
\vs{12}Καὶ ἔσονται οἱ καταλελειμμένοι ἔντιμοι μᾶλλον ἢ τὸ χρυσίον τὸ ἄπυρον· καὶ ἄνθρωπος μᾶλλον ἔντιμος ἔσται ἢ ὁ λίθος ὁ ἐν Σουφείρ.
\vs{13}Ὁ γὰρ οὐρανὸς θυμωθήσεται, καὶ ἡ γῆ σεισθήσεται ἐκ τῶν θεμελίων αὐτῆς, διὰ θυμὸν ὀργῆς Κυρίου σαβαὼθ, ἐν τῇ ἡμέρᾳ ᾗ ἂν ἐπέλθῃ ὁ θυμὸς αὐτοῦ.
\vs{14}Καὶ ἔσονται οἱ καταλελειμμένοι ὡς δορκάδιον φεῦγον, καὶ ὡς πρόβατον πλανώμενον, καὶ οὐκ ἔσται ὁ συνάγων, ὥστε ἄνθρωπον εἰς τὸν λαὸν αὐτοῦ ἀποστραφῆναι, καὶ ἄνθρωπος εἰς τὴν χώραν ἑαυτοῦ διώξεται.
\vs{15}Ὃς γὰρ ἂν ἁλῷ, ἡττηθήσεται, καὶ οἵτινες συνηγμένοι εἰσὶ, πεσοῦνται μαχαίρᾳ.
\vs{16}Καὶ τὰ τέκνα αὐτῶν ῥάξουσιν ἐνώπιον αὐτῶν, καὶ τὰς οἰκίας αὐτῶν προνομεύσουσιν, καὶ τὰς γυναῖκας αὐτῶν ἕξουσιν.

\vs{17}Ἰδοὺ ἐπεγείρω ὑμῖν τοὺς Μήδους, οἳ ἀργύριον οὐ λογίζονται, οὐδὲ χρυσίου χρείαν ἔχουσι.
\vs{18}Τοξεύματα νεανίσκων συντρίψουσι, καὶ τὰ τέκνα ὑμῶν οὐ μὴ ἐλεήσωσιν, οὐδὲ ἐπὶ τοῖς τέκνοις σου φείσονται οἱ ὀφθαλμοὶ αὐτῶν.
\vs{19}Καὶ ἔσται Βαβυλὼν ἣ καλεῖται ἔνδοξος ἀπὸ βασιλέως Χαλδαίων, ὃν τρόπον κατέστρεψεν ὁ Θεὸς Σόδομα καὶ Γόμοῤῥα·
\vs{20}Οὐ κατοικηθήσεται εἰς τὸν αἰῶνα χρόνον, οὐδὲ μὴ εἰσέλθωσιν εἰς αὐτὴν διὰ πολλῶν γενεῶν, οὐδὲ μὴ διέλθωσιν αὐτὴν Ἄραβες, οὐδὲ ποιμένες οὐ μὴ ἀναπαύσονται ἐν αὐτῇ.
\vs{21}Καὶ ἀναπαύσονται ἐκεῖ θηρία, καὶ ἐμπλησθήσονται αἱ οἰκίαι ἤχου· καὶ ἀναπαύσονται ἐκεῖ σειρῆνες, καὶ δαιμόνια ἐκεῖ ὀρχήσονται,
\vs{22}καὶ ὀνοκένταυροι ἐκεῖ κατοικήσουσι, καὶ νοσσοποιήσουσιν ἐχῖνοι ἐν τοῖς οἴκοις αὐτῶν. Ταχὺ ἔρχεται καὶ οὐ χρονιεῖ.

\ch{14}
Καὶ ἐλεήσει Κύριος τὸν Ἰακὼβ, καὶ ἐκλέξεται ἔτι τὸν Ἰσραὴλ, καὶ ἀναπαύσονται ἐπὶ τῆς γῆς αὐτῶν, καὶ ὁ γειώρας προστεθήσεται πρὸς αὐτοὺς, καὶ προστεθήσεται πρὸς τὸν οἶκον Ἰακώβ·
\vs{2}Καὶ λήψονται αὐτοὺς ἔθνη, καὶ εἰσάξουσιν εἰς τὸν τόπον αὐτῶν, καὶ κατακληρονομήσουσι, καὶ πληθυνθήσονται ἐπὶ τῆς γῆς εἰς δούλους καὶ δούλας· καὶ ἔσονται αἰχμάλωτοι οἱ αἰχμαλωτεύσαντες αὐτοὺς, καὶ κυριευθήσονται οἱ κυριεύσαντες αὐτῶν.

\vs{3}Καὶ ἔσται ἐν τῇ ἡμέρᾳ ἐκείνῃ, ἀναπαύσει σε Κύριος ἀπὸ τῆς ὀδύνης καὶ τοῦ θυμοῦ σου, τῆς δουλείας σου τῆς σκληρᾶς, ἧς ἐδούλευσας αὐτοῖς.
\vs{4}Καὶ λήψῃ τὸν θρῆνον τοῦτον ἐπὶ τὸν βασιλέα Βαβυλῶνος,

Πῶς ἀναπέπαυται ὁ ἀπαιτῶν, καὶ ἀναπέπαυται ὁ ἐπισπουδαστής;
\vs{5}Συνέτριψε Κύριος τὸν ζυγὸν τῶν ἁμαρτωλῶν, τὸν ζυγὸν τῶν ἀρχόντων.
\vs{6}Πατάξας ἔθνος θυμῷ, πληγῇ ἀνιάτῳ, παίων ἔθνος πληγὴν θυμοῦ, ἣ οὐκ ἐφείσατο, ἀνεπαύσατο πεποιθώς.
\vs{7}Πᾶσα ἡ γῆ βοᾷ μετʼ εὐφροσύνης,
\vs{8}καὶ τὰ ξύλα τοῦ Λιβάνου εὐφράνθησαν ἐπὶ σοὶ, καὶ ἡ κέδρος τοῦ Λιβάνου, ἀφʼ οὗ σὺ κεκοίμησαι, οὐκ ἀνέβη ὁ κόπτων ἡμᾶς.
\vs{9}Ὁ ᾅδης κάτωθεν ἐπικράνθη συναντήσας σοι· συνηγέρθησάν σοι πάντες οἱ γίγαντες οἱ ἄρξαντες τῆς γῆς, οἱ ἐγείραντες ἐκ τῶν θρόνων αὐτῶν πάντας βασιλεῖς ἐθνῶ.
\vs{10}Πάντες ἀποκριθήσονται, καὶ ἐροῦσί σοι, Καὶ σὺ ἑάλως, ὥσπερ καὶ ἡμεῖς· ἐν ἡμῖν δὲ κατελογίσθης.
\vs{11}Κατέβη εἰς ᾅδου ἡ δόξα σου, ἡ πολλὴ εὐφροσύνη σου· ὑποκάτω σου στρώσουσι σῆψιν, καὶ τὸ κατακάλυμμά σου σκώληξ.
\vs{12}Πῶς ἐξέπεσεν ἐκ τοῦ οὐρανοῦ ὁ Ἑωσφόρος ὁ πρωὶ ἀνατέλλων; συνετρίβη εἰς τὴν γῆν ὁ ἀποστέλλων πρὸς πάντα τὰ ἔθνη.
\vs{13}Σὺ δὲ εἶπας ἐν τῇ διανοίᾳ σου, εἰς τὸν οὐρανὸν ἀναβήσομαι, ἐπάνω τῶν ἀστέρων τοῦ οὐρανοῦ θήσω τὸν θρόνον μου, καθιῶ ἐν ὄρει ὑψηλῷ, ἐπὶ τὰ ὄρη τὰ ὑψηλὰ τὰ πρὸς Βοῤῥᾶν,
\vs{14}ἀναβήσομαι ἐπάνω τῶν νεφῶν, ἔσομαι ὅμοιος τῷ ὑψίστῳ.
\vs{15}Νῦν δὲ εἰς ᾅδην καταβήσῃ, καὶ εἰς τὰ θεμέλια τῆς γῆς.
\vs{16}Οἱ ἰδόντες σε θαυμάσονται ἐπὶ σοὶ, καὶ ἐροῦσιν, οὗτος ὁ ἄνθρωπος ὁ παροξύνων τὴν γῆν, ὁ σείων βασιλεῖς,
\vs{17}ὁ θεὶς τὴν οἰκουμένην ὅλην ἔρημον, καὶ τὰς πόλεις αὐτοῦ καθεῖλε, τοὺς ἐν ἐπαγωγῇ οὐκ ἔλυσε.
\vs{18}Πάντες οἱ βασιλεῖς τῶν ἐθνῶν ἐκοιμήθησαν ἐν τιμῇ, ἄνθρωπος ἐν τῷ οἴκῳ αὐτοῦ.
\vs{19}Σὺ δὲ ῥιφήσῃ ἐν τοῖς ὄρεσιν, ὡς νεκρὸς ἐβδελυγμένος, μετὰ πολλῶν τεθνηκότων ἐκκεκεντημένων μαχαίραις, καταβαινόντων εἰς ᾅδου.
\vs{20}Ὃν τρόπον ἱμάτιον ἐν αἵματι πεφυρμένον οὐκ ἔσται καθαρὸν, οὕτως οὐδὲ σὺ ἔσῃ καθαρός· διότι τὴν γῆν μου ἀπώλεσας, καὶ τὸν λαόν μου ἐπέκτεινας· οὐ μὴ μείνῃς εἰς τὸν αἰῶνα χρόνον, σπέρμα πονηρόν.
\vs{21}Ἑτοίμασον τὰ τέκνα σου σφαγῆναι ταῖς ἁμαρτίαις τοῦ πατρὸς αὐτῶν, ἵνα μὴ ἀναστῶσι καὶ κληρονομήσωσι τὴν γῆν, καὶ ἐμπλήσωσι τὴν γῆν πολέμων.
\vs{22}Καὶ ἐπαναστήσομαι αὐτοῖς, λέγει Κύριος σαβαὼθ, καὶ ἀπολῶ αὐτῶν ὄνομα, καὶ κατάλειμμα, καὶ σπέρμα· τάδε λέγει Κύριος.
\vs{23}Καὶ θήσω τὴν Βαβυλωνίαν ἔρημον, ὥστε κατοικεῖν ἐχίνους, καὶ ἔσται εἰς οὐδέν· καὶ θήσω αὐτὴν πηλοῦ βάραθρον εἰς ἀπώλειαν.

\vs{24}Τάδε λέγει Κύριος σαβαὼθ, ὃν τρόπον εἴρηκα, οὕτως ἔσται, καὶ ὃν τρόπον βεβούλευμαι, οὕτως μενεῖ,
\vs{25}τοῦ ἀπολέσαι τοὺς Ἀσσυρίους ἐπὶ τῆς γῆς τῆς ἐμῆς, καὶ ἐπὶ τῶν ὀρέων μου· καὶ ἔσονται εἰς καταπάτημα, καὶ ἀφαιρεθήσεται ἀπʼ αὐτῶν ὁ ζυγὸς αὐτῶν, καὶ τὸ κῦδος αὐτῶν ἀπὸ τῶν ὤμων ἀφαιρεθήσεται.
\vs{26}Αὕτη ἡ βουλὴ ἣν βεβούλευται Κύριος ἐπὶ τὴν ὅλην οἰκουμένην, καὶ αὕτη ἡ χεὶρ ἡ ὑψηλὴ ἐπὶ πάντα τὰ ἔθνη.
\vs{27}Ἃ γὰρ ὁ Θεὸς ὁ ἅγιος βεβούλευται, τίς διασκεδάσει; καὶ τὴν χεῖρα αὐτοῦ τὴν ὑψηλὴν τίς ἀποστρέψει;

\vs{28}Τοῦ ἔτους οὗ ἀπέθανεν ὁ βασιλεὺς Ἄχαζ, ἐγενήθη τὸ ῥῆμα τοῦτο.

\vs{29}Μὴ εὐφρανθείητε οἱ ἀλλόφυλοι πάντες, συνετρίβη γὰρ ὁ ζυγὸς τοῦ παίοντος ὑμᾶς· ἐκ γὰρ σπέρματος ὄφεως ἐξελεύσεται ἔκγονα ἀσπίδων, καὶ τὰ ἔκγονα αὐτῶν ἐξελεύσονται ὄφεις πετάμενοι.
\vs{30}Καὶ βοσκηθήσονται πτωχοὶ διʼ αὐτοῦ· πτωχοὶ δὲ ἄνθρωποι ἐπὶ εἰρήνης ἀναπαύσονται· ἀνελεῖ δὲ ἐν λιμῷ τὸ σπέρμα σου, καὶ τὸ κατάλειμμά σου ἀνελεῖ.
\vs{31}Ὀλολύξατε πύλαι πόλεων, κεκραγέτωσαν πόλεις τεταραγμέναι, οἱ ἀλλόφυλοι πάντες, ὅτι ἀπὸ βοῤῥᾶ καπνὸς ἔρχεται, καὶ οὐκ ἔστι τοῦ εἶναι.
\vs{32}Καὶ τί ἀποκριθήσονται βασιλεῖς ἐθνῶν; ὅτι Κύριος ἐθεμελίωσε Σιων, καὶ διʼ αὐτοῦ σωθήσονται οἱ ταπεινοὶ τοῦ λαοῦ.

ΤΟ ῬΗΜΑ ΤΟ ΚΑΤΑ ΤΗΣ ΜΩΑΒΙΤΙΔΟΣ.

\ch{15}
Νυκτὸς ἀπολεῖται ἡ Μωαβῖτις, νυκτὸς γὰρ ἀπολεῖται τὸ τεῖχος τῆς Μωαβίτιδος.
\vs{2}Λυπεῖσθε ἐφʼ ἑαυτοὺς, ἀπολεῖται γὰρ καὶ Δηβὼν, οὗ ὁ βωμὸς ὑμῶν· ἐκεῖ ἀναβήσεσθε κλαίειν, ἐπὶ Ναβαῦ τῆς Μωαβίτιδος· ὀλολύξατε, ἐπὶ πάσης κεφαλῆς φαλάκρωμα, πάντες βραχίονες κατατετμημένοι.
\vs{3}Ἐν ταῖς πλατείαις αὐτῆς περιζώσασθε σάκκους, καὶ κόπτεσθε ἐπὶ τῶν δωμάτων αὐτῆς, καὶ ἐν ταῖς ῥύμαις αὐτῆς, πάντες ὀλολύζετε μετὰ κλαυθμοῦ.
\vs{4}Ὅτι κέκραγεν Ἐσεβὼν καὶ Ἐλεαλὴ, ἕως Ἰασσὰ ἠκούσθη ἡ φωνὴ αὐτῶν· διατοῦτο ἡ ὀσφὺς τῆς Μωαβίτιδος βοᾷ, ἡ ψυχὴ αὐτῆς γνώσεται.
\vs{5}Ἡ καρδία τῆς Μωαβείτιδος βοᾷ ἐν αὐτῇ ἕως Σηγώρ· δάμαλις γάρ ἐστι τριετής· ἐπὶ δὲ τῆς ἀναβάσεως Λουεὶθ, πρὸς σὲ κλαίοντες ἀναβήσονται τῇ ὁδῷ Ἀρωνιείμ· βοᾷ, σύντριμμα καὶ σεισμός.
\vs{6}Τὸ ὕδωρ τῆς Νεμηρεὶμ ἔρημον ἔσται, καὶ ὁ χόρτος αὐτῆς ἐκλείψει· χόρτος γὰρ χλωρὸς οὐκ ἔσται.
\vs{7}Μὴ καὶ οὕτως μέλλει σωθῇναι; ἐπάξω γὰρ ἐπὶ τὴν φάραγγα Ἄραβας, καὶ λήψονται αὐτήν.
\vs{8}Συνῆψε γὰρ ἡ βοὴ τὸ ὅριον τῆς Μωαβίτιδος τῆς Ἀγαλεὶμ, καὶ ὀλολυγμὸς αὐτῆς ἕως τοῦ φρέατος τοῦ Αἰλείμ.
\vs{9}Τὸ δὲ ὕδωρ τὸ Δειμὼν πλησθήσεται αἵματος, ἐπάξω γὰρ ἐπὶ Δειμὼν Ἄραβας, καὶ ἀρῶ τὸ σπέρμα Μωὰβ, καὶ Ἀριὴλ, καὶ τὸ κατάλοιπον Ἄδαμὰ.

\ch{16}
Ἀποστέλω ὡς ἑρπετὰ ἐπὶ τὴν γῆν· μὴ πέτρα ἔρημός ἐστι τὸ ὄρος θυγατρὸς Σιών;
\vs{2}Ἔσῃ γὰρ ὡς πετεινοῦ ἀνιπταμένου νοσσὸς ἀφῃρημένος, ἔσῃ θυγάτηρ Μωὰβ, ἔπειτα δέ Ἀρνῶν πλείονα
\vs{3}βουλεύου, ποίει τε σκέπην πένθους αὕτη διαπαντὸς, ἐν μεσημβρινῇ σκοτίᾳ φεύγουσιν, ἐξέστησαν· μὴ ἀχθῇς,
\vs{4}παροικήσουσί σοι οἱ φυγάδες Μωάβ· ἔσονται σκέπη ὑμῖν ἀπὸ προσώπου διώκοντος, ὅτι ἤρθη ἡ συμμαχία σου, καὶ ὁ ἄρχων ἀπώλετο ὁ καταπατῶν ἀπὸ τῆς γῆς.
\vs{5}Καὶ διορθωθήσεται μετʼ ἐλέους θρόνος, καὶ καθιεῖται ἐπʼ αὐτοῦ μετὰ ἀληθείας ἐν σκηνῇ Δαυὶδ, κρίνων καὶ ἐκζητῶν κρίμα καὶ σπεύδων δικαιοσύνην.

\vs{6}Ἠκούσαμεν τὴν ὕβριν Μωὰβ, ὑβριστὴς σφόδρα τὴν ὑπερηφανίαν ἐξῇρα· οὐχ οὕτως ἡ μαντεία σου, οὐχ οὕτως.

\vs{7}Ὀλολύξει Μωὰβ, ἐν γὰρ τῇ Μωαβίτιδι πάντες ὀλολύξουσι τοῖς κατοικοῦσι δέ Σεθ μελετήσεις, καὶ οὐκ ἐντραπήσῃ.
\vs{8}Τὰ πεδία Ἑσεβὼν πενθήσει, ἄμπελος Σεβαμά· καταπίνοντες τὰ ἔθνη, καταπατήστε τὰς ἀμπέλους αὐτῆς, ἕως Ἰαζήρ· οὐ μὴ συνάψητε, πλανήθητε τὴν ἔρημον, οἱ ἀπεσταλμένοι ἐνκατελείφθησαν, διέβησαν γὰρ πρὸς τὴν θάλασσαν.
\vs{9}Διατοῦτο κλαύσομαι ὡς τὸν κλαυθμὸν Ἰαζὴρ ἄμπελον Σεβαμά· τὰ δένδρα σου κατέβαλεν Ἐσεβὼν καὶ Ἐλεαλὴ, ὅτι ἐπὶ τῷ θερισμῷ καὶ ἐπὶ τῷ τρυγητῷ σου καταπατήσω, καὶ πάντα πεσοῦνται.
\vs{10}Καὶ ἀρθήσεται εὐφροσύνη καὶ ἀγαλλίαμα ἐκ τῶν ἀμπελώνων, καὶ ἐν τοῖς ἀμπελῶσί σου οὐ μὴ εὐφρανθήσονται, καὶ οὐ μὴ πατήσουσιν οἶνον εἰς τὰ ὑπολήνια, πέπαυται γάρ.
\vs{11}Διατοῦτο ἡ κοιλία μου ἐπὶ Μωὰβ ὡς κιθάρα ἠχήσει, καὶ τὰ ἐντός μου ὡς τεῖχος ἐνεκαίνισας.
\vs{12}Καὶ ἔσται εἰς τὸ ἐντραπῆναί σε, ὅτι ἐκοπίασε Μωὰβ ἐπὶ τοῖς βωμοῖς, καὶ εἰσελεύσεται εἰς τὰ χειροποίητα αὐτῆς, ὥστε προσεύξασθαι, καὶ οὐ μὴ δύνηται ἐξελέσθαι αὐτόν.

\vs{13}Τοῦτο τὸ ῥῆμα ὃ ἐλάλησε Κύριος ἐπὶ Μωὰβ, ὁπότε ἐλάλησε.
\vs{14}Καὶ νῦν λέγῳ, ἐν τρισὶν ἔτεσιν ἐτῶν μισθωτοῦ ἀτιμασθήσεται ἡ δόξα Μωὰβ παντὶ τῷ πλούτῳ τῷ πολλῷ, καὶ καταλειφθήσεται ὀλιγοστὸς, καὶ οὐκ ἔντιμος.

ΤΟ ῬΗΜΑ ΤΟ ΚΑΤΑ ΔΑΜΑΣΚΟΥ.

\ch{17}
Ἰδοὺ Δαμασκὸς ἀρθήσεται ἀπὸ πόλεων, καὶ ἔσται εἰς πτῶσιν,
\vs{2}καταλελειμμένη εἰς τὸν αἰῶνα, εἰς κοίτην ποιμνίων καὶ ἀνάπαυσιν, καὶ οὐκ ἔσται ὁ διώκων.
\vs{3}Καὶ οὐκέτι ἔσται ὀχυρὰ τοῦ καταφυγεῖν Ἐφραίμ· καὶ οὐκέτι βασιλεία ἐν Δαμασκῷ, καὶ τὸ λοιπὸν τῶν Σύρων οὐ γὰρ σὺ βελτίων εἶ τῶν υἱῶν Ἰσραὴλ, τῆς δόξης αὐτῶν· τάδε λέγει Κύριος σαβαώθ.
\vs{4}Ἔσται ἐν τῇ ἡμέρᾳ ἐκείνῃ ἔκλειψις δόξης Ἰακὼβ, καὶ τὰ πίονα τῆς δόξης αὐτοῦ σεισθήσεται.
\vs{5}Καὶ ἔσται ὃν τρόπον ἐάν τις συναγάγῃ ἀμητὸν ἑστηκότα, καὶ σπέρμα σταχύων ἀμήσῃ· καὶ ἔσται ὃν τρόπον ἐάν τις συναγάγῃ στάχυν ἐν φάραγγι στερεᾷ.
\vs{6}Καὶ καταλειφθῇ ἐν αὐτῇ καλάμη, ἢ ὡς ῥῶγες ἐλαίας δύο ἢ τρεῖς ἐπʼ ἄκρου μετεώρου, ἢ τέσσαρες ἢ πέντε ἀπὶ τῶν κλάδων αὐτῶν καταλειφθῇ· τὰδε λέγει Κύριος ὁ Θεὸς Ἰσραήλ.

\vs{7}Τῇ ἡμέρᾳ ἐκείνῃ πεποιθὼς ἔσται ὁ ἄνθρωπος ἐπὶ τῷ ποιήσαντι αὐτὸν, οἱ δὲ ὀφθαλμοὶ αὐτοῦ εἰς τὸν ἅγιον τοῦ Ἰσραὴλ ἐμβλέψονται,
\vs{8}καὶ οὐ μὴ πεποιθότες ὦσιν ἐπὶ τοῖς βωμοῖς, οὐδὲ ἐπὶ τοῖς ἔργοις τῶν χειρῶν αὐτῶν, ἃ ἐποίησαν οἱ δάκτυλοι αὐτῶν, καὶ οὐκ ὄψονται τὰ δένδρα, οὐδὲ τὰ βδελύγματα αὐτῶν.

\vs{9}Τῇ ἡμέρᾳ ἐκείνῃ ἔσονται αἱ πόλεις σου ἐγκαταλελειμμέναι, ὃν τρόπον κατέλιπον οἱ Ἀμοῤῥαῖοι καὶ οἱ Εὐαῖοι ἀπὸ προσώπου τῶν νἱῶν Ἰσραήλ· καὶ ἔσονται ἔρημοι,
\vs{10}διότι κατέλιπες τὸν Θεὸν τὸν σωτῆρά σου, καὶ Κυρίου τοῦ βοηθοῦ σου οὐκ ἐμνήσθης· διὰτοῦτο φυτεύσεις φύτευμα ἄπιστον, καὶ σπέρμα ἄπιστον.
\vs{11}Τῇ ἡμέρᾳ, ᾗ ἂν φυτεύσῃς, πλανηθήσῃ· τὸ δὲ πρωῒ ἐὰν σπείρῃς, ἀνθήσει εἰς ἀμητὸν ᾗ ἂν ἡμέρᾳ κληρώσῃ, καὶ ὡς πατὴρ ἀνθρώπου κληρώσῃ τοῖς υἱοῖς σου.

\vs{12}Οὐαὶ πλῆθος ἐθνῶν πολλῶν· ὡς θάλασσα κυμαίνουσα, οὕτω ταραχθήσεσθε· καὶ νῶτος ἐθνῶν πολλῶν, ὡς ὕδωρ ἠχήσει.
\vs{13}Ὡς ὕδωρ πολὺ ἔθνη πολλὰ, ὡς ὕδατος πολλοῦ βίᾳ φερομένου· καὶ ἀποσκορακιεῖ αὐτὸν, καὶ πόῤῥω αὐτὸν διώξεται, ὡς χνοῦν ἀχύρου λικμώντων ἀπέναντι ἀνέμου, καὶ ὡς κονιορτὸν τροχοῦ καταιγὶς φέρουσα,

\vs{14}Πρὸς ἑσπέραν καὶ ἔσται πένθος, πριν ἢ πρωῒ, καὶ οὐκ ἔσται· αὕτη ἡ μερὶς τῶν προνομευσάντων ὑμᾶς, καὶ κληρονομία τοῖς ὑμᾶς κληρονομήσασιν.

\ch{18}
Οὐαὶ γῆς πλοίων πτέρυγες, ἐπέκεινα ποταμῶν Αἰθιοπίας·
\vs{2}Ὁ ἀποστέλλων ἐν θαλάσσῃ ὅμηρα, καὶ ἐπιστολὰς βιβλίνας ἐπάνω τοῦ ὕδατος· πορεύσονται γὰρ ἄγγελοι κοῦφοι πρὸς ἔθνος μετέωρον, καὶ ξένον λαὸν καὶ χαλεπόν· τίς αὐτοῦ ἐπέκεινα; ἔθνος ἀνέλπιστον καὶ καταπεπατημένον· νῦν οἱ ποταμοὶ τῆς γῆς
\vs{3}πάντες, ὡς χώρα κατοικουμένη κατοικηθήσεται· ἡ χώρα αὐτῶν ὡσεὶ σημεῖον ἀπὸ ὄρους ἀρθῇ, ὡς σάλπιγγος φωνὴ ἀκουστὸν ἔσται.
\vs{4}Διότι οὕτως εἶπε Κύριός μοι, ἀσφάλεια ἔσται ἐν τῇ ἐμῇ πόλει, ὡς φῶς καύματος μεσημβρίας, καὶ ὡς νεφέλη δρόσου ἡμέρας ἀμητο ἔσται
\vs{5}πρὸ τοῦ θερισμοῦ, ὅταν συντελεσθῇ ἄνθος, καὶ ὄμφαξ ἐξανθήσῃ ἄνθος ὀμφακίζουσα· καὶ ἀφελεῖ τὰ βοτρύδια τὰ μικρὰ τοῖς δρεπάνοις, καὶ τὰς κληματίδας ἀφελεῖ, καὶ ἀποκόψει,
\vs{6}καὶ καταλείψει ἅμα τοῖς πετεινοῖς τοῦ οὐρανοῦ, καὶ τοῖς θηρίοις τῆς γῆς· καὶ συναχθήσεται ἐπʼ αὐτοὺς τὰ πετεινὰ τοῦ οὐρανοῦ, καὶ πάντα τὰ θηρία τῆς γῆς ἐπʼ αὐτὸν ἥξει.
\vs{7}Ἐν τῷ καιρῷ ἐκείνῳ ἀνενεχθήσεται δῶρα Κυρίῳ σαβαὼθ ἐκ λαοῦ τεθλιμμένου καὶ τετιλμένου, καὶ ἀπὸ λαοῦ μεγάλου ἀπὸ τοῦ νῦν καὶ εἰς τὸν αἰῶνα χρόνον· ἔθνος ἐλπίζον καὶ καταπεπατημένον, ὅ ἐστιν ἐν μέρει ποταμοῦ τῆς χώρας αὐτοῦ, εἰς τὸν τόπον οὗ τὸ ὄνομα Κυρίου σαβαὼθ, ὄρος Σιών.

ὍΡΑΣΙΣ ΑΙΓΥΠΤΟΥ.

\ch{19}
Ἰδοὺ Κύριος κάθηται ἐπὶ νεφέλης κούφης, καὶ ἥξει εἰς Αἴγυπτον, καὶ σεισθήσεται τὰ χειροποίητα Αἰγύπτου ἀπὸ προσώπου αὐτοῦ· καὶ ἡ καρδία αὐτῶν ἡττηθήσεται ἐν αὐτοῖς.
\vs{2}Καὶ ἐπεγερθήσονται Αἰγύπτιοι ἐπʼ Αἰγυπτίους, καὶ πολεμήσει ἄνθρωπος τὸν ἀδελφὸν αὐτοῦ, καὶ ἄνθρωπος τὸν πλησίον αὐτοῦ, πόλις ἐπὶ πόλιν, καὶ νὸμος ἐπὶ νὸμον.
\vs{3}Καὶ ταραχθήσεται τὸ πνεῦμα τῶν Αἰγυπτίων ἐν αὐτοῖς, καὶ τὴν βουλὴν αὐτῶν διασκεδάσω, καὶ ἐπερωτήσουσι τοὺς θεοὺς αὐτῶν, καὶ τὰ ἀγάλματα αὐτῶν, καὶ τοὺς ἐκ τῆς γῆς φωνοῦντας, καὶ τοὺς ἐγγαστριμύθους.
\vs{4}Καὶ παραδώσω τὴν Αἴγυπτον εἰς χεῖρας ἀνθρώπων, κυρίων σκληρῶν, καὶ βασιλεῖς σκληροὶ κυριεύσουσιν αὐτῶν. τάδε λέγει Κύριος σαβαώθ.
\vs{5}Καὶ πίονται οἱ Αἰγύπτιοι ὕδωρ τὸ παρὰ θάλασσαν, ὁ δὲ ποταμὸς ἐκλείψει, καὶ ξηρανθήσεται.
\vs{6}Καὶ ἐκλείψουσιν οἱ ποταμοὶ, καὶ αἱ διώρυχες τοῦ ποταμοῦ· καὶ ξηρανθήσεται πᾶσα συναγωγὴ ὕδατος, καὶ ἐν παντὶ ἕλει καλάμου καὶ παπύρου,
\vs{7}καὶ τὸ ἄχι τὸ χλωρὸν πᾶν τὸ κύκλῳ τοῦ ποταμοῦ, καὶ πᾶν τὸ σπειρόμενον διὰ τοῦ ποταμοῦ ξηρανθήσεται ἀνεμόφθορον.
\vs{8}Καὶ στενάξουσιν οἱ ἁλιεῖς, καὶ στενάξουσι πάντες οἱ βάλλοντες ἄγκιστρον εἰς τὸν ποταμὸν, καὶ οἱ βάλλοντες σαγήνας, καὶ οἱ ἀμφιβολεῖς πενθήσουσι·
\vs{9}Καὶ αἰσχύνη λήψεται τοὺς ἐργαζομένους τὸ λίνον τὸ σχιστὸν, καὶ τοὺς ἐργαζομένους τὴν βύσσον.
\vs{10}Καὶ ἔσονται οἱ ἐργαζόμενοι αὐτὰ ἐν ὀδύνῃ, καὶ πάντες οἱ ποιοῦντες τὸν ζῦθον λυπηθήσονται, καὶ τὰς ψυχὰς πονέσουσι.
\vs{11}Καὶ μωροὶ ἔσονται οἱ ἄρχοντες Τάνεως, οἱ σοφοὶ σύμβουλοι τοῦ βασιλέως, ἡ βουλὴ σὐτῶν μωρανθήσεται· πῶς ἐρεῖτε τῷ βασιλεῖ, υἱοὶ συνετῶν ἡμεῖς, υἱοὶ βασιλέων τῶν ἐξ ἀρχῆς;
\vs{12}Ποῦ εἰσι νῦν οἱ σοφοί σου; καὶ ἀναγγειλάτωσάν σοι, καὶ εἰπάτωσαν, τί βεβούλευται Κύριος σαβαὼθ ἐπʼ Αἴγυπτον;
\vs{13}Ἐξέλιπον οἱ ἄρχοντες Τάνεως, καὶ ὑψώθησαν οἱ ἄρχοντες Μέμφεως, καὶ πλανήσουσιν Αἴγυπτον κατὰ φυλάς.
\vs{14}Κύριος γὰρ ἐκέρασεν αὐτοῖς πνεῦμα πλανήσεως, καὶ ἐπλάνησαν Αἴγυπτον ἐν πᾶσι τοῖς ἔργοις αὐτῶν, ὡς πλανᾶται ὁ μεθύων, καὶ ὁ ἐμῶν ἅμα,
\vs{15}καὶ οὐκ ἔσται τοῖς Αἰγυπτίοις ἔργον ὃ ποιήσει κεφαλὴν καὶ οὐρὰν, καὶ ἀρχὴν καὶ τέλος.

\vs{16}Τῇ δὲ ἡμέρᾳ ἐκείνῃ ἔσονται οἱ Αἰγύπτιοι, ὡς γυναῖκες, ἐν φόβῳ καὶ ἐν τρόμῳ ἀπὸ προσώπου τῆς χειρὸς Κυρίου σαβαὼθ, ἣν αὐτὸς ἐπιβαλεῖ αὐτοῖς.
\vs{17}Καὶ ἔσται ἡ χώρα τῶν Ἰουδαίων τοῖς Αἰγυπτίοις εἰς φόβητρον· πᾶς ὃς ἐὰν ὀνομάσῃ αὐτὴν αὐτοῖς, φοβηθήσονται διὰ τὴν βουλὴν ἣν βεβούλευται Κύριος σαβαὼθ ἐπʼ αὐτήν.
\vs{18}Τῇ ἡμέρᾳ ἐκείνῃ ἔσονται πέντε πόλεις ἐν Αἰγύπτῳ λαλοῦσαι τῇ γλώσσῃ τῇ Χαναανίτιδι, καὶ ὀμνῦντες τῷ ὀνόματι Κυρίου σαβαώθ· πόλις ἀσεδὲκ κληθήσεται ἡ μία πόλις.
\vs{19}Τῇ ἡμερᾳ ἐκείνῃ ἔσται θυσιαστήριον τῷ Κυρίῳ ἐν χώρᾳ Αἰγυπτίων, καὶ στήλη πρὸς τὸ ὅριον αὐτῆς τῷ Κυρίῳ.
\vs{20}καὶ ἔσται εἰς σημεῖον εἰς τὸν αἰῶνα Κυρίῳ ἐν χώρᾳ Αἰγύπτου· ὅτι κεκράξονται πρὸς Κύριον διὰ τοὺς θλίβοντας αὐτούς, καὶ ἀποστελεῖ αὐτοῖς ἄνθρωπον ὃς σώσει αὐτοὺς, κρίνων σώσει αὐτούς.
\vs{21}Καὶ γνωστὸς ἔσται Κύριος τοῖς Αἰγυπτίοις· καὶ γνώσονται οἱ Αἰγύπτιοι τὸν Κύριον ἐν τῇ ἡμέρᾳ ἐκείνῃ, καὶ ποιήσουσι θυσίας, καὶ εὔξονται εὐχὰς τῷ Κυρίῳ, καὶ ἀποδώσουσι.
\vs{22}Καὶ πατάξει Κύριος τοὺς Αἰγυπτίους πληγῇ, καὶ ἰάσεται αὐτοὺς ἰάσει, καὶ ἐπιστραφήσονται πρὸς Κύριον, καὶ εἰσακούσεται αὐτῶν, καὶ ἰάσεται αὐτοὺς ἰάσει.
\vs{23}Τῇ ἡμέρᾳ ἐκείνῃ ἔσται ἡ ὁδὸς ἀπὸ Αἰγύπτου πρὸς Ἀσσυρίους, καὶ εἰσελεύσονται Ἀσσύριοι εἰς Αἴγυπτον, καὶ Αἰγύπτιοι πορεύσονται πρὸς Ἀσσυρίους, καὶ δουλεύσουσιν Αἰγύπτιοι τοῖς Ἀσσυρίοις.
\vs{24}Τῇ ἡμέρᾳ ἐκείνῃ ἔσται Ἰσραὴλ τρίτος ἐν τοῖς Αἰγυπτίοις, καὶ ἐν τοῖς Ἀσσυρίοις εὐλογημένος ἐν τῇ γῇ
\vs{25}ἣν εὐλόγησε Κύριος σαβαὼθ, λέγων, εὐλογημένος ὁ λαός μου ὁ ἐν Αἰγύπτῳ, καὶ ὁ ἐν Ἀσσυρίοις, καὶ ἡ κληρονομία μου Ἰσραήλ.

\ch{20}
Τοῦ ἔτους ὅτε εἰσῆλθε Τανὰθαν εἰς Ἄζωτον, ἡνίκα ἀπεστάλη ὑπὸ Ἀρνᾶ βασιλέως Ἀσσυρίων, καὶ ἐπολέμησε τὴν Ἄζωτον, καὶ ἔλαβεν αὐτὴν,
\vs{2}τότε ἐλάλησε Κύριος πρὸς Ἡσαΐαν υἱὸν Ἀμὼς, λέγων, πορεύου καὶ ἄφελε τὸν σάκκον ἀπὸ τῆς ὀσφύος σου, καὶ τὰ σανδάλιά σου ὑπόλυσαι ἀπὸ τῶν ποδῶν σου, καὶ ποίησον οὕτως, πορευόμενος γυμνὸς καὶ ἀνυπόδετος.
\vs{3}Καὶ εἶπε Κύριος, ὃν τρόπον πεπόρευται ὁ παῖς μου Ἡσαΐας γυμνὸς καὶ ἀνυπόδετος τρία ἔτη, τρία ἔτη ἔσται εἰς σημεῖα καὶ τέρατα τοῖς Αἰγυπτίοις καὶ Αἰθίοψιν·
\vs{4}Ὅτι οὕτως ἄξει βασιλεὺς Ἀσσυρίων τὴν αἰχμαλωσίαν Αἰγύπτου καὶ Αἰθιόπων, νεανίσκους καὶ πρεσβύτας, γυμνοὺς καὶ ἀνυποδέτους, ἀνακεκαλυμμένους τὴν αἰσχύνην Αἰγύπτου.
\vs{5}Καὶ αἰσχυνθήσονται ἡττηθέντες ἐπὶ τοῖς Αἰθίοψιν, ἐφʼ οἷς ἦσαν πεποιθότες οἱ Αἰγύπτιοι, ἦσαν γὰρ αὐτοῖς δόξα.
\vs{6}Καὶ ἐροῦσιν οἱ κατοικοῦντες ἐν τῇ νήσῳ ταύτῃ ἐν τῇ ἡμέρᾳ ἐκείνῃ, ἰδοὺ ἡμεῖς ἦμεν πεποιθότες τοῦ φυγεῖν εἰς αὐτοὺς εἰς βοήθειαν, οἳ οὐκ ἠδύναντο σωθῆναι ἀπὸ βασιλέως Ἀσσυρίων, καὶ πῶς ἡμεῖς σωθησόμεθα;

ΤΟ ὍΡΑΜΑ ΤΗΣ ἘΡΗΜΟΥ.

\ch{21}
Ὡς καταιγὶς διʼ ἐρήμου διέλθοι, ἐξ ἐρήμου ἐρχομένη ἐκ γῆς φοβερὸν
\vs{2}τὸ ὅραμα, καὶ σκληρὸν ἀνηγγέλη μοι· ὁ ἀθετῶν ἀθετεῖ, ὁ ἀνομῶν ἀνομεῖ· ἐπʼ ἐμοὶ οἱ Ἐλαμῖται, καὶ οἱ πρέσβεις τῶν Περσῶν ἐπʼ ἐμὲ ἔρχονται· νῦν στενάξω καὶ παρακαλέσω ἐμαυτόν.
\vs{3}Διατοῦτο ἐνεπλήσθη ἡ ὀσφύς μου ἐκλύσεως, καὶ ὠδῖνες ἔλαβόν με ὡς τὴν τίκτουσαν· ἠδίκησα τοῦ μὴ ἀκοῦσαι, ἐσπούδασα τοῦ μὴ βλέπειν.
\vs{4}Ἡ καρδία μου πλανᾶται, καὶ ἡ ἀνομία με βαπτίζει, ἡ ψυχή μου ἐφέστηκεν εἰς φόβον.
\vs{5}Ἑτοίμασον τὴν τράπεζαν, φάγετε, πίετε· ἀναστάντες οἱ ἄρχοντες, ἑτοιμάσατε θυρεοὺς,
\vs{6}ὅτι οὕτως εἶπε πρὸς μὲ Κύριος, Βαδίσας σεαυτῷ στῆσον σκοπὸν, καὶ ὃ ἂν ἴδῃς ἀνάγγειλον.
\vs{7}Καὶ εἶδον ἀναβάτας ἱππεῖς δύο, καὶ ἀναβάτη ὄνου, καὶ ἀναβάτην καμήλου· ἀκρόασαι ἀκρόασιν πολλήν,
\vs{8}καὶ κάλεσον Οὐρίαν εἰς τὴν σκοπιάν· Κύριος εἶπεν, ἔστην διαπαντὸς ἡμέρας, καὶ ἐπὶ τῆς παρεμβολῆς ἐγὼ ἔστην ὅλην τὴν νύκτα,
\vs{9}καὶ ἰδοὺ αὐτὸς ἔρχεται ἀναβάτης ξυνωρίδος· καὶ ἀποκριθεὶς εἶπε, πέπτωκε πέπτωκε Βαβυλὼν, καὶ πάντα τὰ ἀγάλματα αὐτῆς, καὶ τὰ χειροποίητα αὐτῆς συνετρίβησαν εἰς τὴν γῆν.
\vs{10}Ἀκούσατε οἱ καταλελειμμένοι, καὶ οἱ ὀδυνώμενοι ἀκούσατε ἃ ἤκουσα παρὰ Κυρίου σαβαὼθ, ὁ Θεὸς τοῦ Ἰσραὴλ ἀνήγγειλεν ἡμῖν.

ΤΟ ὍΡΑΜΑ ΤΗΣ ἸΔΟΥΜΑΙΑΣ.

\vs{11}Πρὸς ἐμὲ κάλει παρὰ τοῦ Σηεὶρ, φυλάσσετε ἐπάλξεις.
\vs{12}Φυλάσσω τοπρωῒ καὶ τὴν νύκτα· ἐὰν ζητῇς ζήτει, καὶ παρʼ ἐμοὶ οἴκει,
\vs{13}ἐν τῷ δρυμῷ ἑσπέρας κοιμηθῇς, ἢ ἐν τῇ ὁδῷ Δαιδάν.

\vs{14}Εἰς συνάντησιν διψῶντι ὕδωρ φέρετε οἱ ἐνοικοῦντες ἐν χώρᾳ Θαιμὰν, ἄρτοις συναντᾶτε τοῖς φεύγουσι
\vs{15}διὰ τὸ πλῆθος τῶν πεφονευμένων, καὶ διὰ τὸ πλῆθος τῶν πλανωμένων, καὶ διὰ τὸ πλῆθος τῆς μαχαίρας, καὶ διὰ τὸ πλῆθος τῶν τοξευμάτων τῶν διατεταμένων, καὶ διὰ τὸ πλῆθος τῶν πεπτωκότων ἐν τῷ πολέμῳ.
\vs{16}Διότι οὕτως εἶπέ μοι Κύριος, ἔτι ἐνιαυτὸς ὡς ἐνιαυτὸς μισθωτοῦ, ἐκλείψει ἡ δόξα τῶν υἱῶν Κηδὰρ,
\vs{17}καὶ τὸ κατάλοιπον τῶν τοξευμάτων τῶν ἰσχυρῶν υἱῶν Κηδὰρ ἔσται ὀλίγον, ὅτι Κύριος ὁ Θεὸς Ἰσραὴλ ἐλάλησε.

ΤΟ ῬΗΜΑ ΤΗΣ ΦΑΡΑΓΓΟΣ ΣΙΩΝ.

\ch{22}
Τί ἐγένετό σοι, ὅτι νῦν ἀνέβητε πάντες εἰς δώματα μάταια;
\vs{2}Ἐνεπλήσθη ἡ πόλις βοώντων, οἱ τραυματίαι σου οὐ τραυματίαι ἐν μαχαίραις, οὐδὲ οἱ νεκροί σου νεκροὶ πολέμων.
\vs{3}Πάντες οἱ ἄρχοντές σου πεφεύγασι, καὶ οἱ ἁλόντες σκληρῶς δεδεμένοι εἰσί, καὶ οἱ ἰσχύοντες ἐν σοὶ πόῤῥω πεφεύγασι.
\vs{4}Διὰ τοῦτο εἶπα, ἄφετέ με, πικρῶς κλαύσομαι· μὴ κατισχύσητε παρακαλεῖν με ἐπὶ τὸ σύντριμμα τῆς θυγατρὸς τοῦ γένους μου,
\vs{5}ὅτι ἡμέρα ταραχῆς καὶ ἀπωλείας καὶ καταπατήματος, καὶ πλάνησις παρὰ Κυρίου σαβαώθ· ἐν φάραγγι Σιὼν πλανῶνται, ἀπὸ μικροῦ ἕως μεγάλου πλανῶνται ἐπὶ τὰ ὄρη.
\vs{6}Οἱ δὲ Ἐλαμεῖται ἔλαβον φαρέτρας, καὶ ἀναβάται ἄνθρωποι ἐφʼ ἵππους, καὶ συναγωγὴ παρατάξεως.
\vs{7}Καὶ ἔσονται αἱ ἐκλεκταὶ φάραγγές σου, πλησθήσονται ἁρμάτων, οἱ δὲ ἱππεῖς ἐμφράξουσι τὰς πύλας σου,
\vs{8}καὶ ἀνακαλύψουσι τὰς πύλας Ἰούδα· καὶ ἐμβλέψονται τῇ ἡμέρᾳ ἐκείνῃ εἰς τοὺς ἐκλεκτοὺς οἴκους τῆς πόλεως·
\vs{9}Καὶ ἀνακαλύψουσι τὰ κρυπτὰ τῶν οἴκων τῆς ἄκρας Δαυείδ· καὶ εἴδοσαν, ὅτι πλείους εἰσί, καὶ ὅτι ἀπέστρεψε τὸ ὕδωρ τῆς ἀρχαίας κολυμβήθρας εἰς τὴν πόλιν,
\vs{10}καὶ ὅτι καθείλοσαν τοὺς οἴκους Ἱερουσαλὴμ εἰς ὀχυρώματα τείχους τῇ πόλει.
\vs{11}Καὶ ἐποιήσατε ἑαυτοῖς ὕδωρ ἀναμέσον τῶν δύο τειχῶν ἐσώτερον τῆς κολυμβήθρας τῆς ἀρχαίας, καὶ οὐκ ἐνεβλέψατε εἰς τὸν ἀπʼ ἀρχῆς ποιήσαντα αὐτὴν, καὶ τὸν κτίσαντα αὐτὴν οὐκ εἴδετε.
\vs{12}Καὶ ἐκάλεσε Κύριος Κύριος σαβαὼθ ἐν τῇ ἡμέρᾳ ἐκείνῃ κλαυθμὸν καὶ κοπετὸν, καὶ ξύρησιν καὶ ζῶσιν σάκκων,
\vs{13}αὐτοὶ δὲ ἐποιήσαντο εὐφροσύνην καὶ ἀγαλλίαμα, σφάζοντες μόσχους, καὶ θύοντες πρόβατα, ὥστε φαγεῖν κρέατα, καὶ πιεῖν οἶνον, λέγονες, φάγωμεν καὶ πίωμεν, αὔριον γὰρ ἀποθνήσκομεν.
\vs{14}Καὶ ἀνακεκαλυμμένα ταῦτά ἐστιν ἐν τοῖς ὠσὶ Κυρίου σαβαὼθ, ὅτι οὐκ ἀφεθήσεται ὑμῖν αὕτη ἡ ἁμαρτία, ἕως ἂν ἀποθάνητε.

\vs{15}Τάδε λέγει Κύριος σαβαὼθ, πορεύου εἰς τὸ παστοφόριον, πρὸς Σομνᾶν τὸν ταμίαν, καὶ εἰπὸν αὐτῷ,
\vs{16}τί σὺ ὧδε καὶ τί σοί ἐστιν ὧδε; ὅτι ἐλατόμησας σεαυτῷ ὧδε μνημεῖον, καὶ ἐποίησας σεαυτῷ ἐν ὑψηλῷ μνημεῖον, καὶ ἔγραψας σεαυτῷ ἐν πέτρᾳ σκηνήν;
\vs{17}Ἰδοὺ δὴ Κύριος σαβαὼθ ἐκβάλλει καὶ ἐκτρίψει ἄνδρα, καὶ ἀφελεῖ τὴν στολήν σου καὶ τὸν στέφανόν σου τὸν ἔνδοξον,
\vs{18}καὶ ῥίψει σε εἰς χώραν μεγάλην καὶ ἀμέτρητον, καὶ ἐκεῖ ἀποθανῇ· καὶ θήσει τὸ ἅρμα σου τὸ καλὸν εἰς ἀτιμίαν, καὶ τὸν οἶκον τοῦ ἄρχοντός σου εἰς καταπάτημα.
\vs{19}Καὶ ἀφαιρεθήσῃ ἐκ τῆς οἰκονομίας σου καὶ ἐκ τῆς στάσεώς σου.
\vs{20}Καὶ ἔσται ἐν τῇ ἡμέρᾳ ἐκείνῃ καὶ καλέσω τὸν παῖδά μου Ἐλιακεὶμ τὸν τοῦ Χελκίου,
\vs{21}καὶ ἐνδύσω αὐτὸν τὴν στολήν σου, καὶ τὸν στέφανόν σου δώσω αὐτῷ κατὰ κράτος, καὶ τὴν οἰκονομίαν σου δώσω εἰς τὰς χεῖρας αὐτοῦ· καὶ ἔσται ὡς πατὴρ τοῖς ἐνοικοῦσιν ἐν Ἱερουσαλὴμ, καὶ τοῖς ἐνοικοῦσιν ἐν Ἰούδᾳ.
\vs{22}Καὶ δώσω τὴν δόξαν Δαυεὶδ αὐτῷ, καὶ ἄρξει, καὶ οὐκ ἔσται ὁ ἀντιλέγων· καὶ δώσω αὐψῷ τὴν κλεῖδα οἴκου Δαυὶδ ἐπὶ τῷ ὢμῳ αὐτοῦ· καὶ ἀνοίξει, καὶ οὐκ ἔσται ὁ ἀποκλείων· καὶ κλείσει, καὶ οὐκ ἔσται ὁ ἀνοίγων.
\vs{23}Καὶ στήλω αὐτὸν ἄρχοντα ἐν τόπῳ πιστῷ, καὶ ἔσται εἰς θρόνον δόξης τοῦ οἴκου τοῦ πατρὸς αὐτοῦ.
\vs{24}Καὶ ἔσται πεποιθὼς ἐπʼ αὐτὸν πᾶς ἔνδοξος ἐν τῷ οἴκῳ τοῦ πατρὸς αὐτοῦ, ἀπὸ μικροῦ ἕως μεγάλου, καὶ ἔσονται ἐπικρεμάμενοι αὐτῷ
\vs{25}τῇ ἡμέρᾳ ἐκείνῃ· τάδε λέγει Κύριος σαβαὼθ, κινηθήσεται ὁ ἄνθρωπος ὁ ἐστηριγμένος ἐν τόπῳ πιστῷ, καὶ ἀφαιρεθήσεται, καὶ πεσεῖται, καὶ ἐξολεθρευθήσεται ἡ δόξα ἡ ἐπʼ αὐτόν, ὅτι Κύριος ἐλάλησεν.

ΤΟ ῬΗΜΑ ΤΥΡΟΥ.

\ch{23}
Ὀλολύξατε, πλοῖα Καρχηδόνος, ὅτι ἀπώλετο, καὶ οὐκέτι ἔρχονται ἐκ γῆς Κιτιαίων, ἦκται αἰχμάλωτος.
\vs{2}Τίνι ὅμοιοι γεγόνασιν οἱ ἐνοικοῦντες ἐν τῇ νήσῳ, μετάβολοι Φοινίκης, διαπερῶντες τὴν θάλασσαν
\vs{3}ἐν ὕδατι πολλῷ, σπέρμα μετάβολων; ὡς ἀμητοῦ εἰσφερομένου, οἱ μεταβόλοι τῶν ἐθνῶν.
\vs{4}αἰσχύνθητι Σιδὼν, εἶπεν ἡ θάλασσα· ἡ δὲ ἰσχὺς τῆς θαλάσσης εἶπεν, οὐκ ὤδινον, οὐδὲ ἔτεκον, οὐδὲ ἐξέθρεψα νεανίσκους, οὐδὲ ὕψωσα παρθένους.
\vs{5}Ὅταν δὲ ἀκουστὸν γένηται Αἰγύπτῳ, λήμψεται αὐτοὺς ὀδύνη περὶ Τύρου.
\vs{6}Ἀπέλθατε εἰς Καρχηδόνα, ὀλολύξατε οἱ κατοικοῦντες ἐν τῇ νήσῳ ταύτῃ.
\vs{7}Οὐχ αὕτη ἦν ὑμῶν ἡ ὕβρις ἀπʼ ἀρχῆς, πρινὴ παραδοθῆναι αὐτήν;
\vs{8}Τίς ταῦτα ἐβούλευσεν ἐπὶ Τύρον; μὴ ἥσσων ἐστίν, ἢ οὐκ ἰσχύει; οἱ ἔμποροι αὐτῆς ἔνδοξοι ἄρχοντες τῆς γῆς.

\vs{9}Κύριος σαβαὼθ ἐβουλεύσατο παραλῦσαι πᾶσαν τὴν ὕβριν τῶν ἐνδόξων, καὶ ἀτιμάσαι πᾶν ἔνδοξον ἐπὶ τῆς γῆς.
\vs{10}Ἐργάζου τὴν γῆν σου, καὶ γὰρ πλοῖα οὐκέτι ἔρχεται ἐκ Καρχηδόνος.
\vs{11}Ἡ δὲ χείρ σου οὐκέτι ἰσχύει κατὰ θάλασσαν, ἡ παροξύνουσα βασιλεῖς· Κύριος σαβαὼθ ἐνετείλατο περὶ Χαναὰν ἀπολέσαι αὐτῆς τὴν ἰσχύν.
\vs{12}Καὶ ἐροῦσιν, οὐκέτι οὐ μὴ προστεθῆτε τοῦ ὑβρίζειν καὶ ἀδικεῖν τὴν θυγατέρα Σιδῶνος· καὶ ἐὰν ἀπέλθῃς εἰς Κιτιεῖς, οὐδὲ ἐκεῖ ἀνάπαυσις ἔσται σοι·
\vs{13}Καὶ εἰς γῆν Χαλδαίων, καὶ αὕτη ἠρήμωται ἀπὸ τῶν Ἀσσυρίων, ὅτι ὁ τοῖχος αὐτῆς πέπτωκεν.
\vs{14}Ὀλολύξατε πλοῖα Καρχηδόνος, ὅτι ἀπόλωλε τὸ ὀχύρωμα ὑμῶν.

\vs{15}Καὶ ἔσται ἐν τῇ ἡμέρᾳ ἐκείνῃ, καταλειφθήσεται Τύρος ἔτη ἑβδομήκοντα, ὡς χρόνος βασιλέως, ὡς χρόνος ἀνθρώπου· καὶ ἔσται μετὰ ἑβδομήκοντα ἔτη, ἔσται Τύρος ὡς ᾆσμα πόρνης.
\vs{16}Λάβε κιθάραν, ῥέμβευσον πόλις πόρνη ἐπιλελησμένη, καλῶς κιθάρισον, πολλὰ ᾆσον, ἵνα σου μνεία γένηται.
\vs{17}Καὶ ἔσται μετὰ τὰ ἑβδομήκοντα ἔτη, ἐπισκοπὴν ποιήσει ὁ Θεὸς Τύρου, καὶ πάλιν ἀποκαταστήσεται εἰς τὸ ἀρχαῖον, καὶ ἔσται ἐμπόριον πάσαις ταῖς βασιλείαις τῆς οἰκουμένης ἐπὶ πρόσωπον τῆς γῆς.
\vs{18}Καὶ ἔσται αὐτῆς ἡ ἐμπορία καὶ ὁ μισθὸς ἅγιον Κυρίῳ· οὐκ αὐτοῖς συναχθήσεται, ἀλλὰ τοῖς κατοικοῦσιν ἔναντι Κυρίου, πᾶσα ἡ ἐμπορία αὐτῆς, φαγεῖν καὶ πιεῖν καὶ ἐμπλησθῆναι, καὶ εἰς συμβολὴν μνημόσυνον ἔναντι Κυρίου.

\ch{24}
Ἰδοὺ Κύριος καταφθείρει τὴν οἰκουμένην, καὶ ἐρημώσει αὐτὴν, καὶ ἀνακαλύψει τὸ πρόσωπον αὐτῆς, καὶ διασπερεῖ τοὺς ἐνοικοῦντας ἐν αὐτῇ.
\vs{2}Καὶ ἔσται ὁ λαὸς ὡς ἱερεύς, καὶ ὁ παῖς ὡς ὁ κύριος, καὶ ἡ θεράπαινα ὡς ἡ κυρία· ἔσται ὁ ἀγοράζων ὡς ὁ πωλῶν, ὁ δανείζων ὡς ὁ δανειζόμενος, καὶ ὁ ὀφείλων ὡς ᾧ ὀφείλει.

\vs{3}Φθορᾷ φθαρήσεται ἡ γῆ, καὶ προνομῇ προνομευθήσεται ἡ γῆ· τὸ γὰρ στόμα Κυρίου ἐλάλησε ταῦτα.
\vs{4}Ἐπένθησεν ἡ γῆ, καὶ ἐφθάρη ἡ οἰκουμένη, ἐπένθησαν οἱ ὑψηλοὶ τῆς γῆς.
\vs{5}Ἡ δὲ ἠνόμησε διὰ τοὺς κατοικοῦντας αὐτήν, διότι παρήλθοσαν τὸν νόμον, καὶ ἤλλαξαν τὰ προστάγματα διαθήκην αἰώνιον.
\vs{6}Διατοῦτο ἀρὰ ἔδεται τὴν γῆν, ὅτι ἡμάρτοσαν οἱ κατοικοῦντες αὐτήν· διατοῦτο πτωχοὶ ἔσονται οἱ ἐνοικοῦντες ἐν τῇ γῇ, καὶ καταλειφθήσονται ἄνθρωποι ὀλίγοι.
\vs{7}Πενθήσει οἶνος, πενθήσει ἄμπελος, στενάξουσιν πάντες οἱ εὐφραινόμενοι τὴν ψυχήν.
\vs{8}Πέπαυται εὐφροσύνη τυμπάνων, πέπαυται φωνὴ κιθάρας.
\vs{9}Ἠσχύνθησαν, οὐκ ἔπιον οἶνον, πικρὸν ἐγένετο τὸ σίκερα τοῖς πίνουσιν.
\vs{10}Ἠρημώθη πᾶσα πόλις, κλείσει οἰκίαν τοῦ μὴ εἰσελθεῖν.
\vs{11}Ὀλολύζεται περὶ τοῦ οἴνου πανταχῇ, πέπαυται πᾶσα εὐφροσύνη τῆς γῆς, ἀπῆλθε πᾶσα εὐφροσύνη τῆς γῆς.
\vs{12}Καὶ καταλειφθήσονται πόλεις ἔρημοι, καὶ οἶκοι ἐγκαταλελειμμένοι ἀπολοῦνται.

\vs{13}Ταῦτα πάντα ἔσονται ἐν τῇ γῇ ἐν μέσῳ τῶν ἐθνῶν· ὃν τρόπον ἐάν τις καλαμήσηται ἐλαίαν, οὕτως καλαμήσονται αὐτούς· καὶ ἐὰν παύσηται ὁ τρυγητὸς,
\vs{14}οὗτοι βοῇ φωνήσουσιν. οἱ δὲ καταλειφθέντες ἐπὶ τῆς γῆς εὐφρανθήσονται ἅμα τῇ δόξῃ Κυρίου, ταραχθήσεται τὸ ὕδωρ τῆς θαλάσσης.
\vs{15}Διατοῦτο ἡ δόξα Κυρίου ἐν ταῖς νήσοις ἔσται τῆς θαλάσσης, τὸ ὄνομα Κυρίου ἔνδοξον ἔσται.

\vs{16}Κύριε ὁ Θεὸς Ἰσραὴλ, ἀπὸ τῶν πτερύγων τῆς γῆς τέρατα ἠκούσαμεν, ἐλπὶς τῷ εὐσεβεῖ· καὶ ἐροῦσιν, οὐαὶ τοῖς ἀθετοῦσιν, οἱ ἀθετοῦντες τὸν νόμον.
\vs{17}Φόβος καὶ βόθυνος καὶ παγὶς ἐφʼ ὑμᾶς τοὺς ἐνοικοῦντας ἐπὶ τῆς γῆς.
\vs{18}Καὶ ἔσται ὁ φεύγων τὸν φόβον, ἐμπεσεῖται εἰς τὸν βόθυνον· καὶ ὁ ἐκβαίνων ἐκ τοῦ βοθύνου, ἁλώσεται ὑπὸ τῆς παγίδος· ὅτι θυρίδες ἐκ τοῦ οὐρανοῦ ἀνεῴχθησαν, καὶ σεισθήσεται τὰ θεμέλια τῆς γῆς.
\vs{19}Ταραχῇ ταραχθήσεται ἡ γῆ, καὶ ἀπορίᾳ ἀπορηθήσεται ἡ γῆ.
\vs{20}Ἔκλινεν ὡς ὁ μεθύων καὶ κραιπαλῶν, καὶ σεισθήσεται ὡς ὀπωροφυλάκιον ἡ γῆ· κατίσχυσεν γὰρ ἐπʼ αὐτῆς ἡ ἀνομία, καὶ πεσεῖται, καὶ οὐ μὴ δύνηται ἀναστῆναι.

\vs{21}Καὶ ἐπάξει ὁ Θεὸς ἐπὶ τὸν κόσμον τοῦ οὐρανοῦ τὴν χεῖρα, καὶ ἐπὶ τοὺς βασιλεῖς τῆς γῆς.
\vs{22}Καὶ συνάξουσι συναγωγὴν αὐτῆς εἰς δεσμωτήριον, καὶ ἀποκλείσουσιν εἰς ὀχύρωμα· διὰ πολλῶν γενεῶν ἐπισκοπὴ ἔσται αὐτῶν.
\vs{23}Καὶ τακήσεται ἡ πλίνθος, καὶ πεσεῖται τὸ τεῖχος· ὅτι βασιλεύσει Κύριος ἐκ Σιὼν, καὶ ἐξ Ἱερουσαλὴμ, καὶ ἐνώπιον τῶν πρεσβυτέρων δοξασθήσεται.

\ch{25}
Κύριε ὁ Θεὸς δοξάσω σε, ὑμνήσω τὸ ὄνομά σου, ὅτι ἐποίησας θαυμαστὰ πράγματα, βουλὴν ἀρχαίαν ἀληθινήν· γένοιτο.
\vs{2}Ὅτι ἔθηκας πόλεις εἰς χῶμα, πόλεις ὀχυρὰς τοῦ μὴ πεσεῖν αὐτῶν τὰ θεμέλια· τῶν ἀσεβῶν πόλις τὸν αἰῶνα οὐ μὴ οἰκοδομηθῇ.
\vs{3}Διατοῦτο εὐλογήσει σε ὁ λαὸς ὁ πτωχὸς, καὶ πόλεις ἀνθρώπων ἀδικουμένων εὐλογήσουσί σε.
\vs{4}Ἐγένου γὰρ πάσῃ πόλει ταπεινῇ βοηθὸς, καὶ τοῖς ἀθυμήσασιν διʼ ἔνδειαν σκέπη, ἀπὸ ἀνθρώπων πονηρῶν ῥύσῃ αὐτούς· σκέπη διψώντων, καὶ πνεῦμα ἀνθρώπων ἀδικουμένων.

\vs{5}Ὡς ἄνθρωποι ὀλιγόψυχοι διψῶντες ἐν Σιὼν, ἀπὸ ἀνθρώπων ἀσεβῶν, οἷς ἡμᾶς παρέδωκας.
\vs{6}Καὶ ποιήσει Κύριος σαβαὼθ πᾶσι τοῖς ἔθνεσιν· ἐπὶ τὸ ὄρος τοῦτο πίονται εὐφροσύνην, πίονται οἶνον·
\vs{7}Χρίσονται μύρον ἐν τῷ ὄρει τούτῳ· παράδος ταῦτα πάντα τοῖς ἔθνεσιν· ἡ γὰρ βουλὴ αὕτη ἐπὶ πάντα τὰ ἔθνη.
\vs{8}Κατέπιεν ὁ θάνατος ἰσχύσας, καὶ πάλιν ἀφεῖλε Κύριος ὁ Θεὸς πᾶν δάκρυον ἀπὸ παντὸς προσώπου· τὸ ὄνειδος τοῦ λαοῦ ἀφεῖλεν ἀπὸ πάσης τῆς γῆς, τὸ γὰρ στόμα Κυρίου ἐλάλησε.
\vs{9}Καὶ ἐροῦσι τῇ ἡμέρᾳ ἐκείνῃ, ἰδοὺ ὁ Θεὸς ἡμῶν ἐφʼ ᾧ ἠλπίζομεν, καὶ σώσει ἡμᾶς· οὗτος Κύριος, ὑπεμείναμεν αὐτῷ, καὶ ἠγαλλιώμεθα καὶ εὐφρανθησόμεθα ἐπὶ τῇ σωτηρίᾳ ἡμῶν.

\vs{10}Ἀνάπαυσιν δώσει ὁ Θεὸς ἐπὶ τὸ ὄρος τοῦτο, καὶ καταπατηθήσεται ἡ Μωαβίτις, ὃν τρόπον πατοῦσιν ἅλωνα ἐν ἁμάξαις.
\vs{11}Καὶ ἀνήσει τὰς χεῖρας αὐτοῦ, ὃν τρόπον καὶ αὐτὸς ἐταπείνωσε τοῦ ἀπολέσαι· καὶ ταπεινώσει τὴν ὕβριν αὐτοῦ, ἐφʼ ἃ τὰς χεῖρας ἐπέβαλε.
\vs{12}Καὶ τὸ ὕψος τῆς καταφυγῆς τοῦ τοίχου ταπεινώσει, καὶ καταβήσεται ἕως τοῦ ἐδάφους.

\ch{26}
Τῇ ἡμέρᾳ ἐκείνῃ ᾂσονται τὸ ᾆσμα τοῦτο ἐπὶ γῆς τῆς Ἰουδαίας· ἰδοὺ πόλις ἰσχυρά, καὶ σωτήριον θήσει τὸ τεῖχος, καὶ περίτειχος.
\vs{2}Ἀνοίξατε πύλας, εἰσελθέτω λαὸς φυλάσσων δικαιοσύνην, καὶ φυλάσσων ἀλήθειαν,
\vs{3}ἀντιλαμβανόμενος ἀληθείας, καὶ φυλάσσων εἰρήνην· ὅτι ἐπὶ σοὶ ἐλπίδι
\vs{4}ἤλπισαν Κύριε ἕως τοῦ αἰῶνος, ὁ Θεὸς ὁ μέγας, ὁ αἰώνιος,
\vs{5}ὃς ταπεινώσας κατήγαγες τοὺς ἐνοικοῦντας ἐν ὑψηλοῖς· πόλεις ὀχυρὰς καταβαλεῖς, καὶ κατάξεις ἕως ἐδάφους.
\vs{6}Καὶ πατήσουσιν αὐτοὺς πόδες πρᾳέων καὶ ταπεινῶν.

\vs{7}Ὁδὸς εὐσεβῶν εὐθεῖα ἐγένετο, ἡ ὁδὸς τῶν εὐσεβῶν καὶ παρεσκευασμένη.
\vs{8}Ἡ γὰρ ὁδὸς Κυρίου κρίσις· ἠλπίσαμεν ἐπὶ τῷ ὀνόματί σου, καὶ ἐπὶ τῇ μνείᾳ
\vs{9}ᾗ ἐπιθυμεῖ ἡ ψυχὴ ἡμῶν· ἐκ νυκτὸς ὀρθρίζει τὸ πνεῦμά μου πρὸς σὲ ὁ Θεὸς, διότι φῶς τὰ προστάγματά σου ἐπὶ τῆς γῆς· δικαιοσύνην μάθετε οἱ ἐνοικοῦντες ἐπὶ τῆς γῆς.
\vs{10}Πέπαυται γὰρ ὁ ἀσεβής· πᾶς ὃς οὐ μὴ μάθῃ δικαιοσύνην ἐπὶ τῆς γῆς, ἀλήθειαν οὐ μὴ ποιήσει· ἀρθήτω ὁ ἀσεβὴς, ἵνα μὴ ἴδῃ τὴν δόξαν Κυρίου.
\vs{11}Κύριε ὑψηλός σου ὁ βραχίων, καὶ οὐκ ᾔδεισαν, γνόντες δὲ αἰσχυνθήσονται· ζῆλος λήψεται λαὸν ἀπαίδευτον, καὶ νῦν πῦρ τοὺς ὑπεναντίους ἔδεται.
\vs{12}Κύριε ὁ Θεὸς ἡμῶν, εἰρήνην δὸς ἡμῖν, πάντα γὰρ ἀπέδωκας ἡμῖν.
\vs{13}Κύριε ὁ Θεὸς ἡμῶν, κτῆσαι ἡμᾶς· Κύριε ἐκτός σου ἄλλον οὐκ οἴδαμεν· τὸ ὄνομά σου ὀνομάζομεν.

\vs{14}Οἱ δὲ νεκροὶ ζωὴν οὐ μὴ ἴδωσιν, οὐδὲ ἰατροὶ οὐ μὴ ἀναστήσουσι· διατοῦτο ἐπήγαγες, καὶ ἀπώλεσας, καὶ ᾖρας πᾶν ἄρσεν αὐτῶν.
\vs{15}Πρόσθες αὐτοῖς κακὰ Κύριε, πρόσθες κακὰ τοῖς ἐνδόξοις τῆς γῆς.

\vs{16}Κύριε, ἐν θλίψει ἐμνήσθην σου, ἐν θλίψει μικρᾷ ἡ παιδεία σου ἡμῖν.
\vs{17}Καὶ ὡς ἡ ὠδίνουσα ἐγγίζει τεκεῖν, ἐπὶ τῇ ὠδῖνι αὐτῆς ἐκέκραξεν, οὕτως ἐγενήθημεν τῷ ἀγαπητῷ σου.
\vs{18}Διὰ τὸν φόβον σου Κύριε ἐν γαστρὶ ἐλάβομεν, καὶ ὠδινήσαμεν, καὶ ἐτέκομεν πνεῦμα σωτηρίας σου, ὃ ἐποιήσαμεν ἐπὶ τῆς γῆς· οὐ πεσούμεθα, ἀλλὰ πεσοῦνται πάντες οἱ ἐνοικοῦντες ἐπὶ τῆς γῆς.
\vs{19}Ἀναστήσονται οἱ νεκροὶ, καὶ ἐγερθήσονται οἱ ἐν τοῖς μνημείοις, καὶ εὐφρανθήσονται οἱ ἐν τῇ γῇ· ἡ γὰρ δρόσος ἡ παρὰ σοῦ ἴαμα αὐτοῖς ἐστιν, ἡ δὲ γῆ τῶν ἀσεβῶν πεσεῖται.
\vs{20}Βάδιζε λαός μου, εἴσελθε εἰς τὰ ταμεῖά σου, ἀπόκλεισον τὴν θύραν σου, ἀποκρύβηθι μικρὸν ὅσον ὅσον, ἕως ἂν παρέλθῃ ἡ ὀργὴ Κυρίου.
\vs{21}Ἰδοὺ γὰρ Κύριος ἀπὸ τοῦ ἁγίου ἐπάγει τὴν ὀργὴν ἐπὶ τοὺς ἐνοικοῦντασἐ πὶ τῆς γῆς· καὶ ἀνακαλύψει ἡ γῆ τὸ αἷμα αὐτῆς, καὶ οὐ κατακαλύψει τοὺς ἀνῃρημένους.

\ch{27}
Ἐν τῇ ἡμέρᾳ ἐκείνῃ ἐπάξει ὁ Θεὸς τὴν μάχαιραν τὴν ἁγίαν, καὶ τὴν μεγάλην, καὶ τὴν ἰσχυρὰν ἐπὶ τὸν δράκοντα ὄφιν φεύγοντα, ἐπὶ τὸν δράκοντα ὄφιν σκολιόν· ἀνελεῖ τὸν δράκοντα.
\vs{2}Τῇ ἡμέρᾳ ἐκείνῃ ἀμπελὼν καλός, ἐπιθύμημα ἐξάρχειν κατʼ αὐτῆς.
\vs{3}Ἐγὼ πόλις ὀχυρά, πόλις πολιορκουμένη, μάτην ποτιῶ αὐτήν· ἁλώσεται γὰρ νυκτὸς, ἡμέρας δὲ πεσεῖται τεῖχος.
\vs{4}Οὐκ ἔστιν, ἣ οὐκ ἐπελάβετο αὐτῆς· τίς με θήσει φυλάσσειν καλάμην ἐν ἀγρῷ; διὰ τὴν πολεμίαν ταύτην ἠθέτηκα αὐτήν· τοίνυν διατοῦτο ἐποίησε Κύριος πάντα, ὅσα συνέταξε· κατακέκαυμαι
\vs{5}Βοήσονται οἱ ἐνοικοῦντες ἐν αὐτῇ, ποιήσωμεν εἰρήνην αὐτῷ, ποιήσωμεν εἰρήνην,
\vs{6}οἱ ἐρχόμενοι τέκνα Ἰακώβ· βλαστήσει καὶ ἐξανθήσει Ἰσραὴλ, καὶ ἐμπλησθήσεται ἡ οἰκουμένη τοῦ καρποῦ αὐτοῦ.

\vs{7}Μὴ ὡς αὐτὸς ἐπάταξε, καὶ αὐτὸς οὕτως πληγήσεται; καὶ ὡς αὐτὸς ἀνεῖλεν, οὕτως ἀναιρεθήσεται;
\vs{8}Μαχόμενος καὶ ὀνειδίζων ἐξαποστελεῖ αὐτούς· οὐ σὺ ἦσθα μελετῶν τῷ πνεύματι τῷ σκληρῷ, ἀνελεῖν αὐτοὺς πνεύματι θυμοῦ;
\vs{9}Διατοῦτο ἀφαιρεθήσεται ἀνομία Ἰακὼβ, καὶ τοῦτό ἐστιν ἡ εὐλογία αὐτοῦ, ὅταν ἀφέλωμαι τὴν ἁμαρτίαν αὐτοῦ, ὅταν θῶσι πάντας τοὺς λίθους τῶν βωμῶν κατακεκομμένους, ὡς κονίαν λεπτήν· καὶ οὐ μὴ μείνῃ τὰ δένδρα αὐτῶν, καὶ τὰ εἴδωλα αὐτῶν ἐκκεκομμένα, ὥσπερ δρυμὸς μακράν.
\vs{10}Τὸ κατοικούμενον ποίμνιον ἀνειμένον ἔσται, ὡς ποίμνιον καταλελειμμένον· καὶ ἔσται πολὺν χρόνον εἰς βόσκημα, καὶ ἐκεῖ ἀναπαύσονται ποίμνια.
\vs{11}Καὶ μετὰ χρόνον οὐκ ἔσται ἐν αὐτῇ πᾶν χλωρὸν διὰ τὸ ξηρανθῆναι· γυναῖκες ἐρχόμεναι ἀπὸ θέας δεῦτε· οὐ γὰρ λαός ἐστιν ἔχων σύνεσιν, διατοῦτο οὐ μὴ οἰκτειρήσῃ ὁ ποιήσας αὐτοὺς, οὐδὲ ὁ πλάσας αὐτοὺς οὐ μὴ ἐλεήσῃ.

\vs{12}Καὶ ἔσται ἐν τῇ ἡμέρᾳ ἐκείνῃ, συμφράξει ὁ Θεὸς ἀπὸ τῆς διώρυχος τοῦ ποταμοῦ ἕως Ῥινοκορούρων· ὑμεῖς δὲ συναγάγετε κατὰ ἕνα τοὺς υἱοὺς Ἰσραήλ.
\vs{13}Καὶ ἔσται ἐν τῇ ἡμέρᾳ ἐκείνῃ, σαλπιοῦσι τῇ σάλπιγγι τῇ μεγάλῃ, καὶ ἥξουσιν οἱ ἀπολόμενοι ἐν τῇ χώρᾳ τῶν Ἀσσυρίων, καὶ οἱ ἀπολόμενοι ἐν Αἰγύπτῳ, καὶ προσκυνήσουσι τῷ κυρίῳ ἐπὶ τὸ ὄρος τὸ ἅγιον ἐν Ἱερουσαλήμ.

\ch{28}
Οὐαὶ τῷ στεφάνῳ τῆς ὕβρεως, οἱ μισθωτοὶ Ἐφραῒμ, τὸ ἄνθος τὸ ἐκπεσὸν ἐκ τῆς δόξης ἐπὶ τῆς κορυφῆς τοῦ ὄρους τοῦ παχέος, οἱ μεθύοντες ἄνευ οἴνου.
\vs{2}Ἰδοὺ ἰσχυρὸν καὶ σκληρὸν ὁ θυμὸς Κυρίου, ὡς χάλαζα καταφερομένη οὐκ ἔχουσα σκέπην, βίᾳ καταφερομένη· ὡς ὕδατος πολὺ πλῆθος σῦρον χώραν, τῇ γῇ ποιήσει ἀνάπαυμα· ταῖς χερσὶ,
\vs{3}καὶ τοῖς ποσὶ καταπατηθήσεται ὁ στέφανος τῆς ὕβρεως, οἱ μισθωτοὶ τοῦ Ἐφραΐμ.
\vs{4}Καὶ ἔσται τὸ ἄνθος τὸ ἐκπεσὸν τῆς ἐλπίδος τῆς δόζῆς, ἐπʼ ἄκρου τοῦ ὄρους τοῦ ὑψηλοῦ· ὡς πρόδρομος σύκου, ὁ ἰδὼν αὐτό, πρὶν εἰς τὴν χεῖρα αὐτοῦ λαβεῖν αὐτό, θελήσει αὐτὸ καταπιεῖν.

\vs{5}Τῇ ἡμέρᾳ ἐκείνῃ ἔσται Κύριος σαβαὼθ ὁ στέφανος τῆς ἐλπίδος, ὁ πλεκεὶς τῆς δόξης, τῷ καταλειφθέντι τοῦ λαοῦ.
\vs{6}Καταλειφθήσονται ἐπὶ πνεύματι κρίσεως ἐπὶ κρίσιν, καὶ ἰσχὺν κωλύων ἀνελεῖν.
\vs{7}Οὗτοι γὰρ οἴνῳ πεπλημμελημένοι εἰσίν· ἐπλανήθησαν διὰ τὸ σίκερα, ἱερεὺς καὶ προφήτης ἐξέστησαν διὰ τὸ σίκερα, κατεπόθησαν διὰ τὸν οἶνον, ἐσείσθησαν ἀπὸ τῆς μέθης, ἐπλανήθησαν· τοῦτέστι φάσμα.
\vs{8}Ἀρὰ ἔδεται ταύτην τὴν βουλήν, αὕτη γὰρ ἡ βουλὴ ἕνεκα πλεονεξίας.
\vs{9}Τίνι ἀνηγγείλαμεν κακά, καὶ τίνι ἀνηγγείλαμεν ἀγγελίαν; οἱ ἀπογεγαλακτισμένοι ἀπὸ γάλακτος, οἱ ἀπεσπασμένοι ἀπὸ μαστοῦ.
\vs{10}Θλίψιν ἐπὶ θλίψιν προσδέχου, ἐλπίδα ἐπʼ ἐλπίδι, ἐτι μικρὸν ἔτι μικρόν,
\vs{11}διὰ φαυλισμὸν χειλέων, διὰ γλώσσης ἑτέρας, ὅτι λαλήσουσι τῷ λαῷ τούτῳ,
\vs{12}λέγοντες αὐτοῖς, τοῦτο τὸ ἀνάπαυμα τῷ πεινῶντι, καὶ τοῦτο τὸ σύντριμμα· καὶ οὐκ ἠθέλησαν ἀκούειν.

\vs{13}Καὶ ἔσται αὐτοῖς τὸ λόγιον τοῦ Θεοῦ, θλίψις ἐπὶ θλίψιν, ἐλπὶς ἐπʼ ἐλπίδι, ἔτι μικρὸν ἔτι μικρὸν, ἵνα πορεύσωσι καὶ πέσωσιν ὀπίσω· καὶ συντριβήσονται, καὶ κινδυνεύσουσι, καὶ ἁλώσονται.

\vs{14}Διατοῦτο ἀκούσατε λόγον Κυρίου ἄνδρες τεθλιμμένοι, καὶ οἱ ἄρχοντες τοῦ λαοῦ τούτου τοῦ ἐν Ἰερουσαλήμ·
\vs{15}Ὅτι εἴπατε, ἐποιήσαμεν διαθήκην μετὰ τοῦ ᾅδου, καὶ μετὰ τοῦ θανάτου συνθήκας· καταιγὶς φερομένη ἐὰν παρέλθῃ, οὐ μὴ ἔλθῃ ἐφʼ ἡμᾶς· ἐθήκαμεν ψεῦδος τὴν ἐλπίδα ἡμῶν, καὶ τῷ ψεύδει σκεπασθησόμεθα.
\vs{16}Διατοῦτο οὕτω λέγει κύριος Κύριος,

Ἰδοὺ ἐγὼ ἐμβάλλω εἰς τὰ θεμέλια Σιὼν λίθον πολυτελῆ, ἐκλεκτὸν, ἀκρογωνιαῖον, ἔντιμον, εἰς τὰ θεμέλια αὐτῆς, καὶ ὁ πιστεύων οὐ μὴ καταισχυνθῇ.
\vs{17}Καὶ θήσω κρίσιν εἰς ἐλπίδα, ἡ δὲ ἐλεημοσύνη μου εἰς σταθμούς, καὶ οἱ πεποιθότες μάτην ψεύδει· ὅτι οὐ μὴ παρέλθῃ ὑμᾶς καταιγίς,
\vs{18}μὴ καὶ ἀφέλῃ ὑμῶν τὴν διαθήκην τοῦ θανάτου, καὶ ἡ ἐλπὶς ὑμῶν ἡ πρὸς τὸν ᾅδην οὐ μὴ ἐμμείνῃ· καταιγὶς φερομένη ἐὰν ἐπέλθῃ, ἔσεσθε αὐτῇ εἰς καταπάτημα.
\vs{19}Ὅταν παρέλθῃ, λήμψεται ὑμᾶς, πρωῒ πρωῒ παρελεύσεται ἡμέρας, καὶ ἐν νυκτὶ ἔσται ἐλπὶς πονηρά.

\vs{20}Μάθετε ἀκούειν στενοχωρούμενοι· Οὐ δυνάμεθα μάχεσθαι, αὐτοὶ δὲ ἀσθενοῦμεν τοῦ ὑμᾶς συναχθῆναι.
\vs{21}Ὥσπερ ὄρος ἀσεβῶν ἀναστήσεται Κύριος, καὶ ἔσται ἐν τῇ φάραγγι Γαβαὼν, μετὰ θυμοῦ ποιήσει τὰ ἔργα αὐτοῦ, πικρίας ἔργον· ὁ δὲ θυμὸς αὐτοῦ ἀλλοτρίως χρήσεται, καὶ ἡ σαπρία αὐτοῦ ἀλλοτρία.
\vs{22}Καὶ ὑμεῖς μὴ εὐφρανθείητε, μηδὲ ἰσχυσάτωσαν ὑμῶν οἱ δεσμοί· διότι συντετελεσμένα καὶ συντετμημένα πράγματα ἤκουσα παρὰ Κυρίου σαβαὼθ, ἃ ποιήσει ἐπὶ πᾶσαν τὴν γῆν.

\vs{23}Ἐνωτίζεσθε καὶ ἀκούετε τῆς φωνῆς μου, προσέχετε καὶ ἀκούετε τοὺς λόγους μου.
\vs{24}Μὴ ὅλην τὴν ἡμέραν ἀροτριάσει ὁ ἀροτριῶν; ἢ σπόρον προετοιμάσει, πρὶν ἐργάσασθαι τὴν γῆν;
\vs{25}Οὐχ ὅταν ὁμαλίσῃ τὸ πρόσωπον αὐτῆς, τότε σπείρει μικρὸν μελάνθιον ἢ κύμινον, καὶ πάλιν σπείρει πυρὸν, καὶ κριθὴν, καὶ κέγχρον καὶ ζέαν ἐν τοῖς ὁρίοις σου;
\vs{26}Καὶ παιδευθήσῃ κρίματι Θεοῦ σου, καὶ εὐφρανθήσῃ.
\vs{27}Οὐ γὰρ μετὰ σκληρότητος καθαίρεται τὸ μελάνθιον, οὐδὲ τροχὸς ἁμάξης περιάξει ἐπὶ τὸ κύμινον· ἀλλὰ ῥάβδῳ τινάσσεται τὸ μελάνθιον, τὸ δὲ κύμινον
\vs{28}μετὰ ἄρτου βρωθήσεται· οὐ γὰρ εἰς τὸν αἰῶνα ἐγώ εἰμι ὑμῖν ὀργισθήσομαι, οὐδὲ φωνὴ τῆς πικρίας μου καταπατήσει ὑμᾶς.
\vs{29}Καὶ ταῦτα παρὰ Κυρίου σαβαὼθ ἐξῆλθε τὰ τέρατα· βουλεύσασθε, ὑψώσατε ματαίαν παράκλησιν.

\ch{29}
Οὐαὶ Ἀριὴλ πόλις, ἣν ἐπολέμησε Δαυείδ· συναγάγετε γενήματα ἐνιαυτὸν ἐπὶ ἐνιαυτὸν, φάγεσθε, φάγεσθε γὰρ σὺν Μωάβ,
\vs{2}ἐκθλίψω γὰρ Ἀριήλ· καὶ ἔσται αὐτῆς ἡ ἰσχὺς καὶ ὁ πλοῦτος ἐμοί.
\vs{3}Καὶ κυκλώσω ὡς Δαυὶδ ἐπὶ σὲ, καὶ βαλῶ περὶ σὲ χάρακα, καὶ θήσω περὶ σὲ πύργους,
\vs{4}καὶ ταπεινωθήσονται εἰς τὴν γῆν οἱ λόγοι σου, καὶ εἰς τὴν γῆν οἱ λόγοι σου δύσονται· καὶ ἔσονται ὡς οἱ φωνοῦντες ἐκ τῆς γῆς ἡ φωνή σου, καὶ πρὸς τὸ ἔδαφος ἡ φωνή σου ἀσθενήσει.
\vs{5}Καὶ ἔσται ὡς κονιορτὸς ἀπὸ τροχοῦ ὁ πλοῦτος τῶν ἀσεβῶν, καὶ ὡς χνοῦς φερόμενος τὸ πλῆθος τῶν καταδυναστευόντων σε, καὶ ἔσται ὡς στιγμὴ παραχρῆμα
\vs{6}παρὰ Κυρίου σαβαώθ· ἐπισκοπὴ γὰρ ἔσται μετὰ βροντῆς καὶ σεισμοῦ καὶ φωνὴς μεγάλης, καταιγὶς φερομένη, καὶ φλὸξ πυρὸς κατεσθίουσα.
\vs{7}Καὶ ἔσται ὡς ἐνυπνιαζόμενος καθʼ ὕπνους νυκτὸς, ὁ πλοῦτος ἁπάντων τῶν ἐθνῶν, ὅσοι ἐπεστράτευσαν ἐπὶ Ἀριὴλ, καὶ πάντες οἱ στρατευόμενοι ἐπὶ Ἰερουσαλὴμ, καὶ πάντες οἱ συνηγμένοι ἐπʼ αὐτὴν, καὶ οἱ θλίβοντες αὐτήν.
\vs{8}Καὶ ὡς οἱ ἐν τῷ ὕπνῳ πίνοντες καὶ ἔσθοντες, καὶ ἐξαναστάντων, μάταιον τὸ ἐνύπνιον· καὶ ὃν τρόπον ἐνυπνιάζεται ὁ διψῶν ὡς ὁ πίνων, καὶ ἐξαναστὰς ἔτι διψᾷ, ἡ δὲ ψυχὴ αὐτοῦ εἰς κενὸν ἤλπισεν· οὕτως ἔσται ὁ πλοῦτος τῶν ἐθνῶν πάντων, ὅσοι ἐπεστράτευσαν ἐπὶ τὸ ὄρος Σιών.

\vs{9}Ἐκλύθητε καὶ ἔκστητε, καὶ κραιπαλήσατε οὐκ ἀπὸ σίκερα οὐδὲ ἀπὸ οἴνου.
\vs{10}ὅτι πεπότικεν ὑμᾶς Κύριος πνεύματι κατανύξεως, καὶ καμμύσει τοὺς ὀφθαλμοὺς αὐτῶν, καὶ τῶν προφητῶν αὐτῶν, καὶ τῶν ἀρχόντων αὐτῶν, οἱ ὁρῶντες τὰ κρυπτά.
\vs{11}Καὶ ἔσται ὑμῖν τὰ ῥήματα πάντα ταῦτα, ὡς οἱ λόγοι τοῦ βιβλίου τοῦ ἐσφραγισμένου τσύτου, ὃ ἐὰν δῶσιν αὐτὸ ἀνθρώπῳ ἐπισταμένῳ γράμματα, λέγοντες, ἀνάγνωθι ταῦτα· καὶ ἐρεῖ, οὐ δύναμαι ἀναγνῶναι, ἐσφράγισται γάρ.
\vs{12}καὶ δοθήσεται τὸ βιβλίον τοῦτο εἰς χεῖρας ἀνθρώπου μὴ ἐπισταμένου γράμματα, καὶ ἐρεῖ αὐτῷ, ἀνάγνωθι τοῦτο· καὶ ἐρεῖ, οὐκ ἐπίσταμαι γράμματα.

\vs{13}Καὶ εἶπε Κύριος, ἐγγίζει μοι ὁ λαὸς οὗτος ἐν τῷ στόματι αὐτοῦ, καὶ ἐν τοῖς χείλεσιν αὐτῶν τιμῶσί με, ἡ δὲ καρδία αὐτῶν πόῤῥω ἀπέχει ἀπʼ ἐμοῦ· μάτην δὲ σέβονταί με, διδάσκοντες ἐντάλματα ἀνθρώπων καὶ διδασκαλίας.
\vs{14}Διατοῦτο ἰδοὺ προσθήσω τοῦ μεταθεῖναι τὸν λαὸν τοῦτον, καὶ μεταθήσω αὐτούς, καὶ ἀπολῶ τὴν σοφίαν τῶν σοφῶν, καὶ τὴν σύνεσιν τῶν συνετῶν κρύψω.
\vs{15}Οὐαὶ οἱ βαθέως βουλὴν ποιοῦντες, καὶ οὐ διὰ Κυπίυο· οὐαὶ οἱ ἐν κρυφῇ βουλὴν ποιοῦντες, καὶ ἔσται ἐν σκότει τὰ ἔργα αὐτῶν· καὶ ἐροῦσι, τίς ἑόρακεν ἡμᾶς; καὶ τίς ἡμᾶς γνώσεται, ἢ ἃ ἡμεῖς ποιοῦμεν;
\vs{16}Οὐχ ὡς πηλὸς τοῦ κεραμέως λογισθήσεσθε; μὴ ἐρεῖ τὸ πλάσμα τῷ πλάσαντι αὐτό, οὐ σύ με ἔπλασας; ἢ τὸ ποίημα τῷ ποιήσαντι, οὐ συνετῶς με ἐποίησας;
\vs{17}Οὐκέτι μικρὸν καὶ μετατεθήσεται ὁ Λίβανος, ὡς τὸ ὄρος τὸ Χέρμελ, καὶ τὸ Χέρμελ εἰς δρυμὸν λογισθήσεται;
\vs{18}Καὶ ἀκούσονται ἐν τῇ ἡμέρᾳ ἐκείνῃ κωφοὶ λόγους βιβλίου, καὶ οἱ ἐν τῷ σκότει, καὶ οἱ ἐν τῇ ὁμίχλῃ· ὀφθαλμοὶ τυφλῶν ὄψονται,
\vs{19}καὶ ἀγαλλιάσονται πτωχοὶ διὰ Κύριον ἐν εὐφροσύνῃ, καὶ οἱ ἀπηλπισμένοι τῶν ἀνθρώπων ἐμπλησθήσονται εὐφροσύνης.
\vs{20}Ἐξέλιπεν ἄνομος, καὶ ἀπώλετο ὑπερήφανος, καὶ ἐξωλοθρεύθησαν οἱ ἀνομοῦντες ἐπὶ κακίᾳ,
\vs{21}καὶ οἱ ποιοῦντες ἁμαρτεῖν ἀνθρώπους ἐν λόγῳ· πάντας δὲ τοὺς ἐλέγχοντας ἐν πύλαις πρόσκομμα θήσουσιν, ὅτι ἐπλαγίασαν ἐπʼ ἀδίκοις δίκαιον.

\vs{22}Διατοῦτο τάδε λέγει Κύριος ἐπὶ τὸν οἶκον Ἰακώβ, ὃν ἀφώρισεν ἐξ Ἀβραὰμ, οὐ νῦν αἰσχυνθήσεται Ἰακώβ, οὐδὲ νῦν τὸ πρόσωπον μεταβαλεῖ.
\vs{23}Ἀλλʼ ὅταν ἴδωσιν τὰ τέκνα αὐτῶν τὰ ἔργα μου, διʼ ἐμὲ ἁγιάσουσι τὸ ὄνομά μου, καὶ ἁγιάσουσι τὸν ἅγιον Ἰακὼβ, καὶ τὸν Θεὸν τοῦ Ἰσραὴλ φοβηθήσονται.
\vs{24}Καὶ γνώσονται οἱ πλανώμενοι τῷ πνεύματι σύνεσιν, οἱ δὲ γογγύζοντες μαθήσονται ὑπακούειν, καὶ αἱ γλῶσσαι αἱ ψελλίζουσαι μαθήσονται λαλεῖν εἰρήνην.

\ch{30}
Οὐαὶ τέκνα ἀποστάται, λέγει Κύριος· ἐποιήσατε βουλὴν οὐ διʼ ἐμοῦ, καὶ συνθήκας οὐ διὰ τοῦ πνεύματός μου, προσθεῖναι ἁμαρτίας ἐφʼ ἁμαρτίας,
\vs{2}οἱ πορευόμενοι καταβῆναι εἰς Αἴγυπτον, ἐμὲ δὲ οὐκ ἐπερώτησαν τοῦ βοηθηθῆναι ὑπὸ Φαραὼ, καὶ σκεπασθῆναι ὑπὸ Αἰγυπτίων.
\vs{3}Ἔσται γὰρ ὑμῖν σκέπη Φαραὼ εἰς αἰσχύνην, καὶ τοῖς πεποιθόσιν ἐπʼ Αἴγυπτον ὄνειδος.
\vs{4}Ὅτι εἰσὶν ἐν Τάνει ἀρχηγοὶ ἄγγελοι πονηροί.
\vs{5}Μάτην κοπιάσουσι πρὸς λαὸν, ὃς οὐκ ὠφελήσει αὐτοὺς εἰς βοήθειαν, ἀλλὰ εἰς αἰσχύνην καὶ ὄνειδος.

Ἡ ὍΡΑΣΙΣ ΤΩΝ ΤΕΤΡΑΠΟΔΩΝ ΤΩΝ ἘΝ ΤΗ ἘΡΗΜΩ.

\vs{6}Ἐν τῇ θλίψει καὶ τῇ στενοχωρίᾳ, λέων καὶ σκύμνος λέοντος, ἐκεῖθεν καὶ ἀσπίδες, καὶ ἔκγονα ἀσπίδων πετομένων, οἳ ἔφερον ἐπὶ ὄνων καὶ καμήλων τὸν πλοῦτον αὐτῶν πρὸς ἔθνος, ὃ οὐκ ὠφελήσει αὐτούς.
\vs{7}Αἰγύπτιοι μάταια καὶ κενὰ ὠφελήσουσιν ὑμᾶς· ἀπάγγειλον αὐτοῖς, ὅτι ματαία ἡ παράκλησις ὑμῶν αὕτη.

\vs{8}Νῦν οὖν καθίσας γράψον ἐπὶ πυξίου ταῦτα καὶ εἰς βιβλίον, ὅτι ἔσται εἰς ἡμέρας ταῦτα καιρῷ, καὶ ἕως εἰς τὸν αἰῶνα.
\vs{9}Ὅτι ὁ λαὸς ἀπειθής ἐστιν, υἱοὶ ψευδεῖς, οἳ οὐκ ἠβούλοντο ἀκούειν τὸν νόμον τοῦ Θεοῦ·
\vs{10}Οἱ λέγοντες τοῖς προφήταις, μὴ ἀναγγέλλετε ἡμῖν, καὶ τοῖς τὰ ὁράματα ὁρῶσι, μὴ λαλεῖτε ἡμῖν, ἀλλὰ ἡμῖν λαλεῖτε καὶ ἀναγγέλλετε ἡμῖν ἑτέραν πλάνησιν,
\vs{11}καὶ ἀποστρέψατε ἡμᾶς ἀπὸ τῆς ὁδοῦ ταύτης· ἀφέλετε ἀφʼ ἡμῶν τὸν τρίβον τοῦτον, καὶ ἀφέλετε ἀφʼ ἡμῶν τὸ λόγιον τοῦ Ἰσραήλ.

\vs{12}Διατοῦτο τάδε λέγει ὁ ἅγιος τοῦ Ἰσραήλ, ὅτι ἠπειθήσατε τοῖς λόγοις τούτοις, καὶ ἠλπίσατε ἐπὶ ψεύδει, καὶ ὅτι ἐγόγγυσας, καὶ πεποιθὼς ἐγένου ἐπὶ τῷ λόγῳ τούτῳ,
\vs{13}διατοῦτο ἔσται ὑμῖν ἡ ἁμαρτία αὕτη, ὡς τεῖχος πίπτον παραχρῆμα πόλεως ὀχυρᾶς ἑαλωκυίας, ἧς παραχρῆμα πάρεστι τὸ πτῶμα·
\vs{14}Καὶ τὸ πτῶμα αὐτῆς ἔσται ὡς σύντριμμα ἀγγείου ὀστρακίνου, ἐκ κεραμίου λεπτὰ, ὥστε μὴ εὑρεῖν ἐν αὐτοῖς ὄστρακον, ἐν ᾧ πῦρ ἀρεῖς, καὶ ἐν ᾧ ἀποσυριεῖς ὕδωρ μικρόν.

\vs{15}Οὕτω λέγει Κύριος, Κύριος ὁ ἅγιος τοῦ Ἰσραήλ, ὅταν ἀποστραφεὶς στενάξῃς, τότε σωθήσῃ, καὶ γνώσῃ ποῦ ἦσθα, ὅτε ἐπεποίθεις ἐπὶ τοῖς ματαίοις, ματαία ἡ ἰσχὺς ὑμῶν ἐγενήθη· καὶ οὐκ ἠβούλεσθε ἀκούειν,
\vs{16}ἀλλʼ εἴπατε, ἐφʼ ἵππων φευξόμεθα· διατοῦτο φεύξεσθε· καὶ ἐπὶ κούφοις ἀναβάταις ἐσόμεθα· διατοῦτο κοῦφοι ἔσονται οἱ διώκοντες ὑμᾶς.
\vs{17}Χίλιοι διὰ φωνὴν ἑνὸς φεύξονται, καὶ διὰ φωνὴν πέντε φεύξονται πολλοί, ἕως ἂν καταλειφθῆτε ὡς ἱστὸς ἐπʼ ὄρους, καὶ ὡς σημαῖαν φέρων ἐπὶ βουνοῦ.

\vs{18}Καὶ πάλιν μενεῖ ὁ Θεὸς τοῦ οἰκτειρῆσαι ὑμᾶς, καὶ διατοῦτο ὑψωθήσεται τοῦ ἐλεῆσαι ὑμᾶς· διότι κριτὴς Κύριος ὁ Θεὸς ὑμῶν· μακάριοι οἱ ἐμμένοντες ἐπʼ αὐτῷ.

\vs{19}Διότι λαὸς ἅγιος ἐν Σιὼν οἰκήσει· καὶ Ἰερουσαλὴμ κλαυθμῷ ἔκλαυσεν, ἐλέησόν με· ἐλεήσει σε, τὴν φωνὴν τῆς κραυγῆς σου ἡνίκα εἶδεν, ἐπήκουσέ σου.
\vs{20}Καὶ δώσει Κύριος ὑμῖν ἄρτον θλίψεως, καὶ ὕδωρ στενὸν, καὶ οὐκ ἔτι μὴ ἐγγίσωσί σοι οἱ πλανῶντές σε· ὅτι οἱ ὀφθαλμοί σου ὄψονται τοὺς πλανῶντάς σε,
\vs{21}καὶ τὰ ὦτά σου ἀκούσονται τοὺς λόγους τῶν ὀπίσω σε πλανησάντων, οἱ λέγοντες, αὕτη ἡ ὁδὸς, πορευθῶμεν ἐν αὐτῇ, εἴτε δεξιὰ εἴτε ἀριστερά.
\vs{22}Καὶ μιανεῖς τὰ εἴδωλα τὰ περιηργυρωμένα, καὶ περικεχρυσωμένα λεπτὰ ποιήσῃς, καὶ λικμήσῃς ὡς ὕδωρ ἀποκαθημένης, καὶ ὡς κόπρον ὤσεις αὐτά.
\vs{23}Τότε ἔσται ὁ ὑετὸς τῷ σπέρματι τῆς γῆς σου, καὶ ὁ ἄρτος τοῦ γενήματος τῆς γῆς σου ἔσται πλησμονὴ καὶ λιπαρός· καὶ βοσκηθήσεταί σου τὰ κτήνη τῇ ἡμέρᾳ ἐκείνῃ τόπον πίονα καὶ εὐρύχωρον.
\vs{24}Οἱ ταῦροι ὑμῶν καὶ οἱ βόες οἱ ἐργαζόμενοι τὴν γῆν, φάγονται ἄχυρα ἀναπεποιημένα ἐν κριθῇ λελικμημένῃ.
\vs{25}Καὶ ἔσται ἐπὶ παντὸς ὄρους ὑψηλοῦ, καὶ ἐπὶ παντὸς βουνοῦ μετεώρου ὕδωρ διαπορευόμενον ἐν τῇ ἡμέρᾳ ἐκείνῃ, ὅταν ἀπόλωνται πολλοὶ, καὶ ὅταν πέσωσι πύργοι.
\vs{26}Καὶ ἔσται τὸ φῶς τῆς σελήνης ὡς τὸ φῶς τοῦ ἡλίου, καὶ τὸ φῶς τοῦ ἡλίου ἔσται ἑπταπλάσιον ἐν τῇ ἡμέρα, ὅταν ἰάσηται Κύριος τὸ σύντριμμα τοῦ λαοῦ αὐτοῦ, καὶ τὴν ὀδύνην τῆς πληγῆς σου ἰάσεται.

\vs{27}Ἰδοὺ τὸ ὄνομα Κυρίου ἔρχεται διὰ χρόνου, καιόμενος θυμός· μετὰ δόξης τὸ λόγιον τῶν χειλέων αὐτοῦ, λόγιον ὀργῆς πλῆρες, καὶ ἡ ὀργὴ τοῦ θυμοῦ ὡς πῦρ ἔδεται.
\vs{28}Καὶ τὸ πνεῦμα αὐτοῦ ὡς ὕδωρ ἐν φάραγγι σὺρον, ἥξει ἕως τοῦ τραχήλου, καὶ διαιρεθήσεται, τοῦ ταράξαι ἔθνη ἐπὶ πλανήσει ματαίᾳ, καὶ διώξεται αὐτοὺς πλάνησις, καὶ λήψεται αὐτοὺς κατὰ πρόσωπον αὐτῶν.
\vs{29}Μὴ διαπαντὸς δεῖ ὑμᾶς εὐφραίνεσθαι, καὶ εἰσπορεύεσθαι εἰς τὰ ἅγιά μου διαπαντὸς, ὡσεὶ ἑορτάζοντας, καὶ ὡσεὶ εὐφραινομένους εἰσελθεῖν μετὰ αὐλοῦ εἰς τὸ ὄρος Κυρίου πρὸς τὸν Θεὸν τοῦ Ἰσραήλ;
\vs{30}Καὶ ἀκουστὴν ποιήσει Κύριος τὴν δόξαν τῆς φωνῆς αὐτοῦ, καὶ τὸν θυμὸν τοῦ βραχίονος αὐτοῦ, δεῖξαι μετὰ θυμοῦ καὶ ὀργῆς, καὶ φλογὸς κατεσθιούσης, κεραυνώσει βιαίως, καὶ ὡς ὕδωρ καὶ χάλαζα συγκαταφερομένη βίᾳ.
\vs{31}Διὰ γὰρ τῆς φωνῆς Κυρίου ἡττηθήσονται Ἀσσύριοι, τῇ πληγῇ ᾗ ἂν πατάξῃ αὐτούς.
\vs{32}Καὶ ἔσται αὐτῷ κυκλόθεν, ὅθεν ἦν αὐτῶν ἡ ἐλπὶς τῆς βοηθείας, ἐφʼ ᾗ αὐτὸς ἐπεποίθει, αὐτοὶ μετὰ τυμπάνων καὶ κιθάρας πολεμήσουσιν αὐτὸν ἐκ μεταβολῆς.
\vs{33}Σὺ γὰρ πρὸ ἡμερῶν ἀπαιτηθήσῃ· μὴ καί σοὶ ἡτοιμάσθη βασιλεύειν; φάραγγα βαθείαν, ξύλα κείμενα, πῦρ καὶ ξύλα πολλὰ, ὁ θυμὸς Κυρίου ὡς φάραγξ ὑπὸ θείου καιομένη.

\ch{31}
Οὐαὶ οἱ καταβαίνοντες εἰς Αἴγυπτον ἐπὶ βοήθειαν, οἱ ἐφʼ ἵπποις πεποιθότες καὶ ἐφʼ ἅρμασιν, ἔστι γὰρ πολλὰ, καὶ ἐφʼ ἵπποις πλῆθος σφόδρα· καὶ οὐκ ἦσαν πεποιθότες ἐπὶ τὸν ἅγιον τοῦ Ἰσραὴλ, καὶ τὸν κύριον οὐκ ἐζήτησαν·
\vs{2}Καὶ αὐτὸς σοφῶς ἦγεν ἐπʼ αὐτοὺς κακά, καὶ ὁ λόγος αὐτοῦ οὐ μὴ ἀθετηθῇ, καὶ ἐπαναστήσεται ἐπʼ οἴκους ἀνθρώπων πονηρῶν, καὶ ἐπὶ τὴν ἐλπίδα αὐτῶν τὴν ματαίαν,
\vs{3}Αἰγύπτιον, ἄνθρωπον καὶ οὐ Θεὸν, ἵππων σάρκας, καὶ οὐκ ἔστι βοήθεια· ὁ δὲ Κύριος ἐπάξει τὴν χεῖρα αὐτοῦ ἐπʼ αὐτούς· καὶ κοπιάσουσιν οἱ βοηθοῦντες, καὶ ἅμα πάντες ἀπολοῦνται.
\vs{4}Ὅτι οὕτως εἶπέ μοι Κύριος, ὃν τρόπον βοήσῃ ὁ λέων, ἢ ὁ σκύμνος ἐπὶ τῇ θήρᾳ ᾗ ἔλαβε, καὶ κεκράξῃ ἐπʼ αὐτῇ, ἕως ἂν ἐμπλησθῇ τὰ ὄρη τῆς φωνῆς αὐτοῦ, καὶ ἡττήθησαν, καὶ τὸ πλῆθος τοῦ θυμοῦ ἐπτοήθησαν, οὕτως καταβήσεται Κύριος σαβαὼθ ἐπιστρατεῦσαι ἐπὶ τὸ ὄρος τὸ Σιὼν, ἐπὶ τὰ ὄρη αὐτῆς.
\vs{5}Ὡς ὄρνεα πετόμενα, οὕτως ὑπερασπιεῖ Κύριος σαβαὼθ, ὑπὲρ Ἱερουσαλὴμ ὑπερασπιεῖ, καὶ ἐξελεῖται, καὶ περιποιήσεται, καὶ σώσει.
\vs{6}Ἐπιστράφητε οἱ τὴν βαθείαν βουλὴν βουλευόμενοι καὶ ἄνομον, υἱοὶ Ἰσραήλ.
\vs{7}Ὅτι τῇ ἡμέρᾳ ἐκείνῃ ἀπαρνήσονται οἱ ἄνθρωποι τὰ χειροποίητα αὐτῶν τὰ ἀργυρᾶ, καὶ τὰ χειροποίητα τὰ χρυσᾶ, ἃ ἐποίησαν αἱ χεῖρες αὐτῶν.
\vs{8}Καὶ πεσεῖται Ασσούρ· οὐ μάχαιρα ἀνδρὸς, οὐδὲ μάχαιρα ἀνθρώπου καταφάγεται αὐτὸν, καὶ φεύξεται οὐκ ἀπὸ προσώπου μαχαίρας· οἱ δὲ νεανίσκοι ἔσονται εἰς ἥττημα,
\vs{9}πέτρᾳ γὰρ περιληφθήσονται ὡς χάρακι, καὶ ἡττηθήσονται, ὁ δὲ φεύγων ἁλώσεται· τάδε λέγει Κύριος, μακάριος ὃς ἔχει ἐν Σιὼν σπέρμα, καὶ οἰκείους ἐν Ἱερουσαλήμ.

\ch{32}
Ἰδοὺ γὰρ βασιλεὺς δίκαιος βασιλεύσει, καὶ ἄρχοντες μετὰ κρίσεως ἄρξουσι.
\vs{2}Καὶ ἔσται ὁ ἄνθρωπος κρύπτων τοὺς λόγους αὐτοῦ, καὶ κρυβήσεται, ὡς ἀφʼ ὕδατος φερομένου· καὶ φανήσεται ἐν Σειὼν ὡς ποταμὸς φερόμενος ἔνδοξος ἐν γῇ διψώσῃ.
\vs{3}Καὶ οὐκέτι ἔσονται πεποιθότες ἐπʼ ἀνθρώποις, ἀλλὰ τὰ ὦτα ἀκούειν δώσουσι.
\vs{4}Καὶ ἡ καρδία τῶν ἀσθενούντων προσέξει τῷ ἀκούειν, καὶ αἱ γλῶσσαι αἱ ψελλίζουσαι ταχὺ μαθήσονται λαλεῖν εἰρήνην.
\vs{5}Καὶ οὐκέτι μὴ εἴπωσι τῷ μωρῷ ἄρχειν, καὶ οὐκέτι μὴ εἴπωσιν οἱ ὑπηρέται σου, σίγα.
\vs{6}Ὁ γὰρ μωρὸς μωρὰ λαλήσει, καὶ ἡ καρδία αὐτοῦ μάταια νοήσει, τοῦ συντελεῖν ἄνομα, καὶ λαλεῖν πρὸς Κύριον πλάνησιν, τοῦ διασπεῖραι ψυχὰς πεινώσας, καὶ τὰς ψυχὰς τὰς διψώσας κενὰς ποιήσει.
\vs{7}Ἡ γὰρ βουλὴ τῶν πονηρῶν ἄνομα βουλεύσεται, καταφθεῖραι ταπεινοὺς ἐν λόγοις ἀδίκοις, καὶ διασκεδάσαι λόγους ταπεινῶν ἐν κρίσει.
\vs{8}Οἱ δὲ εὐσεβεῖς συνετὰ ἐβουλεύσαντο, καὶ αὕτη ἡ βουλὴ μενεῖ.

\vs{9}Γυναῖκες πλούσιαι ἀνάστητε, καὶ ἀκούσατε τῆς φωνῆς μου· θυγατέρες ἐν ἐλπίδι εἰσακούσατε λόγους μου.
\vs{10}Ἡμέρας ἐνιαυτοῦ μνείαν ποιήσασθε ἐν ὀδύνῃ μετʼ ἐλπίδος· ἀνήλωται ὁ τρυγητὸς, πέπαυται, οὐκέτι μὴ ἔλθῃ.
\vs{11}Ἔκστητε, λυπήθητε αἱ πεποιθυῖαι, ἐκδύσασθε, γυμναὶ γένεσθε, περιζώσασθε τὰς ὀσφῦας,
\vs{12}καὶ ἐπὶ τῶν μαστῶν κόπτεσθε, ἀπὸ ἀγροῦ ἐπιθυμήματος, καὶ ἀμπέλου γεννήματος.
\vs{13}Ἡ γῆ τοῦ λαοῦ μου, ἄκανθα καὶ χόρτος ἀναβήσεται, καὶ ἐκ πάσης οἰκίας εὐφροσύνη ἀρθήσεται.
\vs{14}Πόλις πλουσία, οἶκοι ἐγκαταλελειμμένοι πλοῦτον πόλεως ἀφήσουσιν, οἴκους ἐπιθυμήματος· καὶ ἔσονται αἱ κῶμαι σπήλαια ἕως τοῦ αἰῶνος, εὐφροσύνη ὄνων ἀγρίων, βοσκήματα ποιμένων,
\vs{15}ἕως ἂν ἔλθῃ ἐφʼ ὑμᾶς πνεῦμα ἀφʼ ὑψηλοῦ· καὶ ἔσται ἔρημος ὁ Χέρμελ, καὶ ὁ Χέρμελ εἰς δρυμὸν λογισθήσεται.
\vs{16}Καὶ ἀναπαύσεται ἐν τῇ ἐρήμῳ κρίμα, καὶ δικαιοσύνη ἐν τῷ Καρμήλῳ κατοικήσει.
\vs{17}Καὶ ἔσται τὰ ἔργα τῆς δικαιοσύνης, εἰρήνη· καὶ κρατήσει ἡ δικαιοσύνη ἀνάπαυσιν, καὶ πεποιθότες ἕως τοῦ αἰῶνος.
\vs{18}Καὶ κατοικήσει ὁ λαὸς αὐτοῦ ἐν πόλει εἰρήνης, καὶ ἐνοικήσει πεποιθὼς, καὶ ἀναπαύσονται μετὰ πλούτου.
\vs{19}Ἡ δὲ χάλαζα ἐὰν καταβῇ, οὐκ ἐφʼ ὑμᾶς ἥξει· καὶ ἔσονται οἱ ἐνοικοῦντες ἐν τοῖς δρυμοῖς πεποιθότες, ὡς οἱ ἐν τῷ πεδινῇ.
\vs{20}Μακάριοι οἱ σπείροντες ἐπὶ πᾶν ὕδωρ, οὗ βοῦς καὶ ὄνος πατεῖ.

\ch{33}
Οὐαὶ τοῖς ταλαιπωροῦσιν ὑμᾶς, ὑμᾶς δὲ οὐδεὶς ποιεῖ ταλαιπώρους, καὶ ὁ ἀθετῶν ὑμᾶς οὐκ ἀθετεῖ· ἁλώσονται οἱ ἀθετοῦντες, καὶ παραδοθήσονται, καὶ ὡς σὴς ἐφʼ ἱματίου, οὕτως ἡττηθήσονται.

\vs{2}Κύριε ἐλέησον ἡμᾶς, ἐπὶ σοὶ γὰρ πεποίθαμεν· ἐγενήθη τὸ σπέρμα τῶν ἀπειθούντων εἰς ἀπώλειαν, ἡ δὲ σωτηρία ἡμῶν ἐν καιρῷ θλίψεως.
\vs{3}Διὰ φωνὴν τοῦ φόβου ἐξέστησαν λαοὶ ἀπὸ τοῦ φόβου σου, καὶ διεσπάρησαν τὰ ἔθνη.

\vs{4}Νῦν δὲ συναχθήσεται τὰ σκῦλα ὑμῶν μικροῦ καὶ μεγάλου· ὃν τρόπον ἐάν τις συναγάγῃ ἀκρίδας, οὕτως ἐμπαίξουσιν ὑμῖν.
\vs{5}Ἅγιος ὁ Θεὸς ὁ κατοικῶν ἐν ὑψηλῷ, ἔνεπλήσθη Σιὼν κρίσεως καὶ δικαιοσύνης.
\vs{6}Ἐν νόμῳ παραδοθήσονται, ἐν θησαυροῖς ἡ σωτηρία ἡμῶν, ἐκεῖ σοφία καὶ ἐπιστήμη καὶ εὐσέβεια πρὸς τὸν κύριον· οὗτοί εἰσι θησαυροὶ δικαιοσύνης.

\vs{7}Ἰδοὺ δὴ ἐν τῷ φόβῳ ὑμῶν οὗτοι φοβηθήσονται· οὓς ἐφοβεῖσθε, βοήσονται ἀφʼ ὑμῶν· ἄγγελοι ἀποσταλήσονται, πικρῶς κλαίοντες, παρακαλοῦντες εἰρήνην.
\vs{8}Ἐρημωθήσονται γὰρ αἱ τούτων ὁδοί· πέπαυται ὁ φόβος τῶν ἐθνῶν, καὶ ἡ πρὸς τούτους διαθήκη αἴρεται, καὶ οὐ μὴ λογίσησθε αὐτοὺς ἀνθρώπους.
\vs{9}Ἐπένθησεν ἡ γῆ, ᾐσχύνθη ὁ Λίβανος, ἕλη ἐγένετο ὁ Σάρων· φανερὰ ἔσται ἡ Γαλιλαία, καὶ ὁ Χέρμελ.

\vs{10}Νῦν ἀναστήσομαι, λέγει Κύριος, νῦν δοξασθήσομαι, νῦν ὑψωθήσομαι.
\vs{11}Νῦν ὄψεσθε, νῦν αἰσθηθήσεσθε, ματαία ἔσται ἡ ἰσχὺς τοῦ πνεύματος ὑμῶν· πῦρ κατέδεται ὑμᾶς.
\vs{12}Καὶ ἔσονται ἔθνη κατακεκαυμένα, ὡς ἄκανθα ἐν ἀγρῷ ἐῤῥιμμένη καὶ κατακεκαυμένη.

\vs{13}Ἀκούσονται οἱ πόῤῥωθεν ἃ ἐποίησα, γνώσονται οἱ ἐγγίζοντες τὴν ἰσχύν μου.
\vs{14}Ἀπέστησαν οἱ ἐν Σιὼν ἄνομοι, λήψεται τρόμος τοὺς ἀσεβεῖς· τίς ἀναγγελεῖ ὑμῖν, ὅτι πῦρ καίεται; τίς ἀναγγελεῖ ὑμῖν τὸν τόπον τὸν αἰώνιον;

\vs{15}Πορευόμενος ἐν δικαιοσύνῃ, λαλῶν εὐθεῖαν ὁδόν, μισῶν ἀνομίαν καὶ ἀδικίαν, καὶ τὰς χεῖρας ἀποσειόμενος ἀπὸ δώρων· βαρύνων τὰ ὦτα, ἵνα μὴ ἀκούσῃ κρίσιν αἵματος· καμμύων τοὺς ὀφθαλμοὺς, ἵνα μὴ ἴδῃ ἀδικίαν,
\vs{16}οὗτος οἰκήσει ἐν ὑψηλῷ σπηλαίῳ πέτρας ἰσχυρᾶς· ἄρτος αὐτῷ δοθήσεται, καὶ τὸ ὕδωρ αὐτοῦ πιστόν.
\vs{17}Βασιλέα μετὰ δόξης ὄψεσθε, οἱ ὀφθαλμοὶ ὑμῶν ὄψονται γῆν πόῤῥωθεν,
\vs{18}ἡ ψυχὴ ἡμῶν μελετήσει φόβον· ποῦ εἰσιν οἱ γραμματικοί; ποῦ εἰσιν οἱ συμβουλεύοντες; ποῦ ἐστιν ὁ ἀριθμῶν τοὺς τρεφομένους
\vs{19}μικρὸν καὶ μέγαν λαόν; ᾧ οὐ συνεβουλεύσατο, οὐσὲ ᾔδει βαθύφωνον, ὥστε μὴ ἀκοῦσαι λαὸς πεφαυλισμένος, καὶ οὐκ ἔστι τῷ ἀκούοντι σύνεσις.

\vs{20}Ἰδοὺ Σιὼν ἡ πόλις, τὸ σωτήριον ἡμῶν, οἱ ὀφθαλμοί σου ὄψονται Ἱερουσαλὴμ, πόλις πλουσία, σκηναὶ αἳ οὐ μὴ σεισθῶσιν, οὐδὲ μὴ κινηθῶσιν οἱ πάσσαλοι τῆς σκηνῆς αὐτῆς εἰς τὸν αἰῶνα χρόνον, οὐδὲ τὰ σχοινία αὐτῆς οὐ μὴ διαῤῥαγῶσιν·
\vs{21}ὅτι τὸ ὄνομα Κυρίου μέγα ὑμῖν· τόπος ὑμῖν ἔσται, ποταμοὶ καὶ διώρυχες πλατεῖς καὶ εὐρύχωροι· οὐ πορεύσῃ ταύτην τὴν ὁδὸν, οὐδὲ πορεύσεται πλοῖον ἐλαύνον.
\vs{22}Ὁ γὰρ Θεός μου μέγας ἐστίν· οὐ παρελεύσεταί με Κύριος κριτὴς ἡμῶν, Κύριος ἄρχων ἡμῶν, Κύριος βασιλεὺς ἡμῶν, Κύριος οὗτος ἡμᾶς σώσει.

\vs{23}Ἐῤῥάγησαν τὰ σχοινία σου, ὅτι οὐκ ἐνίσχυσαν· ὁ ἱστός σου ἔκλινεν, οὐ χαλάσει τὰ ἱστία, οὐκ ἀρεῖ σημεῖον, ἕως οὗ παραδοθῇ εἰς προνομήν· τοίνυν πολλοὶ χωλοὶ προνομὴν ποιήσουσι,
\vs{24}καὶ οὐ μὴ εἴπωσι, Κοπιῶ ὁ λαὸς ἐνοικῶν ἐν αὐτοῖς· ἀφεθῆ γὰρ αὐτοῖς ἡ ἁμαρτία.

\ch{34}
Προσαγάγετε ἔθνη, καὶ ἀκούσατε ἄρχοντες· ἀκουσάτω ἡ γῆ, καὶ οἱ ἐν αὐτῇ, ἡ οἰκουμένη, καὶ ὁ λαὸς ὁ ἐν αὐτῇ.
\vs{2}Διότι θυμὸς Κυρίου ἐπὶ πάντα τὰ ἔθνη καὶ ὀργὴ ἐπὶ τὸν ἀριθμὸν αὐτῶν, τοῦ ἀπολέσαι αὐτοὺς, καὶ παραδοῦναι αὐτοὺς εἰς σφαγήν.
\vs{3}Οἱ δὲ τραυματίαι αὐτῶν ῥιφήσονται, καὶ οἱ νεκροί, καὶ ἀναβήσεται αὐτῶν ἡ ὀσμή, καὶ βραχήσεται τὰ ὄρη ἀπὸ τοῦ αἵματος αὐτῶν.
\vs{4}Καὶ τακήσονται πᾶσαι αἱ δυνάμεις τῶν οὐρανῶν, καὶ ἑλιγήσεται ὁ οὐρανός ὡς βιβλίον, καὶ πάντα τὰ ἄστρα πεσεῖται ὡς φύλλα ἐξ ἀμπέλου, καὶ ὡς πίπτει φύλλα ἀπὸ συκῆς.

\vs{5}Ἐμεθύσθη ἡ μάχαιρὰ μου ἐν τῷ οὐρανῷ· ἰδοὺ ἐπὶ τὴν Ἰδουμαίαν καταβήσεται, καὶ ἐπὶ τὸν λαὸν τῆς ἀπωλείας μετὰ κρίσεως.
\vs{6}Ἡ μάχαιρα τοῦ Κυρίου ἐνεπλήσθη αἵματος, ἐπαχύνθη ἀπὸ στέατος, ἀπὸ αἵματος τράγων καὶ ἀμνῶν, καὶ ἀπὸ στέατος τράγων καὶ κριῶν· ὅτι θυσία τῷ κυρίῳ ἐν Βοσὸρ, καὶ σφαγὴ μεγάλη ἐν τῇ Ἰδουμαίᾳ.
\vs{7}Καὶ συνπεσοῦνται οἱ ἁδροὶ μετʼ αὐτῶν, καὶ οἱ κριοὶ καὶ οἱ ταῦροι, καὶ μεθυσθήσεται ἡ γῆ ἀπὸ τοῦ αἵματος, καὶ ἀπὸ τοῦ στέατος αὐτῶν ἐμπλησθήσεται.
\vs{8}Ἡμέρα γὰρ κρίσεως Κυρίου, καὶ ἐνιαυτὸς ἀνταποδόσεως κρίσεως Σιών.
\vs{9}Καὶ στραφήσονται αἱ φάραγγες αὐτῆς εἰς πίσσαν, καὶ ἡ γῆ αὐτῆς εἰς θεῖον· καὶ ἔσται ἡ γῆ αὐτῆς ὡς πίσσα καιομένη
\vs{10}νυκτὸς καὶ ἡμέρας, καὶ οὐ σβεσθήσεται εἰς τὸν αἰῶνα χρόνον, καὶ ἀναβήσεται ὁ καπνὸς αὐτῆς ἄνω, εἰς γενεὰς αὐτῆς ἐρημωθήσεται, καὶ εἰς χρόνον πολὺν
\vs{11}ὄρνεα καὶ ἐχῖνοι, καὶ ἴβεις καὶ κόρακες κατοικήσουσιν ἐν αὐτῇ· καὶ ἐπιβληθήσεται ἐπʼ αὐτὴν σπαρτίον γεωμετρίας ἐρήμου, καὶ ὀνοκένταυροι οἰκήσουσιν ἐν αὐτῇ.
\vs{12}Οἱ ἄρχοντες αὐτῆς οὐκ ἔσονται· οἱ γὰρ βασιλεῖς καὶ οἱ μεγιστᾶνες αὐτῆς ἔσονται εἰς ἀπώλειαν.
\vs{13}Καὶ ἀναφυὴσει εἰς τὰς πόλεις αὐτῶν ἀκάνθινα ξύλα, καὶ εἰς τὰ ὀχυρώματα αὐτῆς· καὶ ἔσται ἐπαύλεις σειρήνων, καὶ αὐλὴ στρουθῶν.
\vs{14}Καὶ συναντήσουσι δαιμόνια ὀνοκενταύροις, καὶ βοήσονται ἕτερος πρὸς τὸν ἕτερον, ἐκεῖ ἀναπαύσονται ὀνοκένταυροι, εὑρόντες αὑτοῖς ἀνάπαυσιν.
\vs{15}Ἐκεῖ ἐνόσσευσεν ἐχῖνος, καὶ ἔσωσεν ἡ γῆ τὰ παιδία αὐτῆς μετὰ ἀσφαλείας· ἐκεῖ συνήντησαν ἔλαφοι καὶ εἶδον τὰ πρόσωπα ἀλλήλων.
\vs{16}Ἀριθμῷ παρῆλθον, καὶ μία αὐτῶν οὐκ ἀπώλετο· ἑτέρα τὴν ἑτέραν οὐκ ἐζήτησαν, ὅτι ὁ Κύριος αὐτοῖς ἐνετείλατο, καὶ τὸ πνεῦμα αὐτοῦ συνήγαγεν αὐτά.
\vs{17}Καὶ αὐτὸς ἐπιβαλεῖ αὐτοῖς κλήρους, καὶ ἡ χεὶρ αὐτοῦ διεμέρισε βόσκεσθαι· εἰς τὸν αἰῶνα χρόνον κληρονομήσετε, γενεὰς γενεῶν ἀναπαύσονται ἐπʼ αὐτῆς.

\ch{35}
Εὐφράνθητι ἔρημος διψῶσα, ἀγαλλιάσθω ἔρημος, καὶ ἀνθείτω ὡς κρίνον.
\vs{2}Καὶ ἐξανθήσει καὶ ἀγαλλιάσεται τὰ ἔρημα τοῦ Ἰορδανου, ἡ δόξα τοῦ Λιβάνου ἐδόθη αὐτῇ, καὶ ἡ τιμὴ τοῦ Καρμήλου, καὶ ὁ λαός μου ὄψεται τὴν δόξαν Κυρίου, καὶ τὸ ὕψος τοῦ Θεοῦ.

\vs{3}Ἰσχύσατε χεῖρες ἀνειμέναι, καὶ γόνατα παραλελυμένα.
\vs{4}Παρακαλέσατε οἱ ὀλιγόψυχοι τῇ διανοίᾳ· ἰσχύσατε, μὴ φοβεῖσθε· ἰδοὺ ὁ Θεὸς ἡμῶν κρίσιν ἀνταποδίδωσι, καὶ ἀνταποδώσει, αὐτὸς ἥξει καὶ σώσει ἡμᾶς.
\vs{5}Τότε ἀνοιχθήσονται ὀφθαλμοὶ τυφλῶν, καὶ ὦτα κωφῶν ἀκούσονται.
\vs{6}Τότε ἁλεῖται ὡς ἔλαφος ὁ χωλὸς, τρανὴ δὲ ἔσται γλῶσσα μογιλάλων, ὅτι ἐῤῥάγη ἐν τῇ ἐρήμῳ ὕδωρ, καὶ φάραγξ ἐν γῇ διψώσῃ.
\vs{7}Καὶ ἔσται ἡ ἄνυδρος εἰς ἕλη, καὶ εἰς τὴν διψῶσαν γῆν πηγὴ ὕδατος ἔσται· ἐκεῖ εὐφροσύνη ὀρνέων, ἐπαύλεις καλάμου καὶ ἕλη.
\vs{8}Ἔσται ἐκεῖ ὁδὸς καθαρά, καὶ ὁδὸς ἁγία κληθήσεται, καὶ οὐ μὴ παρέλθῃ ἐκεῖ ἀκάθαρτος, οὐδὲ ἔσται ἐκεῖ ὁδὸς ἀκάθαρτος· οἱ δὲ διεσπαρμένοι πορεύσονται ἐπʼ αὐτῆς, καὶ οὐ μὴ πλανηθῶσι.
\vs{9}Καὶ οὐκ ἔσται ἐκεῖ λέων, οὐδὲ τῶν πονηρῶν θηρίων οὐ μὴ ἀναβῇ εἰς αὐτὴν, οὐδὲ μὴ εὑρεθῇ ἐκεῖ, ἀλλὰ πορεύσονται ἐν αὐτῇ λελυτρωμένοι,
\vs{10}Καὶ συνηγμένοι διὰ Κύριον, καὶ ἀποστραφήσονται, καὶ ἥξουσιν εἰς Σιὼν μετʼ εὐφροσύνης, καὶ εὐφροσύνη αἰώνιος ὑπὲρ κεφαλῆς αὐτῶν· ἐπὶ γὰρ τῆς κεφαλῆς αὐτῶν αἴνεσις καὶ ἀγαλλίαμα, καὶ εὐφροσύνη καταλήμψεται αὐτοὺς, ἀπέδρα ὀδύμη καὶ λύπη καὶ στεναγμός.

\ch{36}
Καὶ ἐγένετο τοῦ τεσσαρεσκαιδεκάτου ἔτους βασιλεύοντος Ἑζεκίου, ἀνέβη Σενναχηρεὶμ βασιλεὺς Ἀσσυρίων ἐπὶ τὰς πόλεις τῆς Ἰουδαίας τὰς ὀχυρὰς, καὶ ἔλαβεν αὐτάς.
\vs{2}Καὶ ἀπέστειλε βασιλεὺς Ἀσσυρίων τὸν Ῥαβσάκην ἐκ Λάχης εἰς Ἱερουσαλὴμ πρὸς τὸν βασιλέα Ἐζεκίαν μετὰ δυνάμεως πολλῆς· καὶ ἔστη ἐν τῷ ὑδραγωγῷ τῆς κολυμβήθρας τῆς ἄνω ἐν τῇ ὁδῷ τοῦ ἀγροῦ τοῦ κναφέως.
\vs{3}Καὶ ἐξῆλθε πρὸς αὐτὸν Ἑλιακεὶμ ὁ τοῦ Χελκίου ὁ οἰκονόμος, καὶ Σόμνᾶς ὁ γραμματεὺς, καὶ Ἰωὰχ ὁ τοῦ Ἀσὰφ ὁ ὑπομνηματογράφος.

\vs{4}Καὶ εἶπεν αὐτοῖς Ῥαβσάκης, εἴπατε Ἐζεκίᾳ, Τάδε λέγει ὁ βασιλεὺς ὁ μέγας, βασιλεὺς Ἀσσυρίων, τί πεποιθὼς εἶ;
\vs{5}Μὴ ἐν βουλῇ καὶ λόγοις χειλέων παράταξις γίνεται; καὶ νῦν ἐπὶ τίνα πέποιθας, ὅτι ἀπειθεῖς μοι;
\vs{6}Ἰδοὺ πεποιθὼς εἶ ἐπὶ τὴν ῥάβδον τὴν καλαμίνην τὴν τεθλασμένην ταύτην, ἐπʼ Αἴγυπτον· ὡς ἂν ἐπιστηρισθῇ ἀνὴρ ἐπʼ αὐτὴν, εἰσελεύσεται εἰς τὴν χεῖρα αὐτοῦ, καὶ τρήσει αὐτὴν. οὕτως ἐστὶ Φαραὼ βασιλεὺς Αἰγύπτου, καὶ πάντες οἱ πεποιθότες ἐπʼ αὐτῷ.
\vs{7}Εἰ δὲ λέγετε, ἐπὶ Κύριον τὸν Θεὸν ἡμῶν πεποίθαμεν,
\vs{8}νῦν μίχθητε τῷ κυρίῳ μου τῷ βασιλεῖ Ἀσσυρίων, καὶ δώσω ὑμῖν δισχιλίαν ἵππον, εἰ δυνήσεσθε δοῦναι ἀναβάτας ἐπʼ αὐτούς.
\vs{9}Καὶ πῶς δύνασθε ἀποστρέψαι εἰς πρόσωπον τῶν τοπαρχῶν· οἰκέται εἰσὶν, οἱ πεποιθότες ἐπʼ Αἴγυπτίοις, εἰς ἵππον καὶ ἀναβάτην.
\vs{10}Καὶ νῦν μὴ ἄνευ Κυρίου ἀνέβημεν ἐπὶ τὴν χώραν ταύτην πολεμῆσαι αὐτήν; Κύριος εἶπε πρὸς μὲ, ἀνάβηθι ἐπὶ τὴν γῆν ταύτην, καὶ διάφθειρον αὐτήν.

\vs{11}Καὶ εἶπε πρὸς αὐτὸν Ἑλιακεὶμ, καὶ Σομνᾶς, καὶ Ἰωὰχ, λάλησον πρὸς τοὺς παῖδάς σου Συριστί· ἀκούομεν γὰρ ἡμεῖς· καὶ μὴ λάλει πρὸς ἡμᾶς Ἰουδαϊστί· καὶ ἱνατί λαλεῖς εἰς τὰ ὦτα τῶν ἀνθρώπων ἐπὶ τῷ τείχει;
\vs{12}Καὶ εἶπε πρὸς αὐτοὺς Ῥαβσάκης, μὴ πρὸς τὸν κύριον ὑμῶν ἢ πρὸς ὑμᾶς ἀπέσταλκέ με ὁ κύριός μου, λαλῆσαι τοὺς λόγους τούτους; οὐχὶ πρὸς τοὺς ἀνθρώπους τοὺς καθημένους ἐπὶ τῷ τείχει, ἵνα φάγωσι κόπρον, καὶ πίωσιν οὖρον μεθʼ ὑμῶν ἅμα;

\vs{13}Καὶ ἔστη Ῥαβσάκης, καὶ ἀνεβόησε φωνῇ μεγάλῃ Ἰουδαϊστὶ, καὶ εἶπεν, ἀκούσατε τοὺς λόγους τοῦ βασιλέως τοῦ μεγάλου, βασιλέως Ἀσσυρίων.
\vs{14}Τάδε λέγει ὁ βασιλεύς, μὴ ἀπατάτω ὑμᾶς Ἐζεκίας λόγοις, οὐ δύνηται ῥύσασθαι ὑμᾶς.
\vs{15}Καὶ μὴ λεγέτω ὑμῖν Ἐζεκίας, ὅτι ῥύσεται ὑμᾶς ὁ Θεὸς, καὶ οὐ μὴ παραδοθῇ ἡ πόλις αὕτη ἐν χειρὶ βασιλέως Ἀσσυρίων.
\vs{16}Μὴ ἀκούετε Ἐζεκίου· τάδε λέγει ὁ βασιλεὺς Ἀσσυρίων, εἰ βούλεσθε εὐλογηθῆναι, ἐκπορεύεσθε πρὸς μέ, καὶ φάγεσθε ἕκαστος τὴν ἄμπελον αὐτοῦ καὶ τὰς συκὰς, καὶ πίεσθε ὕδωρ ἐκ τοῦ λάκκου ὑμῶν,
\vs{17}ἕως ἂν ἔλθω, καὶ λάβω ὑμᾶς εἰς γῆν, ὡς ἡ γῆ ὑμῶν, γῆ σίτου καὶ οἴνου καὶ ἄρτων καὶ ἀμπελώνων.
\vs{18}Μὴ ἀπατάτω ὑμᾶς Ἐζεκίας, λέγων, ὁ Θεὸς ῥύσεται ὑμᾶς· μὴ ἐῤῥύσαντο οἱ θεοὶ τῶν ἐθνῶν, ἕκαστος τὴν ἑαῦτοῦ χώραν ἐκ χειρὸς βασιλέως Ἀσσυρίων;
\vs{19}Ποῦ ἐστιν ὁ θεὸς Ἐμὰθ καὶ Ἀρφάθ; καὶ ποῦ ὁ θεὸς τῆς πόλεως Ἐπφαρουαίμ; μὴ ἐδύναντο ῥύσασθαι Σαμάρειαν ἐκ χειρός μου;
\vs{20}Τίς τῶν θεῶν πάντων τῶν ἐθνῶν τούτων, ὅστις ἐῤῥύσατο τὴν γῆν αὐτοῦ ἐκ χειρός μου, ὅτι ῥύσεται ὁ Θεὸς τὴν Ἱερουσαλὴμ ἐκ χειρός μου;
\vs{21}Καὶ ἐσιώπησαν, καὶ οὐδεὶς ἀπεκρίθη αὐτῷ λόγον, διὰ τὸ προστάξαι τὸν βασιλέα μηδένα ἀποκριθῆναι.

\vs{22}Καὶ εἰσῆλθεν Ἑλιακεὶμ ὁ τοῦ Χελκίου, οἰκονόμος, καὶ Σομνᾶς ὁ γραμματεὺς τῆς δυνάμεως, καὶ Ἰωὰχ ὁ τοῦ Ἀσὰφ ὁ ὑπομνηματογράφος, πρὸς Ἐζεκίαν, ἐσχισμένοι τοὺς χιτῶνας, καὶ ἀνήγγειλαν αὐτῷ τοὺς λόγους Ῥαβσάκου.

\ch{37}
Καὶ ἐγένετο ἐν τῷ ἀκοῦσαι τὸν βασιλέα Ἐζεκίαν, ἔσχισε τὰ ἱμάτια, καὶ περιεβάλετο σάκκον, καὶ ἀνέβη εἰς τὸν οἶκον Κυρίου.

\vs{2}Καὶ ἀπέστειλεν Ἑλιακεὶμ τὸν οἰκονόμον, καὶ Σόμνᾶν τὸν γραμματέα, καὶ τοὺς πρεσβυτέρους τῶν ἱερέων περιβεβλημένους σάκκους, πρὸς Ἠσαΐαν υἱὸν Ἀμὼς τὸν προφήτην.
\vs{3}Καὶ εἶπαν αὐτῷ, τάδε λέγει Ἐζεκίας, ἡμέρα, θλίψεως καὶ ὀνειδισμοῦ καὶ ἐλεγμοῦ καὶ ὀργῆς ἡ σήμερον ἡμέρα, ὅτι ἥκει ἡ ὠδὶν τῇ τικτούσῃ, ἰσχὺν δὲ οὐκ ἔχει τοῦ τεκεῖν.
\vs{4}Εἰσακούσαι Κύριος ὁ Θεός σου τοὺς λόγους Ῥαβσάκου, οὓς ἀπέστειλε βασιλεὺς Ἀσσυρίων, ὀνειδίζειν Θεὸν ζῶντα, καὶ ὀνειδίζειν λόγους οὓς ἤκουσε Κύριος ὁ Θεός σου, καὶ δεηθήσῃ πρὸς Κύριόν σου περὶ τῶν καταλελειμμένων τούτων.

\vs{5}Καὶ ἦλθον οἱ παῖδες τοῦ βασιλέως Ἐζεκίου πρὸς Ἡσαΐαν.
\vs{6}Καὶ εἶπεν αὐτοῖς Ἡσαΐας, οὕτως ἐρεῖτε πρὸς τὸν κύριον ὑμῶν, τάδε λέγεἱ Κύριος, μὴ φοβηθῇς ἀπὸ τῶν λόγων ὧν ἤκουσας, οὓς ὠνείδισάν με οἱ πρέσβεις βασιλέως Ἀσσυρίων.
\vs{7}Ἰδοὺ ἐγὼ ἐμβάλλω εἰς αὐτὸν πνεῦμα, καὶ ἀκούσας ἀγγελίαν, ἀποστραφήσεται εἰς τὴν χώραν αὐτοῦ, καὶ πεσεῖται μαχαίρᾳ ἐν τῇ γῇ αὐτοῦ.

\vs{8}Καὶ ἀπέστρεψε Ῥαβσάκης, καὶ κατέλαβε τὸν βασιλέα Ἀσσυρίων πολιορκοῦντα Λοβνάν· καὶ ἤκουσεν ὅτι ἀπῇρεν ἀπὸ Λαχίς.
\vs{9}Καὶ ἐξῆλθε Θαρακὰ βασιλεὺς Αἰθιόπων πολιορκῆσαι αὐτόν· καὶ ἀκούσας ἀπέστρεψε, καὶ ἀπέστειλεν ἀγγέλους πρὸς Ἐζεκίαν, λέγων,
\vs{10}οὕτως ἐρεῖτε Ἐζεκίᾳ βασιλεῖ τῆς Ἰουδαίας, μή σε ἀπατάτω ὁ Θεός σου, ἐφʼ ᾧ πέποιθας ἐπʼ αὐτῷ, λέγων, οὐ μὴ παραδοθῇ Ἱερουσαλὴμ ἐν χειρὶ βασιλέως Ἀσσυρίων.

\vs{11}Σὺ οὐκ ἤκουσας ἃ ἐποίησαν βασιλεῖς Ἀσσυρίων, πᾶσαν τὴν γῆν ὡς ἀπώλεσαν; καὶ σὺ ῥυσθήσῃ;
\vs{12}Μὴ ἐῤῥύσαντο αὐτοὺς οἱ θεοὶ τῶν ἐθνῶν, οὓς ἀπώλεσαν οἱ πατέρες μου, τήν τε Γωζᾶν, καὶ Χαῤῥὰν, καὶ Ῥαφὲθ, αἵ εἰσιν ἐν χώρᾳ Θεεμάθ;
\vs{13}Ποῦ εἰσι βασιλεῖς Ἐμάθ; καὶ ποῦ Ἀρφάθ; καὶ ποῦ πόλεως Ἐπφαρουαὶμ, Ἀναγονγάυα;

\vs{14}Καὶ ἔλαβεν Ἐζεκίας τὸ βιβλίον παρὰ τῶν ἀγγέλων, καὶ ἀνέγνω αὐτὸ, καὶ ἀνέβη εἰς οἶκον Κυρίου, καὶ ἤνοιξεν αὐτὸ ἐναντίον Κυρίου.
\vs{15}Καὶ προσηύξατο Ἐζεκίας πρὸς Κύριον, λέγων,

\vs{16}Κύριος σαβαὼθ ὁ Θεὸς Ἰσραὴλ, ὁ καθήμενος ἐπὶ τῶν χερουβὶμ, σὺ εἶ ὁ Θεὸς μόνος πάσης βασιλείας τῆς οἰκουμένης, σὺ ἐποίησας τὸν οὐρανὸν καὶ τὴν γῆν·
\vs{17}Κλῖνον Κύριε τὸ οὖς σου, εἰσάκουσον Κύριε, ἄνοιξον Κύριε τοὺς ὀφθαλμούς σου, εἴσβλεψον Κύριε, καὶ ἴδε τοὺς λόγους Σενναχηρεὶμ, οὓς ἀπέστειλεν ὀνειδίζειν Θεὸν ζῶντα.
\vs{18}Ἐπʼ ἀληθείας γὰρ Κύριε ἠρήμωσαν βασιλεῖς Ἀσσυρίων τὴν οἰκουμένην ὅλην, καὶ τὴν χώραν αὐτῶν,
\vs{19}καὶ ἀνέβαλον τὰ εἴδωλα αὐτῶν εἰς τὸ πῦρ· οὐ γὰρ θεοὶ ἦσαν, ἀλλὰ ἔργα χειρῶν ἀνθρώπων, ξύλα καὶ λίθοι· καὶ ἀπώσαντο αὐτούς.
\vs{20}Νῦν δὲ Κύριε ὁ Θεὸς ἡμῶν σῶσον ἡμᾶς ἐκ χειρὸς αὐτοῦ, ἵνα γνῷ πᾶσα βασιλεία τῆς γῆς, ὅτι σῦ εἶ ὁ Θεὸς μόνος.

\vs{21}Καὶ ἀπεστάλη Ἡσαΐας υἱὸς Ἀμὼς πρὸς Ἐζεκίαν, καὶ εἶπεν αὐτῷ, τάδε λέγει Κύριος ὁ Θεὸς Ἰσραὴλ, ἤκουσα ἃ προσηύξω πρὸς μὲ περὶ Σενναχηρεὶμ βασιλέως Ἀσσυρίων.
\vs{22}Οὗτος ὁ λόγος ὃν ἐλάλησε περὶ αὐτοῦ ὁ Θεὸς, ἐφαύλισέ σε, καὶ ἐμυκτήρισέ σε παρθένος θυγάτηρ Σιὼν, ἐπὶ σοὶ κεφαλὴν ἐκίνησε θυγάτηρ Ἱερουσαλήμ.
\vs{23}Τίνα ὠνείδισας καὶ παρώξυνας; ἢ πρὸς τίνα ὕψωσας τὴν φωνήν σου; καὶ οὐκ ᾖρας εἰς ὕψος τοὺς ὀφθαλμούς σου πρὸς τὸν ἅγιον τοῦ Ἰσραήλ;
\vs{24}Ὅτι διʼ ἀγγέλων ὠνείδισας Κύριον· σὺ γὰρ εἶπας, τῷ πλήθει τῶν ἁρμάτων ἐγὼ ἀνέβην εἰς ὕψος ὀρέων, καὶ εἰς τὰ ἔσχατα τοῦ Λιβάνου, καὶ ἔκοψα τὸ ὕψος τῆς κέδρου αὐτοῦ, καὶ τὸ κάλλος τῆς κυπαρίσσου, καὶ εἰσῆλθον εἰς ὕψος μέρους τοῦ δρυμοῦ,
\vs{25}καὶ ἔθηκα γέφυραν, καὶ ἠρήμωσα ὕδατα καὶ πᾶσαν συναγωγὴν ὕδατος.

\vs{26}Οὐ ταῦτα ἤκουσας πάλαὶ ἃ ἐγὼ ἐποίησα; ἐξ ἡμερῶν ἀρχαίων συνέταξα, νῦν δὲ ἐπέδειξα ἐξερημῶσαι ἔθνη ἐν ὀχυροῖς, καὶ οἰκοῦντας ἐν πόλεσιν ὀχυραῖς.
\vs{27}Ἀνῆκα τὰς χεῖρας, καὶ ἐξηράνθησαν, καὶ ἐγένοντο ὡς χόρτος ξηρὸς ἐπὶ δωμάτων, καὶ ὡς ἄγρωστις.
\vs{28}Νῦν δὲ τὴν ἀνάπαυσίν σου, καὶ τὴν ἔξοδόν σου, καὶ τὴν εἴσοδόν σου ἐγὼ ἐπίσταμαι.
\vs{29}Ὁ δὲ θυμός σου ὃν ἐθυμώθης, καὶ ἡ πικρία σου ἀνέβη πρὸς μὲ, καὶ ἐμβαλῶ φιμὸν εἰς τὴν ῥῖνά σου, καὶ χαλινὸν εἰς τὰ χείλη σου, καὶ ἀποστρέψω σε τῇ ὁδῷ ᾗ ἦλθες ἐν αὐτῇ.

\vs{30}Τοῦτο δέ σοι τὸ σημεῖον· φάγε τοῦτον τὸν ἐνιαυτὸν ἃ ἔσπαρκας, τῷ δὲ ἐνιαὐτῷ τῷ δευτέρῳ τὸ κατάλειμμα, τῷ δὲ τρίτῳ σπείραντες ἀμήσατε, καὶ φυτεύσατε ἀμπελῶνας, καὶ φάγεσθε τὸν καρπὸν αὐτῶν.
\vs{31}Καὶ ἔσονται οἱ καταλελιμμένοι ἐν τῇ Ἰουδαίᾳ, φυήσουσι ῥίζαν κάτω, καὶ ποιήσουσι σπέρμα ἄνω·
\vs{32}Ὅτι ἐξ Ἱερουσαλὴμ ἔσονται οἱ καταλελειμμένοι, καὶ οἱ σωζόμενοι ἐξ ὄρους Σιών· ὁ ζῆλος Κυρίου σαβαὼθ ποιήσει ταῦτα.
\vs{33}Διατοῦτο οὕτως λέγει Κύριος ἐπὶ βασιλέα Ἀσσυρίων, οὐ μὴ εἰσέλθῃ εἰς τὴν πόλιν ταύτην, οὐδὲ μὴ βάλῃ ἐπʼ αὐτὴν βέλος, οὐδὲ μὴ ἐπιβάλῃ ἐπʼ αὐτὴν θυρεὸν, οὐδὲ μὴ κυκλώσῃ ἐπʼ αὐτὴν χάρακα·
\vs{34}Ἀλλὰ τῇ ὁδῷ ᾗ ἦλθεν, ἐν αὐτῇ ἀποστραφήσεται, καὶ εἰς τὴν πόλιν ταύτην οὐ μὴ εἰσέλθῃ· τάδε λέγει Κύριος.
\vs{35}Ὑπερασπιῶ ὑπὲρ τῆς πόλεως ταύτης τοῦ σῶσαι αὐτὴν διʼ ἐμὲ, καὶ διὰ Δαυὶδ τὸν παῖδά μου.

\vs{36}Καὶ ἐξῆλθεν ἄγγελος Κυρίου, καὶ ἀνεῖλεν ἐκ τῆς παρεμβολῆς τῶν Ἀσσυρίων ἑκατὸν ὀγδοηκονταπέντε χιλιάδας· καὶ ἀναστάντες τοπρωῒ, εὗρον πάντα τὰ σώματα νεκρά.
\vs{37}Καὶ ἀπῆλθεν ἀποστραφεὶς Σενναχηρεὶμ βασιλεὺς Ἀσσυρίων, καὶ ᾤκησεν ἐν Νινευῇ.
\vs{38}Καὶ ἐν τῷ αὐτὸν προσκυνεῖν ἐν τῷ οἴκῳ Νασαρὰχ τὸν πάτραρχον αὐτοῦ, Ἀδραμέλεχ καὶ Σαρασὰρ οἱ υἱοὶ αὐτοῦ ἐπάταξαν αὐτὸν μαχαίραις, αὐτοὶ δὲ διεσώθησαν εἰς Ἀρμενίαν, καὶ ἐβασίλευσεν Ἀσορδὰν ὁ υἱὸς αὐτοῦ ἀντʼ αὐτοῦ.

\ch{38}
Ἐγένετο δὲ ἐν τῷ καιρῷ ἐκείνῳ, ἐμαλακίσθη Ἐζεκίας ἕως θανάτου· καὶ ἦλθε πρὸς αὐτὸν Ἡσαΐας υἱὸς Ἀμὼς ὁ προφήτης, καὶ εἶπε πρὸς αὐτὸν, τάδε λέγει Κύριος, τάξαι περὶ τοῦ οἴκου σου, ἀποθνήσκεις γὰρ σὺ, καὶ οὐ ζήσῃ.
\vs{2}Καὶ ἀπέστρεψεν Ἐζεκίας τὸ πρόσωπον αὐτοῦ πρὸς τὸν τεῖχον, καὶ προσηύξατο πρὸς Κύριον,
\vs{3}λέγων, μνήσθητι Κύριε, ὡς ἐπορεύθην ἐνώπιόν σου μετὰ ἀληθείας, ἐν καρδίᾳ ἀληθινῇ, καὶ τὰ ἀρεστὰ ἐνώπιόν σου ἐποίησα· καὶ ἔκλαυσεν Ἐζεκίας κλαυθμῷ μεγάλῳ.
\vs{4}Καὶ ἐγένετο λόγος Κυρίου πρὸς Ἡσαΐαν, λέγων,
\vs{5}πορεύθητι, καὶ εἰπὸν Ἐζεκίᾳ, τάδε λέγει Κύριος ὁ Θεὸς Δαυὶδ τοῦ πατρός σου, ἤκουσα τῆς προσευχῆς σου, καὶ εἶδον τὰ δάκρυά σου· ἰδοὺ προστίθημι πρὸς τὸν χρόνον σου δεκαπέντε ἔτη,
\vs{6}καὶ ἐκ χειρὸς βασιλέως Ἀσσυρίων ῥύσομαί σε καὶ τὴν πόλιν ταύτην, καὶ ὑπερασπιῶ ὑπὲρ τῆς πόλεως ταύτης.
\vs{7}Τοῦτο δέ σοι τὸ σημεῖον παρὰ Κυρίου, ὅτι ποιήσει ὁ Θεὸς τὸ ῥῆμα τοῦτο·
\vs{8}Ἰδοὺ ἐγὼ στρέφω τὴν σκιὰν τῶν ἀναβαθμῶν οὓς κατέβη τοὺς δέκα ἀναβαθμοὺς τοῦ οἴκου τοῦ πατρός σου ὁ ἥλιος, ἀποστρέψω τὸν ἥλιον τοὺς δέκα ἀναβαθμούς· καὶ ἀνέβη ὁ ἥλιος τοὺς δέκα ἀναβαθμοὺς, οὓς κατέβη ἡ σκιά.

\vs{9}ΠΡΟΣΕΥΧΗ ἘΖΕΚΙΟΥ ΒΑΣΙΛΕΩΣ ΤΗΣ ἸΟΥΔΑΙΑΣ, ἩΝΙΚΑ ἘΜΑΛΑΚΙΣΘΗ, ΚΑΙ ἈΝΕΣΤΗ ἘΚ ΤΗΣ ΜΑΛΑΚΙΑΣ ΑΥΤΟΥ.

\vs{10}Ἐγὼ εἶπα ἐν τῷ ὕψει τῶν ἡμερῶν μου, πορεύσομαι ἐν πύλαις ἅδου, καταλείψω τὰ ἔτη τὰ ἐπίλοιπα.
\vs{11}Εἶπα, οὐκέτι οὐ μὴ ἴδω τὸ σωτήριον τοῦ Θεοῦ ἐπὶ γῆς ζώντων, οὐκέτι μὴ ἴδω τὸ σωτήριον τοῦ Ἰσραὴλ ἐπὶ γῆς, οὐκέτι μὴ ἴδω ἄνθρωπον.
\vs{12}Ἐξέλιπεν ἐκ τῆς συγγενείας μου, κατέλιπον τὸ ἐπίλοιπον τῆς ζωῆς μου, ἐξῆλθε καὶ ἀπῆλθεν ἀπʼ ἐμοῦ ὥσπερ ὁ σκηνὴν καταλύων πήξας· ὡς ἱστὸς τὸ πνεῦμά μου παρʼ ἐμοὶ ἐγένετο, ἐρίθου ἐγγιζούσης ἐκτεμεῖν.
\vs{13}Ἐν τῇ ἡμέρᾳ ἐκείνῃ παρεδόθην ἕως πρωῒ ὡς λέοντι, οὕτως συνέτριψε πάντα τὰ ὀστᾶ μου· ἀπὸ γὰρ τῆς ἡμέρας ἕως νυκτὸς παρεδόθην.
\vs{14}Ὡς χελιδὼν, οὕτω φωνήσω, καὶ ὡς περιστερὰ, οὕτω μελετῶ· ἐξέλιπον γάρ μου οἱ ὀφθαλμοὶ τοῦ βλέπειν εἰς τὸ ὕψος τοῦ οὐρανοῦ πρὸς τὸν Κύριον, ὃς ἐξείλατό με,
\vs{15}καὶ ἀφείλατό μου τὴν ὀδύνην τῆς ψυχῆς.
\vs{16}Κύριε, περὶ αὐτῆς γὰρ ἀνηγγέλη σοι, καὶ ἐξήγειράς μου τὴν πνοὴν, καὶ παρακληθεὶς ἔζησα.
\vs{17}Εἵλου γάρ μου τὴν ψυχὴν, ἴνα μὴ ἀπόληται, καὶ ἀπέῤῥιψας ὀπίσω μου πάσας τὰς ἁμαρτίας.
\vs{18}Οὐ γὰρ οἱ ἐν ᾅδου αἰνέσουσί σε, οὐδὲ οἱ ἀποθανόντες εὐλογήσουσί σε, οὐδὲ ἐλπιοῦσιν οἱ ἐν ᾄδου τὴν ἐλεημοσύνην σου.
\vs{19}Οἱ ζῶντες εὐλογήσουσί σε ὃν τρόπον κᾀγώ· ἀπὸ γὰρ τῆς σήμερον παιδία ποιήσω, ἅ ἀναγγελοῦσι τὴν δικαιοσύνην σου
\vs{20}Θεὲ τῆς σωτηρίας μου, καὶ οὐ παύσομαι εὐλογῶν σε μετὰ ψαλτηρίου πάσας τὰς ἡμέρας τῆς ζωῆς μου, κατέναντι τοῦ οἴκου τοῦ Θεοῦ.

\vs{21}Καὶ εἶπεν Ἡσαΐας πρὸς Ἐζεκίαν, λάβε παλάθην ἐκ σύκων, καὶ τρίψον, καὶ κατάπλασαι, καὶ ὑγιὴς ἔσῃ.
\vs{22}Καὶ εἶπεν Ἐζεκίας, τοῦτο σημεῖον πρὸς Ἐζεκίαν, ὅτι ἀναβήσομαι εἰς τὸν οἶκον τοῦ Θεοῦ.

\ch{39}
Ἐν τῷ καιρῷ ἐκείνῳ ἀπέστειλε Μαρωδὰχ Βαλαδὰν ὁ υἱὸς τοῦ Βαλαδὰν, ὁ βασιλεὺς τῆς Βαβυλωνίας, ἐπιστολὰς καὶ πρέσβεις καὶ δῶρα Ἐζεκίᾳ· ἤκουσε γὰρ, ὅτι ἐμαλακίσθῃ ἕως θανάτου, καὶ ἀνέστη.
\vs{2}Καὶ ἐχάρη ἐπʼ αὐτοῖς Ἐξεκίας, καὶ ἔδειξεν αὐτοῖς τὸν οἶκον τοῦ νεχωθᾶ, καὶ τοῦ ἀργυρίου, καὶ τοῦ χρυσίου, καὶ τῆς στακτῆς, καὶ τῶν θυμιαμάτων, καὶ τοῦ μύρου, καὶ πάντας τοὺς οἴκους τῶν σκευῶν τῆς γάζης, καὶ πάντα ὅσα ἦν ἐν τοῖς θησαυροῖς αὐτοῦ· καὶ οὐκ ἦν οὐθὲν ὃ οὐκ ἔδειξεν Ἐζεκίας ἐν τῷ οἴκῳ αὐτοῦ, καὶ ἐν πάσῃ τῇ ἐξουσίᾳ αὐτοῦ.

\vs{3}Καὶ ἦλθεν Ἡσαΐας ὁ προφήτης πρὸς τὸν βασιλέα Ἐζεκίαν, καὶ εἶπε πρὸς αὐτὸν, τί λέγουσιν οἱ ἄνθρωποι οὗτοι; καὶ πόθεν ἥκασι πρὸς σέ; καὶ εἶπεν Ἐζεκίας, ἐκ γῆς πόῤῥωθεν ἥκασι πρὸς μέ, ἐκ Βαβυλῶνος.
\vs{4}Καὶ εἶπεν Ἠσαΐας, τί εἴδοσαν ἐν τῷ οἴκῳ σου; καὶ εἶπεν Ἐζεκίας, πάντα τὰ ἐν τῷ οἴκῳ μου εἴδοσαν, καὶ οὐκ ἔστιν ἐν τῷ οἴκῳ μου ὃ οὐκ εἴδοσαν, ἀλλὰ καὶ τὰ ἐν τοῖς θησαυροῖς μου.
\vs{5}Καὶ εἶπεν Ἡσαΐας αὐτῷ, ἄκουσον τὸν λόγον Κυρίου σαβαώθ.
\vs{6}Ἰδοὺ ἡμέραι ἔρχονται, καὶ λήψονται πάντα τὰ ἐν τῷ οἴκῳ σου, καὶ ὅσα συνήγαγον οἱ πατέρες σου ἕως τῆς ἡμέρας ταύτης, εἰς Βαβυλῶνα ἥξει, καὶ οὐδὲν οὐ μὴ καταλείπωσιν· εἶπε δὲ ὁ Θεὸς,
\vs{7}ὅτι καὶ ἀπὸ τῶν τέκνων σου ὧν γεννήσεις, λήψονται, καὶ ποιήσουσι σπάδοντας ἐν τῷ οἴκῳ τοῦ βασιλέως τῶν Βαβυλωνίων.
\vs{8}Καὶ εἶπεν Ἐζεκίας Ἡσαΐᾳ, ἀγαθὸς ὁ λόγος Κυρίου, ὃν ἐλάλησε· γενέσθω δὴ εἰρήνη καὶ δικαιοσύνη ἐν ταῖς ἡμέραις μου.

\ch{40}
Παρακαλεῖτε παρακαλεῖτε τὸν λαόν μου, λέγει ὁ Θεός.
\vs{2}Ἱερεῖς λαλήσατε εἰς τὴν καρδίαν Ἱερουσαλὴμ, παρακαλέσατε αὐτὴν, ὅτι ἐπλήσθη ἡ ταπείνωσις αὐτῆς, λέλυται αὐτῆς ἡ ἁμαρτία, ὅτι ἐδέξατο ἐκ χειρὸς Κυρίου διπλᾶ τὰ ἁμαρτήματα αὐτῆς.

\vs{3}Φωνὴ βοῶντος ἐν τῇ ἐρήμῳ, ἑτοιμάσατε τὴν ὁδὸν Κυρίου, εὐθείας ποιεῖτε τὰς τρίβους τοῦ Θεοῦ ἡμῶν.
\vs{4}Πᾶσα φάραγξ πληρωθήσεται, καὶ πᾶν ὄρος καὶ βουνὸς ταπεινωθήσεται· καὶ ἔσται πάντα τὰ σκολιὰ εἰς εὐθεῖαν, καὶ ἡ τραχεῖα εἰς πεδία.
\vs{5}Καὶ ὀφθήσεται ἡ δόξα Κυρίου, καὶ ὄψεται πᾶσα σὰρξ τὸ σωτήριον τοῦ Θεοῦ, ὅτι Κύριος ἐλάλησε.

\vs{6}Φωνὴ λέγοντος, βόησον· καὶ εἶπα, τί βοήσω; πᾶσα σὰρξ χορτὸς, καὶ πᾶσα δόξα ἀνθρώπου ὡς ἄνθος χόρτου·
\vs{7}Ἐξηράνθη ὁ χόρτος, καὶ τὸ ἄνθος ἐξέπεσε·
\vs{8}τὸ δὲ ῥῆμα τοῦ Θεοῦ ἡμῶν μένει εἰς τὸν αἰῶνα.

\vs{9}Ἐπʼ ὄρος ὑψηλὸν ἀνάβηθι ὁ εὐαγγελιζόμενος Σιὼν, ὕψωσον τῇ ἰσχύϊ τὴν φωνήν σου ὁ εὐαγγελιζόμενος Ἱερουσαλήμ· ὑψώσατε, μὴ φοβεῖσθε· εἰπὸν ταῖς πόλεσιν Ἰούδα, ἰδοὺ ὁ Θεὸς ὑμῶν, ἰδοὺ Κύριος·
\vs{10}Κύριος μετὰ ἰσχύος ἔρχεται, καὶ ὁ βραχίων μετὰ κυρίας· ἰδοὺ ὁ μισθὸς αὐτοῦ μετʼ αὐτοῦ, καὶ τὸ ἔργον ἐναντίον αὐτοῦ.
\vs{11}Ὡς ποιμὴν ποιμανεῖ τὸ ποίμνιον αὐτοῦ, καὶ τῷ βραχίονι αὐτοῦ συνάξει ἄρνας, καὶ ἐν γαστρὶ ἐχούσας παρακαλέσει.
\vs{12}Τίς ἐμέτρησε τῇ χειρὶ τὸ ὕδωρ, καὶ τὸν οὐρανὸν σπιθαμῇ, καὶ πᾶσαν τὴν γῆν δρακί; τίς ἔστησε τὰ ὄρη σταθμῷ, καὶ τὰς νάπας ζυγῷ;
\vs{13}Τίς ἔγνω νοῦν Κυρίου; καὶ τίς αὐτοῦ σύμβουλος ἐγένετο, ὃς συμβιβᾷ αὐτόν;
\vs{14}Ἢ πρὸς τίνα συνεβουλεύσατο, καὶ συνεβίβασεν αὐτόν; ἢ τίς ἔδειξεν αὐτῷ κρίσιν; ἢ ὁδὸν συνέσεως τίς ἔδειξεν αὐτῷ;
\vs{15}Εἰ πάντα τὰ ἔθνη ὡς σταγὼν ἀπὸ κάδου, καὶ ὡς ῥοπὴ ζυγοῦ ἐλογίσθησαν, ὡς σίελος λογισθήσονται;
\vs{16}Ὁ δὲ Λίβανος οὐχ ἱκανὸς εἰς καῦσιν, καὶ πάντα τὰ τετράποδα οὐχ ἱκανὰ εἰς ὁλοκάρπωσιν,
\vs{17}καὶ πάντα τὰ ἔθνη ὡς οὐδέν εἰσι, καὶ εἰς οὐθὲν ἐλογίσθησαν.

\vs{18}Τίνι ὡμοιώσατε Κύριον; καὶ τίνι ὁμοιώματι ὡμοιώσατε αὐτόν;
\vs{19}Μὴ εἰκόνα ἐποίησε τέκτων, ἢ χρυσοχόος χωνεύσας χρυσίον περιεχρύσωσεν αὐτὸν, ὁμοίωμα κατεσκεύασεν αὐτόν;
\vs{20}Ξύλον γὰρ ἄσηπτον ἐκλέγεται τέκτων, καὶ σοφῶς ζητήσει πῶς στήσει εἰκόνα αὐτοῦ, καὶ ἵνα μὴ σαλεύηται.
\vs{21}Οὐ γνώσεσθε; οὐκ ἀκούσεσθε; οὐκ ἀνηγγέλη ἐξ ἀρχῆς ὑμῖν; οὐκ ἔγνωτε τὰ θεμέλια τῆς γῆς;
\vs{22}Ὁ κατέχων τὸν γῦρον τῆς γῆς, καὶ οἱ ἐνοικοῦντες ἐν αὐτῇ ὡς ἀκρίδες· ὁ στήσας ὡς καμάραν τὸν οὐρανὸν, καὶ διατείνας ὡς σκηνὴν κατοικεῖν·
\vs{23}Ὁ διδοὺς ἄρχοντας ὡς οὐδὲν ἄρχειν, τὴν δὲ γῆν ὡς οὐδὲν ἐποίησεν.
\vs{24}Οὐ γὰρ μὴ φυτεύσωσιν, οὐδὲ μὴ σπείρωσιν, οὐδὲ μὴ ῥιζωθῇ εἰς τὴν γῆν ἡ ῥίζα αὐτῶν· ἔπνευσεν ἐπʼ αὐτοὺς, καὶ ἐξηράνθησαν, καὶ καταιγὶς ὡς φρύγανα λήψεται αὐτούς.

\vs{25}Νῦν οὖν τίνι με ὡμοιώσατε, καὶ ὑψωθήσομαι; εἶπεν ὁ ἅγιος.
\vs{26}Ἀναβλέψατε εἰς ὕψος τοὺς ὀφθαλμοὺς ὑμῶν, καὶ ἴδετε, τίς κατέδειξε ταῦτα πάντα; ὁ ἐκφέρων κατʼ ἀριθμὸν τὸν κόσμον αὐτοῦ, πάντας ἐπʼ ὀνόματι καλέσει ἀπὸ πολλῆς δόξης, καὶ ἐν κράτει ἰσχύος αὐτοῦ· οὐδέν σε ἔλαθε.

\vs{27}Μὴ γὰρ εἴπῃς Ἰακὼβ, καὶ τί ἐλάλησας Ἰσραήλ; ἀπεκρύβη ἡ ὁδός μου ἀπὸ τοῦ Θεοῦ, καὶ ὁ Θεός μου τὴν κρίσιν ἀφεῖλε, καὶ ἀπέστη.
\vs{28}Καὶ νῦν οὐκ ἔγνως; εἰ μὴ ἤκουσας; Θεὸς αἰώνιος, ὁ Θεὸς ὁ κατασκευάσας τὰ ἄκρα τῆς γῆς· οὐ πεινάσει, οὐδὲ κοπιάσει, οὐδὲ ἔστιν ἐξεύρεσις τῆς φρονήσεως αὐτοῦ,
\vs{29}διδοὺς τοῖς πεινῶσιν ἰσχὺν, καὶ τοῖς μὴ ὀδυνωμένοις λύπην.
\vs{30}Πεινάσουσι γὰρ νεώτεροι, καὶ κοπιάσουσι νεανίσκοι, καὶ ἐκλεκτοὶ ἀνίσχυες ἔσονται.
\vs{31}Οἱ δὲ ὑπομένοντες τὸν Θεὸν, ἀλλάξουσιν ἰσχὺν, πτεροφυήσουσιν ὡς ἀετοὶ, δραμοῦνται καὶ οὐ κοπιάσουσι, βαδιοῦνται καὶ οὐ πεινάσουσιν.

\ch{41}
Ἐγκαινίζεσθε πρὸς μὲ νῆσοι, οἱ γὰρ ἄρχοντες ἀλλάξουσιν ἰσχύν· ἐγγισάτωσαν καὶ λαλησάτωσαν ἅμα, τότε κρίσιν ἀναγγειλάτωσαν.

\vs{2}Τίς ἐξήγειρεν ἀπὸ ἀνατολῶν δικαιοσύνην, ἐκάλεσεν αὐτὴν κατὰ πόδας αὐτοῦ, καὶ πορεύσεται; δώσει ἐναντίον ἐθνῶν, καὶ βασιλεῖς ἐκστήσει· καὶ δώσει εἰς γῆν τὰς μαχαίρας αὐτῶν, καὶ ὡς φρύγανα ἐξωσμένα τὰ τόξα αὐτῶν.
\vs{3}Καὶ διώξεται αὐτοὺς, διελεύσεται ἐν εἰρήνῃ ἡ ὁδὸς τῶν ποδῶν αὐτοῦ.
\vs{4}Τίς ἐνήργησε, καὶ ἐποίησε ταῦτα; ἐκάλεσεν αὐτὴν ὁ καλῶν αὐτὴν ἀπὸ γενεῶν ἀρχῆς· ἐγὼ Θεὸς πρῶτος, καὶ εἰς τὰ ἐπερχόμενα ἐγώ εἰμι.

\vs{5}Εἴδοσαν ἔθνη καὶ ἐφοβηθήσαν, τὰ ἄκρα τῆς γῆς ἤγγισαν, καὶ ἦλθον ἅμα,
\vs{6}κρίνων ἕκαστος τῷ πλησίον, καὶ τῷ ἀδελφῷ βοηθῆσαι·
\vs{7}καὶ ἐρεῖ, ἴσχυσεν ἀνὴρ τέκτων, καὶ χαλκεὺς τύπτων σφύρῃ, ἅμα ἐλαύνων· πότε μὲν ἐρεῖ, σύμβλημα καλόν ἐστιν, ἰσχύρωσαν αὐτὰ ἐν ἥλοις, θήσουσιν αὐτὰ, καὶ οὐ κινηθήσονται.

\vs{8}Σὺ δὲ, Ἰσραὴλ παῖς μου Ἰακὼβ, καὶ ὃν ἐξελεξάμην, σπέρμα Ἁβραὰμ, ὃν ἠγάπησα·
\vs{9}Οὗ ἀντελαβόμην ἀπʼ ἄκρων τῆς γῆς, καὶ ἐκ τῶν σκοπιῶν αὐτῆς ἐκάλεσά σε, καὶ εἶπά σοι, παῖς μου εἶ, ἐξελεξάμην σε, καὶ οὐκ ἐγκατέλιπόν σε.
\vs{10}Μὴ φοβοῦ, μετὰ σοῦ γάρ εἰμι, μὴ πλανῶ· ἐγὼ γάρ εἰμι ὁ Θεός σου, ὁ ἐνισχύσας σε, καὶ ἐβοήθησά σοι, καὶ ἠσφαλισάμην σε τῇ δεξιᾷ τῇ δικαίᾳ μου.

\vs{11}Ἰδοὺ αἰσχυνθήσονται καὶ ἐντραπήσονται πάντες οἱ ἀντικείμενοί σοι· ἔσονται γὰρ ὡς οὐκ ὄντες, καὶ ἀπολοῦνται πάντες οἱ ἀντίδικοί σου.
\vs{12}Ζητήσεις αὐτοὺς, καὶ οὐ μὴ εὕρῃς τοὺς ἀνθρώπους οἳ παροινήσουσιν εἰς σέ· ἔσονται γὰρ ὡς οὐκ ὄντες, καὶ οὐκ ἔσονται οἱ ἀντιπολεμοῦντές σε·
\vs{13}Ὅτι ἐγὼ ὁ Θεός σου, ὁ κρατῶν τῆς δεξιᾶς σου, ὁ λέγων σοι,
\vs{14}μὴ φοβοῦ Ἰακὼβ ὀλιγοστὸς Ἰσραὴλ· ἐγὼ ἐβοήθησά σοι, λέγει ὁ Θεὸς σου, ὁ λυτρούμενὸσσε Ἰσραήλ.
\vs{15}Ἰδοὺ ἐποίησά σε ὡς τροχοὺς ἁμάξης ἀλοῶντας καινοὺς πριστηροειδεῖς, καὶ ἀλοήσεις ὄρη, καὶ λεπτυνεῖς βουνοὺς, καὶ ὡς χνοῦν θήσεις,
\vs{16}καὶ λικήσεις, καὶ ἄνεμος λήμψεται αὐτοὺς, καὶ καταιγὶς διασπερεῖ αὐτούς· σὺ δὲ εὐφρανθήσῃ ἐν τοῖς ἁγίοις Ἰσραήλ.

\vs{17}Καὶ ἀγαλλιάσονται οἱ πτωχοὶ καὶ οἱ ἐνδεεῖς· ζητήσουσι γὰρ ὕδωρ, καὶ οὐκ ἔσται, ἡ γλῶσσα αὐτῶν ἀπὸ τῆς δίψης ἐξηράνθη· ἐγὼ Κύριος ὁ Θεὸς, ἐγὼ ἐπακούσομαι ὁ Θεὸς Ἰσραὴλ, καὶ οὐκ ἐγκαταλείψω αὐτοὺς,
\vs{18}ἀλλὰ ἀνοίξω ἐπὶ τῶν ὀρέων ποταμοὺς, καὶ ἐν μέσῳ πεδίων πηγάς· ποιήσω τὴν ἔρημον εἰς ἕλη ὑδάτων, καὶ τὴν διψῶσαν γῆν ἐν ὑδραγωγοῖς.
\vs{19}Θήσω εἰς τὴν ἄνυδρον γῆν, κέδρον καὶ πύξον, μυρσίνην καὶ κυπάρισσον, καὶ λεύκην·
\vs{20}Ἵνα ἴδωσι καὶ γνῶσι, καὶ ἐννοηθῶσι καὶ ἐπιστῶνται ἅμα, ὅτι χεὶρ Κυρίου ἐποίησε ταῦτα, καὶ ὁ ἅγιος τοῦ Ἰσραὴλ κατέδειξεν.

\vs{21}Ἐγγίζει ἡ κρίσις ὑμῶν, λέγει Κύριος ὁ Θεός· ἤγγισαν αἱ βουλαὶ ὑμῶν, λέγει ὁ βασιλεὺς Ἰακώβ.
\vs{22}Ἐγγισάτωσαν, καὶ ἀναγγειλάτωσαν ὑμῖν ἃ συμβήσεται, ἢ τὰ πρότερον τίνα ἦν, εἴπατε, καὶ ἐπιστήσομεν τὸν νοῦν, καὶ γνωσόμεθα τί τὰ ἔσχατα καὶ τὰ ἐπερχόμενα·
\vs{23}εἴπατε ἡμῖν, ἀναγγείλατε ἡμῖν τὰ ἐπερχόμενα ἐπʼ ἐσχάτου, καὶ γνωσόμεθα ὅτι θεοί ἐστε· εὐποιήσατε καὶ κακώσατε, καὶ θαυμασόμεθα, καὶ ὀψόμεθα ἅμα
\vs{24}ὅτι πόθεν ἐστὲ ὑμεῖς, καὶ πόθεν ἡ ἐργασία ὑμῶν· ἐκ γῆς βδέλυγμα ἐξελέξαντο ὑμᾶς.

\vs{25}Ἐγὼ δὲ ἤγειρα τὸν ἀπὸ Βοῤῥᾶ, καὶ τὸν ἀφʼ ἡλίου ἀνατολῶν· κληθήσονται τῷ ὀνόματί μου· ἐρχέσθωσαν ἄρχοντες, καὶ ὡς πηλὸς κεραμέως, καὶ ὡς κεραμεὺς καταπατῶν τὸν πηλὸν, οὕτω καταπατηθήσεσθε.
\vs{26}Τίς γὰρ ἀναγγελεῖ τὰ ἐξ ἀρχῆς, ἵνα γνῶμεν καὶ τὰ ἔμπροσθεν, καἰ ἐροῦμεν ὅτι ἀληθῆ ἐστιν; οὐκ ἔστιν ὁ προλέγων, οὐδὲ ὁ ἀκούων ὑμῶν τοὺς λόγους.
\vs{27}Ἀρχὴν Σιὼν δώσω, καὶ Ἱερουσαλὴμ παρακαλέσω εἰς ὁδόν.
\vs{28}Ἀπὸ γὰρ τῶν ἐθνῶν, ἰδοὺ οὐδείς· καὶ ἀπὸ τῶν εἰδώλων αὐτῶν οὐκ ἦν ὁ ἀναγγέλλων· καὶ ἐὰν ἐρωτήσω αὐτοὺς, πόθεν ἐστέ; οὐ μὴ ἀποκριθῶσί μοι.
\vs{29}Εἰσὶ γὰρ οἱ ποιοῦντες ὑμᾶς, καὶ μάτην οἱ πλανῶντες ὑμᾶς.

\ch{42}
Ἰακὼβ ὁ παῖς μου, ἀντιλήψομαι αὐτοῦ· Ἰσραὴλ ὁ ἐκλεκτός μου, προσεδέξατο αὐτὸν ἡ ψυχή μου· ἔδωκα τὸ πνεῦμά μου ἐπʼ αὐτὸν, κρίσιν τοῖς ἔθνεσιν ἐξοίσει.
\vs{2}Οὐ κεκράξεται, οὐδὲ ἀνήσει, οὐδὲ ἀκουσθήσεται ἔξω ἡ φωνὴ αὐτοῦ.
\vs{3}Κάλαμον τεθλασμένον οὐ συντρίψει, καὶ λίνον καπνιζόμενον οὐ σβέσει, ἀλλὰ εἰς ἀλήθειαν ἐξοίσει κρίσιν.
\vs{4}Ἀναλάμψει, καὶ οὐ θραυσθήσεται, ἕως ἂν θῇ ἐπὶ τῆς γῆς κρίσιν, καὶ ἐπὶ τῷ ὀνόματι αὐτοῦ ἔθνη ἐλπιοῦσιν.

\vs{5}Οὕτω λέγει Κύριος ὁ Θεὸς, ὁ ποιήσας τὸν οὐρανὸν, καὶ πήξας αὐτὸν, ὁ στερεώσας τὴν γῆν, καὶ τὰ ἐν αὐτῇ, καὶ διδοὺς πνοὴν τῷ λαῷ τῷ ἐπʼ αὐτῆς, καὶ πνεῦμα τοῖς πατοῦσιν αὐτήν·
\vs{6}Ἐγὼ Κύριος ὁ Θεὸς ἐκάλεσά σε ἐν δικαιοσύνῃ, καὶ κρατήσω τῆς χειρός σου, καὶ ἐνισχύσω σε, καὶ ἔδωκά σε εἰς διαθήκην γένους, εἰς φῶς ἐθνῶν,
\vs{7}ἀνοῖξαι ὀφθαλμοὺς τυφλῶν, ἐξαγαγεῖν ἐκ δεσμῶν δεδεμένους καὶ ἐξ οἴκου φυλακῆς, καὶ καθημένους ἐν σκότει.

\vs{8}Ἐγὼ Κύριος ὁ Θεός, τοῦτό μου ἐστὶ τὸ ὄνομα, τὴν δόξαν μου ἑτέρῳ οὐ δώσω, οὐδὲ τὰς ἀρετάς μου τοῖς γλυπτοῖς.
\vs{9}Τὰ ἀπʼ ἀρχῆς ἰδοὺ ἥκασι, καὶ καινὰ ἃ ἐγὼ ἀναγγέλλω, καὶ πρὸ τοῦ ἀναγγεῖλαι ἐδηλώθη ὑμῖν.

\vs{10}Ὑμνήσατε τῷ Κυρίῳ ὕμνον καινόν· ἡ ἀρχὴ αὐτοῦ, δοξάζετε τὸ ὄνομα αὐτοῦ ἀπʼ ἄκρου τῆς γῆς, οἱ καταβαίνοντες εἰς τὴν θάλασσαν, καὶ πλέοντες αὐτὴν, αἱ νῆσοι καὶ οἱ κατοικοῦντες αὐτάς.
\vs{11}Εὐφράνθητι ἔρημος, καὶ αἱ κῶμαι αὐτῆς, ἐπαύλεις, καὶ οἱ κατοικοῦντες Κηδάρ· εὐφρανθήσονται οἱ κατοικοῦντες πέτραν, ἀπʼ ἄκρου τῶν ὀρέων βοήσουσι,
\vs{12}δώσουσι τῷ Θεῷ δόξαν, τὰς ἀρετὰς αὐτοῦ ἐν ταῖς νήσοις ἀναγγελοῦσι.

\vs{13}Κύριος ὁ Θεὸς τῶν δυνάμεων ἐξελεύσεται, καὶ συντρίψει πόλεμον, ἐπεγερεῖ ζῆλον, καὶ βοήσεται ἐπὶ τοὺς ἐχθροὺς αὐτοῦ μετὰ ἰσχύος.
\vs{14}Ἐσιώπησα, μὴ καὶ ἀεὶ σιωπήσομαι καὶ ἀνέξομαι; ὡς ἡ τίκτουσα ἐκαρτέρησα, ἐκστήσω καὶ ξηρανῶ ἅμα·
\vs{15}Ἐρημώσω ὄρη καὶ βουνοὺς, καὶ πάντα χόρτον αὐτῶν ξηρανῶ· καὶ θήσω ποταμοὺς εἰς νήσους, καὶ ἕλη ξηρανῶ.
\vs{16}Καὶ ἄξω τυφλοὺς ἐν ὁδῷ ᾗ οὐκ ἔγνωσαν, καὶ τρίβους ἃς οὐκ ᾔδεισαν, πατῆσαι ποιήσω αὐτούς· ποιήσω αὐτοῖς τὸ σκότος εἰς φῶς, καὶ τὰ σκολιὰ εἰς εὐθεῖαν· ταῦτα τὰ ῥήματα ποιήσω, καὶ οὐκ ἐγκαταλείψω αὐτούς·
\vs{17}Αὐτοὶ δὲ ἀπεστράφησαν εἰς τὰ ὀπίσω· αἰσχύνθητε αἰσχύνην οἱ πεποιθότες ἐπὶ τοῖς γλυπτοῖς, οἱ λέγοντες τοῖς χωνευτοῖς, ὑμεῖς ἐστε θεοὶ ἡμῶν.

\vs{18}Οἱ κωφοὶ ἀκούσατε, καὶ οἱ τυφλοὶ ἀναβλέψατε ἰδεῖν.
\vs{19}Καὶ τίς τυφλὸς ἀλλʼ ἢ οἱ παῖδές μου, καὶ κωφοὶ ἀλλʼ ἢ οἱ κυριεύοντες αὐτῶν; καὶ ἐτυφλώθησαν οἱ δοῦλοι τοῦ Θεοῦ.
\vs{20}Εἴδετε πλεονάκις, καὶ οὐκ ἐφυλάξασθε· ἠνοιγμένα τὰ ὦτα. καὶ οὐκ ἠκούσατε.
\vs{21}Κύριος ὁ Θεὸς ἐβουλεύσατο ἵνα δικαιωθῇ, καὶ μεγαλύνῃ αἴνεσιν.
\vs{22}Καὶ εἶδον, καὶ ἐγένετο ὁ λαὸς πεπρονομευμένος, καὶ διηρπασμένος· ἡ γὰρ παγὶς ἐν τοῖς ταμείοις πανταχοῦ, καὶ ἐν οἴκοις ἅμα, ὅπου ἔκρυψαν αὐτούς· ἐγένοντο εἰς προνομὴν, καὶ οὐκ ἦν ἐξαιρούμενος ἅρπαγμα, καὶ οὐκ ἦν ὁ λέγων, ἀπόδος.

\vs{23}Τίς ἐν ὑμῖν ὃς ἐνωτιεῖται ταῦτα; εἰσακούσατε εἰς τὰ ἐπερχόμενα.
\vs{24}Οἷς ἔδωκεν εἰς διαρπαγὴν Ἰακὼβ καὶ Ἰσραὴλ τοῖς προνομεύουσιν αὐτόν; οὐχὶ ὁ Θεὸς ᾧ ἡμάρτοσαν αὐτῷ, καὶ οὐκ ἠβούλοντο ἐν ταῖς ὁδοῖς αὐτοῦ πορεύεσθαι, οὐδὲ ἀκούειν τοῦ νόμου αὐτοῦ;
\vs{25}Καὶ ἐπήγαγεν ἐπʼ αὐτοὺς ὀργὴν θυμοῦ αὐτοῦ, καὶ κατίσχυσεν αὐτοὺς πόλεμος, καὶ οἱ συμφλέγοντες αὐτοὺς κύκλῳ, καὶ οὐκ ἔγνωσαν ἕκαστος αὐτῶν, οὐδὲ ἔθεντο ἐπὶ ψυχήν.

\ch{43}
Καὶ νῦν οὕτως λέγει Κύριος ὁ Θεὸς ὁ ποιήσας σε Ἰακὼβ, καὶ ὁ πλάσας σε Ἰσραὴλ, μὴ φοβοῦ, ὅτι ἐλυτρωσάμην σε, ἐκάλεσά σε τὸ ὄνομά σου· ἐμὸς εἶ σύ.
\vs{2}Καὶ ἐὰν διαβαίνῃς διʼ ὕδατος, μετὰ σοῦ εἰμι, καὶ ποταμοὶ οὐ συνκλύσουσί σε· καὶ ἐὰν διέλθῃς διὰ πυρὸς, οὐ μὴ κατακαυθῇς, φλὸξ οὐ κατακαύσει σε.
\vs{3}Ὅτι ἐγὼ Κύριος ὁ Θεός σου ὁ ἅγιος Ἰσραὴλ, ὁ σώζων σε· ἐποίησα ἄλλαγμά σου Αἴγυπτον καὶ Αἰθιοπίαν, καὶ Σοήνην ὑπὲρ σοῦ.
\vs{4}Αφʼ οὗ ἔντιμος ἐγένου ἐναντίον ἐμοῦ, ἐδοξάσθης, καὶ ἐγώ σε ἠγάπησα, καὶ δώσω ἀνθρώπους ὑπὲρ σοῦ, καὶ ἄρχοντας ὑπὲρ τῆς κεφαλῆς σου.
\vs{5}Μὴ φοβοῦ, ὅτι μετὰ σοῦ εἰμι· ἀπὸ ἀνατολῶν ἄξω τὸ σπέρμα σου, καὶ ἀπὸ δυσμῶν συνάξω σε.
\vs{6}Ἐρῶ τῷ Βοῤῥᾷ, ἄγε, καὶ τῷ Λιβὶ, μὴ κώλυε· ἄγε τοὺς υἱούς μου ἀπὸ τῆς πόῤῥωθεν, καὶ τὰς θυγατέρας μου ἀπʼ ἄκρων τῆς γῆς,
\vs{7}πάντας ὅσοι ἐπικέκληνται τῷ ὀνόματί μου· ἐν γὰρ τῇ δόξῃ μου κατεσκεύασα αὐτὸν, καὶ ἔπλασα αὐτὸν, καὶ ἐποίησα αὐτὸν,
\vs{8}καὶ ἐξήγαγον λαὸν τυφλὸν, καὶ ὀφθαλμοί εἰσιν ὡσαύτως τυφλοὶ, καὶ κωφοὶ τὰ ὦτα ἔχοντες.

\vs{9}Πάντα τὰ ἔθνη συνήχθησαν ἅμα, καὶ συναχθήσονται ἄρχοντες ἐξ αὐτῶν· τίς ἀναγγελεῖ ταῦτα; ἢ τὰ ἐξ ἀρχῆς τίς ἀναγγελεῖ ὑμῖν; ἀγαγέτωσαν τοῦς μάρτυρας αὐτῶν καὶ δικαιωθήτωσαν, καὶ ἀκουσάτωσαν, καὶ εἰπάτωσαν ἀληθῆ.

\vs{10}Γένεσθέ μοι μάρτυρες, καὶ ἐγὼ μάρτυς, λέγει Κύριος ὁ Θεὸς, καὶ ὁ παῖς μου ὃν ἐξελεξάμην, ἵνα γνῶτε καὶ πιστεύσητε, καὶ συνῆτε ὅτι ἐγώ εἰμι· ἔμπροσθέν μου οὐκ ἐγένετο ἄλλος Θεὸς, καὶ μετʼ ἐμὲ οὐκ ἔσται.
\vs{11}Ἐγὼ ὁ Θεὸς, καὶ οὐκ ἔστι πάρεξ ἐμοῦ σώζων.
\vs{12}Ἐγὼ ἀνήγγειλα καὶ ἔσωσα, ὠνείδισα καὶ οὐκ ἦν ἐν ὑμῖν ἀλλότριος· ὑμεῖς ἐμοὶ μάρτυρες, καὶ ἐγὼ Κύριος ὁ Θεὸς
\vs{13}ἔτι ἀπʼ ἀρχῆς, καὶ οὐκ ἔστιν ὁ ἐκ τῶν χειρῶν μου ὁ ἐξαιρούμενος· ποιήσω, καὶ τίς ἀποστρέψει αὐτό;

\vs{14}Οὕτως λέγει Κύριος ὁ Θεὸς ὁ λυτρούμενος ὑμᾶς, ὁ ἅγιος τοῦ Ἰσραὴλ, ἕνεκεν ὑμῶν ἀποστελῶ εἰς Βαβυλῶνα, καὶ ἐπεγερῶ φεύγοντας πάντας, καὶ Χαλδαῖοι ἐν πλοίοις δεθήσονται.
\vs{15}Ἐγὼ Κύριος ὁ Θεὸς ὁ ἅγιος ὑμῶν, ὁ καταδείξας Ἰσραὴλ βασιλέα ὑμῶν.

\vs{16}Οὕτως λέγει Κύριος, ὁ διδοὺς ἐν θαλάσσῃ ὁδὸν, καὶ ἐν ὕδατι ἰσχυρῷ τρίβον,
\vs{17}ὁ ἐξαγαγὼν ἅρματα καὶ ἵππον καὶ ὄχλον ἰσχυρον· ἀλλʼ ἐκοιμήθησαν, καὶ οὐκ ἀναστήσονται, ἐσβέσθησαν ὡς λῖνον ἐσβεσμένον.

\vs{18}Μὴ μνημονεύετε τὰ πρῶτα, καὶ τὰ ἀρχαῖα μὴ συλλογίζεσθε.
\vs{19}Ἰδοὺ ἐγὼ ποιῶ καινὰ, ἃ νῦν ἀνατελεῖ, καὶ γνώσεσθε αὐτά· καὶ ποιήσω ἐν τῇ ἐρήμῳ ὁδὸν, καὶ ἐν τῇ ἀνύδρῳ ποταμούς.
\vs{20}Εὐλογήσουσί με τὰ θηρία τοῦ ἀγροῦ, σειρῆνες, καὶ θυγατέρες στρουθῶν, ὅτι ἔδωκα ἐν τῇ ἐρήμῳ ὕδωρ, καὶ ποταμοὺς ἐν τῇ ἀνύδρῳ, ποτίσαι τὸ γένος μου τὸ ἐκλεκτὸν,
\vs{21}λαόν μου ὃν περιεποιησάμην τὰς ἀρετάς μου διηγεῖσθαι.

\vs{22}Οὐ νῦν ἐκάλεσά σε Ἰακὼβ, οὐδὲ κοπιάσαι σε ἐποίησα Ἰσραήλ.
\vs{23}Οὐκ ἤνεγκάς μοι πρόβατά σου τῆς ὁλοκαρπώσεώς σου, οὐδὲ ἐν ταῖς θυσίαις σου ἐδόξασάς με· οὐκ ἐδούλωσά σε ἐν θυσίαις, οὐδὲ ἔγκοπον ἐποίησά σε ἐν λιβάνῳ,
\vs{24}οὐδὲ ἐκτήσω μοι ἀργυρίου θυσίασμα, οὐδὲ τὸ στέαρ τῶν θυσιῶν σου ἐπεθύμησα· ἀλλὰ ἐν ταῖς ἁμαρτίαις σου προέστης μου, καὶ ἐν ταῖς ἀδικίαις σου.
\vs{25}Ἐγώ εἰμι ἐγώ εἰμι ὁ ἐξαλείφων τὰς ἀνομίας σου ἕνεκεν ἐμοῦ, καὶ τὰς ἁμαρτίας σου, καὶ οὐ μὴ μνησθήσομαι.
\vs{26}Σὺ δὲ μνήσθητι, καὶ κριθῶμεν· λέγε σὺ τὰς ἀνομίας σου πρῶτος, ἵνα δικαιωθῇς.
\vs{27}Οἱ πατέρες ὑμῶν πρῶτοι, καὶ οἱ ἄρχοντες ὑμῶν ἠνόμησαν εἰς ἐμέ.
\vs{28}Καὶ ἐμίαναν οἱ ἄρχοντες τὰ ἅγιά μου· καὶ ἔδωκα ἀπωλέσαι Ἰακὼβ, καὶ Ἰσραὴλ εἰς ὀνειδισμόν.

\ch{44}
Νῦν δὲ ἄκουσον Ἰακὼβ ὁ παῖς μου, καὶ Ἰσραὴλ ὃν ἐξελεξάμην.
\vs{2}Οὕτω λέγει Κύριος ὁ Θεὸς ὁ ποιήσας σε, καὶ ὁ πλάσας σε ἐκ κοιλίας, ἔτι βοηθηθήσῃ· μὴ φοβοῦ παῖς μου Ἰακὼβ, καὶ ἠγαπημένος Ἰσραὴλ ὃν ἐξελεξάμην.
\vs{3}Ὅτι ἐγὼ δώσω ὕδωρ ἐν δίψει τοῖς πορευομένοις ἐν ἀνύδρῳ, ἐπιθήσω τὸ πνεῦμά μου ἐπὶ τὸ σπέρμα σου, καὶ τὰς εὐλογίας μου ἐπὶ τὰ τέκνα σου,
\vs{4}καὶ ἀνατελοῦσιν ὡς ἀναμέσον ὕδατος χόρτος, καὶ ὡς ἰτέα ἐπὶ παραῤῥέον ὕδωρ.
\vs{5}Οὗτος ἐρεῖ, τοῦ Θεοῦ εἰμι, καὶ οὗτος βοήσεται ἐπὶ τῷ ὀνόματι Ἰακώβ· καὶ ἕτερος ἐπιγράψει χειρὶ αὐτοῦ, τοῦ Θεοῦ εἰμι, καὶ ἐπὶ τῷ ὀνόματι Ἰσραὴλ βοήσεται.

\vs{6}Οὕτως λέγει ὁ Θεὸς ὁ βασιλεὺς Ἰσραὴλ, καὶ ῥυσάμενος αὐτὸν Θεὸς σαβαὼθ, ἐγὼ πρῶτος, καὶ ἐγὼ μετὰ ταῦτα· πλὴν ἐμοῦ οὐκ ἔστι Θεός.
\vs{7}Τίς ὥσπερ ἐγὼ; στήτω, καὶ καλεσάτω, καὶ ἀναγγειλάτω, καὶ ἑτοιμασάτω μοι ἀφʼ οὗ ἐποίησα ἄνθρωπον εἰς τὸν αἰῶνα, καὶ τὰ ἐπερχόμενα πρὸ τοῦ ἐλθεῖν ἀναγγειλάτωσαν ὑμῖν.
\vs{8}Μὴ παρακαλύπτεσθε, μηδὲ πλανᾶσθε· οὐκ ἀπʼ ἀρχῆς ἠνωτίσασθε, καὶ ἀπήγγειλα ὑμῖν; μάρτυρες ὑμεῖς ἐστε, εἰ ἕστι Θεὸς πλὴν ἐμοῦ.

\vs{9}Καὶ οὐκ ἤκουσαν τότε οἱ πλάσσοντες· καὶ οἱ γλύφοντες, πάντες μάταιοι, ποιοῦντες τὰ καταθύμια αὐτῶν, ἃ οὐκ ὠφελήσει αὐτούς· ἀλλὰ αἰσχυνθήσονται
\vs{10}οἱ πλάσσοντες θεὸν, καὶ γλύφοντες πάντες ἀνωφελῆ,
\vs{11}καὶ πάντες ὅθεν ἐγένοντο ἐξηράνθησαν· καὶ κωφοὶ ἀπὸ ἀνθρώπων συναχθήτωσαν πάντες, καὶ στησάτωσαν ἅμα· καὶ ἐντραπήτωσαν, καὶ αἰσχυνθήτωσαν ἅμα·

\vs{12}Ὅτι ὤξυνε τέκτων σίδηρον· σκεπάρνῳ εἰργάσατο αὐτὸ, καὶ ἐν τερέτρῳ ἔστησεν αὐτὸ, καὶ εἰργάσατο αὐτὸ ἐν τῷ βραχίονι τῆς ἰσχύος αὐτοῦ, καὶ πεινάσει, καὶ ἀσθενήσει, καὶ οὐ μὴ πίῃ ὕδωρ.
\vs{13}Ἐκλεξάμενος τέκτων ξύλον, ἔστησεν αὐτὸ ἐν μέτρῳ, καὶ ἐν κόλλῃ ἐῤῥύθμισεν αὐτὸ, καὶ ἐποίησεν αὐτὸ ὡς μορφὴν ἀνδρὸς, καὶ ὡς ὡραιότητα ἀνθρώπου, στῆσαι αὐτὸ ἐν οἴκῳ.
\vs{14}Ἔκοψε ξύλον ἐκ τοῦ δρυμοῦ, ὃ ἐφύτευσε Κύριος, πίτυν, καὶ ὑετὸς ἐμήκυνεν,
\vs{15}ἵνα ᾖ ἀνθρώποις εἰς καῦσιν· καὶ λαβῶν ἀπʼ αὐτοῦ, ἐθερμάνθη, καὶ καύσαντες ἔπεψαν ἄρτους ἐπʼ αὐτῶν· τὸ δὲ λοιπὸν εἰργάσαντο θεοὺς, καὶ προσκυνοῦσιν αὐτοῖς·
\vs{16}Οὗ τὸ ἥμισυ αὐτοῦ κατέκαυσεν ἐν πυρὶ, καὶ ἐπὶ τοῦ ἡμίσους αὐτοῦ ἔπεψεν ἐν τοῖς ἄνθραξιν ἄρτους, καὶ ἐπʼ αὐτοῦ κρέας ὀπτήσας ἔφαγε, καὶ ἐνεπλήσθη, καὶ θερμανθεὶς εἶπεν, ἡδύ μοι, ὅτι ἐθερμάνθην, καὶ εἶδον πῦρ.
\vs{17}Τὸ δὲ λοιπὸν ἐποίησεν εἰς θεὸν γλυπτὸν, καὶ προσκυνεῖ, καὶ προσεύχεται λέγων, ἐξελοῦ με, ὅτι θεός μου εἶ σύ.

\vs{18}Οὐκ ἔγνωσαν φρονῆσαι, ὅτι ἀπημαυρώθησαν τοῦ βλέπειν τοῖς ὀφθαλμοῖς αὐτῶν, καὶ τοῦ νοῆσαι τῇ καρδίᾳ αὐτῶν.
\vs{19}Καὶ οὐκ ἐλογίσατο τῇ ψυχῇ αὐτοῦ, οὐδὲ ἔγνω τῇ φρονήσει, ὅτι τὸ ἥμισυ αὐτοῦ κατέκαυσεν ἐν πυρὶ, καὶ ἔπεψεν ἐπὶ τῶν ἀνθράκων αὐτοῦ ἄρτους, καὶ ὀπτήσας κρέα ἔφαγε, καὶ τὸ λοιπὸν αὐτοῦ εἰς βδέλυγμα ἐποίησε, καὶ προσκυνοῦσιν αὐτῷ.
\vs{20}Γνῶθι ὅτι σποδὸς ἡ καρδία αὐτῶν, καὶ πλανῶνται, καὶ οὐδεὶς δύναται ἐξελέσθαι τὴν ψυχὴν αὐτοῦ· ἴδετε, οὐκ ἐρεῖτε, ὅτι ψεῦδος ἐν τῇ δεξιᾷ μου.

\vs{21}Μνήσθητι ταῦτα Ἰακὼβ καὶ Ἰσραὴλ, ὅτι παῖς μου εἶ σὺ, ἔπλασά σε παῖδά μου, καὶ σὺ Ἰσραὴλ μὴ ἐπιλανθάνου μοῦ.
\vs{22}Ἰδοὺ γὰρ ἀπήλειψα ὡς νεφέλην τὰς ἀνομίας σου, καὶ ὡς γνόφον τὴν ἁμαρτίαν σου· ἐπιστράφηθι πρὸς μὲ, καὶ λυτρώσομαί σε.

\vs{23}Εὐφράνθητε οὐρανοὶ, ὅτι ἠλέησεν ὁ Θεὸς τὸν Ἰσραήλ· σαλπίσατε τὰ θεμέλια τῆς γῆς, βοήσατε ὄρη εὐφροσύνην, οἱ βουνοὶ καὶ πάντα τὰ ξύλα τὰ ἐν αὐτοῖς, ὅτι ἐλυτρώσατο ὁ Θεὸς τὸν Ἰακὼβ, καὶ Ἰσραὴλ δοξασθήσεται.

\vs{24}Οὕτω λέγει Κύριος ὁ λυτρούμενός σε, καὶ πλάσσων σε ἐκ κοιλίας, ἐγὼ Κύριος ὁ συντελῶν πάντα, ἐξέτεινα τὸν οὐρανὸν μόνος, καὶ ἐστερέωσα τὴν γῆν.
\vs{25}Τίς ἕτερος διασκεδάσει σημεῖα ἐγγαστριμύθων, καὶ μαντείας ἀπὸ καρδίας; ἀποστρέφων φρονίμους εἰς τὰ ὀπίσω, καὶ τὴν βουλὴν αὐτῶν μωραίνων,
\vs{26}καὶ ἱστῶν ῥὴμα παιδὸς αὐτοῦ, καὶ τὴν βουλὴν τῶν ἀγγέλων αὐτοῦ ἀληθεύων· ὁ λέγων τῇ Ἱερουσαλὴμ, κατοικηθήσῃ, καὶ ταῖς πόλεσι τῆς Ἰδουμαίας, οἰκοδομηθήσεσθε, καὶ τὰ ἔρημα αὐτῆς ἀνατελεῖ·
\vs{27}Ὁ λέγων τῇ ἀβύσσῳ, ἐρημωθήσῃ, καὶ τοὺς ποταμούς σου ξηρανῶ·
\vs{28}Ὁ λέγων Κύρῳ φρονεῖν, καὶ πάντα τὰ θελήματά μου ποιήσει· ὁ λέγων Ἱερουσαλὴμ, οἰκοδομηθήσῃ, καὶ τὸν οἶκον τὸν ἅγιόν μου θεμελιώσω.

\ch{45}
Οὕτῳ λέγει Κύριος ὁ Θεὸς τῷ χριστῷ μου Κύρῳ, οὗ ἐκράτησα τῆς δεξιᾶς, ἐπακοῦσαι ἔμπροσθεν αὐτοῦ ἔθνη, καὶ ἰσχὺν βασιλέων διαῤῥήξω, ἀνοίξω ἔμπροσθεν αὐτοῦ θύρας, καὶ πόλεις οὐ συγκλεισθήσονται.
\vs{2}Ἐγὼ ἔμπροσθέν σου πορεύσομαι, καὶ ὄρη ὁμαλιῶ, θύρας χαλκᾶς συντρίψω, καὶ μοχλοὺς σιδηροῦς συνκλάσω.
\vs{3}Καὶ δώσω σοι θησαυροὺς σκοτεινοὺς, ἀποκρύφους, ἀοράτους ἀνοίξω σοι, ἵνα γνῷς ὅτι ἐγὼ Κύριος ὁ Θεός σου ὁ καλῶν τὸ ὄνομά σουὁ, Θεὸς Ἰσραήλ.
\vs{4}Ἕνεκεν τοῦ παιδός μου Ἰακὼβ, καὶ Ἰσραὴλ τοῦ ἐκλεκτοῦ μου, ἐγὼ καλέσω σε τῷ ὀνόματί σου, καὶ προσδέξομαί σε· σὺ δὲ οὐκ ἔγνως με,
\vs{5}ὅτι ἐγὼ Κύριος ὁ Θεὸς, καὶ οὐκ ἔστιν ἔτι πλὴν ἐμοῦ Θεός· ἐνίσχυσά σε, καὶ οὐκ ᾔδεις με,
\vs{6}ἵνα γνῶσιν οἱ ἀπʼ ἀνατολῶν ἡλίου καὶ οἱ ἀπὸ δυσμῶν, ὅτι οὐκ ἔστι Θεὸς πλὴν ἐμοῦ· ἐγὼ Κύριος ὁ Θεὸς, καὶ οὐκ ἔστιν ἔτι.
\vs{7}Ἐγὼ ἡ κατασκευάσας φῶς, καὶ ποιήσας σκότος, ὁ ποιῶν εἰρήνην, καὶ κτίζων κακά· ἐγὼ Κύριος ὁ Θεὸς, ὁ ποιῶν πάντα ταῦτα.

\vs{8}Εὐφρανθήτω ὁ οὐρανὸς ἄνωθεν, καὶ αἱ νεφέλαι ῥανάτωσαν δικαιοσύνην· ἀνατειλάτω ἡ γῆ, καὶ βλαστησάτω ἔλεος, καὶ δικαιοσύνην ἀνατειλάτω ἅμα· ἐγώ εἰμι Κύριος ὁ κτίσας σε.

\vs{9}Ποῖον βέλτιον κατεσκεύασα ὡς πηλὸν κεραμέως; μὴ ὁ ἀροτριῶν ἀροτριάσει τὴν γῆν ὅλην τὴν ἡμέραν; μὴ ἐρεῖ ὁ πηλὸς τῷ κεραμεῖ, τί ποιεῖς ὅτι οὐκ ἐργάζῃ, οὐδὲ ἔχεις χεῖρας; μὴ ἀποκριθήσεται τὸ πλάσμα πρὸς τὸν πλάσαντα αὐτό;
\vs{10}Ὁ λέγων τῷ πατρὶ, τί γεννήσεις; καὶ τῇ μητρί, τί ὠδίνεις;

\vs{11}Ὅτι οὕτω λέγει Κύριος ὁ Θεὸς ὁ ἅγιος Ἰσραὴλ, ὁ ποιήσας τὰ ἐπερχόμενα, ἐρωτήσατέ με περὶ τῶν υἱῶν μου, καὶ περὶ τῶν ἔργων τῶν χειρῶν μου ἐντείλασθέ μοι.
\vs{12}Ἐγὼ ἐποίησα γῆν, καὶ ἄνθρωπον ἐπʼ αὐτῆς, ἐγὼ τῇ χειρί μου ἐστερέωσα τὸν οὐρανὸν, ἐγὼ πᾶσι τοῖς ἄστροις ἐνετειλάμην.
\vs{13}Ἐγὼ ἤγειρα αὐτὸν μετὰ δικαιοσύνης βασιλέα, καὶ πᾶσαι αἱ ὁδοὶ αὐτοῦ εὐθεῖαι· οὗτος οἰκοδομήσει τὴν πόλιν μου, καὶ τὴν αἰχμαλωσίαν τοῦ λαοῦ μου ἐπιστρέψει, οὐ μετὰ λύτρων, οὐδὲ μετὰ δώρων, εἶπε Κύριος σαβαώθ.

\vs{14}Οὕτω λέγει Κύριος σαβαὼθ, ἐκοπίασεν Αἴγυπτος, καὶ ἐμπορία Αἰθιόπων, καὶ οἱ Σαβαεὶμ ἄνδρες ὑψηλοὶ ἐπὶ σὲ διαβήσονται, καὶ σοὶ ἔσονται δοῦλοι, καὶ ὀπίσω σου ἀκολουθήσουσι δεδεμένοι χειροπέδαις, καὶ διαβήσονται πρὸς σὲ, καὶ προσκυνήσουσί σοι, καὶ ἐν σοὶ προσεύξονται· ὅτι ἐν σοὶ ὁ Θεός ἐστι, καὶ οὐκ ἔστι Θεὸς πλήν σου.
\vs{15}Σὺ γὰρ εἶ Θεὸς, καὶ οὐκ ᾔδειμεν, ὁ Θεὸς τοῦ Ἰσραὴλ σωρήρ.
\vs{16}Αἰσχυνθήσονται καὶ ἐντραπήσονται πάντες οἱ ἀντικείμενοι αὐτῷ, καὶ πορεύσονται ἐν αἰσχύνῃ· ἐγκαινίζεσθε πρὸς μὲ νῆσοι.
\vs{17}Ἰσραὴλ σώζεται ὑπὸ Κυρίου σωτηρίαν αἰώνιον· οὐκ αἰσχυνθήσονται, οὐδὲ μὴ ἐντραπῶσιν ἕως τοῦ αἰῶνος ἔτι.

\vs{18}Οὕτως λέγει Κύριος ὁ ποιήσας τὸν οὐρανὸν, οὗτος ὁ Θεὸς ὁ καταδείξας τὴν γῆν, καὶ ποιήσας αὐτὴν, αὐτὸς διώρισεν αὐτὴν, οὐκ εἰς κενὸν ἐποίησεν αὐτὴν, ἀλλὰ κατοικεῖσθαι ἔπλασεν αὐτὴν, ἐγώ εἰμι Κύριος, καὶ οὐκ ἔστιν ἔτι.
\vs{19}Οὐκ ἐν κρυφῇ λελάληκα, οὐδὲ ἐν τόπῳ γῆς σκοτεινῷ· οὐκ εἶπα τῷ σπέρματι Ἰακὼβ, μάταιον ζητήσατε· ἐγώ εἰμι ἐγώ εἰμι Κύριος ὁ λαλῶν δικαιοσύνην, καὶ ἀναγγέλλων ἀλήθειαν.

\vs{20}Συνάχθητε, καὶ ἥκετε, βουλεύσασθε ἅμα οἱ σωζόμενοι ἀπὸ τῶν ἐθνῶν· οὐκ ἔγνωσαν οἱ αἴροντες τὸ ξύλον γλύμμα αὐτῶν, καὶ οἱ προσευχόμενοι πρὸς θεοὺς, οἳ οὐ σώζουσιν.
\vs{21}Εἰ ἀναγγελοῦσιν, ἐγγισάτωσαν, ἵνα γνῶσιν ἅμα, τίς ἀκουστὰ ἐποίησε ταῦτα ἀπʼ ἀρχῆς· τότε ἀνηγγέλη ὑμῖν· ἐγὼ ὁ Θεὸς, καὶ οὐκ ἔστιν ἄλλος πλὴν ἐμοῦ· δίκαιος καὶ σωτὴρ, οὐκ ἔστι πάρεξ ἐμοῦ.
\vs{22}Ἐπιστράφητε ἐπʼ ἐμὲ, καὶ σωθήσεσθε, οἱ ἀπʼ ἐσχάτου τῆς γῆς· ἐγώ εἰμι ὁ Θεὸς, καὶ οὐκ ἔστιν ἄλλος.
\vs{23}Κατʼ ἐμαυτοῦ ὀμνύω, εἰ μὴ ἐξελεύσεται ἐκ τοῦ στόματός μου δικαιοσύνη, οἱ λόγοι μου οὐκ ἀποστραφήσονται· Ὅτι ἐμοὶ κάμψει πᾶν γόνυ, καὶ ὀμεῖται πᾶσα γλῶσσα τὸν Θεὸν,
\vs{24}λέγων, δικαιοσύνη καὶ δόξα πρὸς αὐτὸν ἥξει, καὶ αἰσχυνθήσονται πάντες οἱ διορίζοντες αὐτούς.
\vs{25}Ἀπὸ Κυρίου δικαιωθήσονται, καὶ ἐν τῷ Θεῷ ἐνδοξασθήσεται πᾶν τὸ σπέρμα τῶν υἱῶν Ἰσραήλ.

\ch{46}
Ἔπεσε Βὴλ, συνετρίβη Ναβὼ, ἐγένετο τὰ γλυπτὰ αὐτῶν εἰς θηρία, καὶ τὰ κτήνη· αἴρετε αὐτὰ καταδεδεμένα ὡς φορτίον κοπιῶντι ἐκλευμένῳ,
\vs{2}καὶ πεινῶντι, οὐκ ἰσχύοντι, ἅμα, οἳ οὐ δυνήσονται σωθῆναι ἀπὸ πολέμου, αὐτοὶ δὲ αἰχμάλωτοι ἤχθησαν.

\vs{3}Ἀκούετέ μου οἶκος τοῦ Ἰακὼβ, καὶ πᾶν τὸ κατάλοιπον τοῦ Ἰσραὴλ, οἱ αἰρόμενοι ἐκ κοιλίας, καὶ παιδευόμενοι ἐκ παιδίου ἕως γήρως·
\vs{4}ἐγώ εἰμι, καὶ ἕως ἂν καταγηράσητε, ἐγώ εἰμι, ἐγὼ ἀνέχομαι ὑμῶν, ἐγὼ ἐποίησα, καὶ ἐγὼ ἀνήσω, ἐγὼ ἀναλήμψομαι, καὶ σώσω ὑμᾶς.

\vs{5}Τίνι με ὡμοιώσατε; ἴδετε, τεχνάσασθε, οἱ πλανώμενοι.
\vs{6}Οἱ συμβαλλόμενοι χρυσίον ἐκ μαρσυππίου, καὶ ἀργύριον ἐν ζυγῷ, στήσουσιν ἐν σταθμῷ, καὶ μισθωσάμενοι χρυσοχόον ἐποίησαν χειροποίητα, καὶ κύψαντες προσκυνοῦσιν αὐτοῖς.
\vs{7}Αἴρουσιν αὐτὸ ἐπὶ τοῦ ὤμου, καὶ πορεύονται· ἐὰν δὲ θῶσιν αὐτό ἐπὶ τοῦ τόπου αὐτοῦ, μένει, οὐ μὴ κινηθῇ· καὶ ὃς ἐὰν βοήσῃ πρὸς αὐτὸν, οὐ μὴ εἰσακούσῃ, ἀπὸ κακῶν οὐ μὴ σώσῃ αὐτόν.

\vs{8}Μνήσθητε ταῦτα, καὶ στενάξατε, μετανοήσατε οἱ πεπλανημένοι, ἐπιστρέψατε τῇ καρδίᾳ,
\vs{9}καὶ μνήσθητε τὰ πρότερα ἀπὸ τοῦ αἰῶνος, ὅτι ἐγώ εἰμι ὁ Θεὸς, καὶ οὐκ ἔστιν ἔτι πλὴν ἐμοῦ,
\vs{10}ἀναγγέλλων πρότερον τὰ ἔσχατα πρὶν γενέσθαι, καὶ ἅμα συνετελέσθη· καὶ εἶπα, πᾶσα ἡ βουλή μου στήσεται, καὶ πάντα ὅσα βεβούλευμαι, ποιήσω·
\vs{11}Καλῶν ἀπὸ ἀνατολῶν πετεινὸν, καὶ ἀπὸ γῆς πόῤῥωθεν περὶ ὧν βεβούλευμαι, ἐλάλησα, καὶ ἤγαγον, ἔκτισα καὶ ἐποίησα, ἤγαγον αὐτὸν, καὶ εὐώδωσα τὴν ὁδὸν αὐτοῦ.

\vs{12}Ἀκούσατέ μου οἱ ἀπολωλεκότες τὴν καρδίαν, οἱ μακρὰν ἀπὸ τῆς δικαιοσύνης.
\vs{13}Ἤγγισα τὴν δικαιοσύνην μου, καὶ τὴν σωτηρίαν τὴν παρʼ ἐμοῦ οὐ βραδυνῶ· δέδωκα ἐν Σιὼν σωτηρίαν τῷ Ἰσραὴλ εἰς δόξασμα.

\ch{47}
Κατάβηθι, κάθισον ἐπὶ τὴν γῆν παρθένος θυγάτηρ Βαβυλῶνος, κάθισον εἰς τὴν γῆν θυγάτηρ Χαλδαίων, ὅτι οὐκέτι προστεθήσῃ κληθῆναι ἁπαλὴ καὶ τρυφερά.
\vs{2}Λάβε μύλον, ἄλεσον ἄλευρον, ἀποκὸλυψαι τὸ κατακάλυμμά σου, ἀνακάλυψαι τὰς πολιὰς, ἀνάσυρε τὰς κνήμας, διάβηθι ποταμούς.
\vs{3}Ἀνακαλυφθήσεται ἡ αἰσχύνη σου, φανήσονται οἱ ὀνειδισμοί σου· τὸ δίκαιον ἐκ σοῦ λήμψομαι, οὐκέτι μὴ παραδῶ ἀνθρώποις.

\vs{4}Ὁ ῥυσάμενός σε Κύριος σαβαὼθ, ὄνομα αὐτῷ Ἅγιος Ἰσραήλ.

\vs{5}Κάθισον κατανενυγμένη, εἴσελθε εἰς τὸ σκότος θύγατεηρ Χαλδαίων, οὐκέτι μὴ κληθήσῃ ἰσχὺς βασιλείας.
\vs{6}Παρωξύνθην ἐπὶ τῷ λαῷ μου, ἐμίανας τὴν κληρονομίαν μου· ἐγὼ ἔδωκα αὐτοὺς εἰς τὴν χεῖρά σου, σὺ δὲ οὐκ ἔδωκας αὐτοῖς ἔλεος, τοῦ πρεσβυτέρου ἐβάρυνας τὸν ζυγὸν σφόδρα,
\vs{7}καὶ εἶπας, εἰς τὸν αἰῶνα ἔσομαι ἄρχουσα· οὐκ ἐνόησας ταῦτα ἐν τῇ καρδίᾳ σου, οὐδὲ ἐμνήσθης τὰ ἔσχατα.

\vs{8}Νῦν δὲ ἄκουε ταῦτα τρυφερὰ, ἡ καθημένη, ἡ πεποιθυῖα, ἡ λέγουσα ἐν καρδίᾳ αὐτῆς, ἐγώ εἰμι, καὶ οὐκ ἔστιν ἑτέρα, οὐ καθιῶ χήρα, οὐδὲ γνώσομαι ὀρφανίαν.
\vs{9}Νῦν δὲ ἥξει ἐπὶ σὲ τὰ δύο ταῦτα ἐξαίφνης ἐν ἡμέρᾳ μιᾷ, ἀτεκνία καὶ χηρεία ἥξει ἐξαίφνης ἐπὶ σὲ, ἐν τῇ φαρμακείᾳ σου, ἐν τῇ ἰσχύϊ τῶν ἐπαοιδῶν σου σφόδρα,
\vs{10}τῇ ἐλπίδι τῆς πονηρίας σου· σὺ γὰρ εἶπας, ἐγώ εἰμι, καὶ οὐκ ἔστιν ἑτέρα· γνῶθι, ἡ σύνεσις τούτων ἔσται, καὶ ἡ πορνεία σου σοὶ αἰσχύνη· καὶ εἶπας τῇ καρδίᾳ σου, ἐγώ εἰμι, καὶ οὐκ ἔστιν ἑτέρα.

\vs{11}Καὶ ἥξει ἐπὶ σὲ ἀπώλεια, καὶ οὐ μὴ γνῷς· βόθυνος, καὶ ἐμπεσῇ εἰς αὐτόν· καὶ ἥξει ἐπὶ σὲ ταλαιπωρία, καὶ οὐ μὴ δυνήσῃ καθαρὰ γενέσθαι· καὶ ἥξει ἐπὶ σὲ ἐξαπίνης ἀπώλεια, καὶ οὐ μὴ γνώσῃ.
\vs{12}Στῆθι νῦν ἐν ταῖς ἐπαοιδαῖς σου, καὶ ἐν τῇ πολλῇ φαρμακείᾳ σου, ἃ ἐμάνθανες ἐκ νεότητός σου, εἰ δυνήσῃ ὠφεληθῆναι.
\vs{13}Κεκοπίακας ἐν ταῖς βουλαῖς σου· στήτωσαν δὴ καὶ σωσάτωσάν σε οἱ ἀστρολόγοι τοῦ οὐρανοῦ, οἱ ὁρῶντες τοὺς ἀστέρας ἀναγγειλάτωσάν σοι, τί μέλλει ἐπὶ σὲ ἔρχεσθαι.
\vs{14}Ἰδοὺ πάντες ὡς φρύγανα ἐπὶ πυρὶ κατακαυθήσονται, καὶ οὐ μὴ ἐξέλωνται τὴν ψυχὴν αὐτῶν ἐκ φλογός· ὅτι ἔχεις ἄνθρακας πυρός· κάθισαι ἐπʼ αὐτοὺς,
\vs{15}οὗτοι ἔσονταί σοι βοήθεια· ἐκοπίασας ἐν τῇ μεταβολῇ ἐκ νεότητος, ἄνθρωπος καθʼ ἑαυτὸν ἐπλανήθη, σοὶ δὲ οὐκ ἔσται σωτηρία.

\ch{48}
Ἀκούσατε ταῦτα οἶκος Ἰακὼβ, οἱ κεκλημένοι ἐπὶ τῷ ὀνόματι Ἰσραὴλ, καὶ ἐκ Ἰούδα ἐξελθόντες, οἱ ὀμνύοντες τῷ ὀνόματι Κυρίου Θεοῦ Ἰσραὴλ, μιμνησκόμενοι οὐ μετὰ ἀληθείας, οὐδὲ μετὰ δικαιοσύνης,
\vs{2}καὶ ἀντεχόμενοι τῷ ὀνόματι τῆς πόλεως τῆς ἁγίας, καὶ ἐπὶ τῷ Θεῷ Ἰσραὴλ ἀντιστηριζόμενοι· Κύριος σαβαὼθ ὄνομα αὐτῷ.
\vs{3}Τὰ πρότερα ἔτι ἀνήγγειλα, καὶ ἐκ τοῦ στόματός μου ἐξῆλθε, καὶ ἀκουστὸν ἐγένετο· ἐξάπινα ἐποίησα, καὶ ἐπῆλθε.

\vs{4}Γινώσκω ὅτι σκληρὸς εἶ, καὶ νεῦρον σιδηροῦν ὁ τράχηλός σου, καὶ τὸ μέτωπόν σου χαλκοῦν.
\vs{5}Καὶ ἀνήγγειλά σοι πάλαι ἃ πρὶν ἐλθεῖν ἐπὶ σέ· ἀκουστόν σοι ἐποίησα, μή ποτε εἴπῃς, ὅτι τὰ εἴδωλά μοι ἐποίησε, καὶ εἴπῃς, ὅτι τὰ γλυπτὰ καὶ τὰ χωνευτὰ ἐνετείλατό μοι.
\vs{6}Ἠκούσατε πάντα, καὶ ὑμεῖς οὐκ ἔγνωτε· ἀλλὰ ἀκουστά σοι ἐποίησα τὰ καινὰ ἀπὸ τοῦ νῦν, ἃ μέλλει γίνεσθαι· καὶ οὐκ εἶπας,
\vs{7}νῦν γίνεται, καὶ οὐ πάλαι, καὶ οὐ προτέραις ἡμέραις ἤκουσας αὐτά· μὴ εἴπῃς, ναὶ, γινώσκω αὐτά.
\vs{8}Οὔτε ἔγνως, οὔτε ἠπίστω, οὔτε ἀπʼ ἀρχῆς ἤνοιξά σου τὰ ὦτα· ἔγνων γὰρ, ὅτι ἀθετῶν ἀθετήσεις, καὶ ἄνομος ἔτι ἐκ κοιλίας κληθήσῃ.

\vs{9}Ἕνεκεν τοῦ ἐμοῦ ὀνόματος δείξω σοι τὸν θυμόν μου, καὶ τὰ ἔνδοξά μου ἐπάξω ἐπὶ σὲ, ἵνα μὴ ἐξολοθρεύσω σε.
\vs{10}Ἰδοὺ πέπρακά σε, οὐχ ἕνεκεν ἀργυρίου· ἐξειλάμην δέ σε ἐκ καμίνου πτωχείας·
\vs{11}Ἕνεκεν ἐμοῦ ποιήσω σοι, ὅτι τὸ ἐμὸν ὄνομα βεβηλοῦται, καὶ τὴν δόξαν μου ἑτέρῳ οὐ δώσω.

\vs{12}Ἄκουέ μου Ἰακὼβ, καὶ Ἰσραὴλ ὃν ἐγὼ καλῶ· ἐγώ εἰμι πρῶτος, καὶ ἐγώ εἰμι εἰς τὸν αἰῶνα.
\vs{13}Καὶ ἡ χείρ μου ἐθεμελίωσε τὴν γῆν, καὶ ἡ δεξιά μου ἐστερέωσε τὸν οὐρανόν· καλέσω αὐτοὺς, καὶ στήσονται ἅμα,
\vs{14}καὶ συναχθήσονται πάντες, καὶ ἀκούσονται· τίς αὐτοῖς ἀνήγγειλε ταῦτα; ἀγαπῶν σε ἐποίησα τὸ θέλημά σου ἐπὶ Βαβυλῶνα, τοῦ ἆραι σπέρμα Χαλδαίων.
\vs{15}Ἐγὼ ἐλάλησα, ἐγὼ ἐκάλεσα, ἤγαγον αὐτὸν, καὶ εὐώδωσα τὴν ὁδὸν αὐτοῦ.

\vs{16}Προσαγάγετε πρὸς μὲ, καὶ ἀκούσατε ταῦτα· οὐκ ἀπʼ ἀρχῆς ἐν κρυφῇ λελάληκα· ἡνίκα ἐγένετο, ἐκεῖ ἤμην, καὶ νῦν Κύριος Κύριος ἀπέστειλέ με, καὶ τὸ πνεῦμα αὐτοῦ.
\vs{17}Οὕτω λέγει Κύριος; ὁ ῥυσάμενος σε, ἅγιος Ἰσραὴλ, ἐγώ εἰμι ὁ Θεός σου, δέδειχά σοι τοῦ εὑρεῖν σε τὴν ὁδὸν ἐν ᾗ πορεύσῃ ἐν αὐτῇ.
\vs{18}Καὶ εἰ ἤκουσας τῶν ἐντολῶν μου, ἐγένετο ἂν ὡσεὶ ποταμὸς ἡ εἰρήνη σου, καὶ ἡ δικαιοσύνη σου ὡς κῦμα θαλάσσης·
\vs{19}Καὶ ἐγένετο ἂν ὡς ἡ ἄμμος τὸ σπέρμα σου, καὶ τὰ ἔκγονα τῆς κοιλίας σου ὡς ὁ χοῦς τῆς γῆς· οὐδὲ νῦν οὐ μὴ ἐξολοθρευθῇς, οὐδὲ ἀπολεῖται τὸ ὄνομά σου ἐνώπιον ἐμοῦ.

\vs{20}Ἔξελθε ἐκ Βαβυλῶνος φεύγων ἀπὸ τῶν Χαλδαίων, φωνὴν εὐφροσύνης ἀναγγειλατε, καὶ ἀκουστὸν γενέσθω τοῦτο, ἀναγγείλατε ἕως ἐσχάτου τῆς γῆς· λέγετε, Ἐῤῥύσατο Κύριος τὸν δοῦλον αὐτοῦ Ἰακώβ·
\vs{21}Καὶ ἐὰν διψήσωσι, διʼ ἐρήμου ἄξει αὐτοὺς, ὕδωρ ἐκ πέτρας ἐξάξει αὐτοῖς· σχισθήσεται πέτρα, καὶ ῥυήσεται ὕδωρ, καὶ πίεται ὁ λαός μου.
\vs{22}Οὐκ ἔστι χαίρειν, λέγει Κύριος, τοῖς ἀσεβέσιν.

\ch{49}
Ἀκούσατε μου νῆσοι, καὶ προσέχετε ἔθνη, διὰ χρόνου πολλοῦ στήσεται, λέγει Κύριος, ἐκ κοιλίας μητρός μου ἐκάλεσε τὸ ὄνομά μου.
\vs{2}Καὶ ἔθηκε τὸ στόμα μου ὡς μάχαιραν ὀξεῖαν, καὶ ὑπὸ τὴν σκέπην τῆς χειρὸς αὐτοῦ ἔκρυψέ με· ἔθηκέ με ὡς βέλος ἐκλεκτὸν, καὶ ἐν τῇ φαρέτρᾳ αὐτοῦ ἔκρυψέ με,
\vs{3}καὶ εἶπέ μοι, δοῦλός μου εἶ σὺ Ἰσραὴλ, καὶ ἐν σοὶ ἐνδοξασθήσομαι.
\vs{4}Καὶ ἐγὼ εἶπα, κενῶς ἐκοπίασα, εἰς μάταιον καὶ εἰς οὐδὲν ἔδωκα τὴν ἰσχύν μου· διατοῦτο ἡ κρίσις μου παρὰ· Κυρίῳ, καὶ ὁ πόνος μου ἐναντίον τοῦ Θεοῦ μου.
\vs{5}Καὶ νῦν οὕτω λέγει Κύριος, ὁ πλάσας με ἐκ κοιλίας δοῦλον ἑαυτῷ, τοῦ συναγαγεῖν τὸν Ἰακὼβ πρὸς αὐτὸν καὶ Ἰσραήλ· συναχθήσομαι καὶ δοξασθήσομαι ἐναντίον Κυρίου, καὶ ὁ Θεὸς μου ἔσται μοι ἰσχύς.
\vs{6}Καὶ εἶπέ μοι, μέγα σοι ἐστὶ τοῦ κληθῆναί σε παῖδά μου, τοῦ στῆσαι τὰς φυλὰς Ἰακὼβ, καὶ τὴν διασπορὰν τοῦ Ἰσραὴλ ἐπιστρέψαι· ἰδοὺ δέδωκά σε εἰς διαθήκην γένους, εἰς φῶς ἐθνῶν, τοῦ εἶναί σε εἰς σωτηρίαν ἕως ἐσχάτου τῆς γῆς.

\vs{7}Οὕτως λέγει Κύριος, ὁ ῥυσάμενός σε, ὁ Θεὸς Ἰσραὴλ, ἁγιάσατε τὸν φαυλίζοντα τὴν ψυχὴν αὐτοῦ, τὸν βδελυσσόμενον ὑπὸ τῶν ἐθνῶν τῶν δούλων τῶν ἀρχόντων· βασιλεῖς ὄψονται αὐτὸν, καὶ ἀναστήσονται ἄρχοντες, καὶ προσκυνήσουσιν αὐτῷ, ἕνεκεν Κυρίου· ὅτι πιστός ἐστιν ὁ ἅγιος Ἰσραὴλ, καὶ ἐξελεξάμην σε.

\vs{8}Οὕτως λέγει Κύριος, καιρῷ δεκτῷ ἐπήκουσά σου, καὶ ἐν ἡμέρᾳ σωτηρίας ἐβοήθησά σοι, καὶ ἔπλασά σε, καὶ ἔδωκά σε εἰς διαθήκην ἐθνῶν τοῦ καταστῆσαι τὴν γῆν, καὶ κληρονομῆσαι κληρονομίας ἐρήμους,
\vs{9}λέγοντα τοῖς ἐν δεσμοῖς, ἐξέλθατε, καὶ τοῖς ἐν τῷ σκότει, ἀνακαλυφθῆναι· ἐν πάσαις ταῖς ὁδοῖς βοσκηθήσονται, καὶ ἐν πάσαις ταῖς τρίβοις ἡ νομὴ αὐτῶν·
\vs{10}Οὐ πεινάσουσιν, οὐδὲ διψήσουσιν, οὐδὲ πατάξει αὐτοὺς ὁ καύσων, οὐδὲ ὁ ἥλιος, ἀλλʼ ὁ ἐλεῶν αὐτοὺς παρακαλέσει, καὶ διὰ πηγῶν ὑδάτων ἄξει αὐτούς·
\vs{11}Καὶ θήσω πᾶν ὄρος εἰς ὁδὸν, καὶ πᾶσαν τρίβον εἰς βόσκημα αὐτοῖς.
\vs{12}Ἰδοὺ οὗτοι πόῤῥωθεν ἥξουσιν, οὗτοι ἀπὸ Βοῤῥᾶ καὶ θαλάσσης, ἄλλοι δὲ ἐκ γῆς Περσῶν.

\vs{13}Εὐφραίνεσθε οὐρανοί, καὶ ἀγαλλιάσθω ἡ γῆ, ῥηξάτωσαν τὰ ὄρη εὐφροσύνην, ὅτι ἠλέησεν ὁ Θεὸς τὸν λαὸν αὐτοῦ, καὶ τοὺς ταπεινοὺς τοῦ λαοῦ αὐτοῦ παρεκάλεσεν.

\vs{14}Εἶπεν δὲ Σιὼν, ἐγκατέλιπέ με Κύριος, καὶ ὅτι Κύριος ἐπελάθετό μου.
\vs{15}Μὴ ἐπιλήσεται γυνὴ τοῦ παιδίου αὐτῆς, ἢ τοῦ μὴ ἐλεῆσαι τὰ ἔκγονα τῆς κοιλίας αὐτῆς; εἰ δὲ καὶ ταῦτα ἐπιλάθοιτο γυνὴ, ἀλλʼ ἐγὼ οὐκ ἐπιλήσομαί σου, εἶπεν Κύριος.

\vs{16}Ἰδοὺ ἐπὶ τῶν χειρῶν μου ἐζωγράφηκά σου τὰ τείχη, καὶ ἐνώπιόν μου εἶ διαπαντὸς,
\vs{17}καὶ ταχὺ οἰκοδομηθήσῃ ὑφʼ ὧν καθῃρέθης, καὶ οἱ ἐρημώσαντές σε ἐξελεύσονται ἐκ σοῦ.

\vs{18}Ἆρον κύκλῳ τοὺς ὀφθαλμούς σου, καὶ ἴδε πάντας, ἰδοὺ συνήχθησαν καὶ ἤλθοσαν πρὸς σέ. Ζῶ ἐγὼ, λέγει Κύριος, ὅτι πάντας αὐτοὺς ὡς κόσμον ἐνδύσῃ, καὶ περιθήσεις αὐτοὺς ὡς κόσμον νύμφη.
\vs{19}Ὅτι τὰ ἔρημά σου καὶ τὰ κατεφθαρμένα καὶ πεπτωκότα, ὅτι νῦν στενοχωρήσει ἀπὸ τῶν κατοικούντων, καὶ μακρυνθήσονται ἀπὸ σοῦ οἱ καταπίνοντές σε.
\vs{20}Ἐροῦσι γὰρ εἰς τὰ ὦτά σου οἱ υἱοί σου, οὓς ἀπολώλεκας, στενός μοι ὁ τόπος, ποίησόν μοι τόπον ἵνα κατοικήσω.
\vs{21}Καὶ ἐρεῖς ἐν τῇ καρδίᾳ σου, τίς ἐγέννησέ μοι τούτους; ἐγὼ δὲ ἄτεκνος καὶ χήρα, τούτους δὲ τίς ἐξέθρεψέ μοι; ἐγὼ δὲ κατελείφθην μόνη, οὗτοι δέ μοι ποῦ ἦσαν;

\vs{22}Οὕτως λέγει Κύριος Κύριος, ἰδοὺ αἴρω εἰς τὰ ἔθνη τὴν χεῖρά μου, καὶ εἰς τὰς νήσους ἀρῶ σύσσημόν μου, καὶ ἄξουσι τοὺς υἱούς σου ἐν κόλπῳ, τὰς δὲ θυγατέρας σου ἐπʼ ὤμων ἀροῦσι.
\vs{23}Καὶ ἔσονται βασιλεῖς τιθηνοί σου, αἱ δὲ ἄρχουσαι αὐτῶν τροφοί σου· ἐπὶ πρόσωπον τῆς γῆς προσκυνήσουσί σε, καὶ τὸν χοῦν τῶν ποδῶν σου λείξουσι, καὶ γνώσῃ, ὅτι ἐγὼ Κύριος, καὶ οὐκ αἰσχυνθήσονται οἱ ὑπομένοντές με.

\vs{24}Μὴ λήμψεταί τις παρὰ γίγαντος σκῦλα; καὶ ἐὰν αἰχμαλωτεύσῃ τις ἀδίκως, σωθήσεται;
\vs{25}Ὅτι οὕτω λέγει Κύριος, ἐάν τις αἰχμαλωτεύσῃ γίγαντα, λήψεται σκῦλα· λαμβάνων δὲ παρὰ ἰσχύοντος σωθήσεται· ἐγὼ δὲ τὴν κρίσιν σου κρινῶ, καὶ ἐγὼ τοὺς υἱούς σου ῥύσομαι·
\vs{26}Καὶ φάγονται οἱ θλίψαντές σε τὰς σάρκας αὐτῶν, καὶ πίονται ὡς οἶνον νέον τὸ αἷμα αὐτῶν, καὶ μεθυσθήσονται· καὶ αἰσθανθήσεται πᾶσα σὰρξ, ὅτι ἐγὼ Κύριος ὁ ῥυσάμενός σε, καὶ ἀντιλαμβανόμενος ἰσχύος Ἰακώβ.

\ch{50}
Οὕτως λέγει Κύριος, ποῖον τὸ βιβλίον τοῦ ἀποστασίου τῆς μητρὸς ὑμῶν, ᾧ ἐξαπέστειλα αὐτήν; ἢ τίνι ὑπόχρεῳ πέπρακα ὑμᾶς; ἰδοὺ ταῖς ἁμαρτίαις ὑμῶν ἐπράθητε, καὶ ταῖς ἀνομίαις ὑμῶν ἐξαπέστειλα τὴν μητέρα ὑμῶν.
\vs{2}Τί ὅτι ἦλθον, καὶ οὐκ ἦν ἄνθρωπος; ἐκάλεσα, καὶ οὐκ ἦν ὁ ὑπακούων; μὴ οὐκ ἰσχύει ἡ χείρ μου τοῦ ῥύσασθαι; ἢ οὐκ ἰσχύω τοῦ ἐξελέσθαι; ἰδοὺ τῷ ἐλεγμῷ μου ἐξερημώσω τὴν θάλασσαν, καὶ θήσω ποταμοὺς ἐρήμους, καὶ ξηρανθήσονται οἱ ἰχθύες αὐτῶν ἀπὸ τοῦ μὴ εἶναι ὕδωρ, καὶ ἀποθανοῦνται ἐν δίψει.
\vs{3}Ἐνδύσω τὸν οὐρανὸν σκότος, καὶ ὡς σάκκον θήσω τὸ περιβόλαιον αὐτοῦ.

\vs{4}Κύριος Κύριος δίδωσί μοι γλῶσσαν παιδείας, τοῦ γνῶναι ἡνίκα δεῖ εἰπεῖν λόγον· ἔθηκέ μοι πρωῒ, προσέθηκέ μοι ὠτίον ἀκούειν,
\vs{5}καὶ ἡ παιδεία Κυρίου Κυρίου ἀνοίγει μου τὰ ὦτα· ἐγὼ δὲ οὐκ ἀπειθῶ, οὐδὲ ἀντιλέγω.
\vs{6}Τὸν νῶτόν μου ἔδωκα εἰς μάστιγας, τὰς δὲ σιαγόνας μου εἰς ῥαπίσματα, τὸ δὲ πρόσωπόν μου οὐκ ἀπέστρεψα ἀπὸ αἰσχύνης ἐμπτυσμάτων,
\vs{7}καὶ Κύριος Κύριος βοηθός μοι ἐγενήθη· διατοῦτο οὐκ ἐνετράπην, ἀλλὰ ἔθηκα τὸ πρόσωπόν μου ὡς στερεὰν πέτραν, καὶ ἔγνων ὅτι οὐ μὴ αἰσχυνθῶ,
\vs{8}ὅτι ἐγγίζει ὁ δικαιώσας με· τίς ὁ κρινόμενός μοι; ἀντιστήτω μοι ἅμα· καὶ τις ὁ κρινόμενός μοι; ἐγγισάτω μοι.
\vs{9}Ἰδοὺ Κύριος Κύριος βοηθήσει μοι· τίς κακώσει με; ἰδοὺ πάντες ὑμεῖς ὡς ἱμάτιον παλαιωθήσεσθε, καὶ σὴς καταφάγεται ὑμᾶς.

\vs{10}Τίς ἐν ὑμῖν ὁ φοβούμενος τὸν Κύριον; ὑπακουσάτω τῆς φωνῆς τοῦ παιδὸς αὐτοῦ· οἱ πορευόμενοι ἐν σκότει, καὶ οὐκ ἔστιν αὐτοῖς φῶς, πεποίθατε ἐπὶ τῷ ὀνόματι Κυρίου, καὶ ἀντιστηρίσασθε ἐπὶ τῷ Θεῷ.
\vs{11}Ἰδοὺ πάντες ὑμεῖς πῦρ καίετε, καὶ κατισχύετε φλόγα· πορεύεσθε τῷ φωτὶ τοῦ πυρὸς ὑμῶν, καὶ τῇ φλογὶ ᾗ ἐξεκαύσατε· διʼ ἐμὲ ἐγένετο ταῦτα ὑμῖν, ἐν λύπῃ κοιμηθήσεσθε.

\ch{51}
Ἀκούσατε μου οἱ διώκοντες τὸ δίκαιον, καὶ ζητοῦντες τὸν Κύριον, ἐμβλέψατε εἰς τὴν στερεὰν πέτραν, ἣν ἐλατομήσατε, καὶ εἰς τὸν βόθυνον τοῦ λάκκου, ὃν ὠρύξατε.
\vs{2}Ἐμβλέψατε εἰς Ἁβραὰμ τὸν πατέρα ὑμῶν, καὶ εἰς Σάῤῥαν τὴν ὠδίνουσαν ὑμᾶς· ὅτι εἷς ἦν, καὶ ἐκάλεσα αὐτὸν, καὶ εὐλόγησα αὐτὸν, καὶ ἠγάπησα αὐτὸν, καὶ ἐπλήθυνα αὐτόν.
\vs{3}Καὶ σὲ νῦν παρακαλέσω, Σιὼν, καὶ παρεκάλεσα πάντα τὰ ἔρημα αὐτῆς, καὶ θήσω τὰ ἔρημα αὐτῆς ὡς παράδεισον, καὶ τὰ πρὸς δυσμὰς αὐτῆς ὡς παράδεισον Κυρίου· εὐφροσύνην καὶ ἀγαλλίαμα εὑρήσουσιν ἐν αὐτῇ, ἐξομολόγησιν καὶ φωνὴν αἰνέσεως.

\vs{4}Ἀκούσατέ μου, ἀκούσατέ μου λαός μου, καὶ οἱ βασιλεῖς πρὸς μὲ ἐνωτίσασθε, ὅτι νόμος παρʼ ἐμοῦ ἐξελεύσεται, καὶ ἡ κρίσις μου εἰς φῶς ἐθνῶν.
\vs{5}Ἐγγίζει ταχὺ ἡ δικαιοσύνη μου, καὶ ἐξελεύσεται ὡς φῶς τὸ σωτήριόν μου, καὶ εἰς τὸν βραχίονά μου ἔθνη ἐλπιοῦσιν· ἐμὲ νῆσοι ὑπομενοῦσι, καὶ εἰς τὸν βραχίονά μου ἐλπιοῦσιν.
\vs{6}Ἄρατε εἰς τὸν οὐρανὸν τοὺς ὀφθαλμοὺς ὑμῶν, καὶ ἐμβλέψατε εἰς τὴν γῆν κάτω· ὅτι ὁ οὐρανὸς ὡς καπνὸς ἐστερεώθη, ἡ δὲ γῆ ὡς ἱμάτιον παλαιωθήσεται, οἱ δὲ κατοικοῦντες ὥσπερ ταῦτα ἀποθανοῦνται, τὸ δὲ σωτήριόν μου εἰς τὸν αἰῶνα ἔσται, ἡ δὲ δικαιοσύνη μου οὐ μὴ ἐκλίπῃ.

\vs{7}Ἀκούσατέ μου οἱ εἰδότες κρίσιν, λαὸς οὗ ὁ νόμος μου ἐν τῇ καρδίᾳ ὑμῶν· μὴ φοβεῖσθε ὀνειδισμὸν ἀνθρώπων, καὶ τῷ φαυλισμῷ αὐτῶν μὴ ἡττᾶσθε·
\vs{8}Ὡς γὰρ ἱμάτιον βρωθήσεται ὑπὸ χρόνου, καὶ ὡς ἔρια βρωθήσεται ὑπὸ σητός· ἡ δὲ δικαιοσύνη μου εἰς τὸν αἰῶνα ἔσται, τὸ δὲ σωτήριόν μου εἰς γενεὰς γενεῶν.

\vs{9}Ἐξεγείρου ἐξεγείρου Ἱερουσαλὴμ, καὶ ἔνδυσαι τὴν ἰσχὺν τοῦ βραχίονός σου, ἐξεγείρου ὡς ἐν ἀρχῇ ἡμέρας, ὡς γενεὰ αἰῶνος.
\vs{10}Οὐ σὺ εἶ ἡ ἐρημοῦσα θάλασσαν, ὕδωρ ἀβύσσου πλῆθος; ἡ θεῖσα τὰ βάθη τῆς θαλάσσης ὁδὸν διαβάσεως ῥυομένοις καὶ λελυτρωμένοις;
\vs{11}ὑπὸ γὰρ Κυρίου ἀποστραφήσονται, καὶ ἥξουσιν εἰς Σιὼν μετʼ εὐφροσύνης καὶ ἀγαλλιάματος αἰωνίου· ἐπὶ κεφαλῆς γὰρ αὐτῶν αἴνεσις καὶ εὐφροσύνη καταλήμψεται αὐτούς· ἀπέδρα ὀδύνη καὶ λύπη καὶ στεναγμός.

\vs{12}Ἐγώ εἰμι, ἐγώ εἰμι ὁ παρακαλῶν σε· γνῶθι τίς οὖσα ἐφοβήθης ἀπὸ ἀνθρώπου θνητοῦ, καὶ ἀπὸ υἱοῦ ἀνθρώπου, οἳ ὡσεὶ χόρτος ἐξηράνθησαν.
\vs{13}Καὶ ἐπελάθου Θεὸν τὸν ποιήσαντά σε, τὸν ποιήσαντα τὸν οὐρανὸν καὶ θεμελιώσαντα τὴν γῆν· καὶ ἐφόβου ἀεὶ πάσας τὰς ἡμέρας τὸ πρόσωπον τοῦ θυμοῦ τοῦ θλίβοντός σε· ὃν τρόπον γὰρ ἐβουλεύσατο τοῦ ἆραί σε, καὶ νῦν ποῦ ὁ θυμὸς τοῦ θλίβοντός σε;
\vs{14}Ἐν γὰρ τῷ σώζεσθαί σε οὐ στήσεται, οὐδὲ χρονιεῖ,
\vs{15}ὅτι ἐγὼ ὁ Θεός σου, ὁ ταράσσων τὴν θάλασσαν, καὶ ἠχῶν τὰ κύματα αὐτῆς· Κύριος σαβαὼθ ὄνομά μοι.
\vs{16}Θήσω τοὺς λόγους μου εἰς τὸ στόμα σου, καὶ ὑπὸ τὴν σκιὰν τῆς χειρός μου σκεπάσω σε, ἐν ᾗ ἔστησα τὸν οὐρανὸν, καὶ ἐθεμελίωσα τὴν γῆν· καὶ ἐρεῖ Σιὼν, λαός μου εἶ σύ.

\vs{17}Ἐξεγείρου ἐξεγείρου, ἀνάστηθι Ἱερουσαλὴμ, ἡ πιοῦσα ἐκ χειρὸς Κυρίου τὸ ποτήριον τοῦ θυμοῦ αὐτοῦ· τὸ ποτήριον γὰρ τῆς πτώσεως, τὸ κόνδυ τοῦ θυμοῦ ἐξέπιες καὶ ἐξεκένωσας,
\vs{18}καὶ οὐκ ἦν ὁ παρακαλῶν σε ἀπὸ πάντων τῶν τέκνων σου ὧν ἔτεκες, καὶ οὐκ ἦν ὁ ἀντιλαμβανόμενος τῆς χειρός σου, οὐδὲ ἀπὸ πάντων τῶν υἱῶν σου ὧν ὕψωσας.
\vs{19}Διὸ ταῦτα ἀντικείμενά σοι, τίς συλλυπηθήσεταί σοι; πτῶμα καὶ σύντριμμα, λιμὸς καὶ μάχαιρα, τίς παρακαλέσει σε;
\vs{20}Οἱ υἱοί σου οἱ ἀπορούμενοι, οἱ καθεύδοντες ἐπʼ ἄκρου πάσης ἐξόδου ὡς σευτλίον ἡμίεφθον, οἱ πλήρεις θυμοῦ Κυρίου, ἐκλελυμένοι διὰ Κυρίου τοῦ Θεοῦ.

\vs{21}Διατοῦτο ἄκουε τεταπεινωμένη, καὶ μεθύουσα οὐκ ἀπὸ οἴνου.
\vs{22}Οὕτως λέγει Κύριος ὁ Θεὸς ὁ κρίνων τὸν λαὸν αὐτοῦ, ἰδοὺ εἴληφα ἐκ τῆς χειρός σου τὸ ποτήριον τῆς πτώσεως, τὸ κόνδυ τοῦ θυμοῦ μου, καὶ οὐ προσθήσῃ ἔτι πιεῖν αὐτό·
\vs{23}καὶ δώσω αὐτὸ εἰς τὰς χεῖρας τῶν ἀδικησάντων σε καὶ τῶν ταπεινωσάντων σε, οἳ εἶπαν τῇ ψυχῇ σου, κύψον, ἵνα παρέλθωμεν· καὶ ἔθηκας ἴσα τῇ γῇ τὰ μέσα σου ἔξω τοῖς παραπορευομένοις.

\ch{52}
Ἐξεγείρου ἐξεγείρου Σιὼν, ἔνδυσαι τὴν ἰσχύν σου Σιὼν, καὶ σὺ ἔνδυσαι τὴν δόξαν σου Ἱερουσαλὴμ πόλις ἡ ἁγία· οὐκέτι προστεθήσεται διελθεῖν διὰ σοῦ ἀπερίτμητος καὶ ἀκάθαρτος.
\vs{2}Ἐκτίναξαι τὸν χοῦν καὶ ἀνάστηθι, κάθισον Ἱερουσαλὴμ, ἔκδυσαι τὸν δεσμὸν τοῦ τραχήλου σου ἡ αἰχμάλωτος θυγάτηρ Σιών.

\vs{3}Ὅτι τάδε λέγει Κύριος, δωρεὰν ἐπράθητε, καὶ οὐ μετὰ ἀργυρίου λυτρωθήσεσθε.
\vs{4}Οὕτως λέγει Κύριος, εἰς Αἴγυπτον κατέβη ὁ λαός μου τὸ πρότερον παροικῆσαι ἐκεῖ, καὶ εἰς Ἀσσυρίους βίᾳ ἤχθησαν.
\vs{5}Καὶ νῦν τί ἐστὲ ὧδε; τάδε λέγει Κύριος, ὅτι ἐλήφθη ὁ λαός μου δωρεὰν, θαυμάζετε καὶ ὀλολύζετε· τάδε λέγει Κύριος, διʼ ὑμᾶς διαπαντὸς τὸ ὄνομά μου βλασφημεῖται ἐν τοῖς ἔθνεσι.
\vs{6}Διατοῦτο γνώσεται ὁ λαός μου τὸ ὄνομά μου ἐν τῇ ἡμέρᾳ ἐκείνῃ, ὅτι ἐγώ εἰμι αὐτὸς ὁ λαλῶν, πάρειμι ὡς ὥρα ἐπὶ τῶν ὀρέων,
\vs{7}ὡς πόδες εὐαγγελιζομένου ἀκοὴν εἰρήνης, ὡς εὐαγγελιζόμενος ἀγαθὰ, ὅτι ἀκουστὴν ποιήσω τὴν σωτηρίαν σου, λέγων, Σιὼν βασιλεύσει σου ὁ Θεός.
\vs{8}Ὅτι φωνὴ τῶν φυλασσόντων σε ὑψώθη, καὶ τῇ φωνῇ ἅμα εὐφρανθήσονται· ὅτι ὀφθαλμοὶ πρὸς ὀφθαλμοὺς ὄψονται, ἡνίκα ἂν ἐλεήσῃ Κύριος τὴν Σιών.
\vs{9}Ῥηξάτω εὐφροσύνην ἅμα τὰ ἔρημα Ἱερουσαλὴμ, ὅτι ἠλέησε Κύριος αὐτὴν, καὶ ἐῤῥύσατο Ἱερουσαλήμ.
\vs{10}Καὶ ἀποκαλύψει Κύριος τὸν βραχίονα τὸν ἅγιον αὐτοῦ ἐνώπιον πάντων τῶν ἐθνῶν, καὶ ὄψονται πάντα ἄκρα τῆς γῆς τὴν σωτηρίαν τὴν παρὰ τοῦ Θεοῦ ἡμῶν.

\vs{11}Ἀπόστητε, ἀπόστητε, ἐξέλθατε ἐκεῖθεν, καὶ ἀκαθάρτου μὴ ἅψησθε, ἐξέλθατε ἐκ μέσου αὐτῆς, ἀφορίσθητε οἱ φέροντες τὰ σκεύη Κυρίου·
\vs{12}Ὅτι οὐ μετὰ ταραχῆς ἐξελεύσεσθε, οὐδὲ φυγῇ πορεύσεσθε· προπορεύσεται γὰρ πρότερος ὑμῶν Κύριος, καὶ ὁ ἐπισυνάγων ὑμᾶς Θεὸς Ἰσραήλ.

\vs{13}Ἰδοὺ, συνήσει ὁ παῖς μου, καὶ ὑψωθήσεται, καὶ δοξασθήσεται σφόδρα.
\vs{14}Ὃν τρόπον ἐκστήσονται ἐπὶ σὲ πολλοὶ, οὕτως ἀδοξήσει ἀπὸ τῶς ἀνθρώπων τὸ εἶδός σου, καὶ ἡ δόξα σου ἀπὸ υἱῶν ἀνθρώπων.
\vs{15}Οὕτω θαυμάσονται ἔθνη πολλὰ ἐπʼ αὐτῷ, καὶ συνέξουσι βασιλεῖς τὸ στόμα αὐτῶν· ὅτι οἷς οὐκ ἀνηγγέλη περὶ αὐτοῦ, ὄψονται, καὶ οἳ οὐκ ἀκηκόασι, συνήσουσι.

\ch{53}
Κύριε τίς ἐπίστευσε τῇ ἀκοῇ ἡμῶν; καὶ ὁ βραχίων Κυρίου τίνι ἀπεκαλύφθη;
\vs{2}Ἀνηγγείλαμεν ὡς παιδίον ἐναντίον αὐτοῦ, ὡς ῥίζα ἐν γῇ διψώσῃ· οὐκ ἔστιν εἶδος αὐτῷ, οὐδὲ δόξα· καὶ εἴδομεν αὐτὸν, καὶ οὐκ εἶχεν εἶδος οὐδὲ κάλλος·
\vs{3}Ἀλλὰ τὸ εἶδος αὐτοῦ ἄτιμον, καὶ ἐκλεῖπον παρὰ τοὺς υἱοὺς τῶν ἀνθρώπων· ἄνθρωπος ἐν πληγῇ ὢν, καὶ εἰδὼς φέρειν μαλακίαν, ὅτι ἀπέστραπται τὸ πρόσωπον αὐτοῦ, ἠτιμάσθη, καὶ οὐκ ἐλογίσθη.
\vs{4}Οὗτος τὰς ἁμαρτίας ἡμῶν φέρει, καὶ περὶ ἡμῶν ὀδυνᾶται, καὶ ἡμεῖς ἐλογισάμεθα αὐτὸν εἶναι ἐν πόνῳ, καὶ ἐν πληγῇ, καὶ ἐν κακώσει.
\vs{5}Αὐτὸς δὲ ἐτραυματίσθη διὰ τὰς ἁμαρτίας ἡμῶν, καὶ μεμαλάκισται διὰ τὰς ἀνομίας ἡμῶν· παιδεία εἰρήνης ἡμῶν ἐπʼ αὐτὸν, τῷ μώλωπι αὐτοῦ ἡμεῖς ἰάθημεν.
\vs{6}Πάντες ὡς πρόβατα ἐπλανήθημεν· ἄνθρωπος τῇ ὁδῷ αὐτοῦ ἐπλανήθη· καὶ Κύριος παρέδωκεν αὐτὸν ταῖς ἁμαρτίαις ἡμῶν.

\vs{7}Καὶ αὐτὸς διὰ τὸ κεκακῶσθαι οὐκ ἀνοίγει τὸ στόμα αὐτοῦ· ὡς πρόβατον ἐπὶ σφαγὴν ἤχθη, καὶ ὡς ἀμνὸς ἐναντίον τοῦ κείροντος ἄφωνος, οὕτως οὐκ ἀνοίγει τὸ στόμα.
\vs{8}Ἐν τῇ ταπεινώσει ἡ κρίσις αὐτοῦ ᾔρθη, τὴν γενεὰν αὐτοῦ τίς διηγήσεται; ὅτι αἴρεται ἀπὸ τῆς γῆς ἡ ζωὴ αὐτοῦ, ἀπὸ τῶν ἀνομιῶν τοῦ λαοῦ μου ἤχθη εἰς θάνατον.
\vs{9}Καὶ δώσω τοὺς πονηροὺς ἀντὶ τῆς ταφῆς αὐτοῦ, καὶ τοὺς πλουσίους ἀντὶ τοῦ θανάτου αὐτοῦ· ὅτι ἀνομίαν οὐκ ἐποίησεν, οὐδὲ δόλον ἐν τῷ στόματι αὐτοῦ.
\vs{10}Καὶ Κύριος βούλεται καθαρίσαι αὐτὸν τῆς πληγῆς· ἐὰν δῶτε περὶ ἁμαρτίας, ἡ ψυχὴ ὑμῶν ὄψεται σπέρμα μακρόβιον·
\vs{11}καὶ βούλεται Κύριος ἀφελεῖν ἀπὸ τοῦ πόνου τῆς ψυχῆς αὐτοῦ, δεῖξαι αὐτῷ φῶς, καὶ πλάσαι τῇ συνέσει, δικαιῶσαι δίκαιον εὖ δουλεύοντα πολλοῖς, καὶ τὰς ἁμαρτίας αὐτῶν αὐτὸς ἀνοίσει.
\vs{12}Διατοῦτο αὐτὸς κληρονομήσει πολλοὺς, καὶ τῶν ἰσχυρῶν μεριεῖ σκῦλα· ἀνθʼ ὧν παρεδόθη εἰς θάνατον ἡ ψυχὴ αὐτοῦ, καὶ ἐν τοῖς ἀνόμοις ἐλογίσθη, καὶ αὐτὸς ἁμαρτίας πολλῶν ἀνήνεγκε, καὶ διὰ τὰς ἀνομίας αὐτῶν παρεδόθη.

\ch{54}
Εὐφράνθητι στεῖρα ἡ οὐ τίκτουσα, ῥῆξον καὶ βόησον ἡ οὐκ ὠδίνουσα, ὅτι πολλὰ τὰ τέκνα τῆς ἐρήμου, μᾶλλον ἢ τῆς ἐχούσης τὸν ἄνδρα· εἶπε γὰρ Κύριος,
\vs{2}πλάτυνον τὸν τόπον τῆς σκηνῆς σου, καὶ τῶν αὐλαιῶν σου, πῆξον, μὴ φείσῃ, μάκρυνον τὰ σχοινίσματά σου, καὶ τοὺς πασσάλους σου κατίσχυσον,
\vs{3}ἔτι εἰς τὰ δεξιὰ καὶ τὰ ἀριστερὰ ἐκπέτασον· καὶ τὸ σπέρμα σου ἔθνη κληρονομήσει, καὶ πόλεις ἠρημωμένας κατοικιεῖς.
\vs{4}Μὴ φοβοῦ, ὅτι κατῃσχύνθης, μηδὲ ἐντραπῇς, ὅτι ὠνειδίσθης, ὅτι αἰσχύνην αἰώνιον ἐπιλήσῃ, καὶ ὄνειδος τῆς χηρείας σου οὐ μὴ μνησθήσῃ ἔτι.
\vs{5}Ὅτι Κύριος ὁ ποιῶν σε, Κύριος σαβαὼθ ὄνομα αὐτῷ· καὶ ὁ ῥυσάμενός σε, αὐτὸς Θεὸς Ἰσραὴλ, πάσῃ τῇ γῇ κληθήσεται.
\vs{6}Οὐχ ὡς γυναῖκα καταλελειμμένην καὶ ὀλιγόψυχον κέκληκέ σε ὁ Κύριος, οὐδʼ ὡς γυναῖκα ἐκ νεότητος μεμισημένην, εἶπεν ὁ Θεός σου.

\vs{7}Χρόνον μικρὸν κατέλιπόν σε, καὶ μετʼ ἐλέους μεγάλου ἐλεήσω σε.
\vs{8}Ἐν θυμῷ μικρῷ ἀπέστρεψα τὸ πρόσωπόν μου ἀπὸ σοῦ, καὶ ἐν ἐλέει αἰωνίῳ ἐλεήσω σε, εἶπεν ὁ ῥυσάμενός σε Κύριος·

\vs{9}Ἀπὸ τοῦ ὕδατος τοῦ ἐπὶ Νῶε τοῦτό μοι ἐστί· καθότι ὤμοσα αὐτῷ ἐν τῷ χρόνῳ ἐκείνῳ, τῇ γῇ μὴ θυμωθήσεσθαι ἐπὶ σοὶ ἔτι, μηδὲ ἐν ἀπειλῇ σου τὰ ὄρη μεταστήσεσθαι,
\vs{10}οὐδʼ οἱ βουνοί σου μετακινηθήσονται· οὕτως οὐδὲ τὸ παρʼ ἐμοῦ σοι ἔλεος ἐκλείψει, οὐδὲ ἡ διαθήκη τῆς εἰρήνης σου οὐ μὴ μεταστῇ· εἶπε γὰρ ἵλεώς σοι Κύριε.

\vs{11}Ταπεινὴ καὶ ἀκατάστατος οὐ παρεκλήθης· ἰδοὺ, ἐγὼ ἑτοιμάζω σοι ἄνθρακα τὸν λίθον σου, καὶ τὰ θεμέλιά σου σάπφειρον,
\vs{12}καὶ θήσω τὰς ἐπάλξεις σου ἴασπιν, καὶ τὰς πύλας σου λίθους κρυστάλλου, καὶ τὸν περίβολόν σου λίθους ἐκλεκτούς·
\vs{13}καὶ πάντας τοὺς υἱοὺς σου διδακτοὺς Θεοῦ, καὶ ἐν πολλῇ εἰρήνῃ τὰ τέκνα σου.
\vs{14}Καὶ ἐν δικαιοσύνῃ οἰκοδομηθήσῃ· ἀπέχου ἀπὸ ἀδίκου, καὶ οὐ φοβηθήσῃ, καὶ τρόμος οὐκ ἐγγιεῖ σοι.
\vs{15}Ἰδοὺ προσήλυτοι προσελεύσονταί σοι διʼ ἐμοῦ, καὶ παροικήσουσί σοι, καὶ ἐπὶ σὲ καταφεύξονται.

\vs{16}Ἰδοὺ ἐγὼ ἔκτισά σε, οὐχ ὡς χαλκεὺς φυσῶν ἄνθρακας, καὶ ἐκφέρων σκεῦος εἰς ἔργον· ἐγὼ δὲ ἔκτισά σε, οὐκ εἰς ἀπώλειαν φθεῖραι.
\vs{17}Πᾶν σκεῦος σκευαστὸν ἐπὶ σὲ, οὐκ εὐοδώσω· καὶ πᾶσα φωνὴ ἀναστήσεται ἐπὶ σὲ εἰς κρίσιν, πάντας αὐτοὺς ἡττήσεις, οἱ δὲ ἔνοχοί σου ἔσονται ἐν αὐτῇ. Ἔστι κληρονομία τοῖς θεραπεύουσι Κύριον· καὶ ὑμεῖς ἔσεσθέ μοι δίκαιοι, λέγει Κύριος.

\ch{55}
Οἱ διψῶντες πορεύεσθε ἐφʼ ὕδωρ, καὶ ὅσοι μὴ ἔχετε ἀργύριον, βαδίσαντες ἀγοράσατε, καὶ φάγετε ἄνευ ἀργυρίου καὶ τιμῆς οἶνον καὶ στέαρ.
\vs{2}Ἱνατί τιμᾶσθε ἀργυρίου, καὶ τὸν μόχθον ὑμῶν οὐκ εἰς πλησμονήν; ἀκούσατέ μου, καὶ φάγεσθε ἀγαθὰ, καὶ ἐντρυφήσει ἐν ἀγαθοῖς ἡ ψυχὴ ὑμῶν.

\vs{3}Προσέχετε τοῖς ὠσὶν ὑμῶν, καὶ ἐπακολουθήσατε ταῖς ὁδοῖς μου· εἰσακούσατέ μου, καὶ ζήσεται ἐν ἀγαθοῖς ἡ ψυχὴ ὑμῶν, καὶ διαθήσομαι ὑμῖν διαθήκην αἰώνιον, τὰ ὅσια Δαυὶδ τὰ πιστά.
\vs{4}Ἰδοὺ, μαρτύριον ἐν ἔθνεσιν ἔδωκα αὐτὸν, ἄρχοντα καὶ προστάσσοντα ἔθνεσιν.
\vs{5}Ἔθνη ἃ οὐκ οἴδασί σε, ἐπικαλέσονταί σε, καὶ λαοὶ οἳ οὐκ ἐπίστανταί σε, ἐπὶ σὲ καταφεύξονται, ἕνεκεν Κυρίου τοῦ Θεοῦ σου τοῦ ἁγίου Ἰσραὴλ, ὅτι ἐδόξασέν σε.

\vs{6}Ζητήσατε τὸν Κύριον, καὶ ἐν τῷ εὑρίσκειν αὐτὸν, ἐπικαλέσασθε· ἡνίκα δʼ ἂν ἐγγίζῃ ὑμῖν,
\vs{7}ἀπολιπέτω ὁ ἀσεβὴς τὰς ὁδοὺς αὐτοῦ, καὶ ἀνὴρ ἄνομος τὰς βουλὰς αὐτοῦ, καὶ ἐπιστραφήτω ἐπὶ Κύριον, καὶ ἐλεηθήσεται, ὅτι ἐπὶ πολὺ ἀφήσει τὰς ἁμαρτίας ὑμῶν.
\vs{8}Οὐ γάρ εἰσιν αἱ βουλαί μου ὥσπερ αἱ βουλαὶ ὑμῶν, οὐδʼ ὥσπερ αἱ ὁδοὶ ὑμῶν αἱ ὁδοί μου, λέγει Κύριος.
\vs{9}Ἀλλʼ ὡς ἀπέχει ὁ οὐρανὸς ἀπὸ τῆς γῆς, οὕτως ἀπέχει ἡ ὁδός μου ἀπὸ τῶν ὁδῶν ὑμῶν, καὶ τὰ διανοήματα ὑμῶν ἀπὸ τῆς διανοίας μου.
\vs{10}Ὡς γὰρ ἂν καταβῇ ὁ ὑετὸς ἢ χιὼν ἐκ τοῦ οὐρανοῦ, καὶ οὐ μὴ ἀποστραφῇ ἕως ἂν μεθύσῃ τὴν γῆν, καὶ ἐκτέκῃ, καὶ ἐκβλαστήσῃ, καὶ δῷ σπέρμα τῷ σπείροντι, καὶ ἄρτον εἰς βρῶσιν·
\vs{11}οὕτως ἔσται τὸ ῥῆμά μου, ὃ ἐὰν ἐξέλθῃ ἐκ τοῦ στόματός μου, οὐ μὴ ἀποστραφῇ, ἕως ἂν τελεσθῇ ὅσα ἂν ἠθέλησα, καὶ εὐοδώσω τὰς ὁδούς σου, καὶ τὰ ἐντάλματά μου.
\vs{12}Ἐν γὰρ εὐφροσύνῃ ἐξελεύσεσθε, καὶ ἐν χαρᾷ διδαχθήσεσθε· τὰ γὰρ ὄρη καὶ οἱ βουνοὶ ἐξαλοῦνται προσδεχόμενοι ὑμᾶς ἐν χαρᾷ, καὶ πάντα τὰ ξύλα τοῦ ἀγροῦ ἐπικροτήσει τοῖς κλάδοις.
\vs{13}Καὶ ἀντὶ τῆς στοιβῆς ἀναβήσεται κυπάρισσος, ἀντὶ δὲ τῆς κονύζης ἀναβήσεται μυρσίνη· καὶ ἔσται Κύριος εἰς ὄνομα, καὶ εἰς σημεῖον αἰώνιον, καὶ οὐκ ἐκλείψει.

\ch{56}
Τάδε λέγει Κύριος, φυλάσσεσθε κρίσιν, καὶ ποιήσατε δικαιοσύνην· ἤγγικε γὰρ τὸ σωτήριόν μου παραγίνεσθαι, καὶ τὸ ἔλεός μου ἀποκαλυφθῆναι.
\vs{2}Μακάριος ἀνὴρ ὁ ποιῶν ταῦτα, καὶ ἄνθρωπος ὁ ἀντεχόμενος αὐτῶν, καὶ φυλάσσων τὰ σάββατα μὴ βεβηλοῦν, καὶ διατηρῶν τὰς χεῖρας αὐτοῦ μὴ ποιεῖν ἄδικα.

\vs{3}Μὴ λεγέτω ὁ ἀλλογενὴς ὁ προσκείμενος πρὸς Κύριον, ἀφοριεῖ με ἄρα Κύριος ἀπὸ τοῦ λαοῦ αὐτοῦ· καὶ μὴ λεγέτω ὁ εὐνοῦχος, ὅτι ξύλον ἐγώ εἰμι ξηρόν.
\vs{4}Τάδε λέγει Κύριος τοῖς εὐνούχοις, ὅσοι ἄν φυλάξωνται τὰ σάββατά μου, καὶ ἐκλέξωνται ἃ ἐγὼ θέλω, καὶ ἀντέχωνται τῆς διαθήκης μου,
\vs{5}δώσω αὐτοῖς ἐν τῷ οἴκῳ μου καὶ ἐν τῷ τείχει μου τόπον ὀνομαστὸν, κρείττω υἱῶν καὶ θυγατέρων· ὄνομα αἰώνιον δώσω αὐτοῖς, καὶ οὐκ ἐκλείψει·
\vs{6}καὶ τοῖς ἀλλογενέσι τοῖς προσκειμένοις Κυρίῳ δουλεύειν αὐτῷ, καὶ ἀγαπᾷν τὸ ὄνομα Κυρίου, τοῦ εἶναι αὐτῷ εἰς δούλους καὶ δούλας· καὶ πάντας τοὺς φυλασσομένους τὰ σάββατά μου μὴ βεβηλοῦν, καὶ ἀντεχομένους τῆς διαθήκης μου,
\vs{7}εἰσάξω αὐτοὺς εἰς τὸ ὄρος τὸ ἅγιόν μου, καὶ εὐφρανῶ αὐτοὺς ἐν τῷ οἴκῳ τῆς προσευχῆς μου· τὰ ὁλοκαυτώματα αὐτῶν, καὶ αἱ θυσίαι αὐτῶν ἔσονται δεκταὶ ἐπὶ τὸ θυσιαστήριόν μου· ὁ γάρ οἶκός μου, οἶκος προσευχῆς κληθήσεται πᾶσι τοῖς ἔθνεσιν,
\vs{8}εἶπε Κύριος ὁ συνάγων τοὺς διεσπαρμένους Ἰσραὴλ, ὅτι συνάξω ἐπʼ αὐτὸν συναγωγήν.

\vs{9}Πάντα τὰ θηρία τὰ ἄγρια, δεῦτε, φάγετε, πάντα τὰ θηρία τοῦ δρυμοῦ.
\vs{10}Ἴδετε, ὅτι ἐκτετύφλωνται πάντες, οὐκ ἔγνωσαν, κύνες ἐνεοί οὐ δυνήσονται ὑλακτεῖν, ἐνυπνιαζόμενοι κοίτην, φιλοῦντες νυστάξαι.
\vs{11}Καὶ οἱ κύνες ἀναιδεῖς τῇ ψυχῇ, οὐκ εἰδότες πλησμονήν· καί εἰσι πονηροὶ οὐκ εἰδότες σύνεσιν, πάντες ταῖς ὁδοῖς αὐτων ἐξηκολούθησαν, ἕκαστος κατὰ τὸ ἑαυτοῦ.

\ch{57}
Ἴδετε ὡς ὁ δίκαιος ἀπώλετο, καὶ οὐδεὶς ἐκδέχεται τῇ καρδίᾳ, καὶ ἄνδρες δίκαιοι αἴρονται, καὶ οὐδεὶς κατανοεῖ· ἀπὸ γὰρ προσώπου ἀδικίας ᾖρται ὁ δίκαιος.
\vs{2}Ἔσται ἐν εἰρήνῃ ἡ ταφὴ αὐτοῦ, ᾖρται ἐκ τοῦ μέσου.

\vs{3}Ὑμεῖς δὲ προσαγάγετε ὧδε υἱοὶ ἄνομοι, σπέρμα μοιχῶν καὶ πόρνης.
\vs{4}Ἐν τίνι ἐνετρυφήσατε; καὶ ἐπὶ τίνα ἠνοίξατε τὸ στόμα ὑμῶν; καὶ ἐπὶ τίνα ἐχαλάσατε τὴν γλῶσσαν ὑμῶν; οὐχ ὑμεῖς ἐστέ τέκνα ἀπωλείας; σπέρμα ἄνομον;
\vs{5}Οἱ παρακαλοῦντες εἴδωλα ὑπὸ δένδρα δασέα, σφάζοντες τὰ τέκνα αὐτῶν ἐν ταῖς φάραγξιν ἀναμέσον τῶν πετρῶν;
\vs{6}Ἐκείνη σου ἡ μερίς, οὗτός σου ὁ κλῆρος, κᾀκείνοις ἐξέχεας σπονδὰς, καὶ τούτοις ἀνήνεγκας θυσίας· ἐπὶ τούτοις οὖν οὐκ ὀργισθήσομαι;

\vs{7}Ἐπʼ ὄρος ὑψηλὸν καὶ μετέωρον, ἐκεῖ σου ἡ κοίτη, καὶ ἐκεῖ ἀνεβίβασας θυσίας σου,
\vs{8}καὶ ὀπίσω τῶν σταθμῶν τῆς θύρας σου ἔθηκας μνημόσυνά σου· ᾤου, ὅτι ἐὰν ἀπʼ ἐμοῦ ἀποστῇς, πλεῖόν τι ἕξεις; ἠγάπησας τοὺς κοιμωμένους μετὰ σοῦ,
\vs{9}καὶ ἐπλήθυνας τὴν πορνείαν σου μετʼ αὐτῶν, καὶ πολλοὺς ἐποίησας τοὺς μακρὰν ἀπὸ σοῦ, καὶ ἀπέστειλας πρέσβεις ὑπὲρ τὰ ὅριά σου, καὶ ἐταπεινώθης ἕως ᾅδου.
\vs{10}Ταῖς πολυοδίαις σου ἐκοπίασας, καὶ οὐκ εἶπας παύσομαι ἐνισχύουσα· ὅτι ἔπραξας ταῦτα, διατοῦτο οὐ κατεδεήθης μου σύ.

\vs{11}Τίνα εὐλαβηθεῖσα ἐφοβήθης, καὶ ἐψεύσω με, καὶ οὐκ ἐμνήσθης, οὐδὲ ἔλαβές με εἰς τὴν διάνοιαν, οὐδὲ εἰς τὴν καρδίαν σου; καὶ ἐγώ σε ἰδὼν παρορῶ, καὶ ἐμὲ οὐκ ἐφοβήθης.

\vs{12}Καὶ ἐγὼ ἀπαγγελῶ τὴν δικαιοσύνην σου, καὶ τὰ κακά σου, ἃ οὐκ ὠφελήσει σε,
\vs{13}ὅταν ἀναβοήσῃς ἐξελέσθωσάν σε ἐν τή θλίψει σου· τούτους γὰρ πάντας ἄνεμος λήψεται, καὶ ἀποίσει καταιγίς· οἱ δὲ ἀντεχόμενοί μου κτήσονται γῆν, καὶ κληρονομήσουσι τὸ ὄρος τὸ ἅγιόν μου·
\vs{14}Καὶ ἐροῦσι, καθαρίσατε ἀπὸ προσώπου αὐτοῦ ὁδοὺς, καὶ ἄρατε σκῶλα ἀπὸ τῆς ὁδοῦ τοῦ λαοῦ μου.

\vs{15}Τάδε λέγει ὁ ὕψιστος, ἐν ὑψηλοῖς κατοικῶν τὸν αἰῶνα, ἅγιος ἐν ἁγίοις, ὄνομα αὐτῷ, ὕψιστος ἐν ἁγίοις ἀναπαυόμενος, καὶ ὀλιγοψύχοις διδοὺς μακροθυμίαν, καὶ διδοὺς ζωὴν τοῖς συντετριμμένοις τὴν καρδίαν.
\vs{16}Οὐκ εἰς τὸν αἰῶνα ἐκδικήσω ὑμᾶς, οὐδὲ διαπαντὸς ὀργισθήσομαι ὑμῖν· πνεῦμα γὰρ παρʼ ἐμοῦ ἐξελεύσεται, καὶ πνοὴν πᾶσαν ἐγὼ ἐποίησα.
\vs{17}Διʼ ἁμαρτίαν βραχύ τι ἐλύπησα αὐτὸν, καὶ ἐπάταξα αὐτὸν, καὶ ἀπέστρεψα τὸ πρόσωπόν μου ἀπʼ αὐτοῦ, καὶ ἐλυπήθη, καὶ ἐπορεύθη στυγνὸς ἐν ταῖς ὁδοῖς αὐτοῦ.
\vs{18}Τὰς ὁδοὺς αὐτοῦ ἑώρακα, καὶ ἰασάμην αὐτὸν, καὶ παρεκάλεσα αὐτὸν, καὶ ἔδωκα αὐτῷ παράκλησιν ἀληθινὴν,
\vs{19}εἰρήνην ἐπʼ εἰρήνῃ τοῖς μακρὰν καὶ τοῖς ἐγγὺς οὖσι· καὶ εἶπε Κύριος, ἰάσομαι αὐτοὺς.

\vs{20}Οἱ δὲ ἄδικοι κλυδωνισθήσονται, καὶ ἀναπαύσασθαι οὐ δυνήσονται.
\vs{21}Οὐκ ἔστι χαίρειν τοῖς ἀσεβέσιν, εἶπεν ὁ Θεός.

\ch{58}
Ἀναβοήσον ἐν ἰσχύϊ, καὶ μὴ φείσῃ, ὡς σάλπιγγι ὕψωσον τὴν φωνήν σου καὶ ἀνάγγειλον τῷ λαῷ μου τὰ ἁμαρτήματα αὐτῶν, καὶ τῷ οἴκῳ Ἰακὼβ τὰς ἀνομίας αὐτῶν.
\vs{2}Ἐμὲ ἡμέραν ἐξ ἡμέρας ζητοῦσι, καὶ γνῶναί μου τὰς ὁδοὺς ἐπιθυμοῦσιν, ὡς λαὸς δικαιοσύνην πεποιηκὼς καὶ κρίσιν Θεοῦ αὐτοῦ μὴ ἐγκαταλελοιπώς· αἰτοῦσί με νῦν κρίσιν δικαίαν, καὶ ἐγγίζειν Θεῷ ἐπιθυμοῦσι,
\vs{3}λέγοντες, τί ὅτι ἐνηστεύσαμεν, καὶ οὐκ εἶδες; ἐταπεινώσαμεν τὰς ψυχὰς ἡμῶν, καὶ οὐκ ἔγνως;

Ἐν γὰρ ταῖς ἡμέραις τῶν νηστειῶν ὑμῶν εὑρίσκετε τὰ θελήματα ὑμῶν, καὶ πάντας τοὺς ὑποχειρίους ὑμῶν ὑπονύσσετε.
\vs{4}Εἰ εἰς κρίσεις καὶ μάχας νηστεύετε, καὶ τύπτετε πυγμαῖς ταπεινὸν, ἱνατί μοι νηστεύετε ὡς σήμερον, ἀκουσθῆναι ἐν κραυγῇ τὴν φωνὴν ὑμῶν;
\vs{5}Οὐ ταύτην τὴν νηστείαν ἐξελεξάμην, καὶ ἡμέραν ταπεινοῦν ἄνθρωπον τὴν ψυχὴν αὐτοῦ· οὐδʼ ἂν κάμψῃς ὡς κρίκον τὸν τράχηλόν σου, καὶ σάκκον καὶ σποδὸν ὑποστρώσῃ, οὐδʼ οὕτω καλέσετε νηστείαν δεκτήν.
\vs{6}Οὐχὶ τοιαύτην νηστείαν ἐξελεξάμην, λέγει Κύριος· ἀλλὰ λύε πάντα σύνδεσμον ἀδικίας, διάλυε στραγγαλιὰς βιαίων συναλλαγμάτων, ἀπόστελλε τεθραυσμένους ἐν ἀφέσει, καὶ πᾶσαν συγγραφὴν ἄδικον διάσπα.
\vs{7}Διάθρυπτε πεινῶντι τὸν ἄρτον σου, καὶ πτωχοὺς ἀστέγους εἴσαγε εἰς τὸν οἶκόν σου· ἐὰν ἴδῃς γυμνὸν, περίβαλε, καὶ ἀπὸ τῶν οἰκείων τοῦ σπέρματός σου οὐχ ὑπερόψει.

\vs{8}Τότε ῥαγήσεται πρώϊμον τὸ φῶς σου, καὶ τὰ ἰάματά σου ταχὺ ἀνατελεῖ· καὶ προπορεύσεται ἔμπροσθέν σου ἡ δικαιοσύνη σου, καὶ ἡ δόξα τοῦ Θεοῦ περιστελεῖ σε.
\vs{9}Τότε βοήσῃ, καὶ ὁ Θεὸς εἰσακούσεταί σου, ἔτι λαλοῦντός σου ἐρεῖ, ἰδοὺ, πάρειμι· ἐὰν ἀφέλῃς ἀπὸ σοῦ σύνδεσμον, καὶ χειροτονίαν, καὶ ῥῆμα γογγυσμοῦ,
\vs{10}καὶ δῷς πεινῶντι τὸν ἄρτον ἐκ ψυχῆς σου, καὶ ψυχὴν τεταπεινωμένην ἐμπλήσῃς, τότε ἀνατελεῖ ἐν τῷ σκότει τὸ φῶς σου, καὶ τὸ σκότος σου ὡς μεσημβρία,
\vs{11}καὶ ἔσται ὁ Θεός σου μετὰ σοῦ διαπαντός· καὶ ἐμπλησθήσῃ καθάπερ ἐπιθυμεῖ ἡ ψυχή σου, καὶ τὰ ὀστᾶ σου πιανθήσεται· καὶ ἔσται ὡς κῆπος μεθύων, καὶ ὡς πηγὴ ἣν μὴ ἐξέλιπεν ὕδωρ·
\vs{12}Καὶ οἰκοδομηθήσονταί σου αἱ ἔρημοι αἰώνιοι, καὶ ἔσται τὰ θεμέλιά σου αἰώνια γενεῶν γενεαῖς, καὶ κληθήσῃ οἰκοδόμος φραγμῶν, καὶ τὰς τρίβους σου ἀναμέσον παύσεις.

\vs{13}Ἐὰν ἀποστρέψῃς τὸν πόδα σου ἀπὸ τῶν σαββάτων τοῦ μὴ ποιεῖν τὰ θελήματά σου ἐν τῇ ἡμέρᾳ τῇ ἁγίᾳ, καὶ καλέσεις τὰ σάββατα τρυφερὰ, ἅγια τῷ Θεῷ· οὐκ ἀρεῖς τὸν πόδα σου ἐπʼ ἔργῳ, οὐδὲ λαλήσεις λόγον ἐν ὀργῇ ἐκ τοῦ στόματός σου,
\vs{14}καὶ ἔσῃ πεποιθὼς ἐπὶ Κύριον, καὶ ἀναβιβάσει σε ἐπὶ τὰ ἀγαθὰ τῆς γῆς, καὶ ψωμιεῖ σε τὴν κληρονομίαν Ἰακὼβ τοῦ πατρός σου· τὸ γὰρ στόμα Κυρίου ἐλάλησε ταῦτα.

\ch{59}
Μὴ οὐκ ἰσχύει ἡ χεὶρ Κυρίου τοῦ σῶσαι; ἢ ἐβάρυνε τὸ οὖς αὐτοῦ τοῦ μὴ εἰσακοῦσαι;
\vs{2}Ἀλλὰ τὰ ἁμαρτήματα ὑμῶν διϊστῶσιν ἀναμέσον ὑμῶν καὶ ἀναμέσον τοῦ Θεοῦ, καὶ διὰ τὰς ἁμαρτίας ὑμῶν ἀπέστρεψεν τὸ πρόσωπον ἀφʼ ὑμῶν τοῦ μὴ ἐλεῆσαι.
\vs{3}Αἱ γὰρ χεῖρες ὑμῶν μεμολυσμέναι αἵματι, καὶ οἱ δάκτυλοι ὑμῶν ἐν ἁμαρτίαις, τὰ δὲ χείλη ὑμῶν ἐλάλησεν ἀνομίαν, καὶ ἡ γλῶσσα ὑμῶν ἀδικίαν μελετᾷ.

\vs{4}Οὐθεὶς λαλεῖ δίκαια, οὐδὲ ἐστι κρίσις ἀληθινή· πεποίθασιν ἐπὶ ματαίοις, καὶ λαλοῦσι κενὰ, ὅτι κύουσι πόνον, καὶ τίκτουσιν ἀνομίαν.
\vs{5}Ὠὰ ἀσπίδων ἔῤῥηξαν, καὶ ἱστὸν ἀράχνης ὑφαίνουσι, καὶ ὁ μέλλων τῶν ὠῶν αὐτῶν φαγεῖν, συντρίψας οὔριον, εὗρε καὶ ἐν αὐτῷ βασιλίσκου.
\vs{6}Ὁ ἱστὸς αὐτῶν οὐκ ἔσται εἰς ἱμάτιον, οὐδὲ μὴ περιβάλωνται ἀπὸ τῶν ἔργων αὐτῶν· τὰ γὰρ ἔργα αὐτῶν, ἔργα ἀνομίας.
\vs{7}Οἱ δὲ πόδες αὐτῶν ἐπὶ πονηρίαν τρέχουσι, ταχινοὶ ἐκχέαι αἷμα, καὶ οἱ διαλογισμοὶ αὐτὼν, διαλογισμοὶ ἀπὸ φόνων· σύντριμμα καὶ ταλαιπωρία ἐν ταῖς ὁδοῖς αὐτῶν,
\vs{8}καὶ ὁδὸν εἰρήνης οὐκ οἴδασι, καὶ οὐκ ἔστι κρίσις ἐν ταῖς ὁδοῖς αὐτῶν· αἱ γὰρ τρίβοι αὐτῶν διεστραμμέναι ἃς διοδεύουσι, καὶ οὐκ οἴδασιν εἰρήνην.

\vs{9}Διατοῦτο ἀπέστη ἡ κρίσις ἀπʼ αὐτῶν, καὶ οὐ μὴ καταλάβῃ αὐτοὺς δικαιοσύνη· ὑπομεινάντων αὐτῶν φῶς ἐγένετο αὐτοῖς σκότος, μείναντες αὐγὴν ἐν ἀωρίᾳ περιεπάτησαν.
\vs{10}Ψηλαφήσουσιν ὡς τυφλοὶ τοῖχον, καὶ ὡς οὐχ ὑπαρχόντων ὀφθαλμῶν ψηλαφήσουσι· καὶ πεσοῦνται ἐν μεσημβρίᾳ ὡς ἐν μεσονυκτίῳ, ὡς ἀποθνήσκοντες στενάξουσιν·
\vs{11}Ὡς ἄρκος καὶ ὡς περιστερὰ ἅμα πορεύσονται· ἀνεμείναμεν κρίσιν, καὶ οὐκ ἔστι σωτηρία, μακρὰν ἀφέστηκεν ἀφʼ ἡμῶν.

\vs{12}Πολλὴ γὰρ ἡμῶν ἡ ἀνομία ἐναντίον σου, καὶ αἱ ἁμαρτίαι ἡμῶν ἀντέστησαν ἡμῖν· αἱ γὰρ ἀνομίαι ἡμῶν ἐν ἡμῖν, καὶ τὰ ἀδικήματα ἡμῶν ἔγνωμεν.
\vs{13}Ἠσεβήσαμεν καὶ ἐψευσάμεθα, καὶ ἀπέστημεν ὄπισθεν τοῦ Θεοῦ ἡμῶν· ἐλαλήσαμεν ἄδικα, καὶ ἠπειθήσαμεν· ἐκύομεν, καὶ ἐμελετήσαμεν ἀπὸ καρδίας ἡμῶν λόγους ἀδίκους·
\vs{14}Καὶ ἀπεστήσαμεν ὀπίσω τὴν κρίσιν, καὶ ἡ δικαιοσύνη μακρὰν ἀφέστηκεν· ὅτι κατηναλώθη ἐν ταῖς ὁδοῖς αὐτῶν ἡ ἀλήθεια, καὶ διʼ εὐθείας οὐκ ἐδύναντο διελθεῖν.
\vs{15}Καὶ ἡ ἀλήθεια ᾖρται, καὶ μετέστησαν τὴν διάνοιαν τοῦ συνιέναι.

Καὶ εἶδε Κύριος, καὶ οὐκ ἤρεσεν αὐτῷ, ὅτι οὐκ ἦν κρίσις.
\vs{16}Καὶ εἶδε, καὶ οὐκ ἦν ἀνήρ, καὶ κατενοήσε, καὶ οὐκ ἦν ὁ ἀντιληψόμενος· καὶ ἠμύνατο αὐτοὺς τῷ βραχίονι αὐτοῦ, καὶ τῇ ἐλεημοσύνῃ ἐστηρίσατο.
\vs{17}Καὶ ἐνεδύσατο δικαιοσύνην ὡς θώρακα, καὶ περιέθετο περικεφαλαίαν σωτηρίου ἐπὶ τῆς κεφαλῆς, καὶ περιεβάλετο ἱμάτιον ἐκδικήσεως, καὶ τὸ περιβόλαιον αὐτοῦ,
\vs{18}ὡς ἀνταποδώσων ἀνταπόδοσιν ὄνειδος τοῖς ὑπεναντίοις.
\vs{19}Καὶ φοβηθήσονται οἱ ἀπὸ δυσμῶν τὸ ὄνομα Κυρίου, καὶ οἱ ἀπʼ ἀνατολῶν ἡλίου τὸ ὄνομα τὸ ἔνδοξον· ἥξει γὰρ ὡς ποταμὸς βίαιος ἡ ὀργὴ παρὰ Κυρίου, ἥξει μετὰ θυμοῦ.

\vs{20}Καὶ ἥξει ἕνεκεν Σιὼν ὁ ῥυόμενος, καὶ ἀποστρέψει ἀσεβείας ἀπὸ Ἰακώβ.
\vs{21}Καὶ αὕτη αὐτοῖς ἡ παρʼ ἐμοῦ διαθήκη, εἶπε Κύριος· τὸ πνεῦμα τὸ ἐμὸν, ὅ ἐστιν ἐπὶ σοί, καὶ τὰ ῥήματα, ἃ ἔδωκα εἰς τὸ στόμα σου, οὐ μὴ ἐκλίπῃ ἐκ τοῦ στόματός σου, καὶ ἐκ τοῦ στόματος τοῦ σπέρματός σου· εἶπε γὰρ Κύριος ἀπὸ τοῦ νῦν καὶ εἰς τὸν αἰῶνα.

\ch{60}
Φωτίζου φωτίζου Ἱερουσαλὴμ, ἥκει γάρ σου τὸ φῶς, καὶ ἡ δόξα Κυρίου ἐπὶ σὲ ἀνατέταλκεν.
\vs{2}Ἰδοὺ, σκότος καλύψει γῆν, καὶ γνόφος ἐπʼ ἔθνη, ἐπὶ δὲ σὲ φανήσεται Κύριος, καὶ ἡ δόξα αὐτοῦ ἐπὶ σὲ ὀφθήσεται.
\vs{3}Καὶ πορεύσονται βασιλεῖς τῷ φωτί σου, καὶ ἔθνη τῇ λαμπρότητί σου.

\vs{4}Ἆρον κύκλῳ τοὺς ὀφθαλμούς σου, καὶ ἴδε συνηγμένα τὰ τέκνα σου· ἥκασι πάντες οἱ υἱοί σου μακρόθεν, καὶ αἱ θυγατέρες σου ἐπʼ ὤμων ἀρθήσονται.
\vs{5}Τότε ὄψῃ, καὶ φοβηθήσῃ, καὶ ἐκστήσῃ τῇ καρδίᾳ, ὅτι μεταβαλεῖ εἰς σὲ πλοῦτος θαλάσσης, καὶ ἐθνῶν καὶ λαῶν, καὶ ἥξουσί σοι ἀγέλαι καμήλων,
\vs{6}καὶ καλύψουσί σε κάμηλοι Μαδιὰμ καὶ Γαιφά· πάντες ἐκ Σαβὰ ἥξουσι φέροντες χρυσίον καὶ λίβανον οἴσουσιν, καὶ τὸ σωτήριον Κυρίου εὐαγγελιοῦνται.
\vs{7}Καὶ πάντα τὰ πρόβατα Κηδὰρ συναχθήσονται, καὶ κριοὶ Ναβαιὼθ ἥξουσι, καὶ ἀνενεχθήσεται δεκτὰ ἐπὶ τὸ θυσιαστήριόν μου, καὶ ὁ οἶκος τῆς προσευχῆς μου δοξασθήσεται.

\vs{8}Τίνες οἵδε, ὡς νεφέλαι πέτονται, καὶ ὡσεὶ περιστεραὶ σὺν νοσσοῖς ἐπʼ ἐμέ;
\vs{9}Ἐμὲ αἱ νῆσοι ὑπέμειναν, καὶ πλοῖα Θαρσεὶς ἐν πρώτοις, ἀγαγεῖν τὰ τέκνα σου μακρόθεν, καὶ τὸν ἄργυρον καὶ τὸν χρυσὸν αὐτῶν μετʼ αὐτῶν, καὶ διὰ τὸ ὄνομα Κυρίου τὸ ἅγιον, καὶ διὰ τὸ τὸν ἅγιον τοῦ Ἰσραὴλ ἔνδοξον εἶναι.
\vs{10}Καὶ οἰκοδομήσουσιν ἀλλογενεῖς τὰ τείχη σου, καὶ οἱ βασιλεῖς αὐτῶν παραστήσονταί σοι· διὰ γὰρ ὀργήν μου ἐπάταξά σε, καὶ διὰ ἔλεον ἠγάπησά σε.
\vs{11}Καὶ ἀνοιχθήσονται αἱ πύλαι σου διαπαντὸς, ἡμέρας καὶ νυκτὸς οὐ κλεισθήσονται, εἰσαγαγεῖν πρὸς σὲ δύναμιν ἐθνῶν, καὶ βασιλεῖς αὐτῶν ἀγομένους.
\vs{12}Τὰ γὰρ ἔθνη καὶ οἱ βασιλεῖς, οἵτινες οὐ δουλεύσουσί σοι, ἀπολοῦνται, καὶ τὰ ἔθνη ἐρημίᾳ ἐρημωθήσεται.

\vs{13}Καὶ ἡ δόξα τοῦ Λιβάνου πρὸς σὲ ἥξει, ἐν κυπαρίσσῳ καὶ πεύκῃ καὶ κέδρῳ ἅμα, δοξάσαι τὸν τόπον τὸν ἅγιόν μου.

\vs{14}Καὶ πορεύσονται πρὸς σὲ δεδοικότες υἱοὶ ταπεινωσάντων σε, καὶ παροξυνάντων σε, καὶ κληθήσῃ Πόλις Σιὼν ἀγίου Ἰσραήλ.
\vs{15}Διὰ τὸ γεγενῆσθαί σε ἐνκαταλελειμμένην καὶ μεμισημένην, καὶ οὐκ ἦν ὁ βοηθῶν· καὶ θήσω σε ἀγαλλίαμα αἰώνιον, εὐφροσύνην γενεῶν γενεαῖς.

\vs{16}Καὶ θηλάσεις γάλα ἐθνῶν, καὶ πλοῦτον βασιλέων φάγεσαι, καὶ γνώσῃ ὅτι ἐγὼ Κύριος ὁ σώζων σε, καὶ ἐξαιρούμενός σε Θεὸς Ἰσραήλ.
\vs{17}Καὶ ἀντὶ χαλκοῦ οἴσω σοι χρυσίον, ἀντὶ δὲ σιδήρου οἴσω σοι ἀργύριον, ἀντὶ δὲ ξύλων οἴσω σοι χαλκόν, ἀντὶ δὲ λίθων, σιδήρον· καὶ δώσω τοὺς ἄρχοντάς σου ἐν εἰρήνῃ, καὶ τοὺς ἐπισκόπους σου ἐν δικαιοσύνῃ·
\vs{18}Καὶ οὐκ ἀκουσθήσεται ἔτι ἀδικία ἐν τῇ γῇ σου, οὐδὲ σύντριμμα, οὐδὲ ταλαιπωρία ἐν τοῖς ὁρίοις σου, ἀλλὰ κληθήσεται Σωτήριον τὰ τείχη σου, καὶ αἱ πύλαι σου Γλύμμα.
\vs{19}Καὶ οὐκ ἔσται σοι ἔτι ὁ ἥλιος εἰς φῶς ἡμέρας, οὐδὲ ἀνατολὴ σελήνης φωτιεῖ σοι τὴν νύκτα, ἀλλʼ ἔσται σοι Κύριος φῶς αἰώνιον, καὶ ὁ Θεὸς δόξα σου.
\vs{20}Οὐ γὰρ δύσεται ὁ ἥλιος σοι, καὶ ἡ σελήνη σοι οὐκ ἐκλείψει· ἔσται γὰρ σοι Κύριός φῶς αἰώνιον, καὶ ἀναπληρωθήσονται αἱ ἡμέραι τοῦ πένθους σου.
\vs{21}Καὶ ὁ λαός σου πᾶς δίκαιος, διʼ αἰῶνος κληρονομήσουσι τὴν γῆν, φυλάσσων τὸ φύτευμα, ἔργα χειρῶν αὐτοῦ, εἰς δόξαν.
\vs{22}Ὁ ὀλιγοστὸς ἔσται εἰς χιλιάδας, καὶ ὁ ἐλάχιστος εἰς ἔθνος μέγα· ἐγὼ Κύριος κατὰ καιρὸν συνάξω αὐτούς.

\ch{61}
Πνεῦμα Κυρίου ἐπʼ ἐμὲ, οὗ εἵνεκε ἔχρισέν με, εὐαγγελίσασθαι πτωχοῖς ἀπέσταλκέ με, ἰάσασθαι τοὺς συντετριμμένους τὴν καρδίαν, κηρῦξαι αἰχμαλώτοις ἄφεσιν, καὶ τυφλοῖς ἀνάβλεψιν,
\vs{2}καλέσαι ἐνιαυτὸν Κυρίου δεκτὸν, καὶ ἡμέραν ἀνταποδόσεως, παρακαλέσαι πάντας τοὺς πενθοῦντας,
\vs{3}δοθῆναι τοῖς πενθοῦσι Σιὼν αὐτοῖς δόξαν ἀντὶ σποδοῦ, ἄλειμμα εὐφροσύνης τοῖς πενθοῦσι, κατὰ στολὴν δόξης ἀντὶ πνεύματος ἀκηδίας· καὶ κληθήσονται γενεαὶ δικαιοσύνης, φύτευμα Κυρίου εἰς δόξαν.

\vs{4}Καὶ οἰκοδομήσουσιν ἐρήμους αἰωνίας, ἐξηρημωμένας πρότερον ἐξαναστήσουσι, καὶ καινιοῦσι πόλεις ἐρήμους, ἐξηρημωμένας εἰς γενεάς.
\vs{5}Καὶ ἥξουσιν ἀλλογενεῖς ποιμαίνοντες τὰ πρόβατά σου, καὶ ἀλλόφυλοι ἀροτῆρες, καὶ ἀμπελουργοί.
\vs{6}Ὑμεῖς δὲ ἱερεῖς Κυρίου κληθήσεσθε, λειτουργοὶ Θεοῦ, ἰσχὺν ἐθνῶν κατέδεσθε, καὶ ἐν τῷ πλούτῳ αὐτῶν θαυμασθήσεσθε.
\vs{7}Οὕτως ἐκ δευτέρας κληρονομήσουσι τὴν γῆν, καὶ εὐφροσύνη αἰώνιος ὑπὲρ κεφαλῆς αὐτῶν.
\vs{8}Ἐγὼ γάρ εἰμι Κύριος ὁ ἀγαπῶν δικαιοσύνην, καὶ μισῶν ἁρπάγματα ἐξ ἀδικίας· καὶ δώσω τὸν μόχθον αὐτῶν δικαίοις, καὶ διαθήκην αἰώνιον διαθήσομαι αὐτοῖς.
\vs{9}Καὶ γνωσθήσεται ἐν τοῖς ἔθνεσι τὸ σπέρμα αὐτῶν, καὶ τὰ ἔκγονα αὐτῶν ἐν μέσῳ τῶν λαῶν· πᾶς ὁ ὁρῶν αὐτοὺς ἐπιγνώσεται αὐτοὺς, ὅτι οὗτοί εἰσι σπέρμα ηὐλογημένον ὑπὸ Θεοῦ,
\vs{10}καὶ εὐφροσύνῃ εὐφρανθήσονται ἐπὶ Κύριον.

Ἀγαλλιάσθω ἡ ψυχή μου ἐπὶ τῷ Κυρίῳ, ἐνέδυσε γάρ με ἱμάτιον σωτηρίου, καὶ χιτῶνα εὐφροσύνης, ὡς νυμφίῳ περιέθηκέ μοι μίτραν, καὶ ὡς νύμφην κατεκόσμησέ με κόσμῳ.

\vs{11}Καὶ ὡς γῆν αὔξουσαν τὸ ἄνθος αὐτῆς, καὶ ὡς κῆπος τὰ σπέρματα αὐτοῦ, οὕτως ἀνατελεῖ Κύριος Κύριος δικαιοσύνην, καὶ ἀγαλλίαμα ἐναντίον πάντων τῶν ἐθνῶν.

\ch{62}
Διὰ Σιὼν οὐ σιωπήσομαι, καὶ διὰ Ἱερουσαλὴμ οὐκ ἀνήσω, ἕως ἂν ἐξέλθῃ ὡς φῶς ἡ δικαιοσύνη αὐτῆς, τὸ δὲ σωτήριόν μου ὡς λαμπὰς καυθήσεται.
\vs{2}Καὶ ὄψονται ἔθνη τὴν δικαιοσύνην σου, καὶ βασιλεῖς τὴν δόξαν σου, καὶ καλέσει σε τὸ ὄνομα τὸ καινὸν, ὃ ὁ Κύριος ὀνομάσει αὐτό.
\vs{3}Καὶ ἔσῃ στέφανος κάλλους ἐν χειρὶ Κυρίου, καὶ διάδημα βασιλείας ἐν χειρὶ Θεοῦ σου.
\vs{4}Καὶ οὐκέτι κληθήσῃ Καταλελειμμένη, καὶ ἡ γῆ σου οὐ κληθήσεται ἔτι Ἔρημος· σοὶ γὰρ κληθήσεται, Θέλημα ἐμὸν, καὶ τῇ γῇ σου, Οἰκουμένη, ὅτι εὐδόκησε Κύριος ἐν σοὶ, καὶ ἡ γῆ σου συνοικισθήσεται.

\vs{5}Καὶ ὡς συνοικῶν νεανίσκος παρθένῳ, οὕτω κατοικήσουσιν οἱ υἱοί σου· καὶ ἔσται ὃν τρόπον εὐφρανθήσεται νυμφίος ἐπὶ νύμφῃ, οὕτως εὐφρανθήσεται Κύριος ἐπὶ σοί.

\vs{6}Καὶ ἐπὶ τῶν τειχῶν σου Ἱερουσαλὴμ κατέστησα φύλακας ὅλην τὴν ἡμέραν καὶ ὅλην τὴν νύκτα, οἳ διὰ τέλους οὐ σιωπήσονται μιμνησκόμενοι Κυρίου.
\vs{7}Οὐκ ἔστι γὰρ ὑμῖν ὅμοιος· ἐὰν διορθώσῃ, καὶ ποιήσῃ Ἱερουσαλὴμ γαυρίαμα ἐπὶ τῆς γῆς.
\vs{8}Ὤμοσε Κύριος κατὰ τῆς δόξης αὐτοῦ, καὶ κατὰ τῆς ἰσχύος τοῦ βραχίονος αὐτοῦ, εἰ ἔτι δώσω τὸν σῖτόν σου, καὶ τὰ βρώματα σου τοῖς ἐχθροῖς σου, καὶ εἰ ἔτι πίονται υἱοὶ ἀλλότριοι τὸν οἶνόν σου, ἐφʼ ᾧ ἐμόχθησας.
\vs{9}Ἀλλʼ οἱ συναγαγόντες φάγονται αὐτὰ, καὶ αἰνέσουσι Κύριον, καὶ οἱ συναγαγόντες πίονται αὐτὰ ἐν ταῖς ἐπαύλεσι ταῖς ἁγίαις μου.

\vs{10}Πορεύεσθε διὰ τῶν πυλῶν μοῦ, καὶ ὁδοποιήσατε τῷ λαῷ μου, καὶ τοὺς λίθους ἐκ τῆς ὁδοῦ διαῤῥίψατε, ἐξάρατε σύσσημον εἰς τὰ ἔθνη.
\vs{11}Ἰδοὺ γὰρ Κύριος ἐποίησεν ἀκουστὸν ἕως ἐσχάτου τῆς γῆς· εἴπατε τῇ θυγατρὶ Σιὼν, ἰδοὺ ὁ σωτήρ σοι παραγέγονεν ἔχων τὸν ἑαυτοῦ μισθὸν, καὶ τὸ ἔργον αὐτοῦ πρὸ προσώπου αὐτοῦ.
\vs{12}Καὶ καλέσει αὐτὸν Λαὸν ἅγιον, λελυτρωμένον ὑπὸ Κυρίου· σὺ δὲ κληθήσῃ Ἐπιζητουμένη πόλις, καὶ οὐκ ἐγκαταλελιμμένη.

\ch{63}
Τίς οὗτος ὁ παραγνόμενος ἐξ Ἐδὼμ, ἐρύθημα ἱματίων ἐκ Βοσόρ; οὕτως ὡραῖος ἐν στολῇ, βίᾳ μετὰ ἰσχύος; ἐγὼ διαλέγομαι δικαιοσύνην καὶ κρίσιν σωτηρίου.

\vs{2}Διατί σου ἐρυθρὰ τὰ ἱμάτια, καὶ τὰ ἐνδύματά σου ὡς ἀπὸ πατητοῦ ληνοῦ;
\vs{3}Πλήρης καταπεπατημένης, καὶ τῶν ἐθνῶν οὐκ ἔστιν ἀνὴρ μετʼ ἐμοῦ, καὶ κατεπάτησα αὐτοὺς ἐν θυμῷ μου, καὶ κατέθλασα αὐτοὺς ὡς γῆν, καὶ κατήγαγον τὸ αἷμα αὐτῶν εἰς γῆν.
\vs{4}Ἡμέρα γὰρ ἀνταποδόσεως ἐπῆλθεν αὐτοῖς, καὶ ἐνιαυτὸς λυτρώσεως πάρεστι.
\vs{5}Καὶ ἐπέβλεψα, καὶ οὐκ ἦν βοηθός· καὶ προσενόησα, καὶ οὐθεὶς ἀντελαμβάνετο· καὶ ἐῤῥύσατο αὐτοὺς ὁ βραχίων μου, καὶ ὁ θυμός μου ἐπέστη.
\vs{6}Καὶ κατεπάτησα αὐτοὺς τῇ ὀργῇ μου, καὶ κατήγαγον τὸ αἷμα αὐτῶν εἰς γῆν.

\vs{7}Τὸν ἔλεον Κυρίου ἐμνήσθην, τὰς ἀρετὰς Κυρίου ἐν πᾶσιν οἷς ἡμῖν ἀνταποδίδωσι· Κύριος κριτὴν ἀγαθὸς τῷ οἴκῳ Ἰσραὴλ, ἐπάγει ἡμῖν κατὰ τὸ ἔλεος αὐτοῦ, καὶ κατὰ τὸ πλῆθος τῆς δικαιοσύνης αὐτοῦ.

\vs{8}Καὶ εἶπεν, οὐχ ὁ λαός μου; τέκνα οὐ μὴ ἀθετήσωσι· καὶ ἐγένετο αὐτοῖς εἰς σωτηρίαν
\vs{9}ἐκ πάσης θλίψεως αὐτῶν· οὐ πρέσβυς, οὐδὲ ἄγγελος, ἀλλʼ αὐτὸς ἔσωσεν αὐτοὺς, διὰ τὸ ἀγαπᾷν αὐτοὺς καὶ φείδεσθαι αὐτῶν· αὐτὸς ἐλυτρώσατο αὐτοὺς, καὶ ἀνέλαβεν αὐτοὺς, καὶ ὕψωσεν αὐτοὺς πάσας τὰς ἡμέρας τοῦ αἰῶνος·

\vs{10}Αὐτοὶ δὲ ἠπείθησαν, καὶ παρώξυναν τὸ πνεῦμα τὸ ἅγιον αὐτοῦ· καὶ ἐστράφη αὐτοῖς εἰς ἔχθραν, αὐτὸς ἐπολέμησεν αὐτούς.
\vs{11}Καὶ ἐμνήσθη ἡμερῶν αἰωνίων· ποῦ ὁ ἀναβιβάσας ἐκ τῆς θαλάσσης τὸν ποιμένα τῶν προβάτων; ποῦ ἐστιν ὁ θεὶς ἐν αὐτοῖς τὸ πνεῦμα τὸ ἅγιον;
\vs{12}ὁ ἀγαγὼν τῇ δεξιᾷ Μωσῆν, ὁ βραχίων τῆς δόξης αὐτοῦ; κατίσχυσεν ὕδωρ ἀπὸ προσώπου αὐτοῦ, ποιῆσαι ἑαυτῷ ὄνομα αἰώνιον.
\vs{13}Ἤγαγεν αὐτοὺς διʼ ἀβύσσου, ὡς ἵππον διʼ ἐρήμου, καὶ οὐκ ἐκοπίασαν,
\vs{14}καὶ ὡς κτήνη διὰ πεδίου· κατέβη πνεῦμα παρὰ Κυρίου, καὶ ὡδήγησεν αὐτούς· οὕτως ἤγαγες τὸν λαόν σου ποιῆσαι σεαυτῷ ὄνομα δόξης.

\vs{15}Ἐπίστρεψον ἐκ τοῦ οὐρανοῦ, καὶ ἴδε ἐκ τοῦ οἴκου τοῦ ἁγίου σου, καὶ δόξης σου· ποῦ ἐστιν ὁ ζῆλός σου καὶ ἡ ἰσχύς σου; ποῦ ἐστι τὸ πλῆθος τοῦ ἐλέους σου, καὶ οἰκτιρμῶν σου, ὅτι ἀνέσχου ἡμῶν;
\vs{16}Σὺ γὰρ εἶ πατὴρ ἡμῶν, ὅτι Ἁβραὰμ οὐκ ἔγνω ἡμᾶς, καὶ Ἰσραὴλ οὐκ ἐπέγνω ἡμᾶς, ἀλλὰ σὺ Κύριε πατὴρ ἡμῶν ῥῦσαι ἡμᾶς, ἀπʼ ἀρχῆς τὸ ὄνομά σου ἐφʼ ἡμᾶς ἐστι.

\vs{17}Τί ἐπλάνησας ἡμᾶς, Κύριε, ἀπὸ τῆς ὁδοῦ σου; ἐσκλήρυνας τὰς καρδίας ἡμῶν, τοῦ μὴ φοβεῖσθαί σε; ἐπίστρεψον διὰ τοὺς δούλους σου, διὰ τὰς φυλὰς τῆς κληρονομίας σου,
\vs{18}ἵνα μικρὸν κληρονομήσωμεν τοῦ ὄρους τοῦ ἁγίου σου.
\vs{19}Ἐγενόμεθα ὡς τὸ ἀπʼ ἀρχῆς ὅτε οὐκ ἦρξας ἡμῶν, οὐδὲ ἐκλήθη τὸ ὄνομά σου ἐφʼ ἡμᾶς.

Ἐὰν ἀνοίξῃς τὸν οὐρανὸν, τρόμος λήψεται ἀπὸ σοῦ ὄρη, καὶ τακήσονται,

\ch{64}ὡς κηρὸς ἀπὸ προσώπου πυρὸς τήκεται, καὶ κατακαύσει πῦρ τοὺς ὑπεναντίους, καὶ φανερὸν ἔσται τὸ ὄνομά σου ἐν τοῖς ὑπεναντίοις· ἀπὸ προσώπου σου ἔθνη ταραχθήσονται,
\vs{2}ὅταν ποιῇς τὰ ἔνδοξα· τρόμος λήψεται ἀπὸ σοῦ ὄρη.

\vs{3}Ἀπὸ τοῦ αἰῶνος οὐκ ἠκούσαμεν, οὐδὲ οἱ ὀφθαλμοὶ ἡμῶν εἶδον Θεὸν πλήν σου, καὶ τὰ ἔργα σου, ἃ ποιήσεις τοῖς ὑπομένουσιν ἔλεον.
\vs{4}Συναντήσεται γὰρ τοῖς ποιοῦσι τὸ δίκαιον, καὶ τῶν ὁδῶν σου μνησθήσονται· ἰδοὺ σὺ ὠργίσθης, καὶ ἡμεῖς ἡμάρτομεν· διατοῦτο ἐπλανήθημεν,
\vs{5}καὶ ἐγενήθημεν ὡς ἀκάθαρτοι πάντες ἡμεῖς, ὡς ῥάκος ἀποκαθημένης πᾶσα ἡ δικαιοσύνη ἡμῶν· καὶ ἐξεῤῥύημεν ὡς φύλλα διὰ τὰς ἀνομίας ἡμῶν· οὕτως ἄνεμος οἴσει ἡμᾶς.
\vs{6}Καὶ οὐκ ἔστιν ὁ ἐπικαλούμενος τὸ ὄνομά σου, καὶ ὁ μνησθεὶς ἀντιλαβέσθαι σου· ὅτι ἀπέστρεψας τὸ πρόσωπόν σου ἀφʼ ἡμῶν, καὶ παρέδωκας ἡμᾶς διὰ τὰς ἁμαρτίας ἡμῶν.

\vs{7}Καὶ νῦν Κύριε, πατὴρ ἡμῶν σὺ, ἡμεῖς δὲ πηλός, ἔργα τῶν χειρῶν σου πάντες.
\vs{8}Μὴ ὀργίζου ἡμῖν σφόδρα, καὶ μὴ ἐν καιρῷ μνησθῇς ἁμαρτιῶν ἡμῶν· καὶ νῦν ἐπίβλεψον, ὅτι λαός σου πάντες ἡμεῖς.
\vs{9}Πόλις τοῦ ἁγίου σου ἐγενήθη ἔρημος, Σιὼν ὡς ἔρημος ἐγενήθη, Ἱερουσαλὴμ εἰς κατάραν.
\vs{10}Ὁ οἶκος τὸ ἅγιον ἡμῶν, καὶ ἡ δόξα ἣν εὐλόγησαν οἱ πατέρες ἡμῶν, ἑγενήθη πυρίκαυστος, καὶ πάντα ἔνδοξα ἡμῶν συνέπεσε.
\vs{11}Καὶ ἐπὶ πᾶσι τούτοις ἀνέσχου Κύριε, καὶ ἐσιώπησας, καὶ ἐταπείνωσας ἡμᾶς σφόδρα.

\ch{65}
Ἐμφανῆς ἐγενήθην τοῖς ἐμὲ μὴ ἐπερωτῶσιν, εὑρέθην τοῖς ἐμὲ μὴ ζητοῦσιν· εἶπα, ἰδοὺ εἰμὶ τῷ ἔθνει, οἳ οὐκ ἐκάλεσάν μου τὸ ὄνομα.
\vs{2}Ἐξεπέτασα τὰς χεῖράς μου ὅλην τὴν ἡμέραν πρὸς λαὸν ἀπειθοῦντα καὶ ἀντιλέγοντα, τοῖς πορευομένοις ὁδῷ οὐ καλῇ, ἀλλʼ ὀπίσω τῶν ἁμαρτιῶν αὐτῶν.
\vs{3}Ὁ λαὸς οὗτος ὁ παροξύνων με ἐναντίον ἐμοῦ διαπαντός· αὐτοὶ θυσιάζουσιν ἐν τοῖς κήποις, καὶ θυμιῶσιν ἐπὶ ταῖς πλίνθοις τοῖς δαιμονίοις, ἃ οὐκ ἔστιν·
\vs{4}Ἐν τοῖς μνήμασι, καὶ ἐν τοῖς σπηλαίοις κοιμῶνται διὰ ἐνὺπνια, οἱ ἔσθοντες κρέας ὕειον, καὶ ζωμὸν θυσιῶν, μεμολυμμένα πάντα τὰ σκεύη αὐτῶν,
\vs{5}οἱ λέγοντες, πόῤῥω ἀπʼ ἐμοῦ, μὴ ἐγγίσῃς μοι, ὅτι καθαρός εἰμι·

Οὗτος καπνὸς τοῦ θυμοῦ μου, πῦρ καίεται ἐν αὐτῷ πάσας τὰς ἡμέρας.
\vs{6}Ἰδοὺ, γέγραπται ἐνώπιόν μου, οὐ σιωπήσω ἕως ἂν ἀποδώσω εἰς τὸν κόλπον αὐτῶν
\vs{7}τὰς ἁμαρτίας αὐτῶν, καὶ τῶν πατέρων αὐτῶν, λέγει Κύριος· οἳ ἐθυμίασαν ἐπὶ τῶν ὀρέων, καὶ ἐπὶ τῶν βουνῶν ὠνείδισάν με, ἀποδώσω τὰ ἔργα αὐτῶν εἰς τὸν κόλπον αὐτῶν.

\vs{8}Οὕτως λέγει Κύριος, ὃν τρόπον εὑρεθήσεται ὁ ῥὼξ ἐν τῷ βότρυϊ, καὶ ἐροῦσι, μὴ λυμήνῃ αὐτὸν, ὅτι εὐλογία ἐστὶν ἐν αὐτῷ, οὕτως ποιήσω ἕνεκεν τοῦ δουλεύοντός μοι, τούτου ἕνεκεν οὐ μὴ ἀπολέσω πάντας.
\vs{9}Καὶ ἐξάξω τὸ ἐξ Ἰακὼβ σπέρμα καὶ ἐξ Ἰούδα, καὶ κληρονομήσει τὸ ὄρος τὸ ἅγιόν μου, καὶ κληρονομήσουσιν οἱ ἐκλεκτοί μου, καὶ οἱ δοῦλοί μου, καὶ κατοικήσουσιν ἐκεῖ.
\vs{10}Καὶ ἔσονται ἐν τῷ δρυμῷ ἐπαύλεις ποιμνίων, καὶ φάραγξ Ἀχὼρ εἰς ἀνάπαυσιν βουκολίων τῷ λαῷ μου, οἳ ἐζήτησάν με.

\vs{11}Ὑμεῖς δὲ οἱ ἐγκαταλιπόντες με, καὶ ἐπιλανθανόμενοι τὸ ὄρος τὸ ὄρος τὸ ἅγιόν μου, καὶ ἐτοιμάζοντες τῷ δαιμονίῳ τράπεζαν, καὶ πληροῦντες τῇ τύχῃ κέρασμα,
\vs{12}ἐγὼ παραδώσω ὑμᾶς εἰς μάχαιραν, πάντες ἐν σφαγῇ πεσεῖσθε· ὅτι ἐκάλεσα ὑμᾶς, καὶ οὐχ ὑπηκούσατε· ἐλάλησα, καὶ παρηκούσατε, καὶ ἐποιήσατε τὸ πονηρὸν ἐναντίον ἐμοῦ, καὶ ἃ οὐκ ἐβουλόμην, ἐξελέξασθε.
\vs{13}Διατοῦτο τάδε λέγει Κύριος, ἰδοὺ, οἱ δουλεύοντές μοι φάγονται, ὑμεῖς δὲ πεινάσετε· ἰδοὺ, οἱ δουλεύοντές μοι πίονται, ὑμεῖς δὲ διψήσετε· ἰδοὺ, οἱ δουλεύοντές μοι εὐφρανθήσονται, ὑμεῖς δὲ αἰσχυνθήσεσθε·
\vs{14}ἰδοὺ, οἱ δουλεύοντές μοι ἀγαλλιάσονται ἐν εὐφροσύνῃ, ὑμεῖς δὲ κεκράξεσθε διὰ τὸν πόνον τῆς καρδίας ὑμῶν, καὶ ἀπὸ συντριβῆς πνεύματος ὑμῶν ὀλολύξετε.
\vs{15}Καταλείψετε γὰρ τὸ ὄνομα ὑμῶν εἰς πλησμονὴν τοῖς ἐκλεκτοῖς μου, ὑμᾶς δὲ ἀνελεῖ Κύριος, τοῖς δὲ δουλεύουσί μοι κληθήσεται ὄνομα καινὸν,
\vs{16}ὃ εὐλογηθήσεται ἐπὶ τῆς γῆς, εὐλογήσουσι γὰρ τὸν Θεὸν τὸν ἀληθινόν· καὶ οἱ ὀμνύοντες ἐπὶ τῆς γῆς, ὀμοῦνται τὸν Θεὸν τὸν ἀληθινόν· ἐπιλήσονται γὰρ τὴν θλίψιν τὴν πρώτην, καὶ οὐκ ἀναβήσεται αὐτῶν ἐπὶ τὴν καρδίαν.

\vs{17}Ἔσται γὰρ ὁ οὐρανὸς καινὸς, καὶ ἡ γῆ καινὴ, καὶ οὐ μὴ μνησθῶσι τῶν προτέρων, οὐδʼ οὐ μὴ ἐπέλθῃ αὐτῶν ἐπὶ τὴν καρδίαν,
\vs{18}ἀλλʼ εὐφροσύνην καὶ ἀγαλλίαμα εὑρήσουσιν ἐν αὐτῇ· ὅτι ἰδοὺ ἐγὼ ποιῶ ἀγαλλίαμα Ἱερουσαλὴμ, καὶ τὸν λαόν μου εὐφροσύνην.
\vs{19}Καὶ ἀγαλλιάσομαι ἐπὶ Ἱερουσαλὴμ, καὶ εὐφρανθήσομαι ἐπὶ τῷ λαῷ μου· καὶ οὐκέτι μὴ ἀκουσθῇ ἐν αὐτῇ φωνὴ κλαῦθμου, καὶ φωνὴ κραυγῆς,
\vs{20}οὐδʼ οὐ μὴ γένηται ἔτι ἐκεῖ ἄωρος καὶ πρεσβύτης, ὃς οὐκ ἐμπλήσει τὸν χρόνον αὐτοῦ· ἔσται γὰρ ὁ νέος ἐκατὸν ἐτῶν, ὁ δὲ ἀποθνήσκων ἁμαρτωλὸς ἑκατὸν ἐτῶν, καὶ ἐπικατάρατος ἔσται.
\vs{21}Καὶ οἰκοδομήσουσιν οἰκίας, καὶ αὐτοὶ ἐνοικήσουσι· καὶ καταφυτεύσουσιν ἀμπελῶνας, καὶ αὐτοὶ φάγονται τὰ γεννήματα αὐτῶν.
\vs{22}Οὐ μὴ οἰκοδομήσουσι, καὶ ἄλλοι ἐνοικήσουσι, καὶ οὐ μὴ φυτεύσουσι, καὶ ἄλλοι φάγονται· κατὰ γὰρ τὰς ἡμέρας τοῦ ξύλου τῆς ζωῆς ἔσονται αἱ ἡμέραι τοῦ λαοῦ μου· τὰ γὰρ ἔργα τῶν πόνων αὐτῶν παλαιώσουσιν.
\vs{23}Οἱ ἐκλεκτοί μου οὐ κοπιάσουσιν εἰς κενὸν, οὐδὲ τεκνοποιήσουσιν εἰς κατάραν, ὅτι σπέρμα εὐλογημένον ὑπὸ Θεοῦ ἐστι, καὶ τὰ ἔκγονα αὐτῶν μετʼ αὐτῶν.

\vs{24}Καὶ ἔσται πρὶν ἢ κεκράξαι αὐτοὺς, ἐγὼ ὑπακούσομαι αὐτῶν· ἔτι λαλούντων αὐτῶν, ἐρῶ, τί ἐστι;
\vs{25}Τότε λύκοι καὶ ἄρνες βοσκηθήσονται ἅμα, καὶ λέων ὡς βοῦς φάγεται ἄχυρα, ὄφις δὲ γῆν ὡς ἄρτον· οὐκ ἀδικήσουσιν, οὐδὲ λυμανοῦνται ἐπὶ τῷ ὄρει τῷ ἁγίῳ μου, λέγει Κύριος.

\ch{66}
Οὕτως λέγει Κύριος, ὁ οὐρανός μου θρόνος, καὶ ἡ γῆ ὑποπόδιον τῶν ποδῶν μου· ποῖον οἶκον οἰκοδομήσετέ μοι; καὶ ποῖος τόπος τῆς καταπαύσεώς μου;
\vs{2}Πάντα γὰρ ταῦτα ἐποίησεν ἡ χείρ μου, καὶ ἔστιν ἐμὰ πάντα ταῦτα, λέγει Κύριος· καὶ ἐπὶ τίνα ἐπιβλέψω, ἀλλʼ ἢ ἐπὶ τὸν ταπεινὸν καὶ ἡσύχιον, καὶ τρέμοντα τοὺς λόγους μου;

\vs{3}Ὁ δὲ ἄνομος ὁ θύων μοι μόσχον, ὡς ὁ ἀποκτέννων κύνα· ὁ δὲ ἀναφέρων σεμίδαλιν, ὡς αἷμα ὕειον· ὁ διδοὺς λίβανον εἰς μνημόσυνον, ὡς βλάσφημος.

Καὶ αὐτοὶ ἐξελέξαντο τὰς ὁδοὺς αὐτῶν, καὶ τὰ βδελύγματα αὐτῶν ἡ ψυχὴ αὐτῶν ἠθέλησε.
\vs{4}Καὶ ἐγὼ ἐκλέξομαι τὰ ἐμπαίγματα αὐτῶν, καὶ τὰς ἁμαρτίας ἀνταποδώσω αὐτοῖς· ὅτι ἐκάλεσα αὐτοὺς, καὶ οὐχ ὑπήκουσάν μου· ἐλάλησα καὶ οὐκ ἤκουσαν, καὶ ἐποίησαν τὸ πονηρὸν ἐναντίον ἐμοῦ, καὶ ἃ οὐκ ἐβουλόμην, ἐξελέξαντο.

\vs{5}Ἀκούσατε ῥήματα Κυρίου οἱ τρέμοντες τὸν λόγον αὐτοῦ· εἴπατε ἀδελφοὶ ὑμῶν τοῖς μισοῦσιν ὑμᾶς καὶ βδελυσσομένοις, ἵνα τὸ ὄνομα Κυρίου δοξασθῇ, καὶ ὀφθῇ ἐν τῇ εὐφροσύνῃ αὐτῶν, καὶ ἐκεῖνοι αἰσχυνθήσονται.

\vs{6}Φωνὴ κραυγῆς ἐκ πόλεως, φωνὴ ἐκ ναοῦ, φωνὴ Κυρίου ἀνταποδιδόντος ἀνταπόδοσιν τοῖς ἀντικειμένοις.
\vs{7}Πρὶν τὴν ὠδίνουσαν τεκεῖν, πρὶν ἐλθεῖν τὸν πόνον τῶν ὠδίνων, ἐξέφυγε καὶ ἔτεκεν ἄρσεν.
\vs{8}Τίς ἤκουσε τοιοῦτο, καὶ τίς ἑὼρακεν οὕτως; εἰ ὤδινε γῆ ἐν ἡμέρᾳ μιᾷ, ἢ καὶ ἐτέχθη ἔθνος εἰς ἅπαξ, ὅτι ὤδινε καὶ ἔτεκε Σιὼν τὰ παιδία αὐτῆς;
\vs{9}Ἐγὼ δὲ ἔδωκα τὴν προσδοκίαν ταύτην, καὶ οὐκ ἐμνήσθης μου, εἶπε Κύριος· οὐκ ἰδοὺ ἐγὼ γεννῶσαν καὶ στεῖραν ἐποίησα; εἶπεν ὁ Θεός σου.

\vs{10}Εὐφράνθητι Ἱερουσαλὴμ, καὶ πανηγυρίσατε ἐν αὐτῇ πάντες οἱ ἀγαπῶντες αὐτὴν, χάρητε ἅμα αὐτῇ χαρᾷ πάντες ὅσοι πενθεῖτε ἐπʼ αὐτῇ,
\vs{11}ἵνα θηλάσητε, καὶ ἐμπλησθῆτε ἀπὸ μαστοῦ παρακλήσεως αὐτῆς, ἵνα ἐκθηλάσαντες τρυφήσητε ἀπὸ εἰσόδου δόξης αὐτῆς.

\vs{12}Ὅτι τάδε λέγει Κύριος, ἰδοὺ ἐγὼ ἐκκλίνω εἰς αὐτοὺς ὡς ποταμὸς εἰρήνης, καὶ ὡς χειμάῤῥους ἐπικλύζων δόξαν ἐθνῶν· τὰ παιδία αὐτῶν ἐπʼ ὤμων ἀρθήσονται, καὶ ἐπὶ γονάτων παρακληθήσονται.
\vs{13}Ὡς εἴ τινα μήτηρ παρακαλέσει, οὕτω κᾀγὼ παρακαλέσω ὑμᾶς, καὶ ἐν Ἱερουσαλὴμ παρακληθήσεσθε.
\vs{14}Καὶ ὄψεσθε, καὶ χαρήσεται ἡ καρδία ὑμῶν, καὶ τὰ ὀστᾶ ὑμῶν ὡς βοτάνη ἀνατελεῖ· καὶ γνωσθήσεται ἡ χεὶρ Κυρίου τοῖς φοβουμένοις αὐτὸν, καὶ ἀπειλήσει τοῖς ἀπειθοῦσιν.

\vs{15}Ἰδοὺ γὰρ Κύριος ὡς πῦρ ἥξει, καὶ ὡς καταιγὶς τὰ ἅρματα αὐτοῦ, ἀποδοῦναι ἐν θυμῷ ἐκδίκησιν αὐτοῦ, καὶ ἀποσκορακισμὸν αὐτοῦ ἐν φλογὶ πυρός.
\vs{16}Ἐν γὰρ τῷ πυρὶ Κυρίου κριθήσεται πᾶσα ἡ γῆ, καὶ ἐν τῇ ῥομφαίᾳ αὐτοῦ πᾶσα σάρξ· πολλοὶ τραυματίαι ἔσονται ὑπὸ Κυρίου.

\vs{17}Οἱ ἁγνιζόμενοι καὶ καθαριζόμενοι εἰς τοὺς κήπους, καὶ ἐν τοῖς προθύροις ἔσθοντες κρέας ὕειον, καὶ τὰ βδελύγματα, καὶ τὸν μῦν, ἐπιτοαυτὸ ἀναλωθήσονται, εἶπε Κύριος.
\vs{18}Κᾀγὼ τὰ ἔργα αὐτῶν καὶ τὸν λογισμὸν αὐτῶν· ἔρχομαι συναγαγεῖν πάντα τὰ ἔθνη καὶ τὰς γλώσσας, καὶ ἥξουσι καὶ ὄψονται τὴν δόξαν μου.
\vs{19}Καὶ καταλείψω ἐπʼ αὐτῶν σημεῖον, καὶ ἐξαποστελῶ ἐξ αὐτῶν σεσωσμένους εἰς τὰ ἔθνη, εἰς Θαρσὶς, καὶ Φοὺδ, καὶ Λοὺδ, καὶ Μοσόχ, καὶ εἰς Θοβὲλ, καὶ εἰς τὴν Ἑλλάδα, καὶ εἰς τὰς νήσους τὰς πόῤῥω, οἳ οὐκ ἀκηκόασί μου τὸ ὄνομα, οὔτε ἑωράκασί μου τὴν δόξαν· καὶ ἀναγγελοῦσι τὴν δόξαν μου ἐν τοῖς ἔθνεσι,
\vs{20}καὶ ἄξουσι τοὺς ἀδελφοὺς ὑμῶν ἐκ πάντων τῶν ἐθνῶν δῶρον Κυρίῳ, μεθʼ ἵππων καὶ ἁρματων ἐν λαμπήναις ἡμιόνων μετὰ σκιαδίων εἰς τὴν ἁγίαν πόλιν Ἱερουσαλήμ, εἶπε Κύριος, ὡς ἀνενέγκαισαν οἱ υἱοὶ Ἰσραὴλ τὰς θυσίας αὐτῶν ἐμοὶ μετὰ ψαλμῶν εἰς τὸν οἶκον Κυρίου.
\vs{21}Καὶ ἀπʼ αὐτῶν λήμψομαι ἱερεῖς καὶ Λευίτας, εἶπε Κύριος.

\vs{22}Ὃν τρόπον γὰρ ὁ οὐρανὸς καινὸς καὶ ἡ γῆ καινὴ, ἃ ἐγὼ ποιῶ, μένει ἐνώπιον ἐμοῦ, λέγει Κύριος, οὕτω στήσεται τὸ σπέρμα ὑμῶν, καὶ τὸ ὄνομα ὑμῶν.
\vs{23}Καὶ ἔσται μὴν ἐκ μηνὸς, καὶ σάββατον ἐκ σαββάτου, ἥξει πᾶσα σὰρξ τοῦ προσκυνῆσαι ἐνώπιον ἐμοῦ ἐν Ἱερουσαλὴμ, εἶπεν Κύριος.
\vs{24}Καὶ ἐξελεύσονται καὶ ὄψονται τὰ κῶλα τῶν ἀνθρώπων τῶν παραβεβηκότων ἐν ἐμοί· ὁ γὰρ σκώληξ αὐτῶν οὐ τελευτήσει, καὶ τὸ πῦρ αὐτῶν οὐ σβεσθήσεται, καὶ ἔσονται εἰς ὅρασιν πάσῃ σαρκί.


\def\book{ΙΕΡΕΜΙΑΣ}
\biblebook{ΙΕΡΕΜΙΑΣ}


\lettrine[lines=2, loversize=0.2, nindent=0em, findent=.25em]{\textcolor{bookheadingcolor}{Τ}}{Ο} ῥῆμα τοῦ Θεοῦ ὃ ἐγένετο ἐπὶ Ἱερεμίαν τὸν του Χελκίου, ἐκ τῶν ἱερέων, ὃς κατῴκει ἐν Ἀναθὼθ ἐν γῇ Βενιαμὶν,
\vs{2}ὡς ἐγενήθη λόγος τοῦ Θεοῦ πρὸς αὐτὸν, ἐν ταῖς ἡμέραις Ἰωσία υἱοῦ Ἀμὼς βασιλέως Ἰούδα, ἔτους τρισκαιδεκάτου ἐν τῇ βασιλείᾳ αὐτοῦ.
\vs{3}Καὶ ἐγένετο ἐν ταῖς ἡμέραις Ἰωακεὶμ υἱοῦ Ἰωσία βασιλέως Ἰούδα, ἕως ἑνδεκάτου ἔτους τοῦ Σεδεκία υἱοῦ Ἰωσία βασιλέως Ἰούδα, ἕως τῆς αἰχμαλωσίας Ἱερουσαλὴμ ἐν τῷ πέμπτῳ μηνί.

\vs{4}Καὶ ἐγένετο λόγος Κυρίου πρὸς αὐτόν·
\vs{5}Πρὸ τοῦ με πλάσαι σε ἐν κοιλίᾳ, ἐπίσταμαί σε, καὶ πρὸ τοῦ σε ἐξελθεῖν ἐκ μήτρας, ἡγίακά σε, προφήτην εἰς ἔθνη τέθεικά σε.

\vs{6}Καὶ εἶπα, ὁ ὢν δέσποτα, Κύριε, ἰδοὺ οὐκ ἐπίσταμαι λαλεῖν, ὅτι νεώτερος ἐγώ εἰμι.
\vs{7}Καὶ εἶπε Κύριος πρὸς μὲ, μὴ λέγε, ὅτι νεώτερος ἐγώ εἰμι, ὅτι πρὸς πάντας οὓς ἐὰν ἐξαποστείλω σε, πορεύσῃ, καὶ κατὰ πάντα ὅσα ἐὰν ἐντείλωμαί σοι, λαλήσεις.
\vs{8}Μὴ φοβηθῇς ἀπὸ προσώπου αὐτῶν, ὅτι μετὰ σοῦ ἐγώ εἰμι τοῦ ἐξαιρεῖσθαί σε, λέγει Κύριος.
\vs{9}Καὶ ἐξέτεινε Κύριος τὴν χεῖρα αὐτοῦ πρὸς μὲ, καὶ ἥψατο τοῦ στόματός μου, καὶ εἶπε Κύριος πρὸς μὲ, ἰδοὺ δέδωκα τοὺς λόγους μου εἰς τὸ στόμα σου.

\vs{10}Ἰδοὺ καθέστακά σε σήμερον ἐπὶ ἔθνη καὶ ἐπὶ βασιλείας, ἐκριζοῦν, καὶ κατασκάπτειν, καὶ ἀπολύειν, καὶ ἀνοικοδομεῖν, καὶ καταφυτεύειν.

\vs{11}Καὶ ἐγένετο λόγος Κυρίου πρὸς μὲ, λέγων, τί σὺ ὁρᾷς; καὶ εἶπα, βακτηρίαν καρυΐνην.
\vs{12}Καὶ εἶπε Κύριος πρὸς μὲ, καλῶς ἑώρακας, διότι ἐγρήγορα ἐγὼ ἐπὶ τοὺς λόγους μου τοῦ ποιῆσαι αὐτούς.
\vs{13}Καὶ ἐγένετο λόγος Κυρίου ἐκ δευτέρου πρὸς μὲ, λέγων, τί σὺ ὁρᾷς; καὶ εἶπα, λέβητα ὑποκαιόμενον, καὶ τὸ πρόσωπον αὐτοῦ ἀπὸ προσώπου Βοῤῥᾶ.
\vs{14}Καὶ εἶπε Κύριος πρὸς μὲ, ἀπὸ προσώπου Βοῤῥᾶ ἐκκαυθήσεται τὰ κακὰ ἐπὶ πάντας τοὺς κατοικοῦντας τὴν γῆν.
\vs{15}Διότι ἰδοὺ ἐγὼ συγκαλῶ πάσας τὰς βασιλείας τῆς γῆς ἀπὸ Βοῤῥᾶ, λέγει Κύριος· καὶ ἥξουσι καὶ θήσουσιν ἕκαστος τὸν θρόνον αὐτοῦ ἐπὶ τὰ πρόθυρα τῶν πυλῶν Ἱερουσαλὴμ, καὶ ἐπὶ πάντα τὰ τείχη τὰ κύκλῳ αὐτῆς, καὶ ἐπὶ πάσας τὰς πόλεις Ἰούδα.
\vs{16}Καὶ λαλήσω πρὸς αὐτοὺς μετὰ κρίσεως, περὶ πάσης τῆς κακίας αὐτῶν, ὡς ἐγκατέλιπόν με, καὶ ἔθυσαν θεοῖς ἀλλοτρίοις, καὶ προσεκύνησαν τοῖς ἔργοις τῶν χειρῶν αὐτῶν.

\vs{17}Καὶ σὺ περίζωσαι τὴν ὀσφύν σου, καὶ ἀνάστηθι, καὶ εἰπὸ πάντα ὅσα ἄν ἐντείλωμαί σοι· μὴ φοβηθῇς ἀπὸ προσώπου αὐτῶν, μηδὲ πτοηθῇς ἐναντίον αὐτῶν, ὅτι μετὰ σοῦ εἰμι, τοῦ ἐξαιρεῖσθαί σε, λέγει Κύριος.
\vs{18}Ἰδοὺ τέθεικά σε ἐν τῇ σήμερον ἡμέρᾳ ὡς πόλιν ὀχυρὰν καὶ ὡς τεῖχος χαλκοῦν, ὀχυρὸν πᾶσι τοῖς βασιλεύσιν Ἰούδα, καὶ τοῖς ἄρχουσιν αὐτοῦ, καὶ τῷ λαῷ τῆς γῆς.
\vs{19}Καὶ πολεμήσουσί σε, καὶ οὐ μὴ δύνωνται πρὸς σὲ διότι μετὰ σοῦ ἐγώ εἰμι, τοῦ ἐξαιρεῖσθαί σε, εἶπε Κύριος.

\ch{2}
\vs{2}Καὶ εἶπε, τάδε λέγει Κύριος, ἐμνήσθην ἐλέους νεότητός σου, καὶ ἀγάπης τελειώσεώς σου, τοῦ ἐξακολουθῆσαί σε τῷ ἁγίῳ Ἰσραὴλ, λέγει Κύριος.
\vs{3}Ὁ ἅγιος Ἰσραὴλ τῷ Κυρίῳ, ἀρχὴ γεννημάτων αὐτοῦ· πάντες οἱ ἔσθοντες αὐτὸν πλημμελήσουσι, κακὰ ἥξει ἐπʼ αὐτοὺς, φησὶ Κύριος.

\vs{4}Ἀκούσατε λόγον Κυρίου οἶκος Ἰακὼβ, καὶ πᾶσα πατριὰ οἴκου Ἰσραήλ.
\vs{5}Τάδε λέγει Κύριος, τί εὕροσαν οἱ πατέρες ὑμῶν ἐν ἐμοὶ πλημμέλημα, ὅτι ἀπέστησαν μακρὰν ἀπʼ ἐμοῦ, καὶ ἐπορεύθησαν ὀπίσω τῶν ματαίων, καὶ ἐματαιώθησαν;
\vs{6}Καὶ οὐκ εἶπαν, ποῦ ἐστι Κύριος, ὁ ἀναγαγὼν ἡμᾶς ἐκ γῆς Αἰγύπτου, ὁ καθοδηγήσας ἡμᾶς ἐν τῇ ἐρήμῳ ἐν γῇ ἀπείρῳ καὶ ἀβάτῳ, ἐν γῇ ἀνύδρῳ καὶ ἀκάρπῳ, ἐν γῇ ᾗ οὐ διώδευσεν ἐν αὐτῇ ἀνὴρ εὐθὲν, καὶ οὐ κατῴκησεν ἄνθρωπος ἐκεῖ;
\vs{7}Καὶ ἤγαγον ὑμᾶς εἰς τὸν Κάρμηλον, τοῦ φαγεῖν ὑμᾶς τοὺς καρποὺς αὐτοῦ, καὶ τὰ ἀγαθὰ αὐτοῦ· καὶ εἰσήλθετε, καὶ ἐμιάνατε τὴν γῆν μου, καὶ τὴν κληρονομίαν μου ἔθεσθε εἰς βδέλυγμα.
\vs{8}Οἱ ἱερεῖς οὐκ εἶπαν, ποῦ ἐστι Κύριος; καὶ οἱ ἀντεχόμενοι τοῦ νόμου, οὐκ ἠπίσταντό με, καὶ οἱ ποιμένες ἠσέβουν εἰς ἐμὲ, καὶ οἱ προφῆται ἐπροφήτευον τῇ Βάαλ, καὶ ὀπίσω ἀνωφελοῦς ἐπορεύθησαν.

\vs{9}Διατοῦτο ἔτι κριθήσομαι πρὸς ὑμᾶς, καὶ πρὸς τοὺς υἱοὺς τῶν υἱῶν ὑμῶν κριθήσομαι.
\vs{10}Διότι ἔλθετε εἰς νήσους Χεττιεὶμ, καὶ ἴδετε, καὶ εἰς Κηδὰρ ἀποστείλατε, καὶ νοήσατε σφόδρα, καὶ ἴδετε εἰ γέγονε τοιαῦτα·
\vs{11}εἰ ἀλλάξωνται ἔθνη θεοὺς αὐτῶν, καὶ οὗτοι οὐκ εἰσὶ θεοί· ὁ δὲ λαός μου ἠλλάξατο τὴν δόξαν αὐτοῦ, ἐξ ἧς οὐκ ὠφεληθήσονται.
\vs{12}Ἐξέστη ὁ οὐρανὸς ἐπὶ τούτῳ, καὶ ἔφριξεν ἐπὶ πλεῖον σφόδρα, λέγει Κύριος.
\vs{13}Ὅτι δύο καὶ πονηρὰ ἐποίησεν ὁ λαός μου· ἐμὲ ἐγκατέλιπον πηγὴν ὕδατος ζωῆς, καὶ ὤρυξαν ἑαυτοῖς λάκκους συντετριμμένους, οἳ οὐ δυνήσονται ὕδωρ συνέχειν.

\vs{14}Μὴ δοῦλός ἐστιν Ἰσραὴλ, ἢ οἰκογενής ἐστι; διατί εἰς προνομὴν ἐγένετο;
\vs{15}Ἐπʼ αὐτὸν ὠρύοντο λέοντες, καὶ ἔδωκαν τὴν φωνὴν αὐτῶν, οἳ ἔταξαν τὴν γῆν αὐτοῦ εἰς ἔρημον, καὶ αἱ πόλεις αὐτοῦ κατεσκάφησαν, παρὰ τὸ μὴ κατοικεῖσθαι.
\vs{16}Καὶ υἱοὶ Μέμφεως καὶ Τάφνας ἔγνωσάν σε, καὶ κατέπαιζόν σου.
\vs{17}Οὐχὶ ταῦτα ἐποίησέ σοι τὸ καταλιπεῖν σε ἐμέ; λέγει Κύριος ὁ Θεός σου.

\vs{18}Καὶ νῦν τί σοι καὶ τῇ ὁδῷ Αἰγύπτου τοῦ πιεῖν ὕδωρ Γηῶν; καὶ τί σοι καὶ τῇ ὁδῷ Ἀσσυρίων τοῦ πιεῖν ὕδωρ ποταμῶν;
\vs{19}Παιδεύσει σε ἡ ἀποστασία σου, καὶ ἡ κακία σου ἐλέγξει σε· καὶ γνῶθι, καὶ ἴδε, ὅτι πικρόν σοι τὸ καταλιπεῖν σε ἐμὲ, λέγει Κύριος ὁ Θεός σου· καὶ οὐκ εὐδόκησα ἐπὶ σοὶ, λέγει Κύριος ὁ Θεός σου.

\vs{20}Ὅτι ἀπʼ αἰῶνος συνέτριψας τὸν ζυγόν σου, καὶ διέσπασας τοὺς δεσμούς σου, καὶ εἶπας, οὐ δουλεύσω σοι ἀλλὰ πορεύσομαι ἐπὶ πάντα βουνὸν ὑψηλὸν, καὶ ὑποκάτω παντὸς ξύλου κατασκίου, ἐκεῖ διαχυθήσομαι ἐν τῇ πορνείᾳ μου.
\vs{21}Ἐγὼ δὲ ἐφύτευσά σε ἄμπελον καρποφόρον πᾶσαν ἀληθινήν· πῶς ἐστράφης εἰς πικρίαν ἡ ἄμπελος ἡ ἀλλοτρία;
\vs{22}Ἐὰν ἀποπλύνῃ ἐν νίτρῳ, καὶ πληθυνῇς σεαυτῇ ποίαν, κεκηλίδωσαι ἐν ταῖς ἀδικίαις σου ἐναντίον ἐμοῦ, λέγει Κύριος.

\vs{23}Πῶς ἐρεῖς, οὐκ ἐμιάνθην, καὶ ὀπίσω τῆς Βάαλ οὐκ ἐπορεύθην; ἴδε τὰς ὁδούς σου ἐν τῷ πολυανδρίῳ, καὶ γνῶθι τί ἐποίησας· ὀψὲ φωνὴ αὐτῆς ὠλόλυξε·
\vs{24}Τὰς ὁδοὺς αὐτῆς ἐπλάτυνεν ἐφʼ ὕδατα ἐρήμου, ἐν ἐπιθυμίαις ψυχῆς αὐτῆς ἐπνευματοφορεῖτο, παρεδόθη, τίς ἐπιστρέψει αὐτήν; πάντες οἱ ζητοῦντες αὐτὴν οὐ κοπιάσουσιν, ἐν τῇ ταπεινώσει αὐτῆς εὑρήσουσιν αὐτήν.
\vs{25}Ἀπόστρεψον τὸν πόδα σου ἀπὸ ὁδοῦ τραχείας, καὶ τὸν φάρυγγά σου ἀπὸ δίψους· ἡ δὲ εἶπεν ἀνδριοῦμαι, ὅτι ἠγαπήκει ἀλλοτρίους, καὶ ὀπίσω αὐτῶν ἐπορεύετο.

\vs{26}Ὡς αἰσχύνη κλέπτου ὅταν ἁλῷ, οὕτως αἰσχυνθήσονται οἱ υἱοὶ Ἰσραὴλ, αὐτοὶ καὶ οἱ βασιλεῖς αὐτῶν, καὶ οἱ ἄρχοντες αὐτῶν, καὶ οἱ ἱερεῖς αὐτῶν, καὶ οἱ προφῆται αὐτῶν.
\vs{27}Τῷ ξύλῳ εἶπαν, ὅτι πατήρ μου εἶ σὺ, καὶ τῷ λίθῳ, σὺ ἐγέννησάς με· καὶ ἔστρεψαν ἐπʼ ἐμὲ νῶτα, καὶ οὐ πρόσωπα αὐτῶν· καὶ ἐν τῷ καιρῷ τῶν κακῶν αὐτῶν ἐροῦσιν, ἀνάστα καὶ σῶσον ἡμᾶς.
\vs{28}Καὶ ποῦ εἰσὶν οἱ θεοί σου, οὓς ἐποίησας σεαυτῷ; εἰ ἀναστήσονται καὶ σώσουσιν ἐν καιρῷ τῆς κακώσεώς σου; ὅτι κατʼ ἀριθμὸν τῶν πόλεών σου ἦσαν θεοί σου Ἰούδα, καὶ κατʼ ἀριθμὸν διόδων τῆς Ἱερουσαλὴμ ἔθυον τῇ Βάαλ.
\vs{29}Ἱνατί λαλεῖτε πρὸς μέ; πάντες ὑμεῖς ἠσεβήσατε, καὶ πάντες ὑμεῖς ἠνομήσατε εἰς ἐμὲ, λέγει Κύριος.
\vs{30}Μάτην ἐπάταξα τὰ τέκνα ὑμῶν, παιδείαν οὐκ ἐδέξασθε, μάχαιρα κατέφαγε τοὺς προφήτας ὑμῶν ὡς λέων ὀλοθρεύων, καὶ οὐκ ἐφοβήθητε.

\vs{31}Ἀκούσατε λόγον Κυρίου· τάδε λέγει Κύριος, μὴ ἔρημος ἐγενόμην τῷ Ἰσραὴλ ἢ γῆ κεχερσωμένη; διατί εἶπεν ὁ λαός μου, οὐ κυριευθησόμεθα, καὶ οὐχ ἥξομεν πρὸς σὲ ἔτι;
\vs{32}Μὴ ἐπιλήσεται νύμφη τὸν κόσμον αὐτῆς, καὶ παρθένος τὴν στηθοδεσμίδα αὐτῆς; ὁ δὲ λαός μου ἐπελάθετό μου ἡμέρας ὧν οὐκ ἔστιν ἀριθμός.
\vs{33}Τί ἔτι καλὸν ἐπιτηδεύσεις ἐν ταῖς ὁδοῖς σου, τοῦ ζητῆσαι ἀγάπησιν; οὐχ οὕτως· ἀλλὰ καὶ σὺ ἐπονηρεύσω τοῦ μιάναι τὰς ὁδούς σου,
\vs{34}καὶ ἐν ταῖς χερσί σου εὑρέθησαν αἵματα ψυχῶν ἀθώων· οὐκ ἐν διορύγμασιν εὗρον αὐτοὺς, ἀλλʼ ἐπὶ πάσῃ δρυΐ.
\vs{35}Καὶ εἶπας, ἀθῶός εἰμι, ἀλλὰ ἀποστραφήτω ὁ θυμὸς αὐτοῦ ἀπʼ ἐμοῦ.

Ἰδοὺ ἐγὼ κρίνομαι πρὸς σὲ, ἐν τῷ λέγειν σε, οὐχ ἥμαρτον·
\vs{36}Ὅτι κατεφρόνησας σφόδρα τοῦ δευτερῶσαι τὰς ὁδούς σου· καὶ ἀπὸ Αἰγύπτου καταισχυνθήσῃ, καθὼς κατῃσχύνθης ἀπὸ Ἀσσούρ·
\vs{37}ὅτι καὶ ἐντεῦθεν ἐξελεύσῃ, καὶ αἱ χεῖρές σου ἐπὶ τῆς κεφαλῆς σου· ὅτι ἀπώσατο Κύριος τὴν ἐλπίδα σου, καὶ οὐκ εὐοδωθήσῃ ἐν αὐτῇ.

\ch{3}
Ἐὰν ἐξαποστείλῃ ἀνὴρ τὴν γυναῖκα αὐτοῦ, καὶ ἀπέλθῃ ἀπʼ αὐτοῦ, καὶ γένηται ἀνδρὶ ἑτέρῳ, μὴ ἀνακάμπτουσα ἀνακάμψει πρὸς αὐτὸν ἔτι; οὐ μιαινομένη μιανθήσεται ἡ γυνὴ ἐκείνη; καὶ σὺ ἐξεπόρνευσας ἐν ποιμέσι πολλοῖς, καὶ ἀνέκαμπτες πρὸς μὲ, λέγει Κύριος.
\vs{2}Ἆρον τοὺς ὀφθαλμούς σου εἰς εὐθεῖαν, καὶ ἴδε, ποῦ οὐχὶ ἐξεφύρθης· ἐπὶ ταῖς ὁδοῖς ἐκάθισας αὐτοῖς ὡσεὶ κορώνη ἐρημουμένη, καὶ ἐμίανας τὴν γῆν ἐν ταῖς πορνείαις σου καὶ ἐν ταῖς κακίαις σου,
\vs{3}καὶ ἔσχες ποιμένας πολλοὺς εἰς πρόσκομμα σεαυτῇ· ὄψις πόρνης ἐγένετό σου, ἀπηναισχύντησας πρὸς πάντας.

\vs{4}Οὐχ ὡς οἶκόν με ἐκαλέσας, καὶ πατέρα καὶ ἀρχηγὸν τῆς παρθενίας σου;
\vs{5}Μὴ διαμενεῖ εἰς τὸν αἰῶνα, ἢ φυλαχθήσεται εἰς νῖκος; ἰδοὺ ἐλάλησας, καὶ ἐποίησας τὰ πονηρὰ ταῦτα, καὶ ἠδυνάσθης.

\vs{6}Καὶ εἶπε Κύριος πρὸς μὲ ἐν ταῖς ἡμέραις Ἰωσείου τοῦ βασιλέως, εἶδες ἃ ἐποίησέ μοι ἡ κατοικία τοῦ Ἰσραήλ; ἐπορεύθησαν ἐπὶ πᾶν ὄρος ὑψηλὸν, καὶ ὑποκάτω παντὸς ξύλου ἀλσώδους, καὶ ἐπόρνευσαν ἐκεῖ.
\vs{7}Καὶ εἶπα, μετὰ τὸ πορνεῦσαι αὐτὴν ταῦτα πάντα, πρὸς μὲ ἀνάστρεψον· καὶ οὐκ ἀνέστρεψε· καὶ εἶδε τὴν ἀσυνθεσίαν αὐτῆς ἡ ἀσύνθετος Ἰούδα.
\vs{8}Καὶ εἶδον, ὅτι περὶ πάντων ὧν κατελήφθη ἐν οἷς ἐμοιχᾶτο ἡ κατοικία Ἰσραὴλ, καὶ ἐξαπέστειλα αὐτὴν, καὶ ἔδωκα αὐτῇ βιβλίον ἀποστασίου εἰς τὰς χεῖρας αὐτῆς· καὶ οὐκ ἐφοβήθη ἡ ἀσύνθετος Ἰούδα, καὶ ἐπορεύθη, καὶ ἐπόρνευσε καὶ αὐτὴ,
\vs{9}και ἐγένετο εἰς οὐθὲν ἡ πορνεία αὐτῆς, καὶ ἐμοίχευσε τὸ ξύλον καὶ τὸν λίθον.
\vs{10}Καὶ ἐν πᾶσι τούτοις οὐκ ἐπεστράφη πρὸς μὲ ἡ ἀσύνθετος Ἰούδα ἐξ ὅλης τῆς καρδίας αὐτῆς, ἀλλʼ ἐπὶ ψεύδει.

\vs{11}Καὶ εἶπε Κύριος πρὸς μὲ, ἐδικαίωσε τὴν ψυχὴν αὐτοῦ Ἰσραὴλ ἀπὸ τῆς ἀσυνθέτου Ἰούδα.
\vs{12}Πορεύου καὶ ἀνάγνωθι τοὺς λόγους τούτους πρὸς Βοῤῥᾶν, καὶ ἐρεῖς, ἐπιστράφηθι πρὸς μὲ ἡ κατοικία τοῦ Ἰσραὴλ, λέγει Κύριος· καὶ μὴ στηριῶ τὸ πρόσωπόν μου ἐφʼ ὑμᾶς, ὅτι ἐλεήμων ἐγώ εἰμι, λέγει Κύριος, καὶ οὐ μηνιῶ ὑμῖν εἰς τὸν αἰῶνα.
\vs{13}Πλὴν, γνῶθι τὴν ἀδικίαν σου, ὅτι εἰς Κύριον τὸν Θεόν σου ἠσέβησας, καὶ διέχεας τὰς ὁδούς σου εἰς ἀλλοτρίους ὑποκάτω παντὸς ξύλου ἀλσώδους, τῆς δὲ φωνῆς μου οὐχ ὑπήκουσας, λέγει Κύριος.
\vs{14}Ἐπιστράφητε υἱοὶ ἀφεστηκότες, λέγει Κύριος, διότι ἐγὼ κατακυριεύσω ὑμῶν, καὶ λήψομαι ὑμᾶς ἕνα ἐκ πόλεως καὶ δύο ἐκ πατριᾶς, καὶ εἰσάξω ὑμᾶς εἰς Σιὼν,
\vs{15}καὶ δώσω ὑμῖν ποιμένας κατὰ τὴν καρδίαν μου, καὶ ποιμανοῦσιν ὑμᾶς ποιμαίνοντες μετʼ ἐπιστήμης.

\vs{16}Καὶ ἔσται ἐὰν πληθυνθῆτε, καὶ αὐξηθῆτε ἐπὶ τῆς γῆς, λέγει Κύριος, ἐν ταῖς ἡμέραις ἐκείναις οὐκ ἐροῦσιν ἔτι, κιβωτὸς διαθήκης ἁγίου Ἰσραὴλ, οὐκ ἀναβήσεται ἐπὶ καρδίαν, οὐκ ὀνομασθήσεται, οὐδὲ ἐπισκεφθήσεται, καὶ οὐ ποιηθήσεται ἔτι.
\vs{17}Ἐν ταῖς ἡμέραις ἐκείναις καὶ ἐν τῷ καιρῷ ἐκείνῳ καλέσουσι τὴν Ἱερουσαλὴμ, Θρόνον Κυρίου· καὶ συναχθήσονται πάντα τὰ ἔθνη εἰς αὐτὴν, καὶ οὐ πορεύσονται ἔτι ὀπίσω τῶν ἐνθυμημάτων τῆς καρδίας αὐτῶν τῆς πονηρᾶς.

\vs{18}Ἐν ταῖς ἡμέραις ἐκείναις συνελεύσονται ὁ οἶκος Ἰούδα ἐπὶ τὸν οἶκον τοῦ Ἰσραὴλ, καὶ ἥξουσιν ἐπιτοαυτὸ ἀπὸ γῆς Βοῤῥᾶ, καὶ ἀπὸ πασῶν τῶν χωρῶν ἐπὶ τὴν γῆν, ἣν κατεκληρονόμησα τοὺς πατέρας αὐτῶν.
\vs{19}Καὶ ἐγὼ εἶπα, γένοιτο Κύριε· ὅτι τάξω σε εἰς τέκνα, καὶ δώσω σοι γῆν ἐκλεκτὴν, κληρονομίαν Θεοῦ παντοκράτορος ἐθνῶν· καὶ εἶπα, πατέρα καλέσετέ με, καὶ ἀπʼ ἐμοῦ οὐκ ἀποστραφήσεσθε.
\vs{20}Πλὴν ὡς ἀθετεῖ γυνὴ εἰς τὸν συνόντα αὐτῇ, οὕτως ἠθέτησεν εἰς ἐμὲ ὁ οἶκος Ἰσραὴλ, λέγει Κύριος.

\vs{21}Φωνὴ ἐκ χειλέων ἠκούσθη κλαυθμοῦ καὶ δεήσεως υἱῶν Ἰσραὴλ, ὅτι ἠδίκησαν ἐν ταῖς ὁδοῖς αὐτῶν, ἐπελάθοντο Θεοῦ ἁγίου αὐτῶν.
\vs{22}Ἐπιστράφητε υἱοὶ ἐπιστρέφοντες, καὶ ἰάσομαι τὰ συντρίμματα ὑμῶν.

Ἰδοὺ δοῦλοι ἡμεῖς ἐσόμεθά σοι· ὅτι σὺ Κύριος ὁ Θεὸς ἡμῶν εἶ.
\vs{23}Ὄντως εἰς ψεῦδος ἦσαν οἱ βουνοὶ, καὶ ἡ δύναμις τῶν ὀρέων, πλὴν διὰ Κυρίου Θεοῦ ἡμῶν ἡ σωτηρία τοῦ Ἰσραήλ.
\vs{24}Ἡ δὲ αἰσχύνη κατηνάλωσε τοὺς μόχθους τῶν πατέρων ἡμῶν, ἀπὸ νεότητος ἡμῶν, τὰ πρόβατα αὐτῶν καὶ τοὺς μόσχους αὐτῶν καὶ τοὺς υἱοὺς αὐτῶν καὶ τὰς θυγατέρας αὐτῶν.
\vs{25}Ἐκοιμήθημεν ἐν τῇ αἰσχύνῃ ἡμῶν, καὶ ἐπεκάλυψεν ἡμᾶς ἡ ἀτιμία ἡμῶν, διότι ἔναντι τοῦ Θεοῦ ἡμῶν ἡμάρτομεν ἡμεῖς, καὶ οἱ πατέρες ἡμῶν, ἀπὸ νεότητος ἡμῶν ἕως τῆς ἡμέρας ταύτης· καὶ οὐχ ὑπηκούσαμεν τῆς φωνῆς Κυρίου τοῦ Θεοῦ ἡμῶν.

\ch{4}
Ἐὰν ἐπιστραφῇ Ἰσραὴλ, λέγει Κύριος, πρὸς μὲ, ἐπιστραφήσεται· καὶ ἐὰν περιέλῃ τὰ βδελύγματα αὐτοῦ ἐκ στόματος αὐτοῦ, καὶ ἀπὸ τοῦ προσώπου μου εὐλαβηθῇ,
\vs{2}καὶ ὀμόσῃ, ζῇ Κύριος, μετὰ ἀληθείας ἐν κρίσει καὶ ἐν δικαιοσύνῃ, καὶ εὐλογήσουσιν ἐν αὐτῷ ἔθνη, καὶ ἐν αὐτῷ αἰνέσουσι τῷ Θεῷ ἐν Ἰερουσαλήμ.

\vs{3}Ὅτι τάδε λέγει Κύριος τοῖς ἀνδράσιν Ἰούδα, καὶ τοῖς κατοικοῦσιν Ἱερουσαλὴμ, νεώσατε ἑαυτοῖς νεώματα, καὶ μὴ σπείρητε ἐπʼ ἀκάνθαις.
\vs{4}Περιτμήθητε τῷ Θεῷ ὑμῶν, καὶ περιτέμεσθε τὴν σκληροκαρδίαν ὑμῶν ἄνδρες Ἰούδα, καὶ οἱ κατοικοῦντες Ἱερουσαλὴμ, μὴ ἐξέλθῃ ὡς πῦρ ὁ θυμός μου, καὶ ἐκκαυθήσεται, καὶ οὐκ ἔσται ὁ σβέσων, ἀπὸ προσώπου πονηρίας ἐπιτηδευμάτων ὑμῶν.

\vs{5}Ἀναγγείλατε ἐν τῷ Ἰούδᾳ, καὶ ἀκουσθήτω ἐν Ἱερουσαλήμ· εἴπατε, σημάνατε ἐπὶ τῆς γῆς σάλπιγγι, κεκράξατε μέγα· εἴπατε, συνάχθητε, καὶ εἰσέλθωμεν εἰς τὰς πόλεις τὰς τειχήρεις.
\vs{6}Ἀναλαβόντες φεύγετε εἰς Σιών· σπεύσατε, μὴ στῆτε, ὅτι κακὰ ἐγὼ ἐπάγω ἀπὸ Βοῤῥα, καὶ συντριβὴν μεγάλην.
\vs{7}Ἀνέβη λέων ἐκ τῆς μάνδρας αὐτοῦ, ἐξολοθρεύων ἔθνη ἐξῇρε, καὶ ἐξῆλθεν ἐκ τοῦ τόπου αὐτοῦ, τοῦ θεῖναι τὴν γῆν εἰς ἐρήμωσιν· καὶ αἱ πόλεις καθαιρεθήσονται, παρὰ τὸ μὴ κατοικεῖσθαι αὐτάς.
\vs{8}Ἐπὶ τούτοις περιζώσασθε σάκκους, καὶ κόπτεσθε, καὶ ἀλαλάξατε, διότι οὐκ ἀπεστράφη ὁ θυμὸς Κυρίου ἀφʼ ὑμῶν.
\vs{9}Καὶ ἔσται ἐν ἐκείνῃ τῇ ἡμέρᾳ, λέγει Κύριος, ἀπολεῖται ἡ καρδία τοῦ βασιλέως, καὶ ἡ καρδία τῶν ἀρχόντων, καὶ οἱ ἱερεῖς ἐκστήσονται, καὶ οἱ προφῆται θαυμάσονται.

\vs{10}Καὶ εἶπα, ὦ δέσποτα Κύριε, ἆρά γε ἀπατῶν ἠπάτησας τὸν λαὸν τοῦτον καὶ τὴν Ἱερουσαλὴμ, λέγων, εἰρήνη ἔσται, καὶ ἰδοὺ ἥψατο ἡ μάχαιρα ἕως τῆς ψυχῆς αὐτῶν.

\vs{11}Ἐν τῷ καιρῷ ἐκείνῳ ἐροῦσι τῷ λαῷ τούτῳ καὶ τῇ Ἱερουσαλὴμ, πνεῦμα πλανήσεως ἐν τῇ ἐρήμῳ, ὁδὸς τῆς θυγατρὸς τοῦ λαοῦ μου, οὐκ εἰς καθαρὸν, οὐδʼ εἰς ἅγιον.
\vs{12}Πνεῦμα πληρώσεως ἥξει μοι· νῦν δὲ ἐγὼ λαλῶ κρίματά μου πρὸς αὐτούς.
\vs{13}Ἰδοὺ ὡς νεφέλη ἀναβήσεται, καὶ ὡς καταιγὶς τὰ ἄρματα αὐτοῦ, κουφότεροι ἀετῶν οἱ ἵπποι αὐτοῦ· οὐαὶ ἡμῖν, ὅτι ταλαιπωροῦμεν.

\vs{14}Ἀπόπλυνε ἀπὸ κακίας τὴν καρδίαν σου Ἱερουσαλὴμ, ἵνα σωθῇς· ἕως πότε ὑπάρχουσιν ἐν σοὶ διαλογισμοὶ πόνων σου;
\vs{15}Διότι φωνὴ ἀγγέλλοντος ἐκ Δὰν ἥξει, καὶ ἀκουσθήσεται πόνος ἐξ ὄρους Ἐφραίμ.
\vs{16}Ἀναμνήσατε ἔθνη, ἰδοὺ ἥκασιν· ἀναγγείλατε ἐν Ἱερουσαλὴμ, συστροφαὶ ἔρχονται ἐκ γῆς μακρόθεν, καὶ ἔδωκαν ἐπὶ τὰς πόλεις Ἰούδα φωνὴν αὐτῶν.
\vs{17}Ὡς φυλάσσοντες ἀγρὸν, ἐγένοντο ἐπʼ αὐτὴν κύκλῳ, ὅτι ἐμοῦ ἠμέλησας, λέγει Κύριος.
\vs{18}Αἱ ὁδοί σου καὶ τὰ ἐπιτηδεύματά σου ἐποίησαν ταῦτά σοι· αὕτη ἡ κακία σου, ὅτι πικρὰ, ὅτι ἥψατο ἕως τῆς καρδίας σου.

\vs{19}Τὴν κοιλίαν μου, τὴν κοιλίαν μου ἀλγῶ, καὶ τὰ αἰσθητήρια τῆς καρδίας μου, μαιμάσσει ἡ ψυχή μου, σπαράσσεται ἡ καρδία μου· οὐ σιωπήσομαι, ὅτι φωνὴν σάλπιγγος ἤκουσεν ἡ ψυχή μου, κραυγὴν πολέμου
\vs{20}καὶ ταλαιπωρίας συντριμμὸν ἐπικαλεῖται, ὅτι τεταλαιπώρηκε πᾶσα ἡ γῆ, ἄφνω τεταλαιπώρηκεν ἡ σκηνὴ, διεσπάσθησαν αἱ δέῤῥεις μου.
\vs{21}Ἕως πότε ὄψομαι φεύγοντας ἀκούων φωνὴν σαλπίγγων;

\vs{22}Διότι οἱ ἡγούμενοι τοῦ λαοῦ μου ἐμὲ οὐκ ᾔδεισαν· υἱοὶ ἄφρονές εἰσι καὶ οὐ συνετοὶ, σοφοί εἰσι τοῦ κακοποιῆσαι, τὸ δὲ καλῶς ποιῆσαι οὐκ ἐπέγνωσαν.

\vs{23}Ἐπέβλεψα ἐπὶ τὴν γῆν, καὶ ἰδοὺ οὐθὲν, καὶ εἰς τὸν οὐρανὸν, καὶ οὐκ ἦν τὰ φῶτα αὐτοῦ.
\vs{24}Εἶδον τὰ ὄρη, καὶ ἦν τρέμοντα, καὶ πάντας τοὺς βουνοὺς ταρασσομένους.
\vs{25}Ἐπέβλεψα, καὶ ἰδοὺ οὐκ ἦν ἄνθρωπος, καὶ πάντα τὰ πετεινὰ τοῦ οὐρανοῦ ἐπτοεῖτο.
\vs{26}Εἶδον, καὶ ἰδοὺ ὁ Κάρμηλος ἔρημος, καὶ πᾶσαι αἱ πόλεις ἐμπεπυρισμέναι πυρὶ ἀπὸ προσώπου Κυρίου, καὶ ἀπὸ προσώπου ὀργῆς θυμοῦ αὐτοῦ ἠφανίσθησαν.

\vs{27}Τάδε λέγει Κύριος, ἔρημος ἔσται πᾶσα ἡ γῆ, συντέλειαν δὲ οὐ μὴ ποιήσω.
\vs{28}Ἐπὶ τούτοις πενθείτω ἡ γῆ, καὶ συσκοτασάτω ὁ οὐρανὸς ἄνωθεν· διότι ἐλάλησα, καὶ οὐ μετανοήσω, ὥρμησα, καὶ οὐκ ἀποστρέψω ἀπʼ αὐτῆς.
\vs{29}Ἀπὸ φωνῆς ἱππέως, καὶ ἐντεταμένου τόξου ἀνεχώρησε πᾶσα ἡ χώρα· εἰσέδυσαν εἴς τὰ σπήλαια, καὶ εἰς τὰ ἄλση ἐκρύβησαν, καὶ ἐπὶ τὰς πέτρας ἀνέβησαν· πᾶσα πόλις ἐγκατελείφθη, οὐ κατῴκει ἐν αὐταῖς ἄνθρωπος.
\vs{30}Καὶ σὺ τί ποιήσεις; ἐὰν περιβάλῃ κόκκινον, καὶ κοσμήσῃ κόσμῳ χρυσῷ· ἐὰν ἐγχρίσῃ στίβι τοὺς ὀφθαλμούς σου, εἰς μάταιον ὡραϊσμός σου· ἀπώσαντό σε οἱ ἐρασταί σου, τὴν ψυχήν σου ζητοῦσιν.

\vs{31}Ὅτι φωνὴν ὡς ὠδινούσης ἤκουσα τοῦ στεναγμοῦ σου, ὡς πρωτοτοκούσης· φωνὴ θυγατρὸς Σιὼν ἐκλυθήσεται, καὶ παρήσει τὰς χεῖρας αὐτῆς· οἴμοι ἐγὼ, ὅτι ἐκλείπει ἡ ψυχή μου ἐπὶ τοῖς ἀνῃρημένοις.

\ch{5}
Περιδράμετε ἐν ταῖς ὁδοῖς Ἱερουσαλὴμ, καὶ ἴδετε, καὶ γνῶτε, καὶ ζητήσατε ἐν ταῖς πλατείαις αὐτῆς, ἐὰν εὕρητε, εἰ ἔστι ποιῶν κρίμα, καὶ ζητῶν πίστιν· καὶ ἵλεως ἔσομαι αὐτοῖς, λέγει Κύριος.
\vs{2}Ζῇ Κύριος, λέγουσι, διατοῦτο οὐκ ἐν ψεύδεσιν ὀμνύουσι;
\vs{3}Κύριε, οἱ ὀφθαλμοί σου εἰς πίστιν· ἐμαστίγωσας αὐτοὺς, καὶ οὐκ ἐπόνεσαν, συνετέλεσας αὐτοὺς, καὶ οὐκ ἠθέλησαν δέξασθαι παιδείαν· ἐστερέωσαν τὰ πρόσωπα αὐτῶν ὑπὲρ πέτραν, καὶ οὐκ ἠθέλησαν ἐπιστραφῆναι.
\vs{4}Καὶ ἐγὼ εἶπα, ἴσως πτωχοί εἰσι, διότι οὐκ ἐδυνάσθησαν, ὅτι οὐκ ἔγνωσαν ὁδὸν Κυρίου, καὶ κρίσιν Θεοῦ.
\vs{5}Πορεύσομαι πρὸς τοὺς ἁδροὺς, καὶ λαλήσω αὐτοῖς, ὅτι αὐτοὶ ἐπέγνωσαν ὁδὸν Κυρίου, καὶ κρίσιν Θεοῦ· καὶ ἰδοὺ ὁμοθυμαδὸν συνέτριψαν ζυγὸν, διέῤῥηξαν δεσμούς.

\vs{6}Διατοῦτο ἔπαισεν αὐτοὺς λέων ἐκ τοῦ δρυμοῦ, καὶ λύκος ἕως τῶν οἰκιῶν ὠλόθρευσεν αὐτοὺς, καὶ πάρδαλις ἐγρηγόρησεν ἐπὶ τὰς πόλεις αὐτῶν· πάντες οἱ ἐκπορευόμενοι ἀπʼ αὐτῶν θηρευθήσονται, ὅτι ἐπλήθυναν ἀσεβείας αὐτῶν, ἴσχυσαν ἐν ταῖς ἀποστροφαῖς αὐτῶν.
\vs{7}Ποίᾳ τούτων ἵλεως γένωμαί σοι; οἱ υἱοί σου ἐγκατέλιπόν με, καὶ ὤμνυον ἐν τοῖς οὐκ οὖσι θεοῖς· καὶ ἐχόρτασα αὐτοὺς, καὶ ἐμοιχῶντο, καὶ ἐν οἴκοις πορνῶν κατέλυον.
\vs{8}Ἵπποι θηλυμανεῖς ἐγενήθησαν, ἕκαστος ἐπὶ τὴν γυναῖκα τοῦ πλησίον αὐτοῦ ἐχρεμέτιζον.
\vs{9}Μὴ ἐπὶ τούτοις οὐκ ἐπισκέψομαι; λέγει Κύριος· ἢ ἐν ἔθνει τοιούτῳ οὐκ ἐκδικήσει ἡ ψυχή μου;

\vs{10}Ἀνάβητε ἐπὶ τοὺς προμαχῶνας αὐτῆς, καὶ κατασκάψατε, συντέλειαν δὲ μὴ ποιήσητε· ὑπολίπεσθε τὰ ὑποστηρίγματα αὐτῆς, ὅτι τοῦ Κυρίου εἰσίν.
\vs{11}Ὅτι ἀθετῶν ἠθέτησεν εἰς ἐμὲ, λέγει Κύριος, οἶκος Ἰσραὴλ, καὶ οἶκος Ἰούδα
\vs{12}ἐψεύσατο τῷ Κυρίῳ αὐτῶν, καὶ εἶπαν, οὐκ ἔστι ταῦτα· οὐχ ἥξει ἐφʼ ἡμᾶς κακὰ, καὶ μάχαιραν καὶ λιμὸν οὐκ ὀψόμεθα.
\vs{13}Οἱ προφῆται ἡμῶν ἦσαν εἰς ἄνεμον, καὶ λόγος Κυρίου οὐχ ὑπῆρχεν ἐν αὐτοῖς.

\vs{14}Διατοῦτο τάδε λέγει Κύριος παντοκράτωρ, ἀνθʼ ὧν ἐλαλήσατε τὸ ῥῆμα τοῦτο, ἰδοὺ ἐγὼ δέδωκα τοὺς λόγους μου εἰς τὸ στόμα σου πῦρ, καὶ τὸν λαὸν τοῦτον ξύλα, καὶ καταφάγεται αὐτούς.
\vs{15}Ἰδοὺ ἐγὼ ἐπάγω ἐφʼ ὑμᾶς ἔθνος πόῤῥωθεν, οἶκος Ἰσραὴλ, λέγει Κύριος, ἔθνος οὗ οὐκ ἀκούσει τῆς φωνῆς τῆς γλώσσης αὐτοῦ.
\vs{16}Πάντες ἰσχυροὶ, καὶ κατέδονται τὸν θερισμὸν ὑμῶν,
\vs{17}καὶ τοὺς ἄρτους ὑμῶν, καὶ κατέδονται τοὺς υἱοὺς ὑμῶν, καὶ τὰς θυγατέρας ὑμῶν, καὶ κατέδονται τὰ πρόβατα ὑμῶν, καὶ τοὺς μόσχους ὑμῶν, καὶ κατέδονται τοὺς ἀμπελῶνας ὑμῶν, καὶ τοὺς συκῶνας ὑμῶν, καὶ τοὺς ἐλαιῶνας ὑμῶν· καὶ ἀλοήσουσι τὰς πόλεις τὰς ὀχυρὰς ὑμῶν, ἐφʼ αἷς ὑμεῖς πεποίθατε ἐπʼ αὐταῖς, ἐν ῥομφαίᾳ.
\vs{18}Καὶ ἔσται ἐν ταῖς ἡμέραις ἐκείναις, λέγει Κύριος ὁ Θεός σου, οὐ μὴ ποιήσω ὑμᾶς εἰς συντέλειαν.

\vs{19}Καὶ ἔσται ὅταν εἴπητε, τίνος ἕνεκεν ἐποίησε Κύριος ὁ Θεὸς ἡμῶν ἡμῖν πάντα ταῦτα; καὶ ἐρεῖς αὐτοῖς, ἀνθʼ ὧν ἐδουλεύσατε θεοῖς ἀλλοτρίοις ἐν τῇ γῇ ὑμῶν, οὕτως δουλεύσετε ἀλλοτρίοις ἐν τῇ γῇ οὐχ ὑμῶν.

\vs{20}Ἀναγγείλατε ταῦτα εἰς τὸν οἶκον Ἰακὼβ, καὶ ἀκουσθήτω ἐν τῷ οἴκῳ Ἰούδα.
\vs{21}Ἀκούσατε δὴ ταῦτα λαὸς μωρὸς καὶ ἀκάρδιος, ὀφθαλμοὶ αὐτοῖς καὶ οὐ βλέπουσιν, ὦτα αὐτοῖς καὶ οὐκ ἀκούουσι.
\vs{22}Μὴ ἐμὲ οὐ φοβηθήσεσθε; λέγει Κύριος· ἢ ἀπὸ προσώπου μου οὐκ εὐλαβηθήσεσθε; τὸν τάξαντα ἄμμον ὅριον τῇ θαλάσσῃ, πρόσταγμα αἰώνιον, καὶ οὐχ ὑπερβήσεται αὐτὸ, καὶ ταραχθήσεται, καὶ οὐ δυνήσεται· καὶ ἠχήσουσι τὰ κύματα αὐτῆς, καὶ οὐχ ὑπερβήσεται αὐτό.

\vs{23}Τῷ δὲ λαῷ τούτῳ ἐγενήθη καρδία ἀνήκοος καὶ ἀπειθὴς, καὶ ἐξέκλιναν καὶ ἀπήλθοσαν,
\vs{24}καὶ οὐκ εἶπον ἐν τῇ καρδίᾳ αὐτῶν, φοβηθῶμεν δὴ Κύριον τὸν Θεὸν ἡμῶν, τὸν διδόντα ἡμῖν ὑετὸν πρώϊμον καὶ ὄψιμον, κατὰ καιρὸν πληρώσεως προστάγματος θερισμοῦ, καὶ ἐφύλαξεν ἡμῖν.
\vs{25}Αἱ ἀνομίαι ὑμῶν ἐξέκλιναν ταῦτα, καὶ αἱ ἁμαρτίαι ὑμῶν ἀπέστησαν τὰ ἀγαθὰ ἀφʼ ὑμῶν.
\vs{26}Ὅτι εὑρέθησαν ἐν τῷ λαῷ μου ἀσεβεῖς, καὶ παγίδας ἔστησαν τοῦ διαφθεῖραι ἄνδρας, καὶ συνελαμβάνοσαν.

\vs{27}Ὡς παγὶς ἐφεσταμένη πλήρης πετεινῶν, οὕτως οἱ οἶκοι αὐτῶν πλήρεις δόλου· διατοῦτο ἐμεγαλύνθησαν, καὶ ἐπλούτησαν,
\vs{28}καὶ παρέβησαν κρίσιν, οὐκ ἔκριναν κρίσιν ὀρφανοῦ, καὶ κρίσιν χήρας οὐκ ἐκρίνοσαν.
\vs{29}Μὴ ἐπὶ τούτοις οὐκ ἐπισκέψομαι; λέγει Κύριος· ἢ ἐν ἔθνει τῷ τοιούτῳ οὐκ ἐκδικήσει ἡ ψυχή μου;

\vs{30}Ἔκστασις καὶ φρικτὰ ἐγενήθη ἐπὶ τῆς γῆς.
\vs{31}Οἱ προφῆται προφητεύουσιν ἄδικα, καὶ οἱ ἱερεῖς ἐπεκρότησαν ταῖς χερσὶν αὐτῶν, καὶ ὁ λαός μου ἠγάπησεν οὕτως· καὶ τί ποιήσετε εἰς τὰ μετὰ ταῦτα;

\ch{6}
Ἐνισχύσατε υἱοὶ Βενιαμὶν ἐκ μέσου τῆς Ἱερουσαλὴμ, καὶ ἐν Θεκουὲ σημάνατε σάλπιγγι, καὶ ὑπὲρ Βαιθαχαρμὰ ἄρατε σημεῖον, ὅτι κακὰ ἐκκέκυφεν ἀπὸ Βοῤῥᾶ, καὶ συντριβὴ μεγάλη γίνεται·
\vs{2}καὶ ἀφαιρεθήσεται τὸ ὕψος θύγατερ Σιών.
\vs{3}Εἰς αὐτὴν ἥξουσι ποιμένες, καὶ τὰ ποίμνια αὐτῶν· καὶ πήξουσιν ἐπʼ αὐτὴν σκηνὰς κύκλῳ, καὶ ποιμανοῦσιν ἕκαστος τῇ χειρὶ αὐτοῦ.

\vs{4}Παρασκευάσασθε ἐπʼ αὐτὴν εἰς πόλεμον, ἀνάστητε, καὶ ἀναβῶμεν ἐπʼ αὐτὴν μεσημβρίας· οὐαὶ ἡμῖν, ὅτι κέκλικεν ἡ ἡμέρα, ὅτι ἐκλείπουσιν αἱ σκιαὶ τῆς ἡμέρας.
\vs{5}Ἀνάστητε, καὶ ἀναβῶμεν ἐπʼ αὐτὴν νυκτὶ, καὶ διαφθείρωμεν τὰ θεμέλια αὐτῆς.

\vs{6}Ὅτι τάδε λέγει Κύριος, ἔκκοψον τὰ ξύλα αὐτῆς, ἔκχεον ἐπὶ Ἱερουσαλὴμ δύναμιν· ὦ πόλις ψευδὴς, ὅλη καταδυναστεία ἐν αὐτῇ.
\vs{7}Ὡς ψύχει λάκκος ὕδωρ, οὕτω ψύχει κακία αὐτῆς· ἀσέβεια καὶ ταλαιπωρία ἀκουσθήσεται ἐν αὐτῇ ἐπὶ πρόσωπον αὐτῆς διαπαντός.

\vs{8}Πόνῳ καὶ μάστιγι παιδευθήσῃ Ἱερουσαλὴμ, μὴ ἀποστῇ ἡ ψυχή μου ἀπὸ σοῦ, μὴ ποιήσω σε ἄβατον γῆν, ἥτις οὐ κατοικισθῇ.

\vs{9}Ὅτι τάδε λέγει Κύριος, καλαμᾶσθε, καλαμᾶσθε ὡς ἄμπελον τὰ κατάλοιπα τοῦ Ἰσραὴλ, ἐπιστρέψατε ὡς ὁ τρυγῶν ἐπὶ τὸν κάρταλλον αὐτοῦ.
\vs{10}Πρὸς τίνα λαλήσω, καὶ διαμαρτύρωμαι, καὶ εἰσακούσεται; ἰδοὺ ἀπερίτμητα τὰ ὦτα αὐτῶν, καὶ οὐ δυνήσονται ἀκούειν· ἰδοὺ τὸ ῥῆμα Κυρίου ἐγένετο αὐτοῖς εἰς ὀνειδισμὸν, οὐ μὴ βουληθῶσιν αὐτό.

\vs{11}Καὶ τὸν θυμόν μου ἔπλησα, καὶ ἐπέσχον, καὶ οὐ συνετέλεσα αὐτούς· ἐκχεῶ ἐπὶ νήπια ἔξωθεν, καὶ ἐπὶ συναγωγὴν νεανίσκων ἅμα, ὅτι ἀνὴρ καὶ γυνὴ συλληφθήσονται, πρεσβύτερος μετὰ πλήρους ἡμερῶν.
\vs{12}Καὶ μεταστραφήσονται αἱ οἰκίαι αὐτῶν εἰς ἑτέρους, ἀγροὶ καὶ αἱ γυναῖκες αὐτῶν ἐπιτοαυτό· ὅτι ἐκτενῶ τὴν χεῖρά μου ἐπὶ τοὺς κατοικοῦντας τὴν γῆν ταύτην, λέγει Κύριος.

\vs{13}Ὅτι ἀπὸ μικροῦ αὐτῶν καὶ ἕως μεγάλου πάντες συνετελέσαντο ἄνομα, ἀπὸ ἱερέως καὶ ἕως ψευδοπροφήτου πάντες ἐποίησαν ψευδῆ.
\vs{14}Καὶ ἰῶντο σύντριμμα τοῦ λαοῦ μου, ἐξουθενοῦντες καὶ λέγοντες, εἰρήνη εἰρήνη· καὶ ποῦ ἐστιν εἰρήνη;
\vs{15}Κατῃσχύνθησαν, ὅτι ἐξελίποσαν· καὶ οὐδʼ ὡς καταισχυνόμενοι κατῃσχύνθησαν, καὶ τὴν ἀτιμίαν αὐτῶν οὐκ ἔγνωσαν· διατοῦτο πεσοῦνται ἐν τῇ πτώσει αὐτῶν, καὶ ἐν καιρῷ ἐπισκοπῆς ἀπολοῦνται, εἶπε Κύριος.

\vs{16}Τάδε λέγει Κύριος, στῆτε ἐπὶ ταῖς ὁδοῖς, καὶ ἴδετε, καὶ ἐρωτήσατε τρίβους Κυρίου αἰωνίους· καὶ ἴδετε ποία ἐστὶν ἡ ὁδὸς ἡ ἀγαθὴ, καὶ βαδίζετε ἐν αὐτῇ, καὶ εὑρήσετε ἁγνισμὸν ταῖς ψυχαῖς ὑμῶν· Καὶ εἶπαν, οὐ πορευσόμεθα.
\vs{17}Καθέστακα ἐφʼ ὑμᾶς σκοπούς· ἀκούσατε τῆς φωνῆς τῆς σάλπιγγος· καὶ εἶπαν, οὐκ ἀκουσόμεθα.

\vs{18}Διὰ τοῦτο ἤκουσαν τὰ ἔθνη, καὶ οἱ ποιμαίνοντες τὰ ποίμνια αὐτῶν.
\vs{19}Ἄκουε γῆ, ἰδοὺ ἐγὼ ἐπάγω ἐπὶ τὸν λαὸν τοῦτον κακὰ, τὸν καρπὸν ἀποστροφῆς αὐτῶν, ὅτι τῶν λόγων μου οὐ προσέσχον, καὶ τὸν νόμον μου ἀπώσαντο.
\vs{20}Ἱνατί μοι λίβανον ἐκ Σαβὰ φέρετε, καὶ κινάμωμον ἐκ γῆς μακρόθεν; τὰ ὁλοκαυτώματα ὑμῶν οὐκ εἰσὶ δεκτὰ, καὶ αἱ θυσίαι ὑμῶν οὐχ ἥδυνάν μοι.
\vs{21}Διατοῦτο τάδε λέγει Κύριος, ἰδοὺ ἐγὼ δίδωμι ἐπὶ τὸν λαὸν τοῦτον ἀσθένειαν, καὶ ἀσθενήσουσι πατέρες καὶ υἱοὶ ἅμα, γείτων καὶ ὁ πλησίον αὐτοῦ ἀπολοῦνται.

\vs{22}Τάδε λέγει Κύριος, ἰδοὺ λαὸς ἔρχεται ἀπὸ Βοῤῥᾶ, καὶ ἔθνη ἐξεγερθήσονται ἀπʼ ἐσχάτου τῆς γῆς.
\vs{23}Τόξον καὶ ζιβύνην κρατήσουσιν· ἰταμός ἐστι, καὶ οὐκ ἐλεήσει· φωνὴ αὐτοῦ, ὡς θάλασσα κυμαίνουσα· ἐφʼ ἵπποις καὶ ἅρμασι παρατάξεται ὡς πῦρ εἰς πόλεμον πρὸς σὲ, θύγατερ Σιών.

\vs{24}Ἠκούσαμεν τὴν ἀκοὴν αὐτῶν, παρελύθησαν αἱ χεῖρες ἡμῶν, θλίψις κατέσχεν ἡμᾶς, ὠδῖνες ὡς τικτούσης.
\vs{25}Μὴ ἐκπορεύεσθε εἰς ἀγρὸν, καὶ ἐν ταῖς ὁδοῖς μὴ βαδίζετε, ὅτι ῥομφαία τῶν ἐχθρῶν παροικεῖ κυκλόθεν.
\vs{26}Θύγατερ λαοῦ μου περίζωσαι σάκκον, κατάπασσε ἐν σποδῷ, πένθος ἀγαπητοῦ ποιήσαι σεαυτῇ κοπετὸν οἰκτρὸν, ὅτι ἐξαίφνης ἥξει ταλαιπωρία ἐφʼ ὑμᾶς.

\vs{27}Δοκιμαστὴν δέδωκά σε ἐν λαοῖς δεδοκιμασμένοις, καὶ γνώσῃ με ἐν τῷ δοκιμάσαι με τὴν ὁδὸν αὐτῶν,
\vs{28}πάντες ἀνήκοοι πορεύομενοι σκολιῶς· χαλκὸς καὶ σίδηρος, πάντες διεφθαρμένοι εἰσίν.
\vs{29}Ἐξέλιπε φυσητὴρ ἀπὸ πυρὸς, ἐξέλιπε μόλιβος· εἰς κενὸν ἀργυροκόπος ἀργυροκοπεῖ, πονηρία αὐτῶν οὐκ ἐτάκη.
\vs{30}Ἀργύριον ἀποδεδοκιμασμένον καλέσατε αὐτοὺς, ὅτι ἀπεδοκίμασεν αὐτοὺς Κύριος.

\ch{7}
\vs{2}Ἀκούσατε λόγον Κυρίου πᾶσα Ἰουδαία.
\vs{3}Τάδε λέγει Κύριος ὁ Θεὸς Ἰσραὴλ, διορθώσατε τὰς ὁδοὺς ὑμῶν, καὶ τὰ ἐπιτηδεύματα ὑμῶν, καὶ κατοικιῶ ὑμᾶς ἐν τῷ τόπῳ τούτῳ.
\vs{4}Μὴ πεποίθατε ἐφʼ ἑαυτοῖς ἐπὶ λόγοις ψευδέσιν, ὅτι τὸ παράπαν οὐκ ὠφελήσουσιν ὑμᾶς, λέγοντες, ναὸς Κυρίου, ναὸς Κυρίου ἐστίν.

\vs{5}Ὅτι ἐὰν διορθοῦντες διορθώσητε τὰς ὁδοὺς ὑμῶν καὶ τὰ ἐπιτηδεύματα ὑμῶν, καὶ ποιοῦντες ποιήσητε κρίσιν ἀναμέσον ἀνδρὸς, καὶ ἀναμέσον τοῦ πλησίον αὐτοῦ,
\vs{6}καὶ προσήλυτον καὶ ὀρφανὸν καὶ χήραν μὴ καταδυναστεύσητε, καὶ αἷμα ἀθῶον μὴ ἐκχέητε ἐν τῷ τόπῳ τούτῳ, καὶ ὀπίσω θεῶν ἀλλοτρίων μὴ πορεύησθε εἰς κακὸν ὑμῖν,
\vs{7}καὶ κατοικιῶ ὑμᾶς ἐν τῷ τόπῳ τούτῳ ἐν γῇ ᾗ ἔδωκα τοῖς πατράσιν ὑμῶν ἐξ αἰῶνος καὶ ἕως αἰῶνος.

\vs{8}Εἰ δὲ ὑμεῖς πεποίθατε ἐπὶ λόγοις ψευδέσιν, ὅθεν οὐκ ὠφεληθήσεσθε,
\vs{9}καὶ φονεύετε, καὶ μοιχᾶσθε, καὶ κλέπτετε, καὶ ὀμνύετε ἐπʼ ἀδίκῳ, καὶ θυμιᾶτε τῇ Βάαλ, καὶ ἐπορεύεσθε ὀπίσω θεῶν ἀλλοτρίων, ὧν οὐκ οἴδατε, τοῦ κακῶς εἶναι ὑμῖν,
\vs{10}καὶ ἤλθετε καὶ ἔστητε ἐνώπιον ἐμοῦ ἐν τῷ οἴκῳ οὗ ἐπικέκληται τὸ ὄνομά μου ἐπʼ αὐτῷ, καὶ εἴπατε, ἀπεσχήμεθα τοῦ μὴ ποιεῖν πάντα τὰ βδελύγματα ταῦτα.
\vs{11}Μὴ σπήλαιον λῃστῶν ὁ οἶκός μου, οὗ ἐπικέκληται τὸ ὄνομά μου ἐπʼ αὐτῷ ἐκεῖ ἐνώπιον ὑμῶν; καὶ ἰδοὺ ἐγὼ ἑώρακα, λέγει Κύριος.
\vs{12}Ὅτι πορεύθητε εἰς τὸν τόπον μου τὸν ἐν Σηλὼ, οὗ κατεσκήνωσα τὸ ὄνομά μου ἐκεῖ ἔμπροσθεν, καὶ ἴδετε ἃ ἐποίησα αὐτῷ ἀπὸ προσώπου κακίας λαοῦ μου Ἰσραήλ.

\vs{13}Καὶ νῦν ἀνθʼ ὧν ἐποιήσατε πάντα τὰ ἔργα ταῦτα, καὶ ἐλάλησα πρὸς ὑμᾶς καὶ οὐκ ἠκούσατέ μου, καὶ ἐκάλεσα ὑμᾶς καὶ οὐκ ἀπεκρίθητε,
\vs{14}τοίνυν κᾀγὼ ποιήσω τῷ οἴκῳ ᾧ ἐπικέκληται τὸ ὄνομά μου ἐπʼ αὐτῷ, ἐφʼ ᾧ ὑμεῖς πεποίθατε ἐπʼ αὐτῷ, καὶ τῷ τόπῳ ᾧ ἔδωκα ὑμῖν καὶ τοῖς πατράσιν ὑμῶν, καθὼς ἐποίησα τῇ Σηλώ.
\vs{15}Καὶ ἀποῤῥίψω ὑμᾶς ἀπὸ προσώπου μου, καθὼς ἀπέῤῥιψα τοὺς ἀδελφοὺς ὑμῶν, πᾶν τὸ σπέρμα Ἐφραΐμ.

\vs{16}Καὶ σὺ μὴ προσεύχου περὶ τοῦ λαοῦ τούτου, καὶ μὴ ἀξιοῦ τοῦ ἐλεηθῆναι αὐτοὺς, καὶ μὴ εὔχου, καὶ μὴ προσέλθῃς μοι περὶ αὐτῶν, ὅτι οὐκ εἰσακούσομαι.
\vs{17}Ἢ οὐχ ὁρᾷς, τί αὐτοὶ ποιοῦσιν ἐν ταῖς πόλεσιν Ἰούδα, καὶ ἐν ταῖς ὁδοῖς Ἱερουσαλήμ;
\vs{18}Οἱ υἱοὶ αὐτῶν συλλέγουσι ξύλα, καὶ οἱ πατέρες αὐτῶν καίουσι πῦρ, καὶ αἱ γυναῖκες αὐτῶν τρίβουσι σταῖς, τοῦ ποιῆσαι καυῶνας τῇ στρατιᾷ τοῦ οὐρανοῦ, καὶ ἔσπεισαν σπονδὰς θεοῖς ἀλλοτρίοις, ἵνα παροργίσωσί με.
\vs{19}Μὴ ἐμὲ αὐτοὶ παροργίζουσι; λέγει Κύριος· οὐχὶ ἑαυτοὺς, ὅπως καταισχυνθῇ τὰ πρόσωπα αὐτῶν;

\vs{20}Διατοῦτο τάδε λέγει Κύριος, ἰδοὺ ὀργὴ καὶ θυμός μου χεῖται ἐπὶ τὸν τόπον τοῦτον, καὶ ἐπὶ τοὺς ἀνθρώπους, καὶ ἐπὶ τὰ κτήνη, καὶ ἐπὶ πᾶν ξύλου τοῦ ἀγροῦ αὐτῶν, καὶ ἐπὶ τὰ γεννήματα τῆς γῆς, καὶ καυθήσεται καὶ οὐ σβεσθήσεται.

\vs{21}Τάδε λέγει Κύριος, τὰ ὁλοκαυτώματα ὑμῶν συναγάγετε μετὰ τῶν θυσιῶν ὑμῶν, καὶ φάγετε κρέα.
\vs{22}Ὅτι οὐκ ἐλάλησα πρὸς τοὺς πατέρας ὑμῶν, καὶ οὐκ ἐνετειλάμην αὐτοῖς ἐν ἡμέρᾳ ᾗ ἀνήγαγον αὐτοὺς ἐκ γῆς Αἰγύπτου, περὶ ὁλοκαυτωμάτων καὶ θυσίας.
\vs{23}Ἀλλʼ ἢ τὸ ῥῆμα τοῦτο ἐνετειλάμην αὐτοῖς, λέγων, ἀκούσατε τῆς φωνῆς μου, καὶ ἔσομαι ὑμῖν εἰς Θεὸν, καὶ ὑμεῖς ἔσεσθέ μοι εἰς λαὸν, καὶ πορεύεσθε ἐν πάσαις ταῖς ὁδοῖς μου, αἷς ἂν ἐντείλωμαι ὑμῖν, ὅπως ἂν εὖ ᾖ ὑμῖν.
\vs{24}Καὶ οὐκ ἤκουσάν μου, καὶ οὐ προσέσχε τὸ οὖς αὐτῶν, ἀλλʼ ἐπορεύθησαν ἐν τοῖς ἐνθυμήμασι τῆς καρδίας αὐτῶν τῆς κακῆς, καὶ ἐγενήθησαν εἰς τὰ ὄπισθεν· καὶ οὐκ εἰς τὰ ἔμπροσθεν,
\vs{25}ἀφʼ ἧς ἡμέρας ἐξήλθοσαν οἱ πατέρες αὐτῶν ἐκ γῆς Αἰγύπτου, καὶ ἕως τῆς ἡμέρας ταύτης· καὶ ἐξαπέστειλα πρὸς ὑμᾶς πάντας τοὺς δούλους μου, τοὺς προφήτας, ἡμέρας καὶ ὄρθρου, καὶ ἀπέστειλα,
\vs{26}καὶ οὐκ εἰσήκουσάν μου, καὶ οὐ προσέσχε τὸ οὖς αὐτῶν, καὶ ἐσκλήρυναν τὸν τράχηλον αὐτῶν ὑπὲρ τοὺς πατέρας αὐτῶν.

\vs{27}Καὶ ἐρεῖς αὐτοῖς τοῦτον τὸν λόγον,
\vs{28}τοῦτο τὸ ἔθνος ὃ οὐκ ἤκουσε τῆς φωνῆς Κυρίου, οὐδὲ ἐδέξατο παιδείαν, ἐξέλιπεν ἡ πίστις ἐκ στόματος αὐτῶν.

\vs{29}Κεῖρε τὴν κεφαλήν σου, καὶ ἀπόῤῥιπτε, καὶ ἀνάλαβε ἐπὶ χειλέων θρῆνον, ὅτι ἀπεδοκίμασε Κύριος, καὶ ἀπώσατο τὴν γενεὰν τὴν ποιοῦσαν ταῦτα.
\vs{30}Ὅτι ἐποίησαν οἱ υἱοὶ Ἰούδα τὸ πονηρὸν ἐναντίον ἐμοῦ, λέγει Κύριος· ἔταξαν τὰ βδελύγματα αὐτῶν ἐν τῷ οἴκῳ, οὗ ἐπικέκληται τὸ ὄνομά μου ἐπʼ αὐτὸν, τοῦ μιάναι αὐτόν.
\vs{31}Καὶ ᾠκοδόμησαν τὸν βωμὸν τοῦ Ταφὲθ, ὅς ἐστιν ἐν φάραγγι υἱοῦ Ἐννὸμ, τοῦ κατακαίειν τοὺς υἱοὺς αὐτῶν καὶ τὰς θυγατέρας αὐτῶν ἐν πυρὶ, ὃ οὐκ ἐνετειλάμην αὐτοῖς, καὶ οὐ διενοήθην ἐν τῇ καρδίᾳ μου.

\vs{32}Διατοῦτο ἰδοὺ ἡμέραι ἔρχονται, λέγει Κύριος, καὶ οὐκ ἐροῦσιν ἔτι, Βωμὸς τοῦ Ταφὲθ καὶ φάραγξ υἱοῦ Ἐννὸμ, ἀλλʼ ἡ φάραγξ τῶν ἀνῃρημένων· καὶ θάψουσιν ἐν τῷ Ταφὲθ, διὰ τὸ μὴ ὑπάρχειν τόπον.
\vs{33}Καὶ ἔσονται οἱ νεκροὶ τοῦ λαοῦ τούτου εἰς βρῶσιν τοῖς πετεινοῖς τοῦ οὐρανοῦ, καὶ τοῖς θηρίοις τῆς γῆς, καὶ οὐκ ἔσται ὁ ἀποσοβῶν.
\vs{34}Καὶ καταλύσω ἐκ πόλεων Ἰούδα καὶ ἐκ διόδων Ἱερουσαλὴμ φωνὴν εὐφραινομένων, καὶ φωνὴν χαιρόντων, φωνὴν νυμφίου, καὶ φωνὴν νύμφης, ὅτι εἰς ἐρήμωσιν ἔσται πᾶσα ἡ γῆ.

\ch{8}
Ἐν τῷ καιρῷ ἐκείνῳ, λέγει Κύριος, ἐξοίσουσι τὰ ὀστᾶ τῶν βασιλέων Ἰούδα, καὶ τὰ ὀστᾶ τῶν ἀρχόντων αὐτοῦ, καὶ τὰ ὀστᾶ τῶν ἱερέων, καὶ τὰ ὀστᾶ προφητῶν, καὶ τὰ ὀστᾶ τῶν κατοικούντων ἐν Ἱερουσαλὴμ ἐκ τῶν τάφων αὐτῶν,
\vs{2}καὶ ψύξουσιν αὐτὰ πρὸς τὸν ἥλιον καὶ τὴν σελήνην, καὶ πρὸς πάντας τοὺς ἀστέρας, καὶ πρὸς πᾶσαν τὴν στρατιὰν τοῦ οὐρανοῦ, ἃ ἠγάπησαν, καὶ οἷς ἐδούλευσαν, καὶ ὧν ἐπορεύθησαν ὀπίσω αὐτῶν, καὶ ὧν ἀντείχοντο, καὶ οἷς προσεκύνησαν αὐτοῖς· οὐ κοπήσονται, καὶ οὐ ταφήσονται, καὶ ἔσονται εἰς παράδειγμα ἐπὶ πρόσωπου τῆς γῆς,
\vs{3}ὅτι εἵλοντο τὸν θάνατον ἢ τὴν ζωὴν, καὶ πᾶσι τοῖς καταλοίποις τοῖς καταλειφθεῖσιν ἀπὸ τῆς γενεᾶς ἐκείνης, ἐν παντὶ τόπῳ οὗ ἐὰν ἐξώσω αὐτοὺς ἐκεῖ.

\vs{4}Ὅτι τάδε λέγει Κύριος, μὴ ὁ πίπτων οὐκ ἀνίσταται; ἢ ὁ ἀποστρέφων οὐκ ἀναστρέφει;
\vs{5}Διατί ἀπέστρεψεν ὁ λαός μου οὗτος ἀποστροφὴν ἀναιδῆ, καὶ κατεκρατήθησαν ἐν τῇ προαιρέσει αὐτῶν, καὶ οὐκ ἠθέλησαν τοῦ ἐπιστρέψαι;
\vs{6}ἐνωτίσασθε δὴ, καὶ ἀκούσατε· οὐχ οὕτω λαλήσουσιν, οὐκ ἔστιν ἄνθρωπος ὁ μετανοῶν ἀπὸ τῆς κακίας αὐτοῦ, λέγων, τί ἐποίησα; διέλιπεν ὁ τρέχων ἀπὸ τοῦ δρόμου αὐτοῦ, ὡς ἵππος κάθιδρος ἐν χρεμετισμῷ αὐτοῦ.
\vs{7}Καὶ ἡ ἀσίδα ἐν τῷ οὐρανῷ ἔγνω τὸν καιρὸν αὐτῆς, τρυγὼν καὶ χελιδὼν ἀγροῦ, στρουθία ἐφύλαξαν καιροὺς εἰσόδων ἑαυτῶν, ὁ δὲ λαός μου οὗτος οὐκ ἔγνω τὰ κρίματα Κυρίου.

\vs{8}Πῶς ἐρεῖτε, ὅτι σοφοί ἐσμεν ἡμεῖς, καὶ νόμος Κυρίου μεθʼ ἡμῶν ἐστιν; εἰς μάτην ἐγενήθη σχοῖνος ψευδὴς γραμματεῦσιν.
\vs{9}Ἠσχύνθησαν σοφοὶ, καὶ ἐπτοήθησαν καὶ ἑάλωσαν, ὅτι τὸν νόμον Κυρίου ἀπεδοκίμασαν· σοφία τίς ἐστιν ἐν αὐτοῖς;
\vs{10}Διατοῦτο δώσω τὰς γυναῖκας αὐτῶν ἑτέροις, καὶ τοὺς ἀγροὺς αὐτῶν τοῖς κληρονόμοις,
\vs{13}καὶ συνάξουσι τὰ γεννήματα αὐτῶν, λέγει Κύριος. Οὐκ ἔστι σταφυλὴ ἐν ταῖς ἀμπέλοις, καὶ οὐκ ἐστι σῦκα ἐν ταῖς συκαῖς, καὶ τὰ φύλλα κατεῤῥύηκεν.

\vs{14}Ἐπὶ τί ἡμεῖς καθήμεθα; συνάχθητε, καὶ εἰσέλθωμεν εἰς τὰς πόλεις τὰς ὀχυρὰς, καὶ ἀποῤῥιφῶμεν ἐκεῖ, ὅτι ὁ Θεὸς ἀπέῤῥιψεν ἡμᾶς, καὶ ἐπότισεν ἡμᾶς ὕδωρ χολῆς, ὅτι ἡμάρτομεν ἐναντίον αὐτοῦ.
\vs{15}Συνήχθημεν εἰς εἰρήνην, καὶ οὐκ ἦν ἀγαθὰ, εἰς καιρὸν ἰάσεως, καὶ ἰδοὺ σπουδή.

\vs{16}Ἐκ Δὰν ἀκουσόμεθα φωνὴν ὀξύτητος ἵππων αὐτοῦ· ἀπὸ φωνῆς χρεμετισμοῦ ἱππασίας ἵππων αὐτοῦ ἐσείσθη πᾶσα ἡ γῆ. καὶ ἥξει καὶ καταφάγεται τὴν γῆν, καὶ τὸ πλήρωμα αὐτῆς, πόλιν καὶ τοὺς κατοικοῦντας ἐν αὐτῇ.
\vs{17}Διότι ἰδοὺ ἐγὼ ἐξαποστέλλω εἰς ὑμᾶς ὄφεις θανατοῦντας, οἷς οὐκ ἔστιν ἐπᾴσαι, καὶ δήξονται ὑμᾶς
\vs{18}ἀνίατα μετʼ ὀδύνης καρδίας ὑμῶν ἀπορουμένης.

\vs{19}Ἰδοὺ φωνὴ κραυγῆς θυγατρὸς λαοῦ μου ἀπὸ γῆς μακρόθεν· μὴ Κύριος οὐκ ἔστιν ἐν Σιών; ἢ βασιλεὺς οὐκ ἔστιν ἐκεῖ; διότι παρώργισάν με ἐν τοῖς γλυπτοῖς αὐτῶν, καὶ ἐν ματαίοις ἀλλοτρίοις.
\vs{20}Διῆλθε θέρος, παρῆλθεν ἀμητὸς, καὶ ἡμεῖς οὐ διεσώθημεν.

\vs{21}Ἐπὶ συντρίμματι θυγατρὸς λαοῦ μου ἐσκοτώθην· ἐν ἀπορίᾳ κατίσχυσάν με ὠδῖνες ὡς τικτούσης.
\vs{22}Καὶ μὴ ῥητίνη οὐκ ἔστιν ἐν Γαλαὰδ, ἢ ἰατρὸς οὐκ ἔστιν ἐκεῖ; διατί οὐκ ἀνέβη ἴασις θυγατρὸς λαοῦ μου;

\vs{23}Τίς δώσει κεφαλῇ μου ὕδωρ, καὶ ὀφθαλμοῖς μου πηγὴν δακρύων; καὶ κλαύσομαι τὸν λαόν μου τοῦτον ἡμέρας καὶ νυκτὸς, τοὺς τετραυματισμένους θυγατρὸς λαοῦ μου.

\ch{9}
Τίς δῴη μοι ἐν τῇ ἐρήμῳ σταθμὸν ἔσχατον, καὶ καταλείψω τὸν λαόν μου, καὶ ἀπελεύσομαι ἀπʼ αὐτῶν; ὅτι πάντες μοιχῶνται, σύνοδος ἀθετούντων,
\vs{2}καὶ ἐνέτειναν τὴν γλῶσσαν αὐτῶν ὡς τόξον· ψεῦδος, καὶ οὐ πίστις ἐνίσχυσεν ἐπὶ τῆς γῆς, ὅτι ἐκ κακῶν εἰς κακὰ ἐξήλθοσαν, καὶ ἐμὲ οὐκ ἔγνωσαν, φησὶ Κύριος.
\vs{3}Ἕκαστος ἀπὸ τοῦ πλησίον αὐτοῦ φυλάξασθε, καὶ ἐπʼ ἀδελφοῖς αὐτῶν μὴ πεποίθατε, ὅτι πᾶς ἀδελφὸς πτέρνῃ πτερνιεῖ, καὶ πᾶς φίλος δολίως πορεύσεται.
\vs{4}Ἕκαστος κατὰ τοῦ φίλου αὐτοῦ καταπαίξεται, ἀλήθειαν οὐ μὴ λαλήσωσι· μεμάθηκεν ἡ γλῶσσα αὐτῶν λαλεῖν ψευδῆ, ἠδίκησαν, καὶ οὐ διέλιπον τοῦ ἐπιστρέψαι.
\vs{5}Τόκος ἐπὶ τόκῳ, καὶ δόλος ἐπὶ δόλῳ· οὐκ ἤθελον εἰδέναι με, φησὶ Κύριος.

\vs{6}Διατοῦτο τάδε λέγει Κύριος, ἰδοὺ ἐγὼ πυρώσω αὐτοὺς, καὶ δοκιμῶ αὐτούς· ὅτι ποιήσω ἀπὸ προσώπου πονηρίας θυγατρὸς λαοῦ μου.
\vs{7}Βολὶς τιτρώσκουσα ἡ γλῶσσα αὐτῶν, δόλια τὰ ῥήματα τοῦ στόματος αὐτῶν· τῷ πλησίον αὐτοῦ λαλεῖ εἰρηνικὰ, καὶ ἐν ἑαυτῷ ἔχει τὴν ἔχθραν.
\vs{8}Μὴ ἐπὶ τούτοις οὐκ ἐπισκέψομαι; λέγει Κύριος· ἢ ἐν λαῷ τοιούτῳ οὐκ ἐκδικήσει ἡ ψυχή μου;
\vs{9}Ἐπὶ τὰ ὄρη λάβετε κοπετὸν, καὶ ἐπὶ τὰς τρίβους τῆς ἐρήμου θρῆνον, ὅτι ἐξέλιπον παρὰ τὸ μὴ εἶναι ἀνθρώπους· οὐκ ἤκουσαν φωνὴν ὑπάρξεως ἀπὸ πετεινῶν τοῦ οὐρανοῦ, καὶ ἕως κτηνῶν, ἐξέστησαν, ᾤχοντο.
\vs{10}Καὶ δώσω τὴν Ἱερουσαλὴμ εἰς μετοικίαν, καὶ εἰς κατοικητήριον δρακόντων, καὶ τὰς πόλεις Ἰούδα εἰς ἀφανισμὸν θήσομαι, παρὰ τὸ μὴ κατοικεῖσθαι.

\vs{11}Τίς ὁ ἄνθρωπος ὁ συνετὸς, καὶ συνέτω τοῦτο; καὶ ᾧ λόγος στόματος Κυρίου πρὸς αὐτὸν, ἀναγγειλάτω ὑμῖν, ἕνεκεν τίνος ἀπώλετο ἡ γῆ, ἀνήφθη, ὡς ἔρημος, παρὰ τὸ μὴ διοδεύεσθαι αὐτήν;
\vs{12}Καὶ εἶπε Κύριος πρὸς μὲ, διὰ τὸ ἐγκαταλιπεῖν αὐτοὺς τὸν νόμον μου, ὃν ἔδωκα πρὸ προσώπου αὐτῶν, καὶ οὐκ ἤκουσαν τῆς φωνῆς μου,
\vs{13}ἀλλʼ ἐπορεύθησαν ὀπίσω τῶν ἀρεστῶν τῆς καρδίας αὐτῶν τῆς κακῆς, καὶ ὀπίσω τῶν εἰδώλων ἃ ἐδίδαξαν αὐτοὺς οἱ πατέρες αὐτῶν.
\vs{14}Διατοῦτο τάδε λέγει Κύριος ὁ Θεὸς Ἰσραὴλ, ἰδοὺ ἐγὼ ψωμιῶ αὐτοὺς ἀνάγκας, καὶ ποτιῶ αὐτοὺς ὕδωρ χολῆς,
\vs{15}καὶ διασκορπιῶ αὐτοὺς ἐν τοῖς ἔθνεσιν, εἰς οὓς οὐκ ἐγίνωσκον αὐτοὶ καὶ οἱ πατέρες αὐτῶν, καὶ ἐπαποστελῶ ἐπʼ αὐτοὺς τὴν μάχαιραν, ἕως τοῦ ἐξαναλῶσαι αὐτοὺς ἐν αὐτῇ.

\vs{16}Τάδε λέγει Κύριος, καλέσατε τὰς θρηνούσας, καὶ ἐλθέτωσαν, καὶ πρὸς τὰς σοφὰς ἀποστείλατε, καὶ φθεγξάσθωσαν,
\vs{17}καὶ λαβέτωσαν ἐφʼ ὑμᾶς θρῆνον, καὶ καταγαγέτωσαν οἱ ὀφθαλμοὶ ὑμῶν δάκρυα, καὶ τὰ βλέφαρα ὑμῶν ῥείτω ὕδωρ,
\vs{18}ὅτι φωνὴ οἰκτροῦ ἠκούσθη ἐν Σιών· πῶς ἐταλαιπωρήσαμεν, κατῃσχύνθημεν σφόδρα, ὅτι ἐγκατελίπομεν τὴν γῆν, καὶ ἀπεῤῥίψαμεν τὰ σκηνώματα ἡμῶν;
\vs{19}Ἀκούσατε δὴ γυναῖκες λόγον Θεοῦ, καὶ δεξάσθω τὰ ὦτα ὑμῶν λόγους στόματος αὐτοῦ, καὶ διδάξατε τὰς θυγατέρας ὑμῶν οἶκτον, καὶ γυνὴ τὴν πλησίον αὐτῆς θρῆνον.
\vs{20}Ὅτι ἀνέβη θάνατος διὰ τῶν θυρίδων ὑμῶν, εἰσῆλθεν εἰς τὴν γῆν ὑμῶν, τοῦ ἐκτρίψαι νήπια ἔξωθεν, καὶ νεανίσκους ἀπὸ τῶν πλατειῶν.
\vs{21}Καὶ ἔσονται οἱ νεκροὶ τῶν ἀνθρώπων εἰς παράδειγμα ἐπὶ προσώπου τοῦ πεδίου τῆς γῆς ὑμῶν, ὡς χόρτος ὀπίσω θερίζοντος, καὶ οὐκ ἔσται ὁ συνάγων.

\vs{22}Τάδε λέγει Κύριος, μὴ καυχάσθω ὁ σοφὸς ἐν τῇ σοφίᾳ αὐτοῦ, καὶ μὴ καυχάσθω ὁ ἰσχυρὸς ἐν τῆ ἰσχύϊ αὐτοῦ, καὶ μὴ καυχάσθω ὁ πλούσιος ἐν τῷ πλούτῳ αὐτοῦ,
\vs{23}ἀλλʼ ἢ ἐν τούτῳ καυχάσθω ὁ καυχώμενος, συνιεῖν καὶ γινώσκειν, ὅτι ἐγώ εἰμι Κύριος ὁ ποιῶν ἔλεος καὶ κρίμα καὶ δικαιοσύνην ἐπὶ τῆς γῆς, ὅτι ἐν τούτοις τὸ θέλημά μου, λέγει Κύριος.

\vs{24}Ἰδοὺ ἡμέραι ἔρχονται, λέγει Κύριος, καὶ ἐπισκέψομαι ἐπὶ πάντας περιτετμημένους ἀκροβυστίας αὐτῶν.
\vs{25}Ἐπʼ Αἴγυπτον, καὶ ἐπὶ Ἰδουμαίαν, καὶ ἐπὶ Ἐδὼμ, καὶ ἐπὶ υἱοὺς Ἀμμὼν, καὶ ἐπὶ υἱοὺς Μωὰβ, καὶ ἐπὶ πάντα περικειρόμενον τὰ κατὰ πρόσωπον αὐτοῦ, τοὺς κατοικοῦντας ἐν τῇ ἐρήμῳ, ὅτι πάντα τὰ ἔθνη ἀπερίτμητα σαρκὶ, καὶ πᾶς οἶκος Ἰσραὴλ ἀπερίτμητοι καρδίας αὐτῶν.

\ch{10}
Ἀκούσατε τὸν λόγον Κυρίου, ὃν ἐλάλησεν ἐφʼ ὑμᾶς, οἶκος Ἰσραήλ.

\vs{2}Τάδε λέγει Κύριος, κατὰ τὰς ὁδοὺς τῶν ἐθνῶν μὴ μανθάνετε, καὶ ἀπὸ τῶν σημείων τοῦ οὐρανοῦ μὴ φοβεῖσθε, ὅτι φοβοῦνται αὐτὰ τοῖς προσώποις αὐτῶν.
\vs{3}Ὅτι τὰ νόμιμα τῶν ἐθνῶν μάταια· ξύλον ἐστὶν ἐκ τοῦ δρυμοῦ ἐκκεκομμένον, ἔργον τέκτονος, καὶ χώνευμα,
\vs{4}ἀργυρίῳ καὶ χρυσίῳ κεκαλλωπισμένα, ἐν σφύραις καὶ ἥλοις ἐστερέωσαν αὐτά·
\vs{5}Θήσουσιν αὐτὰ, καὶ οὐ κινηθήσονται· ἀργύριον τορευτόν ἐστιν, οὐ πορεύσονται,
\vs{9}ἀργύριον προσβλητόν ἐστιν. Ἀπὸ Θαρσεὶς ἥξει χρυσίον Μωφὰζ, καὶ χεὶρ χρυσοχόων, ἔργα τεχνιτῶν πάντα, ὑάκινθον καὶ πορφύραν ἐνδύσουσιν αὐτά.
\vs{9a}Αἰρόμενα ἀρθήσονται, ὅτι οὐκ ἐπιβήσονται· μὴ φοβηθῆτε αὐτὰ, ὅτι οὐ μὴ κακοποιήσωσι, καὶ ἀγαθὸν οὐκ ἔστιν ἐν αὐτοῖς.

\vs{11}Οὕτως ἐρεῖτε αὐτοῖς, θεοὶ οἳ τὸν οὐρανὸν καὶ τὴν γῆν οὐκ ἐποίησαν, ἀπολέσθωσαν ἐκ τῆς γῆς, καὶ ὑποκάτωθεν τοῦ οὐρανοῦ τούτου.
\vs{12}Κύριος ὁ ποιήσας τὴν γῆν ἐν τῇ ἰσχύϊ αὐτοῦ, ὁ ἀνορθώσας τὴν οἰκουμένην ἐν τῇ σοφίᾳ αὐτοῦ, καὶ τῇ φρονήσει αὐτοῦ ἐξέτεινε τὸν οὐρανὸν,
\vs{13}καὶ πλῆθος ὕδατος ἐν οὐρανῷ· καὶ ἀνήγαγε νεφέλας ἐξ ἐσχάτου τῆς γῆς, ἀστραπὰς εἰς ὑετὸν ἐποίησε, καὶ ἐξήγαγε φῶς ἐκ θησαυρῶν αὐτοῦ.
\vs{14}Ἐμωράνθη πᾶς ἄνθρωπος ἀπὸ γνώσεως, κατῃσχύνθη πᾶς χρυσοχόος ἐπὶ τοῖς γλυπτοῖς αὐτοῦ, ὅτι ψευδῆ ἐχώνευσεν, οὐκ ἔστι πνεῦμα ἐν αὐτοῖς.
\vs{15}Μάταιά ἐστιν ἔργα ἐμπεπαιγμένα, ἐν καιρῷ ἐπισκοπῆς αὐτῶν ἀπολοῦνται.
\vs{16}Οὐκ ἔστι τοιαύτη μερὶς τῷ Ἰακὼβ, ὅτι ὁ πλάσας τὰ πάντα, αὐτὸς κληρονομία αὐτοῦ, Κύριος ὄνομα αὐτῷ.

\vs{17}Συνήγαγεν ἔξωθεν τὴν ὑπόστασίν σου, κατοικοῦσαν ἐν ἐκλεκτοῖς.
\vs{18}Ὅτι τάδε λέγει Κύριος, ἰδοὺ ἐγὼ σκελίζω τοὺς κατοικοῦντας τὴν γῆν ταύτην ἐν θλίψει, ὅπως εὑρεθῇ ἡ πληγή σου.

\vs{19}Οὐαὶ ἐπὶ συντρίμματί σου, ἀλγηρὰ ἡ πληγή σου· κᾀγὼ εἶπα, ὄντως τοῦτο τὸ τραῦμά σου, καὶ κατέλαβέ σε.
\vs{20}Ἡ σκηνή σου ἐταλαιπώρησεν, ὤλετο· καὶ πᾶσαι αἱ δέῤῥεις σου διεσπάσθησαν· οἱ υἱοί μου καὶ τὰ πρόβατά μου οὐκ εἰσὶν, οὐκ ἔστιν ἔτι τόπος τῆς σκηνῆς μου, τόπος τῶν δέῤῥεών μου.
\vs{21}Ὅτι οἱ ποιμένες ἠφρονεύσαντο, καὶ τὸν Κύριον οὐκ ἐζήτησαν· διατοῦτο οὐκ ἐνόησε πᾶσα ἡ νομὴ, καὶ διεσκορπίσθησαν,
\vs{22}φωνὴ ἀκοῆς ἰδοὺ ἔρχεται καὶ σειαμὸς μέγας ἐκ γῆς Βοῤῥᾶ, τοῦ τάξαι τὰς πόλεις Ἰούδα εἰς ἀφανισμὸν, καὶ κοίτην στρουθῶν.

\vs{23}Οἶδα, Κύριε, ὅτι οὐχὶ τοῦ ἀνθρώπου ἡ ὁδὸς αὐτοῦ, οὐδὲ ἀνὴρ πορεύσεται καὶ κατορθώσει πορείαν αὐτοῦ.
\vs{24}Παίδευσον ἡμᾶς Κύριε, πλὴν ἐν κρίσει καὶ μὴ ἐν θυμῷ, ἵνα μὴ ὀλίγους ἡμᾶς ποιήσῃς.
\vs{25}Ἔκχεον τὸν θυμόν σου ἐπὶ ἔθνη τὰ μὴ εἰδότα σε, καὶ ἐπὶ γενεὰς αἳ τὸ ὄνομά σου οὐκ ἐπεκαλέσαντο, ὅτι κατέφαγον τὸν Ἰακὼβ καὶ ἐξανήλωσαν αὐτὸν, καὶ τὴν νομὴν αὐτοῦ ἠρήμωσαν.

\ch{11}
Ὁ ΛΟΓΟΣ Ὁ ΓΕΝΟΜΕΝΟΣ ΠΑΡΑ ΚΥΡΙΟΥ ΠΡΟΣ ἹΕΡΕΜΙΑΝ, ΛΕΓΩΝ,

\vs{2}Ἀκούσατε τοὺς λόγους τῆς διαθήκης ταύτης, καὶ λαλήσεις πρὸς ἄνδρας Ἰούδα, καὶ πρὸς τοὺς κατοικοῦντας ἐν Ἱερουσαλὴμ,
\vs{3}καὶ ἐρεῖς πρὸς αὐτοὺς, τάδε λέγει Κύριος ὁ Θεὸς Ἰσραὴλ, ἐπικατάρατος ὁ ἄνθρωπος, ὃς οὐκ ἀκούσεται τῶν λόγων τῆς διαθήκης ταύτης,
\vs{4}ἧς ἐνετειλάμην τοῖς πατράσιν ὑμῶν, ἐν ἡμέρᾳ ᾗ ἀνήγαγον αὐτοὺς ἐκ γῆς Αἰγύπτου ἐκ καμίνου τῆς σιδηρᾶς, λέγων, ἀκούσατε τῆς φωνῆς μου, καὶ ποιήσατε πάντα ὅσα ἐὰν ἐντείλωμαι ὑμῖν, καὶ ἔσεσθέ μοι εἰς λαὸν, καὶ ἐγὼ ἔσομαι ὑμῖν εἰς Θεόν·
\vs{5}Ὅπως στήσω τὸν ὅρκον μου, ὃν ὤμοσα τοῖς πατράσιν ὑμῶν, τοῦ δοῦναι αὐτοῖς γῆν ῥέουσαν γάλα καὶ μέλι, καθὼς ἡ ἡμέρα αὕτη· καὶ ἀπεκρίθην καὶ εἶπα, γένοιτο Κύριε.
\vs{6}Καὶ εἶπε Κύριος πρὸς μὲ, ἀνάγνωθι τοὺς λόγους τούτους ἐν πόλεσιν Ἰούδα, καὶ ἔξωθεν Ἱερουσαλὴμ, λέγων, ἀκούσατε τοὺς λόγους τῆς διαθήκης ταύτης, καὶ ποιήσατε αὐτούς.
\vs{8}Καὶ οὐκ ἐποίησαν.

\vs{9}Καὶ εἶπε Κύριος πρὸς μὲ, εὑρέθη σύνδεσμος ἐν ἀνδράσιν Ἰούδα, καὶ ἐν τοῖς κατοικοῦσιν ἐν Ἱερουσαλήμ·
\vs{10}Ἐπεστράφησαν ἐπὶ τὰς ἀδικίας τῶν πατέρων αὐτῶν τῶν πρότερον, οἳ οὐκ ἠθέλησαν εἰσακοῦσαι τῶν λόγων μου, καὶ ἰδοὺ αὐτοὶ πορεύονται ὀπίσω θεῶν ἀλλοτρίων, τοῦ δουλεύειν αὐτοῖς· καὶ διεσκέδασεν οἶκος Ἰσραὴλ καὶ οἶκος Ἰούδα τὴν διαθήκην μου, ἣν διεθέμην πρὸς τοὺς πατέρας αὐτῶν.

\vs{11}Διατοῦτο τάδε λέγει Κύριος, ἰδοὺ ἐγὼ ἐπάγω ἐπὶ τὸν λαὸν τοῦτον κακὰ, ἐξ ὧν οὐ δυνήσονται ἐξελθεῖν ἐξ αὐτῶν· καὶ κεκράξονται πρὸς μὲ, καὶ οὐκ εἰσακούσομαι αὐτῶν.
\vs{12}Καὶ πορεύσονται πόλεις Ἰούδα καὶ οἱ κατοικοῦντες Ἱερουσαλὴμ, καὶ κεκράξονται πρὸς τοὺς θεοὺς, οἷς αὐτοὶ θυμιῶσιν αὐτοῖς, οἳ μὴ σώσουσιν αὐτοὺς ἐν τῷ καιρῷ τῶν κακῶν αὐτῶν.
\vs{13}Ὅτι κατʼ ἀριθμὸν τῶν πόλεών σου ἦσαν θεοί σου Ἰούδα, καὶ κατʼ ἀριθμὸν ἐξόδων τῆς Ἱερουσαλὴμ ἐτάξατε βωμοὺς θυμιᾷν τῇ Βάαλ.

\vs{14}Καὶ σὺ μὴ προσεύχου περὶ τοῦ λαοῦ τούτου, καὶ μὴ ἀξίου περὶ αὐτῶν ἐν δεήσει καὶ προσευχῇ, ὅτι οὐκ εἰσακούσομαι ἐν τῷ καιρῷ ἐν ᾧ ἐπικαλοῦνταί με, ἐν καιρῷ κακώσεως αὐτῶν.
\vs{15}Τί ἡ ἠγαπημένη ἐν τῷ οἴκῳ μου ἐποίησε βδέλυγμα; μὴ εὐχαὶ καὶ κρέα ἅγια ἀφελοῦσιν ἀπὸ σοῦ τὰς κακίας σου, ἢ τούτοις διαφεύξῃ;
\vs{16}Ἐλαίαν ὡραίαν εὔσκιον τῷ εἴδει ἐκάλεσε Κύριος τὸ ὄνομά σου, εἰς φωνὴν περιτομῆς αὐτῆς· ἀνήφθη πῦρ ἐπʼ αὐτὴν, μεγάλη ἡ θλίψις ἐπὶ σὲ, ἠχρειώθησαν οἱ κλάδοι αὐτῆς·
\vs{17}Καὶ Κύριος ὁ καταφυτεύσας σε, ἐλάλησεν ἐπὶ σὲ κακὰ ἀντὶ τῆς κακίας οἴκου Ἰσραὴλ καὶ οἴκου Ἰούδα, ὅ, τι ἐποίησαν ἑαυτοῖς τοῦ παροργίσαι με ἐν τῷ θυμιᾷν αὐτοὺς τῇ Βάαλ.

\vs{18}Κύριε γνώρισόν μοι, καὶ γνώσομαι· τότε εἶδον τὰ ἐπιτηδεύματα αὐτῶν.
\vs{19}Ἐγὼ δὲ ὡς ἀρνίον ἄκακον ἀγόμενον τοῦ θύεσθαι, οὐκ ἔγνων· ἐπʼ ἐμὲ ἐλογίσαντο λογισμὸν πονηρὸν, λέγοντες, δεῦτε καὶ ἐμβάλωμεν ξύλον εἰς τὸν ἄρτον αὐτοῦ, καὶ ἐκτρίψωμεν αὐτὸν ἀπὸ γῆς ζώντων, καὶ τὸ ὄνομα αὐτοῦ οὐ μὴ μνησθῇ οὐκέτι.
\vs{20}Κύριε κρίνων δίκαια, δοκιμάζων νεφροὺς καὶ καρδίας, ἴδοιμι τῆν παρὰ σοῦ ἐκδίκησιν ἐξ αὐτῶν, ὅτι πρὸς σὲ ἀπεκάλυψα τὸ δικαίωμά μου.

\vs{21}Διατοῦτο τάδε λέγει Κύριος ἐπὶ τοὺς ἄνδρας Ἀναθὼθ τοὺς ζητοῦντας τὴν ψυχήν μου, τοὺς λέγοντας, οὐ μὴ προφητεύσεις ἐπὶ τῷ ὀνόματι Κυρίου, εἰ δὲ μὴ, ἀποθάνῃ ἐν ταῖς χερσὶν ἡμῶν·
\vs{22}Ἰδοὺ ἐγὼ ἐπισκέψομαι ἐπʼ αὐτούς· οἱ νεανίσκοι αὐτῶν ἐν μαχαίρᾳ ἀποθανοῦνται, καὶ οἱ υἱοὶ αὐτῶν καὶ αἱ θυγατέρες αὐτῶν τελευτήσουσιν ἐν λιμῷ,
\vs{23}καὶ ἐγκατάλειμμα οὐκ ἔσται αὐτῶν, ὅτι ἐπάξω κακὰ ἐπὶ τοὺς κατοικοῦντας ἐν Ἀναθὼθ, ἐν ἐνιαυτῷ ἐπισκέψεως αὐτῶν.

\ch{12}
Δίκαιος εἶ Κύριε, ὅτι ἀπολογήσομαι πρὸς σέ· πλὴν κρίματα λαλήσω πρὸς σέ· τί ὅτι ὁδὸς ἀσεβῶν εὐοδοῦται; εὐθήνησαν πάντες οἱ ἀθετοῦντες ἀθετήματα;
\vs{2}Ἐφύτευσας αὐτοὺς, καὶ ἐῤῥιζώθησαν· ἐτεκνοποιήσαντο, καὶ ἐποίησαν καρπόν· ἐγγὺς εἶ σὺ τοῦ στόματος αὐτῶν, καὶ πόῤῥω ἀπὸ τῶν νεφρῶν αὐτῶν.
\vs{3}Καὶ σύ Κύριε γινώσκεις με, δεδοκίμακας τὴν καρδίαν μου ἐναντίον σου· ἅγνισον αὐτοὺς εἰς ἡμέραν σφαγῆς αὐτῶν.
\vs{4}Ἕως πότε πενθήσει ἡ γῆ, καὶ πᾶς ὁ χόρτος τοῦ ἀγροῦ ξηρανθήσεται ἀπὸ κακίας τῶν κατοικούντων ἐν αὐτῆ; ἠφανίσθησαν κτήνη καὶ πετεινὰ, ὅτι εἶπαν, οὐκ ὄψεται ὁ Θεὸς ὁδοὺς ἡμῶν.

\vs{5}Σοῦ οἱ πόδες τρέχουσι, καὶ ἐκλύουσί σε· πῶς παρασκευάσῃ ἐφʼ ἵπποις; καὶ ἐν γῇ εἰρήνης σου πέποιθας· πῶς ποιήσεις ἐν φρυάγματι τοῦ Ἰορδάνου;
\vs{6}Ὅτι καὶ οἱ ἀδελφοί σου καὶ ὁ οἶκος τοῦ πατρός σου, καὶ οὗτοι ἠθέτησάν σε, καὶ αὐτοὶ ἐβόησαν, ἐκ τῶν ὀπίσω σου ἐπισυνήχθησαν· μὴ πιστεύσῃς ἐν αὐτοῖς, ὅτι λαλήσουσι πρὸς σὲ καλά.

\vs{7}Ἐγκαταλέλοιπα τὸν οἶκόν μου, ἀφῆκα τὴν κληρονομίαν μου, ἔδωκα τὴν ἠγαπημένην ψυχήν μου εἰς χεῖρας ἐχθρῶν αὐτῆς.
\vs{8}Ἐγενήθη ἡ κληρονομία μου ἐμοὶ ὡς λέων ἐν δρυμῷ, ἔδωκεν ἐπʼ ἐμὲ τὴν φωνὴν αὐτῆς, διατοῦτο ἐμίσησα αὐτήν.
\vs{9}Μὴ σπήλαιον ὑαίνης ἡ κληρονομία μου ἐμοὶ, ἢ σπήλαιον κύκλῳ αὐτῆς; βαδίσατε, συναγάγετε πάντα τὰ θηρία τοῦ ἀγροῦ, καὶ ἐλθέτωσαν τοῦ φαγεῖν αὐτήν.

\vs{10}Ποιμένες πολλοὶ διέφθειραν τὸν ἀμπελῶνά μου, ἐμόλυναν τὴν μερίδα μου, ἔδωκαν τὴν μερίδα τὴν ἐπιθυμητήν μου εἰς ἔρημον ἄβατον,
\vs{11}ἐτέθη εἰς ἀφανισμὸν ἀπωλείας· διʼ ἐμὲ ἀφανισμῷ ἠφανίσθη πᾶσα ἡ γῆ, ὅτι οὐκ ἔστιν ἀνὴρ τιθέμενος ἐν καρδίᾳ.
\vs{12}Ἐπὶ πᾶσαν διεκβολὴν ἐν τῇ ἐρήμῳ ἦλθον ταλαιπωροῦντες, ὅτι μάχαιρα τοῦ Κυρίου καταφάγεται ἀπʼ ἄκρου τῆς γῆς ἕως ἄκρου τῆς γῆς· οὐκ ἔστιν εἰρήνη πάσῃ σαρκί.
\vs{13}Σπείρατε πυροὺς, καὶ ἀκάνθας θερίζετε· οἱ κλῆροι αὐτῶν οὐκ ὠφελήσουσιν αὐτούς· αἰσχύνθητε ἀπὸ καυχήσεως ὑμῶν, ἀπὸ ὀνειδισμοῦ ἔναντι Κυρίου.

\vs{14}Ὅτι τάδε λέγει Κύριος περὶ πάντων τῶν γειτόνων τῶν πονηρῶν, τῶν ἁπτομένων τῆς κληρονομίας μου, ἧς ἐμέρισα τῷ λαῷ μου Ἰσραὴλ, ἰδοὺ ἐγὼ ἀποσπῶ αὐτοὺς ἀπὸ τῆς γῆς αὐτῶν, καὶ τὸν Ἰούδαν ἐκβαλῶ ἐκ μέσου αὐτῶν.

\vs{15}Καὶ ἔσται μετὰ τὸ ἐκβαλεῖν με αὐτοὺς, ἐπιστρέψω καὶ ἐλεήσω αὐτοὺς, καὶ κατοικιῶ αὐτοὺς, ἕκαστον εἰς τὴν κληρονομίαν αὐτοῦ, καὶ ἕκαστον εἰς τὴν γῆν αὐτοῦ.
\vs{16}Καὶ ἔσται ἐὰν μαθόντες μάθωσι τὴν ὁδὸν τοῦ λαοῦ μου, τοῦ ὀμνύειν τῷ ὀνόματί μου, ζῇ Κύριος, καθὼς ἐδίδαξαν τὸν λαόν μου ὀμνύειν τῇ Βάαλ, καὶ οἰκοδομηθήσεται ἐν μέσῳ τοῦ λαοῦ μου.
\vs{17}Ἐὰν δὲ μὴ ἐπιστρέψωσι, καὶ ἐξαρῶ τὸ ἔθνος ἐκεῖνο ἐξάρσει καὶ ἀπωλείᾳ.

\ch{13}
Τάδε λέγει Κύριος, βάδισον καὶ κτῆσαι σεαυτῷ περίζωμα λινοῦν, καὶ περίθου περὶ τὴν ὀσφύν σου, καὶ ἐν ὕδατι οὐ διελεύσεται.
\vs{2}Καὶ ἐκτησάμην τὸ περίζωμα κατὰ τὸν λόγον Κυρίου, καὶ περιέθηκα περὶ τὴν ὀσφύν μου.
\vs{3}Καὶ ἐγενήθη λόγος Κυρίου πρὸς μὲ, λέγων,
\vs{4}λάβε τὸ περίζωμα τὸ περὶ τὴν ὀσφύν σου, καὶ ἀνάστηθι, καὶ βάδισον ἐπὶ τὸν Εὐφράτην, καὶ κατάκρυψον αὐτὸ ἐκεῖ ἐν τῇ τρυμαλιᾷ τῆς πέτρας.
\vs{5}Καὶ ἐπορεύθην, καὶ ἔκρυψα αὐτὸ ἐν τῷ Εὐφράτῃ, καθὼς ἐνετείλατό μοι Κύριος.
\vs{6}Καὶ ἐγένετο μεθʼ ἡμέρας πολλὰς, καὶ εἶπε Κύριος πρὸς μὲ, ἀνάστηθι, βάδισον ἐπὶ τὸν Εὐφράτην, καὶ λάβε ἐκεῖθεν τὸ περίζωμα, ὃ ἐνετειλάμην σοι τοῦ κατακρύψαι ἐκεῖ.
\vs{7}Καὶ ἐπορεύθην ἐπὶ τὸν Εὐφράτην ποταμὸν, καὶ ὤρυξα, καὶ ἔλαβον τὸ περίζωμα ἐκ τοῦ τόπου οὗ κατώρυξα αὐτὸ ἐκεῖ, καὶ ἰδοὺ διεφθαρμένον ἦν, ὃ οὐ μὴ χρησθῇ εἰς οὐθέν.

\vs{8}Καὶ ἐγενήθη λόγος Κυρίου πρὸς μὲ, λέγων,
\vs{9}Τάδε λέγει Κύριος, οὕτω φθερῶ τὴν ὕβριν Ἰούδα καὶ τὴν ὕβριν Ἱερουσαλὴμ,
\vs{10}τὴν πολλὴν ταύτην ὕβριν, τοὺς μὴ βουλομένους ὑπακούειν τῶν λόγων μου, καὶ πορευθέντας ὀπίσω θεῶν ἀλλοτρίων τοῦ δουλεύειν αὐτοῖς, καὶ τοῦ προσκυνεῖν αὐτοῖς· καὶ ἔσονται ὥσπερ τὸ περίζωμα τοῦτο, ὃ οὐ χρησθήσεται εἰς οὐθέν.
\vs{11}Ὅτι καθάπερ κολλᾶται τὸ περίζωμα περὶ τὴν ὀσφὺν τοῦ ἀνθρώπου, οὕτως ἐκόλλησα πρὸς ἐμαυτὸν τὸν οἶκον τοῦ Ἰσραὴλ, καὶ πάντα οἶκον Ἰούδα, τοῦ γενέσθαι μοι εἰς λαὸν ὀνομαστὸν, καὶ εἰς καύχημα, καὶ εἰς δόξαν, καὶ οὐκ εἰσήκουσάν μου.

\vs{12}Καὶ ἐρεῖς πρὸς τὸν λαὸν τοῦτον, πᾶς ἀσκὸς πληρωθήσεται οἴνου· καὶ ἔσται ἐὰν εἴπωσι πρὸς σὲ, μὴ γνόντες οὐ γνωσόμεθα, ὅτι πᾶς ἀσκὸς πληρωθήσεται οἴνου;
\vs{13}Καὶ ἐρεῖς πρὸς αὐτοὺς, τάδε λέγει Κύριος, ἰδοὺ ἐγὼ πληρῶ τοὺς κατοικοῦντας τὴν γῆν ταύτην, καὶ τοὺς βασιλεῖς αὐτῶν τοὺς καθημένους υἱοὺς τοῦ Δαυὶδ ἐπὶ τοῦ θρόνου αὐτῶν, καὶ τοὺς ἱερεῖς καὶ τοὺς προφήτας, καὶ τὸν Ἰούδαν καὶ πάντας τοὺς κατοικοῦντας ἐν Ἱερουσαλὴμ, μεθύσματι·
\vs{14}Καὶ διασκορπιῶ αὐτοὺς ἄνδρα καὶ τὸν ἀδελφὸν αὐτοῦ, καὶ τοὺς πατέρας αὐτῶν, καὶ τοὺς υἱοὺς αὐτῶν ἐν τῷ αὐτῷ· οὐκ ἐπιποθήσω, λέγει Κύριος, καὶ οὐ φείσομαι, καὶ οὐκ οἰκτειρήσω ἀπὸ διαφθορᾶς αὐτῶν.

\vs{15}Ἀκούσατε, καὶ ἐνωτίσασθε, καὶ μὴ ἐπαίρεσθε, ὅτι Κύριος ἐλάλησε.
\vs{16}Δότε τῷ Κυρίῳ Θεῷ ὑμῶν δόξαν πρὸ τοῦ συσκοτάσαι, καὶ πρὸ τοῦ προσκόψαι πόδας ὑμῶν ἐπʼ ὄρη σκοτεινά· καὶ ἀναμενεῖτε εἰς φῶς, καὶ ἐκεῖ σκιὰ θανάτου καὶ τεθήσονται εἰς σκότος.
\vs{17}Ἐὰν δὲ μὴ ἀκούσητε, κεκρυμμένως κλαύσεται ἡ ψυχὴ ὑμῶν ἀπὸ προσώπου ὕβρεως, καὶ κατάξουσιν οἱ ὀφθαλμοὶ ὑμῶν δάκρυα, ὅτι συνετρίβη τὸ ποίμνιον Κυρίου.

\vs{18}Εἴπατε τῷ βασιλεῖ καὶ τοῖς δυναστεύουσι, ταπεινώθητε καὶ καθίσατε, ὅτι καθῃρέθη ἀπὸ κεφαλῆς ὑμῶν στέφανος δόξης ὑμῶν.
\vs{19}Πόλεις αἱ πρὸς Νότον, συνεκλείσθησαν, καὶ οὐκ ἦν ὁ ἀνοίγων· ἀποικίσθη Ἰούδας, συνετέλεσαν ἀποικίαν τελείαν.
\vs{20}Ἀνάλαβε ὀφθαλμούς σου Ἱερουσαλὴμ, καὶ ἴδε τοὺς ἐρχομένους ἀπὸ Βοῤῥᾶ· ποῦ ἐστι τὸ ποίμνιον ὃ ἐδόθη σοι, πρόβατα δόξης σου;
\vs{21}Τί ἐρεῖς ὅταν ἐπισκέπτωνταί σε; καὶ σὺ ἐδίδαξας αὐτοὺς ἐπὶ σὲ μαθήματα εἰς ἀρχήν· οὐκ ὠδῖνες καθέξουσί σε καθὼς γυναῖκα τίκτουσαν;
\vs{22}Καὶ ἐὰν εἴπῃς ἐν τῇ καρδίᾳ σου, διατί ἀπήντησέ μοι ταῦτα; διὰ τὸ πλῆθος τῆς ἀδικίας σου ἀνεκαλύφθη τὰ ὀπίσθιά σου, παραδειγματισθῆναι τὰς πτέρνας σου.

\vs{23}Εἰ ἀλλάξεται Αἰθίοψ τὸ δέρμα αὐτοῦ, καὶ πάρδαλις τὰ ποικίλματα αὐτῆς, καὶ ὑμεῖς δυνήσεσθε εὐποιῆσαι μεμαθηκότες τὰ κακά.
\vs{24}Καὶ διέσπειρα αὐτοὺς, ὡς φρύγανα φερόμενα ὑπὸ ἀνέμου εἰς ἔρημον.
\vs{25}Οὕτως ὁ κλῆρός σου, καὶ μερὶς τοῦ ἀπειθεῖν ὑμᾶς ἐμοὶ, λέγει Κύριος· ὡς ἐπελάθου μου, καὶ ἤλπισας ἐπὶ ψεύδεσι,
\vs{26}κᾀγὼ ἀποκαλύψω τὰ ὀπίσω σου ἐπὶ τὸ πρόσωπόν σου, καὶ ὀφθήσεται ἡ ἀτιμία σου,
\vs{27}καὶ ἡ μοιχεία σου, καὶ χρεμετισμός σου, καὶ ἡ ἀπαλλοτρίωσις τῆς πορνείας σου· ἐπὶ τῶν βουνῶν, καὶ ἐν τοῖς ἀγροῖς ἑώρακα τὰ βδελύγματά σου· οὐαί σοι Ἱερουσαλὴμ, ὅτι οὐκ ἐκαθαρίσθης ὀπίσω μου· ἕως τίνος ἔτι;

\ch{14}
ΚΑΙ ἘΓΕΝΕΤΟ ΛΟΓΟΣ ΚΥΡΙΟΥ ΠΡΟΣ ἹΕΡΕΜΙΑΝ ΠΕΡΙ ΤΗΣ ἈΒΡΟΧΙΑΣ.

\vs{2}Ἐπένθησεν ἡ Ἰουδαία, καὶ αἱ πύλαι αὐτῆς ἐκενώθησαν, καὶ ἐσκοτώθησαν ἐπὶ τῆς γῆς, καὶ ἡ κραυγὴ τῆς Ἱερουσαλὴμ ἀνέβη·
\vs{3}Καὶ οἱ μεγιστᾶνες αὐτῆς ἀπέστειλαν τοὺς νεωτέρους αὐτῶν ἐφʼ ὕδωρ· ἤλθοσαν ἐπὶ τὰ φρέατα, καὶ οὐχ εὕροσαν ὕδωρ, καὶ ἀπέστρεψαν τὰ ἀγγεῖα αὐτῶν κενά.
\vs{4}Καὶ τὰ ἔργα τῆς γῆς ἐξέλιπεν, ὅτι οὐκ ἦν ὑετός· ᾐσχύνθησαν οἱ γεωργοὶ, ἐπεκάλυψαν τὰς κεφαλὰς αὐτῶν.
\vs{5}Καὶ ἔλαφοι ἐν ἀγρῷ ἔτεκον, καὶ ἐγκατέλιπον, ὅτι οὐκ ἦν βοτάνη.
\vs{6}Ὄνοι ἄγριοι ἔστησαν ἐπὶ νάπας, καὶ εἵλκυσαν ἄνεμον, ἐξέλιπον οἱ ὀφθαλμοὶ αὐτῶν, ὅτι οὐκ ἦν χόρτος.

\vs{7}Αἱ ἁμαρτίαι ἡμῶν ἀντέστησαν ἡμῖν· Κύριε ποίησον ἡμῖν ἕνεκέν σου, ὅτι πολλαὶ αἱ ἁμαρτίαι ἡμῶν ἐναντίον σου, ὅτι σοι ἡμάρτομεν.
\vs{8}Ὑπομονὴ Ἰσραὴλ Κύριε, καὶ σώζεις ἐν καιρῷ κακῶν· ἱνατί ἑγενήθης ὡσεὶ πάροικος ἐπὶ τῆς γῆς, καὶ ὡς αὐτόχθων ἐκκλίνων εἰς κατάλυμα;
\vs{9}Μὴ ἔσῃ ὥσπερ ἄνθρωπος ὑπνῶν, ἢ ὡς ἀνὴρ οὐ δυνάμενος σώζειν; καὶ σὺ ἐν ἡμῖν εἶ Κύριε, καὶ τὸ ὄνομά σου ἐπικέκληται ἐφʼ ἡμᾶς· μὴ ἐπιλάθῃ ἡμῶν.

\vs{10}Οὕτως λέγει Κύριος τῷ λαῷ τούτῳ, ἠγάπησαν κινεῖν πόδας αὐτῶν, καὶ οὐκ ἐφείσαντο, καὶ ὁ Θεὸς οὐκ εὐώδωσεν ἐν αὐτοῖς· νῦν μνησθήσεται τῆς ἀδικίας αὐτῶν.
\vs{11}Καὶ εἶπε Κύριος πρὸς μὲ, μὴ προσεύχου περὶ τοῦ λαοῦ τούτου εἰς ἀγαθὰ,
\vs{12}ὅτι ἐὰν νηστεύσωσιν, οὐκ εἰσακούσομαι τῆς δεήσεως αὐτῶν· καὶ ἐὰν προσενέγκωσιν ὁλοκαυτώματα καὶ θυσίας, οὐκ εὐδοκήσω ἐν αὐτοῖς, ὅτι ἐν μαχαίρᾳ καὶ ἐν λιμῷ καὶ ἐν θανάτῳ ἐγὼ συντελέσω αὐτούς.

\vs{13}Καὶ εἶπα, ὁ ὢν Κύριε, ἰδοὺ οἱ προφῆται αὐτῶν προφητεύουσι, καὶ λέγουσιν, οὐκ ὄψεσθε μάχαιραν, οὐδὲ λιμὸς ἔσται ἐν ὑμῖν, ὅτι ἀλήθειαν καὶ εἰρήνην δώσω ἐπὶ τῆς γῆς, καὶ ἐν τῷ τόπῳ τούτῳ.

\vs{14}Καὶ εἶπε Κύριος πρὸς μὲ, ψευδῆ οἱ προφῆται προφητεύουσιν ἐπὶ τῷ ὀνόματί μου, οὐκ ἀπέστειλα αὐτοὺς, καὶ οὐκ ἐνετειλάμην αὐτοῖς, καὶ οὐκ ἐλάλησα πρὸς αὐτούς· ὅτι ὁράσεις ψευδεῖς, καὶ μαντείας, καὶ οἰωνίσματα, καὶ προαιρέσεις καρδίας αὐτῶν αὐτοὶ προφητεύουσιν ὑμῖν.
\vs{15}Διατοῦτο τάδε λέγει Κύριος περὶ τῶν προφητῶν τῶν προφητευόντων ἐπὶ τῷ ὀνόματί μου ψευδῆ, καὶ ἐγὼ οὐκ ἀπέστειλα αὐτοὺς, οἳ λέγουσι, μάχαιρα καὶ λιμὸς οὐκ ἔσται ἐπὶ τῆς γῆς ταύτης· ἐν θανάτῳ νοσερῷ ἀποθανοῦνται, καὶ ἐν λιμῷ συντελεσθήσονται οἱ προφῆται,
\vs{16}καὶ ὁ λαὸς, οἷς αὐτοὶ προφητεύουσιν αὐτοῖς, καὶ ἔσονται ἐῤῥιμμένοι ἐν ταῖς ὁδοῖς Ἱερουσαλὴμ, ἀπὸ προσώπου μαχαίρας καὶ τοῦ λιμοῦ, καὶ οὐκ ἔσται ὁ θάπτων αὐτοὺς, καὶ αἱ γυναῖκες αὐτῶν, καὶ οἱ υἱοὶ αὐτῶν, καὶ αἱ θυγατέρες αὐτῶν, καὶ ἐκχεῶ ἐπʼ αὐτοὺς τὰ κακὰ αὐτῶν.

\vs{17}Καὶ ἐρεῖς πρὸς αὐτοὺς τὸν λόγον τοῦτον, καταγάγετε ἐπ ὀφθαλμοὺς ὑμῶν δάκρυα ἡμερας καὶ νυκτὸς, καὶ μὴ διαλιπέτωσαν, ὅτι συντρίμματι συνετρίβη θυγάτηρ λαοῦ μου, καὶ πληγὴ ὀδυνηρὰ σφόδρα.
\vs{18}Ἐὰν ἐξέλθω εἰς τὸ πεδίον, καὶ ἰδοὺ τραυματίαι μαχαίρας· καὶ ἐὰν εἰσέλθω εἰς τὴν πόλιν, καὶ ἰδοὺ πόνος λιμοῦ· ὅτι ἱερεὺς καὶ προφήτης ἐπορεύθησαν εἰς γῆν, ἣν οὐκ ᾔδεισαν.

\vs{19}Μὴ ἀποδοκιμάζων ἀπεδοκίμασας τὸν Ἰούδαν, καὶ ἀπὸ Σιὼν ἀπέστη ἡ ψυχή σου; ἱνατί ἔπαισας ἡμᾶς, καὶ οὐκ ἔστιν ἡμῖν ἴασις; ὑπεμείναμεν εἰς εἰρήνην, καὶ οὐκ ἦν ἀγαθὰ, εἰς καιρὸν ἰάσεως, καὶ ἰδοὺ ταραχή.
\vs{20}Ἔγνωμεν Κύριε ἁμαρτήματα ἡμῶν, ἀδικίας πατέρων ἡμῶν, ὅτι ἡμάρτομεν ἐναντίον σου.
\vs{21}Κόπασον διὰ τὸ ὄνομά σου, μὴ ἀπολέσῃς θρόνον δόξης σου· μνήσθητι, μὴ διασκεδάσῃς τὴν διαθήκην σου τὴν μεθʼ ἡμῶν.
\vs{22}Μὴ ἔστιν ἐν εἰδώλοις τῶν ἐθνῶν ὑετίζων; καὶ εἰ ὁ οὐρανὸς δώσει πλησμονὴν αὐτοῦ, οὐχὶ σὺ εἶ αὐτός; καὶ ὑπομενοῦμέν σε Κύριε, ὅτι σὺ ἐποίησας πάντα ταῦτα.

\ch{15}
Καὶ εἶπε Κύριος πρὸς μὲ, ἐὰν στῇ Μωυσῆς καὶ Σαμουὴλ πρὸ προσώπου μου, οὐκ ἔστιν ἡ ψυχή μου πρὸς αὐτούς· ἐξαπόστειλον τὸν λαὸν τοῦτον, καὶ ἐξελθέτωσαν.
\vs{2}Καὶ ἔσται ἐὰν εἴπωσι πρὸς σὲ, ποῦ ἐξελευσόμεθα; καὶ ἐρεῖς πρὸς αὐτοὺς, τάδε λέγει Κύριος, ὅσοι εἰς θάνατον, εἰς θάνατον· καὶ ὅσοι εἰς μάχαιραν, εἰς μάχαιραν· καὶ ὅσοι εἰς λιμὸν, εἰς λιμόν· καὶ ὅσοι εἰς αἰχμαλωσίαν, εἰς αἰχμαλωσίαν.
\vs{3}Καὶ ἐκδικήσω ἐπʼ αὐτοὺς τέσσαρα εἴδη, λέγει Κύριος· τὴν μάχαιραν εἰς σφαγὴν, καὶ τοὺς κύνας εἰς διασπασμὸν, καὶ τὰ θηρία τῆς γῆς, καὶ τὰ πετεινὰ τοῦ οὐρανοῦ εἰς βρῶσιν καὶ διαφθοράν.
\vs{4}Καὶ παραδώσω αὐτοὺς εἰς ἀνάγκας πάσαις ταῖς βασιλείαις τῆς γῆς, διὰ Μανασσῆ υἱὸν Ἐζεκίου βασιλέως Ἰούδα, περὶ πάντων ὧν ἐποίησεν ἐν Ἱερουσαλήμ.

\vs{5}Τίς φείσεται ἐπὶ σοὶ Ἱερουσαλήμ; καὶ τίς δειλιάσει ἐπὶ σοί; ἢ τίς ἀνακάμψει εἰς εἰρήνην σοι;
\vs{6}Σὺ ἀπεστράφης με, λέγει Κύριος, ὀπίσω πορεύσῃ· καὶ ἐκτενῶ τὴν χεῖρά μου, καὶ διαφθερῶ σε, καὶ οὐκέτι ἀνήσω αὐτούς·
\vs{7}Καὶ διασπερῶ αὐτοὺς ἐν διασπορᾶ ἐν πύλαις λαοῦ μου ἠτεκνώθησαν, ἀπώλεσαν τὸν λαόν μου διὰ τὰς κακίας αὐτῶν.

\vs{8}Ἐπληθύνθησαν αἱ χῆραι αὐτῶν ὑπὲρ τὴν ἄμμον τῆς θαλάσσης· ἐπήγαγον ἐπὶ μητέρα νεανίσκους, ταλαιπωρίαν ἐν μεσημβρίᾳ, ἐπέῤῥιψα ἐπʼ αὐτὴν ἐξαίφνης τρόμον καὶ σπουδήν.
\vs{9}Ἐκενώθη ἡ τίκτουσα ἑπτὰ, ἀπεκάκησεν ἡ ψυχὴ αὐτῆς, ἐπέδυ ὁ ἥλιος αὐτῇ ἔτι μεσούσης τῆς ἡμέρας, κατῃσχύνθη καὶ ὠνειδίσθη· τοὺς καταλοίπους αὐτῶν εἰς μάχαιραν δώσω ἐναντίον τῶν ἐχθρῶν αὐτῶν.

\vs{10}Οἴμοι ἐγὼ, μῆτερ, ὡς τινά με ἔτεκες ἄνδρα δικαζόμενον, καὶ διακρινόμενον πάσῃ τῇ γῇ· οὔτε ὠφέλησα, οὔτε ὠφέλησέ με οὐδείς· ἡ ἰσχύς μου ἐξέλιπεν ἐν τοῖς καταρωμένοις με.
\vs{11}Γένοιτο δέσποτα κατευθυνόντων αὐτῶν· εἰ μὴ παρέστην σοι ἐν καιρῷ τῶν κακῶν αὐτῶν, καὶ ἐν καιρῷ θλίψεως αὐτῶν, εἰς ἀγαθὰ πρὸς τὸν ἐχθρόν.
\vs{12}Εἰ γνωσθήσεται σίδηρος; καὶ περιβόλαιον χαλκοῦν ἡ ἰσχύς σου.
\vs{13}Καὶ τοὺς θησαυρούς σου εἰς προνομὴν δώσω ἀντάλλαγμα, διὰ πάσας τὰς ἁμαρτίας σου, καὶ ἐν πᾶσι τοῖς ὁρίοις σου.
\vs{14}Καὶ καταδουλώσω σε κύκλῳ τοῖς ἐχθροῖς σου, ἐν τῇ γῇ ᾗ οὐκ ᾔδεις· ὅτι πῦρ ἐκκέκαυται ἐκ τοῦ θυμοῦ μου, ἐφʼ ὑμᾶς καυθήσεται.

\vs{15}Κύριε μνήσθητί μου, καὶ ἐπίσκεψαί με, καὶ ἀθώωσόν με ἀπὸ τῶν καταδιωκόντων με, μὴ εἰς μακροθυμίαν· γνῶθι, ὡς ἔλαβον περὶ σοῦ ὀνειδισμὸν
\vs{16}ὑπὸ τῶν ἀθετούντων τοὺς λόγους σου· συντέλεσον αὐτούς· καὶ ἔσται ὁ λόγος σου ἐμοὶ εἰς εὐφροσύνην καὶ χαρὰν καρδίας μου, ὅτι ἐπικέκληται τὸ ὄνομά σου ἐπʼ ἐμοὶ, Κύριε παντοκράτωρ.
\vs{17}Οὐκ ἐκάθισα ἐν συνεδρίῳ αὐτῶν παιζόντων, ἀλλὰ εὐλαβούμην ἀπὸ προσώπου χειρός σου· καταμόνας ἐκαθήμην, ὅτι πικρίας ἐνεπλήσθην.

\vs{18}Ἱνατί οἱ λυποῦντές με κατισχύουσί μου; ἡ πληγή μου στερεὰ, πόθεν ἰαθήσομαι; γινομένη ἐγενήθη μοι ὡς ὕδωρ ψευδὲς, οὐκ ἔχον πίστιν.

\vs{19}Διατοῦτο τάδε λέγει Κύριος, ἐὰν ἐπιστρέψῃς, καὶ ἀποκαταστήσω σε, καὶ πρὸ προσώπου μου στήσῃ· καὶ ἐὰν ἐξαγάγῃς τίμιον ἀπὸ ἀναξίου, ὡς τὸ στόμα μου ἔσῃ· καὶ ἀναστρέψουσιν αὐτοὶ πρὸς σὲ, καὶ σὺ οὐκ ἀναστρέψεις πρὸς αὐτούς.
\vs{20}Καὶ δώσω σε τῷ λαῷ τούτῳ, ὡς τεῖχος ὀχυρὸν, χαλκοῦν· καὶ πολεμήσουσι πρὸς σὲ, καὶ οὐ μὴ δύνωνται πρὸς σὲ, διότι μετὰ σοῦ εἰμι τοῦ σώζειν σε,
\vs{21}καὶ τοῦ ἐξαιρεῖσθαί σε ἐκ χειρὸς πονηρῶν, καὶ λυτρώσομαί σε ἐκ χειρὸς λοιμῶν.

\ch{16}
Καὶ σὺ μὴ λάβῃς γυναῖκα, λέγει Κύριος ὁ Θεὸς Ἰσραὴλ,
\vs{2}καὶ οὐ γεννηθήσεταί σοι υἱὸς, οὐδὲ θυγάτηρ ἐν τῷ τόπῳ τούτῳ.
\vs{3}Ὅτι τάδε λέγει Κύριος περὶ τῶν υἱῶν καὶ περὶ τῶν θυγατέρων τῶν γεννωμένων ἐν τῷ τόπῳ τούτῳ, καὶ περὶ τῶν μητέρων αὐτῶν τῶν τετοκυιῶν αὐτοὺς, καὶ περὶ τῶν πατέρων αὐτῶν τῶν γεγεννηκότων αὐτοὺς ἐν τῇ γῇ ταύτῃ,
\vs{4}ἐν θανάτῳ νοσερῷ ἀποθανοῦνται, οὐ κοπήσονται, καὶ οὐ ταφήσονται· εἰς παράδειγμα ἐπὶ προσώπου τῆς γῆς ἔσονται· καὶ τοῖς θηρίοις τῆς γῆς ἔσονται καὶ τοῖς πετεινοῖς τοῦ οὐρανοῦ· ἐν μαχαίρᾳ πεσοῦνται, καὶ ἐν λιμῷ συντελεσθήσονται.

\vs{5}Τάδε λέγει Κύριος, μὴ εἰσέλθῃς εἰς θίασον αὐτῶν, καὶ μὴ πορευθῇς τοῦ κόψασθαι, καὶ μὴ πενθήσῃς αὐτοὺς, ὅτι ἀφέστακα τὴν εἰρήνην μου ἀπὸ τοῦ λαοῦ τούτου·
\vs{6}Οὐ μὴ κόψονται αὐτοὺς, οὐδὲ ἐντομίδας οὐ μὴ ποιήσουσι, καὶ οὐ ξυρηθήσονται,
\vs{7}καὶ οὐ μὴ κλασθῇ ἄρτος ἐν πένθει αὐτῶν εἰς παράκλησιν ἐπὶ τεθνηκότι· οὐ ποτιοῦσιν αὐτὸν ποτήριον εἰς παράκλησιν ἐπὶ πατρὶ καὶ μητρὶ αὐτοῦ.

\vs{8}Εἰς οἰκίαν πότου οὐκ εἰσελεύσῃ σὺ, συγκαθίσαι μετʼ αὐτῶν τοῦ φαγεῖν καὶ πιεῖν.
\vs{9}Διότι τάδε λέγει Κύριος ὁ Θεὸς Ἰσραὴλ, ἰδοὺ ἐγὼ καταλύω ἐκ τοῦ τόπου τούτου ἐνώπιον τῶν ὀφθαλμῶν ὑμῶν, καὶ ἐν ταῖς ἡμέραις ὑμῶν, φωνὴν χαρᾶς, καὶ φωνὴν εὐφροσύνης, φωνὴν νυμφίου, καὶ φωνὴν νύμφης.

\vs{10}Καὶ ἔσται ὅταν ἀναγγείλῃς τῷ λαῷ τούτῳ ἅπαντα τὰ ῥήματα ταῦτα, καὶ εἴπωσι πρὸς σὲ, διατί ἐλάλησε Κύριος ἐφʼ ἡμᾶς πάντα τὰ κακὰ ταῦτα; τίς ἡ ἀδικία ἡμῶν; καὶ τίς ἡ ἁμαρτία ἡμῶν, ἣν ἡμάρτομεν ἔναντι Κυρίου τοῦ Θεοῦ ἡμῶν;
\vs{11}Καὶ ἐρεῖς πρὸς αὐτοὺς, ἀνθʼ ὧν ἐγκατέλιπόν με οἱ πατέρες ὑμῶν, λέγει Κύριος, καὶ ᾤχοντο ὀπισω θεῶν ἀλλοτρίων, καὶ ἐδούλευσαν αὐτοῖς, καὶ προσεκύνησαν αὐτοῖς, καὶ ἐμὲ ἐγκατέλιπον, καὶ τὸν νόμον μου οὐκ ἐφυλάξαντο·
\vs{12}Καὶ ὑμεῖς ἐπονηρεύσασθε ὑπὲρ τοὺς πατέρας ὑμῶν· καὶ ἰδοὺ ὑμεῖς πορεύεσθε ἕκαστος ὀπίσω τῶν ἀρεστῶν τῆς καρδίας ὑμῶν τῆς πονηρᾶς, τοῦ μὴ ὑπακούειν μου.
\vs{13}Καὶ ἀποῤῥίψω ὑμᾶς ἀπὸ τῆς γῆς ταύτης εἰς τὴν γῆν, ἣν οὐκ ᾔδειτε ὑμεῖς καὶ οἱ πατέρες ὑμῶν, καὶ δουλεύσετε ἐκεῖ θεοῖς ἑτέροις, οἳ οὐ δώσουσιν ὑμῖν ἔλεος.

\vs{14}Διατοῦτο ἰδοὺ ἡμέραι ἔρχονται, λέγει Κύριος, καὶ οὐκ ἐροῦσιν ἔτι, ζῇ Κύριος ὁ ἀναγαγὼν τοὺς υἱοὺς Ἰσραὴλ ἐκ γῆς Αἰγύπτου,
\vs{15}ἀλλὰ, ζῇ Κύριος, ὃς ἀνήγαγε τὸν οἶκον Ἰσραὴλ ἀπὸ γῆς Βοῤῥᾶ, καὶ ἀπὸ πασῶν τῶν χωρῶν οὗ ἐξώσθησαν ἐκεῖ· καὶ ἀποκαταστήσω αὐτοὺς εἰς τὴν γῆν αὐτῶν, ἣν ἔδωκα τοῖς πατράσιν αὐτῶν.

\vs{16}Ἰδοὺ ἐγὼ ἀποστέλλω τοὺς ἁλιεῖς τοὺς πολλοὺς, λέγει Κύριος, καὶ ἁλιευσουσιν αὐτούς· καὶ μετὰ ταῦτα ἀποστέλλω τοὺς πολλοὺς θηρευτὰς, καὶ θηρεύσουσιν αὐτοὺς ἐπάνω παντὸς ὄρους, καὶ ἐπάνω παντὸς βουνοῦ, καὶ ἐκ τῶν τρυμαλιῶν τῶν πετρῶν.
\vs{17}Ὅτι οἱ ὀφθαλμοί μου ἐπὶ πάσας τὰς ὁδοὺς αὐτῶν, καὶ οὐκ ἐκρύβη τὰ ἀδικήματα αὐτῶν ἀπέναντι τῶν ὀφθαλμῶν μου.
\vs{18}Καὶ ἀνταποδώσω διπλὰς τὰς κακίας αὐτῶν, καὶ τὰς ἁμαρτίας αὐτῶν, ἐφʼ αἷς ἐβεβήλωσαν τὴν γῆν μου ἐν τοῖς θνησιμαίοις τῶν βδελυγμάτων αὐτῶν, καὶ ἐν ταῖς ἀνομίαις αὐτῶν, ἐν αἷς ἐπλημμέλησαν τὴν κληρονομίαν μου.

\vs{19}Κύριε σὺ ἰσχύς μου, καὶ βοήθειά μου, καὶ καταφυγή μου ἐν ἡμέραις κακῶν· πρὸς σὲ ἔθνη ἥξουσιν ἀπʼ ἐσχάτου τῆς γῆς, καὶ ἐροῦσιν, ὡς ψευδῆ ἐκτήσαντο οἱ πατέρες ἡμῶν εἴδωλα, καὶ οὐκ ἔστιν ἐν αὐτοῖς ὠφέλημα;
\vs{20}Εἰ ποιήσει ἑαυτῷ ἄνθρωπος θεοὺς, καὶ οὗτοι οὐκ εἰσὶ θεοί.
\vs{21}Διατοῦτο ἰδοὺ ἐγὼ δηλώσω αὐτοῖς ἐν τῷ καιρῷ τούτῳ τὴν χεῖρά μου, καὶ γνωριῶ αὐτοῖς τὴν δύναμίν μου, καὶ γνώσονται, ὅτι ὄνομά μοι Κύριος.

\ch{17}
\vs{5}Ἐπικατάρατος ὁ ἄνθρωπος, ὃς τὴν ἐλπίδα ἔχει ἐπʼ ἄνθρωπον, καὶ στηρίσει σάρκα βραχίονος αὐτοῦ ἐπʼ αὐτὸν, καὶ ἀπὸ Κυρίου ἀποστῇ ἡ καρδία αὐτοῦ.
\vs{6}Καὶ ἔσται ὡς ἡ ἀγριομυρίκη ἡ ἐν τῇ ἐρήμῳ, οὐκ ὄψεται ὅταν ἔλθῃ τὰ ἀγαθὰ, καὶ κατασκηνώσει ἐν ἁλίμοις, καὶ ἐν ἐρήμῳ, ἐν γῇ ἁλμυρᾷ ἥτις οὐ κατοικεῖται.
\vs{7}Καὶ εὐλογημένος ὁ ἄνθρωπος, ὃς πέποιθεν ἐπὶ τῷ Κυρίῳ, καὶ ἔσται Κύριος ἐλπὶς αὐτοῦ·
\vs{8}Καὶ ἔσται ὡς ξύλον εὐθηνοῦν παρʼ ὕδατα, καὶ ἐπὶ ἰκμάδα βαλεῖ ῥίζαν αὐτοῦ· οὐ φοβηθήσεται ὅταν ἔλθῃ καῦμα, καὶ ἔσται ἐπʼ αὐτῷ στελέχη ἀλσώδη, ἐν ἐνιαυτῷ ἀβροχίας οὐ φοβηθήσεται, καὶ οὐ διαλείψει ποιῶν καρπόν.

\vs{9}Βαθεία ἡ καρδία παρὰ πάντα, καὶ ἄνθρωπός ἐστι, καὶ τίς γνώσεται αὐτόν;
\vs{10}Ἐγὼ Κύριος ἐτάζων καρδίας, καὶ δοκιμάζων νεφροὺς, τοῦ δοῦναι ἑκάστῳ κατὰ τὰς ὁδοὺς αὐτοῦ, καὶ κατὰ τοὺς καρποὺς τῶν ἐπιτηδευμάτων αὐτοῦ.

\vs{11}Ἐφώνησε πέρδιξ, συνήγαγεν ἃ οὐκ ἔτεκε, ποιῶν πλοῦτον αὐτοῦ οὐ μετὰ κρίσεως· ἐν ἡμίσει ἡμερῶν αὐτοῦ ἐγκαταλείψουσιν αὐτὸν, καὶ ἐπʼ ἐσχάτων αὐτοῦ ἔσται ἄφρων.

\vs{12}Θρόνος δόξης ὑψωμένος, ἁγίασμα ἡμῶν,
\vs{13}ὑπομονὴ Ἰσραήλ· Κύριε, πάντες οἱ καταλιπόντες σε καταισχυνθήτωσαν, ἀφεστηκότες ἐπὶ τῆς γῆς γραφήτωσαν, ὅτι ἐγκατέλιπον πηγὴν ζωῆς, τὸν Κύριον.

\vs{14}Ἴασαί με Κύριε, καὶ ἰαθήσομαι· σῶσόν με, καὶ σωθήσομαι, ὅτι καύχημά μου σὺ εἶ.

\vs{15}Ἰδοὺ αὐτοὶ λέγουσι πρὸς μὲ, ποῦ ἐστιν ὁ λόγος Κυρίου; ἐλθέτω.
\vs{16}Ἐγὼ δὲ οὐκ ἐκοπίασα κατακολουθῶν ὀπίσω σου, καὶ ἡμέραν ἀνθρώπου οὐκ ἐπεθύμησα· σὺ ἐπίστῃ· τὰ ἐκπορευόμενα διὰ τῶν χειλέων μου, πρὸ προσώπου σου ἐστί.
\vs{17}Μὴ γενηθῇς μοι εἰς ἀλλοτρίωσιν, φειδόμενός μου ἐν ἡμέρᾳ πονηρᾷ.
\vs{18}Καταισχυνθήτωσαν οἱ διώκοντές με, καὶ μὴ καταισχυνθείην ἐγώ· πτοηθείησαν αὐτοὶ, καὶ μὴ πτοηθείην ἐγώ· ἐπάγαγε ἐπʼ αὐτοὺς ἡμέραν πονηρὰν, δισσὸν σύντριμμα σύντριψον αὐτούς.

\vs{19}Τάδε λέγει Κύριος, βάδισον, καὶ στῆθι ἐν ταῖς πύλαις υἱῶν λαοῦ σου, ἐν αἷς εἰσπορεύονται ἐν αὐταῖς βασιλεῖς Ἰούδα, καὶ ἐν αἷς ἐκπορεύονται ἐν αὐταῖς, καὶ ἐν πάσαις ταῖς πύλαις Ἱερουσαλὴμ,
\vs{20}καὶ ἐρεῖς αὐτοῖς, ἀκούσατε τὸν λόγον Κυρίου βασιλεῖς Ἰούδα, καὶ πᾶσα Ἰουδαία, καὶ πᾶσα Ἱερουσαλὴμ, οἱ εἰσπορευόμενοι ἐν ταῖς πύλαις ταύταις·
\vs{21}Τάδε λέγει Κύριος, φυλάσσεσθε τὰς ψυχὰς ὑμῶν, καὶ μὴ αἴρετε βαστάγματα ἐν τῇ ἡμέρᾳ τῶν σαββάτων, καὶ μὴ ἐκπορεύεσθε ταῖς πύλαις Ἱερουσαλὴμ,
\vs{22}καὶ μὴ ἐκφέρετε βαστάγματα ἐξ οἰκιῶν ὑμῶν ἐν τῇ ἡμέρᾳ τῶν σαββάτων, καὶ πᾶν ἔργον οὐ ποιήσετε· ἁγιάσατε τὴν ἡμέραν τῶν σαββάτων, καθὼς ἐνετειλάμην τοῖς πατράσιν ὑμῶν. Καὶ οὐκ ἤκουσαν, καὶ οὐκ ἔκλιναν τὸ οὖς αὐτῶν,
\vs{23}καὶ ἐσκλήρυναν τὸν τράχηλον αὐτῶν ὑπὲρ τοὺς πατέρας αὐτῶν, τοῦ μὴ ἀκοῦσαί μου, καὶ τοῦ μὴ δέξασθαι παιδείαν.

\vs{24}Καὶ ἓσται, ἐὰν εἰσακούσητέ μου, λέγει Κύριος, τοῦ μὴ εἰσφέρειν βαστάγματα διὰ τῶν πυλῶν τῆς πόλεως ταύτης ἐν τῇ ἡμέρᾳ τῶν σαββάτων καὶ ἁγιάζειν τὴν ἡμέραν τῶν σαββάτων, τοῦ μὴ ποιεῖν πᾶν ἔργον,
\vs{25}καὶ εἰσελεύσονται διὰ τῶν πυλῶν τῆς πόλεως ταύτης βασιλεῖς καὶ ἄρχοντες καθήμενοι ἐπὶ θρόνου Δαυὶδ, καὶ ἐπιβεβηκότες ἐφʼ ἅρμασι καὶ ἵπποις αὐτῶν, αὐτοὶ καὶ οἱ ἄρχοντες αὐτῶν, ἄνδρες Ἰούδα καὶ οἱ κατοικοῦντες ἐν Ἱερουσαλήμ· καὶ κατοικισθήσεται ἡ πόλις αὕτη εἰς τὸν αἰῶνα.
\vs{26}Καὶ ἥξουσιν ἐκ τῶν πόλεων Ἰούδα, καὶ κυκλόθεν Ἱερουσαλὴμ, καὶ ἐκ γῆς Βενιαμὶν, καὶ ἐκ γῆς πεδινῆς, καὶ ἐκ τοῦ ὄρους, καὶ ἐκ τῆς πρὸς Νότον, φέροντες ὁλοκαυτώματα καὶ θυσίας καὶ θυμιάματα καὶ μάννα καὶ λίβανον, φέροντες αἴνεσιν εἰς οἶκον Κυρίου.

\vs{27}Καὶ ἔσται ἐὰν μὴ ἀκούσητέ μου τοῦ ἁγιάζειν τὴν ἡμέραν τῶν σαββάτων, τοῦ μὴ αἴρειν βαστάγματα, καὶ μὴ εἰσπορεύεσθαι ταῖς πύλαις Ἱερουσαλὴμ ἐν τῇ ἡμέρᾳ τῶν σαββάτων, καὶ ἀνάψω πῦρ ἐν ταῖς πύλαις αὐτῆς, καὶ καταφάγεται ἄμφοδα Ἱερουσαλὴμ, καὶ οὐ σβεσθήσεται.

\ch{18}
Ὁ λόγος ὁ γενόμενος παρὰ Κυρίου πρὸς Ἱερεμίαν, λέγων,
\vs{2}ἀνάστηθι, καὶ κατάβηθι εἰς οἶκον τοῦ κεραμέως, καὶ ἐκεῖ ἀκούσῃ τοὺς λόγους μου.
\vs{3}Καὶ κατέβην εἰς τὸν οἶκον τοῦ κεραμέως, καὶ ἰδοὺ αὐτὸς ἐποίει ἔργον ἐπὶ τῶν λίθων.
\vs{4}Καὶ ἔπεσε τὸ ἀγγεῖον, ὃ αὐτὸς ἐποίει ἐν ταῖς χερσὶν αὐτοῦ· καὶ πάλιν αὐτὸς ἐποίησεν αὐτὸ ἀγγεῖον ἕτερον, καθὼς ἤρεσεν ἐνώπιον αὐτοῦ ποιῆσαι.
\vs{5}Καὶ ἐγένετο λόγος Κυρίου πρὸς μὲ, λέγων,

\vs{6}Εἰ καθὼς ὁ κεραμεὺς οὗτος οὐ δυνήσομαι τοῦ ποιῆσαι ὑμᾶς οἶκος Ἰσραήλ; ἰδοὺ, ὡς ὁ πηλὸς τοῦ κεραμέως, ὑμεῖς ἐστε ἐν χερσί μου.
\vs{7}Πέρας λαλήσω ἐπὶ ἔθνος, ἢ ἐπὶ βασιλείαν τοῦ ἐξᾶραι αὐτοὺς, καὶ τοῦ ἀπολλύειν,
\vs{8}καὶ ἐπιστραφῇ τὸ ἔθνος ἐκεῖνο ἀπὸ πάντων τῶν κακῶν αὐτῶν, καὶ μετανοήσω περὶ τῶν κακῶν, ὧν ἐλογισάμην, τοῦ ποιῆσαι αὐτοῖς.
\vs{9}Καὶ πέρας λαλήσω ἐπὶ ἔθνος καὶ βασιλείαν, τοῦ ἀνοικοδομεῖσθαι καὶ τοῦ καταφυτεύεσθαι,
\vs{10}καὶ ποιήσωσι τὰ πονηρὰ ἐναντίον μου, τοῦ μὴ ἀκούειν τῆς φωνῆς μου, καὶ μετανοήσω περὶ τῶν ἀγαθῶν ὧν ἐλάλησα, τοῦ ποιῆσαι αὐτοῖς.

\vs{11}Καὶ νῦν εἰπὸν πρὸς ἄνδρας Ἰούδα, καὶ πρὸς τοὺς κατοικοῦντας Ἱερουσαλὴμ, ἰδοὺ ἐγὼ πλάσσω ἐφʼ ὑμᾶς κακὰ, καὶ λογίζομαι ἐφʼ ὑμᾶς λογισμόν· ἀποστραφήτω δὴ ἕκαστος ἀπὸ ὁδοῦ αὐτοῦ τῆς πονηρᾶς, καὶ καλλίονα ποιήσατε τὰ ἐπιτηδεύματα ὑμῶν.
\vs{12}Καὶ εἶπαν, ἀνδριούμεθα, ὅτι ὀπίσω τῶν ἀποστροφῶν ἡμῶν πορευσόμεθα, καὶ ἕκαστος τὰ ἀρεστὰ τῆς καρδίας αὐτοῦ τῆς πονηρᾶς ποιήσομεν.

\vs{13}Διατοῦτο τάδε λέγει Κύριος, ἐρωτήσατε δὴ ἐν ἔθνεσι, τίς ἤκουσε τοιαῦτα φρικτὰ ἃ ἐποίησε σφόδρα παρθένος Ἰσραήλ;
\vs{14}Μὴ ἐκλείψουσιν ἀπὸ πέτρας μαστοὶ, ἢ χιὼν ἀπὸ τοῦ Λιβάνου; μὴ ἐκκλίνῃ ὕδωρ βιαίως ἀνέμῳ φερόμενον;
\vs{15}Ὅτι ἐπελάθοντό μου λαός μου, εἰς κενὸν ἐθυμίασαν καὶ ἀσθενήσουσιν ἐν ταῖς ὁδοῖς αὐτῶν σχοίνους αἰωνίους, τοῦ ἐπιβῆναι τρίβους οὐκ ἔχοντας ὁδὸν εἰς πορείαν,
\vs{16}τοῦ τάξαι τὴν γῆν αὐτῶν εἰς ἀφανισμὸν, καὶ σύριγμα αἰώνιον· πάντες οἱ διαπορευόμενοι διʼ αὐτῆς ἐκστήσονται, καὶ κινήσουσι τὴν κεφαλὴν αὐτῶν.
\vs{17}Ὡς ἄνεμον καύσωνα διασπερῶ αὐτοὺς κατὰ πρόσωπον ἐχθρῶν αὐτῶν, δείξω αὐτοῖς ἡμέραν ἀπωλείας αὐτῶν.

\vs{18}Καὶ εἶπεν, δεῦτε καὶ λογισώμεθα ἐπὶ Ἱερεμίαν λογισμὸν, ὅτι οὐκ ἀπολεῖται νόμος ἀπὸ ἱερέως, καὶ βουλὴ ἀπὸ συνετοῦ, καὶ λόγος ἀπὸ προφήτου· δεῦτε καὶ πατάξωμεν αὐτὸν ἐν γλώσσῃ, καὶ ἀκουσόμεθα πάντας τοὺς λόγους αὐτοῦ.

\vs{19}Εἰσάκουσόν μου Κύριε, καὶ εἰσάκουσον τῆς φωνῆς τοῦ δικαιώματός μου.
\vs{20}Εἰ ἀνταποδίδοται ἀντὶ ἀγαθῶν κακὰ, ὅτι συνελάλησαν ῥήματα κατὰ τῆς ψυχῆς μου, καὶ τὴν κόλασιν αὐτῶν ἔκρυψάν μοι· μνήσθητι ἑστηκότος μου κατὰ πρόσωπόν σου, τοῦ λαλῆσαι ὑπὲρ αὐτῶν ἀγαθὰ, τοῦ ἀποστρέψαι τὸν θυμόν σου ἀπʼ αὐτῶν.
\vs{21}Διατοῦτο δὸς τοὺς υἱοὺς αὐτῶν εἰς λιμὸν, καὶ ἄθροισον αὐτοὺς εἰς χεῖρας μαχαίρας· γενέσθωσαν αἱ γυναῖκες αὐτῶν ἄτεκνοι καὶ χῆραι, καὶ οἱ ἄνδρες αὐτῶν γενέσθωσαν ἀνῃρημένοι θανάτῳ, καὶ οἱ νεανίσκοι αὐτῶν πεπτωκότες μαχαίρᾳ ἐν πολέμῳ.
\vs{22}Γενηθήτω κραυγὴ ἐν ταῖς οἰκίαις αὐτῶν· ἐπάξεις ἐπʼ αὐτοὺς λῃστὰς ἄφνω, ὅτι ἐνεχείρησαν λόγον εἰς σύλληψίν μου, καὶ παγίδας ἔκρυψαν ἐπʼ ἐμέ.

\vs{23}Καὶ σὺ, Κύριε, ἔγνως ἅπασαν τὴν βουλὴν αὐτῶν ἐπʼ ἐμὲ εἰς θάνατον· μὴ ἀθωώσῃς τὰς ἀδικίας αὐτῶν, καὶ τὰς ἁμαρτίας αὐτῶν ἀπὸ προσώπου σου μὴ ἐξαλείψῃς· γενέσθω ἡ ἀσθένεια αὐτῶν ἐναντίον σου, ἐν καιρῷ θυμοῦ σου ποίησον ἐν αὐτοῖς.

\ch{19}
Τότε εἶπε Κύριος πρὸς μὲ, βάδισον, καὶ κτῆσαι βικὸν πεπλασμένον ὀστράκινον, καὶ ἄξεις ἀπὸ τῶν πρεσβυτέρων τοῦ λαοῦ καὶ ἀπὸ τῶν ἱερέων,
\vs{2}καὶ ἐξελεύσῃ εἰς τὸ πολυάνδριον υἱῶν τῶν τέκνων αὐτῶν, ὅ ἐστιν ἐπὶ τῶν προθύρων πύλης τῆς Χαρσείθ· καὶ ἀνάγνωθι ἐκεῖ πάντας τοὺς λόγους τούτους, οὓς ἂν λαλήσω πρὸς σὲ,
\vs{3}καὶ ἐρεῖς αὐτοῖς,

Ἀκούσατε τὸν λόγον Κυρίου, βασιλεῖς Ἰούδα, καὶ ἄνδρες Ἰούδα, καὶ οἱ κατοικοῦντες ἐν Ἱερουσαλὴμ καὶ οἱ εἰσπορευόμενοι ἐν ταῖς πύλαις ταύταις· τάδε λέγει Κύριος ὁ Θεὸς Ἰσραὴλ, ἰδοὺ ἐγὼ ἐπάγω ἐπὶ τὸν τόπον τοῦτον κακὰ, ὥστε παντὸς ἀκούοντος αὐτὰ ἠχήσει τὰ ὦτα αὐτοῦ·
\vs{4}Ἀνθʼ ὧν ἐγκατέλιπόν με, καὶ ἀπηλλοτρίωσαν τὸν τόπον τοῦτον, καὶ ἐθυμίασαν ἐν αὐτῷ θεοῖς ἀλλοτρίοις, οἷς οὐκ ᾔδεισαν αὐτοὶ καὶ οἱ πατέρες αὐτῶν· καὶ οἱ βασιλεῖς Ἰούδα ἔπλησαν τὸν τόπον τοῦτον αἱμάτων ἀθώων,
\vs{5}καὶ ᾠκοδόμησαν ὑψηλὰ τῇ Βάαλ, τοῦ κατακαίειν τοὺς υἱοὺς αὐτῶν ἐν πυρὶ, ἃ οὐκ ἐνετειλάμην, οὐδὲ διενοήθην ἐν τῇ καρδίᾳ μου.

\vs{6}Διατοῦτο ἰδοὺ ἡμέραι ἔρχονται, λέγει Κύριος, καὶ οὐ κληθήσεται τῷ τόπῳ τούτῳ ἔτι Διάπτωσις καὶ Πολυάνδριον υἱοῦ Ἐννὼμ, ἀλλʼ ἢ Πολυάνδριον τῆς σφαγῆς.
\vs{7}Καὶ σφάξω τὴν βουλὴν Ἰούδα, καὶ τὴν βουλὴν Ἱερουσαλὴμ ἐν τῷ τόπῳ τούτῳ, καὶ καταβαλῶ αὐτοὺς ἐν μαχαίρᾳ ἐναντίον τῶν ἐχθρῶν αὐτῶν, καὶ ἐν χερσὶ τῶν ζητούντων τὰς ψυχὰς αὐτῶν· καὶ δώσω τοὺς νεκροὺς αὐτῶν εἰς βρῶσιν τοῖς πετεινοῖς τοῦ οὐρανοῦ καὶ τοῖς θηρίοις τῆς γῆς·
\vs{8}Καὶ κατάξω τὴν πόλιν ταύτην εἰς ἀφανισμὸν, καὶ εἰς συρισμόν· πᾶς ὁ παραπορευόμενος ἐπʼ αὐτῆς σκυθρωπάσει, καὶ συριεῖ ὑπὲρ πάσης τῆς πληγῆς αὐτῆς.
\vs{9}Καὶ ἔδονται τὰς σάρκας τῶν υἱῶν αὐτῶν, καὶ τὰς σάρκας τῶν θυγατέρων αὐτῶν· καὶ ἕκαστος τὰς σάρκας τοῦ πλησίον αὐτοῦ ἔδονται ἐν τῇ περιοχῇ, καὶ ἐν τῇ πολιορκίᾳ ᾗ πολιορκήσουσιν αὐτοὺς οἱ ἐχθροὶ αὐτῶν.

\vs{10}Καὶ συντρίψεις τὸν βικὸν κατʼ ὀφθαλμοὺς τῶν ἀνδρῶν τῶν ἐκπορευομένων μετὰ σοῦ,
\vs{11}καὶ ἐρεῖς, τάδε λέγει Κύριος, οὕτως συντρίψω τὸν λαὸν τοῦτον, καὶ τὴν πόλιν ταύτην, καθὼς συντρίβεται ἄγγος ὀστράκινον, ὃ οὐ δυνήσεται ἰαθῆναι ἔτι.
\vs{12}Οὕτως ποιήσω, λέγει Κύριος, τῷ τόπῳ τούτῳ, καὶ τοῖς κατοικοῦσιν ἐν αὐτῷ, τοῦ δοθῆναι τὴν πόλιν ταύτην, ὡς τὴν διαπίπτουσαν.
\vs{13}Καὶ οἶκοι Ἱερουσαλὴμ, καὶ οἶκοι βασιλέων Ἰούδα ἔσονται καθὼς ὁ τόπος ὁ διαπίπτων, ἀπὸ τῶν ἀκαθαρσιῶν αὐτῶν ἐν πάσαις ταῖς οἰκίαις, ἐν αἷς ἐθυμίασαν ἐπὶ τῶν δωμάτων αὐτῶν πάσῃ τῇ στρατιᾷ τοῦ οὐρανοῦ, καὶ ἔσπεισαν σπονδὰς θεοῖς ἀλλοτρίοις.

\vs{14}Καὶ ἦλθεν Ἱερεμίας ἀπὸ τῆς διαπτώσεως, οὗ ἀπέστειλεν αὐτὸν Κύριος ἐκεῖ, τοῦ προφητεῦσαι· καὶ ἔστη ἐν τῇ αὐλῇ οἴκου Κυρίου, καὶ εἶπε πρὸς πάντα τὸν λαὸν,
\vs{15}τάδε λέγει Κύριος, ἰδοὺ ἐγὼ ἐπάγω ἐπὶ τὴν πόλιν ταύτην, καὶ ἐπὶ πάσας τὰς πόλεις αὐτῆς, καὶ ἐπὶ τὰς κώμας αὐτῆς, ἅπαντα τὰ κακὰ ἃ ἐλάλησα ἐπʼ αὐτὴν, ὅτι ἐσκλήρυναν τὸν τράχηλον αὐτῶν, τοῦ μὴ εἰσακούειν τῶν ἐντολῶν μου.

\ch{20}
Καὶ ἤκουσε Πασχὼρ ὁ υἱὸς Ἑμμὴρ ὁ ἱερεὺς, καὶ οὗτος ἦν καθεσταμένος ἡγούμενος οἴκου Κυρίου, τοῦ Ἱερεμίου προφητεύοντος τοὺς λόγους τούτους.
\vs{2}Καὶ ἐπάταξεν αὐτὸν, καὶ ἐνέβαλεν αὐτὸν εἰς τὸν καταράκτην, ὃς ἦν ἐν πύλῃ οἴκου ἀποτεταγμένου τοῦ ὑπερῴου, ὃς ἦν ἐν οἴκῳ Κυρίου.

\vs{3}Καὶ ἐξήγαγε Πασχὼρ τὸν Ἱερεμίαν ἐκ τοῦ καταράκτου· καὶ εἶπεν αὐτῷ Ἱερεμίας, οὐχὶ Πασχὼρ ἐκάλεσε τὸ ὄνομά σου, ἀλλʼ ἢ Μέτοικον.
\vs{4}Διότι τάδε λέγει Κύριος, ἰδοὺ ἐγὼ δίδωμί σε εἰς μετοικίαν σὺν πᾶσι τοῖς φίλοις σου· καὶ πεσοῦνται ἐν μαχαίρᾳ ἐχθρῶν αὐτῶν, καὶ οἱ ὀφθαλμοί σου ὄψονται· καὶ σὲ καὶ πάντα Ἰούδα δώσω εἰς χεῖρας βασιλέως Βαβυλῶνος, καὶ μετοικιοῦσιν αὐτοὺς, καὶ κατακόψουσιν ἐν μαχαίραις.
\vs{5}Καὶ δώσω τὴν πᾶσαν ἰσχὺν τῆς πόλεως ταύτης, καὶ πάντας τοὺς πόνους αὐτῆς, καὶ πάντας τοὺς θησαυροὺς τοῦ βασιλέως Ἰούδα εἰς χεῖρας ἐχθρῶν αὐτοῦ, καὶ ἄξουσιν αὐτοὺς εἰς Βαβυλῶνα.
\vs{6}Καὶ σὺ καὶ πάντες οἱ κατοικοῦντες ἐν τῷ οἴκῳ σου, πορεύσεσθε ἐν αἰχμαλωσίᾳ, καὶ ἐν Βαβυλῶνι ἀποθανῇ, καὶ ἐκεῖ ταφήσῃ σὺ καὶ πάντες οἱ φίλοι σου, οἷς ἐπροφήτευσας αὐτοῖς ψευδῆ.

\vs{7}Ἠπάτησάς με Κύριε, καὶ ἠπατήθην, ἐκράτησας, καὶ ἠδυνάσθης· ἐγενόμην εἰς γέλωτα, πᾶσαν ἡμέραν διετέλεσα μυκτηριζόμενος·
\vs{8}Ὅτι πικρῷ λόγῳ μου γελάσομαι, ἀθεσίαν καὶ ταλαιπωρίαν ἐπικαλέσομαι, ὅτι ἐγενήθη λόγος Κυρίου εἰς ὀνειδισμὸν ἐμοὶ καὶ εἰς χλευασμὸν πᾶσαν ἡμέραν μου.
\vs{9}Καὶ εἶπα, οὐ μὴ ὀνομάσω τὸ ὄνομα Κυρίου, καὶ οὐ μὴ λαλήσω ἔτι ἐπὶ τῷ ὀνόματι αὐτοῦ· καὶ ἐγένετο ὡς πῦρ καιόμενον φλέγον ἐν τοῖς ὀστοῖς μου, καὶ παρεῖμαι πάντοθεν, καὶ οὐ δύναμαι φέρειν,
\vs{10}ὅτι ἤκουσα ψόγον πολλῶν συναθροιζομένων κυκλόθεν, ἐπισύστητε, καὶ ἐπισυστῶμεν ἐπʼ αὐτῷ πάντες ἄνδρες φιλοι αὐτοῦ· τηρήσατε τὴν ἐπίνοιαν αὐτοῦ, εἰ ἀπατηθήσεται, καὶ δυνησόμεθα αὐτῷ, καὶ ληψόμεθα τὴν ἐκδίκησιν ἡμῶν ἐξ αὐτοῦ.

\vs{11}Ὁ δὲ κύριος μετʼ ἐμοῦ καθὼς μαχητὴς ἰσχύων· διατοῦτο ἐδίωξαν, καὶ νοῆσαι οὐκ ἠδύναντο· ᾐσχύνθησαν σφόδρα, ὅτι οὐκ ἐνόησαν ἀτιμίας αὐτῶν, αἳ διʼ αἰῶνος οὐκ ἐπιλησθήσονται.

\vs{12}Κύριε δοκιμάζων δίκαια, συνιῶν νεφροὺς καὶ καρδίας, ἴδοιμι τὴν παρὰ σοῦ ἐκδίκησιν ἐν αὐτοῖς, ὅτι πρὸς σὲ ἀπεκάλυψα τὰ ἀπολογήματά μου.
\vs{13}Ἄσατε τῷ Κυρίῳ, αἰνέσατε αὐτῷ, ὅτι ἐξείλατο τὴν ψυχὴν πένητος ἐκ χειρὸς πονηρευομένων.

\vs{14}Ἐπικατάρατος ἡ ἡμέρα ἐν ᾗ ἐτέχθην ἐν αὐτῇ· ἡ ἡμέρα ἐν ᾗ ἔτεκέν με ἡ μήτηρ μου, μὴ ἔστω ἐπευκτή.
\vs{15}Ἐπικατάρατος ὁ ἄνθρωπος ὁ εὐαγγελισάμενος τῷ πατρί μου λέγων, ἐτέχθη σοι ἄρσην· εὐφραινόμενος
\vs{16}ἔστω ὁ ἄνθρωπος ἐκεῖνος, ὡς αἱ πόλεις ἃς κατέστρεψε Κύριος ἐν θυμῷ καὶ οὐ μετεμελήθη· ἀκουσάτω κραυγῆς τῷ πρωὶ, καὶ ἀλαλαγμοῦ μεσημβρίας,
\vs{17}ὅτι οὐκ ἀπέκτεινέ με ἐν μήτρᾳ, καὶ ἐγένετό μοι ἡ μήτηρ μου τάφος μου, καὶ ἡ μήτρα συλλήψεως αἰωνίας.
\vs{18}Ἱνατί τοῦτο ἐξῆλθον ἐκ μήτρας, τοῦ βλέπειν κόπους καὶ πόνους, καὶ διετέλεσαν ἐν αἰσχύνῃ αἱ ἡμέραι μου;

\ch{21}
Ὁ ΛΟΓΟΣ Ὁ ΓΕΝΟΜΕΝΟΣ ΠΑΡΑ ΚΥΡΙΟΥ ΠΡΟΣ ἹΕΡΕΜΙΑΝ, ὍΤΕ ἈΠΕΣΤΕΙΛΕ ΠΡΟΣ ΑΥΤΟΝ Ὁ ΒΑΣΙΛΕΥΣ ΣΕΔΕΚΙΑΣ ΤΟΝ ΠΑΣΧΩΡ ΥΙΟΝ ΜΕΓΧΙΟΥ, ΚΑΙ ΣΟΦΟΝΙΑΝ ΥΙΟΝ ΒΑΣΑΙΟΥ ΤΟΝ ἹΕΡΕΑ, ΛΕΓΩΝ,
\vs{2}ἐπερώτησον περὶ ἡμῶν τὸν Κύριον· ὅτι βασιλεὺς Βαβυλῶνος ἐφέστηκεν ἐφʼ ἡμᾶς· εἰ ποιήσει Κύριος κατὰ πάντα τὰ θαυμάσια αὐτοῦ, καὶ ἀπελεύσεται ἀφʼ ἡμῶν.

\vs{3}Καὶ εἶπε πρὸς αὐτοὺς Ἱερεμίας, οὕτως ἐρεῖτε πρὸς Σεδεκίαν βασιλέα Ἰούδα,
\vs{4}τὰδε λέγει Κύριος, ἰδοὺ ἐγὼ μεταστρέφω τὰ ὅπλα τὰ πολεμικὰ, ἐν οἷς ὑμεῖς πολεμεῖτε ἐν αὐτοῖς πρὸς τοὺς Χαλδαίους τοὺς συγκεκλεικότας ὑμᾶς ἔξωθεν τοῦ τείχους· καὶ συνάξω αὐτοὺς εἰς τὸ μέσον τῆς πόλεως ταύτης,
\vs{5}καὶ πολεμήσω ἐγὼ ὑμᾶς ἐν χειρὶ ἐκτεταμένῃ καὶ ἐν βραχίονι κραταιῷ μετὰ θυμοῦ καὶ ὀργῆς μεγάλης.
\vs{6}Καὶ πατάξω πάντας τοὺς κατοικοῦντας ἐν τῇ πόλει ταύτῃ, τοὺς ἀνθρώπους καὶ τὰ κτήνη ἐν θανάτῳ μεγάλῳ, καὶ ἀποθανοῦνται.
\vs{7}Καὶ μετὰ ταῦτα οὕτως λέγει Κύριος, δώσω τὸν Σεδεκίαν βασιλέα Ἰούδα, καὶ τοὺς παῖδας αὐτοῦ, καὶ τὸν λαὸν καταλειφθέντα ἐν τῇ πόλει ταύτῃ ἀπὸ τοῦ θανάτου καὶ ἀπὸ τοῦ λιμοῦ καὶ ἀπὸ τῆς μαχαίρας, εἰς χεῖρας ἐχθρῶν αὐτῶν, τῶν ζητούντων τὰς ψυχὰς αὐτῶν, καὶ κατακόψουσιν αὐτοὺς ἐν στόματι μαχαίρας· οὐ φείσομαι ἐπʼ αὐτοῖς, καὶ οὐ μὴ οἰκτειρήσω αὐτούς.

\vs{8}Καὶ πρὸς τὸν λαὸν τοῦτον ἐρεῖς, τάδε λέγει Κύριος, ἰδοὺ ἐγὼ δέδωκα πρὸ προσώπου ὑμῶν τὴν ὁδὸν τῆς ζωῆς, καὶ τὴν ὁδὸν τοῦ θανάτου.
\vs{9}Ὁ καθήμενος ἐν τῇ πόλει ταύτῃ, ἀποθανεῖται ἐν μαχαίρᾳ καὶ ἐν λιμῷ· καὶ ὁ ἐκπορευόμενος προσχωρῆσαι πρὸς τοὺς Χαλδαίους τοὺς συγκεκλεικότας ὑμᾶς, ζήσεται, καὶ ἔσται ἡ ψυχὴ αὐτοῦ εἰς σκῦλα, καὶ ζήσεται.
\vs{10}Διότι ἐστήρικα τὸ πρόσωπόν μου ἐπὶ τὴν πόλιν ταύτην εἰς κακὰ, καὶ οὐκ εἰς ἀγαθά· εἰς χεῖρας βασιλέως Βαβυλῶνος παραδοθήσεται, καὶ κατακαύσει αὐτὴν ἐν πυρί.

\vs{11}Ὁ οἶκος βασιλέως Ἰούδα, ἀκούσατε λόγον Κυρίου.
\vs{12}Οἶκος Δαυὶδ, τάδε λέγει Κύριος, κρίνατε πρωῒ κρίμα, καὶ κατευθύνατε, καὶ ἐξέλεσθε διηρπασμένον ἐκ χειρὸς ἀδικοῦντος αὐτὸν, ὅπως μὴ ἀναφθῇ ὡς πῦρ ἡ ὀργή μου, καὶ καυθήσεται, καὶ οὐκ ἔσται ὁ σβέσων.
\vs{13}Ἰδοὺ ἐγὼ πρὸς σὲ τὸν κατοικοῦντα τὴν κοιλάδα σὸρ, τὴν πεδεινὴν, τοὺς λέγοντας, τίς πτοήσει ἡμᾶς; ἢ τίς εἰσελεύσεται πρὸς τὸ κατοικητήριον ἡμῶν;
\vs{14}Καὶ ἀνάψω πῦρ ἐν τῷ δρυμῷ αὐτῆς, καὶ ἔδεται πάντα τὰ κύκλῳ αὐτῆς.

\ch{22}
Τάδε λέγει Κύριος, πορεύου καὶ κατάβηθι εἰς τὸν οἶκον τοῦ βασιλέως Ἰούδα, καὶ λαλήσεις ἐκεῖ τὸν λόγον τοῦτον,
\vs{2}καὶ ἐρεῖς,

Ἄκουε λόγον Κυρίου βασιλεῦ Ἰούδα, ὁ καθήμενος ἐπὶ θρόνου Δαυὶδ, σὺ καὶ ὁ οἶκός σου καὶ ὁ λαός σου, καὶ οἱ εἰσπορευόμενοι ταῖς πύλαις ταύταις·
\vs{3}Τάδε λέγει Κύριος, ποιεῖτε κρίσιν καὶ δικαιοσύνην, καὶ ἐξαιρεῖσθε διηρπασμένον ἐκ χειρὸς ἀδικοῦντος αὐτὸν, καὶ προσήλυτον καὶ ὀρφανὸν καὶ χήραν μὴ καταδυναστεύετε, καὶ μὴ ἀσεβεῖτε, καὶ αἷμα ἀθῶον μὴ ἐκχέητε ἐν τῷ τόπῳ τούτῳ.
\vs{4}Διότι ἐὰν ποιοῦντες ποιήσητε τὸν λόγον τοῦτον, καὶ εἰσελεύσονται ἐν ταῖς πύλαις τοῦ οἴκου τούτου βασιλεῖς καθήμενοι ἐπὶ θρόνου Δαυὶδ, καὶ ἐπιβεβηκότες ἐφʼ ἁρμάτων καὶ ἵππων, αὐτοὶ καὶ οἱ παῖδες αὐτῶν, καὶ ὁ λαὸς αὐτῶν.
\vs{5}Ἐὰν δὲ μὴ ποιήσητε τοὺς λόγους τούτους, κατʼ ἐμαυτοῦ ὤμοσα, λέγει Κύριος, ὅτι εἰς ἐρήμωσιν ἔσται ὁ οἶκος οὗτος.

\vs{6}Ὅτι τάδε λέγει Κύριος κατὰ τοῦ οἴκου βασιλέως Ἰούδα, Γαλαὰδ σύ μοι ἀρχὴ τοῦ Λιβάνου, ἐὰν μὴ θῶ σε εἰς ἔρημον, πόλεις μὴ κατοικηθησομένας,
\vs{7}καὶ ἐπάξω ἐπὶ σὲ ὀλοθρεύοντα ἄνδρα, καὶ τὸν πέλεκυν αὐτοῦ, καὶ ἐκκόψουσι τὰς ἐκλεκτὰς κέδρους σου, καὶ ἐμβαλοῦσιν εἰς τὸ πῦρ.
\vs{8}Καὶ διελεύσονται ἔθνη διὰ τῆς πόλεως ταύτης, καὶ ἐρεῖ ἕκαστος πρὸς τὸν πλησίον αὐτοῦ, διατί ἐποίησε Κύριος οὕτως τῇ πόλει ταύτῃ τῇ μεγάλῃ;
\vs{9}Καὶ ἐροῦσιν, ἀνθʼ ὧν ἐγκατέλιπον τὴν διαθήκην Κυρίου Θεοῦ αὐτῶν, καὶ προσεκύνησαν θεοῖς ἀλλοτρίοις, καὶ ἐδούλευσαν αὐτοῖς.

\vs{10}Μὴ κλαίετε τὸν τεθνηκότα, μηδὲ θρηνεῖτε αὐτόν· κλαύσατε κλαυθμῷ τὸν ἐκπορευόμενον, ὅτι οὐκ ἐπιστρέψει ἔτι, οὐδὲ ὄψεται τὴν γῆν πατρίδος αὐτοῦ.
\vs{11}Διότι τάδε λέγει Κύριος ἐπὶ Σελλὴμ υἱὸν Ἰωσία τὸν βασιλεύοντα ἀντὶ Ἰωσίου τοῦ πατρὸς αὐτοῦ, ὃς ἐξῆλθεν ἐκ τοῦ τόπου τούτου· οὐκ ἀναστρέψει ἐκεῖ ἔτι,
\vs{12}ἀλλʼ ἢ ἐν τῷ τόπῳ τούτῳ οὗ μετῴκισα αὐτὸν, ἐκεῖ ἀποθανεῖται, καὶ τὴν γῆν ταύτην οὐκ ὄψεται ἔτι.

\vs{13}Ὁ οἰκοδομῶν οἰκίαν αὐτοῦ οὐ μετὰ δικαιοσύνης, καὶ τὰ ὑπερῷα αὐτοῦ οὐκ ἐν κρίματι, παρὰ τῷ πλησίον αὐτοῦ ἐργᾶται δωρεὰν, καὶ τὸν μισθὸν αὐτοῦ οὐ μὴ ἀποδώσει αὐτῷ.
\vs{14}Ὠκοδόμησας σεαυτῷ οἶκον σύμμετρον, ὑπερῷα ῥιπιστὰ διεσταλμένα θυρίσι, καὶ ἐξυλωμένα ἐν κέδρῳ, καὶ κεχρισμένα ἐν μίλτῳ.
\vs{15}Μὴ βασιλεύσῃς, ὅτι σὺ παροξυνῇ ἐν Ἄχαζ τῷ πατρί σου; οὐ φάγονται, καὶ οὐ πίονται· βέλτιόν σε ποιεῖν κρίμα καὶ δικαιοσύνην.
\vs{16}Οὐκ ἔγνωσαν, οὐκ ἔκριναν κρίσιν ταπεινῷ, οὐδὲ κρίσιν πένητος· οὐ τοῦτό ἐστι τὸ μὴ γηῶναί σε ἐμὲ; λέγει Κύριος.
\vs{17}Ἰδοὺ οὐκ εἰσὶν οἱ ὀφθαλμοί σου, οὐδὲ ἡ καρδία σου καλὴ, ἀλλὰ εἰς τὴν πλεονεξίαν σου, καὶ εἰς τὸ αἷμα τὸ ἀθῶον τοῦ ἐκχέαι αὐτὸ, καὶ εἰς ἀδικήματα καὶ εἰς φόνον, τοῦ ποιεῖν αὐτά.

\vs{18}Διατοῦτο τάδε λέγει Κύριος ἐπὶ Ἰωακεὶμ υἱὸν Ἰωσία βασιλέα Ἰούδα, καὶ ἐπὶ τὸν ἄνδρα τοῦτον, οὐ κόψονται αὐτὸν, ὦ ἀδελφὲ, οὐδὲ μὴ κλαύσονται αὐτὸν, οἴμοι Κύριε.
\vs{19}Ταφὴν ὄνου ταφήσεται, συμψησθεὶς ῥιφήσεται ἐπέκεινα τῆς πύλης Ἱερουσαλήμ.

\vs{20}Ἀνάβηθι εἰς τὸν Λίβανον, καὶ κράξον, καὶ εἰς τὴν Βασὰν δὸς τὴν φωνήν σου, καὶ βόησον εἰς τὸ πέρας τῆς θαλάσσης, ὅτι συνετρίβησαν πάντες οἱ ἐρασταί σου.
\vs{21}Ἐλάλησα πρὸς σὲ ἐν τῇ παραπτώσει σου, καὶ εἶπας, οὐκ ἀκούσομαι· αὕτη ἡ ὁδός σου ἐκ νεότητός σου, οὐκ ἤκουσας τῆς φωνῆς μου.
\vs{22}Πάντας τοὺς ποιμένας σου ποιμανεῖ ἄνεμος, καὶ οἱ ἐρασταί σου ἐν αἰχμαλωσίᾳ ἐξελεύσονται, ὅτι τότε αἰσχυνθήσῃ καὶ ἀτιμωθήσῃ ἀπὸ πάντων τῶν φιλούντων σε.
\vs{23}Κατοικοῦσα ἐν τῷ Λιβάνῳ, ἐννοσσεύουσα ἐν ταῖς κέδροις, καταστενάξεις ἐν τῷ ἐλθεῖν σοι ὀδύνας ὡς τικτούσης.

\vs{24}Ζῶ ἐγὼ, λέγει Κύριος, ἐὰν γενόμενος γένηται Ἰεχονίας υἱὸς Ἰωακεὶμ βασιλεὺς Ἰούδα ἀποσφράγισμα ἐπὶ τῆς χειρὸς τῆς δεξιᾶς μου, ἐκεῖθεν ἐκσπάσω σε,
\vs{25}καὶ παραδώσω σε εἰς χεῖρας τῶν ζητούντων τὴν ψυχήν σου, ὧν σὺ εὐλαβῇ ἀπὸ προσώπου αὐτῶν, εἰς χεῖρας τῶν Χαλδαίων,
\vs{26}καὶ ἀποῤῥίψω σε καὶ τὴν μητέρα σου τὴν τεκοῦσάν σε εἰς γῆν, οὗ οὐκ ἐτέχθης ἐκεῖ, καὶ ἐκεῖ ἀποθανεῖσθε·
\vs{27}Εἰς δὲ τὴν γῆν ἣν αὐτοὶ εὔχονται ταῖς ψυχαῖς αὐτῶν, οὐ μὴ ἀποστρέψωσιν.
\vs{28}Ἠτιμώθη Ἰεχονίας ὡς σκεῦος οὗ οὐκ ἔστι χρεία αὐτοῦ, ὅτι ἐξεῤῥίφη, καὶ ἐξεβλήθη εἰς γῆν ἣν οὐκ ᾔδει.

\vs{29}Γῆ, γῆ ἄκουε λόγον Κυρίου,
\vs{30}γράψον τὸν ἄνδρα τοῦτον ἐκκήρυκτον ἄνθρωπον, ὅτι οὐ μὴ αὐξηθῇ ἐκ τοῦ σπέρματος αὐτοῦ καθήμενος ἐπὶ θρόνου Δαυὶδ, ἄρχων ἔτι ἐν τῷ Ἰούδα.

\ch{23}
Ὢ ποιμένες οἱ ἀπολλύοντες καὶ διασκορπίζοντες τὰ πρόβατα τῆς νομῆς αὐτῶν.
\vs{2}Διατοῦτο τάδε λέγει Κύριος ἐπὶ τοὺς ποιμαίνοντας τὸν λαόν μου, ὑμεῖς διεσκορπίσατε τὰ πρόβατά μου, καὶ ἐξώσατε αὐτὰ, καὶ οὐκ ἐπεσκέψασθε αὐτὰ, ἰδοὺ ἐγὼ ἐκδικῶ ἐφʼ ὑμᾶς κατὰ τὰ πονηρὰ ἐπιτηδεύματα ὑμῶν.
\vs{3}Καὶ ἐγὼ εἰσδέξομαι τοὺς καταλοίπους τοῦ λαοῦ μου ἐπὶ πάσης τῆς γῆς, οὗ ἔξωσα αὐτοὺς ἐκεῖ, καὶ καταστήσω αὐτοὺς εἰς τὴν νομὴν αὐτῶν, καὶ αὐξηθήσονται, καὶ πληθυνθήσονται.
\vs{4}Καὶ ἀναστήσω αὐτοῖς ποιμένας, οἳ ποιμανοῦσιν αὐτοὺς, καὶ οὐ φοβηθήσονται ἔτι, οὐδὲ πτοηθήσονται, λέγει Κύριος.

\vs{5}Ἰδοὺ ἡμέραι ἔρχονται, λέγει Κύριος, καὶ ἀναστήσω τῷ Δαυὶδ ἀνατολὴν δικαίαν, καὶ βασιλεύσει βασιλεὺς, καὶ συνήσει, καὶ ποιήσει κρίμα καὶ δικαιοσύνην ἐπὶ τῆς γῆς.
\vs{6}Ἐν ταῖς ἡμέραις αὐτοῦ καὶ σωθήσεται Ἰούδας, καὶ Ἰσραὴλ κατασκηνώσει πεποιθὼς, καὶ τοῦτο τὸ ὄνομα αὐτοῦ, ὃ καλέσει αὐτὸν Κύριος, Ἰωσεδὲκ ἐν τοῖς προφήταις.

\vs{9}Συνετρίβη ἡ καρδία μου ἐν ἐμοὶ, ἐσαλεύθη πάντα τὰ ὀστᾶ μου, ἐγενήθην ὡς ἀνὴρ συντετριμμένος, καὶ ὡς ἄνθρωπος συνεχόμενος ἀπὸ οἶνου ἀπὸ προσώπου Κυρίου καὶ ἀπὸ προσώπου εὐπρεπείας δόξης αὐτοῦ.
\vs{10}Ὅτι ἀπὸ προσώπου τούτων ἐπένθησεν ἡ γῆ, ἐξηράνθησαν αἱ νομαὶ τῆς ἐρήμου· καὶ ἐγένετο ὁ δρόμος αὐτῶν πονηρὸς, καὶ ἡ ἰσχὺς αὐτῶν οὕτως.
\vs{11}Ὅτι ἱερεὺς καὶ προφήτης ἐμολύνθησαν, καὶ ἐν τῷ οἴκῳ μου εἶδον πονηρίας αὐτῶν.
\vs{12}Διατοῦτο γενέσθω ἡ ὁδὸς αὐτῶν αὐτοῖς εἰς ὀλίσθημα ἐν γνόφῳ, καὶ ὑποσκελισθήσονται, καὶ πεσοῦνται ἐν αὐτῇ· διότι ἐπάξω ἐπʼ αὐτοὺς κακὰ, ἐν ἐνιαυτῷ ἐπισκέψεως αὐτῶν.

\vs{13}Καὶ ἐν τοῖς προφήταις Σαμαρείας εἶδον ἀνομήματα· ἐπροφήτευσαν διὰ τῆς Βάαλ, καὶ ἐπλάνησαν τὸν λαόν μου Ἰσραήλ.
\vs{14}Καὶ ἐν τοῖς προφήταις Ἱερουσαλὴμ ἑώρακα φρικτὰ, μοιχωμένους, καὶ πορευομένους ἐν ψεύδεσι, καὶ ἀντιλαμβανομένους χειρῶν πολλῶν, τοῦ μὴ ἀποστραφῆναι ἕκαστον ἀπὸ τῆς ὁδοῦ αὐτοῦ τῆς πονηρᾶς· ἐγενήθησάν μοι πάντες ὡς Σόδομα, καὶ οἱ κατοικοῦντες αὐτὴν ὥσπερ Γόμοῤῥα.

\vs{15}Διατοῦτο τάδε λέγει Κύριος, ἰδοὺ ἐγὼ ψωμιῶ αὐτοὺς ὀδύνην, καὶ ποτιῶ αὐτοὺς ὕδωρ πικρὸν, ὅτι ἀπὸ τῶν προφητῶν Ἱερουσαλὴμ ἐξῆλθε μολυσμὸς πάσῃ τῇ γῇ.

\vs{16}Οὕτως λέγει Κύριος παντοκράτωρ, μὴ ἀκούετε τοὺς λόγους τῶν προφητῶν, ὅτι ματαιοῦσιν ἑαυτοῖς ὅρασιν, ἀπὸ καρδίας αὐτῶν λαλοῦσι, καὶ οὐκ ἀπὸ στόματος Κυρίου.
\vs{17}Λέγουσι τοῖς ἀπωθουμένοις τὸν λόγον Κυρίου, εἰρήνη ἔσται ὑμῖν, καὶ πᾶσι τοῖς πορευομένοις τοῖς θελήμασιν αὐτῶν, καὶ παντὶ τῷ πορευομένῳ πλάνῃ καρδίας αὐτοῦ, εἶπαν, οὐχ ἥξει ἐπὶ σὲ κακά.
\vs{18}Ὅτι τίς ἔστη ἐν ὑποστήματι Κυρίου, καὶ εἶδε τὸν λόγον αὐτοῦ; τίς ἠνωτίσατο, καὶ ἤκουσεν;
\vs{19}Ἰδοὺ σεισμὸς παρὰ Κυρίου, καὶ ὀργὴ ἐκπορεύεται εἰς συσσεισμὸν, συστρεφομένη ἐπὶ τοὺς ἀσεβεῖς ἥξει.
\vs{20}Καὶ οὐκ ἔτι ἀποστρέψει ὁ θυμὸς Κυρίου, ἕως ποιήσῃ αὐτὸ, καὶ ἕως ἂν στήσῃ αὐτὸ, ἀπὸ ἐγχειρήματος καρδίας αὐτοῦ· ἐπʼ ἐσχάτου τῶν ἡμερῶν νοήσουσιν αὐτό.

\vs{21}Οὐκ ἀπέστελλον τοὺς προφήτας, καὶ αὐτοὶ ἔτρεχον· οὐδὲ ἐλάλησα πρὸς αὐτοὺς, καὶ αὐτοὶ ἐπροφήτευον.
\vs{22}Καὶ εἰ ἔστησαν ἐν τῇ ὑποστάσει μου, καὶ εἰ ἤκουσαν τῶν λόγων μου, καὶ τὸν λαόν μου ἂν ἀπέστρεφον αὐτοὺς ἀπὸ τῶν πονηρῶν ἐπιτηδευμάτων αὐτῶν.

\vs{23}Θεὸς ἐγγίζων ἐγώ εἰμι, λέγει Κύριος, καὶ οὐχὶ Θεὸς πόῤῥωθεν.
\vs{24}Εἰ κρυβήσεταί τις ἐν κρυφαίοις, καὶ ἐγὼ οὐκ ὄψομαι αὐτόν; μὴ οὐχὶ τὸν οὐρανὸν καὶ τὴν γῆν ἐγὼ πληρῶ; λέγει Κύριος.

\vs{25}Ἤκουσα ἃ λαλοῦσιν οἱ προφῆται, ἃ προφητεύουσιν ἐπὶ τῷ ὀνόματί μου, ψευδῆ λέγοντες, ἠνυπνιασάμην ἐνύπνιον.
\vs{26}Ἕως πότε ἔσται ἐν καρδίᾳ τῶν προφητῶν τῶν προφητευόντων ψευδῆ, ἐν τῷ προφητεύειν αὐτοὺς τὰ θελήματα τῆς καρδίας αὐτῶν,
\vs{27}τῶν λογιζομένων τοῦ ἐπιλαθέσθαι τοῦ νόμου μου ἐν τοῖς ἐνυπνίοις αὐτῶν, ἃ διηγοῦντο ἕκαστος τῷ πλησίον αὐτοῦ, καθάπερ ἐπελάθοντο οἱ πατέρες αὐτῶν τοῦ ὀνόματός μου ἐν τῇ Βάαλ;
\vs{28}Ὁ προφήτης ἐν ᾧ τὸ ἐνύπνιόν ἐστι, διηγησάσθω τὸ ἐνύπνιον αὐτοῦ, καὶ ἐν ᾧ ὁ λόγος μου πρὸς αὐτὸν, διηγησάσθω τὸν λόγον μου ἐπʼ ἀληθείας· τί τὸ ἄχυρον πρὸς τὸν σῖτον; οὕτως οἱ λόγοι μου, λέγει Κύριος.
\vs{29}Οὐκ ἰδοὺ οἱ λόγοι μου, ὥσπερ πῦρ, λέγει Κύριος, καὶ ὡς πέλυξ κόπτων πέτραν;

\vs{30}Ἰδοὺ ἐγὼ διατοῦτο πρὸς τοὺς προφήτας, λέγει Κύριος ὁ Θεὸς, τοὺς κλέπτοντας τοὺς λόγους μου ἕκαστον παρὰ τοῦ πλησίον αὐτοῦ.
\vs{31}Ἰδοὺ ἐγὼ πρὸς τοὺς προφήτας τοὺς ἐκβάλλοντας προφητείας γλώσσης, καὶ νυστάζοντας νυσταγμὸν αὐτῶν.
\vs{32}Διατοῦτο ἰδοὺ ἐγὼ πρὸς τοὺς προφήτας τοὺς προφητεύοντας ἐνύπνια ψευδῆ, καὶ οὐ διηγοῦντο αὐτὰ, καὶ ἐπλάνησαν τὸν λαόν μου ἐν τοῖς ψεύδεσιν αὐτῶν, καὶ ἐν τοῖς πλάνοις αὐτῶν, καὶ ἐγὼ οὐκ ἀπέστειλα αὐτοὺς, καὶ οὐκ ἐνετειλάμην αὐτοῖς, καὶ ὠφέλειαν οὐκ ὠφελήσουσι τὸν λαὸν τοῦτον.

\vs{33}Καὶ ἐὰν ἐρωτήσωσιν ὁ λαὸς οὗτος, ἢ ἱερεὺς, ἢ προφήτης, τί τὸ λῆμμα Κυρίου; καὶ ἐρεῖς αὐτοῖς, ὑμεῖς ἐστε τὸ λῆμμα, καὶ ῥάξω ὑμᾶς, λέγει Κύριος.
\vs{34}Ὁ προφήτης, καὶ οἱ ἱερεῖς, καὶ ὁ λαὸς, οἳ ἂν εἴπωσι, λῆμμα Κυρίου, καὶ ἐκδικήσω τὸν ἄνθρωπον ἐκεῖνον, καὶ τὸν οἶκον αὐτοῦ.
\vs{35}Οὕτως ἐρεῖτε ἕκαστος πρὸς τὸν πλησίον αὐτοῦ, καὶ ἕκαστος πρὸς τὸν ἀδελφὸν αὐτοῦ, τί ἀπεκρίθη Κύριος, καὶ τί ἐλάλησε Κύριος;
\vs{36}Καὶ λῆμμα Κυρίου μὴ ὀνομάζετε ἔτι, ὅτι τὸ λῆμμα τῷ ἀνθρώπῳ ἔσται ὁ λόγος αὐτοῦ.
\vs{37}Καὶ διατί ἐλάλησε Κύριος ὁ Θεὸς ἡμῶν;
\vs{38}Διατοῦτο τάδε λέγει Κύριος ὁ Θεὸς ἡμῶν ἄνθʼ ὧν εἴπατε τὸν λόγον τοῦτον, λῆμμα Κυρίου, καὶ ἀπέστειλα πρὸς ὑμᾶς, λέγων, οὐκ ἐρεῖτε, λῆμμα Κυρίου·
\vs{39}διατοῦτο ἰδοὺ ἐγὼ λαμβάνω, καὶ ῥάσσω ὑμᾶς, καὶ τὴν πόλιν ἣν ἔδωκα ὑμῖν καὶ τοῖς πατράσιν ὑμῶν.
\vs{40}Καὶ δώσω ἐφʼ ὑμᾶς ὀνειδισμὸν αἰώνιον, καὶ ἀτιμίαν αἰώνιον, ἥτις οὐκ ἐπιλησθήσεται.

\vs{40a}Διατοῦτο ἰδοὺ ἡμέραι ἔρχονται, λέγει Κύριος, καὶ οὐκ ἐροῦσιν ἔτι, ζῇ Κύριος, ὃς ἀνήγαγε τὸν οἶκον Ἰσραὴλ ἐκ γῆς Αἰγύπτου,
\vs{40b}ἀλλὰ, ζῇ Κύριος, ὃς συνήγαγε πᾶν τὸ σπέρμα Ἰσραὴλ ἀπὸ γῆς Βοῤῥᾶ, καὶ ἀπὸ πασῶν τῶν χωρῶν, οὗ ἔξωσεν αὐτοὺς ἐκεῖ, καὶ ἀπεκατέστησεν αὐτοὺς εἰς τὴν γῆν αὐτῶν.

\ch{24}
Ἔδειξέ μοι Κύριος δύο καλάθους σύκων, κειμένους κατὰ πρόσωπον ναοῦ Κυρίου, μετὰ τὸ ἀποικίσαι Ναβουχοδονόσορ βασιλέα Βαβυλῶνος τὸν Ἰεχονίαν υἱὸν Ἰωακεὶμ βασιλέα Ἰούδα, καὶ τοὺς ἄρχοντας, καὶ τοὺς τεχνίτας, καὶ τοὺς δεσμώτας, καὶ τοὺς πλουσίους ἐξ Ἱερουσαλὴμ, καὶ ἤγαγεν αὐτοὺς εἰς Βαβυλῶνα.
\vs{2}Ὁ κάλαθος ὁ εἷς σύκων χρηστῶν σφόδρα, ὡς τὰ σύκα τὰ πρώϊμα· καὶ ὁ κάλαθος ὁ ἕτερος σύκων πονηρῶν σφόδρα, ἃ οὐ βρωθήσεται ἀπὸ πονηρίας αὐτῶν.
\vs{3}Καὶ εἶπε Κύριος πρὸς μὲ, τί σὺ ὁρᾷς Ἱερεμία; καὶ εἶπα, σύκα· σύκα τὰ χρηστὰ, χρηστὰ λίαν· καὶ τὰ πονηρὰ, πονηρὰ λίαν, ἃ οὐ βρωθήσεται ἀπὸ πονηρίας αὐτῶν.

\vs{4}Καὶ ἐγένετο λόγος Κυρίου πρὸς μὲ, λέγων,
\vs{5}Τάδε λέγει Κύριος ὁ Θεὸς Ἰσραὴλ, ὡς τὰ σύκα τὰ χρηστὰ ταῦτα, οὕτως ἐπιγνώσομαι τοὺς ἀποικισθέντας Ἰουδαίους, οὓς ἐξαπέσταλκα ἐκ τοῦ τόπου τούτου εἰς γῆν Χαλδαίων εἰς ἀγαθά.
\vs{6}Καὶ στηριῶ τοὺς ὀφθαλμούς μου ἐπʼ αὐτοὺς εἰς ἀγαθὰ, καὶ ἀποκαταστήσω αὐτοὺς εἰς τὴν γῆν ταύτην εἰς ἀγαθά· καὶ ἀνοικοδομήσω αὐτοὺς, καὶ οὐ μὴ καθελῶ αὐτοὺς, καὶ καταφυτεύσω αὐτοὺς, καὶ οὐ μὴ ἐκτίλω.

\vs{7}Καὶ δώσω αὐτοῖς καρδίαν τοῦ εἰδέναι αὐτοὺς ἐμὲ, ὅτι ἐγώ εἰμι Κύριος, καὶ ἔσονταί μου εἰς λαὸν, καὶ ἐγὼ ἔσομαι αὐτοῖς εἰς Θεὸν, ὅτι ἐπιστραφήσονται ἐπʼ ἐμὲ ἐξ ὅλης τῆς καρδίας αὐτῶν.

\vs{8}Καὶ ὡς τὰ σύκα τὰ πονηρὰ, ἃ οὐ βρωθήσονται ἀπὸ πονηρίας αὐτῶν, τάδε λέγει Κύριος, οὕτως παραδώσω τὸν Σεδεκίαν βασιλέα Ἰούδα, καὶ τοὺς μεγιστᾶνας αὐτοῦ, καὶ τὸ κατάλοιπον Ἱερουσαλὴμ τοὺς ὑπολελειμμένους ἐν τῇ γῇ ταύτῃ, καὶ τοὺς κατοικοῦντας ἐν Αἰγύπτῳ.
\vs{9}Καὶ δώσω αὐτοὺς εἰς διασκορπισμὸν εἰς πάσας τὰς βασιλείας τῆς γῆς, καὶ ἔσονται εἰς ὀνειδισμὸν, καὶ εἰς παραβολὴν, καὶ εἰς μῖσος, καὶ εἰς κατάραν ἐν παντὶ τόπῳ οὗ ἔξωσα αὐτοὺς ἐκεῖ.
\vs{10}Καὶ ἀποστελῶ εἰς αὐτοὺς τὸν λιμὸν, καὶ τὸν θάνατον, καὶ τὴν μάχαιραν, ἕως ἂν ἐκλείπωσιν ἀπὸ τῆς γῆς ἧς ἔδωκα αὐτοῖς.

\ch{25}
Ὁ ΛΟΓΟΣ Ὁ ΓΕΝΟΜΕΝΟΣ ΠΡΟΣ ἹΕΡΕΜΙΑΝ ἐπὶ πάντα τὸν λαὸν Ἰούδα ἐν τῷ ἔτει τῷ τετάρτῳ τοῦ Ἰωακεὶμ υἱοῦ Ἰωσία βασιλέως Ἰούδα,
\vs{2}ὃν ἐλάλησε πρὸς πάντα τὸν λαὸν Ἰούδα, καὶ πρὸς τοὺς κατοικοῦντας Ἱερουσαλὴμ, λέγων,

\vs{3}Ἐν τρισκαιδεκάτῳ ἔτει Ἰωσία υἱοῦ Ἀμὼς βασιλέως Ἰούδα, καὶ ἕως τῆς ἡμέρας ταύτης εἴκοσι καὶ τρία ἔτη, καὶ ἐλάλησα πρὸς ὑμᾶς ὀρθρίζων, καὶ λέγων,
\vs{4}καὶ ἀπέστελλον πρὸς ὑμᾶς τοὺς δούλους μου τοὺς προφήτας, ὄρθρου ἀποστέλλων, καὶ οὐκ εἰσηκούσατε, καὶ οὐ προσέσχετε τοῖς ὠσὶν ὑμῶν,
\vs{5}λέγων, ἀποστράφητε ἕκαστος ἀπὸ τῆς ὁδοῦ αὐτοῦ τῆς πονηρᾶς, καὶ ἀπὸ τῶν πονηρῶν ἐπιτηδευμάτων ὑμῶν, καὶ κατοικήσετε ἐπὶ τῆς γῆς ἧς ἔδωκα ὑμῖν καὶ τοῖς πατράσιν ὑμῶν, ἀπʼ αἰῶνος καὶ ἕως αἰῶνος.
\vs{6}Μὴ πορεύεσθε ὀπίσω θεῶν ἀλλοτρίων, τοῦ δουλεύειν αὐτοῖς, καὶ τοῦ προσκυνεῖν αὐτοῖς, ὅπως μὴ παροργίζητέ με ἐν τοῖς ἔργοις τῶν χειρῶν ὑμῶν, τοῦ κακῶσαι ὑμᾶς.
\vs{7}Καὶ οὐκ ἠκούσατέ μου.

\vs{8}Διατοῦτο τάδε λέγει Κύριος, ἐπειδὴ οὐκ ἐπιστεύσατε τοῖς λόγοις μου,
\vs{9}ἰδοὺ ἐγὼ ἀποστέλλω, καὶ λήψομαι πατριὰν ἀπὸ Βοῤῥᾶ, καὶ ἄξω αὐτοὺς ἐπὶ τὴν γῆν ταύτην, καὶ ἐπὶ τοὺς κατοικοῦντας αὐτὴν, καὶ ἐπὶ πάντα τὰ ἔθνη τὰ κύκλῳ αὐτῆς, καὶ ἐξερημώσω αὐτοὺς, καὶ δώσω αὐτοὺς εἰς ἀφανισμὸν, καὶ εἰς συριγμὸν, καὶ εἰς ὀνειδισμὸν αἰώνιον.
\vs{10}Καὶ ἀπολῶ ἀπʼ αὐτῶν φωνὴν χαρᾶς, καὶ φωνὴν εὐφροσύνης, φωνὴν νυμφίου καὶ φωνὴν νύμφης, ὀσμὴν μύρου, καὶ φῶς λύχνου.
\vs{11}Καὶ ἔσται πᾶσα ἡ γῆ εἰς ἀφανισμὸν, καὶ δουλεύσουσιν ἐν τοῖς ἔθνεσιν ἑβδομήκοντα ἔτη.

\vs{12}Καὶ ἐν τῷ πληρωθῆναι τὰ ἑβδομήκοντα ἔτη, ἐκδικήσω τὸ ἔθνος ἐκεῖνο, καὶ θήσομαι αὐτοὺς εἰς ἀφανισμὸν αἰώνιον,
\vs{13}καὶ ἐπάξω ἐπὶ τὴν γῆν ἐκείνην πάντας τοὺς λόγους μου, οὓς ἐλάλησα κατʼ αὐτῆς, πάντα τὰ γεγραμμένα ἐν τῷ βιβλίῳ τούτῳ·

\vs{14}Ἃ ἘΠΡΟΦΗΤΕΥΣΕΝ ἹΕΡΕΜΙΑΣ ἘΠΙ ΤΑ ἜΘΝΗ ΤΑ ΑΙΛΑΜ.

\vs{15}Τάδε λέγει Κύριος, συνετρίβη τὸ τόξον Αἰλὰμ, ἀρχὴ δυναστείας αὐτῶν.
\vs{16}Καὶ ἐπάξω ἐπὶ Αἰλὰμ τέσσαρας ἀνέμους ἐκ τῶν τεσσάρων ἄκρων τοῦ οὐρανοῦ, καὶ διασπερῶ αὐτοὺς ἐν πᾶσι τοῖς ἀνέμοις τούτοις, καὶ οὐκ ἔσται ἔθνος ὃ οὐχ ἥξει ἐκεῖ, οἱ ἐξωσμένοι Αἰλάμ.
\vs{17}Καὶ πτοήσω αὐτοὺς ἐναντίον τῶν ἐχθρῶν αὐτῶν, τῶν ζητούντων τὴν ψυχὴν αὐτῶν, καὶ ἐπάξω ἐπʼ αὐτοὺς κατὰ τὴν ὀργὴν τοῦ θυμοῦ μου, καὶ ἐπαποστελῶ ὀπίσω αὐτῶν τὴν μάχαιράν μου, ἕως τοῦ ἐξαναλῶσαι αὐτούς.
\vs{18}Καὶ θήσω τὸν θρόνον μου ἐν Αἰλὰμ, καὶ ἐξαποστελῶ ἐκεῖθεν βασιλέα καὶ μεγιστᾶνας.
\vs{19}Καὶ ἔσται ἐπʼ ἐσχάτου τῶν ἡμερῶν, καὶ ἀποστρέψω τὴν αἰχμαλωσίαν Αἰλὰμ, λέγει Κύριος.

\vs{20}Ἐν ἀρχῇ βασιλεύοντος Σεδεκίου βασιλέως, ἐγένετο ὁ λόγος οὗτος περὶ Αἰλάμ·

\ch{26}
\vs{2}ΤΗ ΑΙΓΥΠΤΩ ἘΠΙ ΔΥΝΑΜΙΝ ΦΑΡΑΩ ΝΕΧΑΩ ΒΑΣΙΛΕΩΣ ΑΙΓΥΠΤΟΥ, ὃς ἦν ἐπὶ τῷ ποταμῷ Εὐφράτῃ ἐν Χαρμεὶς, ὃν ἐπάταξε Ναβουχοδονόσορ βασιλεὺς Βαβυλῶνος, ἐν τῷ ἔτει τῷ τετάρτῳ Ἰωακεὶμ βασιλέως Ἰούδα.

\vs{3}Ἀναλάβετε ὅπλα καὶ ἀσπίδας, καὶ προσαγάγετε εἰς πόλεμον,
\vs{4}καὶ ἐπισάξατε τοὺς ἵππους, ἐπίβητε οἱ ἱππεῖς, καὶ κατάστητε ἐν ταῖς περικεφαλαίαις ὑμῶν, προσβάλετε τὰ δόρατα, καὶ ἐνδύσασθε τοὺς θώρακας ὑμῶν.

\vs{5}Τί ὅτι αὐτοὶ πτοοῦνται, καὶ ἀποχωροῦσιν εἰς τὸ ὀπίσω; διότι οἱ ἰσχυροὶ αὐτῶν κοπήσονται, φυγῇ ἔφυγον, καὶ οὐκ ἀνέστρεψαν περιεχόμενοι κυκλόθεν, λέγει Κύριος.
\vs{6}Μὴ φευγέτω ὁ κοῦφος, καὶ μὴ ἀνασωζέσθω ὁ ἰσχυρὸς ἐπὶ Βοῤῥᾶν· τὰ παρὰ τὸν Εὐφράτην ἠσθένησε, καὶ πεπτώκασι.

\vs{7}Τίς οὗτος ὡς ποταμὸς ἀναβήσεται, καὶ ὡς ποταμοὶ κυμαίνουσιν ὕδωρ;
\vs{8}Ὕδατα Αἰγύπτου ὡς ποταμὸς ἀναβήσεται· καὶ εἶπεν, ἀναβήσομαι, καὶ κατακαλύψω τὴν γῆν, καὶ ἀπολῶ τοὺς κατοικοῦντας ἐν αὐτῇ.
\vs{9}Ἐπίβητε ἐπὶ τοὺς ἵππους, παρασκευάσατε τὰ ἅρματα, ἐξέλθατε οἱ μαχηταὶ Αἰθιόπων, καὶ Λίβυες καθωπλισμένοι ὅπλοις, καὶ Λυδοὶ ἀνάβητε, ἐντείνατε τόξον.
\vs{10}Καὶ ἡ ἡμέρα ἐκείνη Κυρίῳ τῷ Θεῷ ἡμῶν ἡμέρα ἐκδικήσεως, τοῦ ἐκδικῆσαι τοὺς ἐχθροὺς αὐτοῦ· καὶ καταφάγεται ἡ μάχαιρα Κυρίου, καὶ ἐμπλησθήσεται, καὶ μεθυσθήσεται ἀπὸ τοῦ αἵματος αὐτῶν, ὅτι θυσία τῷ Κυρίῳ ἀπὸ γῆς Βοῤῥᾶ ἐπὶ ποταμῷ Εὐφράτῃ.

\vs{11}Ἀνάβηθι Γαλαὰδ, καὶ λάβε ῥητίνην τῇ παρθένῳ θυγατρὶ Αἰγύπτου· εἰς τὸ κενὸν ἐπλήθυνας ἰάματά σου, ὠφέλεια οὐκ ἔστιν ἐν σοί.
\vs{12}Ἤκουσαν ἔθνη φωνήν σου, καὶ τῆς κραυγῆς σου ἐπλήσθη ἡ γῆ, ὅτι μαχητὴς πρὸς μαχητὴν ἠσθένησαν, ἐπιτοαυτὸ ἔπεσαν ἀμφότεροι.

\vs{13}Ἃ ἘΛΑΛΗΣΕ ΚΥΡΙΟΣ ἐν χειρὶ Ἱερεμίου, τοῦ ἐλθεῖν τὸν βασιλέα Βαβυλῶνος τοῦ κόψαι γῆν Αἰγύπτου.

\vs{14}Ἀναγγείλατε εἰς Μαγδωλὸν, καὶ παραγγείλατε εἰς Μέμφιν· εἴπατε, ἐπίστηθι, καὶ ἑτοίμασον, ὅτι κατέφαγε μάχαιρα τὴν σμίλακά σου.

\vs{15}Διατί ἔφυγεν ἀπὸ σοῦ ὁ Ἄπις; ὁ μόσχος ὁ ἐκλεκτός σου οὐκ ἔμεινεν· ὅτι Κύριος παρέλυσεν αὐτόν.
\vs{16}Καὶ τὸ πλῆθός σου ἠσθένησε, καὶ ἔπεσε· καὶ ἕκαστος πρὸς τὸν πλησίον αὐτοῦ ἐλάλει, ἀναστῶμεν, καὶ ἀναστρέψωμεν πρὸς τὸν λαὸν ἡμῶν εἰς τὴν πατρίδα ἡμῶν, ἀπὸ προσώπου μαχαίρας Ἑλληνικῆς.
\vs{17}Καλέσατε τὸ ὄνομα Φαραὼ Νεχαὼ βασιλέως Αἰγύπτου, Σαὼν ἑσβειὲ μωήδ.
\vs{18}Ζῶ ἐγὼ, λέγει Κύριος ὁ Θεὸς, ὅτι ὡς τὸ Ἰταβύριον ἐν τοῖς ὄρεσι, καὶ ὡς ὁ Κάρμηλος ὁ ἐν τῇ θαλάσσῃ, ἥξει.
\vs{19}Σκεύη ἀποικισμοῦ ποίησον σεαυτῇ κατοικοῦσα θύγατερ Αἰγύπτου, ὅτι Μέμφις εἰς ἀφανισμὸν ἔσται, καὶ κληθήσεται Οὐαὶ, διὰ τὸ μὴ ὑπάρχειν κατοικοῦντας ἐν αὐτῇ.

\vs{20}Δάμαλις κεκαλλωπισμένη Αἴγυπτος, ἀπόσπασμα ἀπὸ Βοῤῥᾶ ἦλθεν ἐπʼ αὐτὴν,
\vs{21}καὶ οἱ μισθωτοὶ αὐτῆς ἐν αὐτῇ, ὥσπερ μόσχοι σιτευτοὶ τρεφόμενοι ἐν αὐτῇ· διότι καὶ αὐτοὶ ἐπεστράφησαν, καὶ ἔφυγον ὁμοθυμαδόν· οὐκ ἔστησαν, ὅτι ἡμέρα ἀπωλείας ἦλθεν ἐπʼ αὐτοὺς, καὶ καιρὸς ἐκδικήσεως αὐτῶν.
\vs{22}Φωνὴ αὐτῶν ὡς ὄφεως συρίζοντος, ὅτι ἐν ἄμμῳ πορεύονται, ἐν ἀξίναις ἥξουσιν ἐπʼ αὐτήν·
\vs{23}ὡς κόπτοντες ξύλα ἐκκόψουσι τὸν δρυμὸν αὐτῆς, λέγει Κύριος, ὅτι οὐ μὴ εἰκασθῇ, ὅτι πληθύνει ὑπὲρ ἀκρίδα, καὶ οὐκ ἔστιν αὐτοῖς ἀριθμός.
\vs{24}Κατῃσχύνθη ἡ θυγάτηρ Αἰγύπτου, παρεδόθη εἰς χεῖρας λαοῦ ἀπὸ Βοῤῥᾶ.

\vs{25}Ἰδοὺ ἐγὼ ἐκδικῶ τὸν Ἀμμὼν τὸν υἱὸν αὐτῆς ἐπὶ Φαραὼ, καὶ ἐπὶ τοὺς πεποιθότας ἐπʼ αὐτῷ.

\vs{27}Σὺ δὲ μὴ φοβηθῇς δοῦλός μου Ἰακὼβ, μηδὲ πτοηθῇς Ἰσραήλ· διότι ἐγὼ ἰδοὺ σώζων σε μακρόθεν, καὶ τὸ σπέρμα σου ἐκ τῆς αἰχμαλωσίας αὐτῶν· καὶ ἀναστρέψει Ἰακὼβ, καὶ ἡσυχάσει καὶ ὑπνώσει, καὶ οὐκ ἔσται ὁ παρενοχλῶν αὐτόν.
\vs{28}Μὴ φοβοῦ παῖς μου Ἰακὼβ, λέγει Κύριος, ὅτι μετὰ σοῦ ἐγώ εἰμι· ἡ ἀπτόητος καὶ τρυφερὰ παρεδόθη· ὅτι ποιήσω ἔθνει συντέλειαν ἐν παντὶ ἔθνει εἰς οὓς ἔξωσά σε ἐκεῖ, σὲ δὲ οὐ μὴ ποιήσω ἐκλιπεῖν, καὶ παιδεύσω σε εἰς κρίμα, καὶ ἀθῶον οὐκ ἀθωώσω σε.

\ch{27}
ΛΟΓΟΣ ΚΥΡΙΟΥ ὋΝ ἘΛΑΛΗΣΕΝ ἘΠΙ ΒΑΒΥΛΩΝΑ.

\vs{2}Ἀναγγείλατε ἐν τοῖς ἔθνεσι, καὶ ἀκουστὰ ποιήσατε, καὶ μὴ κρύψητε· εἴπατε, ἑάλωκε Βαβυλὼν, κατῃσχύνθη Βῆλος, ἡ ἀπτόητος, ἡ τρυφερὰ παρεδόθη Μαιρωδὰχ,
\vs{3}ὅτι ἀνέβη ἐπʼ αὐτὴν ἔθνος ἀπὸ Βοῤῥᾶ, οὗτος θήσει τὴν γῆν αὐτῆς εἰς ἀφανισμὸν, καὶ οὐκ ἔσται ὁ κατοικῶν ἐν αὐτῇ ἀπὸ ἀνθρώπου καὶ ἕως κτήνους.

\vs{4}Ἐν ταῖς ἡμέραις ἐκείναις, καὶ ἐν τῷ καιρῷ ἐκείνῳ, ἥξουσιν οἱ υἱοὶ Ἰσραὴλ αὐτοὶ, καὶ οἱ υἱοὶ Ἰούδα ἐπὶ τὸ αὐτὸ, βαδίζοντες καὶ κλαίοντες πορεύσονται, τὸν Κύριον Θεὸν αὐτῶν ζητοῦντες.
\vs{5}Ἕως Σιὼν ἐρωτήσουσι τὴν ὁδὸν, ὧδε γὰρ τὸ πρόσωπον αὐτῶν δώσουσι, καὶ ἥξουσι, καὶ καταφεύξονται πρὸς Κύριον τὸν Θεόν· διαθήκη γὰρ αἰώνιος οὐκ ἐπιλησθήσεται.

\vs{6}Πρόβατα ἀπολωλότα ἐγενήθη ὁ λαός μου, οἱ ποιμένες αὐτῶν ἔξωσαν αὐτοὺς, ἐπὶ τὰ ὄρη ἀπεπλάνησαν αὐτοὺς, ἐξ ὄρους ἐπὶ βουνὸν ᾤχοντο, ἐπελάθοντο κοίτης αὐτῶν.
\vs{7}Πάντες οἱ εὑρίσκοντες αὐτοὺς, ἀνήλισκον αὐτούς· οἱ ἐχθροὶ αὐτῶν εἶπαν, Μὴ ἀνῶμεν αὐτοὺς, ἀνθʼ ὧν ἥμαρτον τῷ Κυρίῳ· νομὴ δικαιοσύνης τῷ συναγαγόντι τοὺς πατέρας αὐτῶν.

\vs{8}Ἀπαλλοτριώθητε ἐκ μέσου Βαβυλῶνος καὶ ἀπὸ γῆς Χαλδαίων, καὶ ἐξέλθατε, καὶ γένεσθε ὥσπερ δράκοντες κατὰ πρόσωπον προβάτων.
\vs{9}Ὅτι ἰδοὺ ἐγὼ ἐγείρω ἐπὶ Βαβυλῶνα συναγωγὰς ἐθνῶν ἐκ γῆς Βοῤῥᾶ, καὶ παρατάξονται αὐτῇ· ἐκεῖθεν ἁλώσεται, ὡς βολὶς μαχητοῦ συνετοῦ οὐκ ἐπιστρέψει κενή.
\vs{10}Καὶ ἔσται ἡ Χαλδαία εἰς προνομὴν, πάντες οἱ προνομεύοντες αὐτὴν ἐμπλησθήσονται.

\vs{11}Ὅτι ηὐφραίνεσθε, καὶ κατεκαυχᾶσθε, διαρπάζοντες τὴν κληρονομίαν μου, διότι ἐσκιρτᾶτε, ὡς βοΐδια ἐν βοτάνῃ, καὶ ἐκερατίζετε ὡς ταῦροι.
\vs{12}Ἠσχύνθη ἡ μήτηρ ὑμῶν σφόδρα, ἐνετράπη ἡ τεκοῦσα ὑμᾶς μήτηρ ἐπʼ ἀγαθὰ, ἐσχάτη ἐθνῶν, ἔρημος
\vs{13}ἀπὸ ὀργῆς Κυρίου, οὐ κατοικηθήσεται καὶ ἔσται εἰς ἀφανισμὸν πᾶσα· καὶ πᾶς ὁ διοδεύων διὰ Βαβυλῶνος σκυθρωπάσει, καὶ συριοῦσιν ἐπὶ πᾶσαν τὴν πληγὴν αὐτῆς.

\vs{14}Παρατάξασθε ἐπὶ Βαβυλῶνα κύκλῳ πάντες τείνοντες τόξον, τοξεύσατε ἐπʼ αὐτὴν, μὴ φείσησθε ἐπὶ τοῖς τοξεύμασιν ὑμῶν,
\vs{15}καὶ κατακρατήσατε αὐτήν· παρελύθησαν αἱ χεῖρες αὐτῆς, ἔπεσαν αἱ ἐπάλξεις αὐτῆς, καὶ κατεσκάφη τὸ τεῖχος αὐτῆς, ὅτι ἐκδίκησις παρὰ Θεοῦ ἐστιν· ἐκδικεῖτε ἐπʼ αὐτὴν, καθὼς ἐποίησε, ποιήσατε αὐτῇ.
\vs{16}Ἐξολοθρεύσασθε σπέρμα ἐκ Βαβυλῶνος, κατέχοντα δρέπανον ἐν καιρῷ θερισμοῦ· ἀπὸ προσώπου μαχαίρας Ἑλληνικῆς, ἕκαστος εἰς τὸν λαὸν αὐτοῦ ἀποστρέψουσι, καὶ ἕκαστος εἰς τὴν γῆν αὐτοῦ φεύξεται.

\vs{17}Πρόβατον πλανώμενον Ἰσραὴλ, λέοντες ἔξωσαν αὐτόν· ὁ πρῶτος ἔφαγεν αὐτὸν βασιλεὺς Ἀσσοὺρ, καὶ οὗτος ὕστερον τὰ ὀστᾶ αὐτοῦ βασιλεὺς Βαβυλῶνος.
\vs{18}Διατοῦτο τάδε λέγει Κύριος, ἰδοὺ ἐγὼ ἐκδικῶ ἐπὶ τὸν βασιλέα Βαβυλῶνος, καὶ ἐπὶ τὴν γῆν αὐτοῦ, καθὼς ἐξεδίκησα ἐπὶ τὸν βασιλέα Ἀσσούρ.
\vs{19}Καὶ ἀποκαταστήσω τὸν Ἰσραὴλ εἰς τὴν νομὴν αὐτοῦ, καὶ νεμήσεται ἐν τῷ Καρμήλῳ, καὶ ἐν ὄρει Ἐφραῒμ, καὶ ἐν τῷ Γαλαὰδ, καὶ πλησθήσεται ἡ ψυχὴ αὐτοῦ.
\vs{20}Ἐν ταῖς ἡμέραις ἐκείναις, καὶ ἐν τῷ καιρῷ ἐκείνῳ, ζητήσουσι τὴν ἀδικίαν Ἰσραὴλ, καὶ οὐχ ὑπάρξει, καὶ τὰς ἁμαρτίας Ἰούδα, καὶ οὐ μὴ εὑρεθῶσιν, ὅτι ἵλεως ἔσομαι τοῖς ὑπολελειμμένοις ἐπὶ τῆς γῆς,
\vs{21}λέγει Κύριος.

Πικρῶς ἐπίβηθι ἐπʼ αὐτὴν, καὶ ἐπὶ τοὺς κατοικοῦντας ἐπʼ αὐτήν· ἐκδίκησον μάχαιρα, καὶ ἀφάνισον, λέγει Κύριος, καὶ ποίει κατὰ πάντα ὅσα ἐντέλλομαί σοι.
\vs{22}Φωνὴ πολέμου καὶ συντριβὴ μεγάλη ἐν γῇ Χαλδαίων.
\vs{23}Πῶς ἐκλάσθη καὶ συνετρίβη ἡ σφυρὰ πάσης τῆς γῆς; πῶς ἐγενήθη εἰς ἀφανισμὸν Βαβυλὼν ἐν ἔθνεσιν;
\vs{24}Ἐπιβήσονταί σοι, καὶ οὐ γνώσῃ, ὡς Βαβυλὼν καὶ ἁλώσῃ· εὑρέθης, καὶ ἐλήφθης, ὅτι τῷ Κυρίῳ ἀντέστης.

\vs{25}Ἤνοιξε Κύριος τὸν θησαυρὸν αὐτοῦ, καὶ ἐξήνεγκε τὰ σκεύη ὀργῆς αὐτοῦ, ὅτι ἔργον τῷ Κυρίῳ Θεῷ ἐν γῇ Χαλδαίων,
\vs{26}ὅτι ἐληλύθασιν οἱ καιροὶ αὐτῆς· ἀνοίξατε τὰς ἀποθήκας αὐτῆς, ἐρευνήσατε αὐτὴν ὡς σπήλαιον, καὶ ἐξολοθρεύσατε αὐτήν· μὴ γενέσθω αὐτῆς κατάλειμμα·
\vs{27}Ἀναξηράνατε αὐτῆς πάντας τοὺς καρποὺς, καὶ καταβήτωσαν εἰς σφαγήν· οὐαὶ αὐτοῖς, ὅτι ἥκει ἡ ἡμέρα αὐτῶν, καὶ καιρὸς ἐκδικήσεως αὐτῶν.
\vs{28}Φωνὴ φευγόντων, καὶ ἀνασωζομένων ἐκ γῆς Βαβυλῶνος, τοῦ ἀναγγεῖλαι εἰς Σιὼν τὴν ἐκδίκησιν παρὰ Κυρίου Θεοῦ ἡμῶν.

\vs{29}Παραγγείλατε ἐπὶ Βαβυλῶνα πολλοῖς, παντὶ ἐντείνοντι τόξον, παρεμβάλλετε ἐπʼ αὐτὴν κυκλόθεν· μὴ ἔστω αὐτῆς ἀνασωζόμενος· ἀνταπόδοτε αὐτῇ κατὰ τὰ ἔργα αὐτῆς, κατὰ πάντα ὅσα ἐποίησε, ποιήσατε αὐτῇ, ὅτι πρὸς Κύριον ἀντέστη Θεὸν ἅγιον τοῦ Ἰσραήλ.
\vs{30}Διατοῦτο πεσοῦνται οἱ νεανίσκοι αὐτῆς ἐν ταῖς πλατείαις αὐτῆς, καὶ πάντες οἱ ἄνδρες οἱ πολεμισταὶ αὐτῆς ῥιφήσονται, εἶπε Κύριος.

\vs{31}Ἰδοὺ ἐγὼ ἐπὶ σὲ τὴν ὑβρίστριαν, λέγει Κύριος· ὅτι ἥκει ἡ ἡμέρα σου, καὶ ὁ καιρὸς ἐκδικήσεως σου.
\vs{32}Καὶ ἀσθενήσει ἡ ὕβρις σου, καὶ πεσεῖται· καὶ οὐδεὶς ἔσται ὁ ἀνιστῶν αὐτήν· καὶ ἀνάψω πῦρ ἐν τῷ δρυμῷ αὐτῆς, καὶ καταφάγεται πάντα τὰ κύκλῳ αὐτῆς.

\vs{33}Τάδε λέγει Κύριος, καταδεδυνάστευνται οἱ υἱοὶ Ἰσραὴλ, καὶ οἱ υἱοὶ Ἰούδα, ἅμα πάντες οἱ αἰχμαλωτεύσαντες αὐτοὺς, κατεδυνάστευσαν αὐτοὺς, ὅτι οὐκ ἠθέλησαν ἐξαποστεῖλαι αὐτούς.
\vs{34}Καὶ ὁ λυτρούμενος αὐτοὺς ἰσχυρὸς, Κύριος παντοκράτωρ ὄνομα αὐτῷ· κρίσιν κρινεῖ πρὸς τοὺς ἀντιδίκους αὐτοῦ, ὅπως ἐξάρῃ τὴν γῆν· καὶ παροξυνεῖ τοῖς κατοικοῦσι Βαβυλῶνα,
\vs{35}μάχαιραν ἐπὶ τοὺς Χαλδαίους, καὶ ἐπὶ τοὺς κατοικοῦντας Βαβυλῶνα, καὶ ἐπὶ τοὺς μεγιστᾶνας αὐτῆς, καὶ ἐπὶ τοὺς συνετοὺς αὐτῆς·
\vs{36}Μάχαιραν ἐπὶ τοὺς μαχητὰς αὐτῆς, καὶ παραλυθήσονται· μάχαιραν ἐπὶ τοὺς ἵππους αὐτῶν, καὶ ἐπὶ τὰ ἅρματα αὐτῶν·
\vs{37}Μάχαιραν ἐπὶ τοὺς μαχητὰς αὐτῶν, καὶ ἐπὶ τὸν σύμμικτον τὸν ἐν μέσῳ αὐτῆς, καὶ ἔσονται ὡσεὶ γυναῖκες· μάχαιραν ἐπὶ τοὺς θησαυροὺς, καὶ διασκορπισθήσονται
\vs{38}ἐπὶ τῷ ὕδατι αὐτῆς, καὶ καταισχυνθήσονται, ὅτι γῆ τῶν γλυπτῶν ἐστι, καὶ ἐν ταῖς νήσοις, οὗ κατεκαυχῶντο.
\vs{39}Διατοῦτο κατοικήσουσιν ἰνδάλματα ἐν ταῖς νήσοις, καὶ κατοικήσουσιν ἐν αὐτῇ θυγατέρες σειρήνων, οὐ μὴ κατοικηθῇ οὐκέτι εἰς τὸν αἰῶνα.
\vs{40}Καθὼς κατέστρεψεν ὁ Θεὸς Σόδομα καὶ Γόμοῤῥα, καὶ τὰς ὁμορούσας αὐταῖς, εἶπε Κύριος, οὐ μὴ κατοικήσει ἐκεῖ ἄνθρωπος, καὶ οὐ μὴ παροικήσει ἐκεῖ υἱὸς ἀνθρώπου.

\vs{41}Ἰδοὺ λαὸς ἔρχεται ἀπὸ Βοῤῥᾶ, καὶ ἔθνος μέγα, καὶ βασιλεῖς πολλοὶ ἐξεγερθήσονται ἀπʼ ἐσχάτου τῆς γῆς,
\vs{42}τόξον καὶ ἐγχειρίδιον ἔχοντες, ἰταμός ἐστι, καὶ οὐ μὴ ἐλεήσῃ· ἡ φωνὴ αὐτῶν ὡς θάλασσα ἠχήσει, ἐφʼ ἵπποις ἱππάσονται, παρεσκευασμένοι, ὥσπερ πῦρ, εἰς πόλεμον, πρὸς σὲ, θύγατερ Βαβυλῶνος.
\vs{43}Ἤκουσε βασιλεὺς Βαβυλῶνος τὴν ἀκοὴν αὐτῶν, καὶ παρελύθησαν αἱ χεῖρες αὐτοῦ· θλῖψις κατεκράτησεν αὐτοῦ, ὠδῖνες ὡς τικτούσης.
\vs{44}Ἰδοὺ ὥσπερ λέων ἀναβήσεται ἀπὸ τοῦ Ἰορδάνου εἰς Γαιθὰν, ὅτι ταχέως ἐκδιώξω αὐτοὺς ἀπʼ αὐτῆς, καὶ πάντα νεανίσκον ἐπʼ αὐτὴν ἐπιστήσω· ὅτι τίς ὥσπερ ἐγώ; καὶ τίς ἀντιστήσεταί μοι; καὶ τίς οὗτος ποιμὴν, ὃς στήσεται κατὰ πρόσωπόν μου;

\vs{45}Διατοῦτο ἀκούσατε τὴν βουλὴν Κυρίου ἣν βεβούλευται ἐπὶ Βαβυλῶνα, καὶ λογισμοὺς αὐτοῦ οὓς ἐλογίσατο ἐπὶ τοὺς κατοικοῦντας Χαλδαίους· ἐὰν μὴ διαφθαρῇ τὰ ἀρνία τῶν προβάτων αὐτῶν, ἐὰν μὴ ἀφανισθῇ νομὴ ἀπʼ αὐτῶν.
\vs{46}Ὅτι ἀπὸ φωνῆς ἁλώσεως Βαβυλῶνος σεισθήσεται ἡ γῆ, καὶ κραυγὴ ἐν ἔθνεσιν ἀκουσθήσεται.

\ch{28}
Τάδε λέγει Κύριος, ἰδοὺ ἐγὼ ἐξεγείρω ἐπὶ Βαβυλῶνα καὶ ἐπὶ τοὺς κατοικοῦντας Χαλδαίους ἄνεμον καύσωνα διαφθείροντα.
\vs{2}Καὶ ἐξαποστελῶ εἰς Βαβυλῶνα ὑβριστὰς, καὶ καθυβρίσουσιν αὐτὴν, καὶ λυμανοῦνται τὴν γῆν αὐτῆς· οὐαὶ ἐπὶ Βαβυλῶνα κυκλόθεν ἐν ἡμέρᾳ κακώσεως αὐτῆς.
\vs{3}Τεινέτω ὁ τείνων τὸ τόξον αὐτοῦ, καὶ περιθέσθω ᾧ ἐστιν ὅπλα αὐτοῦ, καὶ μὴ φείσησθε ἐπὶ τοὺς νεανίσκους αὐτῆς, καὶ ἀφανίσατε πᾶσαν τὴν δύναμιν αὐτῆς.
\vs{4}Καὶ πεσοῦνται τραυματίαι ἐν γῇ Χαλδαίων, καὶ κατακεκεντημένοι ἔξωθεν αὐτῆς.

\vs{5}Διότι οὐκ ἐχήρευσεν Ἰσραὴλ καὶ Ἰούδας ἀπὸ Θεοῦ αὐτῶν, ἀπὸ Κυρίου παντοκράτορος, ὅτι ἡ γῆ αὐτῶν ἐπλήσθη ἀδικίας ἀπὸ τῶν ἁγίων Ἰσραήλ.
\vs{6}Φεύγετε ἐκ μέσου Βαβυλῶνος, καὶ ἀνασώζετε ἕκαστος τὴν ψυχὴν αὐτοῦ, καὶ μὴ ἀποῤῥιφῆτε ἐν τῇ ἀδικίᾳ αὐτῆς, ὅτι καιρὸς ἐκδικήσεως αὐτῆς ἐστι παρὰ Κυρίου, ἀνταπόδομα αὐτὸς ἀνταποδίδωσιν αὐτῇ.
\vs{7}Ποτήριον χρυσοῦν Βαβυλὼν ἐν χειρὶ Κυρίου, μεθύσκον πᾶσαν τὴν γῆν, ἀπὸ τοῦ οἴνου αὐτῆς ἐπίοσαν ἔθνη, διατοῦτο ἐσαλεύθησαν.
\vs{8}Καὶ ἄφνω ἔπεσε Βαβυλὼν, καὶ συνετρίβη· θρηνεῖτε αὐτὴν, λάβετε ῥητίνην τῇ διαφθορᾷ αὐτῆς, εἴπως ἰαθήσεται.
\vs{9}Ἰατρεύσαμεν τὴν Βαβυλῶνα, καὶ οὐκ ἰάθη· ἐγκαταλίπωμεν αὐτὴν, καὶ ἀπέλθωμεν ἕκαστος εἰς τὴν γῆν αὐτοῦ, ὅτι ἤγγικεν εἰς οὐρανὸν τὸ κρίμα αὐτῆς, ἐξῇρεν ἕως τῶν ἄστρων.
\vs{10}Ἐξήνεγκε Κύριος τὸ κρίμα αὐτοῦ· δεῦτε, καὶ ἀναγγείλωμεν ἐν Σιὼν τὰ ἔργα Κυρίου τοῦ Θεοῦ ἡμῶν.

\vs{11}Παρασκευάζετε τὰ τοξεύματα, πληροῦτε τὰς φαρέτρας· ἤγειρε Κύριος τὸ πνεῦμα βασιλέως Μήδων, ὅτι εἰς Βαβυλῶνα ἡ ὀργὴ αὐτοῦ, τοῦ ἐξολοθρεῦσαι αὐτὴν, ὅτι ἐκδίκησις Κυρίου ἐστὶν, ἐκδίκησις λαοῦ αὐτοῦ ἐστιν.
\vs{12}Ἐπὶ τειχέων Βαβυλῶνος ἄρατε σημεῖον, ἐπιστήσατε φαρέτρας, ἐγείρατε φυλακὰς, ἑτοιμάσατε ὅπλα, ὅτι ἐνεχείρισε, καὶ ποιήσει Κύριος ἃ ἐλάλησεν ἐπὶ τοὺς κατοικοῦντας Βαβυλῶνα,
\vs{13}κατασκηνοῦντας ἐφʼ ὕδασι πολλοῖς, καὶ ἐπὶ πλήθει θησαυρῶν αὐτῆς· ἥκει τὸ πέρας σου ἀληθῶς εἰς τὰ σπλάγχνα σου.
\vs{14}Ὅτι ὤμοσε Κύριος κατὰ τοῦ βραχίονος αὐτοῦ, διότι πληρώσω σε ἀνθρώπων ὡσεὶ ἀκρίδων, καὶ φθέγξονται ἐπὶ σὲ οἱ καταβαίνοντες.

\vs{15}Κύριος ποιῶν γῆν ἐν τῇ ἰσχύϊ αὐτοῦ, ἑτοιμάζων οἰκουμένην ἐν τῇ σοφίᾳ αὐτοῦ, ἐν τῇ συνέσει αὐτοῦ ἐξέτεινε τὸν ουρανὸν,
\vs{16}εἰς φωνὴν ἔθετο ἦχος ὕδατος ἐν οὐρανῷ, καὶ ἀνήγαγε νεφέλας ἀπʼ ἐσχάτου τῆς γῆς· ἀστραπὰς εἰς ὑετὸν ἐποίησε, καὶ ἐξήγαγε φῶς ἐκ τῶν θησαυρῶν αὐτοῦ.
\vs{17}Ἐματαιώθη πᾶς ἄνθρωπος ἀπὸ γνώσεως, κατῃσχύνθη πᾶς χρυσοχόος ἀπὸ τῶν γλυπτῶν αὐτοῦ, ὅτι ψευδῆ ἐχώνευσαν, οὐκ ἔστι πνεῦμα ἐν αὐτοῖς.
\vs{18}Μάταιά ἐστιν ἔργα μεμωκημένα, ἐν καιρῷ ἐπισκέψεως αὐτῶν ἀπολοῦνται.
\vs{19}Οὐ τοιαύτη μερὶς τῷ Ἰακὼβ, ὅτι ὁ πλάσας τὰ πάντα, αὐτός ἐστι κληρονομία αὐτοῦ, Κύριος ὄνομα αὐτῷ.

\vs{20}Διασκορπίζεις σύ μοι σκεύη πολέμου, καὶ διασκορπιῶ ἐν σοὶ ἔθνη, καὶ ἐξαρῶ ἐκ σοῦ βασιλεῖς.
\vs{21}Καὶ διασκορπιῶ ἐν σοὶ ἵππον καὶ ἐπιβάτην αὐτοῦ, καὶ διασκορπιῶ ἐν σοὶ ἅρματα καὶ ἀναβάτας αὐτῶν.
\vs{22}Καὶ διασκορπιῶ ἐν σοὶ νεανίσκον καὶ παρθένον, καὶ διασκορπιῶ ἐν σοὶ ἄνδρα καὶ γυναῖκα.
\vs{23}Καὶ διασκορπιῶ ἐν σοὶ ποιμένα καὶ τὸ ποίμνιον αὐτοῦ, καὶ διασκορπιῶ ἐν σοὶ γεωργὸν καὶ τὸ γεώργιον αὐτοῦ, καὶ διασκορπιῶ ἐν σοὶ ἡγεμόνας καὶ στρατηγούς σου.
\vs{24}Καὶ ἀνταποδώσω τῇ Βαβυλῶνι καὶ πᾶσι τοῖς κατοικοῦσι Χαλδαίοις πάσας τὰς κακίας αὐτῶν, ἃς ἐποίησαν ἐπὶ Σιὼν κατʼ ὀφθαλμοὺς ὑμῶν, λέγει Κύριος.

\vs{25}Ἰδοὺ ἐγὼ πρὸς σὲ τὸ ὄρος τὸ διεφθαρμένον, τὸ διαφθεῖρον πᾶσαν τὴν γῆν, καὶ ἐκτενῶ τὴν χεῖρά μου ἐπὶ σὲ, καὶ κατακυλιῶ σε ἐπὶ τῶν πετρῶν, καὶ δώσω σε ὡς ὄρος ἐμπεπυρισμένον·
\vs{26}Καὶ οὐ μὴ λάβωσιν ἀπὸ σοῦ λίθον εἰς γωνίαν, καὶ λίθον εἰς θεμέλιον, ὅτι εἰς ἀφανισμὸν ἔσῃ εἰς τὸν αἰῶνα, λέγει Κύριος.

\vs{27}Ἄρατε σημεῖον ἐπὶ τῆς γῆς, σαλπίσατε ἐν ἔθνεσι σάλπιγγι, ἁγιάσατε ἐπʼ αὐτὴν ἔθνη, παραγγείλατε ἐπʼ αὐτὴν, βασιλεῖς ἄρατε παρʼ ἐμοῦ, καὶ τοῖς Ἀχαναζέοις· ἐπιστήσατε ἐπʼ αὐτὴν βελοστάσεις, ἀναβιβάσατε ἐπʼ αὐτὴν ἵππον ὡς ἀκρίδων πλῆθος.
\vs{28}Ἀναβιβάσατε ἐπʼ αὐτὴν ἔθνη, τὸν βασιλέα τῶν Μήδων καὶ πάσης τῆς γῆς, τοὺς ἡγουμένους αὐτοῦ, καὶ πάντας τοὺς στρατηγοὺς αὐτοῦ.
\vs{29}Ἐσείσθη ἡ γῆ, καὶ ἐπόνεσε, διότα ἐξανέστη ἐπὶ Βαβυλῶνα λογισμὸς Κυρίου, τοῦ θεῖναι τὴν γῆν Βαβυλῶνος εἰς ἀφανισμὸν, καὶ μὴ κατοικεῖσθαι αὐτήν.

\vs{30}Ἐξέλιπε μαχητὴς Βαβυλῶνος τοῦ πολεμεῖν, καθήσονται ἐκεῖ ἐν περιοχῇ, ἐθραύσθη ἡ δυναστεία αὐτῶν, ἐγενήθησαν ὡσεὶ γυναῖκες· ἐνεπυρίσθη τὰ σκηνώματα αὐτῆς, συνετρίβησαν οἱ μοχλοὶ αὐτῆς.
\vs{31}Διώκων εἰς ἀπάντησιν διώκοντος διώξεται, καὶ ἀναγγέλλων εἰς ἀπάντησιν ἀναγγέλλοντος, τοῦ ἀναγγεῖλαι τῷ βασιλεῖ Βαβυλῶνος, ὅτι ἑάλωκεν ἡ πόλις αὐτοῦ.
\vs{32}Ἀπʼ ἐσχάτου τῶν διαβάσεων αὐτοῦ ἐλήφθησαν, καὶ τὰ συστήματα αὐτῶν ἐνέπρησαν ἐν πυρὶ, καὶ ἄνδρες αὐτοῦ οἱ πολεμισταὶ ἐξέρχονται.

\vs{33}Διότι τάδε λέγει Κύριος, οἶκοι βασιλέως Βαβυλῶνος, ὡς ἅλων ὥριμος ἀλοηθήσονται· ἔτι μικρὸν, καὶ ἥξει ὁ ἀμητὸς αὐτῆς.

\vs{34}Κατέφαγέ με, ἐμερίσατό με, κατέλαβέ με σκότος λεπτὸν, Ναβουχοδονόσορ βασιλεὺς Βαβυλῶνος κατέπιέ με, ὡς δράκων ἔπλησε τὴν κοιλίαν αὐτοῦ ἀπὸ τῆς τρυφῆς μου·
\vs{35}Ἔξωσάν με οἱ μόχθοι μου καὶ αἱ ταλαιπωρίαι μου εἰς Βαβυλῶνα, ἐρεῖ κατοικοῦσα Σιὼν, καὶ τὸ αἷμά μου ἐπὶ τοὺς κατοικοῦντας Χαλδαίους, ἐρεῖ Ἱερουσαλήμ.

\vs{36}Διατοῦτο τάδε λέγει Κύριος, ἰδοὺ ἐγὼ κρινῶ τὴν ἀντίδικόν σου, καὶ ἐκδικήσω τὴν ἐκδίκησίν σου, καὶ ἐρημώσω τὴν θάλασσαν αὐτῆς, καὶ ξηρανῶ τὴν πηγὴν αὐτῆς.
\vs{37}Καὶ ἔσται Βαβυλὼν εἰς ἀφανισμὸν, καὶ οὐ κατοικηθήσεται.
\vs{38}Ὅτι ἅμα ὡς λέοντες ἐξηγέρθησαν, καὶ ὡς σκύμνοι λεόντων.
\vs{39}Ἐν τῇ θερμασίᾳ αὐτῶν δώσω πότημα αὐτοῖς, καὶ μεθύσω αὐτοὺς, ὅπως καρωθῶσι, καὶ ὑπνώσωσιν ὕπνον αἰώνιον, καὶ οὐ μὴ ἐξεγερθῶσι, λέγει Κύριος.
\vs{40}Καὶ καταβίβασον αὐτοὺς ὡς ἄρνας εἰς σφαγὴν, καὶ ὡς κριοὺς μετʼ ἐρίφων.

\vs{41}Πῶς ἑάλω καὶ ἐθηρεύθη τὸ καύχημα πάσης τῆς γῆς; πῶς ἐγένετο Βαβυλὼν εἰς ἀφανισμὸν ἐν τοῖς ἔθνεσιν;
\vs{42}Ἀνέβη ἐπὶ Βαβυλῶνα ἡ θάλασσα ἐν ἤχῳ κυμάτων αὐτῆς, καὶ κατεκαλύφθη,
\vs{43}ἐγενήθησαν αἱ πόλεις αὐτῆς ὡς γῆ ἄνυδρος καὶ ἄβατος, οὐ κατοικήσει ἐν αὐτῇ οὐδὲ εἷς, οὐδὲ μὴ καταλύσει ἐν αὐτῇ υἱὸς ἀνθρώπου.
\vs{44}Καὶ ἐκδικήσω ἐπὶ Βαβυλῶνα, καὶ ἐξοίσω ἃ κατέπιεν ἐκ τοῦ στόματος αὐτῆς, καὶ οὐ μὴ συναχθῶσι πρὸς αὐτὴν ἔτι τὰ ἔθνη,
\vs{49}καὶ ἐν Βαβυλῶνι πεσοῦνται τραυματίαι πάσης τῆς γῆς.
\vs{50}Ἀνασωζόμενοι ἐκ γῆς πορεύεσθε, καὶ μὴ ἵστασθε· οἱ μακρόθεν μνήσθητε τοῦ Κυρίου, καὶ Ἱερουσαλὴμ ἀναβήτω ἐπὶ τὴν καρδίαν ὑμῶν.
\vs{51}Ἠσχύνθημεν ὅτι ἠκούσαμεν ὀνειδισμὸν ἡμῶν, κατεκάλυψεν ἀτιμία τὸ πρόσωπον ἡμῶν, εἰσῆλθον ἀλλογενεῖς εἰς τὰ ἅγια ἡμῶν, εἰς οἶκον Κυρίου.

\vs{52}Διατοῦτο ἰδοὺ ἡμέραι ἔρχονται, λέγει Κύριος, καὶ ἐκδικήσω ἐπὶ τὰ γλυπτὰ αὐτῆς, καὶ ἐν πάσῃ τῇ γῇ αὐτῆς πεσοῦνται τραυματίαι.
\vs{53}Ὅτι ἐὰν ἀναβῇ Βαβυλὼν ὡς ὁ οὐρανὸς, καὶ ὅτι ἐὰν ὀχυρώσῃ τὰ τείχη ἰσχύϊ αὐτῆς, παρʼ ἐμοῦ ἥξουσιν ἐξολοθρεύοντες αὐτὴν, λέγει Κύριος.
\vs{54}Φωνὴ κραυγῆς ἐν Βαβυλῶνι, καὶ συντριβὴ μεγάλη ἐν γῇ Χαλδαίων,
\vs{55}ὅτι ἐξωλόθρευσε Κύριος τὴν Βαβυλῶνα, καὶ ἀπώλεσεν ἀπʼ αὐτῆς φωνὴν μεγάλην ἠχοῦσαν ὡς ὕδατα πολλά· ἔδωκεν εἰς ὄλεθρον φωνὴν αὐτῆς.
\vs{56}Ὅτι ἦλθεν ἐπὶ Βαβυλῶνα ταλαιπωρία, ἑάλωσαν οἱ μαχηταὶ αὐτῆς, ἐπτόηται τὸ τόξον αὐτῶν, ὅτι ὁ Θεὸς ἀνταποδίδωσιν αὐτοῖς.
\vs{57}Κύριος ἀνταποδίδωσι, καὶ μεθύσει μέθῃ τοὺς ἡγεμόνας αὐτῆς, καὶ τοὺς σοφοὺς αὐτῆς, καὶ τοὺς στρατηγοὺς αὐτῆς, λέγει ὁ βασιλεὺς, Κύριος παντοκράτωρ ὄνομα αὐτῷ.

\vs{58}Τάδε λέγει Κύριος, τεῖχος Βαβυλῶνος ἐπλατύνθη, κατασκαπτόμενον κατασκαφήσεται, καὶ αἱ πύλαι αὐτῆς αἱ ὑψηλαὶ ἐμπυρισθήσονται, καὶ οὐ κοπιάσουσι λαοὶ εἰς κενὸν, καὶ ἔθνη ἐν ἀρχῇ ἐκλείψουσιν.

\vs{59}Ὁ ΛΟΓΟΣ ὋΝ ἘΝΕΤΕΙΛΑΤΟ ΚΥΡΙΟΣ ἹΕΡΕΜΙΑ ΤΩ ΠΡΟΦΗΤΗ εἰπεῖν τῷ Σαραίᾳ υἱῷ Νηρείου, υἱοῦ Μαασαίου, ὅτε ἐπορεύετο παρὰ Σεδεκίου βασιλέως Ἰούδα εἰς Βαβυλῶνα, ἐν τῷ ἔτει τῷ τετάρτῳ τῆς βασιλείας αὐτοῦ· καὶ Σαραίας ἄρχων δώρων.
\vs{60}Καὶ ἔγραψεν Ἱερεμίας πάντα τὰ κακὰ ἃ ἥξει ἐπὶ Βαβυλῶνα ἐν βιβλίῳ, πάντας τοὺς λόγους τούτους τοὺς γεγραμμένους ἐπὶ Βαβυλῶνα.
\vs{61}Καὶ εἶπεν Ἱερεμίας πρὸς Σαραίαν, ὅταν ἔλθῃς εἰς Βαβυλῶνα, καὶ ὄψῃ καὶ ἀναγνώσῃ πάντας τοὺς λόγους τούτους,
\vs{62}καὶ ἐρεῖς, Κύριε Κύριε, σὺ ἐλάλησας ἐπὶ τὸν τόπον τοῦτον, τοῦ ἐξολοθρεῦσαι αὐτὸν, καὶ τοῦ μὴ εἶναι ἐν αὐτῷ κατοικοῦντας ἀπὸ ἀνθρώπου ἕως κτήνους, ὅτι ἀφανισμὸς εἰς τὸν αἰῶνα ἔσται.
\vs{63}Καὶ ἔσται ὅταν παύσῃ τοῦ ἀναγινώσκειν τὸ βιβλίον τοῦτο, καὶ ἐπιδήσεις ἐπʼ αὐτὸ λίθον, καὶ ῥίψεις αὐτὸ εἰς μέσον τοῦ Εὐφράτου,
\vs{64}καὶ ἐρεῖς, οὕτως καταδύσεται Βαβυλὼν, καὶ οὐ μὴ ἀναστῇ ἀπὸ προσώπου τῶν κακῶν, ὧν ἐγὼ ἐπάγω ἐπʼ αὐτήν.

\ch{29}
ἘΠΙ ΤΟΥΣ ἈΛΛΟΦΥΛΟΥΣ ΤΑΔΕ ΛΕΓΕΙ ΚΥΡΙΟΣ.

\vs{2}Ἰδοὺ ὕδατα ἀναβαίνει ἀπὸ Βοῤῥᾶ, καὶ ἔσται εἰς χειμάῤῥουν κατακλύζοντα, καὶ κατακλύσει γῆν καὶ τὸ πλήρωμα αὐτῆς, πόλιν καὶ τοὺς κατοικοῦντας ἐν αὐτῇ· καὶ κεκράξονται οἱ ἄνθρωποι, καὶ ἀλαλάξουσιν ἅπαντες οἱ κατοικοῦντες τὴν γῆν,
\vs{3}ἀπὸ φωνῆς ὁρμῆς αὐτοῦ, ἀπὸ τῶν ὅπλων τῶν ποδῶν αὐτοῦ, καὶ ἀπὸ σεισμοῦ τῶν ἁρμάτων αὐτοῦ, ἤχου τροχῶν αὐτοῦ· οὐκ ἐπέστρεψαν πατέρες ἐφʼ υἱοὺς αὐτῶν ἀπὸ ἐκλύσεως χειρῶν αὐτῶν
\vs{4}ἐν τῇ ἡμέρᾳ τῇ ἐπερχομένῃ τοῦ ἀπολέσαι πάντας τοὺς ἀλλοφύλους· καὶ ἀφανιῶ τὴν Τύρον, καὶ τὴν Σιδῶνα, καὶ πάντας τοὺς καταλοίπους τῆς βοηθείας αὐτῶν, ὅτι ἐξολοθρεύσει Κύριος τοὺς καταλοίπους τῶν νήσων.
\vs{5}Ἥκει φαλάκρωμα ἐπὶ Γάζαν, ἀπεῤῥίφη Ἀσκάλων, καὶ οἱ κατάλοιποι Ἐνακίμ.

\vs{6}Ἕως τίνος κόψεις ἡ μάχαιρα τοῦ Κυρίου; ἕως τίνος οὐχ ἡσυχάσεις; ἀποκατάστηθι εἰς τὸν κολεόν σου, ἀνάπαυσαι, καὶ ἐπάρθητι.

\vs{7}Πῶς ἡσυχάσει, καὶ Κύριος ἐνετείλατο αὐτῇ ἐπὶ τὴν Ἀσκάλωνα, καὶ ἐπὶ τὰς παραθαλασσίους, ἐπὶ τὰς καταλοίπους ἐπεγερθῆναι;

\ch{30}
ΤΗ ἸΔΟΥΜΑΙΑ, τάδε λέγει Κύριος, οὐκ ἔστιν ἔτι σοφία ἐν Θαιμὰν, ἀπώλετο βουλὴ ἐκ συνετῶν, ᾤχετο σοφία αὐτῶν, ἠπατήθη ὁ τόπος αὐτῶν·
\vs{2}βαθύνατε εἰς κάθισιν οἱ κατοικοῦντες ἐν Δαιδὰμ, ὅτι δύσκολα ἐποίησεν· ἤγαγον ἐπʼ αὐτὸν ἐν χρόνῳ ᾧ ἐπεσκεψάμην ἐπʼ αὐτόν.
\vs{3}Ὅτι τρυγηταὶ ἦλθον, οἳ οὐ καταλείψουσί σοι κατάλειμμα· ὡς κλέπται ἐν νυκτὶ, ἑπιθήσουσι χεῖρα αὐτῶν.

\vs{4}Ὅτι ἐγὼ κατέσυρα τὸν Ἡσαῦ, ἀνεκάλυψα τὰ κρυπτὰ αὐτῶν, κρυβῆναι οὐ μὴ δύνωνται, ὤλοντο διὰ χεῖρα ἀδελφοῦ αὐτοῦ, γείτονός μου, καὶ οὐκ ἔστιν
\vs{5}ὑπολείπεσθαι ὀρφανόν σου, ἵνα ζήσηται· καὶ ἐγὼ ζήσομαι, καὶ αἱ χῆραι ἐπʼ ἐμὲ πεποίθασιν·

\vs{6}Ὅτι τάδε εἶπε Κύριος, οἷς οὐκ ἦν νόμος πιεῖν τὸ ποτήριον, ἔπιον· καὶ σὺ ἀθωωμένη οὐ μὴ ἀθωωθῇς,
\vs{7}ὅτι κατʼ ἐμαυτοῦ ὤμοσα, λέγει Κύριος, ὅτι εἰς ἄβατον καὶ εἰς ὀνειδισμὸν, καὶ εἰς κατάρασιν ἔσῃ ἐν μέσῳ αὐτῆς, καὶ πᾶσαι αἱ πόλεις αὐτῆς ἔσονται ἔρημοι εἰς αἰῶνα.

\vs{8}Ἀκοὴν ἤκουσα παρὰ Κυρίου, καὶ ἀγγέλους εἰς ἔθνη ἀπέστειλε, συνάχθητε, καὶ παραγένεσθε εἰς αὐτὴν, ἀνάστητε εἰς πόλεμον.
\vs{9}Μικρὸν ἔδωκά σε ἐν ἔθνεσιν, εὐκαταφρόνητον ἐν ἀνθρώποις.
\vs{10}Ἡ παιγνία σου ἐνεχείρησέ σοι, ἰταμία καρδίας σου κατέλυσε τρυμαλιὰς πετρῶν, συνέλαβεν ἰσχὺν βουνοῦ ὑψηλοῦ· ὅτι ὕψωσεν ὥσπερ ἀετὸς νοσσιὰν αὐτοῦ, ἐκεῖθεν καθελῶ σε.

\vs{11}Καὶ ἔσται ἡ Ἰδουμαία εἰς ἄβατον, πᾶς ὁ παραπορευόμενος ἐπʼ αὐτὴν συριεῖ.
\vs{12}Ὥσπερ κατεστράφη Σόδομα καὶ Γόμοῤῥα, καὶ αἱ πάροικοι αὐτῆς, εἶπε Κύριος παντοκράτωρ, οὐ μὴ καθίσει ἐκεῖ ἄνθρωπος, καὶ οὐ μὴ κατοικήσει ἐκεῖ υἱὸς ἀνθρώπου.
\vs{13}Ἰδοὺ ὥσπερ λέων ἀναβήσεται ἐκ μέσου τοῦ Ἰορδάνου εἰς τόπον Αἰθὰμ, ὅτι ταχὺ ἐκδιώξω αὐτοὺς ἀπʼ αὐτῆς, καὶ τοὺς νεανίσκους ἐπʼ αὐτὴν ἐπιστήσατε· ὅτι τίς ὥσπερ ἐγώ; καὶ τίς ἀντιστήσεταί μοι; καὶ τίς οὗτος ποιμὴν, ὃς στήσεται κατὰ πρόσωπόν μου;

\vs{14}Διατοῦτο ἀκούσατε βουλὴν Κυρίου, ἣν ἐβουλεύσατο ἐπὶ τὴν Ἰδουμαίαν, καὶ λογισμὸν αὐτοῦ, ὃν ἐλογίσατο ἐπὶ τοὺς κατοικοῦντας Θαιμὰν, ἐὰν μὴ συμψησθῶσι τὰ ἐλάχιστα τῶν προβάτων, ἐὰν μὴ ἀβατωθῇ ἐπʼ αὐτοὺς κατάλυσις αὐτῶν,
\vs{15}ὅτι ἀπὸ φωνῆς πτώσεως αὐτῶν ἐφοβήθη ἡ γῆ, καὶ κραυγὴ θαλάσσης οὐκ ἠκούσθη.
\vs{16}Ἰδοὺ ὥσπερ ἀετὸς ὄψεται, καὶ ἐκτενεῖ τὰς πτέρυγας ἐπʼ ὀχυρώματα αὐτῆς· καὶ ἔσται ἡ καρδία τῶν ἰσχυρῶν τῆς Ἰδουμαίας ἐν τῇ ἡμέρᾳ ἐκείνῃ, ὡς καρδία γυναικὸς ὠδινούσης.

\vs{17}ΤΟΙΣ ΥΙΟΙΣ ἈΜΜΩΝ οὕτως εἶπε Κύριος, μὴ υἱοὶ οὐκ εἰσὶν ἐν Ἰσραὴλ, ἢ παραληψόμενος οὐκ ἔστιν αὐτοῖς; διατί παρέλαβε Μελχὸλ τὴν Γαλαὰδ, καὶ ὁ λαὸς αὐτῶν ἐν πόλεσιν αὐτῶν ἐνοικήσει;
\vs{18}Διατοῦτο ἰδοὺ ἡμέραι ἔρχονται, φησὶ Κύριος, καὶ ἀκουτιῶ ἐπὶ Ῥαββὰθ θόρυβον πολέμων, καὶ ἔσονται εἰς ἄβατον καὶ εἰς ἀπώλειαν, καὶ βωμοὶ αὐτῆς ἐν πυρὶ κατακαυθήσονται, καὶ παραλήψεται Ἰσραὴλ τὴν ἀρχὴν αὐτοῦ.
\vs{19}Ἀλάλαξον Ἐσεβὼν, ὅτι ὤλετο Γαΐ· κεκράξατε θυγατέρες Ῥαββὰθ, περιζώσασθε σάκκους καὶ κόψασθε, ὅτι Μελχὸλ βαδιεῖται ἐν ἀποικίᾳ, οἱ ἱερεῖς αὐτοῦ καὶ οἱ ἄρχοντες αὐτοῦ ἅμα.

\vs{20}Τί ἀγαλλιᾶσθε ἐν τοῖς πεδίοις Ἐνακεὶμ, θύγατερ ἰταμίας, ἡ πεποιθυῖα ἐπὶ θησαυροῖς, ἡ λέγουσα, τίς εἰσελεύσεται ἐπʼ ἐμέ;
\vs{21}Ἰδοὺ ἐγὼ φέρω φόβον ἐπὶ σὲ, εἶπε Κύριος, ἀπὸ πάσης τῆς περιοίκου σου, καὶ διασπαρήσεσθε ἕκαστος εἰς τὸ πρόσωπον αὐτοῦ, καὶ οὐκ ἔστιν ὁ συνάγων.

\vs{23}ΤΗ ΚΗΔΑΡ ΤΗ ΒΑΣΙΛΙΣΣΗ ΤΗΣ ΑΥΛΗΣ, ἫΝ ἘΠΑΤΑΞΕ ΝΑΒΟΥΧΟΔΟΝΟΣΟΡ ΒΑΣΙΛΕΥΣ ΒΑΒΥΛΩΝΟΣ, οὕτως εἶπε Κύριος,

Ἀνάστητε, καὶ ἀνάβητε ἐπὶ Κηδὰρ, καὶ πλήσατε τοὺς υἱοὺς Κεδέμ.
\vs{24}Σκηνὰς αὐτῶν, καὶ τὰ πρόβατα αὐτῶν λήψονται· ἱμάτια αὐτῶν, καὶ πάντα τὰ σκεύη αὐτῶν, καὶ καμήλους αὐτῶν λήψονται ἑαυτοῖς· καὶ καλέσατε ἐπʼ αὐτοὺς ἀπώλειαν κυκλόθεν.
\vs{25}Φεύγετε, λίαν ἐμβαθύνατε εἰς κάθισιν, καθήμενοι ἐν τῇ αὐλῇ, ὅτι ἐβουλεύσατο ἐφʼ ὑμᾶς βασιλεὺς Βαβυλῶνος βουλὴν, καὶ ἐλογίσατο λογισμόν.

\vs{26}Ἀνάστηθι, καὶ ἀνάβηθι ἐπʼ ἔθνος εὐσταθοῦν, καθήμενον εἰς ἀναψυχὴν, οἷς οὐκ εἰσὶ θύραι, οὐ βάλανοι, οὐ μοχλοὶ, μόνοι καταλύουσι.
\vs{27}Καὶ ἔσονται κάμηλοι αὐτῶν εἰς προνομὴν, καὶ πλῆθος κτηνῶν αὐτῶν εἰς ἀπώλειαν, καὶ λικμήσω αὐτοὺς παντὶ πνεύματι κεκαρμένους πρὸ προσώπου αὐτῶν, ἐκ παντὸς πέραν αὐτῶν οἴσω τὴν τροπὴν αὐτῶν, εἶπε Κύριος.
\vs{28}Καὶ ἔσται ἡ αὐλὴ διατριβὴ στρουθῶν, καὶ ἄβατος ἕως αἰῶνος, οὐ μὴ καθίσῃ ἐκεῖ ἄνθρωπος, καὶ οὐ μὴ κατοικήσει ἐκεῖ υἱὸς ἀνθρώπου.

\vs{29}ΤΗ ΔΑΜΑΣΚΩ. Κατῃσχύνθη Ἠμὰθ, καὶ Ἀρφὰθ, ὅτι ἤκουσαν ἀκοὴν πονηρὰν, ἐξέστησαν, ἐθυμώθησαν, ἀναπαύσασθαι οὐ μὴ δύνωνται.
\vs{30}Ἐξελύθη Δαμασκὸς, ἀπεστράφη εἰς φυγὴν, τρόμος ἐπελάβετο αὐτῆς.
\vs{31}Πῶς οὐχὶ ἐγκατέλιπε πόλιν ἐμὴν, κώμην ἠγάπησαν;

\vs{32}Διατοῦτο πεσοῦνται νεανίσκοι ἐν πλατείαις σου, καὶ πάντες οἱ ἄνδρες οἱ πολεμισταί σου πεσοῦνται, φησὶ Κύριος·
\vs{33}Καὶ καύσω πῦρ ἐν τείχει Δαμασκοῦ, καὶ καταφάγεται ἄμφοδα υἱοῦ Ἄδερ.

\ch{31}
ΤΗ ΜΩΑΒ οὕτως εἶπε Κύριος, οὐαὶ ἐπὶ Ναβαῦ, ὅτι ὤλετο, ἐλήφθη Καριαθαὶμ, ᾐσχύνθη Ἀμὰθ καὶ Ἀγάθ.
\vs{2}Οὐκ ἔστιν ἔτι ἰατρεία Μωὰβ, γαυρίαμα ἐν Ἐσεβὼν, ἐλογίσατο ἐπʼ αὐτὴν κακά· ἐκόψαμεν αὐτὴν ἀπὸ ἔθνους, καὶ παῦσιν παύσεται· ὄπισθέν σου βαδιεῖται μάχαιρα,
\vs{3}ὅτι φωνὴ κεκραγότων ἐξ Ὠρωναὶμ, ὄλεθρον καὶ σύντριμμα μέγα·
\vs{4}Συνετρίβη Μωὰβ, ἀναγγείλατε εἰς Ζογόρα,
\vs{5}ὅτι ἐπλήσθη Ἀλὼθ ἐν κλαυθμῷ· ἀναβήσεται κλαίων ἐν ὁδῷ Ὠρωναὶμ, κραυγὴν συντρίμματος ἠκούσατε.

\vs{6}Φεύγετε καὶ σώσατε τὰς ψυχὰς ὑμῶν, καὶ ἔσεσθε ὥσπερ ὄνος ἄγριος ἐν ἐρήμῳ.
\vs{7}Ἐπειδὴ ἐπεποίθεις ἐν ὀχυρώματί σου, καὶ σὺ συλληφθήσῃ· καὶ ἐξελεύσεται Χαμὼς ἐν ἀποικίᾳ, καὶ οἱ ἱερεῖς αὐτοῦ, καὶ οἱ ἄρχοντες αὐτοῦ ἅμα.
\vs{8}Καὶ ἥξει ὄλεθρος ἐπὶ πᾶσαν πόλιν, οὐ μὴ σωθῇ, καὶ ἀπολεῖται ὁ αὐλῶν, καὶ ἐξολοθρευθήσεται ἡ πεδινὴ, καθὼς εἶπε Κύριος.
\vs{9}Δότε σημεῖα τῇ Μωὰβ, ὅτι ἁφῇ ἁφθήσεται, καὶ πᾶσαι αἱ πόλεις αὐτῆς εἰς ἄβατον ἔσονται· πόθεν ἔνοικος αὐτῇ;
\vs{10}Ἐπικατάρατος ὁ ποιῶν τὰ ἔργα Κυρίου ἀμελῶς, ἐξαίρων μάχαιραν αὐτοῦ ἀφʼ αἵματος.

\vs{11}Ἀνεπαύσατο Μωὰβ ἐκ παιδαρίου, καὶ πεποιθὼς ἦν ἐπὶ τῇ δόξῃ αὐτοῦ, οὐκ ἐνέχεεν ἐξ ἀγγείου εἰς ἀγγεῖον, καὶ εἰς ἀποικισμὸν οὐκ ᾤχετο· διατοῦτο ἔστη γεῦμα αὐτοῦ ἐν αὐτῷ, καὶ ὀσμὴ αὐτοῦ οὐκ ἐξέλιπε.
\vs{12}Διατοῦτο ἰδοὺ ἡμέραι αὐτοῦ ἔρχονται, φησὶ Κύριος, καὶ ἀποστελῶ αὐτῷ κλίνοντας, καὶ κλινοῦσιν αὐτὸν, καὶ τὰ σκεύη αὐτοῦ λεπτυνοῦσι, καὶ τὰ κέρατα αὐτοῦ συγκόψουσι.
\vs{13}Καὶ καταισχυνθήσεται Μωὰβ ἀπὸ Χαμὼς, ὥσπερ κατῃσχύνθη οἶκος Ἰσραὴλ ἀπὸ Βαιθὴλ ἐλπίδος αὐτῶν πεποιθότες ἐπʼ αὐτοῖς.

\vs{14}Πῶς ἐρεῖτε, ἰσχυροί ἐσμεν, καὶ ἄνθρωπος ἰσχύων εἰς τὰ πολεμικά;
\vs{15}Ὤλετο Μωὰβ πόλις αὐτοῦ, καὶ ἐκλεκτοὶ νεανίσκοι αὐτοῦ κατέβησαν εἰς σφαγήν.
\vs{16}Ἐγγὺς ἡμέρα Μωὰβ ἐλθεῖν, καὶ πονηρία αὐτοῦ ταχεῖα σφόδρα.
\vs{17}Κινήσατε αὐτῷ πάντες κυκλόθεν αὐτοῦ, πάντες ἔκδοτε ὄνομα αὐτοῦ· εἴπατε, πῶς συνετρίβη βακτηρία εὐκλεὴς, ῥάβδος μεγαλώματος;

\vs{18}Κατάβηθι ἀπὸ δόξης, καὶ κάθισον ἐν ὑγρασίᾳ καθημένη· Δαιβὼν ἐκτριβήσεται, ὅτι ὤλετο Μωὰβ, ἀνέβη εἰς σὲ λυμαινόμενος ὀχύρωμά σου.
\vs{19}Ἐφʼ ὁδοῦ στῆθι, καὶ ἔπιδε καθημένη ἐν Ἀρὴρ, καὶ ἐρώτησον φεύγοντα, καὶ σωζόμενον, καὶ εἰπὸν, τί ἐγένετο;

\vs{20}Κατῃσχύνθη Μωὰβ, ὅτι συνετρίβη· ὀλόλυξον καὶ κέκραξον, ἀνάγγειλον ἐν Ἀρνὼν, ὅτι ὤλετο Μωὰβ,
\vs{21}καὶ κρίσις ἔρχεται εἰς τὴν γῆν Μεισὼρ ἐπὶ Χελὼν, καὶ Ῥεφὰς, καὶ Μωφὰς,
\vs{22}καὶ ἐπὶ Δαιβὼν, καὶ ἐπὶ Ναβαῦ, καὶ ἐπʼ οἶκον Δαιθλαθαὶμ,
\vs{23}καὶ ἐπὶ Καριαθαὶμ, καὶ ἐπʼ οἶκον Γαιμὼλ, καὶ ἐπʼ οἶκον Μαὼν,
\vs{24}καὶ ἐπὶ Καριὼθ, καὶ ἐπὶ Βοσὸρ, καὶ ἐπὶ πάσας τὰς πόλεις Μωὰβ τὰς πόῤῥω καὶ τὰς ἐγγύς.
\vs{25}Κατεάχθη κέρας Μωὰβ, καὶ τὸ ἐπίχειρον αὐτοῦ συνετρίβη.

\vs{26}Μεθύσατε αὐτὸν, ὅτι ἐπὶ Κύριον ἐμεγαλύνθη· καὶ ἐπικρούσει Μωὰβ ἐν χειρὶ αὐτοῦ, καὶ ἔσται εἰς γέλωτα καὶ αὐτός.
\vs{27}Καὶ εἰ μὴ εἰς γελοιασμὸν ἦν σοι Ἰσραὴλ, καὶ ἐν κλοπαῖς σου εὑρέθη, ὅτι ἐπολέμεις αὐτόν.
\vs{28}Κατέλιπον τὰς πόλεις, καὶ ᾤκησαν ἐν πέτραις οἱ κατοικοῦντες Μωάβ· ἐγενήθησαν ὥσπερ περιστεραὶ νοσσεύουσαι ἐν πέτραις, στόματι βοθύνου.

\vs{29}Καὶ ἤκουσα ὕβριν Μωὰβ, ὕβρισε λίαν ὕβριν αὐτοῦ, καὶ ὑπερηφανίαν αὐτοῦ· καὶ ὑψώθη ἡ καρδία αὐτοῦ.
\vs{30}Ἐγὼ δὲ ἔγνων ἔργα αὐτοῦ· οὐχὶ τὸ ἱκανὸν αὐτῷ οὐχ οὕτως ἐποίησε;

\vs{31}Διατοῦτο ἐπὶ Μωὰβ ὀλολύζετε πάντοθεν· βοήσατε ἐπʼ ἄνδρας κειράδας αὐχμοῦ.
\vs{32}Ὡς κλαυθμὸν Ἰαζὴρ ἀποκλαύσομαί σοι ἄμπελος Ἀσερημὰ, κλήματά σου διῆλθε θάλασσαν, πόλεις Ἰαζὴρ ἥψαντο, ἐπὶ ὀπώραν σου, ἐπὶ τρυγηταῖς σου ὄλεθρος ἐπέπεσε.
\vs{33}Συνεψήσθη χαρμοσύνη καὶ εὐφροσύνη ἐκ τῆς Μωαβίτιδος· καὶ οἶνος ἦν ἐπὶ ληνοῖς σου, πρωῒ οὐκ ἐπάτησαν, οὐδὲ δείλης οὐκ ἐποίησαν,
\vs{34}αἱ δὲ ἀπὸ κραυγῆς Ἐσεβὼν ἕως Αἰτὰμ αἱ πόλεις αὐτῶν ἔδωκαν φωνὴν αὐτῶν, ἀπὸ Ζογὸρ ἕως Ὠρωναὶμ, καὶ ἀγγελίαν σαλασία, ὅτι καὶ τὸ ὕδωρ Νεβρεὶν εἰς κατάκαυμα ἔσται.

\vs{35}Καὶ ἀπολῶ τὸν Μωὰβ, φησὶ Κύριος, ἀναβαίνοντα ἐπὶ τὸν βωμὸν, καὶ θυμιῶντα θεοῖς αὐτοῦ.
\vs{36}Διατοῦτο καρδία τοῦ Μωὰβ, ὥσπερ αὐλοὶ βουβήσουσι, καρδία μου ἐπʼ ἀνθρώπους κειράδας ὥσπερ αὐλὸς βομβήσει· διατοῦτο ἃ περιεποιήσατο, ἀπώλετο ἀπὸ ἀνθρώπου.
\vs{37}Πᾶσαν κεφαλὴν ἐν παντὶ τόπῳ ξυρηθήσονται, καὶ πᾶς πώγων ξυρηθήσεται, καὶ πᾶσαι χεῖρες κόψονται, καὶ ἐπὶ πάσης ὀσφύος σάκκος.
\vs{38}Καὶ ἐπὶ πάντων τῶν δωμάτων Μωὰβ, καὶ ἐπὶ ταῖς πλατείαις αὐτῆς, ὅτι συνέτριψα, φησὶ Κύριος, ὡς ἀγγεῖον, οὗ οὐκ ἔστι χρεία αὐτοῦ.
\vs{39}Πῶς κατήλλαξε; πῶς ἔστρεψε νῶτον Μωάβ; ᾐσχύνθη, καὶ ἐγένετο Μωὰβ εἰς γέλωτα, καὶ ἐγκότημα πᾶσι τοῖς κύκλῳ αὐτῆς.

\vs{40}Ὅτι οὕτως εἶπε Κύριος,
\vs{41}Ἐλήφθη Καριὼθ, καὶ τὰ ὀχυρώματα συνελήφθη,
\vs{42}καὶ ἀπολεῖται Μωὰβ ἀπὸ ὄχλου, ὅτι ἐπὶ τὸν Κύριον ἐμεγαλύνθη.
\vs{43}Παγὶς καὶ φόβος καὶ βόθυνος ἐπὶ σὲ καθήμενος Μωάβ.
\vs{44}Ὁ φεύγων ἀπὸ προσώπου τοῦ φόβου, ἐμπεσεῖται εἰς τὸν βόθυνον· καὶ ὁ ἀναβαίνων ἐκ τοῦ βοθύνου, καὶ συλληφθήσεται ἐν τῇ παγίδι· ὅτι ἐπάξω ταῦτα ἐπὶ Μωὰβ ἐν ἐνιαυτῷ ἐπισκέψεως αὐτῶν.

\ch{32}
\vs{15}Οὕτως εἶπε Κύριος ὁ Θεὸς Ἰσραὴλ, λάβε τὸ ποτήριον τοῦ οἴνου τοῦ ἀκράτου τούτου ἐκ χειρός μου, καὶ ποτιεῖς πάντα τὰ ἔθνη, πρὸς ἃ ἐγὼ ἀποστέλλω σε πρὸς αὐτούς.
\vs{16}Καὶ πίονται, καὶ ἐξεμοῦνται, καὶ ἐκμανήσονται ἀπὸ προσώπου τῆς μαχαίρας, ἧς ἐγὼ ἀποστέλλω ἀναμέσον αὐτῶν.

\vs{17}Καὶ ἔλαβον τὸ ποτήριον ἐκ χειρὸς Κυρίου, καὶ ἐπότισα τὰ ἔθνη, πρὸς ἃ ἀπέστειλέ με Κύριος πρὸς αὐτὰ,
\vs{18}τὴν Ἱερουσαλὴμ, καὶ τὰς πόλεις Ἰούδα, καὶ βασιλεῖς Ἰούδα καὶ ἄρχοντας αὐτοῦ, τοῦ θεῖναι αὐτὰς εἰς ἐρήμωσιν, καὶ εἰς ἄβατον, καὶ εἰς συριγμὸν,
\vs{19}καὶ τὸν Φαραὼ βασιλέα Αἰγύπτου, καὶ τοὺς παῖδας αὐτοῦ, καὶ τοὺς μεγιστᾶνας αὐτοῦ,
\vs{20}καὶ πάντα τὸν λαὸν αὐτοῦ, καὶ πάντας τοὺς συμμίκτους, καὶ πάντας τοὺς βασιλεῖς ἀλλοφύλων, καὶ τὴν Ἀσκάλωνα, καὶ τὴν Γάζαν, καὶ τὴν Ἀκκάρων, καὶ τὸ ἐπίλοιπον Ἀζώτου,
\vs{21}καὶ τὴν Ἰδουμαίαν, καὶ τὴν Μωαβῖτιν, καὶ τοὺς υἱοὺς Ἀμμὼν,
\vs{22}καὶ βασιλεῖς Τύρου, καὶ βασιλεῖς Σιδῶνος, καὶ βασιλεῖς τοὺς ἐν τῷ πέραν τῆς θαλάσσης,
\vs{23}καὶ τὴν Δαιδὰν, καὶ τὴν Θαιμὰν, καὶ τὴν Ῥῶς, καὶ πᾶν περικεκαρμένον κατὰ πρόσωπον αὐτοῦ,
\vs{24}καὶ πάντας τοὺς συμμίκτους τοὺς καταλύοντας ἐν τῇ ἐρήμῳ,
\vs{25}καὶ πάντας βασιλεῖς Αἰλὰμ, καὶ πάντας βασιλεῖς Περσῶν,
\vs{26}καὶ πάντας βασιλεῖς ἀπὸ ἀπηλιώτου τοὺς πόῤῥω καὶ τοὺς ἐγγὺς, ἕκαστον πρὸς τὸν ἀδελφὸν αὐτοῦ, καὶ πάσας βασιλείας τὰς ἐπὶ προσώπου τῆς γῆς.

\vs{27}Καὶ ἐρεῖς αὐτοῖς, οὕτως εἶπε Κύριος παντοκράτωρ, πίετε, μεθύσθητε, καὶ ἐξεμέσετε, καὶ πεσεῖσθε, καὶ οὐ μὴ ἀναστῆτε ἀπὸ προσώπου τῆς μαχαίρας, ἧς ἐγὼ ἀποστέλλω ἀναμέσον ὑμῶν.
\vs{28}Καὶ ἔσται ὅταν μὴ βούλωνται δέξασθαι τὸ ποτήριον ἐκ τῆς χειρός σου, ὥστε πιεῖν, καὶ ἐρεῖς, οὕτως εἶπε Κύριος, πιόντες πίεσθε,
\vs{29}ὅτι ἐν πόλει ἐν ᾗ ὠνομάσθη τὸ ὄνομά μου ἐπʼ αὐτὴν, ἐγὼ ἄρχομαι κακῶσαι, καὶ ὑμεῖς καθάρσει οὐ μὴ καθαρισθῆτε, ὅτι μάχαιραν ἐγὼ καλῶ ἐπὶ πάντας τοὺς καθημένους ἐπὶ τῆς γῆς.

\vs{30}Καὶ σὺ προφητεύσεις ἐπʼ αὐτοὺς τοὺς λόγους τούτους, καὶ ἐρεῖς, Κύριος ἀφʼ ὑψηλοῦ χρηματιεῖ, ἀπὸ τοῦ ἁγίου αὐτοῦ δώσει φωνὴν αὐτοῦ, λόγον χρηματιεῖ ἐπὶ τοῦ τόπου αὐτοῦ· καὶ οἵδε ὥσπερ τρυγῶντες ἀποκριθήσονται· καὶ ἐπὶ καθημένους ἐπὶ τὴν γῆν
\vs{31}ἥκει ὄλεθρος ἐπὶ μέρος τῆς γῆς, ὅτι κρίσις τῷ Κυρίῳ ἐν τοῖς ἔθνεσι· κρίνεται αὐτὸς πρὸς πᾶσαν σάρκα, οἱ δὲ ἀσεβεῖς ἐδόθησαν εἰς μάχαιραν, λέγει Κύριος.

\vs{32}Οὕτως εἶπε Κύριος, ἰδοὺ κακὰ ἔρχεται ἀπὸ ἔθνους ἐπὶ ἔθνος, καὶ λαίλαψ μεγάλη ἐκπορεύεται ἀπʼ ἐσχάτου τῆς γῆς.
\vs{33}Καὶ ἔσονται τραυματίαι ὑπὸ Κυρίου ἐν ἡμέρᾳ Κυρίου, ἐκ μέρους τῆς γῆς, καὶ ἕως εἰς μέρος τῆς γῆς· οὐ μὴ κατορυγῶσιν, εἰς κόπρια ἐπὶ προσώπου τῆς γῆς ἔσονται.
\vs{34}Ἀλαλάξατε ποιμένες, καὶ κεκράξατε, καὶ κόπτεσθε οἱ κριοὶ τῶν προβάτων, ὅτι ἐπληρώθησαν αἱ ἡμέραι ὑμῶν εἰς σφαγὴν, καὶ πεσεῖσθε ὥσπερ οἱ κριοὶ οἱ ἐκλεκτοὶ,
\vs{35}καὶ ἀπολεῖται φυγὴ ἀπὸ τῶν ποιμένων, καὶ σωτηρία ἀπὸ τῶν κριῶν τῶν προβάτων.
\vs{36}Φωνὴ κραυγῆς τῶν ποιμένων, καὶ ἀλαλαγμὸς τῶν προβάτων καὶ τῶν κριῶν, ὅτι ὠλόθρευσε Κύριος τὰ βοσκήματα αὐτῶν·
\vs{37}Καὶ παύσεται τὰ κατάλοιπα τῆς εἰρήνης ἀπὸ προσώπου ὀργῆς θυμοῦ μου.
\vs{38}Ἐγκατέλιπεν ὥσπερ λέων κατάλυμα αὐτοῦ, ὅτι ἐγενήθη ἡ γῆ αὐτῶν εἰς ἄβατον ἀπὸ προσώπου τῆς μαχαίρας τῆς μεγάλης.

\ch{33}
ἘΝ ἈΡΧΗ ΒΑΣΙΛΕΩΣ ἸΩΑΚΕΙΜ ΥΙΟΥ ἸΩΣΙΑ, ἘΓΕΝΗΘΗ Ὁ ΛΟΓΟΣ ΟΥΤΟΣ ΠΑΡΑ ΚΥΡΙΟΥ·
\vs{2}Οὕτως εἶπε Κύριος, στῆθι ἐν αὐλῇ οἴκου Κυρίου, καὶ χρηματιεῖς ἅπασι τοῖς Ἰουδαίοις, καὶ πᾶσι τοῖς ἐρχομένοις προσκυνεῖν ἐν οἴκῳ Κυρίου, ἅπαντας τοὺς λόγους οὓς συνέταξά σοι χρηματίσαι αὐτοῖς, μὴ ἀφέλῃς ῥῆμα.
\vs{3}Ἴσως ἀκούσονται, καὶ ἀποστραφήσονται ἕκαστος ἀπὸ τῆς ὁδοῦ αὐτοῦ τῆς πονηρᾶς, καὶ παύσομαι ἀπὸ τῶν κακῶν ὧν ἐγὼ λογίζομαι τοῦ ποιῆσαι αὐτοῖς ἕνεκεν τῶν πονηρῶν ἐπιτηδευμάτων αὐτῶν.
\vs{4}Καὶ ἐρεῖς, οὕτως εἶπε Κύριος, ἐὰν μὴ ἀκούσητέ μου, τοῦ πορεύεσθαι ἐν τοῖς νομίμοις μου οἷς ἔδωκα κατὰ πρόσωπον ὑμῶν,
\vs{5}εἰσακούειν τῶν λόγων τῶν παίδων μου τῶν προφητῶν, οὓς ἐγὼ ἀποστέλλω πρὸς ὑμᾶς ὄρθρου, καὶ ἀπέστειλα, καὶ οὐκ ἠκούσατέ μου,
\vs{6}καὶ δώσω τὸν οἶκον τοῦτον ὥσπερ Σηλὼ, καὶ τὴν πόλιν δώσω εἰς κατάραν πᾶσι τοῖς ἔθνεσι πάσης τῆς γῆς.

\vs{7}Καὶ ἤκουσαν οἱ ἱερεῖς, καὶ οἱ ψευδοπροφῆται, καὶ πᾶς ὁ λαὸς τοῦ Ἱερεμίου λαλοῦντος τοὺς λόγους τούτους ἐν οἴκῳ Κυρίου.
\vs{8}Καὶ ἐγένετο Ἱερεμίου παυσαμένου λαλοῦντος πάντα ἃ συνέταξε Κύριος αὐτῷ λαλῆσαι παντὶ τῷ λαῷ, καὶ συνελάβοσαν αὐτὸν οἱ ἱερεῖς, καὶ οἱ ψευδοπροφῆται, καὶ πᾶς ὁ λαὸς, λέγων, θανάτῳ ἀποθανῇ,
\vs{9}ὅτι ἐπροφήτευσας τῷ ὀνόματι Κυρίου, λέγων, ὥσπερ Σηλὼ ἔσται ὁ οἶκος οὗτος, καὶ ἡ πόλις αὕτη ἐρημωθήσεται ἀπὸ κατοικούντων·

Καὶ ἐξεκκλησιάσθη πᾶς ὁ λαὸς ἐπὶ Ἱερεμίαν ἐν οἴκῳ Κυρίου.
\vs{10}Καὶ ἤκουσαν οἱ ἄρχοντες Ἰούδα τὸν λόγον τοῦτον, καὶ ἀνέβησαν ἐξ οἴκου τοῦ βασιλέως εἰς οἶκον Κυρίου, καὶ ἐκάθισαν ἐν προθύροις πύλης τῆς καινῆς.
\vs{11}Καὶ εἶπαν οἱ ἱερεῖς καὶ οἱ ψευδοπροφῆται πρὸς τοὺς ἄρχοντας, καὶ παντὶ τῷ λαῷ, κρίσις θανάτου τῷ ἀνθρώπῳ τούτῳ, ὅτι ἐπροφήτευσε κατὰ τῆς πόλεως ταύτης, καθὼς ἠκούσατε ἐν τοῖς ὠσὶν ὑμῶν.

\vs{12}Καὶ εἶπεν Ἱερεμίας πρὸς τοὺς ἄρχοντας, καὶ παντὶ τῷ λαῷ, λέγων, Κύριος ἀπέστειλέ με προφητεῦσαι ἐπὶ τὸν οἶκον τοῦτον, καὶ ἐπὶ τὴν πόλιν ταύτην, πάντας τοὺς λόγους οὓς ἠκούσατε.
\vs{13}Καὶ νῦν βελτίους ποιήσατε τὰς ὁδοὺς ὑμῶν, καὶ τὰ ἔργα ὑμῶν, καὶ ἀκούσατε τῆς φωνῆς Κυρίου, καὶ παύσεται Κύριος ἀπὸ τῶν κακῶν ὧν ἐλάλησεν ἐφʼ ὑμᾶς.
\vs{14}Καὶ ἰδοὺ ἐγὼ ἐν χερσὶν ὑμῶν, ποιήσατέ μοι ὡς συμφέρει, καὶ ὡς βέλτιον ὑμῖν.
\vs{15}Ἀλλʼ ἢ γνόντες γνώσεσθε, ὅτι εἰ ἀναιρεῖτέ με, αἷμα ἀθῶον δίδοτε ἐφʼ ὑμᾶς, καὶ ἐπὶ τὴν πόλιν ταύτην, καὶ ἐπὶ τοὺς κατοικοῦντας ἐν αὐτῇ· ὅτι ἐν ἀληθείᾳ ἀπέσταλκέ με Κύριος πρὸς ὑμᾶς λαλῆσαι εἰς τὰ ὦτα ὑμῶν πάντας τοὺς λόγους τούτους.

\vs{16}Καὶ εἶπον οἱ ἄρχοντες καὶ πᾶς ὁ λαὸς πρὸς τοὺς ἱερεῖς καὶ πρὸς τοὺς ψευδοπροφήτας, οὐκ ἔστι τῷ ἀνθρώπῳ τούτῳ κρίσις θανάτου, ὅτι ἐπὶ τῷ ὀνόματι Κυρίου τοῦ Θεοῦ ἡμῶν ἐλάλησε πρὸς ἡμᾶς.
\vs{17}Καὶ ἀνέστησαν ἄνδρες τῶν πρεσβυτέρων τῆς γῆς, καὶ εἶπαν πάσῃ τῇ συναγωγῇ τοῦ λαοῦ,
\vs{18}Μιχαίας ὁ Μωραθίτης ἦν ἐν ταῖς ἡμέραις Ἐζεκίου βασιλέως Ἰούδα, καὶ εἶπε παντὶ τῷ λαῷ Ἰούδα, οὕτως εἶπε Κύριος, Σιὼν ὡς ἀγρὸς ἀροτριαθήσεται, καὶ Ἱερουσαλὴμ εἰς ἄβατον ἔσται, καὶ τὸ ὄρος τοῦ οἴκου εἰς ἄλσος δρυμοῦ.
\vs{19}Μὴ ἀνελὼν ἀνεῖλεν αὐτὸν Ἐζεκίας καὶ πᾶς Ἰούδα; οὐχ ὅτι ἐφοβήθησαν τὸν Κύριον, καὶ ὅτι ἐδεήθησαν τοῦ προσώπου Κυρίου, καὶ ἐπαύσατο Κύριος ἀπὸ τῶν κακῶν ὧν ἐλάλησεν ἐπʼ αὐτούς; καὶ ἡμεῖς ἐποιήσαμεν κακὰ μεγάλα ἐπὶ ψυχὰς ἡμῶν.

\vs{20}Καὶ ἄνθρωπος ἦν προφητεύων τῷ ὀνόματι Κυρίου, Οὐρίας υἱὸς Σαμαίου ἐκ Καριαθιαρὶμ, καὶ ἐπροφήτευσε περὶ τῆς γῆς ταύτης κατὰ πάντας τοὺς λόγους Ἱερεμίου.
\vs{21}Καὶ ἤκουσεν ὁ βασιλεὺς Ἰωακεὶμ καὶ πάντες οἱ ἄρχοντες πάντας τοὺς λόγους αὐτοῦ, καὶ ἐζήτουν ἀποκτεῖναι αὐτόν· καὶ ἤκουσεν Οὐρίας, καὶ εἰσῆλθεν εἰς Αἴγυπτον.
\vs{22}Καὶ ἐξαπέστειλεν ὁ βασιλεὺς ἄνδρας εἰς Αἴγυπτον,
\vs{23}καὶ ἐξηγάγοσαν αὐτὸν ἐκεῖθεν, καὶ εἰσηγάγοσαν αὐτὸν πρὸς τὸν βασιλέα, καὶ ἐπάταξεν αὐτὸν ἐν μαχαίρᾳ, καὶ ἔῤῥιψεν αὐτὸν εἰς τὸ μνῆμα υἱῶν λαοῦ αὐτοῦ.
\vs{24}Πλὴν χεὶρ Ἀχεικὰμ υἱοῦ Σαφὰν ἦν μετὰ Ἱερεμίου, τοῦ μὴ παραδοῦναι αὐτὸν εἰς χεῖρας τοῦ λαοῦ, μὴ ἀνελεῖν αὐτόν.

\ch{34}
\vs{2}Οὕτως εἶπε Κύριος, ποίησον σεαυτῷ δεσμοὺς καὶ κλοιοὺς, καὶ περίθου περὶ τὸν τράχηλόν σου.
\vs{3}Καὶ ἀποστελεῖς αὐτοὺς πρὸς βασιλέα Ἰδουμαίας, καὶ πρὸς βασιλέα Μωὰβ, καὶ πρὸς βασιλέα υἱῶν Ἀμμὼν, καὶ πρὸς τὸν βασιλέα Τύρου, καὶ πρὸς βασιλέα Σιδῶνος, ἐν χερσὶν ἀγγέλων αὐτῶν τῶν ἐρχομένων εἰς ἀπάντησιν αὐτῶν εἰς Ἱερουσαλὴμ πρὸς Σεδεκίαν βασιλέα Ἰούδα.
\vs{4}Καὶ συντάξεις αὐτοῖς πρὸς τοὺς κυρίους αὐτῶν εἰπεῖν, οὕτως εἶπε Κύριος ὁ Θεὸς Ἰσραὴλ, οὕτως ἐρεῖτε πρὸς τοὺς κυρίους ὑμῶν,
\vs{5}ὅτι ἐγὼ ἐποίησα τὴν γῆν ἐν τῇ ἰσχύϊ μου τῇ μεγάλῃ, καὶ ἐν τῷ ἐπιχείρῳ μου τῷ ὑψηλῷ, καὶ δώσω αὐτὴν ᾧ ἐὰν δόξῃ ἐν ὀφθαλμοῖς μου·
\vs{6}Ἔδωκα τὴν γῆν τῷ Ναβουχοδονόσορ βασιλεῖ Βαβυλῶνος δουλεύειν αὐτῷ, καὶ τὰ θηρία τοῦ ἀγροῦ ἐργάζεσθαι αὐτῷ.
\vs{8}Καὶ τὸ ἔθνος καὶ ἡ βασιλεία, ὅσοι ἐὰν μὴ ἐμβάλωσι τὸν τράχηλον αὐτῶν ὑπὸ τὸν ζυγὸν βασιλέως Βαβυλῶνος, ἐν μαχαίρᾳ καὶ ἐν λιμῷ ἐπισκέψομαι αὐτοὺς, εἶπε Κύριος, ἕως ἐκλίπωσιν ἐν χειρὶ αὐτοῦ.

\vs{9}Καὶ ὑμεῖς μὴ ἀκούετε τῶν ψευδοπροφητῶν ὑμῶν, καὶ τῶν μαντευομένων ὑμῖν, καὶ τῶν ἐνυπνιαζομένων ὑμῖν, καὶ τῶν οἰωνισμάτων ὑμῶν, καὶ τῶν φαρμακῶν ὑμῶν, τῶν λεγόντων, οὐ μὴ ἐργάσησθε τῷ βασιλεῖ Βαβυλῶνος·
\vs{10}Ὅτι ψευδῆ αὐτοὶ προφητεύουσιν ὑμῖν, πρὸς τὸ μακρῦναι ὑμᾶς ἀπὸ τῆς γῆς ὑμῶν.
\vs{11}Καὶ τὸ ἔθνος ὃ ἐὰν εἰσαγάγῃ τὸν τράχηλον αὐτοῦ ὑπὸ τὸν ζυγὸν βασιλέως Βαβυλῶνος, καὶ ἐργάσηται αὐτῷ, καὶ καταλείψω αὐτὸν ἐπὶ τῆς γῆς αὐτοῦ, καὶ ἐργᾶται αὐτῷ, καὶ ἐνοικήσει ἐν αὐτῇ.

\vs{12}Καὶ πρὸς Σεδεκίαν βασιλέα Ἰούδα ἐλάλησα κατὰ πάντας τοὺς λόγους τούτους, λέγων, εἰσαγάγετε τὸν τράχηλον ὑμῶν,
\vs{14}καὶ ἐργάσασθε τῷ βασιλεῖ Βαβυλῶνος, ὅτι ἄδικα αὐτοὶ προφητεύουσιν ὑμῖν,
\vs{15}ὅτι οὐκ ἀπέστειλα αὐτοὺς, φησὶ Κύριος, καὶ προφητεύουσι τῷ ὀνόματί μου ἐπ ἀδίκῳ, πρὸς τὸ ἀπολέσαι ὑμᾶς, καὶ ἀπολεῖσθε ὑμεῖς, καὶ οἱ προφῆται ὑμῶν, οἱ προφητεύοντες ὑμῖν ἐπʼ ἀδίκῳ ψευδῆ.

\vs{16}Ὑμῖν, καὶ παντὶ τῷ λαῷ τούτῳ, καὶ τοῖς ἱερεῦσιν ἐλάλησα, λέγων, οὕτως εἶπε Κύριος, μὴ ἀκούετε τῶν λόγων τῶν προφητῶν, τῶν προφητευόντων ὑμῖν, λεγόντων, ἰδοὺ σκεύη οἴκου Κυρίου ἐπιστρέψει ἐκ Βαβυλῶνος· ὅτι ἄδικα αὐτοὶ προφητεύουσιν ὑμῖν. Οὐκ ἀπέστειλα αὐτούς.
\vs{18}Εἰ προφῆταί εἰσι, καὶ εἰ ἔστι λόγος Κυρίου ἐν αὐτοῖς, ἀπαντησάτωσάν μοι, ὅτι οὕτως εἶπε Κύριος.

\vs{19}Καὶ τῶν ἐπιλοίπων σκευῶν,
\vs{20}ὧν οὐκ ἔλαβε βασιλεὺς Βαβυλῶνος, ὅτε ἀπῴκισε τὸν Ἰεχονίαν ἐξ Ἱερουσαλὴμ,
\vs{22}εἰς Βαβυλῶνα εἰσελεύσεται, λέγει Κύριος.

\ch{35}
Καὶ ἐγένετο ἐν τῷ τετάρτῳ ἔτει Σεδεκία βασιλέως Ἰούδα ἐν μηνὶ τῷ πέμπτῳ, εἶπέ μοι Ἀνανίας υἱὸς Ἀζὼρ ὁ ψευδοπροφήτης ἀπὸ Γαβαὼν ἐν οἴκῳ Κυρίου, κατʼ ὀφθαλμοὺς τῶν ἱερέων, καὶ παντὸς τοῦ λαοῦ, λέγων,
\vs{2}Οὕτως εἶπε Κύριος, συνέτριψα τὸν ζυγὸν τοῦ βασιλέως Βαβυλῶνος.
\vs{3}Ἔτι δύο ἔτη ἡμερῶν, καὶ ἐγὼ ἀποστρέψω εἰς τὸν τόπον τοῦτον τὰ σκεύη οἴκου Κυρίου,
\vs{4}καὶ Ἰεχονίαν, καὶ τὴν ἀποικίαν Ἰούδα, ὅτι συντρίψω τὸν ζυγὸν βασιλέως Βαβυλῶνος.

\vs{5}Καὶ εἶπεν Ἱερεμίας πρὸς Ἀνανίαν κατʼ ὀφθαλμοὺς παντὸς τοῦ λαοῦ καὶ κατʼ ὀφθαλμοὺς τῶν ἱερέων τῶν ἑστηκότων ἐν οἴκῳ Κυρίου,
\vs{6}καὶ εἶπεν Ἱερεμίας, ἀληθῶς οὕτω ποιήσαι Κύριος, στήσαι τὸν λόγον σου ὃν σὺ προφητεύεις, τοῦ ἐπιστρέψαι τὰ σκεύη οἴκου Κυρίου καὶ πᾶσαν τὴν ἀποικίαν ἐκ Βαβυλῶνος εἰς τὸν τόπον τοῦτον.
\vs{7}Πλὴν ἀκούσατε τὸν λόγον Κυρίου, ὃν ἐγὼ λέγω εἰς τὰ ὦτα ὑμῶν καὶ εἰς τὰ ὦτα παντὸς τοῦ λαοῦ.
\vs{8}Οἱ προφῆται οἱ γεγονότες πρότεροί μου καὶ πρότεροι ὑμῶν ἀπὸ τοῦ αἰῶνος, καὶ ἐπροφήτευσαν ἐπὶ γῆς πολλῆς, καὶ ἐπὶ βασιλείας μεγάλας εἰς πόλεμον.
\vs{9}Ὁ προφήτης ὁ προφητεύσας εἰς εἰρήνην, ἐλθόντος τοῦ λόγου, γνώσονται τὸν προφήτην ὃν ἀπέστειλεν αὐτοῖς Κύριος ἐν πίστει.

\vs{10}Καὶ ἔλαβεν Ἀνανίας ἐν ὀφθαλμοῖς παντὸς τοῦ λαοῦ τοὺς κλοιοὺς ἀπὸ τοῦ τραχήλου Ἱερεμίου, καὶ συνέτριψεν αὐτούς.
\vs{11}Καὶ εἶπεν Ἀνανίας κατʼ ὀφθαλμοὺς παντὸς τοῦ λαοῦ, λέγων, οὕτως εἶπε Κύριος, οὕτως συντρίψω τὸν ζυγὸν βασιλέως Βαβυλῶνος ἀπὸ τραχήλων πάντων τῶν ἐθνῶν· καὶ ᾤχετο Ἱερεμίας εἰς τὴν ὁδὸν αὐτοῦ.

\vs{12}Καὶ ἐγένετο λόγος Κυρίου πρὸς Ἱερεμίαν, μετὰ τὸ συντρίψαι Ἀνανίαν τοὺς κλοιοὺς ἀπὸ τοῦ τραχήλου αὐτοῦ, λέγων,
\vs{13}Βάδιζε, καὶ εἰπὸν πρὸς Ἀνανίαν, λέγων, οὕτως εἶπε Κύριος, κλοιοὺς ξυλίνους συνέτριψας, καὶ ποιήσω ἀντʼ αὐτῶν κλοιοὺς σιδηροῦς·
\vs{14}Ὅτι οὕτως εἶπε Κύριος, ζυγὸν σιδηροῦν ἔθηκα ἐπὶ τὸν τράχηλον πάντων τῶν ἐθνῶν, ἐργάζεσθαι τῷ βασιλεῖ Βαβυλῶνος.
\vs{15}Καὶ εἶπεν Ἱερεμίας τῷ Ἀνανίᾳ, οὐκ ἀπέσταλκέ σε Κύριος, καὶ πεποιθέναι ἐποίησας τὸν λαὸν τοῦτον ἐπʼ ἀδίκῳ.
\vs{16}Διατοῦτο οὕτως εἶπε Κύριος, ἰδοὺ ἐγὼ ἐξαποστέλλω σε ἀπὸ προσώπου τῆς γῆς, τούτῳ τῷ ἐνιαυτῷ ἀποθανῇ.
\vs{17}Καὶ ἀπέθανεν ἐν τῷ μηνὶ τῷ ἑβδόμῳ.

\ch{36}
Καὶ οὗτοι οἱ λόγοι τῆς βίβλου οὓς ἀπέστειλεν Ἱερεμίας ἐξ Ἱερουσαλὴμ πρὸς τοὺς πρεσβυτέρους τῆς ἀποικίας, καὶ πρὸς τοὺς ἱερεῖς, καὶ πρὸς τοὺς ψευδοπροφήτας, ἐπιστολὴν εἰς Βαβυλῶνα τῇ ἀποικίᾳ, καὶ πρὸς ἅπαντα τὸν λαὸν,
\vs{2}ὕστερον ἐξελθόντος Ἰεχονίου τοῦ βασιλέως, καὶ τῆς βασιλίσσης, καὶ τῶν εὐνούχων, καὶ παντὸς ἐλευθέρου, καὶ δεσμώτου, καὶ τεχνίτου ἐξ Ἱερουσαλὴμ,
\vs{3}ἐν χειρὶ Ἐλεασὰν υἱοῦ Σαφὰν, καὶ Γαμαρίου υἱοῦ Χελκίου, ὃν ἀπέστειλε Σεδεκίας βασιλεὺς Ἰούδα, πρὸς βασιλέα Βαβυλῶνος εἰς Βαβυλῶνα, λέγων,

\vs{4}Οὕτως εἶπε Κύριος ὁ Θεὸς Ἰσραὴλ ἐπὶ τὴν ἀποικίαν ἣν ἀπῴκισα ἀπὸ Ἱερουσαλὴμ,
\vs{5}οἰκοδομήσατε οἴκους, καὶ κατοικήσατε, καὶ φυτεύσατε παραδείσους, καὶ φάγετε τοὺς καρποὺς αὐτῶν,
\vs{6}καὶ λάβετε γυναῖκας, καὶ τεκνοποιήσατε υἱοὺς καὶ θυγατέρας, καὶ λάβετε τοῖς υἱοῖς ὑμῶν γυναῖκας, καὶ τὰς θυγατέρας ὑμῶν δότε ἀνδράσι, καὶ πληθύνεσθε, καὶ μὴ σμικρυνθῆτε·
\vs{7}Καὶ ζητήσατε εἰς εἰρήνην τῆς γῆς, εἰς ἣν ἀπῴκισα ὑμᾶς ἐκεῖ· καὶ προσεύξεσθε περὶ αὐτῶν πρὸς Κύριον, ὅτι ἐν εἰρήνῃ αὐτῆς εἰρήνη ὑμῖν.

\vs{8}Ὅτι οὕτως εἶπε Κύριος, μὴ ἀναπειθέτωσαν ὑμᾶς οἱ ψευδοπροφῆται οἱ ἐν ὑμῖν, καὶ μὴ ἀναπειθέτωσαν ὑμᾶς οἱ μάντεις ὑμῶν, καὶ μὴ ἀκούετε εἰς τὰ ἐνύπνια ὑμῶν, ἃ ὑμεῖς ἐνυπνιάζεσθε,
\vs{9}ὅτι ἄδικα αὐτοὶ προφητεύουσιν ὑμῖν ἐπὶ τῷ ὀνόματί μου, καὶ οὐκ ἀπέστειλα αὐτούς.
\vs{10}Ὅτι οὕτως εἶπε Κύριος, ὅταν μέλλῃ πληροῦσθαι Βαβυλῶνι ἑβδομήκοντα ἔτη, ἐπισκέψομαι ὑμᾶς, καὶ ἐπιστήσω τοὺς λόγους μου ἐφʼ ὑμᾶς, τοῦ ἀποστρέψαι τὸν λαὸν ὑμῶν εἰς τὸν τόπον τοῦτον.
\vs{11}Καὶ λογιοῦμαι ἐφʼ ὑμᾶς λογισμὸν εἰρήνης, καὶ οὐ κακὰ, τοῦ δοῦναι ὑμῖν ταῦτα.
\vs{12}Καὶ προσεύξασθε πρὸς μὲ, καὶ εἰσακούσομαι ὑμῶν.
\vs{13}Καὶ ἐκζητήσατέ με, καὶ εὑρήσετέ με· ὅτι ζητήσετέ με ἐν ὅλῃ καρδίᾳ ὑμῶν,
\vs{14}καὶ ἐπιφανοῦμαι ὑμῖν·
\vs{15}Ὅτι εἴπατε, κατέστησεν ἡμῖν Κύριος προφήτας ἐν Βαβυλῶνι·

\vs{21}Οὕτως εἶπε Κύριος ἐπὶ Ἀχιὰβ, καὶ ἐπὶ Σεδεκίαν, ἰδοὺ ἐγὼ δίδωμι αὐτοὺς εἰς χεῖρας βασιλέως Βαβυλῶνος, καὶ πατάξει αὐτοὺς κατʼ ὀφθαλμοὺς ὑμῶν.
\vs{22}Καὶ λήψονται ἀπʼ αὐτῶν κατάραν ἐν πάσῃ τῇ ἀποικίᾳ Ἰούδα ἐν Βαβυλῶνι, λέγοντες, ποιήσαι σε Κύριος, ὡς Σεδεκίαν ἐποίησε, καὶ ὡς Ἀχιὰβ, οὓς ἀπετηγάνισε βασιλεὺς Βαβυλῶνος ἐν πυρὶ,
\vs{23}διʼ ἣν ἐποίησαν ἀνομίαν ἐν Ἰσραὴλ, καὶ ἐμοιχῶντο τὰς γυναῖκας τῶν πολιτῶν αὐτῶν, καὶ λόγον ἐχρημάτισαν ἐν τῷ ὀνόματί μου, ὃν οὐ συνέταξα αὐτοῖς· καὶ ἐγὼ μάρτυς, φησὶ Κύριος.

\vs{24}Καὶ πρὸς Σαμαίαν τὸν Αἰλαμίτην ἐρεῖς,
\vs{25}οὐκ ἀπέστειλά σε τῷ ὀνόματί μου· καὶ πρὸς Σοφονίαν υἱὸν Μαασαίου τὸν ἱερέα εἰπὲ,
\vs{26}Κύριος ἔδωκέ σε ἱερέα ἀντὶ Ἰωδαὲ τοῦ ἱερέως, γενέσθαι ἐπιστάτην ἐν τῷ οἴκῳ Κυρίου παντὶ ἀνθρώπῳ προφητεύοντι, καὶ παντὶ ἀνθρώπῳ μαινομένῳ, καὶ δώσεις αὐτὸν εἰς τὸ ἀπόκλεισμα, καὶ εἰς τὸν καταράκτην.
\vs{27}Καὶ νῦν διατί συνελοιδορήσατε Ἱερεμίαν τὸν ἐξ Ἀναθὼθ, τὸν προφητεύσαντα ὑμῖν;
\vs{28}Οὐ διατοῦτο ἀπέστειλεν; ὅτι διὰ τοῦ μηνὸς τούτου ἀπέστειλε πρὸς ὑμᾶς εἰς Βαβυλῶνα, λέγων, μακράν ἐστιν, οἰκοδομήσατε οἰκίας, καὶ κατοικήσατε, καὶ φυτεύσατε κήπους, καὶ φάγεσθε τὸν καρπὸν αὐτῶν.
\vs{29}Καὶ ἀνέγνω Σοφονίας τὸ βιβλίον εἰς τὰ ὦτα Ἱερεμίου.

\vs{30}Καὶ ἐγένετο λόγος Κυρίου πρὸς Ἱερεμίαν, λέγων,
\vs{31}ἀπόστειλον πρὸς τὴν ἀποικίαν, λέγων, οὕτως εἶπε Κύριος ἐπὶ Σαμαίαν τὸν Αἰλαμίτην, ἐπειδὴ ἐπροφήτευσεν ὑμῖν Σαμαίας, καὶ ἐγὼ οὐκ ἀπέστειλα αὐτὸν, καὶ πεποιθέναι ἐποίησεν ὑμᾶς ἐπʼ ἀδίκοις,
\vs{32}διατοῦτο οὕτως εἶπε Κύριος, ἰδοὺ ἐγὼ ἐπισκέψομαι ἐπὶ Σαμαίαν, καὶ ἐπὶ τὸ γένος αὐτοῦ, καὶ οὐκ ἔσται αὐτῶν ἄνθρωπος ἐν μέσῳ ὑμῶν, τοῦ ἰδεῖν τὰ ἀγαθὰ, ἃ ἐγὼ ποιήσω ὑμῖν, οὐκ ὄψονται.

\ch{37}
Ὁ ΛΟΓΟΣ Ὁ ΓΕΝΟΜΕΝΟΣ ΠΡΟΣ ἹΕΡΕΜΙΑΝ ΠΑΡΑ ΚΥΡΙΟΥ, ΕΙΠΕΙΝ,
\vs{2}Οὕτως εἶπε Κύριος ὁ Θεὸς Ἰσραὴλ, λέγων,

Γράψον πάντας τοὺς λόγους οὓς ἐχρημάτισα πρὸς σὲ ἐπὶ βιβλίου.
\vs{3}Ὅτι ἰδοὺ ἡμέραι ἔρχονται, φησὶ Κύριος, καὶ ἀποστρέψω τὴν ἀποικίαν λαοῦ μου Ἰσραὴλ καὶ Ἰούδα, εἶπε Κύριος, καὶ ἀποστρέψω αὐτοὺς εἰς τὴν γῆν ἣν ἔδωκα τοῖς πατράσιν αὐτῶν, καὶ κυριεύσουσιν αὐτῆς.

\vs{4}ΚΑΙ ΟΥΤΟΙ ΟΙ ΛΟΓΟΙ ΟΥΣ ἘΛΑΛΗΣΕ ΚΥΡΙΟΣ ἘΠΙ ἸΣΡΑΗΛ ΚΑΙ ἸΟΥΔΑ.

\vs{5}Οὕτως εἶπε Κύριος, φωνὴν φόβου ἀκούσεσθε· φόβος, καὶ οὐκ ἔστιν εἰρήνη.
\vs{6}Ἐρωτήσατε, καὶ ἴδετε, εἰ ἔτεκεν ἄρσεν; καὶ περὶ φόβου, ἐν ᾧ καθέξουσιν ὀσφὺν καὶ σωτηρίαν; διότι ἑώρακα πάντα ἄνθρωπον, καὶ αἱ χεῖρες αὐτοῦ ἐπὶ τῆς ὀσφύος αὐτοῦ· ἐστράφησαν πρόσωπα εἰς ἴκτερον.
\vs{7}Ὅτι ἐγενήθη μεγάλη ἡ ἡμέρα ἐκείνη, καὶ οὐκ ἔστι τοιαύτη· καὶ χρόνος στενός ἐστι τῷ Ἰακὼβ,
\vs{8}καὶ ἀπὸ τούτου σωθήσεται. Ἐν τῇ ἡμέρᾳ ἐκείνῃ, εἶπε Κύριος, συντρίψω τὸν ζυγὸν ἀπὸ τοῦ τραχήλου αὐτῶν, καὶ τοὺς δεσμοὺς αὐτῶν διαῤῥήξω· καὶ οὐκ ἐργῶνται αὐτοὶ ἔτι ἀλλοτρίοις,
\vs{9}καὶ ἐργῶνται τῷ Κυρίῳ Θεῷ αὐτῶν· καὶ τὸν Δαυὶδ βασιλέα αὐτῶν ἀναστήσω αὐτοῖς.

\vs{12}Οὕτως εἶπε Κύριος, ἀνέστησα σύντριμμα, ἀλγηρὰ ἡ πληγή σου,
\vs{13}οὐκ ἔστι κρίνων κρίσιν σου, εἰς ἀλγηρὸν ἰατρεύθης, ὠφέλειά σοι οὐκ ἔστι.
\vs{14}Πάντες οἱ φίλοι σου ἐπελάθοντό σου, οὐ μὴ ἐπερωτήσωσιν· ὅτι πληγὴν ἐχθροῦ ἔπαισά σε, παιδείαν στερεάν· ἐπὶ πᾶσαν ἀδικίαν σου ἐπλήθυναν αἱ ἁμαρτίαι σου.
\vs{16}Διατοῦτο πάντες οἱ ἔσθοντές σε βρωθήσονται, καὶ πάντες οἱ ἐχθροί σου κρέας αὐτῶν πᾶν ἔδονται. Ἐπὶ πλῆθος ἀδικιῶν σου ἐπληθύνθησαν αἱ ἁμαρτίαι σου· ἐποίησαν ταῦτά σοι· Καὶ ἔσονται οἱ διαφοροῦντές σε εἰς διαφόρημα, καὶ πάντας τοὺς προνομεύσαντάς σε δώσω εἰς προνομήν.
\vs{17}Ὅτι ἀνάξω τὸ ἴαμά σου, ἀπὸ πληγῆς ὀδυνηρᾶς ἰατρεύσω σε, φησὶ Κύριος· ὅτι Ἐσπαρμένη ἐκλήθης, θήρευμα ὑμῶν ἐστιν, ὅτι ζητῶν οὐκ ἔστιν αὐτήν.

\vs{18}Οὕτως εἶπε Κύριος, ἰδοὺ ἐγὼ ἀποστρέψω τὴν ἀποικιαν Ἰακὼβ, καὶ τὴν αἰχμαλωσίαν αὐτοῦ ἐλεήσω· καὶ οἰκοδομηθήσεται πόλις ἐπὶ τὸ ὕψος αὐτῆς, καὶ ὁ λαὸς κατὰ τὸ κρίμα αὐτοῦ καθεδεῖται,
\vs{19}καὶ ἐξελεύσονται ἀπʼ αὐτῶν ᾄδοντες, φωνὴ παιζόντων· καὶ πλεονάσω αὐτοὺς, καὶ οὐ μὴ ἐλαττωθῶσι.
\vs{20}Καὶ εἰσελεύσονται οἱ υἱοὶ αὐτῶν ὡς τὸ πρότερον, καὶ τὰ μαρτύρια αὐτῶν κατὰ πρόσωπόν μου ὀρθωθήσεται· καὶ ἐπισκέψομαι τοὺς θλίβοντας αὐτοὺς.
\vs{21}Καὶ ἔσονται ἰσχυρότεροι αὐτοῦ ἐπʼ αὐτοὺς, καὶ ὁ ἄρχων αὐτοῦ ἐξ αὐτοῦ ἐξελεύσεται· καὶ συνάξω αὐτοὺς, καὶ ἀποστρέψουσι πρὸς μέ· ὅτι τίς ἐστιν οὗτος ὃς ἔδωκε τὴν καρδίαν αὐτοῦ, ἀποστρέψαι πρὸς μέ; φησὶ Κύριος·

\vs{23}Ὅτι ὀργὴ Κυρίου ἐξῆλθε θυμώδης, ἐξῆλθεν ὀργὴ στρεφομένη, ἐπʼ ἀσεβεῖς ἥξει.
\vs{24}Οὐ μὴ ἀποστραφῇ ὀργὴ θυμοῦ Κυρίου, ἕως ποιήσει, καὶ ἕως καταστήσῃ ἐγχείρημα καρδίας αὐτοῦ· ἐπʼ ἐσχάτων τῶν ἡμερῶν γνώσεσθε αὐτά.

\ch{38}
Ἐν τῷ χρόνῳ ἐκείνῳ, εἶπε Κύριος, ἔσομαι εἰς Θεὸν τῷ γένει Ἰσραὴλ, καὶ αὐτοὶ ἔσονταί μοι εἰς λαόν.
\vs{2}Οὕτως εἶπε Κύριος, εὗρον θερμὸν ἐν ἐρήμῳ μετὰ ὀλωλότων ἐν μαχαίρᾳ· βαδίσατε, καὶ μὴ ὀλέσητε τὸν Ἰσραὴλ.
\vs{3}Κύριος πόῤῥωθεν ὤφθη αὐτῷ· ἀγάπησιν αἰώνιον ἠγάπησά σε· διατοῦτο εἵλκυσά σε εἰς οἰκτείρημα.
\vs{4}Ὅτι οἰκοδομήσω σε, καὶ οἰκοδομηθήσῃ παρθένος Ἰσραήλ· ἔπι λήψῃ τύμπανόν σου, καὶ ἐξελεύσῃ μετὰ συναγωγῆς παιζόντων.
\vs{5}Ὅτι ἐφυτεύσατε ἀμπελῶνας ἐν ὄρεσι Σαμαρείας, φυτεύσατε καὶ αἰνέσατε,
\vs{6}ὅτι ἐστὶν ἡμέρα κλήσεως ἀπολογουμένων ἐν ὄρεσιν Ἐφραὶμ, ἀνάστητε, καὶ ἀνάβητε εἰς Σιὼν πρὸς Κύριον τὸν Θεὸν ὑμῶν.

\vs{7}Ὅτι οὕτως εἶπε Κύριος τῷ Ἰακὼβ, εὐφράνθητε, καὶ χρεμετίσατε ἐπὶ κεφαλὴν ἐθνῶν· ἀκουστὰ ποιήσατε, καὶ αἰνέσατε· εἴπατε, ἔσωσε Κύριος τὸν λαὸν αὐτοῦ, τὸ κατάλοιπον τοῦ Ἰσραήλ.
\vs{8}Ἰδοὺ ἐγὼ ἄγω αὐτοὺς ἀπὸ Βοῤῥᾶ, καὶ συνάξω αὐτοὺς ἀπʼ ἐσχάτου τῆς γῆς ἐν ἑορτῇ φασέκ· καὶ τεκνοποιήσει ὄχλον πολὺν, καὶ ἀποστρέψουσιν ὧδε.
\vs{9}Ἐν κλαυθμῷ ἐξῆλθον, καὶ ἐν παρακλήσει ἀνάξω αὐτοὺς, αὐλίζων ἐπὶ διώρυγας ὑδάτων ἐν ὁδῷ ὀρθῇ, καὶ οὐ μὴ πλανηθῶσιν ἐν αὐτῇ· ὅτι ἐγενόμην τῷ Ἰσραὴλ εἰς πατέρα, καὶ Ἐφραὶμ πρωτότοκός μου ἐστίν.

\vs{10}Ἀκούσατε λόγους Κυρίου ἔθνη, καὶ ἀναγγείλατε εἰς νήσους τὰς μακρόθεν· εἴπατε, ὁ λικμήσας τὸν Ἰσραὴλ καὶ συνάξει αὐτὸν, καὶ φυλάξει αὐτὸν, ὡς ὁ βόσκων ποίμνιον αὐτοῦ.
\vs{11}Ὅτι ἐλυτρώσατο Κύριος τὸν Ἰακὼβ, ἐξείλατο αὐτὸν ἐκ χειρὸς στερεωτέρων αὐτοῦ.
\vs{12}Καὶ ἥξουσι, καὶ εὐφρανθήσονται ἐν τῷ ὄρει Σιὼν, καὶ ἥξουσιν ἐπʼ ἀγαθὰ Κυρίου, ἐπὶ γῆν σίτου καὶ οἴνου, καὶ καρπῶν, καὶ κτηνῶν, καὶ προβάτων· καὶ ἔσται ἡ ψυχὴ αὐτῶν ὥσπερ ξύλον ἔγκαρπον, καὶ οὐ πεινάσουσιν ἔτι.
\vs{13}Τότε χαρήσονται παρθένοι ἐν συναγωγῇ νεανίσκων, καὶ πρεσβύται χαρήσονται, καὶ στρέψω τὸ πένθος αὐτῶν εἰς χαρμονὴν, καὶ ποιήσω αὐτοὺς εὐφραινομένους.
\vs{14}Μεγαλυνῶ καὶ μεθύσω τὴν ψυχὴν τῶν ἱερεων υἱῶν Λευὶ, καὶ ὁ λαός μου τῶν ἀγαθῶν μου ἐμπλησθήσεται· οὕτως εἶπε Κύριος.

\vs{15}Φωνὴ ἐν Ῥαμᾷ ἠκούσθη θρήνου, καὶ κλαυθμοῦ, καὶ ὀδυρμοῦ· Ῥαχὴλ ἀποκλαιομένη οὐκ ἤθελε παύσασθαι ἐπὶ τοῖς υἱοῖς αὐτῆς, ὅτι οὐκ εἰσίν.

\vs{16}Οὕτως εἶπε Κύριος, διαλειπέτω ἡ φωνή σου ἀπὸ κλαυθμοῦ, καὶ οἱ ὀφθαλμοί σου ἀπὸ δακρύων σου, ὅτι ἔστι μισθὸς τοῖς σοῖς ἔργοις, καὶ ἐπιστρέψουσιν ἐκ γῆς ἐχθρῶν,
\vs{17}μόνιμον τοῖς σοῖς τέκνοις.

\vs{18}Ἀκοὴν ἤκουσα Ἐφραὶμ ὀδυρομένου, ἐπαίδευσάς με, καὶ ἐπαιδεύθην· ἐγὼ ὥσπερ μόσχος οὐκ ἐδιδάχθην· ἐπίστρεψόν με, καὶ ἐπιστρέψω, ὅτι σὺ Κύριος ὁ Θεός μου.
\vs{19}Ὅτι ὕστερον αἰχμαλωσίας μου μετενόησα, καὶ ὕστερον τοῦ γνῶναί με, ἐστέναξα ἐφʼ ἡμέρας αἰσχύνης, καὶ ὑπέδειξά σοι, ὅτι ἔλαβον ὀνειδισμὸν ἐκ νεότητός μου.
\vs{20}Υἱὸς ἀγαπητὸς Ἐφραὶμ, ἐμοὶ παιδίον ἐντρυφῶν, ὅτι ἀνθʼ ὧν οἱ λόγοι μου ἐν αὐτῷ, μνείᾳ μνησθήσομαι αὐτοῦ· διατοῦτο ἔσπευσα ἐπʼ αὐτῷ, ἐλεῶν ἐλεήσω αὐτὸν, φησὶ Κύριος.

\vs{21}Στῆσον σεαυτὴν Σιὼν, ποίησον τιμωρίαν, δὸς καρδίαν σου εἰς τοὺς ὤμους· ὁδὸν ᾗ ἐπορεύθης, ἀποστράφηθι παρθένος Ἰσραὴλ, ἀποστράφητι εἰς τὰς πόλεις σου πενθοῦσα.
\vs{22}Ἕως πότε ἀποστρέψεις θυγάτηρ ἠτιμωμένη; ὅτι ἔκτισε Κύριος σωτηρίαν εἰς καταφύτευσιν καινὴν, ἐν σωτηρίᾳ περιελεύσονται ἄνθρωποι.

\vs{23}Ὅτι οὕτως εἶπε Κύριος, ἔτι ἐροῦσι τὸν λόγον τοῦτον ἐν γῇ Ἰούδα, καὶ ἐν πόλεσιν αὐτοῦ, ὅταν ἀποστρέψω τὴν αἰχμαλωσίαν αὐτοῦ, εὐλογημένος Κύριος ἐπὶ δίκαιον ὄρος τὸ ἅγιον αὐτοῦ·
\vs{24}Καὶ ἐνοικοῦντες ἐν ταῖς πόλεσιν Ἰούδα, καὶ ἐν πάσῃ τῇ γῇ αὐτοῦ, ἅμα γεωργῷ, καὶ ἀρθήσεται ἐν ποιμνίῳ.
\vs{25}Ὅτι ἐμέθυσα πᾶσαν ψυχὴν διψῶσαν, καὶ πᾶσαν ψυχὴν πεινῶσαν ἐνέπλησα.
\vs{26}Διατοῦτο ἐξηγέρθην, καὶ εἶδον, καὶ ὁ ὕπνος μου ἡδύς μοι ἐγενήθη.

\vs{27}Διατοῦτο ἰδοὺ ἡμέραι ἔρχονται, φησὶ Κύριος, καὶ σπερῶ τὸν Ἰσραὴλ καὶ τὸν Ἰούδαν, σπέρμα ἀνθρώπου καὶ σπέρμα κτήνους.
\vs{28}Καὶ ἔσται ὥσπερ ἐγρηγόρουν ἐπʼ αὐτοὺς καθαιρεῖν καὶ κακοῦν, οὕτως γρηγορήσω ἐπʼ αὐτοὺς τοῦ οἰκοδομεῖν καὶ καταφυτεύειν, φησὶ Κύριος.
\vs{29}Ἐν ταῖς ἡμέραις ἐκείναις οὐ μὴ εἴπωσιν, οἱ πατέρες ἔφαγον ὄμφακα, καὶ οἱ ὀδόντες τῶν τέκνων ᾑμωδίασαν.
\vs{30}Ἀλλʼ ἢ ἕκαστος ἐν τῇ ἑαυτοῦ ἁμαρτίᾳ ἀποθανεῖται, καὶ τοῦ φαγόντος τὸν ὄμφακα αἱμωδιάσουσιν οἱ ὀδόντες αὐτοῦ.

\vs{31}Ἰδοὺ ἡμέραι ἔρχονται, φησὶ Κύριος, καὶ διαθήσομαι τῷ οἴκῳ Ἰσραὴλ καὶ τῷ οἴκῳ Ἰούδα διαθήκην καινὴν,
\vs{32}οὐ κατὰ τὴν διαθήκην ἣν διεθέμην τοῖς πατράσιν αὐτῶν, ἐν ἡμέρᾳ ἐπιλαβομένου μου τῆς χειρὸς αὐτῶν, ἐξαγαγεῖν αὐτοὺς ἐκ γῆς Αἰγύπτου, ὅτι αὐτοὶ οὐκ ἐνέμειναν ἐν τῇ διαθήκῃ μου, καὶ ἐγὼ ἠμέλησα αὐτῶν, φησὶ Κύριος.
\vs{33}Ὅτι αὕτη ἡ διαθήκη μου, ἣν διαθήσομαι τῷ οἴκῳ Ἰσραὴλ, μετὰ τὰς ἡμέρας ἐκείνας, φησὶ Κύριος, διδοὺς δώσω νόμους μου εἰς τὴν διάνοιαν αὐτῶν, καὶ ἐπὶ καρδίας αὐτῶν γράψω αὐτοὺς, καὶ ἔσομαι αὐτοῖς εἰς Θεὸν, καὶ αὐτοὶ ἔσονταί μοι εἰς λαόν.
\vs{34}Καὶ οὐ μὴ διδάξωσιν ἕκαστος τὸν πολίτην αὐτοῦ, καὶ ἕκαστος τὸν ἀδελφὸν αὐτοῦ, λέγων, γνῶθι τὸν Κύριον· ὅτι πάντες εἰδήσουσί με ἀπὸ μικροῦ αὐτῶν ἕως μεγάλου αὐτῶν, ὅτι ἵλεως ἔσομαι ταῖς ἀδικίαις αὐτῶν, καὶ τῶν ἁμαρτιῶν αὐτῶν οὐ μὴ μνησθῶ ἔτι.

\vs{35}Ἐὰν ὑψωθῇ ὁ οὐρανὸς εἰς τὸ μετέωρον, φησὶ Κύριος, καὶ ἐὰν ταπεινωθῇ τὸ ἔδαφος τῆς γῆς κάτω, καὶ ἐγὼ οὐκ ἀποδοκιμῶ τὸ γένος Ἰσραὴλ, φησὶ Κύριος, περὶ πάντων ὧν ἐποίησαν.

\vs{36}Οὕτως εἶπε Κύριος, ὁ δοὺς τὸν ἥλιον εἱς φῶς τῆς ἡμέρας, σελήνην καὶ ἀστέρας εἰς φῶς τῆς νυκτὸς, καὶ κραυγὴν ἐν θαλάσσῃ, καὶ ἐβόμβησε τὰ κύματα αὐτῆς, Κύριος παντοκράτωρ ὄνομα αὐτῷ·
\vs{37}Ἐὰν παύσωναι οἱ νόμοι οὗτοι ἀπὸ προσώπου μου, φησὶ Κύριος, καὶ τὸ γένος Ἰσραὴλ παύσεται γενέσθαι ἔθνος κατὰ πρόσωπόν μου πάσας τὰς ἡμέρας.

\vs{38}Ἰδοὺ ἡμέραι ἔρχονται, φησὶ Κύριος, καὶ οἰκοδομηθήσεται πόλις τῷ Κυρίῳ ἀπὸ πύργου Ἁναμεὴλ, ἕως πύλης τῆς γωνίας.
\vs{39}Καὶ ἐξελεύσεται ἡ διαμέτρησις αὐτῆς ἀπέναντι αὐτῶν ἕως βουνῶν Γαρὴβ, καὶ περικυκλωθήσεται κύκλῳ ἐξ ἐκλεκτῶν λίθων,
\vs{40}καὶ πάντες Ἀσαρημὼθ ἕως Νάχαλ Κέδρων, ἕως γωνίας πύλης ἵππων ἀνατολῆς, ἁγίασμα τῷ Κυρίῳ, καὶ οὐκέτι οὐ μὴ ἐκλίπῃ, καὶ οὐ μὴ καθαιρεθῇ ἕως τοῦ αἰῶνος.

\ch{39}
Ὁ ΛΟΓΟΣ ὁ γενόμενος παρὰ Κυρίου πρὸς Ἱερεμίαν ἐν τῷ ἐνιαυτῷ δεκάτῳ βασιλεῖ Σεδεκίᾳ, οὗτος ἐνιαυτὸς ὀκτωκαιδέκατος τῷ βασιλεῖ Ναβουχοδονόσορ βασιλεῖ Βαβυλῶνος.

\vs{2}Καὶ δύναμις βασιλέως Βαβυλῶνος ἐχαράκωσεν ἐπὶ Ἱερουσαλὴμ, καὶ Ἱερεμίας ἐφυλάσσετο ἐν αὐλῇ τῆς φυλακῆς, ἥ ἐστιν ἐν οἴκῳ βασιλέως,
\vs{3}ἐν ᾗ κατέκλεισεν αὐτὸν ὁ βασιλεὺς Σεδεκίας, λέγων, διατί σὺ προφητεύεις, λέγων, οὕτως εἶπε Κύριος, ἰδοὺ ἐγὼ δίδωμι τὴν πόλιν ταύτην ἐν χερσὶ βασιλέως Βαβυλῶνος, καὶ λήψεται αὐτὴν,
\vs{4}καὶ Σεδεκίας οὐ μὴ σωθῇ ἐκ χειρὸς τῶν Χαλδαίων, ὅτι παραδόσει παραδοθήσεται εἰς χεῖρας βασιλέως Βαβυλῶνος, καὶ λαλήσει στόμα αὐτοῦ πρὸς στόμα αὐτοῦ, καὶ οἱ ὀφθαλμοὶ αὐτοῦ τοὺς ὀφθαλμοὺς αὐτοῦ ὄψονται·
\vs{5}Καὶ εἰσελεύσεται Σεδεκίας εἰς Βαβυλῶνα, καὶ ἐκεῖ καθίεται;

\vs{6}ΚΑΙ Ὁ ΛΟΓΟΣ ΚΥΡΙΟΥ ἘΓΕΝΗΘΗ ΠΡΟΣ ἹΕΡΕΜΙΑΝ, ΛΕΓΩΝ,
\vs{7}Ἰδοὺ Ἀναμεὴλ υἱὸς Σαλῶμ ἀδελφοῦ πατρός σου ἔρχεται πρὸς σὲ, λέγων, κτῆσαι σεαυτῷ τὸν ἀγρόν μου τὸν ἐν Ἀναθὼθ, ὅτι σοὶ κρίσις παραλαβεῖν εἰς κτῆσιν.
\vs{8}Καὶ ἦλθε πρὸς μὲ Ἀναμεὴλ υἱὸς Σαλὼμ, ἀδελφοῦ πατρός μου, εἰς τὴν αὐλὴν τῆς φυλακῆς, καὶ εἶπε, κτῆσαι σεαυτῷ τὸν ἀγρόν μου τὸν ἐν γῇ Βενιαμὶν τὸν ἐν Ἀναθὼθ, ὅτι σοὶ κρίμα κτήσασθαι αὐτὸν, καὶ σὺ πρεσβύτερος. καὶ ἔγνων, ὅτι λόγος Κυρίου ἐστὶ,
\vs{9}καὶ ἐκτησάμην τὸν ἀγρὸν Ἀναμεὴλ υἱοῦ ἀδελφοῦ πατρός μου, καὶ ἔστησα αὐτῷ ἑπτὰ σίκλους καὶ δέκα ἀργυρίου,
\vs{10}καὶ ἔγραψα εἰς βιβλίον, καὶ ἐσφραγισάμην, καὶ διεμαρτυράμην μάρτυρας, καὶ ἔστησα τὸ ἀργύριον ἐν ζυγῷ.
\vs{11}Καὶ ἔλαβον τὸ βιβλίον τῆς κτήσεως τὸ ἐσφραγισμένον,
\vs{12}καὶ ἔδωκα αὐτὸ τῷ Βαροὺχ υἱῷ Νηρίου υἱῷ Μαασαίου, κατʼ ὀφθαλμοὺς Ἀναμεὴλ υἱοῦ ἀδελφοῦ πατρός μου, καὶ κατʼ ὀφθαλμοὺς τῶν ἀνδρῶν τῶν παρεστηκότων καὶ γραφόντων ἐν τῷ βιβλίῳ τῆς κτήσεως, καὶ κατʼ ὀφθαλμοὺς τῶν Ἰουδαίων τῶν ἐν τῇ αὐλῇ τῆς φυλακῆς.
\vs{13}Καὶ συνέταξα τῷ Βαροὺχ κατʼ ὀφθαλμοὺς αὐτῶν, λέγων,
\vs{14}Οὕτως εἶπε Κύριος παντοκράτωρ, λάβε τὸ βιβλίον τῆς κτήσεως τοῦτο, καὶ τὸ βιβλίον τὸ ἀνεγνωσμένον, καὶ θήσεις αὐτὸ εἰς ἀγγεῖον ὀστράκινον, ἵνα διαμείνῃ ἡμέρας πλείους.
\vs{15}Ὅτι οὕτως εἶπε Κύριος, ἔτι κτιηθήσονται ἀγροὶ, καὶ οἰκίαι, καὶ ἀμπελῶνες ἐν τῇ γῇ ταύτῃ.

\vs{16}Καὶ προσευξάμην πρὸς Κύριον μετὰ τὸ δοῦναί με τὸ βιβλίον τῆς κτήσεως πρὸς Βαροὺχ υἱὸν Νηρίου, λέγων,

\vs{17}Ὁ ὢν Κύριε, σὺ ἐποίησας τὸν οὐρανὸν καὶ τὴν γῆν τῇ ἰσχύϊ σου τῇ μεγάλῃ, καὶ τῷ βραχίονί σου τῷ ὑψηλῷ καὶ τῷ μετεώρῳ, οὐ μὴ ἀποκρυβῇ ἀπὸ σοῦ οὐθὲν,
\vs{18}ποιῶν ἔλεος εἰς χιλιάδας, καὶ ἀποδιδοὺς ἁμαρτίας πατέρων εἰς κόλπους τέκνων αὐτῶν μετʼ αὐτούς· ὁ Θεὸς ὁ μέγας, ὁ ἰσχυρὸς,
\vs{19}Κύριος μεγάλης βουλῆς, καὶ δυνατὸς τοῖς ἔργοις, ὁ Θεὸς ὁ μέγας ὁ παντοκράτωρ, καὶ μεγαλώνυμος Κύριος· οἱ ὀφθαλμοί σου εἰς τὰς ὁδοὺς τῶν υἱῶν τῶν ἀνθρώπων, δοῦναι ἑκάστῳ κατὰ τὴν ὁδὸν αὐτοῦ·
\vs{20}Ὃς ἐποίησας σημεῖα καὶ τέρατα ἐν γῇ Αἰγύπτῳ ἕως τῆς ἡμέρας ταύτης, καὶ ἐν Ἰσραὴλ, καὶ ἐν τοῖς γηγενέσι· καὶ ἐποίησας σεαυτῷ ὄνομα, ὡς ἡμέρα αὕτη,
\vs{21}καὶ ἐξήγαγες τὸν λαόν σου Ἰσραὴλ ἐκ γῆς Αἰγύπτου ἐν σημείοις καὶ ἐν τέρασιν, ἐν χειρὶ κραταιᾷ,
\vs{22}καὶ ἐν βραχίονι ὑψηλῷ, καὶ ἐν ὁράμασι μεγάλοις, καὶ ἔδωκας αὐτοῖς τὴν γῆν ταύτην, ἣν ὤμοσας τοῖς πατράσιν αὐτῶν, γῆν ῥέουσαν γάλα καὶ μέλι.
\vs{23}Καὶ εἰσήλθοσαν καὶ ἔλαβον αὐτὴν, καὶ οὐκ ἤκουσαν τῆς φωνῆς σου, καὶ ἐν τοῖς προστάγμασί σου οὐκ ἐπορεύθησαν· ἅπαντα ἃ ἐνετείλω αὐτοῖς οὐκ ἐποίησαν, καὶ ἐποὶησαν συμβῆναι αὐτοῖς πάντα τὰ κακὰ ταῦτα.
\vs{24}Ἰδοὺ ὄχλος ἥκει εἰς τὴν πόλιν συλλαβεῖν αὐτὴν, καὶ ἡ πόλις ἐδόθη εἰς χεῖρας Χαλδαίων τῶν πολεμούντων αὐτὴν ἀπὸ προσώπου μαχαίρας, καὶ τοῦ λιμοῦ ὡς ἐλάλησας, οὕτως ἐγένετο.
\vs{25}Καὶ σὺ λέγεις πρὸς μέ κτῆσαι σεαυτῷ τὸν ἀγρὸν ἀργυρίου· καὶ ἔγραψα βιβλίον, καὶ ἐσφραγισάμην, καὶ ἐπεμαρτυράμην μάρτυρας, καὶ ἡ πόλις ἐδόθη εἰς χεῖρας Χαλδαίων.

\vs{26}Καὶ ἐγένετο λόγος Κυρίου πρὸς μὲ, λέγων,

\vs{27}Ἐγὼ Κύριος ὁ Θεὸς πάσης σαρκός, μὴ ἀπʼ ἐμοῦ κρυβήσεταί τι;
\vs{28}Διατοῦτο οὕτως εἶπε Κύριος ὁ Θεὸς Ἰσραὴλ, δοθεῖσα παραδοθήσεται ἡ πόλις αὕτη εἰς χεῖρας βασιλέως Βαβυλῶνος, καὶ λήψεται αὐτὴν,
\vs{29}καὶ ἥξουσιν οἱ Χαλδαῖοι πολεμοῦντες ἐπὶ τὴν πόλιν ταύτην, καὶ καύσουσι τὴν πόλιν ταύτην ἐν πυρί, καὶ κατακαύσουσι τὰς οἰκίας ἐν αἷς ἐθυμιῶσαν ἐπὶ τῶν δωμάτων αὐτῶν τῇ Βάαλ, καὶ ἔσπενδον σπονδὰς θεοῖς ἑτέροις, πρὸς τὸ παραπικράναι με·
\vs{30}Ὅτι ἦσαν οἱ υἱοὶ Ἰσραὴλ καὶ οἱ υἱοὶ Ἰούδα μόνοι ποιοῦντες τὸ πονηρὸν κατʼ ὀφθαλμούς μου ἐκ νεὁτητος αὐτῶν·
\vs{31}Ὅτι ἐπὶ τὴν ὀργήν μου, καὶ ἐπὶ τὸν θυμόν μου ἦν ἡ πόλις αὕτη, ἀφʼ ἧς ἡμέρας ᾠκοδόμησαν αὐτὴν καὶ ἕως τῆς ἡμέρας ταύτης, ἀπαλλάξαι αὐτὴν ἀπὸ προσώπου μου,
\vs{32}διὰ πάσας τὰς πονηρίας τῶν υἱῶν Ἰσραὴλ καὶ Ἰούδα, ὧν ἐποίησαν πικράναι με, αὐτοὶ καὶ οἱ βασιλεῖς αὐτῶν, καὶ οἱ ἄρχοντες αὐτῶν, καὶ οἱ ἱερεῖς αὐτῶν, καὶ οἱ προφῆται αὐτῶν, ἄνδρες Ἰούδα, καὶ οἱ κατοικοῦντες ἐν Ἱερουσαλὴμ,
\vs{33}καὶ ἀπέστρεψαν πρὸς μὲ νῶτον, καὶ οὐ πρόσωπον· καὶ ἐδίδαξα αὐτοὺς ὄρθρου, καὶ οὐκ ἤκουσαν ἔτι λαβεῖν παιδείαν.
\vs{34}Καὶ ἔθηκαν τὰ μιάσματα αὐτῶν ἐν τῷ οἴκῳ, οὗ ἐπεκλήθη τὸ ὄνομά μου ἐπʼ αὐτῷ, ἐν ἀκαθαρσίαις αὐτῶν.
\vs{35}Καὶ ᾠκοδόμησαν τοὺς βωμοὺς τῇ Βάαλ τοὺς ἐν φάραγγι υἱοὺ Ἑννὸμ, τοῦ ἀναφέρειν τοὺς υἱοὺς αὐτῶν καὶ τὰς θυγατέρας αὐτῶν τῷ Μολὸχ βασιλεῖ, ἃ οὐ συνέταξα αὐτοῖς, καὶ οὐκ ἀνέβη ἐπὶ καρδίαν μου τοῦ ποιῆσαι τὸ βδέλυγμα τοῦτο, πρὸς τὸ ἐφαμαρτεῖν τὸν Ἰούδαν.

\vs{36}Καὶ νῦν οὕτως εἶπε Κύριος ὁ Θεὸς Ἰσραὴλ ἐπὶ τὴν πόλιν, ἣν σὺ λέγεις, παραδοθήσεται εἰς χεῖρας βασιλέως Βαβυλῶνος ἐν μαχαίρᾳ, καὶ ἐν λιμῷ καὶ ἐν ἀποστολῇ·
\vs{37}Ἰδοὺ ἐγὼ συνάγω αὐτοὺς ἐκ πάσης τῆς γῆς, οὗ διέσπειρα αὐτοὺς ἐκεῖ ἐν ὀργῇ μου, καὶ τῷ θυμῷ μου, καὶ ἐν παροξυσμῷ μεγάλῳ· καὶ ἐπιστρέψω αὐτοὺς εἰς τὸν τόπον τοῦτον, καὶ καθιῶ αὐτοὺς πεποιθότας·
\vs{38}Καὶ ἔσονταί μοι εἰς λαὸν, καὶ ἐγὼ ἔσομαι αὐτοῖς εἰς Θεόν.
\vs{39}Καὶ δώσω αὐτοῖς ὁδὸν ἑτέραν καὶ καρδίαν ἑτέραν, φοβηθῆναί με πάσας τὰς ἡμέρας, καὶ εἰς ἀγαθὸν αὐτοῖς καὶ τοῖς τέκνοις αὐτῶν μετʼ αὐτούς.
\vs{40}Καὶ διαθήσομαι αὐτοῖς διαθήκην αἰωνίαν, ἣν οὐ μὴ ἀποστρέψω ὄπισθεν αὐτῶν· καὶ τὸν φόβον μου δώσω εἰς τὴν καρδίαν αὐτῶν, πρὸς τὸ μὴ ἀποστῆναι αὐτοὺς ἀπʼ ἐμοῦ.
\vs{41}Καὶ ἐπισκέψομαι τοῦ ἀγαθῶσαι αὐτοὺς, καὶ φυτεύσω αὐτοὺς ἐν τῇ γῇ ταύτῃ ἐν πίστει, καὶ ἐν πάσῃ καρδίᾳ μου, καὶ ἐν πάσῃ ψυχῇ.

\vs{42}Ὅτι οὕτως εἶπε Κύριος, καθὰ ἐπήγαγον ἐπὶ τὸν λαὸν τοῦτον πάντα τὰ κακὰ τὰ μεγάλα ταῦτα, οὕτως ἐγὼ ἐπάξω ἐπʼ αὐτοὺς πάντα τὰ ἀγαθὰ, ἃ ἐλάλησα ἐπʼ αὐτούς.
\vs{43}Καὶ κτηθήσονται ἔτι ἀγροὶ ἐν τῇ γῇ, ᾗ σὺ λέγεις, ἄβατος ἔσται ἀπὸ ἀνθρώπων καὶ κτήνους, καὶ παρεδόθησαν εἰς χεῖρας Χαλδαίων.
\vs{44}Καὶ κτήσονται ἀγροὺς ἐν ἀργυρίῳ· καὶ γράψεις βιβλίον καὶ σφραγιῇ, καὶ διαμαρτύρῃ μάρτυρας ἐν γῇ Βενιαμὶν, καὶ κύκλῳ τῆς Ἱερουσαλὴμ, καὶ ἐν πόλεσιν Ἰούδα, καὶ ἐν πόλεσι τοῦ ὄρους, καὶ ἐν πόλεσι τῆς σεφηλὰ, καὶ ἐν πόλεσι τῆς ναγὲβ, ὅτι ἀποστρέψω τὰς ἀποικίας αὐτῶν.

\ch{40}
Καὶ ἐγένετο λόγος Κυρίου πρὸς Ἱερεμίαν δεύτερον, καὶ αὐτὸς ἦν ἔτι δεδεμένος ἐν τῇ αὐλῇ τῆς φυλακῆς, λέγων,

\vs{2}Οὕτως εἶπε Κύριος, ποιῶν γῆν, καὶ πλάσσων αὐτὴν, τοῦ ἀνορθῶσαι αὐτὴν, Κύριος ὄνομα αὐτῷ·
\vs{3}Κέκραξον πρὸς μὲ, καὶ ἀποκριθήσομαί σοι, καὶ ἀπαγγελῶ σοι μεγάλα καὶ ἰσχυρὰ, ἃ οὐκ ἔγνως αὐτά.
\vs{4}Ὅτι οὕτως εἶπε Κύριος περὶ οἴκων τῆς πόλεως ταύτης, καὶ περὶ οἴκων βασιλεως Ἰούδα τῶν καθῃρημένων εἰς χάρακας καὶ προμαχῶνας,
\vs{5}τοῦ μάχεσθαι πρὸς τοὺς Χαλδαίους, καὶ πληρῶσαι αὐτὴν τῶν νεκρῶν τῶν ἀνθρώπων, οὓς ἐπάταξα ἐν ὀργῇ μου, καὶ ἐν θυμῷ μου· καὶ ἀπέστρεψα τὸ πρόσωπόν μου ἀπʼ αὐτῶν, περὶ πασῶν τῶν πονηριῶν αὐτῶν.
\vs{6}Ἰδοὺ ἐγὼ ἀνάγω αὐτῇ συνούλωσιν καὶ ἴαμα, καὶ φανερώσω αὐτοῖς, καὶ ἰατρεύσω αὐτὴν, καὶ ποιήσω καὶ εἰρήνην καὶ πίστιν.

\vs{7}Καὶ ἀποστρέψω τὴν ἀποικίαν Ἰούδα, καὶ ἀποικίαν Ἰσραὴλ, καὶ οἰκοδομήσω αὐτοὺς καθὼς καὶ τοπρότερον.
\vs{8}Καὶ καθαριῶ αὐτοὺς ἀπὸ πασῶν τῶν ἀδικιῶν αὐτῶν, ὧν ἡμάρτοσάν μοι, καὶ οὐ μὴ μνησθήσομαι ἁμαρτιῶν αὐτῶν, ὧν ἥμαρτόν μοι, καὶ ἀπέστησαν ἀπʼ ἐμοῦ.
\vs{9}Καὶ ἔσται εἰς εὐφροσύνην καὶ αἴνεσιν, καὶ εἰς μεγαλειότητα παντὶ τῷ λαῷ τῆς γῆς, οἵτινες ἀκούσονται πάντα τὰ ἀγαθὰ ἃ ἐγὼ ποιήσω, καὶ φοβηθήσονται καὶ πικρανθήσονται περὶ πάντων τῶν ἀγαθῶν, καὶ περὶ πάσης τῆς εἰρήνης ἧς ἐγὼ ποιήσω αὐτοῖς.

\vs{10}Οὕτως εἶπε Κύριος, ἔτι ἀκουσθήσεται ἐν τῷ τόπῳ τούτῳ, ᾧ ὑμεῖς λέγετε, ἔρημός ἐστιν ἀπὸ ἀνθρώπων καὶ κτηνῶν, ἐν πόλεσιν Ἰούδα, καὶ ἔξωθεν Ἱερουσαλὴμ, ταῖς ἠρημωμέναις, παρὰ τὸ μὴ εἶναι ἄνθρωπον καὶ κτήνη·
\vs{11}Φωνὴ εὐφροσύνης, καὶ φωνὴ χαρμοσύνης, φωνὴ νυμφίου, καὶ φωνὴ νύμφης, φωνὴ λεγόντων, ἐξομολογεῖσθε Κυρίῳ παντοκράτορι, ὅτι χρηστὸς Κύριος, ὅτι εἰς τὸν αἰῶνα τὸ ἔλεος αὐτοῦ· καὶ εἰσοίσουσι δῶρα εἰς οἶκον Κυρίου, ὅτι ἀποστρέψω πᾶσαν τὴν ἀποικίαν τῆς γῆς ἐκείνης κατὰ τοπρότερον, εἶπε Κύριος.
\vs{12}Οὕτως εἶπε Κύριος τῶν δυνάμεων, ἔτι ἔσται ἐν τῷ τόπῳ τούτῳ τῷ ἐρήμῳ, παρὰ τὸ μὴ εἶναι ἄνθρωπον καὶ κτῆνος, ἐν πάσαις ταῖς πόλεσιν αὐτοῦ καταλύματα ποιμένων κοιταζόντων πρόβατα,
\vs{13}ἐν πόλεσι τῆς ὀρεινῆς, καὶ ἐν πόλεσι τῆς σεφηλὰ, καὶ ἐν πόλεσι τῆς ναγὲβ, καὶ ἐν γῇ Βενιαμὶν, καὶ ἐν ταῖς κύκλῳ Ἱερουσαλὴμ, καὶ ἐν πόλεσιν Ἰούδα· ἔτι παρελεύσεται πρόβατα ἐπὶ χεῖρα ἀριθμοῦντος, εἶπε Κύριος.

\ch{41}
Ὁ ΛΟΓΟΣ ὁ γενόμενος πρὸς Ἱερεμίαν παρὰ Κυρίου, (καὶ Ναβουχοδονόσορ βασιλεὺς Βαβυλῶνος, καὶ πᾶν τὸ στρατόπεδον αὐτοῦ, καὶ πᾶσα ἡ γῆ ἀρχῆς αὐτοῦ ἐπολέμουν ἐπὶ Ἱερουσαλὴμ, καὶ ἐπὶ πάσας τὰς πόλεις Ἰούδα,) λέγων,

\vs{2}Οὕτως εἶπε Κύριος, βάδισον πρὸς Σεδεκίαν βασιλέα Ἰούδα, καὶ ἐρεῖς αὐτῷ, οὕτως εἶπε Κύριος, παραδόσει παραδοθήσεται ἡ πόλις αὕτη εἰς χεῖρας βασιλέως Βαβυλῶνος, καὶ συλλήψεται αὐτὴν, καὶ καύσει αὐτὴν ἐν πυρὶ,
\vs{3}καὶ σὺ οὐ μὴ σωθῇς ἐκ χειρὸς αὐτοῦ, καὶ συλλήψει συλληφθήσῃ, καὶ εἰς χεῖρας αὐτοῦ δοθήσῃ, καὶ ὀφθαλμοί σου τοὺς ὀφθαλμοὺς αὐτοῦ ὄψονται, καὶ εἰς Βαβυλῶνα εἰσελεύσῃ.

\vs{4}Ἀλλʼ ἄκουσον τὸν λόγον Κυρίου, Σεδεκία βασιλεῦ Ἰούδα· οὕτως λέγει Κύριος,
\vs{5}ἐν εἰρήνῃ ἀποθανῇ· καὶ ὡς ἔκλαυσαν τοὺς πατέρας σου τοὺς βασιλεύσαντας πρότερόν σου, κλαύσονται καὶ σὲ, οὐαὶ Κύριε, καὶ ἕως ᾅδου κόψονταί σε, ὅτι λόγον ἐγὼ ἐλάλησα, εἶπε Κύριος.

\vs{6}Καὶ ἐλάλησεν Ἱερεμίας πρὸς τὸν βασιλέα Σεδεκίαν πάντας τοὺς λόγους τούτους ἐν Ἱερουσαλήμ.
\vs{7}Καὶ ἡ δύναμις βασιλέως Βαβυλῶνος ἐπολέμει ἐπὶ Ἱερουσαλὴμ, καὶ ἐπὶ τὰς πόλεις Ἰούδα, καὶ ἐπὶ Λαχὶς, καὶ ἐπὶ Ἄζηκα, ὅτι αὗται κατελείφθησαν ἐν πόλεσιν Ἰούδα πόλεις ὀχυραί.

\vs{8}Ὁ λόγος ὁ γενόμενος πρὸς Ἱερεμίαν παρὰ Κυρίου, μετὰ τὸ συντελέσαι τὸν βασιλέα Σεδεκίαν διαθήκην πρὸς τὸν λαὸν, τοῦ καλέσαι ἄφεσιν,
\vs{9}τοῦ ἐξαποστεῖλαι ἕκαστον τὸν παῖδα αὐτοῦ, καὶ ἕκαστον τὴν παιδίσκην αὐτοῦ, τὸν Ἑβραῖον καὶ τὴν Ἑβραίαν ἐλευθέρους, πρὸς τὸ μὴ δουλεύειν ἄνδρα ἐξ Ἰούδα.
\vs{10}Καὶ ἐπεστράφησαν πάντες οἱ μεγιστᾶνες, καὶ πᾶς ὁ λαὸς οἱ εἰσελθόντες ἐν τῇ διαθήκῃ, τοῦ ἀποστεῖλαι ἕκαστον τὸν παῖδα αὐτοῦ, καὶ ἕκαστον τὴν παιδίσκην αὐτοῦ,
\vs{11}καὶ ἐῶσαν αὐτοὺς εἰς παῖδας καὶ παιδίσκας.

\vs{12}Καὶ ἐγενήθη λόγος Κυρίου πρὸς Ἱερεμίαν, λέγων,
\vs{13}Οὕτως εἶπε Κύριος, ἐγὼ διεθέμην διαθήκην πρὸς τοὺς πατέρας ὑμῶν ἐν τῇ ἡμέρᾳ ᾗ ἐξειλάμην αὐτοὺς ἐκ γῆς Αἰγύπτου, ἐκ οἴκου δουλείας, λέγων,
\vs{14}ὅταν πληρωθῇ ἓξ ἔτη, ἀποστελεῖς τὸν ἀδελφόν σου τὸν Ἑβραῖον, ὃς πραθήσεταί σοι, καὶ ἐργᾶταί σοι ἓξ ἔτη, καὶ ἐξαποστελεῖς αὐτὸν ἐλεύθερον· καὶ οὐκ ἤκουσάν μου, καὶ οὐκ ἔκλιναν τὸ οὖς αὐτῶν.
\vs{15}Καὶ ἐπέστρεψαν σήμερον ποιῆσαι τὸ εὐθὲς πρὸ ὀφθαλμῶν μου, τοῦ καλέσαι ἄφεσιν ἕκαστον τοῦ πλησίον αὐτοῦ· καὶ συνετέλεσαν διαθήκην κατὰ πρόσωπόν μου, ἐν τῷ οἴκῳ οὗ ἐπεκλήθη τὸ ὄνομά μου ἐπʼ αὐτῷ.
\vs{16}Καὶ ἐπεστρέψατε, καὶ ἐβεβηλώσατε τὸ ὄνομά μου, τοῦ ἐπιστρέψαι ἕκαστον τὸν παῖδα αὐτοῦ, καὶ ἕκαστον τὴν παιδίσκην αὐτοῦ, οὓς ἐξαπεστείλατε ἐλευθέρους τῇ ψυχῇ αὐτῶν, τοῦ εἶναι ὑμῖν εἰς παῖδας καὶ παιδίσκας.

\vs{17}Διατοῦτο οὕτως εἶπε Κύριος, ὑμεῖς οὐκ ἠκούσατέ μου, τοῦ καλέσαι ἄφεσιν ἕκαστος πρὸς τὸν πλησίον αὐτοῦ· ἰδοὺ ἐγὼ καλῶ ἄφεσιν ὑμῖν εἰς μάχαιραν, καὶ εἰς τὸν θάνατον, καὶ εἰς τὸν λιμὸν, καὶ δώσω ὑμᾶς εἰς διασπορὰν πάσαις ταῖς βασιλείαις τῆς γῆς·
\vs{18}Καὶ δώσω τοὺς ἄνδρας τοὺς παρεληλυθότας τὴν διαθήκην μου, τοὺς μὴ στήσαντας τὴν διαθήκην μου, ἣν ἐποίησαν κατὰ πρόσωπόν μου, τὸν μόσχον ὃν ἐποίησαν ἐργάζεσθαι αὐτῷ,
\vs{19}τοὺς ἄρχοντας Ἰούδα, καὶ τοὺς δυνάστας, καὶ τοὺς ἱερεῖς, καὶ τὸν λαόν·
\vs{20}Καὶ δώσω αὐτοὺς τοῖς ἐχθροῖς αὐτῶν, καὶ ἔσται τὰ θνησιμαῖα αὐτῶν βρῶσις τοῖς πετεινοῖς τοῦ οὐρανοῦ καὶ τοῖς θηρίοις τῆς γῆς.
\vs{21}Καὶ τὸν Σεδεκίαν βασιλέα τῆς Ἰουδαίας, καὶ τοὺς ἄρχοντας αὐτῶν δώσω εἰς χεῖρας ἐχθρῶν αὐτῶν, καὶ δύναμις βασιλέως Βαβυλῶνος τοῖς ἀποτρέχουσιν ἀπʼ αὐτῶν.
\vs{22}Ἰδοὺ ἐγὼ συντάσσω, φησὶ Κύριος, καὶ ἐπιστρέψω αὐτοὺς εἰς τὴν γῆν ταύτην, καὶ πολεμήσουσιν ἐπʼ αὐτὴν, καὶ λήψονται αὐτὴν, καὶ κατακαύσουσιν αὐτὴν ἐν πυρὶ, καὶ τὰς πόλεις Ἰούδα, καὶ δώσω αὐτὰς ἐρήμους ἀπὸ τῶν κατοικούντων.

\ch{42}
Ὁ ΛΟΓΟΣ Ὁ ΓΕΝΟΜΕΝΟΣ ΠΡΟΣ ἹΕΡΕΜΙΑΝ παρὰ Κυρίου ἐν ἡμέραις Ἰωακεὶμ βασιλέως Ἰούδα, λέγων,
\vs{2}Βάδισον εἰς οἶκον Ἀρχαβεὶν, καὶ ἄξεις αὐτοὺς εἰς οἶκον Κυρίου, εἰς μίαν τῶν αὐλῶν, καὶ ποτιεῖς αὐτοὺς οἶνον.

\vs{3}Καὶ ἐξήγαγον τὸν Ἰεχονιαν υἱὸν Ἱερεμὶν υἱοῦ Χαβασὶν, καὶ τοὺς ἀδελφοὺς αὐτοῦ, καὶ τοὺς υἱοὺς αὐτοῦ, καὶ πᾶσαν τὴν οἰκίαν Ἀρχαβεὶν,
\vs{4}καὶ εἰσήγαγον αὐτοὺς εἰς οἶκον Κυρίου εἰς τὸ παστοφόριον υἱῶν Ἰωνὰν, υἱοῦ Ἀνανίου, υἱοῦ Γοδολίου ἀνθρώπου τοῦ Θεοῦ, ὅς ἐστιν ἐγγὺς τοῦ οἴκου τῶν ἀρχόντων τῶν ἐπάνω τοῦ οἴκου Μαασαίου υἱοῦ Σελὼμ, τοῦ φυλάσσοντος τὴν αὐλήν.
\vs{5}Καὶ ἔδωκα κατὰ πρόσωπον αὐτῶν κεράμιον οἴνου, καὶ ποτήρια, καὶ εἶπα, πίετε οἴνον.

\vs{6}Καὶ εἶπον, οὐ μὴ πίωμεν οἶνον, ὅτι Ἰωναδὰβ υἱὸς Ῥηχὰβ ὁ πατὴρ ἡμῶν ἐνετείλατο ἡμῖν, λέγων, οὐ μὴ πίητε οἶνον ὑμεῖς καὶ οἱ υἱοὶ υμῶν ἕως αἰῶνος,
\vs{7}καὶ οἰκίας οὐ μὴ οἰκοδομήσητε, καὶ σπέρμα οὐ μὴ σπείρητε, καὶ ἀμπελὼν οὐκ ἔσται ὑμῖν, ὅτι ἐν σκηναῖς οἰκήσετε πάσας τὰς ἡμέρας ὑμῶν, ὅπως ἂν ζήσητε ἡμέρας πολλὰς ἐπὶ τῆς γῆς, ἐφʼ ἧς διατρίβετε ὑμεῖς ἐπʼ αὐτῆς.
\vs{8}Καὶ ἠκούσαμεν τῆς φωνῆς Ἰωναδὰβ τοῦ πατρὸς ἡμῶν, πρὸς τὸ μὴ πιεῖν οἶνον πάσας τὰς ἡμέρας ἡμῶν, ἡμεῖς καὶ αἱ γυναῖκες ἡμῶν, καὶ οἱ υἱοὶ ἡμῶν, καὶ αἱ θυγατέρες ἡμῶν,
\vs{9}καὶ πρὸς τὸ μὴ οἰκοδομεῖν οἰκίας τοῦ κατοικεῖν ἐκεῖ, καὶ ἀμπελὼν καὶ ἀγρὸς καὶ σπέρμα οὐκ ἐγένετο ἡμῖν.
\vs{10}Καὶ ᾠκήσαμεν ἐν σκηναῖς, καὶ ἠκούσαμεν, καὶ ἐποιήσαμεν κατὰ πάντα ἃ ἐνετείλατο ἡμῖν Ἰωναδὰβ ὁ πατὴρ ἡμῶν.
\vs{11}Καὶ ἐγενήθη ὅτε ἀνέβη Ναβουχοδονόσορ ἐπὶ τὴν γῆν, καὶ εἴπαμεν εἰσελθεῖν, καὶ εἰσήλθομεν εἰς Ἱερουσαλὴμ, ἀπὸ προσώπου τῆς δυνάμεως τῶν Χαλδαίων, καὶ ἀπὸ προσώπου τῆς δυνάμεως τῶν Ἀσσυρίων, καὶ ᾠκοῦμεν ἐκεῖ.

\vs{12}Καὶ ἐγένετο λόγος Κυρίου πρὸς μὲ, λέγων,
\vs{13}οὕτως λέγει Κύριος, πορεύου, καὶ εἰπὸν ἀνθρώπῳ Ἰούδα, καὶ τοῖς κατοικοῦσιν Ἱερουσαλὴμ, οὐ μὴ λάβητε παιδείαν τοῦ ἀκούειν τοὺς λόγους μου;
\vs{14}Ἔστησαν ῥῆμα υἱοὶ Ἰωναδὰβ υἱοῦ Ῥηχὰβ, ὃ ἐνετείλατο τοῖς τέκνοις αὐτοῦ πρὸς τὸ μὴ πιεῖν οἶνον, καὶ οὐκ ἐπίοσαν· καὶ ἐγὼ ἐλάλησα πρὸς ὑμᾶς ὄρθρου, καὶ οὐκ ἠκούσατε.
\vs{15}Καὶ ἀπέστειλα πρὸς ὑμᾶς τοὺς παῖδάς μου τοὺς προφήτας, λέγων, ἀποστράφητε ἕκαστος ἀπὸ τῆς ὁδοῦ αὐτοῦ τῆς πονηρᾶς, καὶ βελτίω ποιήσατε τὰ ἐπιτηδεύματα ὑμῶν, καὶ οὐ πορεύεσθε ὀπίσω θεῶν ἑτέρων τοῦ δουλεύειν αὐτοῖς, καὶ οἰκήσετε ἐπὶ τῆς γῆς, ἧς ἔδωκα ὑμῖν, καὶ τοῖς πατράσιν ὑμῶν· καὶ οὐκ ἐκλίνατε τὰ ὦτα ὑμῶν, καὶ οὐκ εἰσηκούσατε.
\vs{16}Καὶ ἔστησαν υἱοὶ Ἰωναδὰβ υἱοῦ Ῥηχὰβ τὴν ἐντολὴν τοῦ πατρὸς αὐτῶν, ὁ δὲ λαὸς οὗτος οὐκ ἤκουσέ μου.
\vs{17}Διατοῦτο οὕτως εἶπε Κύριος, ἰδοὺ ἐγὼ φέρω ἐπὶ Ἰούδαν καὶ ἐπὶ τοὺς κατοικοῦντας Ἱερουσαλὴμ πάντα τὰ κακὰ ἃ ἐλάλησα ἐπʼ αὐτούς.

\vs{18}Διατοῦτο οὕτως εἶπε Κύριος, ἐπειδὴ ἤκουσαν υἱοὶ Ἰωναδὰβ υἱοῦ Ῥηχὰβ τὴν ἐντολὴν τοῦ πατρὸς αὐτῶν ποιεῖν καθότι ἐνετείλατο αὐτοῖς ὁ πατὴρ αὐτῶν,
\vs{19}οὐ μὴ ἐκλείπῃ ἀνὴρ τῶν υἱῶν Ἰωναδὰβ υἱοῦ Ῥηχὰβ παρεστηκὼς κατὰ πρόσωπόν μου πάσας τὰς ἡμέρας τῆς γῆς.

\ch{43}
ἘΝ ΤΩ ἘΝΙΑΥΤΩ ΤΩ ΤΕΤΑΡΤΩ ἸΩΑΚΕΙΜ υἱοῦ Ἰωσία βασιλέως Ἰούδα, ἐγενήθη λόγος Κυρίου πρὸς μὲ, λέγων,

\vs{2}Λάβε σεαυτῷ χαρτίον βιβλίου, καὶ γράψον ἐπʼ αὐτοῦ πάντας τοὺς λόγους οὓς ἐλάλησα πρὸς σὲ ἐπὶ Ἱερουσαλὴμ, καὶ ἐπὶ Ἰούδα, καὶ ἐπὶ πάντα τὰ ἔθνη, ἀφʼ ἧς ἡμέρας λαλήσαντός μου πρὸς σὲ ἀφʼ ἡμερῶν Ἰωσία βασιλέως Ἰούδα, καὶ ἕως τῆς ἡμέρας ταύτης.
\vs{3}Ἴσως ἀκούσεται ὁ οἶκος Ἰούδα πάντα τὰ κακὰ ἃ ἐγὼ λογίζομαι ποιῆσαι αὐτοῖς, ἵνα ἀποστρέψωσιν ἀπὸ τῆς ὁδοῦ αὐτῶν τῆς πονηρᾶς, καὶ ἵλεως ἔσομαι ταῖς ἀδικίαις αὐτῶν καὶ ταῖς ἁμαρτίαις αὐτῶν.

\vs{4}Καὶ ἐκάλεσεν Ἱερεμίας τὸν Βαροὺχ υἱὸν Νηρίου· καὶ ἔγραψεν ἀπὸ στόματος Ἱερεμίου πάντας τοὺς λόγους Κυρίου, οὓς ἐλάλησε πρὸς αὐτὸν, εἰς χαρτίον βιβλίου.
\vs{5}Καὶ ἐνετείλατο Ἱερεμίας τῷ Βαροὺχ, λέγων, ἐγὼ φυλάσσομαι, οὐ μὴ δύνωμαι εἰσελθεῖν εἰς οἶκον Κυρίου·
\vs{6}Καὶ ἀναγνώσῃ ἐν τῷ χαρτίῳ τούτῳ εἰς τὰ ὦτα τοῦ λαοῦ ἐν οἴκῳ Κυρίου, ἐν ἡμέρᾳ νηστείας, καὶ ἐν ὠσὶ παντὸς Ἰούδα τῶν ἐρχομένων ἐκ πόλεων αὐτῶν, ἀναγνώσῃ αὐτοῖς.
\vs{7}Ἴσως πεσεῖται ἔλεος αὐτῶν κατὰ πρόσωπον Κυρίου, καὶ ἀποστρέψουσιν ἐκ τῆς ὁδοῦ αὐτῶν τῆς πονηρᾶς, ὅτι μέγας ὁ θυμὸς καὶ ἡ ὀργὴ Κυρίου, ἣν ἐλάλησεν ἐπὶ τὸν λαὸν τοῦτον.

\vs{8}Καὶ ἐποίησε Βαροὺχ κατὰ πάντα ἃ ἐνετείλατο αὐτῷ Ἱερεμίας, τοῦ ἀναγνῶναι ἐν τῷ βιβλίῳ τοὺς λόγους Κυρίου ἐν οἴκῳ Κυρίου.
\vs{9}Καὶ ἐγενήθη ἐν τῷ ἔτει τῷ ὀγδόῳ τῷ βασιλεῖ Ἰωακεὶμ ἐν τῷ μηνὶ τῷ ἐννάτῳ, ἐξεκκλησίασαν νηστείαν κατὰ πρόσωπον Κυρίου πᾶς ὁ λαὸς ἐν Ἱερουσαλὴμ, καὶ οἴκος Ἰούδα.
\vs{10}Καὶ ἀνεγίνωσκε Βαροὺχ ἐν τῷ βιβλίῳ τοὺς λόγους Ἱερεμίου ἐν οἴκῳ Κυρίου, ἐν οἴκῳ Γαμαρίου υἱοῦ Σαφὰν τοῦ γραμματέως, ἐν τῇ αὐλῇ τῇ ἐπάνω ἐν προθύροις πύλης τοῦ οἴκου Κυρίου τῆς καινῆς, καὶ ἐν ὠσὶ παντὸς τοῦ λαοῦ.

\vs{11}Καὶ ἤκουσε Μιχαίας υἱὸς Γαμαρίου υἱοῦ Σαφὰν ἅπαντας τοὺς λόγους Κυρίου, ἐκ τοῦ βιβλίου·
\vs{12}Καὶ κατέβη εἰς οἰκίαν τοῦ βασιλέως, εἰς τὸν οἶκον τοῦ γραμματέως, καὶ ἰδοὺ ἐκεῖ πάντες οἱ ἄρχοντες ἐκάθηντο, Ἐλισαμὰ ὁ γραμματεὺς, καὶ Δαλαίας υἱὸς Σελεμίου, καὶ Ἰωνάθαν υἱὸς Ἀκχοβὼρ, καὶ Γαμαρίας υἱὸς Σαφὰν, καὶ Σεδεκίας υἱὸς Ἀνανίου, καὶ πάντες οἱ ἄρχοντες·
\vs{13}Καὶ ἀνήγγειλεν αὐτοῖς Μιχαίας πάντας τοὺς λόγους οὓς ἤκουσεν ἀναγινώσκοντος Βαροὺχ εἰς τὰ ὦτα τοῦ λαοῦ.

\vs{14}Καὶ ἀπέστειλαν πάντες οἱ ἄρχοντες πρὸς Βαροὺχ υἱὸν Νηρίου, τὸν Ἰουδὶν υἱὸν Ναθανίου, υἱοῦ Σελεμίου, υἱοῦ Χουσὶ, λέγοντες, τὸ χαρτίον ἐν ᾧ σὺ ἀναγινώσκεις ἐν αὐτῷ ἐν ὠσὶ τοῦ λαοῦ, λάβε αὐτὸ εἰς τὴν χεῖρά σου, καὶ ἧκε· καὶ ἔλαβε Βαροὺχ τὸ χαρτίον, καὶ κατέβη πρὸς αὐτούς.
\vs{15}Καὶ εἶπον αὐτῷ, πάλιν ἀνάγνωθι εἰς τὰ ὦτα ἡμῶν· καὶ ἀνέγνω Βαρούχ.
\vs{16}Καὶ ἐγενήθη ὡς ἤκουσαν πάντας τοὺς λόγους, συνεβουλεύσαντο ἕκαστος πρὸς τὸν πλησίον αὐτοῦ, καὶ εἶπον, ἀναγγέλλοντες ἀναγγείλωμεν τῷ βασιλεῖ ἅπαντας τοὺς λόγους τούτους.
\vs{17}Καὶ τὸν Βαροὺχ ἠρώτησαν, λέγοντες, ποῦ ἔγραψας πάντας τοὺς λόγους τούτους;
\vs{18}Καὶ εἶπε Βαροὺχ, ἀπὸ στόματος αὐτοῦ ἀνήγγειλέ μοι Ἱερεμίας πάντας τοὺς λόγους τούτους, καὶ ἔγραφον ἐν βιβλίῳ.
\vs{19}Καὶ εἶπον τῷ Βαροὺχ, βάδισον, καὶ κατακρύβηθι σὺ καὶ Ἱερεμίας, ἄνθρωπος μὴ γνώτω ποῦ ὑμεῖς.

\vs{20}Καὶ εἰσῆλθον πρὸς τὸν βασιλέα εἰς τὴν αὐλὴν, καὶ τὸ χαρτίον ἔδωκαν φυλάσσειν ἐν οἴκῳ Ἐλισαμά· καὶ ἀνήγγειλαν τῷ βασιλεῖ πάντας τοὺς λόγους τούτους.
\vs{21}Καὶ ἀπέστειλεν ὁ βασιλεὺς τὸν Ἰουδὶν, λαβεῖν τὸ χαρτίον· καὶ ἔλαβεν αὐτὸ ἐξ οἴκου Ἐλισαμά· καὶ ἀνέγνω Ἰουδὶν εἰς τὰ ὦτα τοῦ βασιλέως, καὶ εἰς τὰ ὦτα πάντων τῶν ἀρχόντων, τῶν ἑστηκότων περὶ τὸν βασιλέα.
\vs{22}Καὶ ὁ βασιλεὺς ἐκάθητο ἐν οἴκῳ χειμερινῷ, καὶ ἐσχάρα πυρὸς κατὰ πρόσωπον αὐτοῦ.
\vs{23}Καὶ ἐγενήθη ἀναγινώσκοντος Ἰουδὶν τρεῖς σελίδας καὶ τέσσαρας, ἀπέτεμεν αὐτὰς τῷ ξυρῷ τοῦ γραμματέως, καὶ ἔῤῥιπτεν εἰς τὸ πῦρ τὸ ἐπὶ τῆς ἐσχάρας, ἕως ἐξέλιπε πᾶς ὁ χάρτης εἰς τὸ πῦρ τὸ ἐπὶ τῆς ἐσχάρας.
\vs{24}Καὶ οὐκ ἐζήτησαν, καὶ οὐ διέῤῥηξαν τὰ ἱμάτια αὐτῶν ὁ βασιλεὺς καὶ οἱ παῖδες αὐτοῦ οἱ ἀκούοντες πάντας τοὺς λόγους τούτους.
\vs{25}Καὶ Ἐλνάθαν καὶ Γοδολίας ὑπέθεντο τῷ βασιλεῖ, πρὸς τὸ κατακαῦσαι τὸ χαρτίον.

\vs{26}Καὶ ἐνετείλατο ὁ βασιλεὺς τῷ Ἱερεμεὴλ υἱῷ τοῦ βασιλέως, καὶ τῷ Σαραίᾳ υἱῷ Ἐσριὴλ, συλλαβεῖν τὸν Βαροὺχ, καὶ τὸν Ἱερεμίαν, καὶ κατεκρύβησαν.

\vs{27}Καὶ ἐγένετο λόγος Κυρίου πρὸς Ἱερεμίαν, μετὰ τὸ κατακαῦσαι τὸν βασιλέα τὸ χαρτίον, πάντας τοὺς λόγους, οὓς ἔγραψε Βαροὺχ ἀπὸ στόματος Ἱερεμίου, λέγων,
\vs{28}πάλιν λάβε σὺ χαρτίον ἕτερον, καὶ γράψον πάντας τοὺς λόγους, τοὺς ὄντας ἐπὶ τοῦ χαρτίου, οὓς κατέκαυσεν ὁ βασιλεὺς Ἰωακεὶμ,
\vs{29}καὶ ἐρεῖς, οὕτως εἶπε Κύριος, σὺ κατέκαυσας τὸ χαρτίον τοῦτο, λέγων, διατί ἔγραψας ἐπʼ αὐτῷ, λέγων, εἰσπορευόμενος εἰσπορεύσεται βασιλεὺς Βαβυλῶνος, καὶ ἐξολοθρεύσει τὴν γῆν ταύτην, καὶ ἐκλείψει ἐπʼ αὐτῆς ἄνθρωπος καὶ κτήνη;

\vs{30}Διατοῦτο οὕτως εἶπε Κύριος ἐπὶ Ἰωακεὶμ βασιλέα Ἰούδα, οὐκ ἔσται αὐτῷ καθήμενος ἐπὶ θρόνου Δαυὶδ, καὶ τὸ θνησιμαῖον αὐτοῦ ἔσται ἐῤῥιμμένον ἐν τῷ καύματι τῆς ἡμέρας, καὶ ἐν τῷ παγετῷ τῆς νυκτός·
\vs{31}Καὶ ἐπισκέψομαι ἐπʼ αὐτὸν, καὶ ἐπὶ τὸ γένος αὐτοῦ, καὶ ἐπὶ τοὺς παῖδας αὐτοῦ, καὶ ἐπάξω ἐπʼ αὐτὸν, καὶ ἐπὶ τοὺς κατοικοῦντας Ἱερουσαλὴμ, καὶ ἐπὶ γῆν Ἰούδα, πάντα τὰ κακὰ ἃ ἐλάλησα πρὸς αὐτοὺς, καὶ οὐκ ἤκουσαν.

\vs{32}Καὶ ἔλαβε Βαροὺχ χαρτίον ἕτερον, καὶ ἔγραψεν ἐπʼ αὐτῷ ἀπὸ στόματος Ἱερεμίου ἅπαντας τοὺς λόγους τοῦ βιβλίου, οὓς κατέκαυσεν Ἰωακείμ· καὶ ἔτι προσετέθησαν αὐτῷ λόγοι πλείονες, ὡς οὗτοι.

\ch{44}
Καὶ ἐβασίλευσε Σεδεκίας υἱὸς Ἰωσεία ἀντὶ Ἰωακείμ, ὃν ἐβασίλευσε Ναβουχοδονόσορ βασιλεύειν τοῦ Ἰούδα.
\vs{2}Καὶ οὐκ ἤκουσαν αὐτὸς καὶ οἱ παῖδες αὐτοῦ καὶ ὁ λαὸς τῆς γῆς τοῦς λόγους Κυρίου, οὓς ἐλάλησεν ἐν χειρὶ Ἱερεμίου.

\vs{3}Καὶ ἀπέστειλεν ὁ βασιλεὺς Σεδεκίας τὸν Ἰωάχαλ υἱὸν Σελεμίου, καὶ τὸν Σοφονίαν υἱὸν Μαασαίου τὸν ἱερέα πρὸς Ἱερεμίαν λέγων, πρόσευξαι δὴ περὶ ἡμῶν πρὸς Κύριον.
\vs{4}Καὶ Ἱερεμίας ἦλθε καὶ διῆλθε διὰ μέσου τῆς πόλεως, καὶ οὐκ ἔδωκαν αὐτὸν εἰς τὸν οἶκον τῆς φυλακῆς.
\vs{5}Καὶ δύναμις Φαραὼ ἐξῆλθεν ἐξ Αἰγύπτου, καὶ ἤκουσαν οἱ Χαλδαῖοι τὴν ἀκοὴν αὐτῶν, καὶ ἀνέβησαν ἐπὶ Ἰερουσαλήμ.

\vs{6}Καὶ ἐγένετο λόγος Κυρίου πρὸς Ἱερεμίαν, λέγων,
\vs{7}οὕτως εἶπε Κύριος, οὕτως ἐρεῖς πρὸς βασιλέα Ἰούδα τὸν ἀποστείλαντα πρὸς σὲ, τοῦ ἐκζητῆσαί με, ἰδοὺ δύναμις· Φαραὼ ἡ ἐξελθοῦσα ὑμῖν εἰς βοήθειαν· ἀποστρέψουσιν εἰς γῆν Αἰγύπτου,
\vs{8}καὶ ἀναστρέψουσιν αὐτοὶ οἱ Χαλδαῖοι, καὶ πολεμήσουσιν ἐπὶ τὴν πόλιν ταύτην, καὶ συλλήψονται αὐτὴν, καὶ καύσουσιν αὐτὴν ἐν πυρί.
\vs{9}Ὅτι οὕτως εἶπε Κύριος, μὴ ὑπολάβητε ταῖς ψυχαῖς ὑμῶν, λέγοντες, ἀποτρέχοντες ἀπελεύσονται ἀφʼ ἡμῶν οἱ Χαλδαῖοι· ὅτι οὐ μὴ ἀπέλθωσι.
\vs{10}Καὶ ἐὰν πατάξητε πᾶσαν δύναμιν τῶν Χαλδαίων τοὺς πολεμοῦντας ὑμᾶς, καὶ καταλειφθῶσί τινες ἐκκεκεντημένοι, ἕκαστος ἐν τῷ τόπῳ αὐτοῦ οὗτοι ἀναστήσονται, καὶ καύσουσι τὴν πόλιν ταύτην ἐν πυρί.

\vs{11}Καὶ ἐγένετο ὅτε ἀνέβη ἡ δύναμις τῶν Χαλδαίων ἀπὸ Ἱερουσαλὴμ ἀπὸ προσώπου τῆς δυνάμεως Φαραὼ,
\vs{12}ἐξῆλθεν Ἱερεμίας ἀπὸ Ἱερουσαλὴμ τοῦ πορευθῆναι εἰς γῆν Βενιαμεὶν, τοῦ ἀγοράσαι ἐκεῖθεν ἐν μέσῳ τοῦ λαοῦ·
\vs{13}Καὶ ἐγένετο αὐτὸς ἐν πύλῃ Βενιαμὶν, καὶ ἐκεῖ ἄνθρωπος παρʼ ᾧ κατέλυε, Σαρουΐα υἱὸς Σελεμίου, υἱοῦ Ἀνανίου, καὶ συνέλαβε τὸν Ἱερεμίαν, λέγων, πρὸς τοὺς Χαλδαίους σὺ φεύγεις.
\vs{14}Καὶ εἶπε, ψεῦδος, οὐκ εἰς τοὺς Χαλδαίους ἐγὼ φεύγω· καὶ οὐκ εἰσήκουσεν αὐτοῦ· καὶ συνέλαβε Σαρουΐα τὸν Ἱερεμίαν, καὶ εἰσήγαγεν αὐτὸν πρὸς τοὺς ἄρχοντας.
\vs{15}Καὶ ἐπικράνθησαν οἱ ἄρχοντες ἐπὶ Ἱερεμίαν, καὶ ἐπάταξαν αὐτὸν, καὶ ἀπέστειλαν αὐτὸν εἰς τὴν οἰκίαν Ἰωνάθαν τοῦ γραμματέως, ὅτι ταύτην ἐποίησαν εἰς οἰκίαν φυλακῆς.

\vs{16}Καὶ ἦλθεν Ἱερεμίας εἰς οἰκίαν τοῦ λάκκου, καὶ εἰς τὴν χερὲθ, καὶ ἐκάθισεν ἐκεῖ ἡμέρας πολλάς.
\vs{17}Καὶ ἀπέστειλε Σεδεκίας, καὶ ἐκάλεσεν αὐτὸν, καὶ ἠρώτα αὐτὸν ὁ βασιλεὺς κρυφαίως εἰπεῖν, εἰ ἔστιν ὁ λόγος παρὰ Κυρίου; καὶ εἶπεν, ἔστιν· εἰς χεῖρας βασιλέως Βαβυλῶνος παραδοθήσῃ.
\vs{18}Καὶ εἶπεν Ἱερεμίας τῷ βασιλεῖ, τί ἠδίκησά σε, καὶ τοὺς παῖδάς σου, καὶ τὸν λαὸν τοῦτον, ὅτι σὺ δίδως με εἰς οἰκίαν φυλακῆς;
\vs{19}Καὶ ποῦ εἰσιν οἱ προφῆται ὑμῶν οἱ προφητεύσαντες ὑμῖν, λέγοντες, ὅτι οὐ μὴ ἔλθῃ βασιλεὺς Βαβυλῶνος ἐπὶ τὴν γῆν ταύτην;
\vs{20}Καὶ νῦν, Κύριε βασιλεῦ, πεσέτω τὸ ἔλεός μου κατὰ πρόσωπόν σου· καὶ τί ἀποστρέφεις με εἰς οἰκίαν Ἰωνάθαν τοῦ γραμματέως; καὶ οὐ μὴ ἀποθάνω ἐκεῖ.
\vs{21}Καὶ συνέταξεν ὁ βασιλεὺς, καὶ ἐνεβάλοσαν αὐτὸν εἰς οἰκίαν τῆς φυλακῆς, καὶ ἐδίδοσαν αὐτῷ ἄρτον ἕνα τῆς ἡμέρας, ἔξωθεν οὗ πέσσουσιν, ἕως ἐξέλιπον οἱ ἄρτοι ἐκ τῆς πόλεως· καὶ ἐκάθισεν Ἱερεμίας ἐν τῇ αὐλῇ τῆς φυλακῆς.

\ch{45}
Καὶ ἤκουσε Σαφανίας υἱὸς Νάθαν, καὶ Γοδολίας υἱὸς Πασχὼρ, καὶ Ἰωάχαλ υἱὸς Σεμελίου, τοὺς λόγους οὓς Ἱερεμίας ἐλάλει ἐπὶ τὸν λαὸν, λέγων,

\vs{2}Οὕτως εἶπε Κύριος, ὁ κατοικῶν ἐν τῇ πόλει ταύτῃ, ἀποθανεῖται ἐν ῥομφαίᾳ, καὶ ἐν λιμῷ· καὶ ὁ ἐκπορευόμενος πρὸς τοὺς Χαλδαίους, ζήσεται, καὶ ἔσται ἡ ψυχὴ αὐτοῦ εἰς εὕρημα, καὶ ζήσεται.
\vs{3}Ὅτι οὕτως εἶπε Κύριος, παραδιδομένη παραδοθήσεται ἡ πόλις αὕτη εἰς χεῖρας δυνάμεως βασιλέως Βαβυλῶνος, καὶ συλλήψεται αὐτήν.
\vs{4}Καὶ εἶπον τῷ βασιλεῖ, ἀναιρεθήτω δὴ ὁ ἄνθρωπος ἐκεῖνος, ὅτι αὐτὸς ἐκλύει τὰς χεῖρας τῶν ἀνθρώπων τῶν πολεμούντων τῶν καταλειπομένων ἐν τῇ πόλει, καὶ τὰς χεῖρας παντὸς τοῦ λαοῦ, λαλῶν πρὸς αὐτοὺς κατὰ τοὺς λόγους τούτους· ὅτι ὁ ἄνθρωπος οὗτος οὐ χρησμολογεῖ εἰρήνην τῷ λαῷ τούτῳ ἀλλʼ ἢ πονηρά.
\vs{5}Καὶ εἶπεν ὁ βασιλεὺς, ἰδοὺ αὐτὸς ἐν χερσὶν ὑμῶν· ὅτι οὐκ ἠδύνατο ὁ βασιλεὺς πρὸς αὐτούς.
\vs{6}Καὶ ἔῤῥιψαν αὐτὸν εἰς λάκκον Μελχίου υἱοῦ τοῦ βασιλέως, ὃς ἦν ἐν τῇ αὐλῇ τῆς φυλακῆς, καὶ ἐχάλασαν αὐτὸν εἰς τὸν λάκκον, καὶ ἐν τῷ λάκκῳ οὐκ ἦν ὕδωρ, ἀλλʼ ἢ βόρβορος, καὶ ἦν ἐν τῷ βορβόρῳ.

\vs{7}Καὶ ἤκουσεν Ἀβδεμέλεχ ὁ Αἰθίοψ, καὶ αὐτὸς ἐν οἰκίᾳ τοῦ βασιλέως, ὅτι ἔδωκαν Ἱερεμίαν εἰς τὸν λάκκον· καὶ ὁ βασιλεὺς ἦν ἐν τῇ πύλῃ Βενιαμὶν,
\vs{8}καὶ ἐξῆλθε πρὸς αὐτὸν, καὶ ἐλάλησε πρὸς τὸν βασιλέα, καὶ εἶπεν,
\vs{9}ἐπονηρεύσω, ἃ ἐποίησας τοῦ ἀποκτεῖναι τὸν ἄνθρωπον τοῦτον, ἀπὸ προσώπου τοῦ λιμοῦ, ὅτι οὐκ εἰσὶν ἔτι ἄρτοι ἐν τῇ πόλει.
\vs{10}Καὶ ἐνετείλατο ὁ βασιλεὺς τῷ Ἀβδεμέλεχ, λέγων, λάβε εἰς τὰς χεῖράς σου ἐντεῦθεν τριάκοντα ἀνθρώπους, καὶ ἀνάγαγε αὐτὸν ἐκ τοῦ λάκκου, ἵνα μὴ ἀποθάνῃ.
\vs{11}Καὶ ἔλαβεν Ἀβδεμέλεχ τοὺς ἀνθρώπους, καὶ εἰσῆλθεν εἰς τὴν οἰκίαν τοῦ βασιλέως τὴν ὑπόγαιον, καὶ ἔλαβεν ἐκεῖθεν παλαιὰ ῥάκη καὶ παλαιὰ σχοινία, καὶ ἔῤῥιψεν αὐτὰ πρὸς Ἱερεμίαν εἰς τὸν λάκκον,
\vs{12}καὶ εἶπε, ταῦτα θὲς ὑποκάτω τῶν σχοινίων· καὶ ἐποίησεν Ἱερεμίας οὕτως.
\vs{13}Καὶ εἵλκυσαν αὐτὸν τοῖς σχοινίοις, καὶ ἀνήγαγον αὐτὸν ἐκ τοῦ λάκκου· καὶ ἐκάθισεν Ἱερεμίας ἐν τῇ αὐλῇ τῆς φυλακῆς.

\vs{14}Καὶ ἀπέστειλεν ὁ βασιλεὺς, καὶ ἐκάλεσεν αὐτὸν προς ἑαυτὸν εἰς οἰκίαν Ἀσελεισὴλ, τὴν ἐν οἴκῳ Κυρίου· καὶ εἶπεν αὐτῷ ὁ βασιλεὺς, ἐρωτήσω σε λόγον, καὶ μὴ δὴ κρύψῃς ἀπʼ ἐμοῦ ῥῆμα.

\vs{15}Καὶ εἶπεν Ἱερεμίας τῷ βασιλεῖ, ἐὰν ἀναγγείλω σοι, οὐχὶ θανάτῳ με θανατώσεις; καὶ ἐὰν συμβουλεύσω σοι, οὐ μὴ ἀκούσῃς μου.
\vs{16}Καὶ ὤμοσεν αὐτῷ ὁ βασιλεὺς, λέγων, ζῇ Κύριος ὃς ἐποίησεν ἡμῖν τὴν ψυχὴν ταύτην, εἰ ἀποκτενῶ σε, καὶ εἰ δώσω σε εἰς χεῖρας τῶν ἀνθρώπων τούτων.

\vs{17}Καὶ εἶπεν αὐτῷ Ἱερεμίας, οὕτως εἶπε Κύριος, ἐὰν ἐξελθὼν ἐξέλθῃς πρὸς ἡγεμόνας βασιλέως Βαβυλῶνος, ζήσεται ἡ ψυχή σου, καὶ ἡ πόλις αὕτη οὐ μὴ κατακαυθῇ ἐν πυρὶ, καὶ ζήσῃ σὺ καὶ ἡ οἰκία σου.
\vs{18}Καὶ ἐὰν μὴ ἐξέλθῃς, δοθήσεται ἡ πόλις αὕτη εἰς χεῖρας τῶν Χαλδαίων, καὶ καύσουσιν αὐτὴν ἐν πυρὶ καὶ σὺ οὐ μὴ σωθῇς.

\vs{19}Καὶ εἶπεν ὁ βασιλεὺς τῷ Ἱερεμίᾳ, ἐγὼ λόγον ἔχω τῶν Ἰουδαίων τῶν πεφευγότων πρὸς τοὺς Χαλδαίους, μὴ δώσειν με εἰς χεῖρας αὐτῶν, καὶ καταμωκήσονταί μου.

\vs{20}Καὶ εἶπεν Ἱερεμίας, οὐ μὴ παραδῶσί σε· ἄκουσον τὸν λόγον Κυρίου, ὃν ἐγὼ λέγω πρὸς σὲ, καὶ βέλτιον ἔσται σοι, καὶ ζήσεται ἡ ψυχή σου.
\vs{21}Καὶ εἰ μὴ θέλῃς σὺ ἐξελθεῖν, οὗτος ὁ λόγος ὃν ἔδειξέ μοι Κύριος·
\vs{22}Καὶ ἰδοὺ πᾶσαι αἱ γυναῖκες αἱ καταλειφθεῖσαι ἐν οἰκίᾳ βασιλέως Ἰούδα, ἐξήγοντο πρὸς ἄρχοντας βασιλέως Βαβυλῶνος· καὶ αὗται ἔλεγον, ἠπάτησάν σε, καὶ δυνήσονταί σοι ἄνδρες εἰρηνικοί σου· καὶ καταλύσουσιν ἐν ὀλισθήμασι πόδα σου, ἀπέστρεψαν ἀπὸ σοῦ,
\vs{23}καὶ τᾶς γυναῖκάς σου καὶ τὰ τέκνα σου ἐξάξουσι πρὸς τοὺς Χαλδαίους· καὶ σὺ οὐ μὴ σωθῇς, ὅτι ἐν χειρὶ βασιλέως Βαβυλῶνος συλληφθήσῃ, καὶ ἡ πόλις αὕτη κατακαυθήσεται.

\vs{24}Καὶ εἶπεν αὐτῷ ὁ βασιλεὺς, ἄνθρωπος μὴ γνώτω ἐκ τῶν λόγων τούτων, καὶ σὺ οὐ μὴ ἀποθάνῃς.
\vs{25}Καὶ ἐὰν οἱ ἄρχοντες ἀκούσωσιν ὅτι ἐλάλησά σοι, καὶ ἔλθωσι πρὸς σὲ, καὶ εἴπωσί σοι, ἀνάγγειλον ἡμῖν, τί ἐλάλησέ σοι ὁ βασιλεύς; μὴ κρύψῃς ἀφʼ ἡμῶν, καὶ οὐ μὴ ἀνέλωμέν σε· καὶ τί ἐλάλησε πρὸς σὲ ὁ βασιλεύς;
\vs{26}Καὶ ἐρεῖς αὐτοῖς, ῥίπτω ἐγὼ τὸ ἔλεός μου κατʼ ὀφθαλμοὺς τοῦ βασιλέως, πρὸς τὸ μὴ ἀποστρέψαι με εἰς οἰκίαν Ἰωνάθαν, ἀποθανεῖν με ἐκεῖ.

\vs{27}Καὶ ἤλθοσαν πάντες οἱ ἄρχοντες πρὸς Ἰερεμίαν, καὶ ἠρώτησαν αὐτόν· καὶ ἀνήγγειλεν αὐτοῖς κατὰ πάντας τοὺς λόγους τούτους, οὓς ἐνετείλατο αὐτῷ ὁ βασιλεύς· καὶ ἀπεσιώπησαν, ὅτι οὐκ ἠκούσθη ὁ λόγος Κυρίου.
\vs{28}Καὶ ἐκάθισεν Ἱερεμίας ἐν τῇ αὐλῇ τῆς φυλακῆς, ἕως χρόνου οὗ συνελήφθη Ἱερουσαλήμ.

\ch{46}
Καὶ ἐγένετο τῷ μηνὶ τῷ ἐννάτῳ τοῦ Σεδεκία βασιλέως Ἰούδα, παρεγένετο Ναβουχοδονόσορ βασιλεὺς Βαβυλῶνος, καὶ πᾶσα ἡ δύναμις αὐτοῦ ἐπὶ Ἱερουσαλὴμ, καὶ ἐπολιόρκουν αὐτήν.
\vs{2}Καὶ ἐν τῷ ἑνδεκάτῳ ἔτει τοῦ Σεδεκία, ἐν τῷ μηνὶ τῷ τετάρτῳ, ἐννάτῃ τοῦ μηνὸς, ἐῤῥάγη ἡ πόλις,
\vs{3}καὶ εἰσῆλθον πάντες οἱ ἡγούμενοι βασιλέως Βαβυλῶνος, καὶ ἐκάθισαν ἐν πύλῃ τῇ μέσῃ Μαργανασὰρ, καὶ Σαμαγὼθ, καὶ Ναβουσάχαρ, καὶ Ναβουσαρεὶς, Ναγαργᾶς, Νασεῤῥαβαμὰθ, καὶ οἱ κατάλοιποι ἡγεμόνες βασιλέως Βαβυλῶνος.
\vs{14}Καὶ ἀπέστειλαν, καὶ ἔλαβον τὸν Ἱερεμίαν ἐξ αὐλῆς τῆς φυλακῆς, καὶ ἔδωκαν αὐτὸν πρὸς τὸν Γοδολίαν υἱὸν Ἀχεικὰμ, υἱοῦ Σαφὰν, καὶ ἐξήγαγον αὐτὸν, καὶ ἐκάθισεν ἐν μέσῳ τοῦ λαοῦ.

\vs{15}Καὶ πρὸς Ἱερεμίαν ἐγένετο λόγος Κυρίου ἐν τῇ αὐλῇ τῆς φυλακῆς, λέγων,
\vs{16}πορεύου καὶ εἰπὲ πρὸς Ἀβδεμέλεχ τὸν Αἰθίοπα, οὕτως εἶπε Κύριος ὁ Θεὸς Ἰσραὴλ, ἰδοὺ ἐγὼ φέρω τοὺς λόγους μου ἐπὶ τὴν πόλιν ταύτην εἰς κακὰ καὶ οὐκ εἰς ἀγαθά.
\vs{17}Καὶ σώσω σε ἐν τῇ ἡμέρᾳ ἐκείνῃ, καὶ οὐ μὴ δώσω σε εἰς χεῖρας τῶν ἀνθρώπων ὧν σὺ φοβῇ ἀπὸ προσώπου αὐτῶν,
\vs{18}ὅτι σώζων σώσω σε, καὶ ἐν ῥομφαίᾳ οὐ μὴ πέσῃς· καὶ ἔσται ἡ ψυχή σου εἰς εὕρημα, ὅτι ἐπεποίθεις ἐπʼ ἐμοὶ, φησὶ Κύριος.

\ch{47}
Ὁ λόγος ὁ γενόμενος παρὰ Κυρίου πρὸς Ἱερεμίαν, μετὰ τὸ ἀποστεῖλαι αὐτὸν Ναβουζαρδὰν τὸν ἀρχιμάγειρον τὸν ἐκ Ῥαμὰ, ἐν τῷ λαβεῖν αὐτὸν ἐν χειροπέδαις, ἐν μέσῳ ἀποικίας Ἰούδα τῶν ἠγμένων εἰς Βαβυλῶνα.

\vs{2}Καὶ ἔλαβεν αὐτὸν ὁ ἀρχιμάγειρος, καὶ εἶπεν αὐτῷ, Κύριος ὁ Θεός σου ἐλάλησε τὰ κακὰ ταῦτα ἐπὶ τὸν τόπον τοῦτον·
\vs{3}Καὶ ἐποίησε Κύριος, ὅτι ἡμάρτετε αὐτῷ, καὶ οὐκ ἠκούσατε τῆς φωνῆς αὐτοῦ.
\vs{4}Ἰδοὺ ἔλυσά σε ἀπὸ τῶν χειροπέδων τῶν ἐπὶ τὰς χεῖράς σου· εἰ καλὸν ἐναντίον σου ἐλθεῖν μετʼ ἐμοῦ εἰς Βαβυλῶνα, καὶ θήσω τοὺς ὀφθαλμούς μου ἐπὶ σέ.
\vs{5}Εἰ δὲ μὴ, ἀπότρεχε, ἀνάστρεψον πρὸς τὸν Γοδολίαν υἱὸν Ἀχεικὰμ, υἱοῦ Σαφὰν, ὃν κατέστησε βασιλεὺς Βαβυλῶνος ἐν γῇ Ἰούδα, καὶ οἴκησον μετʼ αὐτοῦ ἐν μέσῳ τοῦ λαοῦ ἐν γῇ Ἰούδα, εἰς ἅπαντα τὰ ἀγαθὰ ἐν ὀφθαλμοῖς σου τοῦ πορευθῆναι ἐκεῖ, καὶ πορεύου· καὶ ἔδωκεν αὐτῷ ὁ ἀρχιμάγειρος δῶρα, καὶ ἀπέστειλεν αὐτόν.
\vs{6}Καὶ ἦλθε πρὸς Γοδολίαν εἰς Μασσηφὰ, καὶ ἐκάθισεν ἐν μέσῳ τοῦ λαοῦ αὐτοῦ, τοῦ καταλειφθέντος ἐν τῇ γῇ.

\vs{7}Καὶ ἤκουσαν πάντες οἱ ἡγεμόνες τῆς δυνάμεως τῆς ἐν ἀγρῷ, αὐτοὶ καὶ οἱ ἄνδρες αὐτῶν, ὅτι κατέστησε βασιλεὺς Βαβυλῶνος τὸν Γοδολίαν ἐν τῇ γῇ, καὶ παρακατέθεντο αὐτῷ ἄνδρας καὶ γυναῖκας αὐτῶν, οὓς οὐ κατῴκισεν εἰς Βαβυλῶνα.
\vs{8}Καὶ ἦλθε πρὸς Γοδολίαν εἰς Μασσηφὰ Ἰσραὴλ υἱὸς Ναθανίου, καὶ Ἰωάναν υἱὸς Κάρηε, καὶ Σαραίας υἱὸς Θαναεμὲθ, καὶ υἱοὶ Ἰωφὲ τοῦ Νετωφαθὶ, καὶ Ἐζονίας υἱὸς τοῦ Μωχαθὶ, αὐτοὶ, καὶ οἱ ἄνδρες αὐτῶν.

\vs{9}Καὶ ὤμοσεν αὐτοῖς Γοδολίας, καὶ τοῖς ἀνδράσιν αὐτῶν, λέγων, μὴ φοβηθῆτε ἀπὸ προσώπου τῶν παίδων τῶν Χαλδαίων· κατοικήσατε ἐν τῇ γῇ, καὶ ἐργάσασθε τῷ βασιλεῖ Βαβυλῶνος, καὶ βέλτιον ἔσται ὑμῖν.
\vs{10}Καὶ ἰδοὺ ἐγὼ κάθημαι ἐναντίον ὑμῶν εἰς Μασσηφὰ, στῆναι κατὰ πρόσωπον τῶν Χαλδαίων, οἳ ἂν ἔλθωσιν ἐφʼ ὑμᾶς· καὶ ὑμεῖς συνάγετε οἶνον καὶ ὀπώραν καὶ ἔλαιον, καὶ βάλετε εἰς τὰ ἀγγεῖα ὑμῶν, καὶ οἰκήσατε ἐν ταῖς πόλεσιν αἷς κατεκρατήσατε.

\vs{11}Καὶ πάντες οἱ Ἰουδαῖοι οἱ ἐν Μωὰβ, καὶ ἐν υἱοῖς Ἀμμὼν, καὶ οἱ ἐν τῇ Ἰδουμαίᾳ, καὶ οἱ ἐν πάσῃ τῇ γῇ, ἤκουσαν ὅτι ἔδωκε βασιλεὺς Βαβυλῶνος κατάλειμμα τῷ Ἰούδᾳ, καὶ ὅτι κατέστησεν ἐπʼ αὐτοὺς τὸν Γοδολίαν υἱὸν Ἀχεικάμ.
\vs{12}Καὶ ἦλθον πρὸς Γοδολίαν εἰς γῆν Ἰούδα εἰς Μασσηφὰ, καὶ συνήγαγον οἶνον, καὶ ὀπώραν πολλὴν σφόδρα, καὶ ἔλαιον.

\vs{13}Καὶ Ἰωάναν υἱὸς Κάρηε, καὶ πάντες οἱ ἡγεμόνες τῆς δυνάμεως, οἱ ἐν τοῖς ἀγροῖς, ἦλθον πρὸς τὸν Γοδολίαν εἰς Μασσηφὰ,
\vs{14}καὶ εἶπον αὐτῷ, εἰ γνώσει γινώσκεις, ὅτι Βελεισσὰ βασιλεὺς υἱὸς Ἀμμὼν ἀπέστειλε πρὸς σὲ τὸν Ἰσμαὴλ πατάξαι σου ψυχήν; καὶ οὐκ ἐπίστευσεν αὐτοῖς Γοδολίας.
\vs{15}Καὶ εἶπεν Ἰωάναν τῷ Γοδολίᾳ κρυφαίως ἐν Μασσηφᾷ, πορεύσομαι δὴ καὶ πατάξω τὸν Ἰσμαὴλ, καὶ μηδεὶς γνώτω, μὴ πατάξῃ σου ψυχὴν, καὶ διασπαρῇ πᾶς Ἰούδα οἱ συνηγμένοι πρὸς σὲ, καὶ ἀπολοῦνται οἱ κατάλοιποι Ἰούδα.
\vs{16}Καὶ εἶπε Γοδολίας πρὸς Ἰωάναν, Μῆ ποιήσῃς τὸ πρᾶγμα, ὅτι ψευδῆ σὺ λέγεις ὑπὲρ Ἰσμαήλ.

\ch{48}
Καὶ ἐγένετο τῷ μηνὶ τῷ ἑβδόμῳ, ἦλθεν Ἰσμαὴλ υἱὸς Ναθανίου υἱοῦ Ἐλεασὰ, ἀπὸ γένους τοῦ βασιλέως, καὶ δέκα ἄνδρες μετʼ αὐτοῦ πρὸς Γοδολίαν εἰς Μασσηφὰ, καὶ ἔφαγον ἐκεῖ ἄρτον ἅμα.
\vs{2}Καὶ ἀνέστη Ἰσμαὴλ, καὶ οἱ δέκα ἄνδρες οἳ ἦσαν μετʼ αὐτοῦ, καὶ ἐπάταξαν τὸν Γοδολίαν, ὃν κατέστησε βασιλεὺς Βαβυλῶνος ἐπὶ τῆς γῆς,
\vs{3}καὶ πάντας τοὺς Ἰουδαίους τοὺς ὄντας μετʼ αὐτοῦ ἐν Μασσηφὰ, καὶ πάντας τοὺς Χαλδαίους τοὺς εὑρεθέντας ἐκεῖ.

\vs{4}Καὶ ἐγένετο τῇ ἡμέρᾳ τῇ δευτέρᾳ πατάξαντος αὐτοῦ τὸν Γοδολίαν, καὶ ἄνθρωπος οὐκ ἔγνω.
\vs{5}Καὶ ἤλθοσαν ἄνδρες ἀπὸ Συχὲμ, καὶ ἀπὸ Σαλὴμ, καὶ ἀπὸ Σαμαρίας, ὀγδοήκοντα ἄνδρες, ἐξυρημένοι πώγωνας, καὶ διεῤῥηγμένοι τὰ ἱμάτια, καὶ κοπτόμενοι, καὶ μάννα, καὶ λίβανος ἐν χερσὶν αὐτῶν, τοῦ εἰσενεγκεῖν εἰς οἶκον Κυρίου.
\vs{6}Καὶ ἐξῆλθεν εἰς ἀπάντησιν αὐτοῖς Ἰσμαήλ· αὐτοὶ ἐπορεύοντο, καὶ ἔκλαιον· καὶ εἶπεν αὐτοῖς, εἰσέλθετε πρὸς Γοδολίαν.
\vs{7}Καὶ ἐγένετο, εἰσελθόντων αὐτῶν εἰς τὸ μέσον τῆς πόλεως, ἔσφαξεν αὐτοὺς εἰς τὸ φρέαρ.
\vs{8}Καὶ δέκα ἄνδρες εὑρέθησαν ἐκεῖ, καὶ εἶπον τῷ Ἱσμαὴλ, μὴ ἀνέλῃς ἡμᾶς, ὅτι εἰσὶν ἡμῖν θησαυροὶ ἐν ἀγρῷ, πυροὶ καὶ κριθαὶ, μέλι καὶ ἔλαιον· καὶ παρῆλθε, καὶ οὐκ ἀνεῖλεν αὐτοὺς ἐν μέσῳ τῶν ἀδελφῶν αὐτῶν.

\vs{9}Καὶ τὸ φρέαρ εἰς ὃ ἔῤῥιψεν ἐκεῖ Ἰσμαὴλ πάντας οὓς ἐπάταξε, φρέαρ μέγα τοῦτό ἐστιν, ὃ ἐποίησεν ὁ βασιλεὺς Ἀσὰ ἀπὸ προσώπου Βαασὰ βασιλέως Ἰσραὴλ, τοῦτο ἐνέπλησεν Ἰσμαὴλ τραυματιῶν.

\vs{10}Καὶ ἀπέστρεψεν Ἰσμαὴλ πάντα τὸν λαὸν τὸν καταλειφθέντα εἰς Μασσηφὰ, καὶ τὰς θυγατέρας τοῦ βασιλέως, ἃς παρκατέθετο ὁ ἀρχιμάγειρος τῷ Γοδολίᾳ υἱῷ Ἀχεικὰμ, καὶ ᾤχετο εἰς τὸ πέραν υἱῶν Ἀμμών.

\vs{11}Καὶ ἤκουσεν Ἰωάναν υἱὸς Κάρηε, καὶ πάντες οἱ ἡγεμόνες τῆς δυνάμεως οἱ μετʼ αὐτοῦ, πάντα τὰ κακὰ ἃ ἐποίησεν Ἰσμαὴλ,
\vs{12}καὶ ἤγαγον ἅπαν τὸ στρατόπεδον αὐτῶν, καὶ ᾤχοντο πολεμεῖν αὐτὸν, καὶ εὗρον αὐτὸν ἐπὶ ὕδατος πολλοῦ ἐν Γαβαών.
\vs{13}Καὶ ἐγένετο, ὅτε εἶδε πᾶς ὁ λαὸς ὁ μετὰ Ἰσμαὴλ τὸν Ἰωάναν καὶ τοὺς ἡγεμόνας τῆς δυνάμεως τῆς μετʼ αὐτοῦ,
\vs{14}καὶ ἀνέστρεψαν πρὸς Ἰωάναν.
\vs{15}Καὶ Ἰσμαὴλ ἐσώθη σὺν ὀκτὼ ἀνθρώποις, καὶ ᾤχετο πρὸς τοὺς υἱοὺς Ἀμμών.

\vs{16}Καὶ ἔλαβεν Ἰωάναν, καὶ πάντες οἱ ἡγεμόνες τῆς δυνάμεως οἱ μετʼ αὐτοῦ, πάντας τοὺς καταλοίπους τοῦ λαοῦ, οὓς ἀπέστρέψεν ἀπὸ Ἰσμαὴλ, δυνατοὺς ἄνδρας ἐν πολέμῳ, καὶ τὰς γυναῖκας, καὶ τὰ λοιπὰ, καὶ τοὺς εὐνούχους, οὓς ἀπέστρεψαν ἀπὸ Γαβαὼν,
\vs{17}καὶ ᾤχοντο, καὶ ἐκάθισαν ἐν Γαβηρωχαμάα, τῇ πρὸς Βηθλεὲμ, τοῦ πορευθῆναι εἰς Αἴγυπτον
\vs{18}ἀπὸ προσώπου τῶν Χαλδαίων· ὅτι ἐφοβήθησαν ἀπὸ προσώπου αὐτῶν, ὅτι ἐπάταξεν Ἰσμαὴλ τὸν Γοδολίαν, ὃν κατέστησεν ὁ βασιλεὺς Βαβυλῶνος ἐν τῇ γῇ.

\ch{49}
Καὶ προσῆλθον πάντες οἱ ἡγεμόνες τῆς δυνάμεως, καὶ Ἰωάναν, καὶ Ἀζαρίας υἱὸς Μαασαίου, καὶ πᾶς ὁ λαὸς ἀπὸ μικροῦ καὶ ἕως μεγάλου,
\vs{2}πρὸς Ἱερεμίαν τὸν προφήτην, καὶ εἶπαν αὐτῷ, πεσέτω δὴ τὸ ἔλεος ἡμῶν κατὰ πρόσωπόν σου, καὶ πρόσευξαι πρὸς Κύριον τὸν Θεόν σου περὶ τῶν καταλοίπων τούτων· ὅτι κατελείφθημεν ὀλίγοι ἀπὸ πολλῶν, καθὼς οἱ ὀφθαλμοί σου βλέπουσι.
\vs{3}Καὶ ἀναγγειλάτω ἡμῖν Κύριος ὁ Θεός σου τὴν ὁδὸν ᾗ πορευσόμεθα ἐν αὐτῇ, καὶ λόγον ὃν ποιήσομεν.

\vs{4}Καὶ εἶπεν αὐτοῖς Ἱερεμίας, ἤκουσα, ἰδοὺ ἐγὼ προσεύξομαι ὑπὲρ ὑμῶν πρὸς Κύριον τὸν Θεὸν ἡμῶν, κατὰ τοὺς λόγους ὑμῶν· καὶ ἔσται, ὁ λόγος ὃν ἂν ἀποκριθήσεται Κύριος ὁ Θεὸς, ἁναγγελῶ ὑμῖν, οὐ μὴ κρύψω ἀφʼ ὑμῶν ῥῆμα.

\vs{5}Καὶ αὐτοὶ εἶπαν τῷ Ἱερεμίᾳ, ἔστω Κύριος ἐν ἡμῖν εἰς μάρτυρα δίκαιον καὶ πιστὸν, εἰ μὴ κατὰ πάντα τὸν λόγον, ὃν ἐὰν ἀποστείλῃ Κύριος πρὸς ἡμᾶς, οὕτως ποιήσωμεν.
\vs{6}Καὶ ἐὰν ἀγαθὸν καὶ ἐὰν κακὸν, τὴν φωνὴν Κυρίου τοῦ Θεοῦ ἡμῶν, οὗ ἡμεῖς ἀποστέλλομέν σε πρὸς αὐτὸν, ἀκουσόμεθα, ἵνα βέλτιον ἡμῖν γένηται, ὅτι ἀκουσόμεθα τῆς φωνῆς Κυρίου τοῦ Θεοῦ ἡμῶν.

\vs{7}Καὶ ἐγενήθη μεθʼ ἡμέρας δέκα, ἐγενήθη λόγος Κυρίου πρὸς Ἱερεμίαν.
\vs{8}Καὶ ἐκάλεσε τὸν Ἰωάναν, καὶ τοὺς ἡγεμόνας τῆς δυνάμεως, καὶ πάντα τὸν λαὸν ἀπὸ μικροῦ καὶ ἕως μεγάλου,
\vs{9}καὶ εἶπεν αὐτοῖς, οὕτως εἶπε Κύριος,
\vs{10}ἐὰν καθίσαντες καθίσητε ἐν τῇ γῇ ταύτῃ, οἰκοδομήσω ὑμᾶς, καὶ οὐ μὴ καθελῶ, καὶ φυτεύσω ὑμᾶς, καὶ οὐ μὴ ἐκτιλῶ, ὅτι ἀναπέπαυμαι ἐπὶ τοῖς κακοῖς οἷς ἐποίησα ὑμῖν.
\vs{11}Μὴ φοβηθῆτε ἀπὸ προσώπου βασιλέως Βαβυλῶνος, οὗ ὑμεῖς φοβεῖσθε· ἀπὸ προσώπου αὐτοῦ μὴ φοβηθῆτε, φησὶ Κύριος, ὅτι μεθʼ ὑμῶν ἐγὼ, ἐξαιρεῖσθαι ὑμᾶς, καὶ σώζειν ὑμᾶς ἐκ χειρὸς αὐτῶν.
\vs{12}Καὶ δώσω ὑμῖν ἔλεος, καὶ ἐλεήσω ὑμᾶς, καὶ ἐπιστρέψω ὑμᾶς εἰς τὴν γῆν ὑμῶν.

\vs{13}Καὶ εἰ λέγετε ὑμεῖς, οὐ μὴ καθίσωμεν ἐν τῇ γῇ ταύτῃ, πρὸς τὸ μὴ ἀκοῦσαι φωνῆς Κυρίου,
\vs{14}ὅτι εἰς γῆν Αἰγύπτου εἰσελευσόμεθα, καὶ οὐ μὴ ἴδωμεν πόλεμον, καὶ φωνὴν σάλπιγγος οὐ μὴ ἀκούσωμεν, καὶ ἐν ἄρτοις οὐ μὴ πεινάσωμεν, καὶ ἐκεῖ οἰκήσομεν·
\vs{15}διατοῦτο ἀκούσατε λόγον Κυρίου· οὕτως εἶπε Κύριος, ἐὰν ὑμεῖς δῶτε τὸ πρόσωπον ὑμῶν εἰς Αἴγυπτον, καὶ εἰσέλθητε ἐκεῖ κατοικεῖν,
\vs{16}καὶ ἔσται, ἡ ῥομφαία ἣν ὑμεῖς φοβεῖσθε ἀπὸ προσώπου αὐτῆς εὑρήσει ὑμᾶς ἐν γῇ Αἰγύπτου, καὶ ὁ λιμὸς οὗ ὑμεῖς λόγον ἔχετε ἀπὸ προσώπου αὐτοῦ καταλήψεται ὑμᾶς ὀπίσω ὑμῶν ἐν Αἰγύπτῳ, καὶ ἐκεῖ ἀποθανεῖσθε.
\vs{17}Καὶ ἔσονται πάντες οἱ ἄνθρωποι, καὶ πάντες οἱ ἀλλογενεῖς, οἱ θέντες τὸ πρόσωπον αὐτῶν εἰς γῆν Αἰγύπτου ἐνοικεῖν ἐκεῖ, ἐκλείψουσιν ἐν τῇ ῥομφαίᾳ, καὶ ἐν τῷ λιμῷ, καὶ οὐκ ἔσται αὐτῶν οὐδεὶς σωζόμενος ἀπὸ τῶν κακῶν ὧν ἐγὼ ἐπάγω ἐπʼ αὐτούς.

\vs{18}Ὅτι οὕτως εἶπε Κύριος, καθὼς ἔσταξεν ὁ θυμός μου ἐπὶ τοὺς κατοικοῦντας Ἱερουσαλὴμ, οὕτως στάξει ὁ θυμός μου ἐφʼ ὑμᾶς, εἰσελθόντων ὑμῶν εἰς Αἴγυπτον· καὶ ἔσεσθε εἰς ἄβατον, καὶ ὑποχείριοι, καὶ εἰς ἀρὰν, καὶ εἰς ὀνειδισμὸν, καὶ οὐ μὴ ἴδητε οὐκέτι τὸν τόπον τοῦτον.

\vs{19}Ἃ ἐλάλησε Κύριος ἐφʼ ὑμᾶς τοὺς καταλοίπους Ἰούδα· μὴ εἰσέλθητε εἰς Αἴγυπτον· καὶ νῦν γνόντες γνώσεσθε,
\vs{20}ὅτι ἐπονηρεύσασθε ἐν ψυχαῖς ὑμῶν, ἀποστείλαντές με, λέγοντες, πρόσευξαι περὶ ἡμῶν πρὸς Κύριον, καὶ κατὰ πάντα ἃ ἐὰν λαλήσει σοι Κύριος ποιήσομεν.
\vs{21}Καὶ οὐκ ἠκούσατε τῆς φωνῆς Κυρίου, ἧς ἀπέστειλέ με πρὸς ὑμᾶς.
\vs{22}Καὶ νῦν ἐν ῥομφαίᾳ καὶ ἐν λιμῷ ἐκλείψετε ἐν τῷ τόπῳ ᾧ ὑμεῖς βούλεσθε εἰσελθεῖν κατοικεῖν ἐκεῖ.

\ch{50}
Καὶ ἐγενήθη, ὡς ἐπαύσατο Ἱερεμίας λέγων πρὸς τὸν λαὸν πάντας τοὺς λόγους Κυρίου, οὓς ἀπέστειλεν αὐτὸν Κύριος πρὸς αὐτοὺς, πάντας τοὺς λόγους τούτους,
\vs{2}καὶ εἶπεν Ἀζαρίας υἱὸς Μαασαίου, καὶ Ἰωάναν υἱὸς Κάρηε, καὶ πάντες οἱ ἄνδρες, οἱ εἰπόντες τῷ Ἱερεμίᾳ, λέγοντες, ψεύδη, οὐκ ἀπέστειλέ σε Κύριος πρὸς ἡμᾶς, λέγων, μὴ εἰσέλθητε εἰς Αἴγυπτον οἰκεῖν ἐκεῖ,
\vs{3}ἀλλʼ ἢ Βαροὺχ υἱὸς Νηρίου συμβάλλει σε πρὸς ἡμᾶς, ἵνα δῷς ἡμᾶς εἰς χεῖρας τῶν Χαλδαίων, τοῦ θανατῶσαι ἡμᾶς, καὶ ἀποικισθῆναι ἡμᾶς εἰς Βαβυλῶνα.
\vs{4}Καὶ οὐκ ἤκουσεν Ἰωάναν, καὶ πάντες ἡγεμόνες τῆς δυνάμεως, καὶ πᾶς ὁ λαὸς τῆς φωνῆς Κυρίου, κατοικῆσαι ἐν γῇ Ἰούδα.

\vs{5}Καὶ ἔλαβεν Ἰωάναν, καὶ πάντες οἱ ἡγεμόνες τῆς δυνάμεως πάντας τοὺς καταλοίπους Ἰούδα, τοὺς ἀποστρέψαντας κατοικεῖν ἐν τῇ γῃ,
\vs{6}τοὺς δυνατοὺς ἄνδρας, καὶ τὰς γυναῖκας, καὶ τὰ νήπια τὰ λοιπὰ, καὶ τὰς θυγατέρας τοῦ βασιλέως, καὶ τὰς ψυχὰς ἃς κατέλιπε Ναβουζαρδὰν μετὰ Γοδολίου υἱοῦ Ἀχεικὰμ, καὶ Ἱερεμίαν τὸν προφήτην, καὶ Βαροὺχ, υἱὸν Νηρίου,
\vs{7}καὶ εἰσῆλθον εἰς Αἴγυπτον, ὅτι οὐκ ἤκουσαν τῆς φωνῆς Κυρίου, καὶ εἰσῆλθον εἰς Τάφνας.

\vs{8}Καὶ ἐγένετο λόγος Κυρίου πρὸς Ἱερεμίαν ἐν Τάφνας, λέγων,
\vs{9}λάβε σεαυτῷ λίθους μεγάλους, καὶ κατάκρυψον αὐτοὺς ἐν προθύροις, ἐν πύλῃ τῆς οἰκίας Φαραὼ ἐν Τάφνας, κατʼ ὀφθαλμοὺς ἀνδρῶν Ἰούδα,
\vs{10}καὶ ἐρεῖς, οὕτως εἶπε Κύριος, ἰδοὺ ἐγὼ ἀποστέλλω, καὶ ἄξω Ναβουχοδονόσορ βασιλέα Βαβυλῶνος, καὶ θήσει αὐτοῦ τὸν θρόνον ἐπάνω τῶν λίθων τούτων ὧν κατέκρυψας, καὶ ἀρεῖ τᾶ ὅπλα ἐπʼ αὐτοὺς,
\vs{11}καὶ εἰσελεύσεται, καὶ πατάξει γῆν Αἰγύπτου, οὓς εἰς θάνατον εἰς θάνατον, καὶ οὓς εἰς ἀποικισμὸν εἰς ἀποικισμὸν, καὶ οὓς εἰς ῥομφαίαν εἰς ῥομφαίαν.
\vs{12}Καὶ καύσει πῦρ ἐν οἰκίαις τῶν θεῶν αὐτῶν, καὶ ἐμπυριεῖ αὐτὰς, καὶ ἀποικιεῖ αὐτοὺς, καὶ φθειριεῖ γῆν Αἰγύπτου, ὥσπερ φθειρίζει ποιμὴν τὸ ἱμάτιον αὐτοῦ· καὶ ἐξελεύσεται ἐν εἰρήνῃ,
\vs{13}καὶ συντρίψει τοὺς στύλους Ἡλιουπόλεως τοὺς ἐν Ὢν, καὶ τὰς οἰκίας αὐτῶν κατακαύσει ἐν πυρί.

\ch{51}
Ὁ ΛΟΓΟΣ Ὁ ΓΕΝΟΜΕΝΟΣ ΠΡΟΣ ἹΕΡΕΜΙΑΝ ἅπασι τοῖς Ἰουδαίοις τοῖς κατοικοῦσιν ἐν γῇ Αἰγύπτου, καὶ τοῖς καθημένοις ἐν Μαγδωλῷ, καὶ ἐν Τάφνας, καὶ ἐν γῇ Παθούρης, λέγων,

\vs{2}Οὕτως εἶπε Κύριος ὁ Θεὸς Ἰσραὴλ, ὑμεῖς ἑωράκατε πάντα τὰ κακὰ ἃ ἐπήγαγον ἐπὶ Ἱερουσαλὴμ, καὶ ἐπὶ τὰς πόλεις Ἰούδα· καὶ ἰδού εἰσιν ἔρημοι ἀπὸ ἐνοίκων
\vs{3}ἀπὸ προσώπου πονηρίας αὐτῶν, ἧς ἐποίησαν παραπικράναι με, πορευθέντες θυμιᾷν θεοῖς ἑτέροις, οἷς οὐκ ἔγνωτε.
\vs{4}Καὶ ἀπέστειλα πρὸς ὑμᾶς τοὺς παῖδάς μου τοὺς προφήτας ὄρθρου, καὶ ἀπέστειλα, λέγων, μὴ ποιήσητε τὸ πρᾶγμα τῆς μολύνσεως ταύτης, ἧς ἐμίσησα.

\vs{5}Καὶ οὐκ ἤκουσάν μου, καὶ οὐκ ἔκλιναν τὸ οὖς αὐτῶν ἀποστρέψαι ἀπὸ τῶν κακῶν αὐτῶν, πρὸς τὸ μὴ θυμιᾷν θεοῖς ἑτέροις.
\vs{6}Καὶ ἔσταξεν ἡ ὀργή μου, καὶ ὁ θυμός μου, καὶ ἐξεκαύθη ἐν πύλαις Ἰούδα, καὶ ἔξωθεν Ἱερουσαλήμ· καὶ ἐγενήθησαν εἰς ἐρήμωσιν καὶ εἰς ἄβατον ὡς ἡ ἡμέρα αὕτη.

\vs{7}Καὶ νῦν οὕτως εἶπε Κύριος παντοκράτωρ, ἱνατί ὑμεῖς ποιεῖτε κακὰ μεγάλα ἐπὶ ψυχαῖς ὑμῶν; ἐκκόψαι ὑμῶν ἄνθρωπον καὶ γυναῖκα, νήπιον καὶ θηλάζοντα ἐκ μέσου Ἰούδα, πρὸς τὸ μὴ καταλειφθῆναι ὑμῶν μηδένα,
\vs{8}παραπικράναι με ἐν τοῖς ἔργοις τῶν χειρῶν ὑμῶν, θυμιᾷν θεοῖς ἑτέροις ἐν γῇ Αἰγύπτῳ, εἰς ἣν ἤλθετε κατοικεῖν ἐκεῖ, ἵνα κοπῆτε, καὶ ἵνα γένησθε εἰς κατάραν καὶ εἰς ὀνειδισμὸν ἐν πᾶσι τοῖς ἔθνεσι τῆς γῆς;
\vs{9}Μὴ ἐπιλέλησθε ὑμεῖς τῶν κακῶν τῶν πατέρων ὑμῶν, καὶ τῶν κακῶν τῶν βασιλέων Ἰούδα, καὶ τῶν κακῶν τῶν ἀρχόντων ὑμῶν, καὶ τῶν κακῶν τῶν γυναικῶν ὑμῶν, ὧν ἐποίησαν ἐν γῇ Ἰούδα, καὶ ἔξωθεν Ἱερουσαλήμ;
\vs{10}Καὶ οὐκ ἐπαύσαντο ἕως τῆς ἡμέρας ταύτης, καὶ οὐκ ἀντείχοντο τῶν προσταγμάτων μου, ὧν ἔδωκα κατὰ πρόσωπον τῶν πατέρων αὐτῶν.

\vs{11}Διατοῦτο οὕτως εἶπε Κύριος, ἰδοὺ ἐγὼ ἐφίστημι τὸ πρόσωπόν μου,
\vs{12}τοῦ ἀπολέσαι πάντας τοὺς καταλοίπους τοὺς ἐν Αἰγύπτῳ, καὶ πεσοῦνται ἐν ῥομφαίᾳ καὶ ἐν λιμῷ, καὶ ἐκλείψουσιν ἀπὸ μικροῦ ἕως μεγάλου, καὶ ἔσονται εἰς ὀνειδισμὸν, καὶ εἰς ἀπώλειαν, καὶ εἰς κατάραν.
\vs{13}Καὶ ἐπισκέψομαι ἐπὶ τοὺς καθημένους ἐν γῇ Αἰγύπτῳ, ὡς ἐπεσκεψάμην ἐπὶ Ἱερουσαλὴμ, ἐν ῥομφαίᾳ καὶ ἐν λιμῷ,
\vs{14}καὶ οὐκ ἔσται σεσωσμένος οὐδεὶς τῶν ἐπιλοίπων Ἰούδα τῶν παροικούντων ἐν γῇ Αἰγύπτῳ, τοῦ ἐπιστρέψαι εἰς γῆν Ἰούδα, ἐφʼ ἣν αὐτοὶ ἐλπίζουσι ταῖς ψυχαῖς αὐτῶν τοῦ ἐπιστρέψαι ἐκεῖ· οὐ μὴ ἐπιστρέψωσιν, ἀλλʼ ἢ οἱ ἀνασεσωσμένοι.

\vs{15}Καὶ ἀπεκρίθησαν τῷ Ἱερεμίᾳ πάντες οἱ ἄνδρες οἱ γνόντες ὅτι θυμιῶσιν αἱ γυναῖκες αὐτῶν, καὶ πᾶσαι αἱ γυναῖκες, συναγωγὴ μεγάλη, καὶ πᾶς ὁ λαὸς οἱ καθήμενοι ἐν γῇ Αἰγύπτῳ, ἐν Παθουρῇ, λέγοντες,

\vs{16}Ὁ λόγος ὃν ἐλάλησας πρὸς ἡμᾶς τῷ ὀνόματι Κυρίου, οὐκ ἀκούσομέν σου,
\vs{17}ὅτι ποιοῦντες ποιήσομεν πάντα τὸν λόγον ὃς ἐξελεύσεται ἐκ τοῦ στόματος ἡμῶν, θυμιᾷν τῇ βασιλίσσῃ τοῦ οὐρανοῦ, καὶ σπένδειν αὐτῇ σπονδὰς, καθὰ ἐποιήσαμεν ἡμεῖς καὶ οἱ πατέρες ἡμῶν, καὶ οἱ βασιλεῖς ἡμῶν, καὶ οἱ ἄρχοντες ἡμῶν, ἐν πόλεσιν Ἰούδα καὶ ἔξωθεν Ἱερουσαλήμ· καὶ ἐπλήσθημεν ἄρτων, καὶ ἐγενόμεθα χρηστοὶ, καὶ κακὰ οὐκ εἴδομεν.
\vs{18}Καὶ ὡς διελίπομεν θυμιῶντες τῇ βασιλίσσῃ τοῦ οὐρανοῦ, ἠλαττώθημεν πάντες, καὶ ἐν ῥομφαίᾳ καὶ ἐν λιμῷ ἐξελίπομεν.
\vs{19}Καὶ ὅτι ἡμεῖς ἐθυμιῶμεν τῇ βασιλίσσῃ τοῦ οὐρανοῦ, καὶ ἐσπείσαμεν αὐτῇ σπονδὰς, μὴ ἄνευ τῶν ἀνδρῶν ἡμῶν ἐποιήσαμεν αὐτῇ χαυῶνας, καὶ ἐσπείσαμεν αὐτῇ σπονδάς;

\vs{20}Καὶ εἶπεν Ἱερεμίας παντὶ τῷ λαῷ, τοῖς δυνατοῖς καὶ ταῖς γυναιξὶ καὶ παντὶ τῷ λαῷ τοῖς ἀποκριθεῖσιν αὐτῷ λόγους, λέγων,
\vs{21}οὐχὶ τοῦ θυμιάματος οὗ ἐθυμιάσατε ἐν ταῖς πόλεσιν Ἰούδα, καὶ ἔξωθεν Ἱερουσαλὴμ, ὑμεῖς καὶ οἱ πατέρες ὑμῶν, καὶ οἱ βασιλεῖς ὑμῶν, καὶ οἱ ἄρχοντες ὑμῶν, καὶ ὁ λαὸς τῆς γῆς, ἐμνήσθη Κύριος; καὶ ἀνέβη ἐπὶ τὴν καρδίαν αὐτοῦ;
\vs{22}Καὶ οὐκ ἠδύνατο Κύριος ἔτι φέρειν ἀπὸ προσώπου πονηρίας πραγμάτων ὑμῶν, καὶ ἀπὸ τῶν βδελυγμάτων ὑμῶν ὧν ἐποιήσατε· καὶ ἐγενήθη ἡ γῆ ὑμῶν εἰς ἐρήμωσιν, καὶ εἰς ἄβατον, καὶ εἰς ἀρὰν, ὡς ἐν τῇ ἡμέρᾳ ταύτῃ,
\vs{23}ἀπὸ προσώπου ὧν ἐθυμιᾶτε, καὶ ὧν ἡμάρτετε τῷ Κύριῳ· καὶ οὐκ ἠκούσατε τῆς φωνῆς Κυρίου, καὶ ἐν τοῖς προστάγμασιν αὐτοῦ, καὶ ἐν τῷ νόμῳ καὶ ἐν τοῖς μαρτυρίοις αὐτοῦ οὐκ ἐπορεύθητε, καὶ ἐπελάβετο ὑμῶν τὰ κακὰ ταῦτα.

\vs{24}Καὶ εἶπεν Ἱερεμίας τῷ λαῷ, καὶ ταῖς γυναιξὶν, ἀκουσατε λόγον Κυρίου.
\vs{25}Οὕτως εἶπε Κύριος ὁ Θεὸς Ἰσραὴλ, ὑμεῖς γυναῖκες τῷ στόματι ὑμῶν ἐλαλήσατε, καὶ ταῖς χερσὶν ὑμῶν ἐπληρώσατε, λέγουσαι, ποιοῦσαι ποιήσομεν τὰς ὁμολογίας ἡμῶν ἃς ὡμολογήκαμεν, θυμιᾷν τῇ βασιλίσσῃ τοῦ οὐρανοῦ καὶ σπένδειν αὐτῇ σπονδάς· ἐμμείνασαι ἐνεμείνατε ταῖς ὁμολογίαις ὑμῶν, καὶ ποιοῦσαι ἐποιήσατε.
\vs{26}Διατοῦτο ἀκούσατε λόγον Κυρίου, πᾶς Ἰούδα οἱ καθήμενοι ἐν γῇ Αἰγύπτῳ, ἰδοὺ ὤμοσα τῷ ὀνόματί μου τῷ μεγάλῳ, εἶπε Κύριος, ἐὰν γένηται ἔτι ὄνομά μου ἐν τῷ στόματι παντὸς Ἰούδα εἰπεῖν, ζῇ Κύριος, ἐπὶ πάση γῇ Αἰγύπτῳ.
\vs{27}Ὅτι ἐγὼ ἐγρήγορα ἐπʼ αὐτοὺς, τοῦ κακῶσαι αὐτοὺς, καὶ οὐκ ἀγαθῶσαι· καὶ ἐκλείψουσι πᾶς Ἰούδα, οἱ κατοικοῦντες ἐν γῇ Αἰγύπτῳ, ἐν ῥομφαίᾳ καὶ ἐν λιμῷ, ἕως ἂν ἐκλείπωσι.
\vs{28}Καὶ οἱ σεσωσμένοι ἀπὸ ῥομφαίας ἐπιστρέψουσιν εἰς γῆν Ἰούδα ὀλίγοι ἀριθμῷ, καὶ γνώσονται οἱ κατάλοιποι Ἰούδα οἱ καταστάντες ἐν γῇ Αἰγύπτῳ, κατοικῆσαι ἐκεῖ, λόγος τίνος ἐμμενεῖ.

\vs{29}Καὶ τοῦτο τὸ σημεῖον ὑμῖν, ὅτι ἐπισκέψομαι ἐγὼ ἐφʼ ὑμᾶς εἰς πονηρά.
\vs{30}Οὕτως εἶπε Κύριος, ἰδοὺ ἐγὼ δίδωμι τὸν Οὐαφρῆ βασιλέα Αἰγύπτου εἰς χεῖρας ἐχθροῦ αὐτοῦ, καὶ εἰς χεῖρας ζητοῦντος τὴν ψυχὴν αὐτοῦ, καθὰ ἔδωκα τὸν Σεδεκίαν βασιλέα Ἰούδα εἰς χεῖρας Ναβουχοδονόσορ βασιλέως Βαβυλῶνος ἐχθροῦ αὐτοῦ, καὶ ζητοῦντος τὴν ψυχὴν αὐτοῦ.

\vs{31}Ὁ ΛΟΓΟΣ ὋΝ ἘΛΑΛΗΣΕΝ ἹΕΡΕΜΙΑΣ Ὁ ΠΡΟΦΗΤΗΣ πρὸς Βαροὺχ υἱὸν Νηρίου, ὅτε ἔγραφε τοὺς λόγους τούτους ἐν τῷ βιβλίῳ ἀπὸ στόματος Ἱερεμίου, ἐν τῷ ἐνιαυτῷ τῷ τετάρτῳ Ἰωακεὶμ υἱῷ Ἰωσία βασιλέως Ἰούδα.

\vs{32}Οὕτως εἶπε Κύριος ἐπὶ σοὶ Βαροὺχ,
\vs{33}ὅτι εἶπας, οἴμοι οἴμοι, ὅτι προσέθηκε Κύριος κόπον ἐπίπονόν μοι, ἐκοιμήθην ἐν στεναγμοῖς, ἀνάπαυσιν οὐχ εὗρον·
\vs{34}Εἰπὸν αὐτῷ, οὕτως εἶπε Κύριος, ἰδοὺ οὓς ἐγὼ ᾠκοδόμησα, ἐγὼ καθαιρῶ· καὶ οὓς ἐγὼ ἐφύτευσα, ἐγὼ ἐκτίλλω.
\vs{35}Καὶ σὺ ζητήσεις σεαυτῷ μεγάλα; μὴ ζητήσῃς, ὅτι ἰδοὺ ἐγὼ ἐπάγω κακὰ ἐπὶ πᾶσαν σάρκα, λέγει Κύριος· καὶ δώσω τὴν ψυχήν σου εἰς εὕρημα ἐν παντὶ τόπῳ οὗ ἐὰν βαδίσῃς ἐκεῖ.

\ch{52}
Ὄντος εἰκοστοῦ καὶ ἑνὸς ἔτους Σεδεκίου, ἐν τῷ βασιλεύειν αὐτὸν, καὶ ἕνδεκα ἔτη ἐβασίλευσεν ἐν Ἱερουσαλήμ· καὶ ὄνομα τῇ μητρὶ αὐτοῦ Ἀμειτάαλ, θυγάτηρ Ἱερεμίου, ἐκ Λοβενά.

\vs{4}Καὶ ἐγένετο τῷ ἔτει τῷ ἐννάτῳ τῆς βασιλείας αὐτοῦ, ἐν μηνὶ τῷ ἐννάτῳ, δεκάτῃ τοῦ μηνὸς, ἦλθε Ναβουχοδονόσορ βασιλεὺς Βαβυλῶνος, καὶ πᾶσα ἡ δύναμις αὐτοῦ ἐπὶ Ἱερουσαλὴμ, καὶ περιεχαράκωσαν αὐτὴν, καὶ περιῳκοδόμησαν αὐτὴν τετραπέδοις λίθοις κύκλῳ.

\vs{5}Καὶ ἦλθεν ἡ πόλις εἰς συνοχὴν, ἕως ἑνδεκάτου ἔτους τῷ βασιλεῖ Σεδεκίᾳ,
\vs{6}ἐν τῇ ἐννάτῃ τοῦ μηνὸς, καὶ ἐστερεώθη ὁ λιμὸς ἐν τῇ πόλει, καὶ οὐκ ἦσαν ἄρτοι τῷ λαῷ τῆς γῆς.
\vs{7}Καὶ διεκόπη ἡ πόλις, καὶ πάντες οἱ ἄνδρες οἱ πολεμισταὶ ἐξῆλθον νυκτὸς κατὰ τὴν ὁδὸν τῆς πύλης, ἀναμέσον τοῦ τείχους, καὶ τοῦ προτειχίσματος, ὃ ἦν κατὰ τὸν κῆπον τοῦ βασιλέως, καὶ οἱ Χαλδαῖοι ἐπὶ τῆς πόλεως κύκλῳ, καὶ ἐπορεύθησαν ὁδὸν τὴν εἰς ἄραβα,
\vs{8}Καὶ κατεδίωξεν ἡ δύναμις τῶν Χαλδαίων ὀπίσω τοῦ βασιλέως, καὶ κατέλαβον αὐτὸν ἐν τῷ πέραν Ἱερειχὼ, καὶ πάντες οἱ παῖδες αὐτοῦ διεσπάρησαν ἀπʼ αὐτοῦ.
\vs{9}Καὶ συνέλαβον τὸν βασιλέα, καὶ ἤγαγον αὐτὸν πρὸς τὸν βασιλέα Βαβυλῶνος εἰς Δεβλαθὰ, καὶ ἐλάλησεν αὐτῷ μετὰ κρίσεως.
\vs{10}Καὶ ἔσφαξε βασιλεὺς Βαβυλῶνος τοὺς υἱοὺς Σεδεκίου κατʼ ὀφθαλμοὺς αὐτοῦ, καὶ πάντας τοὺς ἄρχοντας Ἰούδα ἔσφαξεν ἐν Δεβλαθά.
\vs{11}Καὶ τοὺς ὀφθαλμοὺς Σεδεκίου ἐξετύφλωσε, καὶ ἔδησεν αὐτὸν ἐν πέδαις· καὶ ἤγαγεν αὐτὸν βασιλεὺς Βαβυλῶνος εἰς Βαβυλῶνα, καὶ ἔδωκεν αὐτὸν εἰς οἰκίαν μύλωνος, ἕως ἡμέρας ἧς ἀπέθανε.

\vs{12}Καὶ ἐν μηνὶ πέμπτῳ, δεκάτῃ τοῦ μηνὸς, ἦλθε Ναβουζαρδὰν ὁ ἀρχιμάγειρος, ἑστηκὼς κατὰ πρόσωπον τοῦ βασιλέως Βαβυλῶνος, εἰς Ἱερουσαλὴμ,
\vs{13}καὶ ἐνέπρησε τὸν οἶκον Κυρίου, καὶ τὸν οἶκον τοῦ βασιλέως, καὶ πάσας τὰς οἰκίας τῆς πόλεως, καὶ πᾶσαν οἰκίαν μεγάλην ἐνέπρησεν ἐν πυρί·
\vs{14}Καὶ πᾶν τεῖχος Ἱερουσαλὴμ κύκλῳ καθεῖλεν ἡ δύναμις τῶν Χαλδαίων, ἡ μετὰ τοῦ ἀρχιμαγείρου.
\vs{16}Καὶ τοὺς καταλοίπους τοῦ λαοῦ κατέλιπεν ὁ ἀρχιμάγειρος εἰς ἀμπελουργοὺς καὶ εἰς γεωργούς.

\vs{17}Καὶ τοὺς στύλους τοὺς χαλκοῦς τοὺς ἐν οἴκῳ Κυρίου, καὶ τὰς βάσεις, καὶ τὴν θάλασσαν τὴν χαλκῆν τὴν ἐν οἴκῳ Κυρίου συνέτριψαν οἱ Χαλδαῖοι, καὶ ἔλαβον τὸν χαλκὸν αὐτῶν, καὶ ἀπήνεγκαν εἰς Βαβυλῶνα.
\vs{18}Καὶ τὴν στεφάνην, καὶ τὰς φιάλας, καὶ τὰς κρεάγρας, καὶ πάντα τὰ σκεύη τὰ χαλκᾶ, ἐν οἷς ἐλειτούργουν ἐν αὐτοῖς,
\vs{19}καὶ τὰς ἀπφὼθ, καὶ τὰς μασμαρὼθ, καὶ τοὺς ὑποχυτῆρας, καὶ τὰς λυχνίας, καὶ τὰς θυΐσκας, καὶ τοὺς κυάθους, ἃ ἦν χρυσᾶ χρυσᾶ, καὶ ἃ ἦν ἀργυρᾶ ἀργυρᾶ, ἔλαβεν ὁ ἀρχιμάγειρος.
\vs{20}Καὶ οἱ στύλοι δύο, καὶ ἡ θάλασσα μία, καὶ οἱ μόσχοι δώδεκα χαλκοῖ ὑποκάτω τῆς θαλάσσης, ἃ ἐποίησεν ὁ βασιλεὺς Σαλωμὼν εἰς οἶκον Κυρίου, οὗ οὐκ ἦν σταθμὸς τοῦ χαλκοῦ αὐτῶν.

\vs{21}Καὶ οἱ στύλοι τριακονταπέντε πηχῶν ὕψος τοῦ στύλου τοῦ ἑνὸς, καὶ σπαρτίον δώδεκα πήχεων περιεκύκλου αὐτὸν, καὶ τὸ πάχος αὐτοῦ δακτύλων τεσσάρων κύκλῳ,
\vs{22}καὶ γεῖσος ἐπʼ αὐτοῖς χαλκοῦν, καὶ πέντε πήχεων τὸ μῆκος, ὑμεροχὴ τοῦ γείσους τοῦ ἑνός, καὶ δίκτυον καὶ ῥοαὶ ἐπὶ τοῦ γείσους κύκλῳ τὰ πάντα χαλκᾶ, καὶ κατὰ ταῦτα τῷ στύλῳ τῷ δευτέρῳ ὀκτὼ ῥοαὶ τῷ πήχει τοῖς δώδεκα πήχεσι.
\vs{23}Καὶ ἦσαν αἱ ῥοαὶ ἐννενηκονταὲξ τὸ ἓν μέρος, καὶ ἦσαν αἱ πᾶσαι ῥοαὶ ἑκατὸν ἐπὶ τοῦ δικτύου κύκλῳ.

\vs{24}Καὶ ἔλαβεν ὁ ἀρχιμάγειρος τὸν ἱερέα τὸν πρῶτον, καὶ τὸν ἱερέα τὸν δευτεροῦντα, καὶ τοὺς φυλάττοντας τὴν ὁδὸν,
\vs{25}καὶ εὐνοῦχον ἕνα ὃς ἦν ἐπιστάτης τῶν ἀνδρῶν τῶν πολεμιστῶν, καὶ εὐνοῦχον ἕνα ὃς ἦν ἐπιστάτης ἀνδρῶν τῶν πολεμιστῶν, καὶ ἑπτὰ ἄνδρας ὀνομαστοὺς, τοὺς ἐν προσώπῳ τοῦ βασιλέως, τοὺς εὑρεθέντας ἐν τῇ πόλει, καὶ τὸν γραμματέα τῶν δυνάμεων, τὸν γραμματεύοντα τῷ λαῷ τῆς γῆς, καὶ ἑξήκοντα ἀνθρώπους ἐκ τοῦ λαοῦ τῆς γῆς, τοὺς εὑρεθέντας ἐν μέσῳ τῆς πόλεως·
\vs{26}Καὶ ἔλαβεν αὐτοὺς Ναβουζαρδὰν ὁ ἀρχιμάγειρος τοῦ βασιλέως, καὶ ἢγαγεν αὐτοὺς πρὸς βασιλέα Βαβυλῶνος εἰς Δεβλαθά.
\vs{27}Καὶ ἐπάταξεν αὐτοὺς βασιλεὺς Βαβυλῶνος ἐν Δεβλαθὰ, ἐν γῇ Αἱμάθ.

\vs{31}Καὶ ἐγένετο ἐν τῷ τριακοστῷ καὶ ἑβδόμῳ ἔτει, ἀποικισθέντος τοῦ Ἰωακεὶμ βασιλέως Ἰούδα, ἐν τῷ δωδεκάτῳ μηνὶ, ἐν τῇ τετράδι καὶ εἰκάδι τοῦ μηνὸς, ἔλαβεν Οὐλαιμαδάχὰρ βασιλεὺς Βαβυλῶνος, ἐν τῷ ἐνιαυτῷ ᾧ ἐβασίλευσε, τὴν κεφαλὴν Ἰωακεὶμ βασιλέως Ἰούδα, καὶ ἔκειρεν αὐτὸν, καὶ ἐξήγαγεν αὐτὸν ἐξ οἰκίας ἧς ἐφυλάσσετο,
\vs{32}καὶ ἐλάλησεν αὐτῷ χρηστὰ, καὶ ἔδωκε τὸν θρόνον αὐτοῦ ἐπάνω τῶν βασιλέων τῶν μετʼ αὐτοῦ ἐν Βαβυλῶνι,
\vs{33}καὶ ἤλλαξε τὴν στολὴν τῆς φυλακῆς αὐτοῦ, καὶ ἤσθιεν ἄρτον διαπαντὸς κατὰ πρόσωπον αὐτοῦ πάσας τὰς ἡμέρας ἃς ἔζησε.
\vs{34}Καὶ ἡ σύνταξις αὐτῷ ἐδίδοτο διαπαντὸς παρὰ τοῦ βασιλέως Βαβυλῶνος ἐξ ἡμέρας εἰς ἡμέραν, ἕως ἡμέρας ἧς ἀπέθανε.


\def\book{ΘΡΗΝΟΙ ΙΕΡΕΜΙΟΥ}
\biblebook{ΘΡΗΝΟΙ ΙΕΡΕΜΙΟΥ}

\lettrine[lines=2, loversize=0.2, nindent=0em, findent=.25em]{\textcolor{bookheadingcolor}{Κ}}{ΑΙ} ἐγένετο μετὰ τὸ αἰχμαλωτισθῆναι τὸν Ἰσραὴλ, καὶ Ἱερουσαλὴμ ἐρημωθῆναι, ἐκάθισεν Ἱερεμίας κλαίων, καὶ ἐθρήνησε τὸν θρῆνον τοῦτον ἐπὶ Ἱερουσαλὴμ, καὶ εἶπε,


\vs{1}ΑΛΕΦ. Πῶς ἐκάθισε μόνη ἡ πόλις ἡ πεπληθυμμένη λαῶν; ἐγενήθη ὡς χήρα, πεπληθυμμένη ἐν ἔθνεσιν, ἄρχουσα ἐν χώραις ἐγενήθη εἰς φόρον.

\vs{2}ΒΗΘ. Κλαίουσα ἔκλαυσεν ἐν νυκτὶ, καὶ τὰ δάκρυα αὐτῆς ἐπὶ τῶν σιαγόνων αὐτῆς, καὶ οὐχ ὑπάρχει ὁ παρακαλῶν αὐτὴν ἀπὸ πάντων τῶν ἀγαπώντων αὐτήν· πάντες οἱ φιλοῦντες αὐτὴν ἠθέτησαν ἐν αὐτῇ, ἐγένοντο αὐτῇ εἰς ἐχθρούς.

\vs{3}ΓΙΜΕΛ. Μετῳκίσθη Ἰουδαία ἀπὸ ταπεινώσεως αὐτῆς, καὶ ἀπὸ πλήθους δουλείας αὐτῆς· ἐκάθισεν ἐν ἔθνεσιν, οὐχ εὗρεν ἀνάπαυσιν· πάντες οἱ καταδιώκοντες αὐτὴν, κατέλαβον αὐτὴν ἀναμέσον τῶν θλιβόντων.

\vs{4}ΔΑΛΕΘ. Ὁδοὶ Σιὼν πενθοῦσι παρὰ τὸ μὴ εἶναι ἐρχομένους ἐν ἑορτῇ· πᾶσαι αἱ πύλαι αὐτῆς ἠφανισμέναι, οἱ ἱερεῖς αὐτῆς ἀναστενάζουσιν, αἱ παρθένοι αὐτῆς ἀγόμεναι, καὶ αὐτὴ πικραινομένη ἐν ἑαυτῇ.

\vs{5}Η. Ἐγένοντο οἱ θλίβοντες αὐτὴν εἰς κεφαλὴν, καὶ οἱ ἐχθροὶ αὐτῆς εὐθηνοῦσαν, ὅτι Κύριος ἐταπείνωσεν αὐτὴν ἐπὶ τὸ πλῆθος τῶν ἀσεβειῶν αὐτῆς· τὰ νήπια αὐτῆς ἐπορεύθησαν ἐν αἰχμαλωσίᾳ κατὰ πρόσωπον θλίβοντος.

\vs{6}ΟΥΑΥ. Καὶ ἐξῄρθη ἐκ θυγατρὸς Σιὼν πᾶσα ἡ εὐπρέπεια αὐτῆς· ἐγένοντο οἱ ἄρχοντες αὐτῆς ὡς κριοὶ οὐχ εὑρίσκοντες νομὴν, καὶ ἐπορεύοντο ἐν οὐκ ἰσχύϊ κατὰ πρόσωπον διώκοντος.

\vs{7}ΖΑΙΝ. Ἐμνήσθη Ἱερουσαλὴμ ἡμερῶν ταπεινώσεως αὐτῆς, καὶ ἀπωσμῶν αὐτῆς· πάντα τὰ ἐπιθυμήματα αὐτῆς ὅσα ἦν ἐξ ἡμερῶν ἀρχαίων, ἐν τῷ πεσεῖν τὸν λαὸν αὐτῆς εἰς χεῖρας θλίβοντος, καὶ οὐκ ἦν ὁ βοηθῶν αὐτῇ. ἰδόντες οἱ ἐχθροὶ αὐτῆς, ἐγέλασαν ἐπὶ κατοικεσίᾳ αὐτῆς.

\vs{8}ΗΘ. Ἁμαρτίαν ἥμαρτεν Ἱερουσαλὴμ, διατοῦτο εἰς σάλον ἐγένετο· πάντες οἱ δοξάζοντες αὐτὴν, ἐταπείνωσαν αὐτήν· εἶδον γὰρ τὴν ἀσχημοσύνην αὐτῆς, καί γε αὐτὴ στενάζουσα καὶ ἀπεστράφη ὀπίσω.

\vs{9}ΤΗΘ. Ἀκαθαρσία αὐτῆς πρὸ ποδῶν αὐτῆς· οὐκ ἐμνήσθη ἔσχατα αὐτῆς, καὶ κατεβίβασεν ὑπέρογκα· οὐκ ἔστιν ὁ παρακαλῶν αὐτήν. ἴδε Κύριε τὴν ταπείνωσίν μου, ὅτι ἐμεγαλύνθη ὁ ἐχθρός.

\vs{10}ΙΩΔ. Χεῖρα αὐτοῦ ἐξεπέτασε θλίβων ἐπὶ πάντα τὰ ἐπιθυμήματα αὐτῆς· εἶδε γὰρ ἔθνη εἰσελθόντα εἰς τὸ ἁγίασμα αὐτῆς, ἃ ἐνετείλω μὴ εἰσελθεῖν αὐτὰ εἰς ἐκκλησίαν σου.

\vs{11}ΧΑΦ. Πᾶς ὁ λαὸς αὐτῆς καταστενάζοντες, ζητοῦντες ἄρτον· ἔδωκαν τὰ ἐπιθυμήματα αὐτῆς ἐν βρώσει, τοῦ ἐπιστρέψαι ψυχήν· ἴδε Κύριε καὶ ἐπίβλεψον, ὅτι ἐγενήθη ἠτιμωμένη.

\vs{12}ΛΑΜΕΔ. Οἱ πρὸς ὑμᾶς πάντες παραπορευόμενοι ὁδὸν, ἐπιστρέψατε καὶ ἴδετε εἰ ἔστιν ἄλγος κατὰ τὸ ἄλγος μου, ὃ ἐγενήθη· φθεγξάμενος ἐν ἐμοὶ ἐταπείνωσέ με Κύριος ἐν ἡμέρᾳ ὀργῆς θυμοῦ αὐτοῦ.

\vs{13}ΜΗΜ. Ἐξ ὕψους αὐτοῦ ἀπέστειλε πῦρ, ἐν τοῖς ὀστέοις μου κατήγαγεν αὐτό· διεπέτασε δίκτυον τοῖς ποσί μου, ἀπέστρεψέ με εἰς τὰ ὀπίσω· ἔδωκέ με ἠφανισμένην, ὅλην τὴν ἡμέραν ὠδυνωμένην.

\vs{14}ΝΟΥΝ. Ἐγρηγορήθη ἐπὶ τὰ ἀσεβήματά μου, ἐν χερσί μου συνεπλάκησαν· ἀνέβησαν ἐπὶ τὸν τράχηλόν μου· ἠσθένησεν ἡ ἰσχύς μου· ὅτι ἔδωκε Κύριος ἐν χερσί μου ὀδύνας, οὐ δυνήσομαι στῆναι.

\vs{15}ΣΑΜΕΧ. Ἐξῇρε πάντας τοὺς ἰσχυρούς μου ὁ Κύριος ἐκ μέσου μου, ἐκάλεσεν ἐπʼ ἐμὲ καιρὸν τοῦ συντρίψαι ἐκλεκτούς μου· ληνὸν ἐπάτησε Κύριος παρθένῳ θυγατρὶ Ἰούδα· ἐπὶ τούτοις ἐγὼ κλαίω.

\vs{16}ΑΙΝ. Ὁ ὀφθαλμός μου κατήγαγεν ὕδωρ, ὅτι ἐμακρύνθη ἀπʼ ἐμοῦ ὁ παρακαλῶν με, ὁ ἐπιστρέφων ψυχήν μου· ἐγένοντο οἱ υἱοί μου ἠφανισμένοι, ὅτι ἐκραταιώθη ὁ ἐχθρός.

\vs{17}ΦΗ. Διεπέτασε Σιὼν χεῖρα αὐτῆς, οὐκ ἔστιν ὁ παρακαλῶν αὐτήν· ἐνετείλατο Κύριος τῷ Ἰακὼβ, κύκλῳ αὐτοῦ οἱ θλίβοντες αὐτόν· ἐγενήθη Ἱερουσαλὴμ εἰς ἀποκαθημένην ἀναμέσον αὐτῶν.

\vs{18}ΤΣΑΔΗ. Δίκαιός ἐστι Κύριος, ὅτι στόμα αὐτοῦ παρεπίκρανα· ἀκούσατε δὴ πάντες οἱ λαοὶ, καὶ ἴδετε τὸ ἄλγος μου· παρθένοι μου καὶ νεανίσκοι μου ἐπορεύθησαν ἐν αἰχμαλωσίᾳ.

\vs{19}ΚΩΦ. Ἐκάλεσα τοὺς ἐραστάς μου, αὐτοὶ δὲ παρελογίσαντό με· οἱ ἱερεῖς μου καὶ οἱ πρεσβύτεροί μου ἐν τῇ πόλει ἐξέλιπον, ὅτι ἐζήτησαν βρῶσιν αὐτοῖς ἵνα ἐπιστρέψωσι ψυχὰς αὐτῶν, καὶ οὐχ εὗρον.

\vs{20}ΡΗΧΣ. Ἴδε, Κύριε, ὅτι θλίβομαι, ἡ κοιλία μου ἐταράχθη, καὶ ἡ καρδία μου ἐστράφη ἐν ἐμοὶ, ὅτι παραπικραίνουσα παρεπικράνθην, ἔξωθεν ἠτέκνωσέ με μάχαιρα, ὥσπερ θάνατος ἐν οἴκῳ.

\vs{21}ΧΣΕΝ. Ακούσατε δὴ, ὅτι στενάζω ἐγὼ, οὐκ ἔστιν ὁ παρακαλῶν με· πάντες οἱ ἐχθροί μου ἤκουσαν τὰ κακά μου, καὶ ἐχάρησαν, ὅτι σὺ ἐποίησας· ἐπήγαγες ἡμέραν, ἐκάλεσας καιρὸν, ἐγένοντο ὅμοιοι ἐμοί.

\vs{22}ΘΑΥ. Εἰσέλθοι πᾶσα ἡ κακία αὐτῶν κακὰ πρόσωπόν σου, καὶ ἐπιφύλλισον αὐτοῖς, ὃν τρόπον ἐποίησαν ἐπιφυλλίδα περὶ πάντων τῶν ἁμαρτημάτων μου, ὅτι πολλοὶ οἱ στεναγμοί μου, καὶ ἡ καρδία μου λυπεῖται.

\ch{2}
ΑΛΕΦ. Πῶς ἐγνόφωσεν ἐν ὀργῇ αὐτοῦ Κύριος τὴν θυγατέρα Σιών; κατέῤῥιψεν ἐξ οὐρανοῦ εἰς γῆν δόξασμα Ἰσραὴλ, καὶ οὐκ ἐμνήσθη ὑποποδίου ποδῶν αὐτοῦ.

\vs{2}ΒΗΘ. Ἐν ἡμέρᾳ ὀργῆς αὐτοῦ κατεπόντισε Κύριος, οὐ φεισάμενος· πάντα τὰ ὡραῖα Ἰακὼβ καθεῖλεν ἐν θυμῷ αὐτοῦ, τὰ ὀχυρώματα τῆς θυγατρὸς Ἰούδα ἐκόλλησεν εἰς τὴν γῆν, ἐβεβήλωσε βασιλέα αὐτῆς, καὶ ἄρχοντας αὐτῆς.

\vs{3}ΓΙΜΕΛ. Συνέκλασεν ἐν ὀργῇ θυμοῦ αὐτοῦ πᾶν κέρας Ἰσραὴλ, ἀπέστρεψεν ὀπίσω δεξιὰν αὐτοῦ ἀπὸ προσώπου ἐχθροῦ, καὶ ἀνῆψεν ἐν Ἰακὼβ ὡς πῦρ φλόγα, καὶ κατέφαγε πάντα τὰ κύκλῳ.

\vs{4}ΔΑΛΕΘ. Ἐνέτεινε τόξον αὐτοῦ ὡς ἐχθρὸς ὑπεναντίος, ἐστερέωσε δεξιὰν αὐτοῦ ὡς ὑπεναντίος, καὶ ἀπέκτεινε πάντα τὰ ἐπιθυμήματα τῶν ὀφθαλμῶν μου ἐν σκηνῇ θυγατρὸς Σιὼν, ἐξέχεεν ὡς πῦρ τὸν θυμὸν αὐτοῦ.

\vs{5}Η. Ἐγενήθη Κύριος ὡς ἐχθρὸς, κατεπόντισεν Ἰσραὴλ, κατεπόντισε τὰς βάρεις αὐτῆς, διέφθειρε τὰ ὀχυρώματα αὐτοῦ, καὶ ἐπλήθυνε τῇ θυγατρὶ Ἰούδα ταπεινουμένην καὶ τεταπεινωμένην.

\vs{6}ΟΥΑΥ. Καὶ διεπέτασεν ὡς ἄμπελον τὸ σκήνωμα αὐτοῦ, διέφθειρεν ἑορτὴν αὐτοῦ· ἐπελάθετο Κύριος ἃ ἐποίησεν ἐν Σιὼν ἑορτῆς καὶ σαββάτου, καὶ παρώξυνεν ἐμβριμήματι ὀργῆς αὐτοῦ βασιλέα καὶ ἱερέα καὶ ἄρχοντα.

\vs{7}ΖΑΙΝ. Ἀπώσατο Κύριος θυσιαστήριον αὐτοῦ, ἀπετίναξεν ἁγίασμα αὐτοῦ, συνέτριψεν ἐν χειρὶ ἐχθροῦ τεῖχος βάρεων αὐτῆς· φωνὴν ἔδωκαν ἐν οἴκῳ Κυρίου ὡς ἐν ἡμέρᾳ ἑορτῆς.

\vs{8}ΗΘ. Καὶ ἐπέστρεψε διαφθεῖραι τείχος θυγατρὸς Σιών· ἐξέτεινε μέτρον, οὐκ ἀπέστρεψε χεῖρα αὐτοῦ ἀπὸ καταπατήματος, καὶ ἐπένθησε τὸ προτείχισμα, καὶ τεῖχος ὁμοθυμαδὸν ἠσθένησεν.

\vs{9}ΤΗΘ. Ἐνεπάγησαν εἰς γῆν πύλαι αὐτῆς, ἀπώλεσε καὶ συνέτριψε μοχλοὺς αὐτῆς, βασιλέα αὐτῆς καὶ ἄρχοντα αὐτῆς ἐν τοῖς ἔθνεσιν· οὐκ ἔστι νόμος, καί γε προφῆται αὐτῆς οὐκ εἶδον ὅρασιν παρὰ Κυρίου.

\vs{10}ΙΩΔ. Ἐκάθισαν εἰς τὴν γῆν, ἐσιώπησαν πρεσβύτεροι θυγατρὸς Σιὼν, ἀνεβίβασαν χοῦν ἐπὶ τὴν κεφαλὴν αὐτῶν, περιεζώσαντο σάκκους, κατήγαγον εἰς γῆν ἀρχηγοὺς παρθένους ἐν Ἱερουσαλήμ.

\vs{11}ΧΑΦ. Ἐξέλιπον ἐν δάκρυσιν οἱ ὀφθαλμοί μου, ἐταράχθη ἡ καρδία μου, ἐξεχύθη εἰς τὴν γῆν ἡ δόξα μου, ἐπὶ τὸ σύντριμμα τῆς θυγατρὸς λαοῦ μου, ἐν τῷ ἐκλείπειν νήπιον καὶ θηλάζοντα ἐν πλατείαις πόλεως.

\vs{12}ΛΑΜΕΔ. Ταῖς μητράσιν αὐτῶν εἶπαν, ποῦ σῖτος, καὶ οἶνος; ἐν τῷ ἐκλύεσθαι αὐτοὺς ὡς τραυματίας ἐν πλατείαις πόλεως, ἐν τῷ ἐκχεῖσθαι ψυχὰς αὐτῶν εἰς κόλπον μητέρων αὐτῶν.

\vs{13}ΜΗΜ. Τί μαρτυρήσω σοι, ἢ τί ὁμοιώσω σοι, θύγατερ Ἱερουσαλήμ; τίς σώσει καὶ παρακαλέσει σε, παρθένος θύγατερ Σιών; ὅτι ἐμεγαλύνθη ποτήριον συντριβῆς σου, τίς ἰάσεταί σε;

\vs{14}ΝΟΥΝ. Προφῆταί σου εἴδοσάν σοι μάταια καὶ ἀφροσύνην, καὶ οὐκ ἀπεκάλυψαν ἐπὶ τὴν ἀδικίαν σου, τοῦ ἐπιστρέψαι αἰχμαλωσίαν σου, καὶ εἴδοσάν σοι λήμματα μάταια καὶ ἐξώσματα.

\vs{15}ΣΑΜΕΧ. Ἐκρότησαν ἐπὶ σὲ χεῖρας πάντες οἱ παραπορευόμενοι ὁδὸν, ἐσύρισαν καὶ ἐκίνησαν τὴν κεφαλὴν αὐτῶν ἐπὶ τὴν θυγατέρα Ἱερουσαλήμ. Αὕτη ἡ πόλις, ἐροῦσι, στέφανος εὐφροσύνης πάσης τῆς γῆς;

\vs{16}ΑΙΝ. Διήνοιξαν ἐπὶ σὲ στόμα αὐτῶν πάντες οἱ ἐχθροί σου, ἐσυρίσαν καὶ ἔβρυξαν ὀδόντας, καὶ εἶπαν, κατεπίομεν αὐτήν· πλὴν αὕτη ἡ ἡμέρα ἣν προσεδοκῶμεν, εὕρομεν αὐτὴν, εἴδομεν.

\vs{17}ΦΗ. Ἐποίησε Κύριος ἃ ἐνεθυμήθη, συνετέλεσε ῥῆμα αὐτοῦ, ἃ ἐνετείλατο ἐξ ἡμερῶν ἀρχαίων· καθεῖλε, καὶ οὐκ ἐφείσατο, καὶ ηὔφρανεν ἐπὶ σὲ ἐχθρὸν, ὕψωσε κέρας θλίβοντός σε.

\vs{18}ΤΣΑΔΗ. Ἐβόησε καρδία αὐτῶν πρὸς Κύριον, τείχη Σιὼν καταγάγετε ὡς χειμάῤῥους δάκρυα ἡμέρας καὶ νυκτός· μὴ δῷς ἔκνηψιν σεαυτῇ, μὴ σιωπήσαιτο θυγάτηρ ὀφθαλμῶν σου.

\vs{19}ΚΩΦ. Ἀνάστα, ἀγαλλίασαι ἐν νυκτὶ εἰς ἀρχὰς φυλακῆς σου, ἔκχεον ὡς ὕδωρ καρδίαν σου ἀπέναντι προσώπου Κυρίου, ἆρον πρὸς αὐτὸν χεῖράς σου περὶ ψυχῆς νηπίων σου, τῶν ἐκλυομένων λιμῷ ἐπʼ ἀρχῆς πασῶν ἐξόδων.

\vs{20}ΡΗΧΣ. Ἴδε Κύριε καὶ ἐπίβλεψον τίνι ἐπεφύλλισας οὕτως· εἰ φάγονται γυναῖκες καρπὸν κοιλίας αὐτῶν; ἐπιφυλλίδα ἐποίησε μάγειρος, φονευθήσονται νήπια θηλάζοντα μαστούς; ἀποκτενεῖς ἐν ἁγιάσματι Κυρίου ἱερέα καὶ προφήτην;

\vs{21}ΧΣΕΝ. Ἐκοιμήθησαν εἰς τὴν ἔξοδον παιδάριον καὶ πρεσβύτης· παρθένοι μου καὶ νεανίσκοι μου ἐπορεύθησαν ἐν αἰχμαλωσίᾳ, ἐν ῥομφαίᾳ καὶ ἐν λιμῷ ἀπέκτεινας, ἐν ἡμέρᾳ ὀργῆς σου ἐμαγείρευσας, οὐκ ἐφείσω.

\vs{22}ΘΑΥ. Ἐκάλεσεν ἡμέραν ἑορτῆς παροικίας μου κυκλόθεν, καὶ οὐκ ἐγένοντο ἐν ἡμέρᾳ ὀργῆς Κυρίου ἀνασωζόμενος καὶ καταλελειμμένος, ὡς ἐπεκράτησα, καὶ ἐπλήθυνα ἐχθρούς μου πάντας.

\ch{3}
ΑΛΕΦ. Ἐγὼ ἀνὴρ ὁ βλέπων πτωχείαν, ἐν ῥάβδῳ θυμοῦ αὐτοῦ ἐπʼ ἐμέ.
\vs{2}Παρέλαβέ με καὶ ἀπήγαγέ με εἰς σκότος, καὶ οὐ φῶς.
\vs{3}Πλὴν ἐν ἐμοὶ ἐπέστρεψε χεῖρα αὐτοῦ ὅλην τὴν ἡμέραν,
\vs{4}ἐπαλαίωσε σάρκα μου καὶ δέρμα μου, ὀστέα μου συνέτριψεν.

\vs{5}ΒΗΘ. Ἀνῳκοδόμησε κατʼ ἐμοῦ, καὶ ἐκύκλωσε κεφαλήν μου, καὶ ἐμόχθησεν·
\vs{6}ἐν σκοτεινοῖς ἐκάθισέ με ὡς νεκροὺς αἰῶνος.
\vs{7}Ἀνῳκοδόμησε κατʼ ἐμοῦ, καὶ οὐκ ἐξελεύσομαι, ἐβάρυνε χαλκόν μου.

\vs{8}ΓΙΜΕΛ. Καί γε κεκράξομαι καὶ βοήσω, ἀπέφραξε προσευχήν μου.

\vs{9}ΔΑΛΕΘ. Ἀνῳκοδόμησεν ὁδούς μου, ἐνέφραξε τρίβους μου,
\vs{10}ἐτάραξεν ἄρκος ἐνεδρεύουσα, αὐτός μοι λέων ἐν κρυφαίοις,
\vs{11}κατεδίωξεν ἀφεστηκότα, καὶ κατέπαυσέ με, ἔθετό με ἠφανισμένην.

\vs{12}Η. Ἐνέτεινε τόξον αὐτοῦ, καὶ ἐστήλωσέ με ὡς σκοπὸν εἰς βέλος.
\vs{13}Εἰσήγαγεν ἐν τοῖς νεφροῖς μου ἰοὺς φαρέτρας αὐτοῦ.
\vs{14}Ἐγενήθην γέλως παντὶ λαῷ μου, ψαλμὸς αὐτῶν ὅλην τὴν ἡμέραν.

\vs{15}ΟΥΑΥ. Ἐχόρτασέ με πικρίας, ἐμέθυσέ με χολῆς,
\vs{16}καὶ ἐξέβαλε ψήφῳ ὀδόντας μου· ἐψώμισέ με σποδὸν,
\vs{17}καὶ ἀπώσατο ἐξ εἰρήνης ψυχήν μου· ἐπελαθόμην ἀγαθά.
\vs{18}Καὶ ἀπώλετο νῖκός μου, καὶ ἡ ἐλπίς μου ἀπὸ Κυρίου.

\vs{19}ΖΑΙΝ. Ἐμνήσθην ἀπὸ πτωχείας μου, καὶ ἐκ διωγμοῦ πικρία καὶ χολή μου μνησθήσεται,
\vs{20}καὶ καταδολεσχήσει ἐπʼ ἐμὲ ἡ ψυχή μου.
\vs{21}Ταύτην τάξω εἰς τὴν καρδίαν μου, διατοῦτο ὑπομενῶ.

\vs{21a}ΗΘ. Τὰ ἐλέη Κυρίου, ὅτι οὐκ ἐξέλιπέ με, ὅτι οὐ συνετελέσθησαν οἱ οἰκτιρμοὶ αὐτοῦ· μῆνας εἰς τὰς πρωΐας ἐλέησον Κύριε, ὅτι οὐ συνετελέσθημεν, ὅτι οὐ συνετελέσθησαν οἱ οἰκτιρμοὶ αὐτοῦ.
\vs{21b}Καινὰ εἰς τὰς πρωΐας, πολλὴ ἡ πίστις σου.
\vs{21c}Μερίς μου Κύριος, εἶπεν ἡ ψυχή μου· διατοῦτο ὑπομενῶ αὐτῷ·

\vs{25}ΤΗΘ. Ἀγαθὸς Κύριος τοῖς ὑπομένουσιν αὐτὸν, ψυχὴ ἣ ζητήσει αὐτὸν, ἀγαθόν·
\vs{26}καὶ ὑπομενεῖ, καὶ ἡσυχάσει εἰς τὸ σωτήριον Κυρίου.

\vs{27}ΤΗΘ. Ἀγαθὸν ἀνδρὶ, ὅταν ἄρῃ ζυγὸν ἐν νεότητι αὐτοῦ,
\vs{28}καθήσεται κατὰ μόνας, καὶ σιωπήσεται, ὅτι ᾖρεν ἐφʼ ἑαυτῷ.

\vs{30}ΙΩΔ. Δώσει τῷ παίοντι αὐτὸν σιαγόνα, χορτασθήσεται ὀνειδισμῶν·
\vs{31}Ὅτι οὐκ εἰς τὸν αἰῶνα ἀπώσεται Κύριος.

\vs{32}ΧΑΦ. Ὅτι ὁ ταπεινώσας οἰκτειρήσει, καὶ κατὰ τὸ πλῆθος τοῦ ἐλέους αὐτοῦ,
\vs{33}οὐκ ἀπεκριθη ἀπὸ καρδίας αὐτοῦ, καὶ ἐταπείνωσεν υἱοὺς ἀνδρός.

\vs{34}ΛΑΜΕΔ. Τοῦ ταπεινῶσαι ὑπὸ τοὺς πόδας αὐτοῦ πάντας δεσμίους γῆς,
\vs{35}τοῦ ἐκκλίναι κρίσιν ἀνδρὸς κατέναντι προσώπου ὑψίστου,
\vs{36}καταδικάσαι ἄνθρωπον ἐν τῷ κρίνεσθαι αὐτὸν, Κύριος οὐκ εἶπε.
\vs{37}Τίς οὕτως εἶπε, καὶ ἐγενήθη; Κύριος οὐκ ἐνετείλατο.
\vs{38}Ἐκ στόματος ὑψίστου οὐκ ἐξελεύσεται τὰ κακὰ καὶ τὸ ἀγαθόν.

\vs{39}ΜΗΜ. Τί γογγύσει ἄνθρωπος ζῶν, ἀνὴρ περὶ τῆς ἁμαρτίας αὐτοῦ;

\vs{40}ΝΟΥΝ. Ἐξηρευνήθη ἡ ὁδὸς ἡμῶν καὶ ἠτάσθη, καὶ ἐπιστρέψομεν ἕως Κυρίου.
\vs{41}Ἀναλάβωμεν καρδίας ἡμῶν ἐπὶ χειρῶν πρὸς ὑψηλὸν ἐν οὐρανῷ.
\vs{42}Ἡμαρτήσαμεν, ἠσεβήσαμεν, καὶ οὐχ ἱλάσθης.

\vs{43}ΣΑΜΕΧ. Ἐπεσκέπασας ἐν ʼθυμῷ, καὶ ἀπεδίωξας ἡμᾶς, ἀπέκτεινας, οὐκ ἐφείσω.
\vs{44}Ἐπεσκέπασας νεφέλην σεαυτῷ ἕινεκεν προσευχῆς,
\vs{45}καμμύσαι με καὶ ἀπωσθῆναι. ΑΙΝ. Ἔθηκας ἡμᾶς ἐν μέσῳ τῶν λαῶν.
\vs{46}Διήνοιξαν ἐφʼ ἡμᾶς τὸ στόμα αὐτῶν πάντες οἱ ἐχθροὶ ἡμῶν.
\vs{47}Φόβος καὶ θυμὸς ἐγενήθη ἡμῖν, ἔπαρσις καὶ συντριβή.
\vs{48}Ἀφέσεις ὑδάτων κατάξει ὁ ὀφθαλμός μου ἐπὶ τὸ σύντριμμα τῆς θυγατρὸς τοῦ λαοῦ μου.

\vs{49}ΦΗ. Ὁ ὀφθαλμός μου κατεπόθη, καὶ οὐ σιγήσομαι τοῦ μὴ εἶναι ἔκνηψιν,
\vs{50}ἕως οὗ διακύψῃ καὶ ἴδῃ Κύριος ἐξ οὐρανοῦ.
\vs{51}Ὁ ὀφθαλμός μου ἐπιφυλλιεῖ ἐπὶ τὴν ψυχήν μου παρὰ πάσας θυγατέρας πόλεως.

\vs{52}ΤΣΑΔΗ. Θηρεύοντες ἐθήρευσάν με ὡς στρουθίον· πάντες οἱ ἐχθροί μου δωρεὰν
\vs{53}ἐθανάτωσαν ἐν λάκκῳ ζωήν μου, καὶ ἐπέθηκαν λίθον ἐπʼ ἐμοί.
\vs{54}Ὑπερεχύθη ὕδωρ ἐπὶ τὴν κεφαλήν μου· εἶπα, ἄπωσμαι.

\vs{55}ΚΩΦ. Ἐπεκαλεσάμην τὸ ὄνομά σου, Κύριε, ἐκ λάκκου κατωτάτου·
\vs{56}Φωνήν μου ἤκουσας· μὴ κρύψῃς τὰ ὦτά σου εἰς τὴν δέησίν μου·
\vs{57}Εἰς τὴν βοήθειάν μου ἤγγισας· ἐν ἡμέρᾳ ᾗ ἐπεκαλεσάμην σε εἶπάς μοι, μὴ φοβοῦ.

\vs{58}ΡΗΧΣ. Ἐδίκασας, Κύριε, τὰς δίκας τῆς ψυχῆς μου, ἐλυτρώσω τὴν ζωήν μου.
\vs{59}Ἴδες, Κύριε, τὰς ταραχάς μου, ἔκρινας τὴν κρίσιν μου.
\vs{60}Εἶδες πᾶσαν τὴν ἐκδίκησιν αὐτῶν, εἰς πάντας διαλογισμοὺς αὐτῶν ἐν ἐμοί.

\vs{61}ΧΣΕΝ. Ἤκουσας τὸν ὀνειδισμὸν αὐτῶν, πάντας τοὺς διαλογισμοὺς αὐτῶν κατʼ ἐμοῦ,
\vs{62}χείλη ἐπανισταμένων μοι, καὶ μελέτας αὐτῶν κατʼ ἐμοῦ ὅλην τὴν ἡμέραν,
\vs{63}καθέδραν αὐτῶν, καὶ ἀνάστασιν αὐτῶν· ἐπίβλεψον ἐπὶ ὀφθαλμοὺς αὐτῶν.
\vs{64}Ἀποδώσεις αὐτοῖς ἀνταπόδομα, Κύριε, κατὰ τὰ ἔργα τῶν χειρῶν αὐτῶν.

\vs{65}ΘΑΥ. Ἀποδώσεις αὐτοῖς ὑπερασπισμὸν, καρδίας μου μόχθον.
\vs{66}Σὺ αὐτοὺς καταδιώξεις ἐν ὀργῇ, καὶ ἐξαναλώσεις αὐτοὺς ὑποκάτωθεν τοῦ οὐρανοῦ, Κύριε.

\ch{4}
ΑΛΕΦ. Πῶς ἀμαυρωθήσεται χρυσίον, ἀλλοιωθήσεται τὸ ἀργύριον τὸ ἀγαθόν; ἐξεχύθησαν λίθοι ἅγιοι ἐπʼ ἀρχῆς πασῶν ἐξόδων.

\vs{2}ΒΗΘ. Οἱ υἱοὶ Σιὼν οἱ τίμιοι, οἱ ἐπῃρμένοι ἐν χρυσίῳ, πῶς ἐλογίσθησαν εἰς ἀγγεῖα ὀστράκινα, ἔργα χειρῶν κεραμέως;

\vs{3}ΓΙΜΕΛ. Καί γε δράκοντες ἐξέδυσαν μαστοὺς, ἐθήλασαν σκύμνοι αὐτῶν θυγατέρας λαοῦ μου εἰς ἀνίατον, ὡς στρουθίον ἐν ἐρήμῳ.

\vs{4}ΔΑΛΕΘ. Ἐκολλήθη ἡ γλῶσσα θηλάζοντος πρὸς τὸν φάρυγγα αὐτοῦ ἐν δίψει, νήπια ᾔτησαν ἄρτον, ὁ διακλῶν οὐκ ἔστιν αὐτοῖς.

\vs{5}Η. Οἱ ἔσθοντες τὰς τρυφὰς ἠφανίσθησαν ἐν ταῖς ἐξόδοις, οἱ τιθηνούμενοι ἐπὶ κόκκων περιεβάλλοντο κοπρίας.

\vs{6}ΟΥΑΥ. Καὶ ἐμεγαλύνθη ἀνομία θυγατρὸς λαοῦ μου ὑπὲρ ἀνομίας Σοδόμων τῆς κατεστραμμένης ὥσπερ σπουδὴ, καὶ οὐκ ἐπόνεσαν ἐν αὐτῇ χεῖρας.

\vs{7}ΖΑΙΝ. Ἐκαθαριώθησαν Ναζιραῖοι αὐτῆς ὑπὲρ χιόνα, ἔλαμψαν ὑπὲρ γάλα, ἐπυρώθησαν, ὑπὲρ λίθου σαπφείρου τὸ ἀπόσπασμα αὐτῶν.

\vs{8}ΗΘ. Ἐσκότασεν ὑπὲρ ἀσβόλην τὸ εἶδος αὐτῶν, οὐκ ἐπεγνώσθησαν ἐν ταῖς ἐξόδοις· ἐπάγη δέρμα αὐτῶν ἐπὶ τὰ ὀστέα αὐτῶν, ἐξηράνθησαν, ἐγενήθησαν ὥσπερ ξύλον.

\vs{9}ΤΗΘ. Καλοὶ ἦσαν οἱ τραυματίαι ῥομφαίας, ἢ οἱ τραυματίαι λιμοῦ· ἐπορεύθησαν ἐκκεκεντημένοι ἀπὸ γεννημάτων ἀγρῶν.

\vs{10}ΙΩΔ. Χεῖρες γυναικῶν οἰκτιρμόνων ἥψησαν τὰ παιδία αὐτῶν, ἐγενήθησαν εἰς βρῶσιν αὐταῖς, ἐν τῷ συντρίμματι τῆς θυγατρὸς τοῦ λαοῦ μου.

\vs{11}ΧΑΦ. Συνετέλεσε Κύριος θυμὸν αὐτοῦ, ἐξέχεε θυμὸν ὀργῆς αὐτοῦ, καὶ ἀνῆψε πῦρ ἐν Σιὼν, καὶ κατέφαγε τὰ θεμέλια αὐτῆς.

\vs{12}ΛΑΜΕΔ. Οὐκ ἐπίστευσαν βασιλεῖς γῆς, πάντες οἱ κατοικοῦντες τὴν οἰκουμένην, ὅτι εἰσελεύσεται ἐχθρὸς καὶ ἐκθλίβων διὰ τῶν πυλῶν Ἱερουσαλήμ.

\vs{13}ΜΗΜ. Ἐξ ἁμαρτιῶν προφητῶν αὐτῆς, ἀδικιῶν ἱερέων αὐτῆς, τῶν ἐκχεόντων αἷμα δίκαιον ἐν μέσῳ αὐτῆς.

\vs{14}ΝΟΥΝ. Ἐσαλεύθησαν ἐγρήγοροι αὐτῆς ἐν ταῖς ἐξόδοις, ἐμολύνθησαν ἐν αἵματι ἐν τῷ μὴ δύνασθαι αὐτοὺς, ἥψαντο ἐνδυμάτων αὐτῶν.

\vs{15}ΣΑΜΕΧ. Ἀπόστητε ἀκαθάρτων, καλέσατε αὐτοὺς, ἀπόστητε, ἀπόστητε, μὴ ἅπτεσθε, ὅτι ἀνήφθησαν, καί γε ἐσαλεύθησαν· εἴπατε ἐν τοῖς ἔθνεσιν, οὐ μὴ προσθῶσι τοῦ παροικεῖν.

\vs{16}ΑΙΝ. Πρόσωπον Κυρίου μερὶς αὐτῶν, οὐ προσθήσει ἐπιβλέψαι αὐτοῖς· πρόσωπον ἱερέων οὐκ ἔλαβον, προφήτας οὐκ ἠλέησαν.

\vs{17}ΦΗ. Ἔτι ὄντων ἡμῶν ἐξέλιπον οἱ ὀφθαλμοὶ ἡμῶν, εἰς τὴν βοήθειαν ἡμῶν μάταια ἀποσκοπευόντων ἡμῶν.
\vs{18}ΤΣΑΔΗ. Ἀπεσκοπεύσαμεν εἰς ἔθνος οὐ σῶζον, ἐθηρεύσαμεν μικροὺς ἡμῶν, τοῦ μὴ πορεύεσθαι ἐν ταῖς πλατείαις ἡμῶν.
\vs{19}ΚΩΦ. Ἤγγικεν ὁ καιρὸς ἡμῶν, ἐπληρώθησαν αἱ ἡμέραι ἡμῶν, πάρεστιν ὁ καιρὸς ἡμῶν. Κοῦφοι ἐγένοντο οἱ διώκοντες ἡμᾶς ὑπὲρ ἀετοὺς οὐρανοῦ, ἐπὶ τῶν ὀρέων ἐξέπτησαν, ἐν ἐρήμῳ ἐνήδρευσαν ἡμᾶς.

\vs{20}ΡΗΧΣ. Πνεῦμα προσώπου ἡμῶν χριστὸς Κύριος συνελήφθη ἐν ταῖς διαφθοραῖς αὐτῶν, οὗ εἴπαμεν, ἐν τῇ σκιᾶ αὐτοῦ ζησόμεθα ἐν τοῖς ἔθνεσι.

\vs{21}ΧΣΕΝ. Χαῖρε καὶ εὐφραίνου θύγατερ Ἰδουμαίας ἡ κατοικοῦσα ἐπὶ γῆς, καί γε ἐπὶ σὲ διελεύσεται τὸ ποτήριον Κυρίου, μεθυσθήσῃ καὶ ἀποχεεῖς.

\vs{22}ΘΑΥ. Ἐξέλιπεν ἡ ἀνομία σου θύγατερ Σιὼν, οὐ προσθήσει τοῦ ἀποικίσαι σε· ἐπεσκέψατο ἀνομίας σου θύγατερ Ἐδὼμ, ἀπεκάλυψεν ἐπὶ τὰ ἀσεβήματά σου.

\ch{5}
Μνήσθητι Κύριε, ὅ, τι ἐγενήθη ἡμῖν· ἐπίβλεψον, καὶ ἴδε τὸν ὀνειδισμὸν ἡμῶν.

\vs{2}Κληρονομία ἡμῶν μετεστράθη ἀλλοτρίοις, οἱ οἶκοι ἡμῶν ξένοις.
\vs{3}Ὀρφανοὶ ἐγενήθημεν, οὐχ ὑπάρχει πατὴρ, μητέρες ἡμῶν ὡς αἱ χῆραι.
\vs{4}Ὕδωρ ἡμῶν ἐν ἀργυρίῳ ἐπίομεν, ξύλα ἡμῶν ἐν ἀλλάγματι ἦλθεν ἐπὶ τὸν τράχηλον ἡμῶν·
\vs{5}ἐδιώχθημεν, ἐκοπιάσαμεν, οὐκ ἀνεπαύθημεν.

\vs{6}Αἴγυπτος ἔδωκε χεῖρα, Ἀσσοὺρ εἰς πλησμονὴν αὐτῶν.
\vs{7}Οἱ πατέρες ἡμῶν ἥμαρτον, οὐχ ὑπάρχουσιν, ἡμεῖς τὰ ἀνομήματα αὐτῶν ὑπέσχομεν.
\vs{8}Δοῦλοι ἐκυρίευσαν ἡμῶν, λυτρούμενος οὐκ ἔστιν ἐκ τῆς χειρὸς αὐτῶν.
\vs{9}Ἐν ταῖς ψυχαῖς ἡμῶν εἰσοίσομεν ἄρτον ἡμῶν, ἀπὸ προσώπου ῥομφαίας τῆς ἐρήμου.
\vs{10}Τὸ δέρμα ἡμῶν ὡς κλίβανος ἐπελιώθη, συνεσπάσθησαν ἀπὸ προσώπου καταιγίδων λιμοῦ.
\vs{11}Γυναῖκας ἐν Σιὼν ἐταπείνωσαν, παρθένους ἐν πόλεσιν Ἰούδα.
\vs{12}Ἄρχοντες ἐν χερσὶν αὐτῶν ἐκρεμάσθησαν, πρεσβύτεροι οὐκ ἐδοξάσθησαν.
\vs{13}Ἐκλεκτοὶ κλαυθμὸν ἀνέλαβον, καὶ νεανίσκοι ἐν ξύλῳ ἠσθένησαν.
\vs{14}Καὶ πρεσβύται ἀπὸ πύλης κατέπαυσαν, ἐκλεκτοὶ ἐκ ψαλμῶν αὐτῶν κατέπαυσαν.
\vs{15}Κατέλυσε χαρὰ καρδίας ἡμῶν, ἐστράφη εἰς πένθος ὁ χορὸς ἡμῶν·
\vs{16}Ἔπεσεν ὁ στέφανος ἡμῶν τῆς κεφαλῆς· οὐαὶ δὲ ἡμῖν, ὅτι ἡμάρτομεν.

\vs{17}Περὶ τούτου ἐγενήθη ὀδύνη, ὀδυνηρὰ ἡ καρδία ἡμῶν, περὶ τούτου ἐσκότασαν οἱ ὀφθαλμοὶ ἡμῶν.
\vs{18}Ἐπʼ ὄρος Σιὼν, ὅτι ἠφανίσθη, ἀλώπεκες διῆλθον ἐν αὐτῇ.

\vs{19}Σὺ δὲ Κύριε εἰς τὸν αἰῶνα κατοικήσεις, ὁ θρόνος σου εἰς γενεὰν καὶ γενεάν.
\vs{20}Ἱνατί εἰς νῖκος ἐπιλήσῃ ἡμῶν, καταλείψεις ἡμᾶς εἰς μακρότητα ἡμερῶν;
\vs{21}Ἐπίστρεψον ἡμᾶς Κύριε πρὸς σὲ, καὶ ἐπιστραφησόμεθα· καὶ ἀνακαίνισον ἡμέρας ἡμῶν καθὼς ἔμπροσθεν.
\vs{22}Ὅτι ἀπωθούμενος ἀπώσω ἡμᾶς, ὠργίσθης ἐφʼ ἡμᾶς ἕως σθόδρα.


\def\book{ΙΕΖΕΚΙΗΛ}
\biblebook{ΙΕΖΕΚΙΗΛ}


\lettrine[lines=2, loversize=0.2, nindent=0em, findent=.25em]{\textcolor{bookheadingcolor}{Κ}}{ΑΙ} ἐγένετο ἐν τῷ τριακοστῷ ἔτει ἐν τῷ τετάρτῳ μηνὶ πέμπτῃ τοῦ μηνὸς, καὶ ἐγὼ ἤμην ἐν μέσῳ τῆς αἰχμαλωσίας ἐπὶ τοῦ ποταμοῦ τοῦ Χοβάρ· καὶ ἠνοίχθησαν οἱ οὐρανοὶ, καὶ ἴδον ὁράσεις Θεοῦ.
\vs{2}Πέμπτῃ τοῦ μηνός· τοῦτο τὸ ἔτος τὸ πέμπτον τῆς αἰχμαλωσίας τοῦ βασιλέως Ἰωακείμ·
\vs{3}Καὶ ἐγένετο λόγος Κυρίου πρὸς Ἰεζεκιὴλ υἱὸν Βουζεὶ, τὸν ἱερέα, ἐν γῇ Χαλδαίων, ἐπὶ τοῦ ποταμοῦ τοῦ Χοβάρ· καὶ ἐγένετο ἐπʼ ἐμὲ χεὶρ Κυρίου.

\vs{4}Καὶ ἴδον, καὶ ἰδοὺ πνεῦμα ἐξαῖρον ἤρχετο ἀπὸ Βοῤῥᾶ, καὶ νεφέλη μεγάλη ἐν αὐτῷ, καὶ φέγγος κύκλῳ αὐτοῦ καὶ πῦρ ἐξαστράπτον· καὶ ἐν τῷ μέσῳ αὐτοῦ ὡς ὅρασις ἠλέκτρου ἐν μέσῳ τοῦ πυρὸς, καὶ φέγγος ἐν αὐτῷ.

\vs{5}Καὶ ἐν τῷ μέσῳ ὡς ὁμοίωμα τεσσάρων ζῴων· καὶ αὕτη ἡ ὅρασις αὐτῶν· ὁμοίωμα ἀνθρώπου ἐπʼ αὐτοῖς.
\vs{6}Καὶ τέσσαρα πρόσωπα τῷ ἑνὶ, καὶ τέσσαρες πτέρυγες τῷ ἑνί.
\vs{7}Καὶ τὰ σκέλη αὐτῶν ὀρθὰ, καὶ πτερωτοὶ οἱ πόδες αὐτῶν, καὶ σπινθῆρες, ὡς ἐξαστράπτων χαλκός· καὶ ἐλαφραὶ αἱ πτέρυγες αὐτῶν.
\vs{8}Καὶ χεὶρ ἀνθρώπου ὑποκάτωθεν τῶν πτερύγων αὐτῶν ἐπὶ τὰ τέσσαρα μέρη αὐτῶν.
\vs{9}Καὶ τὰ πρόσωπα αὐτῶν τῶν τεσσάρων οὐκ ἐπεστρέφοντο ἐν τῷ βαδίζειν αὐτά· ἕκαστον ἀπέναντι τοῦ προσώπου αὐτῶν ἐπορεύοντο.
\vs{10}Καὶ ὁμοίωσις τῶν προσώπων αὐτῶν, πρόσωπον ἀνθρώπου, καὶ πρόσωπον τοῦ λέοντος ἐκ δεξιῶν τοῖς τέσσαρσι, καὶ πρόσωπον μόσχου ἐξ ἀριστερῶν τοῖς τέσσαρσι, καὶ πρόσωπον ἀετοῦ τοῖς τέσσαρσι.
\vs{11}Καὶ αἱ πτέρυγες αὐτῶν ἐκτεταμέναι ἄνωθεν τοῖς τέσσαρσιν, ἑκατέρῳ δύο συνεζευγμέναι πρὸς ἀλλήλας, καὶ δύο ἐπεκάλυπτον ἐπάνω τοῦ σώματος αὐτῶν.
\vs{12}Καὶ ἑκάτερον κατὰ πρόσωπον αὐτοῦ ἐπορεύετο· οὗ ἂν ἦν τὸ πνεῦμα πορευόμενον ἐπορεύοντο, καὶ οὐκ ἐπέστρεφον.

\vs{13}Καὶ ἐν μέσῳ τῶν ζώων ὅρασις ὡς ἀνθράκων πυρὸς καιομένων, ὡς ὄψις λαμπάδων συστρεφομένων ἀναμέσον τῶν ζώων, καὶ φέγγος τοῦ πυρὸς, καὶ ἐκ τοῦ πυρὸς ἐξεπορεύετο ἀστραπή.

\vs{15}Καὶ ἴδον, καὶ ἰδοὺ τροχὸς εἷς ἐπὶ τῆς γῆς ἐχόμενος τῶν ζώων τοῖς τέσσαρσι.
\vs{16}Καὶ τὸ εἶδος τῶν τροχῶν ὡς εἶδος θαρσείς· καὶ ὁμοίωμα ἓν τοῖς τέσσαρσι· καὶ τὸ ἔργον αὐτῶν ἦν καθὼς ἂν εἴη τροχὸς ἐν τροχῷ.
\vs{17}Ἐπὶ τὰ τέσσαρα μέρη αὐτῶν ἐπορεύοντο· οὐκ ἐπέστρεφον ἐν τῷ πορεύεσθαι αὐτὰ,
\vs{18}οὐδʼ οἱ νῶτοι αὐτῶν· καὶ ὕψος ἦν αὐτοῖς· καὶ ἴδον αὐτὰ, καὶ οἱ νῶτοι αὐτῶν πλήρεις ὀφθαλμῶν κυκλόθεν τοῖς τέσσαρσι.

\vs{19}Καὶ ἐν τῷ πορεύεσθαι τὰ ζῶα, ἐπορεύοντο οἱ τροχοὶ ἐχόμενοι αὐτῶν· καὶ ἐν τῷ ἐξαίρειν τὰ ζῶα ἀπὸ τῆς γῆς, ἐξῄροντο οἱ τροχοί.
\vs{20}Οὗ ἂν ἦν ἡ νεφέλη, ἐκεῖ τὸ πνεῦμα τοῦ πορεύεσθαι, ἐπορεύοντο οἱ τροχοὶ καὶ ἐξῄροντο σὺν αὐτοῖς, διότι πνεῦμα ζωῆς ἐν τοῖς τροχοῖς.
\vs{21}Ἐν τῷ πορεύεσθαι αὐτὰ ἐπορεύοντο, καὶ ἐν τῷ ἑστάναι αὐτὰ εἱστήκεισαν· καὶ ἐν τῷ ἐξαίρειν αὐτὰ ἀπὸ τῆς γῆς, ἐξῄροντο σὺν αὐτοῖς, ὅτι πνεῦμα ζωῆς ἦν ἐν τοῖς τροχοῖς.

\vs{22}Καὶ ὁμοίωμα ὑπὲρ κεφαλῆς αὐτῶς τῶν ζώων ὡσεὶ στερέωμα, ὡς ὅρασις κρυστάλλου, ἐκτεταμένον ἐπὶ τῶν πτερύγων αὐτῶν ἐπάνωθεν.
\vs{23}Καὶ ὑποκάτωθεν τοῦ στερεώματος αἱ πτέρυγες αὐτῶν ἐκτεταμέναι, πτερυσσόμεναι ἑτέρα τῇ ἑτέρᾳ, ἑκάστῳ δύο ἐπικαλύπτουσαι τὰ σώματα αὐτῶν.
\vs{24}Καὶ ἤκουον τὴν φωνὴν τῶν πτερύγων αὐτῶν ἐν τῷ πορεύεσθαι αὐτὰ, ὡς φωνὴν ὕδατος πολλοῦ· καὶ ἐν τῷ ἑστάναι αὐτὰ, κατέπαυον αἱ πτέρυγες αὐτῶν.

\vs{25}Καὶ ἰδοὺ φωνὴ ὑπεράνωθεν τοῦ στερεώματος τοῦ ὄντος ὑπὲρ κεφαλῆς αὐτῶν,
\vs{26}ὡς ὅρασις λίθου σαπφείρου, ὁμοίωμα θρόνου ἐπʼ αὐτοῦ, καὶ ἐπὶ τοὺ ὁμοιώματος του θρόνου ὁμοίωμα ὡς εἶδος ἀνθρώπου ἄνωθεν.
\vs{27}Καὶ ἴδον ὡς ὄψιν ἠλέκτρου ἀπὸ ὁράσεως ὀσφύος καὶ ἐπάνω, καὶ ἀπὸ ὁράσεως ὀσφύος καὶ ἕως κάτω ἴδον ὅρασιν πυρὸς, καὶ τὸ φέγγος αὐτοῦ κύκλῳ,
\vs{28}ὡς ὅρασις τόξου ὅταν ᾖ ἐν τῇ νεφέλῃ ἐν ἡμέραις ὑετοῦ, οὕτως ἡ στάσις τοῦ φέγγους κυκλόθεν.

\ch{2}
Αὕτη ἡ ὅρασις ὁμοιώματος δόξης Κυρίου· καὶ ἴδον, καὶ πίπτω ἐπὶ πρόσωπόν μου, καὶ ἤκουσα φωνὴν λαλοῦντος· καὶ εἶπε πρὸς μὲ, υἱὲ ἀνθρώπου, στῆθι ἐπὶ τοὺς πόδας σου, καὶ λαλήσω πρὸς σέ·

\vs{2}Καὶ ἦλθεν ἐπʼ ἐμὲ πνεῦμα, καὶ ἀνέλαβέ με, καὶ ἐξῇρέ με, καὶ ἔστησέ με ἐπὶ τοὺς πόδας μου, καὶ ἤκουον αὐτοῦ λαλοῦντος πρὸς μέ.
\vs{3}Καὶ εἶπε πρὸς μὲ, υἱὲ ἀνθρώπου, ἐξαποστέλλω ἐγώ σε πρὸς τὸν οἶκον τοῦ Ἰσραὴλ, τοὺς παραπικραίνοντάς με, οἵτινες παρεπίκρανάν με, αὐτοὶ καὶ οἱ πατέρες αὐτῶν ἕως τῆς σήμερον ἡμέρας.
\vs{4}Καὶ ἐρεῖς πρὸς αὐτοὺς, τάδε λέγει Κύριος,
\vs{5}ἐὰν ἆρα ἀκούσωσιν ἢ πτοηθῶσι, διότι οἶκος παραπικραίνων ἐστὶ, καὶ γνώσονται ὅτι προφήτης εἶ σὺ ἐν μέσῳ αὐτῶν.

\vs{6}Καὶ σὺ, υἱὲ ἀνθρώπου, μὴ φοβηθῇς αὐτοὺς, μηδὲ ἐκστῇς ἀπὸ προσώπου αὐτῶν· διότι παροιστρήσουσι, καὶ ἐπισυστήσονται ἐπὶ σὲ κύκλῳ, καὶ ἐν μέσῳ σκορπίων σὺ κατοικεῖς· τοὺς λόγους αὐτῶν μὴ φοβηθῇς, καὶ ἀπὸ προσώπου αὐτῶν μὴ ἐκστῇς, διότι οἶκος παραπικραίνων ἐστί.
\vs{7}Καὶ λαλήσεις τοὺς λόγους μου πρὸς αὐτοὺς, ἐὰν ἆρα ἀκούσωσιν ἢ πτοηθῶσιν, ὅτι οἶκος παραπικραίνων ἐστί.

\vs{8}Καὶ σὺ, υἱὲ ἀνθρώπου, ἄκουε τοῦ λαλοῦντος πρὸς σὲ, μὴ γίνου παραπικραίνων, καθὼς ὁ οἶκος ὁ παραπικραίνων· χάνε τὸ στόμα σου, καὶ φάγε ὃ ἐγὼ δίδωμί σοι.
\vs{9}Καὶ ἴδον, καὶ ἰδοὺ χεὶρ ἐκτεταμένη πρὸς μὲ, καὶ ἐν αὐτῇ κεφαλὶς βιβλίου.
\vs{10}Καὶ ἀνείλησεν αὐτὴν ἐνώπιόν μου, καὶ ἦν ἐν αὐτῇ γεγραμμένα τὰ ἔμπροσθεν καὶ τὰ ὀπίσω· καὶ ἐγέγραπτο θρῆνος καὶ μέλος καὶ οὐαί.

\ch{3}
Καὶ εἶπε πρὸς μὲ, υἱὲ ἀνθρώπου, κατάφαγε τὴν κεφαλίδα ταύτην, καὶ πορεύθητι καὶ λάλησον τοῖς υἱοῖς Ἰσραήλ·
\vs{2}Καὶ διήνοιξε τὸ στόμα μου, καὶ ἐψώμισέ με τὴν κεφαλίδα.

\vs{3}Καὶ εἶπε πρὸς μὲ, υἱὲ ἀνθρώπου, τὸ στόμα σου φάγεται, καὶ ἡ κοιλία σου πλησθήσεται τῆς κεφαλίδος ταύτης τῆς δεδομένης εἰς σέ· καὶ ἔφαγον αὐτὴν, καὶ ἐγένετο ἐν τῷ στόματί μου, ὡς μέλι γλυκάζον.

\vs{4}Καὶ εἶπε πρὸς μὲ, υἱὲ ἀνθρώπου, βάδιζε καὶ εἴσελθε πρὸς τὸν οἶκον τοῦ Ἰσραὴλ, καὶ λάλησον τοὺς λόγους μου πρὸς αὐτούς·
\vs{5}Διότι οὐ πρὸς λαὸν βαθύγλωσσον σὺ ἐξαποστέλλῃ πρὸς τὸν οἶκον τοῦ Ἰσραήλ·
\vs{6}Οὐδὲ πρὸς λαοὺς πολλοὺς ἀλλοφώνους ἢ ἀλλογλώσσους, οὐδὲ στιβαροὺς τῇ γλωσσῃ ὄντας, ὧν οὐκ ἀκούσῃ τοὺς λόγους· καὶ εἰ πρὸς τοιούτους ἐξαπέστειλά σε, οὗτοι ἂν εἰσήκουσάν σου.
\vs{7}Ὁ δὲ οἶκος τοῦ Ἰσραὴλ οὐ μὴ θελήσουσιν εἰσακοῦσαί σου, διότι οὐ βούλονται εἰσακούειν μου, ὅτι πᾶς ὁ οἶκος Ἰσραὴλ φιλόνεικοί εἰσι καὶ σκληροκάρδιοι.
\vs{8}Καὶ ἰδοὺ δέδωκα τὸ πρόσωπόν σου δυνατὸν κατέναντι τῶν προσώπων αὐτῶν, καὶ τὸ νῖκός σου κατισχύσω κατέναντι τοῦ νίκους αὐτῶν·
\vs{9}Καὶ ἔσται διαπαντὸς κραταιότερον πέτρας· μὴ φοβηθῇς ἀπʼ αὐτῶν, μηδὲ πτοηθῇς ἀπὸ προσώπου αὐτῶν, διότι οἶκος παραπικραίνων ἐστί.

\vs{10}Καὶ εἶπε πρὸς μὲ, υἱὲ ἀνθρώπου, πάντας τοὺς λόγους οὓς λελάληκα μετὰ σοῦ, λάβε εἰς τὴν καρδίαν σου, καὶ τοῖς ὠσί σου ἄκουε.
\vs{11}Καὶ βάδιζε, εἴσελθε εἰς τὴν αἰχμαλωσίαν πρὸς τοὺς υἱοὺς τοῦ λαοῦ σοῦ, καὶ λαλήσεις πρὸς αὐτοὺς, καὶ ἐρεῖς πρὸς αὐτοὺς, τάδε λέγει Κύριος, ἐὰν ἄρα ἀκούσωσιν, ἐὰν ἄρα ἐνδῶσι.

\vs{12}Καὶ ἀνέλαβέ με πνεῦμα, καὶ ἤκουσα κατόπισθέν μου φωνὴν σεισμοῦ μεγάλου, εὐλογημένη ἡ δόξα Κυρίου ἐκ τοῦ τόπου αὐτοῦ.
\vs{13}Καὶ ἴδον φωνὴν τῶν πτερύγων τῶν ζώων πτερυσσομένων ἑτέρα πρὸς τὴν ἑτέραν, καὶ φωνὴ τῶν τροχῶν ἐχομένη αὐτῶν, καὶ φωνὴ τοῦ σεισμοῦ.
\vs{14}Καὶ τὸ πνεῦμα ἐξῇρέ με, καὶ ἀνέλαβέ με, καὶ ἐπορεύθην ἐν ὁρμῇ τοῦ πνεύματός μου, καὶ χεὶρ Κυρίου ἐγένετο ἐπʼ ἐμὲ κραταιά.

\vs{15}Καὶ εἰσῆλθον εἰς τὴν αἰχμαλωσίαν μετέωρος, καὶ περιῆλθον τοὺς κατοικοῦντας ἐπὶ τοῦ ποταμοῦ τοῦ Χοβὰρ τοὺς ὄντας ἐκεῖ· καὶ ἐκάθισα ἐκεῖ ἑπτὰ ἡμέρας, ἀναστρεφόμενος ἐν μέσῳ αὐτῶν.

\vs{16}Καὶ ἐγένετο μετὰ τὰς ἑπτὰ ἡμέρας λόγος Κυρίου πρὸς μὲ, λέγων,
\vs{17}υἱὲ ἀνθρώπου, σκοπὸν δέδωκά σε τῷ οἴκῳ Ἰσραὴλ, καὶ ἀκούσῃ ἐκ στόματός μου λόγον, καὶ διαπειλήσῃ αὐτοῖς παρʼ ἐμοῦ.
\vs{18}Ἐν τῷ λέγειν με τῷ ἀνόμῳ, θανάτῳ θανατωθήσῃ· καὶ οὐ διεστείλω αὐτῷ τοῦ διαστείλασθαι τῷ ἀνόμῳ, ἀποστρέψαι ἀπὸ τῶν ὁδῶν αὐτοῦ, τοῦ ζῆσαι αὐτόν· ὁ ἄνομος ἐκεῖνος τῇ ἀδικίᾳ αὐτοῦ ἀποθανεῖται, καὶ τὸ αἷμα αὐτοῦ ἐκ τῆς χειρός σου ἐκζητήσω.
\vs{19}Καὶ σὺ ἐὰν διαστείλῃ τῷ ἀνόμῳ, καὶ μὴ ἀποστρέψῃ ἀπὸ τῆς ἀνομίας αὐτοῦ, καὶ ἀπὸ τῆς ὁδοῦ αὐτοῦ, ὁ ἄνομος ἐκεῖνος ἐν τῇ ἀδικίᾳ αὐτοῦ ἀποθανεῖται, καὶ σὺ τὴν ψυχήν σου ῥύσῃ.

\vs{20}Καὶ ἐν τῷ ἀποστρέφειν δίκαιον ἀπὸ τῶν δικαιοσυνῶν αὐτοῦ, καὶ ποιήσει παράπτωμα, καὶ δώσω τὴν βάσανον εἰς πρόσωπον αὐτοῦ, αὐτὸς ἀποθανεῖται, ὅτι οὐ διεστείλω αὐτῷ· καὶ ἐν ταῖς ἁμαρτίαις αὐτοῦ ἀποθανεῖται, διότι οὐ μὴ μνησθῶσιν αἱ δικαιοσύναι αὐτοῦ, καὶ τὸ αἷμα αὐτοῦ ἐκ τῆς χειρός σου ἐκζητήσω.
\vs{21}Σὺ δὲ ἐὰν διαστείλῃ τῷ δικαίῳ τοῦ μὴ ἁμαρτεῖν, καὶ αὐτὸς μὴ ἁμάρτῃ, ὁ δίκαιος ζωῇ ζήσεται, ὅτι διεστείλω αὐτῷ, καὶ σὺ τὴν σεαυτοῦ ψυχὴν ῥύσῃ.

\vs{22}Καὶ ἐγένετο ἐπʼ ἐμὲ χεὶρ Κυρίου· καὶ εἶπε πρὸς μὲ, ἀνάστηθι, καὶ ἔξελθε εἰς τὸ πεδίον, καὶ ἐκεῖ λαληθήσεται πρὸς σέ.

\vs{23}Καὶ ἀνέστην καὶ ἐξῆλθον πρὸς τὸ πεδίον· καὶ ἰδοὺ ἐκεῖ δόξα Κυρίου εἱστήκει, καθὼς ἡ ὅρασις, καὶ καθὼς ἡ δόξα Κυρίου, ἣν ἴδον ἐπὶ τοῦ ποταμοῦ τοῦ Χοβάρ· καὶ πίπτω ἐπὶ πρόσωπόν μου·
\vs{24}Καὶ ἦλθεν ἐπʼ ἐμὲ πνεῦμα, καὶ ἔστησέ με ἐπὶ τοὺς πόδας μου, καὶ ἐλάλησε πρὸς μὲ, καὶ εἶπε μοι, εἴσελθε, καὶ ἐγκλείσθητι ἐν μέσῳ τοῦ οἴκου σου.
\vs{25}Καὶ σὺ, υἱὲ ἀνθρώπου, ἰδοὺ δέδονται ἐπὶ σὲ δεσμοὶ, καὶ δήσουσί σε ἐν αὐτοῖς, καὶ οὐ μὴ ἐξέλθῃς ἐκ μέσου αὐτῶν.
\vs{26}Καὶ τὴν γλῶσσάν σου συνδήσω, καὶ ἀποκωφωθήσῃ, καὶ οὐκ ἔσῃ αὐτοῖς εἰς ἄνδρα ἐλέγχοντα, διότι οἶκος παραπικραίνων ἐστί.
\vs{27}Καὶ ἐν τῷ λαλεῖν με πρὸς σὲ, ἀνοίξω τὸ στόμα σου, καὶ ἐρεῖς πρὸς αὐτοὺς τάδε λέγει Κύριος, ὁ ἀκούων ἀκουέτω, καὶ ὁ ἀπειθῶν ἀπειθείτω, διότι οἶκος παραπικραίνων ἐστί.

\ch{4}
Καὶ σὺ, υἱὲ ἀνθρώπου, λάβε σεαυτῷ πλίνθον, καὶ θήσεις αὐτὴν πρὸ προσώπου σου, καὶ διαγράψεις ἐπʼ αὐτὴν πόλιν τὴν Ἱερουσαλήμ.
\vs{2}Καὶ δώσεις ἐπʼ αὐτὴν περιοχὴν, καὶ οἰκοδομήσεις ἐπʼ αὐτὴν προμαχῶνας, καὶ περιβαλεῖς ἐπʼ αὐτὴν χάρακα, καὶ δώσεις ἐπʼ αὐτὴν παρεμβολὰς, καὶ τάξεις τὰς βελοστάσεις κύκλῳ.
\vs{3}Καὶ σὺ λάβε σεαυτῷ τήγανον σιδηροῦν, καὶ θήσεις αὐτὸ τοῖχον σιδηροῦν ἀναμέσον σου καὶ ἀναμέσον τῆς πόλεως, καὶ ἑτοιμάσεις τὸ πρόσωπόν σου ἐπʼ αὐτὴν, καὶ ἔσται ἐν συνκλεισμῷ, καὶ συγκλείσεις αὐτήν· σημεῖόν ἐστι τοῦτο τοῖς υἱοῖς Ἰσραήλ.

\vs{4}Καὶ σὺ κοιμηθήσῃ ἐπὶ τὸ πλευρόν σου τὸ ἀριστερὸν, καὶ θήσεις τὰς ἀδικίας τοῦ οἴκου Ἰσραὴλ ἐπʼ αὐτοὺ, κατὰ ἀριθμὸν τῶν ἡμερῶν πεντήκοντα καὶ ἑκατὸν ἃς κοιμηθήσῃ ἐπʼ αὐτοῦ, καὶ λήψῃ τὰς ἀδικίας αὐτῶν.
\vs{5}Καὶ ἐγὼ δέδωκά σοι τὰς ἀδικίας αὐτῶν εἰς ἀριθμὸν ἡμερῶν, ἐννενήκοντα καὶ ἑκατὸν ἡμέρας, καὶ λήψῃ τὰς ἀδικίας τοῦ οἴκου Ἰσραήλ.
\vs{6}Καὶ συντελέσεις ταῦτα, καὶ κοιμηθήσῃ ἐπὶ τὸ πλευρόν σου τὸ δεξιὸν, καὶ λήψῃ τὰς ἀδικίας τοῦ οἴκου Ἰούδα τεσσαράκοντα ἡμέρας, ἡμέραν εἰς ἐνιαυτὸν τέθεικά σοι.

\vs{7}Καὶ εἰς τὸν συγκλεισμὸν Ἱερουσαλὴμ ἑτοιμάσεις τὸ πρόσωπόν σου, καὶ τὸν βραχίονά σου στερεώσεις, καὶ προφητεύσεις ἐπʼ αὐτήν.
\vs{8}Καὶ ἐγὼ ἰδοὺ δέδωκα ἐπὶ σὲ δεσμοὺς, καὶ μὴ στραφῇς ἀπὸ τοῦ πλευροῦ σου ἐπὶ τὸ πλευρόν σου, ἕως οὗ συντελεσθῶσιν ἡμέραι τοῦ συγκλεισμοῦ σου.

\vs{9}Καὶ σὺ λάβε σεαυτῷ πυροὺς, καὶ κριθὰς, καὶ κύαμον, καὶ φακὸν, καὶ κέγχρον, καὶ ὀλύραν, καὶ ἐμβαλεῖς αὐτὰ εἰς ἄγγος ἓν ὀστράνικον, καὶ ποιήσεις αὐτὰ σεαυτῷ εἰς ἄρτους· καὶ κατὰ ἀριθμὸν τῶν ἡμερῶν, ἃς σὺ καθεύδεις ἐπὶ τοῦ πλευροῦ σου, ἐννενήκοντα καὶ ἑκατὸν ἡμέρας φάγεσαι αὐτά.
\vs{10}Καὶ τὸ βρῶμά σου φάγεσαι ἐν σταθμῷ, εἴκοσι σίκλους τὴν ἡμέραν, ἀπὸ καιροῦ ἕως καιροῦ φάγεσαι αὐτά.
\vs{11}Καὶ ὕδωρ ἐν μέτρῳ πίεσαι, καὶ τὸ ἕκτον τοῦ εἲν ἀπὸ καιροῦ ἕως καιροῦ πίεσαι.
\vs{12}Καὶ ἐγκρυφίαν κρίθινον φάγεσαι αὐτὰ, ἐν βολβίτοις κόπρου ἀνθρωπίνης ἐγκρύψεις αὐτὰ κατʼ ὀφθαλμοὺς αὐτῶν.

\vs{13}Καὶ ἐρεῖς, τάδε λέγει Κύριος ὁ Θεὸς τοῦ Ἰσραὴλ, οὕτως φάγονται οἱ υἱοὶ τοῦ Ἰσραὴλ ἀκάθαρτα ἐν τοῖς ἔθνεσι.
\vs{14}Καὶ εἶπα, μηδαμῶς Κύριε Θεὲ Ἰσραήλ· εἰ ἡ ψυχή μου οὐ μεμίανται ἐν ἀκαθαρσίᾳ, καὶ θνησιμαῖον καὶ θηριάλωτον οὐ βέβρωκα ἀπὸ γενέσεώς μου ἕως τοῦ νῦν, οὐδὲ εἰσελήλυθεν εἰς τὸ στόμα μου πᾶν κρέας ἕωλον.

\vs{15}Καὶ εἶπε πρὸς μὲ, ἰδοὺ δέδωκά σοι βόλβιτα βοῶν ἀντὶ τῶν βολβίτων τῶν ἀνθρωπίνων, καὶ ποιήσεις τοὺς ἄρτους σου ἐπʼ αὐτῶν.

\vs{16}Καὶ εἶπε πρὸς μὲ, υἱὲ ἀνθρώπου, ἰδοὺ ἐγὼ συντρίβω στήριγμα ἄρτου ἐν Ἱερουσαλὴμ, καὶ φάγονται ἄρτον ἐν σταθμῷ καὶ ἐν ἐνδείᾳ, καὶ ὕδωρ ἐν μέτρῳ, καὶ ἐν ἀφανισμῷ πίονται·
\vs{17}Ὅπως ἐνδεεῖς γένωνται ἄρτου καὶ ὕδατος· καὶ ἀφανισθήσεται ἄνθρωπος καὶ ἀδελφὸς αὐτοῦ, καὶ ἐντακήσονται ἐν ταῖς ἀδικίαις αὐτῶν.

\ch{5}
Καὶ σὺ, υἱὲ ἀνθρώπου, λάβε σεαυτῷ ῥομφαίαν ὀξεῖαν ὑπὲρ ξυρὸν κουρέως, κτήσῃ αὐτὴν σεαυτῷ, καὶ ἐπάξεις αὐτὴν ἐπὶ τὴν κεφαλήν σου, καὶ ἐπὶ τὸν πώγωνά σου· καὶ λήψῃ ζυγὸν σταθμίων, καὶ διαστήσεις αὐτούς.
\vs{2}Τὸ τέταρτον ἐν πυρὶ ἀνακαύσεις ἐν μέσῃ τῇ πόλει κατὰ τὴν πλήρωσιν τῶν ἡμερῶν τοῦ συγκλεισμοῦ· καὶ λήψῃ τὸ τέταρτον, καὶ κατακαύσεις αὐτὸ ἐν μέσῳ αὐτῆς· καὶ τὸ τέταρτον κατακόψεις ἐν ῥομφαίᾳ κύκλῳ αὐτῆς· καὶ τὸ τέταρτον διασκορπιεῖς τῷ πνεύματι· καὶ μάχαιραν ἐκκενώσω ὀπίσω αὐτῶν.

\vs{3}Καὶ λήψῃ ἐκεῖθεν ὀλίγους ἐν ἀριθμῷ, καὶ συμπεριλήψῃ αὐτοὺς τῇ ἀναβολῇ σου.
\vs{4}Καὶ ἐκ τούτων λήψῃ ἔτι, καὶ ῥίψεις αὐτοὺς εἰς μέσον τοῦ πυρὸς, καὶ κατακαύσεις αὐτοὺς ἐν πυρί· ἐξ αὐτῆς ἐξελεύσεται πῦρ· καὶ ἐρεῖς παντὶ οἴκῳ Ἰσραὴλ,

\vs{5}Τάδε λέγει Κύριος, αὕτη ἡ Ἱερουσαλὴμ, ἐν μέσῳ τῶν ἐθνῶν τέθεικα αὐτὴν, καὶ τὰς κύκλῳ αὐτῆς χώρας.
\vs{6}Καὶ ἐρεῖς τὰ δικαιώματά μου τῇ ἀνόμῳ ἐκ τῶν ἐθνῶν, καὶ τὰ νόμιμά μου τῶν χωρῶν τῶν κύκλῳ αὐτῆς· διότι τὰ δικαιώματά μου ἀπώσαντο, καὶ ἐν τοῖς νομίμοις μου οὐκ ἐπορεύθησαν ἐν αὐτοῖς.

\vs{7}Διατοῦτο τάδε λέγει Κύριος, ἀνθʼ ὧν ἡ ἀφορμὴ ὑμῶν ἐκ τῶν ἐθνῶν τῶν κύκλῳ ὑμῶν, καὶ ἐν τοῖς νομίμοις μου οὐκ ἐπορεύθητε, καὶ τὰ δικαιώματά μου οὐκ ἐποιήσατε, ἀλλʼ οὐδὲ κατὰ τὰ δικαιώματα τῶν ἐθνῶν τῶν κύκλῳ ὑμῶν οὐ πεποιήκατε,
\vs{8}διατοῦτο τάδε λέγει Κύριος, ἰδοὺ ἐγὼ ἐπὶ σὲ, καὶ ποιήσω ἐν μέσῳ σου κρίμα ἐνώπιον τῶν ἐθνῶν.
\vs{9}Καὶ ποιήσω ἐν σοὶ ἃ οὐ πεποίηκα, καὶ ἃ οὐ ποιήσω ὅμοια αὐτοῖς ἔτι κατὰ πάντα τὰ βδελύγματά σου.
\vs{10}Διατοῦτο πατέρες φάγονται τέκνα ἐν μέσῳ σου, καὶ τέκνα φάγονται πατέρας· καὶ ποιήσω ἐν σοὶ κρίματα, καὶ διασκορπιῶ πάντας τοὺς καταλοίπους σου εἰς πάντα ἄνεμον.

\vs{11}Διατοῦτο, ζῶ ἐγὼ, λέγει Κύριος, ἦ μὴν ἀνθʼ ὧν τὰ ἅγιά μου ἐμίανας ἐν πᾶσι τοῖς βδελύγμασί σου, κᾀγὼ ἀπώσομαί σε, οὐ φείσεταί μου ὁ ὀφθαλμὸς, κᾀγὼ οὐκ ἐλεήσω.
\vs{12}Τὸ τέταρτόν σου ἐν θανάτῳ ἀναλωθήσεται, καὶ τὸ τέταρτόν σου ἐν λιμῷ συντελεσθήσεται ἐν μέσῳ σου, καὶ τὸ τέταρτόν σου εἰς πάντα ἄνεμον σκορπιῶ αὐτοὺς, καὶ τὸ τέταρτόν σου ἐν ῥομφαίᾳ πεσοῦνται κύκλῳ σου, καὶ μάχαιραν ἐκκενώσω ὀπίσω αὐτῶν.

\vs{13}Καὶ συντελεσθήσεται ὁ θυμός μου, καὶ ἡ ὀργή μου ἐπʼ αὐτοὺς, καὶ ἐπιγνώσῃ, διότι ἐγὼ Κύριος λελάληκα ἐν ζήλῳ μου, ἐν τῷ συντελέσαι με τὴν ὀργήν μου ἐπʼ αὐτούς.

\vs{14}Καὶ θήσομαί σε εἰς ἔρημον, καὶ τὰς θυγατέρας σου κύκλῳ σου ἐνώπιον παντὸς διοδεύοντος.
\vs{15}Καὶ ἔσῃ στενακτὴ καὶ δηλαϊστὴ ἐν τοῖς ἔθνεσι τοῖς κύκλῳ σου, ἐν τῷ ποιῆσαί με ἐν σοὶ κρίματα ἐν ἐκδικήσει θυμοῦ μου· ἐγὼ Κύριος λελάληκα.
\vs{16}Καὶ ἐν τῷ ἀποστεῖλαί με βολίδας τοῦ λιμοῦ ἐπʼ αὐτούς, καὶ ἔσονται εἰς ἔκλειψιν, καὶ συντρίψω στήριγμα ἄρτου σου.
\vs{17}Καὶ ἐξαποστελῶ ἐπὶ σὲ λιμὸν καὶ θηρία πονηρὰ, καὶ τιμωρήσομαί σε, καὶ θάνατος καὶ αἷμα διελεύσονται ἐπὶ σὲ, καὶ ῥομφαίαν ἐπάξω ἐπὶ σὲ κυκλόθεν· ἐγὼ Κύριος λελάληκα.

\ch{6}
Καὶ ἐγένετο λόγος Κυρίου πρὸς μὲ, λέγων,
\vs{2}υἱὲ ἀνθρώπου, στήρισον τὸ πρόσωπόν σου ἐπὶ τὰ ὄρη Ἰσραὴλ, καὶ προφήτευσον ἐπʼ αὐτὰ,
\vs{3}καὶ ἐρεῖς,

Τὰ ὄρη Ἰσραὴλ ἀκούσατε λόγον Κυρίου· τάδε λέγει Κύριος τοῖς ὄρεσι καὶ τοῖς βουνοῖς καὶ ταῖς φάραγξι καὶ ταῖς νάπαις, ἰδοὺ ἐγὼ ἐπάγω ἐφʼ ὑμᾶς ῥομφαίαν, καὶ ἐξολοθρευθήσεται τὰ ὑψηλὰ ὑμῶν.
\vs{4}Καὶ συντριβήσονται τὰ θυσιαστήρια ὑμῶν, καὶ τὰ τεμένη ὑμῶν, καὶ καταβαλῶ τραυματίας ὑμῶν ἐνώπιον τῶν εἰδώλων ὑμῶν,
\vs{5}καὶ διασκορπιῶ τὰ ὀστᾶ ὑμῶν κύκλῳ τῶν θυσιαστηρίων ὑμῶν,
\vs{6}καὶ ἐν πάσῃ τῇ κατοικίᾳ ὑμῶν· αἱ πόλεις ἐξερημωθήσονται, καὶ τὰ ὑψηλὰ ἀφανισθήσεται, ὅπως ἐξολοθρευθῇ τὰ θυσιαστήρια ὑμῶν, καὶ συντριβήσονται τὰ εἴδωλα ὑμῶν, καὶ ἐξαρθῇ τὰ τεμένη ὑμῶν·
\vs{7}Καὶ πεσοῦνται τραυματίαι ἐν μέσῳ ὑμῶν, καὶ ἐπιγνώσεσθε ὅτι ἐγὼ Κύριος.

\vs{8}Ἐν τῷ γενέσθαι ἐξ ὑμῶν ἀνασωζομένους ἐκ ῥομφαίας ἐν τοῖς ἔθνεσι, καὶ ἐν τῷ διασκορπισμῷ ὑμῶν ἐν ταῖς χώραις,
\vs{9}καὶ μνησθήσονταί μου οἱ ἀνασωζόμενοι ἐξ ὑμῶν ἐν τοῖς ἔθνεσιν, οὗ ᾐχμαλωτεύθησαν ἐκεῖ· ὀμώμοκα τῇ καρδίᾳ αὐτῶν τῇ ἐκπορνευούσῃ ἀπʼ ἐμοῦ, καὶ τοῖς ὀφθαλμοῖς αὐτῶν τοῖς ἐκπορνεύουσιν ὀπίσω τῶν ἐπιτηδευμάτων αὐτῶν· καὶ κόψονται πρόσωπα αὐτῶν ἐν πᾶσι τοῖς βδελύγμασιν αὐτῶν·
\vs{10}Καὶ ἐπιγνώσονται διότι ἐγὼ Κύριος λελάληκα.

\vs{11}Τάδε λέγει Κύριος, κρότησον τῇ χειρὶ καὶ ψόφησον τῷ ποδὶ, καὶ εἰπὸν, εὖγε εὖγε, ἐπὶ πᾶσι τοῖς βδελύγμασιν οἴκου Ἰσραήλ· ἐν ῥομφαίᾳ καὶ ἑν θανάτῳ καὶ ἐν λιμῷ πεσοῦνται.
\vs{12}Ὁ ἐγγὺς ἐν ῥομφαίᾳ πεσεῖται, ὁ δὲ μακρὰν ἐν θανάτῳ τελευτήσει· καὶ ὁ περιεχόμενος ἐν λιμῷ συντελεσθήσεται· καὶ συντελέσω τὴν ὀργήν μου ἐπʼ αὐτούς.

\vs{13}Καὶ γνώσεσθε διότι ἐγὼ Κύριος, ἐν τῷ εἶναι τοὺς τραυματίας ὑμῶν ἐν μέσῳ τῶν εἰδώλων ὑμῶν κύκλῳ τῶν θυσιαστηρίων ὑμῶν· ἐπὶ πάντα βουνὸν ὑψηλὸν, καὶ ὑποκάτω δένδρου συσκίου, οὗ ἔδωκαν ἐκεῖ ὀσμὴν εὐωδίας πᾶσι τοῖς εἰδώλοις αὐτῶν.
\vs{14}Καὶ ἐκτενῶ τὴν χεῖρά μου ἐπʼ αὐτούς, καὶ θήσομαι τὴν γῆν εἰς ἀφανισμὸν καὶ εἰς ὄλεθρον ἀπὸ τῆς ἐρήμου Δεβλαθὰ ἐκ πάσης τῆς κατοικεσίας αὐτῶν· ἐπιγνώσεσθε ὅτι ἐγὼ Κύριος.

\ch{7}
Καὶ ἐγένετο λόγος Κυρίου πρὸς μὲ, λέγων,
\vs{2}καὶ σὺ υἱὲ ἀνθρώπου εἰπόν, τάδε λέγει Κύριος,

Τῇ γῇ τοῦ Ἰσραὴλ πέρας ἥκει, τὸ πέρας ἥκει ἐπὶ τὰς τέσσαρας πτέρυγας τῆς γῆς.
\vs{3}Ἥκει τὸ πέρας
\vs{4}ἐπὶ σὲ τὸν κατοικοῦντα τὴν γῆν· ἥκει ὁ καιρὸς, ἤγγικεν ἡ ἡμέρα, οὐ μετὰ θορύβων, οὐδὲ μετὰ ὠδίνων.

\vs{5}Νῦν ἐγγύθεν ἐκχεῶ τὴν ὀργήν μου ἐπὶ σὲ, καὶ συντελέσω τὸν θυμόν μου ἐν σοὶ, καὶ κρινῶ σε ἐν ταῖς ὁδοῖς σου, καὶ δώσω ἐπὶ σὲ πάντα τὰ βδελύγματά σου,
\vs{6}οὐ φείσεται ὁ ὀφθαλμός μου, οὐδὲ μὴ ἐλεήσω· διότι τὰς ὁδούς σου ἐπὶ σὲ δώσω, καὶ τὰ βδελύγματά σου ἐν μέσῳ σου ἔσονται, καὶ ἐπιγνώσῃ, διότι ἐγώ εἰμι Κύριος ὁ τύπτων.
\vs{7}Νῦν τὸ πέρας πρὸς σὲ, καὶ ἀποστελῶ ἐγὼ ἐπὶ σὲ, καὶ ἐκδικήσω ἐν ταῖς ὁδοῖς σου, καὶ δώσω ἐπὶ σὲ πάντα τὰ βδελύγματά σου,
\vs{8}οὐ φείσεται ὁ ὀφθαλμός μου, οὐδὲ μὴ ἐλεήσω· διότι τὴν ὁδόν σου ἐπὶ σὲ δώσω, καὶ τὰ βδελύγματά σου ἐν μέσῳ σου ἔσται, καὶ ἐπιγνώσῃ διότι ἐγὼ Κύριος.

\vs{9}Διότι τάδε λέγει Κύριος, ἰδοὺ τὸ πέρας ἥκει.

\vs{10}Ἰδοὺ ἡ ἡμέρα Κυρίου· εἰ καὶ ἡ ῥάβδος ἤνθηκεν, ἡ ὕβρις ἐξανέστηκε,
\vs{11}καὶ συντρίψει στήριγμα ἀνόμου, καὶ οὐ μετὰ θορύβου, οὐδὲ μετὰ σπουδῆς.
\vs{12}Ἥκει ὁ καιρὸς, ἰδοὺ ἡ ἡμέρα· ὁ κτώμενος μὴ χαιρέτω, καὶ ὁ πωλῶν μὴ θρηνείτω·
\vs{13}Διότι ὁ κτώμενος πρὸς τὸν πωλοῦντα οὐκέτι μὴ ἐπιστρέψει, καὶ ἄνθρωπος ἐν ὀφθαλμῷ ζωῆς αὐτοῦ οὐ κρατήσει.
\vs{14}Σαλπίσατε ἐν σάλπιγγι, καὶ κρίνατε τὰ σύμπαντα.
\vs{15}Ὁ πόλεμος ἐν ῥομφαίᾳ ἔξωθεν, καὶ ὁ λιμὸς καὶ ὁ θάνατος ἔσωθεν· ὁ ἐν τῷ πεδίῳ ἐν ῥομφαίᾳ τελευτήσει, τοὺς δʼ ἐν τῇ πόλει λιμὸς καὶ θάνατος συντελέσει.

\vs{16}Καὶ ἀνασωθήσονται οἱ ἀνασωζόμενοι ἐξ αὐτῶν, καὶ ἔσονται ἐπὶ τῶν ὀρέων· καὶ πάντας ἀποκτενῶ, ἔκαστον ἐν ταῖς ἀδικίαις αὐτοῦ.
\vs{17}Πᾶσαι χεῖρες ἐκλυθήσονται, καὶ πάντες μηροὶ μολυνθήσονται ὑγρασίᾳ.
\vs{18}Καὶ περιζώσονται σάκκους, καὶ καλύψει αὐτοὺς θάμβος· καὶ ἐπὶ πᾶν πρόσωπον αἰσχύνη ἐπʼ αὐτοὺς, καὶ ἐπὶ πᾶσαν κεφαλὴν φαλάκρωμα.
\vs{19}Τὸ ἀργύριον αὐτῶν ῥιφήσεται ἐν ταῖς πλατείαις, καὶ τὸ χρυσίον αὐτῶν ὑπεροφθήσεται· αἱ ψυχαὶ αὐτῶν οὐ μὴ ἐμπλησθῶσι, καὶ αἱ κοιλίαι αὐτῶν οὐ μὴ πληρωθῶσι, διότι βάσανος τῶν ἀδικιῶν αὐτῶν ἐγένετο.
\vs{20}Ἐκλεκτὰ κόσμου εἰς ὑπερηφανίαν ἔθεντο αὐτὰ, καὶ εἰκόνας τῶν βδελυγμάτων αὐτῶν ἐποίησαν ἐξ αὐτῶν· ἕνεκεν τούτου δέδωκα αὐτὰ αὐτοῖς εἰς ἀκαθαρσίαν.

\vs{21}Καὶ παραδώσω αὐτὰ εἰς χεῖρας ἀλλοτρίων, τοῦ διαρπάσαι αὐτὰ, καὶ τοῖς λοιμοῖς τῆς γῆς εἰς σκῦλα, καὶ βεβηλώσουσιν αὐτά.
\vs{22}Καὶ ἀποστρέψω τὸ πρόσωπόν μου ἀπʼ αὐτῶν, καὶ μιανοῦσι τῆν ἐπισκοπήν μου, καὶ εἰσελεύσονται εἰς αὐτὰ ἀφυλάκτως, καὶ βεβηλώσουσιν αὐτά.
\vs{23}Καὶ ποιήσουσι φυρμόν· διότι ἡ γῆ πλήρης λαῶν, καὶ ἡ πόλις πλήρης ἀνομίας.
\vs{24}Καὶ ἀποστρέψω τὸ φρύαγμα τῆς ἰσχύος αὐτῶν, καὶ μιανθήσεται τὰ ἅγια αὐτῶν.
\vs{25}Καὶ ἐξιλασμὸς ἥξει, καὶ ζητήσει εἰρήνην, καὶ οὐκ ἔσται.
\vs{26}Οὐαὶ ἐπὶ οὐαὶ ἔσται· καὶ ἀγγελία ἐπὶ ἀγγελίαν ἔσται· καὶ ζητηθήσεται ὅρασις ἐκ προφήτου, καὶ νόμος ἀπολεῖται ἐξ ἱερέως, καὶ βουλὴ ἐκ πρεσβυτέρων.
\vs{27}Ἄρχων ἐνδύσεται ἀφανισμὸν, καὶ αἱ χεῖρες τοῦ λαοῦ τῆς γῆς παραλυθήσονται· κατὰ τὰς ὁδοὺς αὐτῶν ποιήσω αὐτοῖς, καὶ ἐν τοῖς κρίμασιν αὐτῶν ἐκδικήσω αὐτοὺς, καὶ γνώσονται ὅτι ἐγὼ Κύριος.

\ch{8}
Καὶ ἐγένετο ἐν τῷ ἕκτῳ ἔτει ἐν τῷ πέμπτῳ μηνὶ, πέμπτῃ τοῦ μηνὸς, ἐγὼ ἐκαθήμην ἐν τῷ οἴκῳ, καὶ οἱ πρεσβύτεροι Ἰούδα ἐκάθηντο ἐνώπιόν μου· καὶ ἐγένετο ἐπʼ ἐμὲ χεὶρ Κυρίου.
\vs{2}Καὶ ἴδον, καὶ ἰδοὺ ὁμοίωμα ἀνδρὸς, ἀπὸ τῆς ὀσφύος αὐτοῦ καὶ ἕως κάτω πῦρ, καὶ ἀπὸ τῆς ὀσφύος αὐτοῦ, ὑπεράνω αὐτοῦ ὡς ὅρασις ἠλέκτρου.
\vs{3}Καὶ ἐξέτεινεν ὁμοίωμα χειρὸς, καὶ ἀνέλαβέ με τῆς κορυφῆς μου, καὶ ἀνέλαβέ με πνεῦμα ἀναμέσον τῆς γῆς καὶ ἀναμέσον τοῦ οὐρανοῦ, καὶ ἤγαγέ με εἰς Ἱερουσαλὴμ ἐν ὁράσει Θεοῦ ἐπὶ τὰ πρόθυρα τῆς πύλης τῆς βλεπούσης εἰς Βοῤῥᾶν, οὗ ἦν ἡ στήλη τοῦ κτωμένου·
\vs{4}Καὶ ἰδοὺ ἦν ἐκεῖ δόξα Κυρίου Θεοῦ Ἰσραὴλ κατὰ τὴν ὅρασιν ἣν ἴδον ἐν τῷ πεδίῳ.

\vs{5}Καὶ εἶπε πρὸς μὲ, υἱὲ ἀνθρώπου, ἀνάβλεψον τοῖς ὀφθαλμοῖς σου πρὸς Βοῤῥᾶν· καὶ ἀνέβλεψα τοῖς ὀφθαλμοῖς μου πρὸς Βοῤῥᾶν, καὶ ἰδοὺ ἀπὸ Βοῤῥᾶ ἐπὶ τὴν πύλην τὴν πρὸς ἀνατολάς,
\vs{6}Καὶ εἶπε πρὸς μὲ, υἱὲ ἀνθρώπου, ἑώρακας τί οὗτοι ποιοῦσιν· ἀνομίας μεγάλας ποιοῦσιν ὧδε τοῦ ἀπέχεσθαι ἀπὸ τῶν ἁγίων μου· καὶ ἔτι ὄψει ἀνομίας μείζονας.

\vs{7}Καὶ εἰσήγαγέ με ἐπὶ τὰ πρόθυρα τῆς αὐλῆς.
\vs{8}Καὶ εἶπε πρὸς μὲ, υἱὲ ἀνθρώπου, ὄρυξον· καὶ ὤρυξα, καὶ ἰδοὺ θύρα.
\vs{9}Καὶ εἶπε πρὸς μὲ, εἴσελθε, καὶ ἴδε τὰς ἀνομίας ἃς οὗτοι ποιοῦσιν ὧδε.
\vs{10}Καὶ εἰσῆλθον, καὶ ἴδον, καὶ ἰδοὺ μάταια βδελύγματα, καὶ πάντα τὰ εἴδωλα οἴκου Ἰσραὴλ διαγεγραμμένα ἐπʼ αὐτοὺς κύκλῳ.
\vs{11}Καὶ ἑβδομήκοντα ἄνδρες ἐκ τῶν πρεσβυτέρων οἴκου Ἰσραὴλ, καὶ Ἰεχονίας ὁ τοῦ Σαφὰν ἐν μέσῳ αὐτῶν εἱστήκει πρὸ προσώπου αὐτῶν, καὶ ἕκαστος θυμιατήριον αὐτοῦ εἶχεν ἐν τῇ χειρὶ, καὶ ἡ ἀτμὶς τοῦ θυμιάματος ἀνέβαινε.
\vs{12}Καὶ εἶπε πρὸς μὲ, ἑώρακας, υἱὲ ἀνθρώπου, ἃ οἱ πρεσβύτεροι οἴκου Ἰσραὴλ ποιοῦσιν, ἕκαστος αὐτῶν ἐν τῷ κοιτῶνι τῷ κρυπτῷ αὐτῶν, διότι εἶπαν, οὐχ ὁρᾷ ὁ Κύριος, ἐγκαταλέλοιπε Κύριος τὴν γῆν.

\vs{13}Καὶ εἶπε πρὸς μὲ, ἔτι ὄψει ἀνομίας μείζονας ἃς οὗτοι ποιοῦσι.
\vs{14}Καὶ εἰσήγαγέ με ἐπὶ τὰ πρόθυρα τῆς πύλης οἴκου Κυρίου τῆς βλεπούσης πρὸς Βοῤῥᾶν· καὶ ἰδοὺ ἐκεῖ γυναῖκες καθήμεναι θρηνοῦσαι τὸν Θαμμούζ.
\vs{15}Καὶ εἶπε πρὸς μὲ, υἱὲ ἀνθρώπου ἑώρακας, καὶ ἔτι ὄψει ἐπιτηδεύματα μείζονα τούτων.

\vs{16}Καὶ εἰσήγαγέ με εἰς τὴν αὐλὴν οἴκου Κυρίου τὴν ἐσωτέραν, καὶ ἐπὶ τῶν προθύρων τοῦ ναοῦ Κυρίου ἀναμέσον τῶν αἰλὰμ, καὶ ἀναμέσον τοῦ θυσιαστηρίου, ὡς εἴκοσι ἄνδρες τὰ ὀπίσθια αὐτῶν πρὸς τὸν ναὸν τοῦ Κυρίου, καὶ τὰ πρόσωπα αὐτῶν ἀπέναντι, καὶ οὗτοι προσκυνοῦσι τῷ ἡλίῳ.
\vs{17}Καὶ εἶπε πρὸς μὲ, ἑώρακας υἱὲ ἀνθρώπου· μη μικρα τῷ οἴκῳ Ἰούδα τοῦ ποιεῖν τὰς ἀνομίας ἃς πεποιήκασιν ὧδε; διότι ἔπλησαν τὴν γῆν ἀνομίας· καὶ ἰδοὺ αὐτοὶ ὡς μυκτηρίζοντες.
\vs{18}Καὶ ἐγὼ ποιήσω αὐτοῖς μετὰ θυμοῦ· οὐ φείσεται ὁ ὀφθαλμός μου, οὐδὲ μὴ ἐλεήσω.

\ch{9}
Καὶ ἀνέκραγεν εἰς τὰ ὦτά μου φωνῇ μεγάλῃ, λέγων, ἤγγικεν ἡ ἐκδίκησις τῆς πόλεως· καὶ ἕκαστος εἶχε τὰ σκεύη τῆς ἐξολοθρεύσεως ἐν χειρὶ αὐτοῦ.
\vs{2}Καὶ ἰδοὺ ἓξ ἄνδρες ἤρχοντο ἀπὸ τῆς ὁδοῦ τῆς πύλης τῆς ὑψηλῆς τῆς βλεπούσης πρὸς Βοῤῥᾶν, καὶ ἑκάστου πέλυξ ἐν τῇ χειρὶ αὐτοῦ· καὶ εἷς ἀνὴρ ἐν μέσῳ αὐτῶν ἐνδεδυκὼς ποδήρη, καὶ ζώνη σαπφείρου ἐπὶ τῆς ὀσφύος αὐτοῦ, καὶ εἰσήλθοσαν καὶ ἔστησαν ἐχόμενοι τοῦ θυσιαστηρίου τοῦ χαλκοῦ.
\vs{3}Καὶ δόξα Θεοῦ τοῦ Ἰσραὴλ ἀνέβη ἀπὸ τῶν χερουβὶμ, ἡ οὖσα ἐπʼ αὐτῶν, εἰς τὸ αἴθριον τοῦ οἴκου·

\vs{4}Καὶ ἐκάλεσε τὸν ἄνδρα τὸν ἐνδεδυκότα τὸν ποδήρη, ὃς εἶχεν ἐπὶ τῆς ὀσφύος αὐτοῦ τὴν ζώνην, καὶ εἶπε πρὸς αὐτὸν, δίελθε μέσην Ἱερουσαλὴμ, καὶ δὸς σημεῖον ἐπὶ τὰ μέτωπα τῶν ἀνδρῶν τῶν καταστεναζόντων καὶ τῶν κατοδυνωμένων ἐπὶ πάσαις ταῖς ἀνομίαις ταῖς γινομέναις ἐν μέσῳ αὐτῶν.
\vs{5}Καὶ τούτοις εἶπεν ἀκούοντός μου, πορεύεσθε ὀπίσω αὐτοῦ εἰς τὴν πόλιν, καὶ κόπτετε, καὶ μὴ φείδεσθε τοῖς ὀφθαλμοῖς ὑμῶν, καὶ μὴ ἐλεήσητε.
\vs{6}Πρεσβύτερον καὶ νεανίσκον καὶ παρθένον καὶ νήπια καὶ γυναῖκας ἀποκτείνατε εἰς ἐξάλειψιν· ἐπὶ δὲ πάντας ἐφʼ οὕς ἐστι τὸ σημεῖον, μὴ ἐγγίσητε· ἀπὸ τῶν ἁγίων μου ἄρξασθε.

Καὶ ἤρξαντο ἀπὸ τῶν ἀνδρῶν τῶν πρεσβυτέρων οἳ ἦσαν ἔσω ἐν τῷ οἴκῳ.
\vs{7}Καὶ εἶπε πρὸς αὐτούς, μιάνατε τὸν οἶκον, καὶ πλήσατε τὰς ὁδοὺς νεκρῶν ἐκπορευόμενοι, καὶ κόπτετε.

\vs{8}Καὶ ἐγένετο ἐν τῷ κόπτειν αὐτοὺς, καὶ πίπτω ἐπὶ πρόσωπόν μου, καὶ ἀνεβόησα, καὶ εἶπα, οἴμοι Κύριε, ἐξαλείφεις σὺ τοὺς καταλοίπους τοῦ Ἰσραὴλ, ἐν τῷ ἐκχέαι σε τὸν θυμόν σου ἐπὶ Ἱερουσαλήμ;
\vs{9}Καὶ εἶπε πρὸς μὲ, ἀδικία τοῦ οἴκου Ἰσραὴλ καὶ Ἰούδα μεμεγάλυνται σφόδρα σφόδρα, ὅτι ἐπλήσθη ἡ γῆ λαῶν πολλῶν, καὶ ἡ πόλις ἐπλήσθη ἀδικίας καὶ ἀκαθαρσίας, ὅτι εἶπαν, ἐγκατέλιπε Κύριος τὴν γῆν, οὐκ ἐφορᾷ ὁ Κύριος.
\vs{10}Καὶ οὐ φείσεταί μου ὁ ὀφθαλμὸς, οὐδὲ μὴ ἐλεήσω, τὰς ὁδοὺς αὐτῶν εἰς κεφαλὰς αὐτῶν δέδωκα.

\vs{11}Καὶ ἰδοὺ ὁ ἀνὴρ ὁ ἐνδεδυκὼς τὸν ποδήρη, καὶ ἐζωσμένος τῇ ζώνῃ τὴν ὀσφὺν αὐτοῦ, καὶ ἀπεκρίνατο λέγων, πεποίηκα καθὼς ἐνετείλω μοι.

\ch{10}
Καὶ ἴδον, καὶ ἰδοὺ ἐπάνω τοῦ στερεώματος τοῦ ὑπὲρ κεφαλῆς τῶν χερουβὶμ, ὡς λίθος σαπφείρου ὁμοίωμα θρόνου ἐπʼ αὐτῶν.
\vs{2}Καὶ εἶπε πρὸς τὸν ἄνδρα τὸν ἐνδεδυκότα τὴν στολὴν, εἴσελθε εἰς τὸ μέσον τῶν τροχῶν, τῶν ὑποκάτω τῶν χερουβὶμ, καὶ πλῆσον τὰς δράκας σου ἀνθράκων πυρὸς ἐκ μέσου τῶν χερουβὶμ, καὶ διασκόρπισον ἐπὶ τὴν πόλιν· καὶ εἰσῆλθεν ἐνώπιον ἐμοῦ.

\vs{3}Καὶ τὰ χερουβὶμ εἱστήκει ἐκ δεξιῶν τοῦ οἴκου ἐν τῷ εἰσπορεύεσθαι τὸν ἄνδρα, καὶ ἡ νεφέλη ἔπλησε τὴν αὐλήν τὴν ἐσωτέραν.
\vs{4}Καὶ ἀπῇρεν ἡ δόξα Κυρίου ἀπὸ τῶν χερουβὶμ εἰς τὸ αἴθριον τοῦ οἴκου, καὶ ἐπλησε τὸν οἶκον ἡ νεφέλη· καὶ ἡ αὐλὴ ἐπλήσθη τοῦ φέγγους τῆς δόξης Κυρίου.
\vs{5}Καὶ φωνὴ τῶν πτερύγων τῶν χερουβὶμ ἠκούετο ἕως τῆς αὐλῆς τῆς ἐξωτέρας, ὡς φωνὴ Θεοῦ σαδδαῒ λαλοῦντος.

\vs{6}Καὶ ἐγένετο ἐν τῷ ἐντέλλεσθαι αὐτὸν τῷ ἀνδρὶ τῷ ἐνδεδυκότι τὴν στολὴν τὴν ἁγίαν, λέγων, λάβε πῦρ ἐκ μέσου τῶν τροχῶν ἐκ μέσου τῶν χερουβὶμ, καὶ εἰσῆλθε καὶ ἔστη ἐχόμενος τῶν τροχῶν.
\vs{7}Καὶ ἐξέτεινε τὴν χεῖρα αὐτοῦ εἰς μέσον τοῦ πυρὸς τοῦ ὄντος εἰς μέσον τῶν χερουβὶμ, καὶ ἔλαβε, καὶ ἔδωκεν εἰς τὰς χεῖρας τοῦ ἐνδεδυκότος τὴν στολὴν τὴν ἁγίαν, καὶ ἔλαβε καὶ ἐξῆλθε.

\vs{8}Καὶ ἴδον τὰ χερουβὶμ ὁμοίωμα χειρῶν ἀνθρώπων ὑποκάτωθεν τῶν πτερύγων αὐτῶν.
\vs{9}Καὶ ἴδον, καὶ ἰδοὺ τροχοὶ τέσσαρες εἱστήκεισαν ἐχόμενοι τῶν χερουβὶμ· τροχὸς εἷς ἐχόμενος χεροὺβ ἑνός· καὶ ἡ ὄψις τῶν τροχῶν, ὡς ὄψις λίθου ἄνθρακος.
\vs{10}Καὶ ὄψις αὐτῶν ὁμοίωμα ἓν τοῖς τέσσαρσιν, ὃν τρόπον ὅταν ᾖ τροχὸς ἐν μέσῳ τροχοῦ.
\vs{11}Ἐν τῷ πορεύεσθαι αὐτὰ, εἰς τὰ τέσσαρα μέρη αὐτῶν ἐπορεύοντο, οὐκ ἐπέστρεφον ἐν τῷ πορεύεσθαι αὐτά· ὅτι εἰς ὃν ἂν τόπον ἐπέβλεψεν ἡ ἀρχὴ ἡ μία, ἐπορεύοντο, καὶ οὐκ ἐπέστρεφον ἐν τῷ πορεύεσθαι αὐτά.
\vs{12}Καὶ οἱ νῶτοι αὐτῶν, καὶ αἱ χεῖρες αὐτῶν, καὶ αἱ πτέρυγες αὐτῶν, καὶ οἱ τροχοὶ πλήρεις ὀφθαλμῶν κυκλόθεν τοῖς τέσσαρσι τροχοῖς,
\vs{13}τοῖς δὲ τροχοῖς τούτοις ἐπεκλήθη Γελγὲλ ἀκούοντός μου.
\vs{15}Καὶ τὰ χερουβὶμ ἦσαν τοῦτο τὸ ζῶον ὃ ἴδον ἐπὶ τοῦ ποταμοῦ τοῦ Χοβάρ.

\vs{16}Καὶ ἐν τῷ πορεύεσθαι τὰ χερουβὶμ, ἑπορεύοντο οἱ τροχοὶ, καὶ οὗτοι ἐχόμενοι αὐτῶν· καὶ ἐν τῷ ἐξαίρειν τὰ χερουβὶμ τὰς πτέρυγας αὐτῶν τοῦ μετεωρίζεσθαι ἀπὸ τῆς γῆς, οὐκ ἐπέστρεφον οἱ τροχοὶ αὐτῶν.
\vs{17}Ἐν τῷ ἑστᾶναι αὐτὰ, εἱστήκεισαν· καὶ ἐν τῷ μετεωρίζεσθαι αὐτὰ, μετεωρίζοντο μετʼ αὐτῶν· διότι πνεῦμα ζωῆς ἐν αὐτοῖς ἦν.

\vs{18}Καὶ ἐξῆλθε δόξα Κυρίου ἀπὸ τοῦ οἴκου, καὶ ἐπέβη ἐπὶ τὰ χερουβίμ.
\vs{19}Καὶ ἀνέλαβον τὰ χερουβὶμ τὰς πτέρυγας αὐτῶν, καὶ ἐμετεωρίσθησαν ἀπὸ τῆς γῆς ἐνώπιον ἐμοῦ· ἐν τῷ ἐξελθεῖν αὐτὰ, καὶ οἱ τροχοὶ ἐχόμενοι αὐτῶν· καὶ ἔστησαν ἐπὶ τὰ πρόθυρα τῆς πύλης οἴκου Κυρίου τῆς ἀπέναντι, καὶ δόξα Θεοῦ Ἰσραὴλ ἦν ἐπʼ αὐτῶν ὑπεράνω.

\vs{20}Τοῦτο τὸ ζῶόν ἐστιν ὃ ἴδον ὑποκάτω Θεοῦ Ἰσραὴλ ἐπὶ τοῦ ποταμοῦ τοῦ Χοβὰρ, καὶ ἔγνων ὅτι χερουβίμ ἐστι.
\vs{21}Τέσσαρα πρόσωπα τῷ ἑνὶ, καὶ ὀκτὼ πτέρυγες τῷ ἑνὶ, καὶ ὁμοίωμα χειρῶν ἀνθρώπου ὑποκάτωθεν τῶν πτερύγων αὐτῶν.
\vs{22}Καὶ ὁμοίωσις τῶν προσώπων αὐτῶν, ταῦτα τὰ πρόσωπά ἐστιν ἃ ἴδον ὑποκάτω τῆς δόξης τοῦ Θεοῦ Ἰσραὴλ ἐπὶ τοῦ ποταμοῦ Χοβάρ· καὶ αὐτὰ ἕκαστον κατὰ πρόσωπον αὐτῶυ ἐπορεύοντο.

\ch{11}
Καὶ ἀνέλαβέ με πνεῦμα, καὶ ἤγαγέ με ἐπὶ τὴν πύλην τοῦ οἴκου Κυρίου τὴν κατέναντι, τὴν βλέπουσαν κατὰ ἀνατολάς· καὶ ἰδοὺ ἐπὶ τῶν προθύρων τῆς πύλης ὡς εἴκοσι καὶ πέντε ἄνδρες, καὶ ἴδον ἐν μέσῳ αὐτῶν τὸν Ἰεχονίαν τὸν τοῦ Ἔζερ, καὶ Φαλτίαν τὸν τοῦ Βαναίου, τοὺς ἀφηγουμένους τοῦ λαοῦ.

\vs{2}Καὶ εἶπε Κύριος πρὸς μὲ, υἱὲ ἀνθρώπου, οὗτοι οἱ ἄνδρες οἱ λογιζόμενοι μάταια καὶ βουλευόμενοι βουλὴν πονηρὰν ἐν τῇ πόλει ταύτῃ·
\vs{3}Οἱ λέγοντες, οὐχὶ προσφάτως ᾠκοδόμηνται αἱ οἰκίαι; αὕτη ἐστὶν ὁ λέβης, ἡμεῖς δὲ τὰ κρέα.
\vs{4}Διατοῦτο προφήτευσον ἐπʼ αὐτοὺς, προφήτευσον υἱὲ ἀνθρώπου.
\vs{5}Καὶ ἔπεσεν ἐπʼ ἐμὲ πνεῦμα Κυρίου, καὶ εἶπε πρὸς μὲ, λέγε,

Τάδε λέγει Κύριος, οὕτως εἴπατε οἶκος Ἰσραήλ, καὶ τὰ διαβούλια τοῦ πνεύματος ὑμῶν ἐγὼ ἐπίσταμαι.
\vs{6}Ἐπληθύνατε νεκροὺς ὑμῶν ἐν τῇ πόλει ταύτῃ, καὶ ἐνεπλήσατε τὰς ὁδοὺς αὐτῶν τραυματιῶν.
\vs{7}Διατοῦτο τάδε λέγει Κύριος, τοὺς νεκροὺς ὑμῶν οὓς ἐπατάξατε ἐν μέσῳ αὐτῆς, οὗτοί εἰσι τὰ κρέα, αὕτὴ δὲ ὁ λέβης ἐστὶ, καὶ ὑμᾶς ἐξάξω ἐκ μέσου αὐτῆς.
\vs{8}Ῥομφαίαν φοβεῖσθε, καὶ ῥομφαίαν ἐπάξω ἐφʼ ὑμᾶς, λέγει Κύριος.
\vs{9}Καὶ ἐξάξω ὑμᾶς ἐκ μέσου αὐτῆς, καὶ παραδώσω ὑμᾶς εἰς χεῖρας ἀλλοτρίων, καὶ ποιήσω ἐν ὑμῖν κρίματα,
\vs{10}ἐν ῥομφαίᾳ πεσεῖσθε, ἐπὶ τῶν ὀρέων τοῦ Ἰσραὴλ κρινῶ ὑμᾶς, καὶ ἐπιγνώσεσθε ὅτι ἐγὼ Κύριος.

\vs{13}καὶ ἐγένετο ἐν τῷ προφητεύειν με, καὶ Φαλτίας ὁ τοῦ Βαναίου ἀπέθανε· καὶ πίπτω ἐπὶ πρόσωπόν μου, καὶ ἀνεβόησα φωνῇ μεγάλῃ, καὶ εἶπα, οἴμοι οἴμοι Κύριε, εἰς συντέλειαν ποιεῖς σὺ τοὺς καταλοίπους τοῦ Ἰσραήλ;
\vs{14}Καὶ ἐγένετο λόγος Κυρίου πρὸς μὲ, λέγων,
\vs{15}υἱὲ ἀνθρώπου, οἱ ἀδελφοί σου καὶ οἱ ἄνδρες τῆς αἰχμαλωσίας σου καὶ πᾶς ὁ οἶκος τοῦ Ἰσραὴλ συντετέλεσται, οἷς εἶπαν αὐτοῖς οἱ κατοικοῦντες Ἱερουσαλὴμ, μακρὰν ἀπέχετε ἀπὸ τοῦ Κυρίου, ἡμῖν δέδοται ἡ γῆ εἰς κληρονομίαν.
\vs{16}Διατοῦτο εἰπὸν,

Τάδε λέγει Κύριος, ὅτι ἀπώσομαι αὐτοὺς εἰς τὰ ἔθνη, καὶ διασκορπιῶ αὐτοὺς εἰς πᾶσαν γῆν, καὶ ἔσομαι αὐτοῖς εἰς ἁγίασμα μικρὸν ἐν ταῖς χώραις, οὗ ἐὰν εἰσέλθωσιν ἐκεῖ.
\vs{17}Διατοῦτο εἰπὸν, τάδε λέγει Κύριος, καὶ εἰσδέξομαι αὐτοὺς ἐκ τῶν ἐθνῶν, καὶ συνάξω αὐτοὺς ἐκ τῶν χωρῶν οὗ διέσπειρα αὐτοὺς ἐν αὐταῖς· καὶ δώσω αὐτοῖς τὴν γῆν τοῦ Ἰσραήλ.

\vs{18}Καὶ εἰσελεύσονται ἐκεῖ, καὶ ἐξαροῦσι πάντα τὰ βδελύγματα αὐτῆς καὶ πάσας τὰς ἀνομίας αὐτῆς ἐξ αὐτῆς.
\vs{19}Καὶ δώσω αὐτοῖς καρδίαν ἑτέραν, καὶ πνεῦμα καινὸν δώσω ἐν αὐτοῖς, καὶ ἐκσπάσω τὴν καρδίαν τὴν λιθίνην ἐκ τῆς σαρκὸς αὐτῶν, καὶ δώσω αὐτοῖς καρδίαν σαρκίνην·
\vs{20}Ὅπως ἐν τοῖς προστάγμασί μου πορεύωνται, καὶ τὰ δικαιώματά μου φυλάσσωνται καὶ ποιῶσιν αὐτὰ, καὶ ἔσονταί μοι εἰς λαὸν, καὶ ἐγὼ ἔσομαι αὐτοῖς εἰς Θεόν.

\vs{21}Καὶ εἰς τὴν καρδίαν τῶν βδελυγμάτων αὐτῶν καὶ τῶν ἀνομιῶν αὐτῶν, ὡς ἡ καρδία αὐτῶν ἐπορεύετο, τὰς ὁδοὺς αὐτῶν εἰς τὰς κεφαλὰς αὐτῶν δέδωκα, λέγει Κύριος.

\vs{22}Καὶ ἐξῇραν τὰ χερουβὶμ τὰς πτέρυγας αὐτῶν, καὶ οἱ τροχοὶ ἐχόμενοι αὐτῶν, καὶ ἡ δόξα Θεοῦ Ἰσραὴλ ἐπʼ αὐτὰ ὑπεράνω αὐτῶν.
\vs{23}Καὶ ἀνέβη ἡ δόξα Κυρίου ἐκ μέσης τῆς πόλεως, καὶ ἔστη ἐπὶ τοῦ ὄρους ὃ ἦν ἀπέναντι τῆς πόλεως.

\vs{24}Καὶ ἀνέλαβέ με πνεῦμα, καὶ ἤγαγέ με εἰς γῆν Χαλδαίων εἰς τὴν αἰχμαλωσίαν ἐν ὁράσει ἐν πνεύματι Θεοῦ· καὶ ἀνέβην ἀπὸ τῆς ὁράσεως ἧς ἴδον.
\vs{25}Καὶ ἐλάλησα πρὸς τὴν αἰχμαλωσίαν πάντας τοὺς λόγους τοῦ Κυρίου οὓς ἔδειξέ μοι.

\ch{12}
Καὶ ἐγένετο λόγος Κυρίου πρὸς μὲ, λέγων,
\vs{2}υἱὲ ἀνθρώπου, ἐν μέσῳ τῶν ἀδικιῶν αὐτῶν σὺ κατοικεῖς, οἳ ἔχουσιν ὀφθαλμοὺς τοῦ βλέπειν, καὶ οὐ βλέπουσι, καὶ ὦτα ἔχουσι τοῦ ἀκούειν, καὶ οὐκ ἀκούουσι, διότι οἶκος παραπικραίνων ἐστί.

\vs{3}Καὶ σὺ, υἱὲ ἀνθρώπου, ποίησον σεαυτῷ σκεύη αἰχμαλωσίας ἡμέρας ἐνώπιον αὐτῶν, καὶ αἰχμαλωτευθήσῃ ἐκ τοῦ τόπου σου εἰς ἕτερον τόπον ἐνώπιον αὐτῶν, ὅπως ἴδωσι διότι οἶκος παραπικραίνων ἐστί.
\vs{4}Καὶ ἐξοίσεις τὰ σκεύη σου σκεύη αἰχμαλωσίας ἡμέρας κατʼ ὀφθαλμοὺς αὐτῶν, καὶ σὺ ἐξελεύσῃ ἑσπέρας, ὡς ἐκπορεύεται αἰχμάλωτος ἐνώπιον αὐτῶν.
\vs{5}Διόρυξον σεαυτῷ εἰς τὸν τοῖχον, καὶ διεξελεύσῃ διʼ αὐτοῦ ἐνώπιον αὐτῶν·
\vs{6}ἐπʼ ὤμων ἀναληφθήσῃ, καὶ κεκρυμμένος ἐξελεύσῃ, τὸ πρόσωπόν σου συγκαλύψεις, καὶ οὐ μὴ ἴδῃς τὴν γῆν, διότι τέρας δέδωκά σε τῷ οἴκῳ Ἰσραήλ.

\vs{7}Καὶ ἐποίησα οὕτως κατὰ πάντα ὅσα ἐνετείλατό μοι· καὶ σκεύη ἐξήνεγκα αἰχμαλωσίας ἡμέρας, καὶ ἑσπέρας διώρυξα ἐμαυτῷ τὸν τοῖχον, καὶ κεκρυμμένος ἐξῆλθον, ἐπʼ ὤμων ἀνελήφθην ἐνώπιον αὐτῶν.
\vs{8}Καὶ ἐγένετο λόγος Κυρίου τοπρωῒ πρὸς μὲ, λέγων,
\vs{9}υἱὲ ἀνθρώπου, οὐκ εἶπαν πρὸς σὲ ὁ οἶκος τοῦ Ἰσραὴλ, οἶκος ὁ παραπικραίνων, τί σὺ ποιεῖς;
\vs{10}Εἰπὸν πρὸς αὐτοὺς,

Τάδε λέγει Κύριος Κύριος, ὁ ἄρχων καὶ ὁ ἀφηγούμενος ἐν Ἱερουσαλὴμ, καὶ παντὶ οἴκῳ Ἰσραὴλ, οἵ εἰσιν ἐν μέσῳ αὐτῶν,
\vs{11}εἰπὸν, ὅτι ἐγὼ τέρατα ποιῶ· ὃν τρόπον πεποίηκα, οὕτως ἔσται αὐτῷ· ἐν μετοικεσίᾳ καὶ ἐν αἰχμαλωσίᾳ πορεύσονται.
\vs{12}Καὶ ὁ ἄρχων ἐν μέσῳ αὐτων ἐπʼ ὤμων ἀρθήσεται, καὶ κεκρυμμένος ἐξελεύσεται διὰ τοῦ τοίχου, καὶ διορύξει τοῦ ἐξελθεῖν αὐτὸν διʼ αὐτοῦ, τὸ πρόσωπον αὐτοῦ συγκαλύψει, ὅπως μὴ ὁραθῇ ὀφθαλμῷ, καὶ αὐτὸς τὴν γῆν οὐκ ὄψεται.
\vs{13}Καὶ ἐκπετάσω τὸ δίκτυόν μου ἐπʼ αὐτὸν, καὶ συλληφθήσεται ἐν τῇ περιοχῇ μου, καὶ ἄξω αὐτὸν εἰς Βαβυλῶνα εἰς γῆν Χαλδαίων, καὶ αὐτὴν οὐκ ὄψεται, καὶ ἐκεῖ τελευτήσει.
\vs{14}Καὶ πάντας τοὺς κύκλῳ αὐτοῦ τοὺς βοηθοὺς αὐτοῦ, καὶ πάντας τοὺς ἀντιλαμβανομένους αὐτοῦ, διασπερῶ εἰς πάντα ἄνεμον, καὶ ῥομφαίαν ἐκκενώσω ὀπίσω αὐτῶν.
\vs{15}Καὶ γνώσονται διότι ἐγὼ Κύριος, ἐν τῷ διασκορπίσαι με αὐτοὺς ἐν τοῖς ἔθνεσι, καὶ διασπερῶ αὐτοὺς ἐν ταῖς χώραις.

\vs{16}Καὶ ὑπολείψομαι ἐξ αὐτῶν ἄνδρας ἀριθμῷ ἐκ ῥομφαίας, καὶ ἐκ λιμοῦ, καὶ ἐκ θανάτου, ὅπως ἐκδιηγῶνται πάσας τὰς ἀνομίας αὐτῶν ἐν τοῖς ἔθνεσιν οὗ εἰσήλθοσαν ἐκεῖ, καὶ γνώσονται ὅτι ἐγὼ Κύριος.

\vs{17}Καὶ ἐγένετο λόγος Κυρίου πρὸς μὲ, λέγων,
\vs{18}υἱὲ ἀνθρώπου, τὸν ἄρτον σου μετὰ ὀδύνης φάγεσαι, καὶ τὸ ὕδωρ μετὰ βασάνου καὶ θλίψεως πίεσαι.
\vs{19}Καὶ ἐρεῖς πρὸς τὸν λαὸν τῆς γῆς, τάδε λέγει Κύριος τοῖς κατοικοῦσιν Ἱερουσαλὴμ ἐπὶ τῆς γῆς τοῦ Ἰσραὴλ, τοὺς ἄρτους αὐτῶν μετὰ ἐνδείας φάγονται, καὶ τὸ ὕδωρ αὐτῶν μετὰ ἀφανισμοῦ πίονται, ὅπως ἀφανισθῇ ἡ γῆ σὺν πληρώματι αὐτῆς· ἐν ἀσεβείᾳ γὰρ πάντες οἱ κατοικοῦντες ἐν αὐτῇ.
\vs{20}Καὶ αἱ πόλεις αὐτῶν αἱ κατοικούμεναι ἐξερημωθήσονται, καὶ ἡ γῆ εἰς ἀφανισμὸν ἔσται, καὶ ἐπιγνώσεσθε διότι ἐγὼ Κύριος.

\vs{21}Καὶ ἐγένετο λόγος Κυρίου πρὸς μὲ, λέγων,
\vs{22}υἱὲ ἀνθρώπου, τίς ἡ παραβολὴ ὑμῖν ἐπὶ τῆς γῆς τοῦ Ἰσραὴλ, λέγοντες, μακραὶ αἱ ἡμέραι, ἀπόλωλεν ὅρασις;
\vs{23}Διατοῦτο εἰπὸν πρὸς αὐτούς,

Τάδε λέγει Κύριος, καὶ ἀποστρέψω τὴν παραβολὴν ταύτην, καὶ οὐκέτι μὴ εἴπωσι τὴν παραβολὴν ταύτην οἶκος τοῦ Ἰσραήλ, ὅτι λαλήσεις πρὸς αὐτούς, ἡγγίκασιν αἱ ἡμέραι, καὶ λόγος πάσης ὁράσεως.
\vs{24}Ὅτι οὐκ ἔσται ἔτι πᾶσα ὅρασις ψευδὴς, καὶ μαντευόμενος τὰ πρὸς χάριν ἐν μέσῳ τῶν υἱῶν Ἰσραήλ.
\vs{25}Διότι ἐγὼ Κύριος λαλήσω τοὺς λόγους μου, λαλήσω καὶ ποιήσω, καὶ οὐ μὴ μηκύνω ἔτι· ὅτι ἐν ταῖς ἡμέραις ὑμῶν οἶκος ὁ παραπικραίνων, λαλήσω λόγον καὶ ποιήσω, λέγει Κύριος.

\vs{26}Καὶ ἐγένετο λόγος Κυρίου πρὸς μὲ, λέγων,
\vs{27}υἱὲ ἀνθρώπου, ἰδοὺ ὁ οἶκος Ἰσραὴλ ὁ παραπικραίνων, λέγοντες λέγουσιν, ἡ ὅρασις ἣν οὗτος ὁρᾷ, εἰς ἡμέρας πολλὰς, καὶ εἰς καιροὺς μακροὺς οὗτος προφητεύει.
\vs{28}Διατοῦτο εἰπὸν πρὸς αὐτοὺς,

Τάδε λέγει Κύριος, οὐ μὴ μηκύνωσιν οὐκέτι πάντες οἱ λόγοι μου, οὓς ἂν λαλήσω· λαλήσω καὶ ποιήσω, λέγει Κύριος.

\ch{13}
Καὶ ἐγένετο λόγος Κυρίου πρὸς μὲ, λέγων,
\vs{2}υἱὲ ἀνθρώπου, προφήτευσον ἐπὶ τοὺς προφήτας τοῦ Ἰσραὴλ, καὶ προφητεύσεις, καὶ ἐρεῖς πρὸς αὐτοὺς, ἀκούσατε λόγον Κυρίου·

\vs{3}Τάδε λέγει Κύριος, οὐαὶ τοῖς προφητεύουσιν ἀπὸ καρδίας αὐτῶν, καὶ τὸ καθόλου μὴ βλέπουσιν.
\vs{4}Ὡς ἀλώπεκες ἐν ταῖς ἐρήμοις οἱ προφῆταί σου, Ἰσραήλ·
\vs{5}Οὐκ ἔστησαν ἐν στερεώματι· καὶ συνήγαγον ποίμνια ἐπὶ τὸν οἶκον τοῦ Ἰσραήλ· οὐκ ἀνέστησαν οἱ λέγοντες, ἐν ἡμέρᾳ Κυρίου,
\vs{6}βλέποντες ψευδῆ, μαντευόμενοι μάταια, οἱ λέγοντες, λέγει Κύριος, καὶ Κύριος οὐκ ἀπέσταλκεν αὐτοὺς, καὶ ἤρξαντο τοῦ ἀναστῆσαι λόγον.
\vs{7}Οὐχὶ ὅρασιν ψευδῆ ἑωράκατε; καὶ μαντείας ματαίας εἰρήκατε;
\vs{8}Καὶ διατοῦτο εἰπὸν,

Τάδε λέγει Κύριος, ἄνθʼ ὧν οἱ λόγοι ὑμῶν ψευδεῖς, καὶ αἱ μαντεῖαι ὑμῶν μάταιαι, διατοῦτο ἰδοὺ ἐγὼ ἐφʼ ὑμᾶς, λέγει Κύριος.
\vs{9}Καὶ ἐκτενῶ τὴν χεῖρά μου ἐπὶ τοὺς προφήτας τοὺς ὁρῶντας ψευδῆ, καὶ τοὺς ἀποφθεγγομένους μάταια· ἐν παιδείᾳ τοῦ λαοῦ μου οὐκ ἔσονται, οὐδὲ ἐν γραφῇ οἴκου Ἰσραὴλ οὐ γραφήσονται, καὶ εἰς τὴν γῆν τοῦ Ἰσραὴλ οὐκ εἰσελεύσονται, καὶ γνώσονται διότι ἐγὼ Κύριος·
\vs{10}Ἀνθʼ ὧν ἐπλάνησαν τὸν λαόν μου, λέγοντες, εἰρήνη, καὶ οὐκ ἔστιν εἰρήνη· καὶ οὗτος οἰκοδομεῖ τοῖχον, καὶ αὐτοὶ ἀλείφουσιν αὐτὸν, πεσεῖται.

\vs{11}Εἰπὸν πρὸς τοὺς ἀλείφοντας, πεσεῖται, καὶ ἔσται ὑετὸς κατακλύζων, καὶ δώσω λίθους πετροβόλους εἰς τοὺς ἐνδέσμους αὐτῶν, καὶ πεσοῦνται, καὶ πνεῦμα ἐξαῖρον, καὶ ῥαγήσεται.
\vs{12}Καὶ ἰδοὺ πέπτωκεν ὁ τοῖχος, καὶ οὐκ ἐροῦσι πρὸς ὑμᾶς, ποῦ ἐστιν ἡ ἀλοιφὴ ὑμῶν ἣν ἠλείψατε;
\vs{13}Διατοῦτο τάδε λέγει Κύριος, καὶ ῥήξω πνοὴν ἐξαίρουσαν μετὰ θυμοῦ, καὶ ὑετὸς κατακλύζων ἐν ὀργῇ μου ἔσται· καὶ τοὺς λίθους τοὺς πετροβόλους ἐν θυμῷ ἐπάξω εἰς συντέλειαν·
\vs{14}Καὶ κατασκάψω τὸν τοῖχον ὃν ἠλείψατε, και πεσεῖται· καὶ θήσω αὐτὸν ἐπὶ τὴν γῆν, καὶ ἀποκαλυφθήσεται τὰ θεμέλια αὐτοῦ, καὶ πεσεῖται, καὶ συντελεσθήσεσθε μετʼ ἐλέγχων, καὶ ἐπιγνώσεσθε διότι ἐγὼ Κύριος.

\vs{15}Καὶ συντελέσω τὸν θυμόν μου ἐπὶ τὸν τοῖχον, καὶ ἐπὶ τοὺς ἀλείφοντας αὐτὸν, πεσεῖται· καὶ εἶπα πρὸς ὑμᾶς, οὐκ ἔστιν ὁ τοῖχος, οὐδὲ οἱ ἀλείφοντες αὐτὸν,
\vs{16}προφῆται τοῦ Ἰσραὴλ, οἱ προφητεύοντες ἐπὶ Ἱερουσαλὴμ, καὶ οἱ ὁρῶντες αὐτῇ εἰρήνην, καὶ οὐκ ἔστιν εἰρήνη, λέγει Κύριος.

\vs{17}Καὶ σὺ, υἱὲ ἀνθρώπου, στήρισον τὸ πρόσωπόν σου ἐπὶ τὰς θυγατέρας τοῦ λαοῦ σου, τὰς προφητευούσας ἀπὸ καρδίας αὐτῶν, καὶ προφήτευσον ἐπʼ αὐτὰς,
\vs{18}καὶ ἐρεῖς, τάδε λέγει Κύριος, οὐαὶ ταῖς συῤῥαπτούσαις προσκεφάλαια ὑπὸ πάντα ἀγκῶνα χειρὸς, καὶ ποιούσαις ἐπιβόλαια ἐπὶ πᾶσαν κεφαλὴν πάσης ἡλικίας, τοῦ διαστρέφειν ψυχάς· αἱ ψυχαὶ διεστράφησαν τοῦ λαοῦ μου· καὶ ψυχὰς περιεποιοῦντο.
\vs{19}Καὶ ἐβεβήλουν με πρὸς τὸν λαόν μου, ἕνεκεν δρακὸς κριθῶν, καὶ ἕνεκεν κλασμάτων ἄρτων, τοῦ ἀποκτεῖναι ψυχὰς ἃς οὐκ ἔδει ἀποθανεῖν, καὶ τοῦ περιποιήσασθαι ψυχὰς ἃς οὐκ ἔδει ζῆσαι, ἐν τῷ ἀποφθέγγεσθαι ὑμᾶς λαῷ εἰσακούοντι μάταια ἀποφθέγματα.

\vs{20}Διατοῦτο τάδε λέγει Κύριος Κύριος, ἰδοὺ ἐγὼ ἐπὶ τὰ προσκεφάλαια ὑμῶν, ἐφʼ ἃ ὑμεῖς ἐκεῖ συστρέφετε ψυχάς· καὶ διαῤῥήξω αὐτὰ ἀπὸ τῶν βραχιόνων ὑμῶν, καὶ ἐξαποστελῶ τὰς ψυχὰς ἃς ὑμεῖς ἐκστρέφετε τὰς ψυχὰς αὐτῶν εἰς διασκορπισμόν.
\vs{21}Καὶ διαῤῥήξω τὰ ἐπιβόλαια ὑμῶν, καὶ ῥύσομαι τὸν λαόν μου ἐκ χειρὸς ὑμῶν, καὶ οὐκέτι ἔσονται ἐν χερσὶν ὑμῶν εἰς συστροφήν· καὶ ἐπιγνώσεσθε διότι ἐγὼ Κύριος.

\vs{22}Ἀνθʼ ὧν διεστρέψετε καρδίαν δικαίου, καὶ ἐγὼ οὐ διέστρεφον αὐτὸν, καὶ τοῦ κατισχῦσαι χεῖρας ἀνόμου τὸ καθόλου μὴ ἀποστρέψαι ἀπὸ τῆς ὁδοῦ αὐτοῦ τῆς πονηρᾶς, καὶ ζῆσαι αὐτόν·
\vs{23}Διατοῦτο ψευδῆ οὐ μὴ ἴδητε, καὶ μαντείας οὐ μὴ μαντεύσησθε ἔτι· καὶ ῥύσομαι τὸν λαόν μου ἐκ χειρὸς ὑμῶν, καὶ γνώσεσθε ὅτι ἐγὼ Κύριος.

\ch{14}
Καὶ ἦλθον πρὸς μὲ ἐκ τῶν πρεσβυτέρων ἄνδρες τοῦ λαοῦ Ἰσραὴλ καὶ ἐκάθισαν πρὸ προσώπου μου.
\vs{2}Καὶ ἐγένετο πρὸς μὲ λόγος Κυρίου, λέγων,
\vs{3}υἱὲ ἀνθρώπου, οἱ ἄνδρες οὗτοι ἔθεντο τὰ διανοήματα αὐτῶν ἐπὶ τὰς καρδίας αὐτῶν, καὶ τὴν κόλασιν τῶν ἀδικιῶν αὐτῶν ἔθηκαν πρὸ προσώπου αὐτῶν· εἰ ἀποκρινόμενος ἀποκριθῶ αὐτοῖς;
\vs{4}Διατοῦτο λάλησον αὐτοῖς, καὶ ἐρεῖς πρὸς αὐτοὺς, τάδε λέγει Κύριος, ἄνθρωπος ἄνθρωπος ἐκ τοῦ οἴκου Ἰσραὴλ, ὃς ἂν θῇ τὰ διανοήματα αὐτοῦ ἐπὶ τὴν καρδίαν αὐτοῦ, καὶ τὴν κόλασιν τῆς ἀδικίας αὐτοῦ τάξῃ πρὸ προσώπου αὐτοῦ, καὶ ἔλθῃ πρὸς τὸν προφήτην, ἐγὼ Κύριος ἀποκριθήσομαι αὐτῷ ἐν οἷς ἐνέχεται ἡ διάνοια αὐτοῦ,
\vs{5}ὅπως πλαγιάσῃ τὸν οἶκον τοῦ Ἰσραὴλ κατὰ τὰς καρδίας αὐτῶν τὰς ἀπηλλοτριωμένας ἀπʼ ἐμοῦ ἐν τοῖς ἐνθυμήμασιν αὐτῶν.

\vs{6}Διατοῦτο εἰπὸν εἰς τὸν οἶκον τοῦ Ἰσραὴλ, τάδε λέγει Κύριος Κύριος, ἐπιστράφητε καὶ ἀποστρέψατε ἀπὸ τῶν ἐπιτηδευμάτων ὑμῶν, καὶ ἀπὸ πασῶν τῶν ἀσεβειῶν ὑμῶν, καὶ ἐπιστρέψατε τὰ πρόσωπα ὑμῶν.
\vs{7}Διότι ἄνθρωπος ἄνθρωπος ἐκ τοῦ οἴκου Ἰσραὴλ, καὶ ἐκ τῶν προσηλύτων τῶν προσηλυτευόντων ἐν τῷ Ἰσραὴλ, ὃς ἂν ἀπαλλοτριωθῇ ἀπʼ ἐμοῦ, καὶ θῆται τὰ ἐνθυμήματα αὐτοῦ ἐπὶ τὴν καρδίαν αὐτοῦ, καὶ τὴν κόλασιν τῆς ἀδικίας αὐτοῦ τάξῃ πρὸ προσώπου αὐτοῦ, καὶ ἔλθῃ πρὸς τὸν προφήτην τοῦ ἐπερωτῆσαι αὐτὸν ἐν ἐμοὶ, ἐγὼ Κύριος ἀποκριθήσομαι αὐτῷ, ἐν ᾧ ἐνέχεται ἐν αὐτῷ.
\vs{8}Καὶ στηριῶ τὸ πρόσωπόν μου ἐπὶ τὸν ἄνθρωπον ἐκεῖνον, καὶ θήσομαι αὐτὸν εἰς ἔρημον καὶ εἰς ἀφανισμὸν, καὶ ἐξαρῶ αὐτὸν ἐκ μέσου τοῦ λαοῦ μου, καὶ ἐπιγνώσεσθε ὅτι ἐγὼ Κύριος.

\vs{9}Καὶ ὁ προφήτης ἐὰν πλανήσῃ καὶ λαλήσῃ, ἐγὼ Κύριος πεπλάνηκα τὸν προφήτην ἐκεῖνον, καὶ ἐκτενῶ τὴν χεῖρά μου ἐπʼ αὐτὸν, καὶ ἀφανιῶ αὐτὸν ἐκ μέσου τοῦ λαοῦ μου Ἰσραήλ.
\vs{10}Καὶ λήψονται τὴν ἀδικίαν αὐτῶν κατὰ τὸ ἀδίκημα τοῦ ἐπερωτῶντος, καὶ κατὰ τὸ ἀδίκημα ὁμοίως τῷ προφήτῃ ἔσται·
\vs{11}Ὅπως μὴ πλανᾶται ἔτι ὁ οἶκος τοῦ Ἰσραὴλ ἀπʼ ἐμοῦ, καὶ ἵνα μὴ μιαίνωνται ἔτι ἐν πᾶσι τοῖς παραπτώμασιν αὐτῶν· καὶ ἔσονταί μοι εἰς λαὸν, καὶ ἐγὼ ἔσομαι αὐτοῖς εἰς Θεὸν, λέγει Κύριος.

\vs{12}Καὶ ἐγένετο λόγος Κυρίου πρὸς μὲ, λέγων,
\vs{13}υἱὲ ἀνθρώπου, γῆ ἣ ἐὰν ἁμάρτῃ μοι τοῦ παραπεσεῖν παράπτωμα, καὶ ἐκτενῶ τὴν χεῖρά μου ἐπʼ αὐτὴν, καὶ συντρίψω αὐτῆς στήριγμα ἄρτου, καὶ ἐξαποστελῶ ἐπʼ αὐτὴν λιμὸν, καὶ ἐξαρῶ ἐξ αὐτῆς ἄνθρωπον καὶ κτήνη.
\vs{14}Καὶ ἐὰν ὦσιν οἱ τρεῖς ἄνδρες οὗτοι ἐν μέσῳ αὐτῆς, Νῶε καὶ Δανιὴλ καὶ Ἰὼβ, αὐτοὶ ἐν τῇ δικαιοσύνῃ αὐτῶν σωθήσονται, λέγει Κύριος.

\vs{15}Ἐὰν καὶ θηρία πονηρὰ ἐπάγω ἐπὶ τὴν γῆν, καὶ τιμωρήσομαι αὐτὴν, καὶ ἔσται εἰς ἀφανισμὸν, οὐκ ἔσται ὁ διοδεύων ἀπὸ προσώπου τῶν θηρίων,
\vs{16}καὶ οἱ τρεῖς ἄνδρες οὗτοι ἐν μέσῳ αὐτῆς ὦσι, ζῶ ἐγὼ, λέγει Κύριος, εἰ υἱοὶ ἢ θυγατέρες σωθήσονται, ἀλλʼ ἢ αὐτοὶ μόνοι σωθήσονται, ἡ δὲ γῆ ἔσται εἰς ὄλεθρον.
\vs{17}Ἢ καὶ ῥομφαίαν ἐὰν ἐπάγω ἐπὶ τὴν γῆν ἐκείνην, καὶ εἴπω, ῥομφαία διελθάτω διὰ τῆς γῆς, καὶ ἐξαρῶ ἐξ αὐτῶν ἄνθρωπον καὶ κτῆνος,
\vs{18}καὶ οἱ τρεῖς ἄνδρες οὗτοι ἐν μέσῳ αὐτῆς, ζῶ ἐγὼ, λέγει Κύριος, οὐ μὴ ῥύσονται υἱοὺς οὐδὲ θυγατέρας, ἀλλʼ ἢ αὐτοὶ μόνοι σωθήσονται.

\vs{19}Ἢ καὶ θάνατον ἐπαποστέλλω ἐπὶ τὴν γῆν ἐκείνην, καὶ ἐκχεῶ τὸν θυμόν μου ἐπʼ αὐτὴν ἐν αἵματι τοῦ ἐξολοθρεῦσαι ἐξ αὐτῆς ἄνθρωπον καὶ κτῆνος,
\vs{20}καὶ Νῶε καὶ Δανιὴλ καὶ Ἰὼβ ἐν μέσῳ αὐτῆς, ζῶ ἐγὼ, λέγει Κύριος, ἐὰν υἱοὶ ἢ θυγατέρες ὑπολειφθῶσιν, αὐτοὶ ἐν τῇ δικαιοσύνῃ αὐτῶν ῥύσονται τὰς ψυχὰς αὐτῶν.

\vs{21}Τάδε λέγει Κύριος, ἐὰν δὲ καὶ τὰς τέσσαρας ἐκδικήσεις μου τὰς πονηρὰς, ῥομφαίαν, καὶ λιμὸν, καὶ θηρία πονηρὰ, καὶ θάνατον ἐξαποστείλω ἐπὶ Ἱερουσαλὴμ, τοῦ ἐξολοθρεῦσαι ἐξ αὐτῆς ἄνθρωπον καὶ κτῆνος·
\vs{22}Καὶ ἰδοὺ ὑπολελειμμένοι ἐν αὐτῇ, οἱ ἀνασεσωσμένοι αὐτῆς, οἳ ἐξάγουσιν ἐξ αὐτῆς υἱοὺς καὶ θυγατέρας, ἰδοὺ ἐκπορεύονται πρὸς ὑμᾶς, καὶ ὄψεσθε τὰς ὁδοὺς αὐτῶν καὶ τὰ ἐνθυμήματα αὐτῶν, καὶ μεταμεληθήσεσθε ἐπὶ τὰ κακὰ ἃ ἐπήγαγον ἐπὶ Ἱερουσαλὴμ, πάντα τὰ κακὰ ἃ ἐπήγαγον ἐπʼ αὐτήν.
\vs{23}Καὶ παρακαλέσουσιν ὑμᾶς, διότι ὄψεσθε τὰς ὁδοὺς αὐτῶν καὶ τὰ ἐνθυμήματα αὐτῶν, καὶ ἐπιγνώσεσθε διότι οὐ μάτην πεποίηκα πάντα ὅσα ἐποίησα ἐν αὐτῇ, λέγει Κύριος.

\ch{15}
Καὶ ἐγένετο λόγος Κυρίου πρὸς μὲ, λέγων,

\vs{2}Καὶ σὺ υἱὲ ἀνθρώπου, τί ἂν γένοιτο τὸ ξύλον τῆς ἀμπέλου ἐκ πάντων τῶν ξύλων τῶν κλημάτων τῶν ὄντων ἐν τοῖς ξύλοις τοῦ δρυμοῦ;
\vs{3}Εἰ λήψονται ἐξ αὐτῆς ξύλον τοῦ ποιῆσαι εἰς ἐργασίαν; εἰ λήψονται ἐξ αὐτῆς πάσσαλον τοῦ κρεμᾶσαι ἐπʼ αὐτὸν πᾶν σκεῦος;
\vs{4}Πάρεξ ὃ πυρὶ δέδοται εἰς ἀναλωσιν, τὴν κατʼ ἐνιαυτὸν κάθαρσιν ἀπʼ αὐτῆς ἀναλίσκει τὸ πῦρ, καὶ ἐκλείπει εἰς τέλος· μὴ χρήσιμον ἔσται εἰς ἐργασίαν;
\vs{5}Οὐδὲ ἔτι αὐτοῦ ὄντος ὁλοκλήρου οὐκ ἔσται εἰς ἐργασίαν; μὴ ὅτι ἐὰν καὶ πῦρ αὐτὸ ἀναλώσῃ εἰς τέλος, εἰ ἔτι ἔσται εἰς ἐργασίαν.
\vs{6}Διὰ τοῦτο εἰπὸν,

Τάδε λέγει Κύριος, ὃν τρόπον τὸ ξύλον τῆς ἀμπέλου ἐν τοῖς ξύλοις τοῦ δρυμοῦ ὃ δέδωκα αὐτὸ πυρὶ εἰς ἀνάλωσιν, οὕτως δέδωκα τοὺς κατοικοῦντας Ἰερουσαλήμ.
\vs{7}Καὶ δώσω τὸ πρόσωπόν μου ἐπʼ αὐτούς· ἐκ τοῦ πυρὸς ἐξελεύσονται, καὶ πῦρ αὐτοὺς καταφάγεται, καὶ ἐπιγνώσονται ὅτι ἐγὼ Κύριος ἐν τῷ στηρίσαι με τὸ πρόσωπόν μου ἐπʼ αὐτούς.
\vs{8}Καὶ δώσω τὴν γῆν εἰς ἀφανισμὸν, ἀνθʼ ὧν παρέπεσον παραπτώματι, λέγει Κύριος.

\ch{16}
Καὶ ἐγένετο λόγος Κυρίου πρὸς μὲ, λέγων,
\vs{2}υἱὲ ἀνθρώπου, διαμάρτυραι τῇ Ἱερουσαλὴμ τὰς ἀνομίας αὐτῆς,
\vs{3}καὶ ἐρεῖς,

Τάδε λέγει Κύριος τῇ Ἱερουσαλὴμ, ἡ ῥίζα σου καὶ ἡ γένεσίς σου ἐκ γῆς Χαναάν· ὁ πατήρ σου Ἀμοῤῥαῖος, καὶ ἡ μήτηρ σου Χετταία.
\vs{4}Καὶ ἡ γένεσίς σου ἐν ᾗ ἡμέρᾳ ἐτέχθης, οὐκ ἔδησας τοὺς μαστοὺς σου, καὶ ἐν ὕδατι οὐκ ἐλούσθης, οὐδὲ ἁλὶ ἡλίσθης, καὶ ἐν σπαργάνοις οὐκ ἐσπαργανώθης,
\vs{5}οὐδὲ ἐφείσατο ὁ ὀφθαλμός μου ἐπὶ σοὶ, τοῦ ποιῆσαί σοι ἓν ἐκ πάντων τούτων, τοῦ παθεῖν τι ἐπὶ σοί· καὶ ἀπεῤῥίφης ἐπὶ πρόσωπον τοῦ πεδίου τῇ σκολιότητι τῆς ψυχῆς σου ἐν ἡμέρᾳ ᾗ ἐτέχθης.

\vs{6}Καὶ διῆλθον ἐπὶ σὲ, καὶ ἴδον σε πεφυρμένην ἐν τῷ αἵματί σου· καὶ εἶπά σοι, ἐκ τοῦ αἵματός σου ζωὴ,
\vs{7}πληθύνου, καθὼς ἡ ἀνατολὴ τοῦ ἀγροῦ δέδωκά σε· καὶ ἐπληθύνθης καὶ ἐμεγαλύνθης, καὶ εἰσῆλθες εἰς πόλεις πόλεων· οἱ μαστοί σου ἀνωρθώθησαν, καὶ ἡ θρίξ σου ἀνέτειλε· σὺ δὲ ἦσθα γυμνὴ καὶ ἀσχημονοῦσα.

\vs{8}Καὶ διῆλθον διὰ σοῦ, καὶ ἴδον σε· καὶ ἰδοὺ καιρός σου, καὶ καιρὸς καταλυόντων· καὶ διεπέτασα τὰς πτέρυγάς μου ἐπὶ σὲ, καὶ ἐκάλυψα τὴν ἀσχημοσύνην σου, καὶ ὤμοσά σοι· καὶ εἰσῆλθον ἐν διαθήκῃ μετὰ σοῦ, λέγει Κύριος· καὶ ἐγένου μοι.
\vs{9}Καὶ ἔλουσά σε ἐν ὕδατι, καὶ ἀπέπλυνα τὸ αἷμά σου ἀπὸ σοῦ, καὶ ἔχρισά σε ἐν ἐλαίῳ.
\vs{10}Καὶ ἐνέδυσά σε ποικίλα, καὶ ὑπέδυσά σε ὑάκινθον· καὶ ἔζωσά σε βύσσῳ καὶ περιέβαλόν σε τριχαπτῳ,
\vs{11}καὶ ἐκόσμησά σε κόσμῳ, καὶ περιέθηκα ψέλλια περὶ τὰς χεῖράς σου, καὶ κάθεμα περὶ τὸν τράχηλόν σου·
\vs{12}Καὶ ἔδωκα ἐνώτιον περὶ τὸν μυκτῆρά σου, καὶ τροχίσκους ἐπὶ τὰ ὦτά σου, καὶ στέφανον καυχήσεως ἐπὶ τὴν κεφαλήν σου·
\vs{13}Καὶ ἐκοσμήθης χρυσίῳ καὶ ἀργυρίῳ, καὶ τὰ περιβόλαιά σου βύσσινα, καὶ τριχαπτὰ, καὶ ποικίλα· σεμίδαλιν καὶ ἔλαιον καὶ μέλι ἔφαγες, καὶ ἐγένου καλὴ σφόδρα.
\vs{14}Καὶ ἐξῆλθέ σου ὄνομα ἐν τοῖς ἔθνεσιν ἐν τῷ κάλλει σου, διότι συντετελεσμένον ἦν ἐν εὐπρεπείᾳ ἐν τῇ ὡραιότητι ᾗ ἔταξα ἐπὶ σὲ, λέγει Κύριος.

\vs{15}Κατεπεποίθεις ἐν τῷ κάλλει σου, καὶ ἐπόρνευσας ἐπὶ τῷ ὀνόματί σου, καὶ ἐξέχεας τὴν πορνείαν σου ἐπὶ πάντα πάροδον·
\vs{16}Καὶ ἔλαβες ἐκ τῶν ἱματίων σου, καὶ ἐποίησας σεαυτῇ εἴδωλα ῥαπτὰ, καὶ ἐξεπόρνευσας ἐπʼ αὐτὰ, καὶ οὐ μὴ εἰσέλθῃς οὐδὲ μὴ γένηται.
\vs{17}Καὶ ἔλαβες τὰ σκεύη τῆς καυχήσεώς σου ἐκ τοῦ χρυσίου μου καὶ ἐκ τοῦ ἀργυρίου μου, ἐξ ὧν ἔδωκά σοι, καὶ ἐποίησας σεαυτῇ εἰκόνας ἀρσενικὰς, καὶ ἐξεπόρνευσας ἐν αὐταῖς.
\vs{18}Καὶ ἔλαβες τὸν ἱματισμὸν τὸν ποικίλον σου, καὶ περιέβαλες αὐτὰς, καὶ τὸ ἔλαιόν μου καὶ τὸ θυμίαμά μου ἔθηκας πρὸ προσώπου αὐτῶν.
\vs{19}Καὶ τοὺς ἄρτοὺς μου οὓς ἔδωκά σοι, σεμίδαλιν καὶ ἔλαιον καὶ μέλι ἐψώμισά σε, καὶ ἔθηκας αὐτὰ τρὸ τροσώπου αὐτῶν εἰς ὀσμὴν εὐωδίας· καὶ ἐγένετο, λέγει Κύριος.

\vs{20}Καὶ ἔλαβες τοὺς υἱούς σου καὶ τὰς θυγατέρας σου ἃς ἐγέννησας, καὶ ἔθυσας αὐτοῖς εἰς ἀνάλωσιν· ὡς μικρὰ ἐξεπόρνευσας,
\vs{21}καὶ ἔσφαξας τὰ τέκνα σου, καὶ ἔδωκας αὐτὰ ἐν τῷ ἀποτροπιάζεσθαί σε αὐτὰ αὐτοῖς·
\vs{22}Τοῦτο παρὰ πᾶσαν τὴν πορνείαν σου, καὶ οὐκ ἐμνήσθης τῆς νηπιότητός σου, ὅτε ἦσθα γυμνὴ καὶ ἀσχημονοῦσα, πεφυρμένη ἐν τῷ αἵματί σου ἔζησας.

\vs{23}Καὶ ἐγένετο μετὰ πάσας τὰς κακίας σου, λέγει Κύριος,
\vs{24}καὶ ᾠκοδόμησας σεαυτῇ οἴκημα πορνικὸν, καὶ ἐποίησας σεαυτῇ ἔκθεμα ἐν πάσῃ πλατείᾳ,
\vs{25}καὶ ἐπʼ ἀρχῆς πάσης ὁδοῦ ᾠκοδόμησας τὰ πορνεῖά σου, καὶ ἐλυμῄνω τὸ κάλλος σου· καὶ διήγαγες τὰ σκέλη σου παντὶ παρόδῳ καὶ ἐπλήθυνας τὴν πορνείαν σου,
\vs{26}καὶ ἐξεπόρνευσας ἐπὶ τοὺς υἱοὺς Αἰγύπτου τοὺς ὁμοροῦντάς σοι τοὺς μεγαλοσάρκους, καὶ πολλαχῶς ἐξεπόρνευσας τοῦ παροργίσαι με.

\vs{27}Ἐὰν δὲ ἐκτείνω τὴν χεῖρά μου ἐπὶ σὲ, καὶ ἐξαρῶ τὰ νόμιμά σου, καὶ παραδώσω εἰς ψυχὰς μισούντων σε, θυγατέρας ἀλλοφύλων τὰς ἐκκλινούσας σε ἐκ τῆς ὁδοῦ σου ἧς ἠσέβησας.

\vs{28}Καὶ ἐξεπόρνευσας ἐπὶ τὰς θυγατέρας Ἀσσοὺρ, καὶ οὐδʼ οὕτως ἐνεπλήσθης, καὶ ἐξεπόρνευσας καὶ οὐκ ἐνεπίπλω.
\vs{29}Καὶ ἐπλήθυνας τὰς διαθήκας σου πρὸς γῆν Χαλδαίων, καὶ οὐδε ἐν τούτοις ἐνεπλήσθης.

\vs{30}Τί διαθῶ τὴν θυγατέρα σου, λέγει Κύριος, ἐν τῷ ποιῆσαί σε πάντα ταῦτα ἔργα γυναικὸς πόρνης; καὶ ἐξεπόρνευσας τρισσῶς ἐν ταῖς θυγατράσι σου·
\vs{31}Τὸ πορνεῖον ᾠκοδόμησας ἐν πάσῃ ἀρχῇ ὁδοῦ, καὶ τὴν βάσιν σου ἐποίησας ἐν πάσῃ πλατείᾳ· καὶ ἐγένου ὡς πόρνη συνάγουσα μισθώματα.
\vs{32}Ἡ γυνὴ ἡ μοιχωμένη ὁμοία σοι, παρὰ τοῦ ἀνδρὸς αὐτῆς λαμβάνουσα μισθώματα, πᾶσι τοῖς ἐκπορνεύσασαιν αὐτὴν προσεδίδου μισθώματα·
\vs{33}καὶ σὺ δέδωκας μισθώματα πᾶσι τοῖς ἐρασταῖς σου, καὶ ἐφόρτιζες αὐτοὺς τοῦ ἔρχεσθαι πρὸς σὲ κυκλόθεν ἐν τῇ πορνείᾳ σου.

\vs{34}Καὶ ἐγένετο ἐν σοὶ διεστραμμένον παρὰ τὰς γυναῖκας ἐν τῇ πορνείᾳ σου, καὶ μετὰ σοῦ πεπορνεύκασιν, ἐν τῷ προσδιδόναι σε μισθώματα, καὶ σοὶ μισθώματα οὐκ ἐδόθη, καὶ ἐγένετο ἐν σοὶ διεστραμμένα.

\vs{35}Διατοῦτο πόρνη ἄκουε λόγον Κυρίου.
\vs{36}Τάδε λέγει Κύριος, ἀνθʼ ὧν ἐξέχεας τὸν χαλκόν σου, καὶ ἀποκαλυφθήσεται ἡ αἰσχύνη σου ἐν τῇ πορνείᾳ σου πρὸς τοὺς ἐραστάς σου, καὶ εἰς πάντα τὰ ἐνθυμήματα τῶν ἀνομιῶν σου, καὶ ἐν τοῖς αἵμασι τῶν τέκνων σου ὧν ἔδωκας αὐτοῖς·
\vs{37}Διατοῦτο ἰδοὺ ἐγὼ ἐπισυνάγω πάντας τοὺς ἐραστάς σου ἐν οἷς ἐπεμίγης ἐν αὐτοῖς, καὶ πάντας οὓς ἠλάπησας, σὺν πᾶσιν οἷς ἐμίσεις, καὶ συνάξω αὐτοὺς ἐπὶ σὲ κυκλόθεν, καὶ ἀποκαλύψω τὰς κακίας σου πρὸς αὐτούς, καὶ ὄψονται πᾶσαν τὴν αἰσχύνην σου,
\vs{38}καὶ ἐκδικήσω σε ἐκδικήσει μοιχαλίδος, καὶ θήσω σε ἐν αἵματι θυμοῦ καὶ ζήλου.
\vs{39}Καὶ παραδώσω σε εἰς χεῖρας αὐτῶν, καὶ κατασκάψουσι τὸ πορνεῖόν σου, καὶ καθελοῦσι τὴν βάσιν σου, καὶ ἐκδύσουσί σε τὰ ἱμάτιά σου, καὶ λήψονται τὰ σκεύη τῆς καυχήσεώς σου· καὶ ἀφήσουσί σε γυμνὴν, καὶ ἀσχημονοῦσαν.
\vs{40}Καὶ ἄξουσιν ἐπὶ σὲ ὄχλους, καὶ λιθοβολήσουσί σε ἐν λίθοις, καὶ κατασφάξουσί σε ἐν τοῖς ξίφεσιν αὐτῶν.
\vs{41}Καὶ ἐυπρήσουσι τοὺς οἴκους σου πυρὶ, καὶ ποιήσουσιν ἐν σοὶ ἐκδικήσεις ἐνώπιον γυναικῶν πολλῶν· καὶ ἀποστρέψω σε ἐκ πορνείας, καὶ μισθώματα οὐ μὴ δώσω οὐκέτι.

\vs{42}Καὶ ἐπαφήσω τὸν θυμόν μου ἐπὶ σὲ, καὶ ἐξαρθήσεται ὁ ζῆλός μου ἐκ σοῦ· καὶ ἀναπαύσομαι, καὶ οὐ μὴ μεριμνήσω οὐκέτι.
\vs{43}Ἀνθʼ ὧν οὐκ ἐμνήσθης τῆς νηπιότητός σου, καὶ ἐλύπεις με ἐν πᾶσι τούτοις· καὶ ἰδοὺ ἐγὼ τὰς ὁδούς σου εἰς κεφαλήν σου δέδωκα, λέγει Κύριος· καὶ οὕτως ἐποίησας τὴν ἀσέβειαν ἐπὶ πάσαις ταῖς ἀνομίαις σου.

\vs{44}Ταῦτά ἐστι πάντα ὅσα εἶπαν κατὰ σοῦ ἐν παραβολῇ, λέγοντες, καθὼς ἡ μήτηρ,
\vs{45}καὶ ἡ θυγάτηρ θυγάτηρ τῆς μητρός σου· σὺ εἶ ἡ ἀπωσαμένη τὸν ἄνδρα αὐτῆς καὶ τὰ τέκνα αὐτῆς, καὶ ἀδελφαὶ τῶν ἀδελφῶν σου αἱ ἀπωσάμεναι τοὺς ἄνδρας αὐτῶν καὶ τὰ τέκνα αὐτῶν· ἡ μήτηρ ὑμῶν Χετταία, καὶ ὁ πατὴρ Ἀμοῤῥαῖος.
\vs{46}Ἡ ἀδελφὴ ὑμῶν ἡ πρεσβυτέρα Σαμάρεια, αὐτὴ καὶ αἱ θυγατέρες αὐτῆς, ἡ κατοικοῦσα ἐξ εὐωνύμων σου· καὶ ἡ ἀδελφή σου ἡ νεωτέρα σου, ἡ κατοικοῦσα ἐκ δεξιῶν σου, Σόδομα καὶ αἱ θυγατέρες αὐτῆς.
\vs{47}Καὶ οὐδʼ ὡς ἐν ταῖς ὁδοῖς αὐτῶν ἐπορεύθης, οὐδὲ κατὰ τὰς ἀνομίας αὐτῶν ἐποίησας παρὰ μικρὸν, καὶ ὑπέρκεισαι αὐτὰς ἐν πάσαις ταῖς ὁδοῖς σου.

\vs{48}Ζῶ ἐγὼ, λέγει Κύριος, εἰ πεποίηκε Σόδομα αὕτη καὶ αἱ θυγατέρες αὐτῆς, ὃν τρόπον ἐποίησας σὺ καὶ αἱ θυγατέρες σου.
\vs{49}Πλὴν τοῦτο τὸ ἀνόμημα Σοδόμων τῆς ἀδελφῆς σου, ὑπερηφανία, ἐν πλησμονῇ ἄρτων καὶ ἐν εὐθηνίᾳ ἐσπατάλων αὕτη καὶ αἱ θυγατέρες αὐτῆς· τοῦτο ὑπῆρχεν αὐτῇ καὶ ταῖς θυγατράσιν αὐτῆς, καὶ χεῖρα πτωχοῦ καὶ πένητος οὐκ ἀντελαμβάνοντο.
\vs{50}Καὶ ἐμεγαλαύχουν, καὶ ἐποίησαν ἀνομήματα ἐνώπιον ἐμοῦ· καὶ ἐξῇρα αὐτὰς καθὼς ἴδον.

\vs{51}Καὶ Σαμάρεια κατὰ τὰς ἡμίσεις τῶν ἁμαρτιῶν σου οὐχ ἥμαρτε· καὶ ἐπλήθυνας τὰς ἀνομίας σου ὑπὲρ αὐτὰς, καὶ ἐδικαίωσας τὰς ἀδελφάς σου ἐν πάσαις ταῖς ἀνομίαις σου αἷς ἐποίησας.
\vs{52}Καὶ σὺ κόμισαι βάσανόν σου, ἐν ᾗ ἔφθειρας τὰς ἀδελφάς σου ἐν ταῖς ἁμαρτίαις σου αἷς ἠνόμησας ὑπὲρ αὐτὰς, καὶ ἐδικαίωσας αὐτὰς ὑπὲρ σεαυτήν· καὶ σὺ αἰσχύνθητι, καὶ λάβε τὴν ἀτιμίαν σου ἐν τῷ δικαιῶσαί σε τὰς ἀδελφάς σου.
\vs{53}Καὶ ἀποστρέψω τὰς ἀποστροφὰς αὐτῶν, τὴν ἀποστροφὴν Σοδόμων καὶ τῶν θυγατέρων αὐτῆς· καὶ ἀποστρέψω τὴν ἀποστροφὴν Σαμαρείας καὶ τῶν θυγατέρων αὐτῆς· καὶ ἀποστρέψω τὴν ἀποστροφήν σου ἐν μέσῳ αὐτῶν·
\vs{54}Ὅπως κομίσῃ τὴν βάσανόν σου, καὶ ἀτιμωθήσῃ ἐκ πάντων ὧν ἐποίησας ἐν τῷ παροργίσαι με.

\vs{55}Καὶ ἡ ἀδελφή σου Σόδομα καὶ αἱ θυγατέρες αὐτῆς ἀποκατασταθήσονται καθὼς ἦσαν ἀπʼ ἀρχῆς· καὶ σὺ καὶ αἱ θυγατέρες σου ἀποκατασταθήσεσθε καθὼς ἀπʼ ἀρχῆς ἦτε.

\vs{56}Καὶ εἰ μὴ ἦν Σόδομα ἡ ἀδελφή σου εἰς ἀκοὴν ἐν τῷ στόματί σου ἐν ταῖς ἡμέραις ὑπερηφανίας σου·
\vs{57}πρὸ τοῦ ἀποκαλυφθῆναι τὰς κακίας σου, ὃν τρόπον νῦν ὄνειδος εἶ θυγατέρων Συρίας, καὶ πάντων τῶν κύκλῳ αὐτῆς θυγατέρων ἀλλοφύλων τῶν περιεχουσῶν σε κύκλῳ.
\vs{58}Τὰς ἀσεβείας σου καὶ τὰς ἀνομίας σου σὺ κεκόμισαι αὐτὰς, λέγει Κύριος.

\vs{59}Τάδε λέγει Κύριος, καὶ ποιήσω ἐν σοὶ καθὼς ἐποίησας, ὡς ἠτίμωσας ταῦτα τοῦ παραβῆναι τὴν διαθήκην μου.
\vs{60}Καὶ μνησθήσομαι ἐγὼ τῆς διαθήκης μου τῆς μετὰ σοῦ ἐν ἡμέραις νηπιότητός σου, καὶ ἀναστήσω σοι διαθήκην αἰώνιον.
\vs{61}Καὶ μνησθήσῃ τὴν ὁδόν σου, καὶ ἐξατιμωθήσῃ ἐν τῷ ἀναλαβεῖν σε τὰς ἀδελφάς σου τὰς πρεσβυτέρας σου σὺν ταῖς νεωτέραις σου, καὶ δώσω αὐτάς σοι εἰς οἰκοδομὴν, καὶ οὐκ ἐκ διαθήκης σου.
\vs{62}Καὶ ἀναστήσω ἐγὼ τὴν διαθήκην μου μετὰ σοῦ, καὶ ἐπιγνῶσῃ ὅτι ἐγὼ Κύριος·
\vs{63}Ὅπως μνησθῇς καὶ αἰσχυνθῇς, καὶ μὴ ᾖ σοι ἔτι ἀνοῖξαι τὸ στόμα σου ἀπὸ προσώπου τῆς ἀτιμίας σου, ἐν τῷ ἐξιλάσκεσθαί με σοὶ κατὰ πάντα ὅσα ἐποίησας, λέγει Κύριος.

\ch{17}
Καὶ ἐγένετο λόγος Κυρίου πρὸς μὲ, λέγων,
\vs{2}υἱὲ ἀνθρώπου, διήγησαι διήγημα καὶ εἰπὸν παραβολὴν πρὸς τὸν οἶκον τοῦ Ἰσραὴλ,
\vs{3}καὶ ἐρεῖς, τάδε λέγει Κύριος,

Ἀετὸς ὁ μέγας ὁ μεγαλοπτέρυγος, ὁ μακρὸς τῇ ἐκτάσει, πλήρης ὀνύχων, ὃς ἔχει τὸ ἥγημα εἰσελθεῖν εἰς τὸν Λίβανον, καὶ ἔλαβε τὰ ἐπίλεκτα τῆς κέδρου,
\vs{4}τὰ ἄκρα τῆς ἁπαλότητος ἀπέκνισε, καὶ ἤνεγκεν αὐτὰ εἰς γῆν Χαναὰν, εἰς πόλιν τετειχισμένην ἔθετο αὐτά.
\vs{5}Καὶ ἔλαβεν ἀπὸ τοῦ σπέρματος τῆς γῆς, καὶ ἔδωκεν αὐτὸ εἰς τὸ πεδίον φυτὸν ἐφʼ ὕδατι πολλῷ, ἐπιβλεπόμενον ἐταξεν αὐτό.
\vs{6}Καὶ ἀνέτειλε, καὶ ἐγένετο εἰς ἄμπελον ἀσθενοῦσαν καὶ μικρὰν, τοῦ ἐπιφαίνεσθαι αὐτὴν τὰ κλήματα αὐτῆς ἐπʼ αὐτό· καὶ ῥίζαι αὐτῆς ὑποκάτω αὐτῆς ἦσαν· καὶ ἐγένετο εἰς ἄμπελον, καὶ ἐποίησεν ἀπώρυγας, καὶ ἐξέτεινε τὴν ἀναδενδράδα αὐτῆς.

\vs{7}Καὶ ἐγένετο ἀετὸς ἕτερος μέγας, μεγαλοπτέρυγος πολὺς ὄνυξι· καὶ ἰδοὺ ἡ ἄμπελος αὕτη περιπεπλεγμένη πρὸς αὐτὸν, καὶ ῥίζαι αὐτῆς πρὸς αὐτὸν, καὶ τὰ κλήματα αὐτῆς ἐξαπέστειλεν αὐτῷ, τοῦ ποτίσαι αὐτὴν σὺν τῷ βόλῳ τῆς φυτείας αὐτῆς.
\vs{8}Εἰς πεδίον καλὸν ἐφʼ ὕδατι πολλῶ αὕτη πιαίνεται, τοῦ ποιεῖν βλαστοὺς, καὶ φέρειν καρπὸν, τοῦ εἶναι εἰς ἄμπελον μεγάλην.

\vs{9}Διατοῦτο εἰπὸν, τάδε λέγει Κύριος, εἰ κατευθυνεῖ; οὐχὶ αἱ ῥίζαι τῆς ἁπαλότητος αὐτῆς καὶ ὁ καρπὸς σαπήσεται; καὶ ξηρανθήσεται πάντα τὰ προανατέλλοντα αὐτῆς, καὶ οὐκ ἐν βραχίονι μεγάλῳ, οὐδὲ ἐν λαῷ πολλῷ, τοῦ ἐκσπᾶσαι αὐτὴν ἐκ ῥιζῶν αὐτῆς.
\vs{10}Καὶ ἰδοὺ πιαίνεται· μὴ κατευθυνεῖ; οὐχὶ ἅμα τῷ ἅψεσθαι αὐτῆς ἄνεμον τὸν καύσωνα ξηρανθήσεται; σὺν τῷ βόλῳ ἀνατολῆς αὐτῆς ξηρανθήσεται.

\vs{11}Καὶ ἐγένετο λόγος Κυρίου πρὸς μὲ, λέγων,
\vs{12}υἱὲ ἀνθρώπου, εἰπὸν δὴ πρὸς τὸν οἶκον τὸν παραπικραίνοντα, οὐκ ἐπίστασθε τί ἦν ταῦτα; εἰπὸν, ὅταν ἔλθῃ βασιλεὺς Βαβυλῶνος ἐπὶ Ἱερουσαλὴμ, καὶ λήψεται τὸν βασιλέα αὐτῆς καὶ τοὺς ἄρχοντας αὐτῆς, καὶ ἄξῃ αὐτοὺς πρὸς ἑαυτὸν εἰς Βαβυλῶνα.
\vs{13}Καὶ λήψεται ἐκ τοῦ σπέρματος τῆς βασιλείας, καὶ διαθήσεται πρὸς αὐτὸν διαθήκην, καὶ εἰσάξει αὐτὸν ἐν ἀρᾷ καὶ τοὺς ἡγεμόνας τῆς γῆς λήψεται,
\vs{14}τοῦ γενέσθαι εἰς βασιλείαν ἀσθενῆ, τὸ καθόλου μὴ ἐπαίρεσθαι, τοῦ φυλάσσειν τὴν διαθήκην αὐτοῦ, καὶ ἱστάνειν αὐτήν.
\vs{15}Καὶ ἀποστήσεται ἀπʼ αὐτοῦ, τοῦ ἐξαποστέλλειν ἀγγέλους ἑαυτοῦ εἰς Αἴγυπτον, τοῦ δοῦναι αὐτῷ ἵππους καὶ λαὸν πολύν· εἰ κατευθυνεῖ, εἰ διασωθήσεται ὁ ποιῶν ἐναντία; καὶ παραβαίνων διαθήκην εἰ διασωθήσεται;

\vs{16}Ζῶ ἐγὼ, λέγει Κύριος, ἐὰν μὴ ἐν τόπῳ ὁ βασιλεὺς ὁ βασιλεύσας αὐτὸν, ὃς ἠτίμωσε τὴν ἀράν μου, καὶ ὃς παρέβη τὴν διαθήκην μου, μετʼ αὐτοῦ ἐν μέσῳ Βαβυλῶνος τελευτήσει,
\vs{17}καὶ οὐκ ἐν δυνάμει μεγάλῃ οὐδὲ ἐν ὄχλῳ πολλῷ ποιήσει πρὸς αὐτὸν Φαραὼ πόλεμον, ἐν χαρακοβολίᾳ καὶ ἐν οἰκοδομῇ βελοστάσεων, τοῦ ἐξᾶραι ψυχάς.
\vs{18}Καὶ ἠτίμωσεν ὁρκωμοσίαν τοῦ παραβῆναι διαθήκην· καὶ ἰδοὺ δέδωκα τὴν χεῖρα αὐτοῦ, καὶ πάντα ταῦτα ἐποίησεν αὐτῷ, μὴ σωθήσεται.

\vs{19}Διατοῦτο εἰπὸν, τάδε λέγει Κύριος, ζῶ ἐγὼ, ἐὰν μὴ τὴν ὁρκωμοσίαν μου ἣν ἠτίμωσε, καὶ τὴν διαθήκην μου ἣν παρέβη, καὶ δώσω αὐτὴν εἰς κεφαλὴν αὐτοῦ.
\vs{20}Καὶ ἐκπετάσω ἐπʼ αὐτὸν τὸ δίκτυον, καὶ ἁλώσεται ἐν τῇ περιοχῇ αὐτοῦ.
\vs{21}Ἐν πάσῃ παρατάξει αὐτοῦ ἐν ῥομφαίᾳ πεσοῦνται· καὶ τοὺς καταλοίπους εἰς πάντα ἄνεμον διασπερῶ, καὶ ἐπιγνώσεσθε διότι ἐγὼ Κύριος λελάληκα.

\vs{22}Διότι τάδε λέγει Κύριος, καὶ λήψομαι ἐγὼ ἐκ τῶν ἐκλεκτῶν τῆς κέδρου ἐκ κορυφῆς, καρδίας αὐτῶν ἀποκνιῶ, καὶ καταφυτεύσω ἐγὼ ἐπʼ ὄρος ὑψηλὸν,
\vs{23}καὶ κρεμάσω αὐτὸν ἐν ὄρει μετεώρῳ Ἰσραήλ· καὶ καταφυτεύσω, καὶ ἐξοίσει βλαστὸν, καὶ ποιήσει καρπὸν, καὶ ἔσται εἰς κέδρον μεγάλην· καὶ ἀναπαύσεται ὑποκάτω αὐτοῦ πᾶν ὄρνεον, καὶ πᾶν πετεινὸν ὑπὸ τὴν σκιὰν αὐτοῦ ἀναπαύσεται· τὰ κλήματα αὐτοῦ ἀποκατασταθήσεται.
\vs{24}Καὶ γνώσονται πάντα τὰ ξύλα τοῦ πεδίου, διότι ἐγὼ Κύριος ὁ ταπεινῶν ξύλον ὑψηλὸν, καὶ ὑψῶν ξύλον ταπεινὸν, καὶ ξηραίνων ξύλον χλωρὸν, καὶ ἀναθάλλων ξύλον ξηρόν· ἐγὼ Κύριος λελάληκα, καὶ ποιήσω.

\ch{18}
Καὶ ἐγένετο λόγος Κυρίου πρὸς μὲ, λέγων,
\vs{2}υἱὲ ἀνθρώπου, τί ὑμῖν ἡ παραβολὴ αὕτη ἐν τοῖς υἱοῖς Ἰσραὴλ, λέγοντες, οἱ πατέρες ἔφαγον ὄμφακα, καὶ οἱ ὀδόντες τῶν τέκνων ἐγομφίασαν;

\vs{3}Ζῶ ἐγὼ, λέγει Κύριος, ἐὰν γένηται ἔτι λεγομένη ἡ παραβολὴ αὕτη ἐν τῷ Ἰσραήλ·
\vs{4}Ὅτι πᾶσαι αἱ ψυχαὶ ἐμαί εἰσιν, ὃν τρόπον ἡ ψυχὴ τοῦ πατρὸς, οὕτως καὶ ἡ ψυχὴ τοῦ υἱοῦ, ἐμαί εἰσιν· ἡ ψυχὴ ἡ ἁμαρτάνουσα, αὕτη ἀποθανεῖται.

\vs{5}Ὁ δὲ ἄνθρωπος ὃς ἔσται δίκαιος, ὁ ποιῶν κρίμα καὶ δικαιοσύνην,
\vs{6}ἐπὶ τῶν ὀρέων οὐ φάγεται, καὶ τοὺς ὀφθαλμοὺς αὐτοῦ οὐ μὴ ἐπάρῃ πρὸς τὰ ἐνθυμήματα οἴκου Ἰσραὴλ, καὶ τὴν γυναῖκα τοῦ πλησίον αὐτοῦ οὐ μὴ μιάνῃ, καὶ πρὸς γυναῖκα ἐν ἀφέδρῳ οὖσαν οὐ προσεγγιεῖ,
\vs{7}καὶ ἄνθρωπον οὐ μὴ καταδυναστεύσῃ, ἐνεχυρασμὸν ὀφείλοντος ἀποδώσει, καὶ ἅρπαγμα οὐχ ἁρπᾶται, τὸν ἄρτον αὐτοῦ τῷ πεινῶντι δώσει, καὶ γυμνὸν περιβαλεῖ,
\vs{8}καὶ τὸ ἀργύριον αὐτοῦ ἐπὶ τόκῳ οὐ δώσει, καὶ πλεονασμὸν οὐ λήψεται, καὶ ἐξ ἀδικίας ἀποστρέψει τὴν χεῖρα αὐτοῦ, κρίμα δίκαιον ποιήσει ἀναμέσον ἀνδρὸς καὶ ἀναμέσον τοῦ πλησίον αὐτοῦ,
\vs{9}καὶ τοῖς προστάγμασί μου πεπόρευται, καὶ τὰ δικαιώματά μου πεφύλακται, τοῦ ποιῆσαι αὐτά· δίκαιος οὗτός ἐστι, ζωῇ ζήσεται, λέγει Κύριος.

\vs{10}Καὶ ἐὰν γεννήσῃ υἱὸν λοιμὸν, ἐκχέοντα αἷμα καὶ ποιοῦντα ἁμαρτήματα,
\vs{11}ἐν τῇ ὁδῷ τοῦ πατρὸς αὐτοῦ τοῦ δικαίου οὐκ ἐπορεύθη, ἀλλὰ καὶ ἐπὶ τῶν ὀρέων ἔφαγε, καὶ τὴν γυναῖκα τοῦ πλησίον αὐτοῦ ἐμίανε,
\vs{12}καὶ πτωχὸν καὶ πένητα κατεδυνάστευσε, καὶ ἅρπαγμα ἥρπασε, καὶ ἐνεχυρασμὸν οὐκ ἀπέδωκε, καὶ εἰς τὰ εἴδωλα ἔθετο τοὺς ὀφθαλμοὺς αὐτοῦ, ἀνομίαν πεποίηκε,
\vs{13}μετὰ τόκου ἔδωκε, καὶ πλεονασμὸν ἔλαβεν· οὗτος ζωῇ οὐ ζήσεται, πάσας τὰς ἀνομίας ταύτας ἐποίησε, θανάτῳ θανατωθήσεται· τὸ αἷμα αὐτοῦ ἐπʼ αὐτὸν ἔσται.

\vs{14}Ἐὰν δὲ γεννήσῃ υἱόν, καὶ ἴδῃ πάσας τὰς ἁμαρτίας τοῦ πατρὸς αὐτοῦ ἃς ἐποίησε, καὶ φοβηθῇ, καὶ μὴ ποιήσῃ κατʼ αὐτὰς,
\vs{15}ἐπὶ τῶν ὀρέων οὐ βέβρωκε, καὶ τοὺς ὀφθαλμοὺς αὐτοῦ οὐκ ἔθετο εἰς τὰ ἐνθυμήματα οἴκου Ἰσραὴλ, καὶ τὴν γυναῖκα τοῦ πλησίον αὐτοῦ οὐκ ἐμίανε,
\vs{16}καὶ ἄνθρωπον οὐ κατεδυνάστευσε, καὶ ἐνεχυρασμὸν οὐκ ἐνεχύρασε, καὶ ἅρπαγμα οὐχ ἥρπασε, τὸν ἄρτον αὐτοῦ τῷ πεινῶντι ἔδωκε, καὶ γυμνὸν περιέβαλε,
\vs{17}καὶ ἀπὸ ἀδικίας ἀπέστρεψε τὴν χεῖρα αὐτοῦ, τόκον οὐδὲ πλεονασμὸν οὐκ ἔλαβε, δικαιοσύνην ἐποίησε, καὶ ἐν τοῖς προστάγμασί μου ἐπορεύθη, οὐ τελευτήσει ἐν ἀδικίαις πατρὸς αὐτοῦ, ζωῇ ζήσεται.
\vs{18}Ὁ δὲ πατὴρ αὐτοῦ ἐὰν θλίψει θλίψῃ καὶ ἁρπάσῃ ἅρπαγμα, ἐναντία ἐποίησεν ἐν μέσῳ τοῦ λαοῦ μου, καὶ ἀποθανεῖται ἐν τῇ ἀδικίᾳ αὐτοῦ.

\vs{19}Καὶ ἐρεῖτε, τί ὅτι οὐκ ἔλαβε τὴν ἀδικίαν ὁ υἱὸς τοῦ πατρός; ὅτι ὁ υἱὸς δικαιοσύνην καὶ ἔλεος πεποίηκε, πάντα τὰ νόμιμά μου συνετήρησε, καὶ ἐποίησεν αὐτὰ, ζωῇ ζήσεται.
\vs{20}Ἡ δὲ ψυχὴ ἡ ἁμαρτάνουσα, ἀποθανεῖται· ὁ δὲ υἱὸς οὐ λήψεται τὴν ἀδικίαν τοῦ πατρὸς, οὐδὲ ὁ πατὴρ λήψεται τὴν ἀδικίαν τοῦ υἱοῦ· δικαιοσύνη δικαίῳ ἐπʼ αὐτὸν ἔσται, καὶ ἀνομία ἀνόμῳ ἐπʼ αὐτὸν ἔσται.

\vs{21}Καὶ ὁ ἄνομος ἐὰν ἀποστρέψῃ ἐκ πασῶν τῶν ἀνομιῶν αὐτοῦ ὧν ἐποίησε, καὶ φυλάξηται πάσας τὰς ἐντολάς μου, καὶ ποιήσῃ δικαιοσύνην καὶ ἔλεος, ζωῇ ζήσεται, καὶ οὐ μὴ ἀποθάνῃ.
\vs{22}Πάντα τὰ παραπτώματα αὐτοῦ ὅσα ἐποίησεν, οὐ μνησθήσονται· ἐν τῇ δικαιοσύνῃ αὐτοῦ ᾗ ἐποίησε, ζήσεται.
\vs{23}Μὴ θελήσει θελήσω τὸν θάνατον τοῦ ἀνόμου, λέγει Κύριος, ὡς τὸ ἀποστρέψαι αὐτὸν ἐκ τῆς ὁδοῦ τῆς πονηρᾶς, καὶ ζῇν αὐτόν;

\vs{24}Ἐν δὲ τῶ ἀποστρέψαι δίκαιον ἐκ τῆς δικαιοσύνης αὐτοῦ, καὶ ποιῆσαι ἀδικίαν κατὰ πάσας τὰς ἀνομίας ἃς ἐποίησεν ὁ ἄνομος, πᾶσαι αἱ δικαιοσύναι αὐτοῦ ἃς ἐποίησεν, οὐ μὴ μνησθῶσιν· ἐν τῷ παραπτώματι αὐτοῦ ᾧ παρέπεσε, καὶ ἐν ταῖς ἁμαρτίαις αὐτοῦ αἷς ἥμαρτεν, ἐν αὐταῖς ἀποθανεῖται.

\vs{25}Καὶ εἰπατε, οὐ κατευθύνει ἡ ὁδὸς Κυρίου· ἀκούσατε δὴ πᾶς ὁ οἶκος Ἰσραήλ· μὴ ἡ ὁδός μου οὐ κατευθυνεῖ; οὐχὶ ἡ ὁδὸς ὑμῶν κατευθύνει;
\vs{26}Ἐν τῷ ἀποστρέψαι τὸν δίκαιον ἐκ τῆς δικαιοσύνης αὐτοῦ, καὶ ποιήσει παράπτωμα, καὶ ἀποθάνῃ ἐν τῷ παραπτώματι ᾧ ἐποίησεν, ἐν αὐτῷ ἀποθανεῖται.
\vs{27}Καὶ ἐν τῷ ἀποστρέψαι ἄνομον ἀπὸ τῆς ἀνομίας αὐτοῦ ἧς ἐποίησε, καὶ ποιήσει κρίμα καὶ δικαιοσύνην, οὗτος τὴν ψυχὴν αὐτοῦ ἐφύλαξε,
\vs{28}καὶ ἀπέστρεψεν ἐκ πασῶν ἀσεβειῶν αὐτοῦ ὧν ἐποίησε, ζωῇ ζήσεται, οὐ μὴ ἀποθάνῃ.

\vs{29}Καὶ λέγουσιν ὁ οἶκος τοῦ Ἰσραὴλ, οὐ κατορθοῖ ἡ ὁδὸς Κυρίου· μὴ ἡ ὁδός μου οὐ κατορθοῖ, οἶκος Ἰσραήλ; οὐχὶ ἡ ὁδὸς ὑμῶν οὐ κατορθοῖ;
\vs{30}Ἕκαστον κατὰ τὴν ὁδὸν αὐτοῦ κρινῷ ὑμᾶς οἶκος Ἰσραὴλ, λέγει Κύριος· ἐπιστράφητε καὶ ἀποστρέψατε ἐκ πασῶν τῶν ἀσεβειῶν ὑμῶν, καὶ οὐκ ἔσονται ὑμῖν εἰς κόλασιν ἀδικίας.
\vs{31}Ἀποῤῥίψατε ἀφʼ ἑαυτῶν πάσας τὰς ἀσεβείας ὑμῶν ἃς ἠσεβήσατε εἰς ἐμὲ, καὶ ποιήσατε ἑαυτοῖς καρδίαν καινὴν καὶ πνεῦμα καινόν· καὶ ἱνατί ἀποθνήσκετε, οἶκος Ἰσραήλ;
\vs{32}Διότι οὐ θέλω τὸν θάνατον τοῦ ἀποθνήσκοντος, λέγει Κύριος.

\ch{19}
Καὶ σὺ λάβε θρῆνον ἐπὶ τὸν ἄρχοντα τοῦ Ἰσραὴλ,
\vs{2}καὶ ἐρεῖς, τί ἡ μήτηρ σου σκύμνος ἐν μέσῳ λεόντων ἐγενήθη; ἐν μέσῳ λεόντων ἐπλήθυεν σκύμνους αὐτῆς.
\vs{3}Καὶ ἀπεπήδησεν εἷς τῶν σκύμνων αὐτῆς, λέων ἐγένετο, καὶ ἔμαθε τοῦ ἁρπάζειν ἁρπάγματα, ἀνθρώπους ἔφαγε.
\vs{4}Καὶ ἤκουσαν κατʼ αὐτοῦ ἔθνη, ἐν τῇ διαφθορᾷ αὐτῶν συνελήφθη, καὶ ἤγαγον αὐτὸν ἐν κημῷ εἰς γῆν Αἰγύπτου.

\vs{5}Καὶ εἶδεν ὅτι ἀπῶσται ἀπʼ αὐτῆς, ἀπώλετο ἡ ὑπόστασις αὐτῆς· καὶ ἔλαβεν ἄλλον ἐκ τῶν σκύμνων αὐτῆς· λέοντα ἔταξεν αὐτὸν.
\vs{6}Καὶ ἀνεστρέφετο ἐν μέσῳ λεόντων, λέων ἐγένετο, καὶ ἔμαθεν ἁρπάζειν ἁρπάγματα, ἀνθρώπους ἔφαγε,
\vs{7}καὶ ἐνέμετο τῷ θράσει αὐτοῦ, καὶ τὰς πόλεις αὐτῶν ἐξηρήμωσε, καὶ ἠφάνισε γῆν, καὶ τὸ πλήρωμα αὐτῆς ἀπὸ φωνῆς ὠρυώματος αὐτοῦ.

\vs{8}Καὶ ἔδωκαν ἐπʼ αὐτὸν ἔθνη ἐκ χωρῶν κυκλόθεν, καὶ ἐξεπέτασαν ἐπʼ αὐτὸν δίκτυα αὐτῶν, ἐν διαφθορᾷ αὐτῶν συνελήφθη·
\vs{9}Καὶ ἔθεντο αὐτὸν ἐν κημῷ καὶ ἐν γαλεάγρᾳ, ἦλθε πρὸς βασιλέα Βαβυλῶνος, καὶ εἰσήγαγεν αὐτὸν εἰς φυλακὴν, ὅπως μὴ ἀκουσθῇ ἡ φωνὴ αὐτοῦ ἐπὶ τὰ ὄρη τοῦ Ἰσραήλ.

\vs{10}Ἡ μήτηρ σου ὡς ἄμπελος καὶ ὡς ἄνθος ἐν ῥοᾷ, ἐν ὕδατι πεφυτευμένη, ὁ καρπὸς αὐτῆς καὶ ὁ βλαστὸς αὐτῆς ἐγένετο ἐξ ὕδατος πολλοῦ.
\vs{11}Καὶ ἐγένετο αὕτη ῥάβδος ἐπὶ φυλὴν ἡγουμένων, καὶ ὑψώθη τῷ μεγέθει αὐτῆς ἐν μέσῳ στελεχῶν· καὶ εἶδε τὸ μέγεθος αὐτῆς ἐν πλήθει κλημάτων αὐτῆς.

\vs{12}Καὶ κατεκλάσθη ἐν θυμῷ, ἐπὶ γῆν ἐῤῥίφη, καὶ ἄνεμος ὁ καύσων ἐξήρανε τὰ ἐκλεκτὰ αὐτῆς· ἐξεδικήθησαν, καὶ ἐξηράνθη ἡ ῥάβδος ἰσχύος αὐτῆς· πῦρ ἀνήλωσεν αὐτήν.
\vs{13}Καὶ νῦν πεφύτευκαν αὐτὴν ἐν τῇ ἠρημῳ, ἐν γῇ ἀνύδρῳ.
\vs{14}Καὶ ἐξῆλθε πῦρ ἐκ ῥάβδου ἐκλεκτῶν αὐτῆς, καὶ κατέφαγεν αὐτὴν, καὶ οὐκ ἦν ἐν αὐτῇ ῥάβδος ἰσχύος· φυλὴ εἰς παραβολὴν θρήνου ἐστὶ, καὶ ἐσται εἰς θρῆνον.

\ch{20}
Καὶ ἐγένετο ἐν τῷ ἔτει τῷ ἑβδόμῳ τῇ πεντεκαιδεκάτῃ τοῦ μηνὸς, ἦλθον ἄνδρες ἐκ τῶν πρεσβυτέρων οἴκου Ἰσραὴλ, ἐπερωτῆσαι τὸν Κύριον, καὶ ἐκάθισαν πρὸ προσώπου μου.
\vs{2}Καὶ ἐγένετο λόγος Κυρίου πρὸς μὲ, λέγων,
\vs{3}υἱὲ ἀνθρώπου, λάλησον πρὸς τοὺς πρεσβυτέρους τοῦ οἴκου Ἰσραὴλ, καὶ ἐρεῖς πρὸς αὐτοὺς, τάδε λέγει Κύριος, εἰ ἐπερωτῆσαί με ὑμεῖς ἔρχεσθε; ζῶ ἐγὼ, εἰ ἀποκριθήσομαι ὑμῖν, λέγει Κύριος.
\vs{4}Εἰ ἐκδικήσω αὐτοὺς ἐκδικήσει, υἱὲ ἀνθρώπου; τὰς ἀνομίας τῶν πατέρων αὐτῶν διαμάρτυραι αὐτοῖς,
\vs{5}καὶ ἐρεῖς πρὸς αὐτοὺς, τάδε λέγει Κύριος,

Ἀφʼ ἧς ἡμέρας ᾑρέτισα τὸν οἶκον Ἰσραὴλ, καὶ ἐγνωρίσθην τῷ σπέρματι οἴκου Ἰακώβ, καὶ ἐγνώσθην αὐτοῖς ἐν γῇ Αἰγύπτου, καὶ ἀντελαβόμην τῇ χειρί μου αὐτῶν, λέγων, ἐγὼ Κύριος ὁ Θεὸς ὑμῶν·
\vs{6}Ἐν ἐκείνῃ τῇ ἡμέρᾳ ἀντελαβόμην τῇ χειρί μου αὐτῶν, τοῦ ἐξαγαγεῖν αὐτοὺς ἐκ γῆς Αἰγύπτου εἰς τὴν γῆν ἣν ἡτοίμασα αὐτοῖς, γῆν ῥέουσαν γάλα καὶ μέλι, κηρίον ἐστὶ παρὰ πᾶσαν τὴν γῆν.
\vs{7}Καὶ εἶπα πρὸς αὐτοὺς, ἕκαστος βδελύγματα τῶν ὀφθαλμῶν αὐτοῦ ἀποῤῥιψάτω, καὶ ἐν τοῖς ἐπιτηδεύμασιν Αἰγύπτου μὴ μιαίνεσθε, ἐγὼ Κύριος ὁ Θεὸς ὑμῶν.

\vs{8}Καὶ ἀπέστησαν ἀπʼ ἐμοῦ, καὶ οὐκ ἠθέλησαν εἰσακοῦσαί μου, τὰ βδελύγματα τῶν ὀφθαλμῶν αὐτῶν οὐκ ἀπέῤῥιψαν, καὶ τὰ ἐπιτηδεύματα Αἰγύπτου οὐκ ἐγκατέλιπον· καὶ εἶπα τοῦ ἐκχέαι τὸν θυμόν μου ἐπʼ αὐτοὺς, τοῦ συντελέσαι ὀργήν μου ἐν αὐτοῖς ἐν μέσῳ τῆς Αἰγύπτου.
\vs{9}Καὶ ἐποίησα ὅπως τὸ ὄνομά μου τὸ παράπαν μὴ βεβηλωθῇ ἐνώπιον τῶν ἐθνῶν, ὧν αὐτοί εἰσιν ἐν μέσῳ αὐτῶν, ἐν οἷς ἐγνώσθην πρὸς αὐτοὺς ἐνώπιον αὐτῶν, τοῦ ἐξαγαγεῖν αὐτοὺς ἐκ γῆς Αἰγύπτου.

\vs{10}Καὶ ἤγαγον αὐτοὺς εἰς τὴν ἔρημον,
\vs{11}καὶ ἔδωκα αὐτοῖς τὰ προστάγματά μου, καὶ τὰ δικαιώματά μου ἐγνώρισα αὐτοῖς, ὅσα ποιήσει αὐτὰ ἄνθρωπος, καὶ ζήσεται ἐν αὐτοῖς.
\vs{12}Καὶ τὰ σάββατά μου ἔδωκα αὐτοῖς, τοῦ εἶναι εἰς σημεῖον ἀναμέσον ἐμοῦ καὶ ἀναμέσον αὐτῶν, τοῦ γνῶναι αὐτοὺς διότι ἐγὼ Κύριος ὁ ἁγιάζων αὐτούς.

\vs{13}Καὶ εἶπα πρὸς τὸν οἶκον τοῦ Ἰσραὴλ ἐν τῇ ἐρήμῳ, ἐν τοῖς προστάγμασί μου πορεύεσθε· καὶ οὐκ ἐπορεύθησαν· καὶ τὰ δικαιώματά μου ἀπώσαντο, ἃ ποιήσει αὐτὰ ἄνθρωπος καὶ ζήσεται ἐν αὐτοῖς· καὶ τὰ σάββατά μου ἐβεβήλωσαν σφόδρα· καὶ εἶπα τοῦ ἐκχέαι τὸν θυμόν μου ἐπʼ αὐτοὺς ἐν τῇ ἐρήμῳ, τοῦ ἐξαναλῶσαι αὐτούς.
\vs{14}Καὶ ἐποίησα ὅπως τὸ ὄνομά μου τὸ παράπαν μὴ βεβηλωθῇ ἐνώπιον τῶν ἐθνῶν, ὧν ἐξήγαγον αὐτοὺς κατʼ ὀφθαλμοὺς αὐτῶν.

\vs{15}Καὶ ἐγὼ ἐξῇρα τὴν χεῖρά μου ἐπʼ αὐτοὺς ἐν τῇ ἐρήμῳ τὸ παράπαν, τοῦ μὴ εἰσαγαγεῖν αὐτοὺς εἰς τὴν γῆν ἣν ἔδωκα αὐτοῖς, γῆν ῥέουσαν γάλα καὶ μέλι· κηρίον ἐστὶ παρὰ πᾶσαν τὴν γῆν·
\vs{16}Ἀνθʼ ὧν τὰ δικαιώματά μου ἀπώσαντο, καὶ ἐν τοῖς προστάγμασί μου οὐκ ἐπορεύθησαν ἐν αὐτοῖς, καὶ τὰ σάββατά μου ἐβεβήλουν, καὶ ὀπίσω τῶν ἐνθυμημάτων καρδίας αὐτῶν ἐπορεύοντο.

\vs{17}Καὶ ἐφείσατο ὁ ὀφθαλμός μου ἐπʼ αὐτοὺς, τοῦ ἐξαλεῖψαι αὐτοὺς, καὶ οὐκ ἐποίησα αὐτοὺς εἰς συντέλειαν ἐν τῇ ἐρήμῳ.
\vs{18}Καὶ εἶπα πρὸς τὰ τέκνα αὐτῶν ἐν τῇ ἐρήμῳ, ἐν τοῖς νομίμοις τῶν πατέρων ὑμῶν μὴ πορεύεσθε, καὶ τὰ δικαιώματα αὐτῶν μὴ φυλάσσεσθε, καὶ ἐν τοῖς ἐπιτηδεύμασιν αὐτῶν μὴ συναναμίσγεσθε, καὶ μὴ μιαίνεσθε.
\vs{19}Ἐγὼ Κύριος ὁ Θεὸς ὑμῶν· ἐν τοῖς προστάγμασί μου πορεύεσθε, καὶ τὰ δικαιώματά μου φυλάσσεσθε, καὶ ποιεῖτε αὐτά.
\vs{20}Καὶ τὰ σάββατά μου ἁγιάζετε, καὶ ἔστω εἰς σημεῖον ἀναμέσον ἐμοῦ καὶ ὑμῶν, τοῦ γινώσκειν διότι ἐγὼ Κύριος ὁ Θεὸς ὑμῶν.

\vs{21}Καὶ παρεπίκρανάν με, καὶ τὰ τέκνα αὐτῶν ἐν τοῖς προστάγμασί μου οὐκ ἐπορεύθησαν, καὶ τὰ δικαιώματά μου οὐκ ἐφυλάξαντο τοῦ ποιεῖν αὐτὰ, ἃ ποιήσει ἄνθρωπος καὶ ζήσεται ἐν αὐτοῖς, καὶ τὰ σάββατά μου ἐβεβήλουν· καῖ εἶπα τοῦ ἐκχέαι τὸν θυμόν μου ἐπʼ αὐτοὺς ἐν τῇ ἐρήμῳ, τοῦ συντελέσαι τὴν ὀργήν μου ἐπʼ αὐτούς.
\vs{22}Καὶ ἐποίησα ὅπως τὸ ὄνομά μου τὸ παράπαν μὴ βεβηλωθῇ ἐνώπιον τῶν ἐθνῶν, καὶ ἐξήγαγον αὐτοὺς κατʼ ὀφθαλμοὺς αὐτῶν.

\vs{23}Ἐξῇρα τὴν χεῖρά μου ἐπʼ αὐτοὺς ἐν τῇ ἐρήμῳ τοῦ διασκορπίσαι αὐτοὺς ἐν τοῖς ἔθνεσι, διασπεῖραι αὐτοὺς ἐν ταῖς χώραις,
\vs{24}ἀνθʼ ὧν τὰ δικαιώματά μου οὐκ ἐποίησαν, καὶ τὰ προστάγματά μου ἀπώσαντο, καὶ τὰ σάββατά μου ἐβεβήλουν καὶ ὀπίσω τῶν ἐνθυμημάτων τῶν πατέρων αὐτῶν ἦσαν οἱ ὀφθαλμοὶ αὐτῶν.

\vs{25}Καὶ ἐγὼ ἔδωκα αὐτοῖς προστάγματα οὐ καλὰ, καὶ δικαιώματα ἐν οἷς οὐ ζήσονται ἐν αὐτοῖς.
\vs{26}Καὶ μιανῶ αὐτοὺς ἐν τοῖς δόγμασιν αὐτῶν, ἐν τῷ διαπορεύεσθαί με πᾶν διανοῖγον μήτραν, ὅπως ἀφανίσω αὐτούς.

\vs{27}Διατοῦτο λάλησον πρὸς τὸν οἶκον τοῦ Ἰσραὴλ, υἱὲ ἀνθρώπου, καὶ ἐρεῖς πρὸς αὐτοὺς, τάδε λέγει Κύριος, ἕως τούτου παρώργισάν με οἱ πατέρες ὑμῶν ἐν τοῖς παραπτώμασιν αὐτῶν ἐν οἷς παρέπεσον εἰς ἐμέ.
\vs{28}Καὶ εἰσήγαγον αὐτοὺς εἰς τὴν γῆν ἣν ᾖρα τὴν χεῖρά μου δοῦναι αὐτὴν αὐτοῖς· καὶ ἴδον πάντα βουνὸν ὑψηλὸν, καὶ πᾶν ξύλον κατάσκιον, καὶ ἔθυσαν ἐκεῖ τοῖς θεοῖς αὐτῶν, καὶ ἔταξαν ἐκεῖ ὀσμὴν εὐωδίας, καὶ ἔσπεισαν ἐκεῖ τὰς σπονδὰς αὐτῶν.
\vs{29}Καὶ εἶπον πρὸς αὐτοὺς, τί ἐστιν Ἀβαμὰ, ὅτι ὑμεῖς εἰσπορεύεσθε ἐκεῖ; καὶ ἐπεκάλεσαν τὸ ὄνομα αὐτοῦ Ἀβαμὰ, ἕως τῆς σήμερον ἡμέρας·

\vs{30}Διατοῦτο εἰπὸν πρὸς τὸν οἶκον τοῦ Ἰσραὴλ, τάδε λέγει Κύριος, εἰ ἐν ταῖς ἀνομίαις τῶν πατέρων ὑμῶν ὑμεῖς μιαίνεσθε, καὶ ὀπίσω τῶν βδελυγμάτων αὐτῶν ὑμεῖς ἐκπορνεύετε,
\vs{31}καὶ ἐν ταῖς ἀπαρχαῖς τῶν δομάτων ὑμῶν, ἐν τοῖς ἀφορισμοῖς οἷς ὑμεῖς μιαίνεσθε ἐν πᾶσι τοῖς ἐνθυμήμασιν ὑμῶν ἕως τῆς σήμερον ἡμέρας, καὶ ἐγὼ ἀποκριθῶ ὑμῖν οἶκος τοῦ Ἰσραήλ; Ζῶ ἐγὼ, λέγει Κύριος, εἰ ἀποκριθήσομαι ὑμῖν, καὶ εἰ ἀναβήσεται ἐπὶ τὸ πνεῦμα ὑμῶν τοῦτο.
\vs{32}Καὶ οὐκ ἔσται ὃν τρόπον ὑμεῖς λέγετε, ἐσόμεθα ὡς τὰ ἔθνη, καὶ ὡς αἱ φυλαὶ τῆς γῆς, τοῦ λατρεύειν ξύλοις καὶ λίθοις.

\vs{33}Διατοῦτο, ζῶ ἐγὼ, λέγει Κύριος, ἐν χειρὶ κραταιᾷ καὶ ἐν βραχίονι ὑψηλῷ καὶ ἐν θυμῷ κεχυμένῳ βασιλεύσω ἐφʼ ὑμᾶς,
\vs{34}καὶ ἐξάξω ὑμᾶς ἐκ τῶν λαῶν, καὶ εἰσδέξομαι ὑμᾶς ἐκ τῶν χωρῶν οὗ διεσκορπίσθητε ἐν αὐταῖς, ἐν χειρὶ κραταιᾷ καὶ ἐν βραχίονι ὑψηλῷ καὶ ἐν θυμῷ κεχυμένῳ.
\vs{35}Καὶ ἄξω ὑμᾶς εἰς τὴν ἔρημον τῶν λαῶν, καὶ διακριθήσομαι πρὸς ὑμᾶς ἐκεῖ πρόσωπον κατὰ προσώπον.

\vs{36}Ὃν τρόπον διεκρίθην πρὸς τοὺς πατέρας ὑμῶν ἐν τῇ ἐρήμῳ γῆς Αἰγύπτου, οὕτως κρινῶ ὑμᾶς, λέγει Κύριος.
\vs{37}Καὶ διάξω ὑμᾶς ὑπὸ τὴν ῥάβδον μου, καὶ εἰσάξω ὑμᾶς ἐν ἀριθμῷ.
\vs{38}Καὶ ἐκλέξω ἐξ ὑμῶν τοὺς ἀσεβεῖς καὶ τοὺς ἀφεστηκότας, διότι ἐκ τῆς παροικεσίας αὐτῶν ἐξάξω αὐτοὺς, καὶ εἰς τὴν γῆν τοῦ Ἰσραὴλ οὐκ εἰσελεύσονται· καὶ ἐπιγνώσεσθε διότι ἐγὼ Κύριος Κύριος.

\vs{39}Καὶ ὑμεῖς οἶκος Ἰσραὴλ, τάδε λέγει Κύριος Κύριος, ἕκαστος τὰ ἐπιτηδεύματα αὐτοῦ ἐξάρατε, καὶ μετὰ ταῦτα εἰ ὑμεῖς εἰσακούετέ μου, καὶ τὸ ὄνομά μου τὸ ἅγιον οὐ βεβηλώσετε οὐκέτι ἐν τοῖς δώροις ὑμῶν καὶ ἐν τοῖς ἐπιτηδεύμασιν ὑμῶν·
\vs{40}Διότι ἐπι τοῦ ὄρους τοῦ ἁγίου μου, ἐπʼ ὄρους ὑψηλοῦ, λέγει Κύριος Κύριος, ἐκεῖ δουλεύσουσί μοι πᾶς οἶκος Ἰσραὴλ εἰς τέλος· καὶ ἐκεῖ προσδέξομαι, καὶ ἐκεῖ ἐπισκέψομαι τὰς ἀπαρχὰς ὑμῶν, καὶ τὰς ἀπαρχὰς τῶν ἀφορισμῶν ὑμῶν, ἐν πᾶσι τοῖς ἁγιάσμασιν ὑμῶν.

\vs{41}Ἐν ὀσμῇ εὐωδίας προσδέξομαι ὑμᾶς, ἐν τῷ ἐξαγαγεῖν με ὑμᾶς ἐκ τῶν λαῶν, καὶ εἰσδέχεσθαι ὑμᾶς ἐκ τῶν χωρῶν ἐν αἷς διεσκορπίσθητε ἐν αὐταῖς, καὶ ἁγιασθήσομαι ἐν ὑμῖν κατʼ ὀφθαλμοὺς τῶν λαῶν.
\vs{42}Καὶ ἐπιγνώσεσθε διότι ἐγὼ Κύριος, ἐν τῷ εἰσαγαγεῖν με ὑμᾶς εἰς τὴν γῆν τοῦ Ἰσραὴλ, εἰς τὴν γῆν εἰς ἣν ᾖρα τὴν χεῖρά μου τοῦ δοῦναι αὐτὴν τοῖς πατράσιν ὑμῶν.
\vs{43}Καὶ μνησθήσεσθε ἐκεῖ τὰς ὁδοὺς ὑμῶν, καὶ τὰ ἐπιτηδεύματα ὑμῶν ἐν οἷς ἐμιαίνεσθε ἐν αὐτοῖς, καὶ κόψεσθε τὰ πρόσωπα ὑμῶν ἐν πάσαις ταῖς κακίαις ὑμῶν.
\vs{44}Καὶ ἐπιγνώσεσθε διότι ἐγὼ Κύριος, ἐν τῷ ποιῆσαί με οὕτως ὑμῖν, ὅπως τὸ ὄνομά μου μὴ βεβηλωθῇ κατὰ τὰς ὁδοὺς ὑμῶν τὰς κακὰς, καὶ κατὰ τὰ ἐπιτηδεύματα ὑμῶν τὰ διεφθαρμένα, λέγει Κύριος.

\ch{21}
Καὶ ἐγένετο λόγος Κυρίου πρὸς μὲ, λέγων,
\vs{2}υἱὲ ἀνθρώπου, στήρισον τὸ πρόσωπόν σου ἐπὶ θαιμὰν, καὶ ἐπίβλεψον ἐπὶ Δαρὸμ, καὶ προφήτευσον ἐπὶ δρυμὸν ἡγούμενον Ναγὲβ,
\vs{3}καὶ ἐρεῖς τῷ δρυμῷ Ναγὲβ, ἄκουε λόγον Κυρίου· τάδε λέγει Κύριος Κύριος, ἰδοὺ ἐγὼ ἀνάπτω ἐν σοὶ πῦρ, καὶ καταφάγεται ἐν σοὶ πᾶν ξύλον χλωρὸν καὶ πᾶν ξύλον ξηρὸν, οὐ σβεσθήσεται ἡ φλὸξ ἡ ἐξαφθεῖσα, καὶ κατακαυθήσεται ἐν αὐτῇ πᾶν πρόσωπον ἀπὸ ἀπηλιώτου ἕως Βοῤῥᾶ.
\vs{4}Καὶ ἐπιγνώσεται πᾶσα σὰρξ, ὅτι ἐγὼ Κύριος ἐξέκαυσα αὐτὸ, οὐ σβεσθήσεται.

\vs{5}Καὶ εἶπα, μηδαμῶς Κύριε Κύριε· αὐτοὶ λέγουσι πρὸς μὲ, οὐχὶ παραβολή ἐστι λεγομένη αὕτη;

\vs{6}Καὶ ἐγένετο λόγος Κυρίου πρὸς μὲ, λέγων,

\vs{7}Διατοῦτο προφήτευσον υἱὲ ἀνθρώπου, στήρισον τὸ πρόσωπόν σου ἐπὶ Ἱερουσαλὴμ, καὶ ἐπίβλεψον ἐπὶ τὰ ἅγια αὐτῶν,
\vs{8}καὶ προφητεύσεις ἐπὶ τὴν γῆν τοῦ Ἰσραὴλ, καὶ ἐρεῖς πρὸς τὴν γῆν τοῦ Ἰσραὴλ, τάδε λέγει Κύριος, ἰδοὺ ἐγὼ πρὸς σὲ, καὶ ἐκσπάσω τὸ ἐγχειρίδιόν μου ἐκ τοῦ κολεοῦ αὐτοῦ, καὶ ἐξολοθρεύσω ἐκ σοῦ ἄνομον καὶ ἄδικον·
\vs{9}Ἀνθʼ ὧν ἐξολοθρεύσω ἐκ σοῦ ἄδικον καὶ ἄνομον, οὕτως ἐξελεύσεται τὸ ἐγχειρίδιόν μου ἐκ τοῦ κολεοῦ αὐτοῦ ἐπὶ πᾶσαν σάρκα ἀπὸ ἀπηλιώτου ἔως Βοῤῥᾶ,
\vs{10}καὶ ἐπιγνώσεται πᾶσα σὰρξ, διότι ἐγὼ Κύριος ἐξέσπασα τὸ ἐγχειρίδιόν μου ἐκ τοῦ κολεοῦ αὐτοῦ, οὐκ ἀποστρέψει οὐκέτι.

\vs{11}Καὶ σὺ υἱὲ ἀνθρώπου καταστέναξον ἐν συντριβῇ ὀσφύος σου, καὶ ἐν ὀδύναις στενάξεις κατʼ ὀφθαλμοὺς αὐτῶν.
\vs{12}Καὶ ἔσται ἐὰν εἴπωσι πρὸς σὲ, ἕνεκα τίνος σὺ στενάζεις; καὶ ἐρεῖς, ἐπὶ τῇ ἀγγελίᾳ, διότι ἔρχεται, καὶ θραυσθήσεται πᾶσα καρδία, καὶ πᾶσαι χεῖρες παραλυθήσονται, καὶ ἐκψύξει πᾶσα σὰρξ καὶ πᾶν πνεῦμα, καὶ πάντες μηροὶ μολυνθήσονται ὑγρασίᾳ· ἰδοὺ ἔρχεται, λέγει Κύριος.

\vs{13}Καὶ ἐγένετο λόγος Κυρίου πρὸς μὲ, λέγων,
\vs{14}υἱὲ ἀνθρώπου προφήτευσον, καὶ ἐρεῖς, τάδε λέγει Κύριος, εἰπὸν, ῥομφαία ῥομφαία ὀξύνου καὶ θυμώθητι,
\vs{15}ὅπως σφάξῃς σφάγια, ὀξύνου ὅπως γένῃ εἰς στίλβωσιν, ἑτοίμη εἰς παράλυσιν· σφάζε, ἐξουδένει, ἀπόθου πᾶν ξῦλον.
\vs{16}Καὶ ἔδωκεν αὐτὴν ἑτοίμην τοῦ κρατεῖν χεῖρα αὐτοῦ· ἐξηκονήθη ἡ ῥομφαία, ἐστὶν ἑτοίμη τοῦ δοῦναι αὐτὴν εἰς χεῖρα, ἀποκεντοῦντος.

\vs{17}Ἀνάκραγε καὶ ὀλύλυξον, υἱὲ ἀνθρώπου, ὅτι αὕτη ἐγένετο ἐν τῷ λαῷ μου, αὕτη ἐν πᾶσι τοῖς ἀφηγουμένοις τοῦ Ἰσραήλ· παροικήσουσιν, ἐπὶ ῥομφαίᾳ ἐγένετο ἐν τῷ λαῷ μου· διατοῦτο κρότησον ἐπὶ τὴν χεῖρά σου,
\vs{18}ὅτι δεδικαίωται· καὶ τί εἰ καὶ φυλὴ ἀπωσθῇ; οὐκ ἔσται, λέγει Κύριος Κύριος.

\vs{19}Καὶ σὺ υἱὲ ἀνθρώπου προφήτευσον, καὶ κρότησον χεῖρα ἐπὶ χεῖρα, καὶ διπλασίασον ῥομφαίαν· ἡ τρίτη ῥομφαία τραυματιῶν ἐστι, ῥομφαία τραυματιῶν ἡ μεγάλη· καὶ ἐκστήσεις αὐτοὺς,
\vs{20}ὅπως μὴ θραυσθῇ ἡ καρδία, καὶ πληθυνθῶσιν οἱ ἀσθενοῦντες ἐπὶ πᾶσαν πύλην, παραδέδονται εἰς σφάγια ῥομφαίας· εὖ γέγονεν εἰς σφαγὴν, εὖ γέγονεν εἰς στίλβωσιν.
\vs{21}Καὶ διαπορεύου, ὀξύνου ἐκ δεξιῶν καὶ ἐξ εὐωνύμων, οὗ ἄν τὸ πρόσωπόν σου ἐξεγείρηται.

\vs{22}Καὶ ἐγὼ δὲ κροτήσω χεῖρά μου πρὸς χεῖρά μου, καὶ ἐναφήσω θυμόν μου, ἐγὼ Κύριος λελάληκα.

\vs{23}Καὶ ἐγένετο λόγος Κυρίου πρὸς μὲ, λέγων,
\vs{24}καὶ σὺ, υἱὲ ἀνθρώπου, διάταξον σεαυτῷ δύο ὁδοὺς, τοῦ εἰσελθεῖν ῥομφαίαν βασιλέως Βαβυλῶνος, ἐκ χώρας μιᾶς ἐξελεύσονται αἱ δύο, καὶ χεὶρ ἐν ἀρχῇ ὁδοῦ πόλεως, ἐπʼ ἀρχῆς ὁδοῦ διατάξεις,
\vs{25}τοῦ εἰσελθεῖν ῥομφαίαν ἐπὶ Ῥαββὰθ υἱῶν Ἀμμὼν, καὶ ἐπὶ τὴν Ἰουδαίαν, καὶ ἐπὶ τὴν Ἱερουσαλὴμ, ἐν μέσῳ αὐτῆς.

\vs{26}Διότι στήσεται βασιλεὺς Βαβυλῶνος ἐπὶ τὴν ἀρχαίαν ὁδόν, ἐπʼ ἀρχῆς τῶν δύο ὁδῶν, τοῦ μαντεύσασθαι μαντείαν, τοῦ ἀναβράσαι ῥάβδον, καὶ ἐπερωτῆσαι ἐν τοῖς γλυπτοῖς, καὶ κατασκοπήσασθαι.
\vs{27}Ἐκ δεξιῶν αὐτοῦ ἐγένετο τὸ μαντεῖον ἐπὶ Ἱερουσαλὴμ, τοῦ βαλεῖν χάρακα, τοῦ διανοῖξαι στόμα ἐν βοῇ, ὑψῶσαι φωνὴν μετὰ κραυγῆς, τοῦ βαλεῖν χάρακα ἐπὶ τὰς πύλας αὐτῆς, καὶ βαλεῖν χῶμα, καὶ οἰκοδομῆσαι βελοστάσεις.
\vs{28}Καὶ αὐτός αὐτοῖς ὡς μαντευόμενος μαντείαν ἐνώπιον αὐτῶν, καὶ αὐτὸς ἀναμιμνήσκων ἀδικίας αὐτοῦ μνησθῆναι.

\vs{29}Διατοῦτο τάδε λέγει Κύριος, ἀνθʼ ὧν ἀνεμνήσατε τὰς ἀδικίας ὑμῶν, ἐν τῷ ἀποκαλυφθῆναι τὰς ἀσεβείας ὑμῶν, τοῦ ὀραθῆναι ἁμαρτίας ὑμῶν, ἐν πάσαις ταῖς ἀσεβείαις ὑμῶν καὶ ἐν τοῖς ἐπιτηδεύμασιν ὑμῶν, ἀνθʼ ὧν ἀνεμνήσατε, ἐν τούτοις ἁλώσεσθε.
\vs{30}Καὶ σὺ βέβηλε, ἄνομε, ἀφηγούμενε τοῦ Ἰσραὴλ, οὗ ἥκει ἡ ἡμέρα ἐν καιρῷ ἀδικίας, πέρας,
\vs{31}τάδε λέγει Κύριος, ἀφείλου τὴν κίδαριν, καὶ ἐπέθου τὸν στέφανον, αὐτῇ οὐ τοιαύτη ἔσται· ἐταπείνωσας τὸ ὑψηλὸν, καὶ ὕψωσας τὸ ταπεινόν.
\vs{32}Ἀδικίαν, ἀδικίαν, ἀδικίαν θήσομαι αὐτὴν, οὐαὶ αὐτῇ, τοιαύτη ἔσται ἕως οὗ ἔλθῃ ᾧ καθήκει, καὶ παραδώσω αὐτῷ.

\vs{33}Καὶ σὺ υἱὲ ἀνθρώπου προφήτευσον, καὶ ἐρεῖς, τάδε λέγει Κύριος πρὸς τοὺς υἱοὺς Ἀμμὼν καὶ πρὸς τὸν ὀνειδισμὸν αὐτῶν· καὶ ἐρεῖς, ῥομφαία ῥομφαία ἐσπασμένη εἰς σφάγια, καὶ ἐσπασμένη εἰς συντέλειαν, ἐγείρου ὅπως στίλβῃς·
\vs{34}ἐν τῇ ὁράσει σου τῇ ματαίᾳ, καὶ ἐν τῷ μαντεύεσθαί σε ψευδῆ, τοῦ παραδοῦναί σε ἐπὶ τραχήλους τραυματιῶν ἀνόμων, ἥκει ἡ ἡμέρα ἐν καιρῷ ἀδικίας, πέρας.

\vs{35}Ἀπόστρεφε, μὴ καταλύσῃς ἐν τῷ τόπῳ τούτῳ ᾧ γεγέννησαι, ἐν τῇ γῇ τῇ ἰδίᾳ σου κρινῶ σε.
\vs{36}Καὶ ἐκχεῶ ἐπὶ σὲ ὀργήν μου, ἐν πυρὶ ὀργῆς μου ἐμφυσήσω ἐπὶ σὲ, καὶ παραδώσω σε εἰς χεῖρας ἀνδρῶν βαρβάρων τεκταινόντων διαφθοράς.
\vs{37}Ἐν πυρὶ ἔσῃ κατάβρωμα, τὸ αἷμά σου ἔσται ἐν μέσῳ τῆς γῆς σου· οὐ μὴ γένηταί σου μνεία, διότι ἐγὼ Κύριος λελάληκα.

\ch{22}
Καὶ ἐγένετο λόγος Κυρίου πρὸς μὲ, λέγων,
\vs{2}καὶ σύ υἱὲ ἀνθρώπου, εἰ κρινεῖς τὴν πόλιν τῶν αἱμάτων; καὶ παράδειξον αὐτῇ πάσας τὰς ἀνομίας αὐτῆς,
\vs{3}καὶ ἐρεῖς, τάδε λέγει Κύριος Κύριος, ὦ πόλις ἐκχέουσα αἵματα ἐν μέσῳ αὐτῆς, τοῦ ἐλθεῖν καιρὸν αὐτῆς, καὶ ποιοῦσα ἐνθυμήματα καθʼ ἑαυτῆς, τοῦ μιαίνειν αὐτήν,
\vs{4}ἐν τοῖς αἵμασιν αὐτῶν οἷς ἐξέχεας, παραπέπτωκας, καὶ ἐν τοῖς ἐνθυμήμασί σου οἷς ἐποίεις, ἐμιαίνου· καὶ ἤγγισας τὰς ἡμέρας σου, καὶ ἤγαγες καιρὸν ἐτῶν σου· διατοῦτο δέδωκά σε εἰς ὀνειδισμὸν τοῖς ἔθνεσι, καὶ εἰς ἐμπαιγμὸν πάσαις ταῖς χώραις ταῖς ἐγγιζούσαις πρὸς σὲ,
\vs{5}καὶ ταῖς μακρὰν ἀπεχούσαις ἀπὸ σοῦ, καὶ ἐμπαίξονται ἐν σοὶ ἀκάθαρτος ἡ ἀνομαστὴ, καὶ πολλὴ ἐν ταῖς ἀνομίαις.

\vs{6}Ἰδοὺ οἱ ἀφηγούμενοι οἴκου Ἰσραήλ, ἕκαστος πρὸς τοὺς συγγενεῖς αὐτοῦ, συνεφύροντο ἐν σοὶ, ὅπως ἐκχέωσιν αἷμα.
\vs{7}Πατέρα καὶ μητὲρα ἐκακολόγουν ἐν σοὶ, καὶ πρὸς τὸν προσήλυτον ἀνεστρέφοντο ἐν ἀδικίαις ἐν σοί, ὀρφανὸν καὶ χήραν κατεδυνάστευον.
\vs{8}Καὶ τὰ ἅγιά μου ἐξουδένουν, καὶ τὰ σάββατά μου ἐβεβήλουν ἐν σοί.
\vs{9}Ἄνδρες λῃσταὶ ἐν σοὶ, ὅπως ἐκχέωσιν ἐν σοὶ αἷμα, καὶ ἐπὶ τῶν ὀρέων ἤσθιον ἐπὶ σοὶ, ἀνόσια ἐποίουν ἐν μέσῳ σου.
\vs{10}Αἰσχύνην πατρὸς ἀπεκάλυψαν ἐν σοὶ, καὶ ἐν ἀκαθαρσίαις ἀποκαθημένην ἐταπείνουν ἐν σοί.
\vs{11}Ἕκαστος τὴν γυναῖκα τοῦ πλησίον αὐτοῦ ἠνομοῦσαν, καὶ ἕκαστος τὴν νύμφην αὐτοῦ ἐμίαινεν ἐν ἀσεβείᾳ, καὶ ἕκαστος τὴν ἀδελφὴν αὐτοῦ θυγατέρα τοῦ πατρὸς αὐτοῦ ἐταπείνουν ἐν σοί.

\vs{12}Δῶρα ἐλαμβάνοσαν ἐν σοὶ, ὅπως ἐκχέωσιν αἷμα, τόκον καὶ πλεονασμὸν ἐλαμβάνοσαν ἐν σοὶ, καὶ συνετελέσω συντέλειαν κακίας σου τὴν ἐν καταδυναστείᾳ, ἐμοῦ δὲ ἐπελάθου, λέγει Κύριος.

\vs{13}Ἐὰν δὲ πατάξω χεῖρά μου ἐφʼ οἷς συντετέλεσαι οἷς ἐποίησας, καὶ ἐπὶ τοῖς αἵμασί σου τοῖς γεγενημένοις ἐν μέσῳ σου,
\vs{14}εἰ ὑποστήσεται ἡ καρδία σου; εἰ κρατήσουσιν αἱ χεῖρές σου ἐν ταῖς ἡμέραις αἷς ἐγὼ ποιῶ ἐν σοί; ἐγὼ Κύριος λελάληκα, καὶ ποιήσω.
\vs{15}Καὶ διασκορπιῶ σε ἐν τοῖς ἔθνεσι, καὶ διασπερῶ σε ἐν ταῖς χώραις, καὶ ἐκλείψει ἡ ἀκαθαρσία σου ἐκ σοῦ.
\vs{16}Καὶ κατακληρονομήσω ἐν σοὶ κατʼ ὀφθαλμοὺς τῶν ἐθνῶν, καὶ γνώσεσθε διότι ἐγὼ Κύριος.

\vs{17}Καὶ ἐγένετο λόγος Κυρίου πρὸς μὲ, λέγων,
\vs{18}υἱὲ ἀνθρώπου, ἰδοὺ γεγόνασί μοι ὁ οἶκος Ἰσραὴλ ἀναμεμιγμένοι πάντες χαλκῷ, καὶ σιδήρῳ, καὶ κασσιτέρῳ, καὶ μολίβῳ, ἐν μέσῳ ἀργυρίου ἀναμεμιγμένος ἐστί.

\vs{19}Διατοῦτο εἰπὸν, τάδε λέγει Κύριος Κύριος, ἀνθʼ ὧν ἐγένεσθε εἰς σύγκρασιν μίαν, διατοῦτο ἐγὼ εἰσδέχομαι ὑμᾶς εἰς μέσον Ἱερουσαλήμ·
\vs{20}Καθὼς εἰσδέχεται ἄργυρος, καὶ χαλκὸς, καὶ σίδηρος, καὶ κασσίτερος, καὶ μόλιβος εἰς μέσον καμίνου, τοῦ ἐκφυσῆσαι εἰς αὐτὸ πῦρ τοῦ χωνευθῆναι, οὕτως εἰσδέξομαι ἐν ὀργῇ μου, καὶ συνάξω καὶ χωνεύσω ὑμᾶς,
\vs{21}καὶ ἐκφυσήσω ἐφʼ ὑμᾶς ἐν πυρὶ ὀργῆς μου, καὶ χωνευθήσεσθε ἐν μέσῳ αὐτῆς.
\vs{22}Ὃν τρόπον χωνεύεται ἀργύριον ἐν μέσῳ καμίνου, οὕτως χωνευθήσεσθε ἐν μέσῳ αὐτῆς, καὶ ἐπιγνώσεσθε διότι ἐγὼ Κύριος ἐξέχεα τὸν θυμόν μου ἐφʼ ὑμᾶς.

\vs{23}Καὶ ἐγένετο λόγος Κυρίου πρὸς μὲ, λέγων,
\vs{24}υἱὲ ἀνθρώπου, εἰπὸν αὐτῇ, σὺ εἶ γῆ ἡ οὐ βρεχομένη, οὐδὲ ὑετὸς ἐγένετο ἐπὶ σὲ ἐν ἡμέρᾳ ὀργῆς.
\vs{25}Ἧς οἱ ἀφηγούμενοι ἐν μέσῳ αὐτῆς ὡς λέοντες ὠρυόμενοι ἁρπάζοντες ἁρπάγματα, ψυχὰς κατεσθίοντες ἐν δυναστείᾳ, καὶ τιμὰς λαμβάνοντες· καὶ χῆραί σου ἐπληθύνθησαν ἐν μέσῳ σου.
\vs{26}Καὶ οἱ ἱερεῖς αὐτῆς ἠθέτησαν νόμον μου, καὶ ἐβεβήλουν τὰ ἅγιά μου· ἀναμέσον ἁγίου καὶ βεβήλου οὐ διέστελλον, καὶ ἀναμέσον ἀκαθάρτου καὶ τοῦ καθαροῦ οὐ διέστελλον, καὶ ἀπὸ τῶν σαββάτων μου παρεκάλυπτον τοὺς ὀφθαλμοὺς αὐτῶν, καὶ ἐβεβηλούμην ἐν μέσῳ αὐτῶν.
\vs{27}Οἱ ἄρχοντες αὐτῆς ἐν μέσῳ αὐτῆς ὡς λύκοι ἁρπάζοντες ἁρπάγματα, τοῦ ἐκχέαι αἷμα, ὅπως πλεονεξίᾳ πλεονεκτῶσι.
\vs{28}Καὶ οἱ προφῆται αὐτῆς ἀλείφοντες αὐτοὺς πεσοῦνται, ὁρῶντες μάταια, μαντευόμενοι ψευδῆ, λέγοντες, τάδε λέγει Κύριος· καὶ Κύριος οὐκ ἐλάλησε.
\vs{29}Λαὸν τῆς γῆς ἐκπιεζοῦντες ἀδικίᾳ, καὶ διαρπάζοντες ἁρπάγματα, πτωχὸν καὶ πένητα καταδυναστεύοντες, καὶ πρὸς τὸν προσήλυτον οὐκ ἀναστρεφόμενοι μετὰ κρίματος.

\vs{30}Καὶ ἐζήτουν ἐξ αὐτῶν ἄνδρα ἀναστρεφόμενον ὀρθῶς, καὶ ἑστῶτα πρὸ προσώπου μου ὁλοσχερῶς ἐν καιρῷ τῆς ὀργῆς, τοῦ μὴ εἰς τέλος ἐξαλεῖψαι αὐτὴν, καὶ οὐχ εὗρον.
\vs{31}Καὶ ἐξέχεα ἐπʼ αὐτὴν θυμόν μου ἐν πυρὶ ὀργῆς μου, τοῦ συντελέσαι· τὰς ὁδοὺς αὐτῶν εἰς κεφαλὰς αὐτῶν δέδωκα, λέγει Κύριος Κύριος.

\ch{23}
Καὶ ἐγένετο λόγος Κυρίου πρὸς μὲ, λέγων,
\vs{2}υἱὲ ἀνθρώπου, δύο γυναῖκες ἦσαν θυγατέρες μητρὸς μιᾶς,
\vs{3}καὶ ἐξεπόρνευσαν ἐν Αἰγύπτῳ ἐν τῇ νεότητι αὐτῶν, ἐκεῖ ἔπεσον οἱ μαστοὶ αὐτῶν, ἐκεῖ διεπαρθενεύθησαν.
\vs{4}Καὶ τὰ ὀνόματα αὐτῶν ἦν Ὀολὰ ἡ πρεσβυτέρα, καὶ Ὀολιβὰ ἡ ἀδελφὴ αὐτῆς· καὶ ἐγένοντό μοι, καὶ ἔτεκον υἱοὺς καὶ θυγατέρας, καὶ τὰ ὀνόματα αὐτῶν, Σαμάρεια ἦν Ὀολὰ, καὶ Ἱερουσαλὴμ ἦν Ὀολιβά.

\vs{5}Καὶ ἐξεπόρνευσεν ἡ Ὀολὰ ἀπʼ ἐμοῦ, καὶ ἐπέθετο ἐπὶ τοὺς ἐραστὰς αὐτῆς, ἐπὶ τοὺς Ἀσσυρίους τοὺς ἐγγίζοντας αὐτῇ,
\vs{6}ἐνδεδυκότας ὑακίνθινα, ἡγουμένους καὶ στρατηγούς· νεανίσκοι καὶ ἐπίλεκτοι, πάντες ἱππεῖς ἱππαζόμενοι ἐφʼ ἵππων.
\vs{7}Καὶ ἔδωκε τὴν πορνείαν αὐτῆς ἐπʼ αὐτούς· ἐπίλεκτοι υἱοὶ Ἀσσυρίων πάντες· καὶ ἐπὶ πάντας οὓς ἐπέθετο, ἐν πᾶσι τοῖς ἐνθυμήμασιν αὐτοῖς ἐμιαίνετο.
\vs{8}Καὶ τὴν πορνείαν αὐτῆς ἐξ Αἰγύπτου οὐκ ἐγκατέλιπεν, ὅτι μετʼ αὐτῆς ἐκοιμῶντο ἐν νεότητι αὐτῆς, καὶ αὐτοὶ διεπαρθένευσαν αὐτὴν, καὶ ἐξέχεαν τὴν πορνείαν αὐτῶν ἐπʼ αὐτήν.
\vs{9}Διατοῦτο παρέδωκα αὐτὴν εἰς χεῖρας τῶν ἐραστῶν αὐτῆς, εἰς χεῖρας υἱῶν Ἀσσυρίων, ἐφʼ οὓς ἐπετίθετο.
\vs{10}Αὐτοὶ ἀπεκάλυψαν τὴν αἰσχύνην αὐτῆς, υἱοὺς καὶ θυγατέρας αὐτῆς ἔλαβον, καὶ αὐτὴν ἐν ῥομφαίᾳ ἀπέκτειναν· καὶ ἐγένετο λάλημα εἰς γυναῖκας, καὶ ἐποίησαν ἐκδικήσεις ἐν αὐτῇ εἰς τὰς θυγατέρας.

\vs{11}Καὶ εἶδεν ἡ ἀδελφὴ αὐτῆς Ὀολιβὰ, καὶ διέφθειρε τὴν ἐπίθεσιν αὐτῆς ὑπὲρ ἐπί θεσιν αὐτῆς, καὶ τὴν πορνείαν αὐτῆς ὑπὲρ τὴν πορνείαν τῆς τῆς ἀδελφῆς αὐτῆς.
\vs{12}Ἐπὶ τοὺς υἱοὺς τῶν Ἀσσυρίων ἐπέθετο, ἡγουμένους καὶ στρατηγοὺς, τοὺς ἐγγὺς αὐτῆς, ἐνδεδυκότας εὐπάρυφα, ἱππεῖς ἱππαζομένους ἐφʼ ἵππων· νεανίσκοι ἐπίλεκτοι πάντες·
\vs{13}Καὶ ἴδον ὅτι μεμίανται ὁδὸς μία τῶν δύο.

\vs{14}Καὶ προσέθετο πρὸς τὴν πορνείαν αὐτῆς, καὶ εἶδεν ἄνδρας ἐζωγραφημένους ἐπὶ τοῦ τοίχου, εἰκόνας Χαλδαίων ἐζωγραφημένους ἐν γραφίδι,
\vs{15}ἐζωσμένους ποικίλματα ἐπὶ τὰς ὀσφύας αὐτῶν, παραβαπτὰ καὶ ἐπὶ τῶν κεφαλῶν αὐτῶν, ὄψις τρισσὴ πάντων, ὁμοίωμα υἱῶν Χαλδαίων, γῆς πατρίδος αὐτοῦ·
\vs{16}Καὶ ἐπέθετο ἐπʼ αὐτοὺς τῇ ὁράσει ὀφθαλμῶν αὐτῆς· καὶ ἐξαπέστειλεν ἀγγέλους πρὸς αὐτοὺς εἰς γῆν Χαλδαίων.
\vs{17}Καὶ ἤλθοσαν πρὸς αὐτὴν υἱοὶ Βαβυλῶνος, εἰς κοίτην καταλυόντων, καὶ ἐμίαινον αὐτὴν ἐν τῇ πορνείᾳ αὐτῆς, καὶ ἐμιάνθη ἐν αὐτοῖς, καὶ ἀπέστη ἡ ψυχὴ αὐτῆς ἀπʼ αὐτῶν.
\vs{18}Καὶ ἀπεκάλυψε τὴν πορνείαν αὐτῆς, καὶ ἀπεκάλυψεν αἰσχύνην αὐτῆς· καὶ ἀπέστη ἡ ψυχή μου ἀπʼ αὐτῆς, ὃν τρόπον ἀπέστη ἡ ψυχή μου ἀπὸ τῆς ἀδελφῆς αὐτῆς.

\vs{19}Καὶ ἐπλήθυνας τὴν πορνείαν σου, τοῦ ἀναμνῆσαι ἡμέραν νεότητός σου, ἐν αἷς ἐπόρνευσας ἐν Αἰγύπτῳ,
\vs{20}καὶ ἐπέθου ἐπὶ τοὺς Χαλδαίους, ὧν ὡς ὄνων αἱ σάρκες αὐτῶν, καὶ αἰδοῖα ἵππων τὰ αἰδοῖα αὐτῶν,
\vs{21}καὶ ἐπεσκέψω τὴν ἀνομίαν νεότητός σου, ἃ ἐποίεις ἐν Αἰγύπτῳ ἐν τῷ καταλύματί σου, οὗ οἱ μαστοὶ νεότητός σου.

\vs{22}Διατοῦτο, Ὀολιβὰ, τάδε λέγει Κύριος, ἰδοὺ ἐγὼ ἐξεγείρω τοὺς ἐραστάς σου ἐπὶ σὲ, ἀφʼ ὧν ἀπέστη ἡ ψυχή σου ἀπʼ αὐτῶν, καὶ ἐπάξω αὐτοὺς ἐπὶ σὲ κυκλόθεν,
\vs{23}υἱοὺς Βαβυλῶνος, καὶ πάντας τοὺς Χαλδαίους, Φακοὺκ, καὶ Σουὲ, καὶ Ὑχουὲ, καὶ πάντας υἱοὺς Ἀσσυρίων μετʼ αὐτῶν, νεανίσκους ἐπιλέκτους, ἡγεμόνας καὶ στρατηγοὺς, πάντας τρισσοὺς καὶ ὀνομαστὺος, ἱππεύοντας ἐφʼ ἵππων.
\vs{24}Καὶ πάντες ἥξουσιν ἐπὶ σὲ ἀπὸ Βοῤῥᾶ, ἅρματα καὶ τροχοὶ μετʼ ὄχλου λαῶν, θυρεοὶ καὶ πέλται, καὶ βαλεῖ φυλακὴν ἐπὶ σὲ κύκλῳ·
\vs{25}καὶ δώσω πρὸ προσώπου αὐτῶν κρίμα, καὶ ἐκδικήσουσί σε ἐν τοῖς κρίμασιν αὐτῶν. Καὶ δώσω τὸν ζῆλόν μου ἐν σοὶ, καὶ ποιήσουσι μετὰ σου ἐν ὀργῇ θυμοῦ· μυκτῆρά σου καὶ ὦτά σου ἀφελοῦσι, καὶ τοὺς καταλοίπους σου ἐν ῥομφαίᾳ καταβαλοῦσιν· αὐτοὶ υἱούς σου καὶ θυγατέρας σου λήψονται, καὶ τοὺς καταλοίπους σου πῦρ καταφάγεται.
\vs{26}Καὶ ἐκδύσουσί σε τὸν ἱματισμόν σου, καὶ λήψονται τὰ σκεύη τῆς καυχήσεώς σου.
\vs{27}Καὶ ἀποστρέψω τὰς ἀσεβείας σου ἐκ σοῦ, καὶ τὴν πορνείαν σου ἐκ γῆς Αἰγύπτου· καὶ οὐ μὴ ἄρῃς τοὺς ὀφθαλμούς σου ἐπʼ αὐτοὺς, καὶ Αἰγύπτου οὐ μὴ μνησθῇς οὐκέτι.

\vs{28}Διότι τάδε λέγει Κύριος Κύριος, ἰδοὺ ἐγὼ παραδίδωμί σε εἰς χεῖρας ὧν μισεῖς, ἀφʼ ὧν ἀπέστη ἡ ψυχή σου ἀπʼ αὐτῶν.
\vs{29}Καὶ ποιήσουσιν ἐν σοὶ ἑν μίσει, καὶ λήψονται πάντας τοὺς πόνους σου καὶ τοὺς μόχθους σου, καὶ ἔσῃ γυμνὴ καὶ αἰσχύνουσα, καὶ ἀποκαλυφθήσεται αἰσχύνη πορνείας σου· καὶ ἀσέβειά σου, καὶ ἡ πορνεία σου
\vs{30}ἐποίησε ταῦτά σοι, ἐν τῷ ἐκπορνεῦσαί σε ὀπίσω ἐθνῶν, καὶ ἐμιαίνου ἐν τοῖς ἐνθυμήμασιν αὐτῶν.

\vs{31}Ἐν τῇ ὁδῷ τῆς ἀδελφῆς σου ἐπορεύθης· καὶ δώσω τὸ ποτήριον αὐτὴς εἰς χεῖράς σου.
\vs{32}Τάδε λέγει Κύριος, τὸ ποτήριον τῆς ἀδελφῆς σου πίεσαι, τὸ βαθὺ καὶ τὸ πλατὺ, καὶ τὸ πλεονάζον τοῦ συντελέσαι μέθην.
\vs{33}Καὶ ἐκλύσεως πλησθήσῃ, καὶ τὸ ποτήριον ἀφανισμοῦ, ποτήριον ἀδελφῆς σου Σαμαρείας, πίεσαι αὐτό·
\vs{34}καὶ τὰς ἑορτὰς καὶ τὰς νουμηνίας αὐτῆς ἀποστρέψω· διότι ἐγὼ λελάληκα, λέγει Κύριος.
\vs{35}Διατοῦτο τάδε λέγει Κύριος, ἀνθʼ ὧν ἐπελάθον μου, καὶ ἀπέῤῥιψάς με ὀπίσω τοῦ σώματός σου, καὶ σὺ λάβε τὴν ἀσέβειάν σου καὶ τὴν πορνείαν σου.

\vs{36}Καὶ εἶπε Κύριος πρὸς μὲ, υἱὲ ἀνθρώπου, οὐ κρινεῖς τὴν Ὀολὰν καὶ τὴν Ὀολιβὰν, καὶ ἀναγγελεῖς αὐταῖς τὰς ὀνομίας αὐτῶν;
\vs{37}Ὅτι ἐμοιχῶντο, καὶ αἷμα ἐν χερσὶν αὐτῶν, τὰ ἐνθυμήματα αὐτῶν ἐμοιχῶντο, καὶ τὰ τέκνα αὐτῶν ἃ ἐγέννησάν μοι, διήγαγον αὐτοῖς διʼ ἐμπύρων.
\vs{38}Ἕως καὶ ταῦτα ἐποίησάν μοι, τὰ ἅγιά μου ἐμίαινον, καὶ τὰ σάββατά μου ἐβεβήλουν,
\vs{39}καὶ ἐν τῷ σφάζειν αὐτοὺς τὰ τέκνα αὐτῶν τοῖς εἰδώλοις αὐτῶν, καὶ εἰσεπορεύοντο εἰς τὰ ἅγιά μου, τοῦ βεβηλοῦν αὐτά· καὶ ὅτι οὕτως ἐποίουν ἐν μέσῳ τοῦ οἴκου μου,
\vs{40}καὶ ὅτι τοῖς ἀνδράσι τοῖς ἐρχομένοις μακρόθεν, οἷς ἀγγέλους ἐξαπέστειλαν πρὸς αὐτοὺς, καὶ ἅμα τῷ ἔρχεσθαι αὐτοὺς, εὐθὺς ἐλούου, καὶ ἐστιβίζου τοὺς ὀφθαλμούς σου, καὶ ἐκόσμου κόσμῳ,
\vs{41}καὶ ἐκάθου ἐπὶ κλίνης ἐστρωμένης, καὶ τράπεζα κεκοσμημένη πρὸ προσώπου αὐτῆς, καὶ τὸ θυμίαμα καὶ τὸ ἔλαιόν μου εὐφραίνοντο ἐν αὐτοῖς,
\vs{42}καὶ φωνὴν ἁρμονίας ἀνεκρούοντο, καὶ πρὸς ἄνδρας ἐκ πλήθους ἀνθρώπων ἥκοντας ἐκ τῆς ἐρήμου· καὶ ἐδίδοσαν ψέλλια ἐπὶ τὰς χεῖρας αὐτῶν, καὶ στέφανον καυχήσεως ἐπὶ τὰς κεφαλὰς αὐτῶν.

\vs{43}Καὶ εἶπα, οὐκ ἐν τούτοις μοιχεύουσι; καὶ ἔργα πόρνης καὶ αὐτὴ ἐξεπόρνευσε;
\vs{44}Καὶ εἰσεπορεύοντο πρὸς αὐτὴν, ὃν τρόπον εἰσπορεύονται πρὸς γυναῖκα πόρνην, οὕτως εἰσεπορεύοντο πρὸς Ὀολὰν, καὶ πρὸς Ὀολιβὰν, τοῦ ποιῆσαι ἀνομίαν.
\vs{45}Καὶ ἄνδρες δίκαιοι αὐτοὶ καὶ ἐκδικήσουσιν αὐτὰς ἐκδικήσει μοιχαλίδος, καὶ ἐκδικήσει αἵματος, ὅτι μοιχαλίδες εἰσὶ, καὶ αἷμα ἐν χερσὶν αὐτῶν.

\vs{46}Τάδε λέγει Κύριος Κύριος, ἀνάγαγε ἐπʼ αὐτὰς ὄχλον, καὶ δὸς ἐν αὐταῖς ταραχὴν καὶ διαρπαγὴν,
\vs{47}καὶ λιθοβόλησον ἐπʼ αὐτὰς λίθοις ὄχλων, καὶ κατακέντει αὐτὰς ἐν τοῖς ξίφεσιν αὐτῶν· υἱοὺς αὐτῶν καὶ θυγατέρας αὐτῶν ἀποκτενοὺσι, καὶ τοὺς οἴκους αὐτῶν ἐμπρήσουσι.
\vs{48}Καὶ ἀποστρέψω ἀσεβείαν ἐκ τῆς γῆς, καὶ παιδευθήσονται πᾶσαι αἱ γυναῖκες, καὶ οὐ μὴ ποιήσουσι κατὰ τὰς ἀσεβείας αὐτῶν.
\vs{49}Καὶ δοθήσεται ἡ ἀσέβεια ὑμῶν ἐφʼ ὑμᾶς, καὶ τὰς ἁμαρτίας τῶν ἐνθυμημάτων ὑμῶν λήψεσθε, καὶ γνώσεσθε διότι ἐγὼ Κύριος.

\ch{24}
Καὶ ἐγένετο λόγος Κυρίου πρὸς μὲ, ἐν τῷ ἔτει τῷ ἐννάτῳ, ἐν τῷ μηνὶ τῷ δεκάτῳ, δεκάτῃ τοῦ μηνὸς, λέγων,
\vs{2}υἱὲ ἀνθρώπου γράψον σεαυτῷ εἰς ἡμέραν ἀπὸ τῆς ἡμέρας ταύτης, ἀφʼ ἧς ἀπηρείσατο βασιλεὺς Βαβυλῶνος ἐπὶ Ἰερουσαλήμ, ἀπὸ τῆς ἡμέρας τῆς σήμερον·
\vs{3}Καὶ εἰπὸν ἐπὶ τὸν οἶκον τὸν παραπικραίνοντα παραβολήν, καὶ ἐρεῖς πρὸς αὐτούς

Τάδε λέγει Κύριος, ἐπίστησον τὸν λέβητα, καὶ ἔγχεον εἰς αὐτὸν ὕδωρ,
\vs{4}καὶ ἔμβαλε εἰς αὐτὸν τὰ διχοτομήματα, πᾶν διχοτόμημα καλὸν, σκέλος καὶ ὦμον ἐκσεσαρκισμένα ἀπὸ τῶν ὀστῶν,
\vs{5}ἐξ ἐπιλέκτων κτηνῶν εἰλημμένων· καὶ ὑπόκαιε τὰ ὀστᾶ ὑποκάτω αὐτῶν, ἔζεσε καὶ ἥψηται τὰ ὀστᾶ αὐτῆς ἐν μέσῳ αὐτῆς.

\vs{6}Διατοῦτο τάδε λέγει Κύριος, ὦ πόλις αἱμάτων, λέβης, ἐν ᾧ ἐστιν ἰὸς ἐν αὐτῷ, καὶ ὁ ἰὸς οὐκ ἐξῆλθεν ἐξ αὐτῆς, κατὰ μέλος αὐτῆς ἐξήνεγκεν, οὐκ ἐπεσεν ἐπʼ αὐτὴν κλῆρος.
\vs{7}Ὅτι αἷμα αὐτῆς ἐν μέσῳ αὐτῆς ἐστιν, ἐπὶ λεωπετρίαν τέταχα αὐτό· οὐκ ἐκκέχυκα αὐτὸ ἐπὶ τὴν γῆν, τοῦ καλύψαι ἐπʼ αὐτὸ γῆν,
\vs{8}τοῦ ἀναβῆναι θυμὸν εἰς ἐκδίκησιν ἐκδικηθῆναι· δέδωκα τὸ αἷμα αὐτῆς ἐπὶ λεωπετρίαν, τοῦ μὴ καλύψαι αὐτό.

\vs{9}Διατοῦτο τάδε λέγει Κύριος, κᾀγὼ μεγαλυνῶ τὸν δαλόν·
\vs{10}Καὶ πληθυνῶ τὰ ξύλα, καὶ ἀνακαύσω τὸ πῦρ, ὅπως τακῇ τὰ κρέα, καὶ ἐλαττωθῇ ὁ ζωμὸς,
\vs{11}καὶ στῇ ἐπὶ τοὺς ἄνθρακας, ὅπως προσκαυθῇ καὶ θερμανθῇ ὁ χαλκὸς αὐτῆς, καὶ τακῇ ἐν μέσῳ ἀκαθαρσίας αὐτῆς, καὶ ἐκλίπῃ ὁ ἰὸς αὐτῆς,
\vs{12}καὶ οὐ μὴ ἐξέλθῃ ἐξ αὐτῆς πολὺς ὁ ἰὸς αὐτῆς. Καταισχυνθήσεται ὁ ἰὸς αὐτῆς,
\vs{13}ἀνθʼ ὧν ἐμιαίνου σύ· καὶ τί ἐὰν μὴ καθαρισθῇς ἔτι ἕως οὗ ἐμπλήσω τὸν θυμόν μου;

\vs{14}Ἐγὼ Κύριος λελάληκα, καὶ ἥξει, καὶ ποιήσω, οὐ διαστελῶ, οὐδὲ μὴ ἐλεήσω· κατὰ τὰς ὁδούς σου, καὶ κατὰ τὰ ἐνθυμήματά σου κρινῶ σε, λέγει Κύριος· διατοῦτο ἐγὼ κρινῶ σε κατὰ τὰ αἵματά σου, καὶ κατὰ τὰ ἐνθυμήματά σου κρινῶ σε, ἡ ἀκάθαρτος, ἡ ὀνομαστὴ, καὶ πολλὴ τοῦ παραπικραίνειν.

\vs{15}Καὶ ἐγένετο λόγος Κυρίου πρὸς μὲ, λέγων,
\vs{16}υἱὲ ἀνθρώπου, ἰδοὺ ἐγὼ λαμβάνω ἐκ σοῦ τὰ ἐπιθυμήματα τῶν ὀφθαλμῶν σου ἐν παρατάξει, οὐ μὴ κοπῇς, οὐδʼ οὐ μὴ κλαυσθῇς,
\vs{17}στεναγμὸς αἵματος, ὀσφύος πένθος ἔσῃ· οὐκ ἔσται τὸ τρίχωμά σου συμπεπλεγμένον ἐπὶ σὲ, καὶ τὰ ὑοιδήματά σου ἐν τοῖς ποσί σου· οὐ μὴ παρακληθῇς ἐν χείλεσιν αὐτῶν, καὶ ἄρτον ἀνδρῶν οὐ μὴ φάγῃς.

\vs{18}Καὶ ἐλάλησα πρὸς τὸν λαὸν τοπρωῒ, ὃν τρόπον ἐνετείλατὸ μοι ἑσπέρας, καὶ ἐποίησα τοπρωῒ, ὃν τρόπον ἐπετάγη μοι.
\vs{19}Καὶ εἶπε πρὸς μὲ ὁ λαὸς, οὐκ ἀναγγέλλεις ἡμῖν τί ἐστι ταῦτα ἃ σὺ ποιεῖς;
\vs{20}Καὶ εἶπα πρὸς αὐτοὺς, λόγος Κυρίου ἐγένετο πρὸς μὲ, λέγων,
\vs{21}εἰπὸν πρὸς τὸν οἶκον τοῦ Ἰσραήλ,

Τάδε λέγει Κύριος, ἰδοὺ ἐγὼ βεβηλῶ τὰ ἅγιά μου, φρύαγμα ἰσχύος ὑμῶν, ἐπιθυμήματα ὀφθαλμῶν ὑμῶν, καὶ ὑπὲρ ὧν φείδονται αἱ ψυχαὶ ὑμῶν· καὶ οἱ υἱοὶ ὑμῶν καὶ αἱ θυγατέρες ὑμῶν, οὓς ἐγκατελίπετε, ἐν ῥομφαίᾳ πεσοῦνται.
\vs{22}Καὶ ποιήσετε ὃν τρόπον πεποίηκα· ἀπὸ στόματος αὐτῶν οὐ παρακληθήσεσθε, καὶ ἄρτον ἀνδρῶν οὐ φάγεσθε,
\vs{23}καὶ αἱ κόμαι ὑμῶν ἐπὶ τῆς κεφαλῆς ὑμῶν, καὶ τὰ ὑποδήματα ὑμῶν ἐν τοῖς ποσὶν ὑμῶν· οὔτε μὲ κόψησθε, οὔτε μὴ κλαύσητε, καὶ ἐντακήσεσθε ἐν ταῖς ἀδικίαις ὑμῶν, καὶ παρακαλέσετε ἕκαστος τὸν ἀδελφὸν αὐτοῦ.
\vs{24}Καὶ ἔσται Ἰεζεκιὴλ ὑμῖν εἰς τέρας, κατὰ πάντα ὅσα ἐποίησα ποιήσετε, ὅταν ἔλθῃ ταῦτα, καὶ ἐπιγνώσεσθε διότι ἐγὼ Κύριος.

\vs{25}Καὶ σύ, υἱὲ ἀνθρώπου, οὐχὶ ἐν τῇ ἡμέρᾳ ὅταν λαμβάνω τὴν ἰσχὺν παρʼ αὐτῶν, τὴν ἔπαρσιν τῆς καυχήσεως αὐτῶν, τὰ ἐπιθυμήματα ὀφθαλμῶν αὐτῶν, καὶ τὴν ἔπαρσιν ψυχῆς αὐτῶν, υἱοὺς αὐτῶν καὶ θυγατέρας αὐτῶν,
\vs{26}ἐν τῇ ἡμέρᾳ ἐκείνῃ ἥξει ὁ ἀνασωζόμενος πρὸς σὲ, τοῦ ἀναγγεῖλαί σοι εἰς τὰ ὦτα;
\vs{27}Ἐν τῇ ἡμέρᾳ ἐκείνῃ διανοιχθήσεται τὸ στόμα σου πρὸς τὸν ἀνασωζόμενον· λαλήσεις, καὶ οὐ μὴ ἀποκωφωθῇς οὐκέτι, καὶ ἔσῃ αὐτοῖς εἰς τέρας, καὶ ἐπιγνώσονται διότι ἐγὼ Κύριος.

\ch{25}
Καὶ ἐγένετο λόγος Κυρίου πρὸς μὲ, λέγων,
\vs{2}υἱὲ ἀνθρώπου, στήρισον τὸ πρόσωπόν σου ἐπὶ τοὺς υἱοὺς Ἀμμὼν, καὶ προφήτευσον ἐπʼ αὐτοὺς,
\vs{3}καὶ ἐρεῖς τοῖς υἱοῖς Ἀμμὼν,

Ἀκούσατε λόγον Κυρίου· τάδε λέγει Κύριος, ἀνθʼ ὧν ἐπεχάρητε ἐπὶ τὰ ἅγιά μου, ὅτι ἐβεβηλώθη, καὶ ἐπὶ τὴν γῆν τοῦ Ἰσραὴλ, ὅτι ἠφανίσθη, καὶ ἐπὶ τὸν οἶκον τοῦ Ἰούδα, ὅτι ἐπορεύθησαν ἐν αἰχμαλωσίᾳ,
\vs{4}διατοῦτο ἰδοὺ ἐγὼ παραδίδωμι ὑμᾶς τοῖς υἱοῖς Κεδὲμ εἰς κληρονομίαν, καὶ κατασκηνώσουσιν ἐν τῇ ἀπαρτίᾳ αὐτῶν ἐν σοὶ, καὶ δώσουσιν ἐν σοὶ τὰ σκηνώματα αὐτῶν· αὐτοὶ φάγονται τοὺς καρπούς σου, καὶ αὐτοὶ πίονται τὴν πιότητά σου.
\vs{5}Καὶ δώσω τὴν πόλιν τοῦ Ἀμμὼν εἰς νομὰς καμήλων, καὶ τοὺς υἱοὺς Ἀμμὼν εἰς νομὴν προβάτων, καὶ ἐπιγνώσεσθε διότι ἐγὼ Κύριος.

\vs{6}Διότι τάδε λέγει Κύριος, ἀνθʼ ὧν ἐκρότησας τὴν χεῖρά σου, καὶ ἐψόφησας τῷ ποδί σου, καὶ ἐπέχαρας ἐκ ψυχῆς σου ἐπὶ τὴν γῆν τοῦ Ἰσραὴλ,
\vs{7}διατοῦτο ἐκτενῶ τὴν χεῖρά μου ἐπὶ σὲ, καὶ δώσω σε εἰς διαρπαγὴν ἐν τοῖς ἔθνεσι, καὶ ἐξολεθρεύσω σε ἐκ τῶν λαῶν, καὶ ἀπολῶ σε ἐκ τῶν χωρῶν ἀπωλείᾳ, καὶ ἐπιγνώσῃ διότι ἐγὼ Κύριος.

\vs{8}Τάδε λέγει Κύριος, ἀνθʼ ὧν εἶπε Μωάβ, ἰδοὺ, οὐχ ὃν τρόπον πάντα τὰ ἔθνη, οἶκος Ἰσραὴλ καὶ Ἰούδα;
\vs{9}Διατοῦτο ἰδοὺ ἐγὼ παραλύω τὸν ὦμον Μωάβ ἀπὸ πόλεων ἀκρωτηρίων αὐτοῦ, ἐκλεκτὴν γῆν, οἶκον Βεθασιμούθ ἐπάνω πηγῆς πόλεως παραθαλασσίας,
\vs{10}τοὺς υἱοὺς Κεδὲμ ἑπὶ τοὺς υἱοὺς Ἀμμὼν δέδωκα αὐτῷ εἰς κληρονομίαν, ὅπως μὴ μνεία γένηται τῶν υἱῶν Ἀμμών.
\vs{11}Καὶ εἰς Μωὰβ ποιήσω ἐκδίκησιν, καὶ ἐπιγνώσονται διότι ἐγὼ Κύριος.

\vs{12}Τάδε λέγει Κύριος, ἀνθʼ ὧν ἐποίησεν ἡ Ἰδουμαία ἐν τῷ ἐκδικῆσαι αὐτοὺς ἐκδίκησιν εἰς τὸν οἶκον Ἰούδα, καὶ ἐμνησικάκησαν, καὶ ἐξεδίκησαν δίκην,
\vs{13}διατοῦτο τάδε λέγει Κύριος, καὶ ἐκτενῶ τὴν χεῖρά μου ἐπὶ τὴν Ἰδουμαίαν, καὶ ἐξολοθρεύσω ἐξ αὐτῆς ἄνθρωπον καὶ κτῆνος, καὶ θήσομαι αὐτὴν ἔρημον, καὶ ἐκ Θαιμὰν διωκόμενοι ἐν ῥομφαίᾳ πεσοῦνται.
\vs{14}Καὶ δώσω ἐκδίκησίν μου ἐπὶ τὴν Ἰδουμαίαν ἐν χειρὶ λαοῦ μου Ἰσραὴλ, καὶ ποιήσουσιν ἐν τῇ Ἰδουμαίᾳ κατὰ τὴν ὀργήν μου καὶ κατὰ τὸν θυμόν μου, καὶ ἐπιγνώσονται τὴν ἐκδίκησίν μου, λέγει Κύριος.

\vs{15}Διατοῦτο τάδε λέγει Κύριος, ἀνθʼ ὧν ἐποίησαν οἱ ἀλλόφυλοι ἐν ἐκδικήσει, καὶ ἐξανέστησαν ἐκδίκησιν ἐπιχαίροντες ἐκ ψυχῆς, τοῦ ἐξαλεῖψαι ἕως ἑνὸς,
\vs{16}διατοῦτο τάδε λέγει Κύριος, ἰδοὺ ἐγὼ ἐκτείνω τὴν χεῖρά μου ἐπὶ τοὺς ἀλλοφύλους, καὶ ἐξολοθρεύσω Κρῆτας, καὶ ἀπολῶ τοὺς καταλοίπους τοὺς κατοικοῦντας τὴν παραλίαν·
\vs{17}Καὶ ποιήσω ἐν αὐτοῖς ἐκδικήσεις μεγάλας, καὶ ἐπιγνώσονται διότι ἐγὼ Κύριος, ἐν τῷ δοῦναι τὴν ἐκδίκησίν μου ἐπʼ αὐτούς.

\ch{26}
Καὶ ἐγενήθη ἐν τῷ ἑνδεκάτῳ ἔτει, μιᾷ τοῦ μηνὸς, ἐγένετο λόγος Κυρίου πρὸς μὲ, λέγων,

\vs{2}Υἱὲ ἀνθρώπου, ἀνθʼ οὗ εἶπε Σὸρ ἐπὶ Ἱερουσαλὴμ, εὖγε συνετρίβη, ἀπόλωλε τὰ ἔθνη, ἐπεστράφη πρὸς μὲ, ἡ πλήρης ἠρήμωται·
\vs{3}Διατοῦτο τάδε λέγει Κύριος, ἰδοὺ ἐγὼ ἐπὶ σὲ Σὸρ, καὶ ἀνάξω ἐπὶ σὲ ἔθνη πολλὰ, ὡς ἀναβαίνει ἡ θάλασσα τοῖς κύμασιν αὐτῆς.
\vs{4}Καὶ καταβαλοῦσι τὰ τείχη Σὸρ, καὶ καταβαλοῦσι τοὺς πύργους σου, καὶ λικμήσω τὸν χοῦν αὐτῆς ἀπʼ αὐτῆς, καὶ δώσω αὐτὴν εἰς λεωπετρίαν·
\vs{5}Ψυγμὸς σαγηνῶν ἔσται ἐν μέσῳ θαλάσσης, ὅτι ἐγὼ λελάληκα, λέγει Κύριος· καὶ ἔσται εἰς προνομὴν τοῖς ἔθνεσι.
\vs{6}Καὶ αἱ θυγατέρες αὐτῆς ἐν πεδίῳ μαχαίρᾳ ἀναιρεθήσονται, καὶ γνώσονται ὅτι ἐγὼ Κύριος.

\vs{7}Ὅτι τάδε λέγει Κύριος, ἰδοὺ ἐγὼ ἐπάγω ἐπὶ σὲ Σόρ τὸν Ναβουχοδονόσορ βασιλέα Βαβυλῶνος ἀπὸ τοῦ Βοῤῥᾶ, βασιλεὺς βασιλέων ἐστὶ, μεθʼ ἵππων καὶ ἁρμάτων καὶ ἱππέων καὶ συναγωγῆς ἐθνῶν πολλῶν σφόδρα.
\vs{8}Οὗτος τὰς θυγατέρας σου τὰς ἐν τῷ πεδίῳ μαχαίρᾳ ἀνελεῖ, καὶ δώσει ἐπὶ σὲ προφυλακὴν, καὶ περιοικοδομήσει, καὶ ποιήσει ἐπὶ σὲ κύκλῳ χάρακα, καὶ περίστασιν ὅπλων, καὶ τὰς λόγχας αὐτοῦ ἀπέναντί σου δώσει·
\vs{9}Τὰ τείχη σου, καὶ τοὺς πύργους σου καταβαλεῖ ἐν ταῖς μαχαίραις αὐτοῦ.
\vs{10}Ἀπὸ τοῦ πλήθους τῶν ἵππων αὐτοῦ κατακαλύψει σε ὁ κονιορτὸς αὐτῶν· καὶ ἀπὸ τῆς φωνῆς τῶν ἱππέων αὐτοῦ καὶ τῶν τροχῶν τῶν ἁρμάτων αὐτοῦ σεισθήσεται τὰ τείχη σου, εἰσπορευομένου αὐτοῦ τὰς πύλας σου, ὡς εἰσπορευόμενος εἰς πόλιν ἐκ πεδίου.
\vs{11}Ἐν ταῖς ὁπλαῖς τῶν ἵππων αὐτοῦ καταπατήσουσί σου πάσας τὰς πλατείας· τὸν λαόν σου μαχαίρᾳ ἀνελεῖ, καὶ τὴν ὑπόστασιν τῆς ἰσχύος σου ἐπὶ τὴν γῆν κατάξει.

\vs{12}Καὶ προνομεύσει τὴν δύναμίν σου, καὶ σκυλεύσει τὰ ὑπάρχοντά σου, καὶ καταβαλεῖ τὰ τείχη σου, καὶ τοὺς οἴκους σου τοὺς ἐπιθυμητοὺς καθελεῖ· καὶ τοὺς λίθους σου, καὶ τὰ ξύλα σου, καὶ τὸν χοῦν σου εἰς μέσον τῆς θαλάσσης σου ἐμβαλεῖ.
\vs{13}Καὶ καταλύσει τὸ πλῆθος τῶν μουσικῶν σου, καὶ ἡ φωνὴ τῶν ψαλτηρίων σου οὐ μὴ ἀκουσθῇ ἔτι.
\vs{14}Καὶ δώσω σε λεωπετρίαν, ψυγμὸς σαγηνῶν ἔσῃ, οὐ μὴ οἰκοδομηθῇς ἔτι, ὅτι ἐγὼ Κύριος ἐλάλησα, λέγει Κύριος.

\vs{15}Διότι τάδε λέγει Κύριος Κύριος τῇ Σὸρ, οὐκ ἀπὸ φωνῆς τῆς πτώσεώς σου, ἐν τῷ στενάξαι τραυματίας, ἐν τῷ σπᾶσαι μάχαιραν ἐν μέσῳ σου, σεισθήσονται αἱ νῆσαι;
\vs{16}Καὶ καταβήσονται ἀπὸ τῶν θρόνων αὐτῶν πάντες οἱ ἄρχοντες ἐκ τῶν ἐθνῶν τῆς θαλάσσης, καὶ ἀφελοῦνται τὰς μίτρας ἀπὸ τῶν κεφαλῶν αὐτῶν, καὶ τὸν ἱματισμὸν τὸν ποικίλον αὐτῶν ἐκδύσονται· ἐκστάσει ἐκστήσονται, ἐπὶ γῆν καθεδοῦνται, καὶ φοβηθήσονται τὴν ἀπώλειαν αὐτῶν, καὶ στενάξουσιν ἐπὶ σὲ,
\vs{17}καὶ λήψονται ἐπὶ σὲ θρῆνον, καὶ ἐροῦσί σοι, πῶς κατελύθης ἐκ θαλάσσης ἡ πόλις ἡ ἐπαινετὴ, ἡ δοῦσα τὸν φόβον αὐτῆς πᾶσι τοῖς κατοικοῦσιν αὐτήν;
\vs{18}Καὶ φοβηθήσονται αἱ νῆσοι ἀπὸ ἡμέρας πτώσεώς σου.

\vs{19}Ὅτι τάδε λέγει Κύριος Κύριος, ὅταν δῶ πόλιν ἠρημωμένην ὡς τὰς πόλεις τὰς μὴ κατοικισθησομένας, ἐν τῷ ἀναγαγεῖν με ἐπὶ σὲ τὴν ἄβυσσον, καὶ κατακαλύψει σε ὕδωρ πολὺ,
\vs{20}καὶ καταβιβάσω σε πρὸς τοὺς καταβαίνοντας εἰς βόθρον πρὸς λαὸν αἰῶνος, καὶ κατοικιῶ σε εἰς βάθη τῆς γῆς ὡς ἔρημον αἰώνιον μετὰ καταβαινόντων εἰς βόθρον, ὅπως μὴ κατοικηθῇς, μηδὲ ἀναστῇς ἐπὶ γῆς ζωῆς,
\vs{21}ἀπώλειάν σε δώσω, καὶ οὐχ ὑπάρξεις ἔτι εἰς τὸν αἰῶνα, λέγει Κύριος Κύριος.

\ch{27}
Καὶ ἐγένετο λόγος Κυρίου πρὸς μὲ λέγων,

\vs{2}Καὶ σὺ υἱὲ ἀνθρώπου λάβε ἐπὶ Σὸρ θρῆνον,
\vs{3}καὶ ἐρεῖς τῇ Σὸρ τῇ κατοικούσῃ ἐπὶ τῆς εἰσόδου τῆς θαλάσσης, τῷ ἐμπορίῳ τῶν λαῶν, ἀπὸ νήσων πολλῶν, τάδε λέγει Κύριος τῇ Σόρ,

\vs{4}Σὺ εἶπας, ἐγὼ περιέθηκα ἐμαυτῇ κάλλος μου, ἐν καρδίᾳ θαλάσσης τῷ Βεελεὶμ υἱοί σου περιέθηκάν σοι κάλλος.
\vs{5}Κέδρος ἐν Σενεὶρ ᾠκοδομήθη σοι, ταινίαι σανίδων κυπαρίσσου ἐκ τοῦ Λιβάνου ἐλήφθησαν, τοῦ ποιῆσαί σοι ἱστοὺς ἐλατίνους,
\vs{6}ἐκ τῆς Βασανίτιδος ἐποίησαν τὰς κώπας σου, τὰ ἱερά σου ἐποίησαν ἐξ ἐλέφαντος, οἴκους ἀλσώδεις ἀπὸ νήσων τῶν Χετιείμ.
\vs{7}Βύσσος μετὰ ποικιλίας ἐξ Αἰγύπτου ἐγένετό σοι στρωμνὴ, τοῦ περιθεῖναί σοι δόξαν, καὶ περιβαλεῖν σε ὑάκινθον καὶ πορφύραν ἐκ τῶν νήσων Ἐλεισαὶ, καὶ ἐγένετο περιβόλαιά σου.

\vs{8}Καὶ οἱ ἄρχοντές σου οἱ κατοικοῦντες Σιδῶνα, καὶ Ἀράδιοι ἐγένοντο κωπηλάται σου· οἱ σοφοί σου Σὸρ, οἳ ἦσαν ἐν σοὶ, οὗτοι κυβερνῆταί σου.
\vs{9}Οἱ πρεσβύτεροι Βιβλίων, καὶ οἱ σοφοὶ αὐτῶν, οἳ ἦσαν ἐν σοί, οὗτοι ἐνίσχυον τὴν βουλήν σου· καὶ πάντα τὰ πλοῖα τῆς θαλάσσης καὶ οἱ κωπηλάται αὐτῶν ἐγένοντό σοι ἐπὶ δυσμὰς δυσμῶν.

\vs{10}Πέρσαι, καὶ Λυδοὶ, καὶ Λίβυες ἦσαν ἐν τῇ δυνάμει σου, ἄνδρες πολεμισταί σου πέλτας καὶ περικεφαλαίας ἐκρέμασαν ἐν σοί, οὗτοι ἔδωκαν τὴν δόξαν σου.
\vs{11}Υἱοὶ Ἀραδίων καὶ ἡ δύναμίς σου ἐπὶ τῶν τειχέων σου· φύλακες ἐν τοῖς πύργοις σου ἦσαν, τὰς φαρέτρας αὐτῶν ἐκρέμασαν ἐπὶ τῶν ὅρμων σου κύκλῳ, οὗτοι ἐτελείωσάν σου τὸ κάλλος.

\vs{12}Καρχηδόνιοι ἔμποροί σου ἀπὸ πλήθους πάσης ἰσχύος σου, ἀργύριον καὶ χρυσίον καὶ σίδηρον καὶ κασσίτερον καὶ μόλιβον ἔδωκαν τὴν ἀγοράν σου.
\vs{13}Ἡ Ἑλλὰς, καὶ ἡ σύμπασα, καὶ τὰ παρατείνοντα, οὗτοι ἐνεπορεύοντό σοι ἐν ψυχαῖς ἀνθρώπων, καὶ σκεύη χαλκᾶ ἔδωκαν τὴν ἐμπορίαν σου.

\vs{14}Ἐξ οἴκου Θογαρμὰ ἵπποι καὶ ἱππεῖς ἔδωκαν ἀγοράν σου.
\vs{15}Υἱοὶ Ῥοδίων ἔμποροί σου, ἀπὸ νήσων ἐπλήθυναν τὴν ἐμπορίαν σου ὀδόντας ἐλεφαντίνους, καὶ τοῖς εἰσαγομένοις ἀντεδίδους τοὺς μισθούς σου,
\vs{16}ἀνθρώπους ἐμπορίαν σου ἀπὸ πλήθους τοῦ συμμίκτου σου, στακτὴν καὶ ποικίλματα ἐκ Θαρσὶς, καὶ Ῥαμὸθ, καὶ Χορχὸρ ἔδωκαν τὴν ἀγοράν σου.
\vs{17}Ἰούδας καὶ οἱ υἱοὶ τοῦ Ἰσραὴλ, οὗτοι ἔμποροί σου, ἐν πράσει σίτου καὶ μύρων καὶ κασίας, καὶ πρῶτον μέλι καὶ ἔλαιον καὶ ῥητίνην ἔδωκαν εἰς τὸν σύμμικτόν σου.
\vs{18}Δαμασκὸς ἔμποροί σου, ἐκ πλήθους πάσης δυνάμεώς σου· οἶνος ἐκ Χελβὼν, καὶ ἔρια ἐκ Μιλήτου, καὶ οἶνον εἰς τὴν ἀγοράν σου ἔδωκαν.

\vs{19}Ἐξ Ἀσὴλ σίδηρος εἰργασμένος, καὶ τροχιὰς ἐν τῷ συμμίκτῳ σου ἐστί.
\vs{20}Δαιδὰν ἔμποροί σου, μετὰ κτηνῶν ἐκλεκτῶν εἰς ἅρματα.
\vs{21}Ἡ Ἀραβία καὶ πάντες οἱ ἄρχοντες Κηδὰρ, οὗτοι ἔμποροί σου διὰ χειρός σου, καμήλους καὶ ἀμνοὺς καὶ κριοὺς ἐν οἷς ἐμπορεύονταί σε.
\vs{22}Ἔμποροι Σαββὰ, καὶ Ῥαμμὰ, οὗτοι ἔμποροί σου, μετὰ πρώτων ἡδυσμάτων, καὶ λίθων χρηστῶν, καὶ χρυσὸν ἔδωκαν τὴν ἀγοράν σου.
\vs{23}Χαῤῥὰ, καὶ Χαναὰ, οὗτοι ἔμποροί σου·
\vs{24}Ἀσσοὺρ καὶ Χαρμὰν ἔμποροί σου, φέροντες ἐμπορίαν ὑάκινθον, καὶ θησαυροὺς ἐκλεκτοὺς δεδεμένους σχοίνίοις, καὶ κυπαρίσσινα.
\vs{25}Πλοῖα ἔμποροί σου, ἐν τῷ πλήθει, ἐν τῷ συμμίκτῳ σου· καὶ ἐνεπλήσθης καὶ ἐβαρύνθης σφόδρα ἐν καρδίᾳ θαλάσσης.

\vs{26}Ἐν ὕδατι πολλῷ ἦγόν σε οἱ κωπηλάται σου, τὸ πνεῦμα τοῦ Νότου συνέτριψέ σε ἐν καρδίᾳ θαλάσσης.
\vs{27}Ἦσαν δυνάμεις σου, καὶ ὁ μισθός σου, καὶ τῶν συμμίκτων σου, καὶ οἱ κωπηλάται σου, καὶ οἱ κυβερνῆταί σου, καὶ οἱ σύμβουλοί σου, καὶ οἱ σύμμικτοί σου ἐκ τῶν συμμίκτων σου, καὶ πάντες οἱ ἄνδρες οἱ πολεμισταί σου οἱ ἐν σοί· καὶ πᾶσα συναγωγή σου ἐν μέσῳ σου πεσοῦνται ἐν καρδίᾳ θαλάσσης, ἐν τῇ ἡμέρᾳ τῆς πτώσεώς σου.

\vs{28}Πρὸς τὴν κραυγὴν τῆς φωνῆς σου οἱ κυβερνῆταί σου φόβῳ φοβηθήσονται.
\vs{29}Καὶ καταβήσονται ἀπὸ τῶν πλοίων πάντες οἱ κωπηλάται, καὶ οἱ ἐπιβάται, καὶ οἱ πρωρεῖς τῆς θαλάσσης ἐπὶ τὴν γῆν στήσονται.
\vs{30}Καὶ ἀλαλάξουσιν ἐπὶ σὲ τῇ φωνῇ αὐτῶν, καὶ κεκράξονται πικρὸν, καὶ ἐπιθήσουσι γῆν ἐπὶ τὴν κεφαλὴν αὐτῶν, καὶ σποδὸν στρώσονται.

\vs{32}Καὶ λήψονται οἱ υἱοὶ αὐτῶν ἐπὶ σὲ θρῆνον, θρήνημα Σόρ·
\vs{33}Πόσον τινὰ εὗρες μισθὸν ἀπὸ τῆς θαλάσσης; Ἐνέπλησας ἔθνη ἀπὸ τοῦ πλήθους σου, καὶ ἀπὸ τοῦ συμμίκτου σου ἐπλούτησας πάντας βασιλεῖς τῆς γῆς.
\vs{34}Νῦν συνετρίβης ἐν θαλάσσῃ, ἐν βάθει ὕδατος ὁ σύμμικτός σου, καὶ πᾶσα ἡ συναγωγή σου ἐν μέσῳ σου· ἔπεσον πάντες οἱ κωπηλάται σου.
\vs{35}Πάντες οἱ κατοικοῦντες τὰς νήσους ἐστύγνασαν ἐπὶ σέ, καὶ οἱ βασιλεῖς αὐτῶν ἐκστάσει ἐξέστησαν, καὶ ἐδάκρυσε τὸ πρόσωπον αὐτῶν.
\vs{36}Ἔμποροι ἀπὸ ἐθνῶν ἐσύρισάν σε, ἀπώλεια ἐγένου, καὶ οὐκέτι ἔσῃ εἰς τὸν αἰῶνα.

\ch{28}
Καὶ ἐγένετο λόγος Κυρίου πρὸς μὲ, λέγων,

\vs{2}Καὶ σὺ υἱὲ ἀνθρώπου εἰπὸν τῷ ἄρχοντι Τύρου, τάδε λέλει Κύριος, ἀνθʼ οὗ ὑψώθη σου ἡ καρδία, καὶ εἶπας, Θεός εἰμι ἐγὼ, κατοικίαν Θεοῦ, καὶ ἔδωκας τὴν καρδίαν σου ὡς καρδίαν Θεοῦ.
\vs{3}Μὴ σοφώτερος εἶ σὺ τοῦ Δανιήλ; ἢ σοφοὶ οὐκ ἐπαίδευσάν σε τῇ ἐπιστήμῃ αὐτῶν;
\vs{4}Μὴ ἐν τῇ ἐπιστήμῃ σου ἢ τῇ φρονήσει σου ἐποίησας σεαυτῷ δύναμιν, καὶ χρυσίον καὶ ἀργύριον ἐν τοῖς θησαυροῖς σου;
\vs{5}Ἐν τῇ πολλῇ ἐπιστήμῃ σου καὶ ἐμπορίᾳ σου ἐπλήθυνας δύναμίν σου, ὑψώθη ἡ καρδία σου ἐν τῇ δυνάμει σου.

\vs{6}Διατοῦτο τάδε λέγει Κύριος, ἐπειδὴ δέδωκας τὴν καρδίαν σου ὡς καρδίαν Θεοῦ,
\vs{7}ἀντὶ τούτου ἰδοὺ ἐγὼ ἐπάγω ἐπὶ σὲ ἀλλοτρίους λοιμοὺς ἀπὸ ἐθνῶν, καὶ ἐκκενώσουσι τὰς μαχαίρας αὐτῶν ἐπὶ σὲ, καὶ ἐπὶ τὸ κάλλος τῆς ἐπιστήμης σου, καὶ στρώσουσιν τὸ κάλλος σου εἰς ἀπώλειαν,
\vs{8}καὶ καταβιβάσουσί σε, καὶ ἀποθανῇ θανάτῳ τραυματιῶν ἐν καρδίᾳ θαλάσσης.
\vs{9}Μὴ λέγων ἐρεῖς, Θεός εἰμι ἐγὼ, ἐνώπιον τῶν ἀναιρούντων σε; σὺ δὲ εἶ ἄνθρωπος, καὶ οὐ Θεός.
\vs{10}Ἐν πλήθει ἀπεριτμήτων ἀπολῇ ἐν χερσὶν ἀλλοτρίων, ὅτι ἐγὼ ἐλάλησα, λέγει Κύριος.

\vs{11}Καὶ ἐγένετο λόγος Κυρίου πρὸς μὲ, λέγων,
\vs{12}υἱὲ ἀνθρώπου, λάβε θρῆνον ἐπὶ τὸν ἄρχοντα Τύρου, καὶ εἰπὸν αὐτῷ, τάδε λέγει Κύριος Κύριος, σὺ ἀποσφράγισμα ὁμοιώσεως,
\vs{13}καὶ στέφανος κάλλους ἐν τῇ τρυφῇ τοῦ παραδείσου τοῦ Θεοῦ ἐγενήθης· πᾶντα λίθον χρηστὸν ἐνδέδεσαι, σάρδιον καὶ τοπάζιον, καὶ σμάραγδον καὶ ἄνθρακα, καὶ σάπφειρον καὶ ἴασπιν, καὶ ἀργύριον καὶ χρυσίον, καὶ λιγύριον καὶ ἀχάτην, καὶ ἀμέθυστον καὶ χρυσόλιθον, καὶ βηρύλλιον καὶ ὀνύχιον, καὶ χρυσίου ἐνέπλήσας τοὺς θησαυρούς σου καὶ τὰς ἀποθήκας σου ἐν σοί.
\vs{14}Ἀφʼ ἧς ἡμέρας ἐκτίσθης σὺ μετὰ τοῦ χερούβ, ἔθηκά σε ἐν ὄρει ἁγίῳ Θεοῦ, ἐγενήθης ἐν μέσῳ λίθων πυρίνων.
\vs{15}Ἐγενήθης σὺ ἄμωμος ἐν ταῖς ἡμέραις σου, ἀφʼ ἧς ἡμέρας σὺ ἐκτίσθης ἕως εὑρέθη τὰ ἀδικήματα ἐν σοί.

\vs{16}Ἀπὸ πλήθους τῆς ἐμπορίας σου ἔπλησας τὰ ταμεῖά σου ἀνομίας, καὶ ἥμαρτες, καὶ ἐτραυματίσθης ἀπὸ ὄρους τοῦ Θεοῦ· καὶ ἤγαγέ σε τὸ χεροὺβ ἐκ μέσου λίθων πυρίνων.
\vs{17}Ὑψώθη ἡ καρδία σου ἐπὶ τῷ κάλλει σου, διεφθάρη ἡ ἐπιστήμη σου μετὰ τοῦ κάλλους σου· διὰ πλῆθος ἁμαρτιῶν σου ἐπὶ τὴν γῆν ἔῤῥιψά σε, ἐναντίον βασιλέων ἔδωκά σε παραδειγματισθῆναι.
\vs{18}Διὰ τὸ πλῆθος τῶν ἁμαρτιῶν σου καὶ τῶν ἀδικιῶν τῆς ἐμπορίας σου, ἐβεβήλωσα τὰ ἱερά σου, καὶ ἐξάξω πῦρ ἐκ μέσου σου, τοῦτο καταφάγεταί σε· καὶ δώσω σε σποδὸν ἐπὶ τῆς γῆς σου ἐναντίον πάντων τῶν ὁρώντων σε.
\vs{19}Καὶ πάντες οἱ ἐπιστάμενοί σε ἐν τοῖς ἔθνεσι στενάξουσιν ἐπὶ σέ· ἀπώλεια ἐγένου, καὶ οὐχ ὑπάρξεις ἔτι εἰς τὸν αἰῶνα.

\vs{20}Καὶ ἐγένετο λόγος Κυρίου πρὸς μὲ, λέγων,
\vs{21}υἱὲ ἀνθρώπου, στήρισον τὸ πρόσωπόν σου ἐπὶ Σιδῶνα, καὶ προφήτευσον ἐπʼ αὐτὴν,
\vs{22}καὶ εἰπὸν,

Τάδε λέγει Κύριος, ἰδοὺ ἐγὼ ἐπὶ σέ Σιδὼν, καὶ ἐνδοξασθήσομαι ἐν σοί, καὶ γνώσῃ ὅτι ἐγώ εἰμι Κύριος, ἐν τῷ ποιῆσαί με ἐν σοὶ κρίματα, καὶ ἁγιασθήσομαι ἐν σοί.
\vs{23}Αἷμα καὶ θάνατος ἐν ταῖς πλατείαις σου, καὶ πεσοῦνται τετραυματισμέμοι μαχαίραις ἐν σοὶ περικύκλῳ σου, καὶ γνώσονται διότι ἐγώ εἰμι Κύριος.
\vs{24}Καὶ οὐκ ἔσονται οὐκέτι ἐν τῷ οἴκῳ τοῦ Ἰσραὴλ σκόλοψ πικρίας καὶ ἄκανθα ὀδύνης ἀπὸ τῶν περικύκλῳ αὐτῶν, τῶν ἀτιμασάντων αὐτοὺς, καὶ γνώσονται ὅτι ἐγώ εἰμι Κύριος.

\vs{25}Τάδε λέγει Κύριος Κύριος, καὶ συνάξω τὸν Ἰσραὴλ ἐκ τῶν ἐθνῶν οὗ διεσκορπίσθησαν ἐκεῖ, καὶ ἁγιασθήσομαι ἐν αὐτοῖς, καὶ ἐνώπιον τῶν λαῶν καὶ τῶν ἐθνῶν·
\vs{26}Καὶ κατοικήσουσι ἐπὶ τῆς γῆς αὐτῶν, ἣν δέδωκα τῷ δούλῳ μου Ἰακὼβ, καὶ κατοικήσουσιν ἐπʼ αὐτῆς ἐν ἐλπίδι, καὶ οἰκοδομήσουσιν οἰκίας, καὶ φυτεύσουσιν ἀμπελῶνας, καὶ κατοικήσουσιν ἐν ἐλπίδι, ὅταν ποιήσω κρίμα ἐν πᾶσι τοῖς ἀτιμάσασιν αὐτοὺς ἐν τοῖς κύκλῳ αὐτῶν· καὶ γνώσονται ὅτι ἐγώ εἰμι Κύριος ὁ Θεὸς αὐτῶν καὶ ὁ Θεὸς τῶν πατέρων αὐτῶν.

\ch{29}
Ἐν τῷ ἔτει τῷ δωδεκάτῳ, ἐν τῷ δεκάτῳ μηνὶ, μιᾷ τοῦ μηνὸς, ἐγένετο λόγος Κυρίου πρὸς μὲ, λέγων,
\vs{2}υἱὲ ἀνθρώπου, στήρισον τὸ πρόσωπόν σου ἐπὶ Φαραὼ βασιλέα Αἰγύπτου, καὶ προφήτευσον ἐπʼ αὐτὸν καὶ ἐπʼ Αἴγυπτον ὅλην,
\vs{3}καὶ εἰπόν,

Τάδε λέγει Κύριος, ἰδοὺ ἐγὼ ἐπὶ Φαραὼ, τὸν δράκοντα τὸν μέγαν τὸν ἐγκαθήμενον ἐν μέσῳ ποταμῶν αὐτοῦ, τὸν λέγοντα, ἐμοί εἰσιν οἱ ποταμοὶ, καὶ ἐγὼ ἐποίησα αὐτούς.
\vs{4}Καὶ ἐγὼ δώσω παγίδας εἰς τὰς σιαγόνας σου, καὶ προσκολλήσω τοὺς ἰχθῦας τοῦ ποταμοῦ σου πρὸς τὰς πτέρυγάς σου, καὶ ἀνάξω σε ἐκ μέσου τοῦ ποταμοῦ σου,
\vs{5}καὶ καταβαλῶ σε ἐν τάχει καὶ πάντας τοὺς ἰχθύας τοῦ ποταμοῦ σου· ἐπὶ πρόσωπον τοῦ πεδίου πέσῇ, καὶ οὐ μὴ συναχθῇς, καὶ οὐ μὴ περισταλῇς, τοῖς θηρίοις τῆς γῆς καὶ τοῖς πετεινοῖς τοῦ οὐρανοῦ δέδωκά σε εἰς κατάβρωμα·
\vs{6}Καὶ γνώσονται πάντες οἱ κατοικοῦντες Αἴγυπτον, ὅτι ἐγώ εἰμι Κύριος· ἀνθʼ ὧν ἐγενήθης ῥάβδος καλαμίνη τῷ οἴκῳ Ἰσραὴλ,
\vs{7}ὅτε ἐπελάβετό σου τῇ χειρὶ αὐτῶν, ἐθλάσθης· καὶ ὅτε ἐπεκρότησεν ἐπʼ αὐτοὺς πᾶσα χεὶρ, καὶ ὅτε ἐπανεπαύσαντο ἐπὶ σέ, συνετρίβης, καὶ συνέκλασας αὐτῶν πᾶσαν ὀσφύν.

\vs{8}Διατοῦτο τάδε λέγει Κύριος, ἰδοὺ ἐγὼ ἐπάγω ἐπὶ σὲ ῥομφαίαν, καὶ ἀπολῶ ἀπὸ σοῦ ἀνθρώπους καὶ κτήνη,
\vs{9}καὶ ἔσται ἡ γῆ Αἰγύπτου ἀπώλεια καὶ ἔρημος· καὶ γνώσονται ὅτι ἐγώ εἰμι Κύριος· ἀντὶ τοῦ λέγειν σε, οἱ ποταμοὶ ἐμοί εἰσι, καὶ ἐγὼ ἐποίησα αὐτοὺς,

\vs{10}Διατοῦτο ἰδοὺ ἐγὼ ἐπὶ σὲ, καὶ ἐπὶ πάντας τοὺς ποταμούς σου, καὶ δώσω γῆν Αἰγύπτου εἰς ἔρημον καὶ ῥομφαίαν καὶ ἀπώλειαν ἀπὸ Μαγδωλοῦ καὶ Συήνης καὶ ἕως ὁρίων Αἰθιόπων.
\vs{11}Οὐ μὴ διέλθῃ ἐν αὐτῇ ποὺς ἀνθρώπου, καὶ ποὺς κτήνους οὐ μὴ διέλθῃ αὐτὴν, καὶ οὐ κατοικηθήσεται τεσσαράκοντα ἔτη.

\vs{12}Καὶ δώσω τὴν γῆν αὐτῆς ἀπώλειαν ἐν μέσῳ γῆς ἠρημωμένης, καὶ αἱ πόλεις αὐτῆς ἐν μέσῳ πόλεων ἠρημωμένων ἔσονται τεσσεράκοντα ἔτη· καὶ διασπερῶ Αἴγυπτον ἐν τοῖς ἔθνεσι, καὶ λικμήσω αὐτοὺς εἰς τὰς χώρας.

\vs{13}Τάδε λέγει Κύριος, μετὰ τεσσεράκοντα ἔτη συνάξω Αἰγυπτίους ἀπὸ τῶν ἐθνῶν οὗ διεσκορπίσθησαν ἐκεῖ,
\vs{14}καὶ ἀποστρέψω τὴν αἰχμαλωσίαν τῶν Αἰγυπτίων, καὶ κατοικίσω αὐτοὺς ἐν γῇ Φαθωρῆς, ἐν τῇ γῇ ὅθεν ἐλήμφθησαν, καὶ ἔσται ἀρχὴ ταπεινὴ
\vs{15}παρὰ πάσας τὰς ἀρχάς· οὐ μὴ ὑψωθῇ ἔτι ἐπὶ τὰ ἔθνη· καὶ ὀλιγοστοὺς αὐτοὺς ποιήσω, τοῦ μὴ εἶναι αὐτοὺς πλείονας ἐν τοῖς ἔθνεσι.
\vs{16}Καὶ οὐκέτι ἔσονται τῷ οἴκῳ Ἰσραὴλ εἰς ἐλπίδα ἀναμιμνήσκουσαν ἀνομίαν, ἐν τῷ ἀκολουθῆσαι αὐτοὺς ὀπίσω ἀυτῶν, καὶ γνώσονται ὅτι ἐγώ εἰμι Κύριος.

\vs{17}Καὶ ἐγένετο ἐν τῷ ἑβδόμῳ καὶ εἰκοστῷ ἔτει, μιᾷ τοῦ μηνὸς τοῦ πρώτου, ἐγένετο λόγος Κυρίου πρὸς μὲ, λέγων,

\vs{18}Υἱὲ ἀνθρώπου, Ναβουχοδονόσορ βασιλεὺς Βαβυλῶνος κατεδουλώσατο τὴν δύναμιν αὐτοῦ δουλείᾳ μεγάλῃ ἐπὶ Τύρου, πᾶσα κεφαλὴ φαλακρὰ, καὶ πᾶς ὦμος μαδῶν· καὶ μισθὸς οὐκ ἐγενήθη αὐτῷ καὶ τῇ δυνάμει αὐτοῦ ἐπὶ Τύρον, καὶ τῆς δουλείας ἧς ἐδούλευσαν ἐπʼ αὐτήν.

\vs{19}Τάδε λέγει Κύριος Κύριος, ἰδοὺ ἐγὼ δίδωμι τῷ Ναβουχοδονόσορ βασιλεῖ Βαβυλῶνος γῆν Αἰγύπτου, καὶ προνομεύσει τὴν προνομὴν αὐτῆς, καὶ σκυλεύσει τὰ σκῦλα αὐτῆς· καὶ ἔσται μισθὸς τῇ δυνάμει αὐτοῦ,
\vs{20}ἀντὶ τῆς λειτουργίας αὐτοῦ ἧς ἐδούλευσεν ἐπὶ Τύρον, δέδωκα αὐτῷ γῆν Αἰγύπτου· τάδε λέγει Κύριος Κύριος.

\vs{21}Ἐν τῇ ἡμέρᾳ ἐκείνῃ ἀνατελεῖ κέρας παντὶ τῷ οἴκῳ Ἰσραὴλ, καὶ σοὶ δώσω στόμα ἀνεῳγμένον ἐν μέσῳ αὐτῶν, καὶ γνώσονται ὅτι ἐγώ εἰμι Κύριος.

\ch{30}
Καὶ ἐγένετο λόγος Κυρίου πρὸς μὲ, λέγων,
\vs{2}υἱὲ ἀνθρώπου, προφήτευσον, καὶ εἰπὸν τάδε λέγει Κύριος, ὢ ὢ ἡμέρα,
\vs{3}ὅτι ἐγγὺς ἡμέρα τοῦ Κυρίου, ἡμέρα πέρας ἐθνῶν ἔσται·

\vs{4}Καὶ ἥξει μάχαιρα ἐπʼ Αἰγυπτίους, καὶ ἔσται ταραχὴ ἐν γῇ Αἰθιοπίᾳ, καὶ συμπεσοῦνται τετραυματισμένοι ἐν Αἰγύπτῳ, καὶ συμπεσεῖται τὰ θεμέλια αὐτῆς.
\vs{5}Πέρσαι, καὶ Κρῆτες, καὶ Λυδοὶ, καὶ Λίβυες, καὶ πάντες οἱ ἐπίμικτοι καὶ τῶν υἱῶν τῆς διαθήκης μου, μαχαίρᾳ πεσοῦνται ἐν αὐτῇ.
\vs{6}Καὶ πεσοῦνται τὰ ἀντιστηρίγματα Αἰγύπτου, καὶ καταβήσεται ἡ ὕβρις τῆς ἰσχύος αὐτῆς ἀπὸ Μαγδωλοῦ ἕως Συήνης, μαχαίρᾳ πεσοῦνται ἐν αὐτῇ, λέγει Κύριος.
\vs{7}Καὶ ἐρημωθήσεται ἐν μέσῳ χωρῶν ἐρημωμένων, καὶ αἱ πόλεις αὐτῶν ἐν μέσῳ πόλεων ἠρημωμένων ἔσονται,
\vs{8}καὶ γνώσονται ὅτι ἐγώ εἰμι Κύριος, ὅταν δῶ πῦρ ἐπʼ Αἴγυπτον, καὶ συντριβῶσι πάντες οἱ βοηθοῦντες αὐτῇ.
\vs{9}Ἐν τῇ ἡμέρᾳ ἐκείνῃ ἐξελεύσονται ἄγγελοι σπεύδοντες ἀφανίσαι τὴν Αἰθιοπίαν, καὶ ἔσται ταραχὴ ἐν αὐτοῖς ἐν τῇ ἡμέρᾳ Αἰγύπτου, ὅτι ἰδοὺ ἥκει.

\vs{10}Τάδε λέγει Κύριος Κύριος, καὶ ἀπολῶ πλῆθος Αἰγυπτίων διὰ χειρὸς Ναβουχοδονόσορ βασιλέως Βαβυλῶνος,
\vs{11}αὐτοῦ καὶ τοῦ λαοῦ αὐτοῦ, λοιμοὶ ἀπὸ ἐθνῶν ἀπεσταλμένοι ἀπολέσαι γῆν· καὶ ἐκκενώσουσι πάντες τὰς μαχαίρας αὐτῶν ἐπʼ Αἴγυπτον, καὶ πλησθήσεται ἡ γῆ τραυματιῶν,
\vs{12}καὶ δώσω τοὺς ποταμοὺς αὐτῶν ἐρήμους, καὶ ἀπολῶ τὴν γῆν, καὶ τὸ πλήρωμα αὐτῆς ἐν χερσὶν ἀλλοτρίων, ἐγὼ Κύριος λελάληκα.

\vs{13}Ὅτι τάδε λέγει Κύριος Κύριος, καὶ ἀπολῶ μεγιστᾶνας ἀπὸ Μέμφεως, καὶ ἄρχοντας Μέμφεως ἐκ γῆς Αἰγύπτου, καὶ οὐκ ἔσονται ἔτι.
\vs{14}Καὶ ἀπολῶ γῆν Φαθωρῆς, καὶ δώσω πῦρ ἐπὶ Τάνιν, καὶ ποιήσω ἐκδίκησιν ἐν Διοσπόλει.
\vs{15}Καὶ ἐκχεῶ τὸν θυμόν μου ἐπὶ Σαὶν τὴν ἰσχὺν Αἰγύπτου, καὶ ἀπολῶ τὸ πλῆθος Μέμφεως,
\vs{16}καὶ δώσω πῦρ ἐπʼ Αἴγυπτον, καὶ ταραχῇ ταραχθήσεται ἡ Συήνη, καὶ ἐν Διοσπόλει ἔσται ἔκρηγμα, καὶ διαχυθήσεται ὕδατα.
\vs{17}Νεανίσκοι Ἡλιουπόλεως καὶ Βουβάστου ἐν μαχαίρᾳ πεσοῦνται, καὶ αἱ γυναῖκες ἐν αἰχμαλωσίᾳ πορεύσονται,
\vs{18}καὶ ἐν Τάφναις συσκοτάσει ἡ ἡμέρα, ἐν τῷ συντρίψαι με ἐκεῖ τὰ σκῆπτρα Αἰγύπτου· καὶ ἀπολεῖται ἐκεῖ ἡ ὕβρις ἰσχύος αὐτῆς, καὶ ταύτην νεφέλη καλύψει, καὶ αἱ θυγατέρες αὐτῆς αἰχμάλωτοι ἀρθήσονται.
\vs{19}Καὶ ποιήσω κρίμα ἐν Αἰγύπτῳ, καὶ γνώσονται ὅτι ἐγὼ εἰμι Κύριος.

\vs{20}Καὶ ἐγένετο ἐν τῷ ἑνδεκάτῳ ἔτει, ἐν τῷ πρώτῳ μηνὶ, ἑβδόμῃ τοῦ μηνὸς, ἐγένετο λόγος Κυρίου πρὸς μὲ, λέγων,
\vs{21}υἱὲ ἀνθρώπου, τοὺς βραχίονας Φαραὼ βασιλέως Αἰγύπτου συνέτριψα, καὶ ἰδοὺ οὐ κατεδεήθη τοῦ δοθῆναι ἴασιν, τοῦ δοθῆναι ἐπʼ αὐτὸν μάλαγμα, τοῦ δοθῆναι ἰσχὺν ἐπιλαβέσθαι μαχαίρας.
\vs{22}Διατοῦτο τάδε λέγει Κύριος Κύριος, ἰδοὺ ἐγὼ ἐπὶ Φαραὼ βασιλέα Αἰγύπτου, καὶ συντρίψω τοὺς βραχίονας αὐτοῦ τοὺς ἰσχυροὺς, καὶ τοὺς τεταμένους, καὶ καταβαλῶ τὴν μάχαιραν αὐτοῦ ἐκ τῆς χειρὸς αὐτοῦ,
\vs{23}καὶ διασπερῶ Αἴγυπτον εἰς τὰ ἔθνη, καὶ λικμήσω αὐτοὺς εἰς τὰς χώρας.

\vs{24}Καὶ κατισχύσω τοὺς βραχίονας βασιλέως Βαβυλῶνος, καὶ δώσω τὴν ῥομφαίαν μου εἰς τὴν χεῖρα αὐτοῦ, καὶ ἐπάξει αὐτὴν ἐπʼ Αἴγυπτον, καὶ προνομεύσει τὴν προνομὴν αὐτῆς, καὶ σκυλεύσει τὰ σκῦλα αὐτῆς.
\vs{25}Καὶ ἐνισχύσω τοὺς βραχίονας βασιλέως Βαβυλῶνος, οἱ δὲ βραχίονες Φαραὼ πεσοῦνται· καὶ γνώσονται ὅτι ἐγώ εἰμι Κύριος, ἐν τῷ δοῦναι τὴν ῥομφαίαν μου εἰς χεῖρας βασιλέως Βαβυλῶνος, καὶ ἐκτενεῖ αὐτὴν ἐπὶ γῆν Αἰγύπτου.
\vs{26}Καὶ διασπερῶ Αἴγυπτον εἰς τὰ ἔθνη, καὶ λικμήσω αὐτοὺς εἰς τὰς χώρας, καὶ γνώσονται πάντες ὅτι ἐγώ εἰμι Κύριος.

\ch{31}
Καὶ ἐγένετο ἐν τῷ ἑνδεκάτῳ ἔτει, ἐν τῷ τρίτῳ μηνὶ, μιᾷ τοῦ μηνὸς, ἐγένετο λόγος Κυρίου πρὸς μὲ, λέγων,
\vs{2}υἱὲ ἀνθρώπου, εἰπὸν πρὸς Φαραὼ βασιλέα Αἰγύπτου καὶ τῷ πλήθει αὐτοῦ,

Τίνι ὡμοίωσας σεαυτὸν ἐν τῷ ὕψει σου;
\vs{3}Ἰδοὺ Ἀσσοὺρ κυπάρισσος ἐν τῷ Λιβάνῳ, καὶ καλὸς ταῖς παραφυάσι, καὶ ὑψηλὸς τῷ μεγέθει, εἰς μέσον νεφελῶν ἐγένετο ἡ ἀρχὴ αὐτοῦ.
\vs{4}ὕδωρ ἐξέθρεψεν αὐτὸν, ἡ ἄβυσσος ὕψωσεν αὐτὸν, τοὺς ποταμοὺς αὐτῆς ἤγαγε κύκλῳ τῶν φυτῶν αὐτοῦ, καὶ τὰ συστήματα αὐτῆς ἐξαπέστειλεν εἰς πάντα τὰ ξύλα τοῦ πεδίου.
\vs{5}Ἕνεκεν τούτου ὑψώθη τὸ μέγεθος αὐτοῦ παρὰ πάντα τὰ ξύλα τοῦ πεδίου, καὶ ἐπλατύνθησαν οἱ κλάδοι αὐτοῦ ἀφʼ ὕδατος πολλοῦ.
\vs{6}Ἐν ταῖς παραφυάσιν αὐτοῦ ἐνόσσευσαν πάντα τὰ πετεινὰ τοῦ οὐρανοῦ, καὶ ὑποκάτω τῶν κλάδων αὐτοῦ ἐγεννῶσαν πάντα τὰ θηρία τοῦ πεδίου, ἐν τῇ σκιᾷ αὐτοῦ κατῴκησε πᾶν πλῆθος ἐθνῶν.
\vs{7}Καὶ ἐγένετο καλὸς ἐν τῷ ὕψει αὐτοῦ διὰ τὸ πλῆθος τῶν κλάδων αὐτοῦ, ὅτι ἐγενήθησαν αἱ ῥίζαι αὐτοῦ εἰς ὕδωρ πολύ.
\vs{8}Καὶ κυπάρισσοι τοιαῦται ἐν τῷ παραδείσῳ τοῦ Θεοῦ, καὶ αἱ πίτυες οὐχ ὅμοιαι ταῖς παραφυάσιν αὐτοῦ, καὶ ἐλάται οὐκ ἐγένοντο ὅμοιαι τοῖς κλάδοις αὐτοῦ· πᾶν ξύλον ἐν τῷ παραδείσῳ τοῦ Θεοῦ οὐχ ὡμοιώθη αὐτῷ ἐν τῷ κάλλει αὐτοῦ,
\vs{9}διὰ τὸ πλῆθος τῶν κλάδων αὐτοῦ· καὶ ἐζήλωσεν αὐτὸν τὰ ξύλα τοῦ παραδείσου τῆς τρυφῆς τοῦ Θεοῦ.

\vs{10}Διατοῦτο τάδε λέγει Κύριος, ἀνθʼ ὧν ἐγένου μέγας τῷ μεγέθει, καὶ ἔδωκας τὴν ἀρχήν σου εἰς μέσον νεφελῶν, καὶ εἶδον ἐν τῷ ὑψωθῆναι αὐτόν·
\vs{11}Καὶ παρέδωκα αὐτὸν εἰς χεῖρας ἄρχοντος ἐθνῶν, καὶ ἐποίησε τὴν ἀπώλειαν αὐτοῦ.
\vs{12}Καὶ ἐξωλόθρευσαν αὐτὸν ἀλλότριοι λοιμοὶ ἀπὸ ἐθνῶν, καὶ κατέβαλον αὐτὸν ἐπὶ τῶν ὀρέων· ἐν πάσαις ταῖς φάραγξιν ἔπεσαν οἱ κλάδοι αὐτοῦ, καὶ συνετρίβη τὰ στελέχη αὐτοῦ ἐν παντὶ πεδίῳ τῆς γῆς, καὶ κατέβησαν ἀπὸ τῆς σκέπης αὐτῶν πάντες οἱ λαοὶ τῶν ἐθνῶν, καὶ ἠδάφισαν αὐτόν.

\vs{13}Ἐπὶ τὴν πτῶσιν αὐτοῦ ἀνεπαύσαντο πάντα τὰ πετεινὰ τοῦ οὐρανοῦ, καὶ ἐπὶ τὰ στελέχη αὐτοῦ ἐγίνοντο πάντα τὰ θηρία τοῦ ἀγροῦ,
\vs{14}ὅπως μὴ ὑψωθῶσιν ἐν τῷ μεγέθει αὐτῶν πάντα τὰ ξύλα τὰ ἐν τῷ ὕδατι· καὶ ἔδωκαν τὴν ἀρχὴν αὐτῶν εἰς μέσον νεφελῶν, καὶ οὐκ ἔστησαν ἐν τῷ ὕψει αὐτῶν πρὸς αὐτά· πάντες οἱ πίνοντες ὕδωρ, πάντες ἐδόθησαν εἰς θάνατον, εἰς γῆς βάθος, ἐν μέσῳ υἱῶν ἀνθρώπων πρὸς καταβαίνοντας εἰς βόθρον.

\vs{15}Τάδε λέγει Κύριος Κύριος, ἐν ᾗ ἡμέρᾳ κατέβη εἰς ᾅδου, ἐπένθησεν αὐτὸν ἡ ἄβυσσος· καὶ ἐπέστησα τοὺς ποταμοὺς αὐτῆς, καὶ ἐκώλυσα πλῆθος ὕδατος, καὶ ἐσκότασεν ἐπʼ αὐτὸν ὁ Λίβανος, πάντα τὰ ξύλα τοῦ πεδίου ἐπʼ αὐτῷ ἐξελύθησαν.
\vs{16}Ἀπὸ τῆς φωνῆς τῆς πτώσεως αὐτοῦ ἐσείσθησαν τὰ ἔθνη, ὅτε κατεβίβαζον αὐτὸν εἰς ᾅδου μετὰ τῶν καταβαινόντων εἰς λάκκον, καὶ παρεκάλουν αὐτὸν ἐν γῇ πάντα τὰ ξύλα τῆς τρυφῆς, καὶ τὰ ἐκλεκτὰ τοῦ Λιβάνου πάντα τὰ πίνοντα ὕδωρ.
\vs{17}Καὶ γὰρ αὐτοὶ κατέβησαν μετʼ αὐτοῦ εἰς ᾅδου ἐν τοῖς τραυματίαις ἀπὸ μαχαίρας, καὶ τὸ σπέρμα αὐτοῦ οἱ κατοικοῦντες ὑπὸ τὴν σκέπην αὐτοῦ ἐν μέσῳ τῆς ζωῆς αὐτῶν ἀπώλοντο.

\vs{18}Τίνι ὡμοιώθης; κατάβηθι, καὶ καταβιβάσθητι μετὰ τῶν ξύλων τῆς τρυφῆς εἰς γῆς βάθος· ἐν μέσῳ ἀπεριτμήτων κοιμηθήσῃ μετὰ τραυματιῶν μαχαίρας· οὕτως Φαραὼ καὶ τὸ πλῆθος τῆς ἰσχύος αὐτοῦ, λέγει Κύριος Κύριος.

\ch{32}
Καὶ ἐγένετο ἐν τῷ δεκάτῳ ἔτει, ἐν τῷ δεκάτῳ μηνὶ, μιᾷ τοῦ μηνὸς, ἐγένετο λόγος Κυρίου πρὸς μὲ, λέγων,
\vs{2}υἱὲ ἀνθρώπου, λάβε θρῆνον ἐπὶ Φαραὼ βασιλέα Αἰγύπτου, καὶ ἐρεῖς αὐτῷ,

Λέοντι ἐθνῶν ὡμοιώθης σὺ, καὶ ὡς δράκων ὁ ἐν τῇ θαλάσσῃ, καὶ ἐκεράτιζες τοῖς ποταμοῖς σου, καὶ ἐτάρασσες ὕδωρ τοῖς ποσί σου, καὶ κατεπάτεις τοὺς ποταμούς σου.

\vs{3}Τάδε λέγει Κύριος, καὶ περιβαλῶ ἐπὶ σὲ δίκτυα λαῶν πολλῶν, καὶ ἀνάξω σε ἐν τῷ ἀγκίστρῳ μου,
\vs{4}καὶ ἐκτενῶ σε ἐπὶ τὴν γῆν· πεδία πλησθήσεται, καὶ ἐπικαθιῶ ἐπὶ σὲ πάντα τὰ πετεινὰ τοῦ οὐρανοῦ, καὶ ἐμπλήσω πάντα τὰ θηρία πάσης τῆς γῆς.
\vs{5}Καὶ δώσω τὰς σάρκας σου ἐπὶ τὰ ὄρη, καὶ ἐμπλήσω ἀπὸ τοῦ αἵματός σου.
\vs{6}Καὶ ποτισθήσεται ἡ γῆ ἀπὸ τῶν προχωρημάτων σου, ἀπὸ τοῦ πλήθους σου ἐπὶ τῶν ὀρέων· φάραγγας ἐμπλήσω ἀπὸ σοῦ,
\vs{7}καὶ κατακαλύψω ἐν τῷ σβεσθῆναί σε οὐρανὸν, καὶ συσκοτάσω τὰ ἄστρα αὐτοῦ, ἥλιον ἐν νεφέλῃ καλύψω, καὶ σελήνη οὐ μὴ φάνῃ τὸ φῶς αὐτῆς.
\vs{8}Πάντα τὰ φαίνοντα φῶς ἐν τῷ οὐρανῷ, συσκοτάσουσιν ἐπὶ σέ· καὶ δώσω σκότος ἐπὶ τὴν γῆν, λέγει Κύριος Κύριος.

\vs{9}Καὶ παροργιῶ καρδίαν λαῶν πολλῶν, ἡνίκα ἂν ἄγω αἰχμαλωσίαν σου εἰς τὰ ἔθνη, εἰς γῆν ἣν οὐκ ἔγνως.
\vs{10}Καὶ στυγνάσουσιν ἐπὶ σὲ ἔθνη πολλὰ, καὶ οἱ βασιλεῖς αὐτῶν ἐκστάσει ἐκστήσονται, ἐν τῷ πετᾶσθαι τὴν ῥομφαίαν μου ἐπὶ πρόσωπα αὐτῶν, προσδεχόμενοι τὴν πτῶσιν αὐτῶν ἀφʼ ἡμέρας πτώσεώς σου.

\vs{11}Ὅτι τάδε λέγει Κύριος Κύριος, ῥομφαία βασιλέως Βαβυλῶνος ἥξει σοι
\vs{12}ἐν μαχαίραις γιγάντων, καὶ καταβαλῶ τὴν ἰσχύν σου, λοιμοὶ ἀπὸ ἐθνῶν πάντες, καὶ ἀπολοῦσι τὴν ὕβριν Αἰγύπτου, καὶ συντριβήσεται πᾶσα ἡ ἰσχὺς αὐτῆς.
\vs{13}Καὶ ἀπολῶ πάντα τὰ κτήνη αὐτῆς ἀφʼ ὕδατος πολλοῦ, καὶ οὐ μὴ ταράξῃ αὐτὸ πούς ἀνθρώπου ἔτι, καὶ ἴχνος κτηνῶν οὐ μὴ καταπατήσῃ αὐτό.
\vs{14}Οὕτως τότε ἡσυχάσει τὰ ὕδατα αὐτῶν, καὶ οἱ ποταμοὶ αὐτῶν ὡς ἔλαιον πορεύσονται, λέγει Κύριος,
\vs{15}ὅταν δῶ Αἴγυπτον εἰς ἀπώλειαν, καὶ ἐρημωθῇ ἡ γῆ σὺν· τῇ πληρώσει αὐτῆς, ὅταν διασπερῶ πάντας τοὺς κατοικοῦντας ἐν αὐτῇ, καὶ γνώσονται ὅτι ἐγώ εἰμι Κύριος.
\vs{16}Θρῆνός ἐστι, καὶ θρηνήσεις αὐτὸν, καὶ αἱ θυγατέρες τῶν ἐθνῶν θρηνήσουσιν αὐτὸν, ἐπʼ Αἴγυπτον, καὶ ἐπὶ πᾶσαν τὴν ἰσχὺν αὐτῆς θρηνήσουσιν αὐτὴν, λέγει Κύριος Κύριος.

\vs{17}Καὶ ἐγενήθη ἐν τῷ δωδεκάτῳ ἔτει τοῦ πρώτου μηνὸς, πεντεκαιδεκάτῃ τοῦ μηνὸς, ἐγένετο λόγος Κυρίου πρὸς μὲ, λέγων,

\vs{18}Υἱὲ ἀνθρώπου, θρήνησον ἐπὶ τὴν ἰσχὺν Αἰγύπτου, καὶ καταβιβάσουσιν αὐτῆς τὰς θυγατέρας τὰ ἔθνη νεκρὰς εἰς τὸ βάθος τῆς γῆς, πρὸς τοὺς καταβαίνοντας εἰς βόθρον·
\vs{20}Ἐν μέσῳ μαχαίρας τραυματιῶν πεσοῦνται μετʼ αὐτοῦ, καὶ κοιμηθήσεται πᾶσα ἡ ἰσχὺς αὐτοῦ·
\vs{21}καὶ ἐροῦσί σοι οἱ γίγαντες, ἐν βάθει βόθρου γίνου, τίνος κρείττων εἶ; καὶ κατάβηθι, καὶ κοιμήθητι μετὰ ἀπεριτμήτων, ἐν μέσῳ τραυματιῶν μαχαίρας.

\vs{22}Ἐκεῖ Ἀσσοὺρ καὶ πᾶσα ἡ συναγωγὴ αὐτοῦ, πάντες τραυματίαι ἐκεῖ ἐδόθησαν, καὶ ἡ ταφὴ αὐτῶν ἐν βάθει βόθρου, καὶ ἐγενήθη ἡ συναγωγὴ αὐτοῦ περικύκλῳ τοῦ μνήματος αὐτοῦ, πάντες οἱ τραυματίαι οἱ πεπτωκότες μαχαίρᾳ,
\vs{23}οἱ δόντες τὸν φόβον αὐτῶν ἐπὶ γῆς ζωῆς.

\vs{24}Ἐκεῖ Αἰλὰμ καὶ πᾶσα ἡ δύναμις αὐτοῦ περικύκλῳ τοῦ μνήματος αὐτοῦ, πάντες οἱ τραυματίαι οἱ πεπτωκότες μαχαίρᾳ, καὶ οἱ καταβαίνοντες ἀπερίτμητοι εἰς γῆς βάθος, οἱ δεδωκότες αὐτῶν φόβον ἐπὶ γῆς ζωῆς· καὶ ἐλάβοσαν τὴν βάσανον αὐτῶν μετὰ τῶν καταβαινόντων εἰς βόθρον,
\vs{25}ἐν μέσῳ τραυματιῶν.

\vs{26}Ἐκεῖ ἐδόθησαν Μοσὸχ καὶ Θοβὲλ, καὶ πᾶσα ἡ ἰσχὺς αὐτοῦ περικύκλῳ τοῦ μνήματος αὐτοῦ, πάντες τραυματίαι αὐτοῦ, πάντες ἀπερίτμητοι τραυματίαι ἀπὸ μαχαίρας, οἱ δεδωκότες τὸν φόβον αὐτῶν ἐπὶ γῆς ζωῆς.
\vs{27}Καὶ ἐκοιμήθησαν μετὰ τῶν γιγάντων τῶν πεπτωκότων ἀπʼ αἰῶνος, οἳ κατέβησαν εἰς ᾅδου ἐν ὅπλοις πολεμικοῖς, καὶ ἔθηκαν τὰς μαχαίρας αὐτῶν ὑπὸ τὰς κεφαλὰς αὐτῶν, καὶ ἐγενήθησαν αἱ ἀνομῖαι αὐτῶν ἐπὶ τῶν ὀστέων αὐτῶν, ὅτι ἐξεφόβησαν πάντας ἐν τῇ ζωῇ αὐτῶν.
\vs{28}Καὶ σὺ ἐν μέσῳ ἀπεριτμήτων κοιμηθήσῃ μετὰ τετραυματισμένων μαχαίρᾳ.

\vs{29}Ἐκεῖ ἐδόθησαν οἱ ἄρχοντες Ἀσσοὺρ, οἱ δόντες τὴν ἰσχὺν αὐτοῦ εἰς τραῦμα μαχαίρας· οὗτοι μετὰ τραυματιῶν ἐκοιμήθησαν, μετὰ καταβαινόντων εἰς βόθρον.

\vs{30}Ἐκεῖ οἱ ἄρχοντες τοῦ Βοῤῥᾶ πάντες στρατηγοὶ Ἀσσοὺρ, οἱ καταβαίνοντες τραυματίαι, σὺν τῷ φόβῳ αὐτῶν καὶ τῇ ἰσχύϊ αὐτῶν ἐκοιμήθησαν ἀπερίτμητοι μετὰ τραυματιῶν μαχαίρας, καὶ ἀπήνεγκαν τὴν βάσανον αὐτῶν μετὰ τῶν καταβαινόντων εἰς βόθρον.

\vs{31}Ἐκείνους ὄψεται βασιλεὺς Φαραὼ, καὶ παρακληθήσεται ἐπὶ πᾶσαν τὴν ἰσχὺν αὐτῶν, λέγει Κύριος Κύριος.
\vs{32}Ὅτι δέδωκα τὸν φόβον αὐτοῦ ἐπὶ γῆς ζωῆς, καὶ κοιμηθήσεται ἐν μέσῳ ἀπεριτμήτων μετὰ τραυματιῶν μαχαίρας Φαραὼ, καὶ πᾶν τὸ πλῆθος αὐτοῦ μετʼ αὐτοῦ, λέγει Κύριος Κύριος.

\ch{33}
Καὶ ἐγένετο λόγος Κυρίου πρὸς μὲ, λέγων,
\vs{2}υἱὲ ἀνθρώπου, λάλησον τοῖς υἱοῖς τοῦ λαοῦ σου, καὶ ἐρεῖς πρὸς αὐτοὺς,

Γῆ ἐφʼ ἣν ἂν ἐπάγω ῥομφαίαν, καὶ λάβῃ ὁ λαὸς τῆς γῆς ἄνθρωπον ἕνα ἐξ αὐτῶν, καὶ δῶσιν αὐτὸν ἑαυτοῖς εἰς σκοπὸν,
\vs{3}καὶ ἴδῃ τὴν ῥομφαίαν ἐρχομένην ἐπὶ τὴν γῆν, καὶ σαλπίσῃ τῇ σάλπιγγι, καὶ σημάνῃ τῷ λαῷ,
\vs{4}καὶ ἀκούσῃ ὁ ἀκούσας τῆς φωνῆς τῆς σάλπιγγος, καὶ μὴ φυλάξηται, καὶ ἐπέλθῃ ἡ ῥομφαία, καὶ καταλάβῃ αὐτὸν, τὸ αἷμα αὐτοῦ ἐπὶ τῆς κεφαλῆς αὐτοῦ ἔσται.
\vs{5}Ὅτι τὴν φωνὴν τῆς σάλπιγγος ἀκούσας οὐκ ἐφυλάξατο, τὸ αἷμα αὐτοῦ ἐπʼ αὐτοῦ ἔσται· καὶ οὗτος ὅτι ἐφυλάξατο, τὴν ψυχὴν αὐτοῦ ἐξείλατο.

\vs{6}Καὶ ὁ σκοπὸς ἐὰν ἴδῃ τὴν ῥομφαίαν ἐρχομένην, καὶ μὴ σημάνῃ τῇ σάλπιγγι, καὶ ὁ λαὸς μὴ φυλάξηται, καὶ ἐλθοῦσα ἡ ῥομφαία λάβῃ ἐξ αὐτῶν ψυχὴν, αὕτη διὰ τὴν αὐτῆς ἀνομίαν ἐλήφθη, καὶ τὸ αἷμα ἐκ χειρὸς τοῦ σκοποῦ ἐκζητήσω.

\vs{7}Καὶ σὺ υἱὲ ἀνθρώπου, σκοπὸν δέδωκά σε τῷ οἴκῳ Ἰσραὴλ, καὶ ἀκούσῃ ἐκ στόματός μου λόγον.
\vs{8}Ἐν τῷ εἰπεῖν με τῷ ἁμαρτωλῷ, θανάτῳ θανατωθήσῃ, καὶ μὴ λαλήσῃς τοῦ φυλάξασθαι τὸν ἀσεβῆ ἀπὸ τῆς ὁδοῦ αὐτοῦ, αὐτὸς ὁ ἄνομος τῇ ἀνομίᾳ αὐτοῦ ἀποθανεῖται, τὸ δὲ αἷμα αὐτοῦ ἐκ τῆς χειρός σου ἐκζητήσω.
\vs{9}Σὺ δὲ ἐὰν προαπαγγείλῃς τῷ ἀσεβεῖ τὴν ὁδὸν αὐτοῦ ἀποστρέψαι ἀπʼ αὐτῆς, καὶ μὴ ἀποστρέψῃ ἀπὸ τῆς ὁδοῦ αὐτοῦ, οὗτος τῇ ἀσεβείᾳ αὐτοῦ ἀποθανεῖται, καὶ σὺ τὴν ψυχὴν σαυτοῦ ἐξῄρησαι.

\vs{10}Καὶ σὺ υἱὲ ἀνθρώπου, εἰπὸν τῷ οἴκῳ Ἰσραὴλ, οὓτως ἐλαλήσατε λέγοντες, αἱ πλάναι ἡμῶν, καὶ αἱ ἀνομίαι ἡμῶν ἐφʼ ἡμῖν εἰσι, καὶ ἐν αὐταῖς ἡμεῖς τηκόμεθα, καὶ πῶς ζησόμεθα;
\vs{11}Εἰπὸν αὐτοῖς, ζῶ ἐγὼ, τάδε λέγει Κύριος, οὐ βούλομαι τὸν θάνατον τοῦ ἀσεβοῦς, ὡς ἀποστρέψαι τὸν ἀσεβῆ ἀπὸ τῆς ὁδοῦ αὐτοῦ, καὶ ζῇν αὐτόν· ἀποστροφῇ ἀποστρέψατε ἀπὸ τῆς ὁδοῦ ὑμῶν· καὶ ἱνατί ἀποθνήσκετε οἶκος Ἰσραήλ;

\vs{12}Εἰπὸν πρὸς τοὺς υἱοὺς τοῦ λαοῦ σου, δικαιοσύνη δικαίου οὐ μὴ ἐξελεῖται αὐτὸν, ἐν ᾗ ἂν ἡμέρᾳ πλανηθῇ· καὶ ἀνομία ἀσεβοῦς οὐ μὴ κακώσῃ αὐτὸν, ἐν ᾗ ἂν ἡμέρᾳ ἀποστρέψῃ ἀπὸ τῆς ἀνομίας αὐτοῦ· καὶ δίκαιος οὐ μὴ δύνηται σωθῆναι.

\vs{13}Ἐν τῷ εἰπεῖν με τῷ δικαίῳ, οὗτος πέποιθεν ἐπὶ τῇ δικαιοσύνῃ αὐτοῦ, καὶ ποιήσει ἀνομίαν, πᾶσαι αἱ δικαιοσύναι αὐτοῦ οὐ μὴ ἀναμνησθῶσιν, ἐν τῇ ἀδικίᾳ αὐτοῦ ᾗ ἐποίησεν ἐν αὐτῇ ἀποθανεῖται.

\vs{14}Καὶ ἐν τῷ εἰπεῖν με τῷ ἀσεβεῖ, θανάτῳ θανατωθήσῃ, καὶ ἀποστρέψει ἀπὸ τῆς ἁμαρτίαις αὐτοῦ, καὶ ποιήσει κρίμα καὶ δικαιοσύνην,
\vs{15}καὶ ἐνεχύρασμα ἀποδοῖ, καὶ ἁρπάγματα ἀποτίσει, ἐν προστάγμασιν ζωῆς διαπορεύηται, τοῦ μὴ ποιῆσαι ἄδικον, ζωῇ ζήσεται, καὶ οὐ μὴ ἀποθάνῃ.
\vs{16}Πᾶσαι αἱ ἁμαρτίαι αὐτοῦ ἃς ἥμαρτεν, οὐ μὴ ἀναμνησθῶσιν, ὅτι κρίμα καὶ δικαιοσύνην ἐποίησεν, ἐν αὐταῖς ζήσεται.

\vs{17}Καὶ ἐροῦσιν οἱ υἱοὶ τοῦ λαοῦ σου, οὐκ εὐθεῖα ἡ ὁδὸς τοῦ Κυρίου· καὶ αὕτη ἡ ὁδὸς αὐτῶν οὐκ εὐθεῖα.
\vs{18}Ἐν τῷ ἀποστρέψαι δίκαιον ἀπὸ τῆς δικαιοσύνης αὐτοῦ, καὶ ποιήσει ἀνομίας, καὶ ἀποθανεῖται ἐν αὐταῖς.
\vs{19}Καὶ ἐν τῷ ἀποστρέψαι τὸν ἁμαρτωλὸν ἀπὸ τῆς ἀνομίας αὐτοῦ, καὶ ποιήσει κρίμα καὶ δικαιοσύνην, ἐν αὐτοῖς αὐτὸς ζήσεται.
\vs{20}Καὶ τοῦτό ἐστιν ὃ εἴπατε, οὐκ εὐθεῖα ἡ ὁδὸς Κυρίου· ἕκαστον ἐν ταῖς ὁδοῖς αὐτοῦ κρινῶ ὑμᾶς, οἶκος Ἰσραήλ.

\vs{21}Καὶ ἐγενήθη ἐν τῷ δεκάτῳ ἔτει, ἐν τῷ δωδεκάτῳ μηνὶ, πέμπτῃ τοῦ μηνὸς τῆς αἰχμαλωσίας ἡμῶν, ἦλθε πρὸς μὲ ὁ ἀνασωθεὶς ἀπὸ Ἱερουσαλὴμ, λέγων, ἑάλω ἡ πόλις.
\vs{22}Καὶ χεὶρ Κυρίου ἐγενήθη ἐπʼ ἐμὲ ἑσπέρας πρὶν ἐλθεῖν αὐτὸν, καὶ ἤνοιξέ μου τὸ στόμα, ὡς ἦλθε πρὸς μὲ τὸ πρωΐ· καὶ ἀνοιχθὲν τὸ στόμα μου, οὐ συνεσχέθη ἔτι.

\vs{23}Καὶ ἐγενήθη λόγος Κυρίου πρὸς μὲ, λέγων,
\vs{24}υἱὲ ἀνθρώπου, οἱ κατοικοῦντες τὰς ἠρημωμένας ἐπὶ τῆς γῆς τοῦ Ἰσραὴλ, λέγουσιν, εἷς ἦν Ἁβραὰμ, καὶ κατέσχε τὴν γῆν, καὶ ἡμεῖς πλείους ἐσμὲν, ἡμῖν δέδοται ἡ γῆ εἰς κατάσχεσιν.

\vs{25}Διατοῦτο εἰπὸν αὐτοῖς, τάδε λέγει Κύριος Κύριος,
\vs{27}ζῶ ἐγὼ, εἰ μὴν οἱ ἐν ταῖς ἠρημωμέναις μαχαίραις πεσοῦνται, καὶ οἱ ἐπὶ προσώπου τοῦ πεδίου τοῖς θηρίοις τοῦ ἀγροῦ δοθήσονται εἰς κατάβρωμα, καὶ τοὺς ἐν ταῖς τετειχισμέναις, καὶ τοὺς ἐν τοῖς σπηλαίοις θανάτῳ ἀποκτενῶ.
\vs{28}Καὶ δώσω τὴν γῆν ἔρημον, καὶ ἀπολεῖται ἡ ὕβρις τῆς ἰσχύος αὐτῆς, καὶ ἐρημωθήσεται τὰ ὄρῃ τοῦ Ἰσραὴλ διὰ τὸ μὴ εἶναι διαπορευόμενον.
\vs{29}Καὶ γνώσονται ὅτι ἐγώ εἰμι Κύριος· καὶ ποιήσω τὴν γῆν αὐτῶν ἔρημον, καὶ ἐρημωθήσεται διὰ πάντα τὰ βδελύγματα αὐτῶν ἃ ἐποίησαν.

\vs{30}Καὶ σὺ υἱὲ ἀνθρώπου, οἱ υἱοὶ τοῦ λαοῦ σου οἱ λαλοῦντες περὶ σοῦ παρὰ τὰ τείχη καὶ ἐν τοῖς πυλῶσι τῶν οἰκιῶν, καὶ λαλοῦσιν ἄνθρωπος τῷ ἀδελφῷ αὐτοῦ, λέγοντες, συνέλθωμεν, καὶ ἀκούσωμεν τὰ ἐκπορευόμενα παρὰ Κυρίου.
\vs{31}Ἔρχονται πρὸς σὲ, ὡς συμπορεύεται λαὸς, καὶ κάθηνται ἐναντίον σου, καὶ ἀκούουσι τὰ ῥήματά σου, καὶ αὐτὰ οὐ μὴ ποιήσουσιν, ὅτι ψεῦδος ἐν τῷ στόματι αὐτῶν, καὶ ὀπίσω τῶν μιασμάτων ἡ καρδία αὐτῶν.
\vs{32}Καὶ γίνῃ αὐτοῖς ὡς φωνὴ ψαλτηρίου ἡδυφώνου εὐαρμόστου, καὶ ἀκούσονταί σου τὰ ῥήματα, καὶ οὐ μὴ ποιήσουσιν αὐτά.
\vs{33}Καὶ ἡνίκα ἐὰν ἔλθῃ, ἐροῦσιν, ἰδοὺ ἥκει· καὶ γνώσονται ὅτι προφήτης ἦν ἐν μέσῳ αὐτῶν.

\ch{34}
Καὶ ἐγένετο λόγος Κυρίου πρὸς μὲ, λέγων,
\vs{2}υἱὲ ἀνθρώπου, προφήτευσον ἐπὶ τοὺς ποιμένας τοῦ Ἰσραὴλ, προφήτευσον, καὶ εἰπὸν τοῖς ποιμέσι,

Τάδε λέγει Κύριος Κύριος, ὦ ποιμένες Ἰσραὴλ, μὴ βόσκουσι ποιμένες ἑαυτούς; οὐ τὰ πρόβατα βόσκουσιν οἱ ποιμένες;
\vs{3}Ἰδοὺ τὸ γάλα κατέσθετε, καὶ τὰ ἔρια περιβάλλεσθε, καὶ τὸ παχὺ σφάζετε, καὶ τὰ πρόβατά μου οὐ βόσκετε.
\vs{4}Τὸ ἠσθενηκὸς οὐκ ἐνισχύσατε, καὶ τὸ κακῶς ἔχον οὐκ ἐσωματοποιήσατε, καὶ τὸ συντετριμμένον οὐ κατεδήσατε, καὶ τὸ πλανώμενον οὐκ ἀπεστρέψατε, καὶ τὸ ἀπολωλὸς οὐκ ἐζητήσατε, καὶ τὸ ἰσχυρὸν κατειργάσασθε μόχθῳ.
\vs{5}Καὶ διεσπάρη τὰ πρόβατά μου, διὰ τὸ μὴ εἶναι ποιμένας, καὶ ἐγενήθη εἰς κατάβρωμα πᾶσι τοῖς θηρίοις τοῦ ἀγροῦ.
\vs{6}Καὶ διεσπάρη τὰ πρόβατά μου ἐν παντὶ ὄρει, καὶ ἐπὶ πᾶν βουνὸν ὑψηλὸν, καὶ ἐπὶ προσώπου τῆς γῆς διεσπάρη, καὶ οὐκ ἦν ὁ ἐκζητῶν οὐδὲ ὁ ἀποστρέφων.

\vs{7}Διατοῦτο ποιμένες ἀκούσατε λόγον Κυρίου.
\vs{8}Ζῶ ἐγὼ, λέγει Κύριος Κύριος, εἰ μὴν ἀντὶ τοῦ γενέσθαι τὰ πρόβατά μου εἰς προνομὴν, καὶ γενέσθαι τὰ πρόβατά μου εἰς κατάβρωμα πᾶσι τοῖς θηρίοις τοῦ πεδίου, παρὰ τὸ μὴ εἶναι ποιμένας, καὶ οὐκ ἐξεζήτησαν οἱ ποιμένες τὰ πρόβατά μου, καὶ ἐβόσκησαν οἱ ποιμένες ἑαυτοὺς, τὰ δὲ πρόβατά μου οὐκ ἐβόσκησαν·
\vs{9}Ἀντὶ τούτου ποιμένες,
\vs{10}τάδε λέγει Κύριος Κύριος, ἰδοὺ ἐγὼ ἐπὶ τοὺς ποιμένας, καὶ ἐκζητήσω τὰ πρόβατά μου ἐκ τῶν χειρῶν αὐτῶν, καὶ ἀποστρέψω αὐτοὺς τοῦ μὴ ποιμαίνειν τὰ πρόβατά μου, καὶ οὐ βοσκήσουσιν ἔτι οἱ ποιμένες αὐτά· καὶ ἐξελοῦμαι τὰ πρόβατά μου ἐκ τοῦ στόματος αὐτῶν, καὶ οὐκ ἔσονται αὐτοῖς ἔτι εἰς κατάβρωμα.

\vs{11}Διότι τάδε λέγει Κύριος Κύριος, ἰδοὺ ἐγὼ ἐκζητήσω τὰ πρόβατά μου, καὶ ἐπισκέψομαι αὐτά.
\vs{12}Ὥσπερ ζητεῖ ὁ ποιμὴν τὸ ποίμνιον αὐτοῦ ἐν ἡμέρᾳ ὅτʼ ἄν ᾖ γνόφος καὶ νεφέλη ἐν μέσῳ προβάτων διακεχωρισμένων, οὕτως ἐκζητήσω τὰ πρόβατά μου, καὶ ἀπελάσω αὐτὰ ἀπὸ παντὸς τόπου οὗ διεσπάρησαν ἐκεῖ ἐν ἡμέρᾳ νεφέλης καὶ γνόφου.
\vs{13}Καὶ ἐξάξω αὐτοὺς ἐκ τῶν ἐθνῶν, καὶ συνάξω αὐτοὺς ἀπὸ τῶν χωρῶν, καὶ εἰσάξω αὐτοὺς εἰς τὴν γῆν αὐτῶν, καὶ βοσκήσω αὐτοὺς ἐπὶ τὰ ὄρη Ἰσραὴλ, καὶ ἐν ταῖς φάραγξι, καὶ ἐν πάσῃ κατοικίᾳ τῆς γῆς,
\vs{14}ἐν νομῇ ἀγαθῇ βοσκήσω αὐτοὺς, ἐν τῷ ὄρει τῷ ὑψηλῷ Ἰσραήλ· καὶ ἔσονται αἱ μάνδραι αὐτῶν ἐκεῖ, καὶ κοιμηθήσονται, καὶ ἐκεῖ ἀναπαύσονται ἐν τρυφῇ ἀγαθῇ, καὶ ἐν νομῇ πίονι βοσκηθήσονται ἐπὶ τῶν ὀρέων Ἰσραήλ.
\vs{15}Ἐγὼ βοσκήσω τὰ πρόβατά μου, καὶ ἐγὼ ἀναπαύσω αὐτὰ, καὶ γνώσονται ὅτι ἐγώ εἰμι Κύριος· τάδε λέγει Κύριος Κύριος.
\vs{16}Τὸ ἀπολωλὸς ζητήσω, καὶ τὸ πλανώμενον ἀποστρέψω, καὶ τὸ συντετριμμένον καταδήσω, καὶ τὸ ἐκλεῖπον ἐνισχύσω, καὶ τὸ ἰσχυρὸν φυλάξω, καὶ βοσκήσω αὐτὰ μετὰ κρίματος.

\vs{17}Καὶ ὑμεῖς πρόβατα, τάδε λέγει Κύριος Κύριος, ἰδοὺ ἐγὼ διακρινῶ ἀναμέσον προβάτου καὶ προβάτου, κριῶν καὶ τράγων.
\vs{18}Καὶ οὐχ ἱκανὸν ὑμῖν, ὅτι τὴν καλὴν νομὴν ἐνέμεσθε, καὶ τὰ κατάλοιπα τῆς νομῆς ὑμῶν κατεπατεῖτε τοῖς ποσὶν ὑμῶν; καὶ τὸ καθεστηκὸς ὕδωρ ἐπίνετε, καὶ τὸ λοιπὸν τοῖς ποσὶν ὑμῶν ἐταράσσετε;
\vs{19}Καὶ τὰ πρόβατά μου τὰ πατήματα τῶν ποδῶν ὑμῶν ἐνέμοντο, καὶ τὸ τεταραγμένον ὕδωρ ὑπὸ τῶν ποδῶν ὑμῶν ἔπινον.

\vs{20}Διατοῦτο τάδε λέγει Κύριος Κύριος, ἰδοὺ ἐγὼ διακρινῶ ἀναμέσον προβάτου ἰσχυροῦ καὶ ἀναμέσον προβάτου ἀσθενοῦς·
\vs{21}Ἐπὶ ταῖς πλευραῖς καὶ τοῖς ὤμοις ὑμῶν διωθεῖσθε, καὶ τοῖς κέρασιν ὑμῶν ἐκερατίζετε, καὶ πᾶν τὸ ἐκλεῖπον ἐξεθλίβετε.
\vs{22}Καὶ σώσω τὰ πρόβατά μου, καὶ οὐ μὴ ὦσιν ἔτι εἰς προνομὴν, καὶ κρινῶ ἀναμέσον κριοῦ πρὸς κριόν.

\vs{23}Καὶ ἀναστήσω ἐπʼ αὐτοὺς ποιμένα ἕνα, καὶ ποιμανεῖ αὐτοὺς, τὸν δοῦλόν μου Δαυὶδ, καὶ ἔσται αὐτῶν ποιμὴν,
\vs{24}καὶ ἐγὼ Κύριος ἔσομαι αὐτοῖς εἰς Θεὸν, καὶ Δαυὶδ ἄρχων ἐν μέσῳ αὐτῶν· ἐγὼ Κύριος ἐλάλησα.
\vs{25}Καὶ διαθήσομαι τῷ Δαυὶδ διαθήκην εἰρήνης, καὶ ἀφανιῶ θηρία πονηρὰ ἀπὸ τῆς γῆς, καὶ κατοικήσουσιν ἐν τῇ ἐρήμῳ, καὶ ὑπνώσουσιν ἐν τοῖς δρυμοῖς.
\vs{26}Καὶ δώσω αὐτοὺς περικύκλῳ τοῦ ὄρους μου· καὶ δώσω τὸν ὑετὸν ὑμῖν, ὑετὸν εὐλογίας.
\vs{27}Καὶ τὰ ξύλα τὰ ἐν τῷ πεδίῳ δώσει τὸν καρπὸν αὐτῶν, καὶ ἡ γῆ δώσει τὴν ἰσχὺν αὐτῆς, καὶ κατοικήσουσιν ἐπὶ τῆς γῆς αὐτῶν ἐν ἐλπίδι εἰρήνης, καὶ γνώσονται ὅτι ἐγώ εἰμι Κύριος, ἐν τῷ συντρίψαι με τὸν ζυγὸν αὐτῶν· καὶ ἐξελοῦμαι αὐτοὺς ἐκ χειρὸς τῶν καταδουλωσαμένων αὐτοὺς,
\vs{28}καὶ οὐκ ἔσονται ἔτι ἐν προνομῇ τοῖς ἔθνεσι, καὶ τὰ θηρία τῆς γῆς οὐκέτι μὴ φάγωσιν αὐτούς· καὶ κατοικήσουσιν ἐν ἐλπίδι, καὶ οὐκ ἔσται ὁ ἐκφοβῶν αὐτούς.
\vs{29}Καὶ ἀναστήσω αὐτοῖς φυτὸν εἰρήνης, καὶ οὐκέτι ἔσονται ἀπολλύμενοι λιμῷ ἐπὶ τῆς γῆς, καὶ ὀνειδισμὸν ἐθνῶν οὐ μὴ ἐνέγκωσιν ἔτι.

\vs{30}Καὶ γνώσονται ὅτι ἐγώ εἰμι Κύριος ὁ Θεὸς αὐτῶν, καὶ αὐτοὶ λαός μου. Οἶκος Ἰσραὴλ, λέγει Κύριος Κύριος,
\vs{31}πρόβατά μου καὶ πρόβατα ποιμνίου μου ἐστὲ, καὶ ἐγὼ Κύριος ὁ Θεὸς ὑμῶν, λέγει Κύριος Κύριος.

\ch{35}
Καὶ ἐγένετο λόγος Κυρίου πρὸς μὲ, λέγων,
\vs{2}υἱὲ ἀνθρώπου, ἐπίστρεψον τὸ πρόσωπόν σου ἐπʼ ὄρος Σηεὶρ, καὶ προφήτευσον εἰς αὐτὸ,
\vs{3}καὶ εἰπὸν αὐτῷ,

Τάδε λέγει Κύριος Κύριος, ἰδοὺ ἐγὼ ἐπὶ σὲ ὄρος Σηεὶρ, καὶ ἐκτενῶ τὴν χεῖρά μου ἐπὶ σὲ, καὶ δώσω σε εἰς ἔρημον, καὶ ἐρημωθήσῃ,
\vs{4}καὶ ταῖς πόλεσί σου ἐρημίαν ποιήσω, καὶ σὺ ἔρημος ἔσῃ, καὶ γνώσῃ ὅτι ἐγώ εἰμι Κύριος.
\vs{5}Ἀντὶ τοῦ γενέσθαι σε ἐχθρὰν αἰωνίαν, καὶ ἐνεκάθισας τῷ οἴκῳ Ἰσραὴλ δόλῳ, ἐν χειρὶ ἐχθρῶν μαχαίρᾳ ἐν καιρῷ ἀδικίας, ἐπʼ ἐσχάτων.

\vs{6}Διατοῦτο ζῶ ἐγὼ, λέγει Κύριος Κύριος, εἰ μὴν εἰς αἷμα ἥμαρτες, καὶ αἷμα διώξεταί σε.
\vs{7}Καὶ δώσω ὄρος Σηεὶρ εἰς ἔρημον, καὶ ἠρημωμένον, καὶ ἀπολῶ ἀπʼ αὐτοῦ ἀνθρώπους καὶ κτήνη,
\vs{8}καὶ ἐμπλήσω τῶν τραυματιῶν βουνούς σου καὶ τὰς φάραγγάς σου, καὶ ἐν πᾶσι τοῖς πεδίοις σου τετραυματισμένοι μαχαίρᾳ πεσοῦνται ἐν σοί.
\vs{9}Ἐρημίαν αἰώνιον θήσομαί σε, καὶ αἱ πόλεις σου οὐ μὴ κατοικηθῶσιν ἔτι, καὶ γνώσῃ ὅτι ἐγώ εἰμι Κύριος.

\vs{10}Διὰ τὸ εἰπεῖν σε, τὰ δύο ἔθνη, καὶ αἱ δύο χῶραι ἐμαὶ ἔσονται, καὶ κληρονομήσω αὐτὰς, καὶ Κύριος ἐκεῖ ἐστι·
\vs{11}Διατοῦτο ζῶ ἐγὼ, λέγει Κύριος, καὶ ποιήσω σοι κατὰ τὴν ἔχθραν σου, καὶ γνωσθήσομαί σοι ἡνίκα ἂν κρινῶ σε,
\vs{12}καὶ γνώσῃ ὅτι ἐγώ εἰμι Κύριος· ἤκουσα τῆς φωνῆς τῶν βλασφημιῶν σου, ὅτι εἶπας, τὰ ὄρη Ἰσραὴλ ἔρημα ἡμῖν δέδοται εἰς κατάβρωμα,
\vs{13}καὶ ἐμεγαλοῤῥημόνησας ἐπʼ ἐμὲ τῷ στόματί σου, ἐγὼ ἤκουσα.

\vs{14}Τάδε λέγει Κύριος, ἐν τῇ εὐφροσύνῃ πάσης τῆς γῆς ἔρημον ποιήσω σε.
\vs{15}Ἔρημον ἔσῃ ὄρος Σηεὶρ, καὶ πᾶσα ἡ Ἰδουμαία, καὶ ἐξαναλωθήσεται, καὶ γνώσῃ ὅτι ἐγώ εἰμι Κύριος ὁ Θεὸς αὐτῶν.

\ch{36}
Καὶ σὺ υἱὲ ἀνθρώπου προφήτευσον ἐπὶ τὰ ὄρη Ἰσραὴλ, καὶ εἰπὸν τοῖς ὄρεσι τοῦ Ἰσραὴλ, ἀκούσατε λόγον Κυρίου.

\vs{2}Τάδε λέγει Κύριος Κύριος, ἀνθʼ οὗ εἶπεν ἐφʼ ὑμᾶς ὁ ἐχθρὸς, εὖγε ἔρημα αἰώνια εἰς κατάσχεσιν ἡμῖν ἐγενήθη·
\vs{3}Διατοῦτο προφήτευσον, καὶ εἰπὸν, τάδε λέγει Κύριος Κύριος, ἀντὶ τοῦ ἀτιμασθῆναι ὑμᾶς, καὶ μισηθῆναι ὑμᾶς ὑπὸ τῶν κύκλῳ ὑμῶν, τοῦ εἶναι ὑμᾶς εἰς κατάσχεσιν τοῖς καταλοίποις ἔθνεσι, καὶ ἀνέβητε λάλημα γλώσσῃ, καὶ εἰς ὀνείδισμα ἔθνεσι,
\vs{4}διατοῦτο ὄρη Ἰσραὴλ ἀκούσατε λόγον Κυρίου· τάδε λέγει Κύριος τοῖς ὄρεσι, καὶ τοῖς βουνοῖς, καὶ τοῖς χειμάῤῥοις, καὶ ταῖς φάραγξι, καὶ τοῖς ἐξηρημωμένοις καὶ ἠφανισμένοις, καὶ ταῖς πόλεσι ταῖς ἐγκαταλελειμμέναις, καὶ ἐγένοντο εἰς προνομὴν, καὶ εἰς καταπάτημα τοῖς καταλειφθεῖσιν ἔθνεσι περικύκλῳ.
\vs{5}Διατοῦτο τάδε λέγει Κύριος Κύριος, εἰ μὴν ἐν πυρὶ θυμοῦ μου ἐλάλησα ἐπὶ τὰ λοιπὰ ἔθνη, καὶ ἐπὶ τὴν Ἰδουμαίαν πᾶσαν, ὅτι ἔδωκαν τὴν γῆν μου ἑαυτοῖς εἰς κατάσχεσιν μετʼ εὐφροσύνης, ἀτιμάσαντες ψυχὰς, τοῦ ἀφανίσαι ἐν προνομῇ·
\vs{6}Διατοῦτο προφήτευσον ἐπὶ τὴν γῆν τοῦ Ἰσραὴλ, καὶ εἰπὸν τοῖς ὄρεσι, καὶ τοῖς βουνοῖς, καὶ ταῖς φάραγξι, καὶ ταῖς νάπαις, τάδε λέγει Κύριος, ἰδοὺ ἐγὼ ἐν τῷ ζήλῳ μου, καὶ ἐν τῷ θυμῷ μου ἐλάλησα, ἀντὶ τοῦ ὀνειδισμοὺς ἐθνῶν ἐνέγκαι ὑμᾶς.
\vs{7}Διατοῦτο ἐγὼ ἀρῶ τὴν χεῖρά μου ἐπὶ τὰ ἔθνη τὰ περικύκλῳ ὑμῶν, οὗτοι τὴν ἀτιμίαν αὐτῶν λήψονται.

\vs{8}Ὑμῶν δὲ ὄρη Ἰσραὴλ τὴν σταφυλὴν καὶ τὸν καρπὸν ὑμῶν καταφάγεται ὁ λαός μου, ὅτι ἐλπίζουσι τοῦ ἐλθεῖν·
\vs{9}Ὅτι ἰδοὺ ἐγὼ ἐφʼ ὑμᾶς, καὶ ἐπιβλέψω ἐφʼ ὑμᾶς, καὶ κατεργασθήσεσθε, καὶ σπαρήσεσθε,
\vs{10}καὶ πληθυνῶ ἐφʼ ὑμᾶς ἀνθρώπους, πᾶν οἶκον Ἰσραὴλ εἰς τέλος, καὶ κατοικηθήσονται αἱ πόλεις, καὶ ἡ ἠρημωμένη οἰκοδομηθήσεται.
\vs{11}Καὶ πληθυνῶ ἐφʼ ὑμᾶς ἀνθρώπους, καὶ κτήνη, καὶ κατοικιῶ ὑμᾶς ὡς τὸ ἐν ἀρχῇ ὑμῶν, καὶ εὖ ποιήσω ὑμᾶς, ὥσπερ τὰ ἔμπροσθεν ὑμῶν, καὶ γνώσεσθε ὅτι ἐγώ εἰμι Κύριος.
\vs{12}Καὶ γεννήσω ἐφʼ ὑμᾶς ἀνθρώπους τὸν λαόν μου Ἰσραὴλ, καὶ κληρονομήσουσιν ὑμᾶς, καὶ ἔσεσθε αὐτοῖς εἰς κατάσχεσιν· καὶ οὐ μὴ προστεθῆτε ἔτι ἀτεκνωθῆναι ἀπʼ αὐτῶν.

\vs{13}Τάδε λέγει Κύριος Κύριος, ἀνθʼ ὧν εἶπάν σοι, κατέσθουσα ἀνθρώπους εἶ, καὶ ἠτεκνωμένη ὑπὸ τοῦ ἔθνους σου ἐγένου·
\vs{14}Διατοῦτο ἀνθρώπους οὐκέτι φάγεσαι, καὶ τὸ ἔθνος σου οὐκ ἀτεκνώσεις ἔτι, λέγει Κύριος Κύριος.
\vs{15}Καὶ οὐκ ἀκουσθήσεται οὐκέτι ἐφʼ ὑμᾶς ἀτιμία ἐθνῶν, καὶ ὀνειδισμοὺς λαῶν οὐ μὴ ἀνενέγκητε ἔτι, λέγει Κύριος Κύριος.

\vs{16}Καὶ ἐγένετο λόγος Κυρίου πρὸς μὲ, λέγων,
\vs{17}υἱὲ ἀνθρώπου, οἶκος Ἰσραὴλ κατῴκησεν ἐπὶ τῆς γῆς αὐτῶν, καὶ ἐμίαναν αὐτὴν ἐν τῇ ὁδῷ αὐτῶν, καὶ ἐν τοῖς εἰδώλοις αὐτῶν, καὶ ἐν ταῖς ἀκαθαρσίαις αὐτῶν, καὶ κατὰ τὴν ἀκαθαρσίαν τῆς ἀποκαθημένης ἐγενήθη ἡ ὁδὸς αὐτῶν πρὸ προσώπου μου.
\vs{18}Καὶ ἐξέχεα τὸν θυμόν μου ἐπʼ αὐτοὺς,
\vs{19}καὶ διέσπειρα αὐτοὺς εἰς τὰ ἔθνη, καὶ ἐλίκμησα αὐτοὺς εἰς τὰς χώρας, κατὰ τὴν ὁδὸν αὐτῶν, καὶ κατὰ τὴν ἁμαρτίαν αὐτῶν ἔκρινα αὐτούς.
\vs{20}Καὶ εἰσῆλθον εἰς τὰ ἔθνη, οὗ εἰσῆλθον ἐκεῖ, καὶ ἐβεβήλωσαν τὸ ὄνομά μου τὸ ἅγιον ἐν τῷ λέγεσθαι αὐτοὺς, λαὸς Κυρίου οὗτοι, καὶ ἐκ τῆς γῆς αὐτοῦ ἐξεληλύθασι.
\vs{21}Καὶ ἐπεισάμην αὐτῶν διὰ τὸ ὄνομά μου τὸ ἅγιον, ὃ ἐβεβήλωσαν οἶκος Ἰσραὴλ ἐν τοῖς ἔθνεσιν, οὗ εἰσήλθοσαν ἐκεῖ.

\vs{22}Διατοῦτο εἰπὸν τῷ οἴκῳ Ἰσραὴλ, τάδε λέγει Κύριος, οὐχ ὑμῖν ἐγὼ ποιῶ οἶκος Ἰσραὴλ, ἀλλʼ ἢ διὰ τὸ ὄνομά μου τὸ ἅγιον, ὃ ἐβεβηλώσατε ἐν τοῖς ἔθνεσιν, οὗ εἰσήλθετε ἐκεῖ.
\vs{23}Καὶ ἁγιάσω τὸ ὄνομά μου τὸ μέγα τὸ βεβηλωθὲν ἐν τοῖς ἔθνεσιν, ὃ ἐβεβηλώσατε ἐν μέσῳ αὐτῶν, καὶ γνώσονται τὰ ἔθνη ὅτι ἐγώ εἰμι Κύριος, ἐν τῷ ἁγιασθῆναί με ἐν ὑμῖν κατʼ ὀφθαλμοὺς αὐτῶν.

\vs{24}Καὶ λήψομαι ὑμᾶς ἐκ τῶν ἐθνῶν, καὶ ἀθροίσω ὑμᾶς ἐκ πασῶν τῶν γαιῶν, καὶ εἰσάξω ὑμᾶς εἰς τὴν γῆν ὑμῶν,
\vs{25}καὶ ῥανῶ ἐφʼ ὑμᾶς καθαρὸν ὕδωρ, καὶ καθαρισθήσεσθε ἀπὸ πασῶν τῶν ἀκαθαρσιῶν ὑμῶν, καὶ ἀπὸ πάντων τῶν εἰδώλων ὑμῶν, καὶ καθαριῶ ὑμᾶς·
\vs{26}Καὶ δώσω ὑμῖν καρδίαν καινὴν, καὶ πνεῦμα καινὸν δώσω ἐν ὑμῖν, καὶ ἀφελῶ τὴν καρδίαν τὴν λιθίνην ἐκ τῆς σαρκὸς ὑμῶν, καὶ δώσω ὑμῖν καρδίαν σαρκίνην·
\vs{27}Καὶ τὸ πνεῦμά μου δώσω ἐν ὑμῖν, καὶ ποιήσω ἵνα ἐν τοῖς δικαιώμασί μου πορεύησθε, καὶ τὰ κρίματά μου φυλάξησθε, καὶ ποιήσητε·
\vs{28}Καὶ κατοικήσετε ἐπὶ τῆς γῆς ἧς ἔδωκα τοῖς πατράσιν ὑμῶν, καὶ ἔσεσθέ μοι εἰς λαὸν, καὶ ἐγὼ ἔσομαι ὑμῖν εἰς Θεόν.
\vs{29}Καὶ σώσω ὑμᾶς ἐκ πασῶν τῶν ἀκαθαρσιῶν ὑμῶν, καὶ καλέσω τὸν σῖτον, καὶ πληθυνῶ αὐτὸν, καὶ οὐ δώσω ἐφʼ ὑμᾶς λιμόν.
\vs{30}Καὶ πληθυνῶ τὸν καρπὸν τοῦ ξύλου, καὶ τὰ γεννήματα τοῦ ἀγροῦ, ὅπως ἂν μὴ λάβητε ὀνειδισμὸν λιμοῦ ἐν τοῖς ἔθνεσι.

\vs{31}Καὶ μνησθήσεσθε τὰς ὁδοὺς ὑμῶν τὰς πονηρὰς, καὶ τὰ ἐπιτηδεύματα ὑμῶν τὰ μὴ ἀγαθὰ, καὶ προσοχθιεῖτε κατὰ πρόσωπον αὐτῶν ἐν ταῖς ἀνομίαις ὑμῶν, καὶ ἐπὶ τοῖς βδελύγμασιν αὐτῶν.
\vs{32}Οὐ διʼ ὑμᾶς ἐγὼ ποιῶ, λέγει Κύριος Κύριος, γνωστόν ἐστιν ὑμῖν· αἰσχύνθητε καὶ ἐντράπητε ἐκ τῶν ὁδῶν ὑμῶν οἶκος Ἰσραήλ.

\vs{33}Τάδε λέγει ἀδωναῒ Κύριος, ἐν ἡμέρᾳ ᾗ καθαριῶ ὑμᾶς ἐκ πασῶν ἀνομιῶν ὑμῶν, καὶ κατοικιῶ τὰς πόλεις, καὶ οἰκοδομηθήσονται ἔρημοι,
\vs{34}καὶ ἡ γῆ ἠφανισμένη ἐργασθήσεται ἀνθʼ ὧν ὅτι ἠφανισμένη ἐγενήθη κατʼ ὀφθαλμοὺς παντὸς παροδεύοντος.
\vs{35}Καὶ ἐροῦσιν, ἡ γῆ ἐκείνη ἠφανισμένη ἐγενήθη ὡς κῆπος τρυφῆς, καὶ αἱ πόλεις αἱ ἔρημοι καὶ ἠφανισμέναι καὶ κατεσκαμμέναι ὀχυραὶ ἐκάθισαν.
\vs{36}Καὶ γνώσονται τὰ ἔθνη, ὅσα ἂν καταλειφθῶσιν κύκλῳ ὑμῶν, ὅτι ἐγὼ Κύριος ᾠκοδόμησα τὰς καθῃρημένας, καὶ κατεφύτευσα τὰς ἠφανισμένας· ἐγὼ Κύριος ἐλάλησα, καὶ ποιήσω.

\vs{37}Τάδε λέγει ἀδωναῒ Κύριος, ἔτι τοῦτο ζητηθήσομαι τῷ οἴκῳ Ἰσραὴλ τοῦ ποιῆσαι αὐτούς· πληθυνῶ αὐτοὺς ὡς πρόβατα ἀνθρώπους,
\vs{38}ὡς πρόβατα ἅγια, ὡς πρόβατα Ἱερουσαλὴμ ἐν ταῖς ἑορταῖς αὐτῆς· οὕτως ἔσονται αἱ πόλεις αἱ ἔρημοι πλήρεις προβάτων ἀνθρώπων· καὶ γνώσονται ὅτι ἐγὼ Κύριος.

\ch{37}
Καὶ ἐγένετο ἐπʼ ἐμὲ χεὶρ Κυρίου, καὶ ἐξήγαγέ με ἐν πνεύματι Κύριος, καὶ ἔθηκέ με ἐν μέσῳ τοῦ πεδίου, καὶ τοῦτο ἦν μεστὸν ὀστέων ἀνθρωπίνων,
\vs{2}καὶ περιήγαγέ με ἐπʼ αὐτὰ κυκλόθεν κύκλῳ, καὶ ἰδοὺ πολλὰ σφόδρα ἐπὶ προσώπου τοῦ πεδίου, ξηρὰ σφόδρα.

\vs{3}Καὶ εἶπε πρὸς μὲ, υἱὲ ἀνθρώπου, εἰ ζήσεται τὰ ὀστέα ταῦτα; καὶ εἶπα, Κύριε Κύριε, σὺ ἐπίστῃ ταῦτα.
\vs{4}Καὶ εἶπε πρὸς μὲ, προφήτευσον ἐπὶ τὰ ὀστᾶ ταῦτα, καὶ ἐρεῖς αὐτοῖς, τὰ ὀστᾶ τὰ ξηρὰ, ἀκούσατε λόγον Κυρίου.
\vs{5}Τάδε λέγει Κύριος τοῖς ὀστέοις τούτοις, ἰδοὺ ἐγὼ φέρω ἐφʼ ὑμᾶς πνεῦμα ζωῆς,
\vs{6}καὶ δώσω ἐφʼ ὑμᾶς νεῦρα, καὶ ἀνάξω ἐφʼ ὑμᾶς σάρκα, καὶ ἐκτενῶ ἐφʼ ὑμᾶς δέρμα, καὶ δώσω πνεῦμά μου εἰς ὑμᾶς, καὶ ζήσεσθε, καὶ γνώσεσθε ὅτι ἐγώ εἰμι Κύριος.

\vs{7}Καὶ προεφήτευσα, καθὼς ἐνετείλατό μοι· καὶ ἐγένετο ἐν τῷ ἐμὲ προφητεῦσαι, καὶ ἰδοὺ σεισμὸς, καὶ προσήγαγε τὰ ὀστᾶ ἑκάτερον πρὸς τὴν ἁρμονίαν αὐτοῦ.
\vs{8}Καὶ ἴδον, καὶ ἰδοὺ, ἐπʼ αὐτὰ νεῦρα καὶ σάρκες ἐφύοντο, καὶ ἀνέβαινεν ἐπʼ αὐτὰ δέρματα ἐπάνω, καὶ πνεῦμα οὐκ ἦν ἐπʼ αὐτοῖς.
\vs{9}Καὶ εἶπε πρὸς μὲ, προφήτευσον ἐπὶ τὸ πνεῦμα, προφήτευσον υἱὲ ἀνθρώπου, καὶ εἰπὸν τῷ πνεύματι, τάδε λέγει Κύριος, ἐκ τῶν τεσσάρων πνευμάτων ἐλθὲ, καὶ ἐμφύσησον εἰς τοὺς νεκροὺς τούτους, καὶ ζησάτωσαν.
\vs{10}Καὶ προεφήτευσα καθότι ἐνετείλατό μοι, καὶ εἰσῆλθεν εἰς αὐτοὺς τὸ πνεῦμα, καὶ ἔζησαν, καὶ ἔστησαν ἐπὶ τῶν ποδῶν αὐτῶν, συναγωγὴ πολλὴ σφόδρα.

\vs{11}Καὶ ἐλάλησε Κύριος πρὸς μὲ, λέγων, υἱὲ ἀνθρώπου, τὰ ὀστᾶ ταῦτα πᾶς οἶκος Ἰσραήλ ἐστι, καὶ αὐτοὶ λέγουσι, ξηρὰ γέγονε τὰ ὀστᾶ ἡμῶν, ἀπόλωλεν ἡ ἐλπὶς ἡμῶν, διαπεφωνήκαμεν.
\vs{12}Διατοῦτο προφήτευσον, καὶ εἰπὸν,

Τάδε λέγει Κύριος, ἰδοὺ ἐγὼ ἀνοίγω τὰ μνήματα ὑμῶν, καὶ ἀνάξω ὑμᾶς ἐκ τῶν μνημάτων ὑμῶν, καὶ εἰσάξω ὑμᾶς εἰς τὴν γῆν τοῦ Ἰσραήλ·
\vs{13}Καὶ γνώσεσθε ὅτι ἐγώ εἰμι Κύριος, ἐν τῷ ἀνοῖξαί με τοὺς τάφους ὑμῶν, τοῦ ἀναγαγεῖν με ἐκ τῶν τάφων τὸν λαόν μου.
\vs{14}Καὶ δώσω πνεῦμά μου εἰς ὑμὰς, καὶ ζήσεσθε, καὶ θήσομαι ὑμᾶς ἐπὶ τὴν γῆν ὑμῶν, καὶ γνώσεσθε ὅτι ἐγὼ Κύριος· λελάληκα καὶ ποιήσω, λέγει Κύριος.

\vs{15}Καὶ ἐγένετο λόγος Κυρίου πρὸς μὲ λέγων,
\vs{16}υἱὲ ἀνθρώπου, λάβε σεαυτῷ ῥάβδον, καὶ γράψον ἐπʼ αὐτὴν τὸν Ἰούδαν, καὶ τοὺς υἱοὺς Ἰσραὴλ τοὺς προσκειμένους ἐπʼ αὐτόν· καὶ ῥάβδον δευτέραν λήψῃ σεαυτῷ, καὶ γράψεις αὐτὴν τῷ Ἰωσὴφ, ῥάβδον Ἐφραὶμ, καὶ πάντας τοὺς υἱοὺς Ἰσραὴλ τοὺς προστεθέντας πρὸς αὐτόν.
\vs{17}Καὶ συνάψεις αὐτὰς προσαλλήλας σεαυτῷ, εἰς ῥάβδον μίαν τοῦ δῆσαι ἑαυτὰς, καὶ ἔσονται ἐν τῇ χειρί σου.

\vs{18}Καὶ ἔσται ὅταν λέγωσι πρὸς σὲ οἱ υἱοὶ τοῦ λαοῦ σου, οὐκ ἀναγγέλλεις ἡμῖν, τί ἐστι ταῦτά σοι;
\vs{19}Καὶ ἐρεῖς πρὸς αὐτοὺς, τάδε λέγει Κύριος, ἰδοὺ ἐγὼ λήψομαι τὴν φυλὴν Ἰωσὴφ, τὴν διὰ χειρὸς Ἐφραὶμ, καὶ τὰς φυλὰς Ἰσραὴλ τὰς προσκειμένας πρὸς αὐτὸν, καὶ δώσω αὐτοὺς ἐπὶ τὴν φυλὴν Ἰούδα, καὶ ἔσονται εἰς ῥάβδον μίαν τῇ χειρὶ Ἰούδα.
\vs{20}Καὶ ἔσονται αἱ ῥάβδοι ἐφʼ αἷς σὺ ἔγραψας ἐπʼ αὐταῖς, ἐν τῇ χειρί σου ἐνώπιον αὐτῶν.
\vs{21}Καὶ ἐρεῖς αὐτοῖς,

Τάδε λέγει Κύριος Κύριος, ἰδοὺ ἐγὼ λαμβάνω πάντα οἶκον Ἰσραὴλ ἐκ μέσου τῶν ἐθνῶν, οὗ εἰσήλθοσαν ἐκεῖ, καὶ συνάξω αὐτοὺς ἀπὸ πάντων τῶν περικύκλῳ αὐτῶν, καὶ εἰσάξω αὐτοὺς εἰς τὴν γῆν τοῦ Ἰσραὴλ,
\vs{22}καὶ δώσω αὐτοὺς εἰς ἔθνος ἐν τῇ γῇ μου, καὶ ἐν τοῖς ὄρεσιν Ἰσραήλ· καὶ ἄρχων εἷς ἔσται αὐτῶν, καὶ οὐκ ἔσονται ἔτι εἰς δύο ἔθνη, οὐδὲ, μὴ διαιρεθῶσιν οὐκέτι εἰς δύο βασιλείας,
\vs{23}ἵνα μὴ μιαίνωνται ἔτι ἐν τοῖς εἰδώλοις αὐτῶν· καὶ ῥύσομαι αὐτοὺς ἀπὸ πασῶν τῶν ἀνομιῶν αὐτῶν, ὧν ἡμάρτοσαν ἐν αὐταῖς, καὶ καθαριῶ αὐτοὺς, καὶ ἔσονταί μοι εἰς λαὸν, καὶ ἐγὼ Κύριος ἔσομαι αὐτοῖς εἰς Θεόν.

\vs{24}Καὶ ὁ δοῦλός μου Δαυὶδ ἄρχων ἐν μέσῳ αὐτῶν, ἔσται ποιμὴν εἷς πάντων, ὅτι ἐν τοῖς προστάγμασί μου πορεύσονται, καὶ τὰ κρίματά μου φυλάξονται, καὶ ποιήσουνσιν αὐτά.
\vs{25}Καὶ κατοικήσουσιν ἐπὶ τῆς γῆς αὐτῶν, ἣν ἐγὼ δέδωκα τῷ δούλῳ μου Ἰακὼβ, οὗ κατῴκησαν ἐκεῖ οἱ πατέρες αὐτῶν, καὶ κατοικήσουσιν ἐπʼ αὐτῆς αὐτοί· καὶ Δαυὶδ ὁ δοῦλός μου ἄρχων εἰς τὸν αἰῶνα.

\vs{26}Καὶ διαθήσομαι αὐτοῖς διαθήκην εἰρήνης, διαθήκη αἰωνία ἔσται μετʼ αὐτῶν, καὶ θήσω τὰ ἅγιά μου ἐν μέσῳ αὐτῶν εἰς τὸν αἰῶνα,
\vs{27}καὶ ἔσται ἡ κατασκήνωσίς μου ἐν αὐτοῖς, καὶ ἔσομαι αὐτοῖς Θεὸς, καὶ αὐτοί μου ἔσονται λαός·
\vs{28}Καὶ γνώσονται τὰ ἔθνη ὅτι ἐγώ εἰμι Κύριος ὁ ἁγιάζων αὐτοὺς, ἐν τῷ εἶναι τὰ ἅγιά μου ἐν μέσῳ αὐτῶν εἰς τὸν αἰῶνα.

\ch{38}
Καὶ ἐγένετο λόγος Κυρίου πρὸς μὲ, λέγων,
\vs{2}υἱὲ ἀνθρώπου, στήρισον τὸ πρόσωπόν σου ἐπὶ Γὼγ, καὶ τὴν γῆν τοῦ Μαγὼγ, ἄρχοντα Ῥὼς Μεσὸχ, καὶ Θοβὲλ, καὶ προφήτευσον ἐπʼ αὐτόν.
\vs{3}Καὶ εἰπὸν αὐτῷ, τάδε λέγει Κύριος Κύριος,

\vs{4}Ἰδοὺ ἐγὼ ἐπὶ σὲ ἄρχοντα Ῥὼς Μεσὸχ, καὶ Θοβὲλ, καὶ συνάξω σε, καὶ πᾶσαν τὴν δύναμίν σου, ἵππους καὶ ἱππεῖς ἐνδεδυμένους θώρακας πάντας συναγωγῇ πολλῇ, πέλται καὶ περικεφαλαίαι καὶ μάχαιραι·
\vs{5}Πέρσαι, καὶ Αἰθίοπες, καὶ Λίβυες, πάντες περικεφαλαίαις καὶ πέλταις,
\vs{6}Γομὲρ, καὶ πάντες οἱ περὶ αὐτὸν, οἶκος τοῦ Θοργαμὰ, ἀπʼ ἐσχάτου Βοῤῥᾶ, καὶ πάντες οἱ περὶ αὐτὸν, καὶ ἔθνη πολλὰ μετὰ σοῦ.

\vs{7}Ἑτοιμάσθητι, ἑτοίμασον σεαυτὸν σὺ, καὶ πᾶσα ἡ συναγωγή σου ἡ συνηγμένη μετὰ σοῦ, καὶ ἔσῃ μοι εἰς προφυλακήν.
\vs{8}Ἀφʼ ἡμερῶν πλειόνων ἑτοιμασθήσεται, καὶ ἐπʼ ἐσχάτου ἐτῶν ἐλεύσεται, καὶ ἥξει εἰς τὴν γῆν τὴν ἀπεστραμμένην ἀπὸ μαχαίρας, συνηγμένων ἀπὸ ἐθνῶν πολλῶν ἐπὶ γῆν Ἰσραὴλ, ἣ ἐγενήθη ἔρημος διʼ ὅλου· καὶ οὗτος ἐξ ἐθνῶν ἐξελήλυθε, καὶ κατοικήσουσιν ἐπʼ εἰρήνης ἅπαντες.
\vs{9}Καὶ ἀναβήσῃ ὡς ὑετὸς, καὶ ἥξεις ὡς νεφέλη κατακαλύψαι γῆν, καὶ ἔσῃ σὺ καὶ πάντες οἱ περὶ σὲ, καὶ ἔθνη πολλὰ μετὰ σοῦ.

\vs{10}Τάδε λέγει Κύριος Κύριος, καὶ ἔσται ἐν τῇ ἡμέρᾳ ἐκείνῃ, ἀναβήσεται ῥήματα ἐπὶ τὴν καρδίαν σου, καὶ λογιῇ λογισμοὺς πονηροὺς,
\vs{11}καὶ ἐρεῖς, ἀναβήσομαι ἐπὶ γῆν ἀπεῤῥιμμένην, ἥξω ἐπὶ ἡσυχάζοντας ἐν τῇ ἡσυχίᾳ, καὶ οἰκοῦντας ἐπʼ εἰρήνης, πάντας κατοικοῦντας γῆν, ἐν ᾗ οὐχ ὑπάρχει τεῖχος, οὐδὲ μοχλοὶ, καὶ θύραι οὐκ εἰσὶν αὐτοῖς,
\vs{12}προνομεῦσαι προνομὴν, καὶ σκῦλα σκυλεῦσαι αὐτῶν, τοῦ ἐπιστρέψαι χεῖράς μου εἰς τὴν ἠρημωμένην ἣ κατῳκίσθη, καὶ ἐπʼ ἔθνος συνηγμένον ἀπὸ ἐθνῶν πολλῶν, πεποιηκότας κτήσεις, κατοικοῦντας ἐπὶ τὸν ὀμφαλὸν τῆς γῆς.
\vs{13}Σαββὰ, καὶ Δαιδὰν, καὶ ἔμποροι Καρχηδόνιοι, καὶ πᾶσαι αἱ κῶμαι αὐτῶν ἐροῦσί σοι, εἰς προνομὴν τοῦ προνομεῦσαι σὺ ἔρχῃ, καὶ σκυλεῦσαι σκῦλα, συνήγαγες συναγωγήν σου λαβεῖν ἀργύριον καὶ χρυσίον, ἀπενέγκασθαι κτῆσιν, τοῦ σκυλεῦσαι σκῦλα.

\vs{14}Διατοῦτο προφήτευσον υἱὲ ἀνθρώπου, καὶ εἰπὸν τῷ Γὼγ, τάδε λέγει Κύριος, οὐκ ἐν τῇ ἡμέρᾳ ἐκείνῃ, ἐν τῷ κατοικισθῆναι τὸν λαόν μου Ἰσραὴλ ἐπʼ εἰρήνης, ἐγερθήσῃ,
\vs{15}καὶ ἥξεις ἐκ τοῦ τόπου σου ἀπʼ ἐσχάτου Βοῤῥᾶ, καὶ ἔθνη πολλὰ μετὰ σοῦ; ἀναβάται ἵππων πάντες, συναγωγὴ μεγάλη καὶ δύναμις πολλή.
\vs{16}Καὶ ἀναβήσῃ ἐπὶ τὸν λαόν μου Ἰσραὴλ ὡς νεφέλη καλύψαι γῆν· ἐπʼ ἐσχάτων τῶν ἡμερῶν ἔσται, καὶ ἀνάξω σε ἐπὶ τὴν γῆν μου, ἵνα γνῶσι πάντα τὰ ἔθνη ἐμὲ, ἐν τῷ ἁγιασθῆναί με ἐν σοὶ ἐνώπιον αὐτῶν.

\vs{17}Τάδε λέγει Κύριος Κύριος τῷ Γὼγ, σὺ εἶ περὶ οὗ ἐλάλησα πρὸ ἡμερῶν τῶν ἔμπροσθεν, διὰ χειρὸς τῶν δούλων μου τῶν προφητῶν τοῦ Ἰσραὴλ, ἐν ταῖς ἡμέραις ἐκείναις καὶ ἔτεσι, τοῦ ἀναγαγεῖν σε ἐπʼ αὐτούς.
\vs{18}Καὶ ἔσται ἐν τῇ ἡμέρᾳ ἐκείνῃ, ἐν ἡμέρᾳ ᾗ ἂν ἔλθῃ Γὼγ ἐπὶ τὴν γῆν Ἰσραὴλ, λέγει Κύριος Κύριος, ἀναβήσεται ὁ θυμός μου,
\vs{19}καὶ ὁ ζῆλός μου· ἐν πυρὶ τῆς ὀργῆς μου ἐλάλησα, εἰ μὴν ἐν τῇ ἡμέρᾳ ἐκείνῃ ἔσται σεισμὸς μέγας ἐπὶ γῆς Ἰσραήλ·
\vs{20}Καὶ σεισθήσονται ἀπὸ προσώπου Κυρίου οἱ ἰχθύες τῆς θαλάσσης, καὶ τὰ πετεινὰ τοῦ οὐρανοῦ, καὶ τὰ θηρία τοῦ πεδίου, καὶ πάντα τὰ ἑρπετὰ τὰ ἕρποντα ἐπὶ τῆς γῆς, καὶ πάντες οἱ ἄνθρωποι οἱ ἐπὶ προσώπου τῆς γῆς, καὶ ῥαγήσεται τὰ ὄρη, καὶ πεσοῦνται αἱ φάραγγες, καὶ πᾶν τεῖχος ἐπὶ τὴν γῆν πεσεῖται.
\vs{21}Καὶ καλέσω ἐπʼ αὐτὸ καὶ πᾶν φόβον, λέγει Κύριος· μάχαιρα ἀνθρώπου ἐπὶ τὸν ἀδελφὸν αὐτοῦ ἔσται.
\vs{22}Καὶ κρινῶ αὐτὸν θανάτῳ, καὶ αἵματι, καὶ ὑετῷ κατακλύζοντι, καὶ λίθοις χαλάζης, καὶ πῦρ καὶ θεῖον βρέξω ἐπʼ αὐτὸν, καὶ ἐπὶ πάντας τοὺς μετʼ αὐτοῦ, καὶ ἐπʼ ἔθνη πολλὰ μετʼ αὐτοῦ.

\vs{23}Καὶ μεγαλυνθήσομαι, καὶ ἁγιασθήσομαι, καὶ ἐνδοξασθήσομαι, καὶ γνωσθήσομαι ἐναντίον ἐθνῶν πολλῶν, καὶ γνώσονται ὅτι ἐγώ εἰμι Κύριος.

\ch{39}
Καὶ σὺ υἱὲ ἀνθρώπου προφήτευσον ἐπὶ Γὼγ, καὶ εἰπὸν, τάδε λέγει Κύριος, ἰδοὺ ἐγὼ ἐπὶ σὲ Γὼγ ἄρχοντα Ῥὼς, Μεσὸχ, καὶ Θοβέλ.
\vs{2}Καὶ συνάξω σε, καὶ καθοδηγήσω σε, καὶ ἀναβιβῶ σε ἐπʼ ἐσχάτου τοῦ Βοῤῥᾶ, καὶ ἀνάξω σε ἐπὶ τὰ ὄρη τῷ Ἰσραήλ.
\vs{3}Καὶ ἀπολῶ τὸ τόκον σου ἀπὸ τῆς χειρός σου τῆς ἀριστερᾶς, καὶ τὰ τοξεύματά σου ἀπὸ τῆς χειρός σου τῆς δεξιᾶς, καὶ καταβαλῶ σε ἐπὶ τὰ ὄρη τὰ Ἰσραὴλ,
\vs{4}καὶ πεσῇ σὺ καὶ πάντες οἱ περὶ σὲ, καὶ τὰ ἔθνη τὰ μετὰ σοῦ δοθήσονται εἰς πλήθη ὀρνέων, παντὶ πετεινῷ, καὶ πᾶσι τοῖς θηρίοις τοῦ πεδίου δέδωκά σε καταβρωθῆναι.
\vs{5}Ἐπὶ προσώπου τοῦ πεδίου πεσῇ, ὅτι ἐγὼ ἐλάλησα, λέγει Κύριος.

\vs{6}Καὶ ἀποστελῶ πῦρ ἐπὶ Γὼγ, καὶ κατοικηθήσονται αἱ νῆσοι ἐπʼ εἰρήνης, καὶ γνώσονται ὅτι ἐγώ εἰμι Κύριος.
\vs{7}Καὶ τὸ ὄνομά μου τὸ ἅγιον γνωσθήσεται ἐν μέσῳ λαοῦ μου Ἰσραὴλ, καὶ οὐ βεβηλωθήσεται τὸ ὄνομά μου τὸ ἅγιον οὐκέτι, καὶ γνώσονται τὰ ἔθνη ὅτι ἐγώ εἰμι Κύριος, ἅγιος ἐν Ἰσραήλ.
\vs{8}Ἰδοὺ ἥκει, καὶ γνώσῃ ὅτι ἔσται, λέγει Κύριος Κύριος· αὕτη ἐστὶν ἡ ἡμέρα ἐν ᾗ ἐλάλησα.

\vs{9}Καὶ ἐξελεύσονται οἱ κατοικοῦντες τὰς πόλεις Ἰσραὴλ, καὶ καύσουσιν ἐν τοῖς ὅπλοις, πέλταις καὶ κοντοῖς, καὶ τόξοις καὶ τοξεύμασι, καὶ ῥάβδοις χειρῶν, καὶ λόγχαις, καὶ καύσουσιν ἐν αὐτοῖς πῦρ ἑπτὰ ἔτη·
\vs{10}Καὶ οὐ μὴ λάβωσι ξύλα ἐκ τοῦ πεδίου, οὐδὲ μὴ κόψωσιν ἐκ τῶν δρυμῶν, ἀλλʼ ἢ τὰ ὅπλα κατακαύσουσι πυρί· καὶ προνομεύσουσι τοὺς προνομεύσαντας αὐτοὺς, καὶ σκυλεύσουσι τοὺς σκυλεύσαντας αὐτοὺς, λέγει Κύριος.

\vs{11}Καὶ ἔσται, ἐν τῇ ἡμέρᾳ ἐκείνῃ δώσω τῷ Γὼγ τόπον ὀνομαστὸν, μνημεῖον ἐν Ἰσραὴλ, τὸ πολυάνδριον τῶν ἐπελθόντων πρὸς τῇ θαλάσσῃ καὶ περιοικοδομήσουσι τὸ περιστόμιον τῆς φάραγγος, καὶ κατορύξουσιν ἐκεῖ τὸν Γὼγ καὶ πᾶν τὸ πλῆθος αὐτοῦ, καὶ κληθήσεται τότε Τὸ πολυάνδριον τοῦ Γώγ.
\vs{12}Καὶ κατορύξουσιν αὐτοὺς οἶκος Ἰσραὴλ, ἵνα καθαρισθῇ ἡ γῆ ἐν ἑπταμήνῳ.
\vs{13}Καὶ κατορύξουσιν αὐτοὺς πᾶς ὁ λαὸς τῆς γῆς, καὶ ἔσται αὐτοῖς ὀνομαστὸν, ᾗ ἡμέρᾳ ἐδοξάσθη, λέγει Κύριος.
\vs{14}Καὶ ἄνδρας διαπαντὸς διαστελοῦσιν ἐπιπορευομένους τὴν γῆν, θάψαι τοὺς καταλελειμμένους ἐπὶ προσώπου τῆς γῆς, καθαρίσαι αὐτὴν μετὰ τὴν ἐπτάμηνον, καὶ ἐκζητήσουσι.
\vs{15}Καὶ πᾶς ὁ διαπορευόμενος τὴν γῆν, καὶ ἰδὼν ὀστοῦν ἀνθρώπου, οἰκοδομήσει παρʼ αὐτῷ σημεῖον, ἕως ὅτου θάψωσιν αὐτὸ οἱ θάπτοντες εἰς γαὶ τὸ πολυάνδριον τοῦ Γώγ.
\vs{16}Καὶ γὰρ τὸ ὄνομα τῆς πόλεως, Πολυάνδριον· καὶ καθαρισθήσεται ἡ γῆ.

\vs{17}Καὶ σὺ υἱὲ ἀνθρώπου εἰπὸν, τάδε λέγει Κύριος, εἰπὸν παντὶ ὀρνέῳ πετεινῷ, καὶ πρὸς πάντα τὰ θηρία τοῦ πεδίου,

Συνάχθητε καὶ ἔρχεσθε, συνάχθητε ἀπὸ πάντων τῶν περικύκλῳ ἐπὶ τὴν θυσίαν μου, ἣν τέθυκα ὑμῖν θυσίαν μεγάλην ἐπὶ τὰ ὄρη Ἰσραὴλ, καὶ φάγεσθε κρέα, καὶ πίεσθε αἷμα.
\vs{18}Κρέα γιγάντων φάγεσθε, καὶ αἷμα ἀρχόντων τῆς γῆς πίεσθε· κριοὺς καὶ μόσχους καὶ τράγους, καὶ οἱ μόσχοι ἐστεατωμένοι πάντες.
\vs{19}Καὶ φάγεσθε στέαρ εἰς πλησμονὴν, καὶ πίεσθε αἷμα εἰς μέθην ἀπὸ τῆς θυσίας μου, ἧς ἔθυσα ὑμῖν.
\vs{20}Καὶ ἐμπλησθήσεσθε ἐπὶ τῆς τραπέζης μου, ἵππον καὶ ἀναβάτην καὶ γίγαντα καὶ πάντα ἄνδρα πολεμιστὴν, λέγει Κύριος.

\vs{21}Καὶ δώσω τὴν δόξαν μου ἐν ὑμῖν, καὶ ὄψονται πάντα τὰ ἔθνη τὴν κρίσιν μου ἣν ἐποίησα, καὶ τὴν χεῖρά μου ἣν ἐπήγαγον ἐπʼ αὐτούς.
\vs{22}Καὶ γνώσονται οἶκος Ἰσραὴλ, ὅτι ἐγώ εἰμι Κύριος ὁ Θεὸς αὐτῶν, ἀπὸ τῆς ἡμέρας ταύτης καὶ ἐπέκεινα.

\vs{23}Καὶ γνώσονται πάντα τὰ ἔθνη, ὅτι διὰ τὰς ἁμαρτίας αὐτῶν ᾐχμαλωτεύθησαν οἶκος Ἰσραὴλ, ἀνθʼ ὧν ἠθέτησαν εἰς ἐμὲ, καὶ ἀπέστρεψα τὸ πρόσωπόν μου ἀπʼ αὐτῶν, καὶ παρέδωκα αὐτοὺς εἰς χεῖρας τῶν ἐχθρῶν αὐτῶν, καὶ ἔπεσαν πάντες μαχαίρᾳ.
\vs{24}Κατὰ τὰς ἀκαθαρσίας αὐτῶν, καὶ κατὰ τὰ ἀνομήματα αὐτῶν ἐποίησα αὐτοῖς, καὶ ἀπέστρεψα τὸ πρόσωπόν μου ἀπʼ αὐτῶν.

\vs{25}Διατοῦτο τάδε λέγει Κύριος Κύριος, νῦν ἀποστρέψω αἰχμαλωσίαν ἐν Ἰακὼβ, καὶ ἐλεήσω τὸν οἶκον Ἰσραὴλ, καὶ ζηλώσω διὰ τὸ ὄνομα τὸ ἅγιόν μου.
\vs{26}Καὶ λήψονται τὴν ἀτιμίαν αὐτῶν, καὶ τὴν ἀδικίαν ἣν ἠδίκησαν ἐν τῷ κατοικισθῆναι αὐτοὺς ἐπὶ τὴν γῆν αὐτῶν ἐπʼ εἰρήνης· καὶ οὐκ ἔσται ὁ ἐκφοβῶν,
\vs{27}ἐν τῷ ἀποστρέψαι με αὐτοὺς ἐκ τῶν ἐθνῶν, καὶ συναγαγεῖν με αὐτοὺς ἐκ τῶν χωρῶν τῶν ἐθνῶν· καὶ ἁγιασθήσομαι ἐν αὐτοῖς ἐνώπιον τῶν ἐθνῶν.
\vs{28}Καὶ γνώσονται ὅτι ἐγώ εἰμι Κύριος ὁ Θεὸς αὐτῶν, ἐν τῷ ἐπιφανῆναί με αὐτοῖς ἐν τοῖς ἔθνεσι.
\vs{29}Καὶ οὐκ ἀποστρέψω οὐκέτι τὸ πρόσωπόν μου ἀπʼ αὐτῶν, ἀνθʼ ὧν ἐξέχεα τὸν θυμόν μου ἐπὶ τὸν οἶκον Ἰσραὴλ, λέγει Κύριος Κύριος.

\ch{40}
Καὶ ἐγένετο ἐν τῷ πέμπτῳ καὶ εἰκοστῷ ἔτει τῆς αἰχμαλωσίας ἡμῶν, ἐν τῷ πρώτῳ μηνὶ, δεκάτῃ τοῦ μηνὸς, ἐν τῷ τεσσαρεσκαιδεκάτῳ ἔτει μετὰ τὸ ἁλῶναι τὴν πόλιν, ἐν τῇ ἡμέρᾳ ἐκείνῃ ἐγένετο ἐπʼ ἐμὲ χεὶρ Κυρίου,
\vs{2}καὶ ἤγαγέ με ἐν ὁράσει Θεοῦ εἰς τὴν γῆν Ἰσραὴλ, καὶ ἔθηκέ με ἐπʼ ὄρος ὑψηλὸν σφόδρα, καὶ ἐπʼ αὐτῷ ὡσεὶ οἰκοδομὴ πόλεως ἀπέναντι,

\vs{3}Καὶ εἰσήγαγέ με ἐκεῖ· καὶ ἰδοὺ ἀνὴρ, καὶ ἡ ὅρασις αὐτοῦ ἦν ὡσεὶ ὅρασις χαλκοῦ στίλβοντος, καὶ ἐν τῇ χειρὶ αὐτοῦ ἦν σπαρτίον οἰκοδόμων, καὶ κάλαμος μέτρον, καὶ αὐτὸς εἱστήκει ἐπὶ τῆς πύλης.
\vs{4}Καὶ εἶπε πρὸς μὲ ὁ ἀνὴρ,

Ὃν ἑώρακας υἱὲ ἀνθρώπου ἐν τοῖς ὀφθαλμοῖς σου ἴδε, καὶ ἐν τοῖς ὠσί σου ἄκουε, καὶ τάξον εἰς τὴν καρδίαν σου πάντα ὅσα ἐγὼ δεικνύω σοι, διότι ἕνεκα τοῦ δεῖξαί σοι εἰσελήλυθας ὧδε, καὶ δείξεις πάντα ὅσα σὺ ὁρᾷς τῷ οἴκῳ τοῦ Ἰσραήλ.

\vs{5}Καὶ ἰδοὺ περίβολος ἔξωθεν τοῦ οἴκου κύκλῳ, καὶ ἐν τῇ χειρὶ τοῦ ἀνδρὸς κάλαμος, τὸ μέτρον πηχῶν ἓξ ἐν πήχει καὶ παλαιστῆς· καὶ διεμέτρησε τὸ προτείχισμα, πλάτος ἶσον τῷ καλάμῳ, καὶ τὸ ὕψος αὐτοῦ ἶσον τῷ καλάμῳ.

\vs{6}Καὶ εἰσῆλθεν εἰς τὴν πύλην τὴν βλέπουσαν κατὰ ἀνατολὰς ἐν ἑπτὰ ἀναβαθμοῖς, καὶ διεμέτρησε τὸ αἰλὰμ τῆς πύλης ἶσον τῷ καλάμῳ·
\vs{7}Καὶ τὸ θεὲ ἶσον τῷ καλάμῳ τὸ μῆκος, καὶ ἶσον τῷ καλάμῳ τὸ πλάτος, καὶ τὸ αἰλὰμ ἀναμέσον τοῦ θεηλὰθ πηχῶν ἕξ· καὶ τὸ θεὲ τὸ δεύτερον, ἶσον τῷ καλάμῳ πλάτος, καὶ ἶσον τῷ καλάμῳ μῆκος, καὶ τὸ αἰλὰμ πηχέων πέντε·
\vs{8}Καὶ τὸ θεὲ τὸ τρίτον, ἶσον τῷ καλάμῳ μῆκος, καὶ ἶσον τῷ καλάμῳ πλάτος.
\vs{9}Καὶ τὸ αἰλὰμ τοῦ πυλῶνος, πλησίον τοῦ αἰλὰμ τῆς πύλης πηχῶν ὀκτὼ, καὶ τὰ αἰλεῦ πηχῶν δύο· καὶ τὰ αἰλὰμ τῆς πύλης ἔσωθεν,
\vs{10}καὶ τὰ θεὲ τῆς πύλης τοῦ θεὲ κατέναντι, τρεῖς ἔνθεν καὶ τρεῖς ἔνθεν, καὶ μέτρον ἓν τοῖς τρισί· μέτρον ἓν τοῖς αἰλὰμ ἔνθεν καὶ ἔνθεν.
\vs{11}Καὶ διεμέτρησε τὸ πλάτος τῆς θύρας τοῦ πυλῶνος πηχῶν δέκα, καὶ τὸ εὖρος τοῦ πυλῶνος πηχῶν δεκατριῶν.
\vs{12}Καὶ πῆχυς ἐπισυναγόμενος ἐπὶ πρόσωπον τῶν θεεὶμ ἔνθεν καὶ ἔνθεν, καὶ τὸ θεὲ πηχῶν ἓξ ἔνθεν καὶ πηχῶν ἓξ ἔνθεν.

\vs{13}Καὶ διεμέτρησε τὴν πύλην ἀπὸ τοῦ τοίχου τοῦ θεὲ ἐπὶ τὸν τοῖχον τοῦ θεὲ, πλάτος πήχεις εἴκοσι καὶ πέντε· αὕτη πύλη ἐπὶ πύλην.
\vs{14}Καὶ τὸ αἴθριον τοῦ αἰλὰμ τῆς πύλης ἔξωθεν, πήχεις εἴκοσι θεεὶμ τῆς πύλης κύκλῳ.
\vs{15}Καὶ τὸ αἴθριον τῆς πύλης ἔξωθεν, εἰς τὸ αἴθριον αἰλὰμ τῆς πύλης ἔσωθεν πηχῶν πεντήκοντα.

\vs{16}Καὶ θυρίδες κρυπταὶ ἐπὶ τὸ θεεὶμ, καὶ ἐπὶ τὰ αἰλὰμ ἔσωθεν τῆς πύλης τῆς αὐλῆς κυκλόθεν· καὶ ὡσαύτως τοῖς αἰλὰμ θυρίδας κύκλῳ ἔσωθεν, καὶ ἐπὶ τὸ αἰλὰμ φοίνικες ἔνθεν καὶ ἔνθεν.

\vs{17}Καὶ εἰσήγαγέ με εἰς τὴν αὐλὴν τὴν ἐσωτέραν, καὶ ἰδοὺ παστοφόρια, καὶ περίστυλα κύκλῳ τῆς αὐλῆς, τριάκοντα παστοφόρια ἐν τοῖς περιστύλοις,
\vs{18}καὶ αἱ στοαὶ κατὰ νώτου τῶν πυλῶν κατὰ τὸ μῆκος τῶν πυλῶν τὸ περίστυλον τὸ ὑποκάτω.
\vs{19}Καὶ διεμέτρησε τὸ πλάτος τῆς αὐλῆς, ἀπὸ τοῦ αἰθρίου τῆς πύλης τῆς ἐξωτέρας ἔσωθεν ἐπὶ τὸ αἴθριον τῆς πύλης τῆς βλεπούσης ἔξω, πήχεις ἑκατὸν τῆς βλεπούσης κατὰ ἀνατολάς· καὶ ἤγαγέ με ἐπὶ Βοῤῥᾶν,
\vs{20}καὶ ἰδοὺ πύλη βλέπουσα πρὸς Βοῤῥᾶν τῇ αὐλῇ τῇ ἐξωτέρᾳ, καὶ διεμέτρησεν αὐτὴν, τό, τε μῆκος αὐτῆς καὶ τὸ πλάτος,
\vs{21}καὶ τὸ θεὲ τρεῖς ἔνθεν καὶ τρεῖς ἔνθεν, καὶ τὸ αἰλεῦ, καὶ τὰ αἰλαμμὼν, καὶ τοὺς φοίνικας αὐτῆς· καὶ ἐγένετο κατὰ τὰ μέτρα τῆς πύλης τῆς βλεπούσης κατὰ ἀνατολὰς, πηχῶν πεντήκοντα τὸ μῆκος αὐτῆς, καὶ πηχῶν εἰκοσιπέντε τὸ εὖρος αὐτῆς·
\vs{22}Καὶ αἱ θυρίδες αὐτῆς, καὶ τὰ αἰλαμμὼν, καὶ οἱ φοίνικες αὐτῆς καθὼς ἡ πύλη ἡ βλέπουσα κατὰ ἀνατολάς· καὶ ἐν ἑπτὰ κλημακτῆρσιν ἀνέβαινον ἐπʼ αὐτὸν, καὶ τὰ αἰλαμμὼν ἔσωθεν·
\vs{23}Καὶ πύλη τῇ αὐλῇ τῇ ἐσωτέρᾳ βλέπουσα ἐπὶ πύλην τοῦ Βοῤῥᾶ, ὃν τρόπον τῆς πύλης τῆς βλεπούσης κατὰ ἀνατολάς· καὶ διεμέτρησε τὴν αὐλὴν ἀπὸ πύλης ἐπὶ πύλην, πήχεις ἑκατόν.

\vs{24}Καὶ ἤγαγέ με κατὰ Νότον, καὶ ἰδοὺ πύλη βλέπουσα πρὸς Νότον, καὶ διεμέτρησεν αὐτὴν, καὶ τὰ θεὲ, καὶ τὰ αἰλεῦ, καὶ τὰ αἰλαμμὼν, κατὰ τὰ μέτρα ταῦτα.
\vs{25}Καὶ αἱ θυρίδες αὐτῆς, καὶ τὰ αἰλαμμὼν κυκλόθεν, καθὼς αἱ θυρίδες τοῦ αἰλὰμ, πηχῶν πεντήκοντα τὸ μῆκος αὐτῆς, καὶ πηχῶν εἰκοσιπέντε τὸ εὖρος αὐτῆς,
\vs{26}καὶ ἑπτὰ κλημακτῆρες αὐτῇ, καὶ αἰλαμμὼν ἔσωθεν, καὶ φοίνικες αὐτῇ, εἷς ἔνθεν καὶ εἷς ἔνθεν ἐπὶ τὰ αἰλεῦ.
\vs{27}Καὶ πύλη κατέναντι τῆς πύλης τῆς αὐλῆς τῆς ἐσωτέρας πρὸς Νότον, καὶ διεμέτρησε τὴν αὐλὴν ἀπὸ πύλης ἐπὶ πύλην, πήχεις ἑκατὸν τὸ εὖρος πρὸς Νότον.

\vs{28}Καὶ εἰσήγαγέ με εἰς τὴν αὐλὴν τὴν ἐσωτέραν τῆς πύλης τῆς πρὸς Νότον, καὶ διεμέτρησε τὴν πύλην κατὰ τὰ μέτρα ταῦτα,
\vs{29}καὶ τὰ θεὲ, καὶ τὰ αἰλεῦ, καὶ τὰ αἰλαμμὼν κατὰ τὰ μέτρα ταῦτα, καὶ θυρίδες αὐτῇ, καὶ τῷ αἰλαμμὼν κύκλῳ, πήχεις πεντήκοντα τὸ μῆκος αὐτῆς,
\vs{31}καὶ τὸ εὖρος πήχεις εἰκοσιπέντε τοῦ αἰλὰμ εἰς τὴν αὐλὴν τὴν ἐξωτέραν, καὶ φοίνικες τῷ αἰλεῦ, καὶ ὀκτὼ κλημακτῆρες.

\vs{32}Καὶ εἰσήγαγέ με εἰς τὴν πύλην τὴν βλέπουσαν κατὰ ἀνατολὰς, καὶ διεμέτρησεν αὐτὴν κατὰ τὰ μέτρα ταῦτα,
\vs{33}καὶ τὰ θεὲ, καὶ τὰ αἰλεῦ, καὶ τὰ αἰλαμμὼν κατὰ τὰ μέτρα ταῦτα, καὶ θυρίδες αὐτῇ καὶ αἰλαμμὼν κύκλῳ πήχεις πεντήκοντα μῆκος αὐτῆς, καὶ εὖρος αὐτῆς πήχεις εἰκοσιπέντε.
\vs{34}Καὶ αἰλαμμὼν εἰς τὴν αὐλὴν τὴν ἐσωτέραν, καὶ φοίνικες ἐπὶ τοῦ αἰλεῦ ἔνθεν καὶ ἔνθεν, καὶ ὀκτὼ κλημακτῆρες αὐτῇ.

\vs{35}Καὶ εἰσήγαγέ με εἰς τὴν πύλην τὴν πρὸς βοῤῥᾶν, καὶ διεμέτρησε κατὰ τὰ μέτρα ταῦτα,
\vs{36}καὶ τὰ θεὲ, καὶ τὰ αἰλεῦ, καὶ τὰ αἰλαμμὼν, καὶ θυρίδες αὐτῇ κύκλῳ, καὶ τὸ αἰλαμμὼν αὐτῆς, πήχεις πεντήκοντα μῆκος αὐτῆς, καὶ εὖρος πήχεις εἰκοσιπέντε·
\vs{37}Καὶ τὰ αἰλαμμὼν εἰς τὴν αὐλὴν τὴν ἐξωτέραν, καὶ φοίνικες τῷ αἰλεῦ ἔνθεν καὶ ἔνθεν, καὶ ὀκτὼ κλημακτῆρες αὐτῇ.

\vs{38}Τὰ παστοφόρια αὐτῆς, καὶ τὰ θυρώματα αὐτῆς, καὶ τὰ αἰλαμμὼν αὐτῆς ἐπὶ τῆς πύλης τῆς δευτέρας ἔκρυσις,
\vs{39}ὅπως σφάζωσιν ἐν αὐτῇ τὰ ὑπὲρ ἁμαρτίας, καὶ τὰ ὑπὲρ ἀγνοίας.
\vs{40}Καὶ κατὰ νώτου τοῦ ῥύακος τῶν ὁλοκαυτωμάτων τῆς βλεπούσης πρὸς Βοῤῥᾶν, δύο τράπεζαι πρὸς ἀνατολὰς κατὰ νώτου τῆς δευτέρας, καὶ τοῦ αἰλὰμ τῆς πύλης δύο τράπεζαι κατὰ ἀνατολάς.
\vs{41}Τέσσαρες ἔνθεν, καὶ τέσσαρες ἔνθεν κατὰ νώτου τῆς πύλης, ἐπʼ αὐτὰς σφάζουσι τὰ θύματα· κατέναντι τῶν ὀκτὼ τραπεζῶν τῶν θυμάτων.
\vs{42}Καὶ τέσσαρες τράπεζαι τῶν ὁλοκαυτωμάτων λίθιναι λελαξευμέναι, πήχεως καὶ ἡμίσους τὸ πλάτος, καὶ πήχεων δύο ἡμίσους τὸ μῆκος, καὶ ἐπὶ πῆχυν τὸ ὕψος· ἐπʼ αὐτὰ ἐπιθήσουσι τὰ σκεύη, ἑν οἷς σφάζουσιν ἐκεῖ τὰ ὁλοκαυτώματα καὶ τὰ θύματα.
\vs{43}Καὶ παλαιστὴν ἕξουσι γεῖσος λελαξευμένον ἔσωθεν κύκλῳ, καὶ ἐπὶ τὰς τραπέζας ἐπάνωθεν στέγας, τοῦ καλύπτεσθαι ἀπὸ τοῦ ὑετοῦ, καὶ ἀπὸ τῆς ξηρασίας.

\vs{44}Καὶ εἰσήγαγέ με εἰς τὴν αὐλὴν τὴν ἐσωτέραν· καὶ ἰδοὺ δύο ἐξέδραι ἐν τῇ αὐλῇ τῇ ἐσωτέρᾳ, μία κατὰ νώτου τῆς πύλης τῆς βλεπούσης πρὸς Βοῤῥᾶν, φέρουσα πρὸς Νότον, καὶ μία κατὰ νώτου τῆς πύλης τῆς πρὸς Νότον, βλεπούσης δὲ πρὸς Βοῤῥᾶν.
\vs{45}Καὶ εἶπε πρὸς μὲ, ἡ ἐξέδρα αὕτη ἡ βλέπουσα πρὸς Νότον, τοῖς ἱερεῦσι τοῖς φυλάσσουσι τὴν φυλακὴν τοῦ οἴκου,
\vs{46}καὶ ἡ ἐξέδρα ἡ βλέπουσα πρὸς Βοῤῥᾶν, τοῖς ἱερεῦσι τοῖς φυλάσσουσι τὴν φυλακὴν τοῦ θυσιαστηρίου· ἐκεῖνοί εἰσιν οἱ υἱοὶ Σαδδοὺκ, οἱ ἐγγίζοντες ἐκ τοῦ Λευὶ πρὸς Κύριον, λειτουργεῖν αὐτῷ.

\vs{47}Καὶ διεμέτρησε τὴν αὐλὴν, μῆκος πηχῶν ἑκατὸν, καὶ εὖρος πήχεις ἑκατὸν, ἐπὶ τὰ τέσσαρα μέρη αὐτῆς, καὶ τὸ θυσιαστήριον ἀπέναντι τοῦ οἴκου.
\vs{48}Καὶ εἰσήγαγέ με εἰς τὸ αἰλὰμ τοῦ οἴκου· καὶ διεμέτρησε τὸ αἲλ τοῦ αἰλὰμ πηχῶν πέντε τὸ πλάτος ἔνθεν, καὶ πηχῶν πέντε ἔνθεν, καὶ τὸ εὖρος τοῦ θυρώματος πηχῶν δεκατεσσάρων, καὶ ἐπωμίδες τῆς θύρας τοῦ αἰλὰμ πηχῶν τριῶν ἔνθεν, καὶ πηχῶν τριῶν ἔνθεν.
\vs{49}Καὶ τὸ μῆκος τοῦ αἰλὰμ πηχῶν εἴκοσι, καὶ τὸ εὖρος πηχῶν δώδεκα· καὶ ἐπὶ δέκα ἀναβαθμῶν ἀνέβαινον ἐπʼ αὐτὸ, καὶ στύλοι ἦσαν ἐπὶ τὸ αἰλάμ, εἷς ἔνθεν καὶ εἷς ἔνθεν.

\ch{41}
Καὶ εἰσήγαγέ με εἰς τὸν ναὸν, ᾧ διεμέτρησε τὸ αἰλὰμ πηχῶν ἓξ τὸ πλάτος ἔνθεν,
\vs{2}καὶ πηχῶν ἓξ τὸ εὖρος τοῦ αἰλὰμ ἔνθεν, καὶ τὸ εὖρος τοῦ πυλῶνος πηχῶν δέκα, καὶ ἐπωμίδες τοῦ πυλῶνος πηχῶν πέντε ἔνθεν, καὶ πηχῶν πέντε ἔνθεν· καὶ διεμέτρησε τὸ μῆκος αὐτοῦ πηχῶν τεσσαράκοντα, καὶ τὸ εὖρος πηχῶν εἴκοσι.

\vs{3}Καὶ εἰσῆλθεν εἰς τὴν αὐλὴν τὴν ἐσωτέραν, καὶ διεμέτρησε τὸ αἲλ τοῦ θυρώματος πηχῶν δύο, καὶ τὸ θύρωμα πηχῶν ἓξ, καὶ τὰς ἐπωμίδας τοῦ θυρώματος πηχῶν ἑπτὰ ἔνθεν, καὶ πηχῶν ἑπτὰ ἔνθεν.
\vs{4}Καὶ διεμέτρησε τὸ μῆκος τῶν θυρῶν πηχῶν τεσσαράκοντα, καὶ τὸ εὖρος πηχῶν εἴκοσι, κατὰ πρόσωπον τοῦ ναοῦ· καὶ εἶπε, τοῦτο τὸ ἅγιον τῶν ἁγίων.

\vs{5}Καὶ διεμέτρησε τὸν τοῖχον τοῦ οἴκου πηχῶν ἓξ, καὶ τὸ εὖρος τῆς πλευρᾶς πηχῶν τεσσάρων κυκλόθεν,
\vs{6}καὶ πλευρὰ πλευρὸν ἐπὶ πλευρὸν τριάκοντα τρὶς δίς· καὶ διάστημα ἐν τῷ τοίχῳ τοῦ οἴκου ἐν τοῖς πλευροῖς κύκλῳ, τοῦ εἶναι τοῖς ἐπιλαμβανομένοις ὁρᾷν, ὅπως τὸ παράπαν μὴ ἅπτωνται τῶν τοίχων τοῦ οἴκου.
\vs{7}Καὶ τὸ εὖρος τῆς ἀνωτέρας τῶν πλευρῶν κατὰ τὸ πρόσθεμα ἐκ τοῦ τοίχου, πρὸς τὴν ἀνωτέραν κύκλῳ τοῦ οἴκου, ὅπως διαπλατύνηται ἄνωθεν, καὶ ἐκ τῶν κάτωθεν ἀναβαίνωσιν ἐπὶ τὰ ὑπερῷα, καὶ ἐκ τῶν γεισῶν ἐπὶ τὰ τριώροφα,

\vs{8}Καὶ τὸ θραὲλ τοῦ οἴκου ὕψος κύκλῳ διάστημα τῶν πλευρῶν ἶσον τῷ καλάμῳ πηχῶν ἕξ· διαστήματα
\vs{9}καὶ εὖρος τοῦ τοίχου τῆς πλευρᾶς ἔξωθεν πηχῶν πέντε, καὶ τὰ ἀπόλοιπα τὰ ἀναμέσον τῶν πλευρῶν τοῦ οἴκου,
\vs{10}καὶ ἀναμέσον τῶν ἐξεδρῶν εὖρος πηχῶν εἴκοσι, τὸ περιφερὲς τῷ οἴκῳ κύκλῳ.

\vs{11}Καὶ αἱ θύραι τῶν ἐξεδρῶν ἐπὶ τὸ ἀπόλοιπον τῆς θύρας τῆς μιᾶς τῆς πρὸς Βοῤῥᾶν· καὶ ἡ θύρα ἡ μία πρὸς Νότον· καὶ τὸ εὖρος τοῦ φωτὸς τοῦ ἀπολοίπου πηχῶν πέντε πλάτος κυκλόθεν.

\vs{12}Καὶ τὸ διορίζον κατὰ πρόσωπον τοῦ ἀπολοίπου, ὡς πρὸς θάλασσαν πηχῶν ἑβδομήκοντα πλάτος, τοῦ τοίχου τοῦ διορίζοντος πηχῶν πέντε εὖρος κυκλόθεν, καὶ μῆκος αὐτοῦ πηχῶν ἐννενήκοντα.
\vs{13}Καὶ διεμέτρησε κατέναντι τοῦ οἴκου μῆκος πηχῶν ἑκατὸν, καὶ τὰ ἀπόλοιπα καὶ τὰ διορίζοντα, καὶ οἱ τοῖχοι αὐτῶν μῆκος πηχῶν ἑκατόν.
\vs{14}Καὶ τὸ εὖρος κατὰ πρόσωπον τοῦ οἴκου, καὶ τὰ ἀπόλοιπα κατέναντι πηχῶν ἑκατόν.

\vs{15}Καὶ διεμέτρησε μῆκος τοῦ διορίζοντος κατὰ πρόσωπον τοῦ ἀπολοίπου τῶν κατόπισθεν τοῦ οἴκου ἐκείνου, καὶ τὰ ἀπόλοιπα ἔνθεν καὶ ἔνθεν πηχῶν ἑκατὸν τὸ μῆκος· καὶ ὁ ναὸς, καὶ αἱ γωνίαι, καὶ τὸ αἰλὰμ τὸ ἐξώτερον, πεφατνωμένα.
\vs{16}Καὶ αἱ θυρίδες δικτυωταὶ, ὑποφαύσεις κύκλῳ τοῖς τρισὶν, ὥστε διακύπτειν· καὶ ὁ οἶκος καὶ τὰ πλησίον ἐξυλωμένα κύκλῳ, καὶ τὸ ἔδαφος, καὶ ἐκ τοῦ ἐδάφους ἕως τῶν θυρίδων· καὶ αἱ θυρίδες ἀναπτυσσόμεναι τρισσῶς εἰς τὸ διακύπτειν.
\vs{17}Καὶ ἕως πλησίον τῆς ἐσωτέρας καὶ ἕως τῆς ἐξωτέρας, καὶ ἐφʼ ὅλον τὸν τοῖχον κύκλῳ ἐν τῷ ἔσωθεν καὶ ἐν τῷ ἔξωθεν,
\vs{18}γεγλυμμένα χερουβὶμ, καὶ φοίνικες ἀναμέσον χεροὺβ καὶ ἀναμέσον χερούβ· δύο πρόσωπα τῷ χερούβ·
\vs{19}Πρόσωπον ἀνθρώπου πρὸς τὸν φοίνικα ἔνθεν καὶ ἔνθεν, καὶ πρόσωπον λέοντος πρὸς τὸν φοίνικα ἔνθεν καὶ ἔνθεν· διαγεγλυμμένος ὁ οἶκος κυκλόθεν·
\vs{20}Ἐκ τοῦ ἐδάφους ἕως τοῦ φατνώματος, τὰ χερουβὶμ καὶ οἱ φοίνικες διαγεγλυμμένοι.

\vs{21}Καὶ τὸ ἅγιον καὶ ὁ ναὸς ἀναπτυσσόμενος τετράγωνα, κατὰ πρόσωπον τῶν ἁγίων ὅρασις ὡς ὄψις
\vs{22}θυσιαστηρίου ξυλίνου, πηχῶν τριῶν τὸ ὕψος αὐτοῦ, καὶ τὸ μῆκος πηχῶν δύο, καὶ τὸ εὖρος πηχῶν δύο· καὶ κέρατα εἶχε, καὶ ἡ βάσις αὐτοῦ καὶ οἱ τοῖχοι αὐτοῦ ξύλινοι· καὶ εἶπε πρὸς μὲ, αὕτη ἡ τράπεζα, ἡ πρὸ προσώπου Κυρίου.

\vs{23}Καὶ δύο θυρώματα τῷ ναῷ, καὶ δύο θυρώματα τῷ ἁγίῳ,
\vs{24}τοῖς δυσὶ θυρώμασι τοῖς στροφωτοῖς· δύο θυρώματα τῷ ἑνὶ, καὶ δύο θυρώματα τῇ θύρᾳ τῇ δευτέρᾳ.
\vs{25}Καὶ γλυφὴ ἐπʼ αὐτῶν, καὶ ἐπὶ τὰ θυρώματα τοῦ ναοῦ χερουβίμ· καὶ φοίνικες κατὰ τὴν γλυφὴν τῶν ἁγίων, καὶ σπουδαῖα ξύλα κατὰ πρόσωπον τοῦ αἰλὰμ ἔξωθεν,

\vs{26}Καὶ θυρίδες κρυπταί· καὶ διεμέτρησεν ἔνθεν καὶ ἔνθεν, εἰς τὰ ὀροφώματα τοῦ αἰλὰμ, καὶ τὰ πλευρὰ τοῦ οἴκου ἐζυγωμένα.

\ch{42}
Καὶ εἰσήγαγέ με εἰς τὴν αὐλὴν τὴν ἐσωτέραν κατὰ ἀνατολὰς, κατέναντι τῆς πύλης τῆς πρὸς Βοῤῥᾶν, καὶ εἰσήγαγέ με· καὶ ἰδοὺ ἐξέδραι πέντε, ἐχόμεναι τοῦ ἀπολοίπου, καὶ ἐχόμεναι τοῦ διορίζοντος πρὸς Βοῤῥᾶν,
\vs{2}ἐπὶ πήχεις ἑκατὸν μῆκος πρὸς Βοῤῥᾶν, καὶ τὸ πλάτος πεντήκοντα,
\vs{3}διαγεγραμμέναι ὃν τρόπον αἱ πύλαι τῆς αὐλῆς τῆς ἐσωτέρας, καὶ ὃν τρόπον τὰ περίστυλα τῆς αὐλῆς τῆς ἐξωτέρας ἐστοιχισμέναι, ἀντιπρόσωποι στοαὶ τρισσαί.
\vs{4}Καὶ κατέναντι τῶν ἐξεδρῶν περίπατος πηχῶν δέκα τὸ πλάτος, ἐπὶ πήχεις ἑκατὸν τὸ μῆκος, καὶ τὰ θυρώματα αὐτῶν πρὸς Βοῤῥᾶν,
\vs{5}καὶ οἱ περίπατοι οἱ ὑπερῷοι ὡσαύτως· ὅτι ἐξείχετο τὸ περίστυλον ἐξ αὐτοῦ, ἐκ τοῦ ὑποκάτωθεν περιστύλου, καὶ τὸ διάστημα· οὕτως περίστυλον καὶ διάστημα, καὶ οὕτως στοαὶ δύο.
\vs{6}Διότι τριπλαῖ ἦσαν, καὶ στύλους οὐκ εἶχον καθὼς οἱ στύλοι τῶν ἐξωτέρων· διατοῦτο ἐξείχοντο τῶν ὑποκάτωθεν καὶ τῶν μέσων ἀπὸ τῆς γῆς.

\vs{7}Καὶ φῶς ἔξωθεν, ὃν τρόπον αἱ ἐξέδραι τῆς αὐλῆς τῆς ἐξωτέρας, αἱ βλέπουσαι ἀπέναντι τῶν ἐξεδρῶν τῶν πρὸς Βοῤῥᾶν, μῆκος πηχῶν πεντήκοντα.
\vs{8}Ὅτι τὸ μῆκος τῶν ἐξεδρῶν τῶν βλεπουσῶν εἰς τὴν αὐλὴν τὴν ἐξωτέραν, ἦν πηχῶν πεντήκοντα, καὶ αὗταί εἰσιν αἱ ἀντιπρόσωποι ταύταις· τὸ πᾶν πηχῶν ἑκατόν.

\vs{9}Καὶ αἱ θύραι τῶν ἐξεδρῶν τούτων τῆς εἰσόδου τῆς πρὸς ἀνατολὰς, τοῦ εἰσπορεύεσθαι διʼ αὐτῶν ἐκ τῆς αὐλῆς τῆς ἐξωτέρας,
\vs{10}κατὰ τὸ φῶς τοῦ ἐν ἀρχῇ περιπάτου, καὶ τὰ πρὸς Νότον κατὰ πρόσωπον τοῦ Νότου κατὰ πρόσωπον τοῦ ἀπολοίπου, καὶ κατὰ πρόσωπον τοῦ διορίζοντος, καὶ αἱ ἐξέδραι.
\vs{11}Καὶ ὁ περίπατος κατὰ πρόσωπον αὐτῶν, κατὰ τὰ μέτρα ἐξεδρῶν τῶν πρὸς Βοῤῥᾶν, καὶ κατὰ τὸ μῆκος αὐτῶν, καὶ κατὰ τὸ εὖρος αὐτῶν, καὶ κατὰ πάσας τὰς ἐξόδους αὐτῶν, καὶ κατὰ πάσας τὰς ἐπιστροφὰς αὐτῶν, καὶ κατὰ τὰ φῶτα αὐτῶν, καὶ κατὰ τὰ θυρώματα αὐτῶν,
\vs{12}τῶν ἐξεδρῶν τῶν πρὸς Νότον, καὶ κατὰ τὰ θυρώματα ἀπʼ ἀρχῆς τοῦ περιπάτου, ὡς ἐπὶ φῶς διαστμήατος καλάμου, καὶ κατὰ ἀνατολὰς τοῦ εἰσπορεύεσθαι διʼ αὐτῶν.

\vs{13}Καὶ εἶπε πρὸς μὲ, αἱ ἐξέδραι αἱ πρὸς Βοῤῥᾶν, καὶ αἱ ἐξέδραι αἱ πρὸς Νότον, οὖσαι κατὰ πρόσωπον τῶν διαστημάτων, αὗταί εἰσιν αἱ ἐξέδραι τοῦ ἁγίου, ἐν αἷς φάγονται ἐκεῖ οἱ ἱερεῖς υἱοὶ Σαδδοὺκ, οἱ ἐγγίζοντες πρὸς Κύριον, τὰ ἅγια τῶν ἁγίων, καὶ ἐκεῖ θήσουσι τὰ ἅγια τῶν ἁγίων, καὶ τὴν θυσίαν, καὶ τὰ περὶ ἁμαρτίας, καὶ τὰ περὶ ἀγνοίας, διότι ὁ τόπος ἅγιος.
\vs{14}Οὐκ εἰσελεύσονται ἐκεῖ πάρεξ τῶν ἱερέων, οὐκ ἐξελεύσονται ἐκ τοῦ ἁγίου εἰς τὴν αὐλὴν τὴν ἐξωτέραν, ὅπως διαπαντὸς ἅγιοι ὦσιν οἱ προσάγοντες, καὶ μὴ ἅπτωνται τοῦ στολισμοῦ αὐτῶν, ἐν οἷς λειτουργοῦσιν ἐν αὐτοῖς, διότι ἅγιά ἐστι· καὶ ἐνδύσονται ἱμάτια ἕτερα, ὅταν ἅπτωνται τοῦ λαοῦ.

\vs{15}Καὶ συνετελέσθη ἡ διαμέτρησις τοῦ οἴκου ἔσωθεν· καὶ ἐξήγαγέ με καθʼ ὁδὸν τῆς πύλης τῆς βλεπούσης πρὸς ἀνατολὰς, καὶ διεμέτρησε τὸ ὑπόδειγμα τοῦ οἴκου κυκλόθεν ἐν διατάξει.

\vs{16}Καὶ ἔστη κατὰ νώτου τῆς πύλης τῆς βλεπούσης κατὰ ἀνατολὰς, καὶ διεμέτρησε πεντακοσίους ἐν τῷ καλάμῳ τοῦ μέτρου.
\vs{17}Καὶ ἐπέστρεψε πρὸς Βοῤῥᾶν, καὶ διεμέτρησε τὸ κατὰ πρόσωπον τοῦ Βοῤῥᾶ, πήχεις πεντακοσίους ἐν τῷ καλάμῳ τοῦ μέτρου.
\vs{18}Καὶ ἐπέστρεψε πρὸς θάλασσαν, καὶ διεμέτρησε τὸ κατὰ πρόσωπον θαλάσσης, πεντακοσίους ἐν τῷ καλάμῳ τοῦ μέτρου.
\vs{19}Καὶ ἐπέστρεψε πρὸς Νότον, καὶ διεμέτρησε κατέναντι τοῦ Νότου, πεντακοσίους ἐν τῷ καλάμῳ τοῦ μέτρου,
\vs{20}τὰ τέσσαρα μέρη τοῦ αὐτοῦ καλάμου· καὶ διέταξεν αὐτὸν, καὶ περίβολον αὐτῶν κύκλῳ, πεντακοσίων πρὸς ἀνατολὰς, καὶ πεντακοσίων πηχῶν εὖρος, τοῦ διαστέλλειν ἀναμέσον τῶν ἁγίων, καὶ ἀναμέσον τοῦ προτειχίσματος, τοῦ ἐν διατάξει τοῦ οἴκου.

\ch{43}
Καὶ ἤγαγέ με ἐπὶ τὴν πύλην τὴν βλέπουσαν κατὰ ἀνατολὰς, καὶ ἐξήγαγέ με.
\vs{2}Καὶ ἰδοὺ δόξα Θεοῦ Ἰσραὴλ ἤρχετο κατὰ τὴν ὁδὸν τὴν πρὸς ἀνατολὰς, καὶ φωνὴ τῆς παρεμβολῆς, ὡς φωνὴ διπλασιαζόντων πολλῶν· καὶ ἡ γῆ ἐξέλαμπεν ὡς φέγγος ἀπὸ τῆς δόξης κυκλόθεν.
\vs{3}Καὶ ἡ ὅρασις ἣν ἴδον, κατὰ τὴν ὅρασιν ἣν ἴδον, ὅτε εἰσεπορευόμην τοῦ χρίσαι τὴν πόλιν· καὶ ἡ ὅρασις τοῦ ἅρματος οὗ ἴδον, κατὰ τὴν ὅρασιν ἣν ἴδον ἐπὶ τοῦ ποταμοῦ τοῦ Χοβάρ· καὶ πίπτω ἐπὶ πρόσωπόν μου.

\vs{4}Καὶ δόξα Κυρίου εἰσῆλθεν εἰς τὸν οἶκον, κατὰ τὴν ὁδὸν τῆς πύλης τῆς βλεπούσης κατὰ ἀνατολάς.
\vs{5}Καὶ ἀνέλαβέ με πνεῦμα, καὶ εἰσήγαγέ με εἰς τὴν αὐλὴν τὴν ἐσωτέραν· καὶ ἰδοὺ πλήρης δόξης ὁ Κυρίου οἶκος.
\vs{6}Καὶ ἔστην, καὶ ἰδοὺ φωνὴ ἐκ τοῦ οἴκου λαλοῦντος πρὸς μέ, καὶ ὁ ἀνὴρ εἱστήκει ἐχόμενός μου,
\vs{7}καὶ εἶπε πρὸς μὲ,

Υἱὲ ἀνθρώπου, ἑώρακας τὸν τόπον τοῦ θρόνου μου, καὶ τὸν τόπον τοῦ ἴχνους τῶν ποδῶν μου, ἐν οἷς κατασκηνώσῃ τὸ ὄνομά μου ἐν μέσῳ οἴκου Ἰσραὴλ τὸν αἰῶνα· καὶ οὐ βεβηλώσουσιν οὐκέτι οἶκος Ἰσραὴλ τὸ ὄνομα τὸ ἅγιόν μου, αὐτοὶ καὶ οἱ ἡγούμενοι αὐτῶν ἐν τῇ πορνείᾳ αὐτῶν, καὶ ἐν τοῖς φόνοις τῶν ἡγουμένων ἐν μέσῳ αὐτῶν,
\vs{8}ἐν τῷ τιθέναι αὐτοὺς τὸ πρόθυρόν μου ἐν τοῖς προθύροις αὐτῶν, καὶ τὰς φλιάς μου ἐχομένας τῶν φλιῶν αὐτῶν· καὶ ἔδωκαν τὸν τοῖχόν μου ὡς συνεχόμενον ἐμοῦ καὶ αὐτῶν, καὶ ἐβεβήλωσαν τὸ ὄνομα τὸ ἅγιόν μου ἐν ταῖς ἀνομίαις αὐτῶν αἷς ἐποίουν· καὶ ἐξέτριψα αὐτοὺς ἐν θυμῷ μου, καὶ ἐν φόνῳ.
\vs{9}Καὶ νῦν ἀπωσάσθωσαν τὴν πορνείαν αὐτῶν, καὶ τοὺς φόνους τῶν ἡγουμένων αὐτῶν ἀπʼ ἐμοῦ, καὶ κατασκηνώσω ἐν μέσῳ αὐτῶν τὸν αἰῶνα.

\vs{10}Καὶ σὺ, υἱὲ ἀνθρώπου, δεῖξον τῷ οἴκῳ Ἰσραὴλ τὸν οἶκον· καὶ κοπάσουσιν ἀπὸ τῶν ἁμαρτιῶν αὐτῶν· καὶ τὴν ὅρασιν αὐτοῦ, καὶ τὴν διάταξιν αὐτοῦ.
\vs{11}Καὶ αὐτοὶ λήψονται τὴν κόλασιν αὐτῶν περὶ πάντων ὧν ἐποίησαν· καὶ διαγράψεις τὸν οἶκον, καὶ τὰς ἐξόδους αὐτοῦ, καὶ τὴν ὑπόστασιν αὐτοῦ, καὶ πάντα τὰ προστάγματα αὐτοῦ, καὶ πάντα τὰ νόμιμα αὐτοῦ γνωριεῖς αὐτοῖς, καὶ διαγράψεις ἐναντίον αὐτῶν· καὶ φυλάξονται πάντα τὰ δικαιώματά μου, καὶ πάντα τὰ προστάγματά μου, καὶ ποιήσουσιν αὐτά.

\vs{12}Καὶ τὴν διαγραφὴν τοῦ οἴκου ἐπὶ τῆς κορυφῆς τοῦ ὄρους, πάντα τὰ ὅρια αὐτοῦ κυκλόθεν ἅγια ἁγίων.

\vs{13}Καὶ ταῦτα τὰ μέτρα τοῦ θυσιαστηρίου ἐν πήχει τοῦ πήχεως, καὶ παλαιστῆς, κόλπωμα βάθους ἐπὶ πῆχυν, καὶ πῆχυς τὸ εὖρος, καὶ γεῖσος ἐπὶ τὸ χεῖλος αὐτοῦ κυκλόθεν, σπιθαμῆς· καὶ τοῦτο τὸ ὕψος τοῦ θυσιαστηρίου
\vs{14}ἐκ βάθους τῆς ἀρχῆς τοῦ κοιλώματος αὐτοῦ, πρὸς τὸ ἱλαστήριον τὸ μέγα τοῦτο, ὑποκάτωθεν πηχῶν δύο, καὶ τὸ εὖρος πήχεος· καὶ ἀπὸ τοῦ ἱλαστηρίου τοῦ μικροῦ ἐπὶ τὸ ἱλαστήριον τὸ μέγα, πήχεις τέσσαρες, καὶ εὖρος πῆχυς.
\vs{15}Καὶ τὸ ἀριὴλ πηχῶν τεσσάρων, καὶ ἀπὸ τοῦ ἀριὴλ, καὶ ὑπεράνω τῶν κεράτων πῆχυς.
\vs{16}Καὶ τὸ ἀριὴλ πηχῶν δώδεκα μήκους, ἐπὶ πήχεις δώδεκα, τετράγωνον ἐπὶ τὰ τέσσαρα μέρη αὐτοῦ.

\vs{17}Καὶ τὸ ἱλαστήριον πηχῶν δεκατεσσάρων τὸ μῆκος, ἐπὶ πήχεις δέκατέσσαρας τὸ εὖρος ἐπὶ τέσσαρα μέρη αὐτοῦ, καὶ τὸ γεῖσος αὐτῷ κυκλόθεν κυκλούμενον αὐτῷ ἥμισυ πήχεως· καὶ τὸ κύκλωμα αὐτοῦ πῆχυς κυκλόθεν, καὶ οἱ κλημακτῆρες αὐτοῦ βλέποντες κατὰ ἀνατολάς.

\vs{18}Καὶ εἶπε πρὸς μὲ, υἱὲ ἀνθρώπου, τάδε λέγει Κύριος ὁ Θεὸς Ἰσραὴλ, ταῦτα τὰ προστάγματα τοῦ θυσιαστηρίου ἐν ἡμέρᾳ ποιήσεως αὐτοῦ, τοῦ ἀναφέρειν ἐπʼ αὐτοῦ ὁλοκαυτώματα, καὶ προσχέειν πρὸς αὐτὸ αἷμα.
\vs{19}Καὶ δώσεις τοῖς ἱερεῦσι τοῖς Λευίταις τοῖς ἐκ τοῦ σπέρματος Σαδδοὺκ τοῖς ἐγγίζουσι πρὸς μὲ, λέγει Κύριος ὁ Θεὸς, τοῦ λειτουργεῖν μοι μόσχον ἐκ βοῶν περὶ ἁμαρτίας.
\vs{20}Καὶ λήψονται ἐκ τοῦ αἵματος αὐτοῦ, καὶ ἐπιθήσουσιν ἐπὶ τὰ τέσσαρα κέρατα τοῦ θυσιαστηρίου, καὶ ἐπὶ τὰς τέσσαρας γωνίας τοῦ ἱλαστηρίου καὶ ἐπὶ τὴν βάσιν κύκλῳ, καὶ ἐξιλάσονται αὐτό.
\vs{21}Καὶ λήψονται τὸν μόσχον τὸν περὶ ἁμαρτίας, καὶ κατακαυθήσεται ἐν τῷ ἀποκεχωρισμένῳ τοῦ οἴκου, ἔξωθεν τῶν ἁγίων.
\vs{22}Καὶ τῇ ἡμέρᾳ τῇ δευτέρᾳ λήψονται ἐρίφους δύο αἰγῶν ἀμώμους ὑπὲρ ἁμαρτίας, καὶ ἐξιλάσονται τὸ θυσιαστήριον, καθότι ἐξιλάσαντο ἐν τῷ μόσχῳ.
\vs{23}Καὶ μετὰ τὸ συντελέσαι τὸν ἐξιλασμὸν, προσοίσουσι μόσχον ἐκ βοῶν ἄμωμον, καὶ κριὸν ἐκ προβάτων ἄμωμον,
\vs{24}καὶ προσοίσετε ἐναντίον Κυρίου· καὶ ἐπιῤῥίψουσιν οἱ ἱερεῖς ἐπʼ αὐτὰ ἅλα, καὶ ἀνοίσουσιν αὐτὰ ὁλοκαυτώματα τῷ Κυρίῳ.

\vs{25}Ἑπτὰ ἡμέρας ποιήσεις ἔριφον ὑπὲρ ἁμαρτίας καθʼ ἡμέραν, καὶ μόσχον ἐκ βοῶν, καὶ κριὸν ἐκ προβάτων, ἄμωμα ποιήσουσιν
\vs{26}ἑπτὰ ἡμέρας· καὶ ἐξιλάσονται τὸ θυσιαστήριον, καὶ καθαριοῦσιν αὐτό· καὶ πλήσουσι χεῖρας αὐτῶν.
\vs{27}Καὶ ἔσται ἀπὸ τῆς ἡμέρας τῆς ὀγδόης καὶ ἐπέκεινα, ποιήσουσιν οἱ ἱερεῖς ἐπὶ τὸ θυσιαστήριον τὰ ὁλοκαυτώματα ὑμῶν, καὶ τὰ τοῦ σωτηρίου ὑμῶν, καὶ προσδέξομαι ὑμᾶς, λέγει Κύριος.

\ch{44}
Καὶ ἐπέστρεψέ με κατὰ τὴν ὁδὸν τῆς πύλης τῶν ἁγίων τῆς ἐξωτέρας τῆς βλεπούσης κατὰ ἀνατολάς· καὶ αὕτη ἦν κεκλεισμένη.
\vs{2}Καὶ εἶπε Κύριος πρὸς μὲ, ἡ πύλη αὕτη κεκλεισμένη ἔσται, οὐκ ἀνοιχθήσεται, καὶ οὐδεὶς μὴ διέλθῃ διʼ αὐτῆς· ὅτι Κύριος ὁ Θεὸς Ἰσραὴλ εἰσελεύσεται διʼ αὐτῆς, καὶ ἔσται κεκλεισμένη.
\vs{3}Διότι ὁ ἡγούμενος οὗτος καθήσεται ἐν αὐτῇ, τοῦ φαγεῖν ἄρτον ἐναντίον Κυρίου· κατὰ τὴν ὁδὸν αἰλὰμ τῆς πύλης εἰσελεύσεται, καὶ κατὰ τὴν ὁδὸν αὐτοῦ ἐξελεύσεται.

\vs{4}Καὶ εἰσήγαγέ με κατὰ τὴν ὁδὸν τῆς πύλης τῆς πρὸς Βοῤῥᾶν, κατέναντι τοῦ οἴκου· καὶ ἴδον, καὶ ἰδοὺ πλήρης δόξης ὁ οἶκος τοῦ Κυρίου· καὶ πίπτω ἐπὶ πρόσωπόν μου.
\vs{5}Καὶ εἶπε Κύριος πρὸς μὲ, υἱὲ ἀνθρώπου, τάξον εἰς τὴν καρδίαν σου, καὶ ἴδε τοῖς ὀφθαλμοῖς, καὶ τοῖς ὠσί σου ἄκουε πάντα ὅσα ἐγὼ λαλῶ μετὰ σοῦ, κατὰ πάντα τὰ προστάγματα τοῦ οἴκου Κυρίου, καὶ πάντα τὰ νόμιμα αὐτοῦ· καὶ τάξεις τὴν καρδίαν σου εἰς τὴν εἴσοδον τοῦ οἴκου, κατὰ πάσας τὰς ἐξόδους αὐτοῦ, ἐν πᾶσι τοῖς ἁγίοις.
\vs{6}Καὶ ἐρεῖς πρὸς τὸν οἶκον τὸν παραπικραίνοντα, πρὸς τὸν οἶκον τοῦ Ἰσραὴλ, τάδε λέγει Κύριος ὁ Θεὸς, ἱκανούσθω ὑμῖν ἀπὸ πασῶν τῶν ἀνομιῶν ὑμῶν, οἶκος Ἰσραὴλ,
\vs{7}τοῦ εἰσαγαγεῖν ὑμᾶς υἱοὺς ἀλλογενεῖς, ἀπεριτμήτους καρδίᾳ· καὶ ἀπεριτμήτους σαρκὶ, τοῦ γίνεσθαι ἐν τοῖς ἁγίοις μου, καὶ βεβηλοῦν αὐτὰ, ἐν τῷ προσφέρειν ὑμᾶς ἄρτους, σάρκας, καὶ αἷμα· καὶ παρεβαίνετε τὴν διαθήκην μου ἐν πάσαις ταῖς ἀνομίαις ὑμῶν,
\vs{8}καὶ διετάξατε τοῦ φυλάσσειν φυλακὰς ἐν τοῖς ἁγίοις μου.

\vs{9}Διατοῦτο τάδε λέγει Κύριος ὁ Θεὸς, πᾶς υἱὸς ἀλλογενὴς ἀπερίτμητος καρδίᾳ, καὶ ἀπερίτμητος σαρκὶ, οὐκ εἰσελεύσεται εἰς τὰ ἅγιά μου ἐν πᾶσιν υἱοῖς ἀλλογενῶν, τῶν ὄντων ἐν μέσῳ οἴκου Ἰσραήλ·

\vs{10}Ἀλλʼ ἢ οἱ Λευῖται, οἵτινες ἀφήλαντο ἀπʼ ἐμοῦ ἐν τῷ πλανᾶσθαι τὸν Ἰσραὴλ ἀπʼ ἐμοῦ κατόπισθεν τῶν ἐνθυμημάτων αὐτῶν· καὶ λήψονται ἀδικίαν αὐτῶν,
\vs{11}καὶ ἔσονται ἐν τοῖς ἁγίοις μου λειτουργοῦντες, θυρωροὶ ἐπὶ τῶν πυλῶν τοῦ οἴκου, καὶ λειτουργοῦντες τῷ οἴκῳ· οὗτοι σφάξουσι τὰς θυσίας καὶ τὰ ὁλοκαυτώματα τῷ λαῷ, καὶ οὗτοι στήσονται ἐναντίον τοῦ λαοῦ, τοῦ λειτουργεῖν αὐτοῖς.
\vs{12}Ἀνθʼ ὧν ἐλειτούργουν αὐτοῖς πρὸ προσώπου τῶν εἰδώλων αὐτῶν· καὶ ἐγένετο τῷ οἴκῳ Ἰσραὴλ εἰς κόλασιν ἀδικίας· ἕνεκα τούτου ᾖρα τὴν χεῖρά μου ἐπʼ αὐτοὺς, λέγει Κύριος ὁ Θεὸς,
\vs{13}καὶ οὐκ ἐγγιοῦσι πρὸς μὲ τοῦ ἱερατεύειν μοι, οὐδὲ τοῦ προσάγειν πρὸς τὰ ἅγια υἱῶν τοῦ Ἰσραὴλ, οὐδὲ πρὸς τὰ ἅγια τῶν ἁγίων μου· καὶ λήψονται ἀτιμίαν αὐτῶν ἐν τῇ πλανήσει ᾗ ἐπλανήθησαν.
\vs{14}Κατάξουσιν αὐτοὺς φυλάσσειν φυλακὰς τοῦ οἴκου εἰς πάντα τὰ ἔργα αὐτοῦ, καὶ εἰς πάντα ὅσα ἂν ποιήσωσιν.

\vs{15}Οἱ ἱερεῖς οἱ Λευῖται, οἱ υἱοὶ τοῦ Σαδδοὺκ, οἵτινες ἐφυλάξαντο τὰς φυλακὰς τῶν ἁγίων μου, ἐν τῷ πλανᾶσθαι οἶκον Ἰσραὴλ ἀπʼ ἐμοῦ, οὗτοι προσάξουσι πρὸς μὲ, τοῦ λειτουργεῖν μοι, καὶ στήσονται πρὸ προσώπου μου, τοῦ προσφέρειν μοι θυσίαν, στέαρ καὶ αἷμα, λέγει Κύριος ὁ Θεός.
\vs{16}Οὗτοι εἰσελεύσονται εἰς τὰ ἅγιά μου, καὶ οὗτοι προσελεύσονται πρὸς τὴν τράπεζάν μου, τοῦ λειτουργεῖν μοι, καὶ φυλάξουσι τὰς φυλακάς μου.

\vs{17}Καὶ ἔσται ἐν τῷ εἰσπορεύεσθαι αὐτοὺς τὰς πύλας τῆς αὐλῆς τῆς ἐσωτέρας, στολὰς λινὰς ἐνδύσονται, καὶ οὐκ ἐνδύσονται ἔρια ἐν τῷ λειτουργεῖν αὐτοὺς ἀπὸ τῆς πύλης τῆς ἐσωτέρας αὐλῆς.
\vs{18}Καὶ κιδάρεις λινὰς ἕξουσιν ἐπὶ ταῖς κεφαλαῖς αὐτῶν, καὶ περισκελῆ λινὰ ἕξουσιν ἐπὶ τὰς ὀσφύας αὐτῶν, καὶ οὐ περιζώσονται βίᾳ.
\vs{19}Καὶ ἐν τῷ ἐκπορεύεσθαι αὐτοὺς εἰς τὴν αὐλὴν τὴν ἐξωτέραν πρὸς τὸν λαὸν, ἐκδύσονται τὰς στολὰς αὐτῶν, ἐν αἷς αὐτοὶ λειτουργοῦσιν ἐν αὐταῖς· καὶ θήσουσιν αὐτὰς ἐν ταῖς ἐξέδραις τῶν ἁγίων, καὶ ἐνδύσονται στολὰς ἑτέρας, καὶ οὐ μὴ ἁγιάσωσι τὸν λαὸν ἐν ταῖς στολαῖς αὐτῶν.
\vs{20}Καὶ τὰς κεφαλὰς αὐτῶν οὐ ξυρήσονται, καὶ τὰς κόμας αὐτῶν οὐ ψιλώσουσι, καλύπτοντε καλύψουσιν τὰς κεφαλὰς αὐτῶν.
\vs{21}Καὶ οἶνον οὐ μὴ πίωσι πᾶς ἱερεὺς, ἐν τῷ εἰσπορεύεσθαι αὐτοὺς εἰς τὴν αὐλὴν τὴν ἐσωτέραν.
\vs{22}Καὶ χήραν καὶ ἐκβεβλημένην οὐ λήψονται ἑαυτοῖς εἰς γυναῖκα, ἀλλʼ ἢ παρθένον ἐκ τοῦ σπέρματος Ἰσραήλ· καὶ χήρα ἐὰν γένηται ἐξ ἱερέως, λήψονται.

\vs{23}Καὶ τὸν λαόν μου διδάξουσιν ἀναμέσον ἁγίου καὶ βεβήλου, καὶ ἀναμέσον ἀκαθάρτου καὶ καθαροῦ γνωριοῦσιν αὐτοῖς.
\vs{24}Καὶ ἐπὶ κρίσιν αἵματος οὗτοι ἐπιστήσονται τοῦ διακρίνειν· τὰ δικαιώματά μου δικαιώσουσι, καὶ τὰ κρίματά μου κρινοῦσι, καὶ τὰ νόμιμά μου καὶ τὰ προστάγματά μου ἐν πάσαις ταῖς ἑορταῖς μου φυλάξονται, καὶ τὰ σάββατά μου ἁγιάσουσι.

\vs{25}Καὶ ἐπὶ ψυχὴν ἀνθρώπου οὐκ εἰσελεύσονται τοῦ μιανθῆναι, ἀλλʼ ἢ ἐπὶ πατρὶ, καὶ ἐπὶ μητρὶ, καὶ ἐπὶ υἱῷ, καὶ ἐπὶ θυγατρὶ, καὶ ἐπὶ ἀδελφῷ, καὶ ἐπὶ ἀδελφῇ αὐτοῦ, ἣ οὐ γέγονεν ἀνδρὶ, μιανθήσεται.
\vs{26}Καὶ μετὰ τὸ καθαρισθῆναι αὐτὸν, ἑπτὰ ἡμέρας ἐξαριθμήσῃ αὐτῷ.
\vs{27}Καὶ ᾗ ἂν ἡμέρᾳ εἰσπορεύωνται εἰς τὴν αὐλὴν τὴν ἐσωτέραν τοῦ λειτουργεῖν ἐν τῷ ἁγίῳ, προσοίσουσιν ἱλασμὸν, λέγει Κύριος ὁ Θεός·

\vs{28}Καὶ ἔσται αὐτοῖς εἰς κληρονομίαν· ἐγὼ κληρονομία αὐτοῖς, καὶ κατάσχεσις αὐτοῖς οὐ δοθήσεται ἐν τοῖς υἱοῖς Ἰσραὴλ, ὅτι ἐγὼ κατάσχεσις αὐτῶν.
\vs{29}Καὶ τὰς θυσίας, καὶ τὰ ὑπὲρ ἁμαρτίας, καὶ τὰ ὑπὲρ ἀγνοίας, οὗτοι φάγονται· καὶ πᾶν ἀφόρισμα ἐν τῷ Ἰσραὴλ αὐτοῖς ἔσται,
\vs{30}ἀπαρχαὶ πάντων, καὶ τὰ πρωτότοκα πάντων, καὶ τὰ ἀφαιρέματα πάντα· ἐκ πάντων τῶν ἀπαρχῶν ὑμῶν, τοῖς ἱερεῦσιν ἔσται· καὶ τὰ πρωτογεννήματα ὑμῶν δώσετε τῷ ἱερεῖ, τοῦ θεῖναι εὐλογίας ὑμῶν ἐπὶ τοὺς οἴκους ὑμῶν.
\vs{31}Καὶ πᾶν θνησιμαῖον καὶ θηριάλωτον ἐκ τῶν πετεινῶν, καὶ ἐκ τῶν κτηνῶν, οὐ φάγονται οἱ ἱερεῖς.

\ch{45}
Καὶ ἐν τῷ καταμετρεῖσθαι ὑμᾶς τὴν γῆν ἐν κληρονομίᾳ, ἀφοριεῖτε ἀπαρχὴν τῷ Κυρίῳ, ἅγιον ἀπὸ τῆς γῆς πέντε καὶ εἴκοσι χιλιάδας μῆκος, καὶ εὖρος εἴκοσι χιλιάδας, ἅγιον ἔσται ἐν πᾶσι τοῖς ὁρίοις αὐτοῦ κυκλόθεν.
\vs{2}Καὶ ἔσται ἐκ τούτου ἁγιάσματα, πεντακόσιοι ἐπὶ πεντακοσίους, τετράγωνον κυκλόθεν, καὶ πεντήκοντα πήχεις διάστημα αὐτῶν κυκλόθεν.
\vs{3}Καὶ ἐκ ταύτης τῆς διαμετρήσεως διαμετρήσεις μῆκος πέντε καὶ εἴκοσι χιλιάδας, καὶ εὖρος εἴκοσι χιλιάδας· καὶ ἐν αὐτῇ ἔσται ἅγια τῶν ἁγίων·
\vs{4}Ἀπὸ τῆς γῆς ἔσται τοῖς ἱερεῦσι τοῖς λειτουργοῦσιν ἐν τῷ ἁγίῳ, καὶ ἔσται τοῖς ἐγγίζουσι λειτουργεῖν τῷ Κυρίῳ· καὶ ἔσται αὐτοῖς τόπος εἰς οἴκους ἀφωρισμένους τῷ ἁγιασμῷ αὐτῶν,
\vs{5}εἴκοσι καὶ πέντε χιλιάδας μῆκος, καὶ εὖρος εἴκοσι χιλιάδες· καὶ τοῖς Λευίταις τοῖς λειτουργοῦσι τῷ οἴκῳ, αὐτοῖς εἰς κατάσχεσιν πόλεις τοῦ κατοικεῖν.

\vs{6}Καὶ τὴν κατάσχεσιν τῆς πόλεως δώσεις πέντε χιλιάδας εὖρος, καὶ μῆκος πέντε καὶ εἴκοσι χιλιάδας, ὃν τρόπον ἡ ἀπαρχὴ τῶν ἁγίων παντὶ οἴκῳ Ἰσραὴλ ἔσονται.

\vs{7}Καὶ τῷ ἡγουμένῳ ἐκ τούτου, καὶ ἀπὸ τούτου εἰς τὰς ἀπαρχὰς τῶν ἁγίων, εἰς κατάσχεσιν τῆς πόλεως, κατὰ πρόσωπον τῶν ἀπαρχῶν τῶν ἁγίων, καὶ κατὰ πρόσωπον τῆς κατασχέσεως τῆς πόλεως τὰ πρὸς θάλασσαν, καὶ ἀπὸ τῶν πρὸς θάλασσαν πρὸς ἀνατολάς· καὶ τὸ μῆκος ὡς μία τῶν μερίδων ἀπὸ τῶν ὁρίων τῶν πρὸς θάλασσαν, καὶ τὸ μῆκος, ἐπὶ τὰ ὅρια τὰ πρὸς ἀνατολὰς τῆς γῆς.
\vs{8}Καὶ ἔσται αὐτῷ εἰς κατάσχεσιν ἐν τῷ Ἰσραὴλ, καὶ οὐ καταδυναστεύσουσιν οὐκέτι οἱ ἀφηγούμενοι τοῦ Ἰσραὴλ τὸν λαόν μου, καὶ τὴν γῆν κατακληρονομήσουσιν οἶκος Ἰσραὴλ κατὰ φυλὰς αὐτῶν.

\vs{9}Τάδε λέγει Κύριος Θεὸς, ἱκανούσθω ὑμῖν οἱ ἀφηγούμενοι τοῦ Ἰσραὴλ, ἀδικίαν καὶ ταλαιπωρίαν ἀφέλεσθε, κρίμα καὶ δικαιοσύνην ποιήσατε· ἐξάρατε καταδυναστείαν ἀπὸ τοῦ λαοῦ μου, λέγει Κύριος Θεός.
\vs{10}Ζυγὸς δίκαιος, καὶ μέτρον δίκαιον, καὶ χοῖνιξ δικαία ἔσται ὑμῖν τοῦ μέτρου,
\vs{11}καὶ ἡ χοῖνιξ ὁμοίως μία ἔσται τοῦ λαμβάνειν· τὸ δέκατον τοῦ γομόρ, χοῖνιξ, καὶ τὸ δέκατον τοῦ γομόρ, πρὸς τὸ γομόρ ἔσται τὸ ἴσον.
\vs{12}Καὶ τὰ στάθμια, εἴκοσι ὀβολοὶ, πέντε σίκλοι, πέντε καὶ σίκλοι δέκα, καὶ πεντήκοντα σίκλοι ἡ μνᾶ ἔσται ὑμῖν.

\vs{13}Καὶ αὕτη ἡ ἀπαρχὴ ἣν ἀφοριεῖτε, ἕκτον μέτροο ἀπὸ τοῦ γυμὸρ τοῦ πυροῦ, καὶ τὸ ἕκτον αὐτοῦ τοῦ οἴφι ἀπὸ τοῦ κόρου τῶν κριθῶν.
\vs{14}Καὶ τὸ πρόσταγμα τοῦ ἐλαίου κοτύλην ἐλαίου ἀπὸ δέκα κοτυλῶν, ὅτι αἱ δέκα κοτύλαι εἰσὶ γομόρ.
\vs{15}Καὶ πρόβατον ἀπὸ τῶν προβάτων ἀπὸ δέκα ἀφαίρεμα ἐκ πασῶν τῶν πατριῶν τοῦ Ἰσραὴλ, εἰς θυσίας, καὶ εἰς ὁλοκαυτώματα, καὶ εἰς σωτηρίου, τοῦ ἐξιλάσκεσθαι περὶ ὑμῶν, λέγει Κύριος Θεός.
\vs{16}Καὶ πᾶς ὁ λαὸς δώσει τὴν ἀπαρχὴν ταύτην τῷ ἀφηγουμένῳ τοῦ Ἰσραήλ.

\vs{17}Καὶ διὰ τοῦ ἀφηγουμένου ἔσται τὰ ὁλοκαυτώματα, καὶ αἱ θυσίαι, καὶ αἱ σπονδαὶ ἐν ταῖς ἑορταῖς, καὶ ἐν ταῖς νουμηνίαις, καὶ ἐν τοῖς σαββάτοις, καὶ ἐν πάσαις ταῖς ἑορταῖς οἴκου Ἰσραήλ· αὐτὸς ποιήσει τὰ ὑπὲρ ἁμαρτίας, καὶ τὴν θυσίαν, καὶ τὰ ὁλοκαυτώματα, καὶ τὰ τοῦ σωτηρίου, τοῦ ἐξιλάσκεσθαι ὑπὲρ τοῦ οἴκου Ἰσραήλ.

\vs{18}Τάδε λέγει Κύριος Θεὸς, ἐν τῷ πρώτῳ μηνὶ, μιᾷ τοῦ μηνὸς, λήμψεσθε μόσχον ἐκ βοῶν ἄμωμον, τοῦ ἐξιλάσασθαι τὸ ἅγιον.
\vs{19}Καὶ λήψεται ὁ ἱερεὺς ἀπὸ τοῦ αἵματος τοῦ ἐξιλασμοῦ, καὶ δώσει ἐπὶ τὰς φλιὰς τοῦ οἴκου, καὶ ἐπὶ τὰς τέσσαρας γωνίας τοῦ ἱεροῦ, καὶ ἐπὶ τὸ θυσιαστήριον, καὶ ἐπὶ τὰς φλιὰς τῆς πύλης τῆς αὐλῆς τῆς ἐσωτέρας.
\vs{20}Καὶ οὕτως ποιήσεις ἐν τῷ μηνὶ τῷ ἑβδόμῳ· μιᾷ τοῦ μηνὸς λήψῃ παρʼ ἑκάστου ἀπόμοιραν· καὶ ἐξιλάσεσθε τὸν οἶκον.

\vs{21}Καὶ ἐν τῷ πρώτῳ, τεσσαρεσκαιδεκάτῃ τοῦ μηνὸς ἔσται ὑμῖν τὸ πάσχα ἑορτή· ἑπτὰ ἡμέρας ἄζυμα ἔδεσθε.
\vs{22}Καὶ ποιήσει ὁ ἀφηγούμενος ἐν ἐκείνῃ τῇ ἡμέρᾳ ὑπὲρ αὐτοῦ, καὶ τοῦ οἴκου, καὶ ὑπὲρ παντὸς τοῦ λαοῦ τῆς γῆς, μόσχον ὑπὲρ ἁμαρτίας.
\vs{23}Καὶ τὰς ἑπτὰ ἡμέρας τῆς ἑορτῆς ποιήσει ὁλοκαυτώματα τῷ Κυρίῳ ἑπτὰ μόσχους καὶ ἑπτὰ κριοὺς ἀμώμους καθʼ ἡμέραν, τὰς ἑπτὰ ἡμέρας, καὶ ὑπὲρ ἁμαρτίας ἔριφον αἰγῶν καθʼ ἡμέραν, καὶ θυσίαν.
\vs{24}Καὶ πέμμα τῷ μόσχῳ, καὶ πέμματα τῷ κριῷ ποιήσεις, καὶ ἐλαίου τὸ εἲν τῷ πεμματι.
\vs{25}Καὶ ἐν τῷ ἑβδόμῷ μηνὶ, πεντεκαιδεκάτῃ τοῦ μηνὸς, ἐν τῇ ἑορτῇ ποιήσεις κατὰ τὰ αὐτὰ, ἑπτὰ ἡμέρας, καθὼς τὰ ὑπὲρ τῆς ἁμαρτίας, καὶ καθὼς τὰ ὁλοκαυτώματα, καὶ καθὼς τὸ μαναὰ, καὶ καθὼς τὸ ἔλαιον.

\ch{46}
Τάδε λέγει Κύριος Θεὸς, ἡ πύλη ἡ ἐν τῇ αὐλῇ τῇ ἐσωτέρᾳ, ἡ βλέπουσα πρὸς ἀνατολὰς, ἔσται κεκλεισμένη ἓξ ἡμέρας τὰς ἐνεργούς· ἐν τῇ ἡμέρᾳ τῶν σαββάτων ἀνοιχθῇ, καὶ ἐν τῇ ἡμέρᾳ τῆς νουμηνίας ἀνοιχθήσεται.
\vs{2}Καὶ εἰσελεύσεται ὁ ἀφηγούμενος κατὰ τὴν ὁδὸν τοῦ αἰλὰμ τῆς πύλης τῆς ἔσωθεν, καὶ στήσεται ἐπὶ τὰ πρόθυρα τῆς πύλης, καὶ ποιήσουσιν οἱ ἱερεῖς τὰ ὁλοκαυτώματα αὐτοῦ, καὶ τὰ τοῦ σωτηρίου αὐτοῦ.
\vs{3}Καὶ προσκυνήσει ἡ λαὸς τῆς γῆς κατὰ τὰ πρόθυρα τῆς πύλης ἐκείνης, καὶ ἐν τοῖς σαββάτοις, καὶ ἐν ταῖς νουμηνίαις, ἐναντίον Κυρίου.

\vs{4}Καὶ τὰ ὁλοκαυτώματα προσοίσει ὁ ἀφηγούμενος τῷ Κυρίῳ ἐν τῇ ἡμέρᾳ τῶν σαββάτων, ἓξ ἀμνοὺς ἀμώμους, καὶ κριὸν ἂμωμον,
\vs{5}καὶ μαναὰ, πέμμα τῷ κριῷ, καὶ τοῖς ἀμνοῖς θυσίαν, δόμα χειρὸς αὐτοῦ, καὶ ἐλαίου τὸ εἲν τῷ πέμματι.
\vs{6}Καὶ ἐν τῇ ἡμέρᾳ τῆς νουμηνίας μόσχον ἄμωμον, καὶ ἓξ ἀμνοὺς, καὶ κριὸς ἄμωμος ἔσται,
\vs{7}καὶ πέμμα τῷ κριῷ, καὶ πέμμα τῷ μόσχῳ ἔσται μαναά, καὶ τοῖς ἀμνοῖς, καθὼς ἂν ἐκποιῇ ἡ χεὶρ αὐτοῦ, καὶ ἐλαίου τὸ εἲν τῷ πέμματι.

\vs{8}Καὶ ἐν τῷ εἰσπορεύεσθαι τὸν ἀφηγούμενον, κατὰ τὴν ὁδὸν τοῦ αἰλὰμ τῆς πύλης εἰσελεύσεται, καὶ κατὰ τὴν ὁδὸν τῆς πύλης ἐξελεύσεται.
\vs{9}Καὶ ὅταν εἰσπορεύηται ὁ λαὸς τῆς γῆς ἐναντίον Κυρίου ἐν ταῖς ἑορταῖς, ὁ εἰσπορευόμενος κατὰ τὴν ὁδὸν τῆς πύλης τῆς βλεπούσης πρὸς Βοῤῥᾶν προσκυνεῖν, ἐξελεύσεται κατὰ τὴν ὁδὸν τῆς πύλης τῆς πρὸς Νότον· καὶ ὁ εἰσπορευόμενος κατὰ τὴν ὁδὸν τῆς πύλης τῆς πρὸς Νότον, ἐξελεύσεται κατὰ τὴν ὁδὸν τῆς πύλης τῆς πρὸς Βοῤῥᾶν· οὐκ ἀναστρέψει κατὰ τὴν πύλην εἰς ἣν εἰσελήλυθεν, ἀλλʼ ἢ κατʼ εὐθὺ αὐτῆς ἐξελεύσεται.
\vs{10}Καὶ ὁ ἀφηγούμενος ἐν μέσῳ αὐτῶν, ἐν τῷ εἰσπορεύεσθαι αὐτοὺς εἰσελεύσεται μετʼ αὐτῶν, καὶ ἐν τῷ ἐκπορεύεσθαι αὐτοὺς ἐξελεύσεται.

\vs{11}Καὶ ἐν ταῖς ἑορταῖς καὶ ἐν ταῖς πανηγύρεσιν ἔσται τὸ μαναὰ πέμμα τῷ μόσχῳ, καὶ πέμμα τῷ κριῷ, καὶ τοῖς, ἀμνοῖς καθὼς ἂν ἐκποιῇ ἡ χεὶρ αὐτοῦ, καὶ ἐλαίου τὸ εἲν τῷ πέμματι.
\vs{12}Ἐὰν δὲ ποιήσῃ ὁ ἀφηγούμενος ὁμολογίαν ὁλοκαύτωμα σωτηρίου τῷ Κυρίῳ, καὶ ἀνοίξεῃ ἑαυτῷ τὴν πύλην τὴν βλέπουσαν κατὰ ἀνατολάς, καὶ ποιήσῃ τὸ ὁλοκαύτωμα αὐτοῦ, καὶ τὰ τοῦ σωτηρίου αὐτοῦ, ὃν τρόπον ποιεῖ ἐν τῇ ἡμέρᾳ τῶν σαββάτων· καὶ ἐξελεύσεται, καὶ κλείσει τὰς θύρας μετὰ τὸ ἐξελθεῖν αὐτόν·

\vs{13}Καὶ ἀμνὸν ἐνιαύσιον ἄμωμον ποιήσει εἰς ὁλοκαύτωμα καθʼ ἡμέραν τῷ Κυρίῳ, πρωῒ ποιήσει αὐτόν·
\vs{14}Καὶ μαναὰ ποιήσει ἐπʼ αὐτῷ τὸπρωῒ, ἕκτον τοῦ μέτρου, καὶ ἐλαίου τρίτον τοῦ εἲν τοῦ ἀναμίξαι τὴν σεμίδαλιν μαναὰ τῷ Κυρίῳ, πρόσταγμα διαπαντός.
\vs{15}Ποιήσετε τὸν ἀμνὸν, καὶ τὸ μαναὰ, καὶ τὸ ἔλαιον ποιήσετε τοπρωῒ, ὁλοκαύτωμα διαπαντός.

\vs{16}Τάδε λέγει Κύριος Θεὸς, ἐὰν δῷ ὁ ἀφηγούμενος δόμα ἑνὶ ἐκ τῶν υἱῶν αὐτοῦ ἐκ τῆς κληρονομίας αὐτοῦ, τοῦτο τοῖς υἱοῖς αὐτοῦ ἔσται κατάσχεσις κληρονομία·
\vs{17}Ἐὰν δὲ δῷ δόμα ἑνὶ τῶν παίδων αὐτοῦ, καὶ ἔσται αὐτῷ ἕως τοῦ ἔτους τῆς ἀφέσεως, καὶ ἀποδώσει τῷ ἀφηγουμένῳ· πλὴν τῆς κληρονομίας τῶν υἱῶν αὐτοῦ αὐτοῖς ἔσται.
\vs{18}Καὶ οὐ μὴ λάβῃ ὁ ἀφηγούμενος ἐκ τῆς κληρονομίας τοῦ λαοῦ, καταδυναστεῦσαι αὐτοὺς, ἐκ τῆς κατασχέσεως αὐτοῦ κατακληρονομήσει τοῖς υἱοῖς αὐτοῦ, ὅπως μὴ διασκορπίζηται ὁ λαός μου, ἕκαστος ἐκ τῆς κατασχέσεως αὐτοῦ.

\vs{19}Καὶ εἰσήγαγέ με εἰς τὴν εἴσοδον τῆς κατὰ νώτου τῆς πύλης, εἰς τὴν ἐξέδραν τῶν ἁγίων τῶν ἱερέων, τὴν βλέπουσαν πρὸς Βοῤῥᾶν· καὶ ἰδοὺ ἐκεῖ τόπος κεχωρισμένος.
\vs{20}Καὶ εἶπε πρὸς μὲ, οὗτος ὁ τόπος ἐστὶν, οὗ ἑψήσουσιν ἐκεῖ οἱ ἱερεῖς τὰ ὑπὲρ ἀγνοίας καὶ τὰ ὑπὲρ ἁμαρτίας, καὶ ἐκεῖ πέψουσι τὸ μαναὰ τὸ παράπαν, τοῦ μὴ ἐκφέρειν εἰς τὴν αὐλὴν τὴν ἐξωτέραν, τοῦ ἁγιάζειν τὸν λαόν.

\vs{21}Καὶ ἐξήγαγέ με εἰς τὴν αὐλὴν τὴν ἐξῶτέραν, καὶ περιήγαγέ με ἐπὶ τὰ τέσσαρα μέρη τῆς αὐλῆς· καὶ ἰδοὺ αὐλὴ κατὰ τὰ κλίτη τῆς αὐλῆς,
\vs{22}κατὰ τὸ κλίτος αὐλὴ, αὐλὴ ἐπὶ τὰ τέσσαρα, καὶ τῆς αὐλῆς αὐλὴ μικρὰ μήκους πηχῶν τεσσαράκοντα, καὶ εὖρος πηχῶν τριάκοντα, μέτρον ἓν ταῖς τέσσαρσι.
\vs{23}Καὶ ἐξέδραι κύκλῳ ἐν αὐταῖς, κυκλῳ ταῖς τέσσαρσι· καὶ μαγειρεῖα γεγονότα ὑποκάτω τῶν ἐξεδρῶν κύκλῳ.
\vs{24}Καὶ εἶπε πρὸς μὲ, οὗτοι οἱ οἶκοι τῶν μαγείρων, οὗ ἑψήσουσιν ἐκεῖ οἱ λειτουργοῦντες τῷ οἴκῳ τὰ θύματα τοῦ λαοῦ.

\ch{47}
Καὶ εἰσήγαγέ με ἐπὶ τὰ πρόθυρα τοῦ οἴκου· καὶ ἰδοὺ ὕδωρ ἐξεπορεύετο ὑποκάτωθεν τοῦ αἰθρίου κατὰ ἀνατολὰς, ὅτι τὸ πρόσωπον τοῦ οἴκου ἔβλεπε κατὰ ἀνατολὰς, καὶ τὸ ὕδωρ κατέβαινεν ἀπὸ τοῦ κλίτους τοῦ δεξιοῦ, ἀπὸ Νότου ἐπὶ τὸ θυσιαστήριον.
\vs{2}Καὶ ἐξήγαγέ με κατὰ τὴν ὁδὸν τῆς πύλης τῆς πρὸς Βοῤῥᾶν, καὶ περιήγαγέ με τὴν ὁδὸν ἔξωθεν πρὸς τὴν πύλην τῆς αὐλῆς τῆς βλεπούσης κατὰ ἀνατολάς· καὶ ἰδοὺ τὸ ὕδωρ κατεφέρετο ἀπὸ τοῦ κλίτους τοῦ δεξιοῦ,
\vs{3}καθὼς ἔξοδος ἀνδρὸς ἐξεναντίας· καὶ μέτρον ἐν τῇ χειρὶ αὐτοῦ· καὶ διεμέτρησε χιλίους ἐν τῷ μετρῳ, καὶ διῆλθεν ἐν τῷ ὕδατι ὕδωρ ἀφέσεως·
\vs{4}καὶ διεμέτρησε χιλίους, καὶ διῆλθεν ἐν τῷ ὕδατι ὕδωρ ἕως τῶν μηρῶν· καὶ διεμέτρησε χιλίους, καὶ διῆλθεν ὕδωρ ἕως ὀσφύος.
\vs{5}Καὶ διεμέτρησε χιλίους, καὶ οὐκ ἠδύνατο διελθεῖν, ὅτι ἐξύβριζεν ὡς χειμάῤῥου ὃν οὐ διαβήσονται.

\vs{6}Καὶ εἶπε πρὸς μὲ, ἑὼρακας υἱὲ ἀνθρώπου; καὶ ἤγαγέ με, καὶ ἐπέστρεψέ με ἐπι τὸ χεῖλος τοῦ ποταμοῦ,
\vs{7}ἐν τῇ ἐπιστροφῇ μου· καὶ ἰδοὺ ἐπὶ τοῦ χείλους τοῦ ποταμοῦ δένδρα πολλὰ σφόδρα ἔνθεν καὶ ἔνθεν.
\vs{8}Καὶ εἶπε πρὸν μὲ, τὸ ὕδωρ τοῦτο τὸ ἐκπορευόμενον εἰς τὴν Γαλιλαίαν τὴν πρὸς ἀνατολὰς, καὶ κατέβαινεν ἐπὶ τὴν Ἀραβίαν, καὶ ἤρχετο ἕως ἐπὶ τὴν θάλασσαν ἐπὶ τὸ ὕδωρ τῆς διεκβολῆς, καὶ ὑγιάσει τὰ ὕδατα·
\vs{9}Καὶ ἔσται πᾶσα ψυχὴ τῶν ζώων τῶν ἐκζεόντων, ἐπὶ πάντα ἐφʼ ἃ ἂν ἐπέλθῃ ἐκεῖ ὁ ποταμὸς, ζήσεται· καὶ ἔσται ἐκεῖ ἰχθὺς πολὺς σφόδρα, ὅτι ἥκει ἐκεῖ τὸ ὕδωρ τοῦτο, καὶ ὑγιάσει, καὶ ζήσεται, πᾶν ἐφʼ ὃ ἂν ἔλθῃ ὁ ποταμὸς ἐκεῖ, ζήσεται.

\vs{10}Καὶ στήσονται ἐκεῖ ἁλιεῖς ἀπὸ Ἰνγαδεὶν ἕως Ἐναγαλλείμ· ψυγμὸς σαγηνῶν ἔσται, καθʼ ἑαυτὴν ἔσται· καὶ οἱ ἰχθύες αὐτῆς, ὡς οἱ ἰχθύες τῆς θαλάσσης τῆς μεγάλης, πλῆθος πολὺ σφόδρα.
\vs{11}Καὶ ἐν τῇ διεκβολῇ αὐτοῦ, καὶ ἐν τῇ ἐπιστροφῇ αὐτοῦ, καὶ ἐν τῇ ὑπεράρσει αὐτοῦ, οὐ μὴ ὑγιάσωσιν, εἰς ἅλας δέδονται.
\vs{12}Καὶ ἐπὶ τοῦ ποταμοῦ ἀναβήσεται, ἐπὶ τοῦ χείλους αὐτοῦ ἔνθεν καὶ ἔνθεν, πᾶν ξύλον βρώσιμον, οὐ μὴ παλαιωθῇ ἐπʼ αὐτοῦ, οὐδὲ μὴ ἐκλείπῃ ὁ καρπὸς αὐτοῦ, τῆς καινότητος αὐτοῦ πρωτοβολήσει, ὅτι τὰ ὕδατα αὐτῶν ἐκ τῶν ἁγίων ταῦτα ἐκπορεύεται, καὶ ἔσται ὁ καρπὸς αὐτῶν εἰς βρῶσιν, καὶ ἀνάβασις αὐτῶν εἰς ὑγίειαν.

\vs{13}Τάδε λέγει Κύριος Θεὸς, ταῦτα τὰ ὅρια κατακληρονομήσετε τῆς γῆς, ταῖς δώδεκα φυλαῖς τῶν υἱῶν Ἰσραὴλ πρόσθεσις σχοινίσματος.
\vs{14}Καὶ κατακληρονομήσετε αὐτὴν ἕκαστος καθὼς ὁ ἀδελφὸς αὐτοῦ, εἰς ἣν ᾖρα τὴν χεῖρά μου, τοῦ δοῦναι τοῖς πατράσιν αὐτῶν, καὶ πεσεῖται ἡ γῆ αὕτη ὑμῖν ἐν κληρονομίᾳ.

\vs{15}Καὶ ταῦτα τὰ ὅρια τῆς γῆς τῆς πρὸς Βοῤῥᾶν, ἀπὸ θαλάσσης τῆς μεγάλης τῆς καταβαινούσης, καὶ περισχιζούσης τῆς εἰσόδου, Ἠμασελδὰμ,
\vs{16}Μααβθηρὰς Ἑβραμηλιὰμ ἀναμέσον ὁρίων Δαμασκοῦ καὶ ἀναμέσον ὁρίων Ἠμαθεὶ, αὐλὴ τοῦ Σαυνὰν, αἵ εἰσιν ἐπάνω τῶν ὁρίων Αὐρανίτιδος.
\vs{17}Ταῦτα τὰ ὅρια ἀπὸ τῆς θαλάσσης, ἀπὸ τῆς αὐλῆς τοῦ Αἰνὰν, ὅρια Δαμασκοῦ, καὶ τὰ πρὸς βοῤῥᾶν,
\vs{18}Καὶ τὰ πρὸς ἀνατολὰς ἀναμέσον τῆς Λωρανίτιδος, καὶ ἀναμέσον Δαμασκοῦ, καὶ ἀναμέσον τῆς Γαλααδίτιδος, καὶ ἀναμέσον τῆς γῆς τοῦ Ἰσραὴλ, ὁ Ἰορδάνης διορίζει ἐπὶ τὴν θάλασσαν, τὴν πρὸς ἀνατολὰς φοινικῶνος· ταῦτα τὰ πρὸς ἀνατολάς.
\vs{19}Καὶ τὰ πρὸς Νότον καὶ Λίβα ἀπὸ Θαιμὰν καὶ φοινικῶνος, ἕως ὕδατος Μαριμὼθ Καδὴμ, παρεκτεῖνον ἐπὶ τὴν θάλασσαν τὴν μεγάλην·
\vs{20}τοῦτο τὸ μέρος Νότος καὶ λίψ. Τοῦτο τὸ μέρος τῆς θαλάσσης τῆς μεγάλης ὁρίζει, ἕως κατέναντι τῆς εἰσόδου Ἠμὰθ, ἕως εἰσόδου αὐτοῦ· ταῦτά ἐστι τὰ πρὸς θάλασσαν Ἠμάθ.

\vs{21}Καὶ διαμερίσετε τὴν γῆν ταύτην αὐτοῖς, ταῖς φυλαῖς τοῦ Ἰσραήλ.
\vs{22}Βαλεῖτε αὐτὴν ἐν κλήρῳ, ὑμῖν καὶ τοῖς προσηλύτοις τοῖς παροικοῦσιν ἐν μέσῳ ὑμῶν, οἵτινες ἐγέννησαν υἱοὺς ἐν μέσῳ ὑμῶν, καὶ ἔσονται ὑμῶν ὡς αὐτόχθονες ἐν τοῖς υἱοῖς τοῦ Ἰσραήλ· μεθʼ ὑμῶν φάγονται ἐν κληρονομίᾳ ἐν μέσῳ τῶν φυλῶν τοῦ Ἰσραὴλ,
\vs{23}καὶ ἔσονται ἐν φυλῇ προσηλύτων ἐν τοῖς προσηλύτοις τοῖς μετʼ αὐτῶν· ἐκεῖ δώσετε κληρονομίαν αὐτοῖς, λέγει Κύριος Θεός.

\ch{48}
Καὶ ταῦτα τὰ ὀνόματα τῶν φυλῶν ἀπὸ τῆς ἀρχῆς τῆς πρὸς Βοῤῥᾶν, κατὰ τὸ μέρος τῆς καταβάσεως τοῦ περισχίζοντος ἐπὶ τὴν εἴσοδον τῆς Ἠμὰθ αὐλῆς τοῦ Αἰλὰμ, ὅριον Δαμασκοῦ πρὸς Βοῤῥᾶν κατὰ μέρος Ἠμὰθ αὐλῆς· καὶ ἔσται αὐτοῖς τὰ πρὸς ἀνατολὰς ἕως πρὸς θάλασσαν, Δὰν, μία.
\vs{2}Καὶ ἀπὸ τῶν ὁρίων τοῦ Δὰν τὰ πρὸς ἀνατολὰς, ἕως τῶν πρὸς θάλασσαν, Ἀσσὴρ, μία.
\vs{3}Καὶ ἀπὸ τῶν ὁρίων Ἀσσὴρ, ἀπὸ τῶν πρὸς ἀνατολὰς, ἕως τῶν πρὸς θάλασσαν, Νεφθαλεὶμ, μία.
\vs{4}Καὶ ἀπὸ τῶν ὁρίων Νεφθαλεὶμ, ἀπὸ ἀνατολῶν, ἕως τῶν πρὸς θάλασσαν, Μανασσῆ, μία·
\vs{5}Καὶ ἀπὸ τῶν ὁρίων Μανασσῆ, ἀπὸ τῶν πρὸς ἀνατολὰς, ἕως τῶν πρὸς θάλασσαν, Ἐφραίμ, μία.
\vs{6}Καὶ ἀπὸ τῶν ὁρίων Ἐφραὶμ, ἀπὸ τῶν πρὸς ἀνατολὰς, ἕως τῶν πρὸς θάλασσαν, Ῥουβὴν, μία.
\vs{7}Καὶ ἀπὸ τῶν ὁρίων Ῥουβὴν ἀπὸ τῶν πρὸς ἀνατολὰς, ἕως τῶν πρὸς θάλασσαν, Ἰούδα, μία.

\vs{8}Καὶ ἀπὸ τῶν ὁρίων Ἰούδα, ἀπὸ τῶν πρὸς ἀνατολὰς, ἔσται ἡ ἀπαρχὴ τοῦ ἀφορισμοῦ πέντε καὶ εἴκοσι χιλιάδες εὖρος, καὶ μῆκος, καθὼς μία τῶν μερίδων ἀπὸ τῶν πρὸς ἀνατολὰς, καὶ ἕως τῶν πρὸς θάλασσαν· καὶ ἔσται τὸ ἅγιον ἐν μέσῳ αὐτῶν.
\vs{9}Ἀπαρχὴν, ἣν ἀφοριοῦσι τῷ Κυρίῳ, μῆκος πέντε καὶ εἴκοσι χιλιάδες, καὶ εὖρος εἴκοσι πέντε χιλιάδες.
\vs{10}Τούτων ἔσται ἡ ἀπαρχὴ τῶν ἁγίων τοῖς ἱερεῦσιν πρὸς Βοῤῥᾶν, πέντε καὶ εἴκοσι χιλιάδες· καὶ πρὸς θάλασσαν, δέκα χιλιάδες· καὶ πρὸς Νότον, εἴκοσι καὶ πέντε χιλιάδες· καὶ τὸ ὄρος τῶν ἁγίων ἔσται ἐν μέσῳ αὐτοῦ
\vs{11}τοῖς ἱερεῦσι, τοῖς ἡγιασμένοις υἱοῖς Σαδδοὺκ, τοῖς φυλάσσουσι τὰς φυλακὰς τοῦ οἴκου, οἵτινες οὐκ ἐπλανήθησαν ἐν τῇ πλανήσει υἱῶν Ἰσραὴλ, ὃν τρόπον ἐπλανήθησαν οἱ Λευῖται.
\vs{12}Καὶ ἔσται αὐτοῖς ἡ ἀπαρχὴ δεδομένη ἐκ τῶν ἀπαρχῶν τῆς γῆς, ἅγιον ἁγίων ἀπὸ τῶν ὁρίων τῶν Λευειτῶν.

\vs{13}Τοῖς δὲ Λευίταις τὰ ἐχόμενα τῶν ὁρίων τῶν ἱερέων, μῆκος πέντε καὶ εἴκοσι χιλιάδες, καὶ εὖρος δέκα χιλιάδες· πᾶν τὸ μῆκος πέντε καὶ εἴκοσι χιλιάδες, καὶ εὖρος εἴκοσι χιλιάδες.
\vs{14}Οὐ πραθήσεται ἐξ αὐτοῦ, οὐδὲ καταμετρηθήσεται, οὐδὲ ἀφαιρεθήσεται τὰ πρωτογεννήματα τῆς γῆς, ὅτι ἅγιόν ἐστι τῷ Κυρίῳ.

\vs{15}Τὰς δὲ πέντε χιλιάδας τὰς περισσὰς ἐπὶ τῷ πλάτει ἐπὶ ταῖς πέντε καὶ εἴκοσι χιλιάσι, προτείχισμα ἔσται τῇ πόλει εἰς τὴν κατοικίαν, καὶ εἰς διάστημα αὐτοῦ· καὶ ἔσται ἡ πόλις ἐν μέσῳ αὐτοῦ·
\vs{16}Καὶ ταῦτα τὰ μέτρα αὐτῆς· ἀπὸ τῶν πρὸς Βοῤῥᾶν, πεντακόσιοι καὶ τετρακισχίλιοι, καὶ ἀπὸ τῶν πρὸς Νότον πεντακόσιοι καὶ τέσσαρες χιλιάδες, καὶ ἀπὸ τῶν πρὸς ἀνατολὰς, πεντακόσιοι καὶ τέσσαρες χιλιάδες, καὶ ἀπὸ τῶν πρὸς θάλασσαν, τετρακισχιλίους πεντακοσίους.
\vs{17}Καὶ ἔσται διάστημα τῇ πόλει πρὸς Βοῤῥᾶν διακόσιοι πεντήκοντα, καὶ πρὸς Νότον διακόσιοι καὶ πεντήκοντα, καὶ πρὸς ἀνατολὰς, διακόσιοι πεντήκοντα, καὶ πρὸς θάλασσαν διακόσιοι πεντήκοντα.

\vs{18}Καὶ τὸ περισσὸν τοῦ μηκους τὸ ἐχόμενον τῶν ἀπαρχῶν τῶν ἁγίων, δέκα χιλιάδες πρὸς ἀνατολὰς, καὶ δέκα χιλιάδες πρὸς θάλασσαν· καὶ ἔσονται αἱ ἀπαρχαὶ τοῦ ἁγίου, καὶ ἔσται τὰ γεννήματα αὐτῆς εἰς ἄρτους τοῖς ἐργαζομένοις τὴν πόλιν.
\vs{19}Οἱ δὲ ἐργαζόμενοι τὴν πόλιν ἐργῶνται αὐτὴν ἐκ πασῶν τῶν φυλῶν τοῦ Ἰσραήλ.

\vs{20}Πᾶσα ἡ ἀπαρχὴ, πέντε καὶ εἴκοσι χιλιάδες ἐπὶ πέντε καὶ εἴκοσι χιλιάδας τετράγωνον· ἀφοριεῖτε αὐτοῦ τὴν ἀπαρχὴν τοῦ ἁγίου, ἀπὸ τῆς κατασχέσεως τῆς πόλεως.

\vs{21}Τὸ δὲ περισσὸν τῷ ἀφηγουμένῳ ἐκ τούτου καὶ ἐκ τούτου ἀπὸ τῶν ἀπαρχῶν τοῦ ἁγίου, καὶ εἰς τὴν κατάσχεσιν τῆς πόλεως, ἐπὶ πέντε καὶ εἴκοσι χιλιάδας μῆκος, ἕως τῶν ὁρίων τῶν πρὸς ἀνατολὰς καὶ πρὸς θάλασσαν, ἐπὶ πέντε καὶ εἴκοσι χιλιάδας ἕως τῶν ὁρίων τῶν πρὸς θάλασσαν, ἐχόμενα τῶν μερίδων τοῦ ἀφηγουμένου· καὶ ἔσται ἡ ἀπαρχὴ τῶν ἁγίων καὶ τὸ ἁγίασμα τοῦ οἴκου ἐν μέσῳ αὐτῆς.
\vs{22}Καὶ παρὰ τῶν Λευιτῶν, ἀπὸ τῆς κατασχέσεως τῆς πόλεως ἐν μέσῳ τῶν ἀφηγουμένων ἔσται ἀναμέσον τῶν ὁρίων Ἰούδα, καὶ ἀναμέσον τῶν ὁρίων Βενιαμὶν, καὶ τῶν ἀφηγουμένων ἔσται.

\vs{23}Καὶ τὸ περισσὸν τῶν φυλῶν, ἀπὸ τῶν πρὸς ἀνατολὰς, ἕως τῶν πρὸς θάλασσαν, Βενιαμὶν, μία.
\vs{24}Καὶ ἀπὸ τῶν ὁρίων τῶν Βενιαμὶν, ἀπὸ τῶν πρὸς ἀνατολὰς, ἕως τῶν πρὸς θάλα σσαν, Συμεὼν, μία.
\vs{25}Καὶ ἀπὸ τῶν ὁρίων τῶν Συμεὼν, ἀπὸ τῶν πρὸς ἀνατολὰς, ἕως τῶν πρὸς θάλασσαν, Ἰσσάχαρ, μία.
\vs{26}Καὶ ἀπὸ τῶν ὁρίων τῶν Ἰσσάχαρ, ἀπὸ τῶν πρὸς ἀνατολὰς, ἕως τῶν πρὸς θάλασσαν, Ζαβουλὼν, μία.
\vs{27}Καὶ ἀπὸ τῶν ὁρίων τῶν Ζαβουλὼν, ἀπὸ τῶν πρὸς ἀνατολὰς, ἕως τῶν πρὸς θάλασσαν, Γὰδ, μία,
\vs{28}καὶ ἀπὸ τῶν ὁρίων τῶν Γὰδ, ἀπὸ τῶν πρὸς ἀνατολὰς, ἕως τῶν πρὸς Λίβα· καὶ ἔσται τὰ ὅρια αὐτοῦ ἀπὸ Θαιμὰν, καὶ ὕδατος βαριμὼθ Κάδης, κληρονομίας, ἕως τῆς θαλάσσης τῆς μεγάλης.
\vs{29}Αὕτη ἡ γῆ, ἣν βαλεῖτε ἐν κλήρῳ ταῖς φυλαῖς τοῦ Ἰσραήλ· καὶ οὗτοι οἱδιαμερισμοὶ αὐτῶν, λέγει Κύριος Θεός.

\vs{30}Καὶ αὗται αἱ διεκβολαὶ τῆς πόλεως αἱ πρὸς Βοῤῥᾶν, τετρακισχίλιοι καὶ πεντακόσιοι μέτρῳ.
\vs{31}Καὶ αἱ πύλαι τῆς πόλεως, ἐπʼ ὀνόμασι φυλῶν τοῦ Ἰσραήλ· πύλαι τρεῖς πρὸς Βοῤῥᾶν, πύλη Ῥουβὴν μία, καὶ πύλη Ἰούδα, μία, καὶ πύλη Λευῒ, μία.
\vs{32}Καὶ τὰ πρὸς ἀνατολὰς τετρακισχίλιοι καὶ πεντακόσιοι, καὶ πύλαι τρεῖς, πύλη Ἰωσὴφ, μία, καὶ πύλη Βενιαμεὶν, μία, καὶ πύλη Δὰν, μία.
\vs{33}Καὶ τὰ πρὸς Νότον τετρακισχίλιοι καὶ πεντακόσιοι μέτρῳ· καὶ πύλαι τρεῖς, πύλη Συμεὼν, μία, καὶ πύλη Ἰσσάχαρ, μία, καὶ πύλη Ζαβουλὼν, μία.
\vs{34}Καὶ τὰ πρὸς θάλασσαν τετρακισχίλιοι καὶ πεντακόσιοι μέτρῳ· πύλαι τρεῖς, πύλη Γὰδ, μία, καὶ πύλη Ἀσσὴρ, μία, καὶ πύλη Νεφθαλεὶμ, μία.

\vs{35}Κύκλωμα, δέκα καὶ ὀκτὼ χιλιάδες· καὶ τὸ ὄνομα τῆς πόλεως, ἀφʼ ἧς ἂν ἡμέρας γένηται, ἔσται τὸ ὄνομα αὐτῆς.


\def\book{ΔΑΝΙΗΛ}
\biblebook{ΔΑΝΙΗΛ}


\lettrine[lines=2, loversize=0.2, nindent=0em, findent=.25em]{\textcolor{bookheadingcolor}{Ἐ}}{Ν} ἔτει τρίτῳ τῆς βασιλείας Ἰωακεὶμ βασιλέως Ἰούδα, ἦλθε Ναβουχοδονόσορ ὁ βασιλεὺς Βαβυλῶνος εἰς Ἱερουσαλὴμ, καὶ ἐπολιόρκει αὐτήν.
\vs{2}Καὶ ἔδωκε Κύριος ἐν χειρὶ αὐτοῦ τὸν Ἰωακεὶμ βασιλέα Ἰούδα, καὶ ἀπὸ μέρους τῶν σκευῶν οἴκου τοῦ Θεοῦ· καὶ ἤνεγκεν αὐτὰ εἰς γῆν Σενναὰρ οἴκου τοῦ θεοῦ αὐτοῦ, καὶ τὰ σκεύη εἰσήνεγκεν εἰς τὸν οἶκον θησαυροῦ τοῦ θεοῦ αὐτοῦ.
\vs{3}Καὶ εἶπεν ὁ βασιλεὺς τῷ Ἀσφανὲζ τῷ ἀρχιευνούχῳ αὐτοῦ, εἰσαγαγεῖν ἀπὸ τῶν υἱῶν τῆς αἰχμαλωσίας Ἰσραὴλ, καὶ ἀπὸ τοῦ σπέρματος τῆς βασιλείας, καὶ ἀπὸ τῶν φορθομμὶν,
\vs{4}νεανίσκους, οἷς οὐκ ἔστιν ἐν αὐτοῖς μῶμος, καὶ καλοὺς τῇ ὄψει, καὶ συνιέντας ἐν πάσῃ σοφίᾳ, καὶ γινώσκοντας γνῶσιν, καὶ διανοουμένους φρόνησιν, καὶ οἷς ἐστιν ἰσχὺς ἐν αὐτοῖς ἑστάναι ἐν τῷ οἴκῳ ἐνώπιον τοῦ βασιλέως, καὶ διδάξαι αὐτοὺς γράμματα καὶ γλῶσσαν Χαλδαίων.

\vs{5}Καὶ διέταξεν αὐτοῖς ὁ βασιλεὺς τὸ τῆς ἡμέρας καθʼ ἡμέραν, ἀπὸ τῆς τραπέζης τοῦ βασιλέως, καὶ ἀπὸ τοῦ οἴνου τοῦ ποτοῦ αὐτοῦ, καὶ θρέψαι αὐτοὺς ἔτη τρία, καὶ μετὰ ταῦτα στῆναι ἐνώπιον τοῦ βασιλέως.

\vs{6}Καὶ ἐγένετο ἐν αὐτοῖς ἐκ τῶν υἱῶν Ἰούδα, Δανιὴλ, καὶ Ἀνανίας, καὶ Ἀζαρίας, καὶ Μισαήλ.
\vs{7}Καὶ ἐπέθηκεν αὐτοῖς ὁ ἀρχιευνοῦχος ὀνόματα· τῷ Δανιὴλ Βαλτάσαρ, καὶ τῷ Ἀνανίᾳ Σεδρὰχ, καὶ τῷ Μισαὴλ Μισὰχ, καὶ τῷ Ἀζαρίᾳ Ἀβδεναγώ.
\vs{8}Καὶ ἔθετο Δανιὴλ εἰς τὴν καρδίαν αὐτοῦ, ὡς οὐ μὴ ἁλισγηθῇ ἐν τῇ τραπέζῃ τοῦ βασιλέως, καὶ ἐν τῷ οἴνῳ ἀπὸ τοῦ ποτοῦ αὐτοῦ· καὶ ἠξίωσε τὸν ἀρχιευνοῦχον, ὡς οὐ μὴ ἁλισγηθῇ.
\vs{9}Καὶ ἔδωκεν ὁ Θεὸς τὸν Δανιὴλ εἰς ἔλεον καὶ οἰκτιρμὸν ἐνώπιον τοῦ ἀρχιευνούχου.
\vs{10}Καὶ εἶπεν ὁ ἀρχιευνοῦχος τῷ Δανιὴλ, φοβοῦμαι ἐγὼ τὸν κύριόν μου τὸν βασιλέα, τὸν ἐκτάξαντα τὴν βρῶσιν ὑμῶν καὶ τὴν πόσιν ὑμῶν, μήποτε ἴδῃ τὰ πρόσωπα ὑμῶν σκυθρωπὰ παρὰ τὰ παιδάρια τὰ συνήλικα ὑμῶν, καὶ καταδικάσητε τὴν κεφαλήν μου τῷ βασιλεῖ.
\vs{11}Καὶ εἶπε Δανιὴλ πρὸς Ἀμελσὰδ, ὃν κατέστησεν ὁ ἀρχιευνοῦχος ἐπὶ Δανιὴλ, Ἀνανίαν, Μισαὴλ, Ἀζαρίαν,
\vs{12}πείρασον δὴ τοὺς παῖδάς σου ἡμέρας δέκα, καὶ δότωσαν ἡμῖν ἀπὸ τῶν σπερμάτων, καὶ φαγώμεθα, καὶ ὕδωρ πιώμεθα,
\vs{13}καὶ ὀφθήτωσαν ἐνώπιόν σου αἱ ἰδέαι ἡμῶν, καὶ αἱ ἰδέαι τῶν παιδαρίων τῶν ἐσθόντων τὴν τράπεζαν τοῦ βασιλέως, καὶ καθὼς ἐὰν ἴδῃς, ποίησον μετὰ τῶν παίδων σου.

\vs{14}Καὶ εἰσήκουσεν αὐτῶν, καὶ ἐπείρασεν αὐτοὺς ἡμέρας δέκα.
\vs{15}Καὶ μετὰ τὸ τέλος τῶν δέκα ἡμερῶν, ὡράθησαν αἱ ἰδέαι αὐτῶν ἀγαθαὶ καὶ ἰσχυραὶ ταῖς σαρξὶν ὑπὲρ τὰ παιδάρια τὰ ἔσθοντα τὴν τράπεζαν τοῦ βασιλέως.
\vs{16}Καὶ ἐγένετο Ἀμελσὰδ ἀναιρούμενος τὸ δεῖπνον αὐτῶν, καὶ τὸν οἶνον τοῦ πόματος αὐτῶν, καὶ ἐδίδου αὐτοῖς σπέρματα.

\vs{17}Καὶ τὰ παιδάρια ταῦτα οἱ τέσσαρες αὐτοὶ, ἔδωκεν αὐτοῖς ὁ Θεὸς σύνεσιν καὶ φρόνησιν ἐν πάσῃ γραμματικῇ καὶ σοφίᾳ· καὶ Δανιὴλ συνῆκεν ἐν πάσῃ ὁράσει καὶ ἐνυπνίοις.
\vs{18}Καὶ μετὰ τὸ τέλος τῶν ἡμερῶν, ὧν εἶπεν ὁ βασιλεὺς εἰσαγαγεῖν αὐτοὺς, καὶ εἰσήγαγεν αὐτοὺς ὁ ἀρχιευνοῦχος ἐναντίον Ναβουχοδονόσορ.
\vs{19}Καὶ ἐλάλησε μετʼ αὐτῶν ὁ βασιλεύς· καὶ οὐχ εὑρέθησαν ἐκ πάντων αὐτῶν ὅμοιοι Δανιὴλ, καὶ Ἀνανίᾳ, καὶ Μισαὴλ, καὶ Ἀζαρίᾳ· καὶ ἔστησαν ἐνώπιον τοῦ βασιλέως.
\vs{20}Καὶ ἐν παντὶ ῥήματι σοφίας καὶ ἐπιστήμης ὧν ἐζήτησε παρʼ αὐτῶν ὁ βασιλεύς, εὗρεν αὐτοὺς δεκαπλασίονας παρὰ πάντας τοὺς ἐπαοιδοὺς καὶ τοὺς μάγους τοὺς ὄντας ἐν πάσῃ τῇ βασιλείᾳ αὐτοῦ.
\vs{21}Καὶ ἐγένετο Δανιὴλ ἕως ἔτους ἑνὸς Κύρου τοῦ βασιλέως.

\ch{2}
Ἐν τῷ ἔτει τῷ δευτέρῳ τῆς βασιλείας, ἐνυπνιάσθη Ναβουχοδονόσορ ἐνύπνιον, καὶ ἐξέστη τὸ πνεῦμα αὐτοῦ, καὶ ὁ ὕπνος αὐτοῦ ἐγένετο ἀπʼ αὐτοῦ.
\vs{2}Καὶ εἶπεν ὁ βασιλεὺς καλέσαι τοὺς ἐπαοιδοὺς, καὶ τοὺς μάγους, καὶ τοὺς φαρμακοὺς, καὶ τοὺς Χαλδαίους, τοῦ ἀναγγεῖλαι τῷ βασιλεῖ τὰ ἐνύπνια αὐτοῦ· καὶ ἦλθαν, καὶ ἔστησαν ἐνώπιον τοῦ βασιλέως.

\vs{3}Καὶ εἶπεν αὐτοῖς ὁ βασιλεὺς, ἑνυπνιάσθην, καὶ ἐξέστη τὸ πνεῦμά μου, τοῦ γνῶναι τὸ ἐνύπνιον.
\vs{4}Καὶ ἐλάλησαν οἱ Χαλδαῖοι τῷ βασιλεῖ Συριστὶ, βασιλεῦ, εἰς τοὺς αἰῶνας ζῆθι· σὺ εἰπὸν τὸ ἐνύπνιον τοῖς παισί σου, καὶ τὴν σύγκρισιν ἀναγγελοῦμεν.
\vs{5}Ἀπεκρίθη ὁ βασιλεὺς τοῖς Χαλδαίοις, ὁ λόγος ἀπʼ ἐμοῦ ἀπέστη· ἐὰν μὴ γνωρίσητέ μοι τὸ ἐνύπνιον καὶ τὴν σύγκρισιν, εἰς ἀπώλειαν ἔσεσθε, καὶ οἱ οἶκοι ὑμῶν διαρπαγήσονται.
\vs{6}Ἐὰν δὲ τὸ ἐνύπνιον καὶ τὴν σύγκρισιν αὐτοῦ γνωρίσητέ μοι, δόματα καὶ δωρεὰς καὶ τιμὴν πολλὴν λήψεσθε παρʼ ἐμοῦ· πλὴν τὸ ἐνύπνιον καὶ τὴν σύγκρισιν αὐτοῦ ἀπαγγείλατέ μοι.
\vs{7}Ἀπεκρίθησαν δεύτερον, καὶ εἶπαν, ὁ βασιλεὺς εἰπάτω τὸ ἐνύπνιον τοῖς παισὶν αὐτοῦ, καὶ τὴν σύγκρισιν ἀναγγελοῦμεν.

\vs{8}Καὶ ἀπεκρίθη ὁ βασιλεὺς, καὶ εἶπεν, ἐπʼ ἀληθείας οἶδα ἐγὼ, ὅτι καιρὸν ὑμεῖς ἐξαγοράζετε· καθότι ἴδετε, ὅτι ἀπέστη ἀπʼ ἐμοῦ τὸ ῥῆμα.
\vs{9}Ἐὰν οὖν τὸ ἐνύπνιον μὴ ἀναγγείλητέ μοι, οἶδα ὅτι ῥῆμα ψευδὲς καὶ διεφθαρμένον συνέθεσθε εἰπεῖν ἐνώπιόν μου, ἕως οὗ ὁ καιρὸς παρέλθῃ· τὸ ἐνύπνιόν μου εἴπατέ μοι, καὶ γνώσομαι ὅτι καὶ τὴν σύγκρισιν αὐτοῦ ἀναγγελεῖτέ μοι.
\vs{10}Ἀπεκρίθησαν οἱ Χαλδαῖοι ἐνώπιον τοῦ βασιλέως, καὶ λέγουσιν, οὐκ ἔστιν ἄνθρωπος ἐπὶ τῆς ξηρᾶς, ὅστις τὸ ῥῆμα τοῦ βασιλέως δυνήσεται γνωρίσαι, καθότι πᾶς βασιλεὺς μέγας καὶ ἄρχων ῥῆμα τοιοῦτον οὐκ ἐπερωτᾷ ἐπαοιδὸν μάγον καὶ Χαλδαῖον·
\vs{11}Ὅτι ὁ λόγος ὃν ὁ βασιλεὺς ἐπερωτᾷ, βαρύς, καὶ ἕτερος οὐκ ἔστιν ὃς ἀναγγελεῖ αὐτὸν ἐνώπιον τοῦ βασιλέως, ἀλλʼ οἱ θεοὶ, ὧν οὐκ ἔστιν ἡ κατοικία μετὰ πάσης σαρκός.

\vs{12}Τότε ὁ βασιλεὺς ἐν θυμῷ καὶ ὀργῇ εἶπεν ἀπολέσαι πάντας τοὺς σοφοὺς Βαβυλῶνος.
\vs{13}Καὶ τὸ δόγμα ἐξῆλθε, καὶ οἱ σοφοὶ ἀπεκτέννοντο· καὶ ἐζήτησαν Δανιὴλ καὶ τοὺς φίλους αὐτοῦ ἀνελεῖν.

\vs{14}Τότε Δανιὴλ ἀπεκρίθη βουλὴν καὶ γνώμην τῷ Ἀριὼχ τῷ ἀρχιμαγείρῳ τοῦ βασιλέως, ὃς ἐξῆλθεν ἀναιρεῖν τοὺς σοφοὺς Βαβυλῶνος,
\vs{15}ἄρχων τοῦ βασιλέως, περὶ τίνος ἐξῆλθεν ἡ γνώμη ἡ ἀναιδὴς ἐκ προσώπου τοῦ βασιλέως; ἐγνώρισε δὲ ὁ Ἀριὼχ τὸ ῥῆμα τῷ Δανιήλ.
\vs{16}Καὶ Δανιὴλ ἠξίωσε τὸν βασιλέα ὅπως χρόνον δῷ αὐτῷ, καὶ τὴν σύγκρισιν αὐτοῦ ἀναγγελῇ τῷ βασιλεῖ.
\vs{17}Καὶ εἰσῆλθε Δανιὴλ εἰς τὸν οἶκον αὐτοῦ καὶ τῷ Ἀνανίᾳ καὶ τῷ Μισαὴλ καὶ τῷ Ἀζαρίᾳ τοῖς φίλοις αὐτοῦ τὸ ῥῆμα ἐγνώρισε.
\vs{18}Καὶ οἰκτιρμοὺς ἐζήτουν παρὰ τοῦ Θεοῦ τοῦ οὐρανοῦ ὑπὲρ τοῦ μυστηρίου τούτου, ὅπως ἂν μὴ ἀπόλωνται Δανιὴλ καὶ οἱ φίλοι αὐτοῦ μετὰ τῶν ἐπιλοίπων σοφῶν Βαβυλῶνος.

\vs{19}Τότε τῷ Δανιὴλ ἐν ὁράματι τῆς νυκτὸς τὸ μυστήριον ἀπεκαλύφθη· καὶ εὐλόγησε τὸν Θεὸν τοῦ οὐρανοῦ Δανιὴλ, καὶ εἶπεν,

\vs{20}Εἴη τὸ ὄνομα τοῦ Θεοῦ εὐλογημένον ἀπὸ τοῦ αἰῶνος καὶ ἕως τοῦ αἰῶνος, ὅτι ἡ σοφία καὶ ἡ σύνεσις αὐτοῦ ἐστι.
\vs{21}Καὶ αὐτὸς ἀλλοιοῖ καιροὺς καὶ χρόνους, καθιστᾷ βασιλεῖς, καὶ μεθιστᾷ, διδοὺς σοφίαν τοῖς σοφοῖς, καὶ φρόνησιν τοῖς εἰδόσι σύνεσιν,
\vs{22}αὐτὸς ἀποκαλύπτει βαθέα καὶ ἀπόκρυφα, γινώσκων τὰ ἐν τῷ σκότει, καὶ τὸ φῶς μετʼ αὐτοῦ ἐστι·
\vs{23}Σοὶ ὁ Θεὸς τῶν πατέρων μου ἐξομολογοῦμαι καὶ αἰνῶ, ὅτι σοφίαν καὶ δύναμιν δέδωκάς μοι, καὶ ἐγνώρισάς μοι ἃ ἠξιώσαμεν παρὰ σοῦ, καὶ τὸ ὅραμα τοῦ βασιλέως ἐγνώρισάς μοι.

\vs{24}Καὶ ἦλθε Δανιὴλ πρὸς Ἀριὼχ, ὃν κατέστησεν ὁ βασιλεὺς ἀπολέσαι τοὺς σοφοὺς Βαβυλῶνος, καὶ εἶπεν αὐτῷ, τοὺς σοφοὺς Βαβυλῶνος μὴ ἀπολέσῃς, εἰσάγαγε δέ με ἐνώπιον τοῦ βασιλέως, καὶ τὴν σύγκρισιν τῷ βασιλεῖ ἀναγγελῶ.
\vs{25}Τότε Ἀριὼχ ἐν σπουδῇ εἰσήγαγε τὸν Δανιὴλ ἐνώπιον τοῦ βασιλέως, καὶ εἶπεν αὐτῷ, εὕρηκα ἄνδρα ἐκ τῶν υἱῶν τῆς αἰχμαλωσίας τῆς Ἰουδαίας, ὅστις τὸ σύγκριμα τῷ βασιλεῖ ἀναγγελεῖ.
\vs{26}Καὶ ἀπεκρίθη ὁ βασιλεὺς, καὶ εἶπε τῷ Δανιὴλ, οὗ τὸ ὄνομα Βαλτάσαρ, εἰ δύνασαί μοι ἀναγγεῖλαι τὸ ἐνύπνιον ὃ ἴδον, καὶ τὴν σύγκρισιν αὐτοῦ;

\vs{27}Καὶ ἀπεκρίθη Δανιὴλ ἐνώπιον τοῦ βασιλέως, καὶ εἶπε, τὸ μυστήριον ὃ ὁ βασιλεὺς ἐπερωτᾷ, οὐκ ἔστι σοφῶν, μάγων, ἐπαοιδῶν, γαζαρηνῶν ἀναγγεῖλαι τῷ βασιλεῖ·
\vs{28}Ἀλλʼ ἤ ἐστι Θεὸς ἐν οὐρανῷ ἀποκαλύπτων μυστήρια, καὶ ἐγνώρισε τῷ βασιλεῖ Ναβουχοδονόσορ, ἃ δεῖ γενέσθαι ἐπʼ ἐσχάτων τῶν ἡμερῶν· τὸ ἐνύπνιόν σου καὶ αἱ ὁράσεις τῆς κεφαλῆς σου ἐπὶ τῆς κοίτης σου, τοῦτό ἐστι,
\vs{29}βασιλεῦ· οἱ διαλογισμοί σου ἐπὶ τῆς κοίτης σου ἀνέβησαν τί δεῖ γενέσθαι μετὰ ταῦτα· καὶ ὁ ἀποκαλύπτων μυστήρια ἐγνώρισέ σοι ἃ δεῖ γενέσθαι.
\vs{30}Καὶ ἐμοὶ δὲ οὐκ ἐν σοφίᾳ τῇ οὔσῃ ἐν ἐμοὶ παρὰ πάντας τοὺς ζῶντας τὸ μυστήριον τοῦτο ἀπεκαλύφθη, ἀλλʼ ἕνεκεν τοῦ τὴν σύγκρισιν τῷ βασιλεῖ γνωρίσαι, ἵνα τοὺς διαλογισμοὺς τῆς καρδίας σου γνῷς.

\vs{31}Σὺ βασιλεῦ ἐθεώρεις, καὶ ἰδοὺ εἰκὼν μία, μεγάλη ἡ εἰκὼν ἐκείνη, καὶ ἡ πρόσοψις αὐτῆς ὑπερφερὴς, ἑστῶσα πρὸ προσώπου σου, καὶ ἡ ὅρασις αὐτῆς φοβερά.
\vs{32}Εἰκὼν, ἧς ἡ κεφαλὴ χρυσίου χρηστοῦ, αἱ χεῖρες καὶ τὸ στῆθος καὶ οἱ βραχίονες αὐτῆς ἀργυροῖ, ἡ κοιλία καὶ οἱ μηροὶ χαλκοῖ,
\vs{33}αἱ κνῆμαι σιδηραῖ, οἱ πόδες μέρος μέν τι σιδηροῦν, καὶ μέρος δέ τι ὀστράκινον.
\vs{34}Ἐθεώρεις ἕως ἀπεσχίσθη λίθος ἐξ ὄρους ἄνευ χειρῶν, καὶ ἐπάταξε τὴν εἰκόνα ἐπὶ τοὺς πόδας τοὺς σιδηροῦς καὶ ὀστρακίνους, καὶ ἐλέπτυνεν αὐτοὺς εἰς τέλος.
\vs{35}Τότε ἐλεπτύνθησαν εἰσάπαξ τὸ ὄστρακον, ὁ σίδηρος, ὁ χαλκὸς, ὁ ἄργυρος, ὁ χρυσός· καὶ ἐγένετο ὡσεὶ κονιορτὸς ἀπὸ ἅλωνος θερινῆς· καὶ ἐξῇρεν αὐτὰ τὸ πλῆθος τοῦ πνεύματος, καὶ τόπος οὐχ εὑρέθη αὐτοῖς· καὶ ὁ λίθος ὁ πατάξας τὴν εἰκόνα, ἐγενήθη ὄρος μέγα, καὶ ἐπλήρωσε πᾶσαν τὴν γῆν.
\vs{36}Τοῦτό ἐστι τὸ ἐνύπνιον, καὶ τὴν σύγκρισιν αὐτοῦ ἐροῦμεν ἐνώπιον τοῦ βασιλέως.

\vs{37}Σὺ βασιλεῦ βασιλεὺς βασιλέων, ᾧ ὁ Θεὸς τοῦ οὐρανοῦ βασιλείαν ἰσχυρὰν καὶ κραταιὰν καὶ ἔντιμον ἔδωκεν
\vs{38}ἐν παντὶ τόπῳ, ὅπου κατοικοῦσιν οἱ υἱοὶ τῶν ἀνθρώπων· θηρία τε ἀγροῦ, καὶ πετεινὰ οὐρανοῦ, καὶ ἰχθύας τῆς θαλάσσης ἔδωκεν ἐν τῇ χειρί σου, καὶ κατέστησέ σε κύριον πάντων· σὺ εἶ ἡ κεφαλὴ ἡ χρυσῆ.
\vs{39}Καὶ ὀπίσω σου ἀναστήσεται βασιλεία ἑτέρα ἥττων σου, καὶ βασιλεία τρίτη, ἥτις ἐστὶν ὁ χαλκὸς, ἣ κυριεύσει πάσης τῆς γῆς,
\vs{40}καὶ βασιλεία τετάρτη, ἥτις ἔσται ἰσχυρὰ ὡς σίδηρος· ὃν τρόπον ὁ σίδηρος λεπτύνει καὶ δαμάζει πάντα, οὕτως πάντα λεπτυνεῖ καὶ δαμάσει.
\vs{41}Καὶ ὅτι εἶδες τοὺς πόδας, καὶ τοὺς δακτύλους, μέρος μέν τι ὀστράκινον, μέρος δέ τι σιδηροῦν, βασιλεία διῃρημένη ἔσται, καὶ ἀπὸ τῆς ῥίζης τῆς σιδηρᾶς ἔσται ἐν αὐτῇ, ὃν τρόπον εἶδες τὸν σίδηρον ἀναμεμιγμένον τῷ ὀστράκῳ.
\vs{42}Καὶ οἱ δάκτυλοι τῶν ποδῶν μέρος μέν τι σιδηροῦν, μέρος δέ τι ὀστράκινον, μέρος τι τῆς βασιλείας ἔσται ἰσχυρὸν, καὶ ἀπʼ αὐτῆς ἔσται συντριβόμενον.
\vs{43}Ὅτι εἶδες τὸν σίδηρον ἀναμεμιγμένον τῷ ὀστράκῳ, συμμιγεῖς ἔσονται ἐν σπέρματι ἀνθρώπων, καὶ οὐκ ἔσονται προσκολλώμενοι οὗτος μετὰ τούτου, καθὼς ὁ σίδηρος οὐκ ἀναμίγνυται μετὰ τοῦ ὀστράκου.

\vs{44}Καὶ ἐν ταῖς ἡμέραις τῶν βασιλέων ἐκείνων, ἀναστήσει ὁ Θεὸς τοῦ οὐρανοῦ βασιλείαν, ἥτις εἰς τοὺς αἰῶνας οὐ διαφθαρήσεται, καὶ ἡ βασιλεία αὐτοῦ λαῷ ἑτέρῳ οὐχ ὑπολειφθήσεται, λεπτυνεῖ καὶ λικμήσει πάσας τὰς βασιλείας, καὶ αὕτη ἀναστήσεται εἰς τοὺς αἰῶνας·
\vs{45}Ὃν τρόπον εἶδες, ὅτι ἀπὸ ὄρους ἐτμήθη λίθος ἄνευ χειρῶν, καὶ ἐλέπτυνε τὸ ὄστρακον, τὸν σίδηρον, τὸν χαλκόν, τὸν ἄργυρον, τὸν χρυσόν· ὁ Θεὸς ὁ μέγας ἐγνώρισε τῷ βασιλεῖ ἃ δεῖ γενέσθαι μετὰ ταῦτα· καὶ ἀληθινὸν τὸ ἐνύπνιον, καὶ πιστὴ ἡ σύγκρισις αὐτοῦ.

\vs{46}Τότε ὁ βασιλεὺς Ναβουχοδονόσορ ἔπεσεν ἐπὶ πρόσωπον, καὶ τῷ Δανιὴλ προσεκύνησε, καὶ μαναὰ καὶ εὐωδίας εἶπε σπεῖσαι αὐτῷ.
\vs{47}Καὶ ἀποκριθεὶς ὁ βασιλεὺς, εἶπε τῷ Δανιήλ, ἐπʼ ἀληθείας ὁ Θεὸς ὑμῶν, αὐτός ἐστι Θεὸς θεῶν, καὶ Κύριος τῶν βασιλέων, ὁ ἀποκαλύπτων μυστήρια, ὅτι ἠδυνάσθης ἀποκαλύψαι τὸ μυστήριον τοῦτο.
\vs{48}Καὶ ἐμεγάλυνεν ὁ βασιλεὺς τὸν Δανιὴλ, καὶ δόματα μεγάλα καὶ πολλὰ ἔδωκεν αὐτῷ, καὶ κατέστησεν αὐτὸν ἐπὶ πάσης χώρας Βαβυλῶνος, καὶ ἄρχοντα σατραπῶν ἐπὶ πάντας τοὺς σοφοὺς Βαβυλῶνος.
\vs{49}Καὶ Δανιὴλ ᾐτήσατο παρὰ τοῦ βασιλέως, καὶ κατέστησεν ἐπὶ τὰ ἔργα τῆς χώρας Βαβυλῶνος τὸν Σεδρὰχ, Μισὰχ, καὶ Ἀβδεναγώ· καὶ Δανιὴλ ἦν ἐν τῇ αὐλῇ τοῦ βασιλέως.

\ch{3}
Ἔτους ὀκτωκαιδεκάτου Ναβουχοδονόσορ ὁ βασιλεὺς ἐποίησεν εἰκόνα χρυσῆν, ὕψος αὐτῆς πήχεων ἑξήκοντα, εὖρος αὐτῆς πήχεων ἕξ· καὶ ἔστησεν αὐτὴν ἐν πεδίῳ Δεειρᾷ, ἐν χώρᾳ Βαβυλῶνος.
\vs{2}Καὶ ἀπέστειλε συναγαγεῖν τοὺς ὑπάτους, καὶ τοὺς στρατηγοὺς, καὶ τοὺς τοπάρχας, ἡγουμένους, καὶ τυράννους, καὶ τοὺς ἐπʼ ἐξουσιῶν, καὶ πάντας τοὺς ἄρχοντας τῶν χωρῶν, ἐλθεῖν εἰς τὰ ἐγκαίνια τῆς εἰκόνος·
\vs{3}Καὶ συνήχθησαν οἱ τοπάρχαι, ὕπατοι, στρατηγοὶ, ἡγούμενοι, τύραννοι μεγάλοι, οἱ ἐπʼ ἐξουσιῶν, καὶ πάντες οἱ ἄρχοντες τῶν χωρῶν, εἰς τὸν ἐγκαινισμὸν τῆς εἰκόνος, ἧς ἔστησε Ναβουχοδονόσορ ὁ βασιλεύς· καὶ εἱστήκεισαν ἐνώπιον τῆς εἰκόνος.

\vs{4}Καὶ ὁ κήρυξ ἐβόα ἐν ἰσχύϊ, ὑμῖν λέγεται λαοῖς, φυλαὶ, γλῶσσαι,
\vs{5}ᾗ ἂν ὥρᾳ ἀκούσητε φωνῆς σάλπιγγος, σύριγγός τε, καὶ κιθάρας, σαμβύκης τε, καὶ ψαλτηρίου, καὶ παντὸς γένους μουσικῶν, πίπτοντες προσκυνεῖτε τῇ εἰκόνι τῇ χρυσῇ ᾗ ἔστησε Ναβουχοδονόσορ ὁ βασιλεύς.
\vs{6}Καὶ ὃς ἂν μὴ πεσὼν προσκυνήσῃ, αὐτῇ τῇ ὥρᾳ ἐμβληθήσεται εἰς τὴν κάμινον τοῦ πυρὸς τὴν καιομένην.
\vs{7}Καὶ ἐγένετο ὅταν ἤκουον οἱ λαοὶ τῆς φωνῆς τῆς σάλπιγγος, σύριγγός τε, καὶ κιθάρας, σαμβύκης τε, καὶ ψαλτηρίου, καὶ παντὸς γένους μουσικῶν, πίπτοντες πάντες οἱ λαοὶ, φυλαὶ, γλῶσσαι, προσεκύνουν τῇ εἰκόνι τῇ χρυσῇ ἣν ἔστησε Ναβουχοδονόσορ ὁ βασιλεύς.

\vs{8}Τότε προσήλθοσαν ἄνδρες Χαλδαῖοι, καὶ διέβαλον τοὺς Ἰουδαίους τῷ βασιλεῖ·
\vs{9}βασιλεῦ, εἰς τοὺς αἰῶνας ζῆθι.
\vs{10}Σὺ βασιλεῦ ἔθηκας δόγμα, πάντα ἄνθρωπον ὃς ἂν ἀκούσῃ τῆς φωνῆς τῆς σάλπιγγος, σύριγγός τε, καὶ κιθάρας, σαμβύκης, καὶ ψαλτηρίου, καὶ παντὸς γένους μουσικῶν,
\vs{11}καὶ μὴ πεσὼν προσκυνήσῃ τῇ εἰκόνι τῇ χρυσῇ, ἐμβληθήσηται εἰς τὴν κάμινον τοῦ πυρὸς τὴν καιομένην.
\vs{12}Εἰσὶν ἄνδρες Ἰουδαῖοι, οὓς κατέστησας ἐπὶ τὰ ἔργα τῆς χώρας Βαβυλῶνος, Σεδρὰχ, Μισὰχ, Ἀβδεναγὼ, οἳ οὐχ ὑπήκουσαν βασιλεῦ τῷ δόγματί σου, τοῖς θεοῖς σου οὐ λατρεύουσι, καὶ τῇ εἰκόνι τῇ χρυσῇ τᾗ ἔστησας οὐ προσκυνοῦσι.

\vs{13}Τότε Ναβουχοδονόσορ ἐν θυμῷ καὶ ὀργῇ εἶπεν ἀγαγεῖν τὸν Σεδρὰχ, Μισὰχ, καὶ Ἀβδεναγώ· καὶ ἤχθησαν ἐνώπιον τοῦ βασιλέως.
\vs{14}Καὶ ἀπεκρίθη Ναβουχοδονόσορ, καὶ εἶπεν αὐτοῖς, εἶ ἀληθῶς Σεδρὰχ, Μισὰχ, Ἀβδεναγὼ, τοῖς θεοῖς μου οὐ λατρεύετε, καὶ τῇ εἰκόνι τῇ χρυσῇ ᾗ ἔστησα οὐ προσκυνεῖτε;
\vs{15}Νῦν οὖν εἰ ἔχετε ἑτοίμως, ἵνα ὡς ἂν ἀκούσητε τῆς φωνῆς τῆς σάλπιγγος, σύριγγός τε, καὶ κιθάρας, σαμβύκης τε, καὶ ψαλτηρίου, καὶ συμφωνίας, καὶ παντὸς γένους μουσικῶν, πεσόντες προσκυνήσητε τῇ εἰκόνι τῇ χρυσῇ ᾗ ἐποίησα· ἐὰν δὲ μὴ προσκυνήσητε, αὐτῇ τῇ ὥρᾳ ἐμβληθήσεσθε εἰς τὴν κάμινον τοῦ πυρὸς τὴν καιομένην· καὶ τίς ἐστι Θεὸς ὃς ἐξελεῖται ὑμᾶς ἐκ χειρός μου;

\vs{16}Καὶ ἀπεκρίθησαν Σεδρὰχ, Μισὰχ, Ἀβδεναγὼ, λέγοντες τῷ βασιλεῖ Ναβουχοδονόσορ, οὐ χρείαν ἔχομεν ἡμεῖς περὶ τοῦ ῥήματος τούτου ἀποκριθῆναί σοι.
\vs{17}Ἔστι γὰρ Θεὸς ἡμῶν ἐν οὐρανοῖς, ᾧ ἡμεῖς λατρεύομεν, δυνατὸς ἐξελέσθαι ἡμᾶς ἐκ τῆς καμίνου τοῦ πυρὸς τῆς καιομένης, καὶ ἐκ τῶν χειρῶν σου βασιλεῦ ῥύσεται ἡμᾶς.
\vs{18}Καὶ ἐὰν μὴ, γνωστὸν ἔστω σοι, βασιλεῦ, ὅτι τοῖς θεοῖς σου οὐ λατρεύομεν, καὶ τῇ εἰκόνι ᾗ ἔστησας οὐ προσκυνοῦμεν.

\vs{19}Τότε Ναβουχοδονόσορ ἐπλήσθη θυμοῦ, καὶ ἡ ὄψις τοῦ προσώπου αὐτοῦ ἠλλοιώθη ἐπὶ Σεδρὰχ, Μισὰχ, καὶ Ἀβδεναγὼ, καὶ εἶπεν ἐκκαῦσαι τὴν κάμινον ἐπταπλασίως, ἕως οὗ εἰς τέλος ἐκκαῇ.
\vs{20}Καὶ ἄνδρας ἰσχυροὺς ἰσχύϊ εἶπε, πεδήσαντας τὸν Σεδρὰχ, Μισὰχ, καὶ Ἀβδεναγὼ, ἐμβαλεῖν εἰς τὴν κάμινον τοῦ πυρὸς τὴν καιομένην.
\vs{21}Τότε οἱ ἄνδρες ἐκεῖνοι ἐπεδήθησαν σὺν τοῖς σαραβάροις αὐτῶν, καὶ τιάραις, καὶ περικνημίσι, καὶ ἐβλήθησαν εἰς τὸ μέσον τῆς καμίνου τοῦ πυρὸς τῆς καιομένης,
\vs{22}ἐπεὶ τὸ ῥῆμα τοῦ βασιλέως ὑπερίσχυε· καὶ ἡ κάμινος ἐξεκαύθη ἐκ περισσοῦ.
\vs{23}Καὶ οἱ τρεῖς οὗτοι Σεδρὰχ, Μισὰχ, καὶ Ἀβδεναγὼ, ἔπεσον εἰς μέσον τῆς καμίνου τῆς καιομένης πεπεδημένοι,
\vs{24}καὶ περιεπάτουν ἐν μέσῳ τῆς φλογὸς, ὑμνοῦντες τὸν Θεὸν, καὶ εὐλογοῦντες τὸν Κύριον.

\vs{25}Καὶ συστὰς Ἀζαρίας προσηύξατο οὕτως· καὶ ἀνοίξας τὸ στομά αὐτοῦ ἐν μέσῳ τοῦ πυρὸς, εἶπεν,

\vs{26}Εὐλογητὸς εἶ Κύριε ὁ Θεὸς τῶν πατέρων ἡμῶν, καὶ αἰνετὸς, καὶ δεδοξασμένον τὸ ὄνομά σου εἰς τοὺς αἰῶνας.
\vs{27}Οτι δίκαιος εἶ ἐπὶ πᾶσιν οἶς ἐποίησας, καὶ πάντα τὰ ἔργα σου ἀληθινὰ, καὶ εὐθεῖαι αἱ ὁδοί σου, καὶ πᾶσαι αἱ κρίσεις σου ἀλήθεια.

\vs{28}Καὶ κρίματα ἀληθείας ἐποίησας κατὰ πάντα ἃ ἐπήγαγες ἡμῖν, καὶ ἐπὶ τὴν πόλιν τὴν ἁγίαν τὴν τῶν πατέρων ἡμῶν Ἱερουσαλήμ· ὅτι ἐν ἀληθείᾳ καὶ κρίσει ἐπήγαγες ταῦτα πάντα διὰ τὰς ἁμαρτίας ἡμῶν.
\vs{29}Ὅτι ἡμάρτομεν καὶ ἠνομήσαμεν ἀποστῆναι ἀπὸ σοῦ, καὶ ἐξημάρτομεν ἐν πᾶσι, καὶ τῶν ἐντολῶν σου οὐκ ἠκούσαμεν,
\vs{30}οὐδὲ συνετηρήσαμεν, οὐδὲ ἐποιήσαμεν καθὼς ἐνετείλω ἡμῖν, ἵνα εὖ ἡμῖν γένηται.
\vs{31}Καὶ πάντα ὅσα ἐπήγαγες ἡμῖν, καὶ πάντα ὅσα ἐποίησας ἡμῖν, ἐν ἀληθινῇ κρίσει ἐποίησας.

\vs{32}Καὶ παρέδωκας ἡμᾶς εἰς χεῖρας ἐχθρῶν ἀνόμων, καὶ ἐχθίστων ἀποστατῶν, καὶ βασιλεῖ ἀδίκῳ καὶ πονηροτάτῳ παρὰ πᾶσαν τὴν γῆν.
\vs{33}Καὶ νῦν οὐκ ἔστιν ἡμῖν ἀνοῖξαι τὸ στόμα ἡμῶν· αἰσχύνη καὶ ὄνειδος ἐγενήθημεν τοῖς δούλοις σου, καὶ τοῖς σεβομένοις σε.

\vs{34}Μὴ δὴ παραδῴης ἡμᾶς εἰς τέλος διὰ τὸ ὄνομά σου, καὶ μὴ διασκεδάσῃς τὴν διαθήκην σου,
\vs{35}καὶ μὴ ἀποστήσῃς τὸ ἔλεός σου ἀφʼ ἡμῶν, διὰ Ἁβραὰμ τὸν ἠγαπημένον ὑπὸ σοῦ, καὶ διὰ Ἰσαὰκ τὸν δοῦλόν σου, καὶ Ἰσραὴλ τὸν ἅγιόν σου,
\vs{36}οἷς ἐλάλησας πληθῦναι τὸ σπέρμα αὐτῶν, ὡς τὰ ἄστρα τοῦ οὐρανοῦ, καὶ ὡς τὴν ἄμμον τὴν παρὰ τὸ χεῖλος τῆς θαλάσσης.
\vs{37}Ὅτι, δέσποτα, ἐσμικρύνθημεν παρὰ πάντα τὰ ἔθνη, καί ἐσμὲν ταπεενοὶ ἐν πάσῃ τῇ γῇ σήμερον, διὰ τὰς ἁμαρτίας ἡμῶν.
\vs{38}Καὶ οὐκ ἔστιν ἐν τῷ καιρῷ τούτῳ ἄρχων καὶ προφήτης καὶ ἡγούμενος, οὐδὲ ὁλοκαύτωσις, οὐδὲ θυσία, οὐδὲ προσφορὰ, οὐδὲ θυμίαμα, οὐδὲ τόπος τοῦ καρπῶσαι ἐναντίον σου, καὶ εὑρεῖν ἕλεος.

\vs{39}Ἀλλʼ ἐν ψυχῇ συντετριμμένῃ, καὶ πνεύματι ταπεινώσεως προσδεχθείημεν, ὡς ἐν ὁλοκαυτώσει κριῶν καὶ ταύρων, καὶ ἐν μυριάσιν ἀρνῶν πιόνων,
\vs{40}οὕτως γενέσθω ἡ θυσία ἡμῶν ἐνώπιόν σου σήμερον, καὶ ἐκτελέσαι ὄπισθέν σου· ὅτι οὐκ ἔσται αἰσχύνη τοῖς πεποιθόσιν ἐπὶ σοί.

\vs{41}Καὶ νῦν ἐξακολουθοῦμεν ἐν ὅλῃ καρδίᾳ, καὶ φοβούμεθά σε, καὶ ζητοῦμεν τὸ πρόσωπόν σου. Μὴ καταισχύνῃς ἡμᾶς,
\vs{42}ἀλλὰ ποίησον μεθʼ ἡμῶν κατὰ τὴν ἐπιείκειάν σου, καὶ κατὰ τὸ πλῆθος τοῦ ἐλέους σου.

\vs{43}Καὶ ἐξελοῦ ἡμᾶς κατὰ τὰ θαυμάσιά σου, καὶ δὸς δόξαν τῷ ὀνόματί σου, Κύριε·
\vs{44}καὶ ἐντραπείησαν πάντες οἱ ἐνδεικνύμενοι τοῖς δούλοις σου κακὰ, καὶ καταισχυνθείησαν ἀπὸ πάσης τῆς δυναστείας, καὶ ἡ ἰσχὺς αὐτῶν συντριβείη,
\vs{45}καὶ γνώτωσαν ὅτι σὺ εἶ Κύριος, Θεὸς μόνος, καὶ ἔνδοξος ἐφʼ ὅλην τὴν οἰκουμένην.

\vs{46}Καὶ οὐ διέλιπον οἱ ἐμβάλλοντες αὐτοὺς ὑπηρέται τοῦ βασιλέως, καίοντες τὴν κάμινον νάφθαν καὶ πίσσαν καὶ στιππύον καὶ κληματίδα.
\vs{47}Καὶ διεχεῖτο ἡ φλὸξ ἐπάνω τῆς καμίνου ἐπὶ πήχεις τεσσαρακονταεννέα.
\vs{48}Καὶ διώδευσε, καὶ ἐνεπύρισεν οὓς εὗρε περὶ τὴν κάμινον τῶν Χαλδαίων.

\vs{49}Ὁ δὲ ἄγγελος Κυρίου συγκατέβη ἅμα τοῖς περὶ τὸν Αζαρίαν εἰς τὴν κάμινον, καὶ ἐξετίναξε τὴν φλόγα τοῦ πυρὸς ἐκ τῆς καμίνου,
\vs{50}καὶ ἐποίησε τὸ μέσον τῆς καμίνου, ὡς πνεῦμα δρόσου διασυρίζον· καὶ οὐχ ἥψατο αὐτῶν τὸ καθόλου τὸ πῦρ; καὶ οὐκ ἐλύπησεν, οὐδὲ παρηνώχλησεν αὐτοῖς.

\vs{51}Τότε οἱ τρεῖς ὡς ἐξ ἑνὸς στόματος ὕμνουν, καὶ ἐδόξαζον, καὶ ηὐλόγουν τὸν Θεὸν ἐν τῇ καμίνῳ, λέγοντες,

\vs{52}Εὐλογητὸς εἶ Κύριε ὁ Θεὸς τῶν πατέρων ἡμῶν, καὶ αἰνετὸς, καὶ ὑπερυψούμενος εἰς τοὺς αἰῶνας. Καὶ εὐλογημένον τὸ ὄνομα τῆς δόξης σου τὸ ἅγιον, καὶ ὑπεραινετὸν καὶ ὑπερυψούμενον εἰς πάντας τοὺς αἰῶνας.

\vs{53}Εὐλογημένος εἶ ἐν τῷ ναῷ τῆς ἁγίας δόξης σου, καὶ ὑπερυμνητὸς καὶ ὑπερένδοξος εἰς τοὺς αἰῶνας.
\vs{54}Εὐλογημένος εἶ ὁ ἐπιβλέπων ἀβύσσους, καθήμενος ἐπὶ χερουβὶμ, καὶ αἰνετὸς καὶ ὑπερυψούμενος εἰς τοὺς αἰῶνας.
\vs{55}Εὐλογημένος εἶ ἐπὶ θρόνου τῆς βασιλείας σου, καὶ ὑπερυμνητὸς καὶ ὑπερυμνούμενος εἰς τοῦς αἰῶνας.
\vs{56}Εὐλογητὸς εἶ ἐν τῷ στερεώματι τοῦ οὐρανοῦ, καὶ ὑμνητὸς καὶ δεδοξασμένος εἰς τοὺς αἰῶνας.

\vs{57}Εὐλογεῖτε πάντα τὰ ἔργα Κυρίου τὸν Κύριον, ὑμνεῖτε καὶ ὑπερυψοῦτε αὐτὸν εἰς τοὺς αἰῶνας.
\vs{58}Εὐλογεῖτε οὐρανοὶ τὸν Κύριον, ὑμνεῖτε καὶ ὑπερυψοῦτε αὐτὸν εἰς τοὺς αἰῶνας.
\vs{59}Εὐλογεῖτε ἄγγελοι Κυρίου τὸν Κύριον, ὑμνεῖτε καὶ ὑπερυψοῦτε αὐτὸν εἰς τοὺς αἰῶνας.
\vs{60}Εὐλογεῖτε ὕδατα καὶ πάντα τὰ ὑπεράνω τοῦ οὐρανοῦ τὸν Κύριον, ὑμνεῖτε καὶ ὑπερυψοῦτε αὐτὸν εἰς τοὺς αἰῶνας.
\vs{61}Εὐλογείτω πᾶσα ἡ δύναμις Κυρίου τὸν Κύριον, ὑμνεῖτε καὶ ὑπερυψοῦτε αὐτὸν εἰς τοὺς αἰῶνας.

\vs{62}Εὐλογεῖτε ἥλιος καὶ σελήνη τὸν Κύριον, ὑμνεῖτε καὶ ὑπερυψοῦτε αὐτὸν εἰς τοὺς αἰῶνας.
\vs{63}Εὐλογεῖτε ἄστρα τοῦ οὐρανοῦ τὸν Κύριον, ὑμνεῖτε καὶ ὑπερυψοῦτε αὐτὸν εἰς τοὺς αἰῶνας.
\vs{64}Εὐλογείτω πᾶς ὄμβρος καὶ δρόσος τὸν Κύριον, ὑμνεῖτε καὶ ὑπερυψοῦτε αὐτὸν εἰς τοὺς αἰῶνας.
\vs{65}Εὐλογεῖτε πάντα τὰ πνεύματα τὸν Κύριον, ὑμνεῖτε καὶ ὑπερυψοῦτε αὐτὸν εἰς τοὺς αἰῶνας.
\vs{66}Εὐλογεῖτε πῦρ καὶ καῦμα τὸν Κύριον, ὑμνεῖτε καὶ ὑπερυψοῦτε αὐτὸν εἰς τοὺς αἰῶνας.

\vs{71}Εὐλογεῖτε νύκτες καὶ ἡμέραι τὸν Κύριον, ὑμνεῖτε, καὶ ὑπερυψοῦτε αὐτὸν εἰς τοὺς αἰῶνας.
\vs{72}Εὐλογεῖτε φῶς καὶ σκότος τὸν Κύριον, ὑμνεῖτε καὶ ὑπερυψοῦτε αὐτὸν εἰς τοὺς αἰῶνας.
\vs{72a}Εὐλογεῖτε ψύχος καὶ καῦμα τὸν Κύριον, ὑμνεῖτε καὶ ὑπερυψοῦτε αὐτὸν εἰς τοὺς αἰῶνας.
\vs{72b}Εὐλογεῖτε πάχναι καὶ χιόνες τὸν Κύριον, ὑμνεῖτε καὶ ὑπερυψοῦτε αὐτὸν εἰς τοὺς αἰῶνας.
\vs{73}Εὐλογεῖτε ἀστραπαὶ καὶ νεφέλαι τὸν Κύριον, ὑμνεῖτε καὶ ὑπερυψοῦτε αὐτὸν εἰς τοὺς αἰῶνας.

\vs{74}Εὐλογείτω ἡ γῆ τὸν Κύριον, ὑμνείτω καὶ ὑπερυψούτω αὐτὸν εἰς τοὺς αἰῶνας.
\vs{75}Εὐλογεῖτε ὄρη καὶ βουνοὶ τὸν Κύριον, ὑμνεῖτε καὶ ὑπερυψοῦτε αὐτὸν εἰς τοὺς αἰῶνας.
\vs{76}Εὐλογεῖτε πάντα τὰ φυόμενα ἐν τῇ γῇ τὸν Κύριον, ὑμνεῖτε καὶ ὑπερυψοῦτε αὐτὸν εἰς τοὺς αἰῶνας.

\vs{77}Εὐλογεῖτε θάλασσα καὶ ποταμοὶ τὸν Κύριον, ὑμνεῖτε καὶ ὑπερυψοῦτε αὐτὸν εἰς τοὺς αἰῶνας.
\vs{78}Εὐλογεῖτε αἱ πηγαὶ τὸν Κύριον, ὑμνεῖτε καὶ ὑμερυψοῦτε αὐτὸν εἰς τοὺς αἰῶνας.
\vs{79}Εὐλογεῖτε κήτη καὶ πάντα τὰ κινούμενα ἐν τοῖς ὕδασι τὸν Κύριον, ὑμνεῖτε καὶ ὑπερυψοῦτε αὐτὸν εἰς τοὺς αἰῶνας.
\vs{80}Εὐλογεῖτε πάντα τὰ πετεινὰ τοῦ οὐρανοῦ τὸν Κύριον, ὑμνεῖτε καὶ ὑπερυψοῦτε αὐτὸν εἰς τοὺς αἰῶνας.
\vs{81}Εὐλογεῖτε πάντα τὰ θηρία καὶ τὰ κτήνη τὸν Κύριον, ὑμνεῖτε καὶ ὑπερυψοῦτε αὐτὸν εἰς τοὺς αἰῶνας.

\vs{82}Εὐλογεῖτε υἱοὶ τῶν ἀνθρώπων τὸν Κύριον, ὑμνεῖτε καὶ ὑπερυψοῦτε αὐτὸν εἰς τοὺς αἰῶνας.
\vs{83}Εὐλογεῖτε Ἰσραὴλ τὸν Κύριον, ὑμνεῖτε καὶ ὑπερυψοῦτε αὐτὸν εἰς τοὺς αἰῶνας.

\vs{84}Εὐλογεῖτε ἱερεῖς τὸν Κύριον, ὑμνεῖτε καὶ ὑπερυψοῦτε αὐτὸν εἰς τοὺς αἰῶνας.
\vs{85}Εὐλογεῖτε δοῦλοι τὸν Κύριον, ὑμνεῖτε καὶ ὑπερυψοῦτε αὐτὸν εἰς τοὺς αἰῶνας.
\vs{86}Εὐλογεῖτε πνεύματα καὶ ψυχαὶ δικαίων τὸν Κύριον, ὑμνεῖτε καὶ ὑπερυψοῦτε αὐτὸν εἰς τοὺς αἰῶνας.
\vs{87}Εὐλογεῖτε ὅσιοι καὶ ταπεινοὶ τῇ καρδίᾳ τὸν Κύριον, ὑμνεῖτε καὶ ὑπερυψοῦτε αὐτὸν εἰς τοὺς αἰῶνας.

\vs{88}Εὐλογεῖτε Ἀνανία, Ἀζαρία, Μισαὴλ τὸν Κύριον, ὑμνεῖτε καὶ ὑπερυψοῦτε αὐτὸν εἰς τοὺς αἰῶνας· ὅτι ἐξείλετο ἡμᾶς ἐξ ᾅδου, καὶ ἐκ χειρὸς θανάτου ἔσωσεν ἡμᾶς· καὶ ἐῤῥύσατο ἡμᾶς ἐκ μέσου καμίνου καιομένης φλογὸς, καὶ ἐκ μέσου πυρὸς ἐῤῥύσατο ἡμᾶς.
\vs{89}Ἐξομολογεῖσθε τῷ Κυρίῳ ὅτι χρηστὸς, ὅτι εἰς τὸν αἰῶνα τὸ ἔλεος αὐτοῦ.

\vs{90}Εὐλογεῖτε πάντες οἱ σεβόμενοι τὸν Κύριον τὸν Θεὸν τῶν θεῶν, ὑμνεῖτε καὶ ἐξομολογεῖσθε, ὅτι εἰς τὸν αἰῶνα τὸ ἔλεος αὐτοῦ.

\vs{91}Καὶ Ναβουχοδονόσορ ἤκουσεν ὑμνούντων αὐτῶν, καὶ ἐθαύμασε, καὶ ἐξανέστη ἐν σπουδῇ, καὶ εἶπε τοῖς μεγιστᾶσιν αὐτοῦ, οὐχὶ ἄνδρας τρεῖς ἐβάλομεν εἰς τὸ μέσον τοῦ πυρὸς πεπεδημένους; καὶ εἶπον τῷ βασιλεῖ, ἀληθῶς βασιλεῦ.
\vs{92}Καὶ εἶπεν ὁ βασιλεὺς, ὁ δὲ ἐγὼ ὁρῶ ἄνδρας τέσσαρας λελυμένους, καὶ περιπατοῦντας ἐν μέσῳ τοῦ πυρὸς, καὶ διαφθορὰ οὐκ ἔστιν ἐν αὐτοῖς, καὶ ἡ ὅρασις τοῦ τετάρτου ὁμοία υἱῷ Θεοῦ.
\vs{93}Τότε προσῆλθε Ναβουχοδονόσορ πρὸς τὴν θύραν τῆς καμίνου τοῦ πυρὸς τῆς καιομένης, καὶ εἶπε, Σεδρὰχ, Μισὰχ, Ἀβδεναγὼ, οἱ δοῦλοι τοῦ Θεοῦ τοῦ ὑψίστου, ἐξέλθετε καὶ δεῦτε· καὶ ἐξῆλθον Σεδρὰχ, Μισὰχ, Ἀβδεναγὼ, ἐκ μέσου τοῦ πυρός.
\vs{94}Καὶ συνάγονται οἱ σατράπαι, καὶ οἱ στρατηγοὶ, καὶ οἱ τοπάρχαι, καὶ οἱ δυνάσται τοῦ βασιλέως, καὶ ἐθεώρουν τοὺς ἄνδρας, ὅτι οὐκ ἐκυρίευσε τὸ πῦρ τοῦ σώματος αὐτῶν, καὶ ἡ θρὶξ τῆς κεφαλῆς αὐτῶν οὐκ ἐφλογίσθη, καὶ τὰ σαράβαρα αὐτῶν οὐκ ἠλλοιώθη, καὶ ὀσμὴ πυρὸς οὐκ ἦν ἐν αὐτοῖς.

\vs{95}Καὶ ἀπεκρίθη Ναβουχοδονόσορ ὁ βασιλεὺς, καὶ εἶπεν, εὐλογητὸς ὁ Θεὸς τοῦ Σεδρὰχ, Μισὰχ, Ἀβδεσαγὼ, ὃς ἀπέστειλε τὸν ἄγγελον αὐτοῦ, καὶ ἐξείλατο τοὺς παῖδας αὐτοῦ, ὅτι ἐπεποίθεισαν ἐπʼ αὐτῷ· καὶ τὸ ῥῆμα τοῦ βασιλέως ἠλλοίωσαν, καὶ παρέδωκαν τὰ σὼματα αὐτῶν εἰς πῦρ, ὅπως μὴ λατρεύσωσι μηδὲ προσκυνήσωσι παντὶ θεῷ, ἀλλʼ ἢ τῷ Θεῷ αὐτῶν.
\vs{96}Καὶ ἐγὼ ἐκτίθεμαι τὸ δόγμα· πᾶς λαὸς, φυλὴ, γλῶσσα, ἣ ἐὰν εἴπῃ βλασφημίαν κατὰ τοῦ Θεοῦ Σεδρὰχ, Μισὰχ, Ἀβδεναγὼ, εἰς ἀπώλειαν ἔσονται, καὶ οἱ οἶκοι αὐτῶν εἰς διαρπαγὴν, καθότι οὐκ ἔστι Θεὸς ἕτερος ὅστις δυνήσεται ῥύσασθαι οὕτως.
\vs{97}Τότε ὁ βασιλεὺς κατεύθυνε τὸν Σεδρὰχ, Μισὰχ, Ἀβδεναγὼ, ἐν τῇ χώρᾳ Βαβυλῶνος, καὶ ηὔξησεν αὐτοὺς, καὶ ἠξίωσεν αὐτοὺς ἡγεῖσθαι πάντων τῶν Ἰουδαίων, τῶν ἐν τῇ βασιλείᾳ αὐτοῦ.

\ch{4}
Ναβουχοδονόσορ ὁ βασιλεὺς πᾶσι τοῖς λαοῖς, φυλαῖς, καὶ γλώσσαις, τοῖς οἰκοῦσιν ἐν πάσῃ τῇ γῇ, ἐρήνη ὑμῖν πληθυνθείη.
\vs{2}Τὰ σημεῖα καὶ τὰ τέρατα ἃ ἐποίησε μετʼ ἐμοῦ ὁ Θεὸς ὁ ὕψιστος,
\vs{3}ἤρεσεν ἐναντίον ἐμοῦ ἀναγγεῖλαι ὑμῖν, ὡς μεγάλα καὶ ἰσχυρὰ, ἡ βασιλεία αὐτοῦ βασιλεία αἰώνιος, καὶ ἡ ἐξουσία αὐτοῦ εἰς γενεὰν καὶ γενεάν.

\vs{4}Ἐγὼ Ναβουχοδονόσορ εὐθηνῶν ἤμην ἐν τῷ οἴκῳ μου, καὶ εὐθαλῶν.
\vs{5}Ἐνύπνιον ἴδον, καὶ ἐφοβέρισέ με, καὶ ἐταράχθην ἐπὶ τῆς κοίτης μου, καὶ αἱ ὁράσεις τῆς κεφαλῆς μου ἐτάραξάν με.
\vs{6}Καὶ διʼ ἐμοῦ ἐτέθη δόγμα τοῦ εἰσαγαγεῖν ἐνώπιόν μου πάντας τοὺς σοφοὺς Βαβυλῶνος, ὅπως τὴν σύγκρισιν τοῦ ἐνυπνίου γνωρίσωσί μοι.
\vs{7}Καὶ εἰσπορεύοντο οἱ ἐπαοιδοὶ, μάγοι, γαζαρηνοὶ, Χαλδαῖοι· καὶ τὸ ἐνύπνιον ἐγὼ εἶπα ἐνώπιον αὐτῶν· καὶ τὴν σύγκρισιν αὐτοῦ οὐκ ἐγνώρισάν μοι,
\vs{8}ἕως ἦλθε Δανιὴλ, οὗ τὸ ὄνομα Βαλτάσαρ κατὰ τὸ ὄνομα τοῦ θεοῦ μου, ὃς πνεῦμα Θεοῦ ἅγιον ἐν ἑαυτῷ ἔχει. ᾧ εἶπα,

\vs{9}Βαλτάσαρ ὁ ἄρχων τῶν ἐπαοιδῶν, ὃν ἐγὼ ἔγνων ὅτι πνεῦμα Θεοῦ ἅγιον ἐν σοὶ, καὶ πᾶν μυστήριον οὐκ ἀδυνατεῖ σε, ἄκουσον τὴν ὅρασιν τοῦ ἐνυπνίου μου, οὗ ἴδον, καὶ τὴν σύγκρισιν αὐτοῦ εἰπόν μοι.
\vs{10}Ἐπὶ τῆς κοίτης μου ἐθεώρουν, καὶ ἰδοὺ δένδρον ἐν μέσῳ τῆς γῆς, καὶ τὸ ὕψος αὐτοῦ πολύ.
\vs{11}Ἐμεγαλύνθη τὸ δένδρον καὶ ἴσχυσε, καὶ τὸ ὕψος αὐτοῦ ἔφθασεν ἕως τοῦ οὐρανοῦ, καὶ τὸ κῦτος αὐτοῦ εἰς τὸ πέρας ἁπάσης τῆς γῆς,
\vs{12}τὰ φύλλα αὐτοῦ ὡραῖα, καὶ ὁ καρπὸς αὐτοῦ πολὺς, καὶ τροφὴ πάντων ἐν αὐτῷ, καὶ ὑποκάτω αὐτοῦ κατεσκήνουν τὰ θηρία τὰ ἄγρια, καὶ ἐν τοῖς κλάδοις αὐτοῦ κατῴκουν τὰ ὄρνεα τοῦ οὐρανοῦ, καὶ ἐξ αὐτοῦ ἐτρέφετο πᾶσα σάρξ.

\vs{13}Ἐθεώρουν ἐν ὁράματι τῆς νυκτὸς ἐπὶ τῆς κοίτης μου, καὶ ἰδοὺ εἲρ, καὶ ἅγιος ἀπʼ οὐρανοῦ κατέβη,
\vs{14}καὶ ἐφώνησεν ἐν ἰσχύϊ, καὶ οὕτως εἶπεν, ἐκκόψατε τὸ δένδρον, καὶ ἐκτίλατε τοὺς κλάδους αὐτοῦ, καὶ ἐκτινάξατε τὰ φύλλα αὐτοῦ, καὶ διασκορπίσατε τὸν καρπὸν αὐτοῦ· σαλευθήτωσαν τὰ θηρία ὑποκάτωθεν αὐτοῦ, καὶ τὰ ὄρνεα ἀπὸ τῶν κλάδων αὐτοῦ.
\vs{15}Πλὴν τὴν φυὴν τῶν ῥιζῶν αὐτοῦ ἐν τῇ γῇ ἐάσατε, καὶ ἐν δεσμῷ σιδηρῷ καὶ χαλκῷ, καὶ ἐν τῇ χλόῃ τῇ ἔξω, καὶ ἐν τῇ δρόσῳ τοῦ οὐρανοῦ κοιτασθήσεται, καὶ μετὰ τῶν θηρίων ἡ μερὶς αὐτοῦ ἐν τῷ χόρτῳ τῆς γῆς.
\vs{16}Ἡ καρδία αὐτοῦ ἀπὸ τῶν ἀνθρώπων ἀλλοιωθήσεται, καὶ καρδία θηρίου δοθήσεται αὐτῷ, καὶ ἑπτὰ καιροὶ ἀλλαγήσονται ἐπʼ αὐτόν.
\vs{17}Διὰ συνκρίματος εἲρ ὁ λόγος, καὶ ῥῆμα ἁγίων τὸ ἐπερώτημα, ἵνα γνῶσιν οἱ ζῶντες, ὅτι Κύριός ἐστιν ὁ ὕψιστος τῆς βασιλείας τῶν ἀνθρώπων, καὶ ᾧ ἐὰν δόξῃ δώσει αὐτὴν, καὶ ἐξουδένωμα ἀνθρώπων ἀναστήσει ἐπʼ αὐτήν.
\vs{18}Τοῦτο τὸ ἐνύπνιον ὃ ἴδον ἐγὼ Ναβουχοδονόσορ ὁ βασιλεὺς, καὶ σύ Βαλτάσαρ τὸ σύγκριμα εἰπὸν, ὅτι πάντες οἱ σοφοὶ τῆς βασιλείας μου οὐ δύνανται τὸ σύγκριμα αὐτοῦ δηλῶσαί μοι· σὺ δὲ Δανιὴλ δύνασαι, ὅτι πνεῦμα Θεοῦ ἅγιον ἐν σοί.

\vs{19}Τότε Δανιὴλ, οὗ τὸ ὄνομα Βαλτάσαρ, ἀπηνεώθη ὡσεὶ ὥραν μίαν, καὶ οἱ διαλογισμοὶ αὐτοῦ συνετάρασσον αὐτόν· καὶ ἀπεκρίθη Βαλτάσαρ, καὶ εἶπε, κύριε, τὸ ἐνύπνιον ἔστω τοῖς μισοῦσί σε, καὶ ἡ σύγκρισις αὐτοῦ τοῖς ἐχθροῖς σου.
\vs{20}Τὸ δένδρον ὃ εἶδες τὸ μεγαλυνθὲν καὶ τὸ ἰσχυκὸς, οὗ τὸ ὓψος ἔφθανεν εἰς τὸν οὐρανὸν, καὶ τὸ κῦτος αὐτοῦ εἰς πᾶσαν τὴν γῆν,
\vs{21}καὶ τὰ φύλλα αὐτοῦ εὐθαλῆ, καὶ ὁ καρπὸς αὐτοῦ πολύς, καὶ τροφὴ πᾶσιν ἐν αὐτῷ, ὑποκάτω αὐτοῦ κατῴκουν τὰ θηρία τὰ ἄγρια, καὶ ἐν τοῖς κλάδοις αὐτοῦ κατεσκήνουν τὰ ὄρνεα τοῦ οὐρανοῦ,
\vs{22}σὺ εἶ, βασιλεῦ, ὅτι ἐμεγαλύνθης καὶ ἴσχυσας, καὶ ἡ μεγαλωσύνη σου ἐμεγαλύνθη, καὶ ἔφθασεν εἰς τὸν οὐρανὸν, καὶ ἡ κυρεία σου εἰς τὰ πέρατα τῆς γῆς.
\vs{23}Καὶ ὅτι εἶδεν ὁ βασιλεὺς εἲρ, καὶ ἅγιον καταβαίνοντα ἀπὸ τοῦ οὐρανοῦ, καὶ εἶπεν, ἐκτίλατε τὸ δένδρον, καὶ διαφθείρατε αὐτὸ, πλὴν τὴν φυὴν τῶν ῥιζῶν αὐτοῦ ἐν τῇ γῇ ἐάσατε, καὶ ἐν δεσμῷ σιδηρῷ καὶ ἐν χαλκῷ, καὶ ἐν τῇ χλόῃ τῇ ἔξω, καὶ ἐν τῇ δρόσῳ τοῦ οὐρανοῦ αὐλισθήσεται, καὶ μετὰ θηρίων ἀγρίων ἡ μερὶς αὐτοῦ, ἕως οὗ ἑπτὰ καιροὶ ἀλλοιωθῶσιν ἐπʼ αὐτόν·
\vs{24}Τοῦτο ἡ σύγκρισις αὐτοῦ βασιλεῦ, καὶ σύγκριμα ὑψίστου ἐστὶν, ὃ ἔφθασεν ἐπὶ τὸν κύριόν μου τὸν βασιλέα,
\vs{25}καὶ σὲ ἐκδιώξουσιν ἀπὸ τῶν ἀνθρώπων, καὶ μετὰ θηρίων ἀγρίων ἔσται ἡ κατοικία σου, καὶ χόρτον ὡς βοῦν ψωμιοῦσί σε, καὶ ἀπὸ τῆς δρόσου τοῦ οὐρανοῦ αὐλισθήσῃ, καὶ ἑπτὰ καιροὶ ἀλλαγήσονται ἐπὶ σὲ, ἕως οὗ γνῷς ὅτι κυριεύει ὁ ὕψιστος τῆς βασιλείας τῶν ἀνθρώπων, καὶ ᾧ ἂν δόξῃ δώσει αὐτήν.
\vs{26}Καὶ ὅτι εἶπαν, ἐάσατε τὴν φυὴν τῶν ῥιζῶν τοῦ δένδρου· ἡ βασιλεία σου σοὶ μένει, ἀφʼ ἧς ἂν γνῷς τὴν ἐξουσίαν τὴν οὐράνιον.
\vs{27}Διατοῦτο, βασιλεῦ, ἡ βουλή μου ἀρεσάτω σοι, καὶ τὰς ἁμαρτίας σου ἐν ἐλεημοσύναις λύτρωσαι, καὶ τὰς ἀδικίας, ἐν οἰκτιρμοῖς πενήτων· ἴσως ἔσται μακρόθυμος τοῖς παραπτώμασί σου ὁ Θεός.

\vs{28}Ταῦτα πάντα ἔφθασεν ἐπὶ Ναβουχοδονόσορ τὸν βασιλέα.
\vs{29}Μετὰ δωδεκάμηνον, ἐπὶ τῷ ναῷ τῆς βασιλείας αὐτοῦ ἐν Βαβυλῶνι περιπατῶν,
\vs{30}ἀπεκρίθη ὁ βασιλεὺς, καὶ εἶπεν, οὐχ αὕτη ἐστὶ Βαβυλὼν ἡ μεγάλη, ἣν ἐγὼ ᾠκοδόμησα εἰς οἶκον βασιλείας, ἐν τῷ κράτει τῆς ἰσχύος μου, εἰς τιμὴν τῆς δόξης μου;

\vs{31}Ἔτι τοῦ λόγου ἐν τῷ στόματι τοῦ βασιλέως ὄντος, φωνὴ ἀπʼ οὐρανοῦ ἐγένετο, σοὶ λέγουσι Ναβουχοδονόσορ βασιλεῦ, ἡ βασιλεία παρῆλθεν ἀπὸ σοῦ,
\vs{32}καὶ ἀπὸ τῶν ἀνθρώπων σε ἐκδιώκουσι, καὶ μετὰ θηρίων ἀγρίων ἡ κατοικία σου, καὶ χόρτον ὡς βοῦν ψωμιοῦσί σε, καὶ ἑπτὰ καιροὶ ἀλλαγήσονται ἐπὶ σὲ, ἕως γνῷς ὅτι κυριεύει ὁ ὕψιστος τῆς βασιλείας τῶν ἀνθρώπων, καὶ ᾧ ἂν δόξῃ δώσει αὐτήν.

\vs{33}Αὐτῇ τῇ ὥρᾳ ὁ λόγος συνετελέσθη ἐπὶ Ναβουχοδονόσορ, καὶ ἀπὸ τῶν ἀνθρώπων ἐξεδιώχθη, καὶ χόρτον ὡς βοῦς ἤσθιε, καὶ ἀπὸ τῆς δρόσου τοῦ οὐρανοῦ τὸ σῶμα αὐτοῦ ἐβάφη, ἕως αἱ τρίχες αὐτοῦ ὡς λεόντων ἐμεγαλύνθησαν, καὶ οἱ ὄνυχες αὐτοῦ ὡς ὁρνέων.

\vs{34}Καὶ μετὰ τὸ τέλος τῶν ἡμερῶν ἐγὼ Ναβουχοδονόσορ τοὺς ὀφθαλμούς μου εἰς τὸν οὐρανὸν ἀνέλαβον, καὶ αἱ φρένες μου ἐπʼ ἐμὲ ἐπεστράφησαν, καὶ τῷ ὑψίστῳ ηὐλόγησα, καὶ τῷ ζῶντι εἰς τὸν αἰῶνα ᾔνεσα, καὶ ἐδόξασα, ὅτι ἡ ἐξουσία αὐτοῦ ἐξουσία αἰώνιος, καὶ ἡ βασιλεία αὐτοῦ εἰς γενεὰν καὶ γενεὰν,
\vs{35}καὶ πάντες οἱ κατοικοῦντες τὴν γῆν ὡς οὐδὲν ἐλογίσθησαν· καὶ κατὰ τὸ θέλημα αὐτοῦ ποιεῖ ἐν τῇ δυνάμει τοῦ οὐρανοῦ, καὶ ἐν τῇ κατοικίᾳ τῆς γῆς· καὶ οὐκ ἔστιν ὃς ἀντιποιήσεται τῇ χειρὶ αὐτοῦ, καὶ ἐρεῖ αὐτῷ, τί ἐποίησας;
\vs{36}Αὐτῷ τῷ καιρῷ αἱ φρένες μου ἐπεστράφησαν ἐπʼ ἐμέ, καὶ εἰς τὴν τιμὴν τῆς βασιλείας μου ἦλθον· καὶ ἡ μορφή μου ἐπέστρεψεν ἐπʼ ἐμέ, καὶ οἱ τύραννοί μου, καὶ οἱ μεγιστᾶνές μου ἐζήτουν με, καὶ ἐπὶ τὴν βασιλείαν μου ἐκραταιώθην, καὶ μεγαλωσύνη περισσοτέρα προσετέθη μοι.

\vs{37}Νῦν οὖν ἐγὼ Ναβουχοδονόσορ αἰνῶ καὶ ὑπερυψῶ καὶ δοξάζω τὸν βασιλέα τοῦ οὐρανοῦ, ὅτι πάντα τὰ ἔργα αὐτοῦ ἀληθινὰ, καὶ αἱ τρίβοι αὐτοῦ κρίσεις, καὶ πάντας τοὺς πορευομένους ἐν ὑπερηφανίᾳ δύναται ταπεινῶσαι.

\ch{5}
Βαλτάσαρ ὁ βασιλεὺς ἐποίησε δεῖπνον μέγα τοῖς μεγιστᾶσιν αὐτοῦ χιλίοις, καὶ κατέναντι τῶν χιλίων ὁ οἶνος,
\vs{2}καὶ πίνων Βαλτάσαρ εἶπεν ἐν τῇ γεύσει τοῦ οἴνου, τοῦ ἐνεγκεῖν τὰ σκεύη τὰ χρυσᾶ καὶ τὰ ἀργυρᾶ, ἃ ἐξήνεγκε Ναβουχοδονόσορ ὁ πατὴρ αὐτοῦ ἐκ τοῦ ναοῦ τοῦ ἐν Ἰερουσαλὴμ, καὶ πιέτωσαν ἐν αὐτοῖς ὁ βασιλεὺς, καὶ οἱ μεγιστᾶνες αὐτοῦ, καὶ αἱ παλλακαὶ αὐτοῦ, καὶ αἱ παράκοιτοι αὐτοῦ.
\vs{3}Καὶ ἠνέχθησαν τὰ σκεύη τὰ χρυσᾶ καὶ τὰ ἀργυρᾶ, ἃ ἐξήνεγκεν ἐκ τοῦ ναοῦ τοῦ Θεοῦ τοῦ ἐν Ἰερουσαλὴμ, καὶ ἔπινον ἐν αὐτοῖς ὁ βασιλεὺς, καὶ οἱ μεγιστᾶνες αὐτοῦ, καὶ αἱ παλλακαὶ αὐτοῦ, καὶ αἱ παράκοιτοι αὐτοῦ.
\vs{4}Ἔπινον οἶνον· καὶ ᾔνεσαν τοὺς θεοὺς τοὺς χρυσοῦς, καὶ ἀργυροῦς, καὶ χαλκοῦς, καὶ σιδηροῦς, καὶ ξυλίνους, καὶ λιθίνους.

\vs{5}Ἐν αὐτῇ τῇ ὥρᾳ ἐξῆλθον δάκτυλοι χειρὸς ἀνθρώπου, καὶ ἔγραφον κατέναντι τῆς λαμπάδος ἐπὶ τὸ κονίαμα τοῦ τοίχου τοῦ οἴκου τοῦ βασιλέως, καὶ ὁ βασιλεὺς ἐθεώρει τοὺς ἀστραγάλους τῆς χειρὸς τῆς γραφούσης.
\vs{6}Τότε τοῦ βασιλέως ἡ μορφὴ ἠλλοιώθη, καὶ οἱ διαλογισμοὶ αὐτοῦ συνετάρασσον αὐτὸν, καὶ οἱ σύνδεσμοι τῆς ὀσφύος αὐτοῦ διελύοντο, καὶ τὰ γόνατα αὐτοῦ συνεκροτοῦντο.
\vs{7}Καὶ ἐβόησεν ὁ βασιλεὺς ἐν ἰσχύϊ, τοῦ εἰσαγαγεῖν μάγους, Χαλδαίους, γαζαρηνούς· καὶ εἶπε τοῖς σοφοῖς Βαβυλῶνος, ὃς ἂν ἀναγνῷ τὴν γραφὴν ταύτην, καὶ τὴν σύγκρισιν γνωρίσῃ μοι, πορφύραν ἐνδύσεται, καὶ ὁ μανιάκης ὁ χρυσοῦς ἐπὶ τὸν τράχηλον αὐτοῦ, καὶ τρίτος ἐν τῇ βασιλείᾳ μου ἄρξει.
\vs{8}Καὶ εἰσεπορεύοντο πάντες οἱ σοφοὶ τοῦ βασιλέως, καὶ οὐκ ἠδύναντο τὴν γραφὴν ἀναγνῶναι, οὐδὲ τὴν σύγκρισιν γνωρίσαι τῷ βασιλεῖ.
\vs{9}Καὶ ὁ βασιλεὺς Βαλτάσαρ ἐταράχθη, καὶ ἡ μορφὴ αὐτοῦ ἠλλοιώθη ἐν αὐτῷ, καὶ οἱ μεγιστᾶνες αὐτοῦ συνεταράσσοντο.

\vs{10}Καὶ εἰσῆλθεν ἡ βασίλισσα εἰς τὸν οἶκον τοῦ πότου, καὶ εἶπε, βασιλεῦ, εἰς τὸν αἰῶνα ζῆθι· μὴ ταρασσέτωσάν σε οἱ διαλογισμοί σου, καὶ ἡ μορφή σου μὴ ἀλλοιούσθω.
\vs{11}Ἔστιν ἀνὴρ ἐν τῇ βασιλείᾳ σου, ἐν ᾧ πνεῦμα Θωοῦ· καὶ ἐν ταῖς ἡμέραις τοῦ πατρός σου, γρηγόρησις καὶ σύνεσις εὑρέθη ἐν αὐτῷ, καὶ ὁ βασιλεὺς Ναβουχοδονόσορ ὁ πατήρ σου ἄρχοντα ἐπαοιδῶν, μάγων, Χαλδαίων, γαζαρηνῶν, κατέστησεν αὐτὸν,
\vs{12}ὅτι πνεῦμα περισσὸν ἐν αὐτῷ, καὶ φρόνησις καὶ σύνεσις ἐν αὐτῷ, συγκρίνων ἐνύπνια, καὶ ἀναγγέλλων κρατούμενα, καὶ λύων συνδέσμους, Δανιήλ· καὶ ὁ βασιλεὺς ἐπέθηκεν ὄνομα αὐτῷ, Βαλτάσαρ· νῦν οὖν κληθήτω, καὶ τὴν σύγκρισιν αὐτοῦ ἀναγγελεῖ σοι.

\vs{13}Τότε Δανιὴλ εἰσήχθη ἐνώπιον τοῦ βασιλέως· καὶ εἶπεν ὁ βασιλεὺς τῷ Δανιὴλ, σὺ εἶ Δανιὴλ, ὁ ἀπὸ τῶν υἱῶν τῆς αἰχμαλωσίας τῆς Ἰουδαίας, ἧς ἤγαγεν ὁ βασιλεὺς ὁ πατήρ μου;
\vs{14}Ἤκουσα περὶ σοῦ ὅτι πνεῦμα Θεοῦ ἐν σοί, καὶ γρηγόρησις, καὶ σύνεσις, καὶ σοφία περισσὴ εὑρέθη ἐν σοί.
\vs{15}Καὶ νῦν εἰσῆλθον ἐνώπιόν μου οἱ σοφοὶ, μάγοι, γαζαρηνοὶ, ἵνα τὴν γραφὴν ταύτην ἀναγνῶσι, καὶ τὴν σύγκρισιν γνωρίσωσί μοι, καὶ οὐκ ἠδυνήθησαν ἀναγγεῖλαί μοι.
\vs{16}Καὶ ἐγὼ ἤκουσα περὶ σοῦ, ὅτι δύνασαι κρίματα συγκρῖναι· νῦν οὖν ἐὰν δυνηθῇς τὴν γραφὴν ἀναγνῶναι, καὶ τὴν σύγκρισιν αὐτῆς γνωρίσαι μοι, πορφύραν ἐνδύσῃ, καὶ ὁ μανιάκης ὁ χρυσοῦς ἔσται ἐπὶ τῷ τραχήλῳ σου, καὶ τρίτος ἐν τῇ βασιλείᾳ μου ἄρξεις.

\vs{17}Καὶ εἶπε Δανιὴλ ἐνώπιον τοῦ βασιλέως, τὰ δόματά σου σοὶ ἔστω, καὶ τὴν δωρεὰν τῆς οἰκίας σου ἑτέρῳ δὸς, ἐγὼ δὲ τὴν γραφὴν ἀναγνώσομαι, καὶ τὴν σύγκρισιν αὐτῆς γνωρίσω σοι,
\vs{18}βασιλεῦ· ὁ Θεὸς ὁ ὕψιστος τὴν βασιλείαν, καὶ τὴν μεγαλωσύνην, καὶ τὴν τιμὴν, καὶ τὴν δόξαν ἔδωκε Ναβουχοδονόσορ τῷ πατρί σου·
\vs{19}Καὶ ἀπὸ τῆς μεγαλωσύνης ἧς ἔδωκεν αὐτῷ πάντες οἱ λαοὶ, φυλαὶ, γλῶσσαι ἦσαν τρέμοντες καὶ φοβούμενοι ἀπὸ προσώπου αὐτοῦ· οὓς ἠβούλετο αὐτὸς ἀνῄρει, καὶ οὓς ἠβούλετο αὐτὸς ἔτυπτε, καὶ οὓς ἠβούλετο αὐτὸς ὕψου, καὶ οὓς ἠβούλετο αὐτὸς ἐταπείνου.
\vs{20}Καὶ ὅτε ὑψώθη ἡ καρδία αὐτοῦ, καὶ τὸ πνεῦμα αὐτοῦ ἐκραταιώθη τοῦ ὑπερηφανεύσασθαι, κατηνέχθη ἀπὸ τοῦ θρόνου τῆς βασιλείας, καὶ ἡ τιμὴ ἀφῃρέθη ἀπʼ αὐτοῦ,
\vs{21}καὶ ἀπὸ τῶν ἀνθρώπων ἐξεδιώχθη, καὶ ἡ καρδία αὐτοῦ μετὰ τῶν θηρίων ἐδόθη, καὶ μετὰ τῶν ὀνάγρων ἡ κατοικία αὐτοῦ, καὶ χόρτον ὡς βοῦν ἐψώμιζον αὐτὸν, καὶ ἀπὸ τῆς δρόσου τοῦ οὐρανοῦ τὸ σῶμα αὐτοῦ ἐβάφη, ἕως οὗ ἔγνω ὅτι κυριεύει ὁ Θεὸς ὕψιστος τῆς βασιλείας τῶν ἀνθρώπων, καὶ ᾧ ἂν δόξῃ δώσει αὐτήν.

\vs{22}Καὶ σὺ οὖν ὁ υἱὸς αὐτοῦ Βαλτάσαρ οὐκ ἐταπείνωσας τὴν καρδίαν σου κατενώπιον τοῦ Θεοῦ· οὐ πάντα ταῦτα ἔγνως;
\vs{23}Καὶ ἐπὶ τὸν Κύριον Θεὸν τοῦ οὐρανοῦ ὑψώθης, καὶ τὰ σκεύη τοῦ οἴκου αὐτοῦ ἤνεγκαν ἐνώπιόν σου, καὶ σὺ καὶ οἱ μεγιστᾶνές σου, καὶ αἱ παλλακαί σου, καὶ αἱ παράκοιτοί σου οἶνον ἐπίνετε ἐν αὐτοῖς· καὶ τοὺς θεοὺς τοὺς χρυσοῦς, καὶ ἀργυροῦς, καὶ χαλκοῦς, καὶ σιδηροῦς, καὶ ξυλίνους, καὶ λιθίνους, οἳ οὐ βλέπουσι, καὶ οἳ οὐκ ἀκούουσι, καὶ οὐ γινώσκουσιν, ᾔνεσας, καὶ τὸν Θεὸν οὗ ἡ πνοή σου ἐν χειρὶ αὐτοῦ καὶ πᾶσαι αἱ ὁδοί σου, αὐτὸν οὐκ ἐδόξασας.
\vs{24}Διατοῦτο ἐκ προσώπου αὐτοῦ ἀπεστάλη ἀστράγαλος χειρὸς, καὶ τὴν γραφὴν ταύτην ἐνέταξε.

\vs{25}Καὶ αὕτη ἡ γραφὴ ἐντεταγμένη, Μανὴ, Θεκὲλ, Φάρες.
\vs{26}Τοῦτο τὸ σύγκριμα τοῦ ῥήματος· Μανὴ, ἐμέτρησεν ὁ Θεὸς τὴν βασιλείαν σου, καὶ ἐπλήρωσεν αὐτήν.
\vs{27}Θεκὲλ, ἐστάθη ἐν ζυγῷ, καὶ εὑρέθη ὑστεροῦσα.
\vs{28}Φάρες, διήρηται ἡ βασιλεία σου, καὶ ἐδόθη Μήδοις καὶ Πέρσαις.

\vs{29}Καὶ εἶπε Βαλτάσαρ, καὶ ἐνέδυσαν τον Δανιὴλ πορφύραν, καὶ τὸν μανιάκην τὸν χρυσοῦν περιέθηκαν περὶ τὸν τράχηλον αὐτοῦ, καὶ ἐκήρυξε περὶ αὐτοῦ, εἶναι αὐτὸν ἄρχοντα τρίτον ἐν τῇ βασιλείᾳ.
\vs{30}Ἐν αὐτῇ τῇ νυκτὶ ἀνῃρέθη Βαλτάσαρ ὁ βασιλεὺς ὁ Χαλδαῖος,

\ch{6}καὶ Δαρεῖος ὁ Μῆδος παρέλαβε τὴν βασιλείαν, ὢν ἐτῶν ἑξήκοντα δύο.

\vs{2}Καὶ ἤρεσεν ἐνώπιον Δαρείου, καὶ κατέστησεν ἐπὶ τῆς βασιλείας σατράπας ἑκατὸν εἴκοσι, τοῦ εἶναι αὐτοὺς ἐν ὅλῃ τῇ βασιλείᾳ αὐτοῦ,
\vs{3}καὶ ἐπάνω αὐτῶν τακτικοὺς τρεῖς, ὃς ἦν Δανιὴλ εἷς ἐξ αὐτῶν, τοῦ ἀποδιδόναι αὐτοῖς τοὺς σατράπας λόγον, ὅπως ὁ βασιλεὺς μὴ ἐνοχλῆται.
\vs{4}Καὶ ἦν Δανιὴλ ὑπὲρ αὐτοὺς, ὅτι πνεῦμα περισσὸν ἐν αὐτῷ, καὶ ὁ βασιλεὺς κατέστησεν αὐτὸν ἐφʼ ὅλης τῆς βασιλείας αὐτοῦ.

\vs{5}Καὶ οἱ τακτικοὶ καὶ οἱ σατράπαι ἐζήτουν πρόφασιν εὑρεῖν κατὰ Δανιήλ· καὶ πᾶσαν πρόφασιν καὶ παράπτωμα καὶ ἀμπλάκημα οὐχ εὗρον κατʼ αὐτοῦ, ὅτι πιστὸς ἦν.
\vs{6}Καὶ εἶπον οἱ τακτικοὶ, οὐχ εὑρήσομεν κατὰ Δανιὴλ πρόφασιν, εἰ μὴ ἐν νομίμοις Θεοῦ αὐτοῦ.

\vs{7}Τότε οἱ τακτικοὶ καὶ οἱ σατράπαι παρέστησαν τῷ βασιλεῖ, καὶ εἶπαν αὐτῷ, Δαρεῖε βασιλεῦ, εἰς τοὺς αἰῶνας ζῆθι.
\vs{8}Συνεβουλεύσαντο πάντες οἱ ἐπὶ τῆς βασιλείας σου στρατηγοὶ καὶ σατράπαι, ὕπατοι καὶ τοπάρχαι, τοῦ στῆσαι στάσει βασιλικῇ, καὶ ἐνισχῦσαι ὁρισμόν, ὅπως ὃς ἂν αἰτήσῃ αἴτημα παρὰ παντὸς θεοῦ καὶ ἀνθρώπου, ἕως ἡμερῶν τριάκοντα, ἀλλʼ ἢ παρὰ σοῦ βασιλεῦ, ἐμβληθήσεται εἰς τὸν λάκκον τῶν λεόντων.
\vs{9}Νῦν οὖν βασιλεῦ στῆσον τὸν ὁρισμὸν, καὶ ἔκθες γραφὴν, ὅπως μὴ ἀλλοιωθῇ τὸ δόγμα Περσῶν καὶ Μήδων.
\vs{10}Τότε ὁ βασιλεὺς Δαρεῖος ἐπέταξε γραφῆναι τὸ δόγμα.

\vs{11}Καὶ Δανιὴλ ἡνίκα ἔγνω ὅτι ἐνετάγη τὸ δόγμα, εἰσῆλθεν εἰς τὸν οἶκον αὐτοῦ· καὶ αἱ θυρίδες ἀνεῳγμέναι αὐτῷ ἐν τοῖς ὑπερῴοις αὐτοῦ κατέναντι Ἱερουσαλὴμ, καὶ καιροὺς τρεῖς τῆς ἡμέρας ἦν κάμπτων ἐπὶ τὰ γόνατα αὐτοῦ, καὶ προσευχόμενος καὶ ἐξομολογούμενος ἐναντίον τοῦ Θεοῦ αὐτοῦ, καθὼς ἦν ποιῶν ἔμπροσθεν.

\vs{12}Τότε οἱ ἄνδρες ἐκεῖνοι παρετήρησαν, καὶ εὗρον τὸν Δανιὴλ ἀξιοῦντα καὶ δεόμενον τοῦ Θεοῦ αὐτοῦ.
\vs{13}Καὶ προσελθόντες λέγουσι τῷ βασιλεῖ, βασιλεῦ, οὐχ ὁρισμὸν ἔταξας, ὅπως πᾶς ἄνθρωπος ὃς ἂν αἰτήσῃ παρὰ παντὸς θεοῦ καὶ ἀνθρώπου αἴτημα, ἕως ἡμερῶν τριάκοντα, ἀλλʼ ἢ παρὰ σοῦ βασιλεῦ, ἐμβληθήσεται εἰς τὸν λάκκον τῶν λεόντων; καὶ εἶπεν ὁ βασιλεὺς, ἀληθινὸς ὁ λόγος, καὶ τὸ δόγμα Μήδων καὶ Περσῶν οὐ παρελεύσεται.
\vs{14}Τότε ἀπεκρίθησαν καὶ λέγουσιν ἐνώπιον τοῦ βασιλέως, Δανιὴλ, ὁ ἀπὸ τῶν υἱῶν τῆς αἰχμαλωσίας τῆς Ἰουδαίας, οὐχ ὑπετάγη τῷ δόγματί σου· καὶ καιροὺς τρεῖς τῆς ἡμέρας αἰτεῖ παρὰ τοῦ Θεοῦ αὐτοῦ τὰ αἰτήματα αὐτοῦ.
\vs{15}Τότε ὁ βασιλεὺς, ὡς τὸ ῥῆμα ἤκουσε, πολὺ ἐλυπήθη ἐπʼ αὐτῷ, καὶ περὶ τοῦ Δανιὴλ ἠγωνίσατο τοῦ ἐξελέσθαι αὐτὸν, καὶ ἕως ἑσπέρας ἦν ἀγωνιζόμενος τοῦ ἐξελέσθαι αὐτόν.

\vs{16}Τότε οἱ ἄνδρες ἐκεῖνοι λέγουσι τῷ βασιλεῖ, γνῶθι βασιλεῦ, ὅτι τὸ δόγμα Μήδοις καὶ Πέρσαις, τοῦ πᾶν ὁρισμὸν καὶ στάσιν ἣν ἂν ὁ βασιλεὺς στήσῃ, οὐ δεῖ παραλλάξαι.
\vs{17}Τότε ὁ βασιλεὺς εἶπε· καὶ ἢγαγον τὸν Δανιὴλ, καὶ ἐνέβαλον αὐτὸν εἰς τὸν λάκκον τῶν λεόντων· καὶ εἶπεν ὁ βασιλεὺς τῷ Δανιὴλ, ὁ Θεός σου, ᾧ σὺ λατρεύεις ἐνδελεχῶς, αὐτὸς ἐξελεῖταί σε.
\vs{18}Καὶ ἤνεγκαν λίθον, καὶ ἐπέθηκαν ἐπὶ τὸ στόμα τοῦ λάκκου, καὶ ἐσφραγίσατο ὁ βασιλεὺς ἐν τῷ δακτυλίῳ αὐτοῦ, καὶ ἐν τῷ δακτυλίῳ τῶν μεγιστάνων αὐτοῦ, ὅπως μὴ ἀλλοιωθῇ πρᾶγμα ἐν τῷ Δανιήλ.
\vs{19}Καὶ ἀπῆλθεν ὁ βασιλεὺς εἰς τὸν οἶκον αὐτοῦ, καὶ ἐκοιμήθη ἄδειπνος, καὶ ἐδέσματα οὐκ εἰσήνεγκαν αὐτῷ· καὶ ὁ ὕπνος ἀπέστη ἀπʼ αὐτοῦ· καὶ ἔκλεισεν ὁ Θεὸς τὰ στόματα τῶν λεόντων, καὶ οὐ παρηνώχλησαν τῷ Δανιήλ.

\vs{20}Τότε ὁ βασιλεὺς ἀνέστη τὸ πρωῒ ἐν τῷ φωτὶ, καὶ ἐν σπουδῇ ἦλθεν ἐπὶ τὸν λάκκον τῶν λεόντων.
\vs{21}Καὶ ἐν τῷ ἐγγίζειν αὐτὸν τῷ λάκκῳ, ἐβοήσε φωνῇ ἰσχυρᾷ, Δανιὴλ, ὁ δοῦλος τοῦ Θεοῦ τοῦ ζῶντος, ὁ Θεός σου ᾧ σὺ λατρεύεις ἐνδελεχῶς, εἰ ἠδυνήθη ἐξελέσθαί σε ἐκ στόματος τῶν λεόντων;
\vs{22}Καὶ εἶπε Δανιὴλ τῷ βασιλεῖ, βασιλεῦ, εἰς τοὺς αἰῶνας ζῆθι.
\vs{23}Ὁ Θεός μου ἀπέστειλε τὸν ἄγγελον αὐτοῦ, καὶ ἐνέφραξε τὰ στόματα τῶν λεόντων, καὶ οὐκ ἐλυμήναντό με, ὅτι κατέναντι αὐτοῦ εὐθύτης εὑρέθη ἐμοὶ, καὶ ἐνώπιον δὲ σου βασιλεῦ παράπτωμα οὐκ ἐποίησα.
\vs{24}Τότε ὁ βασιλεὺς πολὺ ἠγαθύνθη ἐπʼ αὐτῷ, καὶ τὸν Δανιὴλ εἶπεν ἀνενέγκαι ἐκ τοῦ λάκκου· καὶ ἀνηνέχθη Δανιὴλ ἐκ τοῦ λάκκου· καὶ πᾶσα διαφθορὰ οὐχ εὑρέθη ἐν αὐτῷ, ὅτι ἐπίστευσεν ἐν τῷ Θεῷ αὐτοῦ.

\vs{25}Καὶ εἶπεν ὁ βασιλεὺς, καὶ ἠγάγοσαν τοὺς ἄνδρας τοὺς διαβαλόντας τὸν Δανιὴλ, καὶ εἰς τὸν λάκκον τῶν λεόντων ἐνεβλήθησαν αὐτοὶ, καὶ οἱ υἱοὶ αὐτῶν, καὶ αἱ γυναῖκες αὐτῶν· καὶ οὐκ ἔφθασαν εἰς τὸ ἔδαφος τοῦ λάκκου, ἕως οὗ ἐκυρίευσαν αὐτῶν οἱ λέοντες, καὶ πάντα τὰ ὀστᾶ αὐτῶν ἐλέπτυναν.

\vs{26}Τότε Δαρεῖος ὁ βασιλεὺς ἔγραψε πᾶσι τοῖς λαοῖς, φυλαῖς, γλώσσαις, τοῖς οἰκοῦσιν ἐν πάσῃ τῇ γῇ· εἰρήνη ὑμῖν πληθυνθείη.
\vs{27}Ἐκ προσώπου μου ἐτέθη δόγμα τοῦτο ἐν πάσῃ ἀρχῇ τῆς βασιλείας μου, εἶναι τρέμοντας καὶ φοβουμένους ἀπὸ προσώπου τοῦ Θεοῦ Δανιήλ, ὅτι αὐτός ἐστι Θεὸς ζῶν, καὶ μένων εἰς τοὺς αἰῶνας, καὶ ἡ βασιλεία αὐτοῦ οὐ διαφθαρήσεται, καὶ ἡ κυρεία αὐτοῦ ἕως τέλους·
\vs{28}Ἀντιλαμβάνεται καὶ ῥύεται, καὶ ποιεῖ σημεῖα καὶ τέρατα ἐν τῷ οὐρανῷ καὶ ἐπὶ τῆς γῆς, ὅστις ἐξείλατο τὸν Δανιὴλ ἐκ χειρὸς τῶν λεόντων.
\vs{29}Καὶ Δανιὴλ κατηύθυνεν ἐν τῇ βασιλείᾳ Δαρείου, καὶ ἐν τῇ βασιλείᾳ Κύρου τοῦ Πέρσου.

\ch{7}
Ἐν ἔτει τρώτῳ τῷ Βαλτάσαρ βασιλέως Χαλδαίων, Δανιὴλ ἐνύπνιον εἶδε, καὶ αἱ ὁράσεις τῆς κεφαλῆς αὐτοῦ ἐπὶ τῆς κοίτης αὐτοῦ· καὶ τὸ ἐνύπνιον αὐτοῦ ἔγραψεν.

\vs{2}Ἐγὼ Δανιὴλ ἐθεώρουν· καὶ ἰδοὺ οἱ τέσσαρες ἄνεμοι τοῦ οὐρανοῦ προσέβαλον εἰς τὴν θάλασσαν τὴν μεγάλην·
\vs{3}Καὶ τέσσαρα θηρία μεγάλα ἀνέβαινον ἐκ τῆς θαλάσσης, διαφέροντα ἀλλήλων.
\vs{4}Τὸ πρῶτον ὡσεὶ λέαινα, καὶ πτερὰ αὐτῆς ὡς ἀετοῦ· ἐθεώρουν ἕως οὗ ἐξετίλη τὰ πτερὰ αὐτῆς· καὶ ἐξῄρθη ἀπὸ τῆς γῆς, καὶ ἐπὶ ποδῶν ἀνθρώπου ἐστάθη, καὶ καρδία ἀνθρώπου ἐδόθη αὐτῇ.
\vs{5}Καὶ ἰδοὺ θηρίον δεύτερον ὅμοιον ἄρκτῳ, καὶ εἰς μέρος ἓν ἐστάθη, καὶ τρεῖς πλευραὶ ἐν τῷ στόματι αὐτῆς, ἀναμέσον τῶν ὀδόντων αὐτῆς· καὶ οὕτως ἔλεγον αὐτῇ, ἀνάστηθι, φάγε σάρκας πολλάς.
\vs{6}Ὀπίσω τούτου ἐθεώρουν, καὶ ἰδοὺ θηρίον ἕτερον ὡσεὶ παρδάλις· καὶ αὐτῇ πτερὰ τέσσαρα πετεινοῦ ὑπεράνω αὐτῆς, καὶ τέσσαρες κεφαλαὶ τῷ θηρίῳ, καὶ ἐξουσία ἐδόθη αὐτῇ.
\vs{7}Ὀπίσω τούτου ἐθεώρουν, καὶ ἰδοὺ θηρίον τέταρτον φοβερὸν καὶ ἔκθαμβον, καὶ ἰσχυρὸν περισσῶς, καὶ οἱ ὀδόντες αὐτοῦ σιδηροῖ, ἐσθίον, καὶ λεπτύνον, καὶ τὰ ἐπίλοιπα τοῖς ποσὶν αὐτοῦ συνεπάτει, καὶ αὐτὸ διάφορον περισσῶς παρὰ πάντα τὰ θηρία τὰ ἔμπροσθεν αὐτοῦ· καὶ κέρατα δέκα αὐτῷ.
\vs{8}Προσενόουν τοῖς κέρασιν αὐτοῦ, καὶ ἰδοὺ κέρας ἕτερον μικρὸν ἀνέβη ἐν μέσῳ αὐτῶν, καὶ τρία κέρατα τῶν ἔμπροσθεν αὐτοῦ ἐξεῤῥιζώθη ἀπὸ προσώπου αὐτοῦ· καὶ ἰδοὺ ὀφθαλμοὶ, ὡσεὶ ὀφθαλμοὶ ἀνθρώπου ἐν τῷ κέρατι τούτῳ, καὶ στόμα λαλοῦν μεγάλα.

\vs{9}Ἐθεώρουν ἕως ὅτου οἱ θρόνοι ἐτέθησαν, καὶ παλαιὸς ἡμερῶν ἐκάθητο, καὶ τὸ ἔνδυμα αὐτοῦ λευκὸν ὡσεὶ χιὼν, καὶ ἡ θρὶξ τῆς κεφαλῆς αὐτοῦ ὡσεὶ ἔριον καθαρόν, ὁ θρόνος αὐτοῦ φλὸξ πυρός, οἱ τροχοὶ αὐτοῦ πῦρ φλέγον.
\vs{10}Ποταμὸς πυρὸς εἷλκεν ἔμπροσθεν αὐτοῦ· χίλιαι χιλιάδες ἐλειτούργουν αὐτῷ, καὶ μύριαι μυριάδες παρειστήκεισαν αὐτῷ· κριτήριον ἐκάθισε, καὶ βίβλοι ἠνεῴχθησαν.
\vs{11}Ἐθεώρουν τότε ἀπὸ φωνῆς τῶν λόγων τῶν μεγάλων, ὧν τὸ κέρας ἐκεῖνο ἐλάλει, ἕως ἀνῃρέθη τὸ θηρίον, καὶ ἀπώλετο, καὶ τὸ σῶμα αὐτοῦ ἐδόθη εἰς καῦσιν πυρός.
\vs{12}Καὶ τῶν λοιπῶν θηρίων μετεστάθη ἡ ἀρχὴ, καὶ μακρότης ζωῆς ἐδόθη αὐτοῖς ἕως καιροῦ καὶ καιροῦ.

\vs{13}Ἐθεώρουν ἐν ὁράματι τῆς νυκτός, καὶ ἰδοὺ μετὰ τῶν νεφελῶν τοῦ οὐρανοῦ, ὡς υἱὸς ἀνθρώπου ἐρχόμενος, καὶ ἕως τοῦ παλαιοῦ τῶν ἡμερῶν ἔφθασε, καὶ προσηνέχθη αὐτῷ.
\vs{14}Καὶ αὐτῷ ἐδόθη ἡ ἀρχὴ καὶ ἡ τιμὴ καὶ ἡ βασιλεία, καὶ πάντες οἱ λαοὶ, φυλαὶ, καὶ γλῶσσαι αὐτῷ δουλεύσουσιν· ἡ ἐξουσία αὐτοῦ, ἐξουσία αἰώνιος, ἥτις οὐ παρελεύσεται, καὶ ἡ βασιλεία αὐτοῦ οὐ διαφθαρήσεται.

\vs{15}Ἔφριξε τὸ πνεῦμά μου ἐν τῇ ἕξει μου, ἐγὼ Δανιὴλ, καὶ αἱ ὁράσεις τῆς κεφαλῆς μου ἐτάρασσόν με.
\vs{16}Καὶ προσῆλθον ἑνὶ τῶν ἑστηκότων, καὶ τὴν ἀκρίβειαν ἐζήτουν παρʼ αὐτοῦ μαθεῖν περὶ πάντων τούτων· καὶ εἶπέ μοι τὴν ἀκρίβειαν, καὶ τὴν σύγκρισιν τῶν λόγων ἐγνώρισέ μοι.
\vs{17}Ταῦτα τὰ θηρία τὰ τέσσαρα, τέσσαρες βασιλεῖαι ἀναστήσονται ἐπὶ τῆς γῆς, αἳ ἀρθήσονται·
\vs{18}καὶ παραλήψονται τὴν βασιλείαν ἅγιοι ὑψίστου, καὶ καθέξουσιν αὐτὴν ἕως αἰῶνος τῶν αἰώνων.

\vs{19}Καὶ ἐζήτουν ἀκριβῶς περὶ τοῦ θηοίου τοῦ τετάρτου· ὃτι ἦν διαφέρον παρὰ πᾶν θηρίον, φοβερὸν περισσῶς, οἱ ὀδόντες αὐτοῦ σιδηροῖ, καὶ ὄνυχες αὐτοῦ χαλκοῖ, ἐσθίον, καὶ λεπτῦνον, καὶ τὰ ἐπίλοιπα τοῖς ποσὶν αὐτοῦ συνεπάτει·
\vs{20}Καὶ περὶ τῶν κεράτων αὐτοῦ τῶν δέκα τῶν ἐν τῇ κεφαλῇ αὐτοῦ, καὶ τοῦ ἐτέρου τοῦ ἀναβάντος, καὶ ἐκτινάξαντος τῶν πρώτων, ᾧ οἱ ὀφθαλμοὶ καὶ στόμα λαλοῦν μεγάλα, καὶ ἡ ὅρασις αὐτοῦ μείζων τῶν λοιπῶν.
\vs{21}Ἐθεώρουν, καὶ τὸ κέρας ἐκεῖνο ἐποίει πόλεμον μετὰ τῶν ἁγίων, καὶ ἴσχυσε πρὸς αὐτοὺς,
\vs{22}ἕως οὗ ἦλθεν ὁ παλαιὸς ἡμερῶν, καὶ τὸ κρίμα ἔδωκεν ἁγίοις ὑψίστου· καὶ ὁ καιρὸς ἔφθασε, καὶ τὴν βασιλείαν κατέσχον οἱ ἅγιοι.
\vs{23}Καὶ εἶπε, τὸ θηρίον τὸ τέταρτον, βασιλεία τετάρτη ἔσται ἐν τῇ γῇ, ἥτις ὑπερέξει πάσας τὰς βασιλείας, καὶ καταφάγεται πᾶσαν τὴν γῆν, καὶ συμπατήσει αὐτὴν καὶ κατακόψει.
\vs{24}Καὶ τὰ δέκα κέρατα αὐτοῦ, δέκα βασιλεῖς ἀναστήσονται, καὶ ὀπίσω αὐτῶν ἀναστήσεται ἕτερος, ὃς ὑπεροίσει κακοῖς πάντας τοὺς ἔμπροσθεν, καὶ τρεῖς βασιλεῖς ταπεινώσει,
\vs{25}καὶ λόγους πρὸς τὸν ὕψιστον λαλήσει, καὶ τοὺς ἁγίους ὑψίστου παλαιώσει, καὶ ὑπονοήσει τοῦ ἀλλοιῶσαι καιροὺς καὶ νόμον, καὶ δοθήσεται ἐν χειρὶ αὐτοῦ ἕως καιροῦ καὶ καιρῶν καί γε ἥμισυ καιροῦ.
\vs{26}Καὶ τὸ κριτήριον ἐκάθισε, καὶ τὴν ἀρχὴν μεταστήσουσι τοῦ ἀφανίσαι, καὶ τοῦ ἀπολέσαι ἕως τέλους.
\vs{27}Καὶ ἡ βασιλεία καὶ ἡ ἐξουσία καὶ ἡ μεγαλωσύνη τῶν βασιλέων τῶν ὑποκάτω παντὸς τοῦ οὐρανοῦ, ἐδόθη ἁγίοις ὑψίστου· καὶ ἡ βασιλεία αὐτοῦ, βασιλεία αἰώνιος, καὶ πᾶσαι αἱ ἀρχαὶ αὐτῷ δουλεύσουσι καὶ ὑπακούσονται.

Ἕως ὧδε τὸ πέρας τοῦ λόγου·
\vs{28}ἐγὼ Δανιήλ, οἱ διαλογισμοί μου ἐπὶ πολὺ συνετάρασσόν με, καὶ ἡ μορφή μου ἠλλοιώθη, καὶ τὸ ῥῆμα ἐν τῇ καρδίᾳ μου διετήρησα.

\ch{8}
Ἐν ἔτει τρίτῳ τῆς βασιλείας Βαλτάσαρ τοῦ βασιλέως ὅρασις ὤφθη πρὸς μέ· ἐγὼ Δανιὴλ μετὰ τὴν ὀφθεῖσάν μοι τὴν ἀρχήν,
\vs{2}καὶ ἤμην ἐν Σούσοις τῇ βάρει, ἥ ἐστιν ἐν χώρᾳ Αἰλάμ· καὶ ἤμην ἐπὶ τοῦ Οὐβάλ.
\vs{3}Καὶ ᾖρα τοὺς ὀφθαλμούς μου, καὶ ἴδον· καὶ ἰδοὺ κριὸς εἷς ἑστηκὼς πρὸ τοῦ Οὐβάλ· καὶ αὐτῷ κέρατα ὑψηλά· καὶ τὸ ἓν ὑψηλότερον τοῦ ἑτέρου, καὶ τὸ ὑψηλὸν ἀνέβαινεν ἐπʼ ἐσχάτων.
\vs{4}Καὶ ἴδον τὸν κριὸν κερατίζοντα κατὰ θάλασσαν, καὶ Βοῤῥᾶν, καὶ Νότον· καὶ πάντα τὰ θηρία οὐ στήσεται ἐνώπιον αὐτοῦ· καὶ οὐκ ἦν ὁ ἐξαιρούμενος ἐκ χειρὸς αὐτοῦ, καὶ ἐποίησε κατὰ τὸ θέλημα αὐτοῦ καὶ ἐμεγαλύνθη.

\vs{5}Καὶ ἐγὼ ἤμην συνιῶν, καὶ ἰδοὺ τράγος αἰγῶν ἤρχετο ἀπὸ Λιβὸς ἐπὶ πρόσωπον πάσης τῆς γῆς, καὶ οὐκ ἦν ἁπτόμενος τῆς γῆς· καὶ τῷ τράγῳ κέρας μέσον τῶν ὀφθαλμῶν αὐτοῦ.
\vs{6}Καὶ ἦλθεν ἕως τοῦ κριοῦ τοῦ τὰ κέρατα ἔχοντος, οὗ ἴδον ἑστὼς ἐνώπιον τοῦ Οὐβάλ, καὶ ἔδραμε πρὸς αὐτὸν ἐν ὁρμῇ τῆς ἰσχύος αὐτοῦ.
\vs{7}Καὶ ἴδον αὐτὸν φθάνοντα ἕως τοῦ κριοῦ, καὶ ἐξηγριάνθη πρὸς αὐτόν, καὶ ἔπαισε τὸν κριὸν, καὶ συνέτριψεν ἀμφότερα τὰ κέρατα αὐτοῦ· καὶ οὐκ ἦν ἰσχὺς τῷ κριῷ, τοῦ στῆναι ἐνώπιον αὐτοῦ· καὶ ἔῤῥιψεν αὐτὸν ἐπὶ τὴν γῆν, καὶ συνεπάτησεν αὐτόν, καὶ οὐκ ἦν ὁ ἐξαιρούμενος τὸν κριὸν ἐκ χειρὸς αὐτοῦ.

\vs{8}Καὶ ὁ τράγος τῶν αἰγῶν ἐμεγαλύνθη ἕως σφόδρα· καὶ ἐν τῷ ἰσχῦσαι αὐτὸν, συνετρίβη τὸ κέρας αὐτοῦ τὸ μέγα· καὶ ἀνέβη ἕτερα τέσσαρα ὑποκάτω αὐτοῦ εἰς τοὺς τέσσαρας ἀνέμους τοῦ οὐρανοῦ.
\vs{9}Καὶ ἐκ τοῦ ἑνὸς αὐτῶν ἐξῆλθε κέρας ἓν ἰσχυρόν, καὶ ἐμεγαλύνθη περισσῶς πρὸς τὸν Νότον, καὶ πρὸς τὴν δύναμιν,
\vs{10}καὶ ἐμεγαλύνθη ἕως τῆς δυνάμεως τοῦ οὐρανοῦ· καὶ ἔπεσεν ἐπὶ τὴν γῆν ἀπὸ τῆς δυνάμεως τοῦ οὐρανοῦ καὶ ἀπὸ τῶν ἄστρων, καὶ συνεπάτησαν αὐτά.
\vs{11}Καὶ ἕως οὗ ὁ ἀρχιστράτηγος ῥύσηται τὴν αἰχμαλωσίαν, καὶ διʼ αὐτὸν θυσία ἐταράχθη, καὶ κατευωδώθη αὐτῷ· καὶ τὸ ἅγιον ἐρημωθήσεται.
\vs{12}Καὶ ἐδόθη ἐπὶ τὴν θυσίαν ἁμαρτία, καὶ ἐῤῥρίφη χαμαὶ ἡ δικαιοσύνη· καὶ ἐποίησε, καὶ εὐωδώθη.
\vs{13}Καὶ ἤκουσα ἑνὸς ἁγίου λαλοῦντος· καὶ εἶπεν εἷς ἅγιος τῷ φελμουνὶ τῷ λαλοῦντι, ἕως πότε ἡ ὅρασις στήσεται, ἡ θυσία ἡ ἀρθεῖσα, καὶ ἡ ἁμαρτία ἐρημώσεως ἡ δοθεῖσα, καὶ τὸ ἅγιον καὶ ἡ δύναμις συμπατηθήσεται;
\vs{14}Καὶ εἶπεν αὐτῷ, ἕως ἑσπέρας καὶ πρωῒ ἡμέραι δισχίλιαι καὶ τετρακόσιαι, καὶ καθαρισθήσεται τὸ ἅγιον.

\vs{15}Καὶ ἐγένετο ἐν τῷ ἰδεῖν με, ἐγὼ Δανιὴλ, τὴν ὅρασιν, καὶ ἐζήτουν σύνεσιν, καὶ ἰδοὺ ἔστη ἐνώπιον ἐμοῦ ὡς ὅρασις ἀνδρός.
\vs{16}Καὶ ἤκουσα φωνὴν ἀνδρὸς ἀναμέσον τοῦ Οὐβὰλ, καὶ ἐκάλεσε, καὶ εἶπε, Γαβριὴλ συνέτισον ἐκεῖνον τὴν ὅρασιν.
\vs{17}Καὶ ἦλθε, καὶ ἔστη ἐχόμενος τῆς στάσεώς μου· καὶ ἐν τῷ ἐλθεῖν αὐτὸν ἐθαμβήθην, καὶ πίπτω ἐπὶ πρόσωπόν μου· καὶ εἶπε πρὸς μὲ, σύνες υἱὲ ἀνθρώπου· ἔτι γὰρ εἰς καιροῦ πέρας ἡ ὅρασις·
\vs{18}Καὶ ἐν τῷ λαλεῖν αὐτὸν μετʼ ἐμοῦ, πίπτω ἐπὶ πρόσωπόν μου ἐπὶ τὴν γῆν, καὶ ἥψατό μου, καὶ ἔστησέ με ἐπὶ πόδας,
\vs{19}καὶ εἶπεν, ἰδοὺ ἐγὼ γνωρίζω σοι τὰ ἐσόμενα ἐπʼ ἐσχάτων τῆς ὀργῆς· ἔτι γὰρ εἰς καιροῦ πέρας ἡ ὅρασις.

\vs{20}Ὁ κριὸς ὃν εἶδες, ὁ ἔχων τὰ κέρατα, βασιλεὺς Μήδων καὶ Περσῶν.
\vs{21}Ὁ τράγος τῶν αἰγῶν, βασιλεὺς Ἑλλήνων· καὶ τὸ κέρας τὸ μέγα ὃ ἦν ἀναμέσον τῶν ὀφθαλμῶν αὐτοῦ, αὐτός ἐστιν ὁ βασιλεὺς ὁ πρῶτος.
\vs{22}Καὶ τοῦ συντριβέντος οὗ ἔστησαν τέσσαρα κέρατα ὑποκάτω, τέσσαρες βασιλεῖς ἐκ τοῦ ἔθνους αὐτοῦ ἀναστήσονται, καὶ οὐκ ἐν τῇ ἰσχύϊ αὐτῶν.
\vs{23}Καὶ ἐπʼ ἐσχάτων τῆς βασιλείας αὐτῶν, πληρουμένων τῶν ἁμαρτιῶν αὐτῶν, ἀναστήσεται βασιλεὺς ἀναιδὴς προσώπῳ, καὶ συνιῶν προβλήματα·
\vs{24}καὶ κραταιὰ ἡ ἰσχὺς αὐτοῦ, καὶ θαυμαστὰ διαφθερεῖ, καὶ κατευθυνεῖ, καὶ ποιήσει, καὶ διαφθερεῖ ἰσχυροὺς, καὶ λαὸν ἅγιον.
\vs{25}Καὶ ζυγὸς τοῦ κλοιοῦ αὐτοῦ κατευθυνεῖ, δόλος ἐν τῇ χειρὶ αὐτοῦ, καὶ ἐν καρδίᾳ αὐτοῦ μεγαλυνθήσεται, καὶ δόλῳ διαφθερεῖ πολλοὺς, καὶ ἐπὶ ἀπωλείας πολλῶν στήσεται· καὶ ὡς ὠὰ χειρὶ συντρίψει.
\vs{26}Καὶ ἡ ὅρασις τῆς ἑσπέρας καὶ τῆς πρωΐας τῆς ῥηθείσης ἀληθῶς ἐστι· καὶ σὺ σφράγισον τὴν ὅρασιν, ὅτι εἰς ἡμέρας πολλάς.

\vs{27}Καὶ ἐγὼ Δανιὴλ ἐκοιμήθην, καὶ ἐμαλακίσθην, καὶ ἀνέστην, καὶ ἐποίουν τὰ ἔργα τοῦ βασιλέως, καὶ ἐθαύμαζον τὴν ὅρασιν, καὶ οὐκ ἦν ὁ συνιῶν.

\ch{9}
Ἐν τῷ πρώτῳ ἔτει Δαρείου τοῦ υἱοῦ Ἀσουήρου, ἀπὸ τοῦ σπέρματος τῶν Μήδων, ὃς ἐβασίλευσεν ἐπὶ βασιλείαν Χαλδαίων,
\vs{2}ἐγὼ Δανιὴλ συνῆκα ἐν ταῖς βίβλοις τὸν ἀριθμὸν τῶν ἐτῶν, ὃς ἐγενήθη λόγος Κυρίου πρὸς Ἱερεμίαν τὸν προφήτην, εἰς συμπλήρωσιν ἐρημώσεως Ἱερουσαλὴμ ἑβδομήκοντα ἔτη.

\vs{3}Καὶ ἔδωκα τὸ πρόσωπόν μου πρὸς Κύριον τὸν Θεὸν, τοῦ ἐκζητῆσαι προσευχὴν καὶ δεήσεις ἐν νηστείαις καὶ σάκκῳ.
\vs{4}Καὶ προσευξάμην πρὸς Κύριον τὸν Θεόν μου, καὶ ἐξωμολογησάμην, καὶ εἶπα, Κύριε ὁ Θεὸς ὁ μέγας καὶ θαυμαστὸς, ὁ φυλάσσων τὴν διαθήκην σου, καὶ τὸ ἔλεός τοῖς ἀγαπῶσί σε σου καὶ τοῖς φυλάσσουσι, τὰς ἐντολάς σου,
\vs{5}ἡμάρτομεν, ἠδικήσαμεν, ἠνομήσαμεν, καὶ ἀπέστημεν καὶ ἐξεκλίναμεν ἀπὸ τῶν ἐντολῶν σου, καὶ ἀπὸ τῶν κριμάτων σου,
\vs{6}καὶ οὐκ εἰσηκούσαμεν τῶν δούλων σου τῶν προφητῶν οἳ ἐλάλουν ἐν τῷ ὀνόματί σου πρὸς τοὺς βασιλεῖς ἡμῶν, καὶ ἄρχοντας ἡμῶν, καὶ πατέρας ἡμῶν, καὶ πρὸς πάντα τὸν λαὸν τῆς γῆς.
\vs{7}Σοὶ Κύριε ἡ δικαιοσύνη καὶ ἡμῖν ἡ αἰσχύνη τοῦ προσώπου, ὡς ἡ ἡμέρα αὕτη, ἀνδρὶ Ἰούδα, καὶ τοῖς ἐνοικοῦσιν ἐν Ἱερουσαλὴμ, καὶ παντὶ Ἰσραὴλ, τοῖς ἐγγὺς καὶ τοῖς μακρὰν ἐν πάσῃ τῇ γῇ, οὗ διέσπειρας αὐτοὺς ἐκεῖ, ἐν ἀθεσίᾳ αὐτῶν ᾗ ἠθέτησαν.
\vs{8}Ἐν σοὶ Κύριέ ἐστιν ἡμῶν ἡ δικαιοσύνη, καὶ ἡμῖν ἡ αἰσχύνη τοῦ προσώπου, καὶ τοῖς βασιλεῦσιν ἡμῶν, καὶ τοῖς ἄρχουσιν ἡμῶν, καὶ τοῖς πατράσιν ἡμῶν, οἵτινες ἡμάρτομεν.
\vs{9}Σοὶ Κυρί τῷ Θεῷ ἡμῶν οἱ οἰκτιρμοὶ καὶ οἱ ἱλασμοὶ, ὅτι ἀπέστημεν,
\vs{10}καὶ οὐκ εἰσηκούσαμεν τῆς φωνῆς Κυρίου τοῦ Θεοῦ ἡμῶν, πορεύεσθαι ἐν τοῖς νόμοις αὐτοῦ, οἷς ἔδωκε κατὰ πρόσωπον ἡμῶν ἐν χερσὶ τῶν δούλων αὐτοῦ τῶν προφητῶν.

\vs{11}Καὶ πᾶς Ἰσραὴλ παρέβησαν τὸν νόμον σου, καὶ ἐξέκλιναν τοῦ μὴ ἀκοῦσαι τῆς φωνῆς σου· καὶ ἐπῆλθεν ἐφʼ ἡμᾶς ἡ κατάρα, καὶ ὁ ὅρκος ὁ γεγραμμένος ἐν νόμῳ Μωυσέως δούλου τοῦ Θεοῦ, ὅτι ἡμάρτομεν αὐτῷ.
\vs{12}Καὶ ἔστησε τοὺς λόγους αὐτοῦ οὓς ἐλάλησεν ἐφʼ ἡμᾶς, καὶ ἐπὶ τοὺς κριτὰς ἡμῶν, οἳ ἔκρινον ἡμᾶς, ἐπαγαγεῖν ἐφʼ ἡμᾶς κακὰ μεγάλα, οἷα οὐ γέγονεν ὑποκάτω παντὸς τοῦ οὐρανοῦ, κατὰ τὰ γενόμενα ἐν Ἱερουσαλὴμ,
\vs{13}καθὼς γέγραπται ἐν τῷ νόμῳ Μωυσῇ· πάντα τὰ κακὰ τοῦτα ἦλθεν ἐφʼ ἡμᾶς· καὶ οὐκ ἐδεήθημεν τοῦ προσώπου Κυρίου τοῦ Θεοῦ ἡμῶν, ἀποστρέψαι ἀπὸ τῶν ἀδικιῶν ἡμῶν, καὶ τοῦ συνιέναι ἐν πάσῃ ἀληθείᾳ σου.
\vs{14}Καὶ ἐγρηγόρησε Κύριος, καὶ ἐπήγαγεν αὐτὰ ἐφʼ ἡμᾶς, ὅτι δίκαιος Κύριος ὁ Θεὸς ἡμῶν ἐπὶ πᾶσαν τὴν ποίησιν αὐτοῦ ἣν ἐποίησε, καὶ οὐκ εἰσηκούσαμεν τῆς φωνῆς αὐτοῦ.
\vs{15}Καὶ νῦν Κύριε ὁ Θεὸς ἡμῶν, ὃς ἐξήγαγες τὸν λαόν σου ἐκ γῆς Αἰγύπτου ἐν χειρὶ κραταιᾷ, καὶ ἐποίησας σεαυτῷ ὄνομα ὡς ἡ ἡμέρα αὕτη, ἡμάρτομεν, ἠνομήσαμεν.

\vs{16}Κύριε ἐν πᾶσιν ἐλεημοσύνη σου, ἀποστραφήτω δὴ ὁ θυμός σου, καὶ ἡ ὀργή σου ἀπὸ τῆς πόλεώς σου Ἱερουσαλὴμ ὄρους ἁγίου σου, ὅτι ἡμάρτομεν, καὶ ἐν ταῖς ἀδικίαις ἡμῶν, καὶ τῶν πατέρων ἡμῶν, Ἱερουσαλὴμ, καὶ ὁ λαός σου εἰς ὀνειδισμὸν ἐγένετο ἐν πᾶσι τοῖς περικύκλῳ ἡμῶν.
\vs{17}Καὶ νῦν εἰσάκουσον Κύριε ὁ Θεὸς ἡμῶν τῆς προσευχῆς τοῦ δούλου σου καὶ τῶν δεήσεων αὐτοῦ, καὶ ἐπίφανον τὸ πρόσωπόν σου ἐπὶ τὸ ἁγίασμά σου τὸ ἔρημον, ἔνεκέν σου Κύριε,
\vs{18}κλῖνον ὁ Θεός μου τὸ οὖς σου, καὶ ἄκουσον· ἄνοιξον τοὺς ὀφθαλμούς σου, καὶ ἴδε τὸν ἀφανισμὸν ἡμῶν, καὶ τῆς πόλεως σου ἐφʼ ἧς ἐπικέκληται τὸ ὄνομά σου ἐπʼ αὐτῆς· ὅτι οὐκ ἐπὶ ταῖς δικαιοσύναις ἡμῶν ῥιπτοῦμεν τὸν οἰκτιρμὸν ἡμῶν ἐνώπιόν σου, ἀλλʼ ἐπὶ τοὺς οἰκτιρμούς σου τοὺς πολλοὺς Κύριε.
\vs{19}Εἰσάκουσον Κύριε, ἱλάσθητι Κύριε, πρόσχες Κύριε· μὴ χρονίσῃς ἕνεκέν σου ὁ Θεός μου, ὅτι τὸ ὄνομά σου ἐπικέκληται ἐπὶ τὴν πόλιν σου, καὶ ἐπὶ τὸν λαόν σου.

\vs{20}Καὶ ἔτι ἐμοῦ λαλοῦντος, καὶ προσευχομένου, καὶ ἐξαγορεύοντος τὰς ἁμαρτίας μου, καὶ τὰς ἁμαρτίας τοῦ λαοῦ μου Ἰσραὴλ, καὶ ῥιπτοῦντος τὸν ἔλεόν μου ἐναντίον τοῦ Κυρίου τοῦ Θεοῦ μου περὶ τοῦ ὄρους τοῦ ἁγίου,
\vs{21}καὶ ἔτι ἐμοῦ λαλοῦντος ἐν τῇ προσευχῇ, καὶ ἰδοὺ ἀνὴρ Γαβριὴλ, ὃν ἴδον ἐν τῇ ὁράσει ἐν τῇ ἀρχῆ, πετόμενος, καὶ ἥψατό μου, ὡσεὶ ὥραν θυσίας ἑσπερινῆς.
\vs{22}Καὶ συνέτισέ με, καὶ ἐλάλησε μετʼ ἐμοῦ, καὶ εἶπε, Δανιὴλ, νῦν ἐξῆλθον συμβιβάσαι σε σύνεσιν
\vs{23}ἐν ἀρχῇ τῆς δεήσεώς σου, ἐξῆλθε λόγος, καὶ ἐγὼ ἦλθον τοῦ ἀναγγεῖλαί σοι, ὅτι ἀνὴρ ἐπιθυμιῶν εἶ σὺ, καὶ ἐννοήθητι ἐν τῷ ῥήματι, καὶ σύνες ἐν τῇ ὀπτασίᾳ.

\vs{24}Ἑβδομήκοντα ἑβδομάδες συνετμήθησαν ἐπὶ τὸν λαόν σου, καὶ ἐπὶ τὴν πόλιν τὴν ἁγίαν, τοῦ συντελεσθῆναι ἁμαρτίαν, καὶ τοῦ σφραγίσαι ἁμαρτίας, καὶ ἀπαλεῖψαι τὰς ἀδικίας, καὶ τοῦ ἐξιλάσασθαι ἀδικίας, καὶ τοῦ ἀγαγεῖν δικαιοσύνην αἰώνιον, καὶ τοῦ σφραγίσαι ὅρασιν καὶ προφήτην, καὶ τοῦ χρίσαι ἅγιον ἁγίων.

\vs{25}Καὶ γνώσῃ καὶ συνήσεις ἀπὸ ἐξόδου λόγου τοῦ ἀποκριθῆναι, καὶ τοῦ οἰκοδομῆσαι Ἱερουσαλὴμ, ἓως Χριστοῦ ἡγουμένου ἑβδομάδες ἑπτὰ, καὶ ἑβδομάδες ἑξηκονταδύο· καὶ ἐπιστρέψει, καὶ οἰκοδομηθήσεται πλατεία, καὶ τεῖχος, καὶ ἐκκενωθήσονται οἱ καιροί.

\vs{26}Καὶ μετὰ τὰς ἑβδομάδας τὰς ἑξηκονταδύο, ἐξολοθρενθήσεται χρίσμα, καὶ κρίμα οὐκ ἔστιν ἐν αὐτῷ· καὶ τὴν πόλιν, καὶ τὸ ἅγιον διαφθερεῖ σὺν τῷ ἡγουμένῳ τῷ ἐρχομένῳ, ἐκκοπήσονται ἐν κατακλυσμῷ, καὶ ἕως τέλους πολέμου συντετμημένου τάξει, ἀφανισμοίς.

\vs{27}Καὶ δυναμώσει διαθήκην πολλοῖς ἑβδομὰς μία· καὶ ἐν τῷ ἡμίσει τῆς ἑβδομάδος ἀρθήσεταί μου θυσία καὶ σπονδὴ, καὶ ἐπὶ τὸ ἱερὸν βδέλυγμα τῶν ἐρημώσεων, καὶ ἕως τῆς συντελείας καιροῦ συντέλεια δοθήσεται ἐπὶ τὴν ἐρήμωσιν.

\ch{10}
Ἐν ἔτει τρίτῳ Κύρου βασιλέως Περσῶν λόγος ἀπεκαλύφθη τῷ Δανιὴλ, οὗ τὸ ὄνομα ἐπεκλήθη Βαλτάσαρ· καὶ ἀληθινὸς ὁ λόγος, καὶ δύναμις μεγάλη καὶ σύνεσις ἐδόθη αὐτῷ ἐν τῇ ὀπτασίᾳ.
\vs{2}Ἐν ταῖς ἡμέραις ἐκείναις ἐγὼ Δανιὴλ ἤμην πενθῶν τρεῖς ἑβδομάδας ἡμερῶν,
\vs{3}ἄρτον ἐπιθυμιῶν οὐκ ἔφαγον, καὶ κρέας καὶ οἶνος οὐκ εἰσῆλθεν εἰς τὸ στόμα μου, καὶ ἄλειμμα οὐκ ἠλειψάμην, ἕως πληρώσεως τριῶν ἑβδομάδων ἡμερῶν.

\vs{4}Ἐν ἡμέρᾳ εἰκοστῇ τετάρτῃ τοῦ μηνὸς τοῦ πρώτου, καὶ ἐγὼ ἤμην ἐχόμενα τοῦ ποταμοῦ τοῦ μεγάλου, αὐτός ἐστι Τίγρις Ἐδδεκέλ.
\vs{5}Καὶ ᾖρα τοὺς ὀφθαλμούς μου, καὶ ἴδον, καὶ ἰδοὺ ἀνὴρ εἷς ἐνδεδυμένος βαδδὶν, καὶ ἡ ὀσφὺς αὐτοῦ περιεζωσμένη ἐν χρυσίῳ Ὠφὰζ,
\vs{6}καὶ τὸ σῶμα αὐτοῦ ὡσεὶ Θαρσὶς, καὶ τὸ πρόσωπον αὐτοῦ ὡς ἡ ὅρασις ἀστραπῆς, καὶ οἱ ὀφθαλμοὶ αὐτοῦ ὡσεὶ λαμπάδες πυρὸς, καὶ οἱ βραχίονες αὐτοῦ καὶ τὰ σκέλη ὡς ὅρασις χαλκοῦ στίλβοντος, καὶ ἡ φωνὴ τῶν λόγων αὐτοῦ ὡς φωνὴ ὄχλου.
\vs{7}Καὶ ἴδον ἐγὼ Δανιὴλ μόνος τὴν ὀπτασίαν, καὶ οἱ ἄνδρες οἱ μετʼ ἐμοῦ οὐκ ἴδον τὴν ὀπτασίαν, ἀλλʼ ἢ ἔκστασις μεγάλη ἐπέπεσεν ἐπʼ αὐτοὺς, καὶ ἔφυγον ἐν φόβῳ.
\vs{8}Καὶ ἐγὼ ὑπελείφθην μόνος, καὶ ἴδον τὴν ὀπτασίαν τὴν μεγάλην ταύτην, καὶ οὐχ ὑπελείφθη ἐν ἐμοὶ ἰσχὺς, καὶ ἡ δόξα μου μετεστράφη εἰς διαφθοράν· καὶ οὐκ ἐκράτησα ἰσχύος.
\vs{9}Καὶ ἤκουσα τὴν φωνὴν τῶν λόγων αὐτοῦ· καὶ ἐν τῷ ἀκοῦσαί με αὐτοῦ, ἤμην κατανενυγμένος, καὶ τὸ πρόσωπόν μου ἐπὶ τὴν γῆν.

\vs{10}Καὶ ἰδοὺ χεὶρ ἁπτομένη μου, καὶ ἤγειρέ με ἐπὶ τὰ γόνατά μου.
\vs{11}Καὶ εἶπε πρὸς μὲ, Δανιὴλ ἀνὴρ ἐπιθυμιῶν, σύνες ἐν τοῖς λόγοις οἷς ἐγὼ λαλῶ πρὸς σὲ, καὶ στῆθι ἐπὶ τῇ στάσει σου, ὅτι νῦν ἀπεστάλην πρὸς σέ· καὶ ἐν τῷ λαλῆσαι αὐτὸν πρὸς μὲ τὸν λόγον τοῦτον, ἀνέστην ἔντρομος.
\vs{12}Καὶ εἶπε πρὸς μὲ, μὴ φοβοῦ Δανιὴλ, ὅτι ἀπὸ τῆς πρώτης ἡμέρας ἧς ἔδωκας τὴν καρδίαν σου τοῦ συνεῖναι, καὶ κακωθῆναι ἐναντίον Κυρίου τοῦ Θεοῦ σου, ἠκούσθησαν οἱ λόγοι σου, καὶ ἐγὼ ἦλθον ἐν τοῖς λόγοις σου.
\vs{13}Καὶ ὁ ἄρχων βασιλείας Περσῶν εἱστήκει ἐξεναντίας μου εἴκοσι καὶ μίαν ἡμέραν· καὶ ἰδοὺ Μιχαὴλ εἷς τῶν ἀρχόντων ἦλθε βοηθῆσαί μοι, καὶ αὐτὸν κατέλιπον ἐκεῖ μετὰ τοῦ ἄρχοντος βασιλείας Περσῶν,
\vs{14}καὶ ἦλθον συνετίσαι σε ὅσα ἀπαντήσεται τῷ λαῷ σου ἐπʼ ἐσχάτων τῶν ἡμερῶν, ὅτι ἔτι ἡ ὅρασις εἰς ἡμέρας.
\vs{15}Καὶ ἐν τῷ λαλῆσαι αὐτὸν μετʼ ἐμοῦ κατὰ τοὺς λόγους τούτους, ἔδωκα τὸ πρόσωπόν μου ἐπὶ τὴν γῆν, καὶ κατενύγην.

\vs{16}Καὶ ἰδοὺ ὡς ὁμοίωσις υἱοῦ ἀνθρώπου ἥψατο τῶν χειλέων μου· καὶ ἤνοιξα τὸ στόμα μου, καὶ ἐλάλησα, καὶ εἶπα πρὸς τὸν ἑστῶτα ἐναντίον μου, κύριε, ἐν τῇ ὀπτασίᾳ σου ἐστράφη τὰ ἐντός μου ἐν ἐμοὶ, καὶ οὐκ ἔσχον ἰσχύν.
\vs{17}Καὶ πῶς δυνήσεται ὁ παῖς σου κύριε λαλῆσαι μετὰ τοῦ κυρίου μου τούτου; καὶ ἐγὼ, ἀπὸ τοῦ νῦν οὐ στήσεται ἐν ἐμοὶ ἰσχὺς, καὶ πνεῦμα οὐχ ὑπελείφθη ἐν ἐμοί.
\vs{18}Καὶ προσέθετο, καὶ ἥψατό μου ὡς ὅρασις ἀνθρώπου, καὶ ἐνίσχυσέ με,
\vs{19}καὶ εἶπέ μοι, μὴ φοβοῦ ἀνὴρ ἐπιθυμιῶν, εἰρήνη σοι· ἀνδρίζου καὶ ἴσχυε· καὶ ἐν τῷ λαλῆσαι αὐτὸν μετʼ ἐμοῦ, ἴσχυσα, καὶ εἶπα λαλείτω ὁ κύριός μου, ὅτι ἐνίσχυσάς με.

\vs{20}Καὶ εἶπεν, εἰ οἶδας, ἱνατί ἦλθον πρὸς σέ; καὶ νῦν ἐπιστρέψω τοῦ πολεμῆσαι μετὰ τοῦ ἄρχοντος Περσῶν· καὶ ἐγὼ εἰσεπορευόμην, καὶ ὁ ἄρχων τῶν Ἑλλήνων ἤρχετο.
\vs{21}Ἀλλʼ ἢ ἀναγγελῶ σοι τὸ ἐντεταγμένον ἐν γραφῇ ἀληθείας, καὶ οὐκ ἔστιν εἷς ἀντεχόμενος μετʼ ἐμοῦ περὶ τούτων, ἀλλʼ ἢ Μιχαὴλ ὁ ἄρχων ὑμῶν.

\ch{11}
Καὶ ἐγὼ ἐν ἔτει πρώτῳ Κύρου ἔστην εἰς κράτος καὶ ἰσχύν.

\vs{2}Καὶ νῦν ἀλήθειαν ἀναγγελῶ σοι· ἰδοὺ ἔτι τρεῖς βασιλεῖς ἀναστήσονται ἐν τῇ Περσίδι, καὶ ὁ τέταρτος πλουτήσει πλοῦτον μέγαν παρὰ πάντας· καὶ μετὰ τὸ κρατῆσαι αὐτὸν τοῦ πλούτου αὐτοῦ, ἐπαναστήσεται πάσαις βασιλείαις Ἑλλήνων.

\vs{3}Καὶ ἀναστήσεται βασιλεὺς δυνατὸς, καὶ κυριεύσει κυρείας πολλῆς, καὶ ποιήσει κατὰ τὸ θέλημα αὐτοῦ.

\vs{4}Καὶ ὡς ἂν στῇ ἡ βασιλεία αὐτοῦ, συντριβήσεται, καὶ διαιρεθήσεται εἰς τοὺς τέσσαρας ἀνέμους τοῦ οὐρανοῦ, καὶ οὐκ εἰς τὰ ἔσχατα αὐτοῦ, οὐδὲ κατὰ τὴν κυρείαν αὐτοῦ, ἣν ἐκυρίευσεν, ὅτι ἐκτιλήσεται ἡ βασιλεία αὐτοῦ, καὶ ἑτέροις ἐκτὸς τούτων.

\vs{5}Καὶ ἐνισχύσει ὁ βασιλεὺς τοῦ Νότου· καὶ εἷς τῶν ἀρχόντων αὐτῶν ἐνισχύσει ἐπʼ αὐτὸν, καὶ κυριεύσει κυρίαν πολλήν.
\vs{6}Καὶ μετὰ τὰ ἔτη αὐτοῦ συμμιγήσονται, καὶ θυγάτηρ βασιλέως τοῦ Νότου εἰσελεύσεται πρὸς βασιλέα τοῦ Βοῤῥᾶ, τοῦ ποιῆσαι συνθήκας μετʼ αὐτοῦ, καὶ οὐ κρατήσει ἰσχύος βραχίονος, καὶ οὐ στήσεται τὸ σπέρμα αὐτοῦ, καὶ παραδοθήσεται αὕτη, καὶ οἱ φέροντες αὐτὴν, καὶ ἡ νεάνις, καὶ ὁ κατισχύων αὐτὴν ἐν τοῖς καιροῖς.

\vs{7}Ἀναστήσεται ἐκ τοῦ ἄνθους τῆς ῥίζης αὐτῆς, τῆς ἑτοιμασίας αὐτοῦ, καὶ ἥξει πρὸς τὴν δύναμιν, καὶ εἰσελεύσεται εἰς τὰ ὑποστηρίγματα τοῦ βασιλέως τοῦ Βοῤῥᾶ, καὶ ποιήσει ἐν αὐτοῖς, καὶ κατισχύσει.
\vs{8}Καί γε τοὺς θεοὺς αὐτῶν μετὰ τῶν χωνευτῶν αὐτῶν, πᾶν σκεῦος ἐπιθυμητὸν αὐτῶν, ἀργυρίου καὶ χρυσίου, μετὰ αἰχμαλωσίας οἴσει εἰς Αἴγυπτον, καὶ αὐτὸς στήσεται ὑπὲρ βασιλέα τοῦ Βοῤῥᾶ.
\vs{9}Καὶ εἰσελεύσεται εἰς τὴν βασιλείαν τοῦ βασιλέως τοῦ Νότου, καὶ ἀναστρέψει εἰς τὴν γῆν αὐτοῦ.

\vs{10}Καὶ οἱ υἱοὶ αὐτοῦ συνάξουσιν ὄχλον ἀναμέσον πολλῶν· καὶ ἐλεύσεται ἐρχόμενος καὶ κατακλύζων, καὶ παρελεύσεται, καὶ καθίεται, καὶ συμπροσπλακήσεται ἕως τῆς ἰσχύος αὐτοῦ.
\vs{11}Καὶ ἀγριανθήσεται βασιλεὺς τοῦ Νότου, καὶ ἐξελεύσεται, καὶ πολεμήσει μετὰ τοῦ βασιλέως τοῦ Βοῤῥᾶ, καὶ στήσει ὄχλον πολύν, καὶ παραδοθήσεται ὁ ὄχλος ἐν χειρὶ αὐτοῦ,
\vs{12}καὶ λήψεται τὸν ὄχλον, καὶ ὑψωθήσεται ἡ καρδία αὐτοῦ, καὶ καταβαλεῖ μυριάδας, καὶ οὐ κατισχύσει.
\vs{13}Καὶ ἐπιστρέψει ὁ βασιλεὺς τοῦ Βοῤῥᾶ, καὶ ἄξει ὄχλον πολὺν ὑπὲρ τὸν πρότερον· καὶ εἰς τὸ τέλος τῶν καιρῶν ἐνιαυτῶν ἐπελεύσεται εἰσόδια ἐν δυνάμει μεγάλῃ, καὶ ἐν ὑπάρξει πολλῇ.

\vs{14}Καὶ ἐν τοῖς καιροῖς ἐκείνοις πολλοὶ ἐπαναστήσονται ἐπὶ βασιλέα τοῦ Νότου, καὶ οἱ υἱοὶ τῶν λοιμῶν τοῦ λαοῦ σου ἐπαρθήσονται, τοῦ στῆσαι ὅρασιν, καὶ ἀσθενήσουσι.
\vs{15}Καὶ εἰσελεύσεται βασιλεὺς τοῦ Βοῤῥᾶ, καὶ ἐκχεεῖ πρόσχωμα, καὶ συλλήψεται πόλεις ὀχυρὰς, καὶ οἱ βραχίονες τοῦ βασιλέως τοῦ Νότου στήσονται, καὶ ἀναστήσονται οἱ ἐκλεκτοὶ αὐτοῦ, καὶ οὐκ ἔσται ἰσχὺς τοῦ στῆναι.
\vs{16}Καὶ ποιήσει ὁ εἰσπορευόμενος πρὸς αὐτὸν κατὰ τὸ θέλημα αὐτοῦ, καὶ οὐκ ἔστιν ἑστὼς κατὰ πρόσωπον αὐτοῦ· καὶ στήσεται ἐν τῇ γῇ τοῦ Σαβεὶ, καὶ τελεσθήσεται ἐν τῇ χειρὶ αὐτοῦ.

\vs{17}Καὶ τάξει τὸ πρόσωπον αὐτοῦ εἰσελθεῖν ἐν ἰσχύϊ πάσης τῆς βασιλείας αὐτοῦ, καὶ εὐθεῖα πάντα μετʼ αὐτοῦ ποιήσει· καὶ θυγατέρα τῶν γυναικῶν δώσει αὐτῷ διαφθεῖραι αὐτὴν, καὶ οὐ μὴ παραμείνῃ, καὶ οὐκ αὐτῷ ἔσται.
\vs{18}Καὶ ἐπιστρέψει τὸ πρόσωπον αὐτοῦ εἰς τὰς νήσους, καὶ συλλήψεται πολλὰς, καὶ καταπαύσει ἄρχοντας ὀνειδισμοῦ αὐτῶν, πλὴν ὀνειδισμὸς αὐτοῦ ἐπιστρέψει αὐτῷ.
\vs{19}Καὶ ἐπιστρέψει τὸ πρόσωπον αὐτοῦ εἰς τὴν ἰσχὺν τῆς γῆς αὐτοῦ, καὶ ἀσθενήσει, καὶ πεσεῖται, καὶ οὐκ εὑρεθήσεται.

\vs{20}Καὶ ἀναστήσεται ἐκ τῆς ῥίζης αὐτοῦ φυτὸν τῆς βασιλείας ἐπὶ τὴν ἑτοιμασίαν αὐτοῦ παραβιβάζων, πράσσων δόξαν βασιλείας· καὶ ἐν ταῖς ἡμέραις ἐκείναις ἔτι συντριβήσεται, καὶ οὐκ ἐν προσώποις, οὐδὲ ἐν πολέμῳ.

\vs{21}Στήσεται ἐπὶ τὴν ἑτοιμασίαν αὐτοῦ, ἐξουδενώθη, καὶ οὐκ ἔδωκαν ἐπʼ αὐτὸν δόξαν βασιλείας· καὶ ἥξει ἐν εὐθηνίᾳ, καὶ κατισχύσει βασιλείας ἐν ὀλισθήμασι,
\vs{22}καὶ βραχίονες τοῦ κατακλύζοντος κατακλυσθήσονται ἀπὸ προσώπου αὐτοῦ, καὶ συντριβήσονται, καὶ ἡγούμενος διαθήκης·
\vs{23}Καὶ ἀπὸ τῶν συναναμίξεων πρὸς αὐτὸν ποιήσει δόλον, καὶ ἀναβήσεται, καὶ ὑπερισχύσει αὐτοὺς ἐν ὀλίγῳ ἔθνει.
\vs{24}Καὶ ἐν εὐθηνίᾳ, καὶ ἐν πίοσι χώραις ἥξει, καὶ ποιήσει ἃ οὐκ ἐποίησαν οἱ πατέρες αὐτοῦ καὶ τατέρες τῶν πατέρων αὐτοῦ· προνομὴν καὶ σκῦλα καὶ ὕπαρξιν αὐτοῖς διασκορπιεῖ, καὶ ἐπʼ Αἴγυπτον λογιεῖται λογισμοὺς καὶ ἕως καιροῦ.
\vs{25}Καὶ ἐξεγερθήσεται ἡ ἰσχὺς αὐτοῦ, καὶ ἡ καρδία αὐτοῦ ἐπὶ βασιλέα τοῦ Νότου ἐν δυνάμει μεγάλῃ· καὶ ὁ βασιλεὺς τοῦ Νότου συνάψει πόλεμον ἐν δυνάμει μεγάλῃ καὶ ἰσχυρᾷ σφόδρα, καὶ οὐ στήσονται, ὅτι λογιοῦνται ἐπʼ αὐτὸν λογισμοὺς,
\vs{26}καὶ φάγονται τὰ δέοντα αὐτοῦ, καὶ συντρίψουσιν αὐτὸν, καὶ δυνάμεις κατακλύσει, καὶ πεσοῦνται τραυματίαι πολλοί.

\vs{27}Καὶ ἀμφότεροι οἱ βασιλεῖς, αἱ καρδίαι αὐτῶν εἰς πονηρίαν, καὶ ἐπὶ τραπέζῃ μιᾷ ψευδῆ λαλήσουσι, καὶ οὐ κατευθυνεῖ, ὅτι ἔτι πέρας εἰς καιρόν.
\vs{28}Καὶ ἐπιστρέψει εἰς τὴν γῆν αὐτοῦ ἐν ὑπάρξει πολλῇ, καὶ ἡ καρδία αὐτοῦ ἐπὶ διαθήκην ἁγίαν, καὶ ποιήσει, καὶ ἐπιστρέψει εἰς τὴν γῆν αὐτοῦ.

\vs{29}Εἰς τὸν καιρὸν ἐπιστρέψει, καὶ ἥξει ἐν τῷ Νότῳ, καὶ οὐκ ἔσται ὡς ἡ πρώτη καὶ ἡ ἐσχάτη.
\vs{30}Καὶ εἰσελεύσονται ἐν αὐτῷ οἱ ἐκπορευόμενοι Κίτιοι, καὶ ταπεινωθήσεται, καὶ ἐπιστρέψει, καὶ θυμωθήσεται ἐπὶ διαθήκην ἁγίαν· καὶ ποιήσει, καὶ ἐπιστρέψει, καὶ συνήσει ἐπὶ τοὺς καταλιπόντας διαθήκην ἁγίαν.

\vs{31}Καὶ σπέρματα ἐξ αὐτοῦ ἀναστήσονται, καὶ βεβηλώσουσι τὸ ἁγίασμα τῆς δυναστείας, καὶ μεταστήσουσι τὸν ἐνδελεχισμὸν, καὶ δώσουσι βδέλυγμα ἠφανισμένον.
\vs{32}Καὶ οἱ ἀνομοῦντες διαθήκην ἐπάξουσιν ἐν ὀλισθήμασι· καὶ λαὸς γινώσκοντες Θεὸν αὐτοῦ κατισχύσουσι, καὶ ποιήσουσι,
\vs{33}καὶ οἱ συνετοὶ τοῦ λαοῦ συνήσουσιν εἰς πολλὰ, καὶ ἀσθενήσουσιν ἐν ῥομφαίᾳ, καὶ ἐν φλογὶ, καὶ ἐν αἰχμαλωσίᾳ, καὶ ἐν διαρπαγῇ ἡμερῶν.
\vs{34}Καὶ ἐν τῷ ἀσθενῆσαι αὐτοὺς, βοηθήσονται βοήθειαν μικρὰν, καὶ προστεθήσονται πρὸς αὐτοὺς πολλοὶ ἐν ὀλισθήμασι.

\vs{35}Καὶ ἀπὸ τῶν συνιέντων ἀσθενήσουσι, τοῦ πυρῶσαι αὐτοὺς, καὶ τοῦ ἐκλέξασθαι, καὶ τοῦ ἀποκαλυφθῆναι ἕως καιροῦ πέρας, ὅτι ἔτι εἰς καιρόν.

\vs{36}Καὶ ποιήσει κατὰ τὸ θέλημα αὐτοῦ· καὶ ὁ βασιλεὺς ὑψωθήσεται καὶ μεγαλυνθήσεται ἐπὶ πάντα θεόν, καὶ λαλήσει ὑπέρογκα, καὶ κατευθυνεῖ μέχρις οὗ συντελεσθῇ ἡ ὀργὴ, εἰς γὰρ συντέλειαν γίνεται.
\vs{37}Καὶ ἐπὶ πάντος θεοὺς τῶν πατέρων αὐτοῦ οὐ συνήσει, καὶ ἐπιθυμία γυναικῶν, καὶ ἐπὶ πᾶν θεὸν οὐ συνήσει, ὅτι ἐπὶ πάντας μεγαλυνθήσεται.
\vs{38}Καὶ θεὸν Μαωζεὶμ ἐπὶ τόπου αὐτοῦ δοξάσει, καὶ θεὸν ὃν οὐκ ἔγνωσαν οἱ πατέρες αὐτοῦ, δοξάσει ἐν χρυσῷ καὶ ἀργύρῳ καὶ λίθῳ τιμίῳ, καὶ ἐν ἐπιθυμήμασι.
\vs{39}Καὶ ποιήσει τοῖς ὀχυρώμασι τῶν καταφυγῶν μετὰ θεοῦ ἀλλοτρίου, καὶ πληθυνεῖ δόξαν, καὶ ὑποτάξει αὐτοῖς πολλοὺς, καὶ γῆν διελεῖ ἐν δώροις.

\vs{40}Καὶ ἐν καιροῦ πέρατι συγκερατισθήσεται μετὰ τοῦ βασιλέως τοῦ Νότου· καὶ συναχθήσεται ἐπʼ αὐτὸν βασιλεὺς τοῦ Βοῤῥᾶ ἐν ἅρμασι καὶ ἐν ἱππεῦσι καὶ ἐν ναυσὶ πολλαῖς, καὶ εἰσελεύσονται εἰς τὴν γῆν, καὶ συντρίψει, καὶ παρελεύσεται,
\vs{41}καὶ εἰσελεύσεται εἰς τὴν γῆν τοῦ Σαβαεὶμ, καὶ πολλοὶ ἀσθενήσουσι· καὶ οὗτοι διασωθήσονται ἐκ χειρὸς αὐτοῦ, Ἐδὼμ, καὶ Μωὰβ, καὶ ἀρχὴ υἱῶν Ἀμμών.
\vs{42}Καὶ ἐκτενεῖ τὴν χεῖρα ἐπὶ τὴν γῆν, καὶ γῆ Αἰγύπτου οὐκ ἔσται εἰς σωτηρίαν.
\vs{43}Καὶ κυριεύσει ἐν τοῖς ἀποκρύφοις τοῦ χρυσοῦ καὶ τοῦ ἀργύρου, καὶ ἐν πᾶσιν ἐπιθυμητοῖς Αἰγύπτου, καὶ Λιβύων, καὶ Αἰθιόπων, ἐν τοῖς ὀχυρώμασιν αὐτῶν.
\vs{44}Καὶ ἀκοαὶ καὶ σπουδαὶ ταράξουσιν αὐτὸν ἐξ ἀνατολῶν καὶ ἀπὸ Βοῤῥᾶ· καὶ ἥξει ἐν θυμῷ πολλῷ, τοῦ ἀφανίσαι πολλούς.
\vs{45}Καὶ πήξει τὴν σκηνὴν αὐτοῦ Ἐφαδανῶ, ἀναμέσον τῶν θαλασσῶν εἰς ὄρος Σαβαεὶν ἅγιον, ἥξει ἕως μέρους αὐτοῦ, καὶ οὐκ ἔστιν ὁ ῥυόμενος αὐτόν.

\ch{12}
Καὶ ἐν τῷ καιρῷ ἐκείνῳ ἀναστήσεται Μιχαὴλ ὁ ἄρχων ὁ μέγας, ὁ ἑστηκὼς ἐπὶ τοὺς υἱοὺς τοῦ λαοῦ σου· καὶ ἔσται καιρὸς θλίψεως, θλίψις οἵα οὐ γέγονεν ἀφʼ οὗ γεγένηται ἔθνος ἐν τῇ γῇ, ἕως τοῦ καιροῦ ἐκείνου· ἐν τῷ καιρῷ ἐκείνῳ σωθήσεται ὁ λαός σου πᾶς ὁ γεγραμμένος ἐν τῇ βίβλῳ.
\vs{2}Καὶ πολλοὶ τῶν καθευδόντων ἐν γῆς χώματι ἐξεγερθήσονται, οὗτοι εἰς ζωὴν αἰώνιον, καὶ οὗτοι εἰς ὀνειδισμὸν καὶ εἰς αἰσχύνην αἰώνιον.
\vs{3}Καὶ οἱ συνιέντες λάμψουσιν ὡς ἡ λαμπρότης τοῦ στερεώματος, καὶ ἀπὸ τῶν δικαίων τῶν πολλῶν ὡς οἱ ἀστέρες εἰς τοὺς αἰῶνας, καὶ ἔτι.

\vs{4}Καὶ σὺ Δανιὴλ ἔμφραξον τοὺς λόγους, καὶ σφράγισον τὸ βιβλίον ἕως καιροῦ συντελείας, ἕως διδαχθῶσι πολλοὶ καὶ πληθυνθῇ ἡ γνῶσις.

\vs{5}Καὶ ἴδον ἐγὼ Δανιὴλ, καὶ ἰδοὺ δύο ἕτεροι εἱστήκεισαν, εἷς ἐντεῦθεν τοῦ χειλους τοῦ ποταμοῦ, καὶ εἷς ἐντεῦθεν τοῦ χείλους τοῦ ποταμοῦ.
\vs{6}Καὶ εἶπε τῷ ἀνδρὶ τῷ ἐνδεδυμένῳ τὰ βαδδὶν, ὃς ἦν ἐπάνω τοῦ ὕδατος τοῦ ποταμοῦ, ἕως πότε τὸ πέρας ὧν εἴρηκας τῶν θαυμασίων;
\vs{7}Καὶ ἤκουσα τοῦ ἀνδρὸς τοῦ ἐνδεδυμένον τὰ βαδδὶν, ὃς ἦν ἐπάνω τοῦ ὕδατος τοῦ ποταμοῦ· καὶ ὕψωσε τὴν δεξιὰν αὐτοῦ καὶ τὴν ἀριστερὰν αὐτοῦ εἰς τὸν οὐρανὸν, καὶ ὤμοσεν ἐν τῷ ζῶντι εἰς τὸν αἰῶνα, ὅτι εἰς καιρὸν καιρῶν καὶ ἥμισυ καιροῦ, ἐν τῷ συντελεσθῆναι διασκορπισμὸν γνώσονται πάντα ταῦτα.

\vs{8}Καὶ ἐγὼ ἤκουσα, καὶ οὐ συνῆκα· καὶ εἶπα, Κύριε, τί τὰ ἔσχατα τούτων;
\vs{9}Καὶ εἶπε, δεῦρο Δανιὴλ, ὅτι ἐμπεφραγμένοι καὶ ἐσφραγισμένοι οἱ λόγοι ἕως καιροῦ πέρας.
\vs{10}Ἐκλεγῶσι, καὶ ἐκλευκανθῶσι, καὶ πυρωθῶσι, καὶ ἁγιασθῶσι πολλό· καὶ ἀνομήσωσιν ἄνομοι, καὶ οὐ συνήσουσι πάντες ἄνομοι, καὶ οἱ νοήμονες συνήσουσι.
\vs{11}Καὶ ἀπὸ καιροῦ παραλλάξεως τοῦ ἐνδελεχισμοῦ, καὶ δοθήσεται τὸ βδέλυγμα ἐρημώσεως, ἡμέραι χίλιαι διακόσιαι ἐννενήκοντα.
\vs{12}Μακάριος ὁ ὑπομένων καὶ φθάσας εἰς ἡμέρας χιλίας τριακοσίας τριακονταπέντε.
\vs{13}Καὶ σὺ δεῦρο, καὶ ἀναπαύου· ἔτι γὰρ ἡμέραι καὶ ὧραι εἰς ἀναπλήρωσιν συντελείας, καὶ ἀναστήσῃ εἰς τὸν κλῆρόν σου εἰς συντέλειαν ἡμερῶν.


\def\book{ΩΣΗΕ}
\biblebook{ΩΣΗΕ}


\lettrine[lines=2, loversize=0.2, nindent=0em, findent=.25em]{\textcolor{bookheadingcolor}{Λ}}{ΟΓΟΣ} Κυρίου, ὃς ἐγενήθη πρὸς Ὠσηὲ τὸν τοῦ Βεηρεὶ, ἐν ἡμέραις Ὀζίου, καὶ Ἰωάθαμ, καὶ Ἄχαζ, καὶ Ἐζεκίου βασιλέων Ἰούδα, καὶ ἐν ἡμέραις Ἱεροβοὰμ υἱοῦ Ἰωὰς βασιλέως Ἰσραήλ.

\vs{2}Ἀρχὴ λόγου Κυρίου ἐν Ὠσῆέ· καὶ εἶπε Κύριος πρὸς Ὠσηὲ, βάδιζε, λάβε σεαυτῷ γυναῖκα πορνείας, καὶ τέκνα πορνείας, διότι ἐκπορνεύουσα ἐκπορνεύσει ἡ γῆ ἀπὸ ὄπισθεν τοῦ Κυρίου.

\vs{3}Καὶ ἐπορεύθη, καὶ ἔλαβε τὴν Γόμερ, θυγατέρα Δεβηλαΐμ· καὶ συνέλαβε καὶ ἔτεκεν αὐτῷ υἱόν.
\vs{4}Καὶ εἶπε Κύριος πρὸς αὐτὸν, κάλεσον τὸ ὄνομα αὐτοῦ Ἰεζραὲλ, διότι ἔτι μικρὸν, καὶ ἐκδικήσω τὸ αἷμα τοῦ Ἰεζραὲλ ἐπὶ τὸν οἶκον Ἰούδα, καὶ καταπαύσω βασιλείαν οἴκου Ἰσραήλ.
\vs{5}Καὶ ἔσται, ἐν τῇ ἡμέρᾳ ἐκείνῃ, συντρίψω τὸ τόξον τοῦ Ἰσραὴλ ἐν κοιλάδι τοῦ Ἰεζραέλ.

\vs{6}Καὶ συνέλαβεν ἔτι, καὶ ἔτεκε θυγατέρα· καὶ εἶπεν αὐτῷ, κάλεσον τὸ ὄνομα αὐτῆς, οὐκ ἠλεημένη· διότι οὐ μὴ προσθήσω ἔτι ἐλεῆσαι τὸν οἶκον Ἰσραὴλ, ἀλλʼ ἢ ἀντιτασσόμενος ἀντιτάξομαι αὐτοῖς.
\vs{7}Τοὺς δὲ υἱοὺς Ἰούδα ἐλεήσω, καὶ σώσω αὐτοὺς ἐν Κυρίῳ Θεῷ αὐτῶν, καὶ οὐ σώσω αὐτοὺς ἐν τόξῳ, οὐδὲ ἐν ῥομφαίᾳ, οὐδὲ ἐν πολέμῳ, οὐδὲ ἐν ἵπποις, οὐδὲ ἐν ἱππεῦσι.
\vs{8}Καὶ ἀπεγαλάκτισε τὴν οὐκ ἠλεημένην· καὶ συνέλαβεν ἔτι, καὶ ἔτεκεν υἱόν.
\vs{9}Καὶ εἶπε, κάλεσον τὸ ὄνομα αὐτοῦ, οὐ λαός μου· διότι ὑμεῖς οὐ λαός μου, καὶ ἐγὼ οὐκ εἰμὶ ὑμῶν.

\ch{2}Καὶ ἦν ὁ ἀριθμὸς τῶν υἱῶν Ἰσραὴλ, ὡς ἡ ἄμμος τῆς θαλάσσης, ἣ οὐκ ἐκμετρηθήσεται, οὐδὲ ἐξαριθμηθήσεται· καὶ ἔσται, ἐν τῷ τόπῳ, οὗ ἐῤῥέθη αὐτοῖς, οὐ λαός μου ὑμεῖς, κληθήσονται καὶ αὐτοὶ υἱοὶ Θεοῦ ζῶντος.
\vs{2}Καὶ συναχθήσονται υἱοὶ Ἰούδα, καὶ οἱ υἱοὶ Ἰσραὴλ ἐπιτοαυτὸ, καὶ θήσονται ἑαυτοῖς ἀρχὴν μίαν, καὶ ἀναβήσονται ἐκ τῆς γῆς, ὅτι μεγάλη ἡ ἡμέρα τοῦ Ἰεζραέλ.

\vs{3}Εἴπατε τῷ ἀδελφῷ ὑμῶν, λαός μου, καὶ τῇ ἀδελφῇ ὑμῶν, ἠλεημένη.
\vs{4}Κρίθητε πρὸς τὴν μητέρα ὑμῶν, κρίθητε, ὅτι αὕτη οὐ γυνή μου, καὶ ἐγὼ οὐκ ἀνὴρ αὐτῆς· καὶ ἐξαρῶ τὴν πορνείαν αὐτῆς ἐκ προσώπου μου, καὶ τὴν μοιχείαν αὐτῆς ἐκ μέσου μαστῶν αὐτῆς,
\vs{5}ὅπως ἂν ἐκδύσω αὐτὴν γυμνὴν, καὶ ἀποκαταστήσω αὐτὴν καθὼς ἡμέρα γενέσεως αὐτῆς· καὶ θήσω αὐτὴν ἔρημον, καὶ τάξω αὐτὴν ὡς γῆν ἄνυδρον, καὶ ἀποκτενῶ αὐτὴν ἐν δίψει.
\vs{6}Καὶ τὰ τέκνα αὐτῆς οὐ μὴ ἐλεήσω, ὅτι τέκνα πορνείας ἐστίν.
\vs{7}Ὅτι ἐξεπόρνευσεν ἡ μήτηρ αὐτῶν, κατῄσχυνεν ἡ τεκοῦσα αὐτά· ὅτι εἶπε, πορεύσομαι ὀπίσω τῶν ἐραστῶν μου, τῶν διδόντων μοι τοὺς ἄρτους μου, καὶ τὸ ὕδωρ μου, καὶ τὰ ἱμάτιά μου, καὶ τὰ ὀθόνιά μου, τὸ ἔλαιόν μου, καὶ πάντα ὅσα μοι καθήκει.

\vs{8}Διὰ τοῦτο ἰδοὺ ἐγὼ φράσσω τὴν ὁδὸν αὐτῆς ἐν σκόλοψι, καὶ ἀνοικοδομήσω τὰς ὁδοὺς, καὶ τὴν τρίβον αὐτῆς οὐ μὴ εὕρῃ·
\vs{9}Καὶ καταδιώξεται τοὺς ἐραστὰς αὐτῆς, καὶ οὐ μὴ καταλάβῃ αὐτούς· καὶ ζητήσει αὐτοὺς, καὶ οὐ μὴ εὕρῃ αὐτούς· καὶ ἐρεῖ, πορεύσομαι, καὶ ἐπιστρέψω πρὸς τὸν ἄνδρα μου τὸν πρότερον, ὅτι καλῶς μοι ἦν τότε, ἢ νῦν.

\vs{10}Καὶ αὕτη οὐκ ἔγνω ὅτι ἐγὼ ἔδωκα αὐτῇ τὸν σῖτον, καὶ τὸν οἶνον, καὶ τὸ ἔλαιον, καὶ ἀργύριον ἐπλήθυνα αὐτῇ· αὕτη δὲ ἀργυρᾶ καὶ χρυσᾶ ἐποίησεν τῇ Βάαλ.
\vs{11}Διὰ τοῦτο ἐπιστρέψω, καὶ κομιοῦμαι τὸν σῖτόν μου καθʼ ὥραν αὐτοῦ, καὶ τὸν οἶνόν μου ἐν καιρῷ αὐτοῦ· καὶ ἀφελοῦμαι τὰ ἱμάτιά μου, καὶ τὰ ὀθόνιά μου, τοῦ μὴ καλύπτειν τὴν ἀσχημοσύνην αὐτῆς·
\vs{12}Καὶ νῦν ἀποκαλύψω τὴν ἀκαθαρσίαν αὐτῆς ἐνώπιον τῶν ἐραστῶν αὐτῆς, καὶ οὐθεὶς οὐ μὴ ἐξέληται αὐτὴν ἐκ χειρός μου.
\vs{13}Καὶ ἀποστρέψω πάσας τὰς εὐφροσύνας αὐτῆς, ἑορτὰς αὐτῆς, καὶ τὰς νουμηνίας αὐτῆς, καὶ τὰ σάββατα αὐτῆς, καὶ πάσας τὰς πανηγύρεις αὐτῆς.
\vs{14}Καὶ ἀφανιῶ ἄμπελον αὐτῆς, καὶ τὰς συκᾶς αὐτῆς, ὅσα εἶπε, μισθώματά μου ταῦτά ἐστιν ἃ ἔδωκάν μοι οἱ ἐρασταί μου· καὶ θήσομαι αὐτὰ εἰς μαρτύριον, καὶ καταφάγεται αὐτὰ τὰ θηρία τοῦ ἀγροῦ, καὶ τὰ πετεινὰ τοῦ οὐρανου, καὶ τὰ ἑρπετὰ τῆς γῆς.
\vs{15}Καὶ ἐκδικήσω ἐπʼ αὐτὴν τὰς ἡμέρας τῶν Βααλεὶμ, ἐν αἷς ἐπέθυεν αὐτοῖς· καὶ περιετίθετο τὰ ἐνώτια αὐτῆς, καὶ τὰ καθόρμια αὐτῆς, καὶ ἐπορεύετο ὀπίσω τῶν ἐραστῶν αὐτῆς, ἐμοῦ δὲ ἐπελάθετο, λέγει Κύριος.

\vs{16}Διὰ τοῦτο ἰδοὺ ἐγὼ πλανῶ αὐτὴν, καὶ τάξω αὐτὴν ὡς ἔρημον, καὶ λαλήσω ἐπὶ τὴν καρδίαν αὐτῆς,
\vs{17}καὶ δώσω αὐτῇ τὰ κτήματα αὐτῆς ἐκεῖθεν, καὶ τὴν κοιλάδα Ἀχὼρ διανοῖξαι σύνεσιν αὐτῆς· καὶ ταπεινωθήσεται ἐκεῖ κατὰ τὰς ἡμέρας νηπιότητος αὐτῆς, καὶ κατὰ τὰς ἡμέρας ἀναβάσεως αὐτῆς ἐκ γῆς Αἰγύπτου.

\vs{18}Καὶ ἔσται ἐν τῇ ἡμέρᾳ ἐκεῖνῃ, λέγει Κύριος, καλέσει με ὁ ἀνήρ μου, καὶ οὐ καλέσει με ἔτι Βααλείμ.
\vs{19}Καὶ ἐξαρῶ τὰ ὀνόματα τῶν Βααλεὶμ ἐκ στόματος αὐτῆς, καὶ οὐ μὴ μνησθῶσιν οὐκέτι τὰ ὀνόματα αὐτῶν.
\vs{20}Καὶ διαθήσομαι αὐτοῖς διαθήκην ἐν τῇ ἡμέρᾳ ἐκείνῃ μετὰ τῶν θηρίων τοῦ ἀγροῦ, καὶ μετὰ τῶν πετεινῶν τοῦ οὐρανοῦ, καὶ τῶν ἑρπετῶν τῆς γῆς· καὶ τόξον, καὶ ῥομφαίαν, καὶ πόλεμον συντρίψω ἀπὸ τῆς γῆς, καὶ κατοικιῶ σε ἐπʼ ἐλπίδι.
\vs{21}Καὶ μνηστεύσομαί σε ἐμαυτῷ εἰς τὸν αἰῶνα· καὶ μνηστεύσομαί σε ἐμαυτῷ ἐν δικαιοσύνῃ καὶ ἐν κρίματι, καὶ ἐν ἐλέει, καὶ ἐν οἰκτιρμοῖς,
\vs{22}καὶ μνηστεύσομαί σε ἐμαυτῷ ἐν πίστει, καὶ ἐπιγνώσῃ τὸν Κύριον.

\vs{23}Καὶ ἔσται ἐν ἐκείνῃ τῇ ἡμέρᾳ, λέγει Κύριος, ἐπακούσομαι τῷ οὐρανῷ, καὶ αὐτὸς ἐπακούσεται τῇ γῇ,
\vs{24}καὶ ἡ γῆ ἐπακούσεται τὸν σῖτον, καὶ τὸν οἶνον, καὶ τὸ ἔλαιον, καὶ αὐτὰ ἐπακούσεται τῷ Ἰεζραέλ.
\vs{25}Καὶ σπερῶ αὐτὴν ἐμαυτῷ ἐπὶ τῆς γῆς, καὶ ἀγαπήσω τὴν οὐκ ἠγαπημένην, καὶ ἐρῶ τῷ οὐ λαῷ μου, λαός μου εἶ σύ· καὶ αὐτος ἐρεὶ, Κύριος ὁ Θεός μου εἶ σύ.

\ch{3}
Καὶ εἶπε, Κύριος πρὸς μὲ, ἔτι πορεύθητι, καὶ ἀγάπησον γυναῖκα ἀγαπῶσαν πονηρὰ, καὶ μοιχαλὶν, καθὼς ἀγαπᾷ ὁ Θεὸς τοὺς υἱοὺς Ἰσραὴλ, καὶ αὐτοὶ ἐπιβλέπουσιν ἐπὶ θεοὺς αλλοτρίους, καὶ φιλοῦσι πέμματα μετὰ σταφίδος.
\vs{2}Καὶ ἐμισθωσάμην ἐμαυτῷ πεντεκαίδεκα ἀργυρίου, καὶ γομὸρ κριθῶν, καὶ νέβελ οἴνου.
\vs{3}Καὶ εἶπα πρὸς αὐτὴν, ἡμέρας πολλὰς καθήσῃ ἐπʼ ἐμοὶ, καὶ οὐ μὴ πορνεύσῃς, οὐδὲ μὴ γένῃ ἀνδρὶ, καὶ ἐγὼ ἐπὶ σοί.

\vs{4}Διότι ἡμέρας πολλὰς καθήσονται οἱ υἱοὶ Ἰσραὴλ, οὐκ ὄντος βασιλέως, οὐδὲ ὄντος ἄρχοντος, οὐδὲ οὔσης θυσίας, οὐδὲ ὄντος θυσιαστηρ ίου, οὐδὲ ἱερατείας, οὐδὲ δήλων.
\vs{5}Καὶ μετὰ ταῦτα ἐπιστρέψουσιν οἱ υἱοὶ Ἰσραὴλ, καὶ ἐπιζητήσουσι Κύριον τὸν Θεὸν αὐτῶν, καὶ Δαυὶδ τὸν βασιλέα αὐτῶν, καὶ ἐκστήσονται ἐπὶ τῷ Κυρίῳ, καὶ ἐπὶ τοῖς ἀγαθοῖς αὐτοῦ ἐπʼ ἐσχάτων τῶν ἡμερῶν.

\ch{4}
Ἀκούσατε λόγον Κυρίου υἱοὶ Ἰσραὴλ, ὅτι κρίσις τῷ Κυρίῳ πρὸς τοὺς κατοικοῦντας τὴν γῆν, διότι οὐκ ἔστιν ἀλήθεια, οὐδὲ ἔλεος, οὐδὲ ἐπίγνωσις Θεοῦ ἐπὶ τῆς γῆς.
\vs{2}Ἀρὰ, καὶ ψεῦδος, καὶ φόνος καὶ κλοπὴ, καὶ μοιχεία κέχυται ἐπὶ τῆς γῆς, καὶ αἵματα ἐφʼ αἵμασι μίσγουσι.
\vs{3}Διὰ τοῦτο πενθήσει ἡ γῆ, καὶ σμικρυνθήσεται σὺν πᾶσι τοῖς κατοικοῦσιν αὐτὴν, σὺν τοῖς θηρίοις τοῦ ἀγροῦ, καὶ σὺν τοῖς ἑρπετοῖς τῆς γῆς, καὶ σὺν τοῖς πετεινοῖς τοῦ οὐρανοῦ, καὶ οἱ ἰχθύες τῆς θαλάσσης ἐκλείψουσιν,
\vs{4}ὅπως μηδεὶς μήτε δικάζηται, μήτε ἐλέγχῃ μηδείς· ὁ δὲ λαός μου ὡς ἀντιλεγόμενος ἱερεύς.
\vs{5}Καὶ ἀσθενήσει ἡμέρας, καὶ ἀσθενήσει ὁ προφήτης μετὰ σοῦ· νυκτὶ ὡμοίωσα τὴν μητέρα σου.

\vs{6}Ὡμοιώθη ὁ λαός μου, ὡς οὐκ ἔχων γνῶσιν· ὅτι σὺ ἐπίγνωσιν ἀπώσω, κᾀγὼ ἀπώσομαί σε, τοῦμὴ ἱερατεύειν μοι· καὶ ἐπελάθου νόμον Θεοῦ σου, κᾀγὼ ἐπιλήσομαι τέκνων σου.
\vs{7}Κατὰ τὸ πλῆθος αὐτῶν, οὕτως ἥμαρτόν μοι· τὴν δόξαν αὐτῶν εἰς ἀτιμίαν θήσομαι.
\vs{8}Ἁμαρτίας λαοῦ μου φάγονται· καὶ ἐν ταῖς ἀδικίαις αὐτῶν λήψονται τὰς ψυχὰς αὐτῶν.
\vs{9}Καὶ ἔσται καθὼς ὁ λαὸς, οὕτως καὶ ὁ ἱερεύς· καὶ ἐκδικήσω ἐπʼ αὐτὸν τὰς ὁδοὺς αὐτοῦ, καὶ τὰ διαβούλια αὐτοῦ ἀνταποδώσω αὐτῷ.
\vs{10}Καὶ φάγονται, καὶ οὐ μὴ ἐμπλησθῶσιν· ἐπόρνευσαν, καὶ οὐ μὴ κατευθύνωσι· διότι τὸν Κύριον ἐγκατέλιπον τοῦ φυλάξαι.

\vs{11}Πορνείαν καὶ οἶνον καὶ μέθυσμα ἐδέξατο καρδία λαοῦ μου·
\vs{12}ἐν συμβόλοις ἐπηρώτων, καὶ ἐν ῥάβδοις αὐτοῦ ἀπήγγελλον αὐτῷ· πνεύματι πορνείας ἐπλανήθησαν, καὶ ἐξεπόρνευσαν ἀπὸ τοῦ Θεοῦ αὐτῶν.
\vs{13}Ἐπὶ τὰς κορυφὰς τῶν ὀρέων ἐθυσίαζον, καὶ ἐπὶ τοὺς βουνοὺς ἔθυον ὑποκάτω δρυὸς, καὶ λεύκης, καὶ δένδρου συσκιάζοντος, ὅτι καλὸν σκέπη. διὰ τοῦτο ἐκπορνεύσουσιν αἱ θυγατέρες ὑμῶν, καὶ αἱ νύμφαι ὑμῶν μοιχεύσουσι.
\vs{14}Καὶ οὐ μὴ ἐπισκέψωμαι ἐπὶ τὰς θυγατέρας ὑμῶν ὅταν πορνεύσωσι, καὶ ἐπὶ τὰς νύμφας ὑμῶν ὅταν μοιχεύωσιν· ὅτι αὐτοὶ μετὰ τῶν πορνῶν συνεφύροντο, καὶ μετὰ τῶν τετελεσμένων ἔθυον, καὶ ὁ λαὸς ὁ μὴ συνιὼν συνεπλέκετο μετὰ πόρνης.

\vs{15}Σὺ δὲ Ἰσραὴλ μὴ ἀγνόει, καὶ Ἰούδα μὴ εἰσπορεύεσθε εἰς Γάλγαλα, καὶ μὴ ἀναβαίνετε εἰς τὸν οἶκον Ὦν, καὶ μὴ ὀμνύετε ζῶντα Κύριον.
\vs{16}Διότι ὡς δάμαλις παροιστρῶσα παροίστρησεν Ἰσραήλ· νῦν νεμήσει αὐτοὺς Κύριος ὡς ἀμνὸν ἐν εὐρυχώρῳ.
\vs{17}Μέτοχος εἰδώλων Ἐφραὶμ ἔθηκεν ἑαυτῷ σκάνδαλα,
\vs{18}ἠρέτισε Χαναναίους· πορνεύοντες ἐξεπόρνευσαν, ἠγάπησαν ἀτιμίαν ἐκ φρυάγματος αὐτῆς.
\vs{19}Συστροφὴ πνεύματος σὺ εἶ ἐν ταῖς πτέρυξιν αὐτῆς, καὶ καταισχυνθήσονται ἐκ τῶν θυσιαστηρίων αὐτῶν.

\ch{5}
Ἀκούσατε ταῦτα οἱ ἱερεῖς, καὶ προσέχετε οἶκος Ἰσραήλ, καὶ ὁ οἶκος τοῦ βασιλέως ἐνωτίζεσθε, διότι πρὸς ὑμᾶς ἐστι τὸ κρίμα, ὅτι παγὶς ἐγενήθητε τῇ Σκοπιᾷ, καὶ ὡς δίκτυον ἐκτεταμένον ἐπὶ τὸ Ἰταβύριον,
\vs{2}ὃ οἱ ἀγρεύοντες τὴν θήραν κατέπηξαν· ἐγὼ δὲ παιδευτὴς ὑμῶν,
\vs{3}ἐγὼ ἔγνων τὸν Ἐφραὶμ, καὶ Ἰσραὴλ οὐκ ἀπέστη ἀπʼ ἐμοῦ· διότι νῦν ἐξεπόρνευσεν Ἐφραὶμ, ἐμιάνθη Ἰσραήλ.
\vs{4}Οὐκ ἔδωκαν τὰ διαβούλια αὐτῶν τοῦ ἐπιστρέψαι πρὸς τὸν Θεὸν αὐτῶν, ὅτι πνεῦμα πορνίας ἐν αὐτοῖς ἐστίν· τὸν δὲ Κύριον οὐκ ἐπέγνωσαν.

\vs{5}Καὶ ταπεινωθήσεται ἡ ὕβρις τοῦ Ἰσραὴλ εἰς πρόσωπον αὐτοῦ· καὶ Ἰσραὴλ καὶ Ἐφραὶμ ἀσθενήσουσιν ἐν ταῖς ἀδικίαις αὐτῶν· καὶ ἀσθενήσει καὶ Ἰούδας μετʼ αὐτῶν.
\vs{6}Μετὰ προβάτων καὶ μόσχων πορεύσονται τοῦ ἐκζητῆσαι τὸν Κύριον, καὶ οὐ μὴ εὕρωσιν αὐτόν· ὅτι ἐκκέκλικεν ἀπʼ αὐτῶν,
\vs{7}ὅτι τὸν Κύριον ἐγκατέλιπον, ὅτι τέκνα ἀλλότρια ἐγεννήθησαν αὐτοῖς· νῦν καταφάγεται αὐτοὺς ἡ ἐρυσίβη, καὶ τοὺς κλήρους αὐτῶν.

\vs{8}Σαλπίσατε σάλπιγγι ἐπὶ τοὺς βουνοὺς, ἠχήσατε ἐπὶ τῶν ὑψηλῶν, κηρύξατε ἐν τῷ οἴκῳ Ὦν, ἐξέστη Βενιαμὶν,
\vs{9}Ἐφραὶμ εἰς ἀφανισμὸν ἐγένετο ἐν ἡμέραις ἐλέγχου· ἐν ταῖς φυλαῖς τοῦ Ἰσραὴλ ἔδειξα πιστά.
\vs{10}Ἐγένοντο οἱ ἄρχοντες Ἰούδα ὡς μετατιθέντες ὅρια, ἐπʼ αὐτοὺς ἐκχεῶ ὡς ὕδωρ τὸ ὅρμημά μου.

\vs{11}Κατεδυνάστευσεν Ἐφραὶμ τὸν ἀντίδικον αὐτοῦ, κατεπάτησε τὸ κρίμα, ὅτι ἤρξατο πορεύεσθαι ὀπίσω τῶν ματαίων.
\vs{12}Καὶ ἐγὼ ὡς ταραχὴ τῷ Ἐφραὶμ, καὶ ὡς κέντρον τῷ οἴκῳ Ἰούδα.
\vs{13}Καὶ εἶδεν Ἐφραὶμ τὴν νόσον αὐτοῦ, καὶ Ἰούδας τὴν ὀδύνην αὐτοῦ· καὶ ἐπορεύθη Ἐφραὶμ πρὸς Ἀσσυρίους, καὶ ἀπέστειλε πρέσβεις πρὸς βασιλέα Ἰαρείμ· καὶ οὗτος οὐκ ἠδυνάσθη ἰάσασθαι ὑμᾶς, καὶ οὐ μὴ διαπαύσῃ ἐξ ὑμῶν ὀδύνη.
\vs{14}Διότι ἐγώ εἰμι ὡς πανθὴρ τῷ Ἐφραὶμ, καὶ ὡς λέων τῷ οἴκῳ Ἰούδα· καὶ ἐγὼ ἁρπῶμαι καὶ πορεύσομαι, καὶ λήψομαι, καὶ οὐκ ἔσται ὁ ἐξαιρούμενος.

\vs{15}Πορεύσομαι καὶ ἐπιστρέψω εἰς τὸν τόπον μου, ἕως οὗ ἀφανισθῶσι, καὶ ζητήσουσι τὸ πρόσωπόν μου.

Ἐν θλίψει αὐτῶν ὀρθριοῦσι πρὸς μὲ, λέγοντες,

\ch{6}πορευθῶμεν, καὶ ἐπιστρέψωμεν πρὸς Κύριον τὸν Θεὸν ἡμῶν, ὅτι αὐτὸς ἥρπακε, καὶ ἰάσεται ἡμᾶς· πατάξει,
\vs{2}καὶ μοτώσει ἡμᾶς, ὑγιάσει ἡμᾶς μετὰ δύο ἡμέρας· ἐν τῇ ἡμέρᾳ τῇ τρίτῃ ἐξαναστησόμεθα,
\vs{3}καὶ ζησόμεθα ἐνώπιον αὐτοῦ, καὶ γνωσόμεθα· διώξωμεν τοῦ γνῶναι τὸν Κύριον· ὡς ὄρθρον ἕτοιμον εὑρήσομεν αὐτὸν, καὶ ἥξει ὡς ὑετὸς ἡμῖν πρώϊμος καὶ ὄψιμος γῇ.

\vs{4}Τί σοι ποιήσω Ἐφραίμ; τί σοι ποιήσω Ἰούδα; τὸ δὲ ἔλεος ὑμῶν ὡς νεφέλη πρωϊνὴ, καὶ ὡς δρόσος ὀρθρινὴ πορευομένη.
\vs{5}Διὰ τοῦτο ἀπεθέρισα τοὺς προφήτας ὑμῶν· ἀπέκτεινα αὐτοὺς ἐν ῥήματι στόματός μου· καὶ τὸ κρίμα μου ὡς φῶς ἐξελεύσεται.

\vs{6}Διότι ἔλεος θέλω ἢ θυσίαν, καὶ ἐπίγνωσιν Θεοῦ ἢ ὁλοκαυτώματα.
\vs{7}Αὐτοὶ δέ εἰσιν ὡς ἄνθρωπος παραβαίνων διαθήκην· ἐκεῖ κατεφρόνησέ μου Γαλαὰδ πόλις,
\vs{8}ἐργαζομένη μάταια, ταράσσουσα ὕδωρ,
\vs{9}καὶ ἡ ἰσχύς σου ἀνδρὸς πειρατοῦ· ἔκρυψαν ἱερεῖς ὁδόν, ἐφόνευσαν Σίκιμα, ὅτι ἀνομίαν ἐποίησαν ἐν τῷ οἴκῳ τοῦ Ἰσραήλ· εἶδον φρικώδη ἐκεῖ,
\vs{10}πορνείαν τοῦ Ἐφραίμ· ἐμιάνθη Ἰσραὴλ καὶ Ἰούδα· ἄρχου τρυγᾷν σεαυτῷ,
\vs{11}ἐν τῷ ἐπιστρέφειν με τὴν αἰχμαλωσίαν τοῦ λαοῦ μου.

\ch{7}
Ἐν τῷ ἰάσασθαί με τὸν Ἰσραὴλ, καὶ ἀποκαλυφθήσεται ἡ ἀδικία Ἐφραὶμ, καὶ ἡ κακία Σαμαρείας, ὅτι εἰργάσαντο ψευδῆ· καὶ κλέπτης πρὸς αὐτὸν εἰσελεύσεται, ἐκδιδύσκων λῃστὴς ἐν τῇ ὁδῷ αὐτοῦ,
\vs{2}ὅπως συνᾴδωσιν ὡς ᾄοντες τῇ καρδίᾳ αὐτῶν· πάσας τὰς κακίας αὐτῶν ἐμνήσθην· νῦν ἐκύκλωσαν αὐτοὺς τὰ διαβούλια αὐτῶν, ἀπέναντι τοῦ προσώπου μου ἐγένοντο.
\vs{3}Ἐν ταῖς κακίαις αὐτῶν εὔφραναν βασιλεῖς, καὶ ἐν τοῖς ψεύδεσιν αὐτῶν ἄρχοντας.
\vs{4}Πάντες μοιχεύοντες ὡς κλίβανος καιόμενος εἰς πέψιν κατακαύματος ἀπὸ τῆς φλογὸς, ἀπὸ φυράσεως στέατος, ἕως τοῦ ζυμωθῆναι αὐτό.
\vs{5}Ἡμέραι τῶν βασιλέων ὑμῶν, ἤρξαντο οἱ ἄρχοντες θυμοῦσθαι ἐξ οἴνου, ἐξέτεινε τὴν χεῖρα αὐτοῦ μετὰ λοιμῶν.
\vs{6}Διότι ἀνεκαύθησαν ὡς κλίβανος αἱ καρδίαι αὐτῶν, ἐν τῷ καταράσσειν αὐτοὺς ὅλην τὴν νύκτα· ὕπνου Ἐφραὶμ ἐνεπλήσθη, πρωῒ ἐνεγενήθη, ἀνεκαύθη ὡς πυρὸς φέγγος.
\vs{7}Πάντες ἐθερμάνθησαν ὡς κλίβανος, καὶ κατέφαγον τοὺς κριτὰς αὐτῶν· πάντες οἱ βασιλεῖς αὐτῶν ἔπεσαν, οὐκ ἦν ἐν αὐτοῖς ὁ ἐπικαλούμενος πρὸς μέ.

\vs{8}Ἐφραὶμ ἐν τοῖς λαοῖς αὐτοῦ συνεμίγνυτο, Ἐφραὶμ ἐγένετο ἐγκρυφίας, οὐ μεταστρεφόμενος.
\vs{9}Κατέφαγον ἀλλότριοι τὴν ἰσχὺν αὐτοῦ, αὐτὸς δὲ οὐκ ἔγνω· καὶ πολιαὶ ἐξήνθησαν αὐτῷ, καὶ αὐτὸς οὐκ ἔγνω.
\vs{10}Καὶ ταπεινωθήσεται ἡ ὕβρις Ἰσραὴλ εἰς πρόσωπον αὐτοῦ· καὶ οὐκ ἐπέστρεψαν πρὸς Κύριον τὸν Θεὸν αὐτῶν, καὶ οὐκ ἐξεζήτησαν αὐτὸν ἐν πᾶσι τούτοις.

\vs{11}Καὶ ἦν Ἐφραὶμ ὡς περιστερὰ ἄνους, οὐκ ἔχουσα καρδίαν· Αἴγυπτον ἐπεκαλεῖτο, καὶ εἰς Ἀσσυρίους ἐπορεύθησαν·
\vs{12}Καθὼς ἂν πορεύωνται, ἐπιβαλῶ ἐπʼ αὐτοὺς τὸ δίκτυόν μου, καθὼς τὰ πετεινὰ τοῦ οὐρανοῦ κατάξω αὐτοὺς, παιδεύσω αὐτοὺς ἐν τῇ ἀκοῇ τῆς θλίψεως αὐτῶν.

\vs{13}Οὐαὶ αὐτοῖς, ὅτι ἀπεπήδησαν ἀπʼ ἐμοῦ· δείλαιοί εἰσιν, ὅτι ἠσέβησαν εἰς ἐμέ· ἐγὼ δὲ ἐλυτρωσάμην αὐτοὺς, αὐτοὶ δὲ κατελάλησαν κατʼ ἐμοῦ ψευδῆ.
\vs{14}Καὶ οὐκ ἐβόησαν πρὸς μὲ αἱ καρδίαι αὐτῶν, ἀλλʼ ἢ ὠλόλυζον ἐν ταῖς κοίταις αὐτῶν· ἐπὶ σίτῳ καὶ οἴνῳ κατετέμνοντο.
\vs{15}Ἐπαιδεύθησαν ἐν ἐμοὶ, κᾀγὼ κατίσχυσα τοὺς βραχίονας αὐτῶν, καὶ εἰς ἐμὲ ἐλογίσαντο πονηρά.
\vs{16}Ἀπεστράφησαν εἰς οὐδὲν, ἐγένοντο ὡς τόξον ἐντεταμένον· πεσοῦνται ἐν ῥομφαίᾳ οἱ ἄρχοντες αὐτῶν διʼ ἀπαιδευσίαν γλώσσης αὐτῶν· οὗτος ὁ φαυλισμὸς αὐτῶν ἐν γῇ Αἰγύπτῳ.

\ch{8}
Εἰς κόλπον αὐτῶν, ὡς γῆ, ὡς ἀετὸς ἐπʼ οἶκον Κυρίου, ἀνθʼ ὧν παρέβησαν τὴν διαθήκην μου, καὶ κατὰ τοῦ νόμου μου ἠσέβησαν.
\vs{2}Ἐμὲ κεκράξονται, ὁ Θεὸς, ἐγνώκαμέν σε·
\vs{3}Ὅτι Ἰσραὴλ ἀπεστρέψατο ἀγαθὰ, ἐχθρὸν κατεδίωξαν.
\vs{4}Ἑαυτοῖς ἐβασίλευσαν, καὶ οὐ διʼ ἐμοῦ, ἦρξαν, καὶ οὐκ ἐγνώρισάν μοι· τὸ ἀργύριον αὐτῶν καὶ τὸ χρυσίον αὐτῶν ἐποίησαν ἑαυτοῖς εἴδωλα, ὅπως ἐξολοθρευθῶσιν.

\vs{5}Ἀπότριψαι τὸν μόσχον σου Σαμάρεια, παρωξύνθη ὁ θυμός μου ἐπʼ αὐτούς· ἕως τίνος οὐ μὴ δύνωνται καθαρισθῆναι
\vs{6}ἐν τῷ Ἰσραήλ; καὶ αὐτὸ τέκτων ἐποίησε, καὶ οὐ Θεός ἐστι· διότι πλανῶν ἦν ὁ μόσχος σου, Σαμάρεια.
\vs{7}Ὅτι ἀνεμόφθορα ἔσπειραν, καὶ ἡ καταστροφὴ αὐτῶν ἐκδέξεται αὐτά· δράγμα οὐκ ἔχον ἰσχὺν τοῦ ποιῆσαι ἄλευρον· ἐὰν δὲ καὶ ποιήσῃ, ἀλλότριοι καταφάγονται αὐτό.
\vs{8}Κατεπόθη Ἰσραὴλ, νῦν ἐγένετο ἐν τοῖς ἔθνεσιν ὡς σκεῦος ἄχρηστον,
\vs{9}ὅτι αὐτοὶ ἀνέβησαν εἰς Ἀσσυρίους· ἀνέθαλε καθʼ ἑαυτὸν Ἐφραίμ· δῶρα ἠγάπησαν,
\vs{10}διὰ τοῦτο παραδοθήσονται ἐν τοῖς ἔθνεσι· νῦν εἰσδέξομαι αὐτοὺς, καὶ κοπάσουσι μικρὸν τοῦ χρίειν βασιλέα καὶ ἄρχοντας.

\vs{11}Ὅτι ἐπλήθυνεν Ἐφραὶμ θυσιαστήρια, εἰς ἁμαρτίας ἐγένοντο αὐτῷ θυσιαστήρια ἠγαπημένα.
\vs{12}Καταγράψω αὐτῷ πλῆθος, καὶ τὰ νόμιμα αὐτοῦ εἰς ἀλλότρια ἐλογίσθησαν, θυσιαστήρια τὰ ἠγαπημένα.
\vs{13}Διότι ἐὰν θύσωσι θυσίαν, καὶ φάγωσι κρέα, Κύριος οὐ προσδέξεται αὐτά· νῦν μνησθήσεται τὰς ἀδικίας αὐτῶν, καὶ ἐκδικήσει τὰς ἁμαρτίας αὐτῶν· αὐτοὶ εἰς Αἴγυπτον ἀπέστρεψαν, καὶ ἐν Ἀσσυρίοις ἀκάθαρτα φάγονται.
\vs{14}Καὶ ἐπελάθετο Ἰσραὴλ τοῦ ποιήσαντος αὐτόν, καὶ ᾠκοδόμησαν τεμένη· καὶ Ἰούδας ἐπλήθυνε πόλεις τετειχισμένας· καὶ ἐξαποστελῶ πῦρ εἰς τὰς πόλεις αὐτοῦ, καὶ καταφάγεται τὰ θεμέλια αὐτῶν.

\ch{9}
Μὴ χαῖρε Ἰσραήλ, μηδὲ εὐφραίνου καθὼς οἱ λαοὶ, διότι ἐπόρνευσας ἀπὸ τοῦ Θεοῦ σου· ἠγάπησας δόματα ἐπὶ πάντα ἅλωνα σίτου.
\vs{2}Ἅλων καὶ ληνὸς οὐκ ἔγνω αὐτούς, καὶ ὁ οἶνος ἐψεύσατο αὐτούς.
\vs{3}Οὐ κατῴκησαν ἐν τῇ γῇ τοῦ Κυρίου· κατῴκησεν Ἐφραὶμ Αἴγυπτον, καὶ ἐν Ἀσσυρίοις ἀκάθαρτα φάγονται.
\vs{4}Οὐκ ἔσπεισαν τῷ Κυρίῳ οἶνον, καὶ οὐχ ἥδυναν αὐτῷ αἱ θυσίαι αὐτῶν, ὡς ἄρτος πένθους αὐτοῖς· πάντες οἱ ἐσθίοντες αὐτὰ μιανθήσονται, διότι οἱ ἄρτοι αὐτῶν ταῖς ψυχαῖς αὐτῶν οὐκ εἰσελεύσονται εἰς τὸν οἶκον Κυρίου.

\vs{5}Τί ποιήσετε ἐν ἡμέραις πανηγύρεως, καὶ ἐν ἡμέρᾳ ἑορτῆς τοῦ Κυρίου;
\vs{6}Διὰ τοῦτο ἰδοὺ πορεύονται ἐκ ταλαιπωρίας Αἰγύπτου, καὶ ἐκδέξεται αὐτοὺς Μέμφις, καὶ θάψει αὐτοὺς Μαχμάς· τὸ ἀργύριον αὐτῶν ὄλεθρος κληρονομήσει αὐτὸ, ἄκανθαι ἐν τοῖς σκηνώμασιν αὐτῶν.

\vs{7}Ἥκασιν αἱ ἡμέραι τῆς ἐκδικήσεως, ἥκασιν αἱ ἡμέραι τῆς ἀνταποδόσεώς σου, καὶ κακωθήσεται Ἰσραὴλ ὥσπερ ὁ προφήτης ὁ παρεξεστηκώς, ἄνθρωπος ὁ πνευματοφόρος· ὑπὸ τοῦ πλήθους τῶν ἀδικιῶν σου ἐπληθύνθη μανία σου.
\vs{8}Σκοπὸς Ἐφραὶμ μετὰ Θεοῦ· προφήτης παγὶς σκολιὰ ἐπὶ πάσας τὰς ὁδοὺς αὐτοῦ, μανίαν ἐν οἴκῳ Θεοῦ κατέπηξαν.
\vs{9}Ἐφθάρησαν κατὰ τὰς ἡμέρας τοῦ βουνοῦ, μνησθήσεται ἀδικίας αὐτῶν, ἐκδικήσει ἁμαρτίας αὐτῶν.

\vs{10}Ὡς σταφυλὴν ἐν ἐρήμῳ εὗρον τὸν Ἰσραὴλ, καὶ ὡς σκοπὸν ἐν συκῇ πρώϊμον πατέρας αὐτῶν εἶδον· αὐτοὶ εἰσῆλθον πρὸς τὸν Βεελφεγὼρ, καὶ ἀπηλλοτριώθησαν εἰς αἰσχύνην, καὶ ἐγένοντο οἱ ἐβδελυγμένοι ὡς οἱ ἠγαπημένοι.
\vs{11}Ἐφραὶμ ὡς ὄρνεον ἐξεπετάσθη, αἱ δόξαι αὐτῶν ἐκ τόκων καὶ ὠδίνων καὶ συλλήψεων·
\vs{12}Διότι καὶ ἐὰν ἐκθρέψωσι τὰ τέκνα αὐτῶν, ἀτεκνωθήσονται ἐξ ἀνθρώπων· διότι καὶ οὐαὶ αὐτοῖς ἐστι· σάρξ μου ἐξ αὐτῶν.
\vs{13}Ἐφραὶμ, ὃν τρόπον εἶδον, εἰς θήραν παρέστησαν τὰ τέκνα αὐτῶν, καὶ Ἐφραὶμ, τοῦ ἐξαγαγεῖν εἰς ἀποκέντησιν τὰ τέκνα αὐτοῦ.

\vs{14}Δὸς αὐτοῖς Κύριε, τί δώσεις αὐτοῖς; μήτραν ἀτεκνοῦσαν, καὶ μαστοὺς ξηρούς.
\vs{15}Πᾶσαι αἱ κακίαι αὐτῶν ἐν Γαλγάλ, ὅτι ἐκεῖ ἐμίσησα αὐτούς· διὰ τὰς κακίας τῶν ἐπιτηδευμάτων αὐτῶν, ἐκ τοῦ οἴκου μου ἐκβαλῶ αὐτοὺς, οὐ μὴ προσθήσω τοῦ ἀγαπῆσαι αὐτούς· πάντες οἱ ἄρχοντες αὐτῶν ἀπειθοῦντες.
\vs{16}Ἐπόνεσεν Ἐφραίμ· τὰς ῥίζας αὐτοῦ ἐξηράνθη, καρπὸν οὐκ ἔτι μὴ ἐνέγκῃ· διότι καὶ ἐὰν γεννήσωσιν, ἀποκτενῶ τὰ ἐπιθυμήματα κοιλίας αὐτῶν.
\vs{17}Ἀπώσεται αὐτοὺς ὁ Θεός, ὅτι οὐκ εἰσήκουσαν αὐτοῦ, καὶ ἔσονται πλανῆται ἐν τοῖς ἔθνεσιν.

\ch{10}
Ἄμπελος εὐκηματοῦσα Ἰσραὴλ, ὁ καρπὸς εὐθηνῶν αὐτῆς· κατὰ τὸ πλῆθος τῶν καρπῶν αὐτῆς, ἐπλήθυνε τὰ θυσιαστήρια· κατὰ τὰ ἀγαθὰ τῆς γῆς αὐτοῦ, ᾠκοδόμησε στήλας.
\vs{2}Ἐμέρισαν καρδίας αὐτῶν, νῦν ἀφανισθήσονται· αὐτὸς κατασκάψει τὰ θυσιαστήρια αὐτῶν, ταλαιπωρήσουσιν αἱ στῆλαι αὐτῶν.

\vs{3}Διότι νῦν ἐροῦσιν, οὐκ ἔστι βασιλεὺς ἡμῖν, ὅτι οὐκ ἐφοβήθημεν τὸν Κύριον· ὁ δὲ βασιλεὺς τί ποιήσει ἡμῖν,
\vs{4}λαλῶν ῥήματα προφάσεις ψευδεῖς; διαθήσεται διαθήκην, ἀνατελεῖ ὡς ἄγρωστις κρίμα ἐπὶ χέρσον ἀγροῦ.
\vs{5}Τῷ μόσχῳ τοῦ οἴκου Ὦν παροικήσουσιν οἱ κατοικοῦντες Σαμάρειαν, ὅτι ἐπένθησε λαὸς αὐτοῦ ἐπʼ αὐτόν· καὶ καθὼς παρεπίκραναν αὐτὸν, ἐπιχαροῦνται ἐπὶ τὴν δόξαν αὐτοῦ, ὅτι μετῳκίσθη ἀπʼ αὐτοῦ.
\vs{6}Καὶ αὐτὸν εἰς Ἀσσυρίους δήσαντες, ἀπήνεγκαν ξένια τῷ βασιλεῖ Ἰαρείμ· ἐν δόματι Ἐφραὶμ δέξεται, καὶ αἰσχυνθήσεται Ἰσραὴλ ἐν τῇ βουλῇ αὐτοῦ.
\vs{7}Ἀπέῤῥιψεν Σαμάρεια βασιλέα αὐτῆς ὡς φρύγανον ἐπὶ προσώπου ὕδατος·
\vs{8}Καὶ ἐξαρθήσονται βωμοὶ Ὦν ἁμαρτήματα τοῦ Ἰσραὴλ, ἄκανθαι καὶ τρίβολοι ἀναβήσονται ἐπὶ τὰ θυσιαστήρια αὐτῶν· καὶ ἐροῦσι τοῖς ὄρεσι, καλύψατε ἡμᾶς, καὶ τοῖς βουνοῖς, πέσατε ἐφʼ ἡμᾶς.

\vs{9}Ἀφʼ οὗ οἱ βουνοὶ, ἥμαρτεν Ἰσραήλ, ἐκεῖ ἔστησαν· οὐ μὴ καταλάβῃ αὐτοὺς ἐν τῷ βουνῷ πόλεμος ἐπὶ τὰ τὲκνα ἀδικίας
\vs{10}παιδεῦσαι αὐτούς· καὶ συναχθήσονται ἐπʼ αὐτοὺς λαοὶ, ἐν τῷ παιδεύεσθαι αὐτοὺς ἐν ταῖς δυσὶν ἀδικίαις αὐτῶν.
\vs{11}Ἐφραὶμ δάμαλις δεδιδαγμένη ἀγαπᾷν νῖκος, ἐγὼ δὲ ἐπελεύσομαι ἐπὶ τὸ κάλλιστον τοῦ τραχήλου αὐτῆς· ἐπιβιβῶ Ἐφραὶμ, παρασιωπήσομαι Ἰούδαν, ἐνισχύσει αὐτῷ Ἰακώβ.

\vs{12}Σπείρατε ἑαυτοῖς εἰς δικαιοσύνην, τρυγήσατε εἰς καρπὸν ζωῆς, φωτίσατε ἑαυτοῖς φῶς γνώσεως, ἐκζητήσατε τὸν Κύριον ἕως τοῦ ἐλθεῖν γεννήματα δικαιοσύνης ὑμῖν.
\vs{13}Ἱνατί παρεσιωπήσατε ἀσέβειαν, καὶ τὰς ἀδικίας αὐτῆς ἐτρυγήσατε; ἐφάγετε καρπὸν ψευδῆ, ὅτι ἤλπισας ἐν τοῖς ἁμαρτήμασί σου, ἐν πλήθει δυνάμεώς σου.
\vs{14}Καὶ ἐξαναστήσεται ἀπώλεια ἐν τῷ λαῷ σου, καὶ πάντα τὰ περιτετειχισμένα σου οἰχήσεται· ὡς ἄρχεν Σαλαμὰν ἐκ τοῦ οἴκου τοῦ Ἱεροβοὰμ, ἐν ἡμέραις πολέμου μητέρα ἐπὶ τέκνοις ἠδάφισαν,
\vs{15}οὕτως ποιήσω ὑμῖν οἶκος τοῦ Ἰσραὴλ ἀπὸ προσώπου ἀδικίας κακιῶν ὑμῶν.

\ch{11}
Ὄρθρου ἀπεῤῥίφησαν, ἀπεῤῥίφη βασιλεὺς Ἰσραήλ· ὅτι νήπιος Ἰσραὴλ, καὶ ἐγὼ ἠγάπησα αὐτόν, καὶ ἐξ Αἰγύπτου μετεκάλεσα τὰ τέκνα αὐτοῦ.
\vs{2}Καθὼς μετεκάλεσα αὐτούς, οὕτως ἀπῴχοντο ἐκ προσώπου μου· αὐτοὶ τοῖς Βααλεὶμ ἔθυον, καὶ τοῖς γλυπτοῖς ἐθυμίων.
\vs{3}Καὶ ἐγὼ συνεπόδισα τὸν Ἐφραὶμ, ἀνέλαβον αὐτὸν ἐπὶ τὸν βραχίονά μου, καὶ οὐκ ἔγνωσαν ὅτι ἴαμαι αὐτούς.
\vs{4}Ἐν διαφθορᾷ ἀνθρώπων ἐξέτεινα αὐτοὺς ἐν δεσμοῖς ἀγαπήσεώς μου, καὶ ἔσομαι αὐτοῖς ὡς ῥαπίζων ἄνθρωπος ἐπὶ τὰς σιαγόνας αὐτοῦ· καὶ ἐπιβλέψομαι πρὸς αὐτόν, δυνήσομαι αὐτῷ.

\vs{5}Κατῴκησεν Ἐφραὶμ ἐν Αἰγύπτῳ, καὶ Ἀσσοὺρ αὐτὸς βασιλεὺς αὐτοῦ· ὅτι οὐκ ἠθέλησεν ἐπιστρέψαι,
\vs{6}καὶ ἠσθένησεν ἐν ῥομφαίᾳ ἐν ταῖς πόλεσιν αὐτοῦ· καὶ κατέπαυσεν ἐν ταῖς χερσὶν αὐτοῦ, καὶ φάγονται ἐκ τῶν διαβουλίων αὐτῶν,
\vs{7}καὶ ὁ λαὸς αὐτοῦ ἐπικρεμάμενος ἐκ τῆς κατοικίας αὐτοῦ· καὶ ὁ Θεὸς ἐπὶ τὰ τίμια αὐτοῦ θυμωθήσεται, καὶ οὐ μὴ ὑψώσῃ αὐτόν.

\vs{8}Τί σε διαθῶμαι Ἐφραίμ; ὑπερασπιῶ σου Ἰσραήλ; τί σε διαθῶ; ὡς Ἀδάμα θήσομαί σε, καὶ ὡς Σεβοείμ· μετεστράφη ἡ καρδία μου ἐν τῷ αὐτῷ, συνεταράχθη ἡ μεταμέλειά μου·
\vs{9}Οὐ μὴ ποιήσω κατὰ τὴν ὀργὴν τοῦ θυμοῦ μου· οὐ μὴ ἐγκαταλίπω τοῦ ἐξαλειφθῆναι τὸν Ἐφραὶμ· διότι Θεὸς ἐγώ εἰμι, καὶ οὐκ ἄνθρωπος, ἐν σοὶ ἅγιος, καὶ οὐκ εἰσελεύσομαι εἰς πόλιν.
\vs{10}Ὀπίσω Κυρίου πορεύσομαι· ὡς λέων ἐρεύξεται, ὅτι αὐτὸς ὠρύσεται, καὶ ἐκστήσονται τέκνα ὑδάτων.
\vs{11}Ἐκστήσονται ὡς ὄρνεον ἐξ Αἰγύπτου, καὶ ὡς περιστερὰ ἐκ γῆς Ἀσσυρίων· καὶ ἀποκαταστήσω αὐτοὺς εἰς τοὺς οἴκους αὐτῶν, λέγει Κύριος.

\ch{12}
Ἐκύκλωσέ με ἐν ψεύδει Ἐφραὶμ, καὶ ἐν ἀσεβείαις οἶκος Ἰσραὴλ, καὶ Ἰούδα· νῦν ἔγνω αὐτοὺς ὁ Θεὸς, καὶ ὁ λαὸς ἅγιος κεκλήσεται Θεοῦ.

\vs{2}Ὁ δὲ Ἐφραὶμ πονηρὸν πνεῦμα, ἐδίωξε καύσωνα ὅλην τὴν ἡμέραν· κενὰ καὶ μάταια ἐπλήθυνε, καὶ διαθήκην μετὰ Ἀσσυρίων διέθετο, καὶ ἔλαιον εἰς Αἴγυπτον ἐνεπορεύετο.
\vs{3}Καὶ κρίσις τῷ Κυρίῳ πρὸς Ἰούδαν, τοῦ ἐκδικῆσαι τὸν Ἰακώβ· κατὰ τὰς ὁδοὺς αὐτοῦ καὶ κατὰ τὰ ἐπιτηδεύματα αὐτοῦ ἀποδώσει αὐτῷ.

\vs{4}Ἐν τῇ κοιλίᾳ ἐπτέρνισε τὸν ἀδελφὸν αὐτοῦ, καὶ ἐν κόποις αὐτοῦ ἐνίσχυσε πρὸς Θεόν.
\vs{5}Καὶ ἐνίσχυσε μετὰ ἀγγέλου, καὶ ἠδυνάσθη· ἔκλαυσαν, καὶ ἐδεήθησάν μου· ἐν τῷ οἴκῳ Ὦν εὕροσάν με, καὶ ἐκεῖ ἐλαλήθη πρὸς αὐτούς.
\vs{6}Ὁ δὲ Κύριος ὁ Θεὸς ὁ παντοκράτωρ ἔσται μνημόσυνον αὐτοῦ.
\vs{7}Καὶ σὺ ἐν Θεῷ σου ἐπιστρέψεις, ἔλεον καὶ κρίμα φυλάσσου, καὶ ἔγγιζε πρὸς τὸν Θεόν σου διαπαντός.

\vs{8}Χαναὰν, ἐν χειρὶ αὐτοῦ ζυγὸς ἀδικίας, καταδυναστεύειν ἠγάπησε.
\vs{9}Καὶ εἶπεν Ἐφραὶμ, πλὴν πεπλούτηκα, εὕρηκα ἀναψυχὴν ἐμαυτῷ· πάντες οἱ πόνοι αὐτοῦ οὐχ εὑρεθήσονται αὐτῷ, διʼ ἀδικίας ἃς ἥμαρτεν.
\vs{10}Ἐγὼ δὲ Κύριος ὁ Θεός σου ἀνήγαγόν σε ἐκ γῆς Αἰγύπτου, ἔτι κατοικιῶ σε ἐν σκηναῖς, καθὼς ἡμέραι ἑορτῆς.
\vs{11}Καὶ λαλήσω πρὸς προφήτας, καὶ ἐγὼ ὁράσεις ἐπλήθυνα, καὶ ἐν χερσὶ προφητῶν ὡμοιώθην.
\vs{12}Εἰ μὴ Γαλαάδ ἐστιν, ἄρα ψευδεῖς ἦσαν ἐν Γαλαὰδ ἄρχοντες θυσιάζοντες, καὶ τὰ θυσιαστήρια αὐτῶν, ὡς χελῶναι ἐπὶ χέρσον ἀγροῦ.

\vs{13}Καὶ ἀνεχώρησεν Ἰακὼβ εἰς πεδίον Συρίας, καὶ ἐδούλευσεν Ἰσραὴλ ἐν γυναικὶ, καὶ γυναικὶ ἐφυλάξατο.
\vs{14}Καὶ ἐν προφήτῃ ἀνήγαγεν Κύριος τὸν Ἰσραὴλ ἐκ γῆς Αἰγύπτου, καὶ ἐν προφήτῃ διεφυλάχθη.
\vs{15}Ἐθύμωσεν Ἐφραὶμ, καὶ παρώργισε, καὶ τὸ αἷμα αὐτοῦ ἐπʼ αὐτὸν ἐκχυθήσεται, καὶ τὸν ὀνειδισμὸν αὐτοῦ ἀνταποδώσει Κύριος αὐτῷ.

\ch{13}
Κατὰ τὸν λόγον Ἐφραὶμ δικαιώματα ἔλαβεν αὐτὸς ἐν τῷ Ἰσραήλ, καὶ ἔθετο αὐτὰ τῇ Βάαλ καὶ ἀπέθανε.
\vs{2}Καὶ νῦν προσέθεντο τοῦ ἁμαρτάνειν, καὶ ἐποίησαν ἑαυτοῖς χώνευμα ἐκ τοῦ ἀργυρίου αὐτῶν, κατʼ εἰκόνα εἰδώλων, ἔργα τεκτόνων συντετελεσμένα αὐτοῖς· αὐτοὶ λέγουσι, θύσατε ἀνθρώπους, μόσχοι γὰρ ἐκλελοίπασι.
\vs{3}Διὰ τοῦτο ἔσονται ὡς νεφέλη πρωϊνὴ καὶ ὡς δρόσος ὀρθρινὴ πορευομένη, ὡς χνοῦς ἀποφυσώμενος ἀφʼ ἅλωνος, καὶ ὡς ἀτμὶς ἀπὸ δακρύων.
\vs{4}Ἐγὼ δὲ Κύριος ὁ Θεός σου ὁ στερεῶν τὸν οὐρανὸν, καὶ κτίζων γῆν, οὗ αἱ χεῖρες ἔκτισαν πᾶσαν τὴν στρατιὰν τοῦ οὐρανοῦ, καὶ οὐ παρέδειξά σοι αὐτὰ τοῦ πορεύεσθαι ὀπίσω αὐτῶν· καὶ ἐγὼ ἀνήγαγόν σε ἐκ γῆς Αἰγύπτου, καὶ θεὸν πλὴν ἐμοῦ οὐ γνώσῃ, καὶ σώζων οὐκ ἔστι πάρεξ ἐμοῦ.
\vs{5}Ἐγὼ ἐποίμαινόν σε ἐν τῇ ἐρήμῳ, ἐν γῇ ἀοικήτῳ
\vs{6}κατὰ τὰς νομὰς αὐτῶν, καὶ ἐνεπλήσθησαν εἰς πλησμονήν, καὶ ὑψώθησαν αἱ καρδίαι αὐτῶν· ἕνεκα τούτου ἐπελάθοντό μου.
\vs{7}Καὶ ἔσομαι αὐτοῖς ὡς πανθὴρ, καὶ ὡς πάρδαλις· κατὰ τὴν ὁδὸν Ἀσσυρίων
\vs{8}ἀπαντήσομαι αὐτοῖς ὡς ἄρκτος ἡ ἀπορουμένη, καὶ διαῤῥήξω συγκλεισμὸν καρδίας αὐτῶν, καὶ καταφάγονται αὐτοὺς ἐκεῖ σκύμνοι δρυμοῦ, θηρία ἀγροῦ διασπάσει αὐτούς.

\vs{9}Τῇ διαφθορᾷ σου Ἰσραήλ τίς βοηθήσει;
\vs{10}Ποῦ ὁ βασιλεύς σου οὗτος; καὶ διασωσάτω σε ἐν πάσαις ταῖς πόλεσί σου· κρινάτω σε ὃν εἶπας, δός μοι βασιλέα καὶ ἄρχοντα.
\vs{11}Καὶ ἔδωκά σοι βασιλέα ἐν ὀργῇ μου, καὶ ἔσχον ἐν τῷ θυμῷ μου.

\vs{12}Συστροφὴν ἀδικίας Ἐφραὶμ, ἐγκεκρυμμένη ἡ ἁμαρτία αὐτοῦ,
\vs{13}ὠδῖνες ὡς τικτούσης ἥξουσιν αὐτῷ· οὗτος ὁ υἱός σου ὁ φρόνιμος, διότι οὐ μὴ ὑποστῇ ἐν συντριβῇ τέκνων.
\vs{14}Ἐκ χειρὸς ᾅδου ῥύσομαι, καὶ ἐκ θανάτου λυτρώσομαι αὐτούς· ποῦ ἡ δίκη σου θάνατε; ποῦ τὸ κέντρον σου ᾅδη; παράκλησις κέκρυπται ἀπὸ ὀφθαλμῶν μου.

\vs{15}Διότι οὗτος ἀναμέσον ἀδελφῶν διαστελεῖ, ἐπάξει καύσωνα ἄνεμον Κύριος ἐκ τῆς ἐρήμου ἐπʼ αὐτόν, καὶ ἀναξηρανεῖ τὰς φλέβας αὐτοῦ, ἐξερημώσει τὰς πηγὰς αὐτοῦ· αὐτὸς καταξηρανεῖ τὴν γῆν αὐτοῦ, καὶ πάντα τὰ σκεύη τὰ ἐπιθυμητὰ αὐτοῦ.

\ch{14}
Ἀφανισθήσεται Σαμάρεια, ὅτι ἀντέστη πρὸς τὸν Θεὸν αὐτῆς· ἐν ῥομφαίᾳ πεσοῦνται αὐτοὶ, καὶ τὰ ὑποτίτθια αὐτῶν ἐδαφισθήσονται, καὶ αἱ ἐν γαστρὶ ἔχουσαι αὐτῶν διαῤῥαγήσονται.

\vs{2}Ἐπιστράφηθι Ἰσραὴλ πρὸς Κύριον τὸν Θεόν σου, διότι ἠσθένησαν ἐν ταῖς ἀδικίαις σου.
\vs{3}Λάβετε μεθʼ ἑαυτῶν λόγους, καὶ ἐπιστράφητε πρὸς Κύριον τὸν Θεὸν ὑμῶν· εἴπατε αὐτῷ, ὅπως μὴ λάβητε ἀδικίαν, καὶ λάβητε ἀγαθὰ, καὶ ἀνταποδώσομεν καρπὸν χειλέων ἡμῶν.
\vs{4}Ἀσσοὺρ οὐ μὴ σώσῃ ἡμᾶς, ἐφʼ ἵππον οὐκ ἀναβησόμεθα· οὐκέτι μὴ εἴπωμεν, θεοὶ ἡμῶν, τοῖς ἔργοις τῶν χειρῶν ἡμῶν· ὁ ἐν σοὶ ἐλεήσει ὀρφανόν.

\vs{5}Ἰάσομαι τὰς κατοικίας αὐτῶν, ἀγαπήσω αὐτοὺς ὁμολόγως, ὅτι ἀπέστρεψε τὴν ὀργήν μου ἀπʼ αὐτοῦ.
\vs{6}Ἔσομαι ὡς δρόσος τῷ Ἰσραὴλ, ἀνθήσει ὡς κρίνον, καὶ βαλεῖ τὰς ῥίζας αὐτοῦ ὡς ὁ Λίβανος·
\vs{7}Πορεύσονται οἱ κλάδοι αὐτοῦ, καὶ ἔσται ὡς ἐλαία κατάκαρπος, καὶ ἡ ὀσφρασία αὐτοῦ ὡς Λιβάνου.
\vs{8}ἐπιστρέψουσι καὶ καθιοῦνται ὑπὸ τὴν σκέπην αὐτοῦ, ζήσονται καὶ μεθυσθήσονται σίτῳ· καὶ ἐξανθήσει ὡς ἄμπελος· μνημόσυνον αὐτοῦ, ὡς οἶνος Λιβάνου
\vs{9}τῷ Ἐφραίμ· τί αὐτῷ ἔτι καὶ εἰδώλοις; ἐγὼ ἐταπείνωσα αὐτὸν, καὶ κατισχύσω αὐτόν· ἐγὼ ὡς ἄρκευθος πυκάζουσα, ἐξ ἐμοῦ ὁ καρπός σου εὕρηται.

\vs{10}Τίς σοφὸς καὶ συνήσει ταῦτα; ἢ συνετὸς καὶ ἐπιγνώσεται αὐτά; ὅτι εὐθεῖαι αἱ ὁδοὶ τοῦ Κυρίου, καὶ δίκαιοι πορεύσονται ἐν αὐταῖς, οἱ δὲ ἀσεβεῖς ἀσθενήσουσιν ἐν αὐταῖς.


\def\book{ΙΩΗΛ}
\biblebook{ΙΩΗΛ}


\lettrine[lines=2, loversize=0.2, nindent=0em, findent=.25em]{\textcolor{bookheadingcolor}{Λ}}{ΟΓΟΣ} Κυρίου ὃς ἐγενήθη πρὸς Ἰωὴλ τὸν τοῦ Βαθουήλ.

\postdropcapindent\vs{2}Ἀκούσατε ταῦτα, οἱ πρεσβύτεροι, καὶ ἑνωτίσασθε, πάντες οἱ κατοικοῦντες τὴν γῆν. εἰ γέγονεν τοιαῦτα ἐν ταῖς ἡμέραις ἡμῶν, ἢ ἐν ταῖς ἡμέραις τῶν πατέρων ὑμῶν;
\vs{3}ὑπὲρ αὐτῶν τοῖς τέκνοις ὑμῶν διηγήσασθε, καὶ τὰ τέκνα ὑμῶν τοῖς τέκνοις αὐτῶν, καὶ τὰ τέκνα αὐτῶν εἰς γενεὰν ἑτέραν.
\vs{4}τὰ κατάλοιπα τῆς κάμπης κατέφαγεν ἡ ἀκρίς, καὶ τὰ κατάλοιπα τῆς ἀκρίδος κατέφαγεν ὁ βροῦχος, καὶ τὰ κατάλοιπα τοῦ βρούχου κατέφαγεν ἡ ἐρυσίβη.

\vs{5}Ἐκνήψατε, οἱ μεθύοντες, ἐξ οἴνου αὐτῶν καὶ κλαύσατε· θρηνήσατε πάντες οἱ πίνοντες οἶνον εἰς μέθην, ὅτι ἐξῄρθη ἐκ στόματος ὑμῶν εὐφροσύνη καὶ χαρά.
\vs{6}Ὅτι ἔθνος ἀνέβη ἐπὶ τὴν γῆν μου ἰσχυρὸν καὶ ἀναρίθμητον, οἱ ὀδόντες αὐτοῦ ὀδόντες λέοντος, καὶ αἱ μῦλαι αὐτοῦ σκύμνου·
\vs{7}Ἔθετο τὴν ἄμπελόν μου εἰς ἀφανισμὸν, καὶ τὰς συκάς μου εἰς συγκλασμόν· ἐρευνῶν ἐξηρεύνησεν αὐτὴν, καὶ ἔῤῥιψεν· ἐλεύκανε τὰ κλήματα αὐτῆς.

\vs{8}Θρήνησον πρὸς μὲ ὑπὲρ νύμφην περιεζωσμένην σάκκον, ἐπὶ τὸν ἄνδρα αὐτῆς τὸν παρθενικόν.
\vs{9}Ἐξῇρται θυσία καὶ σπονδὴ ἐξ οἴκου Κυρίου· πενθεῖτε οἱ ἱερεῖς οἱ λειτουργοῦντες θυσιαστηρίῳ Κυρίου,
\vs{10}ὅτι τεταλαιπώρηκε τὰ πεδία· πενθείτω ἡ γῆ, ὅτι τεταλαιπώρηκεν σῖτος· ἐξηράνθη οἶνος, ὠλιγώθη ἔλαιον·
\vs{11}Ἐξηράνθησαν γεωργοί· θρηνεῖτε κτήματα ὑπὲρ πυροῦ καὶ κριθῆς, ὅτι ἀπόλωλε τρυγητὸς ἐξ ἀγροῦ·
\vs{12}Ἡ ἄμπελος ἐξηράνθη, καὶ αἱ συκαῖ ὠλιγώθησαν· ῥοὰ, καὶ φοῖνιξ, καὶ μῆλον, καὶ πάντα τὰ ξύλα τοῦ ἀγροῦ ἐξηράνθησαν, ὅτι ᾔσχυναν χαρὰν οἱ υἱοὶ τῶν ἀνθρώπων.

\vs{13}Περιζώσασθε καὶ κόπτεσθε οἱ ἱερεῖς· θρηνεῖτε οἱ λειτουργοῦντες θυσιαστηρίῳ· εἰσέλθετε, ὑπνώσατε ἐν σάκκοις λειτουργοῦντες Θεῷ, ὅτι ἀπέσχηκεν ἐξ οἴκου Θεοῦ ὑμῶν θυσία καὶ σπονδή.

\vs{14}Ἁγιάσατε νηστείαν, κηρύξατε θεραπείαν, συναγάγετε πρεσβυτέρους, πάντας κατοικοῦντας γῆν εἰς οἶκον Θεοῦ ὑμῶν, καὶ κεκράξετε πρὸς Κύριον ἐκτενῶς,

\vs{15}Οἴμοι, οἴμοι, οἴμοι εἰς ἡμέραν, ὅτι ἐγγὺς ἡ ἡμέρα Κυρίου, καὶ ὡς ταλαιπωρία ἐκ ταλαιπωρίας ἥξει.
\vs{16}Κατέναντι τῶν ὀφθαλμῶν ὑμῶν βρώματα ἐξωλθρεύθη, ἐξ οἴκου Θεοῦ ὑμῶν εὐφροσύνη καὶ χαρά.
\vs{17}Ἐσκίρτησαν δαμάλεις ἐπὶ ταῖς φάτναις αὐτῶν, ἠφανίσθησαν θησαυροὶ, κατεσκάφησαν ληνοὶ, ὅτι ἐξηράνθη σῖτος.
\vs{18}Τί ἀποθήσομεν ἑαυτοῖς; ἔκλαυσαν βουκόλια βοῶν, ὅτι οὐχ ὑπῆρχε νομὴ αὐτοῖς· καὶ τὰ ποίμνια τῶν προβάτων ἠφανίσθησαν.
\vs{19}Πρὸς σὲ Κύριε βοήσομαι, ὅτι πῦρ ἀνήλωσε τὰ ὡραῖα τῆς ἐρήμου, καὶ φλὸξ ἀνῆψε πάντα τὰ ξύλα τοῦ ἀγροῦ,
\vs{20}καὶ τὰ κτήνη τοῦ πεδίου ἀνέβλεψαν πρὸς σὲ, ὅτι ἐξηράνθησαν ἀφέσεις ὑδάτων, καὶ πῦρ κατέφαγε τὰ ὡραῖα τῆς ἐρήμου.

\ch{2}
Σαλπίσατε σάλπιγγι ἐν Σιὼν, κηρύξατε ἐν ὄρει ἁγίῳ μου, καὶ συγχυθήτωσαν πάντες οἱ κατοικοῦντες τὴν γῆν, διότι πάρεστιν ἡμέρα Κυρίου, ὅτι ἐγγὺς
\vs{2}ἡμέρα σκότους καὶ γνόφου, ἡμέρα νεφέλης καὶ ὁμίχλης· ὡς ὄρθρος χυθήσεται ἐπὶ τὰ ὄρη λαὸς πολὺς καὶ ἰσχυρὸς, ὅμοιος αὐτῷ οὐ γέγονεν ἀπὸ τοῦ αἰῶνος, καὶ μετʼ αὐτὸν οὐ προστεθήσεται ἕως ἐτῶν εἰς γενεὰς γενεῶν.
\vs{3}Τὰ ἔμπροσθεν αὐτοῦ πῦρ ἀναλίσκον, καὶ τὰ ὀπίσω αὐτοῦ ἀναπτομένη φλόξ· ὡς παράδεισος τρυφῆς ἡ γῆ πρὸ προσώπου αὐτοῦ, καὶ τὰ ὄπισθεν αὐτοῦ πεδίον ἀφανισμοῦ, καὶ ἀνασωζόμενος οὐκ ἔσται αὐτῷ.

\vs{4}Ὡς ὅρασις ἵππων ἡ ὅρασις αὐτῶν, καὶ ὡς ἱππεῖς οὕτως καταδιώξονται.
\vs{5}Ὡς φωνὴ ἁρμάτων ἐπὶ τὰς κορυφὰς τῶν ὀρέων ἐξαλοῦνται, καὶ ὡς φωνὴ φλογὸς πυρὸς κατεσθιούσης καλάμην, καὶ ὡς λαὸς πολὺς καὶ ἰσχυρὸς παρατασσόμενος εἰς πόλεμον.
\vs{6}Ἀπὸ προσώπου αὐτοῦ συντριβήσονται λαοὶ, πᾶν πρόσωπον ὡς πρόσκαυμα χύτρας.
\vs{7}Ὡς μαχηταὶ δραμοῦνται καὶ ὡς ἄνδρες πολεμισταὶ ἀναβήσονται ἐπὶ τὰ τείχη, καὶ ἕκαστος ἐν τῇ ὁδῷ αὐτοῦ πορεύσεται, καὶ οὐ μὴ ἐκκλίνωσι τὰς τρίβους αὐτῶν,
\vs{8}καὶ ἕκαστος ἀπὸ τοῦ ἀδελφοῦ αὐτοῦ οὐκ ἀφέξεται· καταβαρυνόμενοι ἐν τοῖς ὅπλοις αὐτῶν πορεύσονται, καὶ ἐν τοῖς βέλεσιν αὐτῶν πεσοῦνται, καὶ οὐ μὴ συντελεσθῶσι.
\vs{9}Τῆς πόλεως ἐπιλήψονται, καὶ ἐπὶ τῶν τειχέων δραμοῦνται, καὶ ἐπὶ ταῖς οἰκίαις ἀναβήσονται, καὶ διὰ θυρίδων εἰσελεύσονται ὡς κλέπται.
\vs{10}Πρὸ προσώπου αὐτοῦ συγχυθήσεται ἡ γῆ, καὶ σεισθήσεται ὁ οὐρανός· ὁ ἥλιος καὶ ἡ σελήνη συσκοτάσουσι, καὶ ἄστρα δύσουσι τὸ φέγγος αὐτῶν.
\vs{11}Καὶ Κύριος δώσει φωνὴν αὐτοῦ πρὸ προσώπου δυνάμεως αὐτοῦ, ὅτι πολλή ἐστι σφόδρα ἡ παρεμβολὴ αὐτοῦ, ὅτι ἰσχυρὰ ἔργα λόγων αὐτοῦ· διότι μεγάλη ἡ ἡμέρα Κυρίου, ἐπιφανὴς σφόδρα, καὶ τίς ἔσται ἱκανὸς αὐτῇ;

\vs{12}Καὶ νῦν λέγει Κύριος ὁ Θεὸς ὑμῶν, ἐπιστράφητε πρὸς μὲ ἐξ ὅλης τῆς καρδίας ὑμῶν, καὶ ἐν νηστείᾳ, καὶ ἐν κλαυθμῷ, καὶ ἐν κοπετῷ,
\vs{13}καὶ διαῤῥήξατε τὰς καρδίας ὑμῶν, καὶ μὴ τὰ ἱμάτια ὑμῶν· καὶ ἐπιστράφητε πρὸς Κύριον τὸν Θεὸν ὑμῶν, ὅτι ἐλεήμων καὶ οἰκτίρμων ἐστὶ, μακρόθυμος καὶ πολυέλεος, καὶ μετανοῶν ἐπὶ ταῖς κακίαις.
\vs{14}Τίς οἶδεν εἰ ἐπιστρέψει, καὶ μετανοήσει, καὶ ὑπολείψεται ὀπίσω αὐτοῦ εὐλογίαν, καὶ θυσίαν, καὶ σπονδὴν Κυρίῳ τῷ Θεῷ ὑμῶν;

\vs{15}Σαλπίσατε σάλπιγγι ἐν Σιὼν, ἁγιάσατε νηστείαν. κηρύξατε θεραπείαν,
\vs{16}συναγάγετε λαὸν, ἁγιάσατε ἐκκλησίαν, ἐκλέξασθε πρεσβυτέρους, συναγάγετε νήπια θηλάζοντα μαστοὺς, ἐξελθέτω νυμφίος ἐκ τοῦ κοιτῶνος αὐτοῦ, καὶ νύμφη ἐκ τοῦ παστοῦ αὐτῆς.
\vs{17}Ἀναμέσον τῆς κρηπίδος τοῦ θυσιαστηρίου, κλαύσονται οἱ ἱερεῖς οἱ λειτουργοῦντες τῷ Κυρίῳ, καὶ ἐροῦσι, φεῖσαι Κύριε, τοῦ λαοῦ σου, καὶ μὴ δῷς τὴν κληρονομίαν σου εἰς ὄνειδος, τοῦ κατάρξαι αὐτῶν ἔθνη, ὅπως μὴ εἴπωσιν ἐν τοῖς ἔθνεσι, ποῦ ἐστιν ὁ Θεὸς αὐτῶν;

\vs{18}Καὶ ἐζήλωσε Κύριος τὴν γῆν αὐτοῦ, καὶ ἐφείσατο τοῦ λαοῦ αὐτοῦ.
\vs{19}Καὶ ἀπεκρίθη Κύριος, καὶ εἶπε τῷ λαῷ αὐτοῦ, ἰδοὺ ἐγὼ ἐξαποστέλλω ὑμῖν τὸν σῖτον καὶ τὸν οἶνον καὶ τὸ ἔλαιον, καὶ ἐμπλησθήσεσθε αὐτῶν, καὶ οὐ δώσω ὑμᾶς οὐκ ἔτι εἰς ὀνειδισμὸν ἐν τοῖς ἔθνεσι.
\vs{20}Καὶ τὸν ἀπὸ Βοῤῥᾶ ἐκδιώξω ἀφʼ ὑμῶν, καὶ ἐξώσω αὐτὸν εἰς γῆν ἄνυδρον, καὶ ἀφανιῶ τὸ πρόσωπον αὐτοῦ εἰς τὴν θάλασσαν τὴν πρώτην, καὶ τὰ ὀπίσω αὐτοῦ εἰς τὴν θάλασσαν τὴν ἐσχάτην· καὶ ἀναβήσεται ἡ σαπρία αὐτοῦ, καὶ ἀναβήσεται ὁ βρόμος αὐτοῦ, ὅτι ἐμεγάλυνε τὰ ἔργα αὐτοῦ.

\vs{21}Θάρσει γῆ, χαῖρε καὶ εὐφραίνου, ὅτι ἐμεγάλυνε Κύριος τοῦ ποιῆσαι.
\vs{22}Θαρσεῖτε κτήνη τοῦ πεδίου, ὅτι βεβλάστηκε τὰ πεδία τῆς ἐρήμου, ὅτι ξύλον ἤνεγκεν τὸν καρπὸν αὐτοῦ, συκῆ καὶ ἄμπελος ἔδωκαν τὴν ἰσχὺν αὐτῶν.
\vs{23}Καὶ τὰ τέκνα Σιὼν χαίρετε καὶ εὐφραίνεσθε ἐπὶ τῷ Κυρίῳ Θεῷ ὑμῶν, διότι ἔδωκεν ὑμῖν τὰ βρώματα εἰς δικαιοσύνην, καὶ βρέξει ὑμῖν ὑετὸν πρώϊμον καὶ ὄψιμον, καθὼς ἔμπροσθεν,
\vs{24}καὶ πλησθήσονται αἱ ἅλωνες σίτου, καὶ ὑπερχυθήσονται αἱ ληνοὶ οἴνου καὶ ἐλαίου.
\vs{25}Καὶ ἀνταποδώσε ὑμῖν ἀντὶ τῶν ἐτῶν ὧν κατέφαγεν ἡ ἀκρὶς, καὶ ὁ βροῦχος, καὶ ἡ ἐρυσίβη, καὶ ἡ κάμπη, ἡ δύναμίς μου ἡ μεγάλη, ἣν ἐξαπέστειλα εἰς ὑμᾶς.
\vs{26}Καὶ φάγεσθε ἐσθίοντες, καὶ ἐμπλησθήσεσθε, καὶ αἰνέσετε τὸ ὄνομα Κυρίου τοῦ Θεοῦ ὑμῶν, ἃ ἐποίησε μεθʼ ὑμῶν εἰς θαυμάσια· καὶ οὐ μὴ καταισχυνθῇ ὁ λαός μου εἰς τὸν αἰῶνα.
\vs{27}Καὶ ἐπιγνώσεσθε ὅτι ἐν μέσῳ τοῦ Ἰσραὴλ ἐγώ εἰμι, καὶ ἐγὼ Κύριος ὁ Θεὸς ὑμῶν, καὶ οὐκ ἕστιν ἔτι πλὴν ἐμοῦ· καὶ οὐ μὴ καταισχυνθῶσιν ἔτι ὁ λαός μου εἰς τὸν αἰῶνα.

\ch{3}
Καὶ ἔσται μετὰ ταῦτα, καὶ ἐκχεῶ ἀπὸ τοῦ πνεύματός μου ἐπὶ πᾶσαν σάρκα, καὶ προφητεύσουσιν οἱ υἱοὶ ὑμῶν, καὶ αἱ θυγατέρες ὑμῶν, καὶ οἱ πρεσβύτεροι ὑμῶν ἐνύπνια ἐνυπνιασθήσονται, καὶ οἱ νεανίσκοι ὑμῶν ὁράσεις ὄψονται.
\vs{2}Καὶ ἐπὶ τοὺς δούλους μου καὶ ἐπὶ τὰς δούλας ἐν ταῖς ἡμέραις ἐκείναις ἐκχεῶ ἀπὸ τοῦ πνεύματός μου.
\vs{3}Καὶ δώσω τέρατα ἐν οὐρανῷ, καὶ ἐπὶ τῆς γῆς αἷμα καὶ πῦρ καὶ ἀτμίδα καπνοῦ.
\vs{4}Ὁ ἥλιος μεταστραφήσεται εἰς σκότος, καὶ ἡ σελήνη εἰς αἷμα, πρὶν ἐλθεῖν τὴν ἡμέραν Κυρίου τὴν μεγάλην, καὶ ἐπιφανῆ.

\vs{5}Καὶ ἔσται πᾶς ὃς ἂν ἐπικαλέσηται τὸ ὄνομα Κυρίου, σωθήσεται· ὅτι ἐν τῷ ὄρει Σιὼν καὶ ἐν Ἱερουσαλὴμ ἔσται ἀνασωζόμενος καθότι εἶπε Κύριος, καὶ εὐαγγελιζόμενοι οὓς Κύριος προσκέκληται.

\ch{4}
Ὅτι ἰδοὺ ἐγὼ ἐν ταῖς ἡμέραις ἐκείναις καὶ ἐν τῷ καιρῷ ἐκείνῳ, ὅταν ἐπιστρέψω τὴν αἰχμαλωσίαν Ἰούδα καὶ Ἱερουσαλὴμ,
\vs{2}καὶ συνάξω πάντα τὰ ἔθνη, καὶ κατάξω αὐτὰ εἰς τὴν κοιλάδα Ἰωσαφὰτ, καὶ διακριθήσομαι πρὸς αὐτοὺς ἐκεῖ ὑπὲρ τοῦ λαοῦ μου καὶ τῆς κληρονομίας μου Ἰσραὴλ, οἳ διεσπάρησαν ἐν τοῖς ἔθνεσι, καὶ τὴν γῆν μου κατεδιείλαντο,
\vs{3}καὶ ἐπὶ τὸν λαόν μου ἔβαλον κλήρους, καὶ ἔδωκαν τὰ παιδάρια πόρναις, καὶ τὰ κοράσια ἐπώλουν ἀντὶ τοῦ οἴνου, καὶ ἔπινον.

\vs{4}Καὶ τί ὑμεῖς ἐμοὶ Τύρος, καὶ Σιδὼν, καὶ πᾶσα Γαλιλαία ἀλλοφύλων; μὴ ἀνταπόδομα ὑμεῖς ἀνταποδίδοτέ μοι; ἢ μνησικακεῖτε ὑμεῖς ἐπʼ ἐμοί; ὀξέως, καὶ ταχέως ἀνταποδώσω τὸ ἀνταπόδομα ὑμῶν εἰς κεφαλὰς ὑμῶν,
\vs{5}ἀνθʼ ὧν τὸ ἀργύριόν μου καὶ τὸ χρυσίον μου ἐλάβετε, καὶ τὰ ἐπίλεκτά μου τὰ καλὰ εἰσηνέγκατε εἰς τοὺς ναοὺς ὑμῶν,
\vs{6}καὶ τοὺς υἱοὺς Ἰούδα καὶ τοὺς υἱοὺς Ἱερουσαλὴμ ἀπέδοσθε τοῖς υἱοῖς τῶν Ἑλλήνων, ὅπως ἐξώσητε αὐτοὺς ἐκ τῶν ὁρίων αὐτῶν.
\vs{7}Καὶ ἰδοὺ ἐγὼ ἐξεγείρω αὐτοὺς ἐκ τοῦ τόπου οὗ ἀπέδοσθε αὐτοὺς ἐκεῖ, καὶ ἀνταποδώσω τὸ ἀνταπόδομα ὑμῶν εἰς κεφαλὰς ὑμῶν,
\vs{8}καὶ ἀποδώσομαι τοὺς υἱοὺς ὑμῶν καὶ τὰς θυγατέρας ὑμῶν εἰς χεῖρας τῶν υἱῶν Ἰούδα, καὶ ἀποδώσονται αὐτοὺς εἰς αἰχμαλωσίαν εἰς ἔθνος μακρὰν ἀπέχον, ὅτι Κύριος ἐλάλησε.

\vs{9}Κηρύξατε ταῦτα ἐν τοῖς ἔθνεσιν, ἁγιάσατε πόλεμον, ἐξεγείρατε τοὺς μαχητὰς, προσαγάγετε καὶ ἀναβαίνετε πάντες ἄνδρες πολεμισταὶ,
\vs{10}συγκόψατε τὰ ἄροτρα ὑμῶν εἰς ῥομφαίας, καὶ τὰ δρέπανα ὑμῶν εἰς σειρομάστας· ὁ ἀδύνατος λεγέτω, ὅτι ἰσχύω ἐγώ.
\vs{11}Συναθροίζεσθε, καὶ εἰσπορεύεσθε πάντα τὰ ἔθνη κυκλόθεν, καὶ συνάχθητε ἐκεῖ· ὁ πρᾳῢς ἔστω μαχητής.
\vs{12}Ἐξεγειρέσθωσαν, ἀναβαινέτωσαν πάντα τὰ ἔθνη εἰς τὴν κοιλάδα Ἰωσαφὰτ, διότι ἐκεῖ καθιῶ τοῦ διακρῖναι πάντα τὰ ἔθνη κυκλόθεν.

\vs{13}Ἐξαποστείλατε δρέπανα, ὅτι παρέστηκεν ὁ τρυγητός· εἰσπορεύεσθε, πατεῖτε, διότι πλήρης ἡ ληνός· ὑπερεκχεῖτε τὰ ὑπολήνια, ὅτι πεπλήθυνται τὰ κακὰ αὐτῶν.
\vs{14}Ἦχοι ἐξήχησαν ἐν τῇ κοιλάδι τῆς δίκης, ὅτι ἐγγὺς ἡμέρα Κυρίου ἐν τῇ κοιλάδι τῆς δίκης.
\vs{15}Ὁ ἥλιος καὶ ἡ σελήνη συσκοτάσουσι, καὶ οἱ ἀστέρες δύσουσι φέγγος αὐτῶν.

\vs{16}Ὁ δὲ Κύριος ἐκ Σιὼν ἀνακεκράξεται, καὶ ἐξ Ἱερουσαλὴμ δώσει φωνὴν αὐτοῦ, καὶ σεισθήσεται ὁ οὐρανὸς καὶ ἡ γῆ· ὁ δὲ Κύριος φείσεται τοῦ λαοῦ αὐτοῦ, καὶ ἐνισχύσει τοὺς υἱοὺς Ἰσραήλ.
\vs{17}Καὶ ἐπιγνώσεσθε διότι ἐγὼ Κύριος ὁ Θεὸς ὑμῶν, ὁ κατασκηνῶν ἐν Σιὼν ὄρει ἁγίῳ μου· καὶ ἔσται Ἱερουσαλὴμ ἁγία, καὶ ἀλλογενεῖς οὐ διελεύσονται διʼ αὐτῆς οὐκέτι.

\vs{18}Καὶ ἔσται ἐν τῇ ἡμέρᾳ ἐκείνῃ, ἀποσταλάξει τὰ ὄρη γλυκασμὸν, καὶ οἱ βουνοὶ ῥυήσονται γάλα, καὶ πᾶσαι αἱ ἀφέσεις Ἰούδα ῥυήσονται ὕδατα, καὶ πηγὴ ἐξ οἴκου Κυρίου ἐξελεύσεται, καὶ ποτιεῖ τὸν χειμάῤῥουν τῶν σχοίνων.

\vs{19}Αἴγυπτος εἰς ἀφανισμὸν ἔσται, καὶ ἡ Ἰδουμαία εἰς πεδίον ἀφανισμοῦ ἔσται, ἐξ ἀδικιῶν υἱῶν Ἰούδα, ἀνθʼ ὧν ἐξέχεαν αἷμα δίκαιον ἐν τῇ γῇ αὐτῶν.
\vs{20}Ἡ δὲ Ἰουδαία εἰς τὸν αἰῶνα κατοικηθήσεται, καὶ Ἱερουσαλὴμ εἰς γενεὰς γενεῶν.
\vs{21}Καὶ ἐκζητήσω τὸ αἷμα αὐτῶν, καὶ οὐ μὴ ἀθωώσω, καὶ Κύριος κατασκηνώσει ἐν Σιών.


\def\book{ΑΜΩΣ}
\biblebook{ΑΜΩΣ}


\lettrine[lines=2, loversize=0.2, nindent=0em, findent=.25em]{\textcolor{bookheadingcolor}{Λ}}{ΟΓΟΙ} Αμὼς οἳ ἐγένοντο ἐν Ἀκκαρεὶμ ἐν Θεκουὲ, οὓς εἶδεν ὑπὲρ Ἱερουσαλὴμ, ἐν ἡμέραις Ὀζίου βασιλέως Ἰούδα, καὶ ἐν ἡμέραις Ἱεροβοὰμ τοῦ Ἰωὰς βασιλέως Ἰσραὴλ, πρὸ δύο ἐτῶν τοῦ σεισμοῦ.

\vs{2}Καὶ εἶπε, Κύριος ἐκ Σιὼν ἐφθέγξατο, καὶ ἐξ Ἱερουσαλὴμ ἔδωκε φωνὴν αὐτοῦ· καὶ ἐπένθησαν αἱ νομαὶ τῶν ποιμένων, καὶ ἐξηράνθη ἡ κορυφὴ τοῦ Καρμήλου.

\vs{3}Καὶ εἶπε Κύριος, ἐπὶ ταῖς τρισὶν ἀσεβείαις Δαμασκοῦ, καὶ ἐπὶ ταῖς τέσσαρσιν οὐκ ἀποστραφήσομαι αὐτὸν, ἀνθʼ ὧν ἔπριζον πριοσι σιδηροῖς τὰς ἐν γαστρὶ ἐχούσας τῶν ἐν Γαλαάδ.
\vs{4}Καὶ ἀποστελῶ πῦρ εἰς τὸν οἶκον Ἀζαὴλ, καὶ καταφάγεται τὰ θεμέλια υἱοῦ Ἄδερ.
\vs{5}Καὶ συντρίψω μοχλοὺς Δαμασκοῦ, καὶ ἐξολοθρεύσω κατοικοῦντας ἐκ πεδίου Ὦν, καὶ κατακόψω φυλὴν ἐξ ἀνδρῶν Χαῤῥὰν, καὶ αἰχμαλωτευθήσεται λαὸς Συρίας ἐπίκλητος, λέγει Κύριος.

\vs{6}Τάδε λέγει Κύριος, ἐπὶ ταῖς τρισὶν ἀσεβείαις Γάζης, καὶ ἐπὶ ταῖς τέσσαρσιν οὐκ ἀποστραφήσομαι αὐτοὺς, ἕνεκεν τοῦ αἰχμαλωτεῦσαι αὐτοὺς αἰχμαλωσίαν τοῦ Σαλωμὼν, τοῦ συγκλεῖσαι εἰς τὴν Ἰδουμαίαν.
\vs{7}Καὶ ἐξαποστελῶ πῦρ ἐπὶ τὰ τείχη Γάζης, καὶ καταφάγεται τὰ θεμέλια αὐτῆς.
\vs{8}Καὶ ἐξολοθρεύσω κατοικοῦντας ἐξ Ἀζώτου, καὶ ἐξαρθήσεται φυλὴ ἐξ Ἀσκάλωνος, καὶ ἐπάξω τὴν χεῖρά μου ἐπὶ Ἀκκάρων, καὶ ἀπολοῦνται οἱ κατάλοιποι τῶν ἀλλοφύλων, λέγει Κύριος.

\vs{9}Τάδε λέγει Κύριος, ἐπὶ ταῖς τρισὶν ἀσεβείαις Τύρου, καὶ ἐπὶ ταῖς τέσσαρσιν οὐκ ἀποστραφήσομαι αὐτὴν, ἀνθʼ ὧν συνέκλεισαν αἰχμαλωσίαν τοῦ Σαλωμὼν εἰς τὴν Ἰδουμαίαν, καὶ οὐκ ἐμνήσθησαν διαθήκης ἀδελφῶν.
\vs{10}Καὶ ἐξαποστελῶ πῦρ ἐπὶ τὰ τείχη Τύρου, καὶ καταφάγεται τὰ θεμέλια αὐτῆς.

\vs{11}Τάδε λέγει Κύριος, ἐπὶ ταῖς τρισὶν ἀσεβείαις τῆς Ἰδουμαίας, καὶ ἐπὶ ταῖς τέσσαρσιν οὐκ ἀποστραφήσομαι αὐτοὺς, ἕνεκα τοῦ διῶξαι αὐτοὺς ἐν ῥομφαίᾳ τὸν ἀδελφὸν αὐτοῦ, καὶ ἐλυμήνατο μητέρα ἐπὶ γῆς, καὶ ἥρπασεν εἰς μαρτύριον φρίκην αὐτοῦ, καὶ τὸ ὅρμημα αὐτοῦ ἐφύλαξεν εἰς νῖκος.
\vs{12}Καὶ ἐξαποστελῶ πῦρ εἰς Θαμὰν, καὶ καταφάγεται θεμέλια τειχέων αὐτῆς.

\vs{13}Τάδε λέγει Κύριος, ἐπὶ ταῖς τρισὶν ἀσεβείαις υἱῶν Ἀμμὼν, καὶ ἐπὶ ταῖς τέσσαρσιν οὐκ ἀποστραφήσομαι αὐτὸν, ἀνθʼ ὧν ἀνέσχιζον τὰς ἐν γαστρὶ ἐχούσας τῶν Γαλααδιτῶν, ὅπως ἐνπλατύνωσι τὰ ὅρια ἑαυτῶν.
\vs{14}Καὶ ἀνάψω πῦρ ἐπὶ τεὶχη Ῥαββὰθ, καὶ καταφάγεται θεμέλια αὐτῆς μετὰ κραυγῆς ἐν ἡμέρᾳ πολέμου, καὶ σεισθήσεται ἐν ἡμέραις συντελείας αὐτῆς,
\vs{15}καὶ πορεύσονται οἱ βασιλεῖς αὐτῆς ἐν αἰχμαλωσίᾳ, οἱ ἱερεῖς αὐτῶν καὶ οἱ ἄρχοντες αὐτῶν ἐπιτοαυτὸ, λέγει Κύριος.

\ch{2}
Τάδε λέγει Κύριος, ἐπὶ ταῖς τρισὶν ἀσεβείαις Μωὰβ, καὶ ἐπὶ ταῖς τέσσαρσιν οὐκ ἀποστραφήσομαι αὐτὸν, ἀνθʼ ὧν κατέκαυσαν τὰ ὀστᾶ βασιλέως τῆς Ἰδουμαίας εἰς κονίαν.
\vs{2}Καὶ ἐξαποστελῶ πῦρ εἰς Μωὰβ, καὶ καταφάγεται τὰ θεμέλια τῶν πόλεων αὐτῆς, καὶ ἀποθανεῖται ἐν ἀδυναμίᾳ Μωὰβ μετὰ κραυγῆς καὶ μετὰ φωνῆς σάλπιγγος.
\vs{3}Καὶ ἐξολοθρεύσω κριτὴν ἐξ αὐτῆς, καὶ πάντας [ἄρχοντας] αὐτῆς ἀποκτενῶ μετʼ αὐτοῦ, λέγει Κύριος.

\vs{4}Τάδε λέγει Κύριος, ἐπὶ ταῖς τρισὶν ἀσεβείαις υἱῶν Ἰούδα, καὶ ἐπὶ ταῖς τέσσαρσιν οὐκ ἀποστραφήσουμαι αὐτὸν, ἕνεκα τοῦ ἀπώσασθαι αὐτοὺς τὸν νόμον τοῦ Κυρίου, καὶ τὰ προστάγματα αὐτοῦ οὐκ ἐφυλάξαντο, καὶ ἐπλάνησεν αὐτοὺς τὰ μάταια αὐτῶν ἃ ἐποίησαν, οἷς ἐξηκολούθησαν οἱ πατέρες αὐτῶν ὀπίσω αὐτῶν.
\vs{5}Καὶ ἐξαποστελῶ πῦρ ἐπὶ Ἰούδαν, καὶ καταφάγεται θεμέλια Ἱερουσαλήμ.

\vs{6}Τάδε λέγει Κύριος, ἐπὶ ταῖς τρισὶν ἀσεβείαις Ἰσραὴλ, καὶ ἐπὶ ταῖς τέσσαρσιν οὐκ ἀποστραφήσομαι αὐτὸν, ἀνθʼ ὧν ἀπέδοντο ἀργυρίου δίκαιον, καὶ πένητα ἕνεκεν ὑποδημάτων,
\vs{7}τὰ πατοῦντα ἐπὶ τὸν χοῦν τῆς γῆς, καὶ ἐκονδύλιζον εἰς κεφαλὰς πτωχῶν, καὶ ὁδὸν ταπεινῶν ἐξέκλιναν, καὶ υἱὸς καὶ πατὴρ αὐτοῦ εἰσεπορεύοντο πρὸς τὴν αὐτὴν παιδίσκην, ὅπως βεβηλῶσι τὸ ὄνομα τοῦ Θεοῦ αὐτῶν.
\vs{8}Καὶ τὰ ἱμάτια αὐτῶν δεσμεύοντες σχοινίοις, παραπετάσματα ἐποίουν ἐχόμενα τοῦ θυσιαστηρίου, καὶ οἶνον ἐκ συκοφαντιῶν ἔπινον ἐν τῷ οἴκῳ τοῦ Θεοῦ αὐτῶν.

\vs{9}Ἐγὼ δὲ ἐξῇρα τὸν Ἀμοῤῥαῖον ἐκ προσώπου αὐτῶν, οὗ ἦν, καθὼς ὕψος κέδρου, τὸ ὕψος αὐτοῦ, καὶ ἰσχυρὸς ἦν ὡς δρῦς, καὶ ἐξήρανα τὸν καρπὸν αὐτοῦ ἐπάνωθεν, καὶ τὰς ῥίζας αὐτοῦ ὑποκάτωθεν.
\vs{10}Καὶ ἐγὼ ἀνήγαγον ὑμᾶς ἐκ γῆς Αἰγύπτου, καὶ περιήγαγον ὑμᾶς ἐν τῇ ἐρήμῳ τεσσαράκοντα ἔτη, τοῦ κατακληρονομῆσαι τὴν γῆν τῶν Ἀμοῤῥαίων.
\vs{11}Καὶ ἔλαβον ἐκ τῶν υἱῶν ὑμῶν εἰς προφήτας, καὶ ἐκ τῶν νεανίσκων ὑμῶν εἰς ἁγιασμόν· μὴ οὐκ ἔστι ταῦτα υἱοὶ Ἰσραήλ; λέγει Κύριος.
\vs{12}Καὶ ἐποτίζετε τοὺς ἡγιασμένους οἶνον, καὶ τοῖς προφήταις ἐνετέλλεσθε λέγοντες, οὐ μὴ προφητεύσητε.

\vs{13}Διατοῦτο ἰδοὺ ἐγὼ κυλίω ὑποκάτω ὑμῶν, ὃν τρόπον κυλίεται ἡ ἅμαξα ἡ γέμουσα καλάμης.
\vs{14}Καὶ ἀπολεῖται φυγὴ ἐκ δρομέως, καὶ ὁ κραταιὸς οὐ μὴ κρατήσῃ τῆς ἰσχύος αὐτοῦ, καὶ ὁ μαχητὴς οὐ μὴ σώσῃ τὴν ψυχὴν αὐτοῦ.
\vs{15}Καὶ ὁ τοξότης οὐ μὴ ὑποστῇ καὶ ὁ ὀξὺς τοῖς ποσὶν αὐτοῦ οὐ μὴ διασωθῇ, καὶ ὁ ἱππεὺς οὐ μὴ σώσῃ τὴν ψυχὴν αὐτοῦ,
\vs{16}καὶ ὁ κραταιὸς οὐ μὴ εὑρήσει τὴν καρδίαν αὐτοῦ ἐν δυναστείαις, ὁ γυμνὸς διώξεται ἐν ἐκείνῃ τῇ ἡμέρᾳ, λέγει Κύριος.

\ch{3}
Ἀκούσατε τὸν λόγον τοῦτον, ὃν ἐλάλησε Κύριος ἐφʼ ὑμᾶς, οἶκος Ἰσραὴλ, καὶ κατὰ πάσης φυλῆς, ἧς ἀνήγαγον ἐκ γῆς Αἰγύπτου, λέγων,
\vs{2}πλὴν ὑμᾶς ἔγνων ἐκ πασῶν τῶν φυλῶν τῆς γῆς, διατοῦτο ἐκδικήσω ἐφʼ ὑμᾶς πάσας τὰς ἁμαρτίας ὑμῶν.

\vs{3}Εἰ πορεύσονται δύο ἐπιτοαυτὸ καθόλου, ἐὰν μὴ γνωρίσωσιν ἑαυτούς;
\vs{4}Εἰ ἐρεύξεται λέων ἐκ τοῦ δρυμοῦ αὐτοῦ θήραν οὐκ ἔχων; εἰ δώσει σκύμνος φωνὴν αὐτοῦ ἐκ τῆς μάνδρας αὐτοῦ καθόλου, ἐὰν μὴ ἁρπάσῃ τί;
\vs{5}Εἰ πεσεῖται ὄρνεον ἐπὶ τῆς γῆς ἄνευ ἰξευτοῦ; εἰ σχασθήσεται παγὶς ἐπὶ τῆς γῆς ἄνευ τοῦ συλλαβεῖν τί;
\vs{6}Εἰ φωνήσει σάλπιγξ ἐν πόλει, καὶ λαὸς οὐ πτοηθήσεται; εἰ ἔσται κακία ἐν πόλει ἣν Κύριος οὐκ ἐποίησε;
\vs{7}Διότι οὐ μὴ ποιήσῃ Κύριος ὁ Θεὸς πρᾶγμα ἐὰν μὴ ἀποκαλύψῃ παιδείαν πρὸς τοὺς δούλους αὐτοῦ τοὺς προφήτας.
\vs{8}Λέων ἐρεύξεται, καὶ τίς οὐ φοβηθήσεται; Κύριος ὁ Θεὸς ἐλάλησε, καὶ τίς οὐ προφητεύσει;

\vs{9}Ἀναγγείλατε χώραις ἐν Ἀσσυρίοις, καὶ ἐπὶ τὰς χώρας τῆς Αἰγύπτου, καὶ εἴπατε, συνάχθητε ἐπὶ τὸ ὄρος Σαμαρείας, καὶ ἴδετε θαυμαστὰ πολλὰ ἐν μέσῳ αὐτῆς, καὶ καταδυναστείαν τὴν ἐν αὐτῇ.
\vs{10}Καὶ οὐκ ἔγνω ἃ ἔσται ἐναντίον αὐτῆς, λέγει Κύριος, οἱ θησαυρίζοντες ἀδικίαν καὶ ταλαιπωρίαν ἐν ταῖς χώραις αὐτῶν.
\vs{11}Διατοῦτο τάδε λέγει Κύριος ὁ Θεὸς, Τύρος κυκλόθεν ἡ γῆ σου ἐρημωθήσεται, καὶ κατάξει ἐκ σοῦ ἰσχύν σου, καὶ διαρπαγήσονται αἱ χῶραί σου.
\vs{12}Τάδε λέγει Κύριος, ὃν τρόπον ὅταν ἐκσπάσῃ ὁ ποιμὴν ἐκ στόματος τοῦ λέοντος δύο σκέλη ἢ λοβὸν ὠτίου, οὕτως ἐκσπασθήσονται οἱ υἱοὶ Ἰσραὴλ οἱ κατοικοῦντες ἐν Σαμαρείᾳ κατέναντι τῆς φυλῆς, καὶ ἐν Δαμασκῷ.

\vs{13}Ἰερεῖς ἀκούσατε, καὶ ἐπιμαρτύρασθε τῷ οἴκῳ Ἰακὼβ, λέγει Κύριος ὁ Θεὸς ὁ παντοκράτωρ.
\vs{14}Διότι ἐν τῇ ἡμέρᾳ ὅταν ἐκδικῶ ἀσεβείας τοῦ Ἰσραὴλ ἐπʼ αὐτὸν, καὶ ἐκδικήσω ἐπὶ τὰ θυσιαστήρια Βαιθήλ· καὶ κατασκαφήσεται τὰ κέρατα τοῦ θυσιαστηρίου, καὶ πεσοῦνται ἐπὶ τὴν γῆν·
\vs{15}Συγχεῶ καὶ πατάξω τὸν οἶκον τὸν περίπτερον ἐπὶ τὸν οἶκον τὸν θερινὸν, καὶ ἀπολοῦνται οἶκοι ἐλεφάντινοι, καὶ προστεθήσονται ἕτεροι οἶκοι πολλοὶ, λέγει Κύριος.

\ch{4}
Ἀκούσατε τὸν λόγον τοῦτον δαμάλεις τῆς Βασανίτιδος, αἱ ἐν τῷ ὄρει τῆς Σαμαρείας, αἱ καταδυναστεύουσαι πτωχοὺς, καὶ καταπατοῦσαι πένητας, αἱ λέγουσαι τοῖς κυρίοις αὐτῶν, ἐπίδοτε ἡμῖν ὅπως πίωμεν.

\vs{2}Ομνύει Κύριος κατὰ τῶν ἁγίων αὐτοῦ, διότι ἰδοὺ ἡμέραι ἔρχονται ἐφʼ ὑμᾶς, καὶ λήψονται ὑμᾶς ἐν ὅπλοις, καὶ τοὺς μεθʼ ὑμῶν εἰς λέβητας ὑποκαιομένους ἐμβαλοῦσιν ἔμπυροι λοιμοὶ,
\vs{3}καὶ ἐξενεχθήσεσθε γυμναὶ κατέναντι ἀλλήλων, καὶ ἀποῤῥιφήσεσθε εἰς τὸ ὄρος τὸ Ῥομμὰν, λέγει Κύριος.

\vs{4}Εἰσήλθατε εἰς Βαιθὴλ, καὶ ἠσεβήσατε, καὶ εἰς Γάλγαλα ἐπληθύνατε τοῦ ἀσεβῆσαι· καὶ ἠνέγκατε εἰς τοπρωῒ θυσίας ὑμῶν, εἰς τὴν τριημερίαν τὰ ἐπιδέκατα ὑμῶν.
\vs{5}Καὶ ἀνέγνωσαν ἔξω νόμον, καὶ ἐπεκαλέσαντο ὁμολογίας· ἀναγγείλατε ὅτι ταῦτα ἠγάπησαν οἱ υἱοὶ Ἰσραὴλ, λέγει Κύριος.

\vs{6}Καὶ ἐγὼ δώσω ὑμῖν γομφιασμὸν ὀδόντων ἐν πάσαις ταῖς πόλεσιν ὑμῶν, καὶ ἔνδειαν ἄρτων ἐν πᾶσι τοῖς τόποις ὑμῶν· καὶ οὐκ ἐπεστρέψατε πρὸς μὲ, λέγει Κύριος.
\vs{7}Καὶ ἐγὼ ἀνέσχον ἐξ ὑμῶν τὸν ὑετὸν πρὸ τριῶν μηνῶν τοῦ τρυγητοῦ, καὶ βρέξω ἐπὶ πόλιν μίαν, ἐπὶ δὲ πόλιν μίαν οὐ βρέξω· μερὶς μία βραχήσεται, καὶ μερὶς, ἐφʼ ἣν οὐ βρέξω, ξηρανθήσεται.
\vs{8}Καὶ συναθροισθήσονται δύο καὶ τρεῖς πόλεις εἰς πόλιν μίαν τοῦ πιεῖν ὕδωρ, καὶ οὐ μὴ ἐμπλησθῶσι· καὶ οὐκ ἐπεστράφητε πρὸς μὲ, λέγει Κύριος.
\vs{9}Ἐπάταξα ὑμᾶς ἐν πυρώσει, καὶ ἐν ἰκτέρῳ· ἐπληθύνατε κήπους ὑμῶν, ἀμπελῶνας ὑμῶν, καὶ συκεῶνας ὑμῶν· καὶ ἐλαιῶνας ὑμῶν κατέφαγεν ἡ κάμπη· καὶ οὐδʼ ὡς ἐπεστρέψατε πρὸς μὲ, λέγει Κύριος.
\vs{10}Ἐξαπέστειλα εἰς ὑμᾶς θάνατον ἐν ὁδῷ Αἰγύπτου, καὶ ἀπέκτεινα ἐν ῥομφαίᾳ τοὺς νεανίσκους ὑμῶν, μετὰ αἰχμαλωσίας ἵππων σου, καὶ ἀνήγαγον ἐν πυρὶ τὰς παρεμβολὰς ἐν τῇ ὀργῇ ὑμῶν· καὶ οὐδʼ ὡς ἐπεστρέψατε πρὸς μὲ, λέγει Κύριος.
\vs{11}Κατέστρεψα ὑμᾶς, καθὼς κατέστρεψεν ὁ Θεὸς Σόδομα καὶ Γόμοῤῥα, καὶ ἐγένεσθε ὡς δαλὸς ἐξέσπασμένος ἐκ πυρός· καὶ οὐδʼ ὡς ἐπεστρέψατε πρὸς μὲ, λέγει Κύριος.

\vs{12}Διατοῦτο οὕτως ποιήσω σοι Ἰσραήλ· πλὴν ὅτι οὕτως ποιήσω σοι, ἑτοιμάζου τοῦ ἐπικαλεῖσθαι τὸν Θεόν σου Ἰσραήλ.
\vs{13}Διότι ἰδοὺ ἐγὼ στερεῶν βροντὴν, καὶ κτίζων πνεῦμα, καὶ ἀπαγγέλλων εἰς ἀνθρώπους τὸν χριστὸν αὐτοῦ, ποιῶν ὄρθρον καὶ ὁμίχλην, καὶ ἐπιβαίνων ἐπὶ τὰ ὑψηλὰ τῆς γῆς· Κύριος ὁ Θεὸς ὁ παντοκράτωρ ὄνομα αὐτῷ.

\ch{5}
Ἀκούσατε τὸν λόγον Κυρίου τοῦτον, ὃν ἐγὼ λαμβάνω ἐφʼ ὑμᾶς, θρῆνον. Οἶκος Ἰσραὴλ
\vs{2}ἔπεσεν, οὐκέτι μὴ προσθήσει τοῦ ἀναστῆναι. Παρθένος τοῦ Ἰσραὴλ ἔσφαλεν ἐπὶ τῆς γῆς αὐτοῦ, οὐκ ἔστιν ὁ ἀναστήσων αὐτήν.
\vs{3}Διατοῦτο τάδε λέγει Κύριος Κύριος, ἡ πόλις ἐξ ἧς ἐξεπορεύοντο χίλιοι, ὑπολειφθήσονται ἑκατόν· καὶ ἐξ ἧς ἐξεπορεύοντο ἑκατὸν, ὑπολειφθήσονται δέκα τῷ οἴκῳ Ἰσραήλ.

\vs{4}Διότι τάδε λέγει Κύριος πρὸς τὸν οἶκον Ἰσραὴλ, ἐκζητήσατέ με, καὶ ζήσεσθε.
\vs{5}Καὶ μὴ ἐκζητεῖτε Βαιθὴλ, καὶ εἰς Γάλγαλα μὴ εἰσπορεύεσθε, καὶ ἐπὶ τὸ φρέαρ τοῦ ὅρκου μὴ διαβαίνετε, ὅτι Γάλγαλα αἰχμαλωτευομένη αἰχμαλωτευθήσεται, καὶ Βαιθὴλ ἔσται ὡς οὐχ ὑπάρχουσα.
\vs{6}Ἐκζητήσατε τὸν Κύριον, καὶ ζήσατε, ὅπως μὴ ἀναλάμψῃ ὡς πῦρ ὁ οἶκος Ἰωσὴφ, καὶ καταφάγῃ αὐτὸν, καὶ οὐκ ἔσται ὁ σβέσων τῷ οἴκῳ Ἰσραήλ.

\vs{7}Ὁ ποιῶν εἰς ὕψος κρίμα, καὶ δικαιοσύνην εἰς γῆν ἔθηκεν·
\vs{8}ὁ ποιῶν πάντα καὶ μετασκευάζων, καὶ ἐκτρέπων εἰς τοπρωῒ σκιάν, καὶ ἡμέραν εἰς νύκτα συσκοτάζων, ὁ προσκαλούμενος τὸ ὕδωρ τῆς θαλάσσης, καὶ ἐκχέων αὐτὸ ἐπὶ πρόσωπον τῆς γῆς· Κύριος ὄνομα αὐτῷ·
\vs{9}ὁ διαιρῶν συντριμμὸν ἐπὶ ἰσχὺν, καὶ ταλαιπωρίαν ἐπὶ ὀχύρωμα ἐπάγων.

\vs{10}Ἐμίσησαν ἐν πύλαις ἐλέγχοντα, καὶ λόγον ὅσιον ἐβδελύξαντο.
\vs{11}Διατοῦτο ἀνθʼ ὧν κατεκονδύλιζον πτωχοὺς, καὶ δῶρα ἐκλεκτὰ ἐδέξασθε παρʼ αὐτῶν, οἴκους ξεστοὺς ᾠκοδομήσατε, καὶ οὐ μὴ κατοικήσητε ἐν αὐτοῖς· ἀμπελῶνας ἐπιθυμητοὺς ἐφυτεύσατε, καὶ οὐ μὴ πίητε τὸν οἶνον αὐτῶν.
\vs{12}Ὅτι ἔγνων πολλὰς ἀσεβείας ὑμῶν, καὶ ἰσχυραὶ αἱ ἁμαρτίαι ὑμῶν, καταπατοῦντες δίκαιον, λαμβάνοντες ἀλλάγματα, καὶ πένητας ἐν πύλαις ἐκκλίνοντες.

\vs{13}Διατοῦτο ὁ συνιὼν ἐν τῷ καιρῷ ἐκείνῳ σιωπήσεται, ὅτι καιρὸς πονηρῶν ἐστιν.
\vs{14}Ἐκζητήσατε τὸ καλὸν, καὶ μὴ πονηρὸν, ὅπως ζήσητε, καὶ ἔσται οὕτως μεθʼ ὑμῶν Κύριος ὁ Θεὸς ὁ παντοκράτωρ, ὃν τρόπον εἴπατε,
\vs{15}μεμισήκαμεν τὰ πονηρὰ, καὶ ἠγαπήσαμεν τὰ καλά· καὶ ἀποκαταστήσατε ἐν πύλαις κρίμα, ὅπως ἐλεήσῃ Κύριος ὁ Θεὸς ὁ παντοκράτωρ τοὺς περιλοίπους τοῦ Ἰωσήφ.

\vs{16}Διατοῦτο τάδε λέγει Κύριος ὁ Θεὸς ὁ παντοκράτωρ, ἐν πάσαις ταῖς πλατείαις κοπετὸς, καὶ ἐν πάσαις ταῖς ὁδοῖς ῥηθήσεται οὐαὶ, οὐαί· κληθήσεται γεωργὸς εἰς πένθος καὶ κοπετὸν, καὶ εἰς εἰδότας θρῆνον.
\vs{17}Καὶ ἐν πάσαις ὁδοῖς κοπετὸς, διότι ἐλεύσομαι διὰ μέσου σου, εἶπε Κύριος.

\vs{18}Οὐαὶ οἱ ἐπιθυμοῦντες τὴν ἡμέραν Κυρίου· ἱνατί αὕτη ὑμῖν ἡ ἡμέρα τοῦ Κυρίου; καὶ αὕτη ἐστὶ σκότος καὶ οὐ φῶς.
\vs{19}Ὃν τρόπον ἐὰν φύγῃ ἄνθρωπος ἐκ προσώπου τοῦ λέοντος, καὶ ἐμπέσῃ αὐτῷ ἡ ἄρκος, καὶ εἰσπηδήσῃ εἰς τὸν οἶκον αὐτοῦ, καὶ ἀπερείσηται τὰς χεῖρας αὐτοῦ ἐπὶ τὸν τοῖχον, καὶ δάκῃ αὐτὸν ὄφις.
\vs{20}Οὐχὶ σκότος ἡ ἡμέρα τοῦ Κυρίου, καὶ οὐ φῶς, καὶ γνόφος οὐκ ἔχων φέγγος αὕτη;

\vs{21}Μεμίσηκα, ἀπῶσμαι ἑορτὰς ὑμῶν, καὶ οὐ μὴ ὀσφρανθῶ θυσίας ἐν ταῖς πανηγύρεσιν ὑμῶν.
\vs{22}Διότι ἐὰν ἐνέγκητέ μοι ὁλοκαυτώματα καὶ θυσίας ὑμῶν, οὐ προσδέξομαι, καὶ σωτηρίους ἐπιφανείας ὑμῶν οὐκ ἐπιβλέψομαι.
\vs{23}Μετάστησον ἀπʼ ἐμοῦ ἦχον ᾠδῶν σου, καὶ ψαλμὸν ὀργάνων σου οὐκ ἀκούσομαι.
\vs{24}Καὶ κυλισθήσεται ὡς ὕδωρ κρίμα, καὶ δικαιοσύνη ὡς χειμάῤῥους ἄβατος.
\vs{25}Μὴ σφάγια καὶ θυσίας προσηνέγκατέ μοι οἶκος Ἰσραὴλ τεσσαράκοντα ἔτη ἐν τῇ ἐρήμῳ;
\vs{26}Καὶ ἀνελάβετε τὴν σκηνὴν τοῦ Μολὸχ, καὶ τὸ ἄστρον τοῦ θεοῦ ὑμῶν Ῥαιφὰν, τοὺς τύπους αὐτῶν οὓς ἐποιήσατε ἑαυτοῖς.
\vs{27}Καὶ μετοικιῶ ὑμᾶς ἐπέκεινα Δαμασκοῦ, λέγει Κύριος· ὁ Θεὸς ὁ παντοκράτωρ ὄνομα αὐτῷ.

\ch{6}
Οὐαὶ τοῖς ἐξουθενοῦσι Σιὼν, καὶ τοῖς πεποιθόσιν ἐπὶ τὸ ὄρος Σαμαρείας, ἀπετρύγησαν ἀρχὰς ἐθνῶν, καὶ εἰσῆλθον αὐτοί. οἶκος τοῦ Ἰσραὴλ
\vs{2}διάβητε πάντες καὶ ἴδετε, καὶ διέλθατε ἐκεῖθεν εἰς Ἐματραββὰ, καὶ κατάβητε ἐκεῖθεν εἰς Γὲθ ἀλλοφύλων, τὰς κρατίστας ἐκ πασῶν τῶν βασιλειῶν τούτων, εἰ πλείονα τὰ ὅρια αὐτῶν ἐστι τῶν ὑμετέρων ὁρίων.

\vs{3}Οἱ ἐρχόμενοι εἰς ἡμέραν κακὴν, οἱ ἐγγίζοντες καὶ ἐφαπτόμενοι σαββάτων ψευδῶν,
\vs{4}οἱ καθεύδοντες ἐπὶ κλινῶν ἐλεφαντίνων, καὶ κατασπαταλῶντες ἐπὶ ταῖς στρωμναῖς αὐτῶν, καὶ ἔσθοντες ἐρίφους ἐκ ποιμνίων, καὶ μοσχάρια ἐκ μέσου βουκολίων γαλαθηνὰ,
\vs{5}οἱ ἐπικροτοῦντες πρὸς τὴν φωνὴν τῶν ὀργάνων, ὡς ἑστηκότα ἐλογίσαντο, καὶ οὐχ ὡς φεύγοντα,
\vs{6}οἱ πίνοντες τὸν διυλισμένον οἶνον, καὶ τὰ πρῶτα μῦρα χριόμενοι, καὶ οὐκ ἔπασχον οὐδὲν ἐπὶ τῇ συντριβῇ Ἰωσήφ.
\vs{7}Διὰ τοῦτο νῦν αἰχμάλωτοι ἔσονται ἀπʼ ἀρχῆς δυναστῶν, καὶ ἐξαρθήσεται χρεμετισμὸς ἵππων ἐξ Ἐφραίμ·

\vs{8}Ὅτι ὤμοσε Κύριος καθʼ ἑαυτοῦ, διότι βδελύσσομαι ἐγὼ πᾶσαν τὴν ὕβριν Ἰακὼβ, καὶ τὰς χώρας αὐτοῦ μεμίσηκα, καὶ ἐξαρῶ πόλιν σὺν πᾶσι τοῖς κατοικοῦσιν αὐτήν.

\vs{9}Καὶ ἔσται, ἐὰν ὑπολειφθῶσι δέκα ἄνδρες ἐν οἰκίᾳ μιᾷ, καὶ ἀποθανοῦνται, καὶ ὑπολειφθήσονται οἱ κατάλοιποι,
\vs{10}καὶ λήψονται οἱ οἰκεῖοι αὐτῶν, καὶ παραβιῶνται τοῦ ἐξενέγκαι τὰ ὀστᾶ αὐτῶν ἐκ τοῦ οἴκου· καὶ ἐρεῖ τοῖς προεστηκόσι τῆς οἰκίας, εἰ ἔτι ὑπάρχει παρὰ σοί; Καὶ ἐρεῖ, οὐκ ἔτι· καὶ ἐρεῖ, σίγα ἕνεκα τοῦ μὴ ὀνομάσαι τὸ ὄνομα Κυρίου.

\vs{11}Διότι ἰδοὺ Κύριος ἐντέλλεται, καὶ πατάξει τὸν οἶκον τὸν μέγαν θλάσμασι, καὶ τὸν οἶκον τὸν μικρὸν ῥάγμασιν.

\vs{12}Εἰ διώξονται ἐν πέτραις ἵπποι; εἰ παρασιωπήσονται ἐν θηλείαις; ὅτι ἐξεστρέψατε εἰς θυμὸν κρίμα, καὶ καρπὸν δικαιοσύνης εἰς πικρίαν,
\vs{13}οἱ εὐφραινόμενοι ἐπʼ οὐδενὶ λόγῳ, οἱ λέγοντες, οὐκ ἐν τῃ ἰσχύϊ ἡμῶν ἔσχομεν κέρατα;
\vs{14}Διότι ἰδοὺ ἐγὼ ἐπεγερῶ ἐφʼ ὑμᾶς οἶκος Ἰσραὴλ ἔθνος, λέγει Κύριος τῶν δυνάμεων, καὶ ἐκθλίψουσιν ὑμᾶς τοῦ μὴ εἰσελθεῖν εἰς Αἰμὰθ, καὶ ὡς τοῦ χειμάῤῥου τῶν δυσμῶν.

\ch{7}
Οὕτως ἔδειξέ μοι Κύριος ὁ Θεός· καὶ ἰδοὺ ἐπιγονὴ ἀκρίδων ἐρχομένη ἑωθινὴ, καὶ ἰδοὺ βροῦχος εἷς, Γὼγ ὁ βασιλεύς.
\vs{2}Καὶ ἔσται ἐὰν συντελέσῃ τοῦ καταφαγεῖν τὸν χόρτον τῆς γῆς, καὶ εἶπα, Κύριε Κύριε, ἵλεως γενοῦ· τίς ἀναστήσει τὸν Ἰακώβ; ὅτι ὀλιγοστός ἐστι.
\vs{3}Μετανόησον Κύριε ἐπὶ τούτῳ. Καὶ τοῦτο οὐκ ἔσται, λέγει Κύριος.

\vs{4}Οὕτως ἔδειξέ μοι Κύριος· καὶ ἰδοὺ ἐκάλεσε τὴν δίκην ἐν πυρὶ Κύριος, καὶ κατέφαγε τὴν ἄβυσσον τὴν πολλὴν, καὶ κατέφαγε τὴν μερίδα Κυρίου.
\vs{5}Καὶ εἶπα, Κύριε κόπασον δὴ, τίς ἀναστήσει τὸν Ἰακώβ; ὅτι ὀλιγοστός ἐστι.
\vs{6}Μετανόησον Κύριε ἐπὶ τούτῳ. Καὶ τοῦτο οὐ μὴ γένηται, λέγει Κύριος.

\vs{7}Οὕτως ἔδειξέ μοι Κύριος· καὶ ἰδοὺ ἑστηκὼς ἐπὶ τείχους ἀδαμαντίνου, καὶ ἐν τῇ χειρὶ αὐτοῦ ἀδάμας.
\vs{8}Καὶ εἶπε Κύριος πρὸς μὲ, τί σὺ ὁρᾷς Ἀμώς; καὶ εἶπα, ἀδάμαντα· καὶ εἶπε Κύριος πρὸς μὲ, ἰδοὺ ἐγὼ ἐντάσσω ἀδάμαντα ἐν μέσῳ λαοῦ μου Ἰσραὴλ, οὐκ ἔτι μὴ προσθῶ τοῦ παρελθεῖν αὐτόν.
\vs{9}Καὶ ἀφανισθήσονται βωμοὶ τοῦ γέλωτος, καὶ αἱ τελεταὶ τοῦ Ἰσραὴλ ἐρημωθήσονται, καὶ ἀναστήσομαι ἐπὶ τὸν οἶκον Ἱεροβοὰμ ἐν ῥομφαίᾳ.

\vs{10}Καὶ ἐξαπέστειλεν Ἀμασίας ὁ ἱερεὺς Βαιθὴλ πρὸς Ἱεροβοὰμ βασιλέα Ἰσραὴλ, λέγων, συστροφὰς ποιεῖται κατὰ σοῦ Ἀμὼς ἐν μέσῳ οἴκου Ἰσραὴλ, οὐ μὴ δύνηται ἡ γῆ ὑπενεγκεῖν πάντας τοὺς λόγους αὐτοῦ.
\vs{11}Διότι τάδε λέγει Ἀμὼς, ἐν ῥομφαίᾳ τελευτήσει Ἱεροβοὰμ, ὁ δὲ Ἰσραὴλ αἰχμάλωτος ἀχθήσεται ἀπὸ τῆς γῆς αὐτοῦ.

\vs{12}Καὶ εἶπεν Ἀμασίας πρὸς Ἀμὼς ὁ ὁρῶν βάδιζε, ἐκχῶρησον σὺ εἰς γῆν Ἰούδα, καὶ ἐκεῖ καταβίου, καὶ ἐκεῖ προφητεύσεις,
\vs{13}εἰς δὲ Βαιθὴλ οὐκ ἔτι προσθήσεις τοῦ προφητεῦσαι, ὅτι ἁγίασμα βασιλέως ἐστὶ, καὶ οἶκος βασιλείας ἐστί.

\vs{14}Καὶ ἀπεκρίθη Ἀμὼς καὶ εἶπε πρὸς Ἀμασίαν, οὐκ ἤμην προφήτης ἐγὼ, οὐδὲ υἱὸς προφήτου, ἀλλʼ ἢ αἰπόλος ἤμην, καὶ κνίζων συκάμινα·
\vs{15}καὶ ἀνέλαβέ με Κύριος ἐκ τῶν προβάτων, καὶ εἶπε Κύριος πρὸς μὲ, βάδιζε, καὶ προφήτευσον ἐπὶ τὸν λαόν μου Ἰσραήλ.
\vs{16}Καὶ νῦν ἄκουε λόγον Κυρίου· σὺ λέγεις, μὴ προφήτευε ἐπὶ τὸν Ἰσραὴλ, καὶ οὐ μὴ ὀχλαγωγήσῃς ἐπὶ τὸν οἶκον Ἰακώβ.
\vs{17}Διὰ τοῦτο τάδε λέγει Κύριος, ἡ γυνή σου ἐν τῇ πόλει πορνεύσει, καὶ οἱ υἱοί σου καὶ αἱ θυγατέρες σου ἐν ῥομφαίᾳ πεσοῦνται, καὶ ἡ γῆ σου ἐν σχοινίῳ καταμετρηθήσεται, καὶ σὺ ἐν γῇ ἀκαθάρτῳ τελευτήσεις, ὁ δὲ Ἰσραὴλ αἰχμάλωτος ἀχθήσεται ἀπὸ τῆς γῆς αὐτοῦ· οὕτως ἔδειξέ μοι Κύριος Κύριος.

\ch{8}
Καὶ ἰδοὺ ἄγγος ἰξευτοῦ. Καὶ εἶπε, τί σὺ βλέπεις Ἀμώς; καὶ εἶπα, ἄγγος ἰξευτοῦ·
\vs{2}καὶ εἶπε Κύριος πρὸς μὲ, ἥκει τὸ πέρας ἐπὶ τὸν λαόν μου Ἰσραὴλ, οὐ προσθήσω ἔτι τοῦ παρελθεῖν αὐτόν.
\vs{3}Καὶ ὀλολύξει τὰ φατνώματα τοῦ ναοῦ ἐν τῇ ἡμέρᾳ ἐκείνῃ, λέγει Κύριος Κύριος· πολὺς ὁ πεπτωκὼς ἐν παντὶ τόπῳ, ἐπιῤῥίψω σιωπήν.

\vs{4}Ἀκούσατε δὴ ταῦτα οἱ ἐκτρίβοντες εἰς τοπρωῒ πένητα, καὶ καταδυναστεύοντες πτωχοὺς ἀπὸ τῆς γῆς,
\vs{5}λέγοντες, πότε διελεύσεται ὁ μὴν, καὶ ἐμπολήσομεν, καὶ τὰ σάββατα, καὶ ἀνοίξομεν θησαυρὸν τοῦ ποιῆσαι μέτρον μικρὸν, καὶ τοῦ μεγαλῦναι στάθμιον, καὶ ποιῆσαι ζυγὸν ἄδικον,
\vs{6}τοῦ κτᾶσθαι ἐν ἀργυρίῳ καὶ πτωχοὺς, καὶ πένητα ἀντὶ ὑποδημάτων, καὶ ἀπὸ παντὸς γεννήματος ἐμπορευσόμεθα.
\vs{7}Ὀμνύει Κύριος κατὰ τῆς ὑπερηφανίας Ἰακὼβ, εἰ ἐπιλησθήσεται εἰς νῖκος πάντα τὰ ἔργα ὑμῶν,
\vs{8}Καὶ ἐπὶ τούτοις οὐ ταραχθήσεται ἡ γῆ, καὶ πενθήσει πᾶς ὁ κατοικῶν ἐν αὐτῇ; καὶ ἀναβήσεται ὡς ποταμὸς συντέλεια, καὶ καταβήσεται ὡς ποταμὸς Αἰγύπτου.

\vs{9}Καὶ ἔσται ἐν τῇ ἡμέρᾳ ἐκείνῃ, λέγει Κύριος Κύριος, δύσεται ὁ ἥλιος μεσημβρίας, καὶ συσκοτάσει ἐπὶ τῆς γῆς ἐν ἡμέρᾳ τὸ φῶς,
\vs{10}καὶ μεταστρέψω τὰς ἑορτὰς ὑμῶν εἰς πένθος, καὶ πάσας τὰς ᾠδὰς ὑμῶν εἰς θρῆνον, καὶ ἀναβιβῶ ἐπὶ πᾶσαν ὀσφὺν σάκκον, καὶ ἐπὶ πᾶσαν κεφαλὴν φαλάκρωμα· καὶ θήσομαι αὐτὸν ὡς πένθος ἀγαπητοῦ, καὶ τοὺς μετʼ αὐτοῦ ὡς ἡμέραν ὀδύνης.

\vs{11}Ἰδοὺ ἡμέραι ἔρχονται, λέγει Κύριος, καὶ ἐξαποστελῶ λιμὸν ἐπὶ τὴν γῆν, οὐ λιμὸν ἄρτων, οὐδὲ δίψαν ὕδατος, ἀλλὰ λιμὸν τοῦ ἀκοῦσαι τὸν λόγον Κυρίου.
\vs{12}Καὶ σαλευθήσονται ὕδατα ἀπὸ τῆς θαλάσσης ἕως θαλάσσης, καὶ ἀπὸ Βοῤῥᾶ ἕως ἀνατολῶν περιδραμοῦνται ζητοῦντες τὸν λόγον τοῦ Κυρίου, καὶ οὐ μὴ εὕρωσιν.
\vs{13}Ἐν τῇ ἡμέρᾳ ἐκείνῃ ἐκλείψουσιν αἱ παρθένοι αἱ καλαὶ, καὶ οἱ νεανίσκοι ἐν δίψει,
\vs{14}οἱ ὀμνύοντες κατὰ τοῦ ἱλασμοῦ Σαμαρείας, καὶ οἱ λέγοντες, ζῇ ὁ θεός σου Δὰν, καὶ ζῇ ὁ θεός σου Βηρσαβεέ· καὶ πεσοῦνται, καὶ οὐ μὴ ἀναστῶσιν ἔτι.

\ch{9}
Εἶδον τὸν Κύριον ἐφεστῶτα ἐπὶ τοῦ θυσιαστηρίου, καὶ εἶπε,

Πάταξον ἐπὶ τὸ ἱλαστήριον, καὶ σεισθήσεται τὰ πρόπυλα, καὶ διάκοψον εἰς κεφαλὰς πάντων· καὶ τοὺς καταλοίπους αὐτῶν ἐν ῥομφαίᾳ ἀποκτενῶ, οὐ μὴ διαφύγῃ ἐξ αὐτῶν φεύγων, καὶ οὐ μὴ διασωθῇ ἐξ αὐτῶν ἀνασωζόμενος.
\vs{2}Ἐὰν κατακρυβῶσιν εἰς ᾅδου, ἐκεῖθεν ἡ χείρ μου ἀνασπάσει αὐτούς· καὶ ἐὰν ἀναβῶσιν εἰς τὸν οὐρανὸν, ἐκεῖθεν κατάξω αὐτούς.
\vs{3}Ἐὰν ἐγκατακρυβῶσιν εἰς τὴν κορυφὴν τοῦ Καρμήλου, ἐκεῖθεν ἐξερευνήσω, καὶ λήψομαι αὐτούς· καὶ ἐὰν καταδύσωσιν ἐξ ὀφθαλμῶν μου εἰς τὰ βάθη τῆς θαλάσσης, ἐκεῖ ἐντελοῦμαι τῷ δράκοντι, καὶ δήξεται αὐτούς.
\vs{4}Καὶ ἐὰν πορευθῶσιν ἐν αἰχμαλωσίᾳ πρὸ προσώπου τῶν ἐχθρῶν αὐτῶν, ἐκεῖ ἐντελοῦμαι τῇ ῥομφαίᾳ, καὶ ἀποκτενεῖ αὐτούς· καὶ στηριῶ τοὺς ὀφθαλμούς μου ἐπʼ αὐτοὺς εἰς κακὰ, καὶ οὐκ εἰς ἀγαθά.

\vs{5}Καὶ Κύριος Κύριος ὁ Θεὸς ὁ παντοκράτωρ, ὁ ἐφαπτόμενος τῆς γῆς, καὶ σαλεύων αὐτὴν, καὶ πενθήσουσι πάντες οἱ κατοικοῦντες αὐτὴν, καὶ ἀναβήσεται ὡς ποταμὸς συντέλεια αὐτῆς, καὶ καταβήσεται ὡς ποταμὸς Αἰγύπτου·
\vs{6}Ὁ οἰκοδομῶν εἰς τὸν οὐρανὸν ἀνάβασιν αὐτοῦ, καὶ τὴν ἐπαγγελίαν αὐτοῦ ἐπὶ τῆς γῆς θεμελιῶν, ὁ προσκαλούμενος τὸ ὕδωρ τῆς θαλάσσης, καὶ ἐκχέων αὐτὸ ἐπὶ πρόσωπον τῆς γῆς· Κύριος παντοκράτωρ ὄνομα αὐτῷ.

\vs{7}Οὐχ ὡς υἱοὶ Αἰθιόπων ὑμεῖς ἐστὲ ἐμοὶ, υἱοὶ Ἰσραὴλ; λέγει Κύριος· οὐ τὸν Ἰσραὴλ ἀνήγαγον ἐκ γῆς Αἰγύπτου, καὶ τοὺς ἀλλοφύλους ἐκ Καππαδοκίας, καὶ τοὺς Σύρους ἐκ βόθρου;
\vs{8}Ἰδοὺ οἱ ὀφθαλμοὶ Κυρίου τοῦ Θεοῦ ἐπὶ τὴν βασιλείαν τῶν ἁμαρτωλῶν, καὶ ἐξαρῶ αὐτὴν ἀπὸ προσώπου τῆς γῆς· πλὴν ὅτι οὐκ εἰς τέλος ἐξαρῶ τὸν οἶκον Ἰακὼβ, λέγει Κύριος.
\vs{9}Διότι ἐγὼ ἐντέλλομαι, καὶ λικμήσω ἐν πᾶσι τοῖς ἔθνεσι τὸν οἶκον Ἰσραὴλ, ὃν τρόπον λικμᾶται ἐν τῷ λικμῷ, καὶ οὐ μὴ πέσῃ σύντριμμα ἐπὶ τὴν γῆν.
\vs{10}Ἐν ῥομφαίᾳ τελευτήσουσι πάντες ἁμαρτωλοὶ λαοῦ μου, οἱ λέγοντες, οὐ μὴ ἐγγίσῃ, οὐδὲ μὴ γένηται ἐφʼ ἡμᾶς τὰ κακά.

\vs{11}Ἐν τῇ ἡμέρᾳ ἐκείνῃ ἀναστήσω τὴν σκηνὴν Δαυὶδ τὴν πεπτωκυῖαν, καὶ ἀνοικοδομήσω τὰ πεπτωκότα αὐτῆς, καὶ τὰ κατεσκαμμένα αὐτῆς ἀναστήσω, καὶ ἀνοικοδομήσω αὐτὴν καθὼς αἱ ἡμέραι τοῦ αἰῶνος.
\vs{12}Ὅπως ἐκζητήσωσιν οἱ κατάλοιποι τῶν ἀνθρώπων καὶ πάντα τὰ ἔθνη, ἐφʼ οὓς ἐπικέκληται τὸ ὄνομά μου ἐπʼ αὐτοὺς, λέγει Κύριος ὁ ποιῶν πάντα ταῦτα.

\vs{13}Ἰδοὺ ἡμέραι ἔρχονται, λέγει Κύριος, καὶ καταλήψεται ὁ ἀμητὸς τὸν τρυγητὸν, καὶ περκάσει ἡ σταφυλὴ ἐν τῷ σπόρῳ, καὶ ἀποσταλάξει τὰ ὄρη γλυκασμὸν, καὶ πάντες οἱ βουνοὶ σύμφυτοι ἔσονται.
\vs{14}Καὶ ἐπιστρέψω τὴν αἰχμαλωσίαν τοῦ λαοῦ μου Ἰσραὴλ, καὶ οἰκοδομήσουσι πόλεις τὰς ἠφανισμένας, καὶ κατοικήσουσι, καὶ φυτεύσουσιν ἀμπελῶνας, καὶ πίονται τὸν οἶνον αὐτῶν, καὶ ποιήσουσι κήπους, καὶ φάγονται τὸν καρπὸν αὐτῶν·
\vs{15}Καὶ καταφυτεύσω αὐτοὺς ἐπὶ τῆς γῆς αὐτῶν, καὶ οὐ μὴ ἐκσπασθῶσιν οὐκέτι ἀπὸ τῆς γῆς, ἧς ἔδωκα αὐτοῖς, λέγει Κύριος ὁ Θεὸς παντοκράτωρ.


\def\book{ΟΒΔΕΙΟΥ}
\biblebook{ΟΒΔΕΙΟΥ}

 
\lettrine[lines=2, loversize=0.2, nindent=0em, findent=.25em]{\textcolor{bookheadingcolor}{Ὅ}}{ΡΑΣΙΣ} Ὀβδίου. Τάδε λέγει Κύριος ὁ Θεὸς τῇ Ἰδουμαίᾳ, ἀκοὴν ἤκουσα παρὰ Κυρίου, καὶ περιοχὴν εἰς τὰ ἔθνη ἐξαπέστειλεν· ἀνάστητε, καὶ ἐξαναστῶμεν ἐπʼ αὐτὴν εἰς πόλεμον.

\vs{2}Ἰδοὺ ὀλιγοστὸν δέδωκά σε ἐν τοῖς ἔθνεσιν, ἠτιμωμένος εἶ σὺ σφόδρα.
\vs{3}Ὑπερηφανία τῆς καρδίας σου ἐπῇρέ σε κατασκηνοῦντα ἐν ταῖς ὀπαῖς τῶν πετρῶν· ὑψῶν κατοικίαν αὐτοῦ, λέγων ἐν καρδίᾳ αὐτοῦ, τίς κατάξει με ἐπὶ τὴν γῆν;
\vs{4}Ἐὰν μετεωρισθῇς, ὡς ἀετὸς, καὶ ἐὰν ἀναμέσον τῶν ἄστρων θῇς νοσσιάν σου, ἐκεῖθεν κατάξω σε, λέγει Κύριος.
\vs{5}Εἰ κλέπται εἰσῆλθον πρὸς σὲ, ἢ λῃσταὶ νυκτὸς, ποῦ ἂν ἀπεῤῥίφης; οὐκ ἂν ἔκλεψαν τὰ ἱκανὰ ἑαυτοῖς; καὶ εἰ τρυγηταὶ εἰσῆλθον πρὸς σὲ, οὐκ ἂν ἐπελείποντο ἐπιφυλλίδα;

\vs{6}Πῶς ἐξηρευνήθη Ἡσαῦ, καὶ κατελήφθη τὰ κεκρυμμένα αὐτοῦ;
\vs{7}Ἕως τῶν ὁρίων ἐξαπέστειλάν σε· πάντες οἱ ἄνδρες τῆς διαθήκης σου ἀντέστησάν σοι, ἠδυνάσθησαν πρὸς σὲ ἄνδρες εἰρηνικοί σου, ἔθηκαν ἔνεδρα ὑποκάτω σου, οὐκ ἔστι σύνεσις αὐτοῖς.

\vs{8}Ἐν τῇ ἡμέρᾳ ἐκείνῃ, λέγει Κύριος, ἀπολῶ σοφους ἐκ τῆς Ἰδουμαίας, καὶ σύνεσιν ἐξ ὄρους Ἡσαῦ.
\vs{9}Καὶ πτοηθήσονται οἱ μαχηταί σου οἱ ἐκ Θαιμὰν, ὅπως ἐξαρθῇ ἄνθρωπος ἐξ ὄρους Ἡσαῦ.
\vs{10}Διὰ τὴν σφαγὴν, καὶ τὴν ἀσέβειαν ἀδελφοῦ σου Ἰακὼβ, καλύψει σε αἰσχύνη, καὶ ἐξαρθήσῃ εἰς τὸν αἰῶνα.
\vs{11}Ἀφʼ ἧς ἡμέρας ἀντέστης ἐξεναντίας, ἐν ἡμέραις αἰχμαλωτευόντων ἀλλογενῶν δύναμιν αὐτοῦ, καὶ ἀλλότριοι εἰσῆλθον εἰς πύλας αὐτοῦ, καὶ ἐπὶ Ἱερουσαλὴμ ἔβαλον κλήρους, καὶ σὺ ἦς ὡς εἷς ἐξ αὐτῶν.

\vs{12}Καὶ μὴ ἐπίδῃς ἡμέραν ἀδελφοῦ σου ἐν ἡμέρᾳ ἀλλοτρίων, καὶ μὴ ἐπιχαρῇς ἐπὶ τοὺς υἱοὺς Ἰούδα ἐν ἡμέρᾳ ἀπωλείας αὐτῶν, καὶ μὴ μεγαλοῤῥημονῇ ἐν ἡμέρᾳ θλίψεως,
\vs{13}μηδὲ εἰσέλθῃς εἰς πύλας λαῶν ἐν ἡμέρᾳ πόνων αὐτῶν, μηδὲ ἐπίδῃς καὶ σὺ τὴν συναγωγὴν αὐτῶν ἐν ἡμέρᾳ ὀλέθρου αὐτῶν, καὶ μὴ συνεπιθῇ ἐπὶ τὴν δύναμιν αὐτῶν ἐν ἡμέρᾳ ἀπωλείας αὐτῶν,
\vs{14}μηδὲ ἐπιστῇς ἐπὶ τὰς διεκβολὰς αὐτῶν, ἐξολοθρεῦσαι τοὺς ἀνασωζομένους αὐτῶν, μηδὲ συγκλείσῃς τοὺς φεύγοντας αὐτοῦ ἐν ἡμέρᾳ θλίψεως.

\vs{15}Διότι ἐγγὺς ἡμέρα Κυρίου ἐπὶ πάντα τὰ ἔθνη· ὃν τρόπον ἐποίησας, οὕτως ἔσται σοι· τὸ ἀνταπόδομά σου ἀνταποδοθήσεται εἰς κεφαλήν σου.
\vs{16}Διότι ὃν τρόπον ἔπιες ἐπὶ τὸ ὄρος τὸ ἅγιόν μου, πίονται πάντα τὰ ἔθνη οἶνον, πίονται καὶ καταβήσονται, καὶ ἔσονται καθὼς οὐχ ὑπάρχοντες.

\vs{17}Ἑν δὲ τῷ ὄρει Σιὼν ἔσται σωτηρία, καὶ ἔσται ἅγιον· καὶ κατακληρονομήσουσιν ὁ οἶκος Ἰακὼβ τοῦς, κατακληρονομήσαντας αὐτούς.
\vs{18}Καὶ ἔσται ὁ οἶκος Ἰακὼβ πῦρ, ὁ δὲ οἶκος Ἰωσὴφ φλὸξ, ὁ δὲ οἶκος Ἡσαῦ εἰς καλάμην, καὶ ἐκκαυθήσονται εἰς αὐτοὺς, καὶ καταφάγονται αὐτοὺς, καὶ οὐκ ἔσται πυροφόρος τῷ οἴκῳ Ἡσαῦ, διότι Κύριος ἐλάλησε.
\vs{19}Καὶ κατακληρονομήσουσιν οἱ ἐν ναγὲβ τὸ ὄρος τὸ Ἡσαῦ, καὶ οἱ ἐν τῇ σεφηλὰ τοὺς ἀλλοφύλους· καὶ κατακληρονομήσουσι τὸ ὄρος Ἐφραὶμ, καὶ τὸ πεδίον Σαμαρείας, καὶ Βενιαμὶν, καὶ τὴν Γαλααδίτιν.

\vs{20}Καὶ τῆς μετοικεσίας ἡ ἀρχὴ αὕτη τοῖς υἱοῖς Ἰσραὴλ, γῆ τῶν Χαναναίων ἕως Σαρεπτῶν· καὶ ἡ μετοικεσία Ἱερουσαλὴμ ἕως Ἐφραθά· κληρονομήσουσι τὰς πόλεις τοῦ ναγέβ.

\vs{21}Καὶ ἀναβήσονται ἀνασωζόμενοι ἐξ ὄρους Σιὼν, τοῦ ἐκδικῆσαι τὸ ὄρος Ἡσαῦ, καὶ ἔσται τῷ Κυρίῳ ἡ βασιλεία.


\def\book{ΙΩΝΑΣ}
\biblebook{ΙΩΝΑΣ}


\lettrine[lines=2, loversize=0.2, nindent=0em, findent=.25em]{\textcolor{bookheadingcolor}{Κ}}{ΑΙ} ἐγένετο λόγος Κυρίου πρὸς Ἰωνᾶν τὸν τοῦ Ἀμαθὶ, λέγων,
\vs{2}ἀνάστηθι, καὶ πορεύθητι εἰς Νινευὴ τὴν πόλιν τὴν μεγάλην, καὶ κήρυξον ἐν αὐτῇ, ὅτι ἀνέβη ἡ κραυγὴ τῆς κακίας αὐτῆς πρὸς μέ.
\vs{3}Καὶ ἀνέστη Ἰωνᾶς τοῦ φυγεῖν εἰς Θαρσὶς ἐκ προσώπου Κυρίου· καὶ κατέβη εἰς Ἰόππην, καὶ εὗρε πλοῖον βαδίζον εἰς Θαρσὶς, καὶ ἔδωκε τὸ ναῦλον αὐτοῦ, καὶ ἀνέβη εἰς αὐτὸ, τοῦ πλεῦσαι μετʼ αὐτῶν εἰς Θαρσὶς ἐκ προσώπου Κυρίου.

\vs{4}Καὶ Κύριος ἐξήγειρε πνεῦμα ἐπὶ τὴν θάλασσαν, καὶ ἐγένετο κλύδων μέγας ἐν τῇ θαλάσσῃ, καὶ τὸ πλοῖον ἐκινδύνευε τοῦ συντριβῆναι.
\vs{5}Καὶ ἐφοβήθησαν οἱ ναυτικοὶ, καὶ ἀνεβόησαν ἕκαστος πρὸς τὸν θεὸν αὐτοῦ, καὶ ἐκβολὴν ἐποιήσαντο τῶν σκευῶν τῶν ἐν τῷ πλοίῳ εἰς τὴν θάλασσαν, τοῦ κουφισθῆναι ἀπʼ αὐτῶν· Ἰωνᾶς δὲ κατέβη εἰς τὴν κοίλην τοῦ πλοίου, καὶ ἐκάθευδε, καὶ ἔρεγχε.

\vs{6}Καὶ προσῆλθε πρὸς αὐτὸν ὁ πρωρεὺς, καὶ εἶπεν αὐτῷ, τί σὺ ῥέγχεις; ἀνάστα, καὶ ἐπικαλοῦ τὸν Θεόν σου, ὅπως διασώσῃ ὁ Θεὸς ἡμᾶς, καὶ οὐ μὴ ἀπολώμεθα.
\vs{7}Καὶ εἶπεν ἕκαστος πρὸς τὸν πλησίον αὐτοῦ, δεῦτε βάλωμεν κλήρους, καὶ ἐπιγνῶμεν, τίνος ἕνεκεν ἡ κακία αὕτη ἐστὶν ἐν ἡμῖν· καὶ ἔβαλον κλήρους, καὶ ἔπεσεν ὁ κλῆρος ἐπὶ Ἰωνᾶν.

\vs{8}Καὶ εἶπον πρὸς αὐτὸν, ἀπάγγειλον ἡμῖν, τίς σου ἡ ἐργασία ἐστὶ, καὶ πόθεν ἔρχῃ, καὶ ἐκ ποίας χώρας, καὶ ἐκ ποίου λαοῦ εἶ σύ;
\vs{9}Καὶ εἶπε πρὸς αὐτοὺς, δοῦλος Κυρίου εἰμὶ ἐγὼ, καὶ τὸν Κύριον Θεὸν τοῦ οὐρανοῦ ἐγὼ σέβομαι, ὃς ἐποίησε τὴν θάλασσαν καὶ τὴν ξηράν.
\vs{10}Καὶ ἐφοβήθησαν οἱ ἄνδρες φόβον μέγαν, καὶ εἶπον πρὸς αὐτὸν, τί τοῦτο ἐποίησας; διότι ἔγνωσαν οἱ ἄνδρες ὅτι ἐκ προσώπου Κυρίου ἦν φεύγων, ὅτι ἀπήγγειλεν αὐτοῖς·
\vs{11}καὶ εἶπον πρὸς αὐτὸν, τί ποιήσομέν σοι, καὶ κοπάσει ἡ θάλασσα ἀφʼ ἡμῶν; ὅτι ἡ θάλασσα ἐπορεύετο καὶ ἐξήγειρε μᾶλλον κλύδωνα.
\vs{12}Καὶ εἶπεν Ἰωνᾶς πρὸς αὐτοὺς, ἄρατέ με, καὶ ἐμβάλετέ με εἰς τὴν θάλασσαν, καὶ κοπάσει ἡ θάλασσα ἀφʼ ὑμῶν· διότι ἔγνωκα ἐγὼ, ὅτι διʼ ἐμὲ ὁ κλύδων ὁ μέγας οὗτος ἐφʼ ὑμᾶς ἐστι.

\vs{13}Καὶ παρεβιάζοντο οἱ ἄνδρες τοῦ ἐπιστρέψαι πρὸς τὴν γῆν, καὶ οὐκ ἠδύναντο, ὅτι ἡ θάλασσα ἐπορεύετο, καὶ ἐξηγείρετο μᾶλλον ἐπʼ αὐτούς.
\vs{14}Καὶ ἀνεβόησαν πρὸς Κύριον, καὶ εἶπαν, μηδαμῶς Κύριε· μὴ ἀπολώμεθα ἕνεκεν τῆς ψυχῆς τοῦ ἀνθρώπου τούτου, καὶ μὴ δῷς ἐφʼ ἡμᾶς αἷμα δίκαιον, διότι σὺ Κύριε, ὃν τρόπον ἐβούλου, πεποίηκας.
\vs{15}Καὶ ἔλαβον τὸν Ἰωνᾶν, καὶ ἐξέβαλον αὐτὸν εἰς τὴν θάλασσαν, καὶ ἔστη ἡ θάλασσα ἐκ τοῦ σάλου αὐτῆς.
\vs{16}Καὶ ἐφοβήθησαν οἱ ἄνδρες φόβῳ μεγάλῳ τὸν Κύριον, καὶ ἔθυσαν θυσίαν τῷ Κυρίῳ, καὶ ηὔξαντο τὰς εὐχάς.

\ch{2}
Καὶ προσέταξε Κύριος κήτει μεγάλῳ καταπιεῖν τὸν Ἰωνᾶν· καὶ ἦν Ἰωνᾶς ἐν τῇ κοιλίᾳ τοῦ κήτους τρεῖς ἡμέρας καὶ τρεῖς νύκτας.

\vs{2}Καὶ προσηύξατο Ἰωνᾶς πρὸς Κύριον τὸν Θεὸν αὐτοῦ ἐκ τῆς κοιλίας τοῦ κήτους,
\vs{3}καὶ εἶπεν,

Ἐβόησα ἐν θλίψει μου πρὸς Κύριον τὸν Θεόν μου, καὶ εἰσήκουσέ μου, ἐκ κοιλίας ᾅδου κραυγῆς μου, ἤκουσας φωνῆς μου, ἀπέῤῥιψάς με εἰς βάθη καρδίας θαλάσσης,
\vs{4}καὶ ποταμοί ἐκύκλωσάν με, πάντες οἱ μετεωρισμοί σου καὶ τὰ κύματά σου ἐπʼ ἐμὲ διῆλθον.
\vs{5}Καὶ ἐγὼ εἶπα, ἀπῶσμαι ἐξ ὀφθαλμῶν σου· ἆρα προσθήσω τοῦ ἐπιβλέψαι με πρὸς ναὸν τὸν ἅγιόν σου;
\vs{6}Περιεχύθη μοι ὕδωρ ἕως ψυχῆς, ἄβυσσος ἐκύκλωσέ με ἐσχάτη, ἔδυ ἡ κεφαλή μου εἰς σχισμὰς ὀρέων,
\vs{7}κατέβην εἰς γῆν, ἧς οἱ μοχλοὶ αὐτῆς κάτοχοι αἰώνιοι· καὶ ἀναβήτω φθορὰ ζωῆς μου Κύριε ὁ Θεός μου.

\vs{8}Ἐν τῷ ἐκλείπειν ἀπʼ ἐμοῦ τὴν ψυχήν μου, τοῦ Κυρίου ἐμνήσθην, καὶ ἔλθοι πρὸς σὲ ἡ προσευχή μου εἰς ναὸν τὸν ἅγιόν σου.
\vs{9}Φυλασσόμενοι μάταια καὶ ψευδῆ, ἔλεος αὐτῶν ἐγκατέλιπον.
\vs{10}Ἐγὼ δὲ μετὰ φωνῆς αἰνέσεως καὶ ἐξομολογήσεως θύσω σοι, ὅσα ηὐξάμην ἀποδώσω σοι σωτηρίου τῷ Κυρίῳ.

\vs{11}Καὶ προσετάγη ἀπὸ Κυρίου τῷ κήτει, καὶ ἐξέβαλε τὸν Ἰωνᾶν ἐπὶ τὴν ξηράν.

\ch{3}
Καὶ ἐγένετο λόγος Κυρίου πρὸς Ἰωνᾶν ἐκ δευτέρου, λέγων,
\vs{2}ἀνάστηθι, πορεύθητι εἰς Νινευὴ τὴν πόλιν τὴν μεγάλην; καὶ κήρυξον ἐν αὐτῇ κατὰ τὸ κήρυγμα τὸ ἔμπροσθεν, ὃ ἐγὼ ἐλάλησα πρὸς σέ.
\vs{3}Καὶ ἀνέστη Ἰωνᾶς, καὶ ἐπορεύθη εἰς Νινευὴ, καθὰ ἐλάλησε Κύριος· ἡ δὲ Νινευὴ ἦν πόλις μεγάλη τῷ Θεῷ, ὡσεὶ πορείας ὁδοῦ τριῶν ἡμερῶν·
\vs{4}καὶ ἤρξατο Ἰωνᾶς τοῦ εἰσελθεῖν εἰς τὴν πόλιν, ὡσεὶ πορείαν ἡμέρας μιᾶς· καὶ ἐκήρυξε, καὶ εἶπεν, ἔτι τρεῖς ἡμέραι, καὶ Νινευὴ καταστραφήσεται.

\vs{5}Καὶ ἐπίστευσαν οἱ ἄνδρες Νινευὴ τῷ Θεῷ, καὶ ἐκήρυξαν νηστείαν, καὶ ἐνεδύσαντο σάκκους ἀπὸ μεγάλου αὐτῶν ἕως μικροῦ αὐτῶν.
\vs{6}Καὶ ἤγγισεν ὁ λόγος πρὸς τὸν βασιλέα τῆς Νινευὴ, καὶ ἐξανέστη ἀπὸ τοῦ θρόνου αὐτοῦ, καὶ περιείλατο τὴν στολὴν αὐτοῦ ἀφʼ ἑαυτοῦ, καὶ περιεβάλετο σάκκον, καὶ ἐκάθισεν ἐπὶ σποδοῦ.
\vs{7}Καὶ ἐκηρύχθη, καὶ ἐῤῥέθη ἐν τῇ Νινευὴ παρὰ τοῦ βασιλέως καὶ παρὰ τῶν μεγιστάνων αὐτοῦ, λέγων, οἱ ἄνθρωποι, καὶ τὰ κτήνη, καὶ οἱ βόες, καὶ τὰ πρόβατα μὴ γευσάσθωσαν, μηδὲ νεμέσθωσαν, μηδὲ ὕδωρ πιέτωσαν.
\vs{8}Καὶ περιεβάλλοντο σάκκους οἱ ἄνθρωποι καὶ τὰ κτήνη, καὶ ἀνεβόησαν πρὸς τὸν Θεὸν ἐκτενῶς· καὶ ἀπέστρεψαν ἕκαστος ἀπὸ τῆς ὁδοῦ αὐτῶν τῆς πονηρᾶς, καὶ ἀπὸ τῆς ἀδικίας τῆς ἐν χερσὶν αὐτῶν, λέγοντες,
\vs{9}τίς οἶδεν εἰ μετανοήσει ὁ Θεὸς, καὶ ἀποστρέψει ἐξ ὀργῆς θυμοῦ αὐτοῦ, καὶ οὐ μὴ ἀπολώμεθα;

\vs{10}Καὶ εἶδεν ὁ Θεὸς τὰ ἔργα αὐτῶν, ὅτι ἀπέστρεψαν ἀπὸ τῶν ὁδῶν αὐτῶν τῶν πονηρῶν, καὶ μετενόησεν ὁ Θεὸς ἐπὶ τῇ κακίᾳ, ᾗ ἐλάλησε τοῦ ποιῆσαι αὐτοῖς, καὶ οὐκ ἐποίησε.

\ch{4}
Καὶ ἐλυπήθη Ἰωνᾶς λύπην μεγάλην· καὶ συνεχύθη,
\vs{2}καὶ προσεύξατο πρὸς Κύριον, καὶ εἶπεν, Κύριε, οὐχ οὗτοι οἱ λόγοι μου, ἔτι ὄντος μου ἐν τῇ γῇ μου; διατοῦτο προέφθασα τοῦ φυγεῖν εἰς Θαρσὶς, διότι ἔγνων ὅτι σὺ ἐλεήμων καὶ οἰκτίρμων, μακρόθυμος καὶ πολυέλεος, καὶ μετανοῶν ἐπὶ ταῖς κακίαις.
\vs{3}Καὶ νῦν, δέσποτα Κύριε, λάβε τὴν ψυχήν μου ἀπʼ ἐμοῦ, ὅτι καλὸν τὸ ἀποθανεῖν με ἢ ζῇν με.
\vs{4}Καὶ εἶπε Κύριος πρὸς Ἰωνᾶν, εἰ σφόδρα λελύπησαι σύ;

\vs{5}Καὶ ἐξῆλθεν Ἰωνᾶς ἐκ τῆς πόλεως, καὶ ἐκάθισεν ἀπέναντι τῆς πόλεως· καὶ ἐποίησεν αὐτῷ ἐκεῖ σκηνὴν, καὶ ἐκάθητο ὑποκάτω αὐτῆς, ἕως οὗ ἀπίδῃ τί ἔσται τῇ πόλει.
\vs{6}Καὶ προσέταξε Κύριος ὁ Θεὸς κολοκύνθῃ, καὶ ἀνέβη ὑπὲρ κεφαλῆς τοῦ Ἰωνᾶ, τοῦ εἶναι σκιὰν ὑπεράνω τῆς κεφαλῆς αὐτοῦ, τοῦ σκιάζειν αὐτῷ ἀπὸ τῶν κακῶν αὐτοῦ· καὶ ἐχάρη Ἰωνᾶς ἐπὶ τῇ κολοκύνθῃ χαρὰν μεγάλην.

\vs{7}Καὶ προσέταξεν ὁ Θεὸς σκώληκι ἑωθινῇ τῇ ἐπαυρίον, καὶ ἐπάταξε τὴν κολόκυνθαν, καὶ ἀπεξηράνθη.
\vs{8}Καὶ ἐγένετο ἅμα τῷ ἀνατεῖλαι τὸν ἥλιον, καὶ προσέταξεν ὁ Θεὸς πνεύματι καύσωνι συγκαίοντι, καὶ ἐπάταξεν ὁ ἥλιος ἐπὶ τὴν κεφαλὴν τοῦ Ἰωνᾶ· καὶ ὠλιγοψύχησε, καὶ ἀπελέγετο τὴν ψυχὴν αὐτοῦ, καὶ εἶπε, καλόν μοι ἀποθανεῖν με ἢ ζῇν.
\vs{9}Καὶ εἶπεν ὁ Θεὸς πρὸς Ἰωνᾶν, εἰ σφόδρα λελύπησαι σὺ ἐπὶ τῇ κολοκύνθῃ; καὶ εἶπε, σφόδρα λελύπημαι ἐγὼ ἕως θανάτου.

\vs{10}Καὶ εἶπε Κύριος, σὺ ἐφείσω ὑπὲρ τῆς κολοκύνθης, ὑπὲρ ἧς οὐκ ἐκακοπάθησας ἐπʼ αὐτὴν, καὶ οὐδὲ ἐξέθρεψας αὐτὴν, ἣ ἐγενήθη ὑπὸ νύκτα, καὶ ὑπὸ νύκτα ἀπώλετο·
\vs{11}ἐγὼ δὲ οὐ φείσομαι ὑπὲρ Νινευὴ τῆς πόλεως τῆς μεγάλης, ἐν ᾗ κατοικοῦσι πλείους ἢ δώδεκα μυριάδες ἀνθρώπων, οἵτινες οὐκ ἔγνωσαν δεξιὰν αὐτῶν ἢ ἀριστερὰν αὐτῶν, καὶ κτήνη πολλά;


\def\book{ΜΙΧΑΙΑΣ}
\biblebook{ΜΙΧΑΙΑΣ}


\lettrine[lines=2, loversize=0.2, nindent=0em, findent=.25em]{\textcolor{bookheadingcolor}{Κ}}{ΑΙ} ἐγένετο λόγος Κυρίου πρὸς Μιχαίαν τὸν τοῦ Μωρασθεὶ, ἐν ἡμέραις Ἰωάθαμ, καὶ Ἄχαζ, καὶ Ἐζεκίου βασιλέων Ἰούδα, ὑπὲρ ὧν εἶδε περὶ Σαμαρείας καὶ περὶ Ἱερουσαλήμ.

\vs{2}Ἀκούσατε λαοὶ λόγους, καὶ προσεχέτω ἡ γῆ, καὶ πάντες οἱ ἐν αὐτῇ· καὶ ἔσται Κύριος Κύριος ἐν ὑμῖν εἰς μαρτύριον, Κύριος ἐξ οἴκου ἁγίου αὐτοῦ.
\vs{3}Διότι ἰδοὺ Κύριος ἐκπορεύεται ἐκ τοῦ τόπου αὐτοῦ, καὶ καταβήσεται, καὶ ἐπιβήσεται ἐπὶ τὰ ὕψη τῆς γῆς,
\vs{4}καὶ σαλευθήσεται τὰ ὄρη ὑποκάτωθεν αὐτοῦ, καὶ αἱ κοιλάδες τακήσονται ὡσεὶ κηρὸς ἀπὸ προσώπου πυρὸς, καὶ ὡς ὕδωρ καταφερόμενον ἐν καταβάσει.

\vs{5}Διʼ ἀσέβειαν Ἰακὼβ πάντα ταῦτα, καὶ διʼ ἁμαρτίαν οἴκου Ἰσραήλ· τίς ἡ ἀσέβεια τοῦ Ἰακώβ; οὐχ ἡ Σαμάρεια; καὶ τίς ἡ ἁμαρτία οἴκου Ἰούδα; οὐχὶ Ἱερουσαλήμ;
\vs{6}Καὶ θήσομαι Σαμάρειαν εἰς ὀπωροφυλάκιον ἀγροῦ, καὶ εἰς φυτείαν ἀμπελῶνος, καὶ κατασπάσω εἰς χάος τοὺς λίθους αὐτῆς, καὶ τὰ θεμέλια αὐτῆς ἀποκαλύψω.
\vs{7}Καὶ πάντα τὰ γλυπτὰ αὐτῆς κατακόψουσι, καὶ πάντα τὰ μισθώματα αὐτῆς ἐμπρήσουσιν ἐν πυρὶ, καὶ πάντα τὰ εἴδωλα αὐτῆς θήσομαι εἰς ἀφανισμόν· διότι ἐκ μισθωμάτων πορνείας συνήγαγε, καὶ ἐκ μισθωμάτων πορνείας συνέστρεψεν.

\vs{8}Ἕνεκεν τούτου κόψεται, καὶ θρηνήσει, πορεύσεται ἀνυπόδετος, καὶ γυμνὴ ποιήσεται κοπετὸν ὡς δρακόντων, καὶ πένθος ὡς θυγατέρων σειρήνων.
\vs{9}Ὅτι κατεκράτησεν ἡ πληγὴ αὐτῆς, διότι ἦλθεν ἕως Ἰούδα, καὶ ἥψατο ἕως πύλης λαοῦ μου, ἕως Ἱερουσαλήμ.

\vs{10}Οἱ ἐν Γὲθ μὴ μεγαλύνεσθε, καὶ οἱ Ἐνακεὶμ μὴ ἀνοικοδομεῖτε ἐξ οἴκου κατὰ γέλωτα, γῆν καταπάσασθε καταγέλωτα ὑμῶν,
\vs{11}κατοικοῦσα καλῶς τὰς πόλεις αὐτῆς, οὐκ ἐξῆλθε κατοικοῦσα Σενναὰρ, κόψασθαι οἶκον ἐχόμενον αὐτῆς, λήψεται ἐξ ὑμῶν πληγὴν ὀδύνης.

\vs{12}Τίς ἥρξατο εἰς ἀγαθὰ κατοικούσῃ ὀδύνας; ὅτι κατέβη κακὰ παρὰ Κυρίου ἐπὶ πύλας Ἱερουσαλὴμ,
\vs{13}ψόφος ἁρμάτων καὶ ἱππευόντων· κατοικοῦσα Λαχεὶς, ἀρχηγὸς ἁμαρτίας αὕτη ἐστὶ τῇ θυγατρὶ Σιὼν, ὅτι ἐν σοὶ εὑρέθησαν ἀσέβειαι τοῦ Ἰσραήλ.
\vs{14}Διατοῦτο δώσει ἐξαποστελλομένους ἕως κληρονομίας Γὲθ, οἴκους ματαίους, εἰς κενὸν ἐγένοντο τοῖς βασιλεῦσι τοῦ Ἰσραὴλ,
\vs{15}ἕως τοὺς κληρονόμους ἀγάγωσι, κατοικοῦσα Λαχείς· κληρονομία ἕως Ὀδολλὰμ ἥξει, ἡ δόξα τὴς θυγατρὸς Ἰσραήλ.
\vs{16}Ξύρησαι, καὶ κεῖραι ἐπὶ τὰ τέκνα τὰ τρυφερά σου, ἐμπλάτυνον τὴν χηρείαν σου ὡς ἀετὸς, ὅτι ᾐχμαλωτεύθησαν ἀπὸ σοῦ.

\ch{2}
Ἐγένοντο λογιζόμενοι κόπους, καὶ ἐργαζόμενοι κακὰ ἐν ταῖς κοίταις αὐτῶν, καὶ ἅμα τῇ ἡμέρᾳ συνετέλουν αὐτὰ, διότι οὐκ ᾖραν πρὸς τὸν Θεὸν χεῖρας αὐτῶν.
\vs{2}Καὶ ἐπεθύμουν ἀγρούς, καὶ διήρπαζον ὀρφανούς, καὶ οἴκους κατεδυνάστευον, καὶ διήρπαζον ἄνδρα καὶ τὸν οἶκον αὐτοῦ, καὶ ἄνδρα καὶ τὴν κληρονομίαν αὐτοῦ.

\vs{3}Διατοῦτο τάδε λέγει Κύριος ἰδοὺ ἐγὼ λογίζομαι ἐπὶ τὴν φυλὴν ταύτην κακὰ, ἐξ ὧν οὐ μὴ ἄρητε τοὺς τραχήλους ὑμῶν, καὶ οὐ μὴ πορευθῆτε ὀρθοὶ ἐξαίφνης, ὅτι καιρὸς πονηρός ἐστιν.

\vs{4}Ἐν τῇ ἡμέρᾳ ἐκείνῃ ληφθήσεται ἐφʼ ὑμᾶς παραβολὴ, καὶ θρηνηθήσεται θρῆνος ἐν μέλει, λέγων, ταλαιπωρίᾳ ἐταλαιπωρήσαμεν· μερὶς λὰοῦ μου κατεμετρήθη ἐν σχοινίῳ, καὶ οὐκ ἦν ὁ κωλύων αὐτὸν τοῦ ἀποστρέψαι· οἱ ἀγροὶ ὑμῶν διεμερίσθησαν.
\vs{5}Διατοῦτο οὐκ ἔσται σοι βάλλων σχοινίον ἐν κλήρῳ· ἐν ἐκκλησίᾳ Κυρίου
\vs{6}μὴ κλαίετε δάκρυσι, μηδὲ δακρυέτωσαν ἐπὶ τούτοις· οὐδὲ γὰρ ἀπώσεται ὀνείδη,
\vs{7}ὁ λέγων, οἶκος Ἰακὼβ παρώργισε πνεῦμα Κυρίου. οὐ ταῦτα τὰ ἐπιτηδεύματα αὐτοῦ ἐστιν; οὐχ οἱ λόγοι αὐτοῦ εἰσι καλοὶ μετʼ αὐτοῦ; καὶ ὀρθοὶ πεπόρευνται;
\vs{8}Καὶ ἔμπροσθεν ὁ λαός μου εἰς ἔχθραν ἀντέστη, κατέναντι τῆς εἰρήνης αὐτοῦ· τὴν δορὰν αὐτοῦ ἐξέδειραν, τοῦ ἀφελέσθαι ἐλπίδας συντριμμὸν πολέμου.
\vs{9}Ἡγούμενοι λαοῦ μου ἀποῤῥιφήσονται ἐκ τῶν οἰκιῶν τρυφῆς αὐτῶν, διὰ τὰ πονηρὰ ἐπιτηδεύματα αὐτῶν ἐξώσθησαν· ἐγγίσατε ὄρεσιν αἰωνίοις.

\vs{10}Ἀνάστηθι καὶ πορεύου, ὅτι οὐκ ἔστι σοι αὕτη ἀνάπαυσις ἕνεκεν ἀκαθαρσίας· διεφθάρητε φθορᾷ,
\vs{11}κατεδιώχθητε οὐδενὸς διώκοντος· πνεῦμα ἔστησε ψεῦδος, ἐστάλαξέ σοι εἰς οἶνον καὶ μέθυσμα· καὶ ἔσται, ἐκ τῆς σταγόνος τοῦ λαοῦ τούτου
\vs{12}συναγόμενος συναχθήσεται Ἰακὼβ σὺν πᾶσιν· ἐκδεχόμενος ἐκδέξομαι τοὺς καταλοίπους τοῦ Ἰσραὴλ, ἐπιτοαυτὸ θήσομαι τὴν ἀποστροφὴν αὐτοῦ· ὡς πρόβατα ἐν θλίψει, ὡς ποίμνιον ἐν μέσῳ κοίτης αὐτῶν· ἐξαλοῦνται ἐξ ἀνθρώπων
\vs{13}διὰ τῆς διακοπῆς πρὸ προσώπου αὐτῶν· διέκοψαν, καὶ διῆλθον πύλην, καὶ ἐξῆλθον διʼ αὐτῆς, καὶ ἐξῆλθεν ὁ βασιλεὺς αὐτῶν πρὸ προσώπου αὐτῶν, ὁ δὲ Κύριος ἡγήσεται αὐτῶν.

\ch{3}
Καὶ ἐρεῖ, ἀκούσατε δὴ ταῦτα αἱ ἀρχαὶ οἴκου Ἰακὼβ, καὶ οἱ κατάλοιποι οἴκου Ἰσραήλ· οὐχ ὑμῖν ἐστι τοῦ γνῶναι τὸ κρίμα;
\vs{2}μισοῦντες τὰ καλὰ, καὶ ζητοῦντες τὰ πονηρὰ, ἁρπάζοντες τὰ δέρματα αὐτῶν ἀπʼ αὐτῶν, καὶ τὰς σάρκας αὐτῶν ἀπὸ τῶν ὀστέων αὐτῶν.
\vs{3}Ὃν τρόπον κατέφαγον τὰς σάρκας τοῦ λαοῦ μου, καὶ τὰ δέρματα αὐτῶν ἀπʼ αὐτῶν ἐξέδειραν, καὶ τὰ ὀστέα αὐτῶν συνέθλασαν, καὶ ἐμέλισαν ὡς σάρκας εἰς λέβητα, καὶ ὡς κρέα εἰς χύτραν,
\vs{4}οὕτως κεκράξονται πρὸς τὸν Κύριον, καὶ οὐκ εἰσακούσεται αὐτῶν· καὶ ἀποστρέψει τὸ πρόσωπον αὐτοῦ ἀπʼ αὐτῶν ἐν τῷ καιρῷ ἐκείνῳ, ἀνθʼ ὧν ἐπονηρεύσαντο ἐν τοῖς ἐπιτηδεύμασιν αὐτῶν ἐπʼ αὐτούς.

\vs{5}Τάδε λέγει Κύριος ἐπὶ τοὺς προφήτας τοὺς πλανῶντας τὸν λαόν μου, τοὺς δάκνοντας ἐν τοῖς ὀδοῦσιν αὐτῶν, καὶ κηρύσσοντας εἰρήνην ἐπʼ αὐτὸν, καὶ οὐκ ἐδόθη εἰς τὸ στόμα αὐτῶν, ἤγειραν ἐπʼ αὐτὸν πόλεμον·
\vs{6}διατοῦτο νὺξ ὑμῖν ἔσται ἐξ ὁράσεως, καὶ σκοτία ἔσται ὑμῖν ἐκ μαντείας, καὶ δύσεται ὁ ἥλιος ἐπὶ τοὺς προφήτας, καὶ συσκοτάσει ἐπʼ αὐτοὺς ἡ ἡμέρα·
\vs{7}Καὶ καταισχυνθήσονται οἱ ὁρῶντες τὰ ἐνύπνια, καὶ καταγελασθήσονται οἱ μάντεις, καὶ καταλαλήσουσι κατʼ αὐτῶν πάντες αὐτοὶ, διότι οὐκ ἔσται ὁ ἐπακούων αὐτῶν.
\vs{8}ἐὰν μὴ ἐγὼ ἐμπλήσω ἰσχὺν ἐν πνεύματι Κυρίου καὶ κρίματος καὶ δυναστείας, τοῦ ἀπαγγεῖλαι τῷ Ἰακὼβ ἀσεβείας αὐτοῦ, καὶ τῷ Ἰσραὴλ ἁμαρτίας αὐτοῦ.

\vs{9}Ἀκούσατε δὴ ταῦτα οἱ ἡγούμενοι οἴκου Ἰακὼβ, καὶ οἱ κατάλοιποι οἴκου Ἰσραὴλ, οἱ βδελυσσόμενοι κρίμα, καὶ πάντα τὰ ὀρθὰ διαστρέφοντες,
\vs{10}οἱ οἰκοδομοῦντες Σιὼν ἐν αἵμασι, καὶ Ἱερουσαλὴμ ἐν ἀδικίαις,
\vs{11}οἱ ἡγούμενοι αὐτῆς μετὰ δώρων ἔκρινον, καὶ οἱ ἱερεῖς αὐτῆς μετὰ μισθοῦ ἀπεκρίνοντο, καὶ οἱ προφῆται αὐτῆς μετὰ ἀργυρίου ἐμαντεύοντο, καὶ ἐπὶ τὸν Κύριον ἐπανεπαύοντο, λέγοντες, οὐχὶ ὁ Κύριος ἐν ἡμῖν ἐστιν; οὐ μὴ ἐπέλθῃ ἐφʼ ἡμᾶς κακά.
\vs{12}Διατοῦτο διʼ ὑμᾶς Σιὼν ὡς ἀγρὸς ἀροτριαθήσεται, καὶ Ἱερουσαλὴμ ὡς ὀπωροφυλάκιον ἔσται, καὶ τὸ ὄρος τοῦ οἴκου εἰς ἄλσος δρυμοῦ.

\ch{4}
Καὶ ἔσται ἐπʼ ἐσχάτων τῶν ἡμερῶν ἐμφανὲς τὸ ὄρος Κυρίου, ἕτοιμον ἐπὶ τὰς κορυφὰς τῶν ὀρέων, καὶ μετεωρισθήσεται ὑπεράνω τῶν βουνῶν· καὶ σπεύσουσι πρὸς αὐτὸ λαοὶ,
\vs{2}καὶ πορεύσονται ἔθνη πολλὰ καὶ ἐροῦσι, δεῦτε, ἀναβῶμεν εἰς τὸ ὄρος Κυρίου, καὶ εἰς τὸν οἶκον τοῦ Θεοῦ Ἰακώβ· καὶ δείξουσιν ἡμῖν τὴν ὁδὸν αὐτοῦ, καὶ πορευσόμεθα ἐν ταῖς τρίβοις αὐτοῦ· ὅτι ἐκ Σιὼν ἐξελεύσεται νόμος, καὶ λόγος Κυρίου ἐξ Ἱερουσαλήμ.
\vs{3}Καὶ κρινεῖ ἀναμέσον λαῶν πολλῶν, καὶ ἐξελέγξει ἔθνη ἰσχυρὰ ἕως εἰς μακράν· καὶ κατακόψουσι τὰς ῥομφαίας αὐτῶν εἰς ἄροτρα, καὶ τὰ δόρατα αὐτῶν εἰς δρέπανα, καὶ οὐκέτι μὴ ἀντάρῃ ἔθνος ἐπʼ ἔθνος ῥομφαίαν, καὶ οὐκέτι μὴ μάθωσι πολεμεῖν·
\vs{4}Καὶ ἀναπαύσεται ἕκαστος ὑποκάτω ἀμπέλου αὐτοῦ, καὶ ἕκαστος ὑποκάτω συκῆς αὐτοῦ, καὶ οὐκ ἔσται ὁ ἐκφοβῶν, διότι τὸ στόμα Κυρίου παντοκράτορος ἐλάλησε ταῦτα·
\vs{5}Ὅτι πάντες οἱ λαοὶ πορεύσονται ἕκαστος τὴν ὁδὸν αὐτοῦ, ἡμεῖς δὲ πορευσόμεθα ἐν ὀνόματι Κυρίου Θεοῦ ἡμῶν εἰς τὸν αἰῶνα, καὶ ἐπέκεινα.

\vs{6}Ἐν τῇ ἡμέρᾳ ἐκείνῃ, λέγει Κύριος, συνάξω τὴν συντετριμμένην, καὶ τὴν ἐξωσμένην εἰσδέξομαι, καὶ οὓς ἀπωσάμην.
\vs{7}Καὶ θήσομαι τὴν συντετριμμένην εἰς ὑπόλειμμα, καὶ τὴν ἀπωσμένην εἰς ἔθνος δυνατόν· καὶ βασιλεύσει Κύριος ἐπʼ αὐτοὺς ἐν ὄρει Σιὼν ἀπὸ τοῦ νῦν ἕως εἰς τὸν αἰῶνα.

\vs{8}Καὶ σὺ πύργος ποιμνίου αὐχμώδης, θυγάτηρ Σιὼν, ἐπὶ σὲ ἥξει, καὶ εἰσελεύσεται ἡ ἀρχὴ, ἡ πρώτη βασιλεία ἐκ Βαβυλῶνος τῇ θυγατρὶ Ἱερουσαλήμ.

\vs{9}Καὶ νῦν ἱνατί ἔγνως κακά; μὴ βασιλεὺς οὐκ ἦν σοι; ἢ ἡ βουλή σου ἀπώλετο, ὅτι κατεκράτησάν σου ὠδῖνες ὡς τικτούσης;
\vs{10}Ὤδινε καὶ ἀνδρίζου, καὶ ἔγγιζε θυγάτηρ Σιὼν ὡς τίκτουσα· διότι νῦν ἐξελεύσῃ ἐκ πόλεως, καὶ κατασκηνώσεις ἐν πεδίῳ, καὶ ἥξεις ἕως Βαβυλῶνος· ἐκεῖθεν ῥύσεταί σε, καὶ ἐκεῖθεν λυτρώσεταί σε Κύριος ὁ Θεός σου ἐκ χειρὸς ἐχθρῶν σου.

\vs{11}Καὶ νῦν ἐπισυνήχθησαν ἐπὶ σὲ ἔθνη πολλὰ, λέγοντες, ἐπιχαρούμεθα, καὶ ἐπόψονται ἐπὶ Σιὼν οἱ ὀφθαλμοὶ ἡμῶν.
\vs{12}Αὐτοὶ δὲ οὐκ ἔγνωσαν τὸν λογισμὸν Κυρίου, καὶ οὐ συνῆκαν τὴν βουλὴν αὐτοῦ, ὅτι συνήγαγεν αὐτοὺς ὡς δράγματα ἅλωνος.
\vs{13}Ἀνάστηθι, καὶ ἀλόα αὐτοὺς θυγάτηρ Σιὼν, ὅτι τὰ κέρατά σου θήσομαι σιδηρᾶ, καὶ τὰς ὁπλάς σου θήσομαι χαλκᾶς· καὶ κατατήξεις λαοὺς πολλοὺς, καὶ ἀναθήσεις τῷ Κυρίῳ τὸ πλῆθος αὐτῶν, καὶ τὴν ἰσχὺν αὐτῶν τῷ Κυρίῳ πάσης τῆς γῆς.

\vs{14}Νῦν ἐμφραχθήσεται θυγάτηρ ἐμφραγμῷ, συνοχὴν ἔταξεν ἐφʼ ἡμᾶς, ἐν ῥάβδῳ πατάξουσιν ἐπὶ σιαγόνα τὰς φυλὰς τοῦ Ἰσραήλ.

\ch{5}
Καὶ σὺ Βηθλεὲμ οἶκος Ἐφραθὰ, ὀλιγοστὸς εἶ τοῦ εἶναι ἐν χιλιάσιν Ἰούδα· ἐκ σοῦ μοι ἐξελεύσεται, τοῦ εἶναι εἰς ἄρχοντα τοῦ Ἰσραὴλ, καὶ ἔξοδοι αὐτοῦ ἀπʼ ἀρχῆς ἐξ ἡμερῶν αἰῶνος.

\vs{2}Διατοῦτο δώσει αὐτοὺς ἕως καιροῦ τικτούσης, τέξεται, καὶ οἱ ἐπίλοιποι τῶν ἀδελφῶν αὐτῶν ἐπιστρέψουσιν ἐπὶ τοὺς υἱοὺς Ἰσραήλ.
\vs{3}Καὶ στήσεται καὶ ὄψεται, καὶ ποιμανεῖ τὸ ποίμνιον αὐτοῦ ἐν ἰσχύϊ Κύριος, καὶ ἐν τῇ δόξῃ ὀνόματος Κυρίου Θεοῦ αὐτῶν ὑπάρξουσι, διότι νῦν μεγαλυνθήσονται ἕως ἄκρων τῆς γῆς.

\vs{4}Καὶ ἔσται αὐτῇ εἰρήνη, Ἀσσοὺρ ὅταν ἐπέλθῃ ἐπὶ τὴν γῆν ὑμῶν, καὶ ὅταν ἐπιβῇ ἐπὶ τὴν χώραν ὑμῶν, καὶ ἐπεγερθήσονται ἐπʼ αὐτὸν ἑπτὰ ποιμένες, καὶ ὀκτὼ δήγματα ἀνθρώπων,
\vs{5}καὶ ποιμανοῦσι τὸν Ἀσσοὺρ ἐν ῥομφαίᾳ, καὶ τὴν γῆν τοῦ Νεβρὼδ ἐν τῇ τάφρῳ αὐτῆς· καὶ ῥύσεται ἐκ τοῦ Ἀσσοὺρ ὅταν ἐπέλθῃ ἐπὶ τὴν γῆν ὑμῶν, καὶ ὅταν ἐπιβῇ ἐπὶ τὰ ὅρια ὑμῶν.

\vs{6}Καὶ ἔσται τὸ ὑπόλειμμα τοῦ Ἰακὼβ ἐν τοῖς ἔθνεσιν ἐν μέσῳ λαῶν πολλῶν, ὡς δρόσος παρὰ Κυρίου πίπτουσα, καὶ ὡς ἄρνες ἐπὶ ἄγρωστιν, ὅπως μὴ συναχθῇ μηδεὶς, μηδὲ ὑποστῇ ἐν υἱοῖς ἀνθρώπων.
\vs{7}Καὶ ἔσται τὸ ὑπόλειμμα Ἰακὼβ ἐν τοῖς ἔθνεσιν ἐν μέσῳ λαῶν πολλῶν, ὡς λέων ἐν κτήνεσιν ἐν τῷ δρυμῷ, καὶ ὡς σκύμνος ἐν ποιμνίοις προβάτων, ὃν τρόπον ὅταν διέλθῃ, καὶ διαστείλας ἁρπάσῃ, καὶ μὴ ᾖ ὁ ἐξαιρούμενος.
\vs{8}Ὑψωθήσεται ἡ χείρ σου ἐπὶ τοὺς θλίβοντάς σε, καὶ πάντες οἱ ἐχθροί σου ἐξολοθρευθήσονται.

\vs{9}Καὶ ἔσται ἐν τῇ ἡμέρᾳ ἐκείνῃ, λέγει Κύριος, ἐξολοθρεύσω τοὺς ἵππους ἐκ μέσου σου, καὶ ἀπολῶ τὰ ἅρματά σου,
\vs{10}καὶ ἐξολεθρεύσω τὰς πόλεις τῆς γῆς σου, καὶ ἐξαρῶ πάντα τὰ ὀχυρώματά σου·
\vs{11}καὶ ἐξολοθρεύσω τὰ φάρμακά σου ἐκ τῶν χειρῶν σου, καὶ ἀποφθεγγόμενοι οὐκ ἔσονται ἐν σοί·
\vs{12}καὶ ἐξολοθρεύσω τὰ γλυπτά σου, καὶ τὰς στηλὰς σου ἐκ μέσου σου, καὶ οὐκ ἔτι μὴ προσκυνήσεις τοῖς ἔργοις τῶν χειρῶν σου.
\vs{13}Καὶ ἐκκόψω τὰ ἄλση ἐκ μέσου σου, καὶ ἀφανιῶ τὰς πόλεις σου.
\vs{14}Καὶ ποιήσω ἐν ὀργῇ καὶ ἐν θυμῷ ἐκδίκησιν ἐν τοῖς ἔθνεσιν, ἀνθʼ ὧν οὐκ εἰσήκουσαν.

\ch{6}
Ἀκούσατε δὴ λόγον· Κύριος Κύριος εἶπεν, ἀνάστηθι, κρίθητι πρὸς τὰ ὄρη, καὶ ἀκουσάτωσαν βουνοὶ φωνήν σου.

\vs{2}Ἀκούσατε ὄρη τὴν κρίσιν τοῦ Κυρίου, καὶ αἱ φάραγγες θεμέλια τῆς γῆς, ὅτι κρίσις τῷ Κυρίῳ πρὸς τὸν λαὸν αὐτοῦ, καὶ μετὰ τοῦ Ἰσραὴλ διελεγχθήσεται.
\vs{3}Λαός μου, τί ἐποίησά σοι, ἢ τί ἐλύπησά σε, ἢ τί παρηνώχλησά σοι; ἀποκρίθητί μοι.
\vs{4}Διότι ἀνήγαγόν σε ἐκ γῆς Αἰγύπτου, καὶ ἐξ οἴκου δουλείας ἐλυτρωσάμην σε, καὶ ἐξαπέστειλα πρὸ προσώπου σου τὸν Μωυσῆν, καὶ Ἀαρὼν, καὶ Μαριάμ.

\vs{5}Λαός μου μνήσθητι δὴ, τί ἐβουλεύσατο κατὰ σοῦ Βαλὰκ βασιλεὺς Μωὰβ, καὶ τί ἀπεκρίθη αὐτῷ Βαλαὰμ, υἱὸς τοῦ Βεὼρ ἀπὸ τῶν σχοίνων ἕως τοῦ Γαλγὰλ, ὅπως γνωσθῇ ἡ δικαιοσύνη τοῦ Κυρίου.

\vs{6}Ἐν τίνι καταλάβω τὸν Κύριον, ἀντιλήψομαι Θεοῦ μου ὑψίστου; εἰ καταλήψομαι αὐτὸν ἐν ὁλοκαυτώμασιν, ἐν μόσχοις ἐνιαυσίοις;
\vs{7}Εἰ προσδέξεται Κύριος ἐν χιλιάσι κριῶν; ἢ ἐν μυριάσι χιμάρων πιόνων; εἰ δῶ πρωτότοκά μου ὑπὲρ ἀσεβείας, καρπὸν κοιλίας μου ὑπὲρ ἁμαρτίας ψυχῆς μου;
\vs{8}Εἰ ἀνηγγέλη σοι ἄνθρωπε τί καλόν; ἢ τί Κύριος ἐκζητεῖ παρὰ σοῦ, ἀλλʼ ἢ τοῦ ποιεῖν κρίμα, καὶ ἀγαπᾷν ἔλεον, καὶ ἕτοιμον εἶναι τοῦ πορεύεσθαι μετὰ Κυρίου Θεοῦ σου;

\vs{9}Φωνὴ Κυρίου τῇ πόλει ἐπικληθήσεται, καὶ σώσει φοβουμένους τὸ ὄνομα αὐτοῦ· ἄκουε φυλὴ, καὶ τίς κοσμήσει πόλιν;
\vs{10}Μὴ πῦρ καὶ οἶκος ἀνόμου θησαυρίζων θησαυροὺς ἀνόμους, καὶ μετὰ ὕβρεως ἀδικίας;
\vs{11}Εἰ δικαιωθήσεται ἐν ζυγῷ ἄνομος, καὶ ἐν μαρσίππῳ στάθμια δόλου,
\vs{12}ἐξ ὧν τὸν πλοῦτον αὐτῶν ἀσεβείας ἔπλησαν, καὶ οἱ κατοικοῦντες αὐτὴν ἐλάλουν ψεύδη, καὶ ἡ γλῶσσα αὐτῶν ὑψώθη ἐν τῷ στόματι αὐτῶν;

\vs{13}Καὶ ἐγὼ ἄρξομαι τοῦ πατάξαι σε, ἀφανιῶ σε ἐν ταῖς ἁμαρτίαις σου.
\vs{14}Σὺ φάγεσαι, καὶ οὐ μὴ ἐμπλησθῇς, καὶ συσκοτάσει ἐν σοὶ καὶ ἐκνεύσει, καὶ οὐ μὴ διασωθῇς, καὶ ὅσοι ἂν διασωθῶσιν, εἰς ῥομφαίαν παραδοθήσονται·
\vs{15}Σὺ σπερεῖς, καὶ οὐ μὴ ἀμήσῃς, σὺ πιέσεις ἐλαίαν, καὶ οὐ μὴ ἀλείψῃ ἔλαιον, καὶ οἶνον, καὶ οὐ μὴ πίητε, καὶ ἀφανισθήσεται νόμιμα λαοῦ μου.
\vs{16}Καὶ ἐφύλαξας τὰ δικαιώματα Ζαμβρὶ, καὶ πάντα τὰ ἔργα οἴκου Ἀχαὰβ, καὶ ἐπορεύθητε ἐν ταῖς ὁδοῖς αὐτῶν, ὅπως παραδῶ σε εἰς ἀφανισμὸν, καὶ τοὺς κατοικοῦντας αὐτὴν εἰς συρισμὸν, καὶ ὀνείδη λαῶν λήψεσθε.

\ch{7}
Οἴμοι, ὅτι ἐγενήθην ὡς συνάγων καλάμην ἐν ἀμητῷ, καὶ ὡς ἐπιφυλλίδα ἐν τρυγητῳ, οὐχ ὑπάρχοντος βότρυος τοῦ φαγεῖν τὰ πρωτόγονα. οἴμοι ψυχὴ,
\vs{2}ὅτι ἀπόλωλεν εὐσεβὴς ἀπὸ τῆς γῆς, καὶ κατορθῶν ἐν ἀνθρώποις οὐχ ὑπάρχει· πάντες εἰς αἵματα δικάζονται, ἕκαστος τὸν πλησίον αὐτοῦ ἐκθλίβουσιν ἐκθλιβῇ,
\vs{3}ἐπὶ τὸ κακὸν τὰς χεῖρας αὐτῶν ἑτοιμάζουσιν· ὁ ἄρχων αἰτεῖ, καὶ ὁ κριτὴς εἰρηνικοὺς λόγους ἐλάλησε, καταθύμιον ψυχῆς αὐτοῦ ἐστιν· καὶ ἐξελοῦμαι τὰ ἀγαθὰ αὐτῶν
\vs{4}ὡς σὴς ἐκτρώγων, καὶ βαδίζων ἐπὶ κανόνος ἐν ἡμέρᾳ σκοπιᾶς· οὐαὶ οὐαὶ, αἱ ἐκδικήσεις σου ἥκασι, νῦν ἔσονται κλαυθμοὶ αὐτῶν.
\vs{5}Μὴ καταπιστεύετε ἐν φίλοις, καὶ μὴ ἐλπίζετε ἐπὶ ἡγουμένοις· ἀπὸ τῆς συγκοίτου σου φύλαξαι, τοῦ ἀναθέσθαι τι αὐτῇ.
\vs{6}Διότι υἱὸς ἀτιμάζει πατέρα, θυγάτηρ ἐπαναστήσεται ἐπὶ τὴν μητέρα αὐτῆς, νύμφη ἐπὶ τὴν πενθερὰν αὐτῆς, ἐχθροὶ πάντες ἀνδρὸς οἱ ἐν τῷ οἴκῳ αὐτοῦ.

\vs{7}Ἐγὼ δὲ ἐπὶ τὸν Κύριον ἐπιβλέψομαι, ὑπομενῶ ἐπὶ τῷ Θεῷ τῷ σωτῆρί μου, εἰσακούσεταί μου ὁ Θεός μου.

\vs{8}Μὴ ἐπίχαιρέ μοι ἡ ἐχθρά μου, ὅτι πέπτωκα, καὶ ἀναστήσομαι· διότι ἐὰν καθίσω ἐν τῷ σκότει, Κύριος φωτιεῖ μοι.
\vs{9}Ὀργὴν Κυρίου ὑποίσω, ὅτι ἥμαρτον αὐτῷ, ἕως τοῦ δικαιῶσαι αὐτὸν τὴν δίκην μου· καὶ ποιήσει τὸ κρίμα μου, καὶ ἐξάξει με εἰς τὸ φῶς· ὄψομαι τὴν δικαιοσύνην αὐτοῦ,
\vs{10}καὶ ὄψεται ἡ ἐχθρά μου καὶ περιβαλεῖται αἰσχύνην, ἡ λέγουσα, ροῦ Κύριος ὁ Θεός σου; οἱ ὀφθαλμοί μου ἐπόψονται αὐτὴν, νῦν ἔσται εἰς καταπάτημα ὡς πηλὸς ἐν ταῖς ὁδοῖς.

\vs{11}Ἡμέρα ἀλοιφῆς πλίνθου, ἐξάλειψίς σου ἡ ἡμέρα ἐκείνη, καὶ ἀποτρίψεται νόμιμά σου ἡ ἡμέρα ἐκείνη.
\vs{12}Καὶ αἱ πόλεις σου ἥξουσιν εἰς ὁμαλισμὸν, καὶ εἰς διαμερισμὸν Ἀσσυρίων, καὶ αἱ πόλεις σου αἱ ὀχυραὶ εἰς διαμερισμὸν ἀπὸ Τύρου ἕως τοῦ ποταμοῦ, καὶ ἀπὸ θαλάσσης ἕως θαλάσσης, καὶ ἀπὸ ὄρους ἕως τοῦ ὄρους.
\vs{13}Καὶ ἔσται ἡ γῆ εἰς ἀφανισμὸν σὺν τοῖς κατοικοῦσιν αὐτὴν, ἀπὸ καρπῶν ἐπιτηδευμάτων αὐτῶν.

\vs{14}Ποίμαινε λαόν σου ἐν ῥάβδῳ σου, πρόβατα κληρονομίας σου, κατασκηνοῦντας καθʼ ἑαυτοὺς δρυμὸν ἐν μέσῳ τοῦ Καρμήλου· νεμήσονται τὴν Βασανίτιν, καὶ τὴν Γαλααδίτιν καθὼς αἱ ἡμέραι τοῦ αἰῶνος.

\vs{15}Καὶ κατὰ τὰς ἡμέρας ἐξοδίας σου ἐξ Αἰγύπτου, ὄψεσθε θαυμαστά.
\vs{16}Ὄψονται ἔθνη καὶ καταισχυνθήσονται, καὶ ἐκ πάσης τῆς ἰσχύος αὐτῶν, ἐπιθήσουσι χεῖρας ἐπὶ τὸ στόμα αὐτῶν, τὰ ὦτα αὐτῶν ἀποκωφωθήσεται,
\vs{17}λείξουσι χοῦν ὡς ὄφεις σύροντες γῆν, συγχυθήσονται ἐν συγκλεισμῷ αὐτῶν· ἐπὶ τῷ Κυρίῳ Θεῷ ἡμῶν ἐκστήσονται, καὶ φοβηθήσονται ἀπὸ σοῦ.

\vs{18}Τίς Θεὸς ὥσπερ σὺ; ἐξαίρων ἀνομίας, καὶ ὑπερβαίνων ἀσεβείας τοῖς καταλοίποις τῆς κληρονομίας αὐτοῦ; καὶ οὐ συνέσχεν εἰς μαρτύριον ὀργὴν αὐτοῦ, ὅτι θελητὴς ἐλέους ἐστίν.
\vs{19}Ἐπιστρέψει καὶ οἰκτειρήσει ἡμᾶς, καταδύσει τὰς ἀδικίας ἡμῶν, καὶ ἀποῤῥιφήσονται εἰς τὰ βάθη τῆς θαλάσσης πάσας τὰς ἁμαρτίας ἡμῶν.
\vs{20}Δώσει εἰς ἀλήθειαν τῷ Ἰακὼβ, ἔλεον τῷ Ἁβραάμ, καθότι ὤμοσας τοῖς πατράσιν ἡμῶν, κατὰ τὰς ἡμέρας τὰς ἔμπροσθεν.


\def\book{ΝΑΟΥΜ}
\biblebook{ΝΑΟΥΜ}


\lettrine[lines=2, loversize=0.2, nindent=0em, findent=.25em]{\textcolor{bookheadingcolor}{Λ}}{ΗΜΜΑ} Νινευὴ, βιβλίον ὁράσεως Ναοὺμ τοῦ Ἐλκεσαίου.

\postdropcapindent\vs{2}Θεὸς ζηλωτὴς, καὶ ἐκδικῶν Κύριος, ἐκδικῶν Κύριος μετὰ θυμοῦ, ἐκδικῶν Κύριος τοὺς ὑπεναντίους αὐτοῦ, καὶ ἐξαίρων αὐτὸς τοὺς ἐχθροὺς αὐτοῦ.
\vs{3}Κύριος μακρόθυμος, καὶ μεγάλη ἡ ἰσχὺς αὐτοῦ, καὶ ἀθῶον οὐκ ἀθωώσει Κύριος· ἐν συντελείᾳ, καὶ ἐν συσσεισμῷ ἡ ὁδὸς αὐτοῦ, καὶ νεφέλαι κονιορτὸς ποδῶν αὐτοῦ.
\vs{4}Ἀπειλῶν θαλάσσῃ, καὶ ξηραίνων αὐτὴν, καὶ πάντας τοὺς ποταμοὺς ἐξερημῶν· ὠλιγώθη ἡ Βασανίτις, καὶ ὁ Κάρμηλος, καὶ τὰ ἐξανθοῦντα τοῦ Λιβάνου ἐξέλιπε.
\vs{5}Τὰ ὄρη ἐσείσθησαν ἀπʼ αὐτοῦ, καὶ οἱ βουνοὶ ἐσαλεύθησαν· καὶ ἀνεστάλη ἡ γῆ ἀπὸ προσώπου αὐτοῦ, ἡ σύμπασα καὶ πάντες οἱ κατοικοῦντες ἐν αὐτῇ.
\vs{6}Ἀπὸ προσώπου ὀργῆς αὐτοῦ τίς ὑποστήσεται; καὶ τίς ἀντιστήσεται ἐν ὀργῇ θυμοῦ αὐτοῦ; ὁ θυμὸς αὐτοῦ τήκει ἀρχὰς, καὶ αἱ πέτραι διεθρύβησαν ἀπʼ αὐτοῦ.

\vs{7}Χρηστὸς Κύριος τοῖς ὑπομένουσιν αὐτὸν ἐν ἡμέρᾳ θλίψεως, καὶ γινώσκων τοὺς εὐλαβουμένους αὐτόν.
\vs{8}Καὶ ἐν κατακλυσμῷ πορείας συντέλειαν ποιήσεται, τοὺς ἐπεγειρομένους καὶ τοὺς ἐχθροὺς αὐτοῦ διώξεται σκότος.
\vs{9}Τί λογίζεσθε ἐπὶ τὸν Κύριον; συντέλειαν αὐτὸς ποιήσεται, οὐκ ἐκδικήσει δὶς ἐπιτοαυτὸ ἐν θλίψει.
\vs{10}Ὅτι ἕως θεμελίου αὐτοῦ χερσωθήσεται· καὶ ὡς σμῖλαξ περιπλεκομένη βρωθήσεται, καὶ ὡς καλάμη ξηρασίας μεστή.

\vs{11}Ἐκ σοῦ ἐξελεύσεται λογισμὸς κατὰ τοῦ Κυρίου, πονηρὰ βουλευόμενος ἐναντία.

\vs{12}Τάδε λέγει Κύριος κατάρχων ὑδάτων πολλῶν, καὶ οὕτως διασταλήσονται, καὶ ἡ ἀκοή σου οὐκ ἐνακουσθήσεται ἔτι.
\vs{13}Καὶ νῦν συντρίψω τὴν ῥάβδον αὐτοῦ ἀπὸ σοῦ, καὶ τοὺς δεσμοὺς διαῤῥήξω.

\vs{14}Καὶ ἐντελεῖται περὶ σοῦ Κύριος, οὐ σπαρήσεται ἐκ τοῦ ὀνόματός σου ἔτι· ἐξ οἴκου θεοῦ σου ἐξολοθρεύσω τὰ γλυπτὰ, καὶ χωνευτὰ, θήσομαι ταφήν σου, ὅτι ταχεῖς.

\ch{2}
Ἰδοὺ ἐπὶ τὰ ὄρη οἱ πόδες εὐαγγελιζομένου, καὶ ἀπαγγέλλοντος εἰρήνην· ἑόρταζε Ἰούδα τὰς ἑορτάς σου, ἀπόδος τὰς εὐχάς σου, διότι οὐ μὴ προσθήσωσιν ἔτι τοῦ διελθεῖν διὰ σοῦ εἰς παλαίωσιν.

\vs{2}Συντετέλεσται, ἐξῇρται· ἀνέβη ἐμφυσῶν εἰς πρόσωπόν σου, ἐξαιρούμενος ἐκ θλίψεως· σκόπευσον ὁδὸν, κράτησον ὀσφύος, ἄνδρισαι τῇ ἰσχύϊ σφόδρα.

\vs{3}Διότι ἀπέστρεψε Κύριος τὴν ὕβριν Ἰακὼβ, καθὼς ὕβριν τοῦ Ἰσραὴλ, διότι ἐκτινάσσοντες ἐξετίναξαν αὐτοὺς, καὶ τὰ κλήματα αὐτῶν διέφθειραν.
\vs{4}Ὅπλα δυναστείας αὐτῶν ἐξ ἀνθρώπων, ἄνδρας δυνατοὺς ἐμπαίζοντας ἐν πυρί· αἱ ἡνίαι τῶν ἁρμάτων αὐτῶν ἐν ἡμέρᾳ ἑτοιμασίας αὐτοῦ, καὶ οἱ ἱππεῖς θορυβηθήσονται
\vs{5}ἐν ταῖς ὁδοῖς, καὶ συγχυθήσονται τὰ ἅρματα, καὶ συμπλακήσονται ἐν ταῖς πλατείαις· ἡ ὅρασις αὐτῶν ὡς λαμπάδες πυρὸς, καὶ ὡς ἀστραπαὶ διατρέχουσαι.

\vs{6}Καὶ μνησθήσονται οἱ μεγιστᾶνες αὐτῶν, καὶ φεύξονται ἡμέρας, καὶ ἀσθενήσουσιν ἐν τῇ πορείᾳ αὐτῶν, καὶ σπεύσουσιν ἐπὶ τὰ τείχη αὐτῆς, καὶ ἑτοιμάσουσι τὰς προφυλακὰς αὐτῶν.
\vs{7}Πύλαι τῶν πόλεων διηνοίχθησαν, καὶ τὰ βασίλεια διέπεσε,
\vs{8}καὶ ἡ ὑπόστασις ἀπεκαλύφθη· καὶ αὕτη ἀνέβαινε, καὶ αἱ δοῦλαι αὐτῆς ἤγοντο, καθὼς περιστεραὶ φθεγγόμεναι ἐν καρδίαις αὐτῶν.
\vs{9}Καὶ Νινευὴ ὡς κολυμβήθρα ὕδατος τὰ ὕδατα αὐτῆς, καὶ αὐτοὶ φεύγοντες οὐκ ἔστησαν, καὶ οὐκ ἦν ὁ ἐπιβλέπων.

\vs{10}Διήρπαζον τὸ ἀργύριον, διήρπαζον τὸ χρυσίον, καὶ οὐκ ἦν πέρας τοῦ κόσμου αὐτῆς· βεβάρυνται ἐπὶ πάντα τὰ σκεὺη τὰ ἐπιθυμητὰ αὐτῆς.
\vs{11}Ἐκτιναγμὸς, καὶ ἀνατιναγμὸς, καὶ ἐκβρασμὸς, καὶ καρδίας θραυσμὸς, καὶ ὑπόλυσις γονάτων, καὶ ὠδίνες ἐπὶ πᾶσαν ὀσφύν· καὶ τὸ πρόσωπον πάντων ὡς πρόσκαυμα χύτρας.

\vs{12}Ποῦ ἐστι τὸ κατοικητήριον τῶν λεόντων, καὶ ἡ νομὴ ἡ οὖσα τοῖς σκύμνοις; ποῦ ἐπορεύθη λέων, τοῦ εἰσελθεῖν ἐκεῖ σκύμνον λέοντος, καὶ οὐκ ἦν ὁ ἐκφοβῶν;
\vs{13}Λέων ἥρπασε τὰ ἱκανὰ τοῖς σκύμνοις αὐτοῦ, καὶ ἀπέπνιξε τοῖς λέουσιν αὐτοῦ, καὶ ἔπλησε θήρας νοσσιὰν αὐτοῦ, καὶ τὸ κατοικητήριον αὐτοῦ ἁρπαγῆς.

\vs{14}Ἰδοὺ ἐγὼ ἐπὶ σὲ, λέγει Κύριος παντοκράτωρ, καὶ ἐκκαύσω ἐν καπνῷ πλῆθός σου, καὶ τοὺς λέοντάς σου καταφάγεται ῥομφαία, καὶ ἐξολοθρεύσω ἐκ τῆς γῆς τὴν θήραν σου, καὶ οὐ μὴ ἀκουσθῇ οὐκέτι τὰ ἔργα σου.

\ch{3}
Ὦ πόλις αἱμάτων, ὅλη ψευδὴς, ἀδικίας πλήρης, οὐ ψηλαφηθήσεται θήρα.
\vs{2}Φωνὴ μαστίγων, καὶ φωνὴ σεισμοῦ τροχῶν, καὶ ἵππου διώκοντος, καὶ ἅρματος ἀναβράσσοντος,
\vs{3}καὶ ἱππέως ἀναβαίνοντος, καὶ στιλβούσης ῥομφαίας, καὶ ἐξαστραπτόντων ὅπλων, καὶ πλήθους τραυματιῶν, καὶ βαρείας πτώσεως, καὶ οὐκ ἦν πέρας τοῖς ἔθνεσιν αὐτῆς· καὶ ἀσθενήσουσιν ἐν τοῖς σώμασιν αὐτῶν ἀπὸ πλήθους πορνείας·
\vs{4}πόρνη καλὴ, καὶ ἐπίχαρις ἡγουμένη φαρμάκων, ἡ πωλοῦσα ἔθνη ἐν τῇ πορνείᾳ αὐτῆς, καὶ λαοὺς ἐν τοῖς φαρμάκοις αὐτῆς.

\vs{5}Ἰδοὺ ἐγὼ ἐπὶ σὲ, λέγει Κύριος ὁ Θεὸς ὁ παντοκράτωρ, καὶ ἀποκαλύψω τὰ ὀπίσω σου ἐπὶ τὸ πρόσωπόν σου, καὶ δείξω ἔθνεσι τὴν αἰσχύνην σου, καὶ βασιλείαις τὴν ἀτιμίαν σου.
\vs{6}Καὶ ἐπιῤῥίψω ἐπὶ σὲ βδελυγμὸν κατὰ τὰς ἀκαθαρσίας σου, καὶ θήσομαί σε εἰς παράδειγμα.
\vs{7}Καὶ ἔσται, πᾶς ὁ ὁρῶν σε καταβήσεται ἀπὸ σοῦ, καὶ ἐρεῖ, δειλαία Νινευή· τίς στενάξει αὐτήν; πόθεν ζητήσω παράκλησιν αὐτῇ;

\vs{8}Ἑτοιμασὰι μερίδα, ἁρμόσαι χορδὴν, ἑτοιμάσαι μερίδα Ἀμμών· ἡ κατοικοῦσα ἐν ποταμοῖς, ὕδωρ κύκλῳ αὐτῆς, ἧς ἡ ἀρχὴ θάλασσα, καὶ ὕδωρ τὰ τείχη αὐτῆς,
\vs{9}καὶ Αἰθιοπία ἰσχὺς αὐτῆς, καὶ Αἴγυπτος· καὶ οὐκ ἔστη πέρας τῆς φυγῆς· καὶ Λίβυες ἐγένοντο βοηθοὶ αὐτῆς.
\vs{10}Καὶ αὐτὴ εἰς μετοικεσίαν πορεύσεται αἰχμάλωτος καὶ τὰ νήπια αὐτῆς ἐδαφιοῦσιν ἐπʼ ἀρχὰς πασῶν τῶν ὁδῶν αὐτῆς, καὶ ἐπὶ πάντα τὰ ἔνδοξα αὐτῆς βαλοῦσι κλήρους· καὶ πάντες οἱ μεγιστᾶνες αὐτῆς δεθήσονται χειροπέδαις.
\vs{11}Καὶ σὺ μεθυσθήσῃ, καὶ ἔσῃ ὑπερεωραμένη, καὶ σὺ ζητήσεις σεαυτῇ στάσιν ἐξ ἐχθρῶν.
\vs{12}Πάντα τὰ ὀχυρώματά σου συκαῖ σκοποὺς ἔχουσαι· ἐὰν σαλευθῶσι, πεσοῦνται εἰς στόμα ἔσθοντος.
\vs{13}Ἰδοὺ ὁ λαός σου ὡς γυναῖκες ἐν σοὶ, τοῖς ἐχθροῖς σου ἀνοιγόμεναι ἀνοιχθήσονται πύλαι τῆς γῆς σου, καταφάγεται πῦρ τοὺς μοχλούς σου.

\vs{14}Ὕδωρ περιοχῆς ἐπίσπασαι σεαυτῇ, καὶ κατακράτησον τῶν ὀχυρωμάτων σου· ἔμβηθι εἰς πηλὸν, καὶ συμπατήθητι ἐν ἀχύροις, κατακράτησον ὑπὲρ πλίνθον.
\vs{15}Ἐκεῖ καταφάγεταί σε πῦρ, ἐξολοθρεύσει σε ῥομφαία, καταφάγεταί σε ὡς ἀκρὶς, καὶ βαρυνθήσῃ ὡς βροῦχος.
\vs{16}Ἐπλήθυνας τὰς ἐμπορίας σου ὑπὲρ τὰ ἄστρα τοῦ οὐρανοῦ· βροῦχος ὥρμησε, καὶ ἐξεπετάσθη.
\vs{17}Ἐξήλατο ὡς ἀττέλεβος ὁ σύμμικτός σου, ὡς ἀκρὶς ἐπιβεβηκυῖα ἐπὶ φραγμὸν ἐν ἡμέρᾳ πάγους· ὁ ἥλιος ἀνέτειλε, καὶ ἀφήλατο, καὶ οὐκ ἔγνω τὸν τόπον αὐτῆς· οὐαὶ αὐτοῖς.

\vs{18}Ἐνύσταξαν οἱ ποιμένες σου, βασιλεὺς Ἀσσύριος ἐκοίμισε τοὺς δυνάστας σου, ἀπῇρεν ὁ λαός σου ἐπὶ τὰ ὄρη, καὶ οὐκ ἦν ὁ ἐκδεχόμενος.

\vs{19}Οὐκ ἔστιν ἴασις τῇ συντριβῇ σου, ἐφλέγμανεν ἡ πληγή σου, πάντες οἱ ἀκούοντες τὴν ἀγγελίαν σου κροτήσουσι χεῖρας ἐπὶ σέ· διότι ἐπὶ τίνα οὐκ ἐπῆλθεν ἡ κακία σου διαπαντός;


\def\book{ΑΜΒΑΚΟΥΜ}
\biblebook{ΑΜΒΑΚΟΥΜ}


\lettrine[lines=2, loversize=0.2, nindent=0em, findent=.25em]{\textcolor{bookheadingcolor}{Τ}}{Ο} λῆμμα ὃ εἶδεν Ἀμβακοὺμ ὁ προφήτης.

\postdropcapindent\vs{2}Ἕως τίνος Κύριε κεκράξομαι, καὶ οὐ μὴ εἰσακουσεις; βοήσομαι πρὸς σὲ ἀδικούμενος, καὶ οὐ σώσεις;
\vs{3}Ἱνατί ἔδειξάς μοι κόπους καὶ πόνους ἐπιβλέπειν, ταλαιπωρίαν καὶ ἀσέβειαν; ἐξεναντίας μου γέγονε κρίσις, καὶ ὁ κριτὴς λαμβάνει·
\vs{4}Διατοῦτο διεσκέδασται νόμος, καὶ οὐ διεξάγεται εἰς τέλος κρίμα, ὅτι ἀσεβὴς καταδυναστεύει τὸν δίκαιον, ἕνεκεν τούτου ἐξελεύσεται τὸ κρίμα διεστραμμένον.

\vs{5}Ἴδετε οἱ καταφρονηταὶ, καὶ ἐπιβλέψατε, καὶ θαυμάσατε θαυμάσια, καὶ ἀφανίσθητε· διότι ἔργον ἐγὼ ἐργάζομαι ἐν ταῖς ἡμέραις ὑμῶν, ὃ οὐ μὴ πιστεύσητε, ἐάν τις ἐκδιηγῆται.
\vs{6}Διότι ἰδοὺ ἐγὼ ἐξεγείρω τοὺς Χαλδαίους, τὸ ἔνθος τὸ πικρὸν, καὶ τὸ ταχινὸν, τὸ πορευόμενον ἐπὶ τὰ πλάτη τῆς γῆς, τοῦ κατακληρονομῆσαι σκηνώματα οὐκ αὐτοῦ.
\vs{7}Φοβερὸς καὶ ἐπιφανής ἐστιν, ἐξ αὐτοῦ τὸ κρίμα αὐτοῦ ἔσται, καὶ τὸ λῆμμα αὐτοῦ ἐξ αὐτοῦ ἐξελεύσεται.
\vs{8}Καὶ ἐξαλοῦνται ὑπὲρ παρδάλεις οἱ ἵπποι αὐτοῦ, καὶ ὀξύτεροι ὑπὲρ τοὺς λύκους τῆς Ἀραβίας· καὶ ἐξιππάσονται οἱ ἱππεῖς αὐτοῦ, καὶ ὁρμήσουσι μακρόθεν, καὶ πετασθήσονται ὡς ἀετὸς πρόθυμος εἰς τὸ φαγεῖν.
\vs{9}Συντέλεια εἰς ἀσεβεῖς ἥξει, ἀνθεστηκότας προσώποις αὐτῶν ἐξεναντίας, καὶ συνάξει ὡς ἄμμον αἰχμαλωσίαν·
\vs{10}καὶ αὐτὸς ἐν βασιλεῦσιν ἐντρυφήσει, καὶ τύραννοι παίγνια αὐτοῦ, καὶ αὐτὸς εἰς πᾶν ὀχύρωμα ἐμπαίξεται, καὶ βαλεῖ χῶμα, καὶ κρατήσει αὐτοῦ·
\vs{11}τότε μεταβαλεῖ τὸ πνεῦμα, καὶ διελεύσεται, καὶ ἐξιλάσεται· αὕτη ἡ ἰσχὺς τῷ θεῷ μου.

\vs{12}Οὐχὶ σὺ ἀπʼ ἀρχῆς Κύριε ὁ Θεὸς ὁ ἅγιός μου; καὶ οὐ μὴ ἀποθάνωμεν· Κύριε εἰς κρίμα τέταχας αὐτὸ, καὶ ἔπλασέ με τοῦ ἐλέγχειν παιδείαν αὐτοῦ.
\vs{13}Καθαρὸς ὀφθαλμὸς τοῦ μὴ ὁρᾷν πονηρὰ, καὶ ἐπιβλέπειν ἐπὶ πόνους ὀδύνης· ἱνατί ἐπιβλέπεις ἐπὶ καταφρονοῦντας; παρασιωπήσῃ ἐν τῷ καταπίνειν ἀσεβῆ τὸν δίκαιον;
\vs{14}Καὶ ποιήσεις τοὺς ἀνθρώπους ὡς τοὺς ἰχθύας τῆς θαλάσσης, καὶ ὡς τὰ ἑρπετὰ τὰ οὐκ ἔχοντα ἡγούμενον;
\vs{15}Συντέλειαν ἐν ἀγκίστρῳ ἀνέσπασεν, καὶ εἵλκυσεν αὐτὸν ἐν ἀμφιβλήστρῳ, καὶ συνήγαγεν αὐτὸν ἐν ταῖς σαγήναις αὐτοῦ· ἕνεκεν τούτου εὐφρανθήσεται καὶ χαρήσεται ἡ καρδία αὐτοῦ.
\vs{16}Ἕνεκεν τούτου θύσει τῇ σαγήνῃ αὐτοῦ, καὶ θυμιάσει τῷ ἀμφιβλήστρῳ αὐτοῦ, ὅτι ἐν αὐτοῖς ἐλίπανε μερίδα αὐτοῦ, καὶ τὰ βρώματα αὐτοῦ ἐκλεκτά.
\vs{17}Διατοῦτο ἀμφιβαλεῖ τὸ ἀμφίβληστρον αὐτοῦ, καὶ διαπαντὸς ἀποκτέννειν ἔθνη οὐ φείσεται.

\ch{2}
Ἐπὶ τῆς φυλακῆς μου στήσομαι, καὶ ἐπιβήσομαι ἐπὶ πέτραν, καὶ ἀποσκοπεύσω τοῦ ἰδεῖν τί λαλήσει ἐν ἐμοὶ, καὶ τί ἀποκριθῶ ἐπὶ τὸν ἔλεγχόν μου.

\vs{2}Καὶ ἀπεκρίθη πρὸς μὲ Κύριος, καὶ εἶπε, γράψον ὅρασιν, καὶ σαφῶς εἰς πυξίον, ὅπως διώκῃ ὁ ἀναγινώσκων αὐτά.
\vs{3}Διότι ἔτι ὅρασις εἰς καιρὸν, καὶ ἀνατελεῖ εἰς πέρας, καὶ οὐκ εἰς κενόν· ἐὰν ὑστερήσῃ, ὑπὸμεινον αὐτὸν, ὅτι ἐρχόμενος ἥξει, καὶ οὐ μὴ χρονίσῃ.

\vs{4}Ἐὰν ὑποστείληται, οὐκ εὐδοκεῖ ἡ ψυχή μου ἐν αὐτῷ· ὁ δὲ δίκαιος ἐκ πίστεώς μου ζήσεται.
\vs{5}Ὁ δὲ κατοιόμενος, καὶ καταφρονητὴς, ἀνὴρ ἀλαζὼν, οὐθὲν μὴ περάνῃ· ὃς ἐπλάτυνε καθὼς ᾅδης τὴν ψυχὴν αὐτοῦ, καὶ οὗτος ὡς θάνατος οὐκ ἐμπιπλάμενος, καὶ ἐπισυνάξει ἐπʼ αὐτὸν πάντα τὰ ἔθνη, καὶ εἰσδέξεται πρὸς αὐτὸν πάντας τοὺς λαούς.
\vs{6}Οὐχὶ ταῦτα πάντα κατʼ αὐτοῦ παραβολὴν λήψονται, καὶ πρόβλημα εἰς διήγησιν αὐτοῦ; καὶ ἐροῦσιν, οὐαὶ ὁ πληθύνων ἑαυτῷ τὰ οὐκ ὄντα αὐτοῦ ἕως τίνος, καὶ βαρύνων τὸν κλοιὸν αὐτοῦ στιβαρῶς.
\vs{7}Ὅτι ἐξαίφνης ἀναστήσονται δάκνοντες αὐτὸν, καὶ ἐκνήψουσιν οἱ ἐπίβουλοί σου, καὶ ἔσῃ εἰς διαρπαγὴν αὐτοῖς,
\vs{8}διότι ἐσκύλευσας ἔθνη πολλὰ, σκυλεύσουσι πάντες οἱ ὑπολελειμμένοι λαοὶ, διʼ αἵματα ἀνθρώπων, καὶ ἀσεβείας γῆς καὶ πόλεως, καὶ πάντων τῶν κατοικούντων αὐτήν.

\vs{9}Ὢ ὁ πλεονεκτῶν πλεονεξίαν κακὴν τῷ οἴκῳ αὐτοῦ, τοῦ τάξαι εἰς ὕψος νοσσιὰν αὐτοῦ, τοῦ ἐκσπασθῆναι ἐκ χειρὸς κακῶν.
\vs{10}Ἐβουλεύσω αἰσχύνην τῷ οἴκῳ σου, συνεπέρανας πολλοὺς λαοὺς, καὶ ἐξήμαρτεν ἡ ψυχή σου.
\vs{11}Διότι λίθος ἐκ τοίχου βοήσεται, καὶ κάνθαρος ἐκ ξύλου φθέγξεται αὐτά.

\vs{12}Οὐαὶ ὁ οἰκοδομῶν πόλιν ἐν αἵμασι, καὶ ἑτοιμάζων πόλιν ἐν ἀδικίαις.
\vs{13}Οὐ ταῦτά ἐστι παρὰ Κυρίου παντοκράτορος; καὶ ἐξέλιπον λαοὶ ἱκανοὶ ἐν πυρὶ, καὶ ἔθνη πολλὰ ὠλιγοψύχησαν.
\vs{14}Ὅτι ἐμπλησθήσεται ἡ γῆ τοῦ γνῶναι τὴν δόξαν Κυρίου, ὡς ὕδωρ κατακαλύψει αὐτούς.

\vs{15}Ὢ ὁ ποτίζων τὸν πλησίον αὐτοῦ ἀνατροπῇ θολερᾷ, καὶ μεθύσκων ὅπως ἐπιβλέπῃ ἐπὶ τὰ σπήλαια αὐτῶν.
\vs{16}Πλησμονὴν ἀτιμίας ἐκ δόξης πίε καὶ σύ· καρδία σαλεύθητι, καὶ σείσθητι· ἐκύκλωσεν ἐπὶ σὲ ποτήριον δεξιᾶς Κυρίου, καὶ συνήχθη ἀτιμία ἐπὶ τὴν δόξαν σου.
\vs{17}Διότι ἀσέβεια τοῦ Λιβάνου καλύψει σε, καὶ ταλαιπωρία θηρίων πτοήσει σε, διʼ αἵματα ἀνθρώπων, καὶ ἀσεβείας γῆς καὶ πόλεως, καὶ πάντων τῶν κατοικούντων αὐτήν.

\vs{18}Τί ὠφελεῖ γλυπτὸν, ὅτι ἔγλυψαν αὐτό; ἔπλασεν αὐτὸ χώνευμα, φαντασίαν ψευδῆ, ὅτι πέποιθεν ὁ πλάσας ἐπὶ τὸ πλάσμα αὐτοῦ, τοῦ ποιῆσαι εἴδωλα κωφά.

\vs{19}Οὐαὶ ὁ λέγων τῷ ξύλῳ, ἔκνηψον, ἐξεγέρθητι· καὶ τῷ λίθῳ, ὑψώθητι· καὶ αὐτό ἐστι φαντασία· τοῦτο δέ ἐστιν ἔλασμα χρυσίου καὶ ἀργυρίου, καὶ πᾶν πνεῦμα οὐκ ἔστιν ἐν αὐτῷ.
\vs{20}Ὁ δὲ Κύριος ἐν ναῷ ἁγίῳ αὐτοῦ· εὐλαβείσθω ἀπὸ προσώπου αὐτοῦ πᾶσα ἡ γῆ.

\ch{3}
ΠΡΟΣΕΥΧΗ ἈΜΒΑΚΟΥΜ ΤΟΥ ΠΡΟΦΗΤΟΥ, ΜΕΤΑ Ὠ̣ΔΗΣ.

\vs{2}Κύριε εἰσακήκοα τὴν ἀκοήν σου, καὶ ἐφοβήθην· κατενόησα τὰ ἔργα σου, καὶ ἐξέστην· ἐν μέσῳ δύο ζώων γνωσθήσῃ, ἐν τῷ ἐγγίζειν τὰ ἔτη ἐπιγνωσθήσῃ· ἐν τῷ παρεῖναι τὸν καιρὸν ἀναδειχθήσῃ· ἐν τῷ ταραχθῆναι τὴν ψυχήν μου, ἐν ὀργῇ ἐλέους μνησθήσῃ.

\vs{3}Ὁ Θεὸς ἐκ Θαιμὰν ἥξει, καὶ ὁ ἅγιος ἐξ ὄρους Φαρὰν κατασκίου δασέος· διάψαλμα· ἐκάλυψεν οὐρανοὺς ἡ ἀρετὴ αὐτοῦ, καὶ αἰνέσεως αὐτοῦ πλήρης ἡ γῆ.
\vs{4}Καὶ φέγγος αὐτοῦ ὡς φῶς ἔσται· κέρατα ἐν χερσὶν αὐτοῦ, καὶ ἔθετο ἀγάπησιν κραταιὰν ἰσχύος αὐτοῦ.
\vs{5}Πρὸ προσώπου αὐτοῦ πορεύσεται λόγος, καὶ ἐξελεύσεται εἰς πεδία· κατὰ πόδας αὐτοῦ
\vs{6}ἔστη, καὶ ἐσαλεύθη ἡ γῆ· ἐπέβλεψε, καὶ διετάκη ἔθνη· διεθρύβη τὰ ὄρη βίᾳ, ἐτάκησαν βουνοὶ αἰώνιοι πορείας αἰωνίας αὐτοῦ.
\vs{7}Ἀντὶ κόπων εἶδον σκηνώματα Αἰθιόπων, πτοηθήσονται καὶ αἱ σκηναὶ γῆς Μαδιάμ.

\vs{8}Μὴ ἐν ποταμοῖς ὠργίσθης Κύριε; ἢ ἐν ποταμοῖς ὁ θυμός σου; ἢ ἐν θαλάσσῃ τὸ ὅρμημά σου; ὅτι ἐπιβήσῃ ἐπὶ τοὺς ἵππους σου, καὶ ἡ ἱππασία σου σωτηρία.
\vs{9}Ἐντείνων ἐνέτεινας τόξον σου ἐπὶ σκῆπτρα, λέγει Κύριος· διάψαλμα· ποταμῶν ῥαγήσεται γῆ.
\vs{10}Ὄψονταί σε, καὶ ὠδινήσουσι λαοὶ, σκορπίζων ὕδατα πορείας· ἔδωκεν ἡ ἄβυσσος φωνὴν αὐτῆς, ὕψος φαντασίας αὐτῆς.
\vs{11}Ἐπῄρθη ὁ ἥλιος, καὶ ἡ σελήνη ἔστη ἐν τῇ τάξει αὐτῆς· εἰς φῶς βολίδες σου πορεύσονται, εἰς φέγγος ἀστραπῆς ὅπλων σου.
\vs{12}Ἐν ἀπειλῇ ὀλιγώσεις γῆν, καὶ ἐν θυμῷ κατάξεις ἔθνη,
\vs{13}ἐξῆλθες εἰς σωτηρίαν λαοῦ σου, τοῦ σῶσαι τὸν χριστόν σου· βαλεῖς εἰς κεφαλὰς ἀνόμων θάνατον, ἐξήγειρας δεσμοὺς ἕως τραχήλου· διάψαλμα.
\vs{14}Διέκοψας ἐν ἐκστάσει κεφαλὰς δυναστῶν, σεισθήσονται ἐν αὐτῇ· διανοίξουσι χαλινοὺς αὐτῶν, ὡς ἔσθων πτωχὸς λάθρα.
\vs{15}Καὶ ἐπιβιβᾷς εἰς θάλασσαν τοὺς ἵππους σου, ταράσσοντας ὕδωρ πολύ.

\vs{16}Ἐφυλαξάμην, καὶ ἐπτοήθη ἡ κοιλία μου ἀπὸ φωνῆς προσευχῆς χειλέων μου, καὶ εἰσῆλθε τρόμος εἰς τὰ ὀστᾶ μου, καὶ ὑποκάτωθέν μου ἐταράχθη ἡ ἕξις μου· ἀναπαύσομαι ἐν ἡμέρᾳ θλίψεως, τοῦ ἀναβῆναι εἰς λαὸν παροικίας μου.

\vs{17}Διότι συκῆ οὐ καρποφορήσει, καὶ οὐκ ἔσται γεννήματα ἐν ταῖς ἀμπέλοις· ψεύσεται ἔργον ἐλαίας, καὶ τὰ πεδία οὐ ποιήσει βρῶσιν· ἐξέλιπεν ἀπὸ βρώσεως πρόβατα, καὶ οὐχ ὑπάρχουσι βόες ἐπὶ φάτναις·
\vs{18}ἐγὼ δὲ ἐν τῷ Κυρίῳ ἀγαλλιάσομαι, χαρήσομαι ἐπὶ τῷ Θεῷ τῷ σωτῆρί μου.
\vs{19}Κύριος ὁ Θεὸς δύναμίς μου, καὶ τάξει τοὺς πόδας μου εἰς συντέλειαν· ἐπὶ τὰ ὑψηλὰ ἐπιβιβᾷ με, τοῦ νικῆσαι ἐν τῇ ᾠδῇ αὐτοῦ.


\def\book{ΣΟΦΟΝΙΑΣ}
\biblebook{ΣΟΦΟΝΙΑΣ}


\lettrine[lines=2, loversize=0.2, nindent=0em, findent=.25em]{\textcolor{bookheadingcolor}{Λ}}{ΟΓΟΣ} Κυρίου, ὃς ἐγενήθη πρὸς Σοφονίαν τὸν τοῦ Χουσὶ, υἱὸν Γοδολίου, τοῦ Ἀμορίου τοῦ Ἐζεκίου, ἐν ἡμέραις Ἰωσίου υἱοῦ Ἀμὼν βασιλέως Ἰούδα.

\vs{2}Ἐκλείψει ἐκλιπέτω ἀπὸ προσώπου τῆς γῆς, λέγει Κύριος.
\vs{3}Ἐκλιπέτω ἄνθρωπος καὶ κτήνη, ἐκλιπέτω τὰ πετεινὰ τοῦ οὐρανοῦ καὶ οἱ ἰχθύες τῆς θαλάσσης· καὶ ἀσθενήσουσιν οἱ ἀσεβεῖς, καὶ ἐξαρῶ τοὺς ἀνόμους ἀπὸ προσώπου τῆς γῆς, λέγει Κύριος.
\vs{4}Καὶ ἐκτενῶ τὴν χεῖρά μου ἐπὶ Ἰούδα, καὶ ἐπὶ πάντας τοὺς κατοικοῦντας Ἱερουσαλήμ· καὶ ἐξαρῶ ἐκ τοῦ τόπου τούτου τὰ ὀνόματα τῆς Βάαλ, καὶ τὰ ὀνόματα τῶν ἱερέων,
\vs{5}καὶ τοὺς προσκυνοῦντας ἐπὶ τὰ δώματα τῇ στρατιᾷ τοῦ οὐρανοῦ, καὶ τοὺς προσκυνοῦντας καὶ τοὺς ὀμνύοντας κατὰ τοῦ Κυρίου, καὶ τοὺς ὀμνύοντας κατὰ τοῦ βασιλέως αὐτῶν,
\vs{6}καὶ τοὺς ἐκκλίνοντας ἀπὸ τοῦ Κυρίου, καὶ τοὺς μὴ ζητοῦντας τὸν Κύριον, καὶ τοὺς μὴ ἀντεχομένους τοῦ Κυρίου.

\vs{7}Εὐλαβεῖσθε ἀπὸ προσώπου Κυρίου τοῦ Θεοῦ· διότι ἐγγὺς ἡμέρα τοῦ Κυρίου, ὅτι ἡτοίμακε Κύριος τὴν θυσίαν αὐτοῦ, καὶ ἡγίακε τοὺς κλητοὺς αὐτοῦ.
\vs{8}Καὶ ἔσται, ἐν ἡμέρᾳ θυσίας Κυρίου, καὶ ἐκδικήσω ἐπὶ τοὺς ἄρχοντας, καὶ ἐπὶ τὸν οἶκον τοῦ βασιλέως, καὶ ἐπὶ πάντας τοὺς ἐνδεδυμένους ἐνδύματα ἀλλότρια.
\vs{9}Καὶ ἐκδικήσω ἐμφανῶς ἐπὶ τὰ πρόπυλα ἐν ἐκείνῃ τῇ ἡμέρᾳ, τοὺς πληροῦντας τὸν οἶκον Κυρίου Θεοῦ αὐτῶν ἀσεβείας καὶ δόλου.

\vs{10}Καὶ ἔσται ἐν τῇ ἡμέρᾳ ἐκείνῃ, λέγει Κύριος, φωνὴ κραυγῆς ἀπὸ πύλης ἀποκεντούντων, καὶ ὀλολυγμὸς ἀπὸ τῆς δευτέρας, καὶ συντριμμὸς μέγας ἀπὸ τῶν βουνῶν.
\vs{11}Θρηνήσατε οἱ κατοικοῦντες τὴν κατακεκομμένην, ὅτι ὡμοιώθη πᾶς ὁ λαὸς Χαναὰν, καὶ ἐξωλοθρεύθησαν πάντες οἱ ἐπῃρμένοι ἀργυρίῳ.

\vs{12}Καὶ ἔσται ἐν τῇ ἡμέρᾳ ἐκείνῃ, ἐξερευνήσω τὴν Ἱερουσαλὴμ μετὰ λύχνου, καὶ ἐκδικήσω ἐπὶ τοὺς ἄνδρας τοὺς καταφρονοῦντας ἐπὶ τὰ φυλάγματα αὐτῶν· οἱ δὲ λέγοντες ἐν ταῖς καρδίαις αὐτῶν, οὐ μὴ ἀγαθοποιήσῃ Κύριος, οὐδὲ μὴ κακώσῃ·
\vs{13}καὶ ἔσται ἡ δύναμις αὐτῶν εἰς διαρπαγὴν, καὶ οἱ οἶκοι αὐτῶν εἰς ἀφανισμόν· καὶ οἰκοδομήσουσιν οἰκίας, καὶ οὐ μὴ κατοικήσουσιν ἐν αὐταῖς· καὶ καταφυτεύσουσιν ἀμπελῶνας, καὶ οὐ μὴ πίωσι τὸν οἶνον αὐτῶν.

\vs{14}Ὅτι ἐγγὺς ἡμέρα Κυρίου ἡ μεγάλη, ἐγγὺς καὶ ταχεῖα σφόδρα· φωνὴ ἡμέρας Κυρίου πικρὰ καὶ σκληρὰ τέτακται·
\vs{15}δυνατὴ ἡμέρα ὀργῆς, ἡ ἡμέρα ἐκείνη, ἡμέρα θλίψεως καὶ ἀνάγκης, ἡμέρα ἀωρίας καὶ ἀφανισμοῦ, ἡμέρα γνόφου καὶ σκότους, ἡμέρα νεφέλης καὶ ὁμίχλης,
\vs{16}ἡμέρα σάλπιγγος καὶ κραυγῆς ἐπὶ τὰς πόλεις τὰς ὀχυρὰς, καὶ ἐπὶ τὰς γωνίας τὰς ὑψηλάς.
\vs{17}Καὶ ἐκθλίψω τοὺς ἀνθρώπους, καὶ πορεύσονται ὡς τυφλοὶ, ὅτι τῷ Κυρίῳ ἐξήμαρτον· καὶ ἐκχεεῖ τὸ αἷμα αὐτῶν ὡς χοῦν, καὶ τὰς σάρκας αὐτῶν ὡς βόλβιτα.
\vs{18}Καὶ τὸ ἀργύριον αὐτῶν καὶ τὸ χρυσίον αὐτῶν οὐ μὴ δύνηται ἐξελέσθαι αὐτοὺς ἐν ἡμέρᾳ ὀργῆς Κυρίου· καὶ ἐν πυρὶ ζήλου αὐτοῦ καταναλωθήσεται πᾶσα ἡ γῆ, διότι συντέλειαν καὶ σπουδὴν ποιήσει ἐπὶ πάντας τοὺς κατοικοῦντας τὴν γῆν.

\ch{2}
Συνάχθητε, καὶ συνδέθητε τὸ ἔθνος τὸ ἀπαίδευτον,
\vs{2}πρὸ τοῦ γενέσθαι ὑμᾶς ὡς ἄνθος παραπορευόμενον, πρὸ τοῦ ἐπελθεῖν ἐφʼ ὑμᾶς ὀργὴν Κυρίου, πρὸ τοῦ ἐπελθεῖν ἐφʼ ὑμᾶς ἡμέραν θυμοῦ Κυρίου.
\vs{3}Ζητήσατε τὸν Κύριον πάντες ταπεινοὶ γῆς, κρίμα ἐργάζεσθε, καὶ δικαιοσύνην ζητήσατε, καὶ ἀποκρίνασθε αὐτὰ, ὅπως σκεπασθῆτε ἐν ἡμέρᾳ ὀργῆς Κυρίου.

\vs{4}Διότι Γάζα διηρπασμένη ἔσται, καὶ Ἀσκάλων εἰς ἀφανισμὸν, καὶ Ἄζωτος μεσημβρίας ἐκριφήσεται, καὶ Ἀκκαρὼν ἐκριζωθήσεται.
\vs{5}Οὐαὶ οἱ κατοικοῦντες τὸ σχοίνισμα τῆς θαλάσσης, πάροικοι Κρητῶν· λόγος Κυρίου ἐφʼ ὑμᾶς Χαναὰν, γῆ ἀλλοφύλων, καὶ ἀπολῶ ὑμᾶς ἐκ κατοικίας.
\vs{6}Καὶ ἔσται Κρήτη νομὴ ποιμνίων, καὶ μάνδρα προβάτων.
\vs{7}Καὶ ἔσται τὸ σχοίνισμα τῆς θαλάσσης τοῖς καταλοίποις οἴκου Ἰούδα, ἐπʼ αὐτοὺς νεμήσονται ἐν τοῖς οἴκοις Ἀσκάλωνος, δείλης καταλύσουσιν ἀπὸ προσώπου υἱῶν Ἰούδα, ὅτι ἐπέσκεπται αὐτοὺς Κύριος ὁ Θεὸς αὐτῶν, καὶ ἀποστρέψει τὴν αἰχμαλωσίαν αὐτῶν.

\vs{8}Ἤκουσα ὀνειδισμοὺς Μωὰβ, καὶ κονδυλισμοὺς υἱῶν Ἀμμὼν, ἐν οἷς ὠνείδιζον τὸν λαόν μου, καὶ ἐμεγαλύνοντο ἐπὶ τὰ ὅριά μου.
\vs{9}Διατοῦτο, ζῶ ἐγὼ, λέγει Κύριος τῶν δυνάμεων ὁ Θεὸς Ἰσραὴλ, διότι Μωὰβ ὡς Σόδομα ἔσται, καὶ υἱοὶ Ἀμμὼν ὡς Γόμοῤῥα, καὶ Δαμασκὸς ἐκλελειμμένη ὡς θιμωνία ἅλωνος, καὶ ἠφανισμένη εἰς τὸν αἰῶνα· καὶ οἱ κατάλοιποι λαοῦ μου διαρπῶνται αὐτοὺς, καὶ οἱ κατάλοιποι ἔθνους μου κληρονομήσουσιν αὐτούς.
\vs{10}Αὕτη αὐτοῖς ἀντὶ τῆς ὕβρεως αὐτῶν, διότι ὠνείδισαν, καὶ ἐμεγαλύνθησαν ἐπὶ τὸν Κύριον τὸν παντοκράτορα.
\vs{11}Ἐπιφανήσεται Κύριος ἐπʼ αὐτοὺς, καὶ ἐξολοθρεύσει πάντας τοὺς θεοὺς τῶν ἐθνῶν τῆς γῆς, καὶ προσκυνήσουσιν αὐτῷ ἔκαστος ἐκ τοῦ τόπου αὐτοῦ, πᾶσαι αἱ νῆσοι τῶν ἐθνῶν.

\vs{12}Καὶ ὑμεῖς Αἰθίοπες τραυματίαι ῥομφαίας μου ἐστέ.

\vs{13}Καὶ ἐκτενεῖ τὴν χεῖρα αὐτοῦ ἐπὶ Βοῤῥᾶν, καὶ ἀπολεῖ τὸν Ἀσσύριον, καὶ θήσει τὴν Νινευὴ εἰς ἀφανισμὸν ἄνυδρον, ὡς ἔρημον.
\vs{14}Καὶ νεμήσονται ἐν μέσῳ αὐτῆς ποίμνια, καὶ πάντα τὰ θηρία τῆς γῆς, καὶ χαμαιλέοντες, καὶ ἐχῖνοι ἐν τοῖς φατνώμασιν αὐτῆς κοιτασθήσονται· καὶ θηρία φωνήσει ἐν τοῖς διορύγμασιν αὐτῆς καὶ κόρακες ἐν τοῖς πυλῶσιν αὐτῆς, διότι κέδρος τὸ ἀνάστημα αὐτῆς.

\vs{15}Αὕτη ἡ πόλις ἡ φαυλίστρια, ἡ κατοικοῦσα ἐπʼ ἐλπίδι, ἡ λέγουσα ἐν καρδίᾳ αὐτῆς, ἐγώ εἰμι, καὶ οὐκ ἔστι μετʼ ἐμὲ ἔτι· πῶς ἐγενήθη εἰς ἀφανισμὸν, νομὴ θηρίων; πᾶς ὁ διαπορευόμενος διʼ αὐτῆς συριεῖ, καὶ κινήσει τὰς χεῖρας αὐτοῦ.

\ch{3}
Ὢ ἡ ἐπιφανὴς καὶ ἀπολελυτρωμένη πόλις,
\vs{2}ἡ περιστερὰ οὐκ εἰσήκουσε φωνῆς· οὐκ ἐδέξατο παιδείαν, ἐπὶ τῷ Κυρίῳ οὐκ ἐπεποίθει, καὶ πρὸς τὸν Θεὸν αὐτῆς οὐκ ἤγγισεν.
\vs{3}Οἱ ἄρχοντες αὐτῆς ἐν αὐτῇ ὡς λέοντες ὠρυόμενοι, οἱ κριταὶ αὐτῆς ὡς λύκοι τῆς Ἀραβίας, οὐχ ὑπελίποντο εἰς τοπρωΐ.
\vs{4}Οἱ προφῆται αὐτῆς πνευματοφόροι, ἄνδρες καταφρονηταί· ἱερεῖς αὐτῆς βεβηλοῦσι τὰ ἅγια, καὶ ἀσεβοῦσι νόμον.

\vs{5}Ὁ δὲ Κύριος δίκαιος ἐν μέσῳ αὐτῆς, καὶ οὐ μὴ ποιήσῃ ἄδικον· πρωῒ πρωῒ δώσει κρίμα αὐτοῦ εἰς φῶς, καὶ οὐκ ἀπεκρύβη, καὶ οὐκ ἔγνω ἀδικίαν ἐν ἀπαιτήσει, καὶ οὐκ εἰς νεῖκος ἀδικίαν.
\vs{6}Ἐν διαφθορᾷ κατέσπασα ὑπερηφάνους, ἠφανίσθησαν γωνίαι αὐτῶν· ἐξερημώσω τὰς ὁδοὺς αὐτῶν τοπαράπαν, τοῦ μὴ διοδεύειν· ἐξέλιπον αἱ πόλεις αὐτῶν, παρὰ τὸ μηδένα ὑπάρχειν, μηδὲ κατοικεῖν.
\vs{7}Εἶπα, πλὴν φοβεῖσθέ με, καὶ δέξασθε παιδείαν, καὶ οὐ μὴ ἐξολοθρευθῆτε ἐξ ὀφθαλμῶν αὐτῆς πάντα ὅσα ἐξεδίκησα ἐπʼ αὐτήν· ἑτοιμάζου, ὄρθρισον, ἔφθαρται πᾶσα ἡ ἐπιφυλλὶς αὐτῶν.

\vs{8}Διατοῦτο ὑπόμεινόν με, λέγει Κύριος, εἰς ἡμέραν ἀναστάσεώς μου εἰς μαρτύριον· διὸ τὸ κρίμα μου εἰς συναγωγὰς ἐθνῶν, τοῦ εἰσδέκασθαι βασιλεῖς, τοῦ ἐκχέαι ἐπʼ αὐτοὺς πᾶσαν ὀργὴν θυμοῦ· διότι ἐν πυρὶ ζήλου μου καταναλωθήσεται πᾶσα ἡ γῆ.

\vs{9}Ὅτι τότε μεταστρέψω ἐπὶ λαοὺς γλῶσσαν εἰς γενεὰν αὐτῆς, τοῦ ἐπικαλεῖσθαι πάντας τὸ ὄνομα Κυρίου, τοῦ δουλεύειν αὐτῷ ὑπὸ ζυγὸν ἕνα.
\vs{10}Ἐκ περάτων ποταμῶν Αἰθιοπίας· προσδέξομαι ἐν διεσπαρμένοις μου, οἴσουσι θυσίας μοι.
\vs{11}Ἐν τῇ ἡμέρᾳ ἐκείνῃ, οὐ μὴ καταισχυνθῇς ἐκ πάντων τῶν ἐπιτηδευμάτων σου, ὧν ἠσέβησας εἰς ἐμέ· ὅτι τότε περιελῶ ἀπὸ σοῦ τὰ φαυλίσματα τῆς ὕβρεώς σου, καὶ οὐκ ἔτι μὴ προσθῇς, τοῦ μεγαλαυχῆσαι ἐπὶ τὸ ὄρος τὸ ἅγιόν μου.
\vs{12}Καὶ ὑπολήψομαι ἐν σοὶ λαὸν πρᾳῢν καὶ ταπεινὸν, καὶ εὐλαβηθήσονται ἀπὸ τοῦ ὀνόματος Κυρίου
\vs{13}οἱ κατάλοιποι τοῦ Ἰσραὴλ, καὶ οὐ ποιήσουσιν ἀδικίαν, καὶ οὐ λαλήσουσι μάταια, καὶ οὐ μὴ εὑρεθῇ ἐν τῷ στόματι αὐτῶν γλῶσσα δολία· διότι αὐτοὶ νεμήσονται, καὶ κοιτασθήσονται, καὶ οὐκ ἔσται ὁ ἐκφοβῶν αὐτούς.

\vs{14}Χαῖρε θύγατερ Σιὼν, κήρυσσε θύγατερ Ἱερουσαλήμ· εὐφραίνου καὶ κατατέρπου ἐξ ὅλης τῆς καρδίας σου θύγατερ Ἱερουσαλήμ.
\vs{15}Περιεῖλε Κύριος τὰ ἀδικήματά σου, λελύτρωταί σε ἐκ χειρὸς ἐχθρῶν σου· βασιλεὺς Ἰσραὴλ Κύριος ἐν μέσῳ σου, οὐκ ὄψῃ κακὰ οὐκέτι.

\vs{16}Ἐν τῷ καιρῷ ἐκείνῳ ἐρεῖ Κύριος τῇ Ἱερουσαλὴμ, θάρσει Σιὼν, μὴ παρείσθωσαν αἱ χεῖρές σου.
\vs{17}Κύριος ὁ Θεός σου ἐν σοὶ, ὁ δυνατὸς σώσει σε, ἐπάξει ἐπὶ δὲ εὐφροσύνην, καὶ καινιεῖ σε ἐν τῇ ἀγαπήσει αὐτοῦ· καὶ εὐφρανθήσεται ἐπὶ σὲ ἐν τέρψει ὡς ἐν ἡμέρᾳ ἑορτῆς.
\vs{18}Καὶ συνάξω τοὺς συντετριμμένους σου· οὐαὶ, τίς ἔλαβεν ἐπʼ αὐτὴν ὀνειδισμόν;

\vs{19}Ἰδοὺ ἐγὼ ποιῶ ἐν σοὶ ἔνεκέν σου ἐν τῷ καιρῷ ἐκείνῳ, λέγει Κύριος, καὶ σώσω τὴν ἐκπεπιεσμένην, καὶ τὴν ἀπωσμένην εἰσδέξομαι, καὶ θήσομαι αὐτοὺς εἰς καύχημα, καὶ ὀνομαστοὺς ἐν πάσῃ τῇ γῇ.
\vs{20}καὶ καταισχυνθήσονται ἐν τῷ καιρῷ ἐκείνῳ, ὅταν καλῶς ὑμῖν ποιήσω, καὶ ἐν τῷ καιρῷ, ὅταν εἰσδέξομαι ὑμᾶς· διότι δώσω ὑμᾶς ὀνομαστοὺς, καὶ εἰς καύχημα ἐν πᾶσι τοῖς λαοῖς τῆς γῆς, ἐν τῷ ἐπιστρέφειν με τὴν αἰχμαλωσίαν ὑμῶν ἐνώπιον ὑμῶν, λέγει Κύριος.


\def\book{ΑΓΓΑΙΟΣ}
\biblebook{ΑΓΓΑΙΟΣ}


\lettrine[lines=2, loversize=0.2, nindent=0em, findent=.25em]{\textcolor{bookheadingcolor}{Ἐ}}{Ν} τῷ δευτέρῳ ἔτει ἐπὶ Δαρείου τοῦ βασιλέως, ἐν τῷ μηνὶ τῷ ἕκτῳ, μιᾷ τοῦ μηνὸς, ἐγένετο λόγος Κυρίου ἐν χειρὶ Ἀγγαίου τοῦ προφήτου, λέγων, εἰπὸν πρὸς Ζοροβάβελ τὸν τοῦ Σαλαθιὴλ ἐκ φυλῆς Ἰούδα, καὶ πρὸς Ἰησοῦν τὸν τοῦ Ἰωσεδὲκ τὸν ἱερέα τὸν μέγαν, λέγων,
\vs{2}τάδε λέγει Κύριος παντοκράτωρ, λέγων, ὁ λαὸς οὗτος λέγουσιν, οὐκ ἧκεν ὁ καιρὸς τοῦ οἰκοδομῆσαι τὸν οἶκον Κυρίου.
\vs{3}Καὶ ἐγένετο λόγος Κυρίου ἐν χειρὶ Ἀγγαίου τοῦ προφήτου, λέγων,

\vs{4}Εἰ καιρὸς μὲν ὑμῖν ἐστι τοῦ οἰκεῖν ἐν οἴκοις ὑμῶν κοιλοστάθμοις, ὁ δὲ οἶκος ἡμῶν ἐξηρήμωται;

\vs{5}Καὶ νῦν τάδε λέγει Κύριος παντοκράτωρ, τάξατε δὴ καρδίας ὑμῶν εἰς τὰς ὁδοὺς ὑμῶν.
\vs{6}Ἐσπείρατε πολλὰ καὶ εἰσηνέγκατε ὀλίγα, ἐφάγετε καὶ οὐκ εἰς πλησμονὴν, ἐπίετε καὶ οὐκ εἰς μέθην, περιεβάλεσθε καὶ οὐκ ἐθερμάνθητε ἐν αὐτοῖς, καὶ ὁ τοὺς μισθοὺς συνάγων, συνήγαγεν εἰς δεσμὸν τετρυπημένον.

\vs{7}Τάδε λέγει Κύριος παντοκράτωρ, θέσθε τὰς καρδίας ὑμῶν εἰς τὰς ὁδοὺς ὑμῶν.
\vs{8}Ἀνάβητε εἰς τὸ ὄρος, καὶ κόψατε ξύλα, οἰκοδομήσατε τὸν οἶκον, καὶ εὐδοκήσω ἐν αὐτῷ, καὶ ἐνδοξασθήσομαι, εἶπε Κύριος.
\vs{9}Ἐπεβλέψατε εἰς πολλὰ, καὶ ἐγένετο ὀλίγα· καὶ εἰσηνέχθη εἰς τὸν οἶκον, καὶ ἐξεφύσησα αὐτά· διατοῦτο τάδε λέγει Κύριος παντοκράτωρ, ἀνθʼ ὧν ὁ οἶκός μου ἐστὶν ἔρημος, ὑμεῖς δὲ διώκετε ἕκαστος εἰς τὸν οἶκον αὐτοῦ,
\vs{10}διατοῦτο ἀνέξει ὁ οὐρανὸς ἀπὸ δρόσου, καὶ ἡ γῆ ὑποστελεῖται τὰ ἐκφόρια αὐτῆς.
\vs{11}Καὶ ἐπάξω ῥομφαίαν ἐπὶ τὴν γῆν, καὶ ἐπὶ τὰ ὄρη, καὶ ἐπὶ τὸν σῖτον, καὶ ἐπὶ τὸν οἶνον, καὶ ἐπὶ τὸ ἔλαιον, καὶ ὅσα ἐκφέρει ἡ γῆ, καὶ ἐπὶ τοὺς ἀνθρώπους, καὶ ἐπὶ τὰ κτήνη, καὶ ἐπὶ πάντας τοὺς πόνους τῶν χειρῶν αὐτῶν.

\vs{12}Καὶ ἤκουσε Ζοροβάβελ ὁ τοῦ Σαλαθιὴλ ἐκ φυλῆς Ἰούδα, καὶ Ἰησοῦς ὁ τοῦ Ἰωσεδὲκ ὁ ἱερεὺς ὁ μέγας, καὶ πάντες οἱ κατάλοιποι τοῦ λαοῦ τῆς φωνῆς Κυρίου τοῦ Θεοῦ αὐτῶν, καὶ τῶν λόγων τοῦ Ἀγγαίου τοῦ προφήτου, καθότι ἐξαπέστειλεν αὐτὸν Κύριος ὁ Θεὸς αὐτῶν πρὸς αὐτοὺς, καὶ ἐφοβήθη ὁ λαὸς ἀπὸ προσώπου Κυρίου.
\vs{13}Καὶ εἶπεν Ἀγγαῖος ἄγγελος Κυρίου ἐν ἀγγέλοις Κυρίου τῷ λαῷ, ἐγώ εἰμι μεθʼ ὑμῶν, λέγει Κύριος.

\vs{14}Καὶ ἐξήγειρε Κύριος τὸ πνεῦμα Ζοροβάβελ τοῦ Σαλαθιὴλ ἐκ φυλῆς Ἰούδα, καὶ τὸ πνεῦμα Ἰησοῦ τοῦ Ἰωσεδὲκ τοῦ ἱερέως τοῦ μεγάλου, καὶ τὸ πνεῦμα τῶν καταλοίπων παντὸς τοῦ λαοῦ, καὶ εἰσῆλθον, καὶ ἐποίουν ἔργα ἐν τῷ οἴκῳ Κυρίου παντοκράτορος Θεοῦ αὐτῶν.

\vs{15}Τῇ τετράδι καὶ εἰκάδι τοῦ μηνὸς τοῦ ἕκτου, τῷ δευτέρῳ ἔτει, ἐπὶ Δαρείου τοῦ βασιλέως.

\ch{2}Τῷ μηνὶ τῷ ἑβδόμῳ, μιᾷ καὶ εἰκάδι τοῦ μηνὸς, ἐλάλησε Κύριος ἐν χειρὶ Ἀγγαίου τοῦ προφήτου, λέγων,
\vs{2}εἰπὸν δὴ πρὸς Ζοροβάβελ τὸν τοῦ Σαλαθιὴλ ἐκ φυλῆς Ἰούδα, καὶ πρὸς Ἰησοῦν τοῦ Ἰωσεδὲκ τὸν ἱερέα τὸν μέγαν, καὶ πρὸς πάντας τοὺς καταλοίπους τοῦ λαοῦ, λέγων,

\vs{3}Τίς ἐξ ὑμῶν, ὃς εἶδε τὸν οἶκον τοῦτον ἐν τῇ δόξῃ αὐτοῦ τῇ ἔμπροσθεν; καὶ πῶς ὑμεῖς βλέπετε αὐτὸν νῦν καθὼς οὐχ ὑπάρχοντα ἐνώπιον ὑμῶν;
\vs{4}Καὶ νῦν κατίσχυε Ζοροβάβελ, λέγει Κύριος, καὶ κατίσχυε Ἰησοῦ ὁ τοῦ Ἰωσεδὲκ ὁ ἱερεὺς ὁ μέγας, καὶ κατισχυέτω πᾶς ὁ λαὸς τῆς γῆς, λέγει Κύριος, καὶ ποιεῖτε, διότι μεθʼ ὑμῶν ἐγώ εἰμι, λέγει Κύριος ὁ παντοκράτωρ,
\vs{5}καὶ τὸ πνεῦμά μου ἐφέστηκεν ἐν μέσῳ ὑμῶν· θαρσεῖτε,

\vs{6}Διότι τάδε λέγει Κύριος παντοκράτωρ, ἔτι ἅπαξ ἐγὼ σείσω τὸν οὐρανὸν καὶ τὴν γῆν, καὶ τὴν θάλασσαν καὶ τὴν ξηρὰν,
\vs{7}καὶ συσσείσω πάντα τὰ ἔθνη, καὶ ἥξει τὰ ἐκλεκτὰ πάντων τῶν ἐθνῶν· καὶ πλήσω τὸν οἶκον τοῦτον δόξης, λέγει Κύριος παντοκράτωρ.
\vs{8}Ἐμὸν τὸ ἀργύριον, καὶ ἐμὸν τὸ χρυσίον, λέγει Κύριος παντοκράτωρ.
\vs{9}Διότι μεγάλη ἔσται ἡ δόξα τοῦ οἴκου τούτου, ἡ ἐσχάτη ὑπὲρ τὴν πρώτην, λέγει Κύριος παντοκράτωρ· καὶ ἐν τῷ τόπῳ τούτῳ δώσω εἰρήνην, λέγει Κύριος παντοκράτωρ, καὶ εἰρήνην ψυχῆς εἰς περιποίησιν παντὶ τῷ κτίζοντι, τοῦ ἀναστῆσαι τὸν ναὸν τοῦτον.

\vs{10}Τετράδι καὶ εἰκάδι τοῦ ἐννάτου μηνὸς, ἔτους δευτέρου, ἐπὶ Δαρείου, ἐγένετο λόγος Κυρίου πρὸς Ἀγγαῖον τὸν προφήτην, λέγων,
\vs{11}τάδε λέγει Κύριος παντοκράτωρ, ἐπερώτησον δὴ τοὺς ἱερεῖς νόμον, λέγων,
\vs{12}ἐὰν λάβῃ ἄνθρωπος κρέας ἅγιον ἐν τῷ ἄκρῳ τοῦ ἱματίου αὐτοῦ, καὶ ἅψηται τὸ ἄκρον τοῦ ἱματίου αὐτοῦ ἄρτου, ἢ ἑψήματος, ἢ οἴνου, ἢ ἐλαίου, ἢ παντὸς βρώματος, εἰ ἁγιασθήσεται; καὶ ἀπεκρίθησαν οἱ ἱερεῖς, καὶ εἶπαν, οὔ.
\vs{13}Καὶ εἶπεν Ἀγγαῖος, ἐὰν ἅψηται μεμιασμένος ἀκάθαρτος ἐπὶ ψυχῇ ἐπὶ παντὸς τούτων, εἰ μιανθήσεται; καὶ ἀπεκρίθησαν οἱ ἱερεῖς, καὶ εἶπαν, μιανθήσεται.
\vs{14}Καὶ ἀπεκρίθη Ἀγγαῖος, καὶ εἶπεν, οὕτως ὁ λαὸς οὗτος, καὶ οὕτως τὸ ἔθνος τοῦτο ἐνώπιον ἐμοῦ, λέγει Κύριος, καὶ οὕτως πάντα τὰ ἔργα τῶν χειρῶν αὐτῶν· καὶ ὃς ἐὰν ἐγγίσῃ ἐκεῖ, μιανθήσεται ἕνεκεν τῶν λημμάτων αὐτῶν τῶν ὀρθρινῶν, ὀδυνηθήσονται ἀπὸ προσώπου πόνων αὐτῶν, καὶ ἐμισεῖτε ἐν πύλαις ἐλέγχοντα.
\vs{15}Καὶ νῦν θέσθε δὴ εἰς τὰς καρδίας ὑμῶν ἀπὸ τῆς ἡμέρας ταύτης καὶ ὑπεράνω, πρὸ τοῦ θεῖναι λίθον ἐπὶ λίθον ἐν τῷ ναῷ Κυρίου
\vs{16}τίνες ἦτε, ὅτε ἐνεβάλλετε εἰς κυψέλην κριθῆς εἴκοσι σάτα, καὶ ἐγένετο κριθῆς δέκα σάτα· καὶ εἰσεπορεύεσθε εἰς τὸ ὑπολήνιον ἐξαντλῆσαι πεντήκοντα μετρητὰς καὶ ἐγένοντο εἴκοσι.
\vs{17}Ἐπάταξα ὑμᾶς ἐν ἀφορίᾳ, καὶ ἐν ἀνεμοφθορίᾳ, καὶ ἐν χαλάζῃ πάντα τὰ ἔργα τῶν χειρῶν ὑμῶν, καὶ οὐκ ἐπεστρέψατε πρὸς μὲ, λέγει Κύριος.

\vs{18}Ὑποτάξατε δὴ τὰς καρδίας ὑμῶν ἀπὸ τῆς ἡμέρας ταύτης, καὶ ἐπέκεινα, ἀπὸ τῆς τετράδος καὶ εἰκάδος τοῦ ἐννάτου μηνὸς, καὶ ἀπὸ τῆς ἡμέρας ἧς τεθεμελίωται ὁ ναὸς Κυρίου· θέσθε ἐν ταῖς καρδίαις ὑμῶν,
\vs{19}εἰ ἐπιγνωσθήσεται ἐπὶ τῆς ἅλω, καὶ εἰ ἔτι ἡ ἄμπελος, καὶ ἡ συκῆ, καὶ ἡ ῥοὰ, καὶ τὰ ξύλα τῆς ἐλαίας τὰ οὐ φέροντα καρπὸν, ἀπὸ τῆς ἡμέρας ταύτης εὐλογήσω.

\vs{20}Καὶ ἐγένετο λόγος Κυρίου ἐκ δευτέρου πρὸς Ἀγγαῖον τὸν προφήτην, τετράδι καὶ εἰκάδι τοῦ μηνὸς, λέγων,
\vs{21}εἰπὸν πρὸς Ζοροβάβελ τὸν τοῦ Σαλαθιὴλ ἐκ φυλῆς Ἰούδα, λέγων,

Ἐγὼ σείω τὸν οὐρανὸν καὶ τὴν γῆν, καὶ τὴν θάλασσαν καὶ τὴν ξηρὰν,
\vs{22}καὶ καταστρέψω θρόνους βασιλέων, καὶ ὀλοθρεύσω δύναμιν βασιλέων τῶν ἐθνῶν, καὶ καταστρέψω ἅρματα καὶ ἀναβάτας, καὶ καταβήσονται ἵπποι καὶ ἀναβάται αὐτῶν, ἕκαστος ἐν ῥομφαίᾳ πρὸς τὸν ἀδελφὸν αὐτοῦ.
\vs{23}Ἐν τῇ ἡμέρᾳ ἐκείνῃ, λέγει Κύριος παντοκράτωρ, λήψομαί σε Ζοροβάβελ τὸν τοῦ Σαλαθιὴλ, τὸν δοῦλόν μου, λέγει Κύριος, καὶ θήσομαί σε ὡς σφραγίδα, διότι σὲ ᾑρέτισα, λέγει Κύριος παντοκράτωρ.


\def\book{ΖΑΧΑΡΙΑΣ}
\biblebook{ΖΑΧΑΡΙΑΣ}


\lettrine[lines=2, loversize=0.2, nindent=0em, findent=.25em]{\textcolor{bookheadingcolor}{Ἐ}}{Ν} τῷ ὀγδόῳ μηνὶ, ἔτους δευτέρου ἐπὶ Δαρείου, ἐγένετο λόγος Κυρίου πρὸς Ζαχαρίαν τὸν τοῦ Βαραχίου υἱὸν Ἀδδὼ τὸν προφήτην, λέγων,

\vs{2}Ὠργίσθη Κύριος ἐπὶ τοὺς πατέρας ὑμῶν ὀργὴν μεγάλην·
\vs{3}καὶ ἐρεῖς πρὸς αὐτοῦς, τάδε λέγει Κύριος παντοκράτωρ, ἐπιστρέψατε πρὸς μὲ, λέγει Κύριος τῶν δυνάμεων, καὶ ἐπιστραφήσομαι πρὸς ὑμᾶς, λέγει Κύριος τῶν δυνάμεων.
\vs{4}Καὶ μὴ γίνεσθε καθὼς οἱ πατέρες ὑμῶν, οἷς ἐνεκάλεσαν αὐτοῖς οἱ προφῆται ἔμπροσθεν λέγοντες, τάδε λέγει Κύριος παντοκράτωρ, ἀποστρέψατε ἀπὸ τῶν ὁδῶν ὑμῶν τῶν πονηρῶν, καὶ ἀπὸ τῶν ἐπιτηδευμάτων ὑμῶν τῶν πονηρῶν· καὶ οὐκ εἰσήκουσαν, καὶ οὐ προσέσχον τοῦ εἰσακοῦσαί μου, λέγει Κύριος.

\vs{5}Οἱ πατέρες ὑμῶν ποῦ εἰσι καὶ οἱ προφῆται; μὴ τὸν αἰῶνα ζήσονται;
\vs{6}Πλὴν τοὺς λόγους μου καὶ τὰ νόμιμά μου δέχεσθε, ὅσα ἐγὼ ἐντέλλομαι ἐν πνεύματί μου τοῖς δούλοις μου τοῖς προφήταις, οἳ κατελάβοσαν τοὺς πατέρας ὑμῶν· καὶ ἀπεκρίθησαν, καὶ εἶπαν, καθὼς παρατέτακται Κύριος παντοκράτωρ τοῦ ποιῆσαι ἡμῖν κατὰ τὰς ὁδοὺς ἡμῶν καὶ κατὰ τὰ ἐπιτηδεύματα ἡμῶν, οὕτως ἐποίησεν ἡμῖν.

\vs{7}Τῇ τετράδι καὶ εἰκάδι, τῷ ἑνδεκάτῳ μηνὶ, οὗτός ἐστιν ὁ μὴν Σαβὰτ, ἐν τῷ δευτέρῳ ἔτει, ἐπὶ Δαρείου, ἐγένετο λόγος Κυρίου πρὸς Ζαχαρίαν τὸν τοῦ Βαραχίου υἱὸν Ἀδδὼ τὸν προφήτην, λέγων,

\vs{8}Ἑώρακα τὴν νύκτα, καὶ ἰδοὺ ἀνὴρ ἐπιβεβηκὼς ἐπὶ ἵππον πυῤῥὸν, καὶ οὗτος εἱστήκει ἀναμέσον τῶν ὀρέων τῶν κατασκίων, καὶ ὀπίσω αὐτοῦ ἵπποι πυῤῥοὶ, καὶ ψαροὶ, καὶ ποικίλοι, καὶ λευκοί.
\vs{9}Καὶ εἶπα, τί οὗτοι κύριε; καὶ εἶπε πρὸς μὲ ὁ ἄγγελος ὁ λαλῶν ἐν ἐμοὶ, ἐγὼ δείξω σοι τί ἐστι ταῦτα.
\vs{10}Καὶ ἀπεκρίθη ὁ ἀνὴρ ὁ ἐφεστηκὼς ἀναμέσον τῶν ὀρέων, καὶ εἶπε πρὸς μὲ, οὗτοί εἰσιν οὓς ἐξαπέστειλε Κύριος, περιοδεῦσαι τὴν γῆν·
\vs{11}καὶ ἀπεκρίθησαν τῷ ἀγγέλῳ Κυρίου τῷ ἐφεστῶτι ἀναμέσον τῶν ὀρέων, καὶ εἶπον, περιωδεύσαμεν πᾶσαν τὴν γῆν, καὶ ἰδοὺ πᾶσα ἡ γῆ κατοικεῖται, καὶ ἡσυχάζει.

\vs{12}Καὶ ἀπεκρίθη ὁ ἄγγελος Κυρίου, καὶ εἶπε, Κύριε παντοκράτωρ ἕως τίνος οὐ μὴ ἐλεήσῃς τὴν Ἱερουσαλὴμ, καὶ τὰς πόλεις Ἰούδα, ἃς ὑπερεῖδες, τοῦτο ἑβδομηκοστὸν ἔτος;
\vs{13}Καὶ ἀπεκρίθη Κύριος παντοκράτωρ τῷ ἀγγέλῳ λαλοῦντι ἐν ἐμοὶ, ῥήματα καλὰ καὶ λόγους παρακλητικούς.
\vs{14}Καὶ εἶπε πρὸς μὲ ὁ ἄγγελος ὁ λαλῶν ἐν ἐμοὶ, ἀνάκραγε λέγων,

Τάδε λέγει Κύριος παντοκράτωρ, ἐζήλωκα τὴν Ἱερουσαλὴμ καὶ τὴν Σιὼν ζῆλον μέγαν,
\vs{15}καὶ ὀργὴν μεγάλην ἐγὼ ὀργίζομαι ἐπὶ τὰ ἔθνη τὰ συνεπιτιθέμενα, ἀνθʼ ὧν μὲν ἐγὼ ὠργίσθην ὀλίγα, αὐτοὶ δὲ συνεπέθεντο εἰς κακά.
\vs{16}Διατοῦτο τάδε λέγει Κύριος, ἐπιστρέψω ἐπὶ Ἱερουσαλὴμ ἐν οἰκτιρμῷ, καὶ ὁ οἶκός μου ἀνοικοδομηθήσεται ἐν αὐτῇ, λέγει Κύριος παντοκράτωρ, καὶ μέτρον ἐκταθήσεται ἐπὶ Ἱερουσαλὴμ ἔτι·
\vs{17}καὶ εἶπε πρὸς μὲ ὁ ἄγγελος ὁ λαλῶν ἐν ἐμοὶ, ἔτι ἀνάκραγε λέγων, τάδε λέγει Κύριος παντοκράτωρ, ἔτι διαχυθήσονται πόλεις ἐν ἀγαθοῖς, καὶ ἐλεήσει Κύριος ἔτι τὴν Σιὼν, καὶ αἱρετιεῖ τὴν Ἱερουσαλήμ.

\ch{2}
Καὶ ᾖρα τοὺς ὀφθαλμούς μου, καὶ ἴδον, καὶ ἰδοὺ τέσσαρα κέρατα.
\vs{2}Καὶ εἶπα πρὸς τὸν ἄγγελον τὸν λαλοῦντα ἐν ἐμοὶ, τί ἐστι ταῦτα κύριε; καὶ εἶπε πρὸς μὲ, ταῦτα τὰ κέρατα τὰ διασκορπίσαντα τὸν Ἰούδαν, καὶ τὸν Ἰσραὴλ, καὶ Ἱερουσαλήμ.
\vs{3}Καὶ ἔδειξέ μοι Κύριος τέσσαρας τέκτονας.
\vs{4}Καὶ εἶπα, τί οὗτοι ἔρχονται ποιῆσαι; καὶ εἶπε ταῦτα τὰ κέρατα τὰ διασκορπίσαντα τὸν Ἰούδα καὶ τὸν Ἰσραὴλ κατέαξαν, καὶ οὐδεὶς αὐτῶν ᾖρε κεφαλήν· καὶ ἐξήλθοσαν οὗτοι τοῦ ὀξῦναι αὐτὰ εἰς χεῖρας αὐτῶν, τὰ τέσσαρα κέρατα, τὰ ἔθνη τὰ ἐπαιρόμενα κέρας ἐπὶ τὴν γῆν Κυρίου, τοῦ διασκορπίσαι αὐτήν.

\vs{5}Καὶ ᾖρα τοὺς ὀφθαλμούς μου, καὶ ἴδον, καὶ ἰδοὺ ἀνὴρ, καὶ ἐν τῇ χειρὶ αὐτοῦ σχοινίον γεωμετρικόν.
\vs{6}Καὶ εἶπα πρὸς αὐτόν, ποῦ σὺ πορεύῃ; καὶ εἶπε πρὸς μὲ, διαμετρῆσαι τὴν Ἱερουσαλὴμ, τοῦ ἰδεῖν πηλίκον τὸ πλάτος αὐτῆς ἐστι, καὶ πηλίκον τὸ μῆκος.
\vs{7}Καὶ ἰδοὺ ὁ ἄγγελος ὁ λαλῶν ἐν ἐμοὶ εἱστήκει, καὶ ἄγγελος ἕτερος ἐξεπορεύετο εἰς συνάντησιν αὐτῷ,
\vs{8}καὶ εἶπε πρὸς αὐτὸν, λέγων, δράμε, καὶ λάλησον πρὸς τὸν νεανίαν ἐκεῖνον, λέγων,

Κατακάρπως κατοικηθήσεται Ἱερουσαλὴμ ἀπὸ πλήθους ἀνθρώπων καὶ κτηνῶν ἐν μέσῳ αὐτῆς·
\vs{9}καὶ ἐγὼ ἔσομαι αὐτῇ, λέγει Κύριος, τεῖχος πυρὸς κυκλόθεν, καὶ εἰς δόξαν ἔσομαι ἐν μέσῳ αὐτῆς.

\vs{10}Ὢ ὢ φεύγετε ἀπὸ γῆς Βοῤῥᾶ, λέγει Κύριος· διότι ἐκ τῶν τεσσάρων ἀνέμων τοῦ οὐρανοῦ συνάξω ὑμᾶς, λέγει Κύριος,
\vs{11}εἰς Σιὼν, ἀνασώζεσθε οἱ κατοικοῦντες θυγατέρα Βαβυλῶνος.
\vs{12}Διότι τάδε λέγει Κύριος παντοκράτωρ, ὀπίσω δόξης ἀπέσταλκέ με ἐπὶ τὰ ἔθνη τὰ σκυλεύσαντα ὑμᾶς, διότι ὁ ἁπτόμενος ὑμῶν ὡς ὁ ἁπτόμενος τῆς κόρης τοῦ ὀφθαλμοῦ αὐτοῦ.
\vs{13}Διότι ἰδοὺ ἐγὼ ἐπιφέρω τὴν χεῖρά μου ἐπʼ αὐτοὺς, καὶ ἔσονται σκῦλα τοῖς δουλεύουσιν αὐτοῖς, καὶ γνώσεσθε ὅτι Κύριος παντοκράτωρ ἀπέσταλκέ με.

\vs{14}Τέρπου καὶ εὐφραίνου θύγατερ Σιὼν, διότι ἰδοὺ ἐγὼ ἔρχομαι, καὶ κατασκηνώσω ἐν μέσῳ σου, λέγει Κύριος.
\vs{15}Καὶ καταφεύξονται ἔθνη πολλὰ ἐπὶ τὸν Κύριον ἐν τῇ ἡμέρᾳ ἐκείνῃ, καὶ ἔσονται αὐτῷ εἰς λαὸν, καὶ κατασκηνώσουσιν ἐν μέσῳ σου, καὶ ἐπιγνώσῃ ὅτι Κύριος παντοκράτωρ ἐξαπέσταλκέ με πρὸς σέ.
\vs{16}Καὶ κατακληρονομήσει Κύριος τὸν Ἰούδαν τὴν μερίδα αὐτοῦ ἐπὶ τὴν ἁγίαν· καὶ αἱρετιεῖ ἔτι τὴν Ἱερουσαλήμ.
\vs{17}Εὐλαβείσθω πᾶσα σὰρξ ἀπὸ προσώπου Κυρίου, ὅτι ἐξεγήγερται ἐκ νεφελῶν ἁγίων αὐτοῦ.

\ch{3}
Καὶ ἔδειξέ μοι Κύριος τὸν Ἰησοῦν τὸν ἱερέα τὸν μέγαν, ἑστῶτα πρὸ προσώπου ἀγγέλου Κυρίου, καὶ ὁ διάβολος εἱστήκει ἐκ δεξιῶν αὐτοῦ, τοῦ ἀντικεῖσθαι αὐτῷ.
\vs{2}Καὶ εἶπε Κύριος πρὸς τὸν διάβολον,

Ἐπιτιμήσαι Κύριος ἐν σοὶ διάβολε, καὶ ἐπιτιμήσαι Κύριος ἐν σοὶ ὁ ἐκλεξάμενος τὴν Ἱερουσαλήμ· οὐκ ἰδοὺ τοῦτο ὡς δαλὸς ἐξεσπασμένος ἐκ πυρός;

\vs{3}Καὶ Ἰησοῦς ἦν ἐνδεδυμένος ἱμάτια ῥυπαρὰ, καὶ εἱστήκει πρὸ προσώπου τοῦ ἀγγέλου.
\vs{4}Καὶ ἀπεκρίθη καὶ εἶπε πρὸς τοὺς ἑστηκότας πρὸ προσώπου αὐτοῦ, λέγων, ἀφέλετε τὰ ἱμάτια τὰ ῥυπαρὰ ἀπʼ αὐτοῦ· καὶ εἶπε πρὸς αὐτὸν, ἰδοὺ ἀφῄρηκα τὰς ἀνομίας σου· καὶ ἐνδύσατε αὐτὸν ποδήρη,
\vs{5}καὶ ἐπίθετε κίδαριν καθαρὰν ἐπὶ τὴν κεφαλὴν αὐτοῦ· καὶ ἐπέθηκαν κίδαριν καθαρὰν ἐπὶ τὴν κεφαλὴν αὐτοῦ, καὶ περιέβαλον αὐτὸν ἱμάτια· καὶ ὁ ἄγγελος Κυρίου εἱστήκει.
\vs{6}Καὶ διεμαρτύρατο ὁ ἄγγελος Κυρίου πρὸς Ἰησοῦν, λέγων,
\vs{7}τάδε λέγει Κύριος παντοκράτωρ,

Ἐὰν ταῖς ὁδοῖς μου πορεύῃ, καὶ ἐν τοῖς προστάγμασί μου φυλάξῃ, καὶ σὺ διακρινεῖς τὸν οἶκόν μου· καὶ ἐὰν διαφυλάσσῃς τὴν αὐλήν μου, καὶ δώσω σοι ἀναστρεφομένους ἐν μέσῳ τῶν ἑστηκότων τούτων.
\vs{8}Ἄκουε δὴ Ἰησοῦ ὁ ἱερεὺς ὁ μέγας, σὺ καὶ οἱ πλησίον σου οἱ καθήμενοι πρὸ προσώπου, διότι ἄνδρες τερατοσκόποι εἰσὶ, διότι ἰδοὺ ἐγὼ ἄγω τὸν δοῦλόν μου Ἀνατολήν.
\vs{9}Διότι ὁ λίθος ὃν ἔδωκα πρὸ προσώπου τοῦ Ἰησοῦ, ἐπὶ τὸν λίθον τὸν ἕνα ἑπτὰ ὀφθαλμοί εἰσιν· ἰδοὺ ἐγὼ ὀρύσσω βόθρον, λέγει Κύριος παντοκράτωρ, καὶ ψηλαφήσω πᾶσαν τὴν ἀδικίαν τῆς γῆς ἐκείνης ἐν ἡμέρᾳ μιᾷ.
\vs{10}Ἐν τῇ ἡμέρᾳ ἐκείνῃ, λέγει Κύριος παντοκράτωρ, συγκαλέσετε ἕκαστος τὸν πλησίον αὐτοῦ ὑποκάτω ἀμπέλου, καὶ ὑποκάτω συκῆς.

\ch{4}
Καὶ ἐπέστρεψεν ὁ ἄγγελος ὁ λαλῶν ἐν ἐμοὶ, καὶ ἐξήγειρέ με ὃν τρόπον ὅταν ἐξεγερθῇ ἄνθρωπος ἐξ ὕπνου αὐτοῦ.

\vs{2}Καὶ εἶπε πρὸς μὲ, τί σὺ βλέπεις; καὶ εἶπα, ἑώρακα, καὶ ἰδοὺ λυχνία χρυσῆ ὅλη, καὶ τὸ λαμπάδιον ἐπάνω αὐτῆς, καὶ ἑπτὰ λύχνοι ἐπάνω αὐτῆς, καὶ ἑπτὰ ἐπαρυστρίδες τοῖς λύχνοις τοῖς ἐπάνω αὐτῆς,
\vs{3}καὶ δύο ἐλαῖαι ἐπάνω αὐτῆς, μία ἐκ δεξιῶν τοῦ λαμπαδίου αὐτῆς, καὶ μία ἐξ εὐωνύμων.
\vs{4}Καὶ ἐπηρώτησα, καὶ εἶπα πρὸς τὸν ἄγγελον τὸν λαλοῦντα ἐν ἐμοὶ, λέγων, τί ἐστι ταῦτα κύριε;
\vs{5}Καὶ ἀπεκρίθη ὁ ἄγγελος ὁ λαλῶν ἐν ἐμοὶ, καὶ εἶπε πρὸς μὲ, λέγων, οὐ γινώσκεις τί ἐστι ταῦτα; καὶ εἶπα, οὐχὶ κύριε.
\vs{6}Καὶ ἀπεκρίθη, καὶ εἶπε πρὸς μὲ, λέγων, οὗτος ὁ λόγος Κυρίου πρὸς Ζοροβάβελ, λέγων,

Οὐκ ἐν δυνάμει μεγάλῃ, οὐδὲ ἐν ἰσχύϊ, ἀλλὰ ἐν πνεύματί μου, λέγει Κύριος παντοκράτωρ.
\vs{7}Τίς εἶ σὺ τὸ ὄρος τὸ μέγα τὸ πρὸ προσώπου Ζοροβάβελ τοῦ κατορθῶσαι; καὶ ἐξοίσω τὸν λίθον τῆς κληρονομίας, ἰσότητα χάριτος χάριτα αὐτῆς.

\vs{8}Καὶ ἐγένετο λόγος Κυρίου πρὸς μὲ, λέγων,
\vs{9}αἱ χεῖρες Ζοροβάβελ ἐθεμελίωσαν τὸν οἶκον τοῦτον, καὶ αἱ χεῖρες αὐτοῦ ἐπιτελέσουσιν αὐτόν· καὶ ἐπιγνώσῃ διότι Κύριος παντοκράτωρ ἐξαπέσταλκέ με πρὸς σέ.
\vs{10}Διότι τίς ἐξουδένωσεν εἰς ἡμέρας μικράς; καὶ χαροῦνται, καὶ ὄψονται τὸν λίθον τὸν κασσιτέρινον ἐν χειρὶ Ζοροβάβελ· ἑπτὰ οὗτοι ὀφθαλμοί εἰσιν οἱ ἐπιβλέποντες ἐπὶ πᾶσαν τὴν γῆν.

\vs{11}Καὶ ἀπεκρίθην, καὶ εἶπα πρὸς αὐτὸν, τί αἱ δύο ἐλαῖαι αὗται, αἱ ἐκ δεξιῶν τῆς λυχνίας καὶ ἐξ εὐωνύμων;
\vs{12}Καὶ ἐπηρώτησα ἐκ δευτέρου, καὶ εἶπα πρὸς αὐτὸν, τί οἱ δύο κλάδοι τῶν ἐλαιῶν οἱ ἐν ταῖς χερσὶ τῶν δύο μυξωτήρων τῶν χρυσῶν τῶν ἐπιχεόντων, καὶ ἐπαναγόντων τὰς ἐπαρυστρίδας τὰς χρυσᾶς;
\vs{13}Καὶ εἶπε πρὸς μὲ, οὐκ οἶδας τί ἐστι ταῦτα; καὶ εἶπα, οὐχὶ κύριε.
\vs{14}Καὶ εἶπεν, οὗτοι οἱ δύο υἱοὶ τῆς πιότητος παρεστήκασι Κυρίῳ πάσης τῆς γῆς.

\ch{5}
Καὶ ἐπέστρεψα, καὶ ᾖρα τοὺς ὀφθαλμούς μου, καὶ ἴδον, καὶ ἰδοὺ δρέπανον πετόμενον.
\vs{2}Καὶ εἶπε πρὸς μὲ, τί σὺ βλέπεις; καὶ εἶπα, ἐγὼ ὁρῶ δρέπανον πετόμενον μήκους πήχεων εἴκοσι, καὶ πλάτους πήχεων δέκα.
\vs{3}Καὶ εἶπε πρὸς μὲ,

Αὕτη ἡ ἀρὰ ἡ ἐκπορευομένη ἐπὶ πρόσωπον πάσης τῆς γῆς· διότι πᾶς ὁ κλέπτης ἐκ τούτου ἕως θανάτου ἐκδικηθήσεται, καὶ πᾶς ὁ ἐπίορκος ἐκ τούτου ἐκδικηθήσεται.
\vs{4}Καὶ ἐξοίσω αὐτὸ, λέγει Κύριος παντοκράτωρ, καὶ εἰσελεύσεται εἰς τὸν οἶκον τοῦ κλέπτου, καὶ εἰς τὸν οἶκον τοῦ ὀμνύοντος τῷ ὀνόματί μου ἐπὶ ψεύδει, καὶ καταλύσει ἐν μέσῳ τοῦ οἴκου αὐτοῦ, καὶ συντελέσει αὐτὸν, καὶ τὰ ξύλα αὐτοῦ, καὶ τοὺς λίθους αὐτοῦ.

\vs{5}Καὶ ἐξῆλθεν ὁ ἄγγελος ὁ λαλῶν ἐν ἐμοὶ, καὶ εἶπε πρὸς μὲ, ἀνάβλεψον τοῖς ὀφθαλμοῖς σου, καὶ ἴδε τὸ ἐκπορευόμενον τοῦτο.
\vs{6}Καὶ εἶπα, τί ἐστι; καὶ εἶπε, τοῦτο τὸ μέτρον τὸ ἐκπορευόμενον· καὶ εἶπεν, αὕτη ἡ ἀδικία αὐτῶν ἐν πάσῃ τῇ γῇ.
\vs{7}Καὶ ἰδοὺ τάλαντον μολίβδου ἐξαιρόμενον· καὶ ἰδοὺ γυνὴ μία ἐκάθητο ἐν μέσῳ τοῦ μέτρου.
\vs{8}Καὶ εἶπεν, αὕτη ἐστὶν ἡ ἀνομία· καὶ ἔῤῥιψεν αὐτὴν εἰς μέσον τοῦ μέτρου, καὶ ἔῤῥιψε τὸν λίθον τοῦ μολίβδου εἰς τὸ στόμα αὐτῆς.
\vs{9}Καὶ ᾖρα τοὺς ὀφθαλμούς μου, καὶ ἴδον, καὶ ἰδοὺ δύο γυναῖκες ἐκπορευόμεναι, καὶ πνεῦμα ἐν ταῖς πτέρυξιν αὐτῶν, καὶ αὗται εἶχον πτέρυγας ἔποπος· καὶ ἀνέλαβον τὸ μέτρον ἀναμέσον τῆς γῆς, καὶ ἀναμέσον τοῦ οὐρανοῦ.
\vs{10}Καὶ εἶπα πρὸς τὸν ἄγγελον τὸν λαλοῦντα ἐν ἐμοὶ, ποῦ αὗται ἀποφέρουσι τὸ μέτρον;
\vs{11}Καὶ εἶπε πρὸς μὲ, οἰκοδομῆσαι αὐτῷ οἰκίαν ἐν γῇ Βαβυλῶνος, καὶ ἑτοιμάσαι, καὶ θήσουσιν αὐτὸ ἐκεῖ ἐπὶ τὴν ἑτοιμασίαν αὐτοῦ.

\ch{6}
Καὶ ἐπέστρεψα, καὶ ᾖρα τοὺς ὀφθαλμούς μου, καὶ ἴδον, καὶ ἰδοὺ τέσσαρα ἅρματα ἐκπορευόμενα ἐκ μέσου δύο ὀρέων, καὶ τὰ ὄρη ἦν ὄρη χαλκᾶ.
\vs{2}Ἐν τῷ ἅρματι τῷ πρώτῳ ἵπποι πυῤῥοὶ, καὶ ἐν τῷ ἅρματι τῷ δευτέρῳ ἵπποι μέλανες,
\vs{3}καὶ ἐν τῷ ἅρματι τῷ τρίτῳ ἵπποι λευκοὶ, καὶ ἐν τῷ ἅρματι τῷ τετάρτῳ ἵπποι ποικίλοι ψαροί.
\vs{4}Καὶ ἀπεκρίθην, καὶ εἶπα πρὸς τὸν ἄγγελον τὸν λαλοῦντα ἐν ἐμοὶ, τί ἐστι ταῦτα κύριε;

\vs{5}Καὶ ἀπεκρίθη ὁ ἄγγελος ὁ λαλῶν ἐν ἐμοὶ, καὶ εἶπε, ταῦτά ἐστιν οἱ τέσσαρες ἄνεμοι τοῦ οὐρανοῦ, ἐκπορεύονται παραστῆναι τῷ Κυρίῳ πάσης τῆς γῆς.
\vs{6}Ἐν ᾧ ἦσαν ἵπποι οἱ μέλανες, ἐξεπορεύοντο ἐπὶ γῆν Βοῤῥᾶ, καὶ οἱ λευκοὶ ἐξεπορεύοντο κατόπισθεν αὐτῶν, καὶ οἱ ποικίλοι ἐξεπορεύοντο ἐπὶ γῆν Νότου,
\vs{7}καὶ οἱ ψαροὶ ἐξεπορεύοντο, καὶ ἐπέβλεπον τοῦ πορεύεσθαι τοῦ περιοδεῦσαι τὴν γῆν, καὶ εἶπε, πορεύεσθε, καὶ περιοδεύσατε τὴν γῆν· καὶ περιώδευσαν τὴν γῆν.

\vs{8}Καὶ ἀνεβόησε, καὶ ἐλάλησε πρὸς μὲ, λέγων, ἰδοὺ οἱ ἐκπορευόμενοι ἐπὶ γῆν Βοῤῥᾶ, καὶ ἀνέπαυσαν τὸν θυμόν μου ἐν γῇ βοῤῥᾶ.

\vs{9}Καὶ ἐγένετο λόγος Κυρίου πρὸς μὲ, λέγων,
\vs{10}λάβε τὰ ἐκ τῆς αἰχμαλωσίας παρὰ τῶν ἀρχόντων, καὶ παρὰ τῶν χρησίμων αὐτῆς, καὶ παρὰ τῶν ἐπεγνωκότων αὐτὴν, καὶ εἰσελεύσῃ σὺ ἐν τῇ ἡμέρᾳ ἐκείνῃ εἰς τὸν οἶκον Ἰωσίου τοῦ Σοφονίου τοῦ ἥκοντος ἐκ Βαβυλῶνος,
\vs{11}καὶ λήψῃ ἀργύριον καὶ χρυσίον, καὶ ποιήσεις στεφάνους, καὶ ἐπιθήσεις ἐπὶ τὴν κεφαλὴν Ἰησοῦ τοῦ Ἰωσεδὲκ τοῦ ἱερέως τοῦ μεγάλου.
\vs{12}Καὶ ἐρεῖς πρὸς αὐτὸν, τάδε λέγει Κύριος παντοκράτωρ,

Ἰδοὺ ἀνὴρ, Ἀνατολὴ ὄνομα αὐτῷ· καὶ ὑποκάτωθεν αὐτοῦ ἀνατελεῖ, καὶ οἰκοδομήσει τὸν οἶκον Κυρίου,
\vs{13}καὶ αὐτὸς λήψεται ἀρετὴν, καὶ καθιεῖται, καὶ κατάρξει ἐπὶ τοῦ θρόνου αὐτοῦ, καὶ ἔσται ἱερεὺς ἐκ δεξιῶν αὐτοῦ, καὶ βουλὴ εἰρηνικὴ ἔσται ἀναμέσον ἀμφοτέρων·

\vs{14}Ὁ δὲ στέφανος ἔσται τοῖς ὑπομένουσι, καὶ τοῖς χρησίμοις αὐτῆς, καὶ τοῖς ἐπεγνωκόσιν αὐτὴν, καὶ εἰς χάριτα υἱοῦ Σοφονίου, καὶ εἰς ψαλμὸν ἐν οἴκῳ Κυρίου.
\vs{15}Καὶ οἱ μακρὰν ἀπʼ αὐτῶν ἥξουσι, καὶ οἰκοδομήσουσιν ἐν τῷ οἴκῳ Κυρίου, καὶ γνώσεσθε διότι Κύριος παντοκράτωρ ἀπέσταλκέ με πρὸς ὑμᾶς· καὶ ἔσται, ἐὰν εἰσακούοντες εἰσακούσητε τῆς φωνῆς Κυρίου τοῦ Θεοῦ ὑμῶν.

\ch{7}
Καὶ ἐγένετο ἐν τῷ τετάρτῳ ἔτει ἐπὶ Δαρείου τοῦ βασιλέως, ἐγένετο λόγος Κυρίου πρὸς Ζαχαρίαν τετράδι τοῦ μηνὸς τοῦ ἐννάτου, ὅς ἐστι Χασελεύ.
\vs{2}Καὶ ἐξαπέστειλεν εἰς Βαιθὴλ, Σαρασὰρ καὶ Ἀρβεσεὲρ ὁ βασιλεὺς, καὶ οἱ ἄνδρες αὐτοῦ, καὶ ἐξιλάσασθαι τὸν Κύριον,
\vs{3}λέγων πρὸς τοὺς ἱερεῖς τοὺς ἐν τῷ οἴκῳ Κυρίου παντοκράτορος, καὶ πρὸς τοὺς προφήτας, λέγων, εἰσελήλυθεν ὧδε ἐν τῷ μηνὶ τῷ πέμπτῳ τὸ ἁγίασμα, καθότι ἐποίησεν ἤδη ἱκανὰ ἔτη.

\vs{4}Καὶ ἐγένετο λόγος Κυρίου τῶν δυνάμεων πρὸς ἐμὲ, λέγων,
\vs{5}εἰπὸν πρὸς ἅπαντα τὸν λαὸν τῆς γῆς, καὶ πρὸς τοὺς ἱερεῖς, λέγων, ἐὰν νηστεύσητε ἢ κόψησθε ἐν ταῖς πέμπταις ἢ ἐν ταῖς ἑβδόμαις, καὶ ἰδοὺ ἑβδομήκοντα ἔτη, μὴ νηστείαν νενηστεύκατέ μοι;
\vs{6}Καὶ ἐὰν φάγητε ἢ πίητε, οὐκ ὑμεῖς ἔσθετε καὶ πίνετε;
\vs{7}Οὐχ οὗτοι οἱ λόγοι, οὓς ἐλάλησε Κύριος ἐν χερσὶ τῶν προφητῶν τῶν ἔμπροσθεν, ὅτε ἦν Ἱερουσαλὴμ κατοικουμένη, καὶ εὐθηνοῦσα, καὶ αἱ πόλεις κυκλόθεν αὐτῆς, καὶ ἡ ὀρεινὴ καὶ ἡ πεδινὴ κατῳκεῖτο;

\vs{8}Καὶ ἐγένετο λόγος Κυρίου πρὸς Ζαχαρίαν, λέγων,
\vs{9}τάδε λέγει Κύριος παντοκράτωρ,

Κρίμα δίκαιον κρίνετε, καὶ ἔλεος καὶ οἰκτιρμὸν ποιεῖτε ἕκαστος πρὸς τὸν ἀδελφὸν αὐτοῦ,
\vs{10}καὶ χήραν, καὶ ὀρφανὸν, καὶ προσήλυτον, καὶ πένητα μὴ καταδυναστεύετε, καὶ κακίαν ἕκαστος τοῦ ἀδελφοῦ αὐτοῦ μὴ μνησικακείτω ἐν ταῖς καρδίαις ὑμῶν.

\vs{11}Καὶ ἠπείθησαν τοῦ προσέχειν, καὶ ἔδωκαν νῶτον παραφρονοῦντα, καὶ τὰ ὦτα αὐτῶν ἐβάρυναν τοῦ μὴ εἰσακούειν.
\vs{12}Καὶ τὴν καρδίαν αὐτῶν ἔταξαν ἀπειθῆ τοῦ μὴ εἰσακούειν τοῦ νόμου μου, καὶ τοὺς λόγους, οὓς ἐξαπέστειλε Κύριος παντοκράτωρ ἐν πνεύματι αὐτοῦ ἐν χερσὶ τῶν προφητῶν τῶν ἔμπροσθεν· καὶ ἐγένετο ὀργὴ μεγάλη παρὰ Κυρίου παντοκράτορος.
\vs{13}Καὶ ἔσται, ὃν τρόπον εἶπε, καὶ οὐκ εἰσήκουσαν, οὕτως κεκράξονται, καὶ οὐ μὴ εἰσακούσω, λέγει Κύριος παντοκράτωρ.
\vs{14}Καὶ ἐκβαλῶ αὐτοὺς εἰς πάντα τὰ ἔθνη, ἃ οὐκ ἔγνωσαν· καὶ ἡ γῆ ἀφανισθήσεται κατόπισθεν αὐτῶν ἐκ διοδεύοντος καὶ ἐξ ἀναστρέφοντος· καὶ ἔταξαν γῆν ἐκλεκτὴν εἰς ἀφανισμόν.

\ch{8}
Καὶ ἐγένετο λόγος Κυρίου παντοκράτορος, λέγων,
\vs{2}τάδε λέγει Κύριος παντοκράτωρ, ἐζήλωκα τὴν Ἱερουσαλὴμ, καὶ τὴν Σιὼν ζῆλον μέγαν, καὶ θυμῷ μεγάλῳ ἐζήλωκα αὐτήν.

\vs{3}Τάδε λέγει Κύριος, ἐπιστρέψω ἐπὶ Σιὼν, καὶ κατασκηνώσω ἐν μέσῳ Ἱερουσαλὴμ, καὶ κληθήσεται ἡ Ἱερουσαλὴμ πόλις ἀληθινὴ, καὶ τὸ ὄρος Κυρίου παντοκράτορος, ὄρος ἅγιον.

\vs{4}Τάδε λέγει Κύριος παντοκράτωρ, ἔτι καθήσονται πρεσβύτεροι καὶ πρεσβύτεραι ἐν ταῖς πλατείαις Ἱερουσαλὴμ, ἕκαστος τὴν ῥάβδον αὐτοῦ ἔχων ἐν τῇ χειρὶ αὐτοῦ, ἀπὸ πλήθους ἡμερῶν.
\vs{5}Καὶ αἱ πλατεῖαι τῆς πόλεως πλησθήσονται παιδαρίων καὶ κορασίων παιζόντων ἐν ταῖς πλατείαις αὐτῆς.

\vs{6}Τάδε λέγει Κύριος παντοκράτωρ, εἰ ἀδυνατήσει ἐνώπιον τῶν καταλοίπων τοῦ λαοῦ τούτου ἐν ταῖς ἡμέραις ἐκείναις, μὴ καὶ ἐνώπιόν μου ἀδυνατήσει; λέγει Κύριος παντοκράτωρ.

\vs{7}Τάδε λέγει Κύριος παντοκράτωρ, ἰδοὺ ἐγὼ σώζω τὸν λαόν μου ἀπὸ γῆς ἀνατολῶν καὶ ἀπὸ γῆς δυσμῶν,
\vs{8}καὶ εἰσάξω αὐτοὺς, καὶ κατασκηνώσω ἐν μέσῳ Ἱερουσαλὴμ, καὶ ἔσονται ἐμοὶ εἰς λαὸν, κᾀγὼ ἔσομαι αὐτοῖς εἰς Θεὸν ἐν ἀληθείᾳ καὶ ἐν δικαιοσύνῃ.

\vs{9}Τάδε λέγει Κύριος παντοκράτωρ, κατισχυέτωσαν αἱ χεῖρες ὑμῶν τῶν ἀκουόντων ἐν ταῖς ἡμέραις ταύταις τοὺς λόγους τούτους ἐκ στόματος τῶν προφητῶν, ἀφʼ ἧς ἡμέρας τεθεμελίωται ὁ οἶκος Κυρίου παντοκράτορος, καὶ ὁ ναὸς ἀφʼ οὗ ᾠκοδόμηται.
\vs{10}Διότι πρὸ τῶν ἡμερῶν ἐκείνων ὁ μισθὸς τῶν ἀνθρώπων οὐκ ἔσται εἰς ὄνησιν, καὶ ὁ μισθὸς τῶν κτηνῶν οὐχ ὑπάρξει, καὶ τῷ ἐκπορευομένῳ καὶ τῷ εἰσπορευομένῳ οὐκ ἔσται εἰρήνη ἀπὸ τῆς θλίψεως· καὶ ἐξαποστελῶ πάντας τοὺς ἀνθρώπους, ἕκαστον ἐπὶ τὸν πλησίον αὐτοῦ.
\vs{11}Καὶ νῦν, οὐ κατὰ τὰς ἡμέρας τὰς ἔμπροσθεν ἐγὼ ποιῶ τοῖς καταλοίποις τοῦ λαοῦ τούτου, λέγει Κύριος παντοκράτωρ,
\vs{12}ἀλλʼ ἢ δείξω εἰρήνην· ἡ ἄμπελος δώσει τὸν καρπὸν αὐτῆς, καὶ ἡ γῆ δώσει τὰ γεννήματα αὐτῆς, καὶ ὁ οὐρανὸς δώσει τὴν δρόσον αὐτοῦ, καὶ κατακληρονομήσω τοῖς καταλοίποις τοῦ λαοῦ μου τούτου ταῦτα πάντα.
\vs{13}Καὶ ἔσται ὃν τρόπον ἦτε ἐν κατάρᾳ ἐν τοῖς ἔθνεσιν ὁ οἶκος Ἰούδα καὶ οἶκος Ἰσραὴλ, οὕτως διασώσω ὑμᾶς, καὶ ἔσεσθε ἐν εὐλογίᾳ· θαρσεῖτε, καὶ κατισχύετε ἐν ταῖς χερσὶν ὑμῶν.

\vs{14}Διότι τάδε λέγει Κύριος παντοκράτωρ, ὃν τρόπον διενοήθην τοῦ κακῶσαι ὑμᾶς ἐν τῷ παροργίσαι με τοὺς πατέρας ὑμῶν, λέγει Κύριος παντοκράτωρ, καὶ οὐ μετενόησα·
\vs{15}οὕτως παρατέταγμαι, καὶ διανενόημαι ἐν ταῖς ἡμέραις ταύταις, τοῦ καλῶς ποιῆσαι τὴν Ἱερουσαλὴμ, καὶ τὸν οἶκον Ἰούδα· θαρσεῖτε.
\vs{16}Οὗτοι οἱ λόγοι οὓς ποιήσετε· λαλεῖτε ἀλήθειαν ἕκαστος πρὸς τὸν πλησίον αὐτοῦ, ἀλήθειαν καὶ κρίμα εἰρηνικὸν κρίνατε ἐν ταῖς πύλαις ὑμῶν,
\vs{17}καὶ ἔκαστος τὴν κακίαν τοῦ πλησίον αὐτοῦ μὴ λογίζεσθε ἐν ταῖς καρδίαις ὑμῶν, καὶ ὅρκον ψευδῆ μὴ ἀγαπᾶτε· διότι ταῦτα πάντα ἐμίσησα, λέγει Κύριος παντοκράτωρ.

\vs{18}Καὶ ἐγένετο λόγος Κυρίου παντοκράτορος πρὸς μὲ, λέγων,
\vs{19}Τάδε λέγει Κύριος παντοκράτωρ, νηστεία ἡ τετρὰς, καὶ νηστεία ἡ πέμπτη, καὶ νηστεία ἡ ἑβδόμη, καὶ νηστεία ἡ δεκάτη ἔσονται τῷ οἴκῳ Ἰούδα εἰς χαρὰν καὶ εὐφροσύνην, καὶ εἰς ἑορτὰς ἀγαθάς· καὶ εὐφρανθήσεσθε, καὶ τὴν ἀλήθειαν καὶ τὴν εἰρήνην ἀγαπήσατε.

\vs{20}Τάδε λέγει Κύριος παντοκράτωρ, ἔτι ἥξουσι λαοὶ πολλοὶ, καὶ κατοικοῦντες πόλεις πολλὰς,
\vs{21}καὶ συνελεύσονται κατοικοῦντες πέντε πόλεις εἰς μίαν πόλιν, λέγοντες, πορευθῶμεν δεηθῆναι τοῦ προσώπου Κυρίου, καὶ ἐκζητῆσαι τὸ πρόσωπον Κυρίου παντοκράτορος· πορεύσομαι κᾀγώ·
\vs{22}Καὶ ἥξουσι λαοὶ πολλοὶ καὶ ἔθνη πολλὰ ἐκζητῆσαι τὸ πρόσωπον Κυρίου παντοκράτορος ἐν Ἱερουσαλὴμ, καὶ ἐξιλάσασθαι τὸ πρόσωπον Κυρίου.

\vs{23}Τάδε λέγει Κύριος παντοκράτωρ, ἐν ταῖς ἡμέραις ἐκείναις, ἐὰν ἐπιλάβωνται δέκα ἄνδρες ἐκ πασῶν τῶν γλῶσσων τῶν ἐθνῶν, καὶ ἐπιλάβωνται τοῦ κρασπέδου ἀνδρὸς Ἰουδαίου, λέγοντες, πορευσόμεθα μετὰ σοῦ, διότι ἀκηκόαμεν ὅτι ὁ Θεὸς μεθʼ ὑμῶν ἐστι.

\ch{9}
Λῆμμα λόγου Κυρίου ἐν γῇ Σεδρὰχ, καὶ Δαμασκοῦ θυσία αὐτοῦ, διότι Κύριος ἐφορᾷ ἀνθρώπους, καὶ πάσας φυλὰς τοῦ Ἰσραήλ.
\vs{2}Καὶ ἐν Ἠμὰθ ἐν τοῖς ὁρίοις αὐτῆς Τύρος καὶ Σιδὼν, διότι ἐφρόνησαν σφόδρα.
\vs{3}Καὶ ᾠκοδόμησε Τύρος ὀχυρώματα αὐτῇ, καὶ ἐθησαύρισεν ἀργύριον ὡς χοῦν, καὶ συνήγαγε χρυσίον ὡς πηλὸν ὁδῶν.

\vs{4}Καὶ διατοῦτο Κύριος κληρονομήσει αὐτοὺς, καὶ πατάξει εἰς θάλασσαν δύναμιν αὐτῆς, καὶ αὕτη ἐν πυρὶ καταναλωθήσεται.
\vs{5}Ὄψεται Ἀσκάλων, καὶ φοβηθήσεται, καὶ Γάζα, καὶ ὀδυνηθήσεται σφόδρα, καὶ Ἀκκάρων, ὅτι ᾐσχύνθη ἐπὶ τῷ παραπτώματι αὐτῆς· καὶ ἀπολεῖται βασιλεὺς ἐκ Γάζης, καὶ Ἀσκάλων οὐ μὴ κατοικηθῇ.
\vs{6}Καὶ κατοικήσουσιν ἀλλογενεῖς ἐν Ἀζώτῳ, καὶ καθελῶ ὕβριν ἀλλοφύλων,
\vs{7}καὶ ἐξαρῶ τὸ αἷμα αὐτῶν ἐκ στόματος αὐτῶν, καὶ τὰ βδελύγματα αὐτῶν ἐκ μέσου ὀδόντων αὐτῶν· καὶ ὑπολειφθήσονται καὶ οὗτοι τῷ Θεῷ ἡμῶν, καὶ ἔσονται ὡς χιλίαρχος ἐν Ἰούδᾳ, καὶ Ἀκκάρων ὡς ὁ Ἰεβουσαῖος.
\vs{8}Καὶ ὑποστήσομαι τῷ οἴκῳ μου ἀνάστημα, τοῦ μὴ διαπορεύεσθαι, μηδὲ ἀνακάμπτειν, καὶ οὐ μὴ ἐπέλθῃ ἐπʼ αὐτοὺς οὐκέτι ἐξελαύνων, διότι νῦν ἑώρακα ἐν τοῖς ὀφθαλμοῖς μου.

\vs{9}Χαῖρε σφόδρα θύγατερ Σιὼν, κήρυσσε θύγατερ Ἱερουσαλήμ· ἰδοὺ ὁ βασιλεὺς ἔρχεταί σοι δίκαιος καὶ σώζων, αὐτὸς πραῢς, καὶ ἐπιβεβηκὼς ἐπὶ ὑποζύγιον καὶ πῶλον νέον.
\vs{10}Καὶ ἐξολοθρεύσει ἅρματα ἐξ Ἐφραὶμ, καὶ ἵππον ἐξ Ἱερουσαλὴμ, καὶ ἐξολοθρεύσεται τόξον πολεμικὸν, καὶ πλῆθος καὶ εἰρήνη ἐξ ἐθνῶν, καὶ κατάρξει ὑδάτων ἕως θαλάσσης, καὶ ποταμῶν διεκβολὰς γῆς.

\vs{11}Καὶ σὺ ἐν αἵματι διαθήκης σου ἐξαπέστειλας δεσμίους σου ἐκ λάκκου οὐκ ἔχοντος ὕδωρ.
\vs{12}Καθήσεσθε ἐν ὀχυρώμασι δέσμιοι τῆς συναγωγῆς, καὶ ἀντὶ μιᾶς ἡμέρας παροικεσίας σου διπλᾶ ἀνταποδώσω σοι,
\vs{13}διότι ἐνέτεινά σε Ἰούδα ἐμαυτῷ τόξον, ἔπλησα τὸν Ἐφραὶμ, καὶ ἐξεγερῶ τὰ τέκνα σου Σιὼν ἐπὶ τὰ τέκνα τῶν Ἑλλήνων, καὶ ψηλαφήσω σε ὡς ῥομφαίαν μαχητοῦ,
\vs{14}καὶ Κύριος ἔσται ἐπʼ αὐτοὺς, καὶ ἐξελεύσεται ὡς ἀστραπὴ βολὶς, καὶ Κύριος παντοκράτωρ ἐν σάλπιγγι σαλπιεῖ, καὶ πορεύσεται ἐν σάλῳ ἀπειλῆς αὐτοῦ.
\vs{15}Κύριος παντοκράτωρ ὑπερασπιεῖ αὐτούς· καὶ καταναλώσουσιν αὐτοὺς, καὶ καταχώσουσιν αὐτοὺς ἐν λίθοις σφενδόνης, καὶ ἐκπίονται αὐτοὺς ὡς οἶνον, καὶ πλήσουσι τὰς φιάλας ὡς θυσιαστήριον.
\vs{16}Καὶ σώσει αὐτοὺς Κύριος ὁ Θεὸς αὐτῶν ἐν τῇ ἡμέρᾳ ἐκείνῃ, ὡς πρόβατα λαὸν αὐτοῦ, διότι λίθοι ἅγιοι κυλίονται ἐπὶ γῆς αὐτοῦ.
\vs{17}Ὅτι εἴ τι ἀγαθὸν αὐτοῦ, καὶ εἴ τι καλὸν αὐτοῦ, σῖτος νεανίσκοις, καὶ οἶνος εὐωδιάζων εἰς παρθένους.

\ch{10}
Αἰτεῖσθε παρὰ Κυρίου ὑετὸν καθʼ ὥραν, πρώϊμον καὶ ὄψιμον· Κύριος ἐποίησε φαντασίας, καὶ ὑετὸν χειμερινὸν δώσει αὐτοῖς, ἑκάστῳ βοτάνην ἐν ἀγρῷ.
\vs{2}Διότι οἱ ἀποφθεγγόμενοι ἐλάλησαν κόπους, καὶ οἱ μάντεις ὁράσεις ψευδεῖς, καὶ τὰ ἐνύπνια ψευδῆ ἐλάλουν, μάταια παρεκάλουν· διατοῦτο ἐξηράνθησαν ὡς πρόβατα, καὶ ἐκακώθησαν, διότι οὐκ ἦν ἴασις.

\vs{3}Ἐπὶ τοὺς ποιμένας παρωξύνθη ὁ θυμός μου, καὶ ἐπὶ τοὺς ἀμνοὺς ἐπισκέψομαι· καὶ ἐπισκέψεται Κύριος ὁ Θεὸς ὁ παντοκράτωρ τὸ ποίμνιον αὐτοῦ, τὸν οἶκον Ἰούδα, καὶ τάξει αὐτοὺς ὡς ἵππον εὐπρεπῆ αὐτοῦ ἐν πολέμῳ,
\vs{4}καὶ ἀπʼ αὐτοῦ ἐπέβλεψε, καὶ ἀπʼ αὐτοῦ ἔταξε, καὶ ἀπʼ αὐτοῦ τόξον ἐν θυμῷ, ἀπʼ αὐτοῦ ἐξελεύσεται πᾶς ὁ ἐξελαύνων ἐν τῷ αὐτῷ.
\vs{5}Καὶ ἔσονται ὡς μαχηταὶ πατοῦντες πηλὸν ἐν ταῖς ὁδοῖς ἐν πολέμῳ, καὶ παρατάξονται, διότι Κύριος μετʼ αὐτῶν· καὶ καταισχυνθήσονται ἀναβάται ἵππων.

\vs{6}Καὶ κατισχύσω τὸν οἶκον Ἰούδα, καὶ τὸν οἶκον Ἰωσὴφ σώσω, καὶ κατοικιῶ αὐτοὺς, ὅτι ἠγάπησα αὐτοὺς, καὶ ἔσονται, ὃν τρόπον οὐκ ἀπεστρεψάμην αὐτούς· διότι ἐγὼ Κύριος ὁ Θεὸς αὐτῶν· καὶ ἐπακούσομαι αὐτοῖς,
\vs{7}καὶ ἔσονται ὡς μαχηταὶ τοῦ Ἐφραὶμ, καὶ χαρήσεται ἡ καρδία αὐτῶν ὡς ἐν οἴνῳ· καὶ τὰ τέκνα αὐτῶν ὄψονται, καὶ εὐφρανθήσονται, καὶ χαρεῖται ἡ καρδία αὐτῶν ἐπὶ τῷ Κυρίῳ.
\vs{8}Σημανῶ αὐτοῖς, καὶ εἰσδέξομαι αὐτοὺς, διότι λυτρώσομαι αὐτοὺς, καὶ πληθυνθήσονται καθότι ἦσαν πολλοί.

\vs{9}Καὶ σπερῶ αὐτοὺς ἐν λαοῖς, καὶ οἱ μακρὰν μνησθήσονταί μου, ἐκθρέψουσι τὰ τέκνα αὐτῶν, καὶ ἐπιστρέψουσι.
\vs{10}Καὶ ἐπιστρέψω αὐτοὺς ἐκ γῆς Αἰγύπτου, καὶ ἐξ Ἀσσυρίων εἰσδέξομαι αὐτοὺς, καὶ εἰς τὴν Γαλααδίτιν, καὶ εἰς τὸν Λίβανον εἰσάξω αὐτοὺς, καὶ οὐ μὴ ὑπολειφθῇ ἐξ αὐτῶν οὐδὲ εἷς.
\vs{11}Καὶ διελεύσονται ἐν θαλάσσῃ στενῇ, πατάξουσιν ἐν θαλάσσῃ κύματα, καὶ ξηρανθήσεται πάντα τὰ βάθη ποταμῶν, καὶ ἀφαιρεθήσεται πᾶσα ὕβρις Ἀσσυρίων, καὶ σκῆπτρον Αἰγύπτου περιαιρεθήσεται.
\vs{12}Καὶ κατισχύσω αὐτοὺς ἐν Κυρίῳ Θεῷ αὐτῶν, καὶ ἐν τῷ ὀνόματι αὐτοῦ κατακαυχήσονται, λέγει Κύριος.

\ch{11}
Διάνοιξον ὁ Λίβανος τὰς θύρας σου, καὶ καταφαγέτω πῦρ τὰς κέδρους σου.
\vs{2}Ὀλολυξάτω πίτυς, διότι πέπτωκε κέδρος, ὅτι μεγάλως μεγιστᾶνες ἐταλαιπώρησαν· ὀλολύξατε δρύες τῆς Βασανίτιδος, ὅτι κατεσπάσθη ὁ δρυμὸς ὁ σύμφυτος.

\vs{3}Φωνὴ θρηνούντων ποιμένων, ὅτι τεταλαιπώρηκεν ἡ μεγαλωσύνη αὐτῶν· φωνὴ ὠρυομένων λεόντων, ὅτι τεταλαιπώρηκε τὸ φρύαγμα τοῦ Ἰορδάνου.

\vs{4}Τάδε λέγει Κύριος παντοκράτωρ, ποιμαίνετε τὰ πρόβατα τῆς σφαγῆς,
\vs{5}ἃ οἱ κτησάμενοι κατέσφαζον, καὶ οὐ μετεμέλοντο, καὶ οἱ πωλοῦντες αὐτὰ ἔλεγον, εὐλογητὸς Κύριος, καὶ πεπλουτήκαμεν· καὶ οἱ ποιμένες αὐτῶν οὐκ ἔπασχον οὐδὲν ἐπʼ αὐτοῖς.
\vs{6}Διατοῦτο οὐ φείσομαι οὐκέτι ἐπὶ τοὺς κατοικοῦντας τὴν γῆν, λέγει Κύριος· καὶ ἰδοὺ ἐγὼ παραδίδωμι τοὺς ἀνθρώπους, ἕκαστον εἰς χεῖρα τοῦ πλησίον αὐτοῦ, καὶ εἰς χεῖρα βασιλέως αὐτοῦ, καὶ κατακόψουσι τὴν γῆν, καὶ οὐ μὴ ἐξέλωμαι ἐκ χειρὸς αὐτῶν.

\vs{7}Καὶ ποιμανῶ τὰ πρόβατα τῆς σφαγῆς εἰς τὴν Χαναανίτιν· καὶ λήψομαι ἐμαυτῷ δύο ῥάβδους, τὴν μὲν μίαν ἐκάλεσα Κάλλος, καὶ τὴν ἑτέραν ἐκάλεσα Σχοίνισμα, καὶ ποιμανῶ τὰ πρόβατα.
\vs{8}Καὶ ἐξαρῶ τοὺς τρεῖς ποιμένας ἐν μηνὶ ἑνὶ, καὶ βαρυνθήσεται ἡ ψυχή μου ἐπʼ αὐτοὺς, καὶ γὰρ αἱ ψυχαὶ αὐτῶν ἐπωρύοντο ἐπʼ ἐμέ.
\vs{9}Καὶ εἶπα, οὐ ποιμανῶ ὑμᾶς· τὸ ἀποθνῆσκον ἀποθνησκέτω, καὶ τὸ ἐκλεῖπον ἐκλιπέτω, καὶ τὰ κατάλοιπα κατεσθιέτωσαν ἕκαστος τὰς σάρκας τοῦ πλησίον αὐτοῦ.

\vs{10}Καὶ λήψομαι τὴν ῥάβδον μου τὴν καλὴν, καὶ ἀποῤῥίψω αὐτὴν, τοῦ διασκεδάσαι τὴν διαθήκην μου, ἣν διεθέμην πρὸς πάντας τοὺς λαοὺς,
\vs{11}καὶ διασκεδασθήσεται ἐν τῇ ἡμέρᾳ ἐκείνῃ, καὶ γνώσονται οἱ Χαναναῖοι τὰ πρόβατα τὰ φυλασσόμενά μοι, διότι λόγος Κυρίου ἐστί.
\vs{12}Καὶ ἐρῶ πρὸς αὐτοὺς, εἰ καλὸν ἐνώπιον ὑμῶν ἐστι, δότε τὸν μισθόν μου, ἢ ἀπείπασθε· καὶ ἔστησαν τὸν μισθόν μου τριακοντα ἀργυροῦς.
\vs{13}Καὶ εἶπε Κύριος πρὸς μὲ, κάθες αὐτοὺς εἰς τὸ χωνευτήριον, καὶ σκέψομαι εἰ δόκιμόν ἐστιν, ὃν τρόπον ἐδοκιμάσθην ὑπὲρ αὐτῶν, καὶ ἔλαβον τοὺς τριάκοντα ἀργυροῦς, καὶ ἐνέβαλον αὐτοὺς εἰς τὸν οἶκον Κυρίου εἰς τὸ χωνευτήριον.

\vs{14}Καὶ ἀπέῤῥιψα τὴν ῥάβδον τὴν δευτέραν τὸ Σχοίνισμα, τοῦ διασκεδάσαι τὴν κατάσχεσιν ἀναμέσον Ἰούδα, καὶ ἀναμέσον Ἰσραήλ.

\vs{15}Καὶ εἶπε Κύριος πρὸς μὲ, ἔτι λάβε σεαυτῷ σκεύη ποιμενικὰ ποιμένος ἀπείρου·
\vs{16}διότι ἰδοὺ ἐγὼ ἐξεγείρω ποιμένα ἐπὶ τὴν γῆν, τὸ ἐκλιμπάνον οὐ μὴ ἐπισκέψηται, καὶ τὸ ἐσκορπισμένον οὐ μὴ ζητήσῃ, καὶ τὸ συντετριμμένον οὐ μὴ ἰάσηται, καὶ τὸ ὁλόκληρον οὐ μὴ κατευθύνῃ, καὶ τὰ κρέα τῶν ἑκλεκτῶν καταφάγεται, καὶ τοὺς ἀστραγάλους αὐτῶν ἐκστρέψει.

\vs{17}Ὢ οἱ ποιμαίνοντες τὰ μάταια, καταλελοιπότες τὰ πρόβατα, μάχαιρα ἐπὶ τοὺς βραχίονας αὐτοῦ, καὶ ἐπὶ τὸν ὀφθαλμὸν τὸν δεξιὸν αὐτοῦ, ὁ βραχίων αὐτοῦ ξηραινόμενος ξηρανθήσεται, καὶ ὁ ὀφθαλμὸς ὁ δεξιὸς αὐτοῦ ἐκτυφλούμενος ἐκτυφλωθήσεται.

\ch{12}
Λῆμμα λόγου Κυρίου ἐπὶ τὸν Ἰσραήλ· λέγει Κύριος, ἐκτείνων οὐρανὸν, καὶ θεμελιῶν γῆν, καὶ πλάσσων πνεῦμα ἀνθρώπου ἐν αὐτῷ,
\vs{2}ἰδοὺ ἐγὼ τίθημι τὴν Ἱερουσαλὴμ ὡς πρόθυρα σαλευόμενα πᾶσι τοῖς λαοῖς κύκλῳ, καὶ ἐν τῇ Ἰουδαίᾳ ἔσται περιοχὴ ἐπὶ Ἱερουσαλήμ.
\vs{3}Καὶ ἔσται ἐν τῇ ἡμέρᾳ ἐκείνῃ θήσομαι τὴν Ἱερουσαλὴμ λίθον καταπατούμενον πᾶσι τοῖς ἔθνεσι· πᾶς ὁ καταπατῶν αὐτὴν ἐμπαίζων ἐμπαίξεται, καὶ ἐπισυναχθήσονται ἐπʼ αὐτὴν πάντα τὰ ἔθνη τῆς γῆς.
\vs{4}Ἐν τῇ ἡμέρᾳ ἐκείνῃ, λέγει Κύριος παντοκράτωρ, πατάξω πάντα ἵππον ἐν ἐκστάσει, καὶ τὸν ἀναβάτην αὐτοῦ ἐν παραφρονήσει, ἐπὶ δὲ τὸν οἶκον Ἰούδα διανοίξω τοὺς ὀφθαλμούς μου, καὶ πάντας τοὺς ἵππους τῶν λαῶν πατάξω ἐν ἀποτυφλώσει.

\vs{5}Καὶ ἐροῦσιν οἱ χιλίαρχοι Ἰούδα ἐν ταῖς καρδίαις αὐτῶν, εὑρήσομεν ἑαυτοῖς τοὺς κατοικοῦντας Ἱερουσαλὴμ ἐν Κυρίῳ παντοκράτορι Θεῷ αὐτῶν.
\vs{6}Ἐν τῇ ἡμέρᾳ ἐκείνῃ θήσομαι τοὺς χιλιάρχους Ἰούδα ὡς δαλὸν πυρὸς ἐν ξύλοις, καὶ ὡς λαμπάδα πυρὸς ἐν καλάμῃ, καὶ καταφάγονται ἐκ δεξιῶν, καὶ ἐξ εὐωνύμων πάντας τοὺς λαοὺς κυκλόθεν, καὶ κατοικήσει Ἱερουσαλὴμ ἔτι καθʼ ἑαυτὴν ἐν Ἱερουσαλήμ.
\vs{7}Καὶ σώσει Κύριος τὰ σκηνώματα Ἰούδα, καθὼς ἀπʼ ἀρχῆς, ὅπως μὴ μεγαλύνηται καύχημα οἴκου Δαυὶδ, καὶ ἔπαρσις τῶν κατοικούντων Ἱερουσαλὴμ ἐπὶ τὸν Ἰούδα.
\vs{8}Καὶ ἔσται ἐν τῇ ἡμέρᾳ ἐκείνῃ ὑπερασπιεῖ Κύριος ὑπὲρ τῶν κατοικούντων Ἱερουσαλὴμ, καὶ ἔσται ὁ ἀσθενῶν ἐν αὐτοῖς ἐν ἐκείνῃ τῇ ἡμέρᾳ ὡς Δαυὶδ, ὁ δὲ οἶκος Δαυὶδ ὡς οἶκος Θεοῦ, ὡς ἄγγελος Κυρίου ἐνώπιον αὐτῶν.
\vs{9}Καὶ ἔσται ἐν τῇ ἡμέρᾳ ἐκείνῃ, ζητήσω ἐξᾷραι πάντα τὰ ἔθνη τὰ ἐρχόμενα ἐπὶ Ἱερουσαλήμ.
\vs{10}Καὶ ἐκχεῶ ἐπὶ τὸν οἶκον Δαυὶδ, καὶ ἐπὶ τοὺς κατοικοῦντας Ἱερουσαλὴμ πνεῦμα χάριτος καὶ οἰκτιρμοῦ· καὶ ἐπιβλέψονται πρὸς μὲ, ἀνθʼ ὧν κατωρχήσαντο· καὶ κόψονται ἐπʼ αὐτὸν κοπετὸν, ὡς ἐπʼ ἀγαπητῷ, καὶ ὀδυνηθήσονται ὀδύνην, ὡς ἐπὶ τῷ πρωτοτόκῳ.

\vs{11}Ἐν τῇ ἡμέρᾳ ἐκείνῃ μεγαλυνθήσεται ὁ κοπετὸς ἐν Ἱερουσαλὴμ, ὡς κοπετὸς ῥοῶνος ἐν πεδίῳ ἐκκοπτομένου.
\vs{12}Καὶ κόψεται ἡ γῆ κατὰ φυλὰς φυλάς· φυλὴ οἴκου Δαυὶδ καθʼ ἑαυτὴν, καὶ αἱ γυναῖκες αὐτῶν καθʼ ἑαυτάς· φυλὴ οἴκου Νάθαν καθʼ ἑαυτὴν, καὶ αἱ γυναῖκες αὐτῶν καθʼ ἑαυτάς·
\vs{13}φυλὴ οἴκου Λευὶ καθʼ ἑαυτὴν, καὶ αἱ γυναῖκες αὐτῶν καθʼ ἑαυτάς· φυλὴ τοῦ Συμεὼν καθʼ ἑαυτὴν, καὶ αἱ γυναῖκες αὐτῶν καθʼ ἑαυτάς.
\vs{14}Πᾶσαι αἱ ὑπολελειμμέναι φυλαὶ, φυλὴ καθʼ ἑαυτὴν, καὶ γυναῖκες αὐτῶν καθʼ ἑαυτάς.

\ch{13}
Ἐν τῇ ἡμέρᾳ ἐκείνῃ ἔσται πᾶς τόπος διανοιγόμενος τῷ οἴκῳ Δαυὶδ, καὶ τοῖς κατοικοῦσιν Ἱερουσαλὴμ εἰς τὴν μετακίνησιν, καὶ εἰς τὸν χωρισμόν.
\vs{2}Καὶ ἔσται ἐν τῇ ἡμέρᾳ ἐκείνῃ, λέγει Κύριος σαβαὼθ, ἐξολοθρεύσω τὰ ὀνόματα τῶν εἰδώλων ἀπὸ τῆς γῆς, καὶ οὐκ ἔτι αὐτῶν ἔσται μνεία· καὶ τοὺς ψευδοπροφήτας, καὶ τὸ πνεῦμα τὸ ἀκάθαρτον ἐξαρῶ ἀπὸ τῆς γῆς.
\vs{3}Καὶ ἔσται ἐὰν προφητεύσῃ ἄνθρωπος ἔτι, καὶ ἐρεῖ πρὸς αὐτὸν ὁ πατὴρ αὐτοῦ, καὶ ἡ μήτηρ αὐτοῦ, οἱ γεννήσαντες αὐτὸν, οὐ ζήσῃ, ὅτι ψευδῆ ἐλάλησας ἐπʼ ὀνόματι Κυρίου· καὶ συμποδιοῦσιν αὐτὸν ὁ πατὴρ αὐτοῦ, καὶ ἡ μήτηρ αὐτοῦ, οἱ γεννήσαντες αὐτὸν, ἐν τῷ προφητεύειν αὐτόν.

\vs{4}Καὶ ἔσται ἐν τῇ ἡμέρᾳ ἐκείνῃ καταισχυνθήσονται οἱ προφῆται, ἕκαστος ἐκ τῆς ὁράσεως αὐτοῦ, ἐν τῷ προφητεύειν αὐτὸν, καὶ ἐνδύσονται δέῤῥιν τριχίνην, ἀνθʼ ὧν ἐψεύσαντο.
\vs{5}Καὶ ἐρεῖ, οὐκ εἰμὶ προφήτης ἐγὼ, διότι ἄνθρωπος ἐργαζόμενος τὴν γῆν ἐγώ εἰμι, ὅτι ἄνθρωπος ἐγέννησέ με ἐκ νεότητός μου.
\vs{6}Καὶ ἐρῶ πρὸς αὐτὸν, τί αἱ πληγαὶ αὗται ἀναμέσον τῶν χειρῶν σου; καὶ ἐρεῖ, ἃς ἐπλήγην ἐν τῷ οἴκῳ τῷ ἀγαπητῷ μου.

\vs{7}Ῥομφαία ἐξεγέρθητι ἐπὶ τοὺς ποιμένας μου, καὶ ἐπὶ ἄνδρα πολίτην μου, λέγει Κύριος παντοκράτωρ, πατάξατε τοὺς ποιμένας, καὶ ἐκσπάσατε τὰ πρόβατα· καὶ ἐπάξω τὴν χεῖρά μου ἐπὶ τοὺς μικρούς.
\vs{8}Καὶ ἔσται ἐν πάσῃ τῇ γῇ, λέγει Κύριος, τὰ δύο μέρη αὐτῆς ἐξολοθρευθήσεται, καὶ ἐκλείψει, τὸ δὲ τρίτον ὑπολειφθήσεται ἐν αὐτῇ.
\vs{9}Καὶ διάξω τὸ τρίτον διὰ πυρὸς, καὶ πυρώσω αὐτοὺς, ὡς πυροῦται τὸ ἀργύριον, καὶ δοκιμῶ αὐτοὺς, ὡς δοκιμάζεται τὸ χρυσίον· αὐτὸς ἐπικαλέσεται τὸ ὄνομά μου, κἀγὼ ἐπακούσομαι αὐτῷ, καὶ ἐρῶ, λαός μου οὗτός ἐστι· καὶ αὐτὸς ἐρεῖ, Κύριος ὁ Θεός μου.

\ch{14}
Ἰδοὺ ἡμέραι ἔρχονται Κυρίου, καὶ διαμερισθήσονται τὰ σκῦλά σου ἐν σοί·
\vs{2}καὶ ἐπισυνάξω πάντα τὰ ἔθνη ἐπὶ Ἱερουσαλὴμ εἰς πόλεμον, καὶ ἁλώσεται ἡ πόλις, καὶ διαρπαγήσονται αἱ οἰκίαι, καὶ αἱ γυναῖκες μολυνθήσονται, καὶ ἐξελεύσεται τὸ ἥμισυ τῆς πόλεως ἐν αἰχμαλωσίᾳ, οἱ δὲ κατάλοιποι τοῦ λαοῦ μου οὐ μὴ ἐξολοθρευθῶσιν ἐκ τῆς πόλεως.

\vs{3}Καὶ ἐξελεύσεται Κύριος, καὶ παρατάξεται ἐν τοῖς ἔθνεσιν ἐκείνοις, καθὼς ἡμέρα παρατάξεως αὐτοῦ ἐν ἡμέρᾳ πολέμου.
\vs{4}Καὶ στήσονται οἱ πόδες αὐτοῦ ἐν τῇ ἡμέρᾳ ἐκείνῃ ἐπὶ τὸ ὄρος τῶν ἐλαιῶν, τὸ κατέναντι Ἱερουσαλὴμ ἐξ ἀνατολῶν· καὶ σχισθήσεται τὸ ὄρος τῶν ἐλαιῶν, τὸ ἥμισυ αὐτοῦ πρὸς ἀνατολὰς καὶ θάλασσαν, χάος μέγα σφόδρα· καὶ κλινεῖ τὸ ἥμισυ τοῦ ὄρους πρὸς τὸν βοῤῥᾶν, καὶ τὸ ἥμισυ αὐτοῦ πρὸς Νότον.
\vs{5}Καὶ φραχθήσεται ἡ φάραγξ τῶν ὀρέων μου, καὶ ἐγκολληθήσεται φάραγξ ὀρέων ἕως Ἰασὸδ, καὶ ἐνφραχθήσεται καθὼς ἐνεφράγη ἐν ταῖς ἡμέραις τοῦ συσσεισμοῦ, ἐν ἡμέραις Ὀζίου βασιλέως Ἰούδα· καὶ ἥξει Κύριος ὁ Θεός μου, καὶ πάντες οἱ ἅγιοι μετʼ αὐτοῦ.
\vs{6}Καὶ ἔσται ἐν ἐκείνῃ τῇ ἡμέρᾳ οὐκ ἔσται φῶς, καὶ ψύχη καὶ πάγος
\vs{7}ἔσται μίαν ἡμέραν, καὶ ἡ ἡμέρα ἐκείνη γνωστὴ τῷ Κυρίῳ, καὶ οὐκ ἡμέρα, καὶ οὐ νὺξ, καὶ πρὸς ἑσπέραν ἔσται φῶς.

\vs{8}Καὶ ἐν τῇ ἡμέρᾳ ἐκείνῃ ἐξελεύσεται ὕδωρ ζῶν ἐξ Ἱερουσαλὴμ, τὸ ἥμισυ αὐτοῦ εἰς τὴν θάλασσαν τὴν πρώτην, καὶ τὸ ἥμισυ αὐτοῦ εἰς τὴν θάλασσαν τὴν ἐσχάτην· καὶ ἐν θέρει καὶ ἐν ἔαρι ἔσται οὕτως.
\vs{9}Καὶ ἔσται Κύριος εἰς βασιλέα ἐπὶ πᾶσαν τὴν γῆν· ἐν τῇ ἡμέρᾳ ἐκείνῃ ἔσται Κύριος εἷς, καὶ τὸ ὄνομα αὐτοῦ ἓν,
\vs{10}κυκλῶν πᾶσαν τὴν γῆν, καὶ τὴν ἔρημον ἀπὸ Γαβὲ ἕως Ῥεμμὼν κατὰ Νότον Ἱερουσαλήμ. Ῥαμὰ δὲ ἐπὶ τόπου μενεῖ· ἀπὸ τῆς πύλης Βενιαμὶν ἕως τοῦ τόπου τῆς πύλης τῆς πρώτης, ἕως τῆς πύλης τῶν γωνιῶν, καὶ ἕως τοῦ πύργου Ἀναμεὴλ, ἕως τῶν ὑποληνίων τοῦ βασιλέως
\vs{11}κατοικήσουσιν ἐν αὐτῇ, καὶ ἀνάθεμα οὐκ ἔσται ἔτι, καὶ κατοικήσει Ἱερουσαλὴμ πεποιθότως.

\vs{12}Καὶ αὕτη ἔσται ἡ πτῶσις ἣν κόψει Κύριος πάντας τοὺς λαοὺς, ὅσοι ἐπεστράτευσαν ἐπὶ Ἱερουσαλήμ· τακήσονται αἱ σάρκες αὐτῶν, ἑστηκότων ἐπὶ τοὺς πόδας αὐτῶν, καὶ οἱ ὀφθαλμοὶ αὐτῶν ῥυήσονται ἐκ τῶν ὀπῶν αὐτῶν, καὶ ἡ γλῶσσα αὐτῶν τακήσεται ἐν τῷ στόματι αὐτῶν.
\vs{13}Καὶ ἔσται ἐν τῇ ἡμέρᾳ ἐκείνῃ ἔκστασις Κυρίου μεγάλη ἐπʼ αὐτούς· καὶ ἐπιλήψονται ἕκαστος τῆς χειρὸς τοῦ πλησίον αὐτοῦ, καὶ συμπλακήσεται ἡ χεὶρ αὐτοῦ πρὸς τὴν χεῖρα τοῦ πλησίον αὐτοῦ.
\vs{14}Καὶ Ἰούδας παρατάξεται ἐν Ἱερουσαλὴμ, καὶ συνάξει τὴν ἰσχὺν πάντων τῶν λαῶν κυκλόθεν, χρυσίον καὶ ἀργύριον καὶ ἱματισμὸν εἰς πλῆθος σφόδρα.
\vs{15}Καὶ αὕτη ἔσται ἡ πτῶσις τῶν ἵππων, καὶ τῶν ἡμιόνων, καὶ τῶν καμήλων, καὶ τῶν ὄνων, καὶ πάντων τῶν κτηνῶν τῶν ὄντων ἐν ταῖς παρεμβολαῖς ἐκείναις, κατὰ τὴν πτῶσιν ταύτην.

\vs{16}Καὶ ἔσται, ὅσοι ἐὰν καταλειφθῶσιν ἐκ πάντων τῶν ἐθνῶν τῶν ἐλθόντων ἐπʼ Ἱερουσαλὴμ, καὶ ἀναβήσονται κατʼ ἐνιαυτὸν, τοῦ προσκυνῆσαι τῷ βασιλεῖ Κυρίῳ παντοκράτορι, καὶ τοῦ ἑορτάσαι τὴν ἑορτὴν τῆς σκηνοπηγίας.
\vs{17}Καὶ ἔσται, ὅσοι ἐὰν μὴ ἀναβῶσιν ἐκ πασῶν τῶν φυλῶν τῆς γῆς εἰς Ἱερουσαλὴμ, τοῦ προσκυνῆσαι τῷ βασιλεῖ Κυρίῳ παντοκράτορι, καὶ οὗτοι ἐκείνοις προστεθήσονται.
\vs{18}Ἐὰν δὲ φυλὴ Αἰγύπτου μὴ ἀναβῇ, μηδὲ ἔλθῃ, καὶ ἐπὶ τούτους ἔσται ἡ πτῶσις, ἣν πατάξει Κύριος πάντα τὰ ἔθνη, ὅσα ἂν μὴ ἀναβῇ, τοῦ ἑορτάσαι τὴν ἑορτὴν τῆς σκηνοπηγίας.
\vs{19}Αὕτη ἔσται ἡ ἁμαρτία Αἰγύπτου, καὶ ἡ ἁμαρτία πάντων τῶν ἐθνῶν, ὃς ἂν μὴ ἀναβῇ ἑορτάσαι τὴν ἑορτὴν τῆς σκηνοπηγίας.

\vs{20}Ἐν τῇ ἡμέρᾳ ἐκείνῃ ἔσται τὸ ἐπὶ τὸν χαλινὸν τοῦ ἵππου ἅγιον τῷ Κυρίῳ παντοκράτορι· καὶ ἔσονται οἱ λέβητες ἐν τῷ οἴκῳ Κυρίου ὡς φιάλαι πρὸ προσώπου τοῦ θυσιαστηρίου.
\vs{21}Καὶ ἔσται πᾶς λέβης ἐν Ἱερουσαλὴμ καὶ ἐν τῷ Ἰούδᾳ ἅγιος τῷ Κυρίῳ παντοκράτορι· καὶ ἥξουσι πάντες οἱ θυσιάζοντες, καὶ λήψονται ἐξ αὐτῶν, καὶ ἑψήσουσιν ἐν αὐτοῖς· καὶ οὐκ ἔσται Χαναναῖος ἔτι ἐν τῷ οἴκῳ Κυρίου παντοκράτορος ἐν τῇ ἡμέρᾳ ἐκείνῃ.


\def\book{ΜΑΛΑΧΙΑΣ}
\biblebook{ΜΑΛΑΧΙΑΣ}


\lettrine[lines=2, loversize=0.2, nindent=0em, findent=.25em]{\textcolor{bookheadingcolor}{Λ}}{ΗΜΜΑ} λόγου Κυρίου ἐπὶ τὸν Ἰσραὴλ ἐν χειρὶ ἀγγέλου αὐτοῦ, θέσθε δὴ ἐπὶ τὰς καρδίας ὑμῶν.

\postdropcapindent\vs{2}Ἠγάπησα ὑμᾶς, λέγει Κύριος· καὶ εἴπατε, ἐν τίνι ἠγάπησας ἡμᾶς; οὐκ ἀδελφὸς ἦν Ἠσαῦ τοῦ Ἰακὼβ, λέγει Κύριος, καὶ ἠγάπησα τὸν Ἰακὼβ,
\vs{3}τὸν δὲ Ἠσαῦ ἐμίσησα, καὶ ἔταξα τὰ ὅρια αὐτοῦ εἰς ἀφανισμὸν, καὶ τὴν κληρονομίαν αὐτοῦ εἰς δώματα ἐρήμου;
\vs{4}Διότι ἐρεῖ, ἡ Ἰδουμαία κατέστραπται, καὶ ἐπιστρέψωμεν, καὶ ἀνοικοδομήσωμεν τὰς ἐρήμους· τάδε λέγει Κύριος παντοκράτωρ, αὐτοὶ οἰκοδομήσουσι, καὶ ἐγὼ καταστρέψω· καὶ ἐπικληθήσεται αὐτοῖς ὅρια ἀνομίας, καὶ λαὸς ἐφʼ ὃν παρατέτακται Κύριος ἕως αἰῶνος.
\vs{5}Καὶ οἱ ὀφθαλμοὶ ὑμῶν ὄψονται, καὶ ὑμεῖς ἐρεῖτε, ἐμεγαλύνθη Κύριος ὑπεράνω τῶν ὁρίων τοῦ Ἰσραήλ.

\vs{6}Υἱὸς δοξάζει πατέρα, καὶ δοῦλος τὸν κύριον ἑαυτοῦ. καὶ εἰ πατήρ εἰμι ἐγὼ, ποῦ ἐστιν ἡ δόξα μου; καὶ εἰ Κύριός εἰμι ἐγὼ, ποῦ ἐστιν ὁ φόβος μου; λέγει Κύριος παντοκράτωρ· ὑμεῖς οἱ ἱερεῖς οἱ φαυλίζοντες τὸ ὄνομά μου, καὶ εἴπατε, ἐν τίνι ἐφαυλίσαμεν τὸ ὄνομά σου;
\vs{7}Προσάγοντες πρὸς τὸ θυσιαστήριόν μου ἄρτους ἠλισγημένους, καὶ εἴπατε, ἐν τίνι ἠλισγήσαμεν αὐτούς; ἐν τῷ λέγειν ὑμᾶς, τράπεζα Κυρίου ἠλισγημένη ἐστί, καὶ τὰ ἐπιτιθέμενα ἐξουδενώσατε.
\vs{8}Διότι ἐὰν προσαγάγητε τυφλὸν εἰς θυσίας, οὐ κακόν; καὶ ἐὰν προσαγάγητε χωλὸν ἢ ἄῤῥωστον, οὐ κακόν; προσάγαγε δὴ αὐτὸ τῷ ἡγουμένῳ σου, εἰ προσδέξεταί σε, εἰ λήψεται πρόσωπόν σου, λέγει Κύριος παντοκράτωρ.

\vs{9}Καὶ νῦν ἐξιλάσκεσθε τὸ πρόσωπον τοῦ Θεοῦ ὑμῶν, καὶ δεήθητε αὐτοῦ. Ἐν χερσὶν ὑμῶν γέγονε ταῦτα, εἰ λήψομαι ἐξ ὑμῶν πρόσωπα ὑμῶν; λέγει Κύριος παντοκράτωρ.
\vs{10}Διότι καὶ ἐν ὑμῖν συγκλεισθήσονται θύραι, καὶ οὐκ ἀνάψεται τὸ θυσιαστήριόν μου δωρεάν· οὐκ ἔστι μου θέλημα ἐν ὑμῖν, λέγει Κύριος παντοκράτωρ, καὶ θυσίαν οὐ προσδέξομαι ἐκ τῶν χειρῶν ὑμῶν.
\vs{11}Διότι ἀπὸ ἀνατολῶν ἡλίου καὶ ἕως δυσμῶν τὸ ὄνομά μου δεδόξασται ἐν τοῖς ἔθνεσι, καὶ ἐν παντὶ τόπῳ θυμίαμα προσάγεται τῷ ὀνόματὶ μου, καὶ θυσία καθαρά· διότι μέγα τὸ ὄνομά μου ἐν τοῖς ἔθνεσι, λέγει Κύριος παντοκράτωρ.

\vs{12}Ὑμεῖς δὲ βεβηλοῦτε αὐτὸ ἐν τῷ λέγειν ὑμᾶς, τράπεζα Κυρίου ἠλισγημένη ἐστὶ, καὶ τὰ ἐπιτιθέμενα ἐξουδένωται βρώματα αὐτοῦ·
\vs{13}Καὶ εἴπατε, ταῦτα ἐν κακοπαθείας ἐστί. καὶ ἐξεφύσησα αὐτὰ, λέγει Κύριος παντοκράτωρ· καὶ εἰσεφέρετε ἁρπάγματα, καὶ τὰ χωλὰ, καὶ τὰ ἐνοχλούμενα· καὶ ἐὰν φέρητε τὴν θυσίαν, εἰ προσδέξομαι αὐτὰ ἐκ τῶν χειρῶν ὑμῶν; λέγει Κύριος παντοκράτωρ.
\vs{14}Καὶ ἐπικατάρατος, ὃς ἦν δυνατὸς, καὶ ὑπῆρχεν ἐν τῷ ποιμνίῳ αὐτοῦ ἄρσεν, καὶ εὐχὴ αὐτοῦ ἐπʼ αὐτῷ, καὶ θύει διεφθαρμένον τῷ Κυρίῳ· διότι βασιλεὺς μέγας ἐγώ εἰμι, λέγει Κύριος παντοκράτωρ, καὶ τὸ ὄνομά μου ἐπιφανὲς ἐν τοῖς ἔθνεσι.

\ch{2}
Καὶ νῦν ἡ ἐντολὴ αὕτη πρὸς ὑμᾶς οἱ ἱερεῖς.
\vs{2}Ἐὰν μὴ ἀκούσητε, καὶ ἐὰν μὴ θῆσθε εἰς τὴν καρδίαν ὑμῶν, τοῦ δοῦναι δόξαν τῷ ὀνόματί μου, λέγει Κύριος παντοκράτωρ· καὶ ἐξαποστελῶ ἐφʼ ὑμᾶς τὴν κατάραν, καὶ ἐπικαταράσομαι τὴν εὐλογίαν ὑμῶν, καὶ καταράσομαι αὐτήν· καὶ διασκεδάσω τὴν εὐλογίαν ὑμῶν, καὶ οὐκ ἔσται ἐν ὑμῖν, ὅτι ὑμεῖς οὐ τίθεσθε εἰς τὴν καρδίαν ὑμῶν.
\vs{3}Ἰδοὺ ἐγὼ ἀφορίζω ὑμῖν τὸν ὦμον, καὶ σκορπιῶ ἔνυστρον ἐπὶ τὰ πρόσωπα ὑμῶν, ἔνυστρον ἑορτῶν, ὑμῶν, καὶ λήψομαι ὑμᾶς εἰς τὸ αὐτό·
\vs{4}Καὶ ἐπιγνώσεσθε διότι ἐγὼ ἐξαπέσταλκα πρὸς ὑμᾶς τὴν ἐντολὴν ταύτην, τοῦ εἶναι τὴν διαθήκην μου πρὸς τοὺς Λευίτας, λέγει Κύριος παντοκράτωρ.

\vs{5}Ἡ διαθήκη μου ἦν μετʼ αὐτοῦ τῆς ζωῆς καὶ τῆς εἰρήνης, καὶ ἔδωκα αὐτῷ ἐν φόβῳ φοβεῖσθαί με, καὶ ἀπὸ προσώπου ὀνόματός μου στέλλεσθαι αὐτόν·
\vs{6}νόμος ἀληθείας ἦν ἐν τῷ στόματι αὐτοῦ, καὶ ἀδικία οὐχ εὑρέθη ἐν χείλεσιν αὐτοῦ· ἐν εἰρήνῃ κατευθύνων ἐπορεύθη μετʼ ἐμοῦ, καὶ πολλοὺς ἐπέστρεψεν ἀπὸ ἀδικίας.
\vs{7}Ὅτι χείλη ἱερέως φυλάξεται γνῶσιν, καὶ νόμον ἐκζητήσουσιν ἐκ στόματος αὐτοῦ, διότι ἄγγελος Κυρίου παντοκράτορός ἐστιν.

\vs{8}Ὑμεῖς δὲ ἐξεκλίνατε ἐκ τῆς ὁδοῦ, καὶ ἠσθενήσατε πολλοὺς ἐν νόμῳ, διεφθείρατε τὴν διαθήκην τοῦ Λευὶ, λέγει Κύριος παντοκράτωρ.
\vs{9}Κᾀγὼ δέδωκα ὑμᾶς ἐξουδενωμένους καὶ ἀπεῤῥιμμένους εἰς πάντα τὰ ἔθνη, ἀνθʼ ὧν ὑμεῖς οὐκ ἐφυλάξασθε τὰς ὁδούς μου, ἀλλὰ ἐλαμβάνετε πρόσωπα ἐν νόμῳ.

\vs{10}Οὐχὶ πατὴρ εἷς πάντων ὑμῶν; οὐχὶ Θεὸς εἷς ἔκτισεν ὑμᾶς; τί ὅτι ἐγκατελίπετε ἕκαστος τὸν ἀδελφὸν αὐτοῦ, τοῦ βεβηλῶσαι τὴν διαθήκην τῶν πατέρων ὑμῶν;

\vs{11}Ἐνκατελείφθη Ἰούδας, καὶ βδέλυγμα ἐγένετο ἐν τῷ Ἰσραὴλ καὶ ἐν Ἱερουσαλὴμ, διότι ἐβεβήλωσεν Ἰούδας τὰ ἅγια Κυρίου, ἐν οἷς ἠγάπησε, καὶ ἐπετήδευσεν εἰς θεοὺς ἀλλοτρίους.
\vs{12}Ἐξολοθρεύσει Κύριος τὸν ἄνθρωπον τὸν ποιοῦντα ταῦτα, ἕως καὶ ταπεινωθῇ ἐκ σκηνωμάτων Ἰακὼβ, καὶ ἐκ προσαγόντων θυσίαν τῷ κυρίῳ παντοκράτορι.
\vs{13}Καὶ ταῦτα, ἃ ἔμισουν, ἐποιεῖτε· ἐκαλύπτετε δάκρυσι τὸ θυσιαστήριον Κυρίου, καὶ κλαυθμῷ καὶ στεναγμῷ ἐκ κόπων· ἔτι ἄξιον ἐπιβλέψαι εἰς θυσίαν, ἢ λαβεῖν δεκτὸν ἐκ τῶν χειρῶν ὑμῶν;

\vs{14}Καὶ εἴπατε, ἕνεκεν τίνος; ὅτι Κύριος διεμαρτύρατο ἀναμέσον σου, καὶ ἀναμέσον γυναικὸς νεότητός σου, ἣν ἐγκατέλιπες, καὶ αὕτη κοινωνός σου, καὶ γυνὴ διαθήκης σου.
\vs{15}Καὶ οὐ καλὸν ἐποίησε; καὶ ὑπόλειμμα πνεύματος αὐτοῦ· καὶ εἴπατε, τί ἄλλο ἢ σπέρμα ζητεῖ ὁ Θεός; καὶ φυλάξασθε ἐν τῷ πνεύματι ὑμῶν, καὶ γυναῖκα νεότητός σου μὴ ἐγκαταλίπῃς.
\vs{16}Ἀλλὰ ἐὰν μισήσας ἐξαποστείλῃς, λέγει Κύριος ὁ Θεὸς τοῦ Ἰσραὴλ, καὶ καλύψει ἀσέβεια ἐπὶ τὰ ἐνθυμήματά σου, λέγει Κύριος παντοκράτωρ. καὶ φυλάξασθε ἐν τῷ πνεύματι ὑμῶν, καὶ οὐ μὴ ἐγκαταλίπητε
\vs{17}οἱ παροξύναντες τὸν Θεὸν ἐν τοῖς λόγοις ὑμῶν· καὶ εἴπατε, ἐν τίνι παρωξύναμεν αὐτόν; ἐν τῷ λέγειν ὑμᾶς, πᾶς ποιῶν πονηρὸν, καλὸν ἐνώπιον Κυρίου, καὶ ἐν αὐτοῖς αὐτὸς εὐδόκησε, καί ποῦ ἐστιν ὁ Θεὸς τῆς δικαιοσύνης;

\ch{3}
Ἰδοὺ ἐξαποστέλλω τὸν ἄγγελόν μου, καὶ ἐπιβλέψεται ὁδὸν πρὸ προσώπου μου, καὶ ἐξαίφνης ἥξει εἰς τὸν ναὸν ἑαυτοῦ Κύριος, ὃν ὑμεῖς ζητεῖτε, καὶ ὁ ἄγγελος τῆς διαθήκης, ὃν ὑμεῖς θέλετε· ἰδοὺ ἔρχεται, λέγει Κύριος παντοκράτωρ,
\vs{2}καὶ τίς ὑπομενεῖ ἡμέραν εἰσόδου αὐτοῦ; ἢ τίς ὑποστήσεται ἐν τῇ ὀπτασίᾳ αὐτοῦ; διότι αὐτὸς εἰσπορεύεται ὡς πῦρ χωνευτηρίου, καὶ ὡς ποὶα πλυνόντων.
\vs{3}Καθιεῖται χωνεύων καὶ καθαρίζων ὡς τὸ ἀργύριον, καὶ ὡς τὸ χρυσίον, καὶ καθαρίσει τοὺς υἱοὺς Λευὶ, καὶ χεεῖ αὐτοὺς ὥσπερ τὸ χρυσίον καὶ τὸ ἀργύριον· καὶ ἔσονται τῷ Κυρίῳ προσάγοντες θυσίαν ἐν δικαιοσύνῃ.

\vs{4}Καὶ ἀρέσει τῷ Κυρίῳ θυσία Ἰούδα καὶ Ἱερουσαλὴμ, καθὼς αἱ ἡμέραι τοῦ αἰῶνος, καὶ καθὼς τὰ ἔτη τὰ ἔμπροσθεν.
\vs{5}Καὶ προσάξω πρὸς ὑμᾶς ἐν κρίσει, καὶ ἔσομαι μάρτυς ταχὺς ἐπὶ τὰς φαρμακοὺς, καὶ ἐπὶ τὰς μοιχαλίδας, καὶ ἐπὶ τοὺς ὀμνύοντας τῷ ὀνόματί μου ἐπὶ ψεύδει, καὶ ἐπὶ τοὺς ἀποστεροῦντας μισθὸν μισθωτοῦ, καὶ τοὺς καταδυναστεύοντας χήραν, καὶ τοὺς κονδυλίζοντας ὀρφανοὺς, καὶ τοὺς ἐκκλίνοντας κρίσιν προσηλύτου, καὶ τοὺς μὴ φοβουμένους με, λέγει Κύριος παντοκράτωρ.
\vs{6}Διότι ἐγὼ Κύριος ὁ Θεὸς ὑμῶν, καὶ οὐκ ἠλλοίωμαι·
\vs{7}καὶ ὑμεῖς οἱ υἱοὶ Ἰακὼβ οὐκ ἀπέχεσθε ἀπὸ τῶν ἀδικιῶν τῶν πατέρων ὑμῶν, ἐξεκλίνατε νόμιμά μου, καὶ οὐκ ἐφυλάξασθε.

Ἐπιστρέψατε πρὸς μὲ, καὶ ἐπιστραφήσομαι πρὸς ὑμᾶς, λέγει Κύριος παντοκράτωρ· καὶ εἴπατε, ἐν τίνι ἐπιστρέψομεν;
\vs{8}Μήτι πτερνιεῖ ἄνθρωπος Θεόν; διότι ὑμεῖς πτερνίζετέ με· καὶ ἐρεῖτε, ἐν τίνι ἐπτερνίσαμέν σε; ὅτι τὰ ἐπιδέκατα, καὶ αἱ ἀπαρχαὶ μεθʼ ὑμῶν εἰσι.
\vs{9}Καὶ ἀποβλέποντες ὑμεῖς ἀποβλέπετε, καὶ ἐμὲ ὑμεῖς πτερνίζετε.

\vs{10}Τὸ ἔτος συνετελέσθη, καὶ εἰσηνέγκατε πάντα τὰ ἐκφόρια εἰς τοὺς θησαυροὺς, καὶ ἔσται ἡ διαρπαγὴ αὐτοῦ ἐν τῷ οἴκῳ αὐτοῦ· ἐπιστρέψατε δὴ ἐν τούτῳ, λέγει Κύριος παντοκράτωρ· ἐὰν μὴ ἀνοίξω ὑμῖν τοὺς καταῤῥάκτας τοῦ οὐρανοῦ, καὶ ἐκχεῶ τὴν εὐλογίαν μου ὑμῖν, ἕως τοῦ ἱκανωθῆναι·
\vs{11}Καὶ διαστελῶ ὑμῖν εἰς βρῶσιν, καὶ οὐ μὴ διαφθείρω ὑμῶν τὸν καρπὸν τῆς γῆς, καὶ οὐ μὴ ἀσθενήσῃ ὑμῶν ἡ ἄμπελος ἡ ἐν τῷ ἀγρῷ, λέγει Κύριος παντοκράτωρ.
\vs{12}Καὶ μακαριοῦσιν ὑμᾶς πάντα τὰ ἔθνη, διότι ἔσεσθε ὑμεῖς γῆ θελητὴ, λέγει Κύριος παντοκράτωρ.

\vs{13}Ἐβαρύνατε ἐπʼ ἐμὲ τοὺς λόγους ὑμῶν, λέγει Κύριος· καὶ εἴπατε, ἐν τίνι κατελαλήσαμεν κατὰ σοῦ;
\vs{14}Εἴπατε, μάταιος ὁ δουλεύων Θεῷ, καὶ τί πλέον, ὅτι ἐφυλάξαμεν τὰ φυλάγματα αὐτοῦ, καὶ διότι ἐπορεύθημεν ἱκέται πρὸ προσώπου Κυρίου παντοκράτορος;
\vs{15}Καὶ νῦν ἡμεῖς μακαρίζομεν ἀλλοτρίους, καὶ ἀνοικοδομοῦνται πάντες ποιοῦντες ἄνομα, καὶ ἀντέστησαν τῷ Θεῷ, καὶ ἐσώθησαν.

\vs{16}Ταῦτα κατελάλησαν οἱ φοβούμενοι τὸν κύριον, ἕκαστος πρὸς τὸν πλησίον αὐτοῦ· καὶ προσέσχε Κύριος, καὶ εἰσήκουσε, καὶ ἔγραψε βιβλίον μνημοσύνου ἐνώπιον αὐτοῦ τοῖς φοβουμένοις τὸν Κύριον, καὶ εὐλαβουμένοις τὸ ὄνομα αὐτοῦ.
\vs{17}Καὶ ἔσονταί μοι, λέγει Κύριος παντοκράτωρ, εἰς ἡμέραν, ἣν ἐγὼ ποιῶ, εἰς περιποίησιν, καὶ αἱρετιῶ αὐτοὺς, ὃν τρόπον αἱρετίζει ἄνθρωπος τὸν υἱὸν αὐτοῦ, τὸν δουλεύοντα αὐτῷ.
\vs{18}Καὶ ἐπιστραφήσεσθε, καὶ ὄψεσθε ἀναμέσον δικαίου, καὶ ἀναμέσον ἀνόμου, καὶ ἀναμέσον τοῦ δουλεύοντος Θεῷ, καὶ τοῦ μὴ δουλεύοντος.

\vs{19}Διότι ἰδοὺ ἡμέρα ἔρχεται καιομένη ὡς κλίβανος καὶ φλέξει αὐτοὺς, καὶ ἔσονται πάντες οἱ ἀλλογενεῖς, καὶ πάντες οἱ ποιοῦντες ἄνομα, καλάμη, καὶ ἀνάψει αὐτοὺς ἡ ἡμέρα ἡ ἐρχομένη, λέγει Κύριος παντοκράτωρ, καὶ οὐ μὴ ὑπολειφθῇ ἐξ αὐτῶν ῥίζα οὐδὲ κλῆμα.

\vs{20}Καὶ ἀνατελεῖ ὑμῖν τοῖς φοβουμένοις τὸ ὄνομά μου ἥλιος δικαιοσύνης, καὶ ἴασις ἐν ταῖς πτέρυξιν αὐτοῦ· καὶ ἐξελεύσεσθε, καὶ σκιρτήσετε ὡς μοσχάρια ἐκ δεσμῶν ἀνειμένα.
\vs{21}Καὶ καταπατήσετε ἀνόμους, διότι ἔσονται σποδὸς ὑποκάτω τῶν ποδῶν ὑμῶν ἐν τῇ ἡμέρᾳ ᾗ ἐγὼ ποιῶ, λέγει Κύριος παντοκράτωρ.
\vs{22}Καὶ ἰδοὺ ἐγὼ ἀποστελῶ ὑμῖν Ἠλίαν τὸν Θεσβίτην, πρὶν ἐλθεῖν τὴν ἡμέραν Κυρίου τὴν μεγάλην καὶ ἐπιφανῆ,
\vs{23}ὃς ἀποκαταστήσει καρδίαν πατρὸς πρὸς υἱὸν, καὶ καρδίαν ἀνθρώπου πρὸς τὸν πλησίον αὐτοῦ, μὴ ἔλθω καὶ πατάξω τὴν γῆν ἄρδην.

\vs{24}Μνήσθητε νόμου Μωυσῆ τοῦ δούλου μου, καθότι ἐνετειλάμην αὐτῷ ἐν Χωρὴβ πρὸς πάντα τὸν Ἰσραὴλ, προστάγματα καὶ δικαιώματα.


\def\book{ΕΣΔΡΑΣ Αʹ}
\biblebook{ΕΣΔΡΑΣ Αʹ}


\lettrine[lines=2, loversize=0.2, nindent=0em, findent=.25em]{\textcolor{bookheadingcolor}{Κ}}{ΑΙ} ἤγαγεν Ἰωσίας τὸ πάσχα ἐν Ἱερουσαλὴμ τῷ Κυρίῳ αὐτοῦ, καὶ ἔθυσε τὸ πάσχα τῇ τεσσαρεσκαιδεκάτῃ ἡμέρᾳ τοῦ μηνὸς τοῦ πρώτου·
\vs{2}στήσας τοὺς ἱερεῖς κατʼ ἐφημερίας ἐστολισμένους ἐν τῷ ἱερῷ τοῦ Κυρίου.

\vs{3}Καὶ εἶπε τοῖς Λευίταις ἱεροδούλοις τοῦ Ἰσραὴλ, ἁγιάσαι ἑαυτοὺς τῷ Κυρίῳ ἐν τῇ θέσει τῆς ἁγίας κιβωτοῦ τοῦ Κυρίου ἐν τῷ οἴκῳ ᾧ ᾠκοδόμησε Σαλωμὼν ὁ τοῦ Δαυὶδ ὁ βασιλεύς·
\vs{4}οὐκ ἔσται ὑμῖν ἆραι ἐπʼ ὤμων αὐτήν· καὶ νῦν λατρεύετε τῷ Κυρίῳ Θεῷ ὑμῶν, καὶ θεραπεύετε τὸ ἔθνος αὐτοῦ Ἰσραήλ, καὶ ἑτοιμάσατε κατὰ τὰς πατριὰς καὶ τὰς φυλὰς ὑμῶν,
\vs{5}κατὰ τὴν γραφὴν Δαυὶδ βασιλέως Ἰσραὴλ, καὶ κατὰ τὴν μεγαλειότητα Σαλωμὼν τοῦ υἱοῦ αὐτοῦ· καὶ στάντες ἐν τῷ ἁγίῳ κατὰ τὴν μεριδαρχίαν τὴν πατρικὴν ὑμῶν τῶν Λευιτῶν, τῶν ἔμπροσθεν τῶν ἀδελφῶν ὑμῶν υἱῶν Ἰσραὴλ,
\vs{6}ἐν τάξει θύσατε τὸ πάσχα, καὶ τὰς θυσίας ἑτοιμάσατε τοῖς ἀδελφοῖς ὑμῶν, καὶ ποιήσατε τὸ πάσχα κατὰ τὸ πρόσταγμα τοῦ Κυρίου τὸ δοθὲν τῷ Μωυσῇ.

\vs{7}Καὶ ἐδωρήσατο Ἰωσίας τῷ λαῷ τῷ εὑρεθέντι ἀρνῶν καὶ ἐρίφων τριάκοντα χιλιάδας, μόσχους τρισχιλίους· ταῦτα ἐκ τῶν βασιλικῶν ἐδόθη κατʼ ἐπαγγελίαν τῷ λαῷ, καὶ τοῖς ἱερεῦσι, καὶ Λευίταις.
\vs{8}Καὶ ἔδωκε Χελκίας, καὶ Ζαχαρίας, καὶ Συῆλος οἱ ἐπιστάται τοῦ ἱεροῦ τοῖς ἱερεῦσιν εἰς πάσχα πρόβατα δισχίλια ἑξακόσια, μόσχους τριακοσίους.
\vs{9}Καὶ Ἰεχονὶας, καὶ Σαμαίας, καὶ Ναθαναὴλ ὁ ἀδελφὸς, καὶ Ἀσαβίας, καὶ Ὀχίηλος, καὶ Ἰωρὰμ χιλίαρχοι ἔδωκαν τοῖς Λευίταις εἰς πάσχα πρόβατα πεντακισχίλια, μόσχους ἑπτακοσίους.

\vs{10}Καὶ ταῦτα τὰ γενόμενα, εὐπρεπῶς ἔστησαν οἱ ἱερεῖς καὶ οἱ Λευῖται,
\vs{11}ἔχοντες τὰ ἄζυμα κατὰ τὰς φυλὰς καὶ κατὰ τὰς μεριδαρχίας τῶν πατέρων ἔμπροσθεν τοῦ λαοῦ, προσενεγκεῖν τῷ Κυρίῳ κατὰ τὰ γεγραμμένα ἐν βιβλίῳ Μωυσῆ· καὶ οὕτως τὸ πρωϊνόν.
\vs{12}Καὶ ὤπτησαν τὸ πάσχα πυρὶ ὡς καθήκει, καὶ τὰς θυσίας ἥψησαν ἐν τοῖς χαλκείοις καὶ λέβησι μετʼ εὐωδίας, καὶ ἀπήνεγκαν πᾶσι τοῖς ἐκ τοῦ λαοῦ·
\vs{13}μετὰ δὲ ταῦτα ἡτοίμασαν ἑαυτοῖς τε καὶ τοῖς ἱερεῦσιν ἀδελφοῖς αὐτῶν υἱοῖς Ἀαρών·
\vs{14}οἱ γὰρ ἱερεῖς ἀνέφερον τὰ στέατα ἕως ἀωρίας· καὶ οἱ Λευῖται ἡτοίμασαν ἑαυτοῖς καὶ τοῖς ἱερεῦσιν ἀδελφοῖς αὐτῶν υἱοῖς Ἀαρών.
\vs{15}Καὶ οἱ ἱεροψάλται υἱοὶ Ἀσὰφ ἦσαν ἐπὶ τῆς τάξεως αὐτῶν, κατὰ τὰ ὑπὸ Δαυὶδ τεταγμένα, καὶ Ἀσὰφ, καὶ Ζαχαρίας, καὶ Ἐδδινοῦς ὁ παρὰ τοῦ βασιλέως.
\vs{16}Καὶ οἱ θυρωροὶ ἐφʼ ἑκάστου πυλῶνος· οὐκ ἔστι παραβῆναι ἕκαστον τὴν ἑαυτοῦ ἐφημερίαν· οἱ γὰρ ἀδελφοὶ αὐτῶν οἱ Λευῖται ἡτοίμασαν αὐτοῖς,
\vs{17}καὶ συνετελέσθη τὰ τῆς θυσίας τοῦ Κυρίου ἐν ἐκείνῃ τῇ ἡμέρᾳ ἀχθῆναι τὸ πάσχα,
\vs{18}καὶ προσαχθῆναι τὰς θυσίας ἐπὶ τὸ τοῦ Κυρίου θυσιαστήριον, κατὰ τὴν ἐπιταγὴν τοῦ βασιλέως Ἰωσίου.

\vs{19}Καὶ ἠγάγοσαν οἱ υἱοὶ Ἰσραὴλ οἱ εὑρεθέντες ἐν τῷ καιρῷ τούτῳ τὸ πάσχα καὶ τὴν ἑορτὴν τῶν ἀζύμων ἡμέρας ἑπτά.
\vs{20}Καὶ οὐκ ἤχθη τὸ πάσχα τοιοῦτον ἐν τῷ Ἰσραὴλ ἀπὸ τῶν χρόνων Σαμουὴλ τοῦ προφήτου.
\vs{21}Καὶ πάντες οἱ βασιλεῖς τοῦ Ἰσραὴλ οὐκ ἠγάγοσαν πάσχα τοιοῦτον, οἷον ἤγαγεν Ἰωσίας, καὶ οἱ ἱερεῖς, καὶ οἱ Λευῖται, καὶ οἱ Ἰουδαῖοι, καὶ πᾶς Ἰσραὴλ ὁ εὑρεθεὶς ἐν τῇ κατοικήσει αὐτῶν ἐν Ἱερουσαλήμ.
\vs{22}Ὀκτωκαιδεκάτῳ ἔτει βασιλεύοντος Ἰωσίου ἤχθη τὸ πάσχα τοῦτο.

\vs{23}Καὶ ὠρθώθη τὰ ἔργα Ἰωσίου ἐνώπιον τοῦ Κυρίου αὐτοῦ ἐν καρδίᾳ πλήρει εὐσεβείας.
\vs{24}Καὶ τὰ κατʼ αὐτὸν δὲ ἀναγέγραπται ἐν τοῖς ἔμπροσθεν χρόνοις, περὶ τῶν ἡμαρτηκότων καὶ ἠσεβηκότων εἰς τὸν Κύριον παρὰ πᾶν ἔθνος καὶ βασιλείαν, καὶ ἃ ἐλύπησαν αὐτὸν, ἔστι, καὶ οἱ λόγοι τοῦ Κυρίου ἀνέστησαν ἐπὶ Ἰσραήλ.

\vs{25}Καὶ μετὰ πᾶσαν τὴν πρᾶξιν ταύτην Ἰωσίου, συνέβη Φαραὼ βασιλέα Αἰγύπτου ἐλθόντα πόλεμον ἐγεῖραι ἐν Χαρκαμὺς ἐπὶ τοῦ Εὐφράτου· καὶ ἐξῆλθεν εἰς ἀπάντησιν αὐτῷ Ἰωσίας.
\vs{26}Καὶ διεπέμψατο πρὸς αὐτὸν βασιλεὺς Αἰγύπτου, λέγων, τί ἐμοὶ καὶ σοί ἐστι, βασιλεῦ τῆς Ἰουδαίας;
\vs{27}Οὐχί πρὸς σὲ ἐξαπέσταλμαι ὑπὸ Κυρίου τοῦ Θεοῦ· ἐπὶ γὰρ τοῦ Εὐφράτου ὁ πόλεμός μου ἐστί· καὶ νῦν Κύριος μετʼ ἑμοῦ ἐστι, καὶ Κύριος μετʼ ἐνοῦ ἑπισπεύδων ἐστίν· ἀπόστηθι, καὶ μὴ ἐναντιοῦ τῷ Κυρίῳ.

\vs{28}Καὶ οὐκ ἀπέστρεψεν ἑαυτὸν Ἰωσίας ἐπὶ τὸ ἅρμα αὐτοῦ, ἀλλὰ πολεμεῖν αὐτὸν ἐπεχείρει, οὐ προσέχων ῥήμασιν Ἰερεμίου προφήτου ἐκ στόματος Κυρίου.
\vs{29}Ἀλλὰ συνεστήσατο πρὸς αὐτὸν πόλεμον ἐν τῷ πεδίῳ Μαγεδδώ· καὶ κατέβησαν οἱ ἄρχοντες πρὸς τὸν βασιλέα Ἰωσίαν.
\vs{30}Καὶ εἶπεν ὁ βασιλεὺς τοῖς παισὶν ἑαυτοῦ, ἀποστήσατέ με ἀπὸ τῆς μάχης, ἠσθένησα γὰρ λίαν· καὶ εὐθέως ἀπέστησαν αὐτὸν οἱ παῖδες αὐτοῦ ἀπὸ τῆς παρατάξεως.
\vs{31}Καὶ ἀνέβη ἐπὶ τὸ ἅρμα τὸ δευτέριον αὐτοῦ, καὶ ἀποκατασταθεὶς, εἰς Ἰερουσαλὴμ, μετήλλαξε τὸν βίον αὐτοῦ, καὶ ἐτάφη ἐν τῷ πατρικῷ τάφῳ.
\vs{32}Καὶ ἐν ὅλῃ τῇ Ἰουδαίᾳ ἐπένθησαν τὸν Ἰωσίαν, καὶ ἐθρήνησεν Ἱερεμίας ὁ προφήτης ὑπὲρ Ἰωσίον, ὑαι οἱ προκαθήμενοι σὺν γυναιξὶν ἐθρηνοῦσαν αὐτὸν ἕως τῆς ἡμέρας ταύτης· καὶ ἐξεδόθη τοῦτο γίνεσθαι ἀεὶ εἰς ἅπαν τὸ γένος Ἰσραήλ.

\vs{33}Ταῦτα δὲ ἀναγέγραπται ἐν τῇ βίβλῳ τῶν ἱστορουμένων περὶ τῶν βασιλέων τῆς Ἰουδαίας, καὶ τὸ καθʼ ἓν πραχθὲν τῆς πράξεως Ἰωσίου, καὶ τῆς δόξης αὐτοῦ, καὶ τῆς συνέσεως αὐτοῦ ἐν τῷ νόμῳ Κυρίου· τά τε προπραχθέντα ὑπʼ αὐτοῦ καὶ τὰ νῦν, ἱστόρηται ἐν τῷ βιβλίῳ τῶν βασιλέων Ἰσραὴλ καὶ Ἰούδα.

\vs{34}Καὶ ἀναλαβόντες οἱ ἐκ τοῦ ἔθνους τὸν Ἰεχονίαν υἱὸν Ἰωσίου, ἀνέδειξαν βασιλέα ἀντὶ Ἰωσείου τοῦ πατρὸς αὐτοῦ, ὄντα ἐτῶν εἴκοσι τριῶν.
\vs{35}Καὶ ἐβασίλευσεν ἐν Ἰσραὴλ καὶ Ἱερουσαλὴμ μῆνας τρεῖς· καὶ ἀπέστησεν αὐτὸν βασιλεὺς Αἰγύπτου τοῦ μὴ βασιλεύειν ἐν Ἱερουσαλὴμ,
\vs{36}καὶ ἐζημίωσε τὸ ἔθνος ἀργυρίου ταλάντοις ἑκατὸν καὶ χρυσίου ταλάντῳ ἑνί.

\vs{37}Καὶ ἀνέδειξε βασιλεὺς Αἰγύπτου βασιλέα Ἰωακὶμ τὸν ἀδελφὸν αὐτοῦ βασιλέα τῆς Ἰουδαίας καὶ Ἱερουσαλήμ.
\vs{38}Καὶ ἔδησεν Ἰωακὶμ τοὺς μεγιστᾶνας, Ζαράκην δὲ τὸν ἀδελφὸν αὐτοῦ συλλαβὼν ἀνήγαγεν ἐξ Αἰγύπτου.
\vs{39}Ἐτῶν δὲ ἦν εἰκοσιπέντε Ἰωακὶμ ὅτε ἐβασίλευσε τῆς Ἰουδαίας καὶ Ἱερουσαλήμ· καὶ ἐποίησε τὸ πονηρὸν ἐνώπιον Κυρίου.
\vs{40}Μετʼ αὐτὸν δὲ ἀνέβη Ναβουχοδονόσορ ὁ βασιλεὺς Βαβυλῶνος, καὶ ἔδησεν αὐτὸν ἐν χαλκείῳ δεσμῷ, καὶ ἀπήγαγεν εἰς Βαβυλῶνα.
\vs{41}Καὶ ἀπὸ τῶν ἱερῶν σκευῶν τοῦ Κυρίου λαβὼν Ναβουχοδονόσορ καὶ ἀπενέγκας, ἀπηρείσατο ἐν τῷ ναῷ αὐτοῦ ἐν Βαβυλῶνι.
\vs{42}Τὰ δὲ ἱστορηθέντα περὶ αὐτοῦ, καὶ τῆς ἀκαθαρσίας αὐτοῦ καὶ δυσσεβείας, ἀναγέγραπται ἐν τῇ βιβλῳ τῶν χρόνων τῶν βασιλέων.

\vs{43}Καὶ ἐβασίλευσεν ἀντʼ αὐτοῦ Ἰωακὶμ ὁ υἱὸς αὐτοῦ· ὅτε γὰρ ἀνεδείχθη, ἦν ἐτῶν ὀκτώ.
\vs{44}Βασιλεύει δὲ μῆνας τρεῖς καὶ ἡμέρας δέκα ἐν Ἱερουσαλὴμ, καὶ ἐποίησε τὸ πονηρὸν ἔναντι Κυρίου.

\vs{45}Καὶ μετʼ ἐνιαυτὸν ἀποστείλας Ναβουχοδονόσορ μετήγαγεν αὐτὸν εἰς Βαβυλῶνα, ἅμα τοῖς ἱεροῖς σκεύεσι τοῦ Κυρίου,
\vs{46}καὶ ἀνέδειξε Σεδεκίαν βασιλέα τῆς Ἰουδαίας καὶ Ἱερουσαλὴμ, ὄντα ἐτῶν εἴκοσι ἑνός· βασιλεύει δὲ ἔτη ἕνδεκα,
\vs{47}καὶ ἐποίησε τὸ πονηρὸν ἐνώπιον Κυρίου, καὶ οὐκ ἐνετράπη ἀπὸ τῶν ῥηθέντων λόγων ὑπὸ Ἱερεμίου τοῦ προφήτου ἐκ στόματος τοῦ Κυρίου.
\vs{48}Καὶ ὁρκισθεὶς ἀπὸ τοῦ βασιλέως Ναβουχοδονόσορ τῷ ὀνόματι Κυρίου, ἐπιορκήσας ἀπέστη· καὶ σκληρύνας αὐτοῦ τὸν τράχηλον καὶ τὴν καρδίαν αὐτοῦ, παρέβη τὰ νόμιμα Κυρίου Θεοῦ Ἰσραήλ.
\vs{49}Καὶ οἱ ἡγούμενοι δὲ τοῦ λαοῦ καὶ τῶν ἱερέων πολλὰ ἠσέβησαν καὶ ὑπὲρ πάσας τὰς ἀκαθαρσίας πάντων τῶν ἐθνῶν, καὶ ἐμίαναν τὸ ἱερὸν τοῦ Κυρίον τὸ ἁγιαζόμενον ἐν Ἱερουσαλήμ.

\vs{50}Καὶ ἀπέστειλεν ὁ Θεὸζ τῶν πατέρων αὐτῶν διὰ τοῦ ἀγγέλου αὐτοῦ μετακαλέσαι αὐτοὺς, καθότι ἐφείδετο αὐτῶν καὶ τοῦ σκηνώματος αὐτοῦ.
\vs{51}Αὐτοὶ δὲ ἐμυκτήρισαν ἐν τοῖς ἀγγέλοις αὐτοῦ·
\vs{52}καὶ ᾗ ἡμέρᾳ ἐλάλησε Κύριος, ἦσαν ἐκπαίζοντες τοὺς προφήτας αὐτοῦ, ἕως οὗ θυμῶντα αὐτὸν ἐπὶ τῷ ἔθνει αὐτοῦ διὰ τὰ δυσσεβήματα, προστάξαι ἀναβιβάσαι ἐπʼ αὐτοὺς τοὺς βασιλεῖς τῶν Χαλδαίων.
\vs{53}Οὗτοι ἀπέκτειναν τοὺς νεανίσκους αὐτῶν ἐν ῥομφαία, περικύκλῳ τοῦ ἁγίου αὐτῶν ἱεροῦ· καὶ οὐκ ἐφείσαντο νεανίσκου καὶ παρθένου, καὶ πρεσβύτου καὶ νεωτέρου, ἀλλὰ πάντας παρέδωκαν εἰς τὰς χεῖρας αὐτῶν.
\vs{54}Καὶ πάντα τὰ ἱερὰ σκεύη τοῦ Κυρίου τὰ μεγάλα καὶ τὰ μικρά, καὶ τὰς κιβωτοὺς τοῦ Κυρίου, καὶ τὰς βασιλικὰς ἀποθήκας ἀναλαβόντες ἀπήνεγκαν εἰς Βαβυλῶνα.
\vs{55}Καὶ ἐνεπύρισαν τὸν οἶκον τοῦ Κυρίου, καὶ ἔλυσαν τὰ τείχη Ἱερουσαλὴμ, καὶ τοὺς πύργους αὐτῆς ἐνεπύρισαν ἐν πυρί,
\vs{56}καὶ συνετέλεσαν πάντα τὰ ἔνδοξα αὐτῆς ἀχρειῶσαι, καὶ τοὺς ἐπιλοίπους ἀπήγαγε μετὰ ῥομφαίας εἰς Βαβυλῶνα.
\vs{57}Καὶ ἦσαν παῖδες αὐτῷ καὶ τοῖς υἱοῖς αὐτοῦ, μέχρις οὗ βασιλεῦσαι Πέρσας, εἰς ἀναπλήρωσιν ῥήματος τοῦ Κυρίου ἐν στόματι Ἱερεμίου·
\vs{58}ἕως τοῦ εὐδοκῆσαι τὴν γῆν τὰ σάββατα αὐτῆς, πάντα τὸν χρόνον τῆς ἐρημώσεως αὐτῆς, σαββατιεῖ εἰς συμπλήρωσιν ἐτῶν ἑβδομήκοντα.

\ch{2}
Βασιλεύοντος Κύρου Περσῶν ἔτους πρώτου, εἰς συντέλειαν ῥήματος Κυρίου ἐν στόματι Ἱερεμίου,
\vs{2}ἤγειρε Κύριος τὸ πνεῦμα Κύρου βασιλέως Περσῶν, καὶ ἐκήρυξεν ἐν ὅλῃ τῇ βασιλείᾳ αὐτοῦ, καὶ ἅμα διὰ γραπτῶν, λέγων,
\vs{3}τάδε λέγει ὁ βασιλεὺς Περσῶν Κῦρος, ἐμὲ ἀνέδειξε βασιλέα τῆς οἰκουμένης ὁ Κύριος τοῦ Ἰσραὴλ, Κύριος ὁ ὕψιστος.
\vs{4}Καὶ ἐσήμῃνέ μοι οἰκοδομῆσαι αὐτῷ οἶκον ἐν Ἱερουσαλὴμ, τῇ ἐν τῇ Ἰουδαίᾳ.

\vs{5}Εἴ τις ἐστὶν οὖν ὑμῶν ἐκ τοῦ ἔθνους αὐτοῦ, ἔστω ὁ Κύριος αὐτοῦ μετʼ αὐτοῦ, καὶ ἀναβὰς εἰς τὴν Ἱερουσαλὴμ τὴν ἐν τῇ Ἰουδαίᾳ, οἰκοδομείτω τὸν οἶκον τοῦ Κυρίου τοῦ Ἰσραήλ· οὗ ὁ Κύριος, ὁ κατασκηνώσας ἐν Ἱερουσαλήμ.
\vs{6}Ὅσοι οὖν κατὰ τοὺς τόπους οἰκοῦσι, βοηθείτωσαν αὐτῷ οἱ ἐν τῷ τόπῳ αὐτοῦ, ἐν χρυσίῳ καὶ ἐν ἀργυρίῳ, ἐν δὀσεσι,
\vs{7}μεθʼ ἵππων καὶ κτηνῶν, σὺν τοῖς ἄλλοις τοῖς κατʼ εὐχὰς προστεθειμένοις εἰς τὸ ἱερὸν τοῦ Κυρίου τὸ ἐν Ἱερουσαλήμ.

\vs{8}Καὶ καταστήσαντες οἱ ἀρχίφυλοι τῶν πατριῶν τῆς Ἰούδα καὶ Βενιαμὶν φυλῆς, καὶ οἱ ἱερεῖς καὶ οἱ Λευῖται, καὶ πάντων ὧν ἤγειρε Κύριος τό πνεῦμα, ἀναβῆναι οἰκοδομῆσαι οἶκον τῷ Κυρίῳ τὸν ἐν Ἱερουσαλήμ·
\vs{9}καὶ οἱ περικύκλῳ αὐτῶν ἐβοήθησαν ἐν πᾶσιν, ἐν ἀργυρίῳ καὶ χρυσίῳ, ἵπποις, κτήνεσι, καὶ εὐχαῖς ὡς πλείσταις πολλῶν, ὧν ὁ νοῦς ἠγέρθη.
\vs{10}Καὶ ὁ βασιλεὺς Κύρος ἐξήνεγκε τὰ ἱερὰ σκεύη τοῦ Κυρίου, ἃ μετήνεγκε Ναβουχοδονόσορ ἐξ Ἱερουσαλὴμ, καὶ ἀπηρείσατο αὐτὰ ἐν τῷ εἰδωλείῳ αὐτοῦ.

\vs{11}Ἐξενέγκας δὲ αὐτὰ Κύρος ὁ βασιλεὺς Περσῶν παρέδωκεν αὐτὰ Μιθρδάτῃ τῷ ἑαυτοῦ γαζοφύλακι·
\vs{12}Διὰ δὲ τούτου παρεδόθησαν Σαμανασσάρῳ προστάτῃ τῆς Ἰουδαίας.
\vs{13}Ὁ δὲ τούτων ἀριθμὸς ἦν, σπονδεῖα χρυσᾶ χίλια, σπονδεῖα ἀργυρᾶ χίλια, θυΐσκαι ἀργυραῖ εἰκοσιεννέα, φιάλαι χρυσαῖ τριάκοντα, ἀργυραῖ διοχίλιαι τετρακόσιαι δέκα, καὶ ἄλλα σκεύη χίλια.
\vs{14}Τὰ δὲ πάντα σκεύη ἐκομίσθη, χρυσᾶ καὶ ἀργυρᾶ, πεντακισχίλια τετρακόσια ἑξηκονταεννέα.
\vs{15}Ἀνηνέχθη δὲ ὑπὸ Σαμανασσάρον ἅμα τοῖς ἐκ τῆς αἰχμαλωσίας ἐκ Βαβυλῶνος εἰς Ἰερουσαλήμ.

\vs{16}Ἐν δὲ τοῖς ἐπὶ Ἀρταξέρξου τῶν Περσῶν βασιλέως χρόνοις κατέγραψαν αὐτῷ κατὰ τῶν κατοικούντων ἐν τῇ Ἰουδαίᾳ καὶ Ἱερουσαλὴμ, Βήλεμος, καὶ Μιθραδάτης, καὶ Ταβέλλιος, καὶ Ῥάθυμος, καὶ Βεέλτεθμος, καὶ Σαμέλλιος ὁ γραμματεὺς, καὶ οἱ λοιποὶ οἱ τούτοις συντασσόμενοι, οἰκοῦντες δὲ ἐν Σαμαρείᾳ καὶ τοῖς ἄλλοις τόποις, τὴν ὑπογεγραμμένην ἐπιστολήν·
\vs{17}Βασιλεῖ Ἀρταξέρξῃ κυρίῳ οἱ παῖδές σου, Ῥάθυμος ὁ τὰ προσπίπτοντα, καὶ Σαμέλλιος ὁ γραμματεὺς, καὶ οἱ ἐπίλοιποι τῆς βουλῆς αὐτῶν, καὶ κριταὶ οἱ ἐν κοίλῃ Συρίᾳ καὶ Φοινίκῃ.
\vs{18}Καὶ νῦν γνωστὸν ἔστω τῷ κυρίῳ βασιλεῖ, ὅτι οἱ Ἰουδαῖοι ἀναβάντες παρʼ ὑμῶν πρὸς ἡμᾶς ἐλθόντες εἰς Ἱερουσαλήμ, τὴν πόλιν τήν ἀποστάτιν καὶ πονηρὰν, οἰκοδομοῦσι τάς τε ἀγορὰς αὐτῆς, καὶ τὰ τείχη θεραπεύουσι, καὶ ναὸν ὑποβάλλονται.
\vs{19}Ἐὰν οὖν ἡ πόλις αὕτη οἰκοδομηθῇ, καὶ τὰ τείχη συντελεσθῇ, φορολογίαν οὐ μὴ ὑπομείνωσι δοῦναι, ἀλλὰ καὶ βασιλεῦσιν ἀντιστήσονται.

\vs{20}Καὶ ἐπεὶ ἐνεργεῖται τὰ κατὰ τὸν ναὸν, καλῶς ἔχειν ὑπολαμβάνομεν μὴ ὑπεριδεῖν τὸ τοιοῦτο,
\vs{21}ἀλλὰ προσφωνῆσαι τῷ κυρίῳ βασιλεῖ, ὅπως ἂν φαίνηταί σοι, ἐπισκεφθῇ ἐν τοῖς ἀπὸ τῶν πατέρων σου βιβλίοις.
\vs{22}Καὶ εὑρήσεις ἐν τοῖς ὑπομνηματισμοῖς γεγραμμένα περὶ τούτων, καὶ γνώσῃ ὅτι ἡ πόλις ἐκείνη ἦν ἀποστάτις,
\vs{23}καὶ βασιλεῖς καὶ πόλεις ἐνοχλοῦσα, καὶ οἱ Ἰουδαῖοι ἀποστάται καὶ πολιορκίας συνιστάμενοι ἐν αὐτῇ ἔτι ἐξ αἰῶνος, διʼ ἣν αἰτίαν καὶ ἡ πόλις αὕτη ἠρημώθη.
\vs{24}Νῦν οὖν ὑποδεικνύομέν σοι, κύριε βασιλεῦ, ὅτι ἐὰν ἡ πόλις αὕτη οἰκοδομηθῇ, καὶ τὰ ταύτης τείχη ἀνασταθῇ, κάθοδος οὐκ ἔτι σοι ἔσται εἰς κοίλην Συρίαν καὶ Φοινίκην.

\vs{25}Τότε ἀντέγραψεν ὁ βασιλεὺς Ῥαθύμῳ τῷ γράφοντι τὰ προσπίπτοντα, καὶ Βεελτέθμῳ, καὶ Σαμελλίῳ γραμματεῖ, καὶ τοῖς λοιποῖς τοῖς συντασσομένοις καὶ οἰκοῦσιν ἐν τῇ Σαμαρείᾳ, καὶ Συρίᾳ, καὶ Φοινίκῃ, τὰ ὑπογεγραμμένα.
\vs{26}Ἀνέγνων τὴν ἐπιστολὴν ἣν πεπόμφατε πρὸς μέ· ἐπέταξα οὖν ἐπισκέψασθαι· καὶ εὑρέθη ὅτι ἡ πόλις ἐκείνη ἐστὶν ἐξ αἰῶνος βασιλεῦσιν ἀντιπαρατάσσουσα,
\vs{27}καὶ οἱ ἄνθρωποι ἀποστάσεις καὶ πολέμους ἐν αὐτῇ συντελοῦντες, καὶ βασιλεῖς ἰσχυροὶ καὶ σκληροὶ ἦσαν ἐν Ἱερουσαλὴμ κυριεύοντες καὶ φορολογοῦντες κοίλην Συρίαν καὶ Φοινίκην.
\vs{28}Νῦν οὖν ἐπέταξα ἀποκωλῦσαι τοὺς ἀνθρώπους ἐκείνους τοῦ οἰκοδομῆσαι τὴν πόλιν, καὶ προνοηθῆναι ὅπως μηδὲν παρὰ ταῦτα γένηται.
\vs{29}Καὶ μὴ προβῇ ἐπὶ πλεῖον τὰ τῆς κακίας εἰς τὸ βασιλεῖς ἐνοχλῆσαι.

\vs{30}Τότε ἀναγνωσθέντων τῶν παρὰ τοῦ βασιλέως Ἀρταξέρξου γραφέντων, Ῥάθυμος, καὶ Σαμέλλιος ὁ γραμματεὺς, καὶ οἱ τούτοις συντασσόμενοι, ἀναζεύξαντες εἰς Ἱερουσαλὴμ κατὰ σπουδὴν μεθʼ ἵππου καὶ ὄχλου παρατάξεως, ἤρξαντο κωλύειν τοὺς οἰκοδομοῦντας, καὶ ἤργει ἡ οἰκοδομὴ τοῦ ἱεροῦ τοῦ ἐν Ἱερουσαλὴμ μέχρι τοῦ δευτέρου ἔτους τῆς βασιλείας Δαρείου τοῦ Περσῶν βασιλέως.

\ch{3}
Καὶ βασιλεὺς Δαρεῖος ἐποίησε δοχὴν μεγάλην πᾶσι τοῖς ὑπʼ αὐτὸν, καὶ πᾶσι τοῖς οἰκογενέσιν αὐτοῦ, καὶ πᾶσι τοῖς μεγιστᾶσι τῆς Μηδίας καὶ τῆς Περσίδος,
\vs{2}καὶ πᾶσι τοῖς σατράπαις καὶ στρατηγοῖς καὶ τοπάρχαις τοῖς ὑπʼ αὐτὸν, ἀπὸ τῆς Ἰνδικῆς μέχρις Αἰθιοπίας, ἐν ταῖς ἑκατὸν εἰκοσιεπτὰ σατραπείαις.
\vs{3}Καὶ ἐφάγοσαν καὶ ἐπίοσαν, καὶ ἐμπλησθέντες ἀνέλυσαν· ὁ δὲ Δαρεῖος ὁ βασιλεὺς ἀνέλυσεν εἰς τὸν κοιτῶνα ἑαυτοῦ, καὶ ἐκοιμήθη, καὶ ἔξυπνος ἐγένετο.

\vs{4}Τότε οἱ τρεῖς νεανίσκοι οἱ σωματοφύλακες οἱ φυλάσσοντες τὸ σῶμα τοῦ βασιλέως,
\vs{5}εἶπαν ἕτερος πρὸς τὸν ἕτερον, εἴπωμεν ἕκαστος ἡμῶν ἕνα λόγον, ὃς ὑπερισχύσει· καὶ οὗ ἐὰν φανῇ τὸ ῥῆμα αὐτοῦ σοφώτερον τοῦ ἑτέρου, δώσει αὐτῷ Δαρεῖος ὁ βασιλεὺς δωρεὰς μεγάλας, καὶ ἐπινίκια μεγάλα,
\vs{6}καὶ πορφύραν περιβαλέσθαι, καὶ ἐν χρυσώμασι πίνειν, καὶ ἐπὶ χρυσῷ καθεύδειν, καὶ ἅρμα χρυσοχάλινον, καὶ κίδαριν βυσσίνην, καὶ μανιάκην περὶ τὸν τράχηλον,
\vs{7}καὶ δεύτερος καθιεῖται Δαρείου διὰ τὴν σοφίαν αὐτοῦ, καὶ συγγενὴς Δαρείου κληθήσεται.

\vs{8}Καὶ τότε γράψαντες ἕκαστος τὸν ἑαυτοῦ λόγον, ἐσφαγίσαντο καὶ ἔθηκαν ὑπὸ τὸ προσκεφάλαιον Δαρείου τοῦ βασιλέως,
\vs{9}καὶ εἶπαν, ὅταν ἐγερθῇ ὁ βασιλεὺς, δώσουσιν αὐτῷ τὸ γράμμα, καὶ ὃν ἂν κρίνῃ ὁ βασιλεὺς καὶ οἱ τρεῖς μεγιστᾶνες τῆς Περσίδος, ὅτι οὗ ὁ λόγος αὐτοῦ σοφώτερος, αὐτῷ δοθήσεται τὸ νῖκος καθὼς γέγραπται.
\vs{10}Ὁ εἷς ἔγραψεν, ὑπερισχύει ὁ οἶνος.
\vs{11}Ὁ ἕτερος ἔγραψεν, ὑπερισχύει ὁ βασιλεύς.
\vs{12}Ὁ τρίτος ἔγραψεν, ὑπερισχύουσιν αἱ γυναῖκες, ὑπὲρ δὲ πάντα νικᾷ ἡ ἀλήθεια.

\vs{13}Καὶ ὅτε ἐξηγέρθη ὁ βασιλεὺς, λαβόντες τὸ γράμμα ἔδωκαν αὐτῷ, καὶ ἀνέγνω.
\vs{14}Καὶ ἐξαποστείλας ἐκάλεσε πάντας τοὺς μεγιστᾶνας τῆς Περσίδος καὶ τῆς Μηδείας, καὶ τοὺς σατράπας καὶ στρατηγοὺς, καὶ τοπάρχας καὶ ὑπάτους,
\vs{15}καὶ ἐκάθισεν ἐν τῷ χρηματιστηρίῳ, καὶ ἀνεγνώσθη τὸ γράμμα ἐνώπιον αὐτῶν.
\vs{16}Καὶ εἶπε, καλέσατε τοὺς νεανίσκους, καὶ αὐτοὶ δηλώσουσι τοὺς λόγους ἑαυτῶν· καὶ ἐκλήθησαν, καὶ εἰσήλθοσαν.
\vs{17}Καὶ εἶπαν αὐτοῖς, ἀπαγγείλατε ἡμῖν περὶ τῶν γεγραμμένων.

\vs{18}Καὶ ἤρξατο ὁ πρῶτος ὁ εἴπας περὶ τῆς ἰσχύος τοῦ οἴνου, καὶ ἔφη οὕτως, ἄνδρες,
\vs{19}πῶς ὑπερισχύει ὁ οἶνος; πάντας τοὺς ἀνθρώπους τοὺς πιόντας αὐτὸν πλανᾷ, τὴν διάνοιαν τοῦ τε βασιλέως καὶ τοῦ ὀρφανοῦ ποιεῖ τὴν διάνοιαν μίαν, τήν τε τοῦ οἰκέτου καὶ τὴν τοῦ ἐλευθέρου, τήν τε τοῦ πένητος καὶ τὴν τοῦ πλουσίου·
\vs{20}καὶ πᾶσαν διάνοιαν μεταστρέφει εἰς εὐωχίαν καὶ εὐφροσύνην, καὶ οὐ μέμνηται πᾶσαν λύπην καὶ πᾶν ὀφείλημα·
\vs{21}καὶ πάσας καρδίας ποιεῖ πλουσιας, καὶ οὐ μέμνηται βασιλέα οὐδὲ σατράπην· καὶ πάντα διὰ ταλάντων ποιεῖ λαλεῖν.
\vs{22}Καὶ οὐ μέμνηνται, ὅταν πίνωσι, φιλιάζειν φίλοις καὶ ἀδελφοῖς, καὶ μετʼ οὐ πολὺ σπῶνται τὰς μαχαίρας.
\vs{23}Καὶ ὅταν ἀπὸ τοῦ οἴνου ἐγερθῶσιν, οὐ μέμνηνται ἃ ἔπραξαν.
\vs{24}Ὦ ἄνδρες, οὐχ ὑπερισχύει ὁ οἶνος, ὅτι οὕτως ἀναγκάζει ποιεῖν; καὶ ἐσίγησεν οὕτως εἴπας.

\ch{4}
Καὶ ἤρξατο ὁ δεύτερος λαλεῖν, ὁ εἴπας περὶ τῆς ἰσχύος τοῦ βασιλέως.
\vs{2}Ὦ ἄνδρες, οὐχ ὑπερισχύουσιν οἱ ἄνθρωποι, τὴν γῆν καὶ τὴν θάλασσαν κατακρατοῦντες καὶ πάντα τὰ ἐν αὐτοῖς;
\vs{3}Ὁ δὲ βασιλεὺς ὑπερισχύει, καὶ κυριεύει αὐτῶν καὶ δεσπόζει αὐτῶν, καὶ πᾶν ὃ ἐὰν εἴπῃ αὐτοῖς, ἐνακούουσιν.
\vs{4}Ἐὰν εἴπῃ αὐτοῖς ποιῆσαι πόλεμον ἕτερος πρὸς τὸν ἕτερον, ποιοῦσιν· ἐὰν δὲ ἐξαποστείλῃ αὐτοὺς πρὸς τοὺς πολεμίους, βαδίζουσι καὶ κατεργάζονται τὰ ὄρη καὶ τὰ τείχη καὶ τοὺς πύργους,
\vs{5}φονεύουσι καὶ φονεύονται, καὶ τὸν λόγον τοῦ βασιλέως οὐ παραβαίνουσιν· ἐὰν δὲ νικήσωσι, τῷ βασιλεῖ κομίζουσι πάντα, καὶ ἐὰν προνομεύσωσι, καὶ τὰ ἄλλα πάντα.

\vs{6}Καὶ ὅσοι οὐ στρατεύονται οὐδὲ πολεμοῦσιν, ἀλλὰ γεωργοῦσι τὴν γῆν, πάλιν ὅταν σπείρωσι θερίσαντε, ἀναφέρουσι τῷ βασιλεῖ· καὶ ἕτερος τὸν ἕτερον ἀναγκάζοντες, ἀναφέρουσι τοὺς φόρους τῷ βασιλεῖ.
\vs{7}Καὶ αὐτὸς εἷς μόνος ἐστίν· ἐὰν εἴπῃ ἀποκτεῖναι, ἀποκτέννουσιν· ἐὰν εἴπῃ ἀφεῖναι, ἀφίουσιν·
\vs{8}Εἶπε πατάξαι, τύπτουσιν· εἶπεν ἐρημῶσαι, ἐρημοῦσιν· εἶπεν οἰκοδομῆσαι, οἰκοδομοῦσιν·
\vs{9}εἶπεν ἐκκόψαι, ἐκκόπτουσιν· εἶπε φυτεῦσαι, φυτεύουσι.
\vs{10}Καὶ πᾶς ὁ λαὸς αὐτοῦ καὶ αἱ δυνάμεις αὐτοῦ ἐνακούουσι· πρὸς δὲ τούτοις αὐτὸς ἀνάκειται, ἐσθίει καὶ πίνει καὶ καθεύδει,
\vs{11}αὐτοὶ δὲ τηροῦσι κύκλῳ περὶ αὐτόν· καὶ οὐ δύνανται ἕκαστος ἀπελθεῖν, καὶ ποιεῖν τὰ ἔργα αὐτοῦ, οὐδὲ παρακούουσιν αὐτοῦ.
\vs{12}Ὦ ἄνδρες, πῶς οὐχ ὑπερισχύει ὁ βασιλεὺς, ὅτι οὕτως ἐπάκουστός ἐστι; καὶ ἐσίγησεν.

\vs{13}Ὁ δὲ τρίτος ὁ εἴπας περὶ τῶν γυναικῶν καὶ τῆς ἀληθείας, οὗτός ἐστι Ζοροβάβελ, ἤρξατο λαλεῖν·
\vs{14}Ἄνδρες, οὐ μέγας ὁ βασιλεὺς, καὶ πολλοὶ οἱ ἄνθρωποι, καὶ ὁ οἶνος ἰσχύει; τίς οὖν ὁ δεσπόζων αὐτῶν, ἢ τίς ὁ κυριεύων αὐτῶν; οὐχ αἱ γυναῖκες;
\vs{15}Αἱ γυναῖκες ἐγέννησαν τὸν βασιλέα καὶ πάντα τὸν λαὸν ὃς κυριεύει τῆς θαλάσσης καὶ τῆς γῆς,
\vs{16}καὶ ἐξ αὐτῶν ἐγένοντο· καὶ αὗται ἐξέθρεψαν αὐτοὺς τοὺς φυτεύσαντας τοὺς ἀμπελῶνας ἐξ ὧν ὁ οἶνος γίνεται.
\vs{17}Καὶ αὗται ποιοῦσι τὰς στολὰς τῶν ἀνθρώπων, καὶ αὗται ποιοῦσι δόξαν τοῖς ἀνθρώποις, καὶ οὐ δύνανται οἱ ἄνθρωποι χωρὶς τῶν γυναικῶν εἶναι.
\vs{18}Ἐὰν δὲ συναγάγωσι χρυσίον καὶ ἀργύριον καὶ πᾶν πρᾶγμα ὡραῖον, καὶ ἴδωσι γυναῖκα μίαν καλὴν τῷ εἴδει καὶ τῷ κάλλει,
\vs{19}ταῦτα πάντα ἀφέντες, εἰς αὐτὴν ἐκκέχῃναν, καὶ χάσκοντες τὸ στόμα θεωροῦσιν αὐτὴν, καὶ πάντες αὐτὴν αἱρετίζουσι μᾶλλον ἢ τὸ χρυσίον καὶ τὸ ἀργύριον καὶ πᾶν πρᾶγμα ὡραῖον.

\vs{20}Ἄνθρωπος τὸν ἑαυτοῦ πατέρα ἐγκαταλείπει ὃς ἐξέθρεψεν αὐτὸν, καὶ τὴν ἰδίαν χώραν, καὶ πρὸς τὴν ἰδίαν γυναῖκα κολλᾶται,
\vs{21}καὶ μετὰ τῆς γυναικὸς ἀφίησι τὴν ψυχήν, καὶ οὔτε τὸν πατέρα μέμνηται, οὔτε τὴν μητέρα, οὔτε τὴν χώραν.
\vs{22}Καὶ ἐντεῦθεν δεῖ ὑμᾶς γνῶναι ὅτι αἱ γυναῖκες κυριεύουσιν ὑμῶν· οὐχὶ πονεῖτε, καὶ μοχθεῖτε, καὶ πάντα ταῖς γυναιξὶ δίδοτε, καὶ φέρετε;
\vs{23}Και λαμβάνει ὁ ἄνθρωπος τὴν ῥομφαίαν αὐτοῦ, καὶ ἐκπορεύεται ἐξοδεύειν καὶ λῃστεύειν καὶ κλέπτειν, καὶ εἰς τὴν θάλασσαν πλεῖν, καὶ ποταμοὺς,
\vs{24}καὶ τὸν λέοντα θεωρεῖ, καὶ ἐν σκότει βαδίζει· καὶ ὅταν κλέψῃ καὶ ἁρπάσῃ καὶ λωποδυτήσῃ, τῇ ἐρωμένῃ ἀποφέρει.
\vs{25}Καὶ πλεῖον ἀγαπᾷ ἄνθρωπος τὴν ἰδίαν γυναῖκα μᾶλλον ἢ τὸν πατέρα καὶ τὴν μητέρα.
\vs{26}Καὶ πολλοὶ ἀπενοήθησαν ταῖς ἰδίαις διανοίαις διὰ τὰς γυναῖκας, καὶ δοῦλοι ἐγένοντο διʼ αὐτάς·
\vs{27}καὶ πολλοὶ ἀπώλοντο καὶ ἐσφάλησαν καὶ ἡμάρτοσαν διὰ τὰς γυναῖκας.

\vs{28}Καὶ νῦν οὐ πιστεύετέ μοι; οὐχὶ μέγας ὁ βασιλεὺς τῇ ἐξουσίᾳ αὐτοῦ; οὐχὶ πᾶσαι αἱ χῶραι εὐλαβοῦνται ἅψασθαι αὐτοῦ;
\vs{29}Ἐθεώρουν αὐτὸν, καὶ Ἀπάμην τὴν θυγατέρα Βαρτάκου τοῦ θαυμαστοῦ, τὴν παλλακὴν τοῦ βασιλέως, καθημένην ἐν δεξιᾷ τοῦ βασιλέως,
\vs{30}καὶ ἀφαιροῦσαν τὸ διάδημα ἀπὸ τῆς κεφαλῆς τοῦ βασιλέως, καὶ ἐπιτιθοῦσαν ἑαυτῇ· καὶ ἐῤῥάπιζε τὸν βασιλέα τῇ ἀριστερᾷ.
\vs{31}Καὶ πρὸς τούτοις ὁ βασιλεὺς χάσκων τὸ στόμα ἐθεώρει αὐτήν· καὶ ἐὰν προσγελάσῃ αὐτῷ, γελᾷ· ἐὰν δὲ πικρανθῇ ἐπʼ αὐτὸν, κολακεύει αὐτὴν, ὅπως διαλλαγῇ αὐτῷ.
\vs{32}Ὦ ἄνδρες, πῶς οὐχὶ ἰσχυραὶ αἱ γυναῖκες, ὅτι οὕτως πράσσουσι;

\vs{33}Καὶ τότε ὁ βασιλεὺς καὶ οἱ μεγιστᾶνες ἔβλεπον εἷς τὸν ἕτερον· καὶ ἤρξατο λαλεῖν περὶ τῆς ἀληθείας·
\vs{34}Ἄνδρες, οὐχὶ ἰσχυραὶ αἱ γυναῖκες; μεγάλη ἡ γῆ, καὶ ὑψηλὸς ὁ οὐρανὸς, καὶ ταχὺς τῷ δρόμῳ ὁ ἥλιος, ὅτι στρέφεται ἐν τῷ κύκλῳ τοῦ οὐρανοῦ, καὶ πάλιν ἀποτρέχει εἰς τὸν ἑαυτοῦ τόπον ἐν μιᾷ ἡμέρᾳ.
\vs{35}Οὐχὶ μέγας ὃς ταῦτα ποιεῖ; καὶ ἡ ἀλήθεια μεγάλη καὶ ἰσχυροτέρα παρὰ πάντα.
\vs{36}Πᾶσα ἡ γῆ τὴν ἀλὴθειαν καλεῖ, καὶ ὁ οὐρανὸς αὐτὴν εὐλογεῖ, καὶ πάντα τὰ ἔργα σείεται καὶ τρέμει, καὶ οὐκ ἔστι μετʼ αὐτῆς ἄδικον οὐδέν.
\vs{37}Ἄδικος ὁ οἶνος, ἄδικος ὁ βασιλεὺς, ἄδικοι αἱ γυναῖκες, ἄδικοι πάντες οἱ υἱοὶ τῶν ἀνθρώπων, καὶ ἄδικα πάντα τὰ ἔργα αὐτῶν τὰ τοιαῦτα, καὶ οὐκ ἔστιν ἐν αὐτοῖς ἀλήθεια, καὶ ἐν τῇ ἀδικίᾳ αὐτῶν ἀπολοῦνται.

\vs{38}Καὶ ἡ ἀλήθεια μένει καὶ ἰσχύει εἰς τὸν αἰῶνα, καὶ ζῇ καὶ κρατεῖ εἰς τὸν αἰῶνα τοῦ αἰῶνος.
\vs{39}Καὶ οὐκ ἔστι παρʼ αὐτὴν λαμβάνειν πρόσωπα, οὐδὲ διάφορα, ἀλλὰ τὰ δίκαια ποιεῖ ἀπὸ πάντων τῶν ἀδίκων καὶ πονηρῶν· καὶ πάντες εὐδοκοῦσι τοῖς ἔργοις αὐτῆς,
\vs{40}καὶ οὐκ ἔστιν ἐν τῇ κρίσει αὐτῆς οὐδὲν ἄδικον· καὶ αὕτη, ἡ ἰσχὺς, καὶ τὸ βασίλειον, καὶ ἡ ἐξουσία, καὶ ἡ μεγαλειότης τῶν πάντων αἰώνων· εὐλογητὸς ὁ Θεὸς τῆς ἀληθείας.

\vs{41}Καὶ ἐσιώπησε τοῦ λαλεῖν· καὶ πᾶς ὁ λαὸς τότε ἐφώνησε· καὶ τότε εἶπον, μεγάλη ἡ ἀλήθεια, καὶ ὑπερισχύει·
\vs{42}τότε ὁ βασιλεὺς εἶπεν αὐτῷ, αἴτησαι ὃ θέλεις πλείω τῶν γεγραμμένων, καὶ δώσομέν σοι ὅν τρόπον εὑρέθης σοφώτερος, καὶ ἐχόμενός μου καθήσῃ, καὶ συγγενής μου κληθήθῃ.
\vs{43}Τότε εἶπε τῷ βασιλεῖ, μνήσθητι τὴν εὐχὴν, ἣν ηὔξω, οἰκοδομῆσαι τὴν Ἱερουσαλὴμ ἐν τῇ ἡμέρᾳ ᾗ τὸ βασίλειόν σου παρέλαβες,
\vs{44}καὶ πάντα τὰ σκεύη τὰ ληφθέντα ἐξ Ἱερουσαλὴμ, καὶ ἐκπέμψαι ἃ ἐχώρισε Κύρος, ὅτε ηὔξατο ἐκκόψαι Βαβυλῶνα, καὶ ηὔξατο ἐξαποστεῖλαι ἐκεῖ.
\vs{45}Καὶ σὺ ηὔξω οἰκοδομῆσαι τὸν ναὸν ὃν ἐνεπύρισαν οἱ Ἰδουμαῖοι, ὅτε ἠρημώθη ἡ Ἰουδαία ὑπὸ τῶν Χαλδαίων.
\vs{46}Καὶ νῦν τοῦτό ἐστιν ὅ σε ἀξιῶ, κύριε βασιλεῦ, καὶ ὃ αἰτοῦμαί σε, καὶ αὕτη ἐστὶν, ἡ μεγαλωσύνη ἡ παρὰ σοῦ· δέομαι οὖν ἵνα ποιήσῃς τὴν εὐχὴν, ἣν ηὔξω τῷ βασιλεῖ τοῦ οὐρανοῦ, ποιῆσαι ἐκ στόματός σου.

\vs{47}Τότε ἀναστὰς Δαρεῖος ὁ βασιλεὺς κατεφίλησεν αὐτὸν, καὶ ἔγραψεν αὐτῷ τὰς ἐπιστολὰς πρὸς πάντας τοὺς οἰκονόμους, καὶ τοπάρχας, καὶ στρατηγοὺς, καὶ σατράπας, ἵνα προπέμψωσιν αὐτὸν καὶ τοὺς μετʼ αὐτοῦ πάντας ἀναβαίνοντας οἰκοδομῆσαι τὴν Ἱερουσαλήμ.
\vs{48}Καὶ πᾶσι τοῖς τοπάρχαις ἐν κοίλῃ Συρίᾳ, καὶ Φοινίκῃ, καὶ τοῖς ἐν τῷ Λιβάνω ἔγραψεν ἐπιστολὰς, μεταφέρειν ξύλα κέδρινα ἀπὸ τοῦ Λιβάνου εἰς Ἱερουσαλὴμ, καὶ ὅπως οἰκοδομήσωσι μετʼ αὐτοῦ τὴν πόλιν.

\vs{49}Καὶ ἔγραψε πᾶσι τοῖς Ἰουδαίοις τοῖς ἀναβαίνουσιν ἀπὸ τῆς βασιλείας εἰς τὴν Ἰουδαίαν ὑπὲρ τῆς ἐλευθερίας, πάντα δυνατὸν, καὶ τοπάρχην, καὶ σατράπην, καὶ οἰκονόμον μὴ ἐπελεύσεσθαι ἑπὶ τὰς θύρας αὐτῶν,
\vs{50}καὶ πᾶσαν τὴν χώραν ἣν κρατοῦσιν, ἀφορολόγητον αὐτοῖς ὑπάρχειν· καὶ ἵνα οἱ Ἰδουμαῖοι ἀφίωσι τὰς κώμας ἃ διακρατοῦσι τῶν Ἰουδαίων·
\vs{51}καὶ εἰς τὴν οἰκοδομὴν τοῦ ἱεροῦ δοθῆναι κατʼ ἐνιαυτὸν τάλαντα εἴκοσι, μέχρι τοῦ οἰκοδομηθήναι·
\vs{52}καὶ ἐπὶ τὸ θυσιαστήριον ὁλοκαυτώματα καρποῦσθαι καθʼ ἡμέραν, καθὰ ἔχουσιν ἐντολὴν, ἑπτακαίδεκα προσφέρειν ἄλλα τάλαντα, δέκα κατʼ ἐνιαυτόν.
\vs{53}καὶ πᾶσι τοῖς προσβαίνουσιν ἀπὸ τῆς Βαβυλωνίας κτίσαι τὴν πόλιν, ὑπάρχειν τὴν ἐλευθερίαν αὐτοῖς τε καὶ τοῖς ἐκγόνοις αὐτῶν, καὶ πᾶσι τοῖς ἱερεῦσι τοῖς προσβαίνουσιν.
\vs{54}Ἔγραψε δὲ καὶ τὴν χορηγίαν καὶ τὴν ἱερατικὴν στολὴν ἐν τίνι λατρεύουσιν ἐν αὐτῇ.
\vs{55}Καὶ τοῖς Λευίταις ἔγραψε δοῦναι τὴν χορηγίαν, ἕως τῆς ἡμέρας ἡς ἐπιτελεσθῇ ὁ οἶκος καὶ Ἱερουσαλὴμ οἰκοδομηθῆναι.
\vs{56}Καὶ πᾶσι τοῖς φρουροῦσι τὴν πόλιν ἔγραψε δοῦναι αὐτοῖς κλήρους καὶ ὀψώνια.
\vs{57}Καὶ ἐξαπέστειλε πάντα τὰ σκεύη ἃ ἐχώρισε Κύρος ἀπὸ Βαβυλῶνος· καὶ πάντα ὅσα εἶπε Κύρος ποιῆσαι, καὶ αὐτὸς ἐπέταξε ποιῆσαι, καὶ ἐξαποστεῖλαι εἰς Ἱερουσαλήμ.

\vs{58}Καὶ ὅτε ἐξῆλθεν ὁ νεανίσκος, ἄρας τὸ πρόσωπον εἰς τὸν οὐρανὸν ἐναντίον Ἱερουσαλὴμ, εὐλόγησε τῷ βασιλεῖ τοῦ οὐρανοῦ, λέγων,
\vs{59}παρὰ σοῦ νίκη, καὶ παρὰ σοῦ ἡ σοφία, καὶ σὴ ἡ δόξα, καὶ ἐγὼ σὸς οἰκέτης.
\vs{60}Εὐλογητὸς εἶ, ὃς ἔδωκάς μοι σοφίαν, καὶ σοὶ ὁμολογῶ, δέσποτα τῶν πατέρων.
\vs{61}Καὶ ἔλαβε τὰς ἐπιστολὰς, καὶ ἐξῆλθε, καὶ ἦλθεν εἰς Βαβυλῶνα, καὶ ἀπήγγειλε τοῖς ἀδελφοῖς αὐτοῦ πᾶσι.
\vs{62}Καὶ εὐλόγησαν τὸν Θεὸν τῶν πατέρων αὐτῶν, ὅτι ἔδωκεν αὐτοῖς ἄνεσιν καὶ ἄφεσιν,
\vs{63}ἀναβῆναι καὶ οἰκοδομῆσαι τὴν Ἱερουσαλὴμ καὶ τὸ ἱερὸν, οὗ ὠνομάσθη τὸ ὄνομα αὐτοῦ ἐπʼ αὐτῷ· καὶ ἐκωθωνίζοντο μετὰ μουσικῶν καὶ χαρᾶς ἡμέρας ἑπτά.

\ch{5}
Μετὰ δὲ ταῦτα ἐξελέγησαν ἀναβῆναι ἀρχηγοὶ οἴκου πατριῶν κατὰ φυλὰς αὐτῶν, καὶ αἱ γυναῖκες αὐτῶν, καὶ οἱ υἱοὶ αὐτῶν, καὶ αἱ θυγατέρες, καὶ οἱ παῖδες αὐτῶν, καὶ αἱ παιδίσκαι, καὶ τὰ κτήνη αὐτῶν.
\vs{2}Καὶ Δαρεῖος συναπέστειλε μετʼ αὐτῶν ἱππεῖς χιλίους, ἕως τοῦ ἀποκαταστῆσαι αὐτοὺς εἰς Ἱερουσαλὴμ μετʼ εἰρήνης, καὶ μετὰ μουσικῶν, τυμπάνων, καὶ αὐλῶν.
\vs{3}Καὶ πάντες οἱ ἀδελφοὶ αὐτῶν παίζοντες, καὶ ἐποίησεν αὐτοὺς συναναβῆναι μετʼ ἐκείνων.

\vs{4}Καὶ ταῦτα τὰ ὀνόματα τῶν ἀνδρῶν τῶν ἀναβαινόντων κατὰ πατριὰς αὐτῶν εἰς τὰς φυλὰς, ἐπὶ τὴν μεριδαρχίαν αὐτῶν.
\vs{5}Οἱ ἱερεῖς υἱοὶ Φινεὲς, υἱοὶ Ἀαρὼν, Ἰησοῦς ὁ τοῦ Ἰωσεδὲκ τοῦ Σαραίου, καὶ Ἰωακὶμ ὁ τοῦ Ζοροβάβελ τοῦ Σαλαθιὴλ ἐκ τοῦ οἴκου τοῦ Δαυὶδ, ἐκ τῆς γενεᾶς Φαρὲς, φυλῆς δὲ Ἰούδα,
\vs{6}ὃς ἐλάλησεν ἐπὶ Δαρείου τοῦ βασιλέως Περσῶν λόγους σοφοὺς ἐν τῷ δευτέρῳ ἔτει τῆς βασιλείας αὐτοῦ, μηνὶ Νισὰν τοῦ πρώτου μηνός.
\vs{7}Εἰσὶ δὲ οὗτοι οἱ ἐκ τῆς γῆς Ἰουδαίας ἀναβάντες ἐκ τῆς αἰχμαλωσίας τῆς παροικίας, οὓς μετῴκισε Ναβουχοδονόσορ βασιλεὺς Βαβυλῶνος εἰς Βαβυλῶνα.
\vs{8}Καὶ ἐπέστρεψαν εἰς Ἱερουσαλὴμ καὶ τὴν λοιπὴν Ἰουδαίαν ἕκαστος εἰς τὴς ἰδίαν πόλιν, οἱ ἐλθόντες μετὰ Ζοροβάβελ, καὶ Ἰησοῦ, Νεεμίου, Ζαραίου, Ῥησαίου, Ἐνηνέος, Μαρδοχαίου, Βεελσάρου, Ἀσφαράσου, Ῥεελίου, Ῥοΐμου, Βαανὰ, τῶν προηγουμένων αὐτῶν.

\vs{9}Ἀριθμὸς τῶν ἀπὸ τοῦ ἔθνους καὶ οἱ προηγούμενοι αὐτῶν· υἱοὶ Φόρος, δύο χιλιάδες καὶ ἑκατὸν ἑβδομηκονταδύο· υἱοὶ Σαφὰτ, τετρακόσιοι ἑβδομηκονταδύο.

\vs{10}Υἱοὶ Ἀρὲς, ἑπτακόσιοι πεντηκονταέξ.

\vs{11}Υἱοὶ Φαὰθ Μωὰβ εἰς τοὺς υἱοὺς Ἰησοῦ καὶ Ἰωὰβ, δισχίλιοι ὀκτακόσιοι δεκαδύο.

\vs{12}Υἱοὶ Ἠλὰμ, χίλιοι διακόσιοι πεντηκοντατέσσαρες· υἱοὶ Ζαθουῒ, ἐννακόσιοι ἑβδομηκονταπέντε· υἱοὶ Χορβὲ, ἑπτακόσιοι πέντε· υἱοὶ Βανί, ἑξακόσιοι τεσσαρακονταοκτώ.

\vs{13}Υἱοὶ Βηβαὶ, ἑξακόσιοι τριακοντατρεῖς· υἱοὶ Ἀργαὶ, χίλιοι τριακόσιοι εἰκοσιδύο.

\vs{14}Υἱοὶ Ἀδωνικὰν, ἑξακόσιοι τριακονταεπτά. υἱοὶ Βαγοῒ, δισχίλιοι ἑξακόσιοι ἕξ· υἱοὶ Ἀδινοὺ, τετρακόσιοι πεντηκοντατέσσαρες·

\vs{15}Υἱοὶ Ἀτὴρ Ἐζεκίου, ἐννενηκονταδύο· υἱοὶ Κιλὰν, καὶ Ἀζηνὰν, ἑξηκονταεπτά· υἱοὶ Ἀζαροὺ, τετρακόσιοι τριακονταδύο.

\vs{16}Υἱοὶ Ἁννὶς, ἑκατὸν εἷς· υἱοὶ Ἀρὸμ, τριακονταδύο· υἱοὶ Βασσαὶ, τριακόσιοι εἰκοσιτρεῖς· υἱοὶ Ἀρσιφουρὶθ, ἑκατὸν δύο.

\vs{17}Υἱοὶ Βαιτηροὺς, τρισχίλιοι πέντε· υἱοὶ ἐκ Βαιθλωμῶν, ἑκατὸν εἰκοσιτρεῖς.

\vs{18}Οἱ ἐκ Νετωφὰς, πεντηκονταπέντε· οἱ ἐξ Ἀναθὼθ, ἑκατὸν πεντηκονταοκτώ· οἱ ἐκ Βαιθασμὼν, τεσσαρακονταδύο.

\vs{19}Οἱ ἐκ Καριαθιρὶ, εἰκοσιπέντε· οἱ ἐκ Καφείρας, καὶ Βηρὼγ, ἑπτακόσιοι τεσσαρακοντατρεῖς.

\vs{20}Οἱ Χαδιασαὶ καὶ Ἁμμίδιοι, τετρακόσιοι εἰκοσιδύο· οι ἐκ Κιραμᾶς καὶ Γαββῆς, ἑξακόσιοι εἴκοσι εἷς.

\vs{21}Οἱ ἐκ Μακαλὼν, ἑκατὸν εἰκοσιδύο· οἱ ἐκ Βετολίῳ, πεντηκονταδύο· υἱοὶ Νιφὶς, ἑκατὸν πεντηκονταέξ.

\vs{22}Υἱοὶ Καλαμωλάλου, καὶ Ὠνοὺς, ἑπτακόσιοι εἰκοσιπέντε· υἱοὶ Ἱερεχοὺ, διακόσιοι τεσσαρακονταπέντε.

\vs{23}Υἱοὶ Σανάας, τρισχίλιοι τριακόσιοι εἷς.

\vs{24}Οἱ ἱερεῖς οἱ υἱοὶ Ἰεδδοὺ τοῦ Ἰησοῦ εἰς τοὺς υἱοὺς Σανασὶβ, ὀκτακόσιοι ἑβδομηκονταδύο· υἱοὶ Ἐμμηροὺθ, διακόσιοι πεντηκονταδύο.

\vs{25}Υἱοὶ Φασσούρου, χίλιοι τεσσαρακονταεπτά· υἱοὶ Χαρμι, διακόσιοι δεκαεπτά.

\vs{26}Οἱ Λευῖται οἱ υἱοὶ Ἰησοῦ, καὶ Καδοήλου, καὶ Βάννου, καὶ Σουδίου, ἑβδομηκοντατέσσαρες·

\vs{27}Οἱ ἱεροψάλται υἱοὶ Ἀσὰρ, ἑκατὸν εἰκοσιοκτώ.

\vs{28}Οἱ θυρωροὶ υἱοὶ Σαλοὺμ, υἱοὶ Ἀτὰρ, υἱοὶ Τολμὰν, υἱοὶ Δακοὺβ, Ἀτητὰ, υἱοὶ Τωβὶς, πάντες ἑκατὸν τριακονταεννέα.

\vs{29}Οἱ ἱερόδουλοι, υἱοὶ Ἡσαὺ, υἱοὶ Ἀσιφὰ, υἱοὶ Ταβαὼθ, υἱοὶ Κηρὰς, υἱοὶ Σουδὰ, υἱοὶ Φαλαίου, υἱοὶ Λαβανὰ, υἱοὶ Ἀγραβὰ,

\vs{30}Υἱοὶ Ἀκοὺδ, υἱοὶ Οὐτὰ, υἱοὶ Κητὰβ, υἱοὶ Ἀκκαβὰ, υἱοὶ Συβαῒ, υἱοὶ Ἀνὰν, υἱοὶ Καθουὰ, υἱοὶ Γεδδοὺρ,

\vs{31}Υἱοὶ Ἰαΐρου, υἱοὶ Δαισὰν, υἱοὶ Νοεβὰ, υἱοὶ Χασεβά, υἱοὶ Καζηρὰ, υἱοὶ Ὀζίου, υἱοὶ Φινοὲ, υἱοὶ Ἀσαρὰ, υἱοὶ Βασθαῒ, υἱοὶ Ἀσσανὰ, υἱοὶ Μανὶ, υἱοὶ Ναφισὶ, υἱοὶ Ἀκοὺφ, υἱοὶ Ἀχιβὰ, υἱοὶ Ἀσοὺβ, υἱοὶ Φαρακὲμ, υἱοὶ Βασαλὲμ,

\vs{32}Υἱοὶ Μεεδδὰ, υἱοὶ Κουθὰ, υἱοὶ Χαρέα, υἱοὶ Βαρχουὲ, υἱοὶ Σερὰρ, υἱοὶ Θομοῒ, υἱοὶ Νασί, υἱοὶ Ἀτεφά·

\vs{33}Υἱοὶ παίδων Σαλωμὼν, υἱοὶ Ἀσσαπφιὼθ, υἱοὶ Φαριρὰ, υἱοὶ Ἰειηλὶ, υἱοὶ Λοζὼν, υἱοὶ Ἰσδαὴλ, υἱοὶ Σαφυῒ,

\vs{34}Υἱοὶ Ἁγιὰ, υἱοὶ Φακαρὲθ, υἱοὶ Σαβιὴ, υἱοὶ Σαρωθὶ, υἱοὶ Μισαίας, υἱοὶ Γὰς, υἱοὶ Ἀδδοὺς, υἱοὶ Σουβὰ, υἱοὶ Ἀφεῤῥὰ, υἱοὶ Βαρωδὶς, υἱοὶ Σαφὰγ, υἱοὶ Ἀλλώμ·

\vs{35}Πάντες οἱ ἱερόδουλοι, καὶ οἱ υἱοὶ τῶν παίδων Σαλωμὼν τριακόσιοι ἑβδομηκονταδύο.

\vs{36}Οὗτοι ἀναβάντες ἀπὸ Θερμελὲθ, καὶ Θελερσὰς, ἡγούμενος αὐτῶν Χαρααθαλὰν, καὶ Ἀλάρ.
\vs{37}Καὶ οὐκ ἠδύναντο ἀπαγγεῖλαι τὰς πατριὰς αὐτῶν καὶ γενεὰς, ὡς ἐκ τοῦ Ἰσραήλ εἰσιν· υἱοὶ Δαλὰν τοῦ υἱοῦ τοῦ Βαενὰν, υἱοὶ Νεκωδὰν, ἑξακόσιοι πεντη· κονταδύο.

\vs{38}Καὶ ἐκ τῶν ἱερέων οἱ ἐμποιούμενοι ἱερωσύνης, καὶ οὐχ εὑρέθησαν, υἱοὶ Ὀβδία, υἱοὶ Ἁκβὼς, υἱοὶ Ἰαδδοὺ τοῦ λαβόντος Αὐγίαν γυναῖκα τῶν θυγατέρων Φαηζελδαίου, καὶ ἐκλήθη ἐπὶ τῷ ὀνόματι αὐτοῦ.

\vs{39}Καὶ τούτων ζητηθείσης τῆς γενικῆς γραφῆς ἐν τῷ καταλοχισμῷ καὶ μὴ εὑρεθείσης, ἐχωρίσθησαν τοῦ ἱερατεύειν.
\vs{40}Καὶ εἶπεν αὐτοῖς Νεεμίας καὶ Ἀτθαρίας, μὴ μετέχειν τῶν ἁγίων ἕως ἀναστῇ ἀρχιερεὺς ἐνδεδυμένος τὴν δήλωσιν καὶ τὴν ἀλήθειαν.

\vs{41}Οἱ δὲ πάντες Ἰσραὴλ ἦσαν ἀπὸ δωδεκαετοῦς καὶ ἐπάνω χωρὶς παίδων καὶ παιδισκῶν, μυριάδες τέσσαρες δισχίλιοι τριακόσιοι ἑξήκοντα.
\vs{42}Παῖδες τούτων καὶ παιδίσκαι, ἑπτακισχίλιοι τριακόσιοι τριακονταεπτά. ψάλται καὶ ψαλτῳδοὶ, διακόσιοι τεσσαρακονταπέντε·
\vs{43}Κάμηλοι τετρακόσιοι τριακονταπέντε, καὶ ἵπποι ἑπτακισχίλιοι τριακονταὲξ, ἡμίονοι διακόσιοι τεσσαράκοντα πέντε, ὑποζύλια πεντακισχίλια πεντακόσια εἰκοσιπέντε.

\vs{44}Καὶ ἐκ τῶν ἡγουμένων κατὰ τὰς πατριὰς ἐν τῷ παραγίνεσθαι αὐτούς εἰς τὸ ἱερὸν τοῦ Θεοῦ τὸ ἐν Ἱερουσαλὴμ, ηὔξαντο ἐγεῖραι τὸν οἶκον ἐπὶ τοὺ τοῦ τόπου αὐτοῦ κατὰ τὴν αὐτῶν δύναμιν,
\vs{45}καὶ δοῦναι εἰς τὸ ἱερὸν γαζοφυλάκιον τῶν ἔργων, χρυσίου μνᾶς χιλίας καὶ ἀργυρίου μνᾶς πεντακισχιλίας, καὶ στολὰς ἱερατικὰς ἑκατόν.
\vs{46}Καὶ κατῳκίσθησαν οἱ ἱερεῖς, καὶ οἱ Λευεῖται, καὶ οἱ ἐκ τοῦ λαοῦ αὐτοῦ ἐν Ἱερουσαλὴμ καὶ τῇ χώρᾳ, οἵ τε ἱεροψάλται, καὶ οἱ θυρωροὶ, καὶ πᾶς Ἰσραὴλ ἐν ταῖς κώμαις αὐτῶν.

\vs{47}Ἐνστάντος δὲ τοῦ ἑβδόμον μηνὸς, καὶ ὄντων τῶν υἱῶν Ἰσραὴλ ἑκάστου ἐν τοῖς ἰδίοις, συνήχθησαν ὁμοθυμαδὸν εἰς τὸ εὐρύχωρον τοῦ πρώτου πυλῶνος τοῦ πρὸς τῇ ἀνατολῇ.
\vs{48}Καὶ καταστὰς Ἰησοῦς ὁ τοῦ Ἰωσεδὲκ καὶ οἱ ἀδελφοὶ αὐτοῦ οἱ ἱερεῖς, καὶ Ζοροβάβελ ὁ τοῦ Σαλαθιὴλ καὶ οἱ τούτου ἀδελφοὶ, ἡτοίμασαν τὸ θυσιαστήριον τοῦ Θεοῦ Ἰσραὴλ,
\vs{49}προσενέγκαι ἐπʼ αὐτοῦ ὁλοκαυτώσεις, ἀκολούθως τοῖς ἐν τῇ Μωσέως βίβλῳ τοῦ ἀνθρώπου τοῦ Θεοῦ διηγορευμένοις.

\vs{50}Καὶ ἐπισυνήχθησαν αὐτοῖς ἐκ τῶν ἄλλων ἐθνῶν τῆς γῆς, καὶ κατώρθωσαν τὸ θυσιαστήριον ἐπὶ τοῦ τόπου αὐτῶν, ὅτι ἐν ἔχθρᾳ ἦσαν αὐτοῖς, καὶ κατίσχυσαν αὐτοὺς πάντα τὰ ἔθνη τὰ ἐπὶ τῆς γῆς· καὶ ἀνέφερον θυσίας κατὰ τὸν καιρὸν, καὶ ὁλοκαυτώματα Κυρίῳ τὸ πρωϊνὸν καὶ τὸ δειλινόν.
\vs{51}Καὶ ἐγάγοσαν τὴν τῆς σκηνοπηγίας ἑορτὴν, ὡς ἐπιτέτακται ἐν τῷ νόμῳ, καὶ θυσίας καθʼ ἡμέραν, ὡς προσῆκον ἦν·
\vs{52}καὶ μετὰ ταῦτα προσφορὰς ἐνδελεχισμοῦ, καὶ θυσίας σαββάτων καὶ νουμηνιῶν καὶ ἑορτῶν πασῶν ἡγιασμένων.

\vs{53}Καὶ ὅσοι ηὔξαντο εὐχὴν τῷ Θεῷ ἀπὸ τῆς νουμηνίας τοῦ ἑβδόμου μηνὸς, ἤρξαντο προσφέρειν θυσίας τῷ Θεῷ, καὶ ὁ ναὸς τοῦ Θεοῦ οὔπω ᾠκοδόμητο.

\vs{54}Καὶ ἔδωκαν ἀργύριον τοῖς λατόμοις καὶ τέκτοσι, καὶ ποτὰ καὶ βρωτὰ,
\vs{55}καὶ χάῤῥα τοῖς Σιδωνίοις καὶ Τυρίοις εἰς τὸ παράγειν αὐτοὺς ἐκ τοῦ Λιβάνου ξύλα κέδρινα, διαφέρειν σχεδίας εἰς τὸν Ἰόππης λιμένα, κατὰ τὸ πρόσταγμα τὸ γραφὲν αὐτοῖς παρὰ Κύρου τοῦ Περσῶν βασιλέως.

\vs{56}Καὶ τῷ δευτέρῳ ἔτει παραγενόμενος εἰς τὸ ἰερὸν τοῦ Θεοῦ εἰς Ἱερονσαλὴμ μηνὸς δευτέρου, ἤρξατο Ζοροβάβελ ὁ τοῦ Σαλαθιὴλ, καὶ Ἰησοῦς ὁ τοῦ Ἰωσεδὲκ, καὶ οἱ ἀδελφοὶ αὐτῶν, καὶ οἱ ἱερεῖς οἱ Λευῖται, καὶ πάντες οἱ παραγενόμενοι ἐκ τῆς αἰχμαλωσίας εἰς Ἱερουσαλὴμ,
\vs{57}καὶ ἐθεμελίωσαν τὸν ναὸν τοῦ Θεοῦ τῇ νουμηνίᾳ τοῦ δευτέρου μηνὸς τοῦ δευτέρου ἔτους, ἐν τῷ ἐλθεῖν εἰς τὴν Ἰουδαίαν καὶ Ἱερουσαλήμ.
\vs{58}Καὶ ἔστησαν τοὺς Λευίτας ἀπὸ εἰκοσαετοῦς ἐπὶ τῶν ἔργων τοῦ Κυρίου· καὶ ἔστη Ἰησοῦς, καὶ οἱ υἱοὶ, καὶ οἱ ἀδελφοὶ, καὶ Καδμιὴλ ὁ ἀδελφὸς, καὶ οἱ υἱοὶ Ἠμαδαβοὺν, καὶ οἱ υἱοὶ Ἰωδὰ τοῦ Ἡλιαδοὺδ σὺν τοῖς υἱοῖς καὶ ἀδελφοῖς, πάντες οἱ Λευῖται ὁμοθυμαδὸν ἐργοδιῶκται, ποιοῦντες εἰς τὰ ἔργα ἐν τῷ οἴκῳ τοῦ Κυρίου· καὶ ᾠκοδόμησαν οἱ οἰκοδόμοι τὸν ναὸν τοῦ Κυριου.

\vs{59}Καὶ ἔστησαν οἱ ἱερεῖς ἐστολισμενοι μετὰ μουσικῶν καὶ σαλπίγγων,
\vs{60}καὶ οἱ Λευῖται υἱοὶ Ἀσὰφ ἔχοντες τὰ κύμβαλα ὑμνοῦντες τῶ Κυρίῳ, καὶ εὐλογοῦντες κατὰ Δαυὶδ βασιλέα τοῦ Ἰσραήλ.

\vs{61}Καὶ ἐφώνησαν διʼ ὕμνων εὐλογοῦντες τῷ Κυρίῳ, ὅτι ἡ χρηστότης αὐτοῦ καὶ ἡ δόξα εἰς τοὺς αἰῶνας ἐν παντὶ Ἰσραήλ.
\vs{62}Καὶ πᾶς ὁ λαὸς ἐσάλπισαν καὶ ἐβόησαν φωνῇ μεγάλῃ, ὑμνοῦντες τῷ Κυρίῳ ἐπὶ τῇ ἐγέρσει τοῦ οἴκου Κυρίου.

\vs{63}Καὶ ἤλθοσαν ἐκ τῶν ἱερεων τῶν Λευιτῶν καὶ τῶν προκαθημένων κατὰ τὰς πατριὰς αὐτῶν, οἱ πρεσβύτεροι οἱ ἑωρακότες τὸν πρὸ τούτου οἶκον, πρὸς τὴν τουτου οἰκοδομὴν μετὰ κλαυθμοῦ καὶ κραυγῆς μεγάλης,
\vs{64}καὶ πολλοὶ διὰ σαλπιγγων καὶ χαρᾶς μεγάλῃ τῇ φωνῇ,
\vs{65}ὥστε τὸν λαὸν μὴ ἀκούειν τῶν σαλπίγγων διὰ τὸν κλαυθμὸν τοῦ λαοῦ· ὁ γὰρ ὄχλος ἦν ὁ σαλπίζων μεγάλως, ὥστε μακρόθεν ἀκούεσθαι.

\vs{66}Καὶ ἀκούσαντες οἱ ἑχθροί τῆς φυλῆς Ἰούδα καὶ Βενιαμὶν, ἤλθοσαν ἐπιγνῶναι τίς ἡ φωνὴ τῶν σαλπίγγων.
\vs{67}Καὶ ἐπέγνωσαν ὅτι οἱ ἐκ τῆς αἰχμαλωσίας οἰκοδομοῦσι τὸν ναὸν τῷ Κυρίῳ Θεῷ Ἰσραήλ.
\vs{68}Καὶ προσελθόντες τῷ Ζοροβάβελ, καὶ Ἰησοῦ, καὶ τοῖς ἡγουμένοις τῶν πατριῶν, λέγουσιν αὐτοῖς, συνοικοδομήσωμεν ὑμῖν.
\vs{69}Ὁμοίως γὰρ ὑμῖν ἀκούομεν τοῦ Κυρίου ὑμῶν, καὶ αὐτῷ ἐπιθύομεν ἀφʼ ἡμερῶν Ἀσβακαφὰς βασιλέως Ἀσσυρίων, ὃς μετήγαγεν ἡμᾶς ἐνταῦθα.

\vs{70}Καὶ εἶπεν αὐτοῖς Ζοροβάβελ καὶ Ἰησοῦς καὶ οἱ ἡγούμενοι τῶν πατριῶν τοῦ Ἰσραὴλ, οὐχ ἡμῖν καὶ ὑμῖν τοῦ οἰκοδομῆσαι τὸν οἶκον Κυρίῳ Θεῷ ἡμῶν.
\vs{71}Ἡμεῖς γὰρ μόνοι οἰκοδομήσωμεν τῷ Κυρίῳ τοῦ Ἰσραὴλ, ἀκολούθως οἷς προσέταξεν ἡμῖν Κύρος ὁ βασιλεὺς Περσῶν.
\vs{72}Τὰ δὲ ἔθνη τῆς γῆς ἐπικοιμώμενα τοῖς ἐν τῇ Ἰουδαίᾳ καὶ πολιορκοῦντες, εἶργον τοῦ οἰκοδομεῖν,
\vs{73}καὶ βουλὰς δημαγωγοῦντες, καὶ συστάσεις ποιούμενοι, ἀπεκώλυσαν τοῦ ἀποτελεσθῆναι τὴν οἰκοδομὴν πάντα τὸν χρόνον τῆς ζωῆς τοῦ βασιλέως Κύρου· καὶ εἴρχθησαν τῆς οἰκοδομῆς ἔτη δύο ἕως τῆς Δαρείου βασιλείας.

\ch{6}
Ἐν δὲ τῷ δευτέρῳ ἔτει τῆς Δαρείου βασιλείας, ἐπροφήτευσεν Ἀγγαῖος καὶ Ζαχαρίας ὁ τοῦ Ἀδδὼ οἱ προφῆται ἐπὶ τοὺς τοὺς Ἰουδαίους τοὺς ἐν τῇ Ἰουδαίᾳ καὶ Ἱερουσαλὴμ, ἐπὶ τῷ ὀνόματι Κυρίου Θεοῦ Ἰσραὴλ ἐπʼ αὐτούς.

\vs{2}Τότε στὰς Ζοροβάβελ ὁ τοῦ Σαλαθιὴλ καὶ Ἰησοῦς ὁ τοῦ Ἰωσεδὲκ, ἥρξαντο οἰκοδομεῖν τὸν οἶκον τοῦ Κυρίου, τὸν ἐν Ἱερουσαλὴμ, συνόντων τῶν προφητῶν τοῦ Κυρίου, βοηθούντων αὐτοῖς.
\vs{3}Ἐν αὐτῷ τῷ χρόνῳ παρῆν πρὸς αὐτοὺς Σισίννης ὁ ἔπαρχος Συρίας καὶ Φοινίκης, καὶ Σαθραβουζάνης καὶ οἱ συνεταῖροι,
\vs{4}καὶ εἶπαν αὐτοῖς, τίνος ὑμῖν συντάξαντος τὸν οἶκον τοῦτον οἰκοδομεῖτε, καὶ τὴν στέγην ταύτην καὶ τὰ ἄλλα πάντα ἐπιτελεῖτε; καὶ τίνες εἰσὶν οἰκοδόμοι οἱ ταῦτα ἐπιτελοῦντες;

\vs{5}Καὶ ἔσχοσαν χάριν, ἐπισκοπῆς γενομένης ἐπὶ τὴν αἰχμαλωσίαν, παρὰ τοῦ Κυρίου οἱ πρεσβύτεροι τῶν Ἰουδαίων,
\vs{6}καὶ οὐκ ἐκωλύθησαν τῆς οἰκοδομῆς, μέχρις οὗ ἀποσημανθῆναι Δαρείῳ περὶ αὐτῶν, καὶ προσφωνηθῆναι.

\vs{7}ἈΝΤΙΓΡΑΦΟΝ ἘΠΙΣΤΟΛΗΣ ἯΣ ἜΓΡΑΨΕΔΑΡΕΙΩ, ΚΑΙ ἈΠΕΣΤΕΙΛΑΝ. Σισίννης ὁ ἔπαρχος Συρίας καὶ Φοινίκης, καὶ Σαθραβουζάνης, καὶ οἱ συνεταῖροι οἱ ἐν Συρίᾳ καὶ Φοινίκῃ ἡγεμόνες, βασιλεῖ Δαρείῳ χαίρειν.
\vs{8}Πάντα γνωστὰ ἔστω τῶ κυρίῳ ἡμῶν τῷ βασιλεῖ, ὅτι παραγενόμενοι εἰς τὴν χώραν τῆς Ἰουδαίας, καὶ ἐλθόντες εἰς Ἱερουσαλὴμ τὴν πόλιν, κατελάβομεν τῆς αἰχμαλωσίας τοὺς πρεσβυτέρους τῶν Ἰουδαίων ἐν Ἱερουσαλὴμ τῇ πόλει οἰκοδομοῦντας οἶκον τῷ Κυρίῳ μέγαν,
\vs{9}καινὸν διὰ λίθων ξυστῶν πολυτελῶν, ξύλων τιθεμένων ἐν τοῖς τοίχοις,
\vs{10}καὶ τὰ ἔργα ἐκεῖνα ἐπὶ σπουδῆς γινόμενα, καὶ εὐοδούμενον τὸ ἔργον ἐν ταῖς χερσὶν αὐτῶν, καὶ ἐν πάσῃ δόξῃ καὶ ἐπιμελείᾳ συντελούμενον.

\vs{11}Τότε ἐπυνθανόμεθα τῶν πρεσβυτέρων τούτων, λέγοντες, τίνος ὑμῖν προστάξαντος οἰκοδομεῖτε τὸν οἶκον τοῦτον, καὶ τὰ ἔργα ταῦτα θεμελιοῦτε;
\vs{12}Ἐπηρωτήσαμεν οὖν αὐτοὺς, εἵνεκεν τοῦ γνωρίσαι σοι, καὶ γράψαι σοι τοὺς ἀνθρώπους τοὺς ἀφηγουμένους, καὶ τὴν ὀνοματογραφίαν ᾐτοῦμεν αὐτοὺς τῶν προκαθηγουμένων.
\vs{13}Οἱ δὲ ἀπεκρίθησαν ἡμῖν, λέγοντες, ἐσμὲν παῖδες τοῦ Κυρίου τοῦ κτίσαντος τὸν οὐρανὸν καὶ τὴν γῆν·
\vs{14}καὶ ᾠκοδόμητο οἶκος ἔμπροσθεν ἐτῶν πλειόνων διὰ βασιλέως τοῦ Ἰσραὴλ μεγάλου καὶ ἰσχυροῦ, καὶ ἐπετελέσθη.
\vs{15}Καὶ ἐπεὶ οἱ πατέρες ἡμῶν παραπικράναντες ἥμαρτον εἰς τὸν Κύριον τοῦ Ἰσραὴλ τὸν οὐράνιον, παρέδωκεν αὐτοὺς εἰς χεῖρας Ναβουχοδονόσορ βασιλέως Βαβυλῶνος βασιλέως τῶν Χαλδαίων.
\vs{16}Τόν τε οἶκον καθελόντες ἐνεπύρισαν, καὶ τὸν λαὸν ᾐχμαλώτευσαν εἰς Βαβυλῶνα.

\vs{17}Ἐν δὲ τῷ πρώτῳ ἔτει βασιλεύοντος Κύρου χώρας Βαβυλωνίας, ἔγραψεν ὁ βασιλεὺς Κύρος τὸν οἶκον τοῦτον οἰκοδομῆσαι·
\vs{18}Καὶ τὰ ἱερὰ σκεύη τὰ χρυσᾶ καὶ τὰ ἀργυρᾶ, ἃ ἐξήνεγκε Ναβουχοδονόσορ ἐκ τοῦ οἴκου τοῦ ἐν Ἱερουσαλὴμ, καὶ ἀπηρείσατο αὐτὰ ἐν τῷ αὐτοῦ ναῷ, πάλιν ἐξήνεγκεν αὐτὰ Κύρος ὁ βασιλεὺς ἐκ τοῦ ναοῦ τοῦ ἐν Βαβυλωνίᾳ, καὶ παρεδόθη Σαβανασσάρῳ Ζοροβάβελ τῷ ἐπάρχῳ,
\vs{19}καὶ ἐπετάγη αὐτῷ, καὶ ἀπήνεγκε πάντα τὰ σκεύη ταῦτα ἀποθεῖναι ἐν τῷ ναῷ τῷ ἐν Ἱερουσαλὴμ, καὶ τὸν ναὸν τοῦ Κυρίου οἰκοδομηθῆναι ἐπὶ τοῦ τόπου.
\vs{20}Τότε ὁ Σαβανάσσαρος παραγενόμενος ἐνεβάλετο τοὺς θεμελίους τοῦ οἴκου Κυρίου τοῦ ἐν Ἱερουσαλὴμ, καὶ ἀπʼ ἐκείνου μέχρι τοῦ νῦν οἰκοδομούμενος οὐκ ἔλαβε συντέλειαν.

\vs{21}Νῦν οὖν εἰ κρίνεται, βασιλεῦ, ἐπισκεπήτω ἐν τοῖς βασιλικοῖς βιβλιοφυλακίοις τοῦ Κύρου,
\vs{22}καὶ ἐὰν εὑρίσκητε; μετὰ τῆς γνώμης Κύρου τοῦ βασιλέως γενομένην τὴν οἰκοδομὴν τοῦ οἴκου Κυρίου τοῦ ἐν Ἱερουσαλὴμ, καὶ κρίνηται τῷ κυρίῳ βασιλεῖ ἡμῶν, προσφωνησάτω ἡμῖν περὶ τοῦτων.

\vs{23}Τότε ὁ βασιλεὺς Δαρεῖος προσέταξεν ἐπισκέψασθαι ἐν τοῖς βιβλιοφυλακίοις τοῖς κειμένοις ἐν Βαβυλῶνι· καὶ εὑρέθη ἐν Ἐκβατάνοις τῇ βάρει τῇ ἐν Μηδίᾳ χώρᾳ τόπος εἷς, ἐν ῷ ὑπομνημάτιστο τάδε.
\vs{24}Ἔτους πρώτου βασιλεύοντος Κύρου, βασιλεὺς Κύρος προσέταξε τὸν οἶκον τοῦ Κυρίου τὸν ἐν Ἱερουσαλὴμ οἰκοδομῆσαι, ὅπου ἐπιθύουσι διὰ πυρὸς ἐνδελεχοῦς, οὗ τὸ ὕψος πηχῶν ἑξήκοντα,
\vs{25}πλάτος πηχῶν ἑξήκοντα διὰ δόμων λιθίνων ξυστῶν τριῶν, καὶ δόμου ξυλίνου ἐγχωρίου καινοῦ ἑνὸς, καὶ τὸ δαπάνημα δοθῆναι ἐκ τοῦ οἴκου Κύρου τοῦ βασιλέως.
\vs{26}Καὶ τὰ ἱερὰ σκεύη τοῦ οἴκου Κυρίου τά τε χρυσᾶ καὶ ἀργυρᾶ, ἃ ἐξήνεγκε Ναβουχοδονόσορ ἐκ τοῦ οἴκου τοῦ ἐν Ἱερουσαλὴμ, καὶ ἀπήνεγκεν εἰς Βαβυλῶνα, ἀποκατασταθῆναι εἰς τὸν οἶκον τὸν ἐν Ἱερουσαλὴμ, οὗ ἦν κείμενα, ὅπως τεθῇ ἐκεῖ.

\vs{27}Προσέταξε δὲ ἐπιμεληθῆναι Σισίννῃ ἐπάρχῳ Συρίας καὶ Φοινίκης, καὶ Σαθραβουζάνῃ, καὶ τοῖς συνεταίροις, καὶ τοῖς ἀποτεταγμένοις ἐν Συρίᾳ καὶ Φοινίκῃ ἡγεμόσιν ἀπέχεσθαι τοῦ τόπου, ἐᾶσαι δὲ τὸν παῖδα Κυρίου Ζοροβάβελ, ἔπαρχον δὲ τῆς Ἰουδαίας, καὶ τοὺς πρεσβυτέρους τῶν Ἰουδαίων, τὸν οἶκον τοῦ Κυρίου ἐκεῖνον οἰκοδομεῖν ἐπὶ τοῦ τόπου.
\vs{28}Καὶ ἐγὼ δὲ ἐπέταξα ὁλοσχερῶς οἰκοδομῆσαι, καὶ ἀτενίσαι ἵνα συμποιῶσι τοῖς ἐκ τῆς αἰχμαλωσίας τῆς Ἰουδαίας, μέχρι τοῦ ἐπιτελεσθῆναι τὸν οἶκον τοῦ Κυρίου·
\vs{29}καὶ ἀπὸ τῆς φορολογίας κοίλης Συρίας καὶ Φοινίκης ἐπιμελῶς σύνταξιν δίδοσθαι τούτοις τοῖς ἀνθρώποις εἰς θυσίαν τῷ Κυρίῳ, Ζοροβάβελ ἐπάρεχῳ εἰς ταύρους, καὶ κριοὺς, καὶ ἄρνας,
\vs{30}ὁμοίως δὲ καὶ πυρὸν, καὶ ἅλα, καὶ οἶνον, καὶ ἔλαιον ἐνδελεχῶς κατʼ ἐνιαυτὸν, καθὼς ἂν οἱ ἱερεῖς οἱ ἐν Ἱερουσαλὴμ ὑπαγορεύσωσιν ἀναλίσκεσθαι καθʼ ἡμέραν, ἀναμφισβητήτως,
\vs{31}ὅπως προσφέρωνται σπονδαὶ τῷ Θεῷ τῷ ὑψίστῳ ὑπὲρ τοῦ βασιλέως καὶ τῶν παίδων, καὶ προσεύχωνται περὶ τῆς αὐτῶν ζωῆς·
\vs{32}καὶ προστάξαι ἵνα ὅσοι ἐὰν παραβῶσί τὶ τῶν γεγραμμένων καὶ ἀκυρώσωσι, ληφθῆναι ξύλον ἐκ τῶν ἰδίων αὐτοῦ, καὶ ἐπʼ αὐτοῦ κρεμασθῆναι, καὶ τὰ ὑπάρχοντα αὐτοῦ εἶναι βασιλικά.

\vs{33}Διὰ ταῦτα καὶ ὁ Κύριος, οὗ τὸ ὄνομα αὐτοῦ ἐπικέκληται ἐκεῖ, ἀφανίσαι πάντα βασιλέα καὶ ἔθνος, ὃς ἐκτενεῖ τὴν χεῖρα αὐτοῦ κωλῦσαι ἢ κακοποιῆσαι τὸν οἶκον Κυρίου ἐκεῖνον τὸν ἐν Ἱερουσαλήμ.
\vs{34}Ἐγὼ βασιλεὺς Δαρεῖος δεδογμάτικα ἐπιμελῶς κατὰ ταῦτα γίνεσθαι.

\ch{7}
Τότε Σισίννης ἔπαρχος κοίλης Συρίας καὶ Φοινίκης, καὶ Σαθραβουζάνης, καὶ οἱ συνεταῖροι κατακολουθήσαντες τοῖς ὑπὸ τοῦ βασιλέως Δαρείου προσταγεῖσιν,
\vs{2}ἐπεστάτουν τῶν ἱερῶν ἔργων ἐπιμελέστερον συνεργοῦντες τοῖς πρεσβυτέροις τῶν Ἰουδαίων καὶ ἱεροστάταις.
\vs{3}Καὶ εὔοδα ἐγίνετο τὰ ἱερὰ ἔργα, προφητευόντων Ἀγγαίου καὶ Ζαχαρίου τῶν προφητῶν.

\vs{4}Καὶ συνετέλεσαν ταῦτα διὰ προστάγματος Κυρίου Θεοῦ Ἰσραήλ· καὶ μετὰ τῆς γνώμης τοῦ Κύρου καὶ Δαρείου καὶ Ἀρταξέρξου βασιλέων Περσῶν,
\vs{5}συνετελέσθη ὁ οἶκος ὁ ἅγιος ἕως τρίτης καὶ εἰκάδος μηνὸς Ἄδαρ, τοῦ ἕκτου ἔτους βασιλέως Δαρείου.

\vs{6}Καὶ ἐποίησαν οἱ υἱοὶ Ἰσραὴλ, καὶ οἱ ἱερεῖς καὶ οἱ Λευῖται καὶ οἱ λοιποὶ οἱ ἐκ τῆς αἰχμαλωσίας οἱ προστεθέντες, ἀκολούθως τοῖς ἐν τῇ Μωυσέως βίβλῳ.
\vs{7}Καὶ προσήνεγκαν εἰς τὸν ἐγκαινισμὸν τοῦ ἱεροῦ τοῦ Κυρίου ταύρους ἑκατὸν, κριοὺς διακοσίους, ἄρνας τετρακοσίους,
\vs{8}χιμάρους ὑπὲρ ἁμαρτίας παντὸς τοῦ Ἰσραὴλ δώδεκα πρὸς ἀριθμὸν, ἐκ τῶν φυλάρχων τοῦ Ἰσραὴλ δώδεκα.
\vs{9}Καὶ ἔστησαν οἱ ἱερεῖς καὶ οἱ Λενῖται κατὰ φυλὰς ἐστολισμένοι ἐπὶ τῶν ἔργων Κυρίου Θεοῦ Ἰσραὴλ ἀκολούθως τῇ Μωυσέως βίβλῳ, καὶ οἱ θυρωροὶ ἐφʼ ἑκάστου πυλῶνος.

\vs{10}Καὶ ἠγάγοσαν οἱ υἱοὶ Ἰσραὴλ τῶν ἐκ τῆς αἰχμαλωσίας τὸ πάσχα ἐν τῇ τεσσαρεσκαιδεκάτῃ τοῦ πρώτου μηνὸς, ὅτε ἡγνίσθησαν οἱ ἱερεῖς καὶ οἱ Λευῖται,
\vs{11}ἅμα καὶ πάντες οἱ υἱοὶ τῆς αἰχμαλωσίας, ὅτι ἡγνίσθησαν· ὅτι οἱ Λευῖται ἅμα πάντες ἡγνίσθησαν.
\vs{12}Καὶ ἔθυσαν τὸ πάσχα πᾶσι τοῖς υἱοῖς τοῖς αἰχμαλωσίας, καὶ τοῖς ἀδελφοῖς αὐτῶν τοῖς ἱερεῦσι, καὶ ἑαυτοῖς.
\vs{13}Καὶ ἐφάγοσαν οἱ υἱοὶ Ἰσραὴλ οἱ ἐκ τῆς αἰχμαλωσίας, πάντες οἱ χωρισθέντες ἀπὸ τῶν βδελυγμάτων τῶν ἐθνῶν τῆς γῆς, ζητοῦντες τὸν Κύριον·
\vs{14}Καὶ ἠγάγοσαν τὴν ἑορτὴν τῶν ἀζύμων ἑπτὰ ἡμέρας εὐφραινόμενοι ἔναντι Κυρίου,
\vs{15}ὅτι μετέστρεψε τὴν βουλὴν τοῦ βασιλέως Ἀσσυρίων ἐπʼ αὐτοὺς, κατισχῦσαι τὰς χεῖρας αὐτῶν ἐπὶ τὰ ἔργα Κυρίου Θεοῦ Ἰσραήλ.

\ch{8}
Καὶ μεταγενέστερος τούτων ἐστὶ, βασιλεύοντος Ἀρταξέρξου τοῦ Περσῶν βασιλέως, προσέβη Ἔσδρας Ἀζαραίου, τοῦ Ζεχρίου, τοῦ Χελκίου, τοῦ Σαλήμου,
\vs{2}τοῦ Σαδδούκου, τοῦ Ἀχιτὼβ, τοῦ Ἀμαρίου, τοῦ Ὀζίου, τοῦ Βοκκὰ, τοῦ Ἀβισαῒ, τοῦ Φινεὲς, τοῦ Ἐλεάζαρ, τοῦ Ἀαρὼν, τοῦ ἱερέως τοῦ πρώτου·
\vs{3}οὗτος Ἔσδρας ἀνέβη ἐκ Βαβυλῶνος ὡς γραμματεὺς εὐφυὴς ὢν ἐν τῷ Μωυσέως νόμω τῷ ἐκδεδομένῳ ὑπὸ τοῦ Θεοῦ τοῦ Ἰσραήλ.
\vs{4}Καὶ ἔδωκεν αὐτῷ ὁ βασιλεὺς δόξαν, εὑρόντος χάριν ἐνώπιον αὐτοῦ ἐπὶ πάντα τὰ ἀξιώματα αὐτοῦ.

\vs{5}Καὶ συνανέβησαν ἐκ τῶν υἱῶν Ἰσραὴλ, καὶ τῶν ἱερέων, καὶ Λευιτῶν, καὶ ἱεροψαλτῶν, καὶ θυρωρῶν, καὶ ἱεροδούλων εἰς Ἱερουσαλὴμ,
\vs{6}ἔτους ἑβδόμου βασιλεύοντος Ἀρταξέρξου ἐν τῷ πέμπτῳ μηνί· οὗτος ἐνιαυτὸς ἕβδομος τῷ βασιλεῖ· ἐξελθόντες γὰρ ἐκ Βαβυλῶνος τῇ νουμηνίᾳ τοῦ πρώτου μηνὸς, παρεγένοντο εἰς Ἱερουσαλὴμ κατὰ τὴν δοθεῖσαν αὐτοῖς εὐοδίαν παρὰ τοῦ Κυρίου ἐπʼ αὐτῷ.
\vs{7}Ὁ γὰρ Ἔσδρας πολλὴν ἐπιστήμην περιεῖχεν εἰς τὸ μηδὲν παραλιπεῖν τῶν ἐκ τοῦ νόμου Κυρίου καὶ ἐκ τῶν ἐντολῶν, διδάξαι πάντα τὸν Ἰσραὴλ δικαιώματα καὶ κρίματα.

\vs{8}Προσπεσόντος δὲ τοῦ γραφέντος προστάγματος παρὰ Αρταξέρξου βασιλέως πρὸς Ἔσδραν τὸν ἱερέα καὶ ἀναγνώστην τοῦ νόμου Κυρίου, οὗ ἐστιν ἀντίγραφον τὸ ὑποκείμενον·

\vs{9}Βασιλεὺς Ἀρταξέρξης Ἔσδρᾳ τῷ ἱερεῖ καὶ ἀναγνώστῃ τοῦ νόμου Κυρίου χαίρειν.
\vs{10}Καὶ τὰ φιλάνθρωπα ἐγὼ κρίνας προσέταξα τοὺς βουλομένους ἐκ τοῦ ἔθνους τῶν Ἰουδαίων αἱρετίζοντας, καὶ τῶν ἱερέων καὶ τῶν Λευιτῶν, καὶ τῶνδε ἐν τῇ ἡμετέρᾳ βασιλείᾳ, συμπορεύεσθαί σοι εἰς Ἱερουσαλήμ.
\vs{11}Ὅσοι οὖν ἐνθυμοῦνται, συνεξορμάσθωσαν καθάπερ δέδοκται ἐμοί τε, καὶ τοῖς ἑπτὰ φίλοις συμβουλευταῖς,
\vs{12}ὅπως ἐπισκέψωνται τὰ κατὰ τὴν Ἰουδαίαν καὶ Ἱερουσαλὴμ ἀκολούθως ᾧ ἔχει ἐν τῷ νόμῳ Κυρίου,
\vs{13}καὶ ἀπενεγκεῖν δῶρα τῷ Κυρίῳ τοῦ Ἰσραὴλ, ἃ ηὐξάμην ἐγώ τε καὶ οἱ φίλοι, εἰς Ἱερουσαλήμ· καὶ πᾶν χρυσίον καὶ ἀργύριον ὃ ἐὰν εὑρεθῇ ἐν τῇ χώρᾳ τῆς Βαβυλωνίας τῷ Κυρίῳ εἰς Ἱερουσαλὴμ
\vs{14}σὺν τῷ δεδωρημένῳ ὑπὸ τοῦ ἔθνους εἰς τὸ ἱερὸν τοῦ Κυρίου Θεοῦ αὐτῶν τὸ ἐν Ἱερουσαλὴμ, συναχθῆναι τό, τε χρυσίον καὶ τὸ ἀργύριον εἰς ταύρους καὶ κριοὺς καὶ ἄρνας, καὶ τὰ τούτοις ἀκόλουθα,
\vs{15}ὥστε προσενεγκεῖν θυσίας τῷ Κυρίῳ ἐπὶ τὸ θυσιαστήριον τοῦ Κυρίον Θεοῦ αὐτῶν τὸ ἐν Ἱερουσαλήμ.

\vs{16}Καὶ πάντα ὅσα ἐὰν βούλῃ μετὰ τῶν ἀδελφῶν σου ποιῆσαι χρυσίῳ καὶ ἀργυρίῳ, ἐπιτέλει κατὰ τὸ θέλημα τοῦ Θεοῦ σου.
\vs{17}Καὶ τὰ ἱερὰ σκεύη τοῦ Κυρίου τὰ διδόμενά σοι εἰς τὴν χρείαν τοῦ ἱεροῦ τοῦ Θεοῦ σου,
\vs{18}δώσεις ἐκ τοῦ βασιλικοῦ γαζοφυλακίου.

\vs{19}Κᾀγὼ ἰδοὺ Ἀρταξέρξης βασιλεὺς προσέταξα τοῖς γαζοφύλαξι Συρίας καὶ Φοινίκης, ἵνα ὅσα ἐὰν ἀποστείλῃ Ἔσδρας ὁ ἱερεὺς καὶ ἀναγνώστης τοῦ νόμου τοῦ Θεοῦ τοῦ ὑψίστου, ἐπιμελῶς διδῶσιν αὐτῷ ἕως ἀργυρίου ταλάντων ἑκατὸν,
\vs{20}ὁμοίως δὲ καὶ ἕως πυροῦ κόρων ἑκατὸν, καὶ οἴνου μετρητῶν ἑκατόν·
\vs{21}καὶ ἄλλα ἐκ πλήθους πάντα κατὰ τὸν τοῦ Θεοῦ νόμον ἐπιτελεσθήτω ἐπιμελῶς τῷ Θεῷ τῷ ὑψίστῳ, ἕνεκα τοῦ μὴ γενέσθαι ὀργὴν εἰς τὴν βασιλείαν τοῦ βασιλέως καὶ τῶν υἱῶν αὐτοῦ.
\vs{22}Καὶ ὑμῖν δὲ λέγεται ὅπως πᾶσι τοῖς ἱερεῦσι, καὶ τοῖς Λευίταις, καὶ ἱεροψάλταις, καὶ θυρωροῖς, καὶ ἱεροδούλοις, καὶ πραγματικοῖς τοῦ ἱεροῦ τούτου μηδὲ μία φορολογία, μηδὲ ἄλλη ἐπιβουλὴ γίνηται, καὶ μηδένα ἔχειν ἐξουσίαν ἐπιβαλεῖν τι τούτοις.

\vs{23}Καὶ σὺ, Ἔσδρα, κατὰ τὴν σοφίαν τοῦ Θεοῦ, ἀνάδειξον κριτὰς καὶ δικαστὰς, ὅπως δικάζωσιν ἐν ὅλῃ Συρίᾳ καὶ Φοινίκῃ πάντας τοὺς ἐπισταμένους τὸν νόμον τοῦ Θεοῦ σου, καὶ τοὺς μὴ ἐπισταμένους διδάξεις.
\vs{24}Καὶ πάντες ὅσοι ἂν παραβαίνωσι τὸν νόμον τοῦ Θεοῦ σου καὶ τὸς βασιλικὸν, ἐπιμελῶς κολασθήσονται, ἐάν τε καὶ θανάτῳ, ἐάν τε καὶ τιμωρίᾳ, ἢ ἀργυρικῇ ζημίᾳ, ἢ ἀπαγωγῇ.

\vs{25}Καὶ εἶπεν Ἔσδρας ὁ γραμματεὺς, εὐλογητὸς μόνος Κύριος ὁ Θεὸς τῶν πατέρων μου, ὁ δοὺς ταῦτα εἰς τὴν καρδίαν τοῦ βασιλέως, δοξάσαι τὸν οἶκον αὐτοῦ τὸν ἐν Ἱερουσαλὴμ,
\vs{26}καὶ ἐμὲ ἐτίμησεν ἐναντίον τοὺ βασιλέως, καὶ τῶν συμβουλευόντων, καὶ πάντων τῶν φίλων, καὶ μεγιστάνων αὐτοῦ.
\vs{27}Καὶ ἐγὼ εὐθαρσὴς ἐγενόμην κατὰ τὴν ἀντίλμψιν Κυρίου τοῦ Θεοῦ μου, καὶ συνήγαγον ἄνδρας ἐκ τοῦ Ἰσραὴλ ὥστε συναναβῆναί μοι.

\vs{28}Καὶ οὗτοι οἱ προηγούμενοι κατὰ τὰς πατριὰς αὐτῶν καὶ τὰς μεριδαρχίας, οἱ ἀναβάντες μετʼ ἐμοῦ ἐκ Βαβυλῶνος ἐν τῇ βασιλείᾳ Ἀρταξέρξου τοῦ βασιλέως.
\vs{29}Ἐκ τῶν υἱῶν Φινεὲς, Γηρσών· ἐκ τῶν υἱῶν Ἰαθαμάρου, Γαμαλιήλ· ἐκ τῶν υἱῶν Λαυὶδ, Λαττοὺς ὁ Σεχενίου·
\vs{30}ἐκ τῶν υἱῶν Φόρος, Ζαχαρίας, καὶ μετʼ αὐτοῦ ἀπεγράφησαν ἄνδρες ἑκατὸν πεντήκοντα·
\vs{31}ἐκ τῶν υἱῶν Φαὰθ Μωὰβ, Ἐλιαωνίας Ζαραίου, καὶ μετʼ αὐτοῦ ἄνδρες διακόσιοι·
\vs{32}ἐκ τῶν υἱῶν Ζαθόης, Ζεχενίας Ἰεζήλου, καὶ μετʼ αὐτοῦ ἄνδρες τριακόσιοι· ἐκ τῶν υἱῶν Ἀδὶν, Ὠβὴθ Ἰωνάθου, καὶ μετʼ αὐτοῦ ἄνδρες διακόσιοι πεντήκοντα·
\vs{33}ἐκ τῶν υἱῶν Ἠλὰμ, Ἰεσίαν Γοθολίου, καὶ μετʼ αὐτοῦ ἄνδρες ἑβδομήκοντα·
\vs{34}ἐκ τῶν υἱῶν Σαφατίου, Ζαραΐας Μιχαήλου, καὶ μετʼ αὐτοῦ ἄνδρες ἑβδομήκοντα·
\vs{35}ἐκ τῶν υἱῶν Ἰωὰβ, Ἀβαδίας Ἰεζήλου, καὶ μετʼ αὐτοῦ ἄνδρες διακόσιοι δεκαδύο·
\vs{36}Ἐκ τῶν υἱῶν Βανίας, Σαλιμὼθ Ἰωσαφίου, καὶ μετʼ αὐτοῦ ἄνδρες ἑξήκοντα καὶ ἑκατόν·
\vs{37}ἐκ τῶν υἱῶν Βαβὶ, Ζαχαρίας Βηβαῒ, καὶ μετʼ αὐτοῦ ἄνδρες εἰκοσιοκτώ·
\vs{38}ἐκ τῶν υἱῶν Ἀστὰθ, Ἰωάννης Ἀκατὰν, καὶ μετʼ αὐτοῦ ἄνδρες ἑκατὸν δέκα·
\vs{39}ἐκ τῶν υἱῶν Ἀδωνικὰμ, οἱ ἔσχατοι· καὶ ταῦτα τὰ ὀνόματα αὐτῶν· Ἐλιφαλὰ τοῦ Γεουὴλ, καὶ Σαμαίας, καὶ μετʼ αὐτῶν ἄνδρες ἑβδομήκοντα·
\vs{40}ἐκ τῶν υἱῶν Βαγὼ, Οὐθὶ ὁ τοῦ Ἱσταλκούρου, καὶ μετʼ αὐτοῦ ἄνδρες ἑβδομήκοντα.

\vs{41}Καὶ συνήγαγον αὐτοὺς ἐπὶ τὸν λεγόμενον Θερὰν ποταμὸν, καὶ παρενεβάλομεν ἡμέρας τρεῖς αὐτόθι, καὶ κατέμαθον αὐτούς.
\vs{42}Καὶ ἐκ τῶν ἱερέων καὶ ἐκ τῶν Λευιτῶν οὐχ εὑρὼν ἐκεῖ,
\vs{43}ἀπέστειλα πρὸς Ἐλεάζαρον, καὶ Ἱδουῆλον, καὶ Μαιὰ, καὶ Μασμὰν, καὶ Ἀλναθὰν, καὶ Σαμαίαν, καὶ Ἰώριβον, Νάθαν,
\vs{44}Ἐννατὰν, Ζαχαρίαν, καὶ Μοσόλλαμον τοὺς ἡγουμένους καὶ ἐπιοτήμονας,
\vs{45}καὶ εἶπα αὐτοῖς ἐλθεῖν πρὸς Λοδδαῖον τὸν ἡγούμενον τὸν ἐν τῷ τόπῳ τοῦ γαζοφυλακίου,
\vs{46}ἐντειλάμενος αὐτοῖς διαλεχθῆναι Λοδδαίῳ, καὶ τοῖς ἀδελφοῖς αὐτοῦ, καὶ τοῖς ἐν τῷ τόπῳ γαζοφύλαξιν, ἀποστεῖλαι ἡμῖν τοὺν ἱερατεύσοντας ἐν τῷ οἴκῳ τοῦ Κυρίου ἡμῶν·

\vs{47}Καὶ ἤγαγον ἡμῖν κατὰ τὴν κραταιὰν χεῖρα τοῦ Κυρίου ἡμῶν ἄνδρας ἐπιστήμονας τῶν υἱῶν Μοολὶ τοῦ Λευὶ τοῦ Ἰσραὴλ, Ἀσεβηβίαν, καὶ τοὺς υἱοὺς αὐτοῦ, καὶ τοὺς ἀδελφοὺς, ὄντας δέκα καὶ ὀκτώ·
\vs{48}καὶ Ἀσεβίαν, καὶ Ἄννουον, καὶ Ὠσαίαν ἀδελφὸν ἐκ τῶν υἱῶν Χανουναίου, καὶ οἱ υἱοὶ αὐτῶν εἴκοσι ἄνδρες·
\vs{49}καὶ ἐκ τῶν ἱεροδούλων ὧν ἔδωκε Δαυὶδ, καὶ οἱ ἡγούμενοι εἰς τὴν ἐργασίαν τῶν Λευιτῶν, ἱεροδούλους διακοσίους καὶ εἴκοσι· πάντων ἐσημάνθη ἡ ὀνοματογραφία.

\vs{50}Καὶ ηὐξάμην ἐκεῖ νηστείαν τοῖς νεανίσκοις ἔναντι Κυρίου ἡμῶν, ζητῆσαι παρʼ αὐτοῦ εὐοδίαν ἡμῖν τε καὶ τοῖς συνοῦσιν ἡμῖν, τέκνοις ἡμῶν, καὶ κτήνεσιν.
\vs{51}Ἐνετράπην γὰρ αἰτῆσαι τὸν βασιλέα, πεζούς τε καὶ ἱππεῖς, και προπομπὴν ἕνεκεν ἀσφαλείας τῆς πρὸς τοὺς ἐναντιουμένους ἡμῖν.
\vs{52}Εἴπαμεν γὰρ τῷ βασιλεῖ, ὅτι ἡ ἰσχὺς τοῦ Κυρίου ἡμῶν ἔσται μετὰ τῶν ἐπιζητούντων αὐτὸν εἰς πᾶσαν ἐπανόρθωσιν.
\vs{53}Καὶ πάλιν ἐδεήθημεν τοῦ Κυρίου ἡμῶν πάντα ταῦτα, καὶ ἐτύχομεν εὐιλάτου.

\vs{54}Καὶ ἐχώρισα τῶν φυλάρχων τῶν ἱερέων ἄνδρας δεκαδύο, καὶ Ἐσερεβίαν καὶ Σαμίαν, καὶ μετʼ αὐτῶν ἐκ τῶν ἀδελφῶν αὐτῶν ἄνδρας δώδεκα.
\vs{55}Καὶ ἔστησα αὐτοῖς τὸ ἀργύριον, καὶ τὸ χρυσίον, καὶ τὰ ἱερὰ σκεύη τοῦ οἴκου τοῦ Κυρίου ἡμῶν, ἃ ἐδωρήσατο ὁ βασιλεὺς, καὶ οἱ σύμβουλοι αὐτοὺ, καὶ οἱ μεγιστᾶνες, καὶ πᾶς Ἰσραήλ.
\vs{56}Καὶ στήσας παρέδωκα αὑτοῖς ἀργυρίου τάλαντα ἑξακόσια πεντήκοντα, καὶ σκεύη ἀργυρᾶ ταλάντων ἑκατὸν, καὶ χρυσίου τάλαντα ἑκατὸν,
\vs{57}καὶ χρυσώματα εἴκοσι, καὶ σκεύη χάλκεα ἀπὸ χρηστοῦ χαλκοῦ στίλβοντα χρυσοειδῆ σκεύη δώδεκα.

\vs{58}Καὶ εἶπα αὐτοῖς καὶ ὑμεῖς ἅγιοι ἐστὲ τῷ Κυρίῳ, καὶ τὰ σκεύη τὰ ἅγια, καὶ τὸ χρυσίον, καὶ τὸ ἀργύριον, εὐχὴ τῷ Κυρίῳ, Κυρίῳ τῶν πατέρων ἡμῶν.
\vs{59}Ἀγρυπνεῖτε, καὶ φυλάσσετε ἕως τοῦ παραδοῦναι ὑμᾶς αὐτὰ τοῖς φυλάρχοις τῶν ἱερέων καὶ τῶν Λευιτῶν, καὶ τοῖς ἡγουμένοις τὼν πατριῶν τοῦ Ἰσραὴλ ἐν Ἰερουσαλὴμ, ἐν τοῖς παστοφορίοις τοῦ οἴκου τοῦ Θεοῦ ἡμῶν.
\vs{60}Καὶ οἱ παραλαβόντες οἱ ἱερεῖς καὶ οἱ Λευῖται τὸ ἀργύριον, καὶ τὸ χρυσίον, καὶ τὰ σκεύη τὰ ἐν Ἱερουσαλὴμ, εἰσήνεγκαν εἰς τὸ ἱερὸν τοῦ Κυρίου.

\vs{61}Καὶ ἀναζεύξαντες ἀπὸ τοῦ ποταμοῦ Θερὰ τῇ δωδεκάτῃ τοῦ πρώτου μηνὸς, ἕως εἰσήλθομεν εἰς Ἱερουσαλὴμ κατὰ τὴν κραταιὰν χεῖρα τοῦ Κυρίου ἡμῶν τὴν ἐφʼ ἡμῖν· καὶ ἐῤῥύσατο ἡμᾶς ἀπὸ τῆς εἰσόδου ἀπὸ παντὸς ἐχθροῦ, καὶ ἦλθομεν εἰς Ἱερουσαλήμ.
\vs{62}Καὶ γενομένης αὐτόθι ἡμέρας τρίτης, τῇ ἡμέρᾳ τῇ τετάρτῃ σταθὲν τὸ ἀργύριον καὶ τὸ χρυσίον παρεδόθη ἐν τῷ οἴκῳ Κυρίου ἡμῶν Μαρμωθὶ Οὐρία ἱερεῖ.
\vs{63}Καὶ μετʼ αὐτοῦ Ἐλεάζαρ ὁ τοῦ Φινεὲς, καὶ ἦσαν μετʼ αὐτοῦ Ἰωσαβδὸς Ἰησοῦ, καὶ Μωὲθ Σαβάννου· οἱ δὲ Λευῖται, πρὸς ἀριθμὸν καὶ ὁλκὴν ἅπαντα.
\vs{64}Καὶ ἐγράφη πᾶσα ἡ ὁλκὴ αὐτῶν αὐτῇ τῇ ὥρᾳ.

\vs{65}Οἱ δὲ παραγενόμενοι ἐκ τῆς αἰχμαλωσίας προσήνεγκαν θυσίας τῷ Θεῷ τοῦ Ἰσραὴλ Κυρίῳ, ταύρους δώδεκα ὑπὲρ παντὸς Ἰσραὴλ,
\vs{66}κριοὺς ἐνενηκονταὲξ, ἄρνας ἑβδομηκονταδύο, τράγους ὑπὲρ σωτηρίου δώδεκα, ἅπαντα θυσίαν τῷ Κυρίῳ.
\vs{67}Καὶ ἀπέδωκαν τὰ προστάγματα τοῦ βασιλέως τοῖς βασιλικοῖς οἰκονόμοις καὶ τοῖς ἐπάρχοις κοίλης Συρίας καὶ Φοινίκης, καὶ ἐδόξασαν τὸ ἔθνος, καὶ τὸ ἱερὸν τοῦ Κυρίου.

\vs{68}Καὶ τούτων τελεσθέντων, προσήλθοσάν μοι οἱ ἡγούμενοι, λέγοντες,
\vs{69}οὐκ ἐχώρισαν τὸ ἔθνος τοῦ Ἰσραὴλ καὶ οἱ ἄρχοντες καὶ οἱ ἱερεῖς καὶ οἱ Λευῖται τὰ ἀλλογενῆ ἔθνη τῆς γῆς καὶ τὰς ἀκαθαρσίας αὐτῶν, ἀπὸ τῶν ἐθνῶν τῶν Χαναναίων, καὶ Χετταίων, καὶ Φερεζαίων, καὶ Ἰεβουσαίων, καὶ Μωαβιτῶν, καὶ Αἰγυπτίων, καὶ Ἰδουμαίων.
\vs{70}Συνῴκησαν γὰρ μετὰ τῶν θυγατέρων αὐτῶν καὶ αὐτοὶ καὶ οἱ υἱοὶ αὐτῶν, καὶ ἐπεμίγη τὸ σπέρμα τὸ ἅγιον εἰς τὰ ἀλλογενῆ ἔθνη τῆς γῆς, καὶ μετεῖχον οἱ προηγούμενοι καὶ οἱ μεγιστᾶνες τῆς ἀνομίας ταύτης ἀπὸ τῆς ἀρχῆς τοῦ πράγματος.

\vs{71}Καὶ ἅμα τῷ ἀκοῦσαί με ταῦτα, διέῤῥηξα τὰ ἱμάτια καὶ τὴν ἱερὰν ἐσθῆτα, καὶ κατέτιλα τοῦ τριχώματος τῆς κεφαλῆς καὶ τοῦ πώγωνος, καὶ ἐκάθισα σύννους καὶ περίλυπος.
\vs{72}Καὶ ἐπισυνήχθησαν πρὸς μὲ ὅσοι ποτὲ ἐπεκινοῦντο ἐπὶ τῷ ῥήματι Κυρίου Θεοῦ τοῦ Ἰσραὴλ, ἐμοῦ πενθοῦντος ἐπὶ τῇ ἀνομίᾳ· καὶ ἐκαθήμην περίλυπος ἕως τῆς δειλινῆς θυσίας.

\vs{73}Καὶ ἐξεγερθεὶς ἐκ τῆς νηστείας διεῤῥηγμένα ἔχων τὰ ἱμάτια καὶ τὴν ἱερὰν ἐσθῆτα, κάμψας τὰ γόνατα, καὶ ἐκτείνας τὰς χεῖρας πρὸς τὸν Κύριον·
\vs{74}ἔλεγον, Κύριε, ᾔσχυμμαι καὶ ἐντέτραμμαι κατὰ πρόσωπόν σου.
\vs{75}Αἱ γὰρ ἁμαρτίαι ἡμῶν ἐπλεόνασαν ὑπὲρ τὰς κεφαλὰς ἡμῶν, καὶ αἱ ἄγνοιαι ἡμῶν ὑπερήνεγκαν ἕως τοῦ οὐρανοῦ,
\vs{76}ἔτι ἀπὸ τῶν χρόνων τῶν πατέρων ἡμῶν, καὶ ἐσμὲν ἐν μεγάλῃ ἁμαρτίᾳ ἕως τῆς ἡμέρας ταύτης.
\vs{77}Καὶ διὰ τὰς ἁμαρτίας ἡμῶν καὶ τῶν πατέρων ἡμῶν παρεδόθημεν σὺν τοῖς ἀδελφοῖς ἡμῶν, καὶ σὺν τοῖς βασιλεῦσιν ἡμῶν, καὶ σὺν τοῖς ἱερεῦσιν ἡμῶν, τοῖς βασιλεῦσι τῆς γῆς εἰς ῥομφαίαν καὶ αἰχμαλωσίαν καὶ προνομὴν μετὰ αἰσχύνης μέχρι τῆς σήμερον ἡμέρας.

\vs{78}Καὶ νῦν κατὰ πόσον τι ἐγενήθη ἡμῖν ἔλεος παρὰ τοῦ Κυρίου Κυρίου, καταλειφθῆναι ἡμῖν ῥίζαν καὶ ὄνομα ἐν τῷ τόπῳ ἁγιάσματός σου,
\vs{79}καὶ τοῦ ἀνακαλύψαι φωστῆρα ἡμῖν ἐν τῷ οἴκῳ Κυρίου τοῦ Θεοῦ ἡμῶν, δοῦναι ἡμῖν τροφὴν ἐν τῷ καιρῷ τῆς δουλείας ἡμῶν;
\vs{80}Καὶ ἐν τῷ δουλεύειν ἡμᾶς οὐκ ἐγκατελείφθημεν ὑπὸ τοῦ Κυρίου ἡμῶν, ἀλλὰ ἐποίησεν ἡμᾶς ἐν χάριτι ἐνώπιον τῶν βασιλέων Περσῶν,
\vs{81}δοῦναι ἡμῖν τροφὴν, καὶ δοξάσαι τὸ ἱερὸν τοῦ Κυρίου ἡμῶν, καὶ ἐγεῖραι τὴν ἔρημον Σιὼν, δοῦναι ἡμῖν στερέωμα ἐν τῇ Ἰουδαίᾳ καὶ Ἱερουσαλήμ.

\vs{82}Καὶ νῦν τί ἐροῦμεν, Κύριε, ἔχοντες ταῦτα; παρέβημεν γὰρ τὰ προστάγματά σου, ἃ ἔδωκας ἐν χειρὶ τῶν παίδων σου τῶν προφητῶν,
\vs{83}λέγων, ὅτι ἡ γῆ, εἰς ἣν εἰσέρχεσθε κληρονομῆσαι, ἔστι γῆ μεμολυσμένη μολυσμῷ τῶν ἀλλογενῶν τῆς γῆς, καὶ τῆς ἀκαθαρσίας αὐτῶν ἐνέπλησαν αὐτήν.
\vs{84}Καὶ νῦν τὰς θυγατέρας ὑμῶν μὴ συνοικήσητε τοῖς υἱοῖς αὐτῶν, καὶ τὰς θυγατέρας αὐτῶν μὴ λάβητε τοῖς υἱοῖς ὑμῶν,
\vs{85}καὶ οὐ ζητήσετε εἰρηνεῦσαι τὰ πρὸς αὐτοὺς τὸν ἅπαντα χρόνον, ἵνα ἰσχύσαντες φάγητε τὰ ἀγαθὰ τῆς γῆς, καὶ κατακληρονομήσητε τοῖς τέκνοις ὑμῶν ἕως αἰῶνος.

\vs{86}Καὶ τὰ συμβαίνοντα πάντα ἡμῖν γίνεται διὰ τὰ ἔργα ἡμῶν τὰ πονηρὰ, καὶ τὰς μεγάλας ἁμαρτίας ἡμῶν·
\vs{87}σὺ γὰρ Κύριε ὁ κουφίσας τὰς ἁμαρτίας ἡμῶν, ἔδωκας ἡμῖν τοιαύτην ῥίζαν· πάλιν ἀνεκάμψαμεν παραβῆναι τὸν νόμον σου εἰς τὸ ἐπιμιγῆναι τῇ ἀκαθαρσίᾳ τῶν ἐθνῶν τῆς γῆς.
\vs{88}Οὐχὶ ὠργίσθης ἡμῖν ἀπολέσαι ἡμᾶς, ἕως τοῦ μὴ καταλιπεῖν ῥίζαν καὶ σπέρμα καὶ ὄνομα ἡμῶν;

\vs{89}Κύριε τοῦ Ἰσραὴλ, ἀληθινὸς εἶ· κατελείφθημεν γὰρ ῥίζα ἐν τῇ σήμερον.
\vs{90}Ἰδοὺ νῦν ἐσμὲν ἐνώπιόν σου ἐν ταῖς ἀνομίαις ἡμῶν· οὐ γὰρ ἐστι στῆναι ἔτι ἔμπροσθέν σου ἐπὶ τούτοις.
\vs{91}Καὶ ὅτε προσευχόμενος Ἔσδρας ἀνθωμολογεῖτο κλαίων χαμαιπετὴς ἔμπροσθεν τοῦ ἱεροῦ, ἐπισυνήχθησαν πρὸς αὐτὸν ἀπὸ Ἱερουσαλὴμ ὄχλος πολὺς σφόδρα, ἄνδρες, καὶ γυναῖκες, καὶ νεανίαι· κλαυθμὸς γὰρ ἦν μέγας ἐν τῷ πλήθει.

\vs{92}Καὶ φωνήσας Ἰεχονίας Ἰεήλου τῶν υἱῶν Ἰσραὴλ, εἶπεν, Ἔσδρα, ἡμεῖς ἡμάρτομεν εἰς τὸν Κύριον· συνῳκίσαμεν γυναῖκας ἀλλογενεῖς ἐκ τῶν ἐθνῶν τῆς γῆς· καὶ νῦν ἐστιν ἐπάνω πᾶς Ἰσραήλ.
\vs{93}Ἐν τούτῳ γινέσθω ἡμῖν ὁρκωμοσία πρὸς τὸν Κύριον, ἐκβαλεῖν πάσας τὰς γυναῖκας ἡμῶν τὰς ἐκ τῶν ἀλλογενῶν σὺν τοῖς τέκνοις αὐτῶν,
\vs{94}ὡς ἐκρίθη σοι, καὶ ὅσοι πειθαρχοῦσι τοῦ νόμου Κυρίου.
\vs{95}Ἀναστὰς ἐπιτέλει· πρὸς σὲ γὰρ τὸ πρᾶγμα, καὶ ἡμεῖς μετὰ σοῦ ἰσχὺν ποιεῖν.
\vs{96}Καὶ ἀναστὰς Ἔσδρας ὥρκισε τοὺς φυλάρχους τῶν ἱερέων καὶ Λευιτῶν παντὸς τοῦ Ἰσραὴλ, ποιῆσαι κατὰ ταῦτα· καὶ ὤμοσαν.

\ch{9}
Καὶ ἀναστὰς Ἔσρας ἀπὸ τῆς αὐλῆς τοῦ ἱεροῦ, ἐπορεύθη εἰς τὸ παστοφόριον Ἰωνὰν τοῦ Ἐλιασίβου.
\vs{2}Καὶ αὐλισθεὶς ἐκεῖ, ἄρτου οὐκ ἐγεύσατο οὐδὲ ὕδωρ ἔπιε, πενθῶν ἐπὶ τῶν ἀνομιῶν τῶν μεγάλων τοῦ πλήθους.
\vs{3}Καὶ ἐγένετο κήρυγμα ἐν ὅλῃ τῇ Ἰουδαίᾳ καὶ Ἱερουσαλὴμ πᾶσι τοῖς ἐκ τῆς αἰχμαλωσίας, συναχθῆναι εἰς Ἱερουσαλήμ.
\vs{4}Καὶ ὅσοι ἂν μὴ ἀπαντήσωσιν ἐν δυσὶν ἢ τρισὶν ἡμέραις, κατὰ τὸ κρίμα τῶν προκαθημένων πρεσβυτέρων, ἀνιερωθήσονται τὰ κτήνη αὐτῶν, καὶ αὐτὸς ἀλλοτριωθήσεται ἀπὸ τοῦ πλήθους τῆς αἰχμαλωσίας.

\vs{5}Καὶ ἐπισυνήχθησαν πάντες οἱ ἐκ τῆς φυλῆς Ἰούδα καὶ Βενιαμεὶν ἐν τρισὶν ἡμέραις εἰς Ἱερουσαλήμ· οὗτος ὁ μὴν ἔννατος, τῇ εἰκάδι τοῦ μηνός.
\vs{6}Καὶ συνεκάθισαν πᾶν τὸ πλῆθος ἐν τῷ εὐρυχώρῳ τοῦ ἱεροῦ, τρέμοντες διὰ τὸν ἐνεστῶτα χειμῶνα.

\vs{7}Καὶ ἀναστὰς Ἔσδρας εἶπεν αὐτοῖς, ὑμεῖς ἠνομήσατε καὶ συνῳκίσατε γυναιξὶν ἀλλογενέσι, τοῦ προσθεῖναι ἁμαρτίαν τῷ Ἰσραήλ.
\vs{8}Καὶ νῦν δότε ὁμολογίαν δόξαν τῷ Κυρίῳ Θεῷ τῶν πατέρων ἡμῶν,
\vs{9}καὶ ποιήσατε τὸ θέλημα αὐτοῦ, καὶ χωρίσθητε ἀπὸ τῶν ἐθνῶν τῆς γῆς, καὶ ἀπὸ τῶν γυναικῶν τῶν ἀλλογενῶν.

\vs{10}Καὶ ἐφώνησεν ἅπαν τὸ πλῆθος, καὶ εἶπον μεγάλῃ τῇ φωνῇ, οὕτως ὡς εἴρηκας, ποιήσομεν.
\vs{11}Ἀλλὰ τὸ πλῆθος πολὺ καὶ ὥρα χειμερινὴ, καὶ οὐκ ἰσχύομεν στῆναι αἴθριοι· καὶ τὸ ἔργον οὐκ ἔστιν ἡμῖν ἡμέρας μιᾶς οὐδὲ δύο, ἐπὶ πλεῖον γὰρ ἡμάρτομεν ἐν τούτοις.
\vs{12}Στήτωσαν δὲ οἱ προηγούμενοι τοῦ πλήθους, καὶ πάντες οἱ ἐκ τῶν κατοικιῶν ἡμῶν ὅσοι ἔχουσι γυναῖκας ἀλλογενεῖς, παραγενηθήτωσαν λαβόντες χρόνον,
\vs{13}ἑκάστου δὲ τόπου τοὺς πρεσβυτέρους καὶ τοὺς κριτὰς, ἕως τοῦ λῦσαι τὴν ὀργὴν Κυρίου ἀφʼ ἡμῶν τοῦ πράγματος τούτου.

\vs{14}Ἰωνάθας Ἀζαήλου, καὶ Ἐζεκίας Θεωκανοῦ ἐπεδέξαντο κατὰ ταῦτα· καὶ Μοσόλλαμος, καὶ Λευὶς, καὶ Σαββαταῖος συνεβράβευσαν αὐτοῖς.
\vs{15}Καὶ ἐποίησαν κατὰ πάντα ταῦτα οἱ ἐκ τῆς αἰχμαλωσίας·
\vs{16}καὶ ἐπελέξατο αὐτῷ Ἔσδρας ὁ ἱερεὺς ἄνδρας ἡγουμένους τῶν πατριῶν αὐτῶν πάντας κατʼ ὄνομα, καὶ συνεκλείσθησαν τῇ νουμηνίᾳ τοῦ μηνὸς τοῦ δεκάτου, ἐτάσαι τὸ πρᾶγμα.
\vs{17}Καὶ ἤχθη ἐπὶ πέρας τὰ κατὰ τοὺς ἄνδρας τοὺς ἐπισυνέχοντας γυναῖκας ἀλλογενεῖς, ἕως τῆς νουμηνίας τοῦ πρώτου μηνός.

\vs{18}Καὶ εὑρέθησαν τῶν ἱερέων οἱ ἐπισυναχθέντες ἀλλογενεῖς γυναῖκας ἔχοντες,
\vs{19}ἐκ τῶν υἱῶν Ἰησοῦ τοῦ Ἰωσεδὲκ, καὶ τῶν ἀδελφῶν αὐτοῦ, Μαθήλας, καὶ Ἐλεάζαρος, καὶ Ἰώριβος, καὶ Ἰωαδάνος.
\vs{20}Καὶ ἐπέβαλον τὰς χεῖρας ἐκβαλεῖν τὰς γυναῖκας αὐτῶν· καὶ εἰς ἐξιλασμὸν κριοὺς ὑπὲρ τῆς ἀγνοίας αὐτῶν.

\vs{21}Καὶ ἐκ τῶν υἱῶν Ἐμμὴρ, Ἀνανίας, καὶ Ζαβδαίος, καὶ Μάνης, καὶ Σαμαῖος, καὶ Ἱερεὴλ, καὶ Ἀζαρίας·
\vs{22}καὶ ἐκ τῶν υἱῶν Φαισοὺρ, Ἐλιωναῒς, Μασσίας, Ἰσμαῆλος, καὶ Ναθαναῆλος, καὶ Ὠκόδηλος, καὶ Σαλόας.

\vs{23}Καὶ ἐκ τῶν Λευιτῶν, Ἰωζαβάδος, καὶ Σεμεῒς, καὶ Κώϊος (οὗτός ἐστι Καλιτὰς), καὶ Παθαῖος, καὶ Ἰούδας, καὶ Ἰωνάς.
\vs{24}Ἐκ τῶν ἱεροψαλτῶν, Ἐλιάσαβος, Βακχοῦρος.
\vs{25}Ἐκ τῶν θυρωρῶν, Σαλοῦμος, καὶ Τολβάνης.

\vs{26}Ἐκ τοῦ Ἰσραὴλ ἐκ τῶν υἱῶν Φόρος, Ἱερμὰς, καὶ Ἰεζίας, καὶ Μελχίας, καὶ Μαῆλος, καὶ Ἐλεάζαρος, καὶ Ἀσεβίας, καὶ Βαναίας.
\vs{27}Ἐκ τῶν υἱῶν Ἠλὰ, Ματθανίας, Ζαχαρίας, καὶ Ἰεζριῆλος, καὶ Ἰωαβδίος, καὶ Ἱερεμὼθ, καὶ Ἀϊδίας.
\vs{28}Καὶ ἐκ τῶν υἱῶν Ζαμὼθ, Ἐλιαδὰς, Ἐλιάσιμος, Ὀθονίας, Ἰαριμὼθ, καὶ Σάβαθος, καὶ Ζεραλίας.
\vs{29}Καὶ ἐκ τῶν υἱῶν Βηβαῒ, Ἰωάννης, καὶ Ἀνανίας, καὶ Ἰωζάβδος, καὶ Ἀμαθίας.
\vs{30}Ἐκ τῶν υἱῶν Μανὶ, Ὠλαμὸς, Μαμοῦχος, Ἰεδαῖος, Ἰασούβος, καὶ Ἰασαῆλος, καὶ Ἱερεμώθ.
\vs{31}Καὶ ἐξ υἱῶν Ἀδδὶ, Νάαθος, καὶ Μοοσίας, Λακκοῦνος, καὶ Ναΐδος, Ματθανίας, καὶ Σεσθὴλ, καὶ Βαλνούος, καὶ Μανασσίας.
\vs{32}Καὶ ἐκ τῶν υἱῶν Ἀνὰν, Ἐλιωνὰς, καὶ Ἀσαΐας, καὶ Μελχίας, καὶ Σαββαῖος, καὶ Σίμων Χοσαμαίος.
\vs{33}Καὶ ἐκ τῶν υἱῶν Ἀσὸμ, Ἀλταναῖος, καὶ Ματταθίας, καὶ Σαβανναῖος, καὶ Ἐλιφαλὰτ, καὶ Μανασσῆς, καὶ Σεμεΐ.
\vs{34}Καὶ ἐκ τῶν υἱῶν Βαανὶ, Ἱερεμίας, Μομδίος, Ἰσμαῆρος, Ἰουὴλ, Μαβδαῒ, καὶ Πεδίας, καὶ Ἄνως, Ῥαβασίων, καὶ Ἐνάσιβος, καὶ Μαμνιτάναιμος, Ἐλίασις, Βαννοὺς, Ἐλιαλὶ, Σομεῒς, Σελεμίας, Ναθανίας· καὶ ἐκ τῶν υἱῶν Ἐζωρὰ, Σεσὶς, Ἐσρὶλ, Ἀζαῆλος, Σαματὸς, Ζαμβρὶ, Ἰώσηφος.
\vs{35}Καὶ ἐκ τῶν υἱῶν Ἐθμὰ, Μαζιτίας, Ζαβαδαίας, Ἠδαῒς, Ἰουὴλ, Βαναίας.

\vs{36}Πάντες οὗτοι συνῴκισαν γυναῖκας ἀλλογενεῖς, καὶ ἀπέλυσαν αὐτὰς σὺν τέκνοις.

\vs{37}Καὶ κατῴκησαν οἱ ἱερεῖς, καὶ οἱ Λευῖται, καὶ οἱ ἐκ τοῦ Ἰσραὴλ ἐν Ἱερουσαλὴμ καὶ ἐν τῇ χώρᾳ τῇ νουμηνίᾳ τοῦ μηνὸς τοῦ ἑβδόμου, καὶ οἱ υἱοὶ Ἰσραὴλ ἐν ταῖς κατοικίαις αὐτων.

\vs{38}Καὶ συνήχθη πᾶν τὸ πλῆθος ὁμοθυμαδὸν ἐπὶ τὸ εὐρύχωρον τοῦ πρὸς ἀνατολὰς τοῦ ἱεροῦ πυλῶνος,
\vs{39}καὶ εἶπεν Ἔσδρᾳ τῷ ἱερεῖ καὶ ἀναγνώστῃ, κόμισαι τὸν νόμον Μωυσῆ, τὸν παραδοθέντα ὑπὸ Κυρίου Θεοῦ Ἰσραήλ.
\vs{40}Καὶ ἐκόμισεν Ἔσδρας ὁ ἀρχιερεὺς τὸν νόμον παντὶ τῷ πλήθει ἀπὸ ἀνθρώπου ἕως γυναικὸς, καὶ πᾶσι τοῖς ἱερεῦσιν, ἀκοῦσαι τοῦ νόμου νουμηνίᾳ τοῦ ἑβδόμου μηνός.
\vs{41}Καὶ ἀνεγίνωσκεν ἐν τῷ πρὸ τοῦ ἱεροῦ πυλῶνος εὐρυχώρῳ, ἐξ ὄρθρου ἕως μέσης ἡμέρας, ἐνώπιον ἀνδρῶν τε καὶ γυναικῶν· καὶ ἐπέδωκαν πᾶν τὸ πλῆθος τὸν νοῦν εἰς τὸν νόμον.

\vs{42}Καὶ ἔστη Ἔσδρας ὁ ἱερεὺς καὶ ἀναγνώστης τοῦ νόμου ἐπὶ τοῦ ξυλίνου βήματος τοῦ κατασκευασθέντος.
\vs{43}Καὶ ἔστησαν παρʼ αὐτῷ Ματταθίας, Σαμμοὺς, Ἀνανίας, Ἀζαρίας, Οὐρίας, Ἐζεκίας, Βαάλσαμος, ἐκ δεξιῶν·
\vs{44}καὶ ἐξ εὐωνύμων Φαλδαῖος, καὶ Μισαὴλ, Μελχίας, Λωθάσουβος, Ναβαρίας, Ζαχαρίας.

\vs{45}Καὶ ἀναλαβὼν Ἔσδρας τὸ βιβλίον ἐνώπιον τοῦ πλήθους, προεκάθητο ἐπιδόξως ἐνώπιον πάντων.
\vs{46}Καὶ ἐν τῷ λῦσαι τὸν νόμον, πάντες ὀρθοὶ ἔστησαν· καὶ εὐλόγησεν Ἔσδρας τῷ Κυρίῳ Θεῷ ὑψίστῳ Θεῷ σαβαὼθ παντοκράτορι.
\vs{47}Καὶ ἐπεφώνησε πᾶν τὸ πλῆθος, ἀμήν· καὶ ἄραντες ἄνω τὰς χεῖρας, προσπεσόντες ἐπὶ τὴν γῆν, προσεκύνησαν τῷ Κυρίῳ.

\vs{48}Ἰησοῦς, καὶ Ἀννιοὺθ, καὶ Σαραβίας, καὶ Ἰαδινὸς, καὶ Ἰάκουβος, Σαβαταῖος, Αὐταίας, Μαιάννας, καὶ Καλίτας, Ἀζαρίας, καὶ Ἰώζαβδος, καὶ Ἀνανίας, Φαλίας, οἱ Λευῖται, ἐδίδασκον τὸν νόμον τοῦ Κυρίου, καὶ πρὸς τὸ πλῆθος ἀνεγίνωσκον τὸν νόμον τοῦ Κυρίου, ἐμφυσιοῦντες ἅμα τὴν ἀνάγνωσιν.

\vs{49}Καὶ εἶπεν Ἀτθαράτης Ἔσδρᾳ τῷ ἀρχιερεῖ καὶ ἀναγνώστῃ, καὶ τοῖς Λευίταις τοῖς διδάσκουσι τὸ πλῆθος ἐπὶ πάντας,
\vs{50}ἡ ἡμέρα αὕτη ἐστὶν ἁγία τῷ Κυρίῳ· καὶ πάντες ἔκλαιον ἐν τῷ ἀκοῦσαι τοῦ νόμου·
\vs{51}βαδίσαντες οὖν φάγετε λιπάσματα, καὶ πίετε γλυκάσματα, καὶ ἀποστείλατε ἀποστολὰς τοῖς μὴ ἔχουσιν·
\vs{52}ἁγία γὰρ ἡ ἡμέρα τῷ Κυρίῳ· καὶ μὴ λυπεῖσθε, ὁ γὰρ Κύριος δοξάσει ὑμᾶς.

\vs{53}Καὶ οἱ Λευῖται ἐκέλευον παντὶ τῷ δήμῳ, λέγοντες ἡ ἡμέρα αὕτη ἁγία, μὴ λυπεῖσθε.
\vs{54}Καὶ ᾤχοντο πάντες φαγεῖν καὶ πιεῖν καὶ εὐφραίνεσθαι, καὶ δοῦναι ἀποστολὰς τοῖς μὴ ἔχουσι, καὶ εὐφρανθῆναι μεγάλως,
\vs{55}ὅτι γὰρ ἐνεφυσιώθησαν ἐν τοῖς ῥήμασιν οἷς ἐδιδάχθησαν, καὶ ἐπισυνήχθησαν.


\def\book{ΤΩΒΙΤ}
\biblebook{ΤΩΒΙΤ}


\lettrine[lines=2, loversize=0.2, nindent=0em, findent=.25em]{\textcolor{bookheadingcolor}{Β}}{ΙΒΛΟΣ} λόγων Τωβὶτ, τοῦ Τωβιὴλ, τοῦ Ανανιὴλ, τοῦ Ἀδουὴλ, τοῦ Γαβαήλ, ἐκ τοῦ σπέρματος Ἀσιήλ, ἐκ τῆς φυλῆς Νεφθαλί,
\vs{2}ὃς ᾐχμαλωτεύθη ἐν ἡμέραις Ἐνεμεσσάρου τοῦ βασιλέως Ἀσσυρίων ἐκ Θίσβης, ἥ ἐστιν ἐκ δεξιῶν κυδίως τῆς Νεφθαλὶ ἐν τῇ Γαλιλαίᾳ ὑπεράνω Ἀσήρ.

\vs{3}Ἐγὼ Τωβὶτ ὁδοῖς ἀληθείας ἐπορευόμην καὶ δικαιοσύνης πάσας τὰς ἡμέρας τῆς ζωῆς μου· καὶ ἐλεημοσύνας πολλὰς ἐποίησα τοῖς ἀδελφοῖς μον, καὶ τῷ ἔθνει, τοῖς προπορευθεῖσι μετʼ ἐμοῦ εἰς χώραν Ἀσσυρίων εἰς Νινευή.
\vs{4}Καὶ ὅτι ἤμην ἐν τῇ χώρᾳ μου ἐν τῇ γῇ Ἰσραὴλ, νεωτέρου μου ὄντος, πᾶσα φυλὴ τοῦ Νεφθαλὶ τοῦ πατρός μου ἀπέστη ἀπὸ τοῦ οἴκου Ἰεροσολύμων, τῆς ἐκλεγείσης ἀπὸ πασῶν τῶν φυλῶν Ἰσραὴλ, εἰς τὸ θυσιάζειν πάσας τὰς φυλάς· καὶ ἡγιάσθη ὁ ναὸς τῆς κατασκηνώσεως τοῦ ὑψίστου, καὶ ᾠκοδομήθη εἰς πάσας τὰς γενεὰς τοῦ αἰῶνος.

\vs{5}Καὶ πᾶσαι αἱ φυλαὶ αἱ συναποστᾶσαι ἔθυον τῇ Βάαλ τῇ δαμάλει, καὶ ὁ οἶκος Νεφθαλὶ τοῦ πατρός μου.
\vs{6}Κᾀγὼ μόνος ἐπορευόμην πλεονάκις εἰς Ἱεροσόλυμα ἐν ταῖς ἑορταῖς, καθὼς γέγραπται παντὶ τῷ Ἰσραὴλ, ἐν προστάγματι αἰωνίῳ, τὰς ἀπαρχὰς, καὶ τὰς δεκάτας τῶν γεννημάτων, καὶ τὰς πρωτοκουρίας ἔχων, καὶ ἐδίδουν αὐτὰς τοῖς ἱερεῦσι τοῖς υἱοῖς ʼΑαρὼν πρὸς τὸ θυσιαστήριον πάντων τῶν γεννημάτων.
\vs{7}Τὴν δεκάτην ἐδίδουν τοῖς υἱοῖς Λευὶ τοῖς θεραπεύουσιν εἰς Ἱερουσαλήμ, καὶ τὴν δευτέραν δεκάτην ἀπεπρατιζόμην, καὶ ἐπορευόμην καὶ ἐδαπάνων· αὐτὰ ἐν Ἱεροσολύμοις καθʼ ἕκαστον ἐνιαυτὸν,
\vs{8}καὶ τὴν τρίτην ἐδίδουν οἷς καθήκει, καθὼς ἐνετείλατο Δεββωρὰ ἡ μήτηρ τοῦ πατρός μου, διότι ὀρφανὸς κατελείφθην ὑπὸ τοῦ πατρός μου.

\vs{9}Καὶ ὅτε ἐγενόμην ἀνήρ, ἔλαβον Ἄνναν γυναῖκα ἐκ τοῦ σπέρματος τῆς πατριᾶς ἡμῶν· καὶ ἐγέννησα ἐξ αὐτῆς Τωβίαν.
\vs{10}Καὶ ὅτε ᾐχμαλωτίσθημεν εἰς Νινευῆ, πάντες οἱ ἀδελφοί μου, καὶ οἱ ἐκ τοῦ γένους μου ἤσθιον ἐκ τῶν ἄρτων τῶν ἐθνῶν·
\vs{11}ἐγὼ δὲ συνετήρησα τὴν ψυχήν μου μὴ φαγεῖν,
\vs{12}καθότι ἐμεμνήμην τοῦ Θεοῦ ἐν ὅλῃ τῇ ψυχῇ μου.
\vs{13}Καὶ ἔδωκεν ὁ ὕψιστος χάριν καὶ μορφὴν ἐνώπιον Ἐνεμεσσάρου, καὶ ἤμην αὐτοῦ ἀγοραστής.

\vs{14}Καὶ ἐπορευόμην εἰς τὴν Μηδίαν, καὶ παρεθέμην Γαβαήλῳ τῷ ἀδελφῷ Γαβρία ἐν Ῥάγοις τῆς Μηδίας, ἀργυρίου τάλαντα δέκα.
\vs{15}Καὶ ὅτε ἀπέθανεν Ἐνεμεσσὰρ, ἐβασίλευσε Σενναχηρὶμ ὁ υἱὸς αὐτοῦ ἀντʼ αὐτοῦ, καὶ αἱ ὁδοὶ αὐτοῦ ἠκαταστάθησαν, καὶ οὐκ ἔτι ἠδυνάσθην πορευθῆναι εἰς τὴν Μηδίαν.

\vs{16}Καὶ ἐν ταῖς ἡμέραις Ἐνεμεσσάρου ἐλεημοσύνας πολλὰς ἐποίουν τοῖς ἀδελφοῖς μου·
\vs{17}τοὺς ἄρτους μου ἐδίδουν τοῖς πεινῶσι, καὶ ἱμάτια τοῖς γυμνοῖς· καὶ εἴ τινα ἐκ τοῦ γένους μου ἐθεώρουν τεθνηκότα καὶ ἐῤῥιμμένον ὀπίσω τοῦ τείχους Νινευή, ἔθαπτον αὐτόν.
\vs{18}Καὶ εἴ τινα ἀπέκτεννε Σενναχηρὶμ ὁ βασιλεὺς, ὅτε ἦλθε φεύγων ἐκ τῆς Ἰιυδαίας, ἔθαψα αὐτοὺς κλέπτων· πολλοὺς γὰρ ἀπέκτεινεν ἐν τῷ θυμῷ αὐτοῦ· καὶ ἐζητήθη ὑπὸ τοῦ βασιλέως τὰ σώματα, καὶ οὐχ εὑρέθη.

\vs{19}Πορευθεὶς δὲ εἷς τῶν ἐν Νινευῇ, ὑπέδειξε τῷ βασιλεῖ περὶ ἐμοῦ ὅτι θάπτω αὐτοὺς, και ἐκρύβην· ἐπιγνοὺς δὲ ὅτι ζητοῦμαι ἀποθανεῖν, φοβηθεὶς ἀνεχώρησα.
\vs{20}Καὶ διηρπάγη πάντα τὰ ὑπάρχοντά μου, καὶ οὐ κατελείφθη μοι οὐδὲν, πλὴν Ἄννας τῆς γυναικός μου, καὶ Τωβὶου τοῦ υἱοῦ μου.
\vs{21}Καὶ οὐ διῆλθον ἡμέρας πεντήκοντα, ἕως οὗ ἀπέκτειναν αὐτὸν οἱ δύο υἱοὶ αὐτοῦ· καὶ ἔφυγον εἰς τὰ ὄρη Ἀραράθ· καὶ ἐβασίλευσε Σαχερδονὸς υἱὸς αὐτοῦ ἀντʼ αὐτοῦ, καὶ ἔταξεν Ἀχιάχαρον τὸν Ἀναὴλ υἱὸν τοῦ ἀδελφοῦ μου ἐπὶ πᾶσαν τὴν ἐκλογιστίαν τῆς βασιλείας αὐτοῦ, καὶ ἐπὶ πᾶσαν τὴν διοίκησιν.

\vs{22}Καὶ ἠξίωσεν Ἀχιάχαρος περὶ ἐμοῦ, καὶ ἦλθον εἰς Νινευή. Ἀχιάχαρος δὲ ἦν ὁ οἰνοχόος, καὶ ἐπὶ τοῦ δακτυλίου, καὶ διοικητὴς, καὶ ἐκλογιστὴς, καὶ κατέστησεν αὐτὸν ὁ Σαχερδονὸς· ἐκ δευτέρας, ἦν δὲ ἐξάδελφός μου.

\ch{2}
Ὅτε δὲ κατῆλθον εἰς τὸν οἶκόν μου, καὶ ἀπεδόθη μοι Ἄννα ἡ γυνή μου, καὶ Τωβίας ὁ υἱός μου, ἐν τῇ πεντηκοστῇ ἑορτῇ, ἥν ἐστιν ἁγία ἑπτὰ ἑβδομάδων, ἐγενήθη ἄριστον καλόν μοι, καὶ ἀνέπεσα τοῦ φαγεῖν.
\vs{2}Καὶ ἐθεασάμην ὄψα πολλά, καὶ εἶπα τῷ υἱῷ μου, βάδισον καὶ ἄγαγε ὃν ἂν εὕρῃς τῶν ἀδελφῶν ἡμῶν ἐνδεῆ, ὃς μέμνηται τοῦ Κυρίου, καὶ ἰδοὺ μένω σε.

\vs{3}Καὶ ἐλθὼν εἶπε, πάτερ, εἷς ἐκ τοῦ γένους ἡμῶν ἐστραγγαλωμένος ἔῤῥιπται ἐν τῇ ἀγορᾷ.
\vs{4}Κᾀγὼ πρινὴ γεύσασθαί με, ἀναπηδήσας ἀνειλόμην αὐτὸν εἴς τι οἴκημα ἕως οὗ ἔδυ ὁ ἥλιος.
\vs{5}Καὶ ἐπιστρέψας ἐλουσάμην, καὶ ἤσθιον τὸν ἄρτον μου ἐν λύπῃ.
\vs{6}Καὶ ἐμνήσθην τῆς προφητείας Ἀμὼς, καθὼς εἶπε, στραφήσονται αἱ ἑορταὶ ὑμῶν εἰς πένθος, καὶ πᾶσαι αἱ εὐφροσύναι ὑμῶν εἰς θρῆνον.
\vs{7}Καὶ ἔκλαυσα· καὶ ὅτε ἔδυ ὁ ἥλιος, ᾠχόμην, καὶ ὀρύξας ἔθαψα αὐτόν.
\vs{8}Καὶ οἱ πλησίον ἐπεγέλων, λέγοντες, οὐκ ἔτι φοβεῖται φονευθῆναι περὶ τοῦ πράγματος τούτου, καὶ ἀπέδρα, καὶ ἰδοὺ πάλιν θάπτει τοὺς νεκρούς.

\vs{9}Καὶ ἐν αὐτῇ τῇ νυκτὶ ἀνέλυσα θάψας, καὶ ἐκοιμήθην μεμιαμμένος παρὰ τὸν τοῖχον τῆς αὐλῆς, καὶ τὸ πρόσωπόν μου ἀκάλυπτον ἦν.
\vs{10}Καὶ οὐκ ᾔδειν ὅτι στρουθία ἐν τῷ τοίχῳ ἐστί· καὶ τῶν ὀφθαλμῶν μου ἀνεῳγότων, ἀφώδευσαν τὰ στρουθία θερμὸν εἰς τοὺς ὀφθαλμούς μου, καὶ ἐγενήθη λευκώματα ἐν τοῖς ὀφθαλμοῖς μου, καὶ ἐπορεύθην πρὸς ἰατροὺς, καὶ οὐκ ὠφέλησάν με· Ἀχιάχαρος δὲ ἔτρεφέ με ἕως οὗ ἐπορεύθην εἰς τὴν Ἐλυμαΐδα.

\vs{11}Καὶ ἡ γυνή μου Ἄννα ἠριθεύετο ἐν τοῖς γυναικείοις, καὶ ἀπέστελλε τοῖς κυρίοις.
\vs{12}Καὶ ἀπέδωκαν αὐτῇ καὶ αὐτοὶ τὸν μισθὸν, προσδόντες καὶ ἔριφον.
\vs{13}Ὅτε δὲ ἦλθε πρὸς μέ, ἤρξατο κράζειν· καὶ εἶπα αὐτῇ, πόθεν τὸ ἐρίφιον; μὴ κλεψιμαῖόν ἐστίν; ἀπόδος αὐτὸ τοῖς κυρίοις· οὐ γὰρ θεμιτόν ἐστι φαγεῖν κλεψιμαῖον.
\vs{14}Ἡ δὲ εἶπε, δῶρον δέδοταί μοι ἐπὶ τῷ μισθῷ· καὶ οὐκ ἐπίστευον αὐτῇ· καὶ ἔλεγον ἀποδιδόναι αὐτὸ τοῖς κυρίοις, καὶ ἠρυθρίων πρὸς αὐτήν· ἡ δὲ ἀποκριθεῖσα εἶπέ μοι, ποῦ εἰσιν αἱ ἐλεημοσύναι σου, καὶ αἱ δικαιοσύναι σου; ἰδοὺ γνωστὰ πάντα μετὰ σοῦ.

\ch{3}
Καὶ λυπηθεὶς ἔκλαυσα, καὶ προσευξάμην μετʼ ὀδύνης, λέγων, Δίκαιος εἶ Κύριε,
\vs{2}καὶ πάντα τὰ ἔργα σου, καὶ πᾶσαι αἱ ὁδοί σου ἐλεημοσύναι καὶ ἀλήθεια, καὶ κρίσιν ἀληθινὴν καὶ δικαίαν σὺ κρίνεις εἰς τὸν αἰῶνα.
\vs{3}Μνήσθητί μου, καὶ ἐπίβλεψον ἐπʼ ἐμέ· μή με ἐκδικῇς ταῖς ἁμαρτίαις μου καὶ τοῖς ἀγνοήμασί μου, καί τῶν πατέρων μου, ἃ ἥμαρτον ἐνώπιόν σου.
\vs{4}Παρήκουσαν γὰρ τῶν ἐντολῶν σου, καὶ ἔδωκας ἡμᾶς εἰς διαρπαγὴν καὶ αἰχμαλωσίαν καὶ θάνατον καὶ παραβολὴν ὀνειδισμοῦ πᾶσι τοῖς ἔθνεσιν ἐν οἷς ἐσκορπίσμεθα.

\vs{5}Καὶ νῦν πολλαὶ αἱ κρίσεις σου εἰσιὶ καὶ ἀληθιναὶ, ἐξ ἐμοῦ ποιῆσαι περὶ τῶν ἁμαρτιῶν μου καὶ τῶν πατέρων μου, ὃτι οὐκ ἐποιήσαμεν τὰς ἐντολάς σου, οὐ γὰρ ἐπορεύθημεν ἐν ἀληθείᾳ ἐνώπιόν σου.
\vs{6}Καὶ νῦν κατὰ τὸ ἀρεστὸν ἐνώπιόν σου ποίησον μετʼ ἐμοῦ· ἐπίταξον ἀναλαβεῖν τὸ πνεῦμά μου, ὅπως ἀπολυθῶ, καὶ γένωμαι γῆ, διότι λυσιτελεῖ μοι ἀποθανεῖν, ἢ ζῇν· ὅτι ὀνειδισμοὺς ψευδεῖς ἤκουσα, καὶ λύπη ἐστὶ πολλὴ ἐν ἐμοί· ἐπίταξον ἀπολυθῆναί με τῆς ἀνάγκης ἤδη εἰς τὸν αἰώνιον τόπον, μὴ ἀποστρέψῃς τὸ πρόσωπόν σου ἀπʼ ἐμοῦ.

\vs{7}Ἐν τῇ αὐτῇ ἡμέρᾳ συνέβη τῇ θυγατρὶ Ῥαγουὴλ Σάῤῥᾳ ἐν Ἐκβατάνοις τῆς Μηδίας, καὶ ταύτην ὀνειδισθῆναι ὑπὸ παιδισκῶν πατρὸς αὐτῆς,
\vs{8}ὅτι ἦν δεδομένη ἀνδράσιν ἑπτὰ, καὶ Ἀσμοδαῖος τὸ πονηρὸν δαιμόνιον ἀπέκτεινεν αὐτοὺς, πρινὴ γενέσθαι αὐτοὺς μετʼ αὐτῆς ὡς ἐν γυναιξί· καὶ εἶπαν αὐτῇ, οὐ συνιεῖς ἀποπνίγουσά σου τοὺς ἄνδρας; ἤδη ἑπτὰ ἔσχες, καὶ ἑνὸς αὐτῶν οὐκ ὠνομάσθης.
\vs{9}Τί ἡμᾶς μαστιγοῖς; εἰ ἀπέθαναν, βάδιζε μετʼ αὐτῶν, μὴ ἴδοιμέν σου υἱὸν ἢ θυγατέρα εἰς τὸν αἰῶνα.
\vs{10}Ταῦτα ἀκούσασα ἐλυπήθη σφόδρα, ὥστε ἀπάγξασθαι· καὶ εἶπε, μία μέν εἰμι τῷ πατρί μου· ἐὰν ποιήσω τοῦτο, ὄνειδος αὐτῷ ἔσται, καὶ τὸ γῆρας αὐτοῦ κατάξω μετʼ ὀδύνης εἰς ᾅδου.

\vs{11}Καὶ ἐδεήθη πρὸς τῇ θυρίδι, καὶ εἶπεν, εὐλογητὸς εἶ Κύριε ὁ Θεός μου, καὶ εὐλογητὸν τὸ ὄνομά σου τὸ ἅγιον καὶ ἔντιμον εἰς τοὺς αἰῶνας· εὐλογήσαισάν σε πάντα τὰ ἔργα σου εἰς τὸν αἰῶνα.
\vs{12}Καὶ νῦν, Κύριε, τοὺς ὀφθαλμούς μου καὶ τὸ πρόσωπόν μου εἰς σὲ δέδωκα.
\vs{13}Εἶπον, ἀπολῦσαί με ἀπὸ τῆς γῆς, καὶ μὴ ἀκοῦσαί με μηκέτι ὀνειδισμόν.
\vs{14}Σὺ γινώσκεις, Κύριε, ὅτι καθαρά εἰμι ἀπὸ πάσης ἁμαρτίας ἀνδρός,
\vs{15}καὶ οὐκ ἐμόλυνα τὸ ὄνομά μου οὐδὲ τὸ ὄνομα τοῦ πατρός μου ἐν τῇ γῇ τῆς αἰχμαλωσίας μου· μονογενής εἰμι τῷ πατρί μου, καὶ οὐχ ὑπάρχει αὐτῷ παιδίον ὃ κληρονομήσει αὐτόν, οὐδὲ ἀδελφὸς ἐγγὺς, οὐδὲ ὑπάρχων αὐτῷ υἱὸς, ἵνα συντηρήσω ἐμαυτὴν αὐτῷ γυναῖκα, ἤδη ἀπώλοντό μοι ἑπτά· ἱνατί μοι ζῇν; καὶ εἰ μὴ δοκεῖ σοι ἀποκτεῖναί με, ἐπίταξον ἐπιβλέψαι ἐπʼ ἐμὲ, καὶ μηκέτι ἐλεῆσαί με, καὶ ἀκοῦσαί με ὀνειδισμόν.

\vs{16}Καὶ εἰσηκούσθη προσευχὴ ἀμφοτέρων ἐνώπιον τῆς δόξης τοῦ μεγάλου,
\vs{17}Ῥαφαήλ καὶ ἀπεστάλη ἰάσασθαι τοὺς δύο, τοῦ Τωβὶτ λεπίσαι τὰ λευκώματα, καὶ Σάῤῥαν τὴν τοῦ Ῥαγουὴλ δοῦναι Τωβίᾳ τῷ υἱῷ Τωβὶτ γυναῖκα, καὶ δῆσαι Ἀσμοδαῖον τὸ πονηρὸν δαιμόνιον, διότι Τωβίᾳ ἐπιβάλλει κληρονομῆσαι αὐτήν. Ἐν αὐτῷ τῷ καιρῷ ἐπιστρέψας Τωβὶτ εἰσῆλθεν εἰς τὸν οἶκον αὐτοῦ, καὶ Σάῤῥα ἡ τοῦ Ῥαγουὴλ κατέβη ἐκ τοῦ ὑπερῴου αὐτῆς.

\ch{4}
Ἐν τῇ ἡμέρᾳ ἐκείνῃ ἐμνήσθη Τωβὶτ περὶ τοῦ ἀργυρίου, οὗ παρέθετο Γαβαὴλ ἐν Ῥάγοις τῆς Μηδίας.
\vs{2}Καὶ εἶπεν ἐν ἑαυτῷ, ἐγὼ ᾐτησάμην θάνατον, τί οὐ καλῶ Τωβίαν τὸν υἱόν μου, ἵνα αὐτῷ ὑποδείξω, πρὶν ἀποθανεῖν με;

\vs{3}Καὶ καλέσας αὐτὸν, εἶπε, παιδίον, ἐὰν ἀποθάνω, θάψον με, καὶ μὴ ὑπερίδῃς τὴν μητέρα σου· τίμα αὐτὴν πάσας τὰς ἡμέρας τῆς ζωῆς σου, καὶ ποίει τὸ ἀρεστὸν αὐτῇ, καὶ μὴ λυπήσῃς αὐτήν.
\vs{4}Μνήσθητι, παιδίον, ὅτι πολλοὺς κινδύνους ἑώρακεν ἐπὶ σοὶ ἐν τῇ κοιλίᾳ· ὅταν ἀποθάνῃ, θάψον αὐτὴν παρʼ ἐμοὶ ἐν ἑνὶ τάφῳ.

\vs{5}Πάσας τὰς ἡμέρας, παιδίον, Κυρίου τοῦ Θεοῦ ἡμῶν μνημόνευε, καὶ μὴ θελήσῃς ἁμαρτάνειν καὶ παραβῆναι τὰς ἐντολὰς αὐτοῦ· δικαιοσύνην ποίει πάσας τὰς ἡμέρας τῆς ζωῆς σου, καὶ μὴ πορευθῇς ταῖς ὁδοῖς τῆς ἀδικίας.
\vs{6}Διότι ποιουντός σου τὴν ἀλήθειαν, εὐοδίαι ἔσονται ἐν τοῖς ἔργοις σου, καὶ πᾶσι τοῖς ποιοῦσι τὴν δικαιοσύνην.
\vs{7}Ἐκ τῶν ὑπαρχόντων σοι ποίει ἐλεημοσύνην, καὶ μὴ φθονεσάτω σου ὁ ὀφθαλμὸς ἐν τῷ ποιεῖν σε ἐλεημοσύνην· μὴ ἀποστρέψῃς τὸ πρόσωπόν σου ἀπὸ παντὸς πτωχοῦ, καὶ ἀπὸ σοῦ οὐ μὴ ἀποστραφῇ τὸ πρόσωπον τοῦ Θεοῦ.
\vs{7a}Ὡς σοὶ ὑπάρχοι κατὰ τὸ πλῆθος, ποίησον ἐξ αὐτῶν ἐλεημοσύνην· ἐὰν ὀλίγον σοι ὑπάρχῃ, κατὰ τὸ ὀλίγον μὴ φοβοῦ ποιεῖν ἐλεημοσύνην.
\vs{7b}Θέμα γὰρ ἀγαθὸν θησαυρίζεις σεαυτῷ εἰς ἡμέραν ἀνάγκης.
\vs{7c}Διότι ἐλεημοσύνη ἐκ θανάτου ῥύεται, καὶ οὐκ ἐᾴ εἰσελθεῖν εἰς τὸ σκότος.
\vs{7d}Δῶρον γὰρ ἀγαθόν ἐστιν ἐλεημοσύνη πᾶσι τοῖς ποιοῦσιν αὐτὴν ἐνώπιον τοῦ ὑψίστου.

\vs{7e}Πρόσεχε σεαυτῷ, παιδίον, ἀπὸ πάσης πορνείας, καὶ γυναῖκα πρῶτον λάβε ἀπὸ τοῦ σπέρματος τῶν πατέρων σου· μὴ λάβῃς γυναῖκα ἀλλοτρίαν, ἣ οὐκ ἔστιν ἐκ τῆς φυλῆς τοῦ πατέρος σου, διότι υἱοὶ προφητῶν ἐσμέν, Νῶε, Ἁβραάμ, Ἰσαάκ, Ἰακώβ. Οἱ πατέρες ἡμῶν ἀπὸ τοῦ αἰῶνος, μνήσθητι, παιδίον, ὅτι αὐτοὶ πάντες ἔλαβον γυναῖκας ἐκ τῶν ἀδελφῶν αὐτῶν, καὶ εὐλογήθησαν ἐν τοῖς τέκνοις αὐτῶν, καὶ τὸ σπέρμα αὐτῶν κληρονομήσει γῆν.

\vs{7f}Καὶ νῦν, παιδίον, ἀγάπα τοὺς ἀδελφούς σου, καὶ μὴ ὑπερηφανεύου τῇ καρδίᾳ σου ἀπὸ τῶν ἀδελφῶν σου, καὶ τῶν υἱῶν καὶ θυγατέρων τοῦ λαοῦ σου, λαβεῖν σεαυτῷ ἐξ αὐτῶν γυναῖκα· διότι ἐν τῇ ὑπερηφανίᾳ ἀπώλεια καὶ ἀκαταστασία πολλή, καὶ ἐν τῇ ἀχρειότητι ἐλάττωσις καὶ ἔνδεια μεγάλη· ἡ γὰρ ἀχρειότης μήτηρ ἐστὶ τοῦ λιμοῦ.
\vs{7g}Μισθὸς παντὸς ἀνθρώπου ὃς ἐὰν ἐργάσηται, παρὰ σοὶ μὴ αὐλισθήτω, ἀλλʼ ἀπόδος αὐτῷ παρʼ αὐτίκα· ἐὰν δουλεύσῃς τῷ Θεῷ, ἀποδοθήσεταί σοι· πρόσεχε σεαυτῷ, παιδίον, ἐν πᾶσι τοῖς ἔργοις σου, καὶ ἴσθι πεπαιδευμένος ἐν πάσῃ ἀναστροφῇ σου.
\vs{7h}Καὶ ὃ μισεῖς, μηδενὶ ποιήσῃς· οἶνον εἰς μέθην μὴ πίῃς, καὶ μὴ πορευθήτω μετὰ σοῦ μέθη ἐν τῇ ὁδῷ σου.

\vs{7i}Ἐκ τοῦ ἄρτου σου δίδου πεινῶντι, καὶ ἐκ τῶν ἱματίων σου τοῖς γυμνοῖς· πᾶν ὃ ἐὰν περισσεύσῃ σοι, ποίει ἐλεημοσύνην, καὶ μὴ φθονεσάτω σου ὁ ὀφθαλμὸς ἐν τῷ ποιεῖν σε ἐλεημοσύνην.
\vs{7k}Εκχεον τοὺς ἄρτους σου ἐπὶ τὸν τάφον τῶν δικαίων, καὶ μὴ δῷς τοῖς ἁμαρτωλοῖς.
\vs{7l}Συμβουλίαν παρὰ παντὸς φρονίμου ζήτησον, καὶ μὴ καταφρονήσῃς ἐπὶ πάσης συμβουλίας χρησίμης.

\vs{19}Καὶ ἐν παντὶ καιρῷ εὐλόγει Κύριον τὸν Θεὸν, καὶ παρʼ αὐτοῦ αἴτησον, ὅπως αἱ ὁδοί σου εὐθεῖαι γένωνται, καὶ πᾶσαι αἱ τρίβοι καὶ βουλαὶ σου εὐοδωθῶσι· διότι πᾶν ἔθνος οὐκ ἔχει βουλὴν, ἀλλʼ αὐτὸς ὁ Κύριος δίδωσι πάντα τὰ ἀγαθὰ, καὶ ὃν ἐὰν θέλῃ, ταπεινοῖ καθὼς βούλεται· καὶ νῦν, παιδίον, μνημόνευε τῶν ἐντολῶν μου, καὶ μὴ ἐξαλειφθήτωσαν ἐκ τῆς καρδίας σου.

\vs{20}Καὶ νῦν ὑποδεικνύω σοι τὰ δέκα τάλαντα τοῦ ἀργυρίου, ἃ παρεθέμην Γαβαήλῳ τῷ τοῦ Γαβρία ἐν Ῥάγοις τῆς Μηδίας.
\vs{21}Καὶ μὴ φοβοῦ, παιδίον, ὅτι ἐπτωχεύσαμεν· ὑπάρχει σοι πολλὰ, ἐὰν φοβηθῇς τὸν Θεὸν, καὶ ἀποστῇς ἀπὸ πάσης ἁμαρτίας, καὶ ποιήσῃς τὸ ἀρεστὸν ἐνώπιον αὐτοῦ.

\ch{5}
Καὶ ἀποκριθεὶς Τωβίας εἶπεν αὐτῷ, πάτερ, ποιήσω πάντα ὅσα ἐντέταλσαί μοι.
\vs{2}Ἀλλὰ πῶς δυνήσομαι λαβεῖν τὸ ἀργύριον, καὶ οὐ γινώσκω αὐτόν;
\vs{3}Καὶ ἔδωκεν αὐτῷ τὸ χειρόγραφον, καὶ εἶπεν αὐτῷ, ζήτησον σεαυτῷ ἄνθρωπον ὃς συμπορεύσεταί σοι, καὶ δώσω αὐτῷ μισθὸν ἕως ζῶ, καὶ λάβε πορευθεὶς τὸ ἀργύριον.

\vs{4}Καὶ ἐπορεύθη ζητῆσαι ἄνθρωπον, καὶ εὗρε Ῥαφαὴλ, ὃς ἦν ἄγγελος, καὶ οὐκ ᾔδει·
\vs{5}καὶ εἶπεν αὐτῷ, εἰ δύναμαι πορευθῆναι μετὰ σοῦ ἐν Ῥάγοις τῆς Μηδίας, καὶ εἰ ἔμπειρος εἶ τῶν τόπων.
\vs{6}Καὶ εἶπεν αὐτῷ ὁ ἄγγελος, πορεύσομαι μετὰ σοῦ, καὶ τῆς ὁδοῦ ἐμπειρῶ, καὶ παρὰ Γαβαὴλ τὸν ἀδελφὸν ἡμῶν ηὐλίσθην.

\vs{7}Καὶ εἶπεν αὐτῷ Τωβίας ὑπόμεινόν με, καὶ ἐρῶ τῷ πατρί.
\vs{8}Καὶ εἶπεν αὐτῷ, πορεύου, καὶ μὴ χρονίσῃς· καὶ εἰσελθὼν, εἶπε τῷ πατρὶ, ἰδοὺ εὕρηκα ὃς συμπορεύσεταί μοι· ὁ δὲ εἶπε, φώνησον αὐτὸν πρὸς μὲ, ἵνα ἐπιγνῶ ποίας φυλῆς ἐστι, καὶ εἰ πιστὸς τοῦ πορευθῆναι μετὰ σοῦ.
\vs{9}Καὶ ἐκάλεσεν αὐτόν· καὶ εἰσῆλθε, καὶ ἠσπάσαντο ἀλλήλους.

\vs{10}Καὶ εἶπεν αὐτῷ Τωβὶτ, ἀδελφὲ, ἐκ ποίας φυλῆς καὶ ἐκ ποίας πατριᾶς εἶ σύ; ὑπόδειξόν μοι.
\vs{11}Καὶ εἶπεν αὐτῷ, φυλὴν καὶ πατριὰν σὺ ζητεῖς; ἢ μίσθιον, ὃς συμπορεύσεται μετὰ τοῦ υἱοῦ σου; καὶ εἶπεν αὐτῷ Τωβὶτ, βούλομαι, ἀδελφὲ, ἐπιγνῶναι τὸ γένος σου, καὶ τὸ ὄνομα.

\vs{12}Ὃς δὲ εἶπεν, ἐγὼ Ἀζαρίας Ἁνανίου τοῦ μεγάλου, τῶν ἀδελφῶν σου.
\vs{13}Καὶ εἶπεν αὐτῷ, ὑγιαίνων ἔλθοις, ἀδελφέ· καὶ μή μοι ὀργισθῇς, ὅτι ἐζήτησα τὴν φυλήν σου, καὶ τὴν πατριάν σου ἐπιγνῶναι· καὶ σὺ τυγχάνεις ἀδελφός μου ἐκ τῆς καλῆς καὶ ἀγαθῆς γενεᾶς· ἐπεγίνωσκον γὰρ ἐγὼ Ἀνανίαν καὶ Ἰωνάθαν τοὺς υἱοὺς Σεμεῒ τοῦ μεγάλου, ὡς ἐπορευόμεθα κοινῶς εἰς Ἱεροσόλυμα προσκυνεῖν, ἀναφέροντες τὰ πρωτότοκα, καὶ τὰς δεκάτας τῶν γενυημάτων, καὶ οὐκ ἐπλανήθησαν ἐν τῇ πλάνῃ τῶν ἀδελφῶν ἡμῶν· ἐκ ῥίζης καλῆς εἶ, ἀδελφέ.
\vs{14}Ἀλλὰ εἶπόν μοι τίνα σοι ἔσομαι μισθὸν διδόναι; δραχμὴν τῆς ἡμέρας, καὶ τὰ δέοντά σοι ὡς καὶ τῷ υἱῷ μου,
\vs{15}καὶ ἔτι προσθήσω σοι ἐπὶ τὸν μισθὸν, ἐὰν ὑγιαίνοντες ἐπιστρέψητε.

\vs{16}Καὶ εὐδόκησαν οὕτως· καὶ εἶπε πρὸς Τωβίαν, ἕτοιμος γίνου πρὸς τὴν ὁδὸν, καὶ εὐοδωθείητε· καὶ ἡτοίμασεν ὁ υἱὸς αὐτοῦ τὰ πρὸς τὴν ὁδόν· καὶ εἶπεν αὐτῷ ὁ πατὴρ αὐτοῦ, πορεύου μετὰ τοῦ ἀνθρώπου τούτου, ὁ δὲ ἐν τῷ οὐρανῷ οἰκῶν Θεὸς εὐοδώσει τὴν ὁδὸν ὑμῶν, καὶ ὁ ἄγγελος αὐτοῦ συμπορευθήτω ὑμῖν· καὶ ἐξῆλθαν ἀμφότεροι ἀπελθεῖν, καὶ ὁ κύων τοῦ παιδαρίου μετʼ αὐτῶν.

\vs{17}Ἔκλαυσε δὲ Ἄννα ἡ μήτηρ αὐτοῦ, καὶ εἶπε πρὸς Τωβὶτ, τί ἐξαπέστειλας τὸ παιδίον ἡμῶν; ἢ οὐχὶ ἡ ῥάβδος τῆς χειρὸς ἡμῶν ἐστιν ἐν τῷ εἰσπορεύεσθαι αὐτὸν καὶ ἐκπορεύεσθαι ἐνώπιον ἡμῶν;
\vs{18}Ἀργύριον τῷ ἀργυρίῳ μὴ φθάσαι, ἀλλὰ περίψημα τοῦ παιδίου ἡμῶν γένοιτο.
\vs{19}Ὡς γὰρ δέδοται ἡμῖν ζῇν παρὰ τοῦ Κυρίου, τοῦτο ἱκανὸν ἡμῖν ὑπάρχει.
\vs{20}Καὶ εἶπεν αὐτῇ Τωβίτ, μὴ λόγον ἔχε ἀδελφὴ, ὑγιαίνων ἐλεύσεται, καὶ οἱ ὀφθαλμοί σου ὄψονται αὐτόν.
\vs{21}Ἄγγελος γὰρ ἀγαθὸς συμπορεύσεται αὐτῷ, καὶ εὐοδωθήσεται ἡ ὁδὸς αὐτοῦ, καὶ ὑποστρέψει ὑγιαίνων.
\vs{22}Καὶ ἐπαύσατο κλαίουσα.

\ch{6}
Οἱ δὲ πορευόμενοι τὴν ὁδὸν, ἦλθον ἑσπέρας ἐπὶ τὸν Τίγριν ποταμόν, καὶ ηὐλίζοντο ἐκεῖ.
\vs{2}Τὸ δὲ παιδάριον κατέβη περικλύσασθαι, καὶ ἀνεπήδησεν ἰχθὺς ἀπὸ τοῦ ποταμοῦ, καὶ ἐβουλήθη καταπιεῖν τὸ παιδάριον.
\vs{3}Ὁ δὲ ἄγγελος εἶπεν αὐτῷ, ἐπιλαβοῦ τοῦ ἰχθύος· καὶ ἐκράτησε τὸν ἰχθῦν τὸ παιδάριον, καὶ ἀνέβαλεν αὐτὸν ἐπὶ τὴν γῆν.
\vs{4}Καὶ εἶπεν αὐτῷ ὁ ἄγγελος, ἀνάτεμε τὸν ἰχθύν, καὶ λαβὼν τὴν καρδίαν καὶ τὸ ἧπαρ καὶ τὴν χολὴν, θὲς ἀσφαλῶς.
\vs{5}Καὶ ἐποίησε τὸ παιδάριον ὡς εἶπεν αὐτῷ ὁ ἄγγελος· τὸν δὲ ἰχθῦν ὀπτήσαντες, ἔφαγον· καὶ ὥδευον ἀμφότεροι, ἕως οὗ ἤγγισαν ἐν Ἐκβατάνοις.

\vs{6}Καὶ εἶπε τὸ παιδάριον τῷ ἀγγέλῳ, Ἀζαρία ἀδελφὲ, τί ἐστιν ἡ καρδία καὶ τὸ ἧπαρ καὶ ἡ χολὴ τοῦ ἰχθύος;
\vs{7}Καὶ εἶπεν αὐτῷ, ἡ καρδία καὶ τὸ ἧπαρ, ἐάν τινα ὀχλῇ δαιμόνιον ἢ πνεῦμα πονηρὸν, ταῦτα δεῖ καπνίσαι ἐνώπιον ἀνθρώπου, ἢ γυναικὸς, καὶ μηκέτι ὀχληθῇ.
\vs{8}Ἡ δὲ χολὴ, ἔγχρισαι ἄνθρωπον ὃς ἔχει λευκώματα ἐν τοῖς ὀφθαλμοῖς, καὶ ἰαθήσεται.

\vs{9}Ὡς δὲ προσήγγισαν τῇ Ῥάγῃ,
\vs{10}εἶπεν ὁ ἄγγελος τῷ παιδαρίῳ, ἀδελφὲ, σήμερον αὐλισθησόμεθα παρὰ Ῥαγουὴλ, καὶ αὐτὸς συγγενής σου ἐστὶ, καὶ ἔστιν αὐτῷ θυγατηρ ὀνόματι Σάῤῥα· λαλήσω περὶ αὐτῆς, τοῦ δοθῆναί σοι αὐτὴν εἰς γυναῖκα,
\vs{11}καὶ ὅτι σοι ἐπιβάλλει ἡ κληρονομία αὐτῆς, καὶ σὺ μόνος εἶ ἐκ τοῦ γένους αὐτῆς·
\vs{12}Καὶ τὸ κοράσιον καλὸν καὶ φρόνιμόν ἐστι· καὶ νῦν ἄκουσόν μου, καὶ λαλήσω τῷ πατρὶ αὐτῆς, καὶ ὅταν ὑποστρέψομεν ἐκ Ῥαγῶν, ποιήσομεν τὸν γάμον· διότι ἐπίσταμαι Ῥαγουὴλ ὅτι οὐ μὴ δῷ αὐτὴν ἀνδρὶ ἑτέρῳ κατὰ τὸν νόμον Μωυσῇ, ἢ ὀφειλήσει θάνατον, ὅτι τὴν κληρονομίαν σοι καθήκει λαβεῖν, ἢ πάντα ἄνθρωπον.

\vs{13}Τότε εἶπε τὸ παιδάριον τῷ ἀγγέλῳ, Ἀζαρία ἀδελφὲ, ἀκήκοα ἐγὼ τὸ κοράσιον δεδόσθαι ἑπτὰ ἀνδράσι, καὶ πάντας ἐν τῷ νυμφῶνι ἀπολωλότας·
\vs{14}καὶ νῦν ἐγὼ μόνος εἰμὶ τῷ πατρὶ, καὶ φοβοῦμαι μὴ εἰσελθὼν ἀποθάνω καθὼς καὶ οἱ πρότεροι, ὅτι δαιμόνιον φιλεῖ αὐτὴν, ὃ οὐκ ἀδικεῖ οὐδένα πλὴν τῶν προσαγόντων αὐτῇ· καὶ νῦν ἐγὼ φοβοῦμαι μὴ ἀποθάνω, καὶ κατάξω τὴν ζωὴν τοῦ πατρός μου καὶ τῆς μητρός μου μετʼ ὀδύνης ἐπʼ ἐμοὶ εἰς τὸν τάφον αὐτῶν, καὶ υἱὸς ἕτερος οὐκ ὑπάρχει αὐτοῖς ὃς θάψει αὐτούς.

\vs{15}Εἶπε δὲ αὐτῷ ὁ ἄγγελος, οὐ μέμνησαι τῶν λόγων ὧν ἐνετείλατό σοι ὁ πατήρ σου, ὑπὲρ τοῦ λαβεῖν σε γυναῖκα ἐκ τοῦ γένους σου; καὶ νῦν ἄκουσόν μου, ἀδελφὲ, διότι σοι ἔσται εἰς γυναῖκα, καὶ τοῦ δαιμονίου μηδένα λόγον ἔχε, ὅτι τὴν νύκτα ταύτην δοθήσεταί σοι αὕτη εἰς γυναῖκα.
\vs{16}Καὶ ἐὰν εἰσέλθῃς εἰς τὸν νυμφῶνα, λήψῃ τέφραν θυμιαμάτων,
\vs{17}καὶ καπνίσεις, καὶ ὀσφρανθήσεται τὸ δαιμόνιον, καὶ φεύξεται, καὶ οὐκ ἐπανελεύσεται εἰς τὸν αἰῶνα τοῦ αἰῶνος. ὅταν δὲ προσπορεύῃ αὐτῇ, ἐγέρθητε ἀμφότεροι, καὶ βοήσατε πρὸς τὸν ἐλεήμονα Θεὸν, καὶ σώσει ὑμᾶς, καὶ ἐλεήσει· μὴ φοβοῦ, ὅτι σοὶ αὕτη ἡτοιμασμένη ἦν ἀπὸ τοῦ αἰῶνος, καὶ σὺ αὐτὴν σώσεις, καὶ πορεύσεται μετὰ σοῦ, καὶ ὑπολαμβάνω ὅτι σοὶ ἔσται ἐξ αὐτῆς παιδία· καὶ ὡς ἤκουσε Τωβίας ταῦτα, ἐφίλησεν αὐτὴν, καὶ ἡ ψυχὴ αὐτοῦ ἐκολλήθη σφόδρα αὐτῇ· καὶ ἦλθεν εἰς Ἐκβάτανα.

\ch{7}
Καὶ παρεγένετο εἰς τὴν οἰκίαν Ῥαγουήλ· καὶ Σάῤῥα δὲ ὑπήντησεν αὐτῷ, καὶ ἐχαιρέτισεν αὐτὸν, καὶ αὐτὸς αὐτούς· καὶ εἰσήγαγεν αὐτοὺς εἰς τὴν οἰκίαν.
\vs{2}Καὶ εἶπεν Ῥαγουὴλ Ἔδνᾳ τῇ γυναικὶ αὐτοῦ, ὡς ὅμοιος ὁ νεανίσκος Τωβὶτ τῷ ἀνεψιῷ μου;

\vs{3}Καὶ ἠρώτησεν αὐτοὺς Ῥαγουὴλ, πόθεν ἐστὲ, ἀδελφοί; καὶ εἶπον αὐτῷ, ἐκ τῶν υἱῶν Νεφθαλὶ τῶν αἰχμαλώτων ἐν Νινευῇ.
\vs{4}Καὶ εἶπεν αὐτοῖς, γινώσκετε Τωβὶτ τὸν ἀδελφὸν ἡμῶν;
\vs{5}Οἱ δὲ εἶπαν, καὶ ζῇ, καὶ ὑγιαίνει· καὶ εἶπε Τωβίας, πατήρ μου ἐστί.
\vs{6}Καὶ ἀνεπήδησε Ῥαγουὴλ, καὶ κατεφίλησεν αὐτὸν, καὶ ἔκλαυσε,
\vs{7}καὶ εὐλόγησεν αὐτὸν, καὶ εἶπεν αὐτῷ, ὁ τοῦ καλοῦ καὶ ἀγαθοῦ ἀνθρώπου υἱός· καὶ ἀκούσας ὅτι Τωβὶτ ἀπώλεσε τοὺς ὀφθαλμοὺς ἑαυτοῦ, ἐλυπήθη καὶ ἔκλαυσε.

\vs{8}Καὶ Ἔδνα ἡ γυνὴ αὐτοῦ καὶ Σάῤῥα ἡ θυγάτηρ αὐτοῦ ἔκλαυσαν, καὶ ὑπεδέξαντο αὐτοὺς προθύμως· καὶ ἔθυσαν κριὸν προβάτων, καὶ παρέθηκαν ὄψα πλείονα· εἶπε δὲ Τωβίας τῷ Ῥαφαὴλ, Ἀζαρία ἀδελφὲ, λάλησον ὑπὲρ ὧν ἔλεγες ἐν τῇ πορείᾳ, καὶ τελεσθήτω τὸ πρᾶγμα.

\vs{9}Καὶ μετέδωκε τὸν λόγον τῷ Ῥαγουήλ·
\vs{10}καὶ εἶπε Ῥαγουὴλ πρὸς Τωβίαν, φάγε, πίε, καὶ ἡδέως γίνου, σοὶ γὰρ καθήκει τὸ παιδίον μου λαβεῖν· πλὴν ὑποδείξω σοι τὴν ἀλήθειαν.
\vs{11}Ἔδωκα τὸ παιδίον μου ἑπτὰ ἀνδράσι, καὶ ὁπότε ἐὰν εἰσεπορεύοντο πρὸς αὐτὴν, ἀπέθνησκον ὑπὸ τὴν νύκτα· ἀλλὰ τὸ νῦν ἔχον. ἡδέως γίνου· καὶ εἶπε Τωβίας, οὐ γεύομαι οὐδὲν ὧδε, ἕως ἂν στήσητε καὶ σταθῆτε πρὸς μέ.
\vs{12}Καὶ εἶπε Ῥαγουὴλ, κομίζου αὐτὴν ἀπὸ τοῦ νῦν κατὰ τὴν κρίσιν· σὺ δὲ ἀδελφὸς εἶ αὐτῆς, καὶ αὐτή σου ἐστίν· ὁ δὲ ἐλεήμων Θεὸς εὐοδώσει ὑμῖν τὰ κάλλιστα.

\vs{13}Καὶ ἐκάλεσε Σάῤῥαν τὴν θυγατέρα αὐτοῦ, καὶ λαβὼν τῆς χειρὸς αὐτῆς, παρέδωκεν αὐτὴν Τωβίᾳ γυναῖκα, καὶ εἶπεν, ἰδοὺ κατὰ τὸν νόμον Μωυσέως κομίζου αὐτὴν, καὶ ἄπαγε πρὸς τὸν πατέρα σου· καὶ εὐλόγησεν αὐτούς.
\vs{14}Καὶ ἐκάλεσεν Ἔδναν τὴν γυναῖκα αὐτοῦ· καὶ λαβὼν βιβλίον, ἔγραψε συγγραφὴν, καὶ ἐσφραγίσατο.
\vs{15}Καὶ ἤρξαντο ἐσθίειν.

\vs{16}Καὶ ἐκάλεσε Ῥαγουὴλ Ἔδναν τὴν γυναῖκα αὐτοῦ, καὶ εἶπεν αὐτῇ, ἀδελφὴ ἑτοίμασον τὸ ἕτερον ταμεῖον, καὶ εἰσάγαγε αὐτήν.
\vs{17}Καὶ ἐποίησεν ὡς εἶπε· καὶ εἰσήγαγεν αὐτὴν ἐκεῖ, καὶ ἔκλαυσε· καὶ ἀπεδέξατο τὰ δάκρυα τῆς θυγατρὸς αὐτῆς, καὶ εἶπεν αὐτῇ,
\vs{18}θάρσει τέκνον, ὁ Κύριος τοῦ οὐρανοῦ καὶ τῆς γῆς δῴη σοι χάριν ἀντὶ τῆς λύπης σου ταύτης, θάρσει θύγατερ.

\ch{8}
Ὅτε δὲ συνετέλεσαν δειπνοῦντες, εἰσήγαγον Τωβίαν πρὸς αὐτήν.
\vs{2}Ὁ δὲ πορεύομενος ἐμνήσθη τῶν λόγων Ῥαφαὴλ, καὶ ἔλαβε τὴν τέφραν τῶν θυμιαμάτων, καὶ ἐπέθηκε τὴν καρδίαν τοῦ ἰχθύος καὶ τὸ ἧπαρ, καὶ ἐκάπνισεν.
\vs{3}Ὅτε δὲ ὠσφράνθη τὸ δαιμόνιον τῆς ὀσμῆς, ἔφυγεν εἰς τὰ ἀνώτατα Αἰγύπτου, καὶ ἔδησεν αὐτὸ ὁ ἄγγελος.

\vs{4}Ὡς δὲ συνεκλείσθησαν ἀμφότεροι, ἀνέστη Τωβίας ἀπὸ τῆς κλίνης, καὶ εἶπεν, ἀνάστηθι ἀδελφὴ, καὶ προσευξώμεθα ἵνα ἐλεήσῃ ἡμᾶς ὁ Κύριος.
\vs{5}Καὶ ἤρξατο Τωβίας λέγειν, εὐλογητὸς εἶ ὁ Θεὸς τῶν πατέρων ἡμῶν, καὶ εὐλογητὸν τὸ ὄνομά σου τὸ ἅγιον καὶ ἔνδοξον εἰς τοὺς αἰῶνας· εὐλογησάτωσάν σε οἱ οὐρανοὶ, καὶ πᾶσαι αἱ κτίσεις σου.
\vs{6}Σὺ ἐποίησας Ἀδάμ, καὶ ἔδωκας αὐτῷ βοηθὸν Εὖαν στήριγμα τὴν γυναῖκα αὐτοῦ· ἐκ τούτων ἐγενήθη τὸ ἀνθρώπων σπέρμα· σὺ εἶπας, οὐ καλὸν εἶναι τὸν ἄνθρωπον μόνον, ποιήσωμεν αὐτῷ βοηθὸν ὅμοιον αὐτῷ.
\vs{7}Καὶ νῦν, Κύριε, οὐ διὰ πορνείαν ἐγὼ λαμβάνω τὴν ἀδελφήν μου ταύτην, ἀλλὰ ἐπʼ ἀληθείας ἐπίταξον ἐλεῆσαί με, καὶ αὐτῇ συγκαταγηρᾶσαι.
\vs{8}Καὶ εἶπε μετʼ αὐτοῦ, ἀμήν.

\vs{9}Καὶ ἐκοιμήθησαν ἀμφότεροι τὴν νύκτα· καὶ ἀναστὰς Ῥαγουὴλ ἐπορεύθη, καὶ ὤρυξε τὰφον,
\vs{10}λέγων, μὴ καὶ οὗτος ἀποθάνῃ;
\vs{11}Καὶ ἦλθε Ῥαγουὴλ εἰς τὴν οἰκίαν ἑαυτοῦ,
\vs{12}καὶ εἶπεν Ἔδνᾳ τῇ γυναικὶ αὐτοῦ, ἀπόστειλον μίαν τῶν παιδισκῶν, καὶ ἰδέτωσαν εἰ ζῇ· εἰ δὲ μὴ, ἵνα θάψωμεν αὐτὸν, καὶ μηδεὶς γνῷ.
\vs{13}Καὶ εἰσῆλθεν ἡ παιδίσκη ἀνοίξασα τὴν θύραν, καὶ εὗρε τοὺς δύο καθεύδοντας,
\vs{14}καὶ ἐξελθοῦσα ἀπήγγειλεν αὐτοῖς, ὅτι ζῇ.

\vs{15}Καὶ εὐλόγησε Ῥαγουὴλ τὸν Θεὸν, λέγων, εὐλογητὸς εἶ σὺ ὁ Θεὸς ἐν πάσῃ εὐλογίᾳ καθαρᾷ καὶ ἁγίᾳ· καὶ εὐλογείτωσάν σε οἱ ἅγιοί σου, καὶ πᾶσαι αἱ κτίσεις σου, καὶ πάντες οἱ ἄγγελοί σου, καὶ οἱ ἐκλεκτοί σου· εὐλογείτωσάν σε εἰς τοὺς αἰῶνας.
\vs{16}Εὐλογητὸς εἶ, ὅτι ηὔφρανάς με, καὶ οὐκ ἐγένετό μοι καθὼς ὑπενόουν, ἀλλὰ κατὰ τὸ πολὺ ἔλεός σου ἐποίησας μεθʼ ἡμῶν.
\vs{17}Εὐλογητὸς εἶ, ὅτι ἠλέησας δύο μονογενεῖς· ποίησον αὐτοῖς, δέσποτα, ἔλεος, συντέλεσον τὴν ζωὴν αὐτῶν ἐν ὑγιείᾳ μετʼ εὐφροσύνης καὶ ἐλέους.
\vs{18}Ἐκέλευσε δὲ τοῖς οἰκέταις χῶσαι τὸν τάφον.

\vs{19}Καὶ ἐποίησεν αὐτοῖς γάμον ἡμερῶν δεκατεσσάρων.
\vs{20}Καὶ εἶπεν αὐτῷ Ῥαγουὴλ, πρινὴ συντελεσθῆναι τὰς ἡμέρας τοῦ γάμου, ἐνόρκως, μὴ ἐξελθεῖν αὐτὸν ἐὰν μὴ πληρωθῶσιν αἱ δεκατέσσαρες ἡμέραι τοῦ γάμου,
\vs{21}καὶ τότε λαβόντα τὸ ἥμισυ τῶν ὑπαρχόντων αὐτοῦ πορεύεσθαι μεθʼ ὑγείας πρὸς τὸν πατέρα, καὶ τὰ λοιπὰ ὅταν ἀποθάνω, καὶ ἡ γυνή μου.

\ch{9}
Καὶ ἐκάλεσε Τωβίας τὸν Ῥαφαὴλ, καὶ εἶπεν αὐτῷ,
\vs{2}Ἀζαρία ἀδελφὲ, λάβε μετὰ σεαυτοῦ παῖδα καὶ δύο καμήλους, καὶ πορεύθητι ἐν Ῥάγοις τῆς Μηδίας παρὰ Γαβαὴλ, καὶ κόμισαί μοι τὸ ἀργύριον, καὶ αὐτὸν ἄγε μοι εἰς τόν γάμον,
\vs{3}διότι ὀμώμοκε Ῥαγουὴλ, μὴ ἐξελθεῖν με.
\vs{4}Καὶ ὁ πατήρ μου ἀριθμεῖ τὰς ἡμέρας, καὶ ἐὰν χρονίσω μέγα, ὀδυνηθήσεται λίαν.
\vs{5}Καὶ ἐπορεύθη Ῥαφαὴλ, καὶ ηὐλίσθη παρὰ Γαβαὴλ, καὶ ἔδωκεν αὐτῷ τὸ χειρόγραφον· ὃς δὲ προήνεγκε τὰ θυλάκια ἐν ταῖς σφραγίσι, καὶ ἔδωκεν αὐτῷ.

\vs{6}Καὶ ὤρθρευσαν κοινῶς, καὶ ἦλθον εἰς τὸν γάμον· καὶ εὐλόγησε Τωβίας τὴν γυναῖκα αὐτοῦ.

\ch{10}
Καὶ Τωβὶτ ὁ πατὴρ αὐτοῦ ἐλογίσατο ἑκάστης ἡμέρας· καὶ ὡς ἐπληρώθησαν αἱ ἡμέραι τῆς πορείας, καὶ οὐκ ἤρχετο,
\vs{2}εἶπε μήποτε κατῄσχυνται; ἢ μήποτε ἀπέθανε Γαβαήλ, καὶ οὐδεὶς αὐτῷ δίδωσι τὸ ἀργύριον;
\vs{3}Καὶ ἐλυπεῖτο λίαν.
\vs{4}Εἶπε δὲ αὐτῷ ἡ γυνὴ, ἀπώλετο τὸ παιδίον, διότι κεχρόνικε· καὶ ἤρξατο θρηνεῖν αὐτὸν, καὶ εἶπεν,
\vs{5}οὐ μέλει μοι, τέκνον, ὅτι ἀφῆκά σε τὸ φῶς τῶν ὀφθαλμῶν μου.

\vs{6}Καὶ Τωβὶτ λέγει αὐτῇ, σίγα, μὴ λόγον ἔχε, ὑγιαίνει.
\vs{7}Καὶ εἶπεν αὐτῷ, σίγα, μὴ πλάνα με, ἀπώλετο τὸ παιδίον μου· καὶ ἐπορεύετο καθʼ ἡμέραν εἰς τὴν ὁδὸν ἔξω, οἵας ἀπῆλθεν· ἡμέρας τε ἄρτον οὐκ ἤσθιε, τὰς δὲ νύκτας οὐ διελίμπανε θρηνοῦσα Τωβίαν τὸν υἱὸν αὐτῆς, ἕως οὗ συνετελέσθησαν αἱ δεκατέσσαρες ἡμέραι τοῦ γάμου, ἃς ὤμοσε Ῥαγουὴλ ποιῆσαι αὐτὸν ἐκεῖ·
\vs{8}εἶπε δὲ Τωβίας τῷ Ῥαγουὴλ, ἐξαπόστειλόν με, ὅτι ὁ πατήρ μου καὶ ἡ μήτηρ μου οὐκέτι ἐλπίζουσιν ὄψεσθαί με.
\vs{9}Εἶπε δὲ αὐτῷ ὁ πενθερὸς, μεῖνον παρʼ ἐμοὶ, κᾀγὼ ἐξαποστελῶ πρὸς τὸν πατέρα σου, καὶ δηλώσουσιν αὐτῷ τὰ κατά σε.
\vs{10}Καὶ Τωβίας λέγει, ἐξαπόστειλόν με πρὸς τὸν πατέρα μου.

\vs{11}Ἀναστὰς δὲ Ῥαγουὴλ, ἔδωκεν αὐτῷ Σάῤῥαν τὴν γυναῖκα αὐτοῦ, καὶ τὸ ἥμισυ τῶν ὑπαρχόντων, σώματα καὶ κτήνη καὶ ἀργύριον,
\vs{12}καὶ εὐλογήσας αὐτοὺς ἐξαπέστειλε, λέγων, εὐοδώσει ὑμᾶς τέκνα ὁ Θεὸς τοῦ οὐρανοῦ πρὸ τοῦ με ἀποθανεῖν.
\vs{13}Καὶ εἶπε τῇ θυγατρὶ αὐτοῦ, τίμα τοὺς πενθερούς σου, αὐτοὶ νῦν γονεῖς σου εἰσὶν, ἀκούσαιμί σου ἀκοὴν καλήν· καὶ ἐφίλησεν αὐτήν· καὶ Ἔδνα εἶπε πρὸς Τωβίαν, ἀδελφὲ ἀγαπητὲ, ἀποκαταστήσαι σε ὁ Κύριος τοῦ οὐρανοῦ, καὶ δῴη μοι ἰδεῖν σου παιδία ἐκ Σάῤῥας τῆς θυγατρός μου, ἵνα εὐφρανθῶ ἐνώπιον τοῦ Κυρίου· καὶ ἰδοὺ παρατίθεμαί σοι τὴν θυγατέρα μου ἐν παρακαταθήκῃ, καὶ μὴ λυπήσῃς αὐτήν.

\ch{11}
Μετὰ ταῦτα ἐπορεύετο καὶ Τωβίας εὐλογῶν τὸν Θεὸν, ὅτι εὐώδωσε τὴν ὁδὸν αὐτοῦ· καὶ κατευλόγει Ῥαγουὴλ, καὶ Ἔδναν τὴν γυναῖκα αὐτοῦ· καὶ ἐπορεύετο μέχρις οὗ ἐγγίσαι αὐτοὺς εἰς Νινευή.

\vs{2}Καὶ εἶπε Ῥαφαὴλ πρὸς Τωβίαν, οὐ γινώσκεις, ἀδελφὲ, πῶς ἀφῆκας τὸν πατέρα σου;
\vs{3}Προδράμωμεν ἔμπροσθεν τῆς γυναικός σου, καὶ ἑτοιμάσωμεν τὴν οἰκίαν·
\vs{4}λάβε δὲ παρὰ χεῖρα τὴν χολὴν τοῦ ἰχθύος· καὶ ἐπορεύθησαν, καὶ συνῆλθεν ὁ κύων ὄπισθεν αὐτῶν.
\vs{5}Καὶ Ἄννα ἐκάθητο περιβλεπομένη εἰς τὴν ὁδὸν τὸν παῖδα αὐτῆς.
\vs{6}Καὶ προσενόησεν αὐτὸν ἐρχόμενον, καὶ εἶπε τῷ πατρὶ αὐτοῦ, ἰδοὺ ὁ υἱὸς μου ἔρχεται, καὶ ὁ ἄνθρωπος ὁ πορευθεὶς μετʼ αὐτοῦ.

\vs{7}Καὶ Ῥαφαὴλ εἶπεν, ἐπίσταμαι ἐγὼ, ὅτι ἀνοίξει τοὺς ὀφθαλμοὺς ὁ πατήρ σου.
\vs{8}Σὺ ἔγχρισον τὴν χολὴν εἰς τοὺς ὀφθαλμοὺς αὐτοῦ, καὶ δηχθεὶς διατρίψει, καὶ ἀποβαλεῖται τὰ λευκώματα, καὶ ὄψεταί σε.

\vs{9}Καὶ προσδραμοῦσα Ἄννα ἐπέπεσεν ἐπὶ τὸν τράχηλον τοῦ υἱοῦ αὐτῆς, καὶ εἶπεν αὐτῷ, εἶδόν σε παιδίον, ἀπὸ τοῦ νῦν ἀποθανοῦμαι· καὶ ἔκλαυσαν ἄμφότεροι.
\vs{10}Καὶ Τωβὶτ ἐξήρχετο πρὸς τὴν θύραν, καὶ προσέκοπτεν· ὁ δὲ υἱὸς αὐτοῦ προσέδραμεν αὐτῷ,
\vs{11}καὶ ἐπελάβετο τοῦ πατρὸς αὐτοῦ, καὶ προσέπασε τὴν χολὴν ἐπὶ τοὺς ὀφθαλμοὺς τοῦ πατρὸς αὐτοῦ, λέγων, θάρσει πάτερ.
\vs{12}Ὡς δὲ συνεδήχθησαν, διέτριψε τοὺς ὀφθαλμοὺς αὐτοῦ,
\vs{13}καὶ ἐλεπίσθη ἀπὸ τῶν κάνθων τῶν ὀφθαλμῶν αὐτοῦ τὰ λευκώματα· καὶ ἰδὼν τὸν υἱὸν αὐτοῦ ἐπέπεσεν ἐπὶ τὸν τράχηλον αὐτοῦ,

\vs{14}Καὶ ἔκλαυσε, καὶ εἶπεν, εὐλογητὸς εἶ ὁ Θεὸς, καὶ εὐλογητὸν τὸ ὄνομά σου εἰς τοὺς αἰῶνας, καὶ εὐλογημένοι πάντες οἱ ἅγιοί σου ἄγγελοι,
\vs{15}ὅτι ἐμαστίγωσας καὶ ἠλέησάς με· ἰδοὺ βλέπω Τωβίαν τὸν υἱόν μου· καὶ εἰσῆλθεν ὁ υἱὸς αὐτοῦ χαίρων, καὶ ἀπήγγειλε τῷ πατρὶ αὐτοῦ τὰ μεγαλεῖα τὰ γενόμενα αὐτῷ ἐν τῇ Μηδίᾳ.

\vs{16}Καὶ ἐξῆλθε Τωβίτ εἰς συνάντησιν τῇ νύμφῃ αὐτοῦ χαίρων καὶ εὐλογῶν τὸν Θεὸν πρὸς τῇ πύλῃ Νινευή· καὶ ἐθαύμαζον οἱ θεωροῦντες αὐτὸν πορευόμενον, ὅτι ἔβλεψε.
\vs{17}Καὶ Τωβὶτ ἐξωμολογεῖτο ἐνώπιον αὐτοῦ, ὅτι ἠλέησεν αὐτοὺς ὁ Θεός· καὶ ὡς ἤγγισε Τωβὶτ Σάῤῥᾳ τῇ νύμφῃ αὐτοῦ, κατευλόγησεν αὐτὴν, λέγων, Ἔλθοις ὑγιαίνουσα θύγατερ· εὐλογητὸς ὁ Θεός, ὃς ἤγαγέ σε πρὸς ἡμᾶς, καὶ ὁ πατήρ σου καὶ ἡ μήτηρ σου· καὶ ἐγένετο χαρὰ πᾶσι τοῖς ἐν Νινευὴ ἀδελφοῖς αὐτοῦ.
\vs{18}Καὶ παρεγένετο Ἀχιάχαρος, καὶ Νασβὰς ὁ ἐξάδελφος αὐτοῦ,
\vs{19}καὶ ἤχθη ὁ γάμος Τωβία μετʼ εὐφροσύνης ἡμέρας ἑπτά.

\ch{12}
Καὶ ἐκάλεσε Τωβὶτ Τωβίαν τὸν υἱὸν αὐτοῦ, καὶ εἶπεν αὐτῷ· ὅρα, τέκνον, μισθὸν τῷ ἀνθρώπῳ τῷ συνελθόντι σοι· καὶ προσθεῖναι αὐτῷ δεῖ.
\vs{2}Καὶ εἶπε, πάτερ, οὐ βλάπτομαι δοὺς αὐτῷ τὸ ἥμισυ ὧν ἐνήνοχα,
\vs{3}ὅτι με ἀγήοχέ σοι ὑγιῆ, καὶ τὴν γυναῖκά μου ἐθεράπευσε, καὶ τὸ ἀργύριόν μου ἤνεγκε, καὶ σὲ ὁμοίως ἐθεράπευσε.

\vs{4}Καὶ εἶπεν ὁ πρεσβύτης, δικαιοῦται αὐτῷ.
\vs{5}Καὶ ἐκάλεσε τὸν ἄγγελον, καὶ εἶπεν αὐτῷ, λάβε τὸ ἥμισυ πάντων ὧν ἐνηνόχατε, καὶ ὕπαγε ὑγιαίνων.
\vs{6}Τότε καλέσας τοὺς δύο κρυπτῶς, εἶπεν αὐτοῖς, εὐλογεῖτε τὸν Θεὸν, καὶ αὐτῷ ἐξομολογεῖσθε, καὶ μεγαλωσύνην δίδοτε αὐτῷ, καὶ ἐξομολογεῖσθε αὐτῷ ἐνώπιον πάντων τῶν ζώντων περὶ ὧν ἐποίησε μεθʼ ὑμῶν· ἀγαθὸν τὸ εὐλογεῖν τὸν Θεὸν, καὶ ὑψοῦν τὸ ὄνομα αὐτοῦ, τοὺς λόγους τῶν ἔργων τοῦ Θεοῦ ἐντίμως ὑποδεικνύοντες· καὶ μὴ ὀκνεῖτε ἐξομολογεῖσθαι αὐτῷ.
\vs{7}Μυστήριον βασιλέως καλὸν κρύψαι, τὰ δὲ ἔργα τοῦ Θεοῦ ἀνακαλύπτειν ἐνδόξως· ἀγαθὸν ποιήσατε, καὶ κακὸν οὐχ εὑρήσει ὑμᾶς.
\vs{8}Ἀγαθὸν προσευχὴ μετὰ νηστείας καὶ ἐλεημοσύνης καὶ δικαιοσύνης· ἀγαθὸν τὸ ὀλίγον μετὰ δικαιοσύνης, ἢ πολὺ μετὰ ἀδικίας· καλὸν ποιῆσαι ἐλεημοσύνην ἢ θησαυρίσαι χρυσίον.
\vs{9}Ἐλεημοσύνη γὰρ ἐκ θανάτου ῥύεται, καὶ αὕτη ἀποκαθαριεῖ πᾶσαν ἁμαρτίαν· οἱ ποιοῦντες ἐλεημοσύνας καὶ δικαιοσύνας πλησθήσονται ζωῆς.
\vs{10}Οἱ δὲ ἁμαρτάνοντες πολέμιοί εἰσι τῆς ἑαυτῶν ζωῆς.

\vs{11}Οὐ μὴ κρύψω ἀφʼ ὑμῶν πᾶν ῥῆμα· εἴρηκα δὴ, μυστήριον βασιλέως κρύψαι καλὸν, τὰ δὲ ἔργα τοῦ Θεοῦ ἀνακαλύπτειν ἐνδόξως.
\vs{12}Καὶ νῦν ὅτε προσηύξω σὺ καὶ ἡ νύμφη σου Σάῤῥα, ἐγὼ προσήγαγον τὸ μνημόσυνον τῆς προσευχῆς ὑμῶν ἐνώπιον τοῦ ἁγίου· καὶ ὅτε ἔθαπτες τοὺς νεκροὺς, ὡσαύτως συμπαρήγμην σοι.
\vs{13}Καὶ ὅτε οὐκ ὤκνησας ἀναστῆναι καὶ καταλιπεῖν τὸ ἄριστόν σου, ὅπως ἀπελθὼν περιστείλῃς τὸν νεκρὸν, οὐκ ἔλαθές με ἀγαθοποιῶν, ἀλλὰ σὺν σοὶ ἤμην.
\vs{14}Καὶ νῦν ἀπέστειλέ με ὁ Θεὸς ἰάσασθαί σε καὶ τὴν νύμφην σου Σάῤῥαν.
\vs{15}Ἐγώ εἰμι Ῥαφαήλ, εἷς ἐκ τῶν ἑπτὰ ἁγίων ἀγγέλων οἳ προσαναφέρουσι τὰς προσευχὰς τῶν ἁγίων, καὶ εἰσπορεύονται ἐνώπιον τῆς δόξης τοῦ ἁγίου.

\vs{16}Καὶ ἐταράχθησαν οἱ δύο, καὶ ἔπεσον ἐπὶ πρόσωπον, ὅτι ἐφοβήθησαν.
\vs{17}Καὶ εἶπεν αὐτοῖς, μὴ φοβεῖσθε, εἰρήνη ὑμῖν ἔσται· τὸν δὲ Θεὸν εὐλογεῖτε εἰς τὸν αἰῶνα,
\vs{18}ὅτι οὐ τῇ ἐμαυτοῦ χάριτι, ἀλλὰ τῇ θελήσει τοῦ Θεοῦ ἡμῶν ἦλθον, ὅθεν εὐλογεῖτε αὐτὸν εἰς τὸν αἰῶνα.
\vs{19}Πάσας τὰς ἡμέρας ὠπτανόμην ὑμῖν, καὶ οὐκ ἔφαγον οὐδὲ ἔπιον, ἀλλὰ ὅρασιν ὑμεῖς ἐθεωρεῖτε.
\vs{20}Καὶ νῦν ἐξομολογεῖσθε τῷ Θεῷ, διότι ἀναβαίνω πρὸς τὸν ἀποστείλαντά με, καὶ γράψατε πάντα τὰ συντελεσθέντα εἰς βιβλίον.
\vs{21}Καὶ ἀνέστησαν, καὶ οὐκ ἔτι εἶδον αὐτόν.
\vs{22}Καὶ ἐξωμολογοῦντο τὰ ἔργα τὰ μεγάλα καὶ θαυμαστὰ αὐτοῦ, ὡς ὤφθη αὐτοῖς ὁ ἄγγελος Κυρίου.

\ch{13}
Καὶ Τωβὶτ ἔγραψε προσευχὴν εἰς ἀγαλλίασιν, καὶ εἶπεν,

Εὐλογητὸς ὁ Θεὸς ὁ ζῶν εἰς τοὺς αἰῶνας, καὶ ἡ βασιλεία αὐτοῦ,
\vs{2}ὅτι αὐτὸς μαστιγοῖ καὶ ἐλεεῖ, κατάγει εἰς ᾅδην καὶ ἀνάγει, καὶ οὐκ ἔστιν ὃς ἐκφεύξεται τὴν χεῖρα αὐτοῦ.
\vs{3}Ἐξομολογεῖσθε αὐτῷ οἱ υἱοὶ Ἰσραὴλ ἐνώπιον τῶν ἐθνῶν, ὅτι αὐτὸς διέσπειρεν ἡμᾶς ἐν αὐτοῖς.
\vs{4}Ἐκεῖ ὑποδείξατε τὴν μεγαλωσύνην αὐτοῦ, ὑψοῦτε αὐτὸν ἐνώπιον παντὸς ζῶντος, καθότι αὐτὸς Κύριος ἡμῶν, καὶ Θεὸς αὐτὸς πατὴρ ἡμῶν εἰς πάντας τοὺς αἰῶνας.
\vs{5}Καὶ μαστιγώσει ἡμᾶς ἐν ταῖς ἀδικίαις ἡμῶν, καὶ πάλιν ἐλεήσει, καὶ συνάξει ἡμᾶς ἐκ πάντων τῶν ἐθνῶν, οὗ ἐὰν σκορπισθῆτε ἐν αὐτοῖς.

\vs{6}Ἐὰν ἐπιστρέψητε πρὸς αὐτὸν ἐν ὅλῃ τῇ καρδίᾳ ὑμῶν, καὶ ἐν ὅλῃ τῇ ψυχῇ ὑμῶν, ποιῆσαι ἐνώπιον αὐτοῦ ἀλήθειαν, τότε ἐπιστρέψει πρὸς ὑμᾶς, καὶ οὐ μὴ κρύψει τὸ πρόσωπον αὐτοῦ ἀφʼ ὑμῶν· καὶ θεάσασθε ἃ ποιήσει μεθʼ ὑμῶν, καὶ ἐξομολογήσασθε αὐτῷ ἐν ὅλῳ τῷ στόματι ὑμῶν, καὶ εὐλογήσατε τὸν Κύριον τῆς δικαιοσύνης, καὶ ὑψώσατε τὸν βασιλέα τῶν αἰώνων· ἐγὼ ἐν τῇ γῇ τῆς αἰχμαλωσίας μου ἐξομολογοῦμαι αὐτῷ, καὶ δεικνύω τὴν ἰσχὺν καὶ τὴν μεγαλωσύνην αὐτοῦ ἔθνει ἁμαρτωλῶν· ἐπιστρέψατε ἁμαρτωλοὶ, καὶ ποιήσατε δικαιοσύνην ἐνώπιον αὐτοῦ· τίς γινώσκει εἰ θελήσει ὑμᾶς, καὶ ποιήσει ἐλεημοσύνην ὑμῖν;

\vs{7}Τὸν Θεόν μου ὑψῶ, καὶ ἡ φυχή μου τῷ βασιλεῖ τοῦ οὐρανοῦ, καὶ ἀγαλλιάσεται τὴν μεγαλωσύνην αὐτοῦ.
\vs{7a}Λεγέτωσαν πάντες, καὶ ἐξομολογείσθωσαν αὐτῷ ἐν Ἱεροσολύμοις.

\vs{7b}Ἱεροσόλυμα πόλις ἁγίου, μαστιγώσει ἐπὶ τὰ ἔργα τῶν υἱῶν σου, καὶ πάλιν ἐλεήσει τοὺς υἱοὺς τῶν δικαίων.
\vs{7c}Ἐξομολογοῦ τῷ Κυρίῳ ἀγαθῶς, καὶ εὐλόγει τὸν βασιλέα τῶν αἰώνων, ἵνα πάλιν ἡ σκηνὴ αὐτοῦ οἰκοδομηθῇ ἐν σοὶ μετὰ χαρᾶς· καὶ εὐφράναι ἐν σοὶ τοὺς αἰχμαλώτους, καὶ ἀγαπήσαι ἐν σοὶ τους ταλαιπώρους, εἰς πάσας τὰς γενεὰς τοῦ αἰῶνος.

\vs{11}Ἔθνη πολλὰ μακρόθεν ἥξει πρὸς τὸ ὄνομα Κυρίου τοῦ Θεοῦ, δῶρα ἐν χερσὶν ἔχοντες, καὶ δῶρα τῷ βασιλεῖ τοῦ οὐρανοῦ· γενεαὶ γενεῶν δώσουσί σοι ἀγαλλίαμα.
\vs{12}Ἐπικατάρατοι πάντες οἱ μισοῦντές σε, εὐλογημένοι ἔσονται πάντες οἱ ἀγαπῶντές σε εἰς τὸν αἰῶνα.
\vs{13}Χάρηθι καὶ ἀγαλλίασαι ἐπὶ τοῖς υἱοῖς τῶν δικαίων, ὅτι συναχθήσονται καὶ εὐλογήσουσι τὸν Κύριον τῶν δικαίων.
\vs{14}Ὦ μακάριοι οἱ ἀγαπῶντές σε, χαρήσονται ἐπὶ τῇ εἰρήνῃ σου· μακάριοι ὅσοι ἐλυπήθησαν ἐπὶ πάσαις ταῖς μάστιξί σου, ὅτι ἐπὶ σοὶ χαρήσονται θεασάμενοι πᾶσαν τὴν δόξαν σου, καὶ εὐφρανθήσονται εἰς τὸν αἰῶνα.

\vs{15}Ἡ ψυχή μου εὐλογεὶτω τὸν Θεὸν τὸν βασιλέα τὸν μέγαν,
\vs{16}ὅτι οἰκοδομηθήσεται Ἱερουσαλὴμ σαπφείρῳ καὶ σμαράγδῳ, καὶ λίθῳ ἐντίμῳ τὰ τείχη σου, καὶ οἱ πύργοι, καὶ οἱ προμαχῶνες ἐν χρυσίῳ καθαρῷ,
\vs{17}καὶ αἱ πλατεῖαι Ἱερουσαλὴμ ἐν βηρύλλῳ, καὶ ἄνθρακι, καὶ λίθῳ ἐκ Σουφεὶρ ψηφολογηθήσονται.
\vs{18}Καὶ ἐροῦσι πᾶσαι αἱ ῥύμαι αὐτῆς ἀλληλούϊα καὶ αἴνεσιν, λέγοντες, εὐλογητὸς ὁ Θεὸς, ὃς ὕψωσε πάντας τοὺς αἰῶνας.

\ch{14}
Καὶ ἐπαύσατο ἐξομολογούμενος Τωβίτ. Καὶ ἦν ἐτῶν πεντηκονταοκτὼ,
\vs{2}ὅτε ἀπώλεσε τὰς ὄψεις, καὶ μετὰ ἔτη ὀκτὼ ἀνέβλεψε· καὶ ἐποίει ἐλεημοσύνας· καὶ προσέθετο φοβεῖσθαι Κύριον τὸν Θεὸν, καὶ ἐξωμολογεῖτο αὐτῷ.

\vs{3}Μεγάλως δὲ ἐγήρασε· καὶ ἐκάλεσε τὸν υἱὸν αὐτοῦ, καὶ τοὺς υἱοὺς αὐτοῦ, καὶ εἶπεν αὐτῷ, τέκνον, λάβε τοὺς υἱούς σου, ἰδοὺ γεγήρακα, καὶ πρὸς τὸ ἀποτρέχειν ἐκ τοῦ ζῇν εἰμι.
\vs{4}Ἄπελθε εἰς τὴν Μηδίαν, τέκνον, ὅτι πέπεισμαι ὅσα ἐλάλησεν Ἰωνὰς ὁ προφήτης περὶ Νινευὴ, ὅτι καταστραφήσεται· ἐν δὲ τῇ Μηδίᾳ ἔσται εἰρήνη μᾶλλον ἕως καιροῦ· καὶ ὅτι οἱ ἀδελφοὶ ἡμῶν ἐν τῇ γῇ σκορπισθήσονται ἀπὸ τῆς ἀγαθῆς γῆς· καὶ Ἱεροσόλυμα ἔσται ἔρημος, καὶ ὁ οἶκος τοῦ Θεοῦ ἐν αὐτῇ κατακαήσεται, καὶ ἔρημος ἔσται μέχρι χρόνου.
\vs{5}Καὶ πάλιν ἐλεήσει αὐτοὺς ὁ Θεὸς, καὶ ἐπιστρέψει αὐτοὺς εἰς τὴν γῆν, καὶ οἰκοδομήσουσι τὸν οἶκον, οὐχ οἷος ὁ πρότερος, ἕως πληρωθῶσι καιροὶ τοῦ αἰῶνος· καὶ μετὰ ταῦτα ἐπιστρέψουσιν ἐκ τῶν αἰχμαλωσιῶν, καὶ οἰκοδομήσουσιν Ἱερουσαλὴμ ἐντίμως· καὶ ὁ οἶκος τοῦ Θεοῦ ἐν αὐτῇ οἰκοδομηθήσεται ἐνδόξως, καθὼς ἐλάλησαν περὶ αὐτῆς οἱ προφῆται.

\vs{6}Καὶ πάντα τὰ ἔθνη ἐπιστρέψουσιν ἀληθινῶς φοβεῖσθαι Κύριον τὸν Θεὸν, καὶ κατορύξουσι τὰ εἴδωλα αὐτῶν,
\vs{7}καὶ εὐλογήσουσι πάντα τὰ ἔθνη Κύριον· καὶ ὁ λαὸς αὐτοῦ ἐξομολογήσεται τῷ Θεῷ· καὶ ὑψώσει Κύριος τὸν λαὸν αὐτοῦ, καὶ χαρήσονται πάντες οἱ ἀγαπῶντες Κύριον τὸν Θεὸν ἐν ἀληθείᾳ καὶ δικαιοσύνῃ, ποιοῦντες ἔλεος τοῖς ἀδελφοῖς ἡμῶν.

\vs{8}Καὶ νῦν, τέκνον, ἄπελθε ἀπὸ Νινευὴ, ὅτι πάντως ἔσται ἃ ἐλάλησεν ὁ προφήτης Ἰωνάς.
\vs{9}Σὺ δὲ τήρησον τὸν νόμον καὶ τὰ προστάγματα, καὶ γενοῦ φιλελεήμων καὶ δίκαιος, ἵνα σοι καλῶς ᾖ.
\vs{10}Καὶ θάψον με καλῶς, καὶ τὴν μητέρα σου μετʼ ἐμοῦ, καὶ μηκέτι αὐλισθῆτε εἰς Νινευή· τεκνον, ἴδε τί ἐποίησεν Ἀμὰν Ἀχιαχάρῳ τῷ θρέψαντι αὐτόν, ὡς ἐκ τοῦ φωτὸς ἤγαγεν αὐτὸν εἰς τὸ σκότος, καὶ ὅσα ἀνταπέδωκεν αὐτῷ· καὶ Ἀχιάχαρον μὲν ἔσωσεν, ἐκείνῳ δὲ τὸ ἀνταπόδομα ἀπεδόθη, καὶ αὐτὸς κατέβη εἰς τὸ σκότος. Μανασσῆς ἐποίησεν ἐλεημοσύνην, καὶ ἐσώθη ἐκ παγίδος θανάτου ἧς ἔπηξεν αὐτῷ· Ἀμὰν δὲ ἐνέπεσεν εἰς τὴν παγίδα, καὶ ἀπώλετο.

\vs{11}Καὶ νῦν, παιδία, ἴδετε τί ἐλεημοσύνη ποιεῖ, καὶ δικαιοσύνη ῥύεται· καὶ ταῦτα αὐτοῦ λέγοντος, ἐξέλιπεν ἡ ψυχὴ αὐτοῦ ἐπὶ τῆς κλίνης· ἦν δὲ ἐτῶν ἑκατὸν πεντηκονταοκτώ· καὶ ἔθαψαν αὐτὸν ἐνδόξως.
\vs{12}Καὶ ὅτε ἀπέθανεν Ἄννα, ἔθαψεν αὐτὴν μετὰ τοῦ πατρὸς αὐτοῦ.

Ἀπῆλθε δὲ Τωβίας μετὰ τῆς γυναικὸς αὐτοῦ καὶ τῶν υἱῶν αὐτοῦ εἰς Ἐκβάτανα πρὸς Ῥαγουὴλ τὸν πενθερὸν αὐτοῦ,
\vs{13}καὶ ἐγήρασεν ἐντίμως· καὶ ἔθαψε τοὺς πενθεροὺς αὐτοῦ ἐνδόξως, καὶ ἐκληρονόμησε τὴν οὐσίαν, καὶ Τωβὶτ τοῦ πατρὸς αὐτοῦ.
\vs{14}καὶ ἀπέθανεν ἐτῶν ἑκατὸν εἱκοσιεπτὰ ἐν Ἐκβατάνοις τῆς Μηδίας.
\vs{15}Καὶ ἤκουσε πρὶνὴ ἀποθανεῖν αὐτὸν, τὴν ἀπώλειαν Νινευὴ, ἣν ᾐχμαλώτισεν Ναβουχοδονόσορ, καὶ Ἀσύηρος, καὶ ἐχάρη πρὸ τοῦ ἀποθανεῖν ἐπὶ Νινευή.


\def\book{ΙΟΥΔΙΘ}
\biblebook{ΙΟΥΔΙΘ}


\lettrine[lines=2, loversize=0.2, nindent=0em, findent=.25em]{\textcolor{bookheadingcolor}{Ἕ}}{ΤΟΥΣ} δωδεκάτου τῆς βασιλείας Ναβουχοδονόσορ, ὃς ἐβασίλευσεν Ἀσσυρίων ἐν Νινευὴ τῇ πόλει τῇ μεγάλῃ, ἐν ταῖς ἡμέραις Ἀρφαξὰδ,
\vs{2}ὃς ἐβασίλευσε Μήδων ἐν Ἐκβατάνοις, καὶ ᾠκοδόμησεν ἐπʼ Ἐκβατάνων, καὶ κύκλῳ τείχη ἐκ λίθων λελαξευμένων, εἰς πλάτος πηχῶν τριῶν, καὶ εἰς μῆκος πηχῶν ἓξ, καὶ ἐποίησε τὸ ὕψος τοῦ τείχους πηχῶν ἑβδομήκοντα, καὶ τὸ πλάτος αὐτοῦ πηχῶν πεντήκοντα,
\vs{3}καὶ τοὺς πύργους αὐτοῦ ἔστησεν ἐπὶ ταῖς πύλαις αὐτῆς πηχῶν ἑκατὸν, καὶ τὸ πλάτος αὐτῆς ἐθεμελίωσεν εἰς πήχεις ἑξήκοντα.
\vs{4}Καὶ ἐποίησε τὰς πύλας αὐτῆς πύλας διεγειρομένας εἰς ὕψος πηχῶν ἑβδομήκοντα, καὶ τὸ πλάτος αὐτῶν πήχεις τεσσράκοντα εἰς ἐξόδους δυνάμεων δυνατῶν αὐτοῦ, καὶ διατάξεις τῶν πεζῶν αὐτοῦ.

\vs{5}Καὶ ἐποίησε πόλεμον ἐν ταῖς ἡμέραις ἐκείναις ὁ βασιλεὺς Ναβουχοδονόσορ πρὸς βασιλέα Ἀρφαξὰδ ἐν τῷ πεδίῳ τῷ μεγάλῳ, τοῦτό ἐστιν ἐν τοῖς ὁρίοις Ῥαγαῦ.
\vs{6}Καὶ συνήντησαν πρὸς αὐτὸν πάντες οἱ κατοικοῦντες τὴν ὀρεινὴν, καὶ πάντες οἱ κατοικοῦντες τὸν Εὐφράτην, καὶ τὸν Τίγριν, καὶ τὸν Ὑδάσπην, καὶ πεδίῳ Εἰριὼχ ὁ βασιλεὺς Ἐλυμαίων· καὶ συνῆλθον ἔθνη πολλὰ σφόδρα εἰς παράταξιν υἱῶν Χελεούλ.

\vs{7}Καὶ ἀπέστειλε Ναβουχοδονόσορ ὁ βασιλεὺς Ἀσσυρίων ἐπὶ πάντας τοὺς κατοικοῦντας τὴν Περσίδα, καὶ ἐπὶ πάντας τοὺς κατοικοῦντας πρὸς δυσμαῖς, τοὺς κατοικοῦντας Κιλικίαν καὶ Δαμασκὸν, τὸν Λίβανον καὶ Ἀντιλίβανον, καὶ πάντας τοὺς κατοικοῦντας κατὰ πρόσωπον παραλίας,
\vs{8}καὶ τοὺς ἐν τοῖς ἔθνεσι τοῦ Καρμήλου, καὶ Γαλαὰδ, καὶ τὴν ἄνω Γαλιλαίαν, καὶ τὸ μέγα πεδίον Ἐσδρηλὼμ,
\vs{9}καὶ πάντας τοὺς ἐν Σαμαρείᾳ καὶ ταῖς πόλεσιν αὐτῆς, καὶ πέραν τοῦ Ἰορδάνου ἕως Ἱερουσαλὴμ, καὶ Βετὰνη, καὶ Χελλοὺς, καὶ Κάδης, καὶ τοῦ ποταμοῦ Αἰγύπτου, καὶ Ταφνὰς, καὶ Ῥαμεσσὴ,
\vs{10}καὶ πᾶσαν γῆν Γεσὲμ ἕως τοῦ ἐλθεῖν ἐπάνω Τάνεως καὶ Μέμφεως, καὶ πάντας τοὺς κατοικοῦντας τὴν Αἴγυπτον ἕως τοῦ ἐλθεῖν ἐπὶ τὰ ὅρια τῆς Αἰθιοπίας.

\vs{11}Καὶ ἐφαύλισαν πάντες οἱ κατοικοῦντες πᾶσαν τὴν γῆν τὸ ῥῆμα Ναβουχοδονὸσορ τοῦ βασιλέως Ἀσσυρίων, καὶ οὐ συνῆλθον αὐτῷ εἰς τὸν πόλεμον, ὅτι οὐκ ἐφοβήθησαν αὐτὸν, ἀλλ ἦν ἐναντίον αὐτῶν ὡς ἀνὴρ ἶσος· καὶ ἀνέστρεψαν τοὺς ἀγγέλους αὐτοῦ κενοὺς ἐν ἀτιμίᾳ πρὸ προσώπου αὐτῶν.
\vs{12}Καὶ ἐθυμώθη Ναβουχοδονόσορ ἐπὶ πᾶσαν τὴν γῆν ταύτην σφόδρα, καὶ ὤμοσε κατὰ τοῦ θρόνου καὶ τῆς βασιλείας αὐτοῦ, εἰ μὴν ἐκδικήσειν πάντα τὰ ὅρια τῆς Κιλικίας καὶ Δαμασκηνῆς καὶ Συρίας, ἀνελεῖν τῇ ῥομφαίᾳ αὐτοῦ καὶ πάντας τοὺς κατοικοῦντας ἐν γῇ Μωὰβ, καὶ τοὺς υἱοὺς Ἀμμὼν, καὶ πᾶσαν τὴν Ἰουδαίαν, καὶ πάντας τοὺς ἐν Αἰγύπτῳ ἕως τοῦ ἐλθεῖν ἐπὶ τὰ ὅρια τῶν δύο θαλασσῶν.

\vs{13}Καὶ παρετάξατο ἐν τῇ δυνάμει αὐτοῦ πρὸς Ἀρφαξὰδ βασιλέα ἐν τῷ ἔτει τῷ ἑπτακαιδεκάτῳ, καὶ ἐκραταιώθη ἐν τῷ πολέμῳ αὐτοῦ, καὶ ἀνέστρεψε πᾶσαν τὴν δύναμιν Ἀρφαξὰδ, καὶ πᾶσαν τὴν ἵππον αὐτοῦ, καὶ πάντα τὰ ἅρματα αὐτοῦ,
\vs{14}καὶ ἐκυρίευσε τῶν πόλεων αὐτοῦ· καὶ ἀφίκετο ἕως Ἐκβατάνων, καὶ ἐκράτησε τῶν πύργων, καὶ ἐπρονόμευσε τὰς πλατείας αὐτῆς, καὶ τὸν κόσμον αὐτῆς ἔθηκεν εἰς ὄνειδος αὐτης·
\vs{15}Καὶ ἔλαβε τὸν Ἀρφαξὰδ ἐν τοῖς ὄρεσι Ῥαγαῦ, καὶ κατηκόντισεν αὐτὸν ἐν ταῖς ζιβύναις αὐτοῦ, καὶ ἐξωλόθρευσεν αὐτὸν ἕως τῆς ἡμέρας ἐκείνης.

\vs{16}Καὶ ἀνέστρεψε μετʼ αὐτῶν αὐτὸς καὶ πᾶς ὁ σύμμικτος αὐτοῦ, πλῆθος ἀνδρῶν πολεμιστῶν πολὺ σφόδρα· καὶ ἦν ἐκεῖ ῥαθυμῶν καὶ εὐωχούμενος αὐτὸς καὶ ἡ δυνάμις αὐτοῦ ἐφʼ ἡμέρας ἑκατὸν εἴκοσι.

\ch{2}
Καὶ ἐν τῷ ἔτει τῷ ὀκτωκαιδεκάτῳ, δευτέρᾳ καὶ εἰκάδι τοῦ πρώτου μηνὸς, ἐγένετο λόγος ἐν οἴκῳ Ναβουχοδονόσορ βασιλέως Ἀσσυρίων, ἐκδικῆσαι πᾶσαν τὴν γῆν καθὼς ἐλάλησε.
\vs{2}Καὶ συνεκάλεσε πάντας τοὺς θεράποντας αὐτοῦ, καὶ πάντας τοὺς μεγιστᾶνας αὐτοῦ, καὶ ἔθετο μετʼ αὐτῶν τὸ μυστήριον τῆς βουλῆς αὐτοῦ· καὶ συνετέλεσε πᾶσαν τὴν κακίαν τῆς γῆς ἐκ τοῦ στόματος αὐτοῦ·
\vs{3}καὶ αὐτοὶ ἔκριναν ὀλοθρεῦσαι πᾶσαν σάρκα, οἳ οὐκ ἠκολούθησαν τῷ λόγῳ τοῦ στόματος αὐτοῦ.

\vs{4}Καὶ ἐγένετο ὡς συνετέλεσε τὴν βουλὴν αὐτοῦ, ἐκάλεσε Ναβουχοδονόσορ βασιλεὺς Ἀσσυρίων τὸν Ὀλοφέρνην ἀρχιστράτηγον τῆς δυνάμεως αὐτοῦ, δεύτερον ὄντα μετʼ αὐτὸν, καὶ εἶπε πρὸς αὐτὸν,
\vs{5}τάδε λέγει ὁ βασιλεὺς ὁ μέγας, ὁ κύριος πάσης τῆς γῆς, ἰδοὺ σὺ ἐξελεύσῃ ἐκ τοῦ προσώπου μου, καὶ λήψῃ μετὰ σεαυτοῦ ἄνδρας πεποιθότας ἐν ἰσχύϊ αὐτῶν, πεζῶν εἰς χιλιάδας ἑκατὸν εἴκοσι, καὶ πλῆθος ἵππων σὺν ἀναβάταις μυριάδων δεκαδύο,
\vs{6}καὶ ἐξελεύσῃ εἰς συνάντησιν πάσῃ τῇ γῇ ἐπὶ δυσμὰς, ὅτι ἠπείθησαν τῷ ῥήματι τοῦ στόματός μου·
\vs{7}καὶ ἀπαγγελεῖς αὐτοῖς ἑτοιμάζειν γῆν καὶ ὕδωρ, ὅτι ἐξελεύσομαι ἐν θυμῷ μου ἐπʼ αὐτοὺς, καὶ καλύψω πᾶν τὸ πρόσωπον τῆς γῆς ἐν τοῖς ποσὶ τῆς δυνάμεώς μου· καὶ δώσω αὐτοὺς εἰς διαρπαγὴν αὐτοῖς,
\vs{8}καὶ οἱ τραυματίαι αὐτῶν πληρώσουσι τὰς φάραγγας καὶ τοὺς χειμάῤῥους αὐτῶν, καὶ ποταμὸς ἐπικλύζων τοῖς νεκροῖς αὐτῶν πληρωθήσεται·
\vs{9}καὶ ἄξω τὴν αἰχμαλωσίαν αὐτῶν ἐπὶ τὰ ἄκρα πάσης τῆς γῆς.
\vs{10}Σὺ δὲ ἐξελθὼν προκαταλήψῃ μοι πᾶν ὅριον αὐτῶν, καὶ ἐκδώσουσί σοι ἑαυτοὺς, καὶ διατηρήσεις ἐμοὶ αὐτοὺς εἰς ἡμέραν ἐλεγμοῦ αὐτῶν.

\vs{11}Ἐπὶ δὲ τοὺς ἀπειθοῦντας οὐ φείσεται ὁ ὀφθαλμός σου, δοῦναι αὐτοὺς εἰς φόνον καὶ ἁρπαγὴν ἐν πάσῃ τῇ γῇ σου.
\vs{12}Ὅτι ζῶν ἐγὼ, καὶ τὸ κράτος τῆς βασιλείας μου, λελάληκα, καὶ ποιήσω ταῦτα ἐν χειρί μου.
\vs{13}Καὶ σὺ δὲ οὐ παραβήσῃ ἕν τι τῶν ῥημάτων τοῦ κυρίου σου, ἀλλʼ ἐπιτελῶν ἐπιτελέσεις, καθότι προστέταχά σοι, καὶ οὐ μακρυνεῖς τοῦ ποιῆσαι αὐτά.

\vs{14}Καὶ ἐξῆλθεν Ὀλοφέρνης ἀπὸ προσώπου τοῦ κυρίου αὐτοῦ, καὶ ἐκάλεσε πάντας τοὺς δυνάστας, καὶ τοὺς στρατηγοὺς, καὶ ἐπιστάτας τῆς δυνάμεως Ἀσσοὺρ,
\vs{15}καὶ ἠρίθμησεν ἐκλεκτοὺς ἄνδρας εἰς παράταξιν, καθότι ἐκέλευσεν αὐτῷ ὁ κύριος αὐτοῦ, εἰς μυριάδας δεκαδύο, καὶ ἱππεῖς τοξότας μυρίους δισχιλίους,
\vs{16}καὶ διέταξεν αὐτοὺς ὃν τρόπον πολέμου πλῆθος συντάσσεται.
\vs{17}Καὶ ἔλαβε καμήλους καὶ ὄνους καὶ ἡμιόνους εἰς τὴν ἀπαρτίαν αὐτῶν, πλῆθος πολὺ σφόδρα, καὶ πρόβατα καὶ βόας καὶ αἶγας εἰς τὴν παρασκευὴν αὐτῶν, ὧν οὐκ ἦν ἀριθμός,
\vs{18}καὶ ἐπισιτισμὸν παντὶ ἀνδρὶ εἰς πλῆθος, καὶ χρυσίον καὶ ἀργύριον ἐξ οἴκου βασιλέως πολὺ σφόδρα.

\vs{19}Καὶ ἐξῆλθεν αὐτὸς καὶ πᾶσα ἡ δύναμις αὐτοῦ εἰς πορείαν τοῦ προελθεῖν βασιλέως Ναβουχοδονόσορ, καὶ καλύψαι πᾶν τὸ πρόσωπον τῆς γῆς πρὸς δυσμαῖς ἐν ἅρμασι καὶ ἱππεῦσι καὶ πεζοῖς ἐπιλέκτοις αὐτῶν.
\vs{20}Καὶ πολὺς ὁ ἐπίμικτος ὡς ἀκρὶς τυνεξῆλθον αὐτοῖς, καὶ ὡς ἡ ἄμμος τῆς γῇς· οὐ γὰρ ἦν ἀριθμὸς ἀπὸ πλήθους αὐτῶν.

\vs{21}Καὶ ἀπῆλθον ἐκ Νινευὴ ὁδὸν τριῶν ἡμερῶν ἐπὶ πρόσωπον τοῦ πεδίου Βαικτιλαὶθ, καὶ ἐπεστρατοπέδευσεν ἀπὸ Βαικτιλαὶθ πλησίον τοῦ ὄρους τοῦ ἐπʼ ἀριστερᾷ τῆς ἄνω Κιλικίας.
\vs{22}Καὶ ἔλαβε πᾶσαν τὴν δύναμιν αὐτοῦ, τοὺς πεζοὺς, καὶ τοὺς ἱππεῖς, καὶ τὰ ἅρματα αὐτοῦ, καὶ ἀπῆλθεν ἐκεῖθεν εἰς τὴν ὀρεινήν.
\vs{23}Καὶ διέκοψε τὸ Φοὺδ καὶ Λοὺδ, καὶ ἐπρονόμευσαν πάντας υἱοὺς Ῥασσὶς, καὶ υἱοὺς Ἰσμαὴλ τοὺς κατὰ πρόσωπον τῆς ἐρήμου πρὸς Νότον τῆς Χελλαίων.

\vs{24}Καὶ παρῆλθε τὸν Εὐφράτην, καὶ διῆλθε τὴν Μεσοποταμίαν, καὶ διέσκαψε πάσας τὰς πόλεις τὰς ὑψηλὰς τὰς ἐπὶ τοῦ χειμαῤῥοῦ Ἀβρωνᾶ ἕως τοῦ ἐλθεῖν ἐπὶ θάλασσαν.

\vs{25}Καὶ κατελάβετο τὰ ὅρια τῆς Κιλίκιας, καὶ κατέκοψε πάντας τοὺς ἀντιστάντας αὐτῷ· καὶ ἦλθεν ἕως ὁρίων Ἰάφεθ, τὰ πρὸς Νότον κατὰ πρόσωπον τῆς Ἀραβίας.

\vs{26}Καὶ ἐκύκλωσε πάντας τοὺς υἱοὺς Μαδιὰν, καὶ ἐνέπρησε τὰ σκηνώματα αὐτῶν, καὶ ἐπρονόμευσε τὰς μάνδρας αὐτῶν.

\vs{27}Καὶ κατέβη εἰς πεδίον Δαμασκοῦ ἐν ἡμέραις θερισμοῦ πυρῶν, καὶ ἐνέπρησε πάντας τοὺς ἀγροὺς αὐτῶν· καὶ τὰ ποίμνια καὶ τὰ βουκόλια ἔδωκεν εἰς ἀφανισμὸν, καὶ τὰς πόλεις αὐτῶν ἐσκύλευσε, καὶ τὰ πεδία αὐτῶν ἐξελίκμησε, καὶ ἐπάταξε πάντας τοὺς νεανίσκους αὐτῶν ἐν στόματι ῥομφαίας.

\vs{28}Καὶ ἐπέπεσεν ὁ φόβος καὶ ὁ τρόμος αὐτοῦ ἐπὶ τοὺς κατοικοῦντας τὴν παραλίαν, τοὺς ὄντας ἐν Σιδῶνι καὶ Τύρῳ, καὶ τοὺς κατοικοῦντας Σοὺρ, καὶ Ὀκινὰ, καὶ πάντας τοὺς κατοικοῦντας Ἱεμναάν· καὶ οἱ κατοικοῦντες ἐν Ἀζώτῳ καὶ Ἀσκάλωνι ἐφοβήθησαν αὐτὸν σφόδρα.

\ch{3}
Καὶ ἀπέστειλαν πρὸς αὐτὸν ἀγγέλους λόγοις εἰρηνικοῖς, λέγοντες,
\vs{2}ἰδοὺ ἡμεῖς οἱ παῖδες Ναβουχοδονόσορ βασιλέως μεγάλου παρακείμεθα ἐνώπιόν σου, χρῆσαι ἡμῖν καθὼς ἀρεστόν ἐστι τῷ προσώπῳ σου.
\vs{3}Ἰδοὺ αἱ ἐπαύλεις ἡμῶν, καὶ πᾶν πεδίον πυρῶν, καὶ τὰ ποίμνια καὶ τὰ βουκόλια, καὶ πᾶσαι αἱ μάνδραι τῶν σκηνῶν ἡμῶν παράκεινται πρὸ προσώπου σου· χρῆσαι καθʼ ὃ ἂν ἀρέσκῃ σοι.
\vs{4}Ἰδοὺ καὶ αἱ πόλεις ἡμῶν, καὶ οἱ κατοικοῦντες ἐν αὐταῖς δοῦλοί σου εἰσίν· ἐλθὼν ἀπάντησον αὐταῖς ὡς ἔστιν ἀγαθὸν ἐν ὀφθαλμοῖς σου.

\vs{5}Καὶ παρεγένοντο οἱ ἄνδρες πρὸς Ὀλοφέρνην, καὶ ἀπήγγειλαν αὐτῷ κατὰ τὰ ῥήματα ταῦτα.
\vs{6}Καὶ κατέβη ἐπὶ τὴν παραλίαν αὐτὸς καὶ ἡ δύναμις αὐτοῦ, καὶ ἐφρούρησε τὰς πόλεις τὰς ὑψηλάς· καὶ ἔλαβεν ἐξ αὐτῶν εἰς συμμαχίαν ἄνδρας ἐπιλέκτους.
\vs{7}Καὶ ἐδέξαντο αὐτὸν αὐτοὶ, καὶ πᾶσα ἡ περίχωρος αὐτῶν μετὰ στεφάνων καὶ χορῶν καὶ τυμπάνων.
\vs{8}Καὶ κατέσκαψε πάντα τὰ ὅρια αὐτῶν, καὶ τὰ ἄλση αὐτῶν ἐξέκοψε· καὶ ἦν δεδογμένον αὐτῷ ἐξολοθρεῦσαι πάντας τοὺς θεοὺς τῆς γῆς, ὅπως αὐτῷ μόνῳ τῷ Ναβουχοδονόσορ λατρεύσωσι πάντα τὰ ἔθνη, καὶ πᾶσαι αἱ γλῶσσαι καὶ πᾶσαι αἱ φυλαὶ αὐτῶν ἐπικαλέσωνται αὐτὸν εἰς θεόν.

\vs{9}Καὶ ἦλθε κατὰ πρόσωπον Ἐσδρηλὼν πλησίον τῆς Δωταίας, ἥ ἐστιν ἀπέναντι τοῦ πρίονος τοῦ μεγάλου τῆς Ἰουδαίας.
\vs{10}Καὶ κατεστρατοπέδευσεν ἀναμέσον Γαιβαὶ καὶ Σκυθῶν πόλεως, καὶ ἦν ἐκεῖ μῆνα ἡμερῶν εἰς τὸ συλλέξαι πᾶσαν τὴν ἀπαρτίαν τῆς δυνάμεως αὐτοῦ.

\ch{4}
Καὶ ἤκουσαν οἱ υἱοὶ Ἰσραὴλ οἱ κατοικοῦντες ἐν τῇ Ἰουδαίᾳ πάντα ὅσα ἐποίησεν Ὀλοφέρνης τοῖς ἔθνεσιν, ὁ ἀρχιστράτηγος Ναβουχοδονόσορ βασιλέως Ἀσσυρίων, καὶ ὃν τρόπον ἐσκύλευσε πάντα τὰ ἱερὰ αὐτῶν, καὶ ἔδωκεν αὐτὰ εἰς ἀφανισμὸν,
\vs{2}καὶ ἐφοβήθησαν σφόδρα σφόδρα ἀπὸ προσώπου αὐτοῦ, καὶ περὶ Ἱερουσαλὴμ καὶ τοῦ ναοῦ Κυρίου Θεοῦ αὐτῶν ἐταράχθησαν·
\vs{3}ὅτι προσφάτως ἦσαν ἀναβεβηκότες ἐκ τῆς αἰχμαλωσίας, καὶ νεωστὶ πᾶς ὁ λαὸς συνελέλεκτο τῆς Ἰουδαίας, καὶ τὰ σκεύη, καὶ τὸ θυσιαστήριον, καὶ ὁ οἶκος, ἐκ τῆς βεβηλώσεως ἡγιασμένα ἦν,
\vs{4}καὶ ἀπέστειλαν εἰς πᾶν ὅριον Σαμαρείας, καὶ Κωνὰς, καὶ Βαιθωρὼν, καὶ Βελμὲν, καὶ Ἱεριχὼ, καὶ εἰς Χωβὰ, καὶ Αἰσωρὰ, καὶ τὸν αὐλῶνα Σαλὴμ,
\vs{5}καὶ προκατελάβοντο πάσας τὰς κορυφὰς τῶν ὀρέων τῶν ὑψηλῶν, καὶ ἐτειχίσαντο τὰς ἐν αὐτοῖς κώμας, καὶ παρέθεντο εἰς ἐπισιτισμον εἰς παρασκευὴν πολέμου, ὅτι προσφάτως ἦν τὰ πεδία αὐτῶν τεθερισμένα.

\vs{6}Καὶ ἔγραψεν Ἰωακὶμ ὁ ἱερεὺς ὁ μέγας, ὃς ἦν ἐν ταῖς ἡμέραις ἐκείναις ἐν Ἱερουσαλὴμ, τοῖς κατοικοῦσι Βετυλούα, καὶ Βετομεσθαὶμ, ἥ ἐστιν ἀπέναντι Ἐσδρηλὼν, κατὰ πρόσωπον τοῦ πεδίου τοῦ πλησίον Δωθαῒμ,
\vs{7}λέγων, διακατασχεῖν τὰς ἀναβάσεις τῆς ὀρεινῆς, ὅτι διʼ αὐτῶν ἦν ἡ εἴσοδος εἰς τὴν Ἰουδαίαν· καὶ ἦν εὐχερῶς διακωλύσαι αὐτοὺς προσβαίνοντας, στενῆς τῆς προσβάσεως οὔσης, ἐπʼ ἄνδρας τοὺς πάντας δύο.
\vs{8}Καὶ ἐποίησαν οἱ υἱοὶ Ἰσραὴλ καθὰ συνέταξεν αὐτοῖς Ἰωακὶμ ὁ ἱερεὺς ὁ μέγας, καὶ ἡ γερουσία παντὸς δήμου Ἰσραὴλ, οἳ ἐκάθηντο ἐν Ἱερουσαλήμ.

\vs{9}Καὶ ἀνεβόησαν πᾶς ἀνὴρ Ἰσραὴλ πρὸς τὸν Θεὸν ἐν ἐκτενίᾳ μεγάλῃ, καὶ ἐταπεινοῦσαν τὰς ψυχὰς αὐτῶν ἐν ἐκτενίᾳ μεγάλῃ,
\vs{10}αὐτοὶ καὶ αἱ γυναῖκες αὐτῶν, καὶ τὰ νήπια αὐτῶν, καὶ τὰ κτήνη αὐτῶν· καὶ πᾶς πάροικος ἢ μισθωτὸς, καὶ ἀργυρώνητος αὐτῶν, ἐπέθεντο σάκκους ἐπὶ τὰς ὀσφύας αὐτῶν.

\vs{11}Καὶ πᾶς ἀνὴρ Ἰσραὴλ καὶ γυνὴ, τὰ παιδία, καὶ οἱ κατοικοῦντες ἐν Ἱερουσαλὴμ ἔπεσον κατὰ πρόσωπον τοῦ ναοῦ, καὶ ἐσποδώσαντο τὰς κεφαλὰς αὐτῶν, καὶ ἐξέτειναν τοὺς σάκκους αὐτῶν κατὰ πρόσωπον Κυρίου· καὶ τὸ θυσιαστήριον σάκκῳ περιέβαλον,
\vs{12}καὶ ἐβόησαν πρὸς τὸν Θεὸν Ἰσραὴλ ὁμοθυμαδὸν ἐκτενῶς, τοῦ μὴ δοῦναι εἰς διαρπαγὴν τὰ νήπια αὐτῶν, καὶ τὰς γυναῖκας εἰς προνομὴν, καὶ τὰς πόλεις τῆς κληρονομίας αὐτῶν εἰς ἀφανισμὸν, καὶ τὰ ἅγια εἰς βεβήλωσιν καὶ ὀνειδισμὸν, ἐπίχαρμα τοῖς ἔθνεσι.

\vs{13}Καὶ εἰσήκουσε Κύριος τῆς φωνῆς αὐτῶν, καὶ εἰσεῖδε τὴν θλίψιν αὐτῶν· καὶ ἦν ὁ λαὸς νηστεύων ἡμέρας πλείους ἐν πάσῃ τῇ Ἰουδαίᾳ καὶ Ἱερουσαλὴμ, κατὰ πρόσωπον τῶν ἁγίων Κυρίου παντοκράτορος.

\vs{14}Καὶ Ἰωακὶμ ὁ ἱερεὺς ὁ μέγας, καὶ πάντες οἱ παρεστηκότες ἐνώπιον Κυρίου, ἱερεῖς καὶ οἱ λειτουργοῦντες Κυρίῳ, σάκκους περιεζωσμένοι τὰς ὀσφύας αὐτῶν, προσέφερον τὴν ὁλοκαύτωσιν τοῦ ἐνδελεχισμοῦ, καὶ τὰς εὐχὰς, καὶ τὰ ἑκουσία δόματα τοῦ λαοῦ,
\vs{15}καὶ ἦν σποδὸς ἐπὶ τὰς κιδάρεις αὐτῶν, καὶ ἐβόων πρὸς Κύριον ἐκ πάσης δυνάμεως εἰς ἀγαθὸν ἐπισκέψασθαι πάντα οἶκον Ἰσραήλ.

\ch{5}
Καὶ ἀνηγγέλλη Ὀλοφέρνῃ ἀρχιστρατήγῳ δυνάμεως Ἀσσοὺρ, διότι οἱ υἱοὶ Ἰσραὴλ παρεσκευάσαντο εἰς πόλεμον, καὶ τὰς διόδους τῆς ὀρεινῆς συνέκλεισαν, καὶ ἐτείχισαν πᾶσαν κορυφὴν ὄρους ὑψηλοῦ, καὶ ἔθηκαν ἐν τοῖς πεδίοις σκάνδαλα.
\vs{2}Καὶ ὠργίσθη θυμῷ σφόδρα, καὶ ἐκάλεσε πάντας τοὺς ἄρχοντας Μωὰβ, καὶ τοὺς στρατηγοὺς Ἀμμὼν, καὶ πάντας σατράπας τῆς παραλίας,
\vs{3}καὶ εἶπεν αὐτοῖς, ἀναγγείλατε δή μοι, υἱοὶ Χαναὰν, τίς ὁ λαὸς οὗτος ὁ καθήμενος ἐν τῇ ὀρεινῇ, καὶ τίνες ἃς κατοικοῦσι πόλεις; καὶ τὸ πλῆθος τῆς δυνάμεως αὐτῶν, καὶ ἐν τίνι τὸ κράτος αὐτῶν, καὶ ἡ ἰσχὺς αὐτῶν, καὶ τίς ἀνέστηκεν ἐπʼ αὐτῶν βασιλεὺς ἡγούμενος στρατηγίας αὐτῶν;
\vs{4}Καὶ διὰ τί κατενωτίσαντο τοῦ μὴ ἐλθεῖν εἰς ἀπάντησίν μοι παρὰ πάντας τοὺς κατοικοῦντας ἐν δυσμαῖς;

\vs{5}Καὶ εἶπε πρὸς αὐτὸν Ἀχιὼρ ὁ ἡγούμενος πάντων υἱῶν Ἀμμὼν, ἀκουσάτω δὴ ὁ κύριός μου λόγον ἐκ στόματος τοῦ δούλου σου, καὶ ἀναγγελῶ σοι τὴν ἀλήθειαν περὶ τοῦ λαοῦ, ὃς κατοικεῖ τὴν ὀρεινὴν ταύτην, πλησίον σοι οἰκοῦντος, καὶ οὐκ ἐξελεύσεται ψεῦδος ἐκ τοῦ στόματος τοῦ δούλου σου.
\vs{6}Ὁ λαὸς οὗτός εἰσιν ἀπογόνοι Χαλδαίων,
\vs{7}καὶ παρῴκησαν τὸ πρότερον ἐν τῇ Μεσοποταμίᾳ, ὅτι οὐκ ἐβουλήθησαν ἀκολουθῆσαι τοῖς θεοῖς τῶν πατέρων αὐτῶν, οἳ ἐγένοντο ἐν γῇ Χαλδαίων·
\vs{8}καὶ ἐξέβησαν ἐξ ὁδοῦ τῶν γονέων αὐτῶν, καὶ προσεκύνησαν τῷ Θεῷ τοῦ οὐρανοῦ, Θεῷ ᾧ ἐπέγνωσαν· καὶ ἐξέβαλον αὐτοὺς ἀπὸ προσώπου τῶν θεῶν αὐτῶν, καὶ ἔφυγον εἰς Μεσοποταμίαν, καὶ παρῴκησαν ἐκεῖ ἡμέρας πολλάς.

\vs{9}Καὶ εἶπεν ὁ Θεὸς αὐτῶν ἐξελθεῖν ἐκ τῆς παροικίας αὐτῶν, καὶ πορευθῆναι εἰς γῆν Χαναάν· καὶ κατῴκησαν ἐκεῖ, καὶ ἐπληθύνθησαν χρυσίῳ καὶ ἀργυρίῳ καὶ ἐν κτήνεσι πολλοῖς σφόδρα.
\vs{10}Καὶ κατέβησαν εἰς Αἴγυπτον, ἐκάλυψε γὰρ τὸ πρόσωπον τῆς γῆς Χαναὰν λιμὸς, καὶ παρῴκησαν ἐκεῖ μέχρις οὗ διετράφησαν· καὶ ἐγένοντο ἐκεῖ εἰς πλῆθος πολύ, καὶ οὐκ ἦν ἀριθμὸς τοῦ γένους αὐτῶν.
\vs{11}Καὶ ἐπανέστη αὐτοῖς ὁ βασιλεὺς Αἰγύπτου, καὶ κατεσοφίσαντο αὐτοὺς ἐν πόνῳ καὶ ἐν πλίνθῳ, καὶ ἐταπείνωσαν αὐτοὺς, καὶ ἔθεντο αὐτοὺς εἰς δούλους.

\vs{12}Καὶ ἀνεβόησαν πρὸς τὸν Θεὸν αὐτῶν, καὶ ἐπάταξε πᾶσαν τὴν γῆν Αἰγύπτου πληγαῖς, ἐν αἷς οὐκ ἦν ἴασις· καὶ ἐξέβαλον αὐτοὺς οἱ Αἰγύπτιοι ἀπὸ προσώπου αὐτῶν.
\vs{13}Καὶ κατεξήρανεν ὁ Θεὸς τὴν ἐρυθρὰν θάλασσαν ἔμπροσθεν αὐτῶν,
\vs{14}καὶ ἤγαγεν αὐτοὺς εἰς ὁδὸν τοῦ Σινὰ, καὶ Κάδης Βαρνὴ, καὶ ἐξέβαλον πάντας τοὺς κατοικοῦντας ἐν τῇ ἐρήμῳ.

\vs{15}Καὶ ᾤκησαν ἐν γῇ Ἀμοῤῥαίων, καὶ πάντας τοὺς Εσεβωνίτας ἐξωλόθρευσαν ἐν τῇ ἰσχύϊ αὐτῶν· καὶ διαβάντες τὸν Ἰορδάνην ἐκληρονόμησαν πᾶσαν τὴν ὀρεινήν.
\vs{16}Καὶ ἐξέβαλον ἐκ προσώπου αὐτῶν τὸν Χανανῖον, καὶ τὸν Φερεζαῖον, καὶ τὸν Ἰεβουσαῖον. καὶ τὸν Συχὲμ, καὶ πάντας τοὺς Γεργεσαίους, καὶ κατῴκησαν ἐν αὐτῇ ἡμέρας πολλάς.

\vs{17}Καὶ ἕως οὐχ ἥμαρτον ἐνώπιον τοῦ Θεοῦ αὐτῶν, ἦν τὰ ἀγαθὰ μετʼ αὐτῶν, ὅτι Θεὸς μισῶν ἀδικίαν μετʼ αὐτῶν ἐστίν.
\vs{18}Ὅτε δὲ ἀπέστησαν ἀπὸ τῆς ὁδοῦ ἧς διέθετο αὐτοῖς, ἐξωλοθρεύθησαν ἐν πολλοῖς πολέμοις ἐπὶ πολὺ σφόδρα, καὶ ᾐχμαλωτεύθησαν εἰς γῆν οὐκ ἰδίαν, καὶ ὁ ναὸς τοῦ Θεοῦ αὐτῶν ἐγενήθη εἰς ἔδαφος, καὶ αἱ πόλεις αὐτῶν ἐκρατήθησαν ὑπὸ τῶν ὑπεναντίων.

\vs{19}Καὶ νῦν ἐπιστρέψαντες ἐπὶ τὸν Θεὸν αὐτῶν, ἀνέβησαν ἐκ τῆς διασπορᾶς οὗ διεσπάρησαν ἐκεῖ, καὶ κατέσχον τὴν Ἱερουσαλὴμ, οὗ τὸ ἁγίασμα αὐτῶν, καὶ κατῳκίσθησαν ἐν τῇ ὀρεινῇ, ὅτι ἦν ἔρημος,
\vs{20}καὶ νῦν, δέσποτα κύριε, εἰ μέν ἐστιν ἀγνόημα ἐν τῷ λαῷ τούτῳ, καὶ ἁμαρτάνουσιν εἰς τὸν Θεὸν αὐτῶν, καὶ ἐπισκεψόμεθα ὅ, τι ἐστὶν ἐν αὐτοῖς σκάνδαλον τοῦτο, καὶ ἀναβησόμεθα, καὶ ἐκπολεμήσομεν αὐτούς.
\vs{21}Εἰ δὲ οὐκ ἔστιν ἀνομία ἐν τῷ ἔθνει αὐτῶν, παρελθέτω δὴ ὁ κύριός μου, μήποτε ὑπερασπίσῃ ὁ Κύριος αὐτῶν καὶ ὁ Θεὸς αὐτῶν ὑπὲρ αὐτῶν, καὶ ἐσόμεθα εἰς ὀνειδισμὸν ἐναντίον πάσης τῆς γῆς.

\vs{22}Καὶ ἐγένετο ὡς ἐπαύσατο Ἀχιὼρ λαλῶν τοὺς λόγους τούτους, καὶ ἐγόγγυσε πᾶς ὁ λαὸς ὁ κυκλῶν τὴν σκηνὴν καὶ περιεστώς· καὶ εἶπαν οἱ μεγιστᾶνες Ὀλοφέρνου, καὶ πάντες οἱ κατοικοῦντες τὴν παραλίαν καὶ τὴν Μωὰβ, συγκόψαι αὐτόν,
\vs{23}οὐ γὰρ φοβηθησόμεθα ἀπὸ υἱῶν Ἰσραήλ· ἰδοὺ γὰρ λαὸς ἐν ᾧ οὐκ ἔστι δύναμις, οὐδὲ κράτος εἰς παράταξιν ἰσχυράν.

\vs{24}Διὸ δὴ ἀναβησόμεθα, καὶ ἔσονται εἰς κατάβρωμα πάσης τῆς στρατιᾶς σου, δέσποτα Ὀλοφέρνη.

\ch{6}
Καὶ ὡς κατέπαυσεν ὁ θόρυβος τῶν ἀνδρῶν τῶν κύκλῳ τῆς συνεδρείας, καὶ εἶπεν Ὀλοφέρνης ὁ ἀρχιστράτηγος δυνάμεως Ἀσσοὺρ πρὸς Ἀχιὼρ ἐναντίον παντὸς τοῦ δήμου ἀλλοφύλων, καὶ πρὸς πάντας υἱοὺς Μωὰβ,
\vs{2}καὶ τίς εἶ σὺ, Ἀχιὼρ, καὶ οἱ μισθωτοὶ τοῦ Ἐφραὶμ, ὅτι προεφήτευσας ἐν ἡμῖν καθὼς σήμερον, καὶ εἶπας τὸ γένος Ἰσραὴλ μὴ πολεμῆσαι, ὅτι ὁ Θεὸς αὐτῶν ὑπερασπιεῖ αὐτῶν; καὶ τίς ὁ Θεὸς εἰ μὴ Ναβουχοδονόσορ;
\vs{3}Οὗτος ἀποστελεῖ τὸ κράτος αὐτοῦ, καὶ ἐξολοθρεύσει αὐτοὺς ἀπὸ προσώπου τῆς γῆς, καὶ οὐ ῥύσεται αὐτοὺς ὁ Θεὸς αὐτῶν· ἀλλʼ ἡμεῖς οἱ δοῦλοι αὐτοῦ πατάξομεν αὐτοὺς ὡς ἄνθρωπον ἕνα, καὶ οὐχ ὑποστήσονται τὸ κράτος τῶν ἵππων ἡμῶν.
\vs{4}Κατακαύσομεν γὰρ αὐτοὺς ἐν αὐτοῖς, καὶ τὰ ὄρη αὐτῶν μεθυσθήσεται ἐν τῷ αἵματι αὐτῶν, καὶ τὰ πεδία αὐτῶν πληρωθήσεται νεκρῶν αὐτῶν· καὶ οὐκ ἀντιστήσεται τὸ ἴχνος τῶν ποδῶν αὐτῶν κατὰ πρόσωπον ἡμῶν, ἀλλὰ ἀπωλείᾳ ἀπολοῦνται, λάλει ὁ βασιλεὺς Ναβουχοδονόσορ ὁ κύριος πάσης τῆς γῆς· εἶπε γὰρ, οὐ ματαιωθήσεται τὰ ῥήματα τῶν λόγων αὐτοῦ.

\vs{5}Σὺ δὲ Ἀχιὼρ μισθωτὲ τοῦ Ἀμμὼν, ὃς ἐλάλησας τοὺς λόγους τούτους ἐν ἡμέρᾳ ἀδικίας σου, οὐκ ὄψει ἔτι τὸ πρόσωπόν μου ἀπὸ τῆς ἡμέρας ταύτης, ἕως οὗ ἐκδικήσω τὸ γένος τῶν ἐκ Αἰγύπτου.
\vs{6}Καὶ τότε διελεύσεται ὁ σίδηρος τῆς στρατιᾶς μου, καὶ ὁ λαὸς τῶν θεραπόντων μου τὰς πλευράς σου, καὶ πεσῇ ἐν τοῖς τραυματίαις αὐτῶν, ὅταν ἐπιστρέψω.
\vs{7}Καὶ ἀποκαταστήσουσί σε οἱ δοῦλοί μου εἰς τὴν ὀρεινὴν, καὶ θήσουσί σε ἐν μιᾷ τῶν πόλεων τῶν ἀναβάσεων,
\vs{8}καὶ οὐκ ἀπολῇ ἕως οὗ ἐξολοθρευθῇς μετʼ αὐτῶν.
\vs{9}Καὶ εἴπερ ἐλπίζεις τῇ καρδίᾳ σον ὅτι οὐ λημφθήσονται, μὴ συμπεσέτω σου τὸ πρόσωπον· ἐλάλησα, καὶ οὐδὲν διαπεσεῖται τῶν ῥημάτων μου.

\vs{10}Καὶ προσέταξεν Ὀλοφέρνης τοῖς δούλοις αὐτοῦ, οἳ ἦσαν παρεστηκότες ἐν τῇ σκηνῇ αὐτοῦ, συλλαβεῖν τὸν Ἀχιὼρ, καὶ ἀποκαταστῆσαι αὐτὸν εἰς Βετυλούα, καὶ πάραδοῦναὶ εἰς χεῖρας υἱῶν Ἰσραήλ.
\vs{11}Καὶ συνέλαβον αὐτὸν οἱ δοῦλοι αὐτοῦ, καὶ ἤγαγον αὐτὸν ἔξω τῆς παρεμβολῆς εἰς τὸ πεδίον, καὶ ἀπῇραν ἐκ μέσου τῆς πεδινῆς εἰς τὴν ὀρεινὴν, καὶ παρεγένοντο ἐπὶ τὰς πηγὰς αἳ ἦσαν ὑποκάτω Βετυλούα.
\vs{12}Καὶ ὡς εἶδαν αὐτοὺς οἱ ἄνδρες τῆς πόλεως ἐπὶ τὴν κορυφὴν τοῦ ὄρους, ἀνέλαβον τὰ ὅπλα αὐτῶν, καὶ ἀπῆλθον ἔξω τῆς πόλεως ἐπὶ τὴν κορυφὴν τοῦ ὄρους· καὶ πᾶς ἀνὴρ σφενδονητὴς διεκράτησαν τὴν ἀνάβασιν αὐτῶν, καὶ ἔβαλον ἐν λίθοις ἐπʼ αὐτούς.
\vs{13}Καὶ ὑποδύσαντες ὑποκάτω τοῦ ὄρους, ἔδησαν τὸν Ἀχιὼρ, καὶ ἀφῆκαν ἐῤῥιμμένον ὑπὸ τὴν ῥίζαν τοῦ ὄρους, καὶ ἀπῴχοντο πρὸς τὸν κύριον αὐτῶν.

\vs{14}Καταβάντες δὲ υἱοὶ Ἰσραὴλ ἐκ τῆς πόλεως αὐτῶν ἐπέστησαν αὐτῷ, καὶ λύσαντες αὐτὸν ἀπήγαγον εἰς τὴν Βετυλούα, καὶ κατέστησαν αὐτὸν ἐπὶ τοὺς ἄρχοντας τῆς πόλεως αὐτῶν,
\vs{15}οἳ ἦσαν ἐν ταῖς ἡμέραις ἐκείναις, Ὀζίας ὁ τοῦ Μιχὰ ἐκ τῆς φυλῆς Συμεὼν, καὶ Ἀβρὶς ὁ τοῦ Γοθονιὴλ, καὶ Χαρμὶς υἱὸς Μελχιήλ.

\vs{16}Καὶ συνεκάλεσαν πάντας τοὺς πρεσβυτέρους τῆς πόλεως· καὶ συνέδραμον πᾶς νεανίσκος αὐτῶν καὶ αἱ γυναῖκες εἰς τὴν ἐκκλησίαν· καὶ ἔστησαν τὸν Ἀχιὼρ ἐν μέσῳ παντὸς τοῦ λαοῦ αὐτῶν· καὶ ἐπηρώτησεν αὐτὸν Ὀζίας τὸ συμβεβηκός.
\vs{17}Καὶ ἀποκριθεὶς ἀπήγγειλεν αὐτοῖς τὰ ῥήματα τῆς συνεδρίας Ὀλοφέρνου, καὶ πάντα τὰ ῥήματα ὅσα ἐλάλησεν ἐν μέσῳ τῶν ἀρχόντων υἱῶν Ἀσσούρ, καὶ ὅσα ἐμεγαλοῤῥημόνησεν Ὀλοφέρνης εἰς τὸν οἶκον Ἰσραήλ.

\vs{18}Καὶ πεσόντες ὁ λαὸς προσεκύνησαν τῷ Θεῷ, καὶ ἐβόησαν λέγοντες,
\vs{19}κύριε ὁ Θεὸς τοῦ οὐρανοῦ, κάτιδε ἐπὶ τὰς ὑπερηφανείας αὐτῶν, καὶ ἐλέησον τὴν ταπείνωσιν τοῦ γένους ἡμῶν, καὶ ἐπίβλεψον ἐπὶ τὸ πρόσωπον τῶν ἡγιασμένων σοι ἐν τῇ ἡμέρᾳ ταύτῃ.

\vs{20}Καὶ παρεκάλεσαν τὸν Ἀχιὼρ, καὶ ἐπῄνεσαν αὐτὸν σφόδρα.
\vs{21}Καὶ παρέλαβεν αὐτὸν Ὀζίας ἐκ τῆς ἐκκλησίας εἰς οἶκον αὐτοῦ, καὶ ἐποίησε πότον τοῖς πρεσβυτέροις· καὶ ἐπεκαλέσαντο τὸν Θεὸν Ἰσραὴλ εἰς βοήθειαν ὅλην τὴν νύκτα ἐκείνην.

\ch{7}
Τῇ δʼ ἐπαύριον παρήγγειλεν Ὀλοφέρνης πᾶσῃ τῇ στρατιᾷ αὐτοῦ, καὶ παντὶ τῷ λαῷ αὐτοῦ, οἳ παρεγένοντο ἐπὶ τὴν συμμαχίαν αὐτοῦ, ἀναζευγνύειν ἐπὶ Βετυλούα, καὶ τὰς ἀναβάσεις τῆς ὀρεινῆς προκαταλαμβάνεσθαι, καὶ ποιεῖν πόλεμον πρὸς τοὺς υἱοὺς Ἰσραήλ.
\vs{2}Καὶ ἀνέζευξεν ἐν τῇ ἡμέρᾳ ἐκείνῃ πᾶς ἀνὴρ δυνατὸς αὐτῶν· καὶ ἡ δύναμις αὐτῶν ἀνδρῶν πολεμιστῶν, χιλιάδες ἀνδρῶν πεζῶν ἐκατὸν ἑβδομήκοντα, καὶ ἱππέων χιλιάδες δεκαδύο, χωρὶς τῆς ἀποσκευῆς, καὶ τῶν ἀνδρῶν οἳ ἦσαν πεζοὶ ἐν αὐτοῖς, πλῆθος πολὺ σφόδρα.
\vs{3}Καὶ παρενέβαλον ἐν τῷ αὐλῶνι πλησίον Βετυλούα ἐπὶ τῆς πηγῆς, καὶ παρέτειναν εἰς εὖρος ἐπὶ Δωθαῒμ καὶ ἕως Βελθὲμ, καὶ εἰς μῆκος ἀπὸ Βετυλούα ἕως Κυαμῶνος, ἥ ἐστιν ἀπέναντι Ἐσδρηλώμ.

\vs{4}Οἱ δὲ υἱοὶ Ἰσραὴλ, ὡς εἶδον αὐτῶν τὸ πλῆθος, ἐταράχθησαν σφόδρα· καὶ εἶπεν ἕκαστος πρὸς τὸν πλησίον αὐτοῦ, νῦν ἐκλείξουσιν οὗτοι τὸ πρόσωπον τῆς γῆς πάσης, καὶ οὔτε τὰ ὄρη τὰ ὑψηλὰ, οὔτε αἱ φάραγγες, οὔτε οἱ βουνοὶ ὑποστήσονται τὸ βάρος αὐτῶν.
\vs{5}Καὶ ἀναλαβόντες ἕκαστος τὰ σκεύη τὰ πολεμικὰ αὐτῶν, καὶ ἀνακαύσαντες πυρὰς ἐπὶ τοὺς πύργους αὐτῶν, ἔμενον φυλάσσοντες ὅλην τὴν νύκτα ἐκείνην.
\vs{6}Τῇ δὲ ἡμέρᾳ τῇ δευτέρᾳ ἐξήγαγεν Ὀλοφέρνης πᾶσαν τὴν ἵππον αὐτοῦ κατὰ πρόσωπον τῶν υἱῶν Ἰσραὴλ οἳ ἦσαν ἐν Βετυλούᾳ,
\vs{7}καὶ ἐπεσκέψατο τὰς ἀναβάσεις τῆς πόλεως αὐτῶν, καὶ τὰς πηγὰς τῶν ὑδάτων αὐτῶν ἐφώδευσε, καὶ προκατελάβετο αὐτὰς, καὶ ἐπέστησεν αὐταῖς παρεμβολὰς ἀνδρῶν πολεμιστῶν, καὶ αὐτὸς ἀνέζευξεν εἰς τὸν λαὸν αὐτοῦ.

\vs{8}Καὶ προσελθόντες αὐτῷ πάντες οἱ ἄρχοντες τῶν υἱῶν Ἡσαῦ, καὶ πάντες οἱ ἡγούμενοι τοῦ λαοῦ Μωὰβ, καὶ οἱ στρατηγοὶ τῆς παραλίας, εἶπαν,
\vs{9}ἀκουσάτω δὴ λόγον ὁ δεσπότης ἡμῶν, ἵνα μὴ γένηται θραῦσμα ἐν τῇ δυνάμει σου.
\vs{10}Ὁ γὰρ λαὸς οὗτος τῶν υἱῶν Ἰσραὴλ οὐ πέποιθαν ἐπὶ τοῖς δόρασιν αὐτῶν, ἀλλʼ ἐπὶ τοῖς ὕψεσι τῶν ὀρέων αὐτῶν, ἐν οἷς αὐτοὶ ἐνοικοῦσιν ἐν αὐτοῖς· οὐ γάρ ἐστιν εὐχερὲς προσβῆναι ταῖς κορυφαῖς τῶν ὀρέων αὐτῶν.

\vs{11}Καὶ νῦν, δέσποτα, μὴ πολέμει πρὸς αὐτοὺς, καθὼς γίνεται πόλεμος παρατάξεως, καὶ οὐ πεσεῖται ἐκ τοῦ λαοῦ σου ἀνὴρ εἷς.
\vs{12}Ἀνάμεινον ἐπὶ τῆς παρεμβολῆς σου, διαφυλάσσων πάντα ἄνδρα ἐκ τῆς δυνάμεώς σου, καὶ ἐπικρατησάτωσαν οἱ παῖδές σου τῆς πηγῆς τοῦ ὕδατος, ἣ ἐκπορεύεται ἐκ τῆς ῥίζης τοῦ ὄρους,
\vs{13}διότι ἐκεῖθεν ὑδρεύονται πάντες οἱ κατοικοῦντες Βετυλούα, καὶ ἀνελεῖ αὐτοὺς ἡ δίψα, καὶ ἐκδώσουσι τὴν πόλιν ἑαυτῶν· καὶ ἡμεῖς καὶ ὁ λαὸς ἡμῶν ἀναβησόμεθα ἐπὶ τὰς πλησίον κορυφὰς τῶν ὀρέων, καὶ παρεμβαλοῦμεν ἐπʼ αὐταῖς εἰς προφυλακὴν, τοῦ μὴ ἐξελθεῖν ἐκ τῆς πόλεως ἄνδρα ἕνα.
\vs{14}Καὶ τακήσονται ἐν τῷ λιμῷ αὐτοὶ, καὶ αἱ γυναῖκες αὐτῶν, καὶ τὰ τέκνα αὐτῶν· καὶ πρὶν ἐλθεῖν τὴν ῥομφαίαν ἐπʼ αὐτοὺς, καταστρωθήσονται ἐν ταῖς πλατείαις τῆς οἰκήσεως αὐτῶν,
\vs{15}καὶ ἀνταποδώσεις αὐτοῖς ἀνταπόδομα πονηρὸν, ἀνθʼ ὧν ἐστασίασαν, καὶ οὐκ ἀπήντησαν τῷ προσώπῳ σου ἐν εἰρήνῃ.

\vs{16}Καὶ ἤρεσαν οἱ λόγοι αὐτῶν ἐνώπιον Ὀλοφέρνου, καὶ ἐνώπιον πάντων τῶν θεραπόντων αὐτοῦ, καὶ συνέταξαν ποιεῖν καθὼς ἐλάλησαν.
\vs{17}Καὶ ἀπῇρε παρεμβολὴ υἱῶν Ἀμμὼν, καὶ μετʼ αὐτῶν χιλιάδες πέντε υἱῶν Ἀσσούρ· καὶ παρενέβαλον ἐν τῷ αὐλῶνι, καὶ προκατελάβοντο τὰ ὕδατα, καὶ τὰς πηγὰς τῶν ὑδάτων τῶν υἱῶν Ἰσραήλ.

\vs{18}Καὶ ἀνέβησαν υἱοὶ Ἡσαῦ, καὶ οἱ υἱοὶ Ἀμμών, καὶ παρενέβαλον ἐν τῇ ὀρεινῇ ἀπέναντι Δωθαῒμ, καὶ ἀπέστειλαν ἐξ αὐτῶν πρὸς Νότον καὶ ἀπηλιώτην ἀπέναντι Ἐκρεβὴλ, ἥ ἐστι πλησίον Χοὺς, ἥ ἐστιν ἐπὶ τοῦ χειμάῤῥου Μοχμούρ· καὶ ἡ λοιπὴ στρατιὰ τῶν Ἀσσυρίων παρενέβαλον ἐν τῷ πεδίῳ, καὶ ἐκάλυψαν πᾶν τὸ πρόσωπον τῆς γῆς· καὶ αἱ σκηναὶ καὶ αἱ ἀπαρτίαι αὐτῶν κατεστρατοπέδευσαν ἐν ὄχλῳ πολλῷ, καὶ ἦσαν εἰς πλῆθος πολὺ σφόδρα.

\vs{19}Καὶ οἱ υἱοὶ Ἰσραὴλ ἀνεβόησαν πρὸς Κύριον Θεὸν αὐτῶν, ὅτι ὠλιγοψύχησε τὸ πνεῦμα αὐτῶν, ὅτι ἐκύκλωσαν πάντες οἱ ἐχθροὶ αὐτῶν, καὶ οὐκ ἦν διαφυγεῖν ἐκ μέσου αὐτῶν.
\vs{20}Καὶ ἔμεινε κύκλῳ αὐτῶν πᾶσα παρεμβολὴ Ἀσσοὺρ, οἱ πεζοὶ καὶ τὰ ἅρματα καὶ οἱ ἱππεῖς αὐτῶν, ἡμέρας τριακοντατέσσαρας· καὶ ἐξέλιπε πάντας τοὺς κατοικοῦντας Βετυλούα πάντα τὰ ἀγγεῖα αὐτῶν τῶν ὑδάτων.
\vs{21}Καὶ οἱ λάκκοι ἐξεκενοῦντο, καὶ οὐκ εἶχον πιεῖν εἰς πλησμονὴν ὕδωρ ἡμέραν μίαν, ὅτι ἐν μέτρῳ ἐδίδοσαν αὐτοῖς πιεῖν.
\vs{22}Καὶ ἠθύμησαν τὰ νήπια αὐτῶν, καὶ αἱ γυναῖκες αὐτῶν καὶ οἱ νεανίσκοι ἐξέλιπον ἀπὸ τῆς δίψης· καὶ ἔπιπτον ἐν ταῖς πλατείαις τῆς πόλεως, καὶ ἐν ταῖς διόδοις τῶν πυλῶν, καὶ οὐκ ἦν κραταίωσις ἔτι ἐν αὐτοῖς.

\vs{23}Καὶ ἐπισυνήχθησαν πᾶς ὁ λαὸς ἐπὶ Ὀζίαν καὶ τοὺς ἄρχοντας τῆς πόλεως, οἱ νεανίσκοι καὶ αἱ γυναῖκες καὶ τὰ παιδία, καὶ ἀνεβόησαν φωνῇ μεγάλῃ, καὶ εἶπαν ἐναντίον πάντων τῶν πρεσβυτέρων,
\vs{24}κρίναι ὁ Θεὸς ἀναμέσον ἡμῶν καὶ ὑμῶν, ὅτι ἐποιήσατε ἐν ἡμῖν ἀδικίαν· μεγάλην, οὐ λαλήσαντες εἰρηνικὰ μετὰ τῶν υἱῶν Ἀσσούρ.
\vs{25}Καὶ νῦν οὐκ ἔστι βοηθὸς ἡμῶν, ἀλλὰ πέπρακεν ἡμᾶς ὁ Θεὸς εἰς τὰς χεῖρας αὐτῶν, τοῦ καταστρωθῆναι ἐναντίον αὐτῶν ἐν δίψῃ καὶ ἀπωλείᾳ μεγάλῃ.

\vs{26}Καὶ νῦν ἐπικαλέσασθε αὐτοὺς, καὶ ἔκδοσθε τὴν πόλιν πᾶσαν εἰς προνομὴν τῷ λαῷ Ὀλοφέρνου, καὶ πάσῃ τῇ δυνάμει αὐτοῦ.
\vs{27}Κρεῖσσον γὰρ ἡμῖν γενηθῆναι αὐτοῖς εἰς διαρπαγήν· ἐσόμεθα γὰρ εἰς δούλους, καὶ ζήσεται ἡ ψυχὴ ἡμῶν, καὶ οὐκ ὀψόμεθα τὸν θάνατον τῶν νηπίων ἡμῶν ἐν ὀφθαλμοῖς ἡμῶν, καὶ τὰς γυναῖκας καὶ τὰ τέκνα ἡμῶν ἐκλειπούσας τὰς ψυχὰς αὐτῶν.
\vs{28}Μαρτυρόμεθα ὑμῖν τὸν οὐρανὸν καὶ τὴν γῆν καὶ τὸν Θεὸν ἡμῶν, καὶ Κύριον τῶν πατέρων ἡμῶν, ὃς ἐκδικεῖ ἡμᾶς κατὰ τὰς ἁμαρτίας ἡμῶν, καὶ κατὰ τὰ ἁμαρτήματα τῶν πατέρων ἡμῶν, ἵνα μὴ ποιήσῃ κατὰ τὰ ῥήματα ταῦτα ἐν τῇ ἡμέρᾳ τῇ σήμερον·
\vs{29}καὶ ἐγένετο κλαυθμὸς μέγας ἐν μέσῳ τῆς ἐκκλησίας πάντων ὁμοθυμαδὸν, καὶ ἐβόησαν πρὸς Κύριον τὸν Θεὸν φωνῇ μεγάλῃ.

\vs{30}Καὶ εἶπε πρὸς αὐτοὺς Ὀζίας, θαρσεῖτε ἀδελφοὶ, διακαρτερήσωμεν ἔτι πέντε ἡμέρας, ἐν αἷς ἐπιστρέψει Κύριος ὁ Θεὸς ἡμῶν τὸ ἔλεος αὐτοῦ ἐφʼ ἡμᾶς· οὐ γὰρ ἐγκαταλείψει ἡμᾶς εἰς τέλος.
\vs{31}Ἐὰν δὲ διέλθωσιν αὗται, καὶ μὴ ἔλθῃ ἐφʼ ἡμᾶς βοήθεια, ποιήσω κατὰ τὰ ῥήματα ὑμῶν.
\vs{32}Καὶ ἐσκόρπισε τὸν λαὸν εἰς τὴν ἑαυτοῦ παρεμβολήν· καὶ ἐπὶ τὰ τείχη καὶ τοὺς πύργους τῆς πόλεως αὐτῶν ἀπῆλθον, καὶ τὰς γυναῖκας καὶ τὰ τέκνα εἰς τοὺς οἴκους αὐτῶν ἐξαπέστειλε· καὶ ἦσαν ἐν ταπεινώσει πολλῇ ἐν τῇ πόλει.

\ch{8}
Καὶ ἤκουσεν ἐν ἐκείναις ταῖς ἡμέραις Ἰουδὶθ, θυγάτηρ Μεραρὶ, υἱοῦ Ὢξ, υἱοῦ Ἰωσὴφ, υἱοῦ Ὀζιὴλ, υἱοὺ Ἑλκία, υἱοῦ Ἠλιοὺ, υἱοῦ Χελκίου, υἱοῦ Ἐλιὰβ, υἱοῦ Ναθαναὴλ, υἱοῦ Σαλαμιὴλ, υἱοῦ Σαρασαδαΐ, υἱοῦ Ἰσραήλ.

\vs{2}Καὶ ὁ ἀνὴρ αὐτῆς Μανασσῆς, τῆς φυλῆς αὐτῆς, καὶ τῆς πατριᾶς αὐτῆς, καὶ ἀπέθανεν ἐν ἡμέραις θερισμοῦ κριθῶν·
\vs{3}ἐπέστη γὰρ ἐπὶ τοῦ δεσμεύοντος τὸ δρᾶγμα ἐν τῷ πεδίῳ, καὶ ὁ καύσων ἦλθεν ἐπὶ τὴν κεφαλὴν αὐτοῦ, καὶ ἔπεσεν ἐπὶ τὴν κλίνην, καὶ ἐτελεύτησεν ἐν Βετυλούᾳ τῇ πόλει αὐτοῦ, καὶ ἔθαψαν αὐτὸν μετὰ τῶν πατέρων αὐτοῦ ἐν τῷ ἀγρῷ τῷ ἀναμέσον Δωθαῒμ καὶ Βελαμών.

\vs{4}Καὶ ἦν Ἰουδὶθ ἐν τῷ οἴκῳ αὐτῆς χηρεύουσα ἔτη τρία καὶ μῆνας τέσσαρας.
\vs{5}Καὶ ἐποίησεν ἑαυτῇ σκηνὴν ἐπὶ τοῦ δώματος τοῦ οἴκου αὐτῆς, καὶ ἐπέθηκεν ἐπὶ τὴν ὀσφῦν αὐτῆς σάκκον· καὶ ἦν ἐπʼ αὐτῆς τὰ ἱμάτια τῆς χηρεύσεως αὐτῆς.
\vs{6}Καὶ ἐνήστευε πάσας τὰς ἡμέρας χηρεύσεως αὐτῆς χωρὶς προσαββάτων, καὶ σαββάτων, καὶ προνουμηνιῶν, καὶ νουμηνιῶν, καὶ ἑορτῶν, καὶ χαρμοσυνῶν οἴκου Ἰσραήλ.
\vs{7}Καὶ ἦν καλὴ τῷ εἴδει, καὶ ὡραία τῇ ὄψει σφόδρα· καὶ ὑπελείπετο αὐτῇ Μανασσῆς ὁ ἀνὴρ αὐτῆς χρυσίον καὶ ἀργύριον, καὶ παῖδας καὶ παιδίσκας, καὶ κτήνη καὶ ἀγροὺς, καὶ ἔμενεν ἐπʼ αὐτῶν.
\vs{8}Καὶ οὐκ ἦν ὃς ἐπήνεγκεν αὐτῇ ῥῆμα πονηρὸν, ὅτι ἐφοβεῖτο τὸν Θεὸν σφόδρα.

\vs{9}Καὶ ἤκουσε τὰ ῥήματα τοῦ λαοῦ τὰ πονηρὰ ἐπὶ τὸν ἄρχοντα, ὅτι ὠλιγοψύχησαν ἐπὶ τῇ σπάνει τῶν ὑδάτων· καὶ ἤκουσε πάντας τοὺς λόγους Ἰουδὶθ οὓς ἐλάλησε πρὸς αὐτοὺς Ὀζίας, ὡς ὤμοσεν αὐτοῖς παραδώσειν τὴν πόλιν μετὰ ἡμέρας πέντε τοῖς Ἀσσυρίοις.
\vs{10}Καὶ ἀποστείλασα τὴν ἅβραν αὐτῆς τὴν ἐφεστῶσαν πᾶσι τοῖς ὑπάρχουσιν αὐτῆς, ἐκάλεσεν Ὀζίαν καὶ Χαβρὶν καὶ Χαρμὶν τοὺς πρεσβυτέρους τῆς πόλεως αὐτῆς.
\vs{11}Καὶ ἦλθον πρὸς αὐτὴν, καὶ εἶπε πρὸς αὐτοὺς, ἀκούσατε δή μου ἄρχοντες τῶν κατοικούντων ἐν Βετυλούᾳ· ὅτι οὐκ εὐθὴς ὁ λόγος ὑμῶν ὃν ἐλαλήσατε ἐναντίον τοῦ λαοῦ ἐν τῇ ἡμέρᾳ ταύτῃ, καὶ ἐστήσατε τὸν ὅρκον ὃν ἐλαλήσατε ἀναμέσον τοῦ Θεοῦ καὶ ὑμῶν, καὶ εἴπατε ἐκδώσειν τὴν πόλιν τοῖς ἐχθροῖς ὑμῶν, ἐὰν μὴ ἐν αὐταῖς ἐπιστρέψῃ ὁ Κύριος βοηθῆσαι ἡμῖν.
\vs{12}Καὶ νῦν τίνες ἐστὲ ὑμεῖς οἳ ἐπειράσατε τὸν Θεὸν ἐν τῇ ἡμέρᾳ τῇ σήμερον, καὶ ἵστασθε ὑπὲρ τοῦ Θεοῦ ἐν μέσῳ υἱῶν ἀνθρώπων;

\vs{13}Καὶ νῦν Κύριον παντοκράτορα ἐξετάζετε, καὶ οὐθὲν ἐπιγνώσεσθε ἕως τοῦ αἰῶνος·
\vs{14}ὅτι βάθος καρδίας ἀνθρώπου οὐχ εὑρήσετε, καὶ λόγους τῆς διανοίας αὐτοῦ οὐ λήψεσθε, καὶ πῶς τὸν Θεὸν ὃς ἐποίησε τὰ πάντα ταῦτα, ἐρευνήσετε, καὶ τὸν νοῦν αὐτοῦ ἐπιγνώσεσθε, καὶ τὸν λογισμὸν αὐτοῦ κατανοήσετε; μηδαμῶς, ἀδελφοὶ, μὴ παροργίζετε Κύριον τὸν Θεὸν ἡμῶν,
\vs{15}ὅτι ἐὰν μὴ βούληται ἐν ταῖς πέντε ἡμέραις βοηθῆσαι ἡμῖν, αὐτὸς ἔχει τὴν ἐξουσίαν ἐν αἷς θέλει σκεπάσαι ἡμέραις, ἢ καὶ ὀλοθρεῦσαι ἡμᾶς πρὸ προσώπου τῶν ἐχθρῶν ἡμῶν.

\vs{16}Ὑμεῖς δὲ μὴ ἐνεχυράζετε τὰς βουλὰς Κυρίου τοῦ Θεοῦ ἡμῶν, ὅτι οὐχ ὡς ἄνθρωπος ὁ Θεὸς ἀπειληθῆναι, οὐδὲ ὡς υἱὸς ἀνθρώπου διαιτηθῆναι.
\vs{17}Διόπερ ἀναμένοντες τὴν παρʼ αὐτοῦ σωτηρίαν, ἐπικαλεσώμεθα αὐτὸν εἰς βοήθειαν ἡμῶν, καὶ εἰσακούσεται τῆς φωνῆς ἡμῶν, ἐὰν ᾖ αὐτῷ ἀρεστόν.

\vs{18}Ὅτι οὐκ ἀνέστη ἐν ταῖς γενεαῖς ἡμῶν, οὐδε ἐστὶν ἐν τῇ ἡμέρᾳ τῇ σήμερον οὔτε φυλὴ, οὔτε πατριὰ, οὔτε δῆμος, οὔτε πόλις ἐξ ἡμῶν, οἳ προσκυνοῦσι θεοῖς χειροποιήτοις, καθάπερ ἐγένετο ἐν ταῖς πρότερον ἡμέραις,
\vs{19}ὧν χάριν ἐδόθησαν εἰς ῥομφαίαν καὶ εἰς διαρπαγὴν οἱ πατέρες ἡμῶν, καὶ ἔπεσον πτῶμα μέγα ἐνώπιον τῶν ἐχθρῶν ἡμῶν.
\vs{20}Ἡμεῖς δὲ ἕτερον θεὸν οὐκ ἐπέγνωμεν πλὴν αὐτοῦ· ὅθεν ἐλπίζομεν ὅτι οὐχ ὑπερόψεται ἡμᾶς, οὐδʼ ἀπὸ τοῦ γένους ἡμῶν.

\vs{21}Ὅτι ἐν τῷ ληφθῆναι ἡμᾶς, οὕτως καθήσεται πᾶσα ἡ Ἰουδαία, καὶ προνομευθήσεται τὰ ἅγια ἡμῶν, καὶ ζητήσει τὴν βεβήλωσιν αὐτῶν ἐκ τοῦ στόματος ἡμῶν,
\vs{22}καὶ τὸν φόνον τῶν ἀδελφῶν ἡμῶν, καὶ τὴν αἰχμαλωσίαν τῆς γῆς, καὶ τὴν ἐρήμωσιν τῆς κληρονομίας ἡμῶν ἐπιστρέψει εἰς κεφαλὴν ἡμῶν ἐν τοῖς ἔθνεσιν, οὗ ἐὰν δουλεύσωμεν ἐκεῖ, καὶ ἐσόμεθα εἰς πρόσκομμα καὶ εἰς ὄνειδος ἐναντίον τῶν κτωμένων ἡμᾶς·
\vs{23}ὅτι οὐ κατευθυνθήσεται ἡ δουλεία ἡμῶν εἰς χάριν. ἀλλʼ εἰς ἀτιμίαν θήσει αὐτὴν Κύριος ὁ Θεὸς ἡμῶν.

\vs{24}Καὶ νῦν ἀδελφοὶ ἐπιδειξώμεθα τοῖς ἀδελφοῖς ἡμῶν, ὅτι ἐξ ἡμῶν κρέμαται ἡ ψυχὴ αὐτῶν, καὶ τὰ ἅγια καὶ ὁ οἶκος καὶ τὸ θυσιαστήριον ἐπεστήρικται ἐφʼ ἡμῖν.

\vs{25}Παρὰ ταῦτα πάντα εὐχαριστήσωμεν Κυρίῳ τῷ Θεῷ ἡμῶν, ὃς πειράζει ἡμᾶς καθὰ καὶ τοὺς πατέρας ἡμῶν.
\vs{26}Μνήσθητε ὅσα ἐποίησε μετὰ Ἁβραὰμ, καὶ ὅσα ἐπείρασε τὸν Ἰσαὰκ, καὶ ὅσα ἐγένετο τῷ Ἰακὼβ ἐν Μεσοποταμίᾳ τῆς Συρίας ποιμαίνοντι τὰ πρόβατα Λάβαν τοῦ ἀδελφοῦ τῆς μητρὸς αὐτοῦ·
\vs{27}ὅτι οὐ καθὼς ἐκείνους ἐπύρωσεν εἰς ἐτασμὸν τῆς καρδίας αὐτῶν, καὶ ἡμᾶς οὐκ ἐξεδίκησεν, ἀλλʼ εἰς νουθέτησιν μαστιγοῖ Κύριος τοὺς ἐγγίζοντας αὐτῷ.

\vs{28}Καὶ εἶπε πρὸς αὐτὴν Ὀζίας, πάντα ὅσα εἶπας, ἀγαθῇ καρδίᾳ ἐλάλησας, καὶ οὐκ ἔστιν ὃς ἀντιστήσεται τοῖς λόγοις σου.
\vs{29}Ὅτι οὐκ ἐν τῇ σήμερον ἡ σοφία σου πρόδηλός ἐστιν, ἀλλὰ ἀπʼ ἀρχῆς ἡμερῶν σου ἔγνω πᾶς ὁ λαὸς τὴν σύνεσίν σου, καθότι ἀγαθόν ἐστι τὸ πλάσμα τῆς καρδίας σου.
\vs{30}Ἀλλʼ ὁ λαὸς ἐδίψησε σφόδρα, καὶ ἠνάγκασαν ποιῆσαι ἡμᾶς καθὰ ἐλαλήσαμεν αὐτοῖς, καὶ ἀπαγαγεῖν ὅρκον ἐφʼ ἡμᾶς, ὃν οὐ παραβησόμεθα.
\vs{31}Καὶ νῦν δεήθητι περὶ ἡμῶν, ὅτι γυνὴ εὐσεβὴς εἶ, καὶ ἀποστελεῖ Κύριος τὸν ὑετὸν εἰς πλήρωσιν τῶν λάκκων ἡμῶν, καὶ οὐκ ἐκλείψομεν ἔτι.

\vs{32}Καὶ εἶπε πρὸς αὐτοὺς Ἰουδὶθ, ἀκούσατέ μου, καὶ ποιήσω πρᾶγμα ὃ ἀφίξεται εἰς γενεὰς γενεῶν υἱοῖς τοῦ γένους ἡμῶν.
\vs{33}Ὑμεῖς στήσεσθε ἐπὶ τῆς πύλης τὴν νύκτα ταύτην, καὶ ἐξελεύσομαι ἐγὼ μετὰ τῆς ἅβρας μου, καὶ ἐν ταῖς ἡμέραις μεθʼ ἃς εἴπατε παραδώσειν τὴν πόλιν τοῖς ἐχθροῖς ἡμῶν, ἐπισκέψεται Κύριος τὸν Ἰσραὴλ ἐν χειρί μου.
\vs{34}Ὑμεῖς δὲ οὐκ ἐξερευνήσετε τὴν πρᾶξίν μου, οὐ γὰρ ἐρῶ ὑμῖν, ἕως τοῦ τελεσθῆναι ἃ ἐγὼ ποιῶ.

\vs{35}Καὶ εἶπεν Ὀζίας καὶ οἱ ἄρχοντες πρὸς αὐτὴν, πορεύου εἰς εἰρήνην, καὶ Κύριος ὁ Θεὸς ἔμπροσθέν σον εἰς ἐκδίκησιν τῶν ἐχθρῶν ἡμῶν.
\vs{36}Καὶ ἀποστρέψαντες ἐκ τῆς σκηνῆς, ἐπορεύθησαν ἐπὶ τὰς διατάξεις αὐτῶν.

\ch{9}
Ἰουδὶθ δὲ ἔπεσεν ἐπὶ πρόσωπον, καὶ ἐπέθετο σποδὸν ἐπὶ τὴν κεφαλὴν αὐτῆς, καὶ ἐγύμνωσεν ὃν ἐνεδιδύσκετο σάκκον· καὶ ἦν ἄρτι προσφερόμενον ἐν Ἱερουσαλὴμ εἰς τὸν οἶκον τοῦ Θεοῦ τὸ θυμίαμα τῆς ἑσπέρας ἐκείνης· καὶ ἐβόησε φωνῇ μεγάλῃ Ἰουδὶθ πρὸς Κύριον, καὶ εἶπε,
\vs{2}κύριε ὁ Θεὸς τοῦ πατρός μου Συμεὼν, ᾧ ἔδωκας ἐν χειρὶ ῥομφαίαν εἰς ἐκδίκησιν ἀλλογενῶν, οἳ ἔλυσαν μήτραν παρθένου εἰς μίασμα, καὶ ἐγύμνωσαν μηρὸν εἰς αἰσχύνην, καὶ ἐβεβήλωσαν μήτραν εἰς ὄνειδος· εἶπας γὰρ, οὐχ οὕτως ἔσται, καὶ ἐποίησαν·
\vs{3}ἀνθʼ ὧν ἔδωκας ἄρχοντας αὐτῶν εἰς φόνον, καὶ τὴν στρωμνὴν αὐτῶν ἣ ᾐδέσατο τὴν ἀπάτην αὐτῶν, εἰς αἷμα, καὶ ἐπάταξας δούλους ἐπὶ δυνάσταις, καὶ δυνάστας ἐπὶ θρόνους αὐτῶν·
\vs{4}καὶ ἔδωκας γυναῖκας αὐτῶν εἰς προνομὴν, καὶ θυγατέρας εἰς αἰχμαλωσίαν, καὶ πάντα τὰ σκῦλα εἰς διαίρεσιν υἱῶν ἠγαπημένων ὑπὸ σοῦ, οἳ καὶ ἐζήλωσαν τὸν ζῆλόν σου, καὶ ἐβδελύξαντο μίασμα αἵματος αὐτῶν, καὶ ἐπεκαλέσαντό σε εἰς βοηθόν· ὁ Θεὸς ὁ Θεὸς ὁ ἐμὸς, καὶ εἰσάκουσον ἐμοῦ τῆς χήρας.

\vs{5}Σὺ γὰρ ἐποίησας τὰ πρότερα ἐκείνων, καὶ ἐκεῖνα, καὶ τὰ μετέπειτα, καὶ τὰ νῦν, καὶ τὰ ἐπερχόμενα διενοήθης, καὶ ἐγενήθησαν ἃ ἐνενοήθης,
\vs{6}καὶ παρέστησαν ἃ ἐβουλεύσω, καὶ εἶπαν, ἰδοὺ πάρεσμεν· πᾶσαι γὰρ αἱ ὁδοί σου ἕτοιμοι, καὶ ἡ κρίσις σου ἐν προγνώσει.

\vs{7}Ἰδοὺ γὰρ Ἀσσύριοι ἐπληθύνθησαν ἐν δυνάμει αὐτῶν, ὑψώθησαν ἐφʼ ἵππῳ καὶ ἀναβάτῃ, ἐγανρίασαν ἐν βραχίονι πεζῶν, ἤλπισαν ἐν ἀσπίδι καὶ ἐν γαισῷ καὶ τόξῳ καὶ σφενδόνῃ, καὶ οὐκ ἔγνωσαν ὅτι σὺ εἶ Κύριος συντρίβων πολέμους· Κύριος ὄνομά σοι.
\vs{8}Σὺ ῥάξον αὐτῶν τὴν ἰσχὺν ἐν δυνάμει σου, καὶ κάταξον τὸ κράτος αὐτῶν ἐν τῷ θυμῷ σου· ἐβουλεύσαντο γὰρ βεβηλῶσαι τὰ ἅγιά σου, μιᾶναι τὸ σκήνωμα τῆς καταπαύσεως τοῦ ὀνόματος τῆς δόξης σου, καὶ καταβαλεῖν σιδήρῳ κέρας θυσιαστηρίου σου.

\vs{9}Βλέψον εἰς ὑπερηφανίαν αὐτῶν, ἀπόστειλον τὴν ὀργήν σου εἰς κεφαλὰς αὐτῶν· δὸς ἐν χειρί μου τῆς χήρας ὃ διενοήθην κράτος,
\vs{10}πάταξον δοῦλον ἐκ χειλέων ἀπάτης μου ἐπʼ ἄρχοντι, καὶ ἄρχοντα ἐπὶ θεράποντι αὐτοῦ, θραῦσον αὐτῶν τὸ ἀνάστεμα ἐν χειρὶ θηλείας.
\vs{11}Οὐ γὰρ ἐν πλήθει τὸ κράτος σου, οὐδὲ ἡ δυναστεία σου ἐν ἰσχύουσιν, ἀλλὰ ταπεινῶν εἶ Θεὸς, ἐλαττόνων εἶ βοηθὸς, ἀντιλήπτωρ ἀσθενούντων, ἀπεγνωσμένων σκεπαστὴς, ἀπηλπισμένων σωτήρ.

\vs{12}Ναὶ ναὶ ὁ Θεὸς τοῦ πατρὸς μοὺ, καὶ Θεὸς κληρονομίας Ἰσραὴλ, δέσποτα τῶν οὐρανῶν καὶ τῆς γῆς, κτίστα τῶν ὑδάτων, βασιλεῦ πάσης κτίσεώς σου, σὺ εἰσάκουσον τῆς δεήσεώς μου,
\vs{13}καὶ δὸς λόγον μου καὶ ἀπάτην εἰς τραῦμα καὶ μώλωπα αὐτῶν, οἳ κατὰ τῆς διαθήκης σου, καὶ οἴκου ἡγιασμένου σου, καὶ κορυφῆς Σιὼν, καὶ οἴκου κατασχέσεως υἱῶν σου ἐβουλεύσαντο σκληρά.

\vs{14}Καὶ ποίησον ἐπὶ πᾶν τὸ ἔθνος σου, καὶ πάσης φυλῆς ἐπίγνωσιν, τοῦ εἰδῆσαι ὅτι σὺ εἶ ὁ Θεὸς πάσης δυνάμεως καὶ κράτους, καὶ οὐκ ἔστιν ἄλλος ὑπερασπίζων τοῦ γένους Ἰσραὴλ, εἰ μὴ σύ.

\ch{10}
Καὶ ἐγένετο ὡς ἐπαύσατο βοῶσα πρὸς τὸν Θεὸν Ἰσραὴλ, καὶ συνετέλεσε πάντα τὰ ῥήματα ταῦτα,
\vs{2}καὶ ἀνέστη ἀπὸ τῆς πτώσεως καὶ ἐκάλεσε τὴν ἅβραν αὐτῆς, καὶ κατέβη εἰς τὸν οἶκον ἐν ᾧ διέτριβεν ἐν αὐτῷ ἐν ταῖς ἡμέραις τῶν σαββάτων, καὶ ἐν ταῖς ἑορταῖς αὐτῆς,
\vs{3}καὶ περιείλατο τὸν σάκκον ὃν ἐνεδεδύκει, καὶ ἐξεδύσατο τὰ ἱμάτια τῆς χηρεύσεως αὐτῆς, καὶ περιεκλύσατο τὸ σῶμα ὕδατι, καὶ ἐχρίσατο μύρῳ παχεῖ, καὶ διέταξε τὰς τρίχας τῆς κεφαλῆς αὐτῆς, καὶ ἐπέθετο μίτραν ἐπʼ αὐτῆς, καὶ ἐνεδύσατο τὰ ἱμάτια τῆς εὐφροσύνης αὐτῆς, ἐν οἷς ἐστολίζετο ἐν ταῖς ἡμέραις τῆς ζωῆς τοῦ ἀνδρὸς αὐτῆς Μανασσῆ·
\vs{4}καὶ ἔλαβε σανδάλια εἰς τοὺς πόδας αὐτῆς, καὶ περιέθετο τοὺς χλιδῶνας, καὶ τὰ ψέλλια, καὶ τοὺς δακτυλίους, καὶ τὰ ἐνώτια, καὶ πάντα τὸν κόσμον αὐτῆς· καὶ ἐκαλλωπίσατο σφόδρα εἰς ἀπάτησιν ὀφθαλμῶν ἀνδρῶν, ὅσοι ἂν ἴδωσιν αὐτήν.

\vs{5}Καὶ ἔδωκε τῇ ἅβρᾳ αὐτῆς ἀσκοπυτίνην οἴνου, καὶ καμψάκην ἐλαίου, καὶ πήραν ἐπλήρωσεν ἀλφίτων καὶ παλάθης καὶ ἄρτων καθαρῶν, καὶ περιεδίπλωσε πάντα τὰ ἀγγεῖα αὐτῆς, καὶ ἐπέθηκεν ἐπʼ αὐτῇ.
\vs{6}Καὶ ἐξήλθοσαν ἐπὶ τὴν πύλην τῆς πόλεως Βετυλούα, καὶ εὕροσαν ἐφεστῶτας ἐπʼ αὐτῇς Ὀζίαν, καὶ τοὺς πρεσβυτέρους τῆς πόλεως Χαβρὶν καὶ Χαρμίν.

\vs{7}Ὡς δὲ εἶδον αὐτὴν, καὶ ἦν ἠλλοιωμένον τὸ πρόσωπον αὐτῆς, καὶ τὴν στολὴν μεταβεβληκυῖαν αὐτῆς, καὶ ἐθαύμασαν ἐπὶ τῷ κάλλει αὐτῆς ἐπὶ πολὺ σφόδρα, καὶ εἶπαν αὐτῇ,
\vs{8}ὁ Θεὸς ὁ Θεὸς τῶν πατέρων ἡμῶν δῴη σε εἰς χάριν, καὶ τελειώσαι τὰ ἐπιτηδεύματά σου εἰς γαυρίαμα υἱῶν Ἱσραὴλ καὶ ὕψωμα Ἰερουσαλήμ· καὶ προσεκύνησε τῷ Θεῷ,

\vs{9}Καὶ εἶπε πρὸς αὐτοὺς, ἐπιτάξατε ἀνοῖξαί μοι τὴν πύλην τῆς πόλεως, καὶ ἐξελεύσομαι εἰς τελείωσιν τῶν λόγων, ὧν ἐλαλήσατε μετʼ ἐμοῦ· καὶ συνέταξαν τοῖς νεανίσκοις ἀνοῖξαι αὐτῇ καθότι ἐλάλησαν.

\vs{10}Καὶ ἐποίησαν οὕτως· καὶ ἐξῆλθεν Ἰουδὶθ, αὐτὴ καὶ ἡ παιδίσκη αὐτῆς μετʼ αὐτῆς· ἀπεσκόπευον δὲ αὐτὴν οἱ ἄνδρες τῆς πόλεως ἕως οὗ κατέβη τὸ ὄρος, ἕως διῆλθε τὸν αὐλῶνα, καὶ οὐκ ἔτι ἐθεώρουν αὐτήν.
\vs{11}Καὶ ἐπορεύοντο ἐν τῷ αὐλῶνι εἰς εὐθεῖαν, καὶ συνήντησεν αὐτῇ προφυλακὴ τῶν Ἀσσυρίων.
\vs{12}Καὶ συνέλαβον αὐτὴν, καὶ ἐπηρώτησαν, τίνων εἶ; καὶ πόθεν ἔρχῃ; καὶ ποῦ πορεύῃ; καὶ εἶπε, θυγάτηρ εἰμὶ τῶν Ἑβραίων, καὶ ἀποδιδράσκω ἀπὸ προσώπου αὐτῶν, ὅτι μέλλουσι δίδοσθαι ὑμῖν εἰς κατάβρωμα.
\vs{13}Κᾀγὼ ἔρχομαι εἰς τὸ πρόσωπον Ὀλοφέρνου ἀρχιστρατήγου δυνάμεως ὑμῶν, τοῦ ἀναγγεῖλαι ῥήματα ἀληθείας, καὶ δείξω πρὸ προσώπου αὐτοῦ ὁδὸν καθʼ ἣν πορεύσεται, καὶ κυριεύσει πάσης τῆς ὀρεινῆς, καὶ οὐ διαφωνήσει τῶν ἀνδρῶν αὐτοῦ σὰρξ μία, οὐδὲ πνεῦμα ζωῆς.

\vs{14}Ὡς δὲ ἤκουσαν οἱ ἄνδρες τὰ ῥήματα αὐτῆς, καὶ κατενόησαν τὸ πρόσωπον αὐτῆς, καὶ ἦν ἐναντίον αὐτῶν θαυμάσιον τῷ κάλλει σφόδρα, καὶ εἶπαν πρὸς αὐτὴν,
\vs{15}σέσωκας τὴν ψυχήν σου, σπεύσασα καταβῆναι εἰς πρόσωπον τοῦ κυρίου ἡμῶν· καὶ νῦν πρόσελθε ἐπὶ τὴν σκηνὴν αὐτοῦ, καὶ ἀφʼ ἡμῶν προπέμψουσί σε ἕως παραδώσουσί σε εἰς τὰς χεῖρας αὐτοῦ.
\vs{16}Ἐὰν δὲ στῇς ἐναντίον αὐτοῦ, μὴ φοβηθῇς τῇ καρδίᾳ σου, ἀλλὰ ἀνάγγειλον κατὰ τὰ ῥήματά σου, καὶ εὖ σε ποιήσει.

\vs{17}Καὶ ἐπέλεξαν ἐξ αὐτῶν ἄνδρας ἑκατὸν, καὶ παρέζευξαν αὐτῇ καὶ τῇ ἅβρᾳ αὐτῆς, καὶ ἤγαγον αὐτὰς ἐπὶ τὴν σκηνὴν Ὀλοφέρνου.
\vs{18}Καὶ ἐγένετο συνδρομὴ ἐν πὰσῃ τῇ παρεμβολῇ· διεβοήθη γὰρ εἰς τὰ σκηνώματα ἡ παρουσία αὐτῆς· καὶ ἐλθόντες ἐκύκλουν αὐτὴν ὡς εἱστήκει ἔξω τῆς σκηνῆς Ὀλοφέρνου, ἕως προσήγγειλαν αὐτῷ περὶ αὐτῆς.
\vs{19}Καὶ ἐθαύμαζον ἐπὶ τῷ κάλλει αὐτῆς, καὶ ἐθαύμαζον τοὺς υἱοὺς Ἰσραὴλ ἀπʼ αὐτῆς· καὶ εἶπεν ἕκαστος πρὸς τὸν πλησίον αὐτοῦ, τίς καταφρονήσει τοῦ λαοῦ τούτου, ὃς ἔχει ἐν ἑαυτῷ γυναῖκας τοιαύτας; ὅτι οὐ καλόν ἐστιν ὑπολείπεσθαι ἐξ αὐτῶν ἄνδρα ἕνα, οἳ ἀφεθέντες δυνήσονται κατασοφίσασθαι πᾶσαν τὴν γῆν·
\vs{20}καὶ ἐξῆλθον οἱ παρακαθεύδοντες Ὀλοφέρνῃ, καὶ πάντες οἱ θεράποντες αὐτοῦ, καὶ εἰσήγαγον αὐτὴν εἰς τὴν σκηνήν.

\vs{21}Καὶ ἦν Ὀλοφέρνης ἀναπαυόμενος ἐπὶ τῆς κλίνης αὐτοῦ ἐν τῷ κωνωπείῳ, ὃ ἦν ἐκ πορφύρας καὶ χρυσίου καὶ σμαράγδου καὶ λίθων πολυτελῶν καθυφασμένων.
\vs{22}Καὶ ἀνήγγειλαν αὐτῷ περὶ αὐτῆς, καὶ ἐξῆλθεν εἰς τὸ προσκήνιον, καὶ λαμπάδες ἀργυραῖ προάγουσαι αὐτοῦ.
\vs{23}Ὡς δὲ ἦλθε κατὰ πρόσωπον αὐτοῦ Ἰουδὶθ, καὶ τῶν θεραπόντων αὐτοῦ, ἐθαύμασαν πάντες ἐπὶ τῷ κάλλει τοῦ προσώπου αὐτῆς· καὶ πεσοῦσα ἐπὶ πρόσωπον προσεκύνησεν αὐτῷ, καὶ ἤγειραν αὐτὴν οἱ δοῦλοι αὐτοῦ.

\ch{11}
Καὶ εἶπε πρὸς αὐτὴν Ὀλοφέρνης, θάρσησον γύναι, μὴ φοβηθῇς τῇ καρδίᾳ σου, ὅτι ἐγὼ οὐκ ἐκάκωσα ἄνθρωπον ὅστις ᾑρέτικε δουλεύειν βασιλεῖ Ναβουχοδονόσορ πάσης τῆς γῆς.
\vs{2}Καὶ νῦν ὁ λαός σου ὁ κατοικῶν τὴν ὀρεινὴν, εἰ μὴ ἐφαύλισάν με, οὐκ ἂν ᾖρα τὸ δόρυ μου ἐπʼ αὐτοὺς, ἀλλʼ αὐτοὶ ἑαυτοῖς ἐποίησαν ταῦτα.
\vs{3}Καὶ νῦν λέγε μοι, τίνος ἕνεκεν ἀπέδρας ἀπʼ αὐτῶν, καὶ ἦλθες πρὸς ἡμᾶς; ἥκεις γὰρ εἰς σωτηρίαν· θάρσει, ἐν τῇ νυκτὶ ταύτῃ ζήσῃ, καὶ εἰς τὸ λοιπόν.
\vs{4}Οὐ γάρ ἐστιν ὃς ἀδικήσει σε, ἀλλʼ εὖ σε ποιήσει, καθὰ γίνεται τοῖς δούλοις τοῦ κυρίου μου βασιλέως Ναβουχοδονόσορ.

\vs{5}Καὶ εἶπε πρὸς αὐτὸν Ἰουδίθ, δέξαι τὰ ῥήματα τῆς δούλης σου, καὶ λαλησάτω ἡ παιδίσκη σου κατὰ πρόσωπόν σου, καὶ οὐκ ἀναγγελῶ ψεῦδος τῷ κυρίῳ μου ἐν τῇ νυκτὶ ταύτῃ.
\vs{6}Καὶ ἐὰν κατακολουθήσῃς τοῖς λόγοις τῆς παιδίσκης σου, τελείως πρᾶγμα ποιήσει μετὰ σοῦ ὁ Θεὸς, καὶ οὐκ ἀποπεσεῖται ὁ κύριός μου τῶν ἐπιτηδευμάτων αὐτοῦ.

\vs{7}Ζῇ γὰρ βασιλεὺς Ναβουχοδονόσορ πάσης τῆς γῆς, καὶ ζῇ τὸ κράτος αὐτοῦ, ὃς ἀπέστειλέ σε εἰς κατόρθωσιν πάσης ψυχῆς, ὅτι οὐ μόνον ἄνθρωποι διὰ σὲ δουλεύουσιν αὐτῷ, ἀλλὰ καὶ τὰ θηρία τοῦ ἀγροῦ, καὶ τὰ κτήνη, καὶ τὰ πετεινὰ τοῦ οὐρανοῦ διὰ τῆς ἰσχύος σου ζήσονται ἐπὶ Ναβουχοδονόσορ, καὶ πάντα τὸν οἶκον αὐτοῦ.
\vs{8}Ἠκούσαμεν γὰρ τὴν σοφίαν σου, καὶ τὰ πανουργεύματα τῆς ψυχῆς σου, καὶ ἀνηγγέλη πάσῃ τῇ γῇ, ὅτι σὺ μόνος ἀγαθὸς ἐν πάσῃ βασιλείᾳ, καὶ δυνατὸς ἐν ἐπιστήμῃ, καὶ θαυμαστὸς ἐν στρατεύμασι πολέμου.

\vs{9}Καὶ νῦν ὁ λόγος ὃν ἐλάλησεν Ἀχιὼρ ἐν τῇ συνεδρείᾳ σου, ἠκούσαμεν τὰ ῥήματα αὐτοῦ, ὅτι περιεποιήσαντο αὐτὸν οἱ ἄνδρες Βετυλούα, καὶ ἀνήγγειλεν αὐτοῖς πάντα ὅσα ἐξελάλησε παρὰ σοί.
\vs{10}Διὸ, δέσποτα κύριε, μὴ παρέλθῃς τὸν λόγον αὐτοῦ, ἀλλὰ κατάθου αὐτὸν ἐν τῇ καρδίᾳ σου, ὅτι ἀληθής ἐστιν· οὐ γὰρ ἐκδικᾶται τὸ γένος ἡμῶν, οὐ κατισχύει ῥομφαία ἐπʼ αὐτοὺς, ἐὰν μὴ ἁμάρτωσιν εἰς τὸν Θεὸν αὐτῶν.

\vs{11}Καὶ νῦν ἵνα μὴ γένηται ὁ κύριός μου εκβολος καὶ ἄπρακτος, καὶ ἐπιπεσεῖται θάνατος ἐπὶ πρόσωπον αὐτῶν, καὶ κατελάβετο αὐτοὺς ἁμάρτημα ἐν ᾧ παροργιοῦσι τὸν Θεὸν αὐτῶν, ὁπηνίκα ἂν ποιήσωσιν ἀτοπίαν.
\vs{12}Ἐπεὶ γὰρ ἐξέλιπεν αὐτοὺς τὰ βρώματα, καὶ ἐσπανίσθη πᾶν ὕδωρ, ἐβουλεύσαντο ἐπιβαλεῖν τοῖς κτήνεσιν αὐτῶν, καὶ πάντα ὅσα διεστείλατο αὐτοῖς ὁ Θεὸς ἐν τοῖς νόμοις αὐτοῦ μὴ φαγεῖν, διέγνωσαν δαπανῆσαι.
\vs{13}Καὶ τὰς ἀπαρχὰς τοῦ σίτου, καὶ τὰς δεκάτας τοῦ οἴνου καὶ τοῦ ἐλαίου, ἃ διεφύλαξαν ἁγιάσαντες τοῖς ἱερεῦσι τοῖς παρεστηκόσιν ἐν Ἱερουσαλὴμ ἀπέναντι τοῦ προσώπου τοῦ Θεοῦ ἡμῶν, κεκρίκασιν ἐξαναλῶσαι, ὧν οὐδὲ ταῖς χερσὶ καθῆκεν ἅψασθαι οὐδένα τῶν ἐκ τοῦ λαοῦ.
\vs{14}Καὶ ἀπεστάλκασιν εἰς Ἱερουσαλὴμ, ὅτι καὶ οἱ ἐκεῖ κατοικοῦντες ἐποίησαν ταῦτα, τοὺς μετοικίσαντας αὐτοῖς τὴν ἄφεσιν παρὰ τῆς γερουσίας.
\vs{15}Καὶ ἔσται ὡς ἂν ἀναγγείλῃ αὐτοῖς καὶ ποιήσωσιν, δοθήσονταί σοι εἰς ὄλεθρον ἐν τῇ ἡμέρᾳ ἐκείνῃ.

\vs{16}Ὅθεν ἐγὼ ἡ δούλη σου ἐπιγνοῦσα ταῦτα πάντα, ἀπέδρων ἀπὸ προσώπου αὐτῶν· καὶ ἀπέστειλέ με ὁ Θεὸς ποιῆσαι μετὰ σοῦ πράγματα, ἐφʼ οἷς ἐκστήσεται πᾶσα ἡ γῆ ὅσοι ἐὰν ἀκούσωσιν αὐτά.
\vs{17}Ὅτι ἡ δούλη σου θεοσεβής ἐστι, καὶ θεραπεύουσα νυκτὸς καὶ ἡμέρας τὸν Θεὸν τοῦ οὐρανοῦ· καὶ νῦν μενῶ παρὰ σοί, κύριέ μου, καὶ ἐξελεύσεται ἡ δούλη σου κατὰ νύκτα εἰς τὴν φάραγγα, καὶ προσεύξομαι πρὸς τὸν Θεόν· καὶ ἐρεῖ μοι πότε ἐποίησαν τὰ ἁμαρτήματα αὐτῶν·
\vs{18}καὶ ἐλθοῦσα προσανοίσω σοι· ἐξελεύσῃ σὺν πάσῃ τῇ δυνάμει σου, καὶ οὐκ ἔστιν ὃς ἀντιστήσεταί σοι ἐξ αὐτῶν.
\vs{19}Καὶ ἄξω σε διὰ μέσου τῆς Ἰουδαίας, ἕως τοῦ ἐλθεῖν ἀπέναντι Ἱερουσαλήμ· καὶ θήσω τὸν δίφρον σου ἐν μέσῳ αὐτῆς, καὶ ἄξεις αὐτοὺς ὡς πρόβατα οἷς οὐκ ἔστι ποιμήν· καὶ οὐ γρύξει κύων τῇ γλώσσῃ αὐτοῦ ἀπέναντί σου· ὅτι ταῦτα ἐλαλήθη μοι κατὰ πρόγνωσίν μου, καὶ ἀπηγγέλη μοι, καὶ ἀπεστάλην ἀναγγεῖλαί σοι.

\vs{20}Καὶ ἤρεσαν οἱ λόγοι αὐτῆς ἐναντίον Ὀλοφέρνου, καὶ ἐναντίον πάντων τῶν θεραπόντων αὐτοῦ, καὶ ἐθαύμασαν ἐπὶ τῇ σοφίᾳ αὐτῆς, καὶ εἶπαν,
\vs{21}οὐκ ἔστι τοιαύτη γυνὴ ἀπʼ ἄκρου ἕως ἄκρου τῆς γῆς, καλῷ προσώπῳ καὶ συνέσει λόγων.
\vs{22}Καὶ εἶπε πρὸς αὐτὴν Ὀλοφέρνης, εὖ ἐποίησεν ὁ Θεὸς ἀποστείλας σε ἔμπροσθεν τοῦ λαοῦ, τοῦ γενηθῆναι ἐν χερσὶν ἡμῶν κράτος· ἐν δὲ τοῖς φαυλίσασι τὸν κύριόν μου, ἀπώλειαν.
\vs{23}Καὶ νῦν ἀστεῖα εἶ σὺ ἐν τῷ εἴδει σου, καὶ ἀγαθὴ ἐν τοῖς λόγοις σου· ὅτι ἐὰν ποιήσῃς καθὰ ἐλάλησας, ὁ Θεός σου ἔσται μου Θεὸς, καὶ σὺ ἐν οἴκῳ βασιλέως Ναβουχοδονόσορ καθήσῃ, καὶ ἔσῃ ὀνομαστὴ παρὰ πᾶσαν τὴν γῆν.

\ch{12}
Καὶ ἐκέλευσεν εἰσαγαγεῖν αὐτὴν οὗ ἐτίθετο τὰ ἀργυρώματα, καὶ αὐτοῦ συνέταξε καταστρῶσαι αὐτῇ ἀπὸ τῶν ὀψοποιημάτων αὐτοῦ, καὶ τοῦ οἴνου αὐτοῦ πίνειν.

\vs{2}Καὶ εἶπεν Ἰουδὶθ, οὐ φάγομαι ἐξ αὐτῶν, ἵνα μὴ γένηται σκάνδαλον, ἀλλʼ ἐκ τῶν ἠκολουθηκότων μοι χορηγηθήσεται.
\vs{3}Καὶ εἶπε πρὸς αὐτὴν Ὀλοφέρνης, ἐὰν δὲ ἐκλίπῃ τὰ ὄντα μετὰ σοῦ, πόθεν ἐξοίσομέν σοι δοῦναι ὅμοια αὐτοῖς; οὐ γάρ ἐστι μεθʼ ἡμῶν ἐκ τοῦ ἔθνους σου.

\vs{4}Καὶ εἶπεν Ἰουδὶθ πρὸς αὐτὸν, ζῇ ἡ ψυχη σου, κύριέ μου, ὅτι οὐ δαπανήσι ἡ δούλη σου τὰ ὄντα μετʼ ἐμοῦ, ἕως ἂν ποιήσῃ Κὺριος ἐν χειρί μου ἃ ἐβουλεύσατο.

\vs{5}Καὶ ἠγάγοσαν αὐτὴν οἱ θεράποντες Ὀλοφέρνου εἰς τὴν σκηνὴν, καὶ ὕπνωσε μέχρι μεσούσης τῆς νυκτός· καὶ ἀνέστη πρὸς τὴν ἑωθινὴν φυλακὴν,
\vs{6}καὶ ἀπέστειλε πρὸς Ὀλοφέρνην, λέγουσα, ἐπιταξάτω δὴ ὁ κύριός μου, ἐᾶσαι τὴν δούλην σου ἐπὶ προσευχὴν ἐξελθεῖν.

\vs{7}Καὶ προσέταξεν Ὀλοφέρνης τοῖς σωματοφύλαξι μὴ διακωλύειν αὐτήν· καὶ παρέμεινεν ἐν τῇ παρεμβολῇ ἡμέρας τρεῖς, καὶ ἐξεπορεύετο κατὰ νύκτα εἰς τὴν φάραγγα Βετυλούα, καὶ ἐβαπτίζετο ἐν τῇ παρεμβολῇ ἐπὶ τῆς πηγῆς τοῦ ὕδατος.
\vs{8}Καὶ ὡς ἀνέβη, ἐδέετο τοῦ Κυρίου Θεοῦ Ἰσραὴλ κατευθύναι τὴν ὁδὸν αὐτῆς εἰς ἀνάστεμα τῶν υἱῶν τοῦ λαοῦ αὐτοῦ.
\vs{9}Καὶ εἰσπορευομένη καθαρὰ παρέμενε τῇ σκηνῇ, μέχρις οὗ προσηνέγκατο τὴν τροφὴν αὐτῆς πρὸς ἑσπέραν.

\vs{10}Καὶ ἐγένετο ἐν τῇ ἡμέρᾳ τῇ τετάρτῃ, ἐποίησεν Ὀλοφέρνης πότον τοῖς δούλοις αὐτοῦ μόνοις, καὶ οὐκ ἐκάλεσεν εἰς τὴν χρῆσιν οὐδένα τῶν πρὸς ταῖς χρείαις.
\vs{11}Καὶ εἶπε Βαγώᾳ τῷ εὐνούχῳ, ὃς ἦν ἀφεστηκὼς ἐπὶ πάντων τῶν αὐτοῦ, πεῖσον δὴ πορευθεὶς τὴν γυναῖκα τὴν Ἑβραίαν ἥ ἐστι παρὰ σοὶ, τοῦ ἐλθεῖν πρὸς ἡμᾶς, καὶ φαγεῖν καὶ πιεῖν μεθʼ ἡμῶν.
\vs{12}Ἰδοὺ γὰρ αἰσχρὸν τῷ προσώπῳ ἡμῶν, εἰ γυναῖκα τοιαύτην παρήσομεν οὐχ ὁμιλήσαντες αὐτῇ, ὅτι ἐὰν ταύτην μὴ ἐπισπασώμεθα, καταγελάσεται ἡμῶν.

\vs{13}Καὶ ἐξῆλθε Βαγώας ἀπὸ προσώπου Ὀλοφέρνου, καὶ εἰσῆλθε πρὸς αὐτὴν, καὶ εἶπε, μὴ ὀκνησάτω δὴ ἡ παιδίσκη ἡ καλὴ αὕτη ἐλθοῦσα πρὸς τὸν κύριόν μου, δοξασθῆναι κατὰ πρόσωπον αὐτοῦ, καὶ πίεσαι μεθʼ ἡμῶν εἰς εὐφροσύνην οἶνον, καὶ γενηθῆναι ἐν τῇ ἡμέρᾳ ταύτῃ ὡς θυγάτηρ μία τῶν υἱῶν Ἀσσοὺρ, αἳ παρεστήκασιν ἐν οἴκῳ Ναβουχοδονόσορ.

\vs{14}Καὶ εἶπε πρὸς αὐτὸν Ἰουδίθ, καὶ τίς εἰμι ἐγὼ ἀντεροῦσα τῷ κυρίῳ μου; ὅτι πᾶν ὃ ἔσται ἐν τοῖς ὀφθαλμοῖς αὐτοῦ ἀρεστὸν, σπεύσασα ποιήσω, καὶ ἔσται τοῦτο ἀγαλλίαμα ἕως ἡμέρας θανάτου μου.
\vs{15}Καὶ διαναστᾶσα ἐκοσμήθη τῷ ἱματισμῷ καὶ παντὶ τῷ κόσμῳ τῷ γυναικείῳ· καὶ προσῆλθεν ἡ δούλη αὐτῆς, καὶ ἔστρωσεν αὐτῇ κατέναντι Ὀλοφέρνου χαμαὶ τὰ κώδια, ἃ ἔλαβε παρὰ Βαγώου εἰς τὴν καθημερινὴν δίαιταν αὐτῆς, εἰς τὸ ἐσθίειν κατακλινομένην ἐπʼ αὐτῶν.

\vs{16}Καὶ εἰσελθοῦσα ἀνέπεσεν Ἰουδίθ, καὶ ἐξέστη ἡ καρδία Ὀλοφέρνου ἐπʼ αὐτὴν, καὶ ἐσαλεύθη ἡ ψυχὴ αὐτοῦ· καὶ ἦν κατεπίθυμος σφόδρα τοῦ συγγενέσθαι μετʼ αὐτῆς· καὶ ἐτήρει καιρὸν τοῦ ἀπατῆσαι αὐτὴν, ἀφʼ ἧς ἡμέρας εἶδεν αὐτήν.

\vs{17}Καὶ εἶπε πρὸς αὐτὴν Ὀλοφέρνης, πίε δὴ, καὶ γενήθητι μεθʼ ἡμῶν εἰς εὐφροσύνην.
\vs{18}Καὶ εἶπεν Ἰουδὶθ, πίομαι δὴ, κύριε, ὅτι ἐμεγαλύνθη τὸ ζῇν μου ἐν ἐμοὶ σήμερον παρὰ πάσας τὰς ἡμέρας τῆς γενέσεώς μου.
\vs{19}Καὶ λαβοῦσα ἔφαγε καὶ ἔπιε κατέναντι αὐτοῦ ἃ ἡτοίμασεν ἡ δούλη αὐτῆς.

\vs{20}Καὶ ηὐφράνθη Ὀλοφέρνης ἀπʼ αὐτῆς, καὶ ἔπιεν οἶνον πολὺν σφόδρα ὅσον οὐκ ἔπιε πώποτε ἐν ἡμέρᾳ μιᾷ ἀφʼ οὗ ἐγεννήθη.

\ch{13}
Ὡς δὲ ὀψία ἐγένετο, ἐσπούδασαν οἱ δοῦλοι αὐτοῦ ἀναλύειν· καὶ Βαγώας συνέκλεισε τὴν σκηνὴν ἔξωθεν, καὶ ἀπέκλεισε τοὺς παρεστῶτας ἐκ προσώπου τοῦ κυρίου αὐτοῦ, καὶ ἀπῴχοντο εἰς τὰς κοίτας αὐτῶν· ἦσαν γὰρ πάντες κεκοπωμένοι, διὰ τὸ ἐπὶ πλεῖον γεγονέναι τὸν πότον.
\vs{2}Ὑπελείφθη δὲ Ἰουδὶθ μόνη ἐν τῇ σκηνῇ, καὶ Ὀλοφέρνης προπεπτωκὼς ἐπὶ τὴν κλίνην αὐτοῦ· ἦν γὰρ περικεχυμένος αὐτῷ ὁ οἶνος.

\vs{3}Καὶ εἶπεν Ιουδὶθ τῇ δούλῃ αὐτῆς στῆναι ἔξω τοῦ κοιτῶνος αὐτῆς, καὶ ἐπιτηρεῖν τὴν ἔξοδον αὐτῆς καθάπερ καθʼ ἡμέραν, ἐξελεύσεσθαι γὰρ ἔφη ἐπὶ τὴν προσευχὴν αὐτῆς· καὶ τῷ Βαγώᾳ ἐλάλησε κατὰ τὰ ῥήματα ταῦτα.

\vs{4}Καὶ ἀπήλθοσαν πάντες ἐκ προσώπου, καὶ οὐδεὶς κατελείφθη ἐν τῷ κοιτῶνι ἀπὸ μικροῦ ἕως μεγάλου· καὶ στᾶσα Ἰουδὶθ παρὰ τὴν κλίνην αὐτοῦ, εἶπεν ἐν τῇ καρδίᾳ αὐτῆς, Κύριε ὁ Θεὸς πάσης δυνάμεως, ἐπίβλεψον ἐν τῇ ὥρᾳ ταύτῃ ἐπὶ τὰ ἔργα τῶν χειρῶν μου, εἰς ὕψωμα Ἱερουσαλήμ·
\vs{5}ὅτι νῦν καιρὸς ἀντιλαβέσθαι τῆς κληρονομίας σου, καὶ ποιῆσαι τὸ ἐπιτήδευμά μου, εἰς θραῦμα ἔχθρων οἳ ἐπανέστησαν ἡμῖν.

\vs{6}Καὶ προσελθοῦσα τῷ κανόνι τῆς κλίνης ὃς ἦν πρὸς κεφαλῆς Ὀλοφέρνου, καθεῖλε τὸν ἀκινάνην αὐτοῦ ἀπʼ αὐτοῦ.
\vs{7}Καὶ ἐγγίσασα τῆς κλίνης, ἐδράξατο τῆς κόμης τῆς κεφαλῆς αὐτοῦ, καὶ εἶπε, κραταίωσόν με, ὁ Θεὸς Ἰσραὴλ, ἐν τῇ ἡμέρᾳ ταύτῃ·

\vs{8}Καὶ ἐπάταξεν εἰς τὸν τράχηλον αὐτοῦ δὶς ἐν τῇ ἰσχύϊ αὐτῆς, καὶ ἀφεῖλε τὴν κεφαλὴν αὐτοῦ ἀπʼ αὐτοῦ,
\vs{9}καὶ ἀπεκύλισε τὸ σῶμα αὐτοῦ ἀπὸ τῆς στρωμνῆς, καὶ ἀφεῖλε τὸ κωνωπεῖον ἀπὸ τῶν στύλων· καὶ μετʼ ὀλίγον ἐξῆλθε, καὶ παρέδωκε τῇ ἅβρᾳ αὐτῆς τὴν κεφαλὴν Ὀλοφέρνου.
\vs{10}Καὶ ἐνέβαλεν αὐτὴν εἰς τὴν πήραν τῶν βρωμάτων αὐτῆς, καὶ ἐξῆλθον αἱ δύο ἅμα κατὰ τὸν ἐθισμὸν αὐτῶν· καὶ διελθοῦσαι τὴν παρεμβολὴν, ἐκύκλωσαν τὴν φάραγγα ἐκείνην, καὶ προσανέβησαν τὸ ὄρος Βετυλούα, καὶ ἤλθοσαν πρὸς τὰς πύλας αὐτῆς.

\vs{11}Καὶ εἶπεν Ἰουδὶθ μακρόθεν τοῖς φυλάσσουσιν ἐπὶ τῶν πυλῶν, ἀνοίξατε, ἀνοίξατε δὴ τὴν πύλην, μεθʼ ἡμῶν ὁ Θεὸς ὁ Θεὸς ἡμῶν, ποιῆσαι ἔτι ἰσχὺν ἐν Ἰσραὴλ καὶ κράτος κατὰ τῶν ἐχθρων, καθὰ καὶ σήμερον ἐποίησε.

\vs{12}Καὶ ἐγένετο ὡς ἤκουσαν οἱ ἄνδρες τῆς πόλεως αὐτῆς τὴν φωνὴν αὐτῆς, ἐσπούδασαν τοῦ καταβῆναι εἰς τὴν πύλην τῆς πόλεως αὐτῶν· καὶ συνεκάλεσαν τοὺς πρεσβυτέρους τῆς πόλεως.
\vs{13}Καὶ συνέδραμον πάντες ἀπὸ μικροῦ ἕως μεγάλου, ὅτι παράδοξον ἦν αὐτοῖς τὸ ἐλθεῖν αὐτὴν, καὶ ἤνοιξαν, τὴν πύλην, καὶ ὑπεδέξαντο αὐτάς· καὶ ἅψαντες πῦρ εἰς φαῦσιν, περιεκύκλωσαν αὐτάς.

\vs{14}Ἡ δὲ εἶπε πρὸς αὐτοὺς φωνῇ μεγάλῃ, αἰνεῖτε τὸν Θεόν, αἰνεῖτε· αἰνεῖτε τὸν Θεὸν ὃς οὐκ ἀπέστησε τὸ ἔλεος αὐτοῦ ἀπὸ τοῦ οἴκου Ἰσραὴλ, ἀλλʼ ἔθραυσε τοὺς ἐχθροὺς ἡμῶν διὰ χειρός μου ἐν τῇ νυκτὶ ταύτῃ.
\vs{15}Καὶ προελοῦσα τὴν κεφαλὴν ἐκ τῆς πήρας, ἔδειξε, καὶ εἶπεν αὐτοῖς, ἰδοὺ ἡ κεφαλὴ Ὀλοφέρνου ἀρχιστρατήγου δυνάμεως Ἀσσοὺρ, καὶ ἰδοὺ τὸ κωνωπεῖον ἐν ᾧ κατέκειτο ἐν ταῖς μέθαις αὐτοῦ, καὶ ἐπάταξεν αὐτὸν ὁ Κύριος ἐν χειρὶ θηλείας.
\vs{16}Καὶ ζῇ Κύριος ὃς διεφύλαξέ με ἐν τῇ ὁδῷ μου ᾗ ἐπορεύθην, ὅτι ἠπάτησεν αὐτὸν τὸ πρόσωπόν μου εἰς ἀπώλειαν αὐτοῦ, καὶ οὐκ ἐποίησεν ἁμάρτημα μετʼ ἐμοῦ εἰς μίασμα καὶ αἰσχύνην.

\vs{17}Καὶ ἐξέστη πᾶς ὁ λαὸς σφόδρα, καὶ κύψαντες προσεκύνησαν τῷ Θεῷ καὶ εἶπαν ὁμοθυμαδὸν, εὐλογητὸς εἶ ὁ Θεὸς ἡμῶν, ὁ ἐξουδενώσας ἐν τῇ ἡμέρᾳ τῇ σήμερον τοὺς ἐχθροὺς τοῦ λαοῦ σου.
\vs{18}Καὶ εἶπεν αὐτῇ Ὀζίας, εὐλογητὴ σὺ, θυγάτηρ τῷ Θεῷ τῷ ὑψίστῳ παρὰ πάσας τὰς γυναῖκας τὰς ἐπὶ τῆς γῆς, καὶ εὐλογημένος Κύριος ὁ Θεὸς, ὃς ἔκτισε τοὺς οὐρανοὺς καὶ τὴν γῆν, ὃς κατεύθυνέ σε εἰς τραῦμα κεφαλῆς ἄρχοντος ἐχθρῶν ἡμῶν,
\vs{19}ὅτι οὐκ ἀποστήσεται ἡ ἐλπίς σου ἀπὸ καρδίας ἀνθρώπων μνημονευόντων ἰσχὺν Θεοῦ ἕως αἰῶνος.
\vs{20}Καὶ ποιήσαι σοι αὐτὰ ὁ Θεὸς εἰς ὕψος αἰώνιον, τοῦ ἐπισκέψασθαί σε ἐν ἀγαθοῖς, ἀνθʼ ὧν οὐκ ἐφείσω τῆς ψυχῆς σου διὰ τὴν ταπείνωσιν τοῦ γένους ἡμῶν, ἀλλʼ ἐπεξῆλθες πτώματι ἡμῶν, ἐπʼ εὐθείαν πορευθεῖσα ἐνώπιον τοῦ Θεοῦ ἡμῶν· καὶ εἶπαν πᾶς ὁ λαός, γένοιτο, γένοιτο.

\ch{14}
Καὶ εἶπε πρὸς αὐτοὺς Ἰουδὶθ, ἀκούσατε δή μου, ἀδελφοὶ, καὶ λαβόντες τὴν κεφαλὴν ταύτην, κρεμάσατε αὐτὴν ἐπὶ τῆς ἐπάλξεως τοῦ τείχους ὑμῶν.

\vs{2}Καὶ ἔσται ἡνίκα ἂν διαφαύσῃ ὁ ὄρθρος, καὶ ἐξέλθῃ ὁ ἥλιος ἐπὶ τὴν γῆν, ἀναλήψεσθε ἕκαστος τὰ σκεύη τὰ πολεμικὰ ὑμῶν, καὶ ἐξελεύσεσθε πᾶς ἀνὴρ ἰσχύων ἔξω τῆς πόλεως, καὶ δώσετε ἀρχηγὸν εἰς αὐτοὺς, ὡς καταβαίνοντες ἐπὶ τὸ πεδίον εἰς τὴν προφυλακὴν υἱῶν Ἀσσοὺρ, καὶ οὐ καταβήσεσθε.
\vs{3}Καὶ ἀναλαβόντες οὗτοι τὰς πανοπλίας αὐτῶν, πορεύσονται εἰς τὴν παρεμβολὴν αὐτῶν, καὶ ἐγεροῦσι τοὺς στρατηγοὺς τῆς δυνάμεως Ἀσσοὺρ, καὶ συνδραμοῦνται ἐπὶ τὴν σκηνὴν Ὀλοφέρνου, καὶ οὐχ εὑρήσουσιν αὐτὸν, καὶ ἐπιπεσεῖται ἐπʼ αὐτοὺς φόβος, καὶ φεύξονται ἀπὸ προσώπου ὑμῶν.
\vs{4}Καὶ ἐπακολουθήσαντες ὑμεῖς, καὶ πάντες οἱ κατοικοῦντες πᾶν ὅριον Ἰσραὴλ, καταστρώσατε αὐτοὺς ἐν ταῖς ὁδοῖς αὐτῶν.
\vs{5}Πρὸ δὲ τοῦ ποιῆσαι ταῦτα, καλέσατέ μοι Ἀχιὼρ τὸν Ἀμμανίτην, ἵνα ἰδὼν ἐπιγνῷ τὸν ἐκφαυλίσαντα τὸν οἶκον τοῦ Ἰσραὴλ, καὶ αὐτὸν ὡς εἰς θάνατον ἀποστείλαντα εἰς ἡμᾶς.

\vs{6}Καὶ ἐκάλεσαν τὸν Ἀχιὼρ ἐκ τοῦ οἴκου Ὀζία· ὡς δὲ ἦλθε, καὶ εἶδε τὴν κεφαλὴν Ὀλοφέρνου ἐν χειρὶ ἀνδρὸς ἑνὸς ἐν τῇ ἐκκλησίᾳ τοῦ λαοῦ, ἔπεσεν ἐπὶ πρόσωπον, καὶ ἐξελύθη τὸ πνεῦμα αὐτοῦ.

\vs{7}Ὡς δὲ ἀνέλαβον αὐτὸν, προσέπεσε τοῖς ποσὶν Ἰουδὶθ, καὶ προσεκύνησε τῷ προσώτῳ αὐτῆς, καὶ εἶπεν, εὐλογημένη σὺ ἐν παντὶ σκηνώματι Ἰούδα, καὶ ἐν παντὶ ἔθνει, οἵτινες ἀκούσαντες τὸ ὄνομά σου ταραχθήσονται.
\vs{8}Καὶ νῦν ἀνάγγειλόν μοι ὅσα ἐποίησας ἐν ταῖς ἡμέραις ταύταις· καὶ ἀπήγγειλεν αὐτῷ Ἰουδὶθ ἐν μέσῳ τοῦ λαοῦ πάντα ὅσα ἦν πεποιηκυῖα, ἀφʼ ἧς ἡμέρας ἐξῆλθεν ἕως οὗ ἐλάλει αὐτοῖς.
\vs{9}Ὡς δὲ ἐπαύσατο λαλοῦσα, ἠλάλαξεν ὁ λαὸς φωνῇ μεγάλῃ, καὶ ἔδωκε φωνὴν εὐφρόσυνον ἐν τῇ πόλει αὐτῶν.

\vs{10}Ἰδὼν δὲ Ἀχιὼρ πάντα ὅσα ἐποίησεν ὁ Θεὸς τοῦ Ἰσραὴλ, ἐπίστευσε τῷ Θεῷ σφόδρα, καὶ περιετέμετο τὴν σάρκα τῆς ἀκροβυστίας αὐτοῦ, καὶ προσετέθη πρὸς τὸν οἶκον Ἰσραὴλ ἕως τῆς ἡμέρας ταύτης.

\vs{11}Ἡνίκα δὲ ὁ ὄρθρος ἀνέβη, καὶ ἐκρέμασαν τὴν κεφαλὴν Ὀλοφέρνου ἐκ τοῦ τείχους, καὶ ἀνέλαβε πᾶς ἀνὴρ Ἰσραὴλ τὰ ὅπλα αὐτοῦ, καὶ ἐξήλθοσαν κατὰ σπείρας ἐπὶ τὰς ἀναβάσεις τοῦ ὄρους.

\vs{12}Οἱ δὲ υἱοὶ Ἀσσοὺρ, ὡς εἶδον αὐτοὺς, διέπεμψαν ἐπὶ τοὺς ἡγουμένους αὐτῶν· οἱ δὲ ἦλθον ἐπὶ στρατηγοὺς καὶ χιλιάρχους καὶ ἐπὶ πάντα ἄρχοντα αὐτῶν·

\vs{13}Καὶ παρεγένοντο ἐπὶ τὴν σκηνὴν Ὀλοφέρνου, καὶ εἶπαν τῷ ὄντι ἐπὶ πάντων τῶν αὐτοῦ, ἔγειρον δὴ τὸν κύριον ἡμῶν, ὅτι ἐτόλμησαν οἱ δοῦλοι καταβαίνειν ἐφʼ ἡμᾶς εἰς πόλεμον, ἵνα ἐξολοθρευθῶσιν εἰς τέλος.

\vs{14}Καὶ εἰσῆλθε Βαγώας, καὶ ἔκρουσε τὴν αὐλαίαν τῆς σκηνῆς· ὑπενοεῖτο γὰρ καθεύδειν αὐτὸν μετὰ Ἰουδίθ.
\vs{15}Ὡς δὲ οὐδεὶς ἐπήκουσε, διαστείλας εἰσῆλθεν εἰς τὸν κοιτῶνα, καὶ εὗρεν αὐτὸν ἐπὶ τῆς χελωνίδος ἐῤῥιμεμένον νεκρὸν, καὶ ἡ κεφαλὴ αὐτοῦ ἀφῄρητο ἀπʼ αὐτοῦ.
\vs{16}Καὶ ἐβόησε φωνῇ μεγάλῃ μετὰ κλαυθμοῦ καὶ στεναγμοῦ καὶ βοῆς ἰσχυρᾶς, καὶ διέῤῥηξε τὰ ἱμάτια αὐτοῦ.

\vs{17}Καὶ εἰσῆλθεν εἰς τὴν σκηνὴν οὗ ἦν Ἰουδὶθ καταλύουσα, καὶ οὐχ εὗρεν αὐτήν· καὶ ἐξεπήδησεν εἰς τὸν λαὸν, κράζων,
\vs{18}ἠθέτησαν οἱ δοῦλοι, ἐποίησεν αἰσχύνην μία γυνὴ τῶν Ἑβραίων εἰς τὸν οἶκον τοῦ βασιλέως Ναβουχοδονόσορ, ὅτι ἰδοὺ Ὀλοφέρνης χαμαί, καὶ ἡ κεφαλὴ οὐκ ἔστιν ἐπʼ αὐτῷ.
\vs{19}Ὡς δὲ ἤκουσαν ταῦτα τὰ ῥήματα οἱ ἄρχοντες τῆς δυνάμεως Ἀσσοὺρ, τοὺς χιτῶνας αὐτῶν διέῤῥηξαν, καὶ ἐταράχθη ἡ ψυχὴ αὐτῶν σφόδρα, καὶ ἐγένετο αὐτῶν κραυγὴ καὶ βοὴ μεγάλη σφόδρα ἐν μέσῳ τῆς παρεμβολῆς.

\ch{15}
Καὶ ὡς ἤκουσαν οἱ ἐν τοῖς σκηνώμασιν ὄντες, ἐξέστησαν ἐπὶ τὸ γεγονός·
\vs{2}καὶ ἐπέπεσεν ἐπʼ αὐτοὺς φόβος καὶ τρόμος, καὶ οὐκ ἦν ἄνθρωπος μένων κατὰ πρόσωπον τοῦ πλησίον ἔτι, ἀλλʼ ἐκχυθέντες ὁμοθυμαδὸν ἔφευγον ἐπὶ πᾶσαν ὁδὸν τοῦ πεδίου καὶ τῆς ὀρεινῆς.
\vs{3}Καὶ οἱ παρεμβεβληκότες ἐν τῇ ὀρεινῇ κύκλῳ Βετυλούα καὶ ἐτράπησαν εἰς φυγήν· καὶ τότε οἱ υἱοὶ Ἰσραὴλ πᾶς ἀνὴρ πολεμιστὴς ἐξ αὐτῶν ἐξεχύθησαν ἐπʼ αὐτούς.

\vs{4}Καὶ ἀπέστειλεν Ὀζίας εἰς Βαιτομασθαῒμ, καὶ Χωβαῒ, καὶ Χωλὰ, καὶ εἰς πᾶν ὅριον Ἰσραὴλ, τοὺς ἀπαγγέλλοντας ὑπὲρ τῶν συντετελεσμένων, καὶ ἵνα πάντες ἐπεκχυθῶσι τοῖς πολεμίοις εἰς τὴν ἀναίρεσιν αὐτῶν.
\vs{5}Ὡς δὲ ἤκουσαν οἱ υἱοὶ Ἰσραὴλ, πάντες ὁμοθυμαδὸν ἐπέπεσον ἐπʼ αὐτοὺς, καὶ ἔκοπτον αὐτοὺς ἕως Χωβά. ὡσαύτως δὲ καὶ οἱ ἐξ Ἱερουσαλὴμ παρεγενήθησαν καὶ ἐκ πάσης τῆς ὀρεινῆς· ἀνήγγειλαν γὰρ αὐτοῖς τὰ γεγονότα τῇ παρεμβολῇ τῶν ἐχθρῶν αὐτῶν· καὶ οἱ ἐν Γαλαὰδ καὶ οἱ ἐν τῇ Γαλιλαίᾳ ὑπερεκέρασαν αὐτοὺς πληγῇ μεγάλῃ, ἕως οὗ παρῆλθον Δαμασκὸν, καὶ τὰ ὅρια αὐτῆς.

\vs{6}Οἱ δὲ λοιποὶ οἱ κατοικοῦντες Βετυλούα, ἐπέπεσαν τῇ παρεμβολῇ Ἀσσοὺρ, καὶ ἐπρονόμευσαν αὐτοὺς, καὶ ἐπλούτησαν σφόδρα.
\vs{7}Οἱ δὲ υἱοὶ Ἰσραὴλ ἀναστρέψαντες ἀπὸ τῆς κοπῆς, ἐκυρίευσαν τῶν λοιπῶν, καὶ αἱ κῶμαι καὶ ἐπαύλεις ἐν τῇ ὀρείνῃ καὶ πεδινῇ ἐκράτησαν πολλῶν λαφύρων· ἦν γὰρ πλῆθος πολὺ σφόδρα.

\vs{8}Καὶ Ἰωακὶμ ὁ ἱερεὺς ὁ μέγας καὶ ἡ γερουσία τῶν υἱῶν Ἰσραὴλ οἱ κατοικοῦντες ἐν Ἱερουσαλὴμ ἦλθον τοῦ θεάσασθαι τὰ ἀγαθὰ ἃ ἐποίησε Κύριος τῷ Ἰσραήλ, καὶ τοῦ ἰδεῖν τὴν Ἰουδὶθ, καὶ λαλῆσαι μετʼ αὐτῆς εἰρήνην.
\vs{9}Ὡς δὲ εἰσῆλθον πρὸς αὐτὴν, εὐλόγησαν αὐτὴν πάντες ὁμοθυμαδὸν, καὶ εἶπαν πρὸς αὐτὴν, σὺ ὕψωμα Ἰσραὴλ, σὺ γαυρίαμα μέγα τοῦ Ἰσραὴλ, σὺ καύχημα μέγα τοῦ γένους ἡμῶν.
\vs{10}Ἐποίησας πάντα ταῦτα ἐν χειρί σου, ἐποίησας τὰ ἀγαθὰ μετὰ Ἰσραήλ· καὶ εὐδοκήσαι ἐπʼ αὐτοῖς ὁ Θεός· εὐλογημένη γίνου παρὰ τῷ παντοκράτορι Κυρίῳ εἰς τὸν αἰῶνα χρόνον· καὶ εἶπε πᾶς ὁ λαὸς, γένοιτο.

\vs{11}Καὶ ἐλαφύρευσε πᾶς ὁ λαὸς τὴν παρεμβολὴν ἐφʼ ἡμέρας τριάκοντα, καὶ ἔδωκαν τῇ Ἰουδὶθ τὴν σκηνὴν Ὀλοφέρνου, καὶ πάντα τὰ ἀργυρώματα, καὶ τὰς κλίνας καὶ τὰ ὅλκια· καὶ πάντα τὰ σκευάσματα αὐτοῦ· καὶ λαβοῦσα αὕτη ἐπέθηκεν ἐπὶ τὴν ἡμίονον αὐτῆς, καὶ ἔζευξε τὰς ἁμάξας αὐτῆς, καὶ ἐσώρευσεν αὐτὰ ἐπʼ αὐτῶν.

\vs{12}Καὶ συνέδραμε πᾶσα γυνὴ Ἰσραὴλ τοῦ ἰδεῖν αὐτὴν, καὶ εὐλόγησαν αὐτήν· καὶ ἐποίησαν αὐτῇ χορὸν ἐξ αὐτῶν· καὶ ἔλαβε θύρσους ἐν ταῖς χερσὶν αὐτῆς, καὶ ἔδωκε ταῖς γυναιξὶ ταῖς μετʼ αὐτῆς,
\vs{13}καὶ ἐστεφανώσαντο τὴν ἐλαίαν αὕτη καὶ αἱ μετʼ αὐτῆς· καὶ προῆλθε παντὸς τοῦ λαοῦ ἐν χορείᾳ ἡγουμένη πασῶν τῶν γυναικῶν, καὶ ἠκολούθει πᾶς ἀνὴρ Ἰσραὴλ ἐνωπλισμένοι μετὰ στεφάνων καὶ ὕμνων ἐν τῷ στόματι αὐτῶν.

\ch{16}
Καὶ ἐξῆρχεν Ἰουδὶθ τὴν ἐξομολόγησιν ταύτην ἐν παντὶ Ἰσραὴλ, καὶ ὑπεφώνει πᾶς ὁ λαὸς τὴν αἴνεσιν ταύτην·

\vs{2}Καὶ εἶπεν Ἰουδὶθ,

Ἐξάρχετε τῷ Θεῷ μου ἐν τυμπάνοις, ᾄσατε τῷ Κυρίῳ μου ἐν κυμβάλοις, ἐναρμόσασθε αὐτῷ ψαλμὸν καινὸν, ὑψοῦτε καὶ ἐπικαλέσασθε τὸ ὄνομα αὐτοῦ.
\vs{3}ὅτι Θεὸς συντρίβων πολέμους Κύριος, ὅτι εἰς παρεμβολὰς αὐτοῦ ἐν μέσῳ λαοῦ ἐξείλατό με ἐκ χειρὸς τῶν καταδιωκόντων με.

\vs{4}Ἦλθεν Ἀσσοὺρ ἐξ ὀρέων ἀπὸ βοῤῥᾶ, ἦλθεν ἐν μυριάσι δυνάμεως αὐτοῦ, ὧν τὸ πλῆθος αὐτῶν ἐνέφραξε χειμάῤῥους, καὶ ἡ ἵππος αὐτῶν ἐκάλυψε βουνούς.
\vs{5}Εἶπεν ἐμπρήσειν τὰ ὅριά μου, καὶ τοὺς νεανίσκους μου ἀνελεῖν ἐν ῥομφαίᾳ, καὶ τὰ θηλάζοντά μου θήσειν εἰς ἔδαφος, καὶ τὰ νήπιά μου δώσειν εἰς προνομὴν, καὶ τὰς παρθένους μου σκυλεῦσαι.

\vs{6}Κύριος Παντοκράτωρ ἠθέτησεν αὐτοὺς ἐν χειρὶ θηλείας.
\vs{7}Οὐ γὰρ ὑπέπεσεν ὁ δυνατὸς αὐτῶν ὑπὸ νεανίσκων, οὐδὲ υἱοὶ Τιτάνων ἐπάταξαν αὐτόν, οὐδὲ ὑψηλοὶ γίγαντες ἐπέθεντο αὐτῷ, ἀλλὰ Ἰουδὶθ θυγάτηρ Μεραρεὶ ἐν κάλλει προσώπου αὐτῆς παρέλυσεν αὐτόν.
\vs{8}Ἐξεδύσατο γὰρ στολὴν χηρεύσεως αὐτῆς εἰς ὕψος τῶν πονούντων ἐν Ἰσραήλ, ἠλείψατο τὸ πρόσωπον αὐτῆς ἐν μυρισμῷ, καὶ ἐδήσατο τὰς τρίχας αὐτῆς ἐν μίτρᾳ, καὶ ἔλαβε στολὴν λινῆν εἰς ἀπάτην αὐτοῦ.
\vs{9}Τὸ σανδάλιον αὐτῆς ἥρπασεν ὀφθαλμὸν αὐτοῦ, καὶ τὸ κάλλος αὐτῆς ᾐχμαλώτισε ψυχὴν αὐτοῦ· διῆλθεν ὁ ἀκινάκης τὸν τράχηλον αὐτοῦ.

\vs{10}Ἔφριξαν Πέρσαι τὴν τόλμαν αὐτῆς, καὶ Μῆδοι τὸ θράσος αὐτῆς ἐῤῥάχθησαν.
\vs{11}Τότε ἠλάλαξαν οἱ ταπεινοί μου, καὶ ἐφοβήθησαν οἱ ἀσθενοῦντές μου, καὶ ἐπτοήθησαν· ὕψωσαν τὴν φωνὴν αὐτῶν, καὶ ἀνετράπησαν.
\vs{12}Υἱοὶ κορασίων κατεκέντησαν αὐτοὺς, καὶ ὡς παῖδας αὐτομολούντων ἐτίτρωσκον αὐτούς· ἀπώλοντο ἐκ παρατάξεως Κυρίου μου.

\vs{13}Ὑμνήσω τῷ Θεῷ μου ὕμνον καινόν· Κύριε, μέγας εἶ, καὶ ἔνδοξος, θαυμαστὸς ἐν ἰσχύϊ, ἀνυπέρβλητος.
\vs{14}Σοὶ δουλευσάτω πᾶσα ἡ κτίσις σου, ὅτι εἶπας, καὶ ἐγενήθησαν· ἀπέστειλας τὸ πνεῦμά σου, καὶ ᾠκοδόμησε· καὶ οὐκ ἔστιν ὃς ἀντιστήσεται τῇ φωνῇ σου.
\vs{15}Ὄρη γὰρ ἐκ θεμελίων σὺν ὕδασι σαλευθήσεται, πέτραι δὲ ἀπὸ προσώπου σου ὡς κηρὸς τακήσονται, ἐπὶ δὲ τοῖς φοβουμένοις σε, σὺ εὐιλατεύεις αὐτοῖς.
\vs{16}Ὅτι μικρὸν πᾶσα θυσία εἰς ὀσμὴν εὐωδίας, καὶ ἐλάχιστον πᾶν στέαρ εἰς ὁλοκαύτωμά σοι· ὁ δὲ φοβούμενος τὸν Κύριον, μέγας διαπαντός.

\vs{17}Οὐαὶ ἔθνεσιν ἐπανισταμένοις τῷ γένει μου· Κύριος παντοκράτωρ ἐκδικήσει αὐτοὺς ἐν ἡμέρᾳ κρίσεως, δοῦναι πῦρ καὶ σκώληκας εἰς σάρκας αὐτῶν, καὶ κλαύσονται ἐν αἰσθήσει ἕως αἰῶνος.

\vs{18}Ὡς δὲ ἤλθοσαν εἰς Ἱερουσαλήμ, προσεκύνησαν τῷ Θεῷ· καὶ ἡνίκα ἐκαθαρίσθη ὁ λαὸς, ἀνήνεγκαν τὰ ὁλοκαυτώματα αὐτῶν, καὶ τὰ ἑκουσία αὐτῶν, καὶ τὰ δόματα.

\vs{19}Καὶ ἀνέθηκεν Ἰουδὶθ πάντα τὰ σκεύη Ὀλοφέρνου ὅσα ἔδωκεν ὁ λαὸς αὐτῇ, καὶ τὸ κωνωπεῖον ὃ ἔλαβεν αὕτς ἐκ τοῦ κοιτῶνος αὐτοῦ, εἰς ἀνάθημα τῷ Θεῷ ἔδωκε.

\vs{20}Καὶ ἦν ὁ λαὸς εὐφραινόμενος ἐν Ἱερουσαλὴμ κατὰ πρόσωπον τῶν ἁγίων ἐπὶ μῆνας τρεῖς, καὶ Ἰουδὶθ μετʼ αὐτῶν κατέμεινε.

\vs{21}Μετὰ δὲ τὰς ἡμέρας ταύτας ἀνέζευξεν ἕκαστος εἰς τὴν κληρονομίαν αὐτοῦ· καὶ Ἰουδὶθ ἀπῆλθεν εἰς Βετυλούα, καὶ κατέμεινεν ἐπὶ τῆς ὑπάρξεως αὐτῆς· καὶ ἐγένετο κατὰ τὸν καιρὸν αὐτῆς ἔνδοξος ἐν πάσῃ τῇ γῇ.
\vs{22}Καὶ πολλοὶ ἐπεθύμησαν αὐτὴν, καὶ οὐκ ἔγνω ἀνὴρ αὐτὴν πάσας τὰς ἡμέρας τῆς ζωῆς αὐτῆς, ἀφʼ ἧς ἡμέρας ἀπέθανε Μανασσῆς ὁ ἀνὴρ αὐτῆς, καὶ προσετέθη πρὸς τὸν λαὸν αὐτοῦ.

\vs{23}Καὶ ἦν προβαίνουσα μεγάλη σφόδρα· καὶ ἐγήρασεν ἐν τῷ οἴκῳ τοῦ ἀνδρὸς αὐτῆς ἔτη ἑκατὸν πέντε, καὶ ἀφῆκε τὴν ἅβραν αὐτῆς ἐλευθέραν, καὶ ἀπέθανεν εἰς Βετυλούα, καὶ ἔθαψαν αὐτὴν ἐν τῷ σπηλαίῳ τοῦ ἀνδρὸς αὐτῆς Μανασσῆ.
\vs{24}καὶ ἐπένθησεν αὐτὴν οἶκος Ἰσραὴλ ἡμέρας ἑπτά· καὶ διεῖλε τὰ ὑπάρχοντα αὐτῆς πρὸ τοῦ ἀποθανεῖν αὐτὴν, πᾶσι τοῖς ἔγγιστα Μανασσῆ τοῦ ἀνδρὸς αὐτῆς, καὶ τοῖς ἔγγιστα τοῦ γένους αὐτῆς.
\vs{25}Καὶ οὐκ ἦν ἔτι ὁ ἐκφοβῶν τοὺς υἱοὺς Ἰσραὴλ ἐν ταῖς ἡμέραις Ἰουδὶθ, καὶ μετὰ τὸ ἀποθανεῖν αὐτὴν, ἡμέρας πολλάς.


\def\book{ΣΟΦΙΑ ΣΑΛΩΜΩΝ}
\biblebook{ΣΟΦΙΑ ΣΑΛΩΜΩΝ}


\lettrine[lines=2, loversize=0.2, nindent=0em, findent=.25em]{\textcolor{bookheadingcolor}{Ἀ}}{ΓΑΠΗΣΑΤΕ} δικαιοσύνην οἱ κρίνοντες τὴν γῆν, φρονήσατε περὶ τοῦ Κυρίου ἐν ἀγαθότητι, καὶ ἐν ἁπλότητι καρδίας ζητήσατε αὐτόν.
\vs{2}Ὅτι εὑρίσκεται τοῖς μὴ πειράζουσιν αὐτὸν, ἐμφανίζεται δὲ τοῖς μὴ ἀπιστοῦσιν αὐτῷ.
\vs{3}Σκολιοὶ γὰρ λογισμοὶ χωρίζουσιν ἀπὸ Θεοῦ, δοκιμαζομένη τε ἡ δύναμις ἐλέγχει τοὺς ἄφρονας.

\vs{4}Ὅτι εἰς κακότεχνον ψυχὴν οὐκ εἰσελεύσεται σοφία, οὐδὲ κατοικήσει ἐν σώματι κατάχρεῳ ἁμαρτίας.
\vs{5}Ἅγιον γὰρ πνεῦμα παιδείας φεύξεται δόλον, καὶ ἀπαναστήσεται ἀπὸ λογισμῶν ἀσυνέτων, καὶ ἐλεγχθήσεται ἐπελθούσης ἀδικίας.

\vs{6}Φιλάνθρωπον γὰρ πνεῦμα σοφία, καὶ οὐκ ἀθῳώσει βλάσφημον ἀπὸ χειλέων αὐτοῦ, ὅτι τῶν νεφρῶν αὐτοῦ μάρτυς ὁ Θεὸς, καὶ τῆς καρδίας αὐτοῦ ἐπίσκοπος ἀληθὴς, καὶ τῆς γλώσσης ἀκουστής·
\vs{7}ὅτι πνεῦμα Κυρίου πεπλήρωκε τὴν οἰκουμένην, καὶ τὸ συνέχον τὰ πάντα γνῶσιν ἔχει φωνῆς.

\vs{8}Διὰ τοῦτο φθεγγόμενος ἄδικα οὐδεὶς μὴ λάθῃ, οὐδὲ μὴν παροδεύσῃ αὐτὸν ἐλέγχουσα ἡ δίκη.
\vs{9}Ἐν γὰρ διαβουλίοις ἀσεβοῦς ἐξέτασις ἔσται, λόγων δὲ αὐτοῦ ἀκοὴ πρὸς Κύριον ἥξει εἰς ἔλεγχον ἀνομημάτων αὐτοῦ.
\vs{10}Ὅτι οὖς ζηλώσεως ἀκροᾶται τὰ πάντα, καὶ θροῦς γογγυσμῶν οὐκ ἀποκρύπτεται.

\vs{11}Φυλάξασθε τοίνυν γογγυσμὸν ἀνωφελῆ, καὶ ἀπὸ καταλαλιᾶς φείσασθε γλώσσης· ὅτι φθέγμα λαθραῖον κενὸν οὐ πορεύσεται, στόμα δὲ καταψευδόμενον ἀναιρεῖ ψυχήν.

\vs{12}Μὴ ζηλοῦτε θάνατον ἐν πλάνῃ ζωῆς ὑμῶν, μηδὲ ἐπισπᾶσθε ὄλεθρον ἔργοις χειρῶν ὑμῶν·
\vs{13}ὅτι ὁ Θεὸς θάνατον οὐκ ἐποίησεν, οὐδὲ τέρπεται ἐπʼ ἀπωλείᾳ ζώντων.
\vs{14}Ἔκτισε γὰρ εἰς τὸ εἶναι τὰ πάντα, καὶ σωτήριοι αἱ γενέσεις τοῦ κόσμου, καὶ οὐκ ἔστιν ἐν αὐταῖς φάρμακον ὀλέθρου, οὔτε ᾅδου βασίλειον ἐπὶ γῆς.
\vs{15}Δικαιοσύνη γὰρ ἀθάνατός ἐστιν·
\vs{16}ἀσεβεῖς δὲ ταῖς χερσὶ καὶ τοῖς λόγοις προσεκαλέσαντο αὐτὸν, φίλον ἡγησάμεμοι αὐτὸν ἐτάκησαν, καὶ συνθήκην ἔθεντο πρὸς αὐτὸν, ὅτι ἄξιοί εἰσι τῆς ἐκείνου μερίδος εἶναι.

\ch{2}
Εἶπον γὰρ ἑαυτοῖς λογισάμενοι οὐκ ὀρθῶς, ὀλίγος ἐστὶ καὶ λυπηρὸς ὁ βίος ἡμῶν, καὶ οὐκ ἔστιν ἴασις ἐν τελευτῇ ἀνθρώπου, καὶ οὐκ ἐγνώσθη ὁ ἀναλύσας ἐξ ᾅδου.
\vs{2}Ὅτι αὐτοσχεδίως ἐγεννήθημεν, καὶ μετὰ τοῦτο ἐσόμεθα ὡς οὐχ ὑπάρξαντες, ὅτι καπνὸς ἡ πνοὴ ἐν ῥισὶν ἡμῶν, καὶ ὁ λόγος σπινθὴρ ἐν κινήσει καρδίας ἡμῶν,
\vs{3}οὗ σβεσθέντος τέφρα ἀποβήσεται τὸ σῶμα, καὶ τὸ πνεῦμα διαχυθήσεται ὡς χαῦνος ἀήρ.
\vs{4}Καὶ τὸ ὄνομα ἡμῶν ἐπιλησθήσεται ἐν χρόνῳ, καὶ οὐθεὶς μνημονεύσει τῶν ἔργων ἡμῶν· καὶ παρελεύσεται ὁ βίος ἡμῶν ὡς ἴχνη νεφέλης, καὶ ὡς ὁμίχλη διασκεδασθήσεται διωχθεῖσα ὑπὸ ἀκτίνων ἡλίου, καὶ ὑπὸ θερμότητος αὐτοῦ βαρυνθεῖσα.

\vs{5}Σκιᾶς γὰρ πάροδος ὁ βίος ἡμῶν, καὶ οὐκ ἔστιν ἀναποδισμὸς τῆς τελευτῆς ἡμῶν, ὅτι κατεσφραγίσθη, καὶ οὐδεὶς ἀναστρέφει.

\vs{6}Δεῦτε οὖν καὶ ἀπολαύσωμεν τῶν ὄντων ἀγαθῶν, καὶ χρησώμεθα τῇ κτίσει ὡς νεότητι σπουδαίως.
\vs{7}Οἴνου πολυτελοῦς καὶ μύρων πλησθῶμεν, καὶ μὴ παροδευσάτω ἡμᾶς ἄνθος ἀέρος.
\vs{8}Στεψώμεθα ῥόδων κάλυξι πρινὴ μαρανθῆναι.
\vs{9}Μηδεὶς ἡμῶν ἄμοιρος ἔστω τῆς ἡμετέρας ἀγερωχίας, πανταχῆ καταλίπωμεν σύμβολα τῆς εὐφροσύνης, ὅτι αὕτη ἡ μερὶς ἡμῶν καὶ ὁ κλῆρος οὗτος.

\vs{10}Καταδυναστεύσωμεν πένητα δίκαιον, μὴ φεισώμεθα χήρας, μηδὲ πρεσβύτου ἐντραπῶμεν πολιὰς πολυχρονίους.
\vs{11}Ἔστω δὲ ἡμῶν ἡ ἰσχὺς νόμος τῆς δικαιοσύνης, τὸ γὰρ ἀσθενὲς ἄχρηστον ἐλέγχεται.

\vs{12}Ἐνεδρεύσωμεν δὲ τὸν δίκαιον, ὅτι δύσχρηστος ἡμῖν ἐστι καὶ ἐναντιοῦται τοῖς ἔργοις ἡμῶν, καὶ ὀνειδίζει ἡμῖν ἁμαρτήματα νόμου, καὶ ἐπιφημίζει ἡμῖν ἁμαρτήματα παιδείας ἡμῶν.
\vs{13}Ἐπαγγέλλεται γνῶσιν ἔχειν Θεοῦ, καὶ παῖδα Κυρίου ἑαυτὸν ὀνομάζει.
\vs{14}Ἐγένετο ἡμῖν εἰς ἔλεγχον ἐννοιῶν ἡμῶν.
\vs{15}Βαρύς ἐστιν ἡμῖν καὶ βλεπόμενος, ὅτι ἀνόμοιος τοῖς ἄλλοις ὁ βίος αὐτοῦ, καὶ ἐξηλλαγμέναι αἱ τρίβοι αὐτοῦ.
\vs{16}Εἰς κίβδηλον ἐλογίσθημεν αὐτῷ, καὶ ἀπέχεται τῶν ὁδῶν ἡμῶν ὡς ἀπὸ ἀκαθαρσιῶν· μακαρίζει ἔσχατα δικαίων, καὶ ἀλαζονεύεται πατέρα Θεόν.

\vs{17}Ἴδωμεν εἰ οἱ λόγοι αὐτοῦ ἀληθεῖς, καὶ πειράσωμεν τὰ ἐν ἐκβάσει αὐτοῦ.
\vs{18}Εἰ γάρ ἐστιν ὁ δίκαιος υἱὸς Θεοῦ, ἀντιλήμψεται αὐτοῦ, καὶ ῥύσεται αὐτὸν ἐκ χειρὸς ἀνθεστηκότων.
\vs{19}Ὕβρει καὶ βασάνῳ ἐτάσωμεν αὐτόν, ἵνα γνῶμεν τὴν ἐπιείκειαν αὐτοῦ, καὶ δοκιμάσωμεν τὴν ἀνεξικακίαν αὐτοῦ.
\vs{20}Θανάτῳ ἀσχήμονι καταδικάσωμεν αὐτόν· ἔσται γὰρ αὐτοῦ ἐπισκοπὴ ἐκ λόγων αὐτοῦ.

\vs{21}Ταῦτα ἐλογίσαντο, καὶ ἐπλανήθησαν· ἀπετύφλωσε γὰρ αὐτοὺς ἡ κακία αὐτῶν, καὶ οὐκ ἔγνωσαν μυστήρια
\vs{22}Θεοῦ, οὐδὲ μισθὸν ἤλπισαν ὁσιότητος, οὐδὲ ἔκριναν γέρας ψυχῶν ἀμώμων.

\vs{23}Ὅτι ὁ Θεὸς ἔκτισε τὸν ἄνθρωπον ἐπ ἀφθαρσίᾳ, καὶ εἰκόνα τῆς ἰδίας ἰδιότητος ἐποίησεν ἀτόν.
\vs{24}Φθόνῳ δὲ διαβόλου θάνατος εἰσῆλθεν εἰς τὸν κόσμον· πειράζουσι δὲ αὐτὸν οἱ τῆς ἐκείνου μερίδος ὄντες.

\ch{3}
Δίκαιων δὲ ψυχαὶ ἐν χειρὶ Θεοῦ, καὶ οὐ μὴ ἅψηται αὐτῶν βάσανος.
\vs{2}Ἔδοξαν ἐν ὀφθαλμοῖς ἀφρόνων τεθνάναι,
\vs{3}καὶ ἐλογίσθη κάκωσις ἡ ἔξοδος αὐτῶν, καὶ ἡ ἀφʼ ἡμῶν πορεία σύντριμμα· οἱ δὲ εἰσιν ἐν εἰρήνῃ·
\vs{4}Καὶ γὰρ ἐν ὄψει ἀνθρώπων ἐὰν κολασθῶσιν, ἡ ἐλπὶς αὐτῶν ἀθανασίας πλήρης.

\vs{5}Καὶ ὀλίγα παιδευθέντες μεγάλα εὐεργετηθήσονται, ὅτι ὁ Θεὸς ἐπείρασεν αὐτοὺς, καὶ εὗρεν αὐτοὺς ἀξίους ἑαυτοῦ.
\vs{6}Ὡς χρυσὸν ἐν χωνευτηρίῳ ἐδοκίμασεν αὐτούς, καὶ ὡς ὁλοκάρπωμα θυσίας προσεδέξατο αὐτούς.

\vs{7}Καὶ ἐν καιρῷ ἐπισκοπῆς αὐτῶν ἀναλάμψουσι, καὶ ὡς σπινθῆρες ἐν καλάμῃ διαδραμοῦνται.
\vs{8}Κρινοῦσιν ἔθνη καὶ κρατήσουσι λαῶν, καὶ βασιλεύσει αὐτῶν Κύριος εἰς τοὺς αἰῶνας.
\vs{9}Οἱ πεποιθότες ἐπʼ αὐτῷ συνήσουσιν ἀλήθειαν, καὶ οἱ πιστοὶ ἐν ἀγάπῃ προσμενοῦσιν αὐτῷ, ὅτι χάρις καὶ ἔλεος τοῖς ἐκλεκτοῖς αὐτοῦ.
\vs{10}Οἱ δὲ ἀσεβεῖς καθὰ ἐλογίσαντο ἕξουσιν ἐπιτιμίαν, οἱ ἀμελήσαντες τοῦ δικαίου καὶ τοῦ Κυρίου ἀποστάντες.

\vs{11}Σοφίαν γὰρ καὶ παιδείαν ὁ ἐξουθενῶν ταλαίπωρος, καὶ κενὴ ἡ ἐλπὶς αὐτῶν, καὶ οἱ κόποι ἀνόνητοι, καὶ ἄχρηστα τὰ ἔργα αὐτῶν.
\vs{12}Αἱ γυναῖκες αὐτῶν ἄφρονες, καὶ πονηρὰ τὰ τέκνα αὐτῶν.

\vs{13}Ἐπικατάρατος ἡ γένεσις αὐτῶν, ὅτι μακαρία στεῖρα ἡ ἀμίαντος, ἥτις οὐκ ἔγνω κοίτης ἐν παραπτώματι, ἕξει καρπὸν ἐν ἐπισκοπῇ ψυχῶν.

\vs{14}Καὶ εὐνοῦχος ὁ μὴ ἐργασάμενος ἐν χειρὶ ἀνόμημα, μηδὲ ἐνθυμηθεὶς κατὰ τοῦ Κυρίου πονηρά· δοθήσεται γὰρ αὐτῷ τῆς πίστεως χάρις ἐκλεκτὴ, καὶ κλῆρος ἐν ναῷ Κυρίου θυμηρέστερος·
\vs{15}Ἀγαθῶν γὰρ πόνων καρπὸς εὐκλεὴς, καὶ ἀδιάπτωτος ἡ ῥίζα τῆς φρονήσεως.

\vs{16}Τέκνα δὲ μοιχῶν ἀτέλεστα ἔσται, καὶ ἐκ παρανόμου κοίτης σπέρμα ἀφανισθήσεται.
\vs{17}Ἐάν τε γὰρ μακρόβιοι γένωνται, εἰς οὐθὲν λογισθήσονται, καὶ ἄτιμον ἐπʼ ἐσχάτον τὸ γῆρας αὐτῶν.
\vs{18}Ἐάν τε ὀξέως τελευτήσωσιν, οὐχ ἕχουσιν ἐλπίδα, οὐδὲ ἐν ἡμέρᾳ διαγνώσεως παραμύθιον·
\vs{19}γενεᾶς γὰρ ἀδίκου χαλεπὰ τὰ τέλη.

\ch{4}
Κρείσσων ἀτεκνία μετὰ ἀρετῆς, ἀθανασία γάρ ἐστιν ἐν μνήμῃ αὐτῆς, ὅτι καὶ παρὰ Θεῷ γινώσκεται καὶ παρὰ ἀνθρώποις·
\vs{2}παροῦσάν τε μιμοῦνται αὐτὴν, καὶ ποθοῦσιν ἀπελθοῦσαν καὶ ἐν τῷ αἰῶνι στεφανηφοροῦσα πομπεύει, τὸν τῶν ἀμιάντων ἄθλων ἀγῶνα νικήσασα.

\vs{3}Πολύγονον δὲ ἀσεβῶν πλῆθος οὐ χρησιμεύσει, καὶ ἐκ νόθων μοσχευμάτων οὐ δώσει ῥίζαν εἰς βάθος, οὐδὲ ἀσφαλῆ βάσιν ἑδράσει.
\vs{4}Κᾂν γὰρ ἐν κλάδοις πρὸς καιρὸν ἀναθάλῃ, ἐπισφαλῶς βεβηκότα ὑπὸ ἀνέμου σαλευθήσεται, καὶ ὑπὸ βίας ἀνέμων ἐκριζωθήσεται.
\vs{5}Περικλασθήσονται κλῶνες ἀτέλεστοι, καὶ ὁ καρπὸς αὐτῶν ἄχρηστος, ἄωρος εἰς βρῶσιν, καὶ εἰς οὐθὲν ἐπιτήδειος.
\vs{6}Ἐκ γὰρ ἀνόμων ὕπνων τέκνα γεννώμενα μάρτυρές εἰσι πονηρίας κατὰ γονέων ἐν ἐξετασμῷ αὐτῶν.
\vs{7}Δίκαιος δὲ ἐὰν φθάσῃ τελευτῆσαι, ἐν ἀναπαύσει ἔσται.

\vs{8}Γῆρας γὰρ τίμιον οὐ τὸ πολυχρόνιον, οὐδὲ ἀριθμῷ ἐτῶν μεμέτρηται.
\vs{9}Πολιὰ δέ ἐστιν φρόνησις ἀνθρώποις, καὶ ἡλικία γήρως βίος ἀκηλίδωτος.
\vs{10}Εὐάρεστος τῷ Θεῷ γενόμενος ἠγαπήθη, καὶ ζῶν μεταξὺ ἁμαρτωλῶν μετετέθη.
\vs{11}Ἡρπάγη μὴ κακία ἀλλάξῃ σύνεσιν αὐτοῦ, ἢ δόλος ἀπατήσῃ ψυχὴν αὐτοῦ.
\vs{12}Βασκανία γὰρ φαυλότητος ἀμαυροῖ τὰ καλὰ, καὶ ῥεμβασμὸς ἐπιθυμίας μεταλλεύει νοῦν ἄκακον.
\vs{13}Τελειωθεὶς ἐν ὀλίγῳ ἐπλήρωσε χρόνους μακρούς.
\vs{14}Ἀρεστὴ γὰρ ἦν Κυρίῳ ἡ ψυχὴ αὐτοῦ· διὰ τοῦτο ἔσπευσεν ἐκ μέσου πονηρίας.
\vs{15}Οἱ δὲ λαοὶ ἰδόντες καὶ μὴ νοήσαντες, μηδὲ θέντες ἐπὶ διανοίᾳ τὸ τοιοῦτο, ὅτι χάρις καὶ ἔλεος ἐν τοῖς ἐκλεκτοῖς αὐτοῦ, καὶ ἐπισκοπὴ ἐν τοῖς ὁσίοις αὐτοῦ.

\vs{16}Κατακρινεῖ δὲ δίκαιος καμὼν τοῦς ζῶντας ἀσεβεῖς, καὶ νεότης τελεσθεῖσα ταχέως πολυετὲς γῆρας ἀδίκου.
\vs{17}Ὄψονται γὰρ τελευτὴν σοφοῦ, καὶ οὐ νοήσουσι τί ἐβουλεύσατο περὶ αὐτοῦ, καὶ εἰς τί ἠσφαλίσατο αὐτὸν ὁ Κύριος.
\vs{18}Ὄψονται καὶ ἐξουθενήσουσιν, αὐτοὺς δὲ ὁ Κύριος ἐκγελάσεται· καὶ ἔσονται μετὰ τοῦτο εἰς πτῶμα ἄτιμον, καὶ εἰς ὕβριν ἐν νεκροῖς διʼ αἰῶνος.
\vs{19}Ὅτι ῥήξει αὐτοὺς ἀφώνους πρηνεῖς, καὶ σαλεύσει αὐτοὺς ἐκ θεμελίων, καὶ ἕως ἐσχάτου χερσωθήσονται, καὶ ἔσονται ἐν ὀδύνῃ, καὶ ἡ μνήμη αὐτῶν ἀπολεῖται.
\vs{20}Ἐλεύσονται ἐν συλλογισμῷ ἁμαρτημάτων αὐτῶν δειλοὶ, καὶ ἐλέγξει αὐτοὺς ἐξεναντίας τὰ ἀνομήματα αὐτῶν.

\ch{5}
Τότε στήσεται ἐν παῤῥησίᾳ πολλῇ ὁ δίκαιος κατὰ πρόσωπον τῶν θλιψάντων αὐτὸν, καὶ τῶν ἀθετούντων τοὺς πόνους αὐτοῦ.
\vs{2}Ἰδόντες ταραχθήσονται φόβῳ δεινῷ καὶ ἐκστήσονται ἐπὶ τῷ παραδόξῳ τῆς σωτηρίας.
\vs{3}Ἐροῦσιν ἑαυτοῖς μετανοοῦντες, καὶ διὰ στενοχωρίαν πνεύματος στενάζοντες,

Οὗτος ἦν ὃν ἔσχομέν ποτε εἰς γέλωτα καὶ εἰς παραβολὴν ὀνειδισμοῦ.
\vs{4}Οἱ ἄφρονες τὸν βίον αὐτοῦ ἐλογισάμεθα μανίαν, καὶ τὴν τελευτὴν αὐτοῦ ἄτιμον.
\vs{5}Πῶς κατελογίσθη ἐν υἱοῖς Θεοῦ, καὶ ἐν ἁγίοις ὁ κλῆρος αὐτοῦ ἐστιν;
\vs{6}Ἄρα ἐπλανήθημεν ἀπὸ ὁδοῦ ἀληθείας, καὶ τὸ τῆς δικαιοσύνης φῶς οὐκ ἔλαμψεν ἡμῖν, καὶ ὁ ἥλιος οὐκ ἀνέτειλεν ἡμῖν.
\vs{7}Ἀνομίας ἐνεπλήσθημεν τρίβοις καὶ ἀπωλείας, καὶ διωδεύσαμεν ἐρήμους ἀβάτους, τὴν δὲ ὁδὸν Κυρίου οὐκ ἔγνωμεν.

\vs{8}Τί ὠφέλησεν ἡμᾶς ἡ ὑπερηφανία; καὶ τὶ πλοῦτος μετὰ ἀλαζονείας συμβέβληται ἡμῖν;
\vs{9}Παρῆλθεν ἐκεῖνα πάντα ὡς σκιὰ, καὶ ὡς ἀγγελία παρατρέχουσα·
\vs{10}ὡς ναῦς δειρχομενη κυμαινόμενον ὕδωρ, ἧς διαβάσης οὐκ ἔστιν ἴχνος εὑρεῖν, οὐδὲ ἀτραπὸν τρόπιος αὐτῆς ἐν κύμασιν·
\vs{11}ἢ ὡς ὀρνέου διϊπτάντος ἀέρα, οὐθὲν εὑρίσκεται τεκμήριον πορείας, πληγῇ δὲ ταρσῶν μαστιζόμενον πνεῦμα κοῦφον καὶ σχιζόμενον βίᾳ ῥοιζου, κινουμένων πτερύγων διωδεύθη, καὶ μετὰ τοῦτο οὐχ εὑρέθη σημεῖον ἐπιβάσεως ἐν αὐτῷ·
\vs{12}ἢ ὡς βέλους βληθέντος ἐπὶ σκοπὸν, τμηθεὶς ὁ ἀὴρ εὐθέως εἰς ἑαυτὸν ἀνελύθη, ὡς ἀγνοῆσαι τὴν δίοδον αὐτοῦ·
\vs{13}οὕτως καὶ ἡμεῖς γεννηθέντες ἐξελίπομεν· καὶ ἀρετῆς μὲν σημεῖον οὐδὲν ἔσχομεν δεῖξαι, ἐν δὲ τῇ κακίᾳ ἡμῶν κατεδαπανήθημεν.

\vs{14}Ὅτι ἐλπὶς ἀσεβοῦς ὡς φερόμενος χοῦς ὑπὸ ἀνέμου, καὶ ὡς πάχνη ὑπὸ λαίλαπος διωχθεῖσα λεπτὴ, καὶ ὡς καπνὸς ὑπὸ ἀνέμου διεχύθη, καὶ ὡς μνεία καταλύτου μονοημέρου παρώδευσε.

\vs{15}Δίκαιοι δὲ εἰς τὸν αἰῶνα ζῶσι, καὶ ἐν Κυρίῳ ὁ μισθὸς αὐτῶν, καὶ ἡ φροντὶς αὐτῶν παρὰ ὑψίστῳ.
\vs{16}Διὰ τοῦτο λήψονται τὸ βασίλειον τῆς εὐπρεπείας, καὶ τὸ διάδημα τοῦ κάλλους ἐκ χειρὸς Κυρίου, ὅτι τῇ δεξιᾷ σκεπάσει αὐτούς, καὶ τῷ βραχίονι ὑπερασπιεῖ αὐτῶν.

\vs{17}Λήψεται πανοπλίαν τὸν ζῆλον αὐτοῦ, καὶ ὁπλοποιήσει τὴν κτίσιν εἰς ἄμυναν ἐχθρῶν.
\vs{18}Ἐνδύσεται θώρακα δικαιοσύνην, καὶ περιθήσεται κόρυθα κρίσιν ἀνυπόκριτον.
\vs{19}Λήψεται ἀσπίδα ἀκαταμάχητον ὁσιότητα,
\vs{20}ὀξυνεῖ δὲ ἀπότομον ὀργὴν εἰς ῥομφαίαν, συνεκπολεμήσει δὲ αὐτῷ ὁ κόσμος ἐπὶ τοὺς παράφρονας.

\vs{21}Πορεύσονται εὔστοχοι βολίδες ἀστραπῶν, καὶ ὡς ἀπὸ εὐκύκλου τόξου τῶν νεφῶν ἐπὶ σκοπὸν ἁλοῦνται.
\vs{22}Καὶ ἐκ πετροβόλου θυμοῦ πλήρεις ῥιφήσονται χάλαζαι· ἀγανακτήσει κατʼ αὐτῶν ὕδωρ θαλάσσης, ποταμοὶ δὲ συγκλύσουσιν ἀποτόμως.

\vs{23}Ἀντιστήσεται αὐτοῖς πνεῦμα δυνάμεως, καὶ ὡς λαίλαψ ἐκλικμήσει αὐτούς· καὶ ἐρημώσει πᾶσαν τὴν γῆν ἀνομία, καὶ ἡ κακοπραγία περιτρέψει θρόνους δυναστῶν.

\ch{6}
Ἀκούσατε οὖν βασιλεῖς καὶ σύνετε, μάθετε δικασταὶ περάτων γῆς.
\vs{2}Ἐνωτίσασθε οἱ κρατοῦντες πλήθους, καὶ γεγαυρομένοι ἐπὶ ὄχλοις ἐθνῶν.
\vs{3}Ὅτι ἐδόθη παρὰ τοῦ Κυρίου ἡ κράτησις ὑμῖν, καὶ ἡ δυναστεία παρὰ ὑψίστου, ὃς ἐξετάσει ὑμῶν τὰ ἔργα, καὶ τὰς βουλὰς διερευνήσει.
\vs{4}Ὅτι ὑπηρέται ὄντες τῆς αὐτοῦ βασιλείας οὐκ ἐκρίνατε ὀρθῶς, οὐδὲ ἐφυλάξατε νόμον, οὐδὲ κατὰ τὴν βουλὴν τοῦ Θεοῦ ἐπορεύθητε·
\vs{5}φρικτῶς καὶ ταχέως ἐπιστήσεται ὑμῖν, ὅτι κρίσις ἀπότομος ἐν τοῖς ὑπερέχουσιν γίνεται.
\vs{6}Ὁ γὰρ ἐλάχιστος συγγνωστός ἐστιν ἐλέους, δυνατοὶ δὲ δυνατῶς ἐτασθήσονται·
\vs{7}οὐ γὰρ ὑποστελεῖται πρόσωπον ὁ πάντων δεσπότης, οὐδὲ ἐντραπήσεται μέγεθος· ὅτι μικρὸν καὶ μέγαν αὐτὸς ἐποίησεν, ὁμοίως τε προνοεῖ περὶ πάντων.
\vs{8}Τοῖς δὲ κραταιοῖς ἰσχυρὰ ἐφίσταται ἔρευνα.

\vs{9}Πρὸς ὑμᾶς οὖν ὦ τύραννοι οἱ λόγοι μου, ἵνα μάθητε σοφίαν καὶ μὴ παραπέσητε.
\vs{10}Οἱ γὰρ φυλάξαντες ὁσίως τὰ ὅσια ὁσιωθήσονται, καὶ οἱ διδαχθέντες αὐτὰ εὑρήσουσιν ἀπολογίαν.
\vs{11}Ἐπιθυμήσατε οὖν τῶν λόγων μου, ποθήσατε καὶ παιδευθήσεσθε.

\vs{12}Λαμπρὰ καὶ ἀμάραντός ἐστιν ἡ σοφία, καὶ εὐχερῶς θεωρεῖται ὑπὸ τῶν ἀγαπώντων αὐτήν, καὶ εὑρίσκεται ὑπὸ τῶν ζητούντων αὐτήν.

\vs{13}Φθάνει τοὺς ἐπιθυμοῦντας προγνωσθῆναι.
\vs{14}Ὁ ὀρθρίσας ἐπʼ αὐτὴν οὐ κοπιάσει, πάρεδρον γὰρ εὑρήσει τῶν πυλῶν αὐτοῦ.
\vs{15}Τὸ γὰρ ἐνθυμηθῆναι περὶ αὐτῆς φρονήσεως τελειότης, καὶ ὁ ἀγρυπνήσας διʼ αὐτὴν ταχέως ἀμέριμνος ἔσται.
\vs{16}Ὅτι τοὺς ἀξίους αὐτῆς αὕτη περιέρχεται ζητοῦσα, καὶ ἐν ταῖς τρίβοις φαντάζεται αὐτοῖς εὐμενῶς, καὶ ἐν πάσῃ ἐπινοίᾳ ὑπαντᾷ αὐτοῖς.
\vs{17}Ἀρχὴ γὰρ αὐτῆς ἡ ἀληθεστάτη παιδείας ἐπιθυμία, φροντὶς δὲ παιδείας ἀγάπη,
\vs{18}ἀγάπη δὲ τήρησις νόμων αὐτῆς, προσοχὴ δὲ νόμων βεβαίωσις ἀφθαρσίας,
\vs{19}ἀφθαρσία δὲ ἐγγὺς εἶναι ποιεῖ Θεοῦ.
\vs{20}Ἐπιθυμία ἄρα σοφίας ἀνάγει ἐπὶ βασιλείαν.

\vs{21}Εἰ οὖν ἥδεσθε ἐπὶ θρόνοις καὶ σκήπτροις τύραννοι λαῶν, τιμήσατε σοφίαν, ἵνα εἰς τὸν αἰῶνα βασιλεύσητε·
\vs{22}Τί δέ ἐστι σοφία καὶ πῶς ἐγένετο, ἀπαγγελῶ, καὶ οὐκ ἀποκρύψω ὑμῖν μυστήρια, ἀλλʼ ἀπʼ ἀρχῆς γενέσεως ἐξιχνιάσω, καὶ θήσω εἰς τὸ ἐμφανὲς τὴν γνῶσιν αὐτῆς, καὶ οὐ μὴ παροδεύσω τὴν ἀλήθειαν·
\vs{23}οὔτε μὴν φθόνῳ τετηκότι συνοδεύσω, ὅτι οὗτος οὐ κοινωνήσει σοφίᾳ.
\vs{24}Πλῆθος δὲ σοφῶν σωτηρία κόσμου, καὶ βασιλεὺς φρόνιμος εὐστάθεια δήμου.
\vs{25}Ὥστε παιδεύεσθε τοῖς ῥήμασί μου, καὶ ὠφεληθήσεσθε.

\ch{7}
Εἰμι μὲν κᾀγὼ θνητὸς ἄνθρωπος, ἶσος ἅπασι, καὶ γηγενοῦς ἀπόγονος πρωτοπλάστου.
\vs{2}Καὶ ἐν κοιλίᾳ μητρὸς ἐγλύφην σὰρξ δεκαμηνιαίῳ χρόνῳ, παγεὶς ἐν αἵματι ἐκ σπέρματος ἀνδρὸς καὶ ἡδονῆς ὕπνῳ συνελθούσης.
\vs{3}Καὶ ἐγὼ δὲ γενόμενος ἔσπασα τὸν κοινὸν ἀέρα, καὶ ἐπὶ τὴν ὁμοιοπαθῆ κατέπεσον γῆν, πρώτην φωνὴν τὴν ὁμοίαν πᾶσιν ἶσα κλαίων.
\vs{4}Ἐν σπαργάνοις ἀνετράφην, καὶ ἐν φροντίσιν.
\vs{5}Οὐδεὶς γὰρ βασιλεὺς ἑτέραν ἔσχε γενέσεως ἀρχήν.
\vs{6}Μία δὲ πάντων εἴσοδος εἰς τὸν βίον, ἔξοδός τε ἴση.

\vs{7}Διὰ τοῦτο ηὐξάμην, καὶ φρόνησις ἐδόθη μοι, ἐπεκαλεσάμην, καὶ ἦλθέ μοι πνεῦμα σοφίας.
\vs{8}Προέκρινα αὐτὴν σκήπτρων καὶ θρόνων, καὶ πλοῦτον οὐδὲν ἡγησάμην ἐν συγκρίσει αὐτῆς.
\vs{9}Οὐδὲ ὁμοίωσα αὐτῇ λίθον ἀτίμητον, ὅτι ὁ πᾶς χρυσὸς ἐν ὄψει αὐτῆς ψάμμος ὀλίγη, καὶ ὡς πηλὸς λογισθήσεται ἄργυρος ἐναντίον αὐτῆς.
\vs{10}Ὑπὲρ ὑγίειαν καὶ εὐμορφίαν ἠγάπησα αὐτὴν, καὶ προειλόμην αὐτὴν ἀντὶ φωτὸς ἔχειν, ὅτι ἀκοίμητον τὸ ἐκ ταύτης φέγγος.

\vs{11}Ἦλθε δέ μοι τὰ ἀγαθὰ ὁμοῦ πάντα μετʼ αὐτῆς, καὶ ἀναρίθμητος πλοῦτος ἐν χερσὶν αὐτῆς.
\vs{12}Εὐφράνθην δὲ ἐπὶ πάντων, ὅτι αὐτῶν ἡγεῖται σοφία, ἠγνόουν δὲ αὐτὴν γενέτιν εἶναι τούτων.

\vs{13}Ἀδόλως τε ἔμαθον, ἀφθόνως τε μεταδίδωμι, τὸν πλοῦτον αὐτῆς οὐκ ἀποκρύπτομαι.
\vs{14}Ἀνεκλιπὴς γὰρ θησαυρός ἐστιν ἀνθρώποις, ὃν οἱ χρησάμενοι πρὸς Θεὸν ἐστείλαντο φιλίαν, διὰ τὰς ἐκ παιδείας δωρεὰς συσταθέντες.

\vs{15}Ἐμοὶ δὲ δῴη ὁ Θεὸς εἰπεῖν κατὰ γνώμην, καὶ ἐνθυμηθῆναι ἀξίως τῶν δεδομένων, ὅτι αὐτὸς καὶ τῆς σοφίας ὁδηγός ἐστι, καὶ τῶν σοφῶν διορθωτής.
\vs{16}Ἐν γὰρ χειρὶ αὐτοῦ καὶ ἡμεῖς καὶ οἱ λόγοι ἡμῶν, πᾶσά τε φρόνησις καὶ ἐργατειῶν ἐπιστήμη.
\vs{17}Αὐτὸς γάρ μοι ἔδωκε τῶν ὄντων γνῶσιν ἀψευδῆ, εἰδέναι σύστασιν κόσμου καὶ ἐνέργειαν στοιχείων,
\vs{18}ἀρχὴν καὶ τέλος καὶ μεσότητα χρόνων, τροπῶν ἀλλαγὰς καὶ μεταβολὰς καιρῶν,
\vs{19}ἐνιαυτῶν κύκλους καὶ ἀστέρων θέσεις,
\vs{20}φύσεις ζώων καὶ θυμοὺς θηρίων, πνευμάτων βίας καὶ διαλογισμοὺς ἀνθρώπων, διαφορὰς φυτῶν καὶ δυνάμεις ῥιζῶν,
\vs{21}ὅσα τέ ἐστι κρυπτὰ καὶ ἐμφανῆ ἔγνων.

\vs{22}Ἡ γὰρ πάντων τεχνίτις ἐδίδαξέ με σοφία· ἔστι γὰρ ἐν αὐτῇ πνεῦμα νοερὸν, ἅγιον, μονογενὲς, πολυμερὲς, λεπτὸν, εὐκίνητον, τρανὸν, ἀμόλυντον, σαφὲς, ἀπήμαντον, φιλάγαθον, ὀξὺ, ἀκώλυτον, εὐεργετικὸν,
\vs{23}φιλάνθρωπον, βέβαιον, ἀσφαλὲς, ἀμέριμνον, παντοδύναμον, πανεπίσκοπον, καὶ διὰ πάντων χωροῦν πνευμάτων νοερῶν, καθαρῶν, λεπτοτάτων.

\vs{24}Πάσης γὰρ κινήσεως κινητικώτερον σοφία, διήκει δὲ καὶ χωρεῖ διὰ πάντων διὰ τὴν καθαρότητα.
\vs{25}Ἀτμὶς γάρ ἐστι τῆς τοῦ Θεοῦ δυνάμεως. καὶ ἀπόῤῥοια τῆς τοῦ παντοκράτορος δόξης εἰλικρινής· διὰ τοῦτο οὐδὲν μεμιαμμένον εἰς αὐτὴν παρεμπίπτει.
\vs{26}Ἀπαύγασμα γάρ ἐστι φωτὸς ἀϊδίου, καὶ ἔσοπτρον ἀκηλίδωτον τῆς τοῦ Θεοῦ ἐνεργείας, καὶ εἰκὼν τῆς ἀγαθότητος αὐτοῦ.
\vs{27}Μία δὲ οὖσα πάντα δύναται, καὶ μένουσα ἐν αὑτῇ τὰ πάντα καινίζει, καὶ κατὰ γενεὰς εἰς ψυχὰς ὁσίας μεταβαίνουσα, φίλους Θεοῦ καὶ προφήτας κατασκευάζει.

\vs{28}Οὐθὲν γὰρ ἀγαπᾷ ὁ Θεὸς, εἰ μὴ τὸν σοφίᾳ συνοικοῦντα.
\vs{29}Ἔστι γὰρ αὕτη εὐπρεπεστέρα ἡλίου, καὶ ὑπὲρ πᾶσαν ἄστρων θέσιν, φωτὶ συγκρινομένη εὑρίσκεται προτέρα.
\vs{30}Τοῦτο μὲν γὰρ διαδέχεται νὺξ, σοφίας δὲ οὐκ ἀντισχύει κακία.

\ch{8}
Διατείνει δὲ ἀπὸ πέρατος εἰς πέρας εὐρώστως, καὶ διοικεῖ τὰ πάντα χρηστῶς.

\vs{2}Ταύτην ἐφίλησα καὶ ἐξεζήτησα ἐκ νεότητός μου, καὶ ἐζήτησα νύμφην ἀγαγέσθαι ἐμαυτῷ, καὶ ἐραστὴς ἐγενόμην τοῦ κάλλους αὐτῆς.
\vs{3}Εὐγένειαν δοξάζει συμβίωσιν Θεοῦ ἔχουσα, καὶ ὁ πάντων δεσπότης ἠγάπησεν αὐτήν.
\vs{4}Μύστις γάρ ἐστι τῆς τοῦ Θεοῦ ἐπιστήμης, καὶ αἱρετὶς τῶν ἔργων αὐτοῦ.

\vs{5}Εἰ δὲ πλοῦτός ἐστιν ἐπιθυμητὸν κτῆμα ἐν βίῳ, τί σοφίας πλουσιώτερον τῆς τὰ πάντα ἐργαζομένης;
\vs{6}Εἰ δὲ φρόνησις ἐργάζεται, τίς αὐτῆς τῶν ὄντων μᾶλλόν ἐστι τεχνίτης;
\vs{7}Καὶ εἰ δικαιοσύνην ἀγαπᾷ τις, οἱ πόνοι ταύτης εἰσὶν ἀρεταί· σωφροσύνην γὰρ καὶ φρόνησιν ἐκδιδάσκει, δικαιοσύνην καὶ ἀνδρίαν, ὧν χρησιμώτερον οὐδέν ἐστιν ἐν βίῳ ἀνθρώποις.
\vs{8}Εἰ δὲ καὶ πολυπειρίαν ποθεῖ τις, οἶδε τὰ ἀρχαῖα καὶ τὰ μέλλοντα εἰκάζειν, ἐπίσταται στροφὰς λόγων καὶ λύσεις αἰνιγμάτων, σημεῖα καὶ τέρατα προγινώσκει, καὶ ἐκβάσεις καιρῶν καὶ χρόνων.

\vs{9}Ἔκρινα τοίνυν ταύτην ἀγαγέσθαι πρὸς συμβίωσιν, εἰδὼς ὅτι ἔσται μοι σύμβουλος ἀγαθῶν, καὶ παραίνεσις φροντίδων καὶ λύπης.
\vs{10}Ἕξω διʼ αὐτὴν δόξαν ἐν ὄχλοις, καὶ τιμὴν παρὰ πρεσβυτέροις ὁ νέος.
\vs{11}Ὀξὺς εὑρεθήσομαι ἐν κρίσει, καὶ ἐν ὄψει δυναστῶν θαυμασθήσομαι.
\vs{12}Σιγῶντά με περιμενοῦσι, καὶ φθεγγομένῳ προσέξουσι, καὶ λαλοῦντος ἐπιπλεῖον, χεῖρα ἐπιθήσουσιν ἐπὶ στόμα αὐτῶν.

\vs{13}Ἕξω διʼ αὐτὴν ἀθανασίαν, καὶ μνήμην αἰώνιον τοῖς μετ ἐμὲ ἀπολείψω.
\vs{14}Διοικήσω λαοὺς, καὶ ἔθνη ὑποταγήσεταί μοι.
\vs{15}Φοβηθήσονταί με ἀκούσαντες τύραννοι φρικτοὶ, ἐν πλήθει φανοῦμαι ἀγαθὸς, καὶ ἐν πολέμῳ ἀνδρεῖος.
\vs{16}Εἰσελθὼν εἰς τὸν οἶκόν μου προσαναπαύσομαι αὐτῇ· οὐ γὰρ ἔχει πικρίαν ἡ συναναστροφὴ αὐτῆς, οὐδὲ ὀδύνην ἡ συμβίωσις αὐτῆς, ἀλλὰ εὐφροσύνην καὶ χαράν.

\vs{17}Ταῦτα λογισάμενος ἐν ἐμαυτῷ, καὶ φροντίσας ἐν καρδίᾳ μου, ὅτι ἐστὶν ἀθανασία ἐν συγγενείᾳ σοφίας,
\vs{18}καὶ ἐν φιλίᾳ αὐτῆς τέρψις ἀγαθὴ, καὶ ἐν πόνοις χειρῶν αὐτῆς πλοῦτος ἀνεκλιπὴς, καὶ ἐν συγγυμνασίᾳ ὁμιλίας αὐτῆς φρόνησις, καὶ εὔκλεια ἐν κοινωνίᾳ λόγων αὐτῆς, περιῄειν ζητῶν ὅπως λάβω αὐτὴν εἰς ἐμαυτόν.

\vs{19}Παῖς δὲ ἤμην εὐφυὴς, ψυχῆς τε ἔλαχον ἀγαθῆς,
\vs{20}μᾶλλον δὲ ἀγαθὸς ὢν ἦλθον εἰς σῶμα ἀμίαντον.
\vs{21}Γνοὺς δὲ ὅτι οὐκ ἄλλως ἔσομαι ἐγκρατὴς, ἐὰν μὴ ὁ Θεὸς δῷ, καὶ τοῦτο δʼ ἦν φρονήσεως τὸ εἰδέναι τίνος ἡ χάρις, ἐνέτυχον τῷ Κυρίῳ, καὶ ἐδεήθην αὐτοῦ, καὶ εἶπον ἐξ ὅλης τῆς καρδίας μου,

\ch{9}
Θεὲ πατέρων καὶ Κύριε τοῦ ἐλέους σου, ὁ ποιήσας τὰ πάντα ἐν λόγῳ σου,
\vs{2}καὶ τῇ σοφίᾳ σου κατεσκεύασας ἄνθρωπον, ἵνα δεσπόζῃ τῶν ὑπὸ σοῦ γενομένων κτισμάτων,
\vs{3}καὶ διέπῃ τὸν κόσμον ἐν ὁσιότητι καὶ δικαιοσύνῃ, καὶ ἐν εὐθύτητι ψυχῆς κρίσιν κρίνῃ·
\vs{4}δός μοι τὴν τῶν σῶν θρόνων πάρεδρον σοφίαν, καὶ μή με ἀποδοκιμάσῃς ἐκ παίδων σου.
\vs{5}Ὅτι ἐγὼ δοῦλος σὸς καὶ υἱὸς τῆς παιδίσκης σου, ἄνθρωπος ἀσθενὴς καὶ ὀλιγοχρόνιος καὶ ἐλάσσων ἐν συνέσει κρίσεως καὶ νόμων.

\vs{6}Κᾂν γάρ τις ᾖ τέλειος ἐν υἱοῖς ἀνθρώπων, τῆς ἀπὸ σοῦ σοφίας ἀπούσης, εἰς οὐδὲν λογισθήσεται.

\vs{7}Σύ με προείλω βασιλέα λαοῦ σου, καὶ δικαστὴν υἱῶν σου καὶ θυγατέρων.
\vs{8}Εἶπας οἰκοδομῆσαι ναὸν ἐν ὄρει ἁγίῳ σου, καὶ ἐν πόλει κατασκηνώσεώς σου θυσιαστήριον, μίμημα σκηνῆς ἁγίας ἣν προητοίμασας ἀπʼ ἀρχῆς.
\vs{9}Καὶ μετὰ σοῦ ἡ σοφία ἡ εἰδυῖα τὰ ἔργα σου, καὶ παροῦσα ὅτε ἐποίεις τὸν κόσμον, καὶ ἐπισταμένη τί ἀρεστὸν ἐν ὀφθαλμοῖς σου, καὶ τί εὐθὲς ἐν ἐντολαῖς σου.
\vs{10}Ἐξαπόστειλον αὐτὴν ἐξ ἁγίων οὐρανῶν, καὶ ἀπὸ θρόνου δόξης σου πέμψον αὐτὴν, ἵνα συμπαροῦσά μοι κοπιάσῃ, καὶ γνῷ τί εὐάρεστόν ἐστι παρὰ σοί.
\vs{11}Οἶδε γὰρ ἐκείνη πάντα καὶ συνιεῖ, καὶ ὁδηγήσει με ἐν ταῖς πράξεσί μου σωφρόνως, καὶ φυλάξει με ἐν τῇ δόξῃ αὐτῆς.
\vs{12}Καὶ ἔσται προσδεκτὰ τὰ ἔργα μου, καὶ διακρινῶ τὸν λαόν σου δικαίως, καὶ ἔσομαι ἄξιος θρόνων πατρός μου.

\vs{13}Τίς γὰρ ἄνθρωπος γνώσεται βουλὴν Θεοῦ; ἢ τίς ἐνθυμηθήσεται τί θέλει ὁ Κύριος;
\vs{14}Λογισμοὶ γὰρ θνητῶν δειλοὶ, καὶ ἐπισφαλεῖς αἱ ἐπίνοιαι ἡμῶν.
\vs{15}Φθαρτὸν γὰρ σῶμα βαρύνει ψυχὴν, καὶ βρίθει τὸ γεῶδες σκῆνος νοῦν πολυφροντίδα.
\vs{16}Καὶ μόλις εἰκάζομεν τὰ ἐπὶ γῆς, καὶ τὰ ἐν χερσὶν εὑρίσκομεν μετὰ πόνου· τὰ δὲ ἐν οὐρανοῖς τίς ἐξιχνίασε;
\vs{17}Βουλὴν δέ σου τίς ἔγνω, εἰ μὴ σὺ ἔδωκας σοφίαν, καὶ ἔπεμψας τὸ ἅγιόν σου πνεῦμα ἀπὸ ὑψίστων;
\vs{18}Καὶ οὕτως διωρθώθησαν αἱ τρίβοι τῶν ἐπὶ γῆς, καὶ τὰ ἀρεστά σου ἐδιδάχθησαν ἄνθρωποι, καὶ τῇ σοφίᾳ ἐσώθησαν.

\ch{10}
Αὕτη πρωτόπλαστον πατέρα κόσμου μόνον κτισθέντα διεφύλαξε, καὶ ἐξείλατο αὐτὸν ἐκ παραπτώματος ἰδίου,
\vs{2}ἔδωκέ τε αὐτῷ ἰσχὺν κρατῆσαι ἁπάντων.

\vs{3}Ἀποστὰς δὲ ἀπʼ αὐτῆς ἄδικος ἐν ὀργῇ αὐτοῦ, ἀδελφοκτόνοις συναπώλετο θυμοῖς,
\vs{4}διʼ ὃν κατακλυζομένην γῆν πάλιν διέσωσε σοφία, διʼ εὐτελοῦς ξύλου τὸν δίκαιον κυβερνήσασα.
\vs{5}Αὕτη καὶ ἐν ὁμονοίᾳ πονηρίας ἐθνῶν συγχυθέντων εὗρε τὸν δίκαιον, καὶ ἐτήρησεν αὐτὸν ἄμεμπτον Θεῷ, καὶ ἐπὶ τέκνου σπλάγχνοις ἰσχυρὸν ἐφύλαξεν.

\vs{6}Αὕτη δίκαιον, ἐξαπολλυμένων ἀσεβῶν, ἐῤῥύσατο φυγόντα πῦρ καταβάσιον Πενταπόλεως·
\vs{7}οἷς ἐτὶ μαρτύριον τῆς πονηρίας καπνιζομένη καθέστηκε χέρσος, καὶ ἀτελέσιν ὥραις καρποφοροῦντα φυτά· ἀπιστοῦσης ψυχῆς μνημεῖον ἑστηκυῖα στήλη ἁλός.
\vs{8}Σοφίαν γὰρ παροδεύσαντες οὐ μόνον ἐβλάβησαν τοῦ μὴ γνῶναι τὰ καλὰ, ἀλλὰ καὶ τῆς ἀφροσύνης ἀπέλιπον τῷ βίῳ μνημόσυνον, ἵνα ἐν οἷς ἐσφάλησαν μηδὲ λαθεῖν δυνηθῶσι.
\vs{9}Σοφία δὲ τοὺς θεραπεύσαντας αὐτὴν ἐκ πόνων ἐῤῥύσατο.

\vs{10}Αὕτη φυγάδα ὀργῆς ἀδελφοῦ δίκαιον ὡδήγησεν ἐν τρίβοις εὐθείαις, ἔδειξεν αὐτῷ βασιλείαν Θεοῦ, καὶ ἔδωκεν αὐτῷ γνῶσιν ἁγίων, εὐπόρησεν αὐτὸν ἐν μόχθοις, καὶ ἐπλήθυνε τοὺς πόνους αὐτοῦ.
\vs{11}Ἐν πλεονεξίᾳ κατισχύοντων αὐτὸν παρέστη, καὶ ἐπλούτισεν αὐτόν.
\vs{12}Διεφύλαξεν αὐτὸν ἀπὸ ἐχθρῶν, καὶ ἀπὸ ἐνεδρευόντων ἠσφαλίσατο, καὶ ἀγῶνα ἰσχυρὸν ἐβράβευσεν αὐτῷ, ἵνα γνῷ, ὅτι παντὸς δυνατωτέρα ἐστὶν εὐσέβεια.

\vs{13}Αὕτη πραθέντα δίκαιον οὐκ ἐγκατέλιπεν, ἀλλὰ ἐξ ἁμαρτίας ἐῤῥύσατο αὐτόν·
\vs{14}συγκατέβη αὐτῷ εἰς λάκκον, καὶ ἐν δεσμοῖς οὐκ ἀφῆκεν αὐτὸν, ἕως ἤνεγκεν αὐτῷ σκῆπτρα βασιλείας καὶ ἐξουσίαν τυραννούντων αὐτοῦ· ψευδεῖς τε ἔδειξε τοὺς μωμησαμένους αὐτὸν, καὶ ἔδωκεν αὐτῷ δόξαν αἰώνιον.

\vs{15}Αὕτη λαὸν ὅσιον καὶ σπέρμα ἄμεμπτον ἐῤῥύσατο ἐξ ἔθνους θλιβόντων.

\vs{16}Εἰσῆλθεν εἰς ψυχὴν θεράποντος Κυρίου, καὶ ἀντέστη βασιλεῦσι φοβεροῖς ἐν τέρασι καὶ σημείοις.
\vs{17}Ἀπέδωκεν ὁσίοις μισθὸν κόπων αὐτῶν, ὡδήγησεν αὐτοὺς ἐν ὁδῷ θαυμαστῇ, καὶ ἐγένετο αὐτοῖς εἰς σκέπην ἡμέρας, καὶ εἰς φλόγα ἄστρων τὴν νύκτα.
\vs{18}Διεβίβασεν αὐτοὺς θάλασσαν ἐρυθρὰν, καὶ διήγαγεν αὐτοὺς διʼ ὕδατος πολλοῦ.
\vs{19}Τοὺς δὲ ἐχθροὺς αὐτῶν κατέκλυσε, καὶ ἐκ βάθους ἀβύσσου ἀνέβρασεν αὐτούς.
\vs{20}Διὰ τοῦτο δίκαιοι ἐσκύλευσαν ἀσεβεῖς, καὶ ὕμνησαν Κύριε τὸ ὄνομα τὸ ἅγιόν σου, τήν τε ὑπέρμαχόν σου χεῖρα ᾔνεσαν ὁμοθυμαδόν.
\vs{21}Ὅτι ἡ σοφία ἤνοιξε στόμα κωφῶν, καὶ γλώσσας νηπίων ἔθηκε τρανάς.

\ch{11}
Εὐωδώσε τὰ ἔργα αὐτῶν ἐν χειρὶ προφήτου ἁγίου.
\vs{2}Διώδευσαν ἔρημον ἀοίκητον, καὶ ἐν ἀβάτοις ἔπηξαν σκηνάς.
\vs{3}Ἀντέστησαν πολεμίοις, καὶ ἠμύναντο ἐχθρούς.
\vs{4}Ἐδίψησαν καὶ ἐπεκαλέσαντό σε, καὶ ἐδόθη αὐτοῖς ἐκ πέτρας ἀκροτόμου ὕδωρ, καὶ ἴαμα δίψης ἐκ λίθου σκληροῦ.
\vs{5}Διʼ ὧν γὰρ ἐκολάσθησαν οἱ ἐχθροὶ αὐτῶν, διὰ τούτων αὐτοὶ ἀποροῦντες εὐεργετήθησαν.
\vs{6}Ἀντὶ μὲν πηγῆς ἀεννάου ποταμοῦ αἵματι λυθρώδει ταραχθέντες
\vs{7}εἰς ἔλεγχον νηπιοκτόνου διατάγματος, ἔδωκας αὐτοῖς δαψιλὲς ὕδωρ ἀνελπίστως·
\vs{8}δείξας διὰ τοῦ τότε δίψους πῶς τοὺς ὑπεναντίους ἐκόλασας.

\vs{9}Ὅτε γὰρ ἐπειράσθησαν, καίπερ ἐν ἐλέει παιδευόμενοι, ἔγνωσαν πῶς ἐν ὀργῇ κρινόμενοι ἀσεβεῖς ἐβασανίζοντο.
\vs{10}Τούτους μὲν γὰρ ὡς πατὴρ νουθεντῶν ἐδοκίμασας, ἐκείνους δὲ ὡς ἀπότομος βασιλεὺς καταδικάζων ἐξήτασας.
\vs{11}Καὶ ἀπόντες δὲ καὶ παρόντες ὁμοίως ἐτρύχοντο.
\vs{12}Διπλῆ γὰρ αὐτοὺς ἔλαβε λύπη, καὶ στεναγμὸς μνημῶν τῶν παρελθουσῶν.
\vs{13}Ὅτε γὰρ ἤκουσαν διὰ τῶν ἰδίων κολάσεων εὐεργετουμένους αὐτοὺς, ᾔσθοντο τοῦ Κυρίου.
\vs{14}Τὸν γὰρ ἐν ἐκθέσει πάλαι ῥιφέντα ἀπεῖπον χλευάζοντες, ἐπὶ τέλει τῶν ἐκβάσεων ἐθαύμασαν, οὐχ ὅμοια δικαίοις διψήσαντες.

\vs{15}Ἀντὶ δὲ λογισμῶν ἀσυνέτων ἀδικίας αὐτῶν, ἐν οἷς πλανηθέντες ἐθρήσκευον ἄλογα ἑρπετὰ καὶ κνώδαλα εὐτελῆ, ἐπαπέστειλας αὐτοῖς πλῆθος ἀλόγων ζώων εἰς ἐκδίκησιν,
\vs{16}ἵνα γνῶσιν ὅτι διʼ ὧν τις ἁμαρτάνει, διὰ τούτων κολάζεται.

\vs{17}Οὐ γὰρ ἠπόρει ἡ παντοδύναμός σου χεὶρ κτίσασα τὸν κόσμον ἐξ ἀμόρφου ὕλης, ἐπιπέμψαι αὐτοῖς πλῆθος ἄρκων, ἢ θρασεῖς λέοντας,
\vs{18}ἢ νεοκτίστους θυμοῦ πλήρεις θῆρας ἀγνώστους, ἤτοι πυρπνόον φυσῶντας ἆσθμα, ἢ βρόμους λικμωμένους καπνοῦ, ἢ δεινοὺς ἀπʼ ὀμμάτων σπινθῆρας ἀστράπτοντας·
\vs{19}ὧν οὐ μόνον ἡ βλάβη ἠδύνατο συνεκτρίψαι αὐτοὺς, ἀλλὰ καὶ ἡ ὄψις ἐκφοβήσασα διολέσαι.
\vs{20}Καὶ χωρὶς δὲ τούτων, ἑνὶ πνεύματι πεσεῖν ἐδύναντο ὑπὸ τῆς δίκης διωχθέντες, καὶ λικμηθέντες ὑπὸ πνεύματος δυνάμεώς σου· ἀλλὰ πάντα μέτρῳ καὶ ἀριθμῷ καὶ σταθμῷ διέταξας.
\vs{21}Τὸ γὰρ μεγάλως ἰσχύειν πάρεστί σοι πάντοτε, καὶ κράτει βραχίονός σου τις ἀντιστήσεται;
\vs{22}Ὅτι ὡς ῥοπὴ ἐκ πλαστίγγων ὅλος ὁ κόσμος ἐναντίον σου, καὶ ὡς ῥανὶς δρόσον ὀρθρινὴ κατελθοῦσα ἐπὶ γῆν.

\vs{23}Ἐλεεῖς δὲ πάντας, ὅτι πάντα δύνασαι, καὶ παρορᾷς ἁμαρτήματα ἀνθρώπων εἰς μετάνοιαν.
\vs{24}Ἀγαπᾷς γὰρ τὰ ὄντα πάντα, καὶ οὐδὲν βδελύσσῃ ὧν ἐποίησας, οὐδὲ γὰρ ἂν μισῶν τι κατεσκεύασας.
\vs{25}Πῶς δὲ ἔμεινεν ἄν τι εἰ μὴ σὺ ἐθέλησας; ἢ τὸ μὴ κληθὲν ὑπὸ σοῦ διετηρήθη;
\vs{26}Φείδῃ δὲ πάντων, ὅτι σά ἐστι, δέσποτα φιλόψυχε.

\ch{12}
Τὸ γὰρ ἄφθαρτόν σου πνεῦμά ἐστιν ἐν πᾶσι.
\vs{2}Διὸ τοὺς παραπίπτοντας κατʼ ὀλίγον ἐλέγχεις, καὶ ἐν οἷς ἁμαρτάνουσιν ὑπομιμνήσκων νουθετεῖς, ἵνα ἀπαλλαγέντες τῆς κακίας πιστεύσωσιν ἐπὶ σὲ Κύριε.
\vs{3}Καὶ γὰρ τοὺς παλαιοὺς οἰκήτορας τῆς ἁγίας σου γῆς
\vs{4}μισήσας, ἐπὶ τῷ ἔχθιστα πράσσειν ἔργα φαρμακειῶν, καὶ τελετὰς ἀνοσίους,
\vs{5}τέκνων τε φονέας ἀνελεήμονας, καὶ σπλαγχνοφάγων ἀνθρωπίνων σαρκῶν θοῖναν,
\vs{6}καὶ αἴματος ἐκ μέσου μυσταθείας σου, καὶ αὐθέντας γονεῖς ψυχῶν ἀβοηθήτων, ἐβουλήθης ἀπολέσαι διὰ χειρῶν πατέρων ἡμῶν·
\vs{7}ἵνα ἀξίαν ἀποικίαν δέξηται Θεοῦ παίδων ἡ παρὰ σοὶ πασῶν τιμιωτάτη γῆ.

\vs{8}Ἀλλὰ καὶ τούτων ὡς ἀνθρώπων ἐφείσω, ἀπέστειλάς τε προδρόμους τοῦ στρατοπέδου σου σφῆκας, ἵνα αὐτοὺς καταβραχὺ ἐξολοθρεύσωσιν.
\vs{9}Οὐκ ἀδυνατῶν ἐν παρατάξει ἀσεβεῖς δικαίοις ὑποχειρίους δοῦναι, ἢ θηρίοις δεινοῖς, ἢ λόγῳ ἀποτόμῳ ὑφʼ ἓν ἐκτρίψαι·
\vs{10}κρίνων δὲ καταβραχὺ ἐδίδους τόπον μετανοίας, οὐκ ἀγνοῶν, ὅτι πονηρὰ ἡ γένεσις αὐτῶν, καὶ ἔμφυτος ἡ κακία αὐτῶν, καὶ ὅτι οὐ μὴ ἀλλαγῇ ὁ λογισμὸς αὐτῶν εἰς τὸν αἰῶνα·
\vs{11}σπέρμα γὰρ ἦν κατηραμένον ἀπʼ ἀρχῆς· οὐδὲ εὐλαβούμενός τινα, ἐφʼ οἷς ἡμάρτανον ἄδειαν ἐδίδους.
\vs{12}Τίς γὰρ ἐρεῖ, τί ἐποίησας; ἢ τίς ἀντιστήσεται τῷ κρίματί σου; τίς δὲ ἐγκαλέσει σοι κατὰ ἐθνῶν ἀπολωλότων, ἃ σὺ ἐποίησας; ἢ τίς εἰς κατάστασίν σοι ἐλεύσεται ἔκδικος κατὰ ἀδίκων ἀνθρώπων;
\vs{13}Οὔτε γὰρ Θεός ἐστι πλὴν σοῦ, ᾧ μέλει περὶ πάντων, ἵνα δείξῃς ὅτι οὐκ ἀδίκως ἔκρινας.

\vs{14}Οὔτε βασιλεὺς ἢ τύραννος ἀντοφθαλμῆσαι δυνήσεταί σοι περὶ ὧν ἀπώλεσας.
\vs{15}Δίκαιος δὲ ὢν δικαίως τὰ πάντα διέπεις, αὐτὸν τὸν μὴ ὀφείλοντα κολασθῆναι καταδικάσαι ἀλλότριαν ἡγούμενος τῆς σῆς δυνάμεως.
\vs{16}Ἡ γὰρ ἰσχύς σου δικαιοσύνης ἀρχὴ, καὶ τὸ πάντων σε δεσπόζειν, πάντων φείδεσθαι ποιεῖ.
\vs{17}Ἰσχὺν γὰρ ἐνδείκνυσαι ἀπιστούμενος ἐπὶ δυνάμεως τελειότητι, καὶ ἐν τοῖς εἰδόσι τὸ θράσος ἐξελέγχεις.
\vs{18}Σὺ δὲ δεσπόζεν ἰσχύος ἐν ἐπιεικείᾳ κρίνεις καὶ μετὰ πολλῆς φειδοῦς διοικεῖς ἡμᾶς· πάρεστι γάρ σοι ὅταν θέλῃς τὸ δύνασθαι.

\vs{19}Ἐδίδαξας δέ σου τὸν λαὸν διὰ τῶν τοιούτων ἔργων, ὅτι δεῖ τὸν δίκαιον εἶναι φιλάνθρωπον· καὶ εὐέλπιδας ἐποίησας τοὺς υἱούς σου, ὅτι δίδως ἑπὶ ἁμαρτήμασι μετάνοιαν.
\vs{20}Εἰ γὰρ ἐχθροὺς παίδων σου καὶ ὀφειλομένους θανάτῳ μετὰ τοσαύτης ἐτιμώρησας προσοχῆς καὶ δεήσεως, δοὺς χρόνους καὶ τόπον διʼ ὧν ἀπαλλαγῶσι τῆς κακίας·
\vs{21}μετὰ πόσης ἀκριβείας ἔκρινας τοὺς υἱούς σου ὧν τοῖς πατράσιν ὅρκους καὶ συνθήκας ἔδωκας ἀγαθῶν ὑποσχέσεων;
\vs{22}Ἡμᾶς οὖν παιδεύων, τοὺς ἐχθροὺς ἡμῶν ἐν μυριότητι μαστιγοῖς, ἵνα σου τὴν ἀγαθότητα μεριμνῶμεν κρίνοντες, κρινόμενοι δὲ προσδοκῶμει ἔλεος.

\vs{23}Ὅθεν καὶ τοὺς ἐν ἀφροσύνῃ ζωῆς βιώσαντας ἀδίκους, διὰ τῶν ἰδίων ἐβασάνισας βδελυγμάτων.
\vs{24}Καὶ γὰρ τῶν πλάνης ὁδῶν μακρότερον ἐπλανήθησαν, θεοὺς ὑπολαμβάνοντες τὰ καὶ ἐν ζώοις τῶν ἐχθρῶν ἄτιμα, νηπίων δίκην ἀφρόνων ψευσθέντες.
\vs{25}Διὰ τοῦτο ὡς παισὶν ἀλογίστοις τὴν κρίσιν εἰς ἐμπαιγμὸν ἔπεμψας.
\vs{26}Οἱ δὲ παιγνίοις ἐπιτιμήσεως μὴ νουθετηθέντες, ἀξίαν Θεοῦ κρίσιν πειράσουσιν.
\vs{27}Ἐφʼ οἷς γὰρ αὐτοὶ πάσχοντες ἠγανάκτουν, ἐπὶ τούτοις οὓς ἐδόκουν θεοὺς, ἐν αὐτοῖς κολαζόμενοι, ἰδόντες ὃν πάλαι ἠρνοῦντο εἰδέναι, Θεὸν ἐπέγνωσαν ἀληθῆ· διὸ καὶ τὸ τέρμα τῆς καταδίκης ἐπʼ αὐτοὺς ἐπῆλθε.

\ch{13}
Μάταιοι μὲν γὰρ πάντες ἄνθρωποι φύσει, οἷς παρῆν Θεοῦ ἀγνωσία, καὶ ἐκ τῶν ὁρομένων ἀγαθῶν οὐκ ἴσχυσαν εἰδέναι τὸν ὄντα, οὔτε τοῖς ἔργοις προσχόντες ἐπέγνωσαν τὸν τεχνίτην.
\vs{2}Ἀλλʼ ἢ πῦρ, ἢ πνεῦμα, ἢ ταχινὸν ἀέρα, ἢ κύκλον ἄστρων, ἢ βίαιον ὕδωρ, ἢ φωστῆρας οὐρανοῦ, πρυτάνεις κόσμου θεοὺς ἐνόμισαν.
\vs{3}Ὧν εἰ μὲν τῇ καλλονῇ τερπόμενοι, θεοὺς ὑπελάμβανον, γνώτωσαν πόσῳ τούτων ὁ δεσπότης ἐστὶ βελτίων· ὁ γὰρ τοῦ κάλλους γενεσιάρχης ἔκτισεν αὐτά.
\vs{4}Εἰ δὲ δύναμιν καὶ ἐνέργειαν ἐκπλαγέντες, νοησάτωσαν ἀπʼ αὐτῶν πόσῳ ὁ κατασκευάσας αὐτὰ δυνατώτερός ἐστιν.

\vs{5}Ἐκ γὰρ μεγέθους καλλονῆς κτισμάτων ἀναλόγως ὁ γενεσιουργὸς αὐτῶν θεωρεῖται.
\vs{6}Ἀλλʼ ὅμως ἐπὶ τούτοις ἐστὶ μέμψις ὀλίγη, καὶ γὰρ αὐτοὶ τάχα πλανῶνται Θεὸν ζητοῦντες, καὶ θέλοντες εὑρεῖν.
\vs{7}Ἐν γὰρ τοῖς ἔργοις αὐτοῦ ἀναστρεφόμενοι διερευνῶσι, καὶ πείθονται τῇ ὄψει, ὅτι καλὰ τὰ βλεπόμενα.
\vs{8}Πάλιν δὲ οὐδʼ αὐτοὶ συγγνωστοί.
\vs{9}Εἰ γὰρ τοσοῦτον ἴσχυσαν εἰδέναι, ἵνα δύνωνται στοχάσασθαι τὸν αἰῶνα, τὸν τούτων δεσπότην πῶς τάχιον οὐχ εὗρον;

\vs{10}Ταλαίπωροι δὲ καὶ ἐν νεκροῖς αἱ ἐλπίδες αὐτων, οἵτινες ἐκάλεσαν θεοὺς ἔργα χειρῶν ἀνθρώπων, χρυσὸν καὶ ἄργυρον τέχνης ἐμμελέτημα, καὶ ἀπεικάσματα ζώων, ἢ λίθον ἄχρηστον χειρὸς ἔργον ἀρχαίας.
\vs{11}Εἰ δὲ καί τις ὑλοτῦμος τέκτων εὐκίνητον φυτὸν ἐκπρίσας, περιέξυσεν εὐμαθῶς πάντα τὸν φλοιὸν αὐτοῦ, καὶ τεχνησέμενος εὐπρεπῶς κατεσκεύασε χρήσιμον σκεῦος εἰς ὑπηρεσίαν ζωῆς,
\vs{12}τὰ δὲ ἀποβλήματα τῆς ἐργασίας εἰς ἑτοιμασίαν τροφῆς ἀναλώσας ἐνεπλήσθη,
\vs{13}τὸ δὲ ἐξ αὐτῶν ἀπόβλημα εἰς οὐθὲν εὔχρηστον, ξύλον σκολιὸν, καὶ ὄζοις συμπεφυκὸς, λαβὼν ἔγλυψεν ἐν ἐπιμελείᾳ ἀργίας αὐτοῦ, καὶ ἐμπειρίᾳ συνέσεως ἐτύπωσεν αὐτὸ, ἀπείκασεν αὐτὸ εἰκόνι ἀνθρώπου,
\vs{14}ἢ ζώῳ τινὶ εὐτελεῖ ὡμοίωσεν αὐτὸ, καταχρίσας μίλτῳ, καὶ φύκει ἐρυθῄνας χρόαν αὐτοῦ, καὶ πᾶσαν κηλίδα τὴν ἐν αὐτῷ καταχρίσας.
\vs{15}Καὶ ποιήσας αὐτῷ αὐτοῦ ἄξιον οἴκημα, ἐν τοίχῳ ἔθηκεν αὐτὸ ἀσφαλισάμενος σιδήρῳ·
\vs{16}ἵνα μὲν οὖν μὴ καταπέσῃ, προενόησεν αὐτοῦ, εἰδὼς ὅτι ἀδυνατεῖ ἑαυτῷ βοηθῆσαι, καὶ γάρ ἐστιν εἰκὼν, καὶ χρείαν ἔχει βοηθείας.

\vs{17}Περὶ δὲ κτημάτων καὶ γάμων αὐτοῦ καὶ τέκνων προσευχόμενος, οὐκ αἰσχύνεται τῷ ἀψύχῳ προσλαλῶν.
\vs{18}Καὶ περὶ μὲν ὑγείας τὸ ἀσθενὲς ἐπικαλεῖται, περὶ δὲ ζωῆς τὸν νεκρὸν ἀξιοῖ, περὶ δὲ ἐπικουρίας τὸν ἀπειρότατον ἱκετεύει, περὶ δὲ ὁδοιπορίας τὸ μηδὲ βάσει χρῆσθαι δυνάμενον,
\vs{19}περὶ δὲ πορισμοῦ καὶ ἐργασίας καὶ χειρῶν ἐπιτυχίας τὸ ἀδρανέστατον ταῖς χερσὶν εὐδράνειαν αἰτεῖται.

\ch{14}
Πλοῦν τις πάλιν στελλόμενος, καὶ ἄγρια μέλλων διοδεύειν κύματα, τοῦ φέροντος αὐτὸν πλοίου σαθρότερον ξύλον ἐπιβοᾶται.
\vs{2}Ἐκεῖνο μὲν γὰρ ὄρεξις πορισμῶν ἐπενόησε, τεχνίτης δὲ σοφίᾳ κατεσκεύασεν·
\vs{3}ἡ δὲ σὴ, Πάτερ, διακυβερνᾷ πρόνοια, ὅτι ἔδωκας καὶ ἐν θαλάσσῃ ὁδὸν καὶ ἐν κύμασι τρίβον ἀσφαλῆ·
\vs{4}δεικνὺς ὅτι δύνασαι ἐκ παντὸς σώζειν, ἵνα κᾂν ἄνευ τέχνης τις ἐπιβῇ.
\vs{5}Θέλεις δὲ μὴ ἀργὰ εἶναι τὰ τῆς σοφίας σου ἔργα, διὰ τοῦτο καὶ ἐλαχίστῳ ξύλῳ πιστεύουσιν ἄνθρωποι ψυχὰς, καὶ διελθόντες κλύδωνα σχεδίᾳ διεσώθησαν.

\vs{6}Καὶ ἀρχῆς γὰρ ἀπολλυμένων ὑπερηφάνων γιγάντων, ἡ ἐλπὶς τοῦ κόσμου ἐπὶ σχεδίας καταφυγοῦσα, ἀπέλιπεν αἰῶνι σπέρμα γενέσεως τῇ σῇ κυβερνηθεῖσα χειρί.
\vs{7}Εὐλόγηται γὰρ ξύλον διʼ οὗ γίνεται δικαιοσύνη.

\vs{8}Τὸ χειροποίητον δέ ἐπικατάρατον αὐτὸ, καὶ ὁ ποιήσας αὐτὸ, ὅτι ὁ μὲν εἰργάζετο, τὸ δὲ φθαρτὸν θεὸς ὠνομάσθη.
\vs{9}Ἐν ἴσῳ γὰρ μισητὰ Θεῷ καὶ ὁ ἀσεβῶν καὶ ἡ ἀσέβεια αὐτοῦ.
\vs{10}Καὶ γὰρ τὸ πραχθὲν σὺν τῷ δράσαντι κολασθησεται.
\vs{11}Διὰ τοῦτο καὶ ἐν εἰδώλοις ἐθνῶν ἐπισκοπὴ ἔσται, ὅτι ἐν κτίσματι Θεοῦ εἰς βδέλυγμα ἐγενήθησαν, καὶ εἰς σκάνδαλα ψυχαῖς ἀνθρώπων, καὶ εἰς παγίδα ποσὶνἀφρόνων.

\vs{12}Ἀρχὴ γὰρ πορνείας ἐπίνοια εἰδώλων, εὕρεσις δὲ αὐτῶν φθορὰ ζωῆς.
\vs{13}Οὔτε γὰρ ἦν ἀπʼ ἀρχῆς, οὔτε εἰς τὸν αἰῶνα ἔσται.
\vs{14}Κενοδοξίᾳ γὰρ ἀνθρώπων εἰσῆλθεν εἰς κόσμον, καὶ διὰ τοῦτο σύντομον αὐτῶν τέλος ἐπενοήθη.
\vs{15}Ἀώρῳ γὰρ πένθει τρυχόμενος πατὴρ, τοῦ ταχέως ἀφαιρεθέντος τέκνου εἰκόνα ποιήσας, τὸν τότε νεκρὸν ἄνθρωπον, νῦν ὡς θεὸν ἐτίμησεν, καὶ παρέδωκε τοῖς ὑποχειρίοις μυστήρια καὶ τελετάς.
\vs{16}Εἶτα ἐν χρόνῳ κρατυνθὲν τὸ ἀσεβὲς ἔθος ὡς νόμος ἐφυλάχθη,
\vs{17}καὶ τυράννων ἐπιταγαῖς ἐθρησκεύετο τὰ γλυπτά· οὓς ἐν ὄψει μὴ δυνάμενοι τιμᾷν ἄνθρωποι διὰ τὸ μακρὰν οἰκεῖν, τὴν πόῤῥωθεν ὄψιν ἀνατυπωσάμενοι, ἐμφανῆ εἰκόνα τοῦ τιμωμένου βασιλέως ἐποίησαν, ἵνα τὸν ἀπόντα ὡς παρόντα κολακεύωσι διὰ τῆς σπουδῆς.

\vs{18}Εἰς ἐπίτασιν δὲ θρησκείας καὶ τοὺς ἀγνοοῦντας ἡ τοῦ τεχνίτου προετρέψατο φιλοτιμία.
\vs{19}Ὁ μὲν γὰρ τάχα τῷ κρατοῦντι βουλόμενος ἀρέσαι, ἐξεβιάσατο τῇ τέχνῃ τὴν ὁμοιότητα ἐπὶ τὸ κάλλιον.
\vs{20}Τὸ δὲ πλῆθος ἐφελκόμενον διὰ τὸ εὔχαρι τῆς ἐργασίας, τὸν πρὸ ὀλίγου τιμηθέντα ἄνθρωπον, νῦν σέβασμα ἐλογίσαντο.
\vs{21}Καὶ τοῦτο ἐγένετο τῷ βίῳ εἰς ἔνεδρον, ὅτι ἢ συμφορᾷ ἢ τυραννίδι δουλεύσαντες ἄνθρωποι, τὸ ἀκοινώνητον ὄνομα λίθοις καὶ ξύλοις περιέθεσαν.

\vs{22}Εἶτʼ οὐκ ἤρκεσε τὸ πλανᾶσθαι περὶ τὴν τοῦ Θεοῦ γνῶσιν, ἀλλὰ καὶ μεγάλῳ ζῶντες ἀγνοίας πολέμῳ, τὰ τοσαῦτα κακὰ εἰρήνην προσαγορεύουσιν.
\vs{23}Ἢ γὰρ τεκνοφόνους τελετὰς, ἢ κρύφια μυστήρια, ἢ ἐμμανεῖς ἐξ ἄλλων θεσμῶν κώμους ἄγοντες,
\vs{24}οὔτε βίους οὔτε γάμους καθαροὺς ἔτι φυλάσσουσιν, ἕτερος δʼ ἕτερον ἢ λοχῶν ἀναιρεῖ, ἢ νοθεύων ὀδυνᾷ.
\vs{25}Πάντα δʼ ἐπιμὶξ ἔχει αἷμα καὶ φόνος, κλοπὴ καὶ δόλος, φθορὰ, ἀπιστία, ταραχὴ, ἐπιορκία,
\vs{26}θόρυβος ἀγαθῶν, χάριτος ἀμνησία, ψυχῶν μιασμὸς, γενέσεως ἐναλλαγὴ, γάμων ἀταξία, μοιχεία, καὶ ἀσέλγεια.
\vs{27}Ἡ γὰρ τῶν ἀνωνύμων εἰδώλων θρησκεία παντὸς ἀρχὴ κακοῦ καὶ αἰτία καὶ πέρας ἐστίν.
\vs{28}Ἢ γὰρ εὐφραινόμενοι μεμῄνασιν, ἢ προφητεύουσι ψευδῆ, ἢ ζῶσιν ἀδίκως, ἢ ἐπιορκοῦσι ταχέως.
\vs{29}Ἀψύχοις γὰρ πεποιθότες εἰδώλοις, κακῶς ὀμόσαντες, ἀδικηθῆναι οὐ προσδέχονται.

\vs{30}Ἀμφότερα δὲ αὐτοὺς μετελεύσεται τὰ δίκαια, ὅτι κακῶς ἐφρόνησαν περὶ Θεοῦ προσχόντες εἰδώλοις, καὶ ἀδίκως ὤμοσαν ἐν δόλῳ καταφρονήσαντες ὁσιότητος.
\vs{31}Οὐ γὰρ ἡ τῶν ὀμνυομένων δύναμις, ἀλλʼ ἡ τῶν ἁμαρτανόντων δίκη ἐπεξέρχεται ἀεὶ τὴν τῶν ἀδίκων παράβασιν.

\ch{15}
Σὺ δὲ ὁ Θεὸς ἡμῶν χρηστὸς καὶ ἀληθὴς, μακρόθυμος καὶ ἐν ἐλέει διοικῶν τὰ πάντα.
\vs{2}Καὶ γὰρ ἐὰν ἁμάρτωμεν, σοί ἐσμεν, εἰδότες σου τὸ κράτος· οὐχ ἁμαρτησόμεθα δὲ, εἰδότες ὅτι σοὶ λελογίσμεθα.
\vs{3}Τὸ γὰρ ἐπίστασθαί σε ὁλόκληρος δικαιοσύνη, καὶ εἰδέναι τὸ κράτος σου ῥίζα ἀθανασίας.
\vs{4}Οὔτε γὰρ ἐπλάνησεν ἡμᾶς ἀνθρώπων κακότεχνος ἐπίνοια, οὐδὲ σκιαγράφων πόνος ἄκαρπος, εἶδος σπιλωθὲν χρώμασι διηλλαγμένοις.
\vs{5}ὧν ὄψις ἄφροσιν εἰς ὄνειδος ἔρχεται, ποθεῖ, τε νεκρᾶς εἰκόνος εἶδος ἄπνουν.

\vs{6}Κακῶν ἐρασταὶ ἄξιοί τε τοιούτων ἐλπίδων, καὶ οἱ δρῶντες, καὶ οἱ ποθοῦντες, καὶ οἱ σεβόμενοι.
\vs{7}Καὶ γὰρ κεραμεὺς ἁπαλὴν γῆν θλίβων ἐπίμοχθον, πλάσσει πρὸς ὑπηρεσίαν ἡμῶν ἕκαστον. ἀλλʼ ἐκ τοῦ αὐτοῦ πηλοῦ ἀνεπλάσατο τά τε τῶν καθαρῶν ἔργων δοῦλα σκεύη, τά τε ἐναντία, πάνθʼ ὁμοίως· τούτων δὲ ἑκατέρου τίς ἑκάστου ἐστὶν ἡ χρῆσις, κριτὴς ὁ πηλουργός.
\vs{8}Καὶ κακόμοχθος θεὸν μάταιον ἐκ τοῦ αὐτοῦ πλάσσει πηλοῦ, ὃς πρὸ μικροῦ ἐκ γῆς γεννηθεὶς μετʼ ὀλίγον πορεύεται ἐξ ἧς ἐλήφθη, τὸ τῆς ψυχῆς ἀπαιτηθεὶς χρέος.

\vs{9}Ἀλλʼ ἔστιν αὐτῷ φροντὶς οὐχ ὅτι μέλλει κάμνειν, οὐδʼ ὅτι βραχυτελῆ βίον ἔχει, ἀλλʼ ἀντερείδεται μὲν χρυσουργοῖς καὶ ἀργυροχόοις, χαλκοπλάστας τε μιμεῖται, καὶ δόξαν ἡγεῖται ὅτι κίβδηλα πλάσσει.
\vs{10}Σποδὸς ἡ καρδία αὐτοῦ, καὶ γῆς εὐτελεστέρα ἡ ἐλπὶς αὐτοῦ, πηδοῦ τε ἀτιμότερος ὁ βίος αὐτον·
\vs{11}ὅτι ἠγνόησε τὸν πλάσαντα αὐτὸν, καὶ τὸν ἐμπνεύσαντα αὐτῷ ψυχὴν ἐνεργοῦσαν, καὶ ἐμφυσήσαντα πνεῦμα ζωτίκον.
\vs{12}Ἀλλʼ ἐλογίσαντο παίγνιον εἶναι τὴν ζωὴν ἡμῶν, καὶ τὸν βίον πανεγυρισμὸν ἐπικερδῆ· δεῖν γάρ φησιν ὅθεν δὴ κᾂν ἐκ κακοῦ πορίζειν.
\vs{13}Οὗτος γὰρ παρὰ πάντας οἶδεν ὅτι ἁμαρτάνει, ὕλης γεώδους εὔθραυστα σκεύη καὶ γλυπτὰ δημιουργῶν.

\vs{14}Πάντες δʼ ἀφρονέστατοι καὶ τάλανες ὑπὲρ ψυχὴν νηπίου, οἱ ἐχθροὶ τοῦ λαοῦ σου καταδυναστεύσαντες αὐτόν.
\vs{15}Ὅτι καὶ πάντα εἴδωλα τῶν ἐθνῶν ἐλογίσαντο θεοὺς, οἷς οὔτε ὁμμάτων χρῆσις εἰς ὅρασιν, οὔτε ῥῖνες εἰς συνολκὴν ἀέρος, οὔτε ὦτα ἀκούειν, οὔτε δάκτυλοι χειρῶν εἰς ψηλάφησιν, καὶ οἱ πόδες αὐτῶν ἀργοὶ πρὸς ἐπίβασιν.
\vs{16}Ἄνθρωπος γὰρ ἐποίησεν αὐτοὺς, καὶ τὸ πνεῦμα δεδανεισμένος ἔπλασεν αὐτούς· οὐδεὶς γὰρ αὐτῷ ὅμοιον ἄνθρωπος ἰσχύει πλάσαι θεόν·
\vs{17}θνητὸς δὲ ὢν νεκρὸν ἐργάζεται χερσὶν ἀνόμοις· κρείττων γάρ ἐστι τῶν σεβασμάτων αὐτοῦ, ὧν αὐτὸς μὲν ἔζησεν, ἐκεῖνα δὲ οὐδέποτε.
\vs{18}Καὶ τὰ ζῶα δὲ τὰ ἔχθιστα σέβονται, ἄνοια γὰρ συγκρινόμενα τῶν ἄλλων ἐστὶ χείρονα.
\vs{19}Οὐδʼ ὅσον ἐπιποθῆσαι ὡς ἐν ζώων ὄψει καλὰ τυγχάνει, ἐκπέφευγε δὲ καὶ τὸν τοῦ Θεοῦ ἔπαινον καὶ τὴν εὐλογίαν αὐτοῦ.

\ch{16}
Διὰ τοῦτο διʼ ὁμοίων ἐκολάσθησαν ἀξίως, καὶ διὰ πλήθους κνωδάλον ἐβασανίσθησαν.
\vs{2}Ἀνθʼ ἧς κολάσεως εὐεργετήσας τὸν λαόν σου, εἰς ἐπιθυμίαν ὀρέξεως ξένην γεῦσιν, τροφὴν ἡτοίμασας ὀρτυγομήτραν,
\vs{3}ἵνα ἐκεῖνοι μὲν ἐπιθυμοῦντες τροφὴν, διὰ τὴν εἰδέχθειαν τῶν ἐπαπεσταλμένων καὶ τὴν ἀναγκαίαν ὄρεξιν ἀποστρέφωνται, αὐτοὶ δὲ ἐπʼ ὀλίγον ἐνδεεῖς γενόμενοι καὶ ξένης μετάσχωσι γεύσεως.
\vs{4}Ἔδει γὰρ ἐκείνοις μὲν ἀπαραίτητον ἔνδειαν ἐπελθεῖν τυραννοῦσι, τούτοις δὲ μόνον δειχθῆναι πῶς οἱ ἐχθροὶ αὐτῶν ἐβασανίζοντο.
\vs{5}Καὶ γὰρ ὅτε αὐτοῖς δεινὸς ἐπῆλθε θηρίων θυμὸς, δήγμασί τε σκολιῶν διεφθείροντο ὄφεων, οὐ μέχρι τέλους ἔμεινεν ἡ ὀργή σου.

\vs{6}Εἰς νουθεσίαν δὲ πρὸς ὀλίγον ἐταράχθησαν, σύμβολον ἔχοντες σωτηρίας, εἰς ἀνάμνησιν ἐντολῆς νόμου σου.
\vs{7}Ὁ γὰρ ἐπιστραφεὶς οὐ διὰ τὸ θεωρούμενον ἐσώζετο, ἀλλὰ διὰ σὲ τὸν πάντων σωτῆρα.
\vs{8}Καὶ ἐν τούτῳ δὲ ἔπεισας τοὺς ἐχθροὺς ἡμῶν, ὅτι σὺ εἶ ὁ ῥυόμενος ἐκ παντὸς κακοῦ.
\vs{9}Οὓς μὲν γὰρ ἀκρίδων καὶ μυιῶν ἀπέκτεινε δήγματα, καὶ οὐχ εὑρέθη ἴαμα τῇ ψυχῇ αὐτῶν, ὅτι ἄξιοι ησαν ὑπὸ τοιούτων κολασθῆναι.
\vs{10}Τοὺς δὲ υἱούς σου οὐδὲ ἰοβόλων δρακόντων ἐνίκησαν ὀδόντες, τὸ ἔλεος γάρ σου ἀντιπαρῆλθε καὶ ἰάσατο αὐτούς.
\vs{11}Εἰς γὰρ ὑπόμνησιν τῶν λογίων σου ἐνεκεντρίζοντο, καὶ ὀξέως διεσώζοντο, ἵνα μὴ εἰς βαθεῖαν ἐμπεσόντες λήθην, ἀπερίσπαστοι γένωνται τῆς σῆς εὐεργεσίας.

\vs{12}Καὶ γὰρ οὔτε βοτάνη οὔτε μάλαγμα ἐθεράπευσεν αὐτοὺς, ἀλλὰ ὁ σός Κύριε λόγος ὁ πάντα ἰώμενος.
\vs{13}Σὺ γὰρ ζωῆς καὶ θανάτου ἐξουσίαν ἔχεις, καὶ κατάγεις εἰς πύλας ᾅδου καὶ ἀνάγεις.
\vs{14}Ἄνθρωπος δὲ ἀποκτέννει μὲν τῇ κακίᾳ αὐτοῦ, ἐξελθὸν δὲ πνεῦμα οὐκ ἀναστρέφει, οὐδὲ ἀναλύει ψυχὴν παραληφθεῖσαν.
\vs{15}Τὴν δὲ σὴν χεῖρα φυγεῖν ἀδύνατόν ἐστιν.

\vs{16}Ἀρνούμενοι γάρ σε εἰδέναι ἀσεβεῖς, ἐν ἰσχύϊ βραχίονός σου ἐμαστιγώθησαν, ξένοις ὑετοῖς καὶ χαλάζαις καὶ ὄμβροις διωκόμενοι ἀπαραιτήτοις, καὶ πυρὶ καταναλισκόμενοι.

\vs{17}Τὸ γὰρ παραδοξότατον, ἐν τῷ πάντα σβεννῦντι ὕδατι πλεῖον ἐνήργει τὸ πῦρ· ὑπέρμαχος γὰρ ὁ κόσμος ἐστὶ δικαίων.
\vs{18}Ποτὲ μὲν γὰρ ἡμεροῦτο φλὸξ, ἵνα μὴ καταφλέξῃ τὰ ἐπʼ ἀσεβεῖς ἀπεσταλμένα ζῶα, ἀλλʼ αὐτοὶ βλέποντες ἴδωσιν, ὅτι Θεοῦ κρίσει ἐλαύνονται.
\vs{19}Ποτὲ δὲ καὶ μεταξὺ ὕδατος ὑπὲρ τὴν πυρὸς δυναμιν φλέγει, ἵνα ἀδίκου γῆς γεννήματα διαφθείρῃ.
\vs{20}Ἀνθʼ ὧν ἀγγέλων τροφὴν ἐψώμισας τὸ λαόν σου, καὶ ἕτοιμον ἄρτον αὐτοῖς ἀπʼ οὐρανοῦ ἔπεμψας ἀκοπιάτως, πᾶσαν ἡδονὴν ἰσχύοντα καὶ πρὸς πᾶσαν ἁρμόνιον γεῦσιν.
\vs{21}Ἡ μὲν γὰρ ὑπόστασίς σου τὴν σὴν γλυκύτητα πρὸς τέκνα ἐνεφάνισε, τῇ δὲ τοῦ προσφερομένου ἐπιθυμίᾳ ὑπηρετῶν, πρὸς ὅ τις ἐβούλετο μετεκιρνᾶτο.
\vs{22}Χιὼν δὲ καὶ κρύσταλλος ὑπέμεινε πῦρ, καὶ οὐκ ἐτήκετο, ἵνα γνῶσιν ὅτι τοὺς τῶν ἐχθρῶν καρποὺς κατέφθειρε πῦρ φλεγομένον, ἐν τῇ χαλάζῃ καὶ ἐν τοῖς ὑετοῖς διαστράπτον.

\vs{23}Τοῦτο πάλιν δʼ ἵνα τραφῶσι δίκαιοι, καὶ τῆς ἰδίας ἐπιλελῆσθαι δυνάμεως.
\vs{24}Ἡ γὰρ κτίσις σοι τῷ ποιήσαντι ὑπηρετοῦσα, ἐπιτείνεται εἰς κόλασιν κατὰ τῶν ἀδίκων, καὶ ἀνίεται εἰς εὐεργεσίαν ὑπὲρ τῶν εἰς σὲ πεποιθότων.

\vs{25}Διὰ τοῦτο καὶ τότε εἰς πάντα μεταλλευομένη, τῇ παντοτρόφῳ σου δωρεᾷ ὑπηρετεῖ, πρὸς τὴν τῶν δεομένων θέλησιν·
\vs{26}ἵνα μάθωσιν οἱ υἱοί σου, οὓς ἠγάπησας, Κύριε, ὅτι οὐχ αἱ γενέσεις τῶν καρπῶν τρέφουσιν ἄνθρωπον, ἀλλὰ τὸ ῥῆμά σου τούς σοι πιστεύοντας διατηρεῖ.

\vs{27}Τὸ γὰρ ὑπὸ πυρὸς μὴ φθειρόμενον, ἁπλῶς ὑπὸ βραχείας ἀκτῖνος ἡλίου θερμαινόμενον ἐτήκετο·
\vs{28}ὅπως γνωστὸν ᾖ, ὅτι δεῖ φθάνειν τὸν ἥλιον ἐπʼ εὐχαριστίαν σου, καὶ πρὸς ἀνατολὴν φωτὸς ἐντυγχάνειν σοι.
\vs{29}Ἀχαρίστου γὰρ ἐλπὶς ὡς χειμέριος πάχνη τακήσεται, καὶ ῥυήσεται ὡς ὕδωρ ἄχρηστον.

\ch{17}
Μεγάλαι γάρ σου αἱ κρίσεις καὶ δυσδιήγητοι· διὰ τοῦτο ἀπαίδευτοι ψυχαὶ ἐπλανήθησαν.
\vs{2}Ὑπειληφότες γὰρ καταδυναστεύειν ἔθνος ἅγιον ἄνομοι, δέσμιοι σκότους καὶ μακρᾶς πεδῆται νυκτὸς, κατακλεισθέντες ὀρόφοις, φυγάδες τῆς αἰωνίου προνοίας ἔκειντο.
\vs{3}Λανθάνειν γὰρ νομίζοντες ἐπὶ κρυφαίοις ἁμαρτήμασιν, ἀφεγγεῖ λήθης παρακαλύμματι ἐσκορπίσθησαν, θαμβούμενοι δεινῶς καὶ ἰνδάλμασιν ἐκταρασσόμενοι.
\vs{4}Οὐδὲ γὰρ ὁ κατέχων αὐτοὺς μυχὸς ἀφόβως διεφύλασσεν, ἦχοι δὲ καταράσσοντες αὐτοὺς περιεκόμπουν, και φάσματα ἀμειδήτοις κατηφῆ προσώποις ἐνεφανίζετο.

\vs{5}Καὶ πυρὸς μὲν οὐδεμία βία κατίσχυε φωτίζειν, οὔτε ἄστρων ἔκλαμπροι φλόγες καταυγάζειν ὑπέμενον τὴν στυγνὴν ἐκείνην νύκτα.
\vs{6}Διεφαίνετο δʼ αὐτοῖς μόνον αὐτομάτη πυρὰ φόβου πλήρης, ἐκδειματούμενοι δὲ τῆς μὴ θεωρουμένης ἐκείνης ὄψεως, ἡγοῦντο χείρω τὰ βλεπόμενα.
\vs{7}Μαγικῆς δὲ ἐμπαίγματα κατέκειτο τέχνης, καὶ τῆς ἐπὶ φρονήσει ἀλαζονείας ἔλεγχος ἐφύβριστος.
\vs{8}Οἱ γὰρ ὑπισχνούμενοι δείματα καὶ ταραχὰς ἀπελαύνειν ψυχῆς νοσούσης, οὗτοι καταγέλαστον εὐλάβειαν ἐνόσουν.

\vs{9}Καὶ γὰρ εἰ μηδὲν αὐτοὺς ταραχῶδες ἐφόβει, κνωδάλων παρόδοις καὶ ἑρπετῶν συριγμοῖς ἐκσεσοβημένοι,
\vs{10}διώλλυντο ἔντρομοι καὶ τὸν μηδαμόθεν φευκτὸν ἀέρα προσιδεῖν ἀρνούμενοι.

\vs{11}Δειλὸν γὰρ ἰδίως πονηρία μαρτυρεῖ καταδικαζομένη, ἀεὶ δὲ προσείληφε τὰ χαλεπὰ συνεχομένη τῇ συνειδήσει.
\vs{12}Οὐθὲν γάρ ἐστι φόβος, εἰ μὴ προδοσία τῶν ἀπὸ λογισμοῦ βοηθημάτων.
\vs{13}Ἔνδοθεν δὲ οὖσα ἥττων ἡ προσδοκία, πλείονα λογίζεται τὴν ἄγνοιαν τῆς παρεχούσης τὴν βάσανον αἰτίας.
\vs{14}Οἱ δὲ τὴν ἀδύνατον ὄντως νύκτα καὶ ἐξ ἀδυνάτου ᾅδου μυχῶν ἐπελθοῦσαν, τὸν αὐτὸν ὕπνον κοιμώμενοι,
\vs{15}τὰ μὲν τέρασιν ἠλαύνοντο φαντασμάτων, τὰ δὲ τῆς ψυχῆς παρελύοντο προδοσίᾳ· αἰφνίδιος γὰρ αὐτοῖς καὶ ἀπροσδόκητος φόβος ἐπῆλθεν.
\vs{16}Εἶθʼ οὕτως, ὃς δήποτʼ οὖν ἦν ἐκεῖ καταπίπτων, ἐφρουρεῖτο εἰς τὴν ἀσίδηρον εἱρκτὴν κατακλεισθείς.
\vs{17}Εἴτε γὰρ γεωργὸς ἦν τις, ἢ ποιμὴν, ἢ τῶν κατʼ ἐρημίαν ἐργάτης μόχθων, προληφθεὶς τὴν δυσάλυκτον ἔμενεν ἀνάγκην· μιᾷ γὰρ ἁλύσει σκότους πάντες ἐδέθησαν.

\vs{18}Εἴτε πνεῦμα συρίζον, ἢ περὶ ἀμφιλαφεῖς κλάδους ὀρνέων ἦχος εὐμελὴς, ἢ ῥυθμὸς ὕδατος πορευομένου βίᾳ, ἢ κτύπος ἀπηνῆς καταῤῥιπτομένων πετρῶν,
\vs{19}ἢ σκιρτώντων ζώων δρόμος ἀθεώρητος, ἢ ὠρυομένων ἀπηνεστάτων θηρίων φωνὴ, ἢ ἀντανακλωμένη ἐκ κοιλοτάτων ὀρέων ἠχὼ, παρέλυσεν αὐτοὺς ἐκφοβοῦντα.
\vs{20}Ὅλος γὰρ ὁ κόσμος λαμπρῷ κατελάμπετο φωτὶ, καὶ ἀνεμποδίστοις συνείχετο ἔργοις.
\vs{21}Μόνοις δὲ ἐκείνοις ἐπετέτατο βαρεῖα νὺξ, εἰκὼν τοῦ μέλλοντος αὐτοὺς διαδέχεσθαι σκότους, ἑαυτοῖς δὲ ἦσαν βαρύτεροι σκότους.

\ch{18}
Τοῖς δὲ ὁσίοις σου μέγιστον ἦν φῶς, ὧν φωνὴν μὲν ἀκούοντες, μορφὴν δὲ οὐχ ὁρῶντες, ὅτι μὲν οὖν κᾀκεῖνοι ἐπεπόνθεισαν, ἐμακάριζον,
\vs{2}ὅτι δὲ οὐ βλάπτουσι προηδικημένοι, εὐχαριστοῦσι, καὶ τοῦ διενεχθῆναι χάριν ἐδέοντο.
\vs{3}Ἀνθʼ ὧν πυριφλεγῆ στύλον, ὁδηγὸν μὲν ἀγνώστου ὁδοιπορίας, ἥλιον δὲ ἀβλαβῆ φιλοτίμου ξενιτείας παρέσχες.
\vs{4}Ἄξιοι μὲν γὰρ ἐκεῖνοι στερηθῆναι φωτὸς, καὶ φυλακισθῆναι ἐν σκότει, οἱ κατακλείστους φυλάξαντες τοὺς υἱούς σου, διʼ ὧν ἤμελλε τὸ ἄφθαρτον νόμου φῶς τῷ αἰῶνι δίδοσθαι.

\vs{5}Βουλευσαμένους δʼ αὐτοὺς τὰ τῶν ὁσιῶν ἀποκτεῖναι νήπια, καὶ ἑνὸς ἐκτεθέντος τέκνου, καὶ σωθέντος, εἰς ἔλεγχον τὸ αὐτῶν ἀφείλω πλῆθος τέκνων, καὶ ὁμοθυμαδὸν ἀπώλεσας ἐν ὕδατι σφοδρῷ.
\vs{6}Ἐκείνη ἡ νὺξ προεγνώσθη πατράσιν ἡμῶν, ἵνα ἀσφαλῶς εἰδότες οἷς ἐπίστευσαν ὅρκοις, ἐπευθυμήσωσι.

\vs{7}Προσεδέχθη δὲ ὑπὸ λαοῦ σου σωτηρία μὲν δικαίων, ἐχθρῶν δὲ ἀπώλεια.
\vs{8}Ὧς γὰρ ἐτιμωρήσω τοὺς ὑπεναντίους, τούτῳ ἡμᾶς προσκαλεσάμενος ἐδόξασας.
\vs{9}Κρυφῇ γὰρ ἐθυσιαζον ὅσιοι παῖδες ἀγαθῶν, καὶ τὸν τῆς θειότητος νόμον ἐν ὁμονοίᾳ διέθεντο, τῶν αὐτῶν ὁμοίως καὶ ἀγαθῶν καὶ κινδύνων μεταλήψεσθαι τοὺς ἁγίους, πατέρων ἤδη προαναμελπόντων αἴνους.

\vs{10}Ἀντήχει δʼ ἀσύμφωνος ἐχθρῶν βοὴ, καὶ οἰκτρὰ διεφέρετο θρηνουμένων παίδων.
\vs{11}Ὁμοίᾳ δὲ δίκῃ δοῦλος ἅμα δεσπότῃ κολασθεὶς, καὶ δημότης βασιλεῖ τὰ αὐτὰ πάσχων.

\vs{12}Ὁμοθυμαδὸν δὲ πάντες ἐν ἑνὶ ὀνόματι θανάτου νεκροὺς εἶχον ἀναριθμήτους, οὐδὲ γὰρ πρὸς τὸ θάψαι οἱ ζῶντες ἦσαν ἱκανοὶ, ἐπεὶ πρὸς μίαν ῥοπὴν ἡ ἐντιμοτέρα γένεσις αὐτῶν διεφθάρη.
\vs{13}Πάντα γὰρ ἀπιστοῦντες διὰ τὰς φαρμακίας, ἐπὶ τῷ τῶν πρωτοτόκων ὀλέθρῳ, ὡμολόγησαν Θεοῦ υἱὸν λαὸν εἶναι.
\vs{14}Ἡσύχου γὰρ σιγῆς περιεχούσης τὰ πάντα, καὶ νυκτὸς ἐν ἰδίῳ τάχει μεσαζούσης,
\vs{15}ὁ παντοδύναμός σου λόγος ἀπʼ οὐρανῶν ἐκ θρόνων βασιλειῶν, ἀπότομος πολεμιστὴς, εἰς μέσον τῆς ὀλεθρίας ἥλατο γῆς, ξίφος ὀξὺ τὴν ἀνυπόκριτον ἐπιταγήν σον φέρων,
\vs{16}καὶ στὰς ἐπλήρωσε τὰ πάντα θανάτου· καὶ οὐρανοῦ μὲν ἥπτετο, βεβήκει δʼ ἐπὶ γῆς.
\vs{17}Τότε παραχρῆμα φαντασίαι μὲν ὀνείρων δεινῶς ἐξετάραξαν αὐτοὺς, φόβοι δὲ ἐπέστησαν ἀδόκητοι
\vs{18}καὶ ἄλλος ἀλλαχῇ ῥιφεὶς ἡμίθνητος, διʼ ἣν ἔθνησκεν αἰτίαν ἐνεφάνιζεν.
\vs{19}Οἱ γὰρ ὄνειροι θορυβήσαντες αὐτοὺς, τοῦτο προεμήνυσαν, ἵνα μὴ ἀγνοοῦντες διʼ ὃ κακῶς πάσχουσιν, ἀπόλωνται.

\vs{20}Ἥψατο δὲ καὶ δικαίων πεῖρα θανάτου, καὶ θραῦσις ἐν ἐρήμῳ ἐγένετο πλήθους· ἀλλʼ οὐκ ἐπὶ πολὺ ἔμεινεν ἡ ὀργή.
\vs{21}Σπεύσας γὰρ ἀνὴρ ἄμεμπτος προεμάχησε· τὸ τῆς ἰδίας λειτουργίας ὅπλον, προσευχὴν καὶ θυμιάματος ἐξιλασμὸν κομίσας, ἀντέστη τῷ θυμῷ, καὶ πέρας ἐπέθηκε τῇ συμφορᾷ, δεικνὺς ὅτι σός ἐστι θεράπων.

\vs{22}Ἐνίκησε δὲ τὸν ὄχλον οὐκ ἰσχύϊ τοῦ σώματος, οὐχ ὅπλων ἐνεργείᾳ, ἀλλὰ λόγῳ τὸν κολάζοντα ὑπέταξεν, ὅρκους πατέρων καὶ διαθήκας ὑπομνήσας.
\vs{23}Σωρηδὸν γὰρ ἤδη πεπτωκότων ἐπʼ ἀλλήλων νεκρῶν, μεταξὺ στὰς, ἀνέκοψε τὴν ὀργὴν, καὶ διέσχισε τὴν πρὸς τοὺς ζῶντας ὁδόν.
\vs{24}Ἐπὶ γὰρ ποδήρους ἐνδύματος ἦν ὅλος ὁ κόσμος, καὶ πατέρων δόξαι ἐπὶ τετραστίχου λίθου γλυφῆς, καὶ μεγαλωσύνη σου ἐπὶ διαδήματος κεφαλῆς αὐτοῦ.
\vs{25}Τούτοις εἶξεν ὁ ὀλοθρεύων, ταῦτα δὲ ἐφοβήθησαν· ἦν γὰρ μόνη ἡ πεῖρα τῆς ὀργῆς ἱκανή.

\ch{19}
Τοῖς δὲ ἀσεβέσι μέχρι τέλους ἀνελεήμων θυμὸς ἐπέστη· προῄδει γὰρ αὐτῶν καὶ τὰ μέλλοντα,
\vs{2}ὅτι αὐτοὶ ἐπιστρέψαντες τοῦ ἀπεῖναι, καὶ μετὰ σπουδῆς προπέμψαντες αὐτοὺς, διώξουσι μεταμεληθέντες.
\vs{3}Ἔτι γὰρ ἐν χερσὶν ἔχοντες τὰ πὲνθη, καὶ προσοδυρόμενοι τάφοις νεκρῶν, ἕτερον ἐπεσπάσαντο λογισμὸν ἀνοίας, καὶ οὓς ἱκετεύοντες ἐξέβαλον, τούτους ὡς φυγάδας ἐδίωκον.
\vs{4}Εἷλκε γὰρ αὐτοὺς ἡ ἀξία ἐπὶ τοῦτο τὸ πέρας ἀνάγκη, καὶ τῶν συμβεβηκότων ἀμνηστίαν ἐνέβαλεν, ἵνα τὴν λείπουσαν ταῖς βασάνοις προαναπληρώσωσιν κόλασιν·
\vs{5}καὶ ὁ μὲν λαός σου παράδοξον ὁδοιπορίαν περάσῃ, ἐκεῖνοι δὲ ξένον εὕρωσι θάνατον.

\vs{6}Ὅλη γὰρ ἡ κτίσις ἐν ἰδίῳ γένει πάλιν ἄνωθεν διετυποῦτο, ὑπηρετοῦσα ταῖς ἰδίαις ἐπιταγαῖς, ἵνα οἱ σοὶ παῖδες φυλαχθῶσιν ἀβλαβεῖς.
\vs{7}Ἡ τὴν παρεμβολὴν σκιάζουσα νεφέλη, ἐκ δὲ προϋφεστῶτος ὕδατος ξηρᾶς ἀνάδυσις γῆς ἐθεωρήθη, ἐξ ἐρυθρᾶς θαλάσσης ὁδὸς ἀνεμπόδιστος, καὶ χλοηφόρον πεδίον ἐκ κλύδωνος βιαίου,
\vs{8}διʼ οὗ πανεθνὶ διῆλθον οἱ τῇ σῇ σκεπαζόμενοι χειρί, θεωρήσαντες θαυμαστὰ τέρατα.
\vs{9}Ὡς γὰρ ἵπποι ἐνεμήθησαν, καὶ ὡς ἀμνοὶ διεσκίρτησαν, αἰνοῦντές σε, Κύριε, τὸν ῥυόμενον αὐτούς.
\vs{10}Ἐμέμνηντο γὰρ ἔτι τῶν ἐν τῇ παροικίᾳ αὐτῶν, πῶς ἀντὶ μὲν γενέσεως ζώων ἐξήγαγεν ἡ γῆ σκνῖπα, ἀντὶ δὲ ἐνύδρων ἐξηρεύξατο ὁ ποταμὸς πλῆθος βατράχων.

\vs{11}Ἐφʼ ὑστέρῳ δὲ εἶδον καὶ νέαν γένεσιν ὀρνέων, ὅτι ἐπιθυμίᾳ προαχθέντες ᾐτήσαντο ἐδέσματα τρυφῆς.
\vs{12}Εἰς γὰρ παραμυθίαν ἀνέβη αὐτοῖς ἀπὸ θαλάσσης ὀρτυγομήτρα,
\vs{13}καὶ αἱ τιμωρίαι τοῖς ἁμαρτωλοῖς ἐπῆλθον, οὐκ ἄνευ τῶν γεγονότων τεκμήριων τῇ βίᾳ τῶν κεραυνῶν· δικαίως γὰρ ἔπασχον ταῖς ἰδίαις αὐτῶν πονηρίαις· καὶ γὰρ χαλεπωτέραν μισοξενίαν ἐπετήδευσαν.
\vs{14}Οἱ μὲν γὰρ τοὺς ἀγνοοῦντας οὐκ ἐδέχοντο παρόντας, οὗτοι δὲ εὐεργέτας ξένους ἐδουλοῦντο.
\vs{15}Καὶ οὐ μόνον, ἀλλʼ ἥτις ἐπισκοπὴ ἔσται αὐτῶν, ἐπεὶ ἀπεχθῶς προσεδέχοντο τοὺς ἀλλοτρίους·
\vs{16}οἱ δὲ μετὰ ἑορτασμάτων εἰσδεξάμενοι τοὺς ἤδη τῶν αὐτῶν μετεσχηκότας δικαίων δεινοῖς ἐκάκωσαν πόνοις.
\vs{17}Ἐπλήγησαν δὲ καὶ ἀορασίᾳ, ὥσπερ ἐκεῖνοι ἐπὶ ταῖς τοῦ δικαίου θύραις, ὅτε ἀχανεῖ περιβληθέντες σκότει, ἕκαστος τῶν αὐτοῦ θυρῶν τὴν δίοδον ἐζήτει.

\vs{18}Διʼ ἑαυτῶν γὰρ τὰ στοιχεῖα μεθαρμοζόμενα, ὥσπερ ἐν ψαλτηρίῳ φθόγγοι τοῦ ῥυθμοῦ τὸ ὄνομα διαλλάσσουσι, πάντοτε μένοντα ἤχῳ, ὅπερ ἐστὶν εἰκάσαι ἐκ τῆς τῶν γεγονότων ὄψεως ἀκριβῶς.
\vs{19}Χερσαῖα γὰρ εἰς ἔνυδρα μετεβάλλετο, καὶ νηκτὰ μετέβαινεν ἐπὶ γῆς.
\vs{20}Πῦρ ἴσχυσεν ἐν ὕδατι τῆς ἰδίας δυνάμεως, καὶ ὕδωρ τῆς σβεστικῆς δυνάμεως ἐπελανθάνετο.
\vs{21}Φλόγες ἀνάπαλιν εὐφθάρτων ζώων οὐκ ἐμάραναν σάρκας ἐμπεριπατούντων, οὐδὲ τηκτὸν εὔτηκτον κρυσταλλοειδὲς γένος ἀμβροσίας τροφῆς.
\vs{22}Κατὰ πάντα γὰρ, Κύριε, ἐμεγάλυνας τὸν λαόν σου, καὶ ἐδόξασας, καὶ οὐχ ὑπερεῖδες, ἐν παντὶ καιρῷ καὶ τόπῳ παριστάμενος.


\def\book{ΣΟΦΙΑ ΣΕΙΡΑΧ}
\biblebook{ΣΟΦΙΑ ΣΕΙΡΑΧ}

\lettrine[lines=2, loversize=0.2, nindent=0em, findent=.25em]{\textcolor{bookheadingcolor}{Π}}{ΡΟΛΟΓΟΣ.}


\postdropcapindent\vs{1}ΠΟΛΛΩΝ καὶ μεγάλων ἡμῖν διὰ τοῦ νόμου καὶ τῶν προφητῶν καὶ τῶν ἄλλων τῶν κατʼ αὐτοὺς ἠκολουθηκότων δεδομένων. ὑπὲρ ὧν δέον ἐστὶν ἐπαινεῖν τὸν Ἰσραὴλ παιδείας καὶ σοφίας, καὶ ὡς οὐ μόνον αὐτοὺς τοὺς ἀναγινώσκοντας δέον ἐστὶν ἐπιστήμονας γίνεσθαι,
\vs{1a}ἀλλὰ καὶ τοῖς ἐκτὸς δύνασθαι τοὺς φιλομαθοῦντας χρησίμους εἶναι καὶ λέγοντας καὶ γράφοντας· ὁ πάππος μου Ἰησοῦς ἐπὶ πλεῖον ἑαυτὸν δοὺς εἴς τε τὴν τοῦ νόμου καὶ τῶν προφητῶν
\vs{1b}καὶ τῶν ἄλλων πατρίων βιβλίων ἀνάγνωσιν, καὶ ἐν τούτοις ἱκανὴν ἕξιν περιποιησάμενος, προήχθη καὶ αὐτὸς συγγράψαι τι τῶν εἰς παιδείαν καὶ σοφίαν ἀνηκόντων, ὅπως οἱ φιλομαθεῖς, καὶ τούτων ἔνοχοι γενόμενοι, πολλῷ μᾶλλον ἐπιπροσθῶσι διὰ τῆς ἐννόμου βιώσεως.

\vs{1c}Παρακέκλησθε οὖν μετʼ εὐνοίας καὶ προσοχῆς τὴν ἀνάγνωσιν ποιεῖσθαι, καὶ συγγνώμην ἔχειν ἐφʼ οἷς ἂν δοκῶμεν
\vs{1d}τῶν κατὰ τὴν ἑρμηνείαν πεφιλοπονημένων τισὶ τῶν λέξεων ἀδυναμεῖν· οὐ γὰρ ἰσοδυναμεῖ αὐτὰ ἐν ἑαυτοῖς Ἑβραϊστὶ λεγόμενα, καὶ ὅταν μεταχθῇ εἰς ἑτέραν γλῶσσαν. Οὐ μόνον δὲ ταῦτα, ἀλλὰ καὶ αὐτὸς ὁ νόμος, καὶ αἱ προφητεῖαι,
\vs{1e}καὶ τὰ λοιπὰ τῶν βιβλίων οὐ μικρὰν ἔχει τὴν διαφορὰν ἐν ἑαυτοῖς λεγόμενα. Ἐν γὰρ τῷ ὀγδόῳ καὶ τριακοστῷ ἔτει ἐπὶ τοῦ Εὐεργέτου βασιλέως παραγενηθεὶς εἰς Αἴγυπτον καὶ συγχρονίσας, εὗρον οὐ μικρᾶς παιδείας ἀφόμοιον·
\vs{1f}ἀναγκαιότατον ἐθέμην αὐτὸς προσενέγκασθαί τινα σπουδὴν καὶ φιλοπονίαν τοῦ μεθερμηνεῦσαι τήνδε τὴν βίβλον· πολλὴν γὰρ ἀγρυπνίαν καὶ ἐπιστήμην προσενεγκάμενος ἐν τῷ διαστήματι τοῦ χρόνου, πρὸς τὸ ἐπὶ πέρας ἄγοντα τὸ βιβλίον ἐκδόσθαι, καὶ τοῖς ἐν τῇ παροικίᾳ βουλομένοις φιλομαθεῖν,
\vs{1g}προκατασκευαζομένοις τὰ ἤθη ἐν νόμῳ βιοτεύειν.
ΣΟΦΙΑ ΣΕΙΡΑΧ.

\vs{1h}ΠΑΣΑ σοφία παρὰ Κυρίου, καὶ μετʼ αὐτοῦ ἐστιν εἰς τὸν αἰῶνα.
\vs{2}Ἄμμον θαλασσῶν καὶ σταγόνας ὑετοῦ καὶ ἡμέρας αἰῶνος τίς ἐξαριθμήσει;
\vs{3}Ὕψος οὐρανοῦ καὶ πλάτος γῆς καὶ ἄβυσσον καὶ σοφίαν τίς ἐξιχνιάσει;

\vs{4}Προτέρα πάντων ἔκτισται σοφία, καὶ σύνεσις φρονήσεως ἐξ αἰῶνος.
\vs{6}Ῥίζα σοφίας τίνι ἀπεκαλύφθη; καὶ τὰ πανουργεύματα αὐτῆς τίς ἔγνω;

\vs{8}Εἷς ἐστι σοφὸς φοβερὸς σφόδρα καθήμενος ἐπὶ τοῦ θρόνου αὐτοῦ·
\vs{9}Κύριος αὐτὸς ἔκτισεν αὐτὴν, καὶ εἶδε καὶ ἐξηρίθμησεν αὐτὴν, καὶ ἐξέχεεν αὐτὴν ἐπὶ πάντα τὰ ἔργα αὐτοῦ.
\vs{10}Μετὰ πάσης σαρκὸς κατὰ τὴν δόσιν αὐτοῦ, καὶ ἐχορήγησεν αὐτὴν τοῖς ἀγαπῶσιν αὐτόν·

\vs{11}Φόβος Κυρίου δόξα καὶ καύχημα καὶ εὐφροσύνη καὶ στέφανος ἀγαλλιάματος.
\vs{12}Φόβος Κυρίου τέρψει καρδίαν, καὶ δώσει εὐφροσύνην καὶ χαρὰν καὶ μακροημέρευσιν.
\vs{13}Τῷ φοβουμένῳ τὸν Κύριον εὖ ἔσται ἐπʼ ἐσχάτων, καὶ ἐν ἡμέρᾳ τελευτῆς αὐτοῦ εὑρήσει χάριν.

\vs{14}Αρχὴ σοφίας φοβεῖσθαι τὸν Θεὸν, καὶ μετὰ πιστῶν ἐν μήτρᾳ συνεκτίσθη αὐτοῖς.
\vs{15}Καὶ μετὰ ἀνθρώπων θεμέλιον αἰῶνος ἐνόσσευσε, καὶ μετὰ τοῦ σπέρματος αὐτῶν ἐμπιστευθήσεται.
\vs{16}Πλησμονὴ σοφίας φοβεῖσθαι τὸν Κύριον, καὶ μεθύσκει αὐτοὺς ἀπὸ τῶν καρπῶν αὐτῆς.
\vs{17}Πάντα τὸν οἶκον αὐτῆς ἐμπλήσει ἐπιθυμημάτων, καὶ τὰ ἀποδοχεῖα ἀπὸ τῶν γεννημάτων αὐτῆς.
\vs{18}Στέφανος σοφίας φόβος Κυρίου, ἀναθάλλων εἰρήνην καὶ ὑγίειαν ἰάσεως·
\vs{19}ἐπιστήμην καὶ γνῶσιν συνέσεως ἐξώμβρησε, καὶ δόξαν κρατούντων αὐτῆς ἀνύψωσε.

\vs{20}Ῥίζα σοφίας φοβεῖσθαι τὸν Κύριον, καὶ οἱ κλάδοι αὐτῆς μακροημέρευσις.
\vs{22}Οὐ δυνήσεται θυμὸς ἄδικος δικαιωθῆναι· ἡ γὰρ ῥοπὴ τοῦ θυμοῦ αὐτοῦ πτῶσις αὐτῷ.
\vs{23}Ἕως καιροῦ ἀνθέξεται μακρόθυμος, καὶ ὕστερον αὐτῷ ἀναδώσει εὐφροσύνη.
\vs{24}Ἕως καιροῦ κρύψει τοὺς λόγους αὐτοῦ, καὶ χείλη πιστῶν ἐκδιηγήσεται σύνεσιν αὐτοῦ.
\vs{25}Ἐν θησαυροῖς σοφίας παραβολὴ ἐπιστήμης, βδέλυγμα δὲ ἁμαρτωλῷ θεοσέβεια.
\vs{26}Ἐπεθύμησας σοφίαν, διατήρησον ἐντολὰς, καὶ Κύριος χορηγήσει σοι αὐτήν.
\vs{27}Σοφία γὰρ καὶ παιδεία φόβος Κυρίου, καὶ ἡ εὐδοκία αὐτοῦ πίστις καὶ πρᾳότης.

\vs{28}Μὴ ἀπειθήσῃς φόβῳ Κυρίου, καὶ μὴ προσέλθῃς αὐτῷ ἐν καρδίᾳ δισσῇ.
\vs{29}Μὴ ὑποκριθῇς ἐν στόμασιν ἀνθρώπων, καὶ ἐν τοῖς χείλεσί σου πρόσεχε.
\vs{30}Μὴ ἐξύψου σεαυτὸν, ἵνα μὴ πέσῃς, καὶ ἐπαγάγῃς τῇ ψυχῇ σου ἀτιμίαν· καὶ ἀποκαλύψει Κύριος τὰ κρυπτά σου, καὶ ἐν μέσῳ συναγωγῆς καταβελεῖ σε· ὅτι οὐ προσῆλθες φόβῳ Κυρίου, καὶ ἡ καρδία σου πλήρης δόλου.

\ch{2}
Τεκνὸν εἰ προσέρχῃ δουλεύειν Κυρίῳ Θεῷ, ἑτοίμασον τὴν ψυχήν σου εἰς πειρασμόν.
\vs{2}Εὔθυνον τὴν καρδίαν σου καὶ καρτέρησον, καὶ μὴ σπεύσῃς ἐν καιρῷ ἐπαγωγῆς.
\vs{3}Κολλήθητι αὐτῷ καὶ μὴ ἀποστῇς, ἵνα αὐξηθῇς ἐπʼ ἐσχάτων σου.
\vs{4}Πᾶν ὃ ἐὰν ἐπαχθῇ σοι, δέξαι, καὶ ἐν ἀλλάγμασι ταπεινώσεώς σου μακροθύμησον.
\vs{5}Ὅτι ἐν πυρὶ δοκιμάζεται χρυσὸς, καὶ ἄνθρωποι δεκτοὶ ἐν καμίνῳ ταπεινώσεως.
\vs{6}Πίστευσον αὐτῷ καὶ ἀντιλήψεταί σου, εὔθυνον τὰς ὁδούς σου καὶ ἔλπισον ἐπʼ αὐτόν.
\vs{7}Οἱ φοβούμενοι τὸν Κύριον, ἀναμείνατε τὸ ἔλεος αὐτοῦ, καὶ μὴ ἐκκλίνητε ἵνα μὴ πέσητε.
\vs{8}Οἱ φοβούμενοι Κύριον πιστεύσατε αὐτῷ, καὶ οὐ μὴ πταίσῃ ὁ μισθὸς ὑμῶν.
\vs{9}Οἱ φοβούμενοι Κύριον ἐλπίσατε εἰς ἀγαθὰ, καὶ εἰς εὐφροσύνην αἰῶνος καὶ ἐλέους.

\vs{10}Ἐμβλέψατε εἰς ἀρχαίας γενεὰς καὶ ἴδετε, τίς ἐνεπίστευσε Κυρίῳ καὶ κατῃσχύνθη; ἢ τίς ἐνέμεινε τῷ φόβῳ αὐτοῦ καὶ ἐγκατελείφθη; ἢ τίς ἐπεκαλέσατο αὐτὸν, καὶ ὑπερεῖδεν αὐτόν;
\vs{11}Διότι οἰκτίρμων καὶ ἐλεήμων ὁ Κύριος, καὶ ἀφίησιν ἁμαρτίας, καὶ σώζει ἐν καιρῷ θλίψεως.
\vs{12}Οὐαὶ καρδίαις δειλαῖς, καὶ χερσὶ παρειμέναις, καὶ ἁμαρτωλῷ ἐπιβαίνοντι ἐπὶ δύο τρίβους.

\vs{13}Οὐαὶ καρδίᾳ παρειμένῃ, ὅτι οὐ πιστεύει, διὰ τοῦτο οὐ σκεπασθήσεται.
\vs{14}Οὐαὶ ὑμῖν τοῖς ἀπολωλεκόσι τὴν ὑπομονὴν, καὶ τί ποιήσετε ὅταν ἐπισκέπτηται ὁ Κύριος;

\vs{15}Οἱ φοβούμενοι Κύριον οὐκ ἀπειθήσουσι ῥημάτων αὐτοῦ, καὶ οἱ ἀγαπῶντες αὐτὸν συντηπήσουσι τὰς ὁδοὺς αὐτοῦ.
\vs{16}Οἱ φοβούμενοι Κύριον ζητήσουσιν εὐδοκίαν αὐτοῦ, καὶ οἱ ἀγαπῶντες αὐτὸν ἐμπλησθήσονται τοῦ νόμου.
\vs{17}Οἱ φοβούμενοι Κύριον ἐτοιμάσουσι καρδίας αὐτῶν, καὶ ἐνώπιον αὐτοῦ παπεινώσουσι τὰς ψυχὰς αὐτῶν.
\vs{18}Ἐμπεσούμεθα εἰς χεῖρας Κυρίου, καὶ οὐκ εἰς χεῖρας ἀνθρώπων· ὡς γὰρ ἡ μεγαλωσύνη αὐτοῦ, οὕτως καὶ τὸ ἔλεος αὐτοῦ.

\ch{3}
Ἐμοῦ τοῦ πατρὸς ἀκούσατε τέκνα, καὶ οὕτως ποιήσατε, ἵνα σωθῆτε.
\vs{2}Ὁ γὰρ Κύριος ἐδόξασε πατέρα ἐπὶ τέκνοις, καὶ κρίσιν μητρὸς ἐστερέωσεν ἐφʼ υἱοῖς.
\vs{3}Ὁ τιμῶν πατέρα ἐξιλάσεται ἁμαρτίαις.
\vs{4}Καὶ ὡς ὁ ἀποθησαυρίζων, ὁ δοξάζων μητέρα αὐτοῦ.

\vs{5}Ὁ τιμῶν πατερα εὐφρανθήσεται ὑπὸ τέκνων, καὶ ἐν ἡμέρᾳ προσευχῆς αὐτοῦ εἰσακουσθήσεται.
\vs{6}Ὁ δοξάζων πατέρα μακροημερεύσει, καὶ ὁ εἰσακούων Κυρίου ἀναπαύσει μητέρα αὐτοῦ,
\vs{7}καὶ ὡς δεσπόταις δουλεύσει ἐν τοῖς γεννήσασιν αὐτόν.

\vs{8}Ἐν ἔργῳ καὶ λόγῳ τίμα τὸν πατέρα σου, ἵνα ἐπέλθῃ σοι εὐλογία παρʼ αὐτοῦ.
\vs{9}Εὐλογία γὰρ πατρὸς στηρίζει οἴκους τέκνων, κατάρα δὲ μητρὸς ἐκριζοῖ θεμέλια.
\vs{10}Μὴ δοξάζου ἐν ἀτιμίᾳ πατρός σου, οὐ γάρ ἐστί σοι δόξα πατρὸς ἀτιμία.
\vs{11}Ἡ γὰρ δόξα ἀνθρώπου ἐκ τιμῆς πατρὸς αὐτοῦ, καὶ ὄνειδος τέκνοις μήτηρ ἐν ἀδοξίᾳ.

\vs{12}Τέκνον, ἀντιλαβοῦ ἐν γήρᾳ πατρός σου, καὶ μὴ λυπήσῃς αὐτὸν ἐν τῇ ζωῇ αὐτοῦ.
\vs{13}Κᾂν ἀπολείπῃ σύνεσιν, συγγνώμην ἔχε, καὶ μὴ ἀτιμάσῃς αὐτὸν ἐν πάσῃ ἰσχύϊ σου.
\vs{14}Ἐλεημοσύνη γὰρ πατρὸς οὐκ ἐπιλησθήσεται, καὶ ἀντὶ ἁμαρτιῶν προσανοικοδομηθήσεταί σοι.
\vs{15}Ἐν ἡμέρᾳ θλίψεώς σου ἀναμνησθήσεταί σου· ὡς εὐδία ἐπὶ παγετῷ, οὕτως ἀναλυθήσονταί σου αἱ ἁμαρτίαι.

\vs{16}Ὡς βλάσφημος ὁ ἐγκαταλιπὼν πατέρα, καὶ κεκατηραμένος ὑπὸ Κυρίου ὁ παροργίζων μητέρα αὐτοῦ.

\vs{17}Τέκνον, ἐν πραΰτητι τὰ ἔργα σου διέξαγε, καὶ ὑπὸ ἀνθρώπου δεκτοῦ ἀγαπηθήσῃ.
\vs{18}Ὅσῳ μέγας εἶ, τοσούτῳ ταπεινοῦ σεαυτὸν, καὶ ἔναντι Κυρίου εὑρήσεις χάριν.
\vs{20}Ὅτι μεγάλη ἡ δυναστεία τοῦ Κυρίου, καὶ ὑπὸ τῶν ταπεινῶν δοξάζεται.

\vs{21}Χαλέπωτερά σου μὴ ζήτει, καὶ ἰσχυρότερά σου μὴ ἐξέταζε,
\vs{22}ἃ προσετάγη σοι, ταῦτα διανοοῦ· οὐ γάρ ἐστί σοι χρεία τῶν κρυπτῶν.
\vs{23}Ἐν τοῖς περισσοῖς τῶν ἔργων σου μὴ περιεργάζου· πλείονα γὰρ συνέσεως ἀνθρώπων ὑπεδείχθη σοι.
\vs{24}Πολλοὺς γὰρ ἐπλάνησεν ἡ ὑπόληψις αὐτῶν, καὶ ὑπόνοια πονηρὰ ὠλίσθησε διανοίας αὐτῶν.

\vs{26}Καὶ ὁ ἀγαπῶν κίνδυνον, ἐν αὐτῷ ἐμπεσεῖται· καρδία σκληρὰ κακωθήσεται ἐπʼ ἐσχάτων.
\vs{27}Καρδία σκληρὰ βαρυνθήσεται πόνοις, καὶ ὁ ἁμαρτωλὸς προσθήσει ἁμαρτίαν ἐφʼ ἁμαρτίαις.
\vs{28}Ἐπαγωγὴ ὑπερηφάνου οὐκ ἔστιν ἴασις, φυτὸν γὰρ πονηρίας ἐῤῥίζωκεν ἐν αὐτῷ·
\vs{29}καρδία συνετοῦ διανοηθήσεται παραβολὴν, καὶ οὖς ἀκροατοῦ ἐπιθυμία σοφοῦ.

\vs{30}Πῦρ φλογιζόμενον ἀποσβέσει ὕδωρ, καὶ ἐλεημοσύνη ἐξιλάσεται ἁμαρτίας.
\vs{31}Ὁ ἀνταποδιδοὺς χάριτας μέμνηται εἰς τὰ μετὰ ταῦτα, καὶ ἐν καιρῷ πτώσεως εὑρήσει στήριγμα.

\ch{4}
Τέκνον, τὴν ζωὴν τοῦ πτωχοῦ μὴ ἀποστερήσῃς, καὶ μὴ παρελκύσῃς ὀφθαλμοὺς ἐπιδεεῖς.
\vs{2}ψυχὴν πεινῶσαν μὴ λυπήσῃς, καὶ μὴ παροργίσῃς ἄνδρα ἐν ἀπορίᾳ αὐτοῦ.
\vs{3}Καρδίαν παροργισμένην μὴ προσταράξῃς, καὶ μὴ παρελκύσῃς δόσιν προσδεομένου.
\vs{4}Ἱκέτην θλιβόμενον μὴ ἀπαναίνου, καὶ μὴ ἀποστρέψῃς τὸ πρόσωπόν σου ἀπὸ πτωχοῦ.
\vs{5}Ἀπὸ δεομένου μὴ ἀποστρέψῃς ὀφθαλμόν, καὶ μὴ δῷς τόπον ἀνθρώπῳ καταράσασθαί σε.
\vs{6}Καταρωμένου γάρ σε ἐν πικρίᾳ ψυχῆς αὐτοῦ, τῆς δεήσεως αὐτοῦ ἐπακούσεται ὁ ποιήσας αὐτόν.
\vs{7}Προσφιλῆ συναγωγῇ σεαυτὸν ποίει, καὶ μεγιστᾶνι ταπεινοῦ τὴν κεφαλήν σου.
\vs{8}Κλῖνον πτωχῷ τὸ οὖς σου, καὶ ἀποκρίθητι αὐτῷ εἰρηνικὰ ἐν πρᾳΰτητι.

\vs{9}Ἐξελοῦ ἀδικούμενον ἐκ χειρὸς ἀδικοῦντος, καὶ μὴ ὀλιγοψυχήσῃς ἐν τῷ κρίνειν σε.
\vs{10}Γίνου ὀρφανοῖς ὡς πατὴρ, καὶ ἀντὶ ἀνδρὸς τῇ μητρὶ αὐτῶν· καὶ ἔσῃ ὡς υἱὸς ὑψίστου, καὶ ἀγαπήσει σε μᾶλλον ἡ μήτηρ σου.

\vs{11}Ἡ σοφία υἱοὺς ἑαυτῇ ἀνύψωσε, καὶ ἐπιλαμβάνεται τῶν ζητούντων αὐτήν·
\vs{12}ὁ ἀγαπῶν αὐτὴν ἀγαπᾷ ζωὴν, καὶ οἱ ὀρθρίζοντες πρὸς αὐτὴν ἐμπλησθήσονται εὐφροσύνης.

\vs{13}Ὁ κρατῶν αὐτῆς κληρονομήσει δόξαν, καὶ οὗ εἰσπορεύεται εὐλογήσει Κύριος.
\vs{14}Οἱ λατρεύοντες αὐτῇ λειτουργήσουσιν Ἁγίῳ, καὶ τοὺς ἀγαπῶντας αὐτὴν ἀγαπᾷ ὁ Κύριος.
\vs{15}Ὁ ὑπακούων αὐτῆς κρινεῖ ἔθνη, καὶ ὁ προσελθὼν αὐτῇ κατασκηνώσει πεποιθώς.
\vs{16}Ἐὰν ἐμπιστεύσῃς, κατακληρονομήσεις αὐτὴν, καὶ ἐν κατασχέσει ἔσονται αἱ γενεαὶ αὐτοῦ.
\vs{17}Ὅτι διεστραμμένως πορεύεται μετʼ αὐτοῦ ἐν πρώτοις· φόβον δὲ καὶ δειλίαν ἐπάξει ἐπʼ αὐτὸν, καὶ βασανίσει αὐτὸν ἐν παιδιᾷ αὐτῆς, ἕως οὗ ἐμπιστεύσῃ τῇ ψυχῇ αὐτοῦ, καὶ πειράσῃ αὐτὸν ἐν τοῖς δικαιώμασιν αὐτῆς.
\vs{18}Καὶ πάλιν ἐπανήξει κατʼ εὐθεῖαν πρὸς αὐτὸν, καὶ εὐφρανεῖ αὐτὸν, καὶ ἀποκαλύψει αὐτῷ τὰ κρυπτὰ αὐτῆς.
\vs{19}Ἐὰν ἀποπλανηθῇ, ἐγκαταλείψει αὐτὸν, καὶ παραδώσει αὐτὸν εἰς χεῖρας πτώσεως αὐτοῦ.

\vs{20}Συντήρησον καιρὸν καὶ φύλαξαι ἀπὸ πονηροῦ, καὶ περὶ τῆς ψυχῆς σου μὴ αἰσχυνθῇς.
\vs{21}Ἔστι γὰρ αἰσχύνη ἐπάγουσα ἁμαρτίαν, καὶ ἔστιν αἰσχύνη δόξα καὶ χάρις.
\vs{22}Μὴ λάβῃς πρόσωπον κατὰ τῆς ψυχῆς σου, καὶ μὴ ἐντραπῇς εἰς πτῶσίν σου.
\vs{23}Μὴ κωλύσῃς λόγον ἐν καιρῷ σωτηρίας,
\vs{24}ἐν γὰρ λόγῳ γνωσθήσεται σοφία, καὶ παιδεία ἐν ῥήματι γλώσσης.
\vs{25}Μὴ ἀντίλεγε τῇ ἀληθείᾳ, καὶ περὶ τῆς ἀπαιδευσίας σου ἐντράπηθι.
\vs{26}Μὴ αἰσχυνθῇς ὁμολογῆσαι ἐφʼ ἁμαρτίαις σου, καὶ μὴ βιάζου ῥοῦν ποταμοῦ.
\vs{27}Καὶ μὴ ὑποστρώσῃς σεαυτὸν ἀνθρώπῳ μωρῷ, καὶ μὴ λάβῃς πρόσωπον δυνάστου.
\vs{28}Ἕως τοῦ θανάτου ἀγώνισαι περὶ τῆς ἀληθείας, καὶ Κύριος ὁ Θεὸς πολεμήσει ὑπὲρ σοῦ.

\vs{29}Μὴ γίνου τραχὺς ἐν γλώσσῃ σου, καὶ νωθρὸς καὶ παρειμένος ἐν τοῖς ἔργοις σου.
\vs{30}Μὴ ἴσθι ὡς λέων ἐν τῷ οἴκῳ σου, καὶ φαντασιοκοπῶν ἐν τοῖς οἰκέταις σου.
\vs{31}Μὴ ἔστω ἡ χείρ σου ἐκτεταμένη εἰς τὸ λαβεῖν, καὶ ἐν τῷ ἀποδιδόναι συνεσταλμένη.

\ch{5}
Μὴ ἔπεχε ἐπὶ τοῖς χρήμασί σου, καὶ μὴ εἴπῃς, αὐτάρκη μοι ἐστί.
\vs{2}Μὴ ἐξακολούθει τῇ ψυχῇ σου καὶ τῇ ἰσχύϊ σου, τοῦ πορεύεσθαι ἐν ἐπιθυμίαις καρδίας σου.
\vs{3}Καὶ μὴ εἴπῃς, τίς με δυναστεύσει; ὁ γὰρ Κύριος ἐκδικῶν ἐκδικήσει σε.
\vs{4}Μὴ εἶπῃς, ἥμαρτον, καὶ τί μοι ἐγένετο; ὁ γὰρ Κυριός ἐστι μακρόθυμος.
\vs{5}Περὶ ἐξιλασμοῦ μὴ ἄφοβος γίνου προσθεῖναι ἁμαρτίαν ἐφʼ ἁμαρτίαις.
\vs{6}Καὶ μὴ εἴπῃς, ὁ οἰκτιρμὸς αὐτοῦ πολὺς, τὸ πλῆθος τῶν ἁμαρτιῶν μου ἐξιλάσεται· ἔλεος γὰρ καὶ ὀργὴ παρʼ αὐτοῦ, καὶ ἐπὶ ἁμαρτωλοὺς καταπαύσει ὁ θυμὸς αὐτοῦ.

\vs{7}Μὴ ἀνάμενε ἐπιστρέψαι πρὸς Κύριον, καὶ μὴ ὑπερβάλλου ἡμέραν ἐξ ἡμέρας· ἐξάπινα γὰρ ἐξελεύσεται ὀργὴ Κυρίου, καὶ ἐν καιρῷ ἐκδικήσεως ἐξολῇ.
\vs{8}Μὴ ἔπεχε ἐπὶ χρήμασιν ἀδίκοις, οὐδὲν γὰρ ὠφελήσεις ἐν ἡμέρᾳ ἐπαγωγῆς.

\vs{9}Μὴ λίκμα ἐν παντὶ ἀνέμῳ, καὶ μὴ πορεύου ἐν πάσῃ ἀτραπῷ· οὕτως ὁ ἁμαρτωλὸς ὁ δίγλωσσος.
\vs{10}Ἴσθι ἐστηριγμένος ἐν συνέσει σου, καὶ εἷς ἔστω σου ὁ λόγος.
\vs{11}Γίνου ταχὺς ἐν ἀκροάσει σου, καὶ ἐν μακροθυμίᾳ φθέγγου ἀπόκρισιν.
\vs{12}Εἰ ἔστι σοι σύνεσις, ἀποκρίθητι τῷ πλησίον· εἰ δὲ μὴ, ἡ χείρ σου ἔστω ἐπὶ στόματί σου.
\vs{13}Δόξα καὶ ἀτιμία ἐν λαλιᾷ, καὶ γλῶσσα ἀνθρώπου πτῶσις αὐτῷ.
\vs{14}Μὴ κληθῇς ψίθυρος, καὶ τῇ γλώσσῃ σου μὴ ἐνέδρευε· ἐπὶ γὰρ τῷ κλέπτῃ ἐστὶν αἰσχύνη, καὶ κατάγνωσις πονηρὰ ἐπὶ διγλώσσου.
\vs{15}Ἐν μεγάλῳ καὶ ἐν μικρῷ μὴ ἀγνόει.

\ch{6}
Καὶ ἀντὶ φίλου μὴ γίνου ἐχθρός· ὄνομα γὰρ πονηρὸν αἰσχύνην καὶ ὄνειδος κληρονομήσει· οὕτως ὁ ἁμαρτωλὸς ὁ δίγλωσσος.
\vs{2b}Μὴ ἐπάρῃς σεαυτὸν ἐν βουλῇ ψυχῆς σου, ἵνα μὴ διαρπαγῇ ὡς ταῦρος ἡ ψυχή σου.
\vs{3}Τὰ φύλλα σου καταφάγεσαι, καὶ τοὺς καρπούς σου ἀπολέσεις, καὶ ἀφήσεις σεαυτὸν ὡς ξύλον ξηρόν.

\vs{4}Ψυχὴ πονηρὰ ἀπολεῖ τὸν κτησάμενον αὐτήν, καὶ ἐπίχαρμα ἐχθρῶν ποιήσει αὐτόν.
\vs{5}Λάρυγξ γλυκὺς πληθυνεῖ φίλους αὐτοῦ, καὶ γλῶσσα εὔλαλος πληθυνεῖ εὐπροσήγορα.
\vs{6}Οἱ εἰρηνεύοντές σοι ἔστωσαν πολλοὶ, οἱ δὲ σύμβουλοί σου εἷς ἀπὸ χιλίων.

\vs{7}Εἰ κτᾶσαι φίλον, ἐν πειρασμῷ κτῆσαι αὐτόν, καὶ μὴ ταχὺ ἐμπιστεύσῃς αὐτῷ.
\vs{8}Ἔστι γὰρ φίλος ἐν καιρῷ αὐτοῦ, καὶ οὐ μὴ παραμείνῃ ἐν ἡμέρᾳ θλίψεώς σου.
\vs{9}Καὶ ἔστι φίλος μετατιθέμενος εἰς ἔχθραν, καὶ μάχην ὀνειδισμοῦ σου ἀποκαλύψει.
\vs{10}Καὶ ἔστι φίλος κοινωνὸς τραπεζῶν, καὶ οὐ μὴ παραμείνῃ ἐν ἡμέρᾳ θλίψεώς σου.
\vs{11}Καὶ ἐν τοῖς ἀγαθοῖς σου ἔσται ὡς σὺ, καὶ ἐπὶ τοὺς οἰκέτας σου παῤῥησιάσεται.
\vs{12}Ἐὰν ταπεινωθῇς, ἔσται κατὰ σοῦ, καὶ ἀπὸ τοῦ προσώπου σου κρυβήσεται.

\vs{13}Ἀπὸ τῶν ἐχθρῶν σου διαχωρίσθητι, καὶ ἀπὸ τῶν φίλων σου πρόσεχε.
\vs{14}Φίλος πιστὸς σκέπη κραταιὰ, ὁ δὲ εὑρὼν αὐτὸν εὗρε θησαυρόν.
\vs{15}Φίλου πιστοῦ οὐκ ἔστιν ἀντάλλαγμα, καὶ οὐκ ἔστι σταθμὸς τῆς καλλονῆς αὐτοῦ.
\vs{16}Φίλος πιστὸς φάρμακον ζωῆς, καὶ οἱ φοβούμενοι Κύριον εὑρήσουσιν αὐτόν.
\vs{17}Ὁ φοβούμενος Κύριον εὐθύνει φιλίαν αὐτοῦ, ὅτι κατʼ αὐτὸν οὕτως καὶ ὁ πλησίον αὐτοῦ.

\vs{18}Τέκνον, ἐκ νεότητός σου ἐπίλεξαι παιδείαν, καὶ ἕως πολιῶν εὑρήσεις σοφίαν.
\vs{19}Ὡς ὁ ἀροτριῶν καὶ ὁ σπείρων πρόσελθε αὐτῇ, καὶ ἀνάμενε τοὺς ἀγαθοὺς καρποὺς αὐτῆς· ἐν γὰρ τῇ ἐργασίᾳ αὐτῆς ὀλίγον κοπιάσεις, καὶ ταχὺ φάγεσαι γεννημάτων αὐτῆς.
\vs{20}Ὡς τραχεῖά ἐστι σφόδρα τοῖς ἀπαιδεύτοις, καὶ οὐκ ἐμμενεῖ ἐν αὐτῇ ἀκάρδιος.
\vs{21}Ὡς λίθος δοκιμασίας ἰσχυρὸς ἔσται ἐπʼ αὐτῷ, καὶ οὐ χρονιεῖ ἀποῤῥίψαι αὐτήν.
\vs{22}Σοφία γὰρ κατὰ τὸ ὄνομα αὐτῆς ἐστι, καὶ οὐ πολλοῖς ἐστι φανερά.

\vs{23}Ἄκουσον, τέκνον, καὶ δέξαι γνώμην μου, καὶ μὴ ἀπαναίνου τὴν συμβουλίαν μου.
\vs{24}Καὶ εἰσένεγκον τοὺς πόδας σου εἰς τὰς πέδας αὐτῆς, καὶ εἰς τὸν κλοίον αὐτῆς τὸν τράχηλόν σου.
\vs{25}Ὑπόθες τὸν ὦμόν σοῦ, καὶ βάσταξον αὐτὴν, καὶ μὴ προσοχθίσῃς τοῖς δεσμοῖς αὐτῆς.
\vs{26}Ἐν πάσῃ ψυχῇ σου πρόσελθε αὐτῇ, καὶ ἐν ὅλῃ δυνάμει σου συντήρησον τὰς ὁδοὺς αὐτῆς.
\vs{27}Ἐξίχνεύσον καὶ ζήτησον, καὶ γνωσθήσεταί σοι, καὶ ἐγκρατὴς γενόμενος μὴ ἀφῇς αὐτήν.
\vs{28}Ἐπʼ ἐσχάτων γὰρ εὑρήσεις τὴν ἀνάπαυσιν αὐτῆς, καὶ στραφήσεταί σοι εἰς εὐφροσύνην.
\vs{29}Καὶ ἔσονταί σοι αἱ πέδαι εἰς σκέπην ἰσχύος, καὶ οἱ κλοιοὶ αὐτῆς εἰς στολὴν δόξης.
\vs{30}Κόσμος γὰρ χρύσεός ἐστιν ἐπʼ αὐτῆς, καὶ οἱ δεσμοὶ αὐτῆς κλῶσμα ὑακίνθινον.
\vs{31}Στολὴν δόξης ἐνδύσῃ αὐτὴν, καὶ στέφανον ἀγαλλιάματος περιθήσεις σεαυτῷ.

\vs{32}Ἐὰν θέλῃς, τέκνον, παιδευθήσῃ, καὶ ἐὰν δῷς τὴν ψυχήν σου, πανοῦργος ἔσῃ.
\vs{33}Ἐὰν ἀγαπήσῃς ἀκούειν ἐκδέξῃ, καὶ ἐὰν κλίνῃς τὸ οὖς σου σοφὸς ἔσῃ.
\vs{34}Ἐν πλήθει πρεσβυτέρων στῆθι, καὶ, τίς σοφός; αὐτῷ προσκολλήθητι.
\vs{35}Πᾶσαν διήγησιν θείαν θέλε ἀκούειν, καὶ παροιμίαι συνέσεως μὴ ἐκφευγέτωσάν σε.
\vs{36}Ἐὰν ἴδῃς συνετὸν, ὄρθριζε πρὸς αὐτὸν, καὶ βαθμοὺς θυρῶν αὐτοῦ ἐκτριβέτω ὁ πούς σου.
\vs{37}Διανοοῦ ἐν τοῖς προστάγμασιν Κυρίου, καὶ ἐν ταῖς ἐντολαῖς αὐτοῦ μελέτα διαπαντός· αὐτὸς στηριεῖ τὴν καρδίαν σου, καὶ ἡ ἐπιθυμία τῆς σοφίας σου δοθήσεταί σοι.

\ch{7}
Μὴ ποίει κακὰ, καὶ οὐ μὴ σε καταλάβῃ κακόν.
\vs{2}Ἀπόστηθι ἀπὸ ἀδίκου, καὶ ἐκκλινεῖ ἀπὸ σοῦ.

\vs{3}Υἱὲ μὴ σπεῖρε ἐπʼ αὔλακας ἀδικίας, καὶ οὐ μὴ θερίσῃς αὐτὰς ἑπταπλασίως.
\vs{4}Μὴ ζήτει παρὰ Κυρίου ἡγεμονίαν, μηδὲ παρὰ βασιλέως καθέδραν δόξης.
\vs{5}Μὴ δικαιοῦ ἔναντι Κυρίου, καὶ παρὰ βασιλεῖ μὴ σοφίζου·
\vs{6}μὴ ζήτει γενέσθαι κριτὴς, μὴ οὐκ ἐξισχύσεις ἐξάραι ἀδικίας. μήποτε εὐλαβηθῇς ἀπὸ προσώπου δυνάστου, καὶ θήσεις σκάνδαλον ἐν εὐθύτητί σου.
\vs{7}Μὴ ἁμάρτανε εἰς πλῆθος πόλεως, καὶ μὴ καταβάλῃς σεαυτὸν ἐν ὄχλῳ.
\vs{8}Μὴ καταδεσμεύσῃς δὶς ἁμαρτίαν, ἐν γὰρ τῇ μιᾷ οὐκ ἀθῶος ἔσῃ.
\vs{9}Μὴ εἴπῃς, τῷ πλήθει τῶν δώρων μου ἐπόψεται, καὶ ἐν τῷ προσενέγκαι με Θεῷ ὑψίστῳ προσδέξεται.
\vs{10}Μὴ ὀλιγοψυχήσῃς ἐν τῇ προσευχῇ σου, καὶ ἐλεημοσύνην ποιῆσαι μὴ παρίδῃς.

\vs{11}Μὴ καταγέλα ἄνθρωπον ὄντα ἐν πικρίᾳ ψυχῆς αὐτοῦ, ἔστι γὰρ ὁ ταπεινῶν καὶ ἀνυψῶν.
\vs{12}Μὴ ἀροτρία ψεῦδος ἐπʼ ἀδελφῷ σου, μηδὲ φίλῳ τὸ ὅμοιον ποίει.
\vs{13}Μὴ θέλε ψεύδεσθαι πᾶν ψεῦδος, ὁ γὰρ ἐνδελεχισμὸς αὐτοῦ οὐκ εἰς ἀγαθόν.
\vs{14}Μὴ ἀδολέσχει ἐν πλήθει πρεσβυτέρων, καὶ μὴ δευτερώσῃς λόγον ἐν προσευχῇ σου.
\vs{15}Μὴ μισήσῃς ἐπίπονον ἐργασίαν, καὶ γεωργίαν ὑπὸ ὑψίστου ἐκτισμένην.
\vs{16}Μὴ προσλογίζου σεαυτὸν ἐν πλήθει ἁμαρτωλῶν·
\vs{17}Ταπείνωσον σφόδρα τὴν ψυχήν σου.
\vs{17a}Μνήσθητι ὅτι ὀργὴ οὐ χρονιεῖ,
\vs{17b}ὅτι ἐκδίκησις ἀσεβοῦς πῦρ καὶ σκώληξ.
\vs{18}Μὴ ἀλλάξῃς φίλον ἕνεκεν ἀδιαφόρου, μήδʼ ἀδελφὸν γνήσιον ἐν χρυσίῳ Σουφείρ.
\vs{19}Μὴ ἀστόχει γυναικὸς σοφῆς καὶ ἀγαθῆς, καὶ γὰρ χάρις αὐτῆς ὑπὲρ τὸ χρυσίον.
\vs{20}Μὴ κακώσῃς οἰκέτην ἐργαζόμενον ἐν ἀληθείᾳ, μηδὲ μίσθιον διδόντα ψυχὴν αὐτοῦ.

\vs{21}Οἰκέτην συνετὸν ἀγαπάτω σου ἡ ψυχὴ, μὴ στερήσῃς αὐτὸν ἐλευθερίας.
\vs{22}Κτήνη σοί ἐστιν; ἐπισκέπτου αὐτά· καὶ εἰ ἔστι σοι χρήσιμα, ἐμμενέτω σοι.
\vs{23}Τέκνα σοί ἐστι, παίδευσον αὐτὰ, καὶ κάμψον ἐκ νεότητος τὸν τράχηλον αὐτῶν.
\vs{24}Θυγατέρες σοί εἰσι; πρόσεχε τῷ σώματι αὐτῶν, καὶ μὴ ἱλαρώσῃς πρὸς αὐτὰς τὸ πρόσωπόν σου.
\vs{25}Ἔκδου θυγατέρα, καὶ ἔσῃ τετελεκὼς ἔργον μέγα, καὶ ἀνδρὶ συνετῷ δώρησαι αὐτήν.
\vs{26}Γυνή σοι ἐστὶ κατὰ ψυχήν; μὴ ἐκβάλῃς αὐτήν.

\vs{27}Ἐν ὅλῃ καρδίᾳ δόξασον τὸν πατέρα σου, καὶ μητρὸς ὠδῖνας μὴ ἐπιλάθῃ.
\vs{28}Μνήσθητι ὅτι διʼ αὐτῶν ἐγενήθης, καὶ τί ἀνταποδώσεις αὐτοῖς καθὼς αὐτοὶ σοί;
\vs{29}Ἐν ὅλῃ ψυχῇ σου εὐλαβοῦ τὸν Κύριον, καὶ τοὺς ἱερεῖς αὐτοῦ θαύμαζε.
\vs{30}Ἐν ὅλῃ δυνάμει ἀγάπησον τὸν ποιήσαντά σε, καὶ τοὺς λειτουργοὺς αὐτοῦ μὴ ἐγκαταλίπῃς.
\vs{31}Φοβοῦ τὸν Κύριον, καὶ δόξασον ἱερέα, καὶ δὸς τὴν μερίδα αὐτῷ, καθὼς ἐντέταλταί σοι, ἀπαρχὴν, καὶ περὶ πλημμελείας, καὶ δόσιν βραχιόνων, καὶ θυσίαν ἁγιασμοῦ, καὶ ἀπαρχὴν ἁγίων.
\vs{32}Καὶ πτωχῷ ἔκτεινον τὴν χεῖρά σου, ἵνα τελειωθῇ ἡ εὐλογία σου.
\vs{33}Χάρις δόματος ἔναντι παντὸς ζῶντος, καὶ ἐπὶ νεκρῷ μὴ ἀποκωλύσῃς χάριν.
\vs{34}Μὴ ὑστέρει ἀπὸ κλαιόντων, καὶ μετὰ πενθούντων πένθησον.
\vs{35}Μὴ ὄκνει ἐπισκέπτεσθαι ἄῤῥωστον, ἐκ γὰρ τῶν τοιούτων ἀγαπηθήσῃ.
\vs{36}Ἐν πᾶσι τοῖς λόγοις σου μιμνήσκου τὰ ἔσχατά σου, καὶ εἰς τὸν αἰῶνα οὐχ ἁμαρτήσεις.

\ch{8}
Μὴ διαμάχου μετὰ ἀνθρώπου δυνάστου, μήποτε ἐμπέσῃς εἰς τὰς χεῖρας αὐτοῦ.
\vs{2}Μὴ ἔριζε μετὰ ἀνθρώπου πλουσίου, μήποτε ἀντιστήσῃ σου τὴν ὁλκήν· πολλους γὰρ ἀπώλεσε τὸ χρυσίον, καὶ καρδίας βασιλέων ἐξέκλινε.
\vs{3}Μὴ διαμάχου μετὰ ἀνθρώπου γλωσσώδους, καὶ μὴ ἐπιστοιβάσῃς ἐπὶ τὸ πῦρ αὐτοῦ ξύλα.
\vs{4}Μὴ πρόσπαιζε ἀπαιδεύτῳ, ἵνα μὴ ἀτιμάζωνται οἱ πρόγονοί σου.
\vs{5}Μὴ ὀνείδιζε ἄνθρωπον ἀποστρέφοντα ἀπὸ ἁμαρτίας, μνήσθητι ὅτι πάντες ἐσμὲν ἐν ἐπιτιμίοις.
\vs{6}Μὴ ἀτιμάσῃς ἄνθρωπον ἐν γήρει αὐτοῦ, καὶ γὰρ ἐξ ἡμῶν γηράσκουσι.
\vs{7}Μὴ ἐπίχαιρε ἐπὶ νεκρῷ, μνήσθητι ὅτι πάντες τελευτῶμεν.
\vs{8}Μὴ παρίδῃς διήγημα σοφῶν, καὶ ἐν ταῖς παροιμίαις αὐτῶν ἀναστρέφου, ὅτι παρʼ αὐτῶν μαθήσῃ παιδείαν, καὶ λειτουργῆσαι μεγιστᾶσι.

\vs{9}Μὴ ἀστόχει διηγήματος γερόντων, καὶ γὰρ αὐτοὶ ἔμαθον παρὰ τῶν πατέρων αὐτῶν· ὅτι παρʼ αὐτῶν μαθήσῃ σύνεσιν, καὶ ἐν καιρῷ χρείας δοῦναι ἀπόκρισιν.
\vs{10}Μὴ ἔκκαιε ἄνθρακας ἁμαρτωλοῦ, μὴ ἐμπυρισθῇς ἐν πυρὶ φλογὸς αὐτοῦ.
\vs{11}Μὴ ἐξαναστῇς ἀπὸ προσώπου ὑβριστοῦ, ἵνα μὴ ἐγκαθίσῃ ὡς ἔνεδρον τῷ στόματί σου.
\vs{12}Μὴ δανείσῃς ἀνθρώπῳ ἰσχυροτέρῳ σου, καὶ ἐὰν δανείσῃς, ὡς ἀπολωλεκὼς γίνου·
\vs{13}Μὴ ἐγγυήσῃ ὑπὲρ δυναμίν σου, καὶ ἐὰν ἐγγυήσῃ, ὡς ἀποτίσων φρόντιζε.
\vs{14}Μὴ δικάζου μετὰ κριτοῦ, κατὰ γὰρ τὴν δόξαν αὐτοῦ κρινοῦσιν αὐτῷ.
\vs{15}Μετὰ τολμηροῦ μὴ πορεύου ἐν ὁδῷ, ἵνα μὴ βαρύνηται κατὰ σοῦ· αὐτὸς γὰρ κατὰ τὸ θέλημα αὐτοῦ ποιήσει, καὶ τῇ ἀφροσύνῃ αὐτοῦ συναπολῇ.

\vs{16}Μετὰ θυμώδους μὴ ποιήσῃς μάχην, καὶ μὴ διαπορεύου μετʼ αὐτοῦ τὴν ἔρημον, ὅτι ὡς οὐδὲν ἐν ὀφθαλμοῖς αὐτοῦ αἷμα, καὶ ὅπου οὐκ ἔστι βοήθεια, καταβαλεῖ σε.
\vs{17}Μετὰ μωροῦ μὴ συμβουλεύου, οὐ γὰρ δυνήσεται λόγον στέξαι.
\vs{18}Ἐνώπιον ἀλλοτρίου μὴ ποιήσῃς κρυπτὸν, οὐ γὰρ γινώσκεις τί τέξεται.
\vs{19}Παντὶ ἀνθρώπῳ μὴ ἔκφαινε σὴν καρδίαν, καὶ μὴ ἀναφερέτω σοι χάριν.

\ch{9}
Μὴ ζήλου γυναῖκα τοῦ κόλπου σου, μηδὲ διδάξῃς ἐπὶ σεαυτὸν παιδείαν πονηράν.
\vs{2}Μὴ δῷς γυναικὶ τὴν ψυχήν σου, ἐπιβῆναι αὐτὴν ἐπὶ τὴν ἰσχύν σου.
\vs{3}Μὴ ὑπάντα γυναικὶ ἑταιριζομένῃ, μήποτε ἐμπέσῃς εἰς τὰς παγίδας αὐτῆς.
\vs{4}Μετὰ ψαλλούσης μὴ ἐνδελέχιζε, μήποτε ἁλῷς ἐν τοῖς ἐπιχειρήμασιν αὐτῆς.
\vs{5}Παρθένον μὴ καταμάνθανε, μήποτε σκανδαλισθῇς ἐν τοῖς ἐπιτιμίοις αὐτῆς.
\vs{6}Μὴ δῷς πόρναις τὴν ψυχήν σου, ἵνα μὴ ἀπολέσῃς τὴν κληρονομίαν σου.
\vs{7}Μὴ περιβλέπου ἐν ῥύμαις πόλεως, καὶ ἐν ταῖς ἐρήμοις αὐτῆς μὴ πλανῶ.
\vs{8}Ἀπόστρεψον ὀφθαλμὸν ἀπὸ γυναικὸς εὐμόρφου, καὶ μὴ καταμάνθανε κάλλος ἀλλότριον· ἐν κάλλει γυναικὸς πολλοὶ ἐπλανήθησαν, καὶ ἐκ τούτου φιλία ὡς πῦρ ἀνακαίεται.

\vs{9}Μετὰ ὑπάνδρου γυναικὸς μὴ κάθου τὸ σύνολον, καὶ μὴ συμβολοκοπήσῃς μετʼ αὐτῆς ἐν οἴνῳ, μήποτε ἐκκλίνῃ ἡ ψυχή σου ἐπʼ αὐτὴν, καὶ τῷ πνεύματί σου ὀλισθήσῃς εἰς ἀπώλειαν.
\vs{10}Μὴ ἐγκαταλίπῃς φίλον ἀρχαῖον, ὁ γὰρ πρόσφατος οὐκ ἔστιν ἔφισος αὐτῷ· οἶνος νέος, φίλος νέος, ἐὰν παλαιωθῇ, μετʼ εὐφροσύνης πίεσαι αὐτόν.
\vs{11}Μὴ ζηλώσῃς δόξαν ἁμαρτωλοῦ, οὐ γὰρ οἶδας τί ἔσται ἡ καταστροφὴ αὐτοῦ.
\vs{12}Μὴ εὐδοκήσῃς ἐν εὐδοκίᾳ ἀσεβῶν, μνήσθητι ὅτι ἕως ᾅδου οὐ μὴ δικαιωθῶσι.
\vs{13}Μακρὰν ἄπεχε ἀπὸ ἀνθρώπου ὃς ἔχει ἐξουσιαν τοῦ φονεύειν, καὶ οὐ μὴ ὑποπτεύσῃς φόβον θανάτου· κᾂν προσέλθῃς, μὴ πλημμελήσῃς, ἵνα μὴ ἀφέληται τὴν ζωήν σου· ἐπίγνωθι ὅτι ἐν μέσῳ παγίδων διαβαίνεις, καὶ ἐπὶ ἐπάλξεων πόλεων περιπατεῖς.

\vs{14}Κατὰ τὴν ἰσχύν σου στόχασαι τοὺς πλησίον, καὶ μετὰ σοφῶν συμβουλεύου.
\vs{15}Καὶ μετὰ συνετῶν ἔστω ὁ διαλογισμός σου, καὶ πᾶσα διήγησίς σου ἐν νόμῳ ὑψίστου.
\vs{16}Ἄνδρες δίκαιοι ἔστωσαν σύνδειπνοί σου, καὶ ἐν φόβῳ Κυρίου ἔστω τὸ καύχημά σου.
\vs{17}Ἐν χειρὶ τεχνιτῶν ἔργον ἐπαινεθήσεται, καὶ ὁ ἡγούμενος λαοῦ σοφὸς ἐν λόγῳ αὐτοῦ.
\vs{18}Φοβερὸς ἐν πόλει αὐτοῦ ἀνὴρ γλωσσώδης, καὶ ὁ προπετὴς ἐν λόγῳ αὐτοῦ μισηθήσεται.

\ch{10}
Κριτὴς σοφὸς παιδεύσει τὸν λαὸν αὐτοῦ, καὶ ἡγεμονία συνετοῦ τεταγμένη ἔσται.
\vs{2}Κατὰ τὸν κριτὴν τοῦ λαοῦ αὐτοῦ οὕτως καὶ οἱ λειτουργοὶ αὐτοῦ, καὶ κατὰ τὸν ἡγούμενον τῆς πόλεως πάντες οἱ κατοικοῦντες αὐτήν.
\vs{3}Βασιλεὺς ἀπαίδευτος ἀπολεῖ τὸν λαὸν αὐτοῦ, καὶ πόλις οἰκισθήσεται ἐν συνέσει δυναστῶν.
\vs{4}Ἐν χειρὶ Κυρίου ἐξουσία τῆς γῆς, καὶ τὸν χρήσιμον ἐγερεῖ εἰς καιρὸν ἐπʼ αὐτῆς.
\vs{5}Ἐν χειρὶ Κυρίου εὐοδία ἀνδρὸς, καὶ προσώπῳ γραμματέως ἐπιθήσει δόξαν αὐτοῦ.

\vs{6}Ἐπὶ παντὶ ἀδικήματι μὴ μηνιάσῃς τῷ πλησίον, καὶ μὴ πράσσε μηδὲν ἐν ἔργοις ὕβρεως.
\vs{7}Μισητὴ ἔναντι Κυρίου καὶ ἀνθρώπων ὑπερηφανία, καὶ ἐξ ἀμφοτέρων πλημμελήσει ἄδικα.
\vs{8}Βασιλεία ἀπὸ ἔθνους εἰς ἔθνος μετάγεται, διὰ ἀδικίας καὶ ὕβρεις καὶ χρήματα.

\vs{9}Τί ὑπερηφανεύεται γῆ καὶ σποδός; ὅτι ἐν ζωῇ ἔῤῥιψα τὰ ἐνδόσθια αὐτοῦ.
\vs{10}Μακρὸν ἀῤῥώστημα σκώπτει ἰατρὸς, καὶ βασιλεὺς σήμερον, καὶ αὔριον τελευτήσει.
\vs{11}Ἐν γὰρ τῷ ἀποθανεῖν ἄνθρωπον, κληρονομήσει ἑρπετὰ καὶ θηρία καὶ σκώληκας.
\vs{12}Ἀρχὴ ὑπερηφανίας, ἀνθρώπου ἀφισταμένου ἀπὸ Κυρίου, καὶ ἀπὸ τοῦ ποιήσαντος αὐτὸν ἀπέστη ἡ καρδία αὐτοῦ.
\vs{13}Ὅτι ἀρχὴ ὑπερηφανίας ἁμαρτία, καὶ ὁ κρατῶν αὐτῆς ἐξομβρήσει βδέλυγμα· διὰ τοῦτο παρεδόξασε Κύριος τὰς ἐπαγωγὰς, καὶ κατέστρεψεν εἰς τέλος αὐτούς.

\vs{14}Θρόνους ἀρχόντων καθεῖλεν ὁ Κύριος, καὶ ἐκάθισε πρᾳεῖς ἀντʼ αὐτῶν.
\vs{15}Ῥίζας ἐθνῶν ἐξέτιλεν ὁ Κύριος, καὶ ἐφύτευσε ταπεινοὺς ἀντʼ αὐτῶν.
\vs{16}χώρας ἐθνῶν κατέστρεψεν ὁ Κύριος, καὶ ἀπώλεσεν αὐτὰς ἕως θεμελίων γῆς.
\vs{17}Ἐξήρανεν ἐξ αὐτῶν καὶ ἀπώλεσεν αὐτοὺς, καὶ κατέπαυσεν ἀπὸ γῆς τὸ μνημόσυνον αὐτῶν.
\vs{18}Οὐκ ἔκτισται ἀνθρώποις ὑπερηφανία, οὐδὲ ὀργὴ θυμοῦ γεννήμασιν γυναικῶν.
\vs{19}Σπέρμα ἔντιμον ποίον; σπέρμα ἀνθρώπου· σπέρμα ἔντιμον ποῖον; οἱ φοβούμενοι τὸν Κύριον· σπέρμα ἄτιμον ποῖον; σπέρμα ἀνθρώπου· σπέρμα ἄτιμον ποῖον; οἱ παραβαίνοντες ἐντολάς.
\vs{20}Ἐν μέσῳ ἀδελφῶν ὁ ἡγούμενος αὐτῶν ἔντιμος, καὶ οἱ φοβούμενοι Κύριον ἐν ὀφθαλμοῖς αὐτοῦ.
\vs{22}Πλούσιος καὶ ἔνδοξος καὶ πτωχὸς, τὸ καύχημα αὐτῶν φόβος Κυρίου.

\vs{23}Οὐ δίκαιον ἀτιμάσαι πτωχὸν συνετὸν, καὶ οὐ καθήκει δοξάσαι ἄνδρα ἁμαρτωλόν.
\vs{24}Μεγιστὰν καὶ κριτὴς καὶ δυνάστης δοξασθήσεται, καὶ οὐκ ἔστιν αὐτῶν τις μείζων τοῦ φοβουμένου τὸν Κύριον.
\vs{25}Οἰκέτῃ σοφῷ ἐλεύθεροι λειτουργήσουσι, καὶ ἀνὴρ ἐπιστήμων οὐ γογγύσει.
\vs{26}Μὴ σοφίζου ποιῆσαι τὸ ἔργον σου, καὶ μὴ δοξάζου ἐν καιρῷ στενοχωρίας σου.
\vs{27}Κρείσσων ἐργαζόμενος ἐν πᾶσιν, ἢ περιπατῶν, ἢ δοξαζόμενος καὶ ἀπορῶν ἄρτων.

\vs{28}Τέκνον, ἐν πρᾳΰτητι δόξασον τὴν ψυχήν σου, καὶ δὸς αὐτῇ τιμὴν κατὰ τὴν ἀξίαν αὐτῆς.
\vs{29}Τὸν ἁμαρτάνοντα εἰς τὴν ψυχὴν αὐτοῦ τίς δικαιώσει; καὶ τίς δοξάσει τὸν ἀτιμάζοντα τὴν ζωὴν αὐτοῦ;
\vs{30}Πτωχὸς δοξάζεται διʼ ἐπιστήμην αὐτοῦ, καὶ πλούσιος δοξάζεται διὰ τὸν πλοῦτον αὐτοῦ.
\vs{31}Ὁ δὲ δοξαζόμενος ἐν πτωχείᾳ, καὶ ἐν πλούτῳ ποσαχῶς; καὶ ὁ ἄδοξος ἐν πλούτῳ, καὶ ἐν πτωχείᾳ ποσαχῶς;

\ch{11}
Σοφία ταπεινοῦ ἀνύψωσε κεφαλὴν, καὶ ἐν μέσῳ μεγιστάνων καθίσει αὐτόν.
\vs{2}Μὴ αἰνέσεις ἄνδρα ἐν κάλλει αὐτοῦ, καὶ μὴ βδελύξῃ ἄνθρωπον ἐν ὁράσει αὐτοῦ.
\vs{3}Μικρὰ ἐν πετεινοῖς μέλισσα, καὶ ἀρχὴ γλυκασμάτων ὁ καρπὸς αὐτῆς.
\vs{4}Ἐν περιβολῇ ἱματίων μὴ καυχήσῃ, καὶ ἐν ἡμέρᾳ δόξης μὴ ἐπαίρου, ὅτι θαυμαστὰ τὰ ἔργα Κυρίου, καὶ κρυπτὰ τὰ ἔργα αὐτοῦ ἐν ἀνθρώποις.
\vs{5}Πολλοὶ τύραννοι ἐκάθισαν ἐπὶ ἐδάφους, ὁ δὲ ἀνυπονόητος ἐφόρεσε διάδημα.
\vs{6}Πολλοὶ δυνάσται ἠτιμάσθησαν σφόδρα, καὶ ἔνδοξοι παρεδόθησαν εἰς χεῖρας ἑτέρων.
\vs{7}Πρὶν ἐξετάσῃς μὴ μέμψῃ· νόησον πρῶτον καὶ τότε ἐπιτιμα.
\vs{8}Πρὶν ἢ ἀκοῦσαι μὴ ἀποκρίνου, καὶ ἐν μέσῳ λόγων μὴ παρεμβάλλου.
\vs{9}Περὶ πράγματος οὗ οὐκ ἔστι σοι χρεία, μὴ ἔριζε, καὶ ἐν κρίσει ἁμαρτωλῶν μὴ συνέδρευε.

\vs{10}Τέκνον, μὴ περὶ πολλὰ ἔστωσαν αἱ πράξεις σου· ἐὰν πληθυνῇς, οὐκ ἀθωωθήσῃ· καὶ ἐὰν διώκῃς, οὐ μὴ καταλάβῃς, καὶ οὐ μὴ ἐκφύγῃς διαδράς.
\vs{11}Ἔστι κοπιῶν καὶ πονῶν καὶ σπεύδων, καὶ τόσῳ μᾶλλον ὑστερεῖται.
\vs{12}Ἔστι νωθρὸς καὶ προσδεόμενος ἀντιλήψεως, ὑστερῶν ἰσχύϊ, καὶ πτωχείᾳ περισσεύει, καὶ οἱ ὀφθαλμοὶ Κυρίου ἐπέβλεψαν αὐτῷ εἰς ἀγαθὰ, καὶ ἀνώρθωσεν αὐτὸν ἐκ ταπεινώσεως αὐτοῦ·
\vs{13}Καὶ ἀνύψωσε κεφαλὴν αὐτοῦ, καὶ ἀπεθαύμασαν ἐπʼ αὐτῷ πολλοί.

\vs{14}Ἀγαθὰ καὶ κακά, ζωὴ καὶ θάνατος, πτωχεία καὶ πλοῦτος παρὰ Κυρίου ἐστί.
\vs{17}Δόσις Κυρίου παραμένει εὐσεβέσι, καὶ ἡ εὐδοκία αὐτοῦ εἰς τὸν αἰῶνα εὐοδωθήσεται.
\vs{18}Ἔστι πλουτῶν ἀπὸ προσοχῆς καὶ σφιγγίας αὐτοῦ, καὶ αὕτη ἡ μερὶς τοῦ μισθοῦ αὐτοῦ.
\vs{19}Ἐν τῷ εἰπεῖν αὐτὸν, εὗρον ἀνάπαυσιν, καὶ νῦν φάγωμαι ἐκ τῶν ἀγαθῶν μου, καὶ οὐκ οἶδε τίς καιρὸς παρελεύσεται, καὶ καταλείψει αὐτὰ ἑτέροις, καὶ ἀποθανεῖται.
\vs{20}Στῆθι ἐν διαθήκῃ σου καὶ ὁμίλει ἐν αὐτῇ, καὶ ἐν τῷ ἔργῳ σου παλαιώθητι.
\vs{21}Μὴ θαύμαζε ἐν ἔργοις ἁμαρτωλοῦ, πίστευε τῷ Κυρίῳ καὶ ἔμμενε τῷ πόνῳ σου· ὅτι κοῦφον ἐν ὀφθαλμοῖς Κυρίου διὰ τάχους ἐξάπινα πλουτίσαι πένητα.

\vs{22}Εὐλογία Κυρίου ἐν μισθῷ εὐσεβοῦς, καὶ ἐν ὥρᾳ ταχινῇ ἀναθάλλει εὐλογίαν αὐτοῦ.
\vs{23}Μὴ εἴπῃς, τίς ἐστί μου χρεία; καὶ τίνα ἀπὸ τοῦ νῦν ἔσται μου τὰ ἀγαθά;
\vs{24}Μὴ εἴπῃς, αὐτάρκη μοι ἐστὶ, καὶ τί ἀπὸ τοῦ νῦν κακωθήσομαι;
\vs{25}Ἐν ἡμέρᾳ ἀγαθῶν ἀμνησία κακῶν, καὶ ἐν ἡμέρᾳ κακῶν οὐ μνησθήσεται ἀγαθῶν·
\vs{26}ὅτι κοῦφον ἔναντι Κυρίου ἐν ἡμέρᾳ τελευτῆς ἀποδοῦναι ἀνθρώπῳ κατὰ τὰς ὁδοὺς αὐτοῦ.
\vs{27}Κάκωσις ὥρας ἐπιλησμονὴν ποιεῖ τρυφῆς, καὶ ἐν συντελείᾳ ἀνθρώπου ἀποκάλυψις ἔργων αὐτοῦ.
\vs{28}Πρὸ τελευτῆς μὴ μακάριζε μηδένα, καὶ ἐν τέκνοις αὐτῶυ γνωσθήσεται ἀνήρ.

\vs{29}Μὴ πάντα ἄνθρωπον εἴσαγε εἰς τὸν οἶκόν σου, πολλὰ γὰρ τὰ ἔνεδρα τοῦ δολίου.
\vs{30}Πέρδιξ θηρευτὴς ἐν καρτάλλῳ, οὕτως καρδία ὑπερηφάνου, καὶ ὡς ὁ κατάσκοπος ἐπιβλέπει πτῶσιν.
\vs{31}Τὰ γὰρ ἀγαθὰ εἰς κακὰ μεταστρέφων ἐνεδρεύει, καὶ ἐν τοῖς αἱρετοῖς ἐπιθήσει μῶμον.
\vs{32}Ἀπὸ σπινθῆρος πυρὸς πληθύνεται ἀνθρακία, καὶ ἄνθρωπος ἁμαρτωλὸς εἰς αἷμα ἐνεδρεύει.
\vs{33}Πρόσεχε ἀπὸ κακούργου, πονηρὰ γὰρ τεκταίνει, μήποτε μῶμον εἰς τὸν αἰῶνα δῷ σοι.
\vs{34}Ἐνοίκισον ἀλλότριον, καὶ διαστρέψει σε ἐν ταραχαῖς, καὶ ἀπαλλοτριώσει σε τῶν ἰδίων σου.

\ch{12}
Ἐὰν εὖ ποιῇς, γνῶθι τίνι ποιεῖς, καὶ ἔσται χάρις τοῖς ἀγαθοῖς σου.
\vs{2}Εὐποίησον εὐσεβεῖ, καὶ εὑρήσεις ἀνταπόδομα, καὶ εἰ μὴ παρʼ αὐτοῦ, ἀλλὰ παρὰ ὑψίστου.
\vs{3}Οὐκ ἔστιν ἀγαθὰ τῷ ἐνελεχίζοντι εἰς κακὰ, καὶ τῷ ἐλεημοσύνην μὴ χαριζομένῳ.
\vs{4}Δὸς τῷ εὐσεβεῖ, καὶ μὴ ἀντιλάβῃ τοῦ ἁμαρτωλοῦ.
\vs{5}Εὐποίησον τῷ ταπεινῷ, καὶ μὴ δῷς ἀσεβεῖ· ἐμπόδισον τοὺς ἄρτους αὐτοῦ, καὶ μὴ δῷς αὐτῷ ἵνα μὴ ἐν αὐτοῖς σε δυναστεύσῃ· διπλάσια γὰρ κακὰ εὑρήσεις ἐν πᾶσιν ἀγαθοῖς οἷς ἂν ποιήσῃς αὐτῷ.
\vs{6}Ὅτι καὶ ὁ ὕψιστος ἐμίσησεν ἁμαρτωλούς, καὶ τοῖς ἀσεβέσιν ἀποδώσει ἐκδίκησιν.

\vs{7}Δὸς τῷ ἀγαθῷ, καὶ μὴ ἀντιλάβῃ τοῦ ἁμαρτωλοῦ.
\vs{8}Οὐκ ἐκδικηθησεται ἐν ἀγαθοῖς ὁ φίλος, καὶ οὐ κρυβήσεται ἐν κακοῖς ὁ ἐχθρός.
\vs{9}Ἐν ἀγαθοῖς ἀνδρὸς οἱ ἐχθροὶ αὐτοῦ ἐν λύπῃ, καὶ ἐν τοῖς κακοῖς αὐτοῦ καὶ ὁ φίλος διαχωρισθήσεται.
\vs{10}Μὴ πιστεύσῃς τῷ ἐχθρῷ σου εἰς τὸν αἰῶνα· ὡς γὰρ ὁ χαλκὸς ἰοῦται, οὕτως ἡ πονηρία αὐτοῦ.
\vs{11}Καὶ ἐὰν ταπεινωθῇ καὶ πορεύηται συγκεκυφὼς, ἐπίστησον τὴν ψυχήν σου καὶ φύλαξαι ἀπʼ αὐτοῦ, καὶ ἔσῃ αὐτῷ ὡς ἐκμεμαχὼς ἔσοπτρον, καὶ γνώσῃ ὅτι οὐκ εἰς τέλος κατίωσε.
\vs{12}Μὴ στήσῃς αὐτὸν παρὰ σεαυτὸν, μὴ ἀνατρέψας σε στῇ ἐπὶ τὸν τόπον σου· μὴ καθίσης αὐτὸν ἐκ δεξιῶν σου, μήποτε ζητήσῃ τὴν καθέδραν σου, καὶ ἐπʼ ἐσχάτῳ ἐπιγνώσῃ τοὺς λόγους μου, καὶ ἐπὶ τῶν ῥημάτων μου κατανυγήσῃ.

\vs{13}Τίς ἐλεήσει ἐπαοιδὸν ὀφιόδηκτον, καὶ πάντας τοὺς προσάγοντας θηρίοις;
\vs{14}Οὕτως τὸν προσπορεύομενον ἀνδρὶ ἁμαρτωλῷ καὶ συμφυρόμενον ἐν ταῖς ἁμαρτίαις αὐτοῦ.
\vs{15}Ὥραν μετὰ σοῦ διαμενεῖ, καὶ ἐὰν ἐκκλίνῃς, οὐ μὴ καρτερήσῃ.
\vs{16}Καὶ ἐν τοῖς χείλεσιν αὐτοῦ γλυκανεῖ ὁ ἐχθρὸς, καὶ ἐν τῇ καρδία αὐτοῦ βουλεύσεται ἀνατρέψαι σε εἰς βόθρον· ἐν ὀφθαλμοῖς αὐτοῦ δακρύσει ὁ ἐχθρὸς, καὶ ἐὰν εὕρῃ καιρὸν, οὐκ ἐμπλησθήσεται ἀφʼ αἵματος.
\vs{17}Κακὰ ἂν ὑπαντήσῃ σοι, εὑρήσεις αὐτὸν ἐκεῖ πρότερόν σου, καὶ ὡς βοηθῶν ὑποσχάσει πτέρναν σου.
\vs{18}Κινήσει τὴν κεφαλὴναὐτοῦ, καὶ ἐπικροτήσει ταῖς χερσὶν αὐτοῦ, καὶ πολλὰ διαψιθυρίσει, καὶ ἀλλοιώσει τὸ πρόσωπον αὐτοῦ.

\ch{13}
Ὁ ἁπτομενος πίσσης μολυνθήσεται, καὶ ὁ κοινωνῶν ὑπερηφάνῳ ὁμοιωθήσεται αὐτῷ.
\vs{2}Βάρος ὑπὲρ σὲ μὴ ἄρῃς, καὶ ἰσχυροτέρῳ σου καὶ πλουσιωτέρῳ μὴ κοινώνει· τί κοινωνήσει χύτρα πρὸς λέβητα; αὕτη προσκρούσει, καὶ αὕτη συντριβήσεται.

\vs{3}Πλούσιος ἠδίκησε, καὶ αὐτὸς προσενεβριμήσατο· πτωχὸς ἠδίκηται, καὶ αὐτὸς προσδεηθήσεται.
\vs{4}Ἐὰν χρησιμεύσῃς, ἐργᾶται ἐν σοί· καὶ ἐὰν ὑστερήσῃς, καταλείψει σε.
\vs{5}Ἐὰν ἔχῃς, συμβιώσεταί σοι, καὶ ἀποκενώσει σε, καὶ αὐτὸς οὐ πονέσει.
\vs{6}Χρείαν ἔσχηκέ σου, καὶ ἀποπλανήσει σε, καὶ προσγελάσεταί σοι, καὶ δώσει σοι ἐλπίδα· λαλήσει σοι καλὰ, καὶ ἐρεῖ, τίς ἡ χρεία σου;
\vs{7}Καὶ αἰσχυνεῖ σε ἐν τοῖς βρώμασιν αὐτοῦ, ἕως οὗ ἀποκενώσῃ σε δὶς ἢ τρὶς, καὶ ἐπʼ ἐσχάτῳ καταμωκήσεταί σου· μετὰ ταῦτα ὄψεταί σε, καὶ καταλείψει σε, καὶ τὴν κεφαλὴν αὐτοῦ κινήσει ἐπὶ σοί.
\vs{8}Πρόσεχε μὴ ἀποπλανηθῇς, καὶ μὴ ταπεινωθῇς ἐν εὐφροσύνῃ σου.

\vs{9}Προσκαλεσαμένου σε δυνάστου, ὑποχωρῶν γίνου, καὶ τόσῳ μᾶλλον προσκαλέσεταί σε.
\vs{10}Μὴ ἔμπιπτε ἵνα μὴ ἀπωσθῇς, καὶ μὴ μακρὰν ἀφιστῶ ἵνα μὴ ἐπιλησθῇς.
\vs{11}Μὴ ἔπεχε εἰσηγορεῖσθαι μετʼ αὐτοῦ, καὶ μὴ πίστευε τοῖς πλείοσι λόγοις αὐτοῦ· ἐκ πολλῆς γὰρ λαλιᾶς πειράσει σε, καὶ ὡς προσγελῶν ἐξετάσει.
\vs{12}Ἀνελεήμων ὁ μὴ συντηρῶν λόγους, καὶ οὐ μὴ φείσηται περὶ κακώσεως καὶ δεσμῶν.
\vs{13}Συντήρησον καὶ πρόσεχε σφοδρῶς, ὅτι μετὰ τῆς πτώσεώς σου περιπατεῖς.

\vs{15}Πᾶν ζῶον ἀγαπᾷ τὸ ὅμοιον αὐτῷ, καὶ πᾶς ἄνθρωπος τὸν πλησίον αὐτοῦ.
\vs{16}Πᾶσα σὰρξ κατὰ γένος συνάγεται, καὶ τῷ ὁμοίῳ αὐτοῦ προσκολληθήσεται ἀνήρ.
\vs{17}Τί κοινωνήσει λύκος ἀμνῷ; οὕτως ἁμαρτωλὸς πρὸς εὐσεβῆ.
\vs{18}Τίς εἰρήνη ὑαίνῃ πρὸς κύνα; καὶ τίς εἰρήνη πλουσίῳ πρὸς πένητα;
\vs{19}Κυνήγια λεόντων ὄναγροι ἐν ἐρήμῳ, οὕτως νομαὶ πλουσίων πτωχοί.
\vs{20}Βδέλυγμα ὑπερηφάνῳ ταπεινότης, οὕτως βδέλυγμα πλουσίῳ πτωχός.

\vs{21}Πλούσιος σαλευόμενος στηρίζεται ὑπὸ φίλων, ταπεινὸς δὲ πεσὼν προσαπωθεῖται ὑπὸ φίλων.
\vs{22}Πλουσίου σφαλέντος πολλοὶ ἀντιλήμπτορες, ἐλάλησεν ἀπόῤῥητα καὶ ἐδικαίωσαν αὐτόν· ταπεινὸς ἔσφαλε καὶ προσεπετίμησαν αὐτῷ, ἐφθέγξατο σύνεσιν καὶ οὐκ ἐδόθη αὐτῷ τόπος.
\vs{23}Πλούσιος ἐλάλησε καὶ πάντες ἐσίγησαν, καὶ τὸν λόγον αὐτοῦ ἀνύψωσαν ἕως τῶν νεφελῶν· πτωχὸς ἐλάλησε, καὶ εἶπαν, τίς οὗτος; κἂν προσκόψῃ, προσανατρέψουσιν αὐτόν.
\vs{24}Ἀγαθὸς ὁ πλοῦτος ᾧ μὴ ἔστιν ἁμαρτία, καὶ πονηρὰ ἡ πτωχεία ἐν στόμασιν ἀσεβοῦς.

\vs{25}Καρδία ἀνθρώπου ἀλλοιοῖ τὸ πρόσωπον αὐτοῦ, ἐὰν εἰς ἀγαθὰ ἐάν τε εἰς κακά.
\vs{26}Ἴχνος καρδίας ἐν ἀγαθοῖς πρόσωπον ἱλαρὸν, καὶ εὕρεσις παραβολῶν διαλογισμοὶ μετὰ κόπου.

\ch{14}
Μακάριος ἀνὴρ ὃς οὐκ ὠλίσθησεν ἐν στόματι αὐτοῦ, καὶ οὐ κατενύγη ἐν λύπῃ ἁμαρτίας.
\vs{2}Μακάριος οὗ οὐ κατέγνω ἡ ψυχὴ αὐτοῦ, καὶ ὃς οὐκ ἔπεσεν ἀπὸ τῆς ἐλπίδος αὐτοῦ.

\vs{3}Ἀνδρὶ μικρολόγῳ οὐ καλὸς ὁ πλοῦτος, καὶ ἀνθρώπῳ βασκάνῳ ἱνατί χρήματα;
\vs{4}Ὁ συνάγων ἀπὸ τῆς ψυχῆς αὐτοῦ, συνάγει ἄλλοις, καὶ ἐν τοῖς ἀγαθοῖς αὐτοῦ τρυφήσουσιν ἕτεροι.
\vs{5}Ὁ πονηρὸς ἑαυτῷ, τίνι ἀγαθὸς ἔσται; καὶ οὐ μὴ εὐφρανθήσεται ἐν τοῖς χρήμασιν αὐτοῦ.
\vs{6}Τοῦ βασκαίνοντος ἑαυτὸν οὐκ ἔστι πονηρότερος, καὶ τοῦτο ἀνταπόδομα τῆς κακίας αὐτοῦ·
\vs{7}κᾂν εὐποιῇ, ἐν λήθῃ ποιεῖ, καὶ ἐπʼ ἐσχάτων ἐκφαίνει τὴν κακίαν αὐτοῦ.
\vs{8}Πονηρὸς ὁ βασκαίνων ὀφθαλμῷ, ἀποστρέφων πρόσωπον καὶ ὁ ὑπερορῶν ψυχάς.
\vs{9}Πλεονέκτου ὀφθαλμὸς οὐκ ἐμπίπλαται μερίδι, καὶ ἀδικία πονηρὰ ἀναξηρανει ψυχήν.
\vs{10}Ὀφθαλμὸς πονηρὸς φθονερὸς ἐπʼ ἄρτῳ, καὶ ἐλλιπὴς ἐπὶ τῆς τραπέζης αὐτοῦ.

\vs{11}Τέκνον, καθὼς ἐὰν ἔχεις εὐποίει σεαυτόν, καὶ προσφορὰς Κυρίῳ ἀξίως πρόσαγε.
\vs{12}Μνήσθητι ὅτι θάνατος οὐ χρονιεῖ, καὶ διαθήκη ᾅδου οὐχ ὑπεδείχθη σοι.
\vs{13}Πρίν σε τελευτῆσαι εὐποίει φίλῳ, καὶ κατὰ τὴν ἰσχύν σου ἔκτεινον καὶ δὸς αὐτῷ.
\vs{14}Μὴ ἀφυστερήσῃς ἀπὸ ἀγαθῆς ἡμέρας, καὶ μερὶς ἐπιθυμίας ἀγαθῆς μή σε παρελθάτω.
\vs{15}Οὐχὶ ἑτέρῳ καταλείψεις τοὺς πόνους σου, καὶ τοὺς κόπους σου εἰς διαίρεσιν κλήρου;
\vs{16}Δὸς καὶ λάβε, καὶ ἀπάτησον τὴν ψυχήν σου, ὅτι οὐκ ἔστιν ἐν ᾅδου ζητῆσαι τρυφήν.
\vs{17}Πᾶσα σὰρξ ὡς ἱμάτιον παλαιοῦται, ἡ γὰρ διαθήκη ἀπʼ αἰῶνος θανάτῳ ἀποθανῇ.
\vs{18}Ὡς φύλλον θάλλον ἐπὶ δένδρου δασέος, τὰ μὲν καταβάλλει, ἀλλα δὲ φύει· οὕτως γενεὰ σαρκὸς καὶ αἵματος, ἡ μὲν τελευτᾷ, ἑτέρα δὲ γεννᾶται.
\vs{19}Πᾶν ἔργον σηπόμενον ἐκλείπει, καὶ ὁ ἐργαζόμενος αὐτὸ μετʼ αὐτοῦ ἀπελεύσεται.

\vs{20}Μακάριος ἀνὴρ ὃς ἐν σοφίᾳ τελευτήσει, καὶ ὃς ἐν συνέσει αὐτοῦ διαλεχθήσεται·
\vs{21}ὡ διανοούμενος τὰς ὁδοὺς αὐτῆς ἐν καρδίᾳ αὐτοῦ, καὶ ἐν τοῖς ἀποκρύφοις αὐτῆς νοηθήσεται.
\vs{22}Ἔξελθε ὀπίσω αὐτῆς ὡς ἰχνευτὴς, καὶ ἐν ταῖς εἰσόδοις αὐτῆς ἐνέδρευε.
\vs{23}Ὁ παρακύπτων διὰ τῶν θυρίδων αὐτῆς, καὶ ἐπὶ τῶν θυρωμάτων αὐτῆς ἀκροάσεται·
\vs{24}ὁ καταλύων σύνεγγυς τοῦ οἴκου αὐτῆς, καὶ πήξει πάσσαλον ἐν τοῖς τοίχοις αὐτῆς.
\vs{25}Στήσει τὴν σκηνὴν αὐτοῦ κατὰ χεῖρας αὐτῆς, καὶ καταλύσει ἐν καταλύματι ἀγαθῶν.
\vs{26}Θήσει τὰ τέκνα αὐτοῦ ἐν τῇ σκέπῃ αὐτῆς, καὶ ὑπὸ τοὺς κλάδους αὐτῆς αὐλισθήσεται.
\vs{27}Σκεπασθήσεται ὑπʼ αὐτῆς ἀπὸ καύματος, καὶ ἐν τῇ δόξῃ αὐτῆς καταλύσει.

\ch{15}
Ὁ φοβούμενος Κύριον ποιήσει αὐτὸ, καὶ ὁ ἐγκρατὴς τοῦ νόμου καταλήψεται αὐτήν·
\vs{2}καὶ ὑπαντήσεται αὐτῷ ὡς μήτηρ, καὶ ὡς γυνὴ παρθενίας προσδέξεται αὐτόν·
\vs{3}ψωμιεῖ αὐτὸν ἄρτον συνέσεως, καὶ ὕδωρ σοφίας ποτίσει αὐτόν.
\vs{4}Στηριχθήσεται ἐπʼ αὐτὴν καὶ οὐ μὴ κλιθῇ, καὶ ἐπʼ αὐτῆς ἐφέξει καὶ οὐ μὴ καταισχυνθῇ.
\vs{5}Καὶ ὑψώσει αὐτὸν παρὰ τοὺς πλησίον αὐτοῦ, καὶ ἐν μέσῳ ἐκκλησίας ἀνοίξει στόμα αὐτοῦ.
\vs{6}Εὐφροσύνην καὶ στέφανον ἀγαλλιάματος καὶ ὄνομα αἰώνιον κατακληρονομήσει.
\vs{7}Οὐ μὴ καταλήψονται αὐτὴν ἄνθρωποι ἀσύνετοι, καὶ ἄνδρες ἁμαρτωλοὶ οὐ μὴ ἴδωσιν αὐτήν.
\vs{8}Μακράν ἐστιν ὑπερηφανίας, καὶ ἄνδρες ψεῦσται οὐ μὴ μνησθήσονται αὐτῆς.

\vs{9}Οὐχ ὡραῖος αἶνος ἐν στόματι ἁμαρτωλοῦ, ὅτι οὐ παρὰ Κυρίου ἀπεστάλη.
\vs{10}Ἐν γὰρ σοφίᾳ ῥηθήσεται αἶνος, καὶ ὁ Κύριος εὐοδώσει αὐτόν.
\vs{11}Μὴ εἴπῃς, ὅτι διὰ Κύριον ἀπέστην· ἃ γὰρ ἐμίσησεν, οὐ ποιήσεις.
\vs{12}Μὴ εἴπῃς, ὅτι αὐτός με ἐπλάνησεν· οὐ γὰρ χρείαν ἔχει ἀνδρὸς ἁμαρτωλοῦ.

\vs{13}Πᾶν βδέλυγμα ἐμίσησε Κύριος, καὶ οὐκ ἔστιν ἀγαπητὸν τοῖς φοβουμένοις αὐτόν.
\vs{14}Αὐτὸς ἐξ ἀρχῆς ἐποίησεν ἄνθρωπον, καὶ ἀφῆκεν αὐτὸν ἐν χειρὶ διαβουλίου αὐτοῦ.
\vs{15}Ἐὰν θέλῃς, συντηρήσεις ἐντολὰς, καὶ πίστιν ποιῆσαι εὐδοκίας.
\vs{16}Παρέθηκέ σοι πῦρ καὶ ὕδωρ, οὗ ἐὰν θέλῃς ἐκτενεῖς τὴν χεῖρά σου.
\vs{17}Ἔναντι ἀνθρώπων ἡ ζωὴ καὶ ὁ θάνατος, καὶ ὃ ἐὰν εὐδοκήσῃ δοθήσεται αὐτῷ.
\vs{18}Ὅτι πολλὴ σοφία τοῦ Κυρίου, ἰσχυρὸς ἐν δυναστείᾳ καὶ βλέπων τὰ πάντα.
\vs{19}Καὶ οἱ ὀφθαλμοὶ αὐτοῦ ἐπὶ τοὺς φοβουμένους αὐτὸν, καὶ αὐτὸς ἐπιγνώσεται πᾶν ἔργον ἀνθρώπου.
\vs{20}Καὶ οὐκ ἐνετείλατο οὐδενὶ ἀσεβεῖν, καὶ οὐκ ἔδωκεν ἄνεσιν οὐδενὶ ἁμαρτάνειν.

\ch{16}
Μὴ ἐπιθύμει τέκνων πλῆθος ἀχρήστων, μὴ εὐφραίνου ἐπὶ υἱοῖς ἀσεβέσιν.
\vs{2}Ἐὰν πληθύνωσι, μὴ εὐφραίνου ἐπʼ αὐτοῖς, εἰ μή ἐστι φόβος Κυρίου μετʼ αὐτῶν.
\vs{3}Μὴ ἐμπιστεύσῃς τῇ ζωῇ αὐτῶν, καὶ μὴ ἔπεχε ἐπὶ τὸν τόπον αὐτῶν· κρείσσων γὰρ εἷς ἢ χίλιοι, καὶ ἀποθανεῖν ἄτεκνον ἢ ἔχειν τέκνα ἀσεβῆ·
\vs{4}ἀπὸ γὰρ ἑνὸς συνετοῦ συνοικισθήσεται πόλις, φυλὴ δὲ ἀνόμων ἐρημωθήσεται.
\vs{5}Πολλὰ τοιαῦτα ἑώρακα ἐν ὀφθαλμοῖς μου, καὶ ἰσχυρότερα τούτων ἀκήκοε τὸ οὖς μου.

\vs{6}Ἐν συναγωγῇ ἁμαρτωλῶν ἐκκαυθήσεται πῦρ, καὶ ἐν ἔθνει ἀπειθεῖ ἐξεκαύθη ὀργή.
\vs{7}Οὐκ ἐξιλάσατο περὶ τῶν ἀρχαίων γιγάντων, οἳ ἀπέστησαν τῇ ἰσχύϊ αὐτῶν.
\vs{8}Οὐκ ἐφείσατο περὶ τῆς παροικίας Λὼτ, οὓς ἐβδελύξατο διὰ τὴν ὑπερηφανίαν αὐτῶν.
\vs{9}Οὐκ ἠλέησεν ἔθνος ἀπωλείας, τοὺς ἐξῃρμένους ἐν ἁμαρτίαις αὐτῶν·
\vs{10}καὶ οὕτως ἑξακοσίας χιλιάδας πεζῶν τοὺς ἐπισυναχθέντας ἐν σκληροκαρδίᾳ αὐτῶν.
\vs{11}Κᾂν ᾖ εἷς σκληροτράχηλος, θαυμαστὸν, τοῦτο εἰ ἀθωωθήσεται, ἔλεος γὰρ καὶ ὀργὴ παρʼ αὐτοῦ· δυνάστης ἐξιλασμῶν· καὶ ἐκχέων ὀργήν.

\vs{12}Κατὰ τὸ πολὺ ἔλεος αὐτοῦ, οὕτως καὶ πολὺς ὁ ἔλεγχος αὐτοῦ· ἄνδρα κατὰ τὰ ἔργα αὐτοῦ κρίνει.
\vs{13}Οὐκ ἐκφεύξεται ἐν ἁρπάγμασιν ἁμαρτωλὸς, καὶ οὐ μὴ καθυστερήσει ὑπομονὴν εὐσεβοῦς.
\vs{14}Πάσῃ ἐλεημοσύνῃ ποιήσει τόπον, ἕκαστος κατὰ τὰ ἔργα αὐτοῦ εὑρήσει.

\vs{17}Μὴ εἴπῃς, ὅτι ἀπὸ Κυρίου κρυβήσομαι, μὴ ἐξ ὕψους τίς μου μνησθήσεται; ἐν λαῷ πλείονι οὐ μὴ μνησθῶ, τίς γὰρ ἡ ψυχή μου ἐν ἀμετρήτῳ κτίσει;
\vs{18}Ἰδοὺ ὁ οὐρανὸς καὶ ὁ οὐρανὸς τοῦ οὐρανοῦ τοῦ Θεοῦ, ἄβυσσος καὶ γῆ σαλευθήσονται ἐν τῇ ἐπισκοπῇ αὐτοῦ·
\vs{19}ἅμα τὰ ὄρη καὶ τὰ θεμέλια τῆς γῆς, ἐν τῷ ἐπιβλέψαι εἰς αὐτὰ, τρόμῳ συσσείονται,
\vs{20}καὶ ἐπʼ αὐτοῖς οὐ διανοηθήσεται καρδία· καὶ τὰς ὁδοὺς αὐτοῦ τίς ἐνθυμηθήσεται;
\vs{21}Καὶ καταιγὶς, ἣν οὐκ ὄψεται ἄνθρωπος, τὰ δὲ πλείονα τῶν ἔργων αὐτοῦ ἐν ἀποκρύφοις.
\vs{22}Ἔργα δικαιοσύνης τίς ἀναγγελεῖ, ἢ τίς ὑπομενεῖ; μακρὰν γὰρ ἡ διαθήκη.
\vs{23}Ἐλαττούμενος καρδίᾳ διανοεῖται ταῦτα, καὶ ἀνὴρ ἄφρων καὶ πλανώμενος διανοεῖται μωρά.

\vs{24}Ἄκουσόν μου, τέκνον, καὶ μάθε ἐπιστήμην, καὶ ἐπὶ τῶν λόγων μου πρόσεχε τῇ καρδίᾳ σου.
\vs{25}Ἐκφαίνω ἐν σταθμῷ παιδείαν, καὶ ἐν ἀκριβείᾳ ἀπαγγέλλω ἐπιστήμην.
\vs{26}Ἐν κρίσει Κυρίου τὰ ἔργα αὐτοῦ ἀπʼ ἀρχῆς, καὶ ἀπὸ ποιήσεως αὐτῶν διέστειλε μερίδας αὐτῶν.
\vs{27}Ἐκόσμησεν εἰς αἰῶνα τὰ ἔργα αὐτοῦ, καὶ τὰς ἀρχὰς αὐτῶν εἰς γενεὰς αὐτῶν· οὔτε ἐπείνασαν, οὔτε ἐκοπίασαν, καὶ οὐκ ἐξέλιπον ἀπὸ τῶν ἔργων αὐτων.
\vs{28}Ἕκαστος τὸν πλησίον αὐτοῦ οὐκ ἔθλιψε, καὶ ἕως αἰῶνος οὐκ ἀπειθήσουσι τοῦ ῥήματος αὐτοῦ·
\vs{29}καὶ μετὰ ταῦτα Κύριος εἰς τὴν γῆν ἐπέβλεψε, καὶ ἐνέπλησεν αὐτὴν τῶν ἀγαθῶν αὐτοῦ.
\vs{30}Ψυχὴν παντὸς ζώου ἐκάλυψε τὸ πρόσωπον αὐτῆς, καὶ εἰς αὐτὴν ἡ ἀποστροφὴ αὐτῶν.

\ch{17}
Κύριος ἔκτισεν ἐκ γῆς ἄνθρωπον, καὶ πάλιν ἀπέστρεψεν αὐτὸν εἰς αὐτήν.
\vs{2}Ἡμέρας ἀριθμοῦ καὶ καιρὸν ἔδωκεν αὐτοῖς, καὶ ἔδωκεν αὐτοῖς ἐξουσίαν τῶν ἐπʼ αὐτῆς.
\vs{3}Καθʼ ἑαυτοὺς ἐνέδυσεν αὐτοὺς ἰσχὺν, καὶ κατʼ εἰκόνα αὐτοῦ ἐποίησεν αὐτούς.
\vs{4}Καὶ ἔθηκε τὸν φόβον αὐτοῦ ἐπὶ πάσης σαρκός, καὶ κατακυριεύειν θηρίων καὶ πετεινῶν.
\vs{6}Διαβούλιον καὶ γλῶσσαν καὶ ὀφθαλμοὺς, ὦτα καὶ καρδίαν ἔδωκε διανοεῖσθαι αὐτοῖς.
\vs{7}Ἐπιστήμην συνέσεως ἐνέπλησεν αὐτοὺς, καὶ ἀγαθὰ καὶ κακὰ ὑπέδειξεν αὐτοῖς.
\vs{8}Ἔθηκε τὸν ὀφθαλμὸν αὐτοῦ ἐπὶ τὰς καρδίας αὐτῶν, δεῖξαι αὐτοῖς τὸ μεγαλεῖον τῶν ἔργων αὐτοῦ·
\vs{10}καὶ ὄνομα ἁγιασμοῦ αἰνέσουσιν, ἵνα διηγῶνται τὰ μεγαλεῖα τῶν ἔργων αὐτοῦ.

\vs{11}Προσέθηκεν αὐτοῖς ἐπιστήμην, καὶ νόμον ζωῆς ἐκληροδότησεν αὐτοῖς.
\vs{12}Διαθήκην αἰῶνος ἔστησε μετʼ αὐτῶν, καὶ τὰ κρίματα αὐτοῦ ὑπέδειξεν αὐτοῖς.
\vs{13}Μεγαλεῖον δόξης εἶδον οἱ ὀφθαλμοὶ αὐτῶν, καὶ δόξαν φωνῆς αὐτῶν ἤκουσε τὸ οὖς αὐτῶν.
\vs{14}Καὶ εἶπεν αὐτοῖς, προσέχετε ἀπὸ παντὸς ἀδίκου· καὶ ἐνετείλατο αὐτοῖς ἑκάστῳ περὶ τοῦ πλησίον.
\vs{15}Αἱ ὁδοὶ αὐτῶν ἐναντίον αὐτοῦ διαπαντὸς, οὐ κρυβήσονται ἀπὸ τῶν ὀφθαλμῶν αὐτοῦ.
\vs{17}Ἑκάστῳ ἔθνει κατέστησεν ἡγούμενον, καὶ μερὶς Κυρίου Ἰσραήλ ἐστιν.
\vs{19}Ἅπαντα τὰ ἔργα αὐτῶν ὡς ὁ ἥλιος ἐναντίον αὐτοῦ, καὶ οἱ ὀφθαλμοὶ αὐτοῦ ἐνδελεχεῖς ἐπὶ τὰς ὁδοὺς αὐτῶν.
\vs{20}Οὐκ ἐκρύβησαν αἱ ἀδικίαι αὐτῶν ἀπʼ αὐτοῦ, καὶ πᾶσαι αἱ ἁμαρτίαι αὐτῶν ἔναντι Κυρίου.

\vs{22}Ἐλεημοσύνη ἀνδρὸς ὡς σφραγὶς μετʼ αὐτοῦ, καὶ χάριν ἀνθρώπου ὡς κόρην συντηρήσει.
\vs{23}Μετὰ ταῦτα ἐξαναστήσεται καὶ ἀνταποδώσει αὐτοῖς, καὶ τὸ ἀνταπόδομα αὐτῶν εἰς κεφαλὴν αὐτῶν ἀποδώσει.
\vs{24}Πλὴν μετανοοῦσιν ἔδωκεν ἐπάνοδον, καὶ παρεκάλεσεν ἐκλείποντας ὑπομονήν.

\vs{25}Ἐπίστρεφε ἐπὶ Κύριον καὶ ἀπόλειπε ἁμαρτίας, δεήθητι κατὰ πρόσωπον καὶ σμίκρυνον πρόσκομμα.
\vs{26}Ἐπάναγε ἐπὶ ὕψιστον, καὶ ἀπόστρεφε ἀπὸ ἀδικίας, καὶ σφόδρα μίσησον βδέλυγμα.
\vs{27}Ὑψίστῳ τίς αἰνέσει ἐν ᾅδου; ἀντὶ ζώντων καὶ ζώντων καὶ διδόντων ἀνθομολόγησιν.
\vs{28}Ἀπὸ νεκροῦ ὡς μηδὲ ὄντος ἀπόλλυται ἐξομολόγησις· ζῶν καὶ ὑγιὴς αἰνέσει τὸν Κύριον.
\vs{29}Ὡς μεγάλη ἡ ἐλεημοσύνη τοῦ Κυρίου, καὶ ἐξιλασμὸς τοῖς ἐπιστρέφουσιν ἐπʼ αὐτόν·
\vs{30}οὐ γὰρ δύναται πάντα εἶναι ἐν ἀνθρώποις, ὅτι οὐκ ἀθάνατος υἱὸς ἀνθρώπου.

\vs{31}Τί φωτεινότερον ἡλίου; καὶ τοῦτο ἐκλείπει, καὶ πονηρὸς ἐνθυμηθήσεται σάρκα καὶ αἷμα.
\vs{32}Δύναμιν ὕψους οὐρανοῦ αὐτὸς ἐπισκέπτεται, καὶ οἱ ἄνθρωποι πάντες γῆ καὶ σποδός.

\ch{18}
Ὁ ζῶν εἰς τὸν αἰῶνα ἔκτισε τὰ πάντα κοινῇ.
\vs{2}Κύριος μόνος δικαιωθήσεται.
\vs{4}Οὐθενὶ ἐξεποίησεν ἐξαγγεῖλαι τὰ ἔργα αὐτοῦ· καὶ τίς ἐξιχνιάσει τὰ μεγαλεῖα αὐτοῦ;
\vs{5}Κράτος μεγαλωσύνης αὐτοῦ τίς ἐξαριθμήσεται; καὶ τίς προσθήσει ἐκδιηγήσασθαι τὰ ἐλέη αὐτοῦ;

\vs{6}Οὐκ ἔστιν ἐλαττῶσαι οὐδὲ προσθεῖναι, καὶ οὐκ ἔστιν ἐξιχνιάσαι τὰ θαυμάσια τοῦ Κυρίου.
\vs{7}Ὅταν συντελέσῃ ἄνθρωπος τότε ἄρχεται, καὶ ὅταν παύσηται τότε ἀπορηθήσεται.
\vs{8}Τί ἄνθρωπος, καὶ τί ἡ χρῆσις αὐτοῦ; τί τὸ ἀγαθὸν αὐτοῦ, καὶ τί τὸ κακὸν αὐτοῦ;
\vs{9}Ἀριθμὸς ἡμερῶν ἀνθρώπου πολλὰ ἔτη ἑκατόν.
\vs{10}Ὡς σταγὼν ὕδατος ἀπὸ θαλάσσης καὶ ψῆφος ἄμμου, οὕτως ὀλίγα ἔτη ἐν ἡμέρᾳ αἰῶνος.
\vs{11}Διὰ τοῦτο ἐμακροθύμησε Κύριος ἐπʼ αὐτοῖς, καὶ ἐξέχεεν ἐπʼ αὐτοὺς τὸ ἔλεος αὐτοῦ.
\vs{12}Εἶδε καὶ ἐπέγνω τὴν καταστροφὴν αὐτῶν ὅτι πονηρὰ, διὰ τοῦτο ἐπλήθυνε τὸν ἐξιλασμὸν αὐτοῦ.
\vs{13}Ἔλεος ἄνθρώπου ἐπὶ τὸν πλησίον αὐτοῦ, ἔλεος δὲ Κυρίου ἐπὶ πᾶσαν σάρκα, ἐλέγχων καὶ παιδεύων καὶ διδάσκων καὶ ἐπιστρέφων ὡς ποιμὴν τὸ ποίμνιον αὐτοῦ.
\vs{14}Τοὺς ἐκδεχομένους παιδείαν ἐλεεῖ, καὶ τοὺς κατασπεύδοντας ἐπὶ τὰ κρίματα αὐτοῦ.

\vs{15}Τέκνον, ἐν ἀγαθοῖς μὴ δῷς μῶμον, καὶ ἐν πάσῃ δόσει λύπην λόγων.
\vs{16}Οὐχὶ καύσωνα ἀναπαύσει δρόσος; οὕτως κρείσσων λόγος ἢ δόσις.
\vs{17}Οὐκ ἰδοὺ λόγος ὑπὲρ δόμα ἀγαθόν; καὶ ἀμφότερα παρὰ ἀνδρὶ κεχαριτωμένῳ.
\vs{18}Μωρὸς ἀχαρίστως ὀνειδιεῖ, καὶ δόσις βασκάνου ἐκτήκει ὀφθαλμούς.
\vs{19}Πρινὴ λαλῆσαι μάνθανε, καὶ πρὸ ἀῤῥωστίας θεραπεύου.
\vs{20}Πρὸ κρίσεως ἐξέταζε σεαυτόν, καὶ ἐν ὥρᾳ ἐπισκοπῆς εὑρήσεις ἐξιλασμόν.
\vs{21}Πρὶν ἀῤῥωστῆσαί σε ταπεινώθητι, καὶ ἐν καιρῷ ἁμαρτημάτων δεῖξον ἐπιστροφήν.

\vs{22}Μὴ ἐμποδισθῇς τοῦ ἀποδοῦναι εὐχὴν εὐκαίρως, καὶ μὴ μείνῃς ἕως θανάτου δικαιωθῆναι.
\vs{23}Πρὶν εὔξασθαι ἑτοίμασον σεαυτὸν, και μὴ γίνου ὡς ἄνθρωπος πειράζων τὸν Κύριον.
\vs{24}Μνήσθητι θυμοῦ ἐν ἡμέραις τελευτῆς, καὶ καιρὸν ἐκδικήσεως ἐν ἀποστροφῇ προσώπου.
\vs{25}Μνήσθητι καιρὸν λιμοῦ ἐν καιρῷ πλησμονῆς, πτωχείαν καὶ ἔνδειαν ἐν ἡμέραις πλούτου.
\vs{26}Ἀπὸ πρωΐθεν ἕως ἑσπέρας μεταβάλλει καιρὸς, καὶ πάντα ἐστὶ ταχινὰ ἔναντι Κυρίου.

\vs{27}Ἄνθρωπος σοφὸς ἐν παντὶ εὐλαβηθήσεται, καὶ ἐν ἡμέραις ἁμαρτιῶν προσέξει ἀπὸ πλημμελείας·
\vs{28}πᾶς συνετὸς ἔγνω σοφίαν, καὶ τῷ εὑρόντι αὐτὴν δώσει ἐξομολόγησιν.
\vs{29}Συνετοὶ ἐν λόγοις καὶ αὐτοὶ ἐσοφίσαντο, καὶ ἀνώμβρησαν παροιμίας ἀκριβεῖς.

ἘΓΚΡΑΤΕΙΑ ΨΥΧΗΣ.

\vs{30}Ὀπίσω τῶν ἐπιθυμιῶν σου μὴ πορεύου, καὶ ἀπὸ τῶν ὀρέξεών σου κωλύου.
\vs{31}Ἐὰν χορηγήσεῃς τῇ ψυχῇ σου εὐδοκίαν ἐπιθυμίας, ποιήσει σε ἐπίχαρμα τῶν ἐχθρῶν σου.
\vs{32}Μὴ εὐφραίνου ἐπὶ πολλῇ τρυφῇ, μηδὲ προσδεθῇς συμβολῇ αὐτῆς.
\vs{33}Μὴ γίνου πτωχὸς συμβολοκοπῶν ἐκ δανεισμοῦ, καὶ οὐδέν σοι ἐστὶν ἐν μαρσυπείῳ.

\ch{19}
Ἐργάτης μέθυσος οὐ πλουτισθήσεται, ὁ ἐξουθενῶν τὰ ὀλίγα κατὰ μικρὸν πεσεῖται.
\vs{2}Οἶνος καὶ γυναῖκες ἀποστήσουσι συνετοὺς, καὶ ὁ κολλώμενος πόρναις τολμηρότερος ἔσται.

\vs{3}Σήτες καὶ σκώληκες κληρονομήσουσιν αὐτόν, καὶ ψυχὴ τολμηρὰ ἐξαρθήσεται.

\vs{4}Ὁ ταχὺ ἐμπιστεύων, κοῦφος καρδίᾳ, καὶ ὁ ἁμαρτάνων εἰς ψυχὴν αὐτοῦ πλημμελήσει.
\vs{5}Ὁ εὐφραινόμενος καρδίᾳ καταγνωσθήσεται,
\vs{6}καὶ ὁ μισῶν λαλιὰν ἐλαττονοῦται κακίᾳ.

\vs{7}Μηδέποτε δευτερώσῃς λόγον, καὶ οὐθέν σοι οὐ μὴ ἐλαττονωθῇ.
\vs{8}Ἐν φίλῳ καὶ ἐν ἐχθρῷ μὴ διηγοῦ, καὶ εἰ μή ἐστί σοι ἁμαρτία, μὴ ἀποκάλυπτε.
\vs{9}Ἀκήκοε γάρ σου καὶ ἐφυλάξατό σε, καὶ ἐν καιρῷ μισήσει σε.
\vs{10}Ἀκήκοας λόγον; συναποθανέτω σοι, θάρσει, οὐ μή σε ῥήξει.
\vs{11}Ἀπὸ προσώπου λόγου ὠδινήσει μωρὸς, ὡς ἀπὸ προσώπου βρέφους ἡ τίκτουσα.
\vs{12}Βέλος πεπηγὸς ἐν μηρῷ σαρκὸς, οὕτως λόγος ἐν κοιλίᾳ μωροῦ.
\vs{13}Ἔλεγξον φίλον, μήποτε οὐκ ἐποίησε, καὶ εἴ τι ἐποίησε, μήποτε προσθῇ.
\vs{14}Ἔλεγξον τὸν φίλον, μήποτε οὐκ εἶπε, καὶ εἰ εἴρηκεν, ἵνα μὴ δευτερώσῃ.
\vs{15}Ἔλεγξον φίλον, πολλάκις γὰρ γίνεται διαβολὴ, καὶ μὴ παντὶ λόγῳ πίστευε.

\vs{16}Ἔστιν ὀλισθαίνων καὶ οὐκ ἀπὸ ψυχῆς· καὶ τίς οὐχ ἡμάρτησεν ἐν τῇ γλώσσῃ αὐτοῦ;

\vs{17}Ἔλεγξον τὸν πλησίον σου πρινη ἀπειλῆσαι, καὶ δὸς τόπον νόμῳ ὑψίστου.
\vs{20}Πᾶσα σοφία φόβος Κυρίου, καὶ ἐν πάσῃ σοφίᾳ ποίησις νόμου.
\vs{22}Καὶ οὐκ ἔστι σοφία πονηρίας ἐπιστήμη, καὶ οὐκ ἔστιν, ὅπου βουλὴ ἁμαρτωλῶν, φρόνησις.
\vs{23}Ἔστι πονηρία καὶ αὕτη βδέλυγμα, καὶ ἔστιν ἄφρων ἐλαττούμενος σοφίᾳ.

\vs{24}Κρείττων ἡττώμενος ἐν συνέσει ἔμφοβος, ἢ περισσεύων ἐν φρονήσει καὶ παραβαίνων νόμον.
\vs{25}Ἔστι πανουργία ἀκριβὴς καὶ αὕτη ἄδικος, καὶ ἔστι διαστρέφων χάριν τοῦ ἐκφᾶναι κρίμα.
\vs{26}Ἔστι πονηρευόμενος συγκεκυφὼς μελανίᾳ, καὶ τὰ ἐντὸς αὐτοῦ πλήρης δόλου.
\vs{27}Συγκύφων πρόσωπον καὶ ἑτεροκωφῶν, ὅπου οὐκ ἐπεγνώσθη, προφθάσει σε.
\vs{28}Καὶ ἐὰν ὑπὸ ἐλαττώματος ἰσχύος κωλυθῇ ἁμαρτεῖν, ἐὰν εὕρῃ καιρὸν, κακοποιήσει.
\vs{29}Ἀπὸ ὁράσεως ἐπιγνωσθήσεται ἀνὴρ, καὶ ἀπὸ ἀπαντήσεως προσώπου ἐπιγνωσθήσεται νοήμων.
\vs{30}Στολισμὸς ἀνδρὸς καὶ γέλως ὀδόντων καὶ βήματα ἀνθρώπου ἀναγγέλλει τὰ περὶ αὐτοῦ.

\ch{20}Ἔστιν ἔλεγχος ὃς οὐκ ἔστιν ὡραῖος, καὶ ἔστι σιωπῶν καὶ αὐτὸς φρόνιμος.

\vs{2}Ὡς καλὸν ἐλέγξαι ἢ θυμοῦσθαι,
\vs{3}καὶ ὁ ἀνθομολογούμενος ἀπὸ ἐλαττώσεως κωλυθήσεται.
\vs{4}Ἐπιθυμία εὐνούχου ἀποπαρθενῶσαι νεανίδα, οὕτως ὁ ποιῶν ἐν βίᾳ κρίματα.

\vs{5}Εστι σιωπῶν εὑρισκόμενος σοφὸς, καὶ ἔστι μισητὸς ἀπὸ πολλῆς λαλιᾶς.
\vs{6}Ἔστι σιωπῶν, οὐ γὰρ ἔχει ἀπόκρισιν, καὶ ἔστι σιωπῶν εἰδὼς καιρόν.
\vs{7}Ἄνθρωπος σοφὸς σιγήσει ἕως καιροῦ, ὁ δὲ λαπιστὴς καὶ ἄφρων ὑπερβήσεται καιρόν.
\vs{8}Ὁ πλεονάζων λόγῳ βδελυχθήσεται, καὶ ὁ ἐνεξουσιαζόμενος μισηθήσεται.
\vs{9}Ἔστιν εὐοδία ἐν κακοῖς ἀνδρί, καὶ ἔστιν εὕρεμα εἰς ἐλάττωσιν.
\vs{10}Ἔστι δόσις ἣ οὐ λυσιτελήσει σοι, καὶ ἔστι δόσις ἧς τὸ ἀνταπόδομα διπλοῦν.
\vs{11}Ἔστιν ἐλάττωσις ἕνεκεν δόξης, καὶ ἔστιν ὃς ἀπὸ ταπεινώσεως ᾖρε κεφαλήν.
\vs{12}Ἔστιν ἀγοράζων πολλὰ ὀλίγου, καὶ ἀποτιννύων αὐτὰ ἑπταπλάσιον.

\vs{13}Ὁ σοφὸς ἐν λόγῳ ἑαυτὸν προσφιλῆ ποιήσει, χάριτες δὲ μωρῶν ἐκχυθήσονται.
\vs{14}Δόσις ἄφρονος οὐ λυσιτελήσει σοι, οἱ γὰρ ὀφθαλμοὶ αὐτοῦ ἀνθʼ ἑνὸς πολλοί.
\vs{15}Ὀλίγα δώσει καὶ πολλὰ ὀνειδίσει, καὶ ἀνοίξει τὸ στόμα αὐτοῦ ὡς κήρυξ· σήμερον δανειεῖ καὶ αὔριον ἀπαιτήσει· μισητὸς ἄνθρωπος ὁ τοιοῦτος.
\vs{16}Μωρὸς ἐρεῖ, οὐχ ὑπάρχει μοι φίλος, καὶ οὐκ ἔστι χάρις τοῖς ἀγαθοῖς μου· οἱ ἔσθοντες τὸν ἄρτον μου, φαῦλοι γλώσσῃ.
\vs{17}Ποσάκις, καὶ ὅσοι καταγελάσονται αὐτοῦ;

\vs{18}Ὀλίσθημα ἀπὸ ἐδάφους μᾶλλον ἢ ἀπὸ γλώσσης, οὕτως πτῶσις κακῶν κατὰ σπουδὴν ἥξει.
\vs{19}Ἄνθρωπος ἄχαρις, μύθος ἄκαιρος, ἐν στόματι ἀπαιδεύτων ἐνδελεχισθήσεται.
\vs{20}Ἀπὸ στόματος μωροῦ ἀποδοκιμασθήσεται παραβολὴ, οὐ γὰρ μὴ εἴπῃ αὐτὴν ἐν καιρῷ αὐτῆς.
\vs{21}Ἔστι κωλυόμενος ἁμαρτάνειν ἀπὸ ἐνδείας, καὶ ἐν τῇ ἀναπαύσει αὐτοῦ οὐ κατανυγήσεται.
\vs{22}Ἔστιν ἀπολλύων τὴν ψυχὴν αὐτοῦ διʼ αἰσχύνην, καὶ ἀπὸ ἄφρονος προσώπου ἀπολεῖ αὐτήν.
\vs{23}Ἔστι χάριν αἰσχύνης ἐπαγγελλόμενος φίλῳ, καὶ ἐκτήσατο αὐτὸν ἐχθρὸν δωρεάν.
\vs{24}Μῶμος πονηρὸς ἐν ἀνθρώπῳ ψεῦδος, ἐν στόματι ἀπαιδεύτων ἐνδελεχισθήσεται.

\vs{25}Αἱρετὸν κλέπτης, ἢ ὁ ἐνδελεχίζων ψεύδει, ἀμφότεροι δὲ ἀπώλειαν κληρονομήσουσιν.
\vs{26}Ἦθος ἀνθρώπου ψευδοῦς ἀτιμία, καὶ ἡ αἰσχύνη αὐτοῦ μετʼ αὐτοῦ ἐνδελεχῶς.

ΛΟΓΟΙ ΠΑΡΑΒΟΛΩΝ.

\vs{27}Ο σοφὸς ἐν λόγοις προάξει ἑαυτὸν, καὶ ἄνθρωπος φρόνιμος ἀρέσει μεγιστᾶσιν.
\vs{28}Ὁ ἐργαζόμενος γῆν ἀνυψώσει θημωνίαν αὐτοῦ, καὶ ὁ ἀρέσκων μεγιστᾶσιν ἐξιλάσεται ἀδικίαν.
\vs{29}Ξένια καὶ δῶρα ἀποτυφλοῖ ὀφθαλμοὺς σοφῶν, καὶ ὡς φιμὸς ἐν στόματι ἀποτρέπει ἐλεγμούς.
\vs{30}Σοφία κεκρυμμένη καὶ θησαυρὸς ἀφανὴς, τίς ὠφέλεια ἐν ἀμφοτέροις;
\vs{31}Κρείσσων ἄνθρωπος ἀποκρύπτων τὴν μωρίαν αὐτοῦ, ἢ ἄνθρωπος ἀποκρύπτων τὴν σοφίαν αὐτοῦ.

\ch{21}
Τέκνον, ἥμαρτες; μὴ προσθῇς μηκέτι, καὶ περὶ τῶν προτέρων σου δεήθητι.
\vs{2}Ὡς ἀπὸ προσώπου ὄφεως, φεῦγε ἀπὸ ἁμαρτίας, ἐὰν γὰρ προσέλθῃς, δήξεταί σε· ὀδόντες λέοντος οἱ ὀδόντες αὐτῆς, ἀναιροῦντες ψυχὰς ἀνθρώπεν.
\vs{3}Ὡς ῥομφαία δίστομος πᾶσα ἀνομία, τῇ πληγῇ αὐτῆς οὐκ ἔστιν ἴασις.
\vs{4}Καταπληγμὸς καὶ ὕβρις ἐρημώσουσι πλοῦτον, οὕτως οἶκος ὑπερηφάνου ἐρημωθήσεται.
\vs{5}Δέησις πτωχοῦ ἐκ στόματος ἕως ὠτίων αὐτοῦ, καὶ τὸ κρίμα αὐτοῦ κατὰ σπουδὴν ἔρχεται.
\vs{6}Μισῶν ἐλεγμὸν, ἐν ἴχνει ἁμαρτωλοῦ, καὶ ὁ φοβούμενος Κύριον ἐπιστρέψει ἐν καρδίᾳ.

\vs{7}Γνωστὸς μακρόθεν ὁ δυνατὸς ἐν γλώσσῃ, ὁ δὲ νοήμων οἶδεν ἐν τῷ ὀλισθαίνειν αὐτόν.
\vs{8}Ὁ οἰκοδομῶν τὴν οἰκίαν αὐτοῦ ἐν χρήμασιν ἀλλοτρίοις, ὡς ὁ συνάγων αὐτοῦ τοὺς λίθους εἰς χειμῶνα.
\vs{9}Στυππεῖον συνηγμένον συναγωγὴ ἀνόμων, καὶ ἡ συντέλεια αὐτῶν φλὸξ πυρός.
\vs{10}Ὁδὸς ἁμαρτωλῶν ὡμαλισμένη ἐκ λίθων, καὶ ἐπʼ ἐσχάτῳ αὐτῆς βόθρος ᾅδου.
\vs{11}Ὁ φύλασσων νόμον κατακρατεῖ τοῦ ἐννοήματος αὐτοῦ, καὶ συντέλεια τοῦ φόβου Κυρίου σοφία.
\vs{12}Οὐ παιδευθήσεται ὃς οὐκ ἔστι πανοῦργος· ἔστι πανουργία πληθύνουσα πικρίαν.
\vs{13}Γνῶσις σοφοῦ ὡς κατακλυσμὸς πληθυνθήσεται, καὶ ἡ βουλὴ αὐτοῦ ὡς πηγὴ ζωῆς.
\vs{14}Ἔγκατα μωροῦ ὡς ἀγγεῖον συντετριμμένον, καὶ πᾶσαν γνῶσιν οὐ κρατήσει.

\vs{15}Λόγον σοφὸν ἐὰν ἀκούσῃ ἐπιστήμων, αἰνέσει αὐτὸν, καὶ ἐπʼ αὐτὸν προσθήσει· ἤκουσεν ὁ σπαταλῶν καὶ ἀπήρεσεν αὐτῷ, καὶ ἀπέστρεψεν αὐτὸν ὀπίσω τοῦ νώτου αὐτοῦ.
\vs{16}Ἐξήγησις μωροῦ ὡς ἐν ὁδῷ φορτίον, ἐπὶ δὲ χείλους συνετοῦ εὑρεθήσεται χάρις.
\vs{17}Στόμα φρονίμου ζητηθήσεται ἐν ἐκκλησίᾳ, καὶ τοὺς λόγους αὐτοῦ διανοηθήσεται ἐν καρδίᾳ.
\vs{18}Ὡς οἶκος ἠφανισμένος οὕτως μωρῷ σοφία· καὶ γνῶσις ἀσυνέτου, ἀδιεξέταστοι λόγοι.
\vs{19}Πέδαι ἐν ποσὶν ἀνόητοις παιδεία, καὶ ὡς χειροπέδαι ἐπὶ χειρὸς δεξιᾶς.
\vs{20}Μωρὸς ἐν γέλωτι ἀνυψοῖ φωνὴν αὐτοῦ, ἀνὴρ δὲ πανοῦργος μόλις ἡσυχῇ μειδιάσει.
\vs{21}Ὡς κόσμος χρυσοῦ φρονίμῳ παιδεία, καὶ ὡς χλιδὼν ἐπὶ βραχίονι δεξιῷ.

\vs{22}Ποὺς μωροῦ ταχὺς εἰς οἰκίαν, ἄνθρωπος δὲ πολύπειρος αἰσχυνθήσεται ἀπὸ προσώπου.
\vs{23}Ἄφρων ἀπὸ θύρας παρακύπτει εἰς οἰκίαν, ἀνὴρ δὲ πεπαιδευμένος ἔξω στήσεται.
\vs{24}Ἀπαιδευσία ἀνθρώπου ἀκροᾶσθαι παρὰ θύραν, ὁ δὲ φρόνιμος βαρυνθήσεται ἀτιμίᾳ.
\vs{25}Χείλη ἀλλοτρίων ἐν τούτοις βαρυνθήσεται, λόγοι δὲ φρονίμων ἐν ζυγῷ σταθήσονται.
\vs{26}Ἐν στόματι μωρῶν ἡ καρδία αὐτῶν, καρδία δὲ σοφῶν στόμα αὐτῶν.
\vs{27}Ἐν τῷ καταρᾶσθαι ἀσεβῆ τὸν Σατανᾶν, αὐτὸς καταρᾶται τὴν ἑαυτοῦ ψυχήν.

\vs{28}Μολύνει τὴν ἑαυτοῦ ψυχὴν ὁ ψιθυρίζων, καὶ ἐν παροικήσει μισηθήσεται.

\ch{22}
Λίθῳ ἠρδαλωμένῳ συνεβλήθη ὀκνηρὸς, καὶ πᾶς ἐκσυριεῖ ἐπὶ τῇ ἀτιμίᾳ αὐτοῦ.
\vs{2}Βολβίτῳ κοπρίων συνεβλήθη ὀκνηρὸς, πᾶς ὁ ἀναιρούμενος αὐτὸν ἐκτινάξει χεῖρα.
\vs{3}Αἰσχύνη πατρὸς ἐν γεννήσει ἀπαιδεύτου, θυγάτηρ δὲ ἐπʼ ἐλαττώσει γίνεται.
\vs{4}Θυγάτηρ φρονίμη κληρονομήσει ἄνδρα αὐτῆς, καὶ ἡ καταισχύνουσα, εἰς λύπην γεννήσαντος.
\vs{5}Πατέρα καὶ ἄνδρα καταισχύνει ἡ θρασεῖα, καὶ ὑπὸ ἀμφοτέρων ἀμιμασθήσεται.

\vs{6}Μουσικὰ ἐν πένθει ἄκαιρος διήγησις, μάστιγες καὶ παιδεία ἐν παντὶ καιρῷ σοφίας.
\vs{9}Συγκολλῶν ὄστρακον ὁ διδάσκων μωρὸν, ἐξεγείρων καθεύδοντα ἐκ βαθέως ὕπνου.
\vs{10}Διηγούμενος νυστάζοντι ὁ διηγούμενος μωρῷ, καὶ ἐπὶ συντελείᾳ ἐρεῖ, τί ἐστιν;
\vs{11}Ἐπὶ νεκρῷ κλαῦσον, ἐξέλιπε γὰρ φῶς· καὶ ἐπὶ μωρῷ κλαῦσον, ἐξέλιπε γὰρ σύνεσις· ἥδιον κλαῦσον ἐπὶ νεκρῷ, ὅτι ἀνεπαύσατο, τοῦ δὲ μωροῦ ὑπὲρ θάνατον ἡ ζωὴ πονηρά.
\vs{12}Πένθος νεκροῦ ἑπτὰ ἡμέραι, μωροῦ δὲ καὶ ἀσεβοῦς πᾶσαι αἱ ἡμέραι τῆς ζωῆς αὐτοῦ.
\vs{13}Μετὰ ἄφρονος μὴ πληθύνῃς λόγον, καὶ πρὸς ἀσύνετον μὴ πορεύου· φύλαξον ἀπʼ αὐτοῦ ἵνα μὴ κόπον ἔχῃς, καὶ οὐ μὴ μολυνθῇς ἐν τῷ ἐντιναγμῷ αὐτοῦ· ἔκκλινον ἀπʼ αὐτοῦ καὶ εὑρήσεις ἀνάπαυσιν, καὶ οὐ μὴ ἀκηδιάσῃς ἐν τῇ ἀπονοίᾳ αὐτοῦ.
\vs{14}Ὑπὲρ μόλυβδον τί βαρυνθήσεται; καὶ τί αὐτῷ ὄνομα, ἀλλʼ ἢ μωρός;
\vs{15}Ἄμμον καὶ ἅλα καὶ βῶλον σιδήρου εὔκοπον ὑπενεγκεῖν, ἢ ἄνθρωπον ἀσύνετον.

\vs{16}Ἱμάντωσις ξυλίνη ἐνδεδεμένη εἰς οἰκοδομὴν ἐν συσσεισμῷ οὐ διαλυθήσεται, οὕτως καρδία ἐστηριγμένη ἐπὶ διανοήματος βουλῆς ἐν καιρῷ οὐ δειλιάσει.
\vs{17}Καρδία ἡδρασμένη ἐπὶ διανοίας συνέσεως, ὡς κόσμος ψαμμωτὸς τοίχου ξυστοῦ.
\vs{18}Χάρακες ἐπὶ μετεώρου κείμενοι κατέναντι ἀνέμου οὐ μὴ ὑπομείνωσιν, οὕτως καρδία δειλὴ ἐπὶ διανοήματος μωροῦ κατέναντι παντὸς φόβου οὐ μὴ ὑπομείνῃ.

\vs{19}Ὁ νύσσων ὀφθαλμὸν κατάξει δάκρυα, καὶ ὁ νύσσων καρδίαν ἐκφαίνει αἴσθησιν.
\vs{20}Βάλλων λίθον ἐπὶ πετεινὰ ἀποσοβεῖ αὐτὰ, καὶ ὁ ὀνειδίζων φίλον διαλύσει φιλίαν.
\vs{21}Ἐπὶ φίλον ἐὰν σπάσῃς ῥομφαίαν, μὴ ἀπελπίσῃς, ἔστι γὰρ ἐπάνοδος.
\vs{22}Ἐπὶ φίλον ἐὰν ἀνοίξῃς στόμα, μὴ εὐλαβηθῇς, ἔστι γὰρ διαλλαγή· πλὴν ὀνειδισμοῦ, καὶ ὑπερηφανίας, καὶ μυστηρίου ἀποκαλύψεως, καὶ πληγῆς δολίας, ἐν τούτοις ἀποφεύξεται πᾶς φίλος.

\vs{23}Πίστιν κτῆσαι ἐν πτωχείᾳ μετὰ τοῦ πλησίον, ἵνα ἐν τοῖς ἀγαθοῖς αὐτοῦ ὁμοῦ πλησθῇς· ἐν καιρῷ θλίψεως διάμενε αὐτῷ, ἵνα ἐν τῇ κληρονομίᾳ αὐτοῦ συνκληρονομήσῃς.
\vs{24}Πρὸ πυρὸς ἀτμὶς καμίνου καὶ καπνὸς, οὕτως πρὸ αἱμάτων λοιδορίαι.
\vs{25}Φίλον σκεπάσαι οὐκ αἰσχυνθήσομαι, καὶ ἀπὸ προσώπου αὐτοῦ οὐ μὴ κρυβῶ,
\vs{26}καὶ εἰ κακά μοι συμβῇ διʼ αὐτόν, πᾶς ὁ ἀκούων φυλάξεται ἀπʼ αὐτοῦ.
\vs{27}Τίς δώσει μοι ἐπὶ στόμα μου φυλακὴν, καὶ ἐπὶ τῶν χειλέων μου σφραγίδα πανοῦργον, ἵνα μὴ πέσω ἀπʼ αὐτῆς, καὶ ἡ γλῶσσά μου ἀπολέσῃ με;

\ch{23}
Κύριε πάτερ καὶ δέσποτα ζωῆς μου, μὴ ἐγκαταλίπῃς με ἐν βουλῇ αὐτῶν, μὴ ἀφῇς με πεσεῖν ἐν αὐτοῖς.
\vs{2}Τίς ἐπιστήσει ἐπὶ τοῦ διανοήματός μου μάστιγας, καὶ ἐπὶ τῆς καρδίας μου παιδείαν σοφίας; ἵνα ἐπὶ τοῖς ἀγνοήμασί μου μὴ φείσωνται, καὶ οὐ μὴ παρῇ τὰ ἁμαρτήματα αὐτῶν,
\vs{3}ὅπως μὴ πληθύνωσιν αἱ ἄγνοιαί μου, καὶ αἱ ἁμαρτίαι μου πλεονάσωσιν, καὶ πεσοῦμαι ἔναντι τῶν ὑπεναντίων, καὶ ἐπιχαρεῖταί μοι ὁ ἐχθρός μου.

\vs{4}Κύριε πάτερ καὶ Θεὲ ζωῆς μου, μετεωρισμὸν ὀφθαλμῶν μὴ δῷς μοι,
\vs{5}καὶ ἐπιθυμίαν ἀπόστρεψον ἀπʼ ἐμοῦ.
\vs{6}Κοιλίας ὄρεξις καὶ συνουσιασμὸς μὴ καταλαβέτωσάν με, καὶ ψυχῇ ἀναιδεῖ μὴ παραδῷς με.

ΠΑΙΔΕΙΑ ΣΤΟΜΑΤΟΣ.

\vs{7}Παιδείαν στόματος ἀκούσατε τέκνα, καὶ ὁ φυλάσσων οὐ μὴ ἁλῷ ἐν τοῖς χείλεσιν αὐτοῦ.
\vs{8}Καταλειφθήσεται ἁμαρτωλὸς, καὶ λοίδορος καὶ ὑπερήφανος σκανδαλισθήσονται ἐν αὐτοῖς.
\vs{9}Ὅρκῳ μὴ ἐθίσῃς τὸ στόμα σου, καὶ ὀνομασίᾳ τοῦ ἁγίου μὴ συνεθισθῇς.
\vs{10}Ὥσπερ γὰρ οἰκέτης ἐξεταζόμενος ἐνδελεχῶς ἀπὸ μώλωπος οὐκ ἐλαττωθήσεται, οὕτως ὁ καὶ ὀμνύων καὶ ὀνομάζων διαπαντὸς ἀπὸ ἁμαρτίας οὐ μὴ καθαρισθῇ.
\vs{11}Ἀνὴρ πολύορκος πλησθήσεται ἀνομίας, καὶ οὐκ ἀποστήσεται ἀπὸ τοῦ οἴκου αὐτοῦ μάστιξ· ἐὰν πλημμελήσῃ, ἁμαρτία αὐτοῦ ἐπʼ αὐτῷ, κᾄν ὑπερίδῃ, ἥμαρτε δισσῶς· καὶ εἰ διακενῆς ὤμοσεν, οὐ δικαιωθήσεται, πλησθήσεται γὰρ ἐπαγωγῶν ὁ οἶκος αὐτοῦ.
\vs{12}Ἔστι λέξις ἀντιπεριβεβλημένη θανάτῳ, μὴ εὑρεθήτω ἐν κληρονομίᾳ Ἰακώβ· ἀπὸ γὰρ εὐσεβῶν ταῦτα πάντα ἀποστήσεται, καὶ ἐν ἁμαρτίαις οὐκ ἐγκυλισθήσονται.
\vs{13}Ἀπαιδευσίαν ἀσυρῆ μὴ συνεθίσῃς τὸ στόμα σου, ἔστι γὰρ ἐν αὐτῇ λόγος ἁμαρτίας.

\vs{14}Μνήσθητι πατρὸς καὶ μητρός σου, ἀναμέσον γὰρ μεγιστάνων συνεδρεύεις· μήποτʼ ἐπιλάθῃ ἐνώπιον αὐτῶν, καὶ τῷ ἐθισμῷ σου μωρανθῇς, καὶ θελήσεις εἰ μὴ ἐγεννήθης, καὶ τὴν ἡμέραν τοῦ τοκετοῦ σου καταράσῃ.
\vs{15}Ἄνθρωπος συνεθιζόμενος λόγοις ὀνειδισμοῦ, ἐν πάσαις ταῖς ἡμέραις αὐτοῦ οὐ μὴ παιδευθῇ.

\vs{16}Δύο εἴδη πληθύνουσιν ἁμαρτίας, καὶ τὸ τρίτον ἐπάξει ὀργήν· ψυχὴ θερμὴ ὡς πῦρ καιόμενον, οὐ μὴ σβεσθῆ ἕως ἂν καταποθῇ· ἄνθρωπος πόρνος ἐν σώματι σαρκὸς αὐτοῦ, οὐ μὴ παύσηται ἕως ἂν ἐκκαύσῃ πῦρ.
\vs{17}Ἀνθρώπῳ πόρνῳ πᾶς ἄρτος ἡδὺς, οὐ μὴ κοπάσῃ ἕως ἂν τελευτήσῃ.

\vs{18}Ἄνθρωπος παραβαίνων ἀπὸ τῆς κλίνης αὐτοῦ, λέγων ἐν τῇ ψυχῇ αὐτοῦ, τίς μὲ ὁρᾷ; σκότος κύκλῳ μου, καὶ οἱ τοῖχοί με καλύπτουσι, καὶ οὐθείς με ὁρᾷ, τί εὐλαβοῦμαι; τῶν ἁμαρτιῶν μου οὐ μὴ μνησθήσεται ὁ ὕψιστος·
\vs{19}καὶ ὀφθαλμοὶ ἀνθρώπων ὁ φόβος αὐτοῦ· καὶ οὐκ ἔγνω ὅτι ὀφθαλμοὶ Κυρίου μυριοπλασίως ἡλίου φωτεινότεροι, ἐπιβλέποντες πάσας ὁδοὺς ἀνθρώπων, καὶ κατανοοῦντες εἰς ἀπόκρυφα μέρη.
\vs{20}Πρινὴ κτισθῆναι τὰ πάντα ἔγνωσται αὐτῷ, οὕτως καὶ μετὰ τὸ συντελεσθῆναι.
\vs{21}Οὗτος ἐν πλατείαις πόλεως ἐκδικηθήσεται, καὶ οὗ οὐχ ὑπενόησε πιασθήσεται.

\vs{22}Οὕτως καὶ γυνὴ καταλιποῦσα τὸν ἄνδρα, καὶ παριστῶσα κληρονόμον ἐξ ἀλλοτρίου.
\vs{23}Πρῶτον μὲν γὰρ ἐν νόμῳ ὑψίστου ἠπείθησε, καὶ δεύτερον εἰς ἄνδρα ἑαυτῆς ἐπλημμέλησε, καὶ τὸ τρίτον ἐν πορνείᾳ ἐμοιχεύθη, ἐξ ἀλλοτρίου ἀνδρὸς τέκνα παρέστησεν.
\vs{24}Αὕτη εἰς ἐκκλησίαν ἐξαχθήσεται, καὶ ἐπὶ τὰ τέκνα αὐτῆς ἐπισκοπὴ ἔσται.
\vs{25}Οὐ διαδώσουσι τὰ τέκνα αὐτῆς εἰς ῥίζαν, καὶ οἱ κλάδοι αὐτῆς οὐ δώσουσι καρπόν.
\vs{26}Καταλείψει εἰς κατάραν τὸ μνημόσυνον αὐτῆς, καὶ τὸ ὄνειδος αὐτῆς οὐκ ἐξαλειφθήσεται.
\vs{27}Καὶ ἐπιγνώσονται οἱ καταλειφθέντες, ὅτι οὐθὲν κρεῖττον φόβου Κυρίου, καὶ οὐθὲν γλυκύτερον τοῦ προσέχειν ἐντολαῖς Κυρίου.

ΑΙΝΕΣΙΣ ΣΟΦΙΑΣ.

\ch{24}
Ἡ σοφία αἰνέσει ψυχὴν αὐτῆς, καὶ ἐν μέσῳ λαοῦ αὐτῆς καυχήσεται.
\vs{2}Ἐν ἐκκλησίᾳ ὑψίστου στόμα αὐτῆς ἀνοίξει, καὶ ἔναντι δυνάμεως αὐτοῦ καυχήσεται.
\vs{3}Ἐγὼ ἀπὸ στόματος ὑψίστου ἐξῆλθον, καὶ ὡς ὁμίχλη κατεκάλυψα γῆν.
\vs{4}Ἐγὼ ἐν ὑψηλοῖς κατεσκήνωσα, καὶ ὁ θρόνος μου ἐν στύλῳ νεφέλης.
\vs{5}Γῦρον οὐρανοῦ ἐκύκλωσα μόνη, καὶ ἐν βάθει ἀβύσσων περιεπάτησα.
\vs{6}Ἐν κύμασι θαλάσσης καὶ ἐν πάσῃ τῇ γῇ, καὶ ἐν παντὶ λαῷ καὶ ἔθνει ἐκτησάμην.

\vs{7}Μετὰ τούτων πάντων ἀνάπαυσιν ἐζήτησα, καὶ ἐν κληρονομίᾳ τίνος αὐλισθήσομαι.
\vs{8}Τότε ἐνετείλατό μοι ὁ κτίστης ἁπάντων, καὶ ὁ κτίσας με κατέπαυσε τὴν σκηνήν μου, καὶ εἶπεν, ἐν Ἰακὼβ κατασκήνωσον, καὶ ἐν Ἰσραὴλ κατακληρονομήθητι.
\vs{9}Πρὸ τοῦ αἰῶνος ἀπʼ ἀρχῆς ἔκτισέ με, καὶ ἕως αἰῶνος οὐ μὴ ἐκλίπω.

\vs{10}Ἐν σκηνῇ ἁγίᾳ ἐνώπιον αὐτοῦ ἐλειτούργησα, καὶ οὕτως ἐν Σιὼν ἐστηρίχθην.
\vs{11}Ἐν πόλει ἠγαπημένῃ ὁμοίως με κατέπαυσε, καὶ ἐν Ἰερουσαλὴμ ἡ ἐξουσία μου.
\vs{12}Καὶ ἐῤῥίζωσα ἐν λαῷ δεδοξασμένῳ, ἐν μερίδι Κυρίου κληρονομίας αὐτοῦ.
\vs{13}Ὡς κέδρος ἀνυψώθην ἐν Λιβάνῳ, καὶ ὡς κυπάρισσος ἐν ὄρεσιν Ἀερμών.
\vs{14}Ὡς φοῖνιξ ἀνυψώθην ἐν αἰγιαλοῖς, καὶ ὡς φυτὰ ῥόδου ἐν Ἱεριχῶ· ὡς ἐλαία εὐπρεπὴς ἐν πεδίῳ, καὶ ἀνυψώθην ὡς πλάτανος.
\vs{15}Ὡς κιννάμωμον καὶ ἀσπάλαθος ἀρωμάτων δέδωκα ὀσμὴν, καὶ ὡς σμύρνα ἐκλεκτὴ διέδωκα εὐωδίαν· ὡς χαλβάνη καὶ ὄνυξ καὶ στακτὴ, καὶ ὡς λιβάνου ἀτμὶς ἐν σκηνῇ.
\vs{16}Ἐγὼ ὡς τερέμινθος ἐξέτεινα κλάδους μου, καὶ οἱ κλάδοι μου κλάδοι δόξης καὶ χάριτος.
\vs{17}Ἐγὼ ὡς ἄμπελος βλαστήσασα χάριν, καὶ τὰ ἄνθη μου καρπὸς δόξης καὶ πλούτου.
\vs{19}Προσέλθετε πρὸς μὲ οἱ ἐπιθυμοῦντές μου, καὶ ἀπὸ τῶν γεννημάτων μου ἐμπλήσθητε.
\vs{20}Τὸ γὰρ μνημόσυνόν μου ὑπὲρ μέλι γλυκύ, καὶ ἡ κληρονομία μου ὑπὲρ μέλιτος κηροῦ.
\vs{21}Οἱ ἐσθίοντές με ἔτι πεινάσουσι, καὶ οἱ πίνοντές με ἔτι διψήσουσιν.
\vs{22}Ὁ ὑπακούων μου οὐκ αἰσχυνθήσεται, καὶ οἱ ἐργαζόμενοι ἐν ἐμοὶ οὐχ ἁμαρτήσουσι.

\vs{23}Ταῦτα πάντα βίβλος διαθήκης Θεοῦ ὑψίστου, νόμον ὃν ἐνετείλατο Μωυσῆς, κληρονομίαν συναγωγαῖς Ἰακώβ.
\vs{25}Ὁ πιμπλῶν ὡς Φεισὼν σοφίαν, καὶ ὡς Τίγρις ἐν ἡμέραις νέων·
\vs{26}ὁ ἀναπληρῶν ὡς Εὐφράτης σύνεσιν, καὶ ὡς Ἰορδάνης ἐν ἡμέραις θερισμοῦ·
\vs{27}ὁ ἐκφαίνων ὡς φῶς παιδείαν, ὡς Γηὼν ἐν ἡμέραις τρυγητοῦ.

\vs{28}Οὐ συνετέλεσεν ὁ πρῶτος γνῶναι αὐτὴν, καὶ οὕτως ὁ ἔσχατος οὐκ ἐξιχνίασεν αὐτήν·
\vs{29}Ἀπὸ γὰρ θαλάσσης ἐπληθύνθη διανόημα αὐτῆς, καὶ ἡ βουλὴ αὐτῆς ἀπὸ ἀβύσσου μεγάλης.
\vs{30}Κᾀγὼ ὡς διώρυξ ἀπὸ ποταμοῦ, καὶ ὡς ὑδραγωγὸς ἐξῆλθον εἰς παράδεισον.
\vs{31}Εἶπα, ποτιῶ μου τὸν κῆπον, καὶ μεθύσω μου τὴν πρασιάν· καὶ ἰδοὺ ἐγένετό μοι ἡ διώρυξ εἰς ποταμὸν, καὶ ὁ ποταμός μου ἐγένετο εἰς θάλασσαν.
\vs{32}Ἔτι παιδείαν ὡς ὄρθρον φωτιῶ, καὶ ἐκφανῶ αὐτὰ ἕως εἰς μακράν.
\vs{33}Ἔτι διδασκαλίαν ὡς προφητείαν ἐκχεῶ, καὶ καταλείψω αὐτὴν εἰς γενεὰς αἰώνων.
\vs{34}Ἴδετε ὅτι οὐκ ἐμοὶ μόνῳ ἐκοπίασα, ἀλλὰ πᾶσι τοῖς ἐκζητοῦσιν αὐτήν.

\ch{25}
Ἐν τρισὶν ὡραΐσθην, καὶ ἀνέστην ὡραία ἔναντι Κυρίου καὶ ἀνθρώπων· ὁμόνοια ἀδελφῶν, καὶ φιλία τῶν πλησίον, καὶ γυνὴ καὶ ἀνὴρ ἑαυτοῖς συπεριφερόμενοι.

\vs{2}Τρία δὲ εἴδη ἐμίσησεν ἡ ψυχή μου, καὶ προσώχθισα σφόδρα τῇ ζωῇ αὐτῶν· πτωχὸν ὑπερήφανον, καὶ πλούσιον ψεύστην, γέροντα μοιχὸν ἐλαττούμενον συνέσει.

\vs{3}Ἐν νεότητι οὐ συναγήοχας, καὶ πῶς ἂν εὕροις ἐν τῷ γήρᾳ σου;

\vs{4}Ὡς ὡραῖον πολιαῖς κρίσις, καὶ πρεσβυτέροις ἐπιγνῶναι βουλην;
\vs{5}Ὡς ὡραία γερόντων σοφία, καὶ δεδοξασμένοις διανόημα καὶ βουλή.
\vs{6}Στέφανος γερόντων πολυπειρία, καὶ τὸ καύχημα αὐτῶν φόβος Κυρίου.

\vs{7}Ἐννέα ὑπονοήματα ἐμακάρισα ἐν καρδίᾳ, καὶ τὸ δέκατον ἐρῶ ἐπὶ γλώσσης· ἄνθρωπος εὐφραινόμενος ἐπὶ τέκνοις, ζῶν καὶ βλέπων ἐπὶ πτώσει ἐχθρῶν.
\vs{8}Μακάριος ὁ συνοικῶν γυναικὶ συνετῇ, καὶ ὃς ἐν γλώσσῃ οὐκ ὠλίσθησε, καὶ ὃς οὐκ ἐδούλευσεν ἀναξίῳ αὐτοῦ.
\vs{9}Μακάριος ὃς εὗρε φρόνησιν, καὶ ὁ διηγούμενος εἰς ὦτα ἀκουόντων.
\vs{10}Ὡς μέγας ὁ εὑρὼν σοφίαν, ἀλλʼ οὐκ ἔστιν ὑπὲρ τὸν φοβούμενον τὸν Κύριον.
\vs{11}Φόβος Κυρίου ὑπὲρ πᾶν ὑπερέβαλεν, ὁ κρατῶν αὐτοῦ τίνι ὁμοιωθήσεται;

\vs{13}Πᾶσαν πληγήν καὶ μὴ πληγὴν καρδίας, καὶ πᾶσαν πονηρίαν καὶ μὴ πονηρίαν γυναικός·
\vs{14}πᾶσαν ἐπαγωγήν καὶ μὴ ἐπαγωγὴν μισούντων, καὶ πᾶσαν ἐκδίκησιν καὶ μὴ ἐκδίκησιν ἐχθρῶν.
\vs{15}Οὐκ ἔστι κεφαλὴ ὑπὲρ κεφαλὴν ὄφεως, καὶ οὐκ ἔστι θυμὸς ὑπὲρ θυμὸν ἐχθροῦ.

\vs{16}Συνοικῆσαι λέοντι καὶ δράκοντι εὐδοκήσω, ἢ ἐνοικῆσαι μετὰ γυναικὸς πονηρᾶς.
\vs{17}Πονηρία γυναικὸς ἀλλοιοῖ τὴν ὅρασιν αὐτῆς, καὶ σκοτοῖ τὸ πρόσωπον αὐτῆς ὡς σάκκον.
\vs{18}Ἀναμέσον τοῦ πλησίον αὐτοῦ ἀναπεσεῖται ὁ ἀνὴρ αὐτῆς, καὶ ἀκούσας ἀνεστέναξε πικρά.
\vs{19}Μικρὰ πᾶσα κακία πρὸς κακίαν γυναικός· κλῆρος ἁμαρτωλοῦ ἐπιπέσοι αὐτῇ.

\vs{20}Ἀνάβασις ἀμμώδης ἐν ποσὶ πρεσβυτέρου· οὕτως γυνὴ γλωσσώδης ἀνδρὶ ἡσύχῳ.
\vs{21}Μὴ προσπέσῃς ἐπὶ κάλλος γυναικὸς, καὶ γυναῖκα μὴ ἐπιποθήσῃς.
\vs{22}Ὀργὴ καὶ ἀναίδεια καὶ αἰσχύνη μεγάλη, γυνὴ ἐὰν ἐπιχορηγῇ τῷ ἀνδρὶ αὐτῆς.
\vs{23}Καρδία ταπεινὴ καὶ πρόσωπον σκυθρωπὸν καὶ πληγὴ καρδίας γυνὴ πονηρά· χεῖρες παρειμέναι καὶ γόνατα παραλελυμένα, ἥτις οὐ μακαριεῖ τὸν ἄνδρα αὐτῆς.
\vs{24}Ἀπὸ γυναικὸς ἀρχὴ ἁμαρτίας, καὶ διʼ αὐτὴν ἀποθνήσκομεν πάντες.
\vs{25}Μὴ δῷς ὕδατι διέξοδον, μηδὲ γυναικὶ πωνηρᾷ ἐξουσίαν.
\vs{26}Εἰ μὴ πορεύεται κατὰ χεῖρά σου, ἀπὸ τῶν σαρκῶν σου ἀπότεμε αὐτήν.

\ch{26}
Γυναικὸς ἀγαθῆς μακάριος ὁ ἀνὴρ, καὶ ἀριθμὸς τῶν ἡμερῶν αὐτοῦ διπλάσιος.
\vs{2}Γυνὴ ἀνδρεία εὐφραίνει τὸν ἄνδρα αὐτῆς, καὶ τὰ ἔτη αὐτοῦ πληρώσει ἐν εἰρήνῃ.
\vs{3}Γυνὴ ἀγαθὴ μερὶς ἀγαθὴ, ἐν μερίδι φοβουμένων Κύριον δοθήσεται.
\vs{4}Πλουσίου δὲ καὶ πτωχοῦ καρδία ἀγαθὴ, ἐν παντὶ καιρῷ πρόσωπον ἱλαρόν.

\vs{5}Ἀπὸ τριῶν εὐλαβήθη ἡ καρδία μου, καὶ ἐπὶ τῷ τετάρτῳ προσώπῳ ἐδεήθην· διαβολὴν πόλεως, καὶ ἐκκλησίαν ὄχλου, καὶ καταψευσμὸν ὑπὲρ θάνατον, πάντα μοχθηρά.
\vs{6}Ἄλγος καρδίας καὶ πένθος γυνὴ ἀντίζηλος ἐπὶ γυναικὶ, καὶ μάστιξ γλώσσης πᾶσιν ἐπικοινωνοῦσα.
\vs{7}Βοοζύγιον σαλευόμενον γυνὴ πονηρὰ, ὁ κρατῶν αὐτῆς ὡς ὁ δρασσόμενος σκορπίου.
\vs{8}Ὀργὴ μεγάλη γυνὴ μέθυσος, καὶ ἀσχημοσύνην αὐτῆς οὐ συγκαλύψει.
\vs{9}Πορνεία γυναικὸς ἐν μετεωρισμοῖς ὀφθαλμῶν, καὶ ἐν τοῖς βλεφάροις αὐτῆς γνωσθήσεται.
\vs{10}Ἐπὶ θυγατρὶ ἀδιατρέπτῳ στερέωσον φυλακὴν, ἵνα μὴ εὑροῦσα ἄνεσιν ἑαυτῇ χρήσηται.
\vs{11}Ὀπίσω ἀναιδοῦς ὀφθαλμοῦ φύλαξαι, καὶ μὴ θαυμάσῃς ἐὰν εἰς σὲ πλημμελήσῃ·
\vs{12}Ὡς διψῶν ὁδοιπόρος τὸ στόμα ἀνοίγει, καὶ ἀπὸ παντὸς ὕδατος τοῦ σύνεγγυς πίεται, κατέναντι παντὸς πασσάλου καθήσεται, καὶ ἔναντι βέλους ἀνοίξει φαρέτραν.

\vs{13}Χάρις γυναικὸς τέρψει τὸν ἄνδρα αὐτῆς, καὶ τὰ ὀστᾶ αὐτοῦ πιανεῖ ἡ ἐπιστήμη αὐτῆς.
\vs{14}Δόσις Κυρίου γυνὴ σιγηρὰ, καὶ οὐκ ἐστιν ἀντάλλαγμα πεπαιδευμένης ψυχῆς.
\vs{15}Χάρις ἐπὶ χάριτι γυνὴ αἰσχυντηρὰ, καὶ οὐκ ἔστι σταθμὸς πᾶς ἄξιος ἐγκρατοῦς ψυχῆς.
\vs{16}Ἥλιος ἀνατέλλων ἐν ὑψίστοις Κυρίου, καὶ κάλλος ἀγαθῆς γυναικὸς ἐν κόσμῳ οἰκίας αὐτοῦ.
\vs{17}Λύχνος ἐκλάμπων ἐπὶ λυχνίας ἁγίας, καὶ κάλλος προσώπου ἐπὶ ἡλικίᾳ στασίμῃ.
\vs{18}Στύλοι χρύσεοι ἐπὶ βάσεως ἀργυρᾶς, καὶ πόδες ὡραῖοι ἐπὶ στέρνοις εὐσταθοῦς.

\vs{28}Επὶ δυσὶ λελύπηται ἡ καρδία μου, καὶ ἐπὶ τῷ τρίτῳ θυμός μοι ἐπῆλθεν· ἀνὴρ πολεμιστὴς ὑστερῶν διʼ ἔνδειαν, καὶ ἄνδρες συνετοὶ ἐὰν σκυβαλισθῶσιν· ἐπανάγων ἀπὸ δικαιοσύνης ἐπὶ ἁμαρτίαν, ὁ Κύριος ἑτοιμάσει εἰς ῥομφαίαν αὐτόν.
\vs{29}Μόλις ἐξελεῖται ἔμπορος ἀπὸ πλημμελείας, καὶ οὐ δικαιωθήσεται κάπηλος ἀπὸ ἁμαρτίας.

\ch{27}
Χαρὶν ἀδιαφόρου πολλοὶ ἥμαρτον, καὶ ὁ ζητῶν πληθύναι ἀποστρέψει ὀφθαλμόν.
\vs{2}Ἀναμέσον ἁρμῶν λίθων παγήσεται πάσσαλος, καὶ ἀναμέσον πράσεως καὶ ἀγορασμοῦ συντριβήσεται ἁμαρτία.
\vs{3}Ἐὰν μὴ ἐν φόβῳ Κυρίου κρατήσῃ κατὰ σπουδὴν, ἐν τάχει καταστραφήσεται αὐτοῦ ὁ οἶκος.
\vs{4}Ἐν σείσματι κοσκίνου διαμένει κοπρία, οὕτως σκύβαλα ἀνθρώπου ἐν λογισμῷ αὐτοῦ.
\vs{5}Σκεύη κεραμέως δοκιμάζει κάμινος, καὶ πειρασμὸς ἀνθρώπου ἐν διαλογισμῷ αὐτοῦ.
\vs{6}Γεώργιον ξύλου ἐκφαίνει ὁ καρπὸς αὐτοῦ, οὕτως λόγος ἐνθυμήματος καρδίας ἀνθρώπου.

\vs{7}Πρὸ λογισμοῦ μὴ ἐπαινέσῃς ἄνδρα, οὗτος γὰρ πειρασμὸς ἀνθρώπων.
\vs{8}Ἐὰν διώκῃς τὸ δίκαιον, καταλήψῃ, καὶ ἐνδύσῃ αὐτὸ ὡς ποδήρη δόξης.
\vs{9}Πετεινὰ πρὸς τὰ ὅμοια αὐτοῖς καταλύσει, καὶ ἀλήθεια πρὸς τοὺς ἐργαζομένους αὐτὴν ἐπανήξει.
\vs{10}Λέων θήραν ἐνεδρεύει, οὕτως ἁμαρτίαι ἐργαζομένους ἄδικα.
\vs{11}Διήγησις εὐσεβοῦς διαπαντὸς σοφία, ὁ δὲ ἄφρων ὡς σελήνη ἀλλοιοῦται.
\vs{12}Εἰς μέσον ἀσυνέτων συντήρησον καιρὸν, εἰς μέσον δὲ διανοουμένων ἐνδελέχιζε.
\vs{13}Διήγησις μωρῶν προσόχθισμα, καὶ ὁ γέλως αὐτῶν ἐν σπατάλῃ ἁμαρτίας.
\vs{14}Λαλιὰ πολυόρκου ὀρθώσει τρίχας, καὶ ἡ μάχη αὐτῶν ἐμφραγμὸς ὠτίων.
\vs{15}Ἔκχυσις αἵματος μάχη ὑπερηφάνων, καὶ ἡ διαλοιδόρησις αὐτῶν ἀκοὴ μοχθηρά.

\vs{16}Ὁ ἀποκαλύπτων μυστήρια ἀπώλεσε πίστιν, καὶ οὐ μὴ εὕρῃ φίλον πρὸς τὴν ψυχὴν αὐτοῦ.
\vs{17}Στέρξον φίλον, καὶ πιστώθητι μετʼ αὐτοῦ· ἐὰν δὲ ἀποκαλύψῃς τὰ μυστήρια αὐτοῦ, οὐ μὴ καταδιώξῃς ὀπίσω αὐτοῦ.
\vs{18}Καθὼς γὰρ ἀπώλεσεν ἄνθρωπος τὸν ἐχθρὸν αὐτοῦ, οὕτως ἀπώλεσας τὴν φιλίαν τοῦ πλησίου·
\vs{19}καὶ ὡς πετεινὸν ἐκ χειρός σου ἀπελύσας, οὕτως ἀφῆκας τὸν πλησίον, καὶ οὐ θηρεύσεις αὐτόν.
\vs{20}Μὴ αὐτὸν διώξῃς, ὅτι μακρὰν ἀπέστῃ, καὶ ἐξέφυγεν ὡς δορκὰς ἐκ παγίδος.
\vs{21}Ὅτι θραῦσμά ἐστι καταδῆσαι, καὶ λοιδορίας ἐστὶ διαλλαγή· ὁ δὲ ἀποκαλύψας μυστήρια ἀπήλπισε.
\vs{22}Διανεύων ὀφθαλμῷ τεκταίνει κακὰ, καὶ οὐδεὶς αὑτὸν ἀποστήσει ἀπʼ αὐτοῦ.
\vs{23}Ἀπέναντι τῶν ὀφθαλμῶν σου γλυκανεῖ στόμα σου, καὶ ἐπὶ τῶν λόγων σου ἐκθαυμάσει, ὕστερον δὲ διαστρέψει τὸ στόμα αὐτοῦ, καὶ ἐν τοῖς λόγοις σου δώσει σκάνδαλον.
\vs{24}Πολλὰ ἐμίσησα καὶ οὐχ ὡμοίωσα αὐτῷ, καὶ ὁ Κύριος μισήσει αὐτόν.

\vs{25}Ὁ βάλλων λίθον εἰς ὕψος ἐπὶ κεφαλὴν αὐτοῦ βάλλει, καὶ πληγὴ δολία διελεῖ τραύματα.
\vs{26}Ὁ ὀρύσσων βόθρον εἰς αὐτὸν ἐμπεσεῖται, καὶ ὁ ἱστῶν παγίδα ἐν αὐτῇ ἁλώσεται.
\vs{27}Ὁ ποιῶν πονηρὰ εἰς αὐτὸν κυλισθήσεται, καὶ οὐ μὴ ἐπιγνῷ πόθεν ἥκει αὐτῷ.
\vs{28}Ἐμπαιγμὸς καὶ ὀνειδισμὸς ὑπερηφάνων, καὶ ἡ ἐκδίκησις ὡς λέων ἐνεδρεύσει αὐτόν.
\vs{29}Παγίδι ἁλώσονται οἱ εὐφραινόμενοι πτώσει εὐσεβῶν, καὶ ὀδύνη καταναλώσει αὐτοὺς πρὸ τοῦ θανάτου αὐτῶν.
\vs{30}Μῆνις καὶ ὀργὴ, καὶ ταῦτά ἐστι βδελύγματα, καὶ ἀνὴρ ἁμαρτωλὸς ἐγκρατὴς ἔσται αὐτῶν.

\ch{28}
Ὁ ἐκδίκων παρὰ Κυρίου εὑρήσει ἐκδίκησιν, καὶ τὰς ἁμαρτίας αὐτοῦ διατηρῶν διατηρήσει.
\vs{2}Ἄφες ἀδίκημα τῷ πλησίον σου, καὶ τότε δεηθέντος σου αἱ ἁμαρτίαι σου λυθήσονται.
\vs{3}Ἄνθρωπος ἀνθρώπῳ συντηρεῖ ὀργὴν, καὶ παρὰ Κυρίου ζητεῖ ἴασιν.
\vs{4}Ἐπʼ ἄνθρωπον ὅμοιον αὐτῷ οὐκ ἔχει ἔλεος, καὶ περὶ τῶν ἁμαρτιῶν αὐτοῦ δεῖται.
\vs{5}Αὐτὸς σὰρξ ὢν διατηρεῖ μῆνιν, τίς ἐξιλάσεται τὰς ἁμαρτίας αὐτοῦ;
\vs{6}Μνήσθητι τὰ ἔσχατα, καὶ παῦσαι ἔχθραίνων· καταφθορὰν καὶ θάνατον, καὶ ἔμμενε ἐντολαῖς.
\vs{7}Μνήσθητι ἐντολῶν, καὶ μὴ μηνίσῃς τῷ πλησίον· καὶ διαθήκην ὑψίστου, καὶ πάριδε ἄγνοιαν.
\vs{8}Ἀπόσχου ἀπὸ μάχης, καὶ ἐλαττώσεις ἁμαρτίας· ἄνθρωπος γὰρ θυμώδης ἐκκαύσει μάχην.

\vs{9}Καὶ ἀνὴρ ἁμαρτωλὸς ταράξει φίλους, καὶ ἀναμέσον εἰρηνευόντων ἐμβάλλει διαβολήν.
\vs{10}Κατὰ τὴν ὕλην πυρὸς οὕτως ἐκκαυθήσεται, κατὰ τὴν ἰσχὺν τοῦ ἀνθρώπου ὁ θυμὸς αὐτοῦ ἔσται, καὶ κατὰ τὸν πλοῦτον ἀνυψώσει ὀργὴν αὐτοῦ, καὶ κατὰ τὴν στερέωσιν τῆς μάχης ἐκκαυθήσεται.
\vs{11}Ἔρις κατασπευδομένη ἐκκαίει πῦρ, καὶ μάχη κατασπεύδουσα ἐκχέει αἷμα.
\vs{12}Ἐὰν φυσήσῃς σπινθῆρα ἐκκαήσεται, καὶ ἐὰν πτύσῃς ἐπʼ αὐτὸν σβεσθήσεται· καὶ ἀμφότερα ἐκ τοῦ στόματός σου ἐκπορεύεται.

\vs{13}Ψίθυρον καὶ δίγλωσσον καταρᾶσθαι, πολλοὺς γὰρ εἰρηνεύοντας ἀπώλεσαν.
\vs{14}Γλῶσσα τρίτη πολλοὺς ἐσάλευσε, καὶ διέστησεν αὐτοὺς ἀπὸ ἔθνους εἰς ἔθνος, καὶ πόλεις ὀχυρὰς καθεῖλε, καὶ οἰκίας μεγιστάνων κατέστρεψε.
\vs{15}Γλῶσσα τρίτη γυναῖκας ἀνδρείας ἐξέβαλε, καὶ ἐστέρησεν αὐτὰς τῶν πόνων αὐτῶν.
\vs{16}Ὁ προσέχων αὐτῇ οὐ μὴ εὕρῃ ἀνάπαυσιν, οὐδὲ κατασκηνώσει μεθʼ ἡσυχίας.

\vs{17}Πληγὴ μάστιγος ποιεῖ μώλωπας, πληγὴ δὲ γλώσσης συγκλάσει ὀστᾶ.
\vs{18}Πολλοὶ ἔπεσαν ἐν στόματι μαχαίρας, καὶ οὐχ ὡς οἱ πεπτωκότες διὰ γλῶσσαν.
\vs{19}Μακάριος ὁ σκεπασθεὶς ἀπʼ αὐτῆς, ὃς οὐ διῆλθεν ἐν τῷ θυμῷ αὐτῆς, ὃς οὐχ εἵλκυσε τὸν ζυγὸν αὐτῆς, καὶ ἐν τοῖς δεσμοῖς αὐτῆς οὐκ ἐδέθη.
\vs{20}Ὁ γὰρ ζυγὸς αὐτῆς ζυγὸς σιδηροῦς, καὶ οἱ δεσμοὶ αὐτῆς δεσμοὶ χάλκεοι.
\vs{21}Θάνατος πονηρὸς ὁ θάνατος αὐτῆς, καὶ λυσιτελὴς μᾶλλον ὁ ᾅδης αὐτῆς.
\vs{22}Οὐ μὴ κρατήσῃ εὐσεβῶν, καὶ ἐν τῇ φλογὶ αὐτῆς οὐ καήσονται.
\vs{23}Οἱ καταλείποντες Κύριον ἐμπεσοῦνται εἰς αὐτὴν, καὶ ἐν αὐτοῖς ἐκκαήσεται, καὶ οὐ μὴ σβεσθῇ· ἐξαποσταλήσεται ἐπʼ αὐτοῖς ὡς λέων, καὶ ὡς πάρδαλις λυμανεῖται αὐτούς.

\vs{24}Ἴδε περίφραξον τὸ κτῆμά σου ἀκάνθαις, τὸ ἀργύριόν σου καὶ τὸ χρυσίον κατάδησον·
\vs{25}καὶ τοῖς λόγοις σου ποίησον ζυγὸν καὶ σταθμὸν, καὶ τῷ στόματί σου ποίησον θύραν καὶ μοχλόν.
\vs{26}Πρόσεχε μήπως ὀλισθήσῃς ἐν αὐτῇ, μὴ πέσῃς κατέναντι ἐνεδρεύοντος.

\ch{29}
Ὁ ποιῶν ἔλεος δανειεῖ τῷ πλησιον, καὶ ὁ ἐπισχύων τῇ χειρὶ αὐτοῦ τηρεῖ ἐντολάς.
\vs{2}Δάνεισον τῷ πλησίον ἐν καιρῷ χρείας αὐτοῦ, καὶ πάλιν ἀπόδος τῷ πλησίον εἰς τὸν καιρόν.
\vs{3}Στερέωσον λόγον, καὶ πιστώθητι μετʼ αὐτοῦ, καὶ ἐν παντὶ καιρῷ εὑρήσεις τὴν χρείαν σου.
\vs{4}Πολλοὶ ὡς εὕρεμα ἐνόμισαν δάνος, καὶ παρέσχον πόνον τοῖς βοηθήσασιν αὐτοῖς.
\vs{5}Ἕως οὗ λάβῃ, καταφιλήσει χεῖρα αὐτοῦ, καὶ ἐπὶ τῶν χρημάτων τοῦ πλησίον ταπεινώσει φωνήν· καὶ ἐν καιρῷ ἀποδόσεως παρελκύσει χρόνον, καὶ ἀποδώσει λόγους ἀκηδίας, καὶ τὸν καιρὸν αἰτιάσεται.
\vs{6}Ἐὰν ἰσχύσῃ, μόλις κομίσεται τὸ ἥμισυ, καὶ λογιεῖται αὐτὸ ὡς εὕρεμα· εἰ δὲ μὴ, ἀπεστέρησεν αὐτὸν τῶν χρημάτων αὐτοῦ, καὶ ἐκτήσατο αὐτὸν ἐχθρὸν δωρεάν· κατάρας καὶ λοιδορίας ἀποδώσει αὐτῷ, καὶ ἀντὶ δόξης ἀποδώσει αὐτῷ ἀτιμίαν.
\vs{7}Πολλοὶ χάριν πονηρίας ἀπέστρεψαν, ἀποστερηθῆναι δωρεὰν εὐλαβήθησαν.
\vs{8}Πλὴν ἐπὶ ταπεινῷ μακροθύμησον, καὶ ἐπʼ ἐλεημοσύνην μὴ παρελκύσῃς αὐτόν·
\vs{9}Χάριν ἐντολῆς ἀντιλαβοῦ πένητος, καὶ κατὰ τὴν ἔνδειαν αὐτοῦ μὴ ἀποστρέψῃς αὐτὸν κενόν.

\vs{10}Ἀπόλεσον ἀργύριον διʼ ἀδελφὸν καὶ φίλον, καὶ μὴ ἰωθήτω ὑπὸ τὸν λίθον εἰς ἀπώλειαν.
\vs{11}Θὲς τὸν θησαυρόν σου κατʼ ἐντολὰς ὑψίστου, καὶ λυσιτελήσει σοι μᾶλλον ἢ τὸ χρυσίον.
\vs{12}Σύγκλεισον ἐλεημοσύνην ἐν τοῖς ταμείοις σου, καὶ αὕτη ἐξελεῖταί σε ἐκ πάσης κακώσεως.
\vs{13}Ὑπὲρ ἀσπίδα κράτους, καὶ ὑπὲρ δόρυ ἀλκῆς κατέναντι ἐχθροῦ πολεμήσει ὑπὲρ σοῦ.
\vs{14}Ἀνὴρ ἀγαθὸς ἐγγυήσεται τὸν πλησίον, καὶ ὁ ἀπολωλεκὼς αἰσχύνην καταλείψει αὐτόν.
\vs{15}Χάριτας ἐγγύου μὴ ἐπιλάθῃ, ἔδωκε γὰρ τὴν ψυχὴν αὐτοῦ ὑπὲρ σοῦ.
\vs{16}Ἀγαθὰ ἐγγύου ἀνατρέψει ἁμαρτωλὸς,
\vs{17}καὶ ἀχάριστος ἐν διανοίᾳ ἐγκαταλείψει ῥυσάμενον.

\vs{18}Ἐγγύη πολλοὺς ἀπώλεσε κατευθύνοντας, καὶ ἐσάλευσεν αὐτοὺς ὡς κῦμα θαλάσσης· ἄνδρας δυνατοὺς ἀπῴκισε, καὶ ἐπλανήθησαν ἐν ἔθνεσιν ἀλλοτρίοις.
\vs{19}Ἁμαρτωλὸς ἐμπεσὼν εἰς ἐγγύην, καὶ διώκων ἐργολαβείας ἐμπεσεῖται εἰς κρίσεις.
\vs{20}Ἀντιλαβοῦ τοῦ πλησίον κατὰ δύναμίν σου, καὶ πρόσεχε σεαυτῷ μὴ ἐμπέσῃς.

\vs{21}Ἀρχὴ ζωῆς ὕδωρ, καὶ ἄρτος, καὶ ἱμάτιον, καὶ οἶκος καλύπτων ἀσχημοσύνην.
\vs{22}Κρείσσων βίος πτωχοῦ ὑπὸ σκέπην δοκῶν, ἢ ἐδέσματα λαμπρὰ ἐν ἀλλοτρίοις.
\vs{23}Ἐπὶ μικρῷ καὶ μεγάλῳ εὐδοκίαν ἔχε.
\vs{24}Ζωὴ πονηρὰ ἐξ οἰκίας εἰς οἰκίαν, καὶ οὗ παροικήσει, οὐκ ἀνοίξει στόμα.
\vs{25}Ξενιεῖς καὶ ποτιεῖς εἰς ἀχάριστα, καὶ πρὸς ἐπὶ τούτοις πικρὰ ἀκούσῃ·
\vs{26}πάρελθε πάροικε, κόσμησον τράπεζαν, καὶ εἴτι ἐν τῇ χειρί σου ψώμισόν με·
\vs{27}ἔξελθε πάροικε ἀπὸ προσώπου δόξης, ἐπεξένωταί μοι ὁ ἀδελφὸς, χρεία τῆς οἰκίας.
\vs{28}Βαρέα ταῦτα ἀνθρώπῳ ἔχοντι φρόνησιν, ἐπιτίμησις οἰκίας καὶ ὀνειδισμὸς δανειστοῦ.

ΠΕΡΙ ΤΕΚΝΩΝ.

\ch{30}
Ὁ ἀγαπῶν τὸν υἱὸν αὐτοῦ, ἐνδελεχήσει μάστιγας αὐτῷ, ἵνα εὐφρανθῇ ἐπʼ ἐσχάτῳ αὐτοῦ.
\vs{2}Ὁ παιδεύων τὸν υἱὸν αὐτοῦ ὀνήσεται ἐπʼ αὐτῷ, καὶ ἀναμέσον γνωρίμων ἐπʼ αὐτῷ καυχήσεται.
\vs{3}Ὁ διδάσκων τὸν υἱὸν αὐτοῦ παραζηλώσει τὸν ἐχθρὸν, καὶ ἔναντι φίλων ἐπʼ αὐτῷ ἀγαλλιάσεται.
\vs{4}Ἐτελεύτησεν αὐτοῦ ὁ πατὴρ, καὶ ὡς οὐκ ἀπέθανεν, ὅμοιον γὰρ αὐτῷ κατέλιπε μετʼ αὐτόν.
\vs{5}Ἐν τῇ ζωῇ αὐτοῦ εἶδε καὶ εὐφράνθη, καὶ ἐν τῇ τελευτῇ αὐτοῦ οὐκ ἐλυπήθη.
\vs{6}Ἐναντίον ἐχθρῶν κατέλιπεν ἔκδικον, καὶ τοῖς φίλοις ἀνταποδιδόντα χάριν.

\vs{7}Περιψύχων υἱὸν καταδεσμεύσει τραύματα αὐτοῦ, καὶ ἐπὶ πάσῃ βοῇ ταραχθήσεται σπλάγχνα αὐτοῦ.
\vs{8}Ἵππος ἀδάμαστος ἀποβαίνει σκληρὸς, καὶ υἱὸς ἀνειμένος ἐκβαίνει προαλής.
\vs{9}Τιθήνησον τέκνον καὶ ἐκθαμβήσει σε, σύμπαιζον αὐτῷ καὶ λυπήσει σε.
\vs{10}Μὴ συγγελάσῃς αὐτῷ ἵνα μὴ συνοδυνηθῇς, καὶ ἐπʼ ἐσχάτῳ γομφιάσεις τοὺς ὀδόντας σου.
\vs{11}Μὴ δῷς αὐτῷ ἐξουσίαν ἐν νεότητι.
\vs{12}Θλάσον τὰς πλευρὰς αὐτοῦ, ὡς ἔστι νήπιος, μήποτε σκληρυνθεὶς ἀπειθήσῃ σοι.
\vs{13}Παίδευσον τὸν υἱόν σου, καὶ ἔργασαι ἐναὐτῷ, ἵνα μὴ ἐν τῇ ἀσχημοσύνῃ σου προσκόψῃ.
\vs{13a}Καὶ μὴ παρίδῃς τὰς ἀγνοίας αὐτοῦ.
\vs{13b}Κάμψον τὸν τράχηλον αὐτοῦ ἐν νεότητι.

ΠΕΡΙ ʼΥΓΙΕΙΑΣ.

\vs{14}Κρείσσων πτωχὸς ὑγιὴς καὶ ἰσχύων τῇ ἕξει, ἢ πλούσιος μεμαστιγωμένος εἰς σῶμα αὐτοῦ.
\vs{15}Ὑγιεία καὶ εὐεξία βέλτιον παντὸς χρυσίου, καὶ σῶμα εὔρωστον ἢ ὄλβος ἀμέτρητος.
\vs{16}Οὐκ ἔστι πλοῦτος βελτίων ὑγιείας σώματος, καὶ οὐκ ἔστιν εὐφροσύνη ὑπὲρ χαρὰν καρδίας.
\vs{17}Κρείσσων θάνατος ὑπὲρ ζωὴν πικρὰν, ἢ ἀῤῥώστημα ἔμμονον.
\vs{18}Ἀγαθὰ ἐκκεχυμένα ἐπὶ στόματι κεκλεισμένῳ, θέματα βρωμάτων παρακείμενα ἐπὶ τάφῳ.

\vs{19}Τί συμφέρει κάρπωσις εἰδώλῳ; οὔτε γὰρ ἔδεται οὔτε μὴ ὀσφρανθῇ· οὕτως ὁ ἐκδιωκόμενος ὑπὸ Κυρίου.
\vs{20}Βλέπων ἐν ἐν ὀφθαλμοῖς καὶ στενάζων, ὥσπερ εὐνοῦχος περιλαμβάνων παρθένον καὶ στενάζων.
\vs{21}Μὴ δῷς εἰς λύπην τὴν ψυχήν σου, καὶ μὴ θλίψῃς σεαυτὸν ἐν βουλῇ σου.
\vs{22}Εὐφροσύνη καρδίας ζωὴ ἀνθρώπου, καὶ ἀγαλλίαμα ἀνδρὸς μακροημέρευσις.
\vs{23}Ἀγάπα τὴν ψυχήν σου, καὶ παρακάλει τὴν καρδίαν σου, καὶ λύπην μακρὰν ἀπόστησον ἀπὸ σοῦ· πολλοὺς γὰρ ἀπέκτεινεν ἡ λύπη, καὶ οὐκ ἔστιν ὠφέλεια ἐν αὐτῇ.
\vs{24}Ζῆλος καὶ θυμὸς ἐλαττοῦσιν ἡμέρας, καὶ πρὸ καιροῦ γῆρας ἄγει μέριμνα.

\vs{25}Ὡς καλαμώμενος ὀπίσω τρυγητῶν, ἐν εὐλογίᾳ Κυρίου ἔφθασα,
\vs{26}καὶ ὡς τρυγῶν ἐπλήρωσα ληνόν. Κατανοήσατε ὅτι οὐκ ἐμοὶ μόνῳ ἐκοπίασα, ἀλλὰ πᾶσι τοῖς ζητοῦσι παιδείαν.
\vs{27}Ακούσατέ μου μεγιστᾶνες λαοῦ, καὶ οἱ ἡγούμενοι ἐκκλησίας ἐνωτίσασθε.
\vs{28}Ὑιῷ καὶ γυναικὶ, ἀδελφῷ καὶ φίλῳ μὴ δῷς ἐξουσίαν ἐπὶ σὲ ἐν ζωῇ σου, καὶ μὴ δῷς ἑτέρῳ τὰ χρήματά σου, ἵνα μὴ μεταμεληθεὶς δέῃ περὶ αὐτῶν.
\vs{29}Ἕως ἔτι ζῇς καὶ πνοὴ ἐν σοὶ, μὴ ἀλλάξῃς σεαυτὸν πάσῃ σαρκί·
\vs{30}κρεῖσσον γάρ ἐστι τὰ τέκνα δεηθῆναί σου, ἢ σὲ ἐμβλέπειν εἰς χεῖρας υἱῶν σου.
\vs{31}Ἐν πᾶσι τοῖς ἔργοις σου γίνου ὑπεράγων· μὴ δῷς μῶμον ἐν τῇ δόξῃ σου.
\vs{32}Ἐν ἡμέρᾳ συντελείας ἡμερῶν ζωῆς σου καὶ ἐν καιρῷ τελευτῆς διάδος κληρονομίαν.

ΠΕΡΙ ΔΟΥΛΩΝ.

\vs{33}Χορτάσματα καί ῥάβδος καὶ φορτία ὄνῳ, ἄρτος καὶ παιδεία καὶ ἔργον οἱκέτῃ.
\vs{34}Ἔργασαι ἐν παιδὶ καὶ εὑρήσεις ἀνάπαυσιν, ἄνες χεῖρας αὐτῷ καὶ ζητήσει ἐλευθερίαν.
\vs{35}Ζυγὸς καὶ ἱμὰς κάμψουσι τράχηλον, καὶ οἰκέτῃ κακούργῳ στρέβλαι καὶ βάσανοι.
\vs{36}Ἔμβαλε αὐτὸν εἰς ἐργασίαν, ἵνα μὴ ἀργῇ, πολλὴν γὰρ κακίαν ἐδίδαξεν ἡ ἀργία.
\vs{37}Εἰς ἔργα κατάστησον καθὼς πρέπει αὐτῷ, κᾂν μὴ πειθαρχῇ, βάρυνον τὰς πέδας αὐτοῦ.

\vs{38}Καὶ μὴ περισσεύσῃς ἐν πάσῃ σαρκὶ, καὶ ἄνευ κρίσεως μὴ ποιήσῃς μηδέν.
\vs{39}Εἰ ἔστι σοι οἰκέτης, ἔστω ὡς σὺ, ὅτι ἐν αἵματι ἐκτήσω αὐτόν.
\vs{40}Εἰ ἔστι σοι οἰκέτῃς, ἄγε αὐτὸν ὡς σεαυτὸν, ὅτι ὡς ἡ ψυχή σου ἐπιδεήσεις αὐτοῦ· ἐὰν κακώσῃς αὐτὸν, καὶ ἀπάρας ἀποδρᾷ, ἐν ποίᾳ ὁδῷ ζητήσεις αὐτόν;

\ch{31}
Κέναι ἐλπίδες καὶ ψευδεῖς ἀσυνέτῳ ἀνδρὶ, καὶ ἐνύπνια ἀναπτεροῦσιν ἄφρονας.
\vs{2}Ὡς δρασσόμενος σκιᾶς καὶ διώκων ἄνεμον, οὕτως ὁ ἐπέχων ἐνυπνίοις.
\vs{3}Τοῦτο κατὰ τούτου ὅρασις ἐνυπνίων, κατέναντι πρωσώπου ὁμοίωμα προσώπου.
\vs{4}Ἀπὸ ἀκαθάρτου τί καθαρισθήσεται; καὶ ἀπὸ ψευδοῦς τί ἀληθεύσει;
\vs{5}Μαντεῖαι καὶ οἰωνισμοὶ καὶ ἐνύπνια, μάταιά ἐστι, καὶ ὡς ὠδινούσης φαντάζεται καρδία.
\vs{6}Ἐὰν μὴ παρὰ ὑψίστου ἀποσταλῇ ἐν ἐπισκοπῇ, μὴ δῷς εἰς αὐτὰ τὴν καρδίαν σου.
\vs{7}Πολλοὺς ἐπλάνησε τὰ ἐνύπνια, καὶ ἐξέπεσον ἐλπίζοντες ἐπʼ αὐτοῖς.

\vs{8}Ἄνευ ψεύδους συντελεσθήσεται νόμος, καὶ σοφία στόματι πιστῷ τελείωσις.
\vs{9}Ἀνὴρ πεπαιδευμένος ἔγνω πολλὰ, καὶ ὁ πολύπειρος ἐκδιηγήσεται σύνεσιν.
\vs{10}Ὃς οὐκ ἐπειράθη ὀλίγα οἶδεν, ὁ δὲ πεπλανημένος πληθυνεῖ πανουργίαν.
\vs{11}Πολλὰ ἑώρακα ἐν τῇ ἀποπλανήσει μου, καὶ πλείονα τῶν λόγων μου, σύνεσίς μου.
\vs{12}Πλεονάκις ἕως θανάτου ἐκινδύνευσα, καὶ διεσώθην τούτων χάριν.

\vs{13}Πνεῦμα φοβουμένων Κυρὶον ζήσεται, ἡ γὰρ ἐλπὶς αὐτῶν ἐπὶ τὸν σώζοντα αὐτούς.
\vs{14}Ὁ φοβούμενος Κύριον οὐ μὴ εὐλαβηθήσεται, καὶ οὐ μὴ δειλιάσῃ, ὅτι αὐτὸς ἐλπὶς αὐτοῦ.
\vs{15}Φοβουμένου τὸν Κύριον μακαρία ἡ ψυχή· τίνι ἐπέχει, καὶ τίς ἀντιστήριγμα αὐτοῦ;
\vs{16}Οἱ ὀφθαλμοὶ Κυρίου ἐπὶ τοὺς ἀγαπῶντας αὐτὸν, ὑπερασπισμὸς δυναστείας καὶ στήριγμα ἰσχύος, σκέπη ἀπὸ καύσωνος καὶ σκέπη ἀπὸ μεσημβρίας, φυλακὴ ἀπὸ προσκόμματος καὶ βοήθεια ἀπὸ πτώματος, ἀνυψῶν ψυχὴν καὶ φωτίζων ὀφθαλμοὺς.
\vs{17}ἴασιν διδοὺς, ζωὴν καὶ εὐλογίαν.

\vs{18}Θυσιάζων ἐξ ἀδίκου, προσφορὰ μεμωκημένη, καὶ οὐκ εἰς εὐδοκίαν μωκήματα ἀνόμων.
\vs{19}Οὐκ εὐδοκεῖ ὁ ὕψιστος ἐν προσφοραῖς ἀσεβῶν, οὐδὲ ἐγ πλήθει θυσιῶν ἐξιλάσκεται ἁμαρτίας.
\vs{20}Θύων υἱὸν ἔναντι τοῦ πατρὸς αὐτοῦ, ὁ προσάγων θυσίαν ἐκ χρημάτων πενήτων.
\vs{21}Ἄρτος ἐπιδεομένων, ζωὴ πτωχῶν, ὁ ἀποστερῶν αὐτὴν ἄνθρωπος αἱμάτων.
\vs{22}Φὸνεύων τὸν πλησίον ὁ ἀφαιρούμενος συμβίωσιν, καὶ ἐκχέων αἷμα ὁ ἀποστερῶν μισθὸν μισθίου.

\vs{23}Εἷς οἰκοδομῶν, καὶ εἷς καθαιρῶν, τί ὠφέλησαν πλεῖον ἢ κόπους;
\vs{24}Εἷς εὐχόμενος, καὶ εἷς καταρώμενος, τίνος φωνῆς εἰσακούσεται ὁ δεστότης;
\vs{25}Βαπτιζόμενος ἀπὸ νεκροῦ, καὶ πάλιν ἁπτόμενος αὐτοῦ, τί ὠφέλησε τῷ λουτρῷ αὐτοῦ;
\vs{26}Οὕτως ἄνθρωπος νηστεύων ἐπὶ τῶν ἁμαρτιῶν αὐτοῦ, καὶ πάλιν πορευόμενος, καὶ τὰ αὐτὰ ποιῶν· τῆς προσευχῆς αὐτοῦ τίς εἰσακούσεται; καὶ τί ὠφέλησεν ἐν τῷ ταπεινωθῆναι αὐτόν;

\ch{32}
Ὁ συντήρων νόμον πλεονάζει προσφοράς· θυσιάζων σωτηρίου ὁ προσέχων ἐντολαῖς.
\vs{2}Ἀνταποδιδοὺς χάριν προσφέρων σεμίδαλιν, καὶ ὁ ποιῶν ἐλεημοσύνην θυσιάζων αἰνέσεως.
\vs{3}Εὐδοκία Κυρίου ἀποστῆναι ἀπὸ πονηρίας, καὶ ἐξιλασμὸς ἀποστῆναι ἀπὸ ἀδικίας.

\vs{4}Μὴ ὀφθῇς ἐν προσώπῳ Κυρίου κενὸς,
\vs{5}πάντα γὰρ ταῦτα χάριν ἐντολῆς.
\vs{6}Προσφορὰ δικαίου λιπαίνει θυσιαστήριον, καὶ ἡ εὐωδία αὐτῆς ἔναντι ὑψίστου.
\vs{7}Θυσία ἀνδρὸς δικαίου δεκτὴ, καὶ τὸ μνημόσυνον αὐτῆς οὐκ ἐπιλησθήσεται.

\vs{8}Ἐν ἀγαθῷ ὀφθαλμῷ δόξασον τὸν Κύριον, καὶ μὴ σμικρύνῃς ἀπαρχὴν χειρῶν σου.
\vs{9}Ἐν πάσῃ δόσει ἱλάρωσον τὸ πρόσωπόν σου, καὶ ἐν εὐφροσύνῃ ἁγίασον δεκάτην.
\vs{10}Δὸς ὑψίστῳ κατὰ τὴν δόσιν αὐτοῦ, καὶ ἐν ἀγαθῷ ὀφθαλμῷ καθεύρεμα χειρός.
\vs{11}Ὅτι Κύριος ἀνταποδιδούς ἐστι, καὶ ἑπταπλάσια ἀνταποδώσει σοι.
\vs{12}Μὴ δωροκόπει, οὐ γὰρ προσδέξεται· καὶ μὴ ἔπεχε θυσίᾳ ἀδίκῳ, ὅτι Κύριος κριτής ἐστι, καὶ οὐκ ἔστι παρʼ αὐτῷ δόξα προσώπου.
\vs{13}Οὐ λήψεται πρόσωπον ἐπὶ πτωχοῦ, καὶ δέησιν ἠδικημένου εἰσακούσεται.
\vs{14}Οὐ μὴ ὑπερίδῃ ἱκετείαν ὀρφανοῦ, καὶ χήραν ἐὰν ἐκχέῃ λαλιάν.
\vs{15}Οὐχὶ δάκρυα χήρας ἐπὶ σιαγόνα καταβαίνει, καὶ ἡ καταβόησις ἐπὶ τῷ καταγαγόντι αὐτά;

\vs{16}Θεραπεύων ἐν εὐδοκίᾳ δεχθήσεται, καὶ ἡ δέησις αὐτοῦ ἕως νεφελῶν συνάψει.
\vs{17}Προσευχὴ ταπεινοῦ νεφέλας διῆλθε, καὶ ἕως συνεγγίσῃ οὐ μὴ παρακληθῇ, καὶ οὐ μὴ ἀποστῇ ἕως ἐπισκέψηται ὁ ὕψιστος· καὶ κρινεῖ δικαίως, καὶ ποιήσει κρίσιν.
\vs{18}Καὶ ὁ Κύριος οὐ μὴ βραδύνῃ, οὐδὲ μὴ μακροθυμήσει ἐπʼ αὐτοῖς ἕως ἄν συντρίψῃ ὀσφὺν ἀνελεημόνων· καὶ τοῖς ἔθνεσιν ἀνταποδώσει ἐκδίκησιν, ἕως ἐξάρῃ πλῆθος ὑβριστῶν, καὶ σκῆπτρα ἀδίκων συντρίψῃ,
\vs{19}ἕως ἀνταποδῷ ἀνθρώπῳ κατὰ τὰς πράξεις αὐτοῦ, καὶ τὰ ἔργα τῶν ἀνθρώπων κατὰ τὰ ἐνθυμήματα αὐτῶν, ἕως κρινῇ τὴν κρίσιν τοῦ λαοῦ αὐτοῦ, καὶ εὐφρανεῖ αὐτοὺς ἐν τῷ ἐλέει αὐτοῦ.
\vs{20}Ὡραῖον ἔλεος ἐν καιρῷ θλίψεως αὐτοῦ, ὡς νεφέλαι ὑετοῦ ἐν καιρῷ ἀβροχίας.

\ch{33}
Ἐλέησον ἡμᾶς, δέσποτα ὁ Θεὸς πάντων, καὶ ἐπίβλεψον·
\vs{2}καὶ ἐπίβαλε τὸν φόβον σου ἐπὶ πάντα τὰ ἔθνη.
\vs{3}Ἔπαρον τὴν χεῖρά σου ἐπὶ ἔθνη ἀλλότρια, καὶ ἰδέτωσαν τὴν δυναστείαν σου.
\vs{4}Ὥσπερ ἐνώπιον αὐτῶν ἡγιάσθης ἐν ἡμῖν, οὕτως ἐνώπιον ἡμῶν μεγαλυνθείης ἐν αὐτοῖς·
\vs{5}Καὶ ἐπιγνώτωσάν σε καθάπερ καὶ ἡμεῖς ἐπέγνωμεν, ὅτι οὐκ ἔστι Θεὸς πλήν σου Κύριε.

\vs{6}Ἐγκαίνισον σημεῖα, καὶ ἀλλοίωσον θαυμάσια·
\vs{7}δόξασον χεῖρα καὶ βραχίονα δεξιόν· ἔγειρον θυμὸν, καὶ ἔκχεον ὀργήν· ἔξαρον ἀντίδικον, καὶ ἔκτριψον ἐχθρόν.
\vs{8}Σπεῦσον καιρὸν, καὶ μνήσθητι ὁρκισμοῦ, καὶ ἐκδιηγησάσθωσαν τὰ μεγαλεῖά σου.
\vs{9}Ἐν ὀργῇ πυρὸς καταβρωθήτω ὁ σωζόμενος, καὶ οἱ κακοῦντες τὸν λαόν σου εὕροισαν ἀπώλειαν.
\vs{10}Σύντριψον κεφαλὰς ἀρχόντων ἐχθρῶν λεγόντων, οὐκ ἔστι πλὴν ἡμῶν.
\vs{11}Σύναγε πάσας φυλὰς Ἰακώβ.

\vs{12}Λαμπρὰ καρδία καὶ ἀγαθὴ ἐπὶ ἐδέσμασιν τῶν βρωμάτων αὐτῆς ἐπιμελήσεται.

\ch{34}
Ἀγρυπνία πλούτου ἐκτήκει σάρκας, καὶ ἡ μέριμνα αὐτοῦ ἀφιστᾷ ὕπνον.
\vs{2}Μέριμνα ἀγρυπνίας ἀπαιτήσει νυσταγμὸν, καὶ ἀῤῥώστημα βαρὺ ἐκνήψει ὕπνος.
\vs{3}Ἐκοπίασε πλούσιος ἐν συναγωγῇ χρημάτων, καὶ ἐν τῇ ἀναπαύσει ἐμπίπλαται τῶν τρυφημάτων αὐτοῦ·
\vs{4}ἐκοπίασε πτωχὸς ἐν ἐλαττώσει βίου, καὶ ἐν τῇ ἀναπαύσει ἐπιδεὴς γίνεται.

\vs{5}Ὁ ἀγαπῶν χρυσίον οὐ δικαιωθήσεται, καὶ ὁ διώκων διαφθορὰν, αὐτὸς πλησθήσεται.
\vs{6}Πολλοὶ ἐδόθησαν εἰς πτῶμα χάριν χρυσίου, καὶ ἐγενήθη ἀπώλεια αὐτῶν κατὰ πρόσωπον αὐτῶν.
\vs{7}Ξύλον προσκόμματός ἐστι τοῖς ἐνθυσιάζουσιν αὐτῷ, καὶ πᾶς ἄφρων ἁλώσεται ἐν αὐτῷ.
\vs{8}Μακάριος πλούσιος ὃς εὑρέθη ἄμωμος, καὶ ὃς ὀπίσω χρυσίου οὐκ ἐπορεύθη.
\vs{9}Τίς ἐστι καὶ μακαριοῦμεν αὐτόν; ἐποίησε γὰρ θαυμάσια ἐν λαῷ αὐτοῦ.
\vs{10}Τίς ἐδοκιμάσθη ἐν αὐτῷ καὶ ἐτελειώθη, καὶ ἔστω εἰς καύχησιν; τίς ἐδύνατο παραβῆναι, καὶ οὐ παρέβη, καὶ ποιῆσαι κακὰ, καὶ οὐκ ἐποίησε;
\vs{11}Στερεωθήσεται τὰ ἀγαθὰ αὐτοῦ, καὶ τὰς ἐλεημοσύνας αὐτοῦ ἐκδιηγήσεται ἐκκλησία.

\vs{12}Ἐπὶ τραπέζης μεγάλης ἐκάθισας; μὴ ἀνοίξῃς ἐπʼ αὐτῆς φάρυγγά σου· καὶ μὴ εἴπῃς, πολλά γε τὰ ἐπʼ αὐτῆς.
\vs{13}Μνήσθητι ὅτι κακὸν ὀφθαλμὸς πονηρὸς, πονηρότερον ὀφθαλμοῦ τί ἔκτισται; διὰ τοῦτο ἀπὸ παντὸς προσώπου δακρύει.
\vs{14}Οὗ ἐὰν ἐπιβλέψῃ, μὴ ἐκτείνῃς χεῖρα, καὶ μὴ συνθλίβου αὐτῷ ἐν τρυβλίῳ.
\vs{15}Νόει τὰ τοῦ πλησίον ἐκ σεαυτοῦ, καὶ ἐπὶ παντὶ πράγματι διανοοῦ.
\vs{16}Φάγε ὡς ἄνθρωπος τὰ παρακείμενά σοι, καὶ μὴ διαμασῷ, μὴ μισηθῇς.
\vs{17}Παῦσαι πρῶτος χάριν παιδείας, καὶ μὴ ἀπληστεύου, μήποτε προσκόψῃς.
\vs{18}Καὶ εἰ ἀναμέσον πλειόνων ἐκάθισας, πρότερος αὐτῶν μὴ ἐκτείνῃς τὴν χεῖρά σου.
\vs{19}Ὡς ἱκανὸν ἀνθρώπῳ πεπαιδευμένῳ τὸ ὀλίγον, καὶ ἐπὶ τῆς κοίτης αὐτοῦ οὐκ ἀσθμαίνει.
\vs{20}Ὕπνος ὑγείας ἐπὶ ἐντέρῳ μετρίῳ, ἀνέστη πρωῒ, καὶ ἡ ψυχὴ αὐτοῦ μετʼ αὐτοῦ· πόνος ἀγρυπνίας καὶ χολέρας καὶ στρόφος μετὰ ἀνδρὸς ἀπλήστου.
\vs{21}Καὶ εἰ ἐβιάσθης ἐν ἐδέσμασιν, ἀνάστα μεσοπωρῶν καὶ ἀναπαύσῃ.

\vs{22}Ἄκουσόν μου τέκνον καὶ μὴ ἐξουδενώσῃς με, καὶ ἐπʼ ἐσχάτῳ εὑρήσεις τοὺς λόγους μου· ἐν πᾶσι τοῖς ἔργοις σου γίνου ἐντρεχὴς, καὶ πᾶν αρρώστημα οὐ μή σοι ἀπαντήσῃ.
\vs{23}Λαμπρὸν ἐπʼ ἄρτοῖς εὐλογήσει χείλη, καὶ μαρτυρία τῆς καλλονῆς αὐτοῦ πιστή.
\vs{24}Πονηρῷ ἐπʼ ἄρτῳ διαγογγύσει πόλις, καὶ ἡ μαρτυρία τῆς πονηρίας αὐτοῦ ἐκριβής.

\vs{25}Ἐν οἴνῳ μὴ ἀνδρίζου, πολλοὺς γὰρ ἀπώλεσεν ὁ οἶνος·
\vs{26}Κάμινος δοκιμάζει στόμωμα ἐν βαφῇ, οὕτως οἶνος καρδίας ἐν μάχῃ ὑπερηφάνων.
\vs{27}Ἐπίσον ζωῆς οἶνος ἀνθρώπῳ, ἐὰν πίνῃς αὐτὸν μέτρῳ αὐτοῦ· τίς ζωὴ ἐλασσουμένῳ οἴνῳ; καὶ αὐτὸς ἔκτισται εἰς εὐφροσύνην ἀνθρώποις.
\vs{28}Ἀγαλλίαμα καρδίας καὶ εὐφροσύνη ψυχῆς οἶνος πινόμενος ἐν καιρῷ αὐτάρκης·
\vs{29}πικρία ψυχῆς οἶνος πινόμενος πολὺς, ἐν ὀρεθισμῷ καὶ ἀντιπτώματι.
\vs{30}Πληθύνει μέθη θυμὸν ἄφρονος εἰς πρόσκομμα, ἐλαττῶν ἰσχὺν καὶ προσποιῶν τραύματα.
\vs{31}Ἐν συμποσίῳ οἴνου μὴ ἐλέγξῃς τὸν πλησίον, καὶ μὴ ἐξουδενώσῃς αὐτὸν ἐν εὐφροσύνῃ αὐτοῦ· λόγον ὀνειδισμοῦ μὴ εἴπῃς αὐτῷ, καὶ μὴ αὐτὸν θλίψῃς ἐν ἀπαιτήσει.

ΠΕΡΙ ἩΓΟΥΜΕΝΩΝ.

\ch{35}
Ἡγούμενων σε κατέστησαν; μὴ ἐπαίρου, γίνου ἐν αὐτοῖς ὡς εἷς ἐξ αὐτῶν· φρόντισον αὐτῶν, καὶ οὕτω κάθισον,
\vs{2}καὶ πᾶσαν τὴν χρείαν σου ποιήσας ἀνάπεσε, ἵνα εὐφρανθῇς διʼ αὐτοὺς, καὶ εὐκοσμίας χάριν λάβῃς στέφανον.
\vs{3}Λάλησον πρεσβύτερε, πρέπει γάρ σοι, ἐν ἀκριβεῖ ἐπιστήμῃ, καὶ μὴ ἐμποδίσῃς μουσικά.
\vs{4}Ὅπου ἀκρόαμα, μὴ ἐκχέῃς λαλιὰν, καὶ ἀκαίρως μὴ σοφίζου.
\vs{5}Σφραγὶς ἄνθρακος ἐπὶ κόσμῳ χρυσῷ, σύγκριμα μουσικῶν ἐν συμποσίῳ οἴνου.
\vs{6}Ἐν κατασκευάσματι χρυσῷ σφραγὶς σμαράγδου μέλος μουσικῶν ἐφʼ ἡδεῖ οἴνῳ.

\vs{7}Λάλησον νεανίσκε εἰ χρεία σου, μόλις δὶς ἐὰν ἐπερωτηθῇς.
\vs{8}Κεφαλαίωσον λόγον, ἐν ὀλίγοις πολλὰ, γίνου ὡς γινώσκων καὶ ἅμα σιωπῶν.
\vs{9}Ἐν μέσῳ μεγιστάνων μὴ ἐξισάζου, καὶ ἑτέρου λέγοντος μὴ πολλὰ ἀδολέσχει.
\vs{10}Πρὸ βροντῆς κατασπεύδει ἀστραπὴ, καὶ πρὸ αἰσχυντηροῦ προελεύσεται χάρις.
\vs{11}Ἐν ὥρᾳ ἐξεγείρου καὶ μὴ οὐράγει, ἀπότρεχε εἰς οἶκον καὶ μὴ ῥαθύμει.
\vs{12}Ἐκεῖ παῖζε καὶ ποίει τὰ ἐνθυμήματά σου, καὶ μὴ ἁμάρτῃς λόγῳ ὑπερηφάνῳ.
\vs{13}Καὶ ἐπὶ τούτοις εὐλόγησον τὸν ποιήσαντά σε, καὶ μεθύσκοντά σε ἀπὸ τῶν ἀγαθῶν αὐτοῦ.

\vs{14}Ὁ φοβούμενος Κύριον ἐκδέξεται παιδείαν, καὶ οἱ ὀρθρίζοντες εὑρήσουσιν εὐδοκίαν.
\vs{15}Ὁ ζητῶν νόμον ἐμπλησθήσεται αὐτοῦ, καὶ ὁ ὑποκρινόμενος σκανδαλισθήσεται ἐν αὐτῷ.
\vs{16}Οἱ φοβούμενοι Κύριον εὑρήσουσι κρίμα, καὶ δικαιώματα ὡς φῶς ἐξάψουσιν.
\vs{17}Ἄνθρωπος ἁμαρτωλὸς ἐκκλίνει ἐλεγμὸν, καὶ κατὰ τὸ θέλημα αὐτοῦ εὑρήσει σύγκριμα.
\vs{18}Ἀνὴρ βουλῆς οὐ μὴ παρίδῃ διανόημα, ἀλλότριος καὶ ὑπερήφανος οὐ καταπτήξει φόβον, καὶ μετὰ τὸ ποιῆσαι μετʼ αὐτοῦ ἄνευ βουλῆς.
\vs{19}Ἄνευ βουλῆς μηθὲν ποιήσῃς, καὶ ἐν τῷ ποιῆσαί σε μὴ μεταμελοῦ.

\vs{20}Ἐν ὁδῷ ἀντιπτώματος μὴ πορεύου, καὶ μὴ προσκόψῃς ἐν λιθώδεσι.
\vs{21}Μὴ πιστεύσῃς ἐν ὁδῷ ἀπροσκόπῳ,
\vs{22}καὶ ἀπὸ τῶν τέκνων σου φύλαξαι.
\vs{23}Ἐν παντὶ ἔργῳ πίστευε τῇ ψυχῇ σου, καὶ γὰρ τοῦτό ἐστι τήρησις ἐντολῶν.
\vs{24}Ὁ πιστεύων νόμῳ προσέχει ἐντολαῖς, καὶ ὁ πεποιθὼς Κυρίῳ οὐκ ἐλαττωθήσεται.

\ch{36}
Τῷ φοβουμένῳ Κύριον οὐκ ἀπαντήσει κακὸν, ἀλλʼ ἐν πειρασμῷ καὶ πάλιν ἐξελεῖται.
\vs{2}Ἀνὴρ σοφὸς οὐ μισήσει νόμον, ὁ δὲ ὑποκρινόμενος ἐν αὐτῷ, ὡς ἐν καταιγίδι πλοῖον.
\vs{3}Ἄνθρωπος συνετὸς ἐμπιστεύσει νόμῳ, καὶ ὁ νόμος αὐτῷ πιστὸς ὡς ἐρώτημα δικαίων.
\vs{4}Ἑτοίμασον λόγον, καὶ οὕτως ἀκουθήσῃ, σύνδησον παιδείαν καὶ ἀποκρίθητι.
\vs{5}Τροχὸς ἁμάξης σπλάγχνα μωροῦ, καὶ ὡς ἄξων στρεφόμενος ὁ διαλογισμὸς αὐτοῦ.
\vs{6}Ἵππος εἰς ὀχείαν ὡς φίλος μῶκος, ὑποκάτω παντὸς ἐπικαθημένου χρεμετίζει.

\vs{7}Διατί ἡμέρα ἡμέρας ὑπερέχει, καὶ πᾶν φῶς ἡμέρας ἐνιαυτοῦ ἀφʼ ἡλίου;
\vs{8}Ἐν γνώσει Κυρίου διεχωρίσθησαν, καὶ ἠλλοίωσε καιροὺς καὶ ἑορτάς·
\vs{9}ἀπʼ αὐτῶν ἀνύψωσε καὶ ἡγίασε, καὶ ἐξ αὐτῶν ἔθηκεν εἰς ἀριθμὸν ἡμερῶν.

\vs{10}Καὶ ἄνθρωποι πάντες ἀπὸ ἐδάφους, καὶ ἐκ γῆς ἐκτίσθη Ἀδάμ.
\vs{11}Ἐν πλήθει ἐπιστήμης Κύριος διεχώρισεν αὐτοὺς, καὶ ἠλλοίωσε τὰς ὁδοὺς αὐτῶν.
\vs{12}Ἐξ αὐτῶν εὐλόγησε καὶ ἀνύψωσε, καὶ ἐξ αὐτῶν ἡγίασε, καὶ πρὸς αὐτὸν ἤγγισεν· ἀπʼ αὐτῶν κατηράσατο, καὶ ἐταπείνωσε, καὶ ἀνέστρεψεν αὐτοὺς ἀπὸ στάσεως αὐτῶν.
\vs{13}Ὡς πηλὸς κεραμέως ἐν χειρὶ αὐτοῦ, πᾶσαι αἱ ὁδοὶ αὐτοῦ κατὰ τὴν εὐδοκίαν αὐτοῦ· οὕτως ἄνθρωποι ἐν χειρὶ τοῦ ποιήσαντος αὐτοὺς, ἀποδοῦναι αὐτοῖς κατὰ τὴν κρίσιν αὐτοῦ.
\vs{14}Ἀπέναντι τοῦ κακοῦ τὸ ἀγαθὸν, καὶ ἀπέναντι τοῦ θανάτου ἡ ζωὴ, οὕτως ἀπέναντι εὐσεβοῦς ἀμαρτωλός.
\vs{15}Καὶ οὕτως ἔμβλεψον εἰς πάντα τὰ ἔργα τοῦ ὑψίστου, δύο δύο ἓν κατέναντι τοῦ ἑνός.

\vs{16}Κᾀγὼ ἔσχατος ἠγρύπνησα, καὶ κατεκληρονόμησα αὐτοὺς καθὼς ἀπʼ ἀρχῆς.
\vs{17}Ἐλέησον λαὸν, Κύριε, κεκλημένον ἐπʼ ὀνόματί σου, καὶ Ἰσραὴλ ὃν πρωτογόνῳ ὡμοίωσας.
\vs{18}Οἰκτείρησον πόλιν ἁγιάσματός σου Ἱερουσαλὴμ, πόλιν καταπαύματός σου.
\vs{19}Πλῆσον Σιὼν ἆραι τὰ λόγιά σου, καὶ ἀπὸ τῆς δόξης σου τὸν λαόν σου.
\vs{20}Δὸς μαρτύριον τοῖς ἐν ἀρχῇ κτίσμασί σου, καὶ ἔγειρον προφητείας τὰς ἐπʼ ὀνόματί σου·
\vs{21}δὸς μισθὸν τοῖς ὑπομένουσί σε, καὶ οἱ προφῆταί σου ἐμπιστευθήτωσαν.

\vs{22}Εἰσάκουσον, Κύριε, δεήσεως τῶν ἱκετῶν σου, κατὰ τὴν εὐλογίαν Ἀαρὼν περὶ τοῦ λαοῦ σου, καὶ γνώσονται πάντες οἱ ἐπὶ τῆς γῆς, ὅτι σὺ Κύριος εἶ ὁ Θεὸς τῶν αἰώνων.
\vs{23}Πᾶν βρῶμα φάγεται κοιλία, ἔστι δὲ βρῶμα βρώματος κάλλιον.
\vs{24}Φάρυγξ γεύεται βρώματα θήρας, οὕτως καρδία συνετὴ λόγους ψευδεῖς.
\vs{25}Καρδία στρεβλὴ δώσει λύπην, καὶ ἄνθρωπος πολύπειρος ἀνταποδώσει αὐτῷ.

\vs{26}Πάντα ἄῤῥενα ἐπιδέξεται γυνὴ, ἔστι δὲ θυγάτηρ θυγατρὸς κρείσσων.
\vs{27}Κάλλος γυναικὸς ἱλαρύνει πρόσωπον, καὶ ὑπὲρ πᾶσαν ἐπιθυμίαν ἀνθρώπου ὑπεράγει.
\vs{28}Εἰ ἔστιν ἐπὶ γλώσσης αὐτῆς ἔλεος καὶ πραΰτης, οὐκ ἔστιν ὁ ἀνὴρ αὐτῆς καθʼ υἱοὺς ἀνθρώπων.

\vs{29}Ὁ κτώμενος γυναῖκα ἐνάρχεται κτήσεως, βοηθὸν κατʼ αὐτὸν καὶ στύλον ἀναπαύσεων.
\vs{30}Οὗ οὐκ ἐστι φραγμὸς, διαρπαγήσεται κτῆμα, καὶ οὗ οὐκ ἔστι γυνὴ, στενάξει πλανώμενος.
\vs{31}Τίς γὰρ πιστεύσει εὐζώνῳ λῃστῇ σφαλλομένῳ ἐκ πόλεως εἰς πόλιν; οὕτως ἀνθρώπῳ μὴ ἔχοντι νοσσιὰν καὶ καταλύοντι οὗ ἐὰν ὀψισῇ.

\ch{37}
Πὰς φίλος ἐρεῖ, ἐφιλίασα αὐτῷ κᾷγώ· ἀλλʼ ἔστι φίλος ὀνόματι μόνον φίλος.
\vs{2}Οὐχὶ λύπη ἔνι ἕως θανάτου, ἑταῖρος καὶ φίλος τρεπόμενος εἰς ἔχθραν;
\vs{3}Ὦ πονηρὸν ἐνθύμημα, πόθεν ἐνεκυλίσθης καλύψαι τὴν ξηρὰν ἐν δολιότητι;
\vs{4}Ἑταῖρος φίλου ἐν εὐφροσύνῃ ἥδεται, καὶ ἐν καιρῷ θλίψεως ἔσται ἀπέναντι·
\vs{5}Ἐταῖρος φίλῳ συμπονεῖ χάριν γαστρὸς, ἔναντι πολέμου λήψεται ἀσπίδα.
\vs{6}Μὴ ἐπιλάθῃ φίλου ἐν τῇ ψυχῇ σου, καὶ μὴ ἀμνημονήσῃς αὐτοῦ ἐν χρήμασί σου.

\vs{7}Πᾶς σύμβουλος ἐξαίρει βουλὴν, ἀλλʼ ἔστι συμβουλεύων εἰς ἑαυτόν.
\vs{8}Ἀπὸ συμβούλου φύλαξον τὴν ψυχήν σου, καὶ γνῶθι πρότερον τίς αὐτοῦ χρεία· καὶ γὰρ αὐτὸς ἑαυτῷ βουλεύσεται· μήποτε βάλῃ ἐπὶ σοὶ κλῆρον,
\vs{9}καὶ εἴπῃ σοι, καλὴ ἡ ὁδός σου· καὶ στήσεται ἐξ ἐναντίας ἰδεῖν τὸ συμβησόμενόν σοι.
\vs{10}Μὴ βουλεύου μετὰ τοῦ ὑποβλεπομένου σε, καὶ ἀπὸ τῶν ζηλούντων σε κρύψον βουλήν.
\vs{11}Μετὰ γυναικὸς περὶ ἀντιζήλου αὐτῆς, καὶ μετὰ δειλοῦ περὶ πολέμου· μετὰ ἐμπόρου περὶ μεταβολίας, καὶ μετὰ ἀγοράζοντος περὶ πράσεως· μετὰ βασκάνου περὶ εὐχαριστίας, καὶ μετὰ ἀνελεήμονος περὶ χρηστοηθείας· μετὰ ὀκνηροῦ περὶ παντὸς ἔργου, καὶ μετὰ μισθίου ἐφεστίου περὶ συντελείας, οἰκέτῃ ἀργῷ περὶ πολλῆς ἐργασίας· μὴ ἔπεχε ἐπὶ τούτοις περὶ πάσης συμβουλίας.

\vs{12}Ἀλλʼ ἢ μετὰ ἀνδρὸς εὐσεβοῦς ἐνδελέχιζε, ὃν ἂν ἐπιγνῷς συντηροῦντα ἐντολάς· ὃς ἐν τῇ ψυχῇ αὐτοῦ κατὰ τὴν ψυχήν σου, καὶ ἐὰν πταίσῃς, συναλγήσει σοι.
\vs{13}Καὶ βουλὴν καρδίας στῆσον, οὐ γάρ ἐστί σοι πιστότερος αὐτῆς·
\vs{14}ψυχὴ γὰρ ἀνδρὸς ἀπαγγέλλειν ἐνίοτε εἴωθεν, ἢ ἑπτὰ σκοποὶ ἐπὶ μετεώρου καθήμενοι ἐπὶ σκοπῆς.
\vs{15}Καὶ ἐπὶ πᾶσι τούτοις δεήθητι ὑψίστου, ἵνα εὐθύνῃ ἐν ἀληθείᾳ τὴν ὁδόν σου.
\vs{16}Ἀρχὴ παντὸς ἔργου λόγος, καὶ πρὸ πάσης πράξεως βουλή.

\vs{17}Ἴχνος ἀλλοιώσεως καρδίας.
\vs{18}Τέσσαρα μέρη ἀνατέλλει, ἀγαθὸν καὶ κακὸν, ζωὴ καὶ θάνατος, καὶ κυριεύουσα ἐνδελεχῶς αὐτῶν γλῶσσά ἐστιν.
\vs{19}Ἔστιν ἀνὴρ πανοῦργος καὶ πολλῶν παιδευτὴς, καὶ τῇ ἰδίᾳ ψυχῇ ἄχρηστός ἐστιν.
\vs{20}Ἔστι σοφιζόμενος ἐν λόγοις μισητὸς, οὗτος πάσης τροφῆς καθυστερήσει,
\vs{21}οὐ γὰρ ἐδόθη αὐτῷ παρὰ Κυρίου χάρις, ὅτι πάσης σοφίας ἐστερήθη.
\vs{22}Ἔστι σοφὸς τῇ ἰδίᾳ ψυχῇ, καὶ οἱ καρποὶ τῆς συνέσεως αὐτοῦ ἐπὶ στόματος πιστοί.

\vs{23}Ἀνὴρ σοφὸς τὸν ἑαυτοῦ λαὸν παιδεύσει, καὶ οἱ καρποὶ τῆς συνέσεως αὐτοῦ πιστοί.
\vs{24}Ἀνὴρ σοφὸς πλησθήσεται εὐλογίας, καὶ μακαριοῦσιν αὐτὸν πάντες οἱ ὁρῶντες.
\vs{25}Ζωὴ ἀνδρὸς ἐν ἀριθμῷ ἡμερῶν, καὶ αἱ ἡμέραι τοῦ Ἰσραὴλ ἀναρίθμητοι.
\vs{26}Ὁ σοφὸς ἐν τῷ λαῷ αὐτοῦ κληρονομήσει πίστιν, καὶ τὸ ὄνομα αὐτοῦ ζήσεται εἰς τὸν αἰῶνα.

\vs{27}Τέκνον, ἐν τῇ ζωῇ σου πείρασον τὴν ψυχήν σου, καὶ ἴδε τί πονηρὸν αὐτῇ, καὶ μὴ δῷς αὐτῇ.
\vs{28}Οὐ γὰρ πάντα πᾶσι συμφέρει, καὶ οὐ πᾶσα ψυχὴ ἐν παντὶ εὐδοκεῖ.
\vs{29}Μὴ ἀπληστεύου ἐν πάσῃ τρυφῇ, καὶ μὴ ἐκχυθῇς ἐπὶ ἐδεσμάτων·
\vs{30}ἐν πολλοῖς γὰρ βρώμασιν ἔσται πόνος, καὶ ἡ ἀπληστία ἐγγιεῖ ἕως χολέρας.
\vs{31}Διʼ ἀπληστίαν πολλοὶ ἐτελεύτησαν, ὁ δὲ προσέχων προσθήσει ζωήν.

\ch{38}
Τίμα ἰατρὸν πρὸς τὰς χρείας τιμαῖς αὐτοῦ, καὶ γὰρ αὐτὸν ἔκτισε Κύριος.
\vs{2}Παρὰ γὰρ ὑψίστου ἐστὶν ἴασις, καὶ παρὰ βασιλέως λήψεται δόμα.
\vs{3}Ἐπιστήμη ἰατροῦ ἀνυψώσει κεφαλὴν αὐτοῦ, καὶ ἔναντι μεγιστάνων θαυμασθήσεται.
\vs{4}Κύριος ἔκτισεν ἐκ γῆς φάρμακα, καὶ ἀνὴρ φρόνιμος οὐ προσοχθιεῖ αὐτοῖς.
\vs{5}Οὐκ ἀπὸ ξύλου ἐγλυκάνθη ὕδωρ, εἰς τὸ γνωσθῆναι τὴν ἰσχὺν αὐτοῦ;
\vs{6}Καὶ αὐτὸς ἔδωκεν ἀνθρώποις ἐπιστήμην ἐνδοξάζεσθαι ἐν τοῖς θαυμασίοις αὐτοῦ.
\vs{7}Ἐν αὐτοῖς ἐθεράπευσε, καὶ ᾖρε τὸν πόνον αὐτοῦ.
\vs{8}Μυρεψὸς ἐν τούτοις ποιήσει μίγμα, καὶ οὐ μὴ συντελέσῃ ἔργα αὐτοῦ, καὶ εἰρήνη παρʼ αὐτοῦ ἐστιν ἐπὶ προσώπου τῆς γῆς.

\vs{9}Τέκνον, ἐν ἀῤῥωστήματί σου μὴ παράβλεπε, ἀλλʼ εὖξαι Κυρίῳ, καὶ αὐτὸς ἰάσεταί σε.
\vs{10}Ἀπόστησον πλημμέλειαν, καὶ εὔθυνον χεῖρας, καὶ ἀπὸ πάσης ἁμαρτίας καθάρισον καρδίαν.
\vs{11}Δὸς εὐωδίαν, καὶ μνημόσυνον σεμιδάλεως, καὶ λίπανον προσφορὰν, ὡς μὴ ὑπάρχων.
\vs{12}Καὶ ἰατρῷ δὸς τόπον, καὶ γὰρ αὐτὸν ἔκτισε Κύριος· καὶ μὴ ἀποστήτω σου, καὶ γὰρ αὐτοῦ χρεία.
\vs{13}Ἔστι καιρὸς ὅτε καὶ ἐν χερσὶν αὐτῶν εὐωδία.
\vs{14}Καὶ γὰρ αὐτοὶ Κυρίου δεηθήσονται, ἵνα εὐοδώσῃ αὐτοῖς ἀνάπαυσιν καὶ ἴασιν χάριν ἐμβιώσεως.
\vs{15}Ὁ ἁμαρτάνων ἔναντι τοῦ ποιήσαντος αὐτὸν, ἐμπέσοι εἰς χεῖρας ἰατροῦ.

\vs{16}Τέκνον, ἐπὶ νεκρῷ κατάγαγε δάκρυα, καὶ ὡς δεινὰ πάσχων ἔναρξε θρήνου· κατὰ δὲ τὴν κρίσιν αὐτοῦ περίστειλον τὸ σῶμα αὐτοῦ, καὶ μὴ ὑπερίδῃς τὴν ταφὴν αὐτοῦ.
\vs{17}Πίκρανον κλαυθμὸν, καὶ θέρμανον κοπετὸν, καὶ ποίησον τὸ πένθος κατὰ τὴν ἀξίαν αὐτοῦ ἡμέραν μίαν καὶ δύο χάριν διαβολῆς, καὶ παρακλήθητι λύπης ἕνεκα·
\vs{18}ἀπὸ λύπης γὰρ ἐκβαίνει θάνατος, καὶ λύπη καρδίας κάμψει ἰσχύν.
\vs{19}Ἐν ἀπαγωγῇ παραβαίνει καὶ λύπη, καὶ βίος πτωχοῦ κατὰ καρδίας.
\vs{20}Μὴ δῷς εἰς λύπην τὴν καρδίαν σου, ἀπόστησον αὐτὴν μνησθεὶς τὰ ἔσχατα.
\vs{21}Μὴ ἐπιλαθῇ, οὐ γάρ ἐστιν ἐπάνοδος, καὶ τοῦτον οὐκ ὠφελήσεις, καὶ σεαυτὸν κακώσεις.
\vs{22}Μνήσθητι τὸ κρίμα αὐτοῦ, ὅτι οὕτω ὡς καὶ τὸ σόν· ἐμοὶ χθὲς, καὶ σοὶ σήμερον.
\vs{23}Ἐν ἀναπαύσει νεκροῦ κατάπαυσον τὸ μνημόσυνον αὐτοῦ, καὶ παρακλήθητι ἐν αὐτῷ ἐν ἐξόδῳ πνεύματος αὐτοῦ.

\vs{24}Σοφία γραμματέως ἐν εὐκαιρίᾳ σχολῆς, καὶ ὁ ἐλασσούμενος πράξει αὐτοῦ σοφισθήσεται.
\vs{25}Τί σοφισθήσεται ὁ κρατῶν ἀροτροῦ, καὶ καυχώμενος ἐν δόρατι κέντρου, βόας ἐλαύνων καὶ ἀναστρεφόμενος ἐν ἔργοις αὐτῶν, καὶ ἡ διήγησις αὐτοῦ ἐν υἱοῖς ταυρῶν;
\vs{26}Καρδίαν αὐτοῦ δώσει ἐκδοῦναι αὔλακας, καὶ ἡ ἀγρυπνία αὐτοῦ εἰς χορτάσματα δαμάλεων.
\vs{27}Οὕτως πᾶς τέκτων καὶ ἀρχιτέκτων, ὅστις νύκτωρ ὡς ἡμέρᾳ διάγει· οἱ γλύφοντες γλύμματα σφραγίδων, καὶ ἡ ὑπομονὴ αὐτοῦ ἀλλοιῶσαι ποικιλίαν· καρδίαν αὐτοῦ δώσει εἰς τὸ ὁμοιῶσαι ζωγραφίαν, καὶ ἡ ἀγρυπνία αὐτοῦ τελέσαι ἔργον.

\vs{28}Οὕτως χαλκεὺς καθήμενος ἐγγὺς ἄκμονος, καὶ καταμανθάνων ἀργῷ σιδήρῳ· ἀτμὶς πυρὸς πήξει σάρκας αὐτοῦ, καὶ ἐν θέρμῃ καμίνου διαμαχήσεται· φωνὴ σφύρης καινιεῖ τὸ οὖς αὐτοῦ, καὶ κατέναντι ὁμοιώματος σκεύους οἱ ὀφθαλμοὶ αὐτοῦ· καρδίαν αὐτοῦ δώσει εἰς συντέλειαν ἔργων, καὶ ἡ ἀγρυπνία αὐτοῦ κοσμῆσαι ἐπὶ συντελείας.

\vs{29}Οὕτω κεραμεὺς καθήμενος ἐν ἔργῳ αὐτοῦ, καὶ συστρέφων ἐν ποσὶν αὐτοῦ τροχὸν, ὃς ἐν μεριμνῃ κεῖται διαπαντὸς ἐπὶ τὸ ἔργον αὐτοῦ, καὶ ἐναρίθμιος πᾶσα ἡ ἐργασία αὐτοῦ.
\vs{30}Ἐν βραχίονι αὐτοῦ τυπώσει πηλὸν, καὶ πρὸ ποδῶν κάμψει ἰσχὺν αὐτοῦ· καρδίαν ἐπιδώσει συντελέσαι τὸ χρίσμα, καὶ ἡ ἀγρυπνία αὐτοῦ καθαρίσαι κάμινον.
\vs{31}Πάντες οὗτοι εἰς χεῖρας αὐτῶν ἐνεπίστευσαν, καὶ ἕκαστος ἐν τῷ ἔργῳ αὐτοῦ σοφίζεται.
\vs{32}Ἄνευ αὐτῶν οὐκ οἰκισθήσεται πόλις, καὶ οὐ παροικήσουσιν, οὐδὲ περιπατήσουσι·
\vs{33}καὶ ἐν ἐκκλησίᾳ οὐχ ὑπεραλοῦνται· ἐπὶ δίφρον δικαστοῦ οὐ καθιοῦνται, καὶ διαθήκην κρίματος οὐ διανοηθήσονται, οὐδὲ μὴ ἐκφανῶσι δικαιοσύνην καὶ κρίμα· καὶ ἐν παραβολαῖς οὐχ εὑρεθήσονται.
\vs{34}Ἀλλὰ κτίσμα αἰῶνος στηρίσουσι· καὶ ἡ δέησις αὐτῶν ἐν ἐργασίᾳ τέχνης.

Πλὴν τοῦ ἐπιδόντος τὴν ψυχὴν αὐτοῦ, καὶ διανοουμένου ἐν νόμῳ ὑψίστου,

\ch{39}
σοφίαν πάντων ἀρχαίων ἐκζητήσει, καὶ ἐν προφητείαις ἀσχοληθήσεται.
\vs{2}Διηγήσεις ἀνδρῶν ὀνομαστῶν συντηρήσει, καὶ ἐν στροφαῖς παραβολῶν συνεισελεύσεται.
\vs{3}Ἀπόκρυφα παροιμιῶν ἐκζητήσει, καὶ ἐν αἰνίγμασι παραβολῶν ἀναστραφήσεται.
\vs{4}Ἀναμέσον μεγιστάνων ὑπηρετήσει, καὶ ἔναντι ἡγουμένου ὀφθήσεται· ἐν γῇ ἀλλοτρίων ἐθνῶν διελεύσεται, ἀγαθὰ γὰρ καὶ κακὰ ἐν ἀνθρώποις ἐπείρασε.
\vs{5}Τὴν καρδίαν αὐτοῦ ἐπιδώσει ὀρθρίσαι πρὸς Κύριον τὸν ποιήσαντα αὐτὸν, καὶ ἔναντι ὑψίστου δεηθήσεται· καὶ ἀνοίξει τὸ στόμα αὐτοῦ ἐν προσευχῇ, καὶ περὶ τῶν ἁμαρτιῶν αὐτοῦ δεηθήσεται.

\vs{6}Ἐὰν Κύριος ὁ μέγας θελήσῃ, πνεύματι συνέσεως ἐμπλησθήσεται· αὐτὸς ἀνομβρήσει ῥήματα σοφίας αὐτοῦ, καὶ ἐν προσευχῇ ἐξομολογήσεται Κυρίῳ.
\vs{7}Αὐτὸς κατευθύνει βουλὴν αὐτοῦ καὶ ἐπιστήμην, καὶ ἐν τοῖς ἀποκρύφοις αὐτοῦ διανοηθήσεται.
\vs{8}Αὐτὸς ἐκφανεῖ παιδείαν διδασκαλίας αὐτοῦ, καὶ ἐν νόμῳ διαθήκης Κυρίου καυχήσεται.
\vs{9}Αἰνέσουσι τὴν σύνεσιν αὐτοῦ πολλοὶ, ἕως τοῦ αἰῶνος οὐκ ἐξαλειφθήσεται· οὐκ ἀποστήσεται τὸ μνημόσυνον αὐτοῦ, καὶ ὄνομα αὐτοῦ ζήσεται εἰς γενεὰς γενεῶν.
\vs{10}Τὴν σοφίαν αὐτοῦ διηγήσονται ἔθνη, καὶ τὸν ἔπαινον αὐτοῦ ἐξαγγελεῖ ἐκκλησία.
\vs{11}Ἐὰν ἐμμείνῃ ὄνομα καταλείψει ἢ χίλιοι, καὶ ἐὰν ἀναπαύσηται ἐμποιεῖ αὐτῷ.
\vs{12}Ἔτι διανοηθεὶς ἐκδιηγήσομαι, καὶ ὡς διχομηνία ἐπληρώθην.

\vs{13}Εἰσακούσετέ μου υἱοὶ ὅσιοι, καὶ βλαστήσατε ὡς ῥόδον φυόμενον ἐπὶ ῥεύματος ἀγροῦ·
\vs{14}καὶ ὡς λίβανος εὐωδιάσατε ὀσμὴν, καὶ ἀνθήσατε ἄνθος ὡς κρίνον· διάδοτε ὀσμὴν καὶ αἰνέσατε ᾆσμα· εὐλογήσατε Κύριον ἐπὶ πᾶσι τοῖς ἔργοις.
\vs{15}Δότε τῷ ὀνόματι αὐτοῦ μεγαλωσύνην, καὶ ἐξομολογήσασθε ἐν αἰνέσει αὐτοῦ, ἐν ᾠδαῖς χειλέων καὶ ἐν κινύραις, καὶ οὕτως ἐρεῖτε ἐν ἐξομολογήσει,
\vs{16}Τὰ ἔργα Κυρίου πάντα ὅτι καλὰ σφόδρα, καὶ πᾶν πρόσταγμα ἐν καιρῷ αὐτοῦ ἔσται.
\vs{17}Πάντα γὰρ ἐν καιρῷ αὐτοῦ ζητηθήσεται· ἐν λόγῳ αὐτοῦ ἔστη ὡς θημωνία ὕδωρ, καὶ ἐν ῥήματι στόματος αὐτοῦ ἀποδοχεῖα ὑδάτων.
\vs{18}Ἐν προστάγματι αὐτοῦ πᾶσα ἡ εὐδοκία, καὶ οὐκ ἔστιν ὃς ἐλαττώσει τὸ σωτήριον αὐτοῦ.
\vs{19}Ἔργα πάσης σαρκὸς ἐνώπιον αὐτοῦ, καὶ οὐκ ἔστι κρυβῆναι ἀπὸ τῶν ὀφθαλμῶν αὐτοῦ.
\vs{20}Ἀπὸ τοῦ αἰῶνος εἰς τὸν αἰῶνα ἐπέβλεψε, καὶ οὐθέν ἐστι θαυμάσιον ἐναντίον αὐτοῦ.

\vs{21}Οὐκ ἔστιν εἰπεῖν, τί τοῦτο; εἰς τί τοῦτο; πάντα γὰρ εἰς χρείας αὐτῶν ἔκτισται.
\vs{22}Ἡ εὐλογία αὐτοῦ ὡς ποταμὸς ἐπεκάλυψε, καὶ ὡς κατακλυσμὸς ξηρὰν ἐμέθυσεν.
\vs{23}Οὕτως ὀργὴν αὐτοῦ ἔθνη κληρονομήσει, ὡς μετέστρεψεν ὕδατα εἰς ἅλμην.

\vs{24}Αἱ ὁδοὶ αὐτοῦ τοῖς ὁσίοις εὐθεῖαι, οὕτως τοῖς ἀνόμοις προσκόμματα.
\vs{25}Ἀγαθὰ τοῖς ἀγαθοῖς ἔκτισται ἀπʼ ἀρχῆς, οὕτως τοῖς ἁμαρτωλοῖς κακά.
\vs{26}Ἀρχὴ πάσης χρείας εἰς ζωὴν ἀνθρώπου, ὕδωρ, πῦρ, καὶ σίδηρος, καὶ ἅλα, καὶ σεμίδαλις πυροῦ, καὶ μέλι, καὶ γάλα, αἷμα σταφυλῆς, καὶ ἔλαιον, καὶ ἱμάτιον.
\vs{27}Ταῦτα πάντα τοῖς εὐσεβέσιν εἰς ἀγαθὰ, οὕτως τοῖς ἁμαρτωλοῖς τραπήσεται εἰς κακά.

\vs{28}Ἔστι πνεύματα ἃ εἰς ἐκδίκησιν ἔκτισται, καὶ ἐν θυμῷ αὐτῶν ἐστερέωσαν μάστιγας αὐτῶν· καὶ ἐν καιρῷ συντελείας ἰσχὺν ἐκχέουσι, καὶ τὸν θυμὸν τοῦ ποιήσαντος αὐτοὺς κοπάσουσι.
\vs{29}Πῦρ καὶ χάλαζα καὶ λιμὸς καὶ θάνατος, πάντα ταῦτα εἰς ἐκδίκησιν ἔκτισται.
\vs{30}Θηρίων ὁδόντες καὶ σκορπίοι καὶ ἔχεις καὶ ῥομφαία ἐκδικοῦσα εἰς ὄλεθρον ἀσεβεῖς,
\vs{31}ἐν τῇ ἐντολῇ αὐτοῦ εὐφρανθήσονται, καὶ ἐπὶ τῆς γῆς εἰς χρείας ἑτοιμασθήσονται, καὶ ἐν τοῖς καιροῖς αὐτῶν οὐ παραβήσονται λόγον.

\vs{32}Διὰ τοῦτο ἐξ ἀρχῆς ἐστηρίχθην, καὶ διενοήθην, καὶ ἐν γραφῇ ἀφῆκα.
\vs{33}Τὰ ἔργα Κυρίου πάντα ἀγαθὰ, καὶ πᾶσαν χρείαν ἐν ὥρᾳ αὐτῆς χορηγήσει.
\vs{34}Καὶ οὐκ ἔστιν εἰπεῖν, τοῦτο τούτου πονηρότερον, πάντα γὰρ ἐν καιρῷ εὐδοκιμηθήσεται.
\vs{35}Καὶ νῦν ἐν πάσῃ καρδίᾳ καὶ στόματι ὑμνήσατε, καὶ εὐλογήσατε τὸ ὄνομα Κυρίου.

\ch{40}
Ἀσχολιὰ μεγάλη ἔκτισται παντὶ ἀνθρώπῳ, καὶ ζυγὸς βαρὺς ἐπὶ υἱοὺς Ἀδὰμ, ἀφʼ ἡμέρας ἐξόδου ἐκ γαστρὸς μητρὸς αὐτῶν ἕως ἡμέρας ἐπὶ ταφῇ εἰς μητέρα πάντων.
\vs{2}Τοὺς διαλογισμοὺς αὐτῶν καὶ φόβον καρδίας, ἐπίνοια προσδοκίας, ἡμέρα τελευτῆς·
\vs{3}ἀπὸ καθημένου ἐπὶ θρόνου ἐν δόξῃ, καὶ ἕως τεταπεινωμένου ἐν γῇ καὶ σποδῷ·
\vs{4}ἀπὸ φοροῦντος ὑάκινθον καὶ στέφανον, καὶ ἕως περιβαλλουμένου ὠμόλινον·
\vs{5}θυμὸς καὶ ζῆλος καὶ ταραχὴ καὶ σάλος καὶ φόβος θανάτου καὶ μηνίαμα καὶ ἔρις, καὶ ἐν καιρῷ ἀναπαύσεως ἐπὶ κοίτης ὕπνος νυκτὸς ἀλλοιοῖ γνῶσιν αὐτοῦ.
\vs{6}Ὀλίγον ὡς οὐδὲν ἐν ἀναπαύσει, καὶ ἀπʼ ἐκείνου ἐν ὕπνοις ὡς ἐν ἡμέρᾳ σκοπιᾶς, τεθορυβημένος ἐν ὁράσει καρδίας αὐτοῦ, ὡς ἐκπέφευγὼς ἀπὸ προσώπου πολέμου.
\vs{7}Ἐν καιρῷ σωτηρίας αὐτοῦ ἐξηγέρθη, καὶ ἀποθαυμάζων εἰς οὐδένα φόβον.

\vs{8}Μετὰ πάσης σαρκὸς ἀπὸ ἀνθρώπου ἕως κτήνους, καὶ ἐπὶ ἁμαρτωλῶν ἑπταπλάσια πρὸς ταῦτα.
\vs{9}Θάνατος καὶ αἷμα καὶ ἔρις καὶ ῥομφαία,. ἐπαγωγαὶ, λιμὸς καὶ σύντριμμα καὶ μάστιξ,
\vs{10}ἐπὶ τοὺς ἀνόμους ἐκτίσθη ταῦτα πάντα, καὶ διʼ αὐτοὺς ἐγένετο ὁ κατακλυσμός.
\vs{11}Πάντα ὅσα ἀπὸ γῆς εἰς γῆν ἀναστρέφει, καὶ ἀπὸ ὑδάτων εἰς θάλασσαν ἀνακάμπτει.
\vs{12}Πᾶν δῶρον καὶ ἀδικία ἐξαλειφθήσεται, καὶ πίστις εἰς τὸν αἰῶνα στήσεται.
\vs{13}Χρήματα ἀδίκων ὡς ποταμὸς ξηρανθήσεται, καὶ ὡς βροντὴ μεγάλη ἐν ὑετῷ ἐξηχήσει.

\vs{14}Ἐν τῷ ἀνοῖξαι αὐτὸν χεῖρας, εὐφρανθήσεται, οὕτως οἱ παραβαίνοντες εἰς συντέλειαν ἐκλείψουσιν.
\vs{15}Ἔκγονα ἀσεβῶν οὐ πληθύνει κλάδους, καὶ ῥίζαι ἀκάθαρτοι ἐπʼ ἀκροτόμου πέτρας.
\vs{16}Ἄχει ἐπὶ παντὸς ὕδατος καὶ χείλους ποταμοῦ πρὸ παντὸς χόρτου ἐκτιλήσεται.

\vs{17}Χάρις ὡς παράδεισος ἐν εὐλογίαις, καὶ ἐλεημοσύνη εἰς τὸν αἰῶνα διαμένει.
\vs{18}Ζωὴ αὐτάρκους ἐργάτου γλυκανθήσεται, καὶ ὑπὲρ ἀμφότερα ὁ εὑρίσκων θησαυρόν.
\vs{19}Τέκνα καὶ οἰκοδομὴ πόλεως στηρίζουσιν ὄνομα, καὶ ὑπὲρ ἀμφότερα γυνὴ ἄμωμος λογίζεται.
\vs{20}Οἶνος καὶ μουσικὰ εὐφραίνουσι καρδίαν, καὶ ὑπὲρ ἀμφότερα ἀγάπησις σοφίας.

\vs{21}Αὐλὸς καὶ ψαλτήριον ἡδύνουσι μέλι, καὶ ὑπὲρ ἀμφότερα γλῶσσα ἡδεῖα.
\vs{22}Χάριν καὶ κάλλος ἐπιθυμήσει ὁ ὀφθαλμός σου, καὶ ὑπὲρ ἀμφότερα χλόην σπόρου.
\vs{23}Φίλος καὶ ἑταῖρος εἰς καιρὸν ἀπαντῶντες, καὶ ὑπὲρ ἀμφότερα γυνὴ μετὰ ἀνδρός.
\vs{24}Ἀδελφοὶ καὶ βοήθεια εἰς καιρὸν θλίψεως, καὶ ὑπὲρ ἀμφότερα ἐλεημοσύνη ῥύσεται.
\vs{25}Χρυσίον καὶ ἀργύριον ἐπιστήσουσι πόδα, καὶ ὑπὲρ ἀπφότερα βουλὴ εὐδοκιμεῖται.
\vs{26}Χρήματα καὶ ἰσχὺς ἀνυψώσουσι καρδίαν, καὶ ὑπὲρ ἀμφότερα φόβος Κυρίου· οὐκ ἔστι φόβῳ Κυρίου ἐλάττωσις, καὶ οὐκ ἔστιν ἐπιζητῆσαι ἐν αὐτῷ βοήθειαν.
\vs{27}Φόβος Κυρίου ὡς παράδεισος εὐλογίας, καὶ ὑπὲρ πᾶσαν δόξαν ἐκάλυψαν αὐτόν.

\vs{28}Τέκνον, ζωὴν ἐπαιτήσεως μὴ βιώσῃς, κρεῖσσον ἀποθανεῖν ἢ ἐπαιτεῖν.
\vs{29}Ἀνὴρ βλέπων εἰς τράπεζαν ἀλλοτρίαν, οὐκ ἔστιν αὐτοῦ ὁ βίος ἐν λογισμῷ ζωῆς, ἀλισγήσει τὴν ψυχὴν αὐτοῦ ἐν ἐδέσμασιν ἀλλοτρίοις· ἀνὴρ δὲ ἐπιστήμων καὶ πεπαιδευμένος φυλάξεται.
\vs{30}Ἐν στόματι ἀναιδοῦς γλυκανθήσεται ἐπαίτησις, καὶ ἐν κοιλίᾳ αὐτοῦ πῦρ καήσεται.

\ch{41}
Ὦ θάνατε, ὡς πικρόν σου τὸ μνημόσυνόν ἐστιν ἀνθρώπῳ εἰρηνεύοντι ἐν τοῖς ὑπάρχουσιν αὐτοῦ, ἀνδρὶ ἀπερισπάστῳ, καὶ εὐοδουμένῳ ἐν πᾶσι, καὶ ἔτι ἰσχύοντι ἐπιδέξασθαι τροφήν;
\vs{2}Ὦ θάνατε, καλόν σου τὸ κρίμα ἐστὶν ἀνθρώπῳ ἐπιδεομένῳ καὶ ἐλασσουμένῳ ἰσχύϊ, ἐσχατογήρῳ, καὶ περισπωμένῳ περὶ πάντων, καὶ ἀπειθοῦντι, καὶ ἀπολωλεκότι ὑπομονήν.
\vs{3}Μὴ εὐλαβοῦ κρίμα θανάτου· μνήσθητι προτέρων σου καὶ ἐσχάτων, τοῦτο τὸ κρίμα παρὰ Κυρίου πάσῃ σαρκί.
\vs{4}Καὶ τί ἀπαναίνῃ ἐν εὐδοκίᾳ ὑψίστου; εἴτε δέκα, εἴτε ἑκατὸν, εἴτε χίλια ἔτη· οὐκ ἔστιν ἐν ᾅδου ἐλεγμὸς ζωῆς.

\vs{5}Τέκνα βδελυκτὰ γίνεται τέκνα ἁμαρτωλῶν, καὶ συναναστρεφόμενα παροικίαις ἀσεβῶν.
\vs{6}Τέκνων ἁμαρτωλῶν ἀπολεῖται κληρονομία, καὶ μετὰ τοῦ σπέρματος αὐτῶν ἐνδελεχιεῖ ὄνειδος.
\vs{7}Πατρὶ ἀσεβεῖ μέμψεται τέκνα, ὅτι διʼ αὐτὸν ὀνειδισθήσονται.
\vs{8}Οὐαὶ ὑμῖν ἄνδρες ἀσεβεῖς, οἵτινες ἐγκατελίπετε νόμον Θεοῦ ὑψίστου.
\vs{9}Καὶ ἐὰν γεννηθῆτε, εἰς κατάραν γεννηθήσεσθε· καὶ ἐὰν ἀποθάνητε, εἰς κατάραν μερισθήσεσθε.

\vs{10}Πάντα ὅσα ἐκ γῆς, εἰς γῆν ἀπελεύσεται· οὕτως ἀσεβεῖς ἀπὸ κατάρας εἰς ἀπώλείαν.
\vs{11}Πένθος ἀνθρώπων ἐν σώμασιν αὐτῶν, ὄνομα δὲ ἁμαρτωλῶν οὐκ ἀγαθὸν, ἐξαλειφθήσεται.
\vs{12}Φρόντισον περὶ ὀνόματος, αὐτὸ γάρ σοι διαμενεῖ, ἢ χίλιοι μεγάλοι θησαυροὶ χρυσίου.
\vs{13}Ἀγαθῆς ζωῆς ἀριθμὸς ἡμερῶν, καὶ ἀγαθὸν ὄνομα εἰς αἰῶνα διαμένει.

\vs{14}Παιδείαν ἐν εἰρήνῃ συντηρήσατε τέκνα, σοφία δὲ κεκρυμμένη καὶ θησαυρὸς ἀφανὴς, τίς ὠφέλεια ἐν ἀμφοτέροις;
\vs{15}Κρείσσων ἄνθρωπος ἀποκρύπτων τὴν μωρίαν αὐτοῦ, ἢ ἄνθρωπος ἀποκρύπτων τὴν σοφίαν αὐτοῦ.
\vs{16}Τοιγαροῦν ἐντράπητε ἐπὶ τῷ ῥήματί μου· οὐ γάρ ἐστι πᾶσαν αἰσχύνην διαφυλάξαι καλὸν, καὶ οὐ πάντα πᾶσιν ἐν πίστει εὐδοκιμεῖται.

\vs{17}Αἰσχύνεσθε ἀπὸ πατρὸς καὶ μητρὸς περὶ πορνείας, καὶ ἀπὸ ἡγουμένου καὶ δυνάστου περὶ ψεύδους·
\vs{18}ἀπὸ κριτοῦ καὶ ἄρχοντος περὶ πλημμελείας, ἀπὸ συναγωγῆς καὶ λαοῦ περὶ ἀνομίας· ἀπὸ κοινωνοῦ καὶ φίλου περὶ ἀδικίας, καὶ ἀπὸ τόπου οὗ παροικεῖς περὶ κλοπῆς·
\vs{19}καὶ ἀπὸ ἀληθείας Θεοῦ καὶ διαθήκης, καὶ ἀπὸ πήξεως ἀγκῶνος ἐπʼ ἄρτοις· ἀπὸ σκορακισμοῦ λήψεως καὶ δόσεως,
\vs{20}καὶ ἀπὸ ἀσπαζομένων περὶ σιωπῆς· ἀπὸ ὁράσεως γυναικὸς ἑταίρας,
\vs{21}καὶ ἀπὸ ἀποστροφῆς προσώπου συγγενοῦς· ἀπὸ ἀφαιρέσεως μερίδος καὶ δόσεως, καὶ ἀπὸ κατανοήσεως γυναικὸς ὑπάνδρου,
\vs{22}ἀπὸ περιεργείας παιδίσκης αὐτοῦ, καὶ μὴ ἐπιστῇς ἐπὶ τὴν κοίτην αὐτῆς, ἀπὸ φίλων περὶ λόγων ὀνειδισμοῦ, καὶ μετὰ τὸ δοῦναι, μὴ ὀνείδιζε.

\vs{23}Ἀπὸ δευτερώσεως καὶ λόγου ἀκοῆς, καὶ ἀπὸ ἀποκαλύψεων λόγων κρυφίων·
\vs{24}καὶ ἔσῃ αἰσχυντηρὸς ἀληθινῶς, καὶ εὑρίσκων χάριν ἔναντι παντὸς ἀνθρώπου·

\ch{42}μὴ περὶ τούτων αἰσχυνθῇς, καὶ μὴ λάβῃς πρόσωπον τοῦ ἁμαρτάνειν·
\vs{2}περὶ νόμου ὑψίστου καὶ διαθήκης, καὶ περὶ κρίματος δικαιῶσαι τὸν ἀσεβῆ·
\vs{3}περὶ λόγου κοινωνοῦ καὶ ὁδοιπόρων, καὶ περὶ δόσεως κληρονομίας ἑταίρων·
\vs{4}περὶ ἀκριβείας ζυγοῦ καὶ σταθμιῶν, περὶ κτήσεως πολλῶν καὶ ὀλίγων·
\vs{5}περὶ ἀδιαφόρου πράσεως, καὶ ἐμπόρων, καὶ περὶ παιδείας τέκνων πολλῆς, καὶ οἰκέτῃ πονηρῷ πλευρὰν αἱμάξαι.

\vs{6}Ἐπὶ γυναικὶ πονηρᾷ καλὸν σφραγίς· καὶ ὅπου χεῖρες πολλαὶ, κλεῖσον.
\vs{7}Ὃ ἐὰν παραδίδως, ἐν ἀριθμῷ καὶ σταθμῷ, καὶ δόσις καὶ λῆψις παντὶ ἐν γραφῇ.
\vs{8}Περὶ παιδείας ἀνοήτου καὶ μωροῦ καὶ ἐσχατογήρου κρινομένου πρὸς νέους, καὶ ἔσῃ πεπαιδευμένος ἀληθινῶς, καὶ δεδοκιμασμένος ἔναντι παντὸς ζῶντος.

\vs{9}Θυγάτηρ πατρὶ ἀπόκρυφος ἀγρυπνία, καὶ ἡ μέριμνα αὐτῆς ἀφιστᾷ ὕπνον· ἐν νεότητι αὐτῆς μήποτε παρακμάσῃ, καὶ συνῳκηκυῖα μήποτε μισηθῇ.
\vs{10}Ἐν παρθενίᾳ μήποτε βεβηλωθῇ, καὶ ἐν τοῖς πατρικοῖς αὐτῆς ἔγκυος γένηται· μετὰ ἀνδρὸς οὖσα μήποτε παραβῇ, καὶ συνῳκηκυῖα μήποτε στειρώσῃ.
\vs{11}Ἐπὶ θυγατρὶ ἀδιατρέπτῳ στερέωσον φυλακὴν, μήποτε ποιήσῃ σε ἐπίχαρμα ἐχθροῖς, λαλιὰν ἐν πόλει, καὶ ἔκκλητον λαοῦ, καὶ καταισχύνῇ σε ἐν πλήθει πολλῶν.
\vs{12}Παντὶ ἀνθρώπῳ μὴ ἔμβλεπε ἐν κάλλει, καὶ ἐν μέσῳ γυναικῶν μὴ συνέδρευε·
\vs{13}ἀπὸ γὰρ ἱματίων ἐκπορεύεται σὴς, καὶ ἀπὸ γυναικὸς πονηρία γυναικός.
\vs{14}Κρείσσων πονηρία ἀνδρὸς ἢ ἀγαθοποιὸς γυνὴ, καὶ γυνὴ καταισχύνουσα εἰς ὀνειδισμόν.

\vs{15}Μνησθήσομαι δὴ τὰ ἔργα Κυρίου, καὶ ἃ ἑώρακα ἐκδιηγήσομαι· ἐν λόγοις Κυρίου τὰ ἔργα αὐτοῦ.
\vs{16}Ἥλιος φωτίζων κατὰ πᾶν ἐπέβλεψε, καὶ τῆς δόξης αὐτοῦ πλῆρες τὸ ἔργον αὐτοῦ.

\vs{17}Οὐκ ἐνεποίησε τοῖς ἁγίοις Κύριος ἐκδιηγήσασθαι πάντα τὰ θαυμάσια αὐτοῦ, ἃ ἐστερέωσε Κύριος ὁ παντοκράτωρ, στηριχθῆναι ἐν δόξῃ αὐτοῦ τὸ πᾶν.
\vs{18}Ἄβυσσον καὶ καρδίαν ἐξίχνευσε, καὶ ἐν πανουργεύμασιν αὐτῶν διενοήθη· ἔγνω γὰρ ὁ Κύριος πᾶσαν εἴδησιν, καὶ ἐνέβλεψεν εἰς σημεῖον αἰῶνος·
\vs{19}ἀπαγγέλλων τὰ παρεληλυθότα καὶ ἐπεσόμενα, καὶ ἀποκαλύπτων ἴχνη ἀποκρύφων.
\vs{20}Οὐ παρῆλθεν αὐτὸν πᾶν διανόημα, οὐκ ἐκρύβη ἀπʼ αὐτοῦ οὐδὲ εἷς λόγος.

\vs{21}Τὰ μεγαλεῖα τῆς σοφίας αὐτοῦ ἐκόσμησε, καὶ ἕως ἐστὶ πρὸ τοῦ αἰῶνος καὶ εἰς τὸν αἰῶνα, οὔτε προσετέθη οὔτε ἠλαττώθη, καὶ οὐδὲ προσεδεήθη οὐδενὸς συμβούλου.
\vs{22}Ὡς πάντα τὰ ἔργα αὐτοῦ ἐπιθυμητὰ, καὶ ὡς σπινθῆρός ἐστι θεωρῆσαι.
\vs{23}Πάντα ταῦτα ζῇ καὶ μένει εἰς τὸν αἰῶνα ἐν πάσαις χρείαις, καὶ πάντα ὑπακούει.
\vs{24}Πάντα δισσὰ ἓν κατέναντι τοῦ ἑνὸς, καὶ οὐκ ἐποίησεν οὐδὲν ἐκλεῖπον.
\vs{25}Ἓν τοῦ ἑνὸς ἐστερέωσε τὰ ἀγαθὰ, καὶ τίς πλησθήσεται ὁρῶν δόξαν αὐτοῦ;

\ch{43}
Γαυρίαμα ὕψους, στερέωμα καθαριότητος, εἶδος οὐρανοῦ ἐν ὁράματι δόξης.
\vs{2}Ἥλιος ἐν ὀπτασίᾳ διαγγέλλων ἐν ἐξόδῳ, σκεῦος θαυμαστὸν, ἔργον ὑψίστου.
\vs{3}Ἐν μεσημβρίᾳ αὐτοῦ ἀναξηραίνει χώραν, καὶ ἐναντίον καύματος αὐτοῦ τίς ὑποστήσεται;
\vs{4}Κάμινον φυσῶν ἐν ἔργοις καύματος, τριπλασίως ἥλιος ἐκκαίων ὄρη· ἀτμίδας πυρώδεις ἐμφυσῶν, καὶ ἐκλάμπων ἀκτίνας ἀμαυροῖ ὀφθαλμούς.
\vs{5}Μέγας Κύριος ὁ ποιήσας αὐτὸν, καὶ ἐν λόγοις αὐτοῦ κατέσπευσε πορείαν.
\vs{6}Καὶ ἡ σελήνη ἐν πᾶσιν εἰς καιρὸν αὐτῆς, ἀνάδειξιν χρόνων, καὶ σημεῖον αἰῶνος.
\vs{7}Ἀπὸ σελήνης σημεῖον ἑορτῆς, φωστὴρ μειούμενος ἐπὶ συντελείας.
\vs{8}Μὴν κατὰ τὸ ὄνομα αὐτῆς ἐστιν, αὐξανομένη θαυμαστῶς ἐν ἀλλοιώσει· σκεῦος παρεμβολῶν ἐν ὕψει, ἐν στερεώματι οὐρανοῦ ἐκλάμπων·
\vs{9}κάλλος οὐρανοῦ, δόξα ἄστρων, κόσμος φωτίζων, ἐν ὑψίστοις Κύριος.
\vs{10}Ἐν λόγοις ἁγίου στήσονται κατὰ κρίμα, καὶ οὐ μὴ ἐκλυθῶσιν ἐν φυλακαῖς αὐτῶν.
\vs{11}Ἴδε τόξον, καὶ εὐλόγησον τὸν ποιήσαντα αὐτὸ, σφόδρα ὡραῖον ἐν τῷ αὐγάσματι αὐτοῦ.
\vs{12}Ἐγύρωσεν οὐρανὸν ἐν κυκλώσει δόξης, χεῖρες ὑψίστου ἐτάννυσαν αὐτό.
\vs{13}Προστάγματι αὐτοῦ κατέσπευσε χιόνα, καὶ ταχύνει ἀστραπὰς κρίματος αὐτοῦ.
\vs{14}Διὰ τοῦτο ἠνεῴχθησαν θησαυροὶ, καὶ ἐξέπτησαν νεφέλαι ὡς πετεινά.
\vs{15}Ἐν μεγαλείῳ αὐτοῦ ἴσχυσε νεφέλας, καὶ διεθρύβησαν λίθοι χαλάζης.

\vs{16}Καὶ ἐν ὀπτασίαις αὐτοῦ σαλευθήσεται ὄρη, ἐν θελήματι πνεύσεται νότος.
\vs{17}Φωνὴ βροντῆς αὐτοῦ ὠδίνησε γῆν, καὶ καταιγὶς Βορέου καὶ συστροφὴ πνεύματος· ὡς πετεινὰ καθιπτάμενα πάσσει χιόνα, καὶ ὡς ἀκρὶς καταλύουσα ἡ κατάβασις αὐτῆς.
\vs{18}Κάλλος λευκότητος αὐτῆς ἐκθαυμάσει ὀφθαλμὸς, καὶ ἐπὶ τοῦ ὑετοῦ αὐτῆς ἐκστήσεται καρδία.
\vs{19}Καὶ πάχνην ὡς ἅλα ἐπὶ γῆς χέει, καὶ παγεῖσα γίνεται σκολόπων ἄκρα.

\vs{20}Ψυχρὸς ἄνεμος Βορέης πνεύσει, καὶ παγήσεται κρύσταλλος ἀφʼ ὕδατος· ἐπὶ πᾶσαν συναγωγὴν ὕδατος καταλύσει, καὶ ὡς θώρακα ἐνδύσεται τὸ ὕδωρ.
\vs{21}Καταφάγεται ὄρη, καὶ ἔρημον ἐκκαύσει, καὶ ἀποσβέσει χλόην ὡς πῦρ.
\vs{22}Ἴασις πάντων κατὰ σπουδὴν ὁμίχλη, δρόσος ἀπαντῶσα ἀπὸ καύσωνος ἱλαρώσει.

\vs{23}Λογισμῷ αὐτοῦ ἐκόπασεν ἄβυσσον, καὶ ἐφύτευσεν αὐτὴν Ἰησοῦς.
\vs{24}Οἱ πλέοντες τὴν θάλασσαν διηγοῦνται τὸν κίνδυνον αὐτῆς, καὶ ἀκοαῖς ὠτίῳν ἡμῶν θαυμάζομεν.
\vs{25}Καὶ ἐκεῖ τὰ παρὰδοξα καὶ θαυμάσια ἔργα, ποικιλία παντὸς ζώου, κτίσις κητῶν.
\vs{26}Διʼ αὐτὸν εὐοδία τέλος αὐτοὺ, καὶ ἐν λόγῳ αὐτοῦ σύγκειται πάντα.

\vs{27}Πολλὰ ἐροῦμεν καὶ οὐ μὴ ἐφικώμεθα, καὶ συντέλεια λόγων τὸ πᾶν ἐστιν αὐτός.
\vs{28}Δοξάζοντες ποῦ ἰσχύσωμεν; αὐτὸς γὰρ ὁ μέγας παρὰ πάντα τὰ ἔργα αὐτοῦ.
\vs{29}Φοβερὸς Κύριος καὶ σφόδρα μέγας, καὶ θαυμαστὴ ἡ δυναστεία αὐτοῦ.
\vs{30}Δοξάζοντες Κύριον ὑψώσατε καθόσον ἂν δύνησθε, ὑπερέξει γὰρ καὶ ἔτι· καὶ ὑψοῦντες αὐτὸν πληθύνατε ἐν ἰσχύϊ, μὴ κοπιᾶτε, οὐ γὰρ μὴ ἐφίκησθε.
\vs{31}Τίς ἑώρακεν αὐτὸν καὶ ἐκδιηγήσεται; καὶ τίς μεγαλύνει αὐτὸν καθώς ἐστι;
\vs{32}Πολλὰ ἀπόκρυφά ἐστι μείζονα τούτων, ὀλίγα γὰρ ἑωράκαμεν τῶν ἔργων αὐτοῦ.
\vs{33}Πάντα γὰρ ἐποίησεν ὁ Κύριος, καὶ τοῖς εὐσεβέσιν ἔδωκε σοφίαν.

ΠΑΤΕΡΩΝ ὝΜΝΟΣ.

\ch{44}
Αἰνέσωμεν δὴ ἄνδρας ἐνδόξους, καὶ τοὺς πατέρας ἡμῶν τῇ γενέσει.
\vs{2}Πολλὴν δόξαν ἔκτισεν ὁ Κύριος, τὴν μεγαλωσύνην αὐτοῦ ἀπʼ αἰῶνος.
\vs{3}Κυριεύοντες ἐν ταῖς βασιλείαις αὐτῶν, καὶ ἄνδρες ὀνομαστοὶ ἐν δυνάμει· βουλεύσονται ἐν συνέσει αὐτῶν, ἀπηγγελκότες ἐν προφητείαις·
\vs{4}ἡγούμενοι λαοῦ ἐν διαβουλίοις, καὶ συνέσει γραμματείας λαοῦ· σοφοὶ λόγοι ἐν παιδείᾳ αὐτῶν·
\vs{5}ἐκζητοῦντες μέλη μουσικῶν, διηγούμενοι ἔπη ἐν γραφῇ·
\vs{6}ἄνδρες πλούσιοι κεχορηγημένοι ἰσχύϊ, εἰρηνεύοντες ἐν παροικίαις αὐτῶν·
\vs{7}πάντες οὗτοι ἐν γενεαῖς ἐδοξάσθησαν, καὶ ἐν ταῖς ἡμέραις αὐτῶν καύχημα.
\vs{8}Εἰσὶν αὐτῶν οἳ κατέλιπον ὄνομα τοῦ ἐκδιηγήσασθαι ἐπαίνους·
\vs{9}καὶ εἰσὶν ὧν οὐκ ἔστι μνημόσυνον, καὶ ἀπώλοντο ὡς οὐχ ὑπάρξαντες, καὶ ἐγένοντο ὡς οὐ γεγονότες, καὶ τὰ τέκνα αὐτῶν μετʼ αὐτούς.

\vs{10}Ἀλλʼ ἢ οὗτοι ἄνδρες ἐλέους, ὧν αἱ δικαιοσύναι οὐκ ἐπελήσθησαν.
\vs{11}Μετὰ τοῦ σπέρματος αὐτῶν διαμενεῖ ἀγαθὴ κληρονομία, ἔκγονα αὐτῶν ἐν ταῖς διαθήκαις.
\vs{12}Ἔστη σπέρμα αὐτῶν καὶ τέκνα αὐτῶν διʼ αὐτούς·
\vs{13}ἕως αἰῶνος μενεῖ σπέρμα αὐτῶν, καὶ ἡ δόξα αὐτῶν οὐκ. ἐξαλειφθήσεται.
\vs{14}Τὸ σῶμα αὐτῶν ἐν εἰρήνῃ ἐτάφη, καὶ τὸ ὄνομα αὐτῶν ζῇ εἰς γενεάς.
\vs{15}Σοφίαν αὐτῶν διηγήσονται λαοὶ, καὶ τὸν ἔπαινον ἐξαγγέλλει ἐκκλησία.

\vs{16}Ἐνὼχ εὐηρέστησε Κυρίῳ, καὶ μετετέθη ὑπόδειγμα μετανοίας ταῖς γενεαῖς.
\vs{17}Νῶε εὑρέθη τέλειος δίκαιος, ἐν καιρῷ ὀργῆς ἐγένετο ἀντάλλαγμα· διὰ τοῦτο ἐγενήθη κατάλειμμα τῇ γῇ, διὰ τοῦτο ἐγένετο κατακλυσμός.
\vs{18}Διαθῆκαι αἰῶνος ἐτέθησαν πρὸς αὐτὸν, ἵνα μὴ ἐξαλειφθῇ κατακλυσμῷ πᾶσα σάρξ.

\vs{19}Ἁβραὰμ μέγας πατὴρ πλήθους ἐθνῶν, καὶ οὐχ εὑρέθη ὅμοιος ἐν τῇ δόξῃ,
\vs{20}ὃς συνετήρησε νόμον ὑψίστου, καὶ ἐγένετο ἐν διαθήκῃ μετʼ αὐτοῦ· καὶ ἐν σαρκὶ αὐτοῦ ἔστησε διαθήκην, καὶ ἐν πειρασμῷ εὑρέθη πιστός.
\vs{21}Διὰ τοῦτο ἐν ὅρκῳ ἔστησεν αὐτῷ, ἐνευλογηθῆναι ἔθνη ἐν τῷ σπέρματι αὐτοῦ, πληθύναι αὐτὸν ὡς χοῦν τῆς γῆς, καὶ ὡς ἄστρα ἀνυψῶσαι τὸ σπέρμα αὐτοῦ, καὶ κατακληρονομῆσαι αὐτοὺς ἀπὸ θαλάσσης ἕως θαλάσσης, καὶ ἀπὸ ποταμοῦ ἕως ἄκρου γῆς.

\vs{22}Καὶ ἐν τῷ Ἰσαὰκ ἔστησεν οὕτως διὰ Ἁβραὰμ τὸν πατέρα αὐτοῦ, εὐλογίαν πάντων ἀνθρώπων καὶ διαθήκην.
\vs{23}Καὶ κατέπαυσεν ἐπὶ κεφαλὴν Ἰακώβ· ἐπέγνω αὐτὸν ἐν εὐλογίαις αὐτοῦ, καὶ ἔδωκεν αὐτῷ ἐν κληρονομίᾳ· καὶ διέστειλε μερίδας αὐτοῦ, ἐν φυλαῖς ἐμέρισε δεκαδύο.

\ch{45}Καὶ ἐξήγαγεν ἐξ αὐτοῦ ἄνδρα ἐλέους, εὑρίσκοντα χάριν ἐν ὀφθαλμοῖς πάσης σαρκός· ἐγαπημένον ὑπὸ Θεοῦ καὶ ἀνθρώπων Μωυσῆν, οὗ τὸ μνημόσυνον ἐν εὐλογίαις.
\vs{2}Ὡμοίωσεν αὐτὸν δόξῃ ἁγίων, καὶ ἐμεγάλυνεν αὐτὸν ἐν φόβοις ἐχθρῶν.
\vs{3}Ἐν λόγοις αὐτοῦ σημεῖα κατέπαυσεν, ἐδόξασεν αὐτὸν κατὰ πρόσωπον βασιλέων· ἐνετείλατο αὐτῷ πρὸς λαὸν αὐτοῦ, καὶ ἔδειξεν αὐτῷ τῆς δόξης αὐτοῦ.
\vs{4}Ἐν πίστει καὶ πραΰτητι αὐτοῦ ἡγίασεν, ἐξελέξατο αὐτὸν ἐκ πάσης σαρκός.
\vs{5}Ἠκούτισεν αὐτὸν τῆς φωνῆς αὐτοῦ, καὶ εἰσήγαγεν αὐτὸν εἰς τὸν γνόφον· καὶ ἔδωκεν αὐτῷ κατὰ πρόσωπον ἐντολὰς, νόμον ζωῆς καὶ ἐπιστημης, διδάξαι τὸν Ἰακὼβ διαθήκην, καὶ κρίματα αὐτοῦ τὸν Ἰσραήλ.

\vs{6}Ἀαρὼν ὕψωσεν ἅγιον ὅμοιον αὐτῷ, ἀδελφὸν αὐτοῦ, ἐκ φυλῆς Λευΐ.
\vs{7}Ἔστησεν αὐτῷ διαθήκην αἰῶνος, καὶ ἔδωκεν αὐτῷ ἱερατείαν λαοῦ· ἐμακάρισεν αὐτὸν ἐν εὐκοσμίᾳ, καὶ περιέζωσεν αὐτὸν στολὴν δόξης.
\vs{8}Ἐνέδυσεν αὐτὸν συντέλειαν καυχήματος, καὶ ἐστερέωσεν αὐτὸν σκεύεσιν ἰσχύος, περισκελῆ καὶ ποδήρη καὶ ἐπωμίδα,
\vs{9}καὶ ἐκύκλωσεν αὐτὸν ῥοΐσκοις χρυσοῖς, κώδωσι πλείστοις κυκλόθεν, ἠχῆσαι φωνὴν ἐν βήμασιν αὐτοῦ, ἀκουστὸν ποιῆσαι ἦχον ἐν ναῷ εἰς μνημόσυνον υἱοῖς λαοῦ αὐτοῦ,
\vs{10}στολῇ ἁγίᾳ, χρυσῷ, καὶ ὑακίνθῳ, καὶ πορφύρᾳ, ἔργῳ ποικιλτοῦ, λογείῳ κρίσεως, δήλοις ἀληθείας,
\vs{11}κεκλωσμένῃ κόκκῳ, ἔργῳ τεχνίτου, λίθοις πολυτελέσι γλύμματος σφραγίδος, ἐν δέσει χρυσίου, ἔργῳ λιθουργοῦ, εἰς μνημόσυνον ἐν γραφῇ κεκολαμμένῃ κατʼ ἀριθμὸν φυλῶν Ἰσραήλ·
\vs{12}στέφανον χρυσοῦν ἐπάνω κιδάρεως, ἐκτύπωμα σφραγίδος ἁγιάσματος, καύχημα τιμῆς, ἔργον ἰσχύος, ἐπιθυμήματα ὀφθαλμῶν κοσμούμενα ὡραῖα.

\vs{13}Πρὸ αὐτοῦ οὐ γέγονε τοιαῦτα ἕως αἰῶνος, οὐκ ἐνεδύσατο ἀλλογενὴς, πλὴν τῶν υἱῶν αὐτοῦ μόνον, καὶ τὰ ἔκγονα αὐτοῦ διαπαντός.
\vs{14}Θυσίαι αὐτοῦ ὁλοκαρπωθήσονται καθημέραν ἐνδελεχῶς δίς.

\vs{15}Ἐπλήρωσε Μωυσῆς τὰς χεῖρας, καὶ ἔχρισεν αὐτὸν ἐν ἐλαίῳ ἁγίῳ· ἐγενήθη αὐτῷ εἰς διαθήκην αἰώνιον, καὶ ἐν τῷ σπέρματι αὐτοῦ ἐν ἡμέραις οὐρανοῦ, λειτουργεῖν αὐτῷ ἅμα καὶ ἱερατεύειν, καὶ εὐλογεῖν τὸν λαὸν αὐτοῦ ἐν τῷ ὀνόματι αὐτοῦ.

\vs{16}Ἐξελέξατο αὐτὸν ἀπὸ παντὸς ζῶντος, προσαγαγεῖν κάρπωσιν Κυρίῳ, θυμίαμα καὶ εὐωδίαν εἰς μνημόσυνον, ἐξιλάσκεσθαι περὶ τοῦ λαοῦ σου.
\vs{17}Ἔδωκεν αὐτὸν ἐν ἐντολαῖς αὐτοῦ, ἐξουσίαν ἐν διαθήκαις κριμάτων, διδάξαι τὸν Ἰακὼβ τὰ μαρτύρια, καὶ ἐν νόμῳ αὐτοῦ φωτίσαι Ἰσραήλ.

\vs{18}Ἐπισυνέστησαν αὐτῷ ἀλλότριοι, καὶ ἐζήλωσαν αὐτὸν ἐν τῇ ἐρήμῳ, ἄνδρες οἱ περὶ Δαθὰν καὶ Ἀβειρὼν, καὶ ἡ συναγωγὴ Κορὲ ἐν θυμῷ καὶ ὀργῇ.

\vs{19}Εἶδε Κύριος καὶ οὐκ εὐδόκησε, καὶ συνετελέσθησαν ἐν θυμῷ ὀργῆς· ἐποίησεν αὐτοῖς τέρατα, καταναλῶσαι ἐν πυρὶ φλογὸς αὐτοῦ.
\vs{20}Καὶ προσέθηκεν Ἀαρὼν δόξαν, καὶ ἔδωκεν αὐτῷ κληρονομίαν· ἀπαρχὰς πρωτογενημάτων ἐμέρισεν αὐτοῖς· ἄρτον ἐν πρώτοις ἡτοίμασε πλησμονήν.
\vs{21}Καὶ γὰρ θυσίας Κυρίου φάγονται, ἃς ἔδωκεν αὐτῷ τε καὶ τῷ σπέρματι αὐτοῦ.
\vs{22}Πλὴν ἐν γῇ λαοῦ οὐ κληρονομήσει, καὶ μερὶς οὐκ ἔστιν αὐτῷ ἐν λαῷ, αὐτὸς γὰρ μερίς σου, κληρονομία.

\vs{23}Καὶ Φινεὲς υἱὸς Ἐλεάζαρ τρίτος εἰς δόξαν, ἐν τῷ ζηλῶσαι αὐτὸν ἐν φόβῳ Κυρίου, καὶ στῆσαι αὐτὸν ἐν τροπῇ λαοῦ, ἐν ἀγαθότητι προθυμίας ψυχῆς αὐτοῦ, καὶ ἐξιλάσατο περὶ τοῦ Ἰσραήλ.
\vs{24}Διὰ τοῦτο ἐστάθη αὐτῷ διαθήκη εἰρήνης, προστάτην ἁγίων καὶ λαῷ αὐτοῦ, ἵνα αὐτῷ ᾖ καὶ τῷ σπέρματι αὐτοῦ ἱερωσύνης μεγαλεῖον εἰς τοὺς αἰῶνας·
\vs{25}καὶ διαθηκήν τῷ Δαυὶδ υἱῷ ἐκ φυλῆς Ἰούδα, κληρονομία βασιλέως υἱοῦ ἐξ υἱοῦ μόνου, κληρονομία Ἀαρὼν καὶ τῷ σπέρματι αὐτοῦ.
\vs{26}Δῴη ὑμῖν σοφίαν ἐν καρδίᾳ ὑμῶν, κρίνειν τὸν λαὸν αὐτοῦ ἐν δικαιοσύνῃ, ἵνα μὴ ἀφανισθῇ τὰ ἀγαθὰ αὐτῶν, καὶ τὴν δόξαν αὐτῶν εἰς γενεὰς αὐτῶν.

\ch{46}
Κράταιος ἐν πολέμοις Ἰησοῦς Ναυῆ, καὶ διάδοχος Μωυσῆ ἐν προφητείαις· ὃς ἐγένετο κατὰ τὸ ὄνομα αὐτοῦ μέγας ἐπὶ σωτηρίᾳ ἐκλεκτῶν αὐτοῦ, ἐκδικῆσαι ἐπεγειρομένους ἐχθροὺς, ὅπως κληρονομήσῃ τὸν Ἰσραήλ.
\vs{2}Ὡς ἐδοξάσθη ἐν τῷ ἐπᾶραι χεῖρας αὐτοῦ, καὶ τῷ ἐκκλῖναι ῥομφαίαν ἐπὶ πόλεις;
\vs{3}Τίς πρότερον αὐτοῦ οὕτως ἔστη; τοὺς γὰρ πολεμίους Κύριος αὐτὸς ἐπήγαγεν.
\vs{4}Οὐχὶ ἐν χειρὶ αὐτοῦ ἀνεπόδισεν ὁ ἥλιος, καὶ μία ἡμέρα ἐγενήθη πρὸς δύο;
\vs{5}Ἐπεκαλέσατο τὸν ὕψιστον δυνάστην, ἐν τῷ θλίψαι αὐτὸν ἐχθροὺς κυκλόθεν· καὶ ἐπήκουσεν αὐτῶν μέγος Κύριος.
\vs{6}Ἐν λίθοις χαλάζης δυνάμεως κραταιᾶς· κατέῤῥαξεν ἐπʼ ἔθνος πόλεμον, καὶ ἐν καταβάσει ἀπώλεσεν ἀνθεστηκότας· ἵνα γνῶσιν ἔθνη πανοπλίαν αὐτῶν, ὅτι ἐναντίον Κυρίου ὁ πόλεμος αὐτοῦ, καὶ γὰρ ἐπηκολούθησεν ὀπίσω δυνάστου.

\vs{7}Καὶ ἐν ἡμέραις Μωυσέως ἐποίησεν ἔλεος, αὐτὸς καὶ Χαλὲβ υἱὸς Ἰεφοννῆ, ἀντιστῆναι ἔναντι ἐχθροῦ, κωλῦσαι λαὸν ἀπὸ ἁμαρτίας, καὶ κοπάσαι γογγυσμὸν πονηρίας.
\vs{8}Καὶ αὐτοὶ δύο ὄντες διεσώθησαν ἀπὸ ἑξακοσίων χιλιάδων πεζῶν, εἰσαγαγεῖν αὐτοὺς εἰς κληρονομίαν, εἰς γῆν ῥέουσαν γάλα καὶ μέλι.

\vs{9}Καὶ ἔδωκεν ὁ Κύριος τῷ Χαλὲβ ἰσχὺν, καὶ ἕως γήρους διέμεινεν αὐτῷ, ἐπιβῆναι αὐτὸν ἐπὶ ὕψος τῆς γῆς, καὶ τὸ σπέρμα αὐτοῦ κατέσχε κληρονομίαν·
\vs{10}ὅπως ἴδωσι πάντες οἱ υἱοὶ Ἰσραὴλ, ὅτι καλὸν τὸ πορεύεσθαι ὀπίσω Κυρίου.
\vs{11}Καὶ οἱ κριταὶ ἕκαστος τῷ αὐτοῦ ὀνόματι, ὅσων οὐκ ἐξεπόρνευσεν ἡ καρδία, καὶ ὅσοι οὐκ ἀπεστράφησαν ἀπὸ Κυρίου, εἴη τὸ μνημόσυνον αὐτῶν ἐν εὐλογίαις·
\vs{12}τὰ ὀστᾶ αὐτῶν ἀναθάλοι ἐκ τοῦ τόπου αὐτῶν, καὶ τὸ ὄνομα αὐτῶν ἀντικαταλλασσόμένον ἐφ υἱοῖς δεδοξασμένων αὐτῶν.

\vs{13}Ἠγαπημένος ὑπὸ Κυρίου αὐτοῦ Σαμουὴλ προφήτης Κυρίου κατέστησε βασιλείαν, καὶ ἔχρισεν ἄρχοντας ἐπὶ τὸν λαὸν αὐτοῦ.
\vs{14}Ἐν νόμῳ Κυρίου ἔκρινε συναγωγὴν, καὶ ἐπεσκέψατο Κύριος τὸν Ἰακώβ.
\vs{15}Ἐν πίστει αὐτοῦ ἠκριβάσθη προφήτης, καὶ ἐγνώσθη ἐν πίστει αὐτοῦ πιστὸς ὁράσεως.
\vs{16}Καὶ ἐπεκαλέσατο τὸν Κύριον δυνάστην, ἐν τῷ θλίψαι ἐχθροὺς αὐτοῦ κυκλόθεν, ἐν προσφορᾷ ἀρνὸς γαλαθηνοῦ.
\vs{17}Καὶ ἐβρόντησεν ἀπʼ οὐρανοῦ Κύριος, καὶ ἐν ἤχῳ μεγάλῳ ἀκουστὴν ἐποίησε τὴν φωνὴν αὐτοῦ.
\vs{18}Καὶ ἐξέτριψεν ἡγουμένους Τυρίων, καὶ πάντας ἄρχοντας Φυλιστιείμ.

\vs{19}Καὶ πρὸ καιροῦ κοιμήσεως αἰῶνος ἐπεμαρτύρατο ἔναντι Κυρίου καὶ χριστοῦ, χρήματα καὶ ἕως ὑποδημάτων ἀπὸ πάσης σαρκὸς οὐκ εἴληφα· καὶ οὐκ ἐνεκάλεσεν αὐτῷ ἄνθρωπος.
\vs{20}Καὶ μετὰ τὸ ὑπνῶσαι αὐτὸν προεφήτευσε, καὶ ὑπέδειξε βασιλεῖ τὴν τελευτὴν αὐτοῦ, καὶ ἀνύψωσεν ἐκ γῆς τὴν φωνὴν αὐτοῦ, ἐν προφητείᾳ ἐξαλεῖψαι ἀνομίαν λαοῦ.

\ch{47}
Καὶ μετὰ τοῦτο ἀνέστη Νάθαν προφητεύειν ἐν ἡμέραις Δαυίδ.

\vs{2}Ὥσπερ στέαρ ἀφωρισμένον ἀπὸ σωτηρίου, οὕτως Δαυὶδ ἀπὸ τῶν υἱῶν Ἰσραήλ.
\vs{3}Ἐν λέουσιν ἔπαισεν ὡς ἐν ἐρίφοις, καὶ ἐν ἄρκοις ὡς ἐν ἄρνασι προβάτων.
\vs{4}Ἐν νεότητι αὐτοῦ οὐχὶ ἀπέκτεινε γίγαντα, καὶ ἐξῇρεν ὀνειδισμὸν ἐκ λαοῦ, ἐν τῷ ἐπᾶραι χεῖρα ἐν λίθῳ σφενδόνης, καὶ καταβαλεῖν γαυρίαμα τοῦ Γολιάθ;
\vs{5}Ἐπεκαλέσατο γὰρ Κύριον τὸν ὕψιστον, καὶ ἔδωκεν ἐν τῇ δεξιᾷ αὐτοῦ κράτος ἐξᾶραι ἄνθρωπον δυνατὸν ἐν πολέμῳ, ἀνυψῶσαι κέρας λαοῦ αὐτοῦ.

\vs{6}Οὕτως ἐν μυριάσιν ἐδόξασεν αὐτὸν, καὶ ᾔνεσεν αὐτὸν ἐν εὐλογίαις Κυρίου, ἐν τῷ φέρεσθαι αὐτῷ διάδημα δόξης.
\vs{7}Ἐξέτριψε γὰρ ἐχθροὺς κυκλόθεν, καὶ ἐξουδένωσε Φυλιστιεὶμ τοὺς ὑπεναντίους· ἕως σήμερον συνετριψεν αὐτῶν κέρας.
\vs{8}Ἐν παντὶ ἔργῳ αὐτοῦ ἔδωκεν ἐξομολόγησιν· ἁγίῳ ὑψίστῳ ῥήματι δόξης ἐν πάσῃ καρδίᾳ αὐτοῦ ὕμνησε, καὶ ἠγάπησε τὸν ποιήσαντα αὐτόν.
\vs{9}Καὶ ἔστησε ψαλτῳδοὺς κατέναντι τοῦ θυσιαστηρίου, καὶ ἐξ ἤχου αὐτῶν γλυκαίνει μέλη.
\vs{10}Ἔδωκεν ἐν ἐορταῖς εὐπρέπειαν, καὶ ἐκόσμησε καιροὺς μέχρι συντελείας· ἐν τῷ αἰνεῖν αὐτοὺς τὸ ἅγιον ὄνομα αὐτοῦ, καὶ ἀπὸ πρωῒ ἠχεῖν τὸ ἁγίασμα.

\vs{11}Κύριος ἀφεῖλε τὰς ἁμαρτίας ἀυτοῦ, καὶ ἀνύψωσεν εἰς αἰῶνα τὸ κέρας αὐτοῦ, καὶ ἔδωκεν αὐτῷ διαθήκην βασιλέων καὶ θρόνον δόξης ἐν τῷ Ἰσραήλ.
\vs{12}Μετὰ τούτου ἀνέστη υἱὸς ἐπιστήμων, καὶ διʼ αὐτὸν κατέλυσεν ἐν πλατυσμῷ.
\vs{13}Σαλωμὼν ἐβασίλευσεν ἐν ἡμέραις εἰρήνης, ᾧ ὁ Θεὸς κατέπαυσε κυκλόθεν, ἵνα στήσῃ οἶκον ἐπʼ ὀνόματι αὐτοῦ, καὶ ἑτοιμάσῃ ἁγίασμα εἰς τὸν αἰῶνα.
\vs{14}Ὡς ἐσοφίσθης ἐν νεότητί σου, καὶ ἐνεπλήσθης ὡς ποταμὸς συνέσεως.
\vs{15}Γῆν ἐπεκάλυψεν ἡ ψυχή σου, καὶ ἐνέπλησας ἐν παραβολαῖς αἰνιγμάτων.

\vs{16}Εἰς νήσους πόῤῥω ἀφίκετο τὸ ὄνομά σου, καὶ ἠγαπήθης ἐν τῇ εἰρήνῃ σου.
\vs{17}Ἐν ᾠδαῖς καὶ παροιμίαις καὶ παραβολαῖς, καὶ ἐν ἑρμηνείαις ἀπεθαύμασάν σε χῶραι.
\vs{18}Ἐν ὀνόματι Κυρίου τοῦ Θεοῦ τοῦ ἐπικεκλημένου Θεοῦ Ἰσραὴλ, συνήγαγες ὡς κασσίτερον τὸ χρυσίον, καὶ ὡς μόλιβδον ἐπλήθυνας ἀργύριον.
\vs{19}Παρενέκλινας τὰς λαγόνας σου γυναιξὶ, καὶ ἐνεξουσιάσθης ἐν τῷ σώματί σου.
\vs{20}Ἔδωκας μῶμον ἐν τῇ δόξῃ σου, καὶ ἐβεβήλωσας τὸ σπέρμα σου, ἐπαγαγεῖν ὀργὴν ἐπὶ τὰ τέκνα σου, καὶ κατενύγην ἐπὶ τῇ ἀφροσύνῃ σου,
\vs{21}γενέσθαι δίχα τυραννίδα, καὶ ἐξ Ἐφραὶμ ἄρξαι βασιλείαν ἀπειθῆ.

\vs{22}Ὁ δὲ Κύριος οὐ μὴ καταλίπῃ τὸ ἔλεος αὐτοῦ, καὶ οὐ μη διαφθαρῇ ἀπὸ τῶν ἔργων αὐτοῦ· οὐδὲ μὴ ἐξαλείψῃ ἐκλεκτοῦ ἔκγονα, καὶ σπέρμα τοῦ ἀγαπήσαντος αὐτὸν οὐ μὴ ἐξάρῃ· καὶ τῷ Ἰακὼβ ἔδωκε κατάλειμμα, καὶ τῷ Δαυὶδ ἐξ αὐτοῦ ῥίζαν.

\vs{23}Καὶ ἀνεπαύσατο Σαλωμὼν μετὰ τῶν πατέρων· καὶ κατέλιπε μετʼ αὐτὸν ἐκ τοῦ σπέρματος αὐτοῦ, λαοῦ ἀφροσύνην καὶ ἐλασσούμενον συνέσει, Ῥοβοὰμ, ὃς ἀπέστησε λαὸν ἐκ βουλῆς αὐτοῦ· καὶ Ἱεροβοὰμ υἱὸν Ναβὰτ, ὃς ἐξήμαρτε τὸν Ἰσραὴλ, καὶ ἔδωκε τῷ Ἐφραὶμ ὁδὸν ἁμαρτίας.
\vs{24}Καὶ ἐπληθύνθησαν αἱ ἁμαρτίαι αὐτῶν σφόδρα, ἀποστῆσαι αὐτοὺς ἀπὸ τῆς γῆς αὐτῶν.
\vs{25}Καὶ πᾶσαν πονηρίαν ἐξεζήτησαν, ἕως ἐκδίκησις ἔλθῃ ἐπʼ αὐτούς.

\ch{48}
Καὶ ἀνέστη Ἠλίας προφήτης ὡς πῦρ, καὶ ὁ λόγος αὐτοῦ ὡς λαμπὰς ἐκαίετο·
\vs{2}ὃς ἐπήγαγεν ἐπʼ αὐτοὺς λιμὸν, καὶ τῷ ζήλῳ αὐτοῦ ὠλιγοποίησεν αὐτούς.
\vs{3}Ἐν λόγῳ Κυρίου ἀνέσχεν οὐρανὸν, κατήγαγεν οὕτως τρὶς πῦρ.
\vs{4}Ὡς ἐδοξάσθης Ἠλία ἐν τοῖς θαυμασίοις σου; καὶ τίς ὅμοιός σοι καυχᾶσθαι;
\vs{5}Ὁ ἐγείρας νεκρὸν ἐκ θυνάτου καὶ ἐξ ᾅδου ἐν λόγῳ ὑψίστου·
\vs{6}ὁ καταγαγὼν βασιλεῖς εἰς ἀπώλειαν, καὶ δεδοξασμένους ἀπὸ κλίνης αὐτῶν·
\vs{7}ὁ ἀκούων ἐν Σινᾷ ἐλεγμὸν, καὶ ἐν Χωρὴβ κρίματα ἐκδικήσεως·
\vs{8}ὁ χρίων βασιλεῖς εἰς ἀνταπόδομα, καὶ προφήτας διαδόχους μετʼ αὐτόν·
\vs{9}ὁ ἀναληφθεὶς ἐν λαίλαπι πυρὸς ἐν ἅρματι ἵππων πυρίνων·
\vs{10}ὁ καταγραφεὶς ἐν ἐλεγμοῖς εἰς καιροὺς, κοπάσαι ὀργὴν πρὸ θυμοῦ, καὶ ἐπιστρέψαι καρδίαν πατρὸς πρὸς υἱὸν, καὶ καταστῆσαι φυλὰς Ἰακώβ.
\vs{11}Μακάριοι οἱ ἰδόντες σε, καὶ οἱ ἐν ἀγαπήσει κεκοσμημένοι· καὶ γὰρ ἡμεῖς ζωῇ ζησόμεθα.

\vs{12}Ἠλίας, ὃς ἐν λαίλαπι ἐσκεπάσθη· καὶ Ἐλισαιὲ ἐνεπλήσθη πνεύματος αὐτοῦ· καὶ ἐν ἡμέραις αὐτοῦ οὐκ ἐσαλεύθη ὑπὸ ἄρχοντος, καὶ οὐ κατεδυνάστευσεν αὐτὸν οὐδείς.
\vs{13}Πᾶς λόγος οὐχ ὑπερῇρεν αὐτὸν, καὶ ἐν κοιμήσει ἐπροφήτευσε τὸ σῶμα αὐτοῦ.
\vs{14}Καὶ ἐν ζωῇ αὐτοῦ ἐποίησε τέρατα. καὶ ἐν τελευτῇ θαυμάσια τὰ ἔργα αὐτοῦ.

\vs{15}Ἐν πᾶσι τούτοις οὐ μετενόησεν ὁ λαὸς, καὶ οὐκ ἀπέστησαν ἀπὸ τῶν ἁμαρτιῶν, ἕως ἐπρονομεύθησαν ἀπὸ τῆς γῆς αὐτῶν, καὶ ἐσκορπίσθησαν ἐν πάσῃ τῇ γῇ· καὶ κατελείφθη ὁ λαὸς ὀλιγοστὸς, καὶ ἄρχων τῷ οἴκῳ Δαυίδ.
\vs{16}Τινὲς μὲν αὐτῶν ἐποίησαν τὸ ἀρεστὸν, τινὲς δὲ ἐπλήθυναν ἁμαρτίας.

\vs{17}Ἐζεκίας ὠχύρωσε τὴν πόλιν αὐτοῦ, καὶ εἰσήγαγεν εἰς μέσον αὐτῶν τὸν Γώγ· ὤρυξε σιδήρῳ ἀκρότομον, καὶ ᾠκοδόμησε κρήνας εἰς ὕδατα.
\vs{18}Ἐν ἡμέραις αὐτοῦ ἀνέβη Σενναχηρὶμ, καὶ ἀπέστειλε Ῥαψάκην, καὶ ἀπῇρε· καὶ ἐπῇρεν ἡ χεὶρ αὐτοῦ ἐπὶ Σιὼν, καὶ ἐμεγαλαύχησεν ὑπερηφανίᾳ αὐτοῦ.
\vs{19}Τότε ἐσαλεύθησαν καρδίαι καὶ χεῖρες αὐτῶν, καὶ ὠδίνησαν ὡς αἱ τίκτουσαι.

\vs{20}Καὶ ἐπεκαλέσαντο τὸν Κύριον τὸν ἐλεήμονα, ἐκπετάσαντες τὰς χεῖρας αὐτῶν πρὸς αὐτόν· καὶ ὁ ἅγιος ἐξ οὐρανοῦ ταχὺ ἐπήκουσεν αὐτῶν, καὶ ἐλυτρώσατο αὐτοὺς ἐν χειρὶ Ἡσαΐου.
\vs{21}Ἐπάταξε τὴν παρεμβολὴν τῶν Ἀσσυρίων, καὶ ἐξέτριψεν αὐτοὺς ὁ ἄγγελος αὐτοῦ.
\vs{22}Ἐποίησε γὰρ Ἐζεκίας τὸ ἀρεστὸν Κυρίῳ, καὶ ἐνίσχυσεν ἐν ὁδοῖς Δαυὶδ τοῦ πατρὸς αὐτοῦ, ἃς ἐνετείλατο Ἡσαΐας ὁ προφήτης ὁ μέγας, καὶ πιστὸς ἐν ὁράσει αὐτοῦ.
\vs{23}Ἐν ταῖς ἡμέραις αὐτοῦ ἀνεπόδισεν ὁ ἥλιος, καὶ προσέθηκε ζωὴν βασιλεῖ·
\vs{24}Πνεύματι μεγάλῳ εἶδε τὰ ἔσχατα, καὶ παρεκάλεσε τοὺς πενθοῦντας ἐν Σιών.
\vs{25}Ἕως τοῦ αἰῶνος ὑπέδειξε τὰ ἐσόμενα, καὶ τὰ ἀπόκρυφα πρινὴ παραγενέσθαι αὐτά.

\ch{49}
Μνημόσυνον Ἰωσίου εἰς σύνθεσιν θυμιάματος, ἐσκευασμένον ἔργῳ μυρεψοῦ, ἐν παντὶ στόματι ὡς μέλι γλυκανθήσεται, καὶ ὡς μουσικὰ ἐν συμποσίῳ οἴνου.
\vs{2}Αὐτὸς κατευθύνθη ἐν ἐπιστροφῇ λαοῦ, καὶ ἐξῇρε βδελύγματα ἀνομίας.
\vs{3}Κατεύθυνε πρὸς Κύριον τὴν καρδίαν αὐτοῦ, ἐν ἡμέραις ἀνόμων κατίσχυσε τὴν εὐσέβειαν.
\vs{4}Πάρεξ Δαυὶδ, καὶ Ἐζεκίου, καὶ Ἰωσίου, πάντες πλημμέλειαν ἐπλημμέλησαν· κατέλιπον γὰρ τὸν νόμον τοῦ ὑψίστου, οἱ βασιλεῖς Ἰούδα ἐξέλιπον.
\vs{5}Ἔδωκαν γὰρ τὸ κέρας αὐτῶν ἑτέροις, καὶ τὴν δόξαν αὐτῶν ἔθνει ἀλλοτρίῳ.

\vs{6}Ἐνεπύρισαν ἐκλεκτὴν πόλιν ἁγιάσματος, καὶ ἠρήμωσαν τὰς ὁδοὺς αὐτῆς ἐν χειρὶ Ἱερεμίου.
\vs{7}Ἐκάκωσαν γὰρ αὐτὸν, καὶ αὐτὸς ἐν μήτρᾳ ἡγιάσθη προφήτης ἐκριζοῦν καὶ κακοῦν καὶ ἀπολλύειν, ὡσαύτως οἰκοδομεῖν καὶ καταφυτεύειν.
\vs{8}Ἰεζεκιὴλ ὃς εἶδεν ὅρασιν δόξης, ἣν ὑπέδειξεν αὐτῷ ἐπὶ ἅρματος χερουβίμ.
\vs{9}Καὶ γὰρ ἐμνήσθη τῶν ἐχθρῶν ἐν ὄμβρῳ, καὶ ἀγαθῶσαι τοὺς εὐθύνοντας ὁδούς.
\vs{10}Καὶ τῶν δώδεκα προφητῶν τὰ ὀστᾶ ἀναθάλοι ἐκ τοῦ τόπου αὐτῶν· παρεκάλεσε δὲ τὸν Ἰακὼβ, καὶ ἐλυτρώσατο αὐτοὺς ἐν πίστει ἐλπίδος.
\vs{11}Πῶς μεγαλύνωμεν τὸν Ζοροβάβελ; καὶ αὐτὸς ὡς σφραγὶς ἐπὶ δεξιᾶς χειρός.

\vs{12}Οὕτως Ἰησοῦς υἱὸς Ἰωσεδέκ· οἳ ἐν ἡμέραις αὐτῶν ᾠκοδόμησαν οἴκον, καὶ ἀνύψωσαν λαὸν ἅγιον Κυρίῳ ἡτοιμασμένον εἰς δόξαν αἰῶνος.
\vs{13}Καὶ Νεεμίου ἐπὶ πολὺ τὸ μνημόσυνον, τοῦ ἐγείραντος ἡμῖν τείχη πεπτωκότα, καὶ στήσαντος πύλας καὶ μοχλοὺς, καὶ ἀνεγείραντος τὰ οἰκόπεδα ἡμῶν.
\vs{14}Οὐδὲ εἷς ἐκτίσθη οἷος Ἐνὼχ τοιοῦτος ἐπὶ τῆς γῆς, καὶ γὰρ αὐτὸς ἀνελήφθη ἀπὸ τῆς γῆς.
\vs{15}Οὐδὲ ὡς Ἰωσὴφ ἐγεννήθη ἀνὴρ, ἡγούμενος ἀδελφῶν, στήριγμα λαοῦ, καὶ τὰ ὀστᾶ αὐτοῦ ἐπεσκέπησαν.
\vs{16}Σὴμ καὶ Σὴθ ἐν ἀνθρώποις ἐδοξάσθησαν, καὶ ὑπὲρ πᾶν ζῶον ἐν τῇ κτίσει Ἀδάμ.

\ch{50}
Σίμων Ὀνίου υἱὸς ἱερεὺς ὁ μέγας, ὃς ἐν ζωῇ αὐτοῦ ὑπέῤῥαψεν οἶκον, καὶ ἐν ἡμέραις αὐτοῦ ἐστερέωσε τὸν ναόν·
\vs{2}καὶ ὑπʼ αὐτοῦ ἐθεμελιώθη ὕψος διπλῆς ἀνάλημμα ὑψηλὸν περιβόλου ἱεροῦ.
\vs{3}Ἐν ἡμέραις αὐτοῦ ἠλαττώθη ἀποδοχεῖον ὑδάτων, χαλκὸς ὡσεὶ θαλάσσης τὸ περίμετρον·
\vs{4}ὁ φροντίζων τοῦ λαοῦ αὐτοῦ ἀπὸ πτώσεως, καὶ ἐνισχύσας πόλιν ἐμπολιορκῆσαι,
\vs{5}ὡς ἐδοξάσθη ἐν περιστροφῇ λαοῦ, ἐν ἐξόδῳ οἴκου καταπετάσματος·
\vs{6}ὡς ἀστὴρ ἑωθινὸς ἐν μέσῳ νεφέλης, ὡς σελήνη πλήρης ἐν ἡμέραις·
\vs{7}ὡς ἥλιος ἐκλάμπων ἐπὶ ναὸν ὑψίστου, καὶ ὡς τόξον φωτίζον ἐν νεφέλαις δόξης·
\vs{8}ὡς ἄνθος ῥόδων ἐν ἡμέραις νέων, ὡς κρίνα ἐπʼ ἐξόδων ὕδατος· ὡς βλαστὸς Λιβάνου ἐν ἡμέραις θέρους,
\vs{9}ὡς πῦρ καὶ λίβανος ἐπὶ πυρείου· ὡς σκεῦος χρυσίου ὁλοσφύρητον κεκοσμημένον παντὶ λίθῳ πολυτελεῖ·
\vs{10}ὡς ἐλαία ἀναθάλλουσα καρποὺς, καὶ ὡς κυπάρισσος ὑψουμένη ἐν νεφέλαις.

\vs{11}Ἐν τῷ ἀναλαμβάνειν αὐτὸν στολὴν δόξης, καὶ ἐνδιδύσκεσθαι αὐτὸν συντέλειαν καυχήματος, ἐν ἀναβάσει θυσιαστηρίου ἁγίου ἐδόξασε περιβολὴν ἁγιάσματος.
\vs{12}Ἐν δὲ τῷ δέχεσθαι μέλη ἐκ χειρῶν ἱερέων, καὶ αὐτὸς ἑστὼς παρʼ ἐσχάρᾳ βωμοῦ, κυκλόθεν αὐτοῦ στέφανος ἀδελφῶν, ὡς βλάστημα κέδρου ἐν τῶ λιβάνῳ· καὶ ἐκύκλωσαν αὐτὸν ὡς στελέχη φοινίκων,
\vs{13}καὶ πάντες οἱ υἱοὶ Ἀαρὼν ἐν δόξῃ αὐτῶν· καὶ προσφορὰ Κυρίου ἐν χερσὶν αὐτῶν ἔναντι πάσης ἐκκλησίας Ἰσραήλ.
\vs{14}Καὶ συντέλειαν λειτουργῶν ἐπὶ βωμῶν, κοσμῆσαι προσφορὰν ὑψίστου παντοκράτορος,
\vs{15}ἐξέτεινεν ἐπὶ σπονδείου χεῖρα αὐτοῦ, καὶ ἔσπεισεν ἐξ αἵματος σταφυλῆς· ἐξέχεεν εἰς θεμέλια θυσιαστηρίου ὀσμὴν εὐωδίας ὑψίστῷ παμβασιλεῖ.

\vs{16}Τότε ἀνέκραγον υἱοὶ Ἀαρὼν, ἐν σάλπιγξιν ἐλαταῖς ἤχησαν· ἀκουστὴν ἐποίησαν φωνὴν μεγάλην εἰς μνημόσυνον ἔναντι ὑψίστου.
\vs{17}Τότε πᾶς ὁ λαὸς κοινῇ κατέσπευσε, καὶ ἔπεσαν ἐπὶ πρόσωπον ἐπὶ τὴν γῆν, προσκυνῆσαι τῷ Κυρίῳ αὐτῶν παντοκράτορι Θεῷ τῷ ὑψίστῳ.
\vs{18}Καὶ ᾔνεσαν οἱ ψαλμῳδοὶ ἐν φωναῖς αὐτῶν, ἐν πλείστῳ οἴκῳ ἐγλυκάνθη μέλος.
\vs{19}Καὶ ἐδεήθη ὁ λαὸς Κυρίου ὑψίστου ἐν προσευχῇ κατέναντι ἐλεήμονος, ἕως συντελεσθῇ κόσμος Κυρίου, καὶ τὴν λειτουργίαν αὐτοῦ ἐτελείωσαν.

\vs{20}Τότε καταβὰς ἐπῇρε χεῖρας αὐτοῦ ἐπὶ πᾶσαν ἐκκλησίαν υἱῶν Ἰσραὴλ, δοῦναι εὐλογίαν Κυρίῳ ἐν χειλέων αὐτοῦ, καὶ ἐν ὀνόματι αὐτοῦ καυχᾶσθαι.
\vs{21}Καὶ ἐδευτέρωσεν ἐν προσκυνήσει ἐπιδείξασθαι τὴν εὐλογίαν παρὰ ὑψίστου.
\vs{22}Καὶ νῦν εὐλογήσατε τῷ Θεῷ πάντες τῷ μεγαλοποιοῦντι πάντη, τὸν ὑψοῦντα ἡμέρας ἡμῶν ἐκ μήτρας, καὶ ποιοῦντα μεθʼ ἡμῶν κατὰ τὸ ἔλεος αὐτοῦ.
\vs{23}Δῴη ἡμῖν εὐφροσύνην καρδίας, καὶ γενέσθαι εἰρήνην ἐν ἡμέραις ἡμῶν ἐν Ἰσραὴλ κατὰ τὰς ἡμέρας τοῦ αἰῶνος,
\vs{24}ἐμπιστεῦσαι μεθʼ ἡμῶν τὸ ἔλεος αὐτοῦ, καὶ ἐν ταῖς ἡμέραις αὐτοῦ λυτρωσάσθω ἡμᾶς.

\vs{25}Ἐν δυσὶν ἔθνεσι προσώχθισεν ἡ ψυχή μου, καὶ τὸ τρίτον οὐκ ἔστιν ἔθνος.
\vs{26}Οἱ καθήμενοι ἐν ὄρει Σαμαρείας, Φυλιστιεὶμ καὶ ὁ λαὸς μωρὸς ὁ κατοικῶν ἐν Σικίμοις.

\vs{27}Παιδείαν συνέσεως καὶ ἐπιστήμης ἐχάραξε ἐν τῷ βιβλίῳ, τούτῳ, Ἰησοῦς υἱὸς Σειρὰχ Ἱεροσολυμίτης, ὃς ἀνώμβρησε σοφίαν ἀπὸ καρδίας αὐτοῦ.
\vs{28}Μακάριος ὃς ἐν τούτοις ἀναστραφήσεται, καὶ ὁ θεὶς αὐτὰ ἐπὶ καρδίαν αὐτοῦ σοφισθήσεται.
\vs{29}Ἐὰν γὰρ αὐτὰ ποιήσῃ, πρὸς πάντα ἰσχύσει, ὅτι φῶς Κυρίου τὸ ἴχνος αὐτοῦ.

ΠΡΟΣΕΥΧΗ ἸΗΣΟΥ ΥΙΟΥ ΣΕΙΡΑΧ.

\ch{51}
Ἐξομολογήσομαι σοι Κύριε βασιλεῦ, καὶ αἰνέσω σε Θεὸν τὸν σωτῆρά μου· ἐξομολογοῦμαι τῷ ὀνόματί σου,
\vs{2}ὅτι σκεπαστὴς καὶ βοηθὸς ἐγένου μοι, καὶ ἐλυτρώσω τὸ σῶμά μου ἐξ ἀπωλείας, καὶ ἐκ παγίδος διαβολῆς γλώσσης· ἀπὸ χειλέων ἐργαζομένων ψεῦδος, καὶ ἔναντι τῶν παρεστηκότων ἐγένου μοι βοηθὸς,
\vs{3}καὶ ἐλυτρώσω με, κατὰ τὸ πλῆθος ἐλέους καὶ ὀνόματός σου, ἐκ βρυγμῶν ἑτοίμων εἰς βρῶμα, ἐκ χειρὸς ζητούντων τὴν ψυχήν μου, ἐκ πλειόνων θλίψεων ὧν ἔσχον,
\vs{4}ἀπὸ πνιγμοῦ πυρὸς κυκλόθεν, καὶ ἐκ μέσου πυρὸς οὗ οὐκ ἐξέκαυσα,
\vs{5}ἐκ βάθους κοιλίας ᾅδου, καὶ ἀπὸ γλώσσης ἀκαθάρτου, καὶ λόγου ψευδοῦς.

\vs{6}Βασιλεῖ διαβολὴ γλώσσης ἀδίκου· ᾔγγισεν ἕως θανάτου ἡ ψυχή μου, καὶ ἡ ζωή μου ἦν σύνεγγυς ᾅδου κάτω.
\vs{7}Περιέσχον με πάντοθεν, καὶ οὐκ ἦν ὁ βοηθῶν· ἐμβλέπων εἰς ἀντίληψιν ἀνθρώπων, καὶ οὐκ ἦν.
\vs{8}Καὶ ἐμνήσθην τοῦ ἐλέους σου Κύριε, καὶ τῆς ἐργασίας σου τῆς ἀπʼ αἰῶνος· ὅτι ἐξελῇ τοὺς ὑπομένοντάς σε, καὶ σώζεις αὐτοὺς ἐκ χειρὸς ἐθνῶν.
\vs{9}Καὶ ἀνύψωσα ἐπὶ γῆς ἱκετείαν μου, καὶ ὑπὲρ θανάτου ῥύσεως ἐδεήθην.
\vs{10}Ἐπεκαλεσάμην Κύριον πατέρα Κυρίου μου, μή με ἐγκαταλιπεῖν ἐν ἡμέραις θλίψεως, ἐν καιρῷ ὑπερηφάνων ἀβοηθησίας.
\vs{11}Αἰνέσω τὸ ὄνομά σου ἐνδελεχῶς, καὶ ὑμνήσω ἐν ἐξομολογήσει· καὶ εἰσηκούσθη ἡ δέησίς μου.
\vs{12}Ἔσωσας γάρ με ἐξ ἀπωλείας, καὶ ἐξείλου με ἐκ καιροῦ πονηροῦ· διὰ τοῦτο ἐξομολογήσομαι καὶ αἰνέσω σοι, καὶ εὐλογήσω τῷ ὀνόματι Κυρίου·

\vs{13}Ἔτι ὢν νεώτερος, πρινὴ πλανηθῆναί με, ἐζήτησα σοφίαν προφανῶς ἐν προσευχῇ μου.
\vs{14}Ἔναντι ναοῦ ἠξίουν περὶ αὐτῆς, καὶ ἕως ἐσχάτων ἐκζητήσω αὐτὴν,
\vs{15}ἐξ ἄνθους ὡς περκαζούσης σταφυλῆς· εὐφράνθη ἡ καρδία μου ἐν αὐτῇ, ἐπέβη ὁ πούς μου ἐν εὐθύτητι, ἐκ νεότητός μου ἴχνευσα αὐτήν.
\vs{16}Ἔκλινα ὀλίγον τὸ οὖς μου, καὶ ἐδεξάμην, καὶ πολλὴν εὗρον ἐμαυτῷ παιδείαν.
\vs{17}Προκοπὴ ἐγένετό μοι ἐν αὐτῇ· τῷ διδόντι μοι σοφίαν, δώσω δόξαν.
\vs{18}Διενοήθην γὰρ τοῦ ποιῆσαι αὐτὴν, καὶ ἐζήλωσα τὸ ἀγαθὸν, καὶ οὐ μὴ αἰσχυνθῶ.

\vs{19}Διαμεμάχηται ἡ ψυχή μου ἐν αὐτῇ, καὶ ἐν ποιήσει λιμοῦ διηκριβωσάμην· τὰς χεῖράς μου ἐξεπέτασα πρός ὕψος, καὶ τὰ ἀγνοήματα αὐτῆς ἐπένθησα,
\vs{20}τὴν ψυχήν μου κατεύθυνα εἰς αὐτὴν, καρδίαν ἐκτησάμην μετʼ αὐτῶν ἀπʼ ἀρχῆς, καὶ ἐν καθαρισμῷ εὗρον αὐτήν· διὰ τοῦτο οὐ μὴ ἐγκαταλειφθῶ.
\vs{21}Καὶ ἡ κοιλία μου ἐταράχθη ἐκζητῆσαι αὐτήν· διὰ τοῦτο ἐκτησάμην ἀγαθὸν κτῆμα.
\vs{22}Ἔδωκε Κύριος γλῶσσάν μοι μισθόν μου, καὶ ἐν αὐτῇ αἰνέσω αὐτόν.

\vs{23}Ἐγγίσατε πρὸς μὲ ἀπαίδευτοι, καὶ αὐλίσθητε ἐν οἴκῳ παιδείας.
\vs{24}Διότι ὑστερεῖτε ἐν τούτοις, καὶ αἱ ψυχαὶ ὑμῶν διψῶσι σφόδρα;
\vs{25}Ἤνοιξα τὸ στὸμα μου, καὶ ἐλάλησα, κτήσασθε ἑαυτοῖς ἄνευ ἀργυριου.
\vs{26}Τὸν τράχηλον ὑμῶν ὑπόθετε ὑπὸ ζυγὸν, καὶ ἐπιδεξάσθω ἡ ψυχὴ ὑμῶν παιδείαν, ἐγγύς ἐστιν εὑρεῖν αὐτήν.
\vs{27}Ἴδετε ἐν ὀφθαλμοῖς ὑμῶν ὅτι ὀλίγον ἐκοπίασα, καὶ εὗρον ἐμαυτῷ πολλὴν ἀνάπαυσιν.

\vs{28}Μετάσχετε παιδείας ἐν πολλῷ ἀριθμῷ ἀργυρίου, καὶ πολὺν χρυσὸν κτήσασθε ἐν αὐτῇ.
\vs{29}Εὐφρανθείη ἡ ψυχὴ ὑμῶν ἐν τῷ ἐλέει αὐτοῦ, καὶ μὴ αἰσχυνθείητε ἐν αἰνέσει αὐτοῦ.
\vs{30}Ἐργάζεσθε τὸ ἔργον ὑμῶν πρὸ καιροῦ, καὶ δώσει τὸν μισθὸν ὑμῶν ἐν καιρῷ αὐτοῦ.


\def\book{ΒΑΡΟΥΧ}
\biblebook{ΒΑΡΟΥΧ}


\lettrine[lines=2, loversize=0.2, nindent=0em, findent=.25em]{\textcolor{bookheadingcolor}{Κ}}{ΑΙ} οὗτοι οἱ λόγοι τοῦ βιβλίου, οὓς ἔγραψε Βαροὺχ υἱὸς Νηρίου, υἱοῦ Μαασαίου, υἱοῦ Σεδεκίου, υἱοῦ Ἀσαδίου, υἱοῦ Χελκίου ἐν Βαβυλῶνι,
\vs{2}ἐν τῷ ἔτει τῷ πέμπτῳ, ἐν ἑβδόμῃ τοῦ μηνὸς, ἐν τῷ καιρῷ ᾧ ἔλαβον οἱ Χαλδαῖοι τὴν Ἱερουσαλὴμ, καὶ ἐνέπρησαν αὐτὴν ἐν πυρί.

\vs{3}Καὶ ἀνέγνω Βαροὺχ τοὺς λόγους τοῦ βιβλίου τούτου ἐν ὠσὶν Ἰεχονίου υἱοῦ Ἰωακεὶμ βασιλέως Ἰούδα, καὶ ἐν ὠσὶ παντὸς τοῦ λαοῦ τῶν ἐρχομένων πρὸς τὴν βίβλον,
\vs{4}καὶ ἐν ὠσὶ τῶν δυνατῶν, καὶ υἱῶν τῶν βασιλέων, καὶ ἐν ὠσὶ τῶν πρεσβυτέρων, καὶ ἐν ὠσὶ παντὸς τοῦ λαοῦ, ἀπὸ μικροῦ ἕως μεγάλου, πάντων τῶν κατοικούντων ἐν Βαβυλῶνι ἐπὶ ποταμοῦ Σούδ.
\vs{5}Καὶ ἔκλαιον, καὶ ἐνήστευον, καὶ ηὔχοντο ἐναντίον Κυρίου.

\vs{6}Καὶ συνήγαγον ἀργύριον, καθὸ ἑκάστου ἠδύνατο ἡ χεὶρ,
\vs{7}καὶ ἀπέστειλαν εἰς Ἱερουσαλὴμ πρὸς Ἰωακεὶμ υἱὸν Χελκίου, υἱοῦ Σαλὼμ, τὸν ἱερέα, καὶ πρὸς τοὺς ἱερεῖς, καὶ πρὸς πάντα τὸν λαὸν, τοὺς εὑρεθέντας μετʼ αὐτοῦ ἐν Ἱερουσαλὴμ,
\vs{8}ἐν τῷ λαβεῖν αὐτὸν τὰ σκεύη οἴκου Κυρίου, τὰ ἐξενεχθέντα ἐκ τοῦ ναοῦ, ἀποστρέψαι εἰς γῆν Ἰούδα, τῇ δεκάτῃ τοῦ Σειουὰλ, σκεύη ἀργυρᾶ, ἃ ἐποίησε Σεδεκίας υἱὸς Ἰωσία βασιλεὺς Ἰούδα,
\vs{9}μετὰ τὸ ἀποικίσαι Ναβουχοδονόσορ βασιλέα Βαβυλῶνος τὸν Ἰεχονίαν, καὶ τοὺς ἄρχοντας, καὶ τοὺς δεσμώτας, καὶ τοὺς δυνατοὺς, καὶ τὸν λαὸν τῆς γῆς ἀπὸ Ἱερουσαλὴμ, καὶ ἤγαγεν αὐτὸν εἰς Βαβυλῶνα.

\vs{10}Καὶ εἶπαν, ἰδοὺ ἀπεστείλαμεν πρὸς ὑμᾶς ἀργύριον, καὶ ἀγοράσατε τοῦ ἀργυρίου ὁλοκαυτώματα, καὶ περὶ ἁμαρτίας, καὶ λίβανον, καὶ ποιήσατε μάννα, καὶ ἀνοίσατε ἐπὶ τὸ θυσιαστήριον Κυρίου τοῦ Θεοῦ ἡμῶν,
\vs{11}καὶ προσεύξασθε περὶ τῆς ζωῆς Ναβουχοδονόσορ βασιλέως Βαβυλῶνος, καὶ εἰς ζωὴν Βαλτάσαρ υἱοῦ αὐτοῦ, ἵνα ὦσιν αἱ ἡμέραι αὐτῶν ὡς αἱ ἡμέραι τοῦ οὐρανοῦ ἐπὶ τῆς γῆς.
\vs{12}Καὶ δώσει Κύριος ἰσχὺν ἡμῖν, καὶ φωτίσει τοὺς ὀφθαλμοὺς ἡμῶν, καὶ ζησόμεθα ὑπὸ τὴν σκιὰν Ναβουχοδονόσορ βασιλέως Βαβυλῶνος, καὶ ὑπὸ τὴν σκιὰν Βαλτάσαρ υἱοῦ αὐτοῦ, καὶ δουλεύσομεν αὐτοῖς ἡμέρας πολλὰς, καὶ εὑρήσομεν χάριν ἐναντίον αὐτῶν.

\vs{13}Καὶ προσεύξασθε περὶ ἡμῶν πρὸς Κύριον τὸν Θεὸν ἡμῶν, ὅτι ἡμάρτομεν τῷ Κυρίῳ Θεῷ ἡμῶν, καὶ οὐκ ἀπέστρεψεν ὁ θυμὸς Κυρίου καὶ ἡ ὀργὴ αὐτοῦ ἀφʼ ἡμῶν, ἕως τῆς ἡμέρας ταύτης.
\vs{14}Καὶ ἀναγνώσεσθε τὸ βιβλίον τοῦτο ὃ ἀπεστείλαμεν πρὸς ὑμᾶς, ἐξαγορεῦσαι ἐν οἴκῳ Κυρίου, ἐν ἡμέρᾳ ἑορτῆς, καὶ ἐν ἡμέραις καιροῦ,

\vs{15}Καὶ ἐρεῖτε, τῷ Κυρίῳ Θεῷ ἡμῶν ἡ δικαιοσύνη, ἡμῖν δὲ αἰσχύνη τῶν προσώπων, ὡς ἡ ἡμέρα αὕτη, ἀνθρώπῳ Ἰούδα, καὶ τοῖς κατοικοῦσιν Ἱερουσαλὴμ,
\vs{16}καὶ τοῖς βασιλεῦσιν ἡμῶν, καὶ τοῖς ἄρχουσιν ἡμῶν, καὶ τοῖς ἱερεῦσιν ἡμῶν, καὶ τοῖς προφήταις ἡμῶν, καὶ τοῖς πατράσιν ἡμῶν,
\vs{17}ὧν ἡμάρτομεν ἔναντι Κυρίου,
\vs{18}καὶ ἠπειθήσαμεν αὐτῷ, καὶ οὐκ ἠκούσαμεν τῆς φωνῆς Κυρίου Θεοῦ ἡμῶν, πορεύεσθαι τοῖς προστάγμασι Κυρίου, οἷς ἔδωκε κατὰ πρόσωπον ἡμῶν,
\vs{19}ἀπὸ τῆς ἡμέρας ἧς ἐξήγαγε Κύριος τοὺς πατέρας ἡμῶν ἐκ γῆς Αἰγύπτου· καὶ ἕως τῆς ἡμέρας ταύτης ἤμεθα ἀπειθοῦντες πρὸς Κύριον Θεὸν ἡμῶν, καὶ ἐσχεδιάζομεν πρὸς τὸ μὴ ἀκούειν τῆς φωνῆς αὐτοῦ.

\vs{20}Καὶ ἐκολλήθη εἰς ἡμᾶς τὰ κακὰ, καὶ ἡ ἀρὰ ἣν συνέταξε Κύριος τῷ Μωυσῇ παιδὶ αὐτοῦ, ἐν ἡμέρᾳ ᾗ ἐξήγαγε τοὺς πατέρας ἡμῶν ἐκ γῆς Αἰγύπτου, δοῦναι ἡμῖν γῆν ῥέουσαν γάλα καὶ μέλι, ὡς ἡ ἡμέρα αὐτη.
\vs{21}Καὶ οὐκ ἠκούσαμεν τῆς φωνῆς Κυρίου τοῦ Θεοῦ ἡμῶν, κατὰ πάντας τοὺς λόγους τῶν προφητῶν, ὧν ἀπέστειλε πρὸς ἡμᾶς.
\vs{22}Καὶ ᾠχόμεθα ἕκαστος ἐν διανοίᾳ καρδίας αὐτοῦ τῆς πονηρᾶς, ἐργάζεσθαι θεοῖς ἑτέροις ποιῆσαι τὰ κακὰ κατʼ ὀφθαλμοὺς Κυρίου Θεοῦ ἡμῶν.

\ch{2}
Καὶ ἔστησε Κύριος τὸν λόγον αὐτοῦ, ὃν ἐλάλησεν ἐφʼ ἡμᾶς, καὶ ἐπὶ τοὺς δικαστὰς ἡμῶν, τοὺς δικάσαντας τὸν Ἰσραὴλ, καὶ ἐπὶ τοὺς βασιλεῖς ἡμῶν, καὶ ἐπὶ τοὺς ἄρχοντας ἡμῶν,
\vs{2}καὶ ἐπὶ ἄνθρωπον Ἰσραὴλ καὶ Ἰούδα, τοῦ ἀγαγεῖν ἐφʼ ἡμᾶς κακὰ μεγάλα, ἃ οὐκ ἐποιήθη ὑποκάτω παντὸς τοῦ οὐρανοῦ, καθὰ ἐποίησεν ἐν Ἱερουσαλὴμ, κατὰ τὰ γεγραμμένα ἐν τῷ νόμῳ Μωυσῆ,
\vs{3}τοῦ φαγεῖν ἡμᾶς, ἄνθρωπον σάρκας υἱοῦ αὐτοῦ, καὶ ἄνθρωπον σάρκας θυγατρὸς αὐτοῦ.
\vs{4}Καὶ ἔδωκεν αὐτοὺς ὑποχειρίους πάσαις ταῖς βασιλείαις ταῖς κύκλῳ ἡμῶν, εἰς ὀνειδισμὸν, καὶ ἄβατον ἐν πᾶσι τοῖς λαοῖς τοῖς κύκλῳ, οὗ διέσπειρεν αὐτοὺς Κύριος ἐκεῖ.

\vs{5}Καὶ ἐγενήθησαν ὑποκάτω καὶ οὐκ ἐπάνω, ὅτι ἡμάρτομεν Κυρίῳ Θεῷ ἡμῶν, πρὸς τὸ μὴ ἀκούειν τῆς φωνῆς αὐτοῦ.

\vs{6}Τῷ Κυρίῳ Θεῷ ἡμῶν ἡ δικαιοσύνη, ἡμῖν δὲ καὶ τοῖς πατράσιν ἡμῶν ἡ αἰσχύνη τῶν προσώπων, ὡς ἡ ἡμέρα αὕτη.
\vs{7}Ἃ ἐλάλησε Κύριος ἐφʼ ἡμᾶς, πάντα τὰ κακὰ ταῦτα ἃ ἦλθεν ἐφʼ ἡμᾶς,
\vs{8}καὶ οὐκ ἐδεήθημεν τοῦ προσώπου Κυρίου, τοῦ ἀποστρέψαι ἕκαστον ἀπὸ τῶν νοημάτων τῆς καρδίας αὐτῶν τῆς πονηρᾶς.
\vs{9}Καὶ ἐγρηγόρησε Κύριος ἐπὶ τοῖς κακοῖς, καὶ ἐπήγαγε Κύριος ἐφʼ ἡμᾶς, ὅτι δίκαιος ὁ Κύριος ἐπὶ πάντα τὰ ἔργα αὐτοῦ, ἃ ἐνετείλατο ἡμῖν.
\vs{10}Καὶ οὐκ ἠκούσαμεν τῆς φωνῆς αὐτοῦ, πορεύεσθαι τοῖς προστάγμασι Κυρίου, οἷς ἔδωκε κατὰ πρόσωπον ἡμῶν·

\vs{11}Καὶ νῦν Κύριε ὁ Θεὸς Ἰσραὴλ, ὃς ἐξήγαγες τὸν λαόν σου ἐκ γῆς Αἰγύπτου, ἐν χειρὶ κραταιᾷ, καὶ ἐν σημείοις, καὶ ἐν τέρασι, καὶ ἐν δυνάμει μεγάλῃ, καὶ ἐν βραχίονι ὑψηλῷ, καὶ ἐποίησας σεαυτῷ ὄνομα, ὡς ἡ ἡμέρα αὕτη,
\vs{12}ἡμάρτομεν, ἠσεβήσαμεν, ἠδικήσαμεν, Κύριε ὁ Θεὸς ἡμῶν, ἐπὶ πᾶσι τοῖς δικαιώμασί σου.

\vs{13}Ἀποστραφήτω ὁ θυμός σου ἀφʼ ἡμῶν, ὅτι κατελείφθημεν ὀλίγοι ἐν τοῖς ἔθνεσιν οὗ διέσπειρας ἡμᾶς ἐκεῖ.

\vs{14}Εἰσάκουσον Κύριε τῆς προσευχῆς ἡμῶν καὶ τῆς δεήσεως ἡμῶν, καὶ ἐξελοῦ ἡμᾶς ἕνεκέν σου, καὶ δὸς ἡμῖν χάριν κατὰ πρόσωπον τῶν ἀποικισάντων ἡμᾶς,
\vs{15}ἵνα γνῷ πᾶσα ἡ γῆ, ὅτι σὺ Κύριος ὁ Θεὸς ἡμῶν, ὅτι τὸ ὄνομά σου ἐπεκλήθη ἐπὶ Ἰσραὴλ, καὶ ἐπὶ τὸ γένος αὐτοῦ.

\vs{16}Κύριε κάτιδε ἐκ τοῦ οἴκου τοῦ ἁγίου σου, καὶ ἐννόησον εἰς ἡμᾶς, καὶ κλῖνον Κύριε τὸ οὖς σου, καὶ ἄκουσον.

\vs{17}Ἄνοιξον ὀφθαλμούς σου, καὶ ἴδε, ὅτι οὐχ οἱ τεθνηκότες ἐν τῷ ᾅδῃ, ὧν ἐλήφθη τὸ πνεῦμα αὐτῶν ἀπὸ τῶν σπλάγχνων αὐτῶν, δώσουσι δόξαν καὶ δικαίωμα τῷ Κυρίῳ·
\vs{18}ἀλλὰ ἡ ψυχὴ ἡ λυπουμένη ἐπὶ τὸ μέγεθος, ὃ βαδίζει κύπτον καὶ ἀσθενοῦν, καὶ οἱ ὀφθαλμοὶ οἱ ἐκλείποντες, καὶ ἡ ψυχὴ ἡ πεινῶσα, δώσουσί σοι δόξαν, καὶ δικαιοσύνην, Κύριε.

\vs{19}Ὅτι οὐκ ἐπὶ τὰ δικαιώματα τῶν πατέρων ἡμῶν καὶ τῶν βασιλέων ἡμῶν ἡμεῖς καταβάλλομεν τὸν ἔλεον κατὰ πρόσωπόν σου, Κύριε ὁ Θεὸς ἡμῶν·
\vs{20}ὅτι ἐνῆκας τὸν θυμόν σου καὶ τὴν ὀργήν σου εἰς ἡμᾶς, καθάπερ ἐλάλησας ἐν χειρὶ τῶν παίδων σου τῶν προφητῶν·

\vs{21}Οὕτως εἶπε Κύριος, κλίνατε τὸν ὦμον ὑμῶν, καὶ ἐργάσασθε τῷ βασιλεῖ Βαβυλῶνος, καὶ καθίσατε ἐπὶ τὴν γῆν, ἣν δέδωκα τοῖς πατράσιν ὑμῶν.
\vs{22}Καὶ ἐὰν μὴ ἀκούσητε τῆς φωνῆς Κυρίου, ἐργάσασθαι τῷ βασιλεῖ Βαβυλῶνος,
\vs{23}ἐκλείψειν ποιήσω ἐκ πόλεων Ἰούδα καὶ ἔξωθεν Ἱερουσαλὴμ φωνὴν εὐφροσύνης, καὶ φωνὴν χαρμοσύνης, φωνὴν νυμφίου, καὶ φωνὴν νύμφης, καὶ ἔσται πᾶσα ἡ γῆ εἰς ἄβατον ἀπὸ ἐνοικούντων.
\vs{24}Καὶ οὐκ ἠκούσαμεν τῆς φωνῆς σου, ἐργάσασθαι τῷ βασιλεῖ Βαβυλῶνος· καὶ ἔστησας τοὺς λόγους σου, οὓς ἐλάλησας ἐν χερσὶ τῶν παίδων σου τῶν προφητῶν, τοῦ ἐξενεχθῆναι τὰ ὀστᾶ βασιλέων ἡμῶν καὶ τὰ ὀστᾶ τῶν πατέρων ἡμῶν ἐκ τοῦ τόπου αὐτῶν.

\vs{25}Καὶ ἰδού ἐστιν ἐξεῤῥιμμένα τῷ καύματι τῆς ἡμέρας, καὶ τῷ παγετῷ τῆς νυκτός· καὶ ἀπεθάνοσαν ἐν πόνοις πονηροῖς, ἐν λιμῷ, καὶ ἐν ῥομφαίᾳ, καὶ ἐν ἀποστολῇ.
\vs{26}Καὶ ἔθηκας τὸν οἶκον, οὗ ἐπεκλήθη τὸ ὄνομά σου ἐπʼ αὐτῷ, ὡς ἡ ἡμέρα αὕτη, διὰ πονηρίαν οἴκου Ἰσραὴλ καὶ οἴκου Ἰούδα.

\vs{27}Καὶ ἐποίησας εἰς ἡμᾶς, Κύριε ὁ Θεὸς ἡμῶν, κατὰ πᾶσαν ἐπιείκειάν σου, καὶ κατὰ πάντα οἰκτιρμόν σου τὸν μέγαν,
\vs{28}καθὰ ἐλάλησας ἐν χειρὶ παιδός σου Μωυσῆ, ἐν ἡμέρᾳ ἐντειλαμένου σου αὐτῷ γράψαι τὸν νόμον σου ἐναντίον υἱῶν Ἰσραὴλ, λέγων.

\vs{29}Ἐὰν μὴ ἀκούσητε τῆς φωνῆς μου, ἦ μὴν ἡ βόμβησις ἡ μεγάλη ἡ πολλὴ αὕτη ἀποστρέψει εἰς μικρὰν ἐν τοῖς ἔθνεσιν, οὗ διασπερῶ αὐτοὺς ἐκεῖ·

\vs{30}Ὅτι ἔγνων ὅτι οὐ μὴ ἀκούσωσί μου, ὅτι λαὸς σκληροτράχηλός ἐστι· καὶ ἐπιστρέψουσιν ἐπὶ καρδίαν αὐτῶν ἐν γῇ ἀποικισμοῦ αὐτῶν,
\vs{31}καὶ γνώσονται ὅτι ἐγὼ Κύριος ὁ Θεὸς αὐτῶν· καὶ δώσω αὐτοῖς καρδίαν καὶ ὦτα ἀκούοντα,
\vs{32}καὶ αἰνέσουσί με ἐν γῇ ἀποικισμοῦ αὐτῶν· καὶ μνησθήσονται τοῦ ὀνόματός μου,
\vs{33}καὶ ἀποστρέψουσιν ἀπὸ τοῦ νώτου αὐτῶν τοῦ σκληροῦ, καὶ ἀπὸ πονηρῶν προσταγμάτων αὐτῶν, ὅτι μνησθήσονται τῆς ὁδοῦ πατέρων αὐτῶν τῶν ἁμαρτόντων ἔναντι Κυρίου.

\vs{34}Καὶ ἀποστρέψω αὐτοὺς εἰς τὴν γῆν, ἣν ὤμοσα τοῖς πατράσιν αὐτῶν, τῷ Ἁβραὰμ, καὶ τῷ Ἰσαὰκ, καὶ τῷ Ἰακὼβ, καὶ κυριεύσουσιν αὐτῆς· καὶ πληθυνῶ αὐτοὺς, καὶ οὐ μὴ σμικρυνθῶσι.
\vs{35}Καὶ στήσω αὐτοῖς διαθήκην αἰώνιον, τοῦ εἶναί με αὐτοῖς εἰς Θεὸν, καὶ αὐτοὶ ἔσονταί μοι εἰς λαόν· καὶ οὐ κινήσω ἔτι τὸν λαόν μου Ἰσραὴλ ἀπὸ τῆς γῆς, ἧς ἔδωκα αὐτοῖς.

\ch{3}
Κύριε παντοκράτωρ ὁ Θεὸς Ἰσραὴλ, ψυχὴ ἐν στενοῖς καὶ πνεῦμα ἀκηδιῶν κέκραγε πρὸς σέ.
\vs{2}Ἄκουσον, Κύριε, καὶ ἐλέησον, ὅτι ἡμάρτομεν ἐναντίον σου·
\vs{3}ὅτι σὺ καθήμενος τὸν αἰῶνα, καὶ ἡμεῖς ἀπολλύμενοι τὸν αἰῶνα.

\vs{4}Κύριε παντοκράτωρ ὁ Θεὸς Ἰσραὴλ, ἄκουσον δὴ τῆς προσευχῆς τῶν τεθνηκότων Ἰσραὴλ, καὶ υἱῶν τῶν ἁμαρτανόντων ἐναντίον σου, οἳ οὐκ ἤκουσαν τῆς φωνῆς σου Θεοῦ αὐτῶν, καὶ ἐκολλήθη ἡμῖν τὰ κακά.
\vs{5}Μὴ μνησθῇς ἀδικιῶν πατέρων ἡμῶν, ἀλλὰ μνήσθητι χειρός σου καὶ ὀνόματός σου ἐν τῷ καιρῷ τούτῳ·
\vs{6}ὅτι σὺ Κύριος ὁ Θεὸς ἡμῶν, καὶ αἰνέσομέν σε Κύριε·

\vs{7}Ὅτι διὰ τοῦτο ἔδωκας τὸν φόβον σου ἐπὶ καρδίαν ἡμῶν, καὶ ἐπικαλεῖσθαι τὸ ὄνομά σου· καὶ αἰνέσομέν σε ἐν τῇ ἀποικίᾳ ἡμῶν, ὅτι ἀπεστρέψαμεν ἀπὸ καρδίας ἡμῶν πᾶσαν ἀδικίαν πατέρων ἡμῶν, τῶν ἡμαρτηκότων ἐναντίον σου.

\vs{8}Ἰδοὺ ἡμεῖς σήμερον ἐν τῇ ἀποικίᾳ ἡμῶν, οὗ διέσπειρας ἡμᾶς ἐκεῖ εἰς ὀνειδισμὸν, καὶ εἰς ἀρὰν, καὶ εἰς ὀφλησιν κατὰ πάσας τὰς ἀδικίας πατέρων ἡμῶν, οἳ ἀπέστησαν ἀπὸ Κυρίου Θεοῦ ἡμῶν.

\vs{9}Ἄκουε Ἰσραὴλ ἐντολὰς ζωῆς, ἐνωτίσασθε γνῶναι φρόνησιν.
\vs{10}Τί ἐστιν Ἰσραήλ; τί ὅτι ἐν γῇ τῶν ἐχθρῶν εἶ; ἐπαλαιώθης ἐν γῇ ἀλλοτρίᾳ, συνεμιάνθης τοῖς νεκροῖς,
\vs{11}προσελογίσθης μετὰ τῶν εἰς ᾅδου,
\vs{12}ἐγκατέλιπες τὴν πηγὴν τῆς σοφίας.
\vs{13}Τῇ ὁδῷ τοῦ Θεοῦ εἰ ἐπορεύθης, κατῴκεις ἂν ἐν εἰρήνῃ τὸν αἰῶνα.

\vs{14}Μάθε ποῦ ἐστι φρόνησις, ποῦ ἐστιν ἰσχὺς, ποῦ ἐστιν σύνεσις, τοῦ γνῶναι ἅμα ποῦ ἐστι μακροβίωσις καὶ ζωὴ, ποῦ ἐστι φῶς ὀφθαλμῶν καὶ εἰρήνη.
\vs{15}Τίς εὗρε τὸν τόπον αὐτῆς, καὶ τίς εἰσῆλθεν εἰς τοὺς θησαυροὺς αὐτῆς;

\vs{16}Ποῦ εἰσιν οἱ ἄρχοντες τῶν ἐθνῶν, καὶ οἱ κυριεύοντες τῶν θηρίων τῶν ἐπὶ τῆς γῆς,
\vs{17}οἱ ἐν τοῖς ὀρνέοις τοῦ οὐρανοῦ ἐμπαίζοντες, καὶ τὸ ἀργύριον θησαυρίζοντες, καὶ τὸ χρυσίον ᾧ ἐπεποίθεισαν ἄνθρωποι, καὶ οὐκ ἔστι τέλος τῆς κτήσεως αὐτῶν;
\vs{18}Ὅτι οἱ τὸ ἀργύριον τεκταίνοντες καὶ μεριμνῶντες, καὶ οὐκ ἔστιν ἐξεύρεσις τῶν ἔργων αὐτῶν.
\vs{19}Ἠφανίσθησαν, καὶ εἰς ᾅδου κατέβησαν, καὶ ἄλλοι ἀνέστησαν ἀντʼ αὐτῶν.

\vs{20}Νεώτεροι εἶδον φῶς, καὶ κατῴκησαν ἐπὶ τῆς γῆς, ὁδὸν δὲ ἐπιστήμης οὐκ ἔγνωσαν,
\vs{21}οὐδὲ συνῆκαν τρίβους αὐτῆς, οὐδὲ ἀντελάβοντο αὐτῆς· οἱ υἱοὶ αὐτῶν ἀπὸ τῆς ὁδοῦ αὐτῶν πόῤῥω ἐγενήθησαν.
\vs{22}Οὐδὲ ἠκούσθη ἐν Χαναὰν, οὐδὲ ὤφθη ἐν Θαιμάν.

\vs{23}Οἵτε υἱοὶ Ἄγαρ οἱ ἐκζητοῦντες τὴν σύνεσιν οἱ ἐπὶ τῆς γῆς, οἱ ἔμποροι τῆς Μεῤῥὰν, καὶ Θαιμὰν, καὶ οἱ μυθολόγοι, καὶ οἱ ἐκζητηταὶ τῆς συνέσεως, ὁδὸν δὲ σοφίας οὐκ ἔγνωσαν, οὐδὲ ἐμνήσθησαν τὰς τρίβους αὐτῆς.

\vs{24}Ὦ Ἰσραὴλ, ὡς μέγας ὁ οἶκος τοῦ Θεοῦ; καὶ ἐπιμήκης ὁ τόπος τῆς κτήσεως αὐτοῦ;
\vs{25}Μέγας, καὶ οὐκ ἔχει τελευτὴν, ὑψηλὸς καὶ ἀμέτρητος.
\vs{26}Ἐκεῖ ἐγεννήθησαν οἱ γίγαντες οἱ ὀνομαστοὶ, ἀπʼ ἀρχῆς γενόμενοι εὐμεγέθεις, ἐπιστάμενοι πόλεμον.
\vs{27}Οὐ τούτους ἐξελέξατο ὁ Θεὸς, οὐδὲ ὁδὸν ἐπιστήμης ἔδωκεν αὐτοῖς.
\vs{28}Καὶ ἀπώλοντο παρὰ τὸ μὴ ἔχειν φρόνησιν, ἀπώλοντο διὰ τὴν ἀβουλίαν αὐτῶν.

\vs{29}Τίς ἀνέβη εἰς τὸν οὐρανὸν, καὶ ἔλαβεν αὐτὴν, καὶ κατεβίβασεν αὐτὴν ἐκ τῶν νεφελῶν;
\vs{30}Τίς διέβη πέραν τῆς θαλάσσης, καὶ εὗρεν αὐτὴν, καὶ οἴσει αὐτὴν χρυσίου ἐκλεκτοῦ;
\vs{31}Οὐκ ἔστιν ὁ γινώσκων τὴν ὁδὸν αὐτῆς, οὐδὲ ὁ ἐνθυμούμενος τὴν τρίβον αὐτῆς.

\vs{32}Ἀλλʼ ὁ εἰδὼς τὰ πάντα γινώσκει αὐτὴν, ἐξεῦρεν αὐτὴν τῇ συνέσει αὐτοῦ· ὁ κατασκευάσας τὴν γῆν εἰς τὸν αἰῶνα χρόνον, ἐνέπλησεν αὐτὴν κτηνῶν τετραπόδων.
\vs{33}Ὁ ἀποστέλλων τὸ φῶς καὶ πορεύεται, ἐκάλεσεν αὐτὸ, καὶ ὑπήκουσεν αὐτῷ τρόμῳ·
\vs{34}Οἱ δὲ ἀστέρες ἔλαμψαν ἐν ταῖς φυλακαῖς αὐτῶν, καὶ εὐφράνθησαν· ἐκάλεσεν αὐτοὺς, καὶ εἶπον, πάρεσμεν· ἔλαμψαν μετʼ εὐφροσύνης τῷ ποιήσαντι αὐτούς.

\vs{35}Οὗτος ὁ Θεὸς ἡμῶν, οὐ λογισθήσεται ἕτερος πρὸς αὐτόν.
\vs{36}Ἐξεῦρε πᾶσαν ὁδὸν ἐπιστήμης, καὶ ἔδωκεν αὐτὴν Ἰακὼβ τῷ παιδὶ αὐτοῦ, καὶ Ἰσραὴλ τῷ ἠγαπημένῳ ὑπʼ αὐτοῦ.
\vs{37}Μετὰ τοῦτο ἐπὶ τῆς γῆς ὤφθη, καὶ ἐν τοῖς ἀνθρώποις συνανεστράφη.

\ch{4}
Αὑτὴ ἡ βίβλος τῶν προσταγμάτων τοῦ Θεοῦ, καὶ ὁ νόμος ὁ ὑπάρχων εἰς τὸν αἰῶνα· πάντες οἱ κρατοῦντες αὐτὴν, εἰς ζωήν· οἱ δὲ καταλείποντες αὐτὴν, ἀποθανοῦνται.
\vs{2}Ἐπιστρέφου Ἱακὼβ, καὶ ἐπιλαβοῦ αὐτῆς, διόδευσον πρὸς τὴν λάμψιν κατέναντι τοῦ φωτὸς αὐτῆς.
\vs{3}Μὴ δῷς ἑτέρῳ τὴν δόξαν σου, καὶ τὰ συμφέροντά σοι ἔθνει ἀλλοτρίῳ.

\vs{4}Μακάριοι ἐσμὲν Ἰσραὴλ, ὅτι τὰ ἀρεστὰ τοῦ Θεοῦ ἡμῖν γνωστά ἐστι.
\vs{5}Θαρσεῖτε λαός μου, μνημόσυνον Ἰσραήλ.
\vs{6}Ἐπράθητε τοῖς ἔθνεσιν οὐκ εἰς ἀπώλειαν, διὰ τὸ παροργίσαι ὑμᾶς τὸν Θεόν· παρεδόθητε τοῖς ὑπεναντίοις.
\vs{7}Παρωξύνατε γὰρ τὸν ποιήσαντα ὑμᾶς, θύσαντες δαιμονίοις, καὶ οὐ Θεῷ.
\vs{8}Ἐπελάθεσθε τὸν τροφεύσαντα ὑμᾶς Θεὸν αἰώνιον, ἐλυπήσατε δὲ καὶ τὴν ἐκθρέψασαν ἡμᾶς Ἱερουσαλήμ.

\vs{9}Εἶδε γὰρ τὴν ἐπελθοῦσαν ὑμῖν ὀργὴν παρὰ τοῦ Θεοῦ, καὶ εἶπεν, ἀκούσατε αἱ πάροικοι Σιὼν, ἐπήγαγέ μοι ὁ Θεὸς πένθος μέγα.
\vs{10}Εἶδον γὰρ τὴν αἰχμαλωσίαν τῶν υἱῶν μου καὶ τῶν θυγατέρων, ἣν ἐπήγαγεν αὐτοῖς ὁ αἰώνιος.
\vs{11}Ἔθρεψα γὰρ αὐτοὺς μετʼ εὐφροσύνης, ἐξαπέστειλα δὲ μετὰ κλαυθμοῦ καὶ πένθους.

\vs{12}Μηδεὶς ἐπιχαιρέτω μοι τῇ χήρᾳ καὶ καταλειφθήσῃ ὑπὸ πολλῶν, ἠρημώθην διὰ τὰς ἁμαρτίας τῶν τέκνων μου, διότι ἐξέκλιναν ἐκ νόμου Θεοῦ,
\vs{13}καὶ δικαιώματα αὐτοῦ οὐκ ἔγνωσαν, οὐδὲ ἐπορεύθησαν ὁδοῖς ἐντολῶν Θεοῦ, οὐδὲ τρίβους παιδείας ἐν δικαιοσύνῃ αὐτοῦ ἐπέβησαν.

\vs{14}Ἐλθέτωσαν αἱ πάροικοι Σιὼν, καὶ μνήσθητε τὴν αἰχμαλωσίαν τῶν υἱῶν μου καὶ θυγατέρων, ἣν ἐπήγαγεν αὐτοῖς ὁ αἰώνιος.
\vs{15}Ἐπήγαγε γὰρ ἐπʼ αὐτοὺς ἔθνος μακρόθεν, ἔθνος ἀναιδὲς καὶ ἀλλόγλωσσον· ὅτι οὐκ ᾐσχύνθησαν πρεσβύτην, οὐδὲ παιδίον ἠλέησαν,
\vs{16}καὶ ἀπήγαγον τοὺς ἀγαπητοὺς τῆς χήρας, καὶ ἀπὸ τῶν θυγατέρων τὴν μόνην ἠρήμωσαν.
\vs{17}Ἐγὼ δὲ τί δυνατὴ βοηθῆσαι ὑμῖν; Ὁ γὰρ ἐπαγαγὼν τὰ κακὰ, ἐξελεῖται ὑμᾶς ἐκ χειρὸς ἐχθρῶν ὑμῶν.
\vs{18}Βαδίζετε τέκνα, βαδίζετε, ἐγὼ γὰρ κατελείφθην ἔρημος.
\vs{19}Ἐξεδυσάμην τὴν στολὴν τῆς εἰρήνης,
\vs{20}ἐνεδυσάμην δὲ σάκκον τῆς δεήσεώς μου· κεκράξομαι πρὸς τὸν αἰώνιον ἐν ταῖς ἡμέραις μου.

\vs{21}Θαῤῥεῖτε τέκνα, βοήσατε πρὸς τὸν Θεὸν, καὶ ἐξελεῖται ὑμᾶς ἐκ δυναστείας, ἐκ χειρος ἐχθρῶν.

\vs{22}Ἐγὼ γὰρ ἤλπισα ἐπὶ τῷ αἰωνίῳ τὴν σωτηρίαν ὑμῶν· καὶ ἦλθέ μοι χαρὰ παρὰ τοῦ ἁγίου ἐπὶ τῇ ἐλεημοσύνῃ, ἣ ἥξει ὑμῖν ἐν τάχει παρὰ τοῦ αἰωνίου σωτῆρος ὑμῶν.

\vs{23}Ἐξέπεμψα γὰρ ὑμᾶς μετὰ κλαυθμοῦ καὶ πένθους, ἀποδώσει δέ μοι ὁ Θεὸς ὑμᾶς μετὰ χαρμοσύνης καὶ εὐφροσύνης εἰς τὸν αἰῶνα.
\vs{24}Ὥσπερ γὰρ νῦν ἑωράκασιν αἱ πάροικοι Σιὼν τὴν ὑμετέραν αἰχμαλωσίαν, οὕτως ὄψονται ἐν τάχει τὴν παρὰ τοῦ Θεοῦ ὑμῶν σωτηρίαν, ἣ ἐπελεύσεται ὑμῖν μετὰ δόξης μεγάλης καὶ λαμπρότητος τοῦ αἰωνίου.

\vs{25}Τέκνα μακροθυμήσατε τὴν παρὰ τοῦ Θεοῦ ἐπελθοῦσαν ὑμῖν ὀργήν, κατεδίωξέ σε ὁ ἐχθρὸς, καὶ ὄψει αὐτοῦ τὴν ἀπώλειαν ἐν τάχει, καὶ ἐπὶ τραχήλους αὐτῶν ἐπιβήσῃ.
\vs{26}Οἱ τρυφεροί μου ἐπορεύθησαν ὁδοὺς τραχείας, ᾔρθησαν ὡς ποίμνιον ἡρπασμένον ὑπὸ ἐχθρῶν.

\vs{27}Θαρσήσατε τέκνα καὶ βοήσατε πρὸς τὸν Θεὸν, ἔσται γὰρ ὑμῶν ὑπὸ τοῦ ἐπάγοντος μνεία.
\vs{28}Ὥσπερ γὰρ ἐγένετο ἡ διάνοια ὑμῶν εἰς τὸ πλανηθῆναι ἀπὸ τοῦ Θεοῦ, δεκαπλασιάσατε ἐπιστραφέντες ζητῆσαι αὐτόν.
\vs{29}Ὁ γὰρ ἐπαγαγὼν ὑμῖν τὰ κακὰ, ἐπάξει ὑμῖν τὴν αἱώνιον εὐφροσύνην μετὰ τῆς σωτηρίας ὑμῶν.

\vs{30}Θάρσει Ἱερουσαλήμ, παρακαλέσει σε ὁ ὀνομάσας σε.
\vs{31}Δείλαιοι οἱ σὲ κακώσαντες, καὶ ἐπιχαρέντες τῇ σῇ πτώσει·
\vs{32}Δείλαιαι αἱ πόλεις αἷς ἐδούλευσαν τὰ τέκνα σου, δειλαία ἡ δεξαμένη τοὺς υἱούς σου.
\vs{33}Ὥσπερ γὰρ ἐχάρη ἐπὶ τῇ σῇ πτώσει, καὶ εὐφράνθη ἐπὶ τῷ πτώματί σου, οὕτως λυπηθήσεται ἐπὶ τῇ ἑαυτῆς ἐρημίᾳ.
\vs{34}Καὶ περιελῶ αὐτῆς τὸ ἀγαλλίαμα τῆς πολυοχλίας καὶ τὸ γαυρίαμα αὐτῆς εἰς πένθος.
\vs{35}Πῦρ γὰρ ἐπελεύσεται αὐτῇ παρὰ τοῦ αἰωνίου εἰς ἡμέρας μακρὰς, καὶ κατοικηθήσεται ὑπὸ δαιμονίων τὸν πλείονα χρόνον.

\vs{36}Περίβλεψον πρὸς ἀνατολὰς Ἱερουσαλὴμ, καὶ ἴδε τὴν εὐφροσύνην τὴν παρὰ τοῦ Θεοῦ σοι ἐρχομένην.
\vs{37}Ἰδοὺ ἔρχονται οἱ υἱοί σου οὓς ἐξαπέστειλας, ἔρχονται συνηγμένοι ἀπὸ ἀνατολῶν ἕως δυσμῶν τῷ ῥήματι τοῦ ἁγίου, χαίροντες τῇ τοῦ Θεοῦ δόξῃ.

\ch{5}
Ἔκδυσαι Ἱερουσαλὴμ, τὴν στολὴν τοῦ πένθους καὶ τῆς κακώσεώς σου, καὶ ἔνδυσαι τὴν εὐπρέπειαν τῆς παρὰ τοῦ Θεοῦ δόξης εἰς τὸν αἰῶνα.

\vs{2}Περιβαλοῦ τὴν διπλοΐδα τῆς παρὰ τοῦ Θεοῦ δικαιοσύνης, ἐπίθου τὴν μίτραν ἐπὶ τὴν κεφαλήν σου τῆς δόξης τοῦ αἰωνίου.
\vs{3}Ὁ γὰρ Θεὸς δείξει τῇ ὑπʼ οὐρανὸν πάσῃ τὴν σὴν λαμπρότητα.
\vs{4}Κληθήσεται γάρ σου τὸ ὄνομα παρὰ τοῦ Θεοῦ εἰς τὸν αἰῶνα, εἰρήνη δικαιοσύνης, καὶ δόξα θεοσεβείας.

\vs{5}Ἀνάστηθι Ἱερουσαλὴμ, καὶ στῆθι ἐπὶ τοῦ ὑψηλοῦ, καὶ περίβλεψαι πρὸς ἀνατολὰς, καὶ ἴδε συνηγμένα τὰ τέκνα σου ἀπὸ ἡλίου δυσμῶν ἕως ἀνατολῶν τῷ ῥήματι τοῦ ἁγίου, χαίροντας τῇ τοῦ Θεοῦ μνείᾳ.
\vs{6}Ἐξῆλθον γὰρ παρὰ σοῦ πεζοὶ ἀγόμενοι ὑπὸ ἐχθρῶν, εἰσάγει δὲ αὐτοὺς ὁ Θεὸς πρὸς σὲ αἰρομένους μετὰ δόξης ὡς θρόνον βασιλείας.

\vs{7}Συνέταξε γὰρ ὁ Θεὸς ταπεινοῦσθαι πᾶν ὄρος ὑψηλὸν, καὶ θῖνας ἀεννάους, καὶ φάραγγας πληροῦσθαι εἰς ὁμαλισμὸν τῆς γῆς, ἵνα βαδίσῃ Ἰσραὴλ ἀσφαλῶς τῇ τοῦ Θεοῦ δόξῃ.
\vs{8}Ἐσκίασαν δὲ καὶ οἱ δρυμοὶ καὶ πᾶν ξύλον εὐωδίας τῷ Ἰσραὴλ προστάγματι τοῦ Θεοῦ.
\vs{9}Ἡγήσεται γὰρ ὁ Θεὸς Ἰσραὴλ μετʼ εὐφροσύνης τῷ φωτὶ τῆς δόξης αὐτοῦ, σὺν ἐλεημοσύνῃ καὶ δικαιοσύνῃ τῇ παρʼ αὐτοῦ.


\def\book{ΕΠΙΣΤΟΛΗ ΙΕΡΕΜΙΟΥ}
\biblebook{ΕΠΙΣΤΟΛΗ ΙΕΡΕΜΙΟΥ}

 
\lettrine[lines=2, loversize=0.2, nindent=0em, findent=.25em]{\textcolor{bookheadingcolor}{Α}}{ΝΤΙΓΡΑΦΟΝ} ἐπιστολῆς ἧς ἀπέστειλεν Ἱερεμίας πρὸς τοὺς ἀχθησομένους αἰχμαλώτους εἰς Βαβυλῶνα ὑπὸ τοῦ βασιλέως τῶν Βαβυλωνίων, ἀναγγεῖλαι αὐτοῖς καθότι ἐπετάγη αὐτῷ ὑπὸ τοῦ Θεοῦ.

\vs{2}Διὰ τὰς ἁμαρτίας ἃς ἡμαρτήκατε ἐναντίον τοῦ Θεοῦ, ἀχθήσεσθε εἰς Βαβυλῶνα αἰχμάλωτοι ὑπὸ Ναβουχοδονόσορ βασιλέως τῶν Βαβυλωνίων.
\vs{3}Εἰσελθόντες οὖν εἰς Βαβυλῶνα, ἔσεσθε ἐκεῖ ἔτη πλείονα καὶ χρόνον μακρὸν, ἕως γενεῶν ἑπτά· μετὰ τοῦτο δὲ ἐξάξω ὑμᾶς ἐκεῖθεν μετʼ εἰρήνης.

\vs{4}Νυνὶ δὲ ὄψεσθε ἐν Βαβυλῶνι θεοὺς ἀργυροῦς καὶ χρυσοῦς καὶ ξυλίνους ἐπʼ ὤμοις αἰρομένους, δεικνύντας φόβον τοῖς ἔθνεσιν.
\vs{5}Εὐλαβήθητε οὖν μὴ καὶ ὑμεῖς ἀφομοιωθέντες τοῖς ἀλλοφύλοις ἀφομοιωθῆτε, καὶ φόβος ὑμᾶς λὰβῃ ἐπʼ αὐτοῖς, ἰδόντας ὄχλον ἔμπροσθεν καὶ ὄπισθεν αὐτῶν προσκυνοῦντας αὐτά.
\vs{6}Εἴπατε δὲ τῇ διανοίᾳ, σοὶ δεῖ προσκυνεῖν, δέσποτα.
\vs{7}Ὁ γὰρ ἄγγελός μου μεθʼ ὑμῶν ἐστιν, αὐτός τε ἐκζητῶν τὰς ψυχὰς ὑμῶν.

\vs{8}Γλῶσσα γὰρ αὐτῶν ἐστι κατεξυσμένη ὑπὸ τέκτονος, αὐτά τε περίχρυσα καὶ περιάργυρα, ψευδῆ δʼ ἐστὶ, καὶ οὐ δύνανται λαλεῖν.
\vs{9}Καὶ ὥσπερ παρθένῳ φιλοκόσμῳ λαμβάνοντες χρυσίον, κατασκευάζουσι στεφάνους ἐπὶ τὰς κεφαλὰς τῶν θεῶν αὐτῶν.
\vs{10}Ἔστι δὲ καὶ ὅτε ὑφαιρούμενοι οἱ ἱερεῖς ἀπὸ τῶν θεῶν αὐτῶν χρυσίον καὶ ἀργύριον εἰς ἑαυτοὺς καταναλοῦσι.
\vs{11}Δώσουσι δὲ ἀπʼ αὐτῶν καὶ ταῖς ἐπὶ τοῦ στέγους πόρναις· κοσμοῦσί τε αὐτοὺς, ὡς ἀνθρώπους, τοῖς ἐνδύμασι, θεοὺς ἀργυροῦς, καὶ θεοὺς χρυσοῦς, καὶ ξυλίνους.

\vs{12}Οὗτοι δὲ οὐ διασώζονται ἀπὸ ἰοῦ καὶ βρωμάτων, περιβεβλημένων αὐτῶν ἱματισμὸν πορφυροῦν.
\vs{13}Ἐκμάσσονται τὸ πρόσωπον αὐτῶν διὰ τὸν ἐκ τῆς οἰκίας κονιορτὸν, ὅς ἐστι πλείων ἐπʼ αὐτοῖς.
\vs{14}Καὶ σκῆπτρον ἔχει ὡς ἄνθρωπος κριτὴς χώρας, ὃς τὸν εἰς αὐτὸν ἁμαρτάνοντα οὐκ ἀνελεῖ.
\vs{15}Ἔχει δὲ ἐγχειρίδιον δεξιᾷ, καὶ πέλεκυν· ἑαυτὸν δὲ ἐκ πολέμον καὶ λῃστῶν οὐκ ἐξελεῖται.
\vs{16}Ὅθεν γνώριμοί εἰσιν οὐκ ὄντες θεοί· μὴ οὖν φοβηθῆτε αὐτούς.

\vs{17}Ὥσπερ γὰρ σκεῦος ἀνθρώπου συντριβὲν ἀχρεῖον γῖνεται, τοιοῦτοι ὑπάρχουσιν οἱ θεοὶ αὐτῶν, καθιδρυμένων αὐτῶν ἐν τοῖς οἴκοις· οἱ ὀφθαλμοὶ αὐτῶν πλήρεις εἰσὶ κονιορτοῦ ἀπὸ τῶν ποδῶν τῶν εἰσπορευομένων.
\vs{18}Καὶ ὥσπερ τινὶ ἠδικηκότι βασιλέα, περιπεφραγμέναι εἰσὶν αἱ αὐλαὶ, ὡς ἐπὶ θανάτῳ ἀπηγμένῳ· τοὺς οἴκους αὐτῶν ὀχυροῦσιν οἱ ἱερεῖς θυρώμασί τε καὶ κλείθροις καὶ μοχλοῖς, ὅπως ὑπὸ τῶν λῃστῶν μὴ συληθῶσι.

\vs{19}Λύχνους καίουσι, καὶ πλείους ἢ ἑαυτοῖς, ὧν οὐδένα δύνανται ἰδεῖν.
\vs{20}Ἔστι μὲν ὥσπερ δοκὸς τῶν ἐκ τῆς οἰκίας, τὰς δὲ καρδίας αὐτῶν φασιν ἐκλείχεσθαι τῶν ἀπὸ τῆς γῆς ἑρπετῶν, κατεσθόντων αὐτούς τε καὶ τὸν ἱματισμὸν αὐτῶν οὐκ αἰσθάνονται·
\vs{21}Μεμελανωμένοι τὸ πρόσωπον αὐτῶν ἀπὸ τοῦ καπνοῦ τοῦ ἐκ τῆς οἰκίας.
\vs{22}Ἐπὶ τὸ σῶμα αὐτῶν καὶ ἐπὶ τὴν κεφαλὴν αὐτῶν ἐφίπτανται νυκτερίδες, χελιδόνες, καὶ τὰ ὄρνεα, ὡσαύτως δὲ καὶ οἱ αἴλουροι.
\vs{23}Ὅθεν γνώσεσθε ὅτι οὐκ εἰσὶ θεοί· μὴ οὖν φοβεῖσθε αὐτά.

\vs{24}Τὸ γὰρ χρυσίον ὃ περίκεινται εἰς κάλλος, ἐὰν μή τις ἐκμάξῃ τὸν ἰὸν, οὐ μὴ στίλψωσιν, οὐδὲ γὰρ ὅτε ἐχωνεύοντο, ᾐσθάνοντο.
\vs{25}Ἐκ πάσης τιμῆς ἠγορασμένα ἐστὶν, ἐν οἷς οὐκ ἔστι πνεῦμα.
\vs{26}Ἄνευ ποδῶν ἐπʼ ὤμοις φέρονταὶ, ἐνδεικνύμενοι τὴν ἑαυτῶν ἀτιμίαν τοῖς ἀνθρώποις.

\vs{27}Αἰσχύνονταί τε καὶ οἱ θεραπεύοντες αὐτὰ, διὰ τὸ, εἴποτε ἐπὶ τὴν γῆν πέσῃ, μὴ διʼ αὐτῶν ἀνίστασθαι, μήτε ἐάν τις αὐτὸ ὀρθὸν στήσῃ, διʼ ἑαυτοῦ κινηθήσεται, μητε ἐὰν κλιθῇ, οὐ μὴ ὀρθωθῇ, ἀλλʼ ὥσπερ νεκροῖς τὰ δῶρα αὐτοῖς παρατίθεται.

\vs{28}Τὰς δὲ θυσίας αὐτῶν ἀποδόμενοι οἱ ἱερεῖς αὐτῶν καταχρῶνται· ὡσαύτως δὲ καὶ αἱ γυναῖκες ἀπʼ αὐτῶν ταριχεύουσαι, οὐτε πτωχῷ οὔτε ἀδυνάτῳ μὴ μεταδῶσι.
\vs{29}Τῶν θυσιῶν αὐτῶν ἀποκαθημένη καὶ λεχὼ ἅπτονται· γνόντες οὖν ἀπὸ τούτων ὅτι οὐκ εἰσὶ θεοὶ, μὴ φοβηθῆτε αὐτούς.
\vs{30}Πόθεν γὰρ κληθείησαν θεοί; ὅτι γυναῖκες παρατιθέασι θεοῖς ἀργυροῖς και χρυσοῖς καὶ ξυλίνοις.
\vs{31}Καὶ ἐν τοῖς οἴκοις αὐτῶν οἱ ἱερεῖς διφρεύουσιν, ἔχοντες τοὺς χιτῶνας διεῤῥωγότας, καὶ τὰς κεφαλὰς καὶ τοὺς πώγωνας ἐξυρημένους, ὧν αἱ κεφαλαὶ ἀκάλυπτοί εἰσιν.
\vs{32}Ὠρύονται δὲ βοῶντες ἐναντίον τῶν θεῶν αὐτῶν, ὥσπερ τινὲς ἐν περιδείπνῳ νεκροῦ.

\vs{33}Ἀπὸ τοῦ ἱματισμοῦ αὐτῶν ἀφελόμενοι οἱ ἱερεῖς, ἐνδύσουσι τὰς γυναῖκας αὐτῶν καὶ τὰ παιδία.
\vs{34}Οὔτε ἐὰν κακὸν πάθωσιν ὑπό τινος, οὔτε ἐὰν ἀγαθὸν, δυνήσονται ἀνταποδοῦναι· οὔτε καταστῆσαι βασιλέα δύνανται, οὔτε ἀφελέσθαι.
\vs{35}Ὡσαύτως οὔτε πλοῦτον οὔτε χαλκὸν οὐ μὴ δύνωνται διδόναι· ἐάν τις εὐχὴν αὐτοῖς εὐξάμενος μὴ ἀποδῷ, οὐ μὴ ἐπιζητήσωσιν.
\vs{36}Ἐκ θανάτου ἄνθρωπον οὐ μὴ ῥύσωνται, οὔτε ἥττονα ἀπὸ ἰσχυροῦ μὴ ἐξέλωνται.
\vs{37}Ἄνθρωπον τυφλὸν εἰς ὅρασιν οὐ μὴ περιστήσωσιν, ἐν ἀνάγκῃ ἄνθρωπον ὄντα οὐ μὴ ἐξέλωνται.
\vs{38}Χήραν οὐ μὴ ἐλεήσωσιν, οὔτε ὀρφανὸν εὖ ποιήσωσι.

\vs{39}Τοῖς ἀπὸ τοῦ ὄρους λίθοις ὡμοιωμὲνοι εἰσὶ τὰ ξύλινα, καὶ τὰ περίχρυσα, καὶ τὰ περιάργυρα, οἱ δὲ θεραπεύοντες αὐτὰ καταισχυνθήσονται.

\vs{40}Πῶς οὖν νομιστέον ἢ κλητέον ὑπάρχειν αὐτοὺς θεοὺς, ἔτι δὲ καὶ αὐτῶν τῶν Χαλδαίων ἀτιμαζόντων αὐτά;
\vs{41}Οἳ ὅταν ἴδωσιν ἐνεὸν μὴ δυνάμενον λαλῆσαι, προσενεγκάμενοι τὸν Βῆλον, ἀξιοῦσι φωνῆσαι, ὠς δυνατοῦ ὄντος αὐτοῦ αἰσθὲσθαι.
\vs{42}Καὶ οὐ δύνανται αὐτοὶ νοήσαντες καταλιπεῖν αὐτὰ, αἴσθησιν γὰρ οὐκ ἐχουσιν.

\vs{43}Αἱ δὲ γυναῖκες περιθέμεναι σχοινὶα, ἐν ταῖς ὁδοῖς ἐγκάθηνται, θυμιῶσαι τὰ πίτυρα· ὅτας δέ τις αὐτῶν ἐφελκυσθεῖσα ὑπό τινος τῶν παραπορευομένων κοιμηθῇ, τὴν πλησίον ὀνειδίζει, ὅτι οὐκ ἠξίωται ὥσπερ καὶ αὐτὴ, οὔτε τὸ σχοινίον αὐτῆς διεῤῥάγη.
\vs{44}Πάντα τὰ γενόμενα ἐν αὐτοῖς ἐστι ψευδῆ· πῶς οὖν νομιστέον ἢ κλητέον ὡς θεοὺς αὐτοὺς ὑπάρχειν;

\vs{45}Ὑπὸ τεκτόνων καὶ χρυσοχόων κατεσκευασμένα εἰσίν· οὐθὲν ἄλλο μὴ γὲνηται, ἢ ὃ βούλονται οἱ τεχνίται αὐτὰ γενέσθαι.
\vs{46}Αὐτοί τε οἱ κατασκευάζοντες αὐτὰ οὐ μὴ γένωνται πολυχρόνιοι· πῶς τε δὴ μέλλει τὰ ὑπʼ αὐτῶν κατασκευασθέντα;

\vs{47}Κατέλιπον γὰρ ψεύδη καὶ ὄνειδος τοῖς ἐπιγινομένοις.
\vs{48}Ὅταν γὰρ ἐπέλθῃ ἐπʼ αὐτὰ πόλεμος καὶ κακὰ, βουλεύονται πρὸς ἑαυτοὺς οἱ ἱερεῖς, ποῦ συναποκρυβῶσι μετʼ αὐτῶν.
\vs{49}Πῶς οὖν οὐκ ἔστιν αἰσθέσθαι ὅτι οὐκ εἰσί θεοί, οἳ οὔτε σώζουσιν ἑαυτοὺς ἐκ πολέμου, οὔτε ἐκ κακῶν;
\vs{50}Ὑπάρχοντα γὰρ ξύλινα καὶ περίχρυσα καὶ περιάργυρα, γνωσθήσεται μετὰ ταῦτα ὅτι ἐστὶ ψευδῆ.
\vs{51}Τοῖς ἔθνεσι πᾶσι τοῖς τε βασιλεῦσι φανερὸν ἔσται ὃτι οὐκ εἰσὶ θεοὶ, ἀλλὰ ἔργα χειρῶν ἀνθρώπων, καὶ οὐδὲν Θεοῦ ἔργον ἐν αὐτοῖς ἐστι.

\vs{52}Τίνι οὖν γνωστέον ἐστὶν ὅτι οὐκ εἰσὶ θεοί;
\vs{53}Βασιλὲα γὰρ χώρας οὐ μὴ ἀναστήσωσιν, οὔτε ὑετὸν ἀνθρώποις οὐ μὴ δῶσι.
\vs{54}Κρίσιν τε οὐ μὴ διακρίνωσιν ἑαυτῶν, οὐδὲ μὴ ῥύσωνται ἀδίκημα, ἀδὺνατοι ὄντες· ὥσπερ γὰρ κορῶναι ἀναμέσον τοῦ οὐρανοῦ καὶ τῆς γῆς.

\vs{55}Καὶ γὰρ ὅταν ἐμπέσῃ εἰς οἰκίαν θεῶν ξυλίνων ἡ περιχρύσων ἢ περιαργύρων πῦρ, οἱ μὲν ἱερεῖν φεύξονται καὶ διασωθήσονται, αὐτοὶ δὲ ὥσπερ δοκοὶ μέσοι κατακαυθήσονται.
\vs{56}Βασιλεῖ δὲ καὶ πολεμίοις οὐ μὴ ἀντιστῶσι· πῶς οὖν ἐκδεκτέον ἢ νομιστέον ὅτι εἰσὶ θεοί;
\vs{57}Οὔτε ἀπὸ κλεπτῶν, οὔτε ἀπὸ λῃστῶν οὐ μὴ διασωθῶσι θεοὶ ξύλινον, καὶ περιάργυροι, καὶ περίχρυσοι·
\vs{58}ὧν οἱ ἰσχύοντες περιελοῦνται τὸ χρυσίον καὶ τὸ ἀργύριον, καὶ τὸν ἱματισμὸν τὸν περικείμενον αὐτοῖς ἀπελεύσονται ἔχοντες, οὔτε ἑαυτοῖς οὐ μὴ βοηθήσωσιν.

\vs{59}Ὥστε κρεῖσσον εἶναι βασιλέα ἐπιδεικνύμενον τὴν ἑαυτοῦ ἀνδρείαν, ἢ σκεῦος ἐν οἰκίᾳ χρήσιμον ἐφʼ ᾧ κεχρήσεται ὁ κεκτημένος, ἢ οἱ ψευδεῖς θεοί· ἢ καὶ θύρα ἐν οἰκίᾳ διασώζουσα τὰ ἐν αὐτῇ ὄντα, ἢ οἱ ψευδεῖς θεοί· καὶ ξύλινος στύλος ἐν βασιλείοις, ἤ οἱ ψευδεῖς θεοί.

\vs{60}Ἥλιος μὲν γὰρ καὶ σελήνη καὶ ἄστρα ὄντα λαμπρὰ, καὶ ἀποστελλόμενα ἐπὶ χρείας, εὐήκοά εἰσιν.
\vs{61}Ὡσαύτως καὶ ἀστραπὴ ὅταν ἐπιφανῇ, εὔοπτός ἐστι· τὸ δʼ αὐτὸ καὶ πνεῦμα ἐν πάσῃ χώρᾳ πνεῖ.
\vs{62}Καὶ νεφέλαις ὅταν ἐπιταγῆ ὑπὸ τοῦ Θεοῦ ἐπιπορεύεσθαι ἐπιπορεύεσθαι ἐφʼ ὅλην τὴν οἰκουμένην, συντελοῦσι τὸ ταχθέν.
\vs{63}Τό, τε πῦρ ἐξαποσταλὲν ἄνωθεν ἐξαναλῶσαι ὄρη καὶ δρυμοὺς, ποιεῖ τὸ συνταχθὲν· ταῦτα δὲ οὔτε ταῖς εἰδέαις οὔτε ταῖς δυνάμεσιν αὐτῶν ἀφωμοιωμένα ἐστίν.

\vs{64}Ὅθεν οὔτε νομιστέον οὔτε κλητέον ὑπάρχειν αὐτοὺς θεοὺς, οὐ δυνατῶν ὄντων αὐτῶν οὔτε κλητέον κρίναι, οὔτε εὖ ποιῆσαι ἀνθρώποις.
\vs{65}Γνόντες οὖν ὅτι οὐκ εἰσὶ θεοί, μὴ φοβηθῆτε αὐτούς·

\vs{66}Οὔτε γὰρ βασιλεῦσιν οὐ μὴ καταράσωνται, οὔτε μὴ εὐλογήσωσι.
\vs{67}Σημεῖά τε ἐν ἔθνεσιν ἐν οὐρανῷ οὐ μὴ δείξωσιν, οὐδὲ ὡς ὁ ἥλιος λάμψουσιν, οὔτε φωτιοῦσιν ὡς ἡ σελήνη.
\vs{68}Τὰ θηρία αὐτῶν ἐστι κρείττω, ἃ δύνανται ἐκφυγόντα εἰς σκέπην ἑαυτὰ ὠφελῆσαι.
\vs{69}Κατʼ οὐδένα οὖν τρόπον ἡμῖν ἑστι φανερὸν ὅτι εἰσὶ θεοί· διὸ μὴ φοβηθῆτε αὐτούς.

\vs{70}Ὥσοερ γὰρ ἐν σικυηράτῳ προβασκάνιον οὐδὲν φυλάσσον, οὕτως οἱ θεοὶ αὐτῶν εἰσι ξύλινοι καὶ περίχρυσον καὶ περιάργυροι.
\vs{71}τὸν αὐτὸν τρόπον καὶ τῇ. ἐν κήπῳ ῥάμνῳ, ἐφʼ ἧς πᾶν ὄρνεον ἐπικάθηται, ὡσαύτως δὲ καὶ νεκρῷ ἐῤῥιμμένῳ ἐν σκότει ἀφωμοίωνται οἱ θεοὶ αὐτῶν ξύλινοι καὶ περίχρυσοι καὶ περιάργυροι.
\vs{72}Ἀπό τε τῆς πορφύρας καὶ τῆς μαρμάρου τῆς ἐπʼ αὐτοὺς σηπομένης γνωσθήσονται ὅτι οὐκ εἰσὶ θεοί· αὐτά τε ἐξ ὑστέρου βρωθήσονται, καὶ ἔσται ὄνειδος ἐν τῇ χώρᾳ.

\vs{73}Κρεῖσσον οὖ ἄνθρωπος δίκαιος οὐκ ἔχων εἴδωλα, ἔσται γὰρ μακρὰν ἀπὸ ὀνειδισμοῦ.


\def\book{ΣΩΣΑΝΝΑ}
\biblebook{ΣΩΣΑΝΝΑ}

 
\lettrine[lines=2, loversize=0.2, nindent=0em, findent=.25em]{\textcolor{bookheadingcolor}{Κ}}{ΑΙ} ἦν ἀνὴρ οἰκῶν ἐν Βαβυλῶνι, καὶ ὄνομα αὐτῷ Ἰωακείμ.
\vs{2}Καὶ ἔλαβε γυναῖκα ᾗ ὄνομα Σωσάννα, θυγάτηρ Χελκείου, καλὴ σφόδρα, καὶ φοβουμένη τὸν Κύριον.
\vs{3}Καὶ οἱ γονεῖς αὐτῆς δίκαιοι, καὶ ἐδίδαξαν τὴν θυγατέρα αὐτῶν κατὰ τὸν νόμον Μωυσῆ.
\vs{4}Καὶ ἦν Ἰωακεὶμ πλούσιος σφόδρα, καὶ ἦν αὐτῷ παράδεισος γειτνιῶν τῷ οἴκῳ αὐτοῦ· καὶ πρὸς αὐτὸν προσήγοντο οἱ Ἰουδαῖοι, διὰ τὸ εἶναι αὐτὸν ἐνδοξότερον πάντων.

\vs{5}Καὶ ἀπεδείχθησαν δύο πρεσβύτεροι ἐκ τοῦ λαοῦ κριταὶ ἐν τῷ ἐνιαυτῷ ἐκείνῳ, περὶ ὧν ἐλάλησεν ὁ δεσπότης, ὅτι ἐξῆλθεν ἀνομία ἐκ Βαβυλῶνος ἐκ πρεσβυτέρων κριτῶν, οἳ ἐδόκουν κυβερνᾷν τὸν λαόν.
\vs{6}Οὗτοι προσεκαρτέρουν ἐν τῇ οἰκίᾳ Ἰωακεὶμ, καὶ ἤρχοντο πρὸς αὐτοὺς πάντες οἱ κρινόμενοι.

\vs{7}Καὶ ἐγένετο ἡνίκα ἀπέτρεχεν ὁ λαὸς μέσον ἡμέρας, εἰσεπορεύετο Σωσάννα, καὶ περιεπάτει ἐν τῷ παραδείσῳ τοῦ ἀνδρὸς αὐτῆς.
\vs{8}Καὶ ἐθεώρουν αὐτὴν οἱ δύο πρεσβύτεροι καθʼ ἡμέραν εἰσπορευομένην, καὶ περιπατοῦσαν, καὶ ἐγένοντο ἐν ἐπιθυμίᾳ αὐτῆς,
\vs{9}καὶ διέστρεψαν τὸν ἑαυτῶν νοῦν, καὶ ἐξέκλιναν τοὺς ὀφθαλμοὺς αὐτῶν, τοῦ μὴ βλέπειν εἰς τὸν οὐρανὸν, μηδὲ μνημονεύειν κριμάτων δικαίων.
\vs{10}Καὶ ἦσαν ἀμφότεροι κατανενυγμένοι περὶ αὐτῆς, καὶ οὐκ ἀνήγγειλαν ἀλλήλοις τὴν ὀδύνην ἑαυτῶν·
\vs{11}Ὅτι ᾐσχύνοντο ἀναγγεῖλαι τὴν ἐπιθυμίαν αὐτῶν, ὅτι ἤθελον συγγενέσθαι αὐτῇ.
\vs{12}Καὶ παρετηροῦσαν φιλοτίμως καθʼ ἡμέραν ὁρᾷν αὐτήν.

\vs{13}Καὶ εἶπαν ἕτερος τῷ ἑτέρῳ, πορευθῶμεν δὴ εἰς οἶκον, ὅτι ἀρίστου ὥρα ἐστί.
\vs{14}Καὶ ἐξελθόντες διεχωρίσθησαν ἀπʼ ἀλλήλων, καὶ ἀνακάμφαντες ἦλθον ἐπιτοαυτὸ, καὶ ἀνετάζοντες ἀλλήλους τὴν αἰτίαν, ὡμολόγησαν τὴν ἐπιθυμίαν αὐτῶν· καὶ τότε κοινῇ συνετάξαντο καιρὸν, ὅτε αὐτὴν δυνήσονται εὑρεῖν μόνην.

\vs{15}Καὶ ἐγένετο ἐν τῷ παρατηρεῖν αὐτοὺς ἡμέραν εὔθετον, εἰσῆλθέ ποτε καθὼς χθὲς καὶ τρίτης ἡμέρας μετὰ δύο μόνων κορασίων, καὶ ἐπεθύμησε λούσασθαι ἐν τῷ παραδείσῳ, ὅτι καῦμα ἦν.
\vs{16}Καὶ οὐκ ἧν οὐδεὶς ἐκεῖ πλὴν οἱ δύο πρεσβύτεροι κεκρυμμένοι, καὶ παρατηροῦντες αὐτήν.
\vs{17}Καὶ εἶπε τοῖς κορασίοις, ἐνέγκατε δή μοι ἔλαιον καὶ σμήγματα, καὶ τας θύρας τοῦ παραδείσου κλείσατε, ὅπως λούσωμαι.

\vs{18}Καὶ ἐποίησαν καθὼς εἶπε, καὶ ἀπέκλεισαν τὰς θύρας τοῦ παραδείσου, καὶ ἐξῆλθαν κατὰ τὰς πλαγίας θύρας, ἐνέγκαι τὰ προστεταγμένα αὐταῖς, καὶ οὐκ εἴδοσαν τοὺς πρεσβυτέρους ὅτι ἦσαν κεκρυμμένοι.

\vs{19}Καὶ ἐγένετο ὡς ἐξήλθοσαν τὰ κοράσια, καὶ ἀνέστησαν οἱ δύο πρεσβύται, καὶ ἐπέδραμον αὐτῇ,
\vs{20}καὶ εἶπον, ἰδοὺ αἱ θύραι τοῦ παραδείσου κέκλεινται, καὶ οὐδεὶς θεωρεῖ ἡμᾶς, καὶ ἐν ἐπιθυμίᾳ σου ἐσμέν· διὸ συγκατάθου ἡμῖν, καὶ γενοῦ μεθʼ ἡμῶν.
\vs{21}Εἰ δὲ μὴ, καταμαρτυρήσομέν σου, ὅτι ἦν μετὰ σοῦ νεανίσκος, καὶ διὰ τοῦτο ἐξαπέστειλας τὰ κοράσια ἀπὸ σοῦ.

\vs{22}Καὶ ἀνεστέναξε Σωσάννα, καὶ εἶπε, στενά μοι πάντοθεν· ἐάν τε γὰρ τοῦτο πράξω, θάνατός μοι ἐστίν· ἐάν τε μὴ πράξω, οὐκ ἐκφεύξομαι τὰς χεῖρας ὑμῶν.
\vs{23}Αἱρετώτερόν μοι ἐστὶ μὴ πράξασαν ἐμπεσεῖν εἰς τὰς χεῖρας ὑμῶν, ἢ ἁμαρτεῖν ἐνώπιον Κυρίου.
\vs{24}Καὶ ἀνεβόησε φωνῇ μεγάλῃ Σωσάννα· ἐβόησαν δὲ καὶ οἱ δύο πρεσβύται κατέναντι αὐτῆς.

\vs{25}Καὶ δραμων ὁ εἷς, ἤνοιξε τὰς θύρας τοῦ παραδείσου.
\vs{26}Ὡς δὲ ἤκουσαν τὴν κραυγὴν ἐν τῷ παραδεισῳ οἱ ἐκ τῆς οἰκίας, εἰσεπήδησαν διὰ τῆς πλαγίας θύρας ἰδεῖν τὸ συμβεβηκὸς αὐτῇ.
\vs{27}Ἡνίκα δὲ εἶπα οἱ πρεσβύται τοὺς λόγους αὐτῶν, κατῃσχύνθησαν οἱ δοῦλοι σφόδρα, ὅτι πώποτε οὐκ ἐῤῥήθη λόγος τοιοῦτος περὶ Σωσάννης.

\vs{28}Καὶ ἐγένετο τῇ ἐπαύριον, ὡς συνῆλθεν ὁ λαὸς πρὸς τὸν ἄνδρα αὐτῆς Ἰωακεὶμ, ἦλθον οἱ δύο πρεσβύται πλήρεις τῆς ἀνόμου ἐννοίας κατὰ Σωσάννης, τοῦ θανατῶσαι αὐτὴν,
\vs{29}καὶ εἶπαν ἔμπροσθεν τοῦ λαοῦ, ἀποστείλατε ἐπὶ Σωσάνναν θυγατέρα Χελκίου, ἥ ἐστι γυνὴ Ἰωακείμ· οἱ δὲ ἀπέστειλαν.
\vs{30}Καὶ ἦλθεν αὐτὴ, καὶ οἱ γονεῖς αὐτῆς, καὶ τὰ τέκνα αὐτῆς, καὶ πάντες οἱ συγγενεῖς αὐτῆς.

\vs{31}Ἡ δὲ Σωσάννα ἦν τρυφερὰ σφόδρα, καὶ καλὴ τῷ εἴδει.
\vs{32}Οἱ δὲ παράνομοι ἐκέλευσαν ἀποκαλυφθῆναι αὐτὴν, ἦν γὰρ κατακεκαλυμμένη, ὅπως ἐμπλησθῶσι τοῦ κάλλους αὐτῆς.
\vs{33}Ἔκλαιον δὲ οἱ παρʼ αὐτῆς, καὶ πάντες οἱ ἰδόντες αὐτήν.

\vs{34}Ἀναστάντες δὲ οἱ δύο πρεσβύται ἐν μέσῳ τῷ λαῷ, ἔθηκαν τὰς χεῖρας ἐπὶ τὴν κεφαλὴν αὐτῆς.
\vs{35}Ἡ δὲ κλαίουσα ἀνέβλεψεν εἰς τὸν οὐρανὸν, ὅτι ἦν ἡ καρδία αὐτῆς πεποιθυῖα ἐπὶ Κυρίῳ.

\vs{36}Εἶπον δὲ οἱ πρεσβύται, περιπατούντων ἡμῶν ἐν τῷ παραδείσῳ μόνων, εἰσῆλθεν αὕτη μετὰ δύο παιδισκῶν, καὶ ἀπέκλεισε τὰς θύρας τοῦ παραδείσου, καὶ ἀπέλυσε τὰς παιδίσκας.
\vs{37}Καὶ ἦλθε πρὸς αὐτὴν νεανίσκος ὃς ἦν κεκρυμμένος, καὶ ἀνέπεσε μετʼ αὐτῆς.
\vs{38}Ἡμεῖς δὲ ὄντες ἐν τῇ γωνίᾳ τοῦ παραδείσου, ἰδόντες τὴν ἀνομίαν, ἐδράμομεν ἐπʼ αὐτούς.

\vs{39}Καὶ ἰδόντες συγγινομένους αὐτοὺς, ἐκείνου μὲν οὐκ ἠδυνήθημεν ἐγκρατεῖς γενέσθαι, διὰ τὸ ἰσχύειν αὐτὸν ὑπὲρ ἡμᾶς, καὶ ἀνοίξαντα τὰς θύρας ἐκπεπηδηκέναι,
\vs{40}Ταύτης δὲ ἐπιλαβόμενοι, ἐπηρωτῶμεν, τίς ἦν ὁ νεανίσκος· καὶ οὐκ ἠθέλησεν ἀγγεῖλαι ἡμῖν· ταῦτα μαρτυροῦμεν.
\vs{41}Καὶ ἐπίστευσεν αὐτοῖς ἡ συναγωγὴ ὡς πρεσβυτέροις τοῦ λαοῦ καὶ κριταῖς· καὶ κατέκριναν αὐτὴν ἀποθανεῖν.

\vs{42}Ἀνεβόησε δὲ φωνῇ μεγάλῃ Σωσάννα, καὶ εἶπεν, ὁ Θεὸς ὁ αἰώνιος, ὁ τῶν κρυπτῶν γνώστης, ὁ εἰδὼς τὰ πάντα πρὶν γενέσεως αὐτῶν,
\vs{43}σὺ ἐπίστασαι ὅτι ψευδῆ μου κατεμαρτύρησαν· καὶ ἰδοὺ ἀποθνήσκω μὴ ποιήσασα μηδὲν ὧν οὗτοι ἐπονηρεύσαντο κατʼ ἐμοῦ.
\vs{44}Καὶ εἰσήκουσε Κύριος τῆς φωνῆς αὐτῆς.

\vs{45}Καὶ ἀπαγομένης αὐτῆς ἀπολέσθαι, ὁ Θεὸς ἐξήγειρε τὸ πνεῦμα τὸ ἅγιον παιδαρίου νεωτέρου ᾧ ὄνομα Δανιήλ.
\vs{46}Καὶ ἐβόησε φωνῇ μεγάλῃ, ἀθῶος ἐγὼ ἀπὸ τοῦ αἵματος ταύτης.

\vs{47}Ἐπέστρεψε δὲ πᾶς ὁ λαὸς πρὸς αὐτὸν, καὶ εἶπαν, τίς ὁ λόγος οὗτος, ὃν σὺ λελάληκας;
\vs{48}Ὁ δὲ στὰς ἐν μέσῳ αὐτῶν, εἶπεν, οὕτως μωροὶ οἱ υἱοὶ Ἰσραήλ; οὐκ ἀνακρίναντες, οὐδὲ τὸ σαφὲς ἐπιγνόντες, κατεκρίνατε θυγατέρα Ἰσραήλ;
\vs{49}Ἀναστρέψατε εἰς τὸ κριτήριον, ψευδῆ γὰρ οὗτοι κατεμαρτύρησαν αὐτῆς.

\vs{50}Καὶ ἀνέστρεψε πᾶς ὁ λαὸς μετὰ σπουδῆς· καὶ εἶπαν αὐτῷ οἱ πρεσβύτεροι, δεῦρο κάθισον ἐν μέσῳ ἡμῶν, καὶ ἀνάγγειλον ἡμῖν, ὅτι σοὶ δέδωκεν ὁ Θεὸς τὸ πρεσβεῖον.
\vs{51}Καὶ εἶπε πρὸς αὐτοὺς Δανιὴλ, διαχωρίσατε αὐτοὺς ἀπʼ ἀλλήλων μακρὰν, καὶ ἀνακρινῶ αὐτούς.

\vs{52}Ὡς δὲ διεχωρίσθησαν εἷς ἀπὸ τοῦ ἑνὸς, ἐκάλεσε τὸν ἕνα αὐτῶν, καὶ εἶπε πρὸς αὐτὸν, πεπαλαιωμένε ἡμερῶν κακῶν, νῦν ἥκασιν αἱ ἁμαρτίαι σου, ἃς ἐποίεις τὸ πρότερον,
\vs{53}κρίνων κρίσεις ἀδίκους· καὶ τοὺς μὲν ἀθώους κατακρίνων, ἀπολύων δὲ τοὺς αἰτίους, λέγοντος τοῦ Θεοῦ, ἀθῶον καὶ δίκαιον οὐκ ἀποκτενεῖς.
\vs{54}Νῦν οὖν ταύτην εἴπερ εἶδες, εἰπὸν, ὑπὸ τί δένδρον εἶδες αὐτοὺς ὁμιλοῦντας ἀλλήλοις; ὁ δὲ εἶπεν, ὑπὸ σχῖνον.

\vs{55}Εἶπε δὲ Δανιὴλ, ὀρθῶς ἔψευσαι εἰς τὴν σεαυτοῦ κεφαλήν· ἤδη γὰρ ἄγγελος φάσιν Θεοῦ λαβὼν παρὰ τοῦ Θεοῦ, σχίσει σε μέσον.
\vs{56}Καὶ μεταστήσας αὐτὸν, ἐκέλευσε προσαγαγεῖν τὸν ἕτερον, καὶ εἶπεν αὐτῷ, σπέρμα Χαναὰν, καὶ οὐκ Ἰούδα, τὸ κάλλος ἐξηπάτησέ σε, καὶ ἐπιθυμία διέστρεψε τὴν καρδίαν σου.
\vs{57}Οὕτως ἑποιεῖτε θυγατράσιν Ἰσραὴλ, καὶ ἐκεῖναι φοβούμεναι ὡμίλουν ὑμῖν· ἀλλʼ οὐ θυγάτηρ Ἰούδα ὑπέμεινε τὴν ἀνομίαν ὑμῶν.
\vs{58}Νῦν οὖν λέγε μοι, ὑπὸ τί δένδρον κατέλαβες αὐτοὺς ὁμιλοῦντας ἀλλήλοις; ὁ δὲ εἶπεν, ὑπὸ πρίνον.

\vs{59}Εἶπε δὲ αὐτῷ Δανιὴλ, ὀρθῶς ἔψευσαι καὶ σὺ εἰς τὴν σεαυτοῦ κεφαλὴν· μένει γὰρ ὁ ἄγγελος τοῦ Θεοῦ, τὴν ῥομφαίαν ἔχων πρίσαι σε μέσον, ὅπως ἐξολοθρεύσῃ ὑμᾶς.

\vs{60}Καὶ ἀνεβόησε πᾶσα ἡ συναγωγὴ φωνῇ μεγάλῃ, καὶ εὐλόγησαν τῷ Θεῷ τῷ σώζοντι τοὺς ἐλπίζοντας ἐπʼ αὐτόν.
\vs{61}Καὶ ἀνέστησαν ἐπὶ τοὺς δύο πρεσβύτας, ὅτι συνέστησεν αὐτοὺς Δανιὴλ ἐκ τοῦ στόματος αὐτῶν ψευδομαρτυρήσαντας.
\vs{62}Καὶ ἐποίησαν αὐτοῖς ὃν τρόπον ἐπονηρεύσαντο τῷ πλησίον· ποιῆσαι κατὰ τὸν νόμον Μωυσῆ· καὶ ἀπέκτειναν αὐτοὺς, καὶ ἐσώθη αἷμα ἀναίτιον ἐν τῇ ἡμέρᾳ ἐκείνῃ.

\vs{63}Χελκίας δὲ καὶ ἡ γυνὴ αὐτοῦ ᾔνεσαν περὶ τῆς θυγατρὸς αὐτῶν μετὰ Ἰωακεὶμ τοῦ ἀνδρὸς αὐτῆς καὶ τῶν συγγενῶν αὐτῶν, ὅτι οὐχ εὑρέθη ἐν αὐτῇ ἄσχημον πρᾶγμα.
\vs{64}Καὶ Δανιὴλ ἐγένετο μέγας ἐνώπιον τοῦ λαοῦ ἀπὸ τῆς ἡμέρας ἐκείνης, καὶ ἐπέκεινα.


\def\book{ΒΗΛ ΚΑΙ ΔΡΑΚΩΝ}
\biblebook{ΒΗΛ ΚΑΙ ΔΡΑΚΩΝ}

 
\lettrine[lines=2, loversize=0.2, nindent=0em, findent=.25em]{\textcolor{bookheadingcolor}{Κ}}{ΑΙ} ὁ βασιλεὺς Ἀστυάγης προσετέθη πρὸς τοὺς πατέρας αὐτοῦ· καὶ παρέλαβε Κύρος ὁ Πέρσης τὴν βασιλείαν αὐτοῦ.
\vs{2}Καὶ ἦν Δανιὴλ συμβιωτὴς τοῦ βασιλέως, καὶ ἔνδοξος ὑπὲρ πάντας τοὺς φίλους αὐτοῦ.

\vs{3}Καὶ ἦν εἴδωλον τοῖς Βαβυλωνίοις ᾧ ὄνομα Βὴλ, καὶ ἐδαπανῶντο εἰς αὐτὸν ἑκάστης ἡμέρας σεμιδάλεως ἀρτάβαι δώδεκα, καὶ πρόβατα τεσσαράκοντα, καὶ οἴνου μετρηταὶ ἕξ.
\vs{4}Καὶ ὁ βασιλεὺς ἐσέβετο αὐτὸν, καὶ ἐπορεύετο καθʼ ἑκάστην ἡμέραν προσκυνεῖν αὐτῷ. Δανιὴλ δὲ προσεκύνει τῷ Θεῷ αὐτοῦ· καὶ εἶπεν αὐτῷ ὁ βασιλεὺς, διατί οὐ προσκυνεῖς τῷ Βήλ;
\vs{5}Ὁ δὲ εἶπεν, ὅτι οὐ σέβομαι εἴδωλα χειροποίητα, ἀλλὰ τὸν ζῶντα Θεὸν, τὸν κτίσαντα τὸν οὐρανὸν καὶ τὴν γῆν καὶ ἔχοντα πάσης σαρκὸς κυρείαν.

\vs{6}Καὶ εἶπεν αὐτῷ ὁ βασιλεὺς, οὐ δοκεῖ σοι Βὴλ εἶναι ζῶν θεός; ἢ οὐχ ὁρᾷς ὅσα ἐσθίει καὶ πίνει καθʼ ἑκάστην ἡμέραν;
\vs{7}Καὶ εἶπε Δανιὴλ γελάσας, μὴ πλανῶ, βασιλεῦ, οὗτος γὰρ ἔσωθεν μέν ἐστι πηλὸς, ἔξωθεν δὲ χαλκὸς, καὶ οὐ βέβρωκεν οὐδέποτε.

\vs{8}Θυμωθεὶς δὲ ὁ βασιλεὺς ἐκάλεσε τοὺς ἱερεῖς αὐτοῦ· καὶ εἶπεν αὐτοῖς ἐὰν μὴ εἴποιτέ μοι τίς ὁ κατέσθων τὴν δαπάνην ταύτην, ἀποθανεῖσθε.
\vs{9}Ἐὰν δὲ δείξητε ὅτι Βὴλ κατεσθίει αὐτὰ, ὁ Δανιὴλ ἀποθανεῖται, ὅτι ἐβλασφήμησεν εἰς τὸν Βήλ· καὶ εἶπε Δανιὴλ τῷ βασιλεῖ, γινέσθω κατὰ τὸ ῥῆμά σου.

\vs{10}Καὶ ἦσαν ἱερεῖς τοῦ Βὴλ ἑβδομήκοντα ἐκτὸς γυναικῶν καὶ τέκνων· καὶ ἦλθεν ὁ βασιλεὺς μετὰ Δανιὴλ εἰς τὸν οἶκον τοῦ Βήλ.
\vs{11}Καὶ εἶπαν οἱ ἱερεῖς τοῦ Βὴλ, ἰδοὺ ἡμεῖς ἀποτρέχομεν ἔξω, σὺ δὲ, βασιλεῦ, παράθες τὰ βρώματα, καὶ τὸν οἶνον κεράσας θὲς, καὶ ἀπόκλεισον τὴν θύραν, καὶ σφράγισον τῷ δακτυλίῳ σου.
\vs{12}Καὶ ἐλθὼν πρωῒ, ἐὰν μὴ εὕρῃς πάντα βεβρωμένα ὑπὸ τοῦ Βὴλ, ἀποθανούμεθα· ἢ Δανιὴλ ὁ ψευδόμενος καθʼ ἡμῶν.
\vs{13}Αὐτοὶ δὲ κατεφρόνουν, ὅτι πεποιήκεισαν ὑπὸ τὴν τράπεζαν κεκρυμμένην εἴσοδον, καὶ διʼ αὐτῆς εἰσεπορεύοντο διόλου, καὶ ἀνήλουν αὐτά.

\vs{14}Καὶ ἐγένετο ὡς ἐξήλθοσαν ἐκεῖνοι, καὶ ὁ βασιλεὺς παρέθηκε τὰ βρώματα τῷ Βήλ· καὶ ἐπέταξε Δανιὴλ τοῖς παιδαρίοις αὐτοῦ, καὶ ἤνεγκαν τέφραν· καὶ κατέσεισαν ὅλον τὸν ναὸν ἐνώπιον τοῦ βασιλέως μόνου· καὶ ἐξελθόντες ἔκλεισαν τὴν θύραν, καὶ ἐσφραγίσαντο ἐν τῷ δακτυλίῳ τοῦ βασιλέως, καὶ ἀπῆλθον.
\vs{15}Οἱ δὲ ἱερεῖς ἦλθον τὴν νύκτα κατὰ τὸ ἔθος αὐτῶν, καὶ αἱ γυναῖκες αὐτῶν, καὶ τὰ τέκνα αὐτῶν, καὶ κατέφαγον πάντα, καὶ ἐξέπιον.

\vs{16}Καὶ ὤρθρισεν ὁ βασιλεὺς τὸ πρωῒ, καὶ Δανιὴλ μετʼ αὐτοῦ.
\vs{17}Καὶ εἶπε, σῶοι αἱ σφραγίδες Δανιήλ; ὁ δὲ εἶπε, σῶοι, βασιλεῦ.
\vs{18}Καὶ ἐγένετο ἅμα τῷ ἀνοῖξαι τὰς θύρας, ἐπιβλέψας ἐπὶ τὴν τράπεζαν ὁ βασιλεὺς, ἐβόησε φωνῇ μεγάλῃ, μέγας εἶ Βὴλ, καὶ οὐκ ἔστι παρὰ σοὶ δόλος οὐδὲ εἷς.

\vs{19}Καὶ ἐγέλασε Δανιὴλ, καὶ ἐκράτησε τὸν βασιλέα, τοῦ μὴ εἰσελθεῖν αὐτὸν ἔσω· καὶ εἶπεν, ἴδε δὴ τὸ ἔδαφος, καὶ γνῶθι τίνος τὰ ἴχνη ταῦτα.
\vs{20}Καὶ εἶπεν ὁ βασιλεὺς, ὁρῶ τὰ ἴχνη ἀνδρῶν, καὶ γυναικῶν, καὶ παιδίων· καὶ ὀργισθεὶς ὁ βασιλεὺς τότε συνέλαβε τοὺς ἱερεῖς, καὶ τὰς γυναῖκας,
\vs{21}καὶ τὰ τέκνα αὐτῶν, καὶ ἔδειξαν αὐτῷ τὰς κρυπτὰς θύρας, διʼ ὧν εἰσεπορεύοντο, καὶ ἐδαπάνων τὰ ἐπὶ τῆς τραπέζης.
\vs{22}Καὶ ἀπέκτεινεν αὐτοὺς ὁ βασιλεὺς, καὶ ἔδωκε τὸν Βὴλ ἔκδοτον τῷ Δανιήλ· καὶ κατέστρεψεν αὐτὸν καὶ τὸ ἱερὸν αὐτοῦ.

\vs{23}Καὶ ἦν Δράκων μέγας, καὶ ἐσέβοντο αὐτὸν οἱ Βαβυλώνιοι.
\vs{24}Καὶ εἶπεν ὁ βασιλεὺς τῷ Δανιὴλ, μὴ καὶ τοῦτον ἐρεῖς ὅτι χαλκοῦς ἐστιν; ἰδοὺ ζῇ, καὶ ἐσθίει, καὶ πίνει· οὐ δύνασαι εἰπεῖν, ὅτι οὐκ ἔστιν οὗτος θεὸς ζῶν· καὶ προσκύνησον αὐτῷ.

\vs{25}Καὶ εἶπε Δανιὴλ, Κυρίῳ τῷ Θεῷ μου προσκυνήσω, ὅτι οὗτός ἐστι Θεὸς ζῶν.
\vs{26}Σὺ δὲ, βασιλεῦ, δός μοι ἐξουσίαν, καὶ ἀποκτενῶ τὸν δράκοντα ἄνευ μαχαίρας καὶ ῥάβδου· καὶ εἶπεν ὁ βασιλεὺς δίδωμί σοι.
\vs{27}Καὶ ἔλαβεν ὁ Δανιὴλ πίσσαν καὶ στέαρ καὶ τρίχας, καὶ ἥψησεν ἐπιτοαυτό· καὶ ἐποίησε μάζας, καὶ ἔδωκεν εἰς τὸ στόμα τοῦ δράκοντος, καὶ φαγὼν διεῤῥάγη ὁ δράκων· καὶ εἶπεν, ἴδετε τὰ σεβάσματα ὑμῶν.

\vs{28}Καὶ ἐγένετο, ὡς ἤκουσαν οἱ Βαβυλώνιοι, ἠγανάκτησαν λίαν, καὶ συνεστράφησαν ἐπὶ τὸν βασιλέα, καὶ εἶπαν, Ἰουδαῖος γέγονεν ὁ βασιλεὺς, τὸν Βὴλ κατέσπασε, καὶ τὸν δράκοντα ἀπέκτεινε, καὶ τοὺς ἱερεῖς κατέσφαξε.
\vs{29}Καὶ εἶπαν ἐλθόντες πρὸς τὸν βασιλέα, παράδος ἡμῖν τὸν Δανιήλ· εἰ δὲ μὴ, ἁποκτενοῦμέν σε, καὶ τὸν οἶκόν σου.

\vs{30}Καὶ εἶδεν ὁ βασιλεὺς ὅτι ἐπείγουσιν αὐτὸν σφόδρα, καὶ ἀναγκασθεὶς ὁ βασιλεὺς παρέδωκεν αὐτοῖς τὸν Δανιήλ.
\vs{31}Οἱ δὲ ἔβαλον αὐτὸν εἰς τὸν λάκκον τῶν λεόντων, καὶ ἦν ἐκεῖ ἡμέρας ἕξ.
\vs{32}Ἦσαν δὲ ἐν τῷ λάκκῳ ἑπτὰ λέοντες, καὶ ἐδίδοτο αὐτοῖς τὴν ἡμέραν δύο σώματα καὶ δύο πρόβατα· τότε δὲ οὐκ ἐδόθη αὐτοῖς, ἵνα καταφάγωσι τὸν Δανιήλ.

\vs{33}Καὶ ἦν Ἀμβακοὺμ ὁ προφήτης ἐν τῇ Ἰουδαίᾳ, καὶ αὐτὸς ἥψησεν ἕψεμα, καὶ ἐνέθρυψεν ἄρτους εἰς σκάφην, καὶ ἑπορεύετο εἰς τὸ πεδίον ἀπενέγκαι τοῖς θερισταῖς.
\vs{34}Καὶ εἶπεν ὁ ἄγγελος Κυρίου τῷ Ἀμβακοὺμ, ἀπένεγκε τὸ ἄριστον ὃ ἔχεις εἰς Βαβυλῶνα τῷ Δανιὴλ εἰς τὸν λάκκον τῶν λεόντων.

\vs{35}Καὶ εἶπεν Ἀμβακοὺμ, Κύριε, Βαβυλῶνα οὐχ ἑώρακα, καὶ τὸν λάκκον οὐ γινώσκω.
\vs{36}Καὶ ἐπελάβετο ὁ ἄγγελος Κυρίου τῆς κορυφῆς αὐτοῦ, καὶ βαστάσας τῆς κόμης τῆς κεφαλῆς αὐτοῦ, ἔθηκεν αὐτὸν εἰς Βαβυῶνα ἐπάνω τοῦ λάκκου, ἐν τῷ ῥοίζῳ τοῦ πνεύματος αὐτοῦ.
\vs{37}Καὶ ἐβόησεν Ἀμβακοὺμ, λέγων, Δανιὴλ, Δανιὴλ, λάβε τὸ ἄριστον ὃ ἀπέστειλέ σοι ὁ Θεός.

\vs{38}Καὶ εἶπε Δανιὴλ, ἐμνήσθης γάρ μου ὁ Θεὸς, καὶ οὐκ ἐγκατέλιπες τοὺς ἀγαπῶντάς σε.
\vs{39}Καὶ ἀναστὰς δανιὴλ, ἔφαγεν· ὁ δὲ ἄγγελος τοῦ ἀπεκατέστησε τὸν Ἀμβακοὺμ παραχρῆμα εἰς τὸν τόπον αὐτοῦ.

\vs{40}Ὁ δὲ βασιλεὺς ἦλθε τῇ ἑβδόμῃ πενθῆσαι τὸν Δανιὴλ, καὶ ἦλθεν ἐπὶ τὸν λάκκον, καὶ ἐνέβλεψε, καὶ ἰδοὺ, Δανιὴλ καθήμενος.
\vs{41}Καὶ ἀναβοήσας φωνῇ μεγάλῃ, εἶπε, μέγας εἰ, Κύριε ὁ Θεὸς τοῦ Δανιὴλ, καὶ οὐκ ἔστιν ἄλλος πλὴν σοῦ.
\vs{42}Καὶ ἀνέσπασεν αὐτόν· τοὺς δὲ αἰτίους τῆς ἀπωλείας αὐτοῦ ἐνέβαλεν εἰς τὸν λάκκον· καὶ κατεβρώθησαν παραχρῆμα ἐνώπιον αὐτοῦ.


\def\book{ΜΑΚΚΑΒΑΙΩΝ Αʹ}
\biblebook{ΜΑΚΚΑΒΑΙΩΝ Αʹ}


\lettrine[lines=2, loversize=0.2, nindent=0em, findent=.25em]{\textcolor{bookheadingcolor}{Κ}}{ΑΙ} ἐγένετο μετὰ τὸ πατάξαι Ἀλέξανδρον τὸν Φιλίππου τὸν Μακεδόνα, ὃς ἐξῆλθεν ἐκ τῆς γῆς Χεττειεὶμ, καὶ ἐπάταξε τὸν Δαρεῖον βασιλέα Περσῶν καὶ Μήδων, καὶ ἐβασίλευσεν ἀντʼ αὐτοῦ πρότερος ἐπὶ τὴν Ἑλλάδα.
\vs{2}Καὶ συνεστήσατο πολέμους πολλοὺς, καὶ ἐκράτησεν ὀχυρωμάτων πολλῶν, καὶ ἔσφαξε βασιλεῖς τῆς γῆς.
\vs{3}Καὶ διῆλθεν ἕως ἄκρων τῆς γῆς, καὶ ἔλαβε σκῦλα πλήθους ἐθνῶν· καὶ ἠσύχασεν ἡ γῆ ἐνώπιον αὐτοῦ· καὶ ὑψώθη, καὶ ἐπῄρθη ἡ καρδία αὐτοῦ.
\vs{4}Καὶ συνήγαγε δύναμιν ἰσχυρὰν σφόδρα, καὶ ἦρξε χωρῶν, καὶ ἐθνῶν, καὶ τυράννων, καὶ ἐγένοντο αὐτῷ εἰς φόρον.

\vs{5}Καὶ μετὰ ταῦτα ἔπεσεν ἐπὶ τὴν κοίτην, καὶ ἔγνω ὅτι ἀποθνήσκει.
\vs{6}Καὶ ἐκάλεσε τοὺς παῖδας αὐτοῦ τοὺς ἐνδόξους τοὺς συντρόφους αὐτοῦ ἀπὸ νεότητος, καὶ διεῖλεν αὐτοῖς τὴν βασιλείαν αὐτοῦ ἔτι ζῶντος αὐτοῦ.
\vs{7}Καὶ ἐβασίλευσεν Ἀλέξανδρος ἔτη δώδεκα, καὶ ἀπέθανε.
\vs{8}Καὶ ἐπεκράτησαν οἱ παῖδες αὐτοῦ ἕκαστος ἐν τῷ τόπῳ αὐτοῦ.
\vs{9}Καὶ ἐπέθεντο πάντες διαδήματα μετὰ τὸ ἀποθανεῖν αὐτὸν, καὶ οἱ υἱοὶ αὐτῶν ὀπίσω αὐτῶν ἔτη πολλὰ, καὶ ἐπλήθυναν κακὰ ἐν τῇ γῇ.

\vs{10}Καὶ ἐξῆλθεν ἐξ αὐτῶν ῥίζα ἁμαρτωλὸς Ἀντίοχος Ἐπιφανὴς, υἱὸς Ἀντιόχου βασιλέως, ὃς ἦν ὅμηρα ἐν τῇ Ῥώμῃ· καὶ ἐβασίλευσεν ἐν ἔτει ἑκατοστῷ καὶ τριακοστῷ καὶ ἑβδόμῳ βασιλείας Ἑλλήνων.

\vs{11}Ἐν ταῖς ἡμέραις ἐκείναις ἐξῆλθον ἐξ Ἰσραὴλ υἱοὶ παράνομοι, καὶ ἀνέπεισαν πολλοὺς, λέγοντες, πορευθῶμεν, καὶ διαθώμεθα διαθήκην μετὰ τῶν ἐθνῶν τῶν κύκλῳ ἡμῶν, ὅτι ἀφʼ ἧς ἐχωρίσθημεν ἀπʼ αὐτῶν, εὗρεν ἡμᾶς κακὰ πολλά.
\vs{12}Καὶ ἠγαθύνθη ὁ λόγος ἐν ὀφθαλμοῖς αὐτῶν.
\vs{13}Καὶ προεθυμήθησάν τινες ἀπὸ τοῦ λαοῦ, καὶ ἐπορεύθησαν πρὸς τὸν βασιλέα· καὶ ἔδωκεν αὐτοῖς ἐξουσίαν ποιῆσαι τὰ δικαιώματα τῶν ἐθνῶν.
\vs{14}Καὶ ᾠκοδόμησαν γυμνάσιον ἐν Ἱεροσολύμοις κατὰ τὰ νόμιμα τῶν ἐθνῶν.
\vs{15}Καὶ ἐποίησαν ἑαυτοῖς ἀκροβυστίας, καὶ ἀπέστησαν ἀπὸ διαθήκης ἁγίας· καὶ ἐζεύχθησαν τοῖς ἔθνεσι, καὶ ἐπράθησαν τοῦ ποιῆσαι τὸ πονηρόν.

\vs{16}Καὶ ἡτοιμάσθη ἡ βασιλεία ἐναντίον Ἀντιόχου· καὶ ὑπέλαβε βασιλεῦσαι τῆς Αἰγύπτου, ὅπως βασιλεύσῃ ἐπὶ τὰς δύο βασιλείας.
\vs{17}Καὶ εἰσῆλθεν εἰς Αἴγυπτον ἐν ὄχλῳ βαρεῖ, ἐν ἅρμασι, καὶ ἐν ἐλέφασι, καὶ ἐν ἱππεῦσι, καὶ ἐν στόλῳ μεγάλῳ.
\vs{18}Καὶ συνεστήσατο πόλεμον πρὸς Πτολεμαῖον βασιλέα Αἰγύπτου· καὶ ἐνετράπη Πτολεμαῖος ἀπὸ προσώπου αὐτοῦ, καὶ ἔφυγε· καὶ ἔπεσον τραυματίαι πολλοί.

\vs{19}Καὶ κατελάβοντο τὰς πόλεις τὰς ὀχυρὰς ἐν γῇ Αἰγύπτῳ· καὶ ἔλαβε τὰ σκῦλα γῆς Αἰγύπτου.

\vs{20}Καὶ ἐπέστρεψεν Ἀντιόχος μετὰ τὸ πατάξαι Αἴγυπτον ἐν τῷ ἑκατοστῷ καὶ τεσσαρακοστῷ καὶ τρίτῳ ἔτει· καὶ ἀνέβη ἐπὶ Ἰσραὴλ, καὶ ἀνέβη εἰς Ἱερουσαλὴμ ἐν ὄχλῳ βαρεῖ.
\vs{21}Καὶ εἰσῆλθεν εἰς τὸ ἁγίασμα ἐν ὑπερηφανείᾳ, καὶ ἔλαβε τὸ θυσιαστήριον τὸ χρυσοῦν, καὶ τὴν λυχνίαν τοῦ φωτὸς, καὶ πάντα τὰ σκεύη αὐτῆς,
\vs{22}καὶ τὴν τράπεζαν τῆς προθέσεως, καὶ τὰ σπονδεῖα, καὶ τὰς φιάλας, καὶ τὰς θυΐσκας τὰς χρυσᾶς, καὶ τὸ καταπέτασμα, καὶ τοὺς στεφάνους, καὶ τὸν κόσμον τὸν χρυσοῦν τὸν κατὰ πρόσωπον τοῦ ναοῦ, καὶ ἐλέπισε πάντα.
\vs{23}Καὶ ἔλαβε τὸ ἀργύριον, καὶ τὸ χρυσίον, καὶ τὰ σκεύη τὰ ἐπιθυμητά· καὶ ἔλαβε τοὺς θησαυροὺς τοὺς ἀποκρύφους οὓς εὗρε.

\vs{24}Καὶ λαβὼν πάντα ἀπῆλθεν εἰς τὴν γῆν αὐτοῦ· καὶ ἐποίησε φονοκτονίαν, καὶ ἐλάλησεν ὑπερηφανείαν μεγάλην.
\vs{25}Καὶ ἐγένετο πένθος μέγα ἐπὶ Ἰσραὴλ ἐν παντὶ τόπῳ αὐτῶν.
\vs{26}Καὶ ἐστέναξαν ἄρχοντες καὶ πρεσβύτεροι, παρθένοι καὶ νεανίσκοι ἠσθένησαν, καὶ τὸ κάλλος τῶν γυναικῶν ἠλλοιώθη.
\vs{27}Πᾶς νυμφίος ἀνέλαβε θρῆνον, καὶ καθημένη ἐν παστῷ ἐγένετο ἐν πένθει.
\vs{28}Καὶ ἐσείσθη ἡ γῆ ἐπὶ τοὺς κατοικοῦντας αὐτήν· καὶ πᾶς ὁ οἶκος Ἰακὼβ ἐνεδύσατο αἰσχύνην.

\vs{29}Καὶ μετὰ δύο ἔτη ἡμερῶν ἀπέστειλεν ὁ βασιλεὺς ἄρχοντα φορολογίας εἰς τὰς πόλεις Ἰούδα· καὶ ἦλθεν εἰς Ἱερουσαλὴμ ἐν ὄχλῳ βαρεῖ.
\vs{30}Καὶ ἐλάλησεν αὐτοῖς λόγους εἰρηνικοὺς ἐν δόλῳ· καὶ ἐνεπίστευσαν αὐτῷ· καὶ ἐπέπεσεν ἐπὶ τὴν πόλιν ἐξάπινα, καὶ ἐπάταξεν αὐτὴν πληγὴν μεγάλην, καὶ ἀπώλεσε λαὸν πολὺν ἐξ Ἰσραήλ.
\vs{31}Καὶ ἔλαβε τὰ σκῦλα τῆς πόλεως, καὶ ἐνεπύρισεν αὐτὴν πυρὶ, καὶ καθεῖλε τοὺς οἴκους αὐτῆς καὶ τὰ τείχη αὐτῆς κύκλῳ.
\vs{32}Καὶ ᾐχμαλώτευσαν τὰς γυναῖκας καὶ τὰ τέκνα, καὶ τὰ κτήνη ἐκληρονόμησαν.

\vs{33}Καὶ ᾠκοδόμησαν τὴν πόλιν Δαυὶδ τείχει μεγάλῳ καὶ ἰσχυρῷ, πύργοις ὀχυροῖς, καὶ ἐγένετο αὐτοῖς εἰς ἄκραν.
\vs{34}καὶ ἔθηκαν ἐκεῖ ἔθνος ἁμαρτωλὸν, ἄνδρας παρανόμους, καὶ ἐνίσχυσαν ἐν αὐτῇ.
\vs{35}Καὶ παρέθεντο ὅπλα καὶ τροφὰς, καὶ συναγαγόντες τὰ σκῦλα Ἱερουσαλὴμ ἀπέθεντο ἐκεῖ· καὶ ἐγένοντο εἰς μεγάλην παγίδα.
\vs{36}Καὶ ἐγένετο εἰς ἔνεδρον τῷ ἁγιάσματι, καὶ εἰς διάβολον πονηρὸν τῷ Ἰσραὴλ διαπαντός.

\vs{37}Καὶ ἐξέχεαν αἷμα ἀθῷον κύκλῳ τοῦ ἁγιάσματος, καὶ ἐμόλυναν τὸ ἁγίασμα.
\vs{38}Καὶ ἔφυγον οἱ κάτοικοι Ἱερουσαλὴμ διʼ αὐτοὺς, καὶ ἐγένετο κατοικία ἀλλοτρίων· καὶ ἐγένετο ἀλλοτρία τοῖς γεννήμασιν αὐτῆς, καὶ τὰ τέκνα αὐτῆς ἐγκατέλιπον αὐτήν.
\vs{39}Τὸ ἁγίασμα αὐτῆς ἠρημώθη ὡς ἔρημος, αἱ ἑορταὶ αὐτῆς ἐστράφησαν εἰς πένθος, τὰ σάββατα αὐτῆς εἰς ὀνειδισμὸν, ἡ τιμὴ αὐτῆς εἱς ἐξουδένωσιν.
\vs{40}Κατὰ τὴν δόξαν αὐτῆς ἐπληθύνθη ἡ ἀτιμία αὐτῆς, καὶ τὸ ὕψος αὐτῆς ἐστράφη εἰς πένθος.

\vs{41}Καὶ ἔγραψεν ὁ βασιλεὺς Ἀντίοχος πάσῃ τῇ βασιλείᾳ αὐτοῦ εἶναι πάντας λαὸν ἕνα,
\vs{42}καὶ ἐγκαταλιπεῖν ἕκαστον τὰ νόμιμα αὐτοῦ· καὶ ἐπεδέξατο πάντα τὰ ἔθνη κατὰ τὸν λόγον τοῦ βασιλέως.
\vs{43}Καὶ πολλοὶ ἀπὸ Ἰσραὴλ εὐδόκησαν τῇ λατρείᾳ αὐτοῦ, καὶ ἔθυσαν τοῖς εἰδώλοις, καὶ ἐβεβήλωσαν τὸ σάββατον.

\vs{44}Καὶ ἀπέστειλεν ὁ βασιλεὺς βιβλία ἐν χειρὶ ἀγγέλων εἰς Ἱερουσαλὴμ καὶ τὰς πόλεις Ἰούδα, πορευθῆναι ὀπίσω νομίμων ἀλλοτρίων τῆς γῆς,
\vs{45}καὶ κωλῦσαι ὁλοκαυτώματα καὶ θυσίαν καὶ σπονδὴν ἐκ τοῦ ἁγιάσματος, καὶ βεβηλῶσαι σάββατα καὶ ἑορτὰς,
\vs{46}καὶ μιᾶναι ἁγίασμα καὶ ἁγίους·
\vs{47}οἰκοδομῆσαι βωμοὺς, καὶ τεμένη, καὶ εἰδωλεῖα, καὶ θύειν ὕεια, καὶ κτήνη κοινὰ,
\vs{48}καὶ ἀφιέναι τοὺς υἱοὺς αὐτῶν ἀπεριτμήτους, βδελύξαι τὰς ψυχὰς αὐτῶν ἐν παντὶ ἀκαθάρτῳ καὶ βεβηλώσει,
\vs{49}ὥστε ἐπιλαθέσθαι τοῦ νόμου, καὶ ἀλλάξαι πάντα τὰ δικαιώματα.

\vs{50}Καὶ ὃς ἂν μὴ ποιήσῃ κατὰ τὸ ῥῆμα τοῦ βασιλέως, ἀποθανεῖται.
\vs{51}Κατὰ πάντας τοὺς λόγους τούτους ἔγραψε πάσῃ τῇ βασιλείᾳ αὐτοῦ, καὶ ἐποίησεν ἐπισκόπους ἐπὶ πάντα τὸν λαόν· καὶ ἐνετείλατο ταῖς πόλεσιν Ἰούδα θυσιάζειν κατὰ πόλιν καὶ πόλιν.
\vs{52}Καὶ συνηθροίσθησαν ἀπὸ τοῦ λαοῦ πρὸς αὐτοὺς πολλοὶ, πᾶς ὁ ἐγκαταλιπὼν τὸν νόμον· καὶ ἐποιησαν κακὰ ἐν τῇ γῇ.
\vs{53}Καὶ ἔθεντο τὸν Ἰσραὴλ ἐν κρύφοις ἐν παντὶ φυγαδευτηρίῳ αὐτῶν.

\vs{54}Καὶ τῇ πεντεκαιδεκάτῃ ἡμέρᾳ Χασελεῦ, τῷ πέμπτῳ καὶ τεσσαρακοστῷ καὶ ἑκατοστῷ ἔτει, ᾠκοδόμησαν βδέλυγμα ἐρημώσεως ἐπὶ τὸ θυσιαστήριον, καὶ ἐν πόλεσιν Ἰούδα κύκλῳ ᾠκοδόμησαν βωμούς.
\vs{55}Καὶ ἐπὶ τῶν θυρῶν τῶν οἰκιῶν, καὶ ἐν ταῖς πλατείαις ἐθυμίων.

\vs{56}Καὶ τὰ βιβλία τοῦ νόμου ἃ εὗρον, ἐνεπύρισαν πυρὶ κατασχίσαντες.
\vs{57}Καὶ ὅπου εὑρίσκετο παρά τινι βιβλίον διαθήκης, καὶ εἴ τις συνευδόκει τῷ νόμῳ, τὸ σύγκριμα τοῦ βασιλέως ἐθανάτου αὐτόν.
\vs{58}Ἐν ἰσχύϊ αὐτῶν ἐποίουν οὕτως τῷ Ἰσραὴλ τοῖς εὑρισκομένοις ἐν παντὶ μηνὶ καὶ μηνὶ ἐν ταῖς πόλεσι.
\vs{59}Καὶ τῇ πέμπτῃ καὶ εἰκάδι τοῦ μηνὸς θυσιάζοντες ἐπὶ τὸν βωμὸν ὃς ἦν ἐπὶ τοῦ θυσιαστηρίου.

\vs{60}Καὶ τὰς γυναῖκας τὰς περιτετμηκυίας τὰ τέκνα αὐτῶν ἐθανάτωσαν, κατὰ τὸ πρόσταγμα.
\vs{61}Καὶ ἐκρέμασαν τὰ βρέφη ἐκ τῶν τραχήλων αὐτῶν, καὶ τοὺς οἴκους αὐτῶν προενόμευσαν, καὶ τοὺς περιτετμηκότας αὐτοὺς ἐθανάτωσαν.
\vs{62}Καὶ πολλοὶ ἐν Ἰσραὴλ ἐκραταιώθησαν, καὶ ὠχυρώθησαν ἐν ἑαυτοῖς τοῦ μὴ φαγεῖν κοινά.
\vs{63}Καὶ ἐπελέξαντο ἀποθανεῖν, ἵνα μὴ μιανθῶσι τοῖς βρώμασι, καὶ μὴ βεβηλώσωσι διαθήκην ἁγίαν· καὶ ἀπέθανον.
\vs{64}Καὶ ἐγένετο ὀργὴ μεγάλη ἐπὶ Ἰσραὴλ σφόδρα.

\ch{2}
Ἐν ταῖς ἡμέραις ἐκείναις ἀνέστη Ματταθίας Ἰωάννου τοῦ Συμεὼν, ἱερεὺς τῶν υἱῶν Ἰωαρὶβ ἀπὸ Ἱερουσαλὴμ, καὶ ἐκάθισεν ἐν Μωδεΐν.
\vs{2}Καὶ αὐτῷ υἱοὶ πέντε, Ἰωαννὰν ὁ ἐπικαλούμενος Καδδὶς,
\vs{3}Σίμων ὁ καλούμενος Θασσι,
\vs{4}Ἰούδας ὁ ἐπικαλούμενος Μακκαβαῖος,
\vs{5}Ἐλεάζαρ ὁ ἐπικαλούμενος Αὐαρὰν, Ἰωνάθαν ὁ ἐπικαλούμενος Ἀπφοῦς.

\vs{6}Καὶ εἶδε τὰς βλασφημίας τὰς γινομένας ἐν Ἰούδα καὶ ἐν Ἱερουσαλὴμ,
\vs{7}καὶ εἶπεν, οἴμοι, ἱνατί τοῦτο ἐγεννήθην ἰδεῖν τὸ σύντριμμα τοῦ λαοῦ μου, καὶ τὸ σύντριμμα τῆς πόλεως τῆς ἁγίας, καὶ καθίσαι ἐκεῖ ἐν τῷ δοθῆναι αὐτὴν ἐν χειρὶ ἐχθρῶν, καὶ τὸ ἁγίασμα ἐν χειρὶ ἀλλοτρίων;

\vs{8}Ἐγένετο ὁ ναὸς αὐτῆς ὡς ἀνὴρ ἄδοξος,
\vs{9}τὰ σκεύη τῆς δόξης αὐτῆς αἰχμάλωτα ἀπήχθη, ἀπεκτάνθη τὰ νήπια αὐτῆς ἐν ταῖς πλατείαις, οἱ νεανίσκοι αὐτῆς ἐν ῥομφαίᾳ ἐχθροῦ.
\vs{10}Ποῖον ἔθνος οὐκ ἐκληρονόμησε βασιλείαν αὐτῆς, καὶ οὐκ ἐκράτησε τῶν σκύλων αὐτῆς;
\vs{11}Πᾶς ὁ κόσμος αὐτῆς ἀφῃρέθη, ἀντὶ ἐλευθήρας ἐγένετο εἰς δούλην.
\vs{12}Καὶ ἰδοὺ τὰ ἅγια ἡμῶν καὶ ἡ καλλονὴ ἡμῶν καὶ ἡ δόξα ἡμῶν ἠρημώθη, καὶ ἐβεβήλωσαν αὐτὰ τὰ ἔθνη.
\vs{13}Ἱνατί ἡμῖν ἔτι ζῇν;

\vs{14}Καὶ διέῤῥηξε Ματταθίας καὶ υἱοὶ αὐτοῦ τὰ ἱμάτια αὐτῶν, καὶ περιεβάλοντο σάκκους, καὶ ἐπένθησαν σφόδρα.

\vs{15}Καὶ ἦλθον οἱ παρὰ τοῦ βασιλέως οἱ καταναγκάζοντες τὴν ἀποστασίαν εἰς Μωδεῒν τὴν πόλιν, ἵνα θυσιάσωσι.
\vs{16}Καὶ πολλοὶ ἀπὸ Ἰσραὴλ πρὸς αὐτοὺς προσῆλθον· καὶ Ματταθίας καὶ οἱ υἱοὶ αὐτοῦ συνήχθησαν.

\vs{17}Καὶ ἀπεκρίθησαν οἱ παρὰ τοῦ βασιλέως, καὶ εἰπον τῷ Ματταθίᾳ, λέγοντες, ἄρχων καὶ ἔνδοξος καὶ μέγας εἶ ἐν τῇ πόλει ταύτῃ, καὶ ἐστηριγμένος ἐν υἱοῖς καὶ ἀδελφοῖς.
\vs{18}Νῦν οὖν πρόσελθε πρῶτος, καὶ ποίησον τὸ πρόσταγμα τοῦ βασιλέως, ὡς ἐποίησαν πάντα τὰ ἔθνη, καὶ οἱ ἄνδρες Ἰούδα, καὶ οἱ καταλειφθέντες ἐν Ἱερουσαλήμ· καὶ ἔσῃ σὺ καὶ ὁ οἶκός σου τῶν φίλων τοῦ βασιλέως, καὶ σὺ καὶ οἱ υἱοί σου δοξασθήσεσθε ἀργυρίῳ, καὶ χρυσίῳ, καὶ ἀποστολαῖς πολλαῖς.

\vs{19}Καὶ ἀπεκρίθη Ματταθίας, καὶ εἶπε φωνῇ μεγάλῃ, εἰ πάντα τὰ ἔθνη τὰ ἐν οἴκῳ τῆς βασιλείας τοῦ βασιλέως ἀκούουσιν αὐτοῦ, ἀποστῆναι ἕκαστος ἀπὸ λατρείας πατέρων αὐτοῦ, καὶ ᾑρετίσαντο ἐν ταῖς ἐντολαῖς αὐτοῦ,
\vs{20}ἀλλʼ ἐγὼ καὶ οἱ υἱοί μου καὶ οἱ ἀδελφοί μου πορευσόμεθα ἐν διαθήκῃ πατέρων ἡμῶν.
\vs{21}Ἵλεως ἡμῖν καταλιπεῖν νόμον καὶ δικαιώματὰ.
\vs{22}Τῶν λόγων τοῦ βασιλέως οὐκ ἀκουσόμεθα, τοῦ παρελθεῖν τὴν λατρείαν ἡμῶν, δεξιὰν ἢ ἀριστεράν.

\vs{23}Καὶ ὡς ἐπαύσατο λαλῶν τοὺς λόγους τούτους, προσῆλθεν ἀνὴρ Ἰουδαῖος ἐν ὀφθαλμοῖς πάντων, θυσιᾶσαι ἐπὶ τοῦ βωμοῦ τοῦ ἐν Μωδεῒν κατὰ τὸ πρόσταγμα τοῦ βασιλέως.
\vs{24}Καὶ εἶδε Ματταθίας καὶ ἐζήλωσε, καὶ ἐτρόμησαν οἱ νεφροὶ αὐτοῦ, καὶ ἀνήνεγκε θυμὸν κατὰ τὸ κρίμα, καὶ δραμῶν ἔσφαξεν αὐτὸν ἐπὶ τὸν βωμόν.

\vs{25}Καὶ τὸν ἄνδρα τοῦ βασιλέως τὸν ἀναγκάζοντα θύειν, ἀπέκτεινεν ἐν τῷ καιρῷ ἐκείνῳ, καὶ τὸν βωμὸν καθεῖλε.
\vs{26}Καὶ ἐζήλωσε τῷ νόμῳ καθὼς ἐποίησε Φινεὲς τῷ Ζαμβρὶ υἱῷ Σαλώμ.

\vs{27}Καὶ ἀνέκραξε Ματταθίας ἐν τῇ πόλει φωνῇ μεγάλῃ, λέγων, πᾶς ὁ ζηλῶν τῷ νόμῳ καὶ ἱστῶν διαθήκην, ἐξελθέτω ὀπίσω μου.
\vs{28}Καὶ ἔφυγον αὐτὸς καὶ οἱ υἱοὶ αὐτοῦ εἰς τὰ ὄρη, καὶ ἐγκατέλιπον ὅσα εἶχον ἐν τῇ πόλει.

\vs{29}Τότε κατέβησαν πολλοὶ ζητοῦντες δικαιοσύνην καὶ κρίμα, εἰς τὴν ἔρημον, καθίσαι ἐκεῖ,
\vs{30}αὐτοὶ καὶ οἱ υἱοὶ αὐτῶν καὶ αἱ γυναῖκες αὐτῶν καὶ τὰ κτήνη αὐτῶν, ὅτι ἐπληθύνθη ἐπʼ αὐτοὺς τὰ κακά.

\vs{31}Καὶ ἀνηγγέλη τοῖς ἀνδράσι τοῦ βασιλέως καὶ ταῖς δυνάμεσιν αἳ ἦσαν ἐν Ἱερουσαλὴμ πόλει Δαυὶδ, ὅτι κατέβησαν ἄνδρες, οἵτινες διεσκέδασαν τὴν ἐντολὴν τοῦ βασιλέως, εἰς τοὺς κρύφους ἐν τῇ ἐρήμῳ.
\vs{32}Καὶ ἔδραμον ὀπίσω αὐτῶν πολλοί· καὶ καταλαβόντες αὐτοὺς παρενέβαλον ἐπʼ αὐτοὺς, καὶ συνεστήσαντο πρὸς αὐτοὺς πόλεμον ἐν τῇ ἡμέρᾳ τῶν σαββάτων,
\vs{33}καὶ εἶπον πρὸς αὐτοὺς, ἕως τοῦ νῦν ἱκανόν· ἐξέλθετε καὶ ποιήσατε κατὰ τὸν λόγον τοῦ βασιλέως, καὶ ζήσεσθε.

\vs{34}Καὶ εἶπον, οὐκ ἐξελευσόμεθα, οὐδὲ ποιήσομεν τὸν λόγον τοῦ βασιλέως, τοῦ βεβηλῶσαι τὴν ἡμέραν τῶν σαββάτων.
\vs{35}καὶ ἐτάχυναν ἐπʼ αὐτοὺς πόλεμον.
\vs{36}Καὶ οὐκ ἀπεκρίθησαν αὐτοῖς, οὐδὲ λίθον ἐνετίναξαν αὐτοῖς, οὐδὲ ἐνέφραξαν τοὺς κρύφους,
\vs{37}λέγοντες, ἀποθάνωμεν πάντες ἐν τῇ ἁπλότητι ἡμῶν· μαρτυρεῖ ἐφʼ ἡμᾶς ὁ οὐρανὸς καὶ ἡ γῆ, ὅτι ἀκρίτως ἀπόλλυτε ἡμᾶς.
\vs{38}Καὶ ἀνέστησαν ἐπʼ αὐτοὺς ἐν τῷ πολέμῳ τοῖς σάββασι, καὶ ἀπέθανον αὐτοὶ καὶ αἱ γυναῖκες αὐτῶν, καὶ τὰ τέκνα αὐτῶν, καὶ τὰ κτήνη αὐτῶν, ἕως χιλίων ψυχῶν ἀνθρώπων.

\vs{39}Καὶ ἔγνω Ματταθίας καὶ οἱ φίλοι αὐτοῦ, καὶ ἐπένθησαν ἐπʼ αὐτοὺς ἕως σφόδρα.
\vs{40}Καὶ εἶπεν ἀνὴρ τῷ πλησίον αὐτοῦ, ἐὰν πάντες ποιήσωμεν ὡς οἱ ἀδελφοὶ ἡμῶν ἐποίησαν, καὶ μὴ πολεμήσωμεν πρὸς τὰ ἔθνη ὑπὲρ τῶν ψυχῶν ἡμῶν καὶ τῶν δικαιωμάτων ἡμῶν, νῦν τάχιον ἡμᾶς ἐξολοθρεύσουσιν ἀπὸ τῆς γῆς.

\vs{41}Καὶ ἐβουλεύσαντο τῇ ἡμέρᾳ ἐκείνῃ, λέγοντες, πᾶς ἄνθρωπος ὃς ἐὰν ἔλθῃ πρὸς ἡμᾶς εἰς πόλεμον τῇ ἡμέρᾳ τῶν σαββάτων, πολεμήσωμεν κατέναντι αὐτοῦ, καὶ οὐ μὴ ἀποθάνωμεν πάντες καθὼς ἀπέθανον οἱ ἀδελφοὶ ἡμῶν ἐν τοῖς κρύφοις.

\vs{42}Τότε συνήχθησαν πρὸς αὐτοὺς συναγωγὴ Ἰουδαίων, ἰσχυροὶ δυνάμει ἀπὸ Ἰσραὴλ, πᾶς ὁ ἑκουσιαζόμενος τῷ νόμῳ.
\vs{43}Καὶ πάντες οἱ φυγαδεύοντες ἀπὸ τῶν κακῶν προσετέθησαν αὐτοῖς, καὶ ἐγένοντο αὐτοῖς εἰς στήριγμα.
\vs{44}Καὶ συνεστήσαντο δύναμιν, καὶ ἐπάταξαν ἁμαρτωλοὺς ἐν ὀργῇ αὐτῶν, καὶ ἄνδρας ἀνόμους ἐν θυμῷ αὐτῶν· καὶ οἱ λοιποὶ ἔφυγον εἰς τὰ ἔθνη σωθῆναι.

\vs{45}Καὶ ἐκύκλωσε Ματταθίας καὶ οἱ φίλοι αὐτοῦ, καὶ καθεῖλον τοὺς βωμοὺς.
\vs{46}Καὶ περιέτεμον τὰ παιδάρια τὰ ἀπερίτμητα ὅσα εὗρον ἐν ὁρίοις Ἰσραὴλ ἐν ἰσχύϊ.
\vs{47}Καὶ ἐδίωξαν τοὺς υἱοὺς τῆς ὑπεπηφανίας, καὶ κατευωδώθη τὸ ἔργον ἐν χειρὶ αὐτῶν.
\vs{48}Καὶ ἀντελάβοντο τοῦ νόμου ἐκ χειρὸς τῶν ἐθνῶν καὶ ἐκ χειρὸς τῶν βασιλέων· καὶ οὐκ ἔδωκαν κέρας τῷ ἁμαρτωλῷ.

\vs{49}Καὶ ἤγγισαν αἱ ἡμέραι τοῦ Ματταθίου ἀποθανεῖν, καὶ εἶπε τοῖς υἱοῖς αὐτοῦ, νῦν ἐστηρίχθη ὑπερηφανία καὶ ἐλεγμὸς καὶ καιρὸς καταστροφῆς καὶ ὀργὴ θυμοῦ.
\vs{50}Καὶ νῦν, τέκνα, ζηλώσατε τῷ νόμῳ, καὶ δότε τὰς ψυχὰς ὑμῶν ὑπὲρ διαθήκης πατέρων ἡμῶν.
\vs{51}Μνήσθητε τῶν πατέρων ἡμῶν τὰ ἔργα ἃ ἐποίησαν ἐν ταῖς γενεαῖς αὐτῶν, καὶ δέξασθε δόξαν μεγάλην καὶ ὄνομα αἰώνιον.
\vs{52}Ἁβραὰμ οὐχὶ ἐν πειρασμῷ εὑρέθη πιστὸς, καὶ ἐλογίσθη αὐτῷ εἰς δικαιοσύνην;
\vs{53}Ἰωσὴφ ἐν καιρῷ στενοχωρίας αὐτοῦ ἐφύλαξεν ἐντολὴν, καὶ ἐγένετο κύριος Αἰγύπτου.
\vs{54}Φινεὲς ὁ πατὴρ ἡμῶν ἐν τῷ ζηλῶσαι ζῆλον, ἔλαβε διαθήκην ἱερωσύνης αἰωνίας.

\vs{55}Ἰησοῦς ἐν τῷ πληρῶσαι λόγον, ἐγένετο κριτὴς ἐν Ἰσραήλ.
\vs{56}Χαλὲβ ἐν τῷ ἐπιμαρτύρασθαι ἐν τῇ ἐκκλησία, ἔλαβε γῆς κληρονομίαν.
\vs{57}Δαυὶδ ἐν τῷ ἐλέῳ αὐτοῦ, ἐκληρονόμησε θρόνον βασιλείας εἰς αἰῶνα αἰῶνος.
\vs{58}Ἡλίας ἐν τῷ ζηλῶσαι ζῆλον νόμου, ἀνελήφθη ἕως εἰς τὸν οὐρανόν.
\vs{59}Ἀνανίας, Ἀζαρίας, Μισαήλ, πιστεύσαντες ἐσώθησαν ἐκ φλογός.
\vs{60}Δανιὴλ ἐν τῇ ἁπλότητι αὐτοῦ ἐῤῥύσθη ἐκ στόματος λεόντων.
\vs{61}Καὶ οὕτως ἐννοήθητε κατὰ γενεὰν καὶ γενεὰν, ὅτι πάντες οἱ ἐλπίζοντες ἐπʼ αὐτὸν οὐκ ἀσθενήσουσι.
\vs{62}Καὶ ἀπὸ λόγων ἀνδρὸς ἁμαρτωλοῦ μὴ φοβηθῆτε, ὅτι ἡ δόξα αὐτοῦ εἰς κοπρίαν καὶ εἰς σκώληκας.
\vs{63}Σήμερον ἐπαρθήσεται, καὶ αὔριον οὐ μὴ εὑρεθῇ, ὅτι ἔστρεψεν εἰς τὸν χοῦν αὐτοῦ, καὶ ὁ διαλογισμὸς αὐτοῦ ἀπώλετο.

\vs{64}Καὶ ὑμεῖς, τέκνα, ἰσχύσατε καὶ ἀνδρίζεσθε ἐν τῷ νόμῳ, ὅτι ἐν αὐτῷ δοξασθήσεσθε.
\vs{65}Καὶ ἰδοὺ Συμεὼν ὁ ἀδελφὸς ὑμῶν, οἶδα ὅτι ἀνὴρ βουλῆς ἐστιν, αὐτοῦ ἀκούετε πάσας τὰς ἡμέρας, αὐτὸς ὑμῖν ἔσται εἰς πατέρα.
\vs{66}Καὶ Ἰούδας Μακκαβαῖος ἰσχυρὸς δυνάμει ἐκ νεότητος αὐτοῦ, οὗτος ὑμῖν ἔσται ἄρχων στρατιᾶς, καὶ πολεμήσει πόλεμον λαῶν.

\vs{67}Καὶ ὑμεῖς προσάξατε πρὸς ὑμᾶς πάντας τοὺς ποιητὰς τοῦ νόμου, καὶ ἐκδικήσατε ἐκδίκησιν τοῦ λαοῦ ὑμῶν.
\vs{68}Ἀνταπόδοτε ἀνταπόδομα τοῖς ἔθνεσι, καὶ προσέχετε εἰς τὰ προστάγματα τοῦ νόμου.
\vs{69}Καὶ εὐλόγησεν αὐτούς· καὶ προσετέθη πρὸς τοὺς πατέρας αὐτοῦ.
\vs{70}Καὶ ἀπέθανεν ἐν τῷ ἕκτῳ καὶ τεσσαρακοστῷ καὶ ἑκατοστῷ ἔτει· καὶ ἔθαψαν αὐτὸν οἱ υἱοὶ αὐτοῦ ἐν τάφοις πατέρων αὐτῶν ἐν Μωδεῒν, καὶ ἐκόψαντο αὐτὸν πᾶς Ἰσραὴλ κοπετὸν μέγαν.

\ch{3}
Καὶ ἀνέστη Ἰούδας ὁ καλούμενος Μακκαβαῖος υἱὸς αὐτοῦ ἀντʼ αὐτοῦ.
\vs{2}Καὶ ἐβοήθουν αὐτῷ πάντες οἱ ἀδελφοὶ αὐτοῦ, καὶ πάντες ὅσοι ἐκολλήθησαν τῷ πατρὶ αὐτοῦ, καὶ ἐπολέμουν τὸν πόλεμον Ἰσραὴλ μετʼ εὐφροσύνης.
\vs{3}Καὶ ἐπλάτυνε δόξαν τῷ λαῷ αὐτοῦ, καὶ ἐνεδύσατο θώρακα ὡς γίγας, καὶ συνεζώσατο τὰ σκεύη αὐτοῦ τὰ πολεμικά· καὶ συνεστήσατο πολέμους σκεπάζων παρεμβολὴν ἐν ῥομφαίᾳ.

\vs{4}Καὶ ὡμοιώθη λέοντι ἐν τοῖς ἔργοις αὐτοῦ, καὶ ὡς σκύμνος ἐρευγόμενος εἰς θήραν.
\vs{5}Καὶ ἐδίωξεν ἀνόμους ἐξερευνῶν, καὶ τοὺς ταράσσοντας τὸν λαὸν αὐτοῦ ἐφλόγισε·
\vs{6}Καὶ συνεστάλησαν οἱ ἄνομοι ἀπὸ τοῦ φόβου αὐτοῦ, καὶ πάντες οἱ ἐργάται τῆς ἀνομίας συνεταράχθησαν, καὶ εὐοδώθη σωτηρία ἐν χειρὶ αὐτοῦ.

\vs{7}Καὶ ἐπίκρανε βασιλεῖς πολλούς, καὶ εὔφρανε τὸν Ἰακὼβ ἐν τοῖς ἔργοις αὐτοῦ, καὶ ἕως τοῦ αἰῶνος τὸ μνημόσυνον αὐτοῦ εἰς εὐλογίαν.
\vs{8}Καὶ διῆλθεν ἐν πόλεσιν Ἰούδα, καὶ ἐξωλόθρευσεν ἀσεβεῖς ἐξ αὐτῆς, καὶ ἀπέστρεψεν ὀργὴν ἀπὸ Ἰσραήλ.
\vs{9}Καὶ ὠνομάσθη ἕως ἐσχάτου τῆς γῆς, καὶ συνήγαγεν ἀπολλυμένους.

\vs{10}Καὶ συνήγαγεν Ἀπολλώνιος ἔθνη, καὶ ἀπὸ Σαμαρείας δύναμιν μεγάλην, τοῦ πολεμῆσαι πρὸς Ἰσραήλ.
\vs{11}Καὶ ἔγνω Ἰούδας, καὶ ἐξῆλθεν εἰς συνάντησιν αὐτῷ, καὶ ἐπάταξεν αὐτὸν, καὶ ἀπέκτεινεν αὐτόν· καὶ ἔπεσον τραυματίαι πολλοὶ, καὶ οἱ ἐπίλοιποι ἔφυγον.
\vs{12}Καὶ ἔλαβε τὰ σκῦλα αὐτῶν, καὶ τὴν μάχαιραν Ἀπολλωνίου ἔλαβεν Ἰούδας, καὶ ἦν πολεμῶν ἐν αὐτῇ πάσας τὰς ἡμέρας.

\vs{13}Καὶ ἤκουσε Σήρων, ὁ ἄρχων τῆς δυνάμεως Συρίας, ὅτι ἤθροισεν Ἰούδας ἄθροισμα, καὶ ἐκκλησίαν πιστῶν μετʼ αὐτοῦ ἐκπορευομένων εἰς πόλεμον·
\vs{14}Καὶ εἶπε, ποιήσω ἐμαυτῷ ὄνομα καὶ δοξασθήσομαι ἐν τῇ βασιλείᾳ, καὶ πολεμήσω τὸν Ἰούδαν καὶ τοὺς σὺν αὐτῷ, τοὺς ἐξουδενοῦντας τὸν λόγον τοῦ βασιλέως.
\vs{15}Καὶ προσέθετο τοῦ ἀναβῆναι· καὶ ἀνέβη μετʼ αὐτοῦ παρεμβολὴ ἀσεβῶν ἰσχυρὰ βοηθῆσαι αὐτῷ, καὶ ποιῆσαι τὴν ἐκδίκησιν ἐν υἱοῖς Ἰσραήλ.

\vs{16}Καὶ ἤγγισαν ἕως ἀναβάσεως Βαιθωρών· καὶ ἐξῆλθεν Ἰούδας εἰς συνάντησιν αὐτῷ ὀλιγοστός.
\vs{17}Ὡς δὲ ἴδον τὴν παρεμβολὴν ἐρχομένην εἰς συνάντησιν αὐτοῖς, εἶπον τῷ Ἰούδα, πῶς δυνησόμεθα ὀλιγοστοὶ ὄντες πολεμῆσαι πρὸς πλῆθος τοσοῦτον ἰσχυρόν; καὶ ἡμεῖς ἐκλελύμεθα ἀσιτοῦντες σήμερον.
\vs{18}Καὶ εἶπεν Ἰούδας, εὔκοπόν ἐστι συγκλεισθῆναι πολλοὺς ἐν χερσὶν ὀλίγων· καὶ οὐκ ἔστι διαφορὰ ἐναντίον τοῦ Θεοῦ τοῦ οὐρανοῦ σώζειν ἐν πολλοῖς ἢ ἐν ὀλίγοις.
\vs{19}Ὅτι οὐκ ἐν πλήθει δυνάμεως νίκη πολέμου ἐστὶν, ἀλλʼ ἢ ἐκ τοῦ οὐρανοῦ ἡ ἰσχύς.
\vs{20}Αὐτοὶ ἔρχονται πρὸς ἡμᾶς ἐν πλήθει ὕβρεως καὶ ἀνομίας, τοῦ ἐξᾶραι ἡμᾶς καὶ τὰς γυναῖκας ἡμῶν, καὶ τὰ τέκνα ἡμῶν, τοῦ σκυλεῦσαι ἡμᾶς.
\vs{21}Ἡμεῖς δὲ πολεμοῦμεν περὶ τῶν ψυχῶν ἡμῶν καὶ τῶν νομίμων ἡμῶν.
\vs{22}Καὶ αὐτὸς συντρίψει αὐτοὺς πρὸ προσώπου ἡμῶν· ὑμεῖς δὲ μὴ φοβηθῆτε ἀπʼ αὐτῶν.

\vs{23}Ὡς δὲ ἐπαύσατο λαλῶν, ἐνήλατο εἰς αὐτοὺς ἄφνω, καὶ συνετρίβη Σήρων καὶ ἡ παρεμβολὴ αὐτοῦ ἐνώπιον αὐτοῦ.
\vs{24}Καὶ ἐδίωκον αὐτὸν ἐν τῇ καταβάσει Βαιθωρῶν ἕως τοῦ πεδίου· καὶ ἔπεσον ἀπʼ αὐτῶν εἰς ἄνδρας ὀκτακοσίους· οἱ δὲ λοιποὶ ἔφυγον εἰς γῆν Φυλιστιείμ.
\vs{25}Καὶ ἤρξατο ὁ φόβος Ἰούδα καὶ τῶν ἀδελφῶν αὐτοῦ καὶ ἡ πτόησις ἐπιπίπτειν ἐπὶ τὰ ἔθνη τὰ κύκλῳ αὐτῶν.
\vs{26}Καὶ ἤγγισεν ἕως τοῦ βασιλέως τὸ ὄνομα αὐτοῦ, καὶ ὑπὲρ τῶν παρατάξεων Ἰούδα ἐξηγεῖτο πᾶν ἔθνος.

\vs{27}Ὡς δὲ ἤκουσεν Ἀντίοχος ὁ βασιλεὺς τοὺς λόγους τούτους, ὠργίσθη θυμῷ· καὶ ἀπέστειλε καὶ συνήγαγε τὰς δυνάμεις πάσας τῆς βασιλείας αὐτοῦ, παρεμβολὴν ἰσχυρὰν σφόδρα.
\vs{28}Καὶ ἤνοιξε τὸ γαζοφυλάκιον αὐτοῦ, καὶ ἔδωκεν ὀψώνια ταῖς δυνάμεσιν αὐτοῦ εἰς ἐνιαυτόν· καὶ ἐνετείλατο εἶναι αὐτοὺς ἑτοίμους εἰς πᾶσαν χρείαν.

\vs{29}Καὶ εἶδεν ὅτι ἐξέλιπε τὸ ἀργύριον ἀπὸ τῶν θησαυρῶν, καὶ οἱ φορολόγοι τῆς χώρας ὀλίγοι, χάριν τῆς διχοστασίας καὶ πληγῆς ἧς κατεσκεύασεν ἐν τῇ γῇ, τοῦ ἆραι τὰ νόμιμα ἃ ἦσαν ἀφʼ ἡμερῶν τῶν πρώτων.
\vs{30}Καὶ εὐλαβήθη μὴ οὐκ ἔχῃ ὡς ἅπαξ καὶ δὶς εἰς τὰς δαπάνας καὶ τὰ δόματα ἃ ἐδίδου ἔμπροσθεν δαψιλεῖ χειρὶ, καὶ ἐπερίσσευσεν ὑπέρ τοὺς βασιλεῖς τοὺς ἔμπροσθεν.

\vs{31}Καὶ ἠπορεῖτο τῇ ψυχῇ αὐτοῦ σφόδρα, καὶ ἐβουλεύσατο τοῦ πορευθῆναι εἰς τὴν Περσίδα, καὶ λαβεῖν τοὺς φόρους τῶν χωρῶν, καὶ συναγαγεῖν ἀργύριον πολύ.
\vs{32}Καὶ κατέλιπε Λυσίαν ἄνθρωπον ἔνδοξον καὶ ἀπὸ γένους τῆς βασιλείας, ἐπὶ τῶν πραγμάτων τοῦ βασιλέως ἀπὸ τοῦ ποταμοῦ Εὐφράτου ἕως τῶν ὁρίων Αἰγύπτου,
\vs{33}καὶ τρέφειν Ἀντίοχον τὸν υἱὸν αὐτοῦ ἕως τοῦ ἐπιστρέψαι αὐτόν.
\vs{34}Καὶ παρέδωκεν αὐτῷ τὰς ἡμίσεις τῶν δυνάμεων καὶ τοὺς ἐλέφαντας· καὶ ἐνετείλατο αὐτῷ περὶ πάντων ὧν ἐβούλετο, καὶ περὶ τῶν κατοικούντων τὴν Ἰουδαίαν καὶ Ἱερουσαλὴμ,
\vs{35}ἀποστεῖλαι ἐπʼ αὐτοὺς δύναμιν, τοῦ ἐκτρίψαι καὶ ἐξᾶραι τὴν ἰσχὺν Ἰσραὴλ, καὶ τὸ κατάλειμμα Ἱερουσαλὴμ, καὶ ἆραι τὸ μνημόσυνον αὐτῶν ἀπὸ τοῦ τόπου,
\vs{36}καὶ κατοικῆσαι υἱοὺς ἀλλογενεῖς ἐν πᾶσι τοῖς ὁρίοις αὐτῶν, καὶ κατακληροδοτῆσαι τὴν γῆν αὐτῶν.
\vs{37}Καὶ ὁ βασιλεὺς παρέλαβε τὰς ἡμίσεις τῶν δυνάμεων τὰς καταλειφθείσας, καὶ ἀπῇρεν ἀπὸ Ἀντιοχείας ἀπὸ πόλεως βασιλείας αὐτοῦ, ἔτους ἑβδόμου καὶ τεσσαρακοστοῦ καὶ ἑκατοστοῦ· καὶ διεπέρασε τὸν Εὐφράτην ποταμὸν, καὶ διεπορεύετο τὰς ἐπάνω χώρας.

\vs{38}Καὶ ἐπέλεξε Λυσίας Πτολεμαῖον τὸν Δορυμένους, καὶ Νικάνορα, καὶ Γοργίαν, ἄνδρας δυνατοὺς τῶν φίλων τοῦ βασιλέως.
\vs{39}Καὶ ἀπέστειλε μετʼ αὐτῶν τεσσαράκοντα χιλιάδας ἀνδρῶν καὶ ἑπτακισχιλίαν ἵππον, τοῦ ἐλθεῖν εἰς γῆν Ἰούδα, καὶ καταφθεῖραι αὐτὴν, κατὰ τὸν λόγον τοῦ βασιλέως.
\vs{40}Καὶ ἀπῇραν σὺν πάσῃ τῇ δυνάμει αὐτῶν, καὶ ἦλθον, καὶ παρενέβαλον πλησίον Ἐμμαοὺμ ἐν τῇ γῇ τῇ πεδινῇ.

\vs{41}Καὶ ἤκουσαν οἱ ἔμποροι τῆς χώρας τὸ ὄνομα αὐτῶν, καὶ ἔλαβον ἀργύριον καὶ χρυσίον πολὺ σφόδρα καὶ παῖδας, καὶ ἦλθον εἰς τὴν παρεμβολὴν τοῦ λαβεῖν τοὺς υἱοὺς Ἰσραὴλ εἰς παῖδας· καὶ προσετέθησαν πρὸς αὐτοὺς δύναμις Συρίας καὶ γῆς ἀλλοφύλων.

\vs{42}Καὶ εἶδεν Ἰούδας καὶ οἱ ἀδελφοὶ αὐτοῦ ὅτι ἐπληθύνθη τὰ κακὰ, καὶ αἱ δυνάμεις παρεμβάλλουσιν ἐν τοῖς ὁρίοις αὐτῶν· καὶ ἐπέγνωσαν τοὺς λόγους τοῦ βασιλέως οὓς ἐνετείλατο ποιῆσαι τῷ λαῷ εἰς ἀπώλειαν καὶ συντέλειαν·
\vs{43}καὶ εἶπεν ἕκαστος πρὸς τὸν πλησίον αὐτοῦ, ἀναστήσωμεν τὴν καθαίρεσιν τοῦ λαοῦ ἡμῶν, καὶ πολεμήσωμεν περὶ τοῦ λαοῦ ἡμῶν καὶ τῶν ἁγίων.

\vs{44}Καὶ συνηθροίσθη ἡ συναγωγὴ τοῦ εἶναι ἑτοίμους εἰς πόλεμον, καὶ τοῦ προσεύξασθαι, καὶ αἰτῆσαι ἔλεον καὶ οἰκτιρμούς.

\vs{45}Καὶ Ἱερουσαλὴμ ἦν ἀοίκητος ὡς ἔρημος, οὐκ ἦν ὁ εἰσπορευόμενος καὶ ἐκπορευόμενος ἐκ τῶν γεννημάτων αὐτῆς· καὶ τὸ ἁγίασμα καταπατούμενον, καὶ υἱοὶ ἀλλογενῶν ἐν τῇ ἄκρᾳ, κατάλυμα τοῖς ἔθνεσι· καὶ ἐξῄρθη τέρψις ἐξ Ἰακὼβ, καὶ ἐξέλιπεν αὐλὸς καὶ κινύρα.
\vs{46}Καὶ συνήχθησαν, καὶ ἤλθοσαν εἰς Μασσηφὰ κατέναντι Ἱερουσαλὴμ, ὅτι τόπος προσευχῆς εἰς Μασσηφὰ τὸ πρότερον τῷ Ἰσραήλ.

\vs{47}Καὶ ἐνήστευσαν τῇ ἡμέρᾳ ἐκείνῃ, καὶ περιεβάλοντο σάκκους καὶ σποδὸν ἐπὶ τὰς κεφαλὰς αὐτῶν, καὶ διέῤῥηξαν τὰ ἱμάτια αὐτῶν.
\vs{48}Καὶ ἐξεπέτασαν τὸ βιβλίον τοῦ νόμου, περὶ ὧν ἐξηρεύνων τὰ ἔθνη τὰ ὁμοιώματα τῶν εἰδώλων αὐτῶν.
\vs{49}Καὶ ἤνεγκαν τὰ ἱμάτια τῆς ἱερωσύνης, καὶ τὰ πρωτογεννήματα, καὶ τὰς δεκάτας· καὶ ἤγειραν τοὺς Ναζαραίους, οἳ ἐπλήρωσαν τὰς ἡμέρας.

\vs{50}Καὶ ἐβόησαν φωνῇ εἰς τὸν οὐρανὸν, λέγοντες, τί ποιήσωμεν τούτοις, καὶ ποῦ αὐτοὺς ἀπαγάγωμεν;
\vs{51}Καὶ τὰ ἅγιά σου καταπεπάτηται, καὶ βεβήλωται· καὶ οἱ ἱερεῖς σου ἐν πένθει καὶ ταπεινώσει.
\vs{52}Καὶ ἰδοὺ τὰ ἔθνη συνῆκται ἐφʼ ἡμᾶς τοῦ ἐξᾶραι ἡμᾶς· σὺ οἶδας ἃ λογίζονται ἐφʼ ἡμᾶς.
\vs{53}Πῶς δυνησόμεθα ὑποστῆναι κατὰ πρόσωπον αὐτῶν, ἐὰν μὴ σὺ βοηθήσῃς ἡμῖν;
\vs{54}Καὶ ἐσάλπισαν ταῖς σάλπιγξι, καὶ ἐβόησαν φωνῇ μεγάλῃ.

\vs{55}Καὶ μετὰ τοῦτο κατέστησεν Ἰούδας ἡγουμένους τοῦ λαοῦ, χιλιάρχους, καὶ ἑκατοντάρχους, καὶ πεντηκοντάρχους, καὶ δεκάρχους.
\vs{56}Καὶ εἶπον τοῖς οἰκοδομοῦσιν οἰκίας, καὶ μνηστευομένους γυναῖκας, καὶ φυτεύουσιν ἀμπελῶνας, καὶ δειλοῖς, ἀποστρέφειν ἕκαστον εἰς τὸν οἶκον αὐτοῦ, κατὰ τὸν νόμον.

\vs{57}Καὶ ἀπῇρεν ἡ παρεμβολὴ, καὶ παρενέβαλε κατὰ Νότον Ἐμμαούμ.
\vs{58}Καὶ εἶπεν Ἰούδας, περιζώσασθε, καὶ γίνεσθε εἰς υἱοὺς δυνατοὺς, καὶ γίνεσθε ἕτοιμοι εἰς τοπρωῒ τοῦ πολεμῆσαι ἐν τοῖς ἔθνεσι τούτοις, τοῖς ἐπισυνηγμένοις ἐφʼ ἡμᾶς ἐξᾶραι ἡμᾶς καὶ τὰ ἅγια ἡμῶν.
\vs{59}Ὅτι κρεῖσσον ἡμᾶς ἀποθανεῖν ἐν τῷ πολέμῳ, ἢ ἐπιδεῖν ἐπὶ τὰ κακὰ τοῦ ἔθνους ἡμῶν καὶ τῶν ἁγίων·
\vs{60}Ὡς δʼ ἂν ᾖ θέλημα ἐν οὐρανῷ, οὕτω ποιήσει.

\ch{4}
Καὶ παρέλαβε Γοργίας πεντακισχιλίους ἄνδρας καὶ χιλίαν ἵππον ἐκλεκτὴν, καὶ ἀπῇρεν ἡ παρεμβολὴ νυκτὸς,
\vs{2}ὥστε ἐπιβαλεῖν ἐπὶ τὴν παρεμβολὴν τῶν Ἰουδαίων, καὶ πατάξαι αὐτοὺς ἄφνω· καὶ οἱ υἱοὶ τῆς ἄκρας ἦσαν αὐτῷ ὁδηγοί.
\vs{3}Καὶ ἤκουσεν Ἰούδας, καὶ ἀπῇρεν αὐτὸς καὶ οἱ δυνατοὶ πατάξαι τὴν δύναμιν τοῦ βασιλέως τὴν ἐν Ἐμμαοὺμ,
\vs{4}ἕως ἔτι αἱ δυνάμεις ἐσκορπισμέναι ἦσαν ἀπὸ τῆς παρεμβολῆς.

\vs{5}Καὶ ἦλθε Γοργίας εἰς τὴν παρεμβολὴν Ἰούδα νυκτὸς, καὶ οὐδένα εὗρε· καὶ ἐζήτει αὐτοὺς ἐν τοῖς ὄρεσιν, ὅτι εἶπε, φεύγουσιν οὗτοι ἀφʼ ἡμῶν.

\vs{6}Καὶ ἅμα τῇ ἡμέρᾳ, ὤφθη Ἰούδας ἐν τῷ πεδίῳ ἐν τρισχιλίοις ἀνδράσι· πλὴν καλύμματα καὶ μαχαίρας οὐκ εἶχον καθὼς ἠβούλοντο.
\vs{7}Καὶ εἶδον παρεμβολὴν ἐθνῶν ἰσχυρὰν, τεθωρακισμένην, καὶ ἵππον κυκλοῦσαν αὐτὴν, καὶ οὗτοι διδακτοὶ πολέμου.

\vs{8}Καὶ εἶπεν Ἰούδας τοῖς ἀνδράσι τοῖς μετʼ αὐτοῦ, μὴ φοβεῖσθε τὸ πλῆθος αὐτῶν, καὶ τὸ ὅρμημα αὐτῶν μὴ δειλωθῆτε.
\vs{9}Μνήσθητε πῶς ἐσώθησαν οἱ πατέρες ἡμῶν ἐν θαλάσσῃ ἐρυθρᾷ, ὅτε ἐδίωξεν αὐτοὺς Φαραὼ ἐν δυνάμει.
\vs{10}Καὶ νῦν βοήσωμεν εἰς τὸν οὐρανὸν, εἴπως ἐλεήσει ἡμᾶς, καὶ μνησθήσεται διαθήκης πατέρων ἡμῶν, καὶ συντρίψει τὴν παρεμβολὴν ταύτην κατὰ πρόσωπον ἡμῶν σήμερον.
\vs{11}Καὶ γνώσεται πάντα τὰ ἔθνη, ὅτι ἐστὶν ὁ λυτρούμενος καὶ σώζων τὸν Ἰσραήλ.

\vs{12}Καὶ ᾖραν οἱ ἀλλόφυλοι τοὺς ὀφθαλμοὺς αὐτῶν. καὶ ἴδον αὐτοὺς ἐρχομένους ἐξεναντίας,
\vs{13}καὶ ἐξῆλθον ἐκ τῆς παρεμβολῆς εἰς πόλεμον· καὶ ἐσάλπισαν οἱ μετὰ Ἰούδα.
\vs{14}Καὶ συνῆψαν, καὶ συνετρίβησαν τὰ ἔθνη, καὶ ἔφυγον εἰς τὸ πεδίον.
\vs{15}Οἱ δὲ ἔσχατοι πάντες ἔπεσον ἐν ῥομφαίᾳ· καὶ ἐδίωξαν αὐτοὺς ἕως Γαζηρὼν καὶ ἕως τῶν πεδίων τῆς Ἰδουμαίας καὶ Ἀζώτου καὶ Ἰαμνίας, καὶ ἔπεσον ἐξ αὐτῶν εἰς ἄνδρας τρισχιλίους.

\vs{16}Καὶ ἐπέστρεψεν Ἰούδας καὶ ἡ δύναμις ἀπὸ τοῦ διώκειν ὄπισθεν αὐτῶν,
\vs{17}καὶ εἶπε πρὸς τὸν λαὸν, μὴ ἐπιθυμήσητε τῶν σκύλων, ὅτι πόλεμος ἐξεναντίας ἡμῶν,
\vs{18}καὶ Γοργίας καὶ ἡ δύναμις ἐν τῷ ὄρει ἐγγὺς ἡμῶν· ἀλλὰ στῆτε νῦν ἐναντίον τῶν ἐχθρῶν ἡμῶν, καὶ πολεμήσατε αὐτοὺς, καὶ μετὰ ταῦτα λήψετε τὰ σκύλα μετὰ παῤῥησίας.

\vs{19}Ἔτι λαλοῦντος Ἰούδα ταῦτα, ὤφθη μέρος τι ἐκκύπτον ἐκ τοῦ ὄρους.
\vs{20}Καὶ εἶδεν ὅτι τετρόπωνται, καὶ ἐμπυρίζουσι τὴν παρεμβολὴν, ὁ γὰρ καπνὸς θεωρούμενος ἐνεφάνιζε τὸ γεγονός.
\vs{21}Οἱ δὲ ταῦτα συνιδόντες ἐδειλώθησαν σφόδρα· συνιδόντες δὲ καὶ τὴν Ἰούδα παρεμβολὴν ἐν τῷ πεδίῳ ἑτοίμην εἰς παράταξιν,
\vs{22}ἔφυγον πάντες εἰς γῆν ἀλλοφύλων.
\vs{23}Καὶ ἀνέστρεψεν Ἰούδας ἐπὶ τὴν σκυλείαν τῆς παρεμβολῆς· καὶ ἔλαβον χρυσίον πολὺ καὶ ἀργύριον καὶ ὑάκινθον καὶ πορφύραν θαλασσίαν καὶ πλοῦτον μέγαν.
\vs{24}Καὶ ἐπιστραφέντες ὕμνουν καὶ εὐλόγουν εἰς οὐρανὸν τὸν Κύριον, ὅτι καλὸν, ὅτι εἰς τὸν αἰῶνα τὸ ἔλεος αὐτοῦ.
\vs{25}Καὶ ἐγένετο σωτηρία μεγάλη τῷ Ἰσραὴλ ἐν τῇ ἡμέρᾳ ἐκείνῃ.

\vs{26}Ὅσοι δὲ τῶν ἀλλοφύλων διεσώθησαν, παραγενηθέντες ἀπήγγειλαν τῷ Λυσίᾳ πάντα τὰ συμβεβηκότα.
\vs{27}Ὁ δὲ ἀκούσας συνεχύθη καὶ ἠθύμει, ὅτι οὐχ οἷα ἤθελε, τοιαῦτα γεγόνει τῷ Ἰσραὴλ, καὶ οὐχ οἷα ἐνετείλατο αὐτῷ ὁ βασιλεὺς, τοιαῦτα ἐξέβη.

\vs{28}Καὶ ἐν τῷ ἐχομένῳ ἐνιαυτῷ συνελόχησεν ὁ Λυσίας ἀνδρῶν ἐπιλέκτων ἑξήκοντα χιλιάδας καὶ πεντακισχιλίαν ἵππον, ὥστε ἐκπολεμῆσαι αὐτούς.
\vs{29}Καὶ ἦλθον εἰς τὴν Ἰδουμαίαν, καὶ παρενέβαλον ἐν Βαιθσούροις, καὶ συνήντησεν αὐτοῖς Ἰούδας ἐν δέκα χιλιάσιν ἀνδρῶν.

\vs{30}Καὶ εἶδε τὴν παρεμβολὴν ἰσχυρὰν, καὶ προσηύξατο, καὶ εἶπεν, εὐλογητὸς εἶ, ὁ σωτὴρ τοῦ Ἰσραὴλ, ὁ συντρίψας τὸ ὅρμημα τοῦ δυνατοῦ ἐν χειρὶ τοῦ δούλου σου Δαυὶδ, καὶ παρέδωκας τὴν παρεμβολὴν τῶν ἀλλοφύλων εἰς χεῖρας Ἰωνάθαν υἱοῦ Σαοὺλ, καὶ τοῦ αἴροντος τὰ σκεύη αὐτοῦ.
\vs{31}Σύγκλεισον τὴν παρεμβολὴν ταύτην ἐν χειρὶ λαοῦ σου Ἰσραὴλ, καὶ αἰσχυνθήτωσαν ἐπὶ τῇ δυνάμει καὶ τῇ ἵππῳ αὐτῶν.
\vs{32}Δὸς αὐτοῖς δειλίαν, καὶ τῆξον θράσος ἰσχύος αὐτῶν, καὶ σαλευθήτωσαν τῇ συντριβῇ αὐτῶν.
\vs{33}Κατάβαλε αὐτοὺς ῥομφαίᾳ ἀγαπώντων σε, καὶ αἰνεσάτωσάν σε πάντες οἱ εἰδότες τὸ ὄνομά σου ἐν ὕμνοις.

\vs{34}Καὶ συνέβαλον ἀλλήλοις, καὶ ἔπεσον ἐκ τῆς παρεμβολῆς Λυσίου εἰς πεντακισχιλίους ἄνδρας, καὶ ἔπεσον ἐξ ἐναντίας αὐτῶν.

\vs{35}Ἰδὼν δὲ Λυσίας τὴν γενομένην τροπὴν, τῆς αὐτοῦ συντάξεως, τῆς δὲ Ἰούδα τὸ γεγενημένον θάρσος, καὶ ὡς ἕτοιμοί εἰσιν ἢ ζῇν ἢ τεθνάναι γενναίως, ἀπῇρεν εἰς Ἀντιόχειαν, καὶ ἐξενολόγει· καὶ πλεονάσας τὸν γενηθέντα στρατὸν, ἐλογίζετο πάλιν παραγενέσθαι εἰς τὴν Ἰουδαίαν.

\vs{36}Εἶπε δὲ Ἰούδας καὶ οἱ ἀδελφοὶ αὐτοῦ, Ἰδοὺ συνετρίβησαν οἱ ἐχθροὶ ἡμῶν, ἀναβῶμεν καθαρίσαι τὰ ἅγια καὶ ἐγκαινίσαι.
\vs{37}Καὶ συνήχθη ἡ παρεμβολὴ πᾶσα, καὶ ἀνέβησαν εἰς ὄρος Σιών.
\vs{38}Καὶ ἴδον τὸ ἁγίασμα ἠρημωμένον, καὶ τὸ θυσιαστήριον βεβηλωμένον, καὶ τὰς πύλας κατακεκαυμένας, καὶ ἐν ταῖς αὐλαῖς φυτὰ πεφυκότα ὡς ἐν δρυμῷ ἢ ὡς ἐν ἑνὶ τῶν ὀρέων, καὶ τὰ παστοφόρια καθῃρημένα.
\vs{39}Καὶ διέῤῥηξαν τὰ ἱμάτια αὐτῶν, καὶ ἐκόψαντο κοπετὸν μέγαν, καὶ ἐπέθεντο σποδὸν ἐπὶ τὴν κεφαλὴν αὐτῶν.
\vs{40}Καὶ ἔπεσον ἐπὶ πρόσωπον ἐπὶ τὴν γῆν, καὶ ἐσάλπισαν ταῖς σάλπιγξι τῶν σημασιῶν, καὶ ἐβόησαν εἰς τὸν οὐρανόν.

\vs{41}Τότε ἐπέταξεν Ἰούδας ἄνδρας πολεμεῖν τοὺς ἐν τῇ ἄκρᾳ, ἕως ἂν καθαρίσῃ τὰ ἅγια.
\vs{42}Καὶ ἐπέλεξεν ἱερεῖς ἀμώμους, θελητὰς νόμου.
\vs{43}Καὶ ἐκαθάρισαν τὰ ἅγια, καὶ ᾖραν τοὺς λίθους τοῦ μιασμοῦ εἰς τόπον ἀκάθαρτον.
\vs{44}Καὶ ἐβουλεύσαντο περὶ τοῦ θυσιαστηρίου τῆς ὁλοκαυτώσεως τοῦ βεβηλωμένου, τί αὐτῷ ποιήσωσι.
\vs{45}Καὶ ἐπέπεσεν αὐτοῖς βουλὴ ἀγαθὴ, καθελεῖν αὐτὸ, μήποτε γένηται αὐτοῖς εἰς ὄνειδος, ὅτι ἐμίαναν τὰ ἔθνη αὐτό· καὶ καθεῖλον τὸ θυσιαστήριον,
\vs{46}καὶ ἀπέθεντο τοὺς λίθους ἐν τῷ ὄρει τοῦ οἴκου, ἐν τόπῳ ἐπιτηδείῳ, μέχρι τοῦ παραγενηθῆναι προφήτην τοῦ ἀποκριθῆναι περὶ αὐτῶν.

\vs{47}Καὶ ἔλαβον λίθους ὁλοκλήρους κατὰ τὸν νόμον, καὶ ᾠκοδόμησαν τὸ θυσιαστήριον καινὸν κατὰ τὸ πρότερον.
\vs{48}Καὶ ᾠκοδόμησαν τὰ ἅγια καὶ τὰ ἐντὸς τοῦ οἴκου, καὶ τὰς αὐλὰς ἡγίασαν.
\vs{49}Καὶ ἐποίησαν σκεύη ἅγια καινὰ, καὶ εἰσήνεγκαν τὴν λυχνίαν καὶ τὸ θυσιαστήριον τῶν θυμιαμάτων καὶ τὴν τράπεζαν εἰς τὸν ναόν.

\vs{50}Καὶ ἐθυμίασαν ἐπὶ τὸ θυσιαστήριον, καὶ ἐξῆψαν τοὺς λύχνους τοὺς ἐπὶ τῆς λυχνίας, καὶ ἐφαίνοσαν ἐν τῷ ναῷ.
\vs{51}Καὶ ἐπέθηκαν ἐπὶ τὴν τράπεζαν ἄρτους, καὶ ἐξεπέτασαν τὰ καταπετάσματα· καὶ ἐτέλεσαν πάντα τὰ ἔργα ἃ ἐποίησαν.

\vs{52}Καὶ ὤρθρισαν τοπρωῒ τῇ πέμπτῃ καὶ εἰκάδι τοῦ μηνὸς τοῦ ἐννάτου· οὗτος ὁ μὴν Χασελεῦ τοῦ ὀγδόου καὶ τεσσαρακοστοῦ καὶ ἑκατοστοῦ ἔτους.
\vs{53}Καὶ ἀνήνεγκαν θυσίαν κατὰ τὸν νόμον ἐπὶ τὸ θυσιαστήριον τῶν ὁλοκαυτωμάτων τὸ καινὸν ὃ ἐποίησαν.
\vs{54}Κατὰ τὸν καιρὸν καὶ κατὰ τὴν ἡμέραν ἐν ᾗ ἐβεβήλωσαν αὐτὸ τὰ ἔθνη, ἐν ἐκείνῃ ἐνεκαινίσθη ἐν ᾠδαῖς καὶ κιθάραις καὶ κινύραις, καὶ ἐν κυμβάλοις.
\vs{55}Καὶ ἔπεσον πᾶς ὁ λαὸς ἐπὶ πρόσωπον, καὶ προσεκύνησαν, καὶ εὐλόγησαν εἰς οὐρανὸν τὸν εὐοδώσαντα αὐτοῖς.

\vs{56}Καὶ ἐποίησαν τὸν ἐγκαινισμὸν τοῦ θυσιαστηρίου ἡμέρας ὀκτὼ, καὶ προσήνεγκαν ὁλοκαυτώματα μετʼ εὐφροσύνης, καὶ ἔθυσαν θυσίαν σωτηρίου καὶ αἰνέσεως.
\vs{57}Καὶ κατεκόσμησαν τὸ κατὰ πρόσωπον τοῦ ναοῦ στεφάνοις χρυσοῖς καὶ ἀσπιδίσκαις, καὶ ἐνεκαίνισν τὰς πύλας καὶ τὰ παστοφόρια, καὶ ἐθύρωσαν αὐτά.
\vs{58}Καὶ ἐγενήθη εὐφροσύνη μεγάλη ἐν τῷ λαῷ σφόδρα, καὶ ἀπεστράφη ὄνειδος ἐθνῶν.

\vs{59}Καὶ ἔστησεν Ἰούδας καὶ οἱ ἀδελφοὶ αὐτοῦ καὶ πᾶσα ἡ ἐκκλησία Ἰσραὴλ, ἵνα ἄγωνται αἱ ἡμέραι ἐγκαινισμοῦ τοῦ θυσιαστηρίου ἐν τοῖς καιροῖς αὐτῶν ἐνιαυτὸν κατʼ ἐνιαυτὸν ἡμέρας ὀκτὼ, ἀπὸ τῆς πέμπτης καὶ εἰκάδος τοῦ μηνὸς Χασελεῦ, μετʼ εὐφροσύνης καὶ χαρᾶς.
\vs{60}Καὶ ᾠκοδόμησαν ἐν τῷ καιρῷ ἐκείνῳ τὸ ὄρος Σιὼν, κυκλόθεν τείχη ὑψηλὰ καὶ πύργους ὀχυροὺς, μήποτε παραγενηθέντα τὰ ἔθνη καταπατήσωσιν αὐτὰ, ὡς ἐποίησαν τοπρότερον.
\vs{61}Καὶ ἐπέταξεν ἐκεῖ δύναμιν τηρεῖν αὐτὸ, καὶ ὠχύρωσαν αὐτὸ τηρεῖν τὴν Βαιθσούραν, τοῦ ἔχειν τὸν λαὸν ὀχύρωμα κατὰ πρόσωπον τῆς Ἰδουμαίας.

\ch{5}
Καὶ ἐγένετο ὅτε ἤκουσαν τὰ ἔθνη κυκλόθεν ὅτι ᾠκοδομήθη τὸ θυσιαστήριον, καὶ ἐνεκαινίσθη τὸ ἁγίασμα ὡς τοπρότερον, καὶ ὠργίσθησαν σφόδρα.
\vs{2}Καὶ ἐβουλεύσαντο τοῦ ἆραι τὸ γένος Ἰακὼβ τοὺς ὄντας ἐν μέσῳ αὐτῶν, καὶ ἤρξαντο τοῦ θανατοῦν ἐν τῷ λαῷ καὶ ἐξαίρειν.

\vs{3}Καὶ ἐπολέμει Ἰούδας πρὸς τοὺς υἱοὺς Ἡσαῦ ἐν τῇ Ἰδουμαίᾳ τὴν Ἀκραβαττίνην, ὅτι περιεκάθηντο τὸν Ἰσραὴλ, καὶ ἐπάταξεν αὐτοὺς πληγὴν μεγάλην, καὶ συνέστειλεν αὐτοὺς, καὶ ἔλαβε τὰ σκῦλα αὐτῶν.
\vs{4}Καὶ ἐμνήσθη τῆς κακίας υἱῶν Βαιὰν, οἳ ἦσαν τῷ λαῷ εἰς παγίδα καὶ εἰς σκάνδαλον ἐν τῷ ἐνεδρεύειν αὐτοὺς ἐν ταῖς ὁδοῖς.
\vs{5}Καὶ συνεκλείσθησαν ὑπʼ αὐτοῦ ἐν τοῖς πύργοις, καὶ παρενέβαλεν ἐπʼ αὐτοὺς, καὶ ἀνεθεμάτισεν αὐτοὺς, καὶ ἐνεπύρισε τοὺς πύργους αὐτῆς ἐν πυρὶ σὺν πᾶσι τοῖς ἐνοῦσι.

\vs{6}Καὶ διεπέρασεν ἐπὶ τοὺς υἱοὺς Ἀμμὼν, καὶ εὗρε χεῖρα κραταιὰν καὶ λαὸν πολὺν, καὶ Τιμόθεον ἡγούμενον αὐτῶν.
\vs{7}καὶ συνῆψε πρὸς αὐτοὺς πολέμους πολλοὺς, καὶ συνετρίβησαν πρὸ προσώπου αὐτοῦ, καὶ ἐπάταξεν αὐτούς.
\vs{8}Καὶ προκατελάβετο τὴν Ἰαζὴρ καὶ τὰς θυγατέρας αὐτῆς, καὶ ἀνέστρεψεν εἰς τὴν Ἰούδαίαν.

\vs{9}Καὶ ἐπισυνήχθησαν τὰ ἔθνη τὰ ἐν τῇ Γαλαὰδ ἐπὶ τὸν Ἰσραὴλ τοὺς ὄντας ἐπὶ τοῖς ὁρίοις αὐτῶν τοῦ ἐξᾶραι αὐτούς· καὶ ἔφυγον εἰς Δάθεμα τὸ ὀχύρωμα.
\vs{10}Καὶ ἀπέστειλαν γράμματα πρὸς Ἰούδαν καὶ τοὺς ἀδελφοὺς αὐτοῦ, λέγοντες, ἐπισυνηγμένα ἐστὶν ἐφʼ ἡμᾶς τὰ ἔθνη τὰ κύκλῳ ἡμῶν τοῦ ἐξᾶραι ἡμᾶς.
\vs{11}Καὶ ἑτοιμάζονται ἐλθεῖν καὶ προκαταλαβέσθαι τὸ ὀχύρωμα εἰς ὃ κατεφύγομεν, καὶ Τιμόθεος ἡγεῖται τῆς δυνάμεως αὐτῶν.

\vs{12}Νῦν οὖν ἐλθὼν ἐξελοῦ ἡμᾶς ἐκ χειρὸς αὐτῶν, ὅτι πέπτωκεν ἐξ ἡμῶν πλῆθος.
\vs{13}Καὶ πάντες οἱ ἀδελφοὶ ἡμῶν οἱ ὄντες ἐν τοῖς Τωβίου τεθανάτωνται, καὶ ᾐχμαλωτίκασι τὰς γυναῖκας αὐτῶν καὶ τὰ τέκνα αὐτῶν καὶ τὴν ἀποσκευὴν, καὶ ἀπώλεσαν ἐκεῖ ὡς μίαν χιλιαρχίαν ἀνδρῶν.

\vs{14}Ἔτι αἱ ἐπιστολαὶ ἀνεγινώσκοντο, καὶ ἰδοὺ ἄγγελοι ἕτεροι παρεγένοντο ἐκ τῆς Γαλιλαίας διεῤῥηχότες τὰ ἱμάτια, ἀπαγγέλλοντες κατὰ τὰ ῥήματα ταῦτα,
\vs{15}λέγοντες ἐπισυνῆχθαι ἐπʼ αὐτοὺς ἐκ Πτολεμαΐδος καὶ Τύρου καὶ Σιδῶνος καὶ πάσης Γαλιλαίας ἀλλοφύλων, τοῦ ἐξαναλῶσαι ἡμᾶς.

\vs{16}Ὡς δὲ ἤκουσεν Ἰούδας καὶ ὁ λαὸς τοὺς λόγους τούτους, ἐπισυνήχθη ἐκκλησία μεγάλη, βουλεύσασθαι τί ποιήσωσι τοῖς ἀδελφοῖς αὐτῶν τοῖς οὖσιν ἐν θλίψει, καὶ πολεμουμένοις ὑπʼ αὐτῶν.
\vs{17}Καὶ εἶπεν Ἰούδας Σίμωνι τῷ ἀδελφῷ αὐτοῦ, ἐπίλεξον σεαυτῷ ἄνδρας, καὶ πορεύου καὶ ῥῦσαι τοὺς ἀδελφούς σου τοὺς ἐν τῇ Γαλιλαίᾳ· ἐγὼ δὲ καὶ Ἰωνάθαν ὁ ἀδελφός μου πορευσόμεθα εἰς τὴν Γαλααδίτιν.
\vs{18}Καὶ κατέλιπεν Ἰώσηφον τὸν τοῦ Ζαχαρίου, καὶ Ἀζαρίαν, ἡγουμένους τοῦ λαοῦ, μετὰ τῶν ἐπιλοίπων τῆς δυνάμεως, ἐν τῇ Ἰουδαίᾳ εἰς τήρησιν.
\vs{19}Καὶ ἐνετείλατο αὐτοῖς, λέγων, πρόστητε τοῦ λαοῦ τούτου, καὶ μὴ συνάψητε πόλεμον πρὸς τὰ ἔθνη ἕως τοῦ ἐπιστρέψαι ἡμᾶς.
\vs{20}Καὶ ἐμερίσθησαν Σίμωνι ἄνδρες τρισχίλιοι τοῦ πορευθῆναι εἰς τὴν Γαλιλαίαν, Ἰούδᾳ δὲ ἄνδρες ὀκτακισχίλιοι εἰς τὴν Γαλααδίτιν.

\vs{21}Καὶ ἐπορεύθη Σίμων εἰς τὴν Γαλιλαίαν, καὶ συνῆψε πολέμους πολλοὺς πρὸς τὰ ἔθνη, καὶ συνετρίβη τὰ ἔθνη ἀπὸ προσώπου αὐτοῦ,
\vs{22}καὶ ἐδίωξεν αὐτοὺς ἕως τῆς πύλης Πτολεμαΐδος· καὶ ἔπεσον ἐκ τῶν ἐθνῶν εἰς τρισχιλίους ἄνδρας, καὶ ἔλαβε τὰ σκῦλα αὐτῶν.
\vs{23}Καὶ παρέλαβε τοὺς ἐν τῇ Γαλιλαίᾳ καὶ ἐν Ἀρβάττοις σὺν ταῖς γυναιξὶ καὶ τοῖς τέκνοις, καὶ πάντα ὅσα ἦν αὐτοῖς, καὶ ἤγαγεν εἰς τὴν Ἰουδαίαν μετʼ εὐφροσύνης μεγάλης.

\vs{24}Καὶ Ἰούδας ὁ Μακκαβαῖος καὶ Ἰωνάθαν ὁ ἀδελφὸς αὐτοῦ διέβησαν τὸν Ἰορδάνην, καὶ ἐπορεύθησαν ὁδὸν τριῶν ἡμερῶν ἐν τῷ ἐρήμῳ.
\vs{25}Καὶ συνήντησαν τοῖς Ναβαταίοις, καὶ ἀπήντησαν αὐτοῖς εἰρηνικῶς, καὶ διηγήσαντο αὐτοῖς ἅπαντα τὰ συμβάντα τοῖς ἀδελφοῖς αὐτῶν ἐν τῇ Γαλααδίτιδι.
\vs{26}Καὶ ὅτι πολλοὶ ἐξ αὐτῶν συνειλημμένοι εἰσὶν εἰς Βόσσορα, καὶ Βοσὸρ, ἐν Ἀλέμοις, Χασφὼρ, Μακὲδ, καὶ Καρναΐν· πᾶσαι αἱ πόλεις αὗται ὀχυραὶ καὶ μεγάλαι·
\vs{27}καὶ ἐν ταῖς λοιπαῖς πόλεσι τῆς Γαλααδίτιδός εἰσι συνειλημμένοι, καὶ εἰς αὔριον τάσσονται παρεμβάλλειν ἐπὶ τὰ ὀχυρώματα, καὶ καταλαβέσθαι, καὶ ἐξᾶραι πάντας τούτους ἐν ἡμέρᾳ μιᾷ.

\vs{28}Καὶ ἀπέστρεψεν Ἰούδας καὶ ἡ παρεμβολὴ αὐτοῦ ὁδὸν εἰς τὴν ἔρημον εἰς Βοσὸρ, ἄφνω· καὶ κατελάβετο τὴν πόλιν, καὶ ἀπέκτεινε πᾶν ἀρσενικὸν ἐν στόματι ῥομφαίας, καὶ ἔλαβε πάντα τὰ σκῦλα αὐτῶν, καὶ ἐνέπρησεν αὐτὴν πυρί.
\vs{29}Καὶ ἀπῇρεν ἐκεῖθεν νυκτὸς, καὶ ἐπορεύετο ἕως ἐπὶ τὸ ὀχύρωμα.

\vs{30}Καὶ ἐγένετο ἑωθινὴ, καὶ ᾖραν τοὺς ὀφθαλμοὺς αὐτῶν, καὶ ἰδοὺ λαὸς πολὺς οὗ οὐκ ἦν ἀριθμὸς, αἴροντες κλίμακας καὶ μηχανὰς καταλαβέσθαι τὸ ὀχύρωμα, καὶ ἐπολέμουν αὐτούς.
\vs{31}Καὶ εἶδεν Ἰούδας ὅτι ἦρκται ὁ πόλεμος, καὶ ἡ κραυγὴ τῆς πόλεως ἀνέβη εἰς τὸν οὐρανὸν σάλπιγξι καὶ φωνῇ μεγάλῃ.
\vs{32}Καὶ εἶπε τοῖς ἀνδράσι τῆς δυνάμεως, πολεμήσατε σήμερον ὑπὲρ τῶν ἀδελφῶν ὑμῶν.
\vs{33}Καὶ ἐξῆλθεν ἐν τρισὶν ἀρχαῖς ἐξ ὄπισθεν αὐτῶν· καὶ ἐσάλπισαν ταῖς σάλπιγξι, καὶ ἐβόησαν ἐν προσευχῇ.

\vs{34}Καὶ ἐπέγνω ἡ παρεμβολὴ Τιμοθέου ὅτι Μακκαβαῖός ἐστι, καὶ ἔφυγον ἀπὸ προσώπου αὐτοῦ, καὶ ἐπάταξεν αὐτοὺς πληγὴν μεγάλην, καὶ ἔπεσον ἐξ αὐτῶν ἐν ἐκείνῃ τῇ ἡμέρᾳ εἰς ὀκτακισχιλίους ἄνδρας.

\vs{35}Καὶ ἀπέκλινεν εἰς Μασφὰ, καὶ ἐπολέμησεν αὐτὴν, καὶ προκατελάβετο αὐτὴν, καὶ ἀπέκτεινε πᾶν ἀρσενικὸν αὐτῆς, καὶ ἔλαβε τὰ σκῦλα αὐτῆς, καὶ ἐνέπρησεν αὐτὴν πυρί.
\vs{36}Ἐκεῖθεν ἀπῇρε, καὶ προκατελάβετο τὴν Χασφὼν, Μακὲδ, Βοσὸρ, καὶ τὰς λοιπὰς πόλεις τῆς Γαλααδίτιδος.

\vs{37}Μετὰ δὲ τὰ ῥήματα ταῦτα συνήγαγε Τιμόθεος παρεμβολὴν ἄλλην, καὶ παρενέβαλε κατὰ πρόσωπον Ῥαφὼν ἐκ πέραν τοῦ χειμάῤῥου.
\vs{38}Καὶ ἀπέστειλεν Ἰούδας κατασκοπεῦσαι τὴν παρεμβολήν, καὶ ἀπήγγειλαν αὐτῷ, λέγοντες, ἐπισυνηγμένα εἰσὶ πρὸς αὐτοὺς πάντα τὰ ἔθνη τὰ κύκλῳ ἡμῶν, δύναμις πολλὴ σφόδρα.
\vs{39}Καὶ Ἄραβας μεμίσθωται εἰς βοήθειαν αὐτοῖς, καὶ παρενέβαλον πέραν τοῦ χειμάῤῥου ἕτοιμοι τοῦ ἐλθεῖν ἐπὶ σὲ εἰς πόλεμον· καὶ ἐπορεύθη Ἰούδας εἰς συνάντησιν αὐτῶν.

\vs{40}Καὶ εἶπε Τιμόθεος τοῖς ἄρχουσι τῆς δυνάμεως αὐτοῦ, ἐν τῷ ἐγγίζειν Ἰούδαν καὶ τὴν παρεμβολὴν αὐτοῦ ἐπὶ τὸν χειμάῤῥουν τοῦ ὕδατος, ἐὰν διαβῇ πρὸς ἡμᾶς πρότερος, οὐ δυνησόμεθα ὑποστῆναι αὐτὸν, ὅτι δυνάμενος δυνήσεται πρὸς ἡμᾶς.
\vs{41}Ἐὰν δὲ δειλωθῇ, καὶ παρεμβάλῃ πέραν τοῦ ποταμοῦ, διαπεράσομεν πρὸς αὐτὸν, καὶ δυνησόμεθα πρὸς αὐτόν.

\vs{42}Ὡς δὲ ἤγγισεν Ἰούδας ἐπὶ τὸν χειμάῤῥουν τοῦ ὕδατος, ἔστησε τοὺς γραμματεῖς τοῦ λαοῦ ἐπὶ τοῦ χειμάῤῥου, καὶ ἐνετείλατο αὐτοῖς, λέγων, μὴ ἀφῆτε πάντα ἄνθρωπον παρεμβαλεῖν, ἀλλʼ ἐρχέσθωσαν πάντες εἰς τὸν πόλεμον.
\vs{43}Καὶ διεπέρασεν ἐπʼ αὐτοὺς πρότερος, καὶ πᾶς ὁ λαὸς ὅπισθεν αὐτοῦ· καὶ συνετρίβησαν πρὸ προσώπου αὐτοῦ πάντα τὰ ἔθνη, καὶ ἔῤῥιψαν τὰ ὅπλα αὐτῶν, καὶ ἔφυγον εἰς τὸ τέμενος ἐν Καρναΐν.
\vs{44}Καὶ προκατελάβοντο τὴν πόλιν, καὶ τὸ τέμενος ἐνεπύρισαν ἐν πυρὶ σὺν πᾶσι τοῖς ἐν αὐτῷ· καὶ ἐτροπώθη ἡ Καρναῒν, καὶ οὐκ ἐδύναντο ἔτι ὑποστῆναι κατὰ πρόσωπον Ἰούδα.

\vs{45}Καὶ συνήγαγεν Ἰούδας πάντα Ἰσραὴλ τοὺς ἐν τῇ Γαλααδίτιδι ἀπὸ μικροῦ ἕως μεγάλου, καὶ τὰς γυναῖκας αὐτῶν, καὶ τὰ τέκνα αὐτῶν, καῖ τὴν ἀποσκευὴν, παρεμβολὴν μεγάλην σφόδρα, ἐλθεῖν εἰς γῆν Ἰούδα.
\vs{46}Καὶ ἦλθον ἕως Ἐφρών· καὶ αὕτη ἡ πόλις μεγάλη ἐπὶ τῆς εἰσόδου ὀχυρὰ σφόδρα· οὐκ ἦν ἐκκλῖναι ἀπʼ αὐτῆς δεξιὰν ἢ ἀριστερὰν, ἀλλʼ ἢ διὰ μέσου αὐτῆς πορεύεσθαι.
\vs{47}Καὶ ἀπέκλεισαν αὐτοὺς οἱ ἐκ τῆς πόλεως, καὶ ἐνέφραξαν τὰς πύλας λίθοις.
\vs{48}Καὶ ἀπέστειλε πρὸς αὐτοὺς Ἰούδας λόγοις εἰρηνικοῖς, λέγων, διελευσόμεθα διὰ τῆς γῆς σου τοῦ ἀπελθεῖν εἰς τὴν γῆν ἡμῶν, καὶ οὐδεὶς κακοποιήσει ὑμᾶς, πλὴν τοῖς ποσὶ παρελευσόμεθα· καὶ οὐκ ἠβούλοντο ἀνοῖξαι αὐτῷ.

\vs{49}Καὶ ἐπέταξεν Ἰούδας κηρύξαι ἐν τῇ παρεμβολῇ, τοῦ παρεμβαλεῖν ἕκαστον ἐν ᾧ ἐστι τόπῳ.
\vs{50}Καὶ παρενέβαλον οἱ ἄνδρες τῆς δυνάμεως, καὶ ἐπολέμησαν τὴν πόλιν ὅλην τὴν ἡμέραν ἐκείνην καὶ ὅλην τὴν νύκτα, καὶ παρεδόθη ἡ πόλις ἐν χερσὶν αὐτοῦ.
\vs{51}Καὶ ἀπώλεσε πᾶν ἀρσενικὸν ἐν στόματι ῥομφαίας, καὶ ἐξεῤῥίζωσεν αὐτὴν, καὶ ἔλαβε τὰ σκῦλα αὐτῆς, καὶ διῆλθε διὰ τῆς πόλεως ἐπάνω τῶν ἀπεκταμμένων.

\vs{52}Καὶ διέβησαν τὸν Ἰορδάνην εἰς τὸ πεδίον τὸ μέγα κατὰ πρόσωπον Βαιθσάν.
\vs{53}Καὶ ἦν Ἰούδας ἐπισυνάγων τοὺς ἐσχατίζοντας, καὶ παρακαλῶν τὸν λαὸν κατὰ πᾶσαν τὴν ὁδὸν, ἕως οὗ ἦλθον εἰς γῆν Ἰούδα.
\vs{54}Καὶ ἀνέβησαν εἰς τὸ ὄρος Σιὼν ἐν εὐφροσύνῃ καὶ χαρᾷ· καὶ προσήγαγον ὁλοκαυτώματα, ὅτι οὐκ ἔπεσεν ἐξ αὐτῶν οὐθεὶς ἕως τοῦ ἐπιστρέψαι ἐν εἰρήνῃ.

\vs{55}Καὶ ἐν ταῖς ἡμέραις αἷς ἦν Ἰούδας καὶ Ἰωνάθαν ἐν τῇ Γαλαὰδ, καὶ Σίμων ὁ ἀδελφὸς αὐτοῦ ἐν τῇ Γαλιλαίᾳ κατὰ πρόσωπον Πτολεμαΐδος,
\vs{56}ἤκουσεν Ἰωσὴφ ὁ τοῦ Ζαχαρίου, καὶ Ἀζαρίας, ἄρχοντες τῆς δυνάμεως, τῶν ἀνδραγαθιῶν καὶ τοῦ πολέμου οἷα ἐποίησαν,
\vs{57}καὶ εἶπε, ποιήσωμεν καὶ αὐτοὶ ἑαυτοῖς ὄνομα, καὶ πορευθῶμεν πολεμῆσαι πρὸς τὰ ἔθνη τὰ κύκλῳ ἡμῶν.

\vs{58}Καὶ παρήγγειλαν τοῖς ἀπὸ τῆς δυνάμεως τῆς μετʼ αὐτῶν, καὶ ἐπορεύθησαν ἐπὶ Ἰάμνειαν.
\vs{59}Καὶ ἐξῆλθε Γοργίας ἐκ τῆς πόλεως, καὶ οἱ ἄνδρες αὐτοῦ, εἰς συνάντησιν αὐτοῖς εἰς πόλεμον.
\vs{60}Καὶ ἐτροπώθη Ἰώσηφος καὶ Ἀζαρίας, καὶ ἐδιώχθησαν ἕως τῶν ὁρίων τῆς Ἰουδαίας· καὶ ἔπεσον ἐν τῇ ἡμέρᾳ ἐκείνῃ ἐκ τοῦ λαοῦ τοῦ Ἰσραὴλ εἰς δισχιλίους ἄνδρας.
\vs{61}Καὶ ἐγενήθη τροπὴ μεγάλη ἐν τῷ λαῷ Ἰσραὴλ, ὅτι οὐκ ἤκουσαν Ἰούδα καὶ τῶν ἀδελφῶν αὐτοῦ, οἰόμενοι ἀνδραγαθῆσαι.
\vs{62}Αὐτοὶ δὲ οὐκ ἦσαν ἐκ τοῦ σπέρματος τῶν ἀνδρῶν ἐκείνων, οἷς ἐδόθη σωτηρία Ἰσραὴλ διὰ χειρὸς αὐτῶν.
\vs{63}Καὶ ὁ ἀνὴρ Ἰούδας καὶ οἱ ἀδελφοὶ αὐτοῦ ἐδοξάσθησαν σφόδρα ἑναντίον παντὸς Ἰσραὴλ, καὶ τῶν ἐθνῶν πάντων, οὗ ἠκούετο τὸ ὄνομα αὐτῶν.
\vs{64}Καὶ ἐπισυνήγοντο πρὸς αὐτοὺς εὐφημοῦντες.

\vs{65}Καὶ ἐξῆλθεν Ἰούδας καὶ οἱ ἀδελφοὶ αὐτοῦ, καὶ ἐπολέμουν τοὺς υἱοὺς Ἡσαῦ ἐν τῇ γῇ πρὸς Νότον, καὶ ἐπάταξε τὴν Χεβρὼν καὶ τὰς θυγατέρας αὐτῆς, καὶ καθεῖλε τὸ ὀχύρωμα αὐτῆς, καὶ τοὺς πύργους αὐτῆς ἐνέπρησε κυκλόθεν.
\vs{66}Καὶ ἀπῇρε τοῦ πορευθῆναι εἰς γῆν ἀλλοφύλων, καὶ διεπορεύετο τὴν Σαμάρειαν.

\vs{67}Ἐν τῇ ἡμέρᾳ ἐκείνῃ ἔπεσον ἱερεῖς ἐν πολέμῳ βουλόμενοι ἀνδραγαθῆσαι, ἐν τῷ αὐτοὺς ἐξελθεῖν εἰς πόλεμον ἀβουλεύτως.
\vs{68}Καὶ ἐξέκλινεν Ἰούδας εἰς Ἄζωτον γῆν ἀλλοφύλων, καὶ καθεῖλε τοὺς βωμοὺς αὐτῶν, καὶ τὰ γλυπτὰ τῶν θεῶν αὐτῶν κατέκαυσε πυρὶ, καὶ ἐσκύλευσε τὰ σκῦλα τῶν πόλεων, καὶ ἐπέστρεψεν εἰς τὴν γῆν Ἰούδα.

\ch{6}
Καὶ ὁ βασιλεὺς Ἀντίοχος διεπορεύετο τὰς ἐνάνω χώρας, καὶ ἤκουσεν ὅτι ἐστὶν Ἐλυμαῒς ἐν τῇ Περσίδι πόλις ἔνδοξος πλούτῳ, ἀργυρίῳ τε καὶ χρυσίῳ,
\vs{2}καὶ τὸ ἱερὸν τὸ ἐν αὐτῇ πλούσιον σφόδρα, καὶ ἐκεῖ καλύμματα χρυσᾶ, καὶ θώρακες, καὶ ὅπλα ἃ κατέλιπεν ἐκεῖ Ἀλέξανδρος ὁ Φιλίππου, βασιλεὺς ὁ Μακεδὼν, ὃς ἐβασίλευσε πρῶτος ἐν τοῖς Ἕλλησι.
\vs{3}Καὶ ἦλθε καὶ ἐζήτει καταλαβέσθαι τὴν πόλιν, καὶ προνομεῦσαι αὐτὴν, καὶ οὐκ ἠδυνάσθη, ὅτι ἐγνώσθη ὁ λόγος τοῖς ἐκ τῆς πόλεως.
\vs{4}Καὶ ἀνέστησαν αὐτῷ εἰς πόλεμον, καὶ ἔφυγε καὶ ἀπῇρεν ἐκεῖθεν μετὰ λύπης μεγάλης, ἀποστρέψαι εἰς Βαβυλῶνα.

\vs{5}Καὶ ἦλθεν ἀπαγγέλλων τις αὐτῷ εἰς τὴν Περσίδα, ὅτι τετρόπωνται αἱ παρεμβολαὶ αἱ πορευθεῖσαι εἰς γῆν Ἰούδα.
\vs{6}Καὶ ἐπορεύθη Λυσίας δυνάμει ἰσχυρᾷ ἐν πρώτοις, καὶ ἀνετράπη ἀπὸ προσώπου αὐτῶν, καὶ ἐπίσχυσαν ὅπλοις καὶ δυνάμει καὶ σκύλοις πολλοῖς οἷς ἔλαβον ἀπὸ τῶν παρεμβολῶν ὧν ἐξέκοψαν.
\vs{7}Καὶ καθεῖλον τὸ βδέλυγμα ὃ ᾠκοδόμησεν ἐπὶ τὸ θυσιαστήριον τὸ ἐν Ἱερουσαλὴμ, καὶ τὸ ἁγίασμα καθὼς τὸ πρότερον ἐκύκλωσαν τείχεσιν ὑψηλοῖς, καὶ τὴν Βαιθσούραν πόλιν αὐτοῦ.

\vs{8}Καὶ ἐγένετο ὡς ἤκουσεν ὁ βασιλεὺς τοὺς λόγους τούτους, ἐθαμβήθη καὶ ἐσαλεύθη σφόδρα· καὶ ἔπεσεν ἐπὶ τὴν κοίτην, καὶ ἐνέπεσεν εἰς ἀῤῥωστίαν ἀπὸ τῆς λύπης, ὅτι οὐκ ἐγένετο αὐτῷ καθὼς ἐνεθυμεῖτο.
\vs{9}Καὶ ἦν ἐκεῖ ἡμέρας πλείους, ὅτι ἀνεκαινίσθη ἐπʼ αὐτὸν λύπη μεγάλη, καὶ ἐλογίσατο ὅτι ἀποθνήσκει.
\vs{10}Καὶ ἐκάλεσε πάντας τοὺς φίλους αὐτοῦ, καὶ εἶπε πρὸς αὐτοὺς, ἀφίσταται ὁ ὕπνος ἀπὸ τῶν ὀφθαλμῶν μου, καὶ συμπέπτωκα τῇ καρδίᾳ ἀπὸ τῆς μερίμνης.
\vs{11}Καὶ εἶπα τῇ καρδίᾳ μου, ἕως τίνος θλίψεως ἦλθον καὶ κλύδωνος μεγάλου, ἐν ᾧ νῦν εἰμι; ὅτι χρηστὸς καὶ ἀγαπώμενος ἤμην ἐν τῇ ἐξουσίᾳ μου.
\vs{12}Νῦν δὲ μιμνήσκομαι τῶν κακῶν ὧν ἐποίησα ἐν Ἱερουσαλὴμ, καὶ ἔλαβον πάντα τὰ σκεύη τὰ χρυσᾶ καὶ τὰ ἀργυγᾶ τὰ ἐν αὐτῇ, καὶ ἐξαπέστειλα ἐξᾶραι τοὺς κατοικοῦντας Ἰούδα διακενῆς.
\vs{13}Ἔγνων οὖν ὅτι χάριν τούτων εὗρόν με τὰ κακὰ ταῦτα· καὶ ἰδοὺ ἀπόλλυμαι λύπῃ μεγάλῃ ἐν γῇ ἀλλοτρίᾳ.

\vs{14}Καὶ ἐκάλεσε Φίλιππον ἕνα τῶν φίλων αὐτοῦ, καὶ κατέστησεν αὐτὸν ἐπὶ πάσης τῆς βασιλείας αὐτοῦ.
\vs{15}Καὶ ἔδωκεν αὐτῷ τὸ διάδημα καὶ τὴν στολὴν αὐτοῦ καὶ τὸν δακτύλιον, τοῦ ἀγαγεῖν Ἀντίοχον τὸν υἱὸν αὐτοῦ, καὶ ἐκθρέψαι αὐτὸν τοῦ βασιλεύειν.
\vs{16}Καὶ ἀπέθανεν ἐκεῖ Ἀντίοχος ὁ βασιλεὺς ἔτους ἐννάτου καὶ τεσσαρακοστοῦ καὶ ἑκατοστοῦ.
\vs{17}Καὶ ἐπέγνω Λυσίας ὅτι τέθνηκεν ὁ βασιλεύς, καὶ κατέστησε βασιλεύειν Ἀντίοχον τὸν υἱὸν αὐτοῦ ἀντʼ αὐτοῦ, ὃν ἐξέθρεψε νεώτερον, καὶ ἐκάλεσε τὸ ὄνομα αὐτοῦ Εὐπάτορα.

\vs{18}Καὶ οἱ ἐκ τῆς ἄκρας ἦσαν συγκλείοντες τὸν Ἰσραὴλ κύκλῳ τῶν ἁγίων, καὶ ζητοῦντες τὰ κακὰ διʼ ὅλου, καὶ στήριγμα τοῖς ἔθνεσι
\vs{19}Καὶ ἐλογίσατο Ἰούδας ἐξᾶραι αὐτούς· καὶ ἐξεκκλησίασε πάντα τὸν λαὸν τοῦ περικαθίσαι ἐπʼ αὐτοὺς.
\vs{20}Καὶ συνήχθησαν ἅμα, καὶ περιεκάθισαν ἐπʼ αὐτοὺς ἔτους πεντηκοστοῦ καὶ ἑκατοστοῦ, καὶ ἐποίησεν ἐπʼ αὐτοὺς βελοστάσεις καὶ μηχανάς.

\vs{21}Καὶ ἐξῆλθον ἐξ αὐτῶν ἐκ τοῦ συγκλεισμοῦ, καὶ ἐκολλήθησαν αὐτοῖς τινὲς τῶν ἀσεβῶν ἐξ Ἰσραήλ,
\vs{22}καὶ ἐπορεύθησαν πρὸς τὸν βασιλέα, καὶ εἶπον, ἕως πότε οὐ ποιήσῃ κρίσιν, καὶ ἐκδικήσεις τοὺς ἀδελφοὺς ἡμῶν;
\vs{23}Ἡμεῖς εὐδοκοῦμεν δουλεύειν τῷ πατρί σου, καὶ πορεύεσθαι τοῖς ὑπʼ αὐτοῦ λεγομένοις, καὶ κατακολουθεῖν τοῖς προστάγμασιν αὐτοῦ.
\vs{24}Καὶ περικάθηνται εἰς τὴν ἄκραν υἱοὶ τοῦ λαοῦ ἡμῶν, χάριν τούτου καὶ ἀλλοτριοῦνται ἀφʼ ἡμῶν· πλὴν ὅσοι εὑρίσκοντο ἀφʼ ἡμῶν ἐθανατοῦντο, καὶ αἱ κληρονομίαι ἡμῶν διηρπάζοντο.

\vs{25}Καὶ οὐκ ἐφʼ ἡμᾶς μόνον ἐξέτειναν χεῖρα, ἀλλὰ καὶ ἐπὶ πάντα τὰ ὅρια αὐτῶν.
\vs{26}Καὶ ἰδοὺ παρεμβεβλήκασι σήμερον ἐπὶ τὴν ἄκραν ἐν Ἱερουσαλὴμ, τοῦ καταλαβέσθαι αὐτὴν, καὶ τὸ ἁγίασμα, καὶ τὴν Βαιθσούραν ὠχύρωσαν.
\vs{27}Καὶ ἐὰν μὴ προκαταλάβῃ αὐτοὺς διατάχους, μείζονα τούτων ποιήσουσι, καὶ οὐ δυνήσῃ τοῦ κατασχεῖν αὐτῶν.

\vs{28}Καὶ ὠργίσθη ὁ βασιλεὺς ὅτε ἤκουσε, καὶ συνήγαγε πάντας τοὺς φίλους αὐτοῦ, καὶ τοὺς ἄρχοντας τῆς δυνάμεως αὐτοῦ, καὶ τοὺς ἐπὶ τῶν ἡνιῶν.
\vs{29}Καὶ ἀπὸ βασιλειῶν ἑτέρων καὶ ἀπὸ νήσων θαλασσῶν ἦλθον πρὸς αὐτὸν δυνάμεις μισθωταί.
\vs{30}Καὶ ἦν ὁ ἀριθμὸς τῶν δυνάμεων αὐτοῦ ἑκατὸν χιλιάδες τῶν πεζῶν, καὶ εἴκοσι χιλιάδες ἵππων, καὶ ἐλέφαντες δύο καὶ τριάκοντα εἰδότες πόλεμον.
\vs{31}Καὶ ἤλθοσαν διὰ τῆς Ἰδουμαίας, καὶ παρενεβάλοσαν ἐπὶ Βαιθσούραν, καὶ ἐπολέμησαν ἡμέρας πολλὰς, καὶ ἐποίησαν μηχανάς· καὶ ἐξῆλθον, καὶ ἐνεπύρισαν αὐτὰς ἐν πυρὶ, καὶ ἐπολέμησαν ἀνδρωδῶς.

\vs{32}Καὶ ἀπῇρεν Ἰούδας ἀπὸ τῆς ἄκρας, καὶ παρενέβαλεν εἰς Βαιθζαχαρία ἀπέναντι τῆς παρεμβολῆς τοῦ βασιλέως.
\vs{33}Καὶ ὤρθρισεν ὁ βασιλεὺς τοπρωῒ, καὶ ἀπῇρε τὴν παρεμβολὴν ἐν ὁρμήματι αὐτῆς κατὰ τὴν ὁδὸν Βαιθζαχαρία, καὶ διεσκευάσθησαν αἱ δυνάμεις εἰς τὸν πόλεμον, καὶ ἐσάλπισαν ταῖς σάλπιγξι.

\vs{34}Καὶ τοῖς ἐλέφασιν ἔδειξαν αἷμα σταφυλῆς καὶ μόρων, τοῦ παραστῆσαι αὐτοὺς εἰς τὸν πόλεμον.
\vs{35}Καὶ διεῖλον τὰ θηρία εἰς τὰς φάλαγγας, καὶ παρέστησαν ἑκάστῳ ἐλέφαντι χιλίους ἄνδρας τεθωρακισμένους ἐν ἁλυσιδωτοῖς, καὶ περικεφαλαῖαι χαλκαῖ ἐπὶ τῶν κεφαλῶν αὐτῶν, καὶ πεντακόσιοι ἵπποι διατεταγμένοι ἑκάστῳ θηρίῳ ἐκλελεγμένοι.
\vs{36}Οὗτοι πρὸ καιροῦ, οὗ ἐὰν ἦν τὸ θηρίον, ἦσαν, καὶ οὗ ἐὰν ἐπορεύετο, ἐπορεύοντο ἅμα, οὐκ ἀφίσταντο ἀπʼ αὐτοῦ.
\vs{37}Καὶ πύργοι ξύλινοι ἐπʼ αὐτοὺς ὀχυροὶ σκεπαζόμενοι ἐφʼ ἑκάστου θηρίου, ἐζωσμένοι ἐπʼ αὐτοῦ μηχαναῖς· καὶ ἐφʼ ἑκάστου ἄνδρες δυνάμεως δύο καὶ τριάκοντα οἱ πολεμοῦντες ἐπʼ αὐτοῖς, καὶ ὁ Ἰνδὸς αὐτοῦ.

\vs{38}Καὶ τὴν ἐπίλοιπον ἵππον ἔνθεν καὶ ἔνθεν ἔστησαν ἐπὶ τὰ δύο μέρη τῆς παρεμβολῆς, κατασείοντες καὶ καταφρασσόμενοι ἐν ταῖς φάραγξιν.
\vs{39}Ὡς δὲ ἔστιλβεν ὁ ἥλιος ἐπὶ τὰς χρυσᾶς καὶ χαλκᾶς ἀσπίδας, ἔστιλβε τὰ ὄρη ἀπʼ αὐτῶν, καὶ κατηύγαζεν ὡς λαμπάδες πυρός.
\vs{40}Καὶ ἐξετάθη μέρος τι τῆς παρεμβολῆς τοῦ βασιλέως ἐπὶ τὰ ὑψηλὰ ὄρη, καί τινες ἐπὶ ταπεινά· καὶ ἤρχοντο ἀσφαλῶς καὶ τεταγμένως.
\vs{41}Καὶ ἐσαλεύοντο πάντες οἱ ἀκούοντες φωνῆς πλήθους αὐτῶν, καὶ ὁδοιπαρίας τοῦ πλήθους, καὶ συγκρουσμοῦ τῶν ὅπλων· ἦν γὰρ ἡ παρεμβολὴ μεγάλη σφόδρα καὶ ἰσχυρά.

\vs{42}Καὶ ἤγγισεν Ἰούδας καὶ ἡ παρεμβολὴ αὐτοῦ εἰς παράταξιν· καὶ ἔπεσον ἀπὸ τῆς παρεμβολῆς τοῦ βασιλέως ἑξακόσιοι ἄνδρες.
\vs{43}Καὶ εἶδεν Ἐλεάζαρ ὁ Σαυαρὰν ἓν τῶν θηρίων τεθωρακισμένον θώρακι βασιλικῷ, καὶ ἦν ὑπεράγον πάντα τὰ θηρία, καὶ ὤφθη ὅτι ἐν αὐτῷ ἐστιν ὁ βασιλεύς.
\vs{44}Καὶ ἔδωκεν ἑαυτὸν τοῦ σῶσαι τὸν λαὸν αὐτοῦ, καὶ περιποιῆσαι ἑαυτῷ ὄνομα αἰώνιον.
\vs{45}Καὶ ἐπέδραμεν αὐτῷ θράσει εἰς μέσον τῆς φάλαγγος, καὶ ἐθανάτου δεξιὰ καὶ εὐώνυμα καὶ ἐσχίζοντο ἀπʼ αὐτοῦ ἔνθα καὶ ἔνθα.
\vs{46}Καὶ εἰσέδυ ὑπὸ τὸν ἐλέφαντα, καὶ ὑπέθηκεν αὐτῷ, καὶ ἀνεῖλεν αὐτὸν, καὶ ἔπεσεν ἐπὶ τὴν γῆν ἐπάνω αὐτοῦ, καὶ ἀπέθανεν ἐκεῖ.
\vs{47}Καὶ ἴδον τὴν ἰσχὺν τῆς βασιλείας καὶ τὸ ὅρμημα τῶν δυνάμεων, καὶ ἐξέκλιναν ἀπʼ αὐτῶν.

\vs{48}Οἱ δὲ ἐκ τῆς παρεμβολῆς τοῦ βασιλέως ἀνέβαινον εἰς συνάντησιν αὐτῶν εἰς Ἱερουσαλήμ· καὶ παρενέβαλεν ὁ βασιλεὺς εἰς τὴν Ἰουδαίαν καὶ εἰς τὸ ὄρος Σιὼν,
\vs{49}καὶ ἐποίησεν εἰρήνην μετὰ τῶν ἐκ Βαιθσούρων· καὶ ἐξῆλθον ἐκ τῆς πόλεως, ὅτι οὐκ ἦν αὐτοῖς ἐκεῖ διατροφὴ τοῦ συγκεκλεῖσθαι ἐν αὐτῇ, ὅτι σάββατον ἦν τῇ γῇ.

\vs{50}Καὶ κατελάβετο βασιλεὺς τὴν Βαιθσούραν, καὶ ἀπέταξεν ἐκεῖ φρουρὰν τηρεῖν αὐτὴν,
\vs{51}καὶ παρενέβαλεν ἐπὶ τὸ ἁγίασμα ἡμέρας πολλὰς, καὶ ἔστησεν ἐκεῖ βελοστάσεις καὶ μηχανὰς καὶ πυρόβολα καὶ λιθόβολα καὶ σκορπίδια εἰς τὸ βάλλεσθαι βέλη, καὶ σφενδόνας.
\vs{52}Καὶ ἐποίησαν καὶ αὐτοὶ μηχανὰς πρὸς τὰς μηχανὰς αὐτῶν, καὶ ἐπολέμησαν ἡμέρας πολλάς.
\vs{53}Βρώματα δὲ οὐκ ἦν ἐν τοῖς ἀγγείοις, διὰ τὸ ἕβδομον ἔτος εἶναι, καὶ οἱ ἀνασωζόμενοι εἰς τὴν Ἰουδαίαν ἀπὸ τῶν ἐθνῶν κατέφαγον τὸ ὑπόλειμμα τῆς παραθέσεως.
\vs{54}Καὶ ὑπελείφθησαν ἐν τοῖς ἁγίοις ἄνδρες ὀλίγοι, ὅτι κατεκράτησεν αὐτῶν ὁ λιμός· καὶ ἐσκορπίσθησαν ἕκαστος εἰς τὸν τόπον αὐτοῦ.

\vs{55}Καὶ ἤκουσε Λυσίας, ὅτι Φίλιππος, ὃν κατέστησεν ὁ βασιλεὺς Ἀντίοχος ἔτι ζῶν, ἐκθρέψαι Ἀντίοχον τὸν υἱὸν αὐτοῦ εἰς τὸ βασιλεῦσαι αὐτὸν,
\vs{56}ἀπέστρεψεν ἀπὸ τῆς Περσίδος καὶ Μηδείας, καὶ αἱ δυνάμεις αἱ πορευθεῖσαι τοῦ βασιλέως μετʼ αὐτοῦ, καὶ ὅτι ζητεῖ παραλαβεῖν τὰ πράγματα.
\vs{57}Καὶ κατέσπευσε τοῦ ἀπελθεῖν, καὶ εἰπεῖν πρὸς τὸν βασιλέα καὶ τοὺς ἡγεμόνας τῆς δυνάμεως καὶ τοὺς ἄνδρας, ἐκλείπομεν καθʼ ἡμέραν, καὶ ἡ τροφὴ ἡμῖν ὀλίγη, καὶ ὁ τόπος οὗ παρεμβάλλομεν ἐστιν ὀχυρὸς, καὶ ἐπίκειται ἡμῖν τὰ τῆς βασιλείας.
\vs{58}Νῦν οὖν δῶμεν δεξιὰν τοῖς ἀνθρώποις τούτοις, καὶ ποιήσωμεν μετʼ αὐτῶν εἰρήνην καὶ μετὰ παντὸς ἔθνους αὐτῶν,
\vs{59}καὶ στήσωμεν αὐτοῖς τοῦ πορεύεσθαι τοῖς νομίμοις αὐτῶν, ὡς τοπρότερον· χάριν γὰρ τῶν νομίμων αὐτῶν ὧν διεσκεδάσαμεν, ὠργίσθησαν, καὶ ἐποίησαν ταῦτα πάντα.

\vs{60}Καὶ ἤρεσεν ὁ λόγος ἐναντίον τοῦ βασιλέως καὶ τῶν ἀρχόντων, καὶ ἀπέστειλε πρὸς αὐτοὺς εἰρηνεῦσαι, καὶ ἐπεδέξαντο.
\vs{61}Καὶ ὤμοσεν αὐτοῖς ὁ βασιλεὺς καὶ οἱ ἄρχοντες· ἐπὶ τούτοις ἐξῆλθον ἐκ τοῦ ὀχυρώματος.
\vs{62}Καὶ εἰσῆλθεν ὁ βασιλεὺς εἰς τὸ ὄρος Σιὼν, καὶ εἶδε τὸ ὀχύρωμα τοῦ τόπου· καὶ ἠθέτησε τὸν ὁρκισμὸν ὃν ὤμοσε, καὶ ἐνετείλατο καθελεῖν τὸ τεῖχος κυκλόθεν.
\vs{63}Καὶ ἀπῇρε κατὰ σπουδὴν, καὶ ἀπέστρεψεν εἰς Ἀντιόχειαν, καὶ εὗρε Φίλιππον κυριεύοντα τῆς πόλεως, καὶ ἐπολέμησε πρὸς αὐτὸν, καὶ κατελάβετο τὴν πόλιν βίᾳ.

\ch{7}
Ἔτους ἑνὸς καὶ πεντηκοστοῦ καὶ ἐκατοστοῦ ἐξῆλθε Δημήτριος ὁ τοῦ Σελεύκου ἐκ Ῥώμης, καὶ ἀνέβη σὺν ἀνδράσιν ὀλίγοις εἰς πόλιν παραθαλασσίαν, καὶ ἐβασίλευσεν ἐκεῖ.

\vs{2}Καὶ ἐγένετο ὡς εἰσεπορεύετο εἰς οἶκον βασιλείας πατέρων αὐτοῦ, συνέλαβον αἱ δυνάμεις τὸν Ἀντίοχον καὶ τὸν Λυσίαν ἄγειν αὐτοὺς αὐτῷ.
\vs{3}Καὶ ἐγνώσθη αὐτῷ τὸ πρᾶγμα, καὶ εἶπε, μή μοι δείξητε τὰ πρόσωπα αὐτῶν.
\vs{4}Καὶ ἀπέκτειναν αὐτοὺς αἱ δυνάμεις, καὶ ἐκάθισε Δημήτριος ἐπὶ θρόνου βασιλείας αὐτοῦ.
\vs{5}Καὶ ἦλθον πρὸς αὐτὸν πάντες ἄνδρες ἄνομοι καὶ ἀσεβεῖς ἐξ Ἰσραήλ, καὶ Ἄλκιμος ἡγεῖτο αὐτῶν, βουλόμενος ἱερατεύειν.
\vs{6}Καὶ κατηγόρησαν τοῦ λαοῦ πρὸς τὸν βασιλέα, λέγοντες, ἀπώλεσεν Ἰούδας καὶ οἱ ἀδελφοὶ αὐτοῦ τοὺς φίλους σου, καὶ ἡμᾶς ἐσκόρπισαν ἀπὸ τῆς γῆς ἡμῶν.
\vs{7}Νῦν οὖν ἀπόστειλον ἄνδρα ᾧ πιστεύεις, καὶ πορευθεὶς ἰδέτω τὴν ἐξολόθρευσιν πᾶσαν ἣν ἐποίησεν ἡμῖν καὶ τῇ χώρᾳ τοῦ βασιλέως, καὶ κολασάτω αὐτοὺς καὶ πάντας τοὺς ἐπιβοηθοῦντας αὐτοῖς.

\vs{8}Καὶ ἐπέλεξεν ὁ βασιλεὺς τὸν Βακχίδην τῶν φίλων τοῦ βασιλέως, κυριεύοντα ἐν τῷ πέραν τοῦ ποταμοῦ, καὶ μέγαν ἐν τῇ βασιλείᾳ, καὶ πιστὸν τῷ βασιλεῖ.
\vs{9}Καὶ ἀπέστειλεν αὐτὸν καὶ Ἄλκιμον τὸν ἀσεβῆ, καὶ ἔστησεν αὐτῷ τὴν ἱερωσύνην, καὶ ἐνετείλατο αὐτῷ ποιῆσαι τὴν ἐκδίκησιν ἐν τοῖς υἱοῖς Ἰσραήλ.
\vs{10}Καὶ ἀπῇραν, καὶ ἦλθον μετὰ δυνάμεως πολλῆς εἰς γῆν Ἰούδα· καὶ ἀπέστειλεν ἀγγέλους πρὸς Ἰούδαν, καὶ τοὺς ἀδελφοὺς αὐτοῦ, λόγοις εἰρηνικοῖς μετὰ δόλου.
\vs{11}Καὶ οὐ προσέσχον τοῖς λόγοις αὐτῶν, ἴδον γὰρ ὅτι ἦλθον μετὰ δυνάμεως πολλῆς.

\vs{12}Καὶ ἐπισυνήχθησαν πρὸς Ἄλκιμον καὶ Βακχίδην συναγωγὴ γραμματέων ἐκζητῆσαι δίκαια.
\vs{13}Καὶ πρῶτοι οἱ Ἀσιδαῖοι ἦσαν ἐν υἱοῖς Ἰσραὴλ, καὶ ἐπεζήτουν παρὰ αὐτῶν εἰρήνην·
\vs{14}Εἶπαν γὰρ, ἄνθρωπος ἱερεὺς ἐκ τοῦ σπέρματος Ἀαρὼν ἦλθεν ἐν ταῖς δυνάμεσι, καὶ οὐκ ἀδικήσει ἡμᾶς.
\vs{15}Καὶ ἐλάλησε μετʼ αὐτῶν λόγους εἰρηνικοὺς, καὶ ὤμοσεν αὐτοῖς, λέγων, οὐκ ἐκζητήσομεν ὑμῖν κακὸν, καὶ τοῦ φίλοις ὑμῶν.
\vs{16}Καὶ ἐνεπίστευσαν αὐτῷ· καὶ συνέλαβεν ἐξ αὐτῶν ἑξήκοντα ἄνδρας, καὶ ἀπέκτεινεν αὐτοὺς ἐν ἡμέρᾳ μιᾷ, κατὰ τὸν λόγον ὃν ἔγραψε,
\vs{17}σάρκας ὁσίων σου καὶ αἵματα αὐτῶν ἐξέχεαν κύκλῳ Ἱερουσαλὴμ, καὶ οὐκ ἦν αὐτοῖς ὁ θάπτων.
\vs{18}Καὶ ἐπέπεσεν αὐτῶν ὁ φόβος καὶ ὁ τρόμος ἐπὶ πάντα τὸν λαὸν, ὅτι εἶπαν, οὐκ ἔστιν ἐν αὐτοῖς ἀλήθεια καὶ κρίσις· παρέβησαν γὰρ τὴν στάσιν καὶ τὸν ὅρκον ὃν ὤμοσαν.

\vs{19}Καὶ ἀπῇρε Βακχίδης ἀπὸ Ἱερουσαλὴμ, καὶ παρενέβαλεν ἐν Βηζὲθ, καὶ ἀπέστειλε καὶ συνέλαβε πολλοὺς ἀπὸ τῶν ἀπʼ αὐτοῦ αὐτομολησάντων ἀνδρῶν, καί τινας τοῦ λαοῦ, καὶ ἔθυσεν αὐτοὺς εἰς τὸ φρέαρ τὸ μέγα.
\vs{20}Καὶ κατέστησε τὴν χώραν τῷ Ἀλκίμῳ, καὶ ἀφῆκε μετʼ αὐτοῦ δύναμιν τοῦ βοηθεῖν αὐτῷ· καὶ ἀπῆλθε Βακχίδης πρὸς τὸν βασιλέα.
\vs{21}Καὶ ἠγωνίσατο Ἄλκιμος περὶ τῆς ἀρχιερωσύνης.
\vs{22}Καὶ συνήχθησαν πρὸς αὐτὸν πάντες οἱ ταράσσοντες τὸν λαὸν αὐτῶν, καὶ κατεκράτησαν γῆν Ἰούδα, καὶ ἐποίησαν πληγὴν μεγάλην ἐν Ἰσραήλ.

\vs{23}Καὶ εἶδεν Ἰούδας πᾶσαν τὴν κακίαν ἣν ἐποίησεν Ἄλκιμος καὶ οἱ μετʼ αὐτοῦ ἐν υἱοῖς Ἰσραὴλ ὑπὲρ τὰ ἔθνη·
\vs{24}καὶ ἐξῆλθεν εἰς πάντα τὰ ὅρια τῆς Ἰουδαίας κυκλόθεν, καὶ ἐποίησεν ἐκδίκησιν ἐν τοῖς ἀνδράσι τοῖς αὐτομολήσασι, καὶ ἀνεστάλησαν τοῦ πορεύεσθαι εἰς τὴν χώραν.

\vs{25}Ὡς δὲ εἶδεν Ἄλκιμος ὅτι ἐνίσχυσεν Ἰούδας καὶ οἱ μετʼ αὐτοῦ, καὶ ἔγνω ὅτι οὐ δύναται ὑποστῆναι αὐτούς, καὶ ἐπέστρεψε πρὸς τὸν βασιλέα, καὶ κατηγόρησεν αὐτῶν πονηρά.

\vs{26}Καὶ ἀπέστειλεν ὁ βασιλεὺς Νικάνορα, ἕνα τῶν ἀρχόντων αὐτοῦ τῶν ἐνδόξων, καὶ μισοῦντα καὶ ἐχθραίνοντα τῷ Ἰσραήλ, καὶ ἐνετείλατο αὐτῷ ἐξᾶραι τὸν λαόν.
\vs{27}Καὶ ἦλθε Νικάνωρ εἰς Ἰερουσαλὴμ δυνάμει πολλῇ, καὶ ἀπέστειλε πρὸς Ἰούδαν καὶ τοὺς ἀδελφοὺς αὐτοῦ μετὰ δόλου λόγοις εἰρηνικοῖς, λέγων,
\vs{28}μὴ ἔστω μάχη ἀναμέσον ἐμοῦ καὶ ὑμῶν· ἥξω ἐν ἀνδράσιν ὀλίγοις, ἵνα ὑμῶν ἴδω τὰ πρόσωπα μετʼ εἰρήνης.
\vs{29}Καὶ ἦλθε πρὸς Ἰούδαν, καὶ ἠσπάσαντο ἀλλήλους εἰρηνικῶς· καὶ οἱ πολέμιοι ἦσαν ἕτοιμοι ἐξαρπάσαι τὸν Ἰούδαν.
\vs{30}Καὶ ἐγνώσθη ὁ λόγος τῷ Ἰούδᾳ, ὅτι μετὰ δόλου ἦλθεν ἐπʼ αὐτόν· καὶ ἐπτοήθη ἀπʼ αὐτοῦ, καὶ οὐκ ἐβουλήθη ἔτι ἰδεῖν τὸ πρόσωπον αὐτοῦ.

\vs{31}Καὶ ἔγνω Νικάνωρ ὅτι ἀπεκαλύφθη ἡ βουλὴ αὐτοῦ, καὶ ἐξῆλθεν εἰς συνάντησιν τῷ Ἰούδᾳ ἐν πολέμῳ κατὰ Χαφαρσαλαμά.
\vs{32}Καὶ ἔπεσον τῶν παρὰ Νικάνορος ὡσεὶ πεντακισχίλιοι ἄνδρες, καὶ ἔφυγον εἰς τὴν πόλιν Δαυίδ.

\vs{33}Καὶ μετὰ τοὺς λόγους τούτους ἀνέβη Νικάνωρ εἰς τὸ ὄρος Σιών· καὶ ἐξῆλθον ἀπὸ τῶν ἱερέων ἐκ τῶν ἁγίων καὶ ἀπὸ τῶν πρεσβυτέρων τοῦ λαοῦ ἀσπάσασθαι αὐτὸν εἰρηνικῶς, καὶ δεῖξαι αὐτῷ τὴν ὁλοκαύτωσιν τὴν προσφερομένην ὑπὲρ τοῦ βασιλέως.
\vs{34}Καὶ ἐμυκτήρισεν αὐτοὺς, καὶ κατεγέλασεν αὐτῶν, καὶ ἐμίανεν αὐτούς, καὶ ἐλάλησεν ὑπερηφάνως.
\vs{35}Καὶ ὤμοσε μετὰ θυμοῦ, λέγων, ἐὰν μὴ παραδοθῇ Ἰούδας καὶ ἡ παρεμβολὴ αὐτοῦ εἰς χεῖράς μου τὸ νῦν, καὶ ἔσται ἐὰν ἐπιστρέψω ἐν εἰρήνῃ, ἐμπυριῶ τὸν οἶκον τοῦτον· καὶ ἐξῆλθε μετὰ θυμοῦ μεγάλου.

\vs{36}Καὶ εἰσῆλθον οἱ ἱερεῖς, καὶ ἔστησαν κατὰ πρόσωπον τοῦ θυσιαστηρίου καὶ τοῦ ναοῦ, καὶ ἔκλαυσαν, καὶ εἶπον,
\vs{37}σὺ, Κύριε, ἐξελέξω τὸν οἶκον τοῦτον ἐπικληθῆναι τὸ ὄνομά σου ἐπʼ αὐτῷ, εἶναι οἶκον προσευχῆς καὶ δεήσεως τῷ λαῷ σου.
\vs{38}Ποίησον ἐκδίκησιν ἐν τῷ ἀνθρώπῳ τούτῳ καὶ ἐν τῇ παρεμβολῇ αὐτοῦ, καὶ πεσέτωσαν ἐν ῥομφαίᾳ· μνήσθητι τῶν δυσφημιῶν αὐτῶν, καὶ μὴ δῷς αὐτοῖς μονήν.

\vs{39}Καὶ ἐξῆλθε Νικάνωρ ἐξ Ἱερουσαλὴμ, καὶ παρενέβαλεν ἐν Βαιθωρὼν, καὶ συνήντησεν αὐτῷ δύναμις Συρίας.
\vs{40}Καὶ Ἰούδας παρενέβαλεν ἐν Ἀδασὰ ἐν τρισχιλίοις ἀνδράσι· καὶ προσηύξατο Ἰούδας, καὶ εἶπεν,
\vs{41}οἱ παρὰ τοῦ βασιλέως Ἀσσυρίων ὅτε ἐδυσφήμησαν, ἐξῆλθεν ὁ ἄγγελός σου, Κύριε, καὶ ἐπάταξεν ἐν αὐτοῖς ἑκατὸν ὀγδοηκονταπέντε χιλιάδας.
\vs{42}Οὕτω σύντριψον τὴν παρεμβολὴν ταύτην ἐνώπιον ἡμῶν σήμερον, καὶ γνώτωσαν οἱ ἐπίλοιποι, ὅτι κακῶς ἐλάλησαν ἐπὶ τὰ ἅγιά σου, καὶ κρῖνον αὐτὸν κατὰ τὴν κακίαν αὐτοῦ.

\vs{43}Καὶ συνῆψαν αἱ παρεμβολαὶ εἰς πόλεμον τῇ τρισκαιδεκάτῃ τοῦ μηνὸς Ἄδαρ, καὶ συνετρίβη ἡ παρεμβολὴ Νικάνορος, καὶ ἔπεσεν αὐτὸς πρῶτος ἐν τῷ πολέμῳ.

\vs{44}Ὡς δὲ εἶδεν ἡ παρεμβολὴ αὐτοῦ ὅτι ἔπεσε Νικάνωρ, ῥίψαντες τὰ ὅπλα αὐτῶν ἔφυγον.
\vs{45}Καὶ κατεδίωκον αὐτοὺς ὁδὸν ἡμέρας μιᾶς ἀπὸ Ἀδασὰ ἕως τοῦ ἐλθεῖν εἰς Γάζηρα, καὶ ἐσάλπισαν ὀπίσω αὐτῶν ταῖς σάλπιγξι τῶν σημασιῶν.
\vs{46}Καὶ ἐξῆλθον ἐκ πασῶν τῶν κωμῶν τῆς Ἰουδαίας κυκλόθεν, καὶ ὑπερεκέρων αὐτοὺς, καὶ ἀνέστρεφον οὗτοι πρὸς τούτους· καὶ ἔπεσον πάντες ῥομφαίᾳ, καὶ οὐ κατελείφθη ἐξ αὐτῶν οὐδὲ εἷς.

\vs{47}Καὶ ἔλαβον τὰ σκῦλα καὶ τὴν προνομὴν, καὶ τὴν κεφαλὴν Νικάνορος ἀφεῖλον, καὶ τὴν δεξιὰν αὐτοῦ ἣν ἐξέτεινεν ὑπερηφάνως, καὶ ἤνεγκαν, καὶ ἐξέτειναν παρὰ τὴν Ἱερουσαλήμ.
\vs{48}Καὶ εὐφράνθη δ λαὸς σφόδρα, καὶ ἤγαγον τὴν ἡμέραν ἐκείνην ἡμέραν εὐφροσύνης μεγάλης.
\vs{49}Καὶ ἔστησαν τοῦ ἄγειν κατὰ ἐνιαυτὸν τὴν ἡμέραν ταύτην τὴν τρισκαιδεκάτην τοῦ Ἄδαρ.
\vs{50}Καὶ ἡσύχασεν ἡ γῆ Ἰούδα ἡμέρας ὀλίγας.

\ch{8}
Καὶ ἤκουσεν Ἰούδας τὸ ὄνομα τῶν Ῥωμαίων, ὅτι εἰσὶ δυνατοὶ ἰσχύϊ· καὶ αὐτοὶ εὐδοκοῦσιν ἐν πᾶσι τοῖς προστιθεμένοις αὐτοῖς, καὶ ὅσοι ἂν προσέλθωσιν αὐτοῖς, ἱστῶσιν αὐτοῖς φιλίαν, καὶ ὅτι εἰσὶ δυνατοὶ ἰσχύϊ·
\vs{2}καὶ διηγήσαντο αὐτῷ τοὺς πολέμους αὐτῶν, καὶ τὰς ἀνδραγαθίας ἃς ποιοῦσιν ἐν τοῖς Γαλάταις, καὶ ὅτι κατεκράτησαν αὐτῶν καὶ ἤγαγον αὐτοὺς ὑπὸ φόρον,
\vs{3}καὶ ὅσα ἐποίησαν ἐν χώρᾳ Ἱσπανιας, του κατακρατῆσαι τῶν μετάλλων τοῦ ἀργυρίου καὶ τοῦ χρυσίου τοῦ ἐκεῖ·
\vs{4}καὶ κατεκράτησαν τοῦ τόπου παντὸς τῇ βουλῇ αὐτῶν καὶ τῇ μακροθυμίᾳ, καὶ ὁ τόπος ἦν μακρὰν ἀπέχων ἀπʼ αὐτῶν σφόδρα· καὶ τῶν βασιλέων τῶν ἐπελθόντων ἐπʼ αὐτοὺς ἀπʼ ἄκρου τῆς γῆς ἕως συνέτριψαν αὐτοὺς, καὶ ἐπάταξαν ἐν αὐτοῖς πληγὴν μεγάλην, καὶ οἱ ἐπίλοιποι διδόασιν αὐτοῖς φόρον κατʼ ἐνιαυτόν·

\vs{5}Καὶ τὸν Φίλιππον καὶ τὸν Περσέα Κιτιέων βασιλέα, καὶ τοὺς ἐπῃρμένους ἐπʼ αὐτοὺς, συνέτριψαν αὐτοὺς ἐν πολέμῳ, καὶ κατεκράτησαν αὐτῶν·
\vs{6}καὶ Ἀντίοχον τὸν μέγαν βασιλέα τῆς Ἀσίας, τὸν πορευθέντα ἐπʼ αὐτοὺς εἰς πόλεμον ἔχοντα ἑκατὸν εἴκοσι ἐλέφαντας καὶ ἵππον καὶ ἅρματα καὶ δύναμιν πολλὴν σφόδρα, καὶ συνετρίβη ἀπʼ αὐτῶν·
\vs{7}καὶ ἔλαβον αὐτὸν ζῶντα, καὶ ἔστησαν αὐτοῖς διδόναι αὐτόν τε καὶ τοὺς βασιλεύοντας μετʼ αὐτὸν φόρον μέγαν, διδόναι ὅμηρα καὶ διαστολήν,
\vs{8}καὶ χώραν τὴν Ἰνδικὴν, καὶ Μήδειαν, καὶ Αυδίαν, καὶ ἀπὸ τῶν καλλίστων χωρῶν αὐτῶν, καὶ λαβόντες αὐτὰς παρʼ αὐτοῦ ἔδωκαν αὐτὰς Εὐμένει τῷ βασιλεῖ.

\vs{9}Καὶ ὅτι οἱ ἐκ τῆς Ἑλλάδος ἐβουλεύσαντο ἐλθεῖν καὶ ἐξᾶραι αὐτούς,
\vs{10}καὶ ἐγνώσθη ὁ λόγος αὐτοῖς, καὶ ἀπέστειλαν ἐπʼ αὐτοὺς στρατηγὸν ἕνα, καὶ ἐπολέμησαν πρὸς αὐτοὺς, καὶ ἔπεσον ἐξ αὐτῶν τραυματίαι πολλοὶ, καὶ ᾐχμαλώτευσαν τὰς γυναῖκας αὐτῶν καὶ τὰ τέκνα αὐτῶν, καὶ προενόμευσαν αὐτοὺς, καὶ κατεκράτησαν τῆς γῆς αὐτῶν, καὶ καθεῖλον τὰ ὀχυρώματα αὐτῶν, καὶ κατεδουλώσαντο αὐτοὺς ἕως τῆς ἡμέρας ταύτης.

\vs{11}Καὶ τὰς ἐπιλοίπους βασιλείας, καὶ τὰς νήσους, ὅσοι ποτὲ ἀντέστησαν αὐτοῖς, κατέφθειραν, καὶ ἐδούλωσαν αὐτούς· μετὰ δὲ τῶν φίλων αὐτῶν καὶ τῶν ἐπαναπαυομένων αὐτοῖς συνετήρησαν φιλίαν,
\vs{12}καὶ κατεκράτησαν τῶν βασιλειῶν τῶν ἐγγὺς καὶ τῶν μακρὰν, καὶ ὅσοι ἤκουον τὸ ὄνομα αὐτῶν ἐφοβοῦντο ἀπʼ αὐτῶν·
\vs{13}ὅσοις δʼ ἂν βούλωνται βοηθεῖν καὶ βασιλεύειν, βασιλεύουσιν· οὓς δʼ ἂν βούλωνται, μεθιστῶσι, καὶ ὑψώθησαν σφόδρα·
\vs{14}καὶ ἐν πᾶσι τούτοις οὐκ ἐπέθετο οὐδεὶς αὐτῶν διάδημα, καὶ οὐ περιεβάλοντο πορφύραν, ὥστε ἁδρυνθῆναι ἐν αὐτῇ.
\vs{15}Καὶ βουλευτήριον ἐποίησαν ἑαυτοῖς, καὶ καθʼ ἡμέραν ἐβουλεύοντο τριακόσιοι καὶ εἴκοσι βουλευόμενοι διαπαντὸς περὶ τοῦ πλήθους, τοῦ εὐκοσμεῖν αὐτούς·
\vs{16}καὶ πιστεύουσιν ἑνὶ ἀνθρώπῳ τὴν ἀρχὴν αὐτῶν κατʼ ἐνιαυτὸν, καὶ κυριεύειν πάσης τῆς γῆς αὐτῶν, καὶ πάντες ἀκούουσι τοῦ ἑνὸς, καὶ οὐκ ἔστι φθόνος οὐδὲ ζῆλος ἐν αὐτοῖς·

\vs{17}Καὶ ἐπέλεξεν Ἰούδας τὸν Εὐπόλεμον υἱὸν Ἰωάννου τοῦ Ἀκκὼς, καὶ Ἰάσονα υἱὸν Ἐλεαζάρου, καὶ ἀπέστειλεν αὐτοὺς εἰς Ῥώμην, στῆσαι αὐτοῖς φιλίαν καὶ συμμαχίαν,
\vs{18}καὶ τοῦ ἆραι τὸν ζυγὸν ἀπʼ αὐτῶν, ὅτι ἴδον τὴν βασιλείαν τῶν Ἑλλήνων καταδουλουμένους τὸν Ἰσραὴλ δουλείαν.

\vs{19}Καὶ ἐπορεύθησαν εἰς Ῥώμην, καὶ ἡ ὁδὸς πολλὴ σφόδρα, καὶ εἰσῆλθον εἰς τὸ βουλευτήριον, καὶ ἀπεκρίθησαν καὶ εἶπον,
\vs{20}Ἰούδας ὁ Μακκαβαῖος καὶ οἱ ἀδελφοὶ αὐτοῦ καὶ τὸ πλῆθος τῶν Ἰουδαίων ἀπέστειλαν ἡμᾶς πρὸς ὑμᾶς, στῆσαι μεθʼ ὑμῶν συμμαχίαν καὶ εἰρήνην, καὶ γραφῆναι ἡμᾶς συμμάχους καὶ φίλους ὑμῶν.
\vs{21}Καὶ ἤρεσεν ὁ λόγος ἐνώπιον αὐτῶν.

\vs{22}Καὶ τοῦτο τὸ ἀντίγραφον τῆς ἐπιστολῆς ἧς ἀντέγραψεν ἐπὶ δέλτοῖς χαλκαῖς, καὶ ἀπέστειλεν εἰς Ἱερουσαλὴμ εἶναι παρʼ αὐτοῖς ἐκεῖ μνημόσυνον εἰρήνης καὶ συμμαχίας·
\vs{23}καλῶς γένοιτο Ῥωμαίοις καὶ τῷ ἔθνει Ἰουδαίων ἐν τῇ θαλάσσῃ καὶ ἐπὶ τῆς ξηρᾶς εἰς τὸν αἰῶνα, καὶ ῥομφαία καὶ ἐχθρὸς μακρυνθείη ἀπʼ αὐτῶν.

\vs{24}Ἐὰν δὲ ἐνστῇ πόλεμος ἐν Ῥώμῃ προτέρᾳ ἢ πᾶσι τοῖς συμμάχοις αὐτῶν ἐν πάσῃ κυρείᾳ αὐτῶν,
\vs{25}συμμαχήσει τὸ ἔθνος τῶν Ἰουδαίων, ὡς ἂν ὁ καιρὸς ὑπογραφῇ αὐτοῖς, καρδίᾳ πλήρει.
\vs{26}Καὶ τοῖς πολεμοῦσιν οὐ δώσουσιν οὐδὲ ἐπαρκέσουσι σῖτον, ὅπλα, ἀργύριον, πλοῖα, ὡς ἔδοξε Ῥωμαίοις· καὶ φυλάξονται τὰ φυλάγματα αὐτῶν οὐθὲν λαβόντες·
\vs{27}κατὰ τὰ αὐτὰ δὲ ἐὰν ἔθνει Ἰουδαίων συμβῇ προτέροις πόλεμος, συμμαχήσουσιν οἱ Ῥωμαῖοι ἐκ ψυχῆς, ὡς ἂν αὐτοῖς ὁ καιρὸς ὑπογραφῇ.
\vs{28}Καὶ τοῖς συμμαχοῦσιν οὐ δοθήσεται σῖτος, ὅπλα, ἀργύριον, πλοῖα, ὡς ἔδοξε Ῥώμῃ· καὶ φυλάξονται τὰ φυλάγματα αὐτῶν, καὶ οὐ μετὰ δόλου.

\vs{29}Κατὰ τοὺς λόγους τούτους ἔστησαν Ῥωμαῖοι τῷ δήμῳ τῶν Ἰουδαίων.
\vs{30}Ἐὰν δὲ μετὰ τοὺς λόγους τούτους βουλεύσωνται οὗτοι καὶ οὗτοι προσθεῖναι ἢ ἀφελεῖν, ποιήσονται ἐξ αἱρέσεως αὐτῶν, καὶ ὃ ἐὰν προσθῶσιν ἢ ἀφέλωσιν, ἔσται κύρια.

\vs{31}Καὶ περὶ τῶν κακῶν ὧν ὁ βασιλεὺς Δημήτριος συντελεῖται εἰς αὐτούς, ἐγράψαμεν αὐτῷ, λέγοντες, διατί ἐβάρυνας τὸν ζυγόν σου ἐπὶ τοὺς φίλους ἡμῶν τοὺς συμμάχους Ἰουδαίους;
\vs{32}Ἐὰν οὖν ἔτι ἐντύχωσιν κατὰ σοῦ, ποιήσομεν αὐτοῖς τὴν κρίσιν, καὶ πολεμήσομέν σε διὰ τῆς θαλάσσης καὶ διὰ τῆς ξηρᾶς.

\ch{9}
Καὶ ἤκουσε Δημήτριος ὅτι ἔπεσε Νικάνωρ καὶ αἱ δυνάμεις αὐτοῦ ἐν πολέμῳ, καὶ προσέθετο τὸν Βακχίδην καὶ τὸν Ἄλκιμον ἐκ δευτέρου ἀποστεῖλαι εἰς γῆν Ἰούδα, καὶ τὸ δεξιὸν κέρας μετʼ αὐτῶν.
\vs{2}Καὶ ἐπορεύθησαν ὁδὸν τὴν εἰς Γάλγαλα, καὶ παρενέβαλον ἐπὶ Μαισαλὼθ τὴν ἐν Ἀρβήλοις, καὶ προκατελάβοντο αὐτὴν, καὶ ἀπώλεσαν ψυχὰς ἀνθρώπων πολλάς.
\vs{3}Καὶ τοῦ μηνὸς τοῦ πρώτου ἔτους τοῦ δευτέρου καὶ πεντηκοστοῦ καὶ ἑκατοστοῦ παρενέβαλον ἐπὶ Ἱερουσαλήμ.
\vs{4}Καὶ ἀπῇραν καὶ ἐπορεύθησαν εἰς Βερέαν ἐν εἴκοσι χιλιάσιν ἀνδρῶν καὶ δισχιλίᾳ ἵππῳ.

\vs{5}Καὶ Ἰούδας ἦν παρεμβεβληκὼς ἐν Ἐλεασὰ, καὶ τρισχίλιοι ἄνδρες ἐκλεκτοὶ μετʼ αὐτοῦ.
\vs{6}Καὶ ἴδον τὸ πλῆθος τῶν δυνάμεων ὅτι πολλοί εἰσι, καὶ ἐφοβήθησαν σφόδρα· καὶ ἐξεῤῥύησαν πολλοὶ ἀπὸ τῆς παρεμβολῆς, οὐ κατελείφθησαν ἐξ αὐτῶν ἀλλʼ ἢ ὀκτακόσιοι ἄνδρες.

\vs{7}Καὶ εἶδεν Ἰούδας ὅτι ἀπεῤῥύη ἡ παρεμβολὴ αὐτοῦ, καὶ ὁ πόλεμος ἔθλιβεν αὐτόν· καὶ συνετρίβη τῇ καρδίᾳ, ὅτι οὐκ εἶχε καιρὸν συναγαγεῖν αὐτούς.
\vs{8}Καὶ ἐξελύθη, καὶ εἶπε τοῖς καταλειφθεῖσιν, ἀναστῶμεν καὶ ἀναβῶμεν ἐπὶ τοὺς ὑπεναντίους ἡμῶν, ἐὰν ἄρα δυνώμεθα πολεμῆσαι αὐτούς.
\vs{9}Καὶ ἀπέστρεψαν αὐτὸν, λέγοντες, οὐ μὴ δυνώμεθα, ἀλλʼ ἢ σώζωμεν τὰς ἑαυτῶν ψυχὰς τὸ νῦν, καὶ ἐπιστρέψωμεν μετὰ τῶν ἀδελφῶν ἡμῶν, καὶ πολεμήσωμεν πρὸς αὐτοὺς, ἡμεῖς δὲ ὀλίγοι.

\vs{10}Καὶ εἶπεν Ἰούδας, μή μοι γένοιτο ποιῆσαι τὸ πρᾶγμα τοῦτο, φυγεῖν ἀπʼ αὐτῶν, καὶ εἰ ἤγγικεν ὁ καιρὸς ἡμῶν, καὶ ἀποθάνωμεν ἐν ἀνδρείᾳ χάριν τῶν ἀδελφῶν ἡμῶν, καὶ μὴ καταλίπωμεν αἰτίαν τῇ δόξῃ ἡμῶν.
\vs{11}Καὶ ἀπῇρεν ἡ δύναμις ἀπὸ τῆς παρεμβολῆς, καὶ ἔστησαν εἰς συνάντησιν αὐτοῖς, καὶ ἐμερίσθη ἡ ἵππος εἰς δύο μέρη, καὶ οἱ σφενδονηταὶ καὶ οἱ τοξόται προεπορεύοντο τῆς δυνάμεως, καὶ οἱ πρωταγωνισταὶ πάντες οἱ δυνατοί.

\vs{12}Βακχίδης δὲ ἦν ἐν τῷ δεξιῷ κέρατι, καὶ ἤγγισεν ἡ φάλαγξ ἐκ τῶν δύο μερῶν, καὶ ἐφώνουν ταῖς σάλπιγξι. Καὶ ἐσάλπισαν οἱ παρὰ Ἰούδα καὶ αὐτοὶ ταῖς σάλπιγξι,
\vs{13}καὶ ἐσαλεύθη ἡ γῆ ἀπὸ τῆς φωνῆς τῶν παρεμβολῶν· καὶ ἐγένετο ὁ πόλεμος συνημένος ἀπὸ πρωΐθεν ἕως ἑσπέρας.

\vs{14}Καὶ· εἶδεν Ἰούδας ὅτι Βακχίδης καὶ τὸ στερέωμα τῆς παρεμβολῆς ἐν τοῖς δεξιοῖς, καὶ συνῆλθον αὐτῷ πάντες οἱ εὔψυχοι τῇ καρδίᾳ.
\vs{15}Καὶ συνετρίβη τὸ δεξιὸν κέρας ἀπʼ αὐτῶν, καὶ ἐδίωκεν ὀπίσω αὐτῶν ἕως Ἀζώτου ὄρους.
\vs{16}Καὶ οἱ εἰς τὸ ἀριστερὸν κέρας ἴδον ὅτι συνετρίβη τὸ δεξιὸν κέρας, καὶ ἐπέστρεψαν κατὰ πόδας Ἰούδα καὶ τῶν μετʼ αὐτοῦ ἐκ τῶν ὄπισθεν.
\vs{17}Καὶ ἐβαρύνθη ὁ πόλεμος, καὶ ἔπεσον τραυματίαι πολλοὶ ἐκ τούτων καὶ ἐκ τούτων.
\vs{18}Καὶ Ἰούδας ἔπεσε, καὶ οἱ λοιποὶ ἔφυγον.

\vs{19}Καὶ ᾖραν Ἰωνάθαν καὶ Σίμων Ἰούδαν τὸν ἀδελφὸν αὐτῶν, καὶ ἔθαψαν αὐτὸν ἐν τῷ τάφῳ τῶν πατέρων αὐτοῦ ἐν Μωδεεΐμ.
\vs{20}Καὶ ἔκλαυσαν αὐτὸν, καὶ ἐκόψαντο αὐτὸν πᾶς Ἰσραὴλ κοπετὸν μέγαν, καὶ ἐπένθουν ἡμέρας πολλὰς, καὶ εἶπον,
\vs{21}πῶς ἔπεσε δυνατὸς, σώζων τὸν Ἰσραήλ;
\vs{22}Καὶ τὰ περισσὰ τῶν λόγων Ἰούδα, καὶ τῶν πολέμων, καὶ τῶν ἀνδραγαθιῶν ὧν ἐποίησε, καὶ τῆς μεγαλωσύνης αὐτῶν, οὐ κατεγράφη, πολλὰ γὰρ ἦν σφόδρα.

\vs{23}Καὶ ἐγένετο μετὰ τὴν τελευτὴν Ἰούδα, ἐξέκυψαν οἱ ἄνομοι ἐν πᾶσι τοῖς ὁρίοις Ἰσραὴλ, καὶ ἀνέτειλαν πάντες οἱ ἐργαζόμενοι τὴν ἀδικίαν.
\vs{24}Ἐν ταῖς ἡμέραις ἐκείναις ἐγενήθη λιμὸς μέγας σφόδρα, καὶ ηὐτομόλησεν ἡ χώρα μετʼ αὐτῶν.

\vs{25}Καὶ ἐξέλεξε Βακχίδης τοὺς ἀσεβεῖς ἄνδρας, καὶ κατέστησεν αὐτοὺς κυρίους τῆς χώρας.
\vs{26}Καὶ ἐξεζήτουν καὶ ἐξηρεύνων τοὺς φίλους Ἰούδα, καὶ ἦγον αὐτοὺς πρὸς Βακχίδην· καὶ ἐξεδίκει ἐν αὐτοῖς, καὶ ἐνέπαιζεν αὐτοῖς.
\vs{27}Καὶ ἐγένετο θλίψις μεγάλη ἐν τῷ Ἰσραὴλ, ἥτις οὐκ ἐγένετο ἀφʼ ἧς ἡμέρας οὐκ ὤφθη προφήτης ἐν αὐτοῖς.

\vs{28}Καὶ ἠθροίσθησαν πάντες οἱ φίλοι Ἱούδα, καὶ εἶπον τῷ Ἰωνάθαν,
\vs{29}ἀφʼ οὗ ὁ ἀδελφός σου Ἰούδας τετελεύτηκε, καὶ ἀνὴρ ὅμοιος αὐτῷ οὐκ ἔστιν ἐξελθεῖν πρὸς τοὺς ἐχθροὺς καὶ Βακχίδην, καὶ ἐν τοῖς ἐχθραίνουσι τοῦ ἔθνους ἡμῶν.
\vs{30}Νῦν οὖν σε ᾑρετισάμεθα σήμερον, τοῦ εἶναι ἀντʼ αὐτοῦ ἡμῖν εἰς ἄρχοντα καὶ ἡγούμενον, τοῦ πολεμῆσαι τὸν πόλεμον ἡμῶν.
\vs{31}Καὶ ἐπεδέξατο Ἰωνάθαν ἐν τῷ καιρῷ ἐκείνῳ τὴν ἥγησιν, καὶ ἀνέστη ἀντὶ Ἰούδα τοῦ ἀδελφοῦ αὐτοῦ.
\vs{32}Καὶ ἔγνω Βακχίδης, καὶ ἐζήτει αὐτὸν ἀποκτεῖναι.

\vs{33}Καὶ ἔγνω Ἰωνάθαν, καὶ Σίμων ὁ ἀδελφὸς αὐτοῦ, καὶ πάντες οἱ μετʼ αὐτοῦ, καὶ ἔφυγον εἰς τὴν ἔρημον Θεκωὲ, καὶ παρενέβαλον ἐπὶ τὸ ὕδωρ λάκκου Ἀσφάρ.
\vs{34}Καὶ ἔγνω Βακχίδης τῇ ἡμέρᾳ τῶν σαββάτων, καὶ ἦλθεν αὐτὸς καὶ πᾶν τὸ στράτευμα αὐτοῦ πέραν τοῦ Ἰορδάνου.
\vs{35}Καὶ ἀπέστειλεν Ἰωνάθαν τὸν ἀδελφὸν αὐτοῦ ἡγούμενον τοῦ ὄχλου, καὶ παρεκάλεσε τοὺς Ναυαταίους φίλους αὐτοῦ παραθέσθαι αὐτοῖς τὴν ἀποσκευὴν αὐτῶν τὴν πολλήν.
\vs{36}Καὶ ἐξῆλθον υἱοὶ Ἰαμβρὶ ἐκ Μηδαβὰ, καὶ συνέλαβον Ἰωάννην, καὶ πάντα ὅσα εἶχε, καὶ ἀπῆλθον ἔχοντες.

\vs{37}Μετὰ δὲ τοὺς λόγους τούτους ἀπήγγειλαν τῷ Ἰωνάθαν καὶ Σίμωνι τῷ ἀδελφῷ αὐτοῦ, ὅτι οἱ υἱοὶ Ἰαμβρὶ ποιοῦσι γάμον μέγαν, καὶ ἄγουσι τὴν νύμφην ἀπὸ Ναδαβὰθ, θυγατέρα ἑνὸς τῶν μεγιστάνων μεγάλων τῶν Χαναὰν, μετὰ παραπομπῆς μεγάλης.
\vs{38}Καὶ ἐμνήσθησαν Ἰωάννου τοῦ ἀδελφοῦ αὐτῶν, καὶ ἀνέβησαν, καὶ ἐκρύβησαν ὑπὸ τὴν σκέπην τοῦ ὄρους.
\vs{39}Καὶ ᾖραν τοὺς ὀφθαλμοὺς αὐτῶν, καὶ ἴδον, καὶ ἰδοὺ θροῦς, καὶ ἀποσκευὴ πολλὴ, καὶ ὁ νυμφίος ἐξῆλθε καὶ οἱ φίλοι αὐτοῦ καὶ οἱ ἀδελφοὶ αὐτοῦ εἰς συνάντησιν αὐτῶν μετὰ τυμπάνων, και μουσικῶν, καὶ ὅπλων πολλῶν.

\vs{40}Καὶ ἐξανέστησαν ἐπʼ αὐτοὺς ἀπὸ τοῦ ἐνέδρου οἱ περὶ τὸν Ἰωνάθαν, καὶ ἀπέκτειναν αὐτοὺς, καὶ ἔπεσον τραυματίαι πολλοὶ, καὶ οἱ ἐπίλοιποι ἔφυγον εἰς τὸ ὄρος· καὶ ἔλαβον πάντα τὰ σκῦλα αὐτῶν.
\vs{41}Καὶ μετεστράφη ὁ γάμος εἰς πένθος, καὶ ἡ φωνὴ μουσικῶν αὐτῶν εἰς θρῆνον.
\vs{42}Καὶ ἐξεδίκησαν τὴν ἐκδίκησιν αἵματος ἀδελφοῦ αὐτῶν, καὶ ἀπέστρεψαν εἰς τὸ ἕλος τοῦ Ἰορδάνου.

\vs{43}Καὶ ἤκουσε Βακχίδης, καὶ ἦλθε τῇ ἡμέρᾳ τῶν σαββάτων ἕως τῶν κρηπίδων τοῦ Ἰορδάνου ἐν δυνάμει πολλῇ.
\vs{44}Καὶ εἶπεν Ἰωνάθαν τοῖς παρʼ αὐτοῦ, ἀναστῶμεν νῦν καὶ πολεμήσωμεν ὑπὲρ τῶν ψυχῶν ἡμῶν, οὐ γὰρ ἐστι σήμερον ὡς ἐχθὲς καὶ τρίτην ἡμέραν.
\vs{45}Ἰδοὺ γὰρ ὁ πόλεμος ἐξεναντίας ἡμῶν καὶ ἐξόπισθεν ἡμῶν· τὸ δὲ ὕδωρ τοῦ Ἰορδάνου ἔνθεν καὶ ἔνθεν, καὶ ἕλος καὶ δρυμὸς, οὐκ ἔστι τόπος τοῦ ἐκκλῖναι.
\vs{46}Νῦν οὖν κεκράξατε εἰς οὐρανὸν, ὅπως διασωθῆτε ἐκ χειρὸς ἐχθρῶν ὑμῶν.
\vs{47}Καὶ συνῆψεν ὁ πόλεμος· καὶ ἐξέτεινεν Ἰωνάθαν τὴν χεῖρα αὐτοῦ πατάξαι τὸν Βακχίδην, καὶ ἐξέκλινεν ἀπʼ αὐτοῦ εἰς τὰ ὀπίσω.
\vs{48}Καὶ ἐνεπήδησεν Ἰωνάθαν καὶ οἱ μετʼ αὐτοῦ εἰς τὸν Ἰορδάνην, καὶ διεκολύμβησαν εἰς τὸ πέραν· καὶ οὐ διέβησαν ἐπʼ αὐτοὺς τὸν Ἰορδάνην.
\vs{49}Καὶ διέπεσον παρὰ Βακχίδου τῇ ἡμέρᾳ ἐκείνῃ εἰς χιλίους ἄνδρας.

\vs{50}Καὶ ἐπέστρεψεν εἰς Ἱερουσαλὴμ, καὶ ᾠκοδόμησε πόλεις ὀχυρὰς ἐν τῇ Ἰουδαίᾳ, τὸ ὀχύρωμα τὸ ἐν Ἱεριχὼ, καὶ τὴν Ἐμμαοὺμ, καὶ τὴν Βαιθωρῶν, καὶ τὴν Βαιθὴλ, καὶ τὴν Θαμναθὰ, Φαραθωνὶ, καὶ τὴν Τεφὼν ἐν τείχεσιν ὑψηλοῖς καὶ πύλαις καὶ μοχλοῖς.
\vs{51}Καὶ ἔθετο φρουρὰν ἐν αὐτοῖς τοῦ ἐχθραίνειν τῷ Ἰσραήλ.
\vs{52}Καὶ ὠχύρωσε τὴν πόλιν τὴν ἐν Βαιθσούρᾳ, καὶ τὴν Γάζαρα, καὶ τὴν ἄκραν, καὶ ἔθετο ἐν αὐταῖς δυνάμεις καὶ παραθέσεις βρωμάτων.
\vs{53}Καὶ ἔλαβε τοὺς υἱοὺς τῶν ἡγουμένων τῆς χώρας ὅμηρα, καὶ ἔθετο αὐτοὺς ἐν τῇ ἄκρᾳ ἐν Ἱερουσαλὴμ ἐν φυλακῇ.

\vs{54}Καὶ ἐν ἔτει τρίτῳ καὶ πεντηκοστῷ καὶ ἑκατοστῷ, μηνι τῷ δευτέρῳ, ἐπέταξεν Ἄλκιμος καθαιρεῖν τὸ τεῖχος τῆς αὐλῆς τῶν ἁγίων τῆς ἐσωτέρας, καὶ καθεῖλε τὰ ἔργα τῶν προφητῶν, καὶ ἐνήρξατο τοῦ καθαιρεῖν.
\vs{55}Ἐν τῷ καιρῷ ἐκείνῳ ἐπλήγη Ἄλκιμος, καὶ ἐνεποδίσθη τὰ ἔργα αὐτοῦ, καὶ ἀπεφράγη τὸ στόμα αὐτοῦ, καὶ παρελύθη, καὶ οὐκ ἐδύνατο ἔτι λαλῆσαι λόγον καὶ ἐντείλασθαι περὶ τοῦ οἴκου αὐτοῦ.
\vs{56}Καὶ ἀπέθανεν Ἄλκιμος ἐν τῷ καιρῷ ἐκείνῳ μετὰ βασάνου μεγάλης.

\vs{57}Καὶ εἶδε Βακχίδης ὅτι ἀπέθανεν Ἄλκιμος, καὶ ἀπέστρεψε πρὸς τὸν βασιλέα· καὶ ἡσύχασεν ἡ γῆ Ἰούδα ἔτη δύο.
\vs{58}Καὶ ἐβουλεύσαντο πάντες οἱ ἄνομοι, λέγοντες, ἰδοὺ Ἰωνάθαν καὶ οἱ παρʼ αὐτοῦ ἐν ἡσυχίᾳ κατοικοῦσι πεποιθότες· νῦν οὖν ἄξομεν τὸν Βακχίδην, καὶ συλλήμψεται αὐτοὺς πάντας ἐν νυκτὶ μιᾷ.
\vs{59}Καὶ πορευθέντες συνεβουλεύσαντο αὐτῷ.
\vs{60}Καὶ ἀπῇρε τοῦ ἐλθεῖν μετὰ δυνάμεως πολλῆς, καὶ ἀπέστειλεν ἐπιστολὰς λάθρα πᾶσι τοῖς συμμάχοις αὐτοῦ τοῖς ἐν τῇ Ἰουδαίᾳ, ὅπως συλλάβωσι τὸν Ἰωνάθαν, καὶ τοὺς μετʼ αὐτοῦ· καὶ οὐκ ἐδύναντο, ὅτι ἐγνώσθη αὐτοῖς ἡ βουλὴ αὐτῶν.
\vs{61}Καὶ συνελάβοντο ἀπὸ τῶν ἀνδρῶν τῆς χώρας τῶν ἀρχηγῶν τῆς κακίας εἰς πεντήκοντα ἄνδρας, καὶ ἀπέκτειναν αὐτούς.

\vs{62}Καὶ ἐξεχώρησεν Ἰωνάθαν, καὶ Σίμων, καὶ οἱ μετʼ αὐτοῦ εἰς Βαιθβασὶ τὴν ἐν τῇ ἐρήμῳ, καὶ ᾠκοδόμησε τὰ καθῃρημένα αὐτῆς, καὶ ἐστερέωσαν αὐτήν.
\vs{63}Καὶ ἔγνω Βακχίδης, καὶ συνήγαγε πᾶν τὸ πλῆθος αὐτοῦ, καὶ τοῖς ἐκ τῆς Ἰουδαίας παρήγγειλε.

\vs{64}Καὶ ἐλθὼν παρενέβαλεν ἐπὶ Βαιθβασὶ, καὶ ἐπολέμησεν αὐτὴν ἡμέρας πολλὰς, καὶ ἐποίησε μηχανάς.
\vs{65}Καὶ ἀπέλιπεν Ἰωνάθαν Σίμωνα τὸν ἀδελφὸν αὐτοῦ ἐν τῇ πόλει, καὶ ἐξῆλθεν εἰς τὴν χώραν, καὶ ἐξῆλθεν ἐν ἀριθμῷ.
\vs{66}Καὶ ἐπάταξεν Ὀδοαῤῥὴν, καὶ τοὺς ἀδελφοὺς αὐτοῦ, καὶ τοὺς υἱοὺς Φασιρὼν ἐν τῷ σκηνώματι αὐτῶν, καὶ ἐξήρξατο τύπτειν, καὶ ἀναβαίνειν ἐν δυνάμεσι·
\vs{67}καὶ Σίμων, καὶ οἱ μετʼ αὐτοῦ ἐξῆλθον ἐκ τῆς πόλεως, καὶ ἐνεπύρισαν τὰς μηχανάς.
\vs{68}Καὶ ἐπολέμησαν πρὸς τὸν Βακχίδην, καὶ συνετρίβη ὑπʼ αὐτῶν, καὶ ἔθλιβον αὐτὸν σφόδρα, ὅτι ἦν ἡ βουλὴ αὐτοῦ καὶ ἡ ἔφοδος αὐτοῦ κενή.
\vs{69}Καὶ ὠργίσθηθυμῷ τοῖς ἀνδράσι τοῖς ἀνόμοις τοῖς συμβουλεύσασιν αὐτῷ ἐλθεῖν εἰς τὴν χώραν, καὶ ἀπέκτειναν ἐξ αὐτῶν πολλούς, καὶ ἐβουλεύσατο τοῦ ἀπελθεῖν εἰς τὴν γῆν αὐτοῦ.

\vs{70}Καὶ ἐπέγνω Ἰωνάθαν, καὶ ἀπέστειλε πρὸς αὐτὸν πρέσβεις, τοῦ συνθέσθαι πρὸς αὐτὸν εἰρήνην, καὶ ἀποδοῦναι αὐτοῖς τὴν αἰχμαλωσίαν.
\vs{71}Καὶ ἀπεδέξατο, καὶ ἐποίησε κατὰ τοὺς λόγους αὐτοῦ, καὶ ὤμοσεν αὐτῷ μὴ ἐκζητῆσαι αὐτῷ κακὸν πάσας τὰς ἡμέρας τῆς ζωῆς αὐτοῦ.
\vs{72}Καὶ ἀπέδωκεν αὐτῷ τὴν αἰχμαλωσίαν ἣν ᾐχμαλώτευσε τοπρότερον ἐκ γῆς Ἰούδα· καὶ ἀποστρέψας ἀπῆλθεν εἰς τὴν γῆν αὐτοῦ, καὶ οὐ προσέθετο ἔτι ἐλθεῖν εἰς τὰ ὅρια αὐτῶν.
\vs{73}Καὶ κατέπαυσε ῥομφαία ἐξ Ἰσραήλ· καὶ ᾤκησεν Ἰωνάθαν ἐν Μαχμάς· καὶ ἤρξατο Ἰωνάθαν κρίνειν τὸν λαὸν καὶ ἠφάνισε τοὺς ἀσεβεῖς ἐξ Ἰσραήλ.

\ch{10}
Καὶ ἐν ἔτει ἑξηκοστῷ καὶ ἑκατοστῷ ἀνέβη Ἀλέξανδρος ὁ τοῦ Ἀντιόχου ὁ Ἐπιφανής, καὶ κατελάβετο Πτολεμαΐδα, καὶ ἐπεδέξαντο αὐτὸν, καὶ ἐβασίλευσεν ἐκεῖ.
\vs{2}Καὶ ἤκουσε Δημήτριος ὁ βασιλεύς, καὶ συνήγαγε δυνάμεις πολλὰς σφόδρα, καὶ ἐξῆλθεν εἰς συνάντησιν αὐτῷ εἰς πόλεμον.
\vs{3}Καὶ ἀπέστειλε Δημήτριος πρὸς Ἰωνάθαν ἐπιστολὰς λόγοις εἰρηνικοῖς ὥστε μεγαλῦναι αὐτόν.
\vs{4}Εἶπε γάρ, προφθάσωμεν τοῦ εἰρήνην θεῖναι μετʼ αὐτοῦ, πρινὴ θεῖναι αὐτὸν μετὰ Ἀλεξάνδρου καθʼ ἡμῶν.
\vs{5}Μνησθήσεται γὰρ πάντων τῶν κακῶν ὧν συνετελέσαμεν πρὸς αὐτόν, καὶ εἰς τοὺς ἀδελφοὺς αὐτοῦ, καὶ εἰς τὸ ἔθνος αὐτοῦ.
\vs{6}Καὶ ἔδωκεν αὐτῷ ἐξουσίαν συναγαγεῖν δυνάμεις, καὶ κατασκευάζειν ὅπλα, καὶ εἶναι αὐτὸν σύμμαχον αὐτοῦ, καὶ τὰ ὅμηρα τὰ ἐν τῇ ἄκρᾳ εἶπε παραδοῦναι αὐτῷ.

\vs{7}Καὶ ἦλθεν Ἰωνάθαν εἰς Ἱερουσαλὴμ, καὶ ἀνέγνω τὰς ἐπιστολὰς εἰς τὰ ὦτα παντὸς τοῦ λαοῦ, καὶ τῶν ἐκ τῆς ἄκρας.
\vs{8}Καὶ ἐφοβήθησα φόβον μέγαν ὅτε ἤκουσαν ὅτι ἔδωκεν αὐτῷ ὁ βασιλεὺς ἐξουσίαν συναγαγεῖν δυνάμεις.
\vs{9}Καὶ παρέδωκαν οἱ ἐκ τῆς ἄκρας Ἰωνάθαν τὰ ὅμηρα, καὶ ἀπέδωκεν αὐτοὺς τοῖς γονεῦσιν αὐτῶν.

\vs{10}Καὶ ᾤκησεν Ἰωνάθαν ἐν Ἱερουσαλήμ, καὶ ἤρξατο οἰκοδομεῖν καὶ καινίζειν τὴν πόλιν.
\vs{11}Καὶ εἶπε πρὸς τοὺς ποιοῦντας τὰ ἔργα, οἰκοδομεῖν τὰ τείχη, καὶ τὸ ὄρος Σιὼν κυκλόθεν ἐκ λίθων τετραγώνων εἰς ὀχύρωσιν· καὶ ἐποίησαν οὕτως.

\vs{12}Καὶ ἔφυγον οἱ ἀλλογενεῖς οἱ ὄντες ἐν τοῖς ὀχυρώμασιν οἷς ᾠκοδόμησε Βακχίδης.
\vs{13}Καὶ κατέλιπεν ἕκαστος τὸν τόπον αὐτοῦ, καὶ ἀπῆλθεν εἰς τὴν γῆν αὐτοῦ.
\vs{14}Πλὴν ἐν Βαιθσούρᾳ ὑπελείφθησάν τινες τῶν καταλιπόντων τὸν νόμον καὶ τὰ προστάγματα, ἦν γὰρ αὐτοῖς φυγαδευτήριον.

\vs{15}Καὶ ἤκουσεν Ἀλέξανδρος ὁ βασιλεὺς τὰς ἐπαγγελίας ὅσας ἀπέστειλε Δημήτριος τῷ Ἰωνάθαν, καὶ διηγήσαντο αὐτῷ τοὺς πολέμους καὶ τὰς ἀνδραγαθίας ἃς ἐποίησεν αὐτὸς καὶ οἱ ἀδελφοὶ αὐτοῦ, καὶ τοὺς κόπους οὓς ἔσχον,
\vs{16}καὶ εἶπε, μὴ εὑρήσομεν ἄνδρα τοιοῦτον ἕνα; καὶ νῦν ποιήσομεν αὐτὸν φίλον, καὶ σύμμαχον ἡμῶν.

\vs{17}Καὶ ἔγραψεν ἐπιστολὰς, καὶ ἀπέστειλεν αὐτῷ κατὰ τοὺς λόγους τούτους, λέγων,
\vs{18}βασιλεὺς Ἀλέξανδρος τῷ ἀδελφῷ Ἰωνάθαν χαίρειν.
\vs{19}Ἀκηκόαμεν περὶ σοῦ, ὅτι ἀνὴρ δυνατὸς ἰσχύϊ, καὶ ἐπιτήδειος εἶ τοῦ εἶναι ἡμῖν φίλος.
\vs{20}Καὶ νῦν καθεστάκαμέν σε σήμερον ἀρχιερέα τοῦ ἔθνους σου, καὶ φίλον βασιλέως καλεῖσθαι· καὶ ἀπέστειλεν αὐτῷ πορφύραν καὶ στέφανον χρυσοῦν· καὶ φρονεῖν τὰ ἡμῶν, καὶ συντηρεῖν φιλίαν πρὸς ἡμᾶς.
\vs{21}Καὶ ἐνεδύσατο Ἰωνάθαν τὴν ἁγίαν στολὴν τῷ ἑβδόμῳ μηνὶ ἔτους ἑξηκοστοῦ καὶ ἑκατοστοῦ ἐν ἑορτῇ σκηνοπηγίας, καὶ συνήγαγε δυνάμεις, καὶ κατεσκεύασεν ὅπλα πολλά.

\vs{22}Καὶ ἤκουσε Δημήτριος τοὺς λόγους τούτους. καὶ ἐλυπήθη, καὶ εἶπε,
\vs{23}τί τοῦτο ἐποιήσαμεν, ὅτι προέφθακεν ἡμᾶς ὁ Ἀλέξανδρος τοῦ φιλίαν καταθέσθαι τοῖς Ἰουδαίοις εἰς στήριγμα;
\vs{24}Γράψω αὐτοῖς κᾀγὼ λόγους παρακλήσεως, καὶ ὕψους, καὶ δομάτων, ὅπως ὦσι σὺν ἐμοὶ εἰς βοήθειαν.
\vs{25}Καὶ ἀπέστειλεν αὐτοῖς κατὰ τοὺς λόγους τούτους· βασιλεὺς Δημήτριος τῷ ἔθνει τῶν Ἰουδαίων χαίρειν.
\vs{26}Ἐπεὶ συνετηρήσατε τὰς πρὸς ἡμᾶς συνθήκας, καὶ ἐνεμείνατε τῇ φιλίᾳ ἡμῶν, καὶ οὐ προσεχωρήσατε τοῖς ἐχθροῖς ἡμῶν, ἠκούσαμεν, καὶ ἐχάρημεν.
\vs{27}Καὶ νῦν ἐμμείνατε ἔτι τοῦ συντηρῆσαι πρὸς ἡμᾶς πίστιν, καὶ ἀνταποδώσομεν ὑμῖν ἀγαθὰ, ἀνθʼ ὧν ποιεῖτε μεθʼ ἡμῶν.
\vs{28}Καὶ ἀφήσομεν ὑμῖν ἀφέματα πολλὰ, καὶ δώσομεν ὑμῖν δόματα.

\vs{29}Καὶ νῦν ἀπολύω ὑμᾶς, καὶ ἀφίημι πάντας τοὺς Ἰουδαίους ἀπὸ τῶν φόρων, καὶ τῆς τιμῆς τοῦ ἁλὸς, καὶ ἀπὸ τῶν στεφάνων,
\vs{30}καὶ ἀντὶ τοῦ τρίτου τῆς σπορᾶς, καὶ ἀντὶ τοῦ ἡμίσους τοῦ καρποῦ τοῦ ξυλίνου τοῦ ἐπιβάλλοντός μοι λαβεῖν ἀφίημι ἀπὸ τῆς σήμερον καὶ ἐπέκεινα τοῦ λαβεῖν ἀπὸ τῆς γῆς Ἰούδα, καὶ ἀπὸ τῶν τριῶν νομῶν τῶν προστιθεμένων αὐτῇ ἀπὸ τῆς Σαμαρείτιδος καὶ Γαλιλαίας, καὶ ἀπὸ τῆς σήμερον ἡμέρας καὶ εἰς τὸν αἰῶνα χρόνον.

\vs{31}Καὶ Ἱερουσαλὴμ ἤτω ἁγία καὶ ἀφειμένη, καὶ τὰ ὅρια αὐτῆς, αἱ δεκάται καὶ τὰ τέλη.
\vs{32}Ἀφίημι καὶ τὴν ἐξουσίαν τῆς ἄκρας τῆς ἐν Ἱερουσαλὴμ, καὶ δίδωμι τῷ ἀρχιερεῖ, ὅπως ἂν καταστήσῃ ἐν αὐτῇ ἄνδρας οὓς ἂν ἐκλέξηται αὐτὸς τοῦ φυλάσσειν αὐτήν.

\vs{33}Καὶ πᾶσαν φυχὴν Ἰουδαίων τὴν αἰχμαλωτισθεῖσαν ἀπὸ γῆς Ἰούδα εἰς πᾶσαν βασιλείαν μου ἀφίημι ἐλευθέραν δωρέαν· καὶ πάντες ἀφιέτωσαν τοὺς φόρους καὶ τῶν κτηνῶν αὐτῶν.
\vs{34}Καὶ πᾶσαι αἱ ἑορταὶ καὶ τὰ σάββατα καὶ νουμηνίαι, καὶ ἡμέραι ἀποδεδειγμέναι, καὶ τρεῖς ἡμέραι πρὸ ἑορτῆς καὶ τρεῖς ἡμέραι μετὰ ἑορτὴν, ἔστωσαν πᾶσαι αἱ ἡμέραι ἀτελείας καὶ ἀφέσεως πᾶσι τοῖς Ἰουδαίοις τοῖς οὖσιν ἐν τῇ βασιλείᾳ μου.
\vs{35}Καὶ οὐχ ἕξει ἐξουσίαν οὐδεὶς πράσσειν καὶ παρενοχλεῖν τινα αὐτῶν περὶ παντὸς πράγματος.

\vs{36}Καὶ προγραφήτωσαν τῶν Ἰουδαίων εἰς τὰς δυνάμεις τοῦ βασιλέως εἰς τριάκοντα χιλιάδας ἀνδρῶν, καὶ δοθήσεται αὐτοῖς ξένια ὡς καθήκει πάσαις ταῖς δυνάμεσι τοῦ βασιλέως.
\vs{37}Καὶ κατασταθήσεται ἐξ αὐτῶν ἐν τοῖς ὀχυρώμασι τοῦ βασιλέως τοῖς μεγάλοις, καὶ ἐκ τούτων κατασταθήσεται ἐπὶ χρειῶν τῆς βασιλείας τῶν οὐσῶν εἰς πίστιν· καὶ οἱ ἐπʼ αὐτῶν καὶ οἱ ἄρχοντες ἔστωσαν ἐξ αὐτῶν· καὶ πορευέσθωσαν τοῖς νόμοις αὐτῶν, καθὰ καὶ προσέταξεν ὁ βασιλεὺς ἐν γῇ Ἰούδα.

\vs{38}Καὶ τοὺς τρεῖς νομοὺς τοὺς προστεθέντας τῇ Ἰουδαίᾳ ἀπὸ τῆς χώρας Σαμαρείας, προστεθήτω τῇ Ἰουδαίᾳ πρὸς τὸ λογισθῆναι τοῦ γενέσθαι ὑφʼ ἕνα, τοῦ μὴ ὑπακοῦσαι ἄλλης ἐξουσίας ἀλλʼ ἢ τοῦ ἀρχιερέως.

\vs{39}Πτολεμαΐδα καὶ τὴν προσκυροῦσαν αὐτῇ δέδωκα δόμα τοῖς ἁγίοις τοῖς ἐν Ἰερουσαλὴμ εἰς τὴν προσήκουσαν δαπάνην τοῖς ἁγίοις.
\vs{40}Κᾀγὼ δίδωμι κατʼ ἐνιαυτὸν δεκαπέντε χιλιάδας σίκλων ἀργυρίου ἀπὸ τῶν λόγων τοῦ βασιλέως, ἀπὸ τῶν τόπων τῶν ἀνηκόντων.
\vs{41}Καὶ πᾶν τὸ πλεονάζον ὃ οὐκ ἀπεδίδοσαν οἱ ἀπὸ τῶν χρειῶν, ὡς ἐν τοῖς πρῶτοις ἔτεσιν, ἀπὸ τοῦ νῦν δώσουσιν εἰς τὰ ἔργα τοῦ οἴκου.

\vs{42}Καὶ ἐπὶ τούτοις, πεντακισχιλίους σίκλους ἀργυρίου, οὓς ἐλάμβανον ἀπὸ τῶν χρειῶν τοῦ ἁγίου ἀπὸ τοῦ λόγου κατʼ ἐνιαυτὸν, καὶ ταῦτα ἀφίεται διὰ τὸ ἀνήκειν αὐτὰ τοῖς ἱερεῦσι τοῖς λειτουργοῦσι.
\vs{43}Καὶ ὅσοι ἐὰν φύγωσιν εἰς τὸ ἱερὸν τὸ ἐν Ἱεροσολύμοις καὶ ἐν πᾶσι τοῖς ὁρίοις αὐτοῦ, ὀφείλοντες βασιλικὰ καὶ πᾶν πρᾶγμα, ἀπολελύσθωσαν, καὶ πάντα ὅσα ἐστὶν αὐτοῖς ἐν τῇ βασιλείᾳ μου.
\vs{44}Καὶ τοῦ οἰκοδομηθῆναι καὶ τοῦ ἐπικαινισθῆναι τὰ ἔργα τῶν ἁγίων, καὶ ἡ δαπάνη δοθήσεται ἐκ τοῦ λόγου τοῦ βασιλέως.
\vs{45}Καὶ τοῦ οἰκοδομηθῆναι τὰ τείχη Ἱερουσαλὴμ καὶ ὀχυρῶσαι κυκλόθεν, καὶ ἡ δαπάνη δοθήσεται ἐκ τοῦ λόγου τοῦ βασιλέως, καὶ τοῦ οἰκοδομηθῆναι τὰ τείχη τὰ ἐν τῇ Ἰουδαίᾳ.

\vs{46}Ὡς δὲ ἤκουσεν Ἰωνάθαν καὶ ὁ λαὸς τοὺς λόγους τούτους, οὐκ ἐπίστευσαν αὐτοῖς οὐδὲ ἐπεδέξαντο, ὅτι ἐπεμνήσθησαν τῆς κακίας τῆς μεγάλης ἧς ἐποίησεν ἐν Ἰσραὴλ, καὶ ἔθλιψεν αὐτοὺς σφόδρα.
\vs{47}Καὶ εὐδόκησαν ἐν Ἀλεξάνδρῳ, ὅτι αὐτὸς ἐγένετο αὐτοῖς ἀρχηγὸς λόγων εἰρηνικῶν, καὶ συνεμάχουν αὐτῷ πάσας τὰς ἡμέρας.

\vs{48}Καὶ συνήγαγεν Ἀλέξανδρος ὁ βασιλεὺς δυνάμεις μεγάλας, καὶ παρενέβαλεν ἐξεναντίας Δημητρίου.
\vs{49}Καὶ συνῆψαν πόλεμον οἱ δύο βασιλεῖς, καὶ ἔφυγεν ἡ παρεμβολὴ Δημητρίου, καὶ ἐδίωξεν αὐτὸν ὁ Ἀλέξανδρος, καὶ ἴσχυσεν ἐπʼ αὐτούς.
\vs{50}Καὶ ἐστερέωσε τὸν πόλεμον σφόδρα ἕως ἔδυ ὁ ἥλιος, καὶ ἔπεσεν ὁ Δημήτριος ἐν τῇ ἡμέρᾳ ἐκείνῃ.

\vs{51}Καὶ ἀπέστειλεν Ἀλέξανδρος πρὸς Πτολεμαῖον βασιλέα Αἰγύπτου πρέσβεις κατὰ τοὺς λόγους τούτους, λέγων,
\vs{52}ἐπεὶ ἀνέστρεψα εἰς γῆν βασιλείας μου, καὶ ἐκάθισα ἐπὶ θρόνου πατέρων μου, καὶ ἐκράτησα τῆς ἀρχῆς, καὶ συνέτριψα τὸν Δημήτριον, καὶ ἐπεκράτησα τῆς χώρας ἡμῶν·
\vs{53}καὶ συνῆψα πρὸς αὐτὸν μάχην, καὶ συνετρίβη αὐτὸς καὶ ἡ παρεμβολὴ αὐτοῦ ὑφʼ ἡμῶν, καὶ ἐκαθίσαμεν ἐπὶ θρόνου βασιλείας αὐτοῦ·
\vs{54}καὶ νῦν στήσωμεν πρὸς ἑαυτοὺς φιλίαν, καὶ νῦν δός μοι τὴν θυγατέρα σου εἰς γυναῖκα, καὶ ἐπιγαμβρεύσω σοι, καὶ δώσω σοι δόματα, καὶ αὐτῇ. ἄξιά σου.

\vs{55}Καὶ ἀπεκρίθη Πτολεμαῖος ὁ βασιλεὺς, λέγων, ἀγαθὴ ἡμέρα ἐν ᾗ ἀνέστρεψας εἰς γῆν πατέρων σου, καὶ ἐκάθισας ἐπὶ θρόνου βασιλείας αὐτῶν.
\vs{56}Καὶ νῦν ποιήσω σοι ἃ ἔγραψας, ἀλλʼ ἀπάντησον εἰς Πτολεμαΐδα, ὅπως ἴδωμεν ἀλλήλους, καὶ ἐπιγαμβρεύσω σοι καθὼς εἴρηκας.

\vs{57}Καὶ ἐξῆλθε Πτολεμαῖος ἐξ Αἰγύπτου αὐτὸς καὶ Κλεοπάτρα ἡ θυγάτηρ αὐτοῦ, καὶ εἰσῆλθον εἰς Πτολεμαΐδα ἔτους δευτέρου καὶ ἑξηκοστοῦ καὶ ἑκατοστοῦ.
\vs{58}Καὶ ἀπήντησεν αὐτῷ Ἀλεξανδρος ὁ βασιλεὺς, καὶ ἐξέδοτο αὐτῷ Κλεοπάτραν τὴν θυγατέρα αὐτοῦ, καὶ ἐποίησε τὸν γάμον αὐτῆς ἐν Πτολεμαΐδι, καθὼς οἱ βασιλεῖς, ἐν δόξῃ μεγάλῃ.

\vs{59}Καὶ ἔγραψεν Ἀλέξανδρος ὁ βασιλεὺς τῷ Ἰωνάθαν ἐλθεῖν εἰς συνάντησιν αὐτῷ.
\vs{60}Καὶ ἐπορεύθη μετὰ δόξης εἰς Πτολεμαΐδα, καὶ ἀπήντησε τοῖς δυσὶ βασιλεῦσι· καὶ ἔδωκεν αὐτοῖς ἀργύριον καὶ χρυσίον, καὶ τοῖς φίλοις αὐτῶν, καὶ δόματα πολλὰ, καὶ εὗρε χάριν ἐναντίον αὐτῶν.

\vs{61}Καὶ ἐπισυνήχθησαν πρὸς αὐτὸν ἄνδρες λοιμοὶ ἐξ Ἰσραὴλ, ἄνδρες παράνομοι, ἐντυχεῖν κατʼ αὐτοῦ, καὶ οὐ προσέσχεν αὐτοῖς ὁ βασιλεύς.
\vs{62}Καὶ προσέταξεν ὁ βασιλεὺς, καὶ ἐξέδυσαν Ἰωνάθαν τὰ ἱμάτια αὐτοῦ, καὶ ἐνέδυσαν αὐτὸν πορφύραν, καὶ ἐποίησαν οὕτως.
\vs{63}Καὶ ἐκάθισεν αὐτὸν ὁ βασιλεὺς μετʼ αὐτοῦ, καὶ εἶπε τοῖς ἄρχουσιν αὐτοῦ, ἐξέλθετε μετʼ αὐτοῦ εἰς μέσον τῆς πόλεως, καὶ κηρύξατε τοῦ μηδένα ἐντυγχάνειν κατʼ αὐτοῦ περὶ μηδενὸς πράγματος, καὶ μηδεὶς αὐτῷ παρενοχλείτω περὶ παντὸς λόγου.

\vs{64}Καὶ ἐγένετο ὡς ἴδον οἱ ἐντυγχάνοντες τὴν δόξαν αὐτοῦ καθὼς ἐκήρυξαν, καὶ περιβεβλημένον αὐτὸν πορφύραν, καὶ ἔφυγον πάντες.
\vs{65}Καὶ ἐδόξασεν αὐτὸν ὁ βασιλεύς, καὶ ἔγραψεν αὐτὸν τῶν πρώτων φίλων, καὶ ἔθετο αὐτὸν στρατηγὸν καὶ μεριδάρχην.
\vs{66}Καὶ ἐπέστρεψεν Ἰωνάθαν εἰς Ἱερουσαλὴμ μετʼ εἰρήνης καὶ εὐφροσύνης.

\vs{67}Καί ἐν ἔτει πέμπτῳ καὶ ἑξηκοστῷ καὶ ἑκατοστῷ ἦλθε Δημήτριος υἱὸς Δημητρίου ἐκ Κρήτης εἰς τὴν γῆν τῶν πατέρων αὐτοῦ.
\vs{68}Καὶ ἤκουσεν Ἀλέξανδρος ὁ βασιλεὺς, καὶ ἐλυπήθη σφόδρα, καὶ ἀπέστρεψεν εἰς Ἀντιόχειαν.

\vs{69}Καὶ κατέστησε Δημήτριος Ἀπολλώνιον τὸν ὄντα ἐπὶ κοίλης Συρίας, καὶ συνήγαγε δύναμιν μεγάλην, καὶ παρενέβαλεν ἐν Ἰαμνείᾳ·
\vs{70}καὶ ἀπέστειλε πρὸς Ἰωνάθαν τὸν ἀρχιερέα, λέγων, σὺ μονώτατος ἐπαίρῃ ἐφʼ ἡμᾶς, ἐγὼ δὲ ἐγενήθην εἰς καταγέλωτα καὶ ὀνειδισμὸν διὰ σέ· καὶ διατί σὺ ἐξουσιάζῃ ἐφʼ ἡμᾶς ἐν τοῖς ὄρεσι;

\vs{71}Νῦν οὖν εἰ πέποιθας ἐπὶ ταῖς δυνάμεσί σου, κατάβηθι πρὸς ἡμᾶς εἰς τὸ πεδίον, καὶ συγκριθῶμεν ἑαυτοῖς ἐκεῖ, ὅτι μετʼ ἐμοῦ ἐστι δύναμις τῶν πόλεων.
\vs{72}Ἐρώτησον καὶ μάθε τίς εἰμι καὶ οἱ λοιποὶ οἱ βοηθοῦντες ἡμῖν, καὶ λέγουσιν, οὐκ ἔστιν ὑμῖν στάσις ποδὸς κατὰ πρόσωπον ἡμῶν; ὅτι δὶς ἐτροπώθησαν οἱ πατέρες σου ἐν τῇ γῇ αὐτῶν.
\vs{73}Καὶ νῦν οὐ δυνήσῃ ὑποστῆναι τὴν ἵππον καὶ δύναμιν τοιαύτην ἐν τῷ πεδίῳ, ὅπου οὐκ ἔστι λίθος οὐδὲ κόχλαξ οὐδὲ τόπος τοῦ φυγεῖν.

\vs{74}Ὡς δὲ ἤκουσεν Ἰωνάθαν τῶν λόγων Ἀπολλωνίου, ἐκινήθη τῇ διανοίᾳ, καὶ ἐπέλεξε δέκα χιλιάδας ἀνδρῶν, καὶ ἐξῆλθεν ἐξ Ἱερουσαλὴμ, καὶ συνήντησεν αὐτῷ Σίμων ὁ ἀδελφὸς αὐτοῦ ἐπὶ βοήθειαν αὐτοῦ.
\vs{75}Καὶ παρενέβαλεν ἐπὶ Ἰόππην, καὶ ἀπέκλεισαν αὐτὸν ἐκ τῆς πόλεως, ὅτι φρουρὰ Ἀπολλωνίου ἐν Ἰόππῃ, καὶ ἐπολέμησαν αὐτήν.

\vs{76}Καὶ φοβηθέντες ἤνοιξαν οἱ ἐκ τῆς πόλεως, καὶ ἐκυρίευσεν Ἰωνάθαν Ἰόππης.
\vs{77}Καὶ ἤκουσεν Ἀπολλώνιος, καὶ παρενέβαλε τρισχιλίαν ἵππον, καὶ δύναμιν πολλήν· καὶ ἐπορεύθη εἰς Ἄζωτον ὡς διοδεύων, καὶ ἅμα προῆγεν εἰς τὸ πεδίον, διὰ τὸ ἔχειν αὐτὸν πλῆθος ἵππου καὶ πεποιθέναι ἐπʼ αὐτῇ.

\vs{78}Καὶ κατεδίωξεν Ἰωνάθαν ὀπίσω αὐτοῦ εἰς Ἄζωτον, καὶ συνῆψαν αἱ παρεμβολαὶ εἰς πόλεμον.
\vs{79}Καὶ ἀπέλιπεν Ἀπολλώνιος χιλίαν ἵππον ἐν κρυπτῷ κατόπισθεν αὐτῶν.
\vs{80}Καὶ ἔγνω Ἰωνάθαν ὅτι ἐστὶν ἔνεδρον κατόπισθεν αὐτοῦ, καὶ ἐκύκλωσαν αὐτοῦ τὴν παρεμβολὴν, καὶ ἐξετίναξαν τὰς σχίζας εἰς τὸν λαὸν ἐκ πρωΐθεν ἕως ἑσπέρας.

\vs{81}Ὁ δὲ λαὸς εἱστήκει, καθὼς ἐπέταξεν Ἰωνάθαν, καὶ ἐκοπίασαν οἱ ἵπποι αὐτῶν.
\vs{82}Καὶ εἵλκυσε Σίμων τὴν δύναμιν αὐτοῦ, καὶ συνῆψε πρὸς τὴν φάραγγα· ἡ γὰρ ἵππος ἐξελύθη· καὶ συνετρίβησαν ὑπʼ αὐτοῦ, καὶ ἔφυγον.
\vs{83}Καὶ ἡ ἵππος ἐσκορπίσθη ἐν τῷ πεδίῳ, καὶ ἔφυγον εἰς Ἄζωτον, καὶ εἰσῆλθον εἰς Βηθδαγὼν τὸ εἰδωλεῖον αὐτῶν, τοῦ σωθῆναι

\vs{84}Καὶ ἐνεπύρισεν Ἰωνάθαν τὴν Ἄζωτον καὶ τὰς πόλεις τὰς κύκλῳ αὐτῆς, καὶ ἔλαβε τὰ σκῦλα αὐτῶν, καὶ τὸ ἱερὸν Δαγὼν καὶ τοὺς συμφυγόντας εἰς αὐτὸ ἐνεπύρισε πυρί.
\vs{85}Καὶ ἐγένοντο οἱ πεπτωκότες μαχαίρᾳ σὺν τοῖς ἐμπυρισθεῖσιν εἰς ἄνδρας ὀκτακισχιλίους.
\vs{86}Καὶ ἀπῇρεν ἐκεῖθεν Ἰωνάθαν, καὶ παρενέβαλεν ἐπὶ Ἀσκάλωνα, καὶ ἐξῆλθον οἱ ἐκ τῆς πόλεως εἰς συνάντησιν αὐτῷ ἐν δόξῃ μεγάλῃ.
\vs{87}Καὶ ἐπέστρεψεν Ἰωνάθαν εἰς Ἱερουσαλὴμ σὺν τοῖς παρʼ αὐτοῦ, ἔχοντες σκῦλα πολλά.

\vs{88}Καὶ ἐγένετο ὡς ἤκουσεν Ἀλέξανδρος ὁ βασιλεὺς τοὺς λόγους τούτους, καὶ προσέθετο δοξάσαι τὸν Ἰωνάθαν.
\vs{89}Καὶ ἀπέστειλεν αὐτῷ πόρπην χρυσῆν, ὡς ἔθος ἐστὶ δίδοσθαι τοῖς συγγενέσι τῶν βασιλέων· καὶ ἔδωκεν αὐτῷ τὴν Ἀκκαρὼν καὶ πάντα τὰ ὅρια αὐτῆς εἰς κληροδοσίαν.

\ch{11}
Καὶ ὁ βασιλεὺς Αἰγύπτου ἤθροισε δυνάμεις πολλὰς, ὡς τὴν ἄμμον τὴν περὶ τὸ χεῖλος τῆς θαλάσσης, καὶ πλοῖα πολλά· καὶ ἐζήτησε κατακρατῆσαι τῆς βασιλείας Ἀλεξάνδρου δόλῳ, καὶ προσθεῖναι αὐτὴν τῇ βασιλείᾳ αὑτοῦ.
\vs{2}Καὶ ἐξῆλθεν εἰς Συρίαν λόγοις εἰρηνικοῖς, καὶ ἤνοιγον αὐτῷ οἱ ἀπὸ τῶν πόλεων, καὶ συνήντων αὐτῷ, ὅτι ἐντολὴ ἦν Ἀλεξάνδρου τοῦ βασιλέως συναντᾷν αὐτῷ, διὰ τὸ πενθερὸν αὐτοῦ εἶναι,

\vs{3}Ὡς δὲ εἰσεπορεύετο εἰς τὰς πόλεις Πτολεμαῖος, ἀπέτασσε τὰς δυνάμεις φρουρὰν ἐν ἑκάστῃ πόλει.
\vs{4}Ὡς δὲ ἤγγισεν Ἀζώτου, ἔδειξαν αὐτῷ τὸ ἱερὸν Δαγὼν ἐμπεπυρισμένον, καὶ Ἄζωτον, καὶ τὰ περιπόλια αὐτῆς καθῃρημένα, καὶ τὰ σώματα ἐῤῥιμμένα, καὶ τοὺς ἐμπεπυρισμένους οὓς ἐνεπύρισεν ἐν τῷ πολέμῳ· ἐποίησαν γὰρ θημωνίας αὐτῶν ἐν τῇ ὁδῷ αὐτοῦ.
\vs{5}Καὶ διηγήσαντο τῷ βασιλεῖ ἃ ἐποίησεν Ἰωνάθαν, εἰς τὸ ψογῆσαι αὐτόν· καὶ ἐσίγησεν ὁ βασιλεύς.

\vs{6}Καὶ συνήντησεν Ἰωνάθαν τῷ βασιλεῖ εἰς Ἰόππην μετὰ δόξης, καὶ ἠσπάσαντο ἀλλήλους, καὶ ἐκοιμήθησαν ἐκεῖ.
\vs{7}Καὶ ἐπορεύθη Ἰωνάθαν μετὰ τοῦ βασιλέως ἕως τοῦ ποταμοῦ τοῦ καλουμένου Ἐλευθέρου, καὶ ἐπέστρεψεν εἰς Ἱερουσαλήμ.

\vs{8}Ὁ δὲ βασιλεὺς Πτολεμαῖος ἐκυρίευσε τῶν πόλεων τῆς παραλίας ἕως Σελευκίας τῆς παραθαλασσίας, καὶ διελογίζετο περὶ Ἀλεξάνδρου λογισμοὺς πονηρούς.
\vs{9}Καὶ ἀπέστειλε πρέσβεις πρὸς Δημήτριον τὸν βασιλέα, λέγων, δεῦρο συνθώμεθα πρὸς ἑαυτοὺς διαθήκην, καὶ δώσω σοι τὴν θυγατέρα μου ἣν ἔχει Ἀλέξανδρος, καὶ βασιλεύσεις τῆς βασιλείας τοῦ πατρός σου.
\vs{10}Μεταμεμέλημαι γὰρ δοὺς αὐτῷ τὴν θυγατέρα μου, ἐζήτησε γὰρ ἀποκτεῖναί με.
\vs{11}Καὶ ἐψόγησεν αὐτὸν χάριν τοῦ ἐπιθυμῆσαι αὐτὸν τῆς βασιλείας αὐτοῦ·

\vs{12}Καὶ ἀφελόμενος αὐτοῦ τὴν θυγατέρα, ἔδωκεν αὐτὴν τῷ Δημητρίῳ, καὶ ἠλλοιώθη τοῦ Ἀλεξάνδρου, καὶ ἐφάνη ἡ ἔχθρα αὐτῶν.
\vs{13}Καὶ εἰσῆλθε Πτολεμαῖος εἰς Ἀντιόχειαν, καὶ περιέθετο δύο διαδήματα περὶ τὴν κεφαλὴν αὐτοῦ, τὸ τῆς Ἀσίας καὶ Αἰγύπτου.

\vs{14}Ἀλέξανδρος δὲ ὁ βασιλεὺς ἦν ἐν Κιλικίᾳ κατὰ τοὺς καιροὺς ἐκείνους, ὅτι ἀπεστάτουν οἱ ἀπὸ τῶν τόπων ἐκείνων.
\vs{15}Καὶ ἤκουσεν Ἀλέξανδρος, καὶ ἦλθεν ἐπʼ αὐτὸν πολέμῳ· καὶ ἐξήγαγε Πτολεμαῖος τὴν δύναμιν, καὶ ἀπήντησεν αὐτῷ ἐν χειρὶ ἰσχυρᾷ, καὶ ἐτροπώσατο αὐτόν.

\vs{16}Καὶ ἔφυγεν Ἀλέξανδρος εἰς τὴν Ἀραβίαν, τοῦ σκεπασθῆναι αὐτὸν ἐκεῖ· ὁ δὲ βασιλεὺς Πτολεμαῖος ὑψώθη.
\vs{17}Καὶ ἀφειλε Ζαβδιὴλ ὁ Ἄραψ τὴν κεφαλὴν Ἀλεξάνδρου, καὶ ἀπέστειλε τῷ Πτολεμαίῳ.

\vs{18}Καὶ ὁ βασιλεὺς Πτολεμαῖος ἀπέθανεν ἐν τῇ ἡμέρᾳ τῇ τρίτῃ, καὶ οἱ ὄντες ἐν τοῖς ὀχυρώμασιν ἀπώλοντο ὑπὸ τῶν ἐν τοῖς ὀχυρώμασι.
\vs{19}Καὶ ἐβασίλευσε Δημήτριος ἔτους ἑβδόμου καὶ ἑξηκοστοῦ καὶ ἑκατοστοῦ.

\vs{20}Ἐν ταῖς ἡμέραις ἐκείναις συνήγαγεν Ἰωνάθαν τοὺς ἐκ τῆς Ἰουδαίας, τοῦ ἐκπολεμῆσαι τὴν ἄκραν τὴν ἐν Ἱερουσαλήμ, καὶ ἐποίησεν ἐπʼ αὐτὴν μηχανὰς πολλάς.
\vs{21}Καὶ ἐπορεύθησάν τινες μισοῦντες τὸ ἔθνος αὐτῶν, ἄνδρες παράνομοι, πρὸς τὸν βασιλέα, καὶ ἀπήγγειλαν αὐτῷ ὅτι Ἰωνάθαν περικάθηται τὴν ἄκραν.
\vs{22}Καὶ ἀκούσας ὠργίσθη· ὡς δὲ ἤκουσεν, εὐθέως ἀναζεύξας ἦλθεν εἰς Πτολεμαΐδα, καὶ ἔγραψεν Ἰωνάθαν, τοῦ μὴ περικαθῆσθαι τῇ ἄκρᾳ, καὶ τοῦ ἀπαντῆσαι αὐτὸν αὐτῷ συνμίσγειν εἰς Πτολεμαΐδα τὴν ταχίστην.

\vs{23}Ὡς δὲ ἤκουσεν Ἰωνάθαν, ἐκέλευσε περικαθῆσθαι, καὶ ἐπέλεξε τῶν πρεσβυτέρων Ἰσραὴλ καὶ τῶν ἱερέων, καὶ ἔδωκεν ἑαυτὸν τῷ κινδύνῳ.
\vs{24}Καὶ λαβὼν ἀργύριον, καὶ χρυσίον, καὶ ἱματισμὸν, καὶ ἕτερα ξένια πλείονα, ἐπορεύθη πρὸς τὸν βασιλέα εἰς Πτολεμαΐδα, καὶ εὗρε χάριν ἑνώτιον αὐτοῦ.

\vs{25}Καὶ ἐνετύγχανον κατʼ αὐτοῦ τινὲς ἄνομοι τῆς ἐκ τοῦ ἔθνους.
\vs{26}Καὶ ἐποίησεν αὐτῷ ὁ βασιλεὺς καθὼς ἐποίησαν αὐτῷ οἱ πρὸ αὐτοῦ, καὶ ὕψωσεν αὐτὸν ἐναντίον πάντων τῶν φίλων αὐτοῦ.
\vs{27}Καὶ ἔστησεν αὐτῷ τὴν ἀρχιερωσύνην, καὶ ὅσα ἄλλα εἶχε τίμια τοπρότερον, καὶ ἐποίησεν αὐτὸν τῶν πρώτων φίλων ἡγεῖσθαι.

\vs{28}Καὶ ἠξίωσεν Ἰωνάθαν τὸν βασιλέα ποιῆσαι τὴν Ἰουδαίαν ἀφορολόγητον, καὶ τὰς τρεῖς τοπαρχίας, καὶ τὴν Σαμαρεῖτιν, καὶ ἐπηγγείλατο αὐτῷ τάλαντα τριακόσια.
\vs{29}Καὶ εὐδόκησεν ὁ βασιλεύς, καὶ ἔγραψε τῷ Ἰωνάθαν ἐπιστολὰς περὶ πάντων τούτων ἐχούσας τὸν τρόπον τοῦτον·

\vs{30}Βασιλεὺς Δημήτριος Ἰωνάθαν τῷ ἀδελφῷ χαίρειν, καὶ ἔθνει Ἰουδαίων.
\vs{31}Τὸ ἀντίγραφον τῆς ἐπιστολῆς ἧς ἐγράψαμεν Λασθένει τῷ συγγενεῖ ἡμῶν περὶ ὑμῶν, γεγράφαμεν καὶ πρὸς ὑμᾶς. ὅπως εἰδῆτε.

\vs{32}Βασιλεὺς Δημήτριος Λασθένει τῷ πατρὶ χαίρειν.
\vs{33}Τῷ ἔθνει τῶν Ἰουδαίων φίλοις ἡμῶν καὶ συντηροῦσι τὰ πρὸς ἡμᾶς δίκαια ἐκρίναμεν ἀγαθοποιῆσαι, χάριν τῆς ἑξ αὐτῶν εὐνοίας πρὸς ἡμᾶς.
\vs{34}Ἑστάκαμεν οὖν αὐτοῖς τὰ τε ὅρια τῆς Ἰουδαίας, καὶ τοὺς τρεῖς νομοὺς, Ἀφαίρεμα, καὶ Λύδδαν, καὶ Ῥαμαθὲμ, αἵτινες προσετέθησαν τῇ Ἰουδαίᾳ ἀπὸ τῆς Σαμαρείτιδος, καὶ πάντα τὰ συγκυροῦντα αὐτοῖς πᾶσι τοῖς θυσιάζουσιν εἰς Ἱεροσόλυμα, ἀντὶ τῶν βασιλικῶν ὧν ἐλάμβανεν ὁ βασιλεὺς παρʼ αὐτῶν τοπρότερον κατʼ ἐνιαυτὸν ἀπὸ τῶν γεννημάτων τῆς γῆς, καὶ ἀπὸ τῶν ἀκροδρύων.

\vs{35}Καὶ τὰ ἄλλα τὰ ἀνήκοντα ἡμῖν ἀπὸ τοῦ νῦν τῶν δεκατῶν, καὶ τῶν τελῶν τῶν ἀνηκόντων ἡμῖν, καὶ τὰς τοῦ ἁλὸς λίμνας, καὶ τοὺς ἀνήκοντας ἡμῖν στεφάνους, πάντα ἐπαρκῶς παρίεμεν αὐτοῖς.
\vs{36}Καὶ οὐκ ἀθετηθήσεται οὐδὲ ἓν τούτων ἀπὸ τοῦ νῦν καὶ εἰς τὸν ἅπαντα χρόνον.

\vs{37}Νῦν οὖν ἐπιμέλεσθε τοῦ ποιῆσαι τούτων ἀντίγραφον· καὶ δοθήτω Ἰωνάθαν, καὶ τεθήτω ἐν τῷ ὄρει τῷ ἁγίῳ ἐν τόπῳ ἐπισήμῳ.

\vs{38}Καὶ εἶδε Δημήτριος ὁ βασιλεὺς ὅτι ἡσύχασεν ἡ γῆ ἐνώπιον αὐτοῦ, καὶ οὐδὲν αὐτῷ ἀνθειστήκει, καὶ ἀπέλυσε πάσας τὰς δύναμεις αὐτοῦ ἕκαστον εἰς τὸν ἴδιον τόπον, πλὴν τῶν ξένων δυνάμεων ὧν ἐξενολόγησεν ἀπὸ τῶν νήσων τῶν ἐθνῶν· καὶ ἤχθραναν αὐτῷ πᾶσαι αἱ δυνάμεις τῶν πατέρων αὐτοῦ.

\vs{39}Τρύφων δὲ ἦν τῶν παρὰ Ἀλεξάνδρου τοπρότερον, καὶ εἶδεν ὅτι πᾶσαι αἱ δυνάμεις καταγογγύζουσι τοῦ Δημητρίου, καὶ ἐπορεύθη πρὸς Εἰμαλκουαὶ τὸν Ἄραβα, ὃς ἔτρεφε τὸν Ἀντίοχον τὸ παιδάριον τὸ τοῦ Ἀλεξάνδρου·
\vs{40}καὶ προσήδρευεν αὐτῷ, ὅπως παραδοῖ αὐτὸν αὐτῷ, ὅπως βασιλεύσῃ ἀντὶ τοῦ πατρὸς αὐτοῦ· καὶ ἀπήγγειλεν αὐτῷ ὅσα συνετέλεσε Δημήτριος, καὶ τὴν ἔχθραν ἣν ἐχθραίνουσιν αὐτῷ αἱ δυνάμεις αὐτοῦ· καὶ ἔμεινεν ἐκεῖ ἡμέρας πολλάς.

\vs{41}Καὶ ἀπέστειλεν Ἰωνάθαν πρὸς Δημήτριον τὸν βασιλέα, ἵνα ἐκβάλῃ τοὺς ἐκ τῆς ἄκρας ἐξ Ἱερουσαλὴμ, καὶ τοὺς ἐν τοῖς ὀχυρώμασιν, ἦσαν γὰρ πολεμοῦντες τὸν Ἰσραήλ.
\vs{42}Καὶ ἀπέστειλε Δημήτριος πρὸς Ἰωνάθαν, λέγων, οὐ μόνον ταῦτα ποιήσω σοι καὶ τῷ ἔθνει σου, ἀλλὰ δόξῃ δοξάσω σε καὶ τὸ ἔθνος σου, ἐὰν εὐκαιρίας τύχω.
\vs{43}Νῦν οὖν ὀρθῶς ποιήσεις ἀποστείλας μοι ἄνδρας οἳ συμμαχήσουσιν, ὅτι ἀπέστησαν πᾶσαι αἱ δυνάμεις μου.

\vs{44}Καὶ ἀπέστειλεν Ἰωνάθαν ἄνδρας τρισχιλίους δυνατοὺς ἰσχύϊ αὐτῷ εἰς Ἀντιόχειαν, καὶ ἤλθοσαν πρὸς τὸν βασιλέα, καὶ εὐφράνθη ὁ βασιλεὺς ἐπὶ τῇ ἐφόδῳ αὐτῶν.
\vs{45}Καὶ ἐπισυνήχθησαν οἱ ἐκ τῆς πόλεως εἰς μέσον τῆς πόλεως εἰς ἀνδρῶν δώδεκα μυριάδας, καὶ ἠβούλοντο ἀνελεῖν τὸν βασιλέα.
\vs{46}Καὶ ἔφυγεν ὁ βασιλεὺς εἰς τὴν αὐλὴν, καὶ κατελάβοντο οἱ ἐκ τῆς πόλεως τὰς διόδους τῆς πόλεως, καὶ ἤρξαντο πολεμεῖν.

\vs{47}Καὶ ἐκάλεσεν ὁ βασιλεὺς τοὺς Ἰουδαίους ἐπὶ βοήθειαν, καὶ ἐπισυνήχθησαν πρὸς αὐτὸν πάντες ἅμα· καὶ διεσπάρησαν ἐν τῇ πόλει πάντες ἅμα· καὶ ἀπέκτειναν ἐν τῇ πόλει τῇ ἡμέρᾳ ἐκείνῃ εἰς μυριάδας δέκα.
\vs{48}Καὶ ἐνεπύρισαν τὴν πόλιν, καὶ ἐλάβοσαν σκῦλα πολλὰ ἐν ἐκείνῃ τῇ ἡμέρᾳ, καὶ ἔσωσαν τὸν βασιλέα.

\vs{49}Καὶ ἴδον οἱ ἀπὸ τῆς πόλεως ὅτι κατεκράτησαν οἱ Ἰουδαῖοι τῆς πόλεως, ὡς ἠβούλοντο, καὶ ἠσθένησαν ταῖς διανοίαις αὐτῶν, καὶ ἐκέκραξαν πρὸς τὸν βασιλέα μετὰ δεήσεως, λέγοντες,
\vs{50}δὸς ἡμῖν δεξιὰς, καὶ παυσάσθωσαν οἱ Ἰουδαῖοι πολεμοῦντες ἡμᾶς καὶ τὴν πόλιν.
\vs{51}Καὶ ἔῤῥιψαν τὰ ὅπλα, καὶ ἐποίησαν εἰρήνην· καὶ ἐδοξάσθησαν οἱ Ἰουδαῖοι ἐναντίον τοῦ βασιλέως, καὶ ἐνώπιον πάντων τῶν ἐν τῇ βασιλείᾳ αὐτοῦ, καὶ ἐπέστρεψαν εἰς Ἱερουσαλὴμ ἔχοντες σκῦλα πολλά.

\vs{52}Καὶ ἐκάθισε Δημήτριος ὁ βασιλεὺς ἐπὶ θρόνου τῆς βασιλείας αὐτοῦ, καὶ ἡσύχασεν ἡ γῆ ἐνώπιον αὐτοῦ.
\vs{53}Καὶ ἐψεύσατο πάντα ὅσα εἶπε, καὶ ἠλλοτριώθη τῷ Ἰωνάθαν, καὶ οὐκ ἀνταπέδωκε κατὰ τὰς εὐνοίας ἃς ἀνταπέδωκεν αὐτῷ, καὶ ἔθλιβεν αὐτὸν σφόδρα.

\vs{54}Μετὰ δὲ ταῦτα ἀπέστρεψε Τρύφων καὶ Ἀντίοχος μετʼ αὐτοῦ παιδάριον νεώτερον· καὶ ἐβασίλευσε καὶ ἐπέθετο διάδημα.
\vs{55}Καὶ ἐπισυνήχθησαν πρὸς αὐτὸν πᾶσαι αἱ δυνάμεις ἃς ἀπεσκόρπισε Δημήτριος, καὶ ἐπολέμησαν πρὸς αὐτὸν, καὶ ἔφυγε καὶ ἐτροπώθη.
\vs{56}Καὶ ἔλαβε Τρύφων τὰ θηρία, καὶ κατεκράτησεν Ἀντιοχείας

\vs{57}Καὶ ἔγραψεν Ἀντίοχος ὁ νεώτερος τῷ Ἰωνάθαν, λέγων. ἵστημί σοι τὴν ἀρχιερωσύνην, καὶ καθίστημί σε ἐπὶ τῶν τεσσάρων νομῶν, καὶ εἶναί σε τῶν φίλων τοῦ βασιλέως.
\vs{58}Καὶ ἀπέστειλεν αὐτῷ χρυσώματα καὶ διακονίαν, καὶ ἔδωκεν αὐτῷ ἐξουσίαν πίνειν ἐν χρυσώμασι, καὶ εἶναι ἐν πορφύρᾳ, καὶ ἔχειν πόρπην χρυσῆν.
\vs{59}Καὶ Σίμωνα τὸν ἀδελφὸν αὐτοῦ κατέστησε στρατηγὸν ἀπὸ τῆς κλίμακος Τύρου ἕως τῶν ὁρίων Αἰγύπτου.

\vs{60}Καὶ ἐξῆλθεν Ἰωνάθαν, καὶ διεπορεύετο πέραν τοῦ ποταμοῦ, καὶ ἐν ταῖς πόλεσι, καὶ ἠθροίσθησαν πρὸς αὐτὸν πᾶσαι αἱ δυνάμεις Συρίας εἰς συμμαχίαν, καὶ ἦλθεν εἰς Ἀσκάλωνα, καὶ ἀπήντησαν αὐτῷ οἱ ἐκ τῆς πόλεως ἐνδόξως.

\vs{61}Καὶ ἀπῆλθεν ἐκεῖθεν εἰς Γάζαν, καὶ ἀπέκλεισαν οἱ ἀπὸ Γάζης, καὶ περιεκάθισε περὶ αὐτὴν, καὶ ἐνεπύρισε τὰ περιπόλια αὐτῆς πυρὶ, καὶ ἐσκύλευσεν αὐτά.
\vs{62}Καὶ ἠξίωσαν οἱ ἀπὸ Γάζης τὸν Ἰωνάθαν, καὶ ἔδωκεν αὐτοῖς δεξιάς, καὶ ἔλαβε τοὺς υἱοὺς ἀρχόντων αὐτῶν εἰς ὅμηρα, καὶ ἐξαπέστειλεν αὐτοὺς εἰς Ἱερουσαλὴμ, καὶ διῆλθε τὴν χώραν ἕως Δαμασκοῦ.

\vs{63}Καὶ ἤκουσεν Ἰωνάθαν ὅτι παρῆσαν οἱ ἄρχοντες Δημητρίου εἰς Κάδης τὴν ἐν τῇ Γαλιλαίᾳ, μετὰ δυνάμεως πολλῆς, βουλόμενοι μεταστῆσαι αὐτὸν τῆς χρείας.
\vs{64}Καὶ συνήντησεν αὐτοῖς, τὸν δὲ ἀδελφὸν αὐτοῦ Σίμωνα κατέλιπεν ἐν τῇ χώρᾳ.
\vs{65}Καὶ παρενέβαλε Σίμων ἐπὶ Βαιθσούρα, καὶ ἐπολέμει αὐτὴν ἡμέρας πολλὰς, καὶ συνέκλεισεν αὐτήν.
\vs{66}Καὶ ἠξίωσαν αὐτὸν τοῦ δεξιὰς λαβεῖν, καὶ ἔδωκεν αὐτοῖς, καὶ ἐξέβαλεν αὐτοὺς ἐκεῖθεν, καὶ κατελάβετο τὴν πόλιν, καὶ ἔθετο ἐπʼ αὐτῇ φρουράν.

\vs{67}Καὶ Ἰωνάθαν καὶ ἡ παρεμβολὴ αὐτοῦ παρενέβαλον ἐπὶ τὸ ὕδωρ Γεννησάρ, καὶ ὤρθρισαν τοπρωῒ εἰς τὸ πεδίον Νασώρ.
\vs{68}Καὶ ἰδοὺ παρεμβολὴ ἀλλοφύλων ἀπήντα αὐτῷ ἐν τῷ πεδίῳ, καὶ ἐξέβαλον ἔνεδρον ἐπʼ αὐτὸν ἐν τοῖς ὄρεσιν, αὐτοὶ δὲ ἀπήντησαν ἐξεναντίας.

\vs{69}Τὰ δὲ ἔνεδρα ἐξανέστησαν ἐκ τῶν τόπων αὐτῶν, καὶ συνῆψαν πόλεμον· καὶ ἔφυγον οἱ παρὰ Ἰωνάθαν πάντες,
\vs{70}οὐδὲ εἷς κατελείφθη ἀπʼ αὐτῶν, πλὴν Ματταθίας ὁ τοῦ Ἀβεσσαλώμου, καὶ Ἰούδας ὁ τοῦ Χαλφί, ἄρχοντες τῆς στρατιᾶς τῶν δυνάμεων.

\vs{71}Καὶ διέῤῥηξεν Ἰωνάθαν τὰ ἱμάτια αὐτοῦ, καὶ ἐπέθηκε γῆν ἐπὶ τὴν κεφαλὴν αὐτοῦ, καὶ προσηύξατο.
\vs{72}Καὶ ὑπέστρεψε πρὸς αὐτοὺς πολέμῳ, καὶ ἐτροπώσατο αὐτούς, καὶ ἔφυγον.
\vs{73}Καὶ ἴδον οἱ φεύγοντες οἱ παρʼ αὐτοῦ, καὶ ἐπέστρεψαν πρὸς αὐτόν, καὶ ἐδίωκον μετʼ αὐτοῦ ἕως Κάδης ἕως της παρεμβολῆς αὐτῶν, καὶ παρενέβαλον ἐκεῖ.

\vs{74}Καὶ ἔπεσον ἐκ τῶν ἀλλοφύλων ἐν τῇ ἡμέρα ἐκείνῃ εἰς ἄνδρας τρισχιλίους· καὶ ἐπέστρεψεν Ἰωνάθαν εἰς Ἱερουσαλήμ.

\ch{12}
Καὶ εἶδεν Ἰωνάθαν ὅτι ὁ καιρὸς αὐτῷ συνεργεῖ, καὶ ἐπέλεξεν ἄνδρας, καὶ ἀπέστειλεν εἰς Ῥώμην, στῆσαι καὶ ἀνανεώσασθαι τὴν πρὸς αὐτοὺς φιλίαν.
\vs{2}Καὶ πρὸς Σπαρτιάτας, καὶ τόπους ἑτέρους ἀπέστειλεν ἐπιστολὰς κατὰ τὰ αὐτά.

\vs{3}Καὶ ἐπορεύθησαν εἰς Ῥώμην, καὶ εἰσῆλθον εἰς τὸ βουλευτήριον, καὶ εἶπον, Ἰωνάθαν ὁ ἀρχιερεὺς καὶ τὸ ἔθνος τῶν Ἰουδαίων ἀπέστειλεν ἡμᾶς ἀνανεώσασθαι τὴν φιλίαν αὐτοῖς, καὶ τὴν συμμαχίαν κατὰ τὸ πρότερον.
\vs{4}Καὶ ἔδωκαν ἐπιστολὰς αὐτοῖς πρὸς αὐτοὺς κατὰ τόπον, ὅπως προπέμπωσιν αὐτοὺς εἰς γῆν Ἰούδα μετʼ εἰρήνης.
\vs{5}Καὶ τοῦτο τὸ ἀντίγραφον τῶν ἐπιστολῶν ὧν ἔγραψεν Ἰωνάθαν τοῖς Σπαρτιάταις·

\vs{6}Ἰωνάθαν ἀρχιερεὺς, καὶ ἡ γερουσία τοῦ ἔθνους, καὶ οἱ ἱερεῖς, καὶ ὁ λοιπὸς δῆμος τῶν Ἰουδαίων, Σπαρτιάταις τοῖς ἀδελφοῖς χαίρειν.

\vs{7}Ἔτι πρότερον ἀπεστάλησαν ἐπιστολαὶ πρὸς Ὀνίαν τὸν ἀρχιερέα παρὰ Δαρείου τοῦ βασιλεύοντος ἐν ὑμῖν, ὅτι ἐστὲ ἀδελφοὶ ἡμῶν, ὡς τὸ ἀντίγραφον ὑπόκειται.
\vs{8}Καὶ ἐπεδέξατο Ὀνίας τὸν ἄνδρα τὸν ἀπεσταλμένον ἐνδόξως, καὶ ἔλαβε τὰς ἐπιστολάς ἐν αἷς διεσαφεῖτο περὶ συμμαχίας καὶ φιλίας.

\vs{9}Καὶ ἡμεῖς οὖν ἀπροσδεεῖς τούτων ὄντες, παράκλησιν ἔχοντες τὰ βιβλία τὰ ἅγια τὰ ἐν ταῖς χερσὶν ἡμῶν,
\vs{10}ἐπειράθημεν ἀποστεῖλαι τὴν πρὸς ὑμᾶς ἀδελφότητα καὶ φιλίαν ἀνανεώσασθαι, πρὸς τὸ μὴ ἐξαλλοτριωθῆναι ὑμῶν· πολλοὶ γὰρ καιροὶ διῆλθον ἀφʼ οὗ ἀπεστείλατε πρὸς ἡμᾶς.

\vs{11}Ἡμεῖς οὖν ἐν παντὶ καιρῷ ἀδιαλείπτως ἔν τε ταῖς ἑορταῖς καὶ ταῖς λοιπαῖς καθηκούσαις ἡμέραις μιμνησκόμεθα ὑμων, ἐφʼ ὧν προσφέρομεν θυσιῶν, καὶ ἐν ταῖς προσευχαῖς, ὡς δέον ἐστὶ καὶ πρέπον μνημονεύειν ἀδελφῶν.
\vs{12}Εὐφραινόμεθα δὲ ἐπὶ τῇ δόξῃ ὑμῶν.

\vs{13}Ἡμᾶς δὲ ἐκύκλωσαν πολλαὶ θλίψεις, καὶ πόλεμοι πολλοὶ, καὶ ἐπολέμησαν ἡμᾶς οἱ βασιλεῖς οἱ κύκλῳ ἡμῶν.
\vs{14}Καὶ οὐκ ἠβουλόμεθα οὖν παρενοχλεῖν ὑμῖν, καὶ τοῖς λοιποῖς συμμάχοις, καὶ φίλοις ἡμῶν, ἐν τοῖς πολέμοις τούτοις.
\vs{15}Ἔχομεν γὰρ τὴν ἐξ οὐρανοῦ βοήθειαν βοηθοῦσαν ἡμῖν, καὶ ἐῤῥύσθημεν ἀπὸ τῶν ἐχθρῶν ἡμῶν, καὶ ἐταπεινώθησαν οἱ ἐχθροὶ ἡμῶν.
\vs{16}Ἐπελέξαμεν οὖν Νουμήνιον Ἀντιόχου καὶ Ἀντίπατρον Ἰάσωνος, καὶ ἀπεστάλκαμεν πρὸς Ῥωμαίους ἀνενεώσασθαι τὴν πρὸς αὐτοὺς φιλίαν καὶ συμμαχίαν τὴν προτέραν.
\vs{17}Ἐνετειλάμεθα οὖν αὐτοῖς καὶ πρὸς ὑμᾶς πορευθῆναι, καὶ ἀσπάσασθαι ὑμᾶς, καὶ ἀποδοῦναι ὑμῖν τὰς παρʼ ἡμῶν ἐπιστολὰς περὶ τῆς ἀνανεώσεως καὶ τῆς ἀδελφότητος ἡμῶν.
\vs{18}Καὶ νῦν καλῶς ποιήσετε ἀντιφωνήσαντες ἡμῖν πρὸς ταῦτα.

\vs{19}Καὶ τοῦτο τὸ ἀντίγραφον τῶν ἐπιστολῶν ὧν ἀπέστειλεν. Ὀνιάρης
\vs{20}βασιλεὺς Σπαρτιατῶν Ὀνίᾳ ἱερεῖ μεγάλῳ χαίρειν.

\vs{21}Εὑρέθη ἐν γραφῇ περί τε τῶν Σπαρτιατῶν καὶ Ἰουδαίων ὅτι εἰσὶν ἀδελφοὶ, καὶ ὅτι εἰσὶν ἐκ γένους Ἁβραάμ.
\vs{22}Καὶ νῦν ἀφʼ οὗ ἔγνωμεν ταῦτα, καλῶς ποιήσετε γράφοντες ἡμῖν περὶ τῆς εἰρήνης ὑμῶν.
\vs{23}Καὶ ἡμεῖς δὲ ἁντιγράφομεν ὑμῖν, τὰ κτήνη ὑμῶν καὶ ἡ ὕπαρξις ὑμῶν ἡμῖν ἐστι, καὶ τὰ ἡμῶν ὑμῖν ἐστιν· ἐντελλόμεθα οὖν ὅπως ἀπαγγείλωσιν ὑμῖν κατὰ ταῦτα.

\vs{24}Καὶ ἤκουσεν Ἰωνάθαν ὅτι ἐπέστρεψαν οἱ ἄρχοντες Δημητρίου μετὰ δυνάμεως πολλῆς ὑπὲρ τὸ πρότερον τοῦ πολεμῆσαι πρὸς αὐτόν.
\vs{25}Καὶ ἀπῇρεν ἐξ Ἱερουσαλὴμ, καὶ ἀπήντησεν αὐτοῖς εἰς τὴν Ἀμαθῖτιν χώραν· οὐ γὰρ ἔδωκεν αὐτοῖς ἀνοχὴν ἐμβατεῦσαι εἰς τὴν χώραν αὐτοῦ.

\vs{26}Καὶ ἀπέστειλε κατασκόπους εἰς τὴς παρεμβολὴν αὐτῶν, καὶ ἀπέστρεψαν, καὶ ἀπήγγειλαν αὐτῷ, ὅτι οὕτω τάσσονται ἐπιπεσεῖν ἐπʼ αὐτοὺς τὴν νύκτα.
\vs{27}Ὡς δὲ ἔδυ ὁ ἥλιος, ἐπέταζεν Ἰωνάθαν τοῖς παρʼ αὐτοῦ γρηγορεῖν, καὶ εἶναι ἐπὶ τοῖς ὅπλοις, καὶ ἑτοιμάζεσθαι εἰς πόλεμον διʼ ὅλης τῆς νυκτὸς, καὶ ἐξέβαλε προφύλακας κύκλῳ τῆς παρεμβολῆς.

\vs{28}Καὶ ἤκουσαν οἱ ὑπεναντίοι ὅτι ἡτοίμασται Ἰωνάθαν καὶ οἱ παρʼ αὐτοῦ εἰς πόλεμον, καὶ ἐφοβήθησαν καὶ ἔπτηξαν τῇ καρδίᾳ αὐτῶν, καὶ ἀνέκαυσαν πυρὰς ἐν τῇ παρεμβολῇ αὐτῶν.
\vs{29}Ἰωνάθαν δὲ καὶ οἱ παρʼ αὐτοῦ οὐκ ἔγνωσαν ἕως πρωῒ, ἔβλεπον γὰρ τὰ φῶτα καιόμενα.
\vs{30}Καὶ κατεδίωξεν Ἰωνάθαν ὀπίσω αὐτῶν, καὶ οὐ κατέλαβεν αὐτούς, διέβησαν γὰρ τὸν Ἐλεύθερον ποταμόν.
\vs{31}Καὶ ἐξέκλινεν Ἰωνάθαν ἐπὶ τοὺς Ἄραβας τοὺς καλουμένους Ζαβεδαίους, καὶ ἐπάταξεν αὐτούς, καὶ ἔλαβε τὰ σκῦλα αὐτῶν.
\vs{32}Καὶ ἀναζεύξας ἦλθεν εἰς Δαμασκόν, καὶ διώδευσεν ἐν πάσῃ τῇ χώρᾳ.

\vs{33}Καὶ Σίμων ἐξῆλθε, καὶ διώδευσεν ἕως Ἀσκάλωνος, καὶ τῶν πλησίων ὀχυρωμάτων, καὶ ἐξέκλινεν εἰς Ἰόππην, καὶ προκατελάβετο αὐτήν.
\vs{34}Ἤκουσε γὰρ ὅτι βούλονται τὸ ὀχύρωμα παραδοῦναι τοῖς παρὰ Δημητρίου, καὶ ἔθετο ἐκεῖ φρουρὰν, ὅπως φυλάσσωσιν αὐτήν.

\vs{35}Καὶ ἐπέστρεψεν Ἰωνάθαν, καὶ ἐξεκκλησίασε τοὺς πρεσβυτέρους τοῦ λαοῦ, καὶ ἐβουλεύσατο μετʼ αὐτῶν τοῦ οἰκοδομῆσαι ὀξυρώματα ἐν τῇ Ἰουδαίᾳ,
\vs{36}καὶ προσυψῶσαι τὰ τείχη Ἱερουσαλὴμ, καὶ ὑψῶσαι ὕψος μέγα ἀναμέσον τῆς ἄκρας καὶ τῆς πόλεως, εἰς τὸ διαχωρίζειν αὐτὴν τῆς πόλεως, ἵνα ᾖ αὕτη κατὰ μόνας, ὅπως μήτε ἀγοράζωσι μήτε πωλῶσι.
\vs{37}Καὶ συνήχθησαν τοῦ οἰκοδομεῖν τὴν πόλιν, καὶ ἤγγισε τοῦ τείχους τοῦ χειμάῤῥου τοῦ ἐξ ἀπηλιώτου, καὶ ἐπεσκεύασαν τὸ καλούμενον Χαφεναθά.
\vs{38}καὶ Σίμων ᾠκοδόμησε τὴν Ἀδιδὰ ἐν τῇ Σεφήλᾳ, καὶ ὠχύρωσε θύρας καὶ μοχλούς.

\vs{39}Καὶ ἐζήτησε Τρύφων βασιλεῦσαι τῆς Ἀσίας, καὶ περιθέσθαι τὸ διάδημα, καὶ ἐκτεῖναι χεῖρα ἐπὶ Ἀντίοχον τὸν βασιλέα.
\vs{40}Καὶ ἐφοβήθη μήποτε οὐκ ἐάσῃ αὐτὸν Ἰωνάθαν, καὶ μήποτε πολεμήσῃ πρὸς αὐτὸν, καὶ ἐζήτει πόρον τοῦ συλλαβεῖν τὸν Ἰωνάθαν τοῦ ἀπολυσαι αὐτὸν, καὶ ἀπάρας ἦλθεν εἰς Βαιθσάν.

\vs{41}Καὶ ἐξῆλθεν Ἰωνάθαν εἰς ἀπάντησιν αὐτῷ ἐν τεσσαράκοντα χιλιάσιν ἀνδρῶν ἐπιλελεγμέναις εἰς παράταξιν, καὶ ἦλθεν εἰς Βαιθσάν.

\vs{42}Καὶ εἶδε Τρύφων ὅτι πάρεστιν Ἰωνάθαν μετὰ δυνάμεως πολλῆς, καὶ ἐκτεῖναι χεῖρας ἐπʼ αὐτὸν εὐλαβήθη.
\vs{43}Καὶ ἐπεδέξατο αὐτὸν ἐνδόξως, καὶ συνέστησεν αὐτὸν πᾶσι τοῖς φίλοις αὐτοῦ, καὶ ἔδωκεν αὐτῷ δόματα, καὶ ἐπέταξε ταῖς δυνάμεσιν αὐτοῦ ὑπακούειν αὐτῷ ὡς ἑαυτῷ.

\vs{44}Καὶ εἶπε τῷ Ἰωνάθαν, ἱνατί ἔκοψας πάντα τὸν λαὸν τοῦτον, πολέμου μὴ ἐνεστηκότος ἡμῖν;
\vs{45}Καὶ νῦν ἀπόστειλον αὐτοὺς εἰς τοὺς οἴκους αὐτῶν, ἐπίλεξαι δὲ σεαυτῷ ἄνδρας ὀλίγους οἵτινες ἔσονται μετὰ σοῦ, καὶ δεῦρο μετʼ ἐμοῦ εἰς Πτολεμαΐδα, καὶ παραδώσω σοι αὐτὴν καὶ τὰ λοιπὰ ὀχυρώματα καὶ τὰς δυνάμεις τὰς λοιπὰς καὶ πάντας τοὺς ἐπὶ τῶν χειρῶν, καὶ ἐπιστρέψας ἀπελεύσομαι, τούτου γὰρ χάριν πάρειμι.

\vs{46}Καὶ ἐμπιστεύσας αὐτῷ ἐποίησε καθὼς εἶπε, καὶ ἐξαπέστειλε τὰς δυνάμεις, καὶ ἀπῆλθον εἰς γῆν Ἰούδα.
\vs{47}Κατέλιπε δὲ μεθʼ ἑαυτοῦ ἄνδρας τρισχιλίους, ὧν δισχιλίους ἀφῆκεν ἐν τῇ Γαλιλαίᾳ, χίλιοι δὲ συνῆλθον αὐτῷ.

\vs{48}Ὡς δὲ εἰσῆλθεν Ἰωνάθαν εἰς Πτολεμαΐδα, ἀπέκλεισαν οἱ Πτολεμαεῖς τὰς πύλας, καὶ συνέλαβον αὐτὸν, καὶ πάντας τοὺς εἰσελθόντας μετʼ αὐτοῦ ἀπέκτειναν ἐν ῥομφαίᾳ.
\vs{49}Καὶ ἀπέστειλε Τρύφων δυνάμεις, καὶ ἵππον εἰς τὴν Γαλιλαίαν, καὶ τὸ πεδίον τὸ μέγα, τοῦ ἀπολέσαι πάντας τοὺς παρὰ Ἰωνάθαν.
\vs{50}Καὶ ἐπέγνωσαν ὅτι συνελήφθη Ἰωνάθαν καὶ ἀπόλωλε, καὶ οἱ μετʼ αὐτοῦ, καὶ παρεκάλεσαν ἑαυτοὺς, καὶ ἐπορεύοντο συνεστραμμένοι ἕτοιμοι εἰς πόλεμον.

\vs{51}Καὶ ἴδον οἱ διώκοντες ὅτι περὶ ψυχῆς αὐτοῖς ἐστι, καὶ ἐπέστρεψαν.
\vs{52}Καὶ ἦλθον πάντες μετʼ εἰρήνης εἰς γῆν Ἰούδα, καὶ ἐπένθησαν τὸν Ἰωνάθαν, καὶ τοὺς μετʼ αὐτοῦ, καὶ ἐφοβήθησαν σφόδρα, καὶ ἐπένθησε τᾶς Ἰσραὴλ πένθος μέγα.

\vs{53}Καὶ ἐζήτησαν πάντα τὰ ἔθνη τὰ κύκλῳ αὐτῶν ἐκτρίψαι αὐτούς· εἶπαν γάρ, οὐκ ἔχουσιν ἄρχοντα καὶ βοηθοῦντα· νῦν οὖν πολεμήσωμεν αὐτοὺς, καὶ ἐξάρωμεν ἐξ ἀνθρώπων τὸ μνημόσυνον αὐτῶν.

\ch{13}
Καὶ ἤκουσε Σίμων ὅτι συνήγαγε Τρύφων δύναμιν πολλὴν τοῦ ἐλθεῖν εἰς γῆν Ἰούδα, καὶ ἐκ ρίψαι αὐτήν.
\vs{2}Καὶ εἶδε τὸν λαὸν ὅτι ἐστὶν ἔντρομος, καὶ ἔμφοβος, καὶ ἀνέβη εἰς Ἱερουσαλὴμ, καὶ ἤθροισε τὸν λαόν.
\vs{3}Καὶ παρεκάλεσεν αὐτοὺς, καὶ εἶπεν αὐτοῖς, αὐτοὶ οἴδατε ὅσα ἐγὼ, καὶ οἱ ἀδελφοί μου, ὁ καὶ οἶκος τοῦ πατρός μου, ἐποιήσαμεν περὶ τῶν νόμων, καὶ τῶν ἁγίων, καὶ τοὺς πολέμους, καὶ τὰς στενοχωρίας ἃς εἴδομεν.
\vs{4}Τούτου χάριν ἀπώλοντο οἱ ἀδελφοί μου πάντες χάριν τοῦ Ἰσραὴλ, καὶ κατελείφθην ἐγὼ μόνος.

\vs{5}Καὶ νῦν μή μοι γένοιτο φείσασθαί μου τῆς ψυχῆς ἐν παντὶ καιρῷ θλίψεως, οὐ γάρ εἰμι κρείσσων τῶν ἀδελφῶν μου.
\vs{6}Πλὴν ἐκδικήσω περὶ τοῦ ἔθνους μου, καὶ περὶ τῶν ἁγίων, καὶ περὶ τῶν γυναικῶν καὶ τῶν τέκνων ἡμῶν, ὅτι συνήχθησαν πάντα τὰ ἔθνη ἐκτρίψαι ἡμᾶς ἔχθρας χάριν.

\vs{7}Καὶ ἀνεζωοπύρησε τὸ πνεῦμα τοῦ λαοῦ ἅμα τῷ ἀκοῦσαι τῶν λόγων τούτων,
\vs{8}καὶ ἀπεκρίθησαν φωνῇ μεγάλῃ, λέγοντες, σὺ εἶ ἡμῶν ἡγούμενος ἀντὶ Ἰούδα, καὶ Ἰωνάθαν τοῦ ἀδελφοῦ σου.
\vs{9}Πολέμησον τὸν πόλεμον ἡμῶν, καὶ πάντα ὅσα ἂν εἴπῃς ἡμῖν, ποιήσομεν.

\vs{10}Καὶ συνήγαγε πάντας τοὺς ἄνδρας τοὺς πολεμιστὰς, καὶ ἐτάχυνε τοῦ τελέσαι τὰ τείχη Ἱερουσαλήμ, καὶ ὠχύρωσεν αὐτὴν κυκλόθεν.
\vs{11}Καὶ ἀπέστειλεν Ἰωνάθαν τὸν τοῦ Ἀβεσσαλώμου καὶ μετʼ αὐτοῦ δύναμιν ἱκανὴν εἰς Ἰόππην, καὶ ἐξέβαλε τοὺς ὄντας ἐν αὐτῇ, καὶ ἔμεινεν ἐκεῖ ἐν αὐτῇ.

\vs{12}Καὶ ἀπῇρε Τρύφων ἀπὸ Πτολεμαΐδος μετὰ δυνάμεως πολλῆς εἰσελθεῖν εἰς γῆν Ἰούδα, καὶ Ἰωνάθαν μετʼ αὐτοῦ ἐν φυλακῇ.
\vs{13}Σίμων δὲ παρενέβαλεν ἐν Ἀδιδὰ κατὰ πρόσωπον τοῦ πεδίου.

\vs{14}Καὶ ἐπέγνω Τρύφων ὅτι ἀνέστη Σίμων ἀντὶ Ἰωνάθου τοῦ ἀδελφοῦ αὐτοῦ, καὶ ὅτι συνάπτειν αὐτῷ μέλλει πόλεμον, καὶ ἀπέστειλε πρὸς αὐτὸν πρέσβεις, λέγων,
\vs{15}περὶ ἀργυρίου οὗ ὤφειλεν Ἰωνάθαν ὁ ἀδελφός σου εἰς τὸ βασιλικὸν διʼ ἃς εἶχε χρείας συνέχομεν αὐτόν.
\vs{16}Καὶ νῦν ἀπόστειλον ἀργυρίου τάλαντα ἑκατὸν, καὶ δύο τῶν υἱῶν αὐτοῦ ὅμηρα, ὅπως μὴ ἀφεθεὶς ἀποστατήσῃ ἀφʼ ἡμῶν, καὶ ἀφήσομεν αὐτόν.

\vs{17}Καὶ ἔγνω Σίμων ὅτι δόλῳ λαλοῦσι πρὸς αὐτὸν, καὶ πέμπει τὸ ἀργύριον, καὶ τὰ παιδάρια, μήποτε ἔχθραν ἄρῃ μεγάλην πρός τὸν λαόν,
\vs{18}λέγων, ὅτι οὐκ ἀπέστειλα αὐτῷ τὸ ἀργύριον καὶ τὰ παιδάρια, καὶ ἀπώλετο.
\vs{19}Καὶ ἀπέστειλε τὰ παιδάρια, καὶ τὰ ἑκατὸν τάλαντα· καὶ διεψεύσατο, καὶ οὐκ ἀφῆκε τὸν Ἰωνάθαν.

\vs{20}Καὶ μετὰ ταῦτα ἦλθε Τρύφων τοῦ ἐμβατεῦσααι εἰς τὴν χώραν, καὶ ἐκτρίψαι αὐτὴν, καὶ ἐκύκλωσεν ὁδὸν τὴν εἰς Ἄδωρα· καὶ Σίμων καὶ ἡ παρεμβολὴ αὐτοῦ ἀντιπαρῆγεν αὐτῷ εἰς πάντα τόπον οὗ ἂν ἐπορεύετο.

\vs{21}Οἱ δὲ ἐκ τῆς ἄκρας ἀπέστελλον πρὸς Τρύφωνα πρεσβευτὰς κατασπεύδοντας αὐτὸν τοῦ ἐλθεῖν προὺς αὐτὸν διὰ τῆς ἐρήμου, καὶ· ἀποστεῖλαι αὐτοῖς τροφάς.
\vs{22}Καὶ ἡτοίμασε Τρύφων πᾶσαν τὴν ἵππον αὐτοῦ ἐλθεῖν ἐν τῇ νυκτὶ ἐκείνῃ· καὶ ἦν χιὼν πολλὴ σφόδρα, καὶ οὐκ ἦλθε διὰ τὴν χιόνα, καὶ ἀπῇρε, καὶ ἦλθεν εἰς τὴν Γαλααδίτιν.
\vs{23}Ὡς δὲ ἤγγισε τῇ Βασκαμᾷ, ἀπέκτεινε τὸν Ἰωνάθαν, καὶ ἐτάφη ἐκεῖ.
\vs{24}Καὶ ἐπέστρεψε Τρύφων, καὶ ἀπῆλθεν εἰς τὴν γῆν αὐτοῦ.

\vs{25}Καὶ ἀπέστειλε Σίμων, καὶ ἔλαβε τὰ ὀστᾶ Ἰωνάθαν τοῦ ἀδελφοῦ αὐτοῦ, καὶ ἔθαψεν αὐτὰ ἐν Μωδεῒν πόλει τῶν πατέρων αὐτοῦ.
\vs{26}Καὶ ἐκόψαντο αὐτὸν πᾶς Ἰσραὴλ κοπετὸν μέγαν, καὶ ἐπένθησαν αὐτὸν ἡμέρας πολλάς.

\vs{27}Καὶ ᾠκοδόμησε Σίμων ἐπὶ τὸν τάφον τοῦ πατρὸς αὐτοῦ καὶ τῶν ἀδελφῶν αὐτοῦ, καὶ ὕψωσεν αὐτὸν τῇ ὁράσει λίθῳ ξεστῷ ἐκ τῶν ὄπισθεν καὶ ἐκ τῶν ἔμπροσθεν.
\vs{28}Καὶ ἔστησεν ἐπʼ αὐτὰ ἑπτὰ πυραμίδας, μίαν κατέναντι τῆς μιᾶς, τῷ πατρὶ καὶ τῇ μητρὶ καὶ τῇ μητρὶ καὶ τοῖς τέσσαρσιν ἀδελφοῖς.
\vs{29}Καὶ ταύταις ἐποίησε μηχανήματα, περιθεὶς στύλους μεγάλους, καὶ ἐποίησεν ἐπὶ τοῖς στύλοις πανοπλίας εἰς ὄνομα αἰώνιον, καὶ παρὰ ταῖς πανοπλίαις πλοῖα ἐπιγεγλυμμένα, εἰς τὸ θεωρεῖσθαι ὑπὸ πάντων τῶν πλεόντων τὴν θάλασσαν.
\vs{30}Οὗτος ὁ τάφος ὃν ἐποίησεν ἐν Μωδεῒν, ἕως τῆς ἡμέρας ταύτης.

\vs{31}Ὁ δὲ Τρύφων ἐπορεύετο δόλῳ μετὰ Ἀντιόχου τοῦ βασιλέως τοῦ νεωτέρου, καὶ ἀπέκτεινεν αὐτὸν,
\vs{32}καὶ ἐβασίλευσεν ἀντʼ αὐτοῦ, καὶ περιέθετο διάδημα τῆς Ἀσίας, καὶ ἐποίησε πληγὴν μεγάλην ἐπὶ τῆς γῆς.

\vs{33}Καὶ ᾠκοδόμησε Σίμων τὰ ὀχυρώματα τῆς Ἰουδαίας, καὶ περιετείχισε πύργοις ὑψηλοῖς, καὶ τείχεσι μεγάλοις, καὶ πύλαις, καὶ μοχλοῖς, καὶ ἔθετο βρώματα ἐν τοῖς ὀχυρώμασι.
\vs{34}Καὶ ἐπέλεξε Σίμων ἄνδρας, καὶ ἀπέστειλε πρὸς Δημήτριον τὸν βασιλέα τοῦ ποιῆσαι ἄφεσιν τῇ χώρᾳ, ὅτι πᾶσαι αἱ πράξεις Τρύφωνος ἦσαν ἁρπαγαί.

\vs{35}Καὶ ἀπέστειλεν αὐτῷ Δημήτριος ὁ βασιλεὺς κατὰ τοὺς λόγους τούτους· καὶ ἀπεκρίθη αὐτῷ, καὶ ἔγραψεν αὐτῷ ἐπιστολὴν τοιαύτην·
\vs{36}Βασιλεὺς Δημήτριος Σίμονι ἀρχιερεῖ καὶ φίλῳ βασιλέων, καὶ πρεσβυτέροις, καὶ ἔθνει Ἰουδαίων χαίρειν.
\vs{37}Τὸν στέφανον τὸν χρυσοῦν, καὶ τὴν βαΐνην ἣν ἀπεστείλατε, κεκομίσμεθα, καὶ ἕτοιμοί ἐσμεν τοῦ ποιεῖν ὑμῖν εἰρήνην μεγάλην, καὶ γράφειν τοῖς ἐπὶ τῶν χρειῶν τοῦ ἀφιέναι ὑμῖν ἀφέματα.
\vs{38}Καὶ ὅσα ἑστήκαμεν πρὸς ὑμᾶς ἔστηκε, καὶ τὰ ὀχυρώματα ἃ ᾠκοδομήκατε ὑπαρχέτω ὑμῖν.
\vs{39}Ἀφίεμεν δὲ ἀγνοήματα καὶ τὰ ἁμαρτήματα ἕως τῆς σήμερον ἡμέρας, καὶ τὸν στέφανον ὃν ὠφείλετε, καὶ εἴ τι ἄλλο ἐτελωνεῖτο ἐν Ἱερουσαλήμ, μηκέτι τελωνείσθω.
\vs{40}Καὶ εἴ τινες ἐπιτήδειοι ὑμῶν γραφῆναι εἰς τοὺς περὶ ἡμᾶς, ἐγγραφέσθωσαν, καὶ γινέσθω ἀναμέσον ἡμῶν εἰρήνη.

\vs{41}Ἔτους ἑβδομηκοστοῦ καὶ ἑκατοστοῦ ᾔρθη ὁ ζυγὸς τῶν ἐθνῶν ἀπὸ τοῦ Ἰσραήλ.
\vs{42}Καὶ ἤρξατο ὁ λαὸς Ἰσραὴλ γράφειν ἐν ταῖς συγγραφαῖς καὶ συναλλάγμασιν, ἔτους πρώτου ἐπὶ Σίμωνος ἀρχιερέως μεγάλου καὶ στρατηγοῦ καὶ ἡγουμένου Ἰουδαίων.

\vs{43}Ἐν ταῖς ἡμέραις ἐκείναις παρενέβαλε Σίμων ἐπὶ Γάζαν, καὶ ἐκύκλωσεν αὐτὴν παρεμβολαῖς, καὶ ἐποίησεν ἑλεπόλεις καὶ προσήγαγε τῇ πόλει, καὶ ἐπάταξε πύργον ἕνα καὶ κατελάβετο.
\vs{44}Καὶ ἐξήλλοντο οἱ ἐν τῇ ἑλεπόλει εἰς τὴν πόλιν, καὶ ἐγένετο κίνημα μέγα ἐν τῇ πόλει.
\vs{45}Καὶ ἀνέβησαν οἱ ἐν τῇ πόλει σὺν ταῖς γυναιξὶ καὶ τοῖς τέκνοις ἐπὶ τὸ τεῖχος διεῤῥηχότες τὰ ἱμάτια αὐτῶν, καὶ ἐβόησαν φωνῇ μεγάλῃ ἀξιοῦντες Σίμωνα δεξιὰς αὐτοῖς δοῦναι,
\vs{46}καὶ εἶπον, μὴ ἡμῖν χρήσῃ κατὰ τὰς πονηρίας ἡμῶν, ἀλλὰ κατὰ τὸ ἔλεός σου.

\vs{47}Καὶ συνελύθη Σίμων αὐτοῖς, καὶ οὐκ ἐπολέμησεν αὐτούς· καὶ ἐξέβαλεν αὐτοὺς ἐκ τῆς πόλεως, καὶ ἐκαθάρισε τὰς οἰκίας ἐν αἷς ἦν τὰ εἴδωλα, καὶ οὕτως εἰσῆλθεν εἰς αὐτὴς ὑμνῶν καὶ εὐλογῶν.
\vs{48}Καὶ ἐξέβαλεν ἐξ αὐτῆς πᾶσαν ἀκαθαρσίαν, καὶ κατῴκισεν ἐκεῖ ἄνδρας οἵτινες τὸν νόμον ποιοῦσι, καὶ προσωχύρωσεν αὐτὴν, καὶ ᾠκοδόμησεν ἑαυτῷ ἑν αὐτῇ οἴκησιν.

\vs{49}Οἱ δὲ ἐκ τῆς ἄκρας ἐν Ἱερουσαλὴμ ἐκωλύοντο ἐκπορεύεσθαι καὶ εἰσπορεύεσθαι εἰς τὴν χώραν, καὶ ἀγοράζειν καὶ πωλεῖν, καὶ ἐπείνασαν σφόδρα, καὶ ἀπώλοντο ἐξ αὐτῶν ἱκανοὶ τῇ λιμῷ.
\vs{50}Καὶ ἐβόησαν πρὸς Σίμωνα δεξιὰς λαβεῖν, καὶ ἔδωκεν αὐτοῖς, καὶ ἐξέβαλεν αὐτοὺς ἐκεῖθεν, καὶ ἐκαθάρισε τὴν ἄκραν ἀπὸ τῶν μιασμάτων.
\vs{51}Καὶ εἰσῆλθεν εἰς αὐτὴν τῇ τρίτῃ καὶ εἰκάδι τοῦ δευτέρου μηνὸς ἔτους ἑνὸς καὶ ἑβδομηκοστοῦ καὶ ἑκατοστοῦ μετὰ αἰνέσεως καὶ βαΐων, καὶ ἐν κινύραις, καὶ ἐν κυμβάλοις, καὶ ἐν νάβλαις, καὶ ἐν ὕμνοις, καὶ ἐν ᾠδαῖς, ὅτι συνετρίβη ἐχθρὸς μέγας ἐξ Ἰσραήλ.

\vs{52}Καὶ ἔστησε κατʼ ἐνιαυτὸν τοῦ ἄγειν τὴν ἡμέραν ταύτην μετʼ εὐφροσύνης· καὶ προσωχύρωσε τὸ ὄρος τοῦ ἱεροῦ τὸ παρὰ τὴν ἄκραν, καὶ ᾤκει ἐκεῖ αὐτὸς καὶ οἱ παρʼ αὐτοῦ.
\vs{53}Καὶ εἶδε Σίμων τὸν Ἰωάννην υἱὸν αὐτοῦ, ὅτι ἀνήρ ἐστι, καὶ ἔθετο αὐτὸν ἡγούμενον τῶν δυνάμεων πασῶν, καὶ ᾤκει ἐν Γαζάροις.

\ch{14}
Καὶ ἐν ἔτει δευτέρῳ καὶ ἑβδομηκοστῷ καὶ ἑκατοστῷ συνήγαγε Δημήτριος ὁ βασιλεὺς τὰς δυνάμεις αὐτοῦ· καὶ ἐπορεύθη εἰς Μήδειαν τοῦ ἐπισπάσασθαι βοήθειαν αὐτῷ, ὅπως πολεμήσῃ τὸν Τρύφωνα.

\vs{2}Καὶ ἤκουσεν Ἀρσάκης ὁ βασιλεὺς τῆς Περσίδος καὶ Μηδειας ὅτι ἦλθε Δημήτριος εἰς τὰ ὅρια αὐτοῦ, καὶ ἀπέστειλεν ἕνα τῶν ἀρχόντων αὐτοῦ συλλάβεῖν αὐτὸν ζῶντα.
\vs{3}Καὶ ἐπορεύθη καὶ ἐπάταξε τὴν παρεμβολὴν Δημητρίου, καὶ συνέλαβεν αὐτὸν, καὶ ἤγαγεν αὐτὸν πρὸς Ἀρσάκην, καὶ ἔθετο αὐτὸν ἐν φυλακῇ.

\vs{4}Καὶ ἡσύχασεν ἡ γῆ Ἰούδα πάσας τὰς ἡμέρας Σίμωνος· καὶ ἐζήτησεν ἀγαθὰ τῷ ἔθνει αὐτοῦ, καὶ ἤρεσεν αὐτοῖς ἡ ἐξουσία αὐτοῦ καὶ ἡ δόξα αὐτοῦ πάσας τὰς ἡμέρας.
\vs{5}Καὶ μετὰ πάσης τῆς δόξης αὐτοῦ ἔλαβε τὴν Ἰόππην εἰς λιμένα, καὶ ἐποιήσεν εἴσοδον ταῖς νήσοις τῆς θαλάσσης.
\vs{6}Καὶ ἐπλάτυνε τὰ ὅρια τῷ ἔθνει αὐτοῦ, καὶ ἐκράτησε τῆς χώρας.
\vs{7}Καὶ συνήγαγεν αἰχμαλωσίαν πολλὴν, καὶ ἐκυρίευσε Γαζαρῶν καὶ Βαιθσούρων καὶ τῆς ἄκρας· καὶ ἐξῇρε τὰς ἀκαθαρσίας ἐξ αὐτῆς, καὶ οὐκ ἦν ὁ ἀντικείμενος αὐτῷ.

\vs{8}Καὶ ἦσαν γεωργοῦντες τὴν γῆν αὐτῶν μετʼ εἰρήνης, καὶ ἡ γῆ ἐδίδου τὰ γεννήματα αὐτῆς, καὶ τὰ ξύλα τῶν πεδίων τὸν καρπὸν αὐτῶν.
\vs{9}Πρεσβύτεροι ἐν ταῖς πλατείαις ἐκάθηντο, πάντες περὶ ἀγαθῶν ἐκοινολογοῦντο, καὶ οἱ νεανίσκοι ἐνεδύσαντο δόξας καὶ στολὰς πολέμου.
\vs{10}Ταῖς πόλεσιν ἐχορήγησε βρώματα, καὶ ἔταξεν αὐτὰς ἐν σκεύεσιν ὀχυρώσεως, ἕως ὅτου ὠνομάσθη τὸ ὄνομα τῆς δόξης αὐτοῦ ἕως ἄκρου τῆς γῆς.

\vs{11}Ἐποίησε τὴν εἰρήνην ἐπὶ τῆς γῆς, καὶ εὐφράνθη Ἰσραὴλ εὐφροσύνην μεγάλην.
\vs{12}Καὶ ἐκάθισεν ἕκαστος ὑπὸ τὴν ἄμπελον αὐτοῦ καὶ τὴν συκῆν αὐτοῦ, καὶ οὐκ ἦν ὁ ἐκφοβῶν αὐτούς.
\vs{13}Καὶ ἐξέλιπεν ὁ πολεμῶν αὐτοὺς ἐπὶ τῆς γῆς, καὶ οἱ βασιλεῖς συνετρίβησαν ἐν ταῖς ἡμέραις ἐκείναις.
\vs{14}Καὶ ἐστήρισε πάντας τοὺς ταπεινοὺς τοῦ λαοῦ αὐτοῦ· τὸν νόμον ἐξεζήτησε, καὶ ἐξῇρε πάντα ἄνομον καὶ πονηρόν.
\vs{15}Τὰ ἅγια ἐδόξασε, καὶ ἐπλήθυνε τὰ σκεύη τῶν ἁγίων.

\vs{16}Καὶ ἠκούσθη ἐν Ῥώμῃ ὅτι ἀπέθανεν Ἰωνάθαν, καὶ ἕως Σπάρτης, καὶ ἐλυπήθησαν σφόδρα.
\vs{17}Ὡς δὲ ἤκουσαν ὅτι Σίμων ὁ ἀδελφὸς αὐτοῦ γέγονεν ἀντʼ αὐτοῦ ἀρχιερεὺς, καὶ ἐπικρατεῖ τῆς χώρας καὶ τῶν πόλεων τῶν ἐν αὐτῇ.
\vs{18}Ἔγραψαν πρὸς αὐτὸν δέλτοις χαλκαῖς, τοῦ ἀνανεώσασθαι πρὸς αὐτὸν φιλίαν καὶ τὴν συμμαχίαν ἣν ἔστησαν πρὸς Ἰούδαν καὶ Ἰωνάθαν τοὺς ἀδελφοὺς αὐτοῦ.
\vs{19}Καὶ ἀνεγνώσθησαν ἐνώπιον τῆς ἐκκλησίας ἐν Ἱερουσαλήμ.

\vs{20}Καὶ τοῦτο τὸ ἀντίγραφον τῶν ἐπιστολῶν ὧν ἀπέστειλαν οἱ Σπαρτιάται· Σπαρτιατῶν ἄρχοντες καὶ ἡ πόλις Σίμωνι ἱερεῖ μεγάλῳ, καὶ τοῖς πρεσβυτέροις, καὶ τοῖς ἱερεῦσι, καὶ τῷ λοιπῷ δήμῳ τῶν Ἰουδαίων ἀδελφοῖς χαίρειν.
\vs{21}Οἱ πρεσβεύται οἱ ἀποσταλέντες πρὸς τὸν δῆμον ἡμῶν ἀπήγγειλαν ἡμῖν περὶ τῆς δόξης ὑμῶν καὶ τιμῆς, καὶ εὐφράνθημεν ἐπὶ τῆ ἐφόδῳ αὐτῶν.
\vs{22}Καὶ ἀνεγράψαμεν τὰ ὑπʼ αὐτῶν εἰρημένα ἐν ταῖς βουλαῖς τοῦ δήμου οὕτως, Νουμήνιος Ἀντιόχου καὶ Ἀντίπατρος Ἰάσωνος πρεσβευταὶ Ἰουδαίων ἤλθοσαν πρὸς ἡμᾶς ἀνανεούμενοι τὴν πρὸς ἡμᾶς φιλίαν.

\vs{23}Καὶ ἤρεσε τῷ δήμῳ ἐπιδέξασθαι τοὺς ἄνδρας ἐνδόξως, καὶ τοῦ θέσθαι τὸ ἀντίγραφον τῶν λόγων αὐτῶν ἐν τοῖς ἀποδεδειγμένοις τοῦ δήμου βιβλίοις, τοῦ ἔχειν μνημόσυνον τὸν δῆμον τῶν Σπαρτιατῶν· τὸ δὲ ἀντίγραφον τούτων ἐγράψαμεν Σίμωνι τῷ ἀρχιερεῖ.

\vs{24}Μετὰ ταῦτα ἀπέστειλε Σίμων τὸν Νουμήνιον εἰς Ῥώμην ἔχοντα ἀσπίδα χρυσῆν μεγάλην ὁλκῆς μνῶν χιλιων, εἰς τὸ στῆσαι πρὸς αὐτοὺς τὴν συμμαχίαν.
\vs{25}Ὡς δὲ ἤκουσεν ὁ δῆμος τῶν λόγων τούτων, εἶπον, τίνα χάριν ἀποδώσομεν Σίμωνι καὶ τοῖς υἱοῖς αὐτοῦ;
\vs{26}Ἐστήρισε γὰρ αὐτὸς καὶ οἱ ἀδελφοὶ αὐτοῦ, καὶ ὁ οἶκος τοῦ πατρὸς αὐτοῦ, καὶ ἐπολέμησαν τοὺς ἐχθροὺς Ἰσραὴλ ἀπʼ αὐτῶν, καὶ ἔστησαν αὐτῷ ἐλευθερίαν.

\vs{27}Καὶ κατέγραψαν ἐν δέλτοις χαλκαῖς, καὶ ἔθεντο ἐν στήλαις ἐν ὄρει Σιών· καὶ τοῦτο τὸ ἀντίγραφον τῆς γραφῆς· ὀκτωκαιδεκάτῃ Ἐλούλ, ἔτους δευτέρου καὶ ἑβδομηκοστοῦ καὶ ἑκατοστοῦ· καὶ τοῦτο τρίτον ἔτος ἐπὶ Σίμωνος ἀρχιερέως·
\vs{28}ἐν Σαραμὲλ, ἐπὶ συναγωγῆς μεγάλης ἱερέων, καὶ λαοῦ, καὶ ἀρχόντων ἔθνους, καὶ τῶν πρεσβυτέρων τῆς χώρας ἐγνώρισεν ἡμῖν.

\vs{29}Ἐπεὶ πολλάκις ἐγενήθησαν πόλεμοι ἐν τῇ χώρᾳ· Σίμων δὲ ὁ υἱὸς Ματταθίου ὁ υἱὸς τῶν υἱῶν Ἰαρὶβ καὶ οἱ ἀδελφοὶ αὐτοῦ ἔδωκαν ἑαυτοὺς τῷ κινδύνῳ, καὶ ἀντέστησαν τοῖς ὑπεναντίοις τοῦ ἔθνους αὐτῶν, ὅπως σταθῇ τὰ ἅγια αὐτῶν καὶ ὁ νόμος, καὶ δόξῃ μεγάλῃ ἐδόξασαν τὸ ἔθνος αὐτῶν·

\vs{30}Καὶ ἤθροισεν Ἰωνάθαν τὸ ἔθνος αὐτῶν, καὶ ἐγενήθη αὐτοῖς ἀρχιερεὺς, καὶ προσετέθη πρὸς τὸν λαὸν αὐτοῦ.
\vs{31}Καὶ ἐβουλήθησαν οἱ ἐχθροὶ αὐτῶν ἐμβατεῦσαι εἰς τὴν χώραν αὐτῶν, τοῦ ἐκτρίψαι τὴν χώραν αὐτῶν, καὶ ἐκτεῖναι χεῖρας ἐπὶ τὰ ἅγια αὐτῶν·
\vs{32}τότε ἀνέστη Σίμων, καὶ ἐπολέμησε περὶ τοῦ ἔθνους αὐτοῦ, καὶ ἐδαπάνησε χρήματα πολλὰ τῶν ἑαυτοῦ, καὶ ὡπλοδότησε τοὺς ἄνδρας τῆς δυνάμεως τοῦ ἔθνους αὐτοῦ, καὶ ἔδωκεν αὐτοῖς ὀψώνια,
\vs{33}καὶ ὠχύρωσε τὰς πόλεις τῆς Ἰουδαίας, καὶ τὴν Βαιθσούραν τὴν ἐπὶ τῶν ὁρίων τῆς Ἰουδαίας, οὗ ἦν τὰ ὅπλα τῶν πολεμίων τοπρότερον, καὶ ἔθετο ἐκεῖ φρουρὰν ἄνδρας Ἰουδαίους.
\vs{34}Καὶ Ἰόππην ὠχύρωσε τὴν ἐπὶ τῆς θαλάσσης, καὶ τὴν Γάζαρα τὴν ἐπὶ τῶν ὁρίων Ἀζώτου, ἐν ᾗ ᾤκουν οἱ πολέμιοι τοπρότερον ἐκεῖ, καὶ κατῴκισεν ἐκεῖ Ἰουδαίους, καὶ ὅσα ἐπιτήδεια ἦν πρὸς τὴν τούτων ἐπανόρθωσιν ἔθετο ἐν αὐτοῖς.

\vs{35}Καὶ εἶδεν ὁ λαὸς τὴς πρᾶξιν τοῦ Σίμωνος, καὶ τὴν δόξαν ἣν ἐβουλεύσατο ποιῆσαι τῷ ἔθνει αὐτοῦ, καὶ ἔθεντο αὐτὸν ἡγοὺμενον αὐτῶν καὶ ἀρχιερέα, διὰ τὸ αὐτὸν πεποιηκέναι πάντα ταῦτα, καὶ τὴν δικαιοσύνην, καὶ τὴν πίστιν ἣν συνετήρησε τῷ ἔθνει αὐτοῦ, καὶ ἐζήτησε παντὶ τρόπῳ ὑψῶσαι τὸν λαὸν αὐτοῦ.

\vs{36}Καὶ ἐν ταῖς ἡμέραις αὐτοῦ εὐωδώθη ἐν ταῖς χερσὶν αὐτοῦ, τοῦ ἐξαρθῆναι τὰ ἔθνη ἐκ τῆς χώρας αὐτῶν, καὶ τοὺς ἐν τῇ πόλει Δαυὶδ τοὺς ἐν Ἱερουσαλὴμ, οἳ ἐποίησαν ἑαυτοῖς ἄκραν, ἐξ ἧς ἐξεπορεύοντο καὶ ἐμίαινον κύκλῳ τῶν ἁγίων, καὶ ἐποίουν πληγὴν μεγάλην ἐν τῇ ἁγνείᾳ.
\vs{37}Καὶ κατῴκισεν ἐν αὐτῇ ἄνδρας Ἰουδαίους, καὶ ὠχύρωσεν αὐτὴν πρὸς ἀσφάλειαν τῆς χώρας καὶ τῆς πόλεως, καὶ ὕψωσε τὰ τείχη Ἱερουσαλήμ.

\vs{38}Καὶ ὁ βασιλεὺς Δημήτριος ἔστησεν αὐτῷ τὴν ἀρχιερωσύνην κατὰ ταῦτα,
\vs{39}καὶ ἐποίησεν αὐτὸν τῶν φίλων αὐτοῦ, καὶ ἐδόξασεν αὐτὸν δόξῃ μεγάλῃ.

\vs{40}Ἤκουσε γὰρ ὅτι προσηγόρευνται οἱ Ἰουδαῖοι ὑπὸ Ῥωμαίων φίλοι καὶ σύμμαχοι καὶ ἀδελφοὶ, καὶ ὅτι ἀπήντησαν τοῖς πρεσβευταῖς Σίμωνος ἐνδόξως·
\vs{41}καὶ ὅτι εὐδόκησαν οἱ Ἰουδαῖοι, καὶ οἱ ἱερεῖς, τοῦ εἶναι Σίμωνα ἡγουμένον καὶ ἀρχιερέα εἰς τὸν αἰῶνα, ἕως τοῦ ἀναστῆναι προφήτην πιστόν·
\vs{42}καὶ τοῦ εἶναι ἐπʼ αὐτῶν στρατηγόν, καὶ ὅπως μέλοι αὐτῷ περὶ τῶν ἁγίων καθιστάναι αὐτοὺς ἐπὶ τῶν ἔργων αὐτῶν καὶ ἐπὶ τῆς χώρας, καὶ ἐπὶ τῶν ὅπλων, καὶ ἐπὶ τῶν ὀχυρωμάτων·
\vs{43}καὶ ὅπως μέλοι αὐτῷ περὶ τῶν ἁγίων, καὶ ὅπως ἀκούηται ὑπὸ πάντων, καὶ ὅπως γράφωνται ἐπὶ τῷ ὀνόματι αὐτοῦ πᾶσαι συγγραφαὶ ἐν τῇ χώρᾳ, καὶ ὅπως περιβάληται πορφύραν, καὶ χρυσοφορῇ.

\vs{44}Καὶ οὐκ ἐξέσται οὐδενὶ τοῦ λαοῦ καἰ τῶν ἱερέων ἀθετῆσαί τι τούτων, καὶ ἀντειπεῖν τοῖς ὑπʼ αὐτοῦ ῥηθησομένοις, καὶ ἐπισυστρέψαι συστροφὴν ἐν τῇ χώρᾳ ἄνευ αὐτοῦ, καὶ περιβάλλεσθαι πορφύραν, καὶ ἐμπορποῦσθαι πόρπην χρυσῆν·
\vs{45}Ὃς δʼ ἂν παρὰ ταῦτα ποιήσῃ ἢ ἀθετήσῃ τι τούτων, ἔνοχος ἔσται.
\vs{46}Καὶ εὐδόκησε πᾶς ὁ λαὸς θέσθαι Σίμωνι, καὶ ποιῆσαι κατὰ τοὺς λόγους τούτους.
\vs{47}Καὶ ἐπεδέξατο Σίμων, καὶ εὐδόκησεν ἀρχιερατεύειν, καὶ εἶναι στρατηγὸς καὶ ἐθνάρχης τῶν Ἰουδαίων, καὶ ἱερέων, καὶ τοῦ προστατῆσαι πάντων.

\vs{48}Καὶ τὴν γραφὴν ταύτην εἶπον θέσθαι ἐν δέλτοις χαλκαῖς, καὶ στῆσαι αὐτὰς ἐν περιβόλῳ τῶν ἁγίων ἐν τόπῳ ἐπισήμῳ,
\vs{49}τὰ δὲ ἀντίγραφα αὐτῶν θέσθαι ἐν τῷ γαζοφυλακίῳ, ὅπως ἕχῃ Σίμων, καὶ οἱ υἱοὶ αὐτοῦ.

\ch{15}
Καὶ ἀπέστειλεν ὁ Ἀντίοχος υἱὸς Δημητρίου τοῦ βασιλέως ἐπιστολὰς ἀπὸ τῶν νήσων τῆς θαλάσσης Σίμωνι ἱερεῖ καὶ ἐθνάρχῃ τῶν Ἰουδαίων, καὶ παντὶ τῷ ἔθνει.
\vs{2}Καὶ ἦσαν περιέχουσαι τὸν τρόπον τοῦτον· βασιλεὺς Ἀντίοχος Σίμωνι ἱερεῖ μεγάλῳ, καὶ ἐθνάρχῃ, καὶ ἔθνει Ἰουδαίων χαίρειν.

\vs{3}Ἐπειδὴ ἄνδρες λοιμοὶ κατεκράτησαν τῆς βασιλείας τῶν πατέρων ἡμῶν, βούλομαι δὲ ἀντιποιήσασθαι τῆς βασιλείας, ὅπως ἀποκαταστήσω αὐτὴν ὡς ἦν πρότερον, ἐξενολόγησα δὲ πλῆθος δυνάμεων, καὶ κατεσκεύασα πλοῖα πολεμικά,
\vs{4}βούλομαι δὲ ἐκβῆναι κατὰ τὴν χώραν, ὅπως μετέλθω τοὺς κατεφθαρκότας τὴν χώραν ἡμῶν, καὶ τοὺς ἠρημωκότας πόλεις πολλὰς ἐν τῇ βασιλείᾳ·
\vs{5}νῦν οὗν ἵστημί σοι πάντα τὰ ἀφαιρέματα ἃ ἀφῆκάν σοι οἱ πρὸ ἐμοῦ βασιλεῖς, καὶ ὅσα ἄλλα δόματα ἀφῆκάν σοι.

\vs{6}Καὶ ἐπέτρεψά σοι ποιῆσαι κόμμα ἴδιον νόμισμα τῇ χώρᾳ σου,
\vs{7}Ἱερουσαλὴμ δὲ καὶ τὰ ἅγια εἶναι ἐλεύθερα· καὶ πάντα τὰ ὅπλα ὅσα κατεσκεύασας, καὶ τὰ ὀχυρώματα ἃ ᾠκοδόμησας, ὧν κρατεῖς, μενέτω σοι.
\vs{8}Καὶ πᾶν ὀφείλημα βασιλικὸν, καὶ τὰ ἐσόμενα βασιλικὰ, ἀπὸ τοῦ νῦν καὶ εἰς τὸν ἅπαντα χρόνον ἀφιέσθω σοι·
\vs{9}ὡς δʼ ἂν κρατήσωμεν τῆς βασιλείας ἡμῶν, δοξάσομέν σε, καὶ τὸ ἔθνος σου, καὶ τὸ ἱερὸν δόξῃ μεγάλῃ, ὥστε φανερὰν γενέσθαι τὴν δόξαν ὑμῶν ἐν πάσῃ τῇ γῇ.

\vs{10}Ἔτους τετάρτου καὶ ἑβδομηκοστοῦ καὶ ἑκατοστοῦ ἐξῆλθεν Ἀντίοχος εἰς τὴν γῆν πατέρων αὐτοῦ, καὶ συνῆλθον πρὸς αὐτὸν πᾶσαι αἱ δυνάσεις, ὥστε ὀλίγους εἶναι τσὺς καταλειφθέντας σὺν Τρύφωνι.

\vs{11}Καὶ ἐδίωξεν αὐτὸν Ἀντίοχος ὁ βασιλεύς, καὶ ἦλθε φεύγων εἰς Δωρᾶ τὴν ἐπὶ τῆς θαλάσσης.
\vs{12}Εἶδε γὰρ ὅτι συνῆκται ἐπʼ αὐτὸν τὰ κακὰ, καὶ ἀφῆκαν αὐτὸν αἱ δυνάμεις.

\vs{13}Καὶ παρενέβαλεν Ἀντιοχος ἐπὶ Δωρᾶ, καὶ σύν αὐτῷ δώδεκα μυριάδες ἀνδρῶν πολεμιστῶν, καὶ ὀκτακισχιλία ἵππος.
\vs{14}Καὶ ἐκύκλωσεν τὴν πόλιν, καὶ τὰ πλοῖα ἀπὸ θαλάσσης συνῆψαν, καὶ ἔθλιβε τὴν πόλιν ἀπὸ τῆς γῆς, καὶ τῆς θαλάσσης, καὶ οὐκ εἴασεν οὐδένα ἐκπορεύεσθαι καὶ εἰσπορεύεσθαι.

\vs{15}Καὶ ἦλθε Νουμήνιος, καὶ οἱ παρʼ αὐτοῦ, ἐκ Ῥώμης, ἔχοντες ἐπιστολὰς τοῖς βασιλεῦσι, καὶ ταῖς χώραις ἐν αἷς ἐγέγραπτο τάδε·

\vs{16}Λεύκιος ὕπατος Ῥωμαίων Πτολεμαίῳ βασιλεῖ χαίρειν.
\vs{17}Οἱ πρεσβευταὶ τῶν Ἰουδαίων ἦλθον πρὸς ἡμᾶς φίλοι ἡμῶν, καὶ σύμμαχοι, ἀνανεούμενοι τὴν ἐξ ἀρχῆς φιλίαν καὶ συμμαχίαν, ἀπεσταλμένοι ἀπὸ Σίμωνος τοῦ ἀρχιερέως, καὶ τοῦ δήμου τῶν Ἰουδαίων.
\vs{18}Ἤνεγκαν δὲ ἀσπίδα χρυσῆν ἀπὸ μνῶν χιλίων.
\vs{19}Ἤρεσεν οὖν ἡμῖν γράψαι τοῖς βασιλεῦσι, καὶ ταῖς χώραις, ὅπως μὴ ἐκζητήσωσιν αὐτοῖς, κακὰ καὶ μὴ πολεμήσωσιν αὐτοὺς, καὶ τὰς πόλεις αὐτῶν, καὶ τὴν χώραν αὐτῶν, καὶ ἵνα μὴ συμμαχήσωσι τοῖς πολεμοῦσιν αὐτούς.
\vs{20}Ἔδοξε δὲ ἡμῖν δέξασθαι τὴν ἀσπίδα παρʼ αὐτῶν.
\vs{21}Εἴ τινες οὖν λοιμοὶ διαπεφεύγασιν ἐκ τῆς χώρας αὐτῶν πρὸς ὑμᾶς, παράδοτε αὐτοὺς Σίμωνι τῷ ἀρχιερεῖ, ὅπως ἐκδικήσῃ ἐν αὐτοῖς κατὰ τὸν νόμον αὐτῶν.

\vs{22}Καὶ τὰ αὐτὰ ἔγραψε Δημητρίῳ τῷ βασιλεῖ, καὶ Ἀττάλῳ, Ἀριαράθῃ, καὶ Ἀρσάκῃ.
\vs{23}Καὶ εἰς πάσας τὰς χώρας, καὶ Σαμψάκῃ, καὶ Σπαρτιάταις, καὶ εἰς Δῆλον, καὶ εἰς Μύνδον, καὶ εἰς Σικυῶνα, καὶ εἰς τὴν Καρίαν, καὶ εἰς Σάμον, καὶ εἰς τὴν Παμφυλίαν, καὶ εἰς τὴν Λυκίαν, καὶ εἰς Ἀλικαρνασσὸν, καὶ εἰς Ῥόδον, καὶ εἰς Φασηλίδα, καὶ εἰς Κῶ, καὶ εἰς Σίδην, καὶ εἰς Ἄραδον, καὶ εἰς Γόρτυναν, καὶ Κνίδον, καὶ Κύπρον, καὶ Κυρήνην.
\vs{24}Τὸ δὲ ἀντίγραφον αὐτῶν ἔγραψαν Σίμωνι τῷ ἀρχιερεῖ.

\vs{25}Αντιόχος δὲ ὁ βασιλεὺς παρενέβαλεν ἐπὶ Δωρᾶ ἐν τῇ δευτέρᾳ, προσάγων διαπαντὸς αὐτῇ τὰς χεῖρας, καὶ μηχανὰς ποιούμενος, καὶ συνέκλεισε τὸν Τρύφωνα τοῦ μὴ εἰσπορεύεσθαι καὶ ἐκπορεύεσθαι.

\vs{26}Καὶ ἀπέστειλεν αὐτῷ Σίμων δισχιλίους ἄνδρας ἐκλεκτοὺς συμμαχῆσαι αὐτῷ, καὶ ἀργύριον καὶ χρυσίον, καὶ σκεύη ἱκανά.
\vs{27}Καὶ οὐκ ἠβούλετο αὐτὰ δέξασθαι, ἀλλʼ ἠθέτησε πάντα ὅσα συνέθετο αὐτῷ τοπρότερον, καὶ ἠλλοτριοῦτο αὐτῷ.

\vs{28}Καὶ ἀπέστειλε πρὸς αὐτὸν Ἀθηνόβιον ἕνα τῶν φίλων αὐτοῦ κοινολογησόμενον αὐτῷ λέγων, ὑμεῖς κατακρατεῖτε τῆς Ἰόππης καὶ Γαζάρων καὶ τῆς ἄκρας τῆς ἐν Ἱερουσαλὴμ, πόλεις τῆς βασιλείας μου.
\vs{29}Τὰ ὅρια αὐτῶν ἠρημώσατε, καὶ ἐποιήσατε πληγὴν μεγάλην ἐπὶ τῆς γῆς, καὶ ἐκυριεύσατε τόπων πολλῶν ἐν τῇ βασιλείᾳ μου.
\vs{30}Νῦν οὖν παράδοτε τὰς πόλεις ἃς κατελάβεσθε, καὶ τοὺς φόρους τῶν τόπων ὧν κατεκυριεύσατε ἐκτὸς τῶν ὁρίων τῆς Ἰουδαίας.
\vs{31}Εἰ δὲ μὴ, δότε ἀντʼ αὐτῶν πεντακόσια τάλαντα ἀργυρίου, καὶ τῆς καταφθορᾶς ἧς κατεφθάρκατε, καὶ τῶν φόρων τῶν πόλεων ἄλλα τάλαντα πεντακόσια· εἰ δὲ δὴ, παραγενόμενοι ἐκπολεμήσομεν ὑμᾶς.

\vs{32}Καὶ ἦλθεν Ἀθηνόβιος φίλος τοῦ βασιλέως εἰς Ἱερουσαλήμ, καὶ εἶδε τὴν δόξαν Σίμωνος, καὶ κυλικεῖον μετὰ χρυσωμάτων, καὶ ἀργυρωμάτων, καὶ παράστασιν ἱκανὴν, καὶ ἐξίστατο, καὶ ἀπήγγειλεν αὐτῷ τοὺς λόγους τοῦ βασιλέως.

\vs{33}Καὶ ἀποκριθεὶς Σίμων εἶπεν αὐτῷ, οὔτε γῆν ἀλλοτρίαν εἰλήφαμεν, οὔτε ἀλλοτρίων κεκρατήκαμεν, ἀλλὰ τῆς κληρονομίας τῶν πατέρων ἡμῶν, ὑπὸ δὲ ἐχθρῶν ἡμῶν ἔν τινι καιρῷ ἀκρίτως κατεκρατήθη.
\vs{34}Ἡμεῖς δὲ καιρὸν ἔχοντες ἀντεχόμεθα τῆς κληρονομίας τῶν πατέρων ἡμῶν.
\vs{35}Περὶ δὲ Ἰόππης καὶ Γαζάρων ὧν αἰτεῖς, αὗται ἐποίουν ἐν τῷ λαῷ πληγὴν μεγάλην κατὰ τὴν χώραν ἡμῶν, τούτων δώσομεν τάλαντα ἑκατόν·

\vs{36}Καὶ οὐκ ἀπεκρίθη αὐτῷ Ἀθηνόβιος λόγον. Ἀπέστρεψε δὲ μετὰ θυμοῦ πρὸς τὸν βασιλέα, καὶ ἀπήγγειλεν αὐτῷ τοὺς λόγους τούτους, καὶ τὴν δόξαν Σίμωνος, καὶ πάντα ὅσα εἶδε· καὶ ὠργίσθη ὁ βασιλεὺς ὀργῆν μεγάλην.
\vs{37}Τρύφων δὲ ἐμβὰς εἰς πλοῖον ἔφυγεν εἰς Ὀρθωσιάδα.

\vs{38}Καὶ κατέστησεν ὁ βασιλεὺς τὸν Κενδεβαῖον στρατηγὸν τῆς παραλίας, καὶ δυνάμεις πεζικὰς καὶ ἱππικὰς ἔδωκεν αὐτῷ.
\vs{39}Καὶ ἐνετείλατο αὐτῷ παρεμβαλεῖν κατὰ πρόσωπον τῆς Ἰουδαίας· καὶ ἐνετείλατο αὐτῷ οἰκοδομῆσαι τὴν Κεδρὼν, καὶ ὀχυρῶσαι τὰς πύλας, καὶ ὅπως πολεμήσῃ τὸν λαόν· ὁ δὲ βασιλεὺς ἐδίωκε τὸν Τρύφωνα.

\vs{40}Καὶ παρεγενήθη Κενδεβαῖος εἰς Ἰάμνειαν, καὶ ἤρξατο τοῦ ἐρεθίζειν τὸν λαὸν, καὶ ἐμβατεύειν εἰς τὴν Ἰουδαίαν, καὶ αἰχμαλωτίζειν τὸν λαὸν καὶ φονεύειν.
\vs{41}Καὶ ᾠκοδόμησε τὴν Κεδρών· καὶ ἔταξεν ἐκεῖ ἱππεῖς καὶ δυνάμεις, ὅπως ἐκπορευόμενοι ἐξοδεύωσι τὰς ὁδοὺς τῆς Ἰουδαίας, καθὰ συνεταξεν αὐτῷ ὁ βασιλεύς.

\ch{16}
Καὶ ἀνέβη Ἰωάννης ἐκ Γαζάρων, καὶ ἀπήγγειλε Σιμωνι τῷ πατρὶ αὐτοῦ ἃ συνετέλει Κενδεβαῖος.

\vs{2}Καὶ ἐκάλεσε Σίμων τοὺς δύο υἱοὺς αὐτοῦ τοὺς πρεσβυτέρους Ἰούδαν καὶ Ἰωάννην, καὶ εἶπεν αὐτοῖς, ἐγὼ καὶ οἱ ἀδελφοί μου, καὶ ὁ οἶκος τοῦ πατρός μου, ἐπολεμήσαμεν τοὺς πολεμίους Ἰσραὴλ ἀπὸ νεότητος ἕως τῆς σήμερον ἡμέρας, καὶ εὐωδώθη ἐν ταῖς χερσὶν ἡμῶν ῥύσασθαι τὸν Ἰσραὴλ πλεονάκις.
\vs{3}Νῦν δὲ γεγήρακα, καὶ ὑμεῖς δὲ ἐν τῷ ἐλέει ἱκανοί ἐστε ἐν τοῖς ἔτεσι· γίνεσθε ἀντʼ ἐμοῦ, καὶ τοῦ ἀδελφοῦ μου, καὶ ἐξελθόντες ὑπερμαχεῖτε ὑπὲρ τοῦ ἔθνους ἡμῶν, ἡ δὲ ἐκ τοῦ οὐρανοῦ βοήθεια ἔστω μεθʼ ὑμῶν.

\vs{4}Καὶ ἐπέλεξεν ἐκ τῆς χώρας εἴκοσι χιλιάδας ἀνδρῶν πολεμιστῶν, καὶ ἱππεῖς, καὶ ἐπορεύθωσαν ἐπὶ τὸν Κενδεβαῖον, καὶ ἐκοιμήθησαν ἐν Μωδεΐν.

\vs{5}Καὶ ἀναστάντες τοπρωῒ ἐπορεύοντο εἰς τὸ πεδίον, καὶ ἰδοὺ δύναμις πολλὴ εἰς συνάντησιν αὐτοῖς πεζικὴ, καὶ ἱππεῖς, καὶ ἦν χειμάῤῥους ἀναμέσον αὐτῶν.
\vs{6}Καὶ παρενέβαλε κατὰ πρόσωπον αὐτῶν αὐτὸς καὶ ὁ λαὸς αὐτοῦ· καὶ εἶδε τὸν λαὸν δειλούμενον διαπερᾶσαι τὸν χειμάῤῥουν, καὶ διεπέρασε πρῶτος, καὶ ἴδον αὐτὸν οἱ ἄνδρες, καὶ διεπέρασαν κατόπισθεν αὐτοῦ·
\vs{7}Καὶ διεῖλε τὸν λαὸν, καὶ τοὺς ἱππεῖς ἐν μέσῳ τῶν πέζῶν· ἡ δὲ ἵππος τῶν ὑπεναντίων πολλὴ σφόδρα.

\vs{8}Καὶ ἐσάλπισαν ταῖς ἱεραῖς σάλπιγξι, καὶ ἐτροπώθη Κενδεβαῖος καὶ ἡ παρεμβολὴ αὐτοῦ, καὶ ἔπεσον ἐξ αὐτῶν τραυματίαι πολλοί· οἱ δὲ καταλειφθέντες ἔφυγον εἰς τὸ ὀχύρωμα.

\vs{9}Τότε ἐτραυματίσθη Ἰούδας ὁ ἀδελφὸς Ἰωάννου· Ἰωάννης δὲ κατεδίωξεν αὐτοὺς ἕως ἦλθεν εἰς Κεδρών, ἣν ᾠκοδόμησεν
\vs{10}Καὶ ἔφυγον ἕως εἰς τοὺς πύργους τοὺς ἐν τοῖς ἀγροῖς Ἀζώτου, καὶ ἐνεπύρισεν αὐτὴν ἐν πυρί, καὶ ἔπεσον ἐξ αὐτῶν εἰς ἄνδρας δισχιλίους· καὶ ἀπέστρεψεν εἰς γῆν Ἰούδα μετʼ εἰρήνης.

\vs{11}Καὶ Πτολεμαῖος ὁ τοῦ Ἀβούβου ἦν καθεσταμένος στρατηγὸς εἰς τὸ πεδίον Ἱεριχὼ, καὶ ἔσχεν ἀργύριον καὶ χρυσίον πολύ·
\vs{12}ἦν γὰρ γαμβρὸς τοῦ ἀρχιερέως.
\vs{13}Καὶ ὑψώθη ἡ καρδία αὐτοῦ, καὶ ἠβουλήθη κατακρατῆσαι τῆς χώρας, καὶ ἐβουλεύετο δόλῳ κατὰ Σίμωνος, καὶ τῶν υἱῶν αὐτοῦ, ἆραι αὐτούς.

\vs{14}Σίμων δὲ ἦν ἐφοδεύων τὰς πόλεις τὰς ἐν τῇ χώρᾳ, καὶ φροντίζων τῆς ἐπιμελείας αὐτῶν, καὶ κατέβη εἰς Ἱεριχὼ αὐτὸς, καὶ Ματταθίας καὶ Ἰούδας οἱ υἱοὶ αὐτοῦ, ἐτους ἑβδόμου καὶ ἑβδομηκοστοῦ καὶ ἑκατοστοῦ, ἐν μηνὶ ἑνδεκάτῳ, οὗτος ὁ μὴν Σαβάτ.
\vs{15}Καὶ ὑπεδέξατο αὐτοὺς ὁ τοῦ Ἀβούβου εἰς τὸ ὀχυρωμάτιον τὸ καλούμενον Δὼκ, μετὰ δόλου, ὃ ᾠκοδόμησε, καὶ ἐποίησεν αὐτοῖς πότον μέγαν, καὶ ἐνέκρυψεν ἐκεῖ ἄνδρας.

\vs{16}Καὶ ὅτε ἐμεθύσθη Σίμων καὶ οἱ υἱοὶ αὐτοῦ, ἐξανέστη Πτολεμαῖος καὶ οἱ παρʼ αὐτοῦ, καὶ ἐλάβοσαν τὰ ὅπλα αὐτῶν, καὶ ἐπεισήλθοσαν τῷ Σίμωνι εἰς τὸ συμπόσιον, καὶ ἀπέκτειναν αὐτὸν καὶ τοὺς δύο υἱοὺς αὐτοῦ, καί τινας τῶν παιδαρίων αὐτοῦ.
\vs{17}Καὶ ἐποίησεν ἀθεσίαν μεγάλην, καὶ ἀπέδωκε κατὰ ἀντὶ ἀγαθῶν·

\vs{18}Καὶ ἔγραψε ταῦτα Πτολεμαῖος, καὶ ἀπέστειλε τῷ βασιλεῖ ὅπως ἀποστείλῃ αὐτῷ δυνάμεις εἰς βοήθειαν, καὶ παραδῷ αὐτῷ τὴν χώραν αὐτῶν, καὶ τὰς πόλεις.

\vs{19}Καὶ ἀπέστειλεν ἑτέρους εἰς Γάζαρα ἆραι τὸν Ἰωάννην, καὶ τοῖς χιλιάρχοις ἀπέστειλεν ἐπιστολὰς παραγενέσθαι πρὸς αὐτόν, ὅπως δῷ αὐτοῖς ἀργύριον καὶ χρυσίον καὶ δόματα.
\vs{20}Καὶ ἑτέρους ἀπέστειλε καταλαβέσθαι τὴν Ἱερουσαλὴμ, καὶ τὸ ὄρος τοῦ ἱεροῦ.

\vs{21}Καὶ προδραμών τις ἀπήγγειλεν Ἰωάννῃ εἰς Γάζαρα, ὅτι ἀπώλετο ὁ πατὴρ αὐτοῦ καὶ οἱ ἀδελφοὶ αὐτοῦ, καὶ ὅτι ἀπέσταλκε καὶ σὲ ἀποκτεῖναι.
\vs{22}Καὶ ἀκούσας ἐξέστη σφόδρα· καὶ συνέλαβε τοὺς ἄνδρας τοὺς ἐλθόντας ἀπολέσαι αὐτόν, καὶ ἀπέκτεινεν αὐτοὺς, ἐπέγνω γὰρ ὅτι ἐζήτουν αὐτὸν ἀπολέσαι.

\vs{23}Καὶ τὰ λοιπὰ τῶν λόγων Ἰωάννου, καὶ τῶν πολέμων αὐτοῦ, καὶ τῶν ἀνδραγαθιῶν αὐτοῦ ὧν ἠνδραγάθησε, καὶ τῆς οἰκοδομῆς τῶν τειχέων ὧν ᾠκοδόμησε, καὶ τῶν πράξεων αὐτοῦ,
\vs{24}ἰδοὺ ταῦτα γέγραπται ἐπὶ βιβλίῳ ἡμερῶν ἀρχιερωσύνης αὐτοῦ, ἀφʼ οὗ ἐγενήθη ἀρχιερεὺς μετὰ τὸν πατέρα αὐτοῦ.


\def\book{ΜΑΚΚΑΒΑΙΩΝ Βʹ}
\biblebook{ΜΑΚΚΑΒΑΙΩΝ Βʹ}


\lettrine[lines=2, loversize=0.2, nindent=0em, findent=.25em]{\textcolor{bookheadingcolor}{Τ}}{ΟΙΣ} ἀδελφοῖς τοῖς κατʼ Αἴγυπτον Ἰουδαίοις χαίρειν· οἱ ἀδελφοὶ οἱ ἐν Ἱεροσολύμοις Ἰουδαῖοι, καὶ οἱ ἐν τῇ χώρᾳ τῆς Ἰουδαίας, εἰρήνην ἀγαθήν.

\vs{2}Καὶ ἀγαθοποιήσαι ὑμῖν ὁ Θεὸς, καὶ μνησθείη τῆς διαθήκης αὐτοῦ τῆς πρὸς Ἁβραὰμ, καὶ Ἰσαὰκ, καὶ Ἰακὼβ τῶν δούλων αὐτοῦ τῶν πιστῶν.
\vs{3}Καὶ δῴη ὑμῖν καρδίαν πᾶσιν εἰς τὸ σέβεσθαι αὐτὸν, καὶ ποιεῖν αὐτοῦ τὰ θελήματα καρδίᾳ μεγάλῃ, καὶ ψυχῇ βουλομένῃ·
\vs{4}Καὶ διανοίξαι τὴν καρδίαν ὑμῶν ἐν τῷ νόμῳ αὐτοῦ, καὶ ἐν τοῖς προστάγμασι, καὶ εἰρήνην ποιήσαι,
\vs{5}καὶ ἐπακούσαι ὑμῶν τῶν δεήσεων, καὶ καταλλαγείη ὑμῖν, καὶ μὴ ὑμᾶς ἐγκαταλίποι ἐν καιρῷ πονηρῷ.
\vs{6}Καὶ νῦν ὧδέ ἐσμεν προσευχόμενοι περὶ ὑμῶν.

\vs{7}Βασιλεύοντος Δημητρίου ἔτους ἑκατοστοῦ ἑξηκοστοῦ ἐννατου, ἡμεῖς ὁ Ιουδαῖοι γεγραφήκαμεν ὑμῖν ἐν τῇ θλίψει, καὶ ἐν τῇ ἀκμῇ τῇ ἐπελθούσῃ ἡμῖν ἐν τοῖς ἔτεσι τούτοις, ἀφʼ οὗ ἀπέστη Ἰάσων καὶ οἱ μετʼ αὐτοῦ ἀπὸ τῆς ἁγίας γῆς, καὶ τῆς βασιλείας·
\vs{8}καὶ ἐνεπύρισαν τὸν πυλῶνα, καὶ ἐξέχεαν αἷμα ἀθῶον· καὶ ἐδεήθημεν τοῦ Κυρίου, καὶ εἰσηκούσθημεν, καὶ προσηνέγκαμεν θυσίαν, καὶ σεμίδαλιν, καὶ ἐξήψαμεν τοὺς λύχνους, καὶ προεθήκαμεν τοὺς ἄρτους.
\vs{9}Καὶ νῦν ἵνα ἄγητε τὰς ἡμέρας τῆς σκηνοπηγίας τοῦ Χασελεὺ μηνός.

\vs{10}Ἔτους ἑκατοστοῦ ὀγδοηκοστοῦ καὶ ὀγδόου οἱ ἐν Ἱεροσολύμοις, καὶ οἱ ἐν τῇ Ἰουδαίᾳ, καὶ ἡ γερουσία, καὶ Ἰούδας Ἀριστοβούλῳ διδασκάλῳ Πτολεμαίου τοῦ βασιλέως, ὄντι δὲ ἀπὸ τοῦ τῶν χριστῶν ἱερέων γένους, καὶ τοῖς ἐν Αἰγύπτῳ Ἰουδαίοις, χαίρειν καὶ ὑγιαίνειν.

\vs{11}Ἐκ μεγάλων κινδύνων ὑπὸ τοῦ Θεοῦ σεσωσμένοι, μεγάλως εὐχαριστοῦμεν αὐτῷ, ὡς ἂν πρὸς βασιλέα παρατασσόμενοι.
\vs{12}Αὐτὸς γὰρ ἐξέβρασε τοὺς παραταξαμένους ἐν τῇ ἁγίᾳ πόλει.

\vs{13}Εἰς γὰρ τὴν Περσίδα γενόμενος ὁ ἡγεμὼν, καὶ ἡ περὶ αὐτὸν ἀνυπόστατος δοκοῦσα εἶναι δύναμις, κατεκόπησαν ἐν τῷ τῆς Ναναίας ἱερῷ, παραλογισμῷ χρησαμένων τῶν περὶ τὴν Ναναίαν ἱερέων.
\vs{14}Ὡς γὰρ συνοικήσων αὐτῇ παρεγένετο εἰς τὸν τόπον ὅ, τε Ἀντίοχος, καὶ οἱ σὺν αὐτῷ φίλοι, χάριν τοῦ λαβεῖν τὰ χρήματα εἰς φερνῆς λόγον·
\vs{15}Καὶ προθέντων αὐτὰ τῶν ἱερέων τῆς Ναναίος, κᾀκείνου προσελθόντος μετʼ ὀλίγων εἰς τὸν περίβολον τοῦ τεμένους, συγκλείσαντες τὸ ἱερὸν, ὡς εἰσῆλθεν Ἀντίοχος,
\vs{16}ἀνοίξαντες τὴν τοῦ φατνώματος κρυπτὴν θύραν, βάλλοντες πέτρους συνεκεραύνωσαν τὸν ἡγεμόνα, καὶ μέλη ποιήσαντες, καὶ τὰς κεφαλὰς ἀφελόντες, τοῖς ἔξω παρέῤῥιψαν.

\vs{17}Κατὰ πάντα εὐλογητὸς ἡμῶν ὁ Θεὸς, ὃς παρέδωκε τοὺς ἀσεβήσαντας.

\vs{18}Μέλλοντες οὖν ἄγειν ἐν τῷ Χασελεὺ πέμπτῃ καὶ εἰκάδι τὸν καθαρισμὸν τοῦ ἱεροῦ, δεόν ἡγησάμεθα διασαφῆσαι ὑμῖν, ἵνα καὶ αὐτοὶ ἄγητε τῆς σκηνοπηγίας καὶ τοῦ πυρός, ὅτε Νεεμίας οἰκοδομήσας τό, τε ἱερὸν καὶ τὸ θυσιαστήριον, ἀνήνεγκε θυσίαν.
\vs{19}Καὶ γὰρ ὅτε εἰς τὴν Περσικὴν ἤγοντο οἱ πατέρες ἡμῶν, οἵ τότε εὐσεβεῖς ἱερεῖς λαβόντες ἀπὸ τοῦ πυρὸς τοῦ θυσιαστηρίου λαθραίως, κατέκρυψαν ἐν κοιλώματι φρέατος τάξιν ἔχοντος ἀνύδρον, ἐν ᾧ κατησφαλίσαντο, ὥστε πᾶσιν ἄγνωστον εἶναι τὸν τόπον.

\vs{20}Διελθόντων δὲ ἐτῶν ἱκανῶν, ὅτε ἔδοξεν τῷ Θεῷ, ἀποσταλεὶς Νεεμίας ὑπὸ τοῦ βασιλέως τῆς Περσίδος, τοὺς ἐκγόνους τῶν ἱερέων τῶν ἀποκρυψάντων ἔπεμψεν ἐπὶ τὸ πῦρ·
\vs{21}ὡς δὲ διεσάφησαν ἡμῖν μὴ εὑρηκέναι πῦρ, ἀλλὰ ὕδωρ παχύ, ἐκέλευσεν αὐτοὺς ἀποβάψαντας φέρειν· ὡς δὲ ἀνηνέχθη τὰ τῶν θυσιῶν, ἐκέλευσε τοὺς ἱερεῖς Νεεμίας ἐπιῤῥᾶναι τῷ ὕδατι τά τε ξύλα, καὶ τὰ ἐπικείμενα.
\vs{22}Ὡς δὲ ἐγένετο τοῦτο, καὶ χρόνος διῆλθεν ὅτε ἥλιος ἀνέλαμψε πρότερον ἐπινεφὴς ὤν, ἀνήφθη πυρὰ μεγάλη, ὥστε θαυμάσαι πάντας.

\vs{23}Προσευχὴν δὲ ἐποιήσαντο οἱ ἱερεῖς δαπανωμένης τῆς θυσίας, οἵ τε ἱερεῖς, καὶ πάντες, καταρχομένου Ἰωνάθου, τῶν δὲ λοιπῶν ἐπιφωνούντων, ὡς Νεεμίου.

\vs{24}Ἦν δὲ ἡ προσευχὴ τὸν τρόπον ἔχουσα τοῦτον· Κύριε Κύριε ὁ Θεὸς ὁ πάντων κτίστης, ὁ φοβερὸς, καὶ ἰσχυρὸς, καὶ δίκαιος, καὶ ἐλεήμων, ὁ μόνος βασιλεὺς καὶ χρηστὸς,
\vs{25}ὁ μόνος χορηγός, ὁ μόνος δίκαιος, καὶ παντοκράτωρ, καὶ αἰώνιος, ὁ διασώζων τὸν Ἰσραὴλ ἐκ παντος κακοῦ, ὁ ποιήσας τοὺς πατέρας ἐκλεκτούς, καὶ ἁγιάσας αὐτούς,
\vs{26}πρόσδεξαι τὴν θυσίαν ὑπὲρ παντὸς τοῦ λαοῦ σου Ἰσραὴλ, καὶ διαφύλαξον τὴν μερίδα σου καὶ καθαγίασον.
\vs{27}Ἐπισυνάγαγε τὴν διασπορὰν ἡμῶν, ἐλευθέρωσον τοὺς δουλεύοντας ἐν τοῖς ἔθνεσι, τοὺς ἐξουθενημένους καὶ βδελυκτοὺς ἔπιδε, καὶ γνώτωσαν τὰ ἔθνη ὅτι σὺ εἶ ὁ Θεὸς ἡμῶν.
\vs{28}Βασάνισον τοὺς καταδυναστεύοντας, καὶ ἐξυβρίζοντας ἐν ὑπερηφανίᾳ·
\vs{29}Καταφύτευσον τὸν λαόν σου εἰς τὸν τόπον τὸν ἅγιόν σου, καθὼς εἶπε Μωυσῆς.
\vs{30}Οἱ δὲ ἱερεῖς ἐπέψαλλον τοὺς ὕμνους.

\vs{31}Καθὼς δὲ ἀνηλώθη τὰ τῆς θυσίας, καὶ τὸ περιλειπόμενον ὕδωρ, ὁ Νεεμίας ἐκέλευσε λίθους μείζονας καταχεῖν.
\vs{32}Ὡς δὲ τοῦτο ἐγενήθη, φλὸξ ἀνήφθη· τοῦ δὲ ἀπὸ τοῦ θυσιαστηρίου ἀντιλάμψαντος φωτὸς ἐδαπανήθη.

\vs{33}Ὡς δὲ φανερὸν ἐγενήθη τὸ πρᾶγμα, καὶ διηγγέλη τῷ βασιλεῖ τῶν Περσῶν, ὅτι εἰς τὸν τόπον οὗ τὸ πῦρ ἄπέκρυψαν οἱ μεταχθέντες ἱερεῖς, τὸ ὕδωρ ἐφάνη, ἀφʼ οὗ καὶ οἱ περὶ τὸν Νεεμίαν ἥγνισαν τὰ τῆς θυσίας.
\vs{34}Περιφράξας δὲ ὁ βασιλεὺς ἱερὸν ἐποίησε, δοκιμάσας τὸ πρᾶγμα.

\vs{35}Καὶ οἷς ἐχαρίζετο ὁ βασιλεὺς πολλὰ διάφορα ἐλάμβανε καὶ μετεδίδου.
\vs{36}Προσηγόρευσαν δὲ οἱ περὶ τὸν Νεεμίαν τοῦτο Νέφθαρ, ὃ διερμηνεύεται Καθαρισμός· καλεῖται δὲ παρὰ τοῖς πολλοῖς Νεφθαεί.

\ch{2}
Εὑρίσκεται δὲ ἐν ταῖς ἀπογραφαῖς Ἱερεμίας ὁ προφήτης, ὅτι ἐκέλευσε τοῦ πυρὸς λαβεῖν τοὺς μεταγινομένους, ὡς σεσήμανται,
\vs{2}καὶ ὡς ἐνετείλατο τοῖς μεταγενομένοις ὁ προφήτης, δοὺς αὐτοῖς τὸν νόμον, ἵνα μὴ ἐπιλάθωνται τῶν προσταγμάτων τοῦ Κυρίου, καὶ ἵνα μὴ ἀποπλανηθῶσι ταῖς διανοίαις, βλέποντες ἀγάλματα χρυσᾶ καὶ ἀργυρᾶ, καὶ τὸν περὶ αὐτὰ κόσμον.
\vs{3}Καὶ ἕτερα τοιαῦτα λέγων, παρεκάλει μὴ ἀποστῆναι τὸν νόμον ἀπὸ τῆς καρδίας αὐτῶν.

\vs{4}Ἦν δὲ ἐν τῇ γραφῇ, ὡς τὴν σκηνὴν καὶ τὴν κιβωτὸν ἐκέλευσεν ὁ προφήτης, χρηματισμοῦ γενηθέντος, αὐτῷ συνακολουθεῖν, ὡς δὲ ἐξῆλθεν εἰς τὸ ὄρος οὗ ὁ Μωυσῆς ἀναβὰς ἐθεάσατο τὴν τοῦ Θεοῦ κληρονομίαν.
\vs{5}Καὶ ἐλθὼν ὁ Ἱερεμίας εὗρεν οἶκον ἀντρώδη, καὶ τὴν σκηνὴν, καὶ τὴν κιβωτὸν, καὶ τὸ θυσιαστήριον τοῦ θυμιάματος εἰσήνεγκεν ἐκεῖ, καὶ τὴν θύραν ἐνέφραξε.

\vs{6}Καὶ προσελθόντες τινὲς τῶν συνακολουθούντων ὥστε ἐπισημῄνασθαι τὴν ὁδὸν, καὶ οὐκ ἠδυνήθησαν εὑρεῖν.
\vs{7}Ὡς δὲ ὁ Ἱερεμίας ἔγνω, μεμψάμενος αὐτοῖς εἶπεν, ὅτι καὶ ἄγνωστος ὁ τόπος ἔσται ἕως ἂν συναγάγῃ ὁ Θεὸς ἐπισυναγωγὴν τοῦ λαοῦ, καὶ ἵλεως γένηται.
\vs{8}Καὶ τότε ὁ Κύριος ἀναδείξει ταῦτα, καὶ ὀφθήσεται ἡ δόξα τοῦ Κυρίου καὶ ἡ νεφέλη, ὡς καὶ ἐπὶ Μωυσῇ ἐδηλοῦτο, ὡς καὶ ὁ Σαλωμὼν ἠξίωσεν ἵνα ὁ τόπος καθαγιασθῇ μεγάλως.

\vs{9}Διεσαφεῖτο δὲ καὶ ὡς σοφίαν ἔχων ἀνήνεγκε θυσίαν ἐγκαινισμοῦ, καὶ τῆς τελειώσεως τοῦ ἱεροῦ.
\vs{10}Καθὼς καὶ Μωυσῆς προσηύξατο πρὸς Κύριον, καὶ κατέβη πῦρ ἐκ τοῦ οὐρανοῦ, καὶ τὰ τῆς θυσίας ἐδαπάνησεν· οὕτως καὶ Σαλωμὼν προσηύξατο, καὶ καταβὰν τὸ πῦρ ἀνήλωσε τὰ ὁλοκαυτώματα.
\vs{11}Καὶ εἶπε Μωυσῆς, διὰ τὸ μὴ βεβρῶσθαι τὸ περὶ τῆς ἁμαρτίας, ἀνηλώθη.
\vs{12}Ὡσαύτως καὶ ὁ Σαλωμὼν τὰς ὀκτὼ ἡμέρας ἤγαγεν.

\vs{13}Ἐξηγοῦντο δὲ καὶ ἐν ταῖς ἀναγραφαῖς, καὶ ἐν τοῖς ὑπομνηματισμοῖς τοῖς κατὰ τὸν Νεεμίαν τὰ αὐτά, καὶ ὡς καταβαλλόμενος βιβλιοθήκην, ἐπισυνήγαγε τὰ περὶ τῶν βασιλέων καὶ προφητῶν, καὶ τὰ τοῦ Δαυεὶδ, καὶ ἐπιστολὰς βασιλέων περὶ ἀναθημάτων.
\vs{14}Ὡσαύτως δὲ καὶ Ἰούδας, τὰ διαπεπτωκότα διὰ τὸν πόλεμον τὸν γεγονότα ἡμῖν ἐπισυνήγαγε πάντα, καὶ ἔστι παρʼ ἡμῖν.
\vs{15}Ὧν οὖν ἐὰν χρείαν ἔχητε, τοὺς ἀποκομιοῦντας ὑμῖν ἀποστέλλετε.

\vs{16}Μέλλοντες οὖν ἄγειν τὸν καθαρισμὸν, ἐγράψαμεν ὑμῖν· καλῶς οὖν ποιήσετε ἄγοντες τὰς ἡμέρας.
\vs{17}Ὁ δὲ Θεὸς ὁ σώσας τὸν πάντα λαὸν αὐτοῦ, καὶ ἀποδοὺς τὴν κληρονομίαν πᾶσι, καὶ τὸ βασίλειον, καὶ τὸ ἱεράτευμα, καὶ τὸν ἁγιασμὸν.
\vs{18}Καθὼς ἐπηγγείλατο διὰ τοῦ νόμου ἐλπίζομεν γὰρ ἐπὶ τῷ Θεῷ, ὅτι ταχέως ἡμᾶς ἐλεήσει, καὶ ἐπισυνάξει ἐκ τῆς ὑπὸ τὸν οὐρανὸν εἰς τὸν ἅγιον τόπον· ἐξείλετο γὰρ ἡμᾶς ἐκ μεγάλων κακῶν, καὶ τὸν τόπον ἐκαθάρισε.

\vs{19}Τὰ δὲ κατὰ τὸν Ἰούδαν τὸν Μακκαβαῖον, καὶ τοὺς τούτου ἀδελφοὺς, καὶ τὸν τοῦ ἱεροῦ τοῦ μεγάλου καθαρισμὸν, καὶ τὸν τοῦ βωμοῦ ἐγκαινισμὸν,
\vs{20}ἔτι τε τοὺς πρὸς Ἀντιοχον τὸν Ἐπιφανῆ, καὶ τὸν τούτου υἱὸν Εὐπάτορα πολέμους,
\vs{21}καὶ τὰς ἐξ οὐρανοῦ γενομένας ἐπιφανείας τοῖς ὑπὲρ τοῦ Ἰουδαϊσμοῦ φιλοτίμως ἀνδραγαθήσασιν, ὡστε τὴν ὅλην χώραν ὀλίγους ὄντας λεηλατεῖν, καὶ τὰ βάρβαρα πλήθη διώκειν.
\vs{22}Καὶ τὸ περιβόητον καθʼ ὅλην τὴν οἰκουμένην ἱερὸν ἀνακομίσασθαι, καὶ τὴν πόλιν ἐλευθερῶσαι, καὶ τοὺς μέλλοντας καταλύεσθαι νόμους ἐπανορθῶσαι, τοῦ Κυρίου μετὰ πάσης ἐπιεικείας ἵλεω γενομένου αὐτοῖς,
\vs{23}τὰ ὑπὸ Ἰάσωνος τοῦ Κυρηναίου δεδηλωμένα διὰ πέντε βιβλίων, πειρασόμεθα διʼ ἑνὸς συντάγματος ἐπιτεμεῖν.

\vs{24}Συνορῶντες γὰρ τὸ χύμα τῶν ἀριθμῶν, καὶ τὴν οὖσαν δυσχέρειαν τοῖς θέλουσιν εἰσκυκλεῖσθαι τοῖς τῆς ἱστορίας διηγήμασι διὰ τὸ πλῆθος τῆς ὕλης,
\vs{25}ἐφροντίσαμεν τοῖς μὲν βουλομένοις ἀναγινώσκειν ψυχαγωγίαν, τοῖς δέ φιλοφρονοῦσιν εἰς τὸ διὰ μνήμης ἀναλαβεῖν εὐκοπίαν, πᾶσι δὲ τοῖς ἐντυγχάνουσιν ὠφέλειαν.

\vs{26}Καὶ ἡμῖν μὲν τοῖς τὴν κακοπάθειαν ἐπιδεδεγμένοις τῆς ἐπιτομῆς οὐ ῥᾴδιον, ἱδρῶτος δὲ καὶ ἀγρυπνίας τὸ πρᾶγμα·
\vs{27}καθάπερ τῷ παρασκευάζοντι συμπόσιον, καὶ ζητοῦντι τὴν ἑτέρων λυσιτέλειαν οὐκ εὐχερές μὲν, ὅμως διὰ τὴν τῶν πολλῶν εὐχαριστίαν, ἡδέως τὴν κακοπάθειαν ὑποίσομεν,
\vs{28}τὸ μὲν διακριβοῦν περὶ ἑκάστων τῷ συγγραφεῖ παραχωρήσαντες, τὸ δὲ ἐπιπορεύεσθαι τοῖς ὑπογραμμοῖς τῆς ἐπιτομῆς διαπονοῦντες.
\vs{29}Καθάπερ γὰρ τῆς καινῆς οἰκίας ἀρχιτέκτονι τῆς ὅλης καταβολῆς φροντιστέον, τῷ δὲ ἐγκαίειν καὶ ζωγραφεῖν ἐπιχειροῦντι, τὰ ἐπιτήδεια πρὸς διακόσμησιν ἐξεταστέον· οὕτω δοκῶ καὶ ἐπὶ ἡμῖν.
\vs{30}Τὸ μὲν ἐμβατεύειν, καὶ περί πάντων ποιεῖσθαι λόγον καὶ πολυπραγμονεῖν ἐν τοῖς καταμέρος, τῷ τῆς ἱστορίας ἀρχηγέτῃ καθήκει·
\vs{31}Τὸ δὲ σύντομον τῆς λέξεως μεταδιώκειν, καὶ τὸ ἐξεργαστικὸν τῆς πραγματείας παραιτεῖσθαι, τῷ τὴν μετάφρασιν ποιουμένῳ συγχωρητέον.
\vs{32}Ἐντεῦθεν οὖν ἀρξώμεθα τῆς διηγήσεως, τοῖς προειρημένοις τοσοῦτον ἐπιζεύξαντες· εὔηθες γὰρ τὸ μὲν πρὸ τῆς ἱστορίας πλεονάζειν, τὴν δὲ ἱστορίαν ἐπιτεμεῖν.

\ch{3}
Τῆς ἁγίας τοίνυν πόλεως κατοικουμένης μετὰ πάσης εἰρήνης, καὶ τῶν νόμων ἕτι κάλλιστα συντηρουμένων διὰ τὴν Ὀνίου τοῦ ἀρχιερέως εὐσέβειάν τε καὶ μισοπονηρίαν,
\vs{2}συνέβαινε καὶ αὐτοὺς τοὺς βασιλεῖς τιμᾷν τὸν τόπον, καὶ τὸ ἱερὸν ἀποστολαῖς ταῖς κρατίσταις δοξάζειν,
\vs{3}ὥστε καὶ Σέλευκον τὸν τῆς Ἀσίας βασιλέα χορηγεῖν ἐκ τῶν ἰδίων προσόδων πάντα τὰ πρὸς τὰς λειτουργίας τῶν θυσιῶν ἐπιβάλλοντα δαπανήματα.

\vs{4}Σίμων δέ τις ἐκ τῆς Βενιαμὶν φυλῆς προστάτης τοῦ ἱεροῦ καθεσταμένος, διηνέχθη τῷ ἀρχιερεῖ περὶ τῆς κατὰ τὴν πόλιν παρανομίας·
\vs{5}καὶ νικῆσαι τὸν Ὀνίαν μὴ δυνάμενος, ἦλθε πρὸς Ἀπολλώνιον Θρασαίου, τὸν κατʼ ἐκεῖνον τὸν καιρὸν κοιλῆς Συρίας καὶ Φοινίκης στρατηγόν.
\vs{6}Καὶ προσήγγειλε περὶ τοῦ χρημάτων ἀμυθήτων γέμειν τὸ ἐν Ἱεροσολύμοις γαζοφυλάκιον, ὥστε τὸ πλῆθος τῶν διαφόρων ἐναρίθμητον εἶναι, καὶ μὴ προσήκειν αὐτὰ πρὸς τὸν τῶν θυσιῶν λόγον, εἶναι δὲ δυνατὸν ὑπὸ τὴν τοῦ βασιλέως ἐξουσίαν πεσεῖν ἅπαντα ταῦτα.

\vs{7}Συμμίξας δὲ ὁ Ἀπολλώνιος τῷ βασιλεῖ, περὶ τῶν μηνυθέντων αὐτῷ χρημάτων ἐνεφάνισεν· ὁ δὲ προχειρισάμενος Ἡλιόδωρον τὸν ἐπὶ τῶν πραγμάτων, ἀπέστειλε δοὺς ἐντολὰς, τὴν τῶν προειρημένων χρημάτων ἐκκομιδὴν ποιήσασθαι.
\vs{8}Εὐθέως δὲ ὁ Ἡλιόδωρος ἐποιεῖτο τὴν παρείαν, τῇ μὲν ἐμφάσει ὡς τὰς κατὰ κοίλην Συρίαν καὶ Φοινίκην πόλεις ἐφοδεύσων, τῷ πράγματι δὲ τὴν τοῦ βασιλέως πρόθεσιν ἐπιτελέσων.

\vs{9}Παραγενηθεὶς δὲ εἰς Ἱεροσόλυμα, καὶ φιλοφρόνως ὑπὸ τοῦ ἀρχιερέως τῆς πόλεως ἀποδεχθείς, ἀνέθετο περὶ τοῦ γεγονότος ἐμφανισμοῦ, καὶ τίνος ἕνεκεν πάρεστι διεσάφήσεν· ἐπυνθάνετο δὲ εἰ ταῖς ἀληθείαις ταῦτα οὕτως ἔχοντα τυγχάνει.

\vs{10}Τοῦ δὲ ἀρχιερέως ὑποδείξαντος παραθήκας εἶναι χηρῶν τε καὶ ὀρφανῶν,
\vs{11}τινὰ δὲ καὶ Ὑρκανοῦ τοῦ Τωβίου σφόδρα ἀνδρὸς ἐν ὑπεροχῇ κειμένου, οὐχ ὥσπερ ἦν διαβάλλων ὁ δυσσεβὴς Σίμων, τὰ δὲ πάντα ἀργυρίου τετρακόσια τάλαντα, χρυσίου δὲ διακόσια·
\vs{12}ἀδικηθῆναι δὲ τοὺς πεπιστευκότας τῇ τοῦ τόπου ἁγιωσύνῃ, καὶ τῇ τοῦ τετιμημένου κατὰ τὸν σύμπαντα κόσμον ἱεροῦ σεμνότητι καὶ ἀσυλίᾳ, παντελῶς ἀμήχανον εἶναι.

\vs{13}Ὁ δὲ Ἡλιόδωρος διʼ ἃς εἶχε βασιλικὰς ἐντολὰς, πάντως ἔλεγεν εἰς τὸ βασιλικὸν ἀναληπτέα ταῦτα εἶναι.
\vs{14}Ταξάμενος δὲ ἡμέραν εἰσῄει τὴν περὶ τούτων ἐπίσκεψιν οἰκονομήσων· ἦν δὲ οὐ μικρὰ καθʼ ὅλην τὴν πόλιν ἀγωνία.
\vs{15}Οἱ δὲ ἱερεῖς πρὸ τοῦ θυσιαστηρίου ἐν ταῖς ἱερατικαῖς στολαῖς ῥίψαντες ἑαυτοὺς, ἐπεκαλοῦντο εἰς οὐρανὸν τὸν περὶ παραθήκης νομοθετήσαντα τοῖς παρακαταθεμένοις ταῦτα σῶα διαφυλάξαι.

\vs{16}Ἦν δὲ ὁρῶντα τὴν τοῦ ἀρχιερέως ἰδέαν, τιτρώσκεσθαι τὴν διάνοιαν· ἡ γὰρ ὄψις καὶ τὸ τῆς χρόας παρηλλαγμένον ἐνέφαινε τὴν κατὰ ψυχὴν ἀγωνίαν.
\vs{17}Περιεκέχυτο γὰρ περὶ τὸν ἄνδρα δέος τι καὶ φρικασμὸς σώματος, διʼ ὧν πρόδηλον ἐγένετο τοῖς θεωροῦσι τὸ κατὰ καρδίαν ἐνεστὸς ἄλγος.

\vs{18}Οἱ δὲ ἐκ τῶν οἰκιῶν ἀγεληδὸν ἐξεπήδων ἐπὶ πάνδημον ἱκετείαν, διὰ τὸ μέλλειν εἰς καταφρόνησιν ἔρχεσθαι τὸν τόπον.
\vs{19}Ὑπεζωσμέναι δὲ ὑπὸ τοὺς μαστοὺς αἱ γυναῖκες σάκκους κατὰ τὰς ὁδοὺς ἐπλήθυον· αἱ δὲ κατάκλειστοι τῶν παρθένων, αἱ μὲν συνέτρεχον ἐπὶ τοὺς πυλῶνας, αἱ δὲ ἐπὶ τὰ τείχη, τινὲς δὲ διὰ τῶν θυρίδων διεξέκυπτον.
\vs{20}Πᾶσαι δὲ προτείνουσαι τὰς χεῖρας εἰς τὸν οὐρανον, ἐποιοῦντο τὴν λιτανείαν·

\vs{21}Ἐλεεῖν δʼ ἦν τὴν τοῦ πλήθους παμμιγῆ πρόπτωσιν, τήν τε τοῦ μεγάλως διαγωνιῶντος ἀρχιερέως προσδοκίαν.
\vs{22}Οἱ μὲν οὖν ἐπεκαλοῦντο τὸν παντοκράτορα Θεὸν τὰ πεπιστευμένα τοῖς πεπιστευκόσι σῶα διαφυλάγγειν μετὰ πάσης ἀσφαλείας.

\vs{23}Ὁ δὲ Ἡλιόδωρος τὸ διεγνωσμένον ἐπετέλει.

\vs{24}Αὐτόθι δὲ αὐτοῦ σὺν τοῖς δορυφόροις κατὰ τὸ γαζοφυλάκιον ἤδη παρόντος, ὁ τῶν πατέρων Κύριος καὶ πάσης ἐξουσίας δυνάστης ἐπιφάνειαν μεγάλην ἐποίησεν, ὥστε πάντας τοὺς κατατολμήσαντας συνελθεῖν, καταπλαγέντας τὴν τοῦ Θεοῦ δύναμιν, εἰς ἔκλυσιν καὶ δειλίαν τραπῆναι.
\vs{25}Ὤφθη γάρ τις ἵππος αὐτοῖς φοβερὸν ἔχων τὸν ἐπιβάτην, καὶ καλλίστῃ σαγῇ διακεκοσμημένος, φερόμενος δὲ ῥύδην ἐνέσεισε τῷ Ἡλιοδώρῳ τὰς ἐμπροσθίους ὁπλάς· ὁ δὲ ἐπικαθήμενος ἐφαίνετο χρυσῆν πανοπλίαν ἔχων.

\vs{26}Ἕτεροι δὲ δύο προεφάνησαν αὐτῷ νεανίαι, τῇ ῥώμῃ μὲν ἐκπρεπεῖς, κάλλιστοι δὲ τῇ δόξῃ, διαπρεπεῖς δὲ τὴν περιβολήν· οἳ καὶ παρασταντες ἐξ ἑκατέρου μέρους, ἐμαστίγουν αὐτὸν ἀδιαλείπτως, πολλὰς ἐπιῤῥιπτοῦντες αὐτῷ πληγάς.

\vs{27}Ἄφνω δὲ πεσόντα πρὸς τὴν γῆν, καὶ πολλῷ σκότει περιχυθέντα, συναρπάσαντες, καὶ εἰς φορεῖον ἐνθέντες,
\vs{28}τὸν ἄρτι μετὰ πολλῆς παραδρομῆς καὶ πάσης δορυφορίας εἰς τὸ προειρημένον εἰσελθόντα γαζοφυλάκιον, ἔφερον ἀβοήθητον ἑαυτῷ καθεστῶτα, φανερῶς τὴν τοῦ Θεοῦ δυναστείαν ἐπεγνωκότες.
\vs{29}Καὶ ὁ μὲν διὰ τὴν θείαν ἐνέργειαν ἄφωνος καὶ πάσης ἐστερημένος ἐλπίδος καὶ σωτηρίας ἔῤῥιπτο.
\vs{30}Οἱ δὲ τὸν κύριον εὐλόγουν τὸν παραδοξάζοντα τὸν ἑαυτοῦ τόπον· καὶ τὸ μικρῷ πρότερον δέους καὶ ταραχῆς γέμον ἱερὸν, τοῦ παντοκράτορος ἐπιφανέντος Κυρίου, χαρᾶς καὶ εὐφροσύνης ἐπεπλήρωτο.

\vs{31}Ταχὺ δέ τινες τῶν τοῦ Ἡλιοδώρου συνήθων ἠξίουν τὸν Ὀνίαν ἐπικαλέσασθαι τὸν ὕψιστον, καὶ τὸ ζῇν χαρίσασθι τῷ παντελῶς ἐν ἐσχάτῃ πνοῇ κειμένῳ.
\vs{32}Ὕποπτος δὲ γενόμενος ὁ ἀρχιερεῦς, μήποτε διάληψιν ὁ βασιλεὺς σχῇ, κακουργίαν τινὰ περὶ τὸν Ἡλιόδωρον ὑπὸ τῶν Ἰουδαίων συντετελέσθαι, προσήγαγε θυσίαν ὑπὲρ τῆς τοῦ ἀνδρὸς σωτηρίας.

\vs{33}Ποιουμένου δὲ τοῦ ἀρχιερέως τὸν ἱλασμὸν, οἱ αὐτοὶ νεανίαι πάλιν ἐφάνησαν τῷ Ἡλιοδώρῳ ἐν ταῖς αὐταῖς ἐσθήσεσιν ἐστολισμένοι, καὶ στάντες εἶπον, πολλὰς τῷ Ὀνία τῷ ἀρχιερεῖ χάριτας ἔχε, διὰ γὰρ αὐτὸν σοι κεχάρισται τὸ ζῇν ὁ Κύριος.
\vs{34}Σὺ δὲ ὑπʼ αὐτοῦ μεμαστιγωμένος διάγγελε πᾶσι τὸ μεγαλεῖον τοῦ Θεοῦ κράτος· ταῦτα δὲ εἰπόντες ἀφανεῖς ἐγένοντο.

\vs{35}Ὁ δὲ Ἡλιόδωρος θυσίαν ἀνενέγκας τῷ Κυρίῳ, καὶ εὐχὰς μεγίστας εὐξάμενος τῷ τὸ ζῇν περιποιήσαντι, καὶ τὸν Ὀνίαν ἀποδεξάμενος, ἀνεστρατοπέδευσε πρὸς τὸν βασιλέα.
\vs{36}Ἐξεμαρτύρει δὲ πᾶσιν ἅπερ ἦν ὑπʼ ὄψιν τεθεαμένος ἔργα τοῦ μεγίστου Θεοῦ.

\vs{37}Τοῦ δὲ βασιλέως ἐπερωτήσαντος τὸν Ἡλιόδωρον, ποῖός τις εἴη ἐπιτήδειος ἔτι ἅπαξ διαπεμφθῆναι εἰς Ἱεροσόλυμα, ἔφησεν,
\vs{38}εἴ τινα ἔχεις πολέμιον ἢ πραγμάτων ἐπίβουλον, πέμψον αὐτὸν ἐκεῖ, καὶ μεμαστιγωμένον αὐτὸν προσδέξῃ, ἐάνπερ καὶ διασωθείη, διὰ τὸ περὶ τὸν τόπον ἀληθῶς εἶναί τινα Θεοῦ δύναμιν.
\vs{39}Αὐτὸς γὰρ ὁ τὴν κατοικίαν ἐπουράνιον ἔχων, ἐπόπτης ἐστὶ καὶ βοηθὸς ἐκείνου τοῦ τόπου, καὶ τοὺς παραγινομένους ἐπὶ κακώσει, τύπτων ἀπόλλυσι.

\vs{40}Καὶ τὰ μὲν κατὰ Ἡλιόδωρον, καὶ τὴν τοῦ γαζοφυλακίου τήρησιν οὕτως ἐχώρησεν.

\ch{4}
Ὁ δὲ προειρημένος Σίμων ὁ τῶν χρημάτων καὶ τῆς πατρίδος ἐνδείκτης γεγονὼς, ἐκακολόγει τὸν Ὀνίαν, ὡς αὐτός τε εἴη τὸν Ἡλιόδωρον ἐπισεσεικὼς, καὶ τῶν κακῶν δημιουργὸς καθεστηκώς.
\vs{2}Καὶ τὸν εὐεργέτην τῆς πόλεως, καὶ τὸν κηδεμόνα τῶν ὁμοεθνῶν, καὶ ζηλωτὴν τῶν νόμων, ἐπίβουλον τῶν πραγμάτων ἐτόλμα λέγειν.

\vs{3}Τῆς δὲ ἔχθρας ἐπὶ τοσοῦτον προβαινούσης, ὥστε καὶ διά τινος τῶν ὑπὸ τοῦ Σίμωνος δεδοκιμασμένων φόνους συντελεῖσθαι,
\vs{4}συνορῶν ὁ Ὀνίας τὸ χαλεπὸν τῆς φιλονεικίας, καὶ Ἀπολλώνιον μαίνεσθαι, ὡς τὸν κοίλης Συρίας καὶ Φοινίκης στρατηγὸν, συναύξοντα τὴν κακίαν τοῦ Σίμωνος,
\vs{5}ὡς τὸν βασιλέα διεκομίσθη, οὐ γινόμενος τῶν πολιτῶν κατήγορος, τὸ δὲ συμφέρον κοινῇ κατʼ ἰδίαν παντὶ τῷ πλήθει σκοπῶν.
\vs{6}Ἑώρα γὰρ ἄνευ βασιλικῆς προνοίας ἀδύνατον εἶναι τυχεῖν εἰρήνης ἔτι τὰ πράγματα, καὶ τὸν Σίμωνα παῦλαν οὐ ληψόμενον τῆς ἀνοίας.

\vs{7}Μεταλλάξαντος δὲ τὸν βίον Σελεύκου, καὶ παραλαβόντος τὴν βασιλείαν Ἀντιόχου τοῦ προσαγορευθέντος Ἐπιφανοῦς, ὑπενόθευσεν Ἰάσων ὁ ἀδελφὸς Ὀνίου τὴν ἀρχιερωσύνην,
\vs{8}ἐπαγγειλάμενος τῷ βασιλεῖ διʼ ἐντεύξεως ἀργυοίου τάλαντα ἑξήκοντα πρὸς τοῖς τριακοσίοις, καὶ προσόδου τινὸς ἄλλης τάλαντα ὀγδοήκοντα,
\vs{9}πρὸς δὲ τούτοις ὑπισχνεῖτο καὶ ἕτερα διαγράψαι πεντήκοντα πρὸς τοῖς ἑκατόν, ἐὰν συγχωρηθῇ διὰ τῆς ἐξουσίας αὐτοῦ, γυμνάσιον καὶ ἐφηβίαν αὐτῷ συστήσασθαι, καὶ τοὺς ἐν Ἰεροσολύμοις Ἀντιοχεῖς ἀναγράψαι.
\vs{10}Ἐπινεύσαντος δὲ τοῦ βασιλέως, καὶ τῆς ἀρχῆς κρατήσας, εὐθέως ἐπὶ τὸν Ἑλληνικὸν χαρακτῆρα τοὺς ὁμοφύλους μετῆγε.

\vs{11}Καὶ τὰ κείμενα τοῖς Ἰουδαίοις φιλάνθρωπα βασιλικὰ διὰ Ἰωάννου τοῦ πατρὸς Εὐπολέμου, τοῦ ποιησαμένου τὴν πρεσβείαν ὑπὲρ φιλίας καὶ συμμαχίας πρὸς τοὺς Ῥωμαίους, παρώσατο· καὶ τὰς μὲν νομίμους καταλύων πολιτείας, παρανόμους ἐθισμοὺς ἐκαίνιζεν.
\vs{12}Ἀσμένως γὰρ ὑπʼ αὐτὴν τὴν ἀκρόπολιν γυμνάσιον καθίδρυσε, καὶ τοὺς κρατίστους τῶν ἐφήβων ὑποτάσσων, ὑπὸ πέτασον ἦγεν.

\vs{13}Ἦν δʼ οὕτως ἀκμή τις Ἑλληνισμοῦ, καὶ πρόσβασις ἀλλοφυλισμοῦ διὰ τὴν τοῦ ἀσεβοῦς καὶ οὐκ ἀρχιερέως Ἰάσωνος ὑπερβάλλουσαν ἀναγνείαν,
\vs{14}ὥστε μηκέτι περὶ τὰς τοῦ θυσιαστηρίου λειτουργίας προθύμους εἶναι τοὺς ἱερεῖς, ἀλλὰ τοῦ μὲν ναοῦ καταφρονοῦντες, καὶ τῶν θυσιῶν ἀμελοῦντες ἔσπευδον μετέχειν τῆς ἐν παλαίστρᾳ παρανόμου χορηγίας, μετὰ τὴν τοῦ δίσκου πρόκλησιν.
\vs{15}Καὶ τὰς μὲν πατρῴους τιμὰς ἐν οὐδενὶ τιθέμενοι, τὰς δὲ Ἑλληνικὰς δόξας καλλίστας ἡγούμενοι.

\vs{16}Ὧν χάριν περιέσχεν αὐτοὺς χαλεπὴ περίστασις, καὶ ὧν ἐζήλουν τὰς ἀγωγὰς, καὶ καθάπαν ἤθελον ἐξομοιοῦσθαι, τούτους πολεμίους καὶ τιμωρητὰς ἔσχον.
\vs{17}Ἀσεβεῖν γὰρ εἰς τοὺς θείους νόμους οὐ ῥᾴδιον, ἀλλὰ ταῦτα ὁ ἀκόλουθος καιρὸς δηλώσει.

\vs{18}Ἀγομένου δὲ πενταετηρικοῦ ἀγῶνος ἐν Τύρῳ, καὶ τοῦ βασιλέως παρόντος,
\vs{19}ἀπέστειλεν Ἰάσων ὁ μιαρὸς θεωροὺς ἀπὸ Ἱεροσολύμων Ἀντιοχεῖς ὄντας, παρακομίζοντας ἀργυρίου δραχμὰς τριακοσίας εἰς τὴν τοῦ Ἡρακλέους θυσίαν· ἃς καὶ ἠξίωσαν οἱ παρακομίσαντες μὴ χρῆσθαι πρὸς θυσίαν διὰ τὸ μὴ καθήκειν, εἰς ἑτέραν δὲ καταθέσθαι δαπάνην.
\vs{20}Ἔπεμψεν οὖν ταῦτα, διὰ μὲν τὸν ἀποστείλαντα εἰς τὴν τοῦ Ἡρακλέους θυσίαν, ἕνεκεν δὲ τῶν παρακομιζόντων, εἰς τὰς τῶν τριήρων κατασκευάς.

\vs{21}Ἀποσταλέντος δὲ εἰς Αἴγυπτον Ἀπολλωνίου τοῦ Μενεσθέως διὰ τὰ πρωτοκλίσια Πτολεμαίου τοῦ Φιλομήτορος βασιλέως, μεταλαβὼν Ἀντίοχος ἀλλότριον αὐτὸν τῶυ αὐτῶν γεγονέναι πραγμάτων, τῆς κατʼ αὑτὸν ἀσφαλείας ἐφρόντιζεν· ὅθεν εἰς Ἰόππην παραγενόμενος, κατήντησεν εἰς Ἰεροσόλυμα.
\vs{22}Μεγαλοπρεπῶς δὲ ὑπὸ τοῦ Ἰάσωνος καὶ τῆς πόλεως παραδεχθεὶς, μετὰ δᾳδουχίας καὶ βοῶν εἰσπεπόρευται εἶθʼ οὕτως εἰς τὴν Φοινίκην κατεστρατοπέδευσε.

\vs{23}Μετὰ δὲ τριετῆ χρόνον ἀπέστειλεν Ἰάσων Μενέλαον τὸν τοῦ προσημαινομένου Σίμωνος ἀδελφόν, παρακομίζοντα τὰ χρήματα τῷ βασιλεῖ, καὶ περὶ πραγμάτων ἀναγκαίων ὑπομνηματισμοὺς τελέσοντα.
\vs{24}Ὁ δὲ συσταθεὶς τῷ βασιλεῖ, καὶ δοξάσας αὐτὸν τῷ προσώπῳ τῆς ἐξουσίας, εἰς ἑαυτὸν κατήντησε τὴν ἀρχιερωσύνην, ὑπερβαλὼν τὸν Ἰάσωνα τάλαντα ἀργυρίου τριακόσια.
\vs{25}Λαβὼν δὲ τὰς βασιλικὰς ἐντολὰς παρεγένετο, τῆς μὲν ἀρχιερωσύνης οὐδὲν ἄξιον φέρων, θυμοὺς δὲ ὠμοῦ τυράννου, καὶ θηρὸς βαρβάρου ὀργὰς ἔχων.

\vs{26}Καὶ ὁ μὲν Ἰάσων ὁ τὸν ἴδιον ἀδελφὸν ὑπονοθεύσας, ὑπονοθευθεὶς ὑφʼ ἑτέρου φυγὰς εἰς τὴν Ἀμμανῖτιν χώραν συνήλαστο.
\vs{27}Ὁ δὲ Μενέλαος τῆς μὲν ἀρχῆς ἐκράτει, τῶν δὲ ἐπηγγελμένων τῷ βασιλεῖ χρημάτων οὐδὲν εὐτάκτει,
\vs{28}ποιουμένου δὲ τὴν ἀπαίτησιν Σωστράτου τοῦ τῆς ἀκροπόλεως ἐπάρχου· πρὸς τοῦτον γὰρ ἦν ἡ τῶν φόρων πρᾶξις· διʼ ἣν αἰτίαν οἱ δύο ὑπὸ τοῦ βασιλέως προσεκλήθησαν.

\vs{29}Καὶ ὁ μὲν Μενέλαος ἀπέλιπε τῆς ἀρχιερωσύνης διάδοχον Λυσίμαχον τὸν ἑαυτοῦ ἀδελφόν, Σώστρατος δὲ, Κράτητα τὸν ἐπὶ τῶν Κυπρίων.

\vs{30}Τοιούτων δὲ συνεστηκότων, συνέβη Ταρσεῖς, καὶ Μαλλώτας στασιάζειν, διὰ τὸ Ἀντιοχίδι τῇ παλλακῇ τοῦ βασιλέως ἐν δωρεᾷ δεδόσθαι.
\vs{31}Θᾶττον οὖν ὁ βασιλεὺς ἧκε καταστεῖλαι τὰ πράγματα, καταλιπὼν τὸν διαδεχόμενον Ἀνδρόνικον, τῶν ἐν ἀξιώματι κειμένων.

\vs{32}Νομίσας δὲ ὁ Μενέλαος εἰληφέναι καιρὸν εὐφυῆ, χρυσώματά τινα τῶν τοῦ ἱεροῦ νοσφισάμενος ἐχαρίσατο τῷ Ἀνδρονίκῳ, καὶ ἕτερα ἐτύγχανε πεπρακὼς εἴς τε Τύρον καὶ τὰς κύκλῳ πόλεις.
\vs{33}Ἃ καὶ σαφῶς ἐπεγνωκὼς ὁ Ὀνίας, παρήλεγχεν ἀποκεχωρηκὼς εἰς ἄσυλον τόπον, ἐπὶ Δάφνης τῆς πρὸς Ἀντιόχειαν κειμένης.

\vs{34}Ὅθεν ὁ Μενέλαος λαβὼν ἰδίᾳ τὸν Ἀνδρόνικον, παρεκάλει χειρώσασθαι τὸν Ὀνίαν· ὁ δὲ παραγενόμενος ἐπὶ τὸν Ὀνίαν, καὶ πεισθεὶς ἐπὶ δόλῳ, καὶ δεξιὰς μεθʼ ὅρκων δοὺς, καίπερ ἐν ὑποψίᾳ κείμενος ἔπεισεν ἐκ τοῦ ἀσύλου προελθεῖν, ὃν καὶ παραχρῆμα παρέκλεισεν, οὐκ αἰδεσθεὶς τὸ δίκαιον.
\vs{35}Διʼ ἣν αἰτίαν οὐ μόνον Ἰουδαῖοι, πολλοὶ δὲ καὶ τῶν ἄλλων ἐθνῶν ἐδείναζον, καὶ ἐδυσφόρουν ἐπὶ τῷ τοῦ ἀνδρὸς ἀδίκῳ φόνῳ.

\vs{36}Τοῦ δὲ βασιλέως ἐπανελθόντος ἀπὸ τῶν κατὰ Κιλικίαν τόπων, ἐνετύγχανον οἱ κατὰ πόλιν Ἰουδαῖοι συμισοπονηρούντων καὶ τῶν Ἑλλήνων, ὑπὲρ τοῦ παρὰ λόγον τὸν Ὀνίαν ἀπεκτάνθαι.
\vs{37}Ψυχικῶς οὖν ὁ Ἀντίοχος ἐπιλυπηθεὶς, καὶ τραπεὶς εἰς ἔλεον, καὶ δακρύσας διὰ τὴν τοῦ μετηλλαχότος σωφροσύνην, καὶ πολλὴν εὐταξίαν,
\vs{38}καὶ πυρωθεὶς τοῖς θυμοῖς, παραχρῆμα τὴν τοῦ Ἀνδρονίκου πορφύραν περιελόμενος, καὶ τοὺς χιτῶνας περιῤῥήξας, περιαγαγὼν καθʼ ὅλην τὴν πόλιν, ἐπʼ αὐτὸν τὸν τόπον οὗπερ εἰς τὸν Ὀνίαν ἠσέβησεν, ἐκεῖ τὸν μιαιφόνον ἀπεκόσμησε, τοῦ Κυρίου τὴν ἀξίαν αὐτῷ κόλασιν ἀποδόντος.

\vs{39}Γενομένων δὲ πολλῶν ἱεροσυλημάτων κατὰ τὴν πόλιν ὑπὸ τοῦ Λυσιμάχου μετὰ τῆς Μενελάου γνώμης, καὶ διαδοθείσης ἔξω τῆς φήμης, ἐπισυνήχθη τὸ πλῆθος ἐπὶ τὸν Λυσίμαχον, χρυσωμάτων ἤδη πολλῶν διενηνεγμένων.
\vs{40}Ἐπεγειρομένων δὲ τῶν ὄχλων, καὶ ταῖς ὀργαῖς διεμπιπλαμένων, καθοπλίσας ὁ Λυσίμαχος πρὸς τρισχιλίους, κατήρξατο χειρῶν ἀδίκων, προηγησαμένου τινὸς Τυράννου προβεβηκότος τὴν ἡλικίαν, οὐδὲν δὲ ἧττον καὶ τὴν ἄνοιαν.

\vs{41}Συνιδότες δὲ καὶ τὴν ἐπίθεσιν τοῦ Λυσιμάχου, συναρπάσαντες οἱ μὲν πέτρους, οἱ δὲ ξύλων πάχη, τινὲς δὲ ἐκ τῆς παρακειμένης σποδοῦ δρασσόμενοι, φύρδην ἐνετίνασσον εἰς τοὺς περὶ τὸν Λυσίμαχον.
\vs{42}Διʼ ἣν αἰτίαν πολλοὺς μὲν αὐτῶν τραυματίας ἐποίησαν, τινὰς δὲ καὶ κατέβαλον, πάντας δὲ εἰς φυγὴν συνήλασαν, αὐτὸν δὲ τὸν ἱερόσυλον παρὰ τὸ γαζοφυλάκιον ἐχειρώσαντο.

\vs{43}Περὶ δὲ τούτων ἐνέστη κρίσις πρὸς τὸν Μενέλαον.
\vs{44}Καταντήσαντος δὲ τοῦ βασιλέως εἰς Τύρον, ἐπʼ αὐτοῦ τὴν δικαιολογίαν ἐποιήσαντο οἱ πεμφθέντες ἄνδρες τρεῖς ὑπὸ τῆς γερουσίας.
\vs{45}Ἤδη δὲ λελειμμένος ὁ Μενέλαος ἐπηγγείλατο χρήματα ἱκανὰ τῷ Πτολεμαίῳ τῷ Δορυμένους πρὸς τὸ πεῖσαι τὸν βασιλέα.

\vs{46}Ὅθεν ἀπολαβὼν ὁ Πτολεμαῖος εἴς τι περίστυλον ὡς ἀναψύξοντα τὸν βασιλέα, μετέθηκε.
\vs{47}Καὶ τὸν μὲν τῆς ὅλης κακίας αἴτιον Μενέλαον ἀπέλυσε τῶν κατηγορημάτων, τοῖς δὲ ταλαιπώροις, οἵτινες εἰ καὶ ἐπὶ Σκυθῶν ἔλεγον, ἀπελύθησαν ἄν ἀκατάγνωστοι, τούτοις θάνατον ἐπέκρινε.

\vs{48}Ταχέως οὖν τὴν ἄδικον ζημίαν ὑπέσχον οἱ ὑπὲρ πόλεως καὶ δήμων καὶ τῶν ἱερῶν σκευῶν προαγορεύσαντες.
\vs{49}Διʼ ἣν αἰτίαν καὶ Τύριοι μισοπονηρήσαντες τὰ πρὸς τὴν κηδείαν αὐτῶν μεγαλοπρεπῶς ἐχορήγησαν.
\vs{50}Ὁ δὲ Μενέλαος διὰ τὰς τῶν κρατούντων πλεονεξίας, ἔμενεν ἐπὶ τῆς ἀρχῆς, ἐπιφυόμενος τῇ κακίᾳ, μέγας τῶν πολιτῶν ἐπίβουλος καθεστώς.

\ch{5}
Περὶ δὲ τὸν καιρὸν τοῦτον τὴν δευτέραν ἔφοδον ὁ Ἀντίοχος εἰς Αἴγυπτον ἐστείλατο.
\vs{2}Συνέβη δὲ καθʼ ὅλην τὴν πόλιν σχεδὸν ἐφʼ ἡμέρας τεσσαράκοντα φαίνεσθαι διὰ τοῦ ἀέρος τρέχοντας ἱππεῖς διαχρύσους στολὰς ἔχοντας, καὶ λόγχας σπειρηδὸν ἐξωπλισμένους,
\vs{3}καὶ ἴλας ἵππων διατεταγμένας, καὶ προσβολὰς γινομένας, καὶ καταδρομὰς ἑκατέρων, καὶ ἀσπίδων κινήσεις, καὶ καμάκων πλήθη, καὶ μαχαιρῶν σπασμούς, καὶ βελῶν βολὰς, καὶ χρυσῶν κόσμων ἐκλάμψεις, καὶ παντοίους θωρακισμούς.
\vs{4}Διὸ πάντες ἠξίουν ἐπʼ ἀγαθῷ τὴν ἐπιφάνειαν γενέσθαι.

\vs{5}Γενομένης δὲ λαλιᾶς ψευδοῦς, ὡς μετηλλαχότος τὸν βίον Ἀντιόχου, παραλαβὼν ὁ Ἰάσων οὐκ ἐλάττους τῶν χιλίων, αἰφνιδίως ἐπὶ τὴν πόλιν συνετελέσατο ἐπίθεσιν· τῶν δὲ ἐπὶ τῷ τείχει συνελασθέντων, καὶ τέλος ἤδη καταλαμβανομένης τῆς πόλεως, ὁ Μενέλαος εἰς τὴν ἀκρόπολιν ἐφυγάδευσεν.
\vs{6}Ὁ δὲ Ἰάσων ἐποιεῖτο σφαγὰς τῶν πολιτῶν τῶν ἰδίων ἀφειδῶς, οὐ συννοῶν τὴν εἰς τοὺς συγγενεῖς εὐημερίαν, δυσημερίαν εἶναι τὴν μεγίστην· δοκῶν δὲ πολεμίων καὶ οὐχ ὁμοεθνῶν τρόπαια καταβάλλεσθαι,
\vs{7}τῆς μὲν ἀρχῆς οὐκ ἐκράτησε, τὸ δὲ τέλος τῆς ἐπιβουλῆς αἰσχύνην λαβών, φυγὰς πάλιν εἰς τὴν Ἀμμανίτιν ἀπῆλθε.

\vs{8}Πέρας οὖν κακῆς ἀναστροφῆς ἔτυχεν ἐγκλεισθεὶς πρὸς Ἀρέταν τὸν τῶν Ἀράβων τύραννον, πόλιν ἐκ πόλεως φεύγων, διωκόμενος ὑπὸ πάντων, καὶ στυγούμενος ὡς τῶν νόμων ἀποστάτης, καὶ βδελυσσόμενος ὡς πατρίδος καὶ πολιτῶν δήμιος, εἰς Αἴγυπτον συνεβράσθη.
\vs{9}Καὶ ὁ συχνοὺς τῆς πατρίδος ἀποξενώσας, ἐπὶ ξένης ἀπώλετο πρὸς Λακεδαιμονίους ἀναχθείς, ὡς διὰ τὴν συγγένειαν τευξόμενος σκέπης.
\vs{10}Καὶ ὁ πλῆθος ἀτάφων ἐκρίψας ἀπένθητος ἐγενήθη, καὶ κηδείας οὐδʼ ἡστινοσοῦν οὔτε πατρῴου τάφου μετέσχε.

\vs{11}Προσπεσόντων δὲ τῷ βασιλεῖ περὶ τῶν γεγονότων, διέλαβεν ἀποστατεῖν τὴν Ἰουδαίαν· ὅθεν ἀναζεύξας ἐξ Αἰγύπτου τεθηριωμένος τῇ ψυχῇ, ἔλαβε τὴν μὲν πόλιν δορυάλωτον.
\vs{12}Καὶ ἐκέλευσε τοῖς στρατιώταις κόπτειν ἀφειδῶς τοὺς ἐμπίπτοντας, καὶ τοὺς εἰς τὰς οἰκίας ἀναβαίνοντας κατασφάζειν.
\vs{13}Ἐγίνοντο δὲ νέων καὶ πρεσβυτέρων ἀναιρέσιες, ἀνδρῶν τε καὶ γυναικῶν καὶ τέκνων ἀφανισμὸς, παρθένων τε καὶ νηπίων σφαγαί.
\vs{14}Ὀκτὼ δὲ μυριάδες ἐν ταῖς πάσαις ἡμέραις τρισὶ κατεφθάρησαν, τέσσαρες μὲν ἐν χειρῶν νομαῖς, οὐκ ἧττον δὲ τῶν ἐσφαγμένων ἐπράθησαν.

\vs{15}Καὶ οὐκ ἀρκεσθεὶς δὲ τούτοις, κατετόλμησεν εἰς τὸ πάσης τῆς γῆς ἁγιώτατον ἱερὸν εἰσελθεῖν, ὁδηγὸν ἔχων τὸν Μενέλαον, τὸν καὶ τῶν νόμων καὶ τῆς πατρίδος προδότην γεγονότα.
\vs{16}Καὶ ταῖς μιαραῖς χερσὶ τὰ ἱερὰ σκεύη λαμβάνων, καὶ τὰ ὑπʼ ἄλλων βασιλέων ἀνατεθέντα πρὸς αὔξησιν καὶ δόξαν τοῦ τόπου καὶ τιμήν, ταῖς βεβήλοις χερσὶ συσσύρων ἐπεδίδου.

\vs{17}Καὶ ἐμετεωρίζετο τὴν διάνοιαν ὁ Ἀντίοχος, οὐ συνορῶν ὅτι διὰ τὰς ἁμαρτίας τῶν τὴν πόλιν οἰκούντων ἀπώργισται βραχέως ὁ Δεσπότης, διὸ γέγονε περὶ τὸν τόπον παρόρασις.
\vs{18}Εἰ δὲ μὴ συνέβαινε προενέχεσθαι πολλοῖς ἁμαρτήμασι, καθάπερ ὁ Ἡλιόδωρος ὁ πεμφθεὶς ὑπὸ Σελεύκου τοῦ βασιλέως ἐπὶ τὴν ἐπίσκεψιν τοῦ γαζοφυλακίου, οὗτος προαχθεὶς παραχρῆμα μαστιγωθεὶς ἀνετράπη τοῦ θράσους.

\vs{19}Ἀλλʼ οὐ διὰ τὸν τόπον τὸ ἔθνος, ἀλλὰ διὰ τὸ ἔθνος τὸν τόπον ὁ Κύριος ἐξελέξατο.
\vs{20}Διόπερ καὶ αὐτὸς ὁ τόπος συμμετασχὼν τῶν τοῦ ἔθνους δυσπετημάτων γενομένων, ὕστερον εὐεργετημάτων ὑπὸ τοῦ Κυρίου ἐκοινώνησε· καὶ ὁ καταλειφθεὶς ἐν τῇ τοῦ παντοκράτορος ὀργῇ, πάλιν ἐν τῇ τοῦ μεγάλου Δεσπότου καταλλαγῇ μετὰ πάσης δόξης ἐπανωρθώθη.

\vs{21}Ὁ γοῦν Ἀντίοχος ὀκτακόσια πρὸς τοῖς χιλίοις ἀπενεγκάμενος ἐκ τοῦ ἱεροῦ τάλαντα, θᾶττον εἰς Ἀντιόχειαν ἐχωρίσθη, οἰόμενος ἀπὸ τῆς ὑπερηφανίας τὴν μὲν γῆν πλωτὴν, καὶ τὸ πέλαγος πορευτὸν θέσθαι διὰ τὸν μετεωρισμὸν τῆς καρδίας.

\vs{22}Κατέλιπε δὲ καὶ ἐπιστάτας τοῦ κακοῦν τὸ γένος, ἐν μὲν Ἱεροσολύμοις Φίλιππον, τὸ μὲν γένος Φρύγα, τὸν δὲ τρόπον βαρβαρώτερον ἔχοντα τοῦ καταστήσαντος·
\vs{23}ἐν δὲ Γαριζὶν Ἀνδρόνικον, πρὸς δὲ τούτοις Μενέλαον, ὃς χείριστα τῶν ἄλλων ὑπερῄρετο τοῖς πολίταις, ἀπεχθῆ δὲ πρὸς τοὺς πολίτας Ἰουδαίους ἔχων διάθεσιν.

\vs{24}Ἔπεμψε δὲ τὸν μυσάρχην Ἀπολλώνιον μετὰ στρατεύματος δισμυρίων πρὸς τοῖς δισχιλίοις, προστάξας τοὺς ἐν ἡλικίᾳ πάντας κατασφάξαι, τὰς δὲ γυναῖκας καὶ νεωτέρους πωλεῖν.
\vs{25}Οὗτος δὲ παραγενόμενος εἰς Ἱεροσόλυμα, καὶ τὸν εἰρηνικὸν ὑποκριθείς, ἐπέσχεν ἕως τῆς ἁγίας ἡμέρας τοῦ σαββάτου· καὶ λαβὼν ἀργοῦντας τοὺς Ἰουδαίους, τοῖς ὑφʼ ἑαυτὸν ἐξοπλησίαν παρήγγειλε.
\vs{26}Καὶ τοὺς ἐξελθόντας πάντας ἐπὶ τὴν θεωρίαν συνεξεκέντησε, καὶ εἰς τὴν πόλιν σὺν τοῖς ὅπλοις εἰσδραμὼν ἱκανὰ κατέστρωσε πλήθη.

\vs{27}Ἰούδας δὲ ὁ Μακκαβαῖος δέκατός που γενηθεὶς, καὶ ἀναχωρήσας εἰς τὴν ἔρημον, θηρίων τρόπον ἐν τοῖς ὄρεσι διέζη σὺν τοῖς μετʼ αὐτοῦ, καὶ τὴν χορτώδη τροφὴν σιτούμενοι διατέλουν, πρὸς τὸ μὴ μετασχεῖν τοῦ μολυσμοῦ

\ch{6}
Μετʼ οὐ πολὺν δὲ χρόνον ἐξαπέστειλεν ὁ βασιλεὺς γέροντα Ἀθηναῖον, ἀναγκάζειν τοὺς Ἰουδαίους μεταβαίνειν ἐκ τῶν πατρῴων νόμων, καὶ τοῖς τοῦ Θεοῦ νόμοις μὴ πολιτεύεσθαι,
\vs{2}μολῦναι δὲ καὶ τὸν ἐν Ἱεροσολύμοις νεὼν, καὶ προσονομάσαι Διὸς Ὀλυμπίου, καὶ τὸν ἐν Γαριζὶν, καθὼς ἐτύγχανον οἱ τὸν τόπον οἰκοῦντες, Διὸς Ξενίου.

\vs{3}Χαλεπὴ δὲ καὶ τοῖς ὄχλοις ἦν καὶ δυσχερὴς ἡ ἐπίστασις τῆς κακίας.
\vs{4}Τὸ μὲν γὰρ ἱερὸν ἀσωτίας καὶ κώμων ἐπεπλήρωτο ὑπὸ τῶν ἐθνῶν ῥαθυμούντων μεθʼ ἑταιρῶν, καὶ ἐν τοῖς ἱεροῖς περιβόλοις γυναιξὶ πλησιαζόντων, ἔτι δὲ τὰ μὴ καθήκοντα ἔνδον φερόντων.
\vs{5}Τὸ δὲ θυσιαστήριον τοῖς ἀποδιεσταλμένοις ἀπὸ τῶν νόμων ἀθεμίτοις ἐπεπλήρωτο.
\vs{6}Ἦν δʼ οὔτε σαββατίζειν, οὔτε πατρῴους ἑορτὰς διαφυλάττειν, οὔτε ἁπλῶς Ἰουδαῖον ὁμολογεῖν εἶναι.

\vs{7}Ἤγοντο δὲ μετὰ πικρᾶς ἀνάγκης εἰς τὴν κατὰ μῆνα τοῦ βασιλέως γενέθλιον ἡμέραν ἐπὶ σπλαγχνισμόν· γενομένης δὲ Διονυσίων ἑορτῆς, ἠναγκάζοντο οἱ Ἰουδαῖοι κισσοὺς ἔχοντες πομπεύειν τῷ Διονύσῳ.

\vs{8}Ψήφισμα δὲ ἐξέπεσεν εἰς τὰς ἀστυγείτονας πόλεις Ἑλληνίδας, Πτολεμαίων ὑποτιθεμένων τὴν αὐτὴν ἀγωγὴν κατὰ τῶν Ἰουδαίων, ἄγειν καὶ σπλαγχνίζειν·
\vs{9}τοὺς δὲ μὴ προαιρουμένους μεταβαίνειν ἐπὶ τὰ Ἑλληνικὰ, κατασφάζειν· παρῆν οὖν ὁρᾷν τὴν ἐνεστῶσαν ταλαιπωρίαν.

\vs{10}Δύο γὰρ γυναῖκες ἀνηνέγθησαν περιτετμηκυῖαι τὰ τέκνα αὐτῶν· τούτων δὲ ἐκ τῶν μαστῶν κρεμάσαντες τὰ βρέφη, καὶ δημοσίᾳ περιαγαγόντες αὐτὰς τὴν πόλιν, κατὰ τοῦ τείχους ἐκρήμνισαν.
\vs{11}Ἕτεροι δὲ πλησίον συνδραμόντες εἰς τὰ σπήλαια, λεληθότως ἄγειν τὴν ἑβδομάδα, μηνυθέντες τῷ Φιλίππῳ συνεφλογίσθησαν, διὰ τὸ εὐλαβῶς ἔχειν βοηθῆσαι ἑαυτοῖς κατὰ τὴν δόξαν τῆς σεμνοτάτης ἡμέρας.

\vs{12}Παρακαλῶ οὖν τοὺς ἐντυγχάνοντας τῇδε τῇ βίβλῳ, μὴ συστέλλεσθαι διὰ τὰς συμφορὰς, λογίζεσθαι δὲ τὰς τιμωρίας μὴ πρὸς ὄλεθρον, ἀλλὰ πρὸς παιδείαν τοῦ γένους ἡμῶν εὖναι.
\vs{13}Καὶ τὸ μὴ πολὺν χρόνον ἐᾶσθαι τοὺς δυσσεβοῦντας, ἀλλʼ εὐθέως περιπίπτειν ἐπιτιμίοις, μεγάλης εὐεργεσίας σημεῖόν ἐστιν.

\vs{14}Οὐ γὰρ, καθάπερ καὶ ἐπὶ τῶν ἄλλων ἐθνῶν ἀναμένει μακροθυμῶν ὁ Δεσπότης, μέχρι τοῦ καταντήσαντας αὐτοὺς πρὸς ἐκπλήρωσιν ἁμαρτιῶν, κολάσαι, οὕτω καὶ ἐφʼ ἡμῶν ἔκρινεν
\vs{15}εἶναι, ἵνα μὴ πρὸς τέλος ἀφικομένων ἡμῶν τῶν ἁμαρτιῶν, ὕστερον ἡμᾶς ἐκδικᾷ.
\vs{16}Διόπερ οὐδέ ποτε μὲν τὸν ἔλεον αὐτοῦ ἀφʼ ἡμῶν ἀφίστησι· παιδεύων δὲ μετὰ συμφορᾶς, οὐκ ἐγκαταλείπει τὸν ἑαυτοῦ λαόν.
\vs{17}Πλὴν ἕως ὑπομνήσεως ταῦθʼ ἡμῖν εἰρήσθω· διʼ ὀλίγων δʼ ἐλευστέον ἐπὶ τὴν διήγησιν.

\vs{18}Ἐλεάζαρός τις τῶν πρωτευόντων γραμματέων, ἀνὴρ ἤδη προβεβηκὼς τὴν ἡλικίαν, καὶ τὴν πρόσοψιν τοῦ προσώπου κάλλιστος τυγχάνων, ἀναχανὼν ἠναγκάζετο φαγεῖν ὕειον κρέας.
\vs{19}Ὁ δὲ τὸν μετʼ εὐκλείας θάνατον μᾶλλον ἢ τὸν μετὰ μύσους βίον ἀναδεξάμενος, αὐθαιρέτως ἐπὶ τὸ τύμπανον προσῆγε·
\vs{20}προπτύσας δέ, καθʼ ὃν ἔδει τρόπον προσέρχεσθαι τοὺς ὑπομένοντας ἀμύνεσθαι, ὧν οὐ θέμις γεύσασθαι διὰ τὴν πρὸς τὸ ζῇν φιλοστοργίαν.

\vs{21}Οἱ δὲ πρὸς τῷ παρανόμῳ σπλαγχνισμῷ τεταγμένοι, διὰ τὴν ἐκ τῶν παλαιῶν χρόνων πρὸς τὸν ἄνδρα γνῶσιν, ἀπολαβόντες αὐτὸν κατιδίαν παρεκάλουν, ἐνέγκαντα κρέα οἷς καθῆκον αὐτῷ χρήσασθαι διʼ αὐτοῦ παρασκευασθέντα, ὑποκριθῆναι δὲ ὡς ἐσθίοντα τὰ ὑπὸ τοῦ βασιλέως προστεταγμένα τῶν ἀπὸ τῆς θυσίας κρεῶν,
\vs{22}ἵνα τοῦτο πράξας ἀπολυθῇ τοῦ θανάτου, καὶ διὰ τὴν ἀρχαίαν πρὸς αὐτοὺς φιλίαν τύχῃ φιλανθρωπίας.

\vs{23}Ὁ δὲ λογισμὸν ἀστεῖον ἀναλαβὼν καὶ ἄξιον τῆς ἡλικίας, καὶ τῆς τοῦ γήρως ὑπεροχῆς, καὶ τῆς ἐπικτήτου καὶ ἐπιφανοῦς πολιᾶς, καὶ τῆς ἐκ παιδὸς καλλίστης ἀνατροφῆς, μᾶλλον δὲ τῆς ἁγίας καὶ θεοκτίστου νομοθεσίας, ἀκολούθως ἀπεφῄνατο, ταχέως λέγων προπέμπειν εἰς τὸν ᾅδην.

\vs{24}Οὐ γὰρ τῆς ἡμετέρας ἡλικίας ἄξιόν ἐστιν ὑποκριθῆναι, ἵνα πολλοὶ τῶν νέων ὑπολαβόντες Ἐλεάζαρον τὸν ἐννενηκονταετῆ μεταβεβηκέναι εἰς ἀλλοφυλισμόν,
\vs{25}καὶ αὐτοὶ διὰ τὴν ἐμὴν ὑπόκρισιν, καὶ διὰ τὸ μικρὸν καὶ ἀκαριαῖον ζῇν πλανηθῶσι διʼ ἐμέ, καὶ μῦσος καὶ κηλίδα τοῦ γήρως κατακτήσομαι.
\vs{26}Εἰ γὰρ καὶ ἐπὶ τοῦ παρόντος ἐξελοῦμαι τὴν ἐξ ἀνθρώπων τιμωρίαν, ἀλλὰ τὰς τοῦ παντοκράτορος χεῖρας οὔτε ζῶν οὔτε ἀποθανὼν ἐκφεύξομαι.

\vs{27}Διόπερ ἀνδρείως μὲν νῦν διαλλάξας τὸν βίον, τοῦ μὲν γήρως ἄξιος φανήσομαι,
\vs{28}τοῖς δὲ νέοις ὑπόδειγμα γενναῖον καταλελοιπὼς, εἰς τὸ προθύμως καὶ γενναίως ὑπὲρ τῶν σεμνῶν καὶ ἁγίων νόμων ἀπευθανατίζειν· τοσαῦτα δὲ εἰπὼν. ἐπὶ τὸ τύμπανον εὐθέως ἦλθε.
\vs{29}Τῶν δὲ ἀγόντων τὴν μικρῷ πρότερον εὐμένειαν πρὸς αὐτὸν εἰς δυσμένειαν μεταβαλόντων διὰ τὸ τοὺς προειρημένους λόγους, ὡς αὐτοὶ διελάμβανον, ἀπόνοιαν εἶναι·

\vs{30}Μέλλων δὲ ταῖς πληγαῖς τελευτᾷν, ἀναστενάξας εἶπε, τῷ Κυρίῳ τῷ τὴν ἁγίαν γνῶσιν ἔχοντι φανερόν ἐστιν, ὅτι δυνάμενος ἀπολυθῆναι τοῦ θανάτου, σκληρὰς ὑποφέρω κατὰ τὸ σῶμα ἀλγηδόνας μαστιγούμενος, κατὰ ψυχὴν δὲ ἡδέως διὰ τὸν αὐτοῦ φόβον ταῦτα πάσχω.
\vs{31}Καὶ οὗτος οὖν τοῦτον τὸν τρόπον μετήλλαξεν, οὐ μόνον τοῖς νέοις, ἀλλὰ καὶ τοῖς πλείστοις τοῦ ἔθνους τὸν ἑαυτοῦ θάνατον ὑπόδειγμα γενναιότητος καὶ μνημόσυνον ἀρετῆς καταλιπών.

\ch{7}
Συνέβη δὲ καὶ ἑπτὰ ἀδελφοὺς μετὰ τῆς μητρὸν συλληφθέντας ἀναγκάζεσθαι ὑπὸ τοῦ βασιλέως ἀπὸ τῶν ἀθεμίτων ὑείων κρεῶν ἐφάπτεσθαι, μάστιξι καὶ νευραῖς αἰκιζομένους.

\vs{2}Εἷς δὲ αὐτῶν γενόμενος προήγορος, οὕτως ἔφη, τί μέλλεις ἐρωτᾷν, καὶ μανθάνειν παρʼ ἡμῶν; ἕτοιμοι γὰρ ἀποθνήσκειν ἐσμὲν ἢ παραβαίνειν τοὺς πατρίους νόμους.

\vs{3}Ἔκθυμος δὲ γενόμενος ὁ βασιλεὺς, προσέταξε τήγανα, καὶ λέβητας ἐκπυροῦν.
\vs{4}Τῶν δὲ ἐκπυρωθέντων, παραχρῆμα τὸν γενόμενον αὐτῶν προήγορον προσέταξε γλωσσοτομεῖν, καὶ περισκυθίσαντας ἀκρωτηριάζειν, τῶν λοιπῶν ἀδελφῶν, καὶ τῆς μετρὸς, συνορώντων.

\vs{5}Ἄχρηστον δὲ αὐτὸν τοῖς ὅλοις γενόμενον, ἐκέλευσε τῇ πυρᾷ προσάγειν ἔμπνουν, καὶ τηγανίζειν· τῆς δὲ ἀτμίδος ἐφʼ ἱκανὸν διαδιδούσης τοῦ τηγάνου, ἀλλήλους παρεκάλουν σὺν τῇ μητρὶ γενναίως τελευτᾷν, λέγοντες οὕτως.
\vs{6}Ὁ Κύριος ὁ Θεὸς ἐφορᾷ, καὶ ταῖς ἀληθείαις ἐφʼ ἡμῖν παρακαλεῖται, καθάπερ διὰ τῆς κατὰ πρόσωπον ἀντιμαρτυρούσης ᾠδῆς διεσάφησε Μωυσῆς, λέγων, καὶ ἐπὶ τοῖς δούλοις αὐτοῦ παρακληθήσεται.

\vs{7}Μεταλλάξαντος δὲ τοῦ πρώτου τὸν τρόπον τοῦτον, τὸν δεύτερον ἦγον ἐπὶ τὸν ἐμπαιγμόν· καὶ τὸ τῆς κεφαλῆς δέρμα σὺν ταῖς θριξὶ περισύραντες, ἐπηρώτων, εἰ φάγεσαι πρὸ τοῦ τιμωρηθῆναι τὸ σῶμα κατὰ μέλος;

\vs{8}Ὁ δὲ ἀποκριθεὶς τῇ πατρίῳ φωνῇ εἶπεν, οὐχί· διόπερ καὶ
\vs{9}οὗτος τὴν ἑξῆς ἔλαβε βάσανον, ὡς ὁ πρῶτος. Ἐν ἐσχάτῇ δὲ πνοῇ γενόμενος, εἶπε, σὺ μὲν ἀλάστωρ ἐκ τοῦ παρόντος ἡμᾶς ζῇν ἀπολύεις, ὁ δὲ τοῦ κόσμου βασιλεὺς ἀποθανόντας ἡμᾶς ὑπὲρ τῶν αὐτοῦ νόμων εἰς αἰώνιον ἀναβίωσιν ζωῆς ἡμᾶς ἀναστήσει.

\vs{10}Μετὰ δὲ τοῦτον ὁ τρίτος ἐνεπαίζετο, καὶ τὴν γλῶσσαν αἰτηθεὶς ταχέως προέβαλε, καὶ τὰς χεῖρας εὐθαρσῶς προέτεινε,
\vs{11}καὶ γενναίως εἶπεν, ἐξ οὐρανοῦ ταῦτα κέκτημαι, καὶ διὰ τοὺς αὐτοῦ νόμους ὑπερορῶ ταῦτα, καὶ παρʼ αὐτοῦ ταῦτα πάλιν ἐλπίζω κομίσασθαι.
\vs{12}Ὥστε αὐτὸν τὸν βασιλέα καὶ τοὺς σὺν αὐτῷ ἐκπλήσσεσθαι τὴν τοῦ νεανίσκου ψυχὴν, ὡς ἐν οὐδενὶ τὰς ἀλγηδόνας ἐτίθετο.

\vs{13}Καὶ τούτου δὲ μεταλλάξαντος, τὸν τέταρτον ὡσαύτως ἐβασάνιζον αἰκιζόμενοι.
\vs{14}Καὶ γεννόμενος πρὸς τὸ τελευτᾷν, οὕτως ἔφη, αἱρετὸν μεταλλάσσοντας ὑπʼ ἀνθρώπων τὰς ὑπὸ τοῦ Θεοῦ προσδοκᾷν ἐλπίδας, πάλιν ἀναστήσεσθαι ὑπʼ αὐτοῦ· σοὶ μὲν γὰρ ἀνάστασις εἰς ζωὴν οὐκ ἔσται.

\vs{15}Ἐχομένως δὲ τὸν πέμπτον προσάγοντες ᾐκίζοντο.
\vs{16}Ὁ δὲ πρὸς αὐτὸν ἰδὼν, εἶπεν, ἐξουσίαν ἐν ἀνθρώποις ἔχων φθαρτὸς ὤν, ὅ θελεις ποιεῖς· μὴ δόκει δὲ τὸ γένος ἡμῶν ὑπὸ τοῦ Θεοῦ καταλελεῖφθαι.
\vs{17}Σὺ δὲ καρτέρει, καὶ θεώρει τὸ μεγαλεῖον αὐτοῦ κράτος, ὡς σὲ καὶ τὸ σπέρμα σου βασανίσει.

\vs{18}Μετὰ δὲ τοῦτον ἦγον τὸν ἕκτον, καὶ μέλλων ἀποθνήσκειν, ἔφη, μὴ πλανῶ μάτην, ἡμεῖς γὰρ διʼ ἑαυτοὺς ταῦτα πάσχομεν ἁμαρτάνοντες εἰς τὸν ἑαυτῶν Θεὸν, διὸ ἄξια θαυμασμοῦ γέγονε.
\vs{19}Σὺ δὲ μὴ νομίσῃς ἀθῶος ἔσεσθαι, θεομαχεῖν ἐπιχειρήσας.

\vs{20}Ὑπεραγόντως δὲ ἡ μήτηρ θαυμαστὴ καὶ μνήμης ἀγαθῆς ἀξία, ἥτις ἀπολλυμένους υἱοὺς ἑπτὰ συνορῶσα μιᾶς ὑπὸ καιρὸν ἡμέρας, εὐψύχως ἔφερε διὰ τὰς ἐπὶ Κύριον ἐλπίδας.
\vs{21}Ἕκαστον δὲ αὐτῶν παρεκάλει τῇ πατρίῳ φωνῇ, γενναίῳ πεπληρωμένη φρονήματι, καὶ τὸν θῆλυν λογισμὸν ἄρσενι θυμῷ διεγείρασα, λέγουσα πρὸς αὐτοὺς,
\vs{22}οὐδʼ οἶδʼ ὅπως εἰς τὴν ἐμὴν ἐφάνητε κοιλίαν, οὐδὲ ἐγὼ τὸ πνεῦμα καὶ τὴν ζωὴν ὑμῖν ἐχαρισάμην, καὶ τὴν ἑκάστου στοιχείωσιν οὐκ ἐγὼ διερύθμισα.
\vs{23}Τοιγαροῦν ὁ τοῦ κόσμου κτίστης ὁ πλάσας ἀνθρώπου γένεσιν, καὶ πάντων ἐξευρὼν γένεσιν, καὶ τὸ πνεῦμα καὶ τὴν ζωὴν ὑμῖν πάλιν ἀποδώσει μετʼ ἐλέους, ὡς νῦν ὑπερορᾶτε ἑαυτοὺς διὰ τοὺς αὐτοῦνόμους.

\vs{24}Ὁ δὲ Ἀντίοχος οἰόμενος καταφρονεῖσθαι, καὶ τὴν ὀνειδίζουσαν ὑφορώμενος φωνὴν, ἔτι τοῦ νεωτέρου περιόντος, οὐ μόνον διὰ λόγων ἐποιεῖτο τὴν παράκλησιν, ἀλλὰ καὶ διʼ ὅρκων ἐπίστου, ἅμα πλουτιεῖν καὶ μακαριστὸν ποιήσειν μεταθέμενον ἀπὸ τῶν πατρίων νόμων, καὶ φίλον ἕξειν, καὶ χρείας ἐμπιστεύσειν.

\vs{25}Τοῦ δὲ νεανίου μηδαμῶς προσέχοντος, προσκαλεσάμενος ὁ βασιλεὺς τὴν μητέρα, παρήνει τοῦ μειρακίου γενέσθαι σύμβουλον ἐπὶ σωτηρία.
\vs{26}Πολλὰ δὲ αὐτοῦ παραινέσαντος, ἐπεδέξατο πείσειν τὸν υἱόν.

\vs{27}Προσκύψασα δὲ αὐτῷ, χλευάσασα τὸν ὠμὸν τύραννον, οὕτως ἔφησε τῇ πατρῴᾳ φωνῇ, υἱὲ, ἐλέησόν με τὴν ἐν γαστρὶ περιενέγκασάν σε μῆνας ἐννέα, καὶ θηλάσασάν σε ἔτη τρία, καὶ ἐκθρέψασάν σε καὶ ἀγαγοῦσαν εἰς τὴν ἡλικίαν ταύτην, καὶ τροφοφορήσασαν.
\vs{28}Ἀξιῶ σε, τέκνον, ἀναβλέψαντα εἰς τὸν οὐρανὸν καὶ τὴν γῆν, καὶ τὰ ἐν αὐτοῖς πάντα ἰδόντα, γνῶναι ὅτι ἐξ οὐκ ὄντων ἐποίησεν αὐτὰ ὁ Θεὸς, καὶ τὸ τῶν ἀνθρώπων γένος οὕτως γεγένηται,
\vs{29}μὴ φοβηθῇς τὸν δήμιον τοῦτον, ἀλλὰ τῶν ἀδελφῶν ἄξιος γενόμενος, ἐπίδεξαι τὸν θάνατον, ἵνα ἐν τῷ ἐλέει σὺν τοῖς ἀδελφοῖς σου κομίσωμαί σε.

\vs{30}Ἔτι δὲ ταύτης καταλεγούσης ὁ νεανίας εἶπε, τίνα μένετε; οὐχ ὑπακούω τοῦ προστάγματος τοῦ βασιλέως· τοῦ δὲ προστάγματος ἀκούω τοῦ νόμου τοῦ δοθέντος τοῖς πατράσιν ἡμῶν διὰ Μωυσέως.
\vs{31}Σὺ δὲ πάσης κακίας εὑρετὴς γενόμενος εἰς τοὺς Ἑβραίους, οὐ μὴ διαφύγῃς τὰς χεῖρας τοῦ Θεοῦ.

\vs{32}Ἡμεῖς γὰρ διὰ τὰς ἑαυτῶν ἁμαρτίας πάσχομεν.
\vs{33}Εἰ δὲ χάριν ἐπιπλήξεως καὶ παιδείας ὁ ζῶν Κύριος ἡμῶν βραχέως ἐπώργισται, καὶ πάλιν καταλλαγήσεται τοῖς ἑαυτοῦ δούλοις.
\vs{34}Σὺ δὲ ὦ ἀνόσιε, καὶ πάντων ἀνθρώπων μιαρώτατε, μὴ μάτην μετεωρίζου φρυαττόμενος ἀδήλοις ἐλπίσιν, ἐπὶ τοὺς δούλους αὐτοῦ ἐπαιρόμενος χεῖρα.
\vs{35}Οὔπω γὰρ τὴν τοῦ παντοκράτορος ἐπόπτου Θεοῦ κρίσιν ἐκπέφευγας.

\vs{36}Οἱ μὲν γὰρ νῦν ἡμέτεροι ἀδελφοὶ βραχὺν ὑπενέγκαντες πόνον, ἀεννάου ζωῆς ὑπὸ διαθήκην Θεοῦ πεπτώκασι· σὺ δὲ τῇ τοῦ Θεοῦ κρίσει δίκαια τὰ πρόστιμα τῆς ὑπετηφανίας ἀποίσῃ.
\vs{37}Ἐγὼ δὲ καθάπερ οἱ ἀδελφοί μου, καὶ σῶμα καὶ ψυχὴν προδίδωμι περὶ τῶν πατρίων νόμων, ἐπικαλούμενος τὸν Θεὸν ἵλεων ταχὺ τῷ ἔθνει γενέσθαι, καὶ σὲ μετὰ ἐτασμῶν καὶ μαστίγων ἐξομολογήσασθαι, διότι μόνος αὐτὸς Θεός ἐστιν,
\vs{38}ἐν ἐμοὶ δὲ καὶ τοῖς ἀδελφοῖς μου στῆναι τὴν τοῦ παντοκράτορος ὀργὴν τὴν ἐπὶ τὸ σύμπαν ἡμῶν γένος δικαίως ἐπηγμένην.

\vs{39}Ἔκθυμος δὲ γενόμενος ὁ βασιλεὺς, τούτῳ παρὰ τοὺς ἄλλους χειρίστως ἀπήντησε, πικρῶς φέρων ἐπὶ τῷ μυκτηρισμῷ.
\vs{40}Καὶ οὗτος οὖν καθαρὸς τὸν βίον μετήλλαξε, παντελῶς ἐπὶ τῷ Κυρίῳ πεποιθώς.

\vs{41}Ἐσχάτη δὲ τῶν υἱῶν ἡ μήτηρ ἐτελεύτησε.

\vs{42}Τὰ μὲν οὖν περὶ σπλαγχνισμοὺς, καὶ τὰς ὑπερβαλλούσας αἰκίας ἐπὶ τοσοῦτον δεδηλώσθω.

\ch{8}
Ἰούδας δὲ ὁ Μακκαβαῖος καὶ οἱ σὺν αὐτῷ, παρεισπορευόμενοι λεληθότως εἰς τὰς κώμας, προσεκαλοῦντο τοὺς συγγενεῖς, καὶ τοὺς μεμενηκότας ἐν τῷ Ἰουδαϊσμῷ προσλαβόμενοι, συνήγαγον εἰς ἑξακισχιλίους.

\vs{2}Καὶ ἐπεκαλοῦντο τὸν Κύριον ἐπιδεῖν ἐπὶ τὸν ὑπὸ πάντων καταπατούμενον λαόν, οἰκτεῖραι δὲ καὶ τὸν ναὸν, τὸν ὑπὸ τῶν ἀσεβῶν ἀνθρώπων βεβηλωθέντα,
\vs{3}ἐλεῆσαι δὲ καὶ τὴν καταφθειρομένην πόλιν καὶ μέλλουσαν ἰσόπεδον γίνεσθαι, καὶ τῶν καταβοώντων πρὸς αὐτὸν αἱμάτων εἰσακοῦσαι,
\vs{4}μνησθῆναι δὲ καὶ τῆς τῶν ἀναμαρτήτων νηπίων παρανόμου ἀπωλείας, καὶ περὶ τῶν γενομένων εἰς τὸ ὄνομα αὐτοῦ βλασφημιῶν, καὶ μισοπονηρῆσαι.

\vs{5}Γενόμενος δὲ ἐν συστήματι ὁ Μακκαβαῖος, ἀνυπόστατος ἤδη τοῖς ἔθνεσιν ἐγίνετο, τῆς ὀργῆς τοῦ Κυρίου εἰς ἔλεον τραπείσης.
\vs{6}Πόλεις δὲ καὶ κώμας ἀπροσδοκήτως ἐρχόμενος ἐνεπίμπρα, καὶ τοὺς ἐπικαίρους τόπους ἀπολαμβάνων, οὐκ ὀλίγους τῶν πολεμίων ἐνίκα τροπούμενος.
\vs{7}Μάλιστα τὰς νύκτας πρὸς τὰς τοιαύτας ἐπιβουλὰς συνεργοὺς ἐλάμβανε· καὶ λαλιά τις τῆς εὐανδρίας αὐτοῦ διεχεῖτο πανταχῆ.

\vs{8}Συνορῶν δὲ ὁ Φίλιππος κατὰ μικρὸν εἰς προκοπὴν ἐρχόμενον τὸν ἄνδρα πυκνότερον δὲ ἐν ταῖς εὐημερίαις προβαίνοντα, πρὸς Πτολεμαῖον τὸν κοίλης Συρίας καὶ Φοινίκης στρατηγὸν ἔγραψεν ἐπιβοηθεῖν τοῖς τοῦ βασιλέως πράγμασιν.

\vs{9}Ὁ δὲ ταχέως προχειρισάμενος, Νικάνορα τὸν τοῦ Πατρόκλου, τῶν πρώτων φίλως, ἀπέστειλεν, ὑποτάξας παμφύλων ἔθνη οὐκ ἐλάττους τῶν δισμυρίων, τὸ σύμπαν τῶν Ἰουδαίως ἐξᾶραι γένος· συνέστησε δὲ αὐτῷ καὶ Γοργίαν ἄνδρα στρατηγὸν, καὶ ἐν πολεμικαῖς χρείαις πείραν ἔχοντα.

\vs{10}Διεστήσατο δὲ ὁ Νικάνωρ τὸν φόρον τῷ βασιλεῖ τοῖς Ῥωμαίοις ὄντα ταλάντων δισχιλίων ἐκ τῆς τῶν Ἰουδαίων αἰχμαλωσίας ἐκπληρώσειν.
\vs{11}Εὐθέως δὲ εἰς τὰς παραθαλασσίους πόλεις ἀπέστειλε προσκαλούμενος ἐπʼ ἀγορασμὸν Ἰουδαϊκῶν σωμάτων, ὑπισχνούμενον ἐννενήκοντα σώματα ταλάντου παραχωρήσειν· οὐ προσδεχόμενος τὴν παρὰ τοῦ παντοκράτορος μέλλουσαν παρακολουθήσειν ἐπʼ αὐτῷ δίκην.

\vs{12}Τῷ δὲ ʼΙούδᾳ προσέπεσε περὶ τῆς τοῦ Νικάνορος ἐφόδου· καὶ μεταδόντος αὐτοῦ σὺν αὐτῷ τὴν παρουσίαν τοῦ στρατοπέδου,
\vs{13}οἱ δει λανδροῦντες καὶ ἀπιστοῦντες τὴν τοῦ Θεοῦ δίκην, διεδίδρασκον, καὶ ἐξετόπιζον ἑαυτούς.

\vs{14}Οἱ δὲ τὰ περιλελειμμένα πάντα ἐπώλουν, ὁμοῦ δὲ τὸν Κύριον ἠξίουν ῥύσασθαι τοὺς ὑπὸ τοῦ δυσσεβοῦς Νικάνορος πρὶν συντυχεῖν πεπραμένους.
\vs{15}Καὶ εἰ μὴ διʼ αὐτούς, ἀλλὰ διὰ τὰς πρὸς τοὺς πατέρας αὐτῶν διαθήκας, καὶ ἕνεκεν τῆς ἐπʼ αὐτοὺς ἐπικλήσεως τοῦ σεμνοῦ καὶ μεγαλοπρεποῦς ὀνόματος αὐτοῦ.

\vs{16}Συναγαγὼν δὲ ὁ Μακκαβαῖος τοὺς περὶ αὐτὸν ὄντας τὸν ἀριθμὸν ἑξακισχιλίους, παρεκάλει μὴ καταπλαγῆναι τοὺς πολεμίους, μηδὲ εὐλαβεῖσθαι τὴν τῶν ἀδίκως παραγινομένων ἐπʼ αὐτοὺς ἐθνῶν πολυπληθίαν, ἀγωνίσασθαι δὲ γενναίως,
\vs{17}πρὸ ὀφθαλμῶν λαβόντας τὴν ἀνόμως εἰς τὸν ἅγιον τόπον συντετελεσμένην ὑπʼ αὐτῶν ὕβριν, καὶ τὸν τῆς ἐμπεπαιγμένης πόλεως αἰκισμόν, ἔτι δὲ τὴν τῆς προγονικῆς πολιτείας κατάλυσιν.
\vs{18}Οἱ μὲν γὰρ ὅπλοις πεποίθασιν ἅμα καὶ τόλμαις, ἔφησεν, ἡμεῖς δὲ ἐπὶ τῷ παντοκράτορι Θεῷ δυναμένῳ καὶ τοὺς ἐρχομένους ἐφʼ ἡμᾶς, καὶ τὸν ὅλον κόσμον ἐν ἑνὶ νεύματι καταβαλεῖν, πεποίθαμεν.

\vs{19}Προσαναλεξάμενος δὲ αὐτοῖς καὶ τὰς ἐπὶ τῶν προγόνων γενομένας ἀντιλήψεις, καὶ τὴν ἐπὶ Σενναχηρεὶμ τῶν ἑκατὸν ὀγδοήκοντα πέντε χιλιάδων ὡς ἀπώλοντο.
\vs{20}Καὶ τὴν ἐν τῇ Βαβυλωνίᾳ τὴν πρὸς αὐτοὺς Γαλάτας παράταξιν γενομένην, ὡς οἱ πάντες ἐπὶ τὴν χρείαν ἦλθον ὀκτακισχιλιοι σὺν Μακεδόσι τετρακισχιλίοις, τῶν Μακεδόνων ἀπορουμένων, οἱ ὀκτακισχίλιοι τὰς δώδεκα μυρίαδας ἀπώλεσαν διὰ τὴν γενομένην αὐτοῖς ἀπʼ οὐρανοῦ βοήθειαν, καὶ ὠφέλειαν πολλὴν ἔλαβον.

\vs{21}Ἐφʼ οἷς εὐθαρσεῖς αὐτοὺς παραστήσας, καὶ ἑτοίμους ὑπὲρ τῶν νόμων καὶ τῆς πατρίδος ἀποθνήσκειν, τετραμερές τι τὸ στράτευμα ἐποίησε·
\vs{22}τάξας καὶ τοὺς ἀδελφοὺς αὐτοῦ προηγουμένους ἑκατέρας τάξεως, Σίμωνα καὶ Ἰώσηφον καὶ Ἰωνάθαν, ὑποτάξας ἑκάστῳ χιλίους πρὸς τοῖς πεντακοσίοις,
\vs{23}ἔτι δὲ καὶ Ἐλεάζαρον, παραγνοὺς τὴν ἱερὰν βίβλον, καὶ δοὺς σύνθημα Θεοῦ βοηθείας, τῆς πρώτης σπείρας αὐτὸς προηγούμενος, συνέβαλε τῷ Νικάνορι.

\vs{24}Γενομένου δὲ αὐτοῖς τοῦ παντοκράτορος συμμάχου, κατέσφαξαν τῶν πολεμίων ὑπὲρ τοὺς ἐννακισχιλίους, τραυματίας δὲ καὶ τοῖς μέλεσιν ἀναπήρους τὸ πλεῖστον μέρος τῆς τοῦ Νικάνορος στρατιᾶς ἐποίησαν, πάντας δὲ φυγεῖν ἠνάγκασαν.
\vs{25}Τὰ δὲ χρήματα τῶν παραγεγονότων ἐπὶ τὸν ἀγορασμὸν αὐτῶν ἔλαβον· συνδιώξαντες δὲ αὐτοὺς ἐφʼ ἱκανὸν, ἀνέλυσαν ὑπὸ τῆς ὥρας συγκλειόμενοι.
\vs{26}Ἦν γὰρ ἡ πρὸ τοῦ σαββάτου, διʼ ἥν αἰτίαν οὐκ ἐμακροθύμησαν κατατρέχοντες αὐτούς.

\vs{27}Ὁπλολογήσαντες δὲ αὐτοὺς, καὶ τὰ σκῦλα ἐκδύσαντες τῶν πολεμίων, περὶ τὸ σάββατον ἐγίνοντο, περισσῶς εὐλογοῦντες, καὶ ἐξομολογούμενοι τῷ Κυρίῳ τῷ διασώσαντι αὐτοὺς εἰς τὴν ἡμέραν ταύτην, ἀρχὴν ἐλέους τάξαντος αὐτοῖς.

\vs{28}Μετὰ δὲ τὸ σάββατον τοῖς ᾐκισμένοις, καὶ ταῖς χήραις, καὶ ὀρφανοῖς, μερίσαντες ἀπὸ τῶν σκύλων, τὰ λοιπὰ αὐτοὶ καὶ τὰ παιδία ἐμερίσαντο.
\vs{29}Ταῦτα δὲ διαπραξάμενοι, καὶ κοινὴν ἱκετείαν ποιησάμενοι, τὸν ἐλεήμονα Κύριον ἠξίουν εἰς τέλος, καταλλαγῆναι τοῖς αὐτοῦ δούλοις.

\vs{30}Καὶ τῶν περὶ Τιμόθεον καὶ Βακχίδην συνεριζοντων, ὑπὲρ τοὺς δισμυρίους αὐτῶν ἀνεῖλον, καὶ ὀχυρωμάτων ὑψηλῶν εὖ μάλα ἐγκρατεῖς ἐγένοντο· καὶ λάφυρα πλεῖστα ἐμερίσαντο, ἰσομοίρους ἑαυτοὺς καὶ τοῖς ᾐκισμένοις, καὶ ὀρφανοῖς, καὶ χήραις, ἔτι δὲ καὶ πρεσβυτέροις ποιήσαντες.
\vs{31}Ὁπλολογήσαντες δὲ αὐτοὺς, ἐπιμελῶς πάντα συνέθηκαν εἰς τοὺς ἐπικαίρους τόπους, τὰ δὲ λοιπὰ τῶν σκύλων ἤνεγκαν εἰς Ἱεροσόλυμα.

\vs{32}Τὸν δὲ φυλάρχην τῶν περὶ Τιμόθεον ἀνεῖλον, ἀνοσιώτατον ἄνδρα καὶ πολλὰ τοὺς Ἰουδαίους ἐπιλελυπηκότα.
\vs{33}Ἐπινίκια δὲ ἄγοντες ἐν τῇ πατρίδι, τοὺς ἐμπρήσαντας τοὺς ἱεροὺς πυλῶνας, Καλλισθένην, καί τινας ἄλλους ὑφῆψαν εἰς ἓν οἰκίδιον πεφευγότας, οἵ τινες ἄξιον τῆς δυσσεβείας ἐκομίσαντομισθόν.

\vs{34}Ὁ δὲ τρισαλιτήριος Νικάνωρ, ὁ τοὺς χιλίους ἐμπόρους ἐπὶ τὴν πράσιν τῶν Ἰουδαίων ἀγαγὼν,
\vs{35}ταπεινωθεὶς ὑπὸ τῶν κατʼ αὑτὸν νομιζομένων ἐλαχίστων εἶναι, τῇ τοῦ Κυρίου βοηθείᾳ, τὴν δοξικὴν ἀποθέμενος ἐσθῆτα, διὰ τῆς μεσογείου, δραπέτου τρόπον ἔρημον ἑαυτὸν ποιήσας, ἧκεν εἰς Ἀντιόχειαν, ὑπεράγαν δυσημερήσας ἐπὶ τῇ τοῦ στρατοῦ διαφθορᾷ.
\vs{36}Καὶ ὁ τοῖς Ῥωμαίοις ἀναδεξάμενος φόρον ἀπὸ τῆς τῶν ἐν Ἱεροσολύμοις αἰχμαλωσίας κατορθώσασθαι, κατήγγελλεν ὑπέρμαχον ἔχειν τὸν Θεὸν τοὺς Ἰουδαίους, καὶ διὰ τὸν τρόπον τοῦτον ἀτρώτους εἶναι τοὺς Ἰουδαίους, διὰ τὸ ἀκολουθεῖν τοῖς ὑπʼ αὐτοῦ προτεταγμένοις νόμοις.

\ch{9}
Περὶ δὲ τὸν καιρὸν ἐκεῖνον ἐτύγχανεν Ἀντίοχος ἀναλελυκὼς ἀκόσμως ἐκ τῶν κατὰ τὴν Περσίδα τόπων.
\vs{2}Εἰσεληλύθει γὰρ εἰς τὴν λεγομένην Περσέπολιν, καὶ ἐπεχείρησεν ἱεροσυλεῖν, καὶ τὴν πόλιν συνέχειν· διὸ δὴ τῶν πληθῶν ὁρμησάντων, ἐπὶ τὴν τῶν ὅπλων βοήθειαν ἐτράπησαν· καὶ συνέβη τροπωθέντα τὸν Ἀντίοχον ὑπὸ τῶν ἐγχωρίων, ἀσχήμονα τὴν ἀναζυγὴν ποιήσασθαι.

\vs{3}Ὄντι δὲ αὐτῷ κατʼ Ἐκβάτανα, προσέπεσε τὰ κατὰ Νικάνορα, καὶ τοὺς περὶ Τιμόθεον, γεγονότα.
\vs{4}Ἐπαρθεὶς δὲ τῷ θυμῷ, ᾤετο καὶ τὴν τῶν πεφυγαδευκότων αὐτὸν κακίαν εἰς τοὺς Ἰουδαίους ἐναπερείσασθαι· διὸ συνέταξε τὸν ἁρματηλάτην ἀδιαλείπτως ἐλαύνοντα κατανύειν τὴν πορείαν, τῆς ἐξ οὐρανοῦ δὴ κρίσεως συνούσης αὐτῷ· οὕτως γὰρ ὑπερηφάνως εἶπε, πολυάνδριον Ἰουδαίων Ἱεροσόλυμα ποιήσω παραγενόμενος ἐκεῖ.

\vs{5}Ὁ δὲ πανεπόπτης Κύριος ὁ Θεὸς τοῦ Ἰσραὴλ ἐπάταξεν αὐτὸν ἀνιάτῳ καὶ ἀοράτῳ πληγῇ· ἄρτι δὲ αὐτοῦ καταλήξαντος τὸν λόγον, ἔλαβεν αὐτὸν ἀνήκεστος τῶν σπλάγχνων ἀλγηδὼν, καὶ πικραὶ τῶν ἔνδον βάσανοι,
\vs{6}πάνυ δικαίως, τὸν πολλαῖς καὶ ξενιζούσαις συμφοραῖς ἑτέρων σπλάγχνα βασανίσαντα.

\vs{7}Ὁ δʼ οὐδαμῶς τῆς ἀγερωχίας ἔληγεν· ἔτι δὲ καὶ τῆς ὑπερηφανίας ἐπεπλήρωτο, πῦρ πνέων τοῖς θυμοῖς ἐπὶ τοὺς Ἰουδαίους, καὶ κελεύων ἐποξύνειν τὴν πορείαν· συνέβη δὲ καὶ πεσεῖν αὐτὸν ἀπὸ τοῦ ἅρματος φερομένου ῥοίζῳ, καὶ δυσχερεῖ πτώματι περιπεσόντα, πάντα τὰ μέλη τοῦ σώματος ἀποστρεβλοῦσθαι.

\vs{8}Ὁ δʼ ἄρτι δοκῶν τοῖς τῆς θαλάσσης κύμασιν ἐπιτάσσειν, διὰ τὴν ὑπὲρ ἄνθρωπον ἀλαζονείαν, καὶ πλάστιγγι τὰ τῶν ὀρέων οἰόμενος ὕψη στήσειν, κατὰ γῆν γενόμενος, ἐν φορείῳ παρεκομίζετο, φανερὰν τοῦ Θεοῦ πᾶσι τὴν δύναμιν ἐνδεικνύμενος·
\vs{9}ὥστε καὶ ἐκ τοῦ σώματος τοῦ δυσσεβοῦς σκώληκας ἀναζεῖν, καὶ ζῶντος ἐν ὀδύναις καὶ ἀλγηδόσι τὰς σάρκας αὐτοῦ διαπίπτειν, ὑπὸ δὲ τῆς ὀσμῆς αὐτοῦ πᾶν τὸ στρατόπεδον βαρύνεσθαι τῇν σαπρίᾳ.
\vs{10}Καὶ τὸν μικρῷ πρότερον τῶν οὐρανίων ἄστρων ἅπτεσθαι δοκοῦντα, παρακομίζειν οὐδεὶς ἐδύνατο, διὰ τὸ τῆς ὀσμῆς ἀφόρητον βάρος.

\vs{11}Ἐνταῦθα οὖν ἤρξατο τὸ πολὺ τῆς ὑπερηφανίας λήγειν ὑποτεθραυσμένος, καὶ εἰς ἐπίγνωσιν ἔρχεσθαι θείᾳ μάστιγι κατὰ στιγμὴν ἐπιτεινόμενος ταῖς ἀλγηδόσι.
\vs{12}Καὶ μηδὲ τῆς ὀσμῆς αὐτοῦ δυνάμενος ἀνέχεσθαι, ταῦτʼ ἔφη, δίκαιον ὑποτάσσεσθαι τῷ Θεῷ, καὶ μὴ θνητὸν ὄντα ἰσόθεα φρονεῖν ὑπερηφανῶς.

\vs{13}Ηὔχετο δὲ ὁ μιαρὸς πρὸς τὸν οὐκέτι αὐτὸν ἐλεήσοντα δεσπότην, οὕτω λέγων,
\vs{14}τὴν μὲν ἁγίαν πόλιν ἣν σπεύδων παρεγίνετο ἰσόπεδον ποιῆσαι, καὶ πολυάνδριον οἰκοδομῆσαι, ἐλευθέραν ἀναδεῖξαι·
\vs{15}τοὺς δὲ Ἰουδαίους, οὕς διεγνώκει μηδὲ ταφῆς ἀξιῶσαι, οἰωνοβρώτους δὲ σὺν τοῖς νηπίοις ἐκρίψειν θηρίοις, πάντας αὐτοὺς ἴσους Ἀθηναίοις ποιήσειν·
\vs{16}ὃν δὲ πρότερον ἐσκύλευσεν ἅγιον νεὼν, καλλίστοις ἀναθήμασι κοσμήσειν, καὶ τὰ ἱερὰ σκεύη πολυπλάσια πάντα ἀποδώσειν, τὰς δὲ ἐπιβαλλούσας πρὸς τὰς θυσίας συντάξεις ἐκ τῶν ἰδίων προσόδων χορηγήσειν·
\vs{17}πρὸς δὲ τούτοις, καὶ Ἰουδαῖο ἔσεσθαι, καὶ πάντα τόπον οἰκητὸν ἐπελεύσεσθαι καταγγέλλοντα τὸ τοῦ Θεοῦ κράτος.

\vs{18}Οὐδαμῶς δὲ ληγόντων τῶν πόνων, ἐπεληλύθει γὰρ ἐπʼ αὐτὸν δικαία ἡ τοῦ Θεοῦ κρίσις, τὰ κατʼ αὐτὸν ἀπελπίσας, ἔγραψε πρὸς τοὺς Ἰουδαίους τὴν ὑπογεγραμμένην ἐπιστολήν, ἱκετηρίας τάξιν ἔχουσαν, περιέχουσαν δὲ οὕτως·

\vs{19}Τοῖς χρηστοῖς Ἰουδαίοις τοῖς πολίταις πολλὰ χαίρειν, καὶ ὑγιαίνειν, καὶ εὖ πράττειν, βασιλεὺς καὶ στρατηγὸς Ἀντίοχος.
\vs{20}Εἰ ἔῤῥωσθε, καὶ τὰ τέκνα καὶ τὰ ἴδια κατὰ γνώμην ἔστιν ὑμῖν, εὔχομαι μὲν τῷ Θεῷ τὴν μεγίστην χάριν, εἰς οὐρανὸν τὴν ἐλπίδα ἔχων.

\vs{21}Κᾀγὼ δὲ ἀσθενῶς διεκείμην, ὑμῶν τὴν τιμὴν καὶ τὴν εὔνοιαν ἄν ἐμνημόνευον φιλοστόργως· ἐπανάγων ἐκ τῶν περὶ τὴν Περσίδα τόπων, καὶ περιπεσὼν ἀσθενείᾳ δυσχέρειαν ἐχούσῃ, ἀναγκαῖον ἡγησάμην φροντίσαι τῆς κοινῆς πάντων ἀσφαλείας·
\vs{22}Οὐκ ἀπογινώσκων τὰ κατʼ ἐμαυτόν, ἀλλὰ ἔχων πολλὴν ἐλπίδα ἐκφεύξεσθαι τὴν ἀσθένειαν,
\vs{23}θεωρῶν δὲ ὅτι καὶ ὁ πατήρ καθʼ οὓς καιροὺς εἰς τοὺς ἄνω τόπους ἐστρατοπέδευσεν, ἀνέδειξε τὸν διαδεξόμενον,
\vs{24}ὅπως ἐάν τι παράδοξον ἀποβαίνῃ, ἤ καὶ προσαγγελθῇ τι δυοχερὲς, εἰδότες οἱ κατὰ τὴν χώραν ᾧ καταλέλειπται τὰ πράγματα, μὴ ἐπιταράσσωνται·

\vs{25}Πρὸς δὲ τούτοις κατανοῶν τοὺς παρακειμένους δυνάστας, καὶ γειτνιῶντας τῇ βασιλείᾳ τοῖς καιροῖς ἐπέχοντας, προσδεχομένους τὸ ἀποβησόμενον, ἀναδέδειχα τὸν υἱὸν μου Ἀντίοχον βασιλέα, ὃν πολλάκις ἀνατρέχων εἰς τὰς ἐπάνω σατραπείας τοῖς πλείστοις ὑμῶν παρακατετιθέμην καὶ συνίστων· γέγραφα δὲ πρὸς αὐτὸν τὰ ὑπογεγραμμένα.

\vs{26}Παρακαλῶ οὖν ὑμᾶς καὶ ἀξιῶ, μεμνημένους τῶν εὐεργεσιῶν κοινῇ καὶ κατιδίαν, ἕκαστον συντηρεῖν τὴν οὖσαν εὔνοιαν εἰς ἐμὲ καὶ τὸν υἱόν μου.
\vs{27}Πέπεισμαι γὰρ αὐτὸν ἐπιεικῶς καὶ φιλανθρώπως παρακολουθοῦντα τῇ ἐμῇ προαιρέσει, συμπεριενεχθήσεσθαι ὑμῖν.

\vs{28}Ὁ μὲν οὖν ἀνδροφόνος καὶ βλάσφημος τὰ χείριστα παθών, ὡς ἑτέρους διέθηκεν, ἐπὶ ξένης ἐν τοῖς ὄρεσιν οἰκτίστῳ μόρῳ κατέσπρεψε τὸν βίον.
\vs{29}Παρεκομίζετο δὲ τὸ σῶμα Φίλιππος ὁ σύντροφος αὐτοῦ· ὃς καὶ διευλαβηθεὶς τὸν υἱὸν Ἀντιόχου, πρὸς Πτολεμαῖον τὸν Φιλομήτορα εἰς Αἴγυπτον διεκομίσθη.

\ch{10}
Μακκαβαῖος δὲ καὶ οἱ σὺν αὐτῷ, τοῦ Κυρίου προάγοντος αὐτοὺς, τὸ μὲν ἱερὸν ἐκομίσαντο καὶ τὴν πόλιν,
\vs{2}τοὺς δὲ κατὰ τὴν ἀγορὰν βωμοὺς ὑπὸ τῶν ἀλλοφύλων δεδημιουργημένους, ἔτι δὲ τεμένη καθεῖλον.

\vs{3}Καὶ τὸν νεὼν καθαρίσαντες, ἕτερον θυσιαστήριον ἐποίησαν, καὶ πυρώσαντες λίθους, καὶ πῦρ ἐκ τούτων λαβόντες, ἀνήνεγκαν θυσίαν μετὰ διετῆ χρόνον, καὶ θυμίαμα καὶ λύχνους, καὶ τῶν ἄρτων τὴν πρόθεσιν ἐποιήσαντο.
\vs{4}Ταῦτα δὲ ποιήσαντες ἠξίωσαν τὸν Κύριον πεσόντες ἐπὶ κοιλίαν, μηκέτι περιπεσεῖν τοιούτοις κακοῖς, ἀλλʼ ἐάν ποτε καὶ ἁμάρτωσιν, ὑπʼ αὐτοῦ μετʼ ἐπιεικείας παιδεύεσθαι, καὶ μὴ βλασφήμοις καὶ βαρβάροις ἔθνεσι παραδίδοσθαι.

\vs{5}Ἐν ᾗ δὲ ἡμέρᾳ ὁ νεὼς ὑπὸ ἀλλοφύλων ἐβεβηλώθη, συνέβη κατὰ τὴν αὐτὴν ἡμέραν τὸν καθαρισμὸν γενέσθαι τοῦ ναοῦ, τῇ πέμπτῃ καὶ εἰκάδι τοῦ αὐτοῦ μηνὸς, ὅς ἐστι Χασελεῦ.

\vs{6}Καὶ μετʼ εὐφροσύνης ἦγον ἡμέρας ὀκτὼ σκηνωμάτων τρόπον, μνημονεύοντες ὡς πρὸ μικροῦ χρόνου τὴν τῶν σκηνῶν ἑορτὴν ἐν τοῖς ὄρεσι καὶ ἐν τοῖς σπηλαίοις θηρίων τρόπον ἦσαν νεμόμενοι.
\vs{7}Διὸ θύρσους καὶ κλάδους ὡραίους, ἔτι δὲ φοίνικας ἔχοντες, ὕμνους ἀνέφερον τῷ εὐοδώσαντι καθαρισθῆναι τὸν ἑαυτοῦ τόπον.
\vs{8}Ἐδογμάτισαν δὲ μετὰ κοινοῦ προστάγματος καὶ ψηφίσματος παντὶ τῷ τῶν Ἰουδαίων ἔθνει κατʼ ἐνιαυτὸν ἄγειν τάσδε τὰς ἡμέρας.

\vs{9}Καὶ τὰ μὲν τῆς Ἀντιόχου τοῦ προσαγορευθέντος Ἐπιφανοῦς τελευτῆς οὕτως εἶχε.

\vs{10}Νυνὶ δὲ τὰ κατὰ τὸν Εὐπάτορα Ἀντίοχον, υἱὸν δὲ τοῦ ἀσεβοῦς γενόμενον, δηλώσομεν, αὐτά συντέμνοντες τὰ τῶν πολέμων κακά.
\vs{11}Αὐτὸς γὰρ παραλαβὼν βασιλείαν, ἀνέδειξεν ἐπὶ τῶν πραγμάτων Λυσίαν τινά, κοίλης δὲ Συρίας καὶ Φοινίκης στρατηγὸν πρώταρχον.

\vs{12}Πτολεμαῖος γὰρ ὁ καλούμενος Μάκρων τὸ δίκαιον συντηρεῖν προηγούμενος εἰς τοὺς Ἰουδαίους διὰ τὴν γεγονυῖαν εἰς αὐτοὺς ἀδικίαν, καί ἐπειρᾶτο τὰ πρὸς αὐτοὺς εἰρηνικῶς διεξάγειν.
\vs{13}Ὅθεν κατηγορούμενος ὑπὸ τῶν φίλων πρὸς τὸν Εὐπάτορα, καὶ προδότης παρέκαστα ἀκούων, διὰ τὸ τὴν Κύπρον ἐμπιστευθέντα ὑπὸ τοῦ Φιλομήτορος ἐκλιπεῖν, καὶ πρὸς Ἀντίοχον τὸν Ἐπιφανῆ ἀναχωρῆσαι, μήτʼ εὐγενῆ τὴν ἐξουσίαν ἔχων, ὑπʼ ἀθυμίας φαρμακεύσας ἑαυτὸν ἐξέλιπε τὸν βίον.

\vs{14}Γοργίας δὲ γενόμενος στρατηγὸς τῶν τόπων, ἐξενοτρόφει, καὶ παρέκαστα πρὸς τοὺς Ἰουδαίους ἐπολεμοτρόφει,
\vs{15}Ὁμοῦ δὲ τούτῳ καὶ οἱ Ἰδουμαῖοι ἐγκρατεῖς ἐπικαίρων ὀχυρωμάτων ὄντες, ἐγύμναζον τοὺς Ἰουδαίους, καὶ τοὺς φυγαδευθέντας ἀπὸ Ἰεροσολύμων προσλαβόμενοι πολεμοτροφεῖν ἐπεχείρουν.

\vs{16}Οἱ δὲ περὶ τὸν Μακκαβαῖον, ποιησάμενοι λιτανείαν, καὶ ἀξιώσαντες τὸν Θεὸν σύμμαχον αὐτοῖς γενέσθαι, ἐπὶ τὰ τῶν Ἰδουμαίων ὀχυρώματα ὥρμησαν,
\vs{17}οἷς καὶ προσβαλόντες εὐρώστως, ἐγκρατεῖς ἐγένοντο τῶν τόπων, πάντας τε τοὺς ἐπὶ τῷ τείχει μαχομένους ἠμύναντο· κατέσφαζόν δε τοὺς ἐμπίπτοντας, ἀνεῖλον δὲ οὐχ ἧττον τῶν δισμυρίων.

\vs{18}Συμφυγόντων δὲ οὐκ ἔλαττον τῶν ἐννακισχιλίων εἰς δύο πύργους ὀχυροὺς εὖ μάλα, καὶ πάντα τὰ πρὸς πολιορκίαν ἔχοντας,
\vs{19}ὁ Μακκαβαῖος εἰς ἐπείγοντας τόπους ἀπολιπὼν Σίμωνα καὶ Ἰώσηφον, ἔτι δὲ καὶ Ζακχαῖον καὶ τοὺς σὺν αὐτῷ ἱκανοὺς πρὸς τὴν τούτων πολιορκίαν, αὐτὸς ἐχωρίσθη.

\vs{20}Οἱ δὲ περὶ τὸν Σίμωνα φιλαργυρήσαντες ὑπό τινων τῶν ἐν τοῖς πύργοις ἐπείσθησαν ἀργυρίῳ· ἑπτάκις δὲ μυριάδας δραχμὰς λαβόντες, εἴασάν τινας διαῤῥυῆναι.
\vs{21}Προσαγγελθέντος δὲ τῷ Μακκαβαίῳ περὶ τοῦ γεγονότος, συναγαγὼν τοὺς ἡγουμένους τοῦ λαοῦ, κατηγόρησεν ὡς ἀργυρίου πεπράκασι τοὺς ἀδελφοὺς, τοὺς πολεμίους κατʼ αὐτῶν ἀπολύσαντες.
\vs{22}Τούτους μὲν οὖν προδότας γενομένους ἀπέκτεινε, καὶ παραχρῆμα τοὺς δύο πύργους κατελάβετο.
\vs{23}Τοῖς δὲ ὅπλοις τὰ πάντα ἐν ταῖς χερσὶν εὐοδούμενος, ἀπώλεσεν ἐν τοῖς δυσὶν ὀχυρώμασι πλείους τῶν δισμυρίων.

\vs{24}Τιμόθεος δὲ ὁ πρότερον ἡττηθεὶς ὑπὸ τῶν Ἰουδαίων, συναγαγὼν ξένας δυνάμεις παμπληθεῖς, καὶ τοὺς τῆς Ἀσίας γενομένους ἵππους συναθροίσας οὐκ ὀλίγους, παρῆν ὡς δοριάλωτον ληψόμενος τὴν Ἰουδαίαν.
\vs{25}Οἱ δὲ περὶ τὸν Μακκαβαῖον, συνεγγίζοντος αὐτοῦ, πρὸς ἱκετείαν τοῦ Θεοῦ ἐτράπησαν, γῇ τὰς κεφαλὰς καταπάσαντες, καὶ τὰς ὀσφύας σάκκοις ζώσαντες,
\vs{26}ἐπὶ τὴν ἀπέναντι τοῦ θυσιαστηρίου κρηπίδα προσπεσόντες, ἠξίουν ἵλεων αὐτοῖς γενόμενον ἐχθρεῦσαι τοῖς ἐχθροῖς αὐτῶν, καὶ ἀντικεῖσθαι τοῖς ἀντικειμένοις, καθὼς ὁ νόμος διασαφεῖ.
\vs{27}Γενόμενοι δὲ ἀπὸ τῆς δεήσεως, ἀναλαβόντες τὰ ὅπλα, προῆγον ἀπὸ τῆς πόλεως ἐπὶ πλεῖον· συνεγγίσαντες δὲ τοῖς πολεμίοις, ἐφʼ ἑαυτῶν ἦσαν.

\vs{28}Ἄρτι δὲ τῆς ἀνατολῆς διαδεχομένης, προσέβαλον ἑκάτεροι· οἱ μὲν ἔγγυον ἔχοντες εὐημερίας καὶ νίκης μετʼ ἀρετῆς τὴν ἐπὶ τὸν Κύριον καταφυγὴν, οἱ δὲ καθηγεμόνα τῶν ἀγώνων ταττόμενοι τὸν θυμόν.

\vs{29}Γενομένης δὲ καρτερᾶς μάχης, ἐφάνησαν τοῖς ὑπεναντίοις ἐξ οὐρανοῦ ἐφʼ ἵππων χρυσοχαλίνων ἄνδρες πέντε διαπρεπεῖς, καὶ ἀφηγούμενοι τῶν Ἰουδαίων οἱ δύο,
\vs{30}καὶ τὸν Μακκαβαῖον μέσον λαβόντες, καὶ σκεπάζοντες ταῖς ἑαυτῶν πανοπλίαις, ἄτρωτον διεφύλαττον· εἰς δὲ τοὺς ὑπεναντίους τοξεύματα καὶ κεραυνοὺς ἐξεῤῥίπτουν· διὸ συγχυθέντες ἀορασίᾳ, κατεκόπτοντο ταραχῆς πεπληρωμένοι.
\vs{31}Κατεσφάγησαν δὲ δισμύριοι πρὸς τοῖς πεντακοσίοις, ἱππεῖς δὲ ἑξακόσιοι.

\vs{32}Αὐτὸς δὲ ὁ Τιμόθεος συνέφυγεν εἰς Γάζαρα λεγόμενον ὀχύρωμα, εὖ μάλα φρούριον, στρατηγοῦντος ἐκεῖ Χαιρέου.

\vs{33}Οἱ δὲ περὶ τὸν Μακκαβαῖον ἄσμενοι περιεκάθισαν τὸ φρούριον ἡμέρας τέσσαρας.
\vs{34}Οἱ δὲ ἔνδον τῇ ἐρυμνότητι τοῦ τόπου πεποιθότες, ὑπεράγαν ἐβλασφήμουν, καὶ λόγους ἀθεμίτους προΐοντο.

\vs{35}Ὑποφαινούσης δὲ τῆς πέμπτης ἡμέρας, εἴκοσι νεανίαι τῶν περὶ τὸν Μακκαβαῖον πυρωθέντες τοῖς θυμοῖς διὰ τὰς βλασφημίας, προσβαλόντες τῷ τείχει, ἀῤῥενωδῶς καὶ θηριώδει θυμῷ τὸν ἐμπίπτοντα ἔκοπτον,
\vs{36}ἕτεροι δὲ ὁμοίως προσαναβάντες ἐν τῷ περισπασμῷ πρὸς τοὺς ἔνδον, ἐνεπίμπρων τοὺς πύργους, καὶ πυρὰς ἀνάψαντες ζῶντας τοὺς βλασφήμους κατέκαιον· οἱ δὲ τὰς πύλας διέκοπτον, εἰσδεξάμενοι δὲ τὴν λοιπὴν τάξιν, προκατελάβοντο τὴν πόλιν,
\vs{37}καὶ τὸν Τιμόθεον ἀποκεκρυμμένον ἔν τινι λάκκῳ κατέσφαξαν, καὶ τὸν τούτου ἀδελφὸν Χαιρέαν, καὶ τὸν Ἀπολλοφάνη.

\vs{38}Ταῦτα δὲ διαπραξάμενοι, μεθʼ ὕμνων καὶ ἐξομολογήσεων εὐλόγουν τῷ Κυρίῳ τῷ μεγάλως εὐεργετοῦντι τὸν Ἰσραὴλ, καὶ τὸ νῖκος αὐτοῖς διδόντι.

\ch{11}
Μετʼ ὀλίγον δὲ παντελῶς χρόνον Λυσίας ἐπίτροπος τοῦ βασιλέως καὶ συγγενὴς, καὶ ἐπὶ τῶν πραγμάτων, λίαν βαρέως φέρων ἐπὶ τοῖς γεγονόσι,
\vs{2}συναθροίσας περὶ τὰς ὀκτὼ μυριάδας καὶ τὴν ἵππον πᾶσαν, παρεγένετο ἐπὶ τοὺς Ἰουδαίους, λογιζόμενος τὴν μὲν πόλιν Ἕλλησιν οἰκητήριον ποιήσειν,
\vs{3}τὸ δὲ ἱερὸν ἀργυρολόγητον καθὼς τὰ λοιπὰ τῶν ἐθνῶν τεμένη, πρατὴν δὲ τὴν ἀρχιερωσύνην κατʼ ἔτος ποιήσειν,
\vs{4}οὐδαμῶς ἐπιλογιζόμενος τὸ τοῦ Θεοῦ κράτος, πεφρενωμένος δὲ ταῖς μυριάσι τῶν πεζῶν καὶ ταῖς χιλιάσι τῶν ἱππέων καὶ τοῖς ἐλέφασι τοῖς ὀγδοήκοντα.

\vs{5}Εἰσελθὼν δὲ εἰς τὴν Ἰουδαίαν, καὶ συνεγγίσας τῷ Βαιθσούρᾳ, ὄντι μὲν ἐρυμνῷ χωρίῳ, ἀπὸ δὲ Ἱεροσολύμων ἀπέχοντι ὡσεὶ σταδίους πέντε, τοῦτο ἔθλιβεν.

\vs{6}Ὡς δὲ μετέλαβον οἱ περὶ τὸν Μακκαβαῖον πολιορκοῦντα αὐτὸν τὰ ὀχυρώματα, μετʼ ὀδύρμῶν καὶ δακρύων ἱκέτευον σὺν τοῖς ὄχλοις τὸν Κύριον, ἀγαθὸν ἄγγελον ἀποστεῖλαι πρὸς σωτηρίαν τῷ Ἰσραήλ.
\vs{7}Αὐτὸς δὲ πρῶτος ὁ Μακκαβαῖος ἀναλαβὼν τὰ ὅπλα προετρέψατο τοὺς ἄλλους, ἅμα αὐτῷ διακινδυνεύοντας, ἐπιβοηθεῖν τοῖς ἀδελφοῖς αὐτῶν· ὁμοῦ δὲ καὶ προθύμως ἐξώρμησαν.

\vs{8}Αὐτόθι δὲ καὶ πρὸς τοῖς Ἱεροσολύμοις ὄντων, ἐφάνη προηγούμενος αὐτῶν ἔφιππος ἐν λευκῇ ἐσθῆτι, πανοπλίαν χρυσῆν κραδαίνων.
\vs{9}Ὁμοῦ δὲ πάντες εὐλόγησαν τὸν ἐλεήμονα Θεὸν, καὶ ἐπεῤῥώσθησαν ταῖς ψυχαῖς, οὐ μόνον ἀνθρώπους ἀλλὰ καὶ θῆρας τοὺς ἀγριωτάτους, καὶ σιδηρᾶ τείχη τιτρώσκειν ὄντες ἕτοιμοι.
\vs{10}Προσῆγον ἐν διασκευῇ τὸν ἀπʼ οὐρανοῦ σύμμαχον ἔχοντες, ἐλεήσαντος αὐτοὺς τοῦ Κυρίου.
\vs{11}Λεοντηδὸν δὲ ἐντινάξαντες εἰς τοὺς πολεμίους, κατέστρωσαν αὐτῶν χιλίους πρὸς τοῖς μυρίοις, ἱππεῖς δὲ ἑξακοσίους πρὸς τοῖς χιλίοις· τοὺς δὲ πάντας ἠνάγκασαν φυγεῖν.
\vs{12}Οἱ πλείονες δὲ αὐτῶν τραυματίαι γυμνοὶ διεσώθησαν· καὶ αὐτὸς δὲ ὁ Λυσίας αἰσχρῶς φεύγων διεσώθη.

\vs{13}Οὐκ ἄνους δὲ ὑπάρχων, πρὸς ἑαυτὸν ἀντιβάλλων τὸ γεγονὸς περὶ ἐαυτὸν ἐλάσσωμα, καὶ συννοήσας ἀνικήτους εἶναι τοὺς Ἑβραίους, τοῦ πάντα δυναμένου Θεοῦ συμμαχοῦντος αὐτοῖς,
\vs{14}προσαποστείλας ἔπεισε συλλύσεσθαι ἐπὶ πᾶσι τοῖς δικαίοις· καὶ διότι καὶ τὸν βασιλέα πείσειν φίλον αὐτοῖς ἀναγκάζειν γενέσθαι.
\vs{15}Ἐπένευσε δὲ ὁ Μακκαβαῖος ἐπὶ πᾶσιν οἷς ὁ Λυσίας παρεκάλει τοῦ συμφέροντος φροντίζων· ὅσα γὰρ ὁ Μακκαβαῖος ἐπέδωκε τῷ Λυσίᾳ διὰ γραπτῶν περὶ τῶν Ἰουδαίων, συνεχώρησεν ὁ βασιλεύς.

\vs{16}Ἦσαν γὰρ αἱ γεγραμμέναι τοῖς Ἰουδαίοις ἐπιοτολαὶ παρὰ μὲν Λυσίου περιέχουσαι τὸν τρόπον τοῦτον· Λυσίας τῷ πλήθει τῶν Ἰουδαίων χαίρειν.
\vs{17}Ἰωάννης καὶ Ἀβεσσαλὼμ οἱ πεμφθέντες παρʼ ὑμῶν, ἐπιδόντες τὸν ὑπογεγραμμένον χρηματισμὸν, ἠξίουν περὶ τῶν διʼ αὐτοῦ σημαινομένων.
\vs{18}Ὅσα μὲν οὖν ἔδει καὶ τῷ βασιλεῖ προσενεχθῆναι διεσάφησα, ἃ δὲ ἦν ἐνδεχόμενα, συνεχώρησεν.
\vs{19}Ἐὰν μὲν οὖν συντηρήσητε τὴν εἰς τὰ πράγματα εὔνοιαν, καὶ εἰς τὸ λοιπὸν πειράσομαι παραίτιος ὑμῖν ἀγαθῶν γενέσθαι.
\vs{20}Ὑπὲρ δὲ τῶν κατὰ μέρος ἐντέταλμαι τούτοις τε καὶ τοῖς παρʼ ἐμοῦ διαλεχθῆναι ὑμῖν.
\vs{21}Ἔῤῥωσθε· ἔτους ἑκατοστοῦ τεσσαρακοστοῦ ὀγδόου, Διοσκορινθίου εἰκοστῇ τετάρτῃ.

\vs{22}Ἡ δὲ τοῦ βασιλέως ἐπιστολὴ περιεῖχεν οὕτως· βασιλεὺς Ἀντίοχος τῷ ἀδελφῷ Λυσίᾳ χαίρειν.
\vs{23}Τοῦ πατρὸς ἡμῶν εἰς θεοὺς μεταστάντος, βουλόμενοι τοὺς ἐκ τῆς βασιλείας ἀταράχους ὄντας γενέσθαι πρὸς τὴν τῶν ἰδίων ἐπιμέλειαν,
\vs{24}ἀκηκοότες τοὺς Ἰουδαίους μὴ συνευδοκοῦντας τῇ τοῦ πατρὸς ἐπὶ τὰ Ἑλληνικὰ μεταθέσει, ἀλλὰ τὴν ἑαυτῶν ἀγωγὴν αἱρετίζοντας, καὶ διὰ τοῦτο ἀξιοῦντας συγχωρηθῆναι αὐτοῖς τὰ νόμιμα αὐτῶν·
\vs{25}Αἱρούμενοι οὖν καὶ τοῦτο τὸ ἔθνος ἐκτὸς ταραχῆς εἶναι, κρίνομεν τό, τε ἱερὸν αὐτοῖς ἀποκατασταθῆναι, καὶ πολιτεύεσθαι κατὰ τὰ ἐπὶ τῶν προγόνων αὐτῶν ἔθη.
\vs{26}Εὖ οὖν ποιήσεις διαπεμψάμενος πρὸς αὐτοὺς καὶ δοὺς δεξιὰς, ὅπως εἰδότες τὴν ἡμετέραν προαίρεσιν, εὔθυμοί τε ὦσι, καὶ ἡδέως διαγίνωνται πρὸς τὴν τῶν ἰδίων ἀντίληψιν.

\vs{27}Πρὸς δὲ τὸ ἔθνος ἡ τοῦ βασιλέως ἐπιστολὴ τοιαύτη ἦν· βασιλεὺς Ἀντίοχος τῇ γερουσίᾳ τῶν Ἰουδαίων καὶ τοῖς ἄλλοις Ἰουδαίοις χαίρειν.
\vs{28}Εἰ ἔῤῥωσθε, εἴη ἂν ὡς βουλόμεθα· καὶ αὐτοὶ δὲ ὑγιαίνομεν.
\vs{29}Ἐνεφάνισεν ἡμῖν ὁ Μενέλαος βούλεσθαι κατελθόντας ὑμᾶς γίνεσθαι πρὸς τοῖς ἰδίοις.
\vs{30}Τοῖς οὖν καταπορευομένοις μέχρι τριακάδος Ξανθικοῦ ὑπάρξει δεξιὰ μετὰ τῆς ἀδείας,
\vs{31}χρῆσθαι τοὺς Ἰουδαίους τοῖς ἑαυτῶν δαπανήμασι καὶ νόμοις καθὰ καὶ τὸ πρότερον, καὶ οὐδεὶς αὐτῶν κατʼ οὐδένα τρόπον παρενοχληθήσεται περὶ τῶν ἠγνοημένων.
\vs{32}Πέπομφα δὲ καὶ τὸν Μενέλαον παρακαλέσοντα ὑμᾶς.
\vs{33}Ἔῤῥωσθε· ἔτους ἑκατοστοῦ τεσσαρακοστοῦ ὀγδόου, Ξανθικοῦ πέμπτῃ καὶ δεκάτῃ.

\vs{34}Ἔπεμψαν δὲ καὶ οἱ Ῥωμαιοῖ πρὸς αὐτοὺς ἐπιστολὴν ἔχουσαν οὕτως· Κόϊντος Μέμμιος, Τίτος Μάνλιος, πρεσβύται Ῥωμαίων, τῷ δήμῳ τῶν Ἰουδαίων χαίρειν.
\vs{35}Ὑπὲρ ὧν Λυσίας ὁ συγγενὴς τοῦ βασιλέως συνεχώρησεν ὑμῖν, καὶ ἡμεῖς συνευδοκοῦμεν.
\vs{36}Ἃ δὲ ἔκρινε προσανενεχθῆναι τῷ βασιλεῖ, πέμψατέ τινα παραχρῆμα ἐπισκεψάμενοι περὶ τούτων, ἵνα ἐκθῶμεν ὡς καθήκει ὑμῖν· ἡμεῖς γὰρ προσάγομεν πρὸς Ἀντιόχειαν.
\vs{37}Διὸ σπεύσατε, καὶ πέμψατέ τινας, ὅπως καὶ ἡμεῖς ἐπιγνῶμεν ὁποίας ἐστὲ γνώμης.
\vs{38}Ὑγιαίνετε· ἔτους ἑκατοστοῦ τεσσαρακοστοῦ ὀγδόου, Ξανθικοῦ πέμπτῃ καὶ δεκάτῃ.

\ch{12}
Γενομένων τῶν συνθηκῶν τούτων, ὁ μὲν Λυσίας ἀπῄει πρὸς τὸν βασιλέα, οἱ δὲ Ἰουδαῖοι περὶ τὴν γεωργίαν ἐγίνοντο.
\vs{2}Τῶν δὲ κατὰ τόπον στρατηγῶν Τιμόθεος καὶ Ἀπολλώνιος ὁ τοῦ Γενναίου, ἔτι δὲ Ἱερώνυμος καὶ Δημοφὼν, πρὸς δὲ τούτοις Νικάνωρ ὁ Κυπριάρχης, οὐκ εἴων αὐτοὺς εὐσταθεῖς, καὶ τὰ τῆς ἡσυχίας ἄγειν.

\vs{3}Ἰοππίται δὲ τηλικοῦτο συνετέλεσαν τὸ δυσσέβημα· παρακαλέσαντες τοὺς σὺν αὐτοῖς οἰκοῦντας Ἰουδαίους ἐμβῆναι εἰς τὰ παρασταθέντα ὑπʼ αὐτῶν σκάφη σὺν γυναιξὶ καὶ τέκνοις, ὡς μηδεμιᾶς ἐνεστώσης πρὸς αὐτοὺς δυσμενείας,
\vs{4}κατὰ δὲ τὸ κοινὸν τῆς πόλεως ψήφισμα, καὶ τούτων ἐπιδεξαμένων ὡς ἂν εἰρηνεύειν θελόντων, καὶ μηδὲν ὕποπτον ἐχόντων, ἐπαναχθέντας αὐτοὺς ἐβύθισαν, ὄντας οὐκ ἔλαττον τῶν διακοσίων.

\vs{5}Μεταλαβὼν δὲ Ἰούδας τὴν γεγονυῖαν εἰς τοὺς ὁμοεθνεῖς ὠμότητα, παραγγείλας τοῖς περὶ αὐτὸν ἀνδράσι,
\vs{6}καὶ ἐπικαλεσάμενος τὸν δίκαιον κριτὴν Θεὸν, παρεγένετο ἐπὶ τοὺς μιαιφόνους τῶν ἀδελφῶν· καὶ τὸν μὲν λιμένα νύκτωρ ἐνέπρησε, καὶ τὰ σκάφη κατέφλεξε, τοὺς δὲ ἐκεῖ συμφυγόντας ἐξεκέντησε.
\vs{7}Τοῦ δὲ χωρίου συγκλεισθέντος, ἀνέλυσεν, ὡς πάλιν ἥξων καὶ τὸ σύμπαν τῶν Ἰοππιτῶν ἐκριζῶσαι πολίτευμα.

\vs{8}Μεταλαβὼν δὲ καὶ τοὺς ἐν Ἰαμνείᾳ τὸν αὐτὸν ἐπιτελεῖν βουλομένους τρόπον τοῖς παροικοῦσιν Ἰουδαίοις,
\vs{9}καὶ τοῖς Ἰαμνίταις νυκτὸς ἐπιβαλὼν, ὑφῆψε τὸν λιμένα σὺν τῷ στόλῳ, ὥστε φαίνεσθαι τὰς αὐγὰς τοῦ φέγγους εἰς τὰ Ἱεροσόλυμα, σταδίων ὄντων διακοσίων τεσσαράκοντα.

\vs{10}Ἐκεῖθεν δὲ ἀποσπασθέντων σταδίους ἐννέα, ποιουμένων τὴν πορείαν ἐπὶ τὸν Τιμόθεον, προσέβαλον Ἄραβες αὐτῷ οὐκ ἐλάττους τῶν πεντακισχιλίων, ἱππεῖς δὲ πεντακόσιοι.
\vs{11}Γενομένης δὲ καρτερᾶς μάχης, καὶ τῶν περὶ τὸν Ἰούδαν διὰ τὴν παρὰ τοῦ Θεοῦ βοήθειαν εὐημερησάντων, ἐλαττωθέντες οἱ Νομάδες Ἄραβες ἠξίουν δοῦναι τὸν Ἰούδαν δεξιὰν αὐτοῖς, ὑπισχνούμενοι καὶ βοσκήματα δώσειν, καὶ ἐν τοῖς λοιποῖς ὠφελήσειν αὐτούς.

\vs{12}Ἰούδας δὲ ὑπολαβὼν ὡς ἀληθῶς ἐν πολλοῖς αὐτοὺς χρησίμους, ἐπεχώρησεν εἰρήνην ἄξειν πρὸς αὐτούς· καὶ λαβόντες δεξιὰς, εἰς τὰς σκηνὰς αὐτῶν ἐχωρίσθησαν.

\vs{13}Ἐπέβαλε δὲ καὶ ἐπί τινα πόλιν γεφυροῦν ὀχυρὰν καὶ τείχεσι περιπεφραγμένην, καὶ παμμιγέσιν ἔθνεσι κατοικουμένην, ὄνομα δὲ Κάσπιν.
\vs{14}Οἱ δʼ ἔνδον πεποιθότες τῇ τῶν τειχέων ἐρυμνότητι, τῇ τε τῶν βρωμάτων παραθέσει, ἀναγωγότερον ἐχρῶντο, τοῖς περὶ τὸν Ἰούδαν λοιδοροῦντες, καὶ προσέτι βλασφημοῦντες, καὶ λαλοῦντες ἃ μὴ θέμις.
\vs{15}Οἱ δὲ περὶ τὸν Ἰούδαν ἐπικαλεσάμενοι τὸν μέγαν τοῦ κόσμου δυνάστην, τὸν ἄτερ κριῶν καὶ μηχανῶν ὀργανικῶν κατακρημνίσαντα τὴν Ἱεριχὼ κατὰ τοὺς Ἰησοῦ χρόνους, ἐνέσεισαν θηριωδῶς τῷ τείχει.
\vs{16}Καταλαβόμενοί τε τὴν πόλιν τῇ τοῦ Θεοῦ θελήσει, ἀμυθήτους ἐποιήσαντο σφαγὰς, ὥστε τὴν παρακειμένην λίμνην τὸ πλάτος ἔχουσαν σταδίων δύο, κατάῤῥυτον αἵματι πεπληρωμένην φαίνεσθαι.

\vs{17}Ἐκεῖθεν δὲ ἀποσπάσαντες σταδίους ἑπτακοσίους πεντήκοντα διήνυσαν εἰς τὸν Χάρακα, πρὸς τοὺς λεγομένους Τουβιήνους Ἰουδαίους.
\vs{18}Καὶ Τιμόθεον μὲν ἐπὶ τῶν τόπων οὐ κατέλαβον, ἄπρακτόν τε ἀπὸ τῶν τόπων ἐκλελυκότα, καταλελοιπότα δὲ φρουρὰν ἔν τινι τόπῳ, καὶ μάλα ὀχυράν.
\vs{19}Δωσίθεος δὲ καὶ Σωσίπατρος τῶν περὶ τὸν Μακκαβαῖον ἡγεμόνων, ἐξοδεύσαντες ἀπώλεσαν τοὺς ὑπὸ Τιμοθέου καταλειφθέντας ἐν τῷ ὀχυρώματι πλείους τῶν μυρίων ἀνδρῶν.

\vs{20}Ὁ δὲ Μακκαβαῖος διατάξας τὴν ἑαυτοῦ στρατιὰν σπειρηδὸν, κατέστησεν αὐτοὺς ἐπὶ τῶν σπείρων, καὶ ἐπὶ τὸν Τιμόθεον ὥρμησεν ἔχοντα περὶ αὐτὸν μυριάδας δώδεκα πεζῶν, ἱππεῖς δὲ χιλίους πρὸς τοῖς πεντακοσίοις.

\vs{21}Τὴν δὲ ἔφοδον μεταλαβὼν Ἰούδα, ὁ Τιμόθεος προεξαπέστειλε τὰς γυναῖκας, καὶ τὰ τέκνα, καὶ τὴν ἄλλην ἀποσκευὴν εἰς τὸ λεγόμενον Καρνίον· ἦν γὰρ δυσπολιόρκητον καὶ δυσπρόσιτον τὸ χωρίον διὰ τὴν τῶν πάντων τῶν τόπων στενότητα.

\vs{22}Ἐπιφανείσης δὲ τῆς Ἰούδα σπείρας πρώτης, καὶ γενομένου δέους ἐπὶ τοὺς πολεμίους, φόβου τε ἐκ τῆς τοῦ πάντα ἐφορῶντος ἐπιφανείας γενομένου ἐπʼ αὐτοὺς, εἰς φυγὴν ὥρμησαν ἄλλος ἀλλαχῃ φερόμενος, ὥστε πολλάκις ὑπὸ τῶν ἰδίων βλάπτεσθαι, καὶ ταῖς τῶν ξιφῶν ἀκμαῖς ἀναπείρεσθαι.
\vs{23}Ἐποιεῖτο δὲ τὸν διωγμὸν εὐτονώτερον Ἰούδας, συγκεντῶν τοὺς ἀλιτηρίους, διέφθειρέ τε εἰς μυριάδας τρεῖς ἀνδρῶν.

\vs{24}Αὐτὸς δὲ ὁ Τιμόθεος ἐμπεσὼν τοῖς περὶ τὸν Δωσίθεον καὶ Σωσίπατρον, ἠξίου μετὰ πολλῆς γοητείας ἐξαφεῖναι σῶον αὐτόν· διὰ τὸ πλειόνων μὲν γονεῖς, ὧν δὲ ἀδελφοὺς ἔχειν, καὶ τούτους ἀλογηθῆναι συμβήσεται, εἰ ἀποθάνοι.
\vs{25}Πιστώσαντος δὲ αὐτοῦ διὰ πλειόνων τὸν ὁρισμὸν ἀποκαταστήσειν τούτους ἀπημάντους, ἀπέλυσαν αὐτὸν ἕνεκα τῆς τῶν ἀδελφῶν σωτηρίας.

\vs{26}Ἐξελθὼν δὲ ἐπὶ τὸ Καρνίον καὶ τὸ Ἀταργατεῖον, κατέσφαξε μυριάδας σωμάτων δύο καὶ πεντακισχιλίους.

\vs{27}Καὶ μετὰ τὴν τούτων τροπὴν καὶ ἀπώλειαν ἐπεστράτευσεν Ἰούδας καὶ ἐπὶ Ἐφρὼν, πόλιν ὀχυρὰν, ἐν ᾗ κατῴκει Λυσίας, καὶ πάμφυλα πλήθη· νεανίαι δὲ πρὸ τῶν τειχῶν καθεστῶτες ῥωμαλέοι ἀπεμάχοντο εὐρώστως, ἐνθάδε ὀργάνων καὶ βελῶν πολλαὶ παραθέσεις ὑπῆρχον.
\vs{28}Ἐπικαλεσάμενοι δὲ τὸν Δυνάστην τὸν μετὰ κράτους συντρίβοντα τὰς τῶν πολεμίων ἀλκὰς, ἔλαβον τὴν πόλιν ὑποχείριον, καὶ κατέστρωσαν τῶν ἔνδον εἰς μυριάδας δύο καὶ πεντακισχιλίους.

\vs{29}Ἀναζεύξαντες δὲ ἐκεῖθεν, ὥρμησαν ἐπὶ Σκυθῶν πόλιν, ἀπέχουσαν ἀπὸ Ἱεροσολύμων σταδίους ἑξακοσίους.
\vs{30}Ἀπομαρτυρησάντων δὲ τῶν ἐκεῖ κατοικούντων Ἰουδαίων, ἣν οἱ Σκυθοπολίται ἔσχον πρὸς αὐτοὺς εὔνοιαν, καὶ ἐν τοῖς τῆς ἀτυχίας καιροῖς ἥμερον ἀπάντησιν ἐποιοῦντο,
\vs{31}εὐχαριστήσαντες αὐτοῖς, καὶ προσπαρακαλέσαντες καὶ εἰς τὰ λοιπὰ πρὸς τὸ γένος εὐμενεῖς εἶναι, παρεγένοντο εἰς Ἱεροσόλυμα, τῆς τῶν ἑβδομάδων ἑορτῆς οὕσης ὑπογύου.

\vs{32}Μετὰ δὲ τὴν λεγομένην Πεντηκοστὴν, ὥρμησαν ἐπὶ Γοργίαν τὸν τῆς Ἰδουμαίας στρατηγόν.
\vs{33}Ἐξῆλθε δὲ μετὰ πεζῶν τρισχιλίων, ἱππέων δὲ τετρακοσίων.
\vs{34}Καὶ παραταξαμένων συνέβη πεσεῖν ὀλίγους τῶν Ἰουδαίων.
\vs{35}Δωσίθεος δέ τις τῶν τοῦ Βακήνορος, ἔφιππος ἀνὴρ καὶ καρτερὸς, εἴχετο τοῦ Γοργίου, καὶ λαβόμενος τῆς χλαμύδος, ἦγεν αὐτὸν εὐρώστως, καὶ βουλόμενος τὸν κατάρατον λαβεῖν ζωγρίαν, τῶν ἱππέων Θρακῶν τινὸς ἐπενεχθέντος αὐτῷ καὶ τὸν ὦμον καθελόντος, διέφυγεν ὁ Γοργίας εἰς Μαρισά.

\vs{36}Τῶν δὲ περὶ τὸν Ἔσδριν ἐπιπλεῖον μαχομένων, καὶ κατακόπων ὄντων, ἐπικαλεσάμενος ὁ Ἰούδας τὸν Κύριον σύμμαχον φανῆναι καὶ προοδηγὸν τοῦ πολέμου,
\vs{37}καταρξάμενος τῇ πατρίῳ φωνῇ τὴν μεθʼ ὕμνων κραυγὴν, ἀναβοήσας, καὶ ἐνσείσας ἀπροσδοκήτως τοῖς περὶ τὸν Γοργίαν, τροπὴν αὐτῶν ἐποιήσατο.
\vs{38}Ἰούδας δὲ ἀναλαβὼν τὸ στράτευμα, ἦγεν εἰς Ὀδολλὰμ πόλιν· τῆς δὲ ἑβδομάδος ἐπιβαλλούσης, κατὰ τὸν ἐθισμὸν ἁγνισθέντες αὐτόθι τὸ σάββατον διήγαγον.

\vs{39}Τῇ δὲ ἐχομένῃ ἦλθον οἱ περὶ τὸν Ἰούδαν καθʼ ὃν τρόπον τὸ τῆς χρείας ἐγεγόνει, τὰ τῶν προπεπτωκότων σώματα ἀνακομίσασθαι, καὶ μετὰ τῶν συγγενῶν ἀποκαταστῆσαι εἰς τοὺς πατρῴους τάφους.
\vs{40}Εὗρον δὲ ἑκάστου τῶν τεθνηκότων ὑπὸ τοὺς χιτῶνας ἱερώματα τῶν ἀπὸ Ἰαμνείας εἰδώλων, ἀφʼ ὧν ὁ νόμος ἀπείργει τοὺς Ἰουδαίους· τοῖς δὲ πᾶσι σαφὲς ἐγένετο διὰ τήνδε τὴν αἰτίαν τούσδε πεπτωκέναι.
\vs{41}Πάντες οὖν εὐλογήσαντες τοῦ δικαιοκρίτου Κυρίου τοῦ τὰ κεκρυμμένα φανερὰ ποιοῦντος,
\vs{42}εἰς ἱκετείαν ἐτράπησαν, ἀξιώσαντες τὸ γεγονὸς ἁμάρτημα τελείως ἐξαλειφθῆναι· ὁ δὲ γενναῖος Ἰούδας παρεκάλεσε τὸ πλῆθος συντηρεῖν ἑαυτοὺς ἀναμαρτήτους εἶναι, ὑπʼ ὄψιν ἑωρακότας τὰ γεγονότα, διὰ τὴν τῶν προπεπτωκότων ἁμαρτίαν.

\vs{43}Ποιησάμενός τε κατʼ ἀνδραλογίαν κατασκευάσματα εἰς ἀργυρίου δραχμὰς δισχιλίας, ἀπέστειλεν εἰς Ἱεροσόλυμα προσαγαγεῖν περὶ ἁμαρτίας θυσίαν, πάνυ καλῶς καὶ ἀστείως πράττων, ὑπὲρ ἀναστάσεως διαλογιζόμενος·
\vs{44}εἰ γὰρ μὴ τοὺς προπεπτωκότας ἀναστῆναι προσεδόκα, περισσὸν ἂν ἦν καὶ ληρῶδες ὑπὲρ νεκρῶν προσεύχεσθαι·
\vs{45}εἶτʼ ἐμβλέπων τοῖς μετʼ εὐσεβείας κοιμωμένοις κάλλιστον ἀποκείμενον χαριστήριον· ὁσία καὶ εὐσεβὴς ἡ ἐπίνοια· ὅθεν περὶ τῶν τεθνηκότων τὸν ἐξιλασμὸν ἐποιήσατο, τῆς ἁμαρτίας ἀπολυθῆναι.

\ch{13}
Τῷ δὲ ἐννάτῳ καὶ τεσσαρακοστῷ καὶ ἑκατοστῷ ἔτει προσέπεσε τοῖς περὶ τὸν Ἰούδαν, Ἀντίοχον τὸν Εὐπάτορα παραγενέσθαι σὺν πλήθεσιν ἐπὶ τὴν Ἰουδαίαν,
\vs{2}καὶ σὺν αὐτῷ Λυσίαν τὸν ἐπίτροπον καὶ ἐπὶ τῶν πραγμάτων, ἕκαστον ἔχοντα δύναμιν Ἑλληνικὴν πεζῶν μυριάδας ἕνδεκα, καὶ ἱππεῖς πεντακισχιλίους τριακοσίους, καὶ ἐλέφαντας εἰκοσιδύο, ἅρματα δὲ δρεπανηφόρα τριακόσια.

\vs{3}Καὶ Μενέλαος δὲ συνέμιξεν αὐτοῖς, καὶ παρεκάλει μετὰ πολλῆς εἰρωνείας τὸν Ἀντίοχον, οὐκ ἐπὶ σωτηρίᾳ τῆς πατρίδος, οἰόμενος δὲ ἐπὶ τῆς ἀρχῆς κατασταθήσεσθαι.
\vs{4}Ὁ δὲ βασιλεὺς τῶν βασιλέων ἐξήγειρε τὸν θυμὸν τοῦ Ἀντιόχου ἐπὶ τὸν ἀλιτήριον, καὶ Λυσίου ὑποδείξαντος τοῦτον αἴτιον εἶναι πάντων τῶν κακῶν, προσέταξεν, ὡς ἔθος ἐστὶν ἐν τῷ τόπῳ, προσαπολέσαι ἀγαγόντας αὐτὸν εἰς Βέροιαν.

\vs{5}Ἔστι δὲ ἐν τῷ τόπῳ πύργος πεντήκοντα πηχῶν πλήρης σποδοῦ οὗτος δὲ ὄργανον εἶχε περιφερὲς πάντοθεν ἀπόκρημνον εἰς τὴν σποδόν.
\vs{6}Ἐνταῦθα τὸν ἱεροσυλίας ἔνοχον ὄντα, ἢ καί τινων ἄλλων κακῶν ὑπεροχὴν πεποιημένον, ἅπαντες προσωθοῦσιν εἰς ὄλεθρον.
\vs{7}Τοιούτῳ μόρῳ τὸν παράνομον συνέβη θανεῖν, μηδὲ τῆς γῆς τυχόντα Μενέλαον·
\vs{8}πάνυ δικαίως. Ἐπεὶ γὰρ συνετελέσατο πολλὰ περὶ τὸν βωμὸν ἁμαρτήματα, οὗ τὸ πῦρ ἁγνὸν ἦν καὶ ἡ σποδὸς, ἐν σποδῷ τὸν θάνατον ἐκομισατο.

\vs{9}Τοῖς δὲ φρονήμασιν ὁ βασιλεὺς βεβαρβαρωμένος ἤρχετο, τὰ χείριστα τῶν ἐπὶ τοῦ πατρὸς αὐτοῦ γεγονότων ἐνδειξόμενος τοῖς Ἰουδαίοις.
\vs{10}Μεταλαβὼν δὲ Ἰούδας ταῦτα, παρήγγειλε τῷ πλήθει διʼ ἡμέρας καὶ νυκτὸς ἐπικαλεῖσθαι τὸν Κύριον, εἴποτε ἄλλοτε, καὶ νῦν ἐπιβοηθεῖν τοῖς τοῦ νόμου καὶ πατρίδος καὶ ἱεροῦ ἁγίου στερεῖσθαι μέλλουσι,
\vs{11}καὶ τὸν ἄρτι βραχέως ἀνεψυχότα λαὸν μὴ ἐᾶσαι τοῖς δυσφήμοις ἔθνεσιν ὑποχειρίους γενέσθαι.

\vs{12}Πάντων δὲ τὸ αὐτὸ ποιησάντων ὁμοῦ καὶ καταξιωσάντων τὸν ἐλεήμονα Κύριον μετὰ κλαυθμοῦ καὶ νηστειῶν καὶ προπτώσεως ἐφʼ ἡμέρας τρεῖς ἀδιαλείπτως, παρακαλέσας αὐτοὺς ὁ Ἰούδας ἐκέλευσε παραγίνεσθαι.

\vs{13}Καθʼ ἑαυτὸν δὲ σὺν τοῖς πρεσβυτέροις γενόμενος, ἐβουλεύσατο πρὶν εἰσβαλεῖν τοῦ βασιλέως τὸ στράτευμα εἰς τὴν Ἰουδαίαν, καὶ γενέσθαι τῆς πόλεως ἐγκρατεῖς, ἐξελθόντας κρῖναι τὰ πράγματα τῇ τοῦ Κυρίου βοηθείᾳ.

\vs{14}Δοὺς δὲ τὴν ἐπιτροπὴν τῷ κτίστῃ τοῦ κόσμου, παρακαλέσας τοὺς σὺν αὐτῷ γενναίως ἀγωνίσασθαι μέχρι θανάτου περὶ νόμων, περὶ ἱεροῦ, πόλεως, πατρίδος, πολιτείας, ἐποιήσατο περὶ Μωδεῒν τὴν στρατοπεδείαν.
\vs{15}Δοὺς δὲ τοῖς περὶ αὐτὸν σύνθημα Θεοῦ νίκης, μετὰ νεανίσκων ἀρίστων κεκριμένων ἐπιβαλὼν νύκτωρ ἐπὶ τὴν βασιλικὴν αὐλὴν, ἐν τῇ παρεμβολῇ ἀνεῖλεν εἰς ἄνδρας τετρακισχιλίους, καὶ τὸν πρωτεύοντα τῶν ἐλεφάντων σὺν τῷ κατʼ οἰκίαν ὄχλῳ συνέθηκε,
\vs{16}καὶ τὸ τέλος τὴν παρεμβολὴν δέους καὶ ταραχῆς ἐπλήρωσαν, καὶ ἐξέλυσαν εὑημεροῦντες.
\vs{17}Ὑποφαινούσης δὲ ἤδη τῆς ἡμέρας τοῦτʼ ἐγεγόνει, διὰ τὴν ἐπαρήγουσαν αὐτῷ τοῦ Κυρίου σκέπην.

\vs{18}Ὁ δὲ βασιλεὺς εἰληφὼς γεῦσιν τῆς τῶν Ἰουδαίων εὐτολμίας, κατεπείρασε διὰ μεθόδων τοὺς τόπους.
\vs{19}Καὶ ἐπὶ Βαιθσούρᾳ φρούριον ὀχυρὸν τῶν Ἰουδαίων προσῆγεν· καὶ ἐτροποῦτο, προσέκρουεν, ἠλαττονοῦτο.
\vs{20}Τοῖς δὲ ἔνδον Ἰούδας τὰ δέοντα εἰσέπεμψε.

\vs{21}Προσήγγειλε δὲ τὰ μυστήρια τοῖς πολεμίοις Ῥόδοκος ἐκ τῆς Ἰουδαϊκῆς τάξεως· ἀνεζητήθη δὲ, καὶ κατελήφθη, καὶ κατεκλείσθη.

\vs{22}Ἐδευτερολόγησεν ὁ βασιλεὺς τοῖς ἐν Βαιθσούρᾳ δεξιὰν ἔδωκεν, ἔλαβεν, ἀπῄει, προσέβαλε τοῖς περὶ τὸν Ἰούδαν, ἥττων ἐγένετο,
\vs{23}μετέλαβεν ἀπονενοῆσθαι τὸν Φίλιππον ἐν Ἀντιοχείᾳ τὸν ἀπολελειμμένον ἐπὶ τῶν πραγμάτων, συνεχύθη· τοὺς Ἰουδαίους παρεκάλεσεν, ὑπετάγη, καὶ ὤμοσεν ἐπὶ πᾶσι τοῖς δικαίοις· συνελύθη καὶ θυσίαν προσήγαγεν, ἐτίμησε τὸν νεὼν, καὶ τὸν τόπον ἐφιλανθρώπησε,
\vs{24}καὶ τὸν Μακκαβαῖον ἀπεδέξατο· κατέλιπε στρατηγὸν ἀπὸ Πτολεμαΐδος ἕως τῶν Γεῤῥηνῶν ἡγεμονίδην,
\vs{25}ἦλθεν εἰς Πτολεμαΐδα· ἐδυσφόρουν περὶ τῶν συνθηκῶν οἱ Πτολεμαεῖς, ἐδείναζον γὰρ ὑπὲρ ὧν ἠθέλησαν ἀθετεῖν τὰς διαστάλσεις.

\vs{26}Προσῆλθεν ἐπὶ τὸ βῆμα Λυσίας, ἀπελογήσατο ἐνδεχομένως, συνέπεισε, κατεπρᾴυνεν, εὐμενεῖς ἐποίησεν, ἀνέζευξεν εἰς Ἀντιόχειαν· οὕτω τὰ τοῦ βασιλέως τῆς ἐφόδου καὶ τῆς ἀναζυγῆς ἐχώρησε.

\ch{14}
Μετὰ δὲ τριετῆ χρόνον προσέπεσε τοῖς περὶ τὸν Ἰούδαν, Δημήτριον τὸν τοῦ Σελεύκου διὰ τοῦ κατὰ Τρίπολιν λιμένος εἰσπλεύσαντα μετὰ πλήθους ἰσχυροῦ καὶ στόλου,
\vs{2}κεκρατηκέναι τῆς χώρας, ἐπανελόμενον Ἀντίοχον καὶ τὸν τούτου ἐπίτροπον Λυσίαν.

\vs{3}Ἄλκιμος δέ τις προγενόμενος ἀρχιερεὺς, ἑκουσίως δὲ μεμολυμμένος ἐν τοῖς τῆς ἐπιμιξίας χρόνοις, συννοήσας ὅτι καθʼ ὁντιναοῦν τρόπον οὐκ ἔστιν αὐτῷ σωτηρία, οὐδὲ πρὸς ἅγιον θυσιαστήριον ἔτι πρόσοδος,
\vs{4}ἧκε πρὸς τὸν βασιλέα Δημήτριον πρώτῳ καὶ πεντηκοστῷ καὶ ἑκατοστῷ ἔτει, προσάγων αὐτῷ στέφανον χρυσοῦν καὶ φοίνικα, πρὸς δὲ τούτοις τῶν νομιζομένων θαλλῶν τοῦ ἱεροῦ· καὶ τὴν ἡμέραν ἐκείνην ἡσυχίαν ἔσχε.

\vs{5}Καιρὸν δὲ λαβὼν τῆς ἰδίας ἀνοίας συνεργὸν, προσκληθεὶς εἰς συνέδριον ὑπὸ τοῦ Δημητρίου, καὶ ἐπερωτηθεὶς ἐν τίνι διαθέσει καὶ βουλῇ καθεστήκασιν οἱ Ἰουδαῖοι, πρὸς ταῦτα ἔφη,
\vs{6}οἱ λεγόμενοι τῶν Ἰουδαίων Ἀσιδαῖοι, ὧν ἀφηγεῖται Ἰούδας ὁ Μακκαβαῖος, πολεμοτροφοῦσι καὶ στασιάζουσιν, οὐκ ἐῶντες τὴν βασιλείαν εὐσταθείας τυχεῖν.

\vs{7}Ὅθεν ἀφελόμενος τὴν προγονικὴν δόξαν, λέγω δὴ τὴν ἀρχιερωσύνην, δεῦρο νῦν ἐλήλυθα.
\vs{8}Πρῶτον μὲν ὑπὲρ τῶν ἀνηκόντων τῷ βασιλεῖ γνησίως φρονῶν, δεύτερον δὲ καὶ τῶν ἰδίων πολιτῶν στοχαζόμενος· τῇ μὲν γὰρ τῶν προειρημένων ἀλογιστίᾳ τὸ σύμπαν ἡμῶν γένος οὐ μικρῶς ἀκληρεῖ.

\vs{9}Ἕκαστα δὲ τούτων ἐπεγνωκὼς σὺ βασιλεῦ, καὶ τῆς χώρας καὶ τοῦ περιϊσταμένου γένους ἡμῶν προνοήθητι, καθʼ ἣν ἔχεις πρὸς ἅπαντας εὐαπάντητον φιλανθρωπίαν.
\vs{10}Ἄχρι γὰρ Ἰούδας περίεστιν, ἀδύνατον εἰρήνης τυχεῖν τὰ πράγματα.
\vs{11}Τοιούτων δὲ ῥηθέντων ὑπὸ τούτου, θᾶττον οἱ λοιποὶ φίλοι δυσμενῶς ἔχοντες τὰ πρὸς τὸν Ἰούδαν προσεπύρωσαν τὸν Δημήτριον.

\vs{12}Προσκαλεσάμενος δὲ εὐθέως Νικάνορα τὸν γενόμενον ἐλεφαντάρχην, καὶ στρατηγὸν ἀναδείξας τῆς Ἰουδαίας, ἐξαπέστειλε,
\vs{13}δοὺς ἐντολὰς, αὐτὸν μὲν τὸν Ἰούδαν ἐπανελέσθαι, τοὺς δὲ σὺν αὐτῷ σκορπίσαι, καταστῆσαι δὲ Ἄλκιμον ἀρχιερέα τοῦ μεγίστου ἱεροῦ.
\vs{14}Τὰ δὲ ἐκ τῆς Ἰουδαίας πεφυγαδευκότα τὸν Ἰούδαν ἔθνη συνέμισγον ἀγεληδὸν τῷ Νικάνορι, τὰς τῶν Ἰουδαίων ἀτυχίας καὶ συμφορὰς, ἰδίας εὐημερίας δοκοῦντες ἔσεσθαι.

\vs{15}Ἀκούσαντες δὲ τὴν τοῦ Νικάνορος ἔφοδον καὶ τὴν ἐπίθεσιν τῶν ἐθνῶν, καταπασάμενοι γῆν ἐλιτάνευον τὸν ἄχρι αἰῶνος συστήσαντα τὸν ἑαυτοῦ λαὸν, ἀεὶ δὲ μετʼ ἐπιφανείας ἀντιλαμβανόμενον τῆς ἑαυτοῦ μερίδος.
\vs{16}Προστάξαντος δὲ τοῦ ἡγουμένου, ἐκεῖθεν εὐθέως ἀνέζευξαν, καὶ συμμίσγουσιν αὐτοῖς ἐπὶ κώμην Δεσσαού.

\vs{17}Σίμων δὲ ὁ ἀδελφὸς Ἰούδα συμβεβληκῶς ἦν τῷ Νικάνορι, βραχέως δὲ διὰ τὴν αἰφνίδιον τῶν ἀντιπάλων ἀφασίαν ἐπταικώς.
\vs{18}Ὅμως δὲ ἀκούων ὁ Νικάνωρ ἣν εἶχον οἱ περὶ τὸν Ἰούδαν ἀνδραγαθίαν, καὶ ἐν τοῖς ὑπὲρ τῆς πατρίδος ἀγῶσιν εὐψυχίαν, ἐπευλαβεῖτο τὴν κρίσιν διʼ αἱμάτων ποιήσασθαι·
\vs{19}Διόπερ ἔπεμψε Ποσιδώνιον καὶ Θεόδοτον καὶ Ματταθίαν, δοῦναι καὶ λαβεῖν δεξιάς.

\vs{20}Πλείονος δὲ γενομένης περὶ τούτων ἐπισκέψεως, καὶ τοῦ ἡγεμόνος τοῖς πλήθεσιν ἀνακοινωσαμένου, καὶ φανείσης ὁμοψήφου γνώμης, ἐπένευσαν ταῖς συνθήκαις.
\vs{21}Ἐτάξαντο δὲ ἡμέραν ἐν ᾗ κατʼ ἰδίαν ἥξουσιν εἰς τὸ αὐτό· καὶ προῆλθε, καὶ παρʼ ἑκάστου διαφόρους ἔθεσαν δίφρους.
\vs{22}Διέταξεν Ἰούδας ἐνόπλους ἑτοίμους ἐν τοῖς ἐπικαίροις τόποις, μήποτε ἐκ τῶν πολεμίων αἰφνιδίως κακουργία γένηται· τὴν ἁρμόζουσαν ἐποιήσαντο κοινολογίαν.

\vs{23}Διέτριβεν δὲ ὁ Νικάνωρ ἐν Ἱεροσολύμοις, καὶ ἔπραττεν οὐθὲν ἄτοπον· τοὺς δὲ συναχθέντας ἀγελαίους ὄχλους ἀπέλυσε.
\vs{24}Καὶ εἶχε τὸν Ἰούδαν διαπαντὸς ἐν προσώπῳ, ψυχικῶς τῷ ἀνδρὶ προσεκέκλιτο.
\vs{25}Παρεκάλεσεν αὐτὸν γῆμαι καὶ παιδοποιήσασθαι· ἐγάμησεν, εὐστάθησεν, ἐκοινώνησε βίου.

\vs{26}Ὁ δὲ Ἄλκιμος συνιδὼν τὴν πρὸς ἀλλήλους εὔνοιαν καὶ τὰς γενομένας συνθήκας, ἀναλαβὼν, ἧκε πρὸς τὸν Δημήτριον, καὶ ἔλεγε τὸν Νικάνορα ἀλλότρια φρονεῖν τῶν πραγμάτων· τὸν γὰρ ἐπίβουλον τῆς βασιλείας Ἰούδαν διάδοχον ἀναδέδειχεν ἑαυτοῦ.
\vs{27}Ὁ δὲ βασιλεὺς ἔκθυμος γενόμενος, καὶ ταῖς τοῦ παμπονήρου ἐρεθισθεὶς διαβολαῖς, ἔγραψε Νικάνορι φάσκων, ὑπὲρ μὲν τῶν συνθηκῶν βαρέως φέρειν, κελεύων δὲ τὸν Μακκαβαῖον δέσμιον ἐξαποστέλλειν ταχέως εἰς Ἀντιόχειαν.

\vs{28}Προσπεσόντων δὲ τούτων τῷ Νικάνορι, συνεκέχυτο καὶ δυσφόρως ἔφερεν, εἰ τὰ διεσταλμένα ἀθετήσει μηδὲν τʼ ἀνδρὸς ἠδικηκότος.
\vs{29}Ἐπεὶ δὲ τῷ βασιλεῖ ἀντιπράττειν οὐκ ἦν, εὔκαιρον ἐτήρει στρατηγήματι τοῦτʼ ἐπιτελέσαι.

\vs{30}Ὁ δὲ Μακκαβαῖος αὐστηρότερον διεξάγοντα συνιδὼν τὸν Νικάνορα πρὸς αὐτὸν, καὶ τὴν εἰθισμένην ἀπάντησιν ἀγριωτέραν ἐσχηκότα, νοήσας οὐκ ἀπὸ τοῦ βελτίστου τὴν αὐστηρίαν εἶναι, συστρέψας οὐκ ὀλίγους τῶν περὶ ἑαυτὸν, συνεκρύπτετο τὸν Νικάνορα.
\vs{31}Συγγνοὺς δὲ ὁ ἕτερος ὅτι γενναίως ὑπὸ τοῦ ἀνδρὸς ἐστρατήγηται, παραγενόμενος ἐπὶ τὸ μέγιστον καὶ ἅγιον ἱερὸν, τῶν ἱερέων τὰς καθηκούσας θυσίας προσαγόντων, ἐκέλευσε παραδιδόναι τὸν ἄνδρα.
\vs{32}Τῶν δὲ μεθʼ ὅρκων φασκόντων μὴ γινώσκειν ποῦ ποτʼ ἐστὶν ὁ ζητούμενος,
\vs{33}προτείνας τὴν δεξιὰν εἰς τὸν νεὼν, ταῦτα ὤμοσεν, ἐὰν μὴ δέσμιόν μοι τὸν Ἰούδαν παραδῶτε, τόνδε τοῦ Θεοῦ σηκὸν εἰς πεδίον ποιήσω, καὶ τὸ θυσιαστήριον κατασκάφω, καὶ ἱερὸν ἐνταῦθα τῷ Διονύσῳ ἐπιφανὲς ἀναστήσω.

\vs{34}Τοσαῦτα δὲ εἰπὼν ἀπῆλθεν· οἱ δὲ ἱερεῖς προτείναντες τὰς χεῖρας εἰς τὸν οὐρανὸν, ἐπεκαλοῦντο τὸν διαπαντὸς ὑπέρμαχον τοῦ ἔθνους ἡμῶν, ταῦτα λέγοντες,
\vs{35}σὺ, Κύριε, τῶν ὅλων ἀπροσδεὴς ὑπάρχων, εὐδόκησας ναὸν τῆς σῆς κατασκηνώσεως ἐν ἡμῖν γενέσθαι.
\vs{36}Καὶ νῦν, ἅγιε παντὸς ἁγιασμοῦ Κύριε διατήρησον εἰς αἰῶνα ἀμίαντον τόνδε τὸν προσφάτως κεκαθαρισμένον οἶκον.

\vs{37}Ῥαζὶς δέ τις τῶν ἀπὸ Ἱεροσολύμων πρεσβυτέρων, ἐμηνύθη τῷ Νικάνορι, ἀνὴρ φιλοπολίτης καὶ σφόδρα καλῶς ἀκούων, καὶ κατὰ τὴν εὔνοιαν πατὴρ τῶν Ἰουδαίων προσαγορευόμενος.
\vs{38}Ἦν γὰρ ἐν τοῖς ἔμπροσθεν χρόνοις τῆς ἀμιξίας κρίσιν εἰσενηνεγμένος Ἰουδαϊσμοῦ, καὶ σῶμα καὶ ψυχὴν ὑπὲρ τοῦ Ἰουδαϊσμοῦ παραβεβλημένος μετὰ πάσης ἐκτενίας.

\vs{39}Βουλόμενος δὲ Νικάνωρ πρόδηλον ποιῆσαι ἣν εἶχε πρὸς τοὺς Ἰουδαίους δυσμένειαν, ἀπέστειλε στρατιώτας ὑπὲρ τοὺς πεντακοσίους συλλαβεῖν αὐτόν.
\vs{40}Ἔδοξε γὰρ, ἐκεῖνον συλλαβὼν, τούτοις ἐργάσασθαι συμφοράν,
\vs{41}Τῶν δὲ πληθῶν μελλόντων τὸν πύργον καταλαβέσθαι, καὶ τὴν αὐλαίαν θύραν βιαζομένων, καὶ κελευόντων πῦρ προσάγειν καὶ τὰς θύρας ὑφάπτειν, περικατάληπτος γενόμενος ὑπέθηκεν ἑαυτῷ ξίφος,
\vs{42}εὐγενῶς θέλων ἀποθανεῖν, ἤπερ τοῖς ἀλιτηρίοις ὑποχείριος γενέσθαι, καὶ τῆς ἰδίας εὐγενείας ἀναξίως ὑβρισθῆναι.
\vs{43}Τῇ δὲ πληγῇ μὴ κατευθικτήσας διὰ τὴν τοῦ ἀγῶνος σπουδὴν, καὶ τῶν ὄχλων εἴσω τῶν θυρωμάτων εἰσβαλόντων, ἀναδραμὼν γενναίως ἐπὶ τὸ τεῖχος, κατεκρήμνισεν ἑαυτὸν ἀνδρείως εἰς τοὺς ὄχλους.
\vs{44}Τῶν δὲ ταχέως ἀναποδισάντων, γενομένου διαστήματος ἦλθε κατὰ μέσον τὸν κενεῶνα.

\vs{45}Ἔτι δὲ ἔμπνους ὑπάρχων καὶ πεπυρωμένος τοῖς θυμοῖς, ἐξαναστὰς φερομένων κρουνηδὸν τῶν αἱμάτων, καὶ δυσχερῶν ὄντων τῶν τραυμάτων, δρόμῳ τοὺς ὄχλους διελθὼν, καὶ στὰς ἐπί τινος πέτρας ἀποῤῥωγάδος,
\vs{46}παντελῶς ἔξαιμος ἤδη γενόμενος, προβαλὼν τὰ ἔντερα, καὶ λαβὼν ἑκατέραις ταῖς χερσὶν, ἐνέσεισε τοῖς ὄχλοις· καὶ ἐπικαλεσάμενος τὸν δεσπόζοντα τῆς ζωῆς καὶ τοῦ πνεύματος, ταῦτα αὐτῷ πάλιν ἀποδοῦναι, τόνδε τὸν τρόπον μετήλλαξεν.

\ch{15}
Ὁ δὲ Νικάνωρ μεταλαβὼν τοὺς περὶ τὸν Ἰούδαν ὄντας ἐν τοῖς κατὰ Σαμάρειαν τόποις, ἐβουλεύσατο τῇ τῆς καταπαύσεως ἡμέρᾳ μετὰ πάσης ἀσφαλείας αὐτοῖς ἐπιβαλεῖν.

\vs{2}Τῶν δὲ κατʼ ἀνάγκην συνεπομένων αὐτῷ Ἰουδαίων, λεγόντων, μηδαμῶς οὕτως ἀγρίως καὶ βαρβάρως ἀπολέσῃς, δόξαν δὲ ἀπομέρισον τῇ προτετιμημένῃ ὑπὸ τοῦ πάντα ἐθορῶντος μεθʼ ἁγιότητος ἡμέρᾳ.

\vs{3}Ὁ δὲ τρισαλιτήριος ἐπηρώτησεν, εἰ ἔστιν ἐν οὐρανῷ δυνάστης ὁ προστεταχὼς ἄγειν τὴν τῶν σαββάτων ἡμέραν;
\vs{4}Τῶν δὲ ἀποφῃναμένων, ἔστιν ὁ Κύριος ζῶν αὐτὸς ἐν οὐρανῷ δυνάστης, ὁ κελεύσας ἀσκεῖν τὴν ἑβδομάδα.
\vs{5}Ὁ δὲ ἕτερος, κᾀγώ φησι, δυνάστης ἐπὶ τῆς γῆς ὁ προστάσσων αἴρειν ὅπλα, καὶ τὰς βασιλικὰς χρείας ἐπιτελεῖν· ὅμως οὐ κατέσχεν ἐπιτελέσαι τό σχέτλιον αὐτοῦ βούλημα.
\vs{6}Καὶ ὁ μὲν Νικάνωρ μετὰ πάσης ἀλαζονείας ὑψαυχενῶν, διεγνώκει κοινὸν τῶν περὶ τὸν Ἰούδαν συστήσασθαι τρόπαιον.

\vs{7}Ὁ δὲ Μακκαβαῖος ἦν ἀδιαλείπτως πεποιθὼς μετὰ πάσης ἐλπίδος ἀντιλήψεως τεύξασθαι παρὰ τοῦ Κυρίου.
\vs{8}Καὶ παρεκάλει τοὺς σὺν αὐτῷ μὴ δειλιᾷν τὴν τῶν ἔθνῶν ἔφοδον, ἔχοντας δὲ κατὰ νοῦν τὰ προγεγονότα αὐτοῖς ἁπʼ οὐρανοῦ βοηθήματα, καὶ τανῦν προσδοκᾷν τὴν παρὰ τοῦ παντοκράτορος ἐσομένην αὐτοῖς νίκην καὶ βοήθειαν.
\vs{9}Καὶ παραμυθούμενος αὐτοὺς ἐκ τοῦ νόμου καὶ τῶν προφητῶν, προσυπομνήσας δὲ αὐτοὺς καὶ τοὺς ἀγῶνας οὓς ἦσαν ἐκτετελεκότες, προθυμοτέρους αὐτοὺς κατέστησε.

\vs{10}Καὶ τοῖς θυμοῖς διεγείρας αὐτοὺς, παρήγγειλεν, ἅμα παρεπιδεικνὺς τὴν τῶν ἐθνῶν ἀθεσίαν καὶ τὴν τῶν ὅρκων παράβασιν.
\vs{11}Ἕκαστον δὲ αὐτῶν καθοπλίσας, οὐ τὴν ἀσπίδων καὶ λογχῶν ἀσφάλειαν, ὡς τὴν ἐν τοῖς ἀγαθοῖς λόγοις παράκλησιν, καὶ προσεξηγησάμενος ὄνειρον ἀξιόπιστον ὕπαρ τι πάντας εὔφρανεν.

\vs{12}Ἦν δὲ ἡ τούτου θεωρία τοιάδε· Ὀνίαν τὸν γενόμενον ἀρχιερέα, ἄνδρα καλὸν καὶ ἀγαθὸν, αἰδήμονα μὲν τὴν ἀπάντησιν, πρᾷον δὲ τὸν τρόπον, καὶ λαλιὰν προϊέμενον πρεπόντως, καὶ ἐκ παιδὸς ἐκμεμελετηκότα πάντα τὰ τῆς ἀρετῆς οἰκεῖα, τοῦτον τὰς χεῖρας προτείναντα κατεύχεσθαι τῷ παντὶ τῶν Ἰουδαίων συστήματι.
\vs{13}Εἶθʼ οὕτως ἐπιφανῆναι ἄνδρα πολιᾷ καὶ δόξῃ διαφέροντα, θαυμαστὴν δέ τινα καὶ μεγαλοπρεπεστάτην εἶναι τὴν περὶ αὐτὸν ὑπεροχήν.
\vs{14}Ἀποκριθέντα δὲ τὸν Ὀνιαν εἰπεῖν, ὁ φιλάδελφος οὗτός ἐστιν ὁ πολλὰ προσευχόμενος περὶ τοῦ λαοῦ καὶ τῆς ἁγίας πόλεως, Ἱερεμίας ὁ τοῦ Θεοῦ προφήτης.
\vs{15}Προτείναντα δὲ τὸν Ἱερεμίαν τὴν δεξιὰν παραδοῦναι τῷ Ἰούδᾳ ῥομφαίαν χρυσῆν, διδόντα δὲ προσφωνῆσαι τάδε,
\vs{16}λάβε τὴν ἁγίαν ῥομφαίαν δῶρον παρὰ τοῦ Θεοῦ, διʼ ἧς θραύσεις τοὺς ὑπεναντίους.

\vs{17}Παρακληθέντες δὲ τοῖς Ἰούδα λόγοις πάνυ καλοῖς καὶ δυναμένοις ἐπʼ ἀρετὴν παρορμῆσαι, καὶ ψυχὰς νέων ἐπανορθῶσαι, διέγνωσαν μὴ στρατοπεδεύεσθαι, γενναίως δὲ ἐμφέρεσθαι, καὶ μετὰ πάσης εὐανδρίας ἐμπλακέντες κρῖναι τὰ πράγματα, διὰ τὸ καὶ τὴν πόλιν, καὶ τὰ ἅγια, καὶ τὸ ἱερὸν κινδυνεύειν.
\vs{18}Ἦν γὰρ ὁ περὶ γυναικῶν καὶ τέκνων, ἔτι δὲ ἀδελφῶν καὶ συγγενῶν ἐν ἥττονι μέρει κείμενος αὐτοῖς ἀγὼν, μέγιστος δὲ καὶ πρῶτος ὁ περὶ τοῦ καθηγιασμένου ναοῦ φόβος.
\vs{19}Ἦν δὲ καὶ τοῖς ἐν τῇ πόλει κατειλημμένοις οὐ πάρεργος ἀγωνία ταρασσομένοις τῆς ἐν ὑπαίθρῳ προσβολῆς.

\vs{20}Καὶ πάντων ἤδη προσδοκώντων τὴν ἐσομένην κρίσιν, καὶ ἤδη συμμιξάντων τῶν πολεμίων, καὶ τῆς στρατιᾶς ἐκταγείσης, καὶ τῶν θηρίων ἐπὶ μέρος εὔκαιρον ἀποκατασταθέντων, τῆς τε ἵππου κατὰ κέρας τεταγμένης,

\vs{21}Συνιδὼν ὁ Μακκαβαῖος τὴν τῶν πληθῶν παρουσίαν, καὶ τῶν ὅπλων τὴν ποικίλην παρασκευὴν, τήν τε τῶν θηρίων ἀγριότητα, προτείνας τὰς χεῖρας εἰς τὸν οὐρανὸν, ἐπεκαλέσατο τὸν τερατοποιὸν Κύριον τὸν κατόπτην, γινώσκων ὅτι οὐκ ἔστι διʼ ὅπλων ἡ νίκη, καθὼς δὲ ἂν αὐτῷ κριθείη, τοῖς ἀξίοις περιποιεῖται τὴν νίκην.
\vs{22}Ἔλεγε δὲ ἐπικαλούμενος τόνδε τὸν τρόπον, σὺ, Δέσποτα, ἀπέστειλας τὸν ἄγγελόν σου ἐπὶ Ἑζεκίου τοῦ βασιλέως τῆς Ἰουδαίας, καὶ ἀνεῖλες ἐκ τῆς παρεμβολῆς Σενναχηρεὶμ εἰς ἑκατὸν ὀγδοηκονταπέντε χιλιάδας.
\vs{23}Καὶ νῦν, Δυνάστα τῶν οὐρανῶν, ἀπόστειλον ἄγγελον ἀγαθὸν ἔμπροσθεν ἡμῶν εἰς δέος καὶ τρόμον.
\vs{24}Μεγέθει βραχίονός σου καταπλαγείησαν οἱ μετὰ βλασφημίας παραγενόμενοι ἐπὶ τὸν ἅγιόν σου λαόν· καὶ οὗτος μὲν ἐν τούτοις ἔληξεν.

\vs{25}Οἱ δὲ περὶ τὸν Νικάνορα μετὰ σαλπίγγων καὶ παιάνων προσῆγον,
\vs{26}οἱ δὲ περὶ τὸν Ἰούδαν μετʼ ἐπικλήσεως καὶ εὐχῶν συνέμιξαν τοῖς πολεμίοις.
\vs{27}Καὶ ταῖς μὲν χερσὶν ἀγωνιζόμενοι, ταῖς δὲ καρδίαις πρὸς τὸν Θεὸν εὐχόμενοι, κατέστρωσαν οὐδὲν ἧττον μυριάδων τριῶν καὶ πεντακισχιλίων, τῇ τοῦ Θεοῦ μεγάλως εὐφρανθέντες ἐπιφανείᾳ.

\vs{28}Γενόμενοι δὲ ἀπὸ τῆς χρείας, καὶ μετὰ χαρᾶς ἀναλύοντες, ἐπέγνωσαν προπεπτωκότα Νικάνορα σὺν τῇ πανοπλίᾳ.
\vs{29}Γενομένης δὲ κραυγῆς καὶ ταραχῆς, εὐλόγουν τὸν Δυνάστην τῇ πατρίῳ φωνῇ.

\vs{30}Καὶ προσέταξεν ὁ καθʼ ἅπαν σώματι καὶ ψυχῇ πρωταγωνιστὴς ὑπὲρ τῶν πολιτῶν, ὁ τὴν τῆς ἡλικίας εὔνοιαν εἰς ὁμοεθνεῖς διαφυλάξας, τὴν τοῦ Νικάνορος κεφαλὴν ἀποτεμόντας, καὶ τὴν χεῖρα σὺν τῷ ὤμῳ φέρειν εἰς Ἱεροσόλυμα.

\vs{31}Παραγενόμενος δὲ ἐκεῖ, καὶ συγκαλέσας τοὺς ὁμοεθνεῖς, καὶ τοὺς ἱερεῖς πρὸ τοῦ θυσιαστηρίου στήσας, μετεπέμψατο τοὺς ἐκ τῆς ἄκρας.
\vs{32}Καὶ ἐπιδειξάμενος τὴν τοῦ μιαροῦ Νικάνορος κεφαλὴν, καὶ τὴν χεῖρα τοῦ δυσφήμου, ἣν ἐκτείνας ἐπὶ τὸν ἅγιον τοῦ παντοκράτορος οἶκον ἐμεγαλαύχησε.

\vs{33}Καὶ τὴν γλῶσσαν τοῦ δυσσεβοῦς Νικάνορος ἐκτεμὼν, ἔφη κατὰ μέρος δώσειν τοῖς ὀρνέοις, τὰ δὲ ἐπίχειρα τῆς ἀνοίας κατέναντι τοῦ ναοῦ κρεμᾶσαι.
\vs{34}Οἱ δὲ πάντες εἰς τὸν οὐρανὸν εὐλόγησαν τὸν ἐπιφανῆ Κύριον, λέγοντες, εὐλογητὸς ὁ διατηρήσας τὸν ἑαυτοῦ τόπον ἀμίαντον.
\vs{35}Ἐξέδησε δὲ τὴν τοῦ Νικάνορος κεφαλὴν ἐκ τῆς ἄκρας, ἐπίδηλον πᾶσι καὶ φανερὸν τῆς τοῦ Κυρίου βοηθείας σημεῖον.

\vs{36}Καὶ ἐδογμάτισαν πάντες μετὰ κοινοῦ ψηφίσματος μηδαμῶς ἐᾶσαι ἀπαρασήμαντον τὴνδε τὴν ἡμέραν· ἔχειν δὲ ἐπίσημον τὴν τρισκαιδεκάτην τοῦ δωδεκάτου μηνὸς, Ἄδαρ λέγεται τῇ Συριακῇ φωνῇ, πρὸ μιᾶς ἡμέρας τῆς Μαρδοχαϊκῆς ἡμέρας.
\vs{37}Τῶν οὖν κατὰ Νικάνορα χωρησάντων οὕτω, καὶ ἀπʼ ἐκείνων τῶν καιρῶν κρατηθείσης τῆς πόλεως ὑπὸ τῶν Ἑβραίων. Καὶ αὐτὸς αὐτόθι καταπαύσω τὸν λόγον.

\vs{38}Καὶ εἰ μὲν καλῶς καὶ εὐθίκτως τῇ συντάξει, τοῦτο καὶ αὐτὸς ἤθελον· εἰ δὲ εὐτελῶς καὶ μετρίως, τοῦτο ἐφικτὸν ἦν μοι.
\vs{39}Καθάπερ γὰρ οἶνον καταμόνας πίνειν, ὡσαύτως δὲ καὶ ὕδωρ πάλιν, πολέμιον· ὃν δὲ τρόπον οἶνος ὕδατι συγκερασθεὶς ἡδὺς, καὶ ἐπιτερπῆ τὴν χάριν ἀποτελεῖ, οὕτω καὶ τὸ τῆς κατασκευῆς τοῦ λόγου τέρπει τὰς ἀκοὰς τῶν ἐντυγχανόντων τῇ συντάξει· ἐνταῦθα δὲ ἔσται ἡ τελευτή.


\def\book{ΜΑΚΚΑΒΑΙΩΝ Γʹ}
\biblebook{ΜΑΚΚΑΒΑΙΩΝ Γʹ}


\lettrine[lines=2, loversize=0.2, nindent=0em, findent=.25em]{\textcolor{bookheadingcolor}{Ὁ}}{} ΔΕ Φιλοπάτωρ μαθὼν παρὰ τῶν ἀνακομισθέντων τὴν γενομένην τῶν ἐπʼ αὐτοῦ κρατουμένων τόπων ἀφαίρεσιν ὑπὸ Ἀντιόχου, παραγγείλας ταῖς πάσαις δυνάμεσι πεζικαῖς τε καὶ ἱππικαῖς αὐτοῦ, καὶ τὴν ἀδελφὴν Ἀρσινόην συμπαραλαβὼν, ἐξώρμησε μέχρι τῶν κατὰ Ῥαφίαν τόπων, ὅπου παρεμβεβλήκεισαν οἱ περὶ Ἀντίοχον.

\vs{2}Θεόδοτος δέ τις ἐκπληρῶσαι τὴν ἐπιβουλὴν διανοηθεὶς, παραλαβὼν τῶν προϋποτεταγμένων αὐτῷ ὅπλων Πτολεμαϊκῶν τὰ κράτιστα, διεκομίσθη νύκτωρ ἐπὶ τὴν τοῦ Πτολεμαίου σκηνὴν, ὡς μόνος κτεῖναι αὐτὸν, καὶ ἐν τούτῳ διαλῦσαι τὸν πόλεμον.
\vs{3}Τοῦτον δὲ διαγαγὼν Δωσίθεος ὁ Δριμύλου λεγόμενος, τὸ γένος Ἰουδαῖος, ὕστερον δὲ μεταβαλὼν τὰ νόμιμα, καὶ τῶν πατρίων δογμάτων ἀπηλλοτριωμένος, ἄσημόν τινα κατέκλινεν ἐν τῇ σκηνῇ, ὃν συνέβη κομίσασθαι τὴν ἐκείνου κόλασιν.

\vs{4}Γενομένης δὲ καρτερᾶς μάχης, καὶ τῶν πραγμάτων μᾶλλον ἐῤῥωμένων τῷ Ἀντιόχῳ, ἱκανῶς ἡ Ἀρσινόη ἐπιπορευσαμένη τὰς δυνάμεις παρεκάλει, μετὰ οἴκτου καὶ δακρύων, τοὺς πλοκάμους λελυμένη, βοηθεῖν ἑαυτοῖς τε καὶ τοῖς τέκνοις καὶ γυναιξὶ θαῤῥαλέως, ἐπαγγελλομένη δώσειν νικήσασιν ἑκάστῳ δύο μνᾶς χρυσίου.
\vs{5}Καὶ οὕτω συνέβη τοὺς ἀντιπάλους ἐν χειρονομίαις διαφθαρῆναι, πολλοὺς δὲ καὶ δορυαλώτους συλληφθῆναι

\vs{6}Κατακρατήσας δὲ τῆς ἐπιβουλῆς ἔκρινε τὰς πλησίον πόλεις ἐπελθὼν παρακαλέσαι.
\vs{7}Ποιήσας δὲ τοῦτο, καὶ τοῖς τεμένεσι δωρεὰς ἀπονείμας, εὐθαρσεῖς τοὺς ὑποτεταγμένους κατέστησε.
\vs{8}Τῶν δὲ Ἰουδαίων διαπεμψαμένων πρὸς αὐτὸν ἀπὸ τῆς γερουσίας καὶ τῶν πρεσβυτέρων τοὺς ἀσπασομένους αὐτὸν. Καὶ ξένια κομιοῦντας, καὶ ἐπὶ τοῖς συμβεβηκόσι συγχαρησομένους, συνέβη μᾶλλον αὐτὸν προθυμηθῆναι ὡς τάχιστα πρὸς αὐτοὺς παραγενέσθαι.

\vs{9}Διακομισθεὶς δὲ εἰς Ἱεροσόλυμα, καὶ θύσας τῷ μεγίστῳ Θεῷ, καὶ χάριτας ἀποδιδοὺς, καὶ τῶν ἑξῆς τι τῷ τόπῳ ποιήσας, καὶ δὴ παραγενόμενος εἰς τὸν τόπον, καὶ τῇ σπουδαιότητι καὶ εὐπρεπείᾳ καταπλαγείς,
\vs{10}θαυμάσας δὲ καὶ τὴν τοῦ ἱεροῦ εὐταξίαν, ἐνεθυμήθη βουλεύσασθαι εἰσελθεῖν εἰς τὸν ναόν.

\vs{11}Τῶν δὲ εἰπόντων μὴ καθήκειν γίνεσθαι τοῦτο, διὰ τὸ μηδὲ τοῖς ἐκ τοῦ ἔθνους ἐξεῖναι εἰσιέναι, μηδὲ πᾶσι τοῖς ἱερεῦσιν, ἀλλʼ ἢ μονῳ τῷ προηγουμένῳ πάντων ἀρχιερεῖ, καὶ τούτῳ ἅπαξ κατʼ ἐνιαυτὸν, οὐδαμῶς ἠβούλετο πείθεσθαι.
\vs{12}Τοῦ τε νόμου παραναγνωσθέντος, οὐδαμῶς ἀπέλιπε προφερόμενος ἑαυτὸν, δεῖν εἰσελθεῖν, λέγων, καὶ εἰ ἐκεῖνοι ἐστέρηνται ταύτης τῆς τιμῆς, ἐμὲ οὺ δεῖ.
\vs{13}Καὶ ἐπυνθάνετο, διὰ τίνα αἰτίαν εἰσερχόμενον αὐτὸν εἰς πᾶν τέμενος, οὐθεὶς ἐλώλυσε τῶν παρόντων.

\vs{14}Καί τις ἀπρονοήτως ἔφη κακῶς αὐτὸ τοῦτο τερατεύεσθαι.
\vs{15}Γενομένου δέ φησι τούτου διά τινα αἰτίαν, οὐχὶ πάντως εἰσελεύσεσθαι καὶ θελόντων αὐτῶν καὶ μή;

\vs{16}Τῶν δὲ ἱερέων ἐν ταῖς ἁγίαις ἐσθήσεσι προπεσόντων, καὶ δεομένων τοῦ μεγίστου Θεοῦ βοηθεῖν τῇ ἐνεστώσῃ ἀνάγκῃ, καί τὴν ὁρμὴν τοῦ κακῶς ἐπιβαλλομένου μεταθεῖναι, κραυγῆς τε μετὰ δακρύων τὸ ἱερὸν ἐμπλησάνων,
\vs{17}οἱ κατὰ τὴν πόλιν ἀπολιπόμενοι, ταραχθέντες ἐξεπήδησαν, ἄδηλον τιθέμενοι τὸ γινόμενον.

\vs{18}Αἵ τε κατάκλειστοι παρθένοι ἐν θαλάμοις σὺν ταῖς τεκούσαις ἐξώρμησαν· καὶ σποδῷ καὶ κόνει τὰς κεφαλὰς πασάμεναι, γόων τε καὶ στεναγμῶν τὰς πλατείας ἐνεπίμπλων.
\vs{19}Αἱ δὲ καὶ προσαρτίως ἐσταλμέναι, τοὺς πρὸς ἀπάντησιν διατεταγμένους παστοὺς, καὶ τὴν ἁρμόζουσαν αἰδὼ παραλείπουσαι, δρόμον ἄτακτον ἐν τῇ πόλει συνίσταντο.
\vs{20}Τὰ δὲ νεογνὰ τῶν τέκνων, αἵ τε πρὸς τούτοις μητέρες καὶ τιθηνοὶ παραλιποῦσαι ἄλλως καὶ ἄλλως, αἱ μὲν κατʼ οἴκους, αἱ δὲ κατὰ τὰς ἀγυιὰς ἀνεπιτρέπτως εἰς τὸ πανυπέρτατον ἱερὸν ἠθροίζοντο.
\vs{21}Ποικίλη δὲ ἦν τῶν εἰς τοῦτο συλλεγομένων ἡ δέησις ἐπὶ τοῖς ἀνοσίως ὑπʼ ἐκείνου κατεγχειρουμένοις.

\vs{22}Σύν τε τούτοις οἱ τῶν πολιτῶν θρασυνθέντες οὐκ ἠνείχοντο τέλεον αὐτοῦ ἐπικειμένου, καὶ τὸ τῆς προθέσεως αὐτοῦ ἐκπληροῦν διανοουμένου.
\vs{23}Θωνήσαντες δὲ τὴν ὁρμὴν ἐπὶ τὰ ὅπλα ποιήσασθαι, καὶ θαῤῥαλέως ὑπὲρ τοῦ πατρῴου νόμου τελευτᾷν, ἱκανὴν ἐποίησαν ἐν τῷ τόπῳ τραχύτητα, μόλις τε ὑπό τε τῶν γεραιῶν καὶ τῶν πρεσβυτέρων ἀποτραπέντες ἐπὶ τὴν αὐτὴν τῆς δεήσεως ἔστησαν στάσιν.

\vs{24}Καὶ τὸ μὲν πλῆθος, ὡς ἔμπροσθεν, ἐν τούτοις ἀνεστρέφετο δεόμενον.
\vs{25}Οἱ δὲ περὶ τὸν βασιλέα πρεσβύτεροι πολλαχῶς ἐπειρῶντο τὸν ἀγέρωχον αὐτοῦ νοῦν ἐξιστάνειν τῆς ἐντεθυμημένης βουλῆς.
\vs{26}Θρασυνθεὶς δὲ καὶ πάντα παραπέμψας, ἤδη καὶ πρόσβασιν ἀποιεῖτο, τέλος ἐπιθήσειν δοκῶν τῷ προειρημένῳ.

\vs{27}Ταῦτʼ οὖν καὶ οἱ περὶ αὐτὸν ὄντες θεωροῦντες, ἐτράπησαν εἰς τὸ σὺν τοῖς ἡμετέροις ἐπικαλεῖσθαι τὸν πᾶν κράτος ἔχοντα, τοῖς παροῦσιν ἐπαμῦναι, μὴ περιδόντα τὴν ἄνομον καὶ ὑπερήφανον πρᾶξιν.
\vs{28}Ἐκ δὲ τῆς πυκνοτάτης τε καὶ ἐμπόνου τῶν ὄχλων συναγομένης κραυγῆς, ἀνείκαστός τις ἦν βοή.
\vs{29}Δοκεῖν γὰρ ἦν μὴ μόνον τοὺς ἀνθρώπους, ἀλλὰ καὶ τὰ τείχη καὶ τὸ πᾶν ἔδαφος ἠχεῖν, ἅτε δὴ τῶν πάντων τότε θάνατον ἀλλασσομένων ἀντὶ τῆς τοῦ τόπου βεβηλώσεως.

\ch{2}
Ὁ μὲν οὖν ἀρχιερεὺς Σίμων ἐξεναντίας τοῦ ναοῦ κάμψας τὰ γόνατα, καί τὰς χεῖρας προτείνας εὐτάκτως, ἐποιήσατο τὴν δέησιν τοιαύτην·

\vs{2}Κύριε Κύριε βασιλεῦ τῶν οὐρανῶν, καὶ Δέσποτα πάσης κτίσεως, ἅγιε ἐν ἁγίοις, μόναρχε, παντοκράτωρ, πρόσχες ἡμῖν καταπονουμένοις ὑπὸ ἀνοσίου καὶ βεβήλου θράσει καὶ σθένει πεφρυαγμένου.
\vs{3}Σὺ γὰρ ὁ κτίσας τὰ πάντα καὶ τῶν ὅλων ἐπικρατῶν δυνάστης δίκαιος εἶ, καὶ τοὺς ὕβρει καὶ ἀγερωχίᾳ πράσσοντάς τι κρίνεις.

\vs{4}Σὺ τοὺς ἔμπροσθεν ἀδικίαν ποιήσαντας, ἐν οἷς καὶ γίγαντες ἦσαν ῥώμῃ καὶ θράσει πεποιθότες, διέφθειρας, ἐπαγαγὼν αὐτοῖς ἀμέτρητον ὕδωρ.
\vs{5}Σὺ τοὺς ὑπερηφανίαν ἐργαζομένους Σοδομίτας, αιαδήλους ταῖς κακίαις γενομένους, πυρὶ καὶ θείῳ κατέφλεξας, παράδειγμα τοῖς ἐπιγινομένοις καταστήσας.

\vs{6}Σὺ τὸν θρασὺν Φαραὼ καταδουλωσάμενον τὸν λαόν σου τὸν ἅγιον Ἰσραὴλ, ποικίλαις καὶ πολλαῖς δοκιμάσας τιμωρίαις, ἐγνώρισας τὴν σὴν δυναστείαν· ἐφʼ αἷς ἐγνώρισας τὸ μέγα σου κράτος.
\vs{7}Καὶ ἐπιδιώξαντα αὐτὸν σὺν ἅρμασι καὶ ὄχλων πλήθει, ἐπέκλυσας βάθει θαλάσσης, τοὺς δὲ ἐνπιστεύσαντας ἐπὶ σοὶ τῷ τῆς ἁπάσης κτίσεως δυναστεύοντι, σώους διεκόμισας.
\vs{8}Οἳ καὶ συνειδότες ἔργα σῆς χειρός, ᾔνεσάν σε τὸν παντοκράτορα.

\vs{9}Σὺ, βασιλεῦ, κτίσας τὴν ἀπέραντον καὶ ἀμέτρητον γῆν, ἐξελέξω τὴν πόλιν ταύτην, καὶ ἁγιάσας τὸν τόπον τοῦτον εἰς ὄνομά σοι τῷ τῶν ἁπάντων ἀπροσδεεῖ, καὶ παρεδόξασας ἐν ἐπιφανείᾳ μεγαλοπρεπεῖ, σύστασιν ποιησάμενος αὐτοῦ πρὸς δόξαν τοῦ μεγάλου καὶ ἐντίμου ὀνόματός σου.

\vs{10}Καὶ ἀγαπῶν τὸν οἶκον τοῦ Ἰσραὴλ, ἐπηγγείλω δὴ ὃτι ἐὰν γένηται ἡμῶν ἀποστροφὴ, καὶ καταλάβῃ ἡμᾶς στενοχωρία, καὶ ἐλθόντες εἰς τὸν τόπον τοῦτου δεηθῶμεν, εἰσακούσῃ τῆς δεήσεως ἡμῶν.
\vs{11}Καὶ δὴ πιστὸς εἶ καὶ ἀληθινός.

\vs{12}Ἐπεὶ δὲ πλεονάκις θλιβέντων τῶν πατέρων ἡμῶν ἐβοήθησας αὐτοῖς ἐν τῇ ταπεινώσει, καὶ ἐῤῥύσω αὐτοὺς ἐκ μεγάλων κινδύνων,
\vs{13}ἰδοὺ δὴ νῦν, ἅγιε βασιλεῦ, διὰ τὰς πολλὰς καὶ μεγάλας ἡμῶν ἁμαρτίας καταπονούμεθα, καὶ ὑπετάγημεν τοῖς ἐχθροῖς ἡμῶν, καὶ παρείμεθα ἐν ἀδυναμίαις.
\vs{14}Ἐν δὲ τῇ ἡμετέρᾳ καταπτώσει ὁ θρασὺς καὶ βέβηλος οὗτος ἐπιτηδεύει καθυβρίσαι τὸν ἐπὶ τῆς γῆς ἀναδεδειγμένον τῷ ὀνόματι τῆς δόξης σου ἅγιον τόπον.

\vs{15}Τὸ μὲν γὰρ οἰκητήριόν σου οὐρανὸς τοῦ οὐρανοῦ ἀνέφικτος ἀνθρώποις ἐστίν.
\vs{16}Ἀλλαʼ ἐπεὶ εὐδόκησας τὴν δόξαν σου ἐν τῷ λαῷ σου Ἰσραὴλ, ἡγίασας τὸν τόπον τοῦτον.
\vs{17}Μὴ ἐκδικήσῃς ἡμᾶς ἐν τῇ τούτων ἀκαθαρσίᾳ, μηδὲ εὐθύνῃς ἡμᾶς ἐν βεβηλώσει· ἵνα μὴ καυχήσωνται οἱ παράνομοι ἐν θυμῷ αὐτῶν, μηδὲ ἀγαλλιάσωνται ἐν ὑπερηφανίᾳ γλώσσης αὐτῶν, λέγοντες,
\vs{18}ἡμεῖς κατεπατήσαμεν τὸν οἶκον τοῦ ἁγιασμοῦ, ὡς καταπατοῦνται οἱ οἶκοι τῶν προσοχθισμάτων.

\vs{19}Ἀπάλειψον τὰς ἁμαρτίας ἡμῶν, καὶ διασκέδασον τὰς ἀμπλακίας ἡμῶν, καὶ ἐπίφανον τὸ ἔλεός σου κατὰ τὴν ὥραν ταύτην.
\vs{20}Ταχὺ προκαταλαβέτωσαν ἡμᾶς οἱ οἰκτιρμοί σου· καὶ δὸς αἰνέσεις ἐν στόματι τῶν καταπεπτωκότων καὶ συντετριμμένων τὰς ψυχὰς, ποιήσας ἡμῖν εἰρήνην.

\vs{21}Ἐνταῦθα ὁ πάντων ἐπόπτης Θεὸς, καὶ πρὸ πάντωρ ἅγιος ἐν ἁγίοις, εἰσακούσας τῆς ἐνθέσμου λιτανείας, τὸν ὕβρει καὶ θράσει μεγάλως ἐπηρμένον ἐμάστιξεν αὐτόν,
\vs{22}ἔνθεν καὶ ἔνθεν κραδάνας αὐτὸν ὡς κάλαμον ὑπὸ ἀνέμου, ὥστε κατʼ ἐδάφους ἄπρακτον ἔτι καὶ τοῖς μέλεσι παραλελυμένον, μηδὲ φωνῆσαι δύνασθαι δικαίᾳ περιπεπλεγμένον κρίσει.

\vs{23}Ὅθεν οἵ τε φίλοι καὶ οἱ σωματοφύλακες αὐτοῦ ταχεῖαν καὶ ὀξεῖαν ἰδόντες τὴν καταλαβοῦσαν αὐτὸν εὐθύναν, φοβούμενοι μὴ καὶ τὸ ζῇν ἐκλείπῃ, ταχέως αὐτὸν ἐξείλκυσαν ὑπερβάλλοντι καταπεπληγμένοι φόβῳ.
\vs{24}Ἐν χρόνῳ δὲ ὕστερον ἀναλεξάμενος ἑαυτὸν, οὐδαμῶς εἰς μετάμελον ἦλθεν ἐπιτιμηθεὶς, μετʼ ἀπειλῆς δὲ πικρᾶς ἀνέλυσε.
\vs{25}Διακομισθεὶς δὲ εἰς τὴν Αἴγυπτον, καὶ τὰ τῆς κακίας ἐπαύξων, διά δὲ τῶν προαποδεδιγμένων συμποτῶν καὶ ἑταίρων τοῦ παντὸς δικαίου κεχωρισμένων,
\vs{26}οὐ μόνον ταῖς ἀναριθμήτοις ἀσελγείαις διηρκέσθη, ἀλλὰ καὶ ἐπὶ τοσοῦτον θράσους προῆλθεν, ὥστε δυσφημίας ἐν τοῖς τόποις συνίστασθαι, καὶ πολλοὺς τῶν φίλων ἀτενίζοντας εἰς τὴν τοῦ βασιλέως πρόθεσιν καὶ αὐτοὺς ἕπεσθαι τῇ ἐκείνου θελήσει.

\vs{27}Προέθετο δὲ δημοσίᾳ κατὰ τοῦ ἔθνους διαδοῦναι ψόγον· καὶ ἐπὶ τοῦ κατὰ τὴν αὐλὴν πύργου στήλην ἀναστήσας, ἐξεκόλαψε γραφήν,
\vs{28}μηδένα τῶν μὴ θυόντων εἰς τὰ ἱερὰ αὐτῶν εἰσιέναι, πάντας δὲ τοὺς Ἰουδαίους εἰς λαογραφίαν καὶ οἰκετικὴν διάθεσιν ἀχθῆναι, τοὺς δὲ ἀντιλέγοντας βίᾳ φερομένους τοῦ ζῇν μεταστῆσαι,
\vs{29}τούτους τε ἀπογραφομένους χαράσσεσθαι καὶ διὰ πυρὸς εἰς τὸ σῶμα παρασήμῳ Διονύσου κισσοφύλλῳ, οὓς καὶ καταχωρίσαι εἰς τὴν προσυνεσταλμένην αὐθεντίαν.

\vs{30}Ἵνα δὲ μὴ τοῖς πᾶσιν ἀπεχθόμενος φαίνηται, ὑπέγραψεν, ἐὰν δέ τινες ἐξ αὐτῶν προαιρῶνται ἐν τοῖς κατὰ τὰς τελετὰς μεμυημένοις ἀναστρέφεσθαι, τούτους ἰσοπολίτας Ἀλεξανδρεῦσιν εἶναι.

\vs{31}Ἔνιοι μὲν οὖν ἐπι πολεως τὰς τῆς πόλεως εὐσεβείας ἐπιβάθρας στυγοῦντες, εὐχερῶς ἑαυτοὺς ἐδίδοσαν, ὡς μεγάλης τινὸς κοινωνήσοντες εὐκλείας ἀπὸ τῆς ἐσομένης τῷ βασιλεῖ συναναστροφῆς.
\vs{32}Οἱ δὲ πλεῖστοι γενναίᾳ ψυχῇ ἐνίσχυσαν καὶ οὐ διέστησαν τῆς εὐσεβείας· τά τε χρήματα περὶ τοῦ ζῇν ἀντικαταλλασσομενοι, ἀδεῶς ἐπειρῶντο ἑαυτοὺς ῥύσασθαι ἐκ τῶν ἀπογραφῶν.
\vs{33}Εὐέλπιδές δὲ καθειστήκεισαν ἀντιλήμψεως τεύξεασθαι, καὶ τοὺς ἀποχωροῦντας ἐξ αὐτῶν ἐβδελύσσοντο, καὶ ὡς πολεμίους τοῦ ἔθνους ἔκρινον, καὶ τῆς κοινῆς συναναστροφῆς καὶ εὐχρηστίας ἐστέρουν.

\ch{3}
Ἃ καὶ μεταλαμβάνων ὁ δυσσεβὴς ἐπι τοσοῦτον ἐχόλησεν, ὥστε οὐ μόνον τοῖς κατʼ Ἀλεξάνδρειαν διοργίζεσθαι, ἀλλὰ καὶ τοῖς ἐν τῇ χώρᾳ βαρυτέρως ἐναντιωθῆναι, καὶ προστάξαι σπεύσαντας συναγαγεῖν πάντας ἐπιτοαυτὸ, καὶ χειρίστῳ μόρῳ τοῦ ζῇν μεταστῆσαι.

\vs{2}Τούτων δὲ οἰκονομουμένων, φήμη δυσμενὴς ἐξηχεῖτο κατὰ τοῦ γένους ἀνθρώποις συμφρονοῦσιν εἰς κακοποίησιν, ἀφορμῆς διδομένης εἰς διάθεσιν, ὡς ἂν ἀπὸ τῶν νομίμων αὐτοὺς κωλυόντων.
\vs{3}Οἱ δὲ Ἰουδαῖοι τὴν μὲν πρὸς τοὺς βασιλεῖς εὔνοιαν καὶ πίστιν ἀδιάστροφον ἦσαν διαφυλάσσοντες·
\vs{4}σεβόμενοι δὲ τὸν Θεόν καὶ τῷ τούτου νόμῳ πολιτευόμενοι, χωρισμὸν ἐποίουν ἐπὶ τινων καὶ καταστροφάς διʼ ἣν αἰτίαν ἔνιος ἀπεχθεῖς ἐφαίνοντο·
\vs{5}Τῇ δὲ τῶν δικαίων εὐπραξίᾳ κοσμοῦντες τὴν συναναστροφήν, ἅπασιν ἀνθρώποις εὐδόκιμοι καθειστήκεισαν.

\vs{6}Τὴν μὲν οὖν περὶ τοῦ γένους ἐν πᾶσι θρυλλουμένην εὐπραξίαν οἱ ἀλλόφυλοι οὐδαμῶς διηριθμήσαντο.
\vs{7}Τὴν δὲ περὶ τῶν προσκυνήσεων καὶ τροφῶν διάστασιν ἐθρύλλουν, φάσκοντες μήτε τῷ βασιλεῖ μήτε ταῖς δυνάμεσιν ὁμοσπόνδους τοὺς ἀνθρώπους γενέσθαι, δυσμενεῖς δὲ εἶναι καὶ μέγα τι τοῖς πράγμασιν ἐναντιουμένους· καὶ οὐ τῷ τυχόντι περιήψαν ψόγῳ.

\vs{8}Οἱ δὲ κατὰ τὴν πόλιν Ἕλληνες οὐδὲν ἠδικημένοι, ταραχὴν ἀπροσδόκητον περὶ τοὺς ἄνθρώπους θεωροῦντες, καὶ συνδρομὰς ἀπροσκόπους γινομένας βοηθεῖν μὲν οὐκ ἔσθενον· τυραννικὴ γὰρ ἦν ἡ διάθεσις· παρεκάλουν δὲ καὶ δυσφόρως εἶχον, καὶ μεταπεσεῖσθαι ταῦτα ὑπελάμβανον·
\vs{9}Μὴ γὰρ οὕτως παροραθήσεται τηλικοῦτο σύστεημα μηδὲν ἠγνοηκώς.
\vs{10}Ἤδη δὲ καί τινες γείτονές τε καὶ φίλοι καὶ συμπραγματευόμενοι, μυστικῶς τινας ἐπισπώμενοι, πίστεις ἐδίδουν συνασπιεῖν, καὶ πᾶν ἐκτενὲς προσοίσεσθαι πρὸς ἀντίληψιν.

\vs{11}Ἐκεῖνος μὲν οὖν τῇ κατὰ τὸ παρὸν εὐημερίᾳ γεγαυρωμένος, καὶ οὐ καθορῶν τὸ τοῦ μεγίστου Θεοῦ κράτος, ὑπολαμβάνων δὲ διηνεκῶς ἐν τῇ αὐτῇ διαμένειν βουλῇ, ἔγραψε κατʼ αὐτῶν ἐπιστολὴν τήνδε

\vs{12}Βασιλεὺς Πτολεμαῖος Φιλοπάτωρ τοῖς κατʼ Αἴγυπτον, καὶ κατὰ τόπον στρατηγοῖς καὶ στρατιώταις, χαίρειν καὶ ἐῤῥῶσθαι.
\vs{13}Ἔῤῥωμαι δὲ καὶ ἐγὼ αὐτὸς καὶ τὰ πράγματα ἡμῶν.
\vs{14}Ἐκ τῆς εἰς τὴν Ἀσίαν γενομένης ἡμῖν ἐπιστρατίας, ἧς ἴστε καὶ αὐτοί, τῇ τῶν θεῶν πρὸς ἡμᾶς ἀπροπτώτῳ συμμαχίᾳ καὶ τῇ ἡμετέρᾳ δὲ ῥώμῃ κατὰ λόγον ἐπʼ ἄριστον τέλος ἀχθείσης,
\vs{15}ἡγησάμεθα μὴ βίᾳ δόρατος, ἐπιεικείᾳ δὲ καὶ πολλῇ φιλανθρωπίᾳ τιθηνήσασθαι τὰ κατοικοῦντα κοίλην Συρίαν καὶ Φοινίκην ἔθνη, εὖ ποιήσαί τε ἀσμένως.

\vs{16}Καὶ τοῖς κατὰ πόλεσιν ἱεροῖς ἀπονείμαντες προσόδους πλείστας, προήχθημεν καὶ εἰς τὰ Ἰεροσόλυμα, ἀναβάντες τιμῆσαι τὸ ἱερὸν τῶν ἀλιτηρίων καὶ μηδέποτε ληγόντων τῆς ἀνοίας.
\vs{17}Οἱ δὲ λόγῳ μὲν τὴν ἡμετέραν ἀποδεξάμενοι παρουσίαν, τῷ δὲ πράγματι νόθως, προθυμηθέντων ἡμῶν εἰσελθεῖν εἰς τὸν ναὸν αὐτῶν, καὶ τοῖς ἐκπρεπέσιν καὶ καλλίστοις ἀναθήμασι τιμῆσαι,
\vs{18}τύφοις φερόμενοι παλαιοτέροις εἶρξαν ἡμᾶς τῆς εἰσόδου, ἀπολειπόμενοι τῆς ἡμετέρας ἀλκῆς, διʼ ἣν ἔχομεν πρὸς ἅπαντας ἀνθρώπους φιλανθρωπίαν.
\vs{19}Τὴν δὲ αὐτῶν εἰς ἡμᾶς δυσμενειαν ἔκδηλον καθιστάντες, ὡς μονώτατοι τῶν ἐθνῶν βασιλεῦσι καὶ τοῖς ἑαυτῶν εὐεργέταις ὑψαυχενοῦντες οὐδὲν γνήσιον βούλονται φέρειν.

\vs{20}Ἡμεῖς δὲ τῇ τούτων ἀνοίᾳ συμπεριενεχθέντες, καὶ μετὰ νίκης διακομισθέντες, καὶ εἰς τὴν Αἴγυπτον τοῖς πᾶσιν ἔθνεσιν φιλανθρώπως ἀπαντήσαντες, καθὼς ἔπρεπεν ἐποιήσαμεν.
\vs{21}ἐν δὲ τούτοις πρὸς τοὺς ὁμοφύλους αὐτῶν ἀμνησικακίαν ἅπασιν γνωρίζοντες, διά τε τὴν συμμαχίαν καὶ τὰ πεπιστευμένα μετὰ ἁπλότητος αὐτοῖς ἀρχῆθεν μύρια πράγματα, ἐξαλλοιῶσαι, ἐβουλήθημεν καὶ πολιτείας αὐτοὺς Ἀλεξανδρέων καταξιῶσαι, καὶ μετόχους τῶν ἀεὶ ἱερέων καταστῆσαι.

\vs{22}Οἱ δὲ τοὐναντίον ἐκδεχόμενοι, καὶ τῇ συμφύτῳ κακοηθείᾳ τὸ καλὸν ἀπωσάμενοι, διηνεκῶς δὲ εἰς τὸ φαῦλον ἐκνεύοντες,
\vs{23}οὐ μόνον ἀπεστρέψαντο τὴν ἀτίμητον πολιτείαν, ἀλλὰ καὶ βδελύσσονται λόγῳ τε καὶ σιγῇ τοὺς ἐν αὐτοῖς ὀλίγους πρὸς ἡμᾶς γνησίως διακειμένους, παρέκαστα ὑφορώμενοι διὰ τῆς δυσκλεεστάτης ἐμβιώσεως διὰ τάχους ἡμᾶς καταστρέψαι τὰ κατορθώματα.
\vs{24}Διὸ καὶ τεκμηρίοις καλῶς πεπεισμένοι τούτους κατὰ πάντα δυσνοεῖν ἡμῖν τρόπον, καὶ προνοούμενοι μήποτε αἰφνιδίου μετέπειτα ταραχῆς ἐνστάσης ἡμῖν, τοὺς δυσσεβεῖς τούτους κατὰ νώτου προδότας καὶ βαρβάρους ἔχωμεν πολεμίους.

\vs{25}Προστετάχαμεν ἅμα τῷ προσπεσεῖν τὴν ἐπιστολὴν τήνδε, αὐθωρὶ τοὺς ἐννεμομένους σὺν γυναιξὶ καὶ τέκνοις μετὰ ὕβρεων καὶ σκυλμῶν ἀποστεῖλαι πρὸς ἡμᾶς ἐνδεσμοῖς σιδηροῖς πάντοθεν κατακεκλεισμένους, εἰς ἀνήκεστον καὶ δυσκλεῆ πρέποντα δυσμενέσι φόνον.
\vs{26}Τούτων γὰρ ὁμοῦ κολασθέντων, διειλήφαμεν εἰς τὸν ἐπίλοιπον χρόνον τελείως ἡμῖν τὰ πράγματα ἐν εὐσταθείᾳ καὶ βελτίστῃ διαθέσει κατασταθήσεσθαι.

\vs{27}Ὅς δʼ ἂν σκεπάσῃ τινὰ τῶν Ἰουδαίων ἀπὸ γεραιοῦ μέχρι νηπίου μέχρι τῶν ὑπομασθίων, αἰσχίστοις βασάνοις ἀποτυμπανισθήσεται πανοικί.
\vs{28}Μηνύειν δέ τὸν βουλόμενον, ἐφʼ ᾧ τὴν οὐσίαν τοῦ ἐμπίπτοντος ὑπὸ τὴν εὐθύναν λήψεται, καὶ ἐκ τοῦ βασιλικοῦ ἀργυρίου δραχμὰς δισχιλίας, καὶ τῆς ἐλευθερίας τεύξεται καὶ στεφανωθήσεται.

\vs{29}Πᾶς δὲ τὸπος οὗ ἐὰν φωραθῇ τὸ σύνοκον σκεπαζόμενος Ἰουδαῖος, ἄβατος καὶ πυριφλεγὴς γινέσθω, καὶ πάσῃ θνητῇ φύσει κατὰ πάντα ἄχρηστος φανήσεται εἰς τὸν ἀεὶ χρόνον.
\vs{30}Καὶ ὁ μὲν τῆς ἐπιστολῆς τύπος οὕτως ἐγέγραπτο.

\ch{4}
Παντῇ δὲ ὅπου προσέπιπτε τοῦτο τὸ πρόσταγμα, δημοτελὴς συνίστατο τοῖς ἔθνεσιν εὐωχία μετὰ ἀλαλαγμῶν καὶ χαρᾶς, ὡς ἄν τῆς προκατεσκιῤῥωμένης αὐτοῖς πάλαι κατὰ διάνοιαν, μετὰ παῤῥησίας συνεκφαινομένης ἀπεχθείας.

\vs{2}Τοῖς δὲ Ἰουδαίοις ἀνήκεστον πένθος ἦν καὶ πανόδυρτος μετὰ δακρύων βοὴ, στεναγμοῖς πεπυρωμένης τῆς αὐτῶν πάντοθεν καρδίας, ὀλοφυρομένων τὴν ἀπροσδόκητον ἐξαίφνης ἐπικριθεῖσαν αὐτοῖς ὀλεθρίαν.
\vs{3}Τίς νομὸς ἢ πόλις, ἢ τίς τὸ σύνολον οἰκητὸς τὸπος, ἢ τίνες ἀγυιαὶ κοπετοῦ καὶ γόων ἐπʼ αὐτοῖς οὐκ ἐμπιπλῶντο;

\vs{4}Οὕτω γὰρ μετὰ πικρᾶς καὶ ἀνοίκτου ψυχῆς ὑπὸ τῶν κατὰ πόλιν στρατηγῶν ὁμοθυμαδὸν ἐξαπεστέλλοντο, ὥστε ἐπὶ ταῖς ἐξάλλοις τιμωρίαις καί τινας τῶν ἐχθρῶν, λαμβάνοντας πρὸ τῶν ὀφθαλμὼν τὸν κοινὸν ἔλεον, καὶ λογιζομένους τὴν ἄδηλον τοῦ βίου καταστροφὴν, δακρύειν αὐτῶν τρισάθλιον ἐξαποστολήν.
\vs{5}Ἤγετο γὰρ γεραιῶν πλῆθος πολιᾷ πεπυκασμένων, τὴν ἐκ τοῦ γήρως νωθρότητα ποδῶν ἐπικύφων, ἀνατροπῆς ὁρμῇ βιαίας, ἁπάσης αἰδοῦς ἄνευ πρὸς ὀξείαν καταχρωμένων πορείαν.

\vs{6}Αἱ δὲ ἄρτι πρὸς βίου κοινωνίαν γαμικὸν ὑπεληλυθυῖαι παστὸν νεάνιδες, ἀντὶ τέρψεως μεταλαβοῦσαι γόους, καὶ κόνει τὴν μυροβραχῆ πεφυρμέναι κόμην, ἀκαλύπτως δὲ ἀγόμεναι, θρῆνον ἀνθʼ ὑμεναίων ὁμοθυμαδὸν ἐξῆρχον, ὡς ἐσπαραγμέναι σκυλμοῖς ἀλλοεθνέσι.
\vs{7}Δέσμιαι δέ δημόσιαι μέχρι τῆς εἰς τὸ πλοῖον ἐμβολῆς εἵλκοντο μετὰ βίας.

\vs{8}Οἵ τε τούτων συζυγεῖς βρόχοις ἀντὶ στεφέων τοὺς αὐχένας περιπεπλεγμένοι μετὰ ἀκμαίας καὶ νεανικῆς ἡλικίας, ἀντὶ εὐωχίας καὶ νεωτερικῆς ῥαθυμίας τὰς ἐπιλοίπους τῶν γάμων ἡμέρας ἐν θρήνοις διῆγον, παρὰ πόδας ἤδη τὸν ᾅδην ὁρῶντες κείμενον.
\vs{9}Κατήχθησαν δὲ θηρίων τρόπον ἀγόμενοι σιδηροδέσμοις ἀνάγκαις· οἱ μὲν τοῖς ζυγοῖς τῶν πλοίων προσηλωμένοι τοὺς τραχήλους, οἱ δὲ τοὺς πόδας ἀῤῥήκτοις κατησθφαλισμένοι πέδαις,
\vs{10}ἔτι καὶ τῷ καθύπερθε πυκνῷ σανιδώματι διακειμένῳ τὸ φέγγος ἀποκλειόμενοι, ὅπως πάντοθεν ἐσκοτισμένοι τοὺς ὀφθαλμοὺς, ἀγωγὴν ἐπιβούλων ἐν παντὶ τῷ κατάπλῳ λαμβάνωσι.

\vs{11}Τούτων δὲ ἐπὶ τὴν λεγομένην Σχεδὶαν ἀχθέντων, καὶ τοῦ παράπλου περανθέντος, καθὼς ἦν δεδογματισμένον τῷ βασιλεῖ, προσέταξεν αὐτοὺς ἐν τῷ πρὸ τῆς πόλεως ἱπποδρόμῳ παρεμβαλεῖν ἀπλέτῳ καθεστῶτι περιμέτρῳ, καὶ πρὸς παραδειγματισμὸν ἄγαν εὐκαιροτάτῳ καθεστῶτι πᾶσι τοῖς καταπορευομένοις εἰς τὴν πόλιν, καὶ τοῖς ἐκ τούτων εἰς τὴν χώραν στελλομένοις πρὸς ἐκδημίαν· πρὸς τὸ μηδὲ ταῖς δυνάμεσιν αὐτοῦ κοινωνεῖν, μηδὲ τὸ σύνολον καταξιῶσαι περιβόλων.

\vs{12}Ὡς δὲ τοῦτο ἐγενήθη, ἀκούσας τοὺς ἐκ τῆς πόλεως ὁμοεθνεῖς κρυβῇ ἐκπορευομένους πυκνότερον ἀποδύρεσθαι τὴν ἀκλεᾶ τῶν ἀδελφῶν ταλαιπωρίαν,
\vs{13}διοργισθεὶς προσέταξε καὶ τούτοις ὁμοῦ τὸν αὐτὸν τρόπον ἐπιμελῶς ὡς ἐκείνοις ποιῆσαι, μὴ λειπομένοις κατὰ μηδένα τρόπον τῆς ἐκείνων τιμωρίας, Ἀπογραφῆναι δὲ πᾶν τὸ φῦλον ἐξ ὀνόματος·
\vs{14}οὐ γὰρ τὴν ἔμπροσθε βραχεῖ προδεδηλωμένην τῶν ἔργων κατὰπονον λατρείαν, στρεβλωθέντας δὲ ταῖς παρηγγελμέναις αἰκίαις τὸ τέλος ἀφανίσαι μιᾶς ὑπὸ καιρὸν ἡμέρας.
\vs{15}Ἐγίνετο μὲν οὖν ἡ τοὺτων ἀπογραφὴ μετὰ πικρᾶς σπουδῆς καὶ φιλοτίμου προσεδρίας ἀπὸ ἀνατολῶν ἡλίου μέχρι δυσμῶν, ἀνήνυτον λαμβάνουσα τὸ τέλος επὶ ἡμέρας τεσσαράκοντα.

\vs{16}Μεγάλως δὲ καὶ διηνεκῶς ὁ βασιλεὺς χαρᾷ πεπληρωμένος, συμπόσια ἐπὶ πάντων τῶν εἰδώλων συνιστάμενος, πεπλανημένῃ, πόῤῥω τῆς ἀληθείας φρενὶ καὶ βεβήλῳ στόματι, τὰ μὲν κωφὰ καὶ μὴ δυνάμενα αὐτοῖς λαλεῖν ἢ ἀρήγειν, ἐπαινῶν, εἰς δὲ τὸν μέγιστον Θεὸν τὰ μὴ καθήκοντα λαλῶν.

\vs{17}Μετὰ δὲ τὸ προειρημένον τοῦ χρόνου διάστημα προσηνέγκαντο οἱ γραμματεῖς τῷ βασιλεῖ, μηκέτι ἰσχύειν τὴν τῶν Ἰουδαίων ἀπογραφὴν ποιεῖσθαι διὰ τὴν ἀμέτρητον αὐτῶν πληθὺν, καί περ ὄντων κατὰ τὴν χώραν ἔτι τῶν πλειόνων,
\vs{18}τῶν μὲν κατὰ τὰς οἰκίας ἔτι συνεστηκότων, τῶν δὲ καὶ κατὰ τὸπον, ὡς ἀδυνάτου καθεστῶτος πᾶσι τοῖς ἐπʼ Αἴγυπτον στρατηγοῖς,
\vs{19}ἀπειλήσαντος δὲ αὐτοῖς σκληρότερον ὡς δεδωροκοπημένοις εἰς μηχανὴν τῆς ἐκφυγῆς, συνέβη σαφῶς αὐτὸν περὶ τούτου πεισθῆναι,
\vs{20}λεγόντων μετὰ ἀποδείξεως, καὶ τὴν χαρτηρίαν ἤδη καὶ τοὺς γραφικοὺς καλάμους ἐν οἷς ἐχρῶντο ἐκλελοιπέναι.
\vs{21}Τοῦτο δὲ ἦν ἐνέργεια τῆς τοῦ βοηθοῦντος τοῖς Ἰουδαίοις ἐξ οὐρανοῦ προνοίας ἀνικήτου.

\ch{5}
Τότε προσκαλεσάμενος Ἕρμωνα τὸν πρὸς τῇ τῶν ἐλεφάντων ἐπιμελείᾳ, βαρείᾳ μεμεστωμένος ὀργῇ καὶ χόλῳ κατὰ πᾶν ἀμετάθετος,
\vs{2}ἐκέλευσεν ὑπὸ τὴν ἐπερχομένην ἡμέραν δαψιλέσι δράκεσι λιβανωτοῦ καὶ οἴνῳ πλείονι ἀκράτῳ ἅπαντας τοὺς ἐλέφαντας ποτίσαι, ὄντας τὸν ἀριθμὸν πεντακοσίους, καὶ ἀγριωθέντας τῇ τοῦ πόματος ἀφθόνῳ χορηγίᾳ, εἰσαγαγεῖν πρὸς συνάντησιν τοῦ μόρου τῶν Ἰουδαίων.
\vs{3}Ὁ μὲν τάδε προστάσσων, ἐτρέπετο πρὸς τὴν εὐωχίαν, συναγαγὼν τοὺς μάλιστα τῶν φίλων καὶ τῆς στρατιᾶς ἀπεχθῶς ἔχοντας πρὸς τοὺς Ἰουδαίους.

\vs{4}Ὁ δὲ ἐλεφαντάρχης τὸ προσταγὲν ἀραρότως Ἕρμων συνετέλει.
\vs{5}Οἵ τε πρὸς τούτοις λειτουργοὶ κατὰ τὴν ἑσπέραν ἐξιόντες τὰς τῶν ταλαιπώρων ἐδέσμευον χεῖρας, τήν τε λοιπὴν ἐμηχανῶντο περὶ αὐτοὺς ἀσφάλειαν, ἔννυχον δόξαντες ὁμοῦ λήψεσθαι τὸ φῦλον πέρας τῆς ὀλεθρίας.

\vs{6}Οἱ δὲ πάσης σκέπης ἔρημοι δοκοῦντες εἶναι τοῖς ἔθνεσιν Ἰουδαῖοι, διὰ τὴν πάντοθεν περιέχουσαν αὐτοὺς μετὰ δεσμῶν ἀνάγκην,
\vs{7}τὸν παντοκράτορα Κύριον καὶ πάσης δυνάμεως δυναστεύοντα, ἐλεήμονα Θεὸν αὐτῶν καὶ πατέρα, δυσκαταπαύστῳ βοῇ πάντες μετὰ δακρύων ἐπεκαλέσαντο δεόμενοι,
\vs{8}τὴν κατʼ αὐτῶν μεταστρέψαι βουλὴν ἀνοσίαν, καὶ ῥύσασθαι αὐτοὺς μετὰ μεγαλομεροῦς ἐπιφανείας ἐκ τοῦ παρὰ πόδας ἐν ἑτοίμῳ μόρου.
\vs{9}Τούτων μὲν οὖν ἐκτενῶς ἡ λιτανεία ἀνέβαινεν εἰς οὐρανόν.

\vs{10}Ὁ δὲ Ἕρμων τοὺς ἀνηλεεῖς ἐλέφαντας ποτίσας πεπληρωμένους τῆς τοῦ οἴνου πολλῆς χορηγίας, καὶ τοῦ λιβάνου μεμεστωμένους, ὄρθριος ἐπὶ τὴν αὐλὴν παρῆν περὶ τούτων προσαγγεῖλαι τῷ βασιλεῖ.
\vs{11}Τοῦτο δʼ ἀπʼ αἰῶνος χρόνου κτίσμα καλὸν ἐν νυκτὶ καὶ ἡμέρᾳ ἐπιβαλλόμενον ὑπὸ τοῦ χαριζομένου πᾶσιν οἷς ἂν αὐτὸς θελήσῃ, ὕπνου μέρος ἀπέστειλε πρὸς τὸν βασιλέα.
\vs{12}Καὶ ἡδίστῳ καὶ βαθεῖ κατεσχέθη τῇ ἐνεργείᾳ τοῦ Δεσπότου, τῆς ἀθέσμου μὲν προθέσεως πολὺ διεσφαλμένος, τοῦ δὲ ἀμεταθέτου λογισμοῦ μεγάως διεψευσμένος.

\vs{13}Οἱ δὲ Ἰουδαῖοι τὴν προσημανθεῖσαν ὥραν διαφυγόντες, τὸν ἅγιον ᾔνουν Θεὸν αὐτῶν· καὶ πάλιν ἠξίουν τὸν εὐκατάλλακτον, δεῖξαι μεγαλοσθενοῦς αὐτοῦ χειρὸς κράτος ἔθνεσιν ὑπερηφάνοις.
\vs{14}Μεσούσης δὲ ἤδη τῆς δεκάτης ὥρας σχεδὸν, ὁ πρὸς ταῖς κλήσεσι τεταγμένος, ἀθρόους τοὺς κλητοὺς ἰδὼν, ἔνυξε προσελθὼν τὸν βασιλέα.
\vs{15}Καὶ μόλις διεγείρας, ὑπέδειξε τὸν τῆς συμποσίας καιρὸν ἤδη παρατρέχοντα, τὸν περὶ τούτων λόγον ποιοῦμενος.

\vs{16}Ὃν ὁ βασιλεὺς λογισάμενος, καὶ τραπεὶς εἰς τὸν πότον, ἐκέλευσε τοὺς παραγεγονότας εἰς τὴν συμποσίαν ἄντικρυς ἀνακλιθῆναι αὐτοῦ.
\vs{17}Οὗ καὶ γενομένου, παρῄνει εἰς εὐωχίαν δόντας ἑαυτοὺς, τὸ παρὸν τῆς συμποσίας ἐπιπολὺ γεραιρομένους εἰς εὐφροσύνην καταθέσθαι μέρος.
\vs{18}Ἐπιπλεῖον δὲ προβαινούσης τῆς ὁμιλίας, τὸν Ἕρμωνα μεταπεμψάμενος ὁ βασιλεὺς, μετὰ πικρὰς ἀπειλῆς ἐπυνθάνετο, τίνος ἕνεκεν αἰτίας εἰάθησαν οἱ Ἰουδαῖοι τὴν περοῦσαν ἡμέραν περιβεβιωκότες.
\vs{19}Τοῦ δὲ ὑποδείξαντος νυκτὸς τὸ προσταγὲν ἐπὶ τέλος ἠγηοχέναι, καὶ τῶν φίλων αὐτῷ προσμαρτυρησάντων,
\vs{20}τὴν ὠμότητα χείρονα Φαλάριδος ἐσχηκὼς ἔφη, τῷ τῆς σήμερον ὕπνῳ χάριν ἔχειν αὐτούς· ἀνυπερθέτως δὲ εἰς τὴν ἐπιτέλλουσαν ἡμέραν κατὰ τὸ ὅμοιον ἑτοίμασον τοὺς ἐλέφαντας ἐπὶ τὸν τῶν ἀθεμίτων Ἰουδαίων ἀφανισμόν.

\vs{21}Εἰπόντος δὲ τοῦ βασιλέως, ἀσμένως πάντες μετὰ χαρᾶς οἱ παρόντες ὁμοῦ συναινέσαντες, εἰς τὸν ἴδιον οἶκον ἕκαστος ἀνέλυσε.
\vs{22}Καὶ οὐχ οὕτως εἰς ὕπνον κατεχρήσαντο τὸν χρόνον τῆς νυκτὸς, ὡς εἰς τὸ παντοίους μηχανᾶσθαι τοῖς ταλαιπώροις δοκοῦσιν ἐμπαιγμούς.

\vs{23}Ἄρτι δὲ ἀλεκτρυὼν ἐκεκράγει ὄρθριος, καὶ τὰ θηρία καθωπλικὼς ὁ Ἕρμων ἐν τῷ μεγάλῳ περιστύλῳ διεκίνει.
\vs{24}Τὰ δὲ κατὰ τὴν πόλιν πλήθη συνήθροιστο πρὸς τὴν οἰκτροτάτην θεωρίαν, προσδοκῶντα τὴν πρωίαν μετὰ σπουδῆς.
\vs{25}Οἱ δὲ Ἰουδαῖοι κατὰ τὸν ἀμερῆ ψυχουλκούμενοι χρόνον, πολυδάκρυον ἱκετίαν ἐν μέλεσι γοεροῖς τείνοντες τὰς χεῖρας εἰς τὸν οὐρανὸν, ἐδέοντο τοῦ μεγίστου Θεοῦ, πάλιν αὐτοῖς βοηθῆσαι συντόμως.

\vs{26}Οὔπω δὲ ἡλίου βολαὶ κατεσπείροντο, καὶ τοῦ βασιλέως τοὺς φίλους ἐκδεχομένου, ὁ Ἕρμων παραστὰς, ἐκάλει πρὸς τὴν ἔξοδον, ὑποδεικνύων τὸ πρόθυμον τοῦ βασιλέως ἐν ἑτοίμῳ κεῖσθαι.
\vs{27}Τοῦ δὲ ἀποδεξαμένου καὶ καταπλαγέντος ἐπὶ τῇ παρανόμῳ ἐξόδῳ, κατὰ πᾶν ἀγνωσίᾳ κεκρατημένος ἐπυνθάνετο, τί τὸ πρᾶγμα ἐφʼ οὗ τοῦτο αὐτῳ μετὰ σπουδῆς τετέλεσται.
\vs{28}Τοῦτο δὲ ἦν ἡ ἐνέργεια τοῦ πάντα δεσποτεύοντος Θεοῦ, τῶν πρὶν αὐτῷ μεμηχανημένων λήθην κατὰ διάνοιαν ἐντεθεικότος.

\vs{29}Ὁ δὲ Ἕρμων ὑπεδείκνυε καὶ πάντες οἱ φίλοι, τὰ θηρία καὶ τὰς δυνάμεις ἡτοιμάσθαι, βασιλεῦ, κατὰ τὴν σὴν ἐκτενῆ πρόθεσιν.
\vs{30}Ὁ δὲ ἐπὶ τοῖς ῥηθεῖσι πληρωθεὶς βαρεῖ χόλῳ, διὰ τὸ περὶ τούτων προνοίᾳ Θεοῦ διεσκεδᾶσθαι πᾶν αὐτοῦ τὸ νόημα, ἐνατενίσας μετὰ ἀπειλῆς εἶπεν,
\vs{31}εἴ σαι γονεῖς παρῆσαν ἢ παίδων γοναὶ, τήνδε θηρσὶν ἀγρίοις ἐσκεύασαν δαψιλῆ θοῖναν, ἀντὶ τῶν ἀνεγκλήτων ἐμοὶ καὶ προγόνοις ἐμοῖς ἀποδεδειγμένων ὁλοσχερῆ βεβαίαν πίστιν ἐξόχως, Ἰουδαίων.
\vs{32}Καί περ εἰ μὴ διὰ τὴν τῆς συστροφίας στοργὴν καὶ τῆς χρείας, τὸ ζῇν ἀντὶ τούτων ἐστερήθης.

\vs{33}Οὕτως ὁ Ἕρμων ἀπροσδόκητον ἐπικίνδυνον ὑπήνεγκεν ἀπειλὴν, καὶ τῇ ὁράσει καὶ τῷ προσώπῳ συνεστάλη.
\vs{34}Ὁ καθεὶς δὲ τῶν φίλων σκυθρωπῶς ὑπεκρέων, τοὺς συνηθροισμένους ἀπέλυσαν ἕκαστον ἐπὶ τὴν ἰδίαν ἀσχολίαν.
\vs{35}Οἵ τε Ἰουδαῖοι τὰ παρὰ τοῦ βασιλέως ἀκούσαντες, τὸν ἐπιφανῆ Θεὸν καὶ βασιλέα τῶν βασιλέων ᾔνουν, καὶ τῆσδε τῆς βοηθείας αὐτοῦ τετευχότες.

\vs{36}Κατὰ δὲ τούτους τοὺς νόμους ὁ βασιλεὺς συστησάμενος πάλιν τὸ συμπόσιον, εἰς εὐφροσύνην τραπῆναι παρεκάλει.
\vs{37}Τὸν δὲ Ἕρμωνα προσκαλεσάμενος μετὰ ἀπειλῆς εἶπε, ποσάκις σοι δὲ περὶ τούτων αὐτῶν προστάττειν, ἀθλιώτατε;
\vs{38}Τοὺς ἐλέφαντας ἔτι καὶ νῦν καθόπλισον εἰς τὴν αὔριον ἐπὶ τὸν τῶν Ἰουδαίων ἀφανισμόν.

\vs{39}Οἱ δὲ συνανακείμενοι συγγενεῖς τὴν ἄστατον διάνοιαν αὐτοῦ θαυμάζοντες, προεφέροντο τάδε
\vs{40}βασιλεῦ, μέχρι τίνος ὡς ἀλόγους ἡμᾶς διαπειράζεις, προστάσσων ἤδη τρίτον αὐτοὺς ἀφανίσαι, καὶ πάλιν ἐπὶ τῶν πραγμάτων ἐκ μεταβολῆς ἀναλύων τὰ σοὶ δεδογμένα;
\vs{41}Ὧν χάριν ἡ πόλις διὰ τὴν προσδοκίαν ὀχλεῖ· καὶ πληθύουσα συστροφαῖς, ἤδη καὶ κινδυνεύει πολλάκις διαρπασθῆναι.

\vs{42}Ὅθεν ὁ κατὰ πάντα Φάλαρις βασιλεὺς ἐμπληθυνθεὶς ἀλογιστίας, καὶ τὰς γινομένας πρὸς ἐπισκοπὴν τῶν Ἰουδαίων ἐν αὐτῷ μεταβολὰς τῆς ψυχῆς παρʼ οὐδὲν ἡγούμενος, ἀτελέστατον ἐβεβαίωσεν ὅρκον, ὁρισάμενος τούτους μὲν ἀνυπερθέτως πέμψειν εἰς ᾅδην, ἐν γόνασι καὶ ποσὶ θηρίων ᾐκισμένους,
\vs{43}ἐπιστρατεύσαντα δὲ ἐπὶ τὴν Ἰουδαίαν, ἰσόπεδον πυρὶ καὶ δόρατι θήσεσθαι διατάχους, καὶ τὸν ἄβατον αὐτῶν ἡμῖν ναὸν πυρὶ πρηνέα ἐν τάχει, καὶ τῶν συντελούντων ἐκεῖ θυσίας ἔρημον τὸν ἅπαντα χρόνον καταστήσειν.

\vs{44}Τότε περιχαρεῖς ἀναλύσαντες οἱ φίλοι καὶ συγγενεῖς, μετὰ πίστεως διέτασσον τὰς δυνάμεις ἐπὶ τοὺς εὐκαιροτάτους τόπους τῆς πόλεως πρὸς τήρησιν.
\vs{45}Ὁ δὲ ἐλεφαντάρχης, τὰ θηρία σχεδὸν εἰπεῖν εἰς κατάστημα μανιῶδες ἀγηοχὼς, εὐωδεστάτοις πόμασιν οἴνῳ λελιβανωμένου φοβεραῖς κατεσκευασμένα σκευαῖς.

\vs{46}Περὶ τὴν ἕω, τῆς πόλεως ἤδη πλήθεσιν ἀναριθμήτοις κατὰ τοῦ ἱπποδρόμου καταμεμεστωμένης, εἰσελθὼν εἰς τὴν αὐλὴν, ἐπὶ τὸ προκείμενον ὤτρυνε τὸν βασιλέα.
\vs{47}Ὁ δὲ ὀργῇ βαρείᾳ γεμίσας δυσσεβῆ φρένα, παντὶ τῷ βάρει σὺν τοῖς θηρίοις ἐξώρμησε, βουλόμενος ἀτρώτῳ καρδίᾳ καὶ κόραις ὀφθαλμῶν θεάσασθαι τὴν ἐπίπονον καὶ ταλαίπωρον τῶν προσεσημαμμένων καταστροφήν.

\vs{48}Ὡς δὲ τῶν ἐλεφάντων ἐξιόντων περὶ πύλην, καὶ τῆς συνεπομένης ἐνόπλου δυνάμεως, τῆς τε τοῦ πλήθους πορείας κονιορτὸν ἰδόντες, καὶ βαρυηχῆ θόρυβον ἀκούσαντες οἱ Ἰουδαῖοι,
\vs{49}ὑστάτην βίου ῥοπὴν αὐτοῖς ἐκείνην δόξαντες εἶναι τὸ τέλος τῆς ἀθλιωτάτης προσδοκίας, εἰς οἶκτον καὶ γόους τραπέντες, κατεφίλουν ἀλλήλους περιπλεκόμενοι τοῖς συγγενέσιν ἐπὶ τοὺς τραχήλους ἐπιπίπτοντες, γονεῖς παισὶ καὶ μητέρες νεάνισιν, ἕτεραι δὲ νεογνὰ πρὸς μαστοὺς ἔχουσαι βρέφη τελευταῖον ἕλκοντα γάλα.

\vs{50}Οὐ μὴν δὲ ἀλλὰ καὶ τὰς ἔμπροσθεν αὐτῶν γεγενημένας ἀντιλήψεις ἐξ οὐρανοῦ συνιδόντες, πρηνεῖς ὁμοθυμαδὸν ῥίψαντες ἑαυτοὺς καὶ τὰ νήπια χωρίσαντες τῶν μαστῶν,
\vs{51}ἀνεβόησαν φωνῇ μεγάλῃ σφόδρα, τὸν τῆς ἁπάσης δυνάστην ἱκετεύοντες, οἰκτεῖραι μετὰ ἐπιφανείας αὐτοὺς ἤδη πρὸς πύλαις ᾅδου καθεστῶτας.

\ch{6}
Ἐλεαζάρος δέ τις ἀνὴρ ἐπίσημος τῶν ἀπὸ τῆς χώρας ἱερέων, ἐν πρεσβείῳ τὴν ἡλικίαν ἤδη λελογχὼς, καὶ πάσῃ τῇ κατὰ τὸν βίον ἀρετῇ κεκοσμημένος, τοὺς περὶ αὐτὸν καταστείλας πρεσβυτέρους ἐπικαλεῖσθαι τὸν ἅγιον Θεὸν προσηύξατο τάδε·

\vs{2}Βασιλεῦ μεγαλοκράτωρ, ὕψιστε, παντοκράτωρ Θεὲ, τὴν πᾶσαν διακυβερνῶν ἐν οἰκτιρμοῖς κτίσιν,
\vs{3}ἔπιδε ἐπὶ Ἀβραὰμ σπέρμα, ἐπὶ ἡγιασμένου τέκνα Ἰακὼβ, μερίδος ἡγιασμένης σου λαὸν ἐν ξένῃ γῇ ξένον ἀδίκως ἀπολλύμενον, πάτερ.

\vs{4}Σὺ Φαραὼ πληθύνοντα ἅρμασι, τὸ πρὶν Αἰγύπτου ταύτης δυνάστην, ἐπαρθέντα ἀνόμῳ θράσει καὶ γλώσσῃ μεγαλοῤῥήμονι, σὺν τῇ ὑπερηφάνῳ στρατιᾷ ποντοβρόχους ἀπώλεσας, φέγγος ἐπιφάνας ἐλέους Ἰσραὴλ γένει.
\vs{5}Σὺ τὸν ἀναριθμήτοις δυνάμεσι γαυρωθέντα Σενναχηρεὶμ βαρὺν Ἀσσυρίων βασιλέα, δόρατι τὴν πᾶσαν ὑποχείριον ἤδη λαβόντα γῆν, καὶ μετεωρισθέντα ἐπὶ τὴν ἁγίαν σου πόλιν, βαρέα λαλοῦντα κόμπῳ καὶ θράσει, Δέσποτα, ἔθραυσας, ἔκδηλον δεικνὺς ἐθνεσι πολλοῖς τὸ σὸν κράτος.
\vs{6}Σὺ τοὺς κατὰ τὴν Βαβυλωνίαν τρεῖς ἑταίρους πυρὶ τὴν ψυχὴν αὐθαιρέτως δεδωκότας εἰς τὸ μὴ λατρεῦσαι τοῖς κενοῖς, διάπυρον δροσίσας κάμινον, ἐῤῥύσω μέχρι τριχὸς ἀπημάντους, φλόγα πᾶσιν ἐπιπέμψας τοῖς ὑπεναντίοις.
\vs{7}Σὺ τὸν διαβολαῖς φθόνου λέουσι κατὰ γῆς ῥιφέντα θηρσὶν βορὰν Δανιὴλ εἰς φῶς ἀνήγαγες ἀσινῆ·
\vs{8}Τόν τε βυθοτρεφοῦς ἐν γαστρὶ κήτους Ἰωνᾶν τηκόμενον ἀφιδῶς, ἀπήμαντον πᾶσιν οἰκείοις ἀναδειξας, πάτερ.

\vs{9}Καὶ νῦν μισούβρι, πολυέλεε, τῶν ὅλων σκεπαστά, τὸ τάχος ἐπιφάνηθι τοῖς ἀπὸ Ἰσραὴλ γένους, ὑπὸ δὲ ἐβδελυγμένων ἀνόμων ἐθνῶν ὑβριζομένοις.
\vs{10}Εἰ δὲ ἀσεβείαις κατὰ τὴν ἀποικίαν ὁ βίος ἡμῶν ἐνέσχηται, ῥυσάμενος ἡμᾶς ἀπὸ ἐχθρῶν χειρός, ὡς προαιρῇ, Δέσποτα, ἀπόλεσον ἡμᾶς μόρῳ.

\vs{11}Μὴ τοῖς ματαίοις οἱ ματαιόφρονες εὐλογησάτωσαν ἐπὶ τῇ τῶν ἠγαπημένων σου ἀπωλείᾳ, λέγοντες, οὐδὲ ὁ Θεὸς αὐτῶν ἐῤῥύσατο αὐτούς.
\vs{12}Σὺ δὲ ὁ πᾶσαν ἀλκὴν καὶ δυναστείαν ἔχων ἅπασαν, αἰώνιε, νῦν ἔπιδε· ἐλέησον ἡμᾶς τοὺς καθʼ ὕβριν ἀνόμων ἀλόγιστον ἐκ τοῦ ζῇν μεθισταμένους ἐν ἐπιβούλων τρόπῳ.
\vs{13}Πτηξάτω δὲ ἔθνη σὴν δύναμιν ἀνίκητον σήμερον, ἔντιμε, δυνάμιν ἔχων, ἐπὶ σωτηρία Ἰακὼβ γένους.
\vs{14}Ἱκετεύει σε τὸ πᾶν πλῆθος τῶν νηπίων καὶ οἱ τούτων γονεῖς μετὰ δακρύων.
\vs{15}Δειχθήτω πᾶσιν ἔθνεσιν, ὅτι μεθʼ ἡμῶν εἶ Κύριε, καὶ οὐκ ἀπέστρεψας τὸ πρόσωπόν σου ἀφʼ ἡμῶν· ἀλλὰ καθὼς εἶπας, ὅτι οὐδʼ ἐν τῇ γῇ τῶν ἐχθρῶν αὐτῶν ὄντων ὑπέριδον αὐτούς, οὕτως ἐπιτέλεσον, Κύριε.

\vs{16}Τοῦ δὲ Ἐλεαζάρου λήγοντος ἄρτι τῆς προσευχῆς, ὁ βασιλεὺς σὺν τοῖς θηρίοις καὶ παντὶ τῷ τῆς δυνάμεως φρυάγματι κατὰ τὸν ἱππόδρομον παρῆγεν.
\vs{17}Καὶ θεωρήσαντες οἱ Ἰουδαῖοι, μέγα εἰς οὐρανὸν ἀνέκραξαν, ὥστε καὶ τοὺς παρακειμένους αὐλῶνας συνηχήσαντας, ἀκατάσχετον οἰμωγὴν ποιῆσαι παντὶ τῷ στρατοπέδῳ.

\vs{18}Τότε ὁ μεγαλόδοξος παντοκράτωρ καὶ ἀληθινὸς Θεός, ἐπιφάνας τὸ ἅγιον αὐτοῦ πρόσωπον, ἠνέῳξε τὰς οὐρανίους πύλας, ἐξ ὧν δεδοξασμένοι δύο φοβεροειδεῖς ἄγγελοι κατέβησαν φανεροὶ πᾶσι πλὴν τοῖς Ἰουδαίοις,
\vs{19}καὶ ἀντέστησαν, καὶ τὴν δύναμιν τῶν ὑπεναντίων ἐπλήρωσαν ταραχῆς καὶ δειλίας, καὶ ἀκινήτοις ἔδησαν πέδαις.
\vs{20}Καὶ ὑπόφρικον καὶ τὸ τοῦ βασιλέως σῶμα ἐγενήθη, καὶ λήθη τὸ θράσος αὐτοῦ τὸ βαρύθυμον ἔλαβε.
\vs{21}Καὶ ἀπέστρεψαν τὰ θηρία ἐπὶ τὰς συνεπομένας ἐνόπλους δυνάμεις, καὶ κατεπάτουν αὐτοὺς καὶ ὠλόθρευον.

\vs{22}Καὶ μετεστράφη τοῦ βασιλέως ἡ ὀργὴ εἰς οἶκτον καὶ δάκρυα ὑπὲρ τῶν ἔμπροσθεν αὐτῷ μεμηχανευμένων.
\vs{23}Ἀκούσας γὰρ τῆς κραυγῆς, καὶ συνιδὼν πρηνεῖς ἅπαντας εἰς τὴν ἀπώλειαν, δακρύσας μετὰ ὀργῆς τοῖς φίλοις διηπειλεῖτο, λέγων,
\vs{24}Παραβασιλεύετε, καὶ τυράννους ὑπερβεβήκατε ὠμότητι· καὶ ἐμὲ αὐτὸν τὸν ὑμῶν εὐεργέτην ἐπιχειρεῖτε τῆς ἀρχῆς ἤδη καὶ τοῦ πνεύματος μεθιστᾷν, λάθρα μηχανώμενοι τὰ μὴ συμφέροντα τῇ βασιλείᾳ.
\vs{25}Τίς τοὺς κρατήσαντας ἡμῶν ἐν πίστει τὰ τῆς χώρας ὀχυρώματα, τῆς οἰκίας ἀποστήσας ἕκαστον ἀλόγως ἤθροισεν ἐνθάδε;
\vs{26}Τίς τοὺς ἐξαρχῆς εὐνοίᾳ πρὸς ἡμᾶς κατὰ πάντα διαφέροντας πάντων ἐθνῶν, καὶ τοὺς χειρίστους πλεονάκις ἀνθρώπων ἐπιδεδεγμένους κινδύνους, οὕτως ἀθέσμοις περιέβαλεν αἰκίαις;

\vs{27}Λύσατε, ἐκλύσατε ἄδικα δεσμά· εἰς τὰ ἴδια μετʼ εἰρήνης ἐξαποστείλατε, τὰ προπεπραγμένα παραιτησάμενοι.
\vs{28}Ἀπολύσατε τοὺς υἱοὺς τοῦ παντοκράτορος ἐπουρανίου Θεοῦ ζῶντος, ὃς ἀφʼ ἡμετέρων μέχρι τοῦ νῦν προγόνων ἀπαραπόδιστον μετὰ δόξης εὐστάθειαν παρέχει τοῖς ἡμετέροις πράγμασιν.

\vs{29}Ὁ μὲν οὖν ταῦτα ἔλεξεν· οἱ δὲ ἐν ἀμερεῖ χρόνῳ λυθέντες, τὸν ἅγιον σωτῆρα Θεὸν αὐτῶν εὐλόγουν, ἄρτι τὸν θάνατον ἐκπεφευγότες.
\vs{30}Εἶτα ὁ βασιλεὺς εἰς τὴν πόλιν ἀπαλλαγεὶς, τὸν ἐπὶ τῶν προσόδων προσκαλεσάμενος, ἐκέλευσεν οἴνους τε καὶ τὰ λοιπὰ πρὸς εὐωχίαν ἐπιτήδεια τοῖς Ἰουδαίοις χορηγεῖν ἐπὶ ἡμέρας ἑπτὰ, κρίνας αὐτοὺς ἐν ᾧ τόπῳ ἔδοξαν τὸν ὄλεθρον ἀναλαμβάνειν, ἐν τούτῳ ἐν εὐφροσύνῃ πάσῃ σωτήρια ἄγειν.

\vs{31}Τότε οἱ πρὶν ἐπονείδιστοι καὶ πλησίον τοῦ ᾅδου, μᾶλλον δὲ ἐπʼ αὐτῷ βεβηκότες, ἀντὶ πικροῦ καὶ δυσαιάκτου μόρου, κώθωνα σωτήριον συστησάμενοι, τὸν εἰς πτῶσιν αὐτοῖς καὶ τάφον ἡτοιμασμένον τόπον κλισίαις κατεμέρισαν πλήρεις χαρμονῆς.
\vs{32}Καταλήξαντες δὲ θρήνου πανόδυρτον μέλος, ἀνέλαβον ᾠδὴν πάτριον, τὸν σωτῆρα καὶ τερατοποιὸν αἰνοῦντες Θεόν· οἰμωγήν τε πᾶσαν καὶ κωκυτὸν ἀπωσάμενοι, χοροὺς συνίσταντο εὐφροσύνης εἰρηνικῆς σημεῖον.

\vs{33}Ὡσαύτως δὲ καὶ ὁ βασιλεὺς περὶ τούτων συμπόσιον βαρὺ συναγαγὼν, ἀδιαλείπτως εἰς οὐρανὸν ἀνθωμολογεῖτο μεγαλομερῶς ἐπὶ τῇ παραδόξῳ γενηθείσῃ αὐτῷ σωτηρίᾳ.
\vs{34}Οἵ τε πρὶν εἰς ὄλεθρον καὶ οἰωνοβρώτους αὐτοὺς ἔσεσθαι τιθέμενοι, μετὰ χαρᾶς ἀπογραψάμενοι, κατεστέναξαν, αἰσχύνην ἐφʼ ἑαυτοῖς περιβαλλόμενοι, καὶ τὴν πυρίπνουν τόλμαν ἀκλεῶς ἐσβεσμένοι.

\vs{35}Οἵ τε Ἰουδαῖοι, καθὼς προειρήκαμεν, συστησάμενοι τὸν προειρημένον χορὸν, μετʼ εὐωχίας ἐν ἐξομολογήσεσιν ἱλαραῖς καὶ ψαλμοῖς διῆγον.
\vs{36}καὶ κοινὸν ὁρισάμενοι περὶ τούτων θεσμὸν ἐπὶ πᾶσαν τὴν παροικίαν αὐτῶν εἰς γενεὰς, τὰς προειρημένας ἡμέρας ἄγειν ἔστησαν εὐφροσύνους, οὐ πότου χάριν καὶ λιχνείας, σωτηρίας δὲ τῆς διὰ Θεὸν γενομένης αὐτοῖς.
\vs{37}Ἐνέτυχον δὲ τῷ βασιλεῖ, τὴν ἀπόλυσιν αὐτῶν εἰς τὰ ἴδια αἰτούμενοι.

\vs{38}Ἀπογράφονται δὲ αὐτοὺς ἀπὸ πέμπτης καὶ εἰκάδος τοῦ Παχὼν ἕως τῆς τετάρτης τοῦ Ἐπιφὶ, ἐπὶ ἡμέρας τεσσαράκοντα· συνίστανται δὲ αὐτῶν τὴν ἀπώλειαν ἀπὸ πέμπτης τοῦ Ἐπιφὶ ἕως ἑβδόμης, ἡμέραις τρισίν.
\vs{39}Ἐν αἷς καὶ μεγαλοδόξως ἐπιφάνας τὸ ἔλεος αὐτοῦ ὁ τῶν ὅλων δυνάστης, ἀπταίστους αὐτοὺς ἐῤῥύσατο ὁμοθυμαδόν.

\vs{40}Εὐωχοῦντο δὲ πάνθʼ ὑπὸ τοῦ βασιλέως χορηγούμενοι μέχρι τῆς τεσσαρεσκαιδεκάτης, ἐν ᾗ καὶ τὴν ἐντυχίαν ἐποιήσαντο περὶ τῆς ἀπολύσεως αὐτῶν.
\vs{41}Συναινέσας τε αὐτοὺς ὁ βασιλεὺς, ἔγραψεν αὐτοῖς τὴν ὑπογεγραμμένην ἐπιστολὴν πρὸς τοὺς κατὰ πόλιν στρατηγοὺς μεγαλοψύχως τὴν ἐκτενίαν ἔχουσαν.

\ch{7}
Βασιλεὺς Πτολεμαῖος ὁ Φιλοπάτωρ τοῖς κατʼ Αἴγυπτον στρατηγοῖς καὶ πᾶσι τοῖς τεταγμένοις ἐπὶ πραγμάτων, χαίρειν καὶ ἐῤῥῶσθαι.
\vs{2}Ἐῤῥώμεθα δὲ καὶ αὐτοὶ καὶ τὰ τέκνα ἡμῶν, κατευθύναντος ἡμῖν τοῦ μεγάλου Θεοῦ τὰ πράγματα καθὼς προαιρούμεθα.

\vs{3}Τῶν φίλων τινὲς κακοηθείᾳ πυκνότερον ἡμῖν παρακείμενοι, συνέπεισαν ἡμᾶς εἰς τὸ τοὺς ὑπὸ τὴν βασιλείαν Ἰουδαίους, συναθροίσαντας σύστημα, κολάσασθαι ξενιζούσαις ἀποστατῶν τιμωρίαις,
\vs{4}προσφερόμενοι μήποτε εὐσταθήσειν τὰ πράγματα ἡμῶν, διʼ ἣν ἔχουσιν οὗτοι πρὸς πάντα τὰ ἔθνη δυσμένειαν, μέχρις ἂν συντελεσθῇ τοῦτο.
\vs{5}Οἳ καὶ δεσμίους καταγαγόντες αὐτοὺς μετὰ σκυλμῶν ὡς ἀνδράποδα, μᾶλλον δὲ ὡς ἐπιβούλους, ἄνευ πάσης ἀνακρίσεως καὶ ἐξετάσεως ἐπεχείρησαν ἀνελεῖν, νόμου Σκυθῶν ἀγριωτέραν ἐμπεπορπημένοι ὠμότητα.

\vs{6}Ἡμεῖς δὲ ἐπὶ τούτοις σκληρότερον διαπειλησάμενοι, καθʼ ἣν ἔχομεν πρὸς ἅπαντας ἀνθρώπους ἐπιείκειαν, μόγις τὸ ζῇν αὐτοῖς χαρισάμενοι, καὶ τὸν ἐπουράνιον Θεὸν ἐγνωκότες ἀσφαλῶς ὑπερησπικότα τῶν Ἰουδαίων, ὡς πατέρα ὑπὲρ υἱῶν διαπαντὸς ὑπερμαχοῦντα,
\vs{7}τήν τε τοῦ φίλου ἣν ἔχουσι πρὸς ἡμᾶς βεβαίαν καὶ τοὺς προγόνους ἡμῶν εὔνοιαν ἀναλογισάμενοι, δικαίως ἀπολελύκαμεν πάσης καθʼ ὁντινοῦν αἰτίας τρόπον·
\vs{8}καὶ προστετάχαμεν ἑκάστῳ πάντας εἰς τὰ ἴδια ἐπιστρέφειν, ἐν παντὶ τόπῳ μηθενὸς αὐτοὺς τὸ σύνολον καταβλάπτοντος, μήτε ὀνειδίζειν περὶ τῶν γεγενημένων παρὰ λόγον.
\vs{9}Γινώσκετε γὰρ ὅτι κατὰ τούτων ἐάν τι κακοτεχνήσωμεν πονηρόν, ἢ ἐπιλυπήσωμεν αὐτοὺς τὸ σύνολον, οὐκ ἄνθρωπον, ἀλλὰ τὸν πάσης δεσπόζοντα δυνάμεως Θεὸν ὕψιστον ἀντικείμενον ἡμῖν ἐπʼ ἐκδικήσει τῶν πραγμάτων κατὰ πᾶν ἀφεύκτως δια παντὸς ἕξομεν· ἔῤῥωσθε.

\vs{10}Λαβόντες δὲ τὴν ἐπιστολὴν ταύτην, οὐκ ἐσπούδασαν εὐθέως γενέσθαι περὶ τὴν ἄφοδον, ἀλλὰ τὸν βασιλέα προσηξίωσαν τοὺς ἐκ τοῦ γένους τῶν Ἰουδαίων τὸν ἅγιον Θεὸν αὐθαιρέτως παραβεβηκότας καὶ τοῦ Θεοῦ τὸν νόμον, τυχεῖν διʼ αὐτῶν τῆς ὀφειλομένης κολάσεως,
\vs{11}προφερόμενοι τοὺς γαστρὸς ἕνεκεν τὰ θεῖα παραβεβηκότας προστάγματα, μηδέποτε εὐνοήσειν μηδὲ τοῖς τοῦ βασιλέως πράγμασιν.

\vs{12}Ὁ δὲ τἀληθὲς αὐτοὺς λέγειν παραδεξάμενος καὶ συναινέσας, ἔδωκεν αὐτοῖς ἄδειαν πάντων, ὅπως τοὺς παραβεβηκότας τοῦ Θεοῦ τὸν νόμον ἐξολοθρεύσωσι κατὰ πάντα τὸν ὑπὸ τὴν βασιλείαν αὐτοῦ τόπον μετὰ παῤῥησίας ἄνευ πάσης βασιλικῆς ἐξουσίας ἢ ἐπισκέψεως.
\vs{13}Τότε κατευφημήσαντες αὐτόν, ὡς πρέπον ἦν, οἱ τούτων ἱερεῖς, καὶ πᾶν τὸ πλῆθος ἐπιφωνήσαντες τὸ ἁλληλούϊα, μετὰ χαρᾶς ἀνέλυσαν.

\vs{14}Τότε τὸν ἐμπεσόντα τῶν μεμιασμένων ὁμοεθνῆ κατὰ τὴν ὁδὸν ἐκολάζοντο, καὶ μετὰ παραδειγματισμῶν ἀνῄρουν.
\vs{15}Ἐκείνῃ δὲ τῇ ἡμέρᾳ ἀνεῖλον ὑπὲρ τοὺς τριακοσίους ἄνδρας, καὶ ἤγαγον εὐφροσύνην μετὰ χαρᾶς τοὺς βεβήλους χειρωσάμενοι.
\vs{16}Αὐτοὶ δὲ οἱ μέχρι θανάτου τὸν Θεὸν ἐσχηκότες, παντελῆ σωτηρίας ἀπόλαυσιν εἰληφότες; ἀνέζευξαν ἐκ τῆς πόλεως παντοίοις εὐωδεστάτοις ἄνθεσι κατεστεμμένοι μετʼ εὐφροσύνης καὶ βοῆς, ἐν αἴνοις καὶ παμμέλεσιν ὕμνοις εὐχαριστοῦντες τῷ Θεῷ τῶν πατέρων αὐτῶν αἰωνίῳ σωτῆρι τοῦ Ἰσραήλ.

\vs{17}Παραγενηθέντες δὲ εἰς Πτολεμαΐδα τὴν ὀνομαζομένην διὰ τὴν τοῦ τόπου ἰδιότητα ῥοδοφόρον, ἐν ᾗ προσέμεινεν αὐτοὺς ὁ στόλος κατὰ κοινὴν αὐτῶν βουλὴν ἡμέρας ἑπτὰ,
\vs{18}ἐκεῖ ἐποίησαν πότον σωτήριον, τοῦ βασιλέως χορηγήσαντος αὐτοῖς εὐψύχως τὰ πρὸς τὴν ἄφιξιν πάντα ἑκάστῳ ἕως εἰς τὴν ἰδίαν οἰκίαν.
\vs{19}Καταχθέντες δὲ μετʼ εἰρήνης ἐν ταῖς πρεπούσαις ἐξομολογήσεσιν, ὡσαύτως κᾀκεῖ ἔστησαν καὶ ταύτας ἄγειν τὰς ἡμέρας ἐπὶ τὸν τῆς παροικίας αὐτῶν χρόνον εὐφροσύνους.
\vs{20}Ἃς καὶ ἀνιερώσαντες ἐν στήλῃ κατὰ τὸν τῆς συμποσίας τόπον προσευχῆς καθιδρύσαντες, ἀνέλυσαν ἀσινεῖς, ἐλεύθεροι, ὑπερχαρεῖς, διά τε γῆς καὶ θαλάσσης καὶ ποταμοῦ ἀνασωζόμενοι τῇ τοῦ βασιλέως ἐπιταγῇ, ἕκαστος εἰς τὴν ἰδίαν.

\vs{21}Καὶ πλείστην ἢ ἔμπροσθεν ἐν τοῖς ἐχθροῖς ἐξουσίαν ἐσχηκότες μετὰ δόκης καὶ φόβου, τὸ σύνολον ὑπὸ μηδενὸς διασεισθέντες τῶν ὑπαρχόντων.
\vs{22}Καὶ πάντα τὰ ἑαυτῶν πάντες ἐκομίσαντο ἐξ ἀπογραφῆς, ὥστε τοὺς ἔχοντάς τι, μετὰ φόβου μεγίστου ἀποδοῦναι αὐτοῖς, τὰ μεγαλεῖα τοῦ μεγίστου Θεοῦ ποιήσαντος τελείως ἐπὶ σωτηρίᾳ αὐτῶν.
\vs{23}Εὐλογητὸς ὁ ῥύστης Ἰσραὴλ εἰς τοὺς ἀεὶ χρόνους. Ἀμήν.


\def\book{ΜΑΚΚΑΒΑΙΩΝ Δʹ}
\biblebook{ΜΑΚΚΑΒΑΙΩΝ Δʹ}


\lettrine[lines=2, loversize=0.2, nindent=0em, findent=.25em]{\textcolor{bookheadingcolor}{Φ}}{ΙΛΟΣΟΦΩΤΑΤΟΝ} λόγον ἐπιδείκνυσθαι μέλλων, εἰ αὐτοδέσποτός ἐστιν τῶν παθῶν ὁ εὐσεβὴς λογισμός· συμβουλεύσαιμʼ ἂν ὑμῖν ὀρθῶς, ὅπως προθύμως προσέχητε τῇ φιλοσοφίᾳ.
\vs{2}Καὶ γὰρ ἀναγκαῖος εἰς ἐπιστήμην παντὶ ὁ λόγος, καὶ ἄλλως τῆς μεγίστης ἀρετῆς ἀρετῆς, λέγω δὴ φρονήσεως, περιέχει ἔπαινον·

\vs{3}Εἰ ἄρα τῶν σωφροσύνης κωλυτικῶν παθῶν ὁ λογισμὸς φαίνεται ἐπικρατεῖν, γαστριμαργίας τε καὶ ἐπιθυμίας·
\vs{4}ἀλλὰ καὶ τῶν τῆς δικαιοσύνης ἐμποδιστικῶν παθῶν κυριεύειν ἀναφαίνεται, οἷον κακοηθείας· καὶ τῶν τῆς ἀνδρείας ἐμποδιστικῶν παθῶν, θυμοῦ τε, καὶ πόνου καὶ φόβου.
\vs{5}Πῶς οὖν, ἴσως εἴποιεν ἄν τινες, εἰ τῶν παθῶν ὁ λογισμὸς κρατεῖ, λήθης καὶ ἀγνοίας οὐ δεσπόζει;
\vs{6}γελοῖον ἐπιχειροῦντες λέγειν· οὐ γὰρ τῶν ἑαυτοῦ παθῶν ὁ λογισμὸς κρατεῖ, ἀλλὰ τῶν τῆς δικαιοσύνης καὶ ἀνδρείας καὶ σωφροσύνης, καὶ φρονήσεως ἐναντίων· καὶ τούτων, οὐχ ὥστε αὐτὰ καταλῦσαι, ἀλλʼ ὥστε αὐτοῖς μὴ εἶξαι.

\vs{7}Πολλαχόθεν μὲν οὖν καὶ ἀλλαχόθεν ἔχοιμʼ ἂν ὑμῖν ἐπιδεῖξαι, ὅτι αὐτοκράτωρ ἐστὶν τῶν παθῶν ὁ εὐσεβὴς λογισμός.
\vs{8}Πολὺ δὲ πλέον τοῦτο ἀποδείξαιμι ἀπὸ τῆς ἀνδραγαθείας τῶν ὑπὲρ ἀρετὴν ἀποθανόντων, Ἐλεαζάρου τε καὶ ἑπτὰ ἀδελφῶν καὶ τῆς τούτων μητρός.
\vs{9}Ἅπαντες γὰρ οὗτοι τῶν ἕως θανάτου πόνων ὑπεριδοντες, ὑπεριδόντες ἐπεδείξαντο ὅτι περικρατεῖ τῶν παθῶν ὁ λογισμός.

\vs{10}Τῶν μὲν οὖν ἀρετῶν, ἔπεστί μοι ἐπαινεῖν τοὺς κατὰ τοῦτον τὸν καιρὸν ὑπὲρ τῆς καλοκᾳγαθίας ἀποθανόντας μετὰ τῆς μητρὸς ἄνδρας·
\vs{11}τῶν δὲ τιμῶν μακαρίσαιμʼ ἄν· θαυμασθέντες γὰρ ἐκεῖνοι οὐ μόνον ὑπὸ πάντων ἀνθρώπων ἐπὶ τῇ ἀνδρείᾳ καὶ τῇ ὑπομονῇ, ἀλλὰ καὶ ὑπὸ τῶν αἰκισαμένων, αἴτιοι κατέστησαν τοῦ καταλυθῆναι τὴν κατὰ τοῦ ἔθνους τυραννίδα, νικήσαντες τὸν τύραννον τῇ ὑπομονῇ, ὥστε διʼ αὐτῶν καθαρισθῆναι τὴν πατρίδα.

\vs{12}Ἀλλὰ καὶ περὶ τούτου νῦν αὐτίκα δὴ λέγειν ἐξέσται, ἀρξαμένων τῆς ὑποθέσεως, ὥσπερ εἴωθα ποιεῖν, καὶ οὕτως εἰς τὸν περὶ αὐτῶν τρέψομαι λόγον, δόξαν διδοὺς τῷ πανσόφῳ Θεῷ.

\vs{13}Ζητοῦμεν δὴ τοίνυν, εἰ αὐτοκράτωρ ἐστὶν παθῶν ὁ λογισμός.
\vs{14}Διακρίνωμεν δὲ, τί ποτέ ἐστιν λογισμός; καὶ τί πάθος; καὶ πόσαι παθῶν ἰδέαι; καὶ εἰ πάντων ἐπικρατεῖ τούτων ὁ λογισμός;

\vs{15}Λογισμὸς μὲν δὴ τοίνυν ἐστὶν νοῦς μετὰ ὀρθοῦς βίου· πρωτιμῶν τὸν σοφίας λόγον.
\vs{16}Σοφία δὴ τοὶνυν ἐστὶν γνῶσις θείων καὶ ἀνθρωπίνων πραγμάτων, καὶ τῶν τούτων αἰτίων.
\vs{17}Αὕτη δὴ τοίνυν ἐστὶν ἠ τοῦ νόμου παιδσεία· διʼ ἧς τὰ θεῖα σεμνῶς, καὶ τὰ ἀνθρώπινα συμφερόντως μανθάνομεν.

\vs{18}Τῆς δὲ σοφίας ἰδέαι καθεστᾶσιν, φρόνησις καὶ δικαιοσύνη καὶ ἀνδρεια καὶ σωφροσύνη.
\vs{19}Κυριωτάτη πάντων ἡ φρόνησις· ἐξ ἧς δὴ τῶν παθῶν ὁ λογισμὸς ἐπικρατεῖ.
\vs{20}Παθῶν δὲ φύσεις εἰσὶν αἱ περιεκτικώταται δύο, ἡδονή τε καὶ πόνος· τούτων δὲ ἐκάτερον καὶ περὶ τὴν ψυχὴν πέφυκεν.
\vs{21}Πολλαὶ δὲ καὶ περὶ τὴν ἡδονὴν καὶ τὸν πόνον παθῶν εἰσὶν ἀκολουθίαι.
\vs{22}Πρὸ μὲν οὖν τῆς ἡδονῆς ἐστιν ἐπιθυμία· μετὰ δὲ τὴν ἡδονὴν, χαρά.
\vs{23}Πρὸ δὲ τοῦ πόνου ἐστὶν φόβος· μετὰ δὲ τὸν πόνον, λύπη.

\vs{24}Θυμὸς δὲ κοινὸν πάθος ἐστὶν ἡδονῆς καὶ πόνου, ἐὰν ἐννοηθῇ τις ὃτε αὐτῷ περιέπεσεν.
\vs{25}Ἐν δὲ τῇ ἡδονῇ ἐστιν καὶ ἡ κακοήθης διάθεσις, πολυτροπωτάτη πάντων τῶν παθῶν οὖσα.
\vs{26}Κατὰ μὲν ψυχῆς ἀλαζονεία, καὶ φιλαργυρία, καὶ φιλοδοξία, καὶ φιλονεικία, ἀπιστία καὶ βασκανία·
\vs{27}κατὰ δὲ τὸ σῶμα, παντοφαγία, καὶ λαιμαργία, καὶ νομοφαγία.

\vs{28}Καθάπερ οὖν δυοῖν τοῦ σώματος καὶ τῆς ψυχῆς φυτῶν ὄντων ἡδονῆς τε καὶ πὸνου, πολλαὶ τούτων τῶν παθῶν εἰσιν παραφυάδες.
\vs{29}Ὧν ἕκαστος ὁ πανγέωργος λογισμὸς περικαθαίρων τε καὶ ἀποκνίζων, καὶ περιπλέκων, καὶ ἐπάρδων, καὶ πάντα τρόπον μεταχέων, ἐξημεροῖ τὰς τῶν ἠθῶν καὶ παθῶν ὕλας.
\vs{30}Ὁ γὰρ λογισμὸς τῶν μὲν ἀρετῶν ἐστιν ἡγεμῶν, τῶν δὲ παθῶν αὐτοκράτωρ. Ἐπιθεώρει γε τοίνυν πρῶτον διʼ αὐτῶν κωλυτικῶν τῆς σωφροσύνης ἔργων, ὅτι αὐτοδέσποτός ἐστιν τῶν παθῶν ὁ λογισμός.

\vs{31}Σωφροσύνη δὴ τοίνυν ἐστὶν ἐπικράτεια τῶν ἐπιθυμιῶν.
\vs{32}Τῶν δὲ ἐπιθυμιῶν αἱ μέν εἰσιν ψυχικαὶ, αἱ δὲ σωματικαί· καὶ τούτων ἀμφοτέρων ὁ λογισμὸς ἐπικρατεῖν φαίνεται.
\vs{33}Ἐπεὶ πόθεν κινούμενοι πρὸς τὰς ἀπειρημένας τοοφὰς, ἀποτρεπόμεθα τὰς ἐξ ἑαυτῶν ἡδονάς; οὐχ ὅτι δύναται τῶν ὀρὲξεων ἐπικρατεῖν ὁ λογισμός; ἐγὼ μὲν οἶμαι.
\vs{34}Τοιγαροῦν ἐνύδρων ἐπιθυμοῦντες καὶ ὀρνέων καὶ τετραπόδων, παντοίων βρωμάτων τὼν ἀπηγορευμένων ἡμῖν κατὰ τὸν νόμον ἀπεχόμεθα διὰ τὴν τοῦ λογισμοῦ ἐπικράτειαν.
\vs{35}Ἀντέχεται γὰρ τὰ τῶν ὀρέξεων πάθη ὑπὸ τοῦ σώφρονος νοὸς ἀνακαμπτόμενα· καὶ φιλοτιμοῦνται πάντα τὰ τοῦ σώματος κινήματα ὑπὸ τοῦ λογισμοῦ.

\ch{2}
Καὶ τὶ θαυμαστὸν; εἰ αἱ τῆς ψυχῆς ἐπιθυμίαι πρὸς τὴν τοῦ κάλλους μετουσίαν ἀκυροῦνται.
\vs{2}Ταύτῃ γοῦν ὁ σώφρων Ἰωσὴφ ἐπαινεῖται, ὅτι τῷ λογισμῷ, διανοίᾳ περιεκράτησεν τῆς ἠδυπαθείας.
\vs{3}Νέος γὰρ ὢν καὶ ἀκμάζων πρὸς συνουσιασμὸν ἠκύρωσεν τῷ λογισμῷ τὸν τῶν παθῶν οἶστρον.

\vs{4}Οὐ μόνον δὲ τὴν τῆς ἡδυπαθείας οἰστρηλασίαν ἐπικρατεῖν ὁ λογισμὸς φαίνεται, ἀλλὰ καὶ πάσης ἐπιθυμίας.
\vs{5}Λέγει γοῦν ὁ νόμος· οὐκ ἐπιθυμήσεις τὴν γυναῖκα τοῦ πλησίον σου, οὐδὲ ὅσα τῷ πλησίον σου ἐστίν.
\vs{6}Καίτοι ὅτε μὴ ἐπιθυμεῖν εἴρηκεν ἡμᾶς ὁ νόμος, πολὺ πλέον πείσαιμʼ ἂν ὑμᾶς, ὅτι τῶν ἐπιθυμιῶν κρατεῖν δύναται ὁ λογισμὸς, ὥσπερ καὶ τῶν κωλυτικῶν τῆς δικαιοσύνης παθῶν.
\vs{7}Ἐπεὶ τίνα τρόπον μονοφάγος τις ὢν τὸ ἦθος, καὶ γαστρίμαργος, καὶ μέθυσος, μεταπαιδεύεται, εἰ μὴ δῆλον, ὅτι κύριός ἐστιν τῶν παθῶν ὁ λογισμός;

\vs{8}Αὐτίκα γοῦν τῷ νόμῳ πολιτευόμενος, κᾂν φιλάργυρός τις εἴη, βιάζεται τὸν ἑαυτοῦ τρόπον, τοῖς δεομένοις δανείζων χωρὶς τόκων, καὶ τὸ δάνειον τῶν ἑβδομάδων ἐντάσσων χρεοκοπούμενος.
\vs{9}Κᾂν φειδωλός τις ᾖ, ὑπὸ τοῦ νόμου κρατεῖται διὰ τὸν λογισμὸν, μήτε ἐπικαρπούμενος τοὺς ἀμητοὺς, μήτε ἐπιῤῥωγολογούμενος τοὺς ἀμπελῶνας, καὶ ἐπὶ τῶν ἐτέρων ἔστιν ἐπιγνῶναι τοῦτο, ὅτι τῶν παθῶν ἐστιν ὁ λογισμὸς κρατῶν.

\vs{10}Ὁ γὰρ νόμος καὶ τῆς πρὸς γονεῖς εὐνοίας κρατεῖ, μὴ καταπροδιδοὺς τὴν ἀρετὴν διʼ αὐτούς·
\vs{11}καὶ τῆς προσγαμετῆς φιλίας ἐπικρατεῖ, διὰ παρανομίαν αὐτὴν ἀπελέγχων.
\vs{12}Καὶ τῆς τέκνων φιλίας κυριεύει, διὰ κακίαν αὐτῶν κολάζων, καὶ τῆς φίλων συνηθείας δεσπόζει, διὰ πονηρίας αὐτοὺς ἐξελέγχων.
\vs{13}Καὶ μὴ νομίσητε παράδοξον εἶναι, ὅπου καὶ ἔχθραν ὁ λογισμὸς ἐπινομίσητε παράδοξον εἶναι, ὅπου καὶ ἔχθραν ὁ λογισμὸς ἐπινομισητε παράδοξον εἶναι, ὅπου καὶ ἔχθραν ὁ λογισμὸς ἐπικρατεῖν δύναται διὰ τὸν νόμον,
\vs{14}μηδὲ δενδροτομῶν τὰ ἥμερα τῶν πολεμίων φυτὰ, τὰ δὲ τῶν ἐχθρῶν τοῖς ἀπολέσασιν διασώζων, καὶ τὰ πεπτωκότα συνεγείρων.

\vs{15}Καὶ τῶν βιοτέρων δὲ παθῶν κρατεῖν ὁ λογισμὸς φαίνεται, φιλαρχίας, καὶ κενοδοξίας, καὶ ἀλαζονείας, καὶ μεγαλαυχίας, καὶ βασκανίας.
\vs{16}Πάντα γὰρ ταῦτα τὰ κακοήθη πάθη ὁ σώφρων νοῦς ἀπωθεῖται, ὥσπερ καὶ τὸν θυμόν· καὶ γὰρ τοῦτο δεσπόζει.

\vs{17}Θυμούμενος γέ τοι Μωσῆς κατὰ Δαθὰν καὶ Ἀβειρῶν, οὐ θυμῷ τι κατʼ αὐτῶν ἐποίησεν, ἀλλὰ λογισμῷ τὸν θυμὸν διῄτησεν.
\vs{18}Δυνατὸς γὰρ ὁ σώφρων νοῦς, ὡς ἔφην, κατὰ τῶν παθῶι ἀριστεῦσαι, καὶ τὰ μὲν αὐτῶν μεταθεῖναι, τὰ δὲ καὶ ἀκυρῶσαι.
\vs{19}Ἐπεὶ διατί ὁ πάνσοφος ἡμῶν πατὴρ Ἰακὼβ τοὺς περὶ Συμεὼν καὶ Λευὶν αἰτιᾶται, μὴ λογισμῷ τοὺς Σικιμίτας ἐθνηδὸν ἀποσφάξαντας, λέγων, ἐπικατάρατος ὁ θυμὸς αὐτῶν;
\vs{20}Εἰ μὴ γὰρ ἐδύνετο τῶν θυμῶν ὁ λογισμὸς κρατεῖν, οὐκ ἂν εἶπεν οὑτως.

\vs{21}Ὁπηνίκα γὰρ ὁ Θεὸς τὸν ἄνθρωπον κατεσκεύαζεν, τὰ πάθη αὐτοῦ καὶ τὰ ἤθη περιεφύτευσεν.
\vs{22}Καὶ τηνικαῦτα δὲ περὶ πάντων τὸν ἱερὸν ἡγεμόνα νοῦν διὰ τῶν αἰσθητηρίων ἐνεθρόνισεν·
\vs{23}καὶ τούτῳ νόμον ἔδωκεν, καθʼ ὃν πολιτευόμενος βασιλεύσει βασιλείαν σώφρονά τε, καὶ δικαίαν, καὶ ἀγαθὴν, καὶ ἀνδρείαν.
\vs{24}Πῶς οὖν, εἴποι τις ἂν, εἰ τῶν παθῶν ὁ λογισμὸς κρατεῖ, λήθης καὶ ἀγνοίας οὐ κρατεῖ;

\ch{3}
Ἐστὶ δὲ κομιδῆ γελοῖος ὁ λογισμός οὐ γὰρ τῶν ἑαυτοῦ παθῶν ὁ λογισμὸς ἐπικρατεῖν φαίνεται, ἀλλὰ τῶν σωματικῶν.
\vs{2}Οἷον ἐπιθυμίαν τις ὑμῶν οὐ δύναται ἐκκόψαι, ἀλλὰ μὴ δουλωθῆναι τῇ ἐπιθυμίᾳ δύναται ὁ λογισμὸς παρασχέσθαι.

\vs{3}Θυμόν τις οὐ δύναται ἐκκόψαι ἡμῶν τῆς ψυχῆς, ἀλλὰ τῷ θυμῷ δυνατὸν βοηθῆσαι.
\vs{4}Κακοήθειάν τις ὑμῶν οὐ δύναται ἐκκόψαι, ἀλλὰ τὸ μὴ καμφθῆναι τῇ κακοηθείᾀ δυνατὸν ὁ λογισμὸς συμμαχῆσαι.
\vs{5}Οὐ γὰρ ἐκριζωτὴς τῶν παθῶν ὁ λογισμός ἐστιν, ἀλλʼ ἀνταγωνιστής.
\vs{6}Ἔστιν γοῦν τοῦτο διὰ τῆς Δαυεὶδ τοῦ βασιλέως δίψης σαφέστερον ἐπιλογίσασθαι.
\vs{7}Ἐπεὶ γὰρ διʼ ὅλης ἡμέρας προσβαλὼν τοῖς ἀλλοφύλοις ὁ Δαυὶδ, πολλοὺς αὐτῶν ἀπέκτεινεν μετὰ τῶν τοῦ ἔθνους στρατιωτῶν·
\vs{8}τότε δὲ γενομένης ἑσπέρας, ὑδρῶν καὶ σφόδρα κεκμηκὼς, ἐπὶ τὴν βασίλειον σκηνὴν ἦλθεν, περὶ ἣν ὁ πᾶς τῶν προγόνων στρατὸς ἐστρατοπέδευκεν.

\vs{9}Οἱ μὲν οὖν ἄλλοι πάντες ἐπὶ τὸ δεῖπνον ἦσαν.
\vs{10}Ὁ δὲ βασιλεὺς ὡς μάλιστα διψῶν, καίπερ ἀφθόνους ἔχων πηγὰς, οὐκ ἠδύνατο διʼ αὐτῶν ἰάσασθαι τὴν δίψαν·
\vs{11}ἀλλά τις αὐτὸν ἀλόγιστος ἐπιθυμία τοῦ παρὰ τοῖς πολεμίοις ὕδατος ἐπιτείνουσα συνέφρυγεν, καὶ λύουσα κατέφλεγεν.

\vs{12}Ὅθεν τῶν ὑπερασπιστῶν ἐπὶ τῇ τοῦ βασιλέως ἐπιθυμία σχετλιαζόντων, δύο νεανίσκοι στρατιῶται καρτεροὶ καταιδεσθέντες τὴν τοῦ βασιλέως ἐπιθυμίαν, τὰς πανοπλίας καθωπλίσαντο, καὶ κάλπην λαβόντες ὑπερέβησαν τοὺς τῶν πολεμίων χάρακας·
\vs{13}καὶ λαθόντες τοὺς τῶν πυλῶν ἀκροφύλακας, διεξῄεσαν εὑράμενοι κατὰ πᾶν τὸ τῶν πολεμίων στρατόπεδον.
\vs{14}Καὶ ἀνευράμενοι θαῤῥαλέως τὴν πηγὴν, ἐξ αὐτῆς ἐγέμισαν τῷ βασιλεῖ τὸ ποτόν.

\vs{15}Ὁ δὲ καὶ περὶ τὴν δίψαν διαπυρούμενος, ἐλογίσατο πάνδεινον εἶναι κίνδυνον τῇ ψυχῇ λογισθὲν ἰσοδύναμον τὸ ποτὸν αἵματι.
\vs{16}Ὅθεν ἀντιθεὶς τῇ ἐπιθυμίᾳ τὸν λογισμὸν, ἔσπεισεν τὸ πόμα τῷ Θεῷ.
\vs{17}Δυνατὸς γὰρ ὁ σώφρων νοῦς νικῆσαι τὰς τῶν παθῶν ἀνάγκας, καὶ σβέσαι τὰς τῶν οἴστρων φλεγμονὰς, καὶ τὰς τῶν σωμάτων ἀλγηδόνας καθʼ ὑπερβολὴν οὔσας καταπαλαῖσαι,
\vs{18}καὶ τῆς καλοκᾳγαθίας τοῦ λογισμοῦ ἀποπτῦσαι πάσας τὰς τῶν παθῶν ἐπικρατείας.

\vs{19}Ἤδη δὲ καὶ ὁ καιρὸς ἡμᾶς καλεῖ ἐπὶ τὴν ἀπόδειξιν τῆς ἱστορίας τοῦ σώφρονος λογισμοῦ.
\vs{20}Ἐπειδὴ γὰρ βαθεῖαν εἰρήνην διὰ τὴν εὐνομίαν οἱ πατέρες ἡμῶν εἶχον, καὶ ἔπραττον καλῶς, ὥστε καὶ τὸν τῆς Ἀσίας βασιλέα Σέλευκον τὸν Νικάνορα καὶ χρήματα εἰς τὴν ἱερουργίαν αὐτοῖς ἀποφορίσαι, καὶ τὴν πολιτείαν αὐτῶν ἀποδέχεσθαι·
\vs{21}τότε δή τινες πρὸς τὴν κοινὴν νεωτερίσαντες ὁμόνοιαν, πυλυτρόπως ἐχρήσαντο συμφοραῖς.

\ch{4}
Σίμων γάρ τις πρὸς Ὀνίαν ἀντιπολιτεύομενος τόν ποτε τὴν ἀρχιερωσύνην ἔχοντα διὰ βίου, καλὸν καὶ ἀγαθὸν ἄνδρα, ἐπειδὴ πάντα τρόπον διαβάλλων ὑπὲρ τοῦ ἔθνους οὐκ ἰσχυσεν κακῶσαι, φυγὰς ᾤχετο, τὴν πατρίδα προδώσων.

\vs{2}Ὅθεν ἥκων πρὸς Ἀπολλώνιος, τὸν Συρίας τε καὶ Φοινίκης καὶ Κιλικίας στρατηγὸν, ἔλεγεν,
\vs{3}εὔνους ὢν τοῖς τοῦ βασιλέως πράγμασιν ἥκω, μηνύων πολλὰς ἰδιωτικῶν χρημάτων μυριάδας ἐν τοῖς Ἱεροσολύμων γαζοφυλακίοις τεθησαύρισται, τῷ ἱερῷ μὴ ἐπικοινωνούσας, ἀλλὰ προσήκειν ταῦτα Σελεύκῳ τῷ βασιλεῖ.

\vs{4}Τούτων ἕκαστα γνοὺς ὁ Ἀπολλώνιος, τὸν μὲν Σίμωνα τῆς εἰς τὸν βασιλέα κηδεμονίας ἐπαινεῖ, πρὸς δὲ τὸν Σέλευκον ἀναβὰς κατεμήνυε τὸν τῶν χρημάτων θησαυρόν·
\vs{5}καὶ λαβὼν τὴν περὶ αὐτῶν ἐξουσίαν, ταχὺ εἰς τὴν πατρίδα ἡμῶν μετὰ τοῦ καταράτου Σίμωνος καὶ βαρυτάτου στρατοῦ προσελθὼν,
\vs{6}ταῖς τοῦ βασιλέως ἐντολαῖς ἥκειν ἔλεγεν, ὅπως τὰ ἰδιωτικὰ τοῦ γαζοφυλακίου λάβοι χρήματα.
\vs{7}Καὶ τοῦ ἔθνους πρὸς τὸν λόγον σχετλιάζοντος, ἀντιλέγοντός τε, πάνδεινον εἶναι νομίσαντες, εἰ οἱ τὰς παρακαταθήκας πιστεύσαντας τῷ ἱερῷ θησαυρῷ στερηθήσονται, ὡς οἷόν τε ἦν ἐκώλυον.
\vs{8}Μετὰ ἀπειλῆς δὲ ὁ Ἀπολλώνιος ἀπῄει εἰς τὸ ἱερόν.

\vs{9}Τῶν δὲ ἱερέων μετὰ γυναικῶν καὶ παιδίων ἐν τῷ ἱερῷ ἱκετευσάντων τὸν Θεὸν ὑπερασπίσαι τοῦ ἱεροῦ καταφρονουμένου τόπου.
\vs{10}Ἀνιόντος τε μετὰ καθωπλισμένης τῆς στρατιᾶς τοῦ Ἀπολλωνίου πρὸς τὴν τῶν χρημάτων ἀρπαγὴν οὐρανόθεν ἔφιπποι προϋφάνησαν ἄγγελοι περιαστράπτοντες τοῖς ὅπλοις, καὶ πολὺν αὐτοῖς φόβον τε καὶ τρόμον ἐνιόντες.
\vs{11}Καταπεσὼν γέ τοι ἡμιθανὴς ὁ Ἀπολλώνιος ἐπὶ τὸν πάμφυλον τοῦ ἱεροῦ περίβολον, τὰς χεῖρας ἐξέτεινεν εἰς τὸν οὐρανὸν, μετὰ διακρύων τοὺς Ἑβραίους παρεκάλει, ὅπως περὶ αὐτοῦ εὐξόμενοι, τὸν ἐπουράνιον ἐξευμενίσωνται στρατόν.
\vs{12}Ἔλεγεν γὰρ ἡμαρτηκὼς, ὥστε καὶ ἀποθανεῖν ἄξιος ὑπάρχειν, πᾶσίν τε ἀνθρώποις ὑμνήσειν σωθεὶς τὴν τοῦ ἱεροῦ τόπου μακαριότητα.

\vs{13}Τούτοις ἐπαχθεὶς τοῖς λόγοις Ὀνίας ὁ ἀρχιερεὺς, καίπερ ἄλλως εὐλαβηθεὶς, μή ποτε νομίσειεν ὁ βασιλεὺς Σέλευκος ἐξ ἀνθρωπίνης ἐπιβουλῆς καὶ μὴ θείας δίκης ἀνῃρήσασθαι τὸν Ἀπωλλώνιον, ηὔξατο περὶ αὐτοῦ.
\vs{14}Καὶ ὁ μὲν παραδὸξως διασωθεὶς ᾤχετο, δηλώσων τῷ βασιλεῖ τὰ συμβάντα αὐτῷ.

\vs{15}Τελευτήσαντος δὲ Σελεύκου τοῦ βασιλέως διαδέχεται τὴν ἀρχὴν ὁ υἱὸς αὐτοῦ Ἀντίοχος Ἐπιφανὴς, ἀνὴρ ὑπερήφανος καὶ δεινὸς.
\vs{16}Ὃς καταλύσας τὸν Ὀνίαν τῆς ἀρχιερωσύνης,
\vs{17}Ἰάσονα τὸν ἀδελφὸν αὐτοῦ κατέστησεν ἀρχιερέα, συνθέμενον δώσειν, εἰ ἐπιτρέψειεν αὐτῷ τὴν ἀρχὴν, κατʼ ἐνιαυτὸν τρισχίλια ἐξακόσια ἑξήκοντα τάλαντα.

\vs{18}Ὁ δὲ ἐπέτρεψεν αὐτῷ ἀρχιερᾶσθαι καὶ τοῦ ἔθνους ἀφηγεῖσθαι.
\vs{19}Ὃς καὶ ἐξεζήτησεν τὸ ἔθνος, καὶ ἐξεπολίτευσεν ἐπὶ πᾶσαν παρανομίαν.
\vs{20}Ὥστε μὴ μόνον ἐπʼ αὐτῇ τῆ ἄκρᾳ τῆς πατρίδος ἡμῶν γυμνάσιον κατασκευάσαι, τὴν τοῦ ἱεροῦ κηδεμονίαν.
\vs{21}Ἐφʼ οἷς ἀγανακτήσασα ἡ θεία δίκη αὐτόν τοι τὸν Ἀντίοχον ἐπολέμησεν.
\vs{22}Ἐπειδὴ γὰρ πολεμῶν ἦν κατʼ Αἴγυπ τον Πτολεμαίῳ, ἤκουσέν τε, ὅτι φήμης διαδοθείσης περὶ τιῦ τεθνάναι αὐτὸν, ὡς ἔνι μάλιστα χαίροιεν οἱ Ἱεροσολυμῖται, ταχέως ἐπʼ αὐτοὺς ἀνέζευξεν.
\vs{23}Καὶ ὡς ἐπόρθσεν αὐτοὺς, δόγμα ἔθετο, ὅπως εἴ τινες αὐτῶν φάνοιεν τῷ πατρίῳ πολιτευόμενοι νόμῳ θάνοιεν.

\vs{24}Καὶ ἐπεὶ κατὰ μηδένα τρόπον ἴσχυεν καταλῦσαι διὰ τῶν δογμάτων τὴν τοῦ ἔθνους εὔνοιαν, ἀλλὰ πάσας τὰς ἑαυτοῦ ἀπειλὰς καὶ τιμωρίας ἑώρα καταλυομένας,
\vs{25}ὥστε καὶ γυναῖκας, ὅτι περιέτεμον τὰ παιδία, μετὰ τῶν βρεφῶν κατακρημνισθῆναι, προειδυίας ὅτι τοῦτο πείσονται·
\vs{26}ἐπεὶ οὖν τὰ δόγματα αὑτοῦ κατεφρονεῖτο ὑπὸ τοῦ λαοῦ, αὐτὸς διὰ βασάνων ἕνα ἕκαστον τούτου ἔθνους ἠνάγκαζεν μικρῶν ἀπογευομένους τροφῶν, ἐξόμνυσθαι τὸν Ἰουδαϊσμόν.

\ch{5}
Προκαθίσας γέ τοι μετὰ τῶν συνέδρων ὁ τύραννος Ἀντίοχος ἐπί τινος ὑψηλοῦ τόπου,
\vs{2}καὶ τῶν στρατευμάτων αὐτῶν ἐνόπλων κυκλόθεν παρεστηκότων παρεκέλευεν τοῖς δορυφόροις ἕνα ἕκαστον τῶν Ἑβραίων περισπᾶσθαι καὶ κρεῶν ὑείων καὶ εἰδωλοθύτων ἀναγκάζειν ἀπογεύεσθαι.
\vs{3}Εἰ δὲ τινες μὴ θέλοιεν μιαροφαγῆσαι, τούτους τροχισθέντας ἀναιρεθῆναι.

\vs{4}Πολλῶν δὲ συναρπασθέντων, εἷς πρῶτος ἐκ τῆς ἀγέλης Ἑβραῖος ὀνόματι Ἐλεάζαρος, τὸ γένος ἱερεὺς, τὴν ἐπιστήμην νομικὸς, καὶ τὴν ἡλικίαν προήκων, καὶ πολλοῖς τῶν περὶ τὸν τύραννον διὰ τὴν ἡλικίαν γνώριμος, παρήχθη πλησίον αὐτοῦ.

\vs{5}Καὶ αὐτὸν ἰδὼν ὁ Ἀντίοχος, ἔφη,
\vs{6}ἐγὼ πρὶν ἂρξασθαι τῶν κατὰ σοῦ βασάνων, ὦ πρεσβύτα, συμβουλεύσαιμʼ ἄν σοι ταῦτα ὅπως ἀπογευσάμενος τῶν ὑείων σώζοιο· αἰδοῦμαι γάρ σου τὴν ἡλικίαν καὶ τὴν πολιὰν, ἥν μετὰ τοσοῦτον ἔχων χρόνον, οὔ μοι δοκεῖς φιλοσοφεῖν, τῇ Ἰουδαιων χρώμενος θρησκείᾳ.
\vs{7}Διατί γὰρ τῆς φύσεως κεχαρισμένης καλλίστην τὴν τοῦδε τοῦ ζώου σαρκοφαγίαν βδελύττῃ;
\vs{8}Καὶ γὰρ ἀνόητον τοῦτο τὸ μὴ ἀπολαύειν τῶν χωρὶς ὀνείδους ἡδέων, καὶ διʼ ἄδικον ἀποστρέφεσθαι τὰς τῆς φύσεως χάριτας.

\vs{9}Σὺ δέ μοι καὶ ἀνοητότερον ποιήσειν δοκεῖς, εἰ κενοδοξῶν περὶ τὸ ἀληθὲς,
\vs{10}ἔτι κᾀμοῦ καταφρονήσεις ἐπὶ τῇ ἰδίᾳ τιμωρίᾳ· οὐκ ἐξυπνώσεις ἀπὸ τῆς φλυάρου φιλοσοφίας ὑμῶν;
\vs{11}Καὶ ἀποσκεδάσεις ψῶν λογισμῶν σου τὸν λῆρον, καὶ ἄξιον τῆς ἡλικίας ἀναλαβῶν νοῦν φιλοσοφήσεις τήν τοῦ συμφέροντος ἀλήθειαν;
\vs{12}καὶ προσκυνήσας μου τὴν φιλάνθρωπον παρηγορίαν οἰκτειρήσεις τὸ σεαυτοῦ γῆρας;
\vs{13}καὶ γὰρ ἐνθυμήθητι, ὡς εἰ καί τις ἐστιν τῆσδε τῆς θρησκείας ἐποπτικὴ δύναμις, συγνωμονήσειν σοι ἐπὶ πᾶσιν διʼ ἀνάγκην παρανομίᾳ γεινομένῃ.

\vs{14}Τοῦτον τὸν τρόπον ἐπὶ τὴν ἔκθεσμον σαρκοφαγίαν ἐποτρύνοντος τοῦ τυράννου, λόγον ᾔτησεν ὁ Ἐλεάζαρος.
\vs{15}Καὶ λαβὼν τοῦ λέγειν ἐξουσίαν, ἤρξατο δημηγορεῖν οὕτως·
\vs{16}ἡμεῖς, Αντίοχε, θείῳ πεπεισμένοι νόμῳ πολιτευεσθαι, οὐδεμίαν ἀνάγκην βιαιοτέραν εἶναι νομίζομεν τῆς πρὸς τὸν νόμον ἡμῶν εὐπειθείας.
\vs{17}Διὸ δὲ κατʼ οὐδένα τρόπον παρανομεῖν ἀξιοῦμεν.
\vs{18}Καί τοι εἰ καὶ κατὰ ἀλήθειαν μὴ ἦν ὁ νόμος ἡμῶν, ὡς σὺ ὑπολαμβάνεις, θεῖος, (ἄλλως δὲ νομίζομεν αὐτὸν εἶναι θεῖον) οὐδὲ οὕτως ἐξὸν ἡμῖν ἦν τὴν ἐπὶ τῇ εὐσεβείᾳ δόκαν ἀκυρῶσαι.
\vs{19}Μὴ μικρὰν οὖν εἶναι νομίσῃς ταύτην, εἰ μιαροφαγήσεμεν, ἁμαρτίαν.
\vs{20}Τὸ γὰρ ἐν μικροῖς καὶ ἐν μεγάλοις παρανομεῖν ἰσοδύναμόν ἐστιν·
\vs{21}διʼ ἑκατέρου γὰρ ὡς ὁμοίως ὁ νόμος ὑπερηφανεῖται.

\vs{22}Χλευάζεις δὲ ἡμῶν τὴν φιλοσοφίαν, ὥσπερ οὐ μετὰ εὐλογιστίας ἐν αὐτῇ βιούντων.
\vs{23}Σωφροσύνην τε γὰρ ἡμᾶς ἐκδιδάσκει, ὥστε πασῶν τῶν ἡδονῶν καὶ ἐπιθυμιῶν κρατεῖν, καὶ ἀνδρείαν ἐξασκεῖν, ὥστε πάντα πόνον ἑκουσίως ὑπομένειν·
\vs{24}καὶ δικαιοσύνην παιδεύει, ὥστε διὰ πάντων τῶν ἠθῶν ἰσονομεῖν καὶ εὐσέβειαν διδάσκειν, ὥστε μόνον τὸν ὄντα Θεὸν σέβειν μεγαλοπρεπῶς.
\vs{25}Διὸ οὐ μιαροφαγοῦμεν· πιστεύοντες γὰρ Θεοῦ καθεστᾶναι τὸν νόμον, οἴδαμεν ὅτι καὶ κατὰ φύσιν ἡμῖν συμπαθεῖ νομοθετῶν ὁ τοῦ κόσμου κτίστης·
\vs{26}τὰ μὲν οἰκειωθωσόμενα ἡμῶν ταῖς ψυχαῖς ἐπέτρεψεν ἐσθίειν, τὰ δὲ ἐναντιωθησόμενα ἐκώλυσεν σαρκοφαγεῖν.

\vs{27}Τυραννικὸν δὲ, οὐ μόνον ἀναγκάζεις ἡμᾶς παρανομεῖν, ἀλλὰ καὶ ἐσθίειν, ὅπως τῇ ἐχθίστῃ ἡμῶν μιαροφαγίᾳ ταύτῃ ἔτι ἐγγελάσῃς.
\vs{28}Ἀλλʼ οὐ γελάσεις κατʼ ἐμοῦ τοῦτον τὸν γέλωτα·
\vs{29}οὔτε τοὺς ἱεροὺς τῶν προγόνων περὶ τοῦ φυλάξαι τὸν νόμον ὅρκους οὐ παρήσω.
\vs{30}Οὐδʼ ἂν ἐκκόψεις μου τὰ ὄμματα, καὶ τὰ σπλάγχνα μου τήξεις.
\vs{31}Οὐχ οὕτως εἰμὶ γέρων ἐγὼ καὶ ἄνανδρος, ὥστε μοι διὰ τὴν εὐσέβειαν μὴ νεάζειν τὸν λογισμόν.

\vs{32}Πρὸς ταῦτα τροχοὺς εὐτρέπιζε, καὶ τὸ πῦρ ἐκφύσα σφοδρότερον.
\vs{33}Οὐχ οὕτως οἰκτειρήσω τὸ ἐμαυτοῦ γῆρας, ὥστε με διʼ ἐμαυτοῦ τὸν πάτριον καταλῦσαι νόμον.
\vs{34}Οὐ ψεύσομαί σε, παιδευτὰ νόμε, οὐδὲ φεύξομαί σε, φίλη ἐγκράτεια.
\vs{35}Οὐδὲ καταισχυνῶ σε, φιλόσοφε λόγε, οὐδὲ ἐξαρνήσεμαί σε, ἱερωσύνη τιμία, καὶ νομοθεσίας ἐπιστήμη·
\vs{36}οὐδὲ μιανεῖς μου τὸ σεμνὸν γήρηως στόμα, οὐδὲ νομίμου βίου ἡλικίαν.

\vs{37}Ἁγνόν με οἱ πατέρες προσδέξονται, μὴ φοβηθέντα σου τὰς μέχρι θανάτου ἀνάγκας.
\vs{38}Ἀσεβῶν μὲν γὰρ τυραννήσεις· τῶν δὲ ἐμῶν περὶ τῆς εὐσεβείας λογισμῶν οὔτε λόγοις δεσπόσεις, οὔτε διʼ ἔργων.

\ch{6}
Τοῦτον τὸν τρόπον ἀντιρητορεύσαντα ταῖς τοῦ τυράννου παρηγορίαις, παραστάντες οἱ δορυφόροι πικρῶς ἔσυραν ἐπὶ τὰ βασανιστήρια τὸν Ἐλεάζαρον.
\vs{2}Καὶ πρῶτον μὲν περιέδυσαν τὸν γηραιὸν ἐκκεκοσμημένον περὶ τὴν εὐσέβειαν εὐσχημοσύνην.
\vs{3}Ἔπειτα περιαγκωνίσαντες ἑκατέρωθεν, μάστιξιν κατῇκιζον·
\vs{4}πείσθητι ταῖς τοῦ βασιλέως ἐντολαῖς, ἑτέρωθεν κήρυκος ἐπιβοῶντος.

\vs{5}Ὁ δὲ μεγαλόφρων καὶ εὐγενὴς ὡς ἀληθῶς Ἐλεάζαρος, ὥσπερ ἐν ὀνείρω βασανιζόμενος κατʼ οὐδένα τρόπον μετετρέπετο.
\vs{6}Ἀλλὰ ὑψηλοὺς ἀνατείνας εἰς τὸν οὐρανὸν τοὺς ὀφθαλμοὺς, ἀπεξαίνετο ταῖς μάστιξιν τὰς σάρκας ὁ γέρων, καὶ κατεῤῥεῖτο τῷ αἵματι,
\vs{7}καὶ τὰ πλευρὰ κατετιτρώσκετο, καὶ πίπτων εἰς τὸ ἔδαφος, ἀπὸ τοῦ μὴ φέρειν τὸ σῶμα τὰς ἀλγηδόνας, ὀρθὸν εἶχεν καὶ ἀκλινῆ τὸν λογισμόν.
\vs{8}Λὰξ γέ τοι τῶν πικρῶν τις δορυφόρων, εἰς τοὺς κενεῶνας ἐναλλόμενος ἔτυπτεν, ὅπως ἐξανίσταιτο πίπτων.

\vs{9}Ὁ δὲ ὑπέμενεν τοὺς πόνους, καὶ περιεφρόνει τῆς ἀνάγκης, καὶ διεκαρτέρει τοὺς αἰκισμοὺς,
\vs{10}καὶ καθάπερ γενναῖος ἀθλητὴς τυπτόμενος ἐνίκα τοὺς βασανίζοντας ὁ γέρων.
\vs{11}Ἱδρῶν γέ τοι τὸ πρόσωπον, καὶ ἐπασθμαίνων σφοδρῶς, καὶ ὑπʼ αὐτῶν τῶν βασανιζόντων ἐθαυμάζετο ἐπὶ τῇ εὐτυχίᾳ.

\vs{12}Ὅθεν τὰ μὲν ἐλεοῦντες τὰ τοῦ γήρως αὐτοῦ, τὰ δὲ ἐν συμπαθείᾳ τῆς συνηθείας ὄντες,
\vs{13}τὰ δὲ ἐν θαυμαστῷ τῆς καρτερίας προσιόντες αὐτῷ τινὲς τῶν τοῦ βασιλέως ἔλεγον,
\vs{14}τί τοῖς κακοῖς τούτοις σεαντὸν ἀλογίστως ἀπολλεῖς,
\vs{15}Ἐλεάζαρ; ἡμεῖς μὲν τῶν ἡψημένων βρωμάτων παραθήσομεν· σὺ δὲ ὑποκρινόμενος τῶν ὑείων ἀπογεύσασθαι, σώθητι.

\vs{16}Καὶ ὁ Ἐλεάζαρος, ὥσπερ πικρότερον διὰ τῆς συμβουλίας αἰκισθεὶς, ἀνεβόησεν,
\vs{17}μὴ οὕτως κακῶς φρονήσαιμεν οἱ Ἁβραὰμ παῖδες, ὥστε μαλακοψυχήσαντας ἀπρεπὲς ἡμῖν δρᾶμα ὑποκρίνασθαι.
\vs{18}Καὶ γὰρ ἀλόγιστον, εἰ πρὸς ἀλήθειαν ζήσαντες τὸν μέχρι γήρως βίον, καὶ τὴν ἐπʼ αὐτῶν δόξαν νομίμως φυλάσσοντες, νῦν μεταβαλοίμεθα,
\vs{19}καὶ αὐτοὶ μὲν ἡμεῖς γενοίμεθα τοῖς νέοις ἀσεβείας τύπος, ἵνα παράδειγμα γενώμεθα τῆς μιεροφαγίας.
\vs{20}Αἰσχρὸν γὰρ εἰ ἐπιβιώσωμεν ἀλίγον χρόνον,
\vs{21}καὶ τοῦτον καταγελώμενοι πρὸς ἁπάντων ἐπὶ δειλίᾳ· καὶ ὑπὸ μὲν τοῦ τυράννου καταφρονηθῶμεν ὡς ἄνανδροι, τὸν δὲ θεῖον ἡμῶν νόμον μέχρι θανάτου μὴ προασπίσαιμεν.
\vs{22}Πρὸς ταῦτα ὑμεῖς μὲν, ὦ Ἁβραὰμ παῖδες, εὐγενῶς ὑπὲρ τῆς εὐσεβείας τελευτᾶτε.
\vs{23}Οἱ δὲ τοῦ τυράννου δορυφόροι, τί μέλλετε;

\vs{24}Πρὸς τὰς ἀνάγκας οὕτως μεγαλοφρονοῦντα αὐτὸν ἰδόντες καὶ μηδὲ πρὸς τὸν οἰκτιρμὸν αὐτῶν μεταβαλλόμενον, ἐπὶ πῦρ αὐτὸν ἤγαγον.
\vs{25}Ἔνθα διὰ κακοτέχνων ὀργάνων καταφλέγοντες αὐτὸν ὑπερέπτοσαν, καὶ δυσώδεις χυλοὺς εἰς τοὺς μυκτῆρας αὐτοῦ κατέχεον.

\vs{26}Ὁ δὲ μέχρι τῶν ὀστέων ἤδη κατακεκαυμένος καὶ μέλλων λιποθυμεῖν, ἀνέτεινεν τὰ ὄμματα πρὸς τὸν Θεὸν, καὶ εἶπεν, σὺ οἶσθα, Θεὲ, παρόν μοι σώζεσθαι,
\vs{27}βασάνοις καυστικαῖς ἀποθνήσκω διὰ τὸν νόμον.
\vs{28}Ἵλεως γενοῦ τῷ ἔθνει σου, ἀρκεσθεὶς τῇ ἡμετέρᾳ περὶ αὐτῶν δίκῃ.
\vs{29}Καθάρσιον αὐτῶν ποίησον τὸ ἐμὸν αἷμα, καὶ ἀντίψυχον αὐτῶν λαβὲ τὴν ἐμὴν ψυχήν.
\vs{30}Καὶ ταῦτα εἰπὼν ὁ ἱερὸς ἀνὴρ εὐγενῶς ταῖς βασάνοις ἐναπέθανεν, καὶ μέχρι τῶν τοῦ θανάτου βασάνων ἀντέστη τῷ λσγισμῷ διὰ τὸν νόμον.

\vs{31}Ὁμολογουμένως οἶν δεσπότης ἐστὶν τῶν παθῶν ὁ εὐσεβὴς λογισμός.
\vs{32}Εἰ γὰρ τὰ πάθη τοῦ λογισμοῦ κεκρατήκει, τούτοις ἂν ἀπεδόμην τὴν τῆς ἐπικρατείας μαρτυρίαν.
\vs{33}Νυνὶ δὲ τοῦ λογισμοῦ τὰ πάθη νικήσαντος, αὐτῷ προσηκόντως τὴν τῆς ἡγεμονίας προσνέμομεν ἐξουσίαν.

\vs{34}Καὶ δίκαιόν ἐστιν ὁμολογεῖν ἡμᾶς, τὸ κράτος εἶναι τοῦ λογισμοῦ, ὅπου γε καὶ τῶν ἔξωθεν ἀλγηδόνων ἐπικρατεῖ.
\vs{35}Ἐπεὶ καὶ γελοῖον· καὶ οὐ μόνον τῶν ἀλγηδόνων ἐπιδείκνυμι κεκρατηκέναι τὸν λογισμὸν, ἀλλὰ καὶ τῶν ἡδονῶν κρατεῖν, μηδὲ αὐταῖς ὑπείκειν.

\ch{7}
Ὥσπερ καὶ ἄριστος κυβερνήτης ὁ τοῦ πατρὸς ἡμῶν Ἐλεαζάρου λογισμὸς, πηδαλιουεχῶν τὴν τῆς εὐσεβείας ναῦν ἐν τῷ τῶν παθῶν πελάγει,
\vs{2}καὶ καταικιζόμενος ταῖς τοῦ τυράννου ἀπειλαῖς, καὶ καταντλούμενος ταῖς τῶν βασάνων τρικυμίαις,
\vs{3}κατʼ οὐδένα τρόπον μετέτρεψεν τοὺς τῆς εὐσεβείας οἴακας, ἕως οὗ ἔπλευσεν ἐπὶ τὸν τῆς θανάτου νίκης λιμένα.

\vs{4}Οὐχ οὕτως πόλις πολλοῖς καὶ ποικίλοις μηχανήμασιν ἀντέσχεν ποτὲ πολιορκουμένη, ὡς ὁ πανάγιος ἐκεῖνος τὴν ἱερὰν ψυχὴν αἰκισμοῖς τε καὶ στρέβλαις πυρπολούμενος, ἐκίνησεν τοὺς πολιορκοῦντας, διὰ τὸν ὑπερασπίζοντα τῆς εὐσεβείας λογισμόν.
\vs{5}Ὥσπερ γὰρ πρόκρημνον ἄκραν, τὴν ἑαυτοῦ διὰνοιαν ὁ πατὴρ Ἐλεάζαρος ἐκτείνας, περιέκλασεν τοὺς μαινομένους τῶν παθῶν κλύδωνας.

\vs{6}Ὦ ἄξιε τῆς ἱερωσύνης ἱερεῦ, οὐκ ἐμίανας τοὺς ἱεροὺς ὀδόντας, οὐδὲ τὴν θεοσέβειαν καὶ καθαρισμὸν χωρήσασαν γαστέρα ἐκοινώνησας μιεροφαγίᾳ·
\vs{7}Ὦ σύμφωνε νόμου, καὶ φιλόσοφε θείου βίου.
\vs{8}Τοίουτους δεῖ εἶναι τοὺς δημιουργοῦντας τὸν νόμον ἰδίῳ αἵματι, καὶ γενναίῳ ἱδρῶτι τοῖς μέχρι θανάτου πάθεσιν ὑπερασπίζοντας.

\vs{9}Σὺ πάτερ, τὴν εὐνομίαν ἡμῶν διὰ τῶν ὑπομονῶν εἰς δόξαν ἐκύρωσας, καὶ τὴν ἁγιαστίαν σεμνολογήσας οὐ κατέλυσας, καὶ διὰ τῶν ἔργων ἐπιστοποίησας τοὺς τῆς φιλοσοφίας λὸγους.
\vs{10}Ὦ βασάνων βιότερε γέρων, πυρὸς εὐτονώτερε πρεσβύτα, καὶ παθῶν μέγιστε βασιλεῦ Ἐλεάζαρ.

\vs{11}Ὥσπερ γὰρ ὁ πατὴρ Ἀαρὸν τῷ θυμιατηρίῳ κατθωπλισμένος, διὰ τοῦ ἐθνοπλήθου ἐπιτρέχων τὸν ἐμπυριστὴν ἐνίκησεν ἄγγελον.
\vs{12}Οὕτως ὁ Ἀαρωνίδης Ἐλεάζαρος διὰ τοῦ πυρὸς ὑπερτηκόμενος οὐ μετετράπη τὸν λογισμόν.
\vs{13}Καίτοι τὸ θαυμασιώτατον, γέρων ὢν, λελυμένων μὲν ἤδη τῶν τοῦ σώματος πόνων, καὶ περιεχαλασμένων δὲ τῶν σαρκῶν, κεκμηκότων δὲ καὶ τῶν νεύρων, ἀνενέασεν.
\vs{14}Τῷ πνεύματι τοῦ λογισμοῦ, καὶ τῷ Ἰσακείῳ λογισμῷ τὴν πολυκέφαλον στρέβλαν ἠκύρωσεν.
\vs{15}Ὦ μακαρίου γήρως, καὶ σεμνῆς πολιᾶς, καὶ βίου νομίμου, ὃν πιστὴ θανάτου σφραγὶς ἐτελείωσεν.
\vs{16}Εἰ δὲ τοίνυν γέρων τῶν μέχρι θανάτου βασάνων περιεφρόνησεν διʼ εὐσέβειαν, ὁμολογουμένως ἡγεμών ἐστιν τῶν παθῶν ὁ εὐσεβὴς λογισμός.

\vs{17}Ἴσως δʼ ἂν εἴποιέν τινες, τῶν παθῶν οὐ πάντες περικρατοῦσιν, ὅτι οὐδὲ πάντες φρόνιμον ἔχουσιν τὸν λογισμόν.
\vs{18}Ἀλλʼ ὅσοι εὑσεβείας προνοοῦσιν ἐξ ὅλης καρδίας, οὗτοι μόνοι δύνανται κρατεῖν τῶν τῆς σαρκὸς παθῶν·
\vs{19}οἱ πιστεύοντες, ὅτι Θεῷ οὐκ ἀποθνήσκουσιν, ὥσπερ γὰρ οἱ πατριάρχαι ἡμῶν Ἁβραὰμ, Ἰσαὰκ, Ἰακὼβ, ζῶσι τῷ Θεῷ.

\vs{20}Οὐδὲν οὖν ἐναντιοῦται τὸ φαίνεσθαί τινας παθοκρατεῖσθαι διὰ τὸν ἀσθενῆ λογισμόν.
\vs{21}Ἐπεὶ τίς πρὸς ὅλον τὸν τῆς φιλοσοφίας κανόνα εὐσεβῶς φιλοσοφῶν, καὶ πεπιστευκὼς Θεῷ,
\vs{22}καὶ εἰδὼς ὅτι διὰ τὴν ἀρετὴν πάντα πόνον ὑπομένειν μακάριόν ἐστιν, οὐκ ἂν περικρατήσειεν τῶν παθῶν διὰ τὴν εὐσέβειαν;
\vs{23}μόνος γὰρ ὁ σοφὸς καὶ σώφρων ἀνδρεῖός ἐστιν τῶν παθῶν κύριος.
\vs{24}Διὰ τοῦτο γέ τοι καὶ μειρακίσκοι τῷ τῆς εὐσεβείας λογισμῷ φιλοσοφουντες χαλεπωτερων βασανιστηρίων ἐπεκράτησαν.
\vs{25}Ἐπειδὴ γὰρ κατὰ τὴν πρώτην πεῖραν ἐνικήθη περιφανὴς ὁ τύραννος, μὴ δυνηθεὶς ἀναγκάσαι γέροντα μιαιροφαγῆσαι.

\ch{8}
Τὸ δὲ δὴ σφόδρα περιπαθῶς ἐκέλευσεν ἄλλους ἐκ τῆς ἠλικίας τῶν Ἑβραίων ἀγαγεῖν· καὶ εἰ μὲν μιεροφαγήσαιεν, ἀπολύειν φάγοντας· εἰ δὲ ἀντιλέγοιεν, πικρότερον βασανίζειν.

\vs{2}Ταῦτα διαδεξαμένου τοῦ τυράννου, παρῆσαν ἀγόμενοι μετὰ γηραιᾶς μητρὸς ἑπτὰ ἀδελφοὶ, καλοί τε καὶ αἰδήμονες καὶ γενναῖοι καὶ ἐν παντὶ χαριέντες.
\vs{3}Οὓς ἰδὼν ὁ τύραννος καθάπερ ἐν χορῷ περιέχοντας μέσην τὴν μητέρα, ἤσθετο ἐπʼ αὐτοῖς, καὶ τῆς εὐπρεπείας ἐκπλαγεὶς καὶ τῆς εὐγενείας προσεμειδίασεν αὐτοῖς, καὶ πλησίον καλέσας, ἔφη,

\vs{4}Ὦ νεανίαι φιλοφρόνως ἐγὼ καθʼ ἐνὸς ἑκάστου ὑμῶν θαυμάζω τὸ κάλλος· καὶ τὸ πλῆθος τοσούτων ἀδελφῶν ὑπερτιμῶν, οὐ μόνον συμβουλεύω μὴ μανῆναι τὴν αὐτὴν τῷ προβασανισθέντι γέροντι μανίαν·
\vs{5}ἀλλὰ καὶ παρακαλῶ συνείξαντας τῆς ἐμῆς ἀπολαῦσαι φιλίας· δυναίμην γὰρ ὥσπερ κολάζειν τοὺς ἐπιτάγμασιν, σὕτως καὶ εὐεργετεῖν τοὺς εὐπειθοῦντάς μοι.

\vs{6}Πιστεύσατε οὖν, καὶ ἀρχὰς ἐπὶ τῶν ἐμῶν πραγμάτων ἡγεμονικὰς λήψεσθε, ἀρνησάμενοι τὸν πάτριον ἡμῶν τῆς πολιτείας θεσμόν·
\vs{7}καὶ μεταλαβόντες Ἑλληνικοῦ βίου, καὶ μεταδιαιτηθέντες ἐντρυφήσατε ταῖς νεότησιν ὑμῶν.
\vs{8}Ἐπεὶ ἐὰν ὀργίλως με διάθησθε διὰ τῆς ἀπειθείας ὑμῶν, ἀναγκάσετέ με ἐπὶ δειναῖς κολάσεσιν ἕνα ἕκαστον ὑμῶν διὰ τῶν βασάνων ἀπολέσαι.
\vs{9}Κατελεήσατε οὖν ἑαυτοὺς, οὕς καὶ ὁ πολέμιος ἔγωγε καὶ τῆς ἡλικίας καὶ τῆς εὐμορφίας οἰκτείρομαι.
\vs{10}Οὐ διαλογιεῖσθε τοῦτο, ὅτι οὐδὲν ὑμῖν ἀπειθήσασιν πλὴν τοῦ μετὰ στρεβλῶν ἀποθανεῖν ἀπόκειται;

\vs{11}Ταῦτα δὲ λέγων, ἐκέλευσεν εἰς τὸ ἔμπροσθεν προτεθῆναι τὰ βασανιστήρια, ὅπως καὶ διὰ τοῦ φόβου πείσειεν αὐτοὺς μιεροφαγῆσαι.
\vs{12}Ὡς δὲ τροχούς τε καὶ ἀρθενβόλους στρεβλωτήρια, καὶ τροχαντῆρας καὶ καταπέλτας καὶ λέβητας, τήγανά τε καὶ δακτυλήθρας, καὶ χεῖρας σιδηρᾶς καὶ σφῆνας, καὶ τὰ ζώπυρα τοῦ πυρὸς οἱ δορυφόροι προέθησαν, ὑπολαβὼν δὲ ὁ τύραννος, ἔφη, μειράκια φοβήθητε,
\vs{13}καὶ ἥν σέβεσθε δίκην, ἵλεως ὑμῖν ἔσται διʼ ἀνάγκην παρανομήσασιν.
\vs{14}Οἱ δὲ ἀκούσαντες ἐπαγωγὰ, καὶ ὁρῶντες δεινὰ, οὐ μόνον οὐκ ἐφοβήθησαν, ἀλλὰ καὶ ἀντεφιλοσόφῃσαν τῷ τυράννῳ, καὶ διὰ τῆς εὐλογιστίας τὴν τυραννίδα αὐτοῦ κατέλυσαν.

\vs{15}Καί τοι λογισώμεθα· εἰ δειλόψυχοί τινες ἦσαν, καὶ ἄνανδροι ἐν αὐτοῖς, ποίοις ἂν ἐχρήσαντο λόγοις; οὐχὶ τούτοις;
\vs{16}Ὦ τάλανες ἡμεῖς, καὶ λίαν ἀνόητοι· βασιλέως ἡμᾶς παρακαλοῦντος, καὶ ἐπὶ εὐεργεσίᾳ φωνοῦντος, μὴ πεισθείημεν αὐτῷ;
\vs{17}Τί βουλήμασιν κενοῖς ἐαυτοὺς εὐφραίνομεν, καὶ θανατηφόρον ἀπείθειαν τολμῶμεν;
\vs{18}Οὐ φοβησόμεθα, ἄνδρες ἀδελφοί, τὰ βασανιστήρια, καὶ λογιούμεθα τὰς τῶν βασάνων ἀπειλὰς, καὶ φευξόμεθα τὴν κενοδοξίαν ταύτην καὶ ὀλεθροφόρον ἀλαζονείαν;
\vs{19}Ἐλεήσωμεν τὰς ἑαυτῶν ἡλικίας, καὶ κατοικτειρήσωμεν τὸ τῆς μητρὸς γῆρας·
\vs{20}καὶ ἐνθυμηθῶμεν, ὅτι ἀπειθοῦντες τεθνηξόμεθα.
\vs{21}Συγγνώσεται δὲ ἡμῖν καὶ ἡ θεία δίκη διʼ ἀνάγκην τὸν βασιλὲα φοβηθεῖσιν.
\vs{22}Τί ἐξάγομεν ἑαυτοὺς τοῦ ἡδίστου βίου, καὶ ἐπιστεροῦμεν ἐαυτοὺς τοῦ γλυκέος κόσμου;
\vs{23}Μὴ βιαζώμεθα τὴν ἀνάγκην, μηδὲ κενοδοξήσωμεν ἐπʼ τῇ ἑαυτῶν στρέβλῃ.
\vs{24}Οὐδὲ αὐτὸς ὁ ναὸς ἑκουσίως ἡμᾶς θανατοῖ φοβηθέντας τὰ βασανιστήρια.
\vs{25}Πόθεν ἡμῖν ἡ τοσαύτη ἐντέηκεν φιλονεικία, καὶ ἡ θανατεφόρος ἀρέσκει καρτερία, παρὸν μετὰ ἀταραξίας χρὴ τῷ βασιλεῖ πεισθέντας;

\vs{26}Ἀλλὰ τούτων οὐδὲν εἶπον οἱ νεανίαι βασανίζεσθαι μέλλοντες, οὐδὲ ἐνεθυμήθησαν.
\vs{27}Ἦσαν γὰρ περιφρονες τῶν παθῶν, καὶ αὐτηκράτορες τῶν ἀλγηδόνων. Ὥστε ἅμα τῷ παύσασθαι τὸν τύραννον συμβουλεύοντα αὐτοῖς μιεροφαγῆσαι, πάντες διὰ μιᾶς φωνῆς ὁμοῦ, ὥσπερ ἀπὸ τῆς αὐτῆς ψυχῆς, εἶπον,

\ch{9}
Τί μέλλεις, ὦ τύραννε; ἕτοιμοι γάρ ἐσμεν ἀποθνήσκειν, ἢ παραβαίνειν τὰς πατρίους ἡμῶν ἐντολάς.
\vs{2}Καὶ αἰσχυνόμεθα γὰρ τοὺς προγόνους εἰκότως, εἰ μὴ τῇ τοῦ νόμου εὐπειθείᾳ καὶ συμβούλῳ γνώσει χρησαίμεθα.

\vs{3}Σύμβουλε τύραννε παρανομίας, μὴ ἡμᾶς μισῶν ὑπὲρ αὐτοὺς ἡμᾶς ἐλέα.
\vs{4}Χαλεπώτερον γὰρ αὐτοὺς τοῦ θανάτου νομίζομεν εἶναί σου τὸν ἐπὶ τῇ παρανόμῳ σωτηρίᾳ ἡμῶν ἔλεον.
\vs{5}Ἐκφοβεῖς δὲ ἡμᾶς, τὸν διὰ τῶν βασάνων ἡμῖν θάνατον ἀπειλῶν, ὥσπερ οὐχὶ πρὸ βραχέως παρὰ Ἐλεαζάρου μαθών.
\vs{6}Εἰ δʼ οἱ γέροντες τῶν Ἑβραίων διὰ τὴν εὐσέβειαν καὶ βασανισμὸυς ὑπομείναντες ἀπέθανον, ἀποθάνοιμεν ἂν δικαιότερον ἡμεῖς οἱ νέοι, τὰς βασάνους τῶν σῶν ἀναγκῶν ὑπεριδόντες, ἃς καὶ ὁ παιδευτὴς γέρων ἐνίκησεν.

\vs{7}Πείραζε γαροῦν τύραννε· καὶ τὰς ἡμῶν ψυχὰς εἰ θανατώσεις διὰ τὴν εὐσέβειαν, μὴ νομίσῃς ἡμᾶς βλάπτειν βασανίζων.
\vs{8}Ἡμεῖς μὲν γὰρ διὰ τῆσδε τῆς κακοπαθείας καὶ ὑπομονῆς, τὰ τῆς ἀρετῆς ἆθλα οἴσομεν.
\vs{9}Σὺ δὲ διὰ τὴν ἡμῶν μιαροφονίαν αὐτάρχη καρτερήσεις περὶ τῆς θείας δίκης αἰώνιον βάσανον διὰ πυρός.

\vs{10}Ταῦτα αὐτῶν εἰπόντων, οὐ μόνον ὡς κατὰ ἀπειθούντων ἐχαλέπαινεν ὁ τύραννος, ἀλλʼ ὡς καὶ κατὰ ἀχαρίστων ὠργίσθη.
\vs{11}Ὅθεν τὸν πρεσβύτατον αὐτῶν κελευθέντες παρήγαγον οἱ μαστισταὶ, καὶ διαῤῥήξαντες τὸν χιτῶνα διέδησαν τὰς χεῖρας αὐτοῦ καὶ τοὺς βραχίονας ἱμᾶσιν ἑκατέρωθεν.
\vs{12}Ὡς δὲ τύπτοντες ταῖς μάστιξιν ἐκοπίασαν, μηδὲν ἀνύοντες, ἀνέβαλον αὐτὸν ἐπὶ τὸν τροχόν.
\vs{13}Περὶ ὃν κατατεινόμενος ὁ εὐγενὴς νεανίας, ἔξαρθρος ἐγίνετο.
\vs{14}Καὶ κατὰ πᾶν μέλος κλώμενος κατηγόρει, λέγων,

\vs{15}Τύραννε μιαιρώτατε, καὶ τῆς οὐρανίου δίκης ἐχθρὲ, καὶ ὠμόφρον, οὐκ ἀνδροφονήσαντά με τοῦτον καταικίζεις τὸν τρόπον, οὐδὲ ἀσεβήσαντα, ἀλλα θείου νόμου προασπίζοντα.
\vs{16}Καὶ τῶν δορυφόρων λεγόντων, ὁμολόγησον φαγεῖν, οὕπως ἀπαλλαγῇς τῶν βασάνων,
\vs{17}ὁ δὲ εἶπεν, οὐχ οὕτως ἰσχυρὸς ὑμῶν ἐστιν ὁ τρόπος, ὦ μιαιροὶ διὰκονοι, ὥστε μου τὸν λογισμὸν ἄξαι· τέμνετέ μου μέλη, καὶ πυροῦτε τὰς σάρκας, καὶ στρεβλοῦτε τὰ ἄρθρα.
\vs{18}Διὰ πασῶν γὰρ ὑμᾶς πείσω τῶν βασάνων· ὅτι μόνοι παῖδες Ἑβραίων ὑπὲρ ἀρετῆς εἰσιν ἀνίκητοι.

\vs{19}Ταῦτα λέγοντες εἰς πῦρ ἐπέτρωσαν, καὶ διερεθίζοντες, τὸν τροχὸν προσεπικατέτεινον.
\vs{20}Ἐμολύνετο δὲ πάντοθεν αἵματι ὁ τρόχος, καὶ ὁ σωρὸς τῆς ἀνθρακιᾶς τοῖς τῶν ἰχώρων ἐσβέννυτο σταλαγμοῖς, καὶ περὶ τοὺς αὔξονας τοῦ ὀργάνου περιέῤῥεον αἱ σάρκες.

\vs{21}Καὶ περιτετηκμένον ἤδη ἔχων τὸ τῶν ὀστέων πῆγμα ὁ μεγαλόφρων καὶ Ἀβραμιαῖος νεανίας οὐκ ἐστέναξεν.
\vs{22}Ἀλλʼ ὥσπερ ἐν πυρὶ μετασχηματιζόμενος εἰς ἀφθαρσίαν, ὑπέμεινεν εὐγενῶς τὰς στρέβλας.
\vs{23}Μιμήσασθέ με, ἀδελφοὶ, λέγων· μή μου τὸν αἰῶνα λειποτακτήσητε, μηδʼ ἐξομόσησθέ μου τὴν τῆς εὐψυχίας ἀδελφότητα· ἱερὰν καὶ εὐγενῆ στρατείαν στρατεύσασθε περὶ τῆς εὐσεβείας.
\vs{24}Διʼ ἧς ἷλεως ἡ δικαία καὶ πάτριος ἡμῶν πρόνοια τῷ ἔθνει γενηθεῖσα τιμωρήσειεν τὸν ἀλάστορα τύραννον.
\vs{25}Καὶ ταῦτα εἰπὼν ὁ ἱεροπρεπὴς νεανίας, ἀπέῤῥηξεν τὴν ψυχήν.

\vs{26}Θαυμασάντων δὲ πάντων τὴν καρτεροψυχίαν αὐτοῦ, ἦγον οἱ δορυφόροι τὸν καθʼ ἡλικίαν τοῶ προτέρου δεύτερον, καὶ σιδηρᾶς ἐναρμοσάμενοί χεῖρας, ὀξέσιν τοῖς ὄνυξιν, τοῖς ὀργάνοις καταπέλτῃ προσέδησαν αὐτόν.
\vs{27}Ὡς δὲ, εἰ φαγεῖν βούλοιτο πρὶν βασανίζεσθαι πυνθανόμενοι, τὴν εὐγενῆ γνώμην ἤκουσαν·
\vs{28}ἀπὸ τῶν τενόντων ταῖς σιδηραῖς χερσὶν ἐπισπασάμενοι, μέχρι γε τῶν γενείων τὴν σάρκα πᾶσαν καὶ τὴν τῆς κεφαλῆς δορὰν οἱ παρδάλειοι θῆρες ἀπέσυραν· ὁ δὲ ταύτην βαρέως τὴν ἀλγηδόνα καρτερῶν, ἔλεγεν,
\vs{29}Ὡς ἡδὺς πᾶς τρόπος θανάτου, διὰ τὴν πάτριον ἡμῶν εὐσέβειαν· ἔφη τε πρὸς τὸν τύραννον,

\vs{30}Οὐ δοκεῖς, πάντων ὠμότατε τύραννε, πλεῖων ἐμοῦ σε νὺν βασανίζεσθαι, ὁρῶν σου νικώμενον τὸν τῆς τυραννίδος ὑπερήφανον λογισμὸν ὑπὸ τῆς διὰ τὴν εὐσέβειαν ἡμῶν ὑπομονῆς.
\vs{31}Ἐγὼ μὲν γὰρ ταῖς διὰ τὴν ἀρετὴν ἡδοναῖς τὸν πόνον ἐπικουφίζομαι.
\vs{32}Σὺ δὲ ἐν ταῖς τῆς ἀσεβείας ἀπειλαῖς βασανίζῃ· οὐκ ἐκφεύξῃ δὲ, μιαιρότατε τύραννε, τὰς τῆς θείας ὀργῆς δίκας.

\ch{10}
Καὶ τούτου τὸν ἀοίδιμον θάνατον καρτερήσαντος, ὁ τρίτος ἤγετο, παρακαλούμενος πολλὰ ὑπὸ πολλῶν ὅπως ἀπογευσάμενος σώζοιτο.
\vs{2}Ὁ δὲ ἀναβοήσας, ἔφη, ἤ ἀγνοεῖτε, ὅτι αὐτός με τοῖς ἀποθανοῦσιν ἔσπειρεν πατὴρ, καὶ ἡ αὐτὴ μήτηρ ἐγέννεσιν, καὶ ἐπὶ τοῖς αὐτοῖς ἀνετράφην δόγμασιν;
\vs{3}Οὐκ ἐξόμνυμαι τὴν εὐγενῆ τῆς ἀδελφότητος συγγένειαν.
\vs{3a}Πρὸς ταῦτα εἴ τι ἔχετε κολαστήριον προσαγάγετε τῷ σώματί μου· τῆς γὰρ ψυχῆς μου, οὐδʼ ἂν θέλητε, ἅψασθαι δύνασθε.

\vs{5}Οἱ δὲ πίκρῶς ἐνέγκαντες τὴν παῤῥησίαν τοῦ ἀνδρὸς, ἀρθρεμβόλοις ὀργάνοις τὰς χεῖρας αὐτοῦ καὶ τοὺς πόδας ἐξήρθρουν, καὶ ἐξ ἁρμῶν ἀναμοχλεύοντες ἐξεμέλιζον·
\vs{6}καὶ τοὺς δακτύλους, καὶ τοὺς βραχίονας, καὶ τὰ σκέλη, καὶ τοὺς ἀγκῶνας περιέλκων.
\vs{7}Καὶ κατὰ μηδένα τρόπον ἰσχύοντες αὐτὸν ἄγξαι, περισύραντες τὸ δέρμα σὺν ἄκραις ταῖς τῶν δακτύλων κορυφαῖς ἀπεσκύθιζον, καὶ εὐθέως ἦγον ἐπὶ τὸν τροχόν.
\vs{8}Περὶ ὃν ἐκ σφονδύλων ἐκμελιζόμενος ἑώρα τὰς ἑαυτοῦ σάρκας περιλακιζομένας καὶ κατὰ σπλάγχνων σταγόνας αἵματος ἀποῤῥεούσας.
\vs{9}Μέλλων δὲ ἀποθνήσκειν,
\vs{10}ἔφη, ἡμεῖς μὲν ὦ μιαιρώτατε τύραννε, διὰ παιδείαν καί ἀρετὴν Θεοῦ ταῦτα πάσχομεν.
\vs{11}Σὺ δὲ διὰ τὴν ἀσέβειαν καὶ μιαιφονίαν, ἀκαταλύτους καρτερήσεις βασάνους.

\vs{12}Καὶ τούτου θανόντος ἀδελφοπρεπῶς, τὸν τέταρτον ἐπεσπῶντο, λέγοντες,
\vs{13}Μὴ μανῄς καὶ σὺ τοῖς ἀδελφοῖς σου τὴν αὐτὴν μανίαν· ἀλλὰ πεισθεὶς τῷ βασιλεῖ, σῶζε σεαυτόν.
\vs{14}Ὁ δὲ αὐτοῖς ἔφη, οὐχ οὕτως καυστικώτερον ἔχετε κατʼ ἐμοῦ τὸ πῦρ, ὥστε με δειλανδρῆσαι.
\vs{15}Μὰ τὸν μακάριον τῶν ἀδελφῶν μου θάνατον, καὶ τὸν αἰώνιον τοῦ τυράννου ὄλεθρον, καὶ τὸν ἀοίδιμον τῶν εὐσεβῶν βίον, οὐκ ἀρνήσομαι τὴν εὐγενῆ ἀδελφότητα.
\vs{16}Ἐπινόει, τύραννε, βασάνους· ἵνα καὶ διὰ τούτων μάθῃς, ὅτι ἀδελφός εἰμι τῶν προβεβανασισθέντων.

\vs{17}Ταῦτα ἀκούσας ὁ αἱμοβόρος καὶ φονώδης καὶ πανμιαιρώτατος Ἀντίοχος, ἐκέλευσεν τὴν γλῶτταν αὐτοῦ ἐκτεμεῖν.
\vs{18}Ὁ δὲ ἔφη, κᾂν ἀφέλῃς τὸ τῆς φωνῆς ὄργανον, καὶ σιωπώντων ἀκούει ὁ Θεός.
\vs{19}Ἰδοὺ κεχάλασται ἡ γλῶσσα· τέμνε· οὐ γὰρ παρὰ τοῦτο τὸν λογισμὸν ἡμῶν γλωσσοτομήσεις.
\vs{20}Ἡδέως ὑπὲρ τοῦ Θεοῦ τὰ τοῦ σώματος μέλη ἀκρωτηριαζόμενα.
\vs{21}Σὲ δὲ ταχέως μετελεύσεται ὁ Θεός· τὴν γὰρ τῶν θείων ὕμνων μελῳδὸν γλῶτταν ἐκτέμνεις.

\ch{11}
Ὡς δὲ καὶ οὗτος ταῖς βασάνοις καταικισθεὶς ἐναπέθανεν, ὁ πέμπτος παρεπήδησεν, λέγων,

\vs{2}Οὐ μέλλω, τύραννε, πρὸς τὸν ὑπὲρ τῆς ἀρετῆς βασανισμὸν παραιτεῖσθαι.
\vs{3}Αὐτὸς δʼ ἀπʼ ἐμαυτοῦ παρῆλθον, ὅπως κᾀμὲ κατακτείνας, περὶ πλειόνων ἀδικημάτων ὀφειλήσῃς τῇ οὐρανίῳ δίκῃ τιμωρίαν.
\vs{4}Ὦ μισάρετε καὶ μισάνθρωπε, τὶ δράσαντας ἡμᾶς τοῦτον πορθεῖς τὸν τρόπον;
\vs{5}Ἢ κακόν σοι δοκεῖ, ὅτι τὸν πάντων κτιστὴν εὐσεβοῦμεν, καὶ κατὰ τὸν ἐνάρετον αὐτοῦ ζῶμεν νόμον;
\vs{6}Ἀλλὰ ταῦτα τιμῶν, οὐ βασάνων ἐστὶν ἄξια.
\vs{6a}Εἴπερ ᾐσθάνου ἀνθρώπου πόθων, καὶ ἐλπίδα εἶχες παρὰ Θεῷ σωτηρίου·
\vs{6b}νῦν ἰδὲ ἀλλότριος ὢν Θεοῦ, πολεμεῖς τοὺς εὐσεβοῦντας εἰς τὸν Θεόν.

\vs{9}Τοιαῦτα λέγοντα οἱ δορυφόροι δήσαντες, αὐτὸν εἷλκον ἐπὶ τὸν καταπέλτην·
\vs{10}ἐφʼ ὃ δήσαντες αὐτὸν ἐπὶ τὰ γόνατα, καὶ ταῦτα ποδάγραις σιδηραῖς ἐφορμάσαντες τὴν ὀσφὺν αὐτοῦ ἐπὶ τὸν τροχιαῖον σφῆνα κατέκαμψαν· περὶ ὃν ὅλος ἐπὶ τὸν τρονὸν σκορπίου τρόπον ἀνακλώμενος ἐξεμελίζετο.
\vs{11}Κατὰ τοῦτον τὸν τρόπον καὶ τὸ πνεῦμα στενοχωρούμενος, καὶ τὸ σῶμα ἀγχόμενος, καλὰς, ἔλεγεν, ἄκων,
\vs{12}ὦ τύραννε, χάριτας ἡμῖν χαρίζῃ διὰ γενναιοτέρων πόνων ἐπιδείξασθαι παρέχων τὴν εἰς τὸν νόμον ἡμῶν καρτερίαν.

\vs{13}Τελευτήσαντος δὲ καὶ τούτου, ὁ ἕκτος ἤγετο μειρακίσκος· ὃς πυνθανομένου τοῦ τύραννου εἰ βούλοιτο φαγὼν ἀπολύεσθαι, ὁ δὲ ἔφη,

\vs{14}Ἐγὼ τῇ μὲν ἡλικίᾳ τῶν ἀδελφῶν μου εἰμὶ νεώτερος, τῇ δὲ διανοίᾳ ἡλικιώτης·
\vs{15}Εἰς τὰ αὐτὰ γὰρ καὶ γεννηθέντες καὶ τραφέντες, ὑπὲρ τῶν αὐτῶν καὶ ἀποθνήσκειν ὀφείλομεν ὁμοίως.
\vs{16}Ὥστε εἰ σοὶ δοκεῖ βασανίξειν, μὴ μιαιροφαγοῦτας βασάνιζε.

\vs{17}Ταῦτα αὐτὸν εἰπόντα παρῆγον ἐπὶ τὸν τροχόν.
\vs{18}Ἐφʼ οὗ κατατεινόμενος εὐμελῶς καὶ ἐκσφονοδυλιζόμενος ὑπεκαίετο.
\vs{19}Καὶ ὀβελίσκους ὀξεῖς πυρώσαντες, τοῖς νότοις προσέφερον· καὶ τὰ πλευπὰ διαπείραντες, ἀπʼ αὐτοῦ σπλάγχνα διέκαιον.

\vs{20}Ὁ δὲ βασανιζόμενος, ὦ ἱεροπρεποῦς αἰῶνος, ἔλεγεν, ἐφʼ ὃν διὰ τὴν εὐσέβειαν εἰς γυμνασίαν πόνων ἀδελφοὶ τοσοῦτοι κληθέντες οὐκ ἐνικήθημεν.
\vs{21}Ἀνίκητος γάρ ἐστιν, ὦ τύραννε, ἡ εὐσεβὴς ἐπιστήμη.
\vs{22}Καλοκᾳγαθίᾳ καθωπλισμένος τεθνήξομαι κἀγὼ μετὰ τῶν ἀδελφῶν μοῦ.
\vs{23}Μέγαν σοὶ προσβάλλων καὶ αὐτὸς ἀλάστορα, καινουργὲ τῶν βασάνων, καὶ πολέμιε τῶν ἀληθῶς εὐσεβούντων.

\vs{24}Ἓξ μειράκια κατελύσαμέν σου τὴν τυραννίδα.
\vs{25}Τὸ γὰρ μὴ δυνηθῆναί σε μεταπεῖσαι τὸν λογισμὸν ἡμῶν, μήτε βιάσασθαι πρὸς τὴν μιαιροφαγίαν, οὐ κατάλυσίς ἐστιν σοῦ;
\vs{26}Τὸ πῦρ σου ψυχρὸν ἡμῖν, καὶ ἄπονοι οἱ καταπέλται, καὶ ἀδύνατος ἡ βία σου.
\vs{27}Οὐ γὰρ τυράννου, ἀλλὰ θείου νόμου προεστήκασιν ἡμῶν οἱ δορυφίροι· διὰ τοῦτο ἀνίκητον ἕχομεν τὸν λογι σμόν.

\ch{12}
Ὡς δὲ καὶ οὗτος μακαρίως ἐναπέθανεν καταβληθεὶς εἰς λέβητα, ὁ ἕβδομος παρεγίνετο, πάντων νεώτερος.
\vs{2}Ὅν κατοικτειρήσας ὁ τύραννος, καίπερ δεινῶς ὑπὸ τῶν ἀδελφῶν αὐτοῦ κακισθεὶς, ὁρῶν ἤση τὰ δεσμὰ περικείμενον,
\vs{3}πλησιέστερον αὐτὸν μετεπέμψατο, καὶ παρηγορεῖν ἐπειρᾶτο, λέγων,

\vs{4}Τῆς μὲν τῶν ἀδελφῶν σου ἀπονοίας τὸ τὲλος ὁρᾷς· διὰ γὰρ ἀπείθειαν στρεβλωθέντες τεθνήκασιν, σὺ, εἰ μὲν μὴ πεισθείης, τὰλας βασανισθεὶς καὶ πασανισθεὶς καὶ αὐτὸς τεθνήξῃ πρὸ ὥρας.
\vs{5}Πεισθείης δὲ φίλος ἔσῃ, καὶ τῶν ἐπὶ τῆς βασιλείας ἀφηγήσῃ πραγμάτων.

\vs{6}Καὶ ταῦτα παρακαλῶν, τὴν μητέρα τοῦ παιδὸς μετεπέμψατο, ὄπως αὐτὴν ἐμεήσαν υἱῶν στερηθῖσαν παρορμήσειεν ἐπὶ τὴν μσωτηρίαν, εὐπειθῆ τὸν περιλειπόμενον.
\vs{7}Ὁ δὲ τῆν μητρὸς τῇ Ἐβραΐδι φωνῇ προτρεψαμένης· αὐτὸν (ὡς ἐροῦμεν μετὰ μικρὸν ὕστερον.) ἀπολύσατε με, φησίν·
\vs{8}εἴπω τῷ βασιλεῖ καὶ τοῖς σὺν αὐτῷ φίλοις πᾶσιν.
\vs{9}Καὶ ἐπιχαρέντες μάλιστα ἐπὶ τῇ ἐπαγγελίᾳ τοῦ παιδὸς ταχέως ἔλυσαν αὐτόν.

\vs{10}Καὶ σραμὼν ἐπὶ πλησίον τῶν τηγάνων, ἔφη, ἀνόσιε,
\vs{11}φησὶν, καὶ πάντων τῶν πονηρῶν ἀσεβέστατε τύραννε, οὐκ ᾐδέσθης παρὰ τοῦ Θεοῦ λαβὼν τὰ ἀγαθὰ καὶ τὴν βασιλείαν, τούς θεράποντας αὐτοῦ κατακτεῖναι, καὶ τοὺς τῆς εὐσεβείας στρεβλῶσαι;
\vs{12}Ἀνθ ὧ ταμιεύεταί σε ἡ θεια δίκη πυκνοτέρῳ καὶ αἰωνίῳ πυεὶ καὶ βασάνοις, αἳ εἰς ὅλον τὸν αιῶνα οὐκ ἀνήσουσίν σε.

\vs{13}Οὐκ ᾐδέσθης ἄνθρωπος ὢν, θηριωδέστατε, τοὺς ὁμοιοπαθεῖς καὶ ἐκ τῶν αὐτῶν γεγονότας στοιξείων γλωττοτομῆσαι, καὶ τοῦτον καταικίσας τὸν τρόπον βασανίσαι;
\vs{14}Ἀλλʼ οἱ μὴν εὐγενῶς ἀποθανόντες ἐπλήρωσαν τὴν εἰν τὸν Θεὸν εὐσέβειαν.
\vs{15}Σὺ δὲ κακὸς κακῶς οἰμώξεις, τοὺς τῆς ἀρετῆν ἀγωνιστὰν ἀναιτίως ἀποκτεῖναι.

\vs{16}Ὅθεν καὶ αὐτὸς ἀποθνήσκειν μέλλων, ἔφη,
\vs{17}οὐκ ἀπαυτομολῶ τῆς τῶν ἀδελφῶν μου μαρτυρίας.
\vs{18}Ἐπικαλοῦμαι δὲ τὸν πατρῷον Θεὸν, ὅπως ἵλεως γένει μου.
\vs{19}Σὲ δὲ καὶ ἐν τῷ νῦν βίῳ καὶ θανόντα τιμωρήσεται.

\vs{20}Καὶ ταῦτα κατευξάμενος, ἐαυτὸν ἔοιψεν κατὰ τῶν τηγάνων· καὶ οὕτως ἀπέδεκεν.

\ch{13}
Εἰ δὲ τοίνυν τῶν μέχρι θανάτου πόνων ὑπερεφρόνησαν οἱ ἑπτὰ ἀδελφοὶ, συνομολογεῖται πανταχόθεν, ὅτι αὐτοδέσποτός ἐστιν τῶν παθῶν ὁ εὐσεβὴς λογισμός.
\vs{2}Ὤσπερ γὰρ εἰ τοῖς πάθεσιν σουλωθέντες ἐμιεροφάγησαν, ἐλέγομεν γὰρ αὐτοὺς τούτοις νενικῆσθαι.
\vs{3}Νυνὶ δὲ οὐχ οὕτως· ἀλλὰ τῷ ἐπαινουμένῳ λογισμῷ παρὰ περιεγένοντο τῶν παθῶν.

\vs{4}Καὶ οὐκ ἐστὶν παριδεῖν τὴν ἡγεμονίαν· ἐπεκρὰτησεν γὰρ καὶ πὰθους καὶ πὸνων.
\vs{5}Πῶς οὖν οὐκ ἐστὶν τούτοις τὴν εὐλογιστίας παθοκρὰτειαν ὁμολογεῖν, οἱ τῶν μὲν διὰ πυρὸς ἀλγησόνων οὐκ ἐπεστράφησαν;
\vs{6}Καθάπερ γὰρ προπλήταις λιμένων πύργοις τὰς κυμάτων ἀπειλὰς ἀνακόπτοντες, γαληνὸν παρὲχουσιν τοῖς εἰσπλέουσιν τὸν ὅρμον.
\vs{7}Οὕτος ἡ ἑπτάπυργος τῶν νεανίσκων εὐλογιστία τὸν τῆς εὐσεβείας ὀχυρώσασα λιμένα τὴν τῶν παθῶν ἐνίκησεν ἀκολασίαν.

\vs{8}Ἱερὸν γὰρ εὐσεβείαν στήσαντες χορὸν παρεθάρσυνον ἀλλήλους, λέγοντες,
\vs{9}ἀδελφικῶς ἀποθάνοιμεν, ἀδελφοὶ, περὶ τοῦ νόμον· μιμησώμεθα τούς τρεῖς τοὺς ἐπὶ Ἀσσυρίας νεανίσκους, οἵ τῆς ἰσεπόλιδος καμίνου κατεφρόνησαν.
\vs{10}Μὴ δειλανδρήσωμεν πρὸς τὴν τῆς εὐσεβείας ἀπόδειξιν.
\vs{11}Καὶ ὁ μὲν, θάῤῥει ἀδελφὲ, ἔλεγεν, ὁ δὲ, εὐγενῶς καρτέρησον.
\vs{12}Ὁ δὲ, ἔλεγεν, μνήσθητε πόθεν ἐστὲ, ἢ τίνος πατρὸς χειρὶ σφαγιασθῆναι διὰ τὴν εὐσέβειαν ὑπέμεινεν ὁ Ἰσαάκ.

\vs{13}Εἶς δὲ ἕκαστος καὶ ἀλλήλους ὁμοῦ πάντες ἐφόρων φαιδροὶ καὶ μάλα θαῤῥαλέοι, ἐαυτοὺς, ἔλεγον, τῷ Θεῷ ἀφιερώσωμεν ἐξ ὅλης τῆς καρδίας τῷ δόντι τὰς ψυχὰς, καὶ χρήσωμεν τῇ περὶ τὸν νόμον φυλακῇ τὰ σώματα.
\vs{14}Μὴ φοβηθῶμεν τὸν δοκοῦντα ἀποκτενεῖν
\vs{15}Μέγας γὰρ ψυξῆς ἀγὼν καὶ κίνδυνος ἐν αἰωνίῳ βασάνῳ κείμενος τοῖς παραβᾶσιν τῆν ἐντολὴν τοῦ Θεοῦ.
\vs{16}Καθοπλισώμεθα τοιγαροῦν τῇ τοῦ θείου λογισμοῦ παθοκρατείᾳ.
\vs{17}Οὕτως παθόντας ἡμᾶς Αβραὰμ καὶ Ἰακὼβ ὑποδέξονται, καὶ πάντες οἱ πατέρες ἐπαινέσουσιν.
\vs{18}Καὶ ἑνὶ ἑκάστῳ τῶν ἀποστωμένων αὐτῶν ἀδελφῶν ἔλεγον οἱ περιλειπόμενοι, μὴ καταισχύνῃς ἡμᾶς ἀδελφὲ, μηδὲ ἀδελφὲ, μηδὲ ψεύσῃ τούς προαποθανόντας.

\vs{19}Οὐκ ἀγνοεῖτε δὲ τὰ τῆς ἀνθρωπότητος φίλτρα, ἅπερ ἡ θεία καὶ πάνσοφος πρόνοια διὰ τῆς μητρῴας φυτεύσασα γαστρός·
\vs{20}ἐν ᾗ τὸν ἶσον ἀδελφοὶ κατοικήσαντες χρόνον, καὶ ἐν τῷ αὐτῳ χρόνῳ πλασθέντες, καὶ ἀπὸ τοῦ αὐτῷ αἵματος ἀξηθέντες, καὶ δια τῆς αὐτῆς ψυχῆς τελεσφορηθέντες,
\vs{21}καὶ διὰ τῶν ἴσων ἀποτεχθέντες χρόνον, καὶ ἀπὸ τῶν αὐτῶν γαλακτοποτοῦντες πηγῶν, ἀφʼ οὗ συντέφονται ἐν ἐναγκαλισμάτων φιλάδελφοι ψυχαί·
\vs{22}καὶ αὔξοντες σφοδρότερον διὰ συντροφίας, καὶ τῆς καθʼ ἠμέραν συνηθείας, καὶ τῆς ἄλλης παιδείας, καὶ τῆν ἡμετέρας ἐν νόμῳ Θεοῦ ἀσκήσεως.

\vs{23}Οὕτως δὲ τοίνυν καθεστηκυίας τῆς φιλαδελφίας συμπαθούσης, οἱ ἑπτὰ ἀδελφοὶ συμπαθέστερον ἔσχον τὴν πρὸς ἀλλήλους ὁμόνοιαν.
\vs{24}Νόμῳ γὰρ τῷ αὐτῷ παιδευθέντες, καὶ τὰς αὐτὰς ἐξασκήσαντες ἀρετὰς, καὶ τῷ δικαίῳ συντραφέντες βίῳ, μᾶλλον ἐπʼ αὐτοὺς ἥγαγον.
\vs{25}Ἡ γὰρ ὁμοζηλία τῆς καλοκᾳγαθίας ἐπέτεινεν αὐτῶν τὴν πρὸς ἀλλήλους ὁμόνοιαν.
\vs{26}Σὺν γάρ τῇ εὐσεβείᾳ ποθεινοτέραν αὐτοῖς κατεσκεύαζεν τὴν φιλαδελφίαν.

\vs{27}Ἀλλʼ ὁμοίως καίπερ τῆς φύσεως καὶ τῆς συνηθείας καὶ τῶν τῆς ἀρετῆς ἠθῶν τὰ τῆς ἀδελφότητος αὐτοῖν φίλτρα συναυξόντων, ἀνέσχοντο διὰ τὴν εὐσέβειαν τοὺν ἀδελφοὺς οἱ ὑπολελειμμένοι τοὺς καταικιζομένους, ὁρῶντες μέχρι θανάτου βασανιζομένους.

\ch{14}
Προσέτι καὶ ἐπὶ τὸν αἰκισμὸν ἐποτρύνοντες, ὡς μὴ μόνον τῶν ἀλγηδόνων περιφρονῆσαι αὐτοὺς, ἀλλὰ καὶ τῆς τῶν ἀδελφῶν φιλαδελφίας παθῶν κρατῆσαι.

\vs{2}Ὦ βασιλέως λογισμοὶ βασιλικώτεροι καὶ ἐλευθέρων ἐλευθερώτεροι.
\vs{3}Ἱερὰς καὶ ἐναρμόστους περὶ τῆς εὐσεβείας τῶν ἑπτὰ ἀδελφῶν συμφωνίας.
\vs{4}Οὐδεὶς ἐκ τῶν ἑπτὰ μειρακίων ἐδειλίασεν, οὐδὲ πρὸς τὸν θάνατον ὤκνησεν.
\vs{5}Ἀλλὰ πάντες, ὥσπερ ἐπʼ ἀθανασίας ὁδὸν τρέχοντες, ἐπὶ τὸν διὰ τῶν βασάνων θάνατον ἔσπευδον.
\vs{6}Καθάπερ γὰρ χεῖρες καὶ πόδες συμφώνως τοῖς τῆς ψυχῆς ἀφηγήμασιν κινοῦνται· οὕτως οἱ ἱεροὶ μείρακες ἐκεῖνοι ὡς ὑπὸ ψυχῆς ἀθανάτου τῆς εὐσεβείας, πρὸς τὸν ὑπὲρ αὐτῆς συνεφώνησαν θάνατον.

\vs{7}Ὦ παναγία ἡ συμφώνον ἀδελφῶν ἐβδομάς· καθάπερ γὰρ ἑπτὰ τῆς κοσμοποιΐας ἡμέραι περὶ τὴν εὐσέβειαν,
\vs{8}οὕτος περὶ τὴν ἑβδομάδα χορεύοντες οἱ μείρακες ἐκύκλουν τὸν τῶν βασάνων φόβον καταλύοντες.
\vs{9}Νῦν ἡμεῖς ἀκούοντες τῆν θλίψιν τῶν νεανίων ἐκείνων, φρίττομεν· οἱ δὲ οὐ μόνον ὁρῶντες, ἀλλʼ οὐδὲ μόνον ἀκούοντες τὸν παραχρῆμα ἀπειλῆς λόγον, ἀλλὰ καὶ πάσχοντες, ἐκαρτέρουν καὶ τοῦτο ταῖς διὰ πυρὸς ὀδύναις.
\vs{10}Ὧν τί γένοιτο ἐπαλγέστερον; ὀξεῖα γὰρ καὶ σύντομος ἡ τοῦ πυρὸς οὖσα δύναμις, ταχέως διέλυσε τὰ σώματα.

\vs{11}Καὶ μὴ θαυμαστὸν ἡγεῖσθε, εἰ ὁ λογισμὸς περιεκράτησεν τῶν ἀνδρῶν ἐκείνων ἐν ταῖς βασάνοις, ὅπου γε καὶ γυναικὸς νοῦς πολυτροπωτέρον ὑπερεφρόνησεν ἀλγηδόνων.
\vs{12}Ἡ μήτηρ γὰρ τῶν ἑπτὰ νεανίσκων ὑπήνεγκεν τὰς ἐφʼ ἑνὶ ἐκάστῳ τῶν τέκνων στρέβλας.

\vs{13}Θεωρεῖτε δὲ πῶς πολύπλοκός ἐστιν ἡ τῆς φιλοτεκνίας στοργὴ, ἕλκουσα πάντα πρὸς τὴν τῶν σπλάγχνων συμπάθειαν.
\vs{14}Ὅπου γε καὶ τὰ ἄλογα ζῶα ὁμοίαν τὴν πρὸς τὰ ἐξ αὐτῶν γεννώμενα συμπάθειαν καὶ στοργὴν ἔχει τοῖς ἀνθρώποις.
\vs{15}Καὶ γὰρ τῶν πετεινῶν, τὰ μὲν ἥμερα κατὰ τὰς οἰκίας ὀροφοιτοῦντα προασπίζει τῶν νεοττῶν.
\vs{16}Τὰ δὲ κατὰ τὰς κορυφὰς ὀρέων καὶ φαράγγων ἀποῤῥῶγας καὶ δένδρων ὀπὰς καὶ τὰς τούτων ἄκρας νοσσοποιησάμενα ἀποτίκτει, καὶ τὸν προσιόντα κωλύει.
\vs{17}Εἰ δὲ καὶ μὴ δύναιντο κωλύειν, περιπτάμενα κυκλόθεν αὐτῶν ἀλγοῦντα τῇ στοργῇ, ἀνακαλούμενα τῇ ἰδίᾳ φωνῇ, καθʼ ὃν δύναται τρόπον βοηθεῖ τοῖς τέκνοις.

\vs{18}Καὶ τί δεῖ τὴν διὰ τῶν ἀλόγων ζώων ἐπιδεικνύναι τὴν πρὸς τὰ τέκνα συμπάθειαν.
\vs{19}Ὅπου γε καὶ μέλισσαι περὶ τὸν τῆς κηρογονίας καιρὸν ἐπαμύνονται τοὺς προσιόντας, καὶ καθάπερ σιδήρῳ τῷ κέντρῳ πλήσσουσι τοὺς προσιόντας τῇ νοσσιᾷ αὐτῶν, καὶ ἐπαμύνονται ἕως θανάτου.

\vs{20}Ἀλλʼ οὐχὶ τὴν Ἁβραὰμ ὁμόψυχον τῶν νεανίων μητέρα μετεκίνησεν συμπάθεια τῆς συμπαθείας τέκνων.

\ch{15}
Ὦ λογίσμε τέκνων, παθῶν τύραννε, καὶ εὐσέβεια μητρὶ τέκνων ποθεινοτέρα.
\vs{2}Μήτηρ δυοῖν προκειμένων εὐσεβείας, καὶ τῆς ἑπτὰ υἱῶν σωτηρίας προκαίρους κατὰ τὴν τοῦ τυράννου ὑπόσχεσιν·
\vs{3}τὴν εὐσέβειαν μᾶλλον ἠγάπησεν τὴν σώζουσαν εἰς αἰώνιον ζωὴν κατὰ Θεόν.

\vs{4}Ὦ τίνα τρόπον ἠθολογήσαιμι φιλότεκνα γονέων πάθη, ψυχῆς τε καὶ μορφῆς ὁμοιότητα εἰς μικρὸν παιδὸς χαρακτῆρα θαυμάσιον ἐναπεσφράγιζον, μάλιστα διὰ τὸν τῶν παθῶν τοῖς γεννηθεῖσιν τὰς μητέρας καθεστάναι συμπαθευτέρας.
\vs{5}Ὅσῳ γὰρ καὶ ἀσθενόψυχοι καὶ πολυγονώτεραι ὑπάρχουσιν μητέρες, τοσούτῳ μᾶλλόν εἰσιν φιλοτεκνότεραι.
\vs{6}Πασῶν δὲ τῶν μητέρων ἐγένετο ἡ τῶν ἑπτὰ μήτηρ φιλοτεκνοτέρα, ἥ τις ἑπτὰ κυοφορίαις τὴν πρὸς αὐτοὺς ἐπιφυτευομένη φιλοστοργία,
\vs{7}καὶ διὰ πολλὰς τὰς καθʼ ἔκαστον αὐτῶν ὠδῖνας ἠναγκασμένην τὴν εἰς αὐτοὺς ἔχειν συμπάθειαν,
\vs{8}διὰ τὸν πρὸς τὸν Θεὸν φόβον ὑπερεῖδεν τὴν τῶν τέκνων πρόσκαιρον σωτηρίαν.

\vs{9}Οὐ μὴν δὲ, ἀλλὰ καὶ διὰ τὴν καλοκᾳγαθίαν τῶν υἱῶν, καὶ τὴν πρὸς τὸν νόμον αὐτῶν εὐπείθειαν, μείζων τὴν ἐν αὐτοῖς ἔσχεν φιλοστοργίαν.
\vs{10}Δίκαιοί τε γὰρ ἦσαν, καὶ σώφρονες, καὶ σώφρονες, καὶ ἀνδρεῖοι, καὶ μεγαλόψυχοι, καὶ φιλάδελφοι, καὶ μεγαλόψυχοι, καὶ μεψαλόψυχοι, καὶ φιλάδελφοι, καὶ φιλομήτορες οὕτως, ὥστε καὶ μέχρι θανάτου τὰ νόμιμα φυλάσσοντες πείθεσθαι αὐτῇ.

\vs{11}Ἀλλʼ ὅμως, καὶ ὑπὲρ τοσούτων ὄντων τῶν περὶ φιλοτεκνίαν εἰς συμπάθειαν ἑλκόντων τὴν μητέρα, ἐπʼ οὐδενὸς αὐτῶν τὸν λογισμὸν αὐτῆς αἱ παμποίκιλοι ἴσχυσαν μετατρέψαι.
\vs{12}Ἀλλὰ καἰ καθʼ ἔνα παῖδα καὶ ὁμοῦ πάντας ἡ μήτηρ ἐπὶ τὸν τῆς εὐσεβείας προετρέπετο θάνατον.
\vs{13}Ὦ φύσις ἱερὰ, καὶ φίλτρα γονέων καὶ γονεῦσιν φιλόστοργε, καὶ τροφεῖα, καὶ μητέρων ἀδάμαστα πάθη.

\vs{14}Καθʼ ἕνα στρεβλούμενον καὶ φλεγόμενον ὁρῶσα υήτηρ, οὐ μετεβάλετο διὰ τὴν εὐσέβεβειαν.
\vs{15}Τὰς σάρκας τῶν τέκνων ἑώρα περὶ τὸ πῦρ τηκομένας, καὶ τοὺς τὼν ποδῶν καὶ χειρῶν δακτύλους ἐπὶ γῆς σπαίροντας, καὶ τὰς τῶν κεφαλῶν μέχρι τῶν περὶ τὰ γένεια σάρκας, ὥσπερ προσωπεῖα προκειμένας.

\vs{16}Ὦ πικροτέρων μὲν νῦν μήτηρ πόνων πειρασθεῖσα, ἤπερ τῶν ἐπʼ αὐτοῖς ὠδίνων.
\vs{17}Ὦ μόνη γυνὴ τὴν εὐσέβειαν ὁλόκληρον ἀποκυήσασα.
\vs{18}Οὐ μετέρεψέν σε πρωτότοκος ἀποπνέων· οὐδὲ δεύτερον εἰς οἶκτρον βλέπων ἐν βασάνοις· οὐδὲ τρίτος ἀποψύχων.
\vs{19}Οὐδὲ τοὺς ὀφθαλμοὺς ἑνὸς ἑκάστου θεωροῦσα ταυρηδὸν ἐπὶ τῶν βασάνων ὁρῶντας τὸν αὐτὸν αἰκισμὸν, καὶ τοὺς μυκτῆρας προσημειουμένους αὐτῶν τὸν θάνατον, οὐκ ἔκλαυσας.
\vs{20}Ἐπὶ σαρξὶν τέκνων ὁρῶσα σάρκας τέκνων ἀποκεκομμένας, καὶ ἐπὶ κεφαλαῖς κεφαλὰς ἀποδειροτομουμένας, καὶ ἐπὶ νεκροῖς νεκροὺς πίπτοντας, καὶ πολυάνδριον ὁρῶσα τῶν τέκνων χορεῖον διὰ τῶν βασάνων, οὐκ ἐδάκρυσας.

\vs{21}Οὐχ οὕτως σειρήνιοι μελῳδίαι, οὐδὲ κύκνειοι πρὸς φιληκοΐαν φωναὶ τοὺς ἀκούοντας ἐφέλκονται, ὦ τέκνων φωναὶ μετὰ βασάνων μητέρα φωνούντων.
\vs{22}Πηλίκαις καὶ πόσαις τότε ἡ μήτηρ, τῶν υἱῶν βασανιζομένων τροχοῖς τε καὶ καυτερίοις ἐβασανίζετο βασάνοις;

\vs{23}Ἀλλὰ τὰ σπλάγχνα αὐτῆς ὁ εὐσεβὴς λογισμὸς ἐν αὐτοῖς τοῖς πάθεσιν ἀνδρειώσας ἐπέτεινεν τὴν πρόσκαιρον φιλοτεκνίαν παριδεῖν.
\vs{24}Καίπερ ἑπτὰ τέκνων ὁρῶσα ἀπώλειαν· ἀσπάσασα ἡ γενναῖα μήτηρ ἐξέδσεν διὰ τὴν πρὸς Θεὸν πίστιν.
\vs{25}Καθάπερ γὰρ ἐν βουλευτηρίῳ τῇ ἑαυτῆς ψυχῇ δεινοὺς ὁρῶσα συμβούλους, φύσιν καὶ γένεσιν καὶ φιλοτεκνίαν καὶ τέκνων στρέβλαν.
\vs{26}Δύο ψήφους κρατοῦσα μήτηρ, θανατηφόρον τε καὶ σωτήριον ὑπὲρ τένων·
\vs{27}Οὐκ ἐπέγνω τὴν σώζουσαν ἑπτὰ υἱοὺς πρὸς ὀλίγον χρόνον σωτηρίαν.
\vs{28}Ἀλλὰ τῆς θεοσεβοῦς Ἁβραὰμ καρτερίας ἡ θυγάτηρ ἐμνήσθη.

\vs{29}Ὦ μήτηρ ἔθνους, ἔκδικε τοῦ νόμου, καὶ ὑπερασπίστεια τῆς εὐσεβείας, καὶ τοῦ διὰ σπλάγχνων ἀγῶνος ἀθλοφόρε.
\vs{30}Ὦ ἀῤῥένων πρὸς καρτερίαν γενναιοτέρα, καὶ ἀνδρῶν πρὸς ὑπομονὴν ἀνδρειοτέρα.
\vs{31}Καθάπερ γὰρ ἡ Νῶε κιβωτὸς ἐν τῷ κοσμοπληθεῖ κατακλυσμῷ κοσμοφοροῦσα καρτεροὺς ὑπήνεγκεν τοὺς κλύδωνας·
\vs{32}οὕτως σὺ, ἡ νομοφύλαξ, πανταχόθεν ἐν τῷ τῶν παθῶν περιαντλουμένη κατακλυσμῷ, καὶ καρτεροῖς ἂν λοιμοῖς ταῖς τῶν υἱῶν βασάνοις συνεχομένη, γενναίως ὑπέμεινας τοὺς τῆς εὐσεβείας χειμῶνας.

\ch{16}
Εἰ δὲ τοίνυν καὶ γυνὴ, καὶ γηραιὰ, καὶ ἑπτὰ παὶδων μήτηρ ὑπέμεινε τὰς μέχρι θανάτου βασάνους ὁρῶσα τῶν τέκνων· ὁμολογουμένως αὐτοκράτωρ ἐστὶν τῶν παθῶν ὁ εὐσεβὴς λογισμός.

\vs{2}Ἀπέδειξα οὖν ὅτι οὐ μόνον τῶν παθῶν ἄνδρες ἐπεκράτησαν, ἀλλὰ καὶ γυνὴ τῶν μεγίστων βασάνων ὑπερεφρόνησεν.
\vs{3}Καὶ οὐχ οὕτως οἱ περὶ Δανιὴλ λέοντες ἦσαν ἄγριοι, οὐδὲ Μισαὴλ ἐκφλεγομένη κάμινος λαβροτάτῳ πυρὶ, ὡς τῆς φιλοτεκνίας περιέκαιεν ἐκείνη φύσις, ὁρῶσα αὑτῆς τοὺς ἑπτὰ υἱοὺς βασανιζομένους.
\vs{4}Ἀλλὰ τῷ λογισμῷ τῆς εὐσεβείας κατέσβεσε τοσαῦτα καὶ τηλικαῦτα πάθη ἡ μήτηρ.

\vs{5}Καὶ γὰρ τοῦτο ἐπιλογίσασθαι, ὅτι εἰ δειλόψυχος ἦν ἡ γυνὴ, καίπερ μήτηρ οὖσα, ὠλοφύρετο ἂν ἐπʼ αὐτοῖς· καὶ ἴσως ἄν ταῦτα οὕτως εἶπεν,

\vs{6}Ὦ μελέα ἔγωγε, καὶ πολλάκις τρισαθλία, ἥτις ἑπτὰ παῖδας τεκοῦσα, οὐδενὸς μήτηρ γεγένημαι.
\vs{7}Ὦ μάταιοι ἐπτὰ κυοφορίαι, καὶ ἀνόνητοι ἐπτὰ δεκάηνοι, καὶ ἄκαρποι τιθηνίαι, καὶ ταλαίπωροι γαλακτοτροφίοι.
\vs{8}Μάτην ἐφʼ ὑμῖν, ὦ παῖδες, πολλὰς ὑπέμεινα ὠδῖνας καὶ χαλεπωτέρας φροντίδας ἀνατροφῆς.
\vs{9}Ὦ τῶν ἐμῶν παίδων, οἱ μὲν ἄγαμοι, οἱ δὲ γαμήσαντες ἀνόνητοι, οὐκ ὄψομαι ὑμῶν τέκνα, οὐδὲ μάμμη κληθεῖσα μακαρισθήσομαι.
\vs{10}Ὦ ἡ πολύπαις καὶ καλλίπαις ἐγὼ γυνὴ χήρα καὶ μόνη πολύθρηνος.
\vs{11}Οὐδʼ ἂν ἀποθάνω, θάπτοντα τῶν υἱῶν ἕξω τινά.

\vs{12}Ἀλλὰ τούτῳ τῷ θρήνῳ οὐδένα ὠλοφύρετο ἡ ἱερὰ καὶ θεοσεβὴς μήτηρ. Οὐδʼ ἵνα μὴ ἀποθάνωσιν ἀπέτρεπεν αὐτῶν τινα, οὐδʼ ὡς ἀποθνησκόντων ἐλυπήθη.
\vs{13}Ἀλλʼ ὥσπερ ἀδαμάντινον ἔχουσα τὸν νοῦν, καὶ εἰς ἀθανασίαν ἀνατίκτουσα τὸν τῶν υἱῶν ἀριθμὸν, μᾶλλον ὑπὲρ τῆς εὐσεβείας ἐπὶ τὸν θάνατον αὐτοὺς προετρέπετο ἱκετεύουσα.

\vs{14}Ὦ διʼ εὐσέβειαν Θεοῦ στρατιῶτι, πρεσβύτι καὶ γυνὴ διὰ καρτερίαν καὶ τύραννον ἐνίκησας, καὶ ἔργοις δυνατωτέρα καὶ λόγοις εὑρέθης ἄνανδρος.
\vs{15}Καὶ γὰν ὃτε συνελήφθης μετὰ τῶν παίδων, εἱστήκεις τὸν Ἐλεάζαρον ὁρῶσα βασανιζόμενον, καὶ ἔλεγες τοῖς παισὶν ἐν τῇ Ἑβοαΐδι φωνῇ,

\vs{16}Ὦ παῖδες, γενναῖος ὁ ἀγών· εφ ὃν κληθέντες ὑπὲρ τῆς διαμαρτυρίας τοῦ ἔθνους, ἐναγωνίσασθε προθύμως ὑπὲρ τοῦ πατρίουνόμου.
\vs{17}Καὶ γὰρ αἰσχρὸν τὸν μὴν γέροντα τοῦτον ὑπομένειν τὰς διὰ τὴν εὐσέβειαν ἀληδόνας, ὑμᾶς δὲ τοὺς νεωτέρους καταπλαγῆναι τὰς βασάνους.

\vs{18}Ἀναμνήσθητε, ὅτι διὰ τὸν Θεὸν τοῦ κόσμου μετελάβετε, καὶ τοῦ βίου ἀπελαύσατέ·
\vs{19}καὶ διὰ τιῦτο ὀφείλετε πάντα πόνον ὑπομένειν διὰ τὸν Θεόν.
\vs{20}Διʼ ὃν καὶ ὁ πατὴρ ἡμῶν Ἁβραὰμ ἔσπευδεν τὸν ἐθνοπάτορα υἱὸν σφαγιάσαι Ἰσαὰκ, καὶ τὴν πατρῷαν χεῖρα ξιφηφόρον καταφερομένην ἐπʼ αὐτὸν ὁρῶν οὐκ ἔπτηξεν.
\vs{21}Καὶ Δανιὴλ ὁ δίκαιος εἰς λέοντας ἐβλήθη· καὶ Ἀνανίας, καὶ Ἀζαρίας, καὶ Μισαὴλ εἰς κάμινον πυρὸς ἀπεσφενδονήθησαν, καὶ ὑπέμειναν, διὰ τὸν Θεόν.
\vs{22}Καὶ ὑμεῖς οὖν τὴν αὐτὴν πίστιν πρὸς τὸν Θεὸν ἔχοντες, μὴ χαλεπαίνητε.
\vs{23}Ἀλόγιστον γὰρ εἰδότας εὐσέβειαν μὴ ἀντιστασθαι τοῖς πόνοις.

\vs{24}Διὰ τούτων τῶν λόγων ἡ ἑπταμήτωρ ἕνα ἕκαστον τῶν υἱῶν παρακαλοῦσα, ἔπεισε μᾶλλον, ἢ παραβῆναι τὴν ἐντολὴν τοῦ Θεοῦ.
\vs{25}Ἔτι δὲ καὶ ταῦτα ἰδόντες, ὅτι διὰ τὸν Θεὸν ἀποθανόντες ζῶσιν τῷ Θεῷ, ὥσπερ Ἁβραὰμ καὶ Ἰσαὰκ καὶ Ἰακὼβ, καὶ πάντες οἱ πατριάρχαι.

\ch{17}
Ἔλεγον δὲ καὶ τῶν δορυφόρων τινὲς, ὡς ὅτε ἔμελλεν καὶ αὐτὴ συλλαμβάνεσθαι πρὸς θάνατον, ἵνα μὴ ψαύσειέν τι τοῦ σώματος ἑαυτῆς, ἑαυτὴν ἔῤῥιψεν κατὰ τῆς πυρᾶς.

\vs{2}Ὦ μήτηρ σὺν ἑπτὰ παισὶν καταλύσασα τὴν τοῦ τυράννου βίαν, καὶ ἀκυρώσασα τὰς κακὰς ἐπινοίας αὐτοῦ, καὶ ἐπιδείξασα τὴν τῆς πίστεως γενναιότητα.
\vs{3}Καθάπερ γὰρ σὺ στέγη ἐπὶ τοῦ στύλου τῶν παίδων γενναίως ἱδρυμένη, ἀκλινῶς ὑπήνεγκας τὸν διὰ τῶν βασάνων σεισμόν.

\vs{4}Θάῤῥει τοιγαροῦν, ὦ μήτηρ ἱερόψυχε, τὴν ἐλπίδα τῆς ὑπομονῆς γενναίως ἔχουσα πρὸς Θεόν.
\vs{5}Οὐχ οὕτω σελήνη κατʼ οὐρανὸν σὺν ἄστροις σεμνὴ καθέστηκεν, ὡς σὺ τοὺς εἰς ἀστέρας ἑπτὰ παῖδας φωταγωγήσασα πρὸς τὴν εὐσέβειαν ἔντιμος καθέστηκας Θεῷ, καὶ ἐστήρισαι ἐν οὐρανῷ σὺν αὐτοῖς.
\vs{6}Ἦν γὰρ ἡ παιδοποιΐα σου ἀπὸ Ἁβραὰμ τοῦ παιδός.

\vs{7}Εἰ δὲ ἐξὸν ἡμῖν ἦν, ὥσπερ τινὸς ζωγραφῆσαι τὴν τῆς ἱστορίας σου εὐσέβειαν, οὐκ ἂν ἔφριττον οἱ θεωροῦντες μητέρα ἑπτὰ τέκνων διʼ εὐσέβειαν ποικίλας βασάνους μέχρι θανάτου ὑπομείνασαν.
\vs{8}Καὶ γὰρ ἄξιον ἦν καὶ ἐπὶ αὐτοῦ τοῦ ἐπιταφίου ἀναγράψαι καὶ ταῦτα τοῖς ἀπὸ τοῦ ἔθνους εἰς μνείαν λεγόμενα.
\vs{9}Ἐνταῦθα γέρων ἱερεὺς, καὶ γυμὴ γεραιὰ καὶ ἑπτὰ παῖδες ἐγκεκήδευνται διὰ τυράννου βίαν, τὴν Ἑβραίων πολιτείαν καταλῦσαι θέλοντος.
\vs{10}Οἳ καὶ ἐξεδίκησαν τὸ ἔθνος εἰς Θεὸν ἀφορῶντες, καὶ μέχρι θανάτου τὰς βασάνους ὑπομείναντες.

\vs{11}Ἀληθῶς γὰρ ἦν ἀγῶν θεῖος ὁ διʼ αὐτῶν γεγενημένος.
\vs{12}Ἠθλότει γὰρ τότε ἀρετὴ διʼ ὑπομονῆς δοκιμάζουσα τὸ νῖκος ἐν ἀφθαρσίᾳ ἐν ζωῇ πολυχρονίῳ.
\vs{13}Ἐλεάζαρ δὲ προηγωνίζετο· ἡ δὲ μήτηρ τῶν ἑπτὰ παίδων ἐνήθλει·
\vs{14}οἱ δὲ ἀδελφοὶ ἠγωνίζοντο· ὁ τύραννος ἀντηγωνίζετο· ὁ δὲ κόσμος καὶ ὁ τῶν ἀνθρώπων βίος ἐθεώρει.
\vs{15}Θεοσέβεια δὲ ἐνίκα, τοὺς ἑαυτῆς ἀθλητὰς στεφανοῦσα.

\vs{16}Τίνες οὐκ ἐθαύμασαν τοὺς τῆς ἀληθείας νομοθεσίας ἀθλητὰς; τίνες οὐκ ἐξεπλάγησαν;
\vs{17}Αὐτός γέ τοι ὁ τύραννος καὶ ὅλον τὸν συνέδριον αὐτῶν ἐξεθαύμασαν αὐτῶν τὴν ὑπομονήν.
\vs{18}Διʼ ἣν καὶ τῷ θείῳ νῦν παρεστήκασιν θρόνῳ, καὶ τὸν μακάριον βιοῦσιν αἰῶνα.
\vs{19}Καὶ γάρ φησιν ὁ Μωσῆς, καὶ πάντες οἱ ἡγιασμένοι ὑπὸ τὰς χεῖράς σου.

\vs{20}Καὶ οὗτοι οὖν ἁγιασθέντες διὰ Θεὸν τετίμηνται οὐ μόνον οὖν ταύτῃ τῇ τιμῇ, ἀλλὰ καὶ τῷ διʼ αὐτοὺς τὸ ἔθνος ἡμῶν τοὺς πολεμίους μὴ ἐπικρατήσας,
\vs{21}καὶ τὸν τύραννον τιμωρηθῆναι, καὶ τὴν πατρίδα καθαρισθῆναι,
\vs{22}ὥσπωρ ἀντίψυχον γεγονότας τῆς τοῦ ἔθνους ἁμαρτίας, καὶ διὰ τοῦ αἵματος τῶν εὐσεβῶν ἐκείνων, καὶ τοῦ ἱλαστηρίου θανάτου αὐτῶν, ἡ θεία πρόνοια τὸν Ἰσραὴλ προκακωθέντα διέσωσεν.

\vs{23}Πρὸς γὰρ τὴν ἀνδρείαν αὐτῶν τῆς ἀρετῆς, καὶ τὴν ἐπὶ ταῖς βασάνοις αὐτῶν ὑπομονὴν ὁ τύραννος ἀφιδὼν Ἀντίοχος ἀνεκήρυξεν τοῖς στρατιώταις αὐτοῦ εἰς ὑπόδειγμα τὴν ἐκείνων ὑπομονήν.
\vs{24}Ἔσχεν τε αὐτοὺς γενναίους καὶ ἀνδρείους εἰς πεζομαχίαν καὶ πολιορκίαν· καὶ ἐκπορθήσας ἐνίκησεν πάντας τοὺς πολεμίους.

\ch{18}
Ὦ τῶν Ἁβραμιαίων σπερμάτων ἀπόγονοι παῖδες Ἰσραηλῖται, πείθεσθε τῷ νόμῳ τούτῳ, καὶ πάντα τρόπον εὐσεβεῖτε·
\vs{2}γινώσκοντες, ὅτι τῶν παθῶν δεσπότης ἐστὶν ὁ εὐσεβὴς λογισμός· καὶ οὐ μόνον τῶν ἔνδοθεν, ἀλλὰ καὶ τῶν ἔξωθεν πόνων·

\vs{3}Ἀνθʼ ὦν διὰ τὴν εὐσέβειαν προϊέμενοι τὰ σώματα τοῖς πόνοις ἐκεῖνοι, οὐ μόνον ὑπὸ τῶν ἀνθρώπων ἐθαυμάσθησαν, ἀλλὰ καὶ θείας μερίδος κατηξιώθησαν.
\vs{4}Καὶ διʼ αὐτοὺς εἰρήνευσεν τὸ ἔθνος, καὶ τὴν εὐνομίαν τὴν ἐπὶ τῆς πατρίδος ἀνανεωσάμενος, ἐκπεπολιόρκηκε τοὺς πολεμίους.
\vs{5}Καὶ ὁ τύραννος Ἀντίοχος καὶ ἐπὶ γῆς τετιμώρηται, καὶ ἀποθανὼν κολάζεται· ὡς γὰρ οὐδὲν οὐδαμῶς ἴσχυσεν ἀναγκάσαι τοὺς Ἱεροσολυμίτας ἀλλοφυλῆσαι, καὶ τῶν πατριῶν ἐθνῶν ἐκδιαιτηθῆναι·
\vs{6}τότε δὴ ἀπάρας ἀπὸ τῶν Ἱεροσολύμων ἐστρατοπέδευσεν ἐπὶ Πέρσας.

\vs{7}Ἔλεγεν δὲ ἡ μήτηρ τῶν ἑπτὰ παίδων καὶ ταῦτα ἡ δικαία τοῖς τέκνοις, ὅτι ἐγὼ ἐγενήθην παρθένος ἁγνὴ, καὶ οὐχ ὑπερέβην πατρικὸν οἶκον· ἐφύλασσον δὲ τὴν ᾠκοδομουμένην πλευράν.
\vs{8}Οὐ διέφθειρέν με λυμεὼν τῆς ἐρημίας φθορεὺς ἐν πεδίῳ· οὐδὲ ἐλυμῄνατό μου τὰ ἁγνὰ τῆς παρθενίας λυμεὼν ἀπατηλὸς ὄφις· ἔμεινα δὲ χρόνον ἀκμῆς σὺν ἀνδρί.

\vs{9}Τούτων δὲ ἐνελίκων γενομένων ἐτελεύτησεν ὁ πατήρ· μακάριος μὲν ἐκεῖνος· τὸν γὰρ τῆς εὐτεκνίας βίον ἐπιζητήσας, τὸν τῆς ἀτεκνίας οὐκ ὠδυνήθη καιρόν.
\vs{10}Ὃς ἐδιδασκεν ὑμᾶς, ἔτι ὢν σὺν ὑμῖν, τὸν νόμον καὶ τοὺς προφήτας.

\vs{11}Τὸν ἀναιρεθέντα Ἀβὲλ ὑπὸ Κάϊν ἀνεγίνωσκεν δὲ ἡμῖν, καὶ τὸν ὁλοκαπούμενον Ἰσαὰκ, καὶ τὸν ἐν φυλακῇ Ἰωσήφ.
\vs{12}Ἔλεγεν δὲ ἡμῖν τὸν ζηλωτὴν Φινεές· ἐδίδασκεν δὲ ὑμᾶς τοὺς ἐν πυρὶ Ανανίαν, καὶ Ἀζαρίαν, καὶ Μισαήλ.
\vs{13}Ἐδόξαζεν δὲ καὶ τὸν ἐν λάκκῳ λεόντων Δανιὴλ, ὃν καὶ ἐμακάριζεν.

\vs{14}Ὑπεμίμνησκεν δὲ ὑμᾶς τῆν Ἠσαΐου γραφὴν τὴν λέγουσαν, κᾂν διὰ πυρὸς διέλθῃς, φλὸξ οὐ κατακαύσει σε.
\vs{15}Τὸν ὑμνογράφον ἐμελῴδει ὑμῖν Δαρίδ τὸν λέγοντα, πολλαὶ αἱ θλίψεις τῶν δικαίων.
\vs{16}Τὸν Σαλομῶντα ἐπαροιμίαζεν ἡμῖν τὸν λέγοντα, ξύλον ζωῆς ἐστιν πᾶσιν τοῖς ποιοῦσιν αὐτοῦ τὸ θέλημα.
\vs{17}Τὸν Ἰεζεκιὴλ ἐπιστοποιεῖτο τὸν λέγοντα, εἰ ζήσεται τὰ ὀστᾶ τὰ ξηρὰ ταῦτα;
\vs{18}Ὧδην μὲν γὰρ ἣν ἑδίδαξεν Μωϋσῆς οὐκ ἐπελάθετο τὴν διδάσκουσαν, ἐγὼ ἀποκτενῶ καὶ ζῇν ποιήσω.
\vs{19}Αὗτη ἡ ζωὴ ἡμῶν καὶ ἡ μακαριότης τῶν ἡμερῶν.

\vs{20}Ὧ πικρᾶς τῆς τότε ἡμέρας, καὶ οὐ πικρᾶς, ὅτε ὁ πικρὸς Ἑλλήμων τύραννος πῦρ φλέξας λέβησιν ὠμοῖς, καὶ ζεουσι θυμοῖς ἀγαγὼν ἐπὶ τὸν καταπέλτην καὶ πάλιν τὰς βασάνους αὐτοῦ τοὺς ἑπτὰ παῖδας τῆς Ἁβρααμίτιδος.
\vs{21}Τὰς τῶν ὀμμάτων κόρας ἐπήρωσεν, καὶ γλώσσας ἐξέτεμεν, καὶ βασάνοις ποικίλαις ἀπέκτεινεν.
\vs{22}Ὑπὲρ ὧν ἡ θεία δίκη μετῆλθεν καὶ μετελεύσεται τὸν ἀλάστορα.

\vs{23}Οἱ δὲ Ἁβραμιαῖοι παῖδες σὺν τῇ ἀθλοφορῳ μητρὶ, εἰς πατέρων χορὸν συναγελάζονται, φυχὰς καὶ ἀθανάτους ἀπειληφότες παρὰ τοῦ Θεοῦ.
\vs{24}Ὦ ἡ δόξα εἰς τοὺς αἰῶνας τῶν αἰώνων. Ἀμήν.


\def\book{ΠΡΟΣΕΥΧΗ ΜΑΝΑΣΣΗ ΥΙΟΥ ΕΖΕΚΙΟΥ}
\biblebook{ΠΡΟΣΕΥΧΗ ΜΑΝΑΣΣΗ ΥΙΟΥ ΕΖΕΚΙΟΥ}

 
\lettrine[lines=2, loversize=0.2, nindent=0em, findent=.25em]{\textcolor{bookheadingcolor}{Κ}}{ΥΡΙΕ} παντοκράτωπ, ὁ Θεὸς τῶν πατέρων ἡμῶν τοῦʼ Αβραὰμ καὶ Ἰσαὰκ καὶ Ἰακὼβ καὶ τοῦ σπέρματος αὐτῶν τοῦ δικαίου.
\vs{2}Ὁ ποιήσας τὸν οὐρανὸν καὶ τὴν γῆν σὺν πάντι τῷ κόσμῳ αὐτῶν.
\vs{3}Ὁ πεδήσας τὴν θάλασσαν τῷ λόγῳ τοῦ προστάγματός σου, ὁ κλείσας τὴν ἄβυσσον καὶ σφραγισάμενος αὐτὴν τῷ φοβερῷ καὶ ἐνδόξῳ ὀνόματί σου·
\vs{4}ὃν πάντα φρίσσει καὶ τρέμει ἀπὸ προσώπου δυνάμεώς σου,
\vs{5}ὅτι ἄστεκτος ἡ μεγαλοπρέπεια τῆς δόξης σου, καὶ ἀνυπόστατος ἡ ὀργὴ τῆς ἐπὶ ἁμαρτωλοὺς ἀπειλῆς σου·
\vs{6}ἀμέτρητόν τε καὶ ἀνεξιχνίαστον τὸ ἔλεος τῆς ἐπαγγελίας σου·
\vs{7}σὺ γὰρ εἶ Κύριος ὕψιστος, εὔσπλαγχνος, μακρόθυμος καὶ πολυέλεος, μετανοῶν ἐπὶ κακίαις ἀνθρώπων. Σὺ, Κύριε, κατὰ τὸ πλῆθος τῆς χρηστότητός σου ἐπηγγείλω μετάνοιαν καὶ ἄφεσιν τοῖς ἡμαρτηκόσιν σοι, καὶ τῷ πλήθει τῶν οἰκτιρμῶν σου ὥρισας μετάνοιαν ἁμαρτωλοῖς εἰς σωτηρίαν.
\vs{8}Σὺ οὖν, Κύριε, ὁ Θεὸς τῶν δικαίων, οὐκ ἔθου μετάνοιαν δικαίοις, τῷ Ἀβραὰμ καὶ Ἰσαὰκ καὶ Ἰακὼβ, τοῖς οὐχ ἡμαρτηκόσιν σοι,
\vs{9}Ἐπλήθυναν αἱ ἀνομίαι μου, Κύριε, ἐπλήθυναν, καὶ οὐκ εἰμὶ ἄξιος ἀτενίσαι καὶ ἰδεῖν τὸ ὕψος τοῦ οὐρανοῦ ἀπὸ πλήθους τῶν ἀδικιῶν μου,
\vs{10}κατακαμπτόμενος πολλῷ δεσμῷ σιδηρῷ εἰς τὸ μὴ ἀνανεῦσαι τὴν κεφαλήν μου, καὶ οὐκ ἔστιν μοι ἄνεσις, διότι παρώργισα τὸν θυμόν σου, καὶ τὸ πονηρὸν ἐνώπιόν σου ἐποίησα, μὴ ποιήσας τὸ θέλημά σου καὶ μὴ φυλάξας τὰ προστάγματά σου, στήσας βδελύγματα καὶ πληθύνας προσοχθίσματα.
\vs{11}Καὶ νῦν κλίνω γόνυ καρδίας μου δεόμενος τῆς παθὰ σοῦ χρηστότητος· ἡμάρτηκα,
\vs{12}Κύριε, ἡμάρτηκα, καὶ τὰς ἀνομίας μου ἐγὼ γινώσκω. Ἀλλʼ αἰτοῦμαι δεόμενός σου· ἄνες μοι,
\vs{13}Κύριε, ἄνες μοι, καὶ μὴ συναπολέσῃς με ταῖς ἀνομίαις μου, μηδὲ εἰς τὸν αἰῶνα μηνίσας τηρήσῃς τὰ κακά μαι, μηδὲ καταδικάσῃς με ἐν τοῖς κατωτάτοις τῆς γῆς, διότι σὺ εἶ Θεὸς, Θεὶς τῶν μετανοούντων.
\vs{14}Καὶ ἐν ἐμοὶ δείξεις πᾶσαν τῆν ἀγαθωσύνην σου, ὁτι ἀνάξιον ὄντα σώσεις με κατὰ τὸ πολὺ ἔλεός σου.
\vs{15}Καὶ αἰνέσω σε διὰ παντὸς ἐν ταῖς ἡμέραις τῆς ζωῆς μου, ὅτι σὲ ὑμνεῖ πᾶσα ἡ δύναμις τῶν οὐρανῶν, καὶ σοῦ ἐστὶν ἡ δόξα εἰς τοὺς αἰῶναι. Ἀμήν

\cleardoublepage
\thispagestyle{empty}
\vspace*{3cm}
\phantomsection
\addcontentsline{toc}{part}{Η ΚΑΙΝΗ ΔΙΑΘΗΚΗ}
\begin{center}
  {\Huge Η ΚΑΙΝΗ ΔΙΑΘΗΚΗ}\\[2em]
\end{center}
\newpage
\thispagestyle{empty}
\null

\def\book{ΚΑΤΑ ΜΑΘΘΑΙΟΝ}
\biblebook{ΚΑΤΑ ΜΑΘΘΑΙΟΝ}


\lettrine[lines=2, loversize=0.2, nindent=0em, findent=.25em]{\textcolor{bookheadingcolor}{Β}}{ίβλος} γενέσεως Ἰησοῦ Χριστοῦ υἱοῦ Δαυὶδ υἱοῦ Ἀβραάμ.

\postdropcapindent\vs{2}Ἀβραὰμ ἐγέννησεν τὸν Ἰσαάκ, Ἰσαὰκ δὲ ἐγέννησεν τὸν Ἰακώβ, Ἰακὼβ δὲ ἐγέννησεν τὸν Ἰούδαν καὶ τοὺς ἀδελφοὺς αὐτοῦ,
\vs{3}Ἰούδας δὲ ἐγέννησεν τὸν Φαρὲς καὶ τὸν Ζαρὰ ἐκ τῆς Θάμαρ, Φαρὲς δὲ ἐγέννησεν τὸν Ἑσρώμ, Ἑσρὼμ δὲ ἐγέννησεν τὸν Ἀράμ,
\vs{4}Ἀρὰμ δὲ ἐγέννησεν τὸν Ἀμιναδάβ, Ἀμιναδὰβ δὲ ἐγέννησεν τὸν Ναασσών, Ναασσὼν δὲ ἐγέννησεν τὸν Σαλμών,
\vs{5}Σαλμὼν δὲ ἐγέννησεν τὸν Βόες ἐκ τῆς Ῥαχάβ, Βόες δὲ ἐγέννησεν τὸν Ἰωβὴδ ἐκ τῆς Ῥούθ, Ἰωβὴδ δὲ ἐγέννησεν τὸν Ἰεσσαί,
\vs{6}Ἰεσσαὶ δὲ ἐγέννησεν τὸν Δαυὶδ τὸν βασιλέα.

Δαυὶδ δὲ ἐγέννησεν τὸν Σολομῶνα ἐκ τῆς τοῦ Οὐρίου,
\vs{7}Σολομὼν δὲ ἐγέννησεν τὸν Ῥοβοάμ, Ῥοβοὰμ δὲ ἐγέννησεν τὸν Ἀβιά, Ἀβιὰ δὲ ἐγέννησεν τὸν Ἀσάφ,
\vs{8}Ἀσὰφ δὲ ἐγέννησεν τὸν Ἰωσαφάτ, Ἰωσαφὰτ δὲ ἐγέννησεν τὸν Ἰωράμ, Ἰωρὰμ δὲ ἐγέννησεν τὸν Ὀζίαν,
\vs{9}Ὀζίας δὲ ἐγέννησεν τὸν Ἰωαθάμ, Ἰωαθὰμ δὲ ἐγέννησεν τὸν Ἄχαζ, Ἄχαζ δὲ ἐγέννησεν τὸν Ἑζεκίαν,
\vs{10}Ἑζεκίας δὲ ἐγέννησεν τὸν Μανασσῆ, Μανασσῆς δὲ ἐγέννησεν τὸν Ἀμώς, Ἀμὼς δὲ ἐγέννησεν τὸν Ἰωσίαν,
\vs{11}Ἰωσίας δὲ ἐγέννησεν τὸν Ἰεχονίαν καὶ τοὺς ἀδελφοὺς αὐτοῦ ἐπὶ τῆς μετοικεσίας Βαβυλῶνος.

\vs{12}Μετὰ δὲ τὴν μετοικεσίαν Βαβυλῶνος Ἰεχονίας ἐγέννησεν τὸν Σαλαθιήλ, Σαλαθιὴλ δὲ ἐγέννησεν τὸν Ζοροβαβέλ,
\vs{13}Ζοροβαβὲλ δὲ ἐγέννησεν τὸν Ἀβιούδ, Ἀβιοὺδ δὲ ἐγέννησεν τὸν Ἐλιακίμ, Ἐλιακὶμ δὲ ἐγέννησεν τὸν Ἀζώρ,
\vs{14}Ἀζὼρ δὲ ἐγέννησεν τὸν Σαδώκ, Σαδὼκ δὲ ἐγέννησεν τὸν Ἀχίμ, Ἀχὶμ δὲ ἐγέννησεν τὸν Ἐλιούδ,
\vs{15}Ἐλιοὺδ δὲ ἐγέννησεν τὸν Ἐλεάζαρ, Ἐλεάζαρ δὲ ἐγέννησεν τὸν Ματθάν, Ματθὰν δὲ ἐγέννησεν τὸν Ἰακώβ,
\vs{16}Ἰακὼβ δὲ ἐγέννησεν τὸν Ἰωσὴφ τὸν ἄνδρα Μαρίας, ἐξ ἧς ἐγεννήθη Ἰησοῦς ὁ λεγόμενος Χριστός.

\vs{17}Πᾶσαι οὖν αἱ γενεαὶ ἀπὸ Ἀβραὰμ ἕως Δαυὶδ γενεαὶ δεκατέσσαρες, καὶ ἀπὸ Δαυὶδ ἕως τῆς μετοικεσίας Βαβυλῶνος γενεαὶ δεκατέσσαρες, καὶ ἀπὸ τῆς μετοικεσίας Βαβυλῶνος ἕως τοῦ Χριστοῦ γενεαὶ δεκατέσσαρες.

\vs{18}Τοῦ δὲ Ἰησοῦ Χριστοῦ ἡ γένεσις οὕτως ἦν. μνηστευθείσης τῆς μητρὸς αὐτοῦ Μαρίας τῷ Ἰωσήφ, πρὶν ἢ συνελθεῖν αὐτοὺς εὑρέθη ἐν γαστρὶ ἔχουσα ἐκ πνεύματος ἁγίου.
\vs{19}Ἰωσὴφ δὲ ὁ ἀνὴρ αὐτῆς, δίκαιος ὢν καὶ μὴ θέλων αὐτὴν δειγματίσαι, ἐβουλήθη λάθρᾳ ἀπολῦσαι αὐτήν.
\vs{20}Ταῦτα δὲ αὐτοῦ ἐνθυμηθέντος ἰδοὺ ἄγγελος Κυρίου κατ᾽ ὄναρ ἐφάνη αὐτῷ λέγων· Ἰωσὴφ υἱὸς Δαυίδ, μὴ φοβηθῇς παραλαβεῖν Μαρίαν τὴν γυναῖκά σου· τὸ γὰρ ἐν αὐτῇ γεννηθὲν ἐκ Πνεύματός ἐστιν Ἁγίου.
\vs{21}τέξεται δὲ υἱὸν, καὶ καλέσεις τὸ ὄνομα αὐτοῦ Ἰησοῦν· αὐτὸς γὰρ σώσει τὸν λαὸν αὐτοῦ ἀπὸ τῶν ἁμαρτιῶν αὐτῶν.
\vs{22}Τοῦτο δὲ ὅλον γέγονεν ἵνα πληρωθῇ τὸ ῥηθὲν ὑπὸ Κυρίου διὰ τοῦ προφήτου λέγοντος·
\begin{poetryblock}

\begin{quote} \vs{23}Ἰδοὺ ἡ παρθένος ἐν γαστρὶ ἕξει καὶ τέξεται υἱόν,\end{quote} 

\begin{quote}καὶ καλέσουσιν τὸ ὄνομα αὐτοῦ Ἐμμανουήλ,\end{quote}
\end{poetryblock}

ὅ ἐστιν μεθερμηνευόμενον Μεθ᾽ ἡμῶν ὁ Θεός.
\vs{24}Ἐγερθεὶς δὲ ὁ Ἰωσὴφ ἀπὸ τοῦ ὕπνου ἐποίησεν ὡς προσέταξεν αὐτῷ ὁ ἄγγελος Κυρίου καὶ παρέλαβεν τὴν γυναῖκα αὐτοῦ,
\vs{25}καὶ οὐκ ἐγίνωσκεν αὐτὴν ἕως οὗ ἔτεκεν υἱόν· καὶ ἐκάλεσεν τὸ ὄνομα αὐτοῦ Ἰησοῦν.

\ch{2}
Τοῦ δὲ Ἰησοῦ γεννηθέντος ἐν Βηθλέεμ τῆς Ἰουδαίας ἐν ἡμέραις Ἡρῴδου τοῦ βασιλέως, ἰδοὺ μάγοι ἀπὸ ἀνατολῶν παρεγένοντο εἰς Ἱεροσόλυμα
\vs{2}λέγοντες· Ποῦ ἐστιν ὁ τεχθεὶς βασιλεὺς τῶν Ἰουδαίων; εἴδομεν γὰρ αὐτοῦ τὸν ἀστέρα ἐν τῇ ἀνατολῇ καὶ ἤλθομεν προσκυνῆσαι αὐτῷ.
\vs{3}Ἀκούσας δὲ ὁ βασιλεὺς Ἡρῴδης ἐταράχθη καὶ πᾶσα Ἱεροσόλυμα μετ᾽ αὐτοῦ,
\vs{4}καὶ συναγαγὼν πάντας τοὺς ἀρχιερεῖς καὶ γραμματεῖς τοῦ λαοῦ ἐπυνθάνετο παρ᾽ αὐτῶν ποῦ ὁ Χριστὸς γεννᾶται.
\vs{5}Οἱ δὲ εἶπαν αὐτῷ· Ἐν Βηθλέεμ τῆς Ἰουδαίας· οὕτως γὰρ γέγραπται διὰ τοῦ προφήτου·
\begin{poetryblock}

\begin{quote} \vs{6}Καὶ σύ Βηθλέεμ, γῆ Ἰούδα,\end{quote} 

\begin{quote}οὐδαμῶς ἐλαχίστη εἶ ἐν τοῖς ἡγεμόσιν Ἰούδα·\end{quote} 

\begin{quote}ἐκ σοῦ γὰρ ἐξελεύσεται ἡγούμενος,\end{quote} 

\begin{quote}ὅστις ποιμανεῖ τὸν λαόν μου τὸν Ἰσραήλ.\end{quote}
\end{poetryblock}

\vs{7}Τότε Ἡρῴδης λάθρᾳ καλέσας τοὺς μάγους ἠκρίβωσεν παρ᾽ αὐτῶν τὸν χρόνον τοῦ φαινομένου ἀστέρος,
\vs{8}καὶ πέμψας αὐτοὺς εἰς Βηθλέεμ εἶπεν· Πορευθέντες ἐξετάσατε ἀκριβῶς περὶ τοῦ παιδίου· ἐπὰν δὲ εὕρητε, ἀπαγγείλατέ μοι, ὅπως κἀγὼ ἐλθὼν προσκυνήσω αὐτῷ.
\vs{9}Οἱ δὲ ἀκούσαντες τοῦ βασιλέως ἐπορεύθησαν καὶ ἰδοὺ ὁ ἀστὴρ, ὃν εἶδον ἐν τῇ ἀνατολῇ, προῆγεν αὐτούς, ἕως ἐλθὼν ἐστάθη ἐπάνω οὗ ἦν τὸ παιδίον.
\vs{10}ἰδόντες δὲ τὸν ἀστέρα ἐχάρησαν χαρὰν μεγάλην σφόδρα.
\vs{11}καὶ ἐλθόντες εἰς τὴν οἰκίαν εἶδον τὸ παιδίον μετὰ Μαρίας τῆς μητρὸς αὐτοῦ, καὶ πεσόντες προσεκύνησαν αὐτῷ καὶ ἀνοίξαντες τοὺς θησαυροὺς αὐτῶν προσήνεγκαν αὐτῷ δῶρα, χρυσὸν καὶ λίβανον καὶ σμύρναν.
\vs{12}Καὶ χρηματισθέντες κατ᾽ ὄναρ μὴ ἀνακάμψαι πρὸς Ἡρῴδην, δι᾽ ἄλλης ὁδοῦ ἀνεχώρησαν εἰς τὴν χώραν αὐτῶν.

\vs{13}Ἀναχωρησάντων δὲ αὐτῶν ἰδοὺ ἄγγελος κυρίου φαίνεται κατ᾽ ὄναρ τῷ Ἰωσὴφ λέγων· Ἐγερθεὶς παράλαβε τὸ παιδίον καὶ τὴν μητέρα αὐτοῦ καὶ φεῦγε εἰς Αἴγυπτον καὶ ἴσθι ἐκεῖ ἕως ἂν εἴπω σοι· μέλλει γὰρ Ἡρῴδης ζητεῖν τὸ παιδίον τοῦ ἀπολέσαι αὐτό.
\vs{14}Ὁ δὲ ἐγερθεὶς παρέλαβεν τὸ παιδίον καὶ τὴν μητέρα αὐτοῦ νυκτὸς καὶ ἀνεχώρησεν εἰς Αἴγυπτον,
\vs{15}καὶ ἦν ἐκεῖ ἕως τῆς τελευτῆς Ἡρῴδου· ἵνα πληρωθῇ τὸ ῥηθὲν ὑπὸ κυρίου διὰ τοῦ προφήτου λέγοντος· 
\begin{poetryblock}

\begin{quote}Ἐξ Αἰγύπτου ἐκάλεσα τὸν υἱόν μου.\end{quote}
\end{poetryblock}

\vs{16}Τότε Ἡρῴδης ἰδὼν ὅτι ἐνεπαίχθη ὑπὸ τῶν μάγων ἐθυμώθη λίαν, καὶ ἀποστείλας ἀνεῖλεν πάντας τοὺς παῖδας τοὺς ἐν Βηθλέεμ καὶ ἐν πᾶσι τοῖς ὁρίοις αὐτῆς ἀπὸ διετοῦς καὶ κατωτέρω, κατὰ τὸν χρόνον ὃν ἠκρίβωσεν παρὰ τῶν μάγων.
\vs{17}τότε ἐπληρώθη τὸ ῥηθὲν διὰ Ἰερεμίου τοῦ προφήτου λέγοντος·
\begin{poetryblock}

\begin{quote} \vs{18}Φωνὴ ἐν Ῥαμὰ ἠκούσθη,\end{quote} 

\begin{quote}κλαυθμὸς καὶ ὀδυρμὸς πολύς·\end{quote} 

\begin{quote}Ῥαχὴλ κλαίουσα τὰ τέκνα αὐτῆς,\end{quote} 

\begin{quote}καὶ οὐκ ἤθελεν παρακληθῆναι,\end{quote} 

\begin{quote}ὅτι οὐκ εἰσίν.\end{quote}
\end{poetryblock}

\vs{19}Τελευτήσαντος δὲ τοῦ Ἡρῴδου ἰδοὺ ἄγγελος Κυρίου φαίνεται κατ᾽ ὄναρ τῷ Ἰωσὴφ ἐν Αἰγύπτῳ
\vs{20}λέγων· Ἐγερθεὶς παράλαβε τὸ παιδίον καὶ τὴν μητέρα αὐτοῦ καὶ πορεύου εἰς γῆν Ἰσραήλ· τεθνήκασιν γὰρ οἱ ζητοῦντες τὴν ψυχὴν τοῦ παιδίου.
\vs{21}Ὁ δὲ ἐγερθεὶς παρέλαβεν τὸ παιδίον καὶ τὴν μητέρα αὐτοῦ καὶ εἰσῆλθεν εἰς γῆν Ἰσραήλ.

\vs{22}ἀκούσας δὲ ὅτι Ἀρχέλαος βασιλεύει τῆς Ἰουδαίας ἀντὶ τοῦ πατρὸς αὐτοῦ Ἡρῴδου ἐφοβήθη ἐκεῖ ἀπελθεῖν· χρηματισθεὶς δὲ κατ᾽ ὄναρ ἀνεχώρησεν εἰς τὰ μέρη τῆς Γαλιλαίας,
\vs{23}καὶ ἐλθὼν κατῴκησεν εἰς πόλιν λεγομένην Ναζαρέτ· ὅπως πληρωθῇ τὸ ῥηθὲν διὰ τῶν προφητῶν ὅτι Ναζωραῖος κληθήσεται.

\ch{3}
Ἐν δὲ ταῖς ἡμέραις ἐκείναις παραγίνεται Ἰωάννης ὁ βαπτιστὴς κηρύσσων ἐν τῇ ἐρήμῳ τῆς Ἰουδαίας
\vs{2}καὶ λέγων· Μετανοεῖτε· ἤγγικεν γὰρ ἡ βασιλεία τῶν οὐρανῶν.

\vs{3}οὗτος γάρ ἐστιν ὁ ῥηθεὶς διὰ Ἠσαΐου τοῦ προφήτου λέγοντος· 
\begin{poetryblock}

\begin{quote}Φωνὴ βοῶντος ἐν τῇ ἐρήμῳ·\end{quote} 

\begin{quote}Ἑτοιμάσατε τὴν ὁδὸν Κυρίου,\end{quote} 

\begin{quote}εὐθείας ποιεῖτε τὰς τρίβους αὐτοῦ.\end{quote}
\end{poetryblock}

\vs{4}Αὐτὸς δὲ ὁ Ἰωάννης εἶχεν τὸ ἔνδυμα αὐτοῦ ἀπὸ τριχῶν καμήλου καὶ ζώνην δερματίνην περὶ τὴν ὀσφὺν αὐτοῦ, ἡ δὲ τροφὴ ἦν αὐτοῦ ἀκρίδες καὶ μέλι ἄγριον.
\vs{5}Τότε ἐξεπορεύετο πρὸς αὐτὸν Ἱεροσόλυμα καὶ πᾶσα ἡ Ἰουδαία καὶ πᾶσα ἡ περίχωρος τοῦ Ἰορδάνου,
\vs{6}καὶ ἐβαπτίζοντο ἐν τῷ Ἰορδάνῃ ποταμῷ ὑπ᾽ αὐτοῦ ἐξομολογούμενοι τὰς ἁμαρτίας αὐτῶν.

\vs{7}Ἰδὼν δὲ πολλοὺς τῶν Φαρισαίων καὶ Σαδδουκαίων ἐρχομένους ἐπὶ τὸ βάπτισμα αὐτοῦ εἶπεν αὐτοῖς· Γεννήματα ἐχιδνῶν, τίς ὑπέδειξεν ὑμῖν φυγεῖν ἀπὸ τῆς μελλούσης ὀργῆς;
\vs{8}ποιήσατε οὖν καρπὸν ἄξιον τῆς μετανοίας
\vs{9}καὶ μὴ δόξητε λέγειν ἐν ἑαυτοῖς· Πατέρα ἔχομεν τὸν Ἀβραάμ. λέγω γὰρ ὑμῖν ὅτι δύναται ὁ Θεὸς ἐκ τῶν λίθων τούτων ἐγεῖραι τέκνα τῷ Ἀβραάμ.
\vs{10}ἤδη δὲ ἡ ἀξίνη πρὸς τὴν ῥίζαν τῶν δένδρων κεῖται· πᾶν οὖν δένδρον μὴ ποιοῦν καρπὸν καλὸν ἐκκόπτεται καὶ εἰς πῦρ βάλλεται.

\vs{11}Ἐγὼ μὲν ὑμᾶς βαπτίζω ἐν ὕδατι εἰς μετάνοιαν, ὁ δὲ ὀπίσω μου ἐρχόμενος ἰσχυρότερός μού ἐστιν, οὗ οὐκ εἰμὶ ἱκανὸς τὰ ὑποδήματα βαστάσαι· αὐτὸς ὑμᾶς βαπτίσει ἐν Πνεύματι Ἁγίῳ καὶ πυρί·
\vs{12}οὗ τὸ πτύον ἐν τῇ χειρὶ αὐτοῦ καὶ διακαθαριεῖ τὴν ἅλωνα αὐτοῦ καὶ συνάξει τὸν σῖτον αὐτοῦ εἰς τὴν ἀποθήκην, τὸ δὲ ἄχυρον κατακαύσει πυρὶ ἀσβέστῳ.

\vs{13}Τότε παραγίνεται ὁ Ἰησοῦς ἀπὸ τῆς Γαλιλαίας ἐπὶ τὸν Ἰορδάνην πρὸς τὸν Ἰωάννην τοῦ βαπτισθῆναι ὑπ᾽ αὐτοῦ.
\vs{14}ὁ δὲ Ἰωάννης διεκώλυεν αὐτὸν λέγων· Ἐγὼ χρείαν ἔχω ὑπὸ σοῦ βαπτισθῆναι, καὶ σὺ ἔρχῃ πρός με;
\vs{15}Ἀποκριθεὶς δὲ ὁ Ἰησοῦς εἶπεν πρὸς αὐτόν· Ἄφες ἄρτι, οὕτως γὰρ πρέπον ἐστὶν ἡμῖν πληρῶσαι πᾶσαν δικαιοσύνην. τότε ἀφίησιν αὐτόν.
\vs{16}Βαπτισθεὶς δὲ ὁ Ἰησοῦς εὐθὺς ἀνέβη ἀπὸ τοῦ ὕδατος· καὶ ἰδοὺ ἠνεῴχθησαν αὐτῷ οἱ οὐρανοί, καὶ εἶδεν τὸ Πνεῦμα τοῦ Θεοῦ καταβαῖνον ὡσεὶ περιστερὰν καὶ ἐρχόμενον ἐπ᾽ αὐτόν·
\vs{17}καὶ ἰδοὺ φωνὴ ἐκ τῶν οὐρανῶν λέγουσα· Οὗτός ἐστιν ὁ Υἱός μου ὁ ἀγαπητός, ἐν ᾧ εὐδόκησα.

\ch{4}
Τότε ὁ Ἰησοῦς ἀνήχθη εἰς τὴν ἔρημον ὑπὸ τοῦ Πνεύματος πειρασθῆναι ὑπὸ τοῦ διαβόλου.
\vs{2}καὶ νηστεύσας ἡμέρας τεσσεράκοντα καὶ νύκτας τεσσεράκοντα, ὕστερον ἐπείνασεν.
\vs{3}Καὶ προσελθὼν ὁ πειράζων εἶπεν αὐτῷ· Εἰ Υἱὸς εἶ τοῦ Θεοῦ, εἰπὲ ἵνα οἱ λίθοι οὗτοι ἄρτοι γένωνται.
\vs{4}Ὁ δὲ ἀποκριθεὶς εἶπεν· Γέγραπται· Οὐκ ἐπ᾽ ἄρτῳ μόνῳ ζήσεται ὁ ἄνθρωπος, Ἀλλ᾽ ἐπὶ παντὶ ῥήματι ἐκπορευομένῳ διὰ στόματος Θεοῦ.

\vs{5}Τότε παραλαμβάνει αὐτὸν ὁ διάβολος εἰς τὴν ἁγίαν πόλιν καὶ ἔστησεν αὐτὸν ἐπὶ τὸ πτερύγιον τοῦ ἱεροῦ
\vs{6}καὶ λέγει αὐτῷ· Εἰ Υἱὸς εἶ τοῦ Θεοῦ, βάλε σεαυτὸν κάτω· γέγραπται γὰρ ὅτι 
\begin{poetryblock}

\begin{quote}Τοῖς ἀγγέλοις αὐτοῦ ἐντελεῖται περὶ σοῦ\end{quote} 

\begin{quote}καὶ ἐπὶ χειρῶν ἀροῦσίν σε,\end{quote} 

\begin{quote}μήποτε προσκόψῃς πρὸς λίθον τὸν πόδα σου.\end{quote}
\end{poetryblock}

\vs{7}Ἔφη αὐτῷ ὁ Ἰησοῦς· Πάλιν γέγραπται· Οὐκ ἐκπειράσεις Κύριον τὸν Θεόν σου.

\vs{8}Πάλιν παραλαμβάνει αὐτὸν ὁ διάβολος εἰς ὄρος ὑψηλὸν λίαν καὶ δείκνυσιν αὐτῷ πάσας τὰς βασιλείας τοῦ κόσμου καὶ τὴν δόξαν αὐτῶν
\vs{9}καὶ εἶπεν αὐτῷ· Ταῦτά σοι πάντα δώσω, ἐὰν πεσὼν προσκυνήσῃς μοι.
\vs{10}Τότε λέγει αὐτῷ ὁ Ἰησοῦς· Ὕπαγε, Σατανᾶ· γέγραπται γάρ· Κύριον τὸν θεόν σου προσκυνήσεις καὶ αὐτῷ μόνῳ λατρεύσεις.

\vs{11}Τότε ἀφίησιν αὐτὸν ὁ διάβολος, καὶ ἰδοὺ ἄγγελοι προσῆλθον καὶ διηκόνουν αὐτῷ.

\vs{12}Ἀκούσας δὲ ὅτι Ἰωάννης παρεδόθη ἀνεχώρησεν εἰς τὴν Γαλιλαίαν.
\vs{13}καὶ καταλιπὼν τὴν Ναζαρὰ ἐλθὼν κατῴκησεν εἰς Καφαρναοὺμ τὴν παραθαλασσίαν ἐν ὁρίοις Ζαβουλὼν καὶ Νεφθαλίμ·
\vs{14}ἵνα πληρωθῇ τὸ ῥηθὲν διὰ Ἠσαΐου τοῦ προφήτου λέγοντος·
\begin{poetryblock}

\begin{quote} \vs{15}Γῆ Ζαβουλὼν καὶ γῆ Νεφθαλίμ,\end{quote} 

\begin{quote}ὁδὸν θαλάσσης, πέραν τοῦ Ἰορδάνου,\end{quote} 

\begin{quote}Γαλιλαία τῶν ἐθνῶν,\end{quote}

\begin{quote} \vs{16}ὁ λαὸς ὁ καθήμενος ἐν σκοτίᾳ\end{quote} 

\begin{quote}φῶς εἶδεν μέγα,\end{quote} 

\begin{quote}καὶ τοῖς καθημένοις ἐν χώρᾳ καὶ σκιᾷ θανάτου\end{quote} 

\begin{quote}φῶς ἀνέτειλεν αὐτοῖς.\end{quote}
\end{poetryblock}

\vs{17}Ἀπὸ τότε ἤρξατο ὁ Ἰησοῦς κηρύσσειν καὶ λέγειν· Μετανοεῖτε· ἤγγικεν γὰρ ἡ βασιλεία τῶν οὐρανῶν.

\vs{18}Περιπατῶν δὲ παρὰ τὴν θάλασσαν τῆς Γαλιλαίας εἶδεν δύο ἀδελφούς, Σίμωνα τὸν λεγόμενον Πέτρον καὶ Ἀνδρέαν τὸν ἀδελφὸν αὐτοῦ, βάλλοντας ἀμφίβληστρον εἰς τὴν θάλασσαν· ἦσαν γὰρ ἁλιεῖς.
\vs{19}καὶ λέγει αὐτοῖς· Δεῦτε ὀπίσω μου, καὶ ποιήσω ὑμᾶς ἁλιεῖς ἀνθρώπων.
\vs{20}οἱ δὲ εὐθέως ἀφέντες τὰ δίκτυα ἠκολούθησαν αὐτῷ.
\vs{21}Καὶ προβὰς ἐκεῖθεν εἶδεν ἄλλους δύο ἀδελφούς, Ἰάκωβον τὸν τοῦ Ζεβεδαίου καὶ Ἰωάννην τὸν ἀδελφὸν αὐτοῦ, ἐν τῷ πλοίῳ μετὰ Ζεβεδαίου τοῦ πατρὸς αὐτῶν καταρτίζοντας τὰ δίκτυα αὐτῶν, καὶ ἐκάλεσεν αὐτούς.
\vs{22}οἱ δὲ εὐθέως ἀφέντες τὸ πλοῖον καὶ τὸν πατέρα αὐτῶν ἠκολούθησαν αὐτῷ.

\vs{23}Καὶ περιῆγεν ἐν ὅλῃ τῇ Γαλιλαίᾳ διδάσκων ἐν ταῖς συναγωγαῖς αὐτῶν καὶ κηρύσσων τὸ εὐαγγέλιον τῆς βασιλείας καὶ θεραπεύων πᾶσαν νόσον καὶ πᾶσαν μαλακίαν ἐν τῷ λαῷ.

\vs{24}καὶ ἀπῆλθεν ἡ ἀκοὴ αὐτοῦ εἰς ὅλην τὴν Συρίαν· καὶ προσήνεγκαν αὐτῷ πάντας τοὺς κακῶς ἔχοντας ποικίλαις νόσοις καὶ βασάνοις συνεχομένους καὶ δαιμονιζομένους καὶ σεληνιαζομένους καὶ παραλυτικούς, καὶ ἐθεράπευσεν αὐτούς.
\vs{25}Καὶ ἠκολούθησαν αὐτῷ ὄχλοι πολλοὶ ἀπὸ τῆς Γαλιλαίας καὶ Δεκαπόλεως καὶ Ἱεροσολύμων καὶ Ἰουδαίας καὶ πέραν τοῦ Ἰορδάνου.

\ch{5}
Ἰδὼν δὲ τοὺς ὄχλους ἀνέβη εἰς τὸ ὄρος, καὶ καθίσαντος αὐτοῦ προσῆλθαν αὐτῷ οἱ μαθηταὶ αὐτοῦ·
\vs{2}καὶ ἀνοίξας τὸ στόμα αὐτοῦ ἐδίδασκεν αὐτοὺς λέγων·
\begin{poetryblock}

\begin{quote} \vs{3}Μακάριοι οἱ πτωχοὶ τῷ πνεύματι,\end{quote} 

\begin{quote}Ὅτι αὐτῶν ἐστιν ἡ βασιλεία τῶν οὐρανῶν.\end{quote}

\begin{quote} \vs{4}Μακάριοι οἱ πενθοῦντες, Ὅτι αὐτοὶ παρακληθήσονται.\end{quote}

\begin{quote} \vs{5}Μακάριοι οἱ πραεῖς,\end{quote} 

\begin{quote}Ὅτι αὐτοὶ κληρονομήσουσιν τὴν γῆν.\end{quote}

\begin{quote} \vs{6}Μακάριοι οἱ πεινῶντες καὶ διψῶντες τὴν δικαιοσύνην,\end{quote} 

\begin{quote}Ὅτι αὐτοὶ χορτασθήσονται.\end{quote}

\begin{quote} \vs{7}Μακάριοι οἱ ἐλεήμονες,\end{quote} 

\begin{quote}Ὅτι αὐτοὶ ἐλεηθήσονται.\end{quote}

\begin{quote} \vs{8}Μακάριοι οἱ καθαροὶ τῇ καρδίᾳ,\end{quote} 

\begin{quote}Ὅτι αὐτοὶ τὸν Θεὸν ὄψονται.\end{quote}

\begin{quote} \vs{9}Μακάριοι οἱ εἰρηνοποιοί,\end{quote} 

\begin{quote}Ὅτι αὐτοὶ υἱοὶ Θεοῦ κληθήσονται.\end{quote}

\begin{quote} \vs{10}Μακάριοι οἱ δεδιωγμένοι ἕνεκεν δικαιοσύνης,\end{quote} 

\begin{quote}Ὅτι αὐτῶν ἐστιν ἡ βασιλεία τῶν οὐρανῶν.\end{quote}

\begin{quote} \vs{11}Μακάριοί ἐστε\end{quote} 

\begin{quote}ὅταν ὀνειδίσωσιν ὑμᾶς καὶ διώξωσιν καὶ εἴπωσιν πᾶν πονηρὸν καθ᾽ ὑμῶν ψευδόμενοι ἕνεκεν ἐμοῦ.\end{quote}

\begin{quote} \vs{12}χαίρετε καὶ ἀγαλλιᾶσθε, ὅτι ὁ μισθὸς ὑμῶν πολὺς ἐν τοῖς οὐρανοῖς· οὕτως γὰρ ἐδίωξαν τοὺς προφήτας τοὺς πρὸ ὑμῶν.\end{quote}
\end{poetryblock}

\vs{13}Ὑμεῖς ἐστε τὸ ἅλας τῆς γῆς· ἐὰν δὲ τὸ ἅλας μωρανθῇ, ἐν τίνι ἁλισθήσεται; εἰς οὐδὲν ἰσχύει ἔτι εἰ μὴ βληθὲν ἔξω καταπατεῖσθαι ὑπὸ τῶν ἀνθρώπων.

\vs{14}Ὑμεῖς ἐστε τὸ φῶς τοῦ κόσμου. οὐ δύναται πόλις κρυβῆναι ἐπάνω ὄρους κειμένη·
\vs{15}οὐδὲ καίουσιν λύχνον καὶ τιθέασιν αὐτὸν ὑπὸ τὸν μόδιον ἀλλ᾽ ἐπὶ τὴν λυχνίαν, καὶ λάμπει πᾶσιν τοῖς ἐν τῇ οἰκίᾳ.
\vs{16}οὕτως λαμψάτω τὸ φῶς ὑμῶν ἔμπροσθεν τῶν ἀνθρώπων, ὅπως ἴδωσιν ὑμῶν τὰ καλὰ ἔργα καὶ δοξάσωσιν τὸν πατέρα ὑμῶν τὸν ἐν τοῖς οὐρανοῖς.

\vs{17}Μὴ νομίσητε ὅτι ἦλθον καταλῦσαι τὸν νόμον ἢ τοὺς προφήτας· οὐκ ἦλθον καταλῦσαι ἀλλὰ πληρῶσαι.
\vs{18}ἀμὴν γὰρ λέγω ὑμῖν· ἕως ἂν παρέλθῃ ὁ οὐρανὸς καὶ ἡ γῆ, ἰῶτα ἓν ἢ μία κεραία οὐ μὴ παρέλθῃ ἀπὸ τοῦ νόμου, ἕως ἂν πάντα γένηται.
\vs{19}Ὃς ἐὰν οὖν λύσῃ μίαν τῶν ἐντολῶν τούτων τῶν ἐλαχίστων καὶ διδάξῃ οὕτως τοὺς ἀνθρώπους, ἐλάχιστος κληθήσεται ἐν τῇ βασιλείᾳ τῶν οὐρανῶν· ὃς δ᾽ ἂν ποιήσῃ καὶ διδάξῃ, οὗτος μέγας κληθήσεται ἐν τῇ βασιλείᾳ τῶν οὐρανῶν.

\vs{20}λέγω γὰρ ὑμῖν ὅτι ἐὰν μὴ περισσεύσῃ ὑμῶν ἡ δικαιοσύνη πλεῖον τῶν γραμματέων καὶ Φαρισαίων, οὐ μὴ εἰσέλθητε εἰς τὴν βασιλείαν τῶν οὐρανῶν.

\vs{21}Ἠκούσατε ὅτι ἐρρέθη τοῖς ἀρχαίοις· Οὐ φονεύσεις· ὃς δ᾽ ἂν φονεύσῃ, ἔνοχος ἔσται τῇ κρίσει.
\vs{22}ἐγὼ δὲ λέγω ὑμῖν ὅτι πᾶς ὁ ὀργιζόμενος τῷ ἀδελφῷ αὐτοῦ ἔνοχος ἔσται τῇ κρίσει· ὃς δ᾽ ἂν εἴπῃ τῷ ἀδελφῷ αὐτοῦ· Ῥακά, ἔνοχος ἔσται τῷ συνεδρίῳ· ὃς δ᾽ ἂν εἴπῃ· Μωρέ, ἔνοχος ἔσται εἰς τὴν γέενναν τοῦ πυρός.
\vs{23}Ἐὰν οὖν προσφέρῃς τὸ δῶρόν σου ἐπὶ τὸ θυσιαστήριον κἀκεῖ μνησθῇς ὅτι ὁ ἀδελφός σου ἔχει τι κατὰ σοῦ,
\vs{24}ἄφες ἐκεῖ τὸ δῶρόν σου ἔμπροσθεν τοῦ θυσιαστηρίου καὶ ὕπαγε πρῶτον διαλλάγηθι τῷ ἀδελφῷ σου, καὶ τότε ἐλθὼν πρόσφερε τὸ δῶρόν σου.
\vs{25}ἴσθι εὐνοῶν τῷ ἀντιδίκῳ σου ταχὺ, ἕως ὅτου εἶ μετ᾽ αὐτοῦ ἐν τῇ ὁδῷ, μήποτέ σε παραδῷ ὁ ἀντίδικος τῷ κριτῇ καὶ ὁ κριτὴς τῷ ὑπηρέτῃ καὶ εἰς φυλακὴν βληθήσῃ·
\vs{26}ἀμὴν λέγω σοι, οὐ μὴ ἐξέλθῃς ἐκεῖθεν, ἕως ἂν ἀποδῷς τὸν ἔσχατον κοδράντην.

\vs{27}Ἠκούσατε ὅτι ἐρρέθη· Οὐ μοιχεύσεις.
\vs{28}ἐγὼ δὲ λέγω ὑμῖν ὅτι πᾶς ὁ βλέπων γυναῖκα πρὸς τὸ ἐπιθυμῆσαι αὐτὴν ἤδη ἐμοίχευσεν αὐτὴν ἐν τῇ καρδίᾳ αὐτοῦ.
\vs{29}εἰ δὲ ὁ ὀφθαλμός σου ὁ δεξιὸς σκανδαλίζει σε, ἔξελε αὐτὸν καὶ βάλε ἀπὸ σοῦ· συμφέρει γάρ σοι ἵνα ἀπόληται ἓν τῶν μελῶν σου καὶ μὴ ὅλον τὸ σῶμά σου βληθῇ εἰς γέενναν.
\vs{30}καὶ εἰ ἡ δεξιά σου χεὶρ σκανδαλίζει σε, ἔκκοψον αὐτὴν καὶ βάλε ἀπὸ σοῦ· συμφέρει γάρ σοι ἵνα ἀπόληται ἓν τῶν μελῶν σου καὶ μὴ ὅλον τὸ σῶμά σου εἰς γέενναν ἀπέλθῃ.

\vs{31}Ἐρρέθη δέ· Ὃς ἂν ἀπολύσῃ τὴν γυναῖκα αὐτοῦ, δότω αὐτῇ ἀποστάσιον.
\vs{32}ἐγὼ δὲ λέγω ὑμῖν ὅτι πᾶς ὁ ἀπολύων τὴν γυναῖκα αὐτοῦ παρεκτὸς λόγου πορνείας ποιεῖ αὐτὴν μοιχευθῆναι, καὶ ὃς ἐὰν ἀπολελυμένην γαμήσῃ, μοιχᾶται.

\vs{33}Πάλιν ἠκούσατε ὅτι ἐρρέθη τοῖς ἀρχαίοις· Οὐκ ἐπιορκήσεις, ἀποδώσεις δὲ τῷ Κυρίῳ τοὺς ὅρκους σου.
\vs{34}ἐγὼ δὲ λέγω ὑμῖν μὴ ὀμόσαι ὅλως· μήτε ἐν τῷ οὐρανῷ, ὅτι θρόνος ἐστὶν τοῦ Θεοῦ,
\vs{35}μήτε ἐν τῇ γῇ, ὅτι ὑποπόδιόν ἐστιν τῶν ποδῶν αὐτοῦ, μήτε εἰς Ἱεροσόλυμα, ὅτι πόλις ἐστὶν τοῦ μεγάλου Βασιλέως,
\vs{36}μήτε ἐν τῇ κεφαλῇ σου ὀμόσῃς, ὅτι οὐ δύνασαι μίαν τρίχα λευκὴν ποιῆσαι ἢ μέλαιναν.
\vs{37}ἔστω δὲ ὁ λόγος ὑμῶν Ναὶ ναί, οὒ Οὔ· τὸ δὲ περισσὸν τούτων ἐκ τοῦ πονηροῦ ἐστιν.

\vs{38}Ἠκούσατε ὅτι ἐρρέθη· Ὀφθαλμὸν ἀντὶ ὀφθαλμοῦ καὶ ὀδόντα ἀντὶ ὀδόντος.
\vs{39}ἐγὼ δὲ λέγω ὑμῖν μὴ ἀντιστῆναι τῷ πονηρῷ· ἀλλ᾽ ὅστις σε ῥαπίζει εἰς τὴν δεξιὰν σιαγόνα σου, στρέψον αὐτῷ καὶ τὴν ἄλλην·
\vs{40}καὶ τῷ θέλοντί σοι κριθῆναι καὶ τὸν χιτῶνά σου λαβεῖν, ἄφες αὐτῷ καὶ τὸ ἱμάτιον·
\vs{41}καὶ ὅστις σε ἀγγαρεύσει μίλιον ἕν, ὕπαγε μετ᾽ αὐτοῦ δύο.
\vs{42}τῷ αἰτοῦντί σε δός, καὶ τὸν θέλοντα ἀπὸ σοῦ δανίσασθαι μὴ ἀποστραφῇς.

\vs{43}Ἠκούσατε ὅτι ἐρρέθη· Ἀγαπήσεις τὸν πλησίον σου καὶ μισήσεις τὸν ἐχθρόν σου.
\vs{44}ἐγὼ δὲ λέγω ὑμῖν· ἀγαπᾶτε τοὺς ἐχθροὺς ὑμῶν καὶ προσεύχεσθε ὑπὲρ τῶν διωκόντων ὑμᾶς,
\vs{45}ὅπως γένησθε υἱοὶ τοῦ Πατρὸς ὑμῶν τοῦ ἐν οὐρανοῖς, ὅτι τὸν ἥλιον αὐτοῦ ἀνατέλλει ἐπὶ πονηροὺς καὶ ἀγαθοὺς καὶ βρέχει ἐπὶ δικαίους καὶ ἀδίκους.
\vs{46}ἐὰν γὰρ ἀγαπήσητε τοὺς ἀγαπῶντας ὑμᾶς, τίνα μισθὸν ἔχετε; οὐχὶ καὶ οἱ τελῶναι τὸ αὐτὸ ποιοῦσιν;
\vs{47}καὶ ἐὰν ἀσπάσησθε τοὺς ἀδελφοὺς ὑμῶν μόνον, τί περισσὸν ποιεῖτε; οὐχὶ καὶ οἱ ἐθνικοὶ τὸ αὐτὸ ποιοῦσιν;
\vs{48}Ἔσεσθε οὖν ὑμεῖς τέλειοι ὡς ὁ Πατὴρ ὑμῶν ὁ οὐράνιος τέλειός ἐστιν.

\ch{6}
Προσέχετε δὲ τὴν δικαιοσύνην ὑμῶν μὴ ποιεῖν ἔμπροσθεν τῶν ἀνθρώπων πρὸς τὸ θεαθῆναι αὐτοῖς· εἰ δὲ μή γε, μισθὸν οὐκ ἔχετε παρὰ τῷ Πατρὶ ὑμῶν τῷ ἐν τοῖς οὐρανοῖς.
\vs{2}Ὅταν οὖν ποιῇς ἐλεημοσύνην, μὴ σαλπίσῃς ἔμπροσθέν σου, ὥσπερ οἱ ὑποκριταὶ ποιοῦσιν ἐν ταῖς συναγωγαῖς καὶ ἐν ταῖς ῥύμαις, ὅπως δοξασθῶσιν ὑπὸ τῶν ἀνθρώπων· ἀμὴν λέγω ὑμῖν, ἀπέχουσιν τὸν μισθὸν αὐτῶν.
\vs{3}σοῦ δὲ ποιοῦντος ἐλεημοσύνην μὴ γνώτω ἡ ἀριστερά σου τί ποιεῖ ἡ δεξιά σου,
\vs{4}ὅπως ᾖ σου ἡ ἐλεημοσύνη ἐν τῷ κρυπτῷ· καὶ ὁ Πατήρ σου ὁ βλέπων ἐν τῷ κρυπτῷ ἀποδώσει σοι.
\vs{5}Καὶ ὅταν προσεύχησθε, οὐκ ἔσεσθε ὡς οἱ ὑποκριταί, ὅτι φιλοῦσιν ἐν ταῖς συναγωγαῖς καὶ ἐν ταῖς γωνίαις τῶν πλατειῶν ἑστῶτες προσεύχεσθαι, ὅπως φανῶσιν τοῖς ἀνθρώποις· ἀμὴν λέγω ὑμῖν, ἀπέχουσιν τὸν μισθὸν αὐτῶν.
\vs{6}σὺ δὲ ὅταν προσεύχῃ, εἴσελθε εἰς τὸ ταμεῖόν σου καὶ κλείσας τὴν θύραν σου πρόσευξαι τῷ Πατρί σου τῷ ἐν τῷ κρυπτῷ· καὶ ὁ Πατήρ σου ὁ βλέπων ἐν τῷ κρυπτῷ ἀποδώσει σοι.
\vs{7}Προσευχόμενοι δὲ μὴ βατταλογήσητε ὥσπερ οἱ ἐθνικοί, δοκοῦσιν γὰρ ὅτι ἐν τῇ πολυλογίᾳ αὐτῶν εἰσακουσθήσονται.
\vs{8}μὴ οὖν ὁμοιωθῆτε αὐτοῖς· οἶδεν γὰρ ὁ Πατὴρ ὑμῶν ὧν χρείαν ἔχετε πρὸ τοῦ ὑμᾶς αἰτῆσαι αὐτόν.

\vs{9}Οὕτως οὖν προσεύχεσθε ὑμεῖς· 
\begin{poetryblock}

\begin{quote}Πάτερ ἡμῶν ὁ ἐν τοῖς οὐρανοῖς·\end{quote} 

\begin{quote}Ἁγιασθήτω τὸ ὄνομά σου·\end{quote}

\begin{quote} \vs{10}Ἐλθέτω ἡ βασιλεία σου·\end{quote} 

\begin{quote}Γενηθήτω τὸ θέλημά σου,\end{quote} 

\begin{quote}Ὡς ἐν οὐρανῷ καὶ ἐπὶ γῆς·\end{quote}

\begin{quote} \vs{11}Τὸν ἄρτον ἡμῶν τὸν ἐπιούσιον δὸς ἡμῖν σήμερον·\end{quote}

\begin{quote} \vs{12}Καὶ ἄφες ἡμῖν τὰ ὀφειλήματα ἡμῶν,\end{quote} 

\begin{quote}Ὡς καὶ ἡμεῖς ἀφήκαμεν τοῖς ὀφειλέταις ἡμῶν·\end{quote}

\begin{quote} \vs{13}Καὶ μὴ εἰσενέγκῃς ἡμᾶς εἰς πειρασμόν,\end{quote} 

\begin{quote}Ἀλλὰ ῥῦσαι ἡμᾶς ἀπὸ τοῦ πονηροῦ.\end{quote}
\end{poetryblock}

\vs{14}Ἐὰν γὰρ ἀφῆτε τοῖς ἀνθρώποις τὰ παραπτώματα αὐτῶν, ἀφήσει καὶ ὑμῖν ὁ Πατὴρ ὑμῶν ὁ οὐράνιος·
\vs{15}ἐὰν δὲ μὴ ἀφῆτε τοῖς ἀνθρώποις, οὐδὲ ὁ Πατὴρ ὑμῶν ἀφήσει τὰ παραπτώματα ὑμῶν.

\vs{16}Ὅταν δὲ νηστεύητε, μὴ γίνεσθε ὡς οἱ ὑποκριταὶ σκυθρωποί, ἀφανίζουσιν γὰρ τὰ πρόσωπα αὐτῶν ὅπως φανῶσιν τοῖς ἀνθρώποις νηστεύοντες· ἀμὴν λέγω ὑμῖν, ἀπέχουσιν τὸν μισθὸν αὐτῶν.
\vs{17}σὺ δὲ νηστεύων ἄλειψαί σου τὴν κεφαλὴν καὶ τὸ πρόσωπόν σου νίψαι,
\vs{18}ὅπως μὴ φανῇς τοῖς ἀνθρώποις νηστεύων ἀλλὰ τῷ Πατρί σου τῷ ἐν τῷ κρυφαίῳ· καὶ ὁ Πατήρ σου ὁ βλέπων ἐν τῷ κρυφαίῳ ἀποδώσει σοι.
\vs{19}Μὴ θησαυρίζετε ὑμῖν θησαυροὺς ἐπὶ τῆς γῆς, ὅπου σὴς καὶ βρῶσις ἀφανίζει καὶ ὅπου κλέπται διορύσσουσιν καὶ κλέπτουσιν·
\vs{20}θησαυρίζετε δὲ ὑμῖν θησαυροὺς ἐν οὐρανῷ, ὅπου οὔτε σὴς οὔτε βρῶσις ἀφανίζει καὶ ὅπου κλέπται οὐ διορύσσουσιν οὐδὲ κλέπτουσιν·
\vs{21}ὅπου γάρ ἐστιν ὁ θησαυρός σου, ἐκεῖ ἔσται καὶ ἡ καρδία σου.

\vs{22}Ὁ λύχνος τοῦ σώματός ἐστιν ὁ ὀφθαλμός. ἐὰν οὖν ᾖ ὁ ὀφθαλμός σου ἁπλοῦς, ὅλον τὸ σῶμά σου φωτεινὸν ἔσται·
\vs{23}ἐὰν δὲ ὁ ὀφθαλμός σου πονηρὸς ᾖ, ὅλον τὸ σῶμά σου σκοτεινὸν ἔσται. εἰ οὖν τὸ φῶς τὸ ἐν σοὶ σκότος ἐστίν, τὸ σκότος πόσον.

\vs{24}Οὐδεὶς δύναται δυσὶ κυρίοις δουλεύειν· ἢ γὰρ τὸν ἕνα μισήσει καὶ τὸν ἕτερον ἀγαπήσει, ἢ ἑνὸς ἀνθέξεται καὶ τοῦ ἑτέρου καταφρονήσει. οὐ δύνασθε Θεῷ δουλεύειν καὶ μαμωνᾷ.

\vs{25}Διὰ τοῦτο λέγω ὑμῖν· μὴ μεριμνᾶτε τῇ ψυχῇ ὑμῶν τί φάγητε ἢ τί πίητε, μηδὲ τῷ σώματι ὑμῶν τί ἐνδύσησθε. οὐχὶ ἡ ψυχὴ πλεῖόν ἐστιν τῆς τροφῆς καὶ τὸ σῶμα τοῦ ἐνδύματος;
\vs{26}ἐμβλέψατε εἰς τὰ πετεινὰ τοῦ οὐρανοῦ ὅτι οὐ σπείρουσιν οὐδὲ θερίζουσιν οὐδὲ συνάγουσιν εἰς ἀποθήκας, καὶ ὁ Πατὴρ ὑμῶν ὁ οὐράνιος τρέφει αὐτά· οὐχ ὑμεῖς μᾶλλον διαφέρετε αὐτῶν;
\vs{27}τίς δὲ ἐξ ὑμῶν μεριμνῶν δύναται προσθεῖναι ἐπὶ τὴν ἡλικίαν αὐτοῦ πῆχυν ἕνα;
\vs{28}Καὶ περὶ ἐνδύματος τί μεριμνᾶτε; καταμάθετε τὰ κρίνα τοῦ ἀγροῦ πῶς αὐξάνουσιν· οὐ κοπιῶσιν οὐδὲ νήθουσιν·
\vs{29}λέγω δὲ ὑμῖν ὅτι οὐδὲ Σολομὼν ἐν πάσῃ τῇ δόξῃ αὐτοῦ περιεβάλετο ὡς ἓν τούτων.
\vs{30}εἰ δὲ τὸν χόρτον τοῦ ἀγροῦ σήμερον ὄντα καὶ αὔριον εἰς κλίβανον βαλλόμενον ὁ Θεὸς οὕτως ἀμφιέννυσιν, οὐ πολλῷ μᾶλλον ὑμᾶς, ὀλιγόπιστοι;
\vs{31}Μὴ οὖν μεριμνήσητε λέγοντες· Τί φάγωμεν; ἤ· Τί πίωμεν; ἤ· Τί περιβαλώμεθα;
\vs{32}πάντα γὰρ ταῦτα τὰ ἔθνη ἐπιζητοῦσιν· οἶδεν γὰρ ὁ Πατὴρ ὑμῶν ὁ οὐράνιος ὅτι χρῄζετε τούτων ἁπάντων.
\vs{33}ζητεῖτε δὲ πρῶτον τὴν βασιλείαν τοῦ θεοῦ καὶ τὴν δικαιοσύνην αὐτοῦ, καὶ ταῦτα πάντα προστεθήσεται ὑμῖν.
\vs{34}Μὴ οὖν μεριμνήσητε εἰς τὴν αὔριον, ἡ γὰρ αὔριον μεριμνήσει ἑαυτῆς· ἀρκετὸν τῇ ἡμέρᾳ ἡ κακία αὐτῆς.

\ch{7}
Μὴ κρίνετε, ἵνα μὴ κριθῆτε·
\vs{2}ἐν ᾧ γὰρ κρίματι κρίνετε κριθήσεσθε, καὶ ἐν ᾧ μέτρῳ μετρεῖτε μετρηθήσεται ὑμῖν.
\vs{3}Τί δὲ βλέπεις τὸ κάρφος τὸ ἐν τῷ ὀφθαλμῷ τοῦ ἀδελφοῦ σου, τὴν δὲ ἐν τῷ σῷ ὀφθαλμῷ δοκὸν οὐ κατανοεῖς;
\vs{4}ἢ πῶς ἐρεῖς τῷ ἀδελφῷ σου· Ἄφες ἐκβάλω τὸ κάρφος ἐκ τοῦ ὀφθαλμοῦ σου, καὶ ἰδοὺ ἡ δοκὸς ἐν τῷ ὀφθαλμῷ σοῦ;
\vs{5}ὑποκριτά, ἔκβαλε πρῶτον ἐκ τοῦ ὀφθαλμοῦ σοῦ τὴν δοκόν, καὶ τότε διαβλέψεις ἐκβαλεῖν τὸ κάρφος ἐκ τοῦ ὀφθαλμοῦ τοῦ ἀδελφοῦ σου.

\vs{6}Μὴ δῶτε τὸ ἅγιον τοῖς κυσίν μηδὲ βάλητε τοὺς μαργαρίτας ὑμῶν ἔμπροσθεν τῶν χοίρων, μήποτε καταπατήσουσιν αὐτοὺς ἐν τοῖς ποσὶν αὐτῶν καὶ στραφέντες ῥήξωσιν ὑμᾶς.

\vs{7}Αἰτεῖτε καὶ δοθήσεται ὑμῖν, ζητεῖτε καὶ εὑρήσετε, κρούετε καὶ ἀνοιγήσεται ὑμῖν·
\vs{8}πᾶς γὰρ ὁ αἰτῶν λαμβάνει καὶ ὁ ζητῶν εὑρίσκει καὶ τῷ κρούοντι ἀνοιγήσεται.
\vs{9}Ἢ τίς ἐστιν ἐξ ὑμῶν ἄνθρωπος, ὃν αἰτήσει ὁ υἱὸς αὐτοῦ ἄρτον, μὴ λίθον ἐπιδώσει αὐτῷ;
\vs{10}ἢ καὶ ἰχθὺν αἰτήσει, μὴ ὄφιν ἐπιδώσει αὐτῷ;
\vs{11}εἰ οὖν ὑμεῖς πονηροὶ ὄντες οἴδατε δόματα ἀγαθὰ διδόναι τοῖς τέκνοις ὑμῶν, πόσῳ μᾶλλον ὁ Πατὴρ ὑμῶν ὁ ἐν τοῖς οὐρανοῖς δώσει ἀγαθὰ τοῖς αἰτοῦσιν αὐτόν.

\vs{12}Πάντα οὖν ὅσα ἐὰν θέλητε ἵνα ποιῶσιν ὑμῖν οἱ ἄνθρωποι, οὕτως καὶ ὑμεῖς ποιεῖτε αὐτοῖς· οὗτος γάρ ἐστιν ὁ νόμος καὶ οἱ προφῆται.

\vs{13}Εἰσέλθατε διὰ τῆς στενῆς πύλης· ὅτι πλατεῖα ἡ πύλη καὶ εὐρύχωρος ἡ ὁδὸς ἡ ἀπάγουσα εἰς τὴν ἀπώλειαν καὶ πολλοί εἰσιν οἱ εἰσερχόμενοι δι᾽ αὐτῆς·
\vs{14}ὅτι στενὴ ἡ πύλη καὶ τεθλιμμένη ἡ ὁδὸς ἡ ἀπάγουσα εἰς τὴν ζωήν καὶ ὀλίγοι εἰσὶν οἱ εὑρίσκοντες αὐτήν.

\vs{15}Προσέχετε ἀπὸ τῶν ψευδοπροφητῶν, οἵτινες ἔρχονται πρὸς ὑμᾶς ἐν ἐνδύμασιν προβάτων, ἔσωθεν δέ εἰσιν λύκοι ἅρπαγες.
\vs{16}ἀπὸ τῶν καρπῶν αὐτῶν ἐπιγνώσεσθε αὐτούς. μήτι συλλέγουσιν ἀπὸ ἀκανθῶν σταφυλὰς ἢ ἀπὸ τριβόλων σῦκα;
\vs{17}οὕτως πᾶν δένδρον ἀγαθὸν καρποὺς καλοὺς ποιεῖ, τὸ δὲ σαπρὸν δένδρον καρποὺς πονηροὺς ποιεῖ.
\vs{18}οὐ δύναται δένδρον ἀγαθὸν καρποὺς πονηροὺς ποιεῖν οὐδὲ δένδρον σαπρὸν καρποὺς καλοὺς ποιεῖν.
\vs{19}πᾶν δένδρον μὴ ποιοῦν καρπὸν καλὸν ἐκκόπτεται καὶ εἰς πῦρ βάλλεται.
\vs{20}ἄρα γε ἀπὸ τῶν καρπῶν αὐτῶν ἐπιγνώσεσθε αὐτούς.

\vs{21}Οὐ πᾶς ὁ λέγων μοι· Κύριε Κύριε, εἰσελεύσεται εἰς τὴν βασιλείαν τῶν οὐρανῶν, ἀλλ᾽ ὁ ποιῶν τὸ θέλημα τοῦ Πατρός μου τοῦ ἐν τοῖς οὐρανοῖς.
\vs{22}πολλοὶ ἐροῦσίν μοι ἐν ἐκείνῃ τῇ ἡμέρᾳ· Κύριε Κύριε, οὐ τῷ σῷ ὀνόματι ἐπροφητεύσαμεν, καὶ τῷ σῷ ὀνόματι δαιμόνια ἐξεβάλομεν, καὶ τῷ σῷ ὀνόματι δυνάμεις πολλὰς ἐποιήσαμεν;
\vs{23}καὶ τότε ὁμολογήσω αὐτοῖς ὅτι Οὐδέποτε ἔγνων ὑμᾶς· ἀποχωρεῖτε ἀπ᾽ ἐμοῦ οἱ ἐργαζόμενοι τὴν ἀνομίαν.

\vs{24}Πᾶς οὖν ὅστις ἀκούει μου τοὺς λόγους τούτους καὶ ποιεῖ αὐτούς, ὁμοιωθήσεται ἀνδρὶ φρονίμῳ, ὅστις ᾠκοδόμησεν αὐτοῦ τὴν οἰκίαν ἐπὶ τὴν πέτραν·
\vs{25}καὶ κατέβη ἡ βροχὴ καὶ ἦλθον οἱ ποταμοὶ καὶ ἔπνευσαν οἱ ἄνεμοι καὶ προσέπεσαν τῇ οἰκίᾳ ἐκείνῃ, καὶ οὐκ ἔπεσεν, τεθεμελίωτο γὰρ ἐπὶ τὴν πέτραν.
\vs{26}καὶ πᾶς ὁ ἀκούων μου τοὺς λόγους τούτους καὶ μὴ ποιῶν αὐτοὺς ὁμοιωθήσεται ἀνδρὶ μωρῷ, ὅστις ᾠκοδόμησεν αὐτοῦ τὴν οἰκίαν ἐπὶ τὴν ἄμμον·
\vs{27}καὶ κατέβη ἡ βροχὴ καὶ ἦλθον οἱ ποταμοὶ καὶ ἔπνευσαν οἱ ἄνεμοι καὶ προσέκοψαν τῇ οἰκίᾳ ἐκείνῃ, καὶ ἔπεσεν καὶ ἦν ἡ πτῶσις αὐτῆς μεγάλη.
\vs{28}Καὶ ἐγένετο ὅτε ἐτέλεσεν ὁ Ἰησοῦς τοὺς λόγους τούτους, ἐξεπλήσσοντο οἱ ὄχλοι ἐπὶ τῇ διδαχῇ αὐτοῦ·
\vs{29}ἦν γὰρ διδάσκων αὐτοὺς ὡς ἐξουσίαν ἔχων καὶ οὐχ ὡς οἱ γραμματεῖς αὐτῶν.

\ch{8}
Καταβάντος δὲ αὐτοῦ ἀπὸ τοῦ ὄρους ἠκολούθησαν αὐτῷ ὄχλοι πολλοί.
\vs{2}καὶ ἰδοὺ λεπρὸς προσελθὼν προσεκύνει αὐτῷ λέγων· Κύριε, ἐὰν θέλῃς δύνασαί με καθαρίσαι.
\vs{3}Καὶ ἐκτείνας τὴν χεῖρα ἥψατο αὐτοῦ λέγων· Θέλω, καθαρίσθητι· καὶ εὐθέως ἐκαθαρίσθη αὐτοῦ ἡ λέπρα.
\vs{4}καὶ λέγει αὐτῷ ὁ Ἰησοῦς· Ὅρα μηδενὶ εἴπῃς, ἀλλὰ ὕπαγε σεαυτὸν δεῖξον τῷ ἱερεῖ καὶ προσένεγκον τὸ δῶρον ὃ προσέταξεν Μωϋσῆς, εἰς μαρτύριον αὐτοῖς.

\vs{5}Εἰσελθόντος δὲ αὐτοῦ εἰς Καφαρναοὺμ προσῆλθεν αὐτῷ ἑκατόνταρχος παρακαλῶν αὐτὸν
\vs{6}καὶ λέγων· Κύριε, ὁ παῖς μου βέβληται ἐν τῇ οἰκίᾳ παραλυτικός, δεινῶς βασανιζόμενος.
\vs{7}Καὶ λέγει αὐτῷ· Ἐγὼ ἐλθὼν θεραπεύσω αὐτόν.
\vs{8}Καὶ ἀποκριθεὶς ὁ ἑκατόνταρχος ἔφη· Κύριε, οὐκ εἰμὶ ἱκανὸς ἵνα μου ὑπὸ τὴν στέγην εἰσέλθῃς, ἀλλὰ μόνον εἰπὲ λόγῳ, καὶ ἰαθήσεται ὁ παῖς μου.
\vs{9}καὶ γὰρ ἐγὼ ἄνθρωπός εἰμι ὑπὸ ἐξουσίαν, ἔχων ὑπ᾽ ἐμαυτὸν στρατιώτας, καὶ λέγω τούτῳ· Πορεύθητι, καὶ πορεύεται, καὶ ἄλλῳ· Ἔρχου, καὶ ἔρχεται, καὶ τῷ δούλῳ μου· Ποίησον τοῦτο, καὶ ποιεῖ.
\vs{10}Ἀκούσας δὲ ὁ Ἰησοῦς ἐθαύμασεν καὶ εἶπεν τοῖς ἀκολουθοῦσιν· Ἀμὴν λέγω ὑμῖν, παρ᾽ οὐδενὶ τοσαύτην πίστιν ἐν τῷ Ἰσραὴλ εὗρον.
\vs{11}λέγω δὲ ὑμῖν ὅτι πολλοὶ ἀπὸ ἀνατολῶν καὶ δυσμῶν ἥξουσιν καὶ ἀνακλιθήσονται μετὰ Ἀβραὰμ καὶ Ἰσαὰκ καὶ Ἰακὼβ ἐν τῇ βασιλείᾳ τῶν οὐρανῶν,
\vs{12}οἱ δὲ υἱοὶ τῆς βασιλείας ἐκβληθήσονται εἰς τὸ σκότος τὸ ἐξώτερον· ἐκεῖ ἔσται ὁ κλαυθμὸς καὶ ὁ βρυγμὸς τῶν ὀδόντων.
\vs{13}Καὶ εἶπεν ὁ Ἰησοῦς τῷ ἑκατοντάρχῃ· Ὕπαγε, ὡς ἐπίστευσας γενηθήτω σοι. καὶ ἰάθη ὁ παῖς αὐτοῦ ἐν τῇ ὥρᾳ ἐκείνῃ.

\vs{14}Καὶ ἐλθὼν ὁ Ἰησοῦς εἰς τὴν οἰκίαν Πέτρου εἶδεν τὴν πενθερὰν αὐτοῦ βεβλημένην καὶ πυρέσσουσαν·
\vs{15}καὶ ἥψατο τῆς χειρὸς αὐτῆς, καὶ ἀφῆκεν αὐτὴν ὁ πυρετός, καὶ ἠγέρθη καὶ διηκόνει αὐτῷ.

\vs{16}Ὀψίας δὲ γενομένης προσήνεγκαν αὐτῷ δαιμονιζομένους πολλούς· καὶ ἐξέβαλεν τὰ πνεύματα λόγῳ καὶ πάντας τοὺς κακῶς ἔχοντας ἐθεράπευσεν,
\vs{17}ὅπως πληρωθῇ τὸ ῥηθὲν διὰ Ἠσαΐου τοῦ προφήτου λέγοντος· 
\begin{poetryblock}

\begin{quote}Αὐτὸς τὰς ἀσθενείας ἡμῶν ἔλαβεν\end{quote} 

\begin{quote}καὶ τὰς νόσους ἐβάστασεν.\end{quote}
\end{poetryblock}

\vs{18}Ἰδὼν δὲ ὁ Ἰησοῦς ὄχλον περὶ αὐτὸν ἐκέλευσεν ἀπελθεῖν εἰς τὸ πέραν.
\vs{19}Καὶ προσελθὼν εἷς γραμματεὺς εἶπεν αὐτῷ· Διδάσκαλε, ἀκολουθήσω σοι ὅπου ἐὰν ἀπέρχῃ.
\vs{20}Καὶ λέγει αὐτῷ ὁ Ἰησοῦς· Αἱ ἀλώπεκες φωλεοὺς ἔχουσιν καὶ τὰ πετεινὰ τοῦ οὐρανοῦ κατασκηνώσεις, ὁ δὲ Υἱὸς τοῦ ἀνθρώπου οὐκ ἔχει ποῦ τὴν κεφαλὴν κλίνῃ.
\vs{21}Ἕτερος δὲ τῶν μαθητῶν αὐτοῦ εἶπεν αὐτῷ· Κύριε, ἐπίτρεψόν μοι πρῶτον ἀπελθεῖν καὶ θάψαι τὸν πατέρα μου.
\vs{22}Ὁ δὲ Ἰησοῦς λέγει αὐτῷ· Ἀκολούθει μοι καὶ ἄφες τοὺς νεκροὺς θάψαι τοὺς ἑαυτῶν νεκρούς.

\vs{23}Καὶ ἐμβάντι αὐτῷ εἰς τὸ πλοῖον ἠκολούθησαν αὐτῷ οἱ μαθηταὶ αὐτοῦ.
\vs{24}καὶ ἰδοὺ σεισμὸς μέγας ἐγένετο ἐν τῇ θαλάσσῃ, ὥστε τὸ πλοῖον καλύπτεσθαι ὑπὸ τῶν κυμάτων, αὐτὸς δὲ ἐκάθευδεν.
\vs{25}καὶ προσελθόντες ἤγειραν αὐτὸν λέγοντες· Κύριε, σῶσον, ἀπολλύμεθα.
\vs{26}Καὶ λέγει αὐτοῖς· Τί δειλοί ἐστε, ὀλιγόπιστοι; τότε ἐγερθεὶς ἐπετίμησεν τοῖς ἀνέμοις καὶ τῇ θαλάσσῃ, καὶ ἐγένετο γαλήνη μεγάλη.
\vs{27}Οἱ δὲ ἄνθρωποι ἐθαύμασαν λέγοντες· Ποταπός ἐστιν οὗτος ὅτι καὶ οἱ ἄνεμοι καὶ ἡ θάλασσα αὐτῷ ὑπακούουσιν;

\vs{28}Καὶ ἐλθόντος αὐτοῦ εἰς τὸ πέραν εἰς τὴν χώραν τῶν Γαδαρηνῶν ὑπήντησαν αὐτῷ δύο δαιμονιζόμενοι ἐκ τῶν μνημείων ἐξερχόμενοι, χαλεποὶ λίαν, ὥστε μὴ ἰσχύειν τινὰ παρελθεῖν διὰ τῆς ὁδοῦ ἐκείνης.
\vs{29}Καὶ ἰδοὺ ἔκραξαν λέγοντες· Τί ἡμῖν καὶ σοί, Υἱὲ τοῦ Θεοῦ; ἦλθες ὧδε πρὸ καιροῦ βασανίσαι ἡμᾶς;
\vs{30}Ἦν δὲ μακρὰν ἀπ᾽ αὐτῶν ἀγέλη χοίρων πολλῶν βοσκομένη.
\vs{31}οἱ δὲ δαίμονες παρεκάλουν αὐτὸν λέγοντες· Εἰ ἐκβάλλεις ἡμᾶς, ἀπόστειλον ἡμᾶς εἰς τὴν ἀγέλην τῶν χοίρων.
\vs{32}Καὶ εἶπεν αὐτοῖς· Ὑπάγετε. οἱ δὲ ἐξελθόντες ἀπῆλθον εἰς τοὺς χοίρους· καὶ ἰδοὺ ὥρμησεν πᾶσα ἡ ἀγέλη κατὰ τοῦ κρημνοῦ εἰς τὴν θάλασσαν καὶ ἀπέθανον ἐν τοῖς ὕδασιν.
\vs{33}Οἱ δὲ βόσκοντες ἔφυγον, καὶ ἀπελθόντες εἰς τὴν πόλιν ἀπήγγειλαν πάντα καὶ τὰ τῶν δαιμονιζομένων.
\vs{34}καὶ ἰδοὺ πᾶσα ἡ πόλις ἐξῆλθεν εἰς ὑπάντησιν τῷ Ἰησοῦ καὶ ἰδόντες αὐτὸν παρεκάλεσαν ὅπως μεταβῇ ἀπὸ τῶν ὁρίων αὐτῶν.

\ch{9}
Καὶ ἐμβὰς εἰς πλοῖον διεπέρασεν καὶ ἦλθεν εἰς τὴν ἰδίαν πόλιν.

\vs{2}Καὶ ἰδοὺ προσέφερον αὐτῷ παραλυτικὸν ἐπὶ κλίνης βεβλημένον. καὶ ἰδὼν ὁ Ἰησοῦς τὴν πίστιν αὐτῶν εἶπεν τῷ παραλυτικῷ· Θάρσει, τέκνον, ἀφίενταί σου αἱ ἁμαρτίαι.
\vs{3}Καὶ ἰδού τινες τῶν γραμματέων εἶπαν ἐν ἑαυτοῖς· Οὗτος βλασφημεῖ.
\vs{4}Καὶ ἰδὼν ὁ Ἰησοῦς τὰς ἐνθυμήσεις αὐτῶν εἶπεν· Ἵνατί ἐνθυμεῖσθε πονηρὰ ἐν ταῖς καρδίαις ὑμῶν;
\vs{5}τί γάρ ἐστιν εὐκοπώτερον, εἰπεῖν· Ἀφίενταί σου αἱ ἁμαρτίαι, ἢ εἰπεῖν· Ἔγειρε καὶ περιπάτει;
\vs{6}ἵνα δὲ εἰδῆτε ὅτι ἐξουσίαν ἔχει ὁ Υἱὸς τοῦ ἀνθρώπου ἐπὶ τῆς γῆς ἀφιέναι ἁμαρτίας— τότε λέγει τῷ παραλυτικῷ· Ἐγερθεὶς ἆρόν σου τὴν κλίνην καὶ ὕπαγε εἰς τὸν οἶκόν σου.
\vs{7}καὶ ἐγερθεὶς ἀπῆλθεν εἰς τὸν οἶκον αὐτοῦ.
\vs{8}Ἰδόντες δὲ οἱ ὄχλοι ἐφοβήθησαν καὶ ἐδόξασαν τὸν Θεὸν τὸν δόντα ἐξουσίαν τοιαύτην τοῖς ἀνθρώποις.

\vs{9}Καὶ παράγων ὁ Ἰησοῦς ἐκεῖθεν εἶδεν ἄνθρωπον καθήμενον ἐπὶ τὸ τελώνιον, Μαθθαῖον λεγόμενον, καὶ λέγει αὐτῷ· Ἀκολούθει μοι. καὶ ἀναστὰς ἠκολούθησεν αὐτῷ.

\vs{10}Καὶ ἐγένετο αὐτοῦ ἀνακειμένου ἐν τῇ οἰκίᾳ, καὶ ἰδοὺ πολλοὶ τελῶναι καὶ ἁμαρτωλοὶ ἐλθόντες συνανέκειντο τῷ Ἰησοῦ καὶ τοῖς μαθηταῖς αὐτοῦ.
\vs{11}καὶ ἰδόντες οἱ Φαρισαῖοι ἔλεγον τοῖς μαθηταῖς αὐτοῦ· Διὰ τί μετὰ τῶν τελωνῶν καὶ ἁμαρτωλῶν ἐσθίει ὁ διδάσκαλος ὑμῶν;
\vs{12}Ὁ δὲ ἀκούσας εἶπεν· Οὐ χρείαν ἔχουσιν οἱ ἰσχύοντες ἰατροῦ ἀλλ᾽ οἱ κακῶς ἔχοντες.
\vs{13}πορευθέντες δὲ μάθετε τί ἐστιν· Ἔλεος θέλω καὶ οὐ θυσίαν· οὐ γὰρ ἦλθον καλέσαι δικαίους ἀλλὰ ἁμαρτωλούς.
\vs{14}Τότε προσέρχονται αὐτῷ οἱ μαθηταὶ Ἰωάννου λέγοντες· Διὰ τί ἡμεῖς καὶ οἱ Φαρισαῖοι νηστεύομεν πολλά, οἱ δὲ μαθηταί σου οὐ νηστεύουσιν;
\vs{15}Καὶ εἶπεν αὐτοῖς ὁ Ἰησοῦς· Μὴ δύνανται οἱ υἱοὶ τοῦ νυμφῶνος πενθεῖν ἐφ᾽ ὅσον μετ᾽ αὐτῶν ἐστιν ὁ νυμφίος; ἐλεύσονται δὲ ἡμέραι ὅταν ἀπαρθῇ ἀπ᾽ αὐτῶν ὁ νυμφίος, καὶ τότε νηστεύσουσιν.
\vs{16}Οὐδεὶς δὲ ἐπιβάλλει ἐπίβλημα ῥάκους ἀγνάφου ἐπὶ ἱματίῳ παλαιῷ· αἴρει γὰρ τὸ πλήρωμα αὐτοῦ ἀπὸ τοῦ ἱματίου καὶ χεῖρον σχίσμα γίνεται.
\vs{17}Οὐδὲ βάλλουσιν οἶνον νέον εἰς ἀσκοὺς παλαιούς· εἰ δὲ μή γε, ῥήγνυνται οἱ ἀσκοί καὶ ὁ οἶνος ἐκχεῖται καὶ οἱ ἀσκοὶ ἀπόλλυνται· ἀλλὰ βάλλουσιν οἶνον νέον εἰς ἀσκοὺς καινούς, καὶ ἀμφότεροι συντηροῦνται.

\vs{18}Ταῦτα αὐτοῦ λαλοῦντος αὐτοῖς, ἰδοὺ ἄρχων εἷς ἐλθὼν προσεκύνει αὐτῷ λέγων ὅτι Ἡ θυγάτηρ μου ἄρτι ἐτελεύτησεν· ἀλλὰ ἐλθὼν ἐπίθες τὴν χεῖρά σου ἐπ᾽ αὐτήν, καὶ ζήσεται.
\vs{19}Καὶ ἐγερθεὶς ὁ Ἰησοῦς ἠκολούθει αὐτῷ καὶ οἱ μαθηταὶ αὐτοῦ.

\vs{20}Καὶ ἰδοὺ γυνὴ αἱμορροοῦσα δώδεκα ἔτη προσελθοῦσα ὄπισθεν ἥψατο τοῦ κρασπέδου τοῦ ἱματίου αὐτοῦ·
\vs{21}ἔλεγεν γὰρ ἐν ἑαυτῇ· Ἐὰν μόνον ἅψωμαι τοῦ ἱματίου αὐτοῦ σωθήσομαι.
\vs{22}Ὁ δὲ Ἰησοῦς στραφεὶς καὶ ἰδὼν αὐτὴν εἶπεν· Θάρσει, θύγατερ· ἡ πίστις σου σέσωκέν σε. καὶ ἐσώθη ἡ γυνὴ ἀπὸ τῆς ὥρας ἐκείνης.

\vs{23}Καὶ ἐλθὼν ὁ Ἰησοῦς εἰς τὴν οἰκίαν τοῦ ἄρχοντος καὶ ἰδὼν τοὺς αὐλητὰς καὶ τὸν ὄχλον θορυβούμενον
\vs{24}ἔλεγεν· Ἀναχωρεῖτε, οὐ γὰρ ἀπέθανεν τὸ κοράσιον ἀλλὰ καθεύδει. καὶ κατεγέλων αὐτοῦ.
\vs{25}Ὅτε δὲ ἐξεβλήθη ὁ ὄχλος εἰσελθὼν ἐκράτησεν τῆς χειρὸς αὐτῆς, καὶ ἠγέρθη τὸ κοράσιον.
\vs{26}καὶ ἐξῆλθεν ἡ φήμη αὕτη εἰς ὅλην τὴν γῆν ἐκείνην.

\vs{27}Καὶ παράγοντι ἐκεῖθεν τῷ Ἰησοῦ ἠκολούθησαν αὐτῷ δύο τυφλοὶ κράζοντες καὶ λέγοντες· Ἐλέησον ἡμᾶς, υἱὸς Δαυίδ.
\vs{28}Ἐλθόντι δὲ εἰς τὴν οἰκίαν προσῆλθον αὐτῷ οἱ τυφλοί, καὶ λέγει αὐτοῖς ὁ Ἰησοῦς· Πιστεύετε ὅτι δύναμαι τοῦτο ποιῆσαι; Λέγουσιν αὐτῷ· Ναί Κύριε.
\vs{29}Τότε ἥψατο τῶν ὀφθαλμῶν αὐτῶν λέγων· Κατὰ τὴν πίστιν ὑμῶν γενηθήτω ὑμῖν.
\vs{30}καὶ ἠνεῴχθησαν αὐτῶν οἱ ὀφθαλμοί. καὶ ἐνεβριμήθη αὐτοῖς ὁ Ἰησοῦς λέγων· Ὁρᾶτε μηδεὶς γινωσκέτω.
\vs{31}οἱ δὲ ἐξελθόντες διεφήμισαν αὐτὸν ἐν ὅλῃ τῇ γῇ ἐκείνῃ.

\vs{32}Αὐτῶν δὲ ἐξερχομένων ἰδοὺ προσήνεγκαν αὐτῷ ἄνθρωπον κωφὸν δαιμονιζόμενον.
\vs{33}καὶ ἐκβληθέντος τοῦ δαιμονίου ἐλάλησεν ὁ κωφός. καὶ ἐθαύμασαν οἱ ὄχλοι λέγοντες· Οὐδέποτε ἐφάνη οὕτως ἐν τῷ Ἰσραήλ.
\vs{34}Οἱ δὲ Φαρισαῖοι ἔλεγον· Ἐν τῷ ἄρχοντι τῶν δαιμονίων ἐκβάλλει τὰ δαιμόνια.

\vs{35}Καὶ περιῆγεν ὁ Ἰησοῦς τὰς πόλεις πάσας καὶ τὰς κώμας διδάσκων ἐν ταῖς συναγωγαῖς αὐτῶν καὶ κηρύσσων τὸ εὐαγγέλιον τῆς βασιλείας καὶ θεραπεύων πᾶσαν νόσον καὶ πᾶσαν μαλακίαν.
\vs{36}Ἰδὼν δὲ τοὺς ὄχλους ἐσπλαγχνίσθη περὶ αὐτῶν, ὅτι ἦσαν ἐσκυλμένοι καὶ ἐρριμμένοι ὡσεὶ πρόβατα μὴ ἔχοντα ποιμένα.
\vs{37}Τότε λέγει τοῖς μαθηταῖς αὐτοῦ· Ὁ μὲν θερισμὸς πολύς, οἱ δὲ ἐργάται ὀλίγοι·
\vs{38}δεήθητε οὖν τοῦ Κυρίου τοῦ θερισμοῦ ὅπως ἐκβάλῃ ἐργάτας εἰς τὸν θερισμὸν αὐτοῦ.

\ch{10}
Καὶ προσκαλεσάμενος τοὺς δώδεκα μαθητὰς αὐτοῦ ἔδωκεν αὐτοῖς ἐξουσίαν πνευμάτων ἀκαθάρτων ὥστε ἐκβάλλειν αὐτὰ καὶ θεραπεύειν πᾶσαν νόσον καὶ πᾶσαν μαλακίαν.
\vs{2}Τῶν δὲ δώδεκα ἀποστόλων τὰ ὀνόματά ἐστιν ταῦτα· πρῶτος Σίμων ὁ λεγόμενος Πέτρος καὶ Ἀνδρέας ὁ ἀδελφὸς αὐτοῦ, καὶ Ἰάκωβος ὁ τοῦ Ζεβεδαίου καὶ Ἰωάννης ὁ ἀδελφὸς αὐτοῦ,
\vs{3}Φίλιππος καὶ Βαρθολομαῖος, Θωμᾶς καὶ Μαθθαῖος ὁ τελώνης, Ἰάκωβος ὁ τοῦ Ἁλφαίου καὶ Θαδδαῖος,
\vs{4}Σίμων ὁ Καναναῖος καὶ Ἰούδας ὁ Ἰσκαριώτης ὁ καὶ παραδοὺς αὐτόν.

\vs{5}Τούτους τοὺς δώδεκα ἀπέστειλεν ὁ Ἰησοῦς παραγγείλας αὐτοῖς λέγων· Εἰς ὁδὸν ἐθνῶν μὴ ἀπέλθητε καὶ εἰς πόλιν Σαμαριτῶν μὴ εἰσέλθητε·
\vs{6}πορεύεσθε δὲ μᾶλλον πρὸς τὰ πρόβατα τὰ ἀπολωλότα οἴκου Ἰσραήλ.
\vs{7}πορευόμενοι δὲ κηρύσσετε λέγοντες ὅτι Ἤγγικεν ἡ βασιλεία τῶν οὐρανῶν.
\vs{8}ἀσθενοῦντας θεραπεύετε, νεκροὺς ἐγείρετε, λεπροὺς καθαρίζετε, δαιμόνια ἐκβάλλετε· δωρεὰν ἐλάβετε, δωρεὰν δότε.
\vs{9}Μὴ κτήσησθε χρυσὸν μηδὲ ἄργυρον μηδὲ χαλκὸν εἰς τὰς ζώνας ὑμῶν,
\vs{10}μὴ πήραν εἰς ὁδὸν μηδὲ δύο χιτῶνας μηδὲ ὑποδήματα μηδὲ ῥάβδον· ἄξιος γὰρ ὁ ἐργάτης τῆς τροφῆς αὐτοῦ.
\vs{11}Εἰς ἣν δ᾽ ἂν πόλιν ἢ κώμην εἰσέλθητε, ἐξετάσατε τίς ἐν αὐτῇ ἄξιός ἐστιν· κἀκεῖ μείνατε ἕως ἂν ἐξέλθητε.
\vs{12}εἰσερχόμενοι δὲ εἰς τὴν οἰκίαν ἀσπάσασθε αὐτήν·
\vs{13}καὶ ἐὰν μὲν ᾖ ἡ οἰκία ἀξία, ἐλθάτω ἡ εἰρήνη ὑμῶν ἐπ᾽ αὐτήν, ἐὰν δὲ μὴ ᾖ ἀξία, ἡ εἰρήνη ὑμῶν πρὸς ὑμᾶς ἐπιστραφήτω.
\vs{14}καὶ ὃς ἂν μὴ δέξηται ὑμᾶς μηδὲ ἀκούσῃ τοὺς λόγους ὑμῶν, ἐξερχόμενοι ἔξω τῆς οἰκίας ἢ τῆς πόλεως ἐκείνης ἐκτινάξατε τὸν κονιορτὸν τῶν ποδῶν ὑμῶν.
\vs{15}ἀμὴν λέγω ὑμῖν, ἀνεκτότερον ἔσται γῇ Σοδόμων καὶ Γομόρρων ἐν ἡμέρᾳ κρίσεως ἢ τῇ πόλει ἐκείνῃ.

\vs{16}Ἰδοὺ ἐγὼ ἀποστέλλω ὑμᾶς ὡς πρόβατα ἐν μέσῳ λύκων· γίνεσθε οὖν φρόνιμοι ὡς οἱ ὄφεις καὶ ἀκέραιοι ὡς αἱ περιστεραί.

\vs{17}Προσέχετε δὲ ἀπὸ τῶν ἀνθρώπων· παραδώσουσιν γὰρ ὑμᾶς εἰς συνέδρια καὶ ἐν ταῖς συναγωγαῖς αὐτῶν μαστιγώσουσιν ὑμᾶς·
\vs{18}καὶ ἐπὶ ἡγεμόνας δὲ καὶ βασιλεῖς ἀχθήσεσθε ἕνεκεν ἐμοῦ εἰς μαρτύριον αὐτοῖς καὶ τοῖς ἔθνεσιν.
\vs{19}ὅταν δὲ παραδῶσιν ὑμᾶς, μὴ μεριμνήσητε πῶς ἢ τί λαλήσητε· δοθήσεται γὰρ ὑμῖν ἐν ἐκείνῃ τῇ ὥρᾳ τί λαλήσητε·
\vs{20}οὐ γὰρ ὑμεῖς ἐστε οἱ λαλοῦντες ἀλλὰ τὸ Πνεῦμα τοῦ Πατρὸς ὑμῶν τὸ λαλοῦν ἐν ὑμῖν.

\vs{21}Παραδώσει δὲ ἀδελφὸς ἀδελφὸν εἰς θάνατον καὶ πατὴρ τέκνον, καὶ ἐπαναστήσονται τέκνα ἐπὶ γονεῖς καὶ θανατώσουσιν αὐτούς.
\vs{22}καὶ ἔσεσθε μισούμενοι ὑπὸ πάντων διὰ τὸ ὄνομά μου· ὁ δὲ ὑπομείνας εἰς τέλος οὗτος σωθήσεται.

\vs{23}Ὅταν δὲ διώκωσιν ὑμᾶς ἐν τῇ πόλει ταύτῃ, φεύγετε εἰς τὴν ἑτέραν· ἀμὴν γὰρ λέγω ὑμῖν, οὐ μὴ τελέσητε τὰς πόλεις τοῦ Ἰσραὴλ ἕως ἂν ἔλθῃ ὁ Υἱὸς τοῦ ἀνθρώπου.

\vs{24}Οὐκ ἔστιν μαθητὴς ὑπὲρ τὸν διδάσκαλον οὐδὲ δοῦλος ὑπὲρ τὸν κύριον αὐτοῦ.
\vs{25}ἀρκετὸν τῷ μαθητῇ ἵνα γένηται ὡς ὁ διδάσκαλος αὐτοῦ καὶ ὁ δοῦλος ὡς ὁ κύριος αὐτοῦ. εἰ τὸν οἰκοδεσπότην Βεελζεβοὺλ ἐπεκάλεσαν, πόσῳ μᾶλλον τοὺς οἰκιακοὺς αὐτοῦ.

\vs{26}Μὴ οὖν φοβηθῆτε αὐτούς· οὐδὲν γάρ ἐστιν κεκαλυμμένον ὃ οὐκ ἀποκαλυφθήσεται καὶ κρυπτὸν ὃ οὐ γνωσθήσεται.
\vs{27}ὃ λέγω ὑμῖν ἐν τῇ σκοτίᾳ εἴπατε ἐν τῷ φωτί, καὶ ὃ εἰς τὸ οὖς ἀκούετε κηρύξατε ἐπὶ τῶν δωμάτων.
\vs{28}Καὶ μὴ φοβεῖσθε ἀπὸ τῶν ἀποκτεννόντων τὸ σῶμα, τὴν δὲ ψυχὴν μὴ δυναμένων ἀποκτεῖναι· φοβεῖσθε δὲ μᾶλλον τὸν δυνάμενον καὶ ψυχὴν καὶ σῶμα ἀπολέσαι ἐν γεέννῃ.
\vs{29}Οὐχὶ δύο στρουθία ἀσσαρίου πωλεῖται; καὶ ἓν ἐξ αὐτῶν οὐ πεσεῖται ἐπὶ τὴν γῆν ἄνευ τοῦ Πατρὸς ὑμῶν.
\vs{30}ὑμῶν δὲ καὶ αἱ τρίχες τῆς κεφαλῆς πᾶσαι ἠριθμημέναι εἰσίν.
\vs{31}μὴ οὖν φοβεῖσθε· πολλῶν στρουθίων διαφέρετε ὑμεῖς.

\vs{32}Πᾶς οὖν ὅστις ὁμολογήσει ἐν ἐμοὶ ἔμπροσθεν τῶν ἀνθρώπων, ὁμολογήσω κἀγὼ ἐν αὐτῷ ἔμπροσθεν τοῦ Πατρός μου τοῦ ἐν τοῖς οὐρανοῖς·
\vs{33}ὅστις δ᾽ ἂν ἀρνήσηταί με ἔμπροσθεν τῶν ἀνθρώπων, ἀρνήσομαι κἀγὼ αὐτὸν ἔμπροσθεν τοῦ Πατρός μου τοῦ ἐν τοῖς οὐρανοῖς.

\vs{34}Μὴ νομίσητε ὅτι ἦλθον βαλεῖν εἰρήνην ἐπὶ τὴν γῆν· οὐκ ἦλθον βαλεῖν εἰρήνην ἀλλὰ μάχαιραν.
\vs{35}ἦλθον γὰρ διχάσαι Ἄνθρωπον κατὰ τοῦ πατρὸς αὐτοῦ Καὶ θυγατέρα κατὰ τῆς μητρὸς αὐτῆς Καὶ νύμφην κατὰ τῆς πενθερᾶς αὐτῆς,
\vs{36}Καὶ ἐχθροὶ τοῦ ἀνθρώπου οἱ οἰκιακοὶ αὐτοῦ.

\vs{37}Ὁ φιλῶν πατέρα ἢ μητέρα ὑπὲρ ἐμὲ οὐκ ἔστιν μου ἄξιος, καὶ ὁ φιλῶν υἱὸν ἢ θυγατέρα ὑπὲρ ἐμὲ οὐκ ἔστιν μου ἄξιος·
\vs{38}καὶ ὃς οὐ λαμβάνει τὸν σταυρὸν αὐτοῦ καὶ ἀκολουθεῖ ὀπίσω μου, οὐκ ἔστιν μου ἄξιος.
\vs{39}ὁ εὑρὼν τὴν ψυχὴν αὐτοῦ ἀπολέσει αὐτήν, καὶ ὁ ἀπολέσας τὴν ψυχὴν αὐτοῦ ἕνεκεν ἐμοῦ εὑρήσει αὐτήν.

\vs{40}Ὁ δεχόμενος ὑμᾶς ἐμὲ δέχεται, καὶ ὁ ἐμὲ δεχόμενος δέχεται τὸν ἀποστείλαντά με.
\vs{41}ὁ δεχόμενος προφήτην εἰς ὄνομα προφήτου μισθὸν προφήτου λήμψεται, καὶ ὁ δεχόμενος δίκαιον εἰς ὄνομα δικαίου μισθὸν δικαίου λήμψεται.
\vs{42}καὶ ὃς ἂν ποτίσῃ ἕνα τῶν μικρῶν τούτων ποτήριον ψυχροῦ μόνον εἰς ὄνομα μαθητοῦ, ἀμὴν λέγω ὑμῖν, οὐ μὴ ἀπολέσῃ τὸν μισθὸν αὐτοῦ.

\ch{11}
Καὶ ἐγένετο ὅτε ἐτέλεσεν ὁ Ἰησοῦς διατάσσων τοῖς δώδεκα μαθηταῖς αὐτοῦ, μετέβη ἐκεῖθεν τοῦ διδάσκειν καὶ κηρύσσειν ἐν ταῖς πόλεσιν αὐτῶν.

\vs{2}Ὁ δὲ Ἰωάννης ἀκούσας ἐν τῷ δεσμωτηρίῳ τὰ ἔργα τοῦ Χριστοῦ πέμψας διὰ τῶν μαθητῶν αὐτοῦ
\vs{3}εἶπεν αὐτῷ· Σὺ εἶ ὁ ἐρχόμενος ἢ ἕτερον προσδοκῶμεν;
\vs{4}Καὶ ἀποκριθεὶς ὁ Ἰησοῦς εἶπεν αὐτοῖς· Πορευθέντες ἀπαγγείλατε Ἰωάννῃ ἃ ἀκούετε καὶ βλέπετε·
\vs{5}τυφλοὶ ἀναβλέπουσιν καὶ χωλοὶ περιπατοῦσιν, λεπροὶ καθαρίζονται καὶ κωφοὶ ἀκούουσιν, καὶ νεκροὶ ἐγείρονται καὶ πτωχοὶ εὐαγγελίζονται·
\vs{6}καὶ μακάριός ἐστιν ὃς ἐὰν μὴ σκανδαλισθῇ ἐν ἐμοί.

\vs{7}Τούτων δὲ πορευομένων ἤρξατο ὁ Ἰησοῦς λέγειν τοῖς ὄχλοις περὶ Ἰωάννου· Τί ἐξήλθατε εἰς τὴν ἔρημον θεάσασθαι; κάλαμον ὑπὸ ἀνέμου σαλευόμενον;
\vs{8}ἀλλὰ τί ἐξήλθατε ἰδεῖν; ἄνθρωπον ἐν μαλακοῖς ἠμφιεσμένον; ἰδοὺ οἱ τὰ μαλακὰ φοροῦντες ἐν τοῖς οἴκοις τῶν βασιλέων εἰσίν.
\vs{9}ἀλλὰ τί ἐξήλθατε ἰδεῖν; προφήτην; ναί λέγω ὑμῖν, καὶ περισσότερον προφήτου.
\vs{10}οὗτός ἐστιν περὶ οὗ γέγραπται· 
\begin{poetryblock}

\begin{quote}Ἰδοὺ ἐγὼ ἀποστέλλω τὸν ἄγγελόν μου πρὸ προσώπου σου,\end{quote} 

\begin{quote}Ὃς κατασκευάσει τὴν ὁδόν σου ἔμπροσθέν σου.\end{quote}
\end{poetryblock}

\vs{11}Ἀμὴν λέγω ὑμῖν· οὐκ ἐγήγερται ἐν γεννητοῖς γυναικῶν μείζων Ἰωάννου τοῦ Βαπτιστοῦ· ὁ δὲ μικρότερος ἐν τῇ βασιλείᾳ τῶν οὐρανῶν μείζων αὐτοῦ ἐστιν.
\vs{12}ἀπὸ δὲ τῶν ἡμερῶν Ἰωάννου τοῦ Βαπτιστοῦ ἕως ἄρτι ἡ βασιλεία τῶν οὐρανῶν βιάζεται καὶ βιασταὶ ἁρπάζουσιν αὐτήν.
\vs{13}πάντες γὰρ οἱ προφῆται καὶ ὁ νόμος ἕως Ἰωάννου ἐπροφήτευσαν·
\vs{14}καὶ εἰ θέλετε δέξασθαι, αὐτός ἐστιν Ἠλίας ὁ μέλλων ἔρχεσθαι.
\vs{15}Ὁ ἔχων ὦτα ἀκουέτω.

\vs{16}Τίνι δὲ ὁμοιώσω τὴν γενεὰν ταύτην; ὁμοία ἐστὶν παιδίοις καθημένοις ἐν ταῖς ἀγοραῖς ἃ προσφωνοῦντα τοῖς ἑτέροις
\vs{17}λέγουσιν· 
\begin{poetryblock}

\begin{quote}Ηὐλήσαμεν ὑμῖν Καὶ οὐκ ὠρχήσασθε,\end{quote} 

\begin{quote}Ἐθρηνήσαμεν Καὶ οὐκ ἐκόψασθε.\end{quote}
\end{poetryblock}

\vs{18}Ἦλθεν γὰρ Ἰωάννης μήτε ἐσθίων μήτε πίνων, καὶ λέγουσιν· Δαιμόνιον ἔχει.
\vs{19}ἦλθεν ὁ Υἱὸς τοῦ ἀνθρώπου ἐσθίων καὶ πίνων, καὶ λέγουσιν· Ἰδοὺ ἄνθρωπος φάγος καὶ οἰνοπότης, τελωνῶν φίλος καὶ ἁμαρτωλῶν. καὶ ἐδικαιώθη ἡ σοφία ἀπὸ τῶν ἔργων αὐτῆς.

\vs{20}Τότε ἤρξατο ὀνειδίζειν τὰς πόλεις ἐν αἷς ἐγένοντο αἱ πλεῖσται δυνάμεις αὐτοῦ, ὅτι οὐ μετενόησαν·
\vs{21}Οὐαί σοι, Χοραζίν, οὐαί σοι, Βηθσαϊδά· ὅτι εἰ ἐν Τύρῳ καὶ Σιδῶνι ἐγένοντο αἱ δυνάμεις αἱ γενόμεναι ἐν ὑμῖν, πάλαι ἂν ἐν σάκκῳ καὶ σποδῷ μετενόησαν.
\vs{22}πλὴν λέγω ὑμῖν, Τύρῳ καὶ Σιδῶνι ἀνεκτότερον ἔσται ἐν ἡμέρᾳ κρίσεως ἢ ὑμῖν.
\vs{23}Καὶ σύ, Καφαρναούμ, μὴ ἕως οὐρανοῦ ὑψωθήσῃ; ἕως ᾅδου καταβήσῃ· ὅτι εἰ ἐν Σοδόμοις ἐγενήθησαν αἱ δυνάμεις αἱ γενόμεναι ἐν σοί, ἔμεινεν ἂν μέχρι τῆς σήμερον.
\vs{24}πλὴν λέγω ὑμῖν ὅτι γῇ Σοδόμων ἀνεκτότερον ἔσται ἐν ἡμέρᾳ κρίσεως ἢ σοί.

\vs{25}Ἐν ἐκείνῳ τῷ καιρῷ ἀποκριθεὶς ὁ Ἰησοῦς εἶπεν· Ἐξομολογοῦμαί σοι, Πάτερ, Κύριε τοῦ οὐρανοῦ καὶ τῆς γῆς, ὅτι ἔκρυψας ταῦτα ἀπὸ σοφῶν καὶ συνετῶν καὶ ἀπεκάλυψας αὐτὰ νηπίοις·
\vs{26}ναί ὁ Πατήρ, ὅτι οὕτως εὐδοκία ἐγένετο ἔμπροσθέν σου.
\vs{27}Πάντα μοι παρεδόθη ὑπὸ τοῦ Πατρός μου, καὶ οὐδεὶς ἐπιγινώσκει τὸν Υἱὸν εἰ μὴ ὁ Πατήρ, οὐδὲ τὸν Πατέρα τις ἐπιγινώσκει εἰ μὴ ὁ Υἱὸς καὶ ᾧ ἐὰν βούληται ὁ Υἱὸς ἀποκαλύψαι.

\vs{28}Δεῦτε πρός με πάντες οἱ κοπιῶντες καὶ πεφορτισμένοι, κἀγὼ ἀναπαύσω ὑμᾶς.
\vs{29}ἄρατε τὸν ζυγόν μου ἐφ᾽ ὑμᾶς καὶ μάθετε ἀπ᾽ ἐμοῦ, ὅτι πραΰς εἰμι καὶ ταπεινὸς τῇ καρδίᾳ, καὶ εὑρήσετε ἀνάπαυσιν ταῖς ψυχαῖς ὑμῶν·
\vs{30}ὁ γὰρ ζυγός μου χρηστὸς καὶ τὸ φορτίον μου ἐλαφρόν ἐστιν.

\ch{12}
Ἐν ἐκείνῳ τῷ καιρῷ ἐπορεύθη ὁ Ἰησοῦς τοῖς σάββασιν διὰ τῶν σπορίμων· οἱ δὲ μαθηταὶ αὐτοῦ ἐπείνασαν καὶ ἤρξαντο τίλλειν στάχυας καὶ ἐσθίειν.
\vs{2}οἱ δὲ Φαρισαῖοι ἰδόντες εἶπαν αὐτῷ· Ἰδοὺ οἱ μαθηταί σου ποιοῦσιν ὃ οὐκ ἔξεστιν ποιεῖν ἐν σαββάτῳ.
\vs{3}Ὁ δὲ εἶπεν αὐτοῖς· Οὐκ ἀνέγνωτε τί ἐποίησεν Δαυὶδ ὅτε ἐπείνασεν καὶ οἱ μετ᾽ αὐτοῦ,
\vs{4}πῶς εἰσῆλθεν εἰς τὸν οἶκον τοῦ Θεοῦ καὶ τοὺς ἄρτους τῆς προθέσεως ἔφαγον, ὃ οὐκ ἐξὸν ἦν αὐτῷ φαγεῖν οὐδὲ τοῖς μετ᾽ αὐτοῦ εἰ μὴ τοῖς ἱερεῦσιν μόνοις;
\vs{5}Ἢ οὐκ ἀνέγνωτε ἐν τῷ νόμῳ ὅτι τοῖς σάββασιν οἱ ἱερεῖς ἐν τῷ ἱερῷ τὸ σάββατον βεβηλοῦσιν καὶ ἀναίτιοί εἰσιν;
\vs{6}λέγω δὲ ὑμῖν ὅτι τοῦ ἱεροῦ μεῖζόν ἐστιν ὧδε.
\vs{7}Εἰ δὲ ἐγνώκειτε τί ἐστιν· Ἔλεος θέλω καὶ οὐ θυσίαν, οὐκ ἂν κατεδικάσατε τοὺς ἀναιτίους.
\vs{8}κύριος γάρ ἐστιν τοῦ σαββάτου ὁ Υἱὸς τοῦ ἀνθρώπου.

\vs{9}Καὶ μεταβὰς ἐκεῖθεν ἦλθεν εἰς τὴν συναγωγὴν αὐτῶν·
\vs{10}καὶ ἰδοὺ ἄνθρωπος χεῖρα ἔχων ξηράν. καὶ ἐπηρώτησαν αὐτὸν λέγοντες· Εἰ ἔξεστιν τοῖς σάββασιν θεραπεῦσαι; ἵνα κατηγορήσωσιν αὐτοῦ.
\vs{11}Ὁ δὲ εἶπεν αὐτοῖς· Τίς ἔσται ἐξ ὑμῶν ἄνθρωπος ὃς ἕξει πρόβατον ἕν καὶ ἐὰν ἐμπέσῃ τοῦτο τοῖς σάββασιν εἰς βόθυνον, οὐχὶ κρατήσει αὐτὸ καὶ ἐγερεῖ;
\vs{12}πόσῳ οὖν διαφέρει ἄνθρωπος προβάτου. ὥστε ἔξεστιν τοῖς σάββασιν καλῶς ποιεῖν.
\vs{13}Τότε λέγει τῷ ἀνθρώπῳ· Ἔκτεινόν σου τὴν χεῖρα. καὶ ἐξέτεινεν καὶ ἀπεκατεστάθη ὑγιὴς ὡς ἡ ἄλλη.
\vs{14}ἐξελθόντες δὲ οἱ Φαρισαῖοι συμβούλιον ἔλαβον κατ᾽ αὐτοῦ ὅπως αὐτὸν ἀπολέσωσιν.
\vs{15}Ὁ δὲ Ἰησοῦς γνοὺς ἀνεχώρησεν ἐκεῖθεν. καὶ ἠκολούθησαν αὐτῷ ὄχλοι πολλοί, καὶ ἐθεράπευσεν αὐτοὺς πάντας
\vs{16}καὶ ἐπετίμησεν αὐτοῖς ἵνα μὴ φανερὸν αὐτὸν ποιήσωσιν,
\vs{17}ἵνα πληρωθῇ τὸ ῥηθὲν διὰ Ἠσαΐου τοῦ προφήτου λέγοντος·
\begin{poetryblock}

\begin{quote} \vs{18}Ἰδοὺ ὁ παῖς μου ὃν ᾑρέτισα,\end{quote} 

\begin{quote}ὁ ἀγαπητός μου εἰς ὃν εὐδόκησεν ἡ ψυχή μου·\end{quote} 

\begin{quote}θήσω τὸ Πνεῦμά μου ἐπ᾽ αὐτόν,\end{quote} 

\begin{quote}καὶ κρίσιν τοῖς ἔθνεσιν ἀπαγγελεῖ.\end{quote}

\begin{quote} \vs{19}οὐκ ἐρίσει οὐδὲ κραυγάσει,\end{quote} 

\begin{quote}οὐδὲ ἀκούσει τις ἐν ταῖς πλατείαις τὴν φωνὴν αὐτοῦ.\end{quote}

\begin{quote} \vs{20}κάλαμον συντετριμμένον οὐ κατεάξει\end{quote} 

\begin{quote}καὶ λίνον τυφόμενον οὐ σβέσει,\end{quote} 

\begin{quote}ἕως ἂν ἐκβάλῃ εἰς νῖκος τὴν κρίσιν.\end{quote}

\begin{quote} \vs{21}καὶ τῷ ὀνόματι αὐτοῦ ἔθνη ἐλπιοῦσιν.\end{quote}
\end{poetryblock}

\vs{22}Τότε προσηνέχθη αὐτῷ δαιμονιζόμενος τυφλὸς καὶ κωφός, καὶ ἐθεράπευσεν αὐτόν, ὥστε τὸν κωφὸν λαλεῖν καὶ βλέπειν.
\vs{23}καὶ ἐξίσταντο πάντες οἱ ὄχλοι καὶ ἔλεγον· Μήτι οὗτός ἐστιν ὁ υἱὸς Δαυίδ;
\vs{24}Οἱ δὲ Φαρισαῖοι ἀκούσαντες εἶπον· Οὗτος οὐκ ἐκβάλλει τὰ δαιμόνια εἰ μὴ ἐν τῷ Βεελζεβοὺλ ἄρχοντι τῶν δαιμονίων.
\vs{25}Εἰδὼς δὲ τὰς ἐνθυμήσεις αὐτῶν εἶπεν αὐτοῖς· Πᾶσα βασιλεία μερισθεῖσα καθ᾽ ἑαυτῆς ἐρημοῦται καὶ πᾶσα πόλις ἢ οἰκία μερισθεῖσα καθ᾽ ἑαυτῆς οὐ σταθήσεται.
\vs{26}καὶ εἰ ὁ Σατανᾶς τὸν Σατανᾶν ἐκβάλλει, ἐφ᾽ ἑαυτὸν ἐμερίσθη· πῶς οὖν σταθήσεται ἡ βασιλεία αὐτοῦ;
\vs{27}Καὶ εἰ ἐγὼ ἐν Βεελζεβοὺλ ἐκβάλλω τὰ δαιμόνια, οἱ υἱοὶ ὑμῶν ἐν τίνι ἐκβάλλουσιν; διὰ τοῦτο αὐτοὶ κριταὶ ἔσονται ὑμῶν.
\vs{28}εἰ δὲ ἐν Πνεύματι Θεοῦ ἐγὼ ἐκβάλλω τὰ δαιμόνια, ἄρα ἔφθασεν ἐφ᾽ ὑμᾶς ἡ βασιλεία τοῦ Θεοῦ.
\vs{29}Ἢ πῶς δύναταί τις εἰσελθεῖν εἰς τὴν οἰκίαν τοῦ ἰσχυροῦ καὶ τὰ σκεύη αὐτοῦ ἁρπάσαι, ἐὰν μὴ πρῶτον δήσῃ τὸν ἰσχυρόν; καὶ τότε τὴν οἰκίαν αὐτοῦ διαρπάσει.
\vs{30}Ὁ μὴ ὢν μετ᾽ ἐμοῦ κατ᾽ ἐμοῦ ἐστιν, καὶ ὁ μὴ συνάγων μετ᾽ ἐμοῦ σκορπίζει.
\vs{31}Διὰ τοῦτο λέγω ὑμῖν, πᾶσα ἁμαρτία καὶ βλασφημία ἀφεθήσεται τοῖς ἀνθρώποις, ἡ δὲ τοῦ Πνεύματος βλασφημία οὐκ ἀφεθήσεται.
\vs{32}καὶ ὃς ἐὰν εἴπῃ λόγον κατὰ τοῦ Υἱοῦ τοῦ ἀνθρώπου, ἀφεθήσεται αὐτῷ· ὃς δ᾽ ἂν εἴπῃ κατὰ τοῦ Πνεύματος τοῦ Ἁγίου, οὐκ ἀφεθήσεται αὐτῷ οὔτε ἐν τούτῳ τῷ αἰῶνι οὔτε ἐν τῷ μέλλοντι.

\vs{33}Ἢ ποιήσατε τὸ δένδρον καλὸν καὶ τὸν καρπὸν αὐτοῦ καλόν, ἢ ποιήσατε τὸ δένδρον σαπρὸν καὶ τὸν καρπὸν αὐτοῦ σαπρόν· ἐκ γὰρ τοῦ καρποῦ τὸ δένδρον γινώσκεται.
\vs{34}γεννήματα ἐχιδνῶν, πῶς δύνασθε ἀγαθὰ λαλεῖν πονηροὶ ὄντες; ἐκ γὰρ τοῦ περισσεύματος τῆς καρδίας τὸ στόμα λαλεῖ.
\vs{35}ὁ ἀγαθὸς ἄνθρωπος ἐκ τοῦ ἀγαθοῦ θησαυροῦ ἐκβάλλει ἀγαθά, καὶ ὁ πονηρὸς ἄνθρωπος ἐκ τοῦ πονηροῦ θησαυροῦ ἐκβάλλει πονηρά.
\vs{36}λέγω δὲ ὑμῖν ὅτι πᾶν ῥῆμα ἀργὸν ὃ λαλήσουσιν οἱ ἄνθρωποι ἀποδώσουσιν περὶ αὐτοῦ λόγον ἐν ἡμέρᾳ κρίσεως·
\vs{37}ἐκ γὰρ τῶν λόγων σου δικαιωθήσῃ, καὶ ἐκ τῶν λόγων σου καταδικασθήσῃ.

\vs{38}Τότε ἀπεκρίθησαν αὐτῷ τινες τῶν γραμματέων καὶ Φαρισαίων λέγοντες· Διδάσκαλε, θέλομεν ἀπὸ σοῦ σημεῖον ἰδεῖν.
\vs{39}Ὁ δὲ ἀποκριθεὶς εἶπεν αὐτοῖς· Γενεὰ πονηρὰ καὶ μοιχαλὶς σημεῖον ἐπιζητεῖ, καὶ σημεῖον οὐ δοθήσεται αὐτῇ εἰ μὴ τὸ σημεῖον Ἰωνᾶ τοῦ προφήτου.
\vs{40}ὥσπερ γὰρ ἦν Ἰωνᾶς ἐν τῇ κοιλίᾳ τοῦ κήτους τρεῖς ἡμέρας καὶ τρεῖς νύκτας, οὕτως ἔσται ὁ Υἱὸς τοῦ ἀνθρώπου ἐν τῇ καρδίᾳ τῆς γῆς τρεῖς ἡμέρας καὶ τρεῖς νύκτας.
\vs{41}Ἄνδρες Νινευῖται ἀναστήσονται ἐν τῇ κρίσει μετὰ τῆς γενεᾶς ταύτης καὶ κατακρινοῦσιν αὐτήν, ὅτι μετενόησαν εἰς τὸ κήρυγμα Ἰωνᾶ, καὶ ἰδοὺ πλεῖον Ἰωνᾶ ὧδε.
\vs{42}βασίλισσα νότου ἐγερθήσεται ἐν τῇ κρίσει μετὰ τῆς γενεᾶς ταύτης καὶ κατακρινεῖ αὐτήν, ὅτι ἦλθεν ἐκ τῶν περάτων τῆς γῆς ἀκοῦσαι τὴν σοφίαν Σολομῶνος, καὶ ἰδοὺ πλεῖον Σολομῶνος ὧδε.

\vs{43}Ὅταν δὲ τὸ ἀκάθαρτον πνεῦμα ἐξέλθῃ ἀπὸ τοῦ ἀνθρώπου, διέρχεται δι᾽ ἀνύδρων τόπων ζητοῦν ἀνάπαυσιν καὶ οὐχ εὑρίσκει.
\vs{44}τότε λέγει· Εἰς τὸν οἶκόν μου ἐπιστρέψω ὅθεν ἐξῆλθον· καὶ ἐλθὸν εὑρίσκει σχολάζοντα σεσαρωμένον καὶ κεκοσμημένον.
\vs{45}τότε πορεύεται καὶ παραλαμβάνει μεθ᾽ ἑαυτοῦ ἑπτὰ ἕτερα πνεύματα πονηρότερα ἑαυτοῦ καὶ εἰσελθόντα κατοικεῖ ἐκεῖ· καὶ γίνεται τὰ ἔσχατα τοῦ ἀνθρώπου ἐκείνου χείρονα τῶν πρώτων. οὕτως ἔσται καὶ τῇ γενεᾷ ταύτῃ τῇ πονηρᾷ.

\vs{46}Ἔτι αὐτοῦ λαλοῦντος τοῖς ὄχλοις ἰδοὺ ἡ μήτηρ καὶ οἱ ἀδελφοὶ αὐτοῦ εἱστήκεισαν ἔξω ζητοῦντες αὐτῷ λαλῆσαι.
\vs{47}εἶπεν δέ τις αὐτῷ· Ἰδοὺ ἡ μήτηρ σου καὶ οἱ ἀδελφοί σου ἔξω ἑστήκασιν ζητοῦντές σοι λαλῆσαι.
\vs{48}Ὁ δὲ ἀποκριθεὶς εἶπεν τῷ λέγοντι αὐτῷ· Τίς ἐστιν ἡ μήτηρ μου καὶ τίνες εἰσὶν οἱ ἀδελφοί μου;
\vs{49}καὶ ἐκτείνας τὴν χεῖρα αὐτοῦ ἐπὶ τοὺς μαθητὰς αὐτοῦ εἶπεν· Ἰδοὺ ἡ μήτηρ μου καὶ οἱ ἀδελφοί μου.
\vs{50}ὅστις γὰρ ἂν ποιήσῃ τὸ θέλημα τοῦ Πατρός μου τοῦ ἐν οὐρανοῖς αὐτός μου ἀδελφὸς καὶ ἀδελφὴ καὶ μήτηρ ἐστίν.

\ch{13}
Ἐν τῇ ἡμέρᾳ ἐκείνῃ ἐξελθὼν ὁ Ἰησοῦς τῆς οἰκίας ἐκάθητο παρὰ τὴν θάλασσαν·
\vs{2}καὶ συνήχθησαν πρὸς αὐτὸν ὄχλοι πολλοί, ὥστε αὐτὸν εἰς πλοῖον ἐμβάντα καθῆσθαι, καὶ πᾶς ὁ ὄχλος ἐπὶ τὸν αἰγιαλὸν εἱστήκει.

\vs{3}Καὶ ἐλάλησεν αὐτοῖς πολλὰ ἐν παραβολαῖς λέγων· Ἰδοὺ ἐξῆλθεν ὁ σπείρων τοῦ σπείρειν.
\vs{4}καὶ ἐν τῷ σπείρειν αὐτὸν ἃ μὲν ἔπεσεν παρὰ τὴν ὁδόν, καὶ ἐλθόντα τὰ πετεινὰ κατέφαγεν αὐτά.
\vs{5}Ἄλλα δὲ ἔπεσεν ἐπὶ τὰ πετρώδη ὅπου οὐκ εἶχεν γῆν πολλήν, καὶ εὐθέως ἐξανέτειλεν διὰ τὸ μὴ ἔχειν βάθος γῆς·
\vs{6}ἡλίου δὲ ἀνατείλαντος ἐκαυματίσθη καὶ διὰ τὸ μὴ ἔχειν ῥίζαν ἐξηράνθη.
\vs{7}Ἄλλα δὲ ἔπεσεν ἐπὶ τὰς ἀκάνθας, καὶ ἀνέβησαν αἱ ἄκανθαι καὶ ἔπνιξαν αὐτά.
\vs{8}Ἄλλα δὲ ἔπεσεν ἐπὶ τὴν γῆν τὴν καλὴν καὶ ἐδίδου καρπόν, ὃ μὲν ἑκατὸν, ὃ δὲ ἑξήκοντα, ὃ δὲ τριάκοντα.
\vs{9}Ὁ ἔχων ὦτα ἀκουέτω.

\vs{10}Καὶ προσελθόντες οἱ μαθηταὶ εἶπαν αὐτῷ· Διὰ τί ἐν παραβολαῖς λαλεῖς αὐτοῖς;
\vs{11}Ὁ δὲ ἀποκριθεὶς εἶπεν αὐτοῖς· Ὅτι Ὑμῖν δέδοται γνῶναι τὰ μυστήρια τῆς βασιλείας τῶν οὐρανῶν, ἐκείνοις δὲ οὐ δέδοται.
\vs{12}ὅστις γὰρ ἔχει, δοθήσεται αὐτῷ καὶ περισσευθήσεται· ὅστις δὲ οὐκ ἔχει, καὶ ὃ ἔχει ἀρθήσεται ἀπ᾽ αὐτοῦ.
\vs{13}διὰ τοῦτο ἐν παραβολαῖς αὐτοῖς λαλῶ, Ὅτι βλέποντες οὐ βλέπουσιν Καὶ ἀκούοντες οὐκ ἀκούουσιν οὐδὲ συνίουσιν,
\vs{14}Καὶ ἀναπληροῦται αὐτοῖς ἡ προφητεία Ἠσαΐου ἡ λέγουσα· 
\begin{poetryblock}

\begin{quote}Ἀκοῇ ἀκούσετε καὶ οὐ μὴ συνῆτε,\end{quote} 

\begin{quote}Καὶ βλέποντες βλέψετε καὶ οὐ μὴ ἴδητε.\end{quote}

\begin{quote} \vs{15}Ἐπαχύνθη γὰρ ἡ καρδία τοῦ λαοῦ τούτου,\end{quote} 

\begin{quote}Καὶ τοῖς ὠσὶν βαρέως ἤκουσαν\end{quote} 

\begin{quote}Καὶ τοὺς ὀφθαλμοὺς αὐτῶν ἐκάμμυσαν,\end{quote} 

\begin{quote}Μήποτε ἴδωσιν τοῖς ὀφθαλμοῖς\end{quote} 

\begin{quote}Καὶ τοῖς ὠσὶν ἀκούσωσιν\end{quote} 

\begin{quote}Καὶ τῇ καρδίᾳ συνῶσιν\end{quote} 

\begin{quote}Καὶ ἐπιστρέψωσιν Καὶ ἰάσομαι αὐτούς.\end{quote}
\end{poetryblock}

\vs{16}Ὑμῶν δὲ μακάριοι οἱ ὀφθαλμοὶ ὅτι βλέπουσιν καὶ τὰ ὦτα ὑμῶν ὅτι ἀκούουσιν.
\vs{17}ἀμὴν γὰρ λέγω ὑμῖν ὅτι πολλοὶ προφῆται καὶ δίκαιοι ἐπεθύμησαν ἰδεῖν ἃ βλέπετε καὶ οὐκ εἶδαν, καὶ ἀκοῦσαι ἃ ἀκούετε καὶ οὐκ ἤκουσαν.

\vs{18}Ὑμεῖς οὖν ἀκούσατε τὴν παραβολὴν τοῦ σπείραντος.
\vs{19}Παντὸς ἀκούοντος τὸν λόγον τῆς βασιλείας καὶ μὴ συνιέντος ἔρχεται ὁ πονηρὸς καὶ ἁρπάζει τὸ ἐσπαρμένον ἐν τῇ καρδίᾳ αὐτοῦ, οὗτός ἐστιν ὁ παρὰ τὴν ὁδὸν σπαρείς.
\vs{20}Ὁ δὲ ἐπὶ τὰ πετρώδη σπαρείς, οὗτός ἐστιν ὁ τὸν λόγον ἀκούων καὶ εὐθὺς μετὰ χαρᾶς λαμβάνων αὐτόν,
\vs{21}οὐκ ἔχει δὲ ῥίζαν ἐν ἑαυτῷ ἀλλὰ πρόσκαιρός ἐστιν, γενομένης δὲ θλίψεως ἢ διωγμοῦ διὰ τὸν λόγον εὐθὺς σκανδαλίζεται.
\vs{22}Ὁ δὲ εἰς τὰς ἀκάνθας σπαρείς, οὗτός ἐστιν ὁ τὸν λόγον ἀκούων, καὶ ἡ μέριμνα τοῦ αἰῶνος καὶ ἡ ἀπάτη τοῦ πλούτου συμπνίγει τὸν λόγον καὶ ἄκαρπος γίνεται.
\vs{23}Ὁ δὲ ἐπὶ τὴν καλὴν γῆν σπαρείς, οὗτός ἐστιν ὁ τὸν λόγον ἀκούων καὶ συνιείς, ὃς δὴ καρποφορεῖ καὶ ποιεῖ ὃ μὲν ἑκατὸν, ὃ δὲ ἑξήκοντα, ὃ δὲ τριάκοντα.

\vs{24}Ἄλλην παραβολὴν παρέθηκεν αὐτοῖς λέγων· Ὡμοιώθη ἡ βασιλεία τῶν οὐρανῶν ἀνθρώπῳ σπείραντι καλὸν σπέρμα ἐν τῷ ἀγρῷ αὐτοῦ.
\vs{25}ἐν δὲ τῷ καθεύδειν τοὺς ἀνθρώπους ἦλθεν αὐτοῦ ὁ ἐχθρὸς καὶ ἐπέσπειρεν ζιζάνια ἀνὰ μέσον τοῦ σίτου καὶ ἀπῆλθεν.
\vs{26}ὅτε δὲ ἐβλάστησεν ὁ χόρτος καὶ καρπὸν ἐποίησεν, τότε ἐφάνη καὶ τὰ ζιζάνια.
\vs{27}Προσελθόντες δὲ οἱ δοῦλοι τοῦ οἰκοδεσπότου εἶπον αὐτῷ· Κύριε, οὐχὶ καλὸν σπέρμα ἔσπειρας ἐν τῷ σῷ ἀγρῷ; πόθεν οὖν ἔχει ζιζάνια;
\vs{28}Ὁ δὲ ἔφη αὐτοῖς· Ἐχθρὸς ἄνθρωπος τοῦτο ἐποίησεν. Οἱ δὲ δοῦλοι λέγουσιν αὐτῷ· Θέλεις οὖν ἀπελθόντες συλλέξωμεν αὐτά;
\vs{29}Ὁ δέ φησιν· Οὔ, μήποτε συλλέγοντες τὰ ζιζάνια ἐκριζώσητε ἅμα αὐτοῖς τὸν σῖτον.
\vs{30}ἄφετε συναυξάνεσθαι ἀμφότερα ἕως τοῦ θερισμοῦ, καὶ ἐν καιρῷ τοῦ θερισμοῦ ἐρῶ τοῖς θερισταῖς· Συλλέξατε πρῶτον τὰ ζιζάνια καὶ δήσατε αὐτὰ εἰς δέσμας πρὸς τὸ κατακαῦσαι αὐτά, τὸν δὲ σῖτον συναγάγετε εἰς τὴν ἀποθήκην μου.

\vs{31}Ἄλλην παραβολὴν παρέθηκεν αὐτοῖς λέγων· Ὁμοία ἐστὶν ἡ βασιλεία τῶν οὐρανῶν κόκκῳ σινάπεως, ὃν λαβὼν ἄνθρωπος ἔσπειρεν ἐν τῷ ἀγρῷ αὐτοῦ·
\vs{32}ὃ μικρότερον μέν ἐστιν πάντων τῶν σπερμάτων, ὅταν δὲ αὐξηθῇ μεῖζον τῶν λαχάνων ἐστὶν καὶ γίνεται δένδρον, ὥστε ἐλθεῖν τὰ πετεινὰ τοῦ οὐρανοῦ καὶ κατασκηνοῦν ἐν τοῖς κλάδοις αὐτοῦ.

\vs{33}Ἄλλην παραβολὴν ἐλάλησεν αὐτοῖς· Ὁμοία ἐστὶν ἡ βασιλεία τῶν οὐρανῶν ζύμῃ, ἣν λαβοῦσα γυνὴ ἐνέκρυψεν εἰς ἀλεύρου σάτα τρία ἕως οὗ ἐζυμώθη ὅλον.
\vs{34}Ταῦτα πάντα ἐλάλησεν ὁ Ἰησοῦς ἐν παραβολαῖς τοῖς ὄχλοις καὶ χωρὶς παραβολῆς οὐδὲν ἐλάλει αὐτοῖς,
\vs{35}ὅπως πληρωθῇ τὸ ῥηθὲν διὰ τοῦ προφήτου λέγοντος· 
\begin{poetryblock}

\begin{quote}Ἀνοίξω ἐν παραβολαῖς τὸ στόμα μου,\end{quote} 

\begin{quote}ἐρεύξομαι κεκρυμμένα ἀπὸ καταβολῆς κόσμου.\end{quote}
\end{poetryblock}

\vs{36}Τότε ἀφεὶς τοὺς ὄχλους ἦλθεν εἰς τὴν οἰκίαν. Καὶ προσῆλθον αὐτῷ οἱ μαθηταὶ αὐτοῦ λέγοντες· Διασάφησον ἡμῖν τὴν παραβολὴν τῶν ζιζανίων τοῦ ἀγροῦ.
\vs{37}Ὁ δὲ ἀποκριθεὶς εἶπεν· Ὁ σπείρων τὸ καλὸν σπέρμα ἐστὶν ὁ Υἱὸς τοῦ ἀνθρώπου,
\vs{38}ὁ δὲ ἀγρός ἐστιν ὁ κόσμος, τὸ δὲ καλὸν σπέρμα οὗτοί εἰσιν οἱ υἱοὶ τῆς βασιλείας· τὰ δὲ ζιζάνιά εἰσιν οἱ υἱοὶ τοῦ πονηροῦ,
\vs{39}ὁ δὲ ἐχθρὸς ὁ σπείρας αὐτά ἐστιν ὁ διάβολος, ὁ δὲ θερισμὸς συντέλεια αἰῶνός ἐστιν, οἱ δὲ θερισταὶ ἄγγελοί εἰσιν.
\vs{40}Ὥσπερ οὖν συλλέγεται τὰ ζιζάνια καὶ πυρὶ κατακαίεται, οὕτως ἔσται ἐν τῇ συντελείᾳ τοῦ αἰῶνος·
\vs{41}ἀποστελεῖ ὁ Υἱὸς τοῦ ἀνθρώπου τοὺς ἀγγέλους αὐτοῦ, καὶ συλλέξουσιν ἐκ τῆς βασιλείας αὐτοῦ πάντα τὰ σκάνδαλα καὶ τοὺς ποιοῦντας τὴν ἀνομίαν
\vs{42}καὶ βαλοῦσιν αὐτοὺς εἰς τὴν κάμινον τοῦ πυρός· ἐκεῖ ἔσται ὁ κλαυθμὸς καὶ ὁ βρυγμὸς τῶν ὀδόντων.
\vs{43}τότε οἱ δίκαιοι ἐκλάμψουσιν ὡς ὁ ἥλιος ἐν τῇ βασιλείᾳ τοῦ Πατρὸς αὐτῶν. Ὁ ἔχων ὦτα ἀκουέτω.

\vs{44}Ὁμοία ἐστὶν ἡ βασιλεία τῶν οὐρανῶν θησαυρῷ κεκρυμμένῳ ἐν τῷ ἀγρῷ, ὃν εὑρὼν ἄνθρωπος ἔκρυψεν, καὶ ἀπὸ τῆς χαρᾶς αὐτοῦ ὑπάγει καὶ πωλεῖ πάντα ὅσα ἔχει καὶ ἀγοράζει τὸν ἀγρὸν ἐκεῖνον.

\vs{45}Πάλιν ὁμοία ἐστὶν ἡ βασιλεία τῶν οὐρανῶν ἀνθρώπῳ ἐμπόρῳ ζητοῦντι καλοὺς μαργαρίτας·
\vs{46}εὑρὼν δὲ ἕνα πολύτιμον μαργαρίτην ἀπελθὼν πέπρακεν πάντα ὅσα εἶχεν καὶ ἠγόρασεν αὐτόν.

\vs{47}Πάλιν ὁμοία ἐστὶν ἡ βασιλεία τῶν οὐρανῶν σαγήνῃ βληθείσῃ εἰς τὴν θάλασσαν καὶ ἐκ παντὸς γένους συναγαγούσῃ·
\vs{48}ἣν ὅτε ἐπληρώθη ἀναβιβάσαντες ἐπὶ τὸν αἰγιαλὸν καὶ καθίσαντες συνέλεξαν τὰ καλὰ εἰς ἄγγη, τὰ δὲ σαπρὰ ἔξω ἔβαλον.
\vs{49}Οὕτως ἔσται ἐν τῇ συντελείᾳ τοῦ αἰῶνος· ἐξελεύσονται οἱ ἄγγελοι καὶ ἀφοριοῦσιν τοὺς πονηροὺς ἐκ μέσου τῶν δικαίων
\vs{50}καὶ βαλοῦσιν αὐτοὺς εἰς τὴν κάμινον τοῦ πυρός· ἐκεῖ ἔσται ὁ κλαυθμὸς καὶ ὁ βρυγμὸς τῶν ὀδόντων.

\vs{51}Συνήκατε ταῦτα πάντα; Λέγουσιν αὐτῷ· Ναί.
\vs{52}Ὁ δὲ εἶπεν αὐτοῖς· Διὰ τοῦτο πᾶς γραμματεὺς μαθητευθεὶς τῇ βασιλείᾳ τῶν οὐρανῶν ὅμοιός ἐστιν ἀνθρώπῳ οἰκοδεσπότῃ, ὅστις ἐκβάλλει ἐκ τοῦ θησαυροῦ αὐτοῦ καινὰ καὶ παλαιά.

\vs{53}Καὶ ἐγένετο ὅτε ἐτέλεσεν ὁ Ἰησοῦς τὰς παραβολὰς ταύτας, μετῆρεν ἐκεῖθεν.
\vs{54}καὶ ἐλθὼν εἰς τὴν πατρίδα αὐτοῦ ἐδίδασκεν αὐτοὺς ἐν τῇ συναγωγῇ αὐτῶν, ὥστε ἐκπλήσσεσθαι αὐτοὺς καὶ λέγειν· Πόθεν τούτῳ ἡ σοφία αὕτη καὶ αἱ δυνάμεις;
\vs{55}οὐχ οὗτός ἐστιν ὁ τοῦ τέκτονος υἱός; οὐχ ἡ μήτηρ αὐτοῦ λέγεται Μαριὰμ καὶ οἱ ἀδελφοὶ αὐτοῦ Ἰάκωβος καὶ Ἰωσὴφ καὶ Σίμων καὶ Ἰούδας;
\vs{56}καὶ αἱ ἀδελφαὶ αὐτοῦ οὐχὶ πᾶσαι πρὸς ἡμᾶς εἰσιν; πόθεν οὖν τούτῳ ταῦτα πάντα;
\vs{57}καὶ ἐσκανδαλίζοντο ἐν αὐτῷ. Ὁ δὲ Ἰησοῦς εἶπεν αὐτοῖς· Οὐκ ἔστιν προφήτης ἄτιμος εἰ μὴ ἐν τῇ πατρίδι καὶ ἐν τῇ οἰκίᾳ αὐτοῦ.
\vs{58}καὶ οὐκ ἐποίησεν ἐκεῖ δυνάμεις πολλὰς διὰ τὴν ἀπιστίαν αὐτῶν.

\ch{14}
Ἐν ἐκείνῳ τῷ καιρῷ ἤκουσεν Ἡρῴδης ὁ τετραάρχης τὴν ἀκοὴν Ἰησοῦ,
\vs{2}καὶ εἶπεν τοῖς παισὶν αὐτοῦ· Οὗτός ἐστιν Ἰωάννης ὁ Βαπτιστής· αὐτὸς ἠγέρθη ἀπὸ τῶν νεκρῶν καὶ διὰ τοῦτο αἱ δυνάμεις ἐνεργοῦσιν ἐν αὐτῷ.

\vs{3}Ὁ γὰρ Ἡρῴδης κρατήσας τὸν Ἰωάννην ἔδησεν αὐτὸν καὶ ἐν φυλακῇ ἀπέθετο διὰ Ἡρῳδιάδα τὴν γυναῖκα Φιλίππου τοῦ ἀδελφοῦ αὐτοῦ·
\vs{4}ἔλεγεν γὰρ ὁ Ἰωάννης αὐτῷ· Οὐκ ἔξεστίν σοι ἔχειν αὐτήν.
\vs{5}καὶ θέλων αὐτὸν ἀποκτεῖναι ἐφοβήθη τὸν ὄχλον, ὅτι ὡς προφήτην αὐτὸν εἶχον.

\vs{6}Γενεσίοις δὲ γενομένοις τοῦ Ἡρῴδου ὠρχήσατο ἡ θυγάτηρ τῆς Ἡρῳδιάδος ἐν τῷ μέσῳ καὶ ἤρεσεν τῷ Ἡρῴδῃ,
\vs{7}ὅθεν μεθ᾽ ὅρκου ὡμολόγησεν αὐτῇ δοῦναι ὃ ἐὰν αἰτήσηται.
\vs{8}Ἡ δὲ προβιβασθεῖσα ὑπὸ τῆς μητρὸς αὐτῆς· Δός μοι, φησίν, Ὧδε ἐπὶ πίνακι τὴν κεφαλὴν Ἰωάννου τοῦ Βαπτιστοῦ.
\vs{9}Καὶ λυπηθεὶς ὁ βασιλεὺς διὰ τοὺς ὅρκους καὶ τοὺς συνανακειμένους ἐκέλευσεν δοθῆναι,
\vs{10}καὶ πέμψας ἀπεκεφάλισεν τὸν Ἰωάννην ἐν τῇ φυλακῇ.
\vs{11}Καὶ ἠνέχθη ἡ κεφαλὴ αὐτοῦ ἐπὶ πίνακι καὶ ἐδόθη τῷ κορασίῳ, καὶ ἤνεγκεν τῇ μητρὶ αὐτῆς.
\vs{12}καὶ προσελθόντες οἱ μαθηταὶ αὐτοῦ ἦραν τὸ πτῶμα καὶ ἔθαψαν αὐτόν καὶ ἐλθόντες ἀπήγγειλαν τῷ Ἰησοῦ.

\vs{13}Ἀκούσας δὲ ὁ Ἰησοῦς ἀνεχώρησεν ἐκεῖθεν ἐν πλοίῳ εἰς ἔρημον τόπον κατ᾽ ἰδίαν· καὶ ἀκούσαντες οἱ ὄχλοι ἠκολούθησαν αὐτῷ πεζῇ ἀπὸ τῶν πόλεων.
\vs{14}Καὶ ἐξελθὼν εἶδεν πολὺν ὄχλον καὶ ἐσπλαγχνίσθη ἐπ᾽ αὐτοῖς καὶ ἐθεράπευσεν τοὺς ἀρρώστους αὐτῶν.

\vs{15}Ὀψίας δὲ γενομένης προσῆλθον αὐτῷ οἱ μαθηταὶ λέγοντες· Ἔρημός ἐστιν ὁ τόπος καὶ ἡ ὥρα ἤδη παρῆλθεν· ἀπόλυσον τοὺς ὄχλους, ἵνα ἀπελθόντες εἰς τὰς κώμας ἀγοράσωσιν ἑαυτοῖς βρώματα.
\vs{16}Ὁ δὲ Ἰησοῦς εἶπεν αὐτοῖς· Οὐ χρείαν ἔχουσιν ἀπελθεῖν, δότε αὐτοῖς ὑμεῖς φαγεῖν.
\vs{17}Οἱ δὲ λέγουσιν αὐτῷ· Οὐκ ἔχομεν ὧδε εἰ μὴ πέντε ἄρτους καὶ δύο ἰχθύας.
\vs{18}Ὁ δὲ εἶπεν· Φέρετέ μοι ὧδε αὐτούς.
\vs{19}καὶ κελεύσας τοὺς ὄχλους ἀνακλιθῆναι ἐπὶ τοῦ χόρτου, λαβὼν τοὺς πέντε ἄρτους καὶ τοὺς δύο ἰχθύας, ἀναβλέψας εἰς τὸν οὐρανὸν εὐλόγησεν καὶ κλάσας ἔδωκεν τοῖς μαθηταῖς τοὺς ἄρτους, οἱ δὲ μαθηταὶ τοῖς ὄχλοις.
\vs{20}Καὶ ἔφαγον πάντες καὶ ἐχορτάσθησαν, καὶ ἦραν τὸ περισσεῦον τῶν κλασμάτων δώδεκα κοφίνους πλήρεις.
\vs{21}οἱ δὲ ἐσθίοντες ἦσαν ἄνδρες ὡσεὶ πεντακισχίλιοι χωρὶς γυναικῶν καὶ παιδίων.

\vs{22}Καὶ εὐθέως ἠνάγκασεν τοὺς μαθητὰς ἐμβῆναι εἰς τὸ πλοῖον καὶ προάγειν αὐτὸν εἰς τὸ πέραν, ἕως οὗ ἀπολύσῃ τοὺς ὄχλους.
\vs{23}καὶ ἀπολύσας τοὺς ὄχλους ἀνέβη εἰς τὸ ὄρος κατ᾽ ἰδίαν προσεύξασθαι. ὀψίας δὲ γενομένης μόνος ἦν ἐκεῖ.
\vs{24}τὸ δὲ πλοῖον ἤδη σταδίους πολλοὺς ἀπὸ τῆς γῆς ἀπεῖχεν βασανιζόμενον ὑπὸ τῶν κυμάτων, ἦν γὰρ ἐναντίος ὁ ἄνεμος.
\vs{25}Τετάρτῃ δὲ φυλακῇ τῆς νυκτὸς ἦλθεν πρὸς αὐτοὺς περιπατῶν ἐπὶ τὴν θάλασσαν.
\vs{26}οἱ δὲ μαθηταὶ ἰδόντες αὐτὸν ἐπὶ τῆς θαλάσσης περιπατοῦντα ἐταράχθησαν λέγοντες ὅτι Φάντασμά ἐστιν, καὶ ἀπὸ τοῦ φόβου ἔκραξαν.
\vs{27}Εὐθὺς δὲ ἐλάλησεν ὁ Ἰησοῦς αὐτοῖς λέγων· Θαρσεῖτε, ἐγώ εἰμι· μὴ φοβεῖσθε.
\vs{28}Ἀποκριθεὶς δὲ αὐτῷ ὁ Πέτρος εἶπεν· Κύριε, εἰ σὺ εἶ, κέλευσόν με ἐλθεῖν πρὸς σὲ ἐπὶ τὰ ὕδατα.
\vs{29}Ὁ δὲ εἶπεν· Ἐλθέ. Καὶ καταβὰς ἀπὸ τοῦ πλοίου ὁ Πέτρος περιεπάτησεν ἐπὶ τὰ ὕδατα καὶ ἦλθεν πρὸς τὸν Ἰησοῦν.
\vs{30}βλέπων δὲ τὸν ἄνεμον ἰσχυρὸν ἐφοβήθη, καὶ ἀρξάμενος καταποντίζεσθαι ἔκραξεν λέγων· Κύριε, σῶσόν με.
\vs{31}Εὐθέως δὲ ὁ Ἰησοῦς ἐκτείνας τὴν χεῖρα ἐπελάβετο αὐτοῦ καὶ λέγει αὐτῷ· Ὀλιγόπιστε, εἰς τί ἐδίστασας;
\vs{32}Καὶ ἀναβάντων αὐτῶν εἰς τὸ πλοῖον ἐκόπασεν ὁ ἄνεμος.
\vs{33}οἱ δὲ ἐν τῷ πλοίῳ προσεκύνησαν αὐτῷ λέγοντες· Ἀληθῶς Θεοῦ Υἱὸς εἶ.

\vs{34}Καὶ διαπεράσαντες ἦλθον ἐπὶ τὴν γῆν εἰς Γεννησαρέτ.
\vs{35}καὶ ἐπιγνόντες αὐτὸν οἱ ἄνδρες τοῦ τόπου ἐκείνου ἀπέστειλαν εἰς ὅλην τὴν περίχωρον ἐκείνην καὶ προσήνεγκαν αὐτῷ πάντας τοὺς κακῶς ἔχοντας
\vs{36}καὶ παρεκάλουν αὐτὸν ἵνα μόνον ἅψωνται τοῦ κρασπέδου τοῦ ἱματίου αὐτοῦ· καὶ ὅσοι ἥψαντο διεσώθησαν.

\ch{15}
Τότε προσέρχονται τῷ Ἰησοῦ ἀπὸ Ἱεροσολύμων Φαρισαῖοι καὶ γραμματεῖς λέγοντες·
\vs{2}Διὰ τί οἱ μαθηταί σου παραβαίνουσιν τὴν παράδοσιν τῶν πρεσβυτέρων; οὐ γὰρ νίπτονται τὰς χεῖρας αὐτῶν ὅταν ἄρτον ἐσθίωσιν.
\vs{3}Ὁ δὲ ἀποκριθεὶς εἶπεν αὐτοῖς· Διὰ τί καὶ ὑμεῖς παραβαίνετε τὴν ἐντολὴν τοῦ Θεοῦ διὰ τὴν παράδοσιν ὑμῶν;
\vs{4}ὁ γὰρ Θεὸς εἶπεν· Τίμα τὸν πατέρα καὶ τὴν μητέρα, καί· Ὁ κακολογῶν πατέρα ἢ μητέρα θανάτῳ τελευτάτω.
\vs{5}ὑμεῖς δὲ λέγετε· Ὃς ἂν εἴπῃ τῷ πατρὶ ἢ τῇ μητρί· Δῶρον ὃ ἐὰν ἐξ ἐμοῦ ὠφεληθῇς,
\vs{6}οὐ μὴ τιμήσει τὸν πατέρα αὐτοῦ· καὶ ἠκυρώσατε τὸν λόγον τοῦ Θεοῦ διὰ τὴν παράδοσιν ὑμῶν.
\vs{7}ὑποκριταί, καλῶς ἐπροφήτευσεν περὶ ὑμῶν Ἠσαΐας λέγων·
\begin{poetryblock}

\begin{quote} \vs{8}Ὁ λαὸς οὗτος τοῖς χείλεσίν με τιμᾷ,\end{quote} 

\begin{quote}Ἡ δὲ καρδία αὐτῶν πόρρω ἀπέχει ἀπ᾽ ἐμοῦ·\end{quote}

\begin{quote} \vs{9}Μάτην δὲ σέβονταί με\end{quote} 

\begin{quote}Διδάσκοντες διδασκαλίας ἐντάλματα ἀνθρώπων.\end{quote}
\end{poetryblock}

\vs{10}Καὶ προσκαλεσάμενος τὸν ὄχλον εἶπεν αὐτοῖς· Ἀκούετε καὶ συνίετε·
\vs{11}οὐ τὸ εἰσερχόμενον εἰς τὸ στόμα κοινοῖ τὸν ἄνθρωπον, ἀλλὰ τὸ ἐκπορευόμενον ἐκ τοῦ στόματος τοῦτο κοινοῖ τὸν ἄνθρωπον.

\vs{12}Τότε προσελθόντες οἱ μαθηταὶ λέγουσιν αὐτῷ· Οἶδας ὅτι οἱ Φαρισαῖοι ἀκούσαντες τὸν λόγον ἐσκανδαλίσθησαν;
\vs{13}Ὁ δὲ ἀποκριθεὶς εἶπεν· Πᾶσα φυτεία ἣν οὐκ ἐφύτευσεν ὁ Πατήρ μου ὁ οὐράνιος ἐκριζωθήσεται.
\vs{14}ἄφετε αὐτούς· τυφλοί εἰσιν ὁδηγοί τυφλῶν· τυφλὸς δὲ τυφλὸν ἐὰν ὁδηγῇ, ἀμφότεροι εἰς βόθυνον πεσοῦνται.

\vs{15}Ἀποκριθεὶς δὲ ὁ Πέτρος εἶπεν αὐτῷ· Φράσον ἡμῖν τὴν παραβολήν ταύτην.
\vs{16}Ὁ δὲ εἶπεν· Ἀκμὴν καὶ ὑμεῖς ἀσύνετοί ἐστε;
\vs{17}οὐ νοεῖτε ὅτι πᾶν τὸ εἰσπορευόμενον εἰς τὸ στόμα εἰς τὴν κοιλίαν χωρεῖ καὶ εἰς ἀφεδρῶνα ἐκβάλλεται;
\vs{18}τὰ δὲ ἐκπορευόμενα ἐκ τοῦ στόματος ἐκ τῆς καρδίας ἐξέρχεται, κἀκεῖνα κοινοῖ τὸν ἄνθρωπον.
\vs{19}ἐκ γὰρ τῆς καρδίας ἐξέρχονται διαλογισμοὶ πονηροί, φόνοι, μοιχεῖαι, πορνεῖαι, κλοπαί, ψευδομαρτυρίαι, βλασφημίαι.
\vs{20}ταῦτά ἐστιν τὰ κοινοῦντα τὸν ἄνθρωπον, τὸ δὲ ἀνίπτοις χερσὶν φαγεῖν οὐ κοινοῖ τὸν ἄνθρωπον.

\vs{21}Καὶ ἐξελθὼν ἐκεῖθεν ὁ Ἰησοῦς ἀνεχώρησεν εἰς τὰ μέρη Τύρου καὶ Σιδῶνος.
\vs{22}καὶ ἰδοὺ γυνὴ Χαναναία ἀπὸ τῶν ὁρίων ἐκείνων ἐξελθοῦσα ἔκραζεν λέγουσα· Ἐλέησόν με, Κύριε υἱὸς Δαυίδ· ἡ θυγάτηρ μου κακῶς δαιμονίζεται.
\vs{23}Ὁ δὲ οὐκ ἀπεκρίθη αὐτῇ λόγον. καὶ προσελθόντες οἱ μαθηταὶ αὐτοῦ ἠρώτουν αὐτὸν λέγοντες· Ἀπόλυσον αὐτήν, ὅτι κράζει ὄπισθεν ἡμῶν.
\vs{24}Ὁ δὲ ἀποκριθεὶς εἶπεν· Οὐκ ἀπεστάλην εἰ μὴ εἰς τὰ πρόβατα τὰ ἀπολωλότα οἴκου Ἰσραήλ.
\vs{25}Ἡ δὲ ἐλθοῦσα προσεκύνει αὐτῷ λέγουσα· Κύριε, βοήθει μοι.
\vs{26}Ὁ δὲ ἀποκριθεὶς εἶπεν· Οὐκ ἔστιν καλὸν λαβεῖν τὸν ἄρτον τῶν τέκνων καὶ βαλεῖν τοῖς κυναρίοις.
\vs{27}Ἡ δὲ εἶπεν· Ναί κύριε, καὶ γὰρ τὰ κυνάρια ἐσθίει ἀπὸ τῶν ψιχίων τῶν πιπτόντων ἀπὸ τῆς τραπέζης τῶν κυρίων αὐτῶν.
\vs{28}Τότε ἀποκριθεὶς ὁ Ἰησοῦς εἶπεν αὐτῇ· Ὦ γύναι, μεγάλη σου ἡ πίστις· γενηθήτω σοι ὡς θέλεις. καὶ ἰάθη ἡ θυγάτηρ αὐτῆς ἀπὸ τῆς ὥρας ἐκείνης.

\vs{29}Καὶ μεταβὰς ἐκεῖθεν ὁ Ἰησοῦς ἦλθεν παρὰ τὴν θάλασσαν τῆς Γαλιλαίας, καὶ ἀναβὰς εἰς τὸ ὄρος ἐκάθητο ἐκεῖ.
\vs{30}καὶ προσῆλθον αὐτῷ ὄχλοι πολλοὶ ἔχοντες μεθ᾽ ἑαυτῶν χωλούς, τυφλούς, κυλλούς, κωφούς, καὶ ἑτέρους πολλούς καὶ ἔρριψαν αὐτοὺς παρὰ τοὺς πόδας αὐτοῦ, καὶ ἐθεράπευσεν αὐτούς·
\vs{31}ὥστε τὸν ὄχλον θαυμάσαι βλέποντας κωφοὺς λαλοῦντας, κυλλοὺς ὑγιεῖς καὶ χωλοὺς περιπατοῦντας καὶ τυφλοὺς βλέποντας· καὶ ἐδόξασαν τὸν Θεὸν Ἰσραήλ.

\vs{32}Ὁ δὲ Ἰησοῦς προσκαλεσάμενος τοὺς μαθητὰς αὐτοῦ εἶπεν· Σπλαγχνίζομαι ἐπὶ τὸν ὄχλον, ὅτι ἤδη ἡμέραι τρεῖς προσμένουσίν μοι καὶ οὐκ ἔχουσιν τί φάγωσιν· καὶ ἀπολῦσαι αὐτοὺς νήστεις οὐ θέλω, μήποτε ἐκλυθῶσιν ἐν τῇ ὁδῷ.
\vs{33}Καὶ λέγουσιν αὐτῷ οἱ μαθηταί· Πόθεν ἡμῖν ἐν ἐρημίᾳ ἄρτοι τοσοῦτοι ὥστε χορτάσαι ὄχλον τοσοῦτον;
\vs{34}Καὶ λέγει αὐτοῖς ὁ Ἰησοῦς· Πόσους ἄρτους ἔχετε; Οἱ δὲ εἶπαν· Ἑπτά καὶ ὀλίγα ἰχθύδια.
\vs{35}Καὶ παραγγείλας τῷ ὄχλῳ ἀναπεσεῖν ἐπὶ τὴν γῆν
\vs{36}ἔλαβεν τοὺς ἑπτὰ ἄρτους καὶ τοὺς ἰχθύας καὶ εὐχαριστήσας ἔκλασεν καὶ ἐδίδου τοῖς μαθηταῖς, οἱ δὲ μαθηταὶ τοῖς ὄχλοις.
\vs{37}Καὶ ἔφαγον πάντες καὶ ἐχορτάσθησαν. καὶ τὸ περισσεῦον τῶν κλασμάτων ἦραν ἑπτὰ σπυρίδας πλήρεις.
\vs{38}οἱ δὲ ἐσθίοντες ἦσαν τετρακισχίλιοι ἄνδρες χωρὶς γυναικῶν καὶ παιδίων.

\vs{39}Καὶ ἀπολύσας τοὺς ὄχλους ἐνέβη εἰς τὸ πλοῖον καὶ ἦλθεν εἰς τὰ ὅρια Μαγαδάν.

\ch{16}
Καὶ προσελθόντες οἱ Φαρισαῖοι καὶ Σαδδουκαῖοι πειράζοντες ἐπηρώτησαν αὐτὸν σημεῖον ἐκ τοῦ οὐρανοῦ ἐπιδεῖξαι αὐτοῖς.
\vs{2}Ὁ δὲ ἀποκριθεὶς εἶπεν αὐτοῖς· Ὀψίας γενομένης λέγετε· Εὐδία, πυρράζει γὰρ ὁ οὐρανός·
\vs{3}καὶ πρωΐ· Σήμερον χειμών, πυρράζει γὰρ στυγνάζων ὁ οὐρανός. τὸ μὲν πρόσωπον τοῦ οὐρανοῦ γινώσκετε διακρίνειν, τὰ δὲ σημεῖα τῶν καιρῶν οὐ δύνασθε;
\vs{4}γενεὰ πονηρὰ καὶ μοιχαλὶς σημεῖον ἐπιζητεῖ, καὶ σημεῖον οὐ δοθήσεται αὐτῇ εἰ μὴ τὸ σημεῖον Ἰωνᾶ. καὶ καταλιπὼν αὐτοὺς ἀπῆλθεν.

\vs{5}Καὶ ἐλθόντες οἱ μαθηταὶ εἰς τὸ πέραν ἐπελάθοντο ἄρτους λαβεῖν.
\vs{6}ὁ δὲ Ἰησοῦς εἶπεν αὐτοῖς· Ὁρᾶτε καὶ προσέχετε ἀπὸ τῆς ζύμης τῶν Φαρισαίων καὶ Σαδδουκαίων.
\vs{7}Οἱ δὲ διελογίζοντο ἐν ἑαυτοῖς λέγοντες Ὅτι Ἄρτους οὐκ ἐλάβομεν.
\vs{8}Γνοὺς δὲ ὁ Ἰησοῦς εἶπεν· Τί διαλογίζεσθε ἐν ἑαυτοῖς, ὀλιγόπιστοι, ὅτι ἄρτους οὐκ ἔχετε;
\vs{9}οὔπω νοεῖτε, οὐδὲ μνημονεύετε τοὺς πέντε ἄρτους τῶν πεντακισχιλίων καὶ πόσους κοφίνους ἐλάβετε;
\vs{10}οὐδὲ τοὺς ἑπτὰ ἄρτους τῶν τετρακισχιλίων καὶ πόσας σπυρίδας ἐλάβετε;
\vs{11}πῶς οὐ νοεῖτε ὅτι οὐ περὶ ἄρτων εἶπον ὑμῖν; προσέχετε δὲ ἀπὸ τῆς ζύμης τῶν Φαρισαίων καὶ Σαδδουκαίων.
\vs{12}Τότε συνῆκαν ὅτι οὐκ εἶπεν προσέχειν ἀπὸ τῆς ζύμης τῶν ἄρτων ἀλλὰ ἀπὸ τῆς διδαχῆς τῶν Φαρισαίων καὶ Σαδδουκαίων.

\vs{13}Ἐλθὼν δὲ ὁ Ἰησοῦς εἰς τὰ μέρη Καισαρείας τῆς Φιλίππου ἠρώτα τοὺς μαθητὰς αὐτοῦ λέγων· Τίνα λέγουσιν οἱ ἄνθρωποι εἶναι τὸν Υἱὸν τοῦ ἀνθρώπου;
\vs{14}Οἱ δὲ εἶπαν· Οἱ μὲν Ἰωάννην τὸν Βαπτιστήν, ἄλλοι δὲ Ἠλίαν, ἕτεροι δὲ Ἰερεμίαν ἢ ἕνα τῶν προφητῶν.
\vs{15}Λέγει αὐτοῖς· Ὑμεῖς δὲ τίνα με λέγετε εἶναι;
\vs{16}Ἀποκριθεὶς δὲ Σίμων Πέτρος εἶπεν· Σὺ εἶ ὁ Χριστὸς ὁ Υἱὸς τοῦ Θεοῦ τοῦ ζῶντος.
\vs{17}Ἀποκριθεὶς δὲ ὁ Ἰησοῦς εἶπεν αὐτῷ· Μακάριος εἶ, Σίμων Βαριωνᾶ, ὅτι σὰρξ καὶ αἷμα οὐκ ἀπεκάλυψέν σοι ἀλλ᾽ ὁ Πατήρ μου ὁ ἐν τοῖς οὐρανοῖς.
\vs{18}κἀγὼ δέ σοι λέγω ὅτι σὺ εἶ Πέτρος, καὶ ἐπὶ ταύτῃ τῇ πέτρᾳ οἰκοδομήσω μου τὴν ἐκκλησίαν καὶ πύλαι ᾅδου οὐ κατισχύσουσιν αὐτῆς.
\vs{19}δώσω σοι τὰς κλεῖδας τῆς βασιλείας τῶν οὐρανῶν, καὶ ὃ ἐὰν δήσῃς ἐπὶ τῆς γῆς ἔσται δεδεμένον ἐν τοῖς οὐρανοῖς, καὶ ὃ ἐὰν λύσῃς ἐπὶ τῆς γῆς ἔσται λελυμένον ἐν τοῖς οὐρανοῖς.
\vs{20}Τότε διεστείλατο τοῖς μαθηταῖς ἵνα μηδενὶ εἴπωσιν ὅτι αὐτός ἐστιν ὁ Χριστός.

\vs{21}Ἀπὸ τότε ἤρξατο ὁ Ἰησοῦς δεικνύειν τοῖς μαθηταῖς αὐτοῦ ὅτι δεῖ αὐτὸν εἰς Ἱεροσόλυμα ἀπελθεῖν καὶ πολλὰ παθεῖν ἀπὸ τῶν πρεσβυτέρων καὶ ἀρχιερέων καὶ γραμματέων καὶ ἀποκτανθῆναι καὶ τῇ τρίτῃ ἡμέρᾳ ἐγερθῆναι.
\vs{22}Καὶ προσλαβόμενος αὐτὸν ὁ Πέτρος ἤρξατο ἐπιτιμᾶν αὐτῷ λέγων· Ἵλεώς σοι, Κύριε· οὐ μὴ ἔσται σοι τοῦτο.
\vs{23}Ὁ δὲ στραφεὶς εἶπεν τῷ Πέτρῳ· Ὕπαγε ὀπίσω μου, Σατανᾶ· σκάνδαλον εἶ ἐμοῦ, ὅτι οὐ φρονεῖς τὰ τοῦ Θεοῦ ἀλλὰ τὰ τῶν ἀνθρώπων.

\vs{24}Τότε ὁ Ἰησοῦς εἶπεν τοῖς μαθηταῖς αὐτοῦ· Εἴ τις θέλει ὀπίσω μου ἐλθεῖν, ἀπαρνησάσθω ἑαυτὸν καὶ ἀράτω τὸν σταυρὸν αὐτοῦ καὶ ἀκολουθείτω μοι.
\vs{25}ὃς γὰρ ἐὰν θέλῃ τὴν ψυχὴν αὐτοῦ σῶσαι ἀπολέσει αὐτήν· ὃς δ᾽ ἂν ἀπολέσῃ τὴν ψυχὴν αὐτοῦ ἕνεκεν ἐμοῦ εὑρήσει αὐτήν.
\vs{26}τί γὰρ ὠφεληθήσεται ἄνθρωπος ἐὰν τὸν κόσμον ὅλον κερδήσῃ τὴν δὲ ψυχὴν αὐτοῦ ζημιωθῇ; ἢ τί δώσει ἄνθρωπος ἀντάλλαγμα τῆς ψυχῆς αὐτοῦ;
\vs{27}μέλλει γὰρ ὁ Υἱὸς τοῦ ἀνθρώπου ἔρχεσθαι ἐν τῇ δόξῃ τοῦ Πατρὸς αὐτοῦ μετὰ τῶν ἀγγέλων αὐτοῦ, καὶ τότε ἀποδώσει ἑκάστῳ κατὰ τὴν πρᾶξιν αὐτοῦ.
\vs{28}Ἀμὴν λέγω ὑμῖν ὅτι εἰσίν τινες τῶν ὧδε ἑστώτων οἵτινες οὐ μὴ γεύσωνται θανάτου ἕως ἂν ἴδωσιν τὸν Υἱὸν τοῦ ἀνθρώπου ἐρχόμενον ἐν τῇ βασιλείᾳ αὐτοῦ.

\ch{17}
Καὶ μεθ᾽ ἡμέρας ἓξ παραλαμβάνει ὁ Ἰησοῦς τὸν Πέτρον καὶ Ἰάκωβον καὶ Ἰωάννην τὸν ἀδελφὸν αὐτοῦ καὶ ἀναφέρει αὐτοὺς εἰς ὄρος ὑψηλὸν κατ᾽ ἰδίαν.
\vs{2}καὶ μετεμορφώθη ἔμπροσθεν αὐτῶν, καὶ ἔλαμψεν τὸ πρόσωπον αὐτοῦ ὡς ὁ ἥλιος, τὰ δὲ ἱμάτια αὐτοῦ ἐγένετο λευκὰ ὡς τὸ φῶς.
\vs{3}Καὶ ἰδοὺ ὤφθη αὐτοῖς Μωϋσῆς καὶ Ἠλίας συλλαλοῦντες μετ᾽ αὐτοῦ.
\vs{4}ἀποκριθεὶς δὲ ὁ Πέτρος εἶπεν τῷ Ἰησοῦ· Κύριε, καλόν ἐστιν ἡμᾶς ὧδε εἶναι· εἰ θέλεις, ποιήσω ὧδε τρεῖς σκηνάς, σοὶ μίαν καὶ Μωϋσεῖ μίαν καὶ Ἠλίᾳ μίαν.
\vs{5}Ἔτι αὐτοῦ λαλοῦντος ἰδοὺ νεφέλη φωτεινὴ ἐπεσκίασεν αὐτούς, καὶ ἰδοὺ φωνὴ ἐκ τῆς νεφέλης λέγουσα· Οὗτός ἐστιν ὁ Υἱός μου ὁ ἀγαπητός, ἐν ᾧ εὐδόκησα· ἀκούετε αὐτοῦ.
\vs{6}καὶ ἀκούσαντες οἱ μαθηταὶ ἔπεσαν ἐπὶ πρόσωπον αὐτῶν καὶ ἐφοβήθησαν σφόδρα.
\vs{7}Καὶ προσῆλθεν ὁ Ἰησοῦς καὶ ἁψάμενος αὐτῶν εἶπεν· Ἐγέρθητε καὶ μὴ φοβεῖσθε.
\vs{8}ἐπάραντες δὲ τοὺς ὀφθαλμοὺς αὐτῶν οὐδένα εἶδον εἰ μὴ αὐτὸν Ἰησοῦν μόνον.

\vs{9}Καὶ καταβαινόντων αὐτῶν ἐκ τοῦ ὄρους ἐνετείλατο αὐτοῖς ὁ Ἰησοῦς λέγων· Μηδενὶ εἴπητε τὸ ὅραμα ἕως οὗ ὁ Υἱὸς τοῦ ἀνθρώπου ἐκ νεκρῶν ἐγερθῇ.
\vs{10}Καὶ ἐπηρώτησαν αὐτὸν οἱ μαθηταὶ λέγοντες· Τί οὖν οἱ γραμματεῖς λέγουσιν ὅτι Ἠλίαν δεῖ ἐλθεῖν πρῶτον;
\vs{11}Ὁ δὲ ἀποκριθεὶς εἶπεν· Ἠλίας μὲν ἔρχεται καὶ ἀποκαταστήσει πάντα·
\vs{12}λέγω δὲ ὑμῖν ὅτι Ἠλίας ἤδη ἦλθεν, καὶ οὐκ ἐπέγνωσαν αὐτὸν ἀλλὰ ἐποίησαν ἐν αὐτῷ ὅσα ἠθέλησαν· οὕτως καὶ ὁ Υἱὸς τοῦ ἀνθρώπου μέλλει πάσχειν ὑπ᾽ αὐτῶν.
\vs{13}τότε συνῆκαν οἱ μαθηταὶ ὅτι περὶ Ἰωάννου τοῦ Βαπτιστοῦ εἶπεν αὐτοῖς.

\vs{14}Καὶ ἐλθόντων πρὸς τὸν ὄχλον προσῆλθεν αὐτῷ ἄνθρωπος γονυπετῶν αὐτὸν
\vs{15}καὶ λέγων· Κύριε, ἐλέησόν μου τὸν υἱόν, ὅτι σεληνιάζεται καὶ κακῶς πάσχει· πολλάκις γὰρ πίπτει εἰς τὸ πῦρ καὶ πολλάκις εἰς τὸ ὕδωρ.
\vs{16}καὶ προσήνεγκα αὐτὸν τοῖς μαθηταῖς σου, καὶ οὐκ ἠδυνήθησαν αὐτὸν θεραπεῦσαι.
\vs{17}Ἀποκριθεὶς δὲ ὁ Ἰησοῦς εἶπεν· Ὦ γενεὰ ἄπιστος καὶ διεστραμμένη, ἕως πότε μεθ᾽ ὑμῶν ἔσομαι; ἕως πότε ἀνέξομαι ὑμῶν; φέρετέ μοι αὐτὸν ὧδε.
\vs{18}καὶ ἐπετίμησεν αὐτῷ ὁ Ἰησοῦς καὶ ἐξῆλθεν ἀπ᾽ αὐτοῦ τὸ δαιμόνιον καὶ ἐθεραπεύθη ὁ παῖς ἀπὸ τῆς ὥρας ἐκείνης.

\vs{19}Τότε προσελθόντες οἱ μαθηταὶ τῷ Ἰησοῦ κατ᾽ ἰδίαν εἶπον· Διὰ τί ἡμεῖς οὐκ ἠδυνήθημεν ἐκβαλεῖν αὐτό;
\vs{20}Ὁ δὲ λέγει αὐτοῖς· Διὰ τὴν ὀλιγοπιστίαν ὑμῶν· ἀμὴν γὰρ λέγω ὑμῖν, ἐὰν ἔχητε πίστιν ὡς κόκκον σινάπεως, ἐρεῖτε τῷ ὄρει τούτῳ· Μετάβα ἔνθεν ἐκεῖ, καὶ μεταβήσεται· καὶ οὐδὲν ἀδυνατήσει ὑμῖν.

\vs{22}Συστρεφομένων δὲ αὐτῶν ἐν τῇ Γαλιλαίᾳ εἶπεν αὐτοῖς ὁ Ἰησοῦς· Μέλλει ὁ Υἱὸς τοῦ ἀνθρώπου παραδίδοσθαι εἰς χεῖρας ἀνθρώπων,
\vs{23}καὶ ἀποκτενοῦσιν αὐτόν, καὶ τῇ τρίτῃ ἡμέρᾳ ἐγερθήσεται. καὶ ἐλυπήθησαν σφόδρα.

\vs{24}Ἐλθόντων δὲ αὐτῶν εἰς Καφαρναοὺμ προσῆλθον οἱ τὰ δίδραχμα λαμβάνοντες τῷ Πέτρῳ καὶ εἶπαν· Ὁ διδάσκαλος ὑμῶν οὐ τελεῖ τὰ δίδραχμα;
\vs{25}Λέγει· Ναί. Καὶ ἐλθόντα εἰς τὴν οἰκίαν προέφθασεν αὐτὸν ὁ Ἰησοῦς λέγων· Τί σοι δοκεῖ, Σίμων; οἱ βασιλεῖς τῆς γῆς ἀπὸ τίνων λαμβάνουσιν τέλη ἢ κῆνσον; ἀπὸ τῶν υἱῶν αὐτῶν ἢ ἀπὸ τῶν ἀλλοτρίων;
\vs{26}Εἰπόντος δέ· Ἀπὸ τῶν ἀλλοτρίων, ἔφη αὐτῷ ὁ Ἰησοῦς· Ἄρα Γε ἐλεύθεροί εἰσιν οἱ υἱοί.
\vs{27}ἵνα δὲ μὴ σκανδαλίσωμεν αὐτούς, πορευθεὶς εἰς θάλασσαν βάλε ἄγκιστρον καὶ τὸν ἀναβάντα πρῶτον ἰχθὺν ἆρον, καὶ ἀνοίξας τὸ στόμα αὐτοῦ εὑρήσεις στατῆρα· ἐκεῖνον λαβὼν δὸς αὐτοῖς ἀντὶ ἐμοῦ καὶ σοῦ.

\ch{18}
Ἐν ἐκείνῃ τῇ ὥρᾳ προσῆλθον οἱ μαθηταὶ τῷ Ἰησοῦ λέγοντες· Τίς ἄρα μείζων ἐστὶν ἐν τῇ βασιλείᾳ τῶν οὐρανῶν;
\vs{2}Καὶ προσκαλεσάμενος παιδίον ἔστησεν αὐτὸ ἐν μέσῳ αὐτῶν
\vs{3}καὶ εἶπεν· Ἀμὴν λέγω ὑμῖν, ἐὰν μὴ στραφῆτε καὶ γένησθε ὡς τὰ παιδία, οὐ μὴ εἰσέλθητε εἰς τὴν βασιλείαν τῶν οὐρανῶν.
\vs{4}ὅστις οὖν ταπεινώσει ἑαυτὸν ὡς τὸ παιδίον τοῦτο, οὗτός ἐστιν ὁ μείζων ἐν τῇ βασιλείᾳ τῶν οὐρανῶν.
\vs{5}καὶ ὃς ἐὰν δέξηται ἓν παιδίον τοιοῦτο ἐπὶ τῷ ὀνόματί μου, ἐμὲ δέχεται.

\vs{6}Ὃς δ᾽ ἂν σκανδαλίσῃ ἕνα τῶν μικρῶν τούτων τῶν πιστευόντων εἰς ἐμέ, συμφέρει αὐτῷ ἵνα κρεμασθῇ μύλος ὀνικὸς περὶ τὸν τράχηλον αὐτοῦ καὶ καταποντισθῇ ἐν τῷ πελάγει τῆς θαλάσσης.
\vs{7}Οὐαὶ τῷ κόσμῳ ἀπὸ τῶν σκανδάλων· ἀνάγκη γὰρ ἐλθεῖν τὰ σκάνδαλα, πλὴν οὐαὶ τῷ ἀνθρώπῳ δι᾽ οὗ τὸ σκάνδαλον ἔρχεται.
\vs{8}Εἰ δὲ ἡ χείρ σου ἢ ὁ πούς σου σκανδαλίζει σε, ἔκκοψον αὐτὸν καὶ βάλε ἀπὸ σοῦ· καλόν σοί ἐστιν εἰσελθεῖν εἰς τὴν ζωὴν κυλλὸν ἢ χωλόν ἢ δύο χεῖρας ἢ δύο πόδας ἔχοντα βληθῆναι εἰς τὸ πῦρ τὸ αἰώνιον.
\vs{9}καὶ εἰ ὁ ὀφθαλμός σου σκανδαλίζει σε, ἔξελε αὐτὸν καὶ βάλε ἀπὸ σοῦ· καλόν σοί ἐστιν μονόφθαλμον εἰς τὴν ζωὴν εἰσελθεῖν ἢ δύο ὀφθαλμοὺς ἔχοντα βληθῆναι εἰς τὴν γέενναν τοῦ πυρός.

\vs{10}Ὁρᾶτε μὴ καταφρονήσητε ἑνὸς τῶν μικρῶν τούτων· λέγω γὰρ ὑμῖν ὅτι οἱ ἄγγελοι αὐτῶν ἐν οὐρανοῖς διὰ παντὸς βλέπουσι τὸ πρόσωπον τοῦ Πατρός μου τοῦ ἐν οὐρανοῖς.

\vs{12}Τί ὑμῖν δοκεῖ; ἐὰν γένηταί τινι ἀνθρώπῳ ἑκατὸν πρόβατα καὶ πλανηθῇ ἓν ἐξ αὐτῶν, οὐχὶ ἀφήσει τὰ ἐνενήκοντα ἐννέα ἐπὶ τὰ ὄρη καὶ πορευθεὶς ζητεῖ τὸ πλανώμενον;
\vs{13}καὶ ἐὰν γένηται εὑρεῖν αὐτό, ἀμὴν λέγω ὑμῖν ὅτι χαίρει ἐπ᾽ αὐτῷ μᾶλλον ἢ ἐπὶ τοῖς ἐνενήκοντα ἐννέα τοῖς μὴ πεπλανημένοις.
\vs{14}οὕτως οὐκ ἔστιν θέλημα ἔμπροσθεν τοῦ Πατρὸς ὑμῶν τοῦ ἐν οὐρανοῖς ἵνα ἀπόληται ἓν τῶν μικρῶν τούτων.

\vs{15}Ἐὰν δὲ ἁμαρτήσῃ εἰς σὲ ὁ ἀδελφός σου, ὕπαγε ἔλεγξον αὐτὸν μεταξὺ σοῦ καὶ αὐτοῦ μόνου. ἐάν σου ἀκούσῃ, ἐκέρδησας τὸν ἀδελφόν σου·
\vs{16}ἐὰν δὲ μὴ ἀκούσῃ, παράλαβε μετὰ σοῦ ἔτι ἕνα ἢ δύο, ἵνα Ἐπὶ στόματος δύο μαρτύρων ἢ τριῶν σταθῇ πᾶν ῥῆμα·
\vs{17}ἐὰν δὲ παρακούσῃ αὐτῶν, εἰπὲ τῇ ἐκκλησίᾳ· ἐὰν δὲ καὶ τῆς ἐκκλησίας παρακούσῃ, ἔστω σοι ὥσπερ ὁ ἐθνικὸς καὶ ὁ τελώνης.
\vs{18}Ἀμὴν λέγω ὑμῖν· ὅσα ἐὰν δήσητε ἐπὶ τῆς γῆς ἔσται δεδεμένα ἐν οὐρανῷ, καὶ ὅσα ἐὰν λύσητε ἐπὶ τῆς γῆς ἔσται λελυμένα ἐν οὐρανῷ.

\vs{19}Πάλιν ἀμὴν λέγω ὑμῖν ὅτι ἐὰν δύο συμφωνήσωσιν ἐξ ὑμῶν ἐπὶ τῆς γῆς περὶ παντὸς πράγματος οὗ ἐὰν αἰτήσωνται, γενήσεται αὐτοῖς παρὰ τοῦ Πατρός μου τοῦ ἐν οὐρανοῖς.
\vs{20}οὗ γάρ εἰσιν δύο ἢ τρεῖς συνηγμένοι εἰς τὸ ἐμὸν ὄνομα, ἐκεῖ εἰμι ἐν μέσῳ αὐτῶν.

\vs{21}Τότε προσελθὼν ὁ Πέτρος εἶπεν αὐτῷ· Κύριε, ποσάκις ἁμαρτήσει εἰς ἐμὲ ὁ ἀδελφός μου καὶ ἀφήσω αὐτῷ; ἕως ἑπτάκις;
\vs{22}Λέγει αὐτῷ ὁ Ἰησοῦς· Οὐ λέγω σοι ἕως ἑπτάκις ἀλλὰ ἕως ἑβδομηκοντάκις ἑπτά.

\vs{23}Διὰ τοῦτο ὡμοιώθη ἡ βασιλεία τῶν οὐρανῶν ἀνθρώπῳ βασιλεῖ, ὃς ἠθέλησεν συνᾶραι λόγον μετὰ τῶν δούλων αὐτοῦ.
\vs{24}ἀρξαμένου δὲ αὐτοῦ συναίρειν προσηνέχθη αὐτῷ εἷς ὀφειλέτης μυρίων ταλάντων.
\vs{25}μὴ ἔχοντος δὲ αὐτοῦ ἀποδοῦναι ἐκέλευσεν αὐτὸν ὁ κύριος πραθῆναι καὶ τὴν γυναῖκα καὶ τὰ τέκνα καὶ πάντα ὅσα ἔχει, καὶ ἀποδοθῆναι.
\vs{26}Πεσὼν οὖν ὁ δοῦλος προσεκύνει αὐτῷ λέγων· Μακροθύμησον ἐπ᾽ ἐμοί, καὶ πάντα ἀποδώσω σοι.
\vs{27}Σπλαγχνισθεὶς δὲ ὁ κύριος τοῦ δούλου ἐκείνου ἀπέλυσεν αὐτόν καὶ τὸ δάνειον ἀφῆκεν αὐτῷ.
\vs{28}Ἐξελθὼν δὲ ὁ δοῦλος ἐκεῖνος εὗρεν ἕνα τῶν συνδούλων αὐτοῦ, ὃς ὤφειλεν αὐτῷ ἑκατὸν δηνάρια, καὶ κρατήσας αὐτὸν ἔπνιγεν λέγων· Ἀπόδος εἴ τι ὀφείλεις.
\vs{29}Πεσὼν οὖν ὁ σύνδουλος αὐτοῦ παρεκάλει αὐτὸν λέγων· Μακροθύμησον ἐπ᾽ ἐμοί, καὶ ἀποδώσω σοι.
\vs{30}Ὁ δὲ οὐκ ἤθελεν ἀλλὰ ἀπελθὼν ἔβαλεν αὐτὸν εἰς φυλακὴν ἕως ἀποδῷ τὸ ὀφειλόμενον.
\vs{31}Ἰδόντες οὖν οἱ σύνδουλοι αὐτοῦ τὰ γενόμενα ἐλυπήθησαν σφόδρα καὶ ἐλθόντες διεσάφησαν τῷ κυρίῳ ἑαυτῶν πάντα τὰ γενόμενα.
\vs{32}Τότε προσκαλεσάμενος αὐτὸν ὁ κύριος αὐτοῦ λέγει αὐτῷ· Δοῦλε πονηρέ, πᾶσαν τὴν ὀφειλὴν ἐκείνην ἀφῆκά σοι, ἐπεὶ παρεκάλεσάς με·
\vs{33}οὐκ ἔδει καὶ σὲ ἐλεῆσαι τὸν σύνδουλόν σου, ὡς κἀγὼ σὲ ἠλέησα;
\vs{34}καὶ ὀργισθεὶς ὁ κύριος αὐτοῦ παρέδωκεν αὐτὸν τοῖς βασανισταῖς ἕως οὗ ἀποδῷ πᾶν τὸ ὀφειλόμενον.
\vs{35}Οὕτως καὶ ὁ Πατήρ μου ὁ οὐράνιος ποιήσει ὑμῖν, ἐὰν μὴ ἀφῆτε ἕκαστος τῷ ἀδελφῷ αὐτοῦ ἀπὸ τῶν καρδιῶν ὑμῶν.

\ch{19}
Καὶ ἐγένετο ὅτε ἐτέλεσεν ὁ Ἰησοῦς τοὺς λόγους τούτους, μετῆρεν ἀπὸ τῆς Γαλιλαίας καὶ ἦλθεν εἰς τὰ ὅρια τῆς Ἰουδαίας πέραν τοῦ Ἰορδάνου.
\vs{2}καὶ ἠκολούθησαν αὐτῷ ὄχλοι πολλοί, καὶ ἐθεράπευσεν αὐτοὺς ἐκεῖ.

\vs{3}Καὶ προσῆλθον αὐτῷ Φαρισαῖοι πειράζοντες αὐτὸν καὶ λέγοντες· Εἰ ἔξεστιν ἀνθρώπῳ ἀπολῦσαι τὴν γυναῖκα αὐτοῦ κατὰ πᾶσαν αἰτίαν;
\vs{4}Ὁ δὲ ἀποκριθεὶς εἶπεν· Οὐκ ἀνέγνωτε ὅτι ὁ κτίσας ἀπ᾽ ἀρχῆς Ἄρσεν καὶ θῆλυ ἐποίησεν αὐτοὺς;
\vs{5}καὶ εἶπεν· Ἕνεκα τούτου καταλείψει ἄνθρωπος τὸν πατέρα καὶ τὴν μητέρα καὶ κολληθήσεται τῇ γυναικὶ αὐτοῦ, καὶ ἔσονται οἱ δύο εἰς σάρκα μίαν.
\vs{6}ὥστε οὐκέτι εἰσὶν δύο ἀλλὰ σὰρξ μία. ὃ οὖν ὁ Θεὸς συνέζευξεν ἄνθρωπος μὴ χωριζέτω.
\vs{7}Λέγουσιν αὐτῷ· Τί οὖν Μωϋσῆς ἐνετείλατο δοῦναι βιβλίον ἀποστασίου καὶ ἀπολῦσαι αὐτήν;
\vs{8}Λέγει αὐτοῖς Ὅτι Μωϋσῆς πρὸς τὴν σκληροκαρδίαν ὑμῶν ἐπέτρεψεν ὑμῖν ἀπολῦσαι τὰς γυναῖκας ὑμῶν, ἀπ᾽ ἀρχῆς δὲ οὐ γέγονεν οὕτως.
\vs{9}λέγω δὲ ὑμῖν ὅτι ὃς ἂν ἀπολύσῃ τὴν γυναῖκα αὐτοῦ μὴ ἐπὶ πορνείᾳ καὶ γαμήσῃ ἄλλην μοιχᾶται.
\vs{10}Λέγουσιν αὐτῷ οἱ μαθηταί αὐτοῦ· Εἰ οὕτως ἐστὶν ἡ αἰτία τοῦ ἀνθρώπου μετὰ τῆς γυναικός, οὐ συμφέρει γαμῆσαι.
\vs{11}Ὁ δὲ εἶπεν αὐτοῖς· Οὐ πάντες χωροῦσιν τὸν λόγον τοῦτον ἀλλ᾽ οἷς δέδοται.
\vs{12}εἰσὶν γὰρ εὐνοῦχοι οἵτινες ἐκ κοιλίας μητρὸς ἐγεννήθησαν οὕτως, καὶ εἰσὶν εὐνοῦχοι οἵτινες εὐνουχίσθησαν ὑπὸ τῶν ἀνθρώπων, καὶ εἰσὶν εὐνοῦχοι οἵτινες εὐνούχισαν ἑαυτοὺς διὰ τὴν βασιλείαν τῶν οὐρανῶν. ὁ δυνάμενος χωρεῖν χωρείτω.

\vs{13}Τότε προσηνέχθησαν αὐτῷ παιδία ἵνα τὰς χεῖρας ἐπιθῇ αὐτοῖς καὶ προσεύξηται· οἱ δὲ μαθηταὶ ἐπετίμησαν αὐτοῖς.
\vs{14}ὁ δὲ Ἰησοῦς εἶπεν· Ἄφετε τὰ παιδία καὶ μὴ κωλύετε αὐτὰ ἐλθεῖν πρός με, τῶν γὰρ τοιούτων ἐστὶν ἡ βασιλεία τῶν οὐρανῶν.
\vs{15}καὶ ἐπιθεὶς τὰς χεῖρας αὐτοῖς ἐπορεύθη ἐκεῖθεν.

\vs{16}Καὶ ἰδοὺ εἷς προσελθὼν αὐτῷ εἶπεν· Διδάσκαλε, τί ἀγαθὸν ποιήσω ἵνα σχῶ ζωὴν αἰώνιον;
\vs{17}Ὁ δὲ εἶπεν αὐτῷ· Τί με ἐρωτᾷς περὶ τοῦ ἀγαθοῦ; εἷς ἐστιν ὁ ἀγαθός· εἰ δὲ θέλεις εἰς τὴν ζωὴν εἰσελθεῖν, τήρησον τὰς ἐντολάς.
\vs{18}Λέγει αὐτῷ· Ποίας; Ὁ δὲ Ἰησοῦς εἶπεν· Τὸ Οὐ φονεύσεις, Οὐ μοιχεύσεις, Οὐ κλέψεις, Οὐ ψευδομαρτυρήσεις,
\vs{19}Τίμα τὸν πατέρα καὶ τὴν μητέρα, καὶ Ἀγαπήσεις τὸν πλησίον σου ὡς σεαυτόν.
\vs{20}Λέγει αὐτῷ ὁ νεανίσκος· Πάντα ταῦτα ἐφύλαξα· τί ἔτι ὑστερῶ;
\vs{21}Ἔφη αὐτῷ ὁ Ἰησοῦς· Εἰ θέλεις τέλειος εἶναι, ὕπαγε πώλησόν σου τὰ ὑπάρχοντα καὶ δὸς τοῖς πτωχοῖς, καὶ ἕξεις θησαυρὸν ἐν οὐρανοῖς, καὶ δεῦρο ἀκολούθει μοι.
\vs{22}Ἀκούσας δὲ ὁ νεανίσκος τὸν λόγον ἀπῆλθεν λυπούμενος· ἦν γὰρ ἔχων κτήματα πολλά.

\vs{23}Ὁ δὲ Ἰησοῦς εἶπεν τοῖς μαθηταῖς αὐτοῦ· Ἀμὴν λέγω ὑμῖν ὅτι πλούσιος δυσκόλως εἰσελεύσεται εἰς τὴν βασιλείαν τῶν οὐρανῶν.
\vs{24}πάλιν δὲ λέγω ὑμῖν, εὐκοπώτερόν ἐστιν κάμηλον διὰ τρυπήματος ῥαφίδος διελθεῖν ἢ πλούσιον εἰσελθεῖν εἰς τὴν βασιλείαν τοῦ Θεοῦ.
\vs{25}Ἀκούσαντες δὲ οἱ μαθηταὶ ἐξεπλήσσοντο σφόδρα λέγοντες· Τίς ἄρα δύναται σωθῆναι;
\vs{26}Ἐμβλέψας δὲ ὁ Ἰησοῦς εἶπεν αὐτοῖς· Παρὰ ἀνθρώποις τοῦτο ἀδύνατόν ἐστιν, παρὰ δὲ Θεῷ πάντα δυνατά.

\vs{27}Τότε ἀποκριθεὶς ὁ Πέτρος εἶπεν αὐτῷ· Ἰδοὺ ἡμεῖς ἀφήκαμεν πάντα καὶ ἠκολουθήσαμέν σοι· τί ἄρα ἔσται ἡμῖν;
\vs{28}Ὁ δὲ Ἰησοῦς εἶπεν αὐτοῖς· Ἀμὴν λέγω ὑμῖν ὅτι ὑμεῖς οἱ ἀκολουθήσαντές μοι ἐν τῇ παλινγενεσίᾳ, ὅταν καθίσῃ ὁ Υἱὸς τοῦ ἀνθρώπου ἐπὶ θρόνου δόξης αὐτοῦ, καθήσεσθε καὶ ὑμεῖς ἐπὶ δώδεκα θρόνους κρίνοντες τὰς δώδεκα φυλὰς τοῦ Ἰσραήλ.
\vs{29}καὶ πᾶς ὅστις ἀφῆκεν οἰκίας ἢ ἀδελφοὺς ἢ ἀδελφὰς ἢ πατέρα ἢ μητέρα ἢ τέκνα ἢ ἀγροὺς ἕνεκεν τοῦ ὀνόματός μου, ἑκατονταπλασίονα λήμψεται καὶ ζωὴν αἰώνιον κληρονομήσει.
\vs{30}Πολλοὶ δὲ ἔσονται πρῶτοι ἔσχατοι καὶ ἔσχατοι πρῶτοι.

\ch{20}
Ὁμοία γάρ ἐστιν ἡ βασιλεία τῶν οὐρανῶν ἀνθρώπῳ οἰκοδεσπότῃ, ὅστις ἐξῆλθεν ἅμα πρωῒ μισθώσασθαι ἐργάτας εἰς τὸν ἀμπελῶνα αὐτοῦ.
\vs{2}συμφωνήσας δὲ μετὰ τῶν ἐργατῶν ἐκ δηναρίου τὴν ἡμέραν ἀπέστειλεν αὐτοὺς εἰς τὸν ἀμπελῶνα αὐτοῦ.
\vs{3}Καὶ ἐξελθὼν περὶ τρίτην ὥραν εἶδεν ἄλλους ἑστῶτας ἐν τῇ ἀγορᾷ ἀργούς
\vs{4}καὶ ἐκείνοις εἶπεν· Ὑπάγετε καὶ ὑμεῖς εἰς τὸν ἀμπελῶνα, καὶ ὃ ἐὰν ᾖ δίκαιον δώσω ὑμῖν.
\vs{5}οἱ δὲ ἀπῆλθον. Πάλιν δὲ ἐξελθὼν περὶ ἕκτην καὶ ἐνάτην ὥραν ἐποίησεν ὡσαύτως.
\vs{6}Περὶ δὲ τὴν ἑνδεκάτην ἐξελθὼν εὗρεν ἄλλους ἑστῶτας καὶ λέγει αὐτοῖς· Τί ὧδε ἑστήκατε ὅλην τὴν ἡμέραν ἀργοί;
\vs{7}Λέγουσιν αὐτῷ· Ὅτι οὐδεὶς ἡμᾶς ἐμισθώσατο. Λέγει αὐτοῖς· Ὑπάγετε καὶ ὑμεῖς εἰς τὸν ἀμπελῶνα.
\vs{8}Ὀψίας δὲ γενομένης λέγει ὁ κύριος τοῦ ἀμπελῶνος τῷ ἐπιτρόπῳ αὐτοῦ· Κάλεσον τοὺς ἐργάτας καὶ ἀπόδος αὐτοῖς τὸν μισθόν ἀρξάμενος ἀπὸ τῶν ἐσχάτων ἕως τῶν πρώτων.
\vs{9}Καὶ ἐλθόντες οἱ περὶ τὴν ἑνδεκάτην ὥραν ἔλαβον ἀνὰ δηνάριον.
\vs{10}καὶ ἐλθόντες οἱ πρῶτοι ἐνόμισαν ὅτι πλεῖον λήμψονται· καὶ ἔλαβον τὸ ἀνὰ δηνάριον καὶ αὐτοί.
\vs{11}Λαβόντες δὲ ἐγόγγυζον κατὰ τοῦ οἰκοδεσπότου
\vs{12}λέγοντες· Οὗτοι οἱ ἔσχατοι μίαν ὥραν ἐποίησαν, καὶ ἴσους ἡμῖν αὐτοὺς ἐποίησας τοῖς βαστάσασι τὸ βάρος τῆς ἡμέρας καὶ τὸν καύσωνα.
\vs{13}Ὁ δὲ ἀποκριθεὶς ἑνὶ αὐτῶν εἶπεν· Ἑταῖρε, οὐκ ἀδικῶ σε· οὐχὶ δηναρίου συνεφώνησάς μοι;
\vs{14}ἆρον τὸ σὸν καὶ ὕπαγε. θέλω δὲ τούτῳ τῷ ἐσχάτῳ δοῦναι ὡς καὶ σοί·
\vs{15}ἢ οὐκ ἔξεστίν μοι ὃ θέλω ποιῆσαι ἐν τοῖς ἐμοῖς; ἢ ὁ ὀφθαλμός σου πονηρός ἐστιν ὅτι ἐγὼ ἀγαθός εἰμι;
\vs{16}Οὕτως ἔσονται οἱ ἔσχατοι πρῶτοι καὶ οἱ πρῶτοι ἔσχατοι.

\vs{17}Καὶ ἀναβαίνων ὁ Ἰησοῦς εἰς Ἱεροσόλυμα παρέλαβεν τοὺς δώδεκα μαθητὰς κατ᾽ ἰδίαν καὶ ἐν τῇ ὁδῷ εἶπεν αὐτοῖς·
\vs{18}Ἰδοὺ ἀναβαίνομεν εἰς Ἱεροσόλυμα, καὶ ὁ Υἱὸς τοῦ ἀνθρώπου παραδοθήσεται τοῖς ἀρχιερεῦσιν καὶ γραμματεῦσιν, καὶ κατακρινοῦσιν αὐτὸν θανάτῳ
\vs{19}καὶ παραδώσουσιν αὐτὸν τοῖς ἔθνεσιν εἰς τὸ ἐμπαῖξαι καὶ μαστιγῶσαι καὶ σταυρῶσαι, καὶ τῇ τρίτῃ ἡμέρᾳ ἐγερθήσεται.

\vs{20}Τότε προσῆλθεν αὐτῷ ἡ μήτηρ τῶν υἱῶν Ζεβεδαίου μετὰ τῶν υἱῶν αὐτῆς προσκυνοῦσα καὶ αἰτοῦσά τι ἀπ᾽ αὐτοῦ.
\vs{21}Ὁ δὲ εἶπεν αὐτῇ· Τί θέλεις; Λέγει αὐτῷ· Εἰπὲ ἵνα καθίσωσιν οὗτοι οἱ δύο υἱοί μου εἷς ἐκ δεξιῶν σου καὶ εἷς ἐξ εὐωνύμων σου ἐν τῇ βασιλείᾳ σου.
\vs{22}Ἀποκριθεὶς δὲ ὁ Ἰησοῦς εἶπεν· Οὐκ οἴδατε τί αἰτεῖσθε. δύνασθε πιεῖν τὸ ποτήριον ὃ ἐγὼ μέλλω πίνειν; Λέγουσιν αὐτῷ· Δυνάμεθα.
\vs{23}Λέγει αὐτοῖς· Τὸ μὲν ποτήριόν μου πίεσθε, τὸ δὲ καθίσαι ἐκ δεξιῶν μου καὶ ἐξ εὐωνύμων οὐκ ἔστιν ἐμὸν τοῦτο δοῦναι, ἀλλ᾽ οἷς ἡτοίμασται ὑπὸ τοῦ Πατρός μου.

\vs{24}Καὶ ἀκούσαντες οἱ δέκα ἠγανάκτησαν περὶ τῶν δύο ἀδελφῶν.
\vs{25}ὁ δὲ Ἰησοῦς προσκαλεσάμενος αὐτοὺς εἶπεν· Οἴδατε ὅτι οἱ ἄρχοντες τῶν ἐθνῶν κατακυριεύουσιν αὐτῶν καὶ οἱ μεγάλοι κατεξουσιάζουσιν αὐτῶν.
\vs{26}οὐχ οὕτως ἔσται ἐν ὑμῖν, ἀλλ᾽ ὃς ἐὰν θέλῃ ἐν ὑμῖν μέγας γενέσθαι ἔσται ὑμῶν διάκονος,
\vs{27}καὶ ὃς ἂν θέλῃ ἐν ὑμῖν εἶναι πρῶτος ἔσται ὑμῶν δοῦλος·
\vs{28}ὥσπερ ὁ Υἱὸς τοῦ ἀνθρώπου οὐκ ἦλθεν διακονηθῆναι ἀλλὰ διακονῆσαι καὶ δοῦναι τὴν ψυχὴν αὐτοῦ λύτρον ἀντὶ πολλῶν.

\vs{29}Καὶ ἐκπορευομένων αὐτῶν ἀπὸ Ἰεριχὼ ἠκολούθησεν αὐτῷ ὄχλος πολύς.
\vs{30}καὶ ἰδοὺ δύο τυφλοὶ καθήμενοι παρὰ τὴν ὁδόν ἀκούσαντες ὅτι Ἰησοῦς παράγει, ἔκραξαν λέγοντες· Ἐλέησον ἡμᾶς, Κύριε, υἱὸς Δαυίδ.
\vs{31}Ὁ δὲ ὄχλος ἐπετίμησεν αὐτοῖς ἵνα σιωπήσωσιν· οἱ δὲ μεῖζον ἔκραξαν λέγοντες· Ἐλέησον ἡμᾶς, Κύριε, υἱὸς Δαυίδ.
\vs{32}Καὶ στὰς ὁ Ἰησοῦς ἐφώνησεν αὐτοὺς καὶ εἶπεν· Τί θέλετε ποιήσω ὑμῖν;
\vs{33}Λέγουσιν αὐτῷ· Κύριε, ἵνα ἀνοιγῶσιν οἱ ὀφθαλμοὶ ἡμῶν.
\vs{34}Σπλαγχνισθεὶς δὲ ὁ Ἰησοῦς ἥψατο τῶν ὀμμάτων αὐτῶν, καὶ εὐθέως ἀνέβλεψαν καὶ ἠκολούθησαν αὐτῷ.

\ch{21}
Καὶ ὅτε ἤγγισαν εἰς Ἱεροσόλυμα καὶ ἦλθον εἰς Βηθφαγὴ εἰς τὸ ὄρος τῶν Ἐλαιῶν, τότε Ἰησοῦς ἀπέστειλεν δύο μαθητὰς
\vs{2}λέγων αὐτοῖς· Πορεύεσθε εἰς τὴν κώμην τὴν κατέναντι ὑμῶν, καὶ εὐθέως εὑρήσετε ὄνον δεδεμένην καὶ πῶλον μετ᾽ αὐτῆς· λύσαντες ἀγάγετέ μοι.
\vs{3}καὶ ἐάν τις ὑμῖν εἴπῃ τι, ἐρεῖτε ὅτι Ὁ Κύριος αὐτῶν χρείαν ἔχει· εὐθὺς δὲ ἀποστελεῖ αὐτούς.
\vs{4}Τοῦτο δὲ γέγονεν ἵνα πληρωθῇ τὸ ῥηθὲν διὰ τοῦ προφήτου λέγοντος·
\begin{poetryblock}

\begin{quote} \vs{5}Εἴπατε τῇ θυγατρὶ Σιών·\end{quote} 

\begin{quote}Ἰδοὺ ὁ Βασιλεύς σου ἔρχεταί σοι\end{quote} 

\begin{quote}πραῢς καὶ ἐπιβεβηκὼς ἐπὶ ὄνον\end{quote} 

\begin{quote}καὶ ἐπὶ πῶλον υἱὸν ὑποζυγίου.\end{quote}
\end{poetryblock}

\vs{6}Πορευθέντες δὲ οἱ μαθηταὶ καὶ ποιήσαντες καθὼς συνέταξεν αὐτοῖς ὁ Ἰησοῦς
\vs{7}ἤγαγον τὴν ὄνον καὶ τὸν πῶλον καὶ ἐπέθηκαν ἐπ᾽ αὐτῶν τὰ ἱμάτια, καὶ ἐπεκάθισεν ἐπάνω αὐτῶν.
\vs{8}Ὁ δὲ πλεῖστος ὄχλος ἔστρωσαν ἑαυτῶν τὰ ἱμάτια ἐν τῇ ὁδῷ, ἄλλοι δὲ ἔκοπτον κλάδους ἀπὸ τῶν δένδρων καὶ ἐστρώννυον ἐν τῇ ὁδῷ.
\vs{9}Οἱ δὲ ὄχλοι οἱ προάγοντες αὐτὸν καὶ οἱ ἀκολουθοῦντες ἔκραζον λέγοντες· 
\begin{poetryblock}

\begin{quote}Ὡσαννὰ τῷ υἱῷ Δαυίδ·\end{quote} 

\begin{quote}Εὐλογημένος ὁ ἐρχόμενος ἐν ὀνόματι Κυρίου·\end{quote} 

\begin{quote}Ὡσαννὰ ἐν τοῖς ὑψίστοις.\end{quote}
\end{poetryblock}

\vs{10}Καὶ εἰσελθόντος αὐτοῦ εἰς Ἱεροσόλυμα ἐσείσθη πᾶσα ἡ πόλις λέγουσα· Τίς ἐστιν οὗτος;
\vs{11}Οἱ δὲ ὄχλοι ἔλεγον· Οὗτός ἐστιν ὁ προφήτης Ἰησοῦς ὁ ἀπὸ Ναζαρὲθ τῆς Γαλιλαίας.

\vs{12}Καὶ εἰσῆλθεν Ἰησοῦς εἰς τὸ ἱερόν καὶ ἐξέβαλεν πάντας τοὺς πωλοῦντας καὶ ἀγοράζοντας ἐν τῷ ἱερῷ, καὶ τὰς τραπέζας τῶν κολλυβιστῶν κατέστρεψεν καὶ τὰς καθέδρας τῶν πωλούντων τὰς περιστεράς,
\vs{13}καὶ λέγει αὐτοῖς· Γέγραπται· 
\begin{poetryblock}

\begin{quote}Ὁ οἶκός μου οἶκος προσευχῆς κληθήσεται,\end{quote} 

\begin{quote}ὑμεῖς δὲ αὐτὸν ποιεῖτε Σπήλαιον λῃστῶν.\end{quote}
\end{poetryblock}

\vs{14}Καὶ προσῆλθον αὐτῷ τυφλοὶ καὶ χωλοὶ ἐν τῷ ἱερῷ, καὶ ἐθεράπευσεν αὐτούς.
\vs{15}ἰδόντες δὲ οἱ ἀρχιερεῖς καὶ οἱ γραμματεῖς τὰ θαυμάσια ἃ ἐποίησεν καὶ τοὺς παῖδας τοὺς κράζοντας ἐν τῷ ἱερῷ καὶ λέγοντας· Ὡσαννὰ τῷ υἱῷ Δαυίδ, ἠγανάκτησαν
\vs{16}καὶ εἶπαν αὐτῷ· Ἀκούεις τί οὗτοι λέγουσιν; Ὁ δὲ Ἰησοῦς λέγει αὐτοῖς· Ναί. οὐδέποτε ἀνέγνωτε ὅτι 
\begin{poetryblock}

\begin{quote}Ἐκ στόματος νηπίων καὶ θηλαζόντων Κατηρτίσω αἶνον;\end{quote}
\end{poetryblock}

\vs{17}Καὶ καταλιπὼν αὐτοὺς ἐξῆλθεν ἔξω τῆς πόλεως εἰς Βηθανίαν καὶ ηὐλίσθη ἐκεῖ.

\vs{18}Πρωῒ δὲ ἐπανάγων εἰς τὴν πόλιν ἐπείνασεν.
\vs{19}καὶ ἰδὼν συκῆν μίαν ἐπὶ τῆς ὁδοῦ ἦλθεν ἐπ᾽ αὐτήν καὶ οὐδὲν εὗρεν ἐν αὐτῇ εἰ μὴ φύλλα μόνον, καὶ λέγει αὐτῇ· μηκέτι ἐκ σοῦ καρπὸς γένηται εἰς τὸν αἰῶνα. καὶ ἐξηράνθη παραχρῆμα ἡ συκῆ.

\vs{20}Καὶ ἰδόντες οἱ μαθηταὶ ἐθαύμασαν λέγοντες· Πῶς παραχρῆμα ἐξηράνθη ἡ συκῆ;
\vs{21}Ἀποκριθεὶς δὲ ὁ Ἰησοῦς εἶπεν αὐτοῖς· Ἀμὴν λέγω ὑμῖν, ἐὰν ἔχητε πίστιν καὶ μὴ διακριθῆτε, οὐ μόνον τὸ τῆς συκῆς ποιήσετε, ἀλλὰ κἂν τῷ ὄρει τούτῳ εἴπητε· Ἄρθητι καὶ βλήθητι εἰς τὴν θάλασσαν, γενήσεται·
\vs{22}καὶ πάντα ὅσα ἂν αἰτήσητε ἐν τῇ προσευχῇ πιστεύοντες λήμψεσθε.

\vs{23}Καὶ ἐλθόντος αὐτοῦ εἰς τὸ ἱερὸν προσῆλθον αὐτῷ διδάσκοντι οἱ ἀρχιερεῖς καὶ οἱ πρεσβύτεροι τοῦ λαοῦ λέγοντες· Ἐν ποίᾳ ἐξουσίᾳ ταῦτα ποιεῖς; καὶ τίς σοι ἔδωκεν τὴν ἐξουσίαν ταύτην;
\vs{24}Ἀποκριθεὶς δὲ ὁ Ἰησοῦς εἶπεν αὐτοῖς· Ἐρωτήσω ὑμᾶς κἀγὼ λόγον ἕνα, ὃν ἐὰν εἴπητέ μοι κἀγὼ ὑμῖν ἐρῶ ἐν ποίᾳ ἐξουσίᾳ ταῦτα ποιῶ·
\vs{25}τὸ βάπτισμα τὸ Ἰωάννου πόθεν ἦν; ἐξ οὐρανοῦ ἢ ἐξ ἀνθρώπων; Οἱ δὲ διελογίζοντο ἐν ἑαυτοῖς λέγοντες· Ἐὰν εἴπωμεν· Ἐξ οὐρανοῦ, ἐρεῖ ἡμῖν· Διὰ τί οὖν οὐκ ἐπιστεύσατε αὐτῷ;
\vs{26}ἐὰν δὲ εἴπωμεν· Ἐξ ἀνθρώπων, φοβούμεθα τὸν ὄχλον, πάντες γὰρ ὡς προφήτην ἔχουσιν τὸν Ἰωάννην.
\vs{27}καὶ ἀποκριθέντες τῷ Ἰησοῦ εἶπαν· Οὐκ οἴδαμεν. Ἔφη αὐτοῖς καὶ αὐτός· Οὐδὲ ἐγὼ λέγω ὑμῖν ἐν ποίᾳ ἐξουσίᾳ ταῦτα ποιῶ.

\vs{28}Τί δὲ ὑμῖν δοκεῖ; ἄνθρωπος εἶχεν τέκνα δύο. καὶ προσελθὼν τῷ πρώτῳ εἶπεν· Τέκνον, ὕπαγε σήμερον ἐργάζου ἐν τῷ ἀμπελῶνι.
\vs{29}Ὁ δὲ ἀποκριθεὶς εἶπεν· Οὐ θέλω, ὕστερον δὲ μεταμεληθεὶς ἀπῆλθεν.
\vs{30}Προσελθὼν δὲ τῷ ἑτέρῳ εἶπεν ὡσαύτως. Ὁ δὲ ἀποκριθεὶς εἶπεν· Ἐγώ, κύριε, καὶ οὐκ ἀπῆλθεν.
\vs{31}Τίς ἐκ τῶν δύο ἐποίησεν τὸ θέλημα τοῦ πατρός; Λέγουσιν· Ὁ πρῶτος. Λέγει αὐτοῖς ὁ Ἰησοῦς· Ἀμὴν λέγω ὑμῖν ὅτι οἱ τελῶναι καὶ αἱ πόρναι προάγουσιν ὑμᾶς εἰς τὴν βασιλείαν τοῦ Θεοῦ.
\vs{32}ἦλθεν γὰρ Ἰωάννης πρὸς ὑμᾶς ἐν ὁδῷ δικαιοσύνης, καὶ οὐκ ἐπιστεύσατε αὐτῷ, οἱ δὲ τελῶναι καὶ αἱ πόρναι ἐπίστευσαν αὐτῷ· ὑμεῖς δὲ ἰδόντες οὐδὲ μετεμελήθητε ὕστερον τοῦ πιστεῦσαι αὐτῷ.

\vs{33}Ἄλλην παραβολὴν ἀκούσατε. Ἄνθρωπος ἦν οἰκοδεσπότης ὅστις ἐφύτευσεν ἀμπελῶνα καὶ φραγμὸν αὐτῷ περιέθηκεν καὶ ὤρυξεν ἐν αὐτῷ ληνὸν καὶ ᾠκοδόμησεν πύργον καὶ ἐξέδετο αὐτὸν γεωργοῖς καὶ ἀπεδήμησεν.
\vs{34}Ὅτε δὲ ἤγγισεν ὁ καιρὸς τῶν καρπῶν, ἀπέστειλεν τοὺς δούλους αὐτοῦ πρὸς τοὺς γεωργοὺς λαβεῖν τοὺς καρποὺς αὐτοῦ.
\vs{35}καὶ λαβόντες οἱ γεωργοὶ τοὺς δούλους αὐτοῦ ὃν μὲν ἔδειραν, ὃν δὲ ἀπέκτειναν, ὃν δὲ ἐλιθοβόλησαν.
\vs{36}Πάλιν ἀπέστειλεν ἄλλους δούλους πλείονας τῶν πρώτων, καὶ ἐποίησαν αὐτοῖς ὡσαύτως.
\vs{37}Ὕστερον δὲ ἀπέστειλεν πρὸς αὐτοὺς τὸν υἱὸν αὐτοῦ λέγων· Ἐντραπήσονται τὸν υἱόν μου.
\vs{38}Οἱ δὲ γεωργοὶ ἰδόντες τὸν υἱὸν εἶπον ἐν ἑαυτοῖς· Οὗτός ἐστιν ὁ κληρονόμος· δεῦτε ἀποκτείνωμεν αὐτὸν καὶ σχῶμεν τὴν κληρονομίαν αὐτοῦ,
\vs{39}καὶ λαβόντες αὐτὸν ἐξέβαλον ἔξω τοῦ ἀμπελῶνος καὶ ἀπέκτειναν.
\vs{40}Ὅταν οὖν ἔλθῃ ὁ κύριος τοῦ ἀμπελῶνος, τί ποιήσει τοῖς γεωργοῖς ἐκείνοις;
\vs{41}Λέγουσιν αὐτῷ· Κακοὺς κακῶς ἀπολέσει αὐτούς καὶ τὸν ἀμπελῶνα ἐκδώσεται ἄλλοις γεωργοῖς, οἵτινες ἀποδώσουσιν αὐτῷ τοὺς καρποὺς ἐν τοῖς καιροῖς αὐτῶν.

\vs{42}Λέγει αὐτοῖς ὁ Ἰησοῦς· Οὐδέποτε ἀνέγνωτε ἐν ταῖς γραφαῖς· 
\begin{poetryblock}

\begin{quote}Λίθον ὃν ἀπεδοκίμασαν οἱ οἰκοδομοῦντες,\end{quote} 

\begin{quote}Οὗτος ἐγενήθη εἰς κεφαλὴν γωνίας·\end{quote} 

\begin{quote}Παρὰ Κυρίου ἐγένετο αὕτη\end{quote} 

\begin{quote}Καὶ ἔστιν θαυμαστὴ ἐν ὀφθαλμοῖς ἡμῶν;\end{quote}
\end{poetryblock}

\vs{43}Διὰ τοῦτο λέγω ὑμῖν ὅτι ἀρθήσεται ἀφ᾽ ὑμῶν ἡ βασιλεία τοῦ Θεοῦ καὶ δοθήσεται ἔθνει ποιοῦντι τοὺς καρποὺς αὐτῆς.
\vs{44}καὶ ὁ πεσὼν ἐπὶ τὸν λίθον τοῦτον συνθλασθήσεται· ἐφ᾽ ὃν δ᾽ ἂν πέσῃ λικμήσει αὐτόν.

\vs{45}Καὶ ἀκούσαντες οἱ ἀρχιερεῖς καὶ οἱ Φαρισαῖοι τὰς παραβολὰς αὐτοῦ ἔγνωσαν ὅτι περὶ αὐτῶν λέγει·
\vs{46}καὶ ζητοῦντες αὐτὸν κρατῆσαι ἐφοβήθησαν τοὺς ὄχλους, ἐπεὶ εἰς προφήτην αὐτὸν εἶχον.

\ch{22}
Καὶ ἀποκριθεὶς ὁ Ἰησοῦς πάλιν εἶπεν ἐν παραβολαῖς αὐτοῖς λέγων·
\vs{2}Ὡμοιώθη ἡ βασιλεία τῶν οὐρανῶν ἀνθρώπῳ βασιλεῖ, ὅστις ἐποίησεν γάμους τῷ υἱῷ αὐτοῦ.
\vs{3}καὶ ἀπέστειλεν τοὺς δούλους αὐτοῦ καλέσαι τοὺς κεκλημένους εἰς τοὺς γάμους, καὶ οὐκ ἤθελον ἐλθεῖν.
\vs{4}Πάλιν ἀπέστειλεν ἄλλους δούλους λέγων· Εἴπατε τοῖς κεκλημένοις· Ἰδοὺ τὸ ἄριστόν μου ἡτοίμακα, οἱ ταῦροί μου καὶ τὰ σιτιστὰ τεθυμένα καὶ πάντα ἕτοιμα· δεῦτε εἰς τοὺς γάμους.
\vs{5}Οἱ δὲ ἀμελήσαντες ἀπῆλθον, ὃς μὲν εἰς τὸν ἴδιον ἀγρόν, ὃς δὲ ἐπὶ τὴν ἐμπορίαν αὐτοῦ·
\vs{6}οἱ δὲ λοιποὶ κρατήσαντες τοὺς δούλους αὐτοῦ ὕβρισαν καὶ ἀπέκτειναν.
\vs{7}Ὁ δὲ βασιλεὺς ὠργίσθη καὶ πέμψας τὰ στρατεύματα αὐτοῦ ἀπώλεσεν τοὺς φονεῖς ἐκείνους καὶ τὴν πόλιν αὐτῶν ἐνέπρησεν.
\vs{8}τότε λέγει τοῖς δούλοις αὐτοῦ· Ὁ μὲν γάμος ἕτοιμός ἐστιν, οἱ δὲ κεκλημένοι οὐκ ἦσαν ἄξιοι·
\vs{9}πορεύεσθε οὖν ἐπὶ τὰς διεξόδους τῶν ὁδῶν καὶ ὅσους ἐὰν εὕρητε καλέσατε εἰς τοὺς γάμους.
\vs{10}Καὶ ἐξελθόντες οἱ δοῦλοι ἐκεῖνοι εἰς τὰς ὁδοὺς συνήγαγον πάντας οὓς εὗρον, πονηρούς τε καὶ ἀγαθούς· καὶ ἐπλήσθη ὁ γάμος ἀνακειμένων.
\vs{11}Εἰσελθὼν δὲ ὁ βασιλεὺς θεάσασθαι τοὺς ἀνακειμένους εἶδεν ἐκεῖ ἄνθρωπον οὐκ ἐνδεδυμένον ἔνδυμα γάμου,
\vs{12}καὶ λέγει αὐτῷ· Ἑταῖρε, πῶς εἰσῆλθες ὧδε μὴ ἔχων ἔνδυμα γάμου; Ὁ δὲ ἐφιμώθη.
\vs{13}Τότε ὁ βασιλεὺς εἶπεν τοῖς διακόνοις· Δήσαντες αὐτοῦ πόδας καὶ χεῖρας ἐκβάλετε αὐτὸν εἰς τὸ σκότος τὸ ἐξώτερον· ἐκεῖ ἔσται ὁ κλαυθμὸς καὶ ὁ βρυγμὸς τῶν ὀδόντων.
\vs{14}Πολλοὶ γάρ εἰσιν κλητοὶ, ὀλίγοι δὲ ἐκλεκτοί.

\vs{15}Τότε πορευθέντες οἱ Φαρισαῖοι συμβούλιον ἔλαβον ὅπως αὐτὸν παγιδεύσωσιν ἐν λόγῳ.
\vs{16}καὶ ἀποστέλλουσιν αὐτῷ τοὺς μαθητὰς αὐτῶν μετὰ τῶν Ἡρῳδιανῶν λέγοντες· Διδάσκαλε, οἴδαμεν ὅτι ἀληθὴς εἶ καὶ τὴν ὁδὸν τοῦ Θεοῦ ἐν ἀληθείᾳ διδάσκεις καὶ οὐ μέλει σοι περὶ οὐδενός· οὐ γὰρ βλέπεις εἰς πρόσωπον ἀνθρώπων,
\vs{17}εἰπὲ οὖν ἡμῖν τί σοι δοκεῖ· ἔξεστιν δοῦναι κῆνσον Καίσαρι ἢ οὔ;
\vs{18}Γνοὺς δὲ ὁ Ἰησοῦς τὴν πονηρίαν αὐτῶν εἶπεν· Τί με πειράζετε, ὑποκριταί;
\vs{19}ἐπιδείξατέ μοι τὸ νόμισμα τοῦ κήνσου. Οἱ δὲ προσήνεγκαν αὐτῷ δηνάριον.
\vs{20}Καὶ λέγει αὐτοῖς· Τίνος ἡ εἰκὼν αὕτη καὶ ἡ ἐπιγραφή;
\vs{21}Λέγουσιν αὐτῷ· Καίσαρος. Τότε λέγει αὐτοῖς· Ἀπόδοτε οὖν τὰ Καίσαρος Καίσαρι καὶ τὰ τοῦ Θεοῦ τῷ Θεῷ.
\vs{22}Καὶ ἀκούσαντες ἐθαύμασαν, καὶ ἀφέντες αὐτὸν ἀπῆλθαν.

\vs{23}Ἐν ἐκείνῃ τῇ ἡμέρᾳ προσῆλθον αὐτῷ Σαδδουκαῖοι, λέγοντες μὴ εἶναι ἀνάστασιν, καὶ ἐπηρώτησαν αὐτὸν
\vs{24}λέγοντες· Διδάσκαλε, Μωϋσῆς εἶπεν· Ἐάν τις ἀποθάνῃ μὴ ἔχων τέκνα, ἐπιγαμβρεύσει ὁ ἀδελφὸς αὐτοῦ τὴν γυναῖκα αὐτοῦ καὶ ἀναστήσει σπέρμα τῷ ἀδελφῷ αὐτοῦ.
\vs{25}ἦσαν δὲ παρ᾽ ἡμῖν ἑπτὰ ἀδελφοί· καὶ ὁ πρῶτος γήμας ἐτελεύτησεν, καὶ μὴ ἔχων σπέρμα ἀφῆκεν τὴν γυναῖκα αὐτοῦ τῷ ἀδελφῷ αὐτοῦ·
\vs{26}ὁμοίως καὶ ὁ δεύτερος καὶ ὁ τρίτος ἕως τῶν ἑπτά.
\vs{27}ὕστερον δὲ πάντων ἀπέθανεν ἡ γυνή.
\vs{28}ἐν τῇ ἀναστάσει οὖν τίνος τῶν ἑπτὰ ἔσται γυνή; πάντες γὰρ ἔσχον αὐτήν·
\vs{29}Ἀποκριθεὶς δὲ ὁ Ἰησοῦς εἶπεν αὐτοῖς· Πλανᾶσθε μὴ εἰδότες τὰς γραφὰς μηδὲ τὴν δύναμιν τοῦ Θεοῦ·
\vs{30}ἐν γὰρ τῇ ἀναστάσει οὔτε γαμοῦσιν οὔτε γαμίζονται, ἀλλ᾽ ὡς ἄγγελοι ἐν τῷ οὐρανῷ εἰσιν.
\vs{31}περὶ δὲ τῆς ἀναστάσεως τῶν νεκρῶν οὐκ ἀνέγνωτε τὸ ῥηθὲν ὑμῖν ὑπὸ τοῦ Θεοῦ λέγοντος·
\vs{32}Ἐγώ εἰμι ὁ Θεὸς Ἀβραὰμ καὶ ὁ Θεὸς Ἰσαὰκ καὶ ὁ Θεὸς Ἰακώβ; οὐκ ἔστιν ὁ Θεὸς νεκρῶν ἀλλὰ ζώντων.
\vs{33}Καὶ ἀκούσαντες οἱ ὄχλοι ἐξεπλήσσοντο ἐπὶ τῇ διδαχῇ αὐτοῦ.

\vs{34}Οἱ δὲ Φαρισαῖοι ἀκούσαντες ὅτι ἐφίμωσεν τοὺς Σαδδουκαίους συνήχθησαν ἐπὶ τὸ αὐτό,
\vs{35}καὶ ἐπηρώτησεν εἷς ἐξ αὐτῶν νομικὸς πειράζων αὐτόν·
\vs{36}Διδάσκαλε, ποία ἐντολὴ μεγάλη ἐν τῷ νόμῳ;
\vs{37}Ὁ δὲ ἔφη αὐτῷ· Ἀγαπήσεις κύριον τὸν Θεόν σου ἐν ὅλῃ τῇ καρδίᾳ σου καὶ ἐν ὅλῃ τῇ ψυχῇ σου καὶ ἐν ὅλῃ τῇ διανοίᾳ σου·
\vs{38}αὕτη ἐστὶν ἡ μεγάλη καὶ πρώτη ἐντολή.
\vs{39}δευτέρα δὲ ὁμοία αὐτῇ· Ἀγαπήσεις τὸν πλησίον σου ὡς σεαυτόν.
\vs{40}ἐν ταύταις ταῖς δυσὶν ἐντολαῖς ὅλος ὁ νόμος κρέμαται καὶ οἱ προφῆται.

\vs{41}Συνηγμένων δὲ τῶν Φαρισαίων ἐπηρώτησεν αὐτοὺς ὁ Ἰησοῦς
\vs{42}λέγων· Τί ὑμῖν δοκεῖ περὶ τοῦ Χριστοῦ; τίνος υἱός ἐστιν; Λέγουσιν αὐτῷ· Τοῦ Δαυίδ.
\vs{43}Λέγει αὐτοῖς· Πῶς οὖν Δαυὶδ ἐν Πνεύματι καλεῖ αὐτὸν Κύριον λέγων·
\begin{poetryblock}

\begin{quote} \vs{44}Εἶπεν Κύριος τῷ Κυρίῳ μου·\end{quote} 

\begin{quote}Κάθου ἐκ δεξιῶν μου,\end{quote} 

\begin{quote}Ἕως ἂν θῶ τοὺς ἐχθρούς σου\end{quote} 

\begin{quote}Ὑποκάτω τῶν ποδῶν σου;\end{quote}
\end{poetryblock}

\vs{45}Εἰ οὖν Δαυὶδ καλεῖ αὐτὸν Κύριον, πῶς υἱὸς αὐτοῦ ἐστιν;
\vs{46}Καὶ οὐδεὶς ἐδύνατο ἀποκριθῆναι αὐτῷ λόγον οὐδὲ ἐτόλμησέν τις ἀπ᾽ ἐκείνης τῆς ἡμέρας ἐπερωτῆσαι αὐτὸν οὐκέτι.

\ch{23}
Τότε ὁ Ἰησοῦς ἐλάλησεν τοῖς ὄχλοις καὶ τοῖς μαθηταῖς αὐτοῦ
\vs{2}λέγων· Ἐπὶ τῆς Μωϋσέως καθέδρας ἐκάθισαν οἱ γραμματεῖς καὶ οἱ Φαρισαῖοι.
\vs{3}πάντα οὖν ὅσα ἐὰν εἴπωσιν ὑμῖν ποιήσατε καὶ τηρεῖτε, κατὰ δὲ τὰ ἔργα αὐτῶν μὴ ποιεῖτε· λέγουσιν γὰρ καὶ οὐ ποιοῦσιν.
\vs{4}δεσμεύουσιν δὲ φορτία βαρέα καὶ δυσβάστακτα καὶ ἐπιτιθέασιν ἐπὶ τοὺς ὤμους τῶν ἀνθρώπων, αὐτοὶ δὲ τῷ δακτύλῳ αὐτῶν οὐ θέλουσιν κινῆσαι αὐτά.
\vs{5}Πάντα δὲ τὰ ἔργα αὐτῶν ποιοῦσιν πρὸς τὸ θεαθῆναι τοῖς ἀνθρώποις· πλατύνουσιν γὰρ τὰ φυλακτήρια αὐτῶν καὶ μεγαλύνουσιν τὰ κράσπεδα,
\vs{6}φιλοῦσιν δὲ τὴν πρωτοκλισίαν ἐν τοῖς δείπνοις καὶ τὰς πρωτοκαθεδρίας ἐν ταῖς συναγωγαῖς
\vs{7}καὶ τοὺς ἀσπασμοὺς ἐν ταῖς ἀγοραῖς καὶ καλεῖσθαι ὑπὸ τῶν ἀνθρώπων Ῥαββί.
\vs{8}Ὑμεῖς δὲ μὴ κληθῆτε Ῥαββί· εἷς γάρ ἐστιν ὑμῶν ὁ διδάσκαλος, πάντες δὲ ὑμεῖς ἀδελφοί ἐστε.
\vs{9}καὶ πατέρα μὴ καλέσητε ὑμῶν ἐπὶ τῆς γῆς, εἷς γάρ ἐστιν ὑμῶν ὁ Πατὴρ ὁ οὐράνιος.
\vs{10}μηδὲ κληθῆτε καθηγηταί, ὅτι καθηγητὴς ὑμῶν ἐστιν εἷς ὁ Χριστός.
\vs{11}ὁ δὲ μείζων ὑμῶν ἔσται ὑμῶν διάκονος.
\vs{12}Ὅστις δὲ ὑψώσει ἑαυτὸν ταπεινωθήσεται καὶ ὅστις ταπεινώσει ἑαυτὸν ὑψωθήσεται.

\vs{13}Οὐαὶ δὲ ὑμῖν, γραμματεῖς καὶ Φαρισαῖοι ὑποκριταί, ὅτι κλείετε τὴν βασιλείαν τῶν οὐρανῶν ἔμπροσθεν τῶν ἀνθρώπων· ὑμεῖς γὰρ οὐκ εἰσέρχεσθε οὐδὲ τοὺς εἰσερχομένους ἀφίετε εἰσελθεῖν.

\vs{15}Οὐαὶ ὑμῖν, γραμματεῖς καὶ Φαρισαῖοι ὑποκριταί, ὅτι περιάγετε τὴν θάλασσαν καὶ τὴν ξηρὰν ποιῆσαι ἕνα προσήλυτον, καὶ ὅταν γένηται ποιεῖτε αὐτὸν υἱὸν γεέννης διπλότερον ὑμῶν.

\vs{16}Οὐαὶ ὑμῖν, ὁδηγοὶ τυφλοὶ οἱ λέγοντες· Ὃς ἂν ὀμόσῃ ἐν τῷ ναῷ, οὐδέν ἐστιν· ὃς δ᾽ ἂν ὀμόσῃ ἐν τῷ χρυσῷ τοῦ ναοῦ, ὀφείλει.
\vs{17}μωροὶ καὶ τυφλοί, τίς γὰρ μείζων ἐστίν, ὁ χρυσὸς ἢ ὁ ναὸς ὁ ἁγιάσας τὸν χρυσόν;
\vs{18}καί· Ὃς ἂν ὀμόσῃ ἐν τῷ θυσιαστηρίῳ, οὐδέν ἐστιν· ὃς δ᾽ ἂν ὀμόσῃ ἐν τῷ δώρῳ τῷ ἐπάνω αὐτοῦ, ὀφείλει.
\vs{19}τυφλοί, τί γὰρ μεῖζον, τὸ δῶρον ἢ τὸ θυσιαστήριον τὸ ἁγιάζον τὸ δῶρον;
\vs{20}ὁ οὖν ὀμόσας ἐν τῷ θυσιαστηρίῳ ὀμνύει ἐν αὐτῷ καὶ ἐν πᾶσι τοῖς ἐπάνω αὐτοῦ·
\vs{21}καὶ ὁ ὀμόσας ἐν τῷ ναῷ ὀμνύει ἐν αὐτῷ καὶ ἐν τῷ κατοικοῦντι αὐτόν,
\vs{22}καὶ ὁ ὀμόσας ἐν τῷ οὐρανῷ ὀμνύει ἐν τῷ θρόνῳ τοῦ Θεοῦ καὶ ἐν τῷ καθημένῳ ἐπάνω αὐτοῦ.

\vs{23}Οὐαὶ ὑμῖν, γραμματεῖς καὶ Φαρισαῖοι ὑποκριταί, ὅτι ἀποδεκατοῦτε τὸ ἡδύοσμον καὶ τὸ ἄνηθον καὶ τὸ κύμινον καὶ ἀφήκατε τὰ βαρύτερα τοῦ νόμου, τὴν κρίσιν καὶ τὸ ἔλεος καὶ τὴν πίστιν· ταῦτα δὲ ἔδει ποιῆσαι κἀκεῖνα μὴ ἀφιέναι.
\vs{24}ὁδηγοὶ τυφλοί, οἱ διϋλίζοντες τὸν κώνωπα, τὴν δὲ κάμηλον καταπίνοντες.

\vs{25}Οὐαὶ ὑμῖν, γραμματεῖς καὶ Φαρισαῖοι ὑποκριταί, ὅτι καθαρίζετε τὸ ἔξωθεν τοῦ ποτηρίου καὶ τῆς παροψίδος, ἔσωθεν δὲ γέμουσιν ἐξ ἁρπαγῆς καὶ ἀκρασίας.
\vs{26}Φαρισαῖε τυφλέ, καθάρισον πρῶτον τὸ ἐντὸς τοῦ ποτηρίου, ἵνα γένηται καὶ τὸ ἐκτὸς αὐτοῦ καθαρόν.

\vs{27}Οὐαὶ ὑμῖν, γραμματεῖς καὶ Φαρισαῖοι ὑποκριταί, ὅτι παρομοιάζετε τάφοις κεκονιαμένοις, οἵτινες ἔξωθεν μὲν φαίνονται ὡραῖοι, ἔσωθεν δὲ γέμουσιν ὀστέων νεκρῶν καὶ πάσης ἀκαθαρσίας.
\vs{28}οὕτως καὶ ὑμεῖς ἔξωθεν μὲν φαίνεσθε τοῖς ἀνθρώποις δίκαιοι, ἔσωθεν δέ ἐστε μεστοὶ ὑποκρίσεως καὶ ἀνομίας.

\vs{29}Οὐαὶ ὑμῖν, γραμματεῖς καὶ Φαρισαῖοι ὑποκριταί, ὅτι οἰκοδομεῖτε τοὺς τάφους τῶν προφητῶν καὶ κοσμεῖτε τὰ μνημεῖα τῶν δικαίων,
\vs{30}καὶ λέγετε· Εἰ ἤμεθα ἐν ταῖς ἡμέραις τῶν πατέρων ἡμῶν, οὐκ ἂν ἤμεθα αὐτῶν κοινωνοὶ ἐν τῷ αἵματι τῶν προφητῶν.
\vs{31}ὥστε μαρτυρεῖτε ἑαυτοῖς ὅτι υἱοί ἐστε τῶν φονευσάντων τοὺς προφήτας.
\vs{32}καὶ ὑμεῖς πληρώσατε τὸ μέτρον τῶν πατέρων ὑμῶν.
\vs{33}ὄφεις, γεννήματα ἐχιδνῶν, πῶς φύγητε ἀπὸ τῆς κρίσεως τῆς γεέννης;

\vs{34}Διὰ τοῦτο ἰδοὺ ἐγὼ ἀποστέλλω πρὸς ὑμᾶς προφήτας καὶ σοφοὺς καὶ γραμματεῖς· ἐξ αὐτῶν ἀποκτενεῖτε καὶ σταυρώσετε καὶ ἐξ αὐτῶν μαστιγώσετε ἐν ταῖς συναγωγαῖς ὑμῶν καὶ διώξετε ἀπὸ πόλεως εἰς πόλιν·
\vs{35}ὅπως ἔλθῃ ἐφ᾽ ὑμᾶς πᾶν αἷμα δίκαιον ἐκχυννόμενον ἐπὶ τῆς γῆς ἀπὸ τοῦ αἵματος Ἅβελ τοῦ δικαίου ἕως τοῦ αἵματος Ζαχαρίου υἱοῦ Βαραχίου, ὃν ἐφονεύσατε μεταξὺ τοῦ ναοῦ καὶ τοῦ θυσιαστηρίου.
\vs{36}ἀμὴν λέγω ὑμῖν, ἥξει ταῦτα πάντα ἐπὶ τὴν γενεὰν ταύτην.

\vs{37}Ἰερουσαλὴμ Ἰερουσαλήμ, ἡ ἀποκτείνουσα τοὺς προφήτας καὶ λιθοβολοῦσα τοὺς ἀπεσταλμένους πρὸς αὐτήν, ποσάκις ἠθέλησα ἐπισυναγαγεῖν τὰ τέκνα σου, ὃν τρόπον ὄρνις ἐπισυνάγει τὰ νοσσία αὐτῆς ὑπὸ τὰς πτέρυγας, καὶ οὐκ ἠθελήσατε.
\vs{38}ἰδοὺ ἀφίεται ὑμῖν ὁ οἶκος ὑμῶν ἔρημος.
\vs{39}λέγω γὰρ ὑμῖν, οὐ μή με ἴδητε ἀπ᾽ ἄρτι ἕως ἂν εἴπητε· 
\begin{poetryblock}

\begin{quote}Εὐλογημένος ὁ ἐρχόμενος ἐν ὀνόματι Κυρίου.\end{quote}
\end{poetryblock}

\ch{24}
Καὶ ἐξελθὼν ὁ Ἰησοῦς ἀπὸ τοῦ ἱεροῦ ἐπορεύετο, καὶ προσῆλθον οἱ μαθηταὶ αὐτοῦ ἐπιδεῖξαι αὐτῷ τὰς οἰκοδομὰς τοῦ ἱεροῦ.
\vs{2}ὁ δὲ ἀποκριθεὶς εἶπεν αὐτοῖς· Οὐ βλέπετε ταῦτα πάντα; ἀμὴν λέγω ὑμῖν, οὐ μὴ ἀφεθῇ ὧδε λίθος ἐπὶ λίθον ὃς οὐ καταλυθήσεται.

\vs{3}Καθημένου δὲ αὐτοῦ ἐπὶ τοῦ ὄρους τῶν Ἐλαιῶν προσῆλθον αὐτῷ οἱ μαθηταὶ κατ᾽ ἰδίαν λέγοντες· Εἰπὲ ἡμῖν, πότε ταῦτα ἔσται καὶ τί τὸ σημεῖον τῆς σῆς παρουσίας καὶ συντελείας τοῦ αἰῶνος;

\vs{4}Καὶ ἀποκριθεὶς ὁ Ἰησοῦς εἶπεν αὐτοῖς· Βλέπετε μή τις ὑμᾶς πλανήσῃ·
\vs{5}πολλοὶ γὰρ ἐλεύσονται ἐπὶ τῷ ὀνόματί μου λέγοντες· Ἐγώ εἰμι ὁ Χριστός, καὶ πολλοὺς πλανήσουσιν.
\vs{6}μελλήσετε δὲ ἀκούειν πολέμους καὶ ἀκοὰς πολέμων· ὁρᾶτε μὴ θροεῖσθε· δεῖ γὰρ γενέσθαι, ἀλλ᾽ οὔπω ἐστὶν τὸ τέλος.
\vs{7}ἐγερθήσεται γὰρ ἔθνος ἐπὶ ἔθνος καὶ βασιλεία ἐπὶ βασιλείαν καὶ ἔσονται λιμοὶ καὶ σεισμοὶ κατὰ τόπους·
\vs{8}πάντα δὲ ταῦτα ἀρχὴ ὠδίνων.

\vs{9}Τότε παραδώσουσιν ὑμᾶς εἰς θλῖψιν καὶ ἀποκτενοῦσιν ὑμᾶς, καὶ ἔσεσθε μισούμενοι ὑπὸ πάντων τῶν ἐθνῶν διὰ τὸ ὄνομά μου.
\vs{10}καὶ τότε σκανδαλισθήσονται πολλοὶ καὶ ἀλλήλους παραδώσουσιν καὶ μισήσουσιν ἀλλήλους·
\vs{11}καὶ πολλοὶ ψευδοπροφῆται ἐγερθήσονται καὶ πλανήσουσιν πολλούς·
\vs{12}Καὶ διὰ τὸ πληθυνθῆναι τὴν ἀνομίαν ψυγήσεται ἡ ἀγάπη τῶν πολλῶν.
\vs{13}ὁ δὲ ὑπομείνας εἰς τέλος οὗτος σωθήσεται.
\vs{14}Καὶ κηρυχθήσεται τοῦτο τὸ εὐαγγέλιον τῆς βασιλείας ἐν ὅλῃ τῇ οἰκουμένῃ εἰς μαρτύριον πᾶσιν τοῖς ἔθνεσιν, καὶ τότε ἥξει τὸ τέλος.

\vs{15}Ὅταν οὖν ἴδητε Τὸ βδέλυγμα τῆς ἐρημώσεως τὸ ῥηθὲν διὰ Δανιὴλ τοῦ προφήτου ἑστὸς ἐν τόπῳ ἁγίῳ, ὁ ἀναγινώσκων νοείτω,
\vs{16}τότε οἱ ἐν τῇ Ἰουδαίᾳ φευγέτωσαν εἰς τὰ ὄρη,
\vs{17}ὁ ἐπὶ τοῦ δώματος μὴ καταβάτω ἆραι τὰ ἐκ τῆς οἰκίας αὐτοῦ,
\vs{18}καὶ ὁ ἐν τῷ ἀγρῷ μὴ ἐπιστρεψάτω ὀπίσω ἆραι τὸ ἱμάτιον αὐτοῦ.
\vs{19}Οὐαὶ δὲ ταῖς ἐν γαστρὶ ἐχούσαις καὶ ταῖς θηλαζούσαις ἐν ἐκείναις ταῖς ἡμέραις.

\vs{20}προσεύχεσθε δὲ ἵνα μὴ γένηται ἡ φυγὴ ὑμῶν χειμῶνος μηδὲ σαββάτῳ.
\vs{21}ἔσται γὰρ τότε θλῖψις μεγάλη οἵα οὐ γέγονεν ἀπ᾽ ἀρχῆς κόσμου ἕως τοῦ νῦν οὐδ᾽ οὐ μὴ γένηται.
\vs{22}καὶ εἰ μὴ ἐκολοβώθησαν αἱ ἡμέραι ἐκεῖναι, οὐκ ἂν ἐσώθη πᾶσα σάρξ· διὰ δὲ τοὺς ἐκλεκτοὺς κολοβωθήσονται αἱ ἡμέραι ἐκεῖναι.

\vs{23}Τότε ἐάν τις ὑμῖν εἴπῃ· Ἰδοὺ ὧδε ὁ Χριστός, ἤ· Ὧδε, μὴ πιστεύσητε·
\vs{24}ἐγερθήσονται γὰρ ψευδόχριστοι καὶ ψευδοπροφῆται καὶ δώσουσιν σημεῖα μεγάλα καὶ τέρατα ὥστε πλανῆσαι, εἰ δυνατὸν, καὶ τοὺς ἐκλεκτούς.
\vs{25}ἰδοὺ προείρηκα ὑμῖν.
\vs{26}Ἐὰν οὖν εἴπωσιν ὑμῖν· Ἰδοὺ ἐν τῇ ἐρήμῳ ἐστίν, μὴ ἐξέλθητε· Ἰδοὺ ἐν τοῖς ταμείοις, μὴ πιστεύσητε·
\vs{27}ὥσπερ γὰρ ἡ ἀστραπὴ ἐξέρχεται ἀπὸ ἀνατολῶν καὶ φαίνεται ἕως δυσμῶν, οὕτως ἔσται ἡ παρουσία τοῦ Υἱοῦ τοῦ ἀνθρώπου·
\vs{28}ὅπου ἐὰν ᾖ τὸ πτῶμα, ἐκεῖ συναχθήσονται οἱ ἀετοί.

\vs{29}Εὐθέως δὲ μετὰ τὴν θλῖψιν τῶν ἡμερῶν ἐκείνων 
\begin{poetryblock}

\begin{quote}Ὁ ἥλιος σκοτισθήσεται,\end{quote} 

\begin{quote}Καὶ ἡ σελήνη οὐ δώσει τὸ φέγγος αὐτῆς,\end{quote} 

\begin{quote}Καὶ οἱ ἀστέρες πεσοῦνται ἀπὸ τοῦ οὐρανοῦ,\end{quote} 

\begin{quote}Καὶ αἱ δυνάμεις τῶν οὐρανῶν σαλευθήσονται.\end{quote}
\end{poetryblock}

\vs{30}Καὶ τότε φανήσεται τὸ σημεῖον τοῦ Υἱοῦ τοῦ ἀνθρώπου ἐν οὐρανῷ, καὶ τότε κόψονται πᾶσαι αἱ φυλαὶ τῆς γῆς καὶ ὄψονται τὸν Υἱὸν τοῦ ἀνθρώπου ἐρχόμενον ἐπὶ τῶν νεφελῶν τοῦ οὐρανοῦ μετὰ δυνάμεως καὶ δόξης πολλῆς·
\vs{31}καὶ ἀποστελεῖ τοὺς ἀγγέλους αὐτοῦ μετὰ σάλπιγγος μεγάλης, καὶ ἐπισυνάξουσιν τοὺς ἐκλεκτοὺς αὐτοῦ ἐκ τῶν τεσσάρων ἀνέμων ἀπ᾽ ἄκρων οὐρανῶν ἕως τῶν ἄκρων αὐτῶν.
\vs{32}Ἀπὸ δὲ τῆς συκῆς μάθετε τὴν παραβολήν· ὅταν ἤδη ὁ κλάδος αὐτῆς γένηται ἁπαλὸς καὶ τὰ φύλλα ἐκφύῃ, γινώσκετε ὅτι ἐγγὺς τὸ θέρος·
\vs{33}οὕτως καὶ ὑμεῖς, ὅταν ἴδητε πάντα ταῦτα, γινώσκετε ὅτι ἐγγύς ἐστιν ἐπὶ θύραις.
\vs{34}ἀμὴν λέγω ὑμῖν ὅτι οὐ μὴ παρέλθῃ ἡ γενεὰ αὕτη ἕως ἂν πάντα ταῦτα γένηται.
\vs{35}ὁ οὐρανὸς καὶ ἡ γῆ παρελεύσεται, οἱ δὲ λόγοι μου οὐ μὴ παρέλθωσιν.

\vs{36}Περὶ δὲ τῆς ἡμέρας ἐκείνης καὶ ὥρας οὐδεὶς οἶδεν, οὐδὲ οἱ ἄγγελοι τῶν οὐρανῶν οὐδὲ ὁ Υἱός, εἰ μὴ ὁ Πατὴρ μόνος.

\vs{37}ὥσπερ γὰρ αἱ ἡμέραι τοῦ Νῶε, οὕτως ἔσται ἡ παρουσία τοῦ Υἱοῦ τοῦ ἀνθρώπου.
\vs{38}ὡς γὰρ ἦσαν ἐν ταῖς ἡμέραις ἐκείναις ταῖς πρὸ τοῦ κατακλυσμοῦ τρώγοντες καὶ πίνοντες, γαμοῦντες καὶ γαμίζοντες, ἄχρι ἧς ἡμέρας εἰσῆλθεν Νῶε εἰς τὴν κιβωτόν,
\vs{39}καὶ οὐκ ἔγνωσαν ἕως ἦλθεν ὁ κατακλυσμὸς καὶ ἦρεν ἅπαντας, οὕτως ἔσται καὶ ἡ παρουσία τοῦ Υἱοῦ τοῦ ἀνθρώπου.
\vs{40}τότε δύο ἔσονται ἐν τῷ ἀγρῷ, εἷς παραλαμβάνεται καὶ εἷς ἀφίεται·
\vs{41}δύο ἀλήθουσαι ἐν τῷ μύλῳ, μία παραλαμβάνεται καὶ μία ἀφίεται.
\vs{42}Γρηγορεῖτε οὖν, ὅτι οὐκ οἴδατε ποίᾳ ἡμέρᾳ ὁ κύριος ὑμῶν ἔρχεται.
\vs{43}ἐκεῖνο δὲ γινώσκετε ὅτι εἰ ᾔδει ὁ οἰκοδεσπότης ποίᾳ φυλακῇ ὁ κλέπτης ἔρχεται, ἐγρηγόρησεν ἂν καὶ οὐκ ἂν εἴασεν διορυχθῆναι τὴν οἰκίαν αὐτοῦ.
\vs{44}διὰ τοῦτο καὶ ὑμεῖς γίνεσθε ἕτοιμοι, ὅτι ᾗ οὐ δοκεῖτε ὥρᾳ ὁ Υἱὸς τοῦ ἀνθρώπου ἔρχεται.

\vs{45}Τίς ἄρα ἐστὶν ὁ πιστὸς δοῦλος καὶ φρόνιμος ὃν κατέστησεν ὁ κύριος ἐπὶ τῆς οἰκετείας αὐτοῦ τοῦ δοῦναι αὐτοῖς τὴν τροφὴν ἐν καιρῷ;
\vs{46}μακάριος ὁ δοῦλος ἐκεῖνος ὃν ἐλθὼν ὁ κύριος αὐτοῦ εὑρήσει οὕτως ποιοῦντα·
\vs{47}ἀμὴν λέγω ὑμῖν ὅτι ἐπὶ πᾶσιν τοῖς ὑπάρχουσιν αὐτοῦ καταστήσει αὐτόν.
\vs{48}Ἐὰν δὲ εἴπῃ ὁ κακὸς δοῦλος ἐκεῖνος ἐν τῇ καρδίᾳ αὐτοῦ· Χρονίζει μου ὁ κύριος,
\vs{49}καὶ ἄρξηται τύπτειν τοὺς συνδούλους αὐτοῦ, ἐσθίῃ δὲ καὶ πίνῃ μετὰ τῶν μεθυόντων,
\vs{50}ἥξει ὁ κύριος τοῦ δούλου ἐκείνου ἐν ἡμέρᾳ ᾗ οὐ προσδοκᾷ καὶ ἐν ὥρᾳ ᾗ οὐ γινώσκει,
\vs{51}καὶ διχοτομήσει αὐτὸν καὶ τὸ μέρος αὐτοῦ μετὰ τῶν ὑποκριτῶν θήσει· ἐκεῖ ἔσται ὁ κλαυθμὸς καὶ ὁ βρυγμὸς τῶν ὀδόντων.

\ch{25}
Τότε ὁμοιωθήσεται ἡ βασιλεία τῶν οὐρανῶν δέκα παρθένοις, αἵτινες λαβοῦσαι τὰς λαμπάδας ἑαυτῶν ἐξῆλθον εἰς ὑπάντησιν τοῦ νυμφίου.
\vs{2}πέντε δὲ ἐξ αὐτῶν ἦσαν μωραὶ καὶ πέντε φρόνιμοι.
\vs{3}αἱ γὰρ μωραὶ λαβοῦσαι τὰς λαμπάδας αὐτῶν οὐκ ἔλαβον μεθ᾽ ἑαυτῶν ἔλαιον.
\vs{4}αἱ δὲ φρόνιμοι ἔλαβον ἔλαιον ἐν τοῖς ἀγγείοις μετὰ τῶν λαμπάδων ἑαυτῶν.
\vs{5}χρονίζοντος δὲ τοῦ νυμφίου ἐνύσταξαν πᾶσαι καὶ ἐκάθευδον.
\vs{6}Μέσης δὲ νυκτὸς κραυγὴ γέγονεν· Ἰδοὺ ὁ νυμφίος, ἐξέρχεσθε εἰς ἀπάντησιν αὐτοῦ.
\vs{7}Τότε ἠγέρθησαν πᾶσαι αἱ παρθένοι ἐκεῖναι καὶ ἐκόσμησαν τὰς λαμπάδας ἑαυτῶν.
\vs{8}αἱ δὲ μωραὶ ταῖς φρονίμοις εἶπαν· Δότε ἡμῖν ἐκ τοῦ ἐλαίου ὑμῶν, ὅτι αἱ λαμπάδες ἡμῶν σβέννυνται.
\vs{9}Ἀπεκρίθησαν δὲ αἱ φρόνιμοι λέγουσαι· Μήποτε οὐ μὴ ἀρκέσῃ ἡμῖν καὶ ὑμῖν· πορεύεσθε μᾶλλον πρὸς τοὺς πωλοῦντας καὶ ἀγοράσατε ἑαυταῖς.
\vs{10}Ἀπερχομένων δὲ αὐτῶν ἀγοράσαι ἦλθεν ὁ νυμφίος, καὶ αἱ ἕτοιμοι εἰσῆλθον μετ᾽ αὐτοῦ εἰς τοὺς γάμους καὶ ἐκλείσθη ἡ θύρα.
\vs{11}Ὕστερον δὲ ἔρχονται καὶ αἱ λοιπαὶ παρθένοι λέγουσαι· Κύριε κύριε, ἄνοιξον ἡμῖν.
\vs{12}Ὁ δὲ ἀποκριθεὶς εἶπεν· Ἀμὴν λέγω ὑμῖν, οὐκ οἶδα ὑμᾶς.
\vs{13}Γρηγορεῖτε οὖν, ὅτι οὐκ οἴδατε τὴν ἡμέραν οὐδὲ τὴν ὥραν.
\vs{14}Ὥσπερ γὰρ ἄνθρωπος ἀποδημῶν ἐκάλεσεν τοὺς ἰδίους δούλους καὶ παρέδωκεν αὐτοῖς τὰ ὑπάρχοντα αὐτοῦ,
\vs{15}καὶ ᾧ μὲν ἔδωκεν πέντε τάλαντα, ᾧ δὲ δύο, ᾧ δὲ ἕν, ἑκάστῳ κατὰ τὴν ἰδίαν δύναμιν, καὶ ἀπεδήμησεν. Εὐθέως
\vs{16}Πορευθεὶς ὁ τὰ πέντε τάλαντα λαβὼν ἠργάσατο ἐν αὐτοῖς καὶ ἐκέρδησεν ἄλλα πέντε·
\vs{17}ὡσαύτως ὁ τὰ δύο ἐκέρδησεν ἄλλα δύο.
\vs{18}ὁ δὲ τὸ ἓν λαβὼν ἀπελθὼν ὤρυξεν γῆν καὶ ἔκρυψεν τὸ ἀργύριον τοῦ κυρίου αὐτοῦ.
\vs{19}Μετὰ δὲ πολὺν χρόνον ἔρχεται ὁ κύριος τῶν δούλων ἐκείνων καὶ συναίρει λόγον μετ᾽ αὐτῶν.
\vs{20}καὶ προσελθὼν ὁ τὰ πέντε τάλαντα λαβὼν προσήνεγκεν ἄλλα πέντε τάλαντα λέγων· Κύριε, πέντε τάλαντά μοι παρέδωκας· ἴδε ἄλλα πέντε τάλαντα ἐκέρδησα.
\vs{21}Ἔφη αὐτῷ ὁ κύριος αὐτοῦ· Εὖ, δοῦλε ἀγαθὲ καὶ πιστέ, ἐπὶ ὀλίγα ἦς πιστός, ἐπὶ πολλῶν σε καταστήσω· εἴσελθε εἰς τὴν χαρὰν τοῦ κυρίου σου.
\vs{22}Προσελθὼν δὲ καὶ ὁ τὰ δύο τάλαντα εἶπεν· Κύριε, δύο τάλαντά μοι παρέδωκας· ἴδε ἄλλα δύο τάλαντα ἐκέρδησα.
\vs{23}Ἔφη αὐτῷ ὁ κύριος αὐτοῦ· Εὖ, δοῦλε ἀγαθὲ καὶ πιστέ, ἐπὶ ὀλίγα ἦς πιστός, ἐπὶ πολλῶν σε καταστήσω· εἴσελθε εἰς τὴν χαρὰν τοῦ κυρίου σου.
\vs{24}Προσελθὼν δὲ καὶ ὁ τὸ ἓν τάλαντον εἰληφὼς εἶπεν· Κύριε, ἔγνων σε ὅτι σκληρὸς εἶ ἄνθρωπος, θερίζων ὅπου οὐκ ἔσπειρας καὶ συνάγων ὅθεν οὐ διεσκόρπισας,
\vs{25}καὶ φοβηθεὶς ἀπελθὼν ἔκρυψα τὸ τάλαντόν σου ἐν τῇ γῇ· ἴδε ἔχεις τὸ σόν.
\vs{26}Ἀποκριθεὶς δὲ ὁ κύριος αὐτοῦ εἶπεν αὐτῷ· Πονηρὲ δοῦλε καὶ ὀκνηρέ, ᾔδεις ὅτι θερίζω ὅπου οὐκ ἔσπειρα καὶ συνάγω ὅθεν οὐ διεσκόρπισα;
\vs{27}ἔδει σε οὖν βαλεῖν τὰ ἀργύριά μου τοῖς τραπεζίταις, καὶ ἐλθὼν ἐγὼ ἐκομισάμην ἂν τὸ ἐμὸν σὺν τόκῳ.
\vs{28}Ἄρατε οὖν ἀπ᾽ αὐτοῦ τὸ τάλαντον καὶ δότε τῷ ἔχοντι τὰ δέκα τάλαντα·
\vs{29}τῷ γὰρ ἔχοντι παντὶ δοθήσεται καὶ περισσευθήσεται, τοῦ δὲ μὴ ἔχοντος καὶ ὃ ἔχει ἀρθήσεται ἀπ᾽ αὐτοῦ.
\vs{30}καὶ τὸν ἀχρεῖον δοῦλον ἐκβάλετε εἰς τὸ σκότος τὸ ἐξώτερον· ἐκεῖ ἔσται ὁ κλαυθμὸς καὶ ὁ βρυγμὸς τῶν ὀδόντων.

\vs{31}Ὅταν δὲ ἔλθῃ ὁ υἱὸς τοῦ ἀνθρώπου ἐν τῇ δόξῃ αὐτοῦ καὶ πάντες οἱ ἄγγελοι μετ᾽ αὐτοῦ, τότε καθίσει ἐπὶ θρόνου δόξης αὐτοῦ·
\vs{32}καὶ συναχθήσονται ἔμπροσθεν αὐτοῦ πάντα τὰ ἔθνη, καὶ ἀφορίσει αὐτοὺς ἀπ᾽ ἀλλήλων, ὥσπερ ὁ ποιμὴν ἀφορίζει τὰ πρόβατα ἀπὸ τῶν ἐρίφων,
\vs{33}καὶ στήσει τὰ μὲν πρόβατα ἐκ δεξιῶν αὐτοῦ, τὰ δὲ ἐρίφια ἐξ εὐωνύμων.
\vs{34}Τότε ἐρεῖ ὁ Βασιλεὺς τοῖς ἐκ δεξιῶν αὐτοῦ· Δεῦτε οἱ εὐλογημένοι τοῦ Πατρός μου, κληρονομήσατε τὴν ἡτοιμασμένην ὑμῖν βασιλείαν ἀπὸ καταβολῆς κόσμου.
\vs{35}ἐπείνασα γὰρ καὶ ἐδώκατέ μοι φαγεῖν, ἐδίψησα καὶ ἐποτίσατέ με, ξένος ἤμην καὶ συνηγάγετέ με,
\vs{36}γυμνὸς καὶ περιεβάλετέ με, ἠσθένησα καὶ ἐπεσκέψασθέ με, ἐν φυλακῇ ἤμην καὶ ἤλθατε πρός με.
\vs{37}Τότε ἀποκριθήσονται αὐτῷ οἱ δίκαιοι λέγοντες· Κύριε, πότε σε εἴδομεν πεινῶντα καὶ ἐθρέψαμεν, ἢ διψῶντα καὶ ἐποτίσαμεν;
\vs{38}πότε δέ σε εἴδομεν ξένον καὶ συνηγάγομεν, ἢ γυμνὸν καὶ περιεβάλομεν;
\vs{39}πότε δέ σε εἴδομεν ἀσθενοῦντα ἢ ἐν φυλακῇ καὶ ἤλθομεν πρός σε;
\vs{40}Καὶ ἀποκριθεὶς ὁ Βασιλεὺς ἐρεῖ αὐτοῖς· Ἀμὴν λέγω ὑμῖν, ἐφ᾽ ὅσον ἐποιήσατε ἑνὶ τούτων τῶν ἀδελφῶν μου τῶν ἐλαχίστων, ἐμοὶ ἐποιήσατε.
\vs{41}Τότε ἐρεῖ καὶ τοῖς ἐξ εὐωνύμων· Πορεύεσθε ἀπ᾽ ἐμοῦ οἱ κατηραμένοι εἰς τὸ πῦρ τὸ αἰώνιον τὸ ἡτοιμασμένον τῷ διαβόλῳ καὶ τοῖς ἀγγέλοις αὐτοῦ.
\vs{42}ἐπείνασα γὰρ καὶ οὐκ ἐδώκατέ μοι φαγεῖν, ἐδίψησα καὶ οὐκ ἐποτίσατέ με,
\vs{43}ξένος ἤμην καὶ οὐ συνηγάγετέ με, γυμνὸς καὶ οὐ περιεβάλετέ με, ἀσθενὴς καὶ ἐν φυλακῇ καὶ οὐκ ἐπεσκέψασθέ με.
\vs{44}Τότε ἀποκριθήσονται καὶ αὐτοὶ λέγοντες· Κύριε, πότε σε εἴδομεν πεινῶντα ἢ διψῶντα ἢ ξένον ἢ γυμνὸν ἢ ἀσθενῆ ἢ ἐν φυλακῇ καὶ οὐ διηκονήσαμέν σοι;
\vs{45}Τότε ἀποκριθήσεται αὐτοῖς λέγων· Ἀμὴν λέγω ὑμῖν, ἐφ᾽ ὅσον οὐκ ἐποιήσατε ἑνὶ τούτων τῶν ἐλαχίστων, οὐδὲ ἐμοὶ ἐποιήσατε.
\vs{46}Καὶ ἀπελεύσονται οὗτοι εἰς κόλασιν αἰώνιον, οἱ δὲ δίκαιοι εἰς ζωὴν αἰώνιον.

\ch{26}
Καὶ ἐγένετο ὅτε ἐτέλεσεν ὁ Ἰησοῦς πάντας τοὺς λόγους τούτους, εἶπεν τοῖς μαθηταῖς αὐτοῦ·
\vs{2}Οἴδατε ὅτι μετὰ δύο ἡμέρας τὸ πάσχα γίνεται, καὶ ὁ Υἱὸς τοῦ ἀνθρώπου παραδίδοται εἰς τὸ σταυρωθῆναι.

\vs{3}Τότε συνήχθησαν οἱ ἀρχιερεῖς καὶ οἱ πρεσβύτεροι τοῦ λαοῦ εἰς τὴν αὐλὴν τοῦ ἀρχιερέως τοῦ λεγομένου Καϊάφα
\vs{4}καὶ συνεβουλεύσαντο ἵνα τὸν Ἰησοῦν δόλῳ κρατήσωσιν καὶ ἀποκτείνωσιν·
\vs{5}ἔλεγον δέ· Μὴ ἐν τῇ ἑορτῇ, ἵνα μὴ θόρυβος γένηται ἐν τῷ λαῷ.

\vs{6}Τοῦ δὲ Ἰησοῦ γενομένου ἐν Βηθανίᾳ ἐν οἰκίᾳ Σίμωνος τοῦ λεπροῦ,
\vs{7}προσῆλθεν αὐτῷ γυνὴ ἔχουσα ἀλάβαστρον μύρου βαρυτίμου καὶ κατέχεεν ἐπὶ τῆς κεφαλῆς αὐτοῦ ἀνακειμένου.
\vs{8}Ἰδόντες δὲ οἱ μαθηταὶ ἠγανάκτησαν λέγοντες· Εἰς τί ἡ ἀπώλεια αὕτη;
\vs{9}ἐδύνατο γὰρ τοῦτο πραθῆναι πολλοῦ καὶ δοθῆναι πτωχοῖς.
\vs{10}Γνοὺς δὲ ὁ Ἰησοῦς εἶπεν αὐτοῖς· Τί κόπους παρέχετε τῇ γυναικί; ἔργον γὰρ καλὸν ἠργάσατο εἰς ἐμέ·
\vs{11}πάντοτε γὰρ τοὺς πτωχοὺς ἔχετε μεθ᾽ ἑαυτῶν, ἐμὲ δὲ οὐ πάντοτε ἔχετε·
\vs{12}βαλοῦσα γὰρ αὕτη τὸ μύρον τοῦτο ἐπὶ τοῦ σώματός μου πρὸς τὸ ἐνταφιάσαι με ἐποίησεν.
\vs{13}ἀμὴν λέγω ὑμῖν, ὅπου ἐὰν κηρυχθῇ τὸ εὐαγγέλιον τοῦτο ἐν ὅλῳ τῷ κόσμῳ, λαληθήσεται καὶ ὃ ἐποίησεν αὕτη εἰς μνημόσυνον αὐτῆς.

\vs{14}Τότε πορευθεὶς εἷς τῶν δώδεκα, ὁ λεγόμενος Ἰούδας Ἰσκαριώτης, πρὸς τοὺς ἀρχιερεῖς
\vs{15}εἶπεν· Τί θέλετέ μοι δοῦναι, κἀγὼ ὑμῖν παραδώσω αὐτόν; οἱ δὲ ἔστησαν αὐτῷ τριάκοντα ἀργύρια.
\vs{16}καὶ ἀπὸ τότε ἐζήτει εὐκαιρίαν ἵνα αὐτὸν παραδῷ.
\vs{17}Τῇ δὲ πρώτῃ τῶν ἀζύμων προσῆλθον οἱ μαθηταὶ τῷ Ἰησοῦ λέγοντες· Ποῦ θέλεις ἑτοιμάσωμέν σοι φαγεῖν τὸ πάσχα;
\vs{18}Ὁ δὲ εἶπεν· Ὑπάγετε εἰς τὴν πόλιν πρὸς τὸν δεῖνα καὶ εἴπατε αὐτῷ· Ὁ Διδάσκαλος λέγει· Ὁ καιρός μου ἐγγύς ἐστιν, πρὸς σὲ ποιῶ τὸ πάσχα μετὰ τῶν μαθητῶν μου.
\vs{19}καὶ ἐποίησαν οἱ μαθηταὶ ὡς συνέταξεν αὐτοῖς ὁ Ἰησοῦς καὶ ἡτοίμασαν τὸ πάσχα.

\vs{20}Ὀψίας δὲ γενομένης ἀνέκειτο μετὰ τῶν δώδεκα.
\vs{21}καὶ ἐσθιόντων αὐτῶν εἶπεν· Ἀμὴν λέγω ὑμῖν ὅτι εἷς ἐξ ὑμῶν παραδώσει με.
\vs{22}Καὶ λυπούμενοι σφόδρα ἤρξαντο λέγειν αὐτῷ εἷς ἕκαστος· Μήτι ἐγώ εἰμι, Κύριε;
\vs{23}Ὁ δὲ ἀποκριθεὶς εἶπεν· Ὁ ἐμβάψας μετ᾽ ἐμοῦ τὴν χεῖρα ἐν τῷ τρυβλίῳ οὗτός με παραδώσει.
\vs{24}ὁ μὲν Υἱὸς τοῦ ἀνθρώπου ὑπάγει καθὼς γέγραπται περὶ αὐτοῦ, οὐαὶ δὲ τῷ ἀνθρώπῳ ἐκείνῳ δι᾽ οὗ ὁ Υἱὸς τοῦ ἀνθρώπου παραδίδοται· καλὸν ἦν αὐτῷ εἰ οὐκ ἐγεννήθη ὁ ἄνθρωπος ἐκεῖνος.
\vs{25}Ἀποκριθεὶς δὲ Ἰούδας ὁ παραδιδοὺς αὐτὸν εἶπεν· Μήτι ἐγώ εἰμι, ῥαββί; Λέγει αὐτῷ· Σὺ εἶπας.

\vs{26}Ἐσθιόντων δὲ αὐτῶν λαβὼν ὁ Ἰησοῦς ἄρτον καὶ εὐλογήσας ἔκλασεν καὶ δοὺς τοῖς μαθηταῖς εἶπεν· Λάβετε φάγετε, τοῦτό ἐστιν τὸ σῶμά μου.
\vs{27}Καὶ λαβὼν ποτήριον καὶ εὐχαριστήσας ἔδωκεν αὐτοῖς λέγων· Πίετε ἐξ αὐτοῦ πάντες,
\vs{28}τοῦτο γάρ ἐστιν τὸ αἷμά μου τῆς διαθήκης τὸ περὶ πολλῶν ἐκχυννόμενον εἰς ἄφεσιν ἁμαρτιῶν.
\vs{29}λέγω δὲ ὑμῖν, οὐ μὴ πίω ἀπ᾽ ἄρτι ἐκ τούτου τοῦ γενήματος τῆς ἀμπέλου ἕως τῆς ἡμέρας ἐκείνης ὅταν αὐτὸ πίνω μεθ᾽ ὑμῶν καινὸν ἐν τῇ βασιλείᾳ τοῦ Πατρός μου.
\vs{30}Καὶ ὑμνήσαντες ἐξῆλθον εἰς τὸ ὄρος τῶν Ἐλαιῶν.

\vs{31}Τότε λέγει αὐτοῖς ὁ Ἰησοῦς· Πάντες ὑμεῖς σκανδαλισθήσεσθε ἐν ἐμοὶ ἐν τῇ νυκτὶ ταύτῃ, γέγραπται γάρ· 
\begin{poetryblock}

\begin{quote}Πατάξω τὸν ποιμένα,\end{quote} 

\begin{quote}Καὶ διασκορπισθήσονται τὰ πρόβατα τῆς ποίμνης.\end{quote}
\end{poetryblock}

\vs{32}Μετὰ δὲ τὸ ἐγερθῆναί με προάξω ὑμᾶς εἰς τὴν Γαλιλαίαν.
\vs{33}Ἀποκριθεὶς δὲ ὁ Πέτρος εἶπεν αὐτῷ· Εἰ πάντες σκανδαλισθήσονται ἐν σοί, ἐγὼ οὐδέποτε σκανδαλισθήσομαι.
\vs{34}Ἔφη αὐτῷ ὁ Ἰησοῦς· Ἀμὴν λέγω σοι ὅτι ἐν ταύτῃ τῇ νυκτὶ πρὶν ἀλέκτορα φωνῆσαι τρὶς ἀπαρνήσῃ με.
\vs{35}Λέγει αὐτῷ ὁ Πέτρος· Κἂν δέῃ με σὺν σοὶ ἀποθανεῖν, οὐ μή σε ἀπαρνήσομαι. ὁμοίως καὶ πάντες οἱ μαθηταὶ εἶπαν.

\vs{36}Τότε ἔρχεται μετ᾽ αὐτῶν ὁ Ἰησοῦς εἰς χωρίον λεγόμενον Γεθσημανὶ καὶ λέγει τοῖς μαθηταῖς· Καθίσατε αὐτοῦ ἕως οὗ ἀπελθὼν ἐκεῖ προσεύξωμαι.
\vs{37}Καὶ παραλαβὼν τὸν Πέτρον καὶ τοὺς δύο υἱοὺς Ζεβεδαίου ἤρξατο λυπεῖσθαι καὶ ἀδημονεῖν.
\vs{38}τότε λέγει αὐτοῖς· Περίλυπός ἐστιν ἡ ψυχή μου ἕως θανάτου· μείνατε ὧδε καὶ γρηγορεῖτε μετ᾽ ἐμοῦ.
\vs{39}Καὶ προελθὼν μικρὸν ἔπεσεν ἐπὶ πρόσωπον αὐτοῦ προσευχόμενος καὶ λέγων· Πάτερ μου, εἰ δυνατόν ἐστιν, παρελθάτω ἀπ᾽ ἐμοῦ τὸ ποτήριον τοῦτο· πλὴν οὐχ ὡς ἐγὼ θέλω ἀλλ᾽ ὡς σύ.
\vs{40}Καὶ ἔρχεται πρὸς τοὺς μαθητὰς καὶ εὑρίσκει αὐτοὺς καθεύδοντας, καὶ λέγει τῷ Πέτρῳ· Οὕτως οὐκ ἰσχύσατε μίαν ὥραν γρηγορῆσαι μετ᾽ ἐμοῦ;
\vs{41}γρηγορεῖτε καὶ προσεύχεσθε, ἵνα μὴ εἰσέλθητε εἰς πειρασμόν· τὸ μὲν πνεῦμα πρόθυμον ἡ δὲ σὰρξ ἀσθενής.
\vs{42}Πάλιν ἐκ δευτέρου ἀπελθὼν προσηύξατο λέγων· Πάτερ μου, εἰ οὐ δύναται τοῦτο παρελθεῖν ἐὰν μὴ αὐτὸ πίω, γενηθήτω τὸ θέλημά σου.
\vs{43}καὶ ἐλθὼν πάλιν εὗρεν αὐτοὺς καθεύδοντας, ἦσαν γὰρ αὐτῶν οἱ ὀφθαλμοὶ βεβαρημένοι.
\vs{44}Καὶ ἀφεὶς αὐτοὺς πάλιν ἀπελθὼν προσηύξατο ἐκ τρίτου τὸν αὐτὸν λόγον εἰπὼν πάλιν.
\vs{45}τότε ἔρχεται πρὸς τοὺς μαθητὰς καὶ λέγει αὐτοῖς· Καθεύδετε τὸ λοιπὸν καὶ ἀναπαύεσθε· ἰδοὺ ἤγγικεν ἡ ὥρα καὶ ὁ Υἱὸς τοῦ ἀνθρώπου παραδίδοται εἰς χεῖρας ἁμαρτωλῶν.
\vs{46}ἐγείρεσθε ἄγωμεν· ἰδοὺ ἤγγικεν ὁ παραδιδούς με.

\vs{47}Καὶ ἔτι αὐτοῦ λαλοῦντος ἰδοὺ Ἰούδας εἷς τῶν δώδεκα ἦλθεν καὶ μετ᾽ αὐτοῦ ὄχλος πολὺς μετὰ μαχαιρῶν καὶ ξύλων ἀπὸ τῶν ἀρχιερέων καὶ πρεσβυτέρων τοῦ λαοῦ.
\vs{48}Ὁ δὲ παραδιδοὺς αὐτὸν ἔδωκεν αὐτοῖς σημεῖον λέγων· Ὃν ἂν φιλήσω αὐτός ἐστιν, κρατήσατε αὐτόν.
\vs{49}καὶ εὐθέως προσελθὼν τῷ Ἰησοῦ εἶπεν· Χαῖρε, ῥαββί, καὶ κατεφίλησεν αὐτόν.
\vs{50}Ὁ δὲ Ἰησοῦς εἶπεν αὐτῷ· Ἑταῖρε, ἐφ᾽ ὃ πάρει. Τότε προσελθόντες ἐπέβαλον τὰς χεῖρας ἐπὶ τὸν Ἰησοῦν καὶ ἐκράτησαν αὐτόν.
\vs{51}καὶ ἰδοὺ εἷς τῶν μετὰ Ἰησοῦ ἐκτείνας τὴν χεῖρα ἀπέσπασεν τὴν μάχαιραν αὐτοῦ καὶ πατάξας τὸν δοῦλον τοῦ ἀρχιερέως ἀφεῖλεν αὐτοῦ τὸ ὠτίον.
\vs{52}Τότε λέγει αὐτῷ ὁ Ἰησοῦς· Ἀπόστρεψον τὴν μάχαιράν σου εἰς τὸν τόπον αὐτῆς· πάντες γὰρ οἱ λαβόντες μάχαιραν ἐν μαχαίρῃ ἀπολοῦνται.
\vs{53}ἢ δοκεῖς ὅτι οὐ δύναμαι παρακαλέσαι τὸν Πατέρα μου, καὶ παραστήσει μοι ἄρτι πλείω δώδεκα λεγιῶνας ἀγγέλων;
\vs{54}πῶς οὖν πληρωθῶσιν αἱ γραφαὶ ὅτι οὕτως δεῖ γενέσθαι;
\vs{55}Ἐν ἐκείνῃ τῇ ὥρᾳ εἶπεν ὁ Ἰησοῦς τοῖς ὄχλοις· Ὡς ἐπὶ λῃστὴν ἐξήλθατε μετὰ μαχαιρῶν καὶ ξύλων συλλαβεῖν με; καθ᾽ ἡμέραν ἐν τῷ ἱερῷ ἐκαθεζόμην διδάσκων καὶ οὐκ ἐκρατήσατέ με.
\vs{56}Τοῦτο δὲ ὅλον γέγονεν ἵνα πληρωθῶσιν αἱ γραφαὶ τῶν προφητῶν. Τότε οἱ μαθηταὶ πάντες ἀφέντες αὐτὸν ἔφυγον.

\vs{57}Οἱ δὲ κρατήσαντες τὸν Ἰησοῦν ἀπήγαγον πρὸς Καϊάφαν τὸν ἀρχιερέα, ὅπου οἱ γραμματεῖς καὶ οἱ πρεσβύτεροι συνήχθησαν.
\vs{58}ὁ δὲ Πέτρος ἠκολούθει αὐτῷ ἀπὸ μακρόθεν ἕως τῆς αὐλῆς τοῦ ἀρχιερέως καὶ εἰσελθὼν ἔσω ἐκάθητο μετὰ τῶν ὑπηρετῶν ἰδεῖν τὸ τέλος.

\vs{59}Οἱ δὲ ἀρχιερεῖς καὶ τὸ συνέδριον ὅλον ἐζήτουν ψευδομαρτυρίαν κατὰ τοῦ Ἰησοῦ ὅπως αὐτὸν θανατώσωσιν,
\vs{60}καὶ οὐχ εὗρον πολλῶν προσελθόντων ψευδομαρτύρων. Ὕστερον δὲ προσελθόντες δύο
\vs{61}εἶπαν· Οὗτος ἔφη· Δύναμαι καταλῦσαι τὸν ναὸν τοῦ Θεοῦ καὶ διὰ τριῶν ἡμερῶν οἰκοδομῆσαι.
\vs{62}Καὶ ἀναστὰς ὁ ἀρχιερεὺς εἶπεν αὐτῷ· Οὐδὲν ἀποκρίνῃ τί οὗτοί σου καταμαρτυροῦσιν;
\vs{63}Ὁ δὲ Ἰησοῦς ἐσιώπα. Καὶ ὁ ἀρχιερεὺς εἶπεν αὐτῷ· Ἐξορκίζω σε κατὰ τοῦ Θεοῦ τοῦ ζῶντος ἵνα ἡμῖν εἴπῃς εἰ σὺ εἶ ὁ Χριστὸς ὁ Υἱὸς τοῦ Θεοῦ.
\vs{64}Λέγει αὐτῷ ὁ Ἰησοῦς· Σὺ εἶπας. πλὴν λέγω ὑμῖν· ἀπ᾽ ἄρτι ὄψεσθε τὸν Υἱὸν τοῦ ἀνθρώπου καθήμενον ἐκ δεξιῶν τῆς δυνάμεως καὶ ἐρχόμενον ἐπὶ τῶν νεφελῶν τοῦ οὐρανοῦ.
\vs{65}Τότε ὁ ἀρχιερεὺς διέρρηξεν τὰ ἱμάτια αὐτοῦ λέγων· Ἐβλασφήμησεν· τί ἔτι χρείαν ἔχομεν μαρτύρων; ἴδε νῦν ἠκούσατε τὴν βλασφημίαν·
\vs{66}τί ὑμῖν δοκεῖ; Οἱ δὲ ἀποκριθέντες εἶπαν· Ἔνοχος θανάτου ἐστίν.

\vs{67}Τότε ἐνέπτυσαν εἰς τὸ πρόσωπον αὐτοῦ καὶ ἐκολάφισαν αὐτόν, οἱ δὲ ἐράπισαν
\vs{68}λέγοντες· Προφήτευσον ἡμῖν, Χριστέ, τίς ἐστιν ὁ παίσας σε;

\vs{69}Ὁ δὲ Πέτρος ἐκάθητο ἔξω ἐν τῇ αὐλῇ· καὶ προσῆλθεν αὐτῷ μία παιδίσκη λέγουσα· Καὶ σὺ ἦσθα μετὰ Ἰησοῦ τοῦ Γαλιλαίου.
\vs{70}Ὁ δὲ ἠρνήσατο ἔμπροσθεν πάντων λέγων· Οὐκ οἶδα τί λέγεις.
\vs{71}Ἐξελθόντα δὲ εἰς τὸν πυλῶνα εἶδεν αὐτὸν ἄλλη καὶ λέγει τοῖς ἐκεῖ· Οὗτος ἦν μετὰ Ἰησοῦ τοῦ Ναζωραίου.
\vs{72}Καὶ πάλιν ἠρνήσατο μετὰ ὅρκου ὅτι Οὐκ οἶδα τὸν ἄνθρωπον.
\vs{73}Μετὰ μικρὸν δὲ προσελθόντες οἱ ἑστῶτες εἶπον τῷ Πέτρῳ· Ἀληθῶς καὶ σὺ ἐξ αὐτῶν εἶ, καὶ γὰρ ἡ λαλιά σου δῆλόν σε ποιεῖ.
\vs{74}Τότε ἤρξατο καταθεματίζειν καὶ ὀμνύειν ὅτι Οὐκ οἶδα τὸν ἄνθρωπον. Καὶ εὐθέως ἀλέκτωρ ἐφώνησεν.

\vs{75}Καὶ ἐμνήσθη ὁ Πέτρος τοῦ ῥήματος Ἰησοῦ εἰρηκότος ὅτι Πρὶν ἀλέκτορα φωνῆσαι τρὶς ἀπαρνήσῃ με· καὶ ἐξελθὼν ἔξω ἔκλαυσεν πικρῶς.

\ch{27}
Πρωΐας δὲ γενομένης συμβούλιον ἔλαβον πάντες οἱ ἀρχιερεῖς καὶ οἱ πρεσβύτεροι τοῦ λαοῦ κατὰ τοῦ Ἰησοῦ ὥστε θανατῶσαι αὐτόν·
\vs{2}καὶ δήσαντες αὐτὸν ἀπήγαγον καὶ παρέδωκαν Πιλάτῳ τῷ ἡγεμόνι.

\vs{3}Τότε ἰδὼν Ἰούδας ὁ παραδιδοὺς αὐτὸν ὅτι κατεκρίθη, μεταμεληθεὶς ἔστρεψεν τὰ τριάκοντα ἀργύρια τοῖς ἀρχιερεῦσιν καὶ πρεσβυτέροις
\vs{4}λέγων· Ἥμαρτον παραδοὺς αἷμα ἀθῷον. Οἱ δὲ εἶπαν· Τί πρὸς ἡμᾶς; σὺ ὄψῃ.
\vs{5}Καὶ ῥίψας τὰ ἀργύρια εἰς τὸν ναὸν ἀνεχώρησεν, καὶ ἀπελθὼν ἀπήγξατο.
\vs{6}Οἱ δὲ ἀρχιερεῖς λαβόντες τὰ ἀργύρια εἶπαν· Οὐκ ἔξεστιν βαλεῖν αὐτὰ εἰς τὸν κορβανᾶν, ἐπεὶ τιμὴ αἵματός ἐστιν.
\vs{7}συμβούλιον δὲ λαβόντες ἠγόρασαν ἐξ αὐτῶν τὸν ἀγρὸν τοῦ κεραμέως εἰς ταφὴν τοῖς ξένοις.
\vs{8}διὸ ἐκλήθη ὁ ἀγρὸς ἐκεῖνος Ἀγρὸς αἵματος ἕως τῆς σήμερον.
\vs{9}τότε ἐπληρώθη τὸ ῥηθὲν διὰ Ἰερεμίου τοῦ προφήτου λέγοντος· Καὶ ἔλαβον τὰ τριάκοντα ἀργύρια, τὴν τιμὴν τοῦ τετιμημένου ὃν ἐτιμήσαντο ἀπὸ υἱῶν Ἰσραήλ,
\vs{10}καὶ ἔδωκαν αὐτὰ εἰς τὸν ἀγρὸν τοῦ κεραμέως, καθὰ συνέταξέν μοι Κύριος.

\vs{11}Ὁ δὲ Ἰησοῦς ἐστάθη ἔμπροσθεν τοῦ ἡγεμόνος· καὶ ἐπηρώτησεν αὐτὸν ὁ ἡγεμὼν λέγων· Σὺ εἶ ὁ Βασιλεὺς τῶν Ἰουδαίων; Ὁ δὲ Ἰησοῦς ἔφη· Σὺ λέγεις.
\vs{12}Καὶ ἐν τῷ κατηγορεῖσθαι αὐτὸν ὑπὸ τῶν ἀρχιερέων καὶ πρεσβυτέρων οὐδὲν ἀπεκρίνατο.
\vs{13}Τότε λέγει αὐτῷ ὁ Πιλᾶτος· Οὐκ ἀκούεις πόσα σου καταμαρτυροῦσιν;
\vs{14}Καὶ οὐκ ἀπεκρίθη αὐτῷ πρὸς οὐδὲ ἓν ῥῆμα, ὥστε θαυμάζειν τὸν ἡγεμόνα λίαν.

\vs{15}Κατὰ δὲ ἑορτὴν εἰώθει ὁ ἡγεμὼν ἀπολύειν ἕνα τῷ ὄχλῳ δέσμιον ὃν ἤθελον.
\vs{16}εἶχον δὲ τότε δέσμιον ἐπίσημον λεγόμενον Ἰησοῦν Βαραββᾶν.
\vs{17}συνηγμένων οὖν αὐτῶν εἶπεν αὐτοῖς ὁ Πιλᾶτος· Τίνα θέλετε ἀπολύσω ὑμῖν, Ἰησοῦν τὸν Βαραββᾶν ἢ Ἰησοῦν τὸν λεγόμενον Χριστόν;
\vs{18}ᾔδει γὰρ ὅτι διὰ φθόνον παρέδωκαν αὐτόν.

\vs{19}Καθημένου δὲ αὐτοῦ ἐπὶ τοῦ βήματος ἀπέστειλεν πρὸς αὐτὸν ἡ γυνὴ αὐτοῦ λέγουσα· Μηδὲν σοὶ καὶ τῷ δικαίῳ ἐκείνῳ· πολλὰ γὰρ ἔπαθον σήμερον κατ᾽ ὄναρ δι᾽ αὐτόν.

\vs{20}Οἱ δὲ ἀρχιερεῖς καὶ οἱ πρεσβύτεροι ἔπεισαν τοὺς ὄχλους ἵνα αἰτήσωνται τὸν Βαραββᾶν, τὸν δὲ Ἰησοῦν ἀπολέσωσιν.
\vs{21}Ἀποκριθεὶς δὲ ὁ ἡγεμὼν εἶπεν αὐτοῖς· Τίνα θέλετε ἀπὸ τῶν δύο ἀπολύσω ὑμῖν; Οἱ δὲ εἶπαν· Τὸν Βαραββᾶν.
\vs{22}Λέγει αὐτοῖς ὁ Πιλᾶτος· Τί οὖν ποιήσω Ἰησοῦν τὸν λεγόμενον Χριστόν; Λέγουσιν πάντες· Σταυρωθήτω.
\vs{23}Ὁ δὲ ἔφη· Τί γὰρ κακὸν ἐποίησεν; Οἱ δὲ περισσῶς ἔκραζον λέγοντες· Σταυρωθήτω.

\vs{24}Ἰδὼν δὲ ὁ Πιλᾶτος ὅτι οὐδὲν ὠφελεῖ ἀλλὰ μᾶλλον θόρυβος γίνεται, λαβὼν ὕδωρ ἀπενίψατο τὰς χεῖρας ἀπέναντι τοῦ ὄχλου λέγων· Ἀθῷός εἰμι ἀπὸ τοῦ αἵματος τούτου· ὑμεῖς ὄψεσθε.
\vs{25}Καὶ ἀποκριθεὶς πᾶς ὁ λαὸς εἶπεν· Τὸ αἷμα αὐτοῦ ἐφ᾽ ἡμᾶς καὶ ἐπὶ τὰ τέκνα ἡμῶν.
\vs{26}Τότε ἀπέλυσεν αὐτοῖς τὸν Βαραββᾶν, τὸν δὲ Ἰησοῦν φραγελλώσας παρέδωκεν ἵνα σταυρωθῇ.

\vs{27}Τότε οἱ στρατιῶται τοῦ ἡγεμόνος παραλαβόντες τὸν Ἰησοῦν εἰς τὸ πραιτώριον συνήγαγον ἐπ᾽ αὐτὸν ὅλην τὴν σπεῖραν.
\vs{28}καὶ ἐκδύσαντες αὐτὸν χλαμύδα κοκκίνην περιέθηκαν αὐτῷ,
\vs{29}καὶ πλέξαντες στέφανον ἐξ ἀκανθῶν ἐπέθηκαν ἐπὶ τῆς κεφαλῆς αὐτοῦ καὶ κάλαμον ἐν τῇ δεξιᾷ αὐτοῦ, καὶ γονυπετήσαντες ἔμπροσθεν αὐτοῦ ἐνέπαιξαν αὐτῷ λέγοντες· Χαῖρε, Βασιλεῦ τῶν Ἰουδαίων,
\vs{30}καὶ ἐμπτύσαντες εἰς αὐτὸν ἔλαβον τὸν κάλαμον καὶ ἔτυπτον εἰς τὴν κεφαλὴν αὐτοῦ.
\vs{31}Καὶ ὅτε ἐνέπαιξαν αὐτῷ, ἐξέδυσαν αὐτὸν τὴν χλαμύδα καὶ ἐνέδυσαν αὐτὸν τὰ ἱμάτια αὐτοῦ καὶ ἀπήγαγον αὐτὸν εἰς τὸ σταυρῶσαι.
\vs{32}Ἐξερχόμενοι δὲ εὗρον ἄνθρωπον Κυρηναῖον ὀνόματι Σίμωνα, τοῦτον ἠγγάρευσαν ἵνα ἄρῃ τὸν σταυρὸν αὐτοῦ.

\vs{33}Καὶ ἐλθόντες εἰς τόπον λεγόμενον Γολγοθᾶ, ὅ ἐστιν κρανίου τόπος λεγόμενος,
\vs{34}ἔδωκαν αὐτῷ πιεῖν οἶνον μετὰ χολῆς μεμιγμένον· καὶ γευσάμενος οὐκ ἠθέλησεν πιεῖν.
\vs{35}Σταυρώσαντες δὲ αὐτὸν διεμερίσαντο τὰ ἱμάτια αὐτοῦ βάλλοντες κλῆρον,
\vs{36}καὶ καθήμενοι ἐτήρουν αὐτὸν ἐκεῖ.
\vs{37}καὶ ἐπέθηκαν ἐπάνω τῆς κεφαλῆς αὐτοῦ τὴν αἰτίαν αὐτοῦ γεγραμμένην· 
\begin{poetryblock}

\begin{quote}ΟΥΤΟΣ ΕΣΤΙΝ ΙΗΣΟΥΣ Ο ΒΑΣΙΛΕΥΣ ΤΩΝ ΙΟΥΔΑΙΩΝ.\end{quote}
\end{poetryblock}

\vs{38}Τότε σταυροῦνται σὺν αὐτῷ δύο λῃσταί, εἷς ἐκ δεξιῶν καὶ εἷς ἐξ εὐωνύμων.
\vs{39}Οἱ δὲ παραπορευόμενοι ἐβλασφήμουν αὐτὸν κινοῦντες τὰς κεφαλὰς αὐτῶν
\vs{40}καὶ λέγοντες· Ὁ καταλύων τὸν ναὸν καὶ ἐν τρισὶν ἡμέραις οἰκοδομῶν, σῶσον σεαυτόν, εἰ Υἱὸς εἶ τοῦ Θεοῦ, καὶ κατάβηθι ἀπὸ τοῦ σταυροῦ.
\vs{41}Ὁμοίως καὶ οἱ ἀρχιερεῖς ἐμπαίζοντες μετὰ τῶν γραμματέων καὶ πρεσβυτέρων ἔλεγον·
\vs{42}Ἄλλους ἔσωσεν, ἑαυτὸν οὐ δύναται σῶσαι· Βασιλεὺς Ἰσραήλ ἐστιν, καταβάτω νῦν ἀπὸ τοῦ σταυροῦ καὶ πιστεύσομεν ἐπ᾽ αὐτόν.
\vs{43}πέποιθεν ἐπὶ τὸν Θεόν, ῥυσάσθω νῦν εἰ θέλει αὐτόν· εἶπεν γὰρ ὅτι Θεοῦ εἰμι Υἱός.
\vs{44}Τὸ δ᾽ αὐτὸ καὶ οἱ λῃσταὶ οἱ συσταυρωθέντες σὺν αὐτῷ ὠνείδιζον αὐτόν.

\vs{45}Ἀπὸ δὲ ἕκτης ὥρας σκότος ἐγένετο ἐπὶ πᾶσαν τὴν γῆν ἕως ὥρας ἐνάτης.
\vs{46}περὶ δὲ τὴν ἐνάτην ὥραν ἀνεβόησεν ὁ Ἰησοῦς φωνῇ μεγάλῃ λέγων· 
\begin{poetryblock}

\begin{quote}Ἠλὶ ἠλὶ λεμὰ σαβαχθάνι;\end{quote}
\end{poetryblock}

τοῦτ᾽ ἔστιν· Θεέ μου θεέ μου, ἵνατί με ἐγκατέλιπες;
\vs{47}Τινὲς δὲ τῶν ἐκεῖ ἑστηκότων ἀκούσαντες ἔλεγον ὅτι Ἠλίαν φωνεῖ οὗτος.
\vs{48}καὶ εὐθέως δραμὼν εἷς ἐξ αὐτῶν καὶ λαβὼν σπόγγον πλήσας τε ὄξους καὶ περιθεὶς καλάμῳ ἐπότιζεν αὐτόν.
\vs{49}Οἱ δὲ λοιποὶ ἔλεγον· Ἄφες ἴδωμεν εἰ ἔρχεται Ἠλίας σώσων αὐτόν.
\vs{50}Ὁ δὲ Ἰησοῦς πάλιν κράξας φωνῇ μεγάλῃ ἀφῆκεν τὸ πνεῦμα.

\vs{51}Καὶ ἰδοὺ τὸ καταπέτασμα τοῦ ναοῦ ἐσχίσθη ἀπ᾽ ἄνωθεν ἕως κάτω εἰς δύο καὶ ἡ γῆ ἐσείσθη καὶ αἱ πέτραι ἐσχίσθησαν,
\vs{52}καὶ τὰ μνημεῖα ἀνεῴχθησαν καὶ πολλὰ σώματα τῶν κεκοιμημένων ἁγίων ἠγέρθησαν,
\vs{53}καὶ ἐξελθόντες ἐκ τῶν μνημείων μετὰ τὴν ἔγερσιν αὐτοῦ εἰσῆλθον εἰς τὴν ἁγίαν πόλιν καὶ ἐνεφανίσθησαν πολλοῖς.

\vs{54}Ὁ δὲ ἑκατόνταρχος καὶ οἱ μετ᾽ αὐτοῦ τηροῦντες τὸν Ἰησοῦν ἰδόντες τὸν σεισμὸν καὶ τὰ γενόμενα ἐφοβήθησαν σφόδρα, λέγοντες· Ἀληθῶς Θεοῦ Υἱὸς ἦν οὗτος.

\vs{55}Ἦσαν δὲ ἐκεῖ γυναῖκες πολλαὶ ἀπὸ μακρόθεν θεωροῦσαι, αἵτινες ἠκολούθησαν τῷ Ἰησοῦ ἀπὸ τῆς Γαλιλαίας διακονοῦσαι αὐτῷ·
\vs{56}ἐν αἷς ἦν Μαρία ἡ Μαγδαληνή καὶ Μαρία ἡ τοῦ Ἰακώβου καὶ Ἰωσὴφ μήτηρ καὶ ἡ μήτηρ τῶν υἱῶν Ζεβεδαίου.

\vs{57}Ὀψίας δὲ γενομένης ἦλθεν ἄνθρωπος πλούσιος ἀπὸ Ἁριμαθαίας, τοὔνομα Ἰωσήφ, ὃς καὶ αὐτὸς ἐμαθητεύθη τῷ Ἰησοῦ·
\vs{58}οὗτος προσελθὼν τῷ Πιλάτῳ ᾐτήσατο τὸ σῶμα τοῦ Ἰησοῦ. τότε ὁ Πιλᾶτος ἐκέλευσεν ἀποδοθῆναι.
\vs{59}καὶ λαβὼν τὸ σῶμα ὁ Ἰωσὴφ ἐνετύλιξεν αὐτὸ ἐν σινδόνι καθαρᾷ
\vs{60}καὶ ἔθηκεν αὐτὸ ἐν τῷ καινῷ αὐτοῦ μνημείῳ ὃ ἐλατόμησεν ἐν τῇ πέτρᾳ καὶ προσκυλίσας λίθον μέγαν τῇ θύρᾳ τοῦ μνημείου ἀπῆλθεν.
\vs{61}Ἦν δὲ ἐκεῖ Μαριὰμ ἡ Μαγδαληνὴ καὶ ἡ ἄλλη Μαρία καθήμεναι ἀπέναντι τοῦ τάφου.

\vs{62}Τῇ δὲ ἐπαύριον, ἥτις ἐστὶν μετὰ τὴν Παρασκευήν, συνήχθησαν οἱ ἀρχιερεῖς καὶ οἱ Φαρισαῖοι πρὸς Πιλᾶτον
\vs{63}λέγοντες· Κύριε, ἐμνήσθημεν ὅτι ἐκεῖνος ὁ πλάνος εἶπεν ἔτι ζῶν· Μετὰ τρεῖς ἡμέρας ἐγείρομαι.
\vs{64}κέλευσον οὖν ἀσφαλισθῆναι τὸν τάφον ἕως τῆς τρίτης ἡμέρας, μήποτε ἐλθόντες οἱ μαθηταὶ αὐτοῦ κλέψωσιν αὐτὸν καὶ εἴπωσιν τῷ λαῷ· Ἠγέρθη ἀπὸ τῶν νεκρῶν, καὶ ἔσται ἡ ἐσχάτη πλάνη χείρων τῆς πρώτης.
\vs{65}Ἔφη αὐτοῖς ὁ Πιλᾶτος· Ἔχετε κουστωδίαν· ὑπάγετε ἀσφαλίσασθε ὡς οἴδατε.
\vs{66}οἱ δὲ πορευθέντες ἠσφαλίσαντο τὸν τάφον σφραγίσαντες τὸν λίθον μετὰ τῆς κουστωδίας.

\ch{28}
Ὀψὲ δὲ σαββάτων, τῇ ἐπιφωσκούσῃ εἰς μίαν σαββάτων ἦλθεν Μαριὰμ ἡ Μαγδαληνὴ καὶ ἡ ἄλλη Μαρία θεωρῆσαι τὸν τάφον.
\vs{2}Καὶ ἰδοὺ σεισμὸς ἐγένετο μέγας· ἄγγελος γὰρ Κυρίου καταβὰς ἐξ οὐρανοῦ καὶ προσελθὼν ἀπεκύλισεν τὸν λίθον καὶ ἐκάθητο ἐπάνω αὐτοῦ.
\vs{3}ἦν δὲ ἡ εἰδέα αὐτοῦ ὡς ἀστραπὴ καὶ τὸ ἔνδυμα αὐτοῦ λευκὸν ὡς χιών.
\vs{4}ἀπὸ δὲ τοῦ φόβου αὐτοῦ ἐσείσθησαν οἱ τηροῦντες καὶ ἐγενήθησαν ὡς νεκροί.
\vs{5}Ἀποκριθεὶς δὲ ὁ ἄγγελος εἶπεν ταῖς γυναιξίν· Μὴ φοβεῖσθε ὑμεῖς, οἶδα γὰρ ὅτι Ἰησοῦν τὸν ἐσταυρωμένον ζητεῖτε·
\vs{6}οὐκ ἔστιν ὧδε, ἠγέρθη γὰρ καθὼς εἶπεν· δεῦτε ἴδετε τὸν τόπον ὅπου ἔκειτο.
\vs{7}καὶ ταχὺ πορευθεῖσαι εἴπατε τοῖς μαθηταῖς αὐτοῦ ὅτι Ἠγέρθη ἀπὸ τῶν νεκρῶν, καὶ ἰδοὺ προάγει ὑμᾶς εἰς τὴν Γαλιλαίαν, ἐκεῖ αὐτὸν ὄψεσθε· ἰδοὺ εἶπον ὑμῖν.

\vs{8}Καὶ ἀπελθοῦσαι ταχὺ ἀπὸ τοῦ μνημείου μετὰ φόβου καὶ χαρᾶς μεγάλης ἔδραμον ἀπαγγεῖλαι τοῖς μαθηταῖς αὐτοῦ.
\vs{9}καὶ ἰδοὺ Ἰησοῦς ὑπήντησεν αὐταῖς λέγων· Χαίρετε. αἱ δὲ προσελθοῦσαι ἐκράτησαν αὐτοῦ τοὺς πόδας καὶ προσεκύνησαν αὐτῷ.
\vs{10}τότε λέγει αὐταῖς ὁ Ἰησοῦς· Μὴ φοβεῖσθε· ὑπάγετε ἀπαγγείλατε τοῖς ἀδελφοῖς μου ἵνα ἀπέλθωσιν εἰς τὴν Γαλιλαίαν, κἀκεῖ με ὄψονται.

\vs{11}Πορευομένων δὲ αὐτῶν ἰδού τινες τῆς κουστωδίας ἐλθόντες εἰς τὴν πόλιν ἀπήγγειλαν τοῖς ἀρχιερεῦσιν ἅπαντα τὰ γενόμενα.
\vs{12}καὶ συναχθέντες μετὰ τῶν πρεσβυτέρων συμβούλιόν τε λαβόντες ἀργύρια ἱκανὰ ἔδωκαν τοῖς στρατιώταις
\vs{13}λέγοντες· Εἴπατε ὅτι Οἱ μαθηταὶ αὐτοῦ νυκτὸς ἐλθόντες ἔκλεψαν αὐτὸν ἡμῶν κοιμωμένων.
\vs{14}καὶ ἐὰν ἀκουσθῇ τοῦτο ἐπὶ τοῦ ἡγεμόνος, ἡμεῖς πείσομεν αὐτὸν καὶ ὑμᾶς ἀμερίμνους ποιήσομεν.
\vs{15}Οἱ δὲ λαβόντες τὰ ἀργύρια ἐποίησαν ὡς ἐδιδάχθησαν. Καὶ διεφημίσθη ὁ λόγος οὗτος παρὰ Ἰουδαίοις μέχρι τῆς σήμερον ἡμέρας.

\vs{16}Οἱ δὲ ἕνδεκα μαθηταὶ ἐπορεύθησαν εἰς τὴν Γαλιλαίαν εἰς τὸ ὄρος οὗ ἐτάξατο αὐτοῖς ὁ Ἰησοῦς,
\vs{17}καὶ ἰδόντες αὐτὸν προσεκύνησαν, οἱ δὲ ἐδίστασαν.
\vs{18}Καὶ προσελθὼν ὁ Ἰησοῦς ἐλάλησεν αὐτοῖς λέγων· Ἐδόθη μοι πᾶσα ἐξουσία ἐν οὐρανῷ καὶ ἐπὶ τῆς γῆς.
\vs{19}πορευθέντες οὖν μαθητεύσατε πάντα τὰ ἔθνη, βαπτίζοντες αὐτοὺς εἰς τὸ ὄνομα τοῦ Πατρὸς καὶ τοῦ Υἱοῦ καὶ τοῦ Ἁγίου Πνεύματος,
\vs{20}διδάσκοντες αὐτοὺς τηρεῖν πάντα ὅσα ἐνετειλάμην ὑμῖν· καὶ ἰδοὺ ἐγὼ μεθ᾽ ὑμῶν εἰμι πάσας τὰς ἡμέρας ἕως τῆς συντελείας τοῦ αἰῶνος.


\def\book{ΚΑΤΑ ΜΑΡΚΟΝ}
\biblebook{ΚΑΤΑ ΜΑΡΚΟΝ}


\lettrine[lines=2, loversize=0.2, nindent=0em, findent=.25em]{\textcolor{bookheadingcolor}{Ἀ}}{ρχὴ} τοῦ εὐαγγελίου Ἰησοῦ Χριστοῦ Υἱοῦ Θεοῦ.
\begin{poetryblock}

\begin{quote} \vs{2}Καθὼς γέγραπται ἐν τῷ Ἠσαΐᾳ τῷ προφήτῃ·\end{quote} 

\begin{quote}Ἰδοὺ ἀποστέλλω τὸν ἄγγελόν μου πρὸ προσώπου σου,\end{quote} 

\begin{quote}ὃς κατασκευάσει τὴν ὁδόν σου·\end{quote}

\begin{quote} \vs{3}Φωνὴ βοῶντος ἐν τῇ ἐρήμῳ·\end{quote} 

\begin{quote}Ἑτοιμάσατε τὴν ὁδὸν Κυρίου,\end{quote} 

\begin{quote}εὐθείας ποιεῖτε τὰς τρίβους αὐτοῦ,\end{quote}
\end{poetryblock}

\vs{4}Ἐγένετο Ἰωάννης ὁ βαπτίζων ἐν τῇ ἐρήμῳ καὶ κηρύσσων βάπτισμα μετανοίας εἰς ἄφεσιν ἁμαρτιῶν.
\vs{5}καὶ ἐξεπορεύετο πρὸς αὐτὸν πᾶσα ἡ Ἰουδαία χώρα καὶ οἱ Ἱεροσολυμῖται πάντες, καὶ ἐβαπτίζοντο ὑπ᾽ αὐτοῦ ἐν τῷ Ἰορδάνῃ ποταμῷ ἐξομολογούμενοι τὰς ἁμαρτίας αὐτῶν.
\vs{6}Καὶ ἦν ὁ Ἰωάννης ἐνδεδυμένος τρίχας καμήλου καὶ ζώνην δερματίνην περὶ τὴν ὀσφὺν αὐτοῦ καὶ ἔσθων ἀκρίδας καὶ μέλι ἄγριον.

\vs{7}καὶ ἐκήρυσσεν λέγων· Ἔρχεται ὁ ἰσχυρότερός μου ὀπίσω μου, οὗ οὐκ εἰμὶ ἱκανὸς κύψας λῦσαι τὸν ἱμάντα τῶν ὑποδημάτων αὐτοῦ.
\vs{8}ἐγὼ ἐβάπτισα ὑμᾶς ὕδατι, αὐτὸς δὲ βαπτίσει ὑμᾶς ἐν Πνεύματι Ἁγίῳ.

\vs{9}Καὶ ἐγένετο ἐν ἐκείναις ταῖς ἡμέραις ἦλθεν Ἰησοῦς ἀπὸ Ναζαρὲτ τῆς Γαλιλαίας καὶ ἐβαπτίσθη εἰς τὸν Ἰορδάνην ὑπὸ Ἰωάννου.
\vs{10}καὶ εὐθὺς ἀναβαίνων ἐκ τοῦ ὕδατος εἶδεν σχιζομένους τοὺς οὐρανοὺς καὶ τὸ Πνεῦμα ὡς περιστερὰν καταβαῖνον εἰς αὐτόν·
\vs{11}καὶ φωνὴ ἐγένετο ἐκ τῶν οὐρανῶν· Σὺ εἶ ὁ Υἱός μου ὁ ἀγαπητός, ἐν σοὶ εὐδόκησα.

\vs{12}Καὶ εὐθὺς τὸ Πνεῦμα αὐτὸν ἐκβάλλει εἰς τὴν ἔρημον.
\vs{13}καὶ ἦν ἐν τῇ ἐρήμῳ τεσσεράκοντα ἡμέρας πειραζόμενος ὑπὸ τοῦ Σατανᾶ, καὶ ἦν μετὰ τῶν θηρίων, καὶ οἱ ἄγγελοι διηκόνουν αὐτῷ.

\vs{14}Μετὰ δὲ τὸ παραδοθῆναι τὸν Ἰωάννην ἦλθεν ὁ Ἰησοῦς εἰς τὴν Γαλιλαίαν κηρύσσων τὸ εὐαγγέλιον τοῦ Θεοῦ
\vs{15}καὶ λέγων ὅτι Πεπλήρωται ὁ καιρὸς καὶ ἤγγικεν ἡ βασιλεία τοῦ Θεοῦ· μετανοεῖτε καὶ πιστεύετε ἐν τῷ εὐαγγελίῳ.

\vs{16}Καὶ παράγων παρὰ τὴν θάλασσαν τῆς Γαλιλαίας εἶδεν Σίμωνα καὶ Ἀνδρέαν τὸν ἀδελφὸν Σίμωνος ἀμφιβάλλοντας ἐν τῇ θαλάσσῃ· ἦσαν γὰρ ἁλιεῖς.
\vs{17}καὶ εἶπεν αὐτοῖς ὁ Ἰησοῦς· Δεῦτε ὀπίσω μου, καὶ ποιήσω ὑμᾶς γενέσθαι ἁλιεῖς ἀνθρώπων.
\vs{18}καὶ εὐθὺς ἀφέντες τὰ δίκτυα ἠκολούθησαν αὐτῷ.
\vs{19}Καὶ προβὰς ὀλίγον εἶδεν Ἰάκωβον τὸν τοῦ Ζεβεδαίου καὶ Ἰωάννην τὸν ἀδελφὸν αὐτοῦ καὶ αὐτοὺς ἐν τῷ πλοίῳ καταρτίζοντας τὰ δίκτυα,
\vs{20}καὶ εὐθὺς ἐκάλεσεν αὐτούς. καὶ ἀφέντες τὸν πατέρα αὐτῶν Ζεβεδαῖον ἐν τῷ πλοίῳ μετὰ τῶν μισθωτῶν ἀπῆλθον ὀπίσω αὐτοῦ.

\vs{21}Καὶ εἰσπορεύονται εἰς Καφαρναούμ· καὶ εὐθὺς τοῖς σάββασιν εἰσελθὼν εἰς τὴν συναγωγὴν ἐδίδασκεν.
\vs{22}καὶ ἐξεπλήσσοντο ἐπὶ τῇ διδαχῇ αὐτοῦ· ἦν γὰρ διδάσκων αὐτοὺς ὡς ἐξουσίαν ἔχων καὶ οὐχ ὡς οἱ γραμματεῖς.

\vs{23}Καὶ εὐθὺς ἦν ἐν τῇ συναγωγῇ αὐτῶν ἄνθρωπος ἐν πνεύματι ἀκαθάρτῳ καὶ ἀνέκραξεν
\vs{24}λέγων· Τί ἡμῖν καὶ σοί, Ἰησοῦ Ναζαρηνέ; ἦλθες ἀπολέσαι ἡμᾶς; οἶδά σε τίς εἶ, ὁ Ἅγιος τοῦ Θεοῦ.
\vs{25}Καὶ ἐπετίμησεν αὐτῷ ὁ Ἰησοῦς λέγων· Φιμώθητι καὶ ἔξελθε ἐξ αὐτοῦ.
\vs{26}καὶ σπαράξαν αὐτὸν τὸ πνεῦμα τὸ ἀκάθαρτον καὶ φωνῆσαν φωνῇ μεγάλῃ ἐξῆλθεν ἐξ αὐτοῦ.
\vs{27}Καὶ ἐθαμβήθησαν ἅπαντες ὥστε συζητεῖν πρὸς ἑαυτοὺς λέγοντας· Τί ἐστιν τοῦτο; διδαχὴ καινή κατ᾽ ἐξουσίαν· καὶ τοῖς πνεύμασι τοῖς ἀκαθάρτοις ἐπιτάσσει, καὶ ὑπακούουσιν αὐτῷ.
\vs{28}καὶ ἐξῆλθεν ἡ ἀκοὴ αὐτοῦ εὐθὺς πανταχοῦ εἰς ὅλην τὴν περίχωρον τῆς Γαλιλαίας.

\vs{29}Καὶ εὐθὺς ἐκ τῆς συναγωγῆς ἐξελθόντες ἦλθον εἰς τὴν οἰκίαν Σίμωνος καὶ Ἀνδρέου μετὰ Ἰακώβου καὶ Ἰωάννου.
\vs{30}ἡ δὲ πενθερὰ Σίμωνος κατέκειτο πυρέσσουσα, καὶ εὐθὺς λέγουσιν αὐτῷ περὶ αὐτῆς.
\vs{31}καὶ προσελθὼν ἤγειρεν αὐτὴν κρατήσας τῆς χειρός· καὶ ἀφῆκεν αὐτὴν ὁ πυρετός, καὶ διηκόνει αὐτοῖς.

\vs{32}Ὀψίας δὲ γενομένης, ὅτε ἔδυ ὁ ἥλιος, ἔφερον πρὸς αὐτὸν πάντας τοὺς κακῶς ἔχοντας καὶ τοὺς δαιμονιζομένους·
\vs{33}καὶ ἦν ὅλη ἡ πόλις ἐπισυνηγμένη πρὸς τὴν θύραν.
\vs{34}καὶ ἐθεράπευσεν πολλοὺς κακῶς ἔχοντας ποικίλαις νόσοις καὶ δαιμόνια πολλὰ ἐξέβαλεν καὶ οὐκ ἤφιεν λαλεῖν τὰ δαιμόνια, ὅτι ᾔδεισαν αὐτόν.

\vs{35}Καὶ πρωῒ ἔννυχα λίαν ἀναστὰς ἐξῆλθεν καὶ ἀπῆλθεν εἰς ἔρημον τόπον κἀκεῖ προσηύχετο.
\vs{36}καὶ κατεδίωξεν αὐτὸν Σίμων καὶ οἱ μετ᾽ αὐτοῦ,
\vs{37}καὶ εὗρον αὐτὸν καὶ λέγουσιν αὐτῷ ὅτι Πάντες ζητοῦσίν σε.
\vs{38}Καὶ λέγει αὐτοῖς· Ἄγωμεν ἀλλαχοῦ εἰς τὰς ἐχομένας κωμοπόλεις, ἵνα καὶ ἐκεῖ κηρύξω· εἰς τοῦτο γὰρ ἐξῆλθον.

\vs{39}καὶ ἦλθεν κηρύσσων εἰς τὰς συναγωγὰς αὐτῶν εἰς ὅλην τὴν Γαλιλαίαν καὶ τὰ δαιμόνια ἐκβάλλων.
\vs{40}Καὶ ἔρχεται πρὸς αὐτὸν λεπρὸς παρακαλῶν αὐτὸν καὶ γονυπετῶν καὶ λέγων αὐτῷ ὅτι Ἐὰν θέλῃς δύνασαί με καθαρίσαι.
\vs{41}Καὶ σπλαγχνισθεὶς ἐκτείνας τὴν χεῖρα αὐτοῦ ἥψατο καὶ λέγει αὐτῷ· Θέλω, καθαρίσθητι·
\vs{42}Καὶ εὐθὺς ἀπῆλθεν ἀπ᾽ αὐτοῦ ἡ λέπρα, καὶ ἐκαθαρίσθη.
\vs{43}Καὶ ἐμβριμησάμενος αὐτῷ εὐθὺς ἐξέβαλεν αὐτόν
\vs{44}καὶ λέγει αὐτῷ· Ὅρα μηδενὶ μηδὲν εἴπῃς, ἀλλὰ ὕπαγε σεαυτὸν δεῖξον τῷ ἱερεῖ καὶ προσένεγκε περὶ τοῦ καθαρισμοῦ σου ἃ προσέταξεν Μωϋσῆς, εἰς μαρτύριον αὐτοῖς.
\vs{45}Ὁ δὲ ἐξελθὼν ἤρξατο κηρύσσειν πολλὰ καὶ διαφημίζειν τὸν λόγον, ὥστε μηκέτι αὐτὸν δύνασθαι φανερῶς εἰς πόλιν εἰσελθεῖν, ἀλλ᾽ ἔξω ἐπ᾽ ἐρήμοις τόποις ἦν· καὶ ἤρχοντο πρὸς αὐτὸν πάντοθεν.

\ch{2}
Καὶ εἰσελθὼν πάλιν εἰς Καφαρναοὺμ δι᾽ ἡμερῶν ἠκούσθη ὅτι ἐν οἴκῳ ἐστίν.
\vs{2}καὶ συνήχθησαν πολλοὶ ὥστε μηκέτι χωρεῖν μηδὲ τὰ πρὸς τὴν θύραν, καὶ ἐλάλει αὐτοῖς τὸν λόγον.
\vs{3}Καὶ ἔρχονται φέροντες πρὸς αὐτὸν παραλυτικὸν αἰρόμενον ὑπὸ τεσσάρων.
\vs{4}καὶ μὴ δυνάμενοι προσενέγκαι αὐτῷ διὰ τὸν ὄχλον ἀπεστέγασαν τὴν στέγην ὅπου ἦν, καὶ ἐξορύξαντες χαλῶσι τὸν κράβαττον ὅπου ὁ παραλυτικὸς κατέκειτο.
\vs{5}Καὶ ἰδὼν ὁ Ἰησοῦς τὴν πίστιν αὐτῶν λέγει τῷ παραλυτικῷ· Τέκνον, ἀφίενταί σου αἱ ἁμαρτίαι.
\vs{6}Ἦσαν δέ τινες τῶν γραμματέων ἐκεῖ καθήμενοι καὶ διαλογιζόμενοι ἐν ταῖς καρδίαις αὐτῶν·
\vs{7}Τί οὗτος οὕτως λαλεῖ; βλασφημεῖ· τίς δύναται ἀφιέναι ἁμαρτίας εἰ μὴ εἷς ὁ Θεός;
\vs{8}Καὶ εὐθὺς ἐπιγνοὺς ὁ Ἰησοῦς τῷ πνεύματι αὐτοῦ ὅτι οὕτως διαλογίζονται ἐν ἑαυτοῖς λέγει αὐτοῖς· Τί ταῦτα διαλογίζεσθε ἐν ταῖς καρδίαις ὑμῶν;
\vs{9}τί ἐστιν εὐκοπώτερον, εἰπεῖν τῷ παραλυτικῷ· Ἀφίενταί σου αἱ ἁμαρτίαι, ἢ εἰπεῖν· Ἔγειρε καὶ ἆρον τὸν κράβαττόν σου καὶ περιπάτει;
\vs{10}ἵνα δὲ εἰδῆτε ὅτι ἐξουσίαν ἔχει ὁ Υἱὸς τοῦ ἀνθρώπου ἀφιέναι ἁμαρτίας ἐπὶ τῆς γῆς— λέγει τῷ παραλυτικῷ·
\vs{11}Σοὶ λέγω, ἔγειρε ἆρον τὸν κράβαττόν σου καὶ ὕπαγε εἰς τὸν οἶκόν σου.
\vs{12}Καὶ ἠγέρθη καὶ εὐθὺς ἄρας τὸν κράβαττον ἐξῆλθεν ἔμπροσθεν πάντων, ὥστε ἐξίστασθαι πάντας καὶ δοξάζειν τὸν Θεὸν λέγοντας ὅτι Οὕτως οὐδέποτε εἴδομεν.

\vs{13}Καὶ ἐξῆλθεν πάλιν παρὰ τὴν θάλασσαν· καὶ πᾶς ὁ ὄχλος ἤρχετο πρὸς αὐτόν, καὶ ἐδίδασκεν αὐτούς.
\vs{14}Καὶ παράγων εἶδεν Λευὶν τὸν τοῦ Ἁλφαίου καθήμενον ἐπὶ τὸ τελώνιον, καὶ λέγει αὐτῷ· Ἀκολούθει μοι. καὶ ἀναστὰς ἠκολούθησεν αὐτῷ.

\vs{15}Καὶ γίνεται κατακεῖσθαι αὐτὸν ἐν τῇ οἰκίᾳ αὐτοῦ, καὶ πολλοὶ τελῶναι καὶ ἁμαρτωλοὶ συνανέκειντο τῷ Ἰησοῦ καὶ τοῖς μαθηταῖς αὐτοῦ· ἦσαν γὰρ πολλοὶ καὶ ἠκολούθουν αὐτῷ.
\vs{16}καὶ οἱ γραμματεῖς τῶν Φαρισαίων ἰδόντες ὅτι ἐσθίει μετὰ τῶν ἁμαρτωλῶν καὶ τελωνῶν ἔλεγον τοῖς μαθηταῖς αὐτοῦ· Ὅτι μετὰ τῶν τελωνῶν καὶ ἁμαρτωλῶν ἐσθίει;
\vs{17}Καὶ ἀκούσας ὁ Ἰησοῦς λέγει αὐτοῖς ὅτι Οὐ χρείαν ἔχουσιν οἱ ἰσχύοντες ἰατροῦ ἀλλ᾽ οἱ κακῶς ἔχοντες· οὐκ ἦλθον καλέσαι δικαίους ἀλλὰ ἁμαρτωλούς.

\vs{18}Καὶ ἦσαν οἱ μαθηταὶ Ἰωάννου καὶ οἱ Φαρισαῖοι νηστεύοντες. καὶ ἔρχονται καὶ λέγουσιν αὐτῷ· Διὰ τί οἱ μαθηταὶ Ἰωάννου καὶ οἱ μαθηταὶ τῶν Φαρισαίων νηστεύουσιν, οἱ δὲ σοὶ μαθηταὶ οὐ νηστεύουσιν;
\vs{19}Καὶ εἶπεν αὐτοῖς ὁ Ἰησοῦς· Μὴ δύνανται οἱ υἱοὶ τοῦ νυμφῶνος ἐν ᾧ ὁ νυμφίος μετ᾽ αὐτῶν ἐστιν νηστεύειν; ὅσον χρόνον ἔχουσιν τὸν νυμφίον μετ᾽ αὐτῶν οὐ δύνανται νηστεύειν.
\vs{20}ἐλεύσονται δὲ ἡμέραι ὅταν ἀπαρθῇ ἀπ᾽ αὐτῶν ὁ νυμφίος, καὶ τότε νηστεύσουσιν ἐν ἐκείνῃ τῇ ἡμέρᾳ.
\vs{21}Οὐδεὶς ἐπίβλημα ῥάκους ἀγνάφου ἐπιράπτει ἐπὶ ἱμάτιον παλαιόν· εἰ δὲ μή, αἴρει τὸ πλήρωμα ἀπ᾽ αὐτοῦ τὸ καινὸν τοῦ παλαιοῦ καὶ χεῖρον σχίσμα γίνεται.
\vs{22}Καὶ οὐδεὶς βάλλει οἶνον νέον εἰς ἀσκοὺς παλαιούς· εἰ δὲ μή, ῥήξει ὁ οἶνος τοὺς ἀσκούς καὶ ὁ οἶνος ἀπόλλυται καὶ οἱ ἀσκοί· ἀλλὰ οἶνον νέον εἰς ἀσκοὺς καινούς.
\vs{23}Καὶ ἐγένετο αὐτὸν ἐν τοῖς σάββασιν παραπορεύεσθαι διὰ τῶν σπορίμων, καὶ οἱ μαθηταὶ αὐτοῦ ἤρξαντο ὁδὸν ποιεῖν τίλλοντες τοὺς στάχυας.
\vs{24}καὶ οἱ Φαρισαῖοι ἔλεγον αὐτῷ· Ἴδε τί ποιοῦσιν τοῖς σάββασιν ὃ οὐκ ἔξεστιν;
\vs{25}Καὶ λέγει αὐτοῖς· Οὐδέποτε ἀνέγνωτε τί ἐποίησεν Δαυὶδ ὅτε χρείαν ἔσχεν καὶ ἐπείνασεν αὐτὸς καὶ οἱ μετ᾽ αὐτοῦ,
\vs{26}πῶς εἰσῆλθεν εἰς τὸν οἶκον τοῦ Θεοῦ ἐπὶ Ἀβιαθὰρ ἀρχιερέως καὶ τοὺς ἄρτους τῆς προθέσεως ἔφαγεν, οὓς οὐκ ἔξεστιν φαγεῖν εἰ μὴ τοὺς ἱερεῖς, καὶ ἔδωκεν καὶ τοῖς σὺν αὐτῷ οὖσιν;
\vs{27}Καὶ ἔλεγεν αὐτοῖς· Τὸ σάββατον διὰ τὸν ἄνθρωπον ἐγένετο καὶ οὐχ ὁ ἄνθρωπος διὰ τὸ σάββατον·
\vs{28}ὥστε κύριός ἐστιν ὁ Υἱὸς τοῦ ἀνθρώπου καὶ τοῦ σαββάτου.

\ch{3}
Καὶ εἰσῆλθεν πάλιν εἰς τὴν συναγωγήν. καὶ ἦν ἐκεῖ ἄνθρωπος ἐξηραμμένην ἔχων τὴν χεῖρα.
\vs{2}καὶ παρετήρουν αὐτὸν εἰ τοῖς σάββασιν θεραπεύσει αὐτόν, ἵνα κατηγορήσωσιν αὐτοῦ.
\vs{3}Καὶ λέγει τῷ ἀνθρώπῳ τῷ τὴν χεῖρα ἔχοντι ξηράν· Ἔγειρε εἰς τὸ μέσον.
\vs{4}καὶ λέγει αὐτοῖς· Ἔξεστιν τοῖς σάββασιν ἀγαθὸν ποιῆσαι ἢ κακοποιῆσαι, ψυχὴν σῶσαι ἢ ἀποκτεῖναι; Οἱ δὲ ἐσιώπων.
\vs{5}Καὶ περιβλεψάμενος αὐτοὺς μετ᾽ ὀργῆς, συλλυπούμενος ἐπὶ τῇ πωρώσει τῆς καρδίας αὐτῶν λέγει τῷ ἀνθρώπῳ· Ἔκτεινον τὴν χεῖρα. καὶ ἐξέτεινεν καὶ ἀπεκατεστάθη ἡ χεὶρ αὐτοῦ.
\vs{6}Καὶ ἐξελθόντες οἱ Φαρισαῖοι εὐθὺς μετὰ τῶν Ἡρῳδιανῶν συμβούλιον ἐδίδουν κατ᾽ αὐτοῦ ὅπως αὐτὸν ἀπολέσωσιν.

\vs{7}Καὶ ὁ Ἰησοῦς μετὰ τῶν μαθητῶν αὐτοῦ ἀνεχώρησεν πρὸς τὴν θάλασσαν, καὶ πολὺ πλῆθος ἀπὸ τῆς Γαλιλαίας ἠκολούθησεν, καὶ ἀπὸ τῆς Ἰουδαίας
\vs{8}καὶ ἀπὸ Ἱεροσολύμων καὶ ἀπὸ τῆς Ἰδουμαίας καὶ πέραν τοῦ Ἰορδάνου καὶ περὶ Τύρον καὶ Σιδῶνα πλῆθος πολύ ἀκούοντες ὅσα ἐποίει ἦλθον πρὸς αὐτόν.
\vs{9}Καὶ εἶπεν τοῖς μαθηταῖς αὐτοῦ ἵνα πλοιάριον προσκαρτερῇ αὐτῷ διὰ τὸν ὄχλον ἵνα μὴ θλίβωσιν αὐτόν·
\vs{10}πολλοὺς γὰρ ἐθεράπευσεν, ὥστε ἐπιπίπτειν αὐτῷ ἵνα αὐτοῦ ἅψωνται ὅσοι εἶχον μάστιγας.
\vs{11}καὶ τὰ πνεύματα τὰ ἀκάθαρτα, ὅταν αὐτὸν ἐθεώρουν, προσέπιπτον αὐτῷ καὶ ἔκραζον λέγοντα ὅτι Σὺ εἶ ὁ Υἱὸς τοῦ Θεοῦ.
\vs{12}καὶ πολλὰ ἐπετίμα αὐτοῖς ἵνα μὴ αὐτὸν φανερὸν ποιήσωσιν.

\vs{13}Καὶ ἀναβαίνει εἰς τὸ ὄρος καὶ προσκαλεῖται οὓς ἤθελεν αὐτός, καὶ ἀπῆλθον πρὸς αὐτόν.
\vs{14}καὶ ἐποίησεν δώδεκα οὓς καὶ ἀποστόλους ὠνόμασεν ἵνα ὦσιν μετ᾽ αὐτοῦ καὶ ἵνα ἀποστέλλῃ αὐτοὺς κηρύσσειν
\vs{15}καὶ ἔχειν ἐξουσίαν ἐκβάλλειν τὰ δαιμόνια·
\vs{16}Καὶ ἐποίησεν τοὺς δώδεκα, καὶ ἐπέθηκεν ὄνομα τῷ Σίμωνι Πέτρον,
\vs{17}καὶ Ἰάκωβον τὸν τοῦ Ζεβεδαίου καὶ Ἰωάννην τὸν ἀδελφὸν τοῦ Ἰακώβου καὶ ἐπέθηκεν αὐτοῖς ὀνόματα Βοανηργές, ὅ ἐστιν Υἱοὶ Βροντῆς·
\vs{18}καὶ Ἀνδρέαν καὶ Φίλιππον καὶ Βαρθολομαῖον καὶ Μαθθαῖον καὶ Θωμᾶν καὶ Ἰάκωβον τὸν τοῦ Ἁλφαίου καὶ Θαδδαῖον καὶ Σίμωνα τὸν Καναναῖον
\vs{19}καὶ Ἰούδαν Ἰσκαριώθ, ὃς καὶ παρέδωκεν αὐτόν.

\vs{20}Καὶ ἔρχεται εἰς οἶκον· καὶ συνέρχεται πάλιν ὁ ὄχλος, ὥστε μὴ δύνασθαι αὐτοὺς μηδὲ ἄρτον φαγεῖν.
\vs{21}καὶ ἀκούσαντες οἱ παρ᾽ αὐτοῦ ἐξῆλθον κρατῆσαι αὐτόν· ἔλεγον γὰρ ὅτι Ἐξέστη.

\vs{22}Καὶ οἱ γραμματεῖς οἱ ἀπὸ Ἱεροσολύμων καταβάντες ἔλεγον ὅτι Βεελζεβοὺλ ἔχει καὶ ὅτι Ἐν τῷ ἄρχοντι τῶν δαιμονίων ἐκβάλλει τὰ δαιμόνια.

\vs{23}Καὶ προσκαλεσάμενος αὐτοὺς ἐν παραβολαῖς ἔλεγεν αὐτοῖς· Πῶς δύναται Σατανᾶς Σατανᾶν ἐκβάλλειν;
\vs{24}καὶ ἐὰν βασιλεία ἐφ᾽ ἑαυτὴν μερισθῇ, οὐ δύναται σταθῆναι ἡ βασιλεία ἐκείνη·
\vs{25}καὶ ἐὰν οἰκία ἐφ᾽ ἑαυτὴν μερισθῇ, οὐ δυνήσεται ἡ οἰκία ἐκείνη σταθῆναι.
\vs{26}καὶ εἰ ὁ Σατανᾶς ἀνέστη ἐφ᾽ ἑαυτὸν καὶ ἐμερίσθη, οὐ δύναται στῆναι ἀλλὰ τέλος ἔχει.
\vs{27}ἀλλ᾽ οὐ δύναται οὐδεὶς εἰς τὴν οἰκίαν τοῦ ἰσχυροῦ εἰσελθὼν τὰ σκεύη αὐτοῦ διαρπάσαι, ἐὰν μὴ πρῶτον τὸν ἰσχυρὸν δήσῃ, καὶ τότε τὴν οἰκίαν αὐτοῦ διαρπάσει.

\vs{28}Ἀμὴν λέγω ὑμῖν ὅτι πάντα ἀφεθήσεται τοῖς υἱοῖς τῶν ἀνθρώπων τὰ ἁμαρτήματα καὶ αἱ βλασφημίαι ὅσα ἐὰν βλασφημήσωσιν·
\vs{29}ὃς δ᾽ ἂν βλασφημήσῃ εἰς τὸ Πνεῦμα τὸ Ἅγιον, οὐκ ἔχει ἄφεσιν εἰς τὸν αἰῶνα, ἀλλὰ ἔνοχός ἐστιν αἰωνίου ἁμαρτήματος.
\vs{30}Ὅτι ἔλεγον· Πνεῦμα ἀκάθαρτον ἔχει.

\vs{31}Καὶ ἔρχονται ἡ μήτηρ αὐτοῦ καὶ οἱ ἀδελφοὶ αὐτοῦ καὶ ἔξω στήκοντες ἀπέστειλαν πρὸς αὐτὸν καλοῦντες αὐτόν.
\vs{32}καὶ ἐκάθητο περὶ αὐτὸν ὄχλος, καὶ λέγουσιν αὐτῷ· Ἰδοὺ ἡ μήτηρ σου καὶ οἱ ἀδελφοί σου καὶ αἱ ἀδελφαί σου ἔξω ζητοῦσίν σε.
\vs{33}Καὶ ἀποκριθεὶς αὐτοῖς λέγει· Τίς ἐστιν ἡ μήτηρ μου καὶ οἱ ἀδελφοί μου;
\vs{34}καὶ περιβλεψάμενος τοὺς περὶ αὐτὸν κύκλῳ καθημένους λέγει· Ἴδε ἡ μήτηρ μου καὶ οἱ ἀδελφοί μου.
\vs{35}ὃς γὰρ ἂν ποιήσῃ τὸ θέλημα τοῦ Θεοῦ, οὗτος ἀδελφός μου καὶ ἀδελφὴ καὶ μήτηρ ἐστίν.

\ch{4}
Καὶ πάλιν ἤρξατο διδάσκειν παρὰ τὴν θάλασσαν· καὶ συνάγεται πρὸς αὐτὸν ὄχλος πλεῖστος, ὥστε αὐτὸν εἰς πλοῖον ἐμβάντα καθῆσθαι ἐν τῇ θαλάσσῃ, καὶ πᾶς ὁ ὄχλος πρὸς τὴν θάλασσαν ἐπὶ τῆς γῆς ἦσαν.
\vs{2}Καὶ ἐδίδασκεν αὐτοὺς ἐν παραβολαῖς πολλά καὶ ἔλεγεν αὐτοῖς ἐν τῇ διδαχῇ αὐτοῦ·

\vs{3}Ἀκούετε. ἰδοὺ ἐξῆλθεν ὁ σπείρων σπεῖραι.
\vs{4}καὶ ἐγένετο ἐν τῷ σπείρειν ὃ μὲν ἔπεσεν παρὰ τὴν ὁδόν, καὶ ἦλθεν τὰ πετεινὰ καὶ κατέφαγεν αὐτό.
\vs{5}Καὶ ἄλλο ἔπεσεν ἐπὶ τὸ πετρῶδες ὅπου οὐκ εἶχεν γῆν πολλήν, καὶ εὐθὺς ἐξανέτειλεν διὰ τὸ μὴ ἔχειν βάθος γῆς·
\vs{6}καὶ ὅτε ἀνέτειλεν ὁ ἥλιος ἐκαυματίσθη καὶ διὰ τὸ μὴ ἔχειν ῥίζαν ἐξηράνθη.
\vs{7}Καὶ ἄλλο ἔπεσεν εἰς τὰς ἀκάνθας, καὶ ἀνέβησαν αἱ ἄκανθαι καὶ συνέπνιξαν αὐτό, καὶ καρπὸν οὐκ ἔδωκεν.
\vs{8}Καὶ ἄλλα ἔπεσεν εἰς τὴν γῆν τὴν καλήν καὶ ἐδίδου καρπὸν ἀναβαίνοντα καὶ αὐξανόμενα καὶ ἔφερεν ἓν τριάκοντα καὶ ἓν ἑξήκοντα καὶ ἓν ἑκατόν.
\vs{9}Καὶ ἔλεγεν· Ὃς ἔχει ὦτα ἀκούειν ἀκουέτω.
\vs{10}Καὶ ὅτε ἐγένετο κατὰ μόνας, ἠρώτων αὐτὸν οἱ περὶ αὐτὸν σὺν τοῖς δώδεκα τὰς παραβολάς.
\vs{11}Καὶ ἔλεγεν αὐτοῖς· Ὑμῖν τὸ μυστήριον δέδοται τῆς βασιλείας τοῦ Θεοῦ· ἐκείνοις δὲ τοῖς ἔξω ἐν παραβολαῖς τὰ πάντα γίνεται,
\begin{poetryblock}

\begin{quote} \vs{12}ἵνα Βλέποντες βλέπωσιν καὶ μὴ ἴδωσιν,\end{quote} 

\begin{quote}Καὶ ἀκούοντες ἀκούωσιν καὶ μὴ συνιῶσιν,\end{quote} 

\begin{quote}Μήποτε ἐπιστρέψωσιν Καὶ ἀφεθῇ αὐτοῖς.\end{quote}
\end{poetryblock}
\vs{13}Καὶ λέγει αὐτοῖς· Οὐκ οἴδατε τὴν παραβολὴν ταύτην, καὶ πῶς πάσας τὰς παραβολὰς γνώσεσθε;
\vs{14}Ὁ σπείρων τὸν λόγον σπείρει.
\vs{15}οὗτοι δέ εἰσιν οἱ παρὰ τὴν ὁδὸν· ὅπου σπείρεται ὁ λόγος καὶ ὅταν ἀκούσωσιν, εὐθὺς ἔρχεται ὁ Σατανᾶς καὶ αἴρει τὸν λόγον τὸν ἐσπαρμένον εἰς αὐτούς.
\vs{16}Καὶ οὗτοί εἰσιν οἱ ἐπὶ τὰ πετρώδη σπειρόμενοι, οἳ ὅταν ἀκούσωσιν τὸν λόγον εὐθὺς μετὰ χαρᾶς λαμβάνουσιν αὐτόν,
\vs{17}καὶ οὐκ ἔχουσιν ῥίζαν ἐν ἑαυτοῖς ἀλλὰ πρόσκαιροί εἰσιν, εἶτα γενομένης θλίψεως ἢ διωγμοῦ διὰ τὸν λόγον εὐθὺς σκανδαλίζονται.
\vs{18}Καὶ ἄλλοι εἰσὶν οἱ εἰς τὰς ἀκάνθας σπειρόμενοι· οὗτοί εἰσιν οἱ τὸν λόγον ἀκούσαντες,
\vs{19}καὶ αἱ μέριμναι τοῦ αἰῶνος καὶ ἡ ἀπάτη τοῦ πλούτου καὶ αἱ περὶ τὰ λοιπὰ ἐπιθυμίαι εἰσπορευόμεναι συμπνίγουσιν τὸν λόγον καὶ ἄκαρπος γίνεται.
\vs{20}Καὶ ἐκεῖνοί εἰσιν οἱ ἐπὶ τὴν γῆν τὴν καλὴν σπαρέντες, οἵτινες ἀκούουσιν τὸν λόγον καὶ παραδέχονται καὶ καρποφοροῦσιν ἓν τριάκοντα καὶ ἓν ἑξήκοντα καὶ ἓν ἑκατόν.

\vs{21}Καὶ ἔλεγεν αὐτοῖς· Μήτι ἔρχεται ὁ λύχνος ἵνα ὑπὸ τὸν μόδιον τεθῇ ἢ ὑπὸ τὴν κλίνην; οὐχ ἵνα ἐπὶ τὴν λυχνίαν τεθῇ;
\vs{22}οὐ γάρ ἐστιν κρυπτὸν ἐὰν μὴ ἵνα φανερωθῇ, οὐδὲ ἐγένετο ἀπόκρυφον ἀλλ᾽ ἵνα ἔλθῃ εἰς φανερόν.
\vs{23}Εἴ τις ἔχει ὦτα ἀκούειν ἀκουέτω.
\vs{24}Καὶ ἔλεγεν αὐτοῖς· Βλέπετε τί ἀκούετε. ἐν ᾧ μέτρῳ μετρεῖτε μετρηθήσεται ὑμῖν καὶ προστεθήσεται ὑμῖν.
\vs{25}ὃς γὰρ ἔχει, δοθήσεται αὐτῷ· καὶ ὃς οὐκ ἔχει, καὶ ὃ ἔχει ἀρθήσεται ἀπ᾽ αὐτοῦ.

\vs{26}Καὶ ἔλεγεν· Οὕτως ἐστὶν ἡ βασιλεία τοῦ Θεοῦ ὡς ἄνθρωπος βάλῃ τὸν σπόρον ἐπὶ τῆς γῆς
\vs{27}καὶ καθεύδῃ καὶ ἐγείρηται νύκτα καὶ ἡμέραν, καὶ ὁ σπόρος βλαστᾷ καὶ μηκύνηται ὡς οὐκ οἶδεν αὐτός.
\vs{28}αὐτομάτη ἡ γῆ καρποφορεῖ, πρῶτον χόρτον εἶτα στάχυν εἶτα πλήρης σῖτον ἐν τῷ στάχυϊ.
\vs{29}ὅταν δὲ παραδοῖ ὁ καρπός, εὐθὺς ἀποστέλλει τὸ δρέπανον, ὅτι παρέστηκεν ὁ θερισμός.

\vs{30}Καὶ ἔλεγεν· Πῶς ὁμοιώσωμεν τὴν βασιλείαν τοῦ Θεοῦ ἢ ἐν τίνι αὐτὴν παραβολῇ θῶμεν;
\vs{31}ὡς κόκκῳ σινάπεως, ὃς ὅταν σπαρῇ ἐπὶ τῆς γῆς, μικρότερον ὂν πάντων τῶν σπερμάτων τῶν ἐπὶ τῆς γῆς,
\vs{32}καὶ ὅταν σπαρῇ, ἀναβαίνει καὶ γίνεται μεῖζον πάντων τῶν λαχάνων καὶ ποιεῖ κλάδους μεγάλους, ὥστε δύνασθαι ὑπὸ τὴν σκιὰν αὐτοῦ τὰ πετεινὰ τοῦ οὐρανοῦ κατασκηνοῦν.

\vs{33}Καὶ τοιαύταις παραβολαῖς πολλαῖς ἐλάλει αὐτοῖς τὸν λόγον καθὼς ἠδύναντο ἀκούειν·
\vs{34}χωρὶς δὲ παραβολῆς οὐκ ἐλάλει αὐτοῖς, κατ᾽ ἰδίαν δὲ τοῖς ἰδίοις μαθηταῖς ἐπέλυεν πάντα.

\vs{35}Καὶ λέγει αὐτοῖς ἐν ἐκείνῃ τῇ ἡμέρᾳ ὀψίας γενομένης· Διέλθωμεν εἰς τὸ πέραν.
\vs{36}καὶ ἀφέντες τὸν ὄχλον παραλαμβάνουσιν αὐτὸν ὡς ἦν ἐν τῷ πλοίῳ, καὶ ἄλλα πλοῖα ἦν μετ᾽ αὐτοῦ.
\vs{37}Καὶ γίνεται λαῖλαψ μεγάλη ἀνέμου καὶ τὰ κύματα ἐπέβαλλεν εἰς τὸ πλοῖον, ὥστε ἤδη γεμίζεσθαι τὸ πλοῖον.
\vs{38}καὶ αὐτὸς ἦν ἐν τῇ πρύμνῃ ἐπὶ τὸ προσκεφάλαιον καθεύδων. καὶ ἐγείρουσιν αὐτὸν καὶ λέγουσιν αὐτῷ· Διδάσκαλε, οὐ μέλει σοι ὅτι ἀπολλύμεθα;
\vs{39}Καὶ διεγερθεὶς ἐπετίμησεν τῷ ἀνέμῳ καὶ εἶπεν τῇ θαλάσσῃ· Σιώπα, πεφίμωσο. καὶ ἐκόπασεν ὁ ἄνεμος καὶ ἐγένετο γαλήνη μεγάλη.
\vs{40}Καὶ εἶπεν αὐτοῖς· Τί δειλοί ἐστε; οὔπω ἔχετε πίστιν;
\vs{41}Καὶ ἐφοβήθησαν φόβον μέγαν καὶ ἔλεγον πρὸς ἀλλήλους· Τίς ἄρα οὗτός ἐστιν ὅτι καὶ ὁ ἄνεμος καὶ ἡ θάλασσα ὑπακούει αὐτῷ;

\ch{5}
Καὶ ἦλθον εἰς τὸ πέραν τῆς θαλάσσης εἰς τὴν χώραν τῶν Γερασηνῶν.
\vs{2}καὶ ἐξελθόντος αὐτοῦ ἐκ τοῦ πλοίου εὐθὺς ὑπήντησεν αὐτῷ ἐκ τῶν μνημείων ἄνθρωπος ἐν πνεύματι ἀκαθάρτῳ,
\vs{3}ὃς τὴν κατοίκησιν εἶχεν ἐν τοῖς μνήμασιν, καὶ οὐδὲ ἁλύσει οὐκέτι οὐδεὶς ἐδύνατο αὐτὸν δῆσαι
\vs{4}διὰ τὸ αὐτὸν πολλάκις πέδαις καὶ ἁλύσεσιν δεδέσθαι καὶ διεσπάσθαι ὑπ᾽ αὐτοῦ τὰς ἁλύσεις καὶ τὰς πέδας συντετρῖφθαι, καὶ οὐδεὶς ἴσχυεν αὐτὸν δαμάσαι·
\vs{5}καὶ διὰ παντὸς νυκτὸς καὶ ἡμέρας ἐν τοῖς μνήμασιν καὶ ἐν τοῖς ὄρεσιν ἦν κράζων καὶ κατακόπτων ἑαυτὸν λίθοις.
\vs{6}Καὶ ἰδὼν τὸν Ἰησοῦν ἀπὸ μακρόθεν ἔδραμεν καὶ προσεκύνησεν αὐτῷ
\vs{7}καὶ κράξας φωνῇ μεγάλῃ λέγει· Τί ἐμοὶ καὶ σοί, Ἰησοῦ Υἱὲ τοῦ Θεοῦ τοῦ Ὑψίστου; ὁρκίζω σε τὸν Θεόν, μή με βασανίσῃς.
\vs{8}ἔλεγεν γὰρ αὐτῷ· Ἔξελθε τὸ πνεῦμα τὸ ἀκάθαρτον ἐκ τοῦ ἀνθρώπου.
\vs{9}Καὶ ἐπηρώτα αὐτόν· Τί ὄνομά σοι; Καὶ λέγει αὐτῷ· Λεγιὼν ὄνομά μοι, ὅτι πολλοί ἐσμεν.
\vs{10}καὶ παρεκάλει αὐτὸν πολλὰ ἵνα μὴ αὐτὰ ἀποστείλῃ ἔξω τῆς χώρας.
\vs{11}Ἦν δὲ ἐκεῖ πρὸς τῷ ὄρει ἀγέλη χοίρων μεγάλη βοσκομένη·
\vs{12}καὶ παρεκάλεσαν αὐτὸν λέγοντες· Πέμψον ἡμᾶς εἰς τοὺς χοίρους, ἵνα εἰς αὐτοὺς εἰσέλθωμεν.
\vs{13}Καὶ ἐπέτρεψεν αὐτοῖς. καὶ ἐξελθόντα τὰ πνεύματα τὰ ἀκάθαρτα εἰσῆλθον εἰς τοὺς χοίρους, καὶ ὥρμησεν ἡ ἀγέλη κατὰ τοῦ κρημνοῦ εἰς τὴν θάλασσαν, ὡς δισχίλιοι, καὶ ἐπνίγοντο ἐν τῇ θαλάσσῃ.

\vs{14}Καὶ οἱ βόσκοντες αὐτοὺς ἔφυγον καὶ ἀπήγγειλαν εἰς τὴν πόλιν καὶ εἰς τοὺς ἀγρούς· καὶ ἦλθον ἰδεῖν τί ἐστιν τὸ γεγονός
\vs{15}καὶ ἔρχονται πρὸς τὸν Ἰησοῦν καὶ θεωροῦσιν τὸν δαιμονιζόμενον καθήμενον ἱματισμένον καὶ σωφρονοῦντα, τὸν ἐσχηκότα τὸν λεγιῶνα, καὶ ἐφοβήθησαν.
\vs{16}καὶ διηγήσαντο αὐτοῖς οἱ ἰδόντες πῶς ἐγένετο τῷ δαιμονιζομένῳ καὶ περὶ τῶν χοίρων.
\vs{17}καὶ ἤρξαντο παρακαλεῖν αὐτὸν ἀπελθεῖν ἀπὸ τῶν ὁρίων αὐτῶν.

\vs{18}Καὶ ἐμβαίνοντος αὐτοῦ εἰς τὸ πλοῖον παρεκάλει αὐτὸν ὁ δαιμονισθεὶς ἵνα μετ᾽ αὐτοῦ ᾖ.
\vs{19}καὶ οὐκ ἀφῆκεν αὐτόν, ἀλλὰ λέγει αὐτῷ· Ὕπαγε εἰς τὸν οἶκόν σου πρὸς τοὺς σούς καὶ ἀπάγγειλον αὐτοῖς ὅσα ὁ Κύριός σοι πεποίηκεν καὶ ἠλέησέν σε.
\vs{20}Καὶ ἀπῆλθεν καὶ ἤρξατο κηρύσσειν ἐν τῇ Δεκαπόλει ὅσα ἐποίησεν αὐτῷ ὁ Ἰησοῦς, καὶ πάντες ἐθαύμαζον.

\vs{21}Καὶ διαπεράσαντος τοῦ Ἰησοῦ ἐν τῷ πλοίῳ πάλιν εἰς τὸ πέραν συνήχθη ὄχλος πολὺς ἐπ᾽ αὐτόν, καὶ ἦν παρὰ τὴν θάλασσαν.
\vs{22}καὶ ἔρχεται εἷς τῶν ἀρχισυναγώγων, ὀνόματι Ἰάϊρος, καὶ ἰδὼν αὐτὸν πίπτει πρὸς τοὺς πόδας αὐτοῦ
\vs{23}καὶ παρακαλεῖ αὐτὸν πολλὰ λέγων ὅτι Τὸ θυγάτριόν μου ἐσχάτως ἔχει, ἵνα ἐλθὼν ἐπιθῇς τὰς χεῖρας αὐτῇ ἵνα σωθῇ καὶ ζήσῃ.
\vs{24}Καὶ ἀπῆλθεν μετ᾽ αὐτοῦ. καὶ ἠκολούθει αὐτῷ ὄχλος πολύς καὶ συνέθλιβον αὐτόν.

\vs{25}Καὶ γυνὴ οὖσα ἐν ῥύσει αἵματος δώδεκα ἔτη
\vs{26}καὶ πολλὰ παθοῦσα ὑπὸ πολλῶν ἰατρῶν καὶ δαπανήσασα τὰ παρ᾽ αὐτῆς πάντα καὶ μηδὲν ὠφεληθεῖσα ἀλλὰ μᾶλλον εἰς τὸ χεῖρον ἐλθοῦσα,
\vs{27}ἀκούσασα περὶ τοῦ Ἰησοῦ, ἐλθοῦσα ἐν τῷ ὄχλῳ ὄπισθεν ἥψατο τοῦ ἱματίου αὐτοῦ·
\vs{28}ἔλεγεν γὰρ ὅτι Ἐὰν ἅψωμαι κἂν τῶν ἱματίων αὐτοῦ σωθήσομαι.
\vs{29}καὶ εὐθὺς ἐξηράνθη ἡ πηγὴ τοῦ αἵματος αὐτῆς καὶ ἔγνω τῷ σώματι ὅτι ἴαται ἀπὸ τῆς μάστιγος.
\vs{30}Καὶ εὐθὺς ὁ Ἰησοῦς ἐπιγνοὺς ἐν ἑαυτῷ τὴν ἐξ αὐτοῦ δύναμιν ἐξελθοῦσαν ἐπιστραφεὶς ἐν τῷ ὄχλῳ ἔλεγεν· Τίς μου ἥψατο τῶν ἱματίων;
\vs{31}Καὶ ἔλεγον αὐτῷ οἱ μαθηταὶ αὐτοῦ· Βλέπεις τὸν ὄχλον συνθλίβοντά σε καὶ λέγεις· Τίς μου ἥψατο;
\vs{32}Καὶ περιεβλέπετο ἰδεῖν τὴν τοῦτο ποιήσασαν.
\vs{33}ἡ δὲ γυνὴ φοβηθεῖσα καὶ τρέμουσα, εἰδυῖα ὃ γέγονεν αὐτῇ, ἦλθεν καὶ προσέπεσεν αὐτῷ καὶ εἶπεν αὐτῷ πᾶσαν τὴν ἀλήθειαν.
\vs{34}Ὁ δὲ εἶπεν αὐτῇ· Θυγάτηρ, ἡ πίστις σου σέσωκέν σε· ὕπαγε εἰς εἰρήνην καὶ ἴσθι ὑγιὴς ἀπὸ τῆς μάστιγός σου.

\vs{35}Ἔτι αὐτοῦ λαλοῦντος ἔρχονται ἀπὸ τοῦ ἀρχισυναγώγου λέγοντες ὅτι Ἡ θυγάτηρ σου ἀπέθανεν· τί ἔτι σκύλλεις τὸν διδάσκαλον;
\vs{36}Ὁ δὲ Ἰησοῦς παρακούσας τὸν λόγον λαλούμενον λέγει τῷ ἀρχισυναγώγῳ· Μὴ φοβοῦ, μόνον πίστευε.
\vs{37}καὶ οὐκ ἀφῆκεν οὐδένα μετ᾽ αὐτοῦ συνακολουθῆσαι εἰ μὴ τὸν Πέτρον καὶ Ἰάκωβον καὶ Ἰωάννην τὸν ἀδελφὸν Ἰακώβου.
\vs{38}Καὶ ἔρχονται εἰς τὸν οἶκον τοῦ ἀρχισυναγώγου, καὶ θεωρεῖ θόρυβον καὶ κλαίοντας καὶ ἀλαλάζοντας πολλά,
\vs{39}καὶ εἰσελθὼν λέγει αὐτοῖς· Τί θορυβεῖσθε καὶ κλαίετε; τὸ παιδίον οὐκ ἀπέθανεν ἀλλὰ καθεύδει.
\vs{40}καὶ κατεγέλων αὐτοῦ. Αὐτὸς δὲ ἐκβαλὼν πάντας παραλαμβάνει τὸν πατέρα τοῦ παιδίου καὶ τὴν μητέρα καὶ τοὺς μετ᾽ αὐτοῦ καὶ εἰσπορεύεται ὅπου ἦν τὸ παιδίον.
\vs{41}καὶ κρατήσας τῆς χειρὸς τοῦ παιδίου λέγει αὐτῇ· Ταλιθὰ κούμ, ὅ ἐστιν μεθερμηνευόμενον· Τὸ Κοράσιον, σοὶ λέγω, ἔγειρε.
\vs{42}καὶ εὐθὺς ἀνέστη τὸ κοράσιον καὶ περιεπάτει· ἦν γὰρ ἐτῶν δώδεκα. καὶ ἐξέστησαν εὐθὺς ἐκστάσει μεγάλῃ.
\vs{43}καὶ διεστείλατο αὐτοῖς πολλὰ ἵνα μηδεὶς γνοῖ τοῦτο, καὶ εἶπεν δοθῆναι αὐτῇ φαγεῖν.

\ch{6}
Καὶ ἐξῆλθεν ἐκεῖθεν καὶ ἔρχεται εἰς τὴν πατρίδα αὐτοῦ, καὶ ἀκολουθοῦσιν αὐτῷ οἱ μαθηταὶ αὐτοῦ.
\vs{2}καὶ γενομένου σαββάτου ἤρξατο διδάσκειν ἐν τῇ συναγωγῇ, καὶ πολλοὶ ἀκούοντες ἐξεπλήσσοντο λέγοντες· Πόθεν τούτῳ ταῦτα, καὶ τίς ἡ σοφία ἡ δοθεῖσα τούτῳ, καὶ αἱ δυνάμεις τοιαῦται διὰ τῶν χειρῶν αὐτοῦ γινόμεναι;
\vs{3}οὐχ οὗτός ἐστιν ὁ τέκτων, ὁ υἱὸς τῆς Μαρίας καὶ ἀδελφὸς Ἰακώβου καὶ Ἰωσῆτος καὶ Ἰούδα καὶ Σίμωνος; καὶ οὐκ εἰσὶν αἱ ἀδελφαὶ αὐτοῦ ὧδε πρὸς ἡμᾶς; καὶ ἐσκανδαλίζοντο ἐν αὐτῷ.
\vs{4}Καὶ ἔλεγεν αὐτοῖς ὁ Ἰησοῦς ὅτι Οὐκ ἔστιν προφήτης ἄτιμος εἰ μὴ ἐν τῇ πατρίδι αὐτοῦ καὶ ἐν τοῖς συγγενεῦσιν αὐτοῦ καὶ ἐν τῇ οἰκίᾳ αὐτοῦ.
\vs{5}καὶ οὐκ ἐδύνατο ἐκεῖ ποιῆσαι οὐδεμίαν δύναμιν, εἰ μὴ ὀλίγοις ἀρρώστοις ἐπιθεὶς τὰς χεῖρας ἐθεράπευσεν.
\vs{6}καὶ ἐθαύμαζεν διὰ τὴν ἀπιστίαν αὐτῶν. Καὶ περιῆγεν τὰς κώμας κύκλῳ διδάσκων.
\vs{7}Καὶ προσκαλεῖται τοὺς δώδεκα καὶ ἤρξατο αὐτοὺς ἀποστέλλειν δύο δύο καὶ ἐδίδου αὐτοῖς ἐξουσίαν τῶν πνευμάτων τῶν ἀκαθάρτων,
\vs{8}καὶ παρήγγειλεν αὐτοῖς ἵνα μηδὲν αἴρωσιν εἰς ὁδὸν εἰ μὴ ῥάβδον μόνον, μὴ ἄρτον, μὴ πήραν, μὴ εἰς τὴν ζώνην χαλκόν,
\vs{9}ἀλλὰ ὑποδεδεμένους σανδάλια, καὶ μὴ ἐνδύσησθε δύο χιτῶνας.
\vs{10}Καὶ ἔλεγεν αὐτοῖς· Ὅπου ἐὰν εἰσέλθητε εἰς οἰκίαν, ἐκεῖ μένετε ἕως ἂν ἐξέλθητε ἐκεῖθεν.
\vs{11}καὶ ὃς ἂν τόπος μὴ δέξηται ὑμᾶς μηδὲ ἀκούσωσιν ὑμῶν, ἐκπορευόμενοι ἐκεῖθεν ἐκτινάξατε τὸν χοῦν τὸν ὑποκάτω τῶν ποδῶν ὑμῶν εἰς μαρτύριον αὐτοῖς.
\vs{12}Καὶ ἐξελθόντες ἐκήρυξαν ἵνα μετανοῶσιν,
\vs{13}καὶ δαιμόνια πολλὰ ἐξέβαλλον, καὶ ἤλειφον ἐλαίῳ πολλοὺς ἀρρώστους καὶ ἐθεράπευον.

\vs{14}Καὶ ἤκουσεν ὁ βασιλεὺς Ἡρῴδης, φανερὸν γὰρ ἐγένετο τὸ ὄνομα αὐτοῦ, καὶ ἔλεγον ὅτι Ἰωάννης ὁ Βαπτίζων ἐγήγερται ἐκ νεκρῶν καὶ διὰ τοῦτο ἐνεργοῦσιν αἱ δυνάμεις ἐν αὐτῷ.
\vs{15}ἄλλοι δὲ ἔλεγον ὅτι Ἠλίας ἐστίν· ἄλλοι δὲ ἔλεγον ὅτι Προφήτης ὡς εἷς τῶν προφητῶν.
\vs{16}Ἀκούσας δὲ ὁ Ἡρῴδης ἔλεγεν· Ὃν ἐγὼ ἀπεκεφάλισα Ἰωάννην, οὗτος ἠγέρθη.

\vs{17}Αὐτὸς γὰρ ὁ Ἡρῴδης ἀποστείλας ἐκράτησεν τὸν Ἰωάννην καὶ ἔδησεν αὐτὸν ἐν φυλακῇ διὰ Ἡρῳδιάδα τὴν γυναῖκα Φιλίππου τοῦ ἀδελφοῦ αὐτοῦ, ὅτι αὐτὴν ἐγάμησεν·
\vs{18}ἔλεγεν γὰρ ὁ Ἰωάννης τῷ Ἡρῴδῃ ὅτι Οὐκ ἔξεστίν σοι ἔχειν τὴν γυναῖκα τοῦ ἀδελφοῦ σου.
\vs{19}Ἡ δὲ Ἡρῳδιὰς ἐνεῖχεν αὐτῷ καὶ ἤθελεν αὐτὸν ἀποκτεῖναι, καὶ οὐκ ἠδύνατο·
\vs{20}ὁ γὰρ Ἡρῴδης ἐφοβεῖτο τὸν Ἰωάννην, εἰδὼς αὐτὸν ἄνδρα δίκαιον καὶ ἅγιον, καὶ συνετήρει αὐτόν, καὶ ἀκούσας αὐτοῦ πολλὰ ἠπόρει, καὶ ἡδέως αὐτοῦ ἤκουεν.

\vs{21}Καὶ γενομένης ἡμέρας εὐκαίρου ὅτε Ἡρῴδης τοῖς γενεσίοις αὐτοῦ δεῖπνον ἐποίησεν τοῖς μεγιστᾶσιν αὐτοῦ καὶ τοῖς χιλιάρχοις καὶ τοῖς πρώτοις τῆς Γαλιλαίας,
\vs{22}καὶ εἰσελθούσης τῆς θυγατρὸς αὐτοῦ Ἡρῳδιάδος καὶ ὀρχησαμένης ἤρεσεν τῷ Ἡρῴδῃ καὶ τοῖς συνανακειμένοις. εἶπεν ὁ βασιλεὺς τῷ κορασίῳ· Αἴτησόν με ὃ ἐὰν θέλῃς, καὶ δώσω σοι·
\vs{23}καὶ ὤμοσεν αὐτῇ Πολλά Ὅ τι ἐάν με αἰτήσῃς δώσω σοι ἕως ἡμίσους τῆς βασιλείας μου.
\vs{24}Καὶ ἐξελθοῦσα εἶπεν τῇ μητρὶ αὐτῆς· Τί αἰτήσωμαι; Ἡ δὲ εἶπεν· Τὴν κεφαλὴν Ἰωάννου τοῦ Βαπτίζοντος.
\vs{25}Καὶ εἰσελθοῦσα εὐθὺς μετὰ σπουδῆς πρὸς τὸν βασιλέα ᾐτήσατο λέγουσα· Θέλω ἵνα ἐξαυτῆς δῷς μοι ἐπὶ πίνακι τὴν κεφαλὴν Ἰωάννου τοῦ Βαπτιστοῦ.
\vs{26}Καὶ περίλυπος γενόμενος ὁ βασιλεὺς διὰ τοὺς ὅρκους καὶ τοὺς ἀνακειμένους οὐκ ἠθέλησεν ἀθετῆσαι αὐτήν·
\vs{27}καὶ εὐθὺς ἀποστείλας ὁ βασιλεὺς σπεκουλάτορα ἐπέταξεν ἐνέγκαι τὴν κεφαλὴν αὐτοῦ. καὶ ἀπελθὼν ἀπεκεφάλισεν αὐτὸν ἐν τῇ φυλακῇ
\vs{28}καὶ ἤνεγκεν τὴν κεφαλὴν αὐτοῦ ἐπὶ πίνακι καὶ ἔδωκεν αὐτὴν τῷ κορασίῳ, καὶ τὸ κοράσιον ἔδωκεν αὐτὴν τῇ μητρὶ αὐτῆς.
\vs{29}καὶ ἀκούσαντες οἱ μαθηταὶ αὐτοῦ ἦλθον καὶ ἦραν τὸ πτῶμα αὐτοῦ καὶ ἔθηκαν αὐτὸ ἐν μνημείῳ.

\vs{30}Καὶ συνάγονται οἱ ἀπόστολοι πρὸς τὸν Ἰησοῦν καὶ ἀπήγγειλαν αὐτῷ πάντα ὅσα ἐποίησαν καὶ ὅσα ἐδίδαξαν.
\vs{31}καὶ λέγει αὐτοῖς· Δεῦτε ὑμεῖς αὐτοὶ κατ᾽ ἰδίαν εἰς ἔρημον τόπον καὶ ἀναπαύσασθε ὀλίγον. ἦσαν γὰρ οἱ ἐρχόμενοι καὶ οἱ ὑπάγοντες πολλοί, καὶ οὐδὲ φαγεῖν εὐκαίρουν.

\vs{32}Καὶ ἀπῆλθον ἐν τῷ πλοίῳ εἰς ἔρημον τόπον κατ᾽ ἰδίαν.
\vs{33}καὶ εἶδον αὐτοὺς ὑπάγοντας καὶ ἐπέγνωσαν πολλοί καὶ πεζῇ ἀπὸ πασῶν τῶν πόλεων συνέδραμον ἐκεῖ καὶ προῆλθον αὐτούς.

\vs{34}Καὶ ἐξελθὼν εἶδεν πολὺν ὄχλον καὶ ἐσπλαγχνίσθη ἐπ᾽ αὐτοὺς, ὅτι ἦσαν ὡς πρόβατα μὴ ἔχοντα ποιμένα, καὶ ἤρξατο διδάσκειν αὐτοὺς πολλά.

\vs{35}Καὶ ἤδη ὥρας πολλῆς γενομένης προσελθόντες αὐτῷ οἱ μαθηταὶ αὐτοῦ ἔλεγον ὅτι Ἔρημός ἐστιν ὁ τόπος καὶ ἤδη ὥρα πολλή·
\vs{36}ἀπόλυσον αὐτούς, ἵνα ἀπελθόντες εἰς τοὺς κύκλῳ ἀγροὺς καὶ κώμας ἀγοράσωσιν ἑαυτοῖς τί φάγωσιν.
\vs{37}Ὁ δὲ ἀποκριθεὶς εἶπεν αὐτοῖς· Δότε αὐτοῖς ὑμεῖς φαγεῖν. Καὶ λέγουσιν αὐτῷ· Ἀπελθόντες ἀγοράσωμεν δηναρίων διακοσίων ἄρτους καὶ δώσομεν αὐτοῖς φαγεῖν;
\vs{38}Ὁ δὲ λέγει αὐτοῖς· Πόσους ἄρτους ἔχετε; ὑπάγετε ἴδετε. Καὶ γνόντες λέγουσιν· Πέντε, καὶ δύο ἰχθύας.
\vs{39}Καὶ ἐπέταξεν αὐτοῖς ἀνακλῖναι πάντας συμπόσια συμπόσια ἐπὶ τῷ χλωρῷ χόρτῳ.
\vs{40}καὶ ἀνέπεσαν πρασιαὶ πρασιαὶ κατὰ ἑκατὸν καὶ κατὰ πεντήκοντα.
\vs{41}Καὶ λαβὼν τοὺς πέντε ἄρτους καὶ τοὺς δύο ἰχθύας ἀναβλέψας εἰς τὸν οὐρανὸν εὐλόγησεν καὶ κατέκλασεν τοὺς ἄρτους καὶ ἐδίδου τοῖς μαθηταῖς αὐτοῦ ἵνα παρατιθῶσιν αὐτοῖς, καὶ τοὺς δύο ἰχθύας ἐμέρισεν πᾶσιν.
\vs{42}Καὶ ἔφαγον πάντες καὶ ἐχορτάσθησαν,
\vs{43}καὶ ἦραν κλάσματα δώδεκα κοφίνων πληρώματα καὶ ἀπὸ τῶν ἰχθύων.
\vs{44}καὶ ἦσαν οἱ φαγόντες τοὺς ἄρτους πεντακισχίλιοι ἄνδρες.

\vs{45}Καὶ εὐθὺς ἠνάγκασεν τοὺς μαθητὰς αὐτοῦ ἐμβῆναι εἰς τὸ πλοῖον καὶ προάγειν εἰς τὸ πέραν πρὸς Βηθσαϊδάν, ἕως αὐτὸς ἀπολύει τὸν ὄχλον.
\vs{46}καὶ ἀποταξάμενος αὐτοῖς ἀπῆλθεν εἰς τὸ ὄρος προσεύξασθαι.
\vs{47}Καὶ ὀψίας γενομένης ἦν τὸ πλοῖον ἐν μέσῳ τῆς θαλάσσης, καὶ αὐτὸς μόνος ἐπὶ τῆς γῆς.
\vs{48}καὶ ἰδὼν αὐτοὺς βασανιζομένους ἐν τῷ ἐλαύνειν, ἦν γὰρ ὁ ἄνεμος ἐναντίος αὐτοῖς, περὶ τετάρτην φυλακὴν τῆς νυκτὸς ἔρχεται πρὸς αὐτοὺς περιπατῶν ἐπὶ τῆς θαλάσσης καὶ ἤθελεν παρελθεῖν αὐτούς.
\vs{49}οἱ δὲ ἰδόντες αὐτὸν ἐπὶ τῆς θαλάσσης περιπατοῦντα ἔδοξαν ὅτι φάντασμά ἐστιν, καὶ ἀνέκραξαν·
\vs{50}πάντες γὰρ αὐτὸν εἶδον καὶ ἐταράχθησαν. ὁ Δὲ εὐθὺς ἐλάλησεν μετ᾽ αὐτῶν, καὶ λέγει αὐτοῖς· Θαρσεῖτε, ἐγώ εἰμι· μὴ φοβεῖσθε.
\vs{51}καὶ ἀνέβη πρὸς αὐτοὺς εἰς τὸ πλοῖον καὶ ἐκόπασεν ὁ ἄνεμος, καὶ λίαν ἐκ περισσοῦ ἐν ἑαυτοῖς ἐξίσταντο·
\vs{52}οὐ γὰρ συνῆκαν ἐπὶ τοῖς ἄρτοις, ἀλλ᾽ ἦν αὐτῶν ἡ καρδία πεπωρωμένη.

\vs{53}Καὶ διαπεράσαντες ἐπὶ τὴν γῆν ἦλθον εἰς Γεννησαρὲτ καὶ προσωρμίσθησαν.
\vs{54}καὶ ἐξελθόντων αὐτῶν ἐκ τοῦ πλοίου εὐθὺς ἐπιγνόντες αὐτὸν
\vs{55}περιέδραμον ὅλην τὴν χώραν ἐκείνην καὶ ἤρξαντο ἐπὶ τοῖς κραβάττοις τοὺς κακῶς ἔχοντας περιφέρειν ὅπου ἤκουον ὅτι ἐστίν.
\vs{56}καὶ ὅπου ἂν εἰσεπορεύετο εἰς κώμας ἢ εἰς πόλεις ἢ εἰς ἀγροὺς, ἐν ταῖς ἀγοραῖς ἐτίθεσαν τοὺς ἀσθενοῦντας καὶ παρεκάλουν αὐτὸν ἵνα κἂν τοῦ κρασπέδου τοῦ ἱματίου αὐτοῦ ἅψωνται· καὶ ὅσοι ἂν ἥψαντο αὐτοῦ ἐσῴζοντο.

\ch{7}
Καὶ συνάγονται πρὸς αὐτὸν οἱ Φαρισαῖοι καί τινες τῶν γραμματέων ἐλθόντες ἀπὸ Ἱεροσολύμων.
\vs{2}καὶ ἰδόντες τινὰς τῶν μαθητῶν αὐτοῦ ὅτι κοιναῖς χερσίν, τοῦτ᾽ ἔστιν ἀνίπτοις, ἐσθίουσιν τοὺς ἄρτους—
\vs{3}Οἱ γὰρ Φαρισαῖοι καὶ πάντες οἱ Ἰουδαῖοι ἐὰν μὴ πυγμῇ νίψωνται τὰς χεῖρας οὐκ ἐσθίουσιν, κρατοῦντες τὴν παράδοσιν τῶν πρεσβυτέρων,
\vs{4}καὶ ἀπ᾽ ἀγορᾶς ἐὰν μὴ βαπτίσωνται οὐκ ἐσθίουσιν, καὶ ἄλλα πολλά ἐστιν ἃ παρέλαβον κρατεῖν, βαπτισμοὺς ποτηρίων καὶ ξεστῶν καὶ χαλκίων καὶ κλινῶν—
\vs{5}Καὶ ἐπερωτῶσιν αὐτὸν οἱ Φαρισαῖοι καὶ οἱ γραμματεῖς· Διὰ τί οὐ περιπατοῦσιν οἱ μαθηταί σου κατὰ τὴν παράδοσιν τῶν πρεσβυτέρων, ἀλλὰ κοιναῖς χερσὶν ἐσθίουσιν τὸν ἄρτον;

\vs{6}Ὁ δὲ εἶπεν αὐτοῖς· Καλῶς ἐπροφήτευσεν Ἠσαΐας περὶ ὑμῶν τῶν ὑποκριτῶν, ὡς γέγραπται ὅτι 
\begin{poetryblock}

\begin{quote}Οὗτος ὁ λαὸς τοῖς χείλεσίν με τιμᾷ,\end{quote} 

\begin{quote}Ἡ δὲ καρδία αὐτῶν πόρρω ἀπέχει ἀπ᾽ ἐμοῦ·\end{quote}

\begin{quote} \vs{7}Μάτην δὲ σέβονταί με\end{quote} 

\begin{quote}Διδάσκοντες διδασκαλίας ἐντάλματα ἀνθρώπων.\end{quote}
\end{poetryblock}

\vs{8}Ἀφέντες τὴν ἐντολὴν τοῦ Θεοῦ κρατεῖτε τὴν παράδοσιν τῶν ἀνθρώπων.
\vs{9}Καὶ ἔλεγεν αὐτοῖς· Καλῶς ἀθετεῖτε τὴν ἐντολὴν τοῦ Θεοῦ, ἵνα τὴν παράδοσιν ὑμῶν τηρήσητε.
\vs{10}Μωϋσῆς γὰρ εἶπεν· Τίμα τὸν πατέρα σου καὶ τὴν μητέρα σου, καί· Ὁ κακολογῶν πατέρα ἢ μητέρα θανάτῳ τελευτάτω.
\vs{11}ὑμεῖς δὲ λέγετε· Ἐὰν εἴπῃ ἄνθρωπος τῷ πατρὶ ἢ τῇ μητρί· Κορβᾶν, ὅ ἐστιν Δῶρον, Ὃ ἐὰν ἐξ ἐμοῦ ὠφεληθῇς,
\vs{12}οὐκέτι ἀφίετε αὐτὸν οὐδὲν ποιῆσαι τῷ πατρὶ ἢ τῇ μητρί,
\vs{13}ἀκυροῦντες τὸν λόγον τοῦ Θεοῦ τῇ παραδόσει ὑμῶν ᾗ παρεδώκατε· καὶ παρόμοια τοιαῦτα πολλὰ ποιεῖτε.

\vs{14}Καὶ προσκαλεσάμενος πάλιν τὸν ὄχλον ἔλεγεν αὐτοῖς· Ἀκούσατέ μου πάντες καὶ σύνετε.
\vs{15}οὐδέν ἐστιν ἔξωθεν τοῦ ἀνθρώπου εἰσπορευόμενον εἰς αὐτὸν ὃ δύναται κοινῶσαι αὐτόν, ἀλλὰ τὰ ἐκ τοῦ ἀνθρώπου ἐκπορευόμενά ἐστιν τὰ κοινοῦντα τὸν ἄνθρωπον.
\vs{17}Καὶ ὅτε εἰσῆλθεν εἰς οἶκον ἀπὸ τοῦ ὄχλου, ἐπηρώτων αὐτὸν οἱ μαθηταὶ αὐτοῦ τὴν παραβολήν.
\vs{18}Καὶ λέγει αὐτοῖς· Οὕτως καὶ ὑμεῖς ἀσύνετοί ἐστε; οὐ νοεῖτε ὅτι πᾶν τὸ ἔξωθεν εἰσπορευόμενον εἰς τὸν ἄνθρωπον οὐ δύναται αὐτὸν κοινῶσαι
\vs{19}ὅτι οὐκ εἰσπορεύεται αὐτοῦ εἰς τὴν καρδίαν ἀλλ᾽ εἰς τὴν κοιλίαν, καὶ εἰς τὸν ἀφεδρῶνα ἐκπορεύεται, καθαρίζων πάντα τὰ βρώματα;
\vs{20}Ἔλεγεν δὲ ὅτι Τὸ ἐκ τοῦ ἀνθρώπου ἐκπορευόμενον, ἐκεῖνο κοινοῖ τὸν ἄνθρωπον.
\vs{21}ἔσωθεν γὰρ ἐκ τῆς καρδίας τῶν ἀνθρώπων οἱ διαλογισμοὶ οἱ κακοὶ ἐκπορεύονται, πορνεῖαι, κλοπαί, φόνοι,
\vs{22}μοιχεῖαι, πλεονεξίαι, πονηρίαι, δόλος, ἀσέλγεια, ὀφθαλμὸς πονηρός, βλασφημία, ὑπερηφανία, ἀφροσύνη·
\vs{23}πάντα ταῦτα τὰ πονηρὰ ἔσωθεν ἐκπορεύεται καὶ κοινοῖ τὸν ἄνθρωπον.

\vs{24}Ἐκεῖθεν δὲ ἀναστὰς ἀπῆλθεν εἰς τὰ ὅρια Τύρου. Καὶ εἰσελθὼν εἰς οἰκίαν οὐδένα ἤθελεν γνῶναι, καὶ οὐκ ἠδυνήθη λαθεῖν·
\vs{25}ἀλλ᾽ εὐθὺς ἀκούσασα γυνὴ περὶ αὐτοῦ, ἧς εἶχεν τὸ θυγάτριον αὐτῆς πνεῦμα ἀκάθαρτον, ἐλθοῦσα προσέπεσεν πρὸς τοὺς πόδας αὐτοῦ·
\vs{26}ἡ δὲ γυνὴ ἦν Ἑλληνίς, Συροφοινίκισσα τῷ γένει· καὶ ἠρώτα αὐτὸν ἵνα τὸ δαιμόνιον ἐκβάλῃ ἐκ τῆς θυγατρὸς αὐτῆς.
\vs{27}Καὶ ἔλεγεν αὐτῇ· Ἄφες πρῶτον χορτασθῆναι τὰ τέκνα, οὐ γάρ ἐστιν καλόν λαβεῖν τὸν ἄρτον τῶν τέκνων καὶ τοῖς κυναρίοις βαλεῖν.
\vs{28}Ἡ δὲ ἀπεκρίθη καὶ λέγει αὐτῷ· Κύριε· καὶ τὰ κυνάρια ὑποκάτω τῆς τραπέζης ἐσθίουσιν ἀπὸ τῶν ψιχίων τῶν παιδίων.
\vs{29}Καὶ εἶπεν αὐτῇ· Διὰ τοῦτον τὸν λόγον ὕπαγε, ἐξελήλυθεν ἐκ τῆς θυγατρός σου τὸ δαιμόνιον.
\vs{30}καὶ ἀπελθοῦσα εἰς τὸν οἶκον αὐτῆς εὗρεν τὸ παιδίον βεβλημένον ἐπὶ τὴν κλίνην καὶ τὸ δαιμόνιον ἐξεληλυθός.

\vs{31}Καὶ πάλιν ἐξελθὼν ἐκ τῶν ὁρίων Τύρου ἦλθεν διὰ Σιδῶνος εἰς τὴν θάλασσαν τῆς Γαλιλαίας ἀνὰ μέσον τῶν ὁρίων Δεκαπόλεως.
\vs{32}Καὶ φέρουσιν αὐτῷ κωφὸν καὶ μογιλάλον καὶ παρακαλοῦσιν αὐτὸν ἵνα ἐπιθῇ αὐτῷ τὴν χεῖρα.
\vs{33}Καὶ ἀπολαβόμενος αὐτὸν ἀπὸ τοῦ ὄχλου κατ᾽ ἰδίαν ἔβαλεν τοὺς δακτύλους αὐτοῦ εἰς τὰ ὦτα αὐτοῦ καὶ πτύσας ἥψατο τῆς γλώσσης αὐτοῦ,
\vs{34}καὶ ἀναβλέψας εἰς τὸν οὐρανὸν ἐστέναξεν καὶ λέγει αὐτῷ· Ἐφφαθά, ὅ ἐστιν Διανοίχθητι.
\vs{35}καὶ εὐθὺς ἠνοίγησαν αὐτοῦ αἱ ἀκοαί, καὶ ἐλύθη ὁ δεσμὸς τῆς γλώσσης αὐτοῦ καὶ ἐλάλει ὀρθῶς.
\vs{36}Καὶ διεστείλατο αὐτοῖς ἵνα μηδενὶ λέγωσιν· ὅσον δὲ αὐτοῖς διεστέλλετο, αὐτοὶ μᾶλλον περισσότερον ἐκήρυσσον.
\vs{37}καὶ ὑπερπερισσῶς ἐξεπλήσσοντο λέγοντες· Καλῶς πάντα πεποίηκεν, καὶ τοὺς κωφοὺς ποιεῖ ἀκούειν καὶ τοὺς ἀλάλους λαλεῖν.

\ch{8}
Ἐν ἐκείναις ταῖς ἡμέραις πάλιν πολλοῦ ὄχλου ὄντος καὶ μὴ ἐχόντων τί φάγωσιν, προσκαλεσάμενος τοὺς μαθητὰς λέγει αὐτοῖς·
\vs{2}Σπλαγχνίζομαι ἐπὶ τὸν ὄχλον, ὅτι ἤδη ἡμέραι τρεῖς προσμένουσίν μοι καὶ οὐκ ἔχουσιν τί φάγωσιν·
\vs{3}καὶ ἐὰν ἀπολύσω αὐτοὺς νήστεις εἰς οἶκον αὐτῶν, ἐκλυθήσονται ἐν τῇ ὁδῷ· καί τινες αὐτῶν ἀπὸ μακρόθεν ἥκασιν.
\vs{4}Καὶ ἀπεκρίθησαν αὐτῷ οἱ μαθηταὶ αὐτοῦ ὅτι Πόθεν τούτους δυνήσεταί τις ὧδε χορτάσαι ἄρτων ἐπ᾽ ἐρημίας;
\vs{5}Καὶ ἠρώτα αὐτούς· Πόσους ἔχετε ἄρτους; Οἱ δὲ εἶπαν· Ἑπτά.
\vs{6}Καὶ παραγγέλλει τῷ ὄχλῳ ἀναπεσεῖν ἐπὶ τῆς γῆς· καὶ λαβὼν τοὺς ἑπτὰ ἄρτους εὐχαριστήσας ἔκλασεν καὶ ἐδίδου τοῖς μαθηταῖς αὐτοῦ ἵνα παρατιθῶσιν, καὶ παρέθηκαν τῷ ὄχλῳ.
\vs{7}καὶ εἶχον ἰχθύδια ὀλίγα· καὶ εὐλογήσας αὐτὰ εἶπεν καὶ ταῦτα παρατιθέναι.
\vs{8}Καὶ ἔφαγον καὶ ἐχορτάσθησαν, καὶ ἦραν περισσεύματα κλασμάτων ἑπτὰ σπυρίδας.
\vs{9}ἦσαν δὲ ὡς τετρακισχίλιοι. καὶ ἀπέλυσεν αὐτούς.

\vs{10}Καὶ εὐθὺς ἐμβὰς εἰς τὸ πλοῖον μετὰ τῶν μαθητῶν αὐτοῦ ἦλθεν εἰς τὰ μέρη Δαλμανουθά.

\vs{11}Καὶ ἐξῆλθον οἱ Φαρισαῖοι καὶ ἤρξαντο συζητεῖν αὐτῷ, ζητοῦντες παρ᾽ αὐτοῦ σημεῖον ἀπὸ τοῦ οὐρανοῦ, πειράζοντες αὐτόν.
\vs{12}Καὶ ἀναστενάξας τῷ πνεύματι αὐτοῦ λέγει· Τί ἡ γενεὰ αὕτη ζητεῖ σημεῖον; ἀμὴν λέγω ὑμῖν, εἰ δοθήσεται τῇ γενεᾷ ταύτῃ σημεῖον.
\vs{13}καὶ ἀφεὶς αὐτοὺς πάλιν ἐμβὰς ἀπῆλθεν εἰς τὸ πέραν.

\vs{14}Καὶ ἐπελάθοντο λαβεῖν ἄρτους καὶ εἰ μὴ ἕνα ἄρτον οὐκ εἶχον μεθ᾽ ἑαυτῶν ἐν τῷ πλοίῳ.
\vs{15}καὶ διεστέλλετο αὐτοῖς λέγων· Ὁρᾶτε, βλέπετε ἀπὸ τῆς ζύμης τῶν Φαρισαίων καὶ τῆς ζύμης Ἡρῴδου.
\vs{16}καὶ διελογίζοντο πρὸς ἀλλήλους ὅτι ἄρτους οὐκ ἔχουσιν.
\vs{17}Καὶ γνοὺς λέγει αὐτοῖς· Τί διαλογίζεσθε ὅτι ἄρτους οὐκ ἔχετε; οὔπω νοεῖτε οὐδὲ συνίετε; πεπωρωμένην ἔχετε τὴν καρδίαν ὑμῶν;
\begin{poetryblock}

\begin{quote} \vs{18}Ὀφθαλμοὺς ἔχοντες οὐ βλέπετε\end{quote} 

\begin{quote}καὶ ὦτα ἔχοντες οὐκ ἀκούετε;\end{quote}
\end{poetryblock}

καὶ οὐ μνημονεύετε,

\vs{19}ὅτε τοὺς πέντε ἄρτους ἔκλασα εἰς τοὺς πεντακισχιλίους, πόσους κοφίνους κλασμάτων πλήρεις ἤρατε; Λέγουσιν αὐτῷ· Δώδεκα.
\vs{20}Ὅτε τοὺς ἑπτὰ εἰς τοὺς τετρακισχιλίους, πόσων σπυρίδων πληρώματα κλασμάτων ἤρατε; Καὶ λέγουσιν αὐτῷ· Ἑπτά.
\vs{21}Καὶ ἔλεγεν αὐτοῖς· Οὔπω συνίετε;

\vs{22}Καὶ ἔρχονται εἰς Βηθσαϊδάν. Καὶ φέρουσιν αὐτῷ τυφλὸν καὶ παρακαλοῦσιν αὐτὸν ἵνα αὐτοῦ ἅψηται.
\vs{23}καὶ ἐπιλαβόμενος τῆς χειρὸς τοῦ τυφλοῦ ἐξήνεγκεν αὐτὸν ἔξω τῆς κώμης καὶ πτύσας εἰς τὰ ὄμματα αὐτοῦ, ἐπιθεὶς τὰς χεῖρας αὐτῷ ἐπηρώτα αὐτόν· Εἴ τι βλέπεις;
\vs{24}Καὶ ἀναβλέψας ἔλεγεν· Βλέπω τοὺς ἀνθρώπους ὅτι ὡς δένδρα ὁρῶ περιπατοῦντας.
\vs{25}Εἶτα πάλιν ἐπέθηκεν τὰς χεῖρας ἐπὶ τοὺς ὀφθαλμοὺς αὐτοῦ, καὶ διέβλεψεν καὶ ἀπεκατέστη καὶ ἐνέβλεπεν τηλαυγῶς ἅπαντα.
\vs{26}καὶ ἀπέστειλεν αὐτὸν εἰς οἶκον αὐτοῦ λέγων· Μηδὲ εἰς τὴν κώμην εἰσέλθῃς.

\vs{27}Καὶ ἐξῆλθεν ὁ Ἰησοῦς καὶ οἱ μαθηταὶ αὐτοῦ εἰς τὰς κώμας Καισαρείας τῆς Φιλίππου· καὶ ἐν τῇ ὁδῷ ἐπηρώτα τοὺς μαθητὰς αὐτοῦ λέγων αὐτοῖς· Τίνα με λέγουσιν οἱ ἄνθρωποι εἶναι;
\vs{28}Οἱ δὲ εἶπαν αὐτῷ λέγοντες ὅτι Ἰωάννην τὸν Βαπτιστήν, καὶ ἄλλοι Ἠλίαν, ἄλλοι δὲ ὅτι εἷς τῶν προφητῶν.
\vs{29}Καὶ αὐτὸς ἐπηρώτα αὐτούς· Ὑμεῖς δὲ τίνα με λέγετε εἶναι; Ἀποκριθεὶς ὁ Πέτρος λέγει αὐτῷ· Σὺ εἶ ὁ Χριστός.
\vs{30}Καὶ ἐπετίμησεν αὐτοῖς ἵνα μηδενὶ λέγωσιν περὶ αὐτοῦ.

\vs{31}Καὶ ἤρξατο διδάσκειν αὐτοὺς ὅτι δεῖ τὸν Υἱὸν τοῦ ἀνθρώπου πολλὰ παθεῖν καὶ ἀποδοκιμασθῆναι ὑπὸ τῶν πρεσβυτέρων καὶ τῶν ἀρχιερέων καὶ τῶν γραμματέων καὶ ἀποκτανθῆναι καὶ μετὰ τρεῖς ἡμέρας ἀναστῆναι·
\vs{32}καὶ παρρησίᾳ τὸν λόγον ἐλάλει. καὶ προσλαβόμενος ὁ Πέτρος αὐτὸν ἤρξατο ἐπιτιμᾶν αὐτῷ.
\vs{33}Ὁ δὲ ἐπιστραφεὶς καὶ ἰδὼν τοὺς μαθητὰς αὐτοῦ ἐπετίμησεν Πέτρῳ καὶ λέγει· Ὕπαγε ὀπίσω μου, Σατανᾶ, ὅτι οὐ φρονεῖς τὰ τοῦ Θεοῦ ἀλλὰ τὰ τῶν ἀνθρώπων.

\vs{34}Καὶ προσκαλεσάμενος τὸν ὄχλον σὺν τοῖς μαθηταῖς αὐτοῦ εἶπεν αὐτοῖς· Εἴ τις θέλει ὀπίσω μου ἀκολουθεῖν, ἀπαρνησάσθω ἑαυτὸν καὶ ἀράτω τὸν σταυρὸν αὐτοῦ καὶ ἀκολουθείτω μοι.
\vs{35}ὃς γὰρ ἐὰν θέλῃ τὴν ψυχὴν αὐτοῦ σῶσαι ἀπολέσει αὐτήν· ὃς δ᾽ ἂν ἀπολέσει τὴν ψυχὴν αὐτοῦ ἕνεκεν ἐμοῦ καὶ τοῦ εὐαγγελίου σώσει αὐτήν.
\vs{36}Τί γὰρ ὠφελεῖ ἄνθρωπον κερδῆσαι τὸν κόσμον ὅλον καὶ ζημιωθῆναι τὴν ψυχὴν αὐτοῦ;
\vs{37}τί γὰρ δοῖ ἄνθρωπος ἀντάλλαγμα τῆς ψυχῆς αὐτοῦ;
\vs{38}ὃς γὰρ ἐὰν ἐπαισχυνθῇ με καὶ τοὺς ἐμοὺς λόγους ἐν τῇ γενεᾷ ταύτῃ τῇ μοιχαλίδι καὶ ἁμαρτωλῷ, καὶ ὁ Υἱὸς τοῦ ἀνθρώπου ἐπαισχυνθήσεται αὐτὸν, ὅταν ἔλθῃ ἐν τῇ δόξῃ τοῦ Πατρὸς αὐτοῦ μετὰ τῶν ἀγγέλων τῶν ἁγίων.

\ch{9}
Καὶ ἔλεγεν αὐτοῖς· Ἀμὴν λέγω ὑμῖν ὅτι εἰσίν τινες ὧδε τῶν ἑστηκότων οἵτινες οὐ μὴ γεύσωνται θανάτου ἕως ἂν ἴδωσιν τὴν βασιλείαν τοῦ Θεοῦ ἐληλυθυῖαν ἐν δυνάμει.

\vs{2}Καὶ μετὰ ἡμέρας ἓξ παραλαμβάνει ὁ Ἰησοῦς τὸν Πέτρον καὶ τὸν Ἰάκωβον καὶ τὸν Ἰωάννην καὶ ἀναφέρει αὐτοὺς εἰς ὄρος ὑψηλὸν κατ᾽ ἰδίαν μόνους. καὶ μετεμορφώθη ἔμπροσθεν αὐτῶν,
\vs{3}καὶ τὰ ἱμάτια αὐτοῦ ἐγένετο στίλβοντα λευκὰ λίαν, οἷα γναφεὺς ἐπὶ τῆς γῆς οὐ δύναται οὕτως λευκᾶναι.
\vs{4}καὶ ὤφθη αὐτοῖς Ἠλίας σὺν Μωϋσεῖ καὶ ἦσαν συλλαλοῦντες τῷ Ἰησοῦ.
\vs{5}Καὶ ἀποκριθεὶς ὁ Πέτρος λέγει τῷ Ἰησοῦ· Ῥαββί, καλόν ἐστιν ἡμᾶς ὧδε εἶναι, καὶ ποιήσωμεν τρεῖς σκηνάς, σοὶ μίαν καὶ Μωϋσεῖ μίαν καὶ Ἠλίᾳ μίαν.
\vs{6}οὐ γὰρ ᾔδει τί ἀποκριθῇ, ἔκφοβοι γὰρ ἐγένοντο.
\vs{7}Καὶ ἐγένετο νεφέλη ἐπισκιάζουσα αὐτοῖς, καὶ ἐγένετο φωνὴ ἐκ τῆς νεφέλης· Οὗτός ἐστιν ὁ Υἱός μου ὁ ἀγαπητός, ἀκούετε αὐτοῦ.
\vs{8}καὶ ἐξάπινα περιβλεψάμενοι οὐκέτι οὐδένα εἶδον ἀλλὰ τὸν Ἰησοῦν μόνον μεθ᾽ ἑαυτῶν.

\vs{9}Καὶ καταβαινόντων αὐτῶν ἐκ τοῦ ὄρους διεστείλατο αὐτοῖς ἵνα μηδενὶ ἃ εἶδον διηγήσωνται, εἰ μὴ ὅταν ὁ Υἱὸς τοῦ ἀνθρώπου ἐκ νεκρῶν ἀναστῇ.
\vs{10}καὶ τὸν λόγον ἐκράτησαν πρὸς ἑαυτοὺς συζητοῦντες τί ἐστιν τὸ ἐκ νεκρῶν ἀναστῆναι.

\vs{11}καὶ ἐπηρώτων αὐτὸν λέγοντες· Ὅτι Λέγουσιν οἱ γραμματεῖς ὅτι Ἠλίαν δεῖ ἐλθεῖν πρῶτον;
\vs{12}Ὁ δὲ ἔφη αὐτοῖς· Ἠλίας μὲν ἐλθὼν πρῶτον ἀποκαθιστάνει πάντα· καὶ πῶς γέγραπται ἐπὶ τὸν Υἱὸν τοῦ ἀνθρώπου ἵνα πολλὰ πάθῃ καὶ ἐξουδενηθῇ;
\vs{13}ἀλλὰ λέγω ὑμῖν ὅτι καὶ Ἠλίας ἐλήλυθεν, καὶ ἐποίησαν αὐτῷ ὅσα ἤθελον, καθὼς γέγραπται ἐπ᾽ αὐτόν.

\vs{14}Καὶ ἐλθόντες πρὸς τοὺς μαθητὰς εἶδον ὄχλον πολὺν περὶ αὐτοὺς καὶ γραμματεῖς συζητοῦντας πρὸς αὐτούς.
\vs{15}καὶ εὐθὺς πᾶς ὁ ὄχλος ἰδόντες αὐτὸν ἐξεθαμβήθησαν καὶ προστρέχοντες ἠσπάζοντο αὐτόν.
\vs{16}Καὶ ἐπηρώτησεν αὐτούς· Τί συζητεῖτε πρὸς αὑτούς;
\vs{17}Καὶ ἀπεκρίθη αὐτῷ εἷς ἐκ τοῦ ὄχλου· Διδάσκαλε, ἤνεγκα τὸν υἱόν μου πρὸς σέ, ἔχοντα πνεῦμα ἄλαλον·
\vs{18}καὶ ὅπου ἐὰν αὐτὸν καταλάβῃ ῥήσσει αὐτόν, καὶ ἀφρίζει καὶ τρίζει τοὺς ὀδόντας καὶ ξηραίνεται· καὶ εἶπα τοῖς μαθηταῖς σου ἵνα αὐτὸ ἐκβάλωσιν, καὶ οὐκ ἴσχυσαν.
\vs{19}Ὁ δὲ ἀποκριθεὶς αὐτοῖς λέγει· Ὦ γενεὰ ἄπιστος, ἕως πότε πρὸς ὑμᾶς ἔσομαι; ἕως πότε ἀνέξομαι ὑμῶν; φέρετε αὐτὸν πρός με.
\vs{20}Καὶ ἤνεγκαν αὐτὸν πρὸς αὐτόν. καὶ ἰδὼν αὐτὸν τὸ πνεῦμα εὐθὺς συνεσπάραξεν αὐτόν, καὶ πεσὼν ἐπὶ τῆς γῆς ἐκυλίετο ἀφρίζων.
\vs{21}Καὶ ἐπηρώτησεν τὸν πατέρα αὐτοῦ· Πόσος χρόνος ἐστὶν ὡς τοῦτο γέγονεν αὐτῷ; Ὁ δὲ εἶπεν· Ἐκ παιδιόθεν·
\vs{22}καὶ πολλάκις καὶ εἰς πῦρ αὐτὸν ἔβαλεν καὶ εἰς ὕδατα ἵνα ἀπολέσῃ αὐτόν· ἀλλ᾽ εἴ τι δύνῃ, βοήθησον ἡμῖν σπλαγχνισθεὶς ἐφ᾽ ἡμᾶς.
\vs{23}Ὁ δὲ Ἰησοῦς εἶπεν αὐτῷ· Τὸ Εἰ δύνῃ, πάντα δυνατὰ τῷ πιστεύοντι.
\vs{24}Εὐθὺς κράξας ὁ πατὴρ τοῦ παιδίου ἔλεγεν· Πιστεύω· βοήθει μου τῇ ἀπιστίᾳ.
\vs{25}Ἰδὼν δὲ ὁ Ἰησοῦς ὅτι ἐπισυντρέχει ὄχλος, ἐπετίμησεν τῷ πνεύματι τῷ ἀκαθάρτῳ λέγων αὐτῷ· Τὸ ἄλαλον καὶ κωφὸν πνεῦμα, ἐγὼ ἐπιτάσσω σοι, ἔξελθε ἐξ αὐτοῦ καὶ μηκέτι εἰσέλθῃς εἰς αὐτόν.
\vs{26}Καὶ κράξας καὶ πολλὰ σπαράξας ἐξῆλθεν· καὶ ἐγένετο ὡσεὶ νεκρὸς, ὥστε τοὺς πολλοὺς λέγειν ὅτι ἀπέθανεν.
\vs{27}ὁ δὲ Ἰησοῦς κρατήσας τῆς χειρὸς αὐτοῦ ἤγειρεν αὐτόν, καὶ ἀνέστη.

\vs{28}Καὶ εἰσελθόντος αὐτοῦ εἰς οἶκον οἱ μαθηταὶ αὐτοῦ κατ᾽ ἰδίαν ἐπηρώτων αὐτόν· Ὅτι ἡμεῖς οὐκ ἠδυνήθημεν ἐκβαλεῖν αὐτό;
\vs{29}Καὶ εἶπεν αὐτοῖς· Τοῦτο τὸ γένος ἐν οὐδενὶ δύναται ἐξελθεῖν εἰ μὴ ἐν προσευχῇ.
\vs{30}Κἀκεῖθεν ἐξελθόντες παρεπορεύοντο διὰ τῆς Γαλιλαίας, καὶ οὐκ ἤθελεν ἵνα τις γνοῖ·
\vs{31}ἐδίδασκεν γὰρ τοὺς μαθητὰς αὐτοῦ καὶ ἔλεγεν αὐτοῖς ὅτι Ὁ Υἱὸς τοῦ ἀνθρώπου παραδίδοται εἰς χεῖρας ἀνθρώπων, καὶ ἀποκτενοῦσιν αὐτόν, καὶ ἀποκτανθεὶς μετὰ τρεῖς ἡμέρας ἀναστήσεται.
\vs{32}οἱ δὲ ἠγνόουν τὸ ῥῆμα, καὶ ἐφοβοῦντο αὐτὸν ἐπερωτῆσαι.

\vs{33}Καὶ ἦλθον εἰς Καφαρναούμ. Καὶ ἐν τῇ οἰκίᾳ γενόμενος ἐπηρώτα αὐτούς· Τί ἐν τῇ ὁδῷ διελογίζεσθε;
\vs{34}οἱ δὲ ἐσιώπων· πρὸς ἀλλήλους γὰρ διελέχθησαν ἐν τῇ ὁδῷ τίς μείζων.
\vs{35}Καὶ καθίσας ἐφώνησεν τοὺς δώδεκα καὶ λέγει αὐτοῖς· Εἴ τις θέλει πρῶτος εἶναι, ἔσται πάντων ἔσχατος καὶ πάντων διάκονος.
\vs{36}Καὶ λαβὼν παιδίον ἔστησεν αὐτὸ ἐν μέσῳ αὐτῶν καὶ ἐναγκαλισάμενος αὐτὸ εἶπεν αὐτοῖς·
\vs{37}Ὃς ἂν ἓν τῶν τοιούτων παιδίων δέξηται ἐπὶ τῷ ὀνόματί μου, ἐμὲ δέχεται· καὶ ὃς ἂν ἐμὲ δέχηται, οὐκ ἐμὲ δέχεται ἀλλὰ τὸν ἀποστείλαντά με.

\vs{38}Ἔφη αὐτῷ ὁ Ἰωάννης· Διδάσκαλε, εἴδομέν τινα ἐν τῷ ὀνόματί σου ἐκβάλλοντα δαιμόνια καὶ ἐκωλύομεν αὐτόν, ὅτι οὐκ ἠκολούθει ἡμῖν.
\vs{39}Ὁ δὲ Ἰησοῦς εἶπεν· Μὴ κωλύετε αὐτόν. οὐδεὶς γάρ ἐστιν ὃς ποιήσει δύναμιν ἐπὶ τῷ ὀνόματί μου καὶ δυνήσεται ταχὺ κακολογῆσαί με·
\vs{40}ὃς γὰρ οὐκ ἔστιν καθ᾽ ἡμῶν, ὑπὲρ ἡμῶν ἐστιν.

\vs{41}Ὃς γὰρ ἂν ποτίσῃ ὑμᾶς ποτήριον ὕδατος ἐν ὀνόματι ὅτι Χριστοῦ ἐστε, ἀμὴν λέγω ὑμῖν ὅτι οὐ μὴ ἀπολέσῃ τὸν μισθὸν αὐτοῦ.

\vs{42}Καὶ ὃς ἂν σκανδαλίσῃ ἕνα τῶν μικρῶν τούτων τῶν πιστευόντων εἰς ἐμέ, καλόν ἐστιν αὐτῷ μᾶλλον εἰ περίκειται μύλος ὀνικὸς περὶ τὸν τράχηλον αὐτοῦ καὶ βέβληται εἰς τὴν θάλασσαν.
\vs{43}Καὶ ἐὰν σκανδαλίζῃ σε ἡ χείρ σου, ἀπόκοψον αὐτήν· καλόν ἐστίν σε κυλλὸν εἰσελθεῖν εἰς τὴν ζωὴν ἢ τὰς δύο χεῖρας ἔχοντα ἀπελθεῖν εἰς τὴν γέενναν, εἰς τὸ πῦρ τὸ ἄσβεστον.
\vs{45}καὶ ἐὰν ὁ πούς σου σκανδαλίζῃ σε, ἀπόκοψον αὐτόν· καλόν ἐστίν σε εἰσελθεῖν εἰς τὴν ζωὴν χωλὸν ἢ τοὺς δύο πόδας ἔχοντα βληθῆναι εἰς τὴν γέενναν.
\vs{47}καὶ ἐὰν ὁ ὀφθαλμός σου σκανδαλίζῃ σε, ἔκβαλε αὐτόν· καλόν σέ ἐστιν μονόφθαλμον εἰσελθεῖν εἰς τὴν βασιλείαν τοῦ Θεοῦ ἢ δύο ὀφθαλμοὺς ἔχοντα βληθῆναι εἰς τὴν γέενναν,
\vs{48}ὅπου Ὁ σκώληξ αὐτῶν οὐ τελευτᾷ καὶ τὸ πῦρ οὐ σβέννυται.

\vs{49}Πᾶς γὰρ πυρὶ ἁλισθήσεται.
\vs{50}καλὸν τὸ ἅλας· ἐὰν δὲ τὸ ἅλας ἄναλον γένηται, ἐν τίνι αὐτὸ ἀρτύσετε; ἔχετε ἐν ἑαυτοῖς ἅλα καὶ εἰρηνεύετε ἐν ἀλλήλοις.

\ch{10}
Καὶ ἐκεῖθεν ἀναστὰς ἔρχεται εἰς τὰ ὅρια τῆς Ἰουδαίας καὶ πέραν τοῦ Ἰορδάνου, καὶ συμπορεύονται πάλιν ὄχλοι πρὸς αὐτόν, καὶ ὡς εἰώθει πάλιν ἐδίδασκεν αὐτούς.

\vs{2}Καὶ προσελθόντες Φαρισαῖοι ἐπηρώτων αὐτὸν εἰ ἔξεστιν ἀνδρὶ γυναῖκα ἀπολῦσαι, πειράζοντες αὐτόν.
\vs{3}Ὁ δὲ ἀποκριθεὶς εἶπεν αὐτοῖς· Τί ὑμῖν ἐνετείλατο Μωϋσῆς;
\vs{4}Οἱ δὲ εἶπαν· Ἐπέτρεψεν Μωϋσῆς βιβλίον ἀποστασίου γράψαι καὶ ἀπολῦσαι.
\vs{5}Ὁ δὲ Ἰησοῦς εἶπεν αὐτοῖς· Πρὸς τὴν σκληροκαρδίαν ὑμῶν ἔγραψεν ὑμῖν τὴν ἐντολὴν ταύτην.
\vs{6}ἀπὸ δὲ ἀρχῆς κτίσεως Ἄρσεν καὶ θῆλυ ἐποίησεν αὐτούς·
\vs{7}Ἕνεκεν τούτου καταλείψει ἄνθρωπος τὸν πατέρα αὐτοῦ καὶ τὴν μητέρα καὶ προσκολληθήσεται πρὸς τὴν γυναῖκα αὐτοῦ,
\vs{8}καὶ ἔσονται οἱ δύο εἰς σάρκα μίαν· ὥστε οὐκέτι εἰσὶν δύο ἀλλὰ μία σάρξ.
\vs{9}ὃ οὖν ὁ Θεὸς συνέζευξεν ἄνθρωπος μὴ χωριζέτω.

\vs{10}Καὶ εἰς τὴν οἰκίαν πάλιν οἱ μαθηταὶ περὶ τούτου ἐπηρώτων αὐτόν.
\vs{11}καὶ λέγει αὐτοῖς· Ὃς ἂν ἀπολύσῃ τὴν γυναῖκα αὐτοῦ καὶ γαμήσῃ ἄλλην μοιχᾶται ἐπ᾽ αὐτήν·
\vs{12}καὶ ἐὰν αὐτὴ ἀπολύσασα τὸν ἄνδρα αὐτῆς γαμήσῃ ἄλλον μοιχᾶται.

\vs{13}Καὶ προσέφερον αὐτῷ παιδία ἵνα αὐτῶν ἅψηται· οἱ δὲ μαθηταὶ ἐπετίμησαν αὐτοῖς.
\vs{14}Ἰδὼν δὲ ὁ Ἰησοῦς ἠγανάκτησεν καὶ εἶπεν αὐτοῖς· Ἄφετε τὰ παιδία ἔρχεσθαι πρός με, μὴ κωλύετε αὐτά, τῶν γὰρ τοιούτων ἐστὶν ἡ βασιλεία τοῦ Θεοῦ.
\vs{15}ἀμὴν λέγω ὑμῖν, ὃς ἂν μὴ δέξηται τὴν βασιλείαν τοῦ Θεοῦ ὡς παιδίον, οὐ μὴ εἰσέλθῃ εἰς αὐτήν.
\vs{16}καὶ ἐναγκαλισάμενος αὐτὰ κατευλόγει τιθεὶς τὰς χεῖρας ἐπ᾽ αὐτά.

\vs{17}Καὶ ἐκπορευομένου αὐτοῦ εἰς ὁδὸν προσδραμὼν εἷς καὶ γονυπετήσας αὐτὸν ἐπηρώτα αὐτόν· Διδάσκαλε ἀγαθέ, τί ποιήσω ἵνα ζωὴν αἰώνιον κληρονομήσω;
\vs{18}Ὁ δὲ Ἰησοῦς εἶπεν αὐτῷ· Τί με λέγεις ἀγαθόν; οὐδεὶς ἀγαθὸς εἰ μὴ εἷς ὁ Θεός.
\vs{19}τὰς ἐντολὰς οἶδας· Μὴ φονεύσῃς, Μὴ μοιχεύσῃς, Μὴ κλέψῃς, Μὴ ψευδομαρτυρήσῃς, Μὴ ἀποστερήσῃς, Τίμα τὸν πατέρα σου καὶ τὴν μητέρα.
\vs{20}Ὁ δὲ ἔφη αὐτῷ· Διδάσκαλε, ταῦτα πάντα ἐφυλαξάμην ἐκ νεότητός μου.
\vs{21}Ὁ δὲ Ἰησοῦς ἐμβλέψας αὐτῷ ἠγάπησεν αὐτὸν καὶ εἶπεν αὐτῷ· Ἕν σε ὑστερεῖ· ὕπαγε, ὅσα ἔχεις πώλησον καὶ δὸς τοῖς πτωχοῖς, καὶ ἕξεις θησαυρὸν ἐν οὐρανῷ, καὶ δεῦρο ἀκολούθει μοι.
\vs{22}Ὁ δὲ στυγνάσας ἐπὶ τῷ λόγῳ ἀπῆλθεν λυπούμενος· ἦν γὰρ ἔχων κτήματα πολλά.

\vs{23}Καὶ περιβλεψάμενος ὁ Ἰησοῦς λέγει τοῖς μαθηταῖς αὐτοῦ· Πῶς δυσκόλως οἱ τὰ χρήματα ἔχοντες εἰς τὴν βασιλείαν τοῦ Θεοῦ εἰσελεύσονται.
\vs{24}οἱ δὲ μαθηταὶ ἐθαμβοῦντο ἐπὶ τοῖς λόγοις αὐτοῦ. ὁ δὲ Ἰησοῦς πάλιν ἀποκριθεὶς λέγει αὐτοῖς· Τέκνα, πῶς δύσκολόν ἐστιν εἰς τὴν βασιλείαν τοῦ Θεοῦ εἰσελθεῖν·
\vs{25}εὐκοπώτερόν ἐστιν κάμηλον διὰ τῆς τρυμαλιᾶς τῆς ῥαφίδος διελθεῖν ἢ πλούσιον εἰς τὴν βασιλείαν τοῦ Θεοῦ εἰσελθεῖν.
\vs{26}Οἱ δὲ περισσῶς ἐξεπλήσσοντο λέγοντες πρὸς ἑαυτούς· Καὶ τίς δύναται σωθῆναι;
\vs{27}Ἐμβλέψας αὐτοῖς ὁ Ἰησοῦς λέγει· Παρὰ ἀνθρώποις ἀδύνατον, ἀλλ᾽ οὐ παρὰ θεῷ· πάντα γὰρ δυνατὰ παρὰ τῷ θεῷ.

\vs{28}Ἤρξατο λέγειν ὁ Πέτρος αὐτῷ· Ἰδοὺ ἡμεῖς ἀφήκαμεν πάντα καὶ ἠκολουθήκαμέν σοι.
\vs{29}Ἔφη ὁ Ἰησοῦς· Ἀμὴν λέγω ὑμῖν, οὐδείς ἐστιν ὃς ἀφῆκεν οἰκίαν ἢ ἀδελφοὺς ἢ ἀδελφὰς ἢ μητέρα ἢ πατέρα ἢ τέκνα ἢ ἀγροὺς ἕνεκεν ἐμοῦ καὶ ἕνεκεν τοῦ εὐαγγελίου,
\vs{30}ἐὰν μὴ λάβῃ ἑκατονταπλασίονα νῦν ἐν τῷ καιρῷ τούτῳ οἰκίας καὶ ἀδελφοὺς καὶ ἀδελφὰς καὶ μητέρας καὶ τέκνα καὶ ἀγροὺς μετὰ διωγμῶν, καὶ ἐν τῷ αἰῶνι τῷ ἐρχομένῳ ζωὴν αἰώνιον.
\vs{31}πολλοὶ δὲ ἔσονται πρῶτοι ἔσχατοι καὶ οἱ ἔσχατοι πρῶτοι.

\vs{32}Ἦσαν δὲ ἐν τῇ ὁδῷ ἀναβαίνοντες εἰς Ἱεροσόλυμα, καὶ ἦν προάγων αὐτοὺς ὁ Ἰησοῦς, καὶ ἐθαμβοῦντο, οἱ δὲ ἀκολουθοῦντες ἐφοβοῦντο. καὶ παραλαβὼν πάλιν τοὺς δώδεκα ἤρξατο αὐτοῖς λέγειν τὰ μέλλοντα αὐτῷ συμβαίνειν
\vs{33}ὅτι Ἰδοὺ ἀναβαίνομεν εἰς Ἱεροσόλυμα, καὶ ὁ Υἱὸς τοῦ ἀνθρώπου παραδοθήσεται τοῖς ἀρχιερεῦσιν καὶ τοῖς γραμματεῦσιν, καὶ κατακρινοῦσιν αὐτὸν θανάτῳ καὶ παραδώσουσιν αὐτὸν τοῖς ἔθνεσιν
\vs{34}καὶ ἐμπαίξουσιν αὐτῷ καὶ ἐμπτύσουσιν αὐτῷ καὶ μαστιγώσουσιν αὐτὸν καὶ ἀποκτενοῦσιν, καὶ μετὰ τρεῖς ἡμέρας ἀναστήσεται.

\vs{35}Καὶ προσπορεύονται αὐτῷ Ἰάκωβος καὶ Ἰωάννης οἱ υἱοὶ Ζεβεδαίου λέγοντες αὐτῷ· Διδάσκαλε, θέλομεν ἵνα ὃ ἐὰν αἰτήσωμέν σε ποιήσῃς ἡμῖν.
\vs{36}Ὁ δὲ εἶπεν αὐτοῖς· Τί θέλετε με ποιήσω ὑμῖν;
\vs{37}Οἱ δὲ εἶπαν αὐτῷ· Δὸς ἡμῖν ἵνα εἷς σου ἐκ δεξιῶν καὶ εἷς ἐξ ἀριστερῶν καθίσωμεν ἐν τῇ δόξῃ σου.
\vs{38}Ὁ δὲ Ἰησοῦς εἶπεν αὐτοῖς· Οὐκ οἴδατε τί αἰτεῖσθε. δύνασθε πιεῖν τὸ ποτήριον ὃ ἐγὼ πίνω ἢ τὸ βάπτισμα ὃ ἐγὼ βαπτίζομαι βαπτισθῆναι;
\vs{39}Οἱ δὲ εἶπαν αὐτῷ· Δυνάμεθα. Ὁ δὲ Ἰησοῦς εἶπεν αὐτοῖς· Τὸ ποτήριον ὃ ἐγὼ πίνω πίεσθε καὶ τὸ βάπτισμα ὃ ἐγὼ βαπτίζομαι βαπτισθήσεσθε,
\vs{40}τὸ δὲ καθίσαι ἐκ δεξιῶν μου ἢ ἐξ εὐωνύμων οὐκ ἔστιν ἐμὸν δοῦναι, ἀλλ᾽ οἷς ἡτοίμασται.
\vs{41}Καὶ ἀκούσαντες οἱ δέκα ἤρξαντο ἀγανακτεῖν περὶ Ἰακώβου καὶ Ἰωάννου.
\vs{42}καὶ προσκαλεσάμενος αὐτοὺς ὁ Ἰησοῦς λέγει αὐτοῖς· Οἴδατε ὅτι οἱ δοκοῦντες ἄρχειν τῶν ἐθνῶν κατακυριεύουσιν αὐτῶν καὶ οἱ μεγάλοι αὐτῶν κατεξουσιάζουσιν αὐτῶν.
\vs{43}οὐχ οὕτως δέ ἐστιν ἐν ὑμῖν, ἀλλ᾽ ὃς ἂν θέλῃ μέγας γενέσθαι ἐν ὑμῖν ἔσται ὑμῶν διάκονος,
\vs{44}καὶ ὃς ἂν θέλῃ ἐν ὑμῖν εἶναι πρῶτος ἔσται πάντων δοῦλος·
\vs{45}καὶ γὰρ ὁ Υἱὸς τοῦ ἀνθρώπου οὐκ ἦλθεν διακονηθῆναι ἀλλὰ διακονῆσαι καὶ δοῦναι τὴν ψυχὴν αὐτοῦ λύτρον ἀντὶ πολλῶν.

\vs{46}Καὶ ἔρχονται εἰς Ἰεριχώ. Καὶ ἐκπορευομένου αὐτοῦ ἀπὸ Ἰεριχὼ καὶ τῶν μαθητῶν αὐτοῦ καὶ ὄχλου ἱκανοῦ ὁ υἱὸς Τιμαίου Βαρτιμαῖος, τυφλὸς προσαίτης, ἐκάθητο παρὰ τὴν ὁδόν.
\vs{47}καὶ ἀκούσας ὅτι Ἰησοῦς ὁ Ναζαρηνός ἐστιν ἤρξατο κράζειν καὶ λέγειν· Υἱὲ Δαυὶδ Ἰησοῦ, ἐλέησόν με.
\vs{48}Καὶ ἐπετίμων αὐτῷ πολλοὶ ἵνα σιωπήσῃ· ὁ δὲ πολλῷ μᾶλλον ἔκραζεν· Υἱὲ Δαυίδ, ἐλέησόν με.
\vs{49}Καὶ στὰς ὁ Ἰησοῦς εἶπεν· Φωνήσατε αὐτόν. Καὶ φωνοῦσιν τὸν τυφλὸν λέγοντες αὐτῷ· Θάρσει, ἔγειρε, φωνεῖ σε.
\vs{50}Ὁ δὲ ἀποβαλὼν τὸ ἱμάτιον αὐτοῦ ἀναπηδήσας ἦλθεν πρὸς τὸν Ἰησοῦν.
\vs{51}Καὶ ἀποκριθεὶς αὐτῷ ὁ Ἰησοῦς εἶπεν· Τί σοι θέλεις ποιήσω; Ὁ δὲ τυφλὸς εἶπεν αὐτῷ· Ραββουνι, ἵνα ἀναβλέψω.
\vs{52}Καὶ ὁ Ἰησοῦς εἶπεν αὐτῷ· Ὕπαγε, ἡ πίστις σου σέσωκέν σε. καὶ εὐθὺς ἀνέβλεψεν καὶ ἠκολούθει αὐτῷ ἐν τῇ ὁδῷ.

\ch{11}
Καὶ ὅτε ἐγγίζουσιν εἰς Ἱεροσόλυμα εἰς Βηθφαγὴ καὶ Βηθανίαν πρὸς τὸ ὄρος τῶν Ἐλαιῶν, ἀποστέλλει δύο τῶν μαθητῶν αὐτοῦ
\vs{2}καὶ λέγει αὐτοῖς· Ὑπάγετε εἰς τὴν κώμην τὴν κατέναντι ὑμῶν, καὶ εὐθὺς εἰσπορευόμενοι εἰς αὐτὴν εὑρήσετε πῶλον δεδεμένον ἐφ᾽ ὃν οὐδεὶς οὔπω ἀνθρώπων ἐκάθισεν· λύσατε αὐτὸν καὶ φέρετε.
\vs{3}καὶ ἐάν τις ὑμῖν εἴπῃ· Τί ποιεῖτε τοῦτο; εἴπατε· Ὁ Κύριος αὐτοῦ χρείαν ἔχει, καὶ εὐθὺς αὐτὸν ἀποστέλλει πάλιν ὧδε.
\vs{4}Καὶ ἀπῆλθον καὶ εὗρον πῶλον δεδεμένον πρὸς θύραν ἔξω ἐπὶ τοῦ ἀμφόδου καὶ λύουσιν αὐτόν.
\vs{5}καί τινες τῶν ἐκεῖ ἑστηκότων ἔλεγον αὐτοῖς· Τί ποιεῖτε λύοντες τὸν πῶλον;
\vs{6}Οἱ δὲ εἶπαν αὐτοῖς καθὼς εἶπεν ὁ Ἰησοῦς, καὶ ἀφῆκαν αὐτούς.
\vs{7}καὶ φέρουσιν τὸν πῶλον πρὸς τὸν Ἰησοῦν καὶ ἐπιβάλλουσιν αὐτῷ τὰ ἱμάτια αὐτῶν, καὶ ἐκάθισεν ἐπ᾽ αὐτόν.
\vs{8}Καὶ πολλοὶ τὰ ἱμάτια αὐτῶν ἔστρωσαν εἰς τὴν ὁδόν, ἄλλοι δὲ στιβάδας κόψαντες ἐκ τῶν ἀγρῶν.
\vs{9}καὶ οἱ προάγοντες καὶ οἱ ἀκολουθοῦντες ἔκραζον· 
\begin{poetryblock}

\begin{quote}Ὡσαννά·\end{quote} 

\begin{quote}Εὐλογημένος ὁ ἐρχόμενος ἐν ὀνόματι Κυρίου·\end{quote}

\begin{quote} \vs{10}Εὐλογημένη ἡ ἐρχομένη βασιλεία τοῦ πατρὸς\end{quote} 

\begin{quote}ἡμῶν Δαυίδ·\end{quote} 

\begin{quote}Ὡσαννὰ ἐν τοῖς ὑψίστοις.\end{quote}
\end{poetryblock}

\vs{11}Καὶ εἰσῆλθεν εἰς Ἱεροσόλυμα εἰς τὸ ἱερόν καὶ περιβλεψάμενος πάντα, ὀψίας ἤδη οὔσης τῆς ὥρας, ἐξῆλθεν εἰς Βηθανίαν μετὰ τῶν δώδεκα.

\vs{12}Καὶ τῇ ἐπαύριον ἐξελθόντων αὐτῶν ἀπὸ Βηθανίας ἐπείνασεν.
\vs{13}καὶ ἰδὼν συκῆν ἀπὸ μακρόθεν ἔχουσαν φύλλα ἦλθεν, εἰ ἄρα τι εὑρήσει ἐν αὐτῇ, καὶ ἐλθὼν ἐπ᾽ αὐτὴν οὐδὲν εὗρεν εἰ μὴ φύλλα· ὁ γὰρ καιρὸς οὐκ ἦν σύκων.
\vs{14}καὶ ἀποκριθεὶς εἶπεν αὐτῇ· Μηκέτι εἰς τὸν αἰῶνα ἐκ σοῦ μηδεὶς καρπὸν φάγοι. καὶ ἤκουον οἱ μαθηταὶ αὐτοῦ.
\vs{15}Καὶ ἔρχονται εἰς Ἱεροσόλυμα. Καὶ εἰσελθὼν εἰς τὸ ἱερὸν ἤρξατο ἐκβάλλειν τοὺς πωλοῦντας καὶ τοὺς ἀγοράζοντας ἐν τῷ ἱερῷ, καὶ τὰς τραπέζας τῶν κολλυβιστῶν καὶ τὰς καθέδρας τῶν πωλούντων τὰς περιστερὰς κατέστρεψεν,
\vs{16}καὶ οὐκ ἤφιεν ἵνα τις διενέγκῃ σκεῦος διὰ τοῦ ἱεροῦ.
\vs{17}καὶ ἐδίδασκεν καὶ ἔλεγεν αὐτοῖς· Οὐ γέγραπται ὅτι 
\begin{poetryblock}

\begin{quote}Ὁ οἶκός μου οἶκος προσευχῆς κληθήσεται\end{quote} 

\begin{quote}πᾶσιν τοῖς ἔθνεσιν; ὑμεῖς δὲ πεποιήκατε αὐτὸν Σπήλαιον λῃστῶν.\end{quote}
\end{poetryblock}

\vs{18}Καὶ ἤκουσαν οἱ ἀρχιερεῖς καὶ οἱ γραμματεῖς καὶ ἐζήτουν πῶς αὐτὸν ἀπολέσωσιν· ἐφοβοῦντο γὰρ αὐτόν, πᾶς γὰρ ὁ ὄχλος ἐξεπλήσσετο ἐπὶ τῇ διδαχῇ αὐτοῦ.

\vs{19}Καὶ ὅταν ὀψὲ ἐγένετο, ἐξεπορεύοντο ἔξω τῆς πόλεως.

\vs{20}Καὶ παραπορευόμενοι πρωῒ εἶδον τὴν συκῆν ἐξηραμμένην ἐκ ῥιζῶν.
\vs{21}καὶ ἀναμνησθεὶς ὁ Πέτρος λέγει αὐτῷ· Ῥαββί, ἴδε ἡ συκῆ ἣν κατηράσω ἐξήρανται.
\vs{22}Καὶ ἀποκριθεὶς ὁ Ἰησοῦς λέγει αὐτοῖς· Ἔχετε πίστιν θεοῦ.
\vs{23}ἀμὴν λέγω ὑμῖν ὅτι ὃς ἂν εἴπῃ τῷ ὄρει τούτῳ· Ἄρθητι καὶ βλήθητι εἰς τὴν θάλασσαν, καὶ μὴ διακριθῇ ἐν τῇ καρδίᾳ αὐτοῦ ἀλλὰ πιστεύῃ ὅτι ὃ λαλεῖ γίνεται, ἔσται αὐτῷ.
\vs{24}διὰ τοῦτο λέγω ὑμῖν, πάντα ὅσα προσεύχεσθε καὶ αἰτεῖσθε, πιστεύετε ὅτι ἐλάβετε, καὶ ἔσται ὑμῖν.
\vs{25}Καὶ ὅταν στήκετε προσευχόμενοι, ἀφίετε εἴ τι ἔχετε κατά τινος, ἵνα καὶ ὁ Πατὴρ ὑμῶν ὁ ἐν τοῖς οὐρανοῖς ἀφῇ ὑμῖν τὰ παραπτώματα ὑμῶν.

\vs{27}Καὶ ἔρχονται πάλιν εἰς Ἱεροσόλυμα. καὶ ἐν τῷ ἱερῷ περιπατοῦντος αὐτοῦ ἔρχονται πρὸς αὐτὸν οἱ ἀρχιερεῖς καὶ οἱ γραμματεῖς καὶ οἱ πρεσβύτεροι
\vs{28}καὶ ἔλεγον αὐτῷ· Ἐν ποίᾳ ἐξουσίᾳ ταῦτα ποιεῖς; ἢ τίς σοι ἔδωκεν τὴν ἐξουσίαν ταύτην ἵνα ταῦτα ποιῇς;
\vs{29}Ὁ δὲ Ἰησοῦς εἶπεν αὐτοῖς· Ἐπερωτήσω ὑμᾶς ἕνα λόγον, καὶ ἀποκρίθητέ μοι καὶ ἐρῶ ὑμῖν ἐν ποίᾳ ἐξουσίᾳ ταῦτα ποιῶ·
\vs{30}τὸ βάπτισμα τὸ Ἰωάννου ἐξ οὐρανοῦ ἦν ἢ ἐξ ἀνθρώπων; ἀποκρίθητέ μοι.
\vs{31}Καὶ διελογίζοντο πρὸς ἑαυτοὺς λέγοντες· Ἐὰν εἴπωμεν· Ἐξ οὐρανοῦ, ἐρεῖ· Διὰ τί οὖν οὐκ ἐπιστεύσατε αὐτῷ;
\vs{32}ἀλλὰ εἴπωμεν· Ἐξ ἀνθρώπων;— ἐφοβοῦντο τὸν ὄχλον· ἅπαντες γὰρ εἶχον τὸν Ἰωάννην ὄντως ὅτι προφήτης ἦν.
\vs{33}καὶ ἀποκριθέντες τῷ Ἰησοῦ λέγουσιν· Οὐκ οἴδαμεν. Καὶ ὁ Ἰησοῦς λέγει αὐτοῖς· Οὐδὲ ἐγὼ λέγω ὑμῖν ἐν ποίᾳ ἐξουσίᾳ ταῦτα ποιῶ.

\ch{12}
Καὶ ἤρξατο αὐτοῖς ἐν παραβολαῖς λαλεῖν· Ἀμπελῶνα ἄνθρωπος ἐφύτευσεν καὶ περιέθηκεν φραγμὸν καὶ ὤρυξεν ὑπολήνιον καὶ ᾠκοδόμησεν πύργον καὶ ἐξέδετο αὐτὸν γεωργοῖς καὶ ἀπεδήμησεν.
\vs{2}Καὶ ἀπέστειλεν πρὸς τοὺς γεωργοὺς τῷ καιρῷ δοῦλον ἵνα παρὰ τῶν γεωργῶν λάβῃ ἀπὸ τῶν καρπῶν τοῦ ἀμπελῶνος·
\vs{3}καὶ λαβόντες αὐτὸν ἔδειραν καὶ ἀπέστειλαν κενόν.
\vs{4}Καὶ πάλιν ἀπέστειλεν πρὸς αὐτοὺς ἄλλον δοῦλον· κἀκεῖνον ἐκεφαλίωσαν καὶ ἠτίμασαν.
\vs{5}καὶ ἄλλον ἀπέστειλεν· κἀκεῖνον ἀπέκτειναν, καὶ πολλοὺς ἄλλους, οὓς μὲν δέροντες, οὓς δὲ ἀποκτέννοντες.
\vs{6}Ἔτι ἕνα εἶχεν υἱὸν ἀγαπητόν· ἀπέστειλεν αὐτὸν ἔσχατον πρὸς αὐτοὺς λέγων ὅτι Ἐντραπήσονται τὸν υἱόν μου.
\vs{7}Ἐκεῖνοι δὲ οἱ γεωργοὶ πρὸς ἑαυτοὺς εἶπαν ὅτι Οὗτός ἐστιν ὁ κληρονόμος· δεῦτε ἀποκτείνωμεν αὐτόν, καὶ ἡμῶν ἔσται ἡ κληρονομία.
\vs{8}καὶ λαβόντες ἀπέκτειναν αὐτόν καὶ ἐξέβαλον αὐτὸν ἔξω τοῦ ἀμπελῶνος.
\vs{9}Τί οὖν ποιήσει ὁ κύριος τοῦ ἀμπελῶνος; ἐλεύσεται καὶ ἀπολέσει τοὺς γεωργούς καὶ δώσει τὸν ἀμπελῶνα ἄλλοις.
\vs{10}οὐδὲ τὴν γραφὴν ταύτην ἀνέγνωτε· 
\begin{poetryblock}

\begin{quote}Λίθον ὃν ἀπεδοκίμασαν οἱ οἰκοδομοῦντες,\end{quote} 

\begin{quote}Οὗτος ἐγενήθη εἰς κεφαλὴν γωνίας·\end{quote}

\begin{quote} \vs{11}Παρὰ Κυρίου ἐγένετο αὕτη\end{quote} 

\begin{quote}Καὶ ἔστιν θαυμαστὴ ἐν ὀφθαλμοῖς ἡμῶν;\end{quote}
\end{poetryblock}

\vs{12}Καὶ ἐζήτουν αὐτὸν κρατῆσαι, καὶ ἐφοβήθησαν τὸν ὄχλον, ἔγνωσαν γὰρ ὅτι πρὸς αὐτοὺς τὴν παραβολὴν εἶπεν. καὶ ἀφέντες αὐτὸν ἀπῆλθον.

\vs{13}Καὶ ἀποστέλλουσιν πρὸς αὐτόν τινας τῶν Φαρισαίων καὶ τῶν Ἡρῳδιανῶν ἵνα αὐτὸν ἀγρεύσωσιν λόγῳ.
\vs{14}καὶ ἐλθόντες λέγουσιν αὐτῷ· Διδάσκαλε, οἴδαμεν ὅτι ἀληθὴς εἶ καὶ οὐ μέλει σοι περὶ οὐδενός· οὐ γὰρ βλέπεις εἰς πρόσωπον ἀνθρώπων, ἀλλ᾽ ἐπ᾽ ἀληθείας τὴν ὁδὸν τοῦ Θεοῦ διδάσκεις· ἔξεστιν δοῦναι κῆνσον Καίσαρι ἢ οὔ; δῶμεν ἢ μὴ δῶμεν;
\vs{15}Ὁ δὲ εἰδὼς αὐτῶν τὴν ὑπόκρισιν εἶπεν αὐτοῖς· Τί με πειράζετε; φέρετέ μοι δηνάριον ἵνα ἴδω.
\vs{16}οἱ δὲ ἤνεγκαν. καὶ λέγει αὐτοῖς· Τίνος ἡ εἰκὼν αὕτη καὶ ἡ ἐπιγραφή; Οἱ δὲ εἶπαν αὐτῷ· Καίσαρος.
\vs{17}Ὁ δὲ Ἰησοῦς εἶπεν αὐτοῖς· Τὰ Καίσαρος ἀπόδοτε Καίσαρι καὶ τὰ τοῦ Θεοῦ τῷ Θεῷ. Καὶ ἐξεθαύμαζον ἐπ᾽ αὐτῷ.

\vs{18}Καὶ ἔρχονται Σαδδουκαῖοι πρὸς αὐτόν, οἵτινες λέγουσιν ἀνάστασιν μὴ εἶναι, καὶ ἐπηρώτων αὐτὸν λέγοντες·
\vs{19}Διδάσκαλε, Μωϋσῆς ἔγραψεν ἡμῖν ὅτι ἐάν τινος ἀδελφὸς ἀποθάνῃ καὶ καταλίπῃ γυναῖκα καὶ μὴ ἀφῇ τέκνον, ἵνα λάβῃ ὁ ἀδελφὸς αὐτοῦ τὴν γυναῖκα καὶ ἐξαναστήσῃ σπέρμα τῷ ἀδελφῷ αὐτοῦ.
\vs{20}ἑπτὰ ἀδελφοὶ ἦσαν· καὶ ὁ πρῶτος ἔλαβεν γυναῖκα καὶ ἀποθνῄσκων οὐκ ἀφῆκεν σπέρμα·
\vs{21}καὶ ὁ δεύτερος ἔλαβεν αὐτήν καὶ ἀπέθανεν μὴ καταλιπὼν σπέρμα· καὶ ὁ τρίτος ὡσαύτως·
\vs{22}καὶ οἱ ἑπτὰ οὐκ ἀφῆκαν σπέρμα. ἔσχατον πάντων καὶ ἡ γυνὴ ἀπέθανεν.
\vs{23}ἐν τῇ ἀναστάσει ὅταν ἀναστῶσιν τίνος αὐτῶν ἔσται γυνή; οἱ γὰρ ἑπτὰ ἔσχον αὐτὴν γυναῖκα.
\vs{24}Ἔφη αὐτοῖς ὁ Ἰησοῦς· Οὐ διὰ τοῦτο πλανᾶσθε μὴ εἰδότες τὰς γραφὰς μηδὲ τὴν δύναμιν τοῦ Θεοῦ;
\vs{25}ὅταν γὰρ ἐκ νεκρῶν ἀναστῶσιν οὔτε γαμοῦσιν οὔτε γαμίζονται, ἀλλ᾽ εἰσὶν ὡς ἄγγελοι ἐν τοῖς οὐρανοῖς.
\vs{26}Περὶ δὲ τῶν νεκρῶν ὅτι ἐγείρονται οὐκ ἀνέγνωτε ἐν τῇ βίβλῳ Μωϋσέως ἐπὶ τοῦ Βάτου πῶς εἶπεν αὐτῷ ὁ Θεὸς λέγων· Ἐγὼ ὁ Θεὸς Ἀβραὰμ καὶ ὁ Θεὸς Ἰσαὰκ καὶ ὁ Θεὸς Ἰακώβ;
\vs{27}οὐκ ἔστιν Θεὸς νεκρῶν ἀλλὰ ζώντων· πολὺ πλανᾶσθε.

\vs{28}Καὶ προσελθὼν εἷς τῶν γραμματέων ἀκούσας αὐτῶν συζητούντων, ἰδὼν ὅτι καλῶς ἀπεκρίθη αὐτοῖς ἐπηρώτησεν αὐτόν· Ποία ἐστὶν ἐντολὴ πρώτη πάντων;
\vs{29}Ἀπεκρίθη ὁ Ἰησοῦς Ὅτι Πρώτη ἐστίν· Ἄκουε, Ἰσραήλ, Κύριος ὁ Θεὸς ἡμῶν Κύριος εἷς ἐστιν,
\vs{30}καὶ ἀγαπήσεις Κύριον τὸν Θεόν σου ἐξ ὅλης τῆς καρδίας σου καὶ ἐξ ὅλης τῆς ψυχῆς σου καὶ ἐξ ὅλης τῆς διανοίας σου καὶ ἐξ ὅλης τῆς ἰσχύος σου.
\vs{31}δευτέρα αὕτη· Ἀγαπήσεις τὸν πλησίον σου ὡς σεαυτόν. μείζων τούτων ἄλλη ἐντολὴ οὐκ ἔστιν.
\vs{32}Καὶ εἶπεν αὐτῷ ὁ γραμματεύς· Καλῶς, Διδάσκαλε, ἐπ᾽ ἀληθείας εἶπες ὅτι εἷς ἐστιν καὶ οὐκ ἔστιν ἄλλος πλὴν αὐτοῦ·
\vs{33}καὶ τὸ ἀγαπᾶν αὐτὸν ἐξ ὅλης τῆς καρδίας καὶ ἐξ ὅλης τῆς συνέσεως καὶ ἐξ ὅλης τῆς ἰσχύος καὶ τὸ ἀγαπᾶν τὸν πλησίον ὡς ἑαυτὸν περισσότερόν ἐστιν πάντων τῶν ὁλοκαυτωμάτων καὶ θυσιῶν.
\vs{34}Καὶ ὁ Ἰησοῦς ἰδὼν αὐτὸν ὅτι νουνεχῶς ἀπεκρίθη εἶπεν αὐτῷ· Οὐ μακρὰν εἶ ἀπὸ τῆς βασιλείας τοῦ Θεοῦ. καὶ οὐδεὶς οὐκέτι ἐτόλμα αὐτὸν ἐπερωτῆσαι.

\vs{35}Καὶ ἀποκριθεὶς ὁ Ἰησοῦς ἔλεγεν διδάσκων ἐν τῷ ἱερῷ· Πῶς λέγουσιν οἱ γραμματεῖς ὅτι ὁ Χριστὸς υἱὸς Δαυίδ ἐστιν;
\vs{36}αὐτὸς Δαυὶδ εἶπεν ἐν τῷ Πνεύματι τῷ Ἁγίῳ· 
\begin{poetryblock}

\begin{quote}Εἶπεν Κύριος τῷ Κυρίῳ μου·\end{quote} 

\begin{quote}Κάθου ἐκ δεξιῶν μου,\end{quote} 

\begin{quote}Ἕως ἂν θῶ τοὺς ἐχθρούς σου\end{quote} 

\begin{quote}Ὑποκάτω τῶν ποδῶν σου.\end{quote}
\end{poetryblock}

\vs{37}Αὐτὸς Δαυὶδ λέγει αὐτὸν Κύριον, καὶ πόθεν αὐτοῦ ἐστιν υἱός;

Καὶ ὁ πολὺς ὄχλος ἤκουεν αὐτοῦ ἡδέως.

\vs{38}Καὶ ἐν τῇ διδαχῇ αὐτοῦ ἔλεγεν· Βλέπετε ἀπὸ τῶν γραμματέων τῶν θελόντων ἐν στολαῖς περιπατεῖν καὶ ἀσπασμοὺς ἐν ταῖς ἀγοραῖς
\vs{39}καὶ πρωτοκαθεδρίας ἐν ταῖς συναγωγαῖς καὶ πρωτοκλισίας ἐν τοῖς δείπνοις,
\vs{40}οἱ κατεσθίοντες τὰς οἰκίας τῶν χηρῶν καὶ προφάσει μακρὰ προσευχόμενοι· οὗτοι λήμψονται περισσότερον κρίμα.

\vs{41}Καὶ καθίσας κατέναντι τοῦ γαζοφυλακίου ἐθεώρει πῶς ὁ ὄχλος βάλλει χαλκὸν εἰς τὸ γαζοφυλάκιον. καὶ πολλοὶ πλούσιοι ἔβαλλον πολλά·
\vs{42}καὶ ἐλθοῦσα μία χήρα πτωχὴ ἔβαλεν λεπτὰ δύο, ὅ ἐστιν κοδράντης.
\vs{43}Καὶ προσκαλεσάμενος τοὺς μαθητὰς αὐτοῦ εἶπεν αὐτοῖς· Ἀμὴν λέγω ὑμῖν ὅτι ἡ χήρα αὕτη ἡ πτωχὴ πλεῖον πάντων ἔβαλεν τῶν βαλλόντων εἰς τὸ γαζοφυλάκιον·
\vs{44}πάντες γὰρ ἐκ τοῦ περισσεύοντος αὐτοῖς ἔβαλον, αὕτη δὲ ἐκ τῆς ὑστερήσεως αὐτῆς πάντα ὅσα εἶχεν ἔβαλεν ὅλον τὸν βίον αὐτῆς.

\ch{13}
Καὶ ἐκπορευομένου αὐτοῦ ἐκ τοῦ ἱεροῦ λέγει αὐτῷ εἷς τῶν μαθητῶν αὐτοῦ· Διδάσκαλε, ἴδε ποταποὶ λίθοι καὶ ποταπαὶ οἰκοδομαί.
\vs{2}Καὶ ὁ Ἰησοῦς εἶπεν αὐτῷ· Βλέπεις ταύτας τὰς μεγάλας οἰκοδομάς; οὐ μὴ ἀφεθῇ ὧδε λίθος ἐπὶ λίθον ὃς οὐ μὴ καταλυθῇ.
\vs{3}Καὶ καθημένου αὐτοῦ εἰς τὸ ὄρος τῶν Ἐλαιῶν κατέναντι τοῦ ἱεροῦ ἐπηρώτα αὐτὸν κατ᾽ ἰδίαν Πέτρος καὶ Ἰάκωβος καὶ Ἰωάννης καὶ Ἀνδρέας·
\vs{4}Εἰπὸν ἡμῖν, πότε ταῦτα ἔσται καὶ τί τὸ σημεῖον ὅταν μέλλῃ ταῦτα συντελεῖσθαι πάντα;
\vs{5}Ὁ δὲ Ἰησοῦς ἤρξατο λέγειν αὐτοῖς· Βλέπετε μή τις ὑμᾶς πλανήσῃ·
\vs{6}πολλοὶ ἐλεύσονται ἐπὶ τῷ ὀνόματί μου λέγοντες ὅτι Ἐγώ εἰμι, καὶ πολλοὺς πλανήσουσιν.
\vs{7}ὅταν δὲ ἀκούσητε πολέμους καὶ ἀκοὰς πολέμων, μὴ θροεῖσθε· δεῖ γενέσθαι, ἀλλ᾽ οὔπω τὸ τέλος.
\vs{8}ἐγερθήσεται γὰρ ἔθνος ἐπ᾽ ἔθνος καὶ βασιλεία ἐπὶ βασιλείαν, ἔσονται σεισμοὶ κατὰ τόπους, ἔσονται λιμοί· ἀρχὴ ὠδίνων ταῦτα.

\vs{9}Βλέπετε δὲ ὑμεῖς ἑαυτούς· παραδώσουσιν ὑμᾶς εἰς συνέδρια καὶ εἰς συναγωγὰς δαρήσεσθε καὶ ἐπὶ ἡγεμόνων καὶ βασιλέων σταθήσεσθε ἕνεκεν ἐμοῦ εἰς μαρτύριον αὐτοῖς.
\vs{10}καὶ εἰς πάντα τὰ ἔθνη πρῶτον δεῖ κηρυχθῆναι τὸ εὐαγγέλιον.
\vs{11}καὶ ὅταν ἄγωσιν ὑμᾶς παραδιδόντες, μὴ προμεριμνᾶτε τί λαλήσητε, ἀλλ᾽ ὃ ἐὰν δοθῇ ὑμῖν ἐν ἐκείνῃ τῇ ὥρᾳ τοῦτο λαλεῖτε· οὐ γάρ ἐστε ὑμεῖς οἱ λαλοῦντες ἀλλὰ τὸ Πνεῦμα τὸ Ἅγιον.
\vs{12}Καὶ παραδώσει ἀδελφὸς ἀδελφὸν εἰς θάνατον καὶ πατὴρ τέκνον, καὶ ἐπαναστήσονται τέκνα ἐπὶ γονεῖς καὶ θανατώσουσιν αὐτούς·
\vs{13}καὶ ἔσεσθε μισούμενοι ὑπὸ πάντων διὰ τὸ ὄνομά μου. ὁ δὲ ὑπομείνας εἰς τέλος οὗτος σωθήσεται.

\vs{14}Ὅταν δὲ ἴδητε τὸ βδέλυγμα τῆς ἐρημώσεως ἑστηκότα ὅπου οὐ δεῖ, ὁ ἀναγινώσκων νοείτω, τότε οἱ ἐν τῇ Ἰουδαίᾳ φευγέτωσαν εἰς τὰ ὄρη,
\vs{15}ὁ δὲ ἐπὶ τοῦ δώματος μὴ καταβάτω μηδὲ εἰσελθάτω ἆραι τι ἐκ τῆς οἰκίας αὐτοῦ,
\vs{16}καὶ ὁ εἰς τὸν ἀγρὸν μὴ ἐπιστρεψάτω εἰς τὰ ὀπίσω ἆραι τὸ ἱμάτιον αὐτοῦ.
\vs{17}Οὐαὶ δὲ ταῖς ἐν γαστρὶ ἐχούσαις καὶ ταῖς θηλαζούσαις ἐν ἐκείναις ταῖς ἡμέραις.

\vs{18}προσεύχεσθε δὲ ἵνα μὴ γένηται χειμῶνος·
\vs{19}ἔσονται γὰρ αἱ ἡμέραι ἐκεῖναι θλῖψις οἵα οὐ γέγονεν τοιαύτη ἀπ᾽ ἀρχῆς κτίσεως ἣν ἔκτισεν ὁ Θεὸς ἕως τοῦ νῦν καὶ οὐ μὴ γένηται.
\vs{20}καὶ εἰ μὴ ἐκολόβωσεν Κύριος τὰς ἡμέρας, οὐκ ἂν ἐσώθη πᾶσα σάρξ· ἀλλὰ διὰ τοὺς ἐκλεκτοὺς οὓς ἐξελέξατο ἐκολόβωσεν τὰς ἡμέρας.

\vs{21}Καὶ τότε ἐάν τις ὑμῖν εἴπῃ· Ἴδε ὧδε ὁ Χριστός, Ἴδε ἐκεῖ, μὴ πιστεύετε·
\vs{22}ἐγερθήσονται γὰρ ψευδόχριστοι καὶ ψευδοπροφῆται καὶ δώσουσιν σημεῖα καὶ τέρατα πρὸς τὸ ἀποπλανᾶν, εἰ δυνατὸν, τοὺς ἐκλεκτούς.
\vs{23}ὑμεῖς δὲ βλέπετε· προείρηκα ὑμῖν πάντα.

\vs{24}Ἀλλὰ ἐν ἐκείναις ταῖς ἡμέραις μετὰ τὴν θλῖψιν ἐκείνην 
\begin{poetryblock}

\begin{quote}Ὁ ἥλιος σκοτισθήσεται,\end{quote} 

\begin{quote}Καὶ ἡ σελήνη οὐ δώσει τὸ φέγγος αὐτῆς,\end{quote}

\begin{quote} \vs{25}Καὶ οἱ ἀστέρες ἔσονται ἐκ τοῦ οὐρανοῦ πίπτοντες,\end{quote} 

\begin{quote}Καὶ αἱ δυνάμεις αἱ ἐν τοῖς οὐρανοῖς σαλευθήσονται.\end{quote}
\end{poetryblock}

\vs{26}Καὶ τότε ὄψονται τὸν Υἱὸν τοῦ ἀνθρώπου ἐρχόμενον ἐν νεφέλαις μετὰ δυνάμεως πολλῆς καὶ δόξης.
\vs{27}καὶ τότε ἀποστελεῖ τοὺς ἀγγέλους καὶ ἐπισυνάξει τοὺς ἐκλεκτοὺς αὐτοῦ ἐκ τῶν τεσσάρων ἀνέμων ἀπ᾽ ἄκρου γῆς ἕως ἄκρου οὐρανοῦ.

\vs{28}Ἀπὸ δὲ τῆς συκῆς μάθετε τὴν παραβολήν· ὅταν ἤδη ὁ κλάδος αὐτῆς ἁπαλὸς γένηται καὶ ἐκφύῃ τὰ φύλλα, γινώσκετε ὅτι ἐγγὺς τὸ θέρος ἐστίν·
\vs{29}οὕτως καὶ ὑμεῖς, ὅταν ἴδητε ταῦτα γινόμενα, γινώσκετε ὅτι ἐγγύς ἐστιν ἐπὶ θύραις.

\vs{30}ἀμὴν λέγω ὑμῖν ὅτι οὐ μὴ παρέλθῃ ἡ γενεὰ αὕτη μέχρις οὗ ταῦτα πάντα γένηται.
\vs{31}ὁ οὐρανὸς καὶ ἡ γῆ παρελεύσονται, οἱ δὲ λόγοι μου οὐ μὴ παρελεύσονται.

\vs{32}Περὶ δὲ τῆς ἡμέρας ἐκείνης ἢ τῆς ὥρας οὐδεὶς οἶδεν, οὐδὲ οἱ ἄγγελοι ἐν οὐρανῷ οὐδὲ ὁ Υἱός, εἰ μὴ ὁ Πατήρ.

\vs{33}Βλέπετε, ἀγρυπνεῖτε· οὐκ οἴδατε γὰρ πότε ὁ καιρός ἐστιν.
\vs{34}Ὡς ἄνθρωπος ἀπόδημος ἀφεὶς τὴν οἰκίαν αὐτοῦ καὶ δοὺς τοῖς δούλοις αὐτοῦ τὴν ἐξουσίαν ἑκάστῳ τὸ ἔργον αὐτοῦ καὶ τῷ θυρωρῷ ἐνετείλατο ἵνα γρηγορῇ.
\vs{35}γρηγορεῖτε οὖν· οὐκ οἴδατε γὰρ πότε ὁ κύριος τῆς οἰκίας ἔρχεται, ἢ ὀψὲ ἢ μεσονύκτιον ἢ ἀλεκτοροφωνίας ἢ πρωΐ,
\vs{36}μὴ ἐλθὼν ἐξαίφνης εὕρῃ ὑμᾶς καθεύδοντας.
\vs{37}ὃ δὲ ὑμῖν λέγω πᾶσιν λέγω, γρηγορεῖτε.

\ch{14}
Ἦν δὲ τὸ πάσχα καὶ τὰ ἄζυμα μετὰ δύο ἡμέρας. καὶ ἐζήτουν οἱ ἀρχιερεῖς καὶ οἱ γραμματεῖς πῶς αὐτὸν ἐν δόλῳ κρατήσαντες ἀποκτείνωσιν·
\vs{2}ἔλεγον γάρ· Μὴ ἐν τῇ ἑορτῇ, μήποτε ἔσται θόρυβος τοῦ λαοῦ.

\vs{3}Καὶ ὄντος αὐτοῦ ἐν Βηθανίᾳ ἐν τῇ οἰκίᾳ Σίμωνος τοῦ λεπροῦ, κατακειμένου αὐτοῦ ἦλθεν γυνὴ ἔχουσα ἀλάβαστρον μύρου νάρδου πιστικῆς πολυτελοῦς, συντρίψασα τὴν ἀλάβαστρον κατέχεεν αὐτοῦ τῆς κεφαλῆς.
\vs{4}Ἦσαν δέ τινες ἀγανακτοῦντες πρὸς ἑαυτούς· Εἰς τί ἡ ἀπώλεια αὕτη τοῦ μύρου γέγονεν;
\vs{5}ἠδύνατο γὰρ τοῦτο τὸ μύρον πραθῆναι ἐπάνω δηναρίων τριακοσίων καὶ δοθῆναι τοῖς πτωχοῖς· καὶ ἐνεβριμῶντο αὐτῇ.
\vs{6}Ὁ δὲ Ἰησοῦς εἶπεν· Ἄφετε αὐτήν· τί αὐτῇ κόπους παρέχετε; καλὸν ἔργον ἠργάσατο ἐν ἐμοί.
\vs{7}πάντοτε γὰρ τοὺς πτωχοὺς ἔχετε μεθ᾽ ἑαυτῶν καὶ ὅταν θέλητε δύνασθε αὐτοῖς εὖ ποιῆσαι, ἐμὲ δὲ οὐ πάντοτε ἔχετε.
\vs{8}ὃ ἔσχεν ἐποίησεν· προέλαβεν μυρίσαι τὸ σῶμά μου εἰς τὸν ἐνταφιασμόν.
\vs{9}ἀμὴν δὲ λέγω ὑμῖν, ὅπου ἐὰν κηρυχθῇ τὸ εὐαγγέλιον εἰς ὅλον τὸν κόσμον, καὶ ὃ ἐποίησεν αὕτη λαληθήσεται εἰς μνημόσυνον αὐτῆς.

\vs{10}Καὶ Ἰούδας Ἰσκαριὼθ ὁ εἷς τῶν δώδεκα ἀπῆλθεν πρὸς τοὺς ἀρχιερεῖς ἵνα αὐτὸν παραδοῖ αὐτοῖς.
\vs{11}οἱ δὲ ἀκούσαντες ἐχάρησαν καὶ ἐπηγγείλαντο αὐτῷ ἀργύριον δοῦναι. καὶ ἐζήτει πῶς αὐτὸν εὐκαίρως παραδοῖ.

\vs{12}Καὶ τῇ πρώτῃ ἡμέρᾳ τῶν ἀζύμων, ὅτε τὸ πάσχα ἔθυον, λέγουσιν αὐτῷ οἱ μαθηταὶ αὐτοῦ· Ποῦ θέλεις ἀπελθόντες ἑτοιμάσωμεν ἵνα φάγῃς τὸ πάσχα;
\vs{13}Καὶ ἀποστέλλει δύο τῶν μαθητῶν αὐτοῦ καὶ λέγει αὐτοῖς· Ὑπάγετε εἰς τὴν πόλιν, καὶ ἀπαντήσει ὑμῖν ἄνθρωπος κεράμιον ὕδατος βαστάζων· ἀκολουθήσατε αὐτῷ
\vs{14}καὶ ὅπου ἐὰν εἰσέλθῃ εἴπατε τῷ οἰκοδεσπότῃ ὅτι Ὁ Διδάσκαλος λέγει· Ποῦ ἐστιν τὸ κατάλυμά μου ὅπου τὸ πάσχα μετὰ τῶν μαθητῶν μου φάγω;
\vs{15}καὶ αὐτὸς ὑμῖν δείξει ἀνάγαιον μέγα ἐστρωμένον ἕτοιμον· καὶ ἐκεῖ ἑτοιμάσατε ἡμῖν.
\vs{16}Καὶ ἐξῆλθον οἱ μαθηταὶ καὶ ἦλθον εἰς τὴν πόλιν καὶ εὗρον καθὼς εἶπεν αὐτοῖς καὶ ἡτοίμασαν τὸ πάσχα.

\vs{17}Καὶ ὀψίας γενομένης ἔρχεται μετὰ τῶν δώδεκα.
\vs{18}Καὶ ἀνακειμένων αὐτῶν καὶ ἐσθιόντων ὁ Ἰησοῦς εἶπεν· Ἀμὴν λέγω ὑμῖν ὅτι εἷς ἐξ ὑμῶν παραδώσει με ὁ ἐσθίων μετ᾽ ἐμοῦ.
\vs{19}ἤρξαντο λυπεῖσθαι καὶ λέγειν αὐτῷ εἷς κατὰ εἷς· Μήτι ἐγώ;
\vs{20}Ὁ δὲ εἶπεν αὐτοῖς· Εἷς τῶν δώδεκα, ὁ ἐμβαπτόμενος μετ᾽ ἐμοῦ εἰς τὸ τρύβλιον.
\vs{21}ὅτι ὁ μὲν Υἱὸς τοῦ ἀνθρώπου ὑπάγει καθὼς γέγραπται περὶ αὐτοῦ, οὐαὶ δὲ τῷ ἀνθρώπῳ ἐκείνῳ δι᾽ οὗ ὁ Υἱὸς τοῦ ἀνθρώπου παραδίδοται· καλὸν αὐτῷ εἰ οὐκ ἐγεννήθη ὁ ἄνθρωπος ἐκεῖνος.

\vs{22}Καὶ ἐσθιόντων αὐτῶν λαβὼν ἄρτον εὐλογήσας ἔκλασεν καὶ ἔδωκεν αὐτοῖς καὶ εἶπεν· Λάβετε, τοῦτό ἐστιν τὸ σῶμά μου.
\vs{23}Καὶ λαβὼν ποτήριον εὐχαριστήσας ἔδωκεν αὐτοῖς, καὶ ἔπιον ἐξ αὐτοῦ πάντες.
\vs{24}καὶ εἶπεν αὐτοῖς· Τοῦτό ἐστιν τὸ αἷμά μου τῆς διαθήκης τὸ ἐκχυννόμενον ὑπὲρ πολλῶν.
\vs{25}ἀμὴν λέγω ὑμῖν ὅτι οὐκέτι οὐ μὴ πίω ἐκ τοῦ γενήματος τῆς ἀμπέλου ἕως τῆς ἡμέρας ἐκείνης ὅταν αὐτὸ πίνω καινὸν ἐν τῇ βασιλείᾳ τοῦ Θεοῦ.

\vs{26}Καὶ ὑμνήσαντες ἐξῆλθον εἰς τὸ ὄρος τῶν Ἐλαιῶν.
\vs{27}Καὶ λέγει αὐτοῖς ὁ Ἰησοῦς ὅτι Πάντες σκανδαλισθήσεσθε, ὅτι γέγραπται· 
\begin{poetryblock}

\begin{quote}Πατάξω τὸν ποιμένα,\end{quote} 

\begin{quote}Καὶ τὰ πρόβατα διασκορπισθήσονται.\end{quote}
\end{poetryblock}

\vs{28}Ἀλλὰ μετὰ τὸ ἐγερθῆναί με προάξω ὑμᾶς εἰς τὴν Γαλιλαίαν.
\vs{29}Ὁ δὲ Πέτρος ἔφη αὐτῷ· Εἰ καὶ πάντες σκανδαλισθήσονται, ἀλλ᾽ οὐκ ἐγώ.
\vs{30}Καὶ λέγει αὐτῷ ὁ Ἰησοῦς· Ἀμὴν λέγω σοι ὅτι σὺ σήμερον ταύτῃ τῇ νυκτὶ πρὶν ἢ δὶς ἀλέκτορα φωνῆσαι τρίς με ἀπαρνήσῃ.
\vs{31}Ὁ δὲ ἐκπερισσῶς ἐλάλει· Ἐὰν δέῃ με συναποθανεῖν σοι, οὐ μή σε ἀπαρνήσομαι. ὡσαύτως δὲ καὶ πάντες ἔλεγον.

\vs{32}Καὶ ἔρχονται εἰς χωρίον οὗ τὸ ὄνομα Γεθσημανί καὶ λέγει τοῖς μαθηταῖς αὐτοῦ· Καθίσατε ὧδε ἕως προσεύξωμαι.
\vs{33}καὶ παραλαμβάνει τὸν Πέτρον καὶ τὸν Ἰάκωβον καὶ τὸν Ἰωάννην μετ᾽ αὐτοῦ καὶ ἤρξατο ἐκθαμβεῖσθαι καὶ ἀδημονεῖν
\vs{34}καὶ λέγει αὐτοῖς· Περίλυπός ἐστιν ἡ ψυχή μου ἕως θανάτου· μείνατε ὧδε καὶ γρηγορεῖτε.
\vs{35}Καὶ προελθὼν μικρὸν ἔπιπτεν ἐπὶ τῆς γῆς καὶ προσηύχετο ἵνα εἰ δυνατόν ἐστιν παρέλθῃ ἀπ᾽ αὐτοῦ ἡ ὥρα,
\vs{36}καὶ ἔλεγεν· Ἀββᾶ ὁ Πατήρ, πάντα δυνατά σοι· παρένεγκε τὸ ποτήριον τοῦτο ἀπ᾽ ἐμοῦ· ἀλλ᾽ οὐ τί ἐγὼ θέλω ἀλλὰ τί σύ.
\vs{37}Καὶ ἔρχεται καὶ εὑρίσκει αὐτοὺς καθεύδοντας, καὶ λέγει τῷ Πέτρῳ· Σίμων, καθεύδεις; οὐκ ἴσχυσας μίαν ὥραν γρηγορῆσαι;
\vs{38}γρηγορεῖτε καὶ προσεύχεσθε, ἵνα μὴ ἔλθητε εἰς πειρασμόν· τὸ μὲν πνεῦμα πρόθυμον ἡ δὲ σὰρξ ἀσθενής.
\vs{39}Καὶ πάλιν ἀπελθὼν προσηύξατο τὸν αὐτὸν λόγον εἰπών.
\vs{40}καὶ πάλιν ἐλθὼν εὗρεν αὐτοὺς καθεύδοντας, ἦσαν γὰρ αὐτῶν οἱ ὀφθαλμοὶ καταβαρυνόμενοι, καὶ οὐκ ᾔδεισαν τί ἀποκριθῶσιν αὐτῷ.
\vs{41}Καὶ ἔρχεται τὸ τρίτον καὶ λέγει αὐτοῖς· Καθεύδετε τὸ λοιπὸν καὶ ἀναπαύεσθε· ἀπέχει· ἦλθεν ἡ ὥρα, ἰδοὺ παραδίδοται ὁ Υἱὸς τοῦ ἀνθρώπου εἰς τὰς χεῖρας τῶν ἁμαρτωλῶν.
\vs{42}ἐγείρεσθε ἄγωμεν· ἰδοὺ ὁ παραδιδούς με ἤγγικεν.

\vs{43}Καὶ εὐθὺς ἔτι αὐτοῦ λαλοῦντος παραγίνεται Ἰούδας εἷς τῶν δώδεκα καὶ μετ᾽ αὐτοῦ ὄχλος μετὰ μαχαιρῶν καὶ ξύλων παρὰ τῶν ἀρχιερέων καὶ τῶν γραμματέων καὶ τῶν πρεσβυτέρων.
\vs{44}Δεδώκει δὲ ὁ παραδιδοὺς αὐτὸν σύσσημον αὐτοῖς λέγων· Ὃν ἂν φιλήσω αὐτός ἐστιν, κρατήσατε αὐτὸν καὶ ἀπάγετε ἀσφαλῶς.
\vs{45}καὶ ἐλθὼν εὐθὺς προσελθὼν αὐτῷ λέγει· Ῥαββί, καὶ κατεφίλησεν αὐτόν·
\vs{46}Οἱ δὲ ἐπέβαλαν τὰς χεῖρας αὐτῷ καὶ ἐκράτησαν αὐτόν.
\vs{47}εἷς δέ τις τῶν παρεστηκότων σπασάμενος τὴν μάχαιραν ἔπαισεν τὸν δοῦλον τοῦ ἀρχιερέως καὶ ἀφεῖλεν αὐτοῦ τὸ ὠτάριον.

\vs{48}Καὶ ἀποκριθεὶς ὁ Ἰησοῦς εἶπεν αὐτοῖς· Ὡς ἐπὶ λῃστὴν ἐξήλθατε μετὰ μαχαιρῶν καὶ ξύλων συλλαβεῖν με;
\vs{49}καθ᾽ ἡμέραν ἤμην πρὸς ὑμᾶς ἐν τῷ ἱερῷ διδάσκων καὶ οὐκ ἐκρατήσατέ με· ἀλλ᾽ ἵνα πληρωθῶσιν αἱ γραφαί.
\vs{50}Καὶ ἀφέντες αὐτὸν ἔφυγον πάντες.
\vs{51}Καὶ νεανίσκος τις συνηκολούθει αὐτῷ περιβεβλημένος σινδόνα ἐπὶ γυμνοῦ, καὶ κρατοῦσιν αὐτόν·
\vs{52}ὁ δὲ καταλιπὼν τὴν σινδόνα γυμνὸς ἔφυγεν.
\vs{53}Καὶ ἀπήγαγον τὸν Ἰησοῦν πρὸς τὸν ἀρχιερέα, καὶ συνέρχονται πάντες οἱ ἀρχιερεῖς καὶ οἱ πρεσβύτεροι καὶ οἱ γραμματεῖς.
\vs{54}καὶ ὁ Πέτρος ἀπὸ μακρόθεν ἠκολούθησεν αὐτῷ ἕως ἔσω εἰς τὴν αὐλὴν τοῦ ἀρχιερέως καὶ ἦν συνκαθήμενος μετὰ τῶν ὑπηρετῶν καὶ θερμαινόμενος πρὸς τὸ φῶς.

\vs{55}Οἱ δὲ ἀρχιερεῖς καὶ ὅλον τὸ συνέδριον ἐζήτουν κατὰ τοῦ Ἰησοῦ μαρτυρίαν εἰς τὸ θανατῶσαι αὐτόν, καὶ οὐχ ηὕρισκον·
\vs{56}πολλοὶ γὰρ ἐψευδομαρτύρουν κατ᾽ αὐτοῦ, καὶ ἴσαι αἱ μαρτυρίαι οὐκ ἦσαν.
\vs{57}Καί τινες ἀναστάντες ἐψευδομαρτύρουν κατ᾽ αὐτοῦ λέγοντες
\vs{58}Ὅτι Ἡμεῖς ἠκούσαμεν αὐτοῦ λέγοντος ὅτι Ἐγὼ καταλύσω τὸν ναὸν τοῦτον τὸν χειροποίητον καὶ διὰ τριῶν ἡμερῶν ἄλλον ἀχειροποίητον οἰκοδομήσω.
\vs{59}καὶ οὐδὲ οὕτως ἴση ἦν ἡ μαρτυρία αὐτῶν.
\vs{60}Καὶ ἀναστὰς ὁ ἀρχιερεὺς εἰς μέσον ἐπηρώτησεν τὸν Ἰησοῦν λέγων· Οὐκ ἀποκρίνῃ οὐδέν τί οὗτοί σου καταμαρτυροῦσιν;
\vs{61}Ὁ δὲ ἐσιώπα καὶ οὐκ ἀπεκρίνατο οὐδέν. Πάλιν ὁ ἀρχιερεὺς ἐπηρώτα αὐτὸν καὶ λέγει αὐτῷ· Σὺ εἶ ὁ Χριστὸς ὁ Υἱὸς τοῦ Εὐλογητοῦ;
\vs{62}Ὁ δὲ Ἰησοῦς εἶπεν· Ἐγώ εἰμι, καὶ ὄψεσθε τὸν Υἱὸν τοῦ ἀνθρώπου ἐκ δεξιῶν καθήμενον τῆς δυνάμεως καὶ ἐρχόμενον μετὰ τῶν νεφελῶν τοῦ οὐρανοῦ.
\vs{63}Ὁ δὲ ἀρχιερεὺς διαρρήξας τοὺς χιτῶνας αὐτοῦ λέγει· Τί ἔτι χρείαν ἔχομεν μαρτύρων;
\vs{64}ἠκούσατε τῆς βλασφημίας· τί ὑμῖν φαίνεται; Οἱ δὲ πάντες κατέκριναν αὐτὸν ἔνοχον εἶναι θανάτου.

\vs{65}Καὶ ἤρξαντό τινες ἐμπτύειν αὐτῷ καὶ περικαλύπτειν αὐτοῦ τὸ πρόσωπον καὶ κολαφίζειν αὐτὸν καὶ λέγειν αὐτῷ· Προφήτευσον, καὶ οἱ ὑπηρέται ῥαπίσμασιν αὐτὸν ἔλαβον.

\vs{66}Καὶ ὄντος τοῦ Πέτρου κάτω ἐν τῇ αὐλῇ ἔρχεται μία τῶν παιδισκῶν τοῦ ἀρχιερέως
\vs{67}καὶ ἰδοῦσα τὸν Πέτρον θερμαινόμενον ἐμβλέψασα αὐτῷ λέγει· Καὶ σὺ μετὰ τοῦ Ναζαρηνοῦ ἦσθα τοῦ Ἰησοῦ.
\vs{68}Ὁ δὲ ἠρνήσατο λέγων· Οὔτε οἶδα οὔτε ἐπίσταμαι σὺ τί λέγεις. καὶ ἐξῆλθεν ἔξω εἰς τὸ προαύλιον καὶ ἀλέκτωρ ἐφώνησεν.

\vs{69}Καὶ ἡ παιδίσκη ἰδοῦσα αὐτὸν ἤρξατο πάλιν λέγειν τοῖς παρεστῶσιν ὅτι Οὗτος ἐξ αὐτῶν ἐστιν.
\vs{70}Ὁ δὲ πάλιν ἠρνεῖτο. Καὶ μετὰ μικρὸν πάλιν οἱ παρεστῶτες ἔλεγον τῷ Πέτρῳ· Ἀληθῶς ἐξ αὐτῶν εἶ, καὶ γὰρ Γαλιλαῖος εἶ.
\vs{71}Ὁ δὲ ἤρξατο ἀναθεματίζειν καὶ ὀμνύναι ὅτι Οὐκ οἶδα τὸν ἄνθρωπον τοῦτον ὃν λέγετε.
\vs{72}καὶ εὐθὺς ἐκ δευτέρου ἀλέκτωρ ἐφώνησεν. Καὶ ἀνεμνήσθη ὁ Πέτρος τὸ ῥῆμα ὡς εἶπεν αὐτῷ ὁ Ἰησοῦς ὅτι Πρὶν ἀλέκτορα φωνῆσαι δὶς τρίς με ἀπαρνήσῃ· καὶ ἐπιβαλὼν ἔκλαιεν.

\ch{15}
Καὶ εὐθὺς πρωῒ συμβούλιον ποιήσαντες οἱ ἀρχιερεῖς μετὰ τῶν πρεσβυτέρων καὶ γραμματέων καὶ ὅλον τὸ συνέδριον, δήσαντες τὸν Ἰησοῦν ἀπήνεγκαν καὶ παρέδωκαν Πιλάτῳ.

\vs{2}Καὶ ἐπηρώτησεν αὐτὸν ὁ Πιλᾶτος· Σὺ εἶ ὁ Βασιλεὺς τῶν Ἰουδαίων; Ὁ δὲ ἀποκριθεὶς αὐτῷ λέγει· Σὺ λέγεις.
\vs{3}Καὶ κατηγόρουν αὐτοῦ οἱ ἀρχιερεῖς πολλά.
\vs{4}Ὁ δὲ Πιλᾶτος πάλιν ἐπηρώτα αὐτὸν λέγων· Οὐκ ἀποκρίνῃ οὐδέν; ἴδε πόσα σου κατηγοροῦσιν.
\vs{5}Ὁ δὲ Ἰησοῦς οὐκέτι οὐδὲν ἀπεκρίθη, ὥστε θαυμάζειν τὸν Πιλᾶτον.

\vs{6}Κατὰ δὲ ἑορτὴν ἀπέλυεν αὐτοῖς ἕνα δέσμιον ὃν παρῃτοῦντο.
\vs{7}ἦν δὲ ὁ λεγόμενος Βαραββᾶς μετὰ τῶν στασιαστῶν δεδεμένος οἵτινες ἐν τῇ στάσει φόνον πεποιήκεισαν.
\vs{8}καὶ ἀναβὰς ὁ ὄχλος ἤρξατο αἰτεῖσθαι καθὼς ἐποίει αὐτοῖς.
\vs{9}Ὁ δὲ Πιλᾶτος ἀπεκρίθη αὐτοῖς λέγων· Θέλετε ἀπολύσω ὑμῖν τὸν Βασιλέα τῶν Ἰουδαίων;
\vs{10}ἐγίνωσκεν γὰρ ὅτι διὰ φθόνον παραδεδώκεισαν αὐτὸν οἱ ἀρχιερεῖς.
\vs{11}Οἱ δὲ ἀρχιερεῖς ἀνέσεισαν τὸν ὄχλον ἵνα μᾶλλον τὸν Βαραββᾶν ἀπολύσῃ αὐτοῖς.
\vs{12}Ὁ δὲ Πιλᾶτος πάλιν ἀποκριθεὶς ἔλεγεν αὐτοῖς· Τί οὖν θέλετε ποιήσω ὃν λέγετε τὸν Βασιλέα τῶν Ἰουδαίων;
\vs{13}Οἱ δὲ πάλιν ἔκραξαν· Σταύρωσον αὐτόν.
\vs{14}Ὁ δὲ Πιλᾶτος ἔλεγεν αὐτοῖς· Τί γὰρ ἐποίησεν κακόν; Οἱ δὲ περισσῶς ἔκραξαν· Σταύρωσον αὐτόν.

\vs{15}Ὁ δὲ Πιλᾶτος βουλόμενος τῷ ὄχλῳ τὸ ἱκανὸν ποιῆσαι ἀπέλυσεν αὐτοῖς τὸν Βαραββᾶν, καὶ παρέδωκεν τὸν Ἰησοῦν φραγελλώσας ἵνα σταυρωθῇ.

\vs{16}Οἱ δὲ στρατιῶται ἀπήγαγον αὐτὸν ἔσω τῆς αὐλῆς, ὅ ἐστιν Πραιτώριον, καὶ συνκαλοῦσιν ὅλην τὴν σπεῖραν.
\vs{17}καὶ ἐνδιδύσκουσιν αὐτὸν πορφύραν καὶ περιτιθέασιν αὐτῷ πλέξαντες ἀκάνθινον στέφανον·
\vs{18}καὶ ἤρξαντο ἀσπάζεσθαι αὐτόν· Χαῖρε, Βασιλεῦ τῶν Ἰουδαίων·
\vs{19}Καὶ ἔτυπτον αὐτοῦ τὴν κεφαλὴν καλάμῳ καὶ ἐνέπτυον αὐτῷ καὶ τιθέντες τὰ γόνατα προσεκύνουν αὐτῷ.
\vs{20}καὶ ὅτε ἐνέπαιξαν αὐτῷ, ἐξέδυσαν αὐτὸν τὴν πορφύραν καὶ ἐνέδυσαν αὐτὸν τὰ ἱμάτια αὐτοῦ.

Καὶ ἐξάγουσιν αὐτὸν ἵνα σταυρώσωσιν αὐτόν.
\vs{21}Καὶ ἀγγαρεύουσιν παράγοντά τινα Σίμωνα Κυρηναῖον ἐρχόμενον ἀπ᾽ ἀγροῦ, τὸν πατέρα Ἀλεξάνδρου καὶ Ῥούφου, ἵνα ἄρῃ τὸν σταυρὸν αὐτοῦ.

\vs{22}Καὶ φέρουσιν αὐτὸν ἐπὶ τὸν Γολγοθᾶν τόπον, ὅ ἐστιν μεθερμηνευόμενον Κρανίου τόπος.
\vs{23}καὶ ἐδίδουν αὐτῷ ἐσμυρνισμένον οἶνον· ὃς δὲ οὐκ ἔλαβεν.
\vs{24}Καὶ σταυροῦσιν αὐτὸν καὶ διαμερίζονται τὰ ἱμάτια αὐτοῦ βάλλοντες κλῆρον ἐπ᾽ αὐτὰ τίς τί ἄρῃ.
\vs{25}Ἦν δὲ ὥρα τρίτη καὶ ἐσταύρωσαν αὐτόν.
\vs{26}καὶ ἦν ἡ ἐπιγραφὴ τῆς αἰτίας αὐτοῦ ἐπιγεγραμμένη· 
\begin{poetryblock}

\begin{quote}Ο ΒΑΣΙΛΕΥΣ ΤΩΝ ΙΟΥΔΑΙΩΝ.\end{quote}
\end{poetryblock}

\vs{27}Καὶ σὺν αὐτῷ σταυροῦσιν δύο λῃστάς, ἕνα ἐκ δεξιῶν καὶ ἕνα ἐξ εὐωνύμων αὐτοῦ.
\vs{29}Καὶ οἱ παραπορευόμενοι ἐβλασφήμουν αὐτὸν κινοῦντες τὰς κεφαλὰς αὐτῶν καὶ λέγοντες· Οὐὰ ὁ καταλύων τὸν ναὸν καὶ οἰκοδομῶν ἐν τρισὶν ἡμέραις,
\vs{30}σῶσον σεαυτὸν καταβὰς ἀπὸ τοῦ σταυροῦ.
\vs{31}Ὁμοίως καὶ οἱ ἀρχιερεῖς ἐμπαίζοντες πρὸς ἀλλήλους μετὰ τῶν γραμματέων ἔλεγον· Ἄλλους ἔσωσεν, ἑαυτὸν οὐ δύναται σῶσαι·
\vs{32}ὁ Χριστὸς ὁ Βασιλεὺς Ἰσραὴλ καταβάτω νῦν ἀπὸ τοῦ σταυροῦ, ἵνα ἴδωμεν καὶ πιστεύσωμεν. καὶ οἱ συνεσταυρωμένοι σὺν αὐτῷ ὠνείδιζον αὐτόν.

\vs{33}Καὶ γενομένης ὥρας ἕκτης σκότος ἐγένετο ἐφ᾽ ὅλην τὴν γῆν ἕως ὥρας ἐνάτης.
\vs{34}καὶ τῇ ἐνάτῃ ὥρᾳ ἐβόησεν ὁ Ἰησοῦς φωνῇ μεγάλῃ· 
\begin{poetryblock}

\begin{quote}Ἐλωῒ Ἐλωῒ λεμὰ σαβαχθάνι;\end{quote}
\end{poetryblock}

ὅ ἐστιν μεθερμηνευόμενον· Ὁ Θεός μου ὁ Θεός μου, εἰς τί ἐγκατέλιπές με;

\vs{35}Καί τινες τῶν παρεστηκότων ἀκούσαντες ἔλεγον· Ἴδε Ἠλίαν φωνεῖ.
\vs{36}Δραμὼν δέ τις καὶ γεμίσας σπόγγον ὄξους περιθεὶς καλάμῳ ἐπότιζεν αὐτόν λέγων· Ἄφετε ἴδωμεν εἰ ἔρχεται Ἠλίας καθελεῖν αὐτόν.
\vs{37}Ὁ δὲ Ἰησοῦς ἀφεὶς φωνὴν μεγάλην ἐξέπνευσεν.

\vs{38}Καὶ τὸ καταπέτασμα τοῦ ναοῦ ἐσχίσθη εἰς δύο ἀπ᾽ ἄνωθεν ἕως κάτω.
\vs{39}Ἰδὼν δὲ ὁ κεντυρίων ὁ παρεστηκὼς ἐξ ἐναντίας αὐτοῦ ὅτι οὕτως ἐξέπνευσεν εἶπεν· Ἀληθῶς οὗτος ὁ ἄνθρωπος Υἱὸς Θεοῦ ἦν.

\vs{40}Ἦσαν δὲ καὶ γυναῖκες ἀπὸ μακρόθεν θεωροῦσαι, ἐν αἷς καὶ Μαρία ἡ Μαγδαληνὴ καὶ Μαρία ἡ Ἰακώβου τοῦ μικροῦ καὶ Ἰωσῆτος μήτηρ καὶ Σαλώμη,
\vs{41}αἳ ὅτε ἦν ἐν τῇ Γαλιλαίᾳ ἠκολούθουν αὐτῷ καὶ διηκόνουν αὐτῷ, καὶ ἄλλαι πολλαὶ αἱ συναναβᾶσαι αὐτῷ εἰς Ἱεροσόλυμα.

\vs{42}Καὶ ἤδη ὀψίας γενομένης, ἐπεὶ ἦν Παρασκευή ὅ ἐστιν προσάββατον,
\vs{43}ἐλθὼν Ἰωσὴφ ὁ ἀπὸ Ἁριμαθαίας εὐσχήμων βουλευτής, ὃς καὶ αὐτὸς ἦν προσδεχόμενος τὴν βασιλείαν τοῦ Θεοῦ, τολμήσας εἰσῆλθεν πρὸς τὸν Πιλᾶτον καὶ ᾐτήσατο τὸ σῶμα τοῦ Ἰησοῦ.
\vs{44}Ὁ δὲ Πιλᾶτος ἐθαύμασεν εἰ ἤδη τέθνηκεν καὶ προσκαλεσάμενος τὸν κεντυρίωνα ἐπηρώτησεν αὐτὸν εἰ πάλαι ἀπέθανεν·
\vs{45}καὶ γνοὺς ἀπὸ τοῦ κεντυρίωνος ἐδωρήσατο τὸ πτῶμα τῷ Ἰωσήφ.
\vs{46}Καὶ ἀγοράσας σινδόνα καθελὼν αὐτὸν ἐνείλησεν τῇ σινδόνι καὶ ἔθηκεν αὐτὸν ἐν μνημείῳ ὃ ἦν λελατομημένον ἐκ πέτρας καὶ προσεκύλισεν λίθον ἐπὶ τὴν θύραν τοῦ μνημείου.
\vs{47}ἡ δὲ Μαρία ἡ Μαγδαληνὴ καὶ Μαρία ἡ Ἰωσῆτος ἐθεώρουν ποῦ τέθειται.

\ch{16}
Καὶ διαγενομένου τοῦ σαββάτου Μαρία ἡ Μαγδαληνὴ καὶ Μαρία ἡ τοῦ Ἰακώβου καὶ Σαλώμη ἠγόρασαν ἀρώματα ἵνα ἐλθοῦσαι ἀλείψωσιν αὐτόν.
\vs{2}καὶ λίαν πρωῒ τῇ μιᾷ τῶν σαββάτων ἔρχονται ἐπὶ τὸ μνημεῖον ἀνατείλαντος τοῦ ἡλίου.
\vs{3}καὶ ἔλεγον πρὸς ἑαυτάς· Τίς ἀποκυλίσει ἡμῖν τὸν λίθον ἐκ τῆς θύρας τοῦ μνημείου;
\vs{4}καὶ ἀναβλέψασαι θεωροῦσιν ὅτι ἀποκεκύλισται ὁ λίθος· ἦν γὰρ μέγας σφόδρα.

\vs{5}Καὶ εἰσελθοῦσαι εἰς τὸ μνημεῖον εἶδον νεανίσκον καθήμενον ἐν τοῖς δεξιοῖς περιβεβλημένον στολὴν λευκήν, καὶ ἐξεθαμβήθησαν.
\vs{6}ὁ δὲ λέγει αὐταῖς· Μὴ ἐκθαμβεῖσθε· Ἰησοῦν ζητεῖτε τὸν Ναζαρηνὸν τὸν ἐσταυρωμένον· ἠγέρθη, οὐκ ἔστιν ὧδε· ἴδε ὁ τόπος ὅπου ἔθηκαν αὐτόν.
\vs{7}ἀλλὰ ὑπάγετε εἴπατε τοῖς μαθηταῖς αὐτοῦ καὶ τῷ Πέτρῳ ὅτι Προάγει ὑμᾶς εἰς τὴν Γαλιλαίαν· ἐκεῖ αὐτὸν ὄψεσθε, καθὼς εἶπεν ὑμῖν.

\vs{8}Καὶ ἐξελθοῦσαι ἔφυγον ἀπὸ τοῦ μνημείου, εἶχεν γὰρ αὐτὰς τρόμος καὶ ἔκστασις· καὶ οὐδενὶ οὐδὲν εἶπαν· ἐφοβοῦντο γάρ.

\vs{9}[[Ἀναστὰς δὲ πρωῒ πρώτῃ σαββάτου ἐφάνη πρῶτον Μαρίᾳ τῇ Μαγδαληνῇ, παρ᾽ ἧς ἐκβεβλήκει ἑπτὰ δαιμόνια.
\vs{10}ἐκείνη πορευθεῖσα ἀπήγγειλεν τοῖς μετ᾽ αὐτοῦ γενομένοις πενθοῦσι καὶ κλαίουσιν·
\vs{11}κἀκεῖνοι ἀκούσαντες ὅτι ζῇ καὶ ἐθεάθη ὑπ᾽ αὐτῆς ἠπίστησαν.

\vs{12}Μετὰ δὲ ταῦτα δυσὶν ἐξ αὐτῶν περιπατοῦσιν ἐφανερώθη ἐν ἑτέρᾳ μορφῇ πορευομένοις εἰς ἀγρόν·
\vs{13}κἀκεῖνοι ἀπελθόντες ἀπήγγειλαν τοῖς λοιποῖς· οὐδὲ ἐκείνοις ἐπίστευσαν.

\vs{14}Ὕστερον δὲ ἀνακειμένοις αὐτοῖς τοῖς ἕνδεκα ἐφανερώθη καὶ ὠνείδισεν τὴν ἀπιστίαν αὐτῶν καὶ σκληροκαρδίαν ὅτι τοῖς θεασαμένοις αὐτὸν ἐγηγερμένον οὐκ ἐπίστευσαν.
\vs{15}Καὶ εἶπεν αὐτοῖς· Πορευθέντες εἰς τὸν κόσμον ἅπαντα κηρύξατε τὸ εὐαγγέλιον πάσῃ τῇ κτίσει.
\vs{16}ὁ πιστεύσας καὶ βαπτισθεὶς σωθήσεται, ὁ δὲ ἀπιστήσας κατακριθήσεται.
\vs{17}σημεῖα δὲ τοῖς πιστεύσασιν ταῦτα παρακολουθήσει· ἐν τῷ ὀνόματί μου δαιμόνια ἐκβαλοῦσιν, γλώσσαις λαλήσουσιν καιναῖς,
\vs{18}και ἐν ταῖς χερσὶν ὄφεις ἀροῦσιν κἂν θανάσιμόν τι πίωσιν οὐ μὴ αὐτοὺς βλάψῃ, ἐπὶ ἀρρώστους χεῖρας ἐπιθήσουσιν καὶ καλῶς ἕξουσιν.

\vs{19}Ὁ μὲν οὖν Κύριος Ἰησοῦς μετὰ τὸ λαλῆσαι αὐτοῖς ἀνελήμφθη εἰς τὸν οὐρανὸν καὶ ἐκάθισεν ἐκ δεξιῶν τοῦ Θεοῦ.
\vs{20}ἐκεῖνοι δὲ ἐξελθόντες ἐκήρυξαν πανταχοῦ, τοῦ Κυρίου συνεργοῦντος καὶ τὸν λόγον βεβαιοῦντος διὰ τῶν ἐπακολουθούντων σημείων.]]

[[πάντα δὲ τὰ παρηγγελμένα τοῖς περὶ τὸν Πέτρον συντόμως ἐξήγγειλαν. μετὰ δὲ ταῦτα καὶ αὐτὸς ὁ Ἰησοῦς ἀπὸ ἀνατολῆς καὶ ἄχρι δύσεως ἐξαπέστειλεν δι᾽ αὐτῶν τὸ ἱερὸν καὶ ἄφθαρτον κήρυγμα τῆς αἰωνίου σωτηρίας. ἀμήν.]]


\def\book{ΚΑΤΑ ΛΟΥΚΑΝ}
\biblebook{ΚΑΤΑ ΛΟΥΚΑΝ}


\lettrine[lines=2, loversize=0.2, nindent=0em, findent=.25em]{\textcolor{bookheadingcolor}{Ἐ}}{πειδήπερ} πολλοὶ ἐπεχείρησαν ἀνατάξασθαι διήγησιν περὶ τῶν πεπληροφορημένων ἐν ἡμῖν πραγμάτων,
\vs{2}καθὼς παρέδοσαν ἡμῖν οἱ ἀπ᾽ ἀρχῆς αὐτόπται καὶ ὑπηρέται γενόμενοι τοῦ λόγου,
\vs{3}ἔδοξε κἀμοὶ παρηκολουθηκότι ἄνωθεν πᾶσιν ἀκριβῶς καθεξῆς σοι γράψαι, κράτιστε Θεόφιλε,
\vs{4}ἵνα ἐπιγνῷς περὶ ὧν κατηχήθης λόγων τὴν ἀσφάλειαν.

\vs{5}Ἐγένετο ἐν ταῖς ἡμέραις Ἡρῴδου βασιλέως τῆς Ἰουδαίας ἱερεύς τις ὀνόματι Ζαχαρίας ἐξ ἐφημερίας Ἀβιά, καὶ γυνὴ αὐτῷ ἐκ τῶν θυγατέρων Ἀαρών καὶ τὸ ὄνομα αὐτῆς Ἐλισάβετ.
\vs{6}ἦσαν δὲ δίκαιοι ἀμφότεροι ἐναντίον τοῦ Θεοῦ, πορευόμενοι ἐν πάσαις ταῖς ἐντολαῖς καὶ δικαιώμασιν τοῦ Κυρίου ἄμεμπτοι.
\vs{7}καὶ οὐκ ἦν αὐτοῖς τέκνον, καθότι ἦν ἡ Ἐλισάβετ στεῖρα, καὶ ἀμφότεροι προβεβηκότες ἐν ταῖς ἡμέραις αὐτῶν ἦσαν.

\vs{8}Ἐγένετο δὲ ἐν τῷ ἱερατεύειν αὐτὸν ἐν τῇ τάξει τῆς ἐφημερίας αὐτοῦ ἔναντι τοῦ Θεοῦ,
\vs{9}κατὰ τὸ ἔθος τῆς ἱερατείας ἔλαχε τοῦ θυμιᾶσαι εἰσελθὼν εἰς τὸν ναὸν τοῦ Κυρίου,
\vs{10}καὶ πᾶν τὸ πλῆθος ἦν τοῦ λαοῦ προσευχόμενον ἔξω τῇ ὥρᾳ τοῦ θυμιάματος.
\vs{11}Ὤφθη δὲ αὐτῷ ἄγγελος Κυρίου ἑστὼς ἐκ δεξιῶν τοῦ θυσιαστηρίου τοῦ θυμιάματος.
\vs{12}καὶ ἐταράχθη Ζαχαρίας ἰδών καὶ φόβος ἐπέπεσεν ἐπ᾽ αὐτόν.
\vs{13}Εἶπεν δὲ πρὸς αὐτὸν ὁ ἄγγελος· Μὴ φοβοῦ, Ζαχαρία, διότι εἰσηκούσθη ἡ δέησίς σου, καὶ ἡ γυνή σου Ἐλισάβετ γεννήσει υἱόν σοι καὶ καλέσεις τὸ ὄνομα αὐτοῦ Ἰωάννην.
\begin{poetryblock}

\begin{quote} \vs{14}καὶ ἔσται χαρά σοι καὶ ἀγαλλίασις\end{quote} 

\begin{quote}καὶ πολλοὶ ἐπὶ τῇ γενέσει αὐτοῦ χαρήσονται.\end{quote}

\begin{quote} \vs{15}ἔσται γὰρ μέγας ἐνώπιον τοῦ Κυρίου,\end{quote} 

\begin{quote}καὶ οἶνον καὶ σίκερα οὐ μὴ πίῃ,\end{quote} 

\begin{quote}καὶ Πνεύματος Ἁγίου πλησθήσεται\end{quote} 

\begin{quote}ἔτι ἐκ κοιλίας μητρὸς αὐτοῦ,\end{quote}

\begin{quote} \vs{16}καὶ πολλοὺς τῶν υἱῶν Ἰσραὴλ ἐπιστρέψει\end{quote} 

\begin{quote}ἐπὶ Κύριον τὸν Θεὸν αὐτῶν.\end{quote}

\begin{quote} \vs{17}καὶ αὐτὸς προελεύσεται ἐνώπιον αὐτοῦ\end{quote} 

\begin{quote}ἐν πνεύματι καὶ δυνάμει Ἠλίου,\end{quote} 

\begin{quote}ἐπιστρέψαι καρδίας πατέρων ἐπὶ τέκνα\end{quote} 

\begin{quote}καὶ ἀπειθεῖς ἐν φρονήσει δικαίων,\end{quote} 

\begin{quote}ἑτοιμάσαι Κυρίῳ λαὸν κατεσκευασμένον.\end{quote}
\end{poetryblock}

\vs{18}Καὶ εἶπεν Ζαχαρίας πρὸς τὸν ἄγγελον· Κατὰ τί γνώσομαι τοῦτο; ἐγὼ γάρ εἰμι πρεσβύτης καὶ ἡ γυνή μου προβεβηκυῖα ἐν ταῖς ἡμέραις αὐτῆς.
\vs{19}Καὶ ἀποκριθεὶς ὁ ἄγγελος εἶπεν αὐτῷ· Ἐγώ εἰμι Γαβριὴλ ὁ παρεστηκὼς ἐνώπιον τοῦ Θεοῦ καὶ ἀπεστάλην λαλῆσαι πρὸς σὲ καὶ εὐαγγελίσασθαί σοι ταῦτα·
\vs{20}καὶ ἰδοὺ ἔσῃ σιωπῶν καὶ μὴ δυνάμενος λαλῆσαι ἄχρι ἧς ἡμέρας γένηται ταῦτα, ἀνθ᾽ ὧν οὐκ ἐπίστευσας τοῖς λόγοις μου, οἵτινες πληρωθήσονται εἰς τὸν καιρὸν αὐτῶν.

\vs{21}Καὶ ἦν ὁ λαὸς προσδοκῶν τὸν Ζαχαρίαν καὶ ἐθαύμαζον ἐν τῷ χρονίζειν ἐν τῷ ναῷ αὐτόν.
\vs{22}ἐξελθὼν δὲ οὐκ ἐδύνατο λαλῆσαι αὐτοῖς, καὶ ἐπέγνωσαν ὅτι ὀπτασίαν ἑώρακεν ἐν τῷ ναῷ· καὶ αὐτὸς ἦν διανεύων αὐτοῖς καὶ διέμενεν κωφός.
\vs{23}καὶ ἐγένετο ὡς ἐπλήσθησαν αἱ ἡμέραι τῆς λειτουργίας αὐτοῦ, ἀπῆλθεν εἰς τὸν οἶκον αὐτοῦ.
\vs{24}Μετὰ δὲ ταύτας τὰς ἡμέρας συνέλαβεν Ἐλισάβετ ἡ γυνὴ αὐτοῦ καὶ περιέκρυβεν ἑαυτὴν μῆνας πέντε λέγουσα
\vs{25}ὅτι Οὕτως μοι πεποίηκεν Κύριος ἐν ἡμέραις αἷς ἐπεῖδεν ἀφελεῖν ὄνειδός μου ἐν ἀνθρώποις.
\vs{26}Ἐν δὲ τῷ μηνὶ τῷ ἕκτῳ ἀπεστάλη ὁ ἄγγελος Γαβριὴλ ἀπὸ τοῦ Θεοῦ εἰς πόλιν τῆς Γαλιλαίας ᾗ ὄνομα Ναζαρὲθ
\vs{27}πρὸς παρθένον ἐμνηστευμένην ἀνδρὶ ᾧ ὄνομα Ἰωσὴφ ἐξ οἴκου Δαυὶδ καὶ τὸ ὄνομα τῆς παρθένου Μαριάμ.
\vs{28}καὶ εἰσελθὼν πρὸς αὐτὴν εἶπεν· Χαῖρε, κεχαριτωμένη, ὁ Κύριος μετὰ σοῦ.
\vs{29}Ἡ δὲ ἐπὶ τῷ λόγῳ διεταράχθη καὶ διελογίζετο ποταπὸς εἴη ὁ ἀσπασμὸς οὗτος.
\vs{30}καὶ εἶπεν ὁ ἄγγελος αὐτῇ· 
\begin{poetryblock}

\begin{quote}Μὴ φοβοῦ, Μαριάμ, εὗρες γὰρ χάριν παρὰ τῷ Θεῷ.\end{quote}

\begin{quote} \vs{31}καὶ ἰδοὺ συλλήμψῃ ἐν γαστρὶ καὶ τέξῃ υἱόν\end{quote} 

\begin{quote}καὶ καλέσεις τὸ ὄνομα αὐτοῦ Ἰησοῦν.\end{quote}

\begin{quote} \vs{32}οὗτος ἔσται μέγας καὶ Υἱὸς Ὑψίστου κληθήσεται\end{quote} 

\begin{quote}καὶ δώσει αὐτῷ Κύριος ὁ Θεὸς τὸν θρόνον Δαυὶδ τοῦ πατρὸς αὐτοῦ,\end{quote}

\begin{quote} \vs{33}καὶ βασιλεύσει ἐπὶ τὸν οἶκον Ἰακὼβ εἰς τοὺς αἰῶνας\end{quote} 

\begin{quote}καὶ τῆς βασιλείας αὐτοῦ οὐκ ἔσται τέλος.\end{quote}
\end{poetryblock}

\vs{34}Εἶπεν δὲ Μαριὰμ πρὸς τὸν ἄγγελον· Πῶς ἔσται τοῦτο, ἐπεὶ ἄνδρα οὐ γινώσκω;
\vs{35}Καὶ ἀποκριθεὶς ὁ ἄγγελος εἶπεν αὐτῇ· 
\begin{poetryblock}

\begin{quote}Πνεῦμα Ἅγιον ἐπελεύσεται ἐπὶ σέ\end{quote} 

\begin{quote}καὶ δύναμις Ὑψίστου ἐπισκιάσει σοι·\end{quote} 

\begin{quote}διὸ καὶ τὸ γεννώμενον ἅγιον κληθήσεται Υἱὸς Θεοῦ.\end{quote}
\end{poetryblock}

\vs{36}καὶ ἰδοὺ Ἐλισάβετ ἡ συγγενίς σου καὶ αὐτὴ συνείληφεν υἱὸν ἐν γήρει αὐτῆς καὶ οὗτος μὴν ἕκτος ἐστὶν αὐτῇ τῇ καλουμένῃ στείρᾳ·
\vs{37}ὅτι οὐκ ἀδυνατήσει παρὰ τοῦ Θεοῦ πᾶν ῥῆμα.
\vs{38}Εἶπεν δὲ Μαριάμ· Ἰδοὺ ἡ δούλη Κυρίου· γένοιτό μοι κατὰ τὸ ῥῆμά σου. καὶ ἀπῆλθεν ἀπ᾽ αὐτῆς ὁ ἄγγελος.

\vs{39}Ἀναστᾶσα δὲ Μαριὰμ ἐν ταῖς ἡμέραις ταύταις ἐπορεύθη εἰς τὴν ὀρεινὴν μετὰ σπουδῆς εἰς πόλιν Ἰούδα,
\vs{40}καὶ εἰσῆλθεν εἰς τὸν οἶκον Ζαχαρίου καὶ ἠσπάσατο τὴν Ἐλισάβετ.
\vs{41}καὶ ἐγένετο ὡς ἤκουσεν τὸν ἀσπασμὸν τῆς Μαρίας ἡ Ἐλισάβετ, ἐσκίρτησεν τὸ βρέφος ἐν τῇ κοιλίᾳ αὐτῆς, καὶ ἐπλήσθη Πνεύματος Ἁγίου ἡ Ἐλισάβετ,
\vs{42}καὶ ἀνεφώνησεν κραυγῇ μεγάλῃ καὶ εἶπεν· 
\begin{poetryblock}

\begin{quote}Εὐλογημένη σὺ ἐν γυναιξίν\end{quote} 

\begin{quote}καὶ εὐλογημένος ὁ καρπὸς τῆς κοιλίας σου.\end{quote}
\end{poetryblock}

\vs{43}καὶ πόθεν μοι τοῦτο ἵνα ἔλθῃ ἡ μήτηρ τοῦ Κυρίου μου πρὸς ἐμέ;
\vs{44}ἰδοὺ γὰρ ὡς ἐγένετο ἡ φωνὴ τοῦ ἀσπασμοῦ σου εἰς τὰ ὦτά μου, ἐσκίρτησεν ἐν ἀγαλλιάσει τὸ βρέφος ἐν τῇ κοιλίᾳ μου.
\vs{45}καὶ μακαρία ἡ πιστεύσασα ὅτι ἔσται τελείωσις τοῖς λελαλημένοις αὐτῇ παρὰ Κυρίου.

\vs{46}Καὶ εἶπεν Μαριάμ· 
\begin{poetryblock}

\begin{quote}Μεγαλύνει ἡ ψυχή μου τὸν Κύριον,\end{quote}

\begin{quote} \vs{47}καὶ ἠγαλλίασεν τὸ πνεῦμά μου ἐπὶ τῷ Θεῷ τῷ Σωτῆρί μου,\end{quote}

\begin{quote} \vs{48}ὅτι ἐπέβλεψεν ἐπὶ τὴν ταπείνωσιν τῆς δούλης αὐτοῦ.\end{quote} 

\begin{quote}ἰδοὺ γὰρ ἀπὸ τοῦ νῦν μακαριοῦσίν με πᾶσαι αἱ γενεαί,\end{quote}

\begin{quote} \vs{49}ὅτι ἐποίησέν μοι μεγάλα ὁ δυνατός.\end{quote} 

\begin{quote}καὶ ἅγιον τὸ ὄνομα αὐτοῦ,\end{quote}

\begin{quote} \vs{50}καὶ τὸ ἔλεος αὐτοῦ εἰς γενεὰς καὶ γενεὰς\end{quote} 

\begin{quote}τοῖς φοβουμένοις αὐτόν.\end{quote}

\begin{quote} \vs{51}Ἐποίησεν κράτος ἐν βραχίονι αὐτοῦ,\end{quote} 

\begin{quote}διεσκόρπισεν ὑπερηφάνους διανοίᾳ καρδίας αὐτῶν·\end{quote}

\begin{quote} \vs{52}καθεῖλεν δυνάστας ἀπὸ θρόνων\end{quote} 

\begin{quote}καὶ ὕψωσεν ταπεινούς,\end{quote}

\begin{quote} \vs{53}πεινῶντας ἐνέπλησεν ἀγαθῶν\end{quote} 

\begin{quote}καὶ πλουτοῦντας ἐξαπέστειλεν κενούς.\end{quote}

\begin{quote} \vs{54}ἀντελάβετο Ἰσραὴλ παιδὸς αὐτοῦ,\end{quote} 

\begin{quote}μνησθῆναι ἐλέους,\end{quote}

\begin{quote} \vs{55}καθὼς ἐλάλησεν πρὸς τοὺς πατέρας ἡμῶν,\end{quote} 

\begin{quote}τῷ Ἀβραὰμ καὶ τῷ σπέρματι αὐτοῦ εἰς τὸν αἰῶνα.\end{quote}
\end{poetryblock}

\vs{56}Ἔμεινεν δὲ Μαριὰμ σὺν αὐτῇ ὡς μῆνας τρεῖς, καὶ ὑπέστρεψεν εἰς τὸν οἶκον αὐτῆς.

\vs{57}Τῇ δὲ Ἐλισάβετ ἐπλήσθη ὁ χρόνος τοῦ τεκεῖν αὐτήν καὶ ἐγέννησεν υἱόν.
\vs{58}καὶ ἤκουσαν οἱ περίοικοι καὶ οἱ συγγενεῖς αὐτῆς ὅτι ἐμεγάλυνεν Κύριος τὸ ἔλεος αὐτοῦ μετ᾽ αὐτῆς καὶ συνέχαιρον αὐτῇ.
\vs{59}Καὶ ἐγένετο ἐν τῇ ἡμέρᾳ τῇ ὀγδόῃ ἦλθον περιτεμεῖν τὸ παιδίον καὶ ἐκάλουν αὐτὸ ἐπὶ τῷ ὀνόματι τοῦ πατρὸς αὐτοῦ Ζαχαρίαν.
\vs{60}καὶ ἀποκριθεῖσα ἡ μήτηρ αὐτοῦ εἶπεν· Οὐχί, ἀλλὰ κληθήσεται Ἰωάννης.
\vs{61}Καὶ εἶπαν πρὸς αὐτὴν ὅτι Οὐδείς ἐστιν ἐκ τῆς συγγενείας σου ὃς καλεῖται τῷ ὀνόματι τούτῳ.
\vs{62}ἐνένευον δὲ τῷ πατρὶ αὐτοῦ τὸ τί ἂν θέλοι καλεῖσθαι αὐτό.
\vs{63}Καὶ αἰτήσας πινακίδιον ἔγραψεν λέγων· Ἰωάννης ἐστὶν ὄνομα αὐτοῦ. καὶ ἐθαύμασαν πάντες.
\vs{64}ἀνεῴχθη δὲ τὸ στόμα αὐτοῦ παραχρῆμα καὶ ἡ γλῶσσα αὐτοῦ, καὶ ἐλάλει εὐλογῶν τὸν Θεόν.
\vs{65}Καὶ ἐγένετο ἐπὶ πάντας φόβος τοὺς περιοικοῦντας αὐτούς, καὶ ἐν ὅλῃ τῇ ὀρεινῇ τῆς Ἰουδαίας διελαλεῖτο πάντα τὰ ῥήματα ταῦτα,
\vs{66}καὶ ἔθεντο πάντες οἱ ἀκούσαντες ἐν τῇ καρδίᾳ αὐτῶν λέγοντες· Τί ἄρα τὸ παιδίον τοῦτο ἔσται; καὶ γὰρ χεὶρ Κυρίου ἦν μετ᾽ αὐτοῦ.

\vs{67}Καὶ Ζαχαρίας ὁ πατὴρ αὐτοῦ ἐπλήσθη Πνεύματος Ἁγίου καὶ ἐπροφήτευσεν λέγων·
\begin{poetryblock}

\begin{quote} \vs{68}Εὐλογητὸς Κύριος ὁ Θεὸς τοῦ Ἰσραήλ,\end{quote} 

\begin{quote}ὅτι ἐπεσκέψατο καὶ ἐποίησεν λύτρωσιν τῷ λαῷ αὐτοῦ,\end{quote}

\begin{quote} \vs{69}καὶ ἤγειρεν κέρας σωτηρίας ἡμῖν\end{quote} 

\begin{quote}ἐν οἴκῳ Δαυὶδ παιδὸς αὐτοῦ,\end{quote}

\begin{quote} \vs{70}καθὼς ἐλάλησεν διὰ στόματος τῶν ἁγίων ἀπ᾽ αἰῶνος προφητῶν αὐτοῦ,\end{quote}

\begin{quote} \vs{71}σωτηρίαν ἐξ ἐχθρῶν ἡμῶν καὶ ἐκ χειρὸς πάντων τῶν μισούντων ἡμᾶς,\end{quote}

\begin{quote} \vs{72}ποιῆσαι ἔλεος μετὰ τῶν πατέρων ἡμῶν\end{quote} 

\begin{quote}καὶ μνησθῆναι διαθήκης ἁγίας αὐτοῦ,\end{quote}

\begin{quote} \vs{73}ὅρκον ὃν ὤμοσεν πρὸς Ἀβραὰμ τὸν πατέρα ἡμῶν,\end{quote} 

\begin{quote}τοῦ δοῦναι ἡμῖν\end{quote}
\end{poetryblock}

\vs{74}ἀφόβως ἐκ χειρὸς ἐχθρῶν ῥυσθέντας 
\begin{poetryblock}

\begin{quote}λατρεύειν αὐτῷ\end{quote}
\end{poetryblock}

\vs{75}ἐν ὁσιότητι καὶ δικαιοσύνῃ 
\begin{poetryblock}

\begin{quote}ἐνώπιον αὐτοῦ πάσαις ταῖς ἡμέραις ἡμῶν.\end{quote}

\begin{quote} \vs{76}Καὶ σὺ δέ, παιδίον, προφήτης Ὑψίστου κληθήσῃ·\end{quote} 

\begin{quote}προπορεύσῃ γὰρ ἐνώπιον Κυρίου ἑτοιμάσαι ὁδοὺς αὐτοῦ,\end{quote}

\begin{quote} \vs{77}τοῦ δοῦναι γνῶσιν σωτηρίας τῷ λαῷ αὐτοῦ\end{quote} 

\begin{quote}ἐν ἀφέσει ἁμαρτιῶν αὐτῶν,\end{quote}

\begin{quote} \vs{78}διὰ σπλάγχνα ἐλέους Θεοῦ ἡμῶν,\end{quote} 

\begin{quote}ἐν οἷς ἐπισκέψεται ἡμᾶς ἀνατολὴ ἐξ ὕψους,\end{quote}

\begin{quote} \vs{79}ἐπιφᾶναι τοῖς ἐν σκότει καὶ σκιᾷ θανάτου καθημένοις,\end{quote} 

\begin{quote}τοῦ κατευθῦναι τοὺς πόδας ἡμῶν εἰς ὁδὸν εἰρήνης.\end{quote}
\end{poetryblock}

\vs{80}Τὸ δὲ παιδίον ηὔξανεν καὶ ἐκραταιοῦτο πνεύματι, καὶ ἦν ἐν ταῖς ἐρήμοις ἕως ἡμέρας ἀναδείξεως αὐτοῦ πρὸς τὸν Ἰσραήλ.

\ch{2}
Ἐγένετο δὲ ἐν ταῖς ἡμέραις ἐκείναις ἐξῆλθεν δόγμα παρὰ Καίσαρος Αὐγούστου ἀπογράφεσθαι πᾶσαν τὴν οἰκουμένην.
\vs{2}αὕτη ἀπογραφὴ πρώτη ἐγένετο ἡγεμονεύοντος τῆς Συρίας Κυρηνίου.
\vs{3}καὶ ἐπορεύοντο πάντες ἀπογράφεσθαι, ἕκαστος εἰς τὴν ἑαυτοῦ πόλιν.
\vs{4}Ἀνέβη δὲ καὶ Ἰωσὴφ ἀπὸ τῆς Γαλιλαίας ἐκ πόλεως Ναζαρὲθ εἰς τὴν Ἰουδαίαν εἰς πόλιν Δαυὶδ ἥτις καλεῖται Βηθλεέμ, διὰ τὸ εἶναι αὐτὸν ἐξ οἴκου καὶ πατριᾶς Δαυίδ,
\vs{5}ἀπογράψασθαι σὺν Μαριὰμ τῇ ἐμνηστευμένῃ αὐτῷ, οὔσῃ ἐγκύῳ.
\vs{6}Ἐγένετο δὲ ἐν τῷ εἶναι αὐτοὺς ἐκεῖ ἐπλήσθησαν αἱ ἡμέραι τοῦ τεκεῖν αὐτήν,
\vs{7}καὶ ἔτεκεν τὸν υἱὸν αὐτῆς τὸν πρωτότοκον, καὶ ἐσπαργάνωσεν αὐτὸν καὶ ἀνέκλινεν αὐτὸν ἐν φάτνῃ, διότι οὐκ ἦν αὐτοῖς τόπος ἐν τῷ καταλύματι.

\vs{8}Καὶ ποιμένες ἦσαν ἐν τῇ χώρᾳ τῇ αὐτῇ ἀγραυλοῦντες καὶ φυλάσσοντες φυλακὰς τῆς νυκτὸς ἐπὶ τὴν ποίμνην αὐτῶν.
\vs{9}καὶ ἄγγελος Κυρίου ἐπέστη αὐτοῖς καὶ δόξα Κυρίου περιέλαμψεν αὐτούς, καὶ ἐφοβήθησαν φόβον μέγαν.
\vs{10}καὶ εἶπεν αὐτοῖς ὁ ἄγγελος· Μὴ φοβεῖσθε, ἰδοὺ γὰρ εὐαγγελίζομαι ὑμῖν χαρὰν μεγάλην ἥτις ἔσται παντὶ τῷ λαῷ,
\vs{11}ὅτι ἐτέχθη ὑμῖν σήμερον Σωτήρ ὅς ἐστιν Χριστὸς Κύριος ἐν πόλει Δαυίδ.
\vs{12}καὶ τοῦτο ὑμῖν τὸ σημεῖον, εὑρήσετε βρέφος ἐσπαργανωμένον καὶ κείμενον ἐν φάτνῃ.
\vs{13}Καὶ ἐξαίφνης ἐγένετο σὺν τῷ ἀγγέλῳ πλῆθος στρατιᾶς οὐρανίου αἰνούντων τὸν Θεὸν καὶ λεγόντων·
\begin{poetryblock}

\begin{quote} \vs{14}Δόξα ἐν ὑψίστοις Θεῷ\end{quote} 

\begin{quote}καὶ ἐπὶ γῆς εἰρήνη\end{quote} 

\begin{quote}ἐν ἀνθρώποις εὐδοκίας.\end{quote}
\end{poetryblock}

\vs{15}Καὶ ἐγένετο ὡς ἀπῆλθον ἀπ᾽ αὐτῶν εἰς τὸν οὐρανὸν οἱ ἄγγελοι, οἱ ποιμένες ἐλάλουν πρὸς ἀλλήλους· Διέλθωμεν δὴ ἕως Βηθλεὲμ καὶ ἴδωμεν τὸ ῥῆμα τοῦτο τὸ γεγονὸς ὃ ὁ Κύριος ἐγνώρισεν ἡμῖν.
\vs{16}Καὶ ἦλθαν σπεύσαντες καὶ ἀνεῦραν τήν τε Μαριὰμ καὶ τὸν Ἰωσὴφ καὶ τὸ βρέφος κείμενον ἐν τῇ φάτνῃ·
\vs{17}ἰδόντες δὲ ἐγνώρισαν περὶ τοῦ ῥήματος τοῦ λαληθέντος αὐτοῖς περὶ τοῦ παιδίου τούτου.
\vs{18}καὶ πάντες οἱ ἀκούσαντες ἐθαύμασαν περὶ τῶν λαληθέντων ὑπὸ τῶν ποιμένων πρὸς αὐτούς·
\vs{19}ἡ δὲ Μαρία πάντα συνετήρει τὰ ῥήματα ταῦτα συμβάλλουσα ἐν τῇ καρδίᾳ αὐτῆς.
\vs{20}Καὶ ὑπέστρεψαν οἱ ποιμένες δοξάζοντες καὶ αἰνοῦντες τὸν Θεὸν ἐπὶ πᾶσιν οἷς ἤκουσαν καὶ εἶδον καθὼς ἐλαλήθη πρὸς αὐτούς.

\vs{21}Καὶ ὅτε ἐπλήσθησαν ἡμέραι ὀκτὼ τοῦ περιτεμεῖν αὐτόν καὶ ἐκλήθη τὸ ὄνομα αὐτοῦ Ἰησοῦς, τὸ κληθὲν ὑπὸ τοῦ ἀγγέλου πρὸ τοῦ συλλημφθῆναι αὐτὸν ἐν τῇ κοιλίᾳ.

\vs{22}Καὶ ὅτε ἐπλήσθησαν αἱ ἡμέραι τοῦ καθαρισμοῦ αὐτῶν κατὰ τὸν νόμον Μωϋσέως, ἀνήγαγον αὐτὸν εἰς Ἱεροσόλυμα παραστῆσαι τῷ Κυρίῳ,
\vs{23}καθὼς γέγραπται ἐν νόμῳ Κυρίου ὅτι Πᾶν ἄρσεν διανοῖγον μήτραν ἅγιον τῷ Κυρίῳ κληθήσεται,
\vs{24}καὶ τοῦ δοῦναι θυσίαν κατὰ τὸ εἰρημένον ἐν τῷ νόμῳ Κυρίου, Ζεῦγος τρυγόνων ἢ δύο νοσσοὺς περιστερῶν.

\vs{25}Καὶ ἰδοὺ ἄνθρωπος ἦν ἐν Ἰερουσαλὴμ ᾧ ὄνομα Συμεών καὶ ὁ ἄνθρωπος οὗτος δίκαιος καὶ εὐλαβής προσδεχόμενος παράκλησιν τοῦ Ἰσραήλ, καὶ Πνεῦμα ἦν Ἅγιον ἐπ᾽ αὐτόν·
\vs{26}καὶ ἦν αὐτῷ κεχρηματισμένον ὑπὸ τοῦ Πνεύματος τοῦ Ἁγίου μὴ ἰδεῖν θάνατον πρὶν ἢ ἂν ἴδῃ τὸν Χριστὸν Κυρίου.
\vs{27}καὶ ἦλθεν ἐν τῷ Πνεύματι εἰς τὸ ἱερόν· καὶ ἐν τῷ εἰσαγαγεῖν τοὺς γονεῖς τὸ παιδίον Ἰησοῦν τοῦ ποιῆσαι αὐτοὺς κατὰ τὸ εἰθισμένον τοῦ νόμου περὶ αὐτοῦ
\vs{28}καὶ αὐτὸς ἐδέξατο αὐτὸ εἰς τὰς ἀγκάλας καὶ εὐλόγησεν τὸν Θεὸν καὶ εἶπεν·
\begin{poetryblock}

\begin{quote} \vs{29}Νῦν ἀπολύεις τὸν δοῦλόν σου, Δέσποτα,\end{quote} 

\begin{quote}κατὰ τὸ ῥῆμά σου ἐν εἰρήνῃ·\end{quote}

\begin{quote} \vs{30}ὅτι εἶδον οἱ ὀφθαλμοί μου τὸ σωτήριόν σου,\end{quote}

\begin{quote} \vs{31}ὃ ἡτοίμασας κατὰ πρόσωπον πάντων τῶν λαῶν,\end{quote}

\begin{quote} \vs{32}φῶς εἰς ἀποκάλυψιν ἐθνῶν\end{quote} 

\begin{quote}καὶ δόξαν λαοῦ σου Ἰσραήλ.\end{quote}
\end{poetryblock}

\vs{33}Καὶ ἦν ὁ πατὴρ αὐτοῦ καὶ ἡ μήτηρ θαυμάζοντες ἐπὶ τοῖς λαλουμένοις περὶ αὐτοῦ.
\vs{34}καὶ εὐλόγησεν αὐτοὺς Συμεὼν καὶ εἶπεν πρὸς Μαριὰμ τὴν μητέρα αὐτοῦ· Ἰδοὺ οὗτος κεῖται εἰς πτῶσιν καὶ ἀνάστασιν πολλῶν ἐν τῷ Ἰσραὴλ καὶ εἰς σημεῖον ἀντιλεγόμενον—
\vs{35}καὶ σοῦ δὲ αὐτῆς τὴν ψυχὴν διελεύσεται ῥομφαία— ὅπως ἂν ἀποκαλυφθῶσιν ἐκ πολλῶν καρδιῶν διαλογισμοί.
\vs{36}Καὶ ἦν Ἅννα προφῆτις, θυγάτηρ Φανουήλ, ἐκ φυλῆς Ἀσήρ· αὕτη προβεβηκυῖα ἐν ἡμέραις πολλαῖς, ζήσασα μετὰ ἀνδρὸς ἔτη ἑπτὰ ἀπὸ τῆς παρθενίας αὐτῆς
\vs{37}καὶ αὐτὴ χήρα ἕως ἐτῶν ὀγδοήκοντα τεσσάρων, ἣ οὐκ ἀφίστατο τοῦ ἱεροῦ νηστείαις καὶ δεήσεσιν λατρεύουσα νύκτα καὶ ἡμέραν.
\vs{38}καὶ αὐτῇ τῇ ὥρᾳ ἐπιστᾶσα ἀνθωμολογεῖτο τῷ Θεῷ καὶ ἐλάλει περὶ αὐτοῦ πᾶσιν τοῖς προσδεχομένοις λύτρωσιν Ἰερουσαλήμ.

\vs{39}Καὶ ὡς ἐτέλεσαν πάντα τὰ κατὰ τὸν νόμον Κυρίου, ἐπέστρεψαν εἰς τὴν Γαλιλαίαν εἰς πόλιν ἑαυτῶν Ναζαρέθ.
\vs{40}Τὸ δὲ παιδίον ηὔξανεν καὶ ἐκραταιοῦτο πληρούμενον σοφίᾳ, καὶ χάρις Θεοῦ ἦν ἐπ᾽ αὐτό.

\vs{41}Καὶ ἐπορεύοντο οἱ γονεῖς αὐτοῦ κατ᾽ ἔτος εἰς Ἰερουσαλὴμ τῇ ἑορτῇ τοῦ πάσχα.
\vs{42}Καὶ ὅτε ἐγένετο ἐτῶν δώδεκα, ἀναβαινόντων αὐτῶν κατὰ τὸ ἔθος τῆς ἑορτῆς
\vs{43}καὶ τελειωσάντων τὰς ἡμέρας, ἐν τῷ ὑποστρέφειν αὐτοὺς ὑπέμεινεν Ἰησοῦς ὁ παῖς ἐν Ἰερουσαλήμ, καὶ οὐκ ἔγνωσαν οἱ γονεῖς αὐτοῦ.
\vs{44}νομίσαντες δὲ αὐτὸν εἶναι ἐν τῇ συνοδίᾳ ἦλθον ἡμέρας ὁδὸν καὶ ἀνεζήτουν αὐτὸν ἐν τοῖς συγγενεῦσιν καὶ τοῖς γνωστοῖς,
\vs{45}καὶ μὴ εὑρόντες ὑπέστρεψαν εἰς Ἰερουσαλὴμ ἀναζητοῦντες αὐτόν.
\vs{46}Καὶ ἐγένετο μετὰ ἡμέρας τρεῖς εὗρον αὐτὸν ἐν τῷ ἱερῷ καθεζόμενον ἐν μέσῳ τῶν διδασκάλων καὶ ἀκούοντα αὐτῶν καὶ ἐπερωτῶντα αὐτούς·
\vs{47}ἐξίσταντο δὲ πάντες οἱ ἀκούοντες αὐτοῦ ἐπὶ τῇ συνέσει καὶ ταῖς ἀποκρίσεσιν αὐτοῦ.
\vs{48}Καὶ ἰδόντες αὐτὸν ἐξεπλάγησαν, καὶ εἶπεν πρὸς αὐτὸν ἡ μήτηρ αὐτοῦ· Τέκνον, τί ἐποίησας ἡμῖν οὕτως; ἰδοὺ ὁ πατήρ σου κἀγὼ ὀδυνώμενοι ἐζητοῦμέν σε.
\vs{49}Καὶ εἶπεν πρὸς αὐτούς· Τί ὅτι ἐζητεῖτέ με; οὐκ ᾔδειτε ὅτι ἐν τοῖς τοῦ Πατρός μου δεῖ εἶναί με;
\vs{50}καὶ αὐτοὶ οὐ συνῆκαν τὸ ῥῆμα ὃ ἐλάλησεν αὐτοῖς.
\vs{51}Καὶ κατέβη μετ᾽ αὐτῶν καὶ ἦλθεν εἰς Ναζαρὲθ καὶ ἦν ὑποτασσόμενος αὐτοῖς. καὶ ἡ μήτηρ αὐτοῦ διετήρει πάντα τὰ ῥήματα ἐν τῇ καρδίᾳ αὐτῆς.
\vs{52}Καὶ Ἰησοῦς προέκοπτεν ἐν τῇ σοφίᾳ καὶ ἡλικίᾳ καὶ χάριτι παρὰ Θεῷ καὶ ἀνθρώποις.

\ch{3}
Ἐν ἔτει δὲ πεντεκαιδεκάτῳ τῆς ἡγεμονίας Τιβερίου Καίσαρος, ἡγεμονεύοντος Ποντίου Πιλάτου τῆς Ἰουδαίας, καὶ τετρααρχοῦντος τῆς Γαλιλαίας Ἡρῴδου, Φιλίππου δὲ τοῦ ἀδελφοῦ αὐτοῦ τετρααρχοῦντος τῆς Ἰτουραίας καὶ Τραχωνίτιδος χώρας, καὶ Λυσανίου τῆς Ἀβιληνῆς τετρααρχοῦντος,
\vs{2}ἐπὶ ἀρχιερέως Ἅννα καὶ Καϊάφα, ἐγένετο ῥῆμα Θεοῦ ἐπὶ Ἰωάννην τὸν Ζαχαρίου υἱὸν ἐν τῇ ἐρήμῳ.
\vs{3}Καὶ ἦλθεν εἰς πᾶσαν τὴν περίχωρον τοῦ Ἰορδάνου κηρύσσων βάπτισμα μετανοίας εἰς ἄφεσιν ἁμαρτιῶν,
\vs{4}ὡς γέγραπται ἐν βίβλῳ λόγων Ἠσαΐου τοῦ προφήτου· 
\begin{poetryblock}

\begin{quote}Φωνὴ βοῶντος ἐν τῇ ἐρήμῳ·\end{quote} 

\begin{quote}Ἑτοιμάσατε τὴν ὁδὸν Κυρίου,\end{quote} 

\begin{quote}εὐθείας ποιεῖτε τὰς τρίβους αὐτοῦ·\end{quote}

\begin{quote} \vs{5}πᾶσα φάραγξ πληρωθήσεται\end{quote} 

\begin{quote}καὶ πᾶν ὄρος καὶ βουνὸς ταπεινωθήσεται,\end{quote} 

\begin{quote}καὶ ἔσται τὰ σκολιὰ εἰς εὐθείαν\end{quote} 

\begin{quote}καὶ αἱ τραχεῖαι εἰς ὁδοὺς λείας·\end{quote}

\begin{quote} \vs{6}καὶ ὄψεται πᾶσα σὰρξ τὸ σωτήριον τοῦ Θεοῦ.\end{quote}
\end{poetryblock}

\vs{7}Ἔλεγεν οὖν τοῖς ἐκπορευομένοις ὄχλοις βαπτισθῆναι ὑπ᾽ αὐτοῦ· Γεννήματα ἐχιδνῶν, τίς ὑπέδειξεν ὑμῖν φυγεῖν ἀπὸ τῆς μελλούσης ὀργῆς;
\vs{8}ποιήσατε οὖν καρποὺς ἀξίους τῆς μετανοίας καὶ μὴ ἄρξησθε λέγειν ἐν ἑαυτοῖς· Πατέρα ἔχομεν τὸν Ἀβραάμ. λέγω γὰρ ὑμῖν ὅτι δύναται ὁ Θεὸς ἐκ τῶν λίθων τούτων ἐγεῖραι τέκνα τῷ Ἀβραάμ.
\vs{9}ἤδη δὲ καὶ ἡ ἀξίνη πρὸς τὴν ῥίζαν τῶν δένδρων κεῖται· πᾶν οὖν δένδρον μὴ ποιοῦν καρπὸν καλὸν ἐκκόπτεται καὶ εἰς πῦρ βάλλεται.

\vs{10}Καὶ ἐπηρώτων αὐτὸν οἱ ὄχλοι λέγοντες· Τί οὖν ποιήσωμεν;
\vs{11}Ἀποκριθεὶς δὲ ἔλεγεν αὐτοῖς· Ὁ ἔχων δύο χιτῶνας μεταδότω τῷ μὴ ἔχοντι, καὶ ὁ ἔχων βρώματα ὁμοίως ποιείτω.
\vs{12}Ἦλθον δὲ καὶ τελῶναι βαπτισθῆναι καὶ εἶπαν πρὸς αὐτόν· Διδάσκαλε, τί ποιήσωμεν;
\vs{13}Ὁ δὲ εἶπεν πρὸς αὐτούς· Μηδὲν πλέον παρὰ τὸ διατεταγμένον ὑμῖν πράσσετε.
\vs{14}Ἐπηρώτων δὲ αὐτὸν καὶ στρατευόμενοι λέγοντες· Τί ποιήσωμεν καὶ ἡμεῖς; Καὶ εἶπεν αὐτοῖς· Μηδένα διασείσητε μηδὲ συκοφαντήσητε καὶ ἀρκεῖσθε τοῖς ὀψωνίοις ὑμῶν.

\vs{15}Προσδοκῶντος δὲ τοῦ λαοῦ καὶ διαλογιζομένων πάντων ἐν ταῖς καρδίαις αὐτῶν περὶ τοῦ Ἰωάννου, μήποτε αὐτὸς εἴη ὁ Χριστός,
\vs{16}ἀπεκρίνατο λέγων πᾶσιν ὁ Ἰωάννης· Ἐγὼ μὲν ὕδατι βαπτίζω ὑμᾶς· ἔρχεται δὲ ὁ ἰσχυρότερός μου, οὗ οὐκ εἰμὶ ἱκανὸς λῦσαι τὸν ἱμάντα τῶν ὑποδημάτων αὐτοῦ· αὐτὸς ὑμᾶς βαπτίσει ἐν Πνεύματι Ἁγίῳ καὶ πυρί·
\vs{17}οὗ τὸ πτύον ἐν τῇ χειρὶ αὐτοῦ διακαθᾶραι τὴν ἅλωνα αὐτοῦ καὶ συναγαγεῖν τὸν σῖτον εἰς τὴν ἀποθήκην αὐτοῦ, τὸ δὲ ἄχυρον κατακαύσει πυρὶ ἀσβέστῳ.

\vs{18}Πολλὰ μὲν οὖν καὶ ἕτερα παρακαλῶν εὐηγγελίζετο τὸν λαόν.
\vs{19}ὁ δὲ Ἡρῴδης ὁ τετραάρχης, ἐλεγχόμενος ὑπ᾽ αὐτοῦ περὶ Ἡρῳδιάδος τῆς γυναικὸς τοῦ ἀδελφοῦ αὐτοῦ καὶ περὶ πάντων ὧν ἐποίησεν πονηρῶν ὁ Ἡρῴδης,
\vs{20}προσέθηκεν καὶ τοῦτο ἐπὶ πᾶσιν καὶ κατέκλεισεν τὸν Ἰωάννην ἐν φυλακῇ.

\vs{21}Ἐγένετο δὲ ἐν τῷ βαπτισθῆναι ἅπαντα τὸν λαὸν καὶ Ἰησοῦ βαπτισθέντος καὶ προσευχομένου ἀνεῳχθῆναι τὸν οὐρανὸν
\vs{22}καὶ καταβῆναι τὸ Πνεῦμα τὸ Ἅγιον σωματικῷ εἴδει ὡς περιστερὰν ἐπ᾽ αὐτόν, καὶ φωνὴν ἐξ οὐρανοῦ γενέσθαι· Σὺ εἶ ὁ Υἱός μου ὁ ἀγαπητός, ἐν σοὶ εὐδόκησα.

\vs{23}Καὶ αὐτὸς ἦν Ἰησοῦς ἀρχόμενος ὡσεὶ ἐτῶν τριάκοντα, Ὢν υἱός, ὡς ἐνομίζετο, Ἰωσὴφ τοῦ Ἠλὶ
\vs{24}τοῦ Μαθθὰτ τοῦ Λευὶ τοῦ Μελχὶ τοῦ Ἰανναὶ τοῦ Ἰωσὴφ
\vs{25}τοῦ Ματταθίου τοῦ Ἀμὼς τοῦ Ναοὺμ τοῦ Ἑσλὶ τοῦ Ναγγαὶ
\vs{26}τοῦ Μαὰθ τοῦ Ματταθίου τοῦ Σεμεῒν τοῦ Ἰωσὴχ τοῦ Ἰωδὰ
\vs{27}τοῦ Ἰωανὰν τοῦ Ῥησὰ τοῦ Ζοροβάβελ τοῦ Σαλαθιὴλ τοῦ Νηρὶ
\vs{28}τοῦ Μελχὶ τοῦ Ἀδδὶ τοῦ Κωσὰμ τοῦ Ἐλμαδὰμ τοῦ Ἢρ
\vs{29}τοῦ Ἰησοῦ τοῦ Ἐλιέζερ τοῦ Ἰωρὶμ τοῦ Μαθθὰτ τοῦ Λευὶ
\vs{30}τοῦ Συμεὼν τοῦ Ἰούδα τοῦ Ἰωσὴφ τοῦ Ἰωνὰμ τοῦ Ἐλιακὶμ
\vs{31}τοῦ Μελεὰ τοῦ Μεννὰ τοῦ Ματταθὰ τοῦ Ναθὰμ τοῦ Δαυὶδ
\vs{32}τοῦ Ἰεσσαὶ τοῦ Ἰωβὴδ τοῦ Βοὸς τοῦ Σαλὰ τοῦ Ναασσὼν
\vs{33}τοῦ Ἀμιναδὰβ τοῦ Ἀδμὶν τοῦ Ἀρνὶ τοῦ Ἑσρὼμ τοῦ Φαρὲς τοῦ Ἰούδα
\vs{34}τοῦ Ἰακὼβ τοῦ Ἰσαὰκ τοῦ Ἀβραὰμ τοῦ Θάρα τοῦ Ναχὼρ
\vs{35}τοῦ Σεροὺχ τοῦ Ῥαγαῦ τοῦ Φάλεκ τοῦ Ἔβερ τοῦ Σαλὰ
\vs{36}τοῦ Καϊνὰμ τοῦ Ἀρφαξὰδ τοῦ Σὴμ τοῦ Νῶε τοῦ Λάμεχ
\vs{37}τοῦ Μαθουσαλὰ τοῦ Ἑνὼχ τοῦ Ἰάρετ τοῦ Μαλελεὴλ τοῦ Καϊνὰμ
\vs{38}τοῦ Ἐνὼς τοῦ Σὴθ τοῦ Ἀδὰμ τοῦ Θεοῦ.

\ch{4}
Ἰησοῦς δὲ πλήρης Πνεύματος Ἁγίου ὑπέστρεψεν ἀπὸ τοῦ Ἰορδάνου καὶ ἤγετο ἐν τῷ Πνεύματι ἐν τῇ ἐρήμῳ
\vs{2}ἡμέρας τεσσεράκοντα πειραζόμενος ὑπὸ τοῦ διαβόλου. Καὶ οὐκ ἔφαγεν οὐδὲν ἐν ταῖς ἡμέραις ἐκείναις καὶ συντελεσθεισῶν αὐτῶν ἐπείνασεν.
\vs{3}Εἶπεν δὲ αὐτῷ ὁ διάβολος· Εἰ Υἱὸς εἶ τοῦ Θεοῦ, εἰπὲ τῷ λίθῳ τούτῳ ἵνα γένηται ἄρτος.
\vs{4}Καὶ ἀπεκρίθη πρὸς αὐτὸν ὁ Ἰησοῦς· Γέγραπται ὅτι Οὐκ ἐπ᾽ ἄρτῳ μόνῳ ζήσεται ὁ ἄνθρωπος.
\vs{5}Καὶ ἀναγαγὼν αὐτὸν ἔδειξεν αὐτῷ πάσας τὰς βασιλείας τῆς οἰκουμένης ἐν στιγμῇ χρόνου
\vs{6}καὶ εἶπεν αὐτῷ ὁ διάβολος· Σοὶ δώσω τὴν ἐξουσίαν ταύτην ἅπασαν καὶ τὴν δόξαν αὐτῶν, ὅτι ἐμοὶ παραδέδοται καὶ ᾧ ἐὰν θέλω δίδωμι αὐτήν·
\vs{7}σὺ οὖν ἐὰν προσκυνήσῃς ἐνώπιον ἐμοῦ, ἔσται σοῦ πᾶσα.
\vs{8}Καὶ ἀποκριθεὶς ὁ Ἰησοῦς εἶπεν αὐτῷ· Γέγραπται· Κύριον τὸν Θεόν σου προσκυνήσεις καὶ αὐτῷ μόνῳ λατρεύσεις.

\vs{9}Ἤγαγεν δὲ αὐτὸν εἰς Ἰερουσαλὴμ καὶ ἔστησεν ἐπὶ τὸ πτερύγιον τοῦ ἱεροῦ καὶ εἶπεν αὐτῷ· Εἰ Υἱὸς εἶ τοῦ Θεοῦ, βάλε σεαυτὸν ἐντεῦθεν κάτω·
\vs{10}γέγραπται γὰρ ὅτι 
\begin{poetryblock}

\begin{quote}Τοῖς ἀγγέλοις αὐτοῦ ἐντελεῖται περὶ σοῦ\end{quote} 

\begin{quote}τοῦ διαφυλάξαι σε\end{quote}
\end{poetryblock}

\vs{11}καὶ ὅτι 
\begin{poetryblock}

\begin{quote}ἐπὶ χειρῶν ἀροῦσίν σε,\end{quote} 

\begin{quote}μήποτε προσκόψῃς πρὸς λίθον τὸν πόδα σου.\end{quote}
\end{poetryblock}

\vs{12}Καὶ ἀποκριθεὶς εἶπεν αὐτῷ ὁ Ἰησοῦς ὅτι Εἴρηται· Οὐκ ἐκπειράσεις Κύριον τὸν Θεόν σου.

\vs{13}Καὶ συντελέσας πάντα πειρασμὸν ὁ διάβολος ἀπέστη ἀπ᾽ αὐτοῦ ἄχρι καιροῦ.

\vs{14}Καὶ ὑπέστρεψεν ὁ Ἰησοῦς ἐν τῇ δυνάμει τοῦ Πνεύματος εἰς τὴν Γαλιλαίαν. καὶ φήμη ἐξῆλθεν καθ᾽ ὅλης τῆς περιχώρου περὶ αὐτοῦ.
\vs{15}καὶ αὐτὸς ἐδίδασκεν ἐν ταῖς συναγωγαῖς αὐτῶν δοξαζόμενος ὑπὸ πάντων.

\vs{16}Καὶ ἦλθεν εἰς Ναζαρά, οὗ ἦν τεθραμμένος, καὶ εἰσῆλθεν κατὰ τὸ εἰωθὸς αὐτῷ ἐν τῇ ἡμέρᾳ τῶν σαββάτων εἰς τὴν συναγωγήν καὶ ἀνέστη ἀναγνῶναι.
\vs{17}καὶ ἐπεδόθη αὐτῷ βιβλίον τοῦ προφήτου Ἠσαΐου καὶ ἀναπτύξας τὸ βιβλίον εὗρεν τὸν τόπον οὗ ἦν γεγραμμένον·
\begin{poetryblock}

\begin{quote} \vs{18}Πνεῦμα Κυρίου ἐπ᾽ ἐμέ\end{quote} 

\begin{quote}Οὗ εἵνεκεν ἔχρισέν με\end{quote} 

\begin{quote}Εὐαγγελίσασθαι πτωχοῖς,\end{quote} 

\begin{quote}Ἀπέσταλκέν με,\end{quote} 

\begin{quote}κηρῦξαι αἰχμαλώτοις ἄφεσιν\end{quote} 

\begin{quote}Καὶ τυφλοῖς ἀνάβλεψιν,\end{quote} 

\begin{quote}Ἀποστεῖλαι τεθραυσμένους ἐν ἀφέσει,\end{quote}

\begin{quote} \vs{19}Κηρῦξαι ἐνιαυτὸν Κυρίου δεκτόν.\end{quote}
\end{poetryblock}

\vs{20}Καὶ πτύξας τὸ βιβλίον ἀποδοὺς τῷ ὑπηρέτῃ ἐκάθισεν· καὶ πάντων οἱ ὀφθαλμοὶ ἐν τῇ συναγωγῇ ἦσαν ἀτενίζοντες αὐτῷ.
\vs{21}ἤρξατο δὲ λέγειν πρὸς αὐτοὺς ὅτι Σήμερον πεπλήρωται ἡ γραφὴ αὕτη ἐν τοῖς ὠσὶν ὑμῶν.
\vs{22}Καὶ πάντες ἐμαρτύρουν αὐτῷ καὶ ἐθαύμαζον ἐπὶ τοῖς λόγοις τῆς χάριτος τοῖς ἐκπορευομένοις ἐκ τοῦ στόματος αὐτοῦ καὶ ἔλεγον· Οὐχὶ υἱός ἐστιν Ἰωσὴφ οὗτος;
\vs{23}Καὶ εἶπεν πρὸς αὐτούς· Πάντως ἐρεῖτέ μοι τὴν παραβολὴν ταύτην· Ἰατρέ, θεράπευσον σεαυτόν· ὅσα ἠκούσαμεν γενόμενα εἰς τὴν Καφαρναοὺμ ποίησον καὶ ὧδε ἐν τῇ πατρίδι σου.
\vs{24}Εἶπεν δέ· Ἀμὴν λέγω ὑμῖν ὅτι οὐδεὶς προφήτης δεκτός ἐστιν ἐν τῇ πατρίδι αὐτοῦ.
\vs{25}ἐπ᾽ ἀληθείας δὲ λέγω ὑμῖν, πολλαὶ χῆραι ἦσαν ἐν ταῖς ἡμέραις Ἠλίου ἐν τῷ Ἰσραήλ, ὅτε ἐκλείσθη ὁ οὐρανὸς ἐπὶ ἔτη τρία καὶ μῆνας ἕξ, ὡς ἐγένετο λιμὸς μέγας ἐπὶ πᾶσαν τὴν γῆν,
\vs{26}καὶ πρὸς οὐδεμίαν αὐτῶν ἐπέμφθη Ἠλίας εἰ μὴ εἰς Σάρεπτα τῆς Σιδωνίας πρὸς γυναῖκα χήραν.
\vs{27}καὶ πολλοὶ λεπροὶ ἦσαν ἐν τῷ Ἰσραὴλ ἐπὶ Ἐλισαίου τοῦ προφήτου, καὶ οὐδεὶς αὐτῶν ἐκαθαρίσθη εἰ μὴ Ναιμὰν ὁ Σύρος.
\vs{28}Καὶ ἐπλήσθησαν πάντες θυμοῦ ἐν τῇ συναγωγῇ ἀκούοντες ταῦτα
\vs{29}καὶ ἀναστάντες ἐξέβαλον αὐτὸν ἔξω τῆς πόλεως καὶ ἤγαγον αὐτὸν ἕως ὀφρύος τοῦ ὄρους ἐφ᾽ οὗ ἡ πόλις ᾠκοδόμητο αὐτῶν ὥστε κατακρημνίσαι αὐτόν·
\vs{30}αὐτὸς δὲ διελθὼν διὰ μέσου αὐτῶν ἐπορεύετο.

\vs{31}Καὶ κατῆλθεν εἰς Καφαρναοὺμ πόλιν τῆς Γαλιλαίας. καὶ ἦν διδάσκων αὐτοὺς ἐν τοῖς σάββασιν·
\vs{32}καὶ ἐξεπλήσσοντο ἐπὶ τῇ διδαχῇ αὐτοῦ, ὅτι ἐν ἐξουσίᾳ ἦν ὁ λόγος αὐτοῦ.

\vs{33}Καὶ ἐν τῇ συναγωγῇ ἦν ἄνθρωπος ἔχων πνεῦμα δαιμονίου ἀκαθάρτου καὶ ἀνέκραξεν φωνῇ μεγάλῃ·
\vs{34}Ἔα, τί ἡμῖν καὶ σοί, Ἰησοῦ Ναζαρηνέ; ἦλθες ἀπολέσαι ἡμᾶς; οἶδά σε τίς εἶ, ὁ Ἅγιος τοῦ Θεοῦ.
\vs{35}Καὶ ἐπετίμησεν αὐτῷ ὁ Ἰησοῦς λέγων· Φιμώθητι καὶ ἔξελθε ἀπ᾽ αὐτοῦ. καὶ ῥίψαν αὐτὸν τὸ δαιμόνιον εἰς τὸ μέσον ἐξῆλθεν ἀπ᾽ αὐτοῦ μηδὲν βλάψαν αὐτόν.
\vs{36}Καὶ ἐγένετο θάμβος ἐπὶ πάντας καὶ συνελάλουν πρὸς ἀλλήλους λέγοντες· Τίς ὁ λόγος οὗτος ὅτι ἐν ἐξουσίᾳ καὶ δυνάμει ἐπιτάσσει τοῖς ἀκαθάρτοις πνεύμασιν καὶ ἐξέρχονται;
\vs{37}καὶ ἐξεπορεύετο ἦχος περὶ αὐτοῦ εἰς πάντα τόπον τῆς περιχώρου.

\vs{38}Ἀναστὰς δὲ ἀπὸ τῆς συναγωγῆς εἰσῆλθεν εἰς τὴν οἰκίαν Σίμωνος. πενθερὰ δὲ τοῦ Σίμωνος ἦν συνεχομένη πυρετῷ μεγάλῳ καὶ ἠρώτησαν αὐτὸν περὶ αὐτῆς.
\vs{39}καὶ ἐπιστὰς ἐπάνω αὐτῆς ἐπετίμησεν τῷ πυρετῷ καὶ ἀφῆκεν αὐτήν· παραχρῆμα δὲ ἀναστᾶσα διηκόνει αὐτοῖς.

\vs{40}Δύνοντος δὲ τοῦ ἡλίου ἅπαντες ὅσοι εἶχον ἀσθενοῦντας νόσοις ποικίλαις ἤγαγον αὐτοὺς πρὸς αὐτόν· ὁ δὲ ἑνὶ ἑκάστῳ αὐτῶν τὰς χεῖρας ἐπιτιθεὶς ἐθεράπευεν αὐτούς.
\vs{41}ἐξήρχετο δὲ καὶ δαιμόνια ἀπὸ πολλῶν κραυγάζοντα καὶ λέγοντα ὅτι Σὺ εἶ ὁ Υἱὸς τοῦ Θεοῦ. καὶ ἐπιτιμῶν οὐκ εἴα αὐτὰ λαλεῖν, ὅτι ᾔδεισαν τὸν Χριστὸν αὐτὸν εἶναι.

\vs{42}Γενομένης δὲ ἡμέρας ἐξελθὼν ἐπορεύθη εἰς ἔρημον τόπον· καὶ οἱ ὄχλοι ἐπεζήτουν αὐτόν καὶ ἦλθον ἕως αὐτοῦ καὶ κατεῖχον αὐτὸν τοῦ μὴ πορεύεσθαι ἀπ᾽ αὐτῶν.
\vs{43}ὁ δὲ εἶπεν πρὸς αὐτοὺς ὅτι Καὶ ταῖς ἑτέραις πόλεσιν εὐαγγελίσασθαί με δεῖ τὴν βασιλείαν τοῦ Θεοῦ, ὅτι ἐπὶ τοῦτο ἀπεστάλην.
\vs{44}Καὶ ἦν κηρύσσων εἰς τὰς συναγωγὰς τῆς Ἰουδαίας.

\ch{5}
Ἐγένετο δὲ ἐν τῷ τὸν ὄχλον ἐπικεῖσθαι αὐτῷ καὶ ἀκούειν τὸν λόγον τοῦ Θεοῦ καὶ αὐτὸς ἦν ἑστὼς παρὰ τὴν λίμνην Γεννησαρέτ
\vs{2}καὶ εἶδεν δύο πλοῖα ἑστῶτα παρὰ τὴν λίμνην· οἱ δὲ ἁλιεῖς ἀπ᾽ αὐτῶν ἀποβάντες ἔπλυνον τὰ δίκτυα.
\vs{3}ἐμβὰς δὲ εἰς ἓν τῶν πλοίων, ὃ ἦν Σίμωνος, ἠρώτησεν αὐτὸν ἀπὸ τῆς γῆς ἐπαναγαγεῖν ὀλίγον· καθίσας δὲ ἐκ τοῦ πλοίου ἐδίδασκεν τοὺς ὄχλους.

\vs{4}Ὡς δὲ ἐπαύσατο λαλῶν, εἶπεν πρὸς τὸν Σίμωνα· Ἐπανάγαγε εἰς τὸ βάθος καὶ χαλάσατε τὰ δίκτυα ὑμῶν εἰς ἄγραν.
\vs{5}Καὶ ἀποκριθεὶς Σίμων εἶπεν· Ἐπιστάτα, δι᾽ ὅλης νυκτὸς κοπιάσαντες οὐδὲν ἐλάβομεν· ἐπὶ δὲ τῷ ῥήματί σου χαλάσω τὰ δίκτυα.
\vs{6}καὶ τοῦτο ποιήσαντες συνέκλεισαν πλῆθος ἰχθύων πολύ, διερρήσσετο δὲ τὰ δίκτυα αὐτῶν.
\vs{7}καὶ κατένευσαν τοῖς μετόχοις ἐν τῷ ἑτέρῳ πλοίῳ τοῦ ἐλθόντας συλλαβέσθαι αὐτοῖς· καὶ ἦλθον καὶ ἔπλησαν ἀμφότερα τὰ πλοῖα ὥστε βυθίζεσθαι αὐτά.
\vs{8}Ἰδὼν δὲ Σίμων Πέτρος προσέπεσεν τοῖς γόνασιν Ἰησοῦ λέγων· Ἔξελθε ἀπ᾽ ἐμοῦ, ὅτι ἀνὴρ ἁμαρτωλός εἰμι, Κύριε.
\vs{9}θάμβος γὰρ περιέσχεν αὐτὸν καὶ πάντας τοὺς σὺν αὐτῷ ἐπὶ τῇ ἄγρᾳ τῶν ἰχθύων ὧν συνέλαβον,
\vs{10}ὁμοίως δὲ καὶ Ἰάκωβον καὶ Ἰωάννην υἱοὺς Ζεβεδαίου, οἳ ἦσαν κοινωνοὶ τῷ Σίμωνι. Καὶ εἶπεν πρὸς τὸν Σίμωνα ὁ Ἰησοῦς· Μὴ φοβοῦ· ἀπὸ τοῦ νῦν ἀνθρώπους ἔσῃ ζωγρῶν.
\vs{11}καὶ καταγαγόντες τὰ πλοῖα ἐπὶ τὴν γῆν ἀφέντες πάντα ἠκολούθησαν αὐτῷ.

\vs{12}Καὶ ἐγένετο ἐν τῷ εἶναι αὐτὸν ἐν μιᾷ τῶν πόλεων καὶ ἰδοὺ ἀνὴρ πλήρης λέπρας· ἰδὼν δὲ τὸν Ἰησοῦν, πεσὼν ἐπὶ πρόσωπον ἐδεήθη αὐτοῦ λέγων· Κύριε, ἐὰν θέλῃς δύνασαί με καθαρίσαι.
\vs{13}Καὶ ἐκτείνας τὴν χεῖρα ἥψατο αὐτοῦ λέγων· Θέλω, καθαρίσθητι· καὶ εὐθέως ἡ λέπρα ἀπῆλθεν ἀπ᾽ αὐτοῦ.
\vs{14}Καὶ αὐτὸς παρήγγειλεν αὐτῷ μηδενὶ εἰπεῖν, Ἀλλὰ ἀπελθὼν δεῖξον σεαυτὸν τῷ ἱερεῖ καὶ προσένεγκε περὶ τοῦ καθαρισμοῦ σου καθὼς προσέταξεν Μωϋσῆς, εἰς μαρτύριον αὐτοῖς.
\vs{15}Διήρχετο δὲ μᾶλλον ὁ λόγος περὶ αὐτοῦ, καὶ συνήρχοντο ὄχλοι πολλοὶ ἀκούειν καὶ θεραπεύεσθαι ἀπὸ τῶν ἀσθενειῶν αὐτῶν·
\vs{16}αὐτὸς δὲ ἦν ὑποχωρῶν ἐν ταῖς ἐρήμοις καὶ προσευχόμενος.

\vs{17}Καὶ ἐγένετο ἐν μιᾷ τῶν ἡμερῶν καὶ αὐτὸς ἦν διδάσκων, καὶ ἦσαν καθήμενοι Φαρισαῖοι καὶ νομοδιδάσκαλοι οἳ ἦσαν ἐληλυθότες ἐκ πάσης κώμης τῆς Γαλιλαίας καὶ Ἰουδαίας καὶ Ἰερουσαλήμ· καὶ δύναμις Κυρίου ἦν εἰς τὸ ἰᾶσθαι αὐτόν.
\vs{18}Καὶ ἰδοὺ ἄνδρες φέροντες ἐπὶ κλίνης ἄνθρωπον ὃς ἦν παραλελυμένος καὶ ἐζήτουν αὐτὸν εἰσενεγκεῖν καὶ θεῖναι αὐτὸν ἐνώπιον αὐτοῦ.
\vs{19}καὶ μὴ εὑρόντες ποίας εἰσενέγκωσιν αὐτὸν διὰ τὸν ὄχλον, ἀναβάντες ἐπὶ τὸ δῶμα διὰ τῶν κεράμων καθῆκαν αὐτὸν σὺν τῷ κλινιδίῳ εἰς τὸ μέσον ἔμπροσθεν τοῦ Ἰησοῦ.
\vs{20}Καὶ ἰδὼν τὴν πίστιν αὐτῶν εἶπεν· Ἄνθρωπε, ἀφέωνταί σοι αἱ ἁμαρτίαι σου.
\vs{21}Καὶ ἤρξαντο διαλογίζεσθαι οἱ γραμματεῖς καὶ οἱ Φαρισαῖοι λέγοντες· Τίς ἐστιν οὗτος ὃς λαλεῖ βλασφημίας; τίς δύναται ἁμαρτίας ἀφεῖναι εἰ μὴ μόνος ὁ Θεός;
\vs{22}Ἐπιγνοὺς δὲ ὁ Ἰησοῦς τοὺς διαλογισμοὺς αὐτῶν ἀποκριθεὶς εἶπεν πρὸς αὐτούς· Τί διαλογίζεσθε ἐν ταῖς καρδίαις ὑμῶν;
\vs{23}τί ἐστιν εὐκοπώτερον, εἰπεῖν· Ἀφέωνταί σοι αἱ ἁμαρτίαι σου, ἢ εἰπεῖν· Ἔγειρε καὶ περιπάτει;
\vs{24}ἵνα δὲ εἰδῆτε ὅτι ὁ Υἱὸς τοῦ ἀνθρώπου ἐξουσίαν ἔχει ἐπὶ τῆς γῆς ἀφιέναι ἁμαρτίας— εἶπεν τῷ παραλελυμένῳ· Σοὶ λέγω, ἔγειρε καὶ ἄρας τὸ κλινίδιόν σου πορεύου εἰς τὸν οἶκόν σου.
\vs{25}Καὶ παραχρῆμα ἀναστὰς ἐνώπιον αὐτῶν, ἄρας ἐφ᾽ ὃ κατέκειτο, ἀπῆλθεν εἰς τὸν οἶκον αὐτοῦ δοξάζων τὸν Θεόν.
\vs{26}καὶ ἔκστασις ἔλαβεν ἅπαντας καὶ ἐδόξαζον τὸν Θεόν καὶ ἐπλήσθησαν φόβου λέγοντες ὅτι Εἴδομεν παράδοξα σήμερον.

\vs{27}Καὶ μετὰ ταῦτα ἐξῆλθεν καὶ ἐθεάσατο τελώνην ὀνόματι Λευὶν καθήμενον ἐπὶ τὸ τελώνιον, καὶ εἶπεν αὐτῷ· Ἀκολούθει μοι.
\vs{28}καὶ καταλιπὼν πάντα ἀναστὰς ἠκολούθει αὐτῷ.
\vs{29}Καὶ ἐποίησεν δοχὴν μεγάλην Λευὶς αὐτῷ ἐν τῇ οἰκίᾳ αὐτοῦ, καὶ ἦν ὄχλος πολὺς τελωνῶν καὶ ἄλλων οἳ ἦσαν μετ᾽ αὐτῶν κατακείμενοι.
\vs{30}καὶ ἐγόγγυζον οἱ Φαρισαῖοι καὶ οἱ γραμματεῖς αὐτῶν πρὸς τοὺς μαθητὰς αὐτοῦ λέγοντες· Διὰ τί μετὰ τῶν τελωνῶν καὶ ἁμαρτωλῶν ἐσθίετε καὶ πίνετε;
\vs{31}Καὶ ἀποκριθεὶς ὁ Ἰησοῦς εἶπεν πρὸς αὐτούς· Οὐ χρείαν ἔχουσιν οἱ ὑγιαίνοντες ἰατροῦ ἀλλὰ οἱ κακῶς ἔχοντες·
\vs{32}οὐκ ἐλήλυθα καλέσαι δικαίους ἀλλὰ ἁμαρτωλοὺς εἰς μετάνοιαν.

\vs{33}Οἱ δὲ εἶπαν πρὸς αὐτόν· Οἱ μαθηταὶ Ἰωάννου νηστεύουσιν πυκνὰ καὶ δεήσεις ποιοῦνται ὁμοίως καὶ οἱ τῶν Φαρισαίων, οἱ δὲ σοὶ ἐσθίουσιν καὶ πίνουσιν.
\vs{34}Ὁ δὲ Ἰησοῦς εἶπεν πρὸς αὐτούς· Μὴ δύνασθε τοὺς υἱοὺς τοῦ νυμφῶνος ἐν ᾧ ὁ νυμφίος μετ᾽ αὐτῶν ἐστιν ποιῆσαι νηστεῦσαι;
\vs{35}ἐλεύσονται δὲ ἡμέραι, καὶ ὅταν ἀπαρθῇ ἀπ᾽ αὐτῶν ὁ νυμφίος, τότε νηστεύσουσιν ἐν ἐκείναις ταῖς ἡμέραις.

\vs{36}Ἔλεγεν δὲ καὶ παραβολὴν πρὸς αὐτοὺς ὅτι Οὐδεὶς ἐπίβλημα ἀπὸ ἱματίου καινοῦ σχίσας ἐπιβάλλει ἐπὶ ἱμάτιον παλαιόν· εἰ δὲ μή γε, καὶ τὸ καινὸν σχίσει καὶ τῷ παλαιῷ οὐ συμφωνήσει τὸ ἐπίβλημα τὸ ἀπὸ τοῦ καινοῦ.
\vs{37}Καὶ οὐδεὶς βάλλει οἶνον νέον εἰς ἀσκοὺς παλαιούς· εἰ δὲ μή γε, ῥήξει ὁ οἶνος ὁ νέος τοὺς ἀσκούς καὶ αὐτὸς ἐκχυθήσεται καὶ οἱ ἀσκοὶ ἀπολοῦνται·
\vs{38}ἀλλὰ οἶνον νέον εἰς ἀσκοὺς καινοὺς βλητέον.
\vs{39}καὶ οὐδεὶς πιὼν παλαιὸν θέλει νέον· λέγει γάρ· Ὁ παλαιὸς χρηστός ἐστιν.

\ch{6}
Ἐγένετο δὲ ἐν σαββάτῳ διαπορεύεσθαι αὐτὸν διὰ σπορίμων, καὶ ἔτιλλον οἱ μαθηταὶ αὐτοῦ καὶ ἤσθιον τοὺς στάχυας ψώχοντες ταῖς χερσίν.
\vs{2}τινὲς δὲ τῶν Φαρισαίων εἶπαν· Τί ποιεῖτε ὃ οὐκ ἔξεστιν τοῖς σάββασιν;
\vs{3}Καὶ ἀποκριθεὶς πρὸς αὐτοὺς εἶπεν ὁ Ἰησοῦς· Οὐδὲ τοῦτο ἀνέγνωτε ὃ ἐποίησεν Δαυὶδ ὅτε ἐπείνασεν αὐτὸς καὶ οἱ μετ᾽ αὐτοῦ ὄντες,
\vs{4}ὡς εἰσῆλθεν εἰς τὸν οἶκον τοῦ Θεοῦ καὶ τοὺς ἄρτους τῆς προθέσεως λαβὼν ἔφαγεν καὶ ἔδωκεν τοῖς μετ᾽ αὐτοῦ, οὓς οὐκ ἔξεστιν φαγεῖν εἰ μὴ μόνους τοὺς ἱερεῖς;
\vs{5}καὶ ἔλεγεν αὐτοῖς· Κύριός ἐστιν τοῦ σαββάτου ὁ Υἱὸς τοῦ ἀνθρώπου.
\vs{6}Ἐγένετο δὲ ἐν ἑτέρῳ σαββάτῳ εἰσελθεῖν αὐτὸν εἰς τὴν συναγωγὴν καὶ διδάσκειν. καὶ ἦν ἄνθρωπος ἐκεῖ καὶ ἡ χεὶρ αὐτοῦ ἡ δεξιὰ ἦν ξηρά.
\vs{7}παρετηροῦντο δὲ αὐτὸν οἱ γραμματεῖς καὶ οἱ Φαρισαῖοι εἰ ἐν τῷ σαββάτῳ θεραπεύει, ἵνα εὕρωσιν κατηγορεῖν αὐτοῦ.
\vs{8}Αὐτὸς δὲ ᾔδει τοὺς διαλογισμοὺς αὐτῶν, εἶπεν δὲ τῷ ἀνδρὶ τῷ ξηρὰν ἔχοντι τὴν χεῖρα· Ἔγειρε καὶ στῆθι εἰς τὸ μέσον· καὶ ἀναστὰς ἔστη.
\vs{9}εἶπεν δὲ ὁ Ἰησοῦς πρὸς αὐτούς· Ἐπερωτῶ ὑμᾶς εἰ ἔξεστιν τῷ σαββάτῳ ἀγαθοποιῆσαι ἢ κακοποιῆσαι, ψυχὴν σῶσαι ἢ ἀπολέσαι;
\vs{10}Καὶ περιβλεψάμενος πάντας αὐτοὺς εἶπεν αὐτῷ· Ἔκτεινον τὴν χεῖρά σου. ὁ δὲ ἐποίησεν καὶ ἀπεκατεστάθη ἡ χεὶρ αὐτοῦ.
\vs{11}αὐτοὶ δὲ ἐπλήσθησαν ἀνοίας καὶ διελάλουν πρὸς ἀλλήλους τί ἂν ποιήσαιεν τῷ Ἰησοῦ.

\vs{12}Ἐγένετο δὲ ἐν ταῖς ἡμέραις ταύταις ἐξελθεῖν αὐτὸν εἰς τὸ ὄρος προσεύξασθαι, καὶ ἦν διανυκτερεύων ἐν τῇ προσευχῇ τοῦ Θεοῦ.
\vs{13}καὶ ὅτε ἐγένετο ἡμέρα, προσεφώνησεν τοὺς μαθητὰς αὐτοῦ, καὶ ἐκλεξάμενος ἀπ᾽ αὐτῶν δώδεκα, οὓς καὶ ἀποστόλους ὠνόμασεν·
\vs{14}Σίμωνα ὃν καὶ ὠνόμασεν Πέτρον, καὶ Ἀνδρέαν τὸν ἀδελφὸν αὐτοῦ, καὶ Ἰάκωβον καὶ Ἰωάννην καὶ Φίλιππον καὶ Βαρθολομαῖον
\vs{15}καὶ Μαθθαῖον καὶ Θωμᾶν καὶ Ἰάκωβον Ἁλφαίου καὶ Σίμωνα τὸν καλούμενον Ζηλωτὴν
\vs{16}καὶ Ἰούδαν Ἰακώβου καὶ Ἰούδαν Ἰσκαριὼθ, ὃς ἐγένετο προδότης.

\vs{17}Καὶ καταβὰς μετ᾽ αὐτῶν ἔστη ἐπὶ τόπου πεδινοῦ, καὶ ὄχλος πολὺς μαθητῶν αὐτοῦ, καὶ πλῆθος πολὺ τοῦ λαοῦ ἀπὸ πάσης τῆς Ἰουδαίας καὶ Ἰερουσαλὴμ καὶ τῆς παραλίου Τύρου καὶ Σιδῶνος,
\vs{18}οἳ ἦλθον ἀκοῦσαι αὐτοῦ καὶ ἰαθῆναι ἀπὸ τῶν νόσων αὐτῶν· καὶ οἱ ἐνοχλούμενοι ἀπὸ πνευμάτων ἀκαθάρτων ἐθεραπεύοντο,
\vs{19}καὶ πᾶς ὁ ὄχλος ἐζήτουν ἅπτεσθαι αὐτοῦ, ὅτι δύναμις παρ᾽ αὐτοῦ ἐξήρχετο καὶ ἰᾶτο πάντας.

\vs{20}Καὶ αὐτὸς ἐπάρας τοὺς ὀφθαλμοὺς αὐτοῦ εἰς τοὺς μαθητὰς αὐτοῦ ἔλεγεν· 
\begin{poetryblock}

\begin{quote}Μακάριοι οἱ πτωχοί, Ὅτι ὑμετέρα ἐστὶν ἡ βασιλεία τοῦ Θεοῦ.\end{quote}

\begin{quote} \vs{21}Μακάριοι οἱ πεινῶντες νῦν,\end{quote} 

\begin{quote}Ὅτι χορτασθήσεσθε.\end{quote} 

\begin{quote}Μακάριοι οἱ κλαίοντες νῦν,\end{quote} 

\begin{quote}Ὅτι γελάσετε.\end{quote}

\begin{quote} \vs{22}Μακάριοί ἐστε ὅταν μισήσωσιν ὑμᾶς οἱ ἄνθρωποι καὶ ὅταν ἀφορίσωσιν ὑμᾶς καὶ ὀνειδίσωσιν καὶ ἐκβάλωσιν τὸ ὄνομα ὑμῶν ὡς πονηρὸν ἕνεκα τοῦ Υἱοῦ τοῦ ἀνθρώπου·\end{quote}
\end{poetryblock}

\vs{23}χάρητε ἐν ἐκείνῃ τῇ ἡμέρᾳ καὶ σκιρτήσατε, ἰδοὺ γὰρ ὁ μισθὸς ὑμῶν πολὺς ἐν τῷ οὐρανῷ· κατὰ τὰ αὐτὰ γὰρ ἐποίουν τοῖς προφήταις οἱ πατέρες αὐτῶν.
\begin{poetryblock}

\begin{quote} \vs{24}Πλὴν οὐαὶ ὑμῖν τοῖς πλουσίοις,\end{quote} 

\begin{quote}Ὅτι ἀπέχετε τὴν παράκλησιν ὑμῶν.\end{quote}

\begin{quote} \vs{25}Οὐαὶ ὑμῖν, οἱ ἐμπεπλησμένοι νῦν,\end{quote} 

\begin{quote}Ὅτι πεινάσετε.\end{quote} 

\begin{quote}Οὐαί, οἱ γελῶντες νῦν,\end{quote} 

\begin{quote}Ὅτι πενθήσετε καὶ κλαύσετε.\end{quote}

\begin{quote} \vs{26}Οὐαὶ ὅταν ὑμᾶς καλῶς εἴπωσιν πάντες οἱ ἄνθρωποι·\end{quote} 

\begin{quote}Κατὰ τὰ αὐτὰ γὰρ ἐποίουν τοῖς ψευδοπροφήταις οἱ πατέρες αὐτῶν.\end{quote}
\end{poetryblock}

\vs{27}Ἀλλὰ ὑμῖν λέγω τοῖς ἀκούουσιν· Ἀγαπᾶτε τοὺς ἐχθροὺς ὑμῶν, καλῶς ποιεῖτε τοῖς μισοῦσιν ὑμᾶς,
\vs{28}εὐλογεῖτε τοὺς καταρωμένους ὑμᾶς, προσεύχεσθε περὶ τῶν ἐπηρεαζόντων ὑμᾶς.
\vs{29}τῷ τύπτοντί σε ἐπὶ τὴν σιαγόνα πάρεχε καὶ τὴν ἄλλην, καὶ ἀπὸ τοῦ αἴροντός σου τὸ ἱμάτιον καὶ τὸν χιτῶνα μὴ κωλύσῃς.
\vs{30}παντὶ αἰτοῦντί σε δίδου, καὶ ἀπὸ τοῦ αἴροντος τὰ σὰ μὴ ἀπαίτει.

\vs{31}καὶ καθὼς θέλετε ἵνα ποιῶσιν ὑμῖν οἱ ἄνθρωποι ποιεῖτε αὐτοῖς ὁμοίως.
\vs{32}Καὶ εἰ ἀγαπᾶτε τοὺς ἀγαπῶντας ὑμᾶς, ποία ὑμῖν χάρις ἐστίν; καὶ γὰρ οἱ ἁμαρτωλοὶ τοὺς ἀγαπῶντας αὐτοὺς ἀγαπῶσιν.
\vs{33}καὶ γὰρ ἐὰν ἀγαθοποιῆτε τοὺς ἀγαθοποιοῦντας ὑμᾶς, ποία ὑμῖν χάρις ἐστίν; καὶ οἱ ἁμαρτωλοὶ τὸ αὐτὸ ποιοῦσιν.
\vs{34}καὶ ἐὰν δανίσητε παρ᾽ ὧν ἐλπίζετε λαβεῖν, ποία ὑμῖν χάρις ἐστίν; καὶ ἁμαρτωλοὶ ἁμαρτωλοῖς δανίζουσιν ἵνα ἀπολάβωσιν τὰ ἴσα.
\vs{35}Πλὴν ἀγαπᾶτε τοὺς ἐχθροὺς ὑμῶν καὶ ἀγαθοποιεῖτε καὶ δανίζετε μηδὲν ἀπελπίζοντες· καὶ ἔσται ὁ μισθὸς ὑμῶν πολύς, καὶ ἔσεσθε υἱοὶ Ὑψίστου, ὅτι αὐτὸς χρηστός ἐστιν ἐπὶ τοὺς ἀχαρίστους καὶ πονηρούς.

\vs{36}Γίνεσθε οἰκτίρμονες καθὼς καὶ ὁ Πατὴρ ὑμῶν οἰκτίρμων ἐστίν.
\vs{37}Καὶ μὴ κρίνετε, καὶ οὐ μὴ κριθῆτε· καὶ μὴ καταδικάζετε, καὶ οὐ μὴ καταδικασθῆτε. ἀπολύετε, καὶ ἀπολυθήσεσθε·
\vs{38}δίδοτε, καὶ δοθήσεται ὑμῖν· μέτρον καλὸν πεπιεσμένον σεσαλευμένον ὑπερεκχυννόμενον δώσουσιν εἰς τὸν κόλπον ὑμῶν· ᾧ γὰρ μέτρῳ μετρεῖτε ἀντιμετρηθήσεται ὑμῖν.

\vs{39}Εἶπεν δὲ καὶ παραβολὴν αὐτοῖς· Μήτι δύναται τυφλὸς τυφλὸν ὁδηγεῖν; οὐχὶ ἀμφότεροι εἰς βόθυνον ἐμπεσοῦνται;
\vs{40}οὐκ ἔστιν μαθητὴς ὑπὲρ τὸν διδάσκαλον· κατηρτισμένος δὲ πᾶς ἔσται ὡς ὁ διδάσκαλος αὐτοῦ.

\vs{41}Τί δὲ βλέπεις τὸ κάρφος τὸ ἐν τῷ ὀφθαλμῷ τοῦ ἀδελφοῦ σου, τὴν δὲ δοκὸν τὴν ἐν τῷ ἰδίῳ ὀφθαλμῷ οὐ κατανοεῖς;
\vs{42}πῶς δύνασαι λέγειν τῷ ἀδελφῷ σου· Ἀδελφέ, ἄφες ἐκβάλω τὸ κάρφος τὸ ἐν τῷ ὀφθαλμῷ σου, αὐτὸς τὴν ἐν τῷ ὀφθαλμῷ σοῦ δοκὸν οὐ βλέπων; ὑποκριτά, ἔκβαλε πρῶτον τὴν δοκὸν ἐκ τοῦ ὀφθαλμοῦ σοῦ, καὶ τότε διαβλέψεις τὸ κάρφος τὸ ἐν τῷ ὀφθαλμῷ τοῦ ἀδελφοῦ σου ἐκβαλεῖν.

\vs{43}Οὐ γάρ ἐστιν δένδρον καλὸν ποιοῦν καρπὸν σαπρόν, οὐδὲ πάλιν δένδρον σαπρὸν ποιοῦν καρπὸν καλόν.
\vs{44}ἕκαστον γὰρ δένδρον ἐκ τοῦ ἰδίου καρποῦ γινώσκεται· οὐ γὰρ ἐξ ἀκανθῶν συλλέγουσιν σῦκα οὐδὲ ἐκ βάτου σταφυλὴν τρυγῶσιν.
\vs{45}ὁ ἀγαθὸς ἄνθρωπος ἐκ τοῦ ἀγαθοῦ θησαυροῦ τῆς καρδίας προφέρει τὸ ἀγαθόν, καὶ ὁ πονηρὸς ἐκ τοῦ πονηροῦ προφέρει τὸ πονηρόν· ἐκ γὰρ περισσεύματος καρδίας λαλεῖ τὸ στόμα αὐτοῦ.

\vs{46}Τί δέ με καλεῖτε· Κύριε κύριε, καὶ οὐ ποιεῖτε ἃ λέγω;
\vs{47}Πᾶς ὁ ἐρχόμενος πρός με καὶ ἀκούων μου τῶν λόγων καὶ ποιῶν αὐτούς, ὑποδείξω ὑμῖν τίνι ἐστὶν ὅμοιος·
\vs{48}ὅμοιός ἐστιν ἀνθρώπῳ οἰκοδομοῦντι οἰκίαν ὃς ἔσκαψεν καὶ ἐβάθυνεν καὶ ἔθηκεν θεμέλιον ἐπὶ τὴν πέτραν· πλημμύρης δὲ γενομένης προσέρηξεν ὁ ποταμὸς τῇ οἰκίᾳ ἐκείνῃ, καὶ οὐκ ἴσχυσεν σαλεῦσαι αὐτὴν διὰ τὸ καλῶς οἰκοδομῆσθαι αὐτήν.
\vs{49}ὁ δὲ ἀκούσας καὶ μὴ ποιήσας ὅμοιός ἐστιν ἀνθρώπῳ οἰκοδομήσαντι οἰκίαν ἐπὶ τὴν γῆν χωρὶς θεμελίου, ᾗ προσέρηξεν ὁ ποταμός, καὶ εὐθὺς συνέπεσεν καὶ ἐγένετο τὸ ῥῆγμα τῆς οἰκίας ἐκείνης μέγα.

\ch{7}
Ἐπειδὴ ἐπλήρωσεν πάντα τὰ ῥήματα αὐτοῦ εἰς τὰς ἀκοὰς τοῦ λαοῦ, εἰσῆλθεν εἰς Καφαρναούμ.
\vs{2}Ἑκατοντάρχου δέ τινος δοῦλος κακῶς ἔχων ἤμελλεν τελευτᾶν, ὃς ἦν αὐτῷ ἔντιμος.
\vs{3}ἀκούσας δὲ περὶ τοῦ Ἰησοῦ ἀπέστειλεν πρὸς αὐτὸν πρεσβυτέρους τῶν Ἰουδαίων ἐρωτῶν αὐτὸν ὅπως ἐλθὼν διασώσῃ τὸν δοῦλον αὐτοῦ.
\vs{4}οἱ δὲ παραγενόμενοι πρὸς τὸν Ἰησοῦν παρεκάλουν αὐτὸν σπουδαίως λέγοντες ὅτι Ἄξιός ἐστιν ᾧ παρέξῃ τοῦτο·
\vs{5}ἀγαπᾷ γὰρ τὸ ἔθνος ἡμῶν καὶ τὴν συναγωγὴν αὐτὸς ᾠκοδόμησεν ἡμῖν.
\vs{6}Ὁ δὲ Ἰησοῦς ἐπορεύετο σὺν αὐτοῖς. ἤδη δὲ αὐτοῦ οὐ μακρὰν ἀπέχοντος ἀπὸ τῆς οἰκίας ἔπεμψεν φίλους ὁ ἑκατοντάρχης λέγων αὐτῷ· Κύριε, μὴ σκύλλου, οὐ γὰρ ἱκανός εἰμι ἵνα ὑπὸ τὴν στέγην μου εἰσέλθῃς·
\vs{7}διὸ οὐδὲ ἐμαυτὸν ἠξίωσα πρὸς σὲ ἐλθεῖν· ἀλλὰ εἰπὲ λόγῳ, καὶ ἰαθήτω ὁ παῖς μου.
\vs{8}καὶ γὰρ ἐγὼ ἄνθρωπός εἰμι ὑπὸ ἐξουσίαν τασσόμενος ἔχων ὑπ᾽ ἐμαυτὸν στρατιώτας, καὶ λέγω τούτῳ· Πορεύθητι, καὶ πορεύεται, καὶ ἄλλῳ· Ἔρχου, καὶ ἔρχεται, καὶ τῷ δούλῳ μου· Ποίησον τοῦτο, καὶ ποιεῖ.
\vs{9}Ἀκούσας δὲ ταῦτα ὁ Ἰησοῦς ἐθαύμασεν αὐτόν καὶ στραφεὶς τῷ ἀκολουθοῦντι αὐτῷ ὄχλῳ εἶπεν· Λέγω ὑμῖν, οὐδὲ ἐν τῷ Ἰσραὴλ τοσαύτην πίστιν εὗρον.
\vs{10}καὶ ὑποστρέψαντες εἰς τὸν οἶκον οἱ πεμφθέντες εὗρον τὸν δοῦλον ὑγιαίνοντα.

\vs{11}Καὶ ἐγένετο ἐν τῷ ἑξῆς ἐπορεύθη εἰς πόλιν καλουμένην Ναΐν καὶ συνεπορεύοντο αὐτῷ οἱ μαθηταὶ αὐτοῦ καὶ ὄχλος πολύς.
\vs{12}ὡς δὲ ἤγγισεν τῇ πύλῃ τῆς πόλεως, καὶ ἰδοὺ ἐξεκομίζετο τεθνηκὼς μονογενὴς υἱὸς τῇ μητρὶ αὐτοῦ καὶ αὐτὴ ἦν χήρα, καὶ ὄχλος τῆς πόλεως ἱκανὸς ἦν σὺν αὐτῇ.
\vs{13}καὶ ἰδὼν αὐτὴν ὁ Κύριος ἐσπλαγχνίσθη ἐπ᾽ αὐτῇ καὶ εἶπεν αὐτῇ· Μὴ κλαῖε.
\vs{14}Καὶ προσελθὼν ἥψατο τῆς σοροῦ, οἱ δὲ βαστάζοντες ἔστησαν, καὶ εἶπεν· Νεανίσκε, σοὶ λέγω, ἐγέρθητι.
\vs{15}καὶ ἀνεκάθισεν ὁ νεκρὸς καὶ ἤρξατο λαλεῖν, καὶ ἔδωκεν αὐτὸν τῇ μητρὶ αὐτοῦ.
\vs{16}Ἔλαβεν δὲ φόβος πάντας καὶ ἐδόξαζον τὸν Θεὸν λέγοντες ὅτι Προφήτης μέγας ἠγέρθη ἐν ἡμῖν καὶ ὅτι Ἐπεσκέψατο ὁ Θεὸς τὸν λαὸν αὐτοῦ.
\vs{17}καὶ ἐξῆλθεν ὁ λόγος οὗτος ἐν ὅλῃ τῇ Ἰουδαίᾳ περὶ αὐτοῦ καὶ πάσῃ τῇ περιχώρῳ.

\vs{18}Καὶ ἀπήγγειλαν Ἰωάννῃ οἱ μαθηταὶ αὐτοῦ περὶ πάντων τούτων. καὶ προσκαλεσάμενος δύο τινὰς τῶν μαθητῶν αὐτοῦ ὁ Ἰωάννης
\vs{19}ἔπεμψεν πρὸς τὸν Κύριον λέγων· Σὺ εἶ ὁ ἐρχόμενος ἢ ἄλλον προσδοκῶμεν;
\vs{20}Παραγενόμενοι δὲ πρὸς αὐτὸν οἱ ἄνδρες εἶπαν· Ἰωάννης ὁ Βαπτιστὴς ἀπέστειλεν ἡμᾶς πρὸς σὲ λέγων· Σὺ εἶ ὁ ἐρχόμενος ἢ ἄλλον προσδοκῶμεν;
\vs{21}Ἐν ἐκείνῃ τῇ ὥρᾳ ἐθεράπευσεν πολλοὺς ἀπὸ νόσων καὶ μαστίγων καὶ πνευμάτων πονηρῶν καὶ τυφλοῖς πολλοῖς ἐχαρίσατο βλέπειν.
\vs{22}καὶ ἀποκριθεὶς εἶπεν αὐτοῖς· Πορευθέντες ἀπαγγείλατε Ἰωάννῃ ἃ εἴδετε καὶ ἠκούσατε· 
\begin{poetryblock}

\begin{quote}τυφλοὶ ἀναβλέπουσιν, χωλοὶ περιπατοῦσιν,\end{quote} 

\begin{quote}λεπροὶ καθαρίζονται καὶ κωφοὶ ἀκούουσιν,\end{quote} 

\begin{quote}νεκροὶ ἐγείρονται, πτωχοὶ εὐαγγελίζονται·\end{quote}
\end{poetryblock}

\vs{23}καὶ μακάριός ἐστιν ὃς ἐὰν μὴ σκανδαλισθῇ ἐν ἐμοί.

\vs{24}Ἀπελθόντων δὲ τῶν ἀγγέλων Ἰωάννου ἤρξατο λέγειν πρὸς τοὺς ὄχλους περὶ Ἰωάννου· Τί ἐξήλθατε εἰς τὴν ἔρημον θεάσασθαι; κάλαμον ὑπὸ ἀνέμου σαλευόμενον;
\vs{25}ἀλλὰ τί ἐξήλθατε ἰδεῖν; ἄνθρωπον ἐν μαλακοῖς ἱματίοις ἠμφιεσμένον; ἰδοὺ οἱ ἐν ἱματισμῷ ἐνδόξῳ καὶ τρυφῇ ὑπάρχοντες ἐν τοῖς βασιλείοις εἰσίν.
\vs{26}Ἀλλὰ τί ἐξήλθατε ἰδεῖν; προφήτην; ναί λέγω ὑμῖν, καὶ περισσότερον προφήτου.
\vs{27}οὗτός ἐστιν περὶ οὗ γέγραπται· 
\begin{poetryblock}

\begin{quote}Ἰδοὺ ἀποστέλλω τὸν ἄγγελόν μου πρὸ προσώπου σου,\end{quote} 

\begin{quote}Ὃς κατασκευάσει τὴν ὁδόν σου ἔμπροσθέν σου.\end{quote}
\end{poetryblock}

\vs{28}Λέγω ὑμῖν, μείζων ἐν γεννητοῖς γυναικῶν Ἰωάννου οὐδείς ἐστιν· ὁ δὲ μικρότερος ἐν τῇ βασιλείᾳ τοῦ Θεοῦ μείζων αὐτοῦ ἐστιν.

\vs{29}Καὶ πᾶς ὁ λαὸς ἀκούσας καὶ οἱ τελῶναι ἐδικαίωσαν τὸν Θεόν βαπτισθέντες τὸ βάπτισμα Ἰωάννου·
\vs{30}οἱ δὲ Φαρισαῖοι καὶ οἱ νομικοὶ τὴν βουλὴν τοῦ Θεοῦ ἠθέτησαν εἰς ἑαυτούς μὴ βαπτισθέντες ὑπ᾽ αὐτοῦ.

\vs{31}Τίνι οὖν ὁμοιώσω τοὺς ἀνθρώπους τῆς γενεᾶς ταύτης καὶ τίνι εἰσὶν ὅμοιοι;
\vs{32}ὅμοιοί εἰσιν παιδίοις τοῖς ἐν ἀγορᾷ καθημένοις καὶ προσφωνοῦσιν ἀλλήλοις ἃ λέγει· 
\begin{poetryblock}

\begin{quote}Ηὐλήσαμεν ὑμῖν καὶ οὐκ ὠρχήσασθε,\end{quote} 

\begin{quote}Ἐθρηνήσαμεν καὶ οὐκ ἐκλαύσατε.\end{quote}
\end{poetryblock}

\vs{33}Ἐλήλυθεν γὰρ Ἰωάννης ὁ Βαπτιστὴς μὴ ἐσθίων ἄρτον μήτε πίνων οἶνον, καὶ λέγετε· Δαιμόνιον ἔχει.
\vs{34}ἐλήλυθεν ὁ Υἱὸς τοῦ ἀνθρώπου ἐσθίων καὶ πίνων, καὶ λέγετε· Ἰδοὺ ἄνθρωπος φάγος καὶ οἰνοπότης, φίλος τελωνῶν καὶ ἁμαρτωλῶν.
\vs{35}καὶ ἐδικαιώθη ἡ σοφία ἀπὸ πάντων τῶν τέκνων αὐτῆς.

\vs{36}Ἠρώτα δέ τις αὐτὸν τῶν Φαρισαίων ἵνα φάγῃ μετ᾽ αὐτοῦ, καὶ εἰσελθὼν εἰς τὸν οἶκον τοῦ Φαρισαίου κατεκλίθη.
\vs{37}καὶ ἰδοὺ γυνὴ ἥτις ἦν ἐν τῇ πόλει ἁμαρτωλός, καὶ ἐπιγνοῦσα ὅτι κατάκειται ἐν τῇ οἰκίᾳ τοῦ Φαρισαίου, κομίσασα ἀλάβαστρον μύρου
\vs{38}καὶ στᾶσα ὀπίσω παρὰ τοὺς πόδας αὐτοῦ κλαίουσα τοῖς δάκρυσιν ἤρξατο βρέχειν τοὺς πόδας αὐτοῦ καὶ ταῖς θριξὶν τῆς κεφαλῆς αὐτῆς ἐξέμασσεν καὶ κατεφίλει τοὺς πόδας αὐτοῦ καὶ ἤλειφεν τῷ μύρῳ.
\vs{39}Ἰδὼν δὲ ὁ Φαρισαῖος ὁ καλέσας αὐτὸν εἶπεν ἐν ἑαυτῷ λέγων· Οὗτος εἰ ἦν προφήτης, ἐγίνωσκεν ἂν τίς καὶ ποταπὴ ἡ γυνὴ ἥτις ἅπτεται αὐτοῦ, ὅτι ἁμαρτωλός ἐστιν.
\vs{40}Καὶ ἀποκριθεὶς ὁ Ἰησοῦς εἶπεν πρὸς αὐτόν· Σίμων, ἔχω σοί τι εἰπεῖν. Ὁ δέ· Διδάσκαλε, εἰπέ, φησίν.
\vs{41}Δύο χρεοφειλέται ἦσαν δανιστῇ τινι· ὁ εἷς ὤφειλεν δηνάρια πεντακόσια, ὁ δὲ ἕτερος πεντήκοντα.
\vs{42}μὴ ἐχόντων αὐτῶν ἀποδοῦναι ἀμφοτέροις ἐχαρίσατο. τίς οὖν αὐτῶν πλεῖον ἀγαπήσει αὐτόν;
\vs{43}Ἀποκριθεὶς Σίμων εἶπεν· Ὑπολαμβάνω ὅτι ᾧ τὸ πλεῖον ἐχαρίσατο. Ὁ δὲ εἶπεν αὐτῷ· Ὀρθῶς ἔκρινας.
\vs{44}Καὶ στραφεὶς πρὸς τὴν γυναῖκα τῷ Σίμωνι ἔφη· Βλέπεις ταύτην τὴν γυναῖκα; εἰσῆλθόν σου εἰς τὴν οἰκίαν, ὕδωρ μοι ἐπὶ πόδας οὐκ ἔδωκας· αὕτη δὲ τοῖς δάκρυσιν ἔβρεξέν μου τοὺς πόδας καὶ ταῖς θριξὶν αὐτῆς ἐξέμαξεν.
\vs{45}φίλημά μοι οὐκ ἔδωκας· αὕτη δὲ ἀφ᾽ ἧς εἰσῆλθον οὐ διέλιπεν καταφιλοῦσά μου τοὺς πόδας.
\vs{46}ἐλαίῳ τὴν κεφαλήν μου οὐκ ἤλειψας· αὕτη δὲ μύρῳ ἤλειψεν τοὺς πόδας μου.
\vs{47}οὗ χάριν λέγω σοι, ἀφέωνται αἱ ἁμαρτίαι αὐτῆς αἱ πολλαί, ὅτι ἠγάπησεν πολύ· ᾧ δὲ ὀλίγον ἀφίεται, ὀλίγον ἀγαπᾷ.
\vs{48}Εἶπεν δὲ αὐτῇ· Ἀφέωνταί σου αἱ ἁμαρτίαι.
\vs{49}Καὶ ἤρξαντο οἱ συνανακείμενοι λέγειν ἐν ἑαυτοῖς· Τίς οὗτός ἐστιν ὃς καὶ ἁμαρτίας ἀφίησιν;
\vs{50}Εἶπεν δὲ πρὸς τὴν γυναῖκα· Ἡ πίστις σου σέσωκέν σε· πορεύου εἰς εἰρήνην.

\ch{8}
Καὶ ἐγένετο ἐν τῷ καθεξῆς καὶ αὐτὸς διώδευεν κατὰ πόλιν καὶ κώμην κηρύσσων καὶ εὐαγγελιζόμενος τὴν βασιλείαν τοῦ Θεοῦ καὶ οἱ δώδεκα σὺν αὐτῷ,
\vs{2}καὶ γυναῖκές τινες αἳ ἦσαν τεθεραπευμέναι ἀπὸ πνευμάτων πονηρῶν καὶ ἀσθενειῶν, Μαρία ἡ καλουμένη Μαγδαληνή, ἀφ᾽ ἧς δαιμόνια ἑπτὰ ἐξεληλύθει,
\vs{3}καὶ Ἰωάννα γυνὴ Χουζᾶ ἐπιτρόπου Ἡρῴδου καὶ Σουσάννα καὶ ἕτεραι πολλαί, αἵτινες διηκόνουν αὐτοῖς ἐκ τῶν ὑπαρχόντων αὐταῖς.

\vs{4}Συνιόντος δὲ ὄχλου πολλοῦ καὶ τῶν κατὰ πόλιν ἐπιπορευομένων πρὸς αὐτὸν εἶπεν διὰ παραβολῆς·
\vs{5}Ἐξῆλθεν ὁ σπείρων τοῦ σπεῖραι τὸν σπόρον αὐτοῦ. καὶ ἐν τῷ σπείρειν αὐτὸν ὃ μὲν ἔπεσεν παρὰ τὴν ὁδόν καὶ κατεπατήθη, καὶ τὰ πετεινὰ τοῦ οὐρανοῦ κατέφαγεν αὐτό.
\vs{6}Καὶ ἕτερον κατέπεσεν ἐπὶ τὴν πέτραν, καὶ φυὲν ἐξηράνθη διὰ τὸ μὴ ἔχειν ἰκμάδα.
\vs{7}Καὶ ἕτερον ἔπεσεν ἐν μέσῳ τῶν ἀκανθῶν, καὶ συμφυεῖσαι αἱ ἄκανθαι ἀπέπνιξαν αὐτό.
\vs{8}Καὶ ἕτερον ἔπεσεν εἰς τὴν γῆν τὴν ἀγαθήν καὶ φυὲν ἐποίησεν καρπὸν ἑκατονταπλασίονα. Ταῦτα λέγων ἐφώνει· Ὁ ἔχων ὦτα ἀκούειν ἀκουέτω.

\vs{9}Ἐπηρώτων δὲ αὐτὸν οἱ μαθηταὶ αὐτοῦ Τίς αὕτη εἴη ἡ παραβολή.
\vs{10}Ὁ δὲ εἶπεν· Ὑμῖν δέδοται γνῶναι τὰ μυστήρια τῆς βασιλείας τοῦ Θεοῦ, τοῖς δὲ λοιποῖς ἐν παραβολαῖς, ἵνα 
\begin{poetryblock}

\begin{quote}Βλέποντες μὴ βλέπωσιν\end{quote} 

\begin{quote}Καὶ ἀκούοντες μὴ συνιῶσιν.\end{quote}
\end{poetryblock}

\vs{11}Ἔστιν δὲ αὕτη ἡ παραβολή· ὁ σπόρος ἐστὶν ὁ λόγος τοῦ Θεοῦ.
\vs{12}οἱ δὲ παρὰ τὴν ὁδόν εἰσιν οἱ ἀκούσαντες, εἶτα ἔρχεται ὁ διάβολος καὶ αἴρει τὸν λόγον ἀπὸ τῆς καρδίας αὐτῶν, ἵνα μὴ πιστεύσαντες σωθῶσιν.
\vs{13}Οἱ δὲ ἐπὶ τῆς πέτρας οἳ ὅταν ἀκούσωσιν μετὰ χαρᾶς δέχονται τὸν λόγον, καὶ οὗτοι ῥίζαν οὐκ ἔχουσιν, οἳ πρὸς καιρὸν πιστεύουσιν καὶ ἐν καιρῷ πειρασμοῦ ἀφίστανται.
\vs{14}Τὸ δὲ εἰς τὰς ἀκάνθας πεσόν, οὗτοί εἰσιν οἱ ἀκούσαντες, καὶ ὑπὸ μεριμνῶν καὶ πλούτου καὶ ἡδονῶν τοῦ βίου πορευόμενοι συμπνίγονται καὶ οὐ τελεσφοροῦσιν.
\vs{15}Τὸ δὲ ἐν τῇ καλῇ γῇ, οὗτοί εἰσιν οἵτινες ἐν καρδίᾳ καλῇ καὶ ἀγαθῇ ἀκούσαντες τὸν λόγον κατέχουσιν καὶ καρποφοροῦσιν ἐν ὑπομονῇ.

\vs{16}Οὐδεὶς δὲ λύχνον ἅψας καλύπτει αὐτὸν σκεύει ἢ ὑποκάτω κλίνης τίθησιν, ἀλλ᾽ ἐπὶ λυχνίας τίθησιν, ἵνα οἱ εἰσπορευόμενοι βλέπωσιν τὸ φῶς.
\vs{17}οὐ γάρ ἐστιν κρυπτὸν ὃ οὐ φανερὸν γενήσεται οὐδὲ ἀπόκρυφον ὃ οὐ μὴ γνωσθῇ καὶ εἰς φανερὸν ἔλθῃ.

\vs{18}Βλέπετε οὖν πῶς ἀκούετε· ὃς ἂν γὰρ ἔχῃ, δοθήσεται αὐτῷ· καὶ ὃς ἂν μὴ ἔχῃ, καὶ ὃ δοκεῖ ἔχειν ἀρθήσεται ἀπ᾽ αὐτοῦ.

\vs{19}Παρεγένετο δὲ πρὸς αὐτὸν ἡ μήτηρ καὶ οἱ ἀδελφοὶ αὐτοῦ καὶ οὐκ ἠδύναντο συντυχεῖν αὐτῷ διὰ τὸν ὄχλον.
\vs{20}ἀπηγγέλη δὲ αὐτῷ· Ἡ μήτηρ σου καὶ οἱ ἀδελφοί σου ἑστήκασιν ἔξω ἰδεῖν θέλοντές σε.
\vs{21}Ὁ δὲ ἀποκριθεὶς εἶπεν πρὸς αὐτούς· Μήτηρ μου καὶ ἀδελφοί μου οὗτοί εἰσιν οἱ τὸν λόγον τοῦ Θεοῦ ἀκούοντες καὶ ποιοῦντες.

\vs{22}Ἐγένετο δὲ ἐν μιᾷ τῶν ἡμερῶν καὶ αὐτὸς ἐνέβη εἰς πλοῖον καὶ οἱ μαθηταὶ αὐτοῦ καὶ εἶπεν πρὸς αὐτούς· Διέλθωμεν εἰς τὸ πέραν τῆς λίμνης, καὶ ἀνήχθησαν.
\vs{23}πλεόντων δὲ αὐτῶν ἀφύπνωσεν. καὶ κατέβη λαῖλαψ ἀνέμου εἰς τὴν λίμνην καὶ συνεπληροῦντο καὶ ἐκινδύνευον.
\vs{24}Προσελθόντες δὲ διήγειραν αὐτὸν λέγοντες· Ἐπιστάτα ἐπιστάτα, ἀπολλύμεθα. Ὁ δὲ διεγερθεὶς ἐπετίμησεν τῷ ἀνέμῳ καὶ τῷ κλύδωνι τοῦ ὕδατος· καὶ ἐπαύσαντο καὶ ἐγένετο γαλήνη.
\vs{25}Εἶπεν δὲ αὐτοῖς· Ποῦ ἡ πίστις ὑμῶν; Φοβηθέντες δὲ ἐθαύμασαν λέγοντες πρὸς ἀλλήλους· Τίς ἄρα οὗτός ἐστιν ὅτι καὶ τοῖς ἀνέμοις ἐπιτάσσει καὶ τῷ ὕδατι, καὶ ὑπακούουσιν αὐτῷ;

\vs{26}Καὶ κατέπλευσαν εἰς τὴν χώραν τῶν Γερασηνῶν, ἥτις ἐστὶν ἀντιπέρα τῆς Γαλιλαίας.
\vs{27}ἐξελθόντι δὲ αὐτῷ ἐπὶ τὴν γῆν ὑπήντησεν ἀνήρ τις ἐκ τῆς πόλεως ἔχων δαιμόνια καὶ χρόνῳ ἱκανῷ οὐκ ἐνεδύσατο ἱμάτιον καὶ ἐν οἰκίᾳ οὐκ ἔμενεν ἀλλ᾽ ἐν τοῖς μνήμασιν.
\vs{28}Ἰδὼν δὲ τὸν Ἰησοῦν ἀνακράξας προσέπεσεν αὐτῷ καὶ φωνῇ μεγάλῃ εἶπεν· Τί ἐμοὶ καὶ σοί, Ἰησοῦ Υἱὲ τοῦ Θεοῦ τοῦ Ὑψίστου; δέομαί σου, μή με βασανίσῃς.
\vs{29}παρήγγειλεν γὰρ τῷ πνεύματι τῷ ἀκαθάρτῳ ἐξελθεῖν ἀπὸ τοῦ ἀνθρώπου. πολλοῖς γὰρ χρόνοις συνηρπάκει αὐτόν καὶ ἐδεσμεύετο ἁλύσεσιν καὶ πέδαις φυλασσόμενος καὶ διαρρήσσων τὰ δεσμὰ ἠλαύνετο ὑπὸ τοῦ δαιμονίου εἰς τὰς ἐρήμους.
\vs{30}Ἐπηρώτησεν δὲ αὐτὸν ὁ Ἰησοῦς· Τί σοι ὄνομά ἐστιν; ὁ Δὲ εἶπεν· Λεγιών, ὅτι εἰσῆλθεν δαιμόνια πολλὰ εἰς αὐτόν.
\vs{31}καὶ παρεκάλουν αὐτὸν ἵνα μὴ ἐπιτάξῃ αὐτοῖς εἰς τὴν ἄβυσσον ἀπελθεῖν.
\vs{32}Ἦν δὲ ἐκεῖ ἀγέλη χοίρων ἱκανῶν βοσκομένη ἐν τῷ ὄρει· καὶ παρεκάλεσαν αὐτὸν ἵνα ἐπιτρέψῃ αὐτοῖς εἰς ἐκείνους εἰσελθεῖν· καὶ ἐπέτρεψεν αὐτοῖς.
\vs{33}Ἐξελθόντα δὲ τὰ δαιμόνια ἀπὸ τοῦ ἀνθρώπου εἰσῆλθον εἰς τοὺς χοίρους, καὶ ὥρμησεν ἡ ἀγέλη κατὰ τοῦ κρημνοῦ εἰς τὴν λίμνην καὶ ἀπεπνίγη.

\vs{34}Ἰδόντες δὲ οἱ βόσκοντες τὸ γεγονὸς ἔφυγον καὶ ἀπήγγειλαν εἰς τὴν πόλιν καὶ εἰς τοὺς ἀγρούς.
\vs{35}ἐξῆλθον δὲ ἰδεῖν τὸ γεγονὸς καὶ ἦλθον πρὸς τὸν Ἰησοῦν καὶ εὗρον καθήμενον τὸν ἄνθρωπον ἀφ᾽ οὗ τὰ δαιμόνια ἐξῆλθεν ἱματισμένον καὶ σωφρονοῦντα παρὰ τοὺς πόδας τοῦ Ἰησοῦ, καὶ ἐφοβήθησαν.
\vs{36}ἀπήγγειλαν δὲ αὐτοῖς οἱ ἰδόντες πῶς ἐσώθη ὁ δαιμονισθείς.
\vs{37}Καὶ ἠρώτησεν αὐτὸν ἅπαν τὸ πλῆθος τῆς περιχώρου τῶν Γερασηνῶν ἀπελθεῖν ἀπ᾽ αὐτῶν, ὅτι φόβῳ μεγάλῳ συνείχοντο· αὐτὸς δὲ ἐμβὰς εἰς πλοῖον ὑπέστρεψεν.
\vs{38}Ἐδεῖτο δὲ αὐτοῦ ὁ ἀνὴρ ἀφ᾽ οὗ ἐξεληλύθει τὰ δαιμόνια εἶναι σὺν αὐτῷ· ἀπέλυσεν δὲ αὐτὸν λέγων·
\vs{39}Ὑπόστρεφε εἰς τὸν οἶκόν σου καὶ διηγοῦ ὅσα σοι ἐποίησεν ὁ Θεός. καὶ ἀπῆλθεν καθ᾽ ὅλην τὴν πόλιν κηρύσσων ὅσα ἐποίησεν αὐτῷ ὁ Ἰησοῦς.

\vs{40}Ἐν δὲ τῷ ὑποστρέφειν τὸν Ἰησοῦν ἀπεδέξατο αὐτὸν ὁ ὄχλος· ἦσαν γὰρ πάντες προσδοκῶντες αὐτόν.
\vs{41}καὶ ἰδοὺ ἦλθεν ἀνὴρ ᾧ ὄνομα Ἰάϊρος καὶ οὗτος ἄρχων τῆς συναγωγῆς ὑπῆρχεν, καὶ πεσὼν παρὰ τοὺς πόδας τοῦ Ἰησοῦ παρεκάλει αὐτὸν εἰσελθεῖν εἰς τὸν οἶκον αὐτοῦ,
\vs{42}ὅτι θυγάτηρ μονογενὴς ἦν αὐτῷ ὡς ἐτῶν δώδεκα καὶ αὐτὴ ἀπέθνῃσκεν. Ἐν δὲ τῷ ὑπάγειν αὐτὸν οἱ ὄχλοι συνέπνιγον αὐτόν.
\vs{43}καὶ γυνὴ οὖσα ἐν ῥύσει αἵματος ἀπὸ ἐτῶν δώδεκα, ἥτις ἰατροῖς προσαναλώσασα ὅλον τὸν βίον οὐκ ἴσχυσεν ἀπ᾽ οὐδενὸς θεραπευθῆναι,
\vs{44}προσελθοῦσα ὄπισθεν ἥψατο τοῦ κρασπέδου τοῦ ἱματίου αὐτοῦ καὶ παραχρῆμα ἔστη ἡ ῥύσις τοῦ αἵματος αὐτῆς.
\vs{45}Καὶ εἶπεν ὁ Ἰησοῦς· Τίς ὁ ἁψάμενός μου; Ἀρνουμένων δὲ πάντων εἶπεν ὁ Πέτρος· Ἐπιστάτα, οἱ ὄχλοι συνέχουσίν σε καὶ ἀποθλίβουσιν.
\vs{46}Ὁ δὲ Ἰησοῦς εἶπεν· Ἥψατό μού τις, ἐγὼ γὰρ ἔγνων δύναμιν ἐξεληλυθυῖαν ἀπ᾽ ἐμοῦ.
\vs{47}Ἰδοῦσα δὲ ἡ γυνὴ ὅτι οὐκ ἔλαθεν, τρέμουσα ἦλθεν καὶ προσπεσοῦσα αὐτῷ δι᾽ ἣν αἰτίαν ἥψατο αὐτοῦ ἀπήγγειλεν ἐνώπιον παντὸς τοῦ λαοῦ καὶ ὡς ἰάθη παραχρῆμα.
\vs{48}Ὁ δὲ εἶπεν αὐτῇ· Θυγάτηρ, ἡ πίστις σου σέσωκέν σε· πορεύου εἰς εἰρήνην.
\vs{49}Ἔτι αὐτοῦ λαλοῦντος ἔρχεταί τις παρὰ τοῦ ἀρχισυναγώγου λέγων ὅτι Τέθνηκεν ἡ θυγάτηρ σου· μηκέτι σκύλλε τὸν Διδάσκαλον.
\vs{50}Ὁ δὲ Ἰησοῦς ἀκούσας ἀπεκρίθη αὐτῷ· Μὴ φοβοῦ, μόνον πίστευσον, καὶ σωθήσεται.
\vs{51}Ἐλθὼν δὲ εἰς τὴν οἰκίαν οὐκ ἀφῆκεν εἰσελθεῖν τινα σὺν αὐτῷ εἰ μὴ Πέτρον καὶ Ἰωάννην καὶ Ἰάκωβον καὶ τὸν πατέρα τῆς παιδὸς καὶ τὴν μητέρα.
\vs{52}ἔκλαιον δὲ πάντες καὶ ἐκόπτοντο αὐτήν. ὁ δὲ εἶπεν· Μὴ κλαίετε, οὐ γὰρ ἀπέθανεν ἀλλὰ καθεύδει.
\vs{53}Καὶ κατεγέλων αὐτοῦ εἰδότες ὅτι ἀπέθανεν.
\vs{54}Αὐτὸς δὲ κρατήσας τῆς χειρὸς αὐτῆς ἐφώνησεν λέγων· Ἡ Παῖς, ἔγειρε.
\vs{55}καὶ ἐπέστρεψεν τὸ πνεῦμα αὐτῆς καὶ ἀνέστη παραχρῆμα καὶ διέταξεν αὐτῇ δοθῆναι φαγεῖν.
\vs{56}καὶ ἐξέστησαν οἱ γονεῖς αὐτῆς· ὁ δὲ παρήγγειλεν αὐτοῖς μηδενὶ εἰπεῖν τὸ γεγονός.

\ch{9}
Συνκαλεσάμενος δὲ τοὺς δώδεκα ἔδωκεν αὐτοῖς δύναμιν καὶ ἐξουσίαν ἐπὶ πάντα τὰ δαιμόνια καὶ νόσους θεραπεύειν
\vs{2}καὶ ἀπέστειλεν αὐτοὺς κηρύσσειν τὴν βασιλείαν τοῦ Θεοῦ καὶ ἰᾶσθαι τοὺς ἀσθενεῖς,
\vs{3}καὶ εἶπεν πρὸς αὐτούς· Μηδὲν αἴρετε εἰς τὴν ὁδόν, μήτε ῥάβδον μήτε πήραν μήτε ἄρτον μήτε ἀργύριον μήτε ἀνὰ δύο χιτῶνας ἔχειν.
\vs{4}καὶ εἰς ἣν ἂν οἰκίαν εἰσέλθητε, ἐκεῖ μένετε καὶ ἐκεῖθεν ἐξέρχεσθε.
\vs{5}καὶ ὅσοι ἂν μὴ δέχωνται ὑμᾶς, ἐξερχόμενοι ἀπὸ τῆς πόλεως ἐκείνης τὸν κονιορτὸν ἀπὸ τῶν ποδῶν ὑμῶν ἀποτινάσσετε εἰς μαρτύριον ἐπ᾽ αὐτούς.
\vs{6}Ἐξερχόμενοι δὲ διήρχοντο κατὰ τὰς κώμας εὐαγγελιζόμενοι καὶ θεραπεύοντες πανταχοῦ.

\vs{7}Ἤκουσεν δὲ Ἡρῴδης ὁ τετραάρχης τὰ γινόμενα πάντα καὶ διηπόρει διὰ τὸ λέγεσθαι ὑπό τινων ὅτι Ἰωάννης ἠγέρθη ἐκ νεκρῶν,
\vs{8}ὑπό τινων δὲ ὅτι Ἠλίας ἐφάνη, ἄλλων δὲ ὅτι προφήτης τις τῶν ἀρχαίων ἀνέστη.
\vs{9}Εἶπεν δὲ Ἡρῴδης· Ἰωάννην ἐγὼ ἀπεκεφάλισα· τίς δέ ἐστιν οὗτος περὶ οὗ ἀκούω τοιαῦτα; καὶ ἐζήτει ἰδεῖν αὐτόν.

\vs{10}Καὶ ὑποστρέψαντες οἱ ἀπόστολοι διηγήσαντο αὐτῷ ὅσα ἐποίησαν. Καὶ παραλαβὼν αὐτοὺς ὑπεχώρησεν κατ᾽ ἰδίαν εἰς πόλιν καλουμένην Βηθσαϊδά.
\vs{11}οἱ δὲ ὄχλοι γνόντες ἠκολούθησαν αὐτῷ· καὶ ἀποδεξάμενος αὐτοὺς ἐλάλει αὐτοῖς περὶ τῆς βασιλείας τοῦ Θεοῦ, καὶ τοὺς χρείαν ἔχοντας θεραπείας ἰᾶτο.

\vs{12}Ἡ δὲ ἡμέρα ἤρξατο κλίνειν· προσελθόντες δὲ οἱ δώδεκα εἶπαν αὐτῷ· Ἀπόλυσον τὸν ὄχλον, ἵνα πορευθέντες εἰς τὰς κύκλῳ κώμας καὶ ἀγροὺς καταλύσωσιν καὶ εὕρωσιν ἐπισιτισμόν, ὅτι ὧδε ἐν ἐρήμῳ τόπῳ ἐσμέν.
\vs{13}Εἶπεν δὲ πρὸς αὐτούς· Δότε αὐτοῖς ὑμεῖς φαγεῖν. Οἱ δὲ εἶπαν· Οὐκ εἰσὶν ἡμῖν πλεῖον ἢ ἄρτοι πέντε καὶ ἰχθύες δύο, εἰ μήτι πορευθέντες ἡμεῖς ἀγοράσωμεν εἰς πάντα τὸν λαὸν τοῦτον βρώματα.
\vs{14}ἦσαν γὰρ ὡσεὶ ἄνδρες πεντακισχίλιοι. Εἶπεν δὲ πρὸς τοὺς μαθητὰς αὐτοῦ· Κατακλίνατε αὐτοὺς κλισίας ὡσεὶ ἀνὰ πεντήκοντα.
\vs{15}καὶ ἐποίησαν οὕτως καὶ κατέκλιναν ἅπαντας.
\vs{16}Λαβὼν δὲ τοὺς πέντε ἄρτους καὶ τοὺς δύο ἰχθύας ἀναβλέψας εἰς τὸν οὐρανὸν εὐλόγησεν αὐτοὺς καὶ κατέκλασεν καὶ ἐδίδου τοῖς μαθηταῖς παραθεῖναι τῷ ὄχλῳ.
\vs{17}Καὶ ἔφαγον καὶ ἐχορτάσθησαν πάντες, καὶ ἤρθη τὸ περισσεῦσαν αὐτοῖς κλασμάτων κόφινοι δώδεκα.

\vs{18}Καὶ ἐγένετο ἐν τῷ εἶναι αὐτὸν προσευχόμενον κατὰ μόνας συνῆσαν αὐτῷ οἱ μαθηταί, καὶ ἐπηρώτησεν αὐτοὺς λέγων· Τίνα με λέγουσιν οἱ ὄχλοι εἶναι;
\vs{19}Οἱ δὲ ἀποκριθέντες εἶπαν· Ἰωάννην τὸν Βαπτιστήν, ἄλλοι δὲ Ἠλίαν, ἄλλοι δὲ ὅτι προφήτης τις τῶν ἀρχαίων ἀνέστη.
\vs{20}Εἶπεν δὲ αὐτοῖς· Ὑμεῖς δὲ τίνα με λέγετε εἶναι; Πέτρος δὲ ἀποκριθεὶς εἶπεν· Τὸν Χριστὸν τοῦ Θεοῦ.
\vs{21}Ὁ δὲ ἐπιτιμήσας αὐτοῖς παρήγγειλεν μηδενὶ λέγειν τοῦτο
\vs{22}εἰπὼν ὅτι Δεῖ τὸν Υἱὸν τοῦ ἀνθρώπου πολλὰ παθεῖν καὶ ἀποδοκιμασθῆναι ἀπὸ τῶν πρεσβυτέρων καὶ ἀρχιερέων καὶ γραμματέων καὶ ἀποκτανθῆναι καὶ τῇ τρίτῃ ἡμέρᾳ ἐγερθῆναι.

\vs{23}Ἔλεγεν δὲ πρὸς πάντας· Εἴ τις θέλει ὀπίσω μου ἔρχεσθαι, ἀρνησάσθω ἑαυτὸν καὶ ἀράτω τὸν σταυρὸν αὐτοῦ καθ᾽ ἡμέραν καὶ ἀκολουθείτω μοι.
\vs{24}ὃς γὰρ ἂν θέλῃ τὴν ψυχὴν αὐτοῦ σῶσαι ἀπολέσει αὐτήν· ὃς δ᾽ ἂν ἀπολέσῃ τὴν ψυχὴν αὐτοῦ ἕνεκεν ἐμοῦ οὗτος σώσει αὐτήν.
\vs{25}Τί γὰρ ὠφελεῖται ἄνθρωπος κερδήσας τὸν κόσμον ὅλον ἑαυτὸν δὲ ἀπολέσας ἢ ζημιωθείς;
\vs{26}ὃς γὰρ ἂν ἐπαισχυνθῇ με καὶ τοὺς ἐμοὺς λόγους, τοῦτον ὁ Υἱὸς τοῦ ἀνθρώπου ἐπαισχυνθήσεται, ὅταν ἔλθῃ ἐν τῇ δόξῃ αὐτοῦ καὶ τοῦ Πατρὸς καὶ τῶν ἁγίων ἀγγέλων.
\vs{27}λέγω δὲ ὑμῖν ἀληθῶς, εἰσίν τινες τῶν αὐτοῦ ἑστηκότων οἳ οὐ μὴ γεύσωνται θανάτου ἕως ἂν ἴδωσιν τὴν βασιλείαν τοῦ Θεοῦ.

\vs{28}Ἐγένετο δὲ μετὰ τοὺς λόγους τούτους ὡσεὶ ἡμέραι ὀκτὼ καὶ παραλαβὼν Πέτρον καὶ Ἰωάννην καὶ Ἰάκωβον ἀνέβη εἰς τὸ ὄρος προσεύξασθαι.
\vs{29}καὶ ἐγένετο ἐν τῷ προσεύχεσθαι αὐτὸν τὸ εἶδος τοῦ προσώπου αὐτοῦ ἕτερον καὶ ὁ ἱματισμὸς αὐτοῦ λευκὸς ἐξαστράπτων.
\vs{30}καὶ ἰδοὺ ἄνδρες δύο συνελάλουν αὐτῷ, οἵτινες ἦσαν Μωϋσῆς καὶ Ἠλίας,
\vs{31}οἳ ὀφθέντες ἐν δόξῃ ἔλεγον τὴν ἔξοδον αὐτοῦ, ἣν ἤμελλεν πληροῦν ἐν Ἰερουσαλήμ.
\vs{32}Ὁ δὲ Πέτρος καὶ οἱ σὺν αὐτῷ ἦσαν βεβαρημένοι ὕπνῳ· διαγρηγορήσαντες δὲ εἶδον τὴν δόξαν αὐτοῦ καὶ τοὺς δύο ἄνδρας τοὺς συνεστῶτας αὐτῷ.
\vs{33}καὶ ἐγένετο ἐν τῷ διαχωρίζεσθαι αὐτοὺς ἀπ᾽ αὐτοῦ εἶπεν ὁ Πέτρος πρὸς τὸν Ἰησοῦν· Ἐπιστάτα, καλόν ἐστιν ἡμᾶς ὧδε εἶναι, καὶ ποιήσωμεν σκηνὰς τρεῖς, μίαν σοὶ καὶ μίαν Μωϋσεῖ καὶ μίαν Ἠλίᾳ, μὴ εἰδὼς ὃ λέγει.
\vs{34}Ταῦτα δὲ αὐτοῦ λέγοντος ἐγένετο νεφέλη καὶ ἐπεσκίαζεν αὐτούς· ἐφοβήθησαν δὲ ἐν τῷ εἰσελθεῖν αὐτοὺς εἰς τὴν νεφέλην.
\vs{35}καὶ φωνὴ ἐγένετο ἐκ τῆς νεφέλης λέγουσα· Οὗτός ἐστιν ὁ Υἱός μου ὁ ἐκλελεγμένος, αὐτοῦ ἀκούετε.
\vs{36}καὶ ἐν τῷ γενέσθαι τὴν φωνὴν εὑρέθη Ἰησοῦς μόνος. καὶ αὐτοὶ ἐσίγησαν καὶ οὐδενὶ ἀπήγγειλαν ἐν ἐκείναις ταῖς ἡμέραις οὐδὲν ὧν ἑώρακαν.

\vs{37}Ἐγένετο δὲ τῇ ἑξῆς ἡμέρᾳ κατελθόντων αὐτῶν ἀπὸ τοῦ ὄρους συνήντησεν αὐτῷ ὄχλος πολύς.
\vs{38}καὶ ἰδοὺ ἀνὴρ ἀπὸ τοῦ ὄχλου ἐβόησεν λέγων· Διδάσκαλε, δέομαί σου ἐπιβλέψαι ἐπὶ τὸν υἱόν μου, ὅτι μονογενής μοί ἐστιν,
\vs{39}καὶ ἰδοὺ πνεῦμα λαμβάνει αὐτόν καὶ ἐξαίφνης κράζει καὶ σπαράσσει αὐτὸν μετὰ ἀφροῦ καὶ μόγις ἀποχωρεῖ ἀπ᾽ αὐτοῦ συντρῖβον αὐτόν·
\vs{40}καὶ ἐδεήθην τῶν μαθητῶν σου ἵνα ἐκβάλωσιν αὐτό, καὶ οὐκ ἠδυνήθησαν.
\vs{41}Ἀποκριθεὶς δὲ ὁ Ἰησοῦς εἶπεν· Ὦ γενεὰ ἄπιστος καὶ διεστραμμένη, ἕως πότε ἔσομαι πρὸς ὑμᾶς καὶ ἀνέξομαι ὑμῶν; προσάγαγε ὧδε τὸν υἱόν σου.
\vs{42}Ἔτι δὲ προσερχομένου αὐτοῦ ἔρρηξεν αὐτὸν τὸ δαιμόνιον καὶ συνεσπάραξεν· ἐπετίμησεν δὲ ὁ Ἰησοῦς τῷ πνεύματι τῷ ἀκαθάρτῳ καὶ ἰάσατο τὸν παῖδα καὶ ἀπέδωκεν αὐτὸν τῷ πατρὶ αὐτοῦ.
\vs{43}Ἐξεπλήσσοντο δὲ πάντες ἐπὶ τῇ μεγαλειότητι τοῦ Θεοῦ.

Πάντων δὲ θαυμαζόντων ἐπὶ πᾶσιν οἷς ἐποίει εἶπεν πρὸς τοὺς μαθητὰς αὐτοῦ·
\vs{44}Θέσθε ὑμεῖς εἰς τὰ ὦτα ὑμῶν τοὺς λόγους τούτους· ὁ γὰρ Υἱὸς τοῦ ἀνθρώπου μέλλει παραδίδοσθαι εἰς χεῖρας ἀνθρώπων.
\vs{45}οἱ δὲ ἠγνόουν τὸ ῥῆμα τοῦτο καὶ ἦν παρακεκαλυμμένον ἀπ᾽ αὐτῶν ἵνα μὴ αἴσθωνται αὐτό, καὶ ἐφοβοῦντο ἐρωτῆσαι αὐτὸν περὶ τοῦ ῥήματος τούτου.

\vs{46}Εἰσῆλθεν δὲ διαλογισμὸς ἐν αὐτοῖς, τὸ τίς ἂν εἴη μείζων αὐτῶν.
\vs{47}ὁ δὲ Ἰησοῦς εἰδὼς τὸν διαλογισμὸν τῆς καρδίας αὐτῶν, ἐπιλαβόμενος παιδίον ἔστησεν αὐτὸ παρ᾽ ἑαυτῷ
\vs{48}καὶ εἶπεν αὐτοῖς· Ὃς ἐὰν δέξηται τοῦτο τὸ παιδίον ἐπὶ τῷ ὀνόματί μου, ἐμὲ δέχεται· καὶ ὃς ἂν ἐμὲ δέξηται, δέχεται τὸν ἀποστείλαντά με· ὁ γὰρ μικρότερος ἐν πᾶσιν ὑμῖν ὑπάρχων οὗτός ἐστιν μέγας.

\vs{49}Ἀποκριθεὶς δὲ Ἰωάννης εἶπεν· Ἐπιστάτα, εἴδομέν τινα ἐν τῷ ὀνόματί σου ἐκβάλλοντα δαιμόνια καὶ ἐκωλύομεν αὐτὸν, ὅτι οὐκ ἀκολουθεῖ μεθ᾽ ἡμῶν.
\vs{50}Εἶπεν δὲ πρὸς αὐτὸν ὁ Ἰησοῦς· Μὴ κωλύετε· ὃς γὰρ οὐκ ἔστιν καθ᾽ ὑμῶν, ὑπὲρ ὑμῶν ἐστιν.

\vs{51}Ἐγένετο δὲ ἐν τῷ συμπληροῦσθαι τὰς ἡμέρας τῆς ἀναλήμψεως αὐτοῦ καὶ αὐτὸς τὸ πρόσωπον ἐστήρισεν τοῦ πορεύεσθαι εἰς Ἰερουσαλήμ.
\vs{52}καὶ ἀπέστειλεν ἀγγέλους πρὸ προσώπου αὐτοῦ. καὶ πορευθέντες εἰσῆλθον εἰς κώμην Σαμαριτῶν ὡς ἑτοιμάσαι αὐτῷ·
\vs{53}καὶ οὐκ ἐδέξαντο αὐτόν, ὅτι τὸ πρόσωπον αὐτοῦ ἦν πορευόμενον εἰς Ἰερουσαλήμ.
\vs{54}Ἰδόντες δὲ οἱ μαθηταὶ Ἰάκωβος καὶ Ἰωάννης εἶπαν· Κύριε, θέλεις εἴπωμεν πῦρ καταβῆναι ἀπὸ τοῦ οὐρανοῦ καὶ ἀναλῶσαι αὐτούς;
\vs{55}Στραφεὶς δὲ ἐπετίμησεν αὐτοῖς.
\vs{56}καὶ ἐπορεύθησαν εἰς ἑτέραν κώμην.

\vs{57}Καὶ πορευομένων αὐτῶν ἐν τῇ ὁδῷ εἶπέν τις πρὸς αὐτόν· Ἀκολουθήσω σοι ὅπου ἐὰν ἀπέρχῃ.
\vs{58}Καὶ εἶπεν αὐτῷ ὁ Ἰησοῦς· Αἱ ἀλώπεκες φωλεοὺς ἔχουσιν καὶ τὰ πετεινὰ τοῦ οὐρανοῦ κατασκηνώσεις, ὁ δὲ Υἱὸς τοῦ ἀνθρώπου οὐκ ἔχει ποῦ τὴν κεφαλὴν κλίνῃ.
\vs{59}Εἶπεν δὲ πρὸς ἕτερον· Ἀκολούθει μοι. ὁ Δὲ εἶπεν· Κύριε, Ἐπίτρεψόν μοι ἀπελθόντι πρῶτον θάψαι τὸν πατέρα μου.
\vs{60}Εἶπεν δὲ αὐτῷ· Ἄφες τοὺς νεκροὺς θάψαι τοὺς ἑαυτῶν νεκρούς, σὺ δὲ ἀπελθὼν διάγγελλε τὴν βασιλείαν τοῦ Θεοῦ.
\vs{61}Εἶπεν δὲ καὶ ἕτερος· Ἀκολουθήσω σοι, Κύριε· πρῶτον δὲ ἐπίτρεψόν μοι ἀποτάξασθαι τοῖς εἰς τὸν οἶκόν μου.
\vs{62}Εἶπεν δὲ πρὸς αὐτὸν ὁ Ἰησοῦς· Οὐδεὶς ἐπιβαλὼν τὴν χεῖρα ἐπ᾽ ἄροτρον καὶ βλέπων εἰς τὰ ὀπίσω εὔθετός ἐστιν τῇ βασιλείᾳ τοῦ Θεοῦ.

\ch{10}
Μετὰ δὲ ταῦτα ἀνέδειξεν ὁ Κύριος ἑτέρους ἑβδομήκοντα δύο καὶ ἀπέστειλεν αὐτοὺς ἀνὰ δύο δύο πρὸ προσώπου αὐτοῦ εἰς πᾶσαν πόλιν καὶ τόπον οὗ ἤμελλεν αὐτὸς ἔρχεσθαι.
\vs{2}ἔλεγεν δὲ πρὸς αὐτούς· Ὁ μὲν θερισμὸς πολύς, οἱ δὲ ἐργάται ὀλίγοι· δεήθητε οὖν τοῦ Κυρίου τοῦ θερισμοῦ ὅπως ἐργάτας ἐκβάλῃ εἰς τὸν θερισμὸν αὐτοῦ.
\vs{3}Ὑπάγετε· ἰδοὺ ἀποστέλλω ὑμᾶς ὡς ἄρνας ἐν μέσῳ λύκων.
\vs{4}μὴ βαστάζετε βαλλάντιον, μὴ πήραν, μὴ ὑποδήματα, καὶ μηδένα κατὰ τὴν ὁδὸν ἀσπάσησθε.
\vs{5}Εἰς ἣν δ᾽ ἂν εἰσέλθητε οἰκίαν, πρῶτον λέγετε· Εἰρήνη τῷ οἴκῳ τούτῳ.
\vs{6}καὶ ἐὰν ἐκεῖ ᾖ υἱὸς εἰρήνης, ἐπαναπαήσεται ἐπ᾽ αὐτὸν ἡ εἰρήνη ὑμῶν· εἰ δὲ μή γε, ἐφ᾽ ὑμᾶς ἀνακάμψει.
\vs{7}ἐν αὐτῇ δὲ τῇ οἰκίᾳ μένετε ἐσθίοντες καὶ πίνοντες τὰ παρ᾽ αὐτῶν· ἄξιος γὰρ ὁ ἐργάτης τοῦ μισθοῦ αὐτοῦ. μὴ μεταβαίνετε ἐξ οἰκίας εἰς οἰκίαν.
\vs{8}Καὶ εἰς ἣν ἂν πόλιν εἰσέρχησθε καὶ δέχωνται ὑμᾶς, ἐσθίετε τὰ παρατιθέμενα ὑμῖν
\vs{9}καὶ θεραπεύετε τοὺς ἐν αὐτῇ ἀσθενεῖς καὶ λέγετε αὐτοῖς· Ἤγγικεν ἐφ᾽ ὑμᾶς ἡ βασιλεία τοῦ Θεοῦ.
\vs{10}Εἰς ἣν δ᾽ ἂν πόλιν εἰσέλθητε καὶ μὴ δέχωνται ὑμᾶς, ἐξελθόντες εἰς τὰς πλατείας αὐτῆς εἴπατε·
\vs{11}Καὶ τὸν κονιορτὸν τὸν κολληθέντα ἡμῖν ἐκ τῆς πόλεως ὑμῶν εἰς τοὺς πόδας ἀπομασσόμεθα ὑμῖν· πλὴν τοῦτο γινώσκετε ὅτι ἤγγικεν ἡ βασιλεία τοῦ Θεοῦ.
\vs{12}λέγω ὑμῖν ὅτι Σοδόμοις ἐν τῇ ἡμέρᾳ ἐκείνῃ ἀνεκτότερον ἔσται ἢ τῇ πόλει ἐκείνῃ.

\vs{13}Οὐαί σοι, Χοραζίν, οὐαί σοι, Βηθσαϊδά· ὅτι εἰ ἐν Τύρῳ καὶ Σιδῶνι ἐγενήθησαν αἱ δυνάμεις αἱ γενόμεναι ἐν ὑμῖν, πάλαι ἂν ἐν σάκκῳ καὶ σποδῷ καθήμενοι μετενόησαν.
\vs{14}πλὴν Τύρῳ καὶ Σιδῶνι ἀνεκτότερον ἔσται ἐν τῇ κρίσει ἢ ὑμῖν.
\vs{15}καὶ σύ, Καφαρναούμ, μὴ ἕως οὐρανοῦ ὑψωθήσῃ; ἕως τοῦ ᾅδου καταβήσῃ.

\vs{16}Ὁ ἀκούων ὑμῶν ἐμοῦ ἀκούει, καὶ ὁ ἀθετῶν ὑμᾶς ἐμὲ ἀθετεῖ· ὁ δὲ ἐμὲ ἀθετῶν ἀθετεῖ τὸν ἀποστείλαντά με.

\vs{17}Ὑπέστρεψαν δὲ οἱ ἑβδομήκοντα δύο μετὰ χαρᾶς λέγοντες· Κύριε, καὶ τὰ δαιμόνια ὑποτάσσεται ἡμῖν ἐν τῷ ὀνόματί σου.
\vs{18}Εἶπεν δὲ αὐτοῖς· Ἐθεώρουν τὸν Σατανᾶν ὡς ἀστραπὴν ἐκ τοῦ οὐρανοῦ πεσόντα.
\vs{19}ἰδοὺ δέδωκα ὑμῖν τὴν ἐξουσίαν τοῦ πατεῖν ἐπάνω ὄφεων καὶ σκορπίων, καὶ ἐπὶ πᾶσαν τὴν δύναμιν τοῦ ἐχθροῦ, καὶ οὐδὲν ὑμᾶς οὐ μὴ ἀδικήσῃ.
\vs{20}πλὴν ἐν τούτῳ μὴ χαίρετε ὅτι τὰ πνεύματα ὑμῖν ὑποτάσσεται, χαίρετε δὲ ὅτι τὰ ὀνόματα ὑμῶν ἐνγέγραπται ἐν τοῖς οὐρανοῖς.

\vs{21}Ἐν αὐτῇ τῇ ὥρᾳ ἠγαλλιάσατο ἐν τῷ Πνεύματι τῷ Ἁγίῳ καὶ εἶπεν· Ἐξομολογοῦμαί σοι, Πάτερ, Κύριε τοῦ οὐρανοῦ καὶ τῆς γῆς, ὅτι ἀπέκρυψας ταῦτα ἀπὸ σοφῶν καὶ συνετῶν καὶ ἀπεκάλυψας αὐτὰ νηπίοις· ναί ὁ Πατήρ, ὅτι οὕτως εὐδοκία ἐγένετο ἔμπροσθέν σου.
\vs{22}Πάντα μοι παρεδόθη ὑπὸ τοῦ Πατρός μου, καὶ οὐδεὶς γινώσκει τίς ἐστιν ὁ Υἱὸς εἰ μὴ ὁ Πατήρ, καὶ τίς ἐστιν ὁ Πατὴρ εἰ μὴ ὁ Υἱὸς καὶ ᾧ ἐὰν βούληται ὁ Υἱὸς ἀποκαλύψαι.

\vs{23}Καὶ στραφεὶς πρὸς τοὺς μαθητὰς κατ᾽ ἰδίαν εἶπεν· Μακάριοι οἱ ὀφθαλμοὶ οἱ βλέποντες ἃ βλέπετε.
\vs{24}λέγω γὰρ ὑμῖν ὅτι πολλοὶ προφῆται καὶ βασιλεῖς ἠθέλησαν ἰδεῖν ἃ ὑμεῖς βλέπετε καὶ οὐκ εἶδαν, καὶ ἀκοῦσαι ἃ ἀκούετε καὶ οὐκ ἤκουσαν.

\vs{25}Καὶ ἰδοὺ νομικός τις ἀνέστη ἐκπειράζων αὐτὸν λέγων· Διδάσκαλε, τί ποιήσας ζωὴν αἰώνιον κληρονομήσω;
\vs{26}Ὁ δὲ εἶπεν πρὸς αὐτόν· Ἐν τῷ νόμῳ τί γέγραπται; πῶς ἀναγινώσκεις;
\vs{27}Ὁ δὲ ἀποκριθεὶς εἶπεν· Ἀγαπήσεις Κύριον τὸν Θεόν σου ἐξ ὅλης τῆς καρδίας σου καὶ ἐν ὅλῃ τῇ ψυχῇ σου καὶ ἐν ὅλῃ τῇ ἰσχύϊ σου καὶ ἐν ὅλῃ τῇ διανοίᾳ σου, καὶ Τὸν πλησίον σου ὡς σεαυτόν.
\vs{28}Εἶπεν δὲ αὐτῷ· Ὀρθῶς ἀπεκρίθης· τοῦτο ποίει καὶ ζήσῃ.
\vs{29}Ὁ δὲ θέλων δικαιῶσαι ἑαυτὸν εἶπεν πρὸς τὸν Ἰησοῦν· Καὶ τίς ἐστίν μου πλησίον;

\vs{30}Ὑπολαβὼν ὁ Ἰησοῦς εἶπεν· Ἄνθρωπός τις κατέβαινεν ἀπὸ Ἰερουσαλὴμ εἰς Ἰεριχὼ καὶ λῃσταῖς περιέπεσεν, οἳ καὶ ἐκδύσαντες αὐτὸν καὶ πληγὰς ἐπιθέντες ἀπῆλθον ἀφέντες ἡμιθανῆ.
\vs{31}Κατὰ συγκυρίαν δὲ ἱερεύς τις κατέβαινεν ἐν τῇ ὁδῷ ἐκείνῃ καὶ ἰδὼν αὐτὸν ἀντιπαρῆλθεν·
\vs{32}Ὁμοίως δὲ καὶ Λευίτης γενόμενος κατὰ τὸν τόπον ἐλθὼν καὶ ἰδὼν ἀντιπαρῆλθεν.
\vs{33}Σαμαρίτης δέ τις ὁδεύων ἦλθεν κατ᾽ αὐτὸν καὶ ἰδὼν ἐσπλαγχνίσθη,
\vs{34}καὶ προσελθὼν κατέδησεν τὰ τραύματα αὐτοῦ ἐπιχέων ἔλαιον καὶ οἶνον, ἐπιβιβάσας δὲ αὐτὸν ἐπὶ τὸ ἴδιον κτῆνος ἤγαγεν αὐτὸν εἰς πανδοχεῖον καὶ ἐπεμελήθη αὐτοῦ.
\vs{35}Καὶ ἐπὶ τὴν αὔριον ἐκβαλὼν ἔδωκεν δύο δηνάρια τῷ πανδοχεῖ καὶ εἶπεν· Ἐπιμελήθητι αὐτοῦ, καὶ ὅ τι ἂν προσδαπανήσῃς ἐγὼ ἐν τῷ ἐπανέρχεσθαί με ἀποδώσω σοι.
\vs{36}Τίς τούτων τῶν τριῶν πλησίον δοκεῖ σοι γεγονέναι τοῦ ἐμπεσόντος εἰς τοὺς λῃστάς;
\vs{37}Ὁ δὲ εἶπεν· Ὁ ποιήσας τὸ ἔλεος μετ᾽ αὐτοῦ. Εἶπεν δὲ αὐτῷ ὁ Ἰησοῦς· Πορεύου καὶ σὺ ποίει ὁμοίως.

\vs{38}Ἐν δὲ τῷ πορεύεσθαι αὐτοὺς αὐτὸς εἰσῆλθεν εἰς κώμην τινά· γυνὴ δέ τις ὀνόματι Μάρθα ὑπεδέξατο αὐτὸν.
\vs{39}καὶ τῇδε ἦν ἀδελφὴ καλουμένη Μαριάμ, ἣ καὶ παρακαθεσθεῖσα πρὸς τοὺς πόδας τοῦ Κυρίου ἤκουεν τὸν λόγον αὐτοῦ.
\vs{40}ἡ δὲ Μάρθα περιεσπᾶτο περὶ πολλὴν διακονίαν· ἐπιστᾶσα δὲ εἶπεν· Κύριε, οὐ μέλει σοι ὅτι ἡ ἀδελφή μου μόνην με κατέλιπεν διακονεῖν; εἰπὲ οὖν αὐτῇ ἵνα μοι συναντιλάβηται.
\vs{41}Ἀποκριθεὶς δὲ εἶπεν αὐτῇ ὁ Κύριος· Μάρθα Μάρθα, μεριμνᾷς καὶ θορυβάζῃ περὶ πολλά,
\vs{42}ἑνός δέ ἐστιν χρεία· Μαριὰμ γὰρ τὴν ἀγαθὴν μερίδα ἐξελέξατο ἥτις οὐκ ἀφαιρεθήσεται αὐτῆς.

\ch{11}
Καὶ ἐγένετο ἐν τῷ εἶναι αὐτὸν ἐν τόπῳ τινὶ προσευχόμενον, ὡς ἐπαύσατο, εἶπέν τις τῶν μαθητῶν αὐτοῦ πρὸς αὐτόν· Κύριε, δίδαξον ἡμᾶς προσεύχεσθαι, καθὼς καὶ Ἰωάννης ἐδίδαξεν τοὺς μαθητὰς αὐτοῦ.
\vs{2}Εἶπεν δὲ αὐτοῖς· Ὅταν προσεύχησθε λέγετε· 
\begin{poetryblock}

\begin{quote}Πάτερ, ἁγιασθήτω τὸ ὄνομά σου· Ἐλθέτω ἡ βασιλεία σου·\end{quote}

\begin{quote} \vs{3}Τὸν ἄρτον ἡμῶν τὸν ἐπιούσιον δίδου ἡμῖν τὸ καθ᾽ ἡμέραν·\end{quote}

\begin{quote} \vs{4}Καὶ ἄφες ἡμῖν τὰς ἁμαρτίας ἡμῶν,\end{quote} 

\begin{quote}Καὶ γὰρ αὐτοὶ ἀφίομεν παντὶ ὀφείλοντι ἡμῖν·\end{quote} 

\begin{quote}Καὶ μὴ εἰσενέγκῃς ἡμᾶς εἰς πειρασμόν.\end{quote}
\end{poetryblock}
\vs{5}Καὶ εἶπεν πρὸς αὐτούς· Τίς ἐξ ὑμῶν ἕξει φίλον καὶ πορεύσεται πρὸς αὐτὸν μεσονυκτίου καὶ εἴπῃ αὐτῷ· Φίλε, χρῆσόν μοι τρεῖς ἄρτους,
\vs{6}ἐπειδὴ φίλος μου παρεγένετο ἐξ ὁδοῦ πρός με καὶ οὐκ ἔχω ὃ παραθήσω αὐτῷ·
\vs{7}Κἀκεῖνος ἔσωθεν ἀποκριθεὶς εἴπῃ· Μή μοι κόπους πάρεχε· ἤδη ἡ θύρα κέκλεισται καὶ τὰ παιδία μου μετ᾽ ἐμοῦ εἰς τὴν κοίτην εἰσίν· οὐ δύναμαι ἀναστὰς δοῦναί σοι.
\vs{8}Λέγω ὑμῖν, εἰ καὶ οὐ δώσει αὐτῷ ἀναστὰς διὰ τὸ εἶναι φίλον αὐτοῦ, διά γε τὴν ἀναίδειαν αὐτοῦ ἐγερθεὶς δώσει αὐτῷ ὅσων χρῄζει.
\vs{9}Κἀγὼ ὑμῖν λέγω, αἰτεῖτε καὶ δοθήσεται ὑμῖν, ζητεῖτε καὶ εὑρήσετε, κρούετε καὶ ἀνοιγήσεται ὑμῖν·
\vs{10}πᾶς γὰρ ὁ αἰτῶν λαμβάνει καὶ ὁ ζητῶν εὑρίσκει καὶ τῷ κρούοντι ἀνοιγήσεται.
\vs{11}Τίνα δὲ ἐξ ὑμῶν τὸν πατέρα αἰτήσει ὁ υἱὸς ἰχθύν, καὶ ἀντὶ ἰχθύος ὄφιν αὐτῷ ἐπιδώσει;
\vs{12}ἢ καὶ αἰτήσει ᾠόν, ἐπιδώσει αὐτῷ σκορπίον;
\vs{13}εἰ οὖν ὑμεῖς πονηροὶ ὑπάρχοντες οἴδατε δόματα ἀγαθὰ διδόναι τοῖς τέκνοις ὑμῶν, πόσῳ μᾶλλον ὁ Πατὴρ ὁ ἐξ οὐρανοῦ δώσει Πνεῦμα Ἅγιον τοῖς αἰτοῦσιν αὐτόν.

\vs{14}Καὶ ἦν ἐκβάλλων δαιμόνιον καὶ αὐτὸ ἦν κωφόν· ἐγένετο δὲ τοῦ δαιμονίου ἐξελθόντος ἐλάλησεν ὁ κωφός καὶ ἐθαύμασαν οἱ ὄχλοι.
\vs{15}τινὲς δὲ ἐξ αὐτῶν εἶπον· Ἐν Βεελζεβοὺλ τῷ ἄρχοντι τῶν δαιμονίων ἐκβάλλει τὰ δαιμόνια·
\vs{16}ἕτεροι δὲ πειράζοντες σημεῖον ἐξ οὐρανοῦ ἐζήτουν παρ᾽ αὐτοῦ.
\vs{17}Αὐτὸς δὲ εἰδὼς αὐτῶν τὰ διανοήματα εἶπεν αὐτοῖς· Πᾶσα βασιλεία ἐφ᾽ ἑαυτὴν διαμερισθεῖσα ἐρημοῦται καὶ οἶκος ἐπὶ οἶκον πίπτει.
\vs{18}εἰ δὲ καὶ ὁ Σατανᾶς ἐφ᾽ ἑαυτὸν διεμερίσθη, πῶς σταθήσεται ἡ βασιλεία αὐτοῦ; ὅτι λέγετε ἐν Βεελζεβοὺλ ἐκβάλλειν με τὰ δαιμόνια.
\vs{19}εἰ δὲ ἐγὼ ἐν Βεελζεβοὺλ ἐκβάλλω τὰ δαιμόνια, οἱ υἱοὶ ὑμῶν ἐν τίνι ἐκβάλλουσιν; διὰ τοῦτο αὐτοὶ ὑμῶν κριταὶ ἔσονται.
\vs{20}εἰ δὲ ἐν δακτύλῳ Θεοῦ ἐγὼ ἐκβάλλω τὰ δαιμόνια, ἄρα ἔφθασεν ἐφ᾽ ὑμᾶς ἡ βασιλεία τοῦ Θεοῦ.
\vs{21}Ὅταν ὁ ἰσχυρὸς καθωπλισμένος φυλάσσῃ τὴν ἑαυτοῦ αὐλήν, ἐν εἰρήνῃ ἐστὶν τὰ ὑπάρχοντα αὐτοῦ·
\vs{22}ἐπὰν δὲ ἰσχυρότερος αὐτοῦ ἐπελθὼν νικήσῃ αὐτόν, τὴν πανοπλίαν αὐτοῦ αἴρει ἐφ᾽ ᾗ ἐπεποίθει καὶ τὰ σκῦλα αὐτοῦ διαδίδωσιν.
\vs{23}Ὁ μὴ ὢν μετ᾽ ἐμοῦ κατ᾽ ἐμοῦ ἐστιν, καὶ ὁ μὴ συνάγων μετ᾽ ἐμοῦ σκορπίζει.
\vs{24}Ὅταν τὸ ἀκάθαρτον πνεῦμα ἐξέλθῃ ἀπὸ τοῦ ἀνθρώπου, διέρχεται δι᾽ ἀνύδρων τόπων ζητοῦν ἀνάπαυσιν καὶ μὴ εὑρίσκον· τότε λέγει· Ὑποστρέψω εἰς τὸν οἶκόν μου ὅθεν ἐξῆλθον·
\vs{25}καὶ ἐλθὸν εὑρίσκει σεσαρωμένον καὶ κεκοσμημένον.
\vs{26}τότε πορεύεται καὶ παραλαμβάνει ἕτερα πνεύματα πονηρότερα ἑαυτοῦ ἑπτά καὶ εἰσελθόντα κατοικεῖ ἐκεῖ· καὶ γίνεται τὰ ἔσχατα τοῦ ἀνθρώπου ἐκείνου χείρονα τῶν πρώτων.

\vs{27}Ἐγένετο δὲ ἐν τῷ λέγειν αὐτὸν ταῦτα ἐπάρασά τις φωνὴν γυνὴ ἐκ τοῦ ὄχλου εἶπεν αὐτῷ· Μακαρία ἡ κοιλία ἡ βαστάσασά σε καὶ μαστοὶ οὓς ἐθήλασας.
\vs{28}Αὐτὸς δὲ εἶπεν· Μενοῦν μακάριοι οἱ ἀκούοντες τὸν λόγον τοῦ Θεοῦ καὶ φυλάσσοντες.

\vs{29}Τῶν δὲ ὄχλων ἐπαθροιζομένων ἤρξατο λέγειν· Ἡ γενεὰ αὕτη γενεὰ πονηρά ἐστιν· σημεῖον ζητεῖ, καὶ σημεῖον οὐ δοθήσεται αὐτῇ εἰ μὴ τὸ σημεῖον Ἰωνᾶ.
\vs{30}καθὼς γὰρ ἐγένετο Ἰωνᾶς τοῖς Νινευίταις σημεῖον, οὕτως ἔσται καὶ ὁ Υἱὸς τοῦ ἀνθρώπου τῇ γενεᾷ ταύτῃ.
\vs{31}Βασίλισσα νότου ἐγερθήσεται ἐν τῇ κρίσει μετὰ τῶν ἀνδρῶν τῆς γενεᾶς ταύτης καὶ κατακρινεῖ αὐτούς, ὅτι ἦλθεν ἐκ τῶν περάτων τῆς γῆς ἀκοῦσαι τὴν σοφίαν Σολομῶνος, καὶ ἰδοὺ πλεῖον Σολομῶνος ὧδε.
\vs{32}ἄνδρες Νινευῖται ἀναστήσονται ἐν τῇ κρίσει μετὰ τῆς γενεᾶς ταύτης καὶ κατακρινοῦσιν αὐτήν· ὅτι μετενόησαν εἰς τὸ κήρυγμα Ἰωνᾶ, καὶ ἰδοὺ πλεῖον Ἰωνᾶ ὧδε.

\vs{33}Οὐδεὶς λύχνον ἅψας εἰς κρύπτην τίθησιν οὐδὲ ὑπὸ τὸν μόδιον ἀλλ᾽ ἐπὶ τὴν λυχνίαν, ἵνα οἱ εἰσπορευόμενοι τὸ φῶς βλέπωσιν.

\vs{34}Ὁ λύχνος τοῦ σώματός ἐστιν ὁ ὀφθαλμός σου. ὅταν ὁ ὀφθαλμός σου ἁπλοῦς ᾖ, καὶ ὅλον τὸ σῶμά σου φωτεινόν ἐστιν· ἐπὰν δὲ πονηρὸς ᾖ, καὶ τὸ σῶμά σου σκοτεινόν.
\vs{35}σκόπει οὖν μὴ τὸ φῶς τὸ ἐν σοὶ σκότος ἐστίν.
\vs{36}εἰ οὖν τὸ σῶμά σου ὅλον φωτεινόν, μὴ ἔχον μέρος τι σκοτεινόν, ἔσται φωτεινὸν ὅλον ὡς ὅταν ὁ λύχνος τῇ ἀστραπῇ φωτίζῃ σε.

\vs{37}Ἐν δὲ τῷ λαλῆσαι ἐρωτᾷ αὐτὸν Φαρισαῖος ὅπως ἀριστήσῃ παρ᾽ αὐτῷ· εἰσελθὼν δὲ ἀνέπεσεν.
\vs{38}ὁ δὲ Φαρισαῖος ἰδὼν ἐθαύμασεν ὅτι οὐ πρῶτον ἐβαπτίσθη πρὸ τοῦ ἀρίστου.
\vs{39}Εἶπεν δὲ ὁ Κύριος πρὸς αὐτόν· Νῦν ὑμεῖς οἱ Φαρισαῖοι τὸ ἔξωθεν τοῦ ποτηρίου καὶ τοῦ πίνακος καθαρίζετε, τὸ δὲ ἔσωθεν ὑμῶν γέμει ἁρπαγῆς καὶ πονηρίας.
\vs{40}ἄφρονες, οὐχ ὁ ποιήσας τὸ ἔξωθεν καὶ τὸ ἔσωθεν ἐποίησεν;
\vs{41}πλὴν τὰ ἐνόντα δότε ἐλεημοσύνην, καὶ ἰδοὺ πάντα καθαρὰ ὑμῖν ἐστιν.

\vs{42}Ἀλλὰ οὐαὶ ὑμῖν τοῖς Φαρισαίοις, ὅτι ἀποδεκατοῦτε τὸ ἡδύοσμον καὶ τὸ πήγανον καὶ πᾶν λάχανον καὶ παρέρχεσθε τὴν κρίσιν καὶ τὴν ἀγάπην τοῦ Θεοῦ· ταῦτα δὲ ἔδει ποιῆσαι κἀκεῖνα μὴ παρεῖναι.

\vs{43}Οὐαὶ ὑμῖν τοῖς Φαρισαίοις, ὅτι ἀγαπᾶτε τὴν πρωτοκαθεδρίαν ἐν ταῖς συναγωγαῖς καὶ τοὺς ἀσπασμοὺς ἐν ταῖς ἀγοραῖς.
\vs{44}οὐαὶ ὑμῖν, ὅτι ἐστὲ ὡς τὰ μνημεῖα τὰ ἄδηλα, καὶ οἱ ἄνθρωποι οἱ περιπατοῦντες ἐπάνω οὐκ οἴδασιν.

\vs{45}Ἀποκριθεὶς δέ τις τῶν νομικῶν λέγει αὐτῷ· Διδάσκαλε, ταῦτα λέγων καὶ ἡμᾶς ὑβρίζεις.
\vs{46}Ὁ δὲ εἶπεν· Καὶ ὑμῖν τοῖς νομικοῖς οὐαί, ὅτι φορτίζετε τοὺς ἀνθρώπους φορτία δυσβάστακτα, καὶ αὐτοὶ ἑνὶ τῶν δακτύλων ὑμῶν οὐ προσψαύετε τοῖς φορτίοις.

\vs{47}οὐαὶ ὑμῖν, ὅτι οἰκοδομεῖτε τὰ μνημεῖα τῶν προφητῶν, οἱ δὲ πατέρες ὑμῶν ἀπέκτειναν αὐτούς.
\vs{48}ἄρα μάρτυρές ἐστε καὶ συνευδοκεῖτε τοῖς ἔργοις τῶν πατέρων ὑμῶν, ὅτι αὐτοὶ μὲν ἀπέκτειναν αὐτοὺς, ὑμεῖς δὲ οἰκοδομεῖτε.
\vs{49}διὰ τοῦτο καὶ ἡ σοφία τοῦ Θεοῦ εἶπεν· Ἀποστελῶ εἰς αὐτοὺς προφήτας καὶ ἀποστόλους, καὶ ἐξ αὐτῶν ἀποκτενοῦσιν καὶ διώξουσιν,
\vs{50}ἵνα ἐκζητηθῇ τὸ αἷμα πάντων τῶν προφητῶν τὸ ἐκκεχυμένον ἀπὸ καταβολῆς κόσμου ἀπὸ τῆς γενεᾶς ταύτης,
\vs{51}ἀπὸ αἵματος Ἅβελ ἕως αἵματος Ζαχαρίου τοῦ ἀπολομένου μεταξὺ τοῦ θυσιαστηρίου καὶ τοῦ οἴκου· ναί λέγω ὑμῖν, ἐκζητηθήσεται ἀπὸ τῆς γενεᾶς ταύτης.

\vs{52}Οὐαὶ ὑμῖν τοῖς νομικοῖς, ὅτι ἤρατε τὴν κλεῖδα τῆς γνώσεως· αὐτοὶ οὐκ εἰσήλθατε καὶ τοὺς εἰσερχομένους ἐκωλύσατε.

\vs{53}Κἀκεῖθεν ἐξελθόντος αὐτοῦ ἤρξαντο οἱ γραμματεῖς καὶ οἱ Φαρισαῖοι δεινῶς ἐνέχειν καὶ ἀποστοματίζειν αὐτὸν περὶ πλειόνων,
\vs{54}ἐνεδρεύοντες αὐτὸν θηρεῦσαί τι ἐκ τοῦ στόματος αὐτοῦ.

\ch{12}
Ἐν οἷς ἐπισυναχθεισῶν τῶν μυριάδων τοῦ ὄχλου, ὥστε καταπατεῖν ἀλλήλους, ἤρξατο λέγειν πρὸς τοὺς μαθητὰς αὐτοῦ πρῶτον· Προσέχετε ἑαυτοῖς ἀπὸ τῆς ζύμης, ἥτις ἐστὶν ὑπόκρισις, τῶν Φαρισαίων.

\vs{2}οὐδὲν δὲ συγκεκαλυμμένον ἐστὶν ὃ οὐκ ἀποκαλυφθήσεται καὶ κρυπτὸν ὃ οὐ γνωσθήσεται.
\vs{3}ἀνθ᾽ ὧν ὅσα ἐν τῇ σκοτίᾳ εἴπατε ἐν τῷ φωτὶ ἀκουσθήσεται, καὶ ὃ πρὸς τὸ οὖς ἐλαλήσατε ἐν τοῖς ταμείοις κηρυχθήσεται ἐπὶ τῶν δωμάτων.

\vs{4}Λέγω δὲ ὑμῖν τοῖς φίλοις μου, μὴ φοβηθῆτε ἀπὸ τῶν ἀποκτεινόντων τὸ σῶμα καὶ μετὰ ταῦτα μὴ ἐχόντων περισσότερόν τι ποιῆσαι.
\vs{5}ὑποδείξω δὲ ὑμῖν τίνα φοβηθῆτε· φοβήθητε τὸν μετὰ τὸ ἀποκτεῖναι ἔχοντα ἐξουσίαν ἐμβαλεῖν εἰς τὴν γέενναν. ναί λέγω ὑμῖν, τοῦτον φοβήθητε.
\vs{6}Οὐχὶ πέντε στρουθία πωλοῦνται ἀσσαρίων δύο; καὶ ἓν ἐξ αὐτῶν οὐκ ἔστιν ἐπιλελησμένον ἐνώπιον τοῦ Θεοῦ.
\vs{7}ἀλλὰ καὶ αἱ τρίχες τῆς κεφαλῆς ὑμῶν πᾶσαι ἠρίθμηνται. μὴ φοβεῖσθε· πολλῶν στρουθίων διαφέρετε.
\vs{8}Λέγω δὲ ὑμῖν, πᾶς ὃς ἂν ὁμολογήσῃ ἐν ἐμοὶ ἔμπροσθεν τῶν ἀνθρώπων, καὶ ὁ Υἱὸς τοῦ ἀνθρώπου ὁμολογήσει ἐν αὐτῷ ἔμπροσθεν τῶν ἀγγέλων τοῦ Θεοῦ·
\vs{9}ὁ δὲ ἀρνησάμενός με ἐνώπιον τῶν ἀνθρώπων ἀπαρνηθήσεται ἐνώπιον τῶν ἀγγέλων τοῦ Θεοῦ.

\vs{10}καὶ πᾶς ὃς ἐρεῖ λόγον εἰς τὸν Υἱὸν τοῦ ἀνθρώπου, ἀφεθήσεται αὐτῷ· τῷ δὲ εἰς τὸ Ἅγιον Πνεῦμα βλασφημήσαντι οὐκ ἀφεθήσεται.

\vs{11}Ὅταν δὲ εἰσφέρωσιν ὑμᾶς ἐπὶ τὰς συναγωγὰς καὶ τὰς ἀρχὰς καὶ τὰς ἐξουσίας, μὴ μεριμνήσητε πῶς ἢ τί ἀπολογήσησθε ἢ τί εἴπητε·
\vs{12}τὸ γὰρ Ἅγιον Πνεῦμα διδάξει ὑμᾶς ἐν αὐτῇ τῇ ὥρᾳ ἃ δεῖ εἰπεῖν.

\vs{13}Εἶπεν δέ τις ἐκ τοῦ ὄχλου αὐτῷ· Διδάσκαλε, εἰπὲ τῷ ἀδελφῷ μου μερίσασθαι μετ᾽ ἐμοῦ τὴν κληρονομίαν.
\vs{14}Ὁ δὲ εἶπεν αὐτῷ· Ἄνθρωπε, τίς με κατέστησεν κριτὴν ἢ μεριστὴν ἐφ᾽ ὑμᾶς;
\vs{15}εἶπεν δὲ πρὸς αὐτούς· Ὁρᾶτε καὶ φυλάσσεσθε ἀπὸ πάσης πλεονεξίας, ὅτι οὐκ ἐν τῷ περισσεύειν τινὶ ἡ ζωὴ αὐτοῦ ἐστιν ἐκ τῶν ὑπαρχόντων αὐτῷ.

\vs{16}Εἶπεν δὲ παραβολὴν πρὸς αὐτοὺς λέγων· Ἀνθρώπου τινὸς πλουσίου εὐφόρησεν ἡ χώρα.
\vs{17}καὶ διελογίζετο ἐν ἑαυτῷ λέγων· Τί ποιήσω, ὅτι οὐκ ἔχω ποῦ συνάξω τοὺς καρπούς μου;
\vs{18}καὶ εἶπεν· Τοῦτο ποιήσω, καθελῶ μου τὰς ἀποθήκας καὶ μείζονας οἰκοδομήσω καὶ συνάξω ἐκεῖ πάντα τὸν σῖτον καὶ τὰ ἀγαθά μου
\vs{19}καὶ ἐρῶ τῇ ψυχῇ μου· Ψυχή, ἔχεις πολλὰ ἀγαθὰ κείμενα εἰς ἔτη πολλά· ἀναπαύου, φάγε, πίε, εὐφραίνου.
\vs{20}Εἶπεν δὲ αὐτῷ ὁ Θεός· Ἄφρων, ταύτῃ τῇ νυκτὶ τὴν ψυχήν σου ἀπαιτοῦσιν ἀπὸ σοῦ· ἃ δὲ ἡτοίμασας, τίνι ἔσται;
\vs{21}Οὕτως ὁ θησαυρίζων ἑαυτῷ καὶ μὴ εἰς Θεὸν πλουτῶν.

\vs{22}Εἶπεν δὲ πρὸς τοὺς μαθητὰς αὐτοῦ· Διὰ τοῦτο λέγω ὑμῖν· μὴ μεριμνᾶτε τῇ ψυχῇ τί φάγητε, μηδὲ τῷ σώματι τί ἐνδύσησθε.
\vs{23}ἡ γὰρ ψυχὴ πλεῖόν ἐστιν τῆς τροφῆς καὶ τὸ σῶμα τοῦ ἐνδύματος.
\vs{24}κατανοήσατε τοὺς κόρακας ὅτι οὐ σπείρουσιν οὐδὲ θερίζουσιν, οἷς οὐκ ἔστιν ταμεῖον οὐδὲ ἀποθήκη, καὶ ὁ Θεὸς τρέφει αὐτούς· πόσῳ μᾶλλον ὑμεῖς διαφέρετε τῶν πετεινῶν.
\vs{25}τίς δὲ ἐξ ὑμῶν μεριμνῶν δύναται ἐπὶ τὴν ἡλικίαν αὐτοῦ προσθεῖναι πῆχυν;
\vs{26}εἰ οὖν οὐδὲ ἐλάχιστον δύνασθε, τί περὶ τῶν λοιπῶν μεριμνᾶτε;
\vs{27}Κατανοήσατε τὰ κρίνα πῶς αὐξάνει· οὐ κοπιᾷ οὐδὲ νήθει· λέγω δὲ ὑμῖν, οὐδὲ Σολομὼν ἐν πάσῃ τῇ δόξῃ αὐτοῦ περιεβάλετο ὡς ἓν τούτων.
\vs{28}εἰ δὲ ἐν ἀγρῷ τὸν χόρτον ὄντα σήμερον καὶ αὔριον εἰς κλίβανον βαλλόμενον ὁ Θεὸς οὕτως ἀμφιέζει, πόσῳ μᾶλλον ὑμᾶς, ὀλιγόπιστοι.
\vs{29}Καὶ ὑμεῖς μὴ ζητεῖτε τί φάγητε καὶ τί πίητε καὶ μὴ μετεωρίζεσθε·
\vs{30}ταῦτα γὰρ πάντα τὰ ἔθνη τοῦ κόσμου ἐπιζητοῦσιν, ὑμῶν δὲ ὁ Πατὴρ οἶδεν ὅτι χρῄζετε τούτων.
\vs{31}πλὴν ζητεῖτε τὴν βασιλείαν αὐτοῦ, καὶ ταῦτα προστεθήσεται ὑμῖν.
\vs{32}Μὴ φοβοῦ, τὸ μικρὸν ποίμνιον, ὅτι εὐδόκησεν ὁ Πατὴρ ὑμῶν δοῦναι ὑμῖν τὴν βασιλείαν.

\vs{33}Πωλήσατε τὰ ὑπάρχοντα ὑμῶν καὶ δότε ἐλεημοσύνην· ποιήσατε ἑαυτοῖς βαλλάντια μὴ παλαιούμενα, θησαυρὸν ἀνέκλειπτον ἐν τοῖς οὐρανοῖς, ὅπου κλέπτης οὐκ ἐγγίζει οὐδὲ σὴς διαφθείρει·
\vs{34}ὅπου γάρ ἐστιν ὁ θησαυρὸς ὑμῶν, ἐκεῖ καὶ ἡ καρδία ὑμῶν ἔσται.

\vs{35}Ἔστωσαν ὑμῶν αἱ ὀσφύες περιεζωσμέναι καὶ οἱ λύχνοι καιόμενοι·
\vs{36}καὶ ὑμεῖς ὅμοιοι ἀνθρώποις προσδεχομένοις τὸν κύριον ἑαυτῶν πότε ἀναλύσῃ ἐκ τῶν γάμων, ἵνα ἐλθόντος καὶ κρούσαντος εὐθέως ἀνοίξωσιν αὐτῷ.
\vs{37}μακάριοι οἱ δοῦλοι ἐκεῖνοι, οὓς ἐλθὼν ὁ κύριος εὑρήσει γρηγοροῦντας· ἀμὴν λέγω ὑμῖν ὅτι περιζώσεται καὶ ἀνακλινεῖ αὐτοὺς καὶ παρελθὼν διακονήσει αὐτοῖς.
\vs{38}κἂν ἐν τῇ δευτέρᾳ κἂν ἐν τῇ τρίτῃ φυλακῇ ἔλθῃ καὶ εὕρῃ οὕτως, μακάριοί εἰσιν ἐκεῖνοι.
\vs{39}Τοῦτο δὲ γινώσκετε ὅτι εἰ ᾔδει ὁ οἰκοδεσπότης ποίᾳ ὥρᾳ ὁ κλέπτης ἔρχεται, οὐκ ἂν ἀφῆκεν διορυχθῆναι τὸν οἶκον αὐτοῦ.
\vs{40}καὶ ὑμεῖς γίνεσθε ἕτοιμοι, ὅτι ᾗ ὥρᾳ οὐ δοκεῖτε ὁ Υἱὸς τοῦ ἀνθρώπου ἔρχεται.

\vs{41}Εἶπεν δὲ ὁ Πέτρος· Κύριε, πρὸς ἡμᾶς τὴν παραβολὴν ταύτην λέγεις ἢ καὶ πρὸς πάντας;
\vs{42}Καὶ εἶπεν ὁ Κύριος· Τίς ἄρα ἐστὶν ὁ πιστὸς οἰκονόμος ὁ φρόνιμος, ὃν καταστήσει ὁ κύριος ἐπὶ τῆς θεραπείας αὐτοῦ τοῦ διδόναι ἐν καιρῷ τὸ σιτομέτριον;
\vs{43}μακάριος ὁ δοῦλος ἐκεῖνος, ὃν ἐλθὼν ὁ κύριος αὐτοῦ εὑρήσει ποιοῦντα οὕτως.
\vs{44}ἀληθῶς λέγω ὑμῖν ὅτι ἐπὶ πᾶσιν τοῖς ὑπάρχουσιν αὐτοῦ καταστήσει αὐτόν.
\vs{45}Ἐὰν δὲ εἴπῃ ὁ δοῦλος ἐκεῖνος ἐν τῇ καρδίᾳ αὐτοῦ· Χρονίζει ὁ κύριός μου ἔρχεσθαι, καὶ ἄρξηται τύπτειν τοὺς παῖδας καὶ τὰς παιδίσκας, ἐσθίειν τε καὶ πίνειν καὶ μεθύσκεσθαι,
\vs{46}ἥξει ὁ κύριος τοῦ δούλου ἐκείνου ἐν ἡμέρᾳ ᾗ οὐ προσδοκᾷ καὶ ἐν ὥρᾳ ᾗ οὐ γινώσκει, καὶ διχοτομήσει αὐτὸν καὶ τὸ μέρος αὐτοῦ μετὰ τῶν ἀπίστων θήσει.

\vs{47}Ἐκεῖνος δὲ ὁ δοῦλος ὁ γνοὺς τὸ θέλημα τοῦ κυρίου αὐτοῦ καὶ μὴ ἑτοιμάσας ἢ ποιήσας πρὸς τὸ θέλημα αὐτοῦ δαρήσεται πολλάς·
\vs{48}ὁ δὲ μὴ γνοὺς, ποιήσας δὲ ἄξια πληγῶν δαρήσεται ὀλίγας. παντὶ δὲ ᾧ ἐδόθη πολύ, πολὺ ζητηθήσεται παρ᾽ αὐτοῦ, καὶ ᾧ παρέθεντο πολύ, περισσότερον αἰτήσουσιν αὐτόν.

\vs{49}Πῦρ ἦλθον βαλεῖν ἐπὶ τὴν γῆν, καὶ τί θέλω εἰ ἤδη ἀνήφθη.
\vs{50}βάπτισμα δὲ ἔχω βαπτισθῆναι, καὶ πῶς συνέχομαι ἕως ὅτου τελεσθῇ.
\vs{51}Δοκεῖτε ὅτι εἰρήνην παρεγενόμην δοῦναι ἐν τῇ γῇ; οὐχί, λέγω ὑμῖν, ἀλλ᾽ ἢ διαμερισμόν.
\vs{52}ἔσονται γὰρ ἀπὸ τοῦ νῦν πέντε ἐν ἑνὶ οἴκῳ διαμεμερισμένοι, τρεῖς ἐπὶ δυσὶν καὶ δύο ἐπὶ τρισίν,
\vs{53}διαμερισθήσονται πατὴρ ἐπὶ υἱῷ καὶ υἱὸς ἐπὶ πατρί, μήτηρ ἐπὶ τὴν θυγατέρα καὶ θυγάτηρ ἐπὶ τὴν μητέρα, πενθερὰ ἐπὶ τὴν νύμφην αὐτῆς καὶ νύμφη ἐπὶ τὴν πενθεράν.
\vs{54}Ἔλεγεν δὲ καὶ τοῖς ὄχλοις· Ὅταν ἴδητε τὴν νεφέλην ἀνατέλλουσαν ἐπὶ δυσμῶν, εὐθέως λέγετε ὅτι Ὄμβρος ἔρχεται, καὶ γίνεται οὕτως·
\vs{55}καὶ ὅταν νότον πνέοντα, λέγετε ὅτι Καύσων ἔσται, καὶ γίνεται.
\vs{56}ὑποκριταί, τὸ πρόσωπον τῆς γῆς καὶ τοῦ οὐρανοῦ οἴδατε δοκιμάζειν, τὸν καιρὸν δὲ τοῦτον πῶς οὐκ οἴδατε δοκιμάζειν;

\vs{57}Τί δὲ καὶ ἀφ᾽ ἑαυτῶν οὐ κρίνετε τὸ δίκαιον;
\vs{58}ὡς γὰρ ὑπάγεις μετὰ τοῦ ἀντιδίκου σου ἐπ᾽ ἄρχοντα, ἐν τῇ ὁδῷ δὸς ἐργασίαν ἀπηλλάχθαι ἀπ᾽ αὐτοῦ, μήποτε κατασύρῃ σε πρὸς τὸν κριτήν, καὶ ὁ κριτής σε παραδώσει τῷ πράκτορι, καὶ ὁ πράκτωρ σε βαλεῖ εἰς φυλακήν.
\vs{59}λέγω σοι, οὐ μὴ ἐξέλθῃς ἐκεῖθεν, ἕως καὶ τὸ ἔσχατον λεπτὸν ἀποδῷς.

\ch{13}
Παρῆσαν δέ τινες ἐν αὐτῷ τῷ καιρῷ ἀπαγγέλλοντες αὐτῷ περὶ τῶν Γαλιλαίων ὧν τὸ αἷμα Πιλᾶτος ἔμιξεν μετὰ τῶν θυσιῶν αὐτῶν.
\vs{2}καὶ ἀποκριθεὶς εἶπεν αὐτοῖς· Δοκεῖτε ὅτι οἱ Γαλιλαῖοι οὗτοι ἁμαρτωλοὶ παρὰ πάντας τοὺς Γαλιλαίους ἐγένοντο, ὅτι ταῦτα πεπόνθασιν;
\vs{3}οὐχί, λέγω ὑμῖν, ἀλλ᾽ ἐὰν μὴ μετανοῆτε πάντες ὁμοίως ἀπολεῖσθε.
\vs{4}ἢ ἐκεῖνοι οἱ δεκαοκτὼ ἐφ᾽ οὓς ἔπεσεν ὁ πύργος ἐν τῷ Σιλωὰμ καὶ ἀπέκτεινεν αὐτούς, δοκεῖτε ὅτι αὐτοὶ ὀφειλέται ἐγένοντο παρὰ πάντας τοὺς ἀνθρώπους τοὺς κατοικοῦντας Ἰερουσαλήμ;
\vs{5}οὐχί, λέγω ὑμῖν, ἀλλ᾽ ἐὰν μὴ μετανοῆτε πάντες ὡσαύτως ἀπολεῖσθε.

\vs{6}Ἔλεγεν δὲ ταύτην τὴν παραβολήν· Συκῆν εἶχέν τις πεφυτευμένην ἐν τῷ ἀμπελῶνι αὐτοῦ, καὶ ἦλθεν ζητῶν καρπὸν ἐν αὐτῇ καὶ οὐχ εὗρεν.
\vs{7}εἶπεν δὲ πρὸς τὸν ἀμπελουργόν· Ἰδοὺ τρία ἔτη ἀφ᾽ οὗ ἔρχομαι ζητῶν καρπὸν ἐν τῇ συκῇ ταύτῃ καὶ οὐχ εὑρίσκω· ἔκκοψον οὖν αὐτήν, ἵνατί καὶ τὴν γῆν καταργεῖ;
\vs{8}Ὁ δὲ ἀποκριθεὶς λέγει αὐτῷ· Κύριε, ἄφες αὐτὴν καὶ τοῦτο τὸ ἔτος, ἕως ὅτου σκάψω περὶ αὐτὴν καὶ βάλω κόπρια,
\vs{9}κἂν μὲν ποιήσῃ καρπὸν εἰς τὸ μέλλον· εἰ δὲ μή γε, ἐκκόψεις αὐτήν.

\vs{10}Ἦν δὲ διδάσκων ἐν μιᾷ τῶν συναγωγῶν ἐν τοῖς σάββασιν.
\vs{11}καὶ ἰδοὺ γυνὴ πνεῦμα ἔχουσα ἀσθενείας ἔτη δεκαοκτώ καὶ ἦν συνκύπτουσα καὶ μὴ δυναμένη ἀνακύψαι εἰς τὸ παντελές.
\vs{12}ἰδὼν δὲ αὐτὴν ὁ Ἰησοῦς προσεφώνησεν καὶ εἶπεν αὐτῇ· Γύναι, ἀπολέλυσαι τῆς ἀσθενείας σου,
\vs{13}καὶ ἐπέθηκεν αὐτῇ τὰς χεῖρας· καὶ παραχρῆμα ἀνωρθώθη καὶ ἐδόξαζεν τὸν Θεόν.
\vs{14}Ἀποκριθεὶς δὲ ὁ ἀρχισυνάγωγος, ἀγανακτῶν ὅτι τῷ σαββάτῳ ἐθεράπευσεν ὁ Ἰησοῦς, ἔλεγεν τῷ ὄχλῳ ὅτι Ἓξ ἡμέραι εἰσὶν ἐν αἷς δεῖ ἐργάζεσθαι· ἐν αὐταῖς οὖν ἐρχόμενοι θεραπεύεσθε καὶ μὴ τῇ ἡμέρᾳ τοῦ σαββάτου.
\vs{15}Ἀπεκρίθη δὲ αὐτῷ ὁ Κύριος καὶ εἶπεν· Ὑποκριταί, ἕκαστος ὑμῶν τῷ σαββάτῳ οὐ λύει τὸν βοῦν αὐτοῦ ἢ τὸν ὄνον ἀπὸ τῆς φάτνης καὶ ἀπαγαγὼν ποτίζει;
\vs{16}ταύτην δὲ θυγατέρα Ἀβραὰμ οὖσαν, ἣν ἔδησεν ὁ Σατανᾶς ἰδοὺ δέκα καὶ ὀκτὼ ἔτη, οὐκ ἔδει λυθῆναι ἀπὸ τοῦ δεσμοῦ τούτου τῇ ἡμέρᾳ τοῦ σαββάτου;
\vs{17}Καὶ ταῦτα λέγοντος αὐτοῦ κατῃσχύνοντο πάντες οἱ ἀντικείμενοι αὐτῷ, καὶ πᾶς ὁ ὄχλος ἔχαιρεν ἐπὶ πᾶσιν τοῖς ἐνδόξοις τοῖς γινομένοις ὑπ᾽ αὐτοῦ.

\vs{18}Ἔλεγεν οὖν· Τίνι ὁμοία ἐστὶν ἡ βασιλεία τοῦ Θεοῦ καὶ τίνι ὁμοιώσω αὐτήν;
\vs{19}ὁμοία ἐστὶν κόκκῳ σινάπεως, ὃν λαβὼν ἄνθρωπος ἔβαλεν εἰς κῆπον ἑαυτοῦ, καὶ ηὔξησεν καὶ ἐγένετο εἰς δένδρον, καὶ τὰ πετεινὰ τοῦ οὐρανοῦ κατεσκήνωσεν ἐν τοῖς κλάδοις αὐτοῦ.

\vs{20}Καὶ πάλιν εἶπεν· Τίνι ὁμοιώσω τὴν βασιλείαν τοῦ Θεοῦ;
\vs{21}ὁμοία ἐστὶν ζύμῃ, ἣν λαβοῦσα γυνὴ ἐνέκρυψεν εἰς ἀλεύρου σάτα τρία ἕως οὗ ἐζυμώθη ὅλον.

\vs{22}Καὶ διεπορεύετο κατὰ πόλεις καὶ κώμας διδάσκων καὶ πορείαν ποιούμενος εἰς Ἱεροσόλυμα.

\vs{23}Εἶπεν δέ τις αὐτῷ· Κύριε, εἰ ὀλίγοι οἱ σῳζόμενοι; Ὁ δὲ εἶπεν πρὸς αὐτούς·
\vs{24}Ἀγωνίζεσθε εἰσελθεῖν διὰ τῆς στενῆς θύρας, ὅτι πολλοί, λέγω ὑμῖν, ζητήσουσιν εἰσελθεῖν καὶ οὐκ ἰσχύσουσιν.
\vs{25}ἀφ᾽ οὗ ἂν ἐγερθῇ ὁ οἰκοδεσπότης καὶ ἀποκλείσῃ τὴν θύραν καὶ ἄρξησθε ἔξω ἑστάναι καὶ κρούειν τὴν θύραν λέγοντες· Κύριε, ἄνοιξον ἡμῖν, Καὶ ἀποκριθεὶς ἐρεῖ ὑμῖν· Οὐκ οἶδα ὑμᾶς πόθεν ἐστέ.
\vs{26}Τότε ἄρξεσθε λέγειν· Ἐφάγομεν ἐνώπιόν σου καὶ ἐπίομεν καὶ ἐν ταῖς πλατείαις ἡμῶν ἐδίδαξας·
\vs{27}Καὶ ἐρεῖ Λέγων ὑμῖν· Οὐκ οἶδα ὑμᾶς πόθεν ἐστέ· ἀπόστητε ἀπ᾽ ἐμοῦ πάντες ἐργάται ἀδικίας.
\vs{28}Ἐκεῖ ἔσται ὁ κλαυθμὸς καὶ ὁ βρυγμὸς τῶν ὀδόντων, ὅταν ὄψησθε Ἀβραὰμ καὶ Ἰσαὰκ καὶ Ἰακὼβ καὶ πάντας τοὺς προφήτας ἐν τῇ βασιλείᾳ τοῦ Θεοῦ, ὑμᾶς δὲ ἐκβαλλομένους ἔξω.
\vs{29}καὶ ἥξουσιν ἀπὸ ἀνατολῶν καὶ δυσμῶν καὶ ἀπὸ βορρᾶ καὶ νότου καὶ ἀνακλιθήσονται ἐν τῇ βασιλείᾳ τοῦ Θεοῦ.
\vs{30}καὶ ἰδοὺ εἰσὶν ἔσχατοι οἳ ἔσονται πρῶτοι καὶ εἰσὶν πρῶτοι οἳ ἔσονται ἔσχατοι.

\vs{31}Ἐν αὐτῇ τῇ ὥρᾳ προσῆλθάν τινες Φαρισαῖοι λέγοντες αὐτῷ· Ἔξελθε καὶ πορεύου ἐντεῦθεν, ὅτι Ἡρῴδης θέλει σε ἀποκτεῖναι.
\vs{32}Καὶ εἶπεν αὐτοῖς· Πορευθέντες εἴπατε τῇ ἀλώπεκι ταύτῃ· Ἰδοὺ ἐκβάλλω δαιμόνια καὶ ἰάσεις ἀποτελῶ σήμερον καὶ αὔριον καὶ τῇ τρίτῃ τελειοῦμαι.
\vs{33}πλὴν δεῖ με σήμερον καὶ αὔριον καὶ τῇ ἐχομένῃ πορεύεσθαι, ὅτι οὐκ ἐνδέχεται προφήτην ἀπολέσθαι ἔξω Ἰερουσαλήμ.

\vs{34}Ἰερουσαλὴμ Ἰερουσαλήμ, ἡ ἀποκτείνουσα τοὺς προφήτας καὶ λιθοβολοῦσα τοὺς ἀπεσταλμένους πρὸς αὐτήν, ποσάκις ἠθέλησα ἐπισυνάξαι τὰ τέκνα σου ὃν τρόπον ὄρνις τὴν ἑαυτῆς νοσσιὰν ὑπὸ τὰς πτέρυγας, καὶ οὐκ ἠθελήσατε.
\vs{35}ἰδοὺ ἀφίεται ὑμῖν ὁ οἶκος ὑμῶν. λέγω δὲ ὑμῖν, οὐ μὴ ἴδητέ με ἕως ἥξει ὅτε εἴπητε· 
\begin{poetryblock}

\begin{quote}Εὐλογημένος ὁ ἐρχόμενος ἐν ὀνόματι Κυρίου.\end{quote}
\end{poetryblock}

\ch{14}
Καὶ ἐγένετο ἐν τῷ ἐλθεῖν αὐτὸν εἰς οἶκόν τινος τῶν ἀρχόντων τῶν Φαρισαίων σαββάτῳ φαγεῖν ἄρτον καὶ αὐτοὶ ἦσαν παρατηρούμενοι αὐτόν.

\vs{2}καὶ ἰδοὺ ἄνθρωπός τις ἦν ὑδρωπικὸς ἔμπροσθεν αὐτοῦ.
\vs{3}καὶ ἀποκριθεὶς ὁ Ἰησοῦς εἶπεν πρὸς τοὺς νομικοὺς καὶ Φαρισαίους λέγων· Ἔξεστιν τῷ σαββάτῳ θεραπεῦσαι ἢ οὔ;
\vs{4}Οἱ δὲ ἡσύχασαν. Καὶ ἐπιλαβόμενος ἰάσατο αὐτὸν καὶ ἀπέλυσεν.
\vs{5}καὶ πρὸς αὐτοὺς εἶπεν· Τίνος ὑμῶν υἱὸς ἢ βοῦς εἰς φρέαρ πεσεῖται, καὶ οὐκ εὐθέως ἀνασπάσει αὐτὸν ἐν ἡμέρᾳ τοῦ σαββάτου;
\vs{6}Καὶ οὐκ ἴσχυσαν ἀνταποκριθῆναι πρὸς ταῦτα.

\vs{7}Ἔλεγεν δὲ πρὸς τοὺς κεκλημένους παραβολήν, ἐπέχων πῶς τὰς πρωτοκλισίας ἐξελέγοντο, λέγων πρὸς αὐτούς·
\vs{8}Ὅταν κληθῇς ὑπό τινος εἰς γάμους, μὴ κατακλιθῇς εἰς τὴν πρωτοκλισίαν, μήποτε ἐντιμότερός σου ᾖ κεκλημένος ὑπ᾽ αὐτοῦ,
\vs{9}καὶ ἐλθὼν ὁ σὲ καὶ αὐτὸν καλέσας ἐρεῖ σοι· Δὸς τούτῳ τόπον, καὶ τότε ἄρξῃ μετὰ αἰσχύνης τὸν ἔσχατον τόπον κατέχειν.
\vs{10}Ἀλλ᾽ ὅταν κληθῇς, πορευθεὶς ἀνάπεσε εἰς τὸν ἔσχατον τόπον, ἵνα ὅταν ἔλθῃ ὁ κεκληκώς σε ἐρεῖ σοι· Φίλε, προσανάβηθι ἀνώτερον· τότε ἔσται σοι δόξα ἐνώπιον πάντων τῶν συνανακειμένων σοι.
\vs{11}ὅτι πᾶς ὁ ὑψῶν ἑαυτὸν ταπεινωθήσεται, καὶ ὁ ταπεινῶν ἑαυτὸν ὑψωθήσεται.

\vs{12}Ἔλεγεν δὲ καὶ τῷ κεκληκότι αὐτόν· Ὅταν ποιῇς ἄριστον ἢ δεῖπνον, μὴ φώνει τοὺς φίλους σου μηδὲ τοὺς ἀδελφούς σου μηδὲ τοὺς συγγενεῖς σου μηδὲ γείτονας πλουσίους, μήποτε καὶ αὐτοὶ ἀντικαλέσωσίν σε καὶ γένηται ἀνταπόδομά σοι.
\vs{13}ἀλλ᾽ ὅταν δοχὴν ποιῇς, κάλει πτωχούς, ἀναπείρους, χωλούς, τυφλούς·
\vs{14}καὶ μακάριος ἔσῃ, ὅτι οὐκ ἔχουσιν ἀνταποδοῦναί σοι, ἀνταποδοθήσεται γάρ σοι ἐν τῇ ἀναστάσει τῶν δικαίων.

\vs{15}Ἀκούσας δέ τις τῶν συνανακειμένων ταῦτα εἶπεν αὐτῷ· Μακάριος ὅστις φάγεται ἄρτον ἐν τῇ βασιλείᾳ τοῦ Θεοῦ.

\vs{16}Ὁ δὲ εἶπεν αὐτῷ· Ἄνθρωπός τις ἐποίει δεῖπνον μέγα, καὶ ἐκάλεσεν πολλούς
\vs{17}καὶ ἀπέστειλεν τὸν δοῦλον αὐτοῦ τῇ ὥρᾳ τοῦ δείπνου εἰπεῖν τοῖς κεκλημένοις· Ἔρχεσθε, ὅτι ἤδη ἕτοιμά ἐστιν.
\vs{18}Καὶ ἤρξαντο ἀπὸ μιᾶς πάντες παραιτεῖσθαι. ὁ πρῶτος εἶπεν αὐτῷ· Ἀγρὸν ἠγόρασα καὶ ἔχω ἀνάγκην ἐξελθὼν ἰδεῖν αὐτόν· ἐρωτῶ σε, ἔχε με παρῃτημένον.
\vs{19}Καὶ ἕτερος εἶπεν· Ζεύγη βοῶν ἠγόρασα πέντε καὶ πορεύομαι δοκιμάσαι αὐτά· ἐρωτῶ σε, ἔχε με παρῃτημένον.
\vs{20}Καὶ ἕτερος εἶπεν· Γυναῖκα ἔγημα καὶ διὰ τοῦτο οὐ δύναμαι ἐλθεῖν.
\vs{21}Καὶ παραγενόμενος ὁ δοῦλος ἀπήγγειλεν τῷ κυρίῳ αὐτοῦ ταῦτα. τότε ὀργισθεὶς ὁ οἰκοδεσπότης εἶπεν τῷ δούλῳ αὐτοῦ· Ἔξελθε ταχέως εἰς τὰς πλατείας καὶ ῥύμας τῆς πόλεως καὶ τοὺς πτωχοὺς καὶ ἀναπείρους καὶ τυφλοὺς καὶ χωλοὺς εἰσάγαγε ὧδε.
\vs{22}Καὶ εἶπεν ὁ δοῦλος· Κύριε, γέγονεν ὃ ἐπέταξας, καὶ ἔτι τόπος ἐστίν.
\vs{23}Καὶ εἶπεν ὁ κύριος πρὸς τὸν δοῦλον· Ἔξελθε εἰς τὰς ὁδοὺς καὶ φραγμοὺς καὶ ἀνάγκασον εἰσελθεῖν, ἵνα γεμισθῇ μου ὁ οἶκος·
\vs{24}λέγω γὰρ ὑμῖν ὅτι οὐδεὶς τῶν ἀνδρῶν ἐκείνων τῶν κεκλημένων γεύσεταί μου τοῦ δείπνου.

\vs{25}Συνεπορεύοντο δὲ αὐτῷ ὄχλοι πολλοί, καὶ στραφεὶς εἶπεν πρὸς αὐτούς·
\vs{26}Εἴ τις ἔρχεται πρός με καὶ οὐ μισεῖ τὸν πατέρα ἑαυτοῦ καὶ τὴν μητέρα καὶ τὴν γυναῖκα καὶ τὰ τέκνα καὶ τοὺς ἀδελφοὺς καὶ τὰς ἀδελφάς ἔτι τε καὶ τὴν ψυχὴν ἑαυτοῦ, οὐ δύναται εἶναί μου μαθητής.
\vs{27}ὅστις οὐ βαστάζει τὸν σταυρὸν ἑαυτοῦ καὶ ἔρχεται ὀπίσω μου, οὐ δύναται εἶναί μου μαθητής.

\vs{28}Τίς γὰρ ἐξ ὑμῶν θέλων πύργον οἰκοδομῆσαι οὐχὶ πρῶτον καθίσας ψηφίζει τὴν δαπάνην, εἰ ἔχει εἰς ἀπαρτισμόν;
\vs{29}ἵνα μήποτε θέντος αὐτοῦ θεμέλιον καὶ μὴ ἰσχύοντος ἐκτελέσαι πάντες οἱ θεωροῦντες ἄρξωνται αὐτῷ ἐμπαίζειν
\vs{30}λέγοντες ὅτι Οὗτος ὁ ἄνθρωπος ἤρξατο οἰκοδομεῖν καὶ οὐκ ἴσχυσεν ἐκτελέσαι.
\vs{31}Ἢ τίς βασιλεὺς πορευόμενος ἑτέρῳ βασιλεῖ συμβαλεῖν εἰς πόλεμον οὐχὶ καθίσας πρῶτον βουλεύσεται εἰ δυνατός ἐστιν ἐν δέκα χιλιάσιν ὑπαντῆσαι τῷ μετὰ εἴκοσι χιλιάδων ἐρχομένῳ ἐπ᾽ αὐτόν;
\vs{32}εἰ δὲ μή γε, ἔτι αὐτοῦ πόρρω ὄντος πρεσβείαν ἀποστείλας ἐρωτᾷ τὰ πρὸς εἰρήνην.
\vs{33}Οὕτως οὖν πᾶς ἐξ ὑμῶν ὃς οὐκ ἀποτάσσεται πᾶσιν τοῖς ἑαυτοῦ ὑπάρχουσιν οὐ δύναται εἶναί μου μαθητής.

\vs{34}Καλὸν οὖν τὸ ἅλας· ἐὰν δὲ καὶ τὸ ἅλας μωρανθῇ, ἐν τίνι ἀρτυθήσεται;
\vs{35}οὔτε εἰς γῆν οὔτε εἰς κοπρίαν εὔθετόν ἐστιν, ἔξω βάλλουσιν αὐτό. Ὁ ἔχων ὦτα ἀκούειν ἀκουέτω.

\ch{15}
Ἦσαν δὲ αὐτῷ ἐγγίζοντες πάντες οἱ τελῶναι καὶ οἱ ἁμαρτωλοὶ ἀκούειν αὐτοῦ.
\vs{2}καὶ διεγόγγυζον οἵ τε Φαρισαῖοι καὶ οἱ γραμματεῖς λέγοντες ὅτι Οὗτος ἁμαρτωλοὺς προσδέχεται καὶ συνεσθίει αὐτοῖς.

\vs{3}Εἶπεν δὲ πρὸς αὐτοὺς τὴν παραβολὴν ταύτην λέγων·
\vs{4}Τίς ἄνθρωπος ἐξ ὑμῶν ἔχων ἑκατὸν πρόβατα καὶ ἀπολέσας ἐξ αὐτῶν ἓν οὐ καταλείπει τὰ ἐνενήκοντα ἐννέα ἐν τῇ ἐρήμῳ καὶ πορεύεται ἐπὶ τὸ ἀπολωλὸς ἕως εὕρῃ αὐτό;
\vs{5}καὶ εὑρὼν ἐπιτίθησιν ἐπὶ τοὺς ὤμους αὐτοῦ χαίρων
\vs{6}καὶ ἐλθὼν εἰς τὸν οἶκον συνκαλεῖ τοὺς φίλους καὶ τοὺς γείτονας λέγων αὐτοῖς· Συνχάρητέ μοι, ὅτι εὗρον τὸ πρόβατόν μου τὸ ἀπολωλός.
\vs{7}λέγω ὑμῖν ὅτι οὕτως χαρὰ ἐν τῷ οὐρανῷ ἔσται ἐπὶ ἑνὶ ἁμαρτωλῷ μετανοοῦντι ἢ ἐπὶ ἐνενήκοντα ἐννέα δικαίοις οἵτινες οὐ χρείαν ἔχουσιν μετανοίας.

\vs{8}Ἢ τίς γυνὴ δραχμὰς ἔχουσα δέκα ἐὰν ἀπολέσῃ δραχμὴν μίαν, οὐχὶ ἅπτει λύχνον καὶ σαροῖ τὴν οἰκίαν καὶ ζητεῖ ἐπιμελῶς ἕως οὗ εὕρῃ;
\vs{9}καὶ εὑροῦσα συνκαλεῖ τὰς φίλας καὶ γείτονας λέγουσα· Συνχάρητέ μοι, ὅτι εὗρον τὴν δραχμὴν ἣν ἀπώλεσα.
\vs{10}οὕτως, λέγω ὑμῖν, γίνεται χαρὰ ἐνώπιον τῶν ἀγγέλων τοῦ Θεοῦ ἐπὶ ἑνὶ ἁμαρτωλῷ μετανοοῦντι.

\vs{11}Εἶπεν δέ· Ἄνθρωπός τις εἶχεν δύο υἱούς.
\vs{12}καὶ εἶπεν ὁ νεώτερος αὐτῶν τῷ πατρί· Πάτερ, δός μοι τὸ ἐπιβάλλον μέρος τῆς οὐσίας. ὁ δὲ διεῖλεν αὐτοῖς τὸν βίον.
\vs{13}Καὶ μετ᾽ οὐ πολλὰς ἡμέρας συναγαγὼν πάντα ὁ νεώτερος υἱὸς ἀπεδήμησεν εἰς χώραν μακράν καὶ ἐκεῖ διεσκόρπισεν τὴν οὐσίαν αὐτοῦ ζῶν ἀσώτως.
\vs{14}Δαπανήσαντος δὲ αὐτοῦ πάντα ἐγένετο λιμὸς ἰσχυρὰ κατὰ τὴν χώραν ἐκείνην, καὶ αὐτὸς ἤρξατο ὑστερεῖσθαι.
\vs{15}καὶ πορευθεὶς ἐκολλήθη ἑνὶ τῶν πολιτῶν τῆς χώρας ἐκείνης, καὶ ἔπεμψεν αὐτὸν εἰς τοὺς ἀγροὺς αὐτοῦ βόσκειν χοίρους,
\vs{16}καὶ ἐπεθύμει χορτασθῆναι ἐκ τῶν κερατίων ὧν ἤσθιον οἱ χοῖροι, καὶ οὐδεὶς ἐδίδου αὐτῷ.
\vs{17}Εἰς ἑαυτὸν δὲ ἐλθὼν ἔφη· Πόσοι μίσθιοι τοῦ πατρός μου περισσεύονται ἄρτων, ἐγὼ δὲ λιμῷ ὧδε ἀπόλλυμαι.
\vs{18}ἀναστὰς πορεύσομαι πρὸς τὸν πατέρα μου καὶ ἐρῶ αὐτῷ· Πάτερ, ἥμαρτον εἰς τὸν οὐρανὸν καὶ ἐνώπιόν σου,
\vs{19}οὐκέτι εἰμὶ ἄξιος κληθῆναι υἱός σου· ποίησόν με ὡς ἕνα τῶν μισθίων σου.
\vs{20}Καὶ ἀναστὰς ἦλθεν πρὸς τὸν πατέρα ἑαυτοῦ. ἔτι δὲ αὐτοῦ μακρὰν ἀπέχοντος εἶδεν αὐτὸν ὁ πατὴρ αὐτοῦ καὶ ἐσπλαγχνίσθη καὶ δραμὼν ἐπέπεσεν ἐπὶ τὸν τράχηλον αὐτοῦ καὶ κατεφίλησεν αὐτόν.
\vs{21}Εἶπεν δὲ ὁ υἱὸς αὐτῷ· Πάτερ, ἥμαρτον εἰς τὸν οὐρανὸν καὶ ἐνώπιόν σου, οὐκέτι εἰμὶ ἄξιος κληθῆναι υἱός σου.
\vs{22}Εἶπεν δὲ ὁ πατὴρ πρὸς τοὺς δούλους αὐτοῦ· Ταχὺ ἐξενέγκατε στολὴν τὴν πρώτην καὶ ἐνδύσατε αὐτόν, καὶ δότε δακτύλιον εἰς τὴν χεῖρα αὐτοῦ καὶ ὑποδήματα εἰς τοὺς πόδας,
\vs{23}καὶ φέρετε τὸν μόσχον τὸν σιτευτόν, θύσατε, καὶ φαγόντες εὐφρανθῶμεν,
\vs{24}ὅτι οὗτος ὁ υἱός μου νεκρὸς ἦν καὶ ἀνέζησεν, ἦν ἀπολωλὼς καὶ εὑρέθη. καὶ ἤρξαντο εὐφραίνεσθαι.
\vs{25}Ἦν δὲ ὁ υἱὸς αὐτοῦ ὁ πρεσβύτερος ἐν ἀγρῷ· καὶ ὡς ἐρχόμενος ἤγγισεν τῇ οἰκίᾳ, ἤκουσεν συμφωνίας καὶ χορῶν,
\vs{26}καὶ προσκαλεσάμενος ἕνα τῶν παίδων ἐπυνθάνετο τί ἂν εἴη ταῦτα.
\vs{27}Ὁ δὲ εἶπεν αὐτῷ ὅτι Ὁ ἀδελφός σου ἥκει, καὶ ἔθυσεν ὁ πατήρ σου τὸν μόσχον τὸν σιτευτόν, ὅτι ὑγιαίνοντα αὐτὸν ἀπέλαβεν.
\vs{28}Ὠργίσθη δὲ καὶ οὐκ ἤθελεν εἰσελθεῖν, ὁ δὲ πατὴρ αὐτοῦ ἐξελθὼν παρεκάλει αὐτόν.
\vs{29}Ὁ δὲ ἀποκριθεὶς εἶπεν τῷ πατρὶ αὐτοῦ· Ἰδοὺ τοσαῦτα ἔτη δουλεύω σοι καὶ οὐδέποτε ἐντολήν σου παρῆλθον, καὶ ἐμοὶ οὐδέποτε ἔδωκας ἔριφον ἵνα μετὰ τῶν φίλων μου εὐφρανθῶ·
\vs{30}ὅτε δὲ ὁ υἱός σου οὗτος ὁ καταφαγών σου τὸν βίον μετὰ πορνῶν ἦλθεν, ἔθυσας αὐτῷ τὸν σιτευτὸν μόσχον.
\vs{31}Ὁ δὲ εἶπεν αὐτῷ· Τέκνον, σὺ πάντοτε μετ᾽ ἐμοῦ εἶ, καὶ πάντα τὰ ἐμὰ σά ἐστιν·
\vs{32}εὐφρανθῆναι δὲ καὶ χαρῆναι ἔδει, ὅτι ὁ ἀδελφός σου οὗτος νεκρὸς ἦν καὶ ἔζησεν, καὶ ἀπολωλὼς καὶ εὑρέθη.

\ch{16}
Ἔλεγεν δὲ καὶ πρὸς τοὺς μαθητάς· Ἄνθρωπός τις ἦν πλούσιος ὃς εἶχεν οἰκονόμον, καὶ οὗτος διεβλήθη αὐτῷ ὡς διασκορπίζων τὰ ὑπάρχοντα αὐτοῦ.
\vs{2}καὶ φωνήσας αὐτὸν εἶπεν αὐτῷ· Τί τοῦτο ἀκούω περὶ σοῦ; ἀπόδος τὸν λόγον τῆς οἰκονομίας σου, οὐ γὰρ δύνῃ ἔτι οἰκονομεῖν.
\vs{3}Εἶπεν δὲ ἐν ἑαυτῷ ὁ οἰκονόμος· Τί ποιήσω, ὅτι ὁ κύριός μου ἀφαιρεῖται τὴν οἰκονομίαν ἀπ᾽ ἐμοῦ; σκάπτειν οὐκ ἰσχύω, ἐπαιτεῖν αἰσχύνομαι.
\vs{4}ἔγνων τί ποιήσω, ἵνα ὅταν μετασταθῶ ἐκ τῆς οἰκονομίας δέξωνταί με εἰς τοὺς οἴκους αὐτῶν.
\vs{5}Καὶ προσκαλεσάμενος ἕνα ἕκαστον τῶν χρεοφειλετῶν τοῦ κυρίου ἑαυτοῦ ἔλεγεν τῷ πρώτῳ· Πόσον ὀφείλεις τῷ κυρίῳ μου;
\vs{6}Ὁ δὲ εἶπεν· Ἑκατὸν βάτους ἐλαίου. Ὁ δὲ εἶπεν αὐτῷ· Δέξαι σου τὰ γράμματα καὶ καθίσας ταχέως γράψον πεντήκοντα.
\vs{7}Ἔπειτα ἑτέρῳ εἶπεν· Σὺ δὲ πόσον ὀφείλεις; Ὁ δὲ εἶπεν· Ἑκατὸν κόρους σίτου. Λέγει αὐτῷ· Δέξαι σου τὰ γράμματα καὶ γράψον ὀγδοήκοντα.
\vs{8}Καὶ ἐπῄνεσεν ὁ κύριος τὸν οἰκονόμον τῆς ἀδικίας ὅτι φρονίμως ἐποίησεν· ὅτι οἱ υἱοὶ τοῦ αἰῶνος τούτου φρονιμώτεροι ὑπὲρ τοὺς υἱοὺς τοῦ φωτὸς εἰς τὴν γενεὰν τὴν ἑαυτῶν εἰσιν.
\vs{9}Καὶ ἐγὼ ὑμῖν λέγω, ἑαυτοῖς ποιήσατε φίλους ἐκ τοῦ μαμωνᾶ τῆς ἀδικίας, ἵνα ὅταν ἐκλίπῃ δέξωνται ὑμᾶς εἰς τὰς αἰωνίους σκηνάς.

\vs{10}Ὁ πιστὸς ἐν ἐλαχίστῳ καὶ ἐν πολλῷ πιστός ἐστιν, καὶ ὁ ἐν ἐλαχίστῳ ἄδικος καὶ ἐν πολλῷ ἄδικός ἐστιν.
\vs{11}εἰ οὖν ἐν τῷ ἀδίκῳ μαμωνᾷ πιστοὶ οὐκ ἐγένεσθε, τὸ ἀληθινὸν τίς ὑμῖν πιστεύσει;
\vs{12}καὶ εἰ ἐν τῷ ἀλλοτρίῳ πιστοὶ οὐκ ἐγένεσθε, τὸ ὑμέτερον τίς ὑμῖν δώσει;

\vs{13}Οὐδεὶς οἰκέτης δύναται δυσὶ κυρίοις δουλεύειν· ἢ γὰρ τὸν ἕνα μισήσει καὶ τὸν ἕτερον ἀγαπήσει, ἢ ἑνὸς ἀνθέξεται καὶ τοῦ ἑτέρου καταφρονήσει. οὐ δύνασθε Θεῷ δουλεύειν καὶ μαμωνᾷ.

\vs{14}Ἤκουον δὲ ταῦτα πάντα οἱ Φαρισαῖοι φιλάργυροι ὑπάρχοντες καὶ ἐξεμυκτήριζον αὐτόν.
\vs{15}καὶ εἶπεν αὐτοῖς· Ὑμεῖς ἐστε οἱ δικαιοῦντες ἑαυτοὺς ἐνώπιον τῶν ἀνθρώπων, ὁ δὲ Θεὸς γινώσκει τὰς καρδίας ὑμῶν· ὅτι τὸ ἐν ἀνθρώποις ὑψηλὸν βδέλυγμα ἐνώπιον τοῦ Θεοῦ.

\vs{16}Ὁ νόμος καὶ οἱ προφῆται μέχρι Ἰωάννου· ἀπὸ τότε ἡ βασιλεία τοῦ Θεοῦ εὐαγγελίζεται καὶ πᾶς εἰς αὐτὴν βιάζεται.
\vs{17}εὐκοπώτερον δέ ἐστιν τὸν οὐρανὸν καὶ τὴν γῆν παρελθεῖν ἢ τοῦ νόμου μίαν κεραίαν πεσεῖν.

\vs{18}Πᾶς ὁ ἀπολύων τὴν γυναῖκα αὐτοῦ καὶ γαμῶν ἑτέραν μοιχεύει, καὶ ὁ ἀπολελυμένην ἀπὸ ἀνδρὸς γαμῶν μοιχεύει.

\vs{19}Ἄνθρωπος δέ τις ἦν πλούσιος, καὶ ἐνεδιδύσκετο πορφύραν καὶ βύσσον εὐφραινόμενος καθ᾽ ἡμέραν λαμπρῶς.
\vs{20}πτωχὸς δέ τις ὀνόματι Λάζαρος ἐβέβλητο πρὸς τὸν πυλῶνα αὐτοῦ εἱλκωμένος
\vs{21}καὶ ἐπιθυμῶν χορτασθῆναι ἀπὸ τῶν πιπτόντων ἀπὸ τῆς τραπέζης τοῦ πλουσίου· ἀλλὰ καὶ οἱ κύνες ἐρχόμενοι ἐπέλειχον τὰ ἕλκη αὐτοῦ.
\vs{22}Ἐγένετο δὲ ἀποθανεῖν τὸν πτωχὸν καὶ ἀπενεχθῆναι αὐτὸν ὑπὸ τῶν ἀγγέλων εἰς τὸν κόλπον Ἀβραάμ· ἀπέθανεν δὲ καὶ ὁ πλούσιος καὶ ἐτάφη.
\vs{23}καὶ ἐν τῷ ᾅδῃ ἐπάρας τοὺς ὀφθαλμοὺς αὐτοῦ, ὑπάρχων ἐν βασάνοις, ὁρᾷ Ἀβραὰμ ἀπὸ μακρόθεν καὶ Λάζαρον ἐν τοῖς κόλποις αὐτοῦ.
\vs{24}Καὶ αὐτὸς φωνήσας εἶπεν· Πάτερ Ἀβραάμ, ἐλέησόν με καὶ πέμψον Λάζαρον ἵνα βάψῃ τὸ ἄκρον τοῦ δακτύλου αὐτοῦ ὕδατος καὶ καταψύξῃ τὴν γλῶσσάν μου, ὅτι ὀδυνῶμαι ἐν τῇ φλογὶ ταύτῃ.
\vs{25}Εἶπεν δὲ Ἀβραάμ· Τέκνον, μνήσθητι ὅτι ἀπέλαβες τὰ ἀγαθά σου ἐν τῇ ζωῇ σου, καὶ Λάζαρος ὁμοίως τὰ κακά· νῦν δὲ ὧδε παρακαλεῖται, σὺ δὲ ὀδυνᾶσαι.
\vs{26}καὶ ἐν πᾶσι τούτοις μεταξὺ ἡμῶν καὶ ὑμῶν χάσμα μέγα ἐστήρικται, ὅπως οἱ θέλοντες διαβῆναι ἔνθεν πρὸς ὑμᾶς μὴ δύνωνται, μηδὲ ἐκεῖθεν πρὸς ἡμᾶς διαπερῶσιν.
\vs{27}Εἶπεν δέ· Ἐρωτῶ σε οὖν, πάτερ, ἵνα πέμψῃς αὐτὸν εἰς τὸν οἶκον τοῦ πατρός μου,
\vs{28}ἔχω γὰρ πέντε ἀδελφούς, ὅπως διαμαρτύρηται αὐτοῖς, ἵνα μὴ καὶ αὐτοὶ ἔλθωσιν εἰς τὸν τόπον τοῦτον τῆς βασάνου.
\vs{29}Λέγει δὲ Ἀβραάμ· Ἔχουσι Μωϋσέα καὶ τοὺς προφήτας· ἀκουσάτωσαν αὐτῶν.
\vs{30}Ὁ δὲ εἶπεν· Οὐχί, πάτερ Ἀβραάμ, ἀλλ᾽ ἐάν τις ἀπὸ νεκρῶν πορευθῇ πρὸς αὐτοὺς μετανοήσουσιν.
\vs{31}Εἶπεν δὲ αὐτῷ· Εἰ Μωϋσέως καὶ τῶν προφητῶν οὐκ ἀκούουσιν, οὐδ᾽ ἐάν τις ἐκ νεκρῶν ἀναστῇ πεισθήσονται.

\ch{17}
Εἶπεν δὲ πρὸς τοὺς μαθητὰς αὐτοῦ· Ἀνένδεκτόν ἐστιν τοῦ τὰ σκάνδαλα μὴ ἐλθεῖν, πλὴν οὐαὶ δι᾽ οὗ ἔρχεται·
\vs{2}λυσιτελεῖ αὐτῷ εἰ λίθος μυλικὸς περίκειται περὶ τὸν τράχηλον αὐτοῦ καὶ ἔρριπται εἰς τὴν θάλασσαν ἢ ἵνα σκανδαλίσῃ τῶν μικρῶν τούτων ἕνα.
\vs{3}Προσέχετε ἑαυτοῖς.

ἐὰν ἁμάρτῃ ὁ ἀδελφός σου ἐπιτίμησον αὐτῷ, καὶ ἐὰν μετανοήσῃ ἄφες αὐτῷ.
\vs{4}καὶ ἐὰν ἑπτάκις τῆς ἡμέρας ἁμαρτήσῃ εἰς σὲ καὶ ἑπτάκις ἐπιστρέψῃ πρὸς σὲ λέγων· Μετανοῶ, ἀφήσεις αὐτῷ.

\vs{5}Καὶ εἶπαν οἱ ἀπόστολοι τῷ Κυρίῳ· Πρόσθες ἡμῖν πίστιν.
\vs{6}Εἶπεν δὲ ὁ Κύριος· Εἰ ἔχετε πίστιν ὡς κόκκον σινάπεως, ἐλέγετε ἂν τῇ συκαμίνῳ ταύτῃ· Ἐκριζώθητι καὶ φυτεύθητι ἐν τῇ θαλάσσῃ· καὶ ὑπήκουσεν ἂν ὑμῖν.

\vs{7}Τίς δὲ ἐξ ὑμῶν δοῦλον ἔχων ἀροτριῶντα ἢ ποιμαίνοντα, ὃς εἰσελθόντι ἐκ τοῦ ἀγροῦ ἐρεῖ αὐτῷ· Εὐθέως παρελθὼν ἀνάπεσε,
\vs{8}ἀλλ᾽ οὐχὶ ἐρεῖ αὐτῷ· Ἑτοίμασον τί δειπνήσω καὶ περιζωσάμενος διακόνει μοι ἕως φάγω καὶ πίω, καὶ μετὰ ταῦτα φάγεσαι καὶ πίεσαι σύ;
\vs{9}μὴ ἔχει χάριν τῷ δούλῳ ὅτι ἐποίησεν τὰ διαταχθέντα;
\vs{10}οὕτως καὶ ὑμεῖς, ὅταν ποιήσητε πάντα τὰ διαταχθέντα ὑμῖν, λέγετε ὅτι Δοῦλοι ἀχρεῖοί ἐσμεν, ὃ ὠφείλομεν ποιῆσαι πεποιήκαμεν.

\vs{11}Καὶ ἐγένετο ἐν τῷ πορεύεσθαι εἰς Ἰερουσαλὴμ καὶ αὐτὸς διήρχετο διὰ μέσον Σαμαρείας καὶ Γαλιλαίας.

\vs{12}καὶ εἰσερχομένου αὐτοῦ εἴς τινα κώμην ἀπήντησαν αὐτῷ δέκα λεπροὶ ἄνδρες, οἳ ἔστησαν πόρρωθεν
\vs{13}καὶ αὐτοὶ ἦραν φωνὴν λέγοντες· Ἰησοῦ Ἐπιστάτα, ἐλέησον ἡμᾶς.
\vs{14}Καὶ ἰδὼν εἶπεν αὐτοῖς· Πορευθέντες ἐπιδείξατε ἑαυτοὺς τοῖς ἱερεῦσιν. καὶ ἐγένετο ἐν τῷ ὑπάγειν αὐτοὺς ἐκαθαρίσθησαν.
\vs{15}Εἷς δὲ ἐξ αὐτῶν, ἰδὼν ὅτι ἰάθη, ὑπέστρεψεν μετὰ φωνῆς μεγάλης δοξάζων τὸν Θεόν,
\vs{16}καὶ ἔπεσεν ἐπὶ πρόσωπον παρὰ τοὺς πόδας αὐτοῦ εὐχαριστῶν αὐτῷ· καὶ αὐτὸς ἦν Σαμαρίτης.
\vs{17}Ἀποκριθεὶς δὲ ὁ Ἰησοῦς εἶπεν· Οὐχὶ οἱ δέκα ἐκαθαρίσθησαν; οἱ δὲ ἐννέα ποῦ;
\vs{18}οὐχ εὑρέθησαν ὑποστρέψαντες δοῦναι δόξαν τῷ Θεῷ εἰ μὴ ὁ ἀλλογενὴς οὗτος;
\vs{19}καὶ εἶπεν αὐτῷ· Ἀναστὰς πορεύου· ἡ πίστις σου σέσωκέν σε.

\vs{20}Ἐπερωτηθεὶς δὲ ὑπὸ τῶν Φαρισαίων πότε ἔρχεται ἡ βασιλεία τοῦ Θεοῦ ἀπεκρίθη αὐτοῖς καὶ εἶπεν· Οὐκ ἔρχεται ἡ βασιλεία τοῦ Θεοῦ μετὰ παρατηρήσεως,
\vs{21}οὐδὲ ἐροῦσιν· Ἰδοὺ ὧδε ἤ· Ἐκεῖ, ἰδοὺ γὰρ ἡ βασιλεία τοῦ Θεοῦ ἐντὸς ὑμῶν ἐστιν.

\vs{22}Εἶπεν δὲ πρὸς τοὺς μαθητάς· Ἐλεύσονται ἡμέραι ὅτε ἐπιθυμήσετε μίαν τῶν ἡμερῶν τοῦ Υἱοῦ τοῦ ἀνθρώπου ἰδεῖν καὶ οὐκ ὄψεσθε.
\vs{23}καὶ ἐροῦσιν ὑμῖν· Ἰδοὺ ἐκεῖ, ἤ· Ἰδοὺ ὧδε· μὴ ἀπέλθητε μηδὲ διώξητε.
\vs{24}ὥσπερ γὰρ ἡ ἀστραπὴ ἀστράπτουσα ἐκ τῆς ὑπὸ τὸν οὐρανὸν εἰς τὴν ὑπ᾽ οὐρανὸν λάμπει, οὕτως ἔσται ὁ Υἱὸς τοῦ ἀνθρώπου ἐν τῇ ἡμέρᾳ αὐτοῦ.
\vs{25}πρῶτον δὲ δεῖ αὐτὸν πολλὰ παθεῖν καὶ ἀποδοκιμασθῆναι ἀπὸ τῆς γενεᾶς ταύτης.
\vs{26}Καὶ καθὼς ἐγένετο ἐν ταῖς ἡμέραις Νῶε, οὕτως ἔσται καὶ ἐν ταῖς ἡμέραις τοῦ Υἱοῦ τοῦ ἀνθρώπου·
\vs{27}ἤσθιον, ἔπινον, ἐγάμουν, ἐγαμίζοντο, ἄχρι ἧς ἡμέρας εἰσῆλθεν Νῶε εἰς τὴν κιβωτόν καὶ ἦλθεν ὁ κατακλυσμὸς καὶ ἀπώλεσεν πάντας.
\vs{28}Ὁμοίως καθὼς ἐγένετο ἐν ταῖς ἡμέραις Λώτ· ἤσθιον, ἔπινον, ἠγόραζον, ἐπώλουν, ἐφύτευον, ᾠκοδόμουν·
\vs{29}ᾗ δὲ ἡμέρᾳ ἐξῆλθεν Λὼτ ἀπὸ Σοδόμων, ἔβρεξεν πῦρ καὶ θεῖον ἀπ᾽ οὐρανοῦ καὶ ἀπώλεσεν πάντας.
\vs{30}Κατὰ τὰ αὐτὰ ἔσται ᾗ ἡμέρᾳ ὁ Υἱὸς τοῦ ἀνθρώπου ἀποκαλύπτεται.
\vs{31}ἐν ἐκείνῃ τῇ ἡμέρᾳ ὃς ἔσται ἐπὶ τοῦ δώματος καὶ τὰ σκεύη αὐτοῦ ἐν τῇ οἰκίᾳ, μὴ καταβάτω ἆραι αὐτά, καὶ ὁ ἐν ἀγρῷ ὁμοίως μὴ ἐπιστρεψάτω εἰς τὰ ὀπίσω.
\vs{32}μνημονεύετε τῆς γυναικὸς Λώτ.
\vs{33}ὃς ἐὰν ζητήσῃ τὴν ψυχὴν αὐτοῦ περιποιήσασθαι ἀπολέσει αὐτήν, ὃς δ᾽ ἂν ἀπολέσῃ ζωογονήσει αὐτήν.
\vs{34}λέγω ὑμῖν, ταύτῃ τῇ νυκτὶ ἔσονται δύο ἐπὶ κλίνης μιᾶς, ὁ εἷς παραλημφθήσεται καὶ ὁ ἕτερος ἀφεθήσεται·
\vs{35}ἔσονται δύο ἀλήθουσαι ἐπὶ τὸ αὐτό, ἡ μία παραλημφθήσεται, ἡ δὲ ἑτέρα ἀφεθήσεται.
\vs{37}Καὶ ἀποκριθέντες λέγουσιν αὐτῷ· Ποῦ, Κύριε; Ὁ δὲ εἶπεν αὐτοῖς· Ὅπου τὸ σῶμα, ἐκεῖ καὶ οἱ ἀετοὶ ἐπισυναχθήσονται.

\ch{18}
Ἔλεγεν δὲ παραβολὴν αὐτοῖς πρὸς τὸ δεῖν πάντοτε προσεύχεσθαι αὐτοὺς καὶ μὴ ἐνκακεῖν,
\vs{2}λέγων· Κριτής τις ἦν ἔν τινι πόλει τὸν Θεὸν μὴ φοβούμενος καὶ ἄνθρωπον μὴ ἐντρεπόμενος.
\vs{3}χήρα δὲ ἦν ἐν τῇ πόλει ἐκείνῃ καὶ ἤρχετο πρὸς αὐτὸν λέγουσα· Ἐκδίκησόν με ἀπὸ τοῦ ἀντιδίκου μου.
\vs{4}Καὶ οὐκ ἤθελεν ἐπὶ χρόνον. μετὰ ταῦτα δὲ εἶπεν ἐν ἑαυτῷ· Εἰ καὶ τὸν Θεὸν οὐ φοβοῦμαι οὐδὲ ἄνθρωπον ἐντρέπομαι,
\vs{5}διά γε τὸ παρέχειν μοι κόπον τὴν χήραν ταύτην ἐκδικήσω αὐτήν, ἵνα μὴ εἰς τέλος ἐρχομένη ὑπωπιάζῃ με.
\vs{6}Εἶπεν δὲ ὁ Κύριος· Ἀκούσατε τί ὁ κριτὴς τῆς ἀδικίας λέγει·
\vs{7}ὁ δὲ Θεὸς οὐ μὴ ποιήσῃ τὴν ἐκδίκησιν τῶν ἐκλεκτῶν αὐτοῦ τῶν βοώντων αὐτῷ ἡμέρας καὶ νυκτός, καὶ μακροθυμεῖ ἐπ᾽ αὐτοῖς;
\vs{8}λέγω ὑμῖν ὅτι ποιήσει τὴν ἐκδίκησιν αὐτῶν ἐν τάχει. πλὴν ὁ Υἱὸς τοῦ ἀνθρώπου ἐλθὼν ἆρα εὑρήσει τὴν πίστιν ἐπὶ τῆς γῆς;

\vs{9}Εἶπεν δὲ καὶ πρός τινας τοὺς πεποιθότας ἐφ᾽ ἑαυτοῖς ὅτι εἰσὶν δίκαιοι καὶ ἐξουθενοῦντας τοὺς λοιποὺς τὴν παραβολὴν ταύτην·
\vs{10}Ἄνθρωποι δύο ἀνέβησαν εἰς τὸ ἱερὸν προσεύξασθαι, ὁ εἷς Φαρισαῖος καὶ ὁ ἕτερος τελώνης.
\vs{11}ὁ Φαρισαῖος σταθεὶς πρὸς ἑαυτὸν ταῦτα προσηύχετο· Ὁ Θεός, εὐχαριστῶ σοι ὅτι οὐκ εἰμὶ ὥσπερ οἱ λοιποὶ τῶν ἀνθρώπων, ἅρπαγες, ἄδικοι, μοιχοί, ἢ καὶ ὡς οὗτος ὁ τελώνης·
\vs{12}νηστεύω δὶς τοῦ σαββάτου, ἀποδεκατῶ πάντα ὅσα κτῶμαι.
\vs{13}Ὁ δὲ τελώνης μακρόθεν ἑστὼς οὐκ ἤθελεν οὐδὲ τοὺς ὀφθαλμοὺς ἐπᾶραι εἰς τὸν οὐρανόν, ἀλλ᾽ ἔτυπτεν τὸ στῆθος αὐτοῦ λέγων· Ὁ Θεός, ἱλάσθητί μοι τῷ ἁμαρτωλῷ.
\vs{14}λέγω ὑμῖν, κατέβη οὗτος δεδικαιωμένος εἰς τὸν οἶκον αὐτοῦ παρ᾽ ἐκεῖνον· ὅτι πᾶς ὁ ὑψῶν ἑαυτὸν ταπεινωθήσεται, ὁ δὲ ταπεινῶν ἑαυτὸν ὑψωθήσεται.

\vs{15}Προσέφερον δὲ αὐτῷ καὶ τὰ βρέφη ἵνα αὐτῶν ἅπτηται· ἰδόντες δὲ οἱ μαθηταὶ ἐπετίμων αὐτοῖς.
\vs{16}Ὁ δὲ Ἰησοῦς προσεκαλέσατο αὐτὰ λέγων· Ἄφετε τὰ παιδία ἔρχεσθαι πρός με καὶ μὴ κωλύετε αὐτά, τῶν γὰρ τοιούτων ἐστὶν ἡ βασιλεία τοῦ Θεοῦ.
\vs{17}ἀμὴν λέγω ὑμῖν, ὃς ἂν μὴ δέξηται τὴν βασιλείαν τοῦ Θεοῦ ὡς παιδίον, οὐ μὴ εἰσέλθῃ εἰς αὐτήν.

\vs{18}Καὶ ἐπηρώτησέν τις αὐτὸν ἄρχων λέγων· Διδάσκαλε ἀγαθέ, τί ποιήσας ζωὴν αἰώνιον κληρονομήσω;
\vs{19}Εἶπεν δὲ αὐτῷ ὁ Ἰησοῦς· Τί με λέγεις ἀγαθόν; οὐδεὶς ἀγαθὸς εἰ μὴ εἷς ὁ Θεός.
\vs{20}τὰς ἐντολὰς οἶδας· Μὴ μοιχεύσῃς, Μὴ φονεύσῃς, Μὴ κλέψῃς, Μὴ ψευδομαρτυρήσῃς, Τίμα τὸν πατέρα σου καὶ τὴν μητέρα.
\vs{21}Ὁ δὲ εἶπεν· Ταῦτα πάντα ἐφύλαξα ἐκ νεότητος.
\vs{22}Ἀκούσας δὲ ὁ Ἰησοῦς εἶπεν αὐτῷ· Ἔτι ἕν σοι λείπει· πάντα ὅσα ἔχεις πώλησον καὶ διάδος πτωχοῖς, καὶ ἕξεις θησαυρὸν ἐν τοῖς οὐρανοῖς, καὶ δεῦρο ἀκολούθει μοι.
\vs{23}Ὁ δὲ ἀκούσας ταῦτα περίλυπος ἐγενήθη· ἦν γὰρ πλούσιος σφόδρα.

\vs{24}Ἰδὼν δὲ αὐτὸν ὁ Ἰησοῦς περίλυπον γενόμενον εἶπεν· Πῶς δυσκόλως οἱ τὰ χρήματα ἔχοντες εἰς τὴν βασιλείαν τοῦ Θεοῦ εἰσπορεύονται·
\vs{25}εὐκοπώτερον γάρ ἐστιν κάμηλον διὰ τρήματος βελόνης εἰσελθεῖν ἢ πλούσιον εἰς τὴν βασιλείαν τοῦ Θεοῦ εἰσελθεῖν.
\vs{26}Εἶπαν δὲ οἱ ἀκούσαντες· Καὶ τίς δύναται σωθῆναι;
\vs{27}Ὁ δὲ εἶπεν· Τὰ ἀδύνατα παρὰ ἀνθρώποις δυνατὰ παρὰ τῷ Θεῷ ἐστιν.

\vs{28}Εἶπεν δὲ ὁ Πέτρος· Ἰδοὺ ἡμεῖς ἀφέντες τὰ ἴδια ἠκολουθήσαμέν σοι.
\vs{29}Ὁ δὲ εἶπεν αὐτοῖς· Ἀμὴν λέγω ὑμῖν ὅτι οὐδείς ἐστιν ὃς ἀφῆκεν οἰκίαν ἢ γυναῖκα ἢ ἀδελφοὺς ἢ γονεῖς ἢ τέκνα ἕνεκεν τῆς βασιλείας τοῦ Θεοῦ,
\vs{30}ὃς οὐχὶ μὴ ἀπολάβῃ πολλαπλασίονα ἐν τῷ καιρῷ τούτῳ καὶ ἐν τῷ αἰῶνι τῷ ἐρχομένῳ ζωὴν αἰώνιον.

\vs{31}Παραλαβὼν δὲ τοὺς δώδεκα εἶπεν πρὸς αὐτούς· Ἰδοὺ ἀναβαίνομεν εἰς Ἰερουσαλήμ, καὶ τελεσθήσεται πάντα τὰ γεγραμμένα διὰ τῶν προφητῶν τῷ Υἱῷ τοῦ ἀνθρώπου·
\vs{32}παραδοθήσεται γὰρ τοῖς ἔθνεσιν καὶ ἐμπαιχθήσεται καὶ ὑβρισθήσεται καὶ ἐμπτυσθήσεται
\vs{33}καὶ μαστιγώσαντες ἀποκτενοῦσιν αὐτόν, καὶ τῇ ἡμέρᾳ τῇ τρίτῃ ἀναστήσεται.
\vs{34}Καὶ αὐτοὶ οὐδὲν τούτων συνῆκαν καὶ ἦν τὸ ῥῆμα τοῦτο κεκρυμμένον ἀπ᾽ αὐτῶν καὶ οὐκ ἐγίνωσκον τὰ λεγόμενα.

\vs{35}Ἐγένετο δὲ ἐν τῷ ἐγγίζειν αὐτὸν εἰς Ἰεριχὼ τυφλός τις ἐκάθητο παρὰ τὴν ὁδὸν ἐπαιτῶν.
\vs{36}ἀκούσας δὲ ὄχλου διαπορευομένου ἐπυνθάνετο τί εἴη τοῦτο.
\vs{37}Ἀπήγγειλαν δὲ αὐτῷ ὅτι Ἰησοῦς ὁ Ναζωραῖος παρέρχεται.
\vs{38}Καὶ ἐβόησεν λέγων· Ἰησοῦ υἱὲ Δαυίδ, ἐλέησόν με.
\vs{39}Καὶ οἱ προάγοντες ἐπετίμων αὐτῷ ἵνα σιγήσῃ, αὐτὸς δὲ πολλῷ μᾶλλον ἔκραζεν· Υἱὲ Δαυίδ, ἐλέησόν με.
\vs{40}Σταθεὶς δὲ ὁ Ἰησοῦς ἐκέλευσεν αὐτὸν ἀχθῆναι πρὸς αὐτόν. ἐγγίσαντος δὲ αὐτοῦ ἐπηρώτησεν αὐτόν·
\vs{41}Τί σοι θέλεις ποιήσω; Ὁ δὲ εἶπεν· Κύριε, ἵνα ἀναβλέψω.
\vs{42}Καὶ ὁ Ἰησοῦς εἶπεν αὐτῷ· Ἀνάβλεψον· ἡ πίστις σου σέσωκέν σε.
\vs{43}καὶ παραχρῆμα ἀνέβλεψεν καὶ ἠκολούθει αὐτῷ δοξάζων τὸν Θεόν. καὶ πᾶς ὁ λαὸς ἰδὼν ἔδωκεν αἶνον τῷ Θεῷ.

\ch{19}
Καὶ εἰσελθὼν διήρχετο τὴν Ἰεριχώ.
\vs{2}Καὶ ἰδοὺ ἀνὴρ ὀνόματι καλούμενος Ζακχαῖος, καὶ αὐτὸς ἦν ἀρχιτελώνης καὶ αὐτὸς πλούσιος·
\vs{3}καὶ ἐζήτει ἰδεῖν τὸν Ἰησοῦν τίς ἐστιν καὶ οὐκ ἠδύνατο ἀπὸ τοῦ ὄχλου, ὅτι τῇ ἡλικίᾳ μικρὸς ἦν.
\vs{4}καὶ προδραμὼν εἰς τὸ ἔμπροσθεν ἀνέβη ἐπὶ συκομορέαν ἵνα ἴδῃ αὐτόν ὅτι ἐκείνης ἤμελλεν διέρχεσθαι.
\vs{5}Καὶ ὡς ἦλθεν ἐπὶ τὸν τόπον, ἀναβλέψας ὁ Ἰησοῦς εἶπεν πρὸς αὐτόν· Ζακχαῖε, σπεύσας κατάβηθι, σήμερον γὰρ ἐν τῷ οἴκῳ σου δεῖ με μεῖναι.
\vs{6}Καὶ σπεύσας κατέβη καὶ ὑπεδέξατο αὐτὸν χαίρων.
\vs{7}καὶ ἰδόντες πάντες διεγόγγυζον λέγοντες ὅτι Παρὰ ἁμαρτωλῷ ἀνδρὶ εἰσῆλθεν καταλῦσαι.
\vs{8}Σταθεὶς δὲ Ζακχαῖος εἶπεν πρὸς τὸν Κύριον· Ἰδοὺ τὰ ἡμίσιά μου τῶν ὑπαρχόντων, Κύριε, τοῖς πτωχοῖς δίδωμι, καὶ εἴ τινός τι ἐσυκοφάντησα ἀποδίδωμι τετραπλοῦν.
\vs{9}Εἶπεν δὲ πρὸς αὐτὸν ὁ Ἰησοῦς ὅτι Σήμερον σωτηρία τῷ οἴκῳ τούτῳ ἐγένετο, καθότι καὶ αὐτὸς υἱὸς Ἀβραάμ ἐστιν·
\vs{10}ἦλθεν γὰρ ὁ Υἱὸς τοῦ ἀνθρώπου ζητῆσαι καὶ σῶσαι τὸ ἀπολωλός.

\vs{11}Ἀκουόντων δὲ αὐτῶν ταῦτα προσθεὶς εἶπεν παραβολὴν διὰ τὸ ἐγγὺς εἶναι Ἰερουσαλὴμ αὐτὸν καὶ δοκεῖν αὐτοὺς ὅτι παραχρῆμα μέλλει ἡ βασιλεία τοῦ Θεοῦ ἀναφαίνεσθαι.
\vs{12}εἶπεν οὖν· Ἄνθρωπός τις εὐγενὴς ἐπορεύθη εἰς χώραν μακρὰν λαβεῖν ἑαυτῷ βασιλείαν καὶ ὑποστρέψαι.
\vs{13}καλέσας δὲ δέκα δούλους ἑαυτοῦ ἔδωκεν αὐτοῖς δέκα μνᾶς καὶ εἶπεν πρὸς αὐτούς· Πραγματεύσασθε ἐν ᾧ ἔρχομαι.
\vs{14}Οἱ δὲ πολῖται αὐτοῦ ἐμίσουν αὐτόν καὶ ἀπέστειλαν πρεσβείαν ὀπίσω αὐτοῦ λέγοντες· Οὐ θέλομεν τοῦτον βασιλεῦσαι ἐφ᾽ ἡμᾶς.
\vs{15}Καὶ ἐγένετο ἐν τῷ ἐπανελθεῖν αὐτὸν λαβόντα τὴν βασιλείαν καὶ εἶπεν φωνηθῆναι αὐτῷ τοὺς δούλους τούτους οἷς δεδώκει τὸ ἀργύριον, ἵνα γνοῖ τί διεπραγματεύσαντο.
\vs{16}Παρεγένετο δὲ ὁ πρῶτος λέγων· Κύριε, ἡ μνᾶ σου δέκα προσηργάσατο μνᾶς.
\vs{17}Καὶ εἶπεν αὐτῷ· Εὖγε, ἀγαθὲ δοῦλε, ὅτι ἐν ἐλαχίστῳ πιστὸς ἐγένου, ἴσθι ἐξουσίαν ἔχων ἐπάνω δέκα πόλεων.
\vs{18}Καὶ ἦλθεν ὁ δεύτερος λέγων· Ἡ μνᾶ σου, κύριε, ἐποίησεν πέντε μνᾶς.
\vs{19}Εἶπεν δὲ καὶ τούτῳ· Καὶ σὺ ἐπάνω γίνου πέντε πόλεων.
\vs{20}Καὶ ὁ ἕτερος ἦλθεν λέγων· Κύριε, ἰδοὺ ἡ μνᾶ σου ἣν εἶχον ἀποκειμένην ἐν σουδαρίῳ·
\vs{21}ἐφοβούμην γάρ σε, ὅτι ἄνθρωπος αὐστηρὸς εἶ, αἴρεις ὃ οὐκ ἔθηκας καὶ θερίζεις ὃ οὐκ ἔσπειρας.
\vs{22}Λέγει αὐτῷ· Ἐκ τοῦ στόματός σου κρίνω σε, πονηρὲ δοῦλε. ᾔδεις ὅτι ἐγὼ ἄνθρωπος αὐστηρός εἰμι, αἴρων ὃ οὐκ ἔθηκα καὶ θερίζων ὃ οὐκ ἔσπειρα;
\vs{23}καὶ διὰ τί οὐκ ἔδωκάς μου τὸ ἀργύριον ἐπὶ τράπεζαν; κἀγὼ ἐλθὼν σὺν τόκῳ ἂν αὐτὸ ἔπραξα.
\vs{24}Καὶ τοῖς παρεστῶσιν εἶπεν· Ἄρατε ἀπ᾽ αὐτοῦ τὴν μνᾶν καὶ δότε τῷ τὰς δέκα μνᾶς ἔχοντι—
\vs{25}Καὶ εἶπαν αὐτῷ· Κύριε, ἔχει δέκα μνᾶς—
\vs{26}Λέγω ὑμῖν ὅτι παντὶ τῷ ἔχοντι δοθήσεται, ἀπὸ δὲ τοῦ μὴ ἔχοντος καὶ ὃ ἔχει ἀρθήσεται.
\vs{27}πλὴν τοὺς ἐχθρούς μου τούτους τοὺς μὴ θελήσαντάς με βασιλεῦσαι ἐπ᾽ αὐτοὺς ἀγάγετε ὧδε καὶ κατασφάξατε αὐτοὺς ἔμπροσθέν μου.

\vs{28}Καὶ εἰπὼν ταῦτα ἐπορεύετο ἔμπροσθεν ἀναβαίνων εἰς Ἱεροσόλυμα.

\vs{29}Καὶ ἐγένετο ὡς ἤγγισεν εἰς Βηθφαγὴ καὶ Βηθανίαν πρὸς τὸ ὄρος τὸ καλούμενον Ἐλαιῶν, ἀπέστειλεν δύο τῶν μαθητῶν
\vs{30}λέγων· Ὑπάγετε εἰς τὴν κατέναντι κώμην, ἐν ᾗ εἰσπορευόμενοι εὑρήσετε πῶλον δεδεμένον, ἐφ᾽ ὃν οὐδεὶς πώποτε ἀνθρώπων ἐκάθισεν, καὶ λύσαντες αὐτὸν ἀγάγετε.
\vs{31}καὶ ἐάν τις ὑμᾶς ἐρωτᾷ· Διὰ τί λύετε; οὕτως ἐρεῖτε· Ὅτι Ὁ Κύριος αὐτοῦ χρείαν ἔχει.
\vs{32}Ἀπελθόντες δὲ οἱ ἀπεσταλμένοι εὗρον καθὼς εἶπεν αὐτοῖς.
\vs{33}λυόντων δὲ αὐτῶν τὸν πῶλον εἶπαν οἱ κύριοι αὐτοῦ πρὸς αὐτούς· Τί λύετε τὸν πῶλον;
\vs{34}Οἱ δὲ εἶπαν· ὅτι Ὁ Κύριος αὐτοῦ χρείαν ἔχει.
\vs{35}Καὶ ἤγαγον αὐτὸν πρὸς τὸν Ἰησοῦν καὶ ἐπιρίψαντες αὐτῶν τὰ ἱμάτια ἐπὶ τὸν πῶλον ἐπεβίβασαν τὸν Ἰησοῦν.
\vs{36}Πορευομένου δὲ αὐτοῦ ὑπεστρώννυον τὰ ἱμάτια ἑαυτῶν ἐν τῇ ὁδῷ.
\vs{37}ἐγγίζοντος δὲ αὐτοῦ ἤδη πρὸς τῇ καταβάσει τοῦ ὄρους τῶν Ἐλαιῶν ἤρξαντο ἅπαν τὸ πλῆθος τῶν μαθητῶν χαίροντες αἰνεῖν τὸν Θεὸν φωνῇ μεγάλῃ περὶ πασῶν ὧν εἶδον δυνάμεων,
\vs{38}λέγοντες· 
\begin{poetryblock}

\begin{quote}Εὐλογημένος ὁ ἐρχόμενος,\end{quote} 

\begin{quote}ὁ Βασιλεὺς ἐν ὀνόματι Κυρίου· Ἐν οὐρανῷ εἰρήνη καὶ δόξα ἐν ὑψίστοις.\end{quote}
\end{poetryblock}
\vs{39}Καί τινες τῶν Φαρισαίων ἀπὸ τοῦ ὄχλου εἶπαν πρὸς αὐτόν· Διδάσκαλε, ἐπιτίμησον τοῖς μαθηταῖς σου.
\vs{40}Καὶ ἀποκριθεὶς εἶπεν· Λέγω ὑμῖν, ἐὰν οὗτοι σιωπήσουσιν, οἱ λίθοι κράξουσιν.

\vs{41}Καὶ ὡς ἤγγισεν ἰδὼν τὴν πόλιν ἔκλαυσεν ἐπ᾽ αὐτήν
\vs{42}λέγων ὅτι Εἰ ἔγνως ἐν τῇ ἡμέρᾳ ταύτῃ καὶ σὺ τὰ πρὸς εἰρήνην· νῦν δὲ ἐκρύβη ἀπὸ ὀφθαλμῶν σου.
\vs{43}ὅτι ἥξουσιν ἡμέραι ἐπὶ σὲ καὶ παρεμβαλοῦσιν οἱ ἐχθροί σου χάρακά σοι καὶ περικυκλώσουσίν σε καὶ συνέξουσίν σε πάντοθεν,
\vs{44}καὶ ἐδαφιοῦσίν σε καὶ τὰ τέκνα σου ἐν σοί, καὶ οὐκ ἀφήσουσιν λίθον ἐπὶ λίθον ἐν σοί, ἀνθ᾽ ὧν οὐκ ἔγνως τὸν καιρὸν τῆς ἐπισκοπῆς σου.

\vs{45}Καὶ εἰσελθὼν εἰς τὸ ἱερὸν ἤρξατο ἐκβάλλειν τοὺς πωλοῦντας
\vs{46}λέγων αὐτοῖς· Γέγραπται· 
\begin{poetryblock}

\begin{quote}Καὶ ἔσται ὁ οἶκός μου οἶκος προσευχῆς, ὑμεῖς δὲ αὐτὸν ἐποιήσατε Σπήλαιον λῃστῶν.\end{quote}
\end{poetryblock}
\vs{47}Καὶ ἦν διδάσκων τὸ καθ᾽ ἡμέραν ἐν τῷ ἱερῷ. οἱ δὲ ἀρχιερεῖς καὶ οἱ γραμματεῖς ἐζήτουν αὐτὸν ἀπολέσαι καὶ οἱ πρῶτοι τοῦ λαοῦ,
\vs{48}καὶ οὐχ εὕρισκον τὸ τί ποιήσωσιν, ὁ λαὸς γὰρ ἅπας ἐξεκρέματο αὐτοῦ ἀκούων.

\ch{20}
Καὶ ἐγένετο ἐν μιᾷ τῶν ἡμερῶν διδάσκοντος αὐτοῦ τὸν λαὸν ἐν τῷ ἱερῷ καὶ εὐαγγελιζομένου ἐπέστησαν οἱ ἀρχιερεῖς καὶ οἱ γραμματεῖς σὺν τοῖς πρεσβυτέροις
\vs{2}καὶ εἶπαν λέγοντες πρὸς αὐτόν· Εἰπὸν ἡμῖν ἐν ποίᾳ ἐξουσίᾳ ταῦτα ποιεῖς, ἢ τίς ἐστιν ὁ δούς σοι τὴν ἐξουσίαν ταύτην;
\vs{3}Ἀποκριθεὶς δὲ εἶπεν πρὸς αὐτούς· Ἐρωτήσω ὑμᾶς κἀγὼ λόγον, καὶ εἴπατέ μοι·
\vs{4}Τὸ βάπτισμα Ἰωάννου ἐξ οὐρανοῦ ἦν ἢ ἐξ ἀνθρώπων;
\vs{5}Οἱ δὲ συνελογίσαντο πρὸς ἑαυτοὺς λέγοντες ὅτι Ἐὰν εἴπωμεν· Ἐξ οὐρανοῦ, ἐρεῖ· Διὰ τί οὐκ ἐπιστεύσατε αὐτῷ;
\vs{6}ἐὰν δὲ εἴπωμεν· Ἐξ ἀνθρώπων, ὁ λαὸς ἅπας καταλιθάσει ἡμᾶς, πεπεισμένος γάρ ἐστιν Ἰωάννην προφήτην εἶναι.
\vs{7}Καὶ ἀπεκρίθησαν μὴ εἰδέναι πόθεν.
\vs{8}Καὶ ὁ Ἰησοῦς εἶπεν αὐτοῖς· Οὐδὲ ἐγὼ λέγω ὑμῖν ἐν ποίᾳ ἐξουσίᾳ ταῦτα ποιῶ.

\vs{9}Ἤρξατο δὲ πρὸς τὸν λαὸν λέγειν τὴν παραβολὴν ταύτην· Ἄνθρωπος τις ἐφύτευσεν ἀμπελῶνα καὶ ἐξέδετο αὐτὸν γεωργοῖς καὶ ἀπεδήμησεν χρόνους ἱκανούς.
\vs{10}καὶ καιρῷ ἀπέστειλεν πρὸς τοὺς γεωργοὺς δοῦλον ἵνα ἀπὸ τοῦ καρποῦ τοῦ ἀμπελῶνος δώσουσιν αὐτῷ· οἱ δὲ γεωργοὶ ἐξαπέστειλαν αὐτὸν δείραντες κενόν.
\vs{11}Καὶ προσέθετο ἕτερον πέμψαι δοῦλον· οἱ δὲ κἀκεῖνον δείραντες καὶ ἀτιμάσαντες ἐξαπέστειλαν κενόν.
\vs{12}Καὶ προσέθετο τρίτον πέμψαι· οἱ δὲ καὶ τοῦτον τραυματίσαντες ἐξέβαλον.
\vs{13}Εἶπεν δὲ ὁ κύριος τοῦ ἀμπελῶνος· Τί ποιήσω; πέμψω τὸν υἱόν μου τὸν ἀγαπητόν· ἴσως τοῦτον ἐντραπήσονται.
\vs{14}Ἰδόντες δὲ αὐτὸν οἱ γεωργοὶ διελογίζοντο πρὸς ἀλλήλους λέγοντες· Οὗτός ἐστιν ὁ κληρονόμος· ἀποκτείνωμεν αὐτόν, ἵνα ἡμῶν γένηται ἡ κληρονομία.
\vs{15}καὶ ἐκβαλόντες αὐτὸν ἔξω τοῦ ἀμπελῶνος ἀπέκτειναν. Τί οὖν ποιήσει αὐτοῖς ὁ κύριος τοῦ ἀμπελῶνος;
\vs{16}ἐλεύσεται καὶ ἀπολέσει τοὺς γεωργοὺς τούτους καὶ δώσει τὸν ἀμπελῶνα ἄλλοις.

Ἀκούσαντες δὲ εἶπαν· Μὴ γένοιτο.
\vs{17}Ὁ δὲ ἐμβλέψας αὐτοῖς εἶπεν· Τί οὖν ἐστιν τὸ γεγραμμένον τοῦτο· 
\begin{poetryblock}

\begin{quote}Λίθον ὃν ἀπεδοκίμασαν οἱ οἰκοδομοῦντες,\end{quote} 

\begin{quote}Οὗτος ἐγενήθη εἰς κεφαλὴν γωνίας;\end{quote}
\end{poetryblock}

\vs{18}Πᾶς ὁ πεσὼν ἐπ᾽ ἐκεῖνον τὸν λίθον συνθλασθήσεται· ἐφ᾽ ὃν δ᾽ ἂν πέσῃ, λικμήσει αὐτόν.

\vs{19}Καὶ ἐζήτησαν οἱ γραμματεῖς καὶ οἱ ἀρχιερεῖς ἐπιβαλεῖν ἐπ᾽ αὐτὸν τὰς χεῖρας ἐν αὐτῇ τῇ ὥρᾳ, καὶ ἐφοβήθησαν τὸν λαόν, ἔγνωσαν γὰρ ὅτι πρὸς αὐτοὺς εἶπεν τὴν παραβολὴν ταύτην.

\vs{20}Καὶ παρατηρήσαντες ἀπέστειλαν ἐνκαθέτους ὑποκρινομένους ἑαυτοὺς δικαίους εἶναι, ἵνα ἐπιλάβωνται αὐτοῦ λόγου, ὥστε παραδοῦναι αὐτὸν τῇ ἀρχῇ καὶ τῇ ἐξουσίᾳ τοῦ ἡγεμόνος.
\vs{21}καὶ ἐπηρώτησαν αὐτὸν λέγοντες· Διδάσκαλε, οἴδαμεν ὅτι ὀρθῶς λέγεις καὶ διδάσκεις καὶ οὐ λαμβάνεις πρόσωπον, ἀλλ᾽ ἐπ᾽ ἀληθείας τὴν ὁδὸν τοῦ Θεοῦ διδάσκεις·
\vs{22}ἔξεστιν ἡμᾶς Καίσαρι φόρον δοῦναι ἢ οὔ;
\vs{23}Κατανοήσας δὲ αὐτῶν τὴν πανουργίαν εἶπεν πρὸς αὐτούς·
\vs{24}Δείξατέ μοι δηνάριον· τίνος ἔχει εἰκόνα καὶ ἐπιγραφήν; Οἱ δὲ εἶπαν· Καίσαρος.
\vs{25}Ὁ δὲ εἶπεν πρὸς αὐτούς· Τοίνυν ἀπόδοτε τὰ Καίσαρος Καίσαρι καὶ τὰ τοῦ Θεοῦ τῷ Θεῷ.
\vs{26}Καὶ οὐκ ἴσχυσαν ἐπιλαβέσθαι αὐτοῦ ῥήματος ἐναντίον τοῦ λαοῦ καὶ θαυμάσαντες ἐπὶ τῇ ἀποκρίσει αὐτοῦ ἐσίγησαν.

\vs{27}Προσελθόντες δέ τινες τῶν Σαδδουκαίων, οἱ ἀντιλέγοντες ἀνάστασιν μὴ εἶναι, ἐπηρώτησαν αὐτὸν
\vs{28}λέγοντες· Διδάσκαλε, Μωϋσῆς ἔγραψεν ἡμῖν, ἐάν τινος ἀδελφὸς ἀποθάνῃ ἔχων γυναῖκα, καὶ οὗτος ἄτεκνος ᾖ, ἵνα λάβῃ ὁ ἀδελφὸς αὐτοῦ τὴν γυναῖκα καὶ ἐξαναστήσῃ σπέρμα τῷ ἀδελφῷ αὐτοῦ.
\vs{29}ἑπτὰ οὖν ἀδελφοὶ ἦσαν· καὶ ὁ πρῶτος λαβὼν γυναῖκα ἀπέθανεν ἄτεκνος·
\vs{30}καὶ ὁ δεύτερος
\vs{31}καὶ ὁ τρίτος ἔλαβεν αὐτήν, ὡσαύτως δὲ καὶ οἱ ἑπτὰ οὐ κατέλιπον τέκνα καὶ ἀπέθανον.
\vs{32}ὕστερον καὶ ἡ γυνὴ ἀπέθανεν.
\vs{33}ἡ γυνὴ οὖν ἐν τῇ ἀναστάσει τίνος αὐτῶν γίνεται γυνή; οἱ γὰρ ἑπτὰ ἔσχον αὐτὴν γυναῖκα.

\vs{34}Καὶ εἶπεν αὐτοῖς ὁ Ἰησοῦς· Οἱ υἱοὶ τοῦ αἰῶνος τούτου γαμοῦσιν καὶ γαμίσκονται,
\vs{35}οἱ δὲ καταξιωθέντες τοῦ αἰῶνος ἐκείνου τυχεῖν καὶ τῆς ἀναστάσεως τῆς ἐκ νεκρῶν οὔτε γαμοῦσιν οὔτε γαμίζονται·
\vs{36}οὐδὲ γὰρ ἀποθανεῖν ἔτι δύνανται, ἰσάγγελοι γάρ εἰσιν καὶ υἱοί εἰσιν Θεοῦ τῆς ἀναστάσεως υἱοὶ ὄντες.
\vs{37}Ὅτι δὲ ἐγείρονται οἱ νεκροὶ, καὶ Μωϋσῆς ἐμήνυσεν ἐπὶ τῆς Βάτου, ὡς λέγει Κύριον Τὸν Θεὸν Ἀβραὰμ καὶ Θεὸν Ἰσαὰκ καὶ Θεὸν Ἰακώβ.
\vs{38}Θεὸς δὲ οὐκ ἔστιν νεκρῶν ἀλλὰ ζώντων, πάντες γὰρ αὐτῷ ζῶσιν.

\vs{39}Ἀποκριθέντες δέ τινες τῶν γραμματέων εἶπαν· Διδάσκαλε, καλῶς εἶπας.
\vs{40}οὐκέτι γὰρ ἐτόλμων ἐπερωτᾶν αὐτὸν οὐδέν.

\vs{41}Εἶπεν δὲ πρὸς αὐτούς· Πῶς λέγουσιν τὸν Χριστὸν εἶναι Δαυὶδ υἱόν;
\vs{42}αὐτὸς γὰρ Δαυὶδ λέγει ἐν βίβλῳ ψαλμῶν· 
\begin{poetryblock}

\begin{quote}Εἶπεν Κύριος τῷ Κυρίῳ μου·\end{quote} 

\begin{quote}Κάθου ἐκ δεξιῶν μου,\end{quote}

\begin{quote} \vs{43}Ἕως ἂν θῶ τοὺς ἐχθρούς σου\end{quote} 

\begin{quote}Ὑποπόδιον τῶν ποδῶν σου.\end{quote}
\end{poetryblock}

\vs{44}Δαυὶδ οὖν Κύριον αὐτὸν καλεῖ, καὶ πῶς αὐτοῦ υἱός ἐστιν;

\vs{45}Ἀκούοντος δὲ παντὸς τοῦ λαοῦ εἶπεν τοῖς μαθηταῖς αὐτοῦ·
\vs{46}Προσέχετε ἀπὸ τῶν γραμματέων τῶν θελόντων περιπατεῖν ἐν στολαῖς καὶ φιλούντων ἀσπασμοὺς ἐν ταῖς ἀγοραῖς καὶ πρωτοκαθεδρίας ἐν ταῖς συναγωγαῖς καὶ πρωτοκλισίας ἐν τοῖς δείπνοις,
\vs{47}οἳ κατεσθίουσιν τὰς οἰκίας τῶν χηρῶν καὶ προφάσει μακρὰ προσεύχονται· οὗτοι λήμψονται περισσότερον κρίμα.

\ch{21}
Ἀναβλέψας δὲ εἶδεν τοὺς βάλλοντας εἰς τὸ γαζοφυλάκιον τὰ δῶρα αὐτῶν πλουσίους.
\vs{2}εἶδεν δέ τινα χήραν πενιχρὰν βάλλουσαν ἐκεῖ λεπτὰ δύο,
\vs{3}Καὶ εἶπεν· Ἀληθῶς λέγω ὑμῖν ὅτι ἡ χήρα αὕτη ἡ πτωχὴ πλεῖον πάντων ἔβαλεν·
\vs{4}πάντες γὰρ οὗτοι ἐκ τοῦ περισσεύοντος αὐτοῖς ἔβαλον εἰς τὰ δῶρα, αὕτη δὲ ἐκ τοῦ ὑστερήματος αὐτῆς πάντα τὸν βίον ὃν εἶχεν ἔβαλεν.

\vs{5}Καί τινων λεγόντων περὶ τοῦ ἱεροῦ ὅτι λίθοις καλοῖς καὶ ἀναθήμασιν κεκόσμηται εἶπεν·
\vs{6}Ταῦτα ἃ θεωρεῖτε ἐλεύσονται ἡμέραι ἐν αἷς οὐκ ἀφεθήσεται λίθος ἐπὶ λίθῳ ὃς οὐ καταλυθήσεται.

\vs{7}Ἐπηρώτησαν δὲ αὐτὸν λέγοντες· Διδάσκαλε, πότε οὖν ταῦτα ἔσται καὶ τί τὸ σημεῖον ὅταν μέλλῃ ταῦτα γίνεσθαι;
\vs{8}Ὁ δὲ εἶπεν· Βλέπετε μὴ πλανηθῆτε· πολλοὶ γὰρ ἐλεύσονται ἐπὶ τῷ ὀνόματί μου λέγοντες· Ἐγώ εἰμι, καί· Ὁ καιρὸς ἤγγικεν. μὴ πορευθῆτε ὀπίσω αὐτῶν.
\vs{9}ὅταν δὲ ἀκούσητε πολέμους καὶ ἀκαταστασίας, μὴ πτοηθῆτε· δεῖ γὰρ ταῦτα γενέσθαι πρῶτον, ἀλλ᾽ οὐκ εὐθέως τὸ τέλος.

\vs{10}Τότε ἔλεγεν αὐτοῖς· Ἐγερθήσεται ἔθνος ἐπ᾽ ἔθνος καὶ βασιλεία ἐπὶ βασιλείαν,
\vs{11}σεισμοί τε μεγάλοι καὶ κατὰ τόπους λιμοὶ καὶ λοιμοὶ ἔσονται, φόβητρά τε καὶ ἀπ᾽ οὐρανοῦ σημεῖα μεγάλα ἔσται.

\vs{12}Πρὸ δὲ τούτων πάντων ἐπιβαλοῦσιν ἐφ᾽ ὑμᾶς τὰς χεῖρας αὐτῶν καὶ διώξουσιν, παραδιδόντες εἰς τὰς συναγωγὰς καὶ φυλακάς, ἀπαγομένους ἐπὶ βασιλεῖς καὶ ἡγεμόνας ἕνεκεν τοῦ ὀνόματός μου·
\vs{13}ἀποβήσεται ὑμῖν εἰς μαρτύριον.
\vs{14}θέτε οὖν ἐν ταῖς καρδίαις ὑμῶν μὴ προμελετᾶν ἀπολογηθῆναι·
\vs{15}ἐγὼ γὰρ δώσω ὑμῖν στόμα καὶ σοφίαν ᾗ οὐ δυνήσονται ἀντιστῆναι ἢ ἀντειπεῖν ἅπαντες οἱ ἀντικείμενοι ὑμῖν.
\vs{16}Παραδοθήσεσθε δὲ καὶ ὑπὸ γονέων καὶ ἀδελφῶν καὶ συγγενῶν καὶ φίλων, καὶ θανατώσουσιν ἐξ ὑμῶν,
\vs{17}καὶ ἔσεσθε μισούμενοι ὑπὸ πάντων διὰ τὸ ὄνομά μου.
\vs{18}καὶ θρὶξ ἐκ τῆς κεφαλῆς ὑμῶν οὐ μὴ ἀπόληται.
\vs{19}ἐν τῇ ὑπομονῇ ὑμῶν κτήσασθε τὰς ψυχὰς ὑμῶν.

\vs{20}Ὅταν δὲ ἴδητε κυκλουμένην ὑπὸ στρατοπέδων Ἰερουσαλήμ, τότε γνῶτε ὅτι ἤγγικεν ἡ ἐρήμωσις αὐτῆς.
\vs{21}τότε οἱ ἐν τῇ Ἰουδαίᾳ φευγέτωσαν εἰς τὰ ὄρη καὶ οἱ ἐν μέσῳ αὐτῆς ἐκχωρείτωσαν καὶ οἱ ἐν ταῖς χώραις μὴ εἰσερχέσθωσαν εἰς αὐτήν,
\vs{22}ὅτι ἡμέραι ἐκδικήσεως αὗταί εἰσιν τοῦ πλησθῆναι πάντα τὰ γεγραμμένα.
\vs{23}Οὐαὶ ταῖς ἐν γαστρὶ ἐχούσαις καὶ ταῖς θηλαζούσαις ἐν ἐκείναις ταῖς ἡμέραις· ἔσται γὰρ ἀνάγκη μεγάλη ἐπὶ τῆς γῆς καὶ ὀργὴ τῷ λαῷ τούτῳ,
\vs{24}καὶ πεσοῦνται στόματι μαχαίρης καὶ αἰχμαλωτισθήσονται εἰς τὰ ἔθνη πάντα, καὶ Ἰερουσαλὴμ ἔσται πατουμένη ὑπὸ ἐθνῶν, ἄχρι οὗ πληρωθῶσιν καιροὶ ἐθνῶν.

\vs{25}Καὶ ἔσονται σημεῖα ἐν ἡλίῳ καὶ σελήνῃ καὶ ἄστροις, καὶ ἐπὶ τῆς γῆς συνοχὴ ἐθνῶν ἐν ἀπορίᾳ ἤχους θαλάσσης καὶ σάλου,
\vs{26}ἀποψυχόντων ἀνθρώπων ἀπὸ φόβου καὶ προσδοκίας τῶν ἐπερχομένων τῇ οἰκουμένῃ, αἱ γὰρ δυνάμεις τῶν οὐρανῶν σαλευθήσονται.
\vs{27}καὶ τότε ὄψονται τὸν Υἱὸν τοῦ ἀνθρώπου ἐρχόμενον ἐν νεφέλῃ μετὰ δυνάμεως καὶ δόξης πολλῆς.
\vs{28}ἀρχομένων δὲ τούτων γίνεσθαι ἀνακύψατε καὶ ἐπάρατε τὰς κεφαλὰς ὑμῶν, διότι ἐγγίζει ἡ ἀπολύτρωσις ὑμῶν.

\vs{29}Καὶ εἶπεν παραβολὴν αὐτοῖς· Ἴδετε τὴν συκῆν καὶ πάντα τὰ δένδρα·
\vs{30}ὅταν προβάλωσιν ἤδη, βλέποντες ἀφ᾽ ἑαυτῶν γινώσκετε ὅτι ἤδη ἐγγὺς τὸ θέρος ἐστίν·
\vs{31}οὕτως καὶ ὑμεῖς, ὅταν ἴδητε ταῦτα γινόμενα, γινώσκετε ὅτι ἐγγύς ἐστιν ἡ βασιλεία τοῦ Θεοῦ.
\vs{32}ἀμὴν λέγω ὑμῖν ὅτι οὐ μὴ παρέλθῃ ἡ γενεὰ αὕτη ἕως ἂν πάντα γένηται.
\vs{33}ὁ οὐρανὸς καὶ ἡ γῆ παρελεύσονται, οἱ δὲ λόγοι μου οὐ μὴ παρελεύσονται.

\vs{34}Προσέχετε δὲ ἑαυτοῖς μήποτε βαρηθῶσιν ὑμῶν αἱ καρδίαι ἐν κραιπάλῃ καὶ μέθῃ καὶ μερίμναις βιωτικαῖς καὶ ἐπιστῇ ἐφ᾽ ὑμᾶς αἰφνίδιος ἡ ἡμέρα ἐκείνη
\vs{35}ὡς παγίς· ἐπεισελεύσεται γὰρ ἐπὶ πάντας τοὺς καθημένους ἐπὶ πρόσωπον πάσης τῆς γῆς.
\vs{36}ἀγρυπνεῖτε δὲ ἐν παντὶ καιρῷ δεόμενοι ἵνα κατισχύσητε ἐκφυγεῖν ταῦτα πάντα τὰ μέλλοντα γίνεσθαι καὶ σταθῆναι ἔμπροσθεν τοῦ Υἱοῦ τοῦ ἀνθρώπου.

\vs{37}Ἦν δὲ τὰς ἡμέρας ἐν τῷ ἱερῷ διδάσκων, τὰς δὲ νύκτας ἐξερχόμενος ηὐλίζετο εἰς τὸ ὄρος τὸ καλούμενον Ἐλαιῶν·
\vs{38}καὶ πᾶς ὁ λαὸς ὤρθριζεν πρὸς αὐτὸν ἐν τῷ ἱερῷ ἀκούειν αὐτοῦ.

\ch{22}
Ἤγγιζεν δὲ ἡ ἑορτὴ τῶν ἀζύμων ἡ λεγομένη Πάσχα.
\vs{2}καὶ ἐζήτουν οἱ ἀρχιερεῖς καὶ οἱ γραμματεῖς τὸ πῶς ἀνέλωσιν αὐτόν, ἐφοβοῦντο γὰρ τὸν λαόν.

\vs{3}Εἰσῆλθεν δὲ Σατανᾶς εἰς Ἰούδαν τὸν καλούμενον Ἰσκαριώτην, ὄντα ἐκ τοῦ ἀριθμοῦ τῶν δώδεκα·
\vs{4}καὶ ἀπελθὼν συνελάλησεν τοῖς ἀρχιερεῦσιν καὶ στρατηγοῖς τὸ πῶς αὐτοῖς παραδῷ αὐτόν.
\vs{5}καὶ ἐχάρησαν καὶ συνέθεντο αὐτῷ ἀργύριον δοῦναι.
\vs{6}καὶ ἐξωμολόγησεν, καὶ ἐζήτει εὐκαιρίαν τοῦ παραδοῦναι αὐτὸν ἄτερ ὄχλου αὐτοῖς.

\vs{7}Ἦλθεν δὲ ἡ ἡμέρα τῶν ἀζύμων, ἐν ᾗ ἔδει θύεσθαι τὸ πάσχα·
\vs{8}καὶ ἀπέστειλεν Πέτρον καὶ Ἰωάννην εἰπών· Πορευθέντες ἑτοιμάσατε ἡμῖν τὸ πάσχα ἵνα φάγωμεν.
\vs{9}Οἱ δὲ εἶπαν αὐτῷ· Ποῦ θέλεις ἑτοιμάσωμεν;
\vs{10}Ὁ δὲ εἶπεν αὐτοῖς· Ἰδοὺ εἰσελθόντων ὑμῶν εἰς τὴν πόλιν συναντήσει ὑμῖν ἄνθρωπος κεράμιον ὕδατος βαστάζων· ἀκολουθήσατε αὐτῷ εἰς τὴν οἰκίαν εἰς ἣν εἰσπορεύεται,
\vs{11}καὶ ἐρεῖτε τῷ οἰκοδεσπότῃ τῆς οἰκίας· Λέγει σοι ὁ Διδάσκαλος· Ποῦ ἐστιν τὸ κατάλυμα ὅπου τὸ πάσχα μετὰ τῶν μαθητῶν μου φάγω;
\vs{12}κἀκεῖνος ὑμῖν δείξει ἀνάγαιον μέγα ἐστρωμένον· ἐκεῖ ἑτοιμάσατε.
\vs{13}Ἀπελθόντες δὲ εὗρον καθὼς εἰρήκει αὐτοῖς καὶ ἡτοίμασαν τὸ πάσχα.

\vs{14}Καὶ ὅτε ἐγένετο ἡ ὥρα, ἀνέπεσεν καὶ οἱ ἀπόστολοι σὺν αὐτῷ.
\vs{15}καὶ εἶπεν πρὸς αὐτούς· Ἐπιθυμίᾳ ἐπεθύμησα τοῦτο τὸ πάσχα φαγεῖν μεθ᾽ ὑμῶν πρὸ τοῦ με παθεῖν·
\vs{16}λέγω γὰρ ὑμῖν ὅτι οὐ μὴ φάγω αὐτὸ ἕως ὅτου πληρωθῇ ἐν τῇ βασιλείᾳ τοῦ Θεοῦ.
\vs{17}Καὶ δεξάμενος ποτήριον εὐχαριστήσας εἶπεν· Λάβετε τοῦτο καὶ διαμερίσατε εἰς ἑαυτούς·
\vs{18}λέγω γὰρ ὑμῖν, ὅτι οὐ μὴ πίω ἀπὸ τοῦ νῦν ἀπὸ τοῦ γενήματος τῆς ἀμπέλου ἕως οὗ ἡ βασιλεία τοῦ Θεοῦ ἔλθῃ.
\vs{19}Καὶ λαβὼν ἄρτον εὐχαριστήσας ἔκλασεν καὶ ἔδωκεν αὐτοῖς λέγων· Τοῦτό ἐστιν τὸ σῶμά μου τὸ ὑπὲρ ὑμῶν διδόμενον· τοῦτο ποιεῖτε εἰς τὴν ἐμὴν ἀνάμνησιν.
\vs{20}καὶ τὸ ποτήριον ὡσαύτως μετὰ τὸ δειπνῆσαι, λέγων· Τοῦτο τὸ ποτήριον ἡ καινὴ διαθήκη ἐν τῷ αἵματί μου τὸ ὑπὲρ ὑμῶν ἐκχυννόμενον.

\vs{21}Πλὴν ἰδοὺ ἡ χεὶρ τοῦ παραδιδόντος με μετ᾽ ἐμοῦ ἐπὶ τῆς τραπέζης.
\vs{22}ὅτι ὁ Υἱὸς μὲν τοῦ ἀνθρώπου κατὰ τὸ ὡρισμένον πορεύεται, πλὴν οὐαὶ τῷ ἀνθρώπῳ ἐκείνῳ δι᾽ οὗ παραδίδοται.
\vs{23}Καὶ αὐτοὶ ἤρξαντο συζητεῖν πρὸς ἑαυτοὺς τὸ τίς ἄρα εἴη ἐξ αὐτῶν ὁ τοῦτο μέλλων πράσσειν.

\vs{24}Ἐγένετο δὲ καὶ φιλονεικία ἐν αὐτοῖς, τὸ τίς αὐτῶν δοκεῖ εἶναι μείζων.
\vs{25}ὁ δὲ εἶπεν αὐτοῖς· Οἱ βασιλεῖς τῶν ἐθνῶν κυριεύουσιν αὐτῶν καὶ οἱ ἐξουσιάζοντες αὐτῶν εὐεργέται καλοῦνται.
\vs{26}ὑμεῖς δὲ οὐχ οὕτως, ἀλλ᾽ ὁ μείζων ἐν ὑμῖν γινέσθω ὡς ὁ νεώτερος καὶ ὁ ἡγούμενος ὡς ὁ διακονῶν.
\vs{27}τίς γὰρ μείζων, ὁ ἀνακείμενος ἢ ὁ διακονῶν; οὐχὶ ὁ ἀνακείμενος; ἐγὼ δὲ ἐν μέσῳ ὑμῶν εἰμι ὡς ὁ διακονῶν.

\vs{28}Ὑμεῖς δέ ἐστε οἱ διαμεμενηκότες μετ᾽ ἐμοῦ ἐν τοῖς πειρασμοῖς μου·
\vs{29}κἀγὼ διατίθεμαι ὑμῖν καθὼς διέθετό μοι ὁ Πατήρ μου βασιλείαν,
\vs{30}ἵνα ἔσθητε καὶ πίνητε ἐπὶ τῆς τραπέζης μου ἐν τῇ βασιλείᾳ μου, καὶ καθήσεσθε ἐπὶ θρόνων τὰς δώδεκα φυλὰς κρίνοντες τοῦ Ἰσραήλ.

\vs{31}Σίμων Σίμων, ἰδοὺ ὁ Σατανᾶς ἐξῃτήσατο ὑμᾶς τοῦ σινιάσαι ὡς τὸν σῖτον·
\vs{32}ἐγὼ δὲ ἐδεήθην περὶ σοῦ ἵνα μὴ ἐκλίπῃ ἡ πίστις σου· καὶ σύ ποτε ἐπιστρέψας στήρισον τοὺς ἀδελφούς σου.
\vs{33}Ὁ δὲ εἶπεν αὐτῷ· Κύριε, μετὰ σοῦ ἕτοιμός εἰμι καὶ εἰς φυλακὴν καὶ εἰς θάνατον πορεύεσθαι.
\vs{34}Ὁ δὲ εἶπεν· Λέγω σοι, Πέτρε, οὐ φωνήσει σήμερον ἀλέκτωρ ἕως τρίς με ἀπαρνήσῃ εἰδέναι.

\vs{35}Καὶ εἶπεν αὐτοῖς· Ὅτε ἀπέστειλα ὑμᾶς ἄτερ βαλλαντίου καὶ πήρας καὶ ὑποδημάτων, μή τινος ὑστερήσατε; Οἱ δὲ εἶπαν· Οὐθενός.
\vs{36}Εἶπεν δὲ αὐτοῖς· Ἀλλὰ νῦν ὁ ἔχων βαλλάντιον ἀράτω, ὁμοίως καὶ πήραν, καὶ ὁ μὴ ἔχων πωλησάτω τὸ ἱμάτιον αὐτοῦ καὶ ἀγορασάτω μάχαιραν.
\vs{37}λέγω γὰρ ὑμῖν ὅτι τοῦτο τὸ γεγραμμένον δεῖ τελεσθῆναι ἐν ἐμοί, Τό· Καὶ μετὰ ἀνόμων ἐλογίσθη· καὶ γὰρ τὸ περὶ ἐμοῦ τέλος ἔχει.
\vs{38}Οἱ δὲ εἶπαν· Κύριε, ἰδοὺ μάχαιραι ὧδε δύο. Ὁ δὲ εἶπεν αὐτοῖς· Ἱκανόν ἐστιν.

\vs{39}Καὶ ἐξελθὼν ἐπορεύθη κατὰ τὸ ἔθος εἰς τὸ ὄρος τῶν Ἐλαιῶν, ἠκολούθησαν δὲ αὐτῷ καὶ οἱ μαθηταί.
\vs{40}γενόμενος δὲ ἐπὶ τοῦ τόπου εἶπεν αὐτοῖς· Προσεύχεσθε μὴ εἰσελθεῖν εἰς πειρασμόν.
\vs{41}Καὶ αὐτὸς ἀπεσπάσθη ἀπ᾽ αὐτῶν ὡσεὶ λίθου βολήν καὶ θεὶς τὰ γόνατα προσηύχετο
\vs{42}λέγων· Πάτερ, εἰ βούλει παρένεγκε τοῦτο τὸ ποτήριον ἀπ᾽ ἐμοῦ· πλὴν μὴ τὸ θέλημά μου ἀλλὰ τὸ σὸν γινέσθω.
\vs{43}Ὤφθη δὲ αὐτῷ ἄγγελος ἀπ᾽ οὐρανοῦ ἐνισχύων αὐτόν.
\vs{44}καὶ γενόμενος ἐν ἀγωνίᾳ ἐκτενέστερον προσηύχετο· καὶ ἐγένετο ὁ ἱδρὼς αὐτοῦ ὡσεὶ θρόμβοι αἵματος καταβαίνοντες ἐπὶ τὴν γῆν.
\vs{45}Καὶ ἀναστὰς ἀπὸ τῆς προσευχῆς ἐλθὼν πρὸς τοὺς μαθητὰς εὗρεν κοιμωμένους αὐτοὺς ἀπὸ τῆς λύπης,
\vs{46}καὶ εἶπεν αὐτοῖς· Τί καθεύδετε; ἀναστάντες προσεύχεσθε, ἵνα μὴ εἰσέλθητε εἰς πειρασμόν.

\vs{47}Ἔτι αὐτοῦ λαλοῦντος ἰδοὺ ὄχλος, καὶ ὁ λεγόμενος Ἰούδας εἷς τῶν δώδεκα προήρχετο αὐτούς καὶ ἤγγισεν τῷ Ἰησοῦ φιλῆσαι αὐτόν.
\vs{48}Ἰησοῦς δὲ εἶπεν αὐτῷ· Ἰούδα, φιλήματι τὸν Υἱὸν τοῦ ἀνθρώπου παραδίδως;
\vs{49}Ἰδόντες δὲ οἱ περὶ αὐτὸν τὸ ἐσόμενον εἶπαν· Κύριε, εἰ πατάξομεν ἐν μαχαίρῃ;
\vs{50}καὶ ἐπάταξεν εἷς τις ἐξ αὐτῶν τοῦ ἀρχιερέως τὸν δοῦλον καὶ ἀφεῖλεν τὸ οὖς αὐτοῦ τὸ δεξιόν.
\vs{51}Ἀποκριθεὶς δὲ ὁ Ἰησοῦς εἶπεν· Ἐᾶτε ἕως τούτου· καὶ ἁψάμενος τοῦ ὠτίου ἰάσατο αὐτόν.

\vs{52}Εἶπεν δὲ Ἰησοῦς πρὸς τοὺς παραγενομένους ἐπ᾽ αὐτὸν ἀρχιερεῖς καὶ στρατηγοὺς τοῦ ἱεροῦ καὶ πρεσβυτέρους· Ὡς ἐπὶ λῃστὴν ἐξήλθατε μετὰ μαχαιρῶν καὶ ξύλων;
\vs{53}καθ᾽ ἡμέραν ὄντος μου μεθ᾽ ὑμῶν ἐν τῷ ἱερῷ οὐκ ἐξετείνατε τὰς χεῖρας ἐπ᾽ ἐμέ, ἀλλ᾽ αὕτη ἐστὶν ὑμῶν ἡ ὥρα καὶ ἡ ἐξουσία τοῦ σκότους.

\vs{54}Συλλαβόντες δὲ αὐτὸν ἤγαγον καὶ εἰσήγαγον εἰς τὴν οἰκίαν τοῦ ἀρχιερέως· ὁ δὲ Πέτρος ἠκολούθει μακρόθεν.
\vs{55}Περιαψάντων δὲ πῦρ ἐν μέσῳ τῆς αὐλῆς καὶ συνκαθισάντων ἐκάθητο ὁ Πέτρος μέσος αὐτῶν.
\vs{56}ἰδοῦσα δὲ αὐτὸν παιδίσκη τις καθήμενον πρὸς τὸ φῶς καὶ ἀτενίσασα αὐτῷ εἶπεν· Καὶ οὗτος σὺν αὐτῷ ἦν.
\vs{57}Ὁ δὲ ἠρνήσατο λέγων· Οὐκ οἶδα αὐτόν, γύναι.
\vs{58}Καὶ μετὰ βραχὺ ἕτερος ἰδὼν αὐτὸν ἔφη· Καὶ σὺ ἐξ αὐτῶν εἶ. Ὁ δὲ Πέτρος ἔφη· Ἄνθρωπε, οὐκ εἰμί.
\vs{59}Καὶ διαστάσης ὡσεὶ ὥρας μιᾶς ἄλλος τις διϊσχυρίζετο λέγων· Ἐπ᾽ ἀληθείας καὶ οὗτος μετ᾽ αὐτοῦ ἦν, καὶ γὰρ Γαλιλαῖός ἐστιν.
\vs{60}Εἶπεν δὲ ὁ Πέτρος· Ἄνθρωπε, οὐκ οἶδα ὃ λέγεις. καὶ παραχρῆμα ἔτι λαλοῦντος αὐτοῦ ἐφώνησεν ἀλέκτωρ.
\vs{61}καὶ στραφεὶς ὁ Κύριος ἐνέβλεψεν τῷ Πέτρῳ, καὶ ὑπεμνήσθη ὁ Πέτρος τοῦ ῥήματος τοῦ Κυρίου ὡς εἶπεν αὐτῷ ὅτι Πρὶν ἀλέκτορα φωνῆσαι σήμερον ἀπαρνήσῃ με τρίς.
\vs{62}καὶ ἐξελθὼν ἔξω ἔκλαυσεν πικρῶς.

\vs{63}Καὶ οἱ ἄνδρες οἱ συνέχοντες αὐτὸν ἐνέπαιζον αὐτῷ δέροντες,
\vs{64}καὶ περικαλύψαντες αὐτὸν ἐπηρώτων λέγοντες· Προφήτευσον, τίς ἐστιν ὁ παίσας σε;
\vs{65}καὶ ἕτερα πολλὰ βλασφημοῦντες ἔλεγον εἰς αὐτόν.

\vs{66}Καὶ ὡς ἐγένετο ἡμέρα, συνήχθη τὸ πρεσβυτέριον τοῦ λαοῦ, ἀρχιερεῖς τε καὶ γραμματεῖς, καὶ ἀπήγαγον αὐτὸν εἰς τὸ συνέδριον αὐτῶν
\vs{67}λέγοντες· Εἰ σὺ εἶ ὁ Χριστός, εἰπὸν ἡμῖν. Εἶπεν δὲ αὐτοῖς· Ἐὰν ὑμῖν εἴπω, οὐ μὴ πιστεύσητε·
\vs{68}ἐὰν δὲ ἐρωτήσω, οὐ μὴ ἀποκριθῆτε.
\vs{69}ἀπὸ τοῦ νῦν δὲ ἔσται ὁ Υἱὸς τοῦ ἀνθρώπου καθήμενος ἐκ δεξιῶν τῆς δυνάμεως τοῦ Θεοῦ.
\vs{70}Εἶπαν δὲ πάντες· Σὺ οὖν εἶ ὁ Υἱὸς τοῦ Θεοῦ; Ὁ δὲ πρὸς αὐτοὺς ἔφη· Ὑμεῖς λέγετε ὅτι ἐγώ εἰμι.
\vs{71}Οἱ δὲ εἶπαν· Τί ἔτι ἔχομεν μαρτυρίας χρείαν; αὐτοὶ γὰρ ἠκούσαμεν ἀπὸ τοῦ στόματος αὐτοῦ.

\ch{23}
Καὶ ἀναστὰν ἅπαν τὸ πλῆθος αὐτῶν ἤγαγον αὐτὸν ἐπὶ τὸν Πιλᾶτον.

\vs{2}ἤρξαντο δὲ κατηγορεῖν αὐτοῦ λέγοντες· Τοῦτον εὕραμεν διαστρέφοντα τὸ ἔθνος ἡμῶν καὶ κωλύοντα φόρους Καίσαρι διδόναι καὶ λέγοντα ἑαυτὸν Χριστὸν βασιλέα εἶναι.
\vs{3}Ὁ δὲ Πιλᾶτος ἠρώτησεν αὐτὸν λέγων· Σὺ εἶ ὁ Βασιλεὺς τῶν Ἰουδαίων; Ὁ δὲ ἀποκριθεὶς αὐτῷ ἔφη· Σὺ λέγεις.
\vs{4}Ὁ δὲ Πιλᾶτος εἶπεν πρὸς τοὺς ἀρχιερεῖς καὶ τοὺς ὄχλους· Οὐδὲν εὑρίσκω αἴτιον ἐν τῷ ἀνθρώπῳ τούτῳ.
\vs{5}Οἱ δὲ ἐπίσχυον λέγοντες ὅτι Ἀνασείει τὸν λαὸν διδάσκων καθ᾽ ὅλης τῆς Ἰουδαίας, καὶ ἀρξάμενος ἀπὸ τῆς Γαλιλαίας ἕως ὧδε.

\vs{6}Πιλᾶτος δὲ ἀκούσας ἐπηρώτησεν εἰ ὁ ἄνθρωπος Γαλιλαῖός ἐστιν,
\vs{7}καὶ ἐπιγνοὺς ὅτι ἐκ τῆς ἐξουσίας Ἡρῴδου ἐστὶν ἀνέπεμψεν αὐτὸν πρὸς Ἡρῴδην, ὄντα καὶ αὐτὸν ἐν Ἱεροσολύμοις ἐν ταύταις ταῖς ἡμέραις.

\vs{8}Ὁ δὲ Ἡρῴδης ἰδὼν τὸν Ἰησοῦν ἐχάρη λίαν, ἦν γὰρ ἐξ ἱκανῶν χρόνων θέλων ἰδεῖν αὐτὸν διὰ τὸ ἀκούειν περὶ αὐτοῦ καὶ ἤλπιζέν τι σημεῖον ἰδεῖν ὑπ᾽ αὐτοῦ γινόμενον.
\vs{9}ἐπηρώτα δὲ αὐτὸν ἐν λόγοις ἱκανοῖς, αὐτὸς δὲ οὐδὲν ἀπεκρίνατο αὐτῷ.
\vs{10}Εἱστήκεισαν δὲ οἱ ἀρχιερεῖς καὶ οἱ γραμματεῖς εὐτόνως κατηγοροῦντες αὐτοῦ.
\vs{11}ἐξουθενήσας δὲ αὐτὸν καὶ ὁ Ἡρῴδης σὺν τοῖς στρατεύμασιν αὐτοῦ καὶ ἐμπαίξας περιβαλὼν ἐσθῆτα λαμπρὰν ἀνέπεμψεν αὐτὸν τῷ Πιλάτῳ.
\vs{12}Ἐγένοντο δὲ φίλοι ὅ τε Ἡρῴδης καὶ ὁ Πιλᾶτος ἐν αὐτῇ τῇ ἡμέρᾳ μετ᾽ ἀλλήλων· προϋπῆρχον γὰρ ἐν ἔχθρᾳ ὄντες πρὸς αὑτούς.

\vs{13}Πιλᾶτος δὲ συνκαλεσάμενος τοὺς ἀρχιερεῖς καὶ τοὺς ἄρχοντας καὶ τὸν λαὸν
\vs{14}εἶπεν πρὸς αὐτούς· Προσηνέγκατέ μοι τὸν ἄνθρωπον τοῦτον ὡς ἀποστρέφοντα τὸν λαόν, καὶ ἰδοὺ ἐγὼ ἐνώπιον ὑμῶν ἀνακρίνας οὐθὲν εὗρον ἐν τῷ ἀνθρώπῳ τούτῳ αἴτιον ὧν κατηγορεῖτε κατ᾽ αὐτοῦ.
\vs{15}ἀλλ᾽ οὐδὲ Ἡρῴδης, ἀνέπεμψεν γὰρ αὐτὸν πρὸς ἡμᾶς, καὶ ἰδοὺ οὐδὲν ἄξιον θανάτου ἐστὶν πεπραγμένον αὐτῷ·
\vs{16}παιδεύσας οὖν αὐτὸν ἀπολύσω.

\vs{18}Ἀνέκραγον δὲ παμπληθεὶ λέγοντες· Αἶρε τοῦτον, ἀπόλυσον δὲ ἡμῖν τὸν Βαραββᾶν·
\vs{19}ὅστις ἦν διὰ στάσιν τινὰ γενομένην ἐν τῇ πόλει καὶ φόνον βληθεὶς ἐν τῇ φυλακῇ.
\vs{20}Πάλιν δὲ ὁ Πιλᾶτος προσεφώνησεν αὐτοῖς θέλων ἀπολῦσαι τὸν Ἰησοῦν.
\vs{21}οἱ δὲ ἐπεφώνουν λέγοντες· Σταύρου σταύρου αὐτόν.
\vs{22}Ὁ δὲ τρίτον εἶπεν πρὸς αὐτούς· Τί γὰρ κακὸν ἐποίησεν οὗτος; οὐδὲν αἴτιον θανάτου εὗρον ἐν αὐτῷ· παιδεύσας οὖν αὐτὸν ἀπολύσω.
\vs{23}Οἱ δὲ ἐπέκειντο φωναῖς μεγάλαις αἰτούμενοι αὐτὸν σταυρωθῆναι, καὶ κατίσχυον αἱ φωναὶ αὐτῶν.
\vs{24}καὶ Πιλᾶτος ἐπέκρινεν γενέσθαι τὸ αἴτημα αὐτῶν·
\vs{25}ἀπέλυσεν δὲ τὸν διὰ στάσιν καὶ φόνον βεβλημένον εἰς φυλακὴν ὃν ᾐτοῦντο, τὸν δὲ Ἰησοῦν παρέδωκεν τῷ θελήματι αὐτῶν.

\vs{26}Καὶ ὡς ἀπήγαγον αὐτόν, ἐπιλαβόμενοι Σίμωνά τινα Κυρηναῖον ἐρχόμενον ἀπ᾽ ἀγροῦ ἐπέθηκαν αὐτῷ τὸν σταυρὸν φέρειν ὄπισθεν τοῦ Ἰησοῦ.

\vs{27}Ἠκολούθει δὲ αὐτῷ πολὺ πλῆθος τοῦ λαοῦ καὶ γυναικῶν αἳ ἐκόπτοντο καὶ ἐθρήνουν αὐτόν.
\vs{28}στραφεὶς δὲ πρὸς αὐτὰς ὁ Ἰησοῦς εἶπεν· Θυγατέρες Ἰερουσαλήμ, μὴ κλαίετε ἐπ᾽ ἐμέ· πλὴν ἐφ᾽ ἑαυτὰς κλαίετε καὶ ἐπὶ τὰ τέκνα ὑμῶν,
\vs{29}ὅτι ἰδοὺ ἔρχονται ἡμέραι ἐν αἷς ἐροῦσιν· Μακάριαι αἱ στεῖραι καὶ αἱ κοιλίαι αἳ οὐκ ἐγέννησαν καὶ μαστοὶ οἳ οὐκ ἔθρεψαν.
\vs{30}τότε Ἄρξονται λέγειν τοῖς ὄρεσιν· Πέσετε ἐφ᾽ ἡμᾶς, Καὶ τοῖς βουνοῖς· Καλύψατε ἡμᾶς·
\vs{31}Ὅτι εἰ ἐν τῷ ὑγρῷ ξύλῳ ταῦτα ποιοῦσιν, ἐν τῷ ξηρῷ τί γένηται;

\vs{32}Ἤγοντο δὲ καὶ ἕτεροι κακοῦργοι δύο σὺν αὐτῷ ἀναιρεθῆναι.

\vs{33}Καὶ ὅτε ἦλθον ἐπὶ τὸν τόπον τὸν καλούμενον Κρανίον, ἐκεῖ ἐσταύρωσαν αὐτὸν καὶ τοὺς κακούργους, ὃν μὲν ἐκ δεξιῶν ὃν δὲ ἐξ ἀριστερῶν.
\vs{34}Ὁ δὲ Ἰησοῦς ἔλεγεν· Πάτερ, ἄφες αὐτοῖς, οὐ γὰρ οἴδασιν τί ποιοῦσιν. διαμεριζόμενοι δὲ τὰ ἱμάτια αὐτοῦ ἔβαλον κλήρους.
\vs{35}Καὶ εἱστήκει ὁ λαὸς θεωρῶν. ἐξεμυκτήριζον δὲ καὶ οἱ ἄρχοντες λέγοντες· Ἄλλους ἔσωσεν, σωσάτω ἑαυτόν, εἰ οὗτός ἐστιν ὁ Χριστὸς τοῦ Θεοῦ ὁ ἐκλεκτός.
\vs{36}Ἐνέπαιξαν δὲ αὐτῷ καὶ οἱ στρατιῶται προσερχόμενοι, ὄξος προσφέροντες αὐτῷ
\vs{37}καὶ λέγοντες· Εἰ σὺ εἶ ὁ Βασιλεὺς τῶν Ἰουδαίων, σῶσον σεαυτόν.
\vs{38}Ἦν δὲ καὶ ἐπιγραφὴ ἐπ᾽ αὐτῷ· 
\begin{poetryblock}

\begin{quote}Ο ΒΑΣΙΛΕΥΣ ΤΩΝ ΙΟΥΔΑΙΩΝ ΟΥΤΟΣ.\end{quote}
\end{poetryblock}

\vs{39}Εἷς δὲ τῶν κρεμασθέντων κακούργων ἐβλασφήμει αὐτόν λέγων· Οὐχὶ σὺ εἶ ὁ Χριστός; σῶσον σεαυτὸν καὶ ἡμᾶς.
\vs{40}Ἀποκριθεὶς δὲ ὁ ἕτερος ἐπιτιμῶν αὐτῷ ἔφη· Οὐδὲ φοβῇ σὺ τὸν Θεόν, ὅτι ἐν τῷ αὐτῷ κρίματι εἶ;
\vs{41}καὶ ἡμεῖς μὲν δικαίως, ἄξια γὰρ ὧν ἐπράξαμεν ἀπολαμβάνομεν· οὗτος δὲ οὐδὲν ἄτοπον ἔπραξεν.
\vs{42}καὶ ἔλεγεν· Ἰησοῦ, μνήσθητί μου ὅταν ἔλθῃς εἰς τὴν βασιλείαν σου.
\vs{43}Καὶ εἶπεν αὐτῷ· Ἀμήν σοι λέγω, σήμερον μετ᾽ ἐμοῦ ἔσῃ ἐν τῷ Παραδείσῳ.

\vs{44}Καὶ ἦν ἤδη ὡσεὶ ὥρα ἕκτη καὶ σκότος ἐγένετο ἐφ᾽ ὅλην τὴν γῆν ἕως ὥρας ἐνάτης
\vs{45}τοῦ ἡλίου ἐκλιπόντος, ἐσχίσθη δὲ τὸ καταπέτασμα τοῦ ναοῦ μέσον.
\vs{46}Καὶ φωνήσας φωνῇ μεγάλῃ ὁ Ἰησοῦς εἶπεν· Πάτερ, εἰς χεῖράς σου παρατίθεμαι τὸ πνεῦμά μου. τοῦτο δὲ εἰπὼν ἐξέπνευσεν.

\vs{47}Ἰδὼν δὲ ὁ ἑκατοντάρχης τὸ γενόμενον ἐδόξαζεν τὸν Θεὸν λέγων· Ὄντως ὁ ἄνθρωπος οὗτος δίκαιος ἦν.
\vs{48}καὶ πάντες οἱ συμπαραγενόμενοι ὄχλοι ἐπὶ τὴν θεωρίαν ταύτην, θεωρήσαντες τὰ γενόμενα, τύπτοντες τὰ στήθη ὑπέστρεφον.
\vs{49}εἱστήκεισαν δὲ πάντες οἱ γνωστοὶ αὐτῷ ἀπὸ μακρόθεν καὶ γυναῖκες αἱ συνακολουθοῦσαι αὐτῷ ἀπὸ τῆς Γαλιλαίας ὁρῶσαι ταῦτα.

\vs{50}Καὶ ἰδοὺ ἀνὴρ ὀνόματι Ἰωσὴφ βουλευτὴς ὑπάρχων καὶ ἀνὴρ ἀγαθὸς καὶ δίκαιος—
\vs{51}οὗτος οὐκ ἦν συνκατατεθειμένος τῇ βουλῇ καὶ τῇ πράξει αὐτῶν— ἀπὸ Ἁριμαθαίας πόλεως τῶν Ἰουδαίων, ὃς προσεδέχετο τὴν βασιλείαν τοῦ Θεοῦ,
\vs{52}οὗτος προσελθὼν τῷ Πιλάτῳ ᾐτήσατο τὸ σῶμα τοῦ Ἰησοῦ
\vs{53}καὶ καθελὼν ἐνετύλιξεν αὐτὸ σινδόνι καὶ ἔθηκεν αὐτὸν ἐν μνήματι λαξευτῷ οὗ οὐκ ἦν οὐδεὶς οὔπω κείμενος.
\vs{54}καὶ ἡμέρα ἦν Παρασκευῆς καὶ σάββατον ἐπέφωσκεν.

\vs{55}Κατακολουθήσασαι δὲ αἱ γυναῖκες, αἵτινες ἦσαν συνεληλυθυῖαι ἐκ τῆς Γαλιλαίας αὐτῷ, ἐθεάσαντο τὸ μνημεῖον καὶ ὡς ἐτέθη τὸ σῶμα αὐτοῦ,
\vs{56}ὑποστρέψασαι δὲ ἡτοίμασαν ἀρώματα καὶ μύρα. Καὶ τὸ μὲν σάββατον ἡσύχασαν κατὰ τὴν ἐντολήν.

\ch{24}
Τῇ δὲ μιᾷ τῶν σαββάτων ὄρθρου βαθέως ἐπὶ τὸ μνῆμα ἦλθον φέρουσαι ἃ ἡτοίμασαν ἀρώματα.
\vs{2}εὗρον δὲ τὸν λίθον ἀποκεκυλισμένον ἀπὸ τοῦ μνημείου,
\vs{3}εἰσελθοῦσαι δὲ οὐχ εὗρον τὸ σῶμα τοῦ Κυρίου Ἰησοῦ.
\vs{4}καὶ ἐγένετο ἐν τῷ ἀπορεῖσθαι αὐτὰς περὶ τούτου καὶ ἰδοὺ ἄνδρες δύο ἐπέστησαν αὐταῖς ἐν ἐσθῆτι ἀστραπτούσῃ.
\vs{5}ἐμφόβων δὲ γενομένων αὐτῶν καὶ κλινουσῶν τὰ πρόσωπα εἰς τὴν γῆν εἶπαν πρὸς αὐτάς· Τί ζητεῖτε τὸν ζῶντα μετὰ τῶν νεκρῶν;
\vs{6}οὐκ ἔστιν ὧδε, ἀλλὰ ἠγέρθη. μνήσθητε ὡς ἐλάλησεν ὑμῖν ἔτι ὢν ἐν τῇ Γαλιλαίᾳ
\vs{7}λέγων Τὸν Υἱὸν τοῦ ἀνθρώπου ὅτι δεῖ παραδοθῆναι εἰς χεῖρας ἀνθρώπων ἁμαρτωλῶν καὶ σταυρωθῆναι καὶ τῇ τρίτῃ ἡμέρᾳ ἀναστῆναι.
\vs{8}Καὶ ἐμνήσθησαν τῶν ῥημάτων αὐτοῦ.

\vs{9}καὶ ὑποστρέψασαι ἀπὸ τοῦ μνημείου ἀπήγγειλαν ταῦτα πάντα τοῖς ἕνδεκα καὶ πᾶσιν τοῖς λοιποῖς.
\vs{10}ἦσαν δὲ ἡ Μαγδαληνὴ Μαρία καὶ Ἰωάννα καὶ Μαρία ἡ Ἰακώβου καὶ αἱ λοιπαὶ σὺν αὐταῖς. ἔλεγον πρὸς τοὺς ἀποστόλους ταῦτα,
\vs{11}καὶ ἐφάνησαν ἐνώπιον αὐτῶν ὡσεὶ λῆρος τὰ ῥήματα ταῦτα, καὶ ἠπίστουν αὐταῖς.
\vs{12}Ὁ δὲ Πέτρος ἀναστὰς ἔδραμεν ἐπὶ τὸ μνημεῖον καὶ παρακύψας βλέπει τὰ ὀθόνια μόνα, καὶ ἀπῆλθεν πρὸς ἑαυτὸν θαυμάζων τὸ γεγονός.

\vs{13}Καὶ ἰδοὺ δύο ἐξ αὐτῶν ἐν αὐτῇ τῇ ἡμέρᾳ ἦσαν πορευόμενοι εἰς κώμην ἀπέχουσαν σταδίους ἑξήκοντα ἀπὸ Ἰερουσαλήμ, ᾗ ὄνομα Ἐμμαοῦς,
\vs{14}καὶ αὐτοὶ ὡμίλουν πρὸς ἀλλήλους περὶ πάντων τῶν συμβεβηκότων τούτων.
\vs{15}καὶ ἐγένετο ἐν τῷ ὁμιλεῖν αὐτοὺς καὶ συζητεῖν καὶ αὐτὸς Ἰησοῦς ἐγγίσας συνεπορεύετο αὐτοῖς,
\vs{16}οἱ δὲ ὀφθαλμοὶ αὐτῶν ἐκρατοῦντο τοῦ μὴ ἐπιγνῶναι αὐτόν.
\vs{17}Εἶπεν δὲ πρὸς αὐτούς· Τίνες οἱ λόγοι οὗτοι οὓς ἀντιβάλλετε πρὸς ἀλλήλους περιπατοῦντες; Καὶ ἐστάθησαν σκυθρωποί.
\vs{18}ἀποκριθεὶς δὲ εἷς ὀνόματι Κλεοπᾶς εἶπεν πρὸς αὐτόν· Σὺ μόνος παροικεῖς Ἰερουσαλὴμ καὶ οὐκ ἔγνως τὰ γενόμενα ἐν αὐτῇ ἐν ταῖς ἡμέραις ταύταις;
\vs{19}Καὶ εἶπεν αὐτοῖς· Ποῖα; οἱ Δὲ εἶπαν αὐτῷ· Τὰ περὶ Ἰησοῦ τοῦ Ναζαρηνοῦ, ὃς ἐγένετο ἀνὴρ προφήτης δυνατὸς ἐν ἔργῳ καὶ λόγῳ ἐναντίον τοῦ Θεοῦ καὶ παντὸς τοῦ λαοῦ,
\vs{20}ὅπως τε παρέδωκαν αὐτὸν οἱ ἀρχιερεῖς καὶ οἱ ἄρχοντες ἡμῶν εἰς κρίμα θανάτου καὶ ἐσταύρωσαν αὐτόν.
\vs{21}ἡμεῖς δὲ ἠλπίζομεν ὅτι αὐτός ἐστιν ὁ μέλλων λυτροῦσθαι τὸν Ἰσραήλ· ἀλλά γε καὶ σὺν πᾶσιν τούτοις τρίτην ταύτην ἡμέραν ἄγει ἀφ᾽ οὗ ταῦτα ἐγένετο.
\vs{22}Ἀλλὰ καὶ γυναῖκές τινες ἐξ ἡμῶν ἐξέστησαν ἡμᾶς, γενόμεναι ὀρθριναὶ ἐπὶ τὸ μνημεῖον,
\vs{23}καὶ μὴ εὑροῦσαι τὸ σῶμα αὐτοῦ ἦλθον λέγουσαι καὶ ὀπτασίαν ἀγγέλων ἑωρακέναι, οἳ λέγουσιν αὐτὸν ζῆν.
\vs{24}καὶ ἀπῆλθόν τινες τῶν σὺν ἡμῖν ἐπὶ τὸ μνημεῖον καὶ εὗρον οὕτως καθὼς καὶ αἱ γυναῖκες εἶπον, αὐτὸν δὲ οὐκ εἶδον.
\vs{25}Καὶ αὐτὸς εἶπεν πρὸς αὐτούς· Ὦ ἀνόητοι καὶ βραδεῖς τῇ καρδίᾳ τοῦ πιστεύειν ἐπὶ πᾶσιν οἷς ἐλάλησαν οἱ προφῆται·
\vs{26}οὐχὶ ταῦτα ἔδει παθεῖν τὸν Χριστὸν καὶ εἰσελθεῖν εἰς τὴν δόξαν αὐτοῦ;
\vs{27}καὶ ἀρξάμενος ἀπὸ Μωϋσέως καὶ ἀπὸ πάντων τῶν προφητῶν διερμήνευσεν αὐτοῖς ἐν πάσαις ταῖς γραφαῖς τὰ περὶ ἑαυτοῦ.

\vs{28}Καὶ ἤγγισαν εἰς τὴν κώμην οὗ ἐπορεύοντο, καὶ αὐτὸς προσεποιήσατο πορρώτερον πορεύεσθαι.
\vs{29}καὶ παρεβιάσαντο αὐτὸν λέγοντες· Μεῖνον μεθ᾽ ἡμῶν, ὅτι πρὸς ἑσπέραν ἐστὶν καὶ κέκλικεν ἤδη ἡ ἡμέρα. Καὶ εἰσῆλθεν τοῦ μεῖναι σὺν αὐτοῖς.
\vs{30}καὶ ἐγένετο ἐν τῷ κατακλιθῆναι αὐτὸν μετ᾽ αὐτῶν λαβὼν τὸν ἄρτον εὐλόγησεν καὶ κλάσας ἐπεδίδου αὐτοῖς,
\vs{31}αὐτῶν δὲ διηνοίχθησαν οἱ ὀφθαλμοὶ καὶ ἐπέγνωσαν αὐτόν· καὶ αὐτὸς ἄφαντος ἐγένετο ἀπ᾽ αὐτῶν.
\vs{32}Καὶ εἶπαν πρὸς ἀλλήλους· Οὐχὶ ἡ καρδία ἡμῶν καιομένη ἦν ἐν ἡμῖν ὡς ἐλάλει ἡμῖν ἐν τῇ ὁδῷ, ὡς διήνοιγεν ἡμῖν τὰς γραφάς;

\vs{33}Καὶ ἀναστάντες αὐτῇ τῇ ὥρᾳ ὑπέστρεψαν εἰς Ἰερουσαλήμ καὶ εὗρον ἠθροισμένους τοὺς ἕνδεκα καὶ τοὺς σὺν αὐτοῖς,
\vs{34}λέγοντας ὅτι Ὄντως ἠγέρθη ὁ Κύριος καὶ ὤφθη Σίμωνι.
\vs{35}Καὶ αὐτοὶ ἐξηγοῦντο τὰ ἐν τῇ ὁδῷ καὶ ὡς ἐγνώσθη αὐτοῖς ἐν τῇ κλάσει τοῦ ἄρτου.

\vs{36}Ταῦτα δὲ αὐτῶν λαλούντων αὐτὸς ἔστη ἐν μέσῳ αὐτῶν καὶ λέγει αὐτοῖς· Εἰρήνη ὑμῖν.
\vs{37}πτοηθέντες δὲ καὶ ἔμφοβοι γενόμενοι ἐδόκουν πνεῦμα θεωρεῖν.
\vs{38}Καὶ εἶπεν αὐτοῖς· Τί τεταραγμένοι ἐστέ καὶ διὰ τί διαλογισμοὶ ἀναβαίνουσιν ἐν τῇ καρδίᾳ ὑμῶν;
\vs{39}ἴδετε τὰς χεῖράς μου καὶ τοὺς πόδας μου ὅτι ἐγώ εἰμι αὐτός· ψηλαφήσατέ με καὶ ἴδετε, ὅτι πνεῦμα σάρκα καὶ ὀστέα οὐκ ἔχει καθὼς ἐμὲ θεωρεῖτε ἔχοντα.
\vs{40}καὶ τοῦτο εἰπὼν ἔδειξεν αὐτοῖς τὰς χεῖρας καὶ τοὺς πόδας.
\vs{41}Ἔτι δὲ ἀπιστούντων αὐτῶν ἀπὸ τῆς χαρᾶς καὶ θαυμαζόντων εἶπεν αὐτοῖς· Ἔχετέ τι βρώσιμον ἐνθάδε;
\vs{42}οἱ δὲ ἐπέδωκαν αὐτῷ ἰχθύος ὀπτοῦ μέρος·
\vs{43}καὶ λαβὼν ἐνώπιον αὐτῶν ἔφαγεν.

\vs{44}Εἶπεν δὲ πρὸς αὐτούς· Οὗτοι οἱ λόγοι μου οὓς ἐλάλησα πρὸς ὑμᾶς ἔτι ὢν σὺν ὑμῖν, ὅτι δεῖ πληρωθῆναι πάντα τὰ γεγραμμένα ἐν τῷ νόμῳ Μωϋσέως καὶ τοῖς προφήταις καὶ ψαλμοῖς περὶ ἐμοῦ.
\vs{45}τότε διήνοιξεν αὐτῶν τὸν νοῦν τοῦ συνιέναι τὰς γραφάς·
\vs{46}Καὶ εἶπεν αὐτοῖς ὅτι Οὕτως γέγραπται παθεῖν τὸν Χριστὸν καὶ ἀναστῆναι ἐκ νεκρῶν τῇ τρίτῃ ἡμέρᾳ,
\vs{47}καὶ κηρυχθῆναι ἐπὶ τῷ ὀνόματι αὐτοῦ μετάνοιαν εἰς ἄφεσιν ἁμαρτιῶν εἰς πάντα τὰ ἔθνη. ἀρξάμενοι ἀπὸ Ἰερουσαλήμ
\vs{48}ὑμεῖς μάρτυρες τούτων.
\vs{49}Καὶ ἰδοὺ ἐγὼ ἀποστέλλω τὴν ἐπαγγελίαν τοῦ Πατρός μου ἐφ᾽ ὑμᾶς· ὑμεῖς δὲ καθίσατε ἐν τῇ πόλει ἕως οὗ ἐνδύσησθε ἐξ ὕψους δύναμιν.

\vs{50}Ἐξήγαγεν δὲ αὐτοὺς ἔξω ἕως πρὸς Βηθανίαν, καὶ ἐπάρας τὰς χεῖρας αὐτοῦ εὐλόγησεν αὐτούς.
\vs{51}καὶ ἐγένετο ἐν τῷ εὐλογεῖν αὐτὸν αὐτοὺς διέστη ἀπ᾽ αὐτῶν καὶ ἀνεφέρετο εἰς τὸν οὐρανόν.

\vs{52}καὶ αὐτοὶ προσκυνήσαντες αὐτὸν ὑπέστρεψαν εἰς Ἰερουσαλὴμ μετὰ χαρᾶς μεγάλης
\vs{53}καὶ ἦσαν διὰ παντὸς ἐν τῷ ἱερῷ εὐλογοῦντες τὸν Θεόν.


\def\book{ΚΑΤΑ ΙΩΑΝΝΗΝ}
\biblebook{ΚΑΤΑ ΙΩΑΝΝΗΝ}


\lettrine[lines=2, loversize=0.2, nindent=0em, findent=.25em]{\textcolor{bookheadingcolor}{Ἐ}}{ν} ἀρχῇ ἦν ὁ Λόγος, καὶ ὁ Λόγος ἦν πρὸς τὸν Θεόν, καὶ Θεὸς ἦν ὁ Λόγος.
\vs{2}Οὗτος ἦν ἐν ἀρχῇ πρὸς τὸν Θεόν.
\vs{3}πάντα δι᾽ αὐτοῦ ἐγένετο, καὶ χωρὶς αὐτοῦ ἐγένετο οὐδὲ ἕν. ὃ γέγονεν
\vs{4}ἐν αὐτῷ ζωὴ ἦν, καὶ ἡ ζωὴ ἦν τὸ φῶς τῶν ἀνθρώπων·
\vs{5}καὶ τὸ φῶς ἐν τῇ σκοτίᾳ φαίνει, καὶ ἡ σκοτία αὐτὸ οὐ κατέλαβεν.

\vs{6}Ἐγένετο ἄνθρωπος, ἀπεσταλμένος παρὰ Θεοῦ, ὄνομα αὐτῷ Ἰωάννης·
\vs{7}οὗτος ἦλθεν εἰς μαρτυρίαν ἵνα μαρτυρήσῃ περὶ τοῦ φωτός, ἵνα πάντες πιστεύσωσιν δι᾽ αὐτοῦ.
\vs{8}οὐκ ἦν ἐκεῖνος τὸ φῶς, ἀλλ᾽ ἵνα μαρτυρήσῃ περὶ τοῦ φωτός.

\vs{9}Ἦν τὸ φῶς τὸ ἀληθινὸν, ὃ φωτίζει πάντα ἄνθρωπον, ἐρχόμενον εἰς τὸν κόσμον.
\vs{10}ἐν τῷ κόσμῳ ἦν, καὶ ὁ κόσμος δι᾽ αὐτοῦ ἐγένετο, καὶ ὁ κόσμος αὐτὸν οὐκ ἔγνω.
\vs{11}εἰς τὰ ἴδια ἦλθεν, καὶ οἱ ἴδιοι αὐτὸν οὐ παρέλαβον.
\vs{12}ὅσοι δὲ ἔλαβον αὐτόν, ἔδωκεν αὐτοῖς ἐξουσίαν τέκνα Θεοῦ γενέσθαι, τοῖς πιστεύουσιν εἰς τὸ ὄνομα αὐτοῦ,
\vs{13}οἳ οὐκ ἐξ αἱμάτων οὐδὲ ἐκ θελήματος σαρκὸς οὐδὲ ἐκ θελήματος ἀνδρὸς ἀλλ᾽ ἐκ Θεοῦ ἐγεννήθησαν.

\vs{14}Καὶ ὁ Λόγος σὰρξ ἐγένετο καὶ ἐσκήνωσεν ἐν ἡμῖν, καὶ ἐθεασάμεθα τὴν δόξαν αὐτοῦ, δόξαν ὡς μονογενοῦς παρὰ Πατρός, πλήρης χάριτος καὶ ἀληθείας.
\vs{15}Ἰωάννης μαρτυρεῖ περὶ αὐτοῦ καὶ κέκραγεν λέγων· Οὗτος ἦν ὃν εἶπον· Ὁ ὀπίσω μου ἐρχόμενος ἔμπροσθέν μου γέγονεν, ὅτι πρῶτός μου ἦν.
\vs{16}Ὅτι ἐκ τοῦ πληρώματος αὐτοῦ ἡμεῖς πάντες ἐλάβομεν καὶ χάριν ἀντὶ χάριτος·
\vs{17}ὅτι ὁ νόμος διὰ Μωϋσέως ἐδόθη, ἡ χάρις καὶ ἡ ἀλήθεια διὰ Ἰησοῦ Χριστοῦ ἐγένετο.
\vs{18}Θεὸν οὐδεὶς ἑώρακεν πώποτε· μονογενὴς Θεὸς ὁ ὢν εἰς τὸν κόλπον τοῦ Πατρὸς ἐκεῖνος ἐξηγήσατο.
\vs{19}Καὶ αὕτη ἐστὶν ἡ μαρτυρία τοῦ Ἰωάννου, ὅτε ἀπέστειλαν πρὸς αὐτὸν οἱ Ἰουδαῖοι ἐξ Ἱεροσολύμων ἱερεῖς καὶ Λευίτας ἵνα ἐρωτήσωσιν αὐτόν· Σὺ τίς εἶ;
\vs{20}καὶ ὡμολόγησεν καὶ οὐκ ἠρνήσατο, καὶ ὡμολόγησεν ὅτι Ἐγὼ οὐκ εἰμὶ ὁ Χριστός.
\vs{21}Καὶ ἠρώτησαν αὐτόν· Τί οὖν; σὺ Ἠλίας εἶ; Καὶ λέγει· Οὐκ εἰμί. Ὁ προφήτης εἶ σύ; Καὶ ἀπεκρίθη· Οὔ.
\vs{22}Εἶπαν οὖν αὐτῷ· Τίς εἶ; ἵνα ἀπόκρισιν δῶμεν τοῖς πέμψασιν ἡμᾶς· τί λέγεις περὶ σεαυτοῦ;

\vs{23}Ἔφη· 
\begin{poetryblock}

\begin{quote}Ἐγὼ φωνὴ βοῶντος ἐν τῇ ἐρήμῳ·\end{quote} 

\begin{quote}Εὐθύνατε τὴν ὁδὸν Κυρίου,\end{quote}
\end{poetryblock}

καθὼς εἶπεν Ἠσαΐας ὁ προφήτης.

\vs{24}Καὶ ἀπεσταλμένοι ἦσαν ἐκ τῶν Φαρισαίων.
\vs{25}καὶ ἠρώτησαν αὐτὸν καὶ εἶπαν αὐτῷ· Τί οὖν βαπτίζεις εἰ σὺ οὐκ εἶ ὁ Χριστὸς οὐδὲ Ἠλίας οὐδὲ ὁ προφήτης;
\vs{26}Ἀπεκρίθη αὐτοῖς ὁ Ἰωάννης λέγων· Ἐγὼ βαπτίζω ἐν ὕδατι· μέσος ὑμῶν ἕστηκεν ὃν ὑμεῖς οὐκ οἴδατε,
\vs{27}ὁ ὀπίσω μου ἐρχόμενος, οὗ οὐκ εἰμὶ ἐγὼ ἄξιος ἵνα λύσω αὐτοῦ τὸν ἱμάντα τοῦ ὑποδήματος.
\vs{28}Ταῦτα ἐν Βηθανίᾳ ἐγένετο πέραν τοῦ Ἰορδάνου, ὅπου ἦν ὁ Ἰωάννης βαπτίζων.

\vs{29}Τῇ ἐπαύριον βλέπει τὸν Ἰησοῦν ἐρχόμενον πρὸς αὐτόν καὶ λέγει· Ἴδε ὁ Ἀμνὸς τοῦ Θεοῦ ὁ αἴρων τὴν ἁμαρτίαν τοῦ κόσμου.
\vs{30}οὗτός ἐστιν ὑπὲρ οὗ ἐγὼ εἶπον· Ὀπίσω μου ἔρχεται ἀνὴρ ὃς ἔμπροσθέν μου γέγονεν, ὅτι πρῶτός μου ἦν.
\vs{31}κἀγὼ οὐκ ᾔδειν αὐτόν, ἀλλ᾽ ἵνα φανερωθῇ τῷ Ἰσραὴλ διὰ τοῦτο ἦλθον ἐγὼ ἐν ὕδατι βαπτίζων.
\vs{32}Καὶ ἐμαρτύρησεν Ἰωάννης λέγων ὅτι Τεθέαμαι τὸ Πνεῦμα καταβαῖνον ὡς περιστερὰν ἐξ οὐρανοῦ καὶ ἔμεινεν ἐπ᾽ αὐτόν.
\vs{33}κἀγὼ οὐκ ᾔδειν αὐτόν, ἀλλ᾽ ὁ πέμψας με βαπτίζειν ἐν ὕδατι ἐκεῖνός μοι εἶπεν· Ἐφ᾽ ὃν ἂν ἴδῃς τὸ Πνεῦμα καταβαῖνον καὶ μένον ἐπ᾽ αὐτόν, οὗτός ἐστιν ὁ βαπτίζων ἐν Πνεύματι Ἁγίῳ.
\vs{34}κἀγὼ ἑώρακα καὶ μεμαρτύρηκα ὅτι οὗτός ἐστιν ὁ Υἱὸς τοῦ Θεοῦ.

\vs{35}Τῇ ἐπαύριον πάλιν εἱστήκει ὁ Ἰωάννης καὶ ἐκ τῶν μαθητῶν αὐτοῦ δύο
\vs{36}καὶ ἐμβλέψας τῷ Ἰησοῦ περιπατοῦντι λέγει· Ἴδε ὁ Ἀμνὸς τοῦ Θεοῦ.
\vs{37}καὶ ἤκουσαν οἱ δύο μαθηταὶ αὐτοῦ λαλοῦντος καὶ ἠκολούθησαν τῷ Ἰησοῦ.
\vs{38}Στραφεὶς δὲ ὁ Ἰησοῦς καὶ θεασάμενος αὐτοὺς ἀκολουθοῦντας λέγει αὐτοῖς· Τί ζητεῖτε; Οἱ δὲ εἶπαν αὐτῷ· Ῥαββί, ὃ λέγεται μεθερμηνευόμενον Διδάσκαλε, Ποῦ μένεις;
\vs{39}Λέγει αὐτοῖς· Ἔρχεσθε καὶ ὄψεσθε. ἦλθαν οὖν καὶ εἶδαν ποῦ μένει καὶ παρ᾽ αὐτῷ ἔμειναν τὴν ἡμέραν ἐκείνην· ὥρα ἦν ὡς δεκάτη.
\vs{40}Ἦν Ἀνδρέας ὁ ἀδελφὸς Σίμωνος Πέτρου εἷς ἐκ τῶν δύο τῶν ἀκουσάντων παρὰ Ἰωάννου καὶ ἀκολουθησάντων αὐτῷ·
\vs{41}εὑρίσκει οὗτος πρῶτον τὸν ἀδελφὸν τὸν ἴδιον Σίμωνα καὶ λέγει αὐτῷ· Εὑρήκαμεν τὸν Μεσσίαν, ὅ ἐστιν μεθερμηνευόμενον Χριστός.
\vs{42}ἤγαγεν αὐτὸν πρὸς τὸν Ἰησοῦν. ἐμβλέψας αὐτῷ ὁ Ἰησοῦς εἶπεν· Σὺ εἶ Σίμων ὁ υἱὸς Ἰωάννου, σὺ κληθήσῃ Κηφᾶς, ὃ ἑρμηνεύεται Πέτρος.

\vs{43}Τῇ ἐπαύριον ἠθέλησεν ἐξελθεῖν εἰς τὴν Γαλιλαίαν καὶ εὑρίσκει Φίλιππον. καὶ λέγει αὐτῷ ὁ Ἰησοῦς· Ἀκολούθει μοι.
\vs{44}ἦν δὲ ὁ Φίλιππος ἀπὸ Βηθσαϊδά, ἐκ τῆς πόλεως Ἀνδρέου καὶ Πέτρου.
\vs{45}Εὑρίσκει Φίλιππος τὸν Ναθαναὴλ καὶ λέγει αὐτῷ· Ὃν ἔγραψεν Μωϋσῆς ἐν τῷ νόμῳ καὶ οἱ προφῆται εὑρήκαμεν, Ἰησοῦν υἱὸν τοῦ Ἰωσὴφ τὸν ἀπὸ Ναζαρέτ.
\vs{46}Καὶ εἶπεν αὐτῷ Ναθαναήλ· Ἐκ Ναζαρὲτ δύναταί τι ἀγαθὸν εἶναι; Λέγει αὐτῷ ὁ Φίλιππος· Ἔρχου καὶ ἴδε.
\vs{47}Εἶδεν ὁ Ἰησοῦς τὸν Ναθαναὴλ ἐρχόμενον πρὸς αὐτὸν καὶ λέγει περὶ αὐτοῦ· Ἴδε ἀληθῶς Ἰσραηλίτης ἐν ᾧ δόλος οὐκ ἔστιν.
\vs{48}Λέγει αὐτῷ Ναθαναήλ· Πόθεν με γινώσκεις; Ἀπεκρίθη Ἰησοῦς καὶ εἶπεν αὐτῷ· Πρὸ τοῦ σε Φίλιππον φωνῆσαι ὄντα ὑπὸ τὴν συκῆν εἶδόν σε.
\vs{49}Ἀπεκρίθη αὐτῷ Ναθαναήλ· Ῥαββί, σὺ εἶ ὁ Υἱὸς τοῦ Θεοῦ, σὺ Βασιλεὺς εἶ τοῦ Ἰσραήλ.
\vs{50}Ἀπεκρίθη Ἰησοῦς καὶ εἶπεν αὐτῷ· Ὅτι εἶπόν σοι ὅτι εἶδόν σε ὑποκάτω τῆς συκῆς, πιστεύεις; μείζω τούτων ὄψῃ.
\vs{51}καὶ λέγει αὐτῷ· Ἀμὴν ἀμὴν λέγω ὑμῖν, ὄψεσθε τὸν οὐρανὸν ἀνεῳγότα καὶ τοὺς ἀγγέλους τοῦ Θεοῦ ἀναβαίνοντας καὶ καταβαίνοντας ἐπὶ τὸν Υἱὸν τοῦ ἀνθρώπου.

\ch{2}
Καὶ τῇ ἡμέρᾳ τῇ τρίτῃ γάμος ἐγένετο ἐν Κανὰ τῆς Γαλιλαίας, καὶ ἦν ἡ μήτηρ τοῦ Ἰησοῦ ἐκεῖ·
\vs{2}ἐκλήθη δὲ καὶ ὁ Ἰησοῦς καὶ οἱ μαθηταὶ αὐτοῦ εἰς τὸν γάμον.
\vs{3}καὶ ὑστερήσαντος οἴνου λέγει ἡ μήτηρ τοῦ Ἰησοῦ πρὸς αὐτόν· Οἶνον οὐκ ἔχουσιν.
\vs{4}Καὶ λέγει αὐτῇ ὁ Ἰησοῦς· Τί ἐμοὶ καὶ σοί, γύναι; οὔπω ἥκει ἡ ὥρα μου.
\vs{5}Λέγει ἡ μήτηρ αὐτοῦ τοῖς διακόνοις· Ὅ τι ἂν λέγῃ ὑμῖν ποιήσατε.
\vs{6}Ἦσαν δὲ ἐκεῖ λίθιναι ὑδρίαι ἓξ κατὰ τὸν καθαρισμὸν τῶν Ἰουδαίων κείμεναι, χωροῦσαι ἀνὰ μετρητὰς δύο ἢ τρεῖς.
\vs{7}λέγει αὐτοῖς ὁ Ἰησοῦς· Γεμίσατε τὰς ὑδρίας ὕδατος. Καὶ ἐγέμισαν αὐτὰς ἕως ἄνω.
\vs{8}Καὶ λέγει αὐτοῖς· Ἀντλήσατε νῦν καὶ φέρετε τῷ ἀρχιτρικλίνῳ· Οἱ δὲ ἤνεγκαν.
\vs{9}ὡς δὲ ἐγεύσατο ὁ ἀρχιτρίκλινος τὸ ὕδωρ οἶνον γεγενημένον καὶ οὐκ ᾔδει πόθεν ἐστίν, οἱ δὲ διάκονοι ᾔδεισαν οἱ ἠντληκότες τὸ ὕδωρ, φωνεῖ τὸν νυμφίον ὁ ἀρχιτρίκλινος
\vs{10}καὶ λέγει αὐτῷ· Πᾶς ἄνθρωπος πρῶτον τὸν καλὸν οἶνον τίθησιν καὶ ὅταν μεθυσθῶσιν τὸν ἐλάσσω· σὺ τετήρηκας τὸν καλὸν οἶνον ἕως ἄρτι.
\vs{11}Ταύτην ἐποίησεν ἀρχὴν τῶν σημείων ὁ Ἰησοῦς ἐν Κανὰ τῆς Γαλιλαίας καὶ ἐφανέρωσεν τὴν δόξαν αὐτοῦ, καὶ ἐπίστευσαν εἰς αὐτὸν οἱ μαθηταὶ αὐτοῦ.

\vs{12}Μετὰ τοῦτο κατέβη εἰς Καφαρναοὺμ αὐτὸς καὶ ἡ μήτηρ αὐτοῦ καὶ οἱ ἀδελφοὶ αὐτοῦ καὶ οἱ μαθηταὶ αὐτοῦ καὶ ἐκεῖ ἔμειναν οὐ πολλὰς ἡμέρας.

\vs{13}Καὶ ἐγγὺς ἦν τὸ πάσχα τῶν Ἰουδαίων, καὶ ἀνέβη εἰς Ἱεροσόλυμα ὁ Ἰησοῦς.

\vs{14}καὶ εὗρεν ἐν τῷ ἱερῷ τοὺς πωλοῦντας βόας καὶ πρόβατα καὶ περιστερὰς καὶ τοὺς κερματιστὰς καθημένους,
\vs{15}καὶ ποιήσας φραγέλλιον ἐκ σχοινίων πάντας ἐξέβαλεν ἐκ τοῦ ἱεροῦ τά τε πρόβατα καὶ τοὺς βόας, καὶ τῶν κολλυβιστῶν ἐξέχεεν τὰ κέρματα καὶ τὰς τραπέζας ἀνέτρεψεν,
\vs{16}καὶ τοῖς τὰς περιστερὰς πωλοῦσιν εἶπεν· Ἄρατε ταῦτα ἐντεῦθεν, μὴ ποιεῖτε τὸν οἶκον τοῦ Πατρός μου οἶκον ἐμπορίου.
\vs{17}Ἐμνήσθησαν οἱ μαθηταὶ αὐτοῦ ὅτι γεγραμμένον ἐστίν· Ὁ ζῆλος τοῦ οἴκου σου καταφάγεταί με.

\vs{18}Ἀπεκρίθησαν οὖν οἱ Ἰουδαῖοι καὶ εἶπαν αὐτῷ· Τί σημεῖον δεικνύεις ἡμῖν ὅτι ταῦτα ποιεῖς;
\vs{19}Ἀπεκρίθη Ἰησοῦς καὶ εἶπεν αὐτοῖς· Λύσατε τὸν ναὸν τοῦτον καὶ ἐν τρισὶν ἡμέραις ἐγερῶ αὐτόν.
\vs{20}Εἶπαν οὖν οἱ Ἰουδαῖοι· Τεσσεράκοντα καὶ ἓξ ἔτεσιν οἰκοδομήθη ὁ ναὸς οὗτος, καὶ σὺ ἐν τρισὶν ἡμέραις ἐγερεῖς αὐτόν;
\vs{21}Ἐκεῖνος δὲ ἔλεγεν περὶ τοῦ ναοῦ τοῦ σώματος αὐτοῦ.
\vs{22}ὅτε οὖν ἠγέρθη ἐκ νεκρῶν, ἐμνήσθησαν οἱ μαθηταὶ αὐτοῦ ὅτι τοῦτο ἔλεγεν, καὶ ἐπίστευσαν τῇ γραφῇ καὶ τῷ λόγῳ ὃν εἶπεν ὁ Ἰησοῦς.

\vs{23}Ὡς δὲ ἦν ἐν τοῖς Ἱεροσολύμοις ἐν τῷ πάσχα ἐν τῇ ἑορτῇ, πολλοὶ ἐπίστευσαν εἰς τὸ ὄνομα αὐτοῦ θεωροῦντες αὐτοῦ τὰ σημεῖα ἃ ἐποίει·
\vs{24}αὐτὸς δὲ Ἰησοῦς οὐκ ἐπίστευεν αὑτὸν αὐτοῖς διὰ τὸ αὐτὸν γινώσκειν πάντας
\vs{25}καὶ ὅτι οὐ χρείαν εἶχεν ἵνα τις μαρτυρήσῃ περὶ τοῦ ἀνθρώπου· αὐτὸς γὰρ ἐγίνωσκεν τί ἦν ἐν τῷ ἀνθρώπῳ.

\ch{3}
Ἦν δὲ ἄνθρωπος ἐκ τῶν Φαρισαίων, Νικόδημος ὄνομα αὐτῷ, ἄρχων τῶν Ἰουδαίων·
\vs{2}οὗτος ἦλθεν πρὸς αὐτὸν νυκτὸς καὶ εἶπεν αὐτῷ· Ῥαββί, οἴδαμεν ὅτι ἀπὸ Θεοῦ ἐλήλυθας διδάσκαλος· οὐδεὶς γὰρ δύναται ταῦτα τὰ σημεῖα ποιεῖν ἃ σὺ ποιεῖς, ἐὰν μὴ ᾖ ὁ Θεὸς μετ᾽ αὐτοῦ.
\vs{3}Ἀπεκρίθη Ἰησοῦς καὶ εἶπεν αὐτῷ· Ἀμὴν ἀμὴν λέγω σοι, ἐὰν μή τις γεννηθῇ ἄνωθεν, οὐ δύναται ἰδεῖν τὴν βασιλείαν τοῦ Θεοῦ.
\vs{4}Λέγει πρὸς αὐτὸν ὁ Νικόδημος· Πῶς δύναται ἄνθρωπος γεννηθῆναι γέρων ὤν; μὴ δύναται εἰς τὴν κοιλίαν τῆς μητρὸς αὐτοῦ δεύτερον εἰσελθεῖν καὶ γεννηθῆναι;
\vs{5}Ἀπεκρίθη Ἰησοῦς· Ἀμὴν ἀμὴν λέγω σοι, ἐὰν μή τις γεννηθῇ ἐξ ὕδατος καὶ Πνεύματος, οὐ δύναται εἰσελθεῖν εἰς τὴν βασιλείαν τοῦ Θεοῦ.
\vs{6}τὸ γεγεννημένον ἐκ τῆς σαρκὸς σάρξ ἐστιν, καὶ τὸ γεγεννημένον ἐκ τοῦ Πνεύματος πνεῦμά ἐστιν.
\vs{7}μὴ θαυμάσῃς ὅτι εἶπόν σοι· Δεῖ ὑμᾶς γεννηθῆναι ἄνωθεν.
\vs{8}τὸ πνεῦμα ὅπου θέλει πνεῖ καὶ τὴν φωνὴν αὐτοῦ ἀκούεις, ἀλλ᾽ οὐκ οἶδας πόθεν ἔρχεται καὶ ποῦ ὑπάγει· οὕτως ἐστὶν πᾶς ὁ γεγεννημένος ἐκ τοῦ Πνεύματος.
\vs{9}Ἀπεκρίθη Νικόδημος καὶ εἶπεν αὐτῷ· Πῶς δύναται ταῦτα γενέσθαι;
\vs{10}Ἀπεκρίθη Ἰησοῦς καὶ εἶπεν αὐτῷ· Σὺ εἶ ὁ διδάσκαλος τοῦ Ἰσραὴλ καὶ ταῦτα οὐ γινώσκεις;
\vs{11}ἀμὴν ἀμὴν λέγω σοι ὅτι ὃ οἴδαμεν λαλοῦμεν καὶ ὃ ἑωράκαμεν μαρτυροῦμεν, καὶ τὴν μαρτυρίαν ἡμῶν οὐ λαμβάνετε.
\vs{12}Εἰ τὰ ἐπίγεια εἶπον ὑμῖν καὶ οὐ πιστεύετε, πῶς ἐὰν εἴπω ὑμῖν τὰ ἐπουράνια πιστεύσετε;
\vs{13}καὶ οὐδεὶς ἀναβέβηκεν εἰς τὸν οὐρανὸν εἰ μὴ ὁ ἐκ τοῦ οὐρανοῦ καταβάς, ὁ Υἱὸς τοῦ ἀνθρώπου.
\vs{14}καὶ καθὼς Μωϋσῆς ὕψωσεν τὸν ὄφιν ἐν τῇ ἐρήμῳ, οὕτως ὑψωθῆναι δεῖ τὸν Υἱὸν τοῦ ἀνθρώπου,
\vs{15}ἵνα πᾶς ὁ πιστεύων ἐν αὐτῷ ἔχῃ ζωὴν αἰώνιον.
\vs{16}Οὕτως γὰρ ἠγάπησεν ὁ Θεὸς τὸν κόσμον, ὥστε τὸν Υἱὸν τὸν μονογενῆ ἔδωκεν, ἵνα πᾶς ὁ πιστεύων εἰς αὐτὸν μὴ ἀπόληται ἀλλ᾽ ἔχῃ ζωὴν αἰώνιον.
\vs{17}οὐ γὰρ ἀπέστειλεν ὁ Θεὸς τὸν Υἱὸν εἰς τὸν κόσμον ἵνα κρίνῃ τὸν κόσμον, ἀλλ᾽ ἵνα σωθῇ ὁ κόσμος δι᾽ αὐτοῦ.
\vs{18}ὁ πιστεύων εἰς αὐτὸν οὐ κρίνεται· ὁ δὲ μὴ πιστεύων ἤδη κέκριται, ὅτι μὴ πεπίστευκεν εἰς τὸ ὄνομα τοῦ μονογενοῦς Υἱοῦ τοῦ Θεοῦ.
\vs{19}Αὕτη δέ ἐστιν ἡ κρίσις ὅτι τὸ φῶς ἐλήλυθεν εἰς τὸν κόσμον καὶ ἠγάπησαν οἱ ἄνθρωποι μᾶλλον τὸ σκότος ἢ τὸ φῶς· ἦν γὰρ αὐτῶν πονηρὰ τὰ ἔργα.
\vs{20}πᾶς γὰρ ὁ φαῦλα πράσσων μισεῖ τὸ φῶς καὶ οὐκ ἔρχεται πρὸς τὸ φῶς, ἵνα μὴ ἐλεγχθῇ τὰ ἔργα αὐτοῦ·
\vs{21}ὁ δὲ ποιῶν τὴν ἀλήθειαν ἔρχεται πρὸς τὸ φῶς, ἵνα φανερωθῇ αὐτοῦ τὰ ἔργα ὅτι ἐν Θεῷ ἐστιν εἰργασμένα.

\vs{22}Μετὰ ταῦτα ἦλθεν ὁ Ἰησοῦς καὶ οἱ μαθηταὶ αὐτοῦ εἰς τὴν Ἰουδαίαν γῆν καὶ ἐκεῖ διέτριβεν μετ᾽ αὐτῶν καὶ ἐβάπτιζεν.

\vs{23}Ἦν δὲ καὶ ὁ Ἰωάννης βαπτίζων ἐν Αἰνὼν ἐγγὺς τοῦ Σαλείμ, ὅτι ὕδατα πολλὰ ἦν ἐκεῖ, καὶ παρεγίνοντο καὶ ἐβαπτίζοντο·
\vs{24}οὔπω γὰρ ἦν βεβλημένος εἰς τὴν φυλακὴν ὁ Ἰωάννης.

\vs{25}Ἐγένετο οὖν ζήτησις ἐκ τῶν μαθητῶν Ἰωάννου μετὰ Ἰουδαίου περὶ καθαρισμοῦ.
\vs{26}καὶ ἦλθον πρὸς τὸν Ἰωάννην καὶ εἶπαν αὐτῷ· Ῥαββί, ὃς ἦν μετὰ σοῦ πέραν τοῦ Ἰορδάνου, ᾧ σὺ μεμαρτύρηκας, ἴδε οὗτος βαπτίζει καὶ πάντες ἔρχονται πρὸς αὐτόν.
\vs{27}Ἀπεκρίθη Ἰωάννης καὶ εἶπεν· Οὐ δύναται ἄνθρωπος λαμβάνειν οὐδὲ ἓν ἐὰν μὴ ᾖ δεδομένον αὐτῷ ἐκ τοῦ οὐρανοῦ.
\vs{28}αὐτοὶ ὑμεῖς μοι μαρτυρεῖτε ὅτι εἶπον ὅτι Οὐκ εἰμὶ ἐγὼ ὁ Χριστός, ἀλλ᾽ ὅτι Ἀπεσταλμένος εἰμὶ ἔμπροσθεν ἐκείνου.
\vs{29}Ὁ ἔχων τὴν νύμφην νυμφίος ἐστίν· ὁ δὲ φίλος τοῦ νυμφίου ὁ ἑστηκὼς καὶ ἀκούων αὐτοῦ χαρᾷ χαίρει διὰ τὴν φωνὴν τοῦ νυμφίου. αὕτη οὖν ἡ χαρὰ ἡ ἐμὴ πεπλήρωται.
\vs{30}ἐκεῖνον δεῖ αὐξάνειν, ἐμὲ δὲ ἐλαττοῦσθαι.

\vs{31}Ὁ ἄνωθεν ἐρχόμενος ἐπάνω πάντων ἐστίν· ὁ ὢν ἐκ τῆς γῆς ἐκ τῆς γῆς ἐστιν καὶ ἐκ τῆς γῆς λαλεῖ. ὁ ἐκ τοῦ οὐρανοῦ ἐρχόμενος ἐπάνω πάντων ἐστίν·
\vs{32}ὃ ἑώρακεν καὶ ἤκουσεν τοῦτο μαρτυρεῖ, καὶ τὴν μαρτυρίαν αὐτοῦ οὐδεὶς λαμβάνει.
\vs{33}ὁ λαβὼν αὐτοῦ τὴν μαρτυρίαν ἐσφράγισεν ὅτι ὁ Θεὸς ἀληθής ἐστιν.
\vs{34}ὃν γὰρ ἀπέστειλεν ὁ Θεὸς τὰ ῥήματα τοῦ Θεοῦ λαλεῖ, οὐ γὰρ ἐκ μέτρου δίδωσιν τὸ Πνεῦμα.
\vs{35}Ὁ Πατὴρ ἀγαπᾷ τὸν Υἱόν καὶ πάντα δέδωκεν ἐν τῇ χειρὶ αὐτοῦ.
\vs{36}ὁ πιστεύων εἰς τὸν Υἱὸν ἔχει ζωὴν αἰώνιον· ὁ δὲ ἀπειθῶν τῷ Υἱῷ οὐκ ὄψεται ζωήν, ἀλλ᾽ ἡ ὀργὴ τοῦ Θεοῦ μένει ἐπ᾽ αὐτόν.

\ch{4}
Ὡς οὖν ἔγνω ὁ Ἰησοῦς ὅτι ἤκουσαν οἱ Φαρισαῖοι ὅτι Ἰησοῦς πλείονας μαθητὰς ποιεῖ καὶ βαπτίζει ἢ Ἰωάννης—
\vs{2}καίτοιγε Ἰησοῦς αὐτὸς οὐκ ἐβάπτιζεν ἀλλ᾽ οἱ μαθηταὶ αὐτοῦ—
\vs{3}ἀφῆκεν τὴν Ἰουδαίαν καὶ ἀπῆλθεν πάλιν εἰς τὴν Γαλιλαίαν.

\vs{4}Ἔδει δὲ αὐτὸν διέρχεσθαι διὰ τῆς Σαμαρείας.
\vs{5}ἔρχεται οὖν εἰς πόλιν τῆς Σαμαρείας λεγομένην Συχὰρ πλησίον τοῦ χωρίου ὃ ἔδωκεν Ἰακὼβ τῷ Ἰωσὴφ τῷ υἱῷ αὐτοῦ·
\vs{6}ἦν δὲ ἐκεῖ πηγὴ τοῦ Ἰακώβ. ὁ οὖν Ἰησοῦς κεκοπιακὼς ἐκ τῆς ὁδοιπορίας ἐκαθέζετο οὕτως ἐπὶ τῇ πηγῇ· ὥρα ἦν ὡς ἕκτη.
\vs{7}Ἔρχεται γυνὴ ἐκ τῆς Σαμαρείας ἀντλῆσαι ὕδωρ. λέγει αὐτῇ ὁ Ἰησοῦς· Δός μοι πεῖν·
\vs{8}οἱ γὰρ μαθηταὶ αὐτοῦ ἀπεληλύθεισαν εἰς τὴν πόλιν ἵνα τροφὰς ἀγοράσωσιν.
\vs{9}Λέγει οὖν αὐτῷ ἡ γυνὴ ἡ Σαμαρῖτις· Πῶς σὺ Ἰουδαῖος ὢν παρ᾽ ἐμοῦ πεῖν αἰτεῖς γυναικὸς Σαμαρίτιδος οὔσης; οὐ γὰρ συνχρῶνται Ἰουδαῖοι Σαμαρίταις.
\vs{10}Ἀπεκρίθη Ἰησοῦς καὶ εἶπεν αὐτῇ· Εἰ ᾔδεις τὴν δωρεὰν τοῦ Θεοῦ καὶ τίς ἐστιν ὁ λέγων σοι· Δός μοι πεῖν, σὺ ἂν ᾔτησας αὐτὸν καὶ ἔδωκεν ἄν σοι ὕδωρ ζῶν.
\vs{11}Λέγει αὐτῷ ἡ γυνή· Κύριε, οὔτε ἄντλημα ἔχεις καὶ τὸ φρέαρ ἐστὶν βαθύ· πόθεν οὖν ἔχεις τὸ ὕδωρ τὸ ζῶν;
\vs{12}μὴ σὺ μείζων εἶ τοῦ πατρὸς ἡμῶν Ἰακώβ, ὃς ἔδωκεν ἡμῖν τὸ φρέαρ καὶ αὐτὸς ἐξ αὐτοῦ ἔπιεν καὶ οἱ υἱοὶ αὐτοῦ καὶ τὰ θρέμματα αὐτοῦ;
\vs{13}Ἀπεκρίθη Ἰησοῦς καὶ εἶπεν αὐτῇ· Πᾶς ὁ πίνων ἐκ τοῦ ὕδατος τούτου διψήσει πάλιν·
\vs{14}ὃς δ᾽ ἂν πίῃ ἐκ τοῦ ὕδατος οὗ ἐγὼ δώσω αὐτῷ, οὐ μὴ διψήσει εἰς τὸν αἰῶνα, ἀλλὰ τὸ ὕδωρ ὃ δώσω αὐτῷ γενήσεται ἐν αὐτῷ πηγὴ ὕδατος ἁλλομένου εἰς ζωὴν αἰώνιον.
\vs{15}Λέγει πρὸς αὐτὸν ἡ γυνή· Κύριε, δός μοι τοῦτο τὸ ὕδωρ, ἵνα μὴ διψῶ μηδὲ διέρχωμαι ἐνθάδε ἀντλεῖν.
\vs{16}Λέγει αὐτῇ· Ὕπαγε φώνησον τὸν ἄνδρα σου καὶ ἐλθὲ ἐνθάδε.
\vs{17}Ἀπεκρίθη ἡ γυνὴ καὶ εἶπεν αὐτῷ· Οὐκ ἔχω ἄνδρα. Λέγει αὐτῇ ὁ Ἰησοῦς· Καλῶς εἶπας ὅτι Ἄνδρα οὐκ ἔχω·
\vs{18}πέντε γὰρ ἄνδρας ἔσχες καὶ νῦν ὃν ἔχεις οὐκ ἔστιν σου ἀνήρ· τοῦτο ἀληθὲς εἴρηκας.
\vs{19}Λέγει αὐτῷ ἡ γυνή· Κύριε, θεωρῶ ὅτι προφήτης εἶ σύ.
\vs{20}οἱ πατέρες ἡμῶν ἐν τῷ ὄρει τούτῳ προσεκύνησαν· καὶ ὑμεῖς λέγετε ὅτι ἐν Ἱεροσολύμοις ἐστὶν ὁ τόπος ὅπου προσκυνεῖν δεῖ.
\vs{21}Λέγει αὐτῇ ὁ Ἰησοῦς· Πίστευέ μοι, γύναι, ὅτι ἔρχεται ὥρα ὅτε οὔτε ἐν τῷ ὄρει τούτῳ οὔτε ἐν Ἱεροσολύμοις προσκυνήσετε τῷ Πατρί.
\vs{22}ὑμεῖς προσκυνεῖτε ὃ οὐκ οἴδατε· ἡμεῖς προσκυνοῦμεν ὃ οἴδαμεν, ὅτι ἡ σωτηρία ἐκ τῶν Ἰουδαίων ἐστίν.
\vs{23}ἀλλὰ ἔρχεται ὥρα καὶ νῦν ἐστιν, ὅτε οἱ ἀληθινοὶ προσκυνηταὶ προσκυνήσουσιν τῷ Πατρὶ ἐν πνεύματι καὶ ἀληθείᾳ· καὶ γὰρ ὁ Πατὴρ τοιούτους ζητεῖ τοὺς προσκυνοῦντας αὐτόν.
\vs{24}Πνεῦμα ὁ Θεός, καὶ τοὺς προσκυνοῦντας αὐτὸν ἐν πνεύματι καὶ ἀληθείᾳ δεῖ προσκυνεῖν.
\vs{25}Λέγει αὐτῷ ἡ γυνή· Οἶδα ὅτι Μεσσίας ἔρχεται ὁ λεγόμενος Χριστός· ὅταν ἔλθῃ ἐκεῖνος, ἀναγγελεῖ ἡμῖν ἅπαντα.
\vs{26}Λέγει αὐτῇ ὁ Ἰησοῦς· Ἐγώ εἰμι, ὁ λαλῶν σοι.

\vs{27}Καὶ ἐπὶ τούτῳ ἦλθαν οἱ μαθηταὶ αὐτοῦ καὶ ἐθαύμαζον ὅτι μετὰ γυναικὸς ἐλάλει· οὐδεὶς μέντοι εἶπεν· Τί ζητεῖς ἢ Τί λαλεῖς μετ᾽ αὐτῆς;
\vs{28}Ἀφῆκεν οὖν τὴν ὑδρίαν αὐτῆς ἡ γυνὴ καὶ ἀπῆλθεν εἰς τὴν πόλιν καὶ λέγει τοῖς ἀνθρώποις·
\vs{29}Δεῦτε ἴδετε ἄνθρωπον ὃς εἶπέν μοι πάντα ὅσα ἐποίησα, μήτι οὗτός ἐστιν ὁ Χριστός;
\vs{30}ἐξῆλθον ἐκ τῆς πόλεως καὶ ἤρχοντο πρὸς αὐτόν.

\vs{31}Ἐν τῷ μεταξὺ ἠρώτων αὐτὸν οἱ μαθηταὶ λέγοντες· Ῥαββί, φάγε.
\vs{32}Ὁ δὲ εἶπεν αὐτοῖς· Ἐγὼ βρῶσιν ἔχω φαγεῖν ἣν ὑμεῖς οὐκ οἴδατε.
\vs{33}Ἔλεγον οὖν οἱ μαθηταὶ πρὸς ἀλλήλους· Μή τις ἤνεγκεν αὐτῷ φαγεῖν;
\vs{34}Λέγει αὐτοῖς ὁ Ἰησοῦς· Ἐμὸν βρῶμά ἐστιν ἵνα ποιήσω τὸ θέλημα τοῦ πέμψαντός με καὶ τελειώσω αὐτοῦ τὸ ἔργον.
\vs{35}οὐχ ὑμεῖς λέγετε ὅτι Ἔτι τετράμηνός ἐστιν καὶ ὁ θερισμὸς ἔρχεται; ἰδοὺ λέγω ὑμῖν, ἐπάρατε τοὺς ὀφθαλμοὺς ὑμῶν καὶ θεάσασθε τὰς χώρας ὅτι λευκαί εἰσιν πρὸς θερισμόν. ἤδη
\vs{36}Ὁ θερίζων μισθὸν λαμβάνει καὶ συνάγει καρπὸν εἰς ζωὴν αἰώνιον, ἵνα ὁ σπείρων ὁμοῦ χαίρῃ καὶ ὁ θερίζων.
\vs{37}ἐν γὰρ τούτῳ ὁ λόγος ἐστὶν ἀληθινὸς ὅτι Ἄλλος ἐστὶν ὁ σπείρων καὶ ἄλλος ὁ θερίζων.
\vs{38}ἐγὼ ἀπέστειλα ὑμᾶς θερίζειν ὃ οὐχ ὑμεῖς κεκοπιάκατε· ἄλλοι κεκοπιάκασιν καὶ ὑμεῖς εἰς τὸν κόπον αὐτῶν εἰσεληλύθατε.

\vs{39}Ἐκ δὲ τῆς πόλεως ἐκείνης πολλοὶ ἐπίστευσαν εἰς αὐτὸν τῶν Σαμαριτῶν διὰ τὸν λόγον τῆς γυναικὸς μαρτυρούσης ὅτι Εἶπέν μοι πάντα ἃ ἐποίησα.
\vs{40}ὡς οὖν ἦλθον πρὸς αὐτὸν οἱ Σαμαρῖται, ἠρώτων αὐτὸν μεῖναι παρ᾽ αὐτοῖς· καὶ ἔμεινεν ἐκεῖ δύο ἡμέρας.
\vs{41}Καὶ πολλῷ πλείους ἐπίστευσαν διὰ τὸν λόγον αὐτοῦ,
\vs{42}τῇ τε γυναικὶ ἔλεγον ὅτι Οὐκέτι διὰ τὴν σὴν λαλιὰν πιστεύομεν, αὐτοὶ γὰρ ἀκηκόαμεν καὶ οἴδαμεν ὅτι οὗτός ἐστιν ἀληθῶς ὁ Σωτὴρ τοῦ κόσμου.
\vs{43}Μετὰ δὲ τὰς δύο ἡμέρας ἐξῆλθεν ἐκεῖθεν εἰς τὴν Γαλιλαίαν·
\vs{44}αὐτὸς γὰρ Ἰησοῦς ἐμαρτύρησεν ὅτι προφήτης ἐν τῇ ἰδίᾳ πατρίδι τιμὴν οὐκ ἔχει.
\vs{45}ὅτε οὖν ἦλθεν εἰς τὴν Γαλιλαίαν, ἐδέξαντο αὐτὸν οἱ Γαλιλαῖοι πάντα ἑωρακότες ὅσα ἐποίησεν ἐν Ἱεροσολύμοις ἐν τῇ ἑορτῇ, καὶ αὐτοὶ γὰρ ἦλθον εἰς τὴν ἑορτήν.

\vs{46}Ἦλθεν οὖν πάλιν εἰς τὴν Κανὰ τῆς Γαλιλαίας, ὅπου ἐποίησεν τὸ ὕδωρ οἶνον.

Καὶ ἦν τις βασιλικὸς οὗ ὁ υἱὸς ἠσθένει ἐν Καφαρναούμ.
\vs{47}οὗτος ἀκούσας ὅτι Ἰησοῦς ἥκει ἐκ τῆς Ἰουδαίας εἰς τὴν Γαλιλαίαν ἀπῆλθεν πρὸς αὐτὸν καὶ ἠρώτα ἵνα καταβῇ καὶ ἰάσηται αὐτοῦ τὸν υἱόν, ἤμελλεν γὰρ ἀποθνήσκειν.
\vs{48}Εἶπεν οὖν ὁ Ἰησοῦς πρὸς αὐτόν· Ἐὰν μὴ σημεῖα καὶ τέρατα ἴδητε, οὐ μὴ πιστεύσητε.
\vs{49}Λέγει πρὸς αὐτὸν ὁ βασιλικός· Κύριε, κατάβηθι πρὶν ἀποθανεῖν τὸ παιδίον μου.
\vs{50}Λέγει αὐτῷ ὁ Ἰησοῦς· Πορεύου, ὁ υἱός σου ζῇ. Ἐπίστευσεν ὁ ἄνθρωπος τῷ λόγῳ ὃν εἶπεν αὐτῷ ὁ Ἰησοῦς καὶ ἐπορεύετο.
\vs{51}ἤδη δὲ αὐτοῦ καταβαίνοντος οἱ δοῦλοι αὐτοῦ ὑπήντησαν αὐτῷ λέγοντες ὅτι ὁ παῖς αὐτοῦ ζῇ.
\vs{52}Ἐπύθετο οὖν τὴν ὥραν παρ᾽ αὐτῶν ἐν ᾗ κομψότερον ἔσχεν· εἶπαν οὖν αὐτῷ ὅτι Ἐχθὲς ὥραν ἑβδόμην ἀφῆκεν αὐτὸν ὁ πυρετός.
\vs{53}Ἔγνω οὖν ὁ πατὴρ ὅτι ἐν ἐκείνῃ τῇ ὥρᾳ ἐν ᾗ εἶπεν αὐτῷ ὁ Ἰησοῦς· Ὁ υἱός σου ζῇ, καὶ ἐπίστευσεν αὐτὸς καὶ ἡ οἰκία αὐτοῦ ὅλη.
\vs{54}Τοῦτο δὲ πάλιν δεύτερον σημεῖον ἐποίησεν ὁ Ἰησοῦς ἐλθὼν ἐκ τῆς Ἰουδαίας εἰς τὴν Γαλιλαίαν.

\ch{5}
Μετὰ ταῦτα ἦν ἑορτὴ τῶν Ἰουδαίων καὶ ἀνέβη Ἰησοῦς εἰς Ἱεροσόλυμα.

\vs{2}Ἔστιν δὲ ἐν τοῖς Ἱεροσολύμοις ἐπὶ τῇ προβατικῇ κολυμβήθρα ἡ ἐπιλεγομένη Ἑβραϊστὶ Βηθζαθά πέντε στοὰς ἔχουσα.
\vs{3}ἐν ταύταις κατέκειτο πλῆθος τῶν ἀσθενούντων, τυφλῶν, χωλῶν, ξηρῶν.
\vs{5}Ἦν δέ τις ἄνθρωπος ἐκεῖ τριάκοντα καὶ ὀκτὼ ἔτη ἔχων ἐν τῇ ἀσθενείᾳ αὐτοῦ·
\vs{6}τοῦτον ἰδὼν ὁ Ἰησοῦς κατακείμενον καὶ γνοὺς ὅτι πολὺν ἤδη χρόνον ἔχει, λέγει αὐτῷ· Θέλεις ὑγιὴς γενέσθαι;
\vs{7}Ἀπεκρίθη αὐτῷ ὁ ἀσθενῶν· Κύριε, ἄνθρωπον οὐκ ἔχω ἵνα ὅταν ταραχθῇ τὸ ὕδωρ βάλῃ με εἰς τὴν κολυμβήθραν· ἐν ᾧ δὲ ἔρχομαι ἐγὼ, ἄλλος πρὸ ἐμοῦ καταβαίνει.
\vs{8}Λέγει αὐτῷ ὁ Ἰησοῦς· Ἔγειρε ἆρον τὸν κράβαττόν σου καὶ περιπάτει.
\vs{9}Καὶ εὐθέως ἐγένετο ὑγιὴς ὁ ἄνθρωπος καὶ ἦρεν τὸν κράβαττον αὐτοῦ καὶ περιεπάτει.

Ἦν δὲ σάββατον ἐν ἐκείνῃ τῇ ἡμέρᾳ.
\vs{10}ἔλεγον οὖν οἱ Ἰουδαῖοι τῷ τεθεραπευμένῳ· Σάββατόν ἐστιν, καὶ οὐκ ἔξεστίν σοι ἆραι τὸν κράβαττον σου.
\vs{11}ὁ δὲ ἀπεκρίθη αὐτοῖς· Ὁ ποιήσας με ὑγιῆ ἐκεῖνός μοι εἶπεν· Ἆρον τὸν κράβαττόν σου καὶ περιπάτει.
\vs{12}Ἠρώτησαν αὐτόν· Τίς ἐστιν ὁ ἄνθρωπος ὁ εἰπών σοι· Ἆρον καὶ περιπάτει;
\vs{13}Ὁ δὲ ἰαθεὶς οὐκ ᾔδει τίς ἐστιν, ὁ γὰρ Ἰησοῦς ἐξένευσεν ὄχλου ὄντος ἐν τῷ τόπῳ.
\vs{14}Μετὰ ταῦτα εὑρίσκει αὐτὸν ὁ Ἰησοῦς ἐν τῷ ἱερῷ καὶ εἶπεν αὐτῷ· Ἴδε ὑγιὴς γέγονας, μηκέτι ἁμάρτανε, ἵνα μὴ χεῖρόν σοί τι γένηται.
\vs{15}ἀπῆλθεν ὁ ἄνθρωπος καὶ ἀνήγγειλεν τοῖς Ἰουδαίοις ὅτι Ἰησοῦς ἐστιν ὁ ποιήσας αὐτὸν ὑγιῆ.
\vs{16}Καὶ διὰ τοῦτο ἐδίωκον οἱ Ἰουδαῖοι τὸν Ἰησοῦν, ὅτι ταῦτα ἐποίει ἐν σαββάτῳ.

\vs{17}ὁ δὲ Ἰησοῦς ἀπεκρίνατο αὐτοῖς· Ὁ Πατήρ μου ἕως ἄρτι ἐργάζεται κἀγὼ ἐργάζομαι·
\vs{18}Διὰ τοῦτο οὖν μᾶλλον ἐζήτουν αὐτὸν οἱ Ἰουδαῖοι ἀποκτεῖναι, ὅτι οὐ μόνον ἔλυεν τὸ σάββατον, ἀλλὰ καὶ Πατέρα ἴδιον ἔλεγεν τὸν Θεόν ἴσον ἑαυτὸν ποιῶν τῷ Θεῷ.

\vs{19}Ἀπεκρίνατο οὖν ὁ Ἰησοῦς καὶ ἔλεγεν αὐτοῖς· Ἀμὴν ἀμὴν λέγω ὑμῖν, οὐ δύναται ὁ Υἱὸς ποιεῖν ἀφ᾽ ἑαυτοῦ οὐδὲν ἐὰν μή τι βλέπῃ τὸν Πατέρα ποιοῦντα· ἃ γὰρ ἂν ἐκεῖνος ποιῇ, ταῦτα καὶ ὁ Υἱὸς ὁμοίως ποιεῖ.
\vs{20}ὁ γὰρ Πατὴρ φιλεῖ τὸν Υἱὸν καὶ πάντα δείκνυσιν αὐτῷ ἃ αὐτὸς ποιεῖ, καὶ μείζονα τούτων δείξει αὐτῷ ἔργα, ἵνα ὑμεῖς θαυμάζητε.
\vs{21}Ὥσπερ γὰρ ὁ Πατὴρ ἐγείρει τοὺς νεκροὺς καὶ ζωοποιεῖ, οὕτως καὶ ὁ Υἱὸς οὓς θέλει ζωοποιεῖ.
\vs{22}οὐδὲ γὰρ ὁ Πατὴρ κρίνει οὐδένα, ἀλλὰ τὴν κρίσιν πᾶσαν δέδωκεν τῷ Υἱῷ,
\vs{23}ἵνα πάντες τιμῶσι τὸν Υἱὸν καθὼς τιμῶσι τὸν Πατέρα. ὁ μὴ τιμῶν τὸν Υἱὸν οὐ τιμᾷ τὸν Πατέρα τὸν πέμψαντα αὐτόν.

\vs{24}Ἀμὴν ἀμὴν λέγω ὑμῖν ὅτι ὁ τὸν λόγον μου ἀκούων καὶ πιστεύων τῷ πέμψαντί με ἔχει ζωὴν αἰώνιον καὶ εἰς κρίσιν οὐκ ἔρχεται, ἀλλὰ μεταβέβηκεν ἐκ τοῦ θανάτου εἰς τὴν ζωήν.
\vs{25}Ἀμὴν ἀμὴν λέγω ὑμῖν ὅτι ἔρχεται ὥρα καὶ νῦν ἐστιν ὅτε οἱ νεκροὶ ἀκούσουσιν τῆς φωνῆς τοῦ Υἱοῦ τοῦ Θεοῦ καὶ οἱ ἀκούσαντες ζήσουσιν.
\vs{26}ὥσπερ γὰρ ὁ Πατὴρ ἔχει ζωὴν ἐν ἑαυτῷ, οὕτως καὶ τῷ Υἱῷ ἔδωκεν ζωὴν ἔχειν ἐν ἑαυτῷ.
\vs{27}καὶ ἐξουσίαν ἔδωκεν αὐτῷ κρίσιν ποιεῖν, ὅτι Υἱὸς ἀνθρώπου ἐστίν.
\vs{28}μὴ θαυμάζετε τοῦτο, ὅτι ἔρχεται ὥρα ἐν ᾗ πάντες οἱ ἐν τοῖς μνημείοις ἀκούσουσιν τῆς φωνῆς αὐτοῦ
\vs{29}καὶ ἐκπορεύσονται οἱ τὰ ἀγαθὰ ποιήσαντες εἰς ἀνάστασιν ζωῆς, οἱ δὲ τὰ φαῦλα πράξαντες εἰς ἀνάστασιν κρίσεως.

\vs{30}Οὐ δύναμαι ἐγὼ ποιεῖν ἀπ᾽ ἐμαυτοῦ οὐδέν· καθὼς ἀκούω κρίνω, καὶ ἡ κρίσις ἡ ἐμὴ δικαία ἐστίν, ὅτι οὐ ζητῶ τὸ θέλημα τὸ ἐμὸν ἀλλὰ τὸ θέλημα τοῦ πέμψαντός με.
\vs{31}Ἐὰν ἐγὼ μαρτυρῶ περὶ ἐμαυτοῦ, ἡ μαρτυρία μου οὐκ ἔστιν ἀληθής·
\vs{32}ἄλλος ἐστὶν ὁ μαρτυρῶν περὶ ἐμοῦ, καὶ οἶδα ὅτι ἀληθής ἐστιν ἡ μαρτυρία ἣν μαρτυρεῖ περὶ ἐμοῦ.
\vs{33}Ὑμεῖς ἀπεστάλκατε πρὸς Ἰωάννην, καὶ μεμαρτύρηκεν τῇ ἀληθείᾳ·
\vs{34}ἐγὼ δὲ οὐ παρὰ ἀνθρώπου τὴν μαρτυρίαν λαμβάνω, ἀλλὰ ταῦτα λέγω ἵνα ὑμεῖς σωθῆτε.
\vs{35}Ἐκεῖνος ἦν ὁ λύχνος ὁ καιόμενος καὶ φαίνων, ὑμεῖς δὲ ἠθελήσατε ἀγαλλιαθῆναι πρὸς ὥραν ἐν τῷ φωτὶ αὐτοῦ.

\vs{36}ἐγὼ δὲ ἔχω τὴν μαρτυρίαν μείζω τοῦ Ἰωάννου· τὰ γὰρ ἔργα ἃ δέδωκέν μοι ὁ Πατὴρ ἵνα τελειώσω αὐτά, αὐτὰ τὰ ἔργα ἃ ποιῶ μαρτυρεῖ περὶ ἐμοῦ ὅτι ὁ Πατήρ με ἀπέσταλκεν.
\vs{37}καὶ ὁ πέμψας με Πατὴρ ἐκεῖνος μεμαρτύρηκεν περὶ ἐμοῦ. οὔτε φωνὴν αὐτοῦ πώποτε ἀκηκόατε οὔτε εἶδος αὐτοῦ ἑωράκατε,
\vs{38}καὶ τὸν λόγον αὐτοῦ οὐκ ἔχετε ἐν ὑμῖν μένοντα, ὅτι ὃν ἀπέστειλεν ἐκεῖνος, τούτῳ ὑμεῖς οὐ πιστεύετε.
\vs{39}Ἐραυνᾶτε τὰς γραφάς, ὅτι ὑμεῖς δοκεῖτε ἐν αὐταῖς ζωὴν αἰώνιον ἔχειν· καὶ ἐκεῖναί εἰσιν αἱ μαρτυροῦσαι περὶ ἐμοῦ·
\vs{40}καὶ οὐ θέλετε ἐλθεῖν πρός με ἵνα ζωὴν ἔχητε.

\vs{41}Δόξαν παρὰ ἀνθρώπων οὐ λαμβάνω,
\vs{42}ἀλλὰ ἔγνωκα ὑμᾶς ὅτι τὴν ἀγάπην τοῦ Θεοῦ οὐκ ἔχετε ἐν ἑαυτοῖς.
\vs{43}ἐγὼ ἐλήλυθα ἐν τῷ ὀνόματι τοῦ Πατρός μου, καὶ οὐ λαμβάνετέ με· ἐὰν ἄλλος ἔλθῃ ἐν τῷ ὀνόματι τῷ ἰδίῳ, ἐκεῖνον λήμψεσθε.
\vs{44}πῶς δύνασθε ὑμεῖς πιστεῦσαι δόξαν παρὰ ἀλλήλων λαμβάνοντες, καὶ τὴν δόξαν τὴν παρὰ τοῦ μόνου Θεοῦ οὐ ζητεῖτε;

\vs{45}Μὴ δοκεῖτε ὅτι ἐγὼ κατηγορήσω ὑμῶν πρὸς τὸν Πατέρα· ἔστιν ὁ κατηγορῶν ὑμῶν Μωϋσῆς, εἰς ὃν ὑμεῖς ἠλπίκατε.
\vs{46}εἰ γὰρ ἐπιστεύετε Μωϋσεῖ, ἐπιστεύετε ἂν ἐμοί· περὶ γὰρ ἐμοῦ ἐκεῖνος ἔγραψεν.
\vs{47}εἰ δὲ τοῖς ἐκείνου γράμμασιν οὐ πιστεύετε, πῶς τοῖς ἐμοῖς ῥήμασιν πιστεύσετε;

\ch{6}
Μετὰ ταῦτα ἀπῆλθεν ὁ Ἰησοῦς πέραν τῆς θαλάσσης τῆς Γαλιλαίας τῆς Τιβεριάδος.
\vs{2}ἠκολούθει δὲ αὐτῷ ὄχλος πολύς, ὅτι ἐθεώρουν τὰ σημεῖα ἃ ἐποίει ἐπὶ τῶν ἀσθενούντων.
\vs{3}ἀνῆλθεν δὲ εἰς τὸ ὄρος Ἰησοῦς καὶ ἐκεῖ ἐκάθητο μετὰ τῶν μαθητῶν αὐτοῦ.
\vs{4}Ἦν δὲ ἐγγὺς τὸ πάσχα, ἡ ἑορτὴ τῶν Ἰουδαίων.

\vs{5}ἐπάρας οὖν τοὺς ὀφθαλμοὺς ὁ Ἰησοῦς καὶ θεασάμενος ὅτι πολὺς ὄχλος ἔρχεται πρὸς αὐτὸν λέγει πρὸς Φίλιππον· Πόθεν ἀγοράσωμεν ἄρτους ἵνα φάγωσιν οὗτοι;
\vs{6}τοῦτο δὲ ἔλεγεν πειράζων αὐτόν· αὐτὸς γὰρ ᾔδει τί ἔμελλεν ποιεῖν.
\vs{7}Ἀπεκρίθη αὐτῷ ὁ Φίλιππος· Διακοσίων δηναρίων ἄρτοι οὐκ ἀρκοῦσιν αὐτοῖς ἵνα ἕκαστος βραχύ τι λάβῃ.
\vs{8}Λέγει αὐτῷ εἷς ἐκ τῶν μαθητῶν αὐτοῦ, Ἀνδρέας ὁ ἀδελφὸς Σίμωνος Πέτρου·
\vs{9}Ἔστιν παιδάριον ὧδε ὃς ἔχει πέντε ἄρτους κριθίνους καὶ δύο ὀψάρια· ἀλλὰ ταῦτα τί ἐστιν εἰς τοσούτους;
\vs{10}Εἶπεν ὁ Ἰησοῦς· Ποιήσατε τοὺς ἀνθρώπους ἀναπεσεῖν. ἦν δὲ χόρτος πολὺς ἐν τῷ τόπῳ. ἀνέπεσαν οὖν οἱ ἄνδρες τὸν ἀριθμὸν ὡς πεντακισχίλιοι.
\vs{11}Ἔλαβεν οὖν τοὺς ἄρτους ὁ Ἰησοῦς καὶ εὐχαριστήσας διέδωκεν τοῖς ἀνακειμένοις ὁμοίως καὶ ἐκ τῶν ὀψαρίων ὅσον ἤθελον.
\vs{12}Ὡς δὲ ἐνεπλήσθησαν, λέγει τοῖς μαθηταῖς αὐτοῦ· Συναγάγετε τὰ περισσεύσαντα κλάσματα, ἵνα μή τι ἀπόληται.
\vs{13}συνήγαγον οὖν καὶ ἐγέμισαν δώδεκα κοφίνους κλασμάτων ἐκ τῶν πέντε ἄρτων τῶν κριθίνων ἃ ἐπερίσσευσαν τοῖς βεβρωκόσιν.
\vs{14}Οἱ οὖν ἄνθρωποι ἰδόντες ὃ ἐποίησεν σημεῖον ἔλεγον ὅτι Οὗτός ἐστιν ἀληθῶς ὁ προφήτης ὁ ἐρχόμενος εἰς τὸν κόσμον.
\vs{15}Ἰησοῦς οὖν γνοὺς ὅτι μέλλουσιν ἔρχεσθαι καὶ ἁρπάζειν αὐτὸν ἵνα ποιήσωσιν βασιλέα, ἀνεχώρησεν πάλιν εἰς τὸ ὄρος αὐτὸς μόνος.
\vs{16}Ὡς δὲ ὀψία ἐγένετο κατέβησαν οἱ μαθηταὶ αὐτοῦ ἐπὶ τὴν θάλασσαν
\vs{17}καὶ ἐμβάντες εἰς πλοῖον ἤρχοντο πέραν τῆς θαλάσσης εἰς Καφαρναούμ. καὶ σκοτία ἤδη ἐγεγόνει καὶ οὔπω ἐληλύθει πρὸς αὐτοὺς ὁ Ἰησοῦς,
\vs{18}ἥ τε θάλασσα ἀνέμου μεγάλου πνέοντος διεγείρετο.
\vs{19}Ἐληλακότες οὖν ὡς σταδίους εἴκοσι πέντε ἢ τριάκοντα θεωροῦσιν τὸν Ἰησοῦν περιπατοῦντα ἐπὶ τῆς θαλάσσης καὶ ἐγγὺς τοῦ πλοίου γινόμενον, καὶ ἐφοβήθησαν.
\vs{20}ὁ δὲ λέγει αὐτοῖς· Ἐγώ εἰμι· μὴ φοβεῖσθε.
\vs{21}ἤθελον οὖν λαβεῖν αὐτὸν εἰς τὸ πλοῖον, καὶ εὐθέως ἐγένετο τὸ πλοῖον ἐπὶ τῆς γῆς εἰς ἣν ὑπῆγον.

\vs{22}Τῇ ἐπαύριον ὁ ὄχλος ὁ ἑστηκὼς πέραν τῆς θαλάσσης εἶδον ὅτι πλοιάριον ἄλλο οὐκ ἦν ἐκεῖ εἰ μὴ ἕν καὶ ὅτι οὐ συνεισῆλθεν τοῖς μαθηταῖς αὐτοῦ ὁ Ἰησοῦς εἰς τὸ πλοῖον ἀλλὰ μόνοι οἱ μαθηταὶ αὐτοῦ ἀπῆλθον·
\vs{23}ἀλλὰ ἦλθεν πλοιάρια ἐκ Τιβεριάδος ἐγγὺς τοῦ τόπου ὅπου ἔφαγον τὸν ἄρτον εὐχαριστήσαντος τοῦ Κυρίου.
\vs{24}ὅτε οὖν εἶδεν ὁ ὄχλος ὅτι Ἰησοῦς οὐκ ἔστιν ἐκεῖ οὐδὲ οἱ μαθηταὶ αὐτοῦ, ἐνέβησαν αὐτοὶ εἰς τὰ πλοιάρια καὶ ἦλθον εἰς Καφαρναοὺμ ζητοῦντες τὸν Ἰησοῦν.
\vs{25}καὶ εὑρόντες αὐτὸν πέραν τῆς θαλάσσης εἶπον αὐτῷ· Ῥαββί, πότε ὧδε γέγονας;

\vs{26}Ἀπεκρίθη αὐτοῖς ὁ Ἰησοῦς καὶ εἶπεν· Ἀμὴν ἀμὴν λέγω ὑμῖν, ζητεῖτέ με οὐχ ὅτι εἴδετε σημεῖα, ἀλλ᾽ ὅτι ἐφάγετε ἐκ τῶν ἄρτων καὶ ἐχορτάσθητε.
\vs{27}ἐργάζεσθε μὴ τὴν βρῶσιν τὴν ἀπολλυμένην ἀλλὰ τὴν βρῶσιν τὴν μένουσαν εἰς ζωὴν αἰώνιον, ἣν ὁ Υἱὸς τοῦ ἀνθρώπου ὑμῖν δώσει· τοῦτον γὰρ ὁ Πατὴρ ἐσφράγισεν ὁ Θεός.
\vs{28}Εἶπον οὖν πρὸς αὐτόν· Τί ποιῶμεν ἵνα ἐργαζώμεθα τὰ ἔργα τοῦ Θεοῦ;
\vs{29}Ἀπεκρίθη ὁ Ἰησοῦς καὶ εἶπεν αὐτοῖς· Τοῦτό ἐστιν τὸ ἔργον τοῦ Θεοῦ, ἵνα πιστεύητε εἰς ὃν ἀπέστειλεν ἐκεῖνος.

\vs{30}Εἶπον οὖν αὐτῷ· Τί οὖν ποιεῖς σὺ σημεῖον, ἵνα ἴδωμεν καὶ πιστεύσωμέν σοι; τί ἐργάζῃ;
\vs{31}οἱ πατέρες ἡμῶν τὸ μάννα ἔφαγον ἐν τῇ ἐρήμῳ, καθώς ἐστιν γεγραμμένον· Ἄρτον ἐκ τοῦ οὐρανοῦ ἔδωκεν αὐτοῖς φαγεῖν.
\vs{32}Εἶπεν οὖν αὐτοῖς ὁ Ἰησοῦς· Ἀμὴν ἀμὴν λέγω ὑμῖν, οὐ Μωϋσῆς δέδωκεν ὑμῖν τὸν ἄρτον ἐκ τοῦ οὐρανοῦ, ἀλλ᾽ ὁ Πατήρ μου δίδωσιν ὑμῖν τὸν ἄρτον ἐκ τοῦ οὐρανοῦ τὸν ἀληθινόν·
\vs{33}ὁ γὰρ ἄρτος τοῦ Θεοῦ ἐστιν ὁ καταβαίνων ἐκ τοῦ οὐρανοῦ καὶ ζωὴν διδοὺς τῷ κόσμῳ.
\vs{34}Εἶπον οὖν πρὸς αὐτόν· Κύριε, πάντοτε δὸς ἡμῖν τὸν ἄρτον τοῦτον.
\vs{35}Εἶπεν αὐτοῖς ὁ Ἰησοῦς· Ἐγώ εἰμι ὁ ἄρτος τῆς ζωῆς· ὁ ἐρχόμενος πρὸς ἐμὲ οὐ μὴ πεινάσῃ, καὶ ὁ πιστεύων εἰς ἐμὲ οὐ μὴ διψήσει πώποτε.

\vs{36}ἀλλ᾽ εἶπον ὑμῖν ὅτι καὶ ἑωράκατέ με καὶ οὐ πιστεύετε.
\vs{37}Πᾶν ὃ δίδωσίν μοι ὁ Πατὴρ πρὸς ἐμὲ ἥξει, καὶ τὸν ἐρχόμενον πρός ἐμὲ οὐ μὴ ἐκβάλω ἔξω,
\vs{38}ὅτι καταβέβηκα ἀπὸ τοῦ οὐρανοῦ οὐχ ἵνα ποιῶ τὸ θέλημα τὸ ἐμὸν ἀλλὰ τὸ θέλημα τοῦ πέμψαντός με.
\vs{39}Τοῦτο δέ ἐστιν τὸ θέλημα τοῦ πέμψαντός με, ἵνα πᾶν ὃ δέδωκέν μοι μὴ ἀπολέσω ἐξ αὐτοῦ, ἀλλὰ ἀναστήσω αὐτὸ ἐν τῇ ἐσχάτῃ ἡμέρᾳ.
\vs{40}τοῦτο γάρ ἐστιν τὸ θέλημα τοῦ Πατρός μου, ἵνα πᾶς ὁ θεωρῶν τὸν Υἱὸν καὶ πιστεύων εἰς αὐτὸν ἔχῃ ζωὴν αἰώνιον, καὶ ἀναστήσω αὐτὸν ἐγὼ ἐν τῇ ἐσχάτῃ ἡμέρᾳ.

\vs{41}Ἐγόγγυζον οὖν οἱ Ἰουδαῖοι περὶ αὐτοῦ ὅτι εἶπεν· Ἐγώ εἰμι ὁ ἄρτος ὁ καταβὰς ἐκ τοῦ οὐρανοῦ,
\vs{42}καὶ ἔλεγον· Οὐχ οὗτός ἐστιν Ἰησοῦς ὁ υἱὸς Ἰωσήφ, οὗ ἡμεῖς οἴδαμεν τὸν πατέρα καὶ τὴν μητέρα; πῶς νῦν λέγει ὅτι Ἐκ τοῦ οὐρανοῦ καταβέβηκα;
\vs{43}Ἀπεκρίθη Ἰησοῦς καὶ εἶπεν αὐτοῖς· Μὴ γογγύζετε μετ᾽ ἀλλήλων.
\vs{44}οὐδεὶς δύναται ἐλθεῖν πρός με ἐὰν μὴ ὁ Πατὴρ ὁ πέμψας με ἑλκύσῃ αὐτόν, κἀγὼ ἀναστήσω αὐτὸν ἐν τῇ ἐσχάτῃ ἡμέρᾳ.
\vs{45}ἔστιν γεγραμμένον ἐν τοῖς προφήταις· Καὶ ἔσονται πάντες διδακτοὶ Θεοῦ· πᾶς ὁ ἀκούσας παρὰ τοῦ Πατρὸς καὶ μαθὼν ἔρχεται πρὸς ἐμέ.
\vs{46}οὐχ ὅτι τὸν Πατέρα ἑώρακέν τις εἰ μὴ ὁ ὢν παρὰ τοῦ Θεοῦ, οὗτος ἑώρακεν τὸν Πατέρα.
\vs{47}Ἀμὴν ἀμὴν λέγω ὑμῖν, ὁ πιστεύων ἔχει ζωὴν αἰώνιον.
\vs{48}ἐγώ εἰμι ὁ ἄρτος τῆς ζωῆς.
\vs{49}οἱ πατέρες ὑμῶν ἔφαγον ἐν τῇ ἐρήμῳ τὸ μάννα καὶ ἀπέθανον·
\vs{50}οὗτός ἐστιν ὁ ἄρτος ὁ ἐκ τοῦ οὐρανοῦ καταβαίνων, ἵνα τις ἐξ αὐτοῦ φάγῃ καὶ μὴ ἀποθάνῃ.
\vs{51}ἐγώ εἰμι ὁ ἄρτος ὁ ζῶν ὁ ἐκ τοῦ οὐρανοῦ καταβάς· ἐάν τις φάγῃ ἐκ τούτου τοῦ ἄρτου ζήσει εἰς τὸν αἰῶνα, καὶ ὁ ἄρτος δὲ ὃν ἐγὼ δώσω ἡ σάρξ μού ἐστιν ὑπὲρ τῆς τοῦ κόσμου ζωῆς.

\vs{52}Ἐμάχοντο οὖν πρὸς ἀλλήλους οἱ Ἰουδαῖοι λέγοντες· Πῶς δύναται οὗτος ἡμῖν δοῦναι τὴν σάρκα αὐτοῦ φαγεῖν;
\vs{53}Εἶπεν οὖν αὐτοῖς ὁ Ἰησοῦς· Ἀμὴν ἀμὴν λέγω ὑμῖν, ἐὰν μὴ φάγητε τὴν σάρκα τοῦ Υἱοῦ τοῦ ἀνθρώπου καὶ πίητε αὐτοῦ τὸ αἷμα, οὐκ ἔχετε ζωὴν ἐν ἑαυτοῖς.
\vs{54}ὁ τρώγων μου τὴν σάρκα καὶ πίνων μου τὸ αἷμα ἔχει ζωὴν αἰώνιον, κἀγὼ ἀναστήσω αὐτὸν τῇ ἐσχάτῃ ἡμέρᾳ.
\vs{55}ἡ γὰρ σάρξ μου ἀληθής ἐστιν βρῶσις, καὶ τὸ αἷμά μου ἀληθής ἐστιν πόσις.
\vs{56}Ὁ τρώγων μου τὴν σάρκα καὶ πίνων μου τὸ αἷμα ἐν ἐμοὶ μένει κἀγὼ ἐν αὐτῷ.
\vs{57}καθὼς ἀπέστειλέν με ὁ ζῶν Πατὴρ κἀγὼ ζῶ διὰ τὸν Πατέρα, καὶ ὁ τρώγων με κἀκεῖνος ζήσει δι᾽ ἐμέ.
\vs{58}οὗτός ἐστιν ὁ ἄρτος ὁ ἐξ οὐρανοῦ καταβάς, οὐ καθὼς ἔφαγον οἱ πατέρες καὶ ἀπέθανον· ὁ τρώγων τοῦτον τὸν ἄρτον ζήσει εἰς τὸν αἰῶνα.

\vs{59}Ταῦτα εἶπεν ἐν συναγωγῇ διδάσκων ἐν Καφαρναούμ.
\vs{60}Πολλοὶ οὖν ἀκούσαντες ἐκ τῶν μαθητῶν αὐτοῦ εἶπαν· Σκληρός ἐστιν ὁ λόγος οὗτος· τίς δύναται αὐτοῦ ἀκούειν;
\vs{61}Εἰδὼς δὲ ὁ Ἰησοῦς ἐν ἑαυτῷ ὅτι γογγύζουσιν περὶ τούτου οἱ μαθηταὶ αὐτοῦ εἶπεν αὐτοῖς· Τοῦτο ὑμᾶς σκανδαλίζει;
\vs{62}ἐὰν οὖν θεωρῆτε τὸν Υἱὸν τοῦ ἀνθρώπου ἀναβαίνοντα ὅπου ἦν τὸ πρότερον;
\vs{63}Τὸ πνεῦμά ἐστιν τὸ ζωοποιοῦν, ἡ σὰρξ οὐκ ὠφελεῖ οὐδέν· τὰ ῥήματα ἃ ἐγὼ λελάληκα ὑμῖν πνεῦμά ἐστιν καὶ ζωή ἐστιν.
\vs{64}ἀλλ᾽ εἰσὶν ἐξ ὑμῶν τινες οἳ οὐ πιστεύουσιν. ᾔδει γὰρ ἐξ ἀρχῆς ὁ Ἰησοῦς τίνες εἰσὶν οἱ μὴ πιστεύοντες καὶ τίς ἐστιν ὁ παραδώσων αὐτόν.
\vs{65}Καὶ ἔλεγεν· Διὰ τοῦτο εἴρηκα ὑμῖν ὅτι οὐδεὶς δύναται ἐλθεῖν πρός με ἐὰν μὴ ᾖ δεδομένον αὐτῷ ἐκ τοῦ Πατρός.
\vs{66}Ἐκ τούτου πολλοὶ ἐκ τῶν μαθητῶν αὐτοῦ ἀπῆλθον εἰς τὰ ὀπίσω καὶ οὐκέτι μετ᾽ αὐτοῦ περιεπάτουν.
\vs{67}εἶπεν οὖν ὁ Ἰησοῦς τοῖς δώδεκα· Μὴ καὶ ὑμεῖς θέλετε ὑπάγειν;
\vs{68}Ἀπεκρίθη αὐτῷ Σίμων Πέτρος· Κύριε, πρὸς τίνα ἀπελευσόμεθα; ῥήματα ζωῆς αἰωνίου ἔχεις,
\vs{69}καὶ ἡμεῖς πεπιστεύκαμεν καὶ ἐγνώκαμεν ὅτι σὺ εἶ ὁ Ἅγιος τοῦ Θεοῦ.
\vs{70}Ἀπεκρίθη αὐτοῖς ὁ Ἰησοῦς· Οὐκ ἐγὼ ὑμᾶς τοὺς δώδεκα ἐξελεξάμην; καὶ ἐξ ὑμῶν εἷς διάβολός ἐστιν.
\vs{71}ἔλεγεν δὲ τὸν Ἰούδαν Σίμωνος Ἰσκαριώτου· οὗτος γὰρ ἔμελλεν παραδιδόναι αὐτόν, εἷς ἐκ τῶν δώδεκα.

\ch{7}
Καὶ μετὰ ταῦτα περιεπάτει ὁ Ἰησοῦς ἐν τῇ Γαλιλαίᾳ· οὐ γὰρ ἤθελεν ἐν τῇ Ἰουδαίᾳ περιπατεῖν, ὅτι ἐζήτουν αὐτὸν οἱ Ἰουδαῖοι ἀποκτεῖναι.

\vs{2}ἦν δὲ ἐγγὺς ἡ ἑορτὴ τῶν Ἰουδαίων ἡ σκηνοπηγία.
\vs{3}εἶπον οὖν πρὸς αὐτὸν οἱ ἀδελφοὶ αὐτοῦ· Μετάβηθι ἐντεῦθεν καὶ ὕπαγε εἰς τὴν Ἰουδαίαν, ἵνα καὶ οἱ μαθηταί σου θεωρήσουσιν σοῦ τὰ ἔργα ἃ ποιεῖς·
\vs{4}οὐδεὶς γάρ τι ἐν κρυπτῷ ποιεῖ καὶ ζητεῖ αὐτὸς ἐν παρρησίᾳ εἶναι. εἰ ταῦτα ποιεῖς, φανέρωσον σεαυτὸν τῷ κόσμῳ.
\vs{5}οὐδὲ γὰρ οἱ ἀδελφοὶ αὐτοῦ ἐπίστευον εἰς αὐτόν.
\vs{6}Λέγει οὖν αὐτοῖς ὁ Ἰησοῦς· Ὁ καιρὸς ὁ ἐμὸς οὔπω πάρεστιν, ὁ δὲ καιρὸς ὁ ὑμέτερος πάντοτέ ἐστιν ἕτοιμος.
\vs{7}οὐ δύναται ὁ κόσμος μισεῖν ὑμᾶς, ἐμὲ δὲ μισεῖ, ὅτι ἐγὼ μαρτυρῶ περὶ αὐτοῦ ὅτι τὰ ἔργα αὐτοῦ πονηρά ἐστιν.
\vs{8}ὑμεῖς ἀνάβητε εἰς τὴν ἑορτήν· ἐγὼ οὐκ ἀναβαίνω εἰς τὴν ἑορτὴν ταύτην, ὅτι ὁ ἐμὸς καιρὸς οὔπω πεπλήρωται.
\vs{9}Ταῦτα δὲ εἰπὼν αὐτὸς ἔμεινεν ἐν τῇ Γαλιλαίᾳ.

\vs{10}Ὡς δὲ ἀνέβησαν οἱ ἀδελφοὶ αὐτοῦ εἰς τὴν ἑορτήν, τότε καὶ αὐτὸς ἀνέβη οὐ φανερῶς ἀλλὰ ὡς ἐν κρυπτῷ.
\vs{11}Οἱ οὖν Ἰουδαῖοι ἐζήτουν αὐτὸν ἐν τῇ ἑορτῇ καὶ ἔλεγον· Ποῦ ἐστιν ἐκεῖνος;
\vs{12}καὶ γογγυσμὸς περὶ αὐτοῦ ἦν πολὺς ἐν τοῖς ὄχλοις· οἱ μὲν ἔλεγον ὅτι Ἀγαθός ἐστιν, Ἄλλοι δὲ ἔλεγον· Οὔ, ἀλλὰ πλανᾷ τὸν ὄχλον.
\vs{13}Οὐδεὶς μέντοι παρρησίᾳ ἐλάλει περὶ αὐτοῦ διὰ τὸν φόβον τῶν Ἰουδαίων.

\vs{14}Ἤδη δὲ τῆς ἑορτῆς μεσούσης ἀνέβη Ἰησοῦς εἰς τὸ ἱερὸν καὶ ἐδίδασκεν.
\vs{15}ἐθαύμαζον οὖν οἱ Ἰουδαῖοι λέγοντες· Πῶς οὗτος γράμματα οἶδεν μὴ μεμαθηκώς;
\vs{16}Ἀπεκρίθη οὖν αὐτοῖς ὁ Ἰησοῦς καὶ εἶπεν· Ἡ Ἐμὴ διδαχὴ οὐκ ἔστιν ἐμὴ ἀλλὰ τοῦ πέμψαντός με·
\vs{17}ἐάν τις θέλῃ τὸ θέλημα αὐτοῦ ποιεῖν, γνώσεται περὶ τῆς διδαχῆς πότερον ἐκ τοῦ Θεοῦ ἐστιν ἢ ἐγὼ ἀπ᾽ ἐμαυτοῦ λαλῶ.
\vs{18}ὁ ἀφ᾽ ἑαυτοῦ λαλῶν τὴν δόξαν τὴν ἰδίαν ζητεῖ· ὁ δὲ ζητῶν τὴν δόξαν τοῦ πέμψαντος αὐτὸν οὗτος ἀληθής ἐστιν καὶ ἀδικία ἐν αὐτῷ οὐκ ἔστιν.

\vs{19}Οὐ Μωϋσῆς δέδωκεν ὑμῖν τὸν νόμον; καὶ οὐδεὶς ἐξ ὑμῶν ποιεῖ τὸν νόμον. τί με ζητεῖτε ἀποκτεῖναι;
\vs{20}Ἀπεκρίθη ὁ ὄχλος· Δαιμόνιον ἔχεις· τίς σε ζητεῖ ἀποκτεῖναι;
\vs{21}Ἀπεκρίθη Ἰησοῦς καὶ εἶπεν αὐτοῖς· Ἓν ἔργον ἐποίησα καὶ πάντες θαυμάζετε.
\vs{22}διὰ τοῦτο Μωϋσῆς δέδωκεν ὑμῖν τὴν περιτομήν— οὐχ ὅτι ἐκ τοῦ Μωϋσέως ἐστὶν ἀλλ᾽ ἐκ τῶν πατέρων— καὶ ἐν σαββάτῳ περιτέμνετε ἄνθρωπον.
\vs{23}εἰ περιτομὴν λαμβάνει ἄνθρωπος ἐν σαββάτῳ ἵνα μὴ λυθῇ ὁ νόμος Μωϋσέως, ἐμοὶ χολᾶτε ὅτι ὅλον ἄνθρωπον ὑγιῆ ἐποίησα ἐν σαββάτῳ;
\vs{24}μὴ κρίνετε κατ᾽ ὄψιν, ἀλλὰ τὴν δικαίαν κρίσιν κρίνετε.

\vs{25}Ἔλεγον οὖν τινες ἐκ τῶν Ἱεροσολυμιτῶν· Οὐχ οὗτός ἐστιν ὃν ζητοῦσιν ἀποκτεῖναι;
\vs{26}καὶ ἴδε παρρησίᾳ λαλεῖ καὶ οὐδὲν αὐτῷ λέγουσιν. Μήποτε ἀληθῶς ἔγνωσαν οἱ ἄρχοντες ὅτι οὗτός ἐστιν ὁ Χριστός;
\vs{27}ἀλλὰ τοῦτον οἴδαμεν πόθεν ἐστίν· ὁ δὲ Χριστὸς ὅταν ἔρχηται οὐδεὶς γινώσκει πόθεν ἐστίν.
\vs{28}Ἔκραξεν οὖν ἐν τῷ ἱερῷ διδάσκων ὁ Ἰησοῦς καὶ λέγων· Κἀμὲ οἴδατε καὶ οἴδατε πόθεν εἰμί· καὶ ἀπ᾽ ἐμαυτοῦ οὐκ ἐλήλυθα, ἀλλ᾽ ἔστιν ἀληθινὸς ὁ πέμψας με, ὃν ὑμεῖς οὐκ οἴδατε·
\vs{29}ἐγὼ οἶδα αὐτόν, ὅτι παρ᾽ αὐτοῦ εἰμι κἀκεῖνός με ἀπέστειλεν.
\vs{30}Ἐζήτουν οὖν αὐτὸν πιάσαι, καὶ οὐδεὶς ἐπέβαλεν ἐπ᾽ αὐτὸν τὴν χεῖρα, ὅτι οὔπω ἐληλύθει ἡ ὥρα αὐτοῦ.

\vs{31}Ἐκ τοῦ ὄχλου δὲ πολλοὶ ἐπίστευσαν εἰς αὐτόν καὶ ἔλεγον· Ὁ Χριστὸς ὅταν ἔλθῃ μὴ πλείονα σημεῖα ποιήσει ὧν οὗτος ἐποίησεν;
\vs{32}Ἤκουσαν οἱ Φαρισαῖοι τοῦ ὄχλου γογγύζοντος περὶ αὐτοῦ ταῦτα, καὶ ἀπέστειλαν οἱ ἀρχιερεῖς καὶ οἱ Φαρισαῖοι ὑπηρέτας ἵνα πιάσωσιν αὐτόν.
\vs{33}εἶπεν οὖν ὁ Ἰησοῦς· Ἔτι χρόνον μικρὸν μεθ᾽ ὑμῶν εἰμι καὶ ὑπάγω πρὸς τὸν πέμψαντά με.
\vs{34}ζητήσετέ με καὶ οὐχ εὑρήσετέ με, καὶ ὅπου εἰμὶ ἐγὼ ὑμεῖς οὐ δύνασθε ἐλθεῖν.
\vs{35}Εἶπον οὖν οἱ Ἰουδαῖοι πρὸς ἑαυτούς· Ποῦ οὗτος μέλλει πορεύεσθαι ὅτι ἡμεῖς οὐχ εὑρήσομεν αὐτόν; μὴ εἰς τὴν Διασπορὰν τῶν Ἑλλήνων μέλλει πορεύεσθαι καὶ διδάσκειν τοὺς Ἕλληνας;
\vs{36}τίς ἐστιν ὁ λόγος οὗτος ὃν εἶπεν· Ζητήσετέ με καὶ οὐχ εὑρήσετέ με, καὶ Ὅπου εἰμὶ ἐγὼ ὑμεῖς οὐ δύνασθε ἐλθεῖν;

\vs{37}Ἐν δὲ τῇ ἐσχάτῃ ἡμέρᾳ τῇ μεγάλῃ τῆς ἑορτῆς εἱστήκει ὁ Ἰησοῦς καὶ ἔκραξεν λέγων· Ἐάν τις διψᾷ ἐρχέσθω πρός με καὶ πινέτω.
\vs{38}ὁ πιστεύων εἰς ἐμέ, καθὼς εἶπεν ἡ γραφή, Ποταμοὶ ἐκ τῆς κοιλίας αὐτοῦ ῥεύσουσιν ὕδατος ζῶντος.
\vs{39}τοῦτο δὲ εἶπεν περὶ τοῦ Πνεύματος οὗ ἔμελλον λαμβάνειν οἱ πιστεύσαντες εἰς αὐτόν· οὔπω γὰρ ἦν Πνεῦμα, ὅτι Ἰησοῦς οὐδέπω ἐδοξάσθη.

\vs{40}Ἐκ τοῦ ὄχλου οὖν ἀκούσαντες τῶν λόγων τούτων ἔλεγον· Οὗτός ἐστιν ἀληθῶς ὁ προφήτης·
\vs{41}Ἄλλοι ἔλεγον· Οὗτός ἐστιν ὁ Χριστός, Οἱ δὲ ἔλεγον· Μὴ γὰρ ἐκ τῆς Γαλιλαίας ὁ Χριστὸς ἔρχεται;
\vs{42}οὐχ ἡ γραφὴ εἶπεν ὅτι ἐκ τοῦ σπέρματος Δαυὶδ καὶ ἀπὸ Βηθλεὲμ τῆς κώμης ὅπου ἦν Δαυὶδ ἔρχεται ὁ Χριστός;
\vs{43}Σχίσμα οὖν ἐγένετο ἐν τῷ ὄχλῳ δι᾽ αὐτόν·
\vs{44}τινὲς δὲ ἤθελον ἐξ αὐτῶν πιάσαι αὐτόν, ἀλλ᾽ οὐδεὶς ἐπέβαλεν ἐπ᾽ αὐτὸν τὰς χεῖρας.

\vs{45}Ἦλθον οὖν οἱ ὑπηρέται πρὸς τοὺς ἀρχιερεῖς καὶ Φαρισαίους, καὶ εἶπον αὐτοῖς ἐκεῖνοι· Διὰ τί οὐκ ἠγάγετε αὐτόν;
\vs{46}Ἀπεκρίθησαν οἱ ὑπηρέται· Οὐδέποτε ἐλάλησεν οὕτως ἄνθρωπος.
\vs{47}Ἀπεκρίθησαν οὖν αὐτοῖς οἱ Φαρισαῖοι· Μὴ καὶ ὑμεῖς πεπλάνησθε;
\vs{48}μή τις ἐκ τῶν ἀρχόντων ἐπίστευσεν εἰς αὐτὸν ἢ ἐκ τῶν Φαρισαίων;
\vs{49}ἀλλὰ ὁ ὄχλος οὗτος ὁ μὴ γινώσκων τὸν νόμον ἐπάρατοί εἰσιν.
\vs{50}Λέγει Νικόδημος πρὸς αὐτούς, ὁ ἐλθὼν πρὸς αὐτὸν τὸ πρότερον, εἷς ὢν ἐξ αὐτῶν·
\vs{51}Μὴ ὁ νόμος ἡμῶν κρίνει τὸν ἄνθρωπον ἐὰν μὴ ἀκούσῃ πρῶτον παρ᾽ αὐτοῦ καὶ γνῷ τί ποιεῖ;
\vs{52}Ἀπεκρίθησαν καὶ εἶπαν αὐτῷ· Μὴ καὶ σὺ ἐκ τῆς Γαλιλαίας εἶ; ἐραύνησον καὶ ἴδε ὅτι ἐκ τῆς Γαλιλαίας προφήτης οὐκ ἐγείρεται.

\vs{53}[[Καὶ ἐπορεύθησαν ἕκαστος εἰς τὸν οἶκον αὐτοῦ,
\inparch{8} \vs{1}Ἰησοῦς δὲ ἐπορεύθη εἰς τὸ ὄρος τῶν ἐλαιῶν.
\vs{2}Ὄρθρου δὲ πάλιν παρεγένετο εἰς τὸ ἱερόν καὶ πᾶς ὁ λαὸς ἤρχετο πρὸς αὐτόν, καὶ καθίσας ἐδίδασκεν αὐτούς.
\vs{3}ἄγουσιν δὲ οἱ γραμματεῖς καὶ οἱ Φαρισαῖοι γυναῖκα ἐπὶ μοιχείᾳ κατειλημμένην καὶ στήσαντες αὐτὴν ἐν μέσῳ
\vs{4}λέγουσιν αὐτῷ· Διδάσκαλε, αὕτη ἡ γυνὴ κατείληπται ἐπ᾽ αυτοφώρῳ μοιχευομένη·
\vs{5}ἐν δὲ τῷ νόμῳ ἡμῖν Μωϋσῆς ἐνετείλατο τὰς τοιαύτας λιθάζειν. σὺ οὖν τί λέγεις;
\vs{6}Τοῦτο δὲ ἔλεγον πειράζοντες αὐτόν, ἵνα ἔχωσιν κατηγορεῖν αὐτοῦ. ὁ δὲ Ἰησοῦς κάτω κύψας τῷ δακτύλῳ κατέγραφεν εἰς τὴν γῆν.
\vs{7}Ὡς δὲ ἐπέμενον ἐρωτῶντες αὐτόν, ἀνέκυψεν καὶ εἶπεν αὐτοῖς· Ὁ ἀναμάρτητος ὑμῶν πρῶτος ἐπ᾽ αὐτῇ βαλέτω λίθον.
\vs{8}καὶ πάλιν κατακύψας ἔγραφεν εἰς τὴν γῆν.
\vs{9}Οἱ δὲ ἀκούσαντες ἐξήρχοντο εἷς καθ εἷς ἀρξάμενοι ἀπὸ τῶν πρεσβυτέρων καὶ κατελείφθη μόνος καὶ ἡ γυνὴ ἐν μέσῳ οὖσα.
\vs{10}ἀνακύψας δὲ ὁ Ἰησοῦς εἶπεν αὐτῇ· Γύναι, ποῦ εἰσιν; οὐδείς σε κατέκρινεν;
\vs{11}Ἡ δὲ εἶπεν· Οὐδείς, κύριε. Εἶπεν δὲ ὁ Ἰησοῦς· Οὐδὲ ἐγώ σε κατακρίνω· πορεύου, καὶ ἀπὸ τοῦ νῦν μηκέτι ἁμάρτανε.]]

\vs{12}Πάλιν οὖν αὐτοῖς ἐλάλησεν ὁ Ἰησοῦς λέγων· Ἐγώ εἰμι τὸ φῶς τοῦ κόσμου· ὁ ἀκολουθῶν ἐμοὶ οὐ μὴ περιπατήσῃ ἐν τῇ σκοτίᾳ, ἀλλ᾽ ἕξει τὸ φῶς τῆς ζωῆς.
\vs{13}Εἶπον οὖν αὐτῷ οἱ Φαρισαῖοι· Σὺ περὶ σεαυτοῦ μαρτυρεῖς· ἡ μαρτυρία σου οὐκ ἔστιν ἀληθής.
\vs{14}Ἀπεκρίθη Ἰησοῦς καὶ εἶπεν αὐτοῖς· Κἂν ἐγὼ μαρτυρῶ περὶ ἐμαυτοῦ, ἀληθής ἐστιν ἡ μαρτυρία μου, ὅτι οἶδα πόθεν ἦλθον καὶ ποῦ ὑπάγω· ὑμεῖς δὲ οὐκ οἴδατε πόθεν ἔρχομαι ἢ ποῦ ὑπάγω.
\vs{15}Ὑμεῖς κατὰ τὴν σάρκα κρίνετε, ἐγὼ οὐ κρίνω οὐδένα.
\vs{16}καὶ ἐὰν κρίνω δὲ ἐγώ, ἡ κρίσις ἡ ἐμὴ ἀληθινή ἐστιν, ὅτι μόνος οὐκ εἰμί, ἀλλ᾽ ἐγὼ καὶ ὁ πέμψας με πατήρ.
\vs{17}Καὶ ἐν τῷ νόμῳ δὲ τῷ ὑμετέρῳ γέγραπται ὅτι δύο ἀνθρώπων ἡ μαρτυρία ἀληθής ἐστιν.
\vs{18}ἐγώ εἰμι ὁ μαρτυρῶν περὶ ἐμαυτοῦ καὶ μαρτυρεῖ περὶ ἐμοῦ ὁ πέμψας με Πατήρ.
\vs{19}Ἔλεγον οὖν αὐτῷ· Ποῦ ἐστιν ὁ Πατήρ σου; Ἀπεκρίθη Ἰησοῦς· Οὔτε ἐμὲ οἴδατε οὔτε τὸν Πατέρα μου· εἰ ἐμὲ ᾔδειτε, καὶ τὸν Πατέρα μου ἂν ᾔδειτε.
\vs{20}Ταῦτα τὰ ῥήματα ἐλάλησεν ἐν τῷ γαζοφυλακίῳ διδάσκων ἐν τῷ ἱερῷ· καὶ οὐδεὶς ἐπίασεν αὐτόν, ὅτι οὔπω ἐληλύθει ἡ ὥρα αὐτοῦ.

\vs{21}Εἶπεν οὖν πάλιν αὐτοῖς· Ἐγὼ ὑπάγω καὶ ζητήσετέ με, καὶ ἐν τῇ ἁμαρτίᾳ ὑμῶν ἀποθανεῖσθε· ὅπου ἐγὼ ὑπάγω ὑμεῖς οὐ δύνασθε ἐλθεῖν.
\vs{22}Ἔλεγον οὖν οἱ Ἰουδαῖοι· Μήτι ἀποκτενεῖ ἑαυτὸν, ὅτι λέγει· Ὅπου ἐγὼ ὑπάγω ὑμεῖς οὐ δύνασθε ἐλθεῖν;
\vs{23}Καὶ ἔλεγεν αὐτοῖς· Ὑμεῖς ἐκ τῶν κάτω ἐστέ, ἐγὼ ἐκ τῶν ἄνω εἰμί· ὑμεῖς ἐκ τούτου τοῦ κόσμου ἐστέ, ἐγὼ οὐκ εἰμὶ ἐκ τοῦ κόσμου τούτου.
\vs{24}εἶπον οὖν ὑμῖν ὅτι ἀποθανεῖσθε ἐν ταῖς ἁμαρτίαις ὑμῶν· ἐὰν γὰρ μὴ πιστεύσητε ὅτι ἐγώ εἰμι, ἀποθανεῖσθε ἐν ταῖς ἁμαρτίαις ὑμῶν.
\vs{25}Ἔλεγον οὖν αὐτῷ· Σὺ τίς εἶ; Εἶπεν αὐτοῖς ὁ Ἰησοῦς· Τὴν ἀρχὴν ὅ τι καὶ λαλῶ ὑμῖν;
\vs{26}πολλὰ ἔχω περὶ ὑμῶν λαλεῖν καὶ κρίνειν, ἀλλ᾽ ὁ πέμψας με ἀληθής ἐστιν, κἀγὼ ἃ ἤκουσα παρ᾽ αὐτοῦ ταῦτα λαλῶ εἰς τὸν κόσμον.
\vs{27}Οὐκ ἔγνωσαν ὅτι τὸν Πατέρα αὐτοῖς ἔλεγεν.
\vs{28}εἶπεν οὖν αὐτοῖς ὁ Ἰησοῦς· Ὅταν ὑψώσητε τὸν Υἱὸν τοῦ ἀνθρώπου, τότε γνώσεσθε ὅτι ἐγώ εἰμι, καὶ ἀπ᾽ ἐμαυτοῦ ποιῶ οὐδέν, ἀλλὰ καθὼς ἐδίδαξέν με ὁ Πατὴρ ταῦτα λαλῶ.
\vs{29}καὶ ὁ πέμψας με μετ᾽ ἐμοῦ ἐστιν· οὐκ ἀφῆκέν με μόνον, ὅτι ἐγὼ τὰ ἀρεστὰ αὐτῷ ποιῶ πάντοτε.

\vs{30}Ταῦτα αὐτοῦ λαλοῦντος πολλοὶ ἐπίστευσαν εἰς αὐτόν.
\vs{31}Ἔλεγεν οὖν ὁ Ἰησοῦς πρὸς τοὺς πεπιστευκότας αὐτῷ Ἰουδαίους· Ἐὰν ὑμεῖς μείνητε ἐν τῷ λόγῳ τῷ ἐμῷ, ἀληθῶς μαθηταί μού ἐστε
\vs{32}καὶ γνώσεσθε τὴν ἀλήθειαν, καὶ ἡ ἀλήθεια ἐλευθερώσει ὑμᾶς.
\vs{33}Ἀπεκρίθησαν πρὸς αὐτόν· Σπέρμα Ἀβραάμ ἐσμεν καὶ οὐδενὶ δεδουλεύκαμεν πώποτε· πῶς σὺ λέγεις ὅτι Ἐλεύθεροι γενήσεσθε;
\vs{34}Ἀπεκρίθη αὐτοῖς ὁ Ἰησοῦς· Ἀμὴν ἀμὴν λέγω ὑμῖν ὅτι πᾶς ὁ ποιῶν τὴν ἁμαρτίαν δοῦλός ἐστιν τῆς ἁμαρτίας.
\vs{35}ὁ δὲ δοῦλος οὐ μένει ἐν τῇ οἰκίᾳ εἰς τὸν αἰῶνα, ὁ υἱὸς μένει εἰς τὸν αἰῶνα.
\vs{36}ἐὰν οὖν ὁ Υἱὸς ὑμᾶς ἐλευθερώσῃ, ὄντως ἐλεύθεροι ἔσεσθε.

\vs{37}Οἶδα ὅτι σπέρμα Ἀβραάμ ἐστε· ἀλλὰ ζητεῖτέ με ἀποκτεῖναι, ὅτι ὁ λόγος ὁ ἐμὸς οὐ χωρεῖ ἐν ὑμῖν.
\vs{38}ἃ ἐγὼ ἑώρακα παρὰ τῷ Πατρὶ λαλῶ· καὶ ὑμεῖς οὖν ἃ ἠκούσατε παρὰ τοῦ πατρὸς ποιεῖτε.
\vs{39}Ἀπεκρίθησαν καὶ εἶπαν αὐτῷ· Ὁ πατὴρ ἡμῶν Ἀβραάμ ἐστιν. Λέγει αὐτοῖς ὁ Ἰησοῦς· Εἰ τέκνα τοῦ Ἀβραάμ ἐστε, τὰ ἔργα τοῦ Ἀβραὰμ ἐποιεῖτε·
\vs{40}νῦν δὲ ζητεῖτέ με ἀποκτεῖναι ἄνθρωπον ὃς τὴν ἀλήθειαν ὑμῖν λελάληκα ἣν ἤκουσα παρὰ τοῦ Θεοῦ· τοῦτο Ἀβραὰμ οὐκ ἐποίησεν.
\vs{41}ὑμεῖς ποιεῖτε τὰ ἔργα τοῦ πατρὸς ὑμῶν. Εἶπαν οὖν αὐτῷ· Ἡμεῖς ἐκ πορνείας οὐ γεγεννήμεθα, ἕνα Πατέρα ἔχομεν τὸν Θεόν.
\vs{42}Εἶπεν αὐτοῖς ὁ Ἰησοῦς· Εἰ ὁ Θεὸς Πατὴρ ὑμῶν ἦν ἠγαπᾶτε ἂν ἐμέ, ἐγὼ γὰρ ἐκ τοῦ Θεοῦ ἐξῆλθον καὶ ἥκω· οὐδὲ γὰρ ἀπ᾽ ἐμαυτοῦ ἐλήλυθα, ἀλλ᾽ ἐκεῖνός με ἀπέστειλεν.
\vs{43}Διὰ τί τὴν λαλιὰν τὴν ἐμὴν οὐ γινώσκετε; ὅτι οὐ δύνασθε ἀκούειν τὸν λόγον τὸν ἐμόν.
\vs{44}ὑμεῖς ἐκ τοῦ πατρὸς τοῦ διαβόλου ἐστὲ καὶ τὰς ἐπιθυμίας τοῦ πατρὸς ὑμῶν θέλετε ποιεῖν. ἐκεῖνος ἀνθρωποκτόνος ἦν ἀπ᾽ ἀρχῆς καὶ ἐν τῇ ἀληθείᾳ οὐκ ἔστηκεν, ὅτι οὐκ ἔστιν ἀλήθεια ἐν αὐτῷ. ὅταν λαλῇ τὸ ψεῦδος, ἐκ τῶν ἰδίων λαλεῖ, ὅτι ψεύστης ἐστὶν καὶ ὁ πατὴρ αὐτοῦ.
\vs{45}ἐγὼ δὲ ὅτι τὴν ἀλήθειαν λέγω, οὐ πιστεύετέ μοι.
\vs{46}Τίς ἐξ ὑμῶν ἐλέγχει με περὶ ἁμαρτίας; εἰ ἀλήθειαν λέγω, διὰ τί ὑμεῖς οὐ πιστεύετέ μοι;
\vs{47}ὁ ὢν ἐκ τοῦ Θεοῦ τὰ ῥήματα τοῦ Θεοῦ ἀκούει· διὰ τοῦτο ὑμεῖς οὐκ ἀκούετε, ὅτι ἐκ τοῦ Θεοῦ οὐκ ἐστέ.

\vs{48}Ἀπεκρίθησαν οἱ Ἰουδαῖοι καὶ εἶπαν αὐτῷ· Οὐ καλῶς λέγομεν ἡμεῖς ὅτι Σαμαρίτης εἶ σὺ καὶ δαιμόνιον ἔχεις;
\vs{49}Ἀπεκρίθη Ἰησοῦς· Ἐγὼ δαιμόνιον οὐκ ἔχω, ἀλλὰ τιμῶ τὸν Πατέρα μου, καὶ ὑμεῖς ἀτιμάζετέ με.
\vs{50}ἐγὼ δὲ οὐ ζητῶ τὴν δόξαν μου· ἔστιν ὁ ζητῶν καὶ κρίνων.
\vs{51}ἀμὴν ἀμὴν λέγω ὑμῖν, ἐάν τις τὸν ἐμὸν λόγον τηρήσῃ, θάνατον οὐ μὴ θεωρήσῃ εἰς τὸν αἰῶνα.
\vs{52}Εἶπον οὖν αὐτῷ οἱ Ἰουδαῖοι· Νῦν ἐγνώκαμεν ὅτι δαιμόνιον ἔχεις. Ἀβραὰμ ἀπέθανεν καὶ οἱ προφῆται, καὶ σὺ λέγεις· Ἐάν τις τὸν λόγον μου τηρήσῃ, οὐ μὴ γεύσηται θανάτου εἰς τὸν αἰῶνα.
\vs{53}μὴ σὺ μείζων εἶ τοῦ πατρὸς ἡμῶν Ἀβραάμ, ὅστις ἀπέθανεν; καὶ οἱ προφῆται ἀπέθανον. τίνα σεαυτὸν ποιεῖς;
\vs{54}Ἀπεκρίθη Ἰησοῦς· Ἐὰν ἐγὼ δοξάσω ἐμαυτόν, ἡ δόξα μου οὐδέν ἐστιν· ἔστιν ὁ Πατήρ μου ὁ δοξάζων με, ὃν ὑμεῖς λέγετε ὅτι Θεὸς ἡμῶν ἐστιν,
\vs{55}καὶ οὐκ ἐγνώκατε αὐτόν, ἐγὼ δὲ οἶδα αὐτόν. κἂν εἴπω ὅτι οὐκ οἶδα αὐτόν, ἔσομαι ὅμοιος ὑμῖν ψεύστης· ἀλλὰ οἶδα αὐτὸν καὶ τὸν λόγον αὐτοῦ τηρῶ.
\vs{56}Ἀβραὰμ ὁ πατὴρ ὑμῶν ἠγαλλιάσατο ἵνα ἴδῃ τὴν ἡμέραν τὴν ἐμήν, καὶ εἶδεν καὶ ἐχάρη.
\vs{57}Εἶπον οὖν οἱ Ἰουδαῖοι πρὸς αὐτόν· Πεντήκοντα ἔτη οὔπω ἔχεις καὶ Ἀβραὰμ ἑώρακας;
\vs{58}Εἶπεν αὐτοῖς Ἰησοῦς· Ἀμὴν ἀμὴν λέγω ὑμῖν, πρὶν Ἀβραὰμ γενέσθαι ἐγὼ εἰμί.
\vs{59}Ἦραν οὖν λίθους ἵνα βάλωσιν ἐπ᾽ αὐτόν. Ἰησοῦς δὲ ἐκρύβη καὶ ἐξῆλθεν ἐκ τοῦ ἱεροῦ.

\ch{9}
Καὶ παράγων εἶδεν ἄνθρωπον τυφλὸν ἐκ γενετῆς.
\vs{2}καὶ ἠρώτησαν αὐτὸν οἱ μαθηταὶ αὐτοῦ λέγοντες· Ῥαββί, τίς ἥμαρτεν, οὗτος ἢ οἱ γονεῖς αὐτοῦ, ἵνα τυφλὸς γεννηθῇ;
\vs{3}Ἀπεκρίθη Ἰησοῦς· Οὔτε οὗτος ἥμαρτεν οὔτε οἱ γονεῖς αὐτοῦ, ἀλλ᾽ ἵνα φανερωθῇ τὰ ἔργα τοῦ Θεοῦ ἐν αὐτῷ.
\vs{4}ἡμᾶς δεῖ ἐργάζεσθαι τὰ ἔργα τοῦ πέμψαντός με ἕως ἡμέρα ἐστίν· ἔρχεται νὺξ ὅτε οὐδεὶς δύναται ἐργάζεσθαι.
\vs{5}ὅταν ἐν τῷ κόσμῳ ὦ, φῶς εἰμι τοῦ κόσμου.
\vs{6}Ταῦτα εἰπὼν ἔπτυσεν χαμαὶ καὶ ἐποίησεν πηλὸν ἐκ τοῦ πτύσματος καὶ ἐπέχρισεν αὐτοῦ τὸν πηλὸν ἐπὶ τοὺς ὀφθαλμούς
\vs{7}καὶ εἶπεν αὐτῷ· Ὕπαγε νίψαι εἰς τὴν κολυμβήθραν τοῦ Σιλωάμ ὃ ἑρμηνεύεται Ἀπεσταλμένος. ἀπῆλθεν οὖν καὶ ἐνίψατο καὶ ἦλθεν βλέπων.
\vs{8}Οἱ οὖν γείτονες καὶ οἱ θεωροῦντες αὐτὸν τὸ πρότερον ὅτι προσαίτης ἦν ἔλεγον· Οὐχ οὗτός ἐστιν ὁ καθήμενος καὶ προσαιτῶν;
\vs{9}Ἄλλοι ἔλεγον ὅτι Οὗτός ἐστιν, ἄλλοι ἔλεγον· Οὐχί, ἀλλὰ ὅμοιος αὐτῷ ἐστιν. Ἐκεῖνος ἔλεγεν ὅτι Ἐγώ εἰμι.
\vs{10}Ἔλεγον οὖν αὐτῷ· Πῶς οὖν ἠνεῴχθησάν σου οἱ ὀφθαλμοί;
\vs{11}Ἀπεκρίθη ἐκεῖνος· Ὁ ἄνθρωπος ὁ λεγόμενος Ἰησοῦς πηλὸν ἐποίησεν καὶ ἐπέχρισέν μου τοὺς ὀφθαλμοὺς καὶ εἶπέν μοι ὅτι Ὕπαγε εἰς τὸν Σιλωὰμ καὶ νίψαι· ἀπελθὼν οὖν καὶ νιψάμενος ἀνέβλεψα.
\vs{12}Καὶ εἶπαν αὐτῷ· Ποῦ ἐστιν ἐκεῖνος; Λέγει· Οὐκ οἶδα.

\vs{13}Ἄγουσιν αὐτὸν πρὸς τοὺς Φαρισαίους τόν ποτε τυφλόν.
\vs{14}ἦν δὲ σάββατον ἐν ᾗ ἡμέρᾳ τὸν πηλὸν ἐποίησεν ὁ Ἰησοῦς καὶ ἀνέῳξεν αὐτοῦ τοὺς ὀφθαλμούς.
\vs{15}πάλιν οὖν ἠρώτων αὐτὸν καὶ οἱ Φαρισαῖοι πῶς ἀνέβλεψεν. Ὁ δὲ εἶπεν αὐτοῖς· Πηλὸν ἐπέθηκέν μου ἐπὶ τοὺς ὀφθαλμούς καὶ ἐνιψάμην καὶ βλέπω.
\vs{16}Ἔλεγον οὖν ἐκ τῶν Φαρισαίων τινές· Οὐκ ἔστιν οὗτος παρὰ Θεοῦ ὁ ἄνθρωπος, ὅτι τὸ σάββατον οὐ τηρεῖ. Ἄλλοι δὲ ἔλεγον· Πῶς δύναται ἄνθρωπος ἁμαρτωλὸς τοιαῦτα σημεῖα ποιεῖν; Καὶ σχίσμα ἦν ἐν αὐτοῖς.
\vs{17}λέγουσιν οὖν τῷ τυφλῷ πάλιν· Τί σὺ λέγεις περὶ αὐτοῦ, ὅτι ἠνέῳξέν σου τοὺς ὀφθαλμούς; Ὁ δὲ εἶπεν ὅτι Προφήτης ἐστίν.

\vs{18}Οὐκ ἐπίστευσαν οὖν οἱ Ἰουδαῖοι περὶ αὐτοῦ ὅτι ἦν τυφλὸς καὶ ἀνέβλεψεν ἕως ὅτου ἐφώνησαν τοὺς γονεῖς αὐτοῦ τοῦ ἀναβλέψαντος
\vs{19}καὶ ἠρώτησαν αὐτοὺς λέγοντες· Οὗτός ἐστιν ὁ υἱὸς ὑμῶν, ὃν ὑμεῖς λέγετε ὅτι τυφλὸς ἐγεννήθη; πῶς οὖν βλέπει ἄρτι;
\vs{20}Ἀπεκρίθησαν οὖν οἱ γονεῖς αὐτοῦ καὶ εἶπαν· Οἴδαμεν ὅτι οὗτός ἐστιν ὁ υἱὸς ἡμῶν καὶ ὅτι τυφλὸς ἐγεννήθη·
\vs{21}πῶς δὲ νῦν βλέπει οὐκ οἴδαμεν, ἢ τίς ἤνοιξεν αὐτοῦ τοὺς ὀφθαλμοὺς ἡμεῖς οὐκ οἴδαμεν· αὐτὸν ἐρωτήσατε, ἡλικίαν ἔχει, αὐτὸς περὶ ἑαυτοῦ λαλήσει.
\vs{22}Ταῦτα εἶπαν οἱ γονεῖς αὐτοῦ ὅτι ἐφοβοῦντο τοὺς Ἰουδαίους· ἤδη γὰρ συνετέθειντο οἱ Ἰουδαῖοι ἵνα ἐάν τις αὐτὸν ὁμολογήσῃ Χριστόν, ἀποσυνάγωγος γένηται.
\vs{23}διὰ τοῦτο οἱ γονεῖς αὐτοῦ εἶπαν ὅτι Ἡλικίαν ἔχει, αὐτὸν ἐπερωτήσατε.

\vs{24}Ἐφώνησαν οὖν τὸν ἄνθρωπον ἐκ δευτέρου ὃς ἦν τυφλὸς καὶ εἶπαν αὐτῷ· Δὸς δόξαν τῷ Θεῷ· ἡμεῖς οἴδαμεν ὅτι οὗτος ὁ ἄνθρωπος ἁμαρτωλός ἐστιν.
\vs{25}Ἀπεκρίθη οὖν ἐκεῖνος· Εἰ ἁμαρτωλός ἐστιν οὐκ οἶδα· ἓν οἶδα ὅτι τυφλὸς ὢν ἄρτι βλέπω.
\vs{26}Εἶπον οὖν αὐτῷ· Τί ἐποίησέν σοι; πῶς ἤνοιξέν σου τοὺς ὀφθαλμούς;
\vs{27}Ἀπεκρίθη αὐτοῖς· Εἶπον ὑμῖν ἤδη καὶ οὐκ ἠκούσατε· τί πάλιν θέλετε ἀκούειν; μὴ καὶ ὑμεῖς θέλετε αὐτοῦ μαθηταὶ γενέσθαι;
\vs{28}Καὶ ἐλοιδόρησαν αὐτὸν καὶ εἶπον· Σὺ μαθητὴς εἶ ἐκείνου, ἡμεῖς δὲ τοῦ Μωϋσέως ἐσμὲν μαθηταί·
\vs{29}ἡμεῖς οἴδαμεν ὅτι Μωϋσεῖ λελάληκεν ὁ Θεός, τοῦτον δὲ οὐκ οἴδαμεν πόθεν ἐστίν.
\vs{30}Ἀπεκρίθη ὁ ἄνθρωπος καὶ εἶπεν αὐτοῖς· Ἐν τούτῳ γὰρ τὸ θαυμαστόν ἐστιν, ὅτι ὑμεῖς οὐκ οἴδατε πόθεν ἐστίν, καὶ ἤνοιξέν μου τοὺς ὀφθαλμούς.
\vs{31}οἴδαμεν ὅτι ἁμαρτωλῶν ὁ Θεὸς οὐκ ἀκούει, ἀλλ᾽ ἐάν τις θεοσεβὴς ᾖ καὶ τὸ θέλημα αὐτοῦ ποιῇ τούτου ἀκούει.
\vs{32}ἐκ τοῦ αἰῶνος οὐκ ἠκούσθη ὅτι ἠνέῳξέν τις ὀφθαλμοὺς τυφλοῦ γεγεννημένου·
\vs{33}εἰ μὴ ἦν οὗτος παρὰ Θεοῦ, οὐκ ἠδύνατο ποιεῖν οὐδέν.
\vs{34}Ἀπεκρίθησαν καὶ εἶπαν αὐτῷ· Ἐν ἁμαρτίαις σὺ ἐγεννήθης ὅλος καὶ σὺ διδάσκεις ἡμᾶς; καὶ ἐξέβαλον αὐτὸν ἔξω.

\vs{35}Ἤκουσεν Ἰησοῦς ὅτι ἐξέβαλον αὐτὸν ἔξω καὶ εὑρὼν αὐτὸν εἶπεν· Σὺ πιστεύεις εἰς τὸν Υἱὸν τοῦ ἀνθρώπου;
\vs{36}Ἀπεκρίθη ἐκεῖνος καὶ εἶπεν· Καὶ τίς ἐστιν, Κύριε, ἵνα πιστεύσω εἰς αὐτόν;
\vs{37}Εἶπεν αὐτῷ ὁ Ἰησοῦς· Καὶ ἑώρακας αὐτὸν καὶ ὁ λαλῶν μετὰ σοῦ ἐκεῖνός ἐστιν.
\vs{38}Ὁ δὲ ἔφη· Πιστεύω, Κύριε· καὶ προσεκύνησεν αὐτῷ.

\vs{39}Καὶ εἶπεν ὁ Ἰησοῦς· Εἰς κρίμα ἐγὼ εἰς τὸν κόσμον τοῦτον ἦλθον, ἵνα οἱ μὴ βλέποντες βλέπωσιν καὶ οἱ βλέποντες τυφλοὶ γένωνται.
\vs{40}ἤκουσαν ἐκ τῶν Φαρισαίων ταῦτα οἱ μετ᾽ αὐτοῦ ὄντες καὶ εἶπον αὐτῷ· Μὴ καὶ ἡμεῖς τυφλοί ἐσμεν;
\vs{41}Εἶπεν αὐτοῖς ὁ Ἰησοῦς· Εἰ τυφλοὶ ἦτε, οὐκ ἂν εἴχετε ἁμαρτίαν· νῦν δὲ λέγετε ὅτι Βλέπομεν, ἡ ἁμαρτία ὑμῶν μένει.

\ch{10}
Ἀμὴν ἀμὴν λέγω ὑμῖν, ὁ μὴ εἰσερχόμενος διὰ τῆς θύρας εἰς τὴν αὐλὴν τῶν προβάτων ἀλλὰ ἀναβαίνων ἀλλαχόθεν ἐκεῖνος κλέπτης ἐστὶν καὶ λῃστής·
\vs{2}ὁ δὲ εἰσερχόμενος διὰ τῆς θύρας ποιμήν ἐστιν τῶν προβάτων.
\vs{3}τούτῳ ὁ θυρωρὸς ἀνοίγει καὶ τὰ πρόβατα τῆς φωνῆς αὐτοῦ ἀκούει καὶ τὰ ἴδια πρόβατα φωνεῖ κατ᾽ ὄνομα καὶ ἐξάγει αὐτά.
\vs{4}Ὅταν τὰ ἴδια πάντα ἐκβάλῃ, ἔμπροσθεν αὐτῶν πορεύεται καὶ τὰ πρόβατα αὐτῷ ἀκολουθεῖ, ὅτι οἴδασιν τὴν φωνὴν αὐτοῦ·
\vs{5}ἀλλοτρίῳ δὲ οὐ μὴ ἀκολουθήσουσιν, ἀλλὰ φεύξονται ἀπ᾽ αὐτοῦ, ὅτι οὐκ οἴδασιν τῶν ἀλλοτρίων τὴν φωνήν.
\vs{6}Ταύτην τὴν παροιμίαν εἶπεν αὐτοῖς ὁ Ἰησοῦς, ἐκεῖνοι δὲ οὐκ ἔγνωσαν τίνα ἦν ἃ ἐλάλει αὐτοῖς.

\vs{7}Εἶπεν οὖν πάλιν ὁ Ἰησοῦς· Ἀμὴν ἀμὴν λέγω ὑμῖν ὅτι ἐγώ εἰμι ἡ θύρα τῶν προβάτων.
\vs{8}πάντες ὅσοι ἦλθον πρὸ ἐμοῦ κλέπται εἰσὶν καὶ λῃσταί, ἀλλ᾽ οὐκ ἤκουσαν αὐτῶν τὰ πρόβατα.
\vs{9}ἐγώ εἰμι ἡ θύρα· δι᾽ ἐμοῦ ἐάν τις εἰσέλθῃ σωθήσεται καὶ εἰσελεύσεται καὶ ἐξελεύσεται καὶ νομὴν εὑρήσει.
\vs{10}ὁ κλέπτης οὐκ ἔρχεται εἰ μὴ ἵνα κλέψῃ καὶ θύσῃ καὶ ἀπολέσῃ· ἐγὼ ἦλθον ἵνα ζωὴν ἔχωσιν καὶ περισσὸν ἔχωσιν.

\vs{11}Ἐγώ εἰμι ὁ ποιμὴν ὁ καλός. ὁ ποιμὴν ὁ καλὸς τὴν ψυχὴν αὐτοῦ τίθησιν ὑπὲρ τῶν προβάτων·
\vs{12}ὁ μισθωτὸς καὶ οὐκ ὢν ποιμήν, οὗ οὐκ ἔστιν τὰ πρόβατα ἴδια, θεωρεῖ τὸν λύκον ἐρχόμενον καὶ ἀφίησιν τὰ πρόβατα καὶ φεύγει— καὶ ὁ λύκος ἁρπάζει αὐτὰ καὶ σκορπίζει—
\vs{13}ὅτι μισθωτός ἐστιν καὶ οὐ μέλει αὐτῷ περὶ τῶν προβάτων.

\vs{14}Ἐγώ εἰμι ὁ ποιμὴν ὁ καλός καὶ γινώσκω τὰ ἐμὰ καὶ γινώσκουσί με τὰ ἐμά,
\vs{15}καθὼς γινώσκει με ὁ Πατὴρ κἀγὼ γινώσκω τὸν Πατέρα, καὶ τὴν ψυχήν μου τίθημι ὑπὲρ τῶν προβάτων.
\vs{16}καὶ ἄλλα πρόβατα ἔχω ἃ οὐκ ἔστιν ἐκ τῆς αὐλῆς ταύτης· κἀκεῖνα δεῖ με ἀγαγεῖν καὶ τῆς φωνῆς μου ἀκούσουσιν, καὶ γενήσονται μία ποίμνη, εἷς ποιμήν.

\vs{17}Διὰ τοῦτό με ὁ Πατὴρ ἀγαπᾷ ὅτι ἐγὼ τίθημι τὴν ψυχήν μου, ἵνα πάλιν λάβω αὐτήν.
\vs{18}οὐδεὶς αἴρει αὐτὴν ἀπ᾽ ἐμοῦ, ἀλλ᾽ ἐγὼ τίθημι αὐτὴν ἀπ᾽ ἐμαυτοῦ. ἐξουσίαν ἔχω θεῖναι αὐτήν, καὶ ἐξουσίαν ἔχω πάλιν λαβεῖν αὐτήν· ταύτην τὴν ἐντολὴν ἔλαβον παρὰ τοῦ Πατρός μου.

\vs{19}Σχίσμα πάλιν ἐγένετο ἐν τοῖς Ἰουδαίοις διὰ τοὺς λόγους τούτους.
\vs{20}ἔλεγον δὲ πολλοὶ ἐξ αὐτῶν· Δαιμόνιον ἔχει καὶ μαίνεται· τί αὐτοῦ ἀκούετε;
\vs{21}Ἄλλοι ἔλεγον· Ταῦτα τὰ ῥήματα οὐκ ἔστιν δαιμονιζομένου· μὴ δαιμόνιον δύναται τυφλῶν ὀφθαλμοὺς ἀνοῖξαι;

\vs{22}Ἐγένετο τότε τὰ ἐνκαίνια ἐν τοῖς Ἱεροσολύμοις, χειμὼν ἦν,
\vs{23}καὶ περιεπάτει ὁ Ἰησοῦς ἐν τῷ ἱερῷ ἐν τῇ στοᾷ τοῦ Σολομῶνος.
\vs{24}ἐκύκλωσαν οὖν αὐτὸν οἱ Ἰουδαῖοι καὶ ἔλεγον αὐτῷ· Ἕως πότε τὴν ψυχὴν ἡμῶν αἴρεις; εἰ σὺ εἶ ὁ Χριστός, εἰπὲ ἡμῖν παρρησίᾳ.
\vs{25}Ἀπεκρίθη αὐτοῖς ὁ Ἰησοῦς· Εἶπον ὑμῖν καὶ οὐ πιστεύετε· τὰ ἔργα ἃ ἐγὼ ποιῶ ἐν τῷ ὀνόματι τοῦ Πατρός μου ταῦτα μαρτυρεῖ περὶ ἐμοῦ·
\vs{26}ἀλλὰ ὑμεῖς οὐ πιστεύετε, ὅτι οὐκ ἐστὲ ἐκ τῶν προβάτων τῶν ἐμῶν.
\vs{27}τὰ πρόβατα τὰ ἐμὰ τῆς φωνῆς μου ἀκούουσιν, κἀγὼ γινώσκω αὐτά καὶ ἀκολουθοῦσίν μοι,
\vs{28}κἀγὼ δίδωμι αὐτοῖς ζωὴν αἰώνιον καὶ οὐ μὴ ἀπόλωνται εἰς τὸν αἰῶνα καὶ οὐχ ἁρπάσει τις αὐτὰ ἐκ τῆς χειρός μου.
\vs{29}ὁ Πατήρ μου ὃ δέδωκέν μοι πάντων μεῖζόν ἐστιν, καὶ οὐδεὶς δύναται ἁρπάζειν ἐκ τῆς χειρὸς τοῦ Πατρός.
\vs{30}ἐγὼ καὶ ὁ Πατὴρ ἕν ἐσμεν.

\vs{31}Ἐβάστασαν πάλιν λίθους οἱ Ἰουδαῖοι ἵνα λιθάσωσιν αὐτόν.
\vs{32}ἀπεκρίθη αὐτοῖς ὁ Ἰησοῦς· Πολλὰ ἔργα καλὰ ἔδειξα ὑμῖν ἐκ τοῦ Πατρός· διὰ ποῖον αὐτῶν ἔργον ἐμὲ λιθάζετε;
\vs{33}Ἀπεκρίθησαν αὐτῷ οἱ Ἰουδαῖοι· Περὶ καλοῦ ἔργου οὐ λιθάζομέν σε ἀλλὰ περὶ βλασφημίας, καὶ ὅτι σὺ ἄνθρωπος ὢν ποιεῖς σεαυτὸν Θεόν.
\vs{34}Ἀπεκρίθη αὐτοῖς ὁ Ἰησοῦς· Οὐκ ἔστιν γεγραμμένον ἐν τῷ νόμῳ ὑμῶν ὅτι Ἐγὼ εἶπα· Θεοί ἐστε;
\vs{35}εἰ ἐκείνους εἶπεν θεοὺς πρὸς οὓς ὁ λόγος τοῦ Θεοῦ ἐγένετο, καὶ οὐ δύναται λυθῆναι ἡ γραφή,
\vs{36}ὃν ὁ Πατὴρ ἡγίασεν καὶ ἀπέστειλεν εἰς τὸν κόσμον ὑμεῖς λέγετε ὅτι Βλασφημεῖς, ὅτι εἶπον· Υἱὸς τοῦ Θεοῦ εἰμι;
\vs{37}Εἰ οὐ ποιῶ τὰ ἔργα τοῦ Πατρός μου, μὴ πιστεύετέ μοι·
\vs{38}εἰ δὲ ποιῶ, κἂν ἐμοὶ μὴ πιστεύητε, τοῖς ἔργοις πιστεύετε, ἵνα γνῶτε καὶ γινώσκητε ὅτι ἐν ἐμοὶ ὁ Πατὴρ κἀγὼ ἐν τῷ Πατρί.
\vs{39}Ἐζήτουν οὖν αὐτὸν πάλιν πιάσαι, καὶ ἐξῆλθεν ἐκ τῆς χειρὸς αὐτῶν.

\vs{40}Καὶ ἀπῆλθεν πάλιν πέραν τοῦ Ἰορδάνου εἰς τὸν τόπον ὅπου ἦν Ἰωάννης τὸ πρῶτον βαπτίζων καὶ ἔμεινεν ἐκεῖ.
\vs{41}καὶ πολλοὶ ἦλθον πρὸς αὐτὸν καὶ ἔλεγον ὅτι Ἰωάννης μὲν σημεῖον ἐποίησεν οὐδέν, πάντα δὲ ὅσα εἶπεν Ἰωάννης περὶ τούτου ἀληθῆ ἦν.
\vs{42}καὶ πολλοὶ ἐπίστευσαν εἰς αὐτὸν ἐκεῖ.

\ch{11}
Ἦν δέ τις ἀσθενῶν, Λάζαρος ἀπὸ Βηθανίας, ἐκ τῆς κώμης Μαρίας καὶ Μάρθας τῆς ἀδελφῆς αὐτῆς.
\vs{2}ἦν δὲ Μαριὰμ ἡ ἀλείψασα τὸν Κύριον μύρῳ καὶ ἐκμάξασα τοὺς πόδας αὐτοῦ ταῖς θριξὶν αὐτῆς, ἧς ὁ ἀδελφὸς Λάζαρος ἠσθένει.
\vs{3}ἀπέστειλαν οὖν αἱ ἀδελφαὶ πρὸς αὐτὸν λέγουσαι· Κύριε, ἴδε ὃν φιλεῖς ἀσθενεῖ.
\vs{4}Ἀκούσας δὲ ὁ Ἰησοῦς εἶπεν· Αὕτη ἡ ἀσθένεια οὐκ ἔστιν πρὸς θάνατον ἀλλ᾽ ὑπὲρ τῆς δόξης τοῦ Θεοῦ, ἵνα δοξασθῇ ὁ Υἱὸς τοῦ Θεοῦ δι᾽ αὐτῆς.
\vs{5}Ἠγάπα δὲ ὁ Ἰησοῦς τὴν Μάρθαν καὶ τὴν ἀδελφὴν αὐτῆς καὶ τὸν Λάζαρον.
\vs{6}ὡς οὖν ἤκουσεν ὅτι ἀσθενεῖ, τότε μὲν ἔμεινεν ἐν ᾧ ἦν τόπῳ δύο ἡμέρας,
\vs{7}ἔπειτα μετὰ τοῦτο λέγει τοῖς μαθηταῖς· Ἄγωμεν εἰς τὴν Ἰουδαίαν πάλιν.
\vs{8}Λέγουσιν αὐτῷ οἱ μαθηταί· Ῥαββί, νῦν ἐζήτουν σε λιθάσαι οἱ Ἰουδαῖοι, καὶ πάλιν ὑπάγεις ἐκεῖ;
\vs{9}Ἀπεκρίθη Ἰησοῦς· Οὐχὶ δώδεκα ὧραί εἰσιν τῆς ἡμέρας; ἐάν τις περιπατῇ ἐν τῇ ἡμέρᾳ, οὐ προσκόπτει, ὅτι τὸ φῶς τοῦ κόσμου τούτου βλέπει·
\vs{10}ἐὰν δέ τις περιπατῇ ἐν τῇ νυκτί, προσκόπτει, ὅτι τὸ φῶς οὐκ ἔστιν ἐν αὐτῷ.

\vs{11}Ταῦτα εἶπεν, καὶ μετὰ τοῦτο λέγει αὐτοῖς· Λάζαρος ὁ φίλος ἡμῶν κεκοίμηται· ἀλλὰ πορεύομαι ἵνα ἐξυπνίσω αὐτόν.
\vs{12}Εἶπαν οὖν οἱ μαθηταὶ αὐτῷ· Κύριε, εἰ κεκοίμηται σωθήσεται.
\vs{13}εἰρήκει δὲ ὁ Ἰησοῦς περὶ τοῦ θανάτου αὐτοῦ, ἐκεῖνοι δὲ ἔδοξαν ὅτι περὶ τῆς κοιμήσεως τοῦ ὕπνου λέγει.
\vs{14}Τότε οὖν εἶπεν αὐτοῖς ὁ Ἰησοῦς παρρησίᾳ· Λάζαρος ἀπέθανεν,
\vs{15}καὶ χαίρω δι᾽ ὑμᾶς ἵνα πιστεύσητε, ὅτι οὐκ ἤμην ἐκεῖ· ἀλλὰ ἄγωμεν πρὸς αὐτόν.
\vs{16}Εἶπεν οὖν Θωμᾶς ὁ λεγόμενος Δίδυμος τοῖς συμμαθηταῖς· Ἄγωμεν καὶ ἡμεῖς ἵνα ἀποθάνωμεν μετ᾽ αὐτοῦ.

\vs{17}Ἐλθὼν οὖν ὁ Ἰησοῦς εὗρεν αὐτὸν τέσσαρας ἤδη ἡμέρας ἔχοντα ἐν τῷ μνημείῳ.
\vs{18}ἦν δὲ ἡ Βηθανία ἐγγὺς τῶν Ἱεροσολύμων ὡς ἀπὸ σταδίων δεκαπέντε.
\vs{19}πολλοὶ δὲ ἐκ τῶν Ἰουδαίων ἐληλύθεισαν πρὸς τὴν Μάρθαν καὶ Μαριὰμ ἵνα παραμυθήσωνται αὐτὰς περὶ τοῦ ἀδελφοῦ.
\vs{20}ἡ οὖν Μάρθα ὡς ἤκουσεν ὅτι Ἰησοῦς ἔρχεται ὑπήντησεν αὐτῷ· Μαριὰμ δὲ ἐν τῷ οἴκῳ ἐκαθέζετο.
\vs{21}Εἶπεν οὖν ἡ Μάρθα πρὸς τὸν Ἰησοῦν· Κύριε, εἰ ἦς ὧδε οὐκ ἂν ἀπέθανεν ὁ ἀδελφός μου·
\vs{22}ἀλλὰ καὶ νῦν οἶδα ὅτι ὅσα ἂν αἰτήσῃ τὸν Θεὸν δώσει σοι ὁ Θεός.
\vs{23}Λέγει αὐτῇ ὁ Ἰησοῦς· Ἀναστήσεται ὁ ἀδελφός σου.
\vs{24}Λέγει αὐτῷ ἡ Μάρθα· Οἶδα ὅτι ἀναστήσεται ἐν τῇ ἀναστάσει ἐν τῇ ἐσχάτῃ ἡμέρᾳ.
\vs{25}Εἶπεν αὐτῇ ὁ Ἰησοῦς· Ἐγώ εἰμι ἡ ἀνάστασις καὶ ἡ ζωή· ὁ πιστεύων εἰς ἐμὲ κἂν ἀποθάνῃ ζήσεται,
\vs{26}καὶ πᾶς ὁ ζῶν καὶ πιστεύων εἰς ἐμὲ οὐ μὴ ἀποθάνῃ εἰς τὸν αἰῶνα. πιστεύεις τοῦτο;
\vs{27}Λέγει αὐτῷ· Ναί Κύριε, ἐγὼ πεπίστευκα ὅτι σὺ εἶ ὁ Χριστὸς ὁ Υἱὸς τοῦ Θεοῦ ὁ εἰς τὸν κόσμον ἐρχόμενος.

\vs{28}Καὶ τοῦτο εἰποῦσα ἀπῆλθεν καὶ ἐφώνησεν Μαριὰμ τὴν ἀδελφὴν αὐτῆς λάθρᾳ εἰποῦσα· Ὁ Διδάσκαλος πάρεστιν καὶ φωνεῖ σε.
\vs{29}ἐκείνη δὲ ὡς ἤκουσεν ἠγέρθη ταχὺ καὶ ἤρχετο πρὸς αὐτόν.
\vs{30}Οὔπω δὲ ἐληλύθει ὁ Ἰησοῦς εἰς τὴν κώμην, ἀλλ᾽ ἦν ἔτι ἐν τῷ τόπῳ ὅπου ὑπήντησεν αὐτῷ ἡ Μάρθα.
\vs{31}οἱ οὖν Ἰουδαῖοι οἱ ὄντες μετ᾽ αὐτῆς ἐν τῇ οἰκίᾳ καὶ παραμυθούμενοι αὐτήν, ἰδόντες τὴν Μαριὰμ ὅτι ταχέως ἀνέστη καὶ ἐξῆλθεν, ἠκολούθησαν αὐτῇ δόξαντες ὅτι ὑπάγει εἰς τὸ μνημεῖον ἵνα κλαύσῃ ἐκεῖ.

\vs{32}ἡ οὖν Μαριὰμ ὡς ἦλθεν ὅπου ἦν Ἰησοῦς ἰδοῦσα αὐτὸν ἔπεσεν αὐτοῦ πρὸς τοὺς πόδας λέγουσα αὐτῷ· Κύριε, εἰ ἦς ὧδε οὐκ ἄν μου ἀπέθανεν ὁ ἀδελφός.
\vs{33}Ἰησοῦς οὖν ὡς εἶδεν αὐτὴν κλαίουσαν καὶ τοὺς συνελθόντας αὐτῇ Ἰουδαίους κλαίοντας, ἐνεβριμήσατο τῷ πνεύματι καὶ ἐτάραξεν ἑαυτόν
\vs{34}καὶ εἶπεν· Ποῦ τεθείκατε αὐτόν; Λέγουσιν αὐτῷ· Κύριε, ἔρχου καὶ ἴδε.
\vs{35}Ἐδάκρυσεν ὁ Ἰησοῦς.
\vs{36}Ἔλεγον οὖν οἱ Ἰουδαῖοι· Ἴδε πῶς ἐφίλει αὐτόν.
\vs{37}Τινὲς δὲ ἐξ αὐτῶν εἶπαν· Οὐκ ἐδύνατο οὗτος ὁ ἀνοίξας τοὺς ὀφθαλμοὺς τοῦ τυφλοῦ ποιῆσαι ἵνα καὶ οὗτος μὴ ἀποθάνῃ;

\vs{38}Ἰησοῦς οὖν πάλιν ἐμβριμώμενος ἐν ἑαυτῷ ἔρχεται εἰς τὸ μνημεῖον· ἦν δὲ σπήλαιον καὶ λίθος ἐπέκειτο ἐπ᾽ αὐτῷ.
\vs{39}λέγει ὁ Ἰησοῦς· Ἄρατε τὸν λίθον. Λέγει αὐτῷ ἡ ἀδελφὴ τοῦ τετελευτηκότος Μάρθα· Κύριε, ἤδη ὄζει, τεταρταῖος γάρ ἐστιν.
\vs{40}Λέγει αὐτῇ ὁ Ἰησοῦς· Οὐκ εἶπόν σοι ὅτι ἐὰν πιστεύσῃς ὄψῃ τὴν δόξαν τοῦ Θεοῦ;
\vs{41}Ἦραν οὖν τὸν λίθον. ὁ δὲ Ἰησοῦς ἦρεν τοὺς ὀφθαλμοὺς ἄνω καὶ εἶπεν· Πάτερ, εὐχαριστῶ σοι ὅτι ἤκουσάς μου.
\vs{42}ἐγὼ δὲ ᾔδειν ὅτι πάντοτέ μου ἀκούεις, ἀλλὰ διὰ τὸν ὄχλον τὸν περιεστῶτα εἶπον, ἵνα πιστεύσωσιν ὅτι σύ με ἀπέστειλας.
\vs{43}Καὶ ταῦτα εἰπὼν φωνῇ μεγάλῃ ἐκραύγασεν· Λάζαρε, δεῦρο ἔξω.
\vs{44}ἐξῆλθεν ὁ τεθνηκὼς δεδεμένος τοὺς πόδας καὶ τὰς χεῖρας κειρίαις καὶ ἡ ὄψις αὐτοῦ σουδαρίῳ περιεδέδετο. Λέγει αὐτοῖς ὁ Ἰησοῦς· Λύσατε αὐτὸν καὶ ἄφετε αὐτὸν ὑπάγειν.

\vs{45}Πολλοὶ οὖν ἐκ τῶν Ἰουδαίων οἱ ἐλθόντες πρὸς τὴν Μαριὰμ καὶ θεασάμενοι ἃ ἐποίησεν ἐπίστευσαν εἰς αὐτόν·
\vs{46}τινὲς δὲ ἐξ αὐτῶν ἀπῆλθον πρὸς τοὺς Φαρισαίους καὶ εἶπαν αὐτοῖς ἃ ἐποίησεν Ἰησοῦς.

\vs{47}Συνήγαγον οὖν οἱ ἀρχιερεῖς καὶ οἱ Φαρισαῖοι συνέδριον καὶ ἔλεγον· Τί ποιοῦμεν ὅτι οὗτος ὁ ἄνθρωπος πολλὰ ποιεῖ σημεῖα;
\vs{48}ἐὰν ἀφῶμεν αὐτὸν οὕτως, πάντες πιστεύσουσιν εἰς αὐτόν, καὶ ἐλεύσονται οἱ Ῥωμαῖοι καὶ ἀροῦσιν ἡμῶν καὶ τὸν τόπον καὶ τὸ ἔθνος.
\vs{49}Εἷς δέ τις ἐξ αὐτῶν Καϊάφας, ἀρχιερεὺς ὢν τοῦ ἐνιαυτοῦ ἐκείνου, εἶπεν αὐτοῖς· Ὑμεῖς οὐκ οἴδατε οὐδέν,
\vs{50}οὐδὲ λογίζεσθε ὅτι συμφέρει ὑμῖν ἵνα εἷς ἄνθρωπος ἀποθάνῃ ὑπὲρ τοῦ λαοῦ καὶ μὴ ὅλον τὸ ἔθνος ἀπόληται.
\vs{51}Τοῦτο δὲ ἀφ᾽ ἑαυτοῦ οὐκ εἶπεν, ἀλλὰ ἀρχιερεὺς ὢν τοῦ ἐνιαυτοῦ ἐκείνου ἐπροφήτευσεν ὅτι ἔμελλεν Ἰησοῦς ἀποθνήσκειν ὑπὲρ τοῦ ἔθνους,
\vs{52}καὶ οὐχ ὑπὲρ τοῦ ἔθνους μόνον ἀλλ᾽ ἵνα καὶ τὰ τέκνα τοῦ Θεοῦ τὰ διεσκορπισμένα συναγάγῃ εἰς ἕν.
\vs{53}Ἀπ᾽ ἐκείνης οὖν τῆς ἡμέρας ἐβουλεύσαντο ἵνα ἀποκτείνωσιν αὐτόν.

\vs{54}Ὁ οὖν Ἰησοῦς οὐκέτι παρρησίᾳ περιεπάτει ἐν τοῖς Ἰουδαίοις, ἀλλὰ ἀπῆλθεν ἐκεῖθεν εἰς τὴν χώραν ἐγγὺς τῆς ἐρήμου, εἰς Ἐφραὶμ λεγομένην πόλιν, κἀκεῖ ἔμεινεν μετὰ τῶν μαθητῶν.

\vs{55}Ἦν δὲ ἐγγὺς τὸ πάσχα τῶν Ἰουδαίων, καὶ ἀνέβησαν πολλοὶ εἰς Ἱεροσόλυμα ἐκ τῆς χώρας πρὸ τοῦ πάσχα ἵνα ἁγνίσωσιν ἑαυτούς.
\vs{56}ἐζήτουν οὖν τὸν Ἰησοῦν καὶ ἔλεγον μετ᾽ ἀλλήλων ἐν τῷ ἱερῷ ἑστηκότες· Τί δοκεῖ ὑμῖν; ὅτι οὐ μὴ ἔλθῃ εἰς τὴν ἑορτήν;
\vs{57}δεδώκεισαν δὲ οἱ ἀρχιερεῖς καὶ οἱ Φαρισαῖοι ἐντολὰς ἵνα ἐάν τις γνῷ ποῦ ἐστιν μηνύσῃ, ὅπως πιάσωσιν αὐτόν.

\ch{12}
Ὁ οὖν Ἰησοῦς πρὸ ἓξ ἡμερῶν τοῦ πάσχα ἦλθεν εἰς Βηθανίαν, ὅπου ἦν Λάζαρος, ὃν ἤγειρεν ἐκ νεκρῶν Ἰησοῦς.
\vs{2}ἐποίησαν οὖν αὐτῷ δεῖπνον ἐκεῖ, καὶ ἡ Μάρθα διηκόνει, ὁ δὲ Λάζαρος εἷς ἦν ἐκ τῶν ἀνακειμένων σὺν αὐτῷ.

\vs{3}ἡ οὖν Μαριὰμ λαβοῦσα λίτραν μύρου νάρδου πιστικῆς πολυτίμου ἤλειψεν τοὺς πόδας τοῦ Ἰησοῦ καὶ ἐξέμαξεν ταῖς θριξὶν αὐτῆς τοὺς πόδας αὐτοῦ· ἡ δὲ οἰκία ἐπληρώθη ἐκ τῆς ὀσμῆς τοῦ μύρου.
\vs{4}Λέγει δὲ Ἰούδας ὁ Ἰσκαριώτης εἷς ἐκ τῶν μαθητῶν αὐτοῦ, ὁ μέλλων αὐτὸν παραδιδόναι·
\vs{5}Διὰ τί τοῦτο τὸ μύρον οὐκ ἐπράθη τριακοσίων δηναρίων καὶ ἐδόθη πτωχοῖς;
\vs{6}εἶπεν δὲ τοῦτο οὐχ ὅτι περὶ τῶν πτωχῶν ἔμελεν αὐτῷ, ἀλλ᾽ ὅτι κλέπτης ἦν καὶ τὸ γλωσσόκομον ἔχων τὰ βαλλόμενα ἐβάσταζεν.
\vs{7}Εἶπεν οὖν ὁ Ἰησοῦς· Ἄφες αὐτήν, ἵνα εἰς τὴν ἡμέραν τοῦ ἐνταφιασμοῦ μου τηρήσῃ αὐτό·
\vs{8}τοὺς πτωχοὺς γὰρ πάντοτε ἔχετε μεθ᾽ ἑαυτῶν, ἐμὲ δὲ οὐ πάντοτε ἔχετε.

\vs{9}Ἔγνω οὖν ὁ ὄχλος πολὺς ἐκ τῶν Ἰουδαίων ὅτι ἐκεῖ ἐστιν καὶ ἦλθον οὐ διὰ τὸν Ἰησοῦν μόνον, ἀλλ᾽ ἵνα καὶ τὸν Λάζαρον ἴδωσιν ὃν ἤγειρεν ἐκ νεκρῶν.
\vs{10}ἐβουλεύσαντο δὲ οἱ ἀρχιερεῖς ἵνα καὶ τὸν Λάζαρον ἀποκτείνωσιν,
\vs{11}ὅτι πολλοὶ δι᾽ αὐτὸν ὑπῆγον τῶν Ἰουδαίων καὶ ἐπίστευον εἰς τὸν Ἰησοῦν.

\vs{12}Τῇ ἐπαύριον ὁ ὄχλος πολὺς ὁ ἐλθὼν εἰς τὴν ἑορτήν, ἀκούσαντες ὅτι ἔρχεται ὁ Ἰησοῦς εἰς Ἱεροσόλυμα
\vs{13}ἔλαβον τὰ βαΐα τῶν φοινίκων καὶ ἐξῆλθον εἰς ὑπάντησιν αὐτῷ καὶ ἐκραύγαζον· Ὡσαννά· Εὐλογημένος ὁ ἐρχόμενος ἐν ὀνόματι Κυρίου, Καὶ ὁ Βασιλεὺς τοῦ Ἰσραήλ.
\vs{14}Εὑρὼν δὲ ὁ Ἰησοῦς ὀνάριον ἐκάθισεν ἐπ᾽ αὐτό, καθώς ἐστιν γεγραμμένον·
\begin{poetryblock}

\begin{quote} \vs{15}Μὴ φοβοῦ, θυγάτηρ Σιών·\end{quote} 

\begin{quote}ἰδοὺ ὁ Βασιλεύς σου ἔρχεται,\end{quote} 

\begin{quote}καθήμενος ἐπὶ πῶλον ὄνου.\end{quote}
\end{poetryblock}

\vs{16}Ταῦτα οὐκ ἔγνωσαν αὐτοῦ οἱ μαθηταὶ τὸ πρῶτον, ἀλλ᾽ ὅτε ἐδοξάσθη Ἰησοῦς τότε ἐμνήσθησαν ὅτι ταῦτα ἦν ἐπ᾽ αὐτῷ γεγραμμένα καὶ ταῦτα ἐποίησαν αὐτῷ.
\vs{17}Ἐμαρτύρει οὖν ὁ ὄχλος ὁ ὢν μετ᾽ αὐτοῦ ὅτε τὸν Λάζαρον ἐφώνησεν ἐκ τοῦ μνημείου καὶ ἤγειρεν αὐτὸν ἐκ νεκρῶν.
\vs{18}διὰ τοῦτο καὶ ὑπήντησεν αὐτῷ ὁ ὄχλος, ὅτι ἤκουσαν τοῦτο αὐτὸν πεποιηκέναι τὸ σημεῖον.
\vs{19}Οἱ οὖν Φαρισαῖοι εἶπαν πρὸς ἑαυτούς· Θεωρεῖτε ὅτι οὐκ ὠφελεῖτε οὐδέν· ἴδε ὁ κόσμος ὀπίσω αὐτοῦ ἀπῆλθεν.

\vs{20}Ἦσαν δὲ Ἕλληνές τινες ἐκ τῶν ἀναβαινόντων ἵνα προσκυνήσωσιν ἐν τῇ ἑορτῇ·
\vs{21}οὗτοι οὖν προσῆλθον Φιλίππῳ τῷ ἀπὸ Βηθσαϊδὰ τῆς Γαλιλαίας καὶ ἠρώτων αὐτὸν λέγοντες· Κύριε, θέλομεν τὸν Ἰησοῦν ἰδεῖν.
\vs{22}ἔρχεται ὁ Φίλιππος καὶ λέγει τῷ Ἀνδρέᾳ, ἔρχεται Ἀνδρέας καὶ Φίλιππος καὶ λέγουσιν τῷ Ἰησοῦ.
\vs{23}Ὁ δὲ Ἰησοῦς ἀποκρίνεται αὐτοῖς λέγων· Ἐλήλυθεν ἡ ὥρα ἵνα δοξασθῇ ὁ Υἱὸς τοῦ ἀνθρώπου.
\vs{24}ἀμὴν ἀμὴν λέγω ὑμῖν, ἐὰν μὴ ὁ κόκκος τοῦ σίτου πεσὼν εἰς τὴν γῆν ἀποθάνῃ, αὐτὸς μόνος μένει· ἐὰν δὲ ἀποθάνῃ, πολὺν καρπὸν φέρει.
\vs{25}ὁ φιλῶν τὴν ψυχὴν αὐτοῦ ἀπολλύει αὐτήν, καὶ ὁ μισῶν τὴν ψυχὴν αὐτοῦ ἐν τῷ κόσμῳ τούτῳ εἰς ζωὴν αἰώνιον φυλάξει αὐτήν.
\vs{26}ἐὰν ἐμοί τις διακονῇ, ἐμοὶ ἀκολουθείτω, καὶ ὅπου εἰμὶ ἐγὼ ἐκεῖ καὶ ὁ διάκονος ὁ ἐμὸς ἔσται· ἐάν τις ἐμοὶ διακονῇ τιμήσει αὐτὸν ὁ Πατήρ.
\vs{27}Νῦν ἡ ψυχή μου τετάρακται, καὶ τί εἴπω; Πάτερ, σῶσόν με ἐκ τῆς ὥρας ταύτης; ἀλλὰ διὰ τοῦτο ἦλθον εἰς τὴν ὥραν ταύτην.
\vs{28}Πάτερ, δόξασόν σου τὸ ὄνομα. Ἦλθεν οὖν φωνὴ ἐκ τοῦ οὐρανοῦ· Καὶ ἐδόξασα καὶ πάλιν δοξάσω.
\vs{29}Ὁ οὖν ὄχλος ὁ ἑστὼς καὶ ἀκούσας ἔλεγεν Βροντὴν γεγονέναι, ἄλλοι ἔλεγον· Ἄγγελος αὐτῷ λελάληκεν.
\vs{30}Ἀπεκρίθη Ἰησοῦς καὶ εἶπεν· Οὐ δι᾽ ἐμὲ ἡ φωνὴ αὕτη γέγονεν ἀλλὰ δι᾽ ὑμᾶς.
\vs{31}νῦν κρίσις ἐστὶν τοῦ κόσμου τούτου, νῦν ὁ ἄρχων τοῦ κόσμου τούτου ἐκβληθήσεται ἔξω·
\vs{32}κἀγὼ ἐὰν ὑψωθῶ ἐκ τῆς γῆς, πάντας ἑλκύσω πρὸς ἐμαυτόν.
\vs{33}τοῦτο δὲ ἔλεγεν σημαίνων ποίῳ θανάτῳ ἤμελλεν ἀποθνήσκειν.

\vs{34}Ἀπεκρίθη οὖν αὐτῷ ὁ ὄχλος· Ἡμεῖς ἠκούσαμεν ἐκ τοῦ νόμου ὅτι ὁ Χριστὸς μένει εἰς τὸν αἰῶνα, καὶ πῶς λέγεις σὺ ὅτι δεῖ ὑψωθῆναι τὸν Υἱὸν τοῦ ἀνθρώπου; τίς ἐστιν οὗτος ὁ Υἱὸς τοῦ ἀνθρώπου;
\vs{35}Εἶπεν οὖν αὐτοῖς ὁ Ἰησοῦς· Ἔτι μικρὸν χρόνον τὸ φῶς ἐν ὑμῖν ἐστιν. περιπατεῖτε ὡς τὸ φῶς ἔχετε, ἵνα μὴ σκοτία ὑμᾶς καταλάβῃ· καὶ ὁ περιπατῶν ἐν τῇ σκοτίᾳ οὐκ οἶδεν ποῦ ὑπάγει.
\vs{36}ὡς τὸ φῶς ἔχετε, πιστεύετε εἰς τὸ φῶς, ἵνα υἱοὶ φωτὸς γένησθε. Ταῦτα ἐλάλησεν Ἰησοῦς, καὶ ἀπελθὼν ἐκρύβη ἀπ᾽ αὐτῶν.

\vs{37}Τοσαῦτα δὲ αὐτοῦ σημεῖα πεποιηκότος ἔμπροσθεν αὐτῶν οὐκ ἐπίστευον εἰς αὐτόν,
\vs{38}ἵνα ὁ λόγος Ἠσαΐου τοῦ προφήτου πληρωθῇ ὃν εἶπεν· 
\begin{poetryblock}

\begin{quote}Κύριε, τίς ἐπίστευσεν τῇ ἀκοῇ ἡμῶν;\end{quote} 

\begin{quote}καὶ ὁ βραχίων Κυρίου τίνι ἀπεκαλύφθη;\end{quote}
\end{poetryblock}

\vs{39}Διὰ τοῦτο οὐκ ἠδύναντο πιστεύειν, ὅτι πάλιν εἶπεν Ἠσαΐας·
\begin{poetryblock}

\begin{quote} \vs{40}Τετύφλωκεν αὐτῶν τοὺς ὀφθαλμοὺς\end{quote} 

\begin{quote}καὶ ἐπώρωσεν αὐτῶν τὴν καρδίαν,\end{quote} 

\begin{quote}ἵνα μὴ ἴδωσιν τοῖς ὀφθαλμοῖς\end{quote} 

\begin{quote}καὶ νοήσωσιν τῇ καρδίᾳ\end{quote} 

\begin{quote}καὶ στραφῶσιν, καὶ ἰάσομαι αὐτούς.\end{quote}
\end{poetryblock}

\vs{41}Ταῦτα εἶπεν Ἠσαΐας ὅτι εἶδεν τὴν δόξαν αὐτοῦ, καὶ ἐλάλησεν περὶ αὐτοῦ.
\vs{42}ὅμως μέντοι καὶ ἐκ τῶν ἀρχόντων πολλοὶ ἐπίστευσαν εἰς αὐτόν, ἀλλὰ διὰ τοὺς Φαρισαίους οὐχ ὡμολόγουν ἵνα μὴ ἀποσυνάγωγοι γένωνται·
\vs{43}ἠγάπησαν γὰρ τὴν δόξαν τῶν ἀνθρώπων μᾶλλον ἤπερ τὴν δόξαν τοῦ Θεοῦ.

\vs{44}Ἰησοῦς δὲ ἔκραξεν καὶ εἶπεν· Ὁ πιστεύων εἰς ἐμὲ οὐ πιστεύει εἰς ἐμὲ ἀλλὰ εἰς τὸν πέμψαντά με,
\vs{45}καὶ ὁ θεωρῶν ἐμὲ θεωρεῖ τὸν πέμψαντά με.
\vs{46}ἐγὼ φῶς εἰς τὸν κόσμον ἐλήλυθα, ἵνα πᾶς ὁ πιστεύων εἰς ἐμὲ ἐν τῇ σκοτίᾳ μὴ μείνῃ.
\vs{47}Καὶ ἐάν τίς μου ἀκούσῃ τῶν ῥημάτων καὶ μὴ φυλάξῃ, ἐγὼ οὐ κρίνω αὐτόν· οὐ γὰρ ἦλθον ἵνα κρίνω τὸν κόσμον, ἀλλ᾽ ἵνα σώσω τὸν κόσμον.
\vs{48}ὁ ἀθετῶν ἐμὲ καὶ μὴ λαμβάνων τὰ ῥήματά μου ἔχει τὸν κρίνοντα αὐτόν· ὁ λόγος ὃν ἐλάλησα ἐκεῖνος κρινεῖ αὐτὸν ἐν τῇ ἐσχάτῃ ἡμέρᾳ.
\vs{49}Ὅτι ἐγὼ ἐξ ἐμαυτοῦ οὐκ ἐλάλησα, ἀλλ᾽ ὁ πέμψας με Πατὴρ αὐτός μοι ἐντολὴν δέδωκεν τί εἴπω καὶ τί λαλήσω.
\vs{50}καὶ οἶδα ὅτι ἡ ἐντολὴ αὐτοῦ ζωὴ αἰώνιός ἐστιν. ἃ οὖν ἐγὼ λαλῶ, καθὼς εἴρηκέν μοι ὁ Πατήρ, οὕτως λαλῶ.

\ch{13}
Πρὸ δὲ τῆς ἑορτῆς τοῦ πάσχα εἰδὼς ὁ Ἰησοῦς ὅτι ἦλθεν αὐτοῦ ἡ ὥρα ἵνα μεταβῇ ἐκ τοῦ κόσμου τούτου πρὸς τὸν Πατέρα, ἀγαπήσας τοὺς ἰδίους τοὺς ἐν τῷ κόσμῳ εἰς τέλος ἠγάπησεν αὐτούς.
\vs{2}καὶ δείπνου γινομένου, τοῦ διαβόλου ἤδη βεβληκότος εἰς τὴν καρδίαν ἵνα παραδοῖ αὐτὸν Ἰούδας Σίμωνος Ἰσκαριώτου,
\vs{3}εἰδὼς ὅτι πάντα ἔδωκεν αὐτῷ ὁ Πατὴρ εἰς τὰς χεῖρας καὶ ὅτι ἀπὸ Θεοῦ ἐξῆλθεν καὶ πρὸς τὸν Θεὸν ὑπάγει,
\vs{4}ἐγείρεται ἐκ τοῦ δείπνου καὶ τίθησιν τὰ ἱμάτια καὶ λαβὼν λέντιον διέζωσεν ἑαυτόν·
\vs{5}εἶτα βάλλει ὕδωρ εἰς τὸν νιπτῆρα καὶ ἤρξατο νίπτειν τοὺς πόδας τῶν μαθητῶν καὶ ἐκμάσσειν τῷ λεντίῳ ᾧ ἦν διεζωσμένος.
\vs{6}Ἔρχεται οὖν πρὸς Σίμωνα Πέτρον· λέγει αὐτῷ· Κύριε, σύ μου νίπτεις τοὺς πόδας;
\vs{7}Ἀπεκρίθη Ἰησοῦς καὶ εἶπεν αὐτῷ· Ὃ ἐγὼ ποιῶ σὺ οὐκ οἶδας ἄρτι, γνώσῃ δὲ μετὰ ταῦτα.
\vs{8}Λέγει αὐτῷ Πέτρος· Οὐ μὴ νίψῃς μου τοὺς πόδας εἰς τὸν αἰῶνα. Ἀπεκρίθη Ἰησοῦς αὐτῷ· Ἐὰν μὴ νίψω σε, οὐκ ἔχεις μέρος μετ᾽ ἐμοῦ.
\vs{9}Λέγει αὐτῷ Σίμων Πέτρος· Κύριε, μὴ τοὺς πόδας μου μόνον ἀλλὰ καὶ τὰς χεῖρας καὶ τὴν κεφαλήν.
\vs{10}Λέγει αὐτῷ ὁ Ἰησοῦς· Ὁ λελουμένος οὐκ ἔχει χρείαν εἰ μὴ τοὺς πόδας νίψασθαι, ἀλλ᾽ ἔστιν καθαρὸς ὅλος· καὶ ὑμεῖς καθαροί ἐστε, ἀλλ᾽ οὐχὶ πάντες.
\vs{11}ᾔδει γὰρ τὸν παραδιδόντα αὐτόν· διὰ τοῦτο εἶπεν ὅτι Οὐχὶ πάντες καθαροί ἐστε.

\vs{12}Ὅτε οὖν ἔνιψεν τοὺς πόδας αὐτῶν καὶ ἔλαβεν τὰ ἱμάτια αὐτοῦ καὶ ἀνέπεσεν πάλιν, εἶπεν αὐτοῖς· Γινώσκετε τί πεποίηκα ὑμῖν;
\vs{13}ὑμεῖς φωνεῖτέ με· Ὁ Διδάσκαλος, καὶ· ὁ Κύριος, καὶ καλῶς λέγετε· εἰμὶ γάρ.
\vs{14}εἰ οὖν ἐγὼ ἔνιψα ὑμῶν τοὺς πόδας ὁ Κύριος καὶ ὁ Διδάσκαλος, καὶ ὑμεῖς ὀφείλετε ἀλλήλων νίπτειν τοὺς πόδας·
\vs{15}ὑπόδειγμα γὰρ ἔδωκα ὑμῖν ἵνα καθὼς ἐγὼ ἐποίησα ὑμῖν καὶ ὑμεῖς ποιῆτε.
\vs{16}ἀμὴν ἀμὴν λέγω ὑμῖν, οὐκ ἔστιν δοῦλος μείζων τοῦ κυρίου αὐτοῦ οὐδὲ ἀπόστολος μείζων τοῦ πέμψαντος αὐτόν.
\vs{17}εἰ ταῦτα οἴδατε, μακάριοί ἐστε ἐὰν ποιῆτε αὐτά.

\vs{18}Οὐ περὶ πάντων ὑμῶν λέγω· ἐγὼ οἶδα τίνας ἐξελεξάμην· ἀλλ᾽ ἵνα ἡ γραφὴ πληρωθῇ· Ὁ τρώγων μου τὸν ἄρτον ἐπῆρεν ἐπ᾽ ἐμὲ τὴν πτέρναν αὐτοῦ.
\vs{19}ἀπ᾽ ἄρτι λέγω ὑμῖν πρὸ τοῦ γενέσθαι, ἵνα πιστεύσητε ὅταν γένηται ὅτι ἐγώ εἰμι.
\vs{20}ἀμὴν ἀμὴν λέγω ὑμῖν, ὁ λαμβάνων ἄν τινα πέμψω ἐμὲ λαμβάνει, ὁ δὲ ἐμὲ λαμβάνων λαμβάνει τὸν πέμψαντά με.

\vs{21}Ταῦτα εἰπὼν ὁ Ἰησοῦς ἐταράχθη τῷ πνεύματι καὶ ἐμαρτύρησεν καὶ εἶπεν· Ἀμὴν ἀμὴν λέγω ὑμῖν ὅτι εἷς ἐξ ὑμῶν παραδώσει με.
\vs{22}Ἔβλεπον εἰς ἀλλήλους οἱ μαθηταὶ ἀπορούμενοι περὶ τίνος λέγει.
\vs{23}ἦν ἀνακείμενος εἷς ἐκ τῶν μαθητῶν αὐτοῦ ἐν τῷ κόλπῳ τοῦ Ἰησοῦ, ὃν ἠγάπα ὁ Ἰησοῦς.
\vs{24}νεύει οὖν τούτῳ Σίμων Πέτρος πυθέσθαι τίς ἂν εἴη περὶ οὗ λέγει.
\vs{25}ἀναπεσὼν οὖν ἐκεῖνος οὕτως ἐπὶ τὸ στῆθος τοῦ Ἰησοῦ λέγει αὐτῷ· Κύριε, τίς ἐστιν;
\vs{26}Ἀποκρίνεται ὁ Ἰησοῦς· Ἐκεῖνός ἐστιν ᾧ ἐγὼ βάψω τὸ ψωμίον καὶ δώσω αὐτῷ. βάψας οὖν τὸ ψωμίον λαμβάνει καὶ δίδωσιν Ἰούδᾳ Σίμωνος Ἰσκαριώτου.
\vs{27}καὶ μετὰ τὸ ψωμίον τότε εἰσῆλθεν εἰς ἐκεῖνον ὁ Σατανᾶς. Λέγει οὖν αὐτῷ ὁ Ἰησοῦς· Ὃ ποιεῖς ποίησον τάχιον.
\vs{28}τοῦτο δὲ οὐδεὶς ἔγνω τῶν ἀνακειμένων πρὸς τί εἶπεν αὐτῷ·
\vs{29}τινὲς γὰρ ἐδόκουν, ἐπεὶ τὸ γλωσσόκομον εἶχεν Ἰούδας, ὅτι λέγει αὐτῷ ὁ Ἰησοῦς· Ἀγόρασον ὧν χρείαν ἔχομεν εἰς τὴν ἑορτήν, ἢ τοῖς πτωχοῖς ἵνα τι δῷ.
\vs{30}λαβὼν οὖν τὸ ψωμίον ἐκεῖνος ἐξῆλθεν εὐθύς. ἦν δὲ νύξ.

\vs{31}Ὅτε οὖν ἐξῆλθεν, λέγει Ἰησοῦς· Νῦν ἐδοξάσθη ὁ Υἱὸς τοῦ ἀνθρώπου καὶ ὁ Θεὸς ἐδοξάσθη ἐν αὐτῷ·
\vs{32}εἰ ὁ Θεὸς ἐδοξάσθη ἐν αὐτῷ, καὶ ὁ Θεὸς δοξάσει αὐτὸν ἐν αὑτῷ, καὶ εὐθὺς δοξάσει αὐτόν.
\vs{33}Τεκνία, ἔτι μικρὸν μεθ᾽ ὑμῶν εἰμι· ζητήσετέ με, καὶ καθὼς εἶπον τοῖς Ἰουδαίοις ὅτι Ὅπου ἐγὼ ὑπάγω ὑμεῖς οὐ δύνασθε ἐλθεῖν, καὶ ὑμῖν λέγω ἄρτι.
\vs{34}Ἐντολὴν καινὴν δίδωμι ὑμῖν, ἵνα ἀγαπᾶτε ἀλλήλους, καθὼς ἠγάπησα ὑμᾶς ἵνα καὶ ὑμεῖς ἀγαπᾶτε ἀλλήλους.
\vs{35}ἐν τούτῳ γνώσονται πάντες ὅτι ἐμοὶ μαθηταί ἐστε, ἐὰν ἀγάπην ἔχητε ἐν ἀλλήλοις.

\vs{36}Λέγει αὐτῷ Σίμων Πέτρος· Κύριε, ποῦ ὑπάγεις; Ἀπεκρίθη αὐτῷ Ἰησοῦς· Ὅπου ὑπάγω οὐ δύνασαί μοι νῦν ἀκολουθῆσαι, ἀκολουθήσεις δὲ ὕστερον.
\vs{37}Λέγει αὐτῷ ὁ Πέτρος· Κύριε, διὰ τί οὐ δύναμαί σοι ἀκολουθῆσαι ἄρτι; τὴν ψυχήν μου ὑπὲρ σοῦ θήσω.
\vs{38}Ἀποκρίνεται Ἰησοῦς· Τὴν ψυχήν σου ὑπὲρ ἐμοῦ θήσεις; ἀμὴν ἀμὴν λέγω σοι, οὐ μὴ ἀλέκτωρ φωνήσῃ ἕως οὗ ἀρνήσῃ με τρίς.

\ch{14}
Μὴ ταρασσέσθω ὑμῶν ἡ καρδία· πιστεύετε εἰς τὸν Θεόν καὶ εἰς ἐμὲ πιστεύετε.
\vs{2}ἐν τῇ οἰκίᾳ τοῦ Πατρός μου μοναὶ πολλαί εἰσιν· εἰ δὲ μή, εἶπον ἂν ὑμῖν ὅτι πορεύομαι ἑτοιμάσαι τόπον ὑμῖν;
\vs{3}καὶ ἐὰν πορευθῶ καὶ ἑτοιμάσω τόπον ὑμῖν, πάλιν ἔρχομαι καὶ παραλήμψομαι ὑμᾶς πρὸς ἐμαυτόν, ἵνα ὅπου εἰμὶ ἐγὼ καὶ ὑμεῖς ἦτε.
\vs{4}καὶ ὅπου ἐγὼ ὑπάγω οἴδατε τὴν ὁδόν.

\vs{5}Λέγει αὐτῷ Θωμᾶς· Κύριε, οὐκ οἴδαμεν ποῦ ὑπάγεις· πῶς δυνάμεθα τὴν ὁδὸν εἰδέναι;
\vs{6}Λέγει αὐτῷ ὁ Ἰησοῦς· Ἐγώ εἰμι ἡ ὁδὸς καὶ ἡ ἀλήθεια καὶ ἡ ζωή· οὐδεὶς ἔρχεται πρὸς τὸν Πατέρα εἰ μὴ δι᾽ ἐμοῦ.
\vs{7}εἰ ἐγνώκατέ με, καὶ τὸν Πατέρα μου γνώσεσθε. καὶ ἀπ᾽ ἄρτι γινώσκετε αὐτὸν καὶ ἑωράκατε αὐτόν.

\vs{8}Λέγει αὐτῷ Φίλιππος· Κύριε, δεῖξον ἡμῖν τὸν Πατέρα, καὶ ἀρκεῖ ἡμῖν.
\vs{9}Λέγει αὐτῷ ὁ Ἰησοῦς· Τοσούτῳ χρόνῳ μεθ᾽ ὑμῶν εἰμι καὶ οὐκ ἔγνωκάς με, Φίλιππε; ὁ ἑωρακὼς ἐμὲ ἑώρακεν τὸν Πατέρα· πῶς σὺ λέγεις· Δεῖξον ἡμῖν τὸν Πατέρα;
\vs{10}οὐ πιστεύεις ὅτι ἐγὼ ἐν τῷ Πατρὶ καὶ ὁ Πατὴρ ἐν ἐμοί ἐστιν; τὰ ῥήματα ἃ ἐγὼ λέγω ὑμῖν ἀπ᾽ ἐμαυτοῦ οὐ λαλῶ, ὁ δὲ Πατὴρ ἐν ἐμοὶ μένων ποιεῖ τὰ ἔργα αὐτοῦ.
\vs{11}πιστεύετέ μοι ὅτι ἐγὼ ἐν τῷ Πατρὶ καὶ ὁ Πατὴρ ἐν ἐμοί· εἰ δὲ μή, διὰ τὰ ἔργα αὐτὰ πιστεύετε.
\vs{12}Ἀμὴν ἀμὴν λέγω ὑμῖν, ὁ πιστεύων εἰς ἐμὲ τὰ ἔργα ἃ ἐγὼ ποιῶ κἀκεῖνος ποιήσει καὶ μείζονα τούτων ποιήσει, ὅτι ἐγὼ πρὸς τὸν Πατέρα πορεύομαι·
\vs{13}καὶ ὅ τι ἂν αἰτήσητε ἐν τῷ ὀνόματί μου τοῦτο ποιήσω, ἵνα δοξασθῇ ὁ Πατὴρ ἐν τῷ Υἱῷ.
\vs{14}ἐάν τι αἰτήσητέ με ἐν τῷ ὀνόματί μου ἐγὼ ποιήσω.

\vs{15}Ἐὰν ἀγαπᾶτέ με, τὰς ἐντολὰς τὰς ἐμὰς τηρήσετε·
\vs{16}Κἀγὼ ἐρωτήσω τὸν Πατέρα καὶ ἄλλον Παράκλητον δώσει ὑμῖν, ἵνα μεθ᾽ ὑμῶν εἰς τὸν αἰῶνα ᾖ,
\vs{17}τὸ Πνεῦμα τῆς ἀληθείας, ὃ ὁ κόσμος οὐ δύναται λαβεῖν, ὅτι οὐ θεωρεῖ αὐτὸ οὐδὲ γινώσκει· ὑμεῖς γινώσκετε αὐτό, ὅτι παρ᾽ ὑμῖν μένει καὶ ἐν ὑμῖν ἔσται.
\vs{18}Οὐκ ἀφήσω ὑμᾶς ὀρφανούς, ἔρχομαι πρὸς ὑμᾶς.
\vs{19}ἔτι μικρὸν καὶ ὁ κόσμος με οὐκέτι θεωρεῖ, ὑμεῖς δὲ θεωρεῖτέ με, ὅτι ἐγὼ ζῶ καὶ ὑμεῖς ζήσετε.
\vs{20}ἐν ἐκείνῃ τῇ ἡμέρᾳ γνώσεσθε ὑμεῖς ὅτι ἐγὼ ἐν τῷ Πατρί μου καὶ ὑμεῖς ἐν ἐμοὶ κἀγὼ ἐν ὑμῖν.
\vs{21}ὁ ἔχων τὰς ἐντολάς μου καὶ τηρῶν αὐτὰς ἐκεῖνός ἐστιν ὁ ἀγαπῶν με· ὁ δὲ ἀγαπῶν με ἀγαπηθήσεται ὑπὸ τοῦ Πατρός μου, κἀγὼ ἀγαπήσω αὐτὸν καὶ ἐμφανίσω αὐτῷ ἐμαυτόν.

\vs{22}Λέγει αὐτῷ Ἰούδας, οὐχ ὁ Ἰσκαριώτης· Κύριε, καὶ τί γέγονεν ὅτι ἡμῖν μέλλεις ἐμφανίζειν σεαυτὸν καὶ οὐχὶ τῷ κόσμῳ;
\vs{23}Ἀπεκρίθη Ἰησοῦς καὶ εἶπεν αὐτῷ· Ἐάν τις ἀγαπᾷ με τὸν λόγον μου τηρήσει, καὶ ὁ Πατήρ μου ἀγαπήσει αὐτόν καὶ πρὸς αὐτὸν ἐλευσόμεθα καὶ μονὴν παρ᾽ αὐτῷ ποιησόμεθα.
\vs{24}ὁ μὴ ἀγαπῶν με τοὺς λόγους μου οὐ τηρεῖ· καὶ ὁ λόγος ὃν ἀκούετε οὐκ ἔστιν ἐμὸς ἀλλὰ τοῦ πέμψαντός με Πατρός.

\vs{25}Ταῦτα λελάληκα ὑμῖν παρ᾽ ὑμῖν μένων·
\vs{26}ὁ δὲ Παράκλητος, τὸ Πνεῦμα τὸ Ἅγιον, ὃ πέμψει ὁ Πατὴρ ἐν τῷ ὀνόματί μου, ἐκεῖνος ὑμᾶς διδάξει πάντα καὶ ὑπομνήσει ὑμᾶς πάντα ἃ εἶπον ὑμῖν ἐγώ.

\vs{27}Εἰρήνην ἀφίημι ὑμῖν, εἰρήνην τὴν ἐμὴν δίδωμι ὑμῖν· οὐ καθὼς ὁ κόσμος δίδωσιν ἐγὼ δίδωμι ὑμῖν. μὴ ταρασσέσθω ὑμῶν ἡ καρδία μηδὲ δειλιάτω.
\vs{28}ἠκούσατε ὅτι ἐγὼ εἶπον ὑμῖν· Ὑπάγω καὶ ἔρχομαι πρὸς ὑμᾶς. εἰ ἠγαπᾶτέ με ἐχάρητε ἄν ὅτι πορεύομαι πρὸς τὸν Πατέρα, ὅτι ὁ Πατὴρ μείζων μού ἐστιν.
\vs{29}καὶ νῦν εἴρηκα ὑμῖν πρὶν γενέσθαι, ἵνα ὅταν γένηται πιστεύσητε.
\vs{30}Οὐκέτι πολλὰ λαλήσω μεθ᾽ ὑμῶν, ἔρχεται γὰρ ὁ τοῦ κόσμου ἄρχων· καὶ ἐν ἐμοὶ οὐκ ἔχει οὐδέν,
\vs{31}ἀλλ᾽ ἵνα γνῷ ὁ κόσμος ὅτι ἀγαπῶ τὸν Πατέρα, καὶ καθὼς ἐνετείλατο μοι ὁ Πατὴρ, οὕτως ποιῶ. Ἐγείρεσθε, ἄγωμεν ἐντεῦθεν.

\ch{15}
Ἐγώ εἰμι ἡ ἄμπελος ἡ ἀληθινή καὶ ὁ Πατήρ μου ὁ γεωργός ἐστιν.
\vs{2}πᾶν κλῆμα ἐν ἐμοὶ μὴ φέρον καρπὸν αἴρει αὐτό, καὶ πᾶν τὸ καρπὸν φέρον καθαίρει αὐτὸ ἵνα καρπὸν πλείονα φέρῃ.
\vs{3}ἤδη ὑμεῖς καθαροί ἐστε διὰ τὸν λόγον ὃν λελάληκα ὑμῖν·
\vs{4}μείνατε ἐν ἐμοί, κἀγὼ ἐν ὑμῖν. καθὼς τὸ κλῆμα οὐ δύναται καρπὸν φέρειν ἀφ᾽ ἑαυτοῦ ἐὰν μὴ μένῃ ἐν τῇ ἀμπέλῳ, οὕτως οὐδὲ ὑμεῖς ἐὰν μὴ ἐν ἐμοὶ μένητε.
\vs{5}Ἐγώ εἰμι ἡ ἄμπελος, ὑμεῖς τὰ κλήματα. ὁ μένων ἐν ἐμοὶ κἀγὼ ἐν αὐτῷ οὗτος φέρει καρπὸν πολύν, ὅτι χωρὶς ἐμοῦ οὐ δύνασθε ποιεῖν οὐδέν.
\vs{6}ἐὰν μή τις μένῃ ἐν ἐμοί, ἐβλήθη ἔξω ὡς τὸ κλῆμα καὶ ἐξηράνθη καὶ συνάγουσιν αὐτὰ καὶ εἰς τὸ πῦρ βάλλουσιν καὶ καίεται.
\vs{7}ἐὰν μείνητε ἐν ἐμοὶ καὶ τὰ ῥήματά μου ἐν ὑμῖν μείνῃ, ὃ ἐὰν θέλητε αἰτήσασθε, καὶ γενήσεται ὑμῖν.
\vs{8}ἐν τούτῳ ἐδοξάσθη ὁ Πατήρ μου, ἵνα καρπὸν πολὺν φέρητε καὶ γένησθε ἐμοὶ μαθηταί.

\vs{9}Καθὼς ἠγάπησέν με ὁ Πατήρ, κἀγὼ ὑμᾶς ἠγάπησα· μείνατε ἐν τῇ ἀγάπῃ τῇ ἐμῇ.
\vs{10}ἐὰν τὰς ἐντολάς μου τηρήσητε, μενεῖτε ἐν τῇ ἀγάπῃ μου, καθὼς ἐγὼ τὰς ἐντολὰς τοῦ Πατρός μου τετήρηκα καὶ μένω αὐτοῦ ἐν τῇ ἀγάπῃ.
\vs{11}Ταῦτα λελάληκα ὑμῖν ἵνα ἡ χαρὰ ἡ ἐμὴ ἐν ὑμῖν ᾖ καὶ ἡ χαρὰ ὑμῶν πληρωθῇ.
\vs{12}Αὕτη ἐστὶν ἡ ἐντολὴ ἡ ἐμὴ, ἵνα ἀγαπᾶτε ἀλλήλους καθὼς ἠγάπησα ὑμᾶς.
\vs{13}μείζονα ταύτης ἀγάπην οὐδεὶς ἔχει, ἵνα τις τὴν ψυχὴν αὐτοῦ θῇ ὑπὲρ τῶν φίλων αὐτοῦ.
\vs{14}Ὑμεῖς φίλοι μού ἐστε ἐὰν ποιῆτε ἃ ἐγὼ ἐντέλλομαι ὑμῖν.
\vs{15}οὐκέτι λέγω ὑμᾶς δούλους, ὅτι ὁ δοῦλος οὐκ οἶδεν τί ποιεῖ αὐτοῦ ὁ κύριος· ὑμᾶς δὲ εἴρηκα φίλους, ὅτι πάντα ἃ ἤκουσα παρὰ τοῦ Πατρός μου ἐγνώρισα ὑμῖν.
\vs{16}οὐχ ὑμεῖς με ἐξελέξασθε, ἀλλ᾽ ἐγὼ ἐξελεξάμην ὑμᾶς καὶ ἔθηκα ὑμᾶς ἵνα ὑμεῖς ὑπάγητε καὶ καρπὸν φέρητε καὶ ὁ καρπὸς ὑμῶν μένῃ, ἵνα ὅ τι ἂν αἰτήσητε τὸν Πατέρα ἐν τῷ ὀνόματί μου δῷ ὑμῖν.
\vs{17}ταῦτα ἐντέλλομαι ὑμῖν, ἵνα ἀγαπᾶτε ἀλλήλους.

\vs{18}Εἰ ὁ κόσμος ὑμᾶς μισεῖ, γινώσκετε ὅτι ἐμὲ πρῶτον ὑμῶν μεμίσηκεν.
\vs{19}εἰ ἐκ τοῦ κόσμου ἦτε, ὁ κόσμος ἂν τὸ ἴδιον ἐφίλει· ὅτι δὲ ἐκ τοῦ κόσμου οὐκ ἐστέ, ἀλλ᾽ ἐγὼ ἐξελεξάμην ὑμᾶς ἐκ τοῦ κόσμου, διὰ τοῦτο μισεῖ ὑμᾶς ὁ κόσμος.
\vs{20}Μνημονεύετε τοῦ λόγου οὗ ἐγὼ εἶπον ὑμῖν· Οὐκ ἔστιν δοῦλος μείζων τοῦ κυρίου αὐτοῦ. εἰ ἐμὲ ἐδίωξαν, καὶ ὑμᾶς διώξουσιν· εἰ τὸν λόγον μου ἐτήρησαν, καὶ τὸν ὑμέτερον τηρήσουσιν.
\vs{21}ἀλλὰ ταῦτα πάντα ποιήσουσιν εἰς ὑμᾶς διὰ τὸ ὄνομά μου, ὅτι οὐκ οἴδασιν τὸν πέμψαντά με.
\vs{22}εἰ μὴ ἦλθον καὶ ἐλάλησα αὐτοῖς, ἁμαρτίαν οὐκ εἴχοσαν· νῦν δὲ πρόφασιν οὐκ ἔχουσιν περὶ τῆς ἁμαρτίας αὐτῶν.
\vs{23}Ὁ ἐμὲ μισῶν καὶ τὸν Πατέρα μου μισεῖ.
\vs{24}εἰ τὰ ἔργα μὴ ἐποίησα ἐν αὐτοῖς ἃ οὐδεὶς ἄλλος ἐποίησεν, ἁμαρτίαν οὐκ εἴχοσαν· νῦν δὲ καὶ ἑωράκασιν καὶ μεμισήκασιν καὶ ἐμὲ καὶ τὸν Πατέρα μου.
\vs{25}ἀλλ᾽ ἵνα πληρωθῇ ὁ λόγος ὁ ἐν τῷ νόμῳ αὐτῶν γεγραμμένος ὅτι Ἐμίσησάν με δωρεάν.

\vs{26}Ὅταν ἔλθῃ ὁ Παράκλητος ὃν ἐγὼ πέμψω ὑμῖν παρὰ τοῦ Πατρός, τὸ Πνεῦμα τῆς ἀληθείας ὃ παρὰ τοῦ Πατρὸς ἐκπορεύεται, ἐκεῖνος μαρτυρήσει περὶ ἐμοῦ·
\vs{27}καὶ ὑμεῖς δὲ μαρτυρεῖτε, ὅτι ἀπ᾽ ἀρχῆς μετ᾽ ἐμοῦ ἐστε.

\ch{16}
Ταῦτα λελάληκα ὑμῖν ἵνα μὴ σκανδαλισθῆτε.
\vs{2}ἀποσυναγώγους ποιήσουσιν ὑμᾶς· ἀλλ᾽ ἔρχεται ὥρα ἵνα πᾶς ὁ ἀποκτείνας ὑμᾶς δόξῃ λατρείαν προσφέρειν τῷ Θεῷ.
\vs{3}καὶ ταῦτα ποιήσουσιν ὅτι οὐκ ἔγνωσαν τὸν Πατέρα οὐδὲ ἐμέ.
\vs{4}ἀλλὰ ταῦτα λελάληκα ὑμῖν ἵνα ὅταν ἔλθῃ ἡ ὥρα αὐτῶν μνημονεύητε αὐτῶν ὅτι ἐγὼ εἶπον ὑμῖν. ταῦτα δὲ ὑμῖν ἐξ ἀρχῆς οὐκ εἶπον, ὅτι μεθ᾽ ὑμῶν ἤμην.
\vs{5}Νῦν δὲ ὑπάγω πρὸς τὸν πέμψαντά με, καὶ οὐδεὶς ἐξ ὑμῶν ἐρωτᾷ με· Ποῦ ὑπάγεις;
\vs{6}ἀλλ᾽ ὅτι ταῦτα λελάληκα ὑμῖν ἡ λύπη πεπλήρωκεν ὑμῶν τὴν καρδίαν.
\vs{7}ἀλλ᾽ ἐγὼ τὴν ἀλήθειαν λέγω ὑμῖν, συμφέρει ὑμῖν ἵνα ἐγὼ ἀπέλθω. ἐὰν γὰρ μὴ ἀπέλθω, ὁ Παράκλητος οὐκ ἐλεύσεται πρὸς ὑμᾶς· ἐὰν δὲ πορευθῶ, πέμψω αὐτὸν πρὸς ὑμᾶς.
\vs{8}Καὶ ἐλθὼν ἐκεῖνος ἐλέγξει τὸν κόσμον περὶ ἁμαρτίας καὶ περὶ δικαιοσύνης καὶ περὶ κρίσεως·
\vs{9}περὶ ἁμαρτίας μέν, ὅτι οὐ πιστεύουσιν εἰς ἐμέ·
\vs{10}περὶ δικαιοσύνης δέ, ὅτι πρὸς τὸν Πατέρα ὑπάγω καὶ οὐκέτι θεωρεῖτέ με·
\vs{11}περὶ δὲ κρίσεως, ὅτι ὁ ἄρχων τοῦ κόσμου τούτου κέκριται.

\vs{12}Ἔτι πολλὰ ἔχω ὑμῖν λέγειν, ἀλλ᾽ οὐ δύνασθε βαστάζειν ἄρτι·
\vs{13}ὅταν δὲ ἔλθῃ ἐκεῖνος, τὸ Πνεῦμα τῆς ἀληθείας, ὁδηγήσει ὑμᾶς ἐν τῇ ἀληθείᾳ πάσῃ· οὐ γὰρ λαλήσει ἀφ᾽ ἑαυτοῦ, ἀλλ᾽ ὅσα ἀκούσει λαλήσει καὶ τὰ ἐρχόμενα ἀναγγελεῖ ὑμῖν.
\vs{14}ἐκεῖνος ἐμὲ δοξάσει, ὅτι ἐκ τοῦ ἐμοῦ λήμψεται καὶ ἀναγγελεῖ ὑμῖν.
\vs{15}πάντα ὅσα ἔχει ὁ Πατὴρ ἐμά ἐστιν· διὰ τοῦτο εἶπον ὅτι ἐκ τοῦ ἐμοῦ λαμβάνει καὶ ἀναγγελεῖ ὑμῖν.

\vs{16}Μικρὸν καὶ οὐκέτι θεωρεῖτέ με, καὶ πάλιν μικρὸν καὶ ὄψεσθέ με.
\vs{17}Εἶπαν οὖν ἐκ τῶν μαθητῶν αὐτοῦ πρὸς ἀλλήλους· Τί ἐστιν τοῦτο ὃ λέγει ἡμῖν· Μικρὸν καὶ οὐ θεωρεῖτέ με, καὶ πάλιν μικρὸν καὶ ὄψεσθέ με; καί· Ὅτι ὑπάγω πρὸς τὸν Πατέρα;
\vs{18}ἔλεγον οὖν· τί ἐστιν Τοῦτο ὃ λέγει Τὸ μικρόν; οὐκ οἴδαμεν τί λαλεῖ.

\vs{19}Ἔγνω ὁ Ἰησοῦς ὅτι ἤθελον αὐτὸν ἐρωτᾶν, καὶ εἶπεν αὐτοῖς· Περὶ τούτου ζητεῖτε μετ᾽ ἀλλήλων ὅτι εἶπον· Μικρὸν καὶ οὐ θεωρεῖτέ με, καὶ πάλιν μικρὸν καὶ ὄψεσθέ με;
\vs{20}ἀμὴν ἀμὴν λέγω ὑμῖν ὅτι κλαύσετε καὶ θρηνήσετε ὑμεῖς, ὁ δὲ κόσμος χαρήσεται· ὑμεῖς λυπηθήσεσθε, ἀλλ᾽ ἡ λύπη ὑμῶν εἰς χαρὰν γενήσεται.
\vs{21}ἡ γυνὴ ὅταν τίκτῃ λύπην ἔχει, ὅτι ἦλθεν ἡ ὥρα αὐτῆς· ὅταν δὲ γεννήσῃ τὸ παιδίον, οὐκέτι μνημονεύει τῆς θλίψεως διὰ τὴν χαρὰν ὅτι ἐγεννήθη ἄνθρωπος εἰς τὸν κόσμον.
\vs{22}καὶ ὑμεῖς οὖν νῦν μὲν λύπην ἔχετε· πάλιν δὲ ὄψομαι ὑμᾶς, καὶ χαρήσεται ὑμῶν ἡ καρδία, καὶ τὴν χαρὰν ὑμῶν οὐδεὶς αἴρει ἀφ᾽ ὑμῶν.

\vs{23}Καὶ ἐν ἐκείνῃ τῇ ἡμέρᾳ ἐμὲ οὐκ ἐρωτήσετε οὐδέν. ἀμὴν ἀμὴν λέγω ὑμῖν, ἄν τι αἰτήσητε τὸν Πατέρα ἐν τῷ ὀνόματί μου δώσει ὑμῖν.
\vs{24}ἕως ἄρτι οὐκ ᾐτήσατε οὐδὲν ἐν τῷ ὀνόματί μου· αἰτεῖτε καὶ λήμψεσθε, ἵνα ἡ χαρὰ ὑμῶν ᾖ πεπληρωμένη.

\vs{25}Ταῦτα ἐν παροιμίαις λελάληκα ὑμῖν· ἔρχεται ὥρα ὅτε οὐκέτι ἐν παροιμίαις λαλήσω ὑμῖν, ἀλλὰ παρρησίᾳ περὶ τοῦ Πατρὸς ἀπαγγελῶ ὑμῖν.
\vs{26}ἐν ἐκείνῃ τῇ ἡμέρᾳ ἐν τῷ ὀνόματί μου αἰτήσεσθε, καὶ οὐ λέγω ὑμῖν ὅτι ἐγὼ ἐρωτήσω τὸν Πατέρα περὶ ὑμῶν·
\vs{27}αὐτὸς γὰρ ὁ Πατὴρ φιλεῖ ὑμᾶς, ὅτι ὑμεῖς ἐμὲ πεφιλήκατε καὶ πεπιστεύκατε ὅτι ἐγὼ παρὰ τοῦ Θεοῦ ἐξῆλθον.
\vs{28}ἐξῆλθον παρὰ τοῦ Πατρὸς καὶ ἐλήλυθα εἰς τὸν κόσμον· πάλιν ἀφίημι τὸν κόσμον καὶ πορεύομαι πρὸς τὸν Πατέρα.

\vs{29}Λέγουσιν οἱ μαθηταὶ αὐτοῦ· Ἴδε νῦν ἐν παρρησίᾳ λαλεῖς καὶ παροιμίαν οὐδεμίαν λέγεις.
\vs{30}νῦν οἴδαμεν ὅτι οἶδας πάντα καὶ οὐ χρείαν ἔχεις ἵνα τίς σε ἐρωτᾷ· ἐν τούτῳ πιστεύομεν ὅτι ἀπὸ Θεοῦ ἐξῆλθες.
\vs{31}Ἀπεκρίθη αὐτοῖς Ἰησοῦς· Ἄρτι πιστεύετε;
\vs{32}ἰδοὺ ἔρχεται ὥρα καὶ ἐλήλυθεν ἵνα σκορπισθῆτε ἕκαστος εἰς τὰ ἴδια κἀμὲ μόνον ἀφῆτε· καὶ οὐκ εἰμὶ μόνος, ὅτι ὁ Πατὴρ μετ᾽ ἐμοῦ ἐστιν.
\vs{33}ταῦτα λελάληκα ὑμῖν ἵνα ἐν ἐμοὶ εἰρήνην ἔχητε. ἐν τῷ κόσμῳ θλῖψιν ἔχετε· ἀλλὰ θαρσεῖτε, ἐγὼ νενίκηκα τὸν κόσμον.

\ch{17}
Ταῦτα ἐλάλησεν Ἰησοῦς καὶ ἐπάρας τοὺς ὀφθαλμοὺς αὐτοῦ εἰς τὸν οὐρανὸν εἶπεν· Πάτερ, ἐλήλυθεν ἡ ὥρα· δόξασόν σου τὸν Υἱόν, ἵνα ὁ Υἱὸς δοξάσῃ σέ,
\vs{2}καθὼς ἔδωκας αὐτῷ ἐξουσίαν πάσης σαρκός, ἵνα πᾶν ὃ δέδωκας αὐτῷ δώσῃ αὐτοῖς ζωὴν αἰώνιον.
\vs{3}αὕτη δέ ἐστιν ἡ αἰώνιος ζωὴ ἵνα γινώσκωσιν σὲ τὸν μόνον ἀληθινὸν Θεὸν καὶ ὃν ἀπέστειλας Ἰησοῦν Χριστόν.
\vs{4}ἐγώ σε ἐδόξασα ἐπὶ τῆς γῆς τὸ ἔργον τελειώσας ὃ δέδωκάς μοι ἵνα ποιήσω·
\vs{5}καὶ νῦν δόξασόν με σύ, Πάτερ, παρὰ σεαυτῷ τῇ δόξῃ ᾗ εἶχον πρὸ τοῦ τὸν κόσμον εἶναι παρὰ σοί.

\vs{6}Ἐφανέρωσά σου τὸ ὄνομα τοῖς ἀνθρώποις οὓς ἔδωκάς μοι ἐκ τοῦ κόσμου. σοὶ ἦσαν κἀμοὶ αὐτοὺς ἔδωκας καὶ τὸν λόγον σου τετήρηκαν.
\vs{7}νῦν ἔγνωκαν ὅτι πάντα ὅσα δέδωκάς μοι παρὰ σοῦ εἰσιν·
\vs{8}ὅτι τὰ ῥήματα ἃ ἔδωκάς μοι δέδωκα αὐτοῖς, καὶ αὐτοὶ ἔλαβον καὶ ἔγνωσαν ἀληθῶς ὅτι παρὰ σοῦ ἐξῆλθον, καὶ ἐπίστευσαν ὅτι σύ με ἀπέστειλας.

\vs{9}Ἐγὼ περὶ αὐτῶν ἐρωτῶ, οὐ περὶ τοῦ κόσμου ἐρωτῶ ἀλλὰ περὶ ὧν δέδωκάς μοι, ὅτι σοί εἰσιν,
\vs{10}καὶ τὰ ἐμὰ πάντα σά ἐστιν καὶ τὰ σὰ ἐμά, καὶ δεδόξασμαι ἐν αὐτοῖς.
\vs{11}καὶ οὐκέτι εἰμὶ ἐν τῷ κόσμῳ, καὶ αὐτοὶ ἐν τῷ κόσμῳ εἰσίν, κἀγὼ πρὸς σὲ ἔρχομαι. Πάτερ ἅγιε, τήρησον αὐτοὺς ἐν τῷ ὀνόματί σου ᾧ δέδωκάς μοι, ἵνα ὦσιν ἓν καθὼς ἡμεῖς.
\vs{12}ὅτε ἤμην μετ᾽ αὐτῶν ἐγὼ ἐτήρουν αὐτοὺς ἐν τῷ ὀνόματί σου ᾧ δέδωκάς μοι, καὶ ἐφύλαξα, καὶ οὐδεὶς ἐξ αὐτῶν ἀπώλετο εἰ μὴ ὁ υἱὸς τῆς ἀπωλείας, ἵνα ἡ γραφὴ πληρωθῇ.
\vs{13}Νῦν δὲ πρὸς σὲ ἔρχομαι καὶ ταῦτα λαλῶ ἐν τῷ κόσμῳ ἵνα ἔχωσιν τὴν χαρὰν τὴν ἐμὴν πεπληρωμένην ἐν ἑαυτοῖς.
\vs{14}ἐγὼ δέδωκα αὐτοῖς τὸν λόγον σου καὶ ὁ κόσμος ἐμίσησεν αὐτούς, ὅτι οὐκ εἰσὶν ἐκ τοῦ κόσμου καθὼς ἐγὼ οὐκ εἰμὶ ἐκ τοῦ κόσμου.
\vs{15}Οὐκ ἐρωτῶ ἵνα ἄρῃς αὐτοὺς ἐκ τοῦ κόσμου, ἀλλ᾽ ἵνα τηρήσῃς αὐτοὺς ἐκ τοῦ πονηροῦ.
\vs{16}ἐκ τοῦ κόσμου οὐκ εἰσὶν καθὼς ἐγὼ οὐκ εἰμὶ ἐκ τοῦ κόσμου.
\vs{17}ἁγίασον αὐτοὺς ἐν τῇ ἀληθείᾳ· ὁ λόγος ὁ σὸς ἀλήθειά ἐστιν.
\vs{18}καθὼς ἐμὲ ἀπέστειλας εἰς τὸν κόσμον, κἀγὼ ἀπέστειλα αὐτοὺς εἰς τὸν κόσμον·
\vs{19}καὶ ὑπὲρ αὐτῶν ἐγὼ ἁγιάζω ἐμαυτόν, ἵνα ὦσιν καὶ αὐτοὶ ἡγιασμένοι ἐν ἀληθείᾳ.

\vs{20}Οὐ περὶ τούτων δὲ ἐρωτῶ μόνον, ἀλλὰ καὶ περὶ τῶν πιστευόντων διὰ τοῦ λόγου αὐτῶν εἰς ἐμέ,
\vs{21}ἵνα πάντες ἓν ὦσιν, καθὼς σύ, πάτερ, ἐν ἐμοὶ κἀγὼ ἐν σοί, ἵνα καὶ αὐτοὶ ἐν ἡμῖν ὦσιν, ἵνα ὁ κόσμος πιστεύῃ ὅτι σύ με ἀπέστειλας.
\vs{22}Κἀγὼ τὴν δόξαν ἣν δέδωκάς μοι δέδωκα αὐτοῖς, ἵνα ὦσιν ἓν καθὼς ἡμεῖς ἕν·
\vs{23}ἐγὼ ἐν αὐτοῖς καὶ σὺ ἐν ἐμοί, ἵνα ὦσιν τετελειωμένοι εἰς ἕν, ἵνα γινώσκῃ ὁ κόσμος ὅτι σύ με ἀπέστειλας καὶ ἠγάπησας αὐτοὺς καθὼς ἐμὲ ἠγάπησας.

\vs{24}Πάτερ, ὃ δέδωκάς μοι, θέλω ἵνα ὅπου εἰμὶ ἐγὼ κἀκεῖνοι ὦσιν μετ᾽ ἐμοῦ, ἵνα θεωρῶσιν τὴν δόξαν τὴν ἐμὴν, ἣν δέδωκάς μοι ὅτι ἠγάπησάς με πρὸ καταβολῆς κόσμου.
\vs{25}Πάτερ δίκαιε, καὶ ὁ κόσμος σε οὐκ ἔγνω, ἐγὼ δέ σε ἔγνων, καὶ οὗτοι ἔγνωσαν ὅτι σύ με ἀπέστειλας·
\vs{26}καὶ ἐγνώρισα αὐτοῖς τὸ ὄνομά σου καὶ γνωρίσω, ἵνα ἡ ἀγάπη ἣν ἠγάπησάς με ἐν αὐτοῖς ᾖ κἀγὼ ἐν αὐτοῖς.

\ch{18}
Ταῦτα εἰπὼν Ἰησοῦς ἐξῆλθεν σὺν τοῖς μαθηταῖς αὐτοῦ πέραν τοῦ χειμάρρου τοῦ Κέδρων ὅπου ἦν κῆπος, εἰς ὃν εἰσῆλθεν αὐτὸς καὶ οἱ μαθηταὶ αὐτοῦ.

\vs{2}ᾔδει δὲ καὶ Ἰούδας ὁ παραδιδοὺς αὐτὸν τὸν τόπον, ὅτι πολλάκις συνήχθη Ἰησοῦς ἐκεῖ μετὰ τῶν μαθητῶν αὐτοῦ.
\vs{3}ὁ οὖν Ἰούδας λαβὼν τὴν σπεῖραν καὶ ἐκ τῶν ἀρχιερέων καὶ ἐκ τῶν Φαρισαίων ὑπηρέτας ἔρχεται ἐκεῖ μετὰ φανῶν καὶ λαμπάδων καὶ ὅπλων.
\vs{4}Ἰησοῦς οὖν εἰδὼς πάντα τὰ ἐρχόμενα ἐπ᾽ αὐτὸν ἐξῆλθεν καὶ λέγει αὐτοῖς· Τίνα ζητεῖτε;
\vs{5}Ἀπεκρίθησαν αὐτῷ· Ἰησοῦν τὸν Ναζωραῖον. Λέγει αὐτοῖς· Ἐγώ εἰμι. Εἱστήκει δὲ καὶ Ἰούδας ὁ παραδιδοὺς αὐτὸν μετ᾽ αὐτῶν.
\vs{6}ὡς οὖν εἶπεν αὐτοῖς· Ἐγώ εἰμι, ἀπῆλθον εἰς τὰ ὀπίσω καὶ ἔπεσαν χαμαί.
\vs{7}Πάλιν οὖν ἐπηρώτησεν αὐτούς· Τίνα ζητεῖτε; Οἱ δὲ εἶπαν· Ἰησοῦν τὸν Ναζωραῖον.
\vs{8}Ἀπεκρίθη Ἰησοῦς· Εἶπον ὑμῖν ὅτι ἐγώ εἰμι. εἰ οὖν ἐμὲ ζητεῖτε, ἄφετε τούτους ὑπάγειν·
\vs{9}ἵνα πληρωθῇ ὁ λόγος ὃν εἶπεν ὅτι Οὓς δέδωκάς μοι οὐκ ἀπώλεσα ἐξ αὐτῶν οὐδένα.
\vs{10}Σίμων οὖν Πέτρος ἔχων μάχαιραν εἵλκυσεν αὐτὴν καὶ ἔπαισεν τὸν τοῦ ἀρχιερέως δοῦλον καὶ ἀπέκοψεν αὐτοῦ τὸ ὠτάριον τὸ δεξιόν· ἦν δὲ ὄνομα τῷ δούλῳ Μάλχος.
\vs{11}εἶπεν οὖν ὁ Ἰησοῦς τῷ Πέτρῳ· Βάλε τὴν μάχαιραν εἰς τὴν θήκην· τὸ ποτήριον ὃ δέδωκέν μοι ὁ Πατήρ οὐ μὴ πίω αὐτό;

\vs{12}Ἡ οὖν σπεῖρα καὶ ὁ χιλίαρχος καὶ οἱ ὑπηρέται τῶν Ἰουδαίων συνέλαβον τὸν Ἰησοῦν καὶ ἔδησαν αὐτὸν
\vs{13}καὶ ἤγαγον πρὸς Ἅνναν πρῶτον· ἦν γὰρ πενθερὸς τοῦ Καϊάφα, ὃς ἦν ἀρχιερεὺς τοῦ ἐνιαυτοῦ ἐκείνου·
\vs{14}ἦν δὲ Καϊάφας ὁ συμβουλεύσας τοῖς Ἰουδαίοις ὅτι συμφέρει ἕνα ἄνθρωπον ἀποθανεῖν ὑπὲρ τοῦ λαοῦ.

\vs{15}Ἠκολούθει δὲ τῷ Ἰησοῦ Σίμων Πέτρος καὶ ἄλλος μαθητής. ὁ δὲ μαθητὴς ἐκεῖνος ἦν γνωστὸς τῷ ἀρχιερεῖ καὶ συνεισῆλθεν τῷ Ἰησοῦ εἰς τὴν αὐλὴν τοῦ ἀρχιερέως,
\vs{16}ὁ δὲ Πέτρος εἱστήκει πρὸς τῇ θύρᾳ ἔξω. ἐξῆλθεν οὖν ὁ μαθητὴς ὁ ἄλλος ὁ γνωστὸς τοῦ ἀρχιερέως καὶ εἶπεν τῇ θυρωρῷ καὶ εἰσήγαγεν τὸν Πέτρον.
\vs{17}Λέγει οὖν τῷ Πέτρῳ ἡ παιδίσκη ἡ θυρωρός· Μὴ καὶ σὺ ἐκ τῶν μαθητῶν εἶ τοῦ ἀνθρώπου τούτου; Λέγει ἐκεῖνος· Οὐκ εἰμί.
\vs{18}Εἱστήκεισαν δὲ οἱ δοῦλοι καὶ οἱ ὑπηρέται ἀνθρακιὰν πεποιηκότες, ὅτι ψῦχος ἦν, καὶ ἐθερμαίνοντο· ἦν δὲ καὶ ὁ Πέτρος μετ᾽ αὐτῶν ἑστὼς καὶ θερμαινόμενος.

\vs{19}Ὁ οὖν ἀρχιερεὺς ἠρώτησεν τὸν Ἰησοῦν περὶ τῶν μαθητῶν αὐτοῦ καὶ περὶ τῆς διδαχῆς αὐτοῦ.
\vs{20}Ἀπεκρίθη αὐτῷ Ἰησοῦς· Ἐγὼ παρρησίᾳ λελάληκα τῷ κόσμῳ, ἐγὼ πάντοτε ἐδίδαξα ἐν συναγωγῇ καὶ ἐν τῷ ἱερῷ, ὅπου πάντες οἱ Ἰουδαῖοι συνέρχονται, καὶ ἐν κρυπτῷ ἐλάλησα οὐδέν.
\vs{21}τί με ἐρωτᾷς; ἐρώτησον τοὺς ἀκηκοότας τί ἐλάλησα αὐτοῖς· ἴδε οὗτοι οἴδασιν ἃ εἶπον ἐγώ.
\vs{22}Ταῦτα δὲ αὐτοῦ εἰπόντος εἷς παρεστηκὼς τῶν ὑπηρετῶν ἔδωκεν ῥάπισμα τῷ Ἰησοῦ εἰπών· Οὕτως ἀποκρίνῃ τῷ ἀρχιερεῖ;
\vs{23}Ἀπεκρίθη αὐτῷ Ἰησοῦς· Εἰ κακῶς ἐλάλησα, μαρτύρησον περὶ τοῦ κακοῦ· εἰ δὲ καλῶς, τί με δέρεις;
\vs{24}Ἀπέστειλεν οὖν αὐτὸν ὁ Ἅννας δεδεμένον πρὸς Καϊάφαν τὸν ἀρχιερέα.

\vs{25}Ἦν δὲ Σίμων Πέτρος ἑστὼς καὶ θερμαινόμενος. εἶπον οὖν αὐτῷ· Μὴ καὶ σὺ ἐκ τῶν μαθητῶν αὐτοῦ εἶ; Ἠρνήσατο ἐκεῖνος καὶ εἶπεν· Οὐκ εἰμί.
\vs{26}Λέγει εἷς ἐκ τῶν δούλων τοῦ ἀρχιερέως, συγγενὴς ὢν οὗ ἀπέκοψεν Πέτρος τὸ ὠτίον· Οὐκ ἐγώ σε εἶδον ἐν τῷ κήπῳ μετ᾽ αὐτοῦ;
\vs{27}Πάλιν οὖν ἠρνήσατο Πέτρος, καὶ εὐθέως ἀλέκτωρ ἐφώνησεν.

\vs{28}Ἄγουσιν οὖν τὸν Ἰησοῦν ἀπὸ τοῦ Καϊάφα εἰς τὸ πραιτώριον· ἦν δὲ πρωΐ· καὶ αὐτοὶ οὐκ εἰσῆλθον εἰς τὸ πραιτώριον, ἵνα μὴ μιανθῶσιν ἀλλὰ φάγωσιν τὸ πάσχα.
\vs{29}Ἐξῆλθεν οὖν ὁ Πιλᾶτος ἔξω πρὸς αὐτοὺς καὶ φησίν· Τίνα κατηγορίαν φέρετε κατὰ τοῦ ἀνθρώπου τούτου;
\vs{30}Ἀπεκρίθησαν καὶ εἶπαν αὐτῷ· Εἰ μὴ ἦν οὗτος κακὸν ποιῶν, οὐκ ἄν σοι παρεδώκαμεν αὐτόν.
\vs{31}Εἶπεν οὖν αὐτοῖς ὁ Πιλᾶτος· Λάβετε αὐτὸν ὑμεῖς καὶ κατὰ τὸν νόμον ὑμῶν κρίνατε αὐτόν. Εἶπον αὐτῷ οἱ Ἰουδαῖοι· Ἡμῖν οὐκ ἔξεστιν ἀποκτεῖναι οὐδένα·
\vs{32}ἵνα ὁ λόγος τοῦ Ἰησοῦ πληρωθῇ ὃν εἶπεν σημαίνων ποίῳ θανάτῳ ἤμελλεν ἀποθνήσκειν.

\vs{33}Εἰσῆλθεν οὖν πάλιν εἰς τὸ πραιτώριον ὁ Πιλᾶτος καὶ ἐφώνησεν τὸν Ἰησοῦν καὶ εἶπεν αὐτῷ· Σὺ εἶ ὁ Βασιλεὺς τῶν Ἰουδαίων;
\vs{34}Ἀπεκρίθη Ἰησοῦς· Ἀπὸ σεαυτοῦ σὺ τοῦτο λέγεις ἢ ἄλλοι εἶπόν σοι περὶ ἐμοῦ;
\vs{35}Ἀπεκρίθη ὁ Πιλᾶτος· Μήτι ἐγὼ Ἰουδαῖός εἰμι; τὸ ἔθνος τὸ σὸν καὶ οἱ ἀρχιερεῖς παρέδωκάν σε ἐμοί· τί ἐποίησας;
\vs{36}Ἀπεκρίθη Ἰησοῦς· Ἡ βασιλεία ἡ ἐμὴ οὐκ ἔστιν ἐκ τοῦ κόσμου τούτου· εἰ ἐκ τοῦ κόσμου τούτου ἦν ἡ βασιλεία ἡ ἐμή, οἱ ὑπηρέται οἱ ἐμοὶ ἠγωνίζοντο ἄν ἵνα μὴ παραδοθῶ τοῖς Ἰουδαίοις· νῦν δὲ ἡ βασιλεία ἡ ἐμὴ οὐκ ἔστιν ἐντεῦθεν.
\vs{37}Εἶπεν οὖν αὐτῷ ὁ Πιλᾶτος· Οὐκοῦν βασιλεὺς εἶ σύ; Ἀπεκρίθη ὁ Ἰησοῦς· Σὺ λέγεις ὅτι βασιλεύς εἰμι. ἐγὼ εἰς τοῦτο γεγέννημαι καὶ εἰς τοῦτο ἐλήλυθα εἰς τὸν κόσμον, ἵνα μαρτυρήσω τῇ ἀληθείᾳ· πᾶς ὁ ὢν ἐκ τῆς ἀληθείας ἀκούει μου τῆς φωνῆς.
\vs{38}Λέγει αὐτῷ ὁ Πιλᾶτος· Τί ἐστιν ἀλήθεια;

Καὶ τοῦτο εἰπὼν πάλιν ἐξῆλθεν πρὸς τοὺς Ἰουδαίους καὶ λέγει αὐτοῖς· Ἐγὼ οὐδεμίαν εὑρίσκω ἐν αὐτῷ αἰτίαν.
\vs{39}ἔστιν δὲ συνήθεια ὑμῖν ἵνα ἕνα ἀπολύσω ὑμῖν ἐν τῷ πάσχα· βούλεσθε οὖν ἀπολύσω ὑμῖν τὸν Βασιλέα τῶν Ἰουδαίων;
\vs{40}Ἐκραύγασαν οὖν πάλιν λέγοντες· Μὴ τοῦτον ἀλλὰ τὸν Βαραββᾶν. ἦν δὲ ὁ Βαραββᾶς λῃστής.

\ch{19}
Τότε οὖν ἔλαβεν ὁ Πιλᾶτος τὸν Ἰησοῦν καὶ ἐμαστίγωσεν.
\vs{2}καὶ οἱ στρατιῶται πλέξαντες στέφανον ἐξ ἀκανθῶν ἐπέθηκαν αὐτοῦ τῇ κεφαλῇ καὶ ἱμάτιον πορφυροῦν περιέβαλον αὐτόν
\vs{3}καὶ ἤρχοντο πρὸς αὐτὸν καὶ ἔλεγον· Χαῖρε ὁ Βασιλεὺς τῶν Ἰουδαίων· καὶ ἐδίδοσαν αὐτῷ ῥαπίσματα.
\vs{4}Καὶ ἐξῆλθεν πάλιν ἔξω ὁ Πιλᾶτος καὶ λέγει αὐτοῖς· Ἴδε ἄγω ὑμῖν αὐτὸν ἔξω, ἵνα γνῶτε ὅτι οὐδεμίαν αἰτίαν εὑρίσκω ἐν αὐτῷ.
\vs{5}ἐξῆλθεν οὖν ὁ Ἰησοῦς ἔξω, φορῶν τὸν ἀκάνθινον στέφανον καὶ τὸ πορφυροῦν ἱμάτιον. καὶ λέγει αὐτοῖς· Ἰδοὺ ὁ ἄνθρωπος.

\vs{6}Ὅτε οὖν εἶδον αὐτὸν οἱ ἀρχιερεῖς καὶ οἱ ὑπηρέται ἐκραύγασαν λέγοντες· Σταύρωσον σταύρωσον. Λέγει αὐτοῖς ὁ Πιλᾶτος· Λάβετε αὐτὸν ὑμεῖς καὶ σταυρώσατε· ἐγὼ γὰρ οὐχ εὑρίσκω ἐν αὐτῷ αἰτίαν.
\vs{7}Ἀπεκρίθησαν αὐτῷ οἱ Ἰουδαῖοι· Ἡμεῖς νόμον ἔχομεν καὶ κατὰ τὸν νόμον ὀφείλει ἀποθανεῖν, ὅτι Υἱὸν Θεοῦ ἑαυτὸν ἐποίησεν.
\vs{8}Ὅτε οὖν ἤκουσεν ὁ Πιλᾶτος τοῦτον τὸν λόγον, μᾶλλον ἐφοβήθη,
\vs{9}καὶ εἰσῆλθεν εἰς τὸ πραιτώριον πάλιν καὶ λέγει τῷ Ἰησοῦ· Πόθεν εἶ σύ; Ὁ δὲ Ἰησοῦς ἀπόκρισιν οὐκ ἔδωκεν αὐτῷ.
\vs{10}Λέγει οὖν αὐτῷ ὁ Πιλᾶτος· Ἐμοὶ οὐ λαλεῖς; οὐκ οἶδας ὅτι ἐξουσίαν ἔχω ἀπολῦσαί σε καὶ ἐξουσίαν ἔχω σταυρῶσαί σε;
\vs{11}Ἀπεκρίθη αὐτῷ Ἰησοῦς· Οὐκ εἶχες ἐξουσίαν κατ᾽ ἐμοῦ οὐδεμίαν εἰ μὴ ἦν δεδομένον σοι ἄνωθεν· διὰ τοῦτο ὁ παραδούς μέ σοι μείζονα ἁμαρτίαν ἔχει.
\vs{12}Ἐκ τούτου ὁ Πιλᾶτος ἐζήτει ἀπολῦσαι αὐτόν· οἱ δὲ Ἰουδαῖοι ἐκραύγασαν λέγοντες· Ἐὰν τοῦτον ἀπολύσῃς, οὐκ εἶ φίλος τοῦ Καίσαρος· πᾶς ὁ βασιλέα ἑαυτὸν ποιῶν ἀντιλέγει τῷ Καίσαρι.
\vs{13}Ὁ οὖν Πιλᾶτος ἀκούσας τῶν λόγων τούτων ἤγαγεν ἔξω τὸν Ἰησοῦν καὶ ἐκάθισεν ἐπὶ βήματος εἰς τόπον λεγόμενον Λιθόστρωτον, Ἑβραϊστὶ δὲ Γαββαθα.
\vs{14}ἦν δὲ Παρασκευὴ τοῦ πάσχα, ὥρα ἦν ὡς ἕκτη. καὶ λέγει τοῖς Ἰουδαίοις· Ἴδε ὁ Βασιλεὺς ὑμῶν.
\vs{15}Ἐκραύγασαν οὖν Ἐκεῖνοι· Ἆρον ἆρον, σταύρωσον αὐτόν. Λέγει αὐτοῖς ὁ Πιλᾶτος· Τὸν Βασιλέα ὑμῶν σταυρώσω; Ἀπεκρίθησαν οἱ ἀρχιερεῖς· Οὐκ ἔχομεν βασιλέα εἰ μὴ Καίσαρα.
\vs{16}Τότε οὖν παρέδωκεν αὐτὸν αὐτοῖς ἵνα σταυρωθῇ.

Παρέλαβον οὖν τὸν Ἰησοῦν,
\vs{17}Καὶ βαστάζων ἑαυτῷ τὸν σταυρὸν ἐξῆλθεν εἰς τὸν λεγόμενον Κρανίου τόπον, ὃ λέγεται Ἑβραϊστὶ Γολγοθᾶ,
\vs{18}ὅπου αὐτὸν ἐσταύρωσαν, καὶ μετ᾽ αὐτοῦ ἄλλους δύο ἐντεῦθεν καὶ ἐντεῦθεν, μέσον δὲ τὸν Ἰησοῦν.
\vs{19}Ἔγραψεν δὲ καὶ τίτλον ὁ Πιλᾶτος καὶ ἔθηκεν ἐπὶ τοῦ σταυροῦ· ἦν δὲ γεγραμμένον· ΙΗΣΟΥΣ Ο ΝΑΖΩΡΑΙΟΣ Ο ΒΑΣΙΛΕΥΣ ΤΩΝ ΙΟΥΔΑΙΩΝ.
\vs{20}Τοῦτον οὖν τὸν τίτλον πολλοὶ ἀνέγνωσαν τῶν Ἰουδαίων, ὅτι ἐγγὺς ἦν ὁ τόπος τῆς πόλεως ὅπου ἐσταυρώθη ὁ Ἰησοῦς· καὶ ἦν γεγραμμένον Ἑβραϊστί, Ῥωμαϊστί, Ἑλληνιστί.
\vs{21}ἔλεγον οὖν τῷ Πιλάτῳ οἱ ἀρχιερεῖς τῶν Ἰουδαίων· Μὴ γράφε· Ὁ Βασιλεὺς τῶν Ἰουδαίων, ἀλλ᾽ ὅτι ἐκεῖνος εἶπεν· Βασιλεύς εἰμι τῶν Ἰουδαίων.
\vs{22}Ἀπεκρίθη ὁ Πιλᾶτος· Ὃ γέγραφα, γέγραφα.

\vs{23}Οἱ οὖν στρατιῶται, ὅτε ἐσταύρωσαν τὸν Ἰησοῦν, ἔλαβον τὰ ἱμάτια αὐτοῦ καὶ ἐποίησαν τέσσαρα μέρη, ἑκάστῳ στρατιώτῃ μέρος, καὶ τὸν χιτῶνα. ἦν δὲ ὁ χιτὼν ἄραφος, ἐκ τῶν ἄνωθεν ὑφαντὸς δι᾽ ὅλου.
\vs{24}εἶπαν οὖν πρὸς ἀλλήλους· Μὴ σχίσωμεν αὐτόν, ἀλλὰ λάχωμεν περὶ αὐτοῦ τίνος ἔσται· ἵνα ἡ γραφὴ πληρωθῇ ἡ λέγουσα· 
\begin{poetryblock}

\begin{quote}Διεμερίσαντο τὰ ἱμάτιά μου ἑαυτοῖς\end{quote} 

\begin{quote}καὶ ἐπὶ τὸν ἱματισμόν μου ἔβαλον κλῆρον.\end{quote}
\end{poetryblock}

Οἱ μὲν οὖν στρατιῶται ταῦτα ἐποίησαν.

\vs{25}Εἱστήκεισαν δὲ παρὰ τῷ σταυρῷ τοῦ Ἰησοῦ ἡ μήτηρ αὐτοῦ καὶ ἡ ἀδελφὴ τῆς μητρὸς αὐτοῦ, Μαρία ἡ τοῦ Κλωπᾶ καὶ Μαρία ἡ Μαγδαληνή.
\vs{26}Ἰησοῦς οὖν ἰδὼν τὴν μητέρα καὶ τὸν μαθητὴν παρεστῶτα ὃν ἠγάπα, λέγει τῇ μητρί· Γύναι, ἴδε ὁ υἱός σου.
\vs{27}εἶτα λέγει τῷ μαθητῇ· Ἴδε ἡ μήτηρ σου. καὶ ἀπ᾽ ἐκείνης τῆς ὥρας ἔλαβεν ὁ μαθητὴς αὐτὴν εἰς τὰ ἴδια.

\vs{28}Μετὰ τοῦτο εἰδὼς ὁ Ἰησοῦς ὅτι ἤδη πάντα τετέλεσται, ἵνα τελειωθῇ ἡ γραφὴ, λέγει· Διψῶ.
\vs{29}σκεῦος ἔκειτο ὄξους μεστόν· σπόγγον οὖν μεστὸν τοῦ ὄξους ὑσσώπῳ περιθέντες προσήνεγκαν αὐτοῦ τῷ στόματι.
\vs{30}ὅτε οὖν ἔλαβεν τὸ ὄξος ὁ Ἰησοῦς εἶπεν· Τετέλεσται, καὶ κλίνας τὴν κεφαλὴν παρέδωκεν τὸ πνεῦμα.

\vs{31}Οἱ οὖν Ἰουδαῖοι, ἐπεὶ Παρασκευὴ ἦν, ἵνα μὴ μείνῃ ἐπὶ τοῦ σταυροῦ τὰ σώματα ἐν τῷ σαββάτῳ, ἦν γὰρ μεγάλη ἡ ἡμέρα ἐκείνου τοῦ σαββάτου, ἠρώτησαν τὸν Πιλᾶτον ἵνα κατεαγῶσιν αὐτῶν τὰ σκέλη καὶ ἀρθῶσιν.
\vs{32}ἦλθον οὖν οἱ στρατιῶται καὶ τοῦ μὲν πρώτου κατέαξαν τὰ σκέλη καὶ τοῦ ἄλλου τοῦ συσταυρωθέντος αὐτῷ·
\vs{33}ἐπὶ δὲ τὸν Ἰησοῦν ἐλθόντες, ὡς εἶδον ἤδη αὐτὸν τεθνηκότα, οὐ κατέαξαν αὐτοῦ τὰ σκέλη,
\vs{34}ἀλλ᾽ εἷς τῶν στρατιωτῶν λόγχῃ αὐτοῦ τὴν πλευρὰν ἔνυξεν, καὶ ἐξῆλθεν εὐθὺς αἷμα καὶ ὕδωρ.
\vs{35}καὶ ὁ ἑωρακὼς μεμαρτύρηκεν, καὶ ἀληθινὴ αὐτοῦ ἐστιν ἡ μαρτυρία, καὶ ἐκεῖνος οἶδεν ὅτι ἀληθῆ λέγει, ἵνα καὶ ὑμεῖς πιστεύητε.
\vs{36}Ἐγένετο γὰρ ταῦτα ἵνα ἡ γραφὴ πληρωθῇ· Ὀστοῦν οὐ συντριβήσεται αὐτοῦ.
\vs{37}καὶ πάλιν ἑτέρα γραφὴ λέγει· Ὄψονται εἰς ὃν ἐξεκέντησαν.

\vs{38}Μετὰ δὲ ταῦτα ἠρώτησεν τὸν Πιλᾶτον Ἰωσὴφ ὁ ἀπὸ Ἁριμαθαίας, ὢν μαθητὴς τοῦ Ἰησοῦ κεκρυμμένος δὲ διὰ τὸν φόβον τῶν Ἰουδαίων, ἵνα ἄρῃ τὸ σῶμα τοῦ Ἰησοῦ· καὶ ἐπέτρεψεν ὁ Πιλᾶτος. ἦλθεν οὖν καὶ ἦρεν τὸ σῶμα αὐτοῦ.
\vs{39}ἦλθεν δὲ καὶ Νικόδημος, ὁ ἐλθὼν πρὸς αὐτὸν νυκτὸς τὸ πρῶτον, φέρων μίγμα σμύρνης καὶ ἀλόης ὡς λίτρας ἑκατόν.
\vs{40}ἔλαβον οὖν τὸ σῶμα τοῦ Ἰησοῦ καὶ ἔδησαν αὐτὸ ὀθονίοις μετὰ τῶν ἀρωμάτων, καθὼς ἔθος ἐστὶν τοῖς Ἰουδαίοις ἐνταφιάζειν.
\vs{41}Ἦν δὲ ἐν τῷ τόπῳ ὅπου ἐσταυρώθη κῆπος, καὶ ἐν τῷ κήπῳ μνημεῖον καινόν ἐν ᾧ οὐδέπω οὐδεὶς ἦν τεθειμένος·
\vs{42}ἐκεῖ οὖν διὰ τὴν Παρασκευὴν τῶν Ἰουδαίων, ὅτι ἐγγὺς ἦν τὸ μνημεῖον, ἔθηκαν τὸν Ἰησοῦν.

\ch{20}
Τῇ δὲ μιᾷ τῶν σαββάτων Μαρία ἡ Μαγδαληνὴ ἔρχεται πρωῒ σκοτίας ἔτι οὔσης εἰς τὸ μνημεῖον καὶ βλέπει τὸν λίθον ἠρμένον ἐκ τοῦ μνημείου.
\vs{2}τρέχει οὖν καὶ ἔρχεται πρὸς Σίμωνα Πέτρον καὶ πρὸς τὸν ἄλλον μαθητὴν ὃν ἐφίλει ὁ Ἰησοῦς καὶ λέγει αὐτοῖς· Ἦραν τὸν Κύριον ἐκ τοῦ μνημείου καὶ οὐκ οἴδαμεν ποῦ ἔθηκαν αὐτόν.
\vs{3}Ἐξῆλθεν οὖν ὁ Πέτρος καὶ ὁ ἄλλος μαθητής καὶ ἤρχοντο εἰς τὸ μνημεῖον.
\vs{4}ἔτρεχον δὲ οἱ δύο ὁμοῦ· καὶ ὁ ἄλλος μαθητὴς προέδραμεν τάχιον τοῦ Πέτρου καὶ ἦλθεν πρῶτος εἰς τὸ μνημεῖον,
\vs{5}καὶ παρακύψας βλέπει κείμενα τὰ ὀθόνια, οὐ μέντοι εἰσῆλθεν.
\vs{6}Ἔρχεται οὖν καὶ Σίμων Πέτρος ἀκολουθῶν αὐτῷ καὶ εἰσῆλθεν εἰς τὸ μνημεῖον, καὶ θεωρεῖ τὰ ὀθόνια κείμενα,
\vs{7}καὶ τὸ σουδάριον, ὃ ἦν ἐπὶ τῆς κεφαλῆς αὐτοῦ, οὐ μετὰ τῶν ὀθονίων κείμενον ἀλλὰ χωρὶς ἐντετυλιγμένον εἰς ἕνα τόπον.
\vs{8}τότε οὖν εἰσῆλθεν καὶ ὁ ἄλλος μαθητὴς ὁ ἐλθὼν πρῶτος εἰς τὸ μνημεῖον καὶ εἶδεν καὶ ἐπίστευσεν·
\vs{9}οὐδέπω γὰρ ᾔδεισαν τὴν γραφὴν ὅτι δεῖ αὐτὸν ἐκ νεκρῶν ἀναστῆναι.
\vs{10}Ἀπῆλθον οὖν πάλιν πρὸς αὑτοὺς οἱ μαθηταί.

\vs{11}Μαρία δὲ εἱστήκει πρὸς τῷ μνημείῳ ἔξω κλαίουσα. ὡς οὖν ἔκλαιεν, παρέκυψεν εἰς τὸ μνημεῖον
\vs{12}καὶ θεωρεῖ δύο ἀγγέλους ἐν λευκοῖς καθεζομένους, ἕνα πρὸς τῇ κεφαλῇ καὶ ἕνα πρὸς τοῖς ποσίν, ὅπου ἔκειτο τὸ σῶμα τοῦ Ἰησοῦ.
\vs{13}Καὶ λέγουσιν αὐτῇ ἐκεῖνοι· Γύναι, τί κλαίεις; Λέγει αὐτοῖς Ὅτι Ἦραν τὸν Κύριόν μου, καὶ οὐκ οἶδα ποῦ ἔθηκαν αὐτόν.
\vs{14}Ταῦτα εἰποῦσα ἐστράφη εἰς τὰ ὀπίσω καὶ θεωρεῖ τὸν Ἰησοῦν ἑστῶτα καὶ οὐκ ᾔδει ὅτι Ἰησοῦς ἐστιν.
\vs{15}λέγει αὐτῇ Ἰησοῦς· Γύναι, τί κλαίεις; τίνα ζητεῖς; Ἐκείνη δοκοῦσα ὅτι ὁ κηπουρός ἐστιν λέγει αὐτῷ· Κύριε, εἰ σὺ ἐβάστασας αὐτόν, εἰπέ μοι ποῦ ἔθηκας αὐτόν, κἀγὼ αὐτὸν ἀρῶ.
\vs{16}Λέγει αὐτῇ Ἰησοῦς· Μαριάμ. Στραφεῖσα ἐκείνη λέγει αὐτῷ Ἑβραϊστί· Ραββουνι, ὃ λέγεται Διδάσκαλε.
\vs{17}Λέγει αὐτῇ Ἰησοῦς· Μή μου ἅπτου, οὔπω γὰρ ἀναβέβηκα πρὸς τὸν Πατέρα· πορεύου δὲ πρὸς τοὺς ἀδελφούς μου καὶ εἰπὲ αὐτοῖς· Ἀναβαίνω πρὸς τὸν Πατέρα μου καὶ Πατέρα ὑμῶν καὶ Θεόν μου καὶ Θεὸν ὑμῶν.
\vs{18}Ἔρχεται Μαριὰμ ἡ Μαγδαληνὴ ἀγγέλλουσα τοῖς μαθηταῖς ὅτι Ἑώρακα τὸν Κύριον, καὶ ταῦτα εἶπεν αὐτῇ.

\vs{19}Οὔσης οὖν ὀψίας τῇ ἡμέρᾳ ἐκείνῃ τῇ μιᾷ σαββάτων καὶ τῶν θυρῶν κεκλεισμένων ὅπου ἦσαν οἱ μαθηταὶ διὰ τὸν φόβον τῶν Ἰουδαίων, ἦλθεν ὁ Ἰησοῦς καὶ ἔστη εἰς τὸ μέσον καὶ λέγει αὐτοῖς· Εἰρήνη ὑμῖν.
\vs{20}καὶ τοῦτο εἰπὼν ἔδειξεν τὰς χεῖρας καὶ τὴν πλευρὰν αὐτοῖς. Ἐχάρησαν οὖν οἱ μαθηταὶ ἰδόντες τὸν Κύριον.
\vs{21}Εἶπεν οὖν αὐτοῖς ὁ Ἰησοῦς πάλιν· Εἰρήνη ὑμῖν· καθὼς ἀπέσταλκέν με ὁ Πατήρ, κἀγὼ πέμπω ὑμᾶς.
\vs{22}καὶ τοῦτο εἰπὼν ἐνεφύσησεν καὶ λέγει αὐτοῖς· Λάβετε Πνεῦμα Ἅγιον·
\vs{23}ἄν τινων ἀφῆτε τὰς ἁμαρτίας ἀφέωνται αὐτοῖς, ἄν τινων κρατῆτε κεκράτηνται.

\vs{24}Θωμᾶς δὲ εἷς ἐκ τῶν δώδεκα, ὁ λεγόμενος Δίδυμος, οὐκ ἦν μετ᾽ αὐτῶν ὅτε ἦλθεν Ἰησοῦς.
\vs{25}ἔλεγον οὖν αὐτῷ οἱ ἄλλοι μαθηταί· Ἑωράκαμεν τὸν Κύριον. Ὁ δὲ εἶπεν αὐτοῖς· Ἐὰν μὴ ἴδω ἐν ταῖς χερσὶν αὐτοῦ τὸν τύπον τῶν ἥλων καὶ βάλω τὸν δάκτυλόν μου εἰς τὸν τύπον τῶν ἥλων καὶ βάλω μου τὴν χεῖρα εἰς τὴν πλευρὰν αὐτοῦ, οὐ μὴ πιστεύσω.
\vs{26}Καὶ μεθ᾽ ἡμέρας ὀκτὼ πάλιν ἦσαν ἔσω οἱ μαθηταὶ αὐτοῦ καὶ Θωμᾶς μετ᾽ αὐτῶν. ἔρχεται ὁ Ἰησοῦς τῶν θυρῶν κεκλεισμένων καὶ ἔστη εἰς τὸ μέσον καὶ εἶπεν· Εἰρήνη ὑμῖν.
\vs{27}εἶτα λέγει τῷ Θωμᾷ· Φέρε τὸν δάκτυλόν σου ὧδε καὶ ἴδε τὰς χεῖράς μου καὶ φέρε τὴν χεῖρά σου καὶ βάλε εἰς τὴν πλευράν μου, καὶ μὴ γίνου ἄπιστος ἀλλὰ πιστός.
\vs{28}Ἀπεκρίθη Θωμᾶς καὶ εἶπεν αὐτῷ· Ὁ Κύριός μου καὶ ὁ Θεός μου.
\vs{29}Λέγει αὐτῷ ὁ Ἰησοῦς· Ὅτι ἑώρακάς με πεπίστευκας; μακάριοι οἱ μὴ ἰδόντες καὶ πιστεύσαντες.

\vs{30}Πολλὰ μὲν οὖν καὶ ἄλλα σημεῖα ἐποίησεν ὁ Ἰησοῦς ἐνώπιον τῶν μαθητῶν αὐτοῦ, ἃ οὐκ ἔστιν γεγραμμένα ἐν τῷ βιβλίῳ τούτῳ·
\vs{31}ταῦτα δὲ γέγραπται ἵνα πιστεύητε ὅτι Ἰησοῦς ἐστιν ὁ Χριστὸς ὁ Υἱὸς τοῦ Θεοῦ, καὶ ἵνα πιστεύοντες ζωὴν ἔχητε ἐν τῷ ὀνόματι αὐτοῦ.

\ch{21}
Μετὰ ταῦτα ἐφανέρωσεν ἑαυτὸν πάλιν ὁ Ἰησοῦς τοῖς μαθηταῖς ἐπὶ τῆς θαλάσσης τῆς Τιβεριάδος· ἐφανέρωσεν δὲ οὕτως.
\vs{2}ἦσαν ὁμοῦ Σίμων Πέτρος καὶ Θωμᾶς ὁ λεγόμενος Δίδυμος καὶ Ναθαναὴλ ὁ ἀπὸ Κανᾶ τῆς Γαλιλαίας καὶ οἱ τοῦ Ζεβεδαίου καὶ ἄλλοι ἐκ τῶν μαθητῶν αὐτοῦ δύο.
\vs{3}λέγει αὐτοῖς Σίμων Πέτρος· Ὑπάγω ἁλιεύειν. Λέγουσιν αὐτῷ· Ἐρχόμεθα καὶ ἡμεῖς σὺν σοί. ἐξῆλθον καὶ ἐνέβησαν εἰς τὸ πλοῖον, καὶ ἐν ἐκείνῃ τῇ νυκτὶ ἐπίασαν οὐδέν.
\vs{4}Πρωΐας δὲ ἤδη γενομένης ἔστη Ἰησοῦς εἰς τὸν αἰγιαλόν, οὐ μέντοι ᾔδεισαν οἱ μαθηταὶ ὅτι Ἰησοῦς ἐστιν.
\vs{5}λέγει οὖν αὐτοῖς ὁ Ἰησοῦς· Παιδία, μή τι προσφάγιον ἔχετε; Ἀπεκρίθησαν αὐτῷ· Οὔ.
\vs{6}Ὁ δὲ εἶπεν αὐτοῖς· Βάλετε εἰς τὰ δεξιὰ μέρη τοῦ πλοίου τὸ δίκτυον, καὶ εὑρήσετε. ἔβαλον οὖν, καὶ οὐκέτι αὐτὸ ἑλκύσαι ἴσχυον ἀπὸ τοῦ πλήθους τῶν ἰχθύων.
\vs{7}Λέγει οὖν ὁ μαθητὴς ἐκεῖνος ὃν ἠγάπα ὁ Ἰησοῦς τῷ Πέτρῳ· Ὁ Κύριός ἐστιν. Σίμων οὖν Πέτρος ἀκούσας ὅτι ὁ Κύριός ἐστιν τὸν ἐπενδύτην διεζώσατο, ἦν γὰρ γυμνός, καὶ ἔβαλεν ἑαυτὸν εἰς τὴν θάλασσαν,
\vs{8}οἱ δὲ ἄλλοι μαθηταὶ τῷ πλοιαρίῳ ἦλθον, οὐ γὰρ ἦσαν μακρὰν ἀπὸ τῆς γῆς ἀλλὰ ὡς ἀπὸ πηχῶν διακοσίων, σύροντες τὸ δίκτυον τῶν ἰχθύων.
\vs{9}Ὡς οὖν ἀπέβησαν εἰς τὴν γῆν βλέπουσιν ἀνθρακιὰν κειμένην καὶ ὀψάριον ἐπικείμενον καὶ ἄρτον.
\vs{10}Λέγει αὐτοῖς ὁ Ἰησοῦς· Ἐνέγκατε ἀπὸ τῶν ὀψαρίων ὧν ἐπιάσατε νῦν.
\vs{11}ἀνέβη οὖν Σίμων Πέτρος καὶ εἵλκυσεν τὸ δίκτυον εἰς τὴν γῆν μεστὸν ἰχθύων μεγάλων ἑκατὸν πεντήκοντα τριῶν· καὶ τοσούτων ὄντων οὐκ ἐσχίσθη τὸ δίκτυον.
\vs{12}Λέγει αὐτοῖς ὁ Ἰησοῦς· Δεῦτε ἀριστήσατε. οὐδεὶς δὲ ἐτόλμα τῶν μαθητῶν ἐξετάσαι αὐτόν· Σὺ τίς εἶ; εἰδότες ὅτι ὁ Κύριός ἐστιν.
\vs{13}ἔρχεται Ἰησοῦς καὶ λαμβάνει τὸν ἄρτον καὶ δίδωσιν αὐτοῖς, καὶ τὸ ὀψάριον ὁμοίως.
\vs{14}Τοῦτο ἤδη τρίτον ἐφανερώθη Ἰησοῦς τοῖς μαθηταῖς ἐγερθεὶς ἐκ νεκρῶν.

\vs{15}Ὅτε οὖν ἠρίστησαν λέγει τῷ Σίμωνι Πέτρῳ ὁ Ἰησοῦς· Σίμων Ἰωάννου, ἀγαπᾷς με πλέον τούτων; Λέγει αὐτῷ· Ναί Κύριε, σὺ οἶδας ὅτι φιλῶ σε. Λέγει αὐτῷ· Βόσκε τὰ ἀρνία μου.
\vs{16}Λέγει αὐτῷ πάλιν δεύτερον· Σίμων Ἰωάννου, ἀγαπᾷς με; Λέγει αὐτῷ· Ναί Κύριε, σὺ οἶδας ὅτι φιλῶ σε. Λέγει αὐτῷ· Ποίμαινε τὰ πρόβατά μου.
\vs{17}Λέγει αὐτῷ τὸ τρίτον· Σίμων Ἰωάννου, φιλεῖς με; Ἐλυπήθη ὁ Πέτρος ὅτι εἶπεν αὐτῷ τὸ τρίτον· Φιλεῖς με; Καὶ λέγει αὐτῷ· Κύριε, πάντα σὺ οἶδας, σὺ γινώσκεις ὅτι φιλῶ σε. Λέγει αὐτῷ ὁ Ἰησοῦς· Βόσκε τὰ πρόβατά μου.
\vs{18}Ἀμὴν ἀμὴν λέγω σοι, ὅτε ἦς νεώτερος, ἐζώννυες σεαυτὸν καὶ περιεπάτεις ὅπου ἤθελες· ὅταν δὲ γηράσῃς, ἐκτενεῖς τὰς χεῖράς σου, καὶ ἄλλος σε ζώσει καὶ οἴσει ὅπου οὐ θέλεις.
\vs{19}τοῦτο δὲ εἶπεν σημαίνων ποίῳ θανάτῳ δοξάσει τὸν Θεόν. Καὶ τοῦτο εἰπὼν λέγει αὐτῷ· Ἀκολούθει μοι.

\vs{20}Ἐπιστραφεὶς ὁ Πέτρος βλέπει τὸν μαθητὴν ὃν ἠγάπα ὁ Ἰησοῦς ἀκολουθοῦντα, ὃς καὶ ἀνέπεσεν ἐν τῷ δείπνῳ ἐπὶ τὸ στῆθος αὐτοῦ καὶ εἶπεν· Κύριε, τίς ἐστιν ὁ παραδιδούς σε;
\vs{21}τοῦτον οὖν ἰδὼν ὁ Πέτρος λέγει τῷ Ἰησοῦ· Κύριε, οὗτος δὲ τί;
\vs{22}Λέγει αὐτῷ ὁ Ἰησοῦς· Ἐὰν αὐτὸν θέλω μένειν ἕως ἔρχομαι, τί πρὸς σέ; σύ μοι ἀκολούθει.
\vs{23}ἐξῆλθεν οὖν οὗτος ὁ λόγος εἰς τοὺς ἀδελφοὺς ὅτι ὁ μαθητὴς ἐκεῖνος οὐκ ἀποθνήσκει· οὐκ εἶπεν δὲ αὐτῷ ὁ Ἰησοῦς ὅτι οὐκ ἀποθνήσκει ἀλλ᾽· Ἐὰν αὐτὸν θέλω μένειν ἕως ἔρχομαι, τί πρὸς σέ;
\vs{24}Οὗτός ἐστιν ὁ μαθητὴς ὁ μαρτυρῶν περὶ τούτων καὶ ὁ γράψας ταῦτα, καὶ οἴδαμεν ὅτι ἀληθὴς αὐτοῦ ἡ μαρτυρία ἐστίν.
\vs{25}Ἔστιν δὲ καὶ ἄλλα πολλὰ ἃ ἐποίησεν ὁ Ἰησοῦς, ἅτινα ἐὰν γράφηται καθ᾽ ἕν, οὐδ᾽ αὐτὸν οἶμαι τὸν κόσμον χωρήσειν τὰ γραφόμενα βιβλία.


\def\book{ΠΡΑΞΕΙΣ ΑΠΟΣΤΟΛΩΝ}
\biblebook{ΠΡΑΞΕΙΣ ΑΠΟΣΤΟΛΩΝ}


\lettrine[lines=2, loversize=0.2, nindent=0em, findent=.25em]{\textcolor{bookheadingcolor}{Τ}}{ὸν} μὲν πρῶτον λόγον ἐποιησάμην περὶ πάντων, ὦ Θεόφιλε, ὧν ἤρξατο ὁ Ἰησοῦς ποιεῖν τε καὶ διδάσκειν,
\vs{2}ἄχρι ἧς ἡμέρας ἐντειλάμενος τοῖς ἀποστόλοις διὰ Πνεύματος Ἁγίου οὓς ἐξελέξατο ἀνελήμφθη.
\vs{3}οἷς καὶ παρέστησεν ἑαυτὸν ζῶντα μετὰ τὸ παθεῖν αὐτὸν ἐν πολλοῖς τεκμηρίοις, δι᾽ ἡμερῶν τεσσεράκοντα ὀπτανόμενος αὐτοῖς καὶ λέγων τὰ περὶ τῆς βασιλείας τοῦ Θεοῦ·
\vs{4}Καὶ συναλιζόμενος παρήγγειλεν αὐτοῖς ἀπὸ Ἱεροσολύμων μὴ χωρίζεσθαι ἀλλὰ περιμένειν τὴν ἐπαγγελίαν τοῦ Πατρὸς Ἣν ἠκούσατέ μου,
\vs{5}ὅτι Ἰωάννης μὲν ἐβάπτισεν ὕδατι, ὑμεῖς δὲ ἐν Πνεύματι βαπτισθήσεσθε Ἁγίῳ οὐ μετὰ πολλὰς ταύτας ἡμέρας.
\vs{6}Οἱ μὲν οὖν συνελθόντες ἠρώτων αὐτὸν λέγοντες· Κύριε, εἰ ἐν τῷ χρόνῳ τούτῳ ἀποκαθιστάνεις τὴν βασιλείαν τῷ Ἰσραήλ;
\vs{7}Εἶπεν δὲ πρὸς αὐτούς· Οὐχ ὑμῶν ἐστιν γνῶναι χρόνους ἢ καιροὺς οὓς ὁ Πατὴρ ἔθετο ἐν τῇ ἰδίᾳ ἐξουσίᾳ,
\vs{8}ἀλλὰ λήμψεσθε δύναμιν ἐπελθόντος τοῦ Ἁγίου Πνεύματος ἐφ᾽ ὑμᾶς καὶ ἔσεσθέ μου μάρτυρες ἔν τε Ἰερουσαλὴμ καὶ ἐν πάσῃ τῇ Ἰουδαίᾳ καὶ Σαμαρείᾳ καὶ ἕως ἐσχάτου τῆς γῆς.

\vs{9}Καὶ ταῦτα εἰπὼν βλεπόντων αὐτῶν ἐπήρθη καὶ νεφέλη ὑπέλαβεν αὐτὸν ἀπὸ τῶν ὀφθαλμῶν αὐτῶν.
\vs{10}καὶ ὡς ἀτενίζοντες ἦσαν εἰς τὸν οὐρανὸν πορευομένου αὐτοῦ, καὶ ἰδοὺ ἄνδρες δύο παρειστήκεισαν αὐτοῖς ἐν ἐσθήσεσι λευκαῖς,
\vs{11}οἳ καὶ εἶπαν· Ἄνδρες Γαλιλαῖοι, τί ἑστήκατε βλέποντες εἰς τὸν οὐρανόν; οὗτος ὁ Ἰησοῦς ὁ ἀναλημφθεὶς ἀφ᾽ ὑμῶν εἰς τὸν οὐρανὸν οὕτως ἐλεύσεται ὃν τρόπον ἐθεάσασθε αὐτὸν πορευόμενον εἰς τὸν οὐρανόν.

\vs{12}Τότε ὑπέστρεψαν εἰς Ἰερουσαλὴμ ἀπὸ ὄρους τοῦ καλουμένου Ἐλαιῶνος, ὅ ἐστιν ἐγγὺς Ἰερουσαλὴμ σαββάτου ἔχον ὁδόν.
\vs{13}καὶ ὅτε εἰσῆλθον, εἰς τὸ ὑπερῷον ἀνέβησαν οὗ ἦσαν καταμένοντες, ὅ τε Πέτρος καὶ Ἰωάννης καὶ Ἰάκωβος καὶ Ἀνδρέας, Φίλιππος καὶ Θωμᾶς, Βαρθολομαῖος καὶ Μαθθαῖος, Ἰάκωβος Ἁλφαίου καὶ Σίμων ὁ Ζηλωτὴς καὶ Ἰούδας Ἰακώβου.
\vs{14}οὗτοι πάντες ἦσαν προσκαρτεροῦντες ὁμοθυμαδὸν τῇ προσευχῇ σὺν γυναιξὶν καὶ Μαριὰμ τῇ μητρὶ τοῦ Ἰησοῦ καὶ τοῖς ἀδελφοῖς αὐτοῦ.

\vs{15}Καὶ ἐν ταῖς ἡμέραις ταύταις ἀναστὰς Πέτρος ἐν μέσῳ τῶν ἀδελφῶν εἶπεν· ἦν τε ὄχλος ὀνομάτων ἐπὶ τὸ αὐτὸ ὡσεὶ ἑκατὸν εἴκοσι·
\vs{16}Ἄνδρες ἀδελφοί, ἔδει πληρωθῆναι τὴν γραφὴν ἣν προεῖπεν τὸ Πνεῦμα τὸ Ἅγιον διὰ στόματος Δαυὶδ περὶ Ἰούδα τοῦ γενομένου ὁδηγοῦ τοῖς συλλαβοῦσιν Ἰησοῦν,
\vs{17}ὅτι κατηριθμημένος ἦν ἐν ἡμῖν καὶ ἔλαχεν τὸν κλῆρον τῆς διακονίας ταύτης.
\vs{18}Οὗτος μὲν οὖν ἐκτήσατο χωρίον ἐκ μισθοῦ τῆς ἀδικίας καὶ πρηνὴς γενόμενος ἐλάκησεν μέσος καὶ ἐξεχύθη πάντα τὰ σπλάγχνα αὐτοῦ·
\vs{19}καὶ γνωστὸν ἐγένετο πᾶσι τοῖς κατοικοῦσιν Ἰερουσαλήμ, ὥστε κληθῆναι τὸ χωρίον ἐκεῖνο τῇ ἰδίᾳ διαλέκτῳ αὐτῶν Ἁκελδαμάχ, τοῦτ᾽ ἔστιν Χωρίον αἵματος.
\vs{20}Γέγραπται γὰρ ἐν βίβλῳ Ψαλμῶν· 
\begin{poetryblock}

\begin{quote}Γενηθήτω ἡ ἔπαυλις αὐτοῦ ἔρημος\end{quote} 

\begin{quote}καὶ μὴ ἔστω ὁ κατοικῶν ἐν αὐτῇ,\end{quote}
\end{poetryblock}

Καί· 
\begin{poetryblock}

\begin{quote}Τὴν ἐπισκοπὴν αὐτοῦ λαβέτω ἕτερος.\end{quote}
\end{poetryblock}

\vs{21}Δεῖ οὖν τῶν συνελθόντων ἡμῖν ἀνδρῶν ἐν παντὶ χρόνῳ ᾧ εἰσῆλθεν καὶ ἐξῆλθεν ἐφ᾽ ἡμᾶς ὁ Κύριος Ἰησοῦς,
\vs{22}ἀρξάμενος ἀπὸ τοῦ βαπτίσματος Ἰωάννου ἕως τῆς ἡμέρας ἧς ἀνελήμφθη ἀφ᾽ ἡμῶν, μάρτυρα τῆς ἀναστάσεως αὐτοῦ σὺν ἡμῖν γενέσθαι ἕνα τούτων.
\vs{23}Καὶ ἔστησαν δύο, Ἰωσὴφ τὸν καλούμενον Βαρσαββᾶν ὃς ἐπεκλήθη Ἰοῦστος, καὶ Μαθθίαν.
\vs{24}καὶ προσευξάμενοι εἶπαν· Σὺ Κύριε καρδιογνῶστα πάντων, ἀνάδειξον ὃν ἐξελέξω ἐκ τούτων τῶν δύο ἕνα
\vs{25}λαβεῖν τὸν τόπον τῆς διακονίας ταύτης καὶ ἀποστολῆς ἀφ᾽ ἧς παρέβη Ἰούδας πορευθῆναι εἰς τὸν τόπον τὸν ἴδιον.
\vs{26}καὶ ἔδωκαν κλήρους αὐτοῖς καὶ ἔπεσεν ὁ κλῆρος ἐπὶ Μαθθίαν καὶ συνκατεψηφίσθη μετὰ τῶν ἕνδεκα ἀποστόλων.

\ch{2}
Καὶ ἐν τῷ συμπληροῦσθαι τὴν ἡμέραν τῆς Πεντηκοστῆς ἦσαν πάντες ὁμοῦ ἐπὶ τὸ αὐτό.
\vs{2}καὶ ἐγένετο ἄφνω ἐκ τοῦ οὐρανοῦ ἦχος ὥσπερ φερομένης πνοῆς βιαίας καὶ ἐπλήρωσεν ὅλον τὸν οἶκον οὗ ἦσαν καθήμενοι
\vs{3}καὶ ὤφθησαν αὐτοῖς διαμεριζόμεναι γλῶσσαι ὡσεὶ πυρός καὶ ἐκάθισεν ἐφ᾽ ἕνα ἕκαστον αὐτῶν,
\vs{4}καὶ ἐπλήσθησαν πάντες Πνεύματος Ἁγίου καὶ ἤρξαντο λαλεῖν ἑτέραις γλώσσαις καθὼς τὸ Πνεῦμα ἐδίδου ἀποφθέγγεσθαι αὐτοῖς.

\vs{5}Ἦσαν δὲ εἰς Ἰερουσαλὴμ κατοικοῦντες Ἰουδαῖοι, ἄνδρες εὐλαβεῖς ἀπὸ παντὸς ἔθνους τῶν ὑπὸ τὸν οὐρανόν.
\vs{6}γενομένης δὲ τῆς φωνῆς ταύτης συνῆλθεν τὸ πλῆθος καὶ συνεχύθη, ὅτι ἤκουον εἷς ἕκαστος τῇ ἰδίᾳ διαλέκτῳ λαλούντων αὐτῶν.
\vs{7}Ἐξίσταντο δὲ καὶ ἐθαύμαζον λέγοντες· Οὐχ ἰδοὺ ἅπαντες οὗτοί εἰσιν οἱ λαλοῦντες Γαλιλαῖοι;
\vs{8}καὶ πῶς ἡμεῖς ἀκούομεν ἕκαστος τῇ ἰδίᾳ διαλέκτῳ ἡμῶν ἐν ᾗ ἐγεννήθημεν;
\vs{9}Πάρθοι καὶ Μῆδοι καὶ Ἐλαμῖται καὶ οἱ κατοικοῦντες τὴν Μεσοποταμίαν, Ἰουδαίαν τε καὶ Καππαδοκίαν, Πόντον καὶ τὴν Ἀσίαν,
\vs{10}Φρυγίαν τε καὶ Παμφυλίαν, Αἴγυπτον καὶ τὰ μέρη τῆς Λιβύης τῆς κατὰ Κυρήνην, καὶ οἱ ἐπιδημοῦντες Ῥωμαῖοι,
\vs{11}Ἰουδαῖοί τε καὶ προσήλυτοι, Κρῆτες καὶ Ἄραβες, ἀκούομεν λαλούντων αὐτῶν ταῖς ἡμετέραις γλώσσαις τὰ μεγαλεῖα τοῦ Θεοῦ.
\vs{12}Ἐξίσταντο δὲ πάντες καὶ διηπόρουν, ἄλλος πρὸς ἄλλον λέγοντες· Τί θέλει τοῦτο εἶναι;
\vs{13}Ἕτεροι δὲ διαχλευάζοντες ἔλεγον ὅτι Γλεύκους μεμεστωμένοι εἰσίν.

\vs{14}Σταθεὶς δὲ ὁ Πέτρος σὺν τοῖς ἕνδεκα ἐπῆρεν τὴν φωνὴν αὐτοῦ καὶ ἀπεφθέγξατο αὐτοῖς· Ἄνδρες Ἰουδαῖοι καὶ οἱ κατοικοῦντες Ἰερουσαλὴμ πάντες, τοῦτο ὑμῖν γνωστὸν ἔστω καὶ ἐνωτίσασθε τὰ ῥήματά μου.
\vs{15}οὐ γὰρ ὡς ὑμεῖς ὑπολαμβάνετε οὗτοι μεθύουσιν, ἔστιν γὰρ ὥρα τρίτη τῆς ἡμέρας,
\vs{16}ἀλλὰ τοῦτό ἐστιν τὸ εἰρημένον διὰ τοῦ προφήτου Ἰωήλ·
\begin{poetryblock}

\begin{quote} \vs{17}Καὶ ἔσται ἐν ταῖς ἐσχάταις ἡμέραις, λέγει ὁ Θεός,\end{quote} 

\begin{quote}ἐκχεῶ ἀπὸ τοῦ Πνεύματός μου ἐπὶ πᾶσαν σάρκα,\end{quote} 

\begin{quote}καὶ προφητεύσουσιν οἱ υἱοὶ ὑμῶν καὶ αἱ θυγατέρες ὑμῶν\end{quote} 

\begin{quote}καὶ οἱ νεανίσκοι ὑμῶν ὁράσεις ὄψονται\end{quote} 

\begin{quote}καὶ οἱ πρεσβύτεροι ὑμῶν ἐνυπνίοις ἐνυπνιασθήσονται·\end{quote}

\begin{quote} \vs{18}καί γε ἐπὶ τοὺς δούλους μου καὶ ἐπὶ τὰς δούλας μου ἐν ταῖς ἡμέραις ἐκείναις\end{quote} 

\begin{quote}ἐκχεῶ ἀπὸ τοῦ Πνεύματός μου, καὶ προφητεύσουσιν.\end{quote}

\begin{quote} \vs{19}καὶ δώσω τέρατα ἐν τῷ οὐρανῷ ἄνω\end{quote} 

\begin{quote}καὶ σημεῖα ἐπὶ τῆς γῆς κάτω,\end{quote} 

\begin{quote}αἷμα καὶ πῦρ καὶ ἀτμίδα καπνοῦ.\end{quote}

\begin{quote} \vs{20}ὁ ἥλιος μεταστραφήσεται εἰς σκότος\end{quote} 

\begin{quote}καὶ ἡ σελήνη εἰς αἷμα,\end{quote} 

\begin{quote}πρὶν ἐλθεῖν ἡμέραν Κυρίου τὴν μεγάλην καὶ ἐπιφανῆ.\end{quote}

\begin{quote} \vs{21}καὶ ἔσται πᾶς ὃς ἂν ἐπικαλέσηται τὸ ὄνομα Κυρίου σωθήσεται.\end{quote}
\end{poetryblock}

\vs{22}Ἄνδρες Ἰσραηλῖται, ἀκούσατε τοὺς λόγους τούτους· Ἰησοῦν τὸν Ναζωραῖον, ἄνδρα ἀποδεδειγμένον ἀπὸ τοῦ Θεοῦ εἰς ὑμᾶς δυνάμεσι καὶ τέρασι καὶ σημείοις οἷς ἐποίησεν δι᾽ αὐτοῦ ὁ Θεὸς ἐν μέσῳ ὑμῶν καθὼς αὐτοὶ οἴδατε,
\vs{23}τοῦτον τῇ ὡρισμένῃ βουλῇ καὶ προγνώσει τοῦ Θεοῦ ἔκδοτον διὰ χειρὸς ἀνόμων προσπήξαντες ἀνείλατε,
\vs{24}ὃν ὁ Θεὸς ἀνέστησεν λύσας τὰς ὠδῖνας τοῦ θανάτου, καθότι οὐκ ἦν δυνατὸν κρατεῖσθαι αὐτὸν ὑπ᾽ αὐτοῦ.
\vs{25}Δαυὶδ γὰρ λέγει εἰς αὐτόν· 
\begin{poetryblock}

\begin{quote}Προορώμην τὸν Κύριον ἐνώπιόν μου διὰ παντός,\end{quote} 

\begin{quote}ὅτι ἐκ δεξιῶν μού ἐστιν ἵνα μὴ σαλευθῶ.\end{quote}

\begin{quote} \vs{26}διὰ τοῦτο ηὐφράνθη ἡ καρδία μου καὶ ἠγαλλιάσατο ἡ γλῶσσά μου,\end{quote} 

\begin{quote}ἔτι δὲ καὶ ἡ σάρξ μου κατασκηνώσει ἐπ᾽ ἐλπίδι,\end{quote}

\begin{quote} \vs{27}ὅτι οὐκ ἐνκαταλείψεις τὴν ψυχήν μου εἰς ᾅδην\end{quote} 

\begin{quote}οὐδὲ δώσεις τὸν Ὅσιόν σου ἰδεῖν διαφθοράν.\end{quote}

\begin{quote} \vs{28}ἐγνώρισάς μοι ὁδοὺς ζωῆς,\end{quote} 

\begin{quote}πληρώσεις με εὐφροσύνης μετὰ τοῦ προσώπου σου.\end{quote}
\end{poetryblock}

\vs{29}Ἄνδρες ἀδελφοί, ἐξὸν εἰπεῖν μετὰ παρρησίας πρὸς ὑμᾶς περὶ τοῦ πατριάρχου Δαυὶδ ὅτι καὶ ἐτελεύτησεν καὶ ἐτάφη, καὶ τὸ μνῆμα αὐτοῦ ἔστιν ἐν ἡμῖν ἄχρι τῆς ἡμέρας ταύτης.
\vs{30}προφήτης οὖν ὑπάρχων καὶ εἰδὼς ὅτι ὅρκῳ ὤμοσεν αὐτῷ ὁ Θεὸς ἐκ καρποῦ τῆς ὀσφύος αὐτοῦ καθίσαι ἐπὶ τὸν θρόνον αὐτοῦ,
\vs{31}προϊδὼν ἐλάλησεν περὶ τῆς ἀναστάσεως τοῦ Χριστοῦ ὅτι οὔτε ἐνκατελείφθη εἰς ᾅδην οὔτε ἡ σὰρξ αὐτοῦ εἶδεν διαφθοράν.
\vs{32}Τοῦτον τὸν Ἰησοῦν ἀνέστησεν ὁ Θεός, οὗ πάντες ἡμεῖς ἐσμεν μάρτυρες·
\vs{33}τῇ δεξιᾷ οὖν τοῦ Θεοῦ ὑψωθεὶς, τήν τε ἐπαγγελίαν τοῦ Πνεύματος τοῦ Ἁγίου λαβὼν παρὰ τοῦ Πατρὸς, ἐξέχεεν τοῦτο ὃ ὑμεῖς καὶ βλέπετε καὶ ἀκούετε.
\vs{34}Οὐ γὰρ Δαυὶδ ἀνέβη εἰς τοὺς οὐρανούς, λέγει δὲ αὐτός· 
\begin{poetryblock}

\begin{quote}Εἶπεν ὁ Κύριος τῷ Κυρίῳ μου· Κάθου ἐκ δεξιῶν μου,\end{quote}

\begin{quote} \vs{35}ἕως ἂν θῶ τοὺς ἐχθρούς σου ὑποπόδιον τῶν ποδῶν σου.\end{quote}
\end{poetryblock}

\vs{36}Ἀσφαλῶς οὖν γινωσκέτω πᾶς οἶκος Ἰσραὴλ ὅτι καὶ Κύριον αὐτὸν καὶ Χριστὸν ἐποίησεν ὁ Θεός, τοῦτον τὸν Ἰησοῦν ὃν ὑμεῖς ἐσταυρώσατε.

\vs{37}Ἀκούσαντες δὲ κατενύγησαν τὴν καρδίαν εἶπόν τε πρὸς τὸν Πέτρον καὶ τοὺς λοιποὺς ἀποστόλους· Τί ποιήσωμεν, ἄνδρες ἀδελφοί;
\vs{38}Πέτρος δὲ πρὸς αὐτούς· Μετανοήσατε, φησίν, Καὶ βαπτισθήτω ἕκαστος ὑμῶν ἐπὶ τῷ ὀνόματι Ἰησοῦ Χριστοῦ εἰς ἄφεσιν τῶν ἁμαρτιῶν ὑμῶν καὶ λήμψεσθε τὴν δωρεὰν τοῦ Ἁγίου Πνεύματος.
\vs{39}ὑμῖν γάρ ἐστιν ἡ ἐπαγγελία καὶ τοῖς τέκνοις ὑμῶν καὶ πᾶσιν τοῖς εἰς μακρὰν, ὅσους ἂν προσκαλέσηται Κύριος ὁ Θεὸς ἡμῶν.
\vs{40}Ἑτέροις τε λόγοις πλείοσιν διεμαρτύρατο καὶ παρεκάλει αὐτοὺς λέγων· Σώθητε ἀπὸ τῆς γενεᾶς τῆς σκολιᾶς ταύτης.
\vs{41}οἱ μὲν οὖν ἀποδεξάμενοι τὸν λόγον αὐτοῦ ἐβαπτίσθησαν καὶ προσετέθησαν ἐν τῇ ἡμέρᾳ ἐκείνῃ ψυχαὶ ὡσεὶ τρισχίλιαι.

\vs{42}Ἦσαν δὲ προσκαρτεροῦντες τῇ διδαχῇ τῶν ἀποστόλων καὶ τῇ κοινωνίᾳ, τῇ κλάσει τοῦ ἄρτου καὶ ταῖς προσευχαῖς.
\vs{43}Ἐγίνετο δὲ πάσῃ ψυχῇ φόβος, πολλά τε τέρατα καὶ σημεῖα διὰ τῶν ἀποστόλων ἐγίνετο.
\vs{44}Πάντες δὲ οἱ πιστεύοντες ἦσαν ἐπὶ τὸ αὐτὸ καὶ εἶχον ἅπαντα κοινά
\vs{45}καὶ τὰ κτήματα καὶ τὰς ὑπάρξεις ἐπίπρασκον καὶ διεμέριζον αὐτὰ πᾶσιν καθότι ἄν τις χρείαν εἶχεν·
\vs{46}Καθ᾽ ἡμέραν τε προσκαρτεροῦντες ὁμοθυμαδὸν ἐν τῷ ἱερῷ, κλῶντές τε κατ᾽ οἶκον ἄρτον, μετελάμβανον τροφῆς ἐν ἀγαλλιάσει καὶ ἀφελότητι καρδίας
\vs{47}αἰνοῦντες τὸν Θεὸν καὶ ἔχοντες χάριν πρὸς ὅλον τὸν λαόν. ὁ δὲ Κύριος προσετίθει τοὺς σῳζομένους καθ᾽ ἡμέραν ἐπὶ τὸ αὐτό.

\ch{3}
Πέτρος δὲ καὶ Ἰωάννης ἀνέβαινον εἰς τὸ ἱερὸν ἐπὶ τὴν ὥραν τῆς προσευχῆς τὴν ἐνάτην.
\vs{2}καί τις ἀνὴρ χωλὸς ἐκ κοιλίας μητρὸς αὐτοῦ ὑπάρχων ἐβαστάζετο, ὃν ἐτίθουν καθ᾽ ἡμέραν πρὸς τὴν θύραν τοῦ ἱεροῦ τὴν λεγομένην Ὡραίαν τοῦ αἰτεῖν ἐλεημοσύνην παρὰ τῶν εἰσπορευομένων εἰς τὸ ἱερόν·
\vs{3}ὃς ἰδὼν Πέτρον καὶ Ἰωάννην μέλλοντας εἰσιέναι εἰς τὸ ἱερὸν, ἠρώτα ἐλεημοσύνην λαβεῖν.
\vs{4}Ἀτενίσας δὲ Πέτρος εἰς αὐτὸν σὺν τῷ Ἰωάννῃ εἶπεν· Βλέψον εἰς ἡμᾶς.
\vs{5}ὁ δὲ ἐπεῖχεν αὐτοῖς προσδοκῶν τι παρ᾽ αὐτῶν λαβεῖν.
\vs{6}εἶπεν δὲ Πέτρος· Ἀργύριον καὶ χρυσίον οὐχ ὑπάρχει μοι, ὃ δὲ ἔχω τοῦτό σοι δίδωμι· ἐν τῷ ὀνόματι Ἰησοῦ Χριστοῦ τοῦ Ναζωραίου ἔγειρε καὶ περιπάτει.
\vs{7}Καὶ πιάσας αὐτὸν τῆς δεξιᾶς χειρὸς ἤγειρεν αὐτόν· παραχρῆμα δὲ ἐστερεώθησαν αἱ βάσεις αὐτοῦ καὶ τὰ σφυδρά,
\vs{8}καὶ ἐξαλλόμενος ἔστη καὶ περιεπάτει καὶ εἰσῆλθεν σὺν αὐτοῖς εἰς τὸ ἱερὸν περιπατῶν καὶ ἁλλόμενος καὶ αἰνῶν τὸν Θεόν.
\vs{9}Καὶ εἶδεν πᾶς ὁ λαὸς αὐτὸν περιπατοῦντα καὶ αἰνοῦντα τὸν Θεόν·
\vs{10}ἐπεγίνωσκον δὲ αὐτὸν ὅτι αὐτὸς ἦν ὁ πρὸς τὴν ἐλεημοσύνην καθήμενος ἐπὶ τῇ Ὡραίᾳ Πύλῃ τοῦ ἱεροῦ καὶ ἐπλήσθησαν θάμβους καὶ ἐκστάσεως ἐπὶ τῷ συμβεβηκότι αὐτῷ.

\vs{11}Κρατοῦντος δὲ αὐτοῦ τὸν Πέτρον καὶ τὸν Ἰωάννην συνέδραμεν πᾶς ὁ λαὸς πρὸς αὐτοὺς ἐπὶ τῇ στοᾷ τῇ καλουμένῃ Σολομῶντος ἔκθαμβοι.
\vs{12}ἰδὼν δὲ ὁ Πέτρος ἀπεκρίνατο πρὸς τὸν λαόν· Ἄνδρες Ἰσραηλῖται, τί θαυμάζετε ἐπὶ τούτῳ ἢ ἡμῖν τί ἀτενίζετε ὡς ἰδίᾳ δυνάμει ἢ εὐσεβείᾳ πεποιηκόσιν τοῦ περιπατεῖν αὐτόν;
\vs{13}Ὁ Θεὸς Ἀβραὰμ καὶ ὁ θεὸς Ἰσαὰκ καὶ ὁ θεὸς Ἰακώβ, ὁ Θεὸς τῶν πατέρων ἡμῶν, ἐδόξασεν τὸν Παῖδα αὐτοῦ Ἰησοῦν ὃν ὑμεῖς μὲν παρεδώκατε καὶ ἠρνήσασθε κατὰ πρόσωπον Πιλάτου, κρίναντος ἐκείνου ἀπολύειν·
\vs{14}ὑμεῖς δὲ τὸν Ἅγιον καὶ Δίκαιον ἠρνήσασθε καὶ ᾐτήσασθε ἄνδρα φονέα χαρισθῆναι ὑμῖν,
\vs{15}τὸν δὲ Ἀρχηγὸν τῆς ζωῆς ἀπεκτείνατε ὃν ὁ Θεὸς ἤγειρεν ἐκ νεκρῶν, οὗ ἡμεῖς μάρτυρές ἐσμεν.
\vs{16}Καὶ ἐπὶ τῇ πίστει τοῦ ὀνόματος αὐτοῦ τοῦτον ὃν θεωρεῖτε καὶ οἴδατε, ἐστερέωσεν τὸ ὄνομα αὐτοῦ, καὶ ἡ πίστις ἡ δι᾽ αὐτοῦ ἔδωκεν αὐτῷ τὴν ὁλοκληρίαν ταύτην ἀπέναντι πάντων ὑμῶν.
\vs{17}Καὶ νῦν, ἀδελφοί, οἶδα ὅτι κατὰ ἄγνοιαν ἐπράξατε ὥσπερ καὶ οἱ ἄρχοντες ὑμῶν·
\vs{18}ὁ δὲ Θεὸς, ἃ προκατήγγειλεν διὰ στόματος πάντων τῶν προφητῶν παθεῖν τὸν Χριστὸν αὐτοῦ, ἐπλήρωσεν οὕτως.
\vs{19}μετανοήσατε οὖν καὶ ἐπιστρέψατε πρὸς τὸ ἐξαλειφθῆναι ὑμῶν τὰς ἁμαρτίας,
\vs{20}ὅπως ἂν ἔλθωσιν καιροὶ ἀναψύξεως ἀπὸ προσώπου τοῦ Κυρίου καὶ ἀποστείλῃ τὸν προκεχειρισμένον ὑμῖν Χριστὸν Ἰησοῦν,
\vs{21}ὃν δεῖ οὐρανὸν μὲν δέξασθαι ἄχρι χρόνων ἀποκαταστάσεως πάντων ὧν ἐλάλησεν ὁ Θεὸς διὰ στόματος τῶν ἁγίων ἀπ᾽ αἰῶνος αὐτοῦ προφητῶν.
\vs{22}Μωϋσῆς μὲν εἶπεν ὅτι Προφήτην ὑμῖν ἀναστήσει Κύριος ὁ Θεὸς ὑμῶν ἐκ τῶν ἀδελφῶν ὑμῶν ὡς ἐμέ· αὐτοῦ ἀκούσεσθε κατὰ πάντα ὅσα ἂν λαλήσῃ πρὸς ὑμᾶς.
\vs{23}ἔσται δὲ πᾶσα ψυχὴ ἥτις ἐὰν μὴ ἀκούσῃ τοῦ προφήτου ἐκείνου ἐξολεθρευθήσεται ἐκ τοῦ λαοῦ.
\vs{24}Καὶ πάντες δὲ οἱ προφῆται ἀπὸ Σαμουὴλ καὶ τῶν καθεξῆς ὅσοι ἐλάλησαν καὶ κατήγγειλαν τὰς ἡμέρας ταύτας.
\vs{25}ὑμεῖς ἐστε οἱ υἱοὶ τῶν προφητῶν καὶ τῆς διαθήκης ἧς διέθετο ὁ Θεὸς πρὸς τοὺς πατέρας ὑμῶν λέγων πρὸς Ἀβραάμ· Καὶ ἐν τῷ σπέρματί σου ἐνευλογηθήσονται πᾶσαι αἱ πατριαὶ τῆς γῆς.
\vs{26}ὑμῖν πρῶτον ἀναστήσας ὁ Θεὸς τὸν Παῖδα αὐτοῦ ἀπέστειλεν αὐτὸν εὐλογοῦντα ὑμᾶς ἐν τῷ ἀποστρέφειν ἕκαστον ἀπὸ τῶν πονηριῶν ὑμῶν.

\ch{4}
Λαλούντων δὲ αὐτῶν πρὸς τὸν λαὸν ἐπέστησαν αὐτοῖς οἱ ἱερεῖς καὶ ὁ στρατηγὸς τοῦ ἱεροῦ καὶ οἱ Σαδδουκαῖοι,
\vs{2}διαπονούμενοι διὰ τὸ διδάσκειν αὐτοὺς τὸν λαὸν καὶ καταγγέλλειν ἐν τῷ Ἰησοῦ τὴν ἀνάστασιν τὴν ἐκ νεκρῶν,
\vs{3}καὶ ἐπέβαλον αὐτοῖς τὰς χεῖρας καὶ ἔθεντο εἰς τήρησιν εἰς τὴν αὔριον· ἦν γὰρ ἑσπέρα ἤδη.
\vs{4}πολλοὶ δὲ τῶν ἀκουσάντων τὸν λόγον ἐπίστευσαν καὶ ἐγενήθη ὁ ἀριθμὸς τῶν ἀνδρῶν ὡς χιλιάδες πέντε.

\vs{5}Ἐγένετο δὲ ἐπὶ τὴν αὔριον συναχθῆναι αὐτῶν τοὺς ἄρχοντας καὶ τοὺς πρεσβυτέρους καὶ τοὺς γραμματεῖς ἐν Ἰερουσαλήμ,
\vs{6}καὶ Ἅννας ὁ ἀρχιερεὺς καὶ Καϊάφας καὶ Ἰωάννης καὶ Ἀλέξανδρος καὶ ὅσοι ἦσαν ἐκ γένους ἀρχιερατικοῦ,
\vs{7}καὶ στήσαντες αὐτοὺς ἐν τῷ μέσῳ ἐπυνθάνοντο· Ἐν ποίᾳ δυνάμει ἢ ἐν ποίῳ ὀνόματι ἐποιήσατε τοῦτο ὑμεῖς;
\vs{8}Τότε Πέτρος πλησθεὶς Πνεύματος Ἁγίου εἶπεν πρὸς αὐτούς· Ἄρχοντες τοῦ λαοῦ καὶ πρεσβύτεροι,
\vs{9}εἰ ἡμεῖς σήμερον ἀνακρινόμεθα ἐπὶ εὐεργεσίᾳ ἀνθρώπου ἀσθενοῦς ἐν τίνι οὗτος σέσωσται,
\vs{10}γνωστὸν ἔστω πᾶσιν ὑμῖν καὶ παντὶ τῷ λαῷ Ἰσραὴλ ὅτι ἐν τῷ ὀνόματι Ἰησοῦ Χριστοῦ τοῦ Ναζωραίου ὃν ὑμεῖς ἐσταυρώσατε, ὃν ὁ Θεὸς ἤγειρεν ἐκ νεκρῶν, ἐν τούτῳ οὗτος παρέστηκεν ἐνώπιον ὑμῶν ὑγιής.
\vs{11}οὗτός ἐστιν Ὁ λίθος, ὁ ἐξουθενηθεὶς ὑφ᾽ ὑμῶν τῶν οἰκοδόμων, ὁ γενόμενος εἰς κεφαλὴν γωνίας.
\vs{12}Καὶ οὐκ ἔστιν ἐν ἄλλῳ οὐδενὶ ἡ σωτηρία, οὐδὲ γὰρ ὄνομά ἐστιν ἕτερον ὑπὸ τὸν οὐρανὸν τὸ δεδομένον ἐν ἀνθρώποις ἐν ᾧ δεῖ σωθῆναι ἡμᾶς.

\vs{13}Θεωροῦντες δὲ τὴν τοῦ Πέτρου παρρησίαν καὶ Ἰωάννου καὶ καταλαβόμενοι ὅτι ἄνθρωποι ἀγράμματοί εἰσιν καὶ ἰδιῶται, ἐθαύμαζον ἐπεγίνωσκόν τε αὐτοὺς ὅτι σὺν τῷ Ἰησοῦ ἦσαν,
\vs{14}τόν τε ἄνθρωπον βλέποντες σὺν αὐτοῖς ἑστῶτα τὸν τεθεραπευμένον οὐδὲν εἶχον ἀντειπεῖν.
\vs{15}κελεύσαντες δὲ αὐτοὺς ἔξω τοῦ συνεδρίου ἀπελθεῖν συνέβαλλον πρὸς ἀλλήλους
\vs{16}λέγοντες· Τί ποιήσωμεν τοῖς ἀνθρώποις τούτοις; ὅτι μὲν γὰρ γνωστὸν σημεῖον γέγονεν δι᾽ αὐτῶν πᾶσιν τοῖς κατοικοῦσιν Ἰερουσαλὴμ φανερόν καὶ οὐ δυνάμεθα ἀρνεῖσθαι·
\vs{17}ἀλλ᾽ ἵνα μὴ ἐπὶ πλεῖον διανεμηθῇ εἰς τὸν λαόν ἀπειλησώμεθα αὐτοῖς μηκέτι λαλεῖν ἐπὶ τῷ ὀνόματι τούτῳ μηδενὶ ἀνθρώπων.
\vs{18}Καὶ καλέσαντες αὐτοὺς παρήγγειλαν τὸ καθόλου μὴ φθέγγεσθαι μηδὲ διδάσκειν ἐπὶ τῷ ὀνόματι τοῦ Ἰησοῦ.
\vs{19}Ὁ δὲ Πέτρος καὶ Ἰωάννης ἀποκριθέντες εἶπον πρὸς αὐτούς· Εἰ δίκαιόν ἐστιν ἐνώπιον τοῦ Θεοῦ ὑμῶν ἀκούειν μᾶλλον ἢ τοῦ Θεοῦ, κρίνατε·
\vs{20}οὐ δυνάμεθα γὰρ ἡμεῖς ἃ εἴδαμεν καὶ ἠκούσαμεν μὴ λαλεῖν.
\vs{21}Οἱ δὲ προσαπειλησάμενοι ἀπέλυσαν αὐτούς, μηδὲν εὑρίσκοντες τὸ πῶς κολάσωνται αὐτούς, διὰ τὸν λαόν, ὅτι πάντες ἐδόξαζον τὸν Θεὸν ἐπὶ τῷ γεγονότι·
\vs{22}ἐτῶν γὰρ ἦν πλειόνων τεσσεράκοντα ὁ ἄνθρωπος ἐφ᾽ ὃν γεγόνει τὸ σημεῖον τοῦτο τῆς ἰάσεως.

\vs{23}Ἀπολυθέντες δὲ ἦλθον πρὸς τοὺς ἰδίους καὶ ἀπήγγειλαν ὅσα πρὸς αὐτοὺς οἱ ἀρχιερεῖς καὶ οἱ πρεσβύτεροι εἶπαν.
\vs{24}οἱ δὲ ἀκούσαντες ὁμοθυμαδὸν ἦραν φωνὴν πρὸς τὸν Θεὸν καὶ εἶπαν· Δέσποτα, σὺ ὁ ποιήσας τὸν οὐρανὸν καὶ τὴν γῆν καὶ τὴν θάλασσαν καὶ πάντα τὰ ἐν αὐτοῖς,
\vs{25}ὁ τοῦ πατρὸς ἡμῶν διὰ Πνεύματος Ἁγίου στόματος Δαυὶδ παιδός σου εἰπών· 
\begin{poetryblock}

\begin{quote}Ἵνατί ἐφρύαξαν ἔθνη\end{quote} 

\begin{quote}καὶ λαοὶ ἐμελέτησαν κενά;\end{quote}

\begin{quote} \vs{26}παρέστησαν οἱ βασιλεῖς τῆς γῆς\end{quote} 

\begin{quote}καὶ οἱ ἄρχοντες συνήχθησαν ἐπὶ τὸ αὐτὸ\end{quote} 

\begin{quote}κατὰ τοῦ Κυρίου καὶ κατὰ τοῦ Χριστοῦ αὐτοῦ.\end{quote}
\end{poetryblock}

\vs{27}Συνήχθησαν γὰρ ἐπ᾽ ἀληθείας ἐν τῇ πόλει ταύτῃ ἐπὶ τὸν ἅγιον Παῖδά σου Ἰησοῦν ὃν ἔχρισας, Ἡρῴδης τε καὶ Πόντιος Πιλᾶτος σὺν ἔθνεσιν καὶ λαοῖς Ἰσραήλ,
\vs{28}ποιῆσαι ὅσα ἡ χείρ σου καὶ ἡ βουλὴ σου προώρισεν γενέσθαι.
\vs{29}καὶ τὰ νῦν, Κύριε, ἔπιδε ἐπὶ τὰς ἀπειλὰς αὐτῶν καὶ δὸς τοῖς δούλοις σου μετὰ παρρησίας πάσης λαλεῖν τὸν λόγον σου,
\vs{30}ἐν τῷ τὴν χεῖρά σου ἐκτείνειν σε εἰς ἴασιν καὶ σημεῖα καὶ τέρατα γίνεσθαι διὰ τοῦ ὀνόματος τοῦ ἁγίου Παιδός σου Ἰησοῦ.
\vs{31}Καὶ δεηθέντων αὐτῶν ἐσαλεύθη ὁ τόπος ἐν ᾧ ἦσαν συνηγμένοι, καὶ ἐπλήσθησαν ἅπαντες τοῦ Ἁγίου Πνεύματος καὶ ἐλάλουν τὸν λόγον τοῦ Θεοῦ μετὰ παρρησίας.
\vs{32}Τοῦ δὲ πλήθους τῶν πιστευσάντων ἦν καρδία καὶ ψυχὴ μία, καὶ οὐδὲ εἷς τι τῶν ὑπαρχόντων αὐτῷ ἔλεγεν ἴδιον εἶναι ἀλλ᾽ ἦν αὐτοῖς πάντα κοινά.
\vs{33}καὶ δυνάμει μεγάλῃ ἀπεδίδουν τὸ μαρτύριον οἱ ἀπόστολοι τῆς ἀναστάσεως τοῦ Κυρίου Ἰησοῦ, χάρις τε μεγάλη ἦν ἐπὶ πάντας αὐτούς.
\vs{34}Οὐδὲ γὰρ ἐνδεής τις ἦν ἐν αὐτοῖς· ὅσοι γὰρ κτήτορες χωρίων ἢ οἰκιῶν ὑπῆρχον, πωλοῦντες ἔφερον τὰς τιμὰς τῶν πιπρασκομένων
\vs{35}καὶ ἐτίθουν παρὰ τοὺς πόδας τῶν ἀποστόλων, διεδίδετο δὲ ἑκάστῳ καθότι ἄν τις χρείαν εἶχεν.
\vs{36}Ἰωσὴφ δὲ ὁ ἐπικληθεὶς Βαρνάβας ἀπὸ τῶν ἀποστόλων, ὅ ἐστιν μεθερμηνευόμενον Υἱὸς παρακλήσεως, Λευίτης, Κύπριος τῷ γένει,
\vs{37}ὑπάρχοντος αὐτῷ ἀγροῦ πωλήσας ἤνεγκεν τὸ χρῆμα καὶ ἔθηκεν πρὸς τοὺς πόδας τῶν ἀποστόλων.

\ch{5}
Ἀνὴρ δέ τις Ἁνανίας ὀνόματι σὺν Σαπφίρῃ τῇ γυναικὶ αὐτοῦ ἐπώλησεν κτῆμα
\vs{2}καὶ ἐνοσφίσατο ἀπὸ τῆς τιμῆς, συνειδυίης καὶ τῆς γυναικός, καὶ ἐνέγκας μέρος τι παρὰ τοὺς πόδας τῶν ἀποστόλων ἔθηκεν.
\vs{3}Εἶπεν δὲ ὁ Πέτρος· Ἁνανία, διὰ τί ἐπλήρωσεν ὁ Σατανᾶς τὴν καρδίαν σου, ψεύσασθαί σε τὸ Πνεῦμα τὸ Ἅγιον καὶ νοσφίσασθαι ἀπὸ τῆς τιμῆς τοῦ χωρίου;
\vs{4}οὐχὶ μένον σοὶ ἔμενεν καὶ πραθὲν ἐν τῇ σῇ ἐξουσίᾳ ὑπῆρχεν; τί ὅτι ἔθου ἐν τῇ καρδίᾳ σου τὸ πρᾶγμα τοῦτο; οὐκ ἐψεύσω ἀνθρώποις ἀλλὰ τῷ Θεῷ.
\vs{5}Ἀκούων δὲ ὁ Ἁνανίας τοὺς λόγους τούτους πεσὼν ἐξέψυξεν, καὶ ἐγένετο φόβος μέγας ἐπὶ πάντας τοὺς ἀκούοντας.
\vs{6}ἀναστάντες δὲ οἱ νεώτεροι συνέστειλαν αὐτὸν καὶ ἐξενέγκαντες ἔθαψαν.

\vs{7}Ἐγένετο δὲ ὡς ὡρῶν τριῶν διάστημα καὶ ἡ γυνὴ αὐτοῦ μὴ εἰδυῖα τὸ γεγονὸς εἰσῆλθεν.
\vs{8}ἀπεκρίθη δὲ πρὸς αὐτὴν Πέτρος· Εἰπέ μοι, εἰ τοσούτου τὸ χωρίον ἀπέδοσθε; Ἡ δὲ εἶπεν· Ναί, τοσούτου.
\vs{9}Ὁ δὲ Πέτρος πρὸς αὐτήν· Τί ὅτι συνεφωνήθη ὑμῖν πειράσαι τὸ Πνεῦμα Κυρίου; ἰδοὺ οἱ πόδες τῶν θαψάντων τὸν ἄνδρα σου ἐπὶ τῇ θύρᾳ καὶ ἐξοίσουσίν σε.
\vs{10}Ἔπεσεν δὲ παραχρῆμα πρὸς τοὺς πόδας αὐτοῦ καὶ ἐξέψυξεν· εἰσελθόντες δὲ οἱ νεανίσκοι εὗρον αὐτὴν νεκράν καὶ ἐξενέγκαντες ἔθαψαν πρὸς τὸν ἄνδρα αὐτῆς,
\vs{11}Καὶ ἐγένετο φόβος μέγας ἐφ᾽ ὅλην τὴν ἐκκλησίαν καὶ ἐπὶ πάντας τοὺς ἀκούοντας ταῦτα.

\vs{12}Διὰ δὲ τῶν χειρῶν τῶν ἀποστόλων ἐγίνετο σημεῖα καὶ τέρατα πολλὰ ἐν τῷ λαῷ. καὶ ἦσαν ὁμοθυμαδὸν ἅπαντες ἐν τῇ στοᾷ Σολομῶντος,
\vs{13}τῶν δὲ λοιπῶν οὐδεὶς ἐτόλμα κολλᾶσθαι αὐτοῖς, ἀλλ᾽ ἐμεγάλυνεν αὐτοὺς ὁ λαός.
\vs{14}μᾶλλον δὲ προσετίθεντο πιστεύοντες τῷ Κυρίῳ, πλήθη ἀνδρῶν τε καὶ γυναικῶν,
\vs{15}ὥστε καὶ εἰς τὰς πλατείας ἐκφέρειν τοὺς ἀσθενεῖς καὶ τιθέναι ἐπὶ κλιναρίων καὶ κραβάττων, ἵνα ἐρχομένου Πέτρου κἂν ἡ σκιὰ ἐπισκιάσῃ τινὶ αὐτῶν.
\vs{16}συνήρχετο δὲ καὶ τὸ πλῆθος τῶν πέριξ πόλεων Ἰερουσαλήμ φέροντες ἀσθενεῖς καὶ ὀχλουμένους ὑπὸ πνευμάτων ἀκαθάρτων, οἵτινες ἐθεραπεύοντο ἅπαντες.

\vs{17}Ἀναστὰς δὲ ὁ ἀρχιερεὺς καὶ πάντες οἱ σὺν αὐτῷ, ἡ οὖσα αἵρεσις τῶν Σαδδουκαίων, ἐπλήσθησαν ζήλου
\vs{18}καὶ ἐπέβαλον τὰς χεῖρας ἐπὶ τοὺς ἀποστόλους καὶ ἔθεντο αὐτοὺς ἐν τηρήσει δημοσίᾳ.
\vs{19}Ἄγγελος δὲ Κυρίου διὰ νυκτὸς ἤνοιξε τὰς θύρας τῆς φυλακῆς ἐξαγαγών τε αὐτοὺς εἶπεν·
\vs{20}Πορεύεσθε καὶ σταθέντες λαλεῖτε ἐν τῷ ἱερῷ τῷ λαῷ πάντα τὰ ῥήματα τῆς Ζωῆς ταύτης.
\vs{21}Ἀκούσαντες δὲ εἰσῆλθον ὑπὸ τὸν ὄρθρον εἰς τὸ ἱερὸν καὶ ἐδίδασκον. Παραγενόμενος δὲ ὁ ἀρχιερεὺς καὶ οἱ σὺν αὐτῷ συνεκάλεσαν τὸ συνέδριον καὶ πᾶσαν τὴν γερουσίαν τῶν υἱῶν Ἰσραήλ καὶ ἀπέστειλαν εἰς τὸ δεσμωτήριον ἀχθῆναι αὐτούς.
\vs{22}οἱ δὲ παραγενόμενοι ὑπηρέται οὐχ εὗρον αὐτοὺς ἐν τῇ φυλακῇ· ἀναστρέψαντες δὲ ἀπήγγειλαν
\vs{23}λέγοντες ὅτι Τὸ δεσμωτήριον εὕρομεν κεκλεισμένον ἐν πάσῃ ἀσφαλείᾳ καὶ τοὺς φύλακας ἑστῶτας ἐπὶ τῶν θυρῶν, ἀνοίξαντες δὲ ἔσω οὐδένα εὕρομεν.
\vs{24}Ὡς δὲ ἤκουσαν τοὺς λόγους τούτους ὅ τε στρατηγὸς τοῦ ἱεροῦ καὶ οἱ ἀρχιερεῖς, διηπόρουν περὶ αὐτῶν τί ἂν γένοιτο τοῦτο.
\vs{25}παραγενόμενος δέ τις ἀπήγγειλεν αὐτοῖς ὅτι Ἰδοὺ οἱ ἄνδρες οὓς ἔθεσθε ἐν τῇ φυλακῇ εἰσὶν ἐν τῷ ἱερῷ ἑστῶτες καὶ διδάσκοντες τὸν λαόν.
\vs{26}Τότε ἀπελθὼν ὁ στρατηγὸς σὺν τοῖς ὑπηρέταις ἦγεν αὐτούς οὐ μετὰ βίας, ἐφοβοῦντο γὰρ τὸν λαόν μὴ λιθασθῶσιν.

\vs{27}ἀγαγόντες δὲ αὐτοὺς ἔστησαν ἐν τῷ συνεδρίῳ. καὶ ἐπηρώτησεν αὐτοὺς ὁ ἀρχιερεὺς
\vs{28}λέγων· Οὐ Παραγγελίᾳ παρηγγείλαμεν ὑμῖν μὴ διδάσκειν ἐπὶ τῷ ὀνόματι τούτῳ, καὶ ἰδοὺ πεπληρώκατε τὴν Ἰερουσαλὴμ τῆς διδαχῆς ὑμῶν καὶ βούλεσθε ἐπαγαγεῖν ἐφ᾽ ἡμᾶς τὸ αἷμα τοῦ ἀνθρώπου τούτου.
\vs{29}Ἀποκριθεὶς δὲ Πέτρος καὶ οἱ ἀπόστολοι εἶπαν· Πειθαρχεῖν δεῖ Θεῷ μᾶλλον ἢ ἀνθρώποις.
\vs{30}ὁ Θεὸς τῶν πατέρων ἡμῶν ἤγειρεν Ἰησοῦν ὃν ὑμεῖς διεχειρίσασθε κρεμάσαντες ἐπὶ ξύλου·
\vs{31}τοῦτον ὁ Θεὸς Ἀρχηγὸν καὶ Σωτῆρα ὕψωσεν τῇ δεξιᾷ αὐτοῦ τοῦ δοῦναι μετάνοιαν τῷ Ἰσραὴλ καὶ ἄφεσιν ἁμαρτιῶν.
\vs{32}καὶ ἡμεῖς ἐσμεν μάρτυρες τῶν ῥημάτων τούτων καὶ τὸ Πνεῦμα τὸ Ἅγιον ὃ ἔδωκεν ὁ Θεὸς τοῖς πειθαρχοῦσιν αὐτῷ.

\vs{33}Οἱ δὲ ἀκούσαντες διεπρίοντο καὶ ἐβούλοντο ἀνελεῖν αὐτούς.
\vs{34}Ἀναστὰς δέ τις ἐν τῷ συνεδρίῳ Φαρισαῖος ὀνόματι Γαμαλιήλ, νομοδιδάσκαλος τίμιος παντὶ τῷ λαῷ, ἐκέλευσεν ἔξω βραχὺ τοὺς ἀνθρώπους ποιῆσαι
\vs{35}Εἶπέν τε πρὸς αὐτούς· Ἄνδρες Ἰσραηλῖται, προσέχετε ἑαυτοῖς ἐπὶ τοῖς ἀνθρώποις τούτοις τί μέλλετε πράσσειν.
\vs{36}πρὸ γὰρ τούτων τῶν ἡμερῶν ἀνέστη Θευδᾶς λέγων εἶναί τινα ἑαυτόν, ᾧ προσεκλίθη ἀνδρῶν ἀριθμὸς ὡς τετρακοσίων· ὃς ἀνῃρέθη, καὶ πάντες ὅσοι ἐπείθοντο αὐτῷ διελύθησαν καὶ ἐγένοντο εἰς οὐδέν.
\vs{37}μετὰ τοῦτον ἀνέστη Ἰούδας ὁ Γαλιλαῖος ἐν ταῖς ἡμέραις τῆς ἀπογραφῆς καὶ ἀπέστησεν λαὸν ὀπίσω αὐτοῦ· κἀκεῖνος ἀπώλετο καὶ πάντες ὅσοι ἐπείθοντο αὐτῷ διεσκορπίσθησαν.
\vs{38}Καὶ τὰ νῦν λέγω ὑμῖν, ἀπόστητε ἀπὸ τῶν ἀνθρώπων τούτων καὶ ἄφετε αὐτούς· ὅτι ἐὰν ᾖ ἐξ ἀνθρώπων ἡ βουλὴ αὕτη ἢ τὸ ἔργον τοῦτο, καταλυθήσεται,
\vs{39}εἰ δὲ ἐκ Θεοῦ ἐστιν, οὐ δυνήσεσθε καταλῦσαι αὐτούς, μήποτε καὶ θεομάχοι εὑρεθῆτε. Ἐπείσθησαν δὲ αὐτῷ
\vs{40}καὶ προσκαλεσάμενοι τοὺς ἀποστόλους δείραντες παρήγγειλαν μὴ λαλεῖν ἐπὶ τῷ ὀνόματι τοῦ Ἰησοῦ καὶ ἀπέλυσαν.
\vs{41}Οἱ μὲν οὖν ἐπορεύοντο χαίροντες ἀπὸ προσώπου τοῦ συνεδρίου, ὅτι κατηξιώθησαν ὑπὲρ τοῦ Ὀνόματος ἀτιμασθῆναι,
\vs{42}πᾶσάν τε ἡμέραν ἐν τῷ ἱερῷ καὶ κατ᾽ οἶκον οὐκ ἐπαύοντο διδάσκοντες καὶ εὐαγγελιζόμενοι τὸν Χριστὸν Ἰησοῦν.

\ch{6}
Ἐν δὲ ταῖς ἡμέραις ταύταις πληθυνόντων τῶν μαθητῶν ἐγένετο γογγυσμὸς τῶν Ἑλληνιστῶν πρὸς τοὺς Ἑβραίους, ὅτι παρεθεωροῦντο ἐν τῇ διακονίᾳ τῇ καθημερινῇ αἱ χῆραι αὐτῶν.
\vs{2}Προσκαλεσάμενοι δὲ οἱ δώδεκα τὸ πλῆθος τῶν μαθητῶν εἶπαν· Οὐκ ἀρεστόν ἐστιν ἡμᾶς καταλείψαντας τὸν λόγον τοῦ Θεοῦ διακονεῖν τραπέζαις.
\vs{3}ἐπισκέψασθε δέ, ἀδελφοί, ἄνδρας ἐξ ὑμῶν μαρτυρουμένους ἑπτὰ, πλήρεις Πνεύματος καὶ σοφίας, οὓς καταστήσομεν ἐπὶ τῆς χρείας ταύτης,
\vs{4}ἡμεῖς δὲ τῇ προσευχῇ καὶ τῇ διακονίᾳ τοῦ λόγου προσκαρτερήσομεν.
\vs{5}Καὶ ἤρεσεν ὁ λόγος ἐνώπιον παντὸς τοῦ πλήθους καὶ ἐξελέξαντο Στέφανον, ἄνδρα πλήρης πίστεως καὶ Πνεύματος Ἁγίου, καὶ Φίλιππον καὶ Πρόχορον καὶ Νικάνορα καὶ Τίμωνα καὶ Παρμενᾶν καὶ Νικόλαον προσήλυτον Ἀντιοχέα,
\vs{6}οὓς ἔστησαν ἐνώπιον τῶν ἀποστόλων, καὶ προσευξάμενοι ἐπέθηκαν αὐτοῖς τὰς χεῖρας.
\vs{7}Καὶ ὁ λόγος τοῦ Θεοῦ ηὔξανεν καὶ ἐπληθύνετο ὁ ἀριθμὸς τῶν μαθητῶν ἐν Ἰερουσαλὴμ σφόδρα, πολύς τε ὄχλος τῶν ἱερέων ὑπήκουον τῇ πίστει.

\vs{8}Στέφανος δὲ πλήρης χάριτος καὶ δυνάμεως ἐποίει τέρατα καὶ σημεῖα μεγάλα ἐν τῷ λαῷ.
\vs{9}ἀνέστησαν δέ τινες τῶν ἐκ τῆς συναγωγῆς τῆς λεγομένης Λιβερτίνων καὶ Κυρηναίων καὶ Ἀλεξανδρέων καὶ τῶν ἀπὸ Κιλικίας καὶ Ἀσίας συζητοῦντες τῷ Στεφάνῳ,
\vs{10}καὶ οὐκ ἴσχυον ἀντιστῆναι τῇ σοφίᾳ καὶ τῷ Πνεύματι ᾧ ἐλάλει.
\vs{11}Τότε ὑπέβαλον ἄνδρας λέγοντας ὅτι Ἀκηκόαμεν αὐτοῦ λαλοῦντος ῥήματα βλάσφημα εἰς Μωϋσῆν καὶ τὸν Θεόν.
\vs{12}συνεκίνησάν τε τὸν λαὸν καὶ τοὺς πρεσβυτέρους καὶ τοὺς γραμματεῖς καὶ ἐπιστάντες συνήρπασαν αὐτὸν καὶ ἤγαγον εἰς τὸ συνέδριον,
\vs{13}Ἔστησάν τε μάρτυρας ψευδεῖς λέγοντας· Ὁ ἄνθρωπος οὗτος οὐ παύεται λαλῶν ῥήματα κατὰ τοῦ τόπου τοῦ ἁγίου τούτου καὶ τοῦ νόμου·
\vs{14}ἀκηκόαμεν γὰρ αὐτοῦ λέγοντος ὅτι Ἰησοῦς ὁ Ναζωραῖος οὗτος καταλύσει τὸν τόπον τοῦτον καὶ ἀλλάξει τὰ ἔθη ἃ παρέδωκεν ἡμῖν Μωϋσῆς.
\vs{15}Καὶ ἀτενίσαντες εἰς αὐτὸν πάντες οἱ καθεζόμενοι ἐν τῷ συνεδρίῳ εἶδον τὸ πρόσωπον αὐτοῦ ὡσεὶ πρόσωπον ἀγγέλου.

\ch{7}
Εἶπεν δὲ ὁ ἀρχιερεύς· Εἰ ταῦτα οὕτως ἔχει;
\vs{2}Ὁ δὲ ἔφη· Ἄνδρες ἀδελφοὶ καὶ πατέρες, ἀκούσατε. Ὁ Θεὸς τῆς δόξης ὤφθη τῷ πατρὶ ἡμῶν Ἀβραὰμ ὄντι ἐν τῇ Μεσοποταμίᾳ πρὶν ἢ κατοικῆσαι αὐτὸν ἐν Χαρράν
\vs{3}καὶ εἶπεν πρὸς αὐτόν· Ἔξελθε ἐκ τῆς γῆς σου καὶ ἐκ τῆς συγγενείας σου, καὶ δεῦρο εἰς τὴν γῆν ἣν ἄν σοι δείξω.
\vs{4}τότε ἐξελθὼν ἐκ γῆς Χαλδαίων κατῴκησεν ἐν Χαρράν. κἀκεῖθεν μετὰ τὸ ἀποθανεῖν τὸν πατέρα αὐτοῦ μετῴκισεν αὐτὸν εἰς τὴν γῆν ταύτην εἰς ἣν ὑμεῖς νῦν κατοικεῖτε,
\vs{5}Καὶ οὐκ ἔδωκεν αὐτῷ κληρονομίαν ἐν αὐτῇ οὐδὲ βῆμα ποδός καὶ ἐπηγγείλατο δοῦναι αὐτῷ εἰς κατάσχεσιν αὐτὴν καὶ τῷ σπέρματι αὐτοῦ μετ᾽ αὐτόν, οὐκ ὄντος αὐτῷ τέκνου.
\vs{6}ἐλάλησεν δὲ οὕτως ὁ Θεὸς ὅτι ἔσται τὸ σπέρμα αὐτοῦ πάροικον ἐν γῇ ἀλλοτρίᾳ καὶ δουλώσουσιν αὐτὸ καὶ κακώσουσιν ἔτη τετρακόσια·
\vs{7}Καὶ τὸ ἔθνος ᾧ ἐὰν δουλεύσουσιν κρινῶ ἐγώ, ὁ Θεὸς εἶπεν, Καὶ μετὰ ταῦτα ἐξελεύσονται καὶ λατρεύσουσίν μοι ἐν τῷ τόπῳ τούτῳ.
\vs{8}Καὶ ἔδωκεν αὐτῷ διαθήκην περιτομῆς· καὶ οὕτως ἐγέννησεν τὸν Ἰσαὰκ καὶ περιέτεμεν αὐτὸν τῇ ἡμέρᾳ τῇ ὀγδόῃ, καὶ Ἰσαὰκ τὸν Ἰακώβ, καὶ Ἰακὼβ τοὺς δώδεκα πατριάρχας.

\vs{9}Καὶ οἱ πατριάρχαι ζηλώσαντες τὸν Ἰωσὴφ ἀπέδοντο εἰς Αἴγυπτον. καὶ ἦν ὁ Θεὸς μετ᾽ αὐτοῦ
\vs{10}καὶ ἐξείλατο αὐτὸν ἐκ πασῶν τῶν θλίψεων αὐτοῦ καὶ ἔδωκεν αὐτῷ χάριν καὶ σοφίαν ἐναντίον Φαραὼ βασιλέως Αἰγύπτου καὶ κατέστησεν αὐτὸν ἡγούμενον ἐπ᾽ Αἴγυπτον καὶ ἐφ᾽ ὅλον τὸν οἶκον αὐτοῦ.
\vs{11}Ἦλθεν δὲ λιμὸς ἐφ᾽ ὅλην τὴν Αἴγυπτον καὶ Χανάαν καὶ θλῖψις μεγάλη, καὶ οὐχ ηὕρισκον χορτάσματα οἱ πατέρες ἡμῶν.
\vs{12}ἀκούσας δὲ Ἰακὼβ ὄντα σιτία εἰς Αἴγυπτον ἐξαπέστειλεν τοὺς πατέρας ἡμῶν πρῶτον.
\vs{13}καὶ ἐν τῷ δευτέρῳ ἀνεγνωρίσθη Ἰωσὴφ τοῖς ἀδελφοῖς αὐτοῦ καὶ φανερὸν ἐγένετο τῷ Φαραὼ τὸ γένος τοῦ Ἰωσήφ.
\vs{14}ἀποστείλας δὲ Ἰωσὴφ μετεκαλέσατο Ἰακὼβ τὸν πατέρα αὐτοῦ καὶ πᾶσαν τὴν συγγένειαν ἐν ψυχαῖς ἑβδομήκοντα πέντε.
\vs{15}Καὶ κατέβη Ἰακὼβ εἰς Αἴγυπτον καὶ ἐτελεύτησεν αὐτὸς καὶ οἱ πατέρες ἡμῶν,
\vs{16}καὶ μετετέθησαν εἰς Συχὲμ καὶ ἐτέθησαν ἐν τῷ μνήματι ᾧ ὠνήσατο Ἀβραὰμ τιμῆς ἀργυρίου παρὰ τῶν υἱῶν Ἑμμὼρ ἐν Συχέμ.
\vs{17}Καθὼς δὲ ἤγγιζεν ὁ χρόνος τῆς ἐπαγγελίας ἧς ὡμολόγησεν ὁ Θεὸς τῷ Ἀβραάμ, ηὔξησεν ὁ λαὸς καὶ ἐπληθύνθη ἐν Αἰγύπτῳ
\vs{18}ἄχρι οὗ ἀνέστη βασιλεὺς ἕτερος ἐπ᾽ Αἴγυπτον ὃς οὐκ ᾔδει τὸν Ἰωσήφ.
\vs{19}οὗτος κατασοφισάμενος τὸ γένος ἡμῶν ἐκάκωσεν τοὺς πατέρας ἡμῶν τοῦ ποιεῖν τὰ βρέφη ἔκθετα αὐτῶν εἰς τὸ μὴ ζωογονεῖσθαι.
\vs{20}Ἐν ᾧ καιρῷ ἐγεννήθη Μωϋσῆς καὶ ἦν ἀστεῖος τῷ Θεῷ· ὃς ἀνετράφη μῆνας τρεῖς ἐν τῷ οἴκῳ τοῦ πατρός,
\vs{21}ἐκτεθέντος δὲ αὐτοῦ ἀνείλατο αὐτὸν ἡ θυγάτηρ Φαραὼ καὶ ἀνεθρέψατο αὐτὸν ἑαυτῇ εἰς υἱόν.
\vs{22}καὶ ἐπαιδεύθη Μωϋσῆς ἐν πάσῃ σοφίᾳ Αἰγυπτίων, ἦν δὲ δυνατὸς ἐν λόγοις καὶ ἔργοις αὐτοῦ.
\vs{23}Ὡς δὲ ἐπληροῦτο αὐτῷ τεσσερακονταέτης χρόνος, ἀνέβη ἐπὶ τὴν καρδίαν αὐτοῦ ἐπισκέψασθαι τοὺς ἀδελφοὺς αὐτοῦ τοὺς υἱοὺς Ἰσραήλ.
\vs{24}καὶ ἰδών τινα ἀδικούμενον ἠμύνατο καὶ ἐποίησεν ἐκδίκησιν τῷ καταπονουμένῳ πατάξας τὸν Αἰγύπτιον.
\vs{25}ἐνόμιζεν δὲ συνιέναι τοὺς ἀδελφοὺς αὐτοῦ ὅτι ὁ Θεὸς διὰ χειρὸς αὐτοῦ δίδωσιν σωτηρίαν αὐτοῖς· οἱ δὲ οὐ συνῆκαν.
\vs{26}Τῇ τε ἐπιούσῃ ἡμέρᾳ ὤφθη αὐτοῖς μαχομένοις καὶ συνήλλασσεν αὐτοὺς εἰς εἰρήνην εἰπών· Ἄνδρες, ἀδελφοί ἐστε· ἱνατί ἀδικεῖτε ἀλλήλους;
\vs{27}Ὁ δὲ ἀδικῶν τὸν πλησίον ἀπώσατο αὐτὸν εἰπών· Τίς σε κατέστησεν ἄρχοντα καὶ δικαστὴν ἐφ᾽ ἡμῶν;
\vs{28}μὴ ἀνελεῖν με σὺ θέλεις ὃν τρόπον ἀνεῖλες ἐχθὲς τὸν Αἰγύπτιον;
\vs{29}ἔφυγεν δὲ Μωϋσῆς ἐν τῷ λόγῳ τούτῳ καὶ ἐγένετο πάροικος ἐν γῇ Μαδιάμ, οὗ ἐγέννησεν υἱοὺς δύο.
\vs{30}Καὶ πληρωθέντων ἐτῶν τεσσεράκοντα ὤφθη αὐτῷ ἐν τῇ ἐρήμῳ τοῦ ὄρους Σινᾶ ἄγγελος ἐν φλογὶ πυρὸς βάτου.
\vs{31}ὁ δὲ Μωϋσῆς ἰδὼν ἐθαύμαζεν τὸ ὅραμα, προσερχομένου δὲ αὐτοῦ κατανοῆσαι ἐγένετο φωνὴ Κυρίου·
\vs{32}Ἐγὼ ὁ Θεὸς τῶν πατέρων σου, ὁ Θεὸς Ἀβραὰμ καὶ Ἰσαὰκ καὶ Ἰακώβ. ἔντρομος δὲ γενόμενος Μωϋσῆς οὐκ ἐτόλμα κατανοῆσαι.
\vs{33}Εἶπεν δὲ αὐτῷ ὁ Κύριος· Λῦσον τὸ ὑπόδημα τῶν ποδῶν σου, ὁ γὰρ τόπος ἐφ᾽ ᾧ ἕστηκας γῆ ἁγία ἐστίν.
\vs{34}ἰδὼν εἶδον τὴν κάκωσιν τοῦ λαοῦ μου τοῦ ἐν Αἰγύπτῳ καὶ τοῦ στεναγμοῦ αὐτῶν ἤκουσα, καὶ κατέβην ἐξελέσθαι αὐτούς· καὶ νῦν δεῦρο ἀποστείλω σε εἰς Αἴγυπτον.
\vs{35}Τοῦτον τὸν Μωϋσῆν ὃν ἠρνήσαντο εἰπόντες· Τίς σε κατέστησεν ἄρχοντα καὶ δικαστήν; τοῦτον ὁ Θεὸς καὶ ἄρχοντα καὶ λυτρωτὴν ἀπέσταλκεν σὺν χειρὶ ἀγγέλου τοῦ ὀφθέντος αὐτῷ ἐν τῇ βάτῳ.
\vs{36}οὗτος ἐξήγαγεν αὐτοὺς ποιήσας τέρατα καὶ σημεῖα ἐν γῇ Αἰγύπτῳ καὶ ἐν Ἐρυθρᾷ Θαλάσσῃ καὶ ἐν τῇ ἐρήμῳ ἔτη τεσσεράκοντα.
\vs{37}Οὗτός ἐστιν ὁ Μωϋσῆς ὁ εἴπας τοῖς υἱοῖς Ἰσραήλ· Προφήτην ὑμῖν ἀναστήσει ὁ Θεὸς ἐκ τῶν ἀδελφῶν ὑμῶν ὡς ἐμέ.
\vs{38}οὗτός ἐστιν ὁ γενόμενος ἐν τῇ ἐκκλησίᾳ ἐν τῇ ἐρήμῳ μετὰ τοῦ ἀγγέλου τοῦ λαλοῦντος αὐτῷ ἐν τῷ ὄρει Σινᾶ καὶ τῶν πατέρων ἡμῶν, ὃς ἐδέξατο λόγια ζῶντα δοῦναι ἡμῖν,
\vs{39}ᾧ οὐκ ἠθέλησαν ὑπήκοοι γενέσθαι οἱ πατέρες ἡμῶν, ἀλλὰ ἀπώσαντο καὶ ἐστράφησαν ἐν ταῖς καρδίαις αὐτῶν εἰς Αἴγυπτον
\vs{40}εἰπόντες τῷ Ἀαρών· Ποίησον ἡμῖν θεοὺς οἳ προπορεύσονται ἡμῶν· ὁ γὰρ Μωϋσῆς οὗτος, ὃς ἐξήγαγεν ἡμᾶς ἐκ γῆς Αἰγύπτου, οὐκ οἴδαμεν τί ἐγένετο αὐτῷ.
\vs{41}Καὶ ἐμοσχοποίησαν ἐν ταῖς ἡμέραις ἐκείναις καὶ ἀνήγαγον θυσίαν τῷ εἰδώλῳ καὶ εὐφραίνοντο ἐν τοῖς ἔργοις τῶν χειρῶν αὐτῶν.
\vs{42}ἔστρεψεν δὲ ὁ Θεὸς καὶ παρέδωκεν αὐτοὺς λατρεύειν τῇ στρατιᾷ τοῦ οὐρανοῦ καθὼς γέγραπται ἐν βίβλῳ τῶν προφητῶν· 
\begin{poetryblock}

\begin{quote}Μὴ σφάγια καὶ θυσίας προσηνέγκατέ μοι\end{quote} 

\begin{quote}ἔτη τεσσεράκοντα ἐν τῇ ἐρήμῳ, οἶκος Ἰσραήλ;\end{quote}

\begin{quote} \vs{43}καὶ ἀνελάβετε τὴν σκηνὴν τοῦ Μολὸχ\end{quote} 

\begin{quote}καὶ τὸ ἄστρον τοῦ θεοῦ ὑμῶν Ῥαιφάν,\end{quote} 

\begin{quote}τοὺς τύπους οὓς ἐποιήσατε προσκυνεῖν αὐτοῖς,\end{quote} 

\begin{quote}καὶ μετοικιῶ ὑμᾶς ἐπέκεινα Βαβυλῶνος.\end{quote}
\end{poetryblock}

\vs{44}Ἡ σκηνὴ τοῦ μαρτυρίου ἦν τοῖς πατράσιν ἡμῶν ἐν τῇ ἐρήμῳ καθὼς διετάξατο ὁ λαλῶν τῷ Μωϋσῇ ποιῆσαι αὐτὴν κατὰ τὸν τύπον ὃν ἑωράκει·
\vs{45}ἣν καὶ εἰσήγαγον διαδεξάμενοι οἱ πατέρες ἡμῶν μετὰ Ἰησοῦ ἐν τῇ κατασχέσει τῶν ἐθνῶν, ὧν ἐξῶσεν ὁ Θεὸς ἀπὸ προσώπου τῶν πατέρων ἡμῶν ἕως τῶν ἡμερῶν Δαυίδ,
\vs{46}ὃς εὗρεν χάριν ἐνώπιον τοῦ Θεοῦ καὶ ᾐτήσατο εὑρεῖν σκήνωμα τῷ οἴκῳ Ἰακώβ.
\vs{47}Σολομῶν δὲ οἰκοδόμησεν αὐτῷ οἶκον.
\vs{48}Ἀλλ᾽ οὐχ ὁ Ὕψιστος ἐν χειροποιήτοις κατοικεῖ, καθὼς ὁ προφήτης λέγει·
\begin{poetryblock}

\begin{quote} \vs{49}Ὁ οὐρανός μοι θρόνος,\end{quote} 

\begin{quote}ἡ δὲ γῆ ὑποπόδιον τῶν ποδῶν μου·\end{quote} 

\begin{quote}ποῖον οἶκον οἰκοδομήσετέ μοι, λέγει Κύριος,\end{quote} 

\begin{quote}ἢ τίς τόπος τῆς καταπαύσεώς μου;\end{quote}

\begin{quote} \vs{50}οὐχὶ ἡ χείρ μου ἐποίησεν ταῦτα πάντα;\end{quote}
\end{poetryblock}

\vs{51}Σκληροτράχηλοι καὶ ἀπερίτμητοι καρδίαις καὶ τοῖς ὠσίν, ὑμεῖς ἀεὶ τῷ Πνεύματι τῷ Ἁγίῳ ἀντιπίπτετε ὡς οἱ πατέρες ὑμῶν καὶ ὑμεῖς.
\vs{52}τίνα τῶν προφητῶν οὐκ ἐδίωξαν οἱ πατέρες ὑμῶν; καὶ ἀπέκτειναν τοὺς προκαταγγείλαντας περὶ τῆς ἐλεύσεως τοῦ Δικαίου, οὗ νῦν ὑμεῖς προδόται καὶ φονεῖς ἐγένεσθε,
\vs{53}οἵτινες ἐλάβετε τὸν νόμον εἰς διαταγὰς ἀγγέλων καὶ οὐκ ἐφυλάξατε.

\vs{54}Ἀκούοντες δὲ ταῦτα διεπρίοντο ταῖς καρδίαις αὐτῶν καὶ ἔβρυχον τοὺς ὀδόντας ἐπ᾽ αὐτόν.
\vs{55}ὑπάρχων δὲ πλήρης Πνεύματος Ἁγίου ἀτενίσας εἰς τὸν οὐρανὸν εἶδεν δόξαν Θεοῦ καὶ Ἰησοῦν ἑστῶτα ἐκ δεξιῶν τοῦ Θεοῦ
\vs{56}καὶ εἶπεν· Ἰδοὺ θεωρῶ τοὺς οὐρανοὺς διηνοιγμένους καὶ τὸν Υἱὸν τοῦ ἀνθρώπου ἐκ δεξιῶν ἑστῶτα τοῦ Θεοῦ.
\vs{57}Κράξαντες δὲ φωνῇ μεγάλῃ συνέσχον τὰ ὦτα αὐτῶν καὶ ὥρμησαν ὁμοθυμαδὸν ἐπ᾽ αὐτόν
\vs{58}καὶ ἐκβαλόντες ἔξω τῆς πόλεως ἐλιθοβόλουν. καὶ οἱ μάρτυρες ἀπέθεντο τὰ ἱμάτια αὐτῶν παρὰ τοὺς πόδας νεανίου καλουμένου Σαύλου,
\vs{59}Καὶ ἐλιθοβόλουν τὸν Στέφανον ἐπικαλούμενον καὶ λέγοντα· Κύριε Ἰησοῦ, δέξαι τὸ πνεῦμά μου.
\vs{60}θεὶς δὲ τὰ γόνατα ἔκραξεν φωνῇ μεγάλῃ· Κύριε, μὴ στήσῃς αὐτοῖς ταύτην τὴν ἁμαρτίαν. καὶ τοῦτο εἰπὼν ἐκοιμήθη.

\ch{8}
Σαῦλος δὲ ἦν συνευδοκῶν τῇ ἀναιρέσει αὐτοῦ.

Ἐγένετο δὲ ἐν ἐκείνῃ τῇ ἡμέρᾳ διωγμὸς μέγας ἐπὶ τὴν ἐκκλησίαν τὴν ἐν Ἱεροσολύμοις, πάντες δὲ διεσπάρησαν κατὰ τὰς χώρας τῆς Ἰουδαίας καὶ Σαμαρείας πλὴν τῶν ἀποστόλων.
\vs{2}συνεκόμισαν δὲ τὸν Στέφανον ἄνδρες εὐλαβεῖς καὶ ἐποίησαν κοπετὸν μέγαν ἐπ᾽ αὐτῷ.
\vs{3}Σαῦλος δὲ ἐλυμαίνετο τὴν ἐκκλησίαν κατὰ τοὺς οἴκους εἰσπορευόμενος, σύρων τε ἄνδρας καὶ γυναῖκας παρεδίδου εἰς φυλακήν.

\vs{4}Οἱ μὲν οὖν διασπαρέντες διῆλθον εὐαγγελιζόμενοι τὸν λόγον.
\vs{5}Φίλιππος δὲ κατελθὼν εἰς τὴν πόλιν τῆς Σαμαρείας ἐκήρυσσεν αὐτοῖς τὸν Χριστόν.
\vs{6}προσεῖχον δὲ οἱ ὄχλοι τοῖς λεγομένοις ὑπὸ τοῦ Φιλίππου ὁμοθυμαδὸν ἐν τῷ ἀκούειν αὐτοὺς καὶ βλέπειν τὰ σημεῖα ἃ ἐποίει.
\vs{7}πολλοὶ γὰρ τῶν ἐχόντων πνεύματα ἀκάθαρτα βοῶντα φωνῇ μεγάλῃ ἐξήρχοντο, πολλοὶ δὲ παραλελυμένοι καὶ χωλοὶ ἐθεραπεύθησαν·
\vs{8}ἐγένετο δὲ πολλὴ χαρὰ ἐν τῇ πόλει ἐκείνῃ.

\vs{9}Ἀνὴρ δέ τις ὀνόματι Σίμων προϋπῆρχεν ἐν τῇ πόλει μαγεύων καὶ ἐξιστάνων τὸ ἔθνος τῆς Σαμαρείας, λέγων εἶναί τινα ἑαυτὸν μέγαν,
\vs{10}ᾧ προσεῖχον πάντες ἀπὸ μικροῦ ἕως μεγάλου λέγοντες· Οὗτός ἐστιν ἡ δύναμις τοῦ Θεοῦ ἡ καλουμένη Μεγάλη.
\vs{11}προσεῖχον δὲ αὐτῷ διὰ τὸ ἱκανῷ χρόνῳ ταῖς μαγείαις ἐξεστακέναι αὐτούς.
\vs{12}Ὅτε δὲ ἐπίστευσαν τῷ Φιλίππῳ εὐαγγελιζομένῳ περὶ τῆς βασιλείας τοῦ Θεοῦ καὶ τοῦ ὀνόματος Ἰησοῦ Χριστοῦ, ἐβαπτίζοντο ἄνδρες τε καὶ γυναῖκες.
\vs{13}ὁ δὲ Σίμων καὶ αὐτὸς ἐπίστευσεν καὶ βαπτισθεὶς ἦν προσκαρτερῶν τῷ Φιλίππῳ, θεωρῶν τε σημεῖα καὶ δυνάμεις μεγάλας γινομένας ἐξίστατο.

\vs{14}Ἀκούσαντες δὲ οἱ ἐν Ἱεροσολύμοις ἀπόστολοι ὅτι δέδεκται ἡ Σαμάρεια τὸν λόγον τοῦ Θεοῦ, ἀπέστειλαν πρὸς αὐτοὺς Πέτρον καὶ Ἰωάννην,
\vs{15}οἵτινες καταβάντες προσηύξαντο περὶ αὐτῶν ὅπως λάβωσιν Πνεῦμα Ἅγιον·
\vs{16}οὐδέπω γὰρ ἦν ἐπ᾽ οὐδενὶ αὐτῶν ἐπιπεπτωκός, μόνον δὲ βεβαπτισμένοι ὑπῆρχον εἰς τὸ ὄνομα τοῦ κυρίου Ἰησοῦ.
\vs{17}τότε ἐπετίθεσαν τὰς χεῖρας ἐπ᾽ αὐτούς καὶ ἐλάμβανον Πνεῦμα Ἅγιον.
\vs{18}Ἰδὼν δὲ ὁ Σίμων ὅτι διὰ τῆς ἐπιθέσεως τῶν χειρῶν τῶν ἀποστόλων δίδοται τὸ Πνεῦμα, προσήνεγκεν αὐτοῖς χρήματα
\vs{19}λέγων· Δότε κἀμοὶ τὴν ἐξουσίαν ταύτην ἵνα ᾧ ἐὰν ἐπιθῶ τὰς χεῖρας λαμβάνῃ Πνεῦμα Ἅγιον.
\vs{20}Πέτρος δὲ εἶπεν πρὸς αὐτόν· Τὸ ἀργύριόν σου σὺν σοὶ εἴη εἰς ἀπώλειαν ὅτι τὴν δωρεὰν τοῦ Θεοῦ ἐνόμισας διὰ χρημάτων κτᾶσθαι·
\vs{21}οὐκ ἔστιν σοι μερὶς οὐδὲ κλῆρος ἐν τῷ λόγῳ τούτῳ, ἡ γὰρ καρδία σου οὐκ ἔστιν εὐθεῖα ἔναντι τοῦ Θεοῦ.
\vs{22}μετανόησον οὖν ἀπὸ τῆς κακίας σου ταύτης καὶ δεήθητι τοῦ Κυρίου, εἰ ἄρα ἀφεθήσεταί σοι ἡ ἐπίνοια τῆς καρδίας σου,
\vs{23}εἰς γὰρ χολὴν πικρίας καὶ σύνδεσμον ἀδικίας ὁρῶ σε ὄντα.
\vs{24}Ἀποκριθεὶς δὲ ὁ Σίμων εἶπεν· Δεήθητε ὑμεῖς ὑπὲρ ἐμοῦ πρὸς τὸν Κύριον ὅπως μηδὲν ἐπέλθῃ ἐπ᾽ ἐμὲ ὧν εἰρήκατε.
\vs{25}Οἱ μὲν οὖν διαμαρτυράμενοι καὶ λαλήσαντες τὸν λόγον τοῦ Κυρίου ὑπέστρεφον εἰς Ἱεροσόλυμα, πολλάς τε κώμας τῶν Σαμαριτῶν εὐηγγελίζοντο.

\vs{26}Ἄγγελος δὲ Κυρίου ἐλάλησεν πρὸς Φίλιππον λέγων· Ἀνάστηθι καὶ πορεύου κατὰ μεσημβρίαν ἐπὶ τὴν ὁδὸν τὴν καταβαίνουσαν ἀπὸ Ἰερουσαλὴμ εἰς Γάζαν, αὕτη ἐστὶν ἔρημος.
\vs{27}καὶ ἀναστὰς ἐπορεύθη. καὶ ἰδοὺ ἀνὴρ Αἰθίοψ εὐνοῦχος δυνάστης Κανδάκης βασιλίσσης Αἰθιόπων, ὃς ἦν ἐπὶ πάσης τῆς γάζης αὐτῆς, ὃς ἐληλύθει προσκυνήσων εἰς Ἰερουσαλήμ,
\vs{28}ἦν τε ὑποστρέφων καὶ καθήμενος ἐπὶ τοῦ ἅρματος αὐτοῦ καὶ ἀνεγίνωσκεν τὸν προφήτην Ἠσαΐαν.
\vs{29}Εἶπεν δὲ τὸ Πνεῦμα τῷ Φιλίππῳ· Πρόσελθε καὶ κολλήθητι τῷ ἅρματι τούτῳ.
\vs{30}Προσδραμὼν δὲ ὁ Φίλιππος ἤκουσεν αὐτοῦ ἀναγινώσκοντος Ἠσαΐαν τὸν προφήτην καὶ εἶπεν· Ἆρά γε γινώσκεις ἃ ἀναγινώσκεις;
\vs{31}Ὁ δὲ εἶπεν· Πῶς γὰρ ἂν δυναίμην ἐὰν μή τις ὁδηγήσει με; παρεκάλεσέν τε τὸν Φίλιππον ἀναβάντα καθίσαι σὺν αὐτῷ.
\vs{32}Ἡ δὲ περιοχὴ τῆς γραφῆς ἣν ἀνεγίνωσκεν ἦν αὕτη· 
\begin{poetryblock}

\begin{quote}Ὡς πρόβατον ἐπὶ σφαγὴν ἤχθη\end{quote} 

\begin{quote}καὶ ὡς ἀμνὸς ἐναντίον τοῦ κείραντος αὐτὸν ἄφωνος,\end{quote} 

\begin{quote}οὕτως οὐκ ἀνοίγει τὸ στόμα αὐτοῦ.\end{quote}

\begin{quote} \vs{33}Ἐν τῇ ταπεινώσει αὐτοῦ ἡ κρίσις αὐτοῦ ἤρθη·\end{quote} 

\begin{quote}τὴν γενεὰν αὐτοῦ τίς διηγήσεται;\end{quote} 

\begin{quote}ὅτι αἴρεται ἀπὸ τῆς γῆς ἡ ζωὴ αὐτοῦ.\end{quote}
\end{poetryblock}

\vs{34}Ἀποκριθεὶς δὲ ὁ εὐνοῦχος τῷ Φιλίππῳ εἶπεν· Δέομαί σου, περὶ τίνος ὁ προφήτης λέγει τοῦτο; περὶ ἑαυτοῦ ἢ περὶ ἑτέρου τινός;
\vs{35}Ἀνοίξας δὲ ὁ Φίλιππος τὸ στόμα αὐτοῦ καὶ ἀρξάμενος ἀπὸ τῆς γραφῆς ταύτης εὐηγγελίσατο αὐτῷ τὸν Ἰησοῦν.
\vs{36}Ὡς δὲ ἐπορεύοντο κατὰ τὴν ὁδόν, ἦλθον ἐπί τι ὕδωρ, καί φησιν ὁ εὐνοῦχος· Ἰδοὺ ὕδωρ, τί κωλύει με βαπτισθῆναι;
\vs{38}καὶ ἐκέλευσεν στῆναι τὸ ἅρμα καὶ κατέβησαν ἀμφότεροι εἰς τὸ ὕδωρ, ὅ τε Φίλιππος καὶ ὁ εὐνοῦχος, καὶ ἐβάπτισεν αὐτόν.
\vs{39}Ὅτε δὲ ἀνέβησαν ἐκ τοῦ ὕδατος, Πνεῦμα Κυρίου ἥρπασεν τὸν Φίλιππον καὶ οὐκ εἶδεν αὐτὸν οὐκέτι ὁ εὐνοῦχος, ἐπορεύετο γὰρ τὴν ὁδὸν αὐτοῦ χαίρων.
\vs{40}Φίλιππος δὲ εὑρέθη εἰς Ἄζωτον· καὶ διερχόμενος εὐηγγελίζετο τὰς πόλεις πάσας ἕως τοῦ ἐλθεῖν αὐτὸν εἰς Καισάρειαν.

\ch{9}
Ὁ Δὲ Σαῦλος ἔτι ἐμπνέων ἀπειλῆς καὶ φόνου εἰς τοὺς μαθητὰς τοῦ Κυρίου, προσελθὼν τῷ ἀρχιερεῖ
\vs{2}ᾐτήσατο παρ᾽ αὐτοῦ ἐπιστολὰς εἰς Δαμασκὸν πρὸς τὰς συναγωγάς, ὅπως ἐάν τινας εὕρῃ τῆς Ὁδοῦ ὄντας, ἄνδρας τε καὶ γυναῖκας, δεδεμένους ἀγάγῃ εἰς Ἰερουσαλήμ.
\vs{3}Ἐν δὲ τῷ πορεύεσθαι ἐγένετο αὐτὸν ἐγγίζειν τῇ Δαμασκῷ, ἐξαίφνης τε αὐτὸν περιήστραψεν φῶς ἐκ τοῦ οὐρανοῦ
\vs{4}καὶ πεσὼν ἐπὶ τὴν γῆν ἤκουσεν φωνὴν λέγουσαν αὐτῷ· Σαοὺλ Σαούλ, τί με διώκεις;
\vs{5}Εἶπεν δέ· Τίς εἶ, Κύριε; Ὁ δέ· Ἐγώ εἰμι Ἰησοῦς ὃν σὺ διώκεις·
\vs{6}ἀλλὰ ἀνάστηθι καὶ εἴσελθε εἰς τὴν πόλιν καὶ λαληθήσεταί σοι ὅ τί σε δεῖ ποιεῖν.
\vs{7}Οἱ δὲ ἄνδρες οἱ συνοδεύοντες αὐτῷ εἱστήκεισαν ἐνεοί, ἀκούοντες μὲν τῆς φωνῆς μηδένα δὲ θεωροῦντες.
\vs{8}ἠγέρθη δὲ Σαῦλος ἀπὸ τῆς γῆς, ἀνεῳγμένων δὲ τῶν ὀφθαλμῶν αὐτοῦ οὐδὲν ἔβλεπεν· χειραγωγοῦντες δὲ αὐτὸν εἰσήγαγον εἰς Δαμασκόν.
\vs{9}καὶ ἦν ἡμέρας τρεῖς μὴ βλέπων καὶ οὐκ ἔφαγεν οὐδὲ ἔπιεν.

\vs{10}Ἦν δέ τις μαθητὴς ἐν Δαμασκῷ ὀνόματι Ἁνανίας, καὶ εἶπεν πρὸς αὐτὸν ἐν ὁράματι ὁ Κύριος· Ἁνανία. Ὁ δὲ εἶπεν· Ἰδοὺ ἐγώ, Κύριε.
\vs{11}Ὁ δὲ Κύριος πρὸς αὐτόν· Ἀναστὰς πορεύθητι ἐπὶ τὴν ῥύμην τὴν καλουμένην Εὐθεῖαν καὶ ζήτησον ἐν οἰκίᾳ Ἰούδα Σαῦλον ὀνόματι Ταρσέα· ἰδοὺ γὰρ προσεύχεται
\vs{12}καὶ εἶδεν ἄνδρα ἐν ὁράματι Ἁνανίαν ὀνόματι εἰσελθόντα καὶ ἐπιθέντα αὐτῷ τὰς χεῖρας ὅπως ἀναβλέψῃ.
\vs{13}Ἀπεκρίθη δὲ Ἁνανίας· Κύριε, ἤκουσα ἀπὸ πολλῶν περὶ τοῦ ἀνδρὸς τούτου ὅσα κακὰ τοῖς ἁγίοις σου ἐποίησεν ἐν Ἰερουσαλήμ·
\vs{14}καὶ ὧδε ἔχει ἐξουσίαν παρὰ τῶν ἀρχιερέων δῆσαι πάντας τοὺς ἐπικαλουμένους τὸ ὄνομά σου.
\vs{15}Εἶπεν δὲ πρὸς αὐτὸν ὁ Κύριος· Πορεύου, ὅτι σκεῦος ἐκλογῆς ἐστίν μοι οὗτος τοῦ βαστάσαι τὸ ὄνομά μου ἐνώπιον ἐθνῶν τε καὶ βασιλέων υἱῶν τε Ἰσραήλ·
\vs{16}ἐγὼ γὰρ ὑποδείξω αὐτῷ ὅσα δεῖ αὐτὸν ὑπὲρ τοῦ ὀνόματός μου παθεῖν.
\vs{17}Ἀπῆλθεν δὲ Ἁνανίας καὶ εἰσῆλθεν εἰς τὴν οἰκίαν καὶ ἐπιθεὶς ἐπ᾽ αὐτὸν τὰς χεῖρας εἶπεν· Σαοὺλ ἀδελφέ, ὁ Κύριος ἀπέσταλκέν με, Ἰησοῦς ὁ ὀφθείς σοι ἐν τῇ ὁδῷ ᾗ ἤρχου, ὅπως ἀναβλέψῃς καὶ πλησθῇς Πνεύματος Ἁγίου.
\vs{18}Καὶ εὐθέως ἀπέπεσαν αὐτοῦ ἀπὸ τῶν ὀφθαλμῶν ὡς λεπίδες, ἀνέβλεψέν τε καὶ ἀναστὰς ἐβαπτίσθη
\vs{19}καὶ λαβὼν τροφὴν ἐνίσχυσεν.

Ἐγένετο δὲ μετὰ τῶν ἐν Δαμασκῷ μαθητῶν ἡμέρας τινὰς
\vs{20}Καὶ εὐθέως ἐν ταῖς συναγωγαῖς ἐκήρυσσεν τὸν Ἰησοῦν ὅτι οὗτός ἐστιν ὁ Υἱὸς τοῦ Θεοῦ.
\vs{21}Ἐξίσταντο δὲ πάντες οἱ ἀκούοντες καὶ ἔλεγον· Οὐχ οὗτός ἐστιν ὁ πορθήσας εἰς Ἰερουσαλὴμ τοὺς ἐπικαλουμένους τὸ ὄνομα τοῦτο, καὶ ὧδε εἰς τοῦτο ἐληλύθει ἵνα δεδεμένους αὐτοὺς ἀγάγῃ ἐπὶ τοὺς ἀρχιερεῖς;
\vs{22}Σαῦλος δὲ μᾶλλον ἐνεδυναμοῦτο καὶ συνέχυννεν τοὺς Ἰουδαίους τοὺς κατοικοῦντας ἐν Δαμασκῷ συμβιβάζων ὅτι οὗτός ἐστιν ὁ Χριστός.

\vs{23}Ὡς δὲ ἐπληροῦντο ἡμέραι ἱκαναί, συνεβουλεύσαντο οἱ Ἰουδαῖοι ἀνελεῖν αὐτόν·
\vs{24}ἐγνώσθη δὲ τῷ Σαύλῳ ἡ ἐπιβουλὴ αὐτῶν. παρετηροῦντο δὲ καὶ τὰς πύλας ἡμέρας τε καὶ νυκτὸς ὅπως αὐτὸν ἀνέλωσιν·
\vs{25}λαβόντες δὲ οἱ μαθηταὶ αὐτοῦ νυκτὸς διὰ τοῦ τείχους καθῆκαν αὐτὸν χαλάσαντες ἐν σπυρίδι.

\vs{26}Παραγενόμενος δὲ εἰς Ἰερουσαλὴμ ἐπείραζεν κολλᾶσθαι τοῖς μαθηταῖς, καὶ πάντες ἐφοβοῦντο αὐτόν μὴ πιστεύοντες ὅτι ἐστὶν μαθητής.
\vs{27}Βαρνάβας δὲ ἐπιλαβόμενος αὐτὸν ἤγαγεν πρὸς τοὺς ἀποστόλους καὶ διηγήσατο αὐτοῖς πῶς ἐν τῇ ὁδῷ εἶδεν τὸν Κύριον καὶ ὅτι ἐλάλησεν αὐτῷ καὶ πῶς ἐν Δαμασκῷ ἐπαρρησιάσατο ἐν τῷ ὀνόματι τοῦ Ἰησοῦ.
\vs{28}Καὶ ἦν μετ᾽ αὐτῶν εἰσπορευόμενος καὶ ἐκπορευόμενος εἰς Ἰερουσαλήμ, παρρησιαζόμενος ἐν τῷ ὀνόματι τοῦ Κυρίου,
\vs{29}ἐλάλει τε καὶ συνεζήτει πρὸς τοὺς Ἑλληνιστάς, οἱ δὲ ἐπεχείρουν ἀνελεῖν αὐτόν.
\vs{30}ἐπιγνόντες δὲ οἱ ἀδελφοὶ κατήγαγον αὐτὸν εἰς Καισάρειαν καὶ ἐξαπέστειλαν αὐτὸν εἰς Ταρσόν.

\vs{31}Ἡ μὲν οὖν ἐκκλησία καθ᾽ ὅλης τῆς Ἰουδαίας καὶ Γαλιλαίας καὶ Σαμαρείας εἶχεν εἰρήνην οἰκοδομουμένη καὶ πορευομένη τῷ φόβῳ τοῦ Κυρίου καὶ τῇ παρακλήσει τοῦ Ἁγίου Πνεύματος ἐπληθύνετο.

\vs{32}Ἐγένετο δὲ Πέτρον διερχόμενον διὰ πάντων κατελθεῖν καὶ πρὸς τοὺς ἁγίους τοὺς κατοικοῦντας Λύδδα.
\vs{33}εὗρεν δὲ ἐκεῖ ἄνθρωπόν τινα ὀνόματι Αἰνέαν ἐξ ἐτῶν ὀκτὼ κατακείμενον ἐπὶ κραβάττου, ὃς ἦν παραλελυμένος.
\vs{34}καὶ εἶπεν αὐτῷ ὁ Πέτρος· Αἰνέα, ἰᾶταί σε Ἰησοῦς Χριστός· ἀνάστηθι καὶ στρῶσον σεαυτῷ. καὶ εὐθέως ἀνέστη.
\vs{35}καὶ εἶδαν αὐτὸν πάντες οἱ κατοικοῦντες Λύδδα καὶ τὸν Σαρῶνα, οἵτινες ἐπέστρεψαν ἐπὶ τὸν Κύριον.

\vs{36}Ἐν Ἰόππῃ δέ τις ἦν μαθήτρια ὀνόματι Ταβιθά, ἣ διερμηνευομένη λέγεται Δορκάς· αὕτη ἦν πλήρης ἔργων ἀγαθῶν καὶ ἐλεημοσυνῶν ὧν ἐποίει.
\vs{37}ἐγένετο δὲ ἐν ταῖς ἡμέραις ἐκείναις ἀσθενήσασαν αὐτὴν ἀποθανεῖν· λούσαντες δὲ ἔθηκαν αὐτὴν ἐν ὑπερῴῳ.
\vs{38}ἐγγὺς δὲ οὔσης Λύδδας τῇ Ἰόππῃ οἱ μαθηταὶ ἀκούσαντες ὅτι Πέτρος ἐστὶν ἐν αὐτῇ ἀπέστειλαν δύο ἄνδρας πρὸς αὐτὸν παρακαλοῦντες· Μὴ ὀκνήσῃς διελθεῖν ἕως ἡμῶν.
\vs{39}Ἀναστὰς δὲ Πέτρος συνῆλθεν αὐτοῖς· ὃν παραγενόμενον ἀνήγαγον εἰς τὸ ὑπερῷον καὶ παρέστησαν αὐτῷ πᾶσαι αἱ χῆραι κλαίουσαι καὶ ἐπιδεικνύμεναι χιτῶνας καὶ ἱμάτια ὅσα ἐποίει μετ᾽ αὐτῶν οὖσα ἡ Δορκάς.
\vs{40}Ἐκβαλὼν δὲ ἔξω πάντας ὁ Πέτρος καὶ θεὶς τὰ γόνατα προσηύξατο καὶ ἐπιστρέψας πρὸς τὸ σῶμα εἶπεν· Ταβιθά, ἀνάστηθι. ἡ δὲ ἤνοιξεν τοὺς ὀφθαλμοὺς αὐτῆς, καὶ ἰδοῦσα τὸν Πέτρον ἀνεκάθισεν.
\vs{41}δοὺς δὲ αὐτῇ χεῖρα ἀνέστησεν αὐτήν· φωνήσας δὲ τοὺς ἁγίους καὶ τὰς χήρας παρέστησεν αὐτὴν ζῶσαν.
\vs{42}Γνωστὸν δὲ ἐγένετο καθ᾽ ὅλης τῆς Ἰόππης καὶ ἐπίστευσαν πολλοὶ ἐπὶ τὸν Κύριον.
\vs{43}Ἐγένετο δὲ ἡμέρας ἱκανὰς μεῖναι ἐν Ἰόππῃ παρά τινι Σίμωνι βυρσεῖ.

\ch{10}
Ἀνὴρ δέ τις ἐν Καισαρείᾳ ὀνόματι Κορνήλιος, ἑκατοντάρχης ἐκ σπείρης τῆς καλουμένης Ἰταλικῆς,
\vs{2}εὐσεβὴς καὶ φοβούμενος τὸν Θεὸν σὺν παντὶ τῷ οἴκῳ αὐτοῦ, ποιῶν ἐλεημοσύνας πολλὰς τῷ λαῷ καὶ δεόμενος τοῦ Θεοῦ διὰ παντός,
\vs{3}εἶδεν ἐν ὁράματι φανερῶς ὡσεὶ περὶ ὥραν ἐνάτην τῆς ἡμέρας ἄγγελον τοῦ Θεοῦ εἰσελθόντα πρὸς αὐτὸν καὶ εἰπόντα αὐτῷ· Κορνήλιε.
\vs{4}Ὁ δὲ ἀτενίσας αὐτῷ καὶ ἔμφοβος γενόμενος εἶπεν· Τί ἐστιν, Κύριε; Εἶπεν δὲ αὐτῷ· Αἱ προσευχαί σου καὶ αἱ ἐλεημοσύναι σου ἀνέβησαν εἰς μνημόσυνον ἔμπροσθεν τοῦ Θεοῦ.
\vs{5}καὶ νῦν πέμψον ἄνδρας εἰς Ἰόππην καὶ μετάπεμψαι Σίμωνά τινα ὃς ἐπικαλεῖται Πέτρος·
\vs{6}οὗτος ξενίζεται παρά τινι Σίμωνι βυρσεῖ, ᾧ ἐστιν οἰκία παρὰ θάλασσαν.
\vs{7}Ὡς δὲ ἀπῆλθεν ὁ ἄγγελος ὁ λαλῶν αὐτῷ, φωνήσας δύο τῶν οἰκετῶν καὶ στρατιώτην εὐσεβῆ τῶν προσκαρτερούντων αὐτῷ
\vs{8}καὶ ἐξηγησάμενος ἅπαντα αὐτοῖς ἀπέστειλεν αὐτοὺς εἰς τὴν Ἰόππην.
\vs{9}Τῇ δὲ ἐπαύριον, ὁδοιπορούντων ἐκείνων καὶ τῇ πόλει ἐγγιζόντων, ἀνέβη Πέτρος ἐπὶ τὸ δῶμα προσεύξασθαι περὶ ὥραν ἕκτην.
\vs{10}ἐγένετο δὲ πρόσπεινος καὶ ἤθελεν γεύσασθαι. παρασκευαζόντων δὲ αὐτῶν ἐγένετο ἐπ᾽ αὐτὸν ἔκστασις
\vs{11}καὶ θεωρεῖ τὸν οὐρανὸν ἀνεῳγμένον καὶ καταβαῖνον σκεῦός τι ὡς ὀθόνην μεγάλην τέσσαρσιν ἀρχαῖς καθιέμενον ἐπὶ τῆς γῆς,
\vs{12}ἐν ᾧ ὑπῆρχεν πάντα τὰ τετράποδα καὶ ἑρπετὰ τῆς γῆς καὶ πετεινὰ τοῦ οὐρανοῦ.
\vs{13}καὶ ἐγένετο φωνὴ πρὸς αὐτόν· Ἀναστάς, Πέτρε, θῦσον καὶ φάγε.
\vs{14}Ὁ δὲ Πέτρος εἶπεν· Μηδαμῶς, Κύριε, ὅτι οὐδέποτε ἔφαγον πᾶν κοινὸν καὶ ἀκάθαρτον.
\vs{15}Καὶ φωνὴ πάλιν ἐκ δευτέρου πρὸς αὐτόν· Ἃ ὁ Θεὸς ἐκαθάρισεν, σὺ μὴ κοίνου.
\vs{16}Τοῦτο δὲ ἐγένετο ἐπὶ τρίς καὶ εὐθὺς ἀνελήμφθη τὸ σκεῦος εἰς τὸν οὐρανόν.
\vs{17}Ὡς δὲ ἐν ἑαυτῷ διηπόρει ὁ Πέτρος τί ἂν εἴη τὸ ὅραμα ὃ εἶδεν, ἰδοὺ οἱ ἄνδρες οἱ ἀπεσταλμένοι ὑπὸ τοῦ Κορνηλίου διερωτήσαντες τὴν οἰκίαν τοῦ Σίμωνος ἐπέστησαν ἐπὶ τὸν πυλῶνα,
\vs{18}καὶ φωνήσαντες ἐπυνθάνοντο εἰ Σίμων ὁ ἐπικαλούμενος Πέτρος ἐνθάδε ξενίζεται.
\vs{19}Τοῦ δὲ Πέτρου διενθυμουμένου περὶ τοῦ ὁράματος εἶπεν αὐτῷ τὸ Πνεῦμα· Ἰδοὺ ἄνδρες τρεῖς ζητοῦντές σε,
\vs{20}ἀλλὰ ἀναστὰς κατάβηθι καὶ πορεύου σὺν αὐτοῖς μηδὲν διακρινόμενος ὅτι ἐγὼ ἀπέσταλκα αὐτούς.
\vs{21}Καταβὰς δὲ Πέτρος πρὸς τοὺς ἄνδρας εἶπεν· Ἰδοὺ ἐγώ εἰμι ὃν ζητεῖτε· τίς ἡ αἰτία δι᾽ ἣν πάρεστε;
\vs{22}Οἱ δὲ εἶπαν· Κορνήλιος ἑκατοντάρχης, ἀνὴρ δίκαιος καὶ φοβούμενος τὸν Θεὸν, μαρτυρούμενός τε ὑπὸ ὅλου τοῦ ἔθνους τῶν Ἰουδαίων, ἐχρηματίσθη ὑπὸ ἀγγέλου ἁγίου μεταπέμψασθαί σε εἰς τὸν οἶκον αὐτοῦ καὶ ἀκοῦσαι ῥήματα παρὰ σοῦ.
\vs{23}Εἰσκαλεσάμενος οὖν αὐτοὺς ἐξένισεν. Τῇ δὲ ἐπαύριον ἀναστὰς ἐξῆλθεν σὺν αὐτοῖς καί τινες τῶν ἀδελφῶν τῶν ἀπὸ Ἰόππης συνῆλθον αὐτῷ.
\vs{24}Τῇ δὲ ἐπαύριον εἰσῆλθεν εἰς τὴν Καισάρειαν. ὁ δὲ Κορνήλιος ἦν προσδοκῶν αὐτοὺς συνκαλεσάμενος τοὺς συγγενεῖς αὐτοῦ καὶ τοὺς ἀναγκαίους φίλους.
\vs{25}Ὡς δὲ ἐγένετο τοῦ εἰσελθεῖν τὸν Πέτρον, συναντήσας αὐτῷ ὁ Κορνήλιος πεσὼν ἐπὶ τοὺς πόδας προσεκύνησεν.
\vs{26}ὁ δὲ Πέτρος ἤγειρεν αὐτὸν λέγων· Ἀνάστηθι· καὶ ἐγὼ αὐτὸς ἄνθρωπός εἰμι.
\vs{27}Καὶ συνομιλῶν αὐτῷ εἰσῆλθεν καὶ εὑρίσκει συνεληλυθότας πολλούς,
\vs{28}ἔφη τε πρὸς αὐτούς· Ὑμεῖς ἐπίστασθε ὡς ἀθέμιτόν ἐστιν ἀνδρὶ Ἰουδαίῳ κολλᾶσθαι ἢ προσέρχεσθαι ἀλλοφύλῳ· κἀμοὶ ὁ Θεὸς ἔδειξεν μηδένα κοινὸν ἢ ἀκάθαρτον λέγειν ἄνθρωπον·
\vs{29}διὸ καὶ ἀναντιρρήτως ἦλθον μεταπεμφθείς. πυνθάνομαι οὖν Τίνι λόγῳ μετεπέμψασθέ με;
\vs{30}Καὶ ὁ Κορνήλιος ἔφη· Ἀπὸ τετάρτης ἡμέρας μέχρι ταύτης τῆς ὥρας ἤμην τὴν ἐνάτην προσευχόμενος ἐν τῷ οἴκῳ μου, καὶ ἰδοὺ ἀνὴρ ἔστη ἐνώπιόν μου ἐν ἐσθῆτι λαμπρᾷ
\vs{31}καὶ φησίν· Κορνήλιε, εἰσηκούσθη σου ἡ προσευχὴ καὶ αἱ ἐλεημοσύναι σου ἐμνήσθησαν ἐνώπιον τοῦ Θεοῦ.
\vs{32}πέμψον οὖν εἰς Ἰόππην καὶ μετακάλεσαι Σίμωνα ὃς ἐπικαλεῖται Πέτρος, οὗτος ξενίζεται ἐν οἰκίᾳ Σίμωνος βυρσέως παρὰ θάλασσαν.
\vs{33}Ἐξαυτῆς οὖν ἔπεμψα πρὸς σέ, σύ τε καλῶς ἐποίησας παραγενόμενος. νῦν οὖν πάντες ἡμεῖς ἐνώπιον τοῦ Θεοῦ πάρεσμεν ἀκοῦσαι πάντα τὰ προστεταγμένα σοι ὑπὸ τοῦ Κυρίου.

\vs{34}Ἀνοίξας δὲ Πέτρος τὸ στόμα εἶπεν· Ἐπ᾽ ἀληθείας καταλαμβάνομαι ὅτι οὐκ ἔστιν προσωπολήμπτης ὁ Θεός,
\vs{35}ἀλλ᾽ ἐν παντὶ ἔθνει ὁ φοβούμενος αὐτὸν καὶ ἐργαζόμενος δικαιοσύνην δεκτὸς αὐτῷ ἐστιν.
\vs{36}τὸν λόγον ὃν ἀπέστειλεν τοῖς υἱοῖς Ἰσραὴλ εὐαγγελιζόμενος εἰρήνην διὰ Ἰησοῦ Χριστοῦ, οὗτός ἐστιν πάντων Κύριος,
\vs{37}Ὑμεῖς οἴδατε τὸ γενόμενον ῥῆμα καθ᾽ ὅλης τῆς Ἰουδαίας, ἀρξάμενος ἀπὸ τῆς Γαλιλαίας μετὰ τὸ βάπτισμα ὃ ἐκήρυξεν Ἰωάννης,
\vs{38}Ἰησοῦν τὸν ἀπὸ Ναζαρέθ, ὡς ἔχρισεν αὐτὸν ὁ Θεὸς Πνεύματι Ἁγίῳ καὶ δυνάμει, ὃς διῆλθεν εὐεργετῶν καὶ ἰώμενος πάντας τοὺς καταδυναστευομένους ὑπὸ τοῦ διαβόλου, ὅτι ὁ Θεὸς ἦν μετ᾽ αὐτοῦ.
\vs{39}Καὶ ἡμεῖς μάρτυρες πάντων ὧν ἐποίησεν ἔν τε τῇ χώρᾳ τῶν Ἰουδαίων καὶ ἐν Ἰερουσαλήμ. ὃν καὶ ἀνεῖλαν κρεμάσαντες ἐπὶ ξύλου,
\vs{40}τοῦτον ὁ Θεὸς ἤγειρεν ἐν τῇ τρίτῃ ἡμέρᾳ καὶ ἔδωκεν αὐτὸν ἐμφανῆ γενέσθαι,
\vs{41}οὐ παντὶ τῷ λαῷ, ἀλλὰ μάρτυσιν τοῖς προκεχειροτονημένοις ὑπὸ τοῦ Θεοῦ, ἡμῖν, οἵτινες συνεφάγομεν καὶ συνεπίομεν αὐτῷ μετὰ τὸ ἀναστῆναι αὐτὸν ἐκ νεκρῶν·
\vs{42}καὶ παρήγγειλεν ἡμῖν κηρύξαι τῷ λαῷ καὶ διαμαρτύρασθαι ὅτι οὗτός ἐστιν ὁ ὡρισμένος ὑπὸ τοῦ Θεοῦ Κριτὴς ζώντων καὶ νεκρῶν.
\vs{43}τούτῳ πάντες οἱ προφῆται μαρτυροῦσιν ἄφεσιν ἁμαρτιῶν λαβεῖν διὰ τοῦ ὀνόματος αὐτοῦ πάντα τὸν πιστεύοντα εἰς αὐτόν.

\vs{44}Ἔτι λαλοῦντος τοῦ Πέτρου τὰ ῥήματα ταῦτα ἐπέπεσεν τὸ Πνεῦμα τὸ Ἅγιον ἐπὶ πάντας τοὺς ἀκούοντας τὸν λόγον.
\vs{45}καὶ ἐξέστησαν οἱ ἐκ περιτομῆς πιστοὶ ὅσοι συνῆλθαν τῷ Πέτρῳ, ὅτι καὶ ἐπὶ τὰ ἔθνη ἡ δωρεὰ τοῦ Ἁγίου Πνεύματος ἐκκέχυται·
\vs{46}ἤκουον γὰρ αὐτῶν λαλούντων γλώσσαις καὶ μεγαλυνόντων τὸν Θεόν. Τότε ἀπεκρίθη Πέτρος·
\vs{47}Μήτι τὸ ὕδωρ δύναται κωλῦσαί τις τοῦ μὴ βαπτισθῆναι τούτους, οἵτινες τὸ Πνεῦμα τὸ Ἅγιον ἔλαβον ὡς καὶ ἡμεῖς;
\vs{48}προσέταξεν δὲ αὐτοὺς ἐν τῷ ὀνόματι Ἰησοῦ Χριστοῦ βαπτισθῆναι. τότε ἠρώτησαν αὐτὸν ἐπιμεῖναι ἡμέρας τινάς.

\ch{11}
Ἤκουσαν δὲ οἱ ἀπόστολοι καὶ οἱ ἀδελφοὶ οἱ ὄντες κατὰ τὴν Ἰουδαίαν ὅτι καὶ τὰ ἔθνη ἐδέξαντο τὸν λόγον τοῦ Θεοῦ.
\vs{2}Ὅτε δὲ ἀνέβη Πέτρος εἰς Ἰερουσαλήμ, διεκρίνοντο πρὸς αὐτὸν οἱ ἐκ περιτομῆς
\vs{3}λέγοντες ὅτι Εἰσῆλθες πρὸς ἄνδρας ἀκροβυστίαν ἔχοντας καὶ συνέφαγες αὐτοῖς.
\vs{4}Ἀρξάμενος δὲ Πέτρος ἐξετίθετο αὐτοῖς καθεξῆς λέγων·
\vs{5}Ἐγὼ ἤμην ἐν πόλει Ἰόππῃ προσευχόμενος καὶ εἶδον ἐν ἐκστάσει ὅραμα, καταβαῖνον σκεῦός τι ὡς ὀθόνην μεγάλην τέσσαρσιν ἀρχαῖς καθιεμένην ἐκ τοῦ οὐρανοῦ, καὶ ἦλθεν ἄχρι ἐμοῦ.
\vs{6}εἰς ἣν ἀτενίσας κατενόουν καὶ εἶδον τὰ τετράποδα τῆς γῆς καὶ τὰ θηρία καὶ τὰ ἑρπετὰ καὶ τὰ πετεινὰ τοῦ οὐρανοῦ.
\vs{7}ἤκουσα δὲ καὶ φωνῆς λεγούσης μοι· Ἀναστάς, Πέτρε, θῦσον καὶ φάγε.
\vs{8}Εἶπον δέ· Μηδαμῶς, Κύριε, ὅτι κοινὸν ἢ ἀκάθαρτον οὐδέποτε εἰσῆλθεν εἰς τὸ στόμα μου.
\vs{9}Ἀπεκρίθη δὲ φωνὴ ἐκ δευτέρου ἐκ τοῦ οὐρανοῦ· Ἃ ὁ Θεὸς ἐκαθάρισεν, σὺ μὴ κοίνου.
\vs{10}Τοῦτο δὲ ἐγένετο ἐπὶ τρίς, καὶ ἀνεσπάσθη πάλιν ἅπαντα εἰς τὸν οὐρανόν.
\vs{11}Καὶ ἰδοὺ ἐξαυτῆς τρεῖς ἄνδρες ἐπέστησαν ἐπὶ τὴν οἰκίαν ἐν ᾗ ἦμεν, ἀπεσταλμένοι ἀπὸ Καισαρείας πρός με.
\vs{12}εἶπεν δὲ τὸ Πνεῦμά μοι συνελθεῖν αὐτοῖς μηδὲν διακρίναντα. ἦλθον δὲ σὺν ἐμοὶ καὶ οἱ ἓξ ἀδελφοὶ οὗτοι καὶ εἰσήλθομεν εἰς τὸν οἶκον τοῦ ἀνδρός.
\vs{13}ἀπήγγειλεν δὲ ἡμῖν πῶς εἶδεν τὸν ἄγγελον ἐν τῷ οἴκῳ αὐτοῦ σταθέντα καὶ εἰπόντα· Ἀπόστειλον εἰς Ἰόππην καὶ μετάπεμψαι Σίμωνα τὸν ἐπικαλούμενον Πέτρον,
\vs{14}ὃς λαλήσει ῥήματα πρὸς σὲ ἐν οἷς σωθήσῃ σὺ καὶ πᾶς ὁ οἶκός σου.
\vs{15}Ἐν δὲ τῷ ἄρξασθαί με λαλεῖν ἐπέπεσεν τὸ Πνεῦμα τὸ Ἅγιον ἐπ᾽ αὐτοὺς ὥσπερ καὶ ἐφ᾽ ἡμᾶς ἐν ἀρχῇ.
\vs{16}ἐμνήσθην δὲ τοῦ ῥήματος τοῦ Κυρίου ὡς ἔλεγεν· Ἰωάννης μὲν ἐβάπτισεν ὕδατι, ὑμεῖς δὲ βαπτισθήσεσθε ἐν Πνεύματι Ἁγίῳ.
\vs{17}εἰ οὖν τὴν ἴσην δωρεὰν ἔδωκεν αὐτοῖς ὁ Θεὸς ὡς καὶ ἡμῖν πιστεύσασιν ἐπὶ τὸν Κύριον Ἰησοῦν Χριστόν, ἐγὼ τίς ἤμην δυνατὸς κωλῦσαι τὸν Θεόν;
\vs{18}Ἀκούσαντες δὲ ταῦτα ἡσύχασαν καὶ ἐδόξασαν τὸν Θεὸν λέγοντες· Ἄρα καὶ τοῖς ἔθνεσιν ὁ Θεὸς τὴν μετάνοιαν εἰς ζωὴν ἔδωκεν.

\vs{19}Οἱ μὲν οὖν διασπαρέντες ἀπὸ τῆς θλίψεως τῆς γενομένης ἐπὶ Στεφάνῳ διῆλθον ἕως Φοινίκης καὶ Κύπρου καὶ Ἀντιοχείας μηδενὶ λαλοῦντες τὸν λόγον εἰ μὴ μόνον Ἰουδαίοις.
\vs{20}Ἦσαν δέ τινες ἐξ αὐτῶν ἄνδρες Κύπριοι καὶ Κυρηναῖοι, οἵτινες ἐλθόντες εἰς Ἀντιόχειαν ἐλάλουν καὶ πρὸς τοὺς Ἑλληνιστάς εὐαγγελιζόμενοι τὸν Κύριον Ἰησοῦν.
\vs{21}καὶ ἦν χεὶρ Κυρίου μετ᾽ αὐτῶν, πολύς τε ἀριθμὸς ὁ πιστεύσας ἐπέστρεψεν ἐπὶ τὸν Κύριον.
\vs{22}Ἠκούσθη δὲ ὁ λόγος εἰς τὰ ὦτα τῆς ἐκκλησίας τῆς οὔσης ἐν Ἰερουσαλὴμ περὶ αὐτῶν καὶ ἐξαπέστειλαν Βαρνάβαν διελθεῖν ἕως Ἀντιοχείας.
\vs{23}ὃς παραγενόμενος καὶ ἰδὼν τὴν χάριν τὴν τοῦ Θεοῦ, ἐχάρη καὶ παρεκάλει πάντας τῇ προθέσει τῆς καρδίας προσμένειν τῷ Κυρίῳ,
\vs{24}ὅτι ἦν ἀνὴρ ἀγαθὸς καὶ πλήρης Πνεύματος Ἁγίου καὶ πίστεως. καὶ προσετέθη ὄχλος ἱκανὸς τῷ Κυρίῳ.
\vs{25}Ἐξῆλθεν δὲ εἰς Ταρσὸν ἀναζητῆσαι Σαῦλον,
\vs{26}καὶ εὑρὼν ἤγαγεν εἰς Ἀντιόχειαν. ἐγένετο δὲ αὐτοῖς καὶ ἐνιαυτὸν ὅλον συναχθῆναι ἐν τῇ ἐκκλησίᾳ καὶ διδάξαι ὄχλον ἱκανόν, χρηματίσαι τε πρώτως ἐν Ἀντιοχείᾳ τοὺς μαθητὰς Χριστιανούς.

\vs{27}Ἐν ταύταις δὲ ταῖς ἡμέραις κατῆλθον ἀπὸ Ἱεροσολύμων προφῆται εἰς Ἀντιόχειαν.
\vs{28}ἀναστὰς δὲ εἷς ἐξ αὐτῶν ὀνόματι Ἅγαβος ἐσήμανεν διὰ τοῦ Πνεύματος λιμὸν μεγάλην μέλλειν ἔσεσθαι ἐφ᾽ ὅλην τὴν οἰκουμένην, ἥτις ἐγένετο ἐπὶ Κλαυδίου.
\vs{29}τῶν δὲ μαθητῶν, καθὼς εὐπορεῖτό τις, ὥρισαν ἕκαστος αὐτῶν εἰς διακονίαν πέμψαι τοῖς κατοικοῦσιν ἐν τῇ Ἰουδαίᾳ ἀδελφοῖς·
\vs{30}ὃ καὶ ἐποίησαν ἀποστείλαντες πρὸς τοὺς πρεσβυτέρους διὰ χειρὸς Βαρνάβα καὶ Σαύλου.

\ch{12}
Κατ᾽ ἐκεῖνον δὲ τὸν καιρὸν ἐπέβαλεν Ἡρῴδης ὁ βασιλεὺς τὰς χεῖρας κακῶσαί τινας τῶν ἀπὸ τῆς ἐκκλησίας.
\vs{2}ἀνεῖλεν δὲ Ἰάκωβον τὸν ἀδελφὸν Ἰωάννου μαχαίρῃ.
\vs{3}ἰδὼν δὲ ὅτι ἀρεστόν ἐστιν τοῖς Ἰουδαίοις, προσέθετο συλλαβεῖν καὶ Πέτρον,— ἦσαν δὲ αἱ ἡμέραι τῶν ἀζύμων—
\vs{4}ὃν καὶ πιάσας ἔθετο εἰς φυλακήν παραδοὺς τέσσαρσιν τετραδίοις στρατιωτῶν φυλάσσειν αὐτόν, βουλόμενος μετὰ τὸ πάσχα ἀναγαγεῖν αὐτὸν τῷ λαῷ.
\vs{5}Ὁ μὲν οὖν Πέτρος ἐτηρεῖτο ἐν τῇ φυλακῇ· προσευχὴ δὲ ἦν ἐκτενῶς γινομένη ὑπὸ τῆς ἐκκλησίας πρὸς τὸν Θεὸν περὶ αὐτοῦ.

\vs{6}Ὅτε δὲ ἤμελλεν προαγαγεῖν αὐτὸν ὁ Ἡρῴδης, τῇ νυκτὶ ἐκείνῃ ἦν ὁ Πέτρος κοιμώμενος μεταξὺ δύο στρατιωτῶν δεδεμένος ἁλύσεσιν δυσίν φύλακές τε πρὸ τῆς θύρας ἐτήρουν τὴν φυλακήν.
\vs{7}καὶ ἰδοὺ ἄγγελος Κυρίου ἐπέστη καὶ φῶς ἔλαμψεν ἐν τῷ οἰκήματι· πατάξας δὲ τὴν πλευρὰν τοῦ Πέτρου ἤγειρεν αὐτὸν λέγων· Ἀνάστα ἐν τάχει. καὶ ἐξέπεσαν αὐτοῦ αἱ ἁλύσεις ἐκ τῶν χειρῶν.
\vs{8}εἶπεν δὲ ὁ ἄγγελος πρὸς αὐτόν· Ζῶσαι καὶ ὑπόδησαι τὰ σανδάλιά σου. ἐποίησεν δὲ οὕτως. καὶ λέγει αὐτῷ· Περιβαλοῦ τὸ ἱμάτιόν σου καὶ ἀκολούθει μοι.
\vs{9}Καὶ ἐξελθὼν ἠκολούθει καὶ οὐκ ᾔδει ὅτι ἀληθές ἐστιν τὸ γινόμενον διὰ τοῦ ἀγγέλου· ἐδόκει δὲ ὅραμα βλέπειν.
\vs{10}διελθόντες δὲ πρώτην φυλακὴν καὶ δευτέραν ἦλθαν ἐπὶ τὴν πύλην τὴν σιδηρᾶν τὴν φέρουσαν εἰς τὴν πόλιν, ἥτις αὐτομάτη ἠνοίγη αὐτοῖς καὶ ἐξελθόντες προῆλθον ῥύμην μίαν, καὶ εὐθέως ἀπέστη ὁ ἄγγελος ἀπ᾽ αὐτοῦ.
\vs{11}Καὶ ὁ Πέτρος ἐν ἑαυτῷ γενόμενος εἶπεν· Νῦν οἶδα ἀληθῶς ὅτι ἐξαπέστειλεν ὁ Κύριος τὸν ἄγγελον αὐτοῦ καὶ ἐξείλατό με ἐκ χειρὸς Ἡρῴδου καὶ πάσης τῆς προσδοκίας τοῦ λαοῦ τῶν Ἰουδαίων.
\vs{12}Συνιδών τε ἦλθεν ἐπὶ τὴν οἰκίαν τῆς Μαρίας τῆς μητρὸς Ἰωάννου τοῦ ἐπικαλουμένου Μάρκου, οὗ ἦσαν ἱκανοὶ συνηθροισμένοι καὶ προσευχόμενοι.
\vs{13}κρούσαντος δὲ αὐτοῦ τὴν θύραν τοῦ πυλῶνος προσῆλθεν παιδίσκη ὑπακοῦσαι ὀνόματι Ῥόδη,
\vs{14}καὶ ἐπιγνοῦσα τὴν φωνὴν τοῦ Πέτρου ἀπὸ τῆς χαρᾶς οὐκ ἤνοιξεν τὸν πυλῶνα, εἰσδραμοῦσα δὲ ἀπήγγειλεν ἑστάναι τὸν Πέτρον πρὸ τοῦ πυλῶνος.
\vs{15}Οἱ δὲ πρὸς αὐτὴν εἶπαν· Μαίνῃ. ἡ δὲ διϊσχυρίζετο οὕτως ἔχειν. οἱ δὲ ἔλεγον· Ὁ ἄγγελός ἐστιν αὐτοῦ.
\vs{16}Ὁ δὲ Πέτρος ἐπέμενεν κρούων· ἀνοίξαντες δὲ εἶδαν αὐτὸν καὶ ἐξέστησαν.
\vs{17}κατασείσας δὲ αὐτοῖς τῇ χειρὶ σιγᾶν διηγήσατο αὐτοῖς πῶς ὁ Κύριος αὐτὸν ἐξήγαγεν ἐκ τῆς φυλακῆς εἶπέν τε· Ἀπαγγείλατε Ἰακώβῳ καὶ τοῖς ἀδελφοῖς ταῦτα. καὶ ἐξελθὼν ἐπορεύθη εἰς ἕτερον τόπον.

\vs{18}Γενομένης δὲ ἡμέρας ἦν τάραχος οὐκ ὀλίγος ἐν τοῖς στρατιώταις τί ἄρα ὁ Πέτρος ἐγένετο.
\vs{19}Ἡρῴδης δὲ ἐπιζητήσας αὐτὸν καὶ μὴ εὑρὼν, ἀνακρίνας τοὺς φύλακας ἐκέλευσεν ἀπαχθῆναι, καὶ κατελθὼν ἀπὸ τῆς Ἰουδαίας εἰς Καισάρειαν διέτριβεν.

\vs{20}Ἦν δὲ θυμομαχῶν Τυρίοις καὶ Σιδωνίοις· ὁμοθυμαδὸν δὲ παρῆσαν πρὸς αὐτόν καὶ πείσαντες Βλάστον, τὸν ἐπὶ τοῦ κοιτῶνος τοῦ βασιλέως, ᾐτοῦντο εἰρήνην διὰ τὸ τρέφεσθαι αὐτῶν τὴν χώραν ἀπὸ τῆς βασιλικῆς.
\vs{21}τακτῇ δὲ ἡμέρᾳ ὁ Ἡρῴδης ἐνδυσάμενος ἐσθῆτα βασιλικὴν καὶ καθίσας ἐπὶ τοῦ βήματος ἐδημηγόρει πρὸς αὐτούς,
\vs{22}ὁ δὲ δῆμος ἐπεφώνει· Θεοῦ φωνὴ καὶ οὐκ ἀνθρώπου.
\vs{23}Παραχρῆμα δὲ ἐπάταξεν αὐτὸν ἄγγελος Κυρίου ἀνθ᾽ ὧν οὐκ ἔδωκεν τὴν δόξαν τῷ Θεῷ, καὶ γενόμενος σκωληκόβρωτος ἐξέψυξεν.

\vs{24}Ὁ δὲ λόγος τοῦ θεοῦ ηὔξανεν καὶ ἐπληθύνετο.
\vs{25}Βαρνάβας δὲ καὶ Σαῦλος ὑπέστρεψαν εἰς Ἰερουσαλὴμ πληρώσαντες τὴν διακονίαν, συμπαραλαβόντες Ἰωάννην τὸν ἐπικληθέντα Μάρκον.

\ch{13}
Ἦσαν δὲ ἐν Ἀντιοχείᾳ κατὰ τὴν οὖσαν ἐκκλησίαν προφῆται καὶ διδάσκαλοι ὅ τε Βαρνάβας καὶ Συμεὼν ὁ καλούμενος Νίγερ καὶ Λούκιος ὁ Κυρηναῖος, Μαναήν τε Ἡρῴδου τοῦ τετραάρχου σύντροφος καὶ Σαῦλος.
\vs{2}Λειτουργούντων δὲ αὐτῶν τῷ Κυρίῳ καὶ νηστευόντων εἶπεν τὸ Πνεῦμα τὸ Ἅγιον· Ἀφορίσατε δή μοι τὸν Βαρνάβαν καὶ Σαῦλον εἰς τὸ ἔργον ὃ προσκέκλημαι αὐτούς.
\vs{3}τότε νηστεύσαντες καὶ προσευξάμενοι καὶ ἐπιθέντες τὰς χεῖρας αὐτοῖς ἀπέλυσαν.

\vs{4}Αὐτοὶ μὲν οὖν ἐκπεμφθέντες ὑπὸ τοῦ Ἁγίου Πνεύματος κατῆλθον εἰς Σελεύκειαν, ἐκεῖθέν τε ἀπέπλευσαν εἰς Κύπρον
\vs{5}καὶ γενόμενοι ἐν Σαλαμῖνι κατήγγελλον τὸν λόγον τοῦ Θεοῦ ἐν ταῖς συναγωγαῖς τῶν Ἰουδαίων. εἶχον δὲ καὶ Ἰωάννην ὑπηρέτην.
\vs{6}Διελθόντες δὲ ὅλην τὴν νῆσον ἄχρι Πάφου εὗρον ἄνδρα τινὰ μάγον ψευδοπροφήτην Ἰουδαῖον ᾧ ὄνομα Βαριησοῦ
\vs{7}ὃς ἦν σὺν τῷ ἀνθυπάτῳ Σεργίῳ Παύλῳ, ἀνδρὶ συνετῷ. οὗτος προσκαλεσάμενος Βαρνάβαν καὶ Σαῦλον ἐπεζήτησεν ἀκοῦσαι τὸν λόγον τοῦ Θεοῦ.
\vs{8}ἀνθίστατο δὲ αὐτοῖς Ἐλύμας ὁ μάγος, οὕτως γὰρ μεθερμηνεύεται τὸ ὄνομα αὐτοῦ, ζητῶν διαστρέψαι τὸν ἀνθύπατον ἀπὸ τῆς πίστεως.
\vs{9}Σαῦλος δέ, ὁ καὶ Παῦλος, πλησθεὶς Πνεύματος Ἁγίου ἀτενίσας εἰς αὐτὸν
\vs{10}εἶπεν· Ὦ πλήρης παντὸς δόλου καὶ πάσης ῥᾳδιουργίας, υἱὲ διαβόλου, ἐχθρὲ πάσης δικαιοσύνης, οὐ παύσῃ διαστρέφων τὰς ὁδοὺς τοῦ Κυρίου τὰς εὐθείας;
\vs{11}καὶ νῦν ἰδοὺ χεὶρ Κυρίου ἐπὶ σέ καὶ ἔσῃ τυφλὸς μὴ βλέπων τὸν ἥλιον ἄχρι καιροῦ. παραχρῆμα δὲ ἔπεσεν ἐπ᾽ αὐτὸν ἀχλὺς καὶ σκότος καὶ περιάγων ἐζήτει χειραγωγούς.
\vs{12}Τότε ἰδὼν ὁ ἀνθύπατος τὸ γεγονὸς ἐπίστευσεν ἐκπλησσόμενος ἐπὶ τῇ διδαχῇ τοῦ Κυρίου.

\vs{13}Ἀναχθέντες δὲ ἀπὸ τῆς Πάφου οἱ περὶ Παῦλον ἦλθον εἰς Πέργην τῆς Παμφυλίας, Ἰωάννης δὲ ἀποχωρήσας ἀπ᾽ αὐτῶν ὑπέστρεψεν εἰς Ἱεροσόλυμα.
\vs{14}Αὐτοὶ δὲ διελθόντες ἀπὸ τῆς Πέργης παρεγένοντο εἰς Ἀντιόχειαν τὴν Πισιδίαν, καὶ εἰσελθόντες εἰς τὴν συναγωγὴν τῇ ἡμέρᾳ τῶν σαββάτων ἐκάθισαν.
\vs{15}μετὰ δὲ τὴν ἀνάγνωσιν τοῦ νόμου καὶ τῶν προφητῶν ἀπέστειλαν οἱ ἀρχισυνάγωγοι πρὸς αὐτοὺς λέγοντες· Ἄνδρες ἀδελφοί, εἴ τίς ἐστιν ἐν ὑμῖν λόγος παρακλήσεως πρὸς τὸν λαόν, λέγετε.
\vs{16}Ἀναστὰς δὲ Παῦλος καὶ κατασείσας τῇ χειρὶ εἶπεν· Ἄνδρες Ἰσραηλῖται καὶ οἱ φοβούμενοι τὸν Θεόν, ἀκούσατε.
\vs{17}ὁ Θεὸς τοῦ λαοῦ τούτου Ἰσραὴλ ἐξελέξατο τοὺς πατέρας ἡμῶν καὶ τὸν λαὸν ὕψωσεν ἐν τῇ παροικίᾳ ἐν γῇ Αἰγύπτου καὶ μετὰ βραχίονος ὑψηλοῦ ἐξήγαγεν αὐτοὺς ἐξ αὐτῆς,
\vs{18}καί ὡς τεσσερακονταετῆ χρόνον ἐτροποφόρησεν αὐτοὺς ἐν τῇ ἐρήμῳ
\vs{19}καὶ καθελὼν ἔθνη ἑπτὰ ἐν γῇ Χανάαν κατεκληρονόμησεν τὴν γῆν αὐτῶν
\vs{20}ὡς ἔτεσιν τετρακοσίοις καὶ πεντήκοντα. Καὶ μετὰ ταῦτα ἔδωκεν κριτὰς ἕως Σαμουὴλ τοῦ προφήτου.
\vs{21}κἀκεῖθεν ᾐτήσαντο βασιλέα καὶ ἔδωκεν αὐτοῖς ὁ Θεὸς τὸν Σαοὺλ υἱὸν Κίς, ἄνδρα ἐκ φυλῆς Βενιαμίν, ἔτη τεσσεράκοντα,
\vs{22}καὶ μεταστήσας αὐτὸν ἤγειρεν τὸν Δαυὶδ αὐτοῖς εἰς βασιλέα ᾧ καὶ εἶπεν μαρτυρήσας· Εὗρον Δαυὶδ τὸν τοῦ Ἰεσσαί, ἄνδρα κατὰ τὴν καρδίαν μου, ὃς ποιήσει πάντα τὰ θελήματά μου.
\vs{23}Τούτου ὁ Θεὸς ἀπὸ τοῦ σπέρματος κατ᾽ ἐπαγγελίαν ἤγαγεν τῷ Ἰσραὴλ Σωτῆρα Ἰησοῦν,
\vs{24}προκηρύξαντος Ἰωάννου πρὸ προσώπου τῆς εἰσόδου αὐτοῦ βάπτισμα μετανοίας παντὶ τῷ λαῷ Ἰσραήλ.
\vs{25}ὡς δὲ ἐπλήρου Ἰωάννης τὸν δρόμον, ἔλεγεν· Τί ἐμὲ ὑπονοεῖτε εἶναι; οὐκ εἰμὶ ἐγώ· ἀλλ᾽ ἰδοὺ ἔρχεται μετ᾽ ἐμὲ οὗ οὐκ εἰμὶ ἄξιος τὸ ὑπόδημα τῶν ποδῶν λῦσαι.
\vs{26}Ἄνδρες ἀδελφοί, υἱοὶ γένους Ἀβραὰμ καὶ οἱ ἐν ὑμῖν φοβούμενοι τὸν Θεόν, ἡμῖν ὁ λόγος τῆς σωτηρίας ταύτης ἐξαπεστάλη.
\vs{27}οἱ γὰρ κατοικοῦντες ἐν Ἰερουσαλὴμ καὶ οἱ ἄρχοντες αὐτῶν τοῦτον ἀγνοήσαντες καὶ τὰς φωνὰς τῶν προφητῶν τὰς κατὰ πᾶν σάββατον ἀναγινωσκομένας κρίναντες ἐπλήρωσαν,
\vs{28}καὶ μηδεμίαν αἰτίαν θανάτου εὑρόντες ᾐτήσαντο Πιλᾶτον ἀναιρεθῆναι αὐτόν.
\vs{29}Ὡς δὲ ἐτέλεσαν πάντα τὰ περὶ αὐτοῦ γεγραμμένα, καθελόντες ἀπὸ τοῦ ξύλου ἔθηκαν εἰς μνημεῖον.
\vs{30}ὁ δὲ Θεὸς ἤγειρεν αὐτὸν ἐκ νεκρῶν,
\vs{31}ὃς ὤφθη ἐπὶ ἡμέρας πλείους τοῖς συναναβᾶσιν αὐτῷ ἀπὸ τῆς Γαλιλαίας εἰς Ἰερουσαλήμ, οἵτινες νῦν εἰσιν μάρτυρες αὐτοῦ πρὸς τὸν λαόν.
\vs{32}Καὶ ἡμεῖς ὑμᾶς εὐαγγελιζόμεθα τὴν πρὸς τοὺς πατέρας ἐπαγγελίαν γενομένην,
\vs{33}ὅτι ταύτην ὁ Θεὸς ἐκπεπλήρωκεν τοῖς τέκνοις αὐτῶν ἡμῶν ἀναστήσας Ἰησοῦν ὡς καὶ ἐν τῷ ψαλμῷ γέγραπται τῷ δευτέρῳ· 
\begin{poetryblock}

\begin{quote}Υἱός μου εἶ σύ,\end{quote} 

\begin{quote}ἐγὼ σήμερον γεγέννηκά σε.\end{quote}
\end{poetryblock}
\vs{34}Ὅτι δὲ ἀνέστησεν αὐτὸν ἐκ νεκρῶν μηκέτι μέλλοντα ὑποστρέφειν εἰς διαφθοράν, οὕτως εἴρηκεν ὅτι Δώσω ὑμῖν τὰ ὅσια Δαυὶδ τὰ πιστά.
\vs{35}Διότι καὶ ἐν ἑτέρῳ λέγει· Οὐ δώσεις τὸν Ὅσιόν σου ἰδεῖν διαφθοράν.
\vs{36}Δαυὶδ μὲν γὰρ ἰδίᾳ γενεᾷ ὑπηρετήσας τῇ τοῦ Θεοῦ βουλῇ ἐκοιμήθη καὶ προσετέθη πρὸς τοὺς πατέρας αὐτοῦ καὶ εἶδεν διαφθοράν·
\vs{37}ὃν δὲ ὁ Θεὸς ἤγειρεν, οὐκ εἶδεν διαφθοράν.
\vs{38}Γνωστὸν οὖν ἔστω ὑμῖν, ἄνδρες ἀδελφοί, ὅτι διὰ τούτου ὑμῖν ἄφεσις ἁμαρτιῶν καταγγέλλεται, καὶ ἀπὸ πάντων ὧν οὐκ ἠδυνήθητε ἐν νόμῳ Μωϋσέως δικαιωθῆναι,
\vs{39}ἐν τούτῳ πᾶς ὁ πιστεύων δικαιοῦται.
\vs{40}βλέπετε οὖν μὴ ἐπέλθῃ τὸ εἰρημένον ἐν τοῖς προφήταις·
\begin{poetryblock}

\begin{quote} \vs{41}Ἴδετε, οἱ καταφρονηταί,\end{quote} 

\begin{quote}καὶ θαυμάσατε καὶ ἀφανίσθητε,\end{quote} 

\begin{quote}ὅτι ἔργον ἐργάζομαι ἐγὼ ἐν ταῖς ἡμέραις ὑμῶν,\end{quote} 

\begin{quote}ἔργον ὃ οὐ μὴ πιστεύσητε ἐάν τις ἐκδιηγῆται ὑμῖν.\end{quote}
\end{poetryblock}

\vs{42}Ἐξιόντων δὲ αὐτῶν παρεκάλουν εἰς τὸ μεταξὺ σάββατον λαληθῆναι αὐτοῖς τὰ ῥήματα ταῦτα.
\vs{43}λυθείσης δὲ τῆς συναγωγῆς ἠκολούθησαν πολλοὶ τῶν Ἰουδαίων καὶ τῶν σεβομένων προσηλύτων τῷ Παύλῳ καὶ τῷ Βαρνάβᾳ, οἵτινες προσλαλοῦντες αὐτοῖς ἔπειθον αὐτοὺς προσμένειν τῇ χάριτι τοῦ Θεοῦ.
\vs{44}Τῷ δὲ ἐρχομένῳ σαββάτῳ σχεδὸν πᾶσα ἡ πόλις συνήχθη ἀκοῦσαι τὸν λόγον τοῦ κυρίου.
\vs{45}ἰδόντες δὲ οἱ Ἰουδαῖοι τοὺς ὄχλους ἐπλήσθησαν ζήλου καὶ ἀντέλεγον τοῖς ὑπὸ Παύλου λαλουμένοις βλασφημοῦντες.
\vs{46}Παρρησιασάμενοί τε ὁ Παῦλος καὶ ὁ Βαρνάβας εἶπαν· Ὑμῖν ἦν ἀναγκαῖον πρῶτον λαληθῆναι τὸν λόγον τοῦ Θεοῦ· ἐπειδὴ ἀπωθεῖσθε αὐτὸν καὶ οὐκ ἀξίους κρίνετε ἑαυτοὺς τῆς αἰωνίου ζωῆς, ἰδοὺ στρεφόμεθα εἰς τὰ ἔθνη.
\vs{47}οὕτως γὰρ ἐντέταλται ἡμῖν ὁ Κύριος· 
\begin{poetryblock}

\begin{quote}Τέθεικά σε εἰς φῶς ἐθνῶν\end{quote} 

\begin{quote}τοῦ εἶναί σε εἰς σωτηρίαν ἕως ἐσχάτου τῆς γῆς.\end{quote}
\end{poetryblock}

\vs{48}Ἀκούοντα δὲ τὰ ἔθνη ἔχαιρον καὶ ἐδόξαζον τὸν λόγον τοῦ Κυρίου καὶ ἐπίστευσαν ὅσοι ἦσαν τεταγμένοι εἰς ζωὴν αἰώνιον·
\vs{49}διεφέρετο δὲ ὁ λόγος τοῦ Κυρίου δι᾽ ὅλης τῆς χώρας.
\vs{50}Οἱ δὲ Ἰουδαῖοι παρώτρυναν τὰς σεβομένας γυναῖκας τὰς εὐσχήμονας καὶ τοὺς πρώτους τῆς πόλεως καὶ ἐπήγειραν διωγμὸν ἐπὶ τὸν Παῦλον καὶ Βαρνάβαν καὶ ἐξέβαλον αὐτοὺς ἀπὸ τῶν ὁρίων αὐτῶν.
\vs{51}οἱ δὲ ἐκτιναξάμενοι τὸν κονιορτὸν τῶν ποδῶν ἐπ᾽ αὐτοὺς ἦλθον εἰς Ἰκόνιον,
\vs{52}οἵ τε μαθηταὶ ἐπληροῦντο χαρᾶς καὶ Πνεύματος Ἁγίου.

\ch{14}
Ἐγένετο δὲ ἐν Ἰκονίῳ κατὰ τὸ αὐτὸ εἰσελθεῖν αὐτοὺς εἰς τὴν συναγωγὴν τῶν Ἰουδαίων καὶ λαλῆσαι οὕτως ὥστε πιστεῦσαι Ἰουδαίων τε καὶ Ἑλλήνων πολὺ πλῆθος.
\vs{2}οἱ δὲ ἀπειθήσαντες Ἰουδαῖοι ἐπήγειραν καὶ ἐκάκωσαν τὰς ψυχὰς τῶν ἐθνῶν κατὰ τῶν ἀδελφῶν.
\vs{3}ἱκανὸν μὲν οὖν χρόνον διέτριψαν παρρησιαζόμενοι ἐπὶ τῷ Κυρίῳ τῷ μαρτυροῦντι ἐπὶ τῷ λόγῳ τῆς χάριτος αὐτοῦ, διδόντι σημεῖα καὶ τέρατα γίνεσθαι διὰ τῶν χειρῶν αὐτῶν.
\vs{4}Ἐσχίσθη δὲ τὸ πλῆθος τῆς πόλεως, καὶ οἱ μὲν ἦσαν σὺν τοῖς Ἰουδαίοις, οἱ δὲ σὺν τοῖς ἀποστόλοις.
\vs{5}ὡς δὲ ἐγένετο ὁρμὴ τῶν ἐθνῶν τε καὶ Ἰουδαίων σὺν τοῖς ἄρχουσιν αὐτῶν ὑβρίσαι καὶ λιθοβολῆσαι αὐτούς,
\vs{6}συνιδόντες κατέφυγον εἰς τὰς πόλεις τῆς Λυκαονίας Λύστραν καὶ Δέρβην καὶ τὴν περίχωρον,
\vs{7}κἀκεῖ εὐαγγελιζόμενοι ἦσαν.
\vs{8}Καί τις ἀνὴρ ἀδύνατος ἐν Λύστροις τοῖς ποσὶν ἐκάθητο, χωλὸς ἐκ κοιλίας μητρὸς αὐτοῦ ὃς οὐδέποτε περιεπάτησεν.
\vs{9}οὗτος ἤκουσεν τοῦ Παύλου λαλοῦντος· ὃς ἀτενίσας αὐτῷ καὶ ἰδὼν ὅτι ἔχει πίστιν τοῦ σωθῆναι,
\vs{10}εἶπεν μεγάλῃ φωνῇ· Ἀνάστηθι ἐπὶ τοὺς πόδας σου ὀρθός. καὶ ἥλατο καὶ περιεπάτει.
\vs{11}Οἵ τε ὄχλοι ἰδόντες ὃ ἐποίησεν Παῦλος ἐπῆραν τὴν φωνὴν αὐτῶν Λυκαονιστὶ λέγοντες· Οἱ θεοὶ ὁμοιωθέντες ἀνθρώποις κατέβησαν πρὸς ἡμᾶς,
\vs{12}ἐκάλουν τε τὸν Βαρνάβαν Δία, τὸν δὲ Παῦλον Ἑρμῆν, ἐπειδὴ αὐτὸς ἦν ὁ ἡγούμενος τοῦ λόγου.
\vs{13}ὅ τε ἱερεὺς τοῦ Διὸς τοῦ ὄντος πρὸ τῆς πόλεως ταύρους καὶ στέμματα ἐπὶ τοὺς πυλῶνας ἐνέγκας σὺν τοῖς ὄχλοις ἤθελεν θύειν.
\vs{14}Ἀκούσαντες δὲ οἱ ἀπόστολοι Βαρνάβας καὶ Παῦλος διαρρήξαντες τὰ ἱμάτια αὐτῶν ἐξεπήδησαν εἰς τὸν ὄχλον κράζοντες
\vs{15}καὶ λέγοντες· Ἄνδρες, τί ταῦτα ποιεῖτε; καὶ ἡμεῖς ὁμοιοπαθεῖς ἐσμεν ὑμῖν ἄνθρωποι εὐαγγελιζόμενοι ὑμᾶς ἀπὸ τούτων τῶν ματαίων ἐπιστρέφειν ἐπὶ θεὸν ζῶντα, ὃς ἐποίησεν τὸν οὐρανὸν καὶ τὴν γῆν καὶ τὴν θάλασσαν καὶ πάντα τὰ ἐν αὐτοῖς·
\vs{16}ὃς ἐν ταῖς παρῳχημέναις γενεαῖς εἴασεν πάντα τὰ ἔθνη πορεύεσθαι ταῖς ὁδοῖς αὐτῶν·
\vs{17}καίτοι οὐκ ἀμάρτυρον αὑτὸν ἀφῆκεν ἀγαθουργῶν, οὐρανόθεν ὑμῖν ὑετοὺς διδοὺς καὶ καιροὺς καρποφόρους, ἐμπιπλῶν τροφῆς καὶ εὐφροσύνης τὰς καρδίας ὑμῶν.
\vs{18}Καὶ ταῦτα λέγοντες μόλις κατέπαυσαν τοὺς ὄχλους τοῦ μὴ θύειν αὐτοῖς.
\vs{19}Ἐπῆλθαν δὲ ἀπὸ Ἀντιοχείας καὶ Ἰκονίου Ἰουδαῖοι καὶ πείσαντες τοὺς ὄχλους καὶ λιθάσαντες τὸν Παῦλον ἔσυρον ἔξω τῆς πόλεως νομίζοντες αὐτὸν τεθνηκέναι.
\vs{20}κυκλωσάντων δὲ τῶν μαθητῶν αὐτὸν ἀναστὰς εἰσῆλθεν εἰς τὴν πόλιν. Καὶ τῇ ἐπαύριον ἐξῆλθεν σὺν τῷ Βαρνάβᾳ εἰς Δέρβην.
\vs{21}Εὐαγγελισάμενοί τε τὴν πόλιν ἐκείνην καὶ μαθητεύσαντες ἱκανοὺς ὑπέστρεψαν εἰς τὴν Λύστραν καὶ εἰς Ἰκόνιον καὶ εἰς Ἀντιόχειαν
\vs{22}ἐπιστηρίζοντες τὰς ψυχὰς τῶν μαθητῶν, παρακαλοῦντες ἐμμένειν τῇ πίστει καὶ ὅτι Διὰ πολλῶν θλίψεων δεῖ ἡμᾶς εἰσελθεῖν εἰς τὴν βασιλείαν τοῦ Θεοῦ.
\vs{23}Χειροτονήσαντες δὲ αὐτοῖς κατ᾽ ἐκκλησίαν πρεσβυτέρους, προσευξάμενοι μετὰ νηστειῶν παρέθεντο αὐτοὺς τῷ Κυρίῳ εἰς ὃν πεπιστεύκεισαν.
\vs{24}Καὶ διελθόντες τὴν Πισιδίαν ἦλθον εἰς τὴν Παμφυλίαν
\vs{25}καὶ λαλήσαντες ἐν Πέργῃ τὸν λόγον κατέβησαν εἰς Ἀττάλειαν
\vs{26}Κἀκεῖθεν ἀπέπλευσαν εἰς Ἀντιόχειαν, ὅθεν ἦσαν παραδεδομένοι τῇ χάριτι τοῦ Θεοῦ εἰς τὸ ἔργον ὃ ἐπλήρωσαν.
\vs{27}Παραγενόμενοι δὲ καὶ συναγαγόντες τὴν ἐκκλησίαν ἀνήγγελλον ὅσα ἐποίησεν ὁ Θεὸς μετ᾽ αὐτῶν καὶ ὅτι ἤνοιξεν τοῖς ἔθνεσιν θύραν πίστεως.
\vs{28}διέτριβον δὲ χρόνον οὐκ ὀλίγον σὺν τοῖς μαθηταῖς.

\ch{15}
Καί τινες κατελθόντες ἀπὸ τῆς Ἰουδαίας ἐδίδασκον τοὺς ἀδελφοὺς ὅτι, Ἐὰν μὴ περιτμηθῆτε τῷ ἔθει τῷ Μωϋσέως, οὐ δύνασθε σωθῆναι.
\vs{2}γενομένης δὲ στάσεως καὶ ζητήσεως οὐκ ὀλίγης τῷ Παύλῳ καὶ τῷ Βαρνάβᾳ πρὸς αὐτοὺς, ἔταξαν ἀναβαίνειν Παῦλον καὶ Βαρνάβαν καί τινας ἄλλους ἐξ αὐτῶν πρὸς τοὺς ἀποστόλους καὶ πρεσβυτέρους εἰς Ἰερουσαλὴμ περὶ τοῦ ζητήματος τούτου.
\vs{3}Οἱ μὲν οὖν προπεμφθέντες ὑπὸ τῆς ἐκκλησίας διήρχοντο τήν τε Φοινίκην καὶ Σαμάρειαν ἐκδιηγούμενοι τὴν ἐπιστροφὴν τῶν ἐθνῶν καὶ ἐποίουν χαρὰν μεγάλην πᾶσιν τοῖς ἀδελφοῖς.
\vs{4}παραγενόμενοι δὲ εἰς Ἱεροσόλυμα παρεδέχθησαν ἀπὸ τῆς ἐκκλησίας καὶ τῶν ἀποστόλων καὶ τῶν πρεσβυτέρων, ἀνήγγειλάν τε ὅσα ὁ Θεὸς ἐποίησεν μετ᾽ αὐτῶν.
\vs{5}Ἐξανέστησαν δέ τινες τῶν ἀπὸ τῆς αἱρέσεως τῶν Φαρισαίων πεπιστευκότες λέγοντες ὅτι Δεῖ περιτέμνειν αὐτοὺς παραγγέλλειν τε τηρεῖν τὸν νόμον Μωϋσέως.

\vs{6}Συνήχθησάν τε οἱ ἀπόστολοι καὶ οἱ πρεσβύτεροι ἰδεῖν περὶ τοῦ λόγου τούτου.
\vs{7}Πολλῆς δὲ ζητήσεως γενομένης ἀναστὰς Πέτρος εἶπεν πρὸς αὐτούς· Ἄνδρες ἀδελφοί, ὑμεῖς ἐπίστασθε ὅτι ἀφ᾽ ἡμερῶν ἀρχαίων ἐν ὑμῖν ἐξελέξατο ὁ Θεὸς διὰ τοῦ στόματός μου ἀκοῦσαι τὰ ἔθνη τὸν λόγον τοῦ εὐαγγελίου καὶ πιστεῦσαι.
\vs{8}καὶ ὁ καρδιογνώστης Θεὸς ἐμαρτύρησεν αὐτοῖς δοὺς τὸ Πνεῦμα τὸ Ἅγιον καθὼς καὶ ἡμῖν
\vs{9}καὶ οὐθὲν διέκρινεν μεταξὺ ἡμῶν τε καὶ αὐτῶν τῇ πίστει καθαρίσας τὰς καρδίας αὐτῶν.
\vs{10}Νῦν οὖν τί πειράζετε τὸν Θεόν ἐπιθεῖναι ζυγὸν ἐπὶ τὸν τράχηλον τῶν μαθητῶν ὃν οὔτε οἱ πατέρες ἡμῶν οὔτε ἡμεῖς ἰσχύσαμεν βαστάσαι;
\vs{11}ἀλλὰ διὰ τῆς χάριτος τοῦ Κυρίου Ἰησοῦ πιστεύομεν σωθῆναι καθ᾽ ὃν τρόπον κἀκεῖνοι.
\vs{12}Ἐσίγησεν δὲ πᾶν τὸ πλῆθος καὶ ἤκουον Βαρνάβα καὶ Παύλου ἐξηγουμένων ὅσα ἐποίησεν ὁ Θεὸς σημεῖα καὶ τέρατα ἐν τοῖς ἔθνεσιν δι᾽ αὐτῶν.
\vs{13}Μετὰ δὲ τὸ σιγῆσαι αὐτοὺς ἀπεκρίθη Ἰάκωβος λέγων· Ἄνδρες ἀδελφοί, ἀκούσατέ μου.
\vs{14}Συμεὼν ἐξηγήσατο καθὼς πρῶτον ὁ Θεὸς ἐπεσκέψατο λαβεῖν ἐξ ἐθνῶν λαὸν τῷ ὀνόματι αὐτοῦ.
\vs{15}καὶ τούτῳ συμφωνοῦσιν οἱ λόγοι τῶν προφητῶν καθὼς γέγραπται·
\begin{poetryblock}

\begin{quote} \vs{16}Μετὰ ταῦτα ἀναστρέψω\end{quote} 

\begin{quote}καὶ ἀνοικοδομήσω τὴν σκηνὴν Δαυὶδ τὴν πεπτωκυῖαν\end{quote} 

\begin{quote}καὶ τὰ κατεσκαμμένα αὐτῆς ἀνοικοδομήσω\end{quote} 

\begin{quote}καὶ ἀνορθώσω αὐτήν,\end{quote}

\begin{quote} \vs{17}ὅπως ἂν ἐκζητήσωσιν οἱ κατάλοιποι τῶν ἀνθρώπων τὸν Κύριον\end{quote} 

\begin{quote}καὶ πάντα τὰ ἔθνη ἐφ᾽ οὓς ἐπικέκληται τὸ ὄνομά μου ἐπ᾽ αὐτούς,\end{quote} 

\begin{quote}λέγει Κύριος ποιῶν ταῦτα\end{quote}
\end{poetryblock}

\vs{18}γνωστὰ ἀπ᾽ αἰῶνος.

\vs{19}Διὸ ἐγὼ κρίνω μὴ παρενοχλεῖν τοῖς ἀπὸ τῶν ἐθνῶν ἐπιστρέφουσιν ἐπὶ τὸν Θεόν,
\vs{20}ἀλλὰ ἐπιστεῖλαι αὐτοῖς τοῦ ἀπέχεσθαι τῶν ἀλισγημάτων τῶν εἰδώλων καὶ τῆς πορνείας καὶ τοῦ πνικτοῦ καὶ τοῦ αἵματος.
\vs{21}Μωϋσῆς γὰρ ἐκ γενεῶν ἀρχαίων κατὰ πόλιν τοὺς κηρύσσοντας αὐτὸν ἔχει ἐν ταῖς συναγωγαῖς κατὰ πᾶν σάββατον ἀναγινωσκόμενος.

\vs{22}Τότε ἔδοξε τοῖς ἀποστόλοις καὶ τοῖς πρεσβυτέροις σὺν ὅλῃ τῇ ἐκκλησίᾳ ἐκλεξαμένους ἄνδρας ἐξ αὐτῶν πέμψαι εἰς Ἀντιόχειαν σὺν τῷ Παύλῳ καὶ Βαρνάβᾳ, Ἰούδαν τὸν καλούμενον Βαρσαββᾶν καὶ Σιλᾶν, ἄνδρας ἡγουμένους ἐν τοῖς ἀδελφοῖς,
\vs{23}γράψαντες διὰ χειρὸς αὐτῶν·

Οἱ ἀπόστολοι καὶ οἱ πρεσβύτεροι ἀδελφοὶ Τοῖς κατὰ τὴν Ἀντιόχειαν καὶ Συρίαν καὶ Κιλικίαν ἀδελφοῖς τοῖς ἐξ ἐθνῶν Χαίρειν.
\vs{24}Ἐπειδὴ ἠκούσαμεν ὅτι τινὲς ἐξ ἡμῶν ἐξελθόντες ἐτάραξαν ὑμᾶς λόγοις ἀνασκευάζοντες τὰς ψυχὰς ὑμῶν οἷς οὐ διεστειλάμεθα,
\vs{25}ἔδοξεν ἡμῖν γενομένοις ὁμοθυμαδὸν ἐκλεξαμένοις ἄνδρας πέμψαι πρὸς ὑμᾶς σὺν τοῖς ἀγαπητοῖς ἡμῶν Βαρνάβᾳ καὶ Παύλῳ,
\vs{26}ἀνθρώποις παραδεδωκόσι τὰς ψυχὰς αὐτῶν ὑπὲρ τοῦ ὀνόματος τοῦ Κυρίου ἡμῶν Ἰησοῦ Χριστοῦ.
\vs{27}ἀπεστάλκαμεν οὖν Ἰούδαν καὶ Σιλᾶν καὶ αὐτοὺς διὰ λόγου ἀπαγγέλλοντας τὰ αὐτά.
\vs{28}Ἔδοξεν γὰρ τῷ Πνεύματι τῷ Ἁγίῳ καὶ ἡμῖν μηδὲν πλέον ἐπιτίθεσθαι ὑμῖν βάρος πλὴν τούτων τῶν ἐπάναγκες,
\vs{29}ἀπέχεσθαι εἰδωλοθύτων καὶ αἵματος καὶ πνικτῶν καὶ πορνείας, ἐξ ὧν διατηροῦντες ἑαυτοὺς εὖ πράξετε. Ἔρρωσθε.

\vs{30}Οἱ μὲν οὖν ἀπολυθέντες κατῆλθον εἰς Ἀντιόχειαν, καὶ συναγαγόντες τὸ πλῆθος ἐπέδωκαν τὴν ἐπιστολήν.
\vs{31}ἀναγνόντες δὲ ἐχάρησαν ἐπὶ τῇ παρακλήσει.
\vs{32}Ἰούδας τε καὶ Σιλᾶς καὶ αὐτοὶ προφῆται ὄντες διὰ λόγου πολλοῦ παρεκάλεσαν τοὺς ἀδελφοὺς καὶ ἐπεστήριξαν,
\vs{33}ποιήσαντες δὲ χρόνον ἀπελύθησαν μετ᾽ εἰρήνης ἀπὸ τῶν ἀδελφῶν πρὸς τοὺς ἀποστείλαντας αὐτούς.
\vs{35}Παῦλος δὲ καὶ Βαρνάβας διέτριβον ἐν Ἀντιοχείᾳ διδάσκοντες καὶ εὐαγγελιζόμενοι μετὰ καὶ ἑτέρων πολλῶν τὸν λόγον τοῦ Κυρίου.

\vs{36}Μετὰ δέ τινας ἡμέρας εἶπεν πρὸς Βαρνάβαν Παῦλος· Ἐπιστρέψαντες δὴ ἐπισκεψώμεθα τοὺς ἀδελφοὺς κατὰ πόλιν πᾶσαν ἐν αἷς κατηγγείλαμεν τὸν λόγον τοῦ Κυρίου πῶς ἔχουσιν.
\vs{37}Βαρνάβας δὲ ἐβούλετο συμπαραλαβεῖν καὶ τὸν Ἰωάννην τὸν καλούμενον Μάρκον·
\vs{38}Παῦλος δὲ ἠξίου, τὸν ἀποστάντα ἀπ᾽ αὐτῶν ἀπὸ Παμφυλίας καὶ μὴ συνελθόντα αὐτοῖς εἰς τὸ ἔργον μὴ συμπαραλαμβάνειν τοῦτον.
\vs{39}Ἐγένετο δὲ παροξυσμός ὥστε ἀποχωρισθῆναι αὐτοὺς ἀπ᾽ ἀλλήλων, τόν τε Βαρνάβαν παραλαβόντα τὸν Μάρκον ἐκπλεῦσαι εἰς Κύπρον,
\vs{40}Παῦλος δὲ ἐπιλεξάμενος Σιλᾶν ἐξῆλθεν παραδοθεὶς τῇ χάριτι τοῦ Κυρίου ὑπὸ τῶν ἀδελφῶν.
\vs{41}διήρχετο δὲ τὴν Συρίαν καὶ τὴν Κιλικίαν ἐπιστηρίζων τὰς ἐκκλησίας.

\ch{16}
Κατήντησεν δὲ καὶ εἰς Δέρβην καὶ εἰς Λύστραν. καὶ ἰδοὺ μαθητής τις ἦν ἐκεῖ ὀνόματι Τιμόθεος, υἱὸς γυναικὸς Ἰουδαίας πιστῆς, πατρὸς δὲ Ἕλληνος,
\vs{2}ὃς ἐμαρτυρεῖτο ὑπὸ τῶν ἐν Λύστροις καὶ Ἰκονίῳ ἀδελφῶν.
\vs{3}τοῦτον ἠθέλησεν ὁ Παῦλος σὺν αὐτῷ ἐξελθεῖν, καὶ λαβὼν περιέτεμεν αὐτὸν διὰ τοὺς Ἰουδαίους τοὺς ὄντας ἐν τοῖς τόποις ἐκείνοις· ᾔδεισαν γὰρ ἅπαντες ὅτι Ἕλλην ὁ πατὴρ αὐτοῦ ὑπῆρχεν.
\vs{4}Ὡς δὲ διεπορεύοντο τὰς πόλεις, παρεδίδοσαν αὐτοῖς φυλάσσειν τὰ δόγματα τὰ κεκριμένα ὑπὸ τῶν ἀποστόλων καὶ πρεσβυτέρων τῶν ἐν Ἱεροσολύμοις.
\vs{5}Αἱ μὲν οὖν ἐκκλησίαι ἐστερεοῦντο τῇ πίστει καὶ ἐπερίσσευον τῷ ἀριθμῷ καθ᾽ ἡμέραν.

\vs{6}Διῆλθον δὲ τὴν Φρυγίαν καὶ Γαλατικὴν χώραν κωλυθέντες ὑπὸ τοῦ Ἁγίου Πνεύματος λαλῆσαι τὸν λόγον ἐν τῇ Ἀσίᾳ·
\vs{7}ἐλθόντες δὲ κατὰ τὴν Μυσίαν ἐπείραζον εἰς τὴν Βιθυνίαν πορευθῆναι, καὶ οὐκ εἴασεν αὐτοὺς τὸ Πνεῦμα Ἰησοῦ·
\vs{8}παρελθόντες δὲ τὴν Μυσίαν κατέβησαν εἰς Τρῳάδα.
\vs{9}Καὶ ὅραμα διὰ τῆς νυκτὸς τῷ Παύλῳ ὤφθη, ἀνὴρ Μακεδών τις ἦν ἑστὼς καὶ παρακαλῶν αὐτὸν καὶ λέγων· Διαβὰς εἰς Μακεδονίαν βοήθησον ἡμῖν.
\vs{10}ὡς δὲ τὸ ὅραμα εἶδεν, εὐθέως ἐζητήσαμεν ἐξελθεῖν εἰς Μακεδονίαν συμβιβάζοντες ὅτι προσκέκληται ἡμᾶς ὁ Θεὸς εὐαγγελίσασθαι αὐτούς.

\vs{11}Ἀναχθέντες δὲ ἀπὸ Τρῳάδος εὐθυδρομήσαμεν εἰς Σαμοθρᾴκην, τῇ δὲ ἐπιούσῃ εἰς Νέαν Πόλιν
\vs{12}κἀκεῖθεν εἰς Φιλίππους, ἥτις ἐστὶν πρώτη μερίδος τῆς Μακεδονίας πόλις, κολωνία. Ἦμεν δὲ ἐν ταύτῃ τῇ πόλει διατρίβοντες ἡμέρας τινάς.
\vs{13}Τῇ τε ἡμέρᾳ τῶν σαββάτων ἐξήλθομεν ἔξω τῆς πύλης παρὰ ποταμὸν οὗ ἐνομίζομεν προσευχὴν εἶναι, καὶ καθίσαντες ἐλαλοῦμεν ταῖς συνελθούσαις γυναιξίν.
\vs{14}Καί τις γυνὴ ὀνόματι Λυδία, πορφυρόπωλις πόλεως Θυατείρων σεβομένη τὸν Θεόν, ἤκουεν, ἧς ὁ Κύριος διήνοιξεν τὴν καρδίαν προσέχειν τοῖς λαλουμένοις ὑπὸ τοῦ Παύλου.
\vs{15}ὡς δὲ ἐβαπτίσθη καὶ ὁ οἶκος αὐτῆς, παρεκάλεσεν λέγουσα· Εἰ κεκρίκατέ με πιστὴν τῷ Κυρίῳ εἶναι, εἰσελθόντες εἰς τὸν οἶκόν μου μένετε· καὶ παρεβιάσατο ἡμᾶς.

\vs{16}Ἐγένετο δὲ πορευομένων ἡμῶν εἰς τὴν προσευχὴν παιδίσκην τινὰ ἔχουσαν πνεῦμα Πύθωνα ὑπαντῆσαι ἡμῖν, ἥτις ἐργασίαν πολλὴν παρεῖχεν τοῖς κυρίοις αὐτῆς μαντευομένη.
\vs{17}αὕτη κατακολουθοῦσα τῷ Παύλῳ καὶ ἡμῖν ἔκραζεν λέγουσα· Οὗτοι οἱ ἄνθρωποι δοῦλοι τοῦ Θεοῦ τοῦ Ὑψίστου εἰσίν, οἵτινες καταγγέλλουσιν ὑμῖν ὁδὸν σωτηρίας.
\vs{18}Τοῦτο δὲ ἐποίει ἐπὶ πολλὰς ἡμέρας. διαπονηθεὶς δὲ Παῦλος καὶ ἐπιστρέψας τῷ πνεύματι εἶπεν· Παραγγέλλω σοι ἐν ὀνόματι Ἰησοῦ Χριστοῦ ἐξελθεῖν ἀπ᾽ αὐτῆς· καὶ ἐξῆλθεν αὐτῇ τῇ ὥρᾳ.
\vs{19}Ἰδόντες δὲ οἱ κύριοι αὐτῆς ὅτι ἐξῆλθεν ἡ ἐλπὶς τῆς ἐργασίας αὐτῶν, ἐπιλαβόμενοι τὸν Παῦλον καὶ τὸν Σιλᾶν εἵλκυσαν εἰς τὴν ἀγορὰν ἐπὶ τοὺς ἄρχοντας
\vs{20}καὶ προσαγαγόντες αὐτοὺς τοῖς στρατηγοῖς εἶπαν· Οὗτοι οἱ ἄνθρωποι ἐκταράσσουσιν ἡμῶν τὴν πόλιν, Ἰουδαῖοι ὑπάρχοντες,
\vs{21}καὶ καταγγέλλουσιν ἔθη ἃ οὐκ ἔξεστιν ἡμῖν παραδέχεσθαι οὐδὲ ποιεῖν Ῥωμαίοις οὖσιν.
\vs{22}Καὶ συνεπέστη ὁ ὄχλος κατ᾽ αὐτῶν καὶ οἱ στρατηγοὶ περιρήξαντες αὐτῶν τὰ ἱμάτια ἐκέλευον ῥαβδίζειν,
\vs{23}πολλάς τε ἐπιθέντες αὐτοῖς πληγὰς ἔβαλον εἰς φυλακήν παραγγείλαντες τῷ δεσμοφύλακι ἀσφαλῶς τηρεῖν αὐτούς.
\vs{24}ὃς παραγγελίαν τοιαύτην λαβὼν ἔβαλεν αὐτοὺς εἰς τὴν ἐσωτέραν φυλακὴν καὶ τοὺς πόδας ἠσφαλίσατο αὐτῶν εἰς τὸ ξύλον.
\vs{25}Κατὰ δὲ τὸ μεσονύκτιον Παῦλος καὶ Σιλᾶς προσευχόμενοι ὕμνουν τὸν Θεόν, ἐπηκροῶντο δὲ αὐτῶν οἱ δέσμιοι.
\vs{26}ἄφνω δὲ σεισμὸς ἐγένετο μέγας ὥστε σαλευθῆναι τὰ θεμέλια τοῦ δεσμωτηρίου· ἠνεῴχθησαν δὲ παραχρῆμα αἱ θύραι πᾶσαι καὶ πάντων τὰ δεσμὰ ἀνέθη.
\vs{27}Ἔξυπνος δὲ γενόμενος ὁ δεσμοφύλαξ καὶ ἰδὼν ἀνεῳγμένας τὰς θύρας τῆς φυλακῆς, σπασάμενος τὴν μάχαιραν ἤμελλεν ἑαυτὸν ἀναιρεῖν νομίζων ἐκπεφευγέναι τοὺς δεσμίους.
\vs{28}ἐφώνησεν δὲ ὁ Παῦλος μεγάλῃ φωνῇ λέγων· Μηδὲν πράξῃς σεαυτῷ κακόν, ἅπαντες γάρ ἐσμεν ἐνθάδε.
\vs{29}Αἰτήσας δὲ φῶτα εἰσεπήδησεν καὶ ἔντρομος γενόμενος προσέπεσεν τῷ Παύλῳ καὶ τῷ Σιλᾷ
\vs{30}καὶ προαγαγὼν αὐτοὺς ἔξω ἔφη· Κύριοι, τί με δεῖ ποιεῖν ἵνα σωθῶ;
\vs{31}Οἱ δὲ εἶπαν· Πίστευσον ἐπὶ τὸν Κύριον Ἰησοῦν καὶ σωθήσῃ σὺ καὶ ὁ οἶκός σου.
\vs{32}καὶ ἐλάλησαν αὐτῷ τὸν λόγον τοῦ κυρίου σὺν πᾶσιν τοῖς ἐν τῇ οἰκίᾳ αὐτοῦ.
\vs{33}καὶ παραλαβὼν αὐτοὺς ἐν ἐκείνῃ τῇ ὥρᾳ τῆς νυκτὸς ἔλουσεν ἀπὸ τῶν πληγῶν, καὶ ἐβαπτίσθη αὐτὸς καὶ οἱ αὐτοῦ πάντες παραχρῆμα,
\vs{34}ἀναγαγών τε αὐτοὺς εἰς τὸν οἶκον παρέθηκεν τράπεζαν καὶ ἠγαλλιάσατο πανοικεὶ πεπιστευκὼς τῷ Θεῷ.
\vs{35}Ἡμέρας δὲ γενομένης ἀπέστειλαν οἱ στρατηγοὶ τοὺς ῥαβδούχους λέγοντες· Ἀπόλυσον τοὺς ἀνθρώπους ἐκείνους.
\vs{36}Ἀπήγγειλεν δὲ ὁ δεσμοφύλαξ τοὺς λόγους τούτους πρὸς τὸν Παῦλον ὅτι Ἀπέσταλκαν οἱ στρατηγοὶ ἵνα ἀπολυθῆτε· νῦν οὖν ἐξελθόντες πορεύεσθε ἐν εἰρήνῃ.
\vs{37}Ὁ δὲ Παῦλος ἔφη πρὸς αὐτούς· Δείραντες ἡμᾶς δημοσίᾳ ἀκατακρίτους, ἀνθρώπους Ῥωμαίους ὑπάρχοντας, ἔβαλαν εἰς φυλακήν, καὶ νῦν λάθρᾳ ἡμᾶς ἐκβάλλουσιν; οὐ γάρ, ἀλλὰ ἐλθόντες αὐτοὶ ἡμᾶς ἐξαγαγέτωσαν.
\vs{38}Ἀπήγγειλαν δὲ τοῖς στρατηγοῖς οἱ ῥαβδοῦχοι τὰ ῥήματα ταῦτα. ἐφοβήθησαν δὲ ἀκούσαντες ὅτι Ῥωμαῖοί εἰσιν,
\vs{39}καὶ ἐλθόντες παρεκάλεσαν αὐτούς καὶ ἐξαγαγόντες ἠρώτων ἀπελθεῖν ἀπὸ τῆς πόλεως.
\vs{40}ἐξελθόντες δὲ ἀπὸ τῆς φυλακῆς εἰσῆλθον πρὸς τὴν Λυδίαν καὶ ἰδόντες παρεκάλεσαν τοὺς ἀδελφοὺς καὶ ἐξῆλθαν.

\ch{17}
Διοδεύσαντες δὲ τὴν Ἀμφίπολιν καὶ τὴν Ἀπολλωνίαν ἦλθον εἰς Θεσσαλονίκην ὅπου ἦν συναγωγὴ τῶν Ἰουδαίων.
\vs{2}κατὰ δὲ τὸ εἰωθὸς τῷ Παύλῳ εἰσῆλθεν πρὸς αὐτοὺς καὶ ἐπὶ σάββατα τρία διελέξατο αὐτοῖς ἀπὸ τῶν γραφῶν,
\vs{3}διανοίγων καὶ παρατιθέμενος ὅτι τὸν Χριστὸν ἔδει παθεῖν καὶ ἀναστῆναι ἐκ νεκρῶν καὶ ὅτι Οὗτός ἐστιν ὁ Χριστός ὁ Ἰησοῦς ὃν ἐγὼ καταγγέλλω ὑμῖν.
\vs{4}καί τινες ἐξ αὐτῶν ἐπείσθησαν καὶ προσεκληρώθησαν τῷ Παύλῳ καὶ τῷ Σιλᾷ, τῶν τε σεβομένων Ἑλλήνων πλῆθος πολὺ, γυναικῶν τε τῶν πρώτων οὐκ ὀλίγαι.
\vs{5}Ζηλώσαντες δὲ οἱ Ἰουδαῖοι καὶ προσλαβόμενοι τῶν ἀγοραίων ἄνδρας τινὰς πονηροὺς καὶ ὀχλοποιήσαντες ἐθορύβουν τὴν πόλιν καὶ ἐπιστάντες τῇ οἰκίᾳ Ἰάσονος ἐζήτουν αὐτοὺς προαγαγεῖν εἰς τὸν δῆμον·
\vs{6}μὴ εὑρόντες δὲ αὐτοὺς ἔσυρον Ἰάσονα καί τινας ἀδελφοὺς ἐπὶ τοὺς πολιτάρχας βοῶντες ὅτι Οἱ τὴν οἰκουμένην ἀναστατώσαντες οὗτοι καὶ ἐνθάδε πάρεισιν,
\vs{7}οὓς ὑποδέδεκται Ἰάσων· καὶ οὗτοι πάντες ἀπέναντι τῶν δογμάτων Καίσαρος πράσσουσι βασιλέα ἕτερον λέγοντες εἶναι Ἰησοῦν.
\vs{8}Ἐτάραξαν δὲ τὸν ὄχλον καὶ τοὺς πολιτάρχας ἀκούοντας ταῦτα,
\vs{9}καὶ λαβόντες τὸ ἱκανὸν παρὰ τοῦ Ἰάσονος καὶ τῶν λοιπῶν ἀπέλυσαν αὐτούς.

\vs{10}Οἱ δὲ ἀδελφοὶ εὐθέως διὰ νυκτὸς ἐξέπεμψαν τόν τε Παῦλον καὶ τὸν Σιλᾶν εἰς Βέροιαν, οἵτινες παραγενόμενοι εἰς τὴν συναγωγὴν τῶν Ἰουδαίων ἀπῄεσαν.
\vs{11}οὗτοι δὲ ἦσαν εὐγενέστεροι τῶν ἐν Θεσσαλονίκῃ, οἵτινες ἐδέξαντο τὸν λόγον μετὰ πάσης προθυμίας καθ᾽ ἡμέραν ἀνακρίνοντες τὰς γραφὰς εἰ ἔχοι ταῦτα οὕτως.
\vs{12}πολλοὶ μὲν οὖν ἐξ αὐτῶν ἐπίστευσαν καὶ τῶν Ἑλληνίδων γυναικῶν τῶν εὐσχημόνων καὶ ἀνδρῶν οὐκ ὀλίγοι.
\vs{13}Ὡς δὲ ἔγνωσαν οἱ ἀπὸ τῆς Θεσσαλονίκης Ἰουδαῖοι ὅτι καὶ ἐν τῇ Βεροίᾳ κατηγγέλη ὑπὸ τοῦ Παύλου ὁ λόγος τοῦ Θεοῦ, ἦλθον κἀκεῖ σαλεύοντες καὶ ταράσσοντες τοὺς ὄχλους.
\vs{14}εὐθέως δὲ τότε τὸν Παῦλον ἐξαπέστειλαν οἱ ἀδελφοὶ πορεύεσθαι ἕως ἐπὶ τὴν θάλασσαν, ὑπέμεινάν τε ὅ τε Σιλᾶς καὶ ὁ Τιμόθεος ἐκεῖ.
\vs{15}οἱ δὲ καθιστάνοντες τὸν Παῦλον ἤγαγον ἕως Ἀθηνῶν, καὶ λαβόντες ἐντολὴν πρὸς τὸν Σιλᾶν καὶ τὸν Τιμόθεον ἵνα ὡς τάχιστα ἔλθωσιν πρὸς αὐτὸν ἐξῄεσαν.

\vs{16}Ἐν δὲ ταῖς Ἀθήναις ἐκδεχομένου αὐτοὺς τοῦ Παύλου παρωξύνετο τὸ πνεῦμα αὐτοῦ ἐν αὐτῷ θεωροῦντος κατείδωλον οὖσαν τὴν πόλιν.
\vs{17}διελέγετο μὲν οὖν ἐν τῇ συναγωγῇ τοῖς Ἰουδαίοις καὶ τοῖς σεβομένοις καὶ ἐν τῇ ἀγορᾷ κατὰ πᾶσαν ἡμέραν πρὸς τοὺς παρατυγχάνοντας.
\vs{18}Τινὲς δὲ καὶ τῶν Ἐπικουρείων καὶ Στοϊκῶν φιλοσόφων συνέβαλλον αὐτῷ, καί τινες ἔλεγον· Τί ἂν θέλοι ὁ σπερμολόγος οὗτος λέγειν; οἱ δέ· Ξένων δαιμονίων δοκεῖ καταγγελεὺς εἶναι, ὅτι τὸν Ἰησοῦν καὶ τὴν ἀνάστασιν εὐηγγελίζετο.
\vs{19}Ἐπιλαβόμενοί τε αὐτοῦ ἐπὶ τὸν Ἄρειον πάγον ἤγαγον λέγοντες· Δυνάμεθα γνῶναι τίς ἡ καινὴ αὕτη ἡ ὑπὸ σοῦ λαλουμένη διδαχή;
\vs{20}ξενίζοντα γάρ τινα εἰσφέρεις εἰς τὰς ἀκοὰς ἡμῶν· βουλόμεθα οὖν γνῶναι τίνα θέλει ταῦτα εἶναι.
\vs{21}Ἀθηναῖοι δὲ πάντες καὶ οἱ ἐπιδημοῦντες ξένοι εἰς οὐδὲν ἕτερον ηὐκαίρουν ἢ λέγειν τι ἢ ἀκούειν τι καινότερον.

\vs{22}Σταθεὶς δὲ ὁ Παῦλος ἐν μέσῳ τοῦ Ἀρείου Πάγου ἔφη· Ἄνδρες Ἀθηναῖοι, κατὰ πάντα ὡς δεισιδαιμονεστέρους ὑμᾶς θεωρῶ.
\vs{23}διερχόμενος γὰρ καὶ ἀναθεωρῶν τὰ σεβάσματα ὑμῶν εὗρον καὶ βωμὸν ἐν ᾧ ἐπεγέγραπτο· 
\begin{poetryblock}

\begin{quote}ΑΓΝΩΣΤΩ ΘΕΩ.\end{quote}
\end{poetryblock}

Ὃ οὖν ἀγνοοῦντες εὐσεβεῖτε, τοῦτο ἐγὼ καταγγέλλω ὑμῖν.
\vs{24}Ὁ Θεὸς ὁ ποιήσας τὸν κόσμον καὶ πάντα τὰ ἐν αὐτῷ, οὗτος οὐρανοῦ καὶ γῆς ὑπάρχων Κύριος οὐκ ἐν χειροποιήτοις ναοῖς κατοικεῖ
\vs{25}οὐδὲ ὑπὸ χειρῶν ἀνθρωπίνων θεραπεύεται προσδεόμενός τινος, αὐτὸς διδοὺς πᾶσι ζωὴν καὶ πνοὴν καὶ τὰ πάντα·
\vs{26}ἐποίησέν τε ἐξ ἑνὸς πᾶν ἔθνος ἀνθρώπων κατοικεῖν ἐπὶ παντὸς προσώπου τῆς γῆς, ὁρίσας προστεταγμένους καιροὺς καὶ τὰς ὁροθεσίας τῆς κατοικίας αὐτῶν
\vs{27}ζητεῖν τὸν Θεὸν, εἰ ἄρα γε ψηλαφήσειαν αὐτὸν καὶ εὕροιεν, καί γε οὐ μακρὰν ἀπὸ ἑνὸς ἑκάστου ἡμῶν ὑπάρχοντα.

\vs{28}Ἐν αὐτῷ γὰρ ζῶμεν καὶ κινούμεθα καὶ ἐσμέν, ὡς καί τινες τῶν καθ᾽ ὑμᾶς ποιητῶν εἰρήκασιν· 
\begin{poetryblock}

\begin{quote}Τοῦ γὰρ καὶ γένος ἐσμέν.\end{quote}
\end{poetryblock}

\vs{29}γένος οὖν ὑπάρχοντες τοῦ Θεοῦ οὐκ ὀφείλομεν νομίζειν χρυσῷ ἢ ἀργύρῳ ἢ λίθῳ, χαράγματι τέχνης καὶ ἐνθυμήσεως ἀνθρώπου, τὸ Θεῖον εἶναι ὅμοιον.
\vs{30}Τοὺς μὲν οὖν χρόνους τῆς ἀγνοίας ὑπεριδὼν ὁ Θεὸς, τὰ νῦν παραγγέλλει τοῖς ἀνθρώποις πάντας πανταχοῦ μετανοεῖν,
\vs{31}καθότι ἔστησεν ἡμέραν ἐν ᾗ μέλλει κρίνειν τὴν οἰκουμένην ἐν δικαιοσύνῃ, ἐν ἀνδρὶ ᾧ ὥρισεν, πίστιν παρασχὼν πᾶσιν ἀναστήσας αὐτὸν ἐκ νεκρῶν.
\vs{32}Ἀκούσαντες δὲ ἀνάστασιν νεκρῶν οἱ μὲν ἐχλεύαζον, οἱ δὲ εἶπαν· Ἀκουσόμεθά σου περὶ τούτου καὶ πάλιν.
\vs{33}οὕτως ὁ Παῦλος ἐξῆλθεν ἐκ μέσου αὐτῶν.
\vs{34}τινὲς δὲ ἄνδρες κολληθέντες αὐτῷ ἐπίστευσαν, ἐν οἷς καὶ Διονύσιος ὁ Ἀρεοπαγίτης καὶ γυνὴ ὀνόματι Δάμαρις καὶ ἕτεροι σὺν αὐτοῖς.

\ch{18}
Μετὰ ταῦτα χωρισθεὶς ἐκ τῶν Ἀθηνῶν ἦλθεν εἰς Κόρινθον.
\vs{2}καὶ εὑρών τινα Ἰουδαῖον ὀνόματι Ἀκύλαν, Ποντικὸν τῷ γένει προσφάτως ἐληλυθότα ἀπὸ τῆς Ἰταλίας καὶ Πρίσκιλλαν γυναῖκα αὐτοῦ, διὰ τὸ διατεταχέναι Κλαύδιον χωρίζεσθαι πάντας τοὺς Ἰουδαίους ἀπὸ τῆς Ῥώμης, προσῆλθεν αὐτοῖς
\vs{3}καὶ διὰ τὸ ὁμότεχνον εἶναι ἔμενεν παρ᾽ αὐτοῖς, καὶ ἠργάζετο· ἦσαν γὰρ σκηνοποιοὶ τῇ τέχνῃ.
\vs{4}Διελέγετο δὲ ἐν τῇ συναγωγῇ κατὰ πᾶν σάββατον ἔπειθέν τε Ἰουδαίους καὶ Ἕλληνας.
\vs{5}Ὡς δὲ κατῆλθον ἀπὸ τῆς Μακεδονίας ὅ τε Σιλᾶς καὶ ὁ Τιμόθεος, συνείχετο τῷ λόγῳ ὁ Παῦλος διαμαρτυρόμενος τοῖς Ἰουδαίοις εἶναι τὸν Χριστὸν Ἰησοῦν.
\vs{6}ἀντιτασσομένων δὲ αὐτῶν καὶ βλασφημούντων ἐκτιναξάμενος τὰ ἱμάτια εἶπεν πρὸς αὐτούς· Τὸ αἷμα ὑμῶν ἐπὶ τὴν κεφαλὴν ὑμῶν· καθαρὸς ἐγώ ἀπὸ τοῦ νῦν εἰς τὰ ἔθνη πορεύσομαι.
\vs{7}Καὶ μεταβὰς ἐκεῖθεν εἰσῆλθεν εἰς οἰκίαν τινὸς ὀνόματι Τιτίου Ἰούστου σεβομένου τὸν Θεόν, οὗ ἡ οἰκία ἦν συνομοροῦσα τῇ συναγωγῇ.
\vs{8}Κρίσπος δὲ ὁ ἀρχισυνάγωγος ἐπίστευσεν τῷ Κυρίῳ σὺν ὅλῳ τῷ οἴκῳ αὐτοῦ, καὶ πολλοὶ τῶν Κορινθίων ἀκούοντες ἐπίστευον καὶ ἐβαπτίζοντο.
\vs{9}Εἶπεν δὲ ὁ Κύριος ἐν νυκτὶ δι᾽ ὁράματος τῷ Παύλῳ· Μὴ φοβοῦ, ἀλλὰ λάλει καὶ μὴ σιωπήσῃς,
\vs{10}διότι ἐγώ εἰμι μετὰ σοῦ καὶ οὐδεὶς ἐπιθήσεταί σοι τοῦ κακῶσαί σε, διότι λαός ἐστί μοι πολὺς ἐν τῇ πόλει ταύτῃ.
\vs{11}Ἐκάθισεν δὲ ἐνιαυτὸν καὶ μῆνας ἓξ διδάσκων ἐν αὐτοῖς τὸν λόγον τοῦ Θεοῦ.

\vs{12}Γαλλίωνος δὲ ἀνθυπάτου ὄντος τῆς Ἀχαΐας κατεπέστησαν ὁμοθυμαδὸν οἱ Ἰουδαῖοι τῷ Παύλῳ καὶ ἤγαγον αὐτὸν ἐπὶ τὸ βῆμα
\vs{13}λέγοντες ὅτι Παρὰ τὸν νόμον ἀναπείθει οὗτος τοὺς ἀνθρώπους σέβεσθαι τὸν Θεόν.
\vs{14}Μέλλοντος δὲ τοῦ Παύλου ἀνοίγειν τὸ στόμα εἶπεν ὁ Γαλλίων πρὸς τοὺς Ἰουδαίους· Εἰ μὲν ἦν ἀδίκημά τι ἢ ῥᾳδιούργημα πονηρόν, ὦ Ἰουδαῖοι, κατὰ λόγον ἂν ἀνεσχόμην ὑμῶν,
\vs{15}εἰ δὲ ζητήματά ἐστιν περὶ λόγου καὶ ὀνομάτων καὶ νόμου τοῦ καθ᾽ ὑμᾶς, ὄψεσθε αὐτοί· κριτὴς ἐγὼ τούτων οὐ βούλομαι εἶναι.
\vs{16}καὶ ἀπήλασεν αὐτοὺς ἀπὸ τοῦ βήματος.
\vs{17}Ἐπιλαβόμενοι δὲ πάντες Σωσθένην τὸν ἀρχισυνάγωγον ἔτυπτον ἔμπροσθεν τοῦ βήματος· καὶ οὐδὲν τούτων τῷ Γαλλίωνι ἔμελεν.

\vs{18}Ὁ δὲ Παῦλος ἔτι προσμείνας ἡμέρας ἱκανὰς τοῖς ἀδελφοῖς ἀποταξάμενος ἐξέπλει εἰς τὴν Συρίαν, καὶ σὺν αὐτῷ Πρίσκιλλα καὶ Ἀκύλας, κειράμενος ἐν Κενχρεαῖς τὴν κεφαλήν, εἶχεν γὰρ εὐχήν.
\vs{19}Κατήντησαν δὲ εἰς Ἔφεσον κἀκείνους κατέλιπεν αὐτοῦ, αὐτὸς δὲ εἰσελθὼν εἰς τὴν συναγωγὴν διελέξατο τοῖς Ἰουδαίοις.
\vs{20}ἐρωτώντων δὲ αὐτῶν ἐπὶ πλείονα χρόνον μεῖναι οὐκ ἐπένευσεν,
\vs{21}ἀλλὰ ἀποταξάμενος καὶ εἰπών· Πάλιν ἀνακάμψω πρὸς ὑμᾶς τοῦ Θεοῦ θέλοντος, ἀνήχθη ἀπὸ τῆς Ἐφέσου,
\vs{22}καὶ κατελθὼν εἰς Καισάρειαν, ἀναβὰς καὶ ἀσπασάμενος τὴν ἐκκλησίαν κατέβη εἰς Ἀντιόχειαν.

\vs{23}Καὶ ποιήσας χρόνον τινὰ ἐξῆλθεν διερχόμενος καθεξῆς τὴν Γαλατικὴν χώραν καὶ Φρυγίαν, στηρίζων πάντας τοὺς μαθητάς.

\vs{24}Ἰουδαῖος δέ τις Ἀπολλῶς ὀνόματι, Ἀλεξανδρεὺς τῷ γένει, ἀνὴρ λόγιος, κατήντησεν εἰς Ἔφεσον, δυνατὸς ὢν ἐν ταῖς γραφαῖς.
\vs{25}οὗτος ἦν κατηχημένος τὴν ὁδὸν τοῦ Κυρίου καὶ ζέων τῷ πνεύματι ἐλάλει καὶ ἐδίδασκεν ἀκριβῶς τὰ περὶ τοῦ Ἰησοῦ, ἐπιστάμενος μόνον τὸ βάπτισμα Ἰωάννου·
\vs{26}οὗτός τε ἤρξατο παρρησιάζεσθαι ἐν τῇ συναγωγῇ. ἀκούσαντες δὲ αὐτοῦ Πρίσκιλλα καὶ Ἀκύλας προσελάβοντο αὐτὸν καὶ ἀκριβέστερον αὐτῷ ἐξέθεντο τὴν ὁδὸν τοῦ Θεοῦ.
\vs{27}Βουλομένου δὲ αὐτοῦ διελθεῖν εἰς τὴν Ἀχαΐαν, προτρεψάμενοι οἱ ἀδελφοὶ ἔγραψαν τοῖς μαθηταῖς ἀποδέξασθαι αὐτόν, ὃς παραγενόμενος συνεβάλετο πολὺ τοῖς πεπιστευκόσιν διὰ τῆς χάριτος·
\vs{28}εὐτόνως γὰρ τοῖς Ἰουδαίοις διακατηλέγχετο δημοσίᾳ ἐπιδεικνὺς διὰ τῶν γραφῶν εἶναι τὸν Χριστὸν Ἰησοῦν.

\ch{19}
Ἐγένετο δὲ ἐν τῷ τὸν Ἀπολλῶ εἶναι ἐν Κορίνθῳ Παῦλον διελθόντα τὰ ἀνωτερικὰ μέρη κατελθεῖν εἰς Ἔφεσον καὶ εὑρεῖν τινας μαθητάς
\vs{2}εἶπέν τε πρὸς αὐτούς· Εἰ Πνεῦμα Ἅγιον ἐλάβετε πιστεύσαντες; Οἱ δὲ πρὸς αὐτόν· Ἀλλ᾽ οὐδ᾽ εἰ Πνεῦμα Ἅγιον ἔστιν ἠκούσαμεν.
\vs{3}Εἶπέν τε· Εἰς τί οὖν ἐβαπτίσθητε; Οἱ δὲ εἶπαν· Εἰς τὸ Ἰωάννου βάπτισμα.
\vs{4}Εἶπεν δὲ Παῦλος· Ἰωάννης ἐβάπτισεν βάπτισμα μετανοίας τῷ λαῷ λέγων εἰς τὸν ἐρχόμενον μετ᾽ αὐτὸν ἵνα πιστεύσωσιν, τοῦτ᾽ ἔστιν εἰς τὸν Ἰησοῦν.
\vs{5}Ἀκούσαντες δὲ ἐβαπτίσθησαν εἰς τὸ ὄνομα τοῦ Κυρίου Ἰησοῦ,
\vs{6}καὶ ἐπιθέντος αὐτοῖς τοῦ Παύλου τὰς χεῖρας ἦλθε τὸ Πνεῦμα τὸ Ἅγιον ἐπ᾽ αὐτούς, ἐλάλουν τε γλώσσαις καὶ ἐπροφήτευον.
\vs{7}ἦσαν δὲ οἱ πάντες ἄνδρες ὡσεὶ δώδεκα.

\vs{8}Εἰσελθὼν δὲ εἰς τὴν συναγωγὴν ἐπαρρησιάζετο ἐπὶ μῆνας τρεῖς διαλεγόμενος καὶ πείθων τὰ περὶ τῆς βασιλείας τοῦ Θεοῦ.
\vs{9}ὡς δέ τινες ἐσκληρύνοντο καὶ ἠπείθουν κακολογοῦντες τὴν Ὁδὸν ἐνώπιον τοῦ πλήθους, ἀποστὰς ἀπ᾽ αὐτῶν ἀφώρισεν τοὺς μαθητάς καθ᾽ ἡμέραν διαλεγόμενος ἐν τῇ σχολῇ Τυράννου.
\vs{10}τοῦτο δὲ ἐγένετο ἐπὶ ἔτη δύο, ὥστε πάντας τοὺς κατοικοῦντας τὴν Ἀσίαν ἀκοῦσαι τὸν λόγον τοῦ Κυρίου, Ἰουδαίους τε καὶ Ἕλληνας.
\vs{11}Δυνάμεις τε οὐ τὰς τυχούσας ὁ Θεὸς ἐποίει διὰ τῶν χειρῶν Παύλου,
\vs{12}ὥστε καὶ ἐπὶ τοὺς ἀσθενοῦντας ἀποφέρεσθαι ἀπὸ τοῦ χρωτὸς αὐτοῦ σουδάρια ἢ σιμικίνθια καὶ ἀπαλλάσσεσθαι ἀπ᾽ αὐτῶν τὰς νόσους, τά τε πνεύματα τὰ πονηρὰ ἐκπορεύεσθαι.

\vs{13}Ἐπεχείρησαν δέ τινες καὶ τῶν περιερχομένων Ἰουδαίων ἐξορκιστῶν ὀνομάζειν ἐπὶ τοὺς ἔχοντας τὰ πνεύματα τὰ πονηρὰ τὸ ὄνομα τοῦ Κυρίου Ἰησοῦ λέγοντες· Ὁρκίζω ὑμᾶς τὸν Ἰησοῦν ὃν Παῦλος κηρύσσει.
\vs{14}ἦσαν δέ τινος Σκευᾶ Ἰουδαίου ἀρχιερέως ἑπτὰ υἱοὶ τοῦτο ποιοῦντες.
\vs{15}ἀποκριθὲν δὲ τὸ πνεῦμα τὸ πονηρὸν εἶπεν αὐτοῖς· Τὸν Μὲν Ἰησοῦν γινώσκω καὶ τὸν Παῦλον ἐπίσταμαι, ὑμεῖς δὲ τίνες ἐστέ;
\vs{16}καὶ ἐφαλόμενος ὁ ἄνθρωπος ἐπ᾽ αὐτοὺς ἐν ᾧ ἦν τὸ πνεῦμα τὸ πονηρὸν, κατακυριεύσας ἀμφοτέρων ἴσχυσεν κατ᾽ αὐτῶν ὥστε γυμνοὺς καὶ τετραυματισμένους ἐκφυγεῖν ἐκ τοῦ οἴκου ἐκείνου.
\vs{17}Τοῦτο δὲ ἐγένετο γνωστὸν πᾶσιν Ἰουδαίοις τε καὶ Ἕλλησιν τοῖς κατοικοῦσιν τὴν Ἔφεσον καὶ ἐπέπεσεν φόβος ἐπὶ πάντας αὐτούς καὶ ἐμεγαλύνετο τὸ ὄνομα τοῦ Κυρίου Ἰησοῦ.
\vs{18}πολλοί τε τῶν πεπιστευκότων ἤρχοντο ἐξομολογούμενοι καὶ ἀναγγέλλοντες τὰς πράξεις αὐτῶν.
\vs{19}ἱκανοὶ δὲ τῶν τὰ περίεργα πραξάντων συνενέγκαντες τὰς βίβλους κατέκαιον ἐνώπιον πάντων, καὶ συνεψήφισαν τὰς τιμὰς αὐτῶν καὶ εὗρον ἀργυρίου μυριάδας πέντε.
\vs{20}Οὕτως κατὰ κράτος τοῦ Κυρίου ὁ λόγος ηὔξανεν καὶ ἴσχυεν.

\vs{21}Ὡς δὲ ἐπληρώθη ταῦτα, ἔθετο ὁ Παῦλος ἐν τῷ πνεύματι διελθὼν τὴν Μακεδονίαν καὶ Ἀχαΐαν πορεύεσθαι εἰς Ἱεροσόλυμα εἰπὼν ὅτι Μετὰ τὸ γενέσθαι με ἐκεῖ δεῖ με καὶ Ῥώμην ἰδεῖν.
\vs{22}ἀποστείλας δὲ εἰς τὴν Μακεδονίαν δύο τῶν διακονούντων αὐτῷ, Τιμόθεον καὶ Ἔραστον, αὐτὸς ἐπέσχεν χρόνον εἰς τὴν Ἀσίαν.

\vs{23}Ἐγένετο δὲ κατὰ τὸν καιρὸν ἐκεῖνον τάραχος οὐκ ὀλίγος περὶ τῆς Ὁδοῦ.
\vs{24}Δημήτριος γάρ τις ὀνόματι, ἀργυροκόπος, ποιῶν ναοὺς ἀργυροῦς Ἀρτέμιδος παρείχετο τοῖς τεχνίταις οὐκ ὀλίγην ἐργασίαν,
\vs{25}οὓς συναθροίσας καὶ τοὺς περὶ τὰ τοιαῦτα ἐργάτας εἶπεν· Ἄνδρες, ἐπίστασθε ὅτι ἐκ ταύτης τῆς ἐργασίας ἡ εὐπορία ἡμῖν ἐστιν
\vs{26}καὶ θεωρεῖτε καὶ ἀκούετε ὅτι οὐ μόνον Ἐφέσου ἀλλὰ σχεδὸν πάσης τῆς Ἀσίας ὁ Παῦλος οὗτος πείσας μετέστησεν ἱκανὸν ὄχλον λέγων ὅτι οὐκ εἰσὶν θεοὶ οἱ διὰ χειρῶν γινόμενοι.
\vs{27}οὐ μόνον δὲ τοῦτο κινδυνεύει ἡμῖν τὸ μέρος εἰς ἀπελεγμὸν ἐλθεῖν ἀλλὰ καὶ τὸ τῆς μεγάλης θεᾶς Ἀρτέμιδος ἱερὸν εἰς οὐθὲν λογισθῆναι, μέλλειν τε καὶ καθαιρεῖσθαι τῆς μεγαλειότητος αὐτῆς ἣν ὅλη ἡ Ἀσία καὶ ἡ οἰκουμένη σέβεται.
\vs{28}Ἀκούσαντες δὲ καὶ γενόμενοι πλήρεις θυμοῦ ἔκραζον λέγοντες· Μεγάλη ἡ Ἄρτεμις Ἐφεσίων.
\vs{29}καὶ ἐπλήσθη ἡ πόλις τῆς συγχύσεως, ὥρμησάν τε ὁμοθυμαδὸν εἰς τὸ θέατρον συναρπάσαντες Γάϊον καὶ Ἀρίσταρχον Μακεδόνας, συνεκδήμους Παύλου.
\vs{30}Παύλου δὲ βουλομένου εἰσελθεῖν εἰς τὸν δῆμον οὐκ εἴων αὐτὸν οἱ μαθηταί·
\vs{31}τινὲς δὲ καὶ τῶν Ἀσιαρχῶν, ὄντες αὐτῷ φίλοι, πέμψαντες πρὸς αὐτὸν παρεκάλουν μὴ δοῦναι ἑαυτὸν εἰς τὸ θέατρον.
\vs{32}Ἄλλοι μὲν οὖν ἄλλο τι ἔκραζον· ἦν γὰρ ἡ ἐκκλησία συγκεχυμένη καὶ οἱ πλείους οὐκ ᾔδεισαν τίνος ἕνεκα συνεληλύθεισαν.
\vs{33}ἐκ δὲ τοῦ ὄχλου συνεβίβασαν Ἀλέξανδρον, προβαλόντων αὐτὸν τῶν Ἰουδαίων· ὁ δὲ Ἀλέξανδρος κατασείσας τὴν χεῖρα ἤθελεν ἀπολογεῖσθαι τῷ δήμῳ.
\vs{34}ἐπιγνόντες δὲ ὅτι Ἰουδαῖός ἐστιν, φωνὴ ἐγένετο μία ἐκ πάντων ὡς ἐπὶ ὥρας δύο κραζόντων· Μεγάλη ἡ Ἄρτεμις Ἐφεσίων.
\vs{35}Καταστείλας δὲ ὁ γραμματεὺς τὸν ὄχλον φησίν· Ἄνδρες Ἐφέσιοι, τίς γάρ ἐστιν ἀνθρώπων ὃς οὐ γινώσκει τὴν Ἐφεσίων πόλιν νεωκόρον οὖσαν τῆς μεγάλης Ἀρτέμιδος καὶ τοῦ διοπετοῦς;
\vs{36}ἀναντιρρήτων οὖν ὄντων τούτων δέον ἐστὶν ὑμᾶς κατεσταλμένους ὑπάρχειν καὶ μηδὲν προπετὲς πράσσειν.
\vs{37}ἠγάγετε γὰρ τοὺς ἄνδρας τούτους οὔτε ἱεροσύλους οὔτε βλασφημοῦντας τὴν θεὸν ἡμῶν.
\vs{38}Εἰ μὲν οὖν Δημήτριος καὶ οἱ σὺν αὐτῷ τεχνῖται ἔχουσι πρός τινα λόγον, ἀγοραῖοι ἄγονται καὶ ἀνθύπατοί εἰσιν, ἐγκαλείτωσαν ἀλλήλοις.
\vs{39}εἰ δέ τι περαιτέρω ἐπιζητεῖτε, ἐν τῇ ἐννόμῳ ἐκκλησίᾳ ἐπιλυθήσεται.
\vs{40}καὶ γὰρ κινδυνεύομεν ἐγκαλεῖσθαι στάσεως περὶ τῆς σήμερον, μηδενὸς αἰτίου ὑπάρχοντος περὶ οὗ οὐ δυνησόμεθα ἀποδοῦναι λόγον περὶ τῆς συστροφῆς ταύτης. Καὶ ταῦτα εἰπὼν ἀπέλυσεν τὴν ἐκκλησίαν.

\ch{20}
Μετὰ δὲ τὸ παύσασθαι τὸν θόρυβον μεταπεμψάμενος ὁ Παῦλος τοὺς μαθητὰς καὶ παρακαλέσας, ἀσπασάμενος ἐξῆλθεν πορεύεσθαι εἰς Μακεδονίαν.
\vs{2}διελθὼν δὲ τὰ μέρη ἐκεῖνα καὶ παρακαλέσας αὐτοὺς λόγῳ πολλῷ ἦλθεν εἰς τὴν Ἑλλάδα
\vs{3}ποιήσας τε μῆνας τρεῖς· γενομένης ἐπιβουλῆς αὐτῷ ὑπὸ τῶν Ἰουδαίων μέλλοντι ἀνάγεσθαι εἰς τὴν Συρίαν, ἐγένετο γνώμης τοῦ ὑποστρέφειν διὰ Μακεδονίας.
\vs{4}Συνείπετο δὲ αὐτῷ Σώπατρος Πύρρου Βεροιαῖος, Θεσσαλονικέων δὲ Ἀρίσταρχος καὶ Σεκοῦνδος, καὶ Γάϊος Δερβαῖος καὶ Τιμόθεος, Ἀσιανοὶ δὲ Τυχικὸς καὶ Τρόφιμος.
\vs{5}οὗτοι δὲ προελθόντες ἔμενον ἡμᾶς ἐν Τρῳάδι,
\vs{6}ἡμεῖς δὲ ἐξεπλεύσαμεν μετὰ τὰς ἡμέρας τῶν ἀζύμων ἀπὸ Φιλίππων καὶ ἤλθομεν πρὸς αὐτοὺς εἰς τὴν Τρῳάδα ἄχρι ἡμερῶν πέντε, ὅπου διετρίψαμεν ἡμέρας ἑπτά.

\vs{7}Ἐν δὲ τῇ μιᾷ τῶν σαββάτων συνηγμένων ἡμῶν κλάσαι ἄρτον, ὁ Παῦλος διελέγετο αὐτοῖς μέλλων ἐξιέναι τῇ ἐπαύριον, παρέτεινέν τε τὸν λόγον μέχρι μεσονυκτίου.
\vs{8}Ἦσαν δὲ λαμπάδες ἱκαναὶ ἐν τῷ ὑπερῴῳ οὗ ἦμεν συνηγμένοι.
\vs{9}καθεζόμενος δέ τις νεανίας ὀνόματι Εὔτυχος ἐπὶ τῆς θυρίδος, καταφερόμενος ὕπνῳ βαθεῖ διαλεγομένου τοῦ Παύλου ἐπὶ πλεῖον, κατενεχθεὶς ἀπὸ τοῦ ὕπνου ἔπεσεν ἀπὸ τοῦ τριστέγου κάτω καὶ ἤρθη νεκρός.
\vs{10}καταβὰς δὲ ὁ Παῦλος ἐπέπεσεν αὐτῷ καὶ συμπεριλαβὼν εἶπεν· Μὴ θορυβεῖσθε, ἡ γὰρ ψυχὴ αὐτοῦ ἐν αὐτῷ ἐστιν.
\vs{11}Ἀναβὰς δὲ καὶ κλάσας τὸν ἄρτον καὶ γευσάμενος ἐφ᾽ ἱκανόν τε ὁμιλήσας ἄχρι αὐγῆς, οὕτως ἐξῆλθεν.
\vs{12}ἤγαγον δὲ τὸν παῖδα ζῶντα καὶ παρεκλήθησαν οὐ μετρίως.

\vs{13}Ἡμεῖς δὲ προελθόντες ἐπὶ τὸ πλοῖον ἀνήχθημεν ἐπὶ τὴν Ἆσσον ἐκεῖθεν μέλλοντες ἀναλαμβάνειν τὸν Παῦλον· οὕτως γὰρ διατεταγμένος ἦν μέλλων αὐτὸς πεζεύειν.
\vs{14}ὡς δὲ συνέβαλλεν ἡμῖν εἰς τὴν Ἆσσον, ἀναλαβόντες αὐτὸν ἤλθομεν εἰς Μιτυλήνην,
\vs{15}κἀκεῖθεν ἀποπλεύσαντες τῇ ἐπιούσῃ κατηντήσαμεν ἄντικρυς Χίου, τῇ δὲ ἑτέρᾳ παρεβάλομεν εἰς Σάμον, τῇ δὲ ἐχομένῃ ἤλθομεν εἰς Μίλητον.
\vs{16}Κεκρίκει γὰρ ὁ Παῦλος παραπλεῦσαι τὴν Ἔφεσον, ὅπως μὴ γένηται αὐτῷ χρονοτριβῆσαι ἐν τῇ Ἀσίᾳ· ἔσπευδεν γὰρ εἰ δυνατὸν εἴη αὐτῷ τὴν ἡμέραν τῆς Πεντηκοστῆς γενέσθαι εἰς Ἱεροσόλυμα.

\vs{17}Ἀπὸ δὲ τῆς Μιλήτου πέμψας εἰς Ἔφεσον μετεκαλέσατο τοὺς πρεσβυτέρους τῆς ἐκκλησίας.
\vs{18}Ὡς δὲ παρεγένοντο πρὸς αὐτὸν εἶπεν αὐτοῖς· Ὑμεῖς ἐπίστασθε, ἀπὸ πρώτης ἡμέρας ἀφ᾽ ἧς ἐπέβην εἰς τὴν Ἀσίαν, πῶς μεθ᾽ ὑμῶν τὸν πάντα χρόνον ἐγενόμην,
\vs{19}δουλεύων τῷ Κυρίῳ μετὰ πάσης ταπεινοφροσύνης καὶ δακρύων καὶ πειρασμῶν τῶν συμβάντων μοι ἐν ταῖς ἐπιβουλαῖς τῶν Ἰουδαίων,
\vs{20}ὡς οὐδὲν ὑπεστειλάμην τῶν συμφερόντων τοῦ μὴ ἀναγγεῖλαι ὑμῖν καὶ διδάξαι ὑμᾶς δημοσίᾳ καὶ κατ᾽ οἴκους,
\vs{21}διαμαρτυρόμενος Ἰουδαίοις τε καὶ Ἕλλησιν τὴν εἰς Θεὸν μετάνοιαν καὶ πίστιν εἰς τὸν Κύριον ἡμῶν Ἰησοῦν.
\vs{22}Καὶ νῦν ἰδοὺ δεδεμένος ἐγὼ τῷ πνεύματι πορεύομαι εἰς Ἰερουσαλήμ τὰ ἐν αὐτῇ συναντήσοντά μοι μὴ εἰδώς,
\vs{23}πλὴν ὅτι τὸ Πνεῦμα τὸ Ἅγιον κατὰ πόλιν διαμαρτύρεταί μοι λέγον ὅτι δεσμὰ καὶ θλίψεις με μένουσιν.
\vs{24}ἀλλ᾽ οὐδενὸς λόγου ποιοῦμαι τὴν ψυχὴν τιμίαν ἐμαυτῷ ὡς τελειῶσαι τὸν δρόμον μου καὶ τὴν διακονίαν ἣν ἔλαβον παρὰ τοῦ Κυρίου Ἰησοῦ, διαμαρτύρασθαι τὸ εὐαγγέλιον τῆς χάριτος τοῦ Θεοῦ.
\vs{25}Καὶ νῦν ἰδοὺ ἐγὼ οἶδα ὅτι οὐκέτι ὄψεσθε τὸ πρόσωπόν μου ὑμεῖς πάντες ἐν οἷς διῆλθον κηρύσσων τὴν βασιλείαν.
\vs{26}διότι μαρτύρομαι ὑμῖν ἐν τῇ σήμερον ἡμέρᾳ ὅτι καθαρός εἰμι ἀπὸ τοῦ αἵματος πάντων·
\vs{27}οὐ γὰρ ὑπεστειλάμην τοῦ μὴ ἀναγγεῖλαι πᾶσαν τὴν βουλὴν τοῦ Θεοῦ ὑμῖν.
\vs{28}Προσέχετε ἑαυτοῖς καὶ παντὶ τῷ ποιμνίῳ, ἐν ᾧ ὑμᾶς τὸ Πνεῦμα τὸ Ἅγιον ἔθετο ἐπισκόπους ποιμαίνειν τὴν ἐκκλησίαν τοῦ Θεοῦ, ἣν περιεποιήσατο διὰ τοῦ αἵματος τοῦ ἰδίου.
\vs{29}ἐγὼ οἶδα ὅτι εἰσελεύσονται μετὰ τὴν ἄφιξίν μου λύκοι βαρεῖς εἰς ὑμᾶς μὴ φειδόμενοι τοῦ ποιμνίου,
\vs{30}καὶ ἐξ ὑμῶν αὐτῶν ἀναστήσονται ἄνδρες λαλοῦντες διεστραμμένα τοῦ ἀποσπᾶν τοὺς μαθητὰς ὀπίσω ἑαυτῶν.
\vs{31}διὸ γρηγορεῖτε μνημονεύοντες ὅτι τριετίαν νύκτα καὶ ἡμέραν οὐκ ἐπαυσάμην μετὰ δακρύων νουθετῶν ἕνα ἕκαστον.
\vs{32}Καὶ τὰ νῦν παρατίθεμαι ὑμᾶς τῷ θεῷ καὶ τῷ λόγῳ τῆς χάριτος αὐτοῦ, τῷ δυναμένῳ οἰκοδομῆσαι καὶ δοῦναι τὴν κληρονομίαν ἐν τοῖς ἡγιασμένοις πᾶσιν.
\vs{33}Ἀργυρίου ἢ χρυσίου ἢ ἱματισμοῦ οὐδενὸς ἐπεθύμησα·
\vs{34}αὐτοὶ γινώσκετε ὅτι ταῖς χρείαις μου καὶ τοῖς οὖσιν μετ᾽ ἐμοῦ ὑπηρέτησαν αἱ χεῖρες αὗται.
\vs{35}πάντα ὑπέδειξα ὑμῖν ὅτι οὕτως κοπιῶντας δεῖ ἀντιλαμβάνεσθαι τῶν ἀσθενούντων, μνημονεύειν τε τῶν λόγων τοῦ Κυρίου Ἰησοῦ ὅτι αὐτὸς εἶπεν· Μακάριόν ἐστιν μᾶλλον διδόναι ἢ λαμβάνειν.
\vs{36}Καὶ ταῦτα εἰπὼν θεὶς τὰ γόνατα αὐτοῦ σὺν πᾶσιν αὐτοῖς προσηύξατο.
\vs{37}ἱκανὸς δὲ κλαυθμὸς ἐγένετο πάντων καὶ ἐπιπεσόντες ἐπὶ τὸν τράχηλον τοῦ Παύλου κατεφίλουν αὐτόν,
\vs{38}ὀδυνώμενοι μάλιστα ἐπὶ τῷ λόγῳ ᾧ εἰρήκει, ὅτι οὐκέτι μέλλουσιν τὸ πρόσωπον αὐτοῦ θεωρεῖν. προέπεμπον δὲ αὐτὸν εἰς τὸ πλοῖον.

\ch{21}
Ὡς δὲ ἐγένετο ἀναχθῆναι ἡμᾶς ἀποσπασθέντας ἀπ᾽ αὐτῶν, εὐθυδρομήσαντες ἤλθομεν εἰς τὴν Κῶ, τῇ δὲ ἑξῆς εἰς τὴν Ῥόδον κἀκεῖθεν εἰς Πάταρα,
\vs{2}καὶ εὑρόντες πλοῖον διαπερῶν εἰς Φοινίκην ἐπιβάντες ἀνήχθημεν.
\vs{3}ἀναφάναντες δὲ τὴν Κύπρον καὶ καταλιπόντες αὐτὴν εὐώνυμον ἐπλέομεν εἰς Συρίαν καὶ κατήλθομεν εἰς Τύρον· ἐκεῖσε γὰρ τὸ πλοῖον ἦν ἀποφορτιζόμενον τὸν γόμον.
\vs{4}Ἀνευρόντες δὲ τοὺς μαθητὰς ἐπεμείναμεν αὐτοῦ ἡμέρας ἑπτά, οἵτινες τῷ Παύλῳ ἔλεγον διὰ τοῦ Πνεύματος μὴ ἐπιβαίνειν εἰς Ἱεροσόλυμα.
\vs{5}ὅτε δὲ ἐγένετο ἡμᾶς ἐξαρτίσαι τὰς ἡμέρας, ἐξελθόντες ἐπορευόμεθα προπεμπόντων ἡμᾶς πάντων σὺν γυναιξὶ καὶ τέκνοις ἕως ἔξω τῆς πόλεως, καὶ θέντες τὰ γόνατα ἐπὶ τὸν αἰγιαλὸν προσευξάμενοι
\vs{6}ἀπησπασάμεθα ἀλλήλους καὶ ἀνέβημεν εἰς τὸ πλοῖον, ἐκεῖνοι δὲ ὑπέστρεψαν εἰς τὰ ἴδια.

\vs{7}Ἡμεῖς δὲ τὸν πλοῦν διανύσαντες ἀπὸ Τύρου κατηντήσαμεν εἰς Πτολεμαΐδα καὶ ἀσπασάμενοι τοὺς ἀδελφοὺς ἐμείναμεν ἡμέραν μίαν παρ᾽ αὐτοῖς.
\vs{8}Τῇ δὲ ἐπαύριον ἐξελθόντες ἤλθομεν εἰς Καισάρειαν καὶ εἰσελθόντες εἰς τὸν οἶκον Φιλίππου τοῦ εὐαγγελιστοῦ, ὄντος ἐκ τῶν ἑπτὰ, ἐμείναμεν παρ᾽ αὐτῷ.
\vs{9}τούτῳ δὲ ἦσαν θυγατέρες τέσσαρες παρθένοι προφητεύουσαι.
\vs{10}Ἐπιμενόντων δὲ ἡμέρας πλείους κατῆλθέν τις ἀπὸ τῆς Ἰουδαίας προφήτης ὀνόματι Ἅγαβος,
\vs{11}καὶ ἐλθὼν πρὸς ἡμᾶς καὶ ἄρας τὴν ζώνην τοῦ Παύλου, δήσας ἑαυτοῦ τοὺς πόδας καὶ τὰς χεῖρας εἶπεν· Τάδε λέγει τὸ Πνεῦμα τὸ Ἅγιον· Τὸν ἄνδρα οὗ ἐστιν ἡ ζώνη αὕτη, οὕτως δήσουσιν ἐν Ἰερουσαλὴμ οἱ Ἰουδαῖοι καὶ παραδώσουσιν εἰς χεῖρας ἐθνῶν.
\vs{12}ὡς δὲ ἠκούσαμεν ταῦτα, παρεκαλοῦμεν ἡμεῖς τε καὶ οἱ ἐντόπιοι τοῦ μὴ ἀναβαίνειν αὐτὸν εἰς Ἰερουσαλήμ.
\vs{13}Τότε ἀπεκρίθη ὁ Παῦλος· Τί ποιεῖτε κλαίοντες καὶ συνθρύπτοντές μου τὴν καρδίαν; ἐγὼ γὰρ οὐ μόνον δεθῆναι ἀλλὰ καὶ ἀποθανεῖν εἰς Ἰερουσαλὴμ ἑτοίμως ἔχω ὑπὲρ τοῦ ὀνόματος τοῦ Κυρίου Ἰησοῦ.
\vs{14}μὴ πειθομένου δὲ αὐτοῦ ἡσυχάσαμεν εἰπόντες· Τοῦ Κυρίου τὸ θέλημα γινέσθω.

\vs{15}Μετὰ δὲ τὰς ἡμέρας ταύτας ἐπισκευασάμενοι ἀνεβαίνομεν εἰς Ἱεροσόλυμα·
\vs{16}συνῆλθον δὲ καὶ τῶν μαθητῶν ἀπὸ Καισαρείας σὺν ἡμῖν, ἄγοντες παρ᾽ ᾧ ξενισθῶμεν Μνάσωνί τινι Κυπρίῳ, ἀρχαίῳ μαθητῇ.
\vs{17}Γενομένων δὲ ἡμῶν εἰς Ἱεροσόλυμα ἀσμένως ἀπεδέξαντο ἡμᾶς οἱ ἀδελφοί.

\vs{18}τῇ δὲ ἐπιούσῃ εἰσῄει ὁ Παῦλος σὺν ἡμῖν πρὸς Ἰάκωβον, πάντες τε παρεγένοντο οἱ πρεσβύτεροι.
\vs{19}καὶ ἀσπασάμενος αὐτοὺς ἐξηγεῖτο καθ᾽ ἓν ἕκαστον, ὧν ἐποίησεν ὁ Θεὸς ἐν τοῖς ἔθνεσιν διὰ τῆς διακονίας αὐτοῦ.
\vs{20}Οἱ δὲ ἀκούσαντες ἐδόξαζον τὸν Θεόν εἶπόν τε αὐτῷ· Θεωρεῖς, ἀδελφέ, πόσαι μυριάδες εἰσὶν ἐν τοῖς Ἰουδαίοις τῶν πεπιστευκότων καὶ πάντες ζηλωταὶ τοῦ νόμου ὑπάρχουσιν·
\vs{21}κατηχήθησαν δὲ περὶ σοῦ ὅτι ἀποστασίαν διδάσκεις ἀπὸ Μωϋσέως τοὺς κατὰ τὰ ἔθνη πάντας Ἰουδαίους λέγων μὴ περιτέμνειν αὐτοὺς τὰ τέκνα μηδὲ τοῖς ἔθεσιν περιπατεῖν.
\vs{22}τί οὖν ἐστιν; πάντως ἀκούσονται ὅτι ἐλήλυθας.
\vs{23}Τοῦτο οὖν ποίησον ὅ σοι λέγομεν· εἰσὶν ἡμῖν ἄνδρες τέσσαρες εὐχὴν ἔχοντες ἐφ᾽ ἑαυτῶν.
\vs{24}τούτους παραλαβὼν ἁγνίσθητι σὺν αὐτοῖς καὶ δαπάνησον ἐπ᾽ αὐτοῖς ἵνα ξυρήσονται τὴν κεφαλήν, καὶ γνώσονται πάντες ὅτι ὧν κατήχηνται περὶ σοῦ οὐδέν ἐστιν ἀλλὰ στοιχεῖς καὶ αὐτὸς φυλάσσων τὸν νόμον.
\vs{25}Περὶ δὲ τῶν πεπιστευκότων ἐθνῶν ἡμεῖς ἐπεστείλαμεν κρίναντες φυλάσσεσθαι αὐτοὺς τό τε εἰδωλόθυτον καὶ αἷμα καὶ πνικτὸν καὶ πορνείαν.
\vs{26}Τότε ὁ Παῦλος παραλαβὼν τοὺς ἄνδρας τῇ ἐχομένῃ ἡμέρᾳ σὺν αὐτοῖς ἁγνισθεὶς, εἰσῄει εἰς τὸ ἱερόν διαγγέλλων τὴν ἐκπλήρωσιν τῶν ἡμερῶν τοῦ ἁγνισμοῦ ἕως οὗ προσηνέχθη ὑπὲρ ἑνὸς ἑκάστου αὐτῶν ἡ προσφορά.

\vs{27}Ὡς δὲ ἔμελλον αἱ ἑπτὰ ἡμέραι συντελεῖσθαι, οἱ ἀπὸ τῆς Ἀσίας Ἰουδαῖοι θεασάμενοι αὐτὸν ἐν τῷ ἱερῷ συνέχεον πάντα τὸν ὄχλον καὶ ἐπέβαλον ἐπ᾽ αὐτὸν τὰς χεῖρας
\vs{28}κράζοντες· Ἄνδρες Ἰσραηλῖται, βοηθεῖτε· οὗτός ἐστιν ὁ ἄνθρωπος ὁ κατὰ τοῦ λαοῦ καὶ τοῦ νόμου καὶ τοῦ τόπου τούτου πάντας πανταχῇ διδάσκων, ἔτι τε καὶ Ἕλληνας εἰσήγαγεν εἰς τὸ ἱερὸν καὶ κεκοίνωκεν τὸν ἅγιον τόπον τοῦτον.
\vs{29}ἦσαν γὰρ προεωρακότες Τρόφιμον τὸν Ἐφέσιον ἐν τῇ πόλει σὺν αὐτῷ, ὃν ἐνόμιζον ὅτι εἰς τὸ ἱερὸν εἰσήγαγεν ὁ Παῦλος.
\vs{30}Ἐκινήθη τε ἡ πόλις ὅλη καὶ ἐγένετο συνδρομὴ τοῦ λαοῦ, καὶ ἐπιλαβόμενοι τοῦ Παύλου εἷλκον αὐτὸν ἔξω τοῦ ἱεροῦ καὶ εὐθέως ἐκλείσθησαν αἱ θύραι.
\vs{31}Ζητούντων τε αὐτὸν ἀποκτεῖναι ἀνέβη φάσις τῷ χιλιάρχῳ τῆς σπείρης ὅτι ὅλη συνχύννεται Ἰερουσαλήμ.
\vs{32}ὃς ἐξαυτῆς παραλαβὼν στρατιώτας καὶ ἑκατοντάρχας κατέδραμεν ἐπ᾽ αὐτούς, οἱ δὲ ἰδόντες τὸν χιλίαρχον καὶ τοὺς στρατιώτας ἐπαύσαντο τύπτοντες τὸν Παῦλον.
\vs{33}Τότε ἐγγίσας ὁ χιλίαρχος ἐπελάβετο αὐτοῦ καὶ ἐκέλευσεν δεθῆναι ἁλύσεσι δυσί, καὶ ἐπυνθάνετο τίς εἴη καὶ τί ἐστιν πεποιηκώς.
\vs{34}Ἄλλοι δὲ ἄλλο τι ἐπεφώνουν ἐν τῷ ὄχλῳ. μὴ δυναμένου δὲ αὐτοῦ γνῶναι τὸ ἀσφαλὲς διὰ τὸν θόρυβον ἐκέλευσεν ἄγεσθαι αὐτὸν εἰς τὴν παρεμβολήν.
\vs{35}ὅτε δὲ ἐγένετο ἐπὶ τοὺς ἀναβαθμούς, συνέβη βαστάζεσθαι αὐτὸν ὑπὸ τῶν στρατιωτῶν διὰ τὴν βίαν τοῦ ὄχλου,
\vs{36}ἠκολούθει γὰρ τὸ πλῆθος τοῦ λαοῦ κράζοντες· Αἶρε αὐτόν.

\vs{37}Μέλλων τε εἰσάγεσθαι εἰς τὴν παρεμβολὴν ὁ Παῦλος λέγει τῷ χιλιάρχῳ· Εἰ ἔξεστίν μοι εἰπεῖν τι πρὸς σέ; Ὁ δὲ ἔφη· Ἑλληνιστὶ γινώσκεις;
\vs{38}οὐκ ἄρα σὺ εἶ ὁ Αἰγύπτιος ὁ πρὸ τούτων τῶν ἡμερῶν ἀναστατώσας καὶ ἐξαγαγὼν εἰς τὴν ἔρημον τοὺς τετρακισχιλίους ἄνδρας τῶν Σικαρίων;
\vs{39}Εἶπεν δὲ ὁ Παῦλος· Ἐγὼ ἄνθρωπος μέν εἰμι Ἰουδαῖος, Ταρσεὺς τῆς Κιλικίας, οὐκ ἀσήμου πόλεως πολίτης· δέομαι δέ σου, ἐπίτρεψόν μοι λαλῆσαι πρὸς τὸν λαόν.
\vs{40}ἐπιτρέψαντος δὲ αὐτοῦ ὁ Παῦλος ἑστὼς ἐπὶ τῶν ἀναβαθμῶν κατέσεισεν τῇ χειρὶ τῷ λαῷ. πολλῆς δὲ σιγῆς γενομένης προσεφώνησεν τῇ Ἑβραΐδι διαλέκτῳ λέγων·

\ch{22}
Ἄνδρες ἀδελφοὶ καὶ πατέρες, ἀκούσατέ μου τῆς πρὸς ὑμᾶς νυνὶ ἀπολογίας.
\vs{2}ἀκούσαντες δὲ ὅτι τῇ Ἑβραΐδι διαλέκτῳ προσεφώνει αὐτοῖς, μᾶλλον παρέσχον ἡσυχίαν. Καὶ φησίν·
\vs{3}Ἐγώ εἰμι ἀνὴρ Ἰουδαῖος, γεγεννημένος ἐν Ταρσῷ τῆς Κιλικίας, ἀνατεθραμμένος δὲ ἐν τῇ πόλει ταύτῃ, παρὰ τοὺς πόδας Γαμαλιήλ πεπαιδευμένος κατὰ ἀκρίβειαν τοῦ πατρῴου νόμου, ζηλωτὴς ὑπάρχων τοῦ Θεοῦ καθὼς πάντες ὑμεῖς ἐστε σήμερον·
\vs{4}ὃς ταύτην τὴν Ὁδὸν ἐδίωξα ἄχρι θανάτου δεσμεύων καὶ παραδιδοὺς εἰς φυλακὰς ἄνδρας τε καὶ γυναῖκας,
\vs{5}ὡς καὶ ὁ ἀρχιερεὺς μαρτυρεῖ μοι καὶ πᾶν τὸ πρεσβυτέριον, παρ᾽ ὧν καὶ ἐπιστολὰς δεξάμενος πρὸς τοὺς ἀδελφοὺς εἰς Δαμασκὸν ἐπορευόμην, ἄξων καὶ τοὺς ἐκεῖσε ὄντας δεδεμένους εἰς Ἰερουσαλὴμ ἵνα τιμωρηθῶσιν.
\vs{6}Ἐγένετο δέ μοι πορευομένῳ καὶ ἐγγίζοντι τῇ Δαμασκῷ περὶ μεσημβρίαν ἐξαίφνης ἐκ τοῦ οὐρανοῦ περιαστράψαι φῶς ἱκανὸν περὶ ἐμέ,
\vs{7}ἔπεσά τε εἰς τὸ ἔδαφος καὶ ἤκουσα φωνῆς λεγούσης μοι· Σαοὺλ Σαούλ, τί με διώκεις;
\vs{8}Ἐγὼ δὲ ἀπεκρίθην· Τίς εἶ, Κύριε; Εἶπέν τε πρὸς ἐμέ· Ἐγώ εἰμι Ἰησοῦς ὁ Ναζωραῖος, ὃν σὺ διώκεις.
\vs{9}οἱ δὲ σὺν ἐμοὶ ὄντες τὸ μὲν φῶς ἐθεάσαντο τὴν δὲ φωνὴν οὐκ ἤκουσαν τοῦ λαλοῦντός μοι.
\vs{10}Εἶπον δέ· Τί ποιήσω, Κύριε; Ὁ δὲ Κύριος εἶπεν πρός με· Ἀναστὰς πορεύου εἰς Δαμασκόν κἀκεῖ σοι λαληθήσεται περὶ πάντων ὧν τέτακταί σοι ποιῆσαι.
\vs{11}Ὡς δὲ οὐκ ἐνέβλεπον ἀπὸ τῆς δόξης τοῦ φωτὸς ἐκείνου, χειραγωγούμενος ὑπὸ τῶν συνόντων μοι ἦλθον εἰς Δαμασκόν.
\vs{12}Ἁνανίας δέ τις, ἀνὴρ εὐλαβὴς κατὰ τὸν νόμον, μαρτυρούμενος ὑπὸ πάντων τῶν κατοικούντων Ἰουδαίων,
\vs{13}ἐλθὼν πρὸς ἐμὲ καὶ ἐπιστὰς εἶπέν μοι· Σαοὺλ ἀδελφέ, ἀνάβλεψον. κἀγὼ αὐτῇ τῇ ὥρᾳ ἀνέβλεψα εἰς αὐτόν.
\vs{14}Ὁ δὲ εἶπεν· Ὁ Θεὸς τῶν πατέρων ἡμῶν προεχειρίσατό σε γνῶναι τὸ θέλημα αὐτοῦ καὶ ἰδεῖν τὸν Δίκαιον καὶ ἀκοῦσαι φωνὴν ἐκ τοῦ στόματος αὐτοῦ,
\vs{15}ὅτι ἔσῃ μάρτυς αὐτῷ πρὸς πάντας ἀνθρώπους ὧν ἑώρακας καὶ ἤκουσας.
\vs{16}καὶ νῦν τί μέλλεις; ἀναστὰς βάπτισαι καὶ ἀπόλουσαι τὰς ἁμαρτίας σου ἐπικαλεσάμενος τὸ ὄνομα αὐτοῦ.
\vs{17}Ἐγένετο δέ μοι ὑποστρέψαντι εἰς Ἰερουσαλὴμ καὶ προσευχομένου μου ἐν τῷ ἱερῷ γενέσθαι με ἐν ἐκστάσει
\vs{18}καὶ ἰδεῖν αὐτὸν λέγοντά μοι· Σπεῦσον καὶ ἔξελθε ἐν τάχει ἐξ Ἰερουσαλήμ, διότι οὐ παραδέξονταί σου μαρτυρίαν περὶ ἐμοῦ.
\vs{19}Κἀγὼ εἶπον· Κύριε, αὐτοὶ ἐπίστανται ὅτι ἐγὼ ἤμην φυλακίζων καὶ δέρων κατὰ τὰς συναγωγὰς τοὺς πιστεύοντας ἐπὶ σέ,
\vs{20}καὶ ὅτε ἐξεχύννετο τὸ αἷμα Στεφάνου τοῦ μάρτυρός σου, καὶ αὐτὸς ἤμην ἐφεστὼς καὶ συνευδοκῶν καὶ φυλάσσων τὰ ἱμάτια τῶν ἀναιρούντων αὐτόν.
\vs{21}Καὶ εἶπεν πρός με· Πορεύου, ὅτι ἐγὼ εἰς ἔθνη μακρὰν ἐξαποστελῶ σε.

\vs{22}Ἤκουον δὲ αὐτοῦ ἄχρι τούτου τοῦ λόγου καὶ ἐπῆραν τὴν φωνὴν αὐτῶν λέγοντες· Αἶρε ἀπὸ τῆς γῆς τὸν τοιοῦτον, οὐ γὰρ καθῆκεν αὐτὸν ζῆν.
\vs{23}Κραυγαζόντων τε αὐτῶν καὶ ῥιπτούντων τὰ ἱμάτια καὶ κονιορτὸν βαλλόντων εἰς τὸν ἀέρα,
\vs{24}ἐκέλευσεν ὁ χιλίαρχος εἰσάγεσθαι αὐτὸν εἰς τὴν παρεμβολήν, εἴπας μάστιξιν ἀνετάζεσθαι αὐτὸν ἵνα ἐπιγνῷ δι᾽ ἣν αἰτίαν οὕτως ἐπεφώνουν αὐτῷ.
\vs{25}Ὡς δὲ προέτειναν αὐτὸν τοῖς ἱμᾶσιν, εἶπεν πρὸς τὸν ἑστῶτα ἑκατόνταρχον ὁ Παῦλος· Εἰ ἄνθρωπον Ῥωμαῖον καὶ ἀκατάκριτον ἔξεστιν ὑμῖν μαστίζειν;
\vs{26}Ἀκούσας δὲ ὁ ἑκατοντάρχης προσελθὼν τῷ χιλιάρχῳ ἀπήγγειλεν λέγων· Τί μέλλεις ποιεῖν; ὁ γὰρ ἄνθρωπος οὗτος Ῥωμαῖός ἐστιν.
\vs{27}Προσελθὼν δὲ ὁ χιλίαρχος εἶπεν αὐτῷ· Λέγε μοι, σὺ Ῥωμαῖος εἶ; Ὁ δὲ ἔφη· Ναί.
\vs{28}Ἀπεκρίθη δὲ ὁ χιλίαρχος· Ἐγὼ πολλοῦ κεφαλαίου τὴν πολιτείαν ταύτην ἐκτησάμην. ὁ Δὲ Παῦλος ἔφη· Ἐγὼ δὲ καὶ γεγέννημαι.
\vs{29}Εὐθέως οὖν ἀπέστησαν ἀπ᾽ αὐτοῦ οἱ μέλλοντες αὐτὸν ἀνετάζειν, καὶ ὁ χιλίαρχος δὲ ἐφοβήθη ἐπιγνοὺς ὅτι Ῥωμαῖός ἐστιν καὶ ὅτι αὐτὸν ἦν δεδεκώς.
\vs{30}Τῇ δὲ ἐπαύριον βουλόμενος γνῶναι τὸ ἀσφαλὲς, τὸ τί κατηγορεῖται ὑπὸ τῶν Ἰουδαίων, ἔλυσεν αὐτόν καὶ ἐκέλευσεν συνελθεῖν τοὺς ἀρχιερεῖς καὶ πᾶν τὸ συνέδριον, καὶ καταγαγὼν τὸν Παῦλον ἔστησεν εἰς αὐτούς.

\ch{23}
Ἀτενίσας δὲ ὁ Παῦλος τῷ συνεδρίῳ εἶπεν· Ἄνδρες ἀδελφοί, ἐγὼ πάσῃ συνειδήσει ἀγαθῇ πεπολίτευμαι τῷ Θεῷ ἄχρι ταύτης τῆς ἡμέρας.
\vs{2}ὁ δὲ ἀρχιερεὺς Ἁνανίας ἐπέταξεν τοῖς παρεστῶσιν αὐτῷ τύπτειν αὐτοῦ τὸ στόμα.
\vs{3}Τότε ὁ Παῦλος πρὸς αὐτὸν εἶπεν· Τύπτειν σε μέλλει ὁ Θεός, τοῖχε κεκονιαμένε· καὶ σὺ κάθῃ κρίνων με κατὰ τὸν νόμον καὶ παρανομῶν κελεύεις με τύπτεσθαι;
\vs{4}Οἱ δὲ παρεστῶτες εἶπαν· Τὸν ἀρχιερέα τοῦ Θεοῦ λοιδορεῖς;
\vs{5}Ἔφη τε ὁ Παῦλος· Οὐκ ᾔδειν, ἀδελφοί, ὅτι ἐστὶν ἀρχιερεύς· γέγραπται γὰρ ὅτι Ἄρχοντα τοῦ λαοῦ σου οὐκ ἐρεῖς κακῶς.
\vs{6}Γνοὺς δὲ ὁ Παῦλος ὅτι τὸ ἓν μέρος ἐστὶν Σαδδουκαίων τὸ δὲ ἕτερον Φαρισαίων ἔκραζεν ἐν τῷ συνεδρίῳ· Ἄνδρες ἀδελφοί, ἐγὼ Φαρισαῖός εἰμι, υἱὸς Φαρισαίων, περὶ ἐλπίδος καὶ ἀναστάσεως νεκρῶν ἐγὼ κρίνομαι.
\vs{7}Τοῦτο δὲ αὐτοῦ λαλοῦντος ἐγένετο στάσις τῶν Φαρισαίων καὶ Σαδδουκαίων καὶ ἐσχίσθη τὸ πλῆθος.
\vs{8}Σαδδουκαῖοι μὲν γὰρ λέγουσιν μὴ εἶναι ἀνάστασιν μήτε ἄγγελον μήτε πνεῦμα, Φαρισαῖοι δὲ ὁμολογοῦσιν τὰ ἀμφότερα.
\vs{9}Ἐγένετο δὲ κραυγὴ μεγάλη, καὶ ἀναστάντες τινὲς τῶν γραμματέων τοῦ μέρους τῶν Φαρισαίων διεμάχοντο λέγοντες· Οὐδὲν κακὸν εὑρίσκομεν ἐν τῷ ἀνθρώπῳ τούτῳ· εἰ δὲ πνεῦμα ἐλάλησεν αὐτῷ ἢ ἄγγελος;
\vs{10}πολλῆς δὲ γινομένης στάσεως φοβηθεὶς ὁ χιλίαρχος μὴ διασπασθῇ ὁ Παῦλος ὑπ᾽ αὐτῶν ἐκέλευσεν τὸ στράτευμα καταβὰν ἁρπάσαι αὐτὸν ἐκ μέσου αὐτῶν ἄγειν τε εἰς τὴν παρεμβολήν.
\vs{11}Τῇ δὲ ἐπιούσῃ νυκτὶ ἐπιστὰς αὐτῷ ὁ Κύριος εἶπεν· Θάρσει· ὡς γὰρ διεμαρτύρω τὰ περὶ ἐμοῦ εἰς Ἰερουσαλὴμ, οὕτω σε δεῖ καὶ εἰς Ῥώμην μαρτυρῆσαι.

\vs{12}Γενομένης δὲ ἡμέρας ποιήσαντες συστροφὴν οἱ Ἰουδαῖοι ἀνεθεμάτισαν ἑαυτοὺς λέγοντες μήτε φαγεῖν μήτε πιεῖν ἕως οὗ ἀποκτείνωσιν τὸν Παῦλον.
\vs{13}ἦσαν δὲ πλείους τεσσεράκοντα οἱ ταύτην τὴν συνωμοσίαν ποιησάμενοι,
\vs{14}οἵτινες προσελθόντες τοῖς ἀρχιερεῦσιν καὶ τοῖς πρεσβυτέροις εἶπαν· Ἀναθέματι ἀνεθεματίσαμεν ἑαυτοὺς μηδενὸς γεύσασθαι ἕως οὗ ἀποκτείνωμεν τὸν Παῦλον.
\vs{15}νῦν οὖν ὑμεῖς ἐμφανίσατε τῷ χιλιάρχῳ σὺν τῷ συνεδρίῳ ὅπως καταγάγῃ αὐτὸν εἰς ὑμᾶς ὡς μέλλοντας διαγινώσκειν ἀκριβέστερον τὰ περὶ αὐτοῦ· ἡμεῖς δὲ πρὸ τοῦ ἐγγίσαι αὐτὸν ἕτοιμοί ἐσμεν τοῦ ἀνελεῖν αὐτόν.
\vs{16}Ἀκούσας δὲ ὁ υἱὸς τῆς ἀδελφῆς Παύλου τὴν ἐνέδραν, παραγενόμενος καὶ εἰσελθὼν εἰς τὴν παρεμβολὴν ἀπήγγειλεν τῷ Παύλῳ.
\vs{17}προσκαλεσάμενος δὲ ὁ Παῦλος ἕνα τῶν ἑκατονταρχῶν ἔφη· Τὸν νεανίαν τοῦτον ἄπαγε πρὸς τὸν χιλίαρχον, ἔχει γὰρ ἀπαγγεῖλαί τι αὐτῷ.
\vs{18}Ὁ μὲν οὖν παραλαβὼν αὐτὸν ἤγαγεν πρὸς τὸν χιλίαρχον καὶ φησίν· Ὁ δέσμιος Παῦλος προσκαλεσάμενός με ἠρώτησεν τοῦτον τὸν νεανίσκον ἀγαγεῖν πρὸς σέ ἔχοντά τι λαλῆσαί σοι.
\vs{19}Ἐπιλαβόμενος δὲ τῆς χειρὸς αὐτοῦ ὁ χιλίαρχος καὶ ἀναχωρήσας κατ᾽ ἰδίαν ἐπυνθάνετο, Τί ἐστιν ὃ ἔχεις ἀπαγγεῖλαί μοι;
\vs{20}Εἶπεν δὲ ὅτι Οἱ Ἰουδαῖοι συνέθεντο τοῦ ἐρωτῆσαί σε ὅπως αὔριον τὸν Παῦλον καταγάγῃς εἰς τὸ συνέδριον ὡς μέλλον τι ἀκριβέστερον πυνθάνεσθαι περὶ αὐτοῦ.
\vs{21}σὺ οὖν μὴ πεισθῇς αὐτοῖς· ἐνεδρεύουσιν γὰρ αὐτὸν ἐξ αὐτῶν ἄνδρες πλείους τεσσεράκοντα, οἵτινες ἀνεθεμάτισαν ἑαυτοὺς μήτε φαγεῖν μήτε πιεῖν ἕως οὗ ἀνέλωσιν αὐτόν, καὶ νῦν εἰσιν ἕτοιμοι προσδεχόμενοι τὴν ἀπὸ σοῦ ἐπαγγελίαν.
\vs{22}Ὁ μὲν οὖν χιλίαρχος ἀπέλυσε τὸν νεανίσκον παραγγείλας Μηδενὶ ἐκλαλῆσαι ὅτι ταῦτα ἐνεφάνισας πρὸς ἐμέ.

\vs{23}Καὶ προσκαλεσάμενός δύο τινας τῶν ἑκατονταρχῶν εἶπεν· Ἑτοιμάσατε στρατιώτας διακοσίους, ὅπως πορευθῶσιν ἕως Καισαρείας, καὶ ἱππεῖς ἑβδομήκοντα καὶ δεξιολάβους διακοσίους ἀπὸ τρίτης ὥρας τῆς νυκτός,
\vs{24}κτήνη τε παραστῆσαι ἵνα ἐπιβιβάσαντες τὸν Παῦλον διασώσωσι πρὸς Φήλικα τὸν ἡγεμόνα,
\vs{25}γράψας ἐπιστολὴν ἔχουσαν τὸν τύπον τοῦτον·
\vs{26}Κλαύδιος Λυσίας Τῷ κρατίστῳ ἡγεμόνι Φήλικι Χαίρειν.
\vs{27}Τὸν ἄνδρα τοῦτον συλλημφθέντα ὑπὸ τῶν Ἰουδαίων καὶ μέλλοντα ἀναιρεῖσθαι ὑπ᾽ αὐτῶν ἐπιστὰς σὺν τῷ στρατεύματι ἐξειλάμην μαθὼν ὅτι Ῥωμαῖός ἐστιν.
\vs{28}βουλόμενός τε ἐπιγνῶναι τὴν αἰτίαν δι᾽ ἣν ἐνεκάλουν αὐτῷ, κατήγαγον εἰς τὸ συνέδριον αὐτῶν
\vs{29}ὃν εὗρον ἐγκαλούμενον περὶ ζητημάτων τοῦ νόμου αὐτῶν, μηδὲν δὲ ἄξιον θανάτου ἢ δεσμῶν ἔχοντα ἔγκλημα.
\vs{30}Μηνυθείσης δέ μοι ἐπιβουλῆς εἰς τὸν ἄνδρα ἔσεσθαι ἐξαυτῆς ἔπεμψα πρὸς σέ παραγγείλας καὶ τοῖς κατηγόροις λέγειν τὰ πρὸς αὐτὸν ἐπὶ σοῦ.
\vs{31}Οἱ μὲν οὖν στρατιῶται κατὰ τὸ διατεταγμένον αὐτοῖς ἀναλαβόντες τὸν Παῦλον ἤγαγον διὰ νυκτὸς εἰς τὴν Ἀντιπατρίδα,
\vs{32}τῇ δὲ ἐπαύριον ἐάσαντες τοὺς ἱππεῖς ἀπέρχεσθαι σὺν αὐτῷ ὑπέστρεψαν εἰς τὴν παρεμβολήν·
\vs{33}οἵτινες εἰσελθόντες εἰς τὴν Καισάρειαν καὶ ἀναδόντες τὴν ἐπιστολὴν τῷ ἡγεμόνι παρέστησαν καὶ τὸν Παῦλον αὐτῷ.
\vs{34}Ἀναγνοὺς δὲ καὶ ἐπερωτήσας ἐκ ποίας ἐπαρχείας ἐστὶν, καὶ πυθόμενος ὅτι ἀπὸ Κιλικίας,
\vs{35}Διακούσομαί σου, ἔφη, Ὅταν καὶ οἱ κατήγοροί σου παραγένωνται· κελεύσας ἐν τῷ πραιτωρίῳ τοῦ Ἡρῴδου φυλάσσεσθαι αὐτόν.

\ch{24}
Μετὰ δὲ πέντε ἡμέρας κατέβη ὁ ἀρχιερεὺς Ἁνανίας μετὰ πρεσβυτέρων τινῶν καὶ ῥήτορος Τερτύλλου τινός, οἵτινες ἐνεφάνισαν τῷ ἡγεμόνι κατὰ τοῦ Παύλου.
\vs{2}Κληθέντος δὲ αὐτοῦ ἤρξατο κατηγορεῖν ὁ Τέρτυλλος λέγων· Πολλῆς εἰρήνης τυγχάνοντες διὰ σοῦ καὶ διορθωμάτων γινομένων τῷ ἔθνει τούτῳ διὰ τῆς σῆς προνοίας,
\vs{3}πάντῃ τε καὶ πανταχοῦ ἀποδεχόμεθα, κράτιστε Φῆλιξ, μετὰ πάσης εὐχαριστίας.
\vs{4}ἵνα δὲ μὴ ἐπὶ πλεῖόν σε ἐνκόπτω, παρακαλῶ ἀκοῦσαί σε ἡμῶν συντόμως τῇ σῇ ἐπιεικείᾳ.
\vs{5}Εὑρόντες γὰρ τὸν ἄνδρα τοῦτον λοιμὸν καὶ κινοῦντα στάσεις πᾶσιν τοῖς Ἰουδαίοις τοῖς κατὰ τὴν οἰκουμένην πρωτοστάτην τε τῆς τῶν Ναζωραίων αἱρέσεως,
\vs{6}ὃς καὶ τὸ ἱερὸν ἐπείρασεν βεβηλῶσαι ὃν καὶ ἐκρατήσαμεν,
\vs{8}παρ᾽ οὗ δυνήσῃ αὐτὸς ἀνακρίνας περὶ πάντων τούτων ἐπιγνῶναι ὧν ἡμεῖς κατηγοροῦμεν αὐτοῦ.
\vs{9}Συνεπέθεντο δὲ καὶ οἱ Ἰουδαῖοι φάσκοντες ταῦτα οὕτως ἔχειν.

\vs{10}Ἀπεκρίθη τε ὁ Παῦλος νεύσαντος αὐτῷ τοῦ ἡγεμόνος λέγειν· Ἐκ πολλῶν ἐτῶν ὄντα σε κριτὴν τῷ ἔθνει τούτῳ ἐπιστάμενος εὐθύμως τὰ περὶ ἐμαυτοῦ ἀπολογοῦμαι,
\vs{11}δυναμένου σου ἐπιγνῶναι ὅτι οὐ πλείους εἰσίν μοι ἡμέραι δώδεκα ἀφ᾽ ἧς ἀνέβην προσκυνήσων εἰς Ἰερουσαλήμ.
\vs{12}καὶ οὔτε ἐν τῷ ἱερῷ εὗρόν με πρός τινα διαλεγόμενον ἢ ἐπίστασιν ποιοῦντα ὄχλου οὔτε ἐν ταῖς συναγωγαῖς οὔτε κατὰ τὴν πόλιν,
\vs{13}οὐδὲ παραστῆσαι δύνανταί σοι περὶ ὧν νυνὶ κατηγοροῦσίν μου.
\vs{14}Ὁμολογῶ δὲ τοῦτό σοι ὅτι κατὰ τὴν Ὁδὸν ἣν λέγουσιν αἵρεσιν, οὕτως λατρεύω τῷ πατρῴῳ Θεῷ πιστεύων πᾶσι τοῖς κατὰ τὸν νόμον καὶ τοῖς ἐν τοῖς προφήταις γεγραμμένοις,
\vs{15}ἐλπίδα ἔχων εἰς τὸν Θεόν ἣν καὶ αὐτοὶ οὗτοι προσδέχονται, ἀνάστασιν μέλλειν ἔσεσθαι δικαίων τε καὶ ἀδίκων.
\vs{16}ἐν τούτῳ καὶ αὐτὸς ἀσκῶ ἀπρόσκοπον συνείδησιν ἔχειν πρὸς τὸν Θεὸν καὶ τοὺς ἀνθρώπους διὰ παντός.
\vs{17}Δι᾽ ἐτῶν δὲ πλειόνων ἐλεημοσύνας ποιήσων εἰς τὸ ἔθνος μου παρεγενόμην καὶ προσφοράς,
\vs{18}ἐν αἷς εὗρόν με ἡγνισμένον ἐν τῷ ἱερῷ οὐ μετὰ ὄχλου οὐδὲ μετὰ θορύβου,
\vs{19}τινὲς δὲ ἀπὸ τῆς Ἀσίας Ἰουδαῖοι, οὓς ἔδει ἐπὶ σοῦ παρεῖναι καὶ κατηγορεῖν εἴ τι ἔχοιεν πρὸς ἐμέ.
\vs{20}ἢ αὐτοὶ οὗτοι εἰπάτωσαν τί εὗρον ἀδίκημα στάντος μου ἐπὶ τοῦ συνεδρίου,
\vs{21}ἢ περὶ μιᾶς ταύτης φωνῆς ἧς ἐκέκραξα ἐν αὐτοῖς ἑστὼς ὅτι Περὶ ἀναστάσεως νεκρῶν ἐγὼ κρίνομαι σήμερον ἐφ᾽ ὑμῶν.

\vs{22}Ἀνεβάλετο δὲ αὐτοὺς ὁ Φῆλιξ, ἀκριβέστερον εἰδὼς τὰ περὶ τῆς Ὁδοῦ εἴπας· Ὅταν Λυσίας ὁ χιλίαρχος καταβῇ, διαγνώσομαι τὰ καθ᾽ ὑμᾶς·
\vs{23}διαταξάμενος τῷ ἑκατοντάρχῃ τηρεῖσθαι αὐτὸν ἔχειν τε ἄνεσιν καὶ μηδένα κωλύειν τῶν ἰδίων αὐτοῦ ὑπηρετεῖν αὐτῷ.

\vs{24}Μετὰ δὲ ἡμέρας τινὰς παραγενόμενος ὁ Φῆλιξ σὺν Δρουσίλλῃ τῇ ἰδίᾳ γυναικὶ οὔσῃ Ἰουδαίᾳ μετεπέμψατο τὸν Παῦλον καὶ ἤκουσεν αὐτοῦ περὶ τῆς εἰς Χριστὸν Ἰησοῦν πίστεως.
\vs{25}διαλεγομένου δὲ αὐτοῦ περὶ δικαιοσύνης καὶ ἐγκρατείας καὶ τοῦ κρίματος τοῦ μέλλοντος, ἔμφοβος γενόμενος ὁ Φῆλιξ ἀπεκρίθη· Τὸ νῦν ἔχον πορεύου, καιρὸν δὲ μεταλαβὼν μετακαλέσομαί σε,
\vs{26}ἅμα καὶ ἐλπίζων ὅτι χρήματα δοθήσεται αὐτῷ ὑπὸ τοῦ Παύλου· διὸ καὶ πυκνότερον αὐτὸν μεταπεμπόμενος ὡμίλει αὐτῷ.
\vs{27}Διετίας δὲ πληρωθείσης ἔλαβεν διάδοχον ὁ Φῆλιξ Πόρκιον Φῆστον, θέλων τε χάριτα καταθέσθαι τοῖς Ἰουδαίοις ὁ Φῆλιξ κατέλιπε τὸν Παῦλον δεδεμένον.

\ch{25}
Φῆστος οὖν ἐπιβὰς τῇ ἐπαρχείᾳ μετὰ τρεῖς ἡμέρας ἀνέβη εἰς Ἱεροσόλυμα ἀπὸ Καισαρείας,
\vs{2}ἐνεφάνισάν τε αὐτῷ οἱ ἀρχιερεῖς καὶ οἱ πρῶτοι τῶν Ἰουδαίων κατὰ τοῦ Παύλου καὶ παρεκάλουν αὐτὸν
\vs{3}αἰτούμενοι χάριν κατ᾽ αὐτοῦ ὅπως μεταπέμψηται αὐτὸν εἰς Ἰερουσαλήμ, ἐνέδραν ποιοῦντες ἀνελεῖν αὐτὸν κατὰ τὴν ὁδόν.
\vs{4}Ὁ μὲν οὖν Φῆστος ἀπεκρίθη τηρεῖσθαι τὸν Παῦλον εἰς Καισάρειαν, ἑαυτὸν δὲ μέλλειν ἐν τάχει ἐκπορεύεσθαι·
\vs{5}Οἱ οὖν ἐν ὑμῖν, φησίν, Δυνατοὶ συνκαταβάντες εἴ τί ἐστιν ἐν τῷ ἀνδρὶ ἄτοπον κατηγορείτωσαν αὐτοῦ.

\vs{6}Διατρίψας δὲ ἐν αὐτοῖς ἡμέρας οὐ πλείους ὀκτὼ ἢ δέκα, καταβὰς εἰς Καισάρειαν, τῇ ἐπαύριον καθίσας ἐπὶ τοῦ βήματος ἐκέλευσεν τὸν Παῦλον ἀχθῆναι.
\vs{7}παραγενομένου δὲ αὐτοῦ περιέστησαν αὐτὸν οἱ ἀπὸ Ἱεροσολύμων καταβεβηκότες Ἰουδαῖοι πολλὰ καὶ βαρέα αἰτιώματα καταφέροντες ἃ οὐκ ἴσχυον ἀποδεῖξαι,
\vs{8}Τοῦ Παύλου ἀπολογουμένου ὅτι Οὔτε εἰς τὸν νόμον τῶν Ἰουδαίων οὔτε εἰς τὸ ἱερὸν οὔτε εἰς Καίσαρά τι ἥμαρτον.
\vs{9}Ὁ Φῆστος δὲ θέλων τοῖς Ἰουδαίοις χάριν καταθέσθαι ἀποκριθεὶς τῷ Παύλῳ εἶπεν· Θέλεις εἰς Ἱεροσόλυμα ἀναβὰς ἐκεῖ περὶ τούτων κριθῆναι ἐπ᾽ ἐμοῦ;
\vs{10}Εἶπεν δὲ ὁ Παῦλος· Ἐπὶ τοῦ βήματος Καίσαρος ἑστὼς εἰμι, οὗ με δεῖ κρίνεσθαι. Ἰουδαίους οὐδὲν ἠδίκησα ὡς καὶ σὺ κάλλιον ἐπιγινώσκεις.
\vs{11}εἰ μὲν οὖν ἀδικῶ καὶ ἄξιον θανάτου πέπραχά τι, οὐ παραιτοῦμαι τὸ ἀποθανεῖν· εἰ δὲ οὐδέν ἐστιν ὧν οὗτοι κατηγοροῦσίν μου, οὐδείς με δύναται αὐτοῖς χαρίσασθαι· Καίσαρα ἐπικαλοῦμαι.
\vs{12}Τότε ὁ Φῆστος συλλαλήσας μετὰ τοῦ συμβουλίου ἀπεκρίθη· Καίσαρα ἐπικέκλησαι, ἐπὶ Καίσαρα πορεύσῃ.

\vs{13}Ἡμερῶν δὲ διαγενομένων τινῶν Ἀγρίππας ὁ βασιλεὺς καὶ Βερνίκη κατήντησαν εἰς Καισάρειαν ἀσπασάμενοι τὸν Φῆστον.
\vs{14}ὡς δὲ πλείους ἡμέρας διέτριβον ἐκεῖ, ὁ Φῆστος τῷ βασιλεῖ ἀνέθετο τὰ κατὰ τὸν Παῦλον λέγων· Ἀνήρ τίς ἐστιν καταλελειμμένος ὑπὸ Φήλικος δέσμιος,
\vs{15}περὶ οὗ γενομένου μου εἰς Ἱεροσόλυμα ἐνεφάνισαν οἱ ἀρχιερεῖς καὶ οἱ πρεσβύτεροι τῶν Ἰουδαίων αἰτούμενοι κατ᾽ αὐτοῦ καταδίκην.
\vs{16}πρὸς οὓς ἀπεκρίθην ὅτι οὐκ ἔστιν ἔθος Ῥωμαίοις χαρίζεσθαί τινα ἄνθρωπον πρὶν ἢ ὁ κατηγορούμενος κατὰ πρόσωπον ἔχοι τοὺς κατηγόρους τόπον τε ἀπολογίας λάβοι περὶ τοῦ ἐγκλήματος.
\vs{17}Συνελθόντων οὖν αὐτῶν ἐνθάδε ἀναβολὴν μηδεμίαν ποιησάμενος τῇ ἑξῆς καθίσας ἐπὶ τοῦ βήματος ἐκέλευσα ἀχθῆναι τὸν ἄνδρα·
\vs{18}περὶ οὗ σταθέντες οἱ κατήγοροι οὐδεμίαν αἰτίαν ἔφερον ὧν ἐγὼ ὑπενόουν πονηρῶν,
\vs{19}ζητήματα δέ τινα περὶ τῆς ἰδίας δεισιδαιμονίας εἶχον πρὸς αὐτὸν καὶ περί τινος Ἰησοῦ τεθνηκότος ὃν ἔφασκεν ὁ Παῦλος ζῆν.
\vs{20}Ἀπορούμενος δὲ ἐγὼ τὴν περὶ τούτων ζήτησιν ἔλεγον εἰ βούλοιτο πορεύεσθαι εἰς Ἱεροσόλυμα κἀκεῖ κρίνεσθαι περὶ τούτων.
\vs{21}τοῦ δὲ Παύλου ἐπικαλεσαμένου τηρηθῆναι αὐτὸν εἰς τὴν τοῦ Σεβαστοῦ διάγνωσιν, ἐκέλευσα τηρεῖσθαι αὐτὸν ἕως οὗ ἀναπέμψω αὐτὸν πρὸς Καίσαρα.
\vs{22}Ἀγρίππας δὲ πρὸς τὸν Φῆστον· Ἐβουλόμην καὶ αὐτὸς τοῦ ἀνθρώπου ἀκοῦσαι. Αὔριον, φησίν, Ἀκούσῃ αὐτοῦ.

\vs{23}Τῇ οὖν ἐπαύριον ἐλθόντος τοῦ Ἀγρίππα καὶ τῆς Βερνίκης μετὰ πολλῆς φαντασίας καὶ εἰσελθόντων εἰς τὸ ἀκροατήριον σύν τε χιλιάρχοις καὶ ἀνδράσιν τοῖς κατ᾽ ἐξοχὴν τῆς πόλεως καὶ κελεύσαντος τοῦ Φήστου ἤχθη ὁ Παῦλος.
\vs{24}Καί φησιν ὁ Φῆστος· Ἀγρίππα βασιλεῦ καὶ πάντες οἱ συμπαρόντες ἡμῖν ἄνδρες, θεωρεῖτε τοῦτον περὶ οὗ ἅπαν τὸ πλῆθος τῶν Ἰουδαίων ἐνέτυχόν μοι ἔν τε Ἱεροσολύμοις καὶ ἐνθάδε βοῶντες μὴ δεῖν αὐτὸν ζῆν μηκέτι.
\vs{25}ἐγὼ δὲ κατελαβόμην μηδὲν ἄξιον αὐτὸν θανάτου πεπραχέναι, αὐτοῦ δὲ τούτου ἐπικαλεσαμένου τὸν Σεβαστὸν ἔκρινα πέμπειν.
\vs{26}περὶ οὗ ἀσφαλές τι γράψαι τῷ κυρίῳ οὐκ ἔχω, διὸ προήγαγον αὐτὸν ἐφ᾽ ὑμῶν καὶ μάλιστα ἐπὶ σοῦ, βασιλεῦ Ἀγρίππα, ὅπως τῆς ἀνακρίσεως γενομένης σχῶ τί γράψω·
\vs{27}ἄλογον γάρ μοι δοκεῖ πέμποντα δέσμιον μὴ καὶ τὰς κατ᾽ αὐτοῦ αἰτίας σημᾶναι.

\ch{26}
Ἀγρίππας δὲ πρὸς τὸν Παῦλον ἔφη· Ἐπιτρέπεταί σοι περὶ σεαυτοῦ λέγειν. Τότε ὁ Παῦλος ἐκτείνας τὴν χεῖρα ἀπελογεῖτο·
\vs{2}Περὶ πάντων ὧν ἐγκαλοῦμαι ὑπὸ Ἰουδαίων, βασιλεῦ Ἀγρίππα, ἥγημαι ἐμαυτὸν μακάριον ἐπὶ σοῦ μέλλων σήμερον ἀπολογεῖσθαι
\vs{3}μάλιστα γνώστην ὄντα σε πάντων τῶν κατὰ Ἰουδαίους ἐθῶν τε καὶ ζητημάτων, διὸ δέομαι μακροθύμως ἀκοῦσαί μου.
\vs{4}Τὴν μὲν οὖν βίωσίν μου τὴν ἐκ νεότητος τὴν ἀπ᾽ ἀρχῆς γενομένην ἐν τῷ ἔθνει μου ἔν τε Ἱεροσολύμοις ἴσασι πάντες οἱ Ἰουδαῖοι
\vs{5}προγινώσκοντές με ἄνωθεν, ἐὰν θέλωσι μαρτυρεῖν, ὅτι κατὰ τὴν ἀκριβεστάτην αἵρεσιν τῆς ἡμετέρας θρησκείας ἔζησα Φαρισαῖος.
\vs{6}Καὶ νῦν ἐπ᾽ ἐλπίδι τῆς εἰς τοὺς πατέρας ἡμῶν ἐπαγγελίας γενομένης ὑπὸ τοῦ Θεοῦ ἕστηκα κρινόμενος,
\vs{7}εἰς ἣν τὸ δωδεκάφυλον ἡμῶν ἐν ἐκτενείᾳ νύκτα καὶ ἡμέραν λατρεῦον ἐλπίζει καταντῆσαι, περὶ ἧς ἐλπίδος ἐγκαλοῦμαι ὑπὸ Ἰουδαίων, βασιλεῦ.
\vs{8}τί ἄπιστον κρίνεται παρ᾽ ὑμῖν εἰ ὁ Θεὸς νεκροὺς ἐγείρει;
\vs{9}Ἐγὼ μὲν οὖν ἔδοξα ἐμαυτῷ πρὸς τὸ ὄνομα Ἰησοῦ τοῦ Ναζωραίου δεῖν πολλὰ ἐναντία πρᾶξαι,
\vs{10}ὃ καὶ ἐποίησα ἐν Ἱεροσολύμοις, καὶ πολλούς τε τῶν ἁγίων ἐγὼ ἐν φυλακαῖς κατέκλεισα τὴν παρὰ τῶν ἀρχιερέων ἐξουσίαν λαβών ἀναιρουμένων τε αὐτῶν κατήνεγκα ψῆφον.
\vs{11}καὶ κατὰ πάσας τὰς συναγωγὰς πολλάκις τιμωρῶν αὐτοὺς ἠνάγκαζον βλασφημεῖν περισσῶς τε ἐμμαινόμενος αὐτοῖς ἐδίωκον ἕως καὶ εἰς τὰς ἔξω πόλεις.
\vs{12}Ἐν οἷς πορευόμενος εἰς τὴν Δαμασκὸν μετ᾽ ἐξουσίας καὶ ἐπιτροπῆς τῆς τῶν ἀρχιερέων
\vs{13}ἡμέρας μέσης κατὰ τὴν ὁδὸν εἶδον, βασιλεῦ, οὐρανόθεν ὑπὲρ τὴν λαμπρότητα τοῦ ἡλίου περιλάμψαν με φῶς καὶ τοὺς σὺν ἐμοὶ πορευομένους.
\vs{14}πάντων τε καταπεσόντων ἡμῶν εἰς τὴν γῆν ἤκουσα φωνὴν λέγουσαν πρός με τῇ Ἑβραΐδι διαλέκτῳ· Σαοὺλ Σαούλ, τί με διώκεις; σκληρόν σοι πρὸς κέντρα λακτίζειν.
\vs{15}Ἐγὼ δὲ εἶπα· Τίς εἶ, Κύριε; Ὁ δὲ Κύριος εἶπεν· Ἐγώ εἰμι Ἰησοῦς ὃν σὺ διώκεις.
\vs{16}ἀλλὰ ἀνάστηθι καὶ στῆθι ἐπὶ τοὺς πόδας σου· εἰς τοῦτο γὰρ ὤφθην σοι, προχειρίσασθαί σε ὑπηρέτην καὶ μάρτυρα ὧν τε εἶδές με ὧν τε ὀφθήσομαί σοι,
\vs{17}ἐξαιρούμενός σε ἐκ τοῦ λαοῦ καὶ ἐκ τῶν ἐθνῶν εἰς οὓς ἐγὼ ἀποστέλλω σε
\vs{18}ἀνοῖξαι ὀφθαλμοὺς αὐτῶν, τοῦ ἐπιστρέψαι ἀπὸ σκότους εἰς φῶς καὶ τῆς ἐξουσίας τοῦ Σατανᾶ ἐπὶ τὸν Θεόν, τοῦ λαβεῖν αὐτοὺς ἄφεσιν ἁμαρτιῶν καὶ κλῆρον ἐν τοῖς ἡγιασμένοις πίστει τῇ εἰς ἐμέ.
\vs{19}Ὅθεν, βασιλεῦ Ἀγρίππα, οὐκ ἐγενόμην ἀπειθὴς τῇ οὐρανίῳ ὀπτασίᾳ
\vs{20}ἀλλὰ τοῖς ἐν Δαμασκῷ πρῶτόν τε καὶ Ἱεροσολύμοις, πᾶσάν τε τὴν χώραν τῆς Ἰουδαίας καὶ τοῖς ἔθνεσιν ἀπήγγελλον μετανοεῖν καὶ ἐπιστρέφειν ἐπὶ τὸν Θεόν, ἄξια τῆς μετανοίας ἔργα πράσσοντας.
\vs{21}ἕνεκα τούτων με Ἰουδαῖοι συλλαβόμενοι ὄντα ἐν τῷ ἱερῷ ἐπειρῶντο διαχειρίσασθαι.
\vs{22}Ἐπικουρίας οὖν τυχὼν τῆς ἀπὸ τοῦ Θεοῦ ἄχρι τῆς ἡμέρας ταύτης ἕστηκα μαρτυρόμενος μικρῷ τε καὶ μεγάλῳ οὐδὲν ἐκτὸς λέγων ὧν τε οἱ προφῆται ἐλάλησαν μελλόντων γίνεσθαι καὶ Μωϋσῆς,
\vs{23}εἰ παθητὸς ὁ Χριστός, εἰ πρῶτος ἐξ ἀναστάσεως νεκρῶν φῶς μέλλει καταγγέλλειν τῷ τε λαῷ καὶ τοῖς ἔθνεσιν.

\vs{24}Ταῦτα δὲ αὐτοῦ ἀπολογουμένου ὁ Φῆστος μεγάλῃ τῇ φωνῇ φησιν· Μαίνῃ, Παῦλε· τὰ πολλά σε γράμματα εἰς μανίαν περιτρέπει.
\vs{25}Ὁ δὲ Παῦλος· Οὐ μαίνομαι, φησίν, Κράτιστε Φῆστε, ἀλλὰ ἀληθείας καὶ σωφροσύνης ῥήματα ἀποφθέγγομαι.
\vs{26}ἐπίσταται γὰρ περὶ τούτων ὁ βασιλεύς πρὸς ὃν καὶ παρρησιαζόμενος λαλῶ, λανθάνειν γὰρ αὐτὸν τι τούτων οὐ πείθομαι οὐθέν· οὐ γάρ ἐστιν ἐν γωνίᾳ πεπραγμένον τοῦτο.
\vs{27}πιστεύεις, βασιλεῦ Ἀγρίππα, τοῖς προφήταις; οἶδα ὅτι πιστεύεις.
\vs{28}Ὁ δὲ Ἀγρίππας πρὸς τὸν Παῦλον· Ἐν ὀλίγῳ με πείθεις Χριστιανὸν ποιῆσαι.
\vs{29}Ὁ δὲ Παῦλος· Εὐξαίμην ἂν τῷ Θεῷ καὶ ἐν ὀλίγῳ καὶ ἐν μεγάλῳ οὐ μόνον σὲ ἀλλὰ καὶ πάντας τοὺς ἀκούοντάς μου σήμερον γενέσθαι τοιούτους ὁποῖος καὶ ἐγώ εἰμι παρεκτὸς τῶν δεσμῶν τούτων.
\vs{30}Ἀνέστη τε ὁ βασιλεὺς καὶ ὁ ἡγεμὼν ἥ τε Βερνίκη καὶ οἱ συνκαθήμενοι αὐτοῖς,
\vs{31}καὶ ἀναχωρήσαντες ἐλάλουν πρὸς ἀλλήλους λέγοντες ὅτι Οὐδὲν θανάτου ἢ δεσμῶν ἄξιον τι πράσσει ὁ ἄνθρωπος οὗτος.
\vs{32}Ἀγρίππας δὲ τῷ Φήστῳ ἔφη· Ἀπολελύσθαι ἐδύνατο ὁ ἄνθρωπος οὗτος εἰ μὴ ἐπεκέκλητο Καίσαρα.

\ch{27}
Ὡς δὲ ἐκρίθη τοῦ ἀποπλεῖν ἡμᾶς εἰς τὴν Ἰταλίαν, παρεδίδουν τόν τε Παῦλον καί τινας ἑτέρους δεσμώτας ἑκατοντάρχῃ ὀνόματι Ἰουλίῳ σπείρης Σεβαστῆς.
\vs{2}ἐπιβάντες δὲ πλοίῳ Ἀδραμυττηνῷ μέλλοντι πλεῖν εἰς τοὺς κατὰ τὴν Ἀσίαν τόπους ἀνήχθημεν ὄντος σὺν ἡμῖν Ἀριστάρχου Μακεδόνος Θεσσαλονικέως.
\vs{3}Τῇ τε ἑτέρᾳ κατήχθημεν εἰς Σιδῶνα, φιλανθρώπως τε ὁ Ἰούλιος τῷ Παύλῳ χρησάμενος ἐπέτρεψεν πρὸς τοὺς φίλους πορευθέντι ἐπιμελείας τυχεῖν.
\vs{4}κἀκεῖθεν ἀναχθέντες ὑπεπλεύσαμεν τὴν Κύπρον διὰ τὸ τοὺς ἀνέμους εἶναι ἐναντίους,
\vs{5}τό τε πέλαγος τὸ κατὰ τὴν Κιλικίαν καὶ Παμφυλίαν διαπλεύσαντες κατήλθομεν εἰς Μύρα τῆς Λυκίας.
\vs{6}Κἀκεῖ εὑρὼν ὁ ἑκατοντάρχης πλοῖον Ἀλεξανδρῖνον πλέον εἰς τὴν Ἰταλίαν ἐνεβίβασεν ἡμᾶς εἰς αὐτό.
\vs{7}Ἐν ἱκαναῖς δὲ ἡμέραις βραδυπλοοῦντες καὶ μόλις γενόμενοι κατὰ τὴν Κνίδον, μὴ προσεῶντος ἡμᾶς τοῦ ἀνέμου ὑπεπλεύσαμεν τὴν Κρήτην κατὰ Σαλμώνην,
\vs{8}μόλις τε παραλεγόμενοι αὐτὴν ἤλθομεν εἰς τόπον τινὰ καλούμενον Καλοὺς Λιμένας ᾧ ἐγγὺς πόλις ἦν Λασαία.

\vs{9}Ἱκανοῦ δὲ χρόνου διαγενομένου καὶ ὄντος ἤδη ἐπισφαλοῦς τοῦ πλοὸς διὰ τὸ καὶ τὴν Νηστείαν ἤδη παρεληλυθέναι παρῄνει ὁ Παῦλος
\vs{10}λέγων αὐτοῖς· Ἄνδρες, θεωρῶ ὅτι μετὰ ὕβρεως καὶ πολλῆς ζημίας οὐ μόνον τοῦ φορτίου καὶ τοῦ πλοίου ἀλλὰ καὶ τῶν ψυχῶν ἡμῶν μέλλειν ἔσεσθαι τὸν πλοῦν.
\vs{11}Ὁ δὲ ἑκατοντάρχης τῷ κυβερνήτῃ καὶ τῷ ναυκλήρῳ μᾶλλον ἐπείθετο ἢ τοῖς ὑπὸ Παύλου λεγομένοις.
\vs{12}ἀνευθέτου δὲ τοῦ λιμένος ὑπάρχοντος πρὸς παραχειμασίαν οἱ πλείονες ἔθεντο βουλὴν ἀναχθῆναι ἐκεῖθεν, εἴ πως δύναιντο καταντήσαντες εἰς Φοίνικα παραχειμάσαι λιμένα τῆς Κρήτης βλέποντα κατὰ λίβα καὶ κατὰ χῶρον.

\vs{13}Ὑποπνεύσαντος δὲ νότου δόξαντες τῆς προθέσεως κεκρατηκέναι, ἄραντες ἆσσον παρελέγοντο τὴν Κρήτην.
\vs{14}μετ᾽ οὐ πολὺ δὲ ἔβαλεν κατ᾽ αὐτῆς ἄνεμος τυφωνικὸς ὁ καλούμενος Εὐρακύλων·
\vs{15}συναρπασθέντος δὲ τοῦ πλοίου καὶ μὴ δυναμένου ἀντοφθαλμεῖν τῷ ἀνέμῳ ἐπιδόντες ἐφερόμεθα.
\vs{16}Νησίον δέ τι ὑποδραμόντες καλούμενον Καῦδα ἰσχύσαμεν μόλις περικρατεῖς γενέσθαι τῆς σκάφης,
\vs{17}ἣν ἄραντες βοηθείαις ἐχρῶντο ὑποζωννύντες τὸ πλοῖον, φοβούμενοί τε μὴ εἰς τὴν Σύρτιν ἐκπέσωσιν, χαλάσαντες τὸ σκεῦος, οὕτως ἐφέροντο.
\vs{18}Σφοδρῶς δὲ χειμαζομένων ἡμῶν τῇ ἑξῆς ἐκβολὴν ἐποιοῦντο
\vs{19}καὶ τῇ τρίτῃ αὐτόχειρες τὴν σκευὴν τοῦ πλοίου ἔρριψαν.
\vs{20}μήτε δὲ ἡλίου μήτε ἄστρων ἐπιφαινόντων ἐπὶ πλείονας ἡμέρας, χειμῶνός τε οὐκ ὀλίγου ἐπικειμένου, λοιπὸν περιῃρεῖτο ἐλπὶς πᾶσα τοῦ σῴζεσθαι ἡμᾶς.

\vs{21}Πολλῆς τε ἀσιτίας ὑπαρχούσης τότε σταθεὶς ὁ Παῦλος ἐν μέσῳ αὐτῶν εἶπεν· Ἔδει μέν, ὦ ἄνδρες, πειθαρχήσαντάς μοι μὴ ἀνάγεσθαι ἀπὸ τῆς Κρήτης κερδῆσαί τε τὴν ὕβριν ταύτην καὶ τὴν ζημίαν.
\vs{22}καὶ τὰ νῦν παραινῶ ὑμᾶς εὐθυμεῖν· ἀποβολὴ γὰρ ψυχῆς οὐδεμία ἔσται ἐξ ὑμῶν πλὴν τοῦ πλοίου.
\vs{23}παρέστη γάρ μοι ταύτῃ τῇ νυκτὶ τοῦ Θεοῦ, οὗ εἰμι ἐγώ ᾧ καὶ λατρεύω, ἄγγελος
\vs{24}λέγων· Μὴ φοβοῦ, Παῦλε, Καίσαρί σε δεῖ παραστῆναι, καὶ ἰδοὺ κεχάρισταί σοι ὁ Θεὸς πάντας τοὺς πλέοντας μετὰ σοῦ.
\vs{25}Διὸ εὐθυμεῖτε, ἄνδρες· πιστεύω γὰρ τῷ Θεῷ ὅτι οὕτως ἔσται καθ᾽ ὃν τρόπον λελάληταί μοι.
\vs{26}εἰς νῆσον δέ τινα δεῖ ἡμᾶς ἐκπεσεῖν.

\vs{27}Ὡς δὲ τεσσαρεσκαιδεκάτη νὺξ ἐγένετο διαφερομένων ἡμῶν ἐν τῷ Ἀδρίᾳ, κατὰ μέσον τῆς νυκτὸς ὑπενόουν οἱ ναῦται προσάγειν τινὰ αὐτοῖς χώραν.
\vs{28}καὶ βολίσαντες εὗρον ὀργυιὰς εἴκοσι, βραχὺ δὲ διαστήσαντες καὶ πάλιν βολίσαντες εὗρον ὀργυιὰς δεκαπέντε·
\vs{29}φοβούμενοί τε μή που κατὰ τραχεῖς τόπους ἐκπέσωμεν, ἐκ πρύμνης ῥίψαντες ἀγκύρας τέσσαρας ηὔχοντο ἡμέραν γενέσθαι.
\vs{30}Τῶν δὲ ναυτῶν ζητούντων φυγεῖν ἐκ τοῦ πλοίου καὶ χαλασάντων τὴν σκάφην εἰς τὴν θάλασσαν προφάσει ὡς ἐκ πρῴρης ἀγκύρας μελλόντων ἐκτείνειν,
\vs{31}εἶπεν ὁ Παῦλος τῷ ἑκατοντάρχῃ καὶ τοῖς στρατιώταις· Ἐὰν μὴ οὗτοι μείνωσιν ἐν τῷ πλοίῳ, ὑμεῖς σωθῆναι οὐ δύνασθε.
\vs{32}τότε ἀπέκοψαν οἱ στρατιῶται τὰ σχοινία τῆς σκάφης καὶ εἴασαν αὐτὴν ἐκπεσεῖν.

\vs{33}Ἄχρι δὲ οὗ ἡμέρα ἤμελλεν γίνεσθαι, παρεκάλει ὁ Παῦλος ἅπαντας μεταλαβεῖν τροφῆς λέγων· Τεσσαρεσκαιδεκάτην σήμερον ἡμέραν προσδοκῶντες ἄσιτοι διατελεῖτε μηθὲν προσλαβόμενοι.
\vs{34}διὸ παρακαλῶ ὑμᾶς μεταλαβεῖν τροφῆς· τοῦτο γὰρ πρὸς τῆς ὑμετέρας σωτηρίας ὑπάρχει, οὐδενὸς γὰρ ὑμῶν θρὶξ ἀπὸ τῆς κεφαλῆς ἀπολεῖται.
\vs{35}Εἴπας δὲ ταῦτα καὶ λαβὼν ἄρτον εὐχαρίστησεν τῷ Θεῷ ἐνώπιον πάντων καὶ κλάσας ἤρξατο ἐσθίειν.
\vs{36}εὔθυμοι δὲ γενόμενοι πάντες καὶ αὐτοὶ προσελάβοντο τροφῆς.
\vs{37}ἤμεθα δὲ αἱ πᾶσαι ψυχαὶ ἐν τῷ πλοίῳ διακόσιαι ἑβδομήκοντα ἕξ.
\vs{38}κορεσθέντες δὲ τροφῆς ἐκούφιζον τὸ πλοῖον ἐκβαλλόμενοι τὸν σῖτον εἰς τὴν θάλασσαν.

\vs{39}Ὅτε δὲ ἡμέρα ἐγένετο, τὴν γῆν οὐκ ἐπεγίνωσκον, κόλπον δέ τινα κατενόουν ἔχοντα αἰγιαλὸν εἰς ὃν ἐβουλεύοντο εἰ δύναιντο ἐξῶσαι τὸ πλοῖον.
\vs{40}καὶ τὰς ἀγκύρας περιελόντες εἴων εἰς τὴν θάλασσαν, ἅμα ἀνέντες τὰς ζευκτηρίας τῶν πηδαλίων καὶ ἐπάραντες τὸν ἀρτέμωνα τῇ πνεούσῃ κατεῖχον εἰς τὸν αἰγιαλόν.
\vs{41}περιπεσόντες δὲ εἰς τόπον διθάλασσον ἐπέκειλαν τὴν ναῦν καὶ ἡ μὲν πρῷρα ἐρείσασα ἔμεινεν ἀσάλευτος, ἡ δὲ πρύμνα ἐλύετο ὑπὸ τῆς βίας τῶν κυμάτων.
\vs{42}Τῶν δὲ στρατιωτῶν βουλὴ ἐγένετο ἵνα τοὺς δεσμώτας ἀποκτείνωσιν, μή τις ἐκκολυμβήσας διαφύγῃ.
\vs{43}ὁ δὲ ἑκατοντάρχης βουλόμενος διασῶσαι τὸν Παῦλον ἐκώλυσεν αὐτοὺς τοῦ βουλήματος, ἐκέλευσέν τε τοὺς δυναμένους κολυμβᾶν ἀπορίψαντας πρώτους ἐπὶ τὴν γῆν ἐξιέναι
\vs{44}καὶ τοὺς λοιποὺς οὓς μὲν ἐπὶ σανίσιν, οὓς δὲ ἐπί τινων τῶν ἀπὸ τοῦ πλοίου. καὶ οὕτως ἐγένετο πάντας διασωθῆναι ἐπὶ τὴν γῆν.

\ch{28}
Καὶ διασωθέντες τότε ἐπέγνωμεν ὅτι Μελίτη ἡ νῆσος καλεῖται.
\vs{2}οἵ τε βάρβαροι παρεῖχον οὐ τὴν τυχοῦσαν φιλανθρωπίαν ἡμῖν, ἅψαντες γὰρ πυρὰν προσελάβοντο πάντας ἡμᾶς διὰ τὸν ὑετὸν τὸν ἐφεστῶτα καὶ διὰ τὸ ψῦχος.
\vs{3}Συστρέψαντος δὲ τοῦ Παύλου φρυγάνων τι πλῆθος καὶ ἐπιθέντος ἐπὶ τὴν πυράν, ἔχιδνα ἀπὸ τῆς θέρμης ἐξελθοῦσα καθῆψεν τῆς χειρὸς αὐτοῦ.
\vs{4}ὡς δὲ εἶδον οἱ βάρβαροι κρεμάμενον τὸ θηρίον ἐκ τῆς χειρὸς αὐτοῦ, πρὸς ἀλλήλους ἔλεγον· Πάντως φονεύς ἐστιν ὁ ἄνθρωπος οὗτος ὃν διασωθέντα ἐκ τῆς θαλάσσης ἡ Δίκη ζῆν οὐκ εἴασεν.
\vs{5}ὁ μὲν οὖν ἀποτινάξας τὸ θηρίον εἰς τὸ πῦρ ἔπαθεν οὐδὲν κακόν,
\vs{6}οἱ δὲ προσεδόκων αὐτὸν μέλλειν πίμπρασθαι ἢ καταπίπτειν ἄφνω νεκρόν. ἐπὶ πολὺ δὲ αὐτῶν προσδοκώντων καὶ θεωρούντων μηδὲν ἄτοπον εἰς αὐτὸν γινόμενον μεταβαλόμενοι ἔλεγον αὐτὸν εἶναι θεόν.

\vs{7}Ἐν δὲ τοῖς περὶ τὸν τόπον ἐκεῖνον ὑπῆρχεν χωρία τῷ πρώτῳ τῆς νήσου ὀνόματι Ποπλίῳ, ὃς ἀναδεξάμενος ἡμᾶς τρεῖς ἡμέρας φιλοφρόνως ἐξένισεν.
\vs{8}ἐγένετο δὲ τὸν πατέρα τοῦ Ποπλίου πυρετοῖς καὶ δυσεντερίῳ συνεχόμενον κατακεῖσθαι, πρὸς ὃν ὁ Παῦλος εἰσελθὼν καὶ προσευξάμενος ἐπιθεὶς τὰς χεῖρας αὐτῷ ἰάσατο αὐτόν.
\vs{9}τούτου δὲ γενομένου καὶ οἱ λοιποὶ οἱ ἐν τῇ νήσῳ ἔχοντες ἀσθενείας προσήρχοντο καὶ ἐθεραπεύοντο,
\vs{10}οἳ καὶ πολλαῖς τιμαῖς ἐτίμησαν ἡμᾶς καὶ ἀναγομένοις ἐπέθεντο τὰ πρὸς τὰς χρείας.

\vs{11}Μετὰ δὲ τρεῖς μῆνας ἀνήχθημεν ἐν πλοίῳ παρακεχειμακότι ἐν τῇ νήσῳ, Ἀλεξανδρινῷ, παρασήμῳ Διοσκούροις.
\vs{12}καὶ καταχθέντες εἰς Συρακούσας ἐπεμείναμεν ἡμέρας τρεῖς,
\vs{13}ὅθεν περιελόντες κατηντήσαμεν εἰς Ῥήγιον. καὶ μετὰ μίαν ἡμέραν ἐπιγενομένου νότου δευτεραῖοι ἤλθομεν εἰς Ποτιόλους,
\vs{14}οὗ εὑρόντες ἀδελφοὺς παρεκλήθημεν παρ᾽ αὐτοῖς ἐπιμεῖναι ἡμέρας ἑπτά· καὶ οὕτως εἰς τὴν Ῥώμην ἤλθαμεν.
\vs{15}Κἀκεῖθεν οἱ ἀδελφοὶ ἀκούσαντες τὰ περὶ ἡμῶν ἦλθαν εἰς ἀπάντησιν ἡμῖν ἄχρι Ἀππίου Φόρου καὶ Τριῶν Ταβερνῶν, οὓς ἰδὼν ὁ Παῦλος εὐχαριστήσας τῷ Θεῷ ἔλαβε θάρσος.
\vs{16}Ὅτε δὲ εἰσήλθομεν εἰς Ῥώμην, ἐπετράπη τῷ Παύλῳ μένειν καθ᾽ ἑαυτὸν σὺν τῷ φυλάσσοντι αὐτὸν στρατιώτῃ.

\vs{17}Ἐγένετο δὲ μετὰ ἡμέρας τρεῖς συνκαλέσασθαι αὐτὸν τοὺς ὄντας τῶν Ἰουδαίων πρώτους· συνελθόντων δὲ αὐτῶν ἔλεγεν πρὸς αὐτούς· Ἐγώ, ἄνδρες ἀδελφοί, οὐδὲν ἐναντίον ποιήσας τῷ λαῷ ἢ τοῖς ἔθεσι τοῖς πατρῴοις δέσμιος ἐξ Ἱεροσολύμων παρεδόθην εἰς τὰς χεῖρας τῶν Ῥωμαίων,
\vs{18}οἵτινες ἀνακρίναντές με ἐβούλοντο ἀπολῦσαι διὰ τὸ μηδεμίαν αἰτίαν θανάτου ὑπάρχειν ἐν ἐμοί.
\vs{19}ἀντιλεγόντων δὲ τῶν Ἰουδαίων ἠναγκάσθην ἐπικαλέσασθαι Καίσαρα οὐχ ὡς τοῦ ἔθνους μου ἔχων τι κατηγορεῖν.
\vs{20}διὰ ταύτην οὖν τὴν αἰτίαν παρεκάλεσα ὑμᾶς ἰδεῖν καὶ προσλαλῆσαι, ἕνεκεν γὰρ τῆς ἐλπίδος τοῦ Ἰσραὴλ τὴν ἅλυσιν ταύτην περίκειμαι.
\vs{21}Οἱ δὲ πρὸς αὐτὸν εἶπαν· Ἡμεῖς οὔτε γράμματα περὶ σοῦ ἐδεξάμεθα ἀπὸ τῆς Ἰουδαίας οὔτε παραγενόμενός τις τῶν ἀδελφῶν ἀπήγγειλεν ἢ ἐλάλησέν τι περὶ σοῦ πονηρόν.
\vs{22}ἀξιοῦμεν δὲ παρὰ σοῦ ἀκοῦσαι ἃ φρονεῖς, περὶ μὲν γὰρ τῆς αἱρέσεως ταύτης γνωστὸν ἡμῖν ἐστιν ὅτι πανταχοῦ ἀντιλέγεται.

\vs{23}Ταξάμενοι δὲ αὐτῷ ἡμέραν ἦλθον πρὸς αὐτὸν εἰς τὴν ξενίαν πλείονες οἷς ἐξετίθετο διαμαρτυρόμενος τὴν βασιλείαν τοῦ Θεοῦ, πείθων τε αὐτοὺς περὶ τοῦ Ἰησοῦ ἀπό τε τοῦ νόμου Μωϋσέως καὶ τῶν προφητῶν, ἀπὸ πρωῒ ἕως ἑσπέρας.
\vs{24}Καὶ οἱ μὲν ἐπείθοντο τοῖς λεγομένοις, οἱ δὲ ἠπίστουν·
\vs{25}ἀσύμφωνοι δὲ ὄντες πρὸς ἀλλήλους ἀπελύοντο εἰπόντος τοῦ Παύλου ῥῆμα ἓν, ὅτι Καλῶς τὸ Πνεῦμα τὸ Ἅγιον ἐλάλησεν διὰ Ἠσαΐου τοῦ προφήτου πρὸς τοὺς πατέρας ὑμῶν
\vs{26}λέγων· 
\begin{poetryblock}

\begin{quote}Πορεύθητι πρὸς τὸν λαὸν τοῦτον καὶ εἰπόν·\end{quote} 

\begin{quote}Ἀκοῇ ἀκούσετε καὶ οὐ μὴ συνῆτε\end{quote} 

\begin{quote}καὶ βλέποντες βλέψετε καὶ οὐ μὴ ἴδητε·\end{quote}

\begin{quote} \vs{27}ἐπαχύνθη γὰρ ἡ καρδία τοῦ λαοῦ τούτου\end{quote} 

\begin{quote}καὶ τοῖς ὠσὶν βαρέως ἤκουσαν\end{quote} 

\begin{quote}καὶ τοὺς ὀφθαλμοὺς αὐτῶν ἐκάμμυσαν·\end{quote} 

\begin{quote}μήποτε ἴδωσιν τοῖς ὀφθαλμοῖς\end{quote} 

\begin{quote}καὶ τοῖς ὠσὶν ἀκούσωσιν\end{quote} 

\begin{quote}καὶ τῇ καρδίᾳ συνῶσιν\end{quote} 

\begin{quote}καὶ ἐπιστρέψωσιν, καὶ ἰάσομαι αὐτούς.\end{quote}
\end{poetryblock}

\vs{28}Γνωστὸν οὖν ἔστω ὑμῖν ὅτι τοῖς ἔθνεσιν ἀπεστάλη τοῦτο τὸ σωτήριον τοῦ Θεοῦ· αὐτοὶ καὶ ἀκούσονται.

\vs{30}Ἐνέμεινεν δὲ διετίαν ὅλην ἐν ἰδίῳ μισθώματι καὶ ἀπεδέχετο πάντας τοὺς εἰσπορευομένους πρὸς αὐτόν,
\vs{31}κηρύσσων τὴν βασιλείαν τοῦ Θεοῦ καὶ διδάσκων τὰ περὶ τοῦ Κυρίου Ἰησοῦ Χριστοῦ μετὰ πάσης παρρησίας ἀκωλύτως.


\def\book{ΠΡΟΣ ΡΩΜΑΙΟΥΣ}
\biblebook{ΠΡΟΣ ΡΩΜΑΙΟΥΣ}


\lettrine[lines=2, loversize=0.2, nindent=0em, findent=.25em]{\textcolor{bookheadingcolor}{Π}}{αῦλος} δοῦλος Χριστοῦ Ἰησοῦ, κλητὸς ἀπόστολος ἀφωρισμένος εἰς εὐαγγέλιον Θεοῦ,
\vs{2}ὃ προεπηγγείλατο διὰ τῶν προφητῶν αὐτοῦ ἐν γραφαῖς ἁγίαις
\vs{3}περὶ τοῦ Υἱοῦ αὐτοῦ τοῦ γενομένου ἐκ σπέρματος Δαυὶδ κατὰ σάρκα,
\vs{4}τοῦ ὁρισθέντος Υἱοῦ Θεοῦ ἐν δυνάμει κατὰ πνεῦμα ἁγιωσύνης ἐξ ἀναστάσεως νεκρῶν, Ἰησοῦ Χριστοῦ τοῦ Κυρίου ἡμῶν,
\vs{5}δι᾽ οὗ ἐλάβομεν χάριν καὶ ἀποστολὴν εἰς ὑπακοὴν πίστεως ἐν πᾶσιν τοῖς ἔθνεσιν ὑπὲρ τοῦ ὀνόματος αὐτοῦ,
\vs{6}ἐν οἷς ἐστε καὶ ὑμεῖς κλητοὶ Ἰησοῦ Χριστοῦ,
\vs{7}Πᾶσιν τοῖς οὖσιν ἐν Ῥώμῃ ἀγαπητοῖς Θεοῦ, κλητοῖς ἁγίοις, Χάρις ὑμῖν καὶ εἰρήνη ἀπὸ Θεοῦ Πατρὸς ἡμῶν καὶ Κυρίου Ἰησοῦ Χριστοῦ.

\vs{8}Πρῶτον μὲν εὐχαριστῶ τῷ Θεῷ μου διὰ Ἰησοῦ Χριστοῦ περὶ πάντων ὑμῶν ὅτι ἡ πίστις ὑμῶν καταγγέλλεται ἐν ὅλῳ τῷ κόσμῳ.
\vs{9}μάρτυς γάρ μού ἐστιν ὁ Θεός, ᾧ λατρεύω ἐν τῷ πνεύματί μου ἐν τῷ εὐαγγελίῳ τοῦ Υἱοῦ αὐτοῦ, ὡς ἀδιαλείπτως μνείαν ὑμῶν ποιοῦμαι
\vs{10}πάντοτε ἐπὶ τῶν προσευχῶν μου δεόμενος εἴ πως ἤδη ποτὲ εὐοδωθήσομαι ἐν τῷ θελήματι τοῦ Θεοῦ ἐλθεῖν πρὸς ὑμᾶς.
\vs{11}ἐπιποθῶ γὰρ ἰδεῖν ὑμᾶς, ἵνα τι μεταδῶ χάρισμα ὑμῖν πνευματικὸν εἰς τὸ στηριχθῆναι ὑμᾶς,
\vs{12}τοῦτο δέ ἐστιν συμπαρακληθῆναι ἐν ὑμῖν διὰ τῆς ἐν ἀλλήλοις πίστεως ὑμῶν τε καὶ ἐμοῦ.
\vs{13}Οὐ θέλω δὲ ὑμᾶς ἀγνοεῖν, ἀδελφοί, ὅτι πολλάκις προεθέμην ἐλθεῖν πρὸς ὑμᾶς, καὶ ἐκωλύθην ἄχρι τοῦ δεῦρο, ἵνα τινὰ καρπὸν σχῶ καὶ ἐν ὑμῖν καθὼς καὶ ἐν τοῖς λοιποῖς ἔθνεσιν.
\vs{14}Ἕλλησίν τε καὶ Βαρβάροις, σοφοῖς τε καὶ ἀνοήτοις ὀφειλέτης εἰμί,
\vs{15}οὕτως τὸ κατ᾽ ἐμὲ πρόθυμον καὶ ὑμῖν τοῖς ἐν Ῥώμῃ εὐαγγελίσασθαι.

\vs{16}Οὐ γὰρ ἐπαισχύνομαι τὸ εὐαγγέλιον, δύναμις γὰρ Θεοῦ ἐστιν εἰς σωτηρίαν παντὶ τῷ πιστεύοντι, Ἰουδαίῳ τε πρῶτον καὶ Ἕλληνι.
\vs{17}δικαιοσύνη γὰρ Θεοῦ ἐν αὐτῷ ἀποκαλύπτεται ἐκ πίστεως εἰς πίστιν, καθὼς γέγραπται· Ὁ δὲ δίκαιος ἐκ πίστεως ζήσεται.

\vs{18}Ἀποκαλύπτεται γὰρ ὀργὴ Θεοῦ ἀπ᾽ οὐρανοῦ ἐπὶ πᾶσαν ἀσέβειαν καὶ ἀδικίαν ἀνθρώπων τῶν τὴν ἀλήθειαν ἐν ἀδικίᾳ κατεχόντων,
\vs{19}διότι τὸ γνωστὸν τοῦ Θεοῦ φανερόν ἐστιν ἐν αὐτοῖς· ὁ θεὸς γὰρ αὐτοῖς ἐφανέρωσεν.
\vs{20}τὰ γὰρ ἀόρατα αὐτοῦ ἀπὸ κτίσεως κόσμου τοῖς ποιήμασιν νοούμενα καθορᾶται, ἥ τε ἀΐδιος αὐτοῦ δύναμις καὶ θειότης, εἰς τὸ εἶναι αὐτοὺς ἀναπολογήτους,
\vs{21}Διότι γνόντες τὸν Θεὸν οὐχ ὡς Θεὸν ἐδόξασαν ἢ ηὐχαρίστησαν, ἀλλὰ ἐματαιώθησαν ἐν τοῖς διαλογισμοῖς αὐτῶν καὶ ἐσκοτίσθη ἡ ἀσύνετος αὐτῶν καρδία.
\vs{22}φάσκοντες εἶναι σοφοὶ ἐμωράνθησαν
\vs{23}καὶ ἤλλαξαν τὴν δόξαν τοῦ ἀφθάρτου Θεοῦ ἐν ὁμοιώματι εἰκόνος φθαρτοῦ ἀνθρώπου καὶ πετεινῶν καὶ τετραπόδων καὶ ἑρπετῶν.
\vs{24}Διὸ παρέδωκεν αὐτοὺς ὁ Θεὸς ἐν ταῖς ἐπιθυμίαις τῶν καρδιῶν αὐτῶν εἰς ἀκαθαρσίαν τοῦ ἀτιμάζεσθαι τὰ σώματα αὐτῶν ἐν αὐτοῖς·
\vs{25}οἵτινες μετήλλαξαν τὴν ἀλήθειαν τοῦ Θεοῦ ἐν τῷ ψεύδει καὶ ἐσεβάσθησαν καὶ ἐλάτρευσαν τῇ κτίσει παρὰ τὸν Κτίσαντα, ὅς ἐστιν εὐλογητὸς εἰς τοὺς αἰῶνας, ἀμήν.
\vs{26}Διὰ τοῦτο παρέδωκεν αὐτοὺς ὁ Θεὸς εἰς πάθη ἀτιμίας, αἵ τε γὰρ θήλειαι αὐτῶν μετήλλαξαν τὴν φυσικὴν χρῆσιν εἰς τὴν παρὰ φύσιν,
\vs{27}ὁμοίως τε καὶ οἱ ἄρσενες ἀφέντες τὴν φυσικὴν χρῆσιν τῆς θηλείας ἐξεκαύθησαν ἐν τῇ ὀρέξει αὐτῶν εἰς ἀλλήλους, ἄρσενες ἐν ἄρσεσιν τὴν ἀσχημοσύνην κατεργαζόμενοι καὶ τὴν ἀντιμισθίαν ἣν ἔδει τῆς πλάνης αὐτῶν ἐν ἑαυτοῖς ἀπολαμβάνοντες.
\vs{28}Καὶ καθὼς οὐκ ἐδοκίμασαν τὸν Θεὸν ἔχειν ἐν ἐπιγνώσει, παρέδωκεν αὐτοὺς ὁ Θεὸς εἰς ἀδόκιμον νοῦν, ποιεῖν τὰ μὴ καθήκοντα,
\vs{29}πεπληρωμένους πάσῃ ἀδικίᾳ πονηρίᾳ πλεονεξίᾳ κακίᾳ, μεστοὺς φθόνου φόνου ἔριδος δόλου κακοηθείας, ψιθυριστάς
\vs{30}καταλάλους θεοστυγεῖς ὑβριστάς ὑπερηφάνους ἀλαζόνας, ἐφευρετὰς κακῶν, γονεῦσιν ἀπειθεῖς,
\vs{31}ἀσυνέτους ἀσυνθέτους ἀστόργους ἀνελεήμονας·
\vs{32}οἵτινες τὸ δικαίωμα τοῦ Θεοῦ ἐπιγνόντες ὅτι οἱ τὰ τοιαῦτα πράσσοντες ἄξιοι θανάτου εἰσίν, οὐ μόνον αὐτὰ ποιοῦσιν ἀλλὰ καὶ συνευδοκοῦσιν τοῖς πράσσουσιν.

\ch{2}
Διὸ ἀναπολόγητος εἶ, ὦ ἄνθρωπε πᾶς ὁ κρίνων· ἐν ᾧ γὰρ κρίνεις τὸν ἕτερον, σεαυτὸν κατακρίνεις, τὰ γὰρ αὐτὰ πράσσεις ὁ κρίνων.
\vs{2}οἴδαμεν δὲ ὅτι τὸ κρίμα τοῦ Θεοῦ ἐστιν κατὰ ἀλήθειαν ἐπὶ τοὺς τὰ τοιαῦτα πράσσοντας.
\vs{3}λογίζῃ δὲ τοῦτο, ὦ ἄνθρωπε ὁ κρίνων τοὺς τὰ τοιαῦτα πράσσοντας καὶ ποιῶν αὐτά, ὅτι σὺ ἐκφεύξῃ τὸ κρίμα τοῦ Θεοῦ;
\vs{4}ἢ τοῦ πλούτου τῆς χρηστότητος αὐτοῦ καὶ τῆς ἀνοχῆς καὶ τῆς μακροθυμίας καταφρονεῖς, ἀγνοῶν ὅτι τὸ χρηστὸν τοῦ Θεοῦ εἰς μετάνοιάν σε ἄγει;
\vs{5}Κατὰ δὲ τὴν σκληρότητά σου καὶ ἀμετανόητον καρδίαν θησαυρίζεις σεαυτῷ ὀργὴν ἐν ἡμέρᾳ ὀργῆς καὶ ἀποκαλύψεως δικαιοκρισίας τοῦ Θεοῦ
\vs{6}ὃς Ἀποδώσει ἑκάστῳ κατὰ τὰ ἔργα αὐτοῦ·
\vs{7}τοῖς μὲν καθ᾽ ὑπομονὴν ἔργου ἀγαθοῦ δόξαν καὶ τιμὴν καὶ ἀφθαρσίαν ζητοῦσιν ζωὴν αἰώνιον,
\vs{8}τοῖς δὲ ἐξ ἐριθείας καὶ ἀπειθοῦσι τῇ ἀληθείᾳ πειθομένοις δὲ τῇ ἀδικίᾳ ὀργὴ καὶ θυμός.
\vs{9}θλῖψις καὶ στενοχωρία ἐπὶ πᾶσαν ψυχὴν ἀνθρώπου τοῦ κατεργαζομένου τὸ κακόν, Ἰουδαίου τε πρῶτον καὶ Ἕλληνος·
\vs{10}δόξα δὲ καὶ τιμὴ καὶ εἰρήνη παντὶ τῷ ἐργαζομένῳ τὸ ἀγαθόν, Ἰουδαίῳ τε πρῶτον καὶ Ἕλληνι·
\vs{11}οὐ γάρ ἐστιν προσωπολημψία παρὰ τῷ Θεῷ.

\vs{12}Ὅσοι γὰρ ἀνόμως ἥμαρτον, ἀνόμως καὶ ἀπολοῦνται, καὶ ὅσοι ἐν νόμῳ ἥμαρτον, διὰ νόμου κριθήσονται·
\vs{13}οὐ γὰρ οἱ ἀκροαταὶ νόμου δίκαιοι παρὰ τῷ Θεῷ, ἀλλ᾽ οἱ ποιηταὶ νόμου δικαιωθήσονται.
\vs{14}Ὅταν γὰρ ἔθνη τὰ μὴ νόμον ἔχοντα φύσει τὰ τοῦ νόμου ποιῶσιν, οὗτοι νόμον μὴ ἔχοντες ἑαυτοῖς εἰσιν νόμος·
\vs{15}οἵτινες ἐνδείκνυνται τὸ ἔργον τοῦ νόμου γραπτὸν ἐν ταῖς καρδίαις αὐτῶν, συμμαρτυρούσης αὐτῶν τῆς συνειδήσεως καὶ μεταξὺ ἀλλήλων τῶν λογισμῶν κατηγορούντων ἢ καὶ ἀπολογουμένων,
\vs{16}ἐν ἡμέρᾳ ὅτε κρίνει ὁ Θεὸς τὰ κρυπτὰ τῶν ἀνθρώπων κατὰ τὸ εὐαγγέλιόν μου διὰ Χριστοῦ Ἰησοῦ.

\vs{17}Εἰ δὲ σὺ Ἰουδαῖος ἐπονομάζῃ καὶ ἐπαναπαύῃ νόμῳ καὶ καυχᾶσαι ἐν Θεῷ
\vs{18}καὶ γινώσκεις τὸ θέλημα καὶ δοκιμάζεις τὰ διαφέροντα κατηχούμενος ἐκ τοῦ νόμου,
\vs{19}πέποιθάς τε σεαυτὸν ὁδηγὸν εἶναι τυφλῶν, φῶς τῶν ἐν σκότει,
\vs{20}παιδευτὴν ἀφρόνων, διδάσκαλον νηπίων, ἔχοντα τὴν μόρφωσιν τῆς γνώσεως καὶ τῆς ἀληθείας ἐν τῷ νόμῳ·
\vs{21}ὁ οὖν διδάσκων ἕτερον σεαυτὸν οὐ διδάσκεις; ὁ κηρύσσων μὴ κλέπτειν κλέπτεις;
\vs{22}ὁ λέγων μὴ μοιχεύειν μοιχεύεις; ὁ βδελυσσόμενος τὰ εἴδωλα ἱεροσυλεῖς;
\vs{23}ὃς ἐν νόμῳ καυχᾶσαι, διὰ τῆς παραβάσεως τοῦ νόμου τὸν Θεὸν ἀτιμάζεις·
\vs{24}τὸ γὰρ Ὄνομα τοῦ Θεοῦ δι᾽ ὑμᾶς βλασφημεῖται ἐν τοῖς ἔθνεσιν, καθὼς γέγραπται.

\vs{25}Περιτομὴ μὲν γὰρ ὠφελεῖ ἐὰν νόμον πράσσῃς· ἐὰν δὲ παραβάτης νόμου ᾖς, ἡ περιτομή σου ἀκροβυστία γέγονεν.
\vs{26}ἐὰν οὖν ἡ ἀκροβυστία τὰ δικαιώματα τοῦ νόμου φυλάσσῃ, οὐχ ἡ ἀκροβυστία αὐτοῦ εἰς περιτομὴν λογισθήσεται;
\vs{27}καὶ κρινεῖ ἡ ἐκ φύσεως ἀκροβυστία τὸν νόμον τελοῦσα σὲ τὸν διὰ γράμματος καὶ περιτομῆς παραβάτην νόμου.
\vs{28}Οὐ γὰρ ὁ ἐν τῷ φανερῷ Ἰουδαῖός ἐστιν οὐδὲ ἡ ἐν τῷ φανερῷ ἐν σαρκὶ περιτομή,
\vs{29}ἀλλ᾽ ὁ ἐν τῷ κρυπτῷ Ἰουδαῖος, καὶ περιτομὴ καρδίας ἐν πνεύματι οὐ γράμματι, οὗ ὁ ἔπαινος οὐκ ἐξ ἀνθρώπων ἀλλ᾽ ἐκ τοῦ Θεοῦ.

\ch{3}
Τί οὖν τὸ περισσὸν τοῦ Ἰουδαίου ἢ τίς ἡ ὠφέλεια τῆς περιτομῆς;
\vs{2}πολὺ κατὰ πάντα τρόπον. πρῶτον μὲν γὰρ ὅτι ἐπιστεύθησαν τὰ λόγια τοῦ Θεοῦ.
\vs{3}Τί γάρ; εἰ ἠπίστησάν τινες, μὴ ἡ ἀπιστία αὐτῶν τὴν πίστιν τοῦ Θεοῦ καταργήσει;
\vs{4}μὴ γένοιτο· γινέσθω δὲ ὁ Θεὸς ἀληθής, πᾶς δὲ ἄνθρωπος ψεύστης, καθὼς γέγραπται· 
\begin{poetryblock}

\begin{quote}Ὅπως ἂν δικαιωθῇς ἐν τοῖς λόγοις σου\end{quote} 

\begin{quote}καὶ νικήσεις ἐν τῷ κρίνεσθαί σε.\end{quote}
\end{poetryblock}

\vs{5}Εἰ δὲ ἡ ἀδικία ἡμῶν Θεοῦ δικαιοσύνην συνίστησιν, τί ἐροῦμεν; μὴ ἄδικος ὁ Θεὸς ὁ ἐπιφέρων τὴν ὀργήν; κατὰ ἄνθρωπον λέγω.
\vs{6}μὴ γένοιτο· ἐπεὶ πῶς κρινεῖ ὁ Θεὸς τὸν κόσμον;
\vs{7}εἰ δὲ ἡ ἀλήθεια τοῦ Θεοῦ ἐν τῷ ἐμῷ ψεύσματι ἐπερίσσευσεν εἰς τὴν δόξαν αὐτοῦ, τί ἔτι κἀγὼ ὡς ἁμαρτωλὸς κρίνομαι;
\vs{8}καὶ μὴ καθὼς βλασφημούμεθα καὶ καθώς φασίν τινες ἡμᾶς λέγειν ὅτι Ποιήσωμεν τὰ κακὰ, ἵνα ἔλθῃ τὰ ἀγαθά; ὧν τὸ κρίμα ἔνδικόν ἐστιν.

\vs{9}Τί οὖν; προεχόμεθα; οὐ πάντως· προῃτιασάμεθα γὰρ Ἰουδαίους τε καὶ Ἕλληνας πάντας ὑφ᾽ ἁμαρτίαν εἶναι,
\vs{10}καθὼς γέγραπται ὅτι 
\begin{poetryblock}

\begin{quote}Οὐκ ἔστιν δίκαιος οὐδὲ εἷς,\end{quote}

\begin{quote} \vs{11}οὐκ ἔστιν ὁ συνίων,\end{quote} 

\begin{quote}οὐκ ἔστιν ὁ ἐκζητῶν τὸν Θεόν.\end{quote}

\begin{quote} \vs{12}πάντες ἐξέκλιναν ἅμα ἠχρεώθησαν·\end{quote} 

\begin{quote}οὐκ ἔστιν ὁ ποιῶν χρηστότητα,\end{quote} 

\begin{quote}οὐκ ἔστιν ἕως ἑνός.\end{quote}

\begin{quote} \vs{13}τάφος ἀνεῳγμένος ὁ λάρυγξ αὐτῶν,\end{quote} 

\begin{quote}ταῖς γλώσσαις αὐτῶν ἐδολιοῦσαν,\end{quote} 

\begin{quote}ἰὸς ἀσπίδων ὑπὸ τὰ χείλη αὐτῶν·\end{quote}

\begin{quote} \vs{14}ὧν τὸ στόμα ἀρᾶς καὶ πικρίας γέμει,\end{quote}

\begin{quote} \vs{15}ὀξεῖς οἱ πόδες αὐτῶν ἐκχέαι αἷμα,\end{quote}

\begin{quote} \vs{16}σύντριμμα καὶ ταλαιπωρία ἐν ταῖς ὁδοῖς αὐτῶν,\end{quote}

\begin{quote} \vs{17}καὶ ὁδὸν εἰρήνης οὐκ ἔγνωσαν.\end{quote}

\begin{quote} \vs{18}οὐκ ἔστιν φόβος Θεοῦ ἀπέναντι τῶν ὀφθαλμῶν αὐτῶν.\end{quote}
\end{poetryblock}

\vs{19}Οἴδαμεν δὲ ὅτι ὅσα ὁ νόμος λέγει τοῖς ἐν τῷ νόμῳ λαλεῖ, ἵνα πᾶν στόμα φραγῇ καὶ ὑπόδικος γένηται πᾶς ὁ κόσμος τῷ Θεῷ·
\vs{20}διότι ἐξ ἔργων νόμου οὐ δικαιωθήσεται πᾶσα σὰρξ ἐνώπιον αὐτοῦ, διὰ γὰρ νόμου ἐπίγνωσις ἁμαρτίας.

\vs{21}Νυνὶ δὲ χωρὶς νόμου δικαιοσύνη Θεοῦ πεφανέρωται μαρτυρουμένη ὑπὸ τοῦ νόμου καὶ τῶν προφητῶν,
\vs{22}δικαιοσύνη δὲ Θεοῦ διὰ πίστεως Ἰησοῦ Χριστοῦ εἰς πάντας τοὺς πιστεύοντας. οὐ γάρ ἐστιν διαστολή,
\vs{23}πάντες γὰρ ἥμαρτον καὶ ὑστεροῦνται τῆς δόξης τοῦ Θεοῦ
\vs{24}δικαιούμενοι δωρεὰν τῇ αὐτοῦ χάριτι διὰ τῆς ἀπολυτρώσεως τῆς ἐν Χριστῷ Ἰησοῦ·
\vs{25}ὃν προέθετο ὁ Θεὸς ἱλαστήριον διὰ τῆς πίστεως ἐν τῷ αὐτοῦ αἵματι εἰς ἔνδειξιν τῆς δικαιοσύνης αὐτοῦ διὰ τὴν πάρεσιν τῶν προγεγονότων ἁμαρτημάτων
\vs{26}ἐν τῇ ἀνοχῇ τοῦ Θεοῦ, πρὸς τὴν ἔνδειξιν τῆς δικαιοσύνης αὐτοῦ ἐν τῷ νῦν καιρῷ, εἰς τὸ εἶναι αὐτὸν δίκαιον καὶ δικαιοῦντα τὸν ἐκ πίστεως Ἰησοῦ.

\vs{27}Ποῦ οὖν ἡ καύχησις; ἐξεκλείσθη. διὰ ποίου νόμου; τῶν ἔργων; οὐχί, ἀλλὰ διὰ νόμου πίστεως.
\vs{28}λογιζόμεθα γὰρ δικαιοῦσθαι πίστει ἄνθρωπον χωρὶς ἔργων νόμου.
\vs{29}Ἢ Ἰουδαίων ὁ Θεὸς μόνον; οὐχὶ καὶ ἐθνῶν; ναὶ καὶ ἐθνῶν,
\vs{30}εἴπερ εἷς ὁ Θεός ὃς δικαιώσει περιτομὴν ἐκ πίστεως καὶ ἀκροβυστίαν διὰ τῆς πίστεως.
\vs{31}Νόμον οὖν καταργοῦμεν διὰ τῆς πίστεως; μὴ γένοιτο· ἀλλὰ νόμον ἱστάνομεν.

\ch{4}
Τί οὖν ἐροῦμεν εὑρηκέναι Ἀβραὰμ τὸν προπάτορα ἡμῶν κατὰ σάρκα;
\vs{2}εἰ γὰρ Ἀβραὰμ ἐξ ἔργων ἐδικαιώθη, ἔχει καύχημα, ἀλλ᾽ οὐ πρὸς Θεόν.
\vs{3}τί γὰρ ἡ γραφὴ λέγει; Ἐπίστευσεν δὲ Ἀβραὰμ τῷ Θεῷ καὶ ἐλογίσθη αὐτῷ εἰς δικαιοσύνην.
\vs{4}Τῷ δὲ ἐργαζομένῳ ὁ μισθὸς οὐ λογίζεται κατὰ χάριν ἀλλὰ κατὰ ὀφείλημα,
\vs{5}τῷ δὲ μὴ ἐργαζομένῳ πιστεύοντι δὲ ἐπὶ τὸν δικαιοῦντα τὸν ἀσεβῆ λογίζεται ἡ πίστις αὐτοῦ εἰς δικαιοσύνην·
\vs{6}καθάπερ καὶ Δαυὶδ λέγει τὸν μακαρισμὸν τοῦ ἀνθρώπου ᾧ ὁ Θεὸς λογίζεται δικαιοσύνην χωρὶς ἔργων·
\begin{poetryblock}

\begin{quote} \vs{7}Μακάριοι ὧν ἀφέθησαν αἱ ἀνομίαι\end{quote} 

\begin{quote}καὶ ὧν ἐπεκαλύφθησαν αἱ ἁμαρτίαι·\end{quote}

\begin{quote} \vs{8}μακάριος ἀνὴρ οὗ οὐ μὴ λογίσηται Κύριος ἁμαρτίαν.\end{quote}
\end{poetryblock}

\vs{9}Ὁ μακαρισμὸς οὖν οὗτος ἐπὶ τὴν περιτομὴν ἢ καὶ ἐπὶ τὴν ἀκροβυστίαν; λέγομεν γάρ· Ἐλογίσθη τῷ Ἀβραὰμ ἡ πίστις εἰς δικαιοσύνην.
\vs{10}πῶς οὖν ἐλογίσθη; ἐν περιτομῇ ὄντι ἢ ἐν ἀκροβυστίᾳ; οὐκ ἐν περιτομῇ ἀλλ᾽ ἐν ἀκροβυστίᾳ·
\vs{11}Καὶ σημεῖον ἔλαβεν περιτομῆς σφραγῖδα τῆς δικαιοσύνης τῆς πίστεως τῆς ἐν τῇ ἀκροβυστίᾳ, εἰς τὸ εἶναι αὐτὸν πατέρα πάντων τῶν πιστευόντων δι᾽ ἀκροβυστίας, εἰς τὸ λογισθῆναι καὶ αὐτοῖς τὴν δικαιοσύνην,
\vs{12}καὶ πατέρα περιτομῆς τοῖς οὐκ ἐκ περιτομῆς μόνον ἀλλὰ καὶ τοῖς στοιχοῦσιν τοῖς ἴχνεσιν τῆς ἐν ἀκροβυστίᾳ πίστεως τοῦ πατρὸς ἡμῶν Ἀβραάμ.

\vs{13}Οὐ γὰρ διὰ νόμου ἡ ἐπαγγελία τῷ Ἀβραὰμ ἢ τῷ σπέρματι αὐτοῦ, τὸ κληρονόμον αὐτὸν εἶναι κόσμου, ἀλλὰ διὰ δικαιοσύνης πίστεως.
\vs{14}εἰ γὰρ οἱ ἐκ νόμου κληρονόμοι, κεκένωται ἡ πίστις καὶ κατήργηται ἡ ἐπαγγελία·
\vs{15}ὁ γὰρ νόμος ὀργὴν κατεργάζεται· οὗ δὲ οὐκ ἔστιν νόμος οὐδὲ παράβασις.
\vs{16}Διὰ τοῦτο ἐκ πίστεως, ἵνα κατὰ χάριν, εἰς τὸ εἶναι βεβαίαν τὴν ἐπαγγελίαν παντὶ τῷ σπέρματι, οὐ τῷ ἐκ τοῦ νόμου μόνον ἀλλὰ καὶ τῷ ἐκ πίστεως Ἀβραάμ, ὅς ἐστιν πατὴρ πάντων ἡμῶν,
\vs{17}καθὼς γέγραπται ὅτι Πατέρα πολλῶν ἐθνῶν τέθεικά σε, κατέναντι οὗ ἐπίστευσεν Θεοῦ τοῦ ζωοποιοῦντος τοὺς νεκροὺς καὶ καλοῦντος τὰ μὴ ὄντα ὡς ὄντα.
\vs{18}ὃς παρ᾽ ἐλπίδα ἐπ᾽ ἐλπίδι ἐπίστευσεν εἰς τὸ γενέσθαι αὐτὸν πατέρα πολλῶν ἐθνῶν κατὰ τὸ εἰρημένον· Οὕτως ἔσται τὸ σπέρμα σου,
\vs{19}καὶ μὴ ἀσθενήσας τῇ πίστει κατενόησεν τὸ ἑαυτοῦ σῶμα ἤδη νενεκρωμένον, ἑκατονταετής που ὑπάρχων, καὶ τὴν νέκρωσιν τῆς μήτρας Σάρρας·
\vs{20}εἰς δὲ τὴν ἐπαγγελίαν τοῦ Θεοῦ οὐ διεκρίθη τῇ ἀπιστίᾳ ἀλλὰ ἐνεδυναμώθη τῇ πίστει, δοὺς δόξαν τῷ Θεῷ
\vs{21}καὶ πληροφορηθεὶς ὅτι ὃ ἐπήγγελται δυνατός ἐστιν καὶ ποιῆσαι.
\vs{22}διὸ καὶ Ἐλογίσθη αὐτῷ εἰς δικαιοσύνην.
\vs{23}Οὐκ ἐγράφη δὲ δι᾽ αὐτὸν μόνον ὅτι Ἐλογίσθη αὐτῷ
\vs{24}ἀλλὰ καὶ δι᾽ ἡμᾶς, οἷς μέλλει λογίζεσθαι, τοῖς πιστεύουσιν ἐπὶ τὸν ἐγείραντα Ἰησοῦν τὸν Κύριον ἡμῶν ἐκ νεκρῶν,
\vs{25}ὃς παρεδόθη διὰ τὰ παραπτώματα ἡμῶν καὶ ἠγέρθη διὰ τὴν δικαίωσιν ἡμῶν.

\ch{5}
Δικαιωθέντες οὖν ἐκ πίστεως εἰρήνην ἔχομεν πρὸς τὸν Θεὸν διὰ τοῦ Κυρίου ἡμῶν Ἰησοῦ Χριστοῦ
\vs{2}δι᾽ οὗ καὶ τὴν προσαγωγὴν ἐσχήκαμεν τῇ πίστει εἰς τὴν χάριν ταύτην ἐν ᾗ ἑστήκαμεν καὶ καυχώμεθα ἐπ᾽ ἐλπίδι τῆς δόξης τοῦ Θεοῦ.
\vs{3}Οὐ μόνον δέ, ἀλλὰ καὶ καυχώμεθα ἐν ταῖς θλίψεσιν, εἰδότες ὅτι ἡ θλῖψις ὑπομονὴν κατεργάζεται,
\vs{4}ἡ δὲ ὑπομονὴ δοκιμήν, ἡ δὲ δοκιμὴ ἐλπίδα.
\vs{5}ἡ δὲ ἐλπὶς οὐ καταισχύνει, ὅτι ἡ ἀγάπη τοῦ Θεοῦ ἐκκέχυται ἐν ταῖς καρδίαις ἡμῶν διὰ Πνεύματος Ἁγίου τοῦ δοθέντος ἡμῖν.
\vs{6}Ἔτι γὰρ Χριστὸς ὄντων ἡμῶν ἀσθενῶν ἔτι κατὰ καιρὸν ὑπὲρ ἀσεβῶν ἀπέθανεν.
\vs{7}μόλις γὰρ ὑπὲρ δικαίου τις ἀποθανεῖται· ὑπὲρ γὰρ τοῦ ἀγαθοῦ τάχα τις καὶ τολμᾷ ἀποθανεῖν·
\vs{8}συνίστησιν δὲ τὴν ἑαυτοῦ ἀγάπην εἰς ἡμᾶς ὁ Θεὸς, ὅτι ἔτι ἁμαρτωλῶν ὄντων ἡμῶν Χριστὸς ὑπὲρ ἡμῶν ἀπέθανεν.
\vs{9}Πολλῷ οὖν μᾶλλον δικαιωθέντες νῦν ἐν τῷ αἵματι αὐτοῦ σωθησόμεθα δι᾽ αὐτοῦ ἀπὸ τῆς ὀργῆς.
\vs{10}εἰ γὰρ ἐχθροὶ ὄντες κατηλλάγημεν τῷ Θεῷ διὰ τοῦ θανάτου τοῦ Υἱοῦ αὐτοῦ, πολλῷ μᾶλλον καταλλαγέντες σωθησόμεθα ἐν τῇ ζωῇ αὐτοῦ·
\vs{11}οὐ μόνον δέ, ἀλλὰ καὶ καυχώμενοι ἐν τῷ Θεῷ διὰ τοῦ Κυρίου ἡμῶν Ἰησοῦ Χριστοῦ δι᾽ οὗ νῦν τὴν καταλλαγὴν ἐλάβομεν.

\vs{12}Διὰ τοῦτο ὥσπερ δι᾽ ἑνὸς ἀνθρώπου ἡ ἁμαρτία εἰς τὸν κόσμον εἰσῆλθεν καὶ διὰ τῆς ἁμαρτίας ὁ θάνατος, καὶ οὕτως εἰς πάντας ἀνθρώπους ὁ θάνατος διῆλθεν, ἐφ᾽ ᾧ πάντες ἥμαρτον·
\vs{13}ἄχρι γὰρ νόμου ἁμαρτία ἦν ἐν κόσμῳ, ἁμαρτία δὲ οὐκ ἐλλογεῖται μὴ ὄντος νόμου,
\vs{14}ἀλλὰ ἐβασίλευσεν ὁ θάνατος ἀπὸ Ἀδὰμ μέχρι Μωϋσέως καὶ ἐπὶ τοὺς μὴ ἁμαρτήσαντας ἐπὶ τῷ ὁμοιώματι τῆς παραβάσεως Ἀδάμ ὅς ἐστιν τύπος τοῦ μέλλοντος.
\vs{15}Ἀλλ᾽ οὐχ ὡς τὸ παράπτωμα, οὕτως καὶ τὸ χάρισμα· εἰ γὰρ τῷ τοῦ ἑνὸς παραπτώματι οἱ πολλοὶ ἀπέθανον, πολλῷ μᾶλλον ἡ χάρις τοῦ Θεοῦ καὶ ἡ δωρεὰ ἐν χάριτι τῇ τοῦ ἑνὸς ἀνθρώπου Ἰησοῦ Χριστοῦ εἰς τοὺς πολλοὺς ἐπερίσσευσεν.
\vs{16}καὶ οὐχ ὡς δι᾽ ἑνὸς ἁμαρτήσαντος τὸ δώρημα· τὸ μὲν γὰρ κρίμα ἐξ ἑνὸς εἰς κατάκριμα, τὸ δὲ χάρισμα ἐκ πολλῶν παραπτωμάτων εἰς δικαίωμα.
\vs{17}εἰ γὰρ τῷ τοῦ ἑνὸς παραπτώματι ὁ θάνατος ἐβασίλευσεν διὰ τοῦ ἑνός, πολλῷ μᾶλλον οἱ τὴν περισσείαν τῆς χάριτος καὶ τῆς δωρεᾶς τῆς δικαιοσύνης λαμβάνοντες ἐν ζωῇ βασιλεύσουσιν διὰ τοῦ ἑνὸς Ἰησοῦ Χριστοῦ.
\vs{18}Ἄρα οὖν ὡς δι᾽ ἑνὸς παραπτώματος εἰς πάντας ἀνθρώπους εἰς κατάκριμα, οὕτως καὶ δι᾽ ἑνὸς δικαιώματος εἰς πάντας ἀνθρώπους εἰς δικαίωσιν ζωῆς·
\vs{19}ὥσπερ γὰρ διὰ τῆς παρακοῆς τοῦ ἑνὸς ἀνθρώπου ἁμαρτωλοὶ κατεστάθησαν οἱ πολλοί, οὕτως καὶ διὰ τῆς ὑπακοῆς τοῦ ἑνὸς δίκαιοι κατασταθήσονται οἱ πολλοί.
\vs{20}Νόμος δὲ παρεισῆλθεν, ἵνα πλεονάσῃ τὸ παράπτωμα· οὗ δὲ ἐπλεόνασεν ἡ ἁμαρτία, ὑπερεπερίσσευσεν ἡ χάρις,
\vs{21}ἵνα ὥσπερ ἐβασίλευσεν ἡ ἁμαρτία ἐν τῷ θανάτῳ, οὕτως καὶ ἡ χάρις βασιλεύσῃ διὰ δικαιοσύνης εἰς ζωὴν αἰώνιον διὰ Ἰησοῦ Χριστοῦ τοῦ Κυρίου ἡμῶν.

\ch{6}
Τί οὖν ἐροῦμεν; ἐπιμένωμεν τῇ ἁμαρτίᾳ, ἵνα ἡ χάρις πλεονάσῃ;
\vs{2}μὴ γένοιτο. οἵτινες ἀπεθάνομεν τῇ ἁμαρτίᾳ, πῶς ἔτι ζήσομεν ἐν αὐτῇ;
\vs{3}ἢ ἀγνοεῖτε ὅτι, ὅσοι ἐβαπτίσθημεν εἰς Χριστὸν Ἰησοῦν, εἰς τὸν θάνατον αὐτοῦ ἐβαπτίσθημεν;
\vs{4}συνετάφημεν οὖν αὐτῷ διὰ τοῦ βαπτίσματος εἰς τὸν θάνατον, ἵνα ὥσπερ ἠγέρθη Χριστὸς ἐκ νεκρῶν διὰ τῆς δόξης τοῦ Πατρός, οὕτως καὶ ἡμεῖς ἐν καινότητι ζωῆς περιπατήσωμεν.
\vs{5}Εἰ γὰρ σύμφυτοι γεγόναμεν τῷ ὁμοιώματι τοῦ θανάτου αὐτοῦ, ἀλλὰ καὶ τῆς ἀναστάσεως ἐσόμεθα·
\vs{6}τοῦτο γινώσκοντες ὅτι ὁ παλαιὸς ἡμῶν ἄνθρωπος συνεσταυρώθη, ἵνα καταργηθῇ τὸ σῶμα τῆς ἁμαρτίας, τοῦ μηκέτι δουλεύειν ἡμᾶς τῇ ἁμαρτίᾳ·
\vs{7}ὁ γὰρ ἀποθανὼν δεδικαίωται ἀπὸ τῆς ἁμαρτίας.
\vs{8}Εἰ δὲ ἀπεθάνομεν σὺν Χριστῷ, πιστεύομεν ὅτι καὶ συζήσομεν αὐτῷ,
\vs{9}εἰδότες ὅτι Χριστὸς ἐγερθεὶς ἐκ νεκρῶν οὐκέτι ἀποθνῄσκει, θάνατος αὐτοῦ οὐκέτι κυριεύει.
\vs{10}ὃ γὰρ ἀπέθανεν, τῇ ἁμαρτίᾳ ἀπέθανεν ἐφάπαξ· ὃ δὲ ζῇ, ζῇ τῷ Θεῷ.
\vs{11}οὕτως καὶ ὑμεῖς λογίζεσθε ἑαυτοὺς εἶναι νεκροὺς μὲν τῇ ἁμαρτίᾳ ζῶντας δὲ τῷ Θεῷ ἐν Χριστῷ Ἰησοῦ.
\vs{12}Μὴ οὖν βασιλευέτω ἡ ἁμαρτία ἐν τῷ θνητῷ ὑμῶν σώματι εἰς τὸ ὑπακούειν ταῖς ἐπιθυμίαις αὐτοῦ,
\vs{13}μηδὲ παριστάνετε τὰ μέλη ὑμῶν ὅπλα ἀδικίας τῇ ἁμαρτίᾳ, ἀλλὰ παραστήσατε ἑαυτοὺς τῷ Θεῷ ὡσεὶ ἐκ νεκρῶν ζῶντας καὶ τὰ μέλη ὑμῶν ὅπλα δικαιοσύνης τῷ Θεῷ.
\vs{14}ἁμαρτία γὰρ ὑμῶν οὐ κυριεύσει· οὐ γάρ ἐστε ὑπὸ νόμον ἀλλὰ ὑπὸ χάριν.
\vs{15}Τί οὖν; ἁμαρτήσωμεν, ὅτι οὐκ ἐσμὲν ὑπὸ νόμον ἀλλὰ ὑπὸ χάριν; μὴ γένοιτο.
\vs{16}οὐκ οἴδατε ὅτι ᾧ παριστάνετε ἑαυτοὺς δούλους εἰς ὑπακοήν, δοῦλοί ἐστε ᾧ ὑπακούετε, ἤτοι ἁμαρτίας εἰς θάνατον ἢ ὑπακοῆς εἰς δικαιοσύνην;
\vs{17}χάρις δὲ τῷ Θεῷ ὅτι ἦτε δοῦλοι τῆς ἁμαρτίας ὑπηκούσατε δὲ ἐκ καρδίας εἰς ὃν παρεδόθητε τύπον διδαχῆς,
\vs{18}ἐλευθερωθέντες δὲ ἀπὸ τῆς ἁμαρτίας ἐδουλώθητε τῇ δικαιοσύνῃ.
\vs{19}Ἀνθρώπινον λέγω διὰ τὴν ἀσθένειαν τῆς σαρκὸς ὑμῶν. ὥσπερ γὰρ παρεστήσατε τὰ μέλη ὑμῶν δοῦλα τῇ ἀκαθαρσίᾳ καὶ τῇ ἀνομίᾳ εἰς τὴν ἀνομίαν, οὕτως νῦν παραστήσατε τὰ μέλη ὑμῶν δοῦλα τῇ δικαιοσύνῃ εἰς ἁγιασμόν.
\vs{20}Ὅτε γὰρ δοῦλοι ἦτε τῆς ἁμαρτίας, ἐλεύθεροι ἦτε τῇ δικαιοσύνῃ.
\vs{21}τίνα οὖν καρπὸν εἴχετε τότε; ἐφ᾽ οἷς νῦν ἐπαισχύνεσθε, τὸ γὰρ τέλος ἐκείνων θάνατος.
\vs{22}νυνὶ δέ ἐλευθερωθέντες ἀπὸ τῆς ἁμαρτίας δουλωθέντες δὲ τῷ Θεῷ ἔχετε τὸν καρπὸν ὑμῶν εἰς ἁγιασμόν, τὸ δὲ τέλος ζωὴν αἰώνιον.
\vs{23}τὰ γὰρ ὀψώνια τῆς ἁμαρτίας θάνατος, τὸ δὲ χάρισμα τοῦ Θεοῦ ζωὴ αἰώνιος ἐν Χριστῷ Ἰησοῦ τῷ Κυρίῳ ἡμῶν.

\ch{7}
Ἢ ἀγνοεῖτε, ἀδελφοί, γινώσκουσιν γὰρ νόμον λαλῶ, ὅτι ὁ νόμος κυριεύει τοῦ ἀνθρώπου ἐφ᾽ ὅσον χρόνον ζῇ;
\vs{2}ἡ γὰρ ὕπανδρος γυνὴ τῷ ζῶντι ἀνδρὶ δέδεται νόμῳ· ἐὰν δὲ ἀποθάνῃ ὁ ἀνήρ, κατήργηται ἀπὸ τοῦ νόμου τοῦ ἀνδρός.
\vs{3}ἄρα οὖν ζῶντος τοῦ ἀνδρὸς μοιχαλὶς χρηματίσει ἐὰν γένηται ἀνδρὶ ἑτέρῳ· ἐὰν δὲ ἀποθάνῃ ὁ ἀνήρ, ἐλευθέρα ἐστὶν ἀπὸ τοῦ νόμου, τοῦ μὴ εἶναι αὐτὴν μοιχαλίδα γενομένην ἀνδρὶ ἑτέρῳ.
\vs{4}Ὥστε, ἀδελφοί μου, καὶ ὑμεῖς ἐθανατώθητε τῷ νόμῳ διὰ τοῦ σώματος τοῦ Χριστοῦ, εἰς τὸ γενέσθαι ὑμᾶς ἑτέρῳ, τῷ ἐκ νεκρῶν ἐγερθέντι, ἵνα καρποφορήσωμεν τῷ Θεῷ.
\vs{5}ὅτε γὰρ ἦμεν ἐν τῇ σαρκί, τὰ παθήματα τῶν ἁμαρτιῶν τὰ διὰ τοῦ νόμου ἐνηργεῖτο ἐν τοῖς μέλεσιν ἡμῶν, εἰς τὸ καρποφορῆσαι τῷ θανάτῳ·
\vs{6}νυνὶ δὲ κατηργήθημεν ἀπὸ τοῦ νόμου ἀποθανόντες ἐν ᾧ κατειχόμεθα, ὥστε δουλεύειν ἡμᾶς ἐν καινότητι πνεύματος καὶ οὐ παλαιότητι γράμματος.

\vs{7}Τί οὖν ἐροῦμεν; ὁ νόμος ἁμαρτία; μὴ γένοιτο· ἀλλὰ τὴν ἁμαρτίαν οὐκ ἔγνων εἰ μὴ διὰ νόμου· τήν τε γὰρ ἐπιθυμίαν οὐκ ᾔδειν εἰ μὴ ὁ νόμος ἔλεγεν· Οὐκ ἐπιθυμήσεις.
\vs{8}ἀφορμὴν δὲ λαβοῦσα ἡ ἁμαρτία διὰ τῆς ἐντολῆς κατειργάσατο ἐν ἐμοὶ πᾶσαν ἐπιθυμίαν· χωρὶς γὰρ νόμου ἁμαρτία νεκρά.
\vs{9}Ἐγὼ δὲ ἔζων χωρὶς νόμου ποτέ, ἐλθούσης δὲ τῆς ἐντολῆς ἡ ἁμαρτία ἀνέζησεν,
\vs{10}ἐγὼ δὲ ἀπέθανον καὶ εὑρέθη μοι ἡ ἐντολὴ ἡ εἰς ζωὴν, αὕτη εἰς θάνατον·
\vs{11}ἡ γὰρ ἁμαρτία ἀφορμὴν λαβοῦσα διὰ τῆς ἐντολῆς ἐξηπάτησέν με καὶ δι᾽ αὐτῆς ἀπέκτεινεν.
\vs{12}Ὥστε ὁ μὲν νόμος ἅγιος καὶ ἡ ἐντολὴ ἁγία καὶ δικαία καὶ ἀγαθή.
\vs{13}Τὸ οὖν ἀγαθὸν ἐμοὶ ἐγένετο θάνατος; μὴ γένοιτο· ἀλλὰ ἡ ἁμαρτία, ἵνα φανῇ ἁμαρτία, διὰ τοῦ ἀγαθοῦ μοι κατεργαζομένη θάνατον, ἵνα γένηται καθ᾽ ὑπερβολὴν ἁμαρτωλὸς ἡ ἁμαρτία διὰ τῆς ἐντολῆς.

\vs{14}Οἴδαμεν γὰρ ὅτι ὁ νόμος πνευματικός ἐστιν, ἐγὼ δὲ σάρκινός εἰμι πεπραμένος ὑπὸ τὴν ἁμαρτίαν.
\vs{15}ὃ γὰρ κατεργάζομαι οὐ γινώσκω· οὐ γὰρ ὃ θέλω τοῦτο πράσσω, ἀλλ᾽ ὃ μισῶ τοῦτο ποιῶ.
\vs{16}εἰ δὲ ὃ οὐ θέλω τοῦτο ποιῶ, σύμφημι τῷ νόμῳ ὅτι καλός.
\vs{17}νυνὶ δὲ οὐκέτι ἐγὼ κατεργάζομαι αὐτὸ ἀλλὰ ἡ οἰκοῦσα ἐν ἐμοὶ ἁμαρτία.
\vs{18}Οἶδα γὰρ ὅτι οὐκ οἰκεῖ ἐν ἐμοί, τοῦτ᾽ ἔστιν ἐν τῇ σαρκί μου, ἀγαθόν· τὸ γὰρ θέλειν παράκειταί μοι, τὸ δὲ κατεργάζεσθαι τὸ καλὸν οὔ·
\vs{19}οὐ γὰρ ὃ θέλω ποιῶ ἀγαθόν, ἀλλὰ ὃ οὐ θέλω κακὸν τοῦτο πράσσω.
\vs{20}εἰ δὲ ὃ οὐ θέλω ἐγὼ τοῦτο ποιῶ, οὐκέτι ἐγὼ κατεργάζομαι αὐτὸ ἀλλὰ ἡ οἰκοῦσα ἐν ἐμοὶ ἁμαρτία.
\vs{21}Εὑρίσκω ἄρα τὸν νόμον, τῷ θέλοντι ἐμοὶ ποιεῖν τὸ καλὸν, ὅτι ἐμοὶ τὸ κακὸν παράκειται·
\vs{22}συνήδομαι γὰρ τῷ νόμῳ τοῦ Θεοῦ κατὰ τὸν ἔσω ἄνθρωπον,
\vs{23}βλέπω δὲ ἕτερον νόμον ἐν τοῖς μέλεσίν μου ἀντιστρατευόμενον τῷ νόμῳ τοῦ νοός μου καὶ αἰχμαλωτίζοντά με ἐν τῷ νόμῳ τῆς ἁμαρτίας τῷ ὄντι ἐν τοῖς μέλεσίν μου.
\vs{24}Ταλαίπωρος ἐγὼ ἄνθρωπος· τίς με ῥύσεται ἐκ τοῦ σώματος τοῦ θανάτου τούτου;
\vs{25}χάρις δὲ τῷ Θεῷ διὰ Ἰησοῦ Χριστοῦ τοῦ Κυρίου ἡμῶν. Ἄρα οὖν αὐτὸς ἐγὼ τῷ μὲν νοῒ δουλεύω νόμῳ Θεοῦ τῇ δὲ σαρκὶ νόμῳ ἁμαρτίας.

\ch{8}
Οὐδὲν ἄρα νῦν κατάκριμα τοῖς ἐν Χριστῷ Ἰησοῦ.
\vs{2}ὁ γὰρ νόμος τοῦ Πνεύματος τῆς ζωῆς ἐν Χριστῷ Ἰησοῦ ἠλευθέρωσέν σε ἀπὸ τοῦ νόμου τῆς ἁμαρτίας καὶ τοῦ θανάτου.
\vs{3}τὸ γὰρ ἀδύνατον τοῦ νόμου ἐν ᾧ ἠσθένει διὰ τῆς σαρκός, ὁ Θεὸς τὸν ἑαυτοῦ Υἱὸν πέμψας ἐν ὁμοιώματι σαρκὸς ἁμαρτίας καὶ περὶ ἁμαρτίας κατέκρινεν τὴν ἁμαρτίαν ἐν τῇ σαρκί,
\vs{4}ἵνα τὸ δικαίωμα τοῦ νόμου πληρωθῇ ἐν ἡμῖν τοῖς μὴ κατὰ σάρκα περιπατοῦσιν ἀλλὰ κατὰ πνεῦμα.
\vs{5}Οἱ γὰρ κατὰ σάρκα ὄντες τὰ τῆς σαρκὸς φρονοῦσιν, οἱ δὲ κατὰ πνεῦμα τὰ τοῦ πνεύματος.
\vs{6}τὸ γὰρ φρόνημα τῆς σαρκὸς θάνατος, τὸ δὲ φρόνημα τοῦ πνεύματος ζωὴ καὶ εἰρήνη·
\vs{7}διότι τὸ φρόνημα τῆς σαρκὸς ἔχθρα εἰς Θεόν, τῷ γὰρ νόμῳ τοῦ Θεοῦ οὐχ ὑποτάσσεται, οὐδὲ γὰρ δύναται·
\vs{8}οἱ δὲ ἐν σαρκὶ ὄντες Θεῷ ἀρέσαι οὐ δύνανται.
\vs{9}Ὑμεῖς δὲ οὐκ ἐστὲ ἐν σαρκὶ ἀλλὰ ἐν πνεύματι, εἴπερ Πνεῦμα Θεοῦ οἰκεῖ ἐν ὑμῖν. εἰ δέ τις Πνεῦμα Χριστοῦ οὐκ ἔχει, οὗτος οὐκ ἔστιν αὐτοῦ.
\vs{10}εἰ δὲ Χριστὸς ἐν ὑμῖν, τὸ μὲν σῶμα νεκρὸν διὰ ἁμαρτίαν τὸ δὲ πνεῦμα ζωὴ διὰ δικαιοσύνην.
\vs{11}εἰ δὲ τὸ Πνεῦμα τοῦ ἐγείραντος τὸν Ἰησοῦν ἐκ νεκρῶν οἰκεῖ ἐν ὑμῖν, ὁ ἐγείρας Χριστὸν ἐκ νεκρῶν ζωοποιήσει καὶ τὰ θνητὰ σώματα ὑμῶν διὰ τοῦ ἐνοικοῦντος αὐτοῦ Πνεύματος ἐν ὑμῖν.

\vs{12}Ἄρα οὖν, ἀδελφοί, ὀφειλέται ἐσμέν οὐ τῇ σαρκὶ τοῦ κατὰ σάρκα ζῆν,
\vs{13}εἰ γὰρ κατὰ σάρκα ζῆτε, μέλλετε ἀποθνήσκειν· εἰ δὲ πνεύματι τὰς πράξεις τοῦ σώματος θανατοῦτε, ζήσεσθε.
\vs{14}ὅσοι γὰρ Πνεύματι Θεοῦ ἄγονται, οὗτοι υἱοί Θεοῦ εἰσιν.
\vs{15}Οὐ γὰρ ἐλάβετε πνεῦμα δουλείας πάλιν εἰς φόβον ἀλλὰ ἐλάβετε πνεῦμα υἱοθεσίας ἐν ᾧ κράζομεν· Ἀββᾶ ὁ Πατήρ.
\vs{16}αὐτὸ τὸ Πνεῦμα συμμαρτυρεῖ τῷ πνεύματι ἡμῶν ὅτι ἐσμὲν τέκνα Θεοῦ.
\vs{17}εἰ δὲ τέκνα, καὶ κληρονόμοι· κληρονόμοι μὲν Θεοῦ, συνκληρονόμοι δὲ Χριστοῦ, εἴπερ συμπάσχομεν ἵνα καὶ συνδοξασθῶμεν.

\vs{18}Λογίζομαι γὰρ ὅτι οὐκ ἄξια τὰ παθήματα τοῦ νῦν καιροῦ πρὸς τὴν μέλλουσαν δόξαν ἀποκαλυφθῆναι εἰς ἡμᾶς.
\vs{19}ἡ γὰρ ἀποκαραδοκία τῆς κτίσεως τὴν ἀποκάλυψιν τῶν υἱῶν τοῦ Θεοῦ ἀπεκδέχεται.
\vs{20}τῇ γὰρ ματαιότητι ἡ κτίσις ὑπετάγη, οὐχ ἑκοῦσα ἀλλὰ διὰ τὸν ὑποτάξαντα, ἐφ᾽ ἑλπίδι
\vs{21}ὅτι καὶ αὐτὴ ἡ κτίσις ἐλευθερωθήσεται ἀπὸ τῆς δουλείας τῆς φθορᾶς εἰς τὴν ἐλευθερίαν τῆς δόξης τῶν τέκνων τοῦ Θεοῦ.
\vs{22}Οἴδαμεν γὰρ ὅτι πᾶσα ἡ κτίσις συστενάζει καὶ συνωδίνει ἄχρι τοῦ νῦν·
\vs{23}οὐ μόνον δέ, ἀλλὰ καὶ αὐτοὶ τὴν ἀπαρχὴν τοῦ Πνεύματος ἔχοντες, ἡμεῖς καὶ αὐτοὶ ἐν ἑαυτοῖς στενάζομεν υἱοθεσίαν ἀπεκδεχόμενοι, τὴν ἀπολύτρωσιν τοῦ σώματος ἡμῶν.
\vs{24}τῇ γὰρ ἐλπίδι ἐσώθημεν· ἐλπὶς δὲ βλεπομένη οὐκ ἔστιν ἐλπίς· ὃ γὰρ βλέπει τις ἐλπίζει;
\vs{25}εἰ δὲ ὃ οὐ βλέπομεν ἐλπίζομεν, δι᾽ ὑπομονῆς ἀπεκδεχόμεθα.
\vs{26}Ὡσαύτως δὲ καὶ τὸ Πνεῦμα συναντιλαμβάνεται τῇ ἀσθενείᾳ ἡμῶν· τὸ γὰρ τί προσευξώμεθα καθὸ δεῖ οὐκ οἴδαμεν, ἀλλὰ αὐτὸ τὸ Πνεῦμα ὑπερεντυγχάνει στεναγμοῖς ἀλαλήτοις·
\vs{27}ὁ δὲ ἐραυνῶν τὰς καρδίας οἶδεν τί τὸ φρόνημα τοῦ Πνεύματος, ὅτι κατὰ Θεὸν ἐντυγχάνει ὑπὲρ ἁγίων.
\vs{28}Οἴδαμεν δὲ ὅτι τοῖς ἀγαπῶσιν τὸν Θεὸν πάντα συνεργεῖ εἰς ἀγαθόν, τοῖς κατὰ πρόθεσιν κλητοῖς οὖσιν.
\vs{29}ὅτι οὓς προέγνω, καὶ προώρισεν συμμόρφους τῆς εἰκόνος τοῦ Υἱοῦ αὐτοῦ, εἰς τὸ εἶναι αὐτὸν πρωτότοκον ἐν πολλοῖς ἀδελφοῖς·
\vs{30}οὓς δὲ προώρισεν, τούτους καὶ ἐκάλεσεν· καὶ οὓς ἐκάλεσεν, τούτους καὶ ἐδικαίωσεν· οὓς δὲ ἐδικαίωσεν, τούτους καὶ ἐδόξασεν.

\vs{31}Τί οὖν ἐροῦμεν πρὸς ταῦτα; εἰ ὁ Θεὸς ὑπὲρ ἡμῶν, τίς καθ᾽ ἡμῶν;
\vs{32}ὅς γε τοῦ ἰδίου Υἱοῦ οὐκ ἐφείσατο ἀλλὰ ὑπὲρ ἡμῶν πάντων παρέδωκεν αὐτόν, πῶς οὐχὶ καὶ σὺν αὐτῷ τὰ πάντα ἡμῖν χαρίσεται;
\vs{33}τίς ἐγκαλέσει κατὰ ἐκλεκτῶν Θεοῦ; Θεὸς ὁ δικαιῶν·
\vs{34}τίς ὁ κατακρινῶν; Χριστὸς Ἰησοῦς ὁ ἀποθανών, μᾶλλον δὲ ἐγερθείς, ὅς καί ἐστιν ἐν δεξιᾷ τοῦ Θεοῦ, ὃς καὶ ἐντυγχάνει ὑπὲρ ἡμῶν.
\vs{35}Τίς ἡμᾶς χωρίσει ἀπὸ τῆς ἀγάπης τοῦ Χριστοῦ; θλῖψις ἢ στενοχωρία ἢ διωγμὸς ἢ λιμὸς ἢ γυμνότης ἢ κίνδυνος ἢ μάχαιρα;
\vs{36}καθὼς γέγραπται ὅτι 
\begin{poetryblock}

\begin{quote}Ἕνεκεν σοῦ θανατούμεθα ὅλην τὴν ἡμέραν,\end{quote} 

\begin{quote}ἐλογίσθημεν ὡς πρόβατα σφαγῆς.\end{quote}
\end{poetryblock}

\vs{37}Ἀλλ᾽ ἐν τούτοις πᾶσιν ὑπερνικῶμεν διὰ τοῦ ἀγαπήσαντος ἡμᾶς.
\vs{38}πέπεισμαι γὰρ ὅτι οὔτε θάνατος οὔτε ζωὴ οὔτε ἄγγελοι οὔτε ἀρχαὶ οὔτε ἐνεστῶτα οὔτε μέλλοντα οὔτε δυνάμεις
\vs{39}οὔτε ὕψωμα οὔτε βάθος οὔτε τις κτίσις ἑτέρα δυνήσεται ἡμᾶς χωρίσαι ἀπὸ τῆς ἀγάπης τοῦ Θεοῦ τῆς ἐν Χριστῷ Ἰησοῦ τῷ Κυρίῳ ἡμῶν.

\ch{9}
Ἀλήθειαν λέγω ἐν Χριστῷ, οὐ ψεύδομαι, συμμαρτυρούσης μοι τῆς συνειδήσεώς μου ἐν Πνεύματι Ἁγίῳ,
\vs{2}ὅτι λύπη μοί ἐστιν μεγάλη καὶ ἀδιάλειπτος ὀδύνη τῇ καρδίᾳ μου.
\vs{3}ηὐχόμην γὰρ ἀνάθεμα εἶναι αὐτὸς ἐγὼ ἀπὸ τοῦ Χριστοῦ ὑπὲρ τῶν ἀδελφῶν μου τῶν συγγενῶν μου κατὰ σάρκα,
\vs{4}οἵτινές εἰσιν Ἰσραηλῖται, ὧν ἡ υἱοθεσία καὶ ἡ δόξα καὶ αἱ διαθῆκαι καὶ ἡ νομοθεσία καὶ ἡ λατρεία καὶ αἱ ἐπαγγελίαι,
\vs{5}ὧν οἱ πατέρες καὶ ἐξ ὧν ὁ Χριστὸς τὸ κατὰ σάρκα, ὁ ὢν ἐπὶ πάντων Θεὸς εὐλογητὸς εἰς τοὺς αἰῶνας, ἀμήν.

\vs{6}Οὐχ οἷον δὲ ὅτι ἐκπέπτωκεν ὁ λόγος τοῦ Θεοῦ. οὐ γὰρ πάντες οἱ ἐξ Ἰσραήλ οὗτοι Ἰσραήλ·
\vs{7}οὐδ᾽ ὅτι εἰσὶν σπέρμα Ἀβραάμ πάντες τέκνα, ἀλλ᾽· Ἐν Ἰσαὰκ κληθήσεταί σοι σπέρμα.
\vs{8}τοῦτ᾽ ἔστιν, οὐ τὰ τέκνα τῆς σαρκὸς ταῦτα τέκνα τοῦ Θεοῦ ἀλλὰ τὰ τέκνα τῆς ἐπαγγελίας λογίζεται εἰς σπέρμα.
\vs{9}ἐπαγγελίας γὰρ ὁ λόγος οὗτος· Κατὰ τὸν καιρὸν τοῦτον ἐλεύσομαι καὶ ἔσται τῇ Σάρρᾳ υἱός.

\vs{10}Οὐ μόνον δέ, ἀλλὰ καὶ Ῥεβέκκα ἐξ ἑνὸς κοίτην ἔχουσα, Ἰσαὰκ τοῦ πατρὸς ἡμῶν·
\vs{11}μήπω γὰρ γεννηθέντων μηδὲ πραξάντων τι ἀγαθὸν ἢ φαῦλον, ἵνα ἡ κατ᾽ ἐκλογὴν πρόθεσις τοῦ Θεοῦ μένῃ,
\vs{12}οὐκ ἐξ ἔργων ἀλλ᾽ ἐκ τοῦ καλοῦντος, ἐρρέθη αὐτῇ ὅτι Ὁ μείζων δουλεύσει τῷ ἐλάσσονι,
\vs{13}καθὼς γέγραπται· Τὸν Ἰακὼβ ἠγάπησα, τὸν δὲ Ἠσαῦ ἐμίσησα.

\vs{14}Τί οὖν ἐροῦμεν; μὴ ἀδικία παρὰ τῷ θεῷ; μὴ γένοιτο.
\vs{15}τῷ Μωϋσεῖ γὰρ λέγει· Ἐλεήσω ὃν ἂν ἐλεῶ καὶ οἰκτιρήσω ὃν ἂν οἰκτίρω.
\vs{16}Ἄρα οὖν οὐ τοῦ θέλοντος οὐδὲ τοῦ τρέχοντος ἀλλὰ τοῦ ἐλεῶντος Θεοῦ.
\vs{17}λέγει γὰρ ἡ γραφὴ τῷ Φαραὼ ὅτι Εἰς αὐτὸ τοῦτο ἐξήγειρά σε ὅπως ἐνδείξωμαι ἐν σοὶ τὴν δύναμίν μου καὶ ὅπως διαγγελῇ τὸ ὄνομά μου ἐν πάσῃ τῇ γῇ.
\vs{18}ἄρα οὖν ὃν θέλει ἐλεεῖ, ὃν δὲ θέλει σκληρύνει.
\vs{19}Ἐρεῖς μοι οὖν· Τί οὖν ἔτι μέμφεται; τῷ γὰρ βουλήματι αὐτοῦ τίς ἀνθέστηκεν;
\vs{20}ὦ ἄνθρωπε, μενοῦνγε σὺ τίς εἶ ὁ ἀνταποκρινόμενος τῷ Θεῷ; μὴ ἐρεῖ τὸ πλάσμα τῷ πλάσαντι· Τί με ἐποίησας οὕτως;
\vs{21}ἢ οὐκ ἔχει ἐξουσίαν ὁ κεραμεὺς τοῦ πηλοῦ ἐκ τοῦ αὐτοῦ φυράματος ποιῆσαι ὃ μὲν εἰς τιμὴν σκεῦος ὃ δὲ εἰς ἀτιμίαν;
\vs{22}Εἰ δὲ θέλων ὁ Θεὸς ἐνδείξασθαι τὴν ὀργὴν καὶ γνωρίσαι τὸ δυνατὸν αὐτοῦ ἤνεγκεν ἐν πολλῇ μακροθυμίᾳ σκεύη ὀργῆς κατηρτισμένα εἰς ἀπώλειαν,
\vs{23}καὶ ἵνα γνωρίσῃ τὸν πλοῦτον τῆς δόξης αὐτοῦ ἐπὶ σκεύη ἐλέους ἃ προητοίμασεν εἰς δόξαν;
\vs{24}οὓς καὶ ἐκάλεσεν ἡμᾶς οὐ μόνον ἐξ Ἰουδαίων ἀλλὰ καὶ ἐξ ἐθνῶν,
\vs{25}ὡς καὶ ἐν τῷ Ὡσηὲ λέγει· 
\begin{poetryblock}

\begin{quote}Καλέσω τὸν οὐ λαόν μου λαόν μου\end{quote} 

\begin{quote}καὶ τὴν οὐκ ἠγαπημένην ἠγαπημένην·\end{quote}

\begin{quote} \vs{26}Καὶ Ἔσται ἐν τῷ τόπῳ οὗ ἐρρέθη αὐτοῖς·\end{quote} 

\begin{quote}Οὐ λαός μου ὑμεῖς,\end{quote} 

\begin{quote}ἐκεῖ κληθήσονται Υἱοὶ Θεοῦ ζῶντος.\end{quote}
\end{poetryblock}

\vs{27}Ἠσαΐας δὲ κράζει ὑπὲρ τοῦ Ἰσραήλ· 
\begin{poetryblock}

\begin{quote}Ἐὰν ᾖ ὁ ἀριθμὸς τῶν υἱῶν Ἰσραὴλ ὡς ἡ ἄμμος τῆς θαλάσσης,\end{quote} 

\begin{quote}τὸ ὑπόλειμμα σωθήσεται·\end{quote}

\begin{quote} \vs{28}λόγον γὰρ συντελῶν καὶ συντέμνων\end{quote} 

\begin{quote}ποιήσει Κύριος ἐπὶ τῆς γῆς.\end{quote}
\end{poetryblock}

\vs{29}Καὶ καθὼς προείρηκεν Ἠσαΐας· 
\begin{poetryblock}

\begin{quote}Εἰ μὴ Κύριος Σαβαὼθ ἐγκατέλιπεν ἡμῖν σπέρμα,\end{quote} 

\begin{quote}ὡς Σόδομα ἂν ἐγενήθημεν καὶ ὡς Γόμορρα ἂν ὡμοιώθημεν.\end{quote}
\end{poetryblock}

\vs{30}Τί οὖν ἐροῦμεν; ὅτι ἔθνη τὰ μὴ διώκοντα δικαιοσύνην κατέλαβεν δικαιοσύνην, δικαιοσύνην δὲ τὴν ἐκ πίστεως,
\vs{31}Ἰσραὴλ δὲ διώκων νόμον δικαιοσύνης εἰς νόμον οὐκ ἔφθασεν.
\vs{32}διὰ τί; ὅτι οὐκ ἐκ πίστεως ἀλλ᾽ ὡς ἐξ ἔργων· προσέκοψαν τῷ λίθῳ τοῦ προσκόμματος,
\vs{33}καθὼς γέγραπται· 
\begin{poetryblock}

\begin{quote}Ἰδοὺ τίθημι ἐν Σιὼν λίθον προσκόμματος καὶ πέτραν σκανδάλου,\end{quote} 

\begin{quote}καὶ ὁ πιστεύων ἐπ᾽ αὐτῷ οὐ καταισχυνθήσεται.\end{quote}
\end{poetryblock}

\ch{10}
Ἀδελφοί, ἡ μὲν εὐδοκία τῆς ἐμῆς καρδίας καὶ ἡ δέησις πρὸς τὸν Θεὸν ὑπὲρ αὐτῶν εἰς σωτηρίαν.
\vs{2}μαρτυρῶ γὰρ αὐτοῖς ὅτι ζῆλον Θεοῦ ἔχουσιν ἀλλ᾽ οὐ κατ᾽ ἐπίγνωσιν·
\vs{3}ἀγνοοῦντες γὰρ τὴν τοῦ Θεοῦ δικαιοσύνην καὶ τὴν ἰδίαν δικαιοσύνην ζητοῦντες στῆσαι, τῇ δικαιοσύνῃ τοῦ Θεοῦ οὐχ ὑπετάγησαν.
\vs{4}τέλος γὰρ νόμου Χριστὸς εἰς δικαιοσύνην παντὶ τῷ πιστεύοντι.
\vs{5}Μωϋσῆς γὰρ γράφει τὴν δικαιοσύνην τὴν ἐκ τοῦ νόμου ὅτι Ὁ ποιήσας αὐτὰ ἄνθρωπος ζήσεται ἐν αὐτῇ.
\vs{6}ἡ δὲ ἐκ πίστεως δικαιοσύνη οὕτως λέγει· Μὴ εἴπῃς ἐν τῇ καρδίᾳ σου· Τίς ἀναβήσεται εἰς τὸν οὐρανόν; τοῦτ᾽ ἔστιν Χριστὸν καταγαγεῖν·
\vs{7}ἤ· Τίς καταβήσεται εἰς τὴν ἄβυσσον; τοῦτ᾽ ἔστιν Χριστὸν ἐκ νεκρῶν ἀναγαγεῖν.
\vs{8}Ἀλλὰ τί λέγει; Ἐγγύς σου τὸ ῥῆμά ἐστιν ἐν τῷ στόματί σου καὶ ἐν τῇ καρδίᾳ σου, τοῦτ᾽ ἔστιν τὸ ῥῆμα τῆς πίστεως ὃ κηρύσσομεν.
\vs{9}ὅτι ἐὰν ὁμολογήσῃς ἐν τῷ στόματί σου Κύριον Ἰησοῦν καὶ πιστεύσῃς ἐν τῇ καρδίᾳ σου ὅτι ὁ Θεὸς αὐτὸν ἤγειρεν ἐκ νεκρῶν, σωθήσῃ·
\vs{10}καρδίᾳ γὰρ πιστεύεται εἰς δικαιοσύνην, στόματι δὲ ὁμολογεῖται εἰς σωτηρίαν.
\vs{11}Λέγει γὰρ ἡ γραφή· Πᾶς ὁ πιστεύων ἐπ᾽ αὐτῷ οὐ καταισχυνθήσεται.
\vs{12}οὐ γάρ ἐστιν διαστολὴ Ἰουδαίου τε καὶ Ἕλληνος, ὁ γὰρ αὐτὸς Κύριος πάντων, πλουτῶν εἰς πάντας τοὺς ἐπικαλουμένους αὐτόν·
\vs{13}Πᾶς γὰρ ὃς ἂν ἐπικαλέσηται τὸ ὄνομα Κυρίου σωθήσεται.

\vs{14}Πῶς οὖν ἐπικαλέσωνται εἰς ὃν οὐκ ἐπίστευσαν; πῶς δὲ πιστεύσωσιν οὗ οὐκ ἤκουσαν; πῶς δὲ ἀκούσωσιν χωρὶς κηρύσσοντος;
\vs{15}πῶς δὲ κηρύξωσιν ἐὰν μὴ ἀποσταλῶσιν; καθὼς γέγραπται· Ὡς ὡραῖοι οἱ πόδες τῶν εὐαγγελιζομένων τὰ ἀγαθά.
\vs{16}Ἀλλ᾽ οὐ πάντες ὑπήκουσαν τῷ εὐαγγελίῳ. Ἠσαΐας γὰρ λέγει· Κύριε, τίς ἐπίστευσεν τῇ ἀκοῇ ἡμῶν;
\vs{17}ἄρα ἡ πίστις ἐξ ἀκοῆς, ἡ δὲ ἀκοὴ διὰ ῥήματος Χριστοῦ.
\vs{18}Ἀλλὰ λέγω, μὴ οὐκ ἤκουσαν; μενοῦνγε· 
\begin{poetryblock}

\begin{quote}Εἰς πᾶσαν τὴν γῆν ἐξῆλθεν ὁ φθόγγος αὐτῶν\end{quote} 

\begin{quote}καὶ εἰς τὰ πέρατα τῆς οἰκουμένης τὰ ῥήματα αὐτῶν.\end{quote}
\end{poetryblock}

\vs{19}Ἀλλὰ λέγω, μὴ Ἰσραὴλ οὐκ ἔγνω; πρῶτος Μωϋσῆς λέγει· 
\begin{poetryblock}

\begin{quote}Ἐγὼ παραζηλώσω ὑμᾶς ἐπ᾽ οὐκ ἔθνει,\end{quote} 

\begin{quote}ἐπ᾽ ἔθνει ἀσυνέτῳ παροργιῶ ὑμᾶς.\end{quote}
\end{poetryblock}

\vs{20}Ἠσαΐας δὲ ἀποτολμᾷ καὶ λέγει· 
\begin{poetryblock}

\begin{quote}Εὑρέθην ἐν τοῖς ἐμὲ μὴ ζητοῦσιν,\end{quote} 

\begin{quote}ἐμφανὴς ἐγενόμην τοῖς ἐμὲ μὴ ἐπερωτῶσιν.\end{quote}
\end{poetryblock}

\vs{21}Πρὸς δὲ τὸν Ἰσραὴλ λέγει· 
\begin{poetryblock}

\begin{quote}Ὅλην τὴν ἡμέραν ἐξεπέτασα τὰς χεῖράς μου\end{quote} 

\begin{quote}πρὸς λαὸν ἀπειθοῦντα καὶ ἀντιλέγοντα.\end{quote}
\end{poetryblock}

\ch{11}
Λέγω οὖν, μὴ ἀπώσατο ὁ Θεὸς τὸν λαὸν αὐτοῦ; μὴ γένοιτο· καὶ γὰρ ἐγὼ Ἰσραηλίτης εἰμί, ἐκ σπέρματος Ἀβραάμ, φυλῆς Βενιαμίν.
\vs{2}οὐκ ἀπώσατο ὁ Θεὸς τὸν λαὸν αὐτοῦ ὃν προέγνω. ἢ οὐκ οἴδατε ἐν Ἠλίᾳ τί λέγει ἡ γραφή, ὡς ἐντυγχάνει τῷ Θεῷ κατὰ τοῦ Ἰσραήλ;
\vs{3}Κύριε, τοὺς προφήτας σου ἀπέκτειναν, τὰ θυσιαστήριά σου κατέσκαψαν, κἀγὼ ὑπελείφθην μόνος καὶ ζητοῦσιν τὴν ψυχήν μου.
\vs{4}Ἀλλὰ τί λέγει αὐτῷ ὁ χρηματισμός; Κατέλιπον ἐμαυτῷ ἑπτακισχιλίους ἄνδρας, οἵτινες οὐκ ἔκαμψαν γόνυ τῇ Βάαλ.
\vs{5}Οὕτως οὖν καὶ ἐν τῷ νῦν καιρῷ λεῖμμα κατ᾽ ἐκλογὴν χάριτος γέγονεν·
\vs{6}εἰ δὲ χάριτι, οὐκέτι ἐξ ἔργων, ἐπεὶ ἡ χάρις οὐκέτι γίνεται χάρις.
\vs{7}Τί οὖν; ὃ ἐπιζητεῖ Ἰσραήλ, τοῦτο οὐκ ἐπέτυχεν, ἡ δὲ ἐκλογὴ ἐπέτυχεν· οἱ δὲ λοιποὶ ἐπωρώθησαν,
\vs{8}καθὼς γέγραπται· 
\begin{poetryblock}

\begin{quote}Ἔδωκεν αὐτοῖς ὁ Θεὸς πνεῦμα κατανύξεως,\end{quote} 

\begin{quote}ὀφθαλμοὺς τοῦ μὴ βλέπειν καὶ ὦτα τοῦ μὴ ἀκούειν,\end{quote} 

\begin{quote}ἕως τῆς σήμερον ἡμέρας.\end{quote}
\end{poetryblock}

\vs{9}Καὶ Δαυὶδ λέγει· 
\begin{poetryblock}

\begin{quote}Γενηθήτω ἡ τράπεζα αὐτῶν εἰς παγίδα καὶ εἰς θήραν\end{quote} 

\begin{quote}καὶ εἰς σκάνδαλον καὶ εἰς ἀνταπόδομα αὐτοῖς,\end{quote}

\begin{quote} \vs{10}σκοτισθήτωσαν οἱ ὀφθαλμοὶ αὐτῶν τοῦ μὴ βλέπειν\end{quote} 

\begin{quote}καὶ τὸν νῶτον αὐτῶν διὰ παντὸς σύνκαμψον.\end{quote}
\end{poetryblock}

\vs{11}Λέγω οὖν, μὴ ἔπταισαν ἵνα πέσωσιν; μὴ γένοιτο· ἀλλὰ τῷ αὐτῶν παραπτώματι ἡ σωτηρία τοῖς ἔθνεσιν εἰς τὸ παραζηλῶσαι αὐτούς.
\vs{12}εἰ δὲ τὸ παράπτωμα αὐτῶν πλοῦτος κόσμου καὶ τὸ ἥττημα αὐτῶν πλοῦτος ἐθνῶν, πόσῳ μᾶλλον τὸ πλήρωμα αὐτῶν.
\vs{13}Ὑμῖν δὲ λέγω τοῖς ἔθνεσιν· ἐφ᾽ ὅσον μὲν οὖν εἰμι ἐγὼ ἐθνῶν ἀπόστολος, τὴν διακονίαν μου δοξάζω,
\vs{14}εἴ πως παραζηλώσω μου τὴν σάρκα καὶ σώσω τινὰς ἐξ αὐτῶν.
\vs{15}εἰ γὰρ ἡ ἀποβολὴ αὐτῶν καταλλαγὴ κόσμου, τίς ἡ πρόσλημψις εἰ μὴ ζωὴ ἐκ νεκρῶν;
\vs{16}εἰ δὲ ἡ ἀπαρχὴ ἁγία, καὶ τὸ φύραμα· καὶ εἰ ἡ ῥίζα ἁγία, καὶ οἱ κλάδοι.
\vs{17}Εἰ δέ τινες τῶν κλάδων ἐξεκλάσθησαν, σὺ δὲ ἀγριέλαιος ὢν ἐνεκεντρίσθης ἐν αὐτοῖς καὶ συνκοινωνὸς τῆς ῥίζης τῆς πιότητος τῆς ἐλαίας ἐγένου,
\vs{18}μὴ κατακαυχῶ τῶν κλάδων· εἰ δὲ κατακαυχᾶσαι οὐ σὺ τὴν ῥίζαν βαστάζεις ἀλλὰ ἡ ῥίζα σέ.
\vs{19}Ἐρεῖς οὖν· Ἐξεκλάσθησαν κλάδοι ἵνα ἐγὼ ἐγκεντρισθῶ.
\vs{20}καλῶς· τῇ ἀπιστίᾳ ἐξεκλάσθησαν, σὺ δὲ τῇ πίστει ἕστηκας. μὴ ὑψηλὰ φρόνει ἀλλὰ φοβοῦ·
\vs{21}εἰ γὰρ ὁ Θεὸς τῶν κατὰ φύσιν κλάδων οὐκ ἐφείσατο, μή πως οὐδὲ σοῦ φείσεται.
\vs{22}ἴδε οὖν χρηστότητα καὶ ἀποτομίαν Θεοῦ· ἐπὶ μὲν τοὺς πεσόντας ἀποτομία, ἐπὶ δὲ σὲ χρηστότης Θεοῦ, ἐὰν ἐπιμένῃς τῇ χρηστότητι, ἐπεὶ καὶ σὺ ἐκκοπήσῃ.
\vs{23}κἀκεῖνοι δέ, ἐὰν μὴ ἐπιμένωσιν τῇ ἀπιστίᾳ, ἐνκεντρισθήσονται· δυνατὸς γάρ ἐστιν ὁ Θεὸς πάλιν ἐνκεντρίσαι αὐτούς.
\vs{24}εἰ γὰρ σὺ ἐκ τῆς κατὰ φύσιν ἐξεκόπης ἀγριελαίου καὶ παρὰ φύσιν ἐνεκεντρίσθης εἰς καλλιέλαιον, πόσῳ μᾶλλον οὗτοι οἱ κατὰ φύσιν ἐνκεντρισθήσονται τῇ ἰδίᾳ ἐλαίᾳ.

\vs{25}Οὐ γὰρ θέλω ὑμᾶς ἀγνοεῖν, ἀδελφοί, τὸ μυστήριον τοῦτο, ἵνα μὴ ἦτε παρ᾽ ἑαυτοῖς φρόνιμοι, ὅτι πώρωσις ἀπὸ μέρους τῷ Ἰσραὴλ γέγονεν ἄχρι οὗ τὸ πλήρωμα τῶν ἐθνῶν εἰσέλθῃ
\vs{26}καὶ οὕτως πᾶς Ἰσραὴλ σωθήσεται, καθὼς γέγραπται· 
\begin{poetryblock}

\begin{quote}Ἥξει ἐκ Σιὼν ὁ Ῥυόμενος,\end{quote} 

\begin{quote}ἀποστρέψει ἀσεβείας ἀπὸ Ἰακώβ.\end{quote}

\begin{quote} \vs{27}καὶ αὕτη αὐτοῖς ἡ παρ᾽ ἐμοῦ διαθήκη,\end{quote} 

\begin{quote}ὅταν ἀφέλωμαι τὰς ἁμαρτίας αὐτῶν.\end{quote}
\end{poetryblock}

\vs{28}Κατὰ μὲν τὸ εὐαγγέλιον ἐχθροὶ δι᾽ ὑμᾶς, κατὰ δὲ τὴν ἐκλογὴν ἀγαπητοὶ διὰ τοὺς πατέρας·
\vs{29}ἀμεταμέλητα γὰρ τὰ χαρίσματα καὶ ἡ κλῆσις τοῦ Θεοῦ.
\vs{30}Ὥσπερ γὰρ ὑμεῖς ποτε ἠπειθήσατε τῷ Θεῷ, νῦν δὲ ἠλεήθητε τῇ τούτων ἀπειθείᾳ,
\vs{31}οὕτως καὶ οὗτοι νῦν ἠπείθησαν τῷ ὑμετέρῳ ἐλέει, ἵνα καὶ αὐτοὶ νῦν ἐλεηθῶσιν.
\vs{32}συνέκλεισεν γὰρ ὁ Θεὸς τοὺς πάντας εἰς ἀπείθειαν, ἵνα τοὺς πάντας ἐλεήσῃ.
\begin{poetryblock}

\begin{quote} \vs{33}Ὦ βάθος πλούτου\end{quote} 

\begin{quote}καὶ σοφίας καὶ γνώσεως Θεοῦ·\end{quote} 

\begin{quote}ὡς ἀνεξεραύνητα τὰ κρίματα αὐτοῦ\end{quote} 

\begin{quote}καὶ ἀνεξιχνίαστοι αἱ ὁδοὶ αὐτοῦ.\end{quote}

\begin{quote} \vs{34}Τίς γὰρ ἔγνω νοῦν Κυρίου;\end{quote} 

\begin{quote}ἢ τίς σύμβουλος αὐτοῦ ἐγένετο;\end{quote}

\begin{quote} \vs{35}Ἢ τίς προέδωκεν αὐτῷ,\end{quote} 

\begin{quote}καὶ ἀνταποδοθήσεται αὐτῷ;\end{quote}

\begin{quote} \vs{36}ὅτι ἐξ αὐτοῦ καὶ δι᾽ αὐτοῦ καὶ εἰς αὐτὸν τὰ πάντα·\end{quote} 

\begin{quote}αὐτῷ ἡ δόξα εἰς τοὺς αἰῶνας, ἀμήν.\end{quote}
\end{poetryblock}

\ch{12}
Παρακαλῶ οὖν ὑμᾶς, ἀδελφοί, διὰ τῶν οἰκτιρμῶν τοῦ Θεοῦ παραστῆσαι τὰ σώματα ὑμῶν θυσίαν ζῶσαν ἁγίαν εὐάρεστον τῷ Θεῷ, τὴν λογικὴν λατρείαν ὑμῶν·
\vs{2}καὶ μὴ συσχηματίζεσθε τῷ αἰῶνι τούτῳ, ἀλλὰ μεταμορφοῦσθε τῇ ἀνακαινώσει τοῦ νοός εἰς τὸ δοκιμάζειν ὑμᾶς τί τὸ θέλημα τοῦ Θεοῦ, τὸ ἀγαθὸν καὶ εὐάρεστον καὶ τέλειον.

\vs{3}Λέγω γὰρ διὰ τῆς χάριτος τῆς δοθείσης μοι παντὶ τῷ ὄντι ἐν ὑμῖν μὴ ὑπερφρονεῖν παρ᾽ ὃ δεῖ φρονεῖν ἀλλὰ φρονεῖν εἰς τὸ σωφρονεῖν, ἑκάστῳ ὡς ὁ Θεὸς ἐμέρισεν μέτρον πίστεως.
\vs{4}καθάπερ γὰρ ἐν ἑνὶ σώματι πολλὰ μέλη ἔχομεν, τὰ δὲ μέλη πάντα οὐ τὴν αὐτὴν ἔχει πρᾶξιν,
\vs{5}οὕτως οἱ πολλοὶ ἓν σῶμά ἐσμεν ἐν Χριστῷ, τὸ δὲ καθ᾽ εἷς ἀλλήλων μέλη.
\vs{6}Ἔχοντες δὲ χαρίσματα κατὰ τὴν χάριν τὴν δοθεῖσαν ἡμῖν διάφορα, εἴτε προφητείαν κατὰ τὴν ἀναλογίαν τῆς πίστεως,
\vs{7}εἴτε διακονίαν ἐν τῇ διακονίᾳ, εἴτε ὁ διδάσκων ἐν τῇ διδασκαλίᾳ,
\vs{8}εἴτε ὁ παρακαλῶν ἐν τῇ παρακλήσει· ὁ μεταδιδοὺς ἐν ἁπλότητι, ὁ προϊστάμενος ἐν σπουδῇ, ὁ ἐλεῶν ἐν ἱλαρότητι.

\vs{9}Ἡ ἀγάπη ἀνυπόκριτος. ἀποστυγοῦντες τὸ πονηρόν, κολλώμενοι τῷ ἀγαθῷ,
\vs{10}τῇ φιλαδελφίᾳ εἰς ἀλλήλους φιλόστοργοι, τῇ τιμῇ ἀλλήλους προηγούμενοι,
\vs{11}τῇ σπουδῇ μὴ ὀκνηροί, τῷ πνεύματι ζέοντες, τῷ Κυρίῳ δουλεύοντες,
\vs{12}τῇ ἐλπίδι χαίροντες, τῇ θλίψει ὑπομένοντες, τῇ προσευχῇ προσκαρτεροῦντες,
\vs{13}ταῖς χρείαις τῶν ἁγίων κοινωνοῦντες, τὴν φιλοξενίαν διώκοντες.
\vs{14}Εὐλογεῖτε τοὺς διώκοντας ὑμᾶς, εὐλογεῖτε καὶ μὴ καταρᾶσθε.
\vs{15}χαίρειν μετὰ χαιρόντων, κλαίειν μετὰ κλαιόντων.
\vs{16}τὸ αὐτὸ εἰς ἀλλήλους φρονοῦντες, μὴ τὰ ὑψηλὰ φρονοῦντες ἀλλὰ τοῖς ταπεινοῖς συναπαγόμενοι. μὴ γίνεσθε φρόνιμοι παρ᾽ ἑαυτοῖς.
\vs{17}μηδενὶ κακὸν ἀντὶ κακοῦ ἀποδιδόντες, προνοούμενοι καλὰ ἐνώπιον πάντων ἀνθρώπων·
\vs{18}εἰ δυνατόν τὸ ἐξ ὑμῶν, μετὰ πάντων ἀνθρώπων εἰρηνεύοντες·
\vs{19}μὴ ἑαυτοὺς ἐκδικοῦντες, ἀγαπητοί, ἀλλὰ δότε τόπον τῇ ὀργῇ, γέγραπται γάρ· Ἐμοὶ ἐκδίκησις, ἐγὼ ἀνταποδώσω, λέγει Κύριος.
\vs{20}Ἀλλὰ Ἐὰν πεινᾷ ὁ ἐχθρός σου, ψώμιζε αὐτόν· ἐὰν διψᾷ, πότιζε αὐτόν· τοῦτο γὰρ ποιῶν ἄνθρακας πυρὸς σωρεύσεις ἐπὶ τὴν κεφαλὴν αὐτοῦ.
\vs{21}Μὴ νικῶ ὑπὸ τοῦ κακοῦ ἀλλὰ νίκα ἐν τῷ ἀγαθῷ τὸ κακόν.

\ch{13}
Πᾶσα ψυχὴ ἐξουσίαις ὑπερεχούσαις ὑποτασσέσθω. οὐ γὰρ ἔστιν ἐξουσία εἰ μὴ ὑπὸ Θεοῦ, αἱ δὲ οὖσαι ὑπὸ Θεοῦ τεταγμέναι εἰσίν.
\vs{2}ὥστε ὁ ἀντιτασσόμενος τῇ ἐξουσίᾳ τῇ τοῦ Θεοῦ διαταγῇ ἀνθέστηκεν, οἱ δὲ ἀνθεστηκότες ἑαυτοῖς κρίμα λήμψονται.
\vs{3}οἱ γὰρ ἄρχοντες οὐκ εἰσὶν φόβος τῷ ἀγαθῷ ἔργῳ ἀλλὰ τῷ κακῷ. θέλεις δὲ μὴ φοβεῖσθαι τὴν ἐξουσίαν· τὸ ἀγαθὸν ποίει, καὶ ἕξεις ἔπαινον ἐξ αὐτῆς·
\vs{4}Θεοῦ γὰρ διάκονός ἐστιν σοὶ εἰς τὸ ἀγαθόν. ἐὰν δὲ τὸ κακὸν ποιῇς, φοβοῦ· οὐ γὰρ εἰκῇ τὴν μάχαιραν φορεῖ· Θεοῦ γὰρ διάκονός ἐστιν ἔκδικος εἰς ὀργὴν τῷ τὸ κακὸν πράσσοντι.
\vs{5}Διὸ ἀνάγκη ὑποτάσσεσθαι, οὐ μόνον διὰ τὴν ὀργὴν ἀλλὰ καὶ διὰ τὴν συνείδησιν.
\vs{6}διὰ τοῦτο γὰρ καὶ φόρους τελεῖτε· λειτουργοὶ γὰρ Θεοῦ εἰσιν εἰς αὐτὸ τοῦτο προσκαρτεροῦντες.
\vs{7}ἀπόδοτε πᾶσιν τὰς ὀφειλάς, τῷ τὸν φόρον τὸν φόρον, τῷ τὸ τέλος τὸ τέλος, τῷ τὸν φόβον τὸν φόβον, τῷ τὴν τιμὴν τὴν τιμήν.

\vs{8}Μηδενὶ μηδὲν ὀφείλετε εἰ μὴ τὸ ἀλλήλους ἀγαπᾶν· ὁ γὰρ ἀγαπῶν τὸν ἕτερον νόμον πεπλήρωκεν.
\vs{9}τὸ γάρ Οὐ μοιχεύσεις, Οὐ φονεύσεις, Οὐ κλέψεις, Οὐκ ἐπιθυμήσεις, καὶ εἴ τις ἑτέρα ἐντολή, ἐν τῷ λόγῳ τούτῳ ἀνακεφαλαιοῦται ἐν τῷ· Ἀγαπήσεις τὸν πλησίον σου ὡς σεαυτόν.
\vs{10}ἡ ἀγάπη τῷ πλησίον κακὸν οὐκ ἐργάζεται· πλήρωμα οὖν νόμου ἡ ἀγάπη.
\vs{11}Καὶ τοῦτο εἰδότες τὸν καιρόν, ὅτι ὥρα ἤδη ὑμᾶς ἐξ ὕπνου ἐγερθῆναι, νῦν γὰρ ἐγγύτερον ἡμῶν ἡ σωτηρία ἢ ὅτε ἐπιστεύσαμεν.
\vs{12}ἡ νὺξ προέκοψεν, ἡ δὲ ἡμέρα ἤγγικεν. ἀποθώμεθα οὖν τὰ ἔργα τοῦ σκότους, ἐνδυσώμεθα δὲ τὰ ὅπλα τοῦ φωτός.
\vs{13}ὡς ἐν ἡμέρᾳ εὐσχημόνως περιπατήσωμεν, μὴ κώμοις καὶ μέθαις, μὴ κοίταις καὶ ἀσελγείαις, μὴ ἔριδι καὶ ζήλῳ,
\vs{14}ἀλλὰ ἐνδύσασθε τὸν Κύριον Ἰησοῦν Χριστόν καὶ τῆς σαρκὸς πρόνοιαν μὴ ποιεῖσθε εἰς ἐπιθυμίας.

\ch{14}
Τὸν δὲ ἀσθενοῦντα τῇ πίστει προσλαμβάνεσθε, μὴ εἰς διακρίσεις διαλογισμῶν.
\vs{2}ὃς μὲν πιστεύει φαγεῖν πάντα, ὁ δὲ ἀσθενῶν λάχανα ἐσθίει.
\vs{3}ὁ ἐσθίων τὸν μὴ ἐσθίοντα μὴ ἐξουθενείτω, ὁ δὲ μὴ ἐσθίων τὸν ἐσθίοντα μὴ κρινέτω, ὁ Θεὸς γὰρ αὐτὸν προσελάβετο.
\vs{4}σὺ τίς εἶ ὁ κρίνων ἀλλότριον οἰκέτην; τῷ ἰδίῳ κυρίῳ στήκει ἢ πίπτει· σταθήσεται δέ, δυνατεῖ γὰρ ὁ Κύριος στῆσαι αὐτόν.
\vs{5}Ὃς μὲν γὰρ κρίνει ἡμέραν παρ᾽ ἡμέραν, ὃς δὲ κρίνει πᾶσαν ἡμέραν· ἕκαστος ἐν τῷ ἰδίῳ νοῒ πληροφορείσθω.
\vs{6}ὁ φρονῶν τὴν ἡμέραν Κυρίῳ φρονεῖ· καὶ ὁ ἐσθίων Κυρίῳ ἐσθίει, εὐχαριστεῖ γὰρ τῷ Θεῷ· καὶ ὁ μὴ ἐσθίων Κυρίῳ οὐκ ἐσθίει καὶ εὐχαριστεῖ τῷ Θεῷ.
\vs{7}Οὐδεὶς γὰρ ἡμῶν ἑαυτῷ ζῇ καὶ οὐδεὶς ἑαυτῷ ἀποθνῄσκει·
\vs{8}ἐάν τε γὰρ ζῶμεν, τῷ Κυρίῳ ζῶμεν, ἐάν τε ἀποθνήσκωμεν, τῷ Κυρίῳ ἀποθνήσκομεν. ἐάν τε οὖν ζῶμεν ἐάν τε ἀποθνήσκωμεν, τοῦ Κυρίου ἐσμέν.
\vs{9}εἰς τοῦτο γὰρ Χριστὸς ἀπέθανεν καὶ ἔζησεν, ἵνα καὶ νεκρῶν καὶ ζώντων κυριεύσῃ.
\vs{10}Σὺ δὲ τί κρίνεις τὸν ἀδελφόν σου; ἢ καὶ σὺ τί ἐξουθενεῖς τὸν ἀδελφόν σου; πάντες γὰρ παραστησόμεθα τῷ βήματι τοῦ Θεοῦ,
\vs{11}γέγραπται γάρ· 
\begin{poetryblock}

\begin{quote}Ζῶ ἐγώ, λέγει Κύριος, ὅτι ἐμοὶ κάμψει πᾶν γόνυ\end{quote} 

\begin{quote}καὶ πᾶσα γλῶσσα ἐξομολογήσεται τῷ Θεῷ.\end{quote}
\end{poetryblock}

\vs{12}Ἄρα οὖν ἕκαστος ἡμῶν περὶ ἑαυτοῦ λόγον δώσει τῷ Θεῷ.

\vs{13}Μηκέτι οὖν ἀλλήλους κρίνωμεν· ἀλλὰ τοῦτο κρίνατε μᾶλλον, τὸ μὴ τιθέναι πρόσκομμα τῷ ἀδελφῷ ἢ σκάνδαλον.
\vs{14}Οἶδα καὶ πέπεισμαι ἐν Κυρίῳ Ἰησοῦ ὅτι οὐδὲν κοινὸν δι᾽ ἑαυτοῦ, εἰ μὴ τῷ λογιζομένῳ τι κοινὸν εἶναι, ἐκείνῳ κοινόν.
\vs{15}εἰ γὰρ διὰ βρῶμα ὁ ἀδελφός σου λυπεῖται, οὐκέτι κατὰ ἀγάπην περιπατεῖς· μὴ τῷ βρώματί σου ἐκεῖνον ἀπόλλυε ὑπὲρ οὗ Χριστὸς ἀπέθανεν.
\vs{16}Μὴ βλασφημείσθω οὖν ὑμῶν τὸ ἀγαθόν.
\vs{17}οὐ γάρ ἐστιν ἡ βασιλεία τοῦ Θεοῦ βρῶσις καὶ πόσις ἀλλὰ δικαιοσύνη καὶ εἰρήνη καὶ χαρὰ ἐν Πνεύματι Ἁγίῳ·
\vs{18}ὁ γὰρ ἐν τούτῳ δουλεύων τῷ Χριστῷ εὐάρεστος τῷ Θεῷ καὶ δόκιμος τοῖς ἀνθρώποις.
\vs{19}Ἄρα οὖν τὰ τῆς εἰρήνης διώκωμεν καὶ τὰ τῆς οἰκοδομῆς τῆς εἰς ἀλλήλους.
\vs{20}μὴ ἕνεκεν βρώματος κατάλυε τὸ ἔργον τοῦ Θεοῦ. πάντα μὲν καθαρά, ἀλλὰ κακὸν τῷ ἀνθρώπῳ τῷ διὰ προσκόμματος ἐσθίοντι.
\vs{21}καλὸν τὸ μὴ φαγεῖν κρέα μηδὲ πιεῖν οἶνον μηδὲ ἐν ᾧ ὁ ἀδελφός σου προσκόπτει.
\vs{22}Σὺ πίστιν ἣν ἔχεις κατὰ σεαυτὸν ἔχε ἐνώπιον τοῦ Θεοῦ. μακάριος ὁ μὴ κρίνων ἑαυτὸν ἐν ᾧ δοκιμάζει·
\vs{23}ὁ δὲ διακρινόμενος ἐὰν φάγῃ κατακέκριται, ὅτι οὐκ ἐκ πίστεως· πᾶν δὲ ὃ οὐκ ἐκ πίστεως ἁμαρτία ἐστίν.

\ch{15}
Ὀφείλομεν δὲ ἡμεῖς οἱ δυνατοὶ τὰ ἀσθενήματα τῶν ἀδυνάτων βαστάζειν καὶ μὴ ἑαυτοῖς ἀρέσκειν.
\vs{2}ἕκαστος ἡμῶν τῷ πλησίον ἀρεσκέτω εἰς τὸ ἀγαθὸν πρὸς οἰκοδομήν·
\vs{3}καὶ γὰρ ὁ Χριστὸς οὐχ ἑαυτῷ ἤρεσεν, ἀλλὰ καθὼς γέγραπται· Οἱ ὀνειδισμοὶ τῶν ὀνειδιζόντων σε ἐπέπεσαν ἐπ᾽ ἐμέ.
\vs{4}ὅσα γὰρ προεγράφη, εἰς τὴν ἡμετέραν διδασκαλίαν ἐγράφη, ἵνα διὰ τῆς ὑπομονῆς καὶ διὰ τῆς παρακλήσεως τῶν γραφῶν τὴν ἐλπίδα ἔχωμεν.
\vs{5}Ὁ δὲ Θεὸς τῆς ὑπομονῆς καὶ τῆς παρακλήσεως δῴη ὑμῖν τὸ αὐτὸ φρονεῖν ἐν ἀλλήλοις κατὰ Χριστὸν Ἰησοῦν,
\vs{6}ἵνα ὁμοθυμαδὸν ἐν ἑνὶ στόματι δοξάζητε τὸν Θεὸν καὶ Πατέρα τοῦ Κυρίου ἡμῶν Ἰησοῦ Χριστοῦ.

\vs{7}Διὸ προσλαμβάνεσθε ἀλλήλους, καθὼς καὶ ὁ Χριστὸς προσελάβετο ὑμᾶς εἰς δόξαν τοῦ Θεοῦ.
\vs{8}λέγω γὰρ Χριστὸν διάκονον γεγενῆσθαι περιτομῆς ὑπὲρ ἀληθείας Θεοῦ, εἰς τὸ βεβαιῶσαι τὰς ἐπαγγελίας τῶν πατέρων,
\vs{9}τὰ δὲ ἔθνη ὑπὲρ ἐλέους δοξάσαι τὸν Θεόν, καθὼς γέγραπται· 
\begin{poetryblock}

\begin{quote}Διὰ τοῦτο ἐξομολογήσομαί σοι ἐν ἔθνεσιν\end{quote} 

\begin{quote}καὶ τῷ ὀνόματί σου ψαλῶ.\end{quote}
\end{poetryblock}

\vs{10}Καὶ πάλιν λέγει· 
\begin{poetryblock}

\begin{quote}Εὐφράνθητε, ἔθνη, μετὰ τοῦ λαοῦ αὐτοῦ.\end{quote}
\end{poetryblock}

\vs{11}Καὶ πάλιν· Αἰνεῖτε, πάντα τὰ ἔθνη, τὸν Κύριον καὶ ἐπαινεσάτωσαν αὐτὸν πάντες οἱ λαοί.
\vs{12}Καὶ πάλιν Ἠσαΐας λέγει· 
\begin{poetryblock}

\begin{quote}Ἔσται ἡ ῥίζα τοῦ Ἰεσσαί\end{quote} 

\begin{quote}καὶ ὁ ἀνιστάμενος ἄρχειν ἐθνῶν,\end{quote} 

\begin{quote}ἐπ᾽ αὐτῷ ἔθνη ἐλπιοῦσιν.\end{quote}
\end{poetryblock}

\vs{13}Ὁ δὲ Θεὸς τῆς ἐλπίδος πληρώσαι ὑμᾶς πάσης χαρᾶς καὶ εἰρήνης ἐν τῷ πιστεύειν, εἰς τὸ περισσεύειν ὑμᾶς ἐν τῇ ἐλπίδι ἐν δυνάμει Πνεύματος Ἁγίου.

\vs{14}Πέπεισμαι δέ, ἀδελφοί μου, καὶ αὐτὸς ἐγὼ περὶ ὑμῶν ὅτι καὶ αὐτοὶ μεστοί ἐστε ἀγαθωσύνης, πεπληρωμένοι πάσης τῆς γνώσεως, δυνάμενοι καὶ ἀλλήλους νουθετεῖν.
\vs{15}τολμηρότερον δὲ ἔγραψα ὑμῖν ἀπὸ μέρους ὡς ἐπαναμιμνήσκων ὑμᾶς διὰ τὴν χάριν τὴν δοθεῖσάν μοι ὑπὸ τοῦ Θεοῦ
\vs{16}εἰς τὸ εἶναί με λειτουργὸν Χριστοῦ Ἰησοῦ εἰς τὰ ἔθνη, ἱερουργοῦντα τὸ εὐαγγέλιον τοῦ Θεοῦ, ἵνα γένηται ἡ προσφορὰ τῶν ἐθνῶν εὐπρόσδεκτος, ἡγιασμένη ἐν Πνεύματι Ἁγίῳ.
\vs{17}Ἔχω οὖν τὴν καύχησιν ἐν Χριστῷ Ἰησοῦ τὰ πρὸς τὸν Θεόν·
\vs{18}οὐ γὰρ τολμήσω τι λαλεῖν ὧν οὐ κατειργάσατο Χριστὸς δι᾽ ἐμοῦ εἰς ὑπακοὴν ἐθνῶν, λόγῳ καὶ ἔργῳ,
\vs{19}ἐν δυνάμει σημείων καὶ τεράτων, ἐν δυνάμει Πνεύματος θεοῦ· ὥστε με ἀπὸ Ἰερουσαλὴμ καὶ κύκλῳ μέχρι τοῦ Ἰλλυρικοῦ πεπληρωκέναι τὸ εὐαγγέλιον τοῦ Χριστοῦ,
\vs{20}οὕτως δὲ φιλοτιμούμενον εὐαγγελίζεσθαι οὐχ ὅπου ὠνομάσθη Χριστός, ἵνα μὴ ἐπ᾽ ἀλλότριον θεμέλιον οἰκοδομῶ,
\vs{21}ἀλλὰ καθὼς γέγραπται· 
\begin{poetryblock}

\begin{quote}Οἷς οὐκ ἀνηγγέλη περὶ αὐτοῦ ὄψονται,\end{quote} 

\begin{quote}καὶ οἳ οὐκ ἀκηκόασιν συνήσουσιν.\end{quote}
\end{poetryblock}

\vs{22}Διὸ καὶ ἐνεκοπτόμην τὰ πολλὰ τοῦ ἐλθεῖν πρὸς ὑμᾶς·
\vs{23}Νυνὶ δὲ μηκέτι τόπον ἔχων ἐν τοῖς κλίμασι τούτοις, ἐπιποθίαν δὲ ἔχων τοῦ ἐλθεῖν πρὸς ὑμᾶς ἀπὸ ἱκανῶν ἐτῶν,
\vs{24}ὡς ἂν πορεύωμαι εἰς τὴν Σπανίαν· ἐλπίζω γὰρ διαπορευόμενος θεάσασθαι ὑμᾶς καὶ ὑφ᾽ ὑμῶν προπεμφθῆναι ἐκεῖ ἐὰν ὑμῶν πρῶτον ἀπὸ μέρους ἐμπλησθῶ.
\vs{25}Νυνὶ δὲ πορεύομαι εἰς Ἰερουσαλὴμ διακονῶν τοῖς ἁγίοις.
\vs{26}εὐδόκησαν γὰρ Μακεδονία καὶ Ἀχαΐα κοινωνίαν τινὰ ποιήσασθαι εἰς τοὺς πτωχοὺς τῶν ἁγίων τῶν ἐν Ἰερουσαλήμ.
\vs{27}εὐδόκησαν γάρ καὶ ὀφειλέται εἰσὶν αὐτῶν· εἰ γὰρ τοῖς πνευματικοῖς αὐτῶν ἐκοινώνησαν τὰ ἔθνη, ὀφείλουσιν καὶ ἐν τοῖς σαρκικοῖς λειτουργῆσαι αὐτοῖς.
\vs{28}Τοῦτο οὖν ἐπιτελέσας καὶ σφραγισάμενος αὐτοῖς τὸν καρπὸν τοῦτον, ἀπελεύσομαι δι᾽ ὑμῶν εἰς Σπανίαν·
\vs{29}οἶδα δὲ ὅτι ἐρχόμενος πρὸς ὑμᾶς ἐν πληρώματι εὐλογίας Χριστοῦ ἐλεύσομαι.

\vs{30}Παρακαλῶ δὲ ὑμᾶς, ἀδελφοί, διὰ τοῦ Κυρίου ἡμῶν Ἰησοῦ Χριστοῦ καὶ διὰ τῆς ἀγάπης τοῦ Πνεύματος συναγωνίσασθαί μοι ἐν ταῖς προσευχαῖς ὑπὲρ ἐμοῦ πρὸς τὸν Θεόν,
\vs{31}ἵνα ῥυσθῶ ἀπὸ τῶν ἀπειθούντων ἐν τῇ Ἰουδαίᾳ καὶ ἡ διακονία μου ἡ εἰς Ἰερουσαλὴμ εὐπρόσδεκτος τοῖς ἁγίοις γένηται,
\vs{32}ἵνα ἐν χαρᾷ ἐλθὼν πρὸς ὑμᾶς διὰ θελήματος Θεοῦ συναναπαύσωμαι ὑμῖν.
\vs{33}Ὁ δὲ Θεὸς τῆς εἰρήνης μετὰ πάντων ὑμῶν, ἀμήν.

\ch{16}
Συνίστημι δὲ ὑμῖν Φοίβην τὴν ἀδελφὴν ἡμῶν, οὖσαν καὶ διάκονον τῆς ἐκκλησίας τῆς ἐν Κενχρεαῖς,
\vs{2}ἵνα αὐτὴν προσδέξησθε ἐν Κυρίῳ ἀξίως τῶν ἁγίων καὶ παραστῆτε αὐτῇ ἐν ᾧ ἂν ὑμῶν χρῄζῃ πράγματι· καὶ γὰρ αὐτὴ προστάτις πολλῶν ἐγενήθη καὶ ἐμοῦ αὐτοῦ.

\vs{3}Ἀσπάσασθε Πρίσκαν καὶ Ἀκύλαν τοὺς συνεργούς μου ἐν Χριστῷ Ἰησοῦ,
\vs{4}οἵτινες ὑπὲρ τῆς ψυχῆς μου τὸν ἑαυτῶν τράχηλον ὑπέθηκαν, οἷς οὐκ ἐγὼ μόνος εὐχαριστῶ ἀλλὰ καὶ πᾶσαι αἱ ἐκκλησίαι τῶν ἐθνῶν,
\vs{5}καὶ τὴν κατ᾽ οἶκον αὐτῶν ἐκκλησίαν. Ἀσπάσασθε Ἐπαίνετον τὸν ἀγαπητόν μου, ὅς ἐστιν ἀπαρχὴ τῆς Ἀσίας εἰς Χριστόν.
\vs{6}Ἀσπάσασθε Μαριάν, ἥτις πολλὰ ἐκοπίασεν εἰς ὑμᾶς.
\vs{7}Ἀσπάσασθε Ἀνδρόνικον καὶ Ἰουνίαν τοὺς συγγενεῖς μου καὶ συναιχμαλώτους μου, οἵτινές εἰσιν ἐπίσημοι ἐν τοῖς ἀποστόλοις, οἳ καὶ πρὸ ἐμοῦ γέγοναν ἐν Χριστῷ.
\vs{8}Ἀσπάσασθε Ἀμπλιᾶτον τὸν ἀγαπητόν μου ἐν Κυρίῳ.
\vs{9}Ἀσπάσασθε Οὐρβανὸν τὸν συνεργὸν ἡμῶν ἐν Χριστῷ καὶ Στάχυν τὸν ἀγαπητόν μου.
\vs{10}Ἀσπάσασθε Ἀπελλῆν τὸν δόκιμον ἐν Χριστῷ. Ἀσπάσασθε τοὺς ἐκ τῶν Ἀριστοβούλου.
\vs{11}Ἀσπάσασθε Ἡρῳδίωνα τὸν συγγενῆ μου. Ἀσπάσασθε τοὺς ἐκ τῶν Ναρκίσσου τοὺς ὄντας ἐν Κυρίῳ.
\vs{12}Ἀσπάσασθε Τρύφαιναν καὶ Τρυφῶσαν τὰς κοπιώσας ἐν Κυρίῳ. Ἀσπάσασθε Περσίδα τὴν ἀγαπητήν, ἥτις πολλὰ ἐκοπίασεν ἐν Κυρίῳ.
\vs{13}Ἀσπάσασθε Ῥοῦφον τὸν ἐκλεκτὸν ἐν Κυρίῳ καὶ τὴν μητέρα αὐτοῦ καὶ ἐμοῦ.
\vs{14}Ἀσπάσασθε Ἀσύνκριτον, Φλέγοντα, Ἑρμῆν, Πατρόβαν, Ἑρμᾶν καὶ τοὺς σὺν αὐτοῖς ἀδελφούς.
\vs{15}Ἀσπάσασθε Φιλόλογον καὶ Ἰουλίαν, Νηρέα καὶ τὴν ἀδελφὴν αὐτοῦ, καὶ Ὀλυμπᾶν καὶ τοὺς σὺν αὐτοῖς πάντας ἁγίους.
\vs{16}Ἀσπάσασθε ἀλλήλους ἐν φιλήματι ἁγίῳ. Ἀσπάζονται ὑμᾶς αἱ ἐκκλησίαι πᾶσαι τοῦ Χριστοῦ.

\vs{17}Παρακαλῶ δὲ ὑμᾶς, ἀδελφοί, σκοπεῖν τοὺς τὰς διχοστασίας καὶ τὰ σκάνδαλα παρὰ τὴν διδαχὴν ἣν ὑμεῖς ἐμάθετε ποιοῦντας, καὶ ἐκκλίνετε ἀπ᾽ αὐτῶν·
\vs{18}οἱ γὰρ τοιοῦτοι τῷ Κυρίῳ ἡμῶν Χριστῷ οὐ δουλεύουσιν ἀλλὰ τῇ ἑαυτῶν κοιλίᾳ, καὶ διὰ τῆς χρηστολογίας καὶ εὐλογίας ἐξαπατῶσιν τὰς καρδίας τῶν ἀκάκων.
\vs{19}Ἡ γὰρ ὑμῶν ὑπακοὴ εἰς πάντας ἀφίκετο· ἐφ᾽ ὑμῖν οὖν χαίρω, θέλω δὲ ὑμᾶς σοφοὺς εἶναι εἰς τὸ ἀγαθόν, ἀκεραίους δὲ εἰς τὸ κακόν.
\vs{20}Ὁ δὲ Θεὸς τῆς εἰρήνης συντρίψει τὸν Σατανᾶν ὑπὸ τοὺς πόδας ὑμῶν ἐν τάχει. Ἡ χάρις τοῦ Κυρίου ἡμῶν Ἰησοῦ μεθ᾽ ὑμῶν.

\vs{21}Ἀσπάζεται ὑμᾶς Τιμόθεος ὁ συνεργός μου καὶ Λούκιος καὶ Ἰάσων καὶ Σωσίπατρος οἱ συγγενεῖς μου.
\vs{22}Ἀσπάζομαι ὑμᾶς ἐγὼ Τέρτιος ὁ γράψας τὴν ἐπιστολὴν ἐν Κυρίῳ.
\vs{23}Ἀσπάζεται ὑμᾶς Γάϊος ὁ ξένος μου καὶ ὅλης τῆς ἐκκλησίας. Ἀσπάζεται ὑμᾶς Ἔραστος ὁ οἰκονόμος τῆς πόλεως καὶ Κούαρτος ὁ ἀδελφός.

\vs{25}Τῷ δὲ δυναμένῳ ὑμᾶς στηρίξαι κατὰ τὸ εὐαγγέλιόν μου καὶ τὸ κήρυγμα Ἰησοῦ Χριστοῦ, κατὰ ἀποκάλυψιν μυστηρίου χρόνοις αἰωνίοις σεσιγημένου,
\vs{26}φανερωθέντος δὲ νῦν διά τε γραφῶν προφητικῶν κατ᾽ ἐπιταγὴν τοῦ αἰωνίου Θεοῦ εἰς ὑπακοὴν πίστεως εἰς πάντα τὰ ἔθνη γνωρισθέντος,
\vs{27}μόνῳ σοφῷ Θεῷ, διὰ Ἰησοῦ Χριστοῦ, ᾧ ἡ δόξα εἰς τοὺς αἰῶνας, ἀμήν.


\def\book{ΠΡΟΣ ΚΟΡΙΝΘΙΟΥΣ Α}
\biblebook{ΠΡΟΣ ΚΟΡΙΝΘΙΟΥΣ Α}


\lettrine[lines=2, loversize=0.2, nindent=0em, findent=.25em]{\textcolor{bookheadingcolor}{Π}}{αῦλος} κλητὸς ἀπόστολος Χριστοῦ Ἰησοῦ διὰ θελήματος Θεοῦ καὶ Σωσθένης ὁ ἀδελφὸς
\vs{2}Τῇ ἐκκλησίᾳ τοῦ Θεοῦ τῇ οὔσῃ ἐν Κορίνθῳ, ἡγιασμένοις ἐν Χριστῷ Ἰησοῦ, κλητοῖς ἁγίοις, σὺν πᾶσιν τοῖς ἐπικαλουμένοις τὸ ὄνομα τοῦ Κυρίου ἡμῶν Ἰησοῦ Χριστοῦ ἐν παντὶ τόπῳ, αὐτῶν καὶ ἡμῶν·
\vs{3}Χάρις ὑμῖν καὶ εἰρήνη ἀπὸ Θεοῦ Πατρὸς ἡμῶν καὶ Κυρίου Ἰησοῦ Χριστοῦ.

\vs{4}Εὐχαριστῶ τῷ Θεῷ μου πάντοτε περὶ ὑμῶν ἐπὶ τῇ χάριτι τοῦ Θεοῦ τῇ δοθείσῃ ὑμῖν ἐν Χριστῷ Ἰησοῦ,
\vs{5}ὅτι ἐν παντὶ ἐπλουτίσθητε ἐν αὐτῷ, ἐν παντὶ λόγῳ καὶ πάσῃ γνώσει,
\vs{6}καθὼς τὸ μαρτύριον τοῦ Χριστοῦ ἐβεβαιώθη ἐν ὑμῖν,
\vs{7}ὥστε ὑμᾶς μὴ ὑστερεῖσθαι ἐν μηδενὶ χαρίσματι ἀπεκδεχομένους τὴν ἀποκάλυψιν τοῦ Κυρίου ἡμῶν Ἰησοῦ Χριστοῦ·
\vs{8}ὃς καὶ βεβαιώσει ὑμᾶς ἕως τέλους ἀνεγκλήτους ἐν τῇ ἡμέρᾳ τοῦ Κυρίου ἡμῶν Ἰησοῦ Χριστοῦ.
\vs{9}πιστὸς ὁ Θεὸς, δι᾽ οὗ ἐκλήθητε εἰς κοινωνίαν τοῦ Υἱοῦ αὐτοῦ Ἰησοῦ Χριστοῦ τοῦ Κυρίου ἡμῶν.

\vs{10}Παρακαλῶ δὲ ὑμᾶς, ἀδελφοί, διὰ τοῦ ὀνόματος τοῦ Κυρίου ἡμῶν Ἰησοῦ Χριστοῦ, ἵνα τὸ αὐτὸ λέγητε πάντες καὶ μὴ ᾖ ἐν ὑμῖν σχίσματα, ἦτε δὲ κατηρτισμένοι ἐν τῷ αὐτῷ νοῒ καὶ ἐν τῇ αὐτῇ γνώμῃ.
\vs{11}ἐδηλώθη γάρ μοι περὶ ὑμῶν, ἀδελφοί μου, ὑπὸ τῶν Χλόης ὅτι ἔριδες ἐν ὑμῖν εἰσιν.
\vs{12}λέγω δὲ τοῦτο ὅτι ἕκαστος ὑμῶν λέγει· Ἐγὼ μέν εἰμι Παύλου, Ἐγὼ δὲ Ἀπολλῶ, Ἐγὼ δὲ Κηφᾶ, Ἐγὼ δὲ Χριστοῦ.
\vs{13}Μεμέρισται ὁ Χριστός; μὴ Παῦλος ἐσταυρώθη ὑπὲρ ὑμῶν, ἢ εἰς τὸ ὄνομα Παύλου ἐβαπτίσθητε;
\vs{14}εὐχαριστῶ τῷ θεῷ ὅτι οὐδένα ὑμῶν ἐβάπτισα εἰ μὴ Κρίσπον καὶ Γάϊον,
\vs{15}ἵνα μή τις εἴπῃ ὅτι εἰς τὸ ἐμὸν ὄνομα ἐβαπτίσθητε.
\vs{16}ἐβάπτισα δὲ καὶ τὸν Στεφανᾶ οἶκον, λοιπὸν οὐκ οἶδα εἴ τινα ἄλλον ἐβάπτισα.
\vs{17}οὐ γὰρ ἀπέστειλέν με Χριστὸς βαπτίζειν ἀλλὰ εὐαγγελίζεσθαι, οὐκ ἐν σοφίᾳ λόγου, ἵνα μὴ κενωθῇ ὁ σταυρὸς τοῦ Χριστοῦ.

\vs{18}Ὁ λόγος γὰρ ὁ τοῦ σταυροῦ τοῖς μὲν ἀπολλυμένοις μωρία ἐστίν, τοῖς δὲ σῳζομένοις ἡμῖν δύναμις Θεοῦ ἐστιν.
\vs{19}γέγραπται γάρ· 
\begin{poetryblock}

\begin{quote}Ἀπολῶ τὴν σοφίαν τῶν σοφῶν\end{quote} 

\begin{quote}καὶ τὴν σύνεσιν τῶν συνετῶν ἀθετήσω.\end{quote}
\end{poetryblock}

\vs{20}Ποῦ σοφός; ποῦ γραμματεύς; ποῦ συζητητὴς τοῦ αἰῶνος τούτου; οὐχὶ ἐμώρανεν ὁ Θεὸς τὴν σοφίαν τοῦ κόσμου;
\vs{21}ἐπειδὴ γὰρ ἐν τῇ σοφίᾳ τοῦ Θεοῦ οὐκ ἔγνω ὁ κόσμος διὰ τῆς σοφίας τὸν Θεόν, εὐδόκησεν ὁ Θεὸς διὰ τῆς μωρίας τοῦ κηρύγματος σῶσαι τοὺς πιστεύοντας·
\vs{22}Ἐπειδὴ καὶ Ἰουδαῖοι σημεῖα αἰτοῦσιν καὶ Ἕλληνες σοφίαν ζητοῦσιν,
\vs{23}ἡμεῖς δὲ κηρύσσομεν Χριστὸν ἐσταυρωμένον, Ἰουδαίοις μὲν σκάνδαλον, ἔθνεσιν δὲ μωρίαν,
\vs{24}αὐτοῖς δὲ τοῖς κλητοῖς, Ἰουδαίοις τε καὶ Ἕλλησιν, Χριστὸν Θεοῦ δύναμιν καὶ Θεοῦ σοφίαν·
\vs{25}Ὅτι τὸ μωρὸν τοῦ Θεοῦ σοφώτερον τῶν ἀνθρώπων ἐστίν καὶ τὸ ἀσθενὲς τοῦ Θεοῦ ἰσχυρότερον τῶν ἀνθρώπων.

\vs{26}Βλέπετε γὰρ τὴν κλῆσιν ὑμῶν, ἀδελφοί, ὅτι οὐ πολλοὶ σοφοὶ κατὰ σάρκα, οὐ πολλοὶ δυνατοί, οὐ πολλοὶ εὐγενεῖς·
\vs{27}ἀλλὰ τὰ μωρὰ τοῦ κόσμου ἐξελέξατο ὁ Θεός, ἵνα καταισχύνῃ τοὺς σοφούς, καὶ τὰ ἀσθενῆ τοῦ κόσμου ἐξελέξατο ὁ Θεός, ἵνα καταισχύνῃ τὰ ἰσχυρά,
\vs{28}καὶ τὰ ἀγενῆ τοῦ κόσμου καὶ τὰ ἐξουθενημένα ἐξελέξατο ὁ Θεός, τὰ μὴ ὄντα, ἵνα τὰ ὄντα καταργήσῃ,
\vs{29}ὅπως μὴ καυχήσηται πᾶσα σὰρξ ἐνώπιον τοῦ Θεοῦ.
\vs{30}Ἐξ αὐτοῦ δὲ ὑμεῖς ἐστε ἐν Χριστῷ Ἰησοῦ, ὃς ἐγενήθη σοφία ἡμῖν ἀπὸ Θεοῦ, δικαιοσύνη τε καὶ ἁγιασμὸς καὶ ἀπολύτρωσις,
\vs{31}ἵνα καθὼς γέγραπται· Ὁ καυχώμενος ἐν Κυρίῳ καυχάσθω.

\ch{2}
Κἀγὼ ἐλθὼν πρὸς ὑμᾶς, ἀδελφοί, ἦλθον οὐ καθ᾽ ὑπεροχὴν λόγου ἢ σοφίας καταγγέλλων ὑμῖν τὸ μυστήριον τοῦ Θεοῦ.
\vs{2}οὐ γὰρ ἔκρινά τι εἰδέναι ἐν ὑμῖν εἰ μὴ Ἰησοῦν Χριστὸν καὶ τοῦτον ἐσταυρωμένον.
\vs{3}κἀγὼ ἐν ἀσθενείᾳ καὶ ἐν φόβῳ καὶ ἐν τρόμῳ πολλῷ ἐγενόμην πρὸς ὑμᾶς,
\vs{4}καὶ ὁ λόγος μου καὶ τὸ κήρυγμά μου οὐκ ἐν πειθοῖς σοφίας λόγοις ἀλλ᾽ ἐν ἀποδείξει Πνεύματος καὶ δυνάμεως,
\vs{5}ἵνα ἡ πίστις ὑμῶν μὴ ᾖ ἐν σοφίᾳ ἀνθρώπων ἀλλ᾽ ἐν δυνάμει Θεοῦ.

\vs{6}Σοφίαν δὲ λαλοῦμεν ἐν τοῖς τελείοις, σοφίαν δὲ οὐ τοῦ αἰῶνος τούτου οὐδὲ τῶν ἀρχόντων τοῦ αἰῶνος τούτου τῶν καταργουμένων·
\vs{7}ἀλλὰ λαλοῦμεν Θεοῦ σοφίαν ἐν μυστηρίῳ τὴν ἀποκεκρυμμένην, ἣν προώρισεν ὁ Θεὸς πρὸ τῶν αἰώνων εἰς δόξαν ἡμῶν,
\vs{8}ἣν οὐδεὶς τῶν ἀρχόντων τοῦ αἰῶνος τούτου ἔγνωκεν· εἰ γὰρ ἔγνωσαν, οὐκ ἂν τὸν Κύριον τῆς δόξης ἐσταύρωσαν.
\vs{9}ἀλλὰ καθὼς γέγραπται· 
\begin{poetryblock}

\begin{quote}Ἃ ὀφθαλμὸς οὐκ εἶδεν καὶ οὖς οὐκ ἤκουσεν\end{quote} 

\begin{quote}καὶ ἐπὶ καρδίαν ἀνθρώπου οὐκ ἀνέβη,\end{quote} 

\begin{quote}ἃ ἡτοίμασεν ὁ Θεὸς τοῖς ἀγαπῶσιν αὐτόν.\end{quote}
\end{poetryblock}

\vs{10}Ἡμῖν γὰρ ἀπεκάλυψεν ὁ Θεὸς διὰ τοῦ Πνεύματος· Τὸ γὰρ Πνεῦμα πάντα ἐραυνᾷ, καὶ τὰ βάθη τοῦ Θεοῦ.
\vs{11}τίς γὰρ οἶδεν ἀνθρώπων τὰ τοῦ ἀνθρώπου εἰ μὴ τὸ πνεῦμα τοῦ ἀνθρώπου τὸ ἐν αὐτῷ; οὕτως καὶ τὰ τοῦ Θεοῦ οὐδεὶς ἔγνωκεν εἰ μὴ τὸ Πνεῦμα τοῦ Θεοῦ.
\vs{12}ἡμεῖς δὲ οὐ τὸ πνεῦμα τοῦ κόσμου ἐλάβομεν ἀλλὰ τὸ πνεῦμα τὸ ἐκ τοῦ Θεοῦ, ἵνα εἰδῶμεν τὰ ὑπὸ τοῦ Θεοῦ χαρισθέντα ἡμῖν·
\vs{13}ἃ καὶ λαλοῦμεν οὐκ ἐν διδακτοῖς ἀνθρωπίνης σοφίας λόγοις ἀλλ᾽ ἐν διδακτοῖς Πνεύματος, πνευματικοῖς πνευματικὰ συνκρίνοντες.
\vs{14}Ψυχικὸς δὲ ἄνθρωπος οὐ δέχεται τὰ τοῦ Πνεύματος τοῦ Θεοῦ· μωρία γὰρ αὐτῷ ἐστίν καὶ οὐ δύναται γνῶναι, ὅτι πνευματικῶς ἀνακρίνεται.
\vs{15}ὁ δὲ πνευματικὸς ἀνακρίνει τὰ πάντα, αὐτὸς δὲ ὑπ᾽ οὐδενὸς ἀνακρίνεται.
\vs{16}Τίς γὰρ ἔγνω νοῦν Κυρίου, ὃς συμβιβάσει αὐτόν; ἡμεῖς δὲ νοῦν Χριστοῦ ἔχομεν.

\ch{3}
Κἀγώ, ἀδελφοί, οὐκ ἠδυνήθην λαλῆσαι ὑμῖν ὡς πνευματικοῖς ἀλλ᾽ ὡς σαρκίνοις, ὡς νηπίοις ἐν Χριστῷ.
\vs{2}γάλα ὑμᾶς ἐπότισα, οὐ βρῶμα· οὔπω γὰρ ἐδύνασθε. Ἀλλ᾽ οὐδὲ ἔτι νῦν δύνασθε,
\vs{3}ἔτι γὰρ σαρκικοί ἐστε. ὅπου γὰρ ἐν ὑμῖν ζῆλος καὶ ἔρις, οὐχὶ σαρκικοί ἐστε καὶ κατὰ ἄνθρωπον περιπατεῖτε;
\vs{4}ὅταν γὰρ λέγῃ τις· Ἐγὼ μέν εἰμι Παύλου, ἕτερος δέ· Ἐγὼ Ἀπολλῶ, οὐκ ἄνθρωποί ἐστε;
\vs{5}Τί οὖν ἐστιν Ἀπολλῶς; τί δέ ἐστιν Παῦλος; διάκονοι δι᾽ ὧν ἐπιστεύσατε, καὶ ἑκάστῳ ὡς ὁ Κύριος ἔδωκεν.
\vs{6}ἐγὼ ἐφύτευσα, Ἀπολλῶς ἐπότισεν, ἀλλὰ ὁ Θεὸς ηὔξανεν·
\vs{7}ὥστε οὔτε ὁ φυτεύων ἐστίν τι οὔτε ὁ ποτίζων ἀλλ᾽ ὁ αὐξάνων Θεός.
\vs{8}ὁ φυτεύων δὲ καὶ ὁ ποτίζων ἕν εἰσιν, ἕκαστος δὲ τὸν ἴδιον μισθὸν λήμψεται κατὰ τὸν ἴδιον κόπον·
\vs{9}Θεοῦ γάρ ἐσμεν συνεργοί, Θεοῦ γεώργιον, Θεοῦ οἰκοδομή ἐστε.
\vs{10}Κατὰ τὴν χάριν τοῦ Θεοῦ τὴν δοθεῖσάν μοι ὡς σοφὸς ἀρχιτέκτων θεμέλιον ἔθηκα, ἄλλος δὲ ἐποικοδομεῖ. ἕκαστος δὲ βλεπέτω πῶς ἐποικοδομεῖ.
\vs{11}θεμέλιον γὰρ ἄλλον οὐδεὶς δύναται θεῖναι παρὰ τὸν κείμενον, ὅς ἐστιν Ἰησοῦς Χριστός.
\vs{12}Εἰ δέ τις ἐποικοδομεῖ ἐπὶ τὸν θεμέλιον χρυσόν, ἄργυρον, λίθους τιμίους, ξύλα, χόρτον, καλάμην,
\vs{13}ἑκάστου τὸ ἔργον φανερὸν γενήσεται, ἡ γὰρ ἡμέρα δηλώσει, ὅτι ἐν πυρὶ ἀποκαλύπτεται· καὶ ἑκάστου τὸ ἔργον ὁποῖόν ἐστιν τὸ πῦρ αὐτὸ δοκιμάσει.
\vs{14}εἴ τινος τὸ ἔργον μενεῖ ὃ ἐποικοδόμησεν, μισθὸν λήμψεται·
\vs{15}εἴ τινος τὸ ἔργον κατακαήσεται, ζημιωθήσεται, αὐτὸς δὲ σωθήσεται, οὕτως δὲ ὡς διὰ πυρός.
\vs{16}Οὐκ οἴδατε ὅτι ναὸς Θεοῦ ἐστε καὶ τὸ Πνεῦμα τοῦ Θεοῦ οἰκεῖ ἐν ὑμῖν;
\vs{17}εἴ τις τὸν ναὸν τοῦ Θεοῦ φθείρει, φθερεῖ τοῦτον ὁ Θεός· ὁ γὰρ ναὸς τοῦ Θεοῦ ἅγιός ἐστιν, οἵτινές ἐστε ὑμεῖς.

\vs{18}Μηδεὶς ἑαυτὸν ἐξαπατάτω· εἴ τις δοκεῖ σοφὸς εἶναι ἐν ὑμῖν ἐν τῷ αἰῶνι τούτῳ, μωρὸς γενέσθω, ἵνα γένηται σοφός.
\vs{19}ἡ γὰρ σοφία τοῦ κόσμου τούτου μωρία παρὰ τῷ Θεῷ ἐστιν. γέγραπται γάρ· Ὁ δρασσόμενος τοὺς σοφοὺς ἐν τῇ πανουργίᾳ αὐτῶν·
\vs{20}καὶ πάλιν· Κύριος γινώσκει τοὺς διαλογισμοὺς τῶν σοφῶν ὅτι εἰσὶν μάταιοι.
\vs{21}Ὥστε μηδεὶς καυχάσθω ἐν ἀνθρώποις· πάντα γὰρ ὑμῶν ἐστιν,
\vs{22}εἴτε Παῦλος εἴτε Ἀπολλῶς εἴτε Κηφᾶς, εἴτε κόσμος εἴτε ζωὴ εἴτε θάνατος, εἴτε ἐνεστῶτα εἴτε μέλλοντα· πάντα ὑμῶν,
\vs{23}ὑμεῖς δὲ Χριστοῦ, Χριστὸς δὲ Θεοῦ.

\ch{4}
Οὕτως ἡμᾶς λογιζέσθω ἄνθρωπος ὡς ὑπηρέτας Χριστοῦ καὶ οἰκονόμους μυστηρίων Θεοῦ.
\vs{2}ὧδε λοιπὸν ζητεῖται ἐν τοῖς οἰκονόμοις, ἵνα πιστός τις εὑρεθῇ.
\vs{3}Ἐμοὶ δὲ εἰς ἐλάχιστόν ἐστιν, ἵνα ὑφ᾽ ὑμῶν ἀνακριθῶ ἢ ὑπὸ ἀνθρωπίνης ἡμέρας· ἀλλ᾽ οὐδὲ ἐμαυτὸν ἀνακρίνω.
\vs{4}οὐδὲν γὰρ ἐμαυτῷ σύνοιδα, ἀλλ᾽ οὐκ ἐν τούτῳ δεδικαίωμαι, ὁ δὲ ἀνακρίνων με Κύριός ἐστιν.
\vs{5}Ὥστε μὴ πρὸ καιροῦ τι κρίνετε ἕως ἂν ἔλθῃ ὁ Κύριος, ὃς καὶ φωτίσει τὰ κρυπτὰ τοῦ σκότους καὶ φανερώσει τὰς βουλὰς τῶν καρδιῶν· καὶ τότε ὁ ἔπαινος γενήσεται ἑκάστῳ ἀπὸ τοῦ Θεοῦ.

\vs{6}Ταῦτα δέ, ἀδελφοί, μετεσχημάτισα εἰς ἐμαυτὸν καὶ Ἀπολλῶν δι᾽ ὑμᾶς, ἵνα ἐν ἡμῖν μάθητε τό Μὴ ὑπὲρ ἃ γέγραπται, ἵνα μὴ εἷς ὑπὲρ τοῦ ἑνὸς φυσιοῦσθε κατὰ τοῦ ἑτέρου.
\vs{7}τίς γάρ σε διακρίνει; τί δὲ ἔχεις ὃ οὐκ ἔλαβες; εἰ δὲ καὶ ἔλαβες, τί καυχᾶσαι ὡς μὴ λαβών;
\vs{8}ἤδη κεκορεσμένοι ἐστέ, ἤδη ἐπλουτήσατε, χωρὶς ἡμῶν ἐβασιλεύσατε· καὶ ὄφελόν γε ἐβασιλεύσατε, ἵνα καὶ ἡμεῖς ὑμῖν συμβασιλεύσωμεν.
\vs{9}δοκῶ γάρ, ὁ Θεὸς ἡμᾶς τοὺς ἀποστόλους ἐσχάτους ἀπέδειξεν ὡς ἐπιθανατίους, ὅτι θέατρον ἐγενήθημεν τῷ κόσμῳ καὶ ἀγγέλοις καὶ ἀνθρώποις.
\vs{10}Ἡμεῖς μωροὶ διὰ Χριστόν, ὑμεῖς δὲ φρόνιμοι ἐν Χριστῷ· ἡμεῖς ἀσθενεῖς, ὑμεῖς δὲ ἰσχυροί· ὑμεῖς ἔνδοξοι, ἡμεῖς δὲ ἄτιμοι.
\vs{11}ἄχρι τῆς ἄρτι ὥρας καὶ πεινῶμεν καὶ διψῶμεν καὶ γυμνιτεύομεν καὶ κολαφιζόμεθα καὶ ἀστατοῦμεν
\vs{12}καὶ κοπιῶμεν ἐργαζόμενοι ταῖς ἰδίαις χερσίν· λοιδορούμενοι εὐλογοῦμεν, διωκόμενοι ἀνεχόμεθα,
\vs{13}δυσφημούμενοι παρακαλοῦμεν· ὡς περικαθάρματα τοῦ κόσμου ἐγενήθημεν, πάντων περίψημα ἕως ἄρτι.

\vs{14}Οὐκ ἐντρέπων ὑμᾶς γράφω ταῦτα ἀλλ᾽ ὡς τέκνα μου ἀγαπητὰ νουθετῶν.
\vs{15}ἐὰν γὰρ μυρίους παιδαγωγοὺς ἔχητε ἐν Χριστῷ ἀλλ᾽ οὐ πολλοὺς πατέρας· ἐν γὰρ Χριστῷ Ἰησοῦ διὰ τοῦ εὐαγγελίου ἐγὼ ὑμᾶς ἐγέννησα.
\vs{16}παρακαλῶ οὖν ὑμᾶς, μιμηταί μου γίνεσθε.
\vs{17}Διὰ τοῦτο ἔπεμψα ὑμῖν Τιμόθεον, ὅς ἐστίν μου τέκνον ἀγαπητὸν καὶ πιστὸν ἐν Κυρίῳ, ὃς ὑμᾶς ἀναμνήσει τὰς ὁδούς μου τὰς ἐν Χριστῷ Ἰησοῦ, καθὼς πανταχοῦ ἐν πάσῃ ἐκκλησίᾳ διδάσκω.
\vs{18}Ὡς μὴ ἐρχομένου δέ μου πρὸς ὑμᾶς ἐφυσιώθησάν τινες·
\vs{19}ἐλεύσομαι δὲ ταχέως πρὸς ὑμᾶς ἐὰν ὁ Κύριος θελήσῃ, καὶ γνώσομαι οὐ τὸν λόγον τῶν πεφυσιωμένων ἀλλὰ τὴν δύναμιν·
\vs{20}οὐ γὰρ ἐν λόγῳ ἡ βασιλεία τοῦ Θεοῦ ἀλλ᾽ ἐν δυνάμει.
\vs{21}τί θέλετε; ἐν ῥάβδῳ ἔλθω πρὸς ὑμᾶς ἢ ἐν ἀγάπῃ πνεύματί τε πραΰτητος;

\ch{5}
Ὅλως ἀκούεται ἐν ὑμῖν πορνεία, καὶ τοιαύτη πορνεία ἥτις οὐδὲ ἐν τοῖς ἔθνεσιν, ὥστε γυναῖκά τινα τοῦ πατρὸς ἔχειν.
\vs{2}καὶ ὑμεῖς πεφυσιωμένοι ἐστέ καὶ οὐχὶ μᾶλλον ἐπενθήσατε, ἵνα ἀρθῇ ἐκ μέσου ὑμῶν ὁ τὸ ἔργον τοῦτο πράξας;
\vs{3}Ἐγὼ μὲν γάρ, ἀπὼν τῷ σώματι παρὼν δὲ τῷ πνεύματι, ἤδη κέκρικα ὡς παρὼν τὸν οὕτως τοῦτο κατεργασάμενον·
\vs{4}ἐν τῷ ὀνόματι τοῦ Κυρίου ἡμῶν Ἰησοῦ συναχθέντων ὑμῶν καὶ τοῦ ἐμοῦ πνεύματος σὺν τῇ δυνάμει τοῦ Κυρίου ἡμῶν Ἰησοῦ,
\vs{5}παραδοῦναι τὸν τοιοῦτον τῷ Σατανᾷ εἰς ὄλεθρον τῆς σαρκός, ἵνα τὸ πνεῦμα σωθῇ ἐν τῇ ἡμέρᾳ τοῦ Κυρίου.
\vs{6}Οὐ καλὸν τὸ καύχημα ὑμῶν. οὐκ οἴδατε ὅτι μικρὰ ζύμη ὅλον τὸ φύραμα ζυμοῖ;
\vs{7}ἐκκαθάρατε τὴν παλαιὰν ζύμην, ἵνα ἦτε νέον φύραμα, καθώς ἐστε ἄζυμοι· καὶ γὰρ τὸ πάσχα ἡμῶν ἐτύθη Χριστός.
\vs{8}ὥστε ἑορτάζωμεν μὴ ἐν ζύμῃ παλαιᾷ μηδὲ ἐν ζύμῃ κακίας καὶ πονηρίας ἀλλ᾽ ἐν ἀζύμοις εἰλικρινείας καὶ ἀληθείας.
\vs{9}Ἔγραψα ὑμῖν ἐν τῇ ἐπιστολῇ μὴ συναναμίγνυσθαι πόρνοις,
\vs{10}οὐ πάντως τοῖς πόρνοις τοῦ κόσμου τούτου ἢ τοῖς πλεονέκταις καὶ ἅρπαξιν ἢ εἰδωλολάτραις, ἐπεὶ ὠφείλετε ἄρα ἐκ τοῦ κόσμου ἐξελθεῖν.
\vs{11}νῦν δὲ ἔγραψα ὑμῖν μὴ συναναμίγνυσθαι ἐάν τις ἀδελφὸς ὀνομαζόμενος ᾖ πόρνος ἢ πλεονέκτης ἢ εἰδωλολάτρης ἢ λοίδορος ἢ μέθυσος ἢ ἅρπαξ, τῷ τοιούτῳ μηδὲ συνεσθίειν.
\vs{12}Τί γάρ μοι τοὺς ἔξω κρίνειν; οὐχὶ τοὺς ἔσω ὑμεῖς κρίνετε;
\vs{13}τοὺς δὲ ἔξω ὁ Θεὸς κρίνει. Ἐξάρατε τὸν πονηρὸν ἐξ ὑμῶν αὐτῶν.

\ch{6}
Τολμᾷ τις ὑμῶν πρᾶγμα ἔχων πρὸς τὸν ἕτερον κρίνεσθαι ἐπὶ τῶν ἀδίκων καὶ οὐχὶ ἐπὶ τῶν ἁγίων;
\vs{2}ἢ οὐκ οἴδατε ὅτι οἱ ἅγιοι τὸν κόσμον κρινοῦσιν; καὶ εἰ ἐν ὑμῖν κρίνεται ὁ κόσμος, ἀνάξιοί ἐστε κριτηρίων ἐλαχίστων;
\vs{3}οὐκ οἴδατε ὅτι ἀγγέλους κρινοῦμεν, μήτι γε βιωτικά;
\vs{4}Βιωτικὰ μὲν οὖν κριτήρια ἐὰν ἔχητε, τοὺς ἐξουθενημένους ἐν τῇ ἐκκλησίᾳ, τούτους καθίζετε;
\vs{5}πρὸς ἐντροπὴν ὑμῖν λέγω. οὕτως οὐκ ἔνι ἐν ὑμῖν οὐδεὶς σοφὸς, ὃς δυνήσεται διακρῖναι ἀνὰ μέσον τοῦ ἀδελφοῦ αὐτοῦ;
\vs{6}ἀλλὰ ἀδελφὸς μετὰ ἀδελφοῦ κρίνεται καὶ τοῦτο ἐπὶ ἀπίστων;
\vs{7}ἤδη μὲν οὖν ὅλως ἥττημα ὑμῖν ἐστιν ὅτι κρίματα ἔχετε μεθ᾽ ἑαυτῶν. διὰ τί οὐχὶ μᾶλλον ἀδικεῖσθε; διὰ τί οὐχὶ μᾶλλον ἀποστερεῖσθε;
\vs{8}ἀλλὰ ὑμεῖς ἀδικεῖτε καὶ ἀποστερεῖτε, καὶ τοῦτο ἀδελφούς.
\vs{9}Ἢ οὐκ οἴδατε ὅτι ἄδικοι Θεοῦ βασιλείαν οὐ κληρονομήσουσιν; μὴ πλανᾶσθε· οὔτε πόρνοι οὔτε εἰδωλολάτραι οὔτε μοιχοὶ οὔτε μαλακοὶ οὔτε ἀρσενοκοῖται
\vs{10}οὔτε κλέπται οὔτε πλεονέκται, οὐ μέθυσοι, οὐ λοίδοροι, οὐχ ἅρπαγες βασιλείαν Θεοῦ κληρονομήσουσιν.
\vs{11}καὶ ταῦτά τινες ἦτε· ἀλλὰ ἀπελούσασθε, ἀλλὰ ἡγιάσθητε, ἀλλὰ ἐδικαιώθητε ἐν τῷ ὀνόματι τοῦ Κυρίου Ἰησοῦ Χριστοῦ καὶ ἐν τῷ Πνεύματι τοῦ Θεοῦ ἡμῶν.

\vs{12}Πάντα μοι ἔξεστιν ἀλλ᾽ οὐ πάντα συμφέρει· Πάντα μοι ἔξεστιν ἀλλ᾽ οὐκ ἐγὼ ἐξουσιασθήσομαι ὑπό τινος.
\vs{13}Τὰ βρώματα τῇ κοιλίᾳ καὶ ἡ κοιλία τοῖς βρώμασιν, ὁ δὲ Θεὸς καὶ ταύτην καὶ ταῦτα καταργήσει. τὸ δὲ σῶμα οὐ τῇ πορνείᾳ ἀλλὰ τῷ Κυρίῳ, καὶ ὁ Κύριος τῷ σώματι·
\vs{14}ὁ δὲ Θεὸς καὶ τὸν Κύριον ἤγειρεν καὶ ἡμᾶς ἐξεγερεῖ διὰ τῆς δυνάμεως αὐτοῦ.
\vs{15}Οὐκ οἴδατε ὅτι τὰ σώματα ὑμῶν μέλη Χριστοῦ ἐστιν; ἄρας οὖν τὰ μέλη τοῦ Χριστοῦ ποιήσω πόρνης μέλη; μὴ γένοιτο.
\vs{16}ἢ οὐκ οἴδατε ὅτι ὁ κολλώμενος τῇ πόρνῃ ἓν σῶμά ἐστιν; Ἔσονται γάρ, φησίν, Οἱ δύο εἰς σάρκα μίαν.
\vs{17}ὁ δὲ κολλώμενος τῷ Κυρίῳ ἓν πνεῦμά ἐστιν.
\vs{18}Φεύγετε τὴν πορνείαν. πᾶν ἁμάρτημα ὃ ἐὰν ποιήσῃ ἄνθρωπος ἐκτὸς τοῦ σώματός ἐστιν· ὁ δὲ πορνεύων εἰς τὸ ἴδιον σῶμα ἁμαρτάνει.
\vs{19}ἢ οὐκ οἴδατε ὅτι τὸ σῶμα ὑμῶν ναὸς τοῦ ἐν ὑμῖν Ἁγίου Πνεύματός ἐστιν οὗ ἔχετε ἀπὸ Θεοῦ, καὶ οὐκ ἐστὲ ἑαυτῶν;
\vs{20}ἠγοράσθητε γὰρ τιμῆς· δοξάσατε δὴ τὸν Θεὸν ἐν τῷ σώματι ὑμῶν.

\ch{7}
Περὶ δὲ ὧν ἐγράψατε, καλὸν ἀνθρώπῳ γυναικὸς μὴ ἅπτεσθαι·
\vs{2}διὰ δὲ τὰς πορνείας ἕκαστος τὴν ἑαυτοῦ γυναῖκα ἐχέτω καὶ ἑκάστη τὸν ἴδιον ἄνδρα ἐχέτω.
\vs{3}Τῇ γυναικὶ ὁ ἀνὴρ τὴν ὀφειλὴν ἀποδιδότω, ὁμοίως δὲ καὶ ἡ γυνὴ τῷ ἀνδρί.
\vs{4}ἡ γυνὴ τοῦ ἰδίου σώματος οὐκ ἐξουσιάζει ἀλλὰ ὁ ἀνήρ, ὁμοίως δὲ καὶ ὁ ἀνὴρ τοῦ ἰδίου σώματος οὐκ ἐξουσιάζει ἀλλὰ ἡ γυνή.
\vs{5}Μὴ ἀποστερεῖτε ἀλλήλους, εἰ μήτι ἂν ἐκ συμφώνου πρὸς καιρὸν, ἵνα σχολάσητε τῇ προσευχῇ καὶ πάλιν ἐπὶ τὸ αὐτὸ ἦτε, ἵνα μὴ πειράζῃ ὑμᾶς ὁ Σατανᾶς διὰ τὴν ἀκρασίαν ὑμῶν.
\vs{6}τοῦτο δὲ λέγω κατὰ συνγνώμην οὐ κατ᾽ ἐπιταγήν.
\vs{7}θέλω δὲ πάντας ἀνθρώπους εἶναι ὡς καὶ ἐμαυτόν· ἀλλὰ ἕκαστος ἴδιον ἔχει χάρισμα ἐκ Θεοῦ, ὁ μὲν οὕτως, ὁ δὲ οὕτως.

\vs{8}Λέγω δὲ τοῖς ἀγάμοις καὶ ταῖς χήραις, καλὸν αὐτοῖς ἐὰν μείνωσιν ὡς κἀγώ·
\vs{9}εἰ δὲ οὐκ ἐγκρατεύονται, γαμησάτωσαν, κρεῖττον γάρ ἐστιν γαμῆσαι ἢ πυροῦσθαι.
\vs{10}Τοῖς δὲ γεγαμηκόσιν παραγγέλλω, οὐκ ἐγὼ ἀλλὰ ὁ Κύριος, γυναῖκα ἀπὸ ἀνδρὸς μὴ χωρισθῆναι,—
\vs{11}ἐὰν δὲ καὶ χωρισθῇ, μενέτω ἄγαμος ἢ τῷ ἀνδρὶ καταλλαγήτω,— καὶ ἄνδρα γυναῖκα μὴ ἀφιέναι.
\vs{12}Τοῖς δὲ λοιποῖς λέγω ἐγώ οὐχ ὁ Κύριος· εἴ τις ἀδελφὸς γυναῖκα ἔχει ἄπιστον καὶ αὕτη συνευδοκεῖ οἰκεῖν μετ᾽ αὐτοῦ, μὴ ἀφιέτω αὐτήν·
\vs{13}καὶ γυνὴ εἴ τις ἔχει ἄνδρα ἄπιστον καὶ οὗτος συνευδοκεῖ οἰκεῖν μετ᾽ αὐτῆς, μὴ ἀφιέτω τὸν ἄνδρα.
\vs{14}ἡγίασται γὰρ ὁ ἀνὴρ ὁ ἄπιστος ἐν τῇ γυναικί καὶ ἡγίασται ἡ γυνὴ ἡ ἄπιστος ἐν τῷ ἀδελφῷ· ἐπεὶ ἄρα τὰ τέκνα ὑμῶν ἀκάθαρτά ἐστιν, νῦν δὲ ἅγιά ἐστιν.
\vs{15}Εἰ δὲ ὁ ἄπιστος χωρίζεται, χωριζέσθω· οὐ δεδούλωται ὁ ἀδελφὸς ἢ ἡ ἀδελφὴ ἐν τοῖς τοιούτοις· ἐν δὲ εἰρήνῃ κέκληκεν ὑμᾶς ὁ Θεός.
\vs{16}τί γὰρ οἶδας, γύναι, εἰ τὸν ἄνδρα σώσεις; ἢ τί οἶδας, ἄνερ, εἰ τὴν γυναῖκα σώσεις;

\vs{17}Εἰ μὴ ἑκάστῳ ὡς ἐμέρισεν ὁ Κύριος, ἕκαστον ὡς κέκληκεν ὁ Θεός, οὕτως περιπατείτω. καὶ οὕτως ἐν ταῖς ἐκκλησίαις πάσαις διατάσσομαι.
\vs{18}περιτετμημένος τις ἐκλήθη, μὴ ἐπισπάσθω· ἐν ἀκροβυστίᾳ κέκληταί τις, μὴ περιτεμνέσθω.
\vs{19}ἡ περιτομὴ οὐδέν ἐστιν καὶ ἡ ἀκροβυστία οὐδέν ἐστιν, ἀλλὰ τήρησις ἐντολῶν Θεοῦ.
\vs{20}Ἕκαστος ἐν τῇ κλήσει ᾗ ἐκλήθη, ἐν ταύτῃ μενέτω.
\vs{21}δοῦλος ἐκλήθης, μή σοι μελέτω· ἀλλ᾽ εἰ καὶ δύνασαι ἐλεύθερος γενέσθαι, μᾶλλον χρῆσαι.
\vs{22}ὁ γὰρ ἐν Κυρίῳ κληθεὶς δοῦλος ἀπελεύθερος Κυρίου ἐστίν, ὁμοίως ὁ ἐλεύθερος κληθεὶς δοῦλός ἐστιν Χριστοῦ.
\vs{23}Τιμῆς ἠγοράσθητε· μὴ γίνεσθε δοῦλοι ἀνθρώπων.
\vs{24}ἕκαστος ἐν ᾧ ἐκλήθη, ἀδελφοί, ἐν τούτῳ μενέτω παρὰ Θεῷ.

\vs{25}Περὶ δὲ τῶν παρθένων ἐπιταγὴν Κυρίου οὐκ ἔχω, γνώμην δὲ δίδωμι ὡς ἠλεημένος ὑπὸ Κυρίου πιστὸς εἶναι.
\vs{26}Νομίζω οὖν τοῦτο καλὸν ὑπάρχειν διὰ τὴν ἐνεστῶσαν ἀνάγκην, ὅτι καλὸν ἀνθρώπῳ τὸ οὕτως εἶναι.
\vs{27}δέδεσαι γυναικί, μὴ ζήτει λύσιν· λέλυσαι ἀπὸ γυναικός, μὴ ζήτει γυναῖκα.
\vs{28}ἐὰν δὲ καὶ γαμήσῃς, οὐχ ἥμαρτες, καὶ ἐὰν γήμῃ ἡ παρθένος, οὐχ ἥμαρτεν· θλῖψιν δὲ τῇ σαρκὶ ἕξουσιν οἱ τοιοῦτοι, ἐγὼ δὲ ὑμῶν φείδομαι.
\vs{29}Τοῦτο δέ φημι, ἀδελφοί, ὁ καιρὸς συνεσταλμένος ἐστίν· τὸ λοιπὸν, ἵνα καὶ οἱ ἔχοντες γυναῖκας ὡς μὴ ἔχοντες ὦσιν
\vs{30}καὶ οἱ κλαίοντες ὡς μὴ κλαίοντες καὶ οἱ χαίροντες ὡς μὴ χαίροντες καὶ οἱ ἀγοράζοντες ὡς μὴ κατέχοντες,
\vs{31}καὶ οἱ χρώμενοι τὸν κόσμον ὡς μὴ καταχρώμενοι· παράγει γὰρ τὸ σχῆμα τοῦ κόσμου τούτου.
\vs{32}Θέλω δὲ ὑμᾶς ἀμερίμνους εἶναι. ὁ ἄγαμος μεριμνᾷ τὰ τοῦ Κυρίου, πῶς ἀρέσῃ τῷ Κυρίῳ·
\vs{33}ὁ δὲ γαμήσας μεριμνᾷ τὰ τοῦ κόσμου, πῶς ἀρέσῃ τῇ γυναικί,
\vs{34}καὶ μεμέρισται. καὶ ἡ γυνὴ ἡ ἄγαμος καὶ ἡ παρθένος μεριμνᾷ τὰ τοῦ Κυρίου, ἵνα ᾖ ἁγία καὶ τῷ σώματι καὶ τῷ πνεύματι· ἡ δὲ γαμήσασα μεριμνᾷ τὰ τοῦ κόσμου, πῶς ἀρέσῃ τῷ ἀνδρί.
\vs{35}Τοῦτο δὲ πρὸς τὸ ὑμῶν αὐτῶν σύμφορον λέγω, οὐχ ἵνα βρόχον ὑμῖν ἐπιβάλω ἀλλὰ πρὸς τὸ εὔσχημον καὶ εὐπάρεδρον τῷ Κυρίῳ ἀπερισπάστως.
\vs{36}Εἰ δέ τις ἀσχημονεῖν ἐπὶ τὴν παρθένον αὐτοῦ νομίζει, ἐὰν ᾖ ὑπέρακμος καὶ οὕτως ὀφείλει γίνεσθαι, ὃ θέλει ποιείτω, οὐχ ἁμαρτάνει, γαμείτωσαν.
\vs{37}ὃς δὲ ἕστηκεν ἐν τῇ καρδίᾳ αὐτοῦ ἑδραῖος μὴ ἔχων ἀνάγκην, ἐξουσίαν δὲ ἔχει περὶ τοῦ ἰδίου θελήματος καὶ τοῦτο κέκρικεν ἐν τῇ ἰδίᾳ καρδίᾳ, τηρεῖν τὴν ἑαυτοῦ παρθένον, καλῶς ποιήσει.
\vs{38}Ὥστε καὶ ὁ γαμίζων τὴν ἑαυτοῦ παρθένον καλῶς ποιεῖ καὶ ὁ μὴ γαμίζων κρεῖσσον ποιήσει.
\vs{39}Γυνὴ δέδεται ἐφ᾽ ὅσον χρόνον ζῇ ὁ ἀνὴρ αὐτῆς· ἐὰν δὲ κοιμηθῇ ὁ ἀνήρ, ἐλευθέρα ἐστὶν ᾧ θέλει γαμηθῆναι, μόνον ἐν Κυρίῳ.
\vs{40}μακαριωτέρα δέ ἐστιν ἐὰν οὕτως μείνῃ, κατὰ τὴν ἐμὴν γνώμην· δοκῶ δὲ κἀγὼ Πνεῦμα Θεοῦ ἔχειν.

\ch{8}
Περὶ δὲ τῶν εἰδωλοθύτων, οἴδαμεν ὅτι πάντες γνῶσιν ἔχομεν. ἡ γνῶσις φυσιοῖ, ἡ δὲ ἀγάπη οἰκοδομεῖ·
\vs{2}εἴ τις δοκεῖ ἐγνωκέναι τι, οὔπω ἔγνω καθὼς δεῖ γνῶναι·
\vs{3}εἰ δέ τις ἀγαπᾷ τὸν Θεόν, οὗτος ἔγνωσται ὑπ᾽ αὐτοῦ.
\vs{4}Περὶ τῆς βρώσεως οὖν τῶν εἰδωλοθύτων, οἴδαμεν ὅτι οὐδὲν εἴδωλον ἐν κόσμῳ καὶ ὅτι οὐδεὶς Θεὸς εἰ μὴ εἷς.
\vs{5}καὶ γὰρ εἴπερ εἰσὶν λεγόμενοι θεοὶ εἴτε ἐν οὐρανῷ εἴτε ἐπὶ γῆς, ὥσπερ εἰσὶν θεοὶ πολλοὶ καὶ κύριοι πολλοί,
\begin{poetryblock}

\begin{quote} \vs{6}ἀλλ᾽ ἡμῖν εἷς Θεὸς ὁ Πατήρ\end{quote} 

\begin{quote}ἐξ οὗ τὰ πάντα καὶ ἡμεῖς εἰς αὐτόν,\end{quote} 

\begin{quote}καὶ εἷς Κύριος Ἰησοῦς Χριστός\end{quote} 

\begin{quote}δι᾽ οὗ τὰ πάντα καὶ ἡμεῖς δι᾽ αὐτοῦ.\end{quote}
\end{poetryblock}
\vs{7}Ἀλλ᾽ οὐκ ἐν πᾶσιν ἡ γνῶσις· τινὲς δὲ τῇ συνηθείᾳ ἕως ἄρτι τοῦ εἰδώλου ὡς εἰδωλόθυτον ἐσθίουσιν, καὶ ἡ συνείδησις αὐτῶν ἀσθενὴς οὖσα μολύνεται.
\vs{8}βρῶμα δὲ ἡμᾶς οὐ παραστήσει τῷ Θεῷ· οὔτε ἐὰν μὴ φάγωμεν ὑστερούμεθα, οὔτε ἐὰν φάγωμεν περισσεύομεν.
\vs{9}Βλέπετε δὲ μή πως ἡ ἐξουσία ὑμῶν αὕτη πρόσκομμα γένηται τοῖς ἀσθενέσιν.
\vs{10}ἐὰν γάρ τις ἴδῃ σὲ τὸν ἔχοντα γνῶσιν ἐν εἰδωλείῳ κατακείμενον, οὐχὶ ἡ συνείδησις αὐτοῦ ἀσθενοῦς ὄντος οἰκοδομηθήσεται εἰς τὸ τὰ εἰδωλόθυτα ἐσθίειν;
\vs{11}ἀπόλλυται γὰρ ὁ ἀσθενῶν ἐν τῇ σῇ γνώσει, ὁ ἀδελφὸς δι᾽ ὃν Χριστὸς ἀπέθανεν.
\vs{12}οὕτως δὲ ἁμαρτάνοντες εἰς τοὺς ἀδελφοὺς καὶ τύπτοντες αὐτῶν τὴν συνείδησιν ἀσθενοῦσαν εἰς Χριστὸν ἁμαρτάνετε.
\vs{13}Διόπερ εἰ βρῶμα σκανδαλίζει τὸν ἀδελφόν μου, οὐ μὴ φάγω κρέα εἰς τὸν αἰῶνα, ἵνα μὴ τὸν ἀδελφόν μου σκανδαλίσω.

\ch{9}
Οὐκ εἰμὶ ἐλεύθερος; οὐκ εἰμὶ ἀπόστολος; οὐχὶ Ἰησοῦν τὸν Κύριον ἡμῶν ἑόρακα; οὐ τὸ ἔργον μου ὑμεῖς ἐστε ἐν Κυρίῳ;
\vs{2}εἰ ἄλλοις οὐκ εἰμὶ ἀπόστολος, ἀλλά γε ὑμῖν εἰμι· ἡ γὰρ σφραγίς μου τῆς ἀποστολῆς ὑμεῖς ἐστε ἐν Κυρίῳ.
\vs{3}Ἡ ἐμὴ ἀπολογία τοῖς ἐμὲ ἀνακρίνουσίν ἐστιν αὕτη.
\vs{4}μὴ οὐκ ἔχομεν ἐξουσίαν φαγεῖν καὶ πεῖν;
\vs{5}μὴ οὐκ ἔχομεν ἐξουσίαν ἀδελφὴν γυναῖκα περιάγειν ὡς καὶ οἱ λοιποὶ ἀπόστολοι καὶ οἱ ἀδελφοὶ τοῦ Κυρίου καὶ Κηφᾶς;
\vs{6}ἢ μόνος ἐγὼ καὶ Βαρνάβας οὐκ ἔχομεν ἐξουσίαν μὴ ἐργάζεσθαι;
\vs{7}Τίς στρατεύεται ἰδίοις ὀψωνίοις ποτέ; τίς φυτεύει ἀμπελῶνα καὶ τὸν καρπὸν αὐτοῦ οὐκ ἐσθίει; ἢ τίς ποιμαίνει ποίμνην καὶ ἐκ τοῦ γάλακτος τῆς ποίμνης οὐκ ἐσθίει;
\vs{8}Μὴ κατὰ ἄνθρωπον ταῦτα λαλῶ ἢ καὶ ὁ νόμος ταῦτα οὐ λέγει;
\vs{9}ἐν γὰρ τῷ Μωϋσέως νόμῳ γέγραπται· Οὐ κημώσεις βοῦν ἀλοῶντα. μὴ τῶν βοῶν μέλει τῷ Θεῷ
\vs{10}ἢ δι᾽ ἡμᾶς πάντως λέγει; δι᾽ ἡμᾶς γὰρ ἐγράφη ὅτι ὀφείλει ἐπ᾽ ἐλπίδι ὁ ἀροτριῶν ἀροτριᾶν καὶ ὁ ἀλοῶν ἐπ᾽ ἐλπίδι τοῦ μετέχειν.
\vs{11}Εἰ ἡμεῖς ὑμῖν τὰ πνευματικὰ ἐσπείραμεν, μέγα εἰ ἡμεῖς ὑμῶν τὰ σαρκικὰ θερίσομεν;
\vs{12}εἰ ἄλλοι τῆς ὑμῶν ἐξουσίας μετέχουσιν, οὐ μᾶλλον ἡμεῖς; ἀλλ᾽ οὐκ ἐχρησάμεθα τῇ ἐξουσίᾳ ταύτῃ, ἀλλὰ πάντα στέγομεν, ἵνα μή τινα ἐνκοπὴν δῶμεν τῷ εὐαγγελίῳ τοῦ Χριστοῦ.
\vs{13}Οὐκ οἴδατε ὅτι οἱ τὰ ἱερὰ ἐργαζόμενοι τὰ ἐκ τοῦ ἱεροῦ ἐσθίουσιν, οἱ τῷ θυσιαστηρίῳ παρεδρεύοντες τῷ θυσιαστηρίῳ συμμερίζονται;
\vs{14}οὕτως καὶ ὁ Κύριος διέταξεν τοῖς τὸ εὐαγγέλιον καταγγέλλουσιν ἐκ τοῦ εὐαγγελίου ζῆν.
\vs{15}ἐγὼ δὲ οὐ κέχρημαι οὐδενὶ τούτων. οὐκ ἔγραψα δὲ ταῦτα, ἵνα οὕτως γένηται ἐν ἐμοί· καλὸν γάρ μοι μᾶλλον ἀποθανεῖν ἢ— τὸ καύχημά μου οὐδεὶς κενώσει.
\vs{16}Ἐὰν γὰρ εὐαγγελίζωμαι, οὐκ ἔστιν μοι καύχημα· ἀνάγκη γάρ μοι ἐπίκειται· οὐαὶ γάρ μοί ἐστιν ἐὰν μὴ εὐαγγελίσωμαι.
\vs{17}εἰ γὰρ ἑκὼν τοῦτο πράσσω, μισθὸν ἔχω· εἰ δὲ ἄκων, οἰκονομίαν πεπίστευμαι·
\vs{18}τίς οὖν μού ἐστιν ὁ μισθός; ἵνα εὐαγγελιζόμενος ἀδάπανον θήσω τὸ εὐαγγέλιον εἰς τὸ μὴ καταχρήσασθαι τῇ ἐξουσίᾳ μου ἐν τῷ εὐαγγελίῳ.

\vs{19}Ἐλεύθερος γὰρ ὢν ἐκ πάντων πᾶσιν ἐμαυτὸν ἐδούλωσα, ἵνα τοὺς πλείονας κερδήσω·
\vs{20}καὶ ἐγενόμην τοῖς Ἰουδαίοις ὡς Ἰουδαῖος, ἵνα Ἰουδαίους κερδήσω· τοῖς ὑπὸ νόμον ὡς ὑπὸ νόμον, μὴ ὢν αὐτὸς ὑπὸ νόμον, ἵνα τοὺς ὑπὸ νόμον κερδήσω·
\vs{21}τοῖς ἀνόμοις ὡς ἄνομος, μὴ ὢν ἄνομος Θεοῦ ἀλλ᾽ ἔννομος Χριστοῦ, ἵνα κερδάνω τοὺς ἀνόμους·
\vs{22}ἐγενόμην τοῖς ἀσθενέσιν ἀσθενής, ἵνα τοὺς ἀσθενεῖς κερδήσω· τοῖς πᾶσιν γέγονα πάντα, ἵνα πάντως τινὰς σώσω.
\vs{23}Πάντα δὲ ποιῶ διὰ τὸ εὐαγγέλιον, ἵνα συνκοινωνὸς αὐτοῦ γένωμαι.

\vs{24}Οὐκ οἴδατε ὅτι οἱ ἐν σταδίῳ τρέχοντες πάντες μὲν τρέχουσιν, εἷς δὲ λαμβάνει τὸ βραβεῖον; οὕτως τρέχετε ἵνα καταλάβητε.
\vs{25}πᾶς δὲ ὁ ἀγωνιζόμενος πάντα ἐγκρατεύεται, ἐκεῖνοι μὲν οὖν ἵνα φθαρτὸν στέφανον λάβωσιν, ἡμεῖς δὲ ἄφθαρτον.
\vs{26}ἐγὼ τοίνυν οὕτως τρέχω ὡς οὐκ ἀδήλως, οὕτως πυκτεύω ὡς οὐκ ἀέρα δέρων·
\vs{27}ἀλλὰ ὑπωπιάζω μου τὸ σῶμα καὶ δουλαγωγῶ, μή πως ἄλλοις κηρύξας αὐτὸς ἀδόκιμος γένωμαι.

\ch{10}
Οὐ θέλω γὰρ ὑμᾶς ἀγνοεῖν, ἀδελφοί, ὅτι οἱ πατέρες ἡμῶν πάντες ὑπὸ τὴν νεφέλην ἦσαν καὶ πάντες διὰ τῆς θαλάσσης διῆλθον
\vs{2}καὶ πάντες εἰς τὸν Μωϋσῆν ἐβαπτίσαντο ἐν τῇ νεφέλῃ καὶ ἐν τῇ θαλάσσῃ
\vs{3}καὶ πάντες τὸ αὐτὸ πνευματικὸν βρῶμα ἔφαγον
\vs{4}καὶ πάντες τὸ αὐτὸ πνευματικὸν ἔπιον πόμα· ἔπινον γὰρ ἐκ πνευματικῆς ἀκολουθούσης πέτρας, ἡ πέτρα δὲ ἦν ὁ Χριστός.
\vs{5}ἀλλ᾽ οὐκ ἐν τοῖς πλείοσιν αὐτῶν εὐδόκησεν ὁ Θεός, κατεστρώθησαν γὰρ ἐν τῇ ἐρήμῳ.
\vs{6}Ταῦτα δὲ τύποι ἡμῶν ἐγενήθησαν, εἰς τὸ μὴ εἶναι ἡμᾶς ἐπιθυμητὰς κακῶν, καθὼς κἀκεῖνοι ἐπεθύμησαν.
\vs{7}μηδὲ εἰδωλολάτραι γίνεσθε καθώς τινες αὐτῶν, ὥσπερ γέγραπται· Ἐκάθισεν ὁ λαὸς φαγεῖν καὶ πεῖν καὶ ἀνέστησαν παίζειν.
\vs{8}μηδὲ πορνεύωμεν, καθώς τινες αὐτῶν ἐπόρνευσαν καὶ ἔπεσαν μιᾷ ἡμέρᾳ εἰκοσι τρεῖς χιλιάδες.
\vs{9}μηδὲ ἐκπειράζωμεν τὸν Χριστόν, καθώς τινες αὐτῶν ἐπείρασαν καὶ ὑπὸ τῶν ὄφεων ἀπώλλυντο.
\vs{10}μηδὲ γογγύζετε, καθάπερ τινὲς αὐτῶν ἐγόγγυσαν καὶ ἀπώλοντο ὑπὸ τοῦ ὀλοθρευτοῦ.
\vs{11}Ταῦτα δὲ τυπικῶς συνέβαινεν ἐκείνοις, ἐγράφη δὲ πρὸς νουθεσίαν ἡμῶν, εἰς οὓς τὰ τέλη τῶν αἰώνων κατήντηκεν.
\vs{12}Ὥστε ὁ δοκῶν ἑστάναι βλεπέτω μὴ πέσῃ.
\vs{13}πειρασμὸς ὑμᾶς οὐκ εἴληφεν εἰ μὴ ἀνθρώπινος· πιστὸς δὲ ὁ Θεός, ὃς οὐκ ἐάσει ὑμᾶς πειρασθῆναι ὑπὲρ ὃ δύνασθε ἀλλὰ ποιήσει σὺν τῷ πειρασμῷ καὶ τὴν ἔκβασιν τοῦ δύνασθαι ὑπενεγκεῖν.

\vs{14}Διόπερ, ἀγαπητοί μου, φεύγετε ἀπὸ τῆς εἰδωλολατρίας.
\vs{15}ὡς φρονίμοις λέγω· κρίνατε ὑμεῖς ὅ φημι.
\vs{16}Τὸ ποτήριον τῆς εὐλογίας ὃ εὐλογοῦμεν, οὐχὶ κοινωνία ἐστὶν τοῦ αἵματος τοῦ Χριστοῦ; τὸν ἄρτον ὃν κλῶμεν, οὐχὶ κοινωνία τοῦ σώματος τοῦ Χριστοῦ ἐστιν;
\vs{17}ὅτι εἷς ἄρτος, ἓν σῶμα οἱ πολλοί ἐσμεν, οἱ γὰρ πάντες ἐκ τοῦ ἑνὸς ἄρτου μετέχομεν.
\vs{18}βλέπετε τὸν Ἰσραὴλ κατὰ σάρκα· οὐχ οἱ ἐσθίοντες τὰς θυσίας κοινωνοὶ τοῦ θυσιαστηρίου εἰσίν;
\vs{19}Τί οὖν φημι; ὅτι εἰδωλόθυτόν τί ἐστιν ἢ ὅτι εἴδωλόν τί ἐστιν;
\vs{20}ἀλλ᾽ ὅτι ἃ θύουσιν, δαιμονίοις καὶ οὐ Θεῷ θύουσιν· οὐ θέλω δὲ ὑμᾶς κοινωνοὺς τῶν δαιμονίων γίνεσθαι.
\vs{21}οὐ δύνασθε ποτήριον Κυρίου πίνειν καὶ ποτήριον δαιμονίων, οὐ δύνασθε τραπέζης Κυρίου μετέχειν καὶ τραπέζης δαιμονίων.
\vs{22}ἢ παραζηλοῦμεν τὸν Κύριον; μὴ ἰσχυρότεροι αὐτοῦ ἐσμεν;

\vs{23}Πάντα ἔξεστιν ἀλλ᾽ οὐ πάντα συμφέρει· Πάντα ἔξεστιν ἀλλ᾽ οὐ πάντα οἰκοδομεῖ.
\vs{24}μηδεὶς τὸ ἑαυτοῦ ζητείτω ἀλλὰ τὸ τοῦ ἑτέρου.
\vs{25}Πᾶν τὸ ἐν μακέλλῳ πωλούμενον ἐσθίετε μηδὲν ἀνακρίνοντες διὰ τὴν συνείδησιν·
\vs{26}Τοῦ Κυρίου γὰρ Ἡ γῆ καὶ τὸ πλήρωμα αὐτῆς.
\vs{27}Εἴ τις καλεῖ ὑμᾶς τῶν ἀπίστων καὶ θέλετε πορεύεσθαι, πᾶν τὸ παρατιθέμενον ὑμῖν ἐσθίετε μηδὲν ἀνακρίνοντες διὰ τὴν συνείδησιν.
\vs{28}ἐὰν δέ τις ὑμῖν εἴπῃ· Τοῦτο ἱερόθυτόν ἐστιν, μὴ ἐσθίετε δι᾽ ἐκεῖνον τὸν μηνύσαντα καὶ τὴν συνείδησιν·
\vs{29}συνείδησιν δὲ λέγω οὐχὶ τὴν ἑαυτοῦ ἀλλὰ τὴν τοῦ ἑτέρου. ἵνατί γὰρ ἡ ἐλευθερία μου κρίνεται ὑπὸ ἄλλης συνειδήσεως;
\vs{30}εἰ ἐγὼ χάριτι μετέχω, τί βλασφημοῦμαι ὑπὲρ οὗ ἐγὼ εὐχαριστῶ;
\vs{31}Εἴτε οὖν ἐσθίετε εἴτε πίνετε εἴτε τι ποιεῖτε, πάντα εἰς δόξαν Θεοῦ ποιεῖτε.
\vs{32}ἀπρόσκοποι καὶ Ἰουδαίοις γίνεσθε καὶ Ἕλλησιν καὶ τῇ ἐκκλησίᾳ τοῦ Θεοῦ,
\vs{33}καθὼς κἀγὼ πάντα πᾶσιν ἀρέσκω μὴ ζητῶν τὸ ἐμαυτοῦ σύμφορον ἀλλὰ τὸ τῶν πολλῶν, ἵνα σωθῶσιν.
\inparch{11} \vs{1}Μιμηταί μου γίνεσθε καθὼς κἀγὼ Χριστοῦ.

\vs{2}Ἐπαινῶ δὲ ὑμᾶς ὅτι πάντα μου μέμνησθε καὶ, καθὼς παρέδωκα ὑμῖν, τὰς παραδόσεις κατέχετε.
\vs{3}Θέλω δὲ ὑμᾶς εἰδέναι ὅτι παντὸς ἀνδρὸς ἡ κεφαλὴ ὁ Χριστός ἐστιν, κεφαλὴ δὲ γυναικὸς ὁ ἀνήρ, κεφαλὴ δὲ τοῦ Χριστοῦ ὁ Θεός.
\vs{4}Πᾶς ἀνὴρ προσευχόμενος ἢ προφητεύων κατὰ κεφαλῆς ἔχων καταισχύνει τὴν κεφαλὴν αὐτοῦ.
\vs{5}πᾶσα δὲ γυνὴ προσευχομένη ἢ προφητεύουσα ἀκατακαλύπτῳ τῇ κεφαλῇ καταισχύνει τὴν κεφαλὴν αὐτῆς· ἓν γάρ ἐστιν καὶ τὸ αὐτὸ τῇ ἐξυρημένῃ.
\vs{6}εἰ γὰρ οὐ κατακαλύπτεται γυνή, καὶ κειράσθω· εἰ δὲ αἰσχρὸν γυναικὶ τὸ κείρασθαι ἢ ξυρᾶσθαι, κατακαλυπτέσθω.
\vs{7}Ἀνὴρ μὲν γὰρ οὐκ ὀφείλει κατακαλύπτεσθαι τὴν κεφαλήν εἰκὼν καὶ δόξα Θεοῦ ὑπάρχων· ἡ γυνὴ δὲ δόξα ἀνδρός ἐστιν.
\vs{8}οὐ γάρ ἐστιν ἀνὴρ ἐκ γυναικός ἀλλὰ γυνὴ ἐξ ἀνδρός·
\vs{9}καὶ γὰρ οὐκ ἐκτίσθη ἀνὴρ διὰ τὴν γυναῖκα ἀλλὰ γυνὴ διὰ τὸν ἄνδρα.
\vs{10}διὰ τοῦτο ὀφείλει ἡ γυνὴ ἐξουσίαν ἔχειν ἐπὶ τῆς κεφαλῆς διὰ τοὺς ἀγγέλους.
\vs{11}Πλὴν οὔτε γυνὴ χωρὶς ἀνδρὸς οὔτε ἀνὴρ χωρὶς γυναικὸς ἐν Κυρίῳ·
\vs{12}ὥσπερ γὰρ ἡ γυνὴ ἐκ τοῦ ἀνδρός, οὕτως καὶ ὁ ἀνὴρ διὰ τῆς γυναικός· τὰ δὲ πάντα ἐκ τοῦ Θεοῦ.
\vs{13}Ἐν ὑμῖν αὐτοῖς κρίνατε· πρέπον ἐστὶν γυναῖκα ἀκατακάλυπτον τῷ Θεῷ προσεύχεσθαι;
\vs{14}οὐδὲ ἡ φύσις αὐτὴ διδάσκει ὑμᾶς ὅτι ἀνὴρ μὲν ἐὰν κομᾷ ἀτιμία αὐτῷ ἐστιν,
\vs{15}γυνὴ δὲ ἐὰν κομᾷ δόξα αὐτῇ ἐστιν; ὅτι ἡ κόμη ἀντὶ περιβολαίου δέδοται αὐτῇ.
\vs{16}Εἰ δέ τις δοκεῖ φιλόνεικος εἶναι, ἡμεῖς τοιαύτην συνήθειαν οὐκ ἔχομεν οὐδὲ αἱ ἐκκλησίαι τοῦ Θεοῦ.

\vs{17}Τοῦτο δὲ παραγγέλλων οὐκ ἐπαινῶ ὅτι οὐκ εἰς τὸ κρεῖσσον ἀλλὰ εἰς τὸ ἧσσον συνέρχεσθε.
\vs{18}πρῶτον μὲν γὰρ συνερχομένων ὑμῶν ἐν ἐκκλησίᾳ ἀκούω σχίσματα ἐν ὑμῖν ὑπάρχειν καὶ μέρος τι πιστεύω.
\vs{19}δεῖ γὰρ καὶ αἱρέσεις ἐν ὑμῖν εἶναι, ἵνα καὶ οἱ δόκιμοι φανεροὶ γένωνται ἐν ὑμῖν.
\vs{20}Συνερχομένων οὖν ὑμῶν ἐπὶ τὸ αὐτὸ οὐκ ἔστιν κυριακὸν δεῖπνον φαγεῖν·
\vs{21}ἕκαστος γὰρ τὸ ἴδιον δεῖπνον προλαμβάνει ἐν τῷ φαγεῖν, καὶ ὃς μὲν πεινᾷ ὃς δὲ μεθύει.
\vs{22}μὴ γὰρ οἰκίας οὐκ ἔχετε εἰς τὸ ἐσθίειν καὶ πίνειν; ἢ τῆς ἐκκλησίας τοῦ Θεοῦ καταφρονεῖτε, καὶ καταισχύνετε τοὺς μὴ ἔχοντας; τί εἴπω ὑμῖν; ἐπαινέσω ὑμᾶς; ἐν τούτῳ οὐκ ἐπαινῶ.

\vs{23}Ἐγὼ γὰρ παρέλαβον ἀπὸ τοῦ Κυρίου, ὃ καὶ παρέδωκα ὑμῖν, ὅτι ὁ Κύριος Ἰησοῦς ἐν τῇ νυκτὶ ᾗ παρεδίδετο ἔλαβεν ἄρτον
\vs{24}καὶ εὐχαριστήσας ἔκλασεν καὶ εἶπεν· Τοῦτό μού ἐστιν τὸ σῶμα τὸ ὑπὲρ ὑμῶν· τοῦτο ποιεῖτε εἰς τὴν ἐμὴν ἀνάμνησιν.
\vs{25}ὡσαύτως καὶ τὸ ποτήριον μετὰ τὸ δειπνῆσαι λέγων· Τοῦτο τὸ ποτήριον ἡ καινὴ διαθήκη ἐστὶν ἐν τῷ ἐμῷ αἵματι· τοῦτο ποιεῖτε, ὁσάκις ἐὰν πίνητε, εἰς τὴν ἐμὴν ἀνάμνησιν.
\vs{26}ὁσάκις γὰρ ἐὰν ἐσθίητε τὸν ἄρτον τοῦτον καὶ τὸ ποτήριον πίνητε, τὸν θάνατον τοῦ Κυρίου καταγγέλλετε ἄχρι οὗ ἔλθῃ.

\vs{27}Ὥστε ὃς ἂν ἐσθίῃ τὸν ἄρτον ἢ πίνῃ τὸ ποτήριον τοῦ Κυρίου ἀναξίως, ἔνοχος ἔσται τοῦ σώματος καὶ τοῦ αἵματος τοῦ Κυρίου.
\vs{28}δοκιμαζέτω δὲ ἄνθρωπος ἑαυτόν καὶ οὕτως ἐκ τοῦ ἄρτου ἐσθιέτω καὶ ἐκ τοῦ ποτηρίου πινέτω·
\vs{29}ὁ γὰρ ἐσθίων καὶ πίνων κρίμα ἑαυτῷ ἐσθίει καὶ πίνει μὴ διακρίνων τὸ σῶμα.
\vs{30}διὰ τοῦτο ἐν ὑμῖν πολλοὶ ἀσθενεῖς καὶ ἄρρωστοι καὶ κοιμῶνται ἱκανοί.
\vs{31}Εἰ δὲ ἑαυτοὺς διεκρίνομεν, οὐκ ἂν ἐκρινόμεθα·
\vs{32}κρινόμενοι δὲ ὑπὸ τοῦ Κυρίου παιδευόμεθα, ἵνα μὴ σὺν τῷ κόσμῳ κατακριθῶμεν.
\vs{33}Ὥστε, ἀδελφοί μου, συνερχόμενοι εἰς τὸ φαγεῖν ἀλλήλους ἐκδέχεσθε.
\vs{34}εἴ τις πεινᾷ, ἐν οἴκῳ ἐσθιέτω, ἵνα μὴ εἰς κρίμα συνέρχησθε. Τὰ δὲ λοιπὰ ὡς ἂν ἔλθω διατάξομαι.

\ch{12}
Περὶ δὲ τῶν πνευματικῶν, ἀδελφοί, οὐ θέλω ὑμᾶς ἀγνοεῖν.
\vs{2}Οἴδατε ὅτι ὅτε ἔθνη ἦτε πρὸς τὰ εἴδωλα τὰ ἄφωνα ὡς ἂν ἤγεσθε ἀπαγόμενοι.
\vs{3}διὸ γνωρίζω ὑμῖν ὅτι οὐδεὶς ἐν Πνεύματι Θεοῦ λαλῶν λέγει· Αναθεμα ΙΗΣΟΥΣ, καὶ οὐδεὶς δύναται εἰπεῖν· Κυριος ΙΗΣΟΥΣ, εἰ μὴ ἐν Πνεύματι Ἁγίῳ.

\vs{4}Διαιρέσεις δὲ χαρισμάτων εἰσίν, τὸ δὲ αὐτὸ Πνεῦμα·
\vs{5}καὶ διαιρέσεις διακονιῶν εἰσιν, καὶ ὁ αὐτὸς Κύριος·
\vs{6}καὶ διαιρέσεις ἐνεργημάτων εἰσίν, ὁ δὲ αὐτὸς Θεός ὁ ἐνεργῶν τὰ πάντα ἐν πᾶσιν.
\vs{7}Ἑκάστῳ δὲ δίδοται ἡ φανέρωσις τοῦ Πνεύματος πρὸς τὸ συμφέρον.
\vs{8}ᾧ μὲν γὰρ διὰ τοῦ Πνεύματος δίδοται λόγος σοφίας, ἄλλῳ δὲ λόγος γνώσεως κατὰ τὸ αὐτὸ Πνεῦμα,
\vs{9}ἑτέρῳ πίστις ἐν τῷ αὐτῷ Πνεύματι, ἄλλῳ δὲ χαρίσματα ἰαμάτων ἐν τῷ ἑνὶ Πνεύματι,
\vs{10}ἄλλῳ δὲ ἐνεργήματα δυνάμεων, ἄλλῳ δὲ προφητεία, ἄλλῳ δὲ διακρίσεις πνευμάτων, ἑτέρῳ γένη γλωσσῶν, ἄλλῳ δὲ ἑρμηνεία γλωσσῶν·
\vs{11}πάντα δὲ ταῦτα ἐνεργεῖ τὸ ἓν καὶ τὸ αὐτὸ Πνεῦμα διαιροῦν ἰδίᾳ ἑκάστῳ καθὼς βούλεται.

\vs{12}Καθάπερ γὰρ τὸ σῶμα ἕν ἐστιν καὶ μέλη πολλὰ ἔχει, πάντα δὲ τὰ μέλη τοῦ σώματος πολλὰ ὄντα ἕν ἐστιν σῶμα, οὕτως καὶ ὁ Χριστός·
\vs{13}καὶ γὰρ ἐν ἑνὶ Πνεύματι ἡμεῖς πάντες εἰς ἓν σῶμα ἐβαπτίσθημεν, εἴτε Ἰουδαῖοι εἴτε Ἕλληνες εἴτε δοῦλοι εἴτε ἐλεύθεροι, καὶ πάντες ἓν Πνεῦμα ἐποτίσθημεν.
\vs{14}Καὶ γὰρ τὸ σῶμα οὐκ ἔστιν ἓν μέλος ἀλλὰ πολλά.
\vs{15}ἐὰν εἴπῃ ὁ πούς· Ὅτι οὐκ εἰμὶ χείρ, οὐκ εἰμὶ ἐκ τοῦ σώματος, οὐ παρὰ τοῦτο οὐκ ἔστιν ἐκ τοῦ σώματος;
\vs{16}καὶ ἐὰν εἴπῃ τὸ οὖς· Ὅτι οὐκ εἰμὶ ὀφθαλμός, οὐκ εἰμὶ ἐκ τοῦ σώματος, οὐ παρὰ τοῦτο οὐκ ἔστιν ἐκ τοῦ σώματος;
\vs{17}εἰ ὅλον τὸ σῶμα ὀφθαλμός, ποῦ ἡ ἀκοή; εἰ ὅλον ἀκοή, ποῦ ἡ ὄσφρησις;
\vs{18}Νυνὶ δὲ ὁ Θεὸς ἔθετο τὰ μέλη, ἓν ἕκαστον αὐτῶν ἐν τῷ σώματι καθὼς ἠθέλησεν.
\vs{19}εἰ δὲ ἦν τὰ πάντα ἓν μέλος, ποῦ τὸ σῶμα;
\vs{20}νῦν δὲ πολλὰ μὲν μέλη, ἓν δὲ σῶμα.
\vs{21}Οὐ δύναται δὲ ὁ ὀφθαλμὸς εἰπεῖν τῇ χειρί· Χρείαν σου οὐκ ἔχω, ἢ πάλιν ἡ κεφαλὴ τοῖς ποσίν· Χρείαν ὑμῶν οὐκ ἔχω·
\vs{22}ἀλλὰ πολλῷ μᾶλλον τὰ δοκοῦντα μέλη τοῦ σώματος ἀσθενέστερα ὑπάρχειν ἀναγκαῖά ἐστιν,
\vs{23}καὶ ἃ δοκοῦμεν ἀτιμότερα εἶναι τοῦ σώματος τούτοις τιμὴν περισσοτέραν περιτίθεμεν, καὶ τὰ ἀσχήμονα ἡμῶν εὐσχημοσύνην περισσοτέραν ἔχει,
\vs{24}τὰ δὲ εὐσχήμονα ἡμῶν οὐ χρείαν ἔχει. Ἀλλὰ ὁ θεὸς συνεκέρασεν τὸ σῶμα τῷ ὑστερουμένῳ περισσοτέραν δοὺς τιμήν,
\vs{25}ἵνα μὴ ᾖ σχίσμα ἐν τῷ σώματι ἀλλὰ τὸ αὐτὸ ὑπὲρ ἀλλήλων μεριμνῶσιν τὰ μέλη.
\vs{26}καὶ εἴτε πάσχει ἓν μέλος, συμπάσχει πάντα τὰ μέλη· εἴτε δοξάζεται ἓν μέλος, συνχαίρει πάντα τὰ μέλη.
\vs{27}Ὑμεῖς δέ ἐστε σῶμα Χριστοῦ καὶ μέλη ἐκ μέρους.
\vs{28}Καὶ οὓς μὲν ἔθετο ὁ Θεὸς ἐν τῇ ἐκκλησίᾳ πρῶτον ἀποστόλους, δεύτερον προφήτας, τρίτον διδασκάλους, ἔπειτα δυνάμεις, ἔπειτα χαρίσματα ἰαμάτων, ἀντιλήμψεις, κυβερνήσεις, γένη γλωσσῶν.
\vs{29}μὴ πάντες ἀπόστολοι; μὴ πάντες προφῆται; μὴ πάντες διδάσκαλοι; μὴ πάντες δυνάμεις;
\vs{30}μὴ πάντες χαρίσματα ἔχουσιν ἰαμάτων; μὴ πάντες γλώσσαις λαλοῦσιν; μὴ πάντες διερμηνεύουσιν;
\vs{31}ζηλοῦτε δὲ τὰ χαρίσματα τὰ μείζονα.

Καὶ ἔτι καθ᾽ ὑπερβολὴν ὁδὸν ὑμῖν δείκνυμι.

\ch{13}
Ἐὰν ταῖς γλώσσαις τῶν ἀνθρώπων λαλῶ καὶ τῶν ἀγγέλων, ἀγάπην δὲ μὴ ἔχω, γέγονα χαλκὸς ἠχῶν ἢ κύμβαλον ἀλαλάζον.
\vs{2}καὶ ἐὰν ἔχω προφητείαν καὶ εἰδῶ τὰ μυστήρια πάντα καὶ πᾶσαν τὴν γνῶσιν καὶ ἐὰν ἔχω πᾶσαν τὴν πίστιν ὥστε ὄρη μεθιστάναι, ἀγάπην δὲ μὴ ἔχω, οὐθέν εἰμι.
\vs{3}κἂν ψωμίσω πάντα τὰ ὑπάρχοντά μου καὶ ἐὰν παραδῶ τὸ σῶμά μου ἵνα καυχήσωμαι, ἀγάπην δὲ μὴ ἔχω, οὐδὲν ὠφελοῦμαι.

\vs{4}Ἡ ἀγάπη μακροθυμεῖ, χρηστεύεται ἡ ἀγάπη, οὐ ζηλοῖ, ἡ ἀγάπη οὐ περπερεύεται, οὐ φυσιοῦται,
\vs{5}οὐκ ἀσχημονεῖ, οὐ ζητεῖ τὰ ἑαυτῆς, οὐ παροξύνεται, οὐ λογίζεται τὸ κακόν,
\vs{6}οὐ χαίρει ἐπὶ τῇ ἀδικίᾳ, συνχαίρει δὲ τῇ ἀληθείᾳ·
\vs{7}πάντα στέγει, πάντα πιστεύει, πάντα ἐλπίζει, πάντα ὑπομένει.

\vs{8}Ἡ ἀγάπη οὐδέποτε πίπτει· εἴτε δὲ προφητεῖαι, καταργηθήσονται· εἴτε γλῶσσαι, παύσονται· εἴτε γνῶσις, καταργηθήσεται.
\vs{9}ἐκ μέρους γὰρ γινώσκομεν καὶ ἐκ μέρους προφητεύομεν·
\vs{10}ὅταν δὲ ἔλθῃ τὸ τέλειον, τὸ ἐκ μέρους καταργηθήσεται.
\vs{11}Ὅτε ἤμην νήπιος, ἐλάλουν ὡς νήπιος, ἐφρόνουν ὡς νήπιος, ἐλογιζόμην ὡς νήπιος· ὅτε γέγονα ἀνήρ, κατήργηκα τὰ τοῦ νηπίου.
\vs{12}βλέπομεν γὰρ ἄρτι δι᾽ ἐσόπτρου ἐν αἰνίγματι, τότε δὲ πρόσωπον πρὸς πρόσωπον· ἄρτι γινώσκω ἐκ μέρους, τότε δὲ ἐπιγνώσομαι καθὼς καὶ ἐπεγνώσθην.
\vs{13}Νυνὶ δὲ μένει πίστις, ἐλπίς, ἀγάπη, τὰ τρία ταῦτα· μείζων δὲ τούτων ἡ ἀγάπη.

\ch{14}
Διώκετε τὴν ἀγάπην, ζηλοῦτε δὲ τὰ πνευματικά, μᾶλλον δὲ ἵνα προφητεύητε.
\vs{2}ὁ γὰρ λαλῶν γλώσσῃ οὐκ ἀνθρώποις λαλεῖ ἀλλὰ Θεῷ· οὐδεὶς γὰρ ἀκούει, πνεύματι δὲ λαλεῖ μυστήρια·
\vs{3}ὁ δὲ προφητεύων ἀνθρώποις λαλεῖ οἰκοδομὴν καὶ παράκλησιν καὶ παραμυθίαν.
\vs{4}ὁ λαλῶν γλώσσῃ ἑαυτὸν οἰκοδομεῖ· ὁ δὲ προφητεύων ἐκκλησίαν οἰκοδομεῖ.
\vs{5}Θέλω δὲ πάντας ὑμᾶς λαλεῖν γλώσσαις, μᾶλλον δὲ ἵνα προφητεύητε· μείζων δὲ ὁ προφητεύων ἢ ὁ λαλῶν γλώσσαις ἐκτὸς εἰ μὴ διερμηνεύῃ, ἵνα ἡ ἐκκλησία οἰκοδομὴν λάβῃ.

\vs{6}Νῦν δέ, ἀδελφοί, ἐὰν ἔλθω πρὸς ὑμᾶς γλώσσαις λαλῶν, τί ὑμᾶς ὠφελήσω ἐὰν μὴ ὑμῖν λαλήσω ἢ ἐν ἀποκαλύψει ἢ ἐν γνώσει ἢ ἐν προφητείᾳ ἢ ἐν διδαχῇ;
\vs{7}ὅμως τὰ ἄψυχα φωνὴν διδόντα, εἴτε αὐλὸς εἴτε κιθάρα, ἐὰν διαστολὴν τοῖς φθόγγοις μὴ δῷ, πῶς γνωσθήσεται τὸ αὐλούμενον ἢ τὸ κιθαριζόμενον;
\vs{8}Καὶ γὰρ ἐὰν ἄδηλον σάλπιγξ φωνὴν δῷ, τίς παρασκευάσεται εἰς πόλεμον;
\vs{9}οὕτως καὶ ὑμεῖς διὰ τῆς γλώσσης ἐὰν μὴ εὔσημον λόγον δῶτε, πῶς γνωσθήσεται τὸ λαλούμενον; ἔσεσθε γὰρ εἰς ἀέρα λαλοῦντες.
\vs{10}Τοσαῦτα εἰ τύχοι γένη φωνῶν εἰσιν ἐν κόσμῳ καὶ οὐδὲν ἄφωνον·
\vs{11}ἐὰν οὖν μὴ εἰδῶ τὴν δύναμιν τῆς φωνῆς, ἔσομαι τῷ λαλοῦντι βάρβαρος καὶ ὁ λαλῶν ἐν ἐμοὶ βάρβαρος.
\vs{12}Οὕτως καὶ ὑμεῖς, ἐπεὶ ζηλωταί ἐστε πνευμάτων, πρὸς τὴν οἰκοδομὴν τῆς ἐκκλησίας ζητεῖτε ἵνα περισσεύητε.

\vs{13}Διὸ ὁ λαλῶν γλώσσῃ προσευχέσθω ἵνα διερμηνεύῃ.
\vs{14}ἐὰν γὰρ προσεύχωμαι γλώσσῃ, τὸ πνεῦμά μου προσεύχεται, ὁ δὲ νοῦς μου ἄκαρπός ἐστιν.
\vs{15}Τί οὖν ἐστιν; προσεύξομαι τῷ πνεύματι, προσεύξομαι δὲ καὶ τῷ νοΐ· ψαλῶ τῷ πνεύματι, ψαλῶ δὲ καὶ τῷ νοΐ.
\vs{16}ἐπεὶ ἐὰν εὐλογῇς ἐν πνεύματι, ὁ ἀναπληρῶν τὸν τόπον τοῦ ἰδιώτου πῶς ἐρεῖ τὸ Ἀμήν ἐπὶ τῇ σῇ εὐχαριστίᾳ; ἐπειδὴ τί λέγεις οὐκ οἶδεν·
\vs{17}σὺ μὲν γὰρ καλῶς εὐχαριστεῖς ἀλλ᾽ ὁ ἕτερος οὐκ οἰκοδομεῖται.
\vs{18}Εὐχαριστῶ τῷ Θεῷ, πάντων ὑμῶν μᾶλλον γλώσσαις λαλῶ·
\vs{19}ἀλλὰ ἐν ἐκκλησίᾳ θέλω πέντε λόγους τῷ νοΐ μου λαλῆσαι, ἵνα καὶ ἄλλους κατηχήσω, ἢ μυρίους λόγους ἐν γλώσσῃ.

\vs{20}Ἀδελφοί, μὴ παιδία γίνεσθε ταῖς φρεσίν ἀλλὰ τῇ κακίᾳ νηπιάζετε, ταῖς δὲ φρεσὶν τέλειοι γίνεσθε.
\vs{21}ἐν τῷ νόμῳ γέγραπται ὅτι 
\begin{poetryblock}

\begin{quote}Ἐν ἑτερογλώσσοις καὶ ἐν χείλεσιν ἑτέρων λαλήσω τῷ λαῷ τούτῳ\end{quote} 

\begin{quote}καὶ οὐδ᾽ οὕτως εἰσακούσονταί μου, λέγει Κύριος.\end{quote}
\end{poetryblock}

\vs{22}Ὥστε αἱ γλῶσσαι εἰς σημεῖόν εἰσιν οὐ τοῖς πιστεύουσιν ἀλλὰ τοῖς ἀπίστοις, ἡ δὲ προφητεία οὐ τοῖς ἀπίστοις ἀλλὰ τοῖς πιστεύουσιν.
\vs{23}Ἐὰν οὖν συνέλθῃ ἡ ἐκκλησία ὅλη ἐπὶ τὸ αὐτὸ καὶ πάντες λαλῶσιν γλώσσαις, εἰσέλθωσιν δὲ ἰδιῶται ἢ ἄπιστοι, οὐκ ἐροῦσιν ὅτι μαίνεσθε;
\vs{24}ἐὰν δὲ πάντες προφητεύωσιν, εἰσέλθῃ δέ τις ἄπιστος ἢ ἰδιώτης, ἐλέγχεται ὑπὸ πάντων, ἀνακρίνεται ὑπὸ πάντων,
\vs{25}τὰ κρυπτὰ τῆς καρδίας αὐτοῦ φανερὰ γίνεται, καὶ οὕτως πεσὼν ἐπὶ πρόσωπον προσκυνήσει τῷ Θεῷ ἀπαγγέλλων ὅτι Ὄντως ὁ Θεὸς ἐν ὑμῖν ἐστιν.

\vs{26}Τί οὖν ἐστιν, ἀδελφοί; ὅταν συνέρχησθε, ἕκαστος ψαλμὸν ἔχει, διδαχὴν ἔχει, ἀποκάλυψιν ἔχει, γλῶσσαν ἔχει, ἑρμηνείαν ἔχει· πάντα πρὸς οἰκοδομὴν γινέσθω.
\vs{27}Εἴτε γλώσσῃ τις λαλεῖ, κατὰ δύο ἢ τὸ πλεῖστον τρεῖς καὶ ἀνὰ μέρος, καὶ εἷς διερμηνευέτω·
\vs{28}ἐὰν δὲ μὴ ᾖ διερμηνευτής, σιγάτω ἐν ἐκκλησίᾳ, ἑαυτῷ δὲ λαλείτω καὶ τῷ Θεῷ.
\vs{29}Προφῆται δὲ δύο ἢ τρεῖς λαλείτωσαν καὶ οἱ ἄλλοι διακρινέτωσαν·
\vs{30}ἐὰν δὲ ἄλλῳ ἀποκαλυφθῇ καθημένῳ, ὁ πρῶτος σιγάτω.
\vs{31}δύνασθε γὰρ καθ᾽ ἕνα πάντες προφητεύειν, ἵνα πάντες μανθάνωσιν καὶ πάντες παρακαλῶνται.
\vs{32}καὶ πνεύματα προφητῶν προφήταις ὑποτάσσεται,
\vs{33}οὐ γάρ ἐστιν ἀκαταστασίας ὁ Θεὸς ἀλλὰ εἰρήνης.

Ὡς ἐν πάσαις ταῖς ἐκκλησίαις τῶν ἁγίων
\vs{34}αἱ γυναῖκες ἐν ταῖς ἐκκλησίαις σιγάτωσαν· οὐ γὰρ ἐπιτρέπεται αὐταῖς λαλεῖν, ἀλλὰ ὑποτασσέσθωσαν, καθὼς καὶ ὁ νόμος λέγει.
\vs{35}εἰ δέ τι μαθεῖν θέλουσιν, ἐν οἴκῳ τοὺς ἰδίους ἄνδρας ἐπερωτάτωσαν· αἰσχρὸν γάρ ἐστιν γυναικὶ λαλεῖν ἐν ἐκκλησίᾳ.
\vs{36}Ἢ ἀφ᾽ ὑμῶν ὁ λόγος τοῦ Θεοῦ ἐξῆλθεν, ἢ εἰς ὑμᾶς μόνους κατήντησεν;

\vs{37}Εἴ τις δοκεῖ προφήτης εἶναι ἢ πνευματικός, ἐπιγινωσκέτω ἃ γράφω ὑμῖν ὅτι Κυρίου ἐστὶν ἐντολή·
\vs{38}εἰ δέ τις ἀγνοεῖ, ἀγνοεῖται.
\vs{39}Ὥστε, ἀδελφοί μου, ζηλοῦτε τὸ προφητεύειν καὶ τὸ λαλεῖν μὴ κωλύετε γλώσσαις·
\vs{40}πάντα δὲ εὐσχημόνως καὶ κατὰ τάξιν γινέσθω.

\ch{15}
Γνωρίζω δὲ ὑμῖν, ἀδελφοί, τὸ εὐαγγέλιον ὃ εὐηγγελισάμην ὑμῖν, ὃ καὶ παρελάβετε, ἐν ᾧ καὶ ἑστήκατε,
\vs{2}δι᾽ οὗ καὶ σῴζεσθε, τίνι λόγῳ εὐηγγελισάμην ὑμῖν εἰ κατέχετε, ἐκτὸς εἰ μὴ εἰκῇ ἐπιστεύσατε.
\vs{3}Παρέδωκα γὰρ ὑμῖν ἐν πρώτοις, ὃ καὶ παρέλαβον, ὅτι Χριστὸς ἀπέθανεν ὑπὲρ τῶν ἁμαρτιῶν ἡμῶν κατὰ τὰς γραφάς
\vs{4}καὶ ὅτι ἐτάφη καὶ ὅτι ἐγήγερται τῇ ἡμέρᾳ τῇ τρίτῃ κατὰ τὰς γραφάς
\vs{5}καὶ ὅτι ὤφθη Κηφᾷ εἶτα τοῖς δώδεκα·
\vs{6}ἔπειτα ὤφθη ἐπάνω πεντακοσίοις ἀδελφοῖς ἐφάπαξ, ἐξ ὧν οἱ πλείονες μένουσιν ἕως ἄρτι, τινὲς δὲ ἐκοιμήθησαν·
\vs{7}ἔπειτα ὤφθη Ἰακώβῳ εἶτα τοῖς ἀποστόλοις πᾶσιν·
\vs{8}ἔσχατον δὲ πάντων ὡσπερεὶ τῷ ἐκτρώματι ὤφθη κἀμοί.
\vs{9}Ἐγὼ γάρ εἰμι ὁ ἐλάχιστος τῶν ἀποστόλων ὃς οὐκ εἰμὶ ἱκανὸς καλεῖσθαι ἀπόστολος, διότι ἐδίωξα τὴν ἐκκλησίαν τοῦ Θεοῦ·
\vs{10}χάριτι δὲ Θεοῦ εἰμι ὅ εἰμι, καὶ ἡ χάρις αὐτοῦ ἡ εἰς ἐμὲ οὐ κενὴ ἐγενήθη, ἀλλὰ περισσότερον αὐτῶν πάντων ἐκοπίασα, οὐκ ἐγὼ δὲ ἀλλὰ ἡ χάρις τοῦ Θεοῦ ἡ σὺν ἐμοί.
\vs{11}εἴτε οὖν ἐγὼ εἴτε ἐκεῖνοι, οὕτως κηρύσσομεν καὶ οὕτως ἐπιστεύσατε.

\vs{12}Εἰ δὲ Χριστὸς κηρύσσεται ὅτι ἐκ νεκρῶν ἐγήγερται, πῶς λέγουσιν ἐν ὑμῖν τινες ὅτι ἀνάστασις νεκρῶν οὐκ ἔστιν;
\vs{13}εἰ δὲ ἀνάστασις νεκρῶν οὐκ ἔστιν, οὐδὲ Χριστὸς ἐγήγερται·
\vs{14}εἰ δὲ Χριστὸς οὐκ ἐγήγερται, κενὸν ἄρα καὶ τὸ κήρυγμα ἡμῶν, κενὴ καὶ ἡ πίστις ὑμῶν·
\vs{15}εὑρισκόμεθα δὲ καὶ ψευδομάρτυρες τοῦ Θεοῦ, ὅτι ἐμαρτυρήσαμεν κατὰ τοῦ Θεοῦ ὅτι ἤγειρεν τὸν Χριστόν, ὃν οὐκ ἤγειρεν εἴπερ ἄρα νεκροὶ οὐκ ἐγείρονται.
\vs{16}Εἰ γὰρ νεκροὶ οὐκ ἐγείρονται, οὐδὲ Χριστὸς ἐγήγερται·
\vs{17}εἰ δὲ Χριστὸς οὐκ ἐγήγερται, ματαία ἡ πίστις ὑμῶν, ἔτι ἐστὲ ἐν ταῖς ἁμαρτίαις ὑμῶν,
\vs{18}ἄρα καὶ οἱ κοιμηθέντες ἐν Χριστῷ ἀπώλοντο.
\vs{19}εἰ ἐν τῇ ζωῇ ταύτῃ ἐν Χριστῷ ἠλπικότες ἐσμὲν μόνον, ἐλεεινότεροι πάντων ἀνθρώπων ἐσμέν.

\vs{20}Νυνὶ δὲ Χριστὸς ἐγήγερται ἐκ νεκρῶν ἀπαρχὴ τῶν κεκοιμημένων.
\vs{21}ἐπειδὴ γὰρ δι᾽ ἀνθρώπου θάνατος, καὶ δι᾽ ἀνθρώπου ἀνάστασις νεκρῶν.
\vs{22}ὥσπερ γὰρ ἐν τῷ Ἀδὰμ πάντες ἀποθνήσκουσιν, οὕτως καὶ ἐν τῷ Χριστῷ πάντες ζωοποιηθήσονται.
\vs{23}Ἕκαστος δὲ ἐν τῷ ἰδίῳ τάγματι· ἀπαρχὴ Χριστός, ἔπειτα οἱ τοῦ Χριστοῦ ἐν τῇ παρουσίᾳ αὐτοῦ,
\vs{24}εἶτα τὸ τέλος, ὅταν παραδιδῷ τὴν βασιλείαν τῷ Θεῷ καὶ Πατρί, ὅταν καταργήσῃ πᾶσαν ἀρχὴν καὶ πᾶσαν ἐξουσίαν καὶ δύναμιν.
\vs{25}δεῖ γὰρ αὐτὸν βασιλεύειν ἄχρι οὗ θῇ πάντας τοὺς ἐχθροὺς ὑπὸ τοὺς πόδας αὐτοῦ.
\vs{26}ἔσχατος ἐχθρὸς καταργεῖται ὁ θάνατος·
\vs{27}Πάντα γὰρ Ὑπέταξεν ὑπὸ τοὺς πόδας αὐτοῦ. ὅταν δὲ εἴπῃ ὅτι πάντα ὑποτέτακται, δῆλον ὅτι ἐκτὸς τοῦ ὑποτάξαντος αὐτῷ τὰ πάντα.
\vs{28}ὅταν δὲ ὑποταγῇ αὐτῷ τὰ πάντα, τότε καὶ αὐτὸς ὁ Υἱὸς ὑποταγήσεται τῷ ὑποτάξαντι αὐτῷ τὰ πάντα, ἵνα ᾖ ὁ Θεὸς τὰ πάντα ἐν πᾶσιν.

\vs{29}Ἐπεὶ τί ποιήσουσιν οἱ βαπτιζόμενοι ὑπὲρ τῶν νεκρῶν; εἰ ὅλως νεκροὶ οὐκ ἐγείρονται, τί καὶ βαπτίζονται ὑπὲρ αὐτῶν;
\vs{30}τί καὶ ἡμεῖς κινδυνεύομεν πᾶσαν ὥραν;
\vs{31}καθ᾽ ἡμέραν ἀποθνῄσκω, νὴ τὴν ὑμετέραν καύχησιν, ἀδελφοί, ἣν ἔχω ἐν Χριστῷ Ἰησοῦ τῷ Κυρίῳ ἡμῶν.
\vs{32}εἰ κατὰ ἄνθρωπον ἐθηριομάχησα ἐν Ἐφέσῳ, τί μοι τὸ ὄφελος; εἰ νεκροὶ οὐκ ἐγείρονται, Φάγωμεν καὶ πίωμεν, αὔριον γὰρ ἀποθνήσκομεν.
\vs{33}Μὴ πλανᾶσθε· 
\begin{poetryblock}

\begin{quote}Φθείρουσιν ἤθη χρηστὰ ὁμιλίαι κακαί.\end{quote}
\end{poetryblock}

\vs{34}ἐκνήψατε δικαίως καὶ μὴ ἁμαρτάνετε, ἀγνωσίαν γὰρ Θεοῦ τινες ἔχουσιν, πρὸς ἐντροπὴν ὑμῖν λαλῶ.

\vs{35}Ἀλλὰ ἐρεῖ τις· Πῶς ἐγείρονται οἱ νεκροί; ποίῳ δὲ σώματι ἔρχονται;
\vs{36}ἄφρων, σὺ ὃ σπείρεις, οὐ ζωοποιεῖται ἐὰν μὴ ἀποθάνῃ·
\vs{37}καὶ ὃ σπείρεις, οὐ τὸ σῶμα τὸ γενησόμενον σπείρεις ἀλλὰ γυμνὸν κόκκον εἰ τύχοι σίτου ἤ τινος τῶν λοιπῶν·
\vs{38}ὁ δὲ Θεὸς δίδωσιν αὐτῷ σῶμα καθὼς ἠθέλησεν, καὶ ἑκάστῳ τῶν σπερμάτων ἴδιον σῶμα.
\vs{39}Οὐ πᾶσα σὰρξ ἡ αὐτὴ σάρξ ἀλλὰ ἄλλη μὲν ἀνθρώπων, ἄλλη δὲ σὰρξ κτηνῶν, ἄλλη δὲ σὰρξ πτηνῶν, ἄλλη δὲ ἰχθύων.
\vs{40}καὶ σώματα ἐπουράνια, καὶ σώματα ἐπίγεια· ἀλλὰ ἑτέρα μὲν ἡ τῶν ἐπουρανίων δόξα, ἑτέρα δὲ ἡ τῶν ἐπιγείων.
\vs{41}ἄλλη δόξα ἡλίου, καὶ ἄλλη δόξα σελήνης, καὶ ἄλλη δόξα ἀστέρων· ἀστὴρ γὰρ ἀστέρος διαφέρει ἐν δόξῃ.
\vs{42}Οὕτως καὶ ἡ ἀνάστασις τῶν νεκρῶν. σπείρεται ἐν φθορᾷ, ἐγείρεται ἐν ἀφθαρσίᾳ·
\vs{43}σπείρεται ἐν ἀτιμίᾳ, ἐγείρεται ἐν δόξῃ· σπείρεται ἐν ἀσθενείᾳ, ἐγείρεται ἐν δυνάμει·
\vs{44}σπείρεται σῶμα ψυχικόν, ἐγείρεται σῶμα πνευματικόν. Εἰ ἔστιν σῶμα ψυχικόν, ἔστιν καὶ πνευματικόν.
\vs{45}οὕτως καὶ γέγραπται· Ἐγένετο ὁ πρῶτος ἄνθρωπος Ἀδὰμ εἰς ψυχὴν ζῶσαν, ὁ ἔσχατος Ἀδὰμ εἰς πνεῦμα ζωοποιοῦν.
\vs{46}Ἀλλ᾽ οὐ πρῶτον τὸ πνευματικὸν ἀλλὰ τὸ ψυχικόν, ἔπειτα τὸ πνευματικόν.
\vs{47}ὁ πρῶτος ἄνθρωπος ἐκ γῆς χοϊκός, ὁ δεύτερος ἄνθρωπος ἐξ οὐρανοῦ.
\vs{48}οἷος ὁ χοϊκός, τοιοῦτοι καὶ οἱ χοϊκοί, καὶ οἷος ὁ ἐπουράνιος, τοιοῦτοι καὶ οἱ ἐπουράνιοι·
\vs{49}καὶ καθὼς ἐφορέσαμεν τὴν εἰκόνα τοῦ χοϊκοῦ, φορέσομεν καὶ τὴν εἰκόνα τοῦ ἐπουρανίου.

\vs{50}Τοῦτο δέ φημι, ἀδελφοί, ὅτι σὰρξ καὶ αἷμα βασιλείαν Θεοῦ κληρονομῆσαι οὐ δύναται οὐδὲ ἡ φθορὰ τὴν ἀφθαρσίαν κληρονομεῖ.
\vs{51}Ἰδοὺ μυστήριον ὑμῖν λέγω· πάντες οὐ κοιμηθησόμεθα, πάντες δὲ ἀλλαγησόμεθα,
\vs{52}ἐν ἀτόμῳ, ἐν ῥιπῇ ὀφθαλμοῦ, ἐν τῇ ἐσχάτῃ σάλπιγγι· σαλπίσει γάρ καὶ οἱ νεκροὶ ἐγερθήσονται ἄφθαρτοι καὶ ἡμεῖς ἀλλαγησόμεθα.
\vs{53}δεῖ γὰρ τὸ φθαρτὸν τοῦτο ἐνδύσασθαι ἀφθαρσίαν καὶ τὸ θνητὸν τοῦτο ἐνδύσασθαι ἀθανασίαν.
\vs{54}Ὅταν δὲ τὸ φθαρτὸν τοῦτο ἐνδύσηται ἀφθαρσίαν καὶ τὸ θνητὸν τοῦτο ἐνδύσηται ἀθανασίαν, τότε γενήσεται ὁ λόγος ὁ γεγραμμένος· 
\begin{poetryblock}

\begin{quote}Κατεπόθη ὁ θάνατος εἰς νῖκος.\end{quote}
\end{poetryblock}

\vs{55}Ποῦ σου, θάνατε, τὸ νῖκος; 
\begin{poetryblock}

\begin{quote}ποῦ σου, θάνατε, τὸ κέντρον;\end{quote}
\end{poetryblock}

\vs{56}Τὸ δὲ κέντρον τοῦ θανάτου ἡ ἁμαρτία, ἡ δὲ δύναμις τῆς ἁμαρτίας ὁ νόμος·
\vs{57}τῷ δὲ Θεῷ χάρις τῷ διδόντι ἡμῖν τὸ νῖκος διὰ τοῦ Κυρίου ἡμῶν Ἰησοῦ Χριστοῦ.
\vs{58}Ὥστε, ἀδελφοί μου ἀγαπητοί, ἑδραῖοι γίνεσθε, ἀμετακίνητοι, περισσεύοντες ἐν τῷ ἔργῳ τοῦ Κυρίου πάντοτε, εἰδότες ὅτι ὁ κόπος ὑμῶν οὐκ ἔστιν κενὸς ἐν Κυρίῳ.

\ch{16}
Περὶ δὲ τῆς λογείας τῆς εἰς τοὺς ἁγίους ὥσπερ διέταξα ταῖς ἐκκλησίαις τῆς Γαλατίας, οὕτως καὶ ὑμεῖς ποιήσατε.
\vs{2}κατὰ μίαν σαββάτου ἕκαστος ὑμῶν παρ᾽ ἑαυτῷ τιθέτω θησαυρίζων ὅ τι ἐὰν εὐοδῶται, ἵνα μὴ ὅταν ἔλθω τότε λογεῖαι γίνωνται.
\vs{3}ὅταν δὲ παραγένωμαι, οὓς ἐὰν δοκιμάσητε, δι᾽ ἐπιστολῶν τούτους πέμψω ἀπενεγκεῖν τὴν χάριν ὑμῶν εἰς Ἰερουσαλήμ·
\vs{4}ἐὰν δὲ ἄξιον ᾖ τοῦ κἀμὲ πορεύεσθαι, σὺν ἐμοὶ πορεύσονται.

\vs{5}Ἐλεύσομαι δὲ πρὸς ὑμᾶς ὅταν Μακεδονίαν διέλθω· Μακεδονίαν γὰρ διέρχομαι,
\vs{6}πρὸς ὑμᾶς δὲ τυχὸν παραμενῶ ἢ καὶ παραχειμάσω, ἵνα ὑμεῖς με προπέμψητε οὗ ἐὰν πορεύωμαι.
\vs{7}οὐ θέλω γὰρ ὑμᾶς ἄρτι ἐν παρόδῳ ἰδεῖν, ἐλπίζω γὰρ χρόνον τινὰ ἐπιμεῖναι πρὸς ὑμᾶς ἐὰν ὁ Κύριος ἐπιτρέψῃ.
\vs{8}ἐπιμενῶ δὲ ἐν Ἐφέσῳ ἕως τῆς Πεντηκοστῆς·
\vs{9}θύρα γάρ μοι ἀνέῳγεν μεγάλη καὶ ἐνεργής, καὶ ἀντικείμενοι πολλοί.

\vs{10}Ἐὰν δὲ ἔλθῃ Τιμόθεος, βλέπετε, ἵνα ἀφόβως γένηται πρὸς ὑμᾶς· τὸ γὰρ ἔργον Κυρίου ἐργάζεται ὡς κἀγώ·
\vs{11}μή τις οὖν αὐτὸν ἐξουθενήσῃ. προπέμψατε δὲ αὐτὸν ἐν εἰρήνῃ, ἵνα ἔλθῃ πρός με· ἐκδέχομαι γὰρ αὐτὸν μετὰ τῶν ἀδελφῶν.
\vs{12}Περὶ δὲ Ἀπολλῶ τοῦ ἀδελφοῦ, πολλὰ παρεκάλεσα αὐτὸν, ἵνα ἔλθῃ πρὸς ὑμᾶς μετὰ τῶν ἀδελφῶν· καὶ πάντως οὐκ ἦν θέλημα ἵνα νῦν ἔλθῃ· ἐλεύσεται δὲ ὅταν εὐκαιρήσῃ.

\vs{13}Γρηγορεῖτε, στήκετε ἐν τῇ πίστει, ἀνδρίζεσθε, κραταιοῦσθε.
\vs{14}πάντα ὑμῶν ἐν ἀγάπῃ γινέσθω.

\vs{15}Παρακαλῶ δὲ ὑμᾶς, ἀδελφοί· οἴδατε τὴν οἰκίαν Στεφανᾶ, ὅτι ἐστὶν ἀπαρχὴ τῆς Ἀχαΐας καὶ εἰς διακονίαν τοῖς ἁγίοις ἔταξαν ἑαυτούς·
\vs{16}ἵνα καὶ ὑμεῖς ὑποτάσσησθε τοῖς τοιούτοις καὶ παντὶ τῷ συνεργοῦντι καὶ κοπιῶντι.
\vs{17}Χαίρω δὲ ἐπὶ τῇ παρουσίᾳ Στεφανᾶ καὶ Φορτουνάτου καὶ Ἀχαϊκοῦ, ὅτι τὸ ὑμέτερον ὑστέρημα οὗτοι ἀνεπλήρωσαν·
\vs{18}ἀνέπαυσαν γὰρ τὸ ἐμὸν πνεῦμα καὶ τὸ ὑμῶν. ἐπιγινώσκετε οὖν τοὺς τοιούτους.

\vs{19}Ἀσπάζονται ὑμᾶς αἱ ἐκκλησίαι τῆς Ἀσίας. Ἀσπάζεται ὑμᾶς ἐν Κυρίῳ πολλὰ Ἀκύλας καὶ Πρίσκα σὺν τῇ κατ᾽ οἶκον αὐτῶν ἐκκλησίᾳ.
\vs{20}Ἀσπάζονται ὑμᾶς οἱ ἀδελφοὶ πάντες. Ἀσπάσασθε ἀλλήλους ἐν φιλήματι ἁγίῳ.

\vs{21}Ὁ ἀσπασμὸς τῇ ἐμῇ χειρὶ Παύλου.
\vs{22}Εἴ τις οὐ φιλεῖ τὸν Κύριον, ἤτω ἀνάθεμα. Μαράνα θά.
\vs{23}Ἡ χάρις τοῦ Κυρίου Ἰησοῦ μεθ᾽ ὑμῶν.
\vs{24}Ἡ ἀγάπη μου μετὰ πάντων ὑμῶν ἐν Χριστῷ Ἰησοῦ.


\def\book{ΠΡΟΣ ΚΟΡΙΝΘΙΟΥΣ Β}
\biblebook{ΠΡΟΣ ΚΟΡΙΝΘΙΟΥΣ Β}


\lettrine[lines=2, loversize=0.2, nindent=0em, findent=.25em]{\textcolor{bookheadingcolor}{Π}}{αῦλος} ἀπόστολος Χριστοῦ Ἰησοῦ διὰ θελήματος Θεοῦ καὶ Τιμόθεος ὁ ἀδελφὸς Τῇ ἐκκλησίᾳ τοῦ Θεοῦ τῇ οὔσῃ ἐν Κορίνθῳ σὺν τοῖς ἁγίοις πᾶσιν τοῖς οὖσιν ἐν ὅλῃ τῇ Ἀχαΐᾳ,
\vs{2}Χάρις ὑμῖν καὶ εἰρήνη ἀπὸ Θεοῦ Πατρὸς ἡμῶν καὶ Κυρίου Ἰησοῦ Χριστοῦ.

\vs{3}Εὐλογητὸς ὁ Θεὸς καὶ Πατὴρ τοῦ Κυρίου ἡμῶν Ἰησοῦ Χριστοῦ, ὁ Πατὴρ τῶν οἰκτιρμῶν καὶ Θεὸς πάσης παρακλήσεως,
\vs{4}ὁ παρακαλῶν ἡμᾶς ἐπὶ πάσῃ τῇ θλίψει ἡμῶν εἰς τὸ δύνασθαι ἡμᾶς παρακαλεῖν τοὺς ἐν πάσῃ θλίψει διὰ τῆς παρακλήσεως ἧς παρακαλούμεθα αὐτοὶ ὑπὸ τοῦ Θεοῦ.
\vs{5}ὅτι καθὼς περισσεύει τὰ παθήματα τοῦ Χριστοῦ εἰς ἡμᾶς, οὕτως διὰ τοῦ Χριστοῦ περισσεύει καὶ ἡ παράκλησις ἡμῶν.
\vs{6}Εἴτε δὲ θλιβόμεθα, ὑπὲρ τῆς ὑμῶν παρακλήσεως καὶ σωτηρίας· εἴτε παρακαλούμεθα, ὑπὲρ τῆς ὑμῶν παρακλήσεως τῆς ἐνεργουμένης ἐν ὑπομονῇ τῶν αὐτῶν παθημάτων ὧν καὶ ἡμεῖς πάσχομεν.
\vs{7}καὶ ἡ ἐλπὶς ἡμῶν βεβαία ὑπὲρ ὑμῶν εἰδότες ὅτι ὡς κοινωνοί ἐστε τῶν παθημάτων, οὕτως καὶ τῆς παρακλήσεως.

\vs{8}Οὐ γὰρ θέλομεν ὑμᾶς ἀγνοεῖν, ἀδελφοί, ὑπὲρ τῆς θλίψεως ἡμῶν τῆς γενομένης ἐν τῇ Ἀσίᾳ, ὅτι καθ᾽ ὑπερβολὴν ὑπὲρ δύναμιν ἐβαρήθημεν ὥστε ἐξαπορηθῆναι ἡμᾶς καὶ τοῦ ζῆν·
\vs{9}ἀλλὰ αὐτοὶ ἐν ἑαυτοῖς τὸ ἀπόκριμα τοῦ θανάτου ἐσχήκαμεν, ἵνα μὴ πεποιθότες ὦμεν ἐφ᾽ ἑαυτοῖς ἀλλ᾽ ἐπὶ τῷ Θεῷ τῷ ἐγείροντι τοὺς νεκρούς·
\vs{10}ὃς ἐκ τηλικούτου θανάτου ἐρρύσατο ἡμᾶς καὶ ῥύσεται, εἰς ὃν ἠλπίκαμεν ὅτι καὶ ἔτι ῥύσεται,
\vs{11}συνυπουργούντων καὶ ὑμῶν ὑπὲρ ἡμῶν τῇ δεήσει, ἵνα ἐκ πολλῶν προσώπων τὸ εἰς ἡμᾶς χάρισμα διὰ πολλῶν εὐχαριστηθῇ ὑπὲρ ἡμῶν.

\vs{12}Ἡ γὰρ καύχησις ἡμῶν αὕτη ἐστίν, τὸ μαρτύριον τῆς συνειδήσεως ἡμῶν, ὅτι ἐν ἁγιότητι καὶ εἰλικρινείᾳ τοῦ Θεοῦ, καὶ οὐκ ἐν σοφίᾳ σαρκικῇ ἀλλ᾽ ἐν χάριτι Θεοῦ, ἀνεστράφημεν ἐν τῷ κόσμῳ, περισσοτέρως δὲ πρὸς ὑμᾶς.
\vs{13}οὐ γὰρ ἄλλα γράφομεν ὑμῖν ἀλλ᾽ ἢ ἃ ἀναγινώσκετε ἢ καὶ ἐπιγινώσκετε· ἐλπίζω δὲ ὅτι ἕως τέλους ἐπιγνώσεσθε,
\vs{14}καθὼς καὶ ἐπέγνωτε ἡμᾶς ἀπὸ μέρους, ὅτι καύχημα ὑμῶν ἐσμεν καθάπερ καὶ ὑμεῖς ἡμῶν ἐν τῇ ἡμέρᾳ τοῦ Κυρίου ἡμῶν Ἰησοῦ.

\vs{15}Καὶ ταύτῃ τῇ πεποιθήσει ἐβουλόμην πρότερον πρὸς ὑμᾶς ἐλθεῖν, ἵνα δευτέραν χάριν σχῆτε,
\vs{16}καὶ δι᾽ ὑμῶν διελθεῖν εἰς Μακεδονίαν καὶ πάλιν ἀπὸ Μακεδονίας ἐλθεῖν πρὸς ὑμᾶς καὶ ὑφ᾽ ὑμῶν προπεμφθῆναι εἰς τὴν Ἰουδαίαν.
\vs{17}Τοῦτο οὖν βουλόμενος μήτι ἄρα τῇ ἐλαφρίᾳ ἐχρησάμην; ἢ ἃ βουλεύομαι κατὰ σάρκα βουλεύομαι, ἵνα ᾖ παρ᾽ ἐμοὶ τό Ναί ναὶ καὶ τὸ Οὔ οὔ;
\vs{18}πιστὸς δὲ ὁ Θεὸς ὅτι ὁ λόγος ἡμῶν ὁ πρὸς ὑμᾶς οὐκ ἔστιν Ναί καὶ Οὔ.
\vs{19}ὁ τοῦ Θεοῦ γὰρ Υἱὸς Ἰησοῦς Χριστὸς ὁ ἐν ὑμῖν δι᾽ ἡμῶν κηρυχθείς, δι᾽ ἐμοῦ καὶ Σιλουανοῦ καὶ Τιμοθέου, οὐκ ἐγένετο Ναί καὶ Οὔ ἀλλὰ Ναί ἐν αὐτῷ γέγονεν.
\vs{20}ὅσαι γὰρ ἐπαγγελίαι Θεοῦ, ἐν αὐτῷ τὸ Ναί· διὸ καὶ δι᾽ αὐτοῦ τὸ Ἀμὴν τῷ Θεῷ πρὸς δόξαν δι᾽ ἡμῶν.
\vs{21}Ὁ δὲ βεβαιῶν ἡμᾶς σὺν ὑμῖν εἰς Χριστὸν καὶ χρίσας ἡμᾶς Θεός,
\vs{22}ὁ καὶ σφραγισάμενος ἡμᾶς καὶ δοὺς τὸν ἀρραβῶνα τοῦ Πνεύματος ἐν ταῖς καρδίαις ἡμῶν.

\vs{23}Ἐγὼ δὲ μάρτυρα τὸν Θεὸν ἐπικαλοῦμαι ἐπὶ τὴν ἐμὴν ψυχήν, ὅτι φειδόμενος ὑμῶν οὐκέτι ἦλθον εἰς Κόρινθον.
\vs{24}οὐχ ὅτι κυριεύομεν ὑμῶν τῆς πίστεως ἀλλὰ συνεργοί ἐσμεν τῆς χαρᾶς ὑμῶν· τῇ γὰρ πίστει ἑστήκατε.
\inparch{2} \vs{1}Ἔκρινα γὰρ ἐμαυτῷ τοῦτο τὸ μὴ πάλιν ἐν λύπῃ πρὸς ὑμᾶς ἐλθεῖν.
\vs{2}εἰ γὰρ ἐγὼ λυπῶ ὑμᾶς, καὶ τίς ὁ εὐφραίνων με εἰ μὴ ὁ λυπούμενος ἐξ ἐμοῦ;
\vs{3}καὶ ἔγραψα τοῦτο αὐτὸ, ἵνα μὴ ἐλθὼν λύπην σχῶ ἀφ᾽ ὧν ἔδει με χαίρειν, πεποιθὼς ἐπὶ πάντας ὑμᾶς ὅτι ἡ ἐμὴ χαρὰ πάντων ὑμῶν ἐστιν.
\vs{4}ἐκ γὰρ πολλῆς θλίψεως καὶ συνοχῆς καρδίας ἔγραψα ὑμῖν διὰ πολλῶν δακρύων, οὐχ ἵνα λυπηθῆτε ἀλλὰ τὴν ἀγάπην ἵνα γνῶτε ἣν ἔχω περισσοτέρως εἰς ὑμᾶς.

\vs{5}Εἰ δέ τις λελύπηκεν, οὐκ ἐμὲ λελύπηκεν, ἀλλὰ ἀπὸ μέρους, ἵνα μὴ ἐπιβαρῶ, πάντας ὑμᾶς.
\vs{6}ἱκανὸν τῷ τοιούτῳ ἡ ἐπιτιμία αὕτη ἡ ὑπὸ τῶν πλειόνων,
\vs{7}ὥστε τοὐναντίον μᾶλλον ὑμᾶς χαρίσασθαι καὶ παρακαλέσαι, μή πως τῇ περισσοτέρᾳ λύπῃ καταποθῇ ὁ τοιοῦτος.
\vs{8}διὸ παρακαλῶ ὑμᾶς κυρῶσαι εἰς αὐτὸν ἀγάπην·
\vs{9}Εἰς τοῦτο γὰρ καὶ ἔγραψα, ἵνα γνῶ τὴν δοκιμὴν ὑμῶν, εἰ εἰς πάντα ὑπήκοοί ἐστε.
\vs{10}ᾧ δέ τι χαρίζεσθε, κἀγώ· καὶ γὰρ ἐγὼ ὃ κεχάρισμαι, εἴ τι κεχάρισμαι, δι᾽ ὑμᾶς ἐν προσώπῳ Χριστοῦ,
\vs{11}ἵνα μὴ πλεονεκτηθῶμεν ὑπὸ τοῦ Σατανᾶ· οὐ γὰρ αὐτοῦ τὰ νοήματα ἀγνοοῦμεν.

\vs{12}Ἐλθὼν δὲ εἰς τὴν Τρῳάδα εἰς τὸ εὐαγγέλιον τοῦ Χριστοῦ καὶ θύρας μοι ἀνεῳγμένης ἐν Κυρίῳ,
\vs{13}οὐκ ἔσχηκα ἄνεσιν τῷ πνεύματί μου τῷ μὴ εὑρεῖν με Τίτον τὸν ἀδελφόν μου, ἀλλὰ ἀποταξάμενος αὐτοῖς ἐξῆλθον εἰς Μακεδονίαν.

\vs{14}Τῷ δὲ Θεῷ χάρις τῷ πάντοτε θριαμβεύοντι ἡμᾶς ἐν τῷ Χριστῷ καὶ τὴν ὀσμὴν τῆς γνώσεως αὐτοῦ φανεροῦντι δι᾽ ἡμῶν ἐν παντὶ τόπῳ·
\vs{15}ὅτι Χριστοῦ εὐωδία ἐσμὲν τῷ Θεῷ ἐν τοῖς σωζομένοις καὶ ἐν τοῖς ἀπολλυμένοις,
\vs{16}οἷς μὲν ὀσμὴ ἐκ θανάτου εἰς θάνατον, οἷς δὲ ὀσμὴ ἐκ ζωῆς εἰς ζωήν. καὶ πρὸς ταῦτα τίς ἱκανός;
\vs{17}Οὐ γάρ ἐσμεν ὡς οἱ πολλοὶ καπηλεύοντες τὸν λόγον τοῦ Θεοῦ, ἀλλ᾽ ὡς ἐξ εἰλικρινείας, ἀλλ᾽ ὡς ἐκ Θεοῦ κατέναντι Θεοῦ ἐν Χριστῷ λαλοῦμεν.

\ch{3}
Ἀρχόμεθα πάλιν ἑαυτοὺς συνιστάνειν; ἢ μὴ χρῄζομεν ὥς τινες συστατικῶν ἐπιστολῶν πρὸς ὑμᾶς ἢ ἐξ ὑμῶν;
\vs{2}ἡ ἐπιστολὴ ἡμῶν ὑμεῖς ἐστε, ἐνγεγραμμένη ἐν ταῖς καρδίαις ἡμῶν, γινωσκομένη καὶ ἀναγινωσκομένη ὑπὸ πάντων ἀνθρώπων,
\vs{3}φανερούμενοι ὅτι ἐστὲ ἐπιστολὴ Χριστοῦ διακονηθεῖσα ὑφ᾽ ἡμῶν, ἐνγεγραμμένη οὐ μέλανι ἀλλὰ Πνεύματι Θεοῦ ζῶντος, οὐκ ἐν πλαξὶν λιθίναις ἀλλ᾽ ἐν πλαξὶν καρδίαις σαρκίναις.

\vs{4}Πεποίθησιν δὲ τοιαύτην ἔχομεν διὰ τοῦ Χριστοῦ πρὸς τὸν Θεόν.
\vs{5}οὐχ ὅτι ἀφ᾽ ἑαυτῶν ἱκανοί ἐσμεν λογίσασθαί τι ὡς ἐξ ἑαυτῶν, ἀλλ᾽ ἡ ἱκανότης ἡμῶν ἐκ τοῦ Θεοῦ,
\vs{6}ὃς καὶ ἱκάνωσεν ἡμᾶς διακόνους καινῆς διαθήκης, οὐ γράμματος ἀλλὰ πνεύματος· τὸ γὰρ γράμμα ἀποκτέννει, τὸ δὲ πνεῦμα ζωοποιεῖ.
\vs{7}Εἰ δὲ ἡ διακονία τοῦ θανάτου ἐν γράμμασιν ἐντετυπωμένη λίθοις ἐγενήθη ἐν δόξῃ, ὥστε μὴ δύνασθαι ἀτενίσαι τοὺς υἱοὺς Ἰσραὴλ εἰς τὸ πρόσωπον Μωϋσέως διὰ τὴν δόξαν τοῦ προσώπου αὐτοῦ τὴν καταργουμένην,
\vs{8}πῶς οὐχὶ μᾶλλον ἡ διακονία τοῦ πνεύματος ἔσται ἐν δόξῃ;
\vs{9}εἰ γὰρ τῇ διακονία τῆς κατακρίσεως δόξα, πολλῷ μᾶλλον περισσεύει ἡ διακονία τῆς δικαιοσύνης δόξῃ.
\vs{10}καὶ γὰρ οὐ δεδόξασται τὸ δεδοξασμένον ἐν τούτῳ τῷ μέρει εἵνεκεν τῆς ὑπερβαλλούσης δόξης.
\vs{11}εἰ γὰρ τὸ καταργούμενον διὰ δόξης, πολλῷ μᾶλλον τὸ μένον ἐν δόξῃ.

\vs{12}Ἔχοντες οὖν τοιαύτην ἐλπίδα πολλῇ παρρησίᾳ χρώμεθα
\vs{13}καὶ οὐ καθάπερ Μωϋσῆς ἐτίθει κάλυμμα ἐπὶ τὸ πρόσωπον αὐτοῦ πρὸς τὸ μὴ ἀτενίσαι τοὺς υἱοὺς Ἰσραὴλ εἰς τὸ τέλος τοῦ καταργουμένου.
\vs{14}Ἀλλὰ ἐπωρώθη τὰ νοήματα αὐτῶν. ἄχρι γὰρ τῆς σήμερον ἡμέρας τὸ αὐτὸ κάλυμμα ἐπὶ τῇ ἀναγνώσει τῆς παλαιᾶς διαθήκης μένει, μὴ ἀνακαλυπτόμενον ὅτι ἐν Χριστῷ καταργεῖται·
\vs{15}ἀλλ᾽ ἕως σήμερον ἡνίκα ἂν ἀναγινώσκηται Μωϋσῆς, κάλυμμα ἐπὶ τὴν καρδίαν αὐτῶν κεῖται·
\vs{16}ἡνίκα δὲ ἐὰν ἐπιστρέψῃ πρὸς Κύριον, περιαιρεῖται τὸ κάλυμμα.
\vs{17}Ὁ δὲ Κύριος τὸ Πνεῦμά ἐστιν· οὗ δὲ τὸ Πνεῦμα Κυρίου, ἐλευθερία.
\vs{18}ἡμεῖς δὲ πάντες ἀνακεκαλυμμένῳ προσώπῳ τὴν δόξαν Κυρίου κατοπτριζόμενοι τὴν αὐτὴν εἰκόνα μεταμορφούμεθα ἀπὸ δόξης εἰς δόξαν καθάπερ ἀπὸ Κυρίου Πνεύματος.

\ch{4}
Διὰ τοῦτο, ἔχοντες τὴν διακονίαν ταύτην καθὼς ἠλεήθημεν, οὐκ ἐγκακοῦμεν
\vs{2}ἀλλὰ ἀπειπάμεθα τὰ κρυπτὰ τῆς αἰσχύνης, μὴ περιπατοῦντες ἐν πανουργίᾳ μηδὲ δολοῦντες τὸν λόγον τοῦ Θεοῦ ἀλλὰ τῇ φανερώσει τῆς ἀληθείας συνιστάνοντες ἑαυτοὺς πρὸς πᾶσαν συνείδησιν ἀνθρώπων ἐνώπιον τοῦ Θεοῦ.
\vs{3}Εἰ δὲ καὶ ἔστιν κεκαλυμμένον τὸ εὐαγγέλιον ἡμῶν, ἐν τοῖς ἀπολλυμένοις ἐστὶν κεκαλυμμένον,
\vs{4}ἐν οἷς ὁ θεὸς τοῦ αἰῶνος τούτου ἐτύφλωσεν τὰ νοήματα τῶν ἀπίστων εἰς τὸ μὴ αὐγάσαι τὸν φωτισμὸν τοῦ εὐαγγελίου τῆς δόξης τοῦ Χριστοῦ, ὅς ἐστιν εἰκὼν τοῦ Θεοῦ.
\vs{5}Οὐ γὰρ ἑαυτοὺς κηρύσσομεν ἀλλὰ Ἰησοῦν Χριστὸν Κύριον, ἑαυτοὺς δὲ δούλους ὑμῶν διὰ Ἰησοῦν.
\vs{6}ὅτι ὁ Θεὸς ὁ εἰπών· Ἐκ σκότους φῶς λάμψει, ὃς ἔλαμψεν ἐν ταῖς καρδίαις ἡμῶν πρὸς φωτισμὸν τῆς γνώσεως τῆς δόξης τοῦ Θεοῦ ἐν προσώπῳ Ἰησοῦ Χριστοῦ.

\vs{7}Ἔχομεν δὲ τὸν θησαυρὸν τοῦτον ἐν ὀστρακίνοις σκεύεσιν, ἵνα ἡ ὑπερβολὴ τῆς δυνάμεως ᾖ τοῦ Θεοῦ καὶ μὴ ἐξ ἡμῶν·
\vs{8}ἐν παντὶ θλιβόμενοι ἀλλ᾽ οὐ στενοχωρούμενοι, ἀπορούμενοι ἀλλ᾽ οὐκ ἐξαπορούμενοι,
\vs{9}διωκόμενοι ἀλλ᾽ οὐκ ἐγκαταλειπόμενοι, καταβαλλόμενοι ἀλλ᾽ οὐκ ἀπολλύμενοι,
\vs{10}πάντοτε τὴν νέκρωσιν τοῦ Ἰησοῦ ἐν τῷ σώματι περιφέροντες, ἵνα καὶ ἡ ζωὴ τοῦ Ἰησοῦ ἐν τῷ σώματι ἡμῶν φανερωθῇ.
\vs{11}ἀεὶ γὰρ ἡμεῖς οἱ ζῶντες εἰς θάνατον παραδιδόμεθα διὰ Ἰησοῦν, ἵνα καὶ ἡ ζωὴ τοῦ Ἰησοῦ φανερωθῇ ἐν τῇ θνητῇ σαρκὶ ἡμῶν.
\vs{12}ὥστε ὁ θάνατος ἐν ἡμῖν ἐνεργεῖται, ἡ δὲ ζωὴ ἐν ὑμῖν.
\vs{13}Ἔχοντες δὲ τὸ αὐτὸ πνεῦμα τῆς πίστεως κατὰ τὸ γεγραμμένον· Ἐπίστευσα, διὸ ἐλάλησα, καὶ ἡμεῖς πιστεύομεν, διὸ καὶ λαλοῦμεν,
\vs{14}εἰδότες ὅτι ὁ ἐγείρας τὸν Κύριον Ἰησοῦν καὶ ἡμᾶς σὺν Ἰησοῦ ἐγερεῖ καὶ παραστήσει σὺν ὑμῖν.
\vs{15}τὰ γὰρ πάντα δι᾽ ὑμᾶς, ἵνα ἡ χάρις πλεονάσασα διὰ τῶν πλειόνων τὴν εὐχαριστίαν περισσεύσῃ εἰς τὴν δόξαν τοῦ Θεοῦ.

\vs{16}Διὸ οὐκ ἐγκακοῦμεν, ἀλλ᾽ εἰ καὶ ὁ ἔξω ἡμῶν ἄνθρωπος διαφθείρεται, ἀλλ᾽ ὁ ἔσω ἡμῶν ἀνακαινοῦται ἡμέρᾳ καὶ ἡμέρᾳ.
\vs{17}τὸ γὰρ παραυτίκα ἐλαφρὸν τῆς θλίψεως ἡμῶν καθ᾽ ὑπερβολὴν εἰς ὑπερβολὴν αἰώνιον βάρος δόξης κατεργάζεται ἡμῖν,
\vs{18}μὴ σκοπούντων ἡμῶν τὰ βλεπόμενα ἀλλὰ τὰ μὴ βλεπόμενα· τὰ γὰρ βλεπόμενα πρόσκαιρα, τὰ δὲ μὴ βλεπόμενα αἰώνια.

\ch{5}
Οἴδαμεν γὰρ ὅτι ἐὰν ἡ ἐπίγειος ἡμῶν οἰκία τοῦ σκήνους καταλυθῇ, οἰκοδομὴν ἐκ Θεοῦ ἔχομεν, οἰκίαν ἀχειροποίητον αἰώνιον ἐν τοῖς οὐρανοῖς.
\vs{2}καὶ γὰρ ἐν τούτῳ στενάζομεν τὸ οἰκητήριον ἡμῶν τὸ ἐξ οὐρανοῦ ἐπενδύσασθαι ἐπιποθοῦντες,
\vs{3}εἴ γε καὶ ἐνδυσάμενοι οὐ γυμνοὶ εὑρεθησόμεθα.
\vs{4}καὶ γὰρ οἱ ὄντες ἐν τῷ σκήνει στενάζομεν βαρούμενοι, ἐφ᾽ ᾧ οὐ θέλομεν ἐκδύσασθαι ἀλλ᾽ ἐπενδύσασθαι, ἵνα καταποθῇ τὸ θνητὸν ὑπὸ τῆς ζωῆς.
\vs{5}ὁ δὲ κατεργασάμενος ἡμᾶς εἰς αὐτὸ τοῦτο Θεός, ὁ δοὺς ἡμῖν τὸν ἀρραβῶνα τοῦ Πνεύματος.
\vs{6}Θαρροῦντες οὖν πάντοτε καὶ εἰδότες ὅτι ἐνδημοῦντες ἐν τῷ σώματι ἐκδημοῦμεν ἀπὸ τοῦ Κυρίου·
\vs{7}διὰ πίστεως γὰρ περιπατοῦμεν, οὐ διὰ εἴδους·
\vs{8}Θαρροῦμεν δὲ καὶ εὐδοκοῦμεν μᾶλλον ἐκδημῆσαι ἐκ τοῦ σώματος καὶ ἐνδημῆσαι πρὸς τὸν Κύριον.
\vs{9}διὸ καὶ φιλοτιμούμεθα, εἴτε ἐνδημοῦντες εἴτε ἐκδημοῦντες, εὐάρεστοι αὐτῷ εἶναι.
\vs{10}τοὺς γὰρ πάντας ἡμᾶς φανερωθῆναι δεῖ ἔμπροσθεν τοῦ βήματος τοῦ Χριστοῦ, ἵνα κομίσηται ἕκαστος τὰ διὰ τοῦ σώματος πρὸς ἃ ἔπραξεν, εἴτε ἀγαθὸν εἴτε φαῦλον.

\vs{11}Εἰδότες οὖν τὸν φόβον τοῦ Κυρίου ἀνθρώπους πείθομεν, Θεῷ δὲ πεφανερώμεθα· ἐλπίζω δὲ καὶ ἐν ταῖς συνειδήσεσιν ὑμῶν πεφανερῶσθαι.
\vs{12}οὐ πάλιν ἑαυτοὺς συνιστάνομεν ὑμῖν ἀλλὰ ἀφορμὴν διδόντες ὑμῖν καυχήματος ὑπὲρ ἡμῶν, ἵνα ἔχητε πρὸς τοὺς ἐν προσώπῳ καυχωμένους καὶ μὴ ἐν καρδίᾳ.
\vs{13}Εἴτε γὰρ ἐξέστημεν, Θεῷ· εἴτε σωφρονοῦμεν, ὑμῖν.
\vs{14}ἡ γὰρ ἀγάπη τοῦ Χριστοῦ συνέχει ἡμᾶς, κρίναντας τοῦτο, ὅτι εἷς ὑπὲρ πάντων ἀπέθανεν, ἄρα οἱ πάντες ἀπέθανον·
\vs{15}καὶ ὑπὲρ πάντων ἀπέθανεν, ἵνα οἱ ζῶντες μηκέτι ἑαυτοῖς ζῶσιν ἀλλὰ τῷ ὑπὲρ αὐτῶν ἀποθανόντι καὶ ἐγερθέντι.
\vs{16}Ὥστε ἡμεῖς ἀπὸ τοῦ νῦν οὐδένα οἴδαμεν κατὰ σάρκα· εἰ καὶ ἐγνώκαμεν κατὰ σάρκα Χριστόν, ἀλλὰ νῦν οὐκέτι γινώσκομεν.
\vs{17}ὥστε εἴ τις ἐν Χριστῷ, καινὴ κτίσις· τὰ ἀρχαῖα παρῆλθεν, ἰδοὺ γέγονεν καινά.
\vs{18}Τὰ δὲ πάντα ἐκ τοῦ Θεοῦ τοῦ καταλλάξαντος ἡμᾶς ἑαυτῷ διὰ Χριστοῦ καὶ δόντος ἡμῖν τὴν διακονίαν τῆς καταλλαγῆς,
\vs{19}ὡς ὅτι Θεὸς ἦν ἐν Χριστῷ κόσμον καταλλάσσων ἑαυτῷ, μὴ λογιζόμενος αὐτοῖς τὰ παραπτώματα αὐτῶν καὶ θέμενος ἐν ἡμῖν τὸν λόγον τῆς καταλλαγῆς.
\vs{20}Ὑπὲρ Χριστοῦ οὖν πρεσβεύομεν ὡς τοῦ Θεοῦ παρακαλοῦντος δι᾽ ἡμῶν· δεόμεθα ὑπὲρ Χριστοῦ, καταλλάγητε τῷ Θεῷ.
\vs{21}τὸν μὴ γνόντα ἁμαρτίαν ὑπὲρ ἡμῶν ἁμαρτίαν ἐποίησεν, ἵνα ἡμεῖς γενώμεθα δικαιοσύνη Θεοῦ ἐν αὐτῷ.

\ch{6}
Συνεργοῦντες δὲ καὶ παρακαλοῦμεν μὴ εἰς κενὸν τὴν χάριν τοῦ Θεοῦ δέξασθαι ὑμᾶς·
\vs{2}λέγει γάρ· 
\begin{poetryblock}

\begin{quote}Καιρῷ δεκτῷ ἐπήκουσά σου\end{quote} 

\begin{quote}καὶ ἐν ἡμέρᾳ σωτηρίας ἐβοήθησά σοι.\end{quote}
\end{poetryblock}

Ἰδοὺ νῦν καιρὸς εὐπρόσδεκτος, ἰδοὺ νῦν ἡμέρα σωτηρίας.

\vs{3}Μηδεμίαν ἐν μηδενὶ διδόντες προσκοπήν, ἵνα μὴ μωμηθῇ ἡ διακονία,
\vs{4}ἀλλ᾽ ἐν παντὶ συνιστάντες ἑαυτοὺς ὡς Θεοῦ διάκονοι, ἐν ὑπομονῇ πολλῇ, ἐν θλίψεσιν, ἐν ἀνάγκαις, ἐν στενοχωρίαις,
\vs{5}ἐν πληγαῖς, ἐν φυλακαῖς, ἐν ἀκαταστασίαις, ἐν κόποις, ἐν ἀγρυπνίαις, ἐν νηστείαις,
\vs{6}ἐν ἁγνότητι, ἐν γνώσει, ἐν μακροθυμίᾳ, ἐν χρηστότητι, ἐν Πνεύματι Ἁγίῳ, ἐν ἀγάπῃ ἀνυποκρίτῳ,
\vs{7}ἐν λόγῳ ἀληθείας, ἐν δυνάμει Θεοῦ· διὰ τῶν ὅπλων τῆς δικαιοσύνης τῶν δεξιῶν καὶ ἀριστερῶν,
\vs{8}διὰ δόξης καὶ ἀτιμίας, διὰ δυσφημίας καὶ εὐφημίας· ὡς πλάνοι καὶ ἀληθεῖς,
\vs{9}ὡς ἀγνοούμενοι καὶ ἐπιγινωσκόμενοι, ὡς ἀποθνήσκοντες καὶ ἰδοὺ ζῶμεν, ὡς παιδευόμενοι καὶ μὴ θανατούμενοι,
\vs{10}ὡς λυπούμενοι ἀεὶ δὲ χαίροντες, ὡς πτωχοὶ πολλοὺς δὲ πλουτίζοντες, ὡς μηδὲν ἔχοντες καὶ πάντα κατέχοντες.

\vs{11}Τὸ στόμα ἡμῶν ἀνέῳγεν πρὸς ὑμᾶς, Κορίνθιοι, ἡ καρδία ἡμῶν πεπλάτυνται·
\vs{12}οὐ στενοχωρεῖσθε ἐν ἡμῖν, στενοχωρεῖσθε δὲ ἐν τοῖς σπλάγχνοις ὑμῶν·
\vs{13}τὴν δὲ αὐτὴν ἀντιμισθίαν, ὡς τέκνοις λέγω, πλατύνθητε καὶ ὑμεῖς.

\vs{14}Μὴ γίνεσθε ἑτεροζυγοῦντες ἀπίστοις· τίς γὰρ μετοχὴ δικαιοσύνῃ καὶ ἀνομίᾳ, ἢ τίς κοινωνία φωτὶ πρὸς σκότος;
\vs{15}τίς δὲ συμφώνησις Χριστοῦ πρὸς Βελιάρ, ἢ τίς μερὶς πιστῷ μετὰ ἀπίστου;
\vs{16}τίς δὲ συνκατάθεσις ναῷ Θεοῦ μετὰ εἰδώλων; ἡμεῖς γὰρ ναὸς Θεοῦ ἐσμεν ζῶντος, καθὼς εἶπεν ὁ Θεὸς ὅτι 
\begin{poetryblock}

\begin{quote}Ἐνοικήσω ἐν αὐτοῖς καὶ ἐμπεριπατήσω\end{quote} 

\begin{quote}καὶ ἔσομαι αὐτῶν Θεός καὶ αὐτοὶ ἔσονταί μου λαός.\end{quote}

\begin{quote} \vs{17}Διὸ ἐξέλθατε ἐκ μέσου αὐτῶν\end{quote} 

\begin{quote}καὶ ἀφορίσθητε, λέγει Κύριος,\end{quote} 

\begin{quote}καὶ ἀκαθάρτου μὴ ἅπτεσθε·\end{quote} 

\begin{quote}κἀγὼ εἰσδέξομαι ὑμᾶς\end{quote}

\begin{quote} \vs{18}Καὶ Ἔσομαι ὑμῖν εἰς Πατέρα\end{quote} 

\begin{quote}καὶ ὑμεῖς ἔσεσθέ μοι εἰς υἱοὺς καὶ θυγατέρας,\end{quote} 

\begin{quote}λέγει Κύριος Παντοκράτωρ.\end{quote}
\end{poetryblock}

\ch{7}
Ταύτας οὖν ἔχοντες τὰς ἐπαγγελίας, ἀγαπητοί, καθαρίσωμεν ἑαυτοὺς ἀπὸ παντὸς μολυσμοῦ σαρκὸς καὶ πνεύματος, ἐπιτελοῦντες ἁγιωσύνην ἐν φόβῳ Θεοῦ.

\vs{2}Χωρήσατε ἡμᾶς· οὐδένα ἠδικήσαμεν, οὐδένα ἐφθείραμεν, οὐδένα ἐπλεονεκτήσαμεν.
\vs{3}πρὸς κατάκρισιν οὐ λέγω· προείρηκα γὰρ ὅτι ἐν ταῖς καρδίαις ἡμῶν ἐστε εἰς τὸ συναποθανεῖν καὶ συζῆν.
\vs{4}πολλή μοι παρρησία πρὸς ὑμᾶς, πολλή μοι καύχησις ὑπὲρ ὑμῶν· πεπλήρωμαι τῇ παρακλήσει, ὑπερπερισσεύομαι τῇ χαρᾷ ἐπὶ πάσῃ τῇ θλίψει ἡμῶν.

\vs{5}Καὶ γὰρ ἐλθόντων ἡμῶν εἰς Μακεδονίαν οὐδεμίαν ἔσχηκεν ἄνεσιν ἡ σὰρξ ἡμῶν ἀλλ᾽ ἐν παντὶ θλιβόμενοι· ἔξωθεν μάχαι, ἔσωθεν φόβοι.
\vs{6}ἀλλ᾽ ὁ παρακαλῶν τοὺς ταπεινοὺς παρεκάλεσεν ἡμᾶς ὁ Θεὸς ἐν τῇ παρουσίᾳ Τίτου,
\vs{7}οὐ μόνον δὲ ἐν τῇ παρουσίᾳ αὐτοῦ ἀλλὰ καὶ ἐν τῇ παρακλήσει ᾗ παρεκλήθη ἐφ᾽ ὑμῖν, ἀναγγέλλων ἡμῖν τὴν ὑμῶν ἐπιπόθησιν, τὸν ὑμῶν ὀδυρμόν, τὸν ὑμῶν ζῆλον ὑπὲρ ἐμοῦ ὥστε με μᾶλλον χαρῆναι.
\vs{8}Ὅτι εἰ καὶ ἐλύπησα ὑμᾶς ἐν τῇ ἐπιστολῇ, οὐ μεταμέλομαι· εἰ καὶ μετεμελόμην, βλέπω γὰρ ὅτι ἡ ἐπιστολὴ ἐκείνη εἰ καὶ πρὸς ὥραν ἐλύπησεν ὑμᾶς,
\vs{9}νῦν χαίρω, οὐχ ὅτι ἐλυπήθητε ἀλλ᾽ ὅτι ἐλυπήθητε εἰς μετάνοιαν· ἐλυπήθητε γὰρ κατὰ Θεόν, ἵνα ἐν μηδενὶ ζημιωθῆτε ἐξ ἡμῶν.
\vs{10}ἡ γὰρ κατὰ Θεὸν λύπη μετάνοιαν εἰς σωτηρίαν ἀμεταμέλητον ἐργάζεται· ἡ δὲ τοῦ κόσμου λύπη θάνατον κατεργάζεται.
\vs{11}Ἰδοὺ γὰρ αὐτὸ τοῦτο τὸ κατὰ Θεὸν λυπηθῆναι πόσην κατειργάσατο ὑμῖν σπουδήν, ἀλλὰ ἀπολογίαν, ἀλλὰ ἀγανάκτησιν, ἀλλὰ φόβον, ἀλλὰ ἐπιπόθησιν, ἀλλὰ ζῆλον, ἀλλὰ ἐκδίκησιν. ἐν παντὶ συνεστήσατε ἑαυτοὺς ἁγνοὺς εἶναι τῷ πράγματι.
\vs{12}ἄρα εἰ καὶ ἔγραψα ὑμῖν, οὐχ ἕνεκεν τοῦ ἀδικήσαντος οὐδὲ ἕνεκεν τοῦ ἀδικηθέντος ἀλλ᾽ ἕνεκεν τοῦ φανερωθῆναι τὴν σπουδὴν ὑμῶν τὴν ὑπὲρ ἡμῶν πρὸς ὑμᾶς ἐνώπιον τοῦ Θεοῦ.
\vs{13}διὰ τοῦτο παρακεκλήμεθα. Ἐπὶ δὲ τῇ παρακλήσει ἡμῶν περισσοτέρως μᾶλλον ἐχάρημεν ἐπὶ τῇ χαρᾷ Τίτου, ὅτι ἀναπέπαυται τὸ πνεῦμα αὐτοῦ ἀπὸ πάντων ὑμῶν·
\vs{14}ὅτι εἴ τι αὐτῷ ὑπὲρ ὑμῶν κεκαύχημαι, οὐ κατῃσχύνθην, ἀλλ᾽ ὡς πάντα ἐν ἀληθείᾳ ἐλαλήσαμεν ὑμῖν, οὕτως καὶ ἡ καύχησις ἡμῶν ἡ ἐπὶ Τίτου ἀλήθεια ἐγενήθη.
\vs{15}καὶ τὰ σπλάγχνα αὐτοῦ περισσοτέρως εἰς ὑμᾶς ἐστιν ἀναμιμνῃσκομένου τὴν πάντων ὑμῶν ὑπακοήν, ὡς μετὰ φόβου καὶ τρόμου ἐδέξασθε αὐτόν.
\vs{16}χαίρω ὅτι ἐν παντὶ θαρρῶ ἐν ὑμῖν.

\ch{8}
Γνωρίζομεν δὲ ὑμῖν, ἀδελφοί, τὴν χάριν τοῦ Θεοῦ τὴν δεδομένην ἐν ταῖς ἐκκλησίαις τῆς Μακεδονίας,
\vs{2}ὅτι ἐν πολλῇ δοκιμῇ θλίψεως ἡ περισσεία τῆς χαρᾶς αὐτῶν καὶ ἡ κατὰ βάθους πτωχεία αὐτῶν ἐπερίσσευσεν εἰς τὸ πλοῦτος τῆς ἁπλότητος αὐτῶν·
\vs{3}ὅτι κατὰ δύναμιν, μαρτυρῶ, καὶ παρὰ δύναμιν, αὐθαίρετοι
\vs{4}μετὰ πολλῆς παρακλήσεως δεόμενοι ἡμῶν τὴν χάριν καὶ τὴν κοινωνίαν τῆς διακονίας τῆς εἰς τοὺς ἁγίους,
\vs{5}καὶ οὐ καθὼς ἠλπίσαμεν ἀλλ᾽ ἑαυτοὺς ἔδωκαν πρῶτον τῷ Κυρίῳ καὶ ἡμῖν διὰ θελήματος Θεοῦ
\vs{6}Εἰς τὸ παρακαλέσαι ἡμᾶς Τίτον, ἵνα καθὼς προενήρξατο οὕτως καὶ ἐπιτελέσῃ εἰς ὑμᾶς καὶ τὴν χάριν ταύτην.
\vs{7}ἀλλ᾽ ὥσπερ ἐν παντὶ περισσεύετε, πίστει καὶ λόγῳ καὶ γνώσει καὶ πάσῃ σπουδῇ καὶ τῇ ἐξ ἡμῶν ἐν ὑμῖν ἀγάπῃ, ἵνα καὶ ἐν ταύτῃ τῇ χάριτι περισσεύητε.
\vs{8}Οὐ κατ᾽ ἐπιταγὴν λέγω ἀλλὰ διὰ τῆς ἑτέρων σπουδῆς καὶ τὸ τῆς ὑμετέρας ἀγάπης γνήσιον δοκιμάζων·
\vs{9}Γινώσκετε γὰρ τὴν χάριν τοῦ Κυρίου ἡμῶν Ἰησοῦ Χριστοῦ, ὅτι δι᾽ ὑμᾶς ἐπτώχευσεν πλούσιος ὤν, ἵνα ὑμεῖς τῇ ἐκείνου πτωχείᾳ πλουτήσητε.
\vs{10}καὶ γνώμην ἐν τούτῳ δίδωμι· τοῦτο γὰρ ὑμῖν συμφέρει, οἵτινες οὐ μόνον τὸ ποιῆσαι ἀλλὰ καὶ τὸ θέλειν προενήρξασθε ἀπὸ πέρυσι·
\vs{11}νυνὶ δὲ καὶ τὸ ποιῆσαι ἐπιτελέσατε, ὅπως καθάπερ ἡ προθυμία τοῦ θέλειν, οὕτως καὶ τὸ ἐπιτελέσαι ἐκ τοῦ ἔχειν.
\vs{12}εἰ γὰρ ἡ προθυμία πρόκειται, καθὸ ἐὰν ἔχῃ εὐπρόσδεκτος, οὐ καθὸ οὐκ ἔχει.
\vs{13}Οὐ γὰρ ἵνα ἄλλοις ἄνεσις, ὑμῖν θλῖψις, ἀλλ᾽ ἐξ ἰσότητος·
\vs{14}ἐν τῷ νῦν καιρῷ τὸ ὑμῶν περίσσευμα εἰς τὸ ἐκείνων ὑστέρημα, ἵνα καὶ τὸ ἐκείνων περίσσευμα γένηται εἰς τὸ ὑμῶν ὑστέρημα, ὅπως γένηται ἰσότης,
\vs{15}καθὼς γέγραπται· Ὁ τὸ πολὺ οὐκ ἐπλεόνασεν, καὶ ὁ τὸ ὀλίγον οὐκ ἠλαττόνησεν.

\vs{16}Χάρις δὲ τῷ Θεῷ τῷ διδόντι τὴν αὐτὴν σπουδὴν ὑπὲρ ὑμῶν ἐν τῇ καρδίᾳ Τίτου,
\vs{17}ὅτι τὴν μὲν παράκλησιν ἐδέξατο, σπουδαιότερος δὲ ὑπάρχων αὐθαίρετος ἐξῆλθεν πρὸς ὑμᾶς.
\vs{18}Συνεπέμψαμεν δὲ μετ᾽ αὐτοῦ τὸν ἀδελφὸν οὗ ὁ ἔπαινος ἐν τῷ εὐαγγελίῳ διὰ πασῶν τῶν ἐκκλησιῶν,
\vs{19}οὐ μόνον δὲ, ἀλλὰ καὶ χειροτονηθεὶς ὑπὸ τῶν ἐκκλησιῶν συνέκδημος ἡμῶν σὺν τῇ χάριτι ταύτῃ τῇ διακονουμένῃ ὑφ᾽ ἡμῶν πρὸς τὴν αὐτοῦ τοῦ Κυρίου δόξαν καὶ προθυμίαν ἡμῶν,
\vs{20}στελλόμενοι τοῦτο, μή τις ἡμᾶς μωμήσηται ἐν τῇ ἁδρότητι ταύτῃ τῇ διακονουμένῃ ὑφ᾽ ἡμῶν·
\vs{21}προνοοῦμεν γὰρ καλὰ οὐ μόνον ἐνώπιον Κυρίου ἀλλὰ καὶ ἐνώπιον ἀνθρώπων.
\vs{22}Συνεπέμψαμεν δὲ αὐτοῖς τὸν ἀδελφὸν ἡμῶν ὃν ἐδοκιμάσαμεν ἐν πολλοῖς πολλάκις σπουδαῖον ὄντα, νυνὶ δὲ πολὺ σπουδαιότερον πεποιθήσει πολλῇ τῇ εἰς ὑμᾶς.
\vs{23}εἴτε ὑπὲρ Τίτου, κοινωνὸς ἐμὸς καὶ εἰς ὑμᾶς συνεργός· εἴτε ἀδελφοὶ ἡμῶν, ἀπόστολοι ἐκκλησιῶν, δόξα Χριστοῦ.
\vs{24}τὴν οὖν ἔνδειξιν τῆς ἀγάπης ὑμῶν καὶ ἡμῶν καυχήσεως ὑπὲρ ὑμῶν εἰς αὐτοὺς ἐνδεικνύμενοι εἰς πρόσωπον τῶν ἐκκλησιῶν.

\ch{9}
Περὶ μὲν γὰρ τῆς διακονίας τῆς εἰς τοὺς ἁγίους περισσόν μοί ἐστιν τὸ γράφειν ὑμῖν·
\vs{2}οἶδα γὰρ τὴν προθυμίαν ὑμῶν ἣν ὑπὲρ ὑμῶν καυχῶμαι Μακεδόσιν, ὅτι Ἀχαΐα παρεσκεύασται ἀπὸ πέρυσι, καὶ τὸ ὑμῶν ζῆλος ἠρέθισεν τοὺς πλείονας.
\vs{3}Ἔπεμψα δὲ τοὺς ἀδελφούς, ἵνα μὴ τὸ καύχημα ἡμῶν τὸ ὑπὲρ ὑμῶν κενωθῇ ἐν τῷ μέρει τούτῳ, ἵνα καθὼς ἔλεγον παρεσκευασμένοι ἦτε,
\vs{4}μή πως ἐὰν ἔλθωσιν σὺν ἐμοὶ Μακεδόνες καὶ εὕρωσιν ὑμᾶς ἀπαρασκευάστους καταισχυνθῶμεν ἡμεῖς, ἵνα μὴ λέγωμεν ὑμεῖς, ἐν τῇ ὑποστάσει ταύτῃ.
\vs{5}ἀναγκαῖον οὖν ἡγησάμην παρακαλέσαι τοὺς ἀδελφοὺς, ἵνα προέλθωσιν εἰς ὑμᾶς καὶ προκαταρτίσωσιν τὴν προεπηγγελμένην εὐλογίαν ὑμῶν, ταύτην ἑτοίμην εἶναι οὕτως ὡς εὐλογίαν καὶ μὴ ὡς πλεονεξίαν.

\vs{6}Τοῦτο δέ, ὁ σπείρων φειδομένως φειδομένως καὶ θερίσει, καὶ ὁ σπείρων ἐπ᾽ εὐλογίαις ἐπ᾽ εὐλογίαις καὶ θερίσει.
\vs{7}ἕκαστος καθὼς προῄρηται τῇ καρδίᾳ, μὴ ἐκ λύπης ἢ ἐξ ἀνάγκης· ἱλαρὸν γὰρ δότην ἀγαπᾷ ὁ Θεός.
\vs{8}δυνατεῖ δὲ ὁ Θεὸς πᾶσαν χάριν περισσεῦσαι εἰς ὑμᾶς, ἵνα ἐν παντὶ πάντοτε πᾶσαν αὐτάρκειαν ἔχοντες περισσεύητε εἰς πᾶν ἔργον ἀγαθόν,
\vs{9}καθὼς γέγραπται· 
\begin{poetryblock}

\begin{quote}Ἐσκόρπισεν, ἔδωκεν τοῖς πένησιν,\end{quote} 

\begin{quote}ἡ δικαιοσύνη αὐτοῦ μένει εἰς τὸν αἰῶνα.\end{quote}
\end{poetryblock}
\vs{10}Ὁ δὲ ἐπιχορηγῶν σπόρον τῷ σπείροντι καὶ ἄρτον εἰς βρῶσιν χορηγήσει καὶ πληθυνεῖ τὸν σπόρον ὑμῶν καὶ αὐξήσει τὰ γενήματα τῆς δικαιοσύνης ὑμῶν.
\vs{11}ἐν παντὶ πλουτιζόμενοι εἰς πᾶσαν ἁπλότητα, ἥτις κατεργάζεται δι᾽ ἡμῶν εὐχαριστίαν τῷ Θεῷ·
\vs{12}ὅτι ἡ διακονία τῆς λειτουργίας ταύτης οὐ μόνον ἐστὶν προσαναπληροῦσα τὰ ὑστερήματα τῶν ἁγίων, ἀλλὰ καὶ περισσεύουσα διὰ πολλῶν εὐχαριστιῶν τῷ Θεῷ.
\vs{13}διὰ τῆς δοκιμῆς τῆς διακονίας ταύτης δοξάζοντες τὸν Θεὸν ἐπὶ τῇ ὑποταγῇ τῆς ὁμολογίας ὑμῶν εἰς τὸ εὐαγγέλιον τοῦ Χριστοῦ καὶ ἁπλότητι τῆς κοινωνίας εἰς αὐτοὺς καὶ εἰς πάντας,
\vs{14}καὶ αὐτῶν δεήσει ὑπὲρ ὑμῶν ἐπιποθούντων ὑμᾶς διὰ τὴν ὑπερβάλλουσαν χάριν τοῦ Θεοῦ ἐφ᾽ ὑμῖν.
\vs{15}Χάρις τῷ Θεῷ ἐπὶ τῇ ἀνεκδιηγήτῳ αὐτοῦ δωρεᾷ.

\ch{10}
Αὐτὸς δὲ ἐγὼ Παῦλος παρακαλῶ ὑμᾶς διὰ τῆς πραΰτητος καὶ ἐπιεικείας τοῦ Χριστοῦ, ὃς κατὰ πρόσωπον μὲν ταπεινὸς ἐν ὑμῖν, ἀπὼν δὲ θαρρῶ εἰς ὑμᾶς·
\vs{2}δέομαι δὲ τὸ μὴ παρὼν θαρρῆσαι τῇ πεποιθήσει ᾗ λογίζομαι τολμῆσαι ἐπί τινας τοὺς λογιζομένους ἡμᾶς ὡς κατὰ σάρκα περιπατοῦντας.
\vs{3}Ἐν σαρκὶ γὰρ περιπατοῦντες οὐ κατὰ σάρκα στρατευόμεθα,
\vs{4}τὰ γὰρ ὅπλα τῆς στρατείας ἡμῶν οὐ σαρκικὰ ἀλλὰ δυνατὰ τῷ Θεῷ πρὸς καθαίρεσιν ὀχυρωμάτων, λογισμοὺς καθαιροῦντες
\vs{5}καὶ πᾶν ὕψωμα ἐπαιρόμενον κατὰ τῆς γνώσεως τοῦ Θεοῦ, καὶ αἰχμαλωτίζοντες πᾶν νόημα εἰς τὴν ὑπακοὴν τοῦ Χριστοῦ,
\vs{6}καὶ ἐν ἑτοίμῳ ἔχοντες ἐκδικῆσαι πᾶσαν παρακοήν, ὅταν πληρωθῇ ὑμῶν ἡ ὑπακοή.

\vs{7}Τὰ κατὰ πρόσωπον βλέπετε. εἴ τις πέποιθεν ἑαυτῷ Χριστοῦ εἶναι, τοῦτο λογιζέσθω πάλιν ἐφ᾽ ἑαυτοῦ, ὅτι καθὼς αὐτὸς Χριστοῦ, οὕτως καὶ ἡμεῖς.
\vs{8}ἐάν τε γὰρ περισσότερόν τι καυχήσωμαι περὶ τῆς ἐξουσίας ἡμῶν ἧς ἔδωκεν ὁ Κύριος εἰς οἰκοδομὴν καὶ οὐκ εἰς καθαίρεσιν ὑμῶν, οὐκ αἰσχυνθήσομαι.
\vs{9}ἵνα μὴ δόξω ὡς ἂν ἐκφοβεῖν ὑμᾶς διὰ τῶν ἐπιστολῶν·
\vs{10}Ὅτι Αἱ ἐπιστολαὶ μέν, φησίν, Βαρεῖαι καὶ ἰσχυραί, ἡ δὲ παρουσία τοῦ σώματος ἀσθενὴς καὶ ὁ λόγος ἐξουθενημένος.
\vs{11}τοῦτο λογιζέσθω ὁ τοιοῦτος, ὅτι οἷοί ἐσμεν τῷ λόγῳ δι᾽ ἐπιστολῶν ἀπόντες, τοιοῦτοι καὶ παρόντες τῷ ἔργῳ.

\vs{12}Οὐ γὰρ τολμῶμεν ἐνκρῖναι ἢ συνκρῖναι ἑαυτούς τισιν τῶν ἑαυτοὺς συνιστανόντων, ἀλλὰ αὐτοὶ ἐν ἑαυτοῖς ἑαυτοὺς μετροῦντες καὶ συνκρίνοντες ἑαυτοὺς ἑαυτοῖς οὐ συνιᾶσιν.
\vs{13}ἡμεῖς δὲ οὐκ εἰς τὰ ἄμετρα καυχησόμεθα ἀλλὰ κατὰ τὸ μέτρον τοῦ κανόνος οὗ ἐμέρισεν ἡμῖν ὁ Θεὸς μέτρου, ἐφικέσθαι ἄχρι καὶ ὑμῶν.
\vs{14}οὐ γὰρ ὡς μὴ ἐφικνούμενοι εἰς ὑμᾶς ὑπερεκτείνομεν ἑαυτούς, ἄχρι γὰρ καὶ ὑμῶν ἐφθάσαμεν ἐν τῷ εὐαγγελίῳ τοῦ Χριστοῦ,
\vs{15}οὐκ εἰς τὰ ἄμετρα καυχώμενοι ἐν ἀλλοτρίοις κόποις, ἐλπίδα δὲ ἔχοντες αὐξανομένης τῆς πίστεως ὑμῶν ἐν ὑμῖν μεγαλυνθῆναι κατὰ τὸν κανόνα ἡμῶν εἰς περισσείαν
\vs{16}εἰς τὰ ὑπερέκεινα ὑμῶν εὐαγγελίσασθαι, οὐκ ἐν ἀλλοτρίῳ κανόνι εἰς τὰ ἕτοιμα καυχήσασθαι.
\vs{17}Ὁ δὲ καυχώμενος ἐν Κυρίῳ καυχάσθω·
\vs{18}οὐ γὰρ ὁ ἑαυτὸν συνιστάνων, ἐκεῖνός ἐστιν δόκιμος, ἀλλὰ ὃν ὁ Κύριος συνίστησιν.

\ch{11}
Ὄφελον ἀνείχεσθέ μου μικρόν τι ἀφροσύνης· ἀλλὰ καὶ ἀνέχεσθέ μου.
\vs{2}ζηλῶ γὰρ ὑμᾶς Θεοῦ ζήλῳ, ἡρμοσάμην γὰρ ὑμᾶς ἑνὶ ἀνδρὶ παρθένον ἁγνὴν παραστῆσαι τῷ Χριστῷ·
\vs{3}Φοβοῦμαι δὲ μή πως, ὡς ὁ ὄφις ἐξηπάτησεν Εὕαν ἐν τῇ πανουργίᾳ αὐτοῦ, φθαρῇ τὰ νοήματα ὑμῶν ἀπὸ τῆς ἁπλότητος καὶ τῆς ἁγνότητος τῆς εἰς τὸν Χριστόν.
\vs{4}εἰ μὲν γὰρ ὁ ἐρχόμενος ἄλλον Ἰησοῦν κηρύσσει ὃν οὐκ ἐκηρύξαμεν, ἢ πνεῦμα ἕτερον λαμβάνετε ὃ οὐκ ἐλάβετε, ἢ εὐαγγέλιον ἕτερον ὃ οὐκ ἐδέξασθε, καλῶς ἀνέχεσθε.

\vs{5}Λογίζομαι γὰρ μηδὲν ὑστερηκέναι τῶν Ὑπερλίαν ἀποστόλων.
\vs{6}εἰ δὲ καὶ ἰδιώτης τῷ λόγῳ, ἀλλ᾽ οὐ τῇ γνώσει, ἀλλ᾽ ἐν παντὶ φανερώσαντες ἐν πᾶσιν εἰς ὑμᾶς.
\vs{7}Ἢ ἁμαρτίαν ἐποίησα ἐμαυτὸν ταπεινῶν ἵνα ὑμεῖς ὑψωθῆτε, ὅτι δωρεὰν τὸ τοῦ Θεοῦ εὐαγγέλιον εὐηγγελισάμην ὑμῖν;
\vs{8}ἄλλας ἐκκλησίας ἐσύλησα λαβὼν ὀψώνιον πρὸς τὴν ὑμῶν διακονίαν,
\vs{9}καὶ παρὼν πρὸς ὑμᾶς καὶ ὑστερηθεὶς οὐ κατενάρκησα οὐθενός· τὸ γὰρ ὑστέρημά μου προσανεπλήρωσαν οἱ ἀδελφοὶ ἐλθόντες ἀπὸ Μακεδονίας, καὶ ἐν παντὶ ἀβαρῆ ἐμαυτὸν ὑμῖν ἐτήρησα καὶ τηρήσω.
\vs{10}ἔστιν ἀλήθεια Χριστοῦ ἐν ἐμοὶ ὅτι ἡ καύχησις αὕτη οὐ φραγήσεται εἰς ἐμὲ ἐν τοῖς κλίμασιν τῆς Ἀχαΐας.
\vs{11}διὰ τί; ὅτι οὐκ ἀγαπῶ ὑμᾶς; ὁ Θεὸς οἶδεν.

\vs{12}Ὃ δὲ ποιῶ, καὶ ποιήσω, ἵνα ἐκκόψω τὴν ἀφορμὴν τῶν θελόντων ἀφορμήν, ἵνα ἐν ᾧ καυχῶνται εὑρεθῶσιν καθὼς καὶ ἡμεῖς.
\vs{13}οἱ γὰρ τοιοῦτοι ψευδαπόστολοι, ἐργάται δόλιοι, μετασχηματιζόμενοι εἰς ἀποστόλους Χριστοῦ.
\vs{14}καὶ οὐ θαῦμα· αὐτὸς γὰρ ὁ Σατανᾶς μετασχηματίζεται εἰς ἄγγελον φωτός.
\vs{15}οὐ μέγα οὖν εἰ καὶ οἱ διάκονοι αὐτοῦ μετασχηματίζονται ὡς διάκονοι δικαιοσύνης· ὧν τὸ τέλος ἔσται κατὰ τὰ ἔργα αὐτῶν.

\vs{16}Πάλιν λέγω, μή τίς με δόξῃ ἄφρονα εἶναι· εἰ δὲ μή γε, κἂν ὡς ἄφρονα δέξασθέ με, ἵνα κἀγὼ μικρόν τι καυχήσωμαι.
\vs{17}ὃ λαλῶ, οὐ κατὰ Κύριον λαλῶ ἀλλ᾽ ὡς ἐν ἀφροσύνῃ, ἐν ταύτῃ τῇ ὑποστάσει τῆς καυχήσεως.
\vs{18}ἐπεὶ πολλοὶ καυχῶνται κατὰ σάρκα, κἀγὼ καυχήσομαι.
\vs{19}ἡδέως γὰρ ἀνέχεσθε τῶν ἀφρόνων φρόνιμοι ὄντες·
\vs{20}ἀνέχεσθε γὰρ εἴ τις ὑμᾶς καταδουλοῖ, εἴ τις κατεσθίει, εἴ τις λαμβάνει, εἴ τις ἐπαίρεται, εἴ τις εἰς πρόσωπον ὑμᾶς δέρει.
\vs{21}κατὰ ἀτιμίαν λέγω, ὡς ὅτι ἡμεῖς ἠσθενήκαμεν.

Ἐν ᾧ δ᾽ ἄν τις τολμᾷ, ἐν ἀφροσύνῃ λέγω, τολμῶ κἀγώ.
\vs{22}Ἑβραῖοί εἰσιν; κἀγώ. Ἰσραηλῖταί εἰσιν; κἀγώ. σπέρμα Ἀβραάμ εἰσιν; κἀγώ.
\vs{23}διάκονοι Χριστοῦ εἰσιν; παραφρονῶν λαλῶ, ὑπὲρ ἐγώ· ἐν κόποις περισσοτέρως, ἐν φυλακαῖς περισσοτέρως, ἐν πληγαῖς ὑπερβαλλόντως, ἐν θανάτοις πολλάκις.
\vs{24}Ὑπὸ Ἰουδαίων πεντάκις τεσσεράκοντα παρὰ μίαν ἔλαβον,
\vs{25}τρὶς ἐραβδίσθην, ἅπαξ ἐλιθάσθην, τρὶς ἐναυάγησα, νυχθήμερον ἐν τῷ βυθῷ πεποίηκα·
\vs{26}ὁδοιπορίαις πολλάκις, κινδύνοις ποταμῶν, κινδύνοις λῃστῶν, κινδύνοις ἐκ γένους, κινδύνοις ἐξ ἐθνῶν, κινδύνοις ἐν πόλει, κινδύνοις ἐν ἐρημίᾳ, κινδύνοις ἐν θαλάσσῃ, κινδύνοις ἐν ψευδαδέλφοις,
\vs{27}κόπῳ καὶ μόχθῳ, ἐν ἀγρυπνίαις πολλάκις, ἐν λιμῷ καὶ δίψει, ἐν νηστείαις πολλάκις, ἐν ψύχει καὶ γυμνότητι·
\vs{28}Χωρὶς τῶν παρεκτὸς ἡ ἐπίστασίς μοι ἡ καθ᾽ ἡμέραν, ἡ μέριμνα πασῶν τῶν ἐκκλησιῶν.
\vs{29}τίς ἀσθενεῖ καὶ οὐκ ἀσθενῶ; τίς σκανδαλίζεται καὶ οὐκ ἐγὼ πυροῦμαι;
\vs{30}Εἰ καυχᾶσθαι δεῖ, τὰ τῆς ἀσθενείας μου καυχήσομαι.
\vs{31}ὁ Θεὸς καὶ Πατὴρ τοῦ Κυρίου Ἰησοῦ οἶδεν, ὁ ὢν εὐλογητὸς εἰς τοὺς αἰῶνας, ὅτι οὐ ψεύδομαι.
\vs{32}ἐν Δαμασκῷ ὁ ἐθνάρχης Ἁρέτα τοῦ βασιλέως ἐφρούρει τὴν πόλιν Δαμασκηνῶν πιάσαι με,
\vs{33}καὶ διὰ θυρίδος ἐν σαργάνῃ ἐχαλάσθην διὰ τοῦ τείχους καὶ ἐξέφυγον τὰς χεῖρας αὐτοῦ.

\ch{12}
Καυχᾶσθαι δεῖ, οὐ συμφέρον μέν, ἐλεύσομαι δὲ εἰς ὀπτασίας καὶ ἀποκαλύψεις Κυρίου.
\vs{2}οἶδα ἄνθρωπον ἐν Χριστῷ πρὸ ἐτῶν δεκατεσσάρων, εἴτε ἐν σώματι οὐκ οἶδα, εἴτε ἐκτὸς τοῦ σώματος οὐκ οἶδα, ὁ Θεὸς οἶδεν, ἁρπαγέντα τὸν τοιοῦτον ἕως τρίτου οὐρανοῦ.
\vs{3}καὶ οἶδα τὸν τοιοῦτον ἄνθρωπον, εἴτε ἐν σώματι εἴτε χωρὶς τοῦ σώματος οὐκ οἶδα, ὁ Θεὸς οἶδεν,
\vs{4}ὅτι ἡρπάγη εἰς τὸν Παράδεισον καὶ ἤκουσεν ἄρρητα ῥήματα ἃ οὐκ ἐξὸν ἀνθρώπῳ λαλῆσαι.
\vs{5}Ὑπὲρ τοῦ τοιούτου καυχήσομαι, ὑπὲρ δὲ ἐμαυτοῦ οὐ καυχήσομαι εἰ μὴ ἐν ταῖς ἀσθενείαις.
\vs{6}ἐὰν γὰρ θελήσω καυχήσασθαι, οὐκ ἔσομαι ἄφρων, ἀλήθειαν γὰρ ἐρῶ· φείδομαι δέ, μή τις εἰς ἐμὲ λογίσηται ὑπὲρ ὃ βλέπει με ἢ ἀκούει τι ἐξ ἐμοῦ
\vs{7}καὶ τῇ ὑπερβολῇ τῶν ἀποκαλύψεων. Διὸ ἵνα μὴ ὑπεραίρωμαι, ἐδόθη μοι σκόλοψ τῇ σαρκί, ἄγγελος Σατανᾶ, ἵνα με κολαφίζῃ, ἵνα μὴ ὑπεραίρωμαι.
\vs{8}ὑπὲρ τούτου τρὶς τὸν Κύριον παρεκάλεσα ἵνα ἀποστῇ ἀπ᾽ ἐμοῦ.
\vs{9}καὶ εἴρηκέν μοι· Ἀρκεῖ σοι ἡ χάρις μου, ἡ γὰρ δύναμις ἐν ἀσθενείᾳ τελεῖται. Ἥδιστα οὖν μᾶλλον καυχήσομαι ἐν ταῖς ἀσθενείαις μου, ἵνα ἐπισκηνώσῃ ἐπ᾽ ἐμὲ ἡ δύναμις τοῦ Χριστοῦ.
\vs{10}διὸ εὐδοκῶ ἐν ἀσθενείαις, ἐν ὕβρεσιν, ἐν ἀνάγκαις, ἐν διωγμοῖς καὶ στενοχωρίαις, ὑπὲρ Χριστοῦ· ὅταν γὰρ ἀσθενῶ, τότε δυνατός εἰμι.

\vs{11}Γέγονα ἄφρων, ὑμεῖς με ἠναγκάσατε. ἐγὼ γὰρ ὤφειλον ὑφ᾽ ὑμῶν συνίστασθαι· οὐδὲν γὰρ ὑστέρησα τῶν Ὑπερλίαν ἀποστόλων εἰ καὶ οὐδέν εἰμι.
\vs{12}τὰ μὲν σημεῖα τοῦ ἀποστόλου κατειργάσθη ἐν ὑμῖν ἐν πάσῃ ὑπομονῇ, σημείοις τε καὶ τέρασιν καὶ δυνάμεσιν.
\vs{13}τί γάρ ἐστιν ὃ ἡσσώθητε ὑπὲρ τὰς λοιπὰς ἐκκλησίας, εἰ μὴ ὅτι αὐτὸς ἐγὼ οὐ κατενάρκησα ὑμῶν; χαρίσασθέ μοι τὴν ἀδικίαν ταύτην.

\vs{14}Ἰδοὺ τρίτον τοῦτο ἑτοίμως ἔχω ἐλθεῖν πρὸς ὑμᾶς, καὶ οὐ καταναρκήσω· οὐ γὰρ ζητῶ τὰ ὑμῶν ἀλλὰ ὑμᾶς. οὐ γὰρ ὀφείλει τὰ τέκνα τοῖς γονεῦσιν θησαυρίζειν ἀλλὰ οἱ γονεῖς τοῖς τέκνοις.
\vs{15}ἐγὼ δὲ ἥδιστα δαπανήσω καὶ ἐκδαπανηθήσομαι ὑπὲρ τῶν ψυχῶν ὑμῶν. εἰ περισσοτέρως ὑμᾶς ἀγαπῶν, ἧσσον ἀγαπῶμαι;
\vs{16}Ἔστω δέ, ἐγὼ οὐ κατεβάρησα ὑμᾶς· ἀλλὰ ὑπάρχων πανοῦργος δόλῳ ὑμᾶς ἔλαβον.
\vs{17}μή τινα ὧν ἀπέσταλκα πρὸς ὑμᾶς, δι᾽ αὐτοῦ ἐπλεονέκτησα ὑμᾶς;
\vs{18}παρεκάλεσα Τίτον καὶ συναπέστειλα τὸν ἀδελφόν· μήτι ἐπλεονέκτησεν ὑμᾶς Τίτος; οὐ τῷ αὐτῷ Πνεύματι περιεπατήσαμεν; οὐ τοῖς αὐτοῖς ἴχνεσιν;

\vs{19}Πάλαι δοκεῖτε ὅτι ὑμῖν ἀπολογούμεθα. κατέναντι Θεοῦ ἐν Χριστῷ λαλοῦμεν· τὰ δὲ πάντα, ἀγαπητοί, ὑπὲρ τῆς ὑμῶν οἰκοδομῆς.
\vs{20}φοβοῦμαι γὰρ μή πως ἐλθὼν οὐχ οἵους θέλω εὕρω ὑμᾶς κἀγὼ εὑρεθῶ ὑμῖν οἷον οὐ θέλετε· μή πως ἔρις, ζῆλος, θυμοί, ἐριθεῖαι, καταλαλιαί, ψιθυρισμοί, φυσιώσεις, ἀκαταστασίαι·
\vs{21}μὴ πάλιν ἐλθόντος μου ταπεινώσῃ με ὁ Θεός μου πρὸς ὑμᾶς καὶ πενθήσω πολλοὺς τῶν προημαρτηκότων καὶ μὴ μετανοησάντων ἐπὶ τῇ ἀκαθαρσίᾳ καὶ πορνείᾳ καὶ ἀσελγείᾳ ᾗ ἔπραξαν.

\ch{13}
Τρίτον τοῦτο ἔρχομαι πρὸς ὑμᾶς· Ἐπὶ στόματος δύο μαρτύρων καὶ τριῶν σταθήσεται πᾶν ῥῆμα.
\vs{2}Προείρηκα καὶ προλέγω, ὡς παρὼν τὸ δεύτερον καὶ ἀπὼν νῦν, τοῖς προημαρτηκόσιν καὶ τοῖς λοιποῖς πᾶσιν, ὅτι ἐὰν ἔλθω εἰς τὸ πάλιν οὐ φείσομαι,
\vs{3}ἐπεὶ δοκιμὴν ζητεῖτε τοῦ ἐν ἐμοὶ λαλοῦντος Χριστοῦ, ὃς εἰς ὑμᾶς οὐκ ἀσθενεῖ ἀλλὰ δυνατεῖ ἐν ὑμῖν.
\vs{4}καὶ γὰρ ἐσταυρώθη ἐξ ἀσθενείας, ἀλλὰ ζῇ ἐκ δυνάμεως Θεοῦ. καὶ γὰρ ἡμεῖς ἀσθενοῦμεν ἐν αὐτῷ, ἀλλὰ ζήσομεν σὺν αὐτῷ ἐκ δυνάμεως Θεοῦ εἰς ὑμᾶς.

\vs{5}Ἑαυτοὺς πειράζετε εἰ ἐστὲ ἐν τῇ πίστει, ἑαυτοὺς δοκιμάζετε· ἢ οὐκ ἐπιγινώσκετε ἑαυτοὺς ὅτι Ἰησοῦς Χριστὸς ἐν ὑμῖν; εἰ μήτι ἀδόκιμοί ἐστε.
\vs{6}ἐλπίζω δὲ ὅτι γνώσεσθε ὅτι ἡμεῖς οὐκ ἐσμὲν ἀδόκιμοι.
\vs{7}Εὐχόμεθα δὲ πρὸς τὸν Θεὸν μὴ ποιῆσαι ὑμᾶς κακὸν μηδέν, οὐχ ἵνα ἡμεῖς δόκιμοι φανῶμεν, ἀλλ᾽ ἵνα ὑμεῖς τὸ καλὸν ποιῆτε, ἡμεῖς δὲ ὡς ἀδόκιμοι ὦμεν.
\vs{8}οὐ γὰρ δυνάμεθά τι κατὰ τῆς ἀληθείας ἀλλὰ ὑπὲρ τῆς ἀληθείας.
\vs{9}χαίρομεν γὰρ ὅταν ἡμεῖς ἀσθενῶμεν, ὑμεῖς δὲ δυνατοὶ ἦτε· τοῦτο καὶ εὐχόμεθα, τὴν ὑμῶν κατάρτισιν.
\vs{10}Διὰ τοῦτο ταῦτα ἀπὼν γράφω, ἵνα παρὼν μὴ ἀποτόμως χρήσωμαι κατὰ τὴν ἐξουσίαν ἣν ὁ Κύριος ἔδωκέν μοι εἰς οἰκοδομὴν καὶ οὐκ εἰς καθαίρεσιν.

\vs{11}Λοιπόν, ἀδελφοί, χαίρετε, καταρτίζεσθε, παρακαλεῖσθε, τὸ αὐτὸ φρονεῖτε, εἰρηνεύετε, καὶ ὁ Θεὸς τῆς ἀγάπης καὶ εἰρήνης ἔσται μεθ᾽ ὑμῶν.
\vs{12}Ἀσπάσασθε ἀλλήλους ἐν ἁγίῳ φιλήματι. Ἀσπάζονται ὑμᾶς οἱ ἅγιοι πάντες.

\vs{13}Ἡ χάρις τοῦ Κυρίου Ἰησοῦ Χριστοῦ καὶ ἡ ἀγάπη τοῦ Θεοῦ καὶ ἡ κοινωνία τοῦ Ἁγίου Πνεύματος μετὰ πάντων ὑμῶν.


\def\book{ΠΡΟΣ ΓΑΛΑΤΑΣ}
\biblebook{ΠΡΟΣ ΓΑΛΑΤΑΣ}


\lettrine[lines=2, loversize=0.2, nindent=0em, findent=.25em]{\textcolor{bookheadingcolor}{Π}}{αῦλος} ἀπόστολος οὐκ ἀπ᾽ ἀνθρώπων οὐδὲ δι᾽ ἀνθρώπου ἀλλὰ διὰ Ἰησοῦ Χριστοῦ καὶ Θεοῦ Πατρὸς τοῦ ἐγείραντος αὐτὸν ἐκ νεκρῶν,
\vs{2}καὶ οἱ σὺν ἐμοὶ πάντες ἀδελφοί Ταῖς ἐκκλησίαις τῆς Γαλατίας,
\vs{3}Χάρις ὑμῖν καὶ εἰρήνη ἀπὸ Θεοῦ Πατρὸς ἡμῶν καὶ Κυρίου Ἰησοῦ Χριστοῦ
\vs{4}τοῦ δόντος ἑαυτὸν ὑπὲρ τῶν ἁμαρτιῶν ἡμῶν, ὅπως ἐξέληται ἡμᾶς ἐκ τοῦ αἰῶνος τοῦ ἐνεστῶτος πονηροῦ κατὰ τὸ θέλημα τοῦ Θεοῦ καὶ Πατρὸς ἡμῶν,
\vs{5}ᾧ ἡ δόξα εἰς τοὺς αἰῶνας τῶν αἰώνων, ἀμήν.

\vs{6}Θαυμάζω ὅτι οὕτως ταχέως μετατίθεσθε ἀπὸ τοῦ καλέσαντος ὑμᾶς ἐν χάριτι Χριστοῦ εἰς ἕτερον εὐαγγέλιον,
\vs{7}ὃ οὐκ ἔστιν ἄλλο, εἰ μή τινές εἰσιν οἱ ταράσσοντες ὑμᾶς καὶ θέλοντες μεταστρέψαι τὸ εὐαγγέλιον τοῦ Χριστοῦ.
\vs{8}Ἀλλὰ καὶ ἐὰν ἡμεῖς ἢ ἄγγελος ἐξ οὐρανοῦ εὐαγγελίζηται ὑμῖν παρ᾽ ὃ εὐηγγελισάμεθα ὑμῖν, ἀνάθεμα ἔστω.
\vs{9}ὡς προειρήκαμεν καὶ ἄρτι πάλιν λέγω· εἴ τις ὑμᾶς εὐαγγελίζεται παρ᾽ ὃ παρελάβετε, ἀνάθεμα ἔστω.

\vs{10}Ἄρτι γὰρ ἀνθρώπους πείθω ἢ τὸν Θεόν; ἢ ζητῶ ἀνθρώποις ἀρέσκειν; εἰ ἔτι ἀνθρώποις ἤρεσκον, Χριστοῦ δοῦλος οὐκ ἂν ἤμην.
\vs{11}γνωρίζω γὰρ ὑμῖν, ἀδελφοί, τὸ εὐαγγέλιον τὸ εὐαγγελισθὲν ὑπ᾽ ἐμοῦ ὅτι οὐκ ἔστιν κατὰ ἄνθρωπον·
\vs{12}οὐδὲ γὰρ ἐγὼ παρὰ ἀνθρώπου παρέλαβον αὐτό οὔτε ἐδιδάχθην, ἀλλὰ δι᾽ ἀποκαλύψεως Ἰησοῦ Χριστοῦ.
\vs{13}Ἠκούσατε γὰρ τὴν ἐμὴν ἀναστροφήν ποτε ἐν τῷ Ἰουδαϊσμῷ, ὅτι καθ᾽ ὑπερβολὴν ἐδίωκον τὴν ἐκκλησίαν τοῦ Θεοῦ καὶ ἐπόρθουν αὐτήν,
\vs{14}καὶ προέκοπτον ἐν τῷ Ἰουδαϊσμῷ ὑπὲρ πολλοὺς συνηλικιώτας ἐν τῷ γένει μου, περισσοτέρως ζηλωτὴς ὑπάρχων τῶν πατρικῶν μου παραδόσεων.
\vs{15}Ὅτε δὲ εὐδόκησεν ὁ θεὸς ὁ ἀφορίσας με ἐκ κοιλίας μητρός μου καὶ καλέσας διὰ τῆς χάριτος αὐτοῦ
\vs{16}ἀποκαλύψαι τὸν Υἱὸν αὐτοῦ ἐν ἐμοὶ, ἵνα εὐαγγελίζωμαι αὐτὸν ἐν τοῖς ἔθνεσιν, εὐθέως οὐ προσανεθέμην σαρκὶ καὶ αἵματι
\vs{17}οὐδὲ ἀνῆλθον εἰς Ἱεροσόλυμα πρὸς τοὺς πρὸ ἐμοῦ ἀποστόλους, ἀλλὰ ἀπῆλθον εἰς Ἀραβίαν καὶ πάλιν ὑπέστρεψα εἰς Δαμασκόν.
\vs{18}Ἔπειτα μετὰ ἔτη τρία ἀνῆλθον εἰς Ἱεροσόλυμα ἱστορῆσαι Κηφᾶν καὶ ἐπέμεινα πρὸς αὐτὸν ἡμέρας δεκαπέντε,
\vs{19}ἕτερον δὲ τῶν ἀποστόλων οὐκ εἶδον εἰ μὴ Ἰάκωβον τὸν ἀδελφὸν τοῦ Κυρίου.
\vs{20}ἃ δὲ γράφω ὑμῖν, ἰδοὺ ἐνώπιον τοῦ Θεοῦ ὅτι οὐ ψεύδομαι.
\vs{21}Ἔπειτα ἦλθον εἰς τὰ κλίματα τῆς Συρίας καὶ τῆς Κιλικίας·
\vs{22}ἤμην δὲ ἀγνοούμενος τῷ προσώπῳ ταῖς ἐκκλησίαις τῆς Ἰουδαίας ταῖς ἐν Χριστῷ.
\vs{23}μόνον δὲ ἀκούοντες ἦσαν ὅτι Ὁ διώκων ἡμᾶς ποτε νῦν εὐαγγελίζεται τὴν πίστιν ἥν ποτε ἐπόρθει,
\vs{24}καὶ ἐδόξαζον ἐν ἐμοὶ τὸν Θεόν.

\ch{2}
Ἔπειτα διὰ δεκατεσσάρων ἐτῶν πάλιν ἀνέβην εἰς Ἱεροσόλυμα μετὰ Βαρνάβα συμπαραλαβὼν καὶ Τίτον·
\vs{2}ἀνέβην δὲ κατὰ ἀποκάλυψιν· καὶ ἀνεθέμην αὐτοῖς τὸ εὐαγγέλιον ὃ κηρύσσω ἐν τοῖς ἔθνεσιν, κατ᾽ ἰδίαν δὲ τοῖς δοκοῦσιν, μή πως εἰς κενὸν τρέχω ἢ ἔδραμον.
\vs{3}ἀλλ᾽ οὐδὲ Τίτος ὁ σὺν ἐμοί, Ἕλλην ὤν, ἠναγκάσθη περιτμηθῆναι·
\vs{4}διὰ δὲ τοὺς παρεισάκτους ψευδαδέλφους, οἵτινες παρεισῆλθον κατασκοπῆσαι τὴν ἐλευθερίαν ἡμῶν ἣν ἔχομεν ἐν Χριστῷ Ἰησοῦ, ἵνα ἡμᾶς καταδουλώσουσιν,
\vs{5}οἷς οὐδὲ πρὸς ὥραν εἴξαμεν τῇ ὑποταγῇ, ἵνα ἡ ἀλήθεια τοῦ εὐαγγελίου διαμείνῃ πρὸς ὑμᾶς.
\vs{6}Ἀπὸ δὲ τῶν δοκούντων εἶναί τι,— ὁποῖοί ποτε ἦσαν οὐδέν μοι διαφέρει· πρόσωπον ὁ Θεὸς ἀνθρώπου οὐ λαμβάνει— ἐμοὶ γὰρ οἱ δοκοῦντες οὐδὲν προσανέθεντο,
\vs{7}ἀλλὰ τοὐναντίον ἰδόντες ὅτι πεπίστευμαι τὸ εὐαγγέλιον τῆς ἀκροβυστίας καθὼς Πέτρος τῆς περιτομῆς,
\vs{8}ὁ γὰρ ἐνεργήσας Πέτρῳ εἰς ἀποστολὴν τῆς περιτομῆς ἐνήργησεν καὶ ἐμοὶ εἰς τὰ ἔθνη,
\vs{9}καὶ γνόντες τὴν χάριν τὴν δοθεῖσάν μοι, Ἰάκωβος καὶ Κηφᾶς καὶ Ἰωάννης, οἱ δοκοῦντες στῦλοι εἶναι, δεξιὰς ἔδωκαν ἐμοὶ καὶ Βαρνάβα κοινωνίας, ἵνα ἡμεῖς εἰς τὰ ἔθνη, αὐτοὶ δὲ εἰς τὴν περιτομήν·
\vs{10}μόνον τῶν πτωχῶν ἵνα μνημονεύωμεν, ὃ καὶ ἐσπούδασα αὐτὸ τοῦτο ποιῆσαι.

\vs{11}Ὅτε δὲ ἦλθεν Κηφᾶς εἰς Ἀντιόχειαν, κατὰ πρόσωπον αὐτῷ ἀντέστην, ὅτι κατεγνωσμένος ἦν.
\vs{12}πρὸ τοῦ γὰρ ἐλθεῖν τινας ἀπὸ Ἰακώβου μετὰ τῶν ἐθνῶν συνήσθιεν· ὅτε δὲ ἦλθον, ὑπέστελλεν καὶ ἀφώριζεν ἑαυτόν φοβούμενος τοὺς ἐκ περιτομῆς.
\vs{13}καὶ συνυπεκρίθησαν αὐτῷ καὶ οἱ λοιποὶ Ἰουδαῖοι, ὥστε καὶ Βαρνάβας συναπήχθη αὐτῶν τῇ ὑποκρίσει.
\vs{14}Ἀλλ᾽ ὅτε εἶδον ὅτι οὐκ ὀρθοποδοῦσιν πρὸς τὴν ἀλήθειαν τοῦ εὐαγγελίου, εἶπον τῷ Κηφᾷ ἔμπροσθεν πάντων· Εἰ σὺ Ἰουδαῖος ὑπάρχων ἐθνικῶς καὶ οὐκ Ἰουδαϊκῶς ζῇς, πῶς τὰ ἔθνη ἀναγκάζεις ἰουδαΐζειν;
\vs{15}Ἡμεῖς φύσει Ἰουδαῖοι καὶ οὐκ ἐξ ἐθνῶν Ἁμαρτωλοί·
\vs{16}εἰδότες δὲ ὅτι οὐ δικαιοῦται ἄνθρωπος ἐξ ἔργων νόμου ἐὰν μὴ διὰ πίστεως Ἰησοῦ Χριστοῦ, καὶ ἡμεῖς εἰς Χριστὸν Ἰησοῦν ἐπιστεύσαμεν, ἵνα δικαιωθῶμεν ἐκ πίστεως Χριστοῦ καὶ οὐκ ἐξ ἔργων νόμου, ὅτι ἐξ ἔργων νόμου οὐ δικαιωθήσεται πᾶσα σάρξ.
\vs{17}Εἰ δὲ ζητοῦντες δικαιωθῆναι ἐν Χριστῷ εὑρέθημεν καὶ αὐτοὶ ἁμαρτωλοί, ἆρα Χριστὸς ἁμαρτίας διάκονος; μὴ γένοιτο.
\vs{18}εἰ γὰρ ἃ κατέλυσα ταῦτα πάλιν οἰκοδομῶ, παραβάτην ἐμαυτὸν συνιστάνω.
\vs{19}ἐγὼ γὰρ διὰ νόμου νόμῳ ἀπέθανον, ἵνα Θεῷ ζήσω. Χριστῷ συνεσταύρωμαι·
\vs{20}ζῶ δὲ οὐκέτι ἐγώ, ζῇ δὲ ἐν ἐμοὶ Χριστός· ὃ δὲ νῦν ζῶ ἐν σαρκί, ἐν πίστει ζῶ τῇ τοῦ Υἱοῦ τοῦ Θεοῦ τοῦ ἀγαπήσαντός με καὶ παραδόντος ἑαυτὸν ὑπὲρ ἐμοῦ.
\vs{21}Οὐκ ἀθετῶ τὴν χάριν τοῦ Θεοῦ· εἰ γὰρ διὰ νόμου δικαιοσύνη, ἄρα Χριστὸς δωρεὰν ἀπέθανεν.

\ch{3}
Ὦ ἀνόητοι Γαλάται, τίς ὑμᾶς ἐβάσκανεν, οἷς κατ᾽ ὀφθαλμοὺς Ἰησοῦς Χριστὸς προεγράφη ἐσταυρωμένος;
\vs{2}τοῦτο μόνον θέλω μαθεῖν ἀφ᾽ ὑμῶν· ἐξ ἔργων νόμου τὸ Πνεῦμα ἐλάβετε ἢ ἐξ ἀκοῆς πίστεως;
\vs{3}Οὕτως ἀνόητοί ἐστε, ἐναρξάμενοι Πνεύματι νῦν σαρκὶ ἐπιτελεῖσθε;
\vs{4}τοσαῦτα ἐπάθετε εἰκῇ; εἴ γε καὶ εἰκῇ.
\vs{5}ὁ οὖν ἐπιχορηγῶν ὑμῖν τὸ Πνεῦμα καὶ ἐνεργῶν δυνάμεις ἐν ὑμῖν, ἐξ ἔργων νόμου ἢ ἐξ ἀκοῆς πίστεως;

\vs{6}Καθὼς Ἀβραὰμ ἐπίστευσεν τῷ Θεῷ, καὶ ἐλογίσθη αὐτῷ εἰς δικαιοσύνην·
\vs{7}Γινώσκετε ἄρα ὅτι οἱ ἐκ πίστεως, οὗτοι υἱοί εἰσιν Ἀβραάμ.
\vs{8}προϊδοῦσα δὲ ἡ γραφὴ ὅτι ἐκ πίστεως δικαιοῖ τὰ ἔθνη ὁ Θεὸς, προευηγγελίσατο τῷ Ἀβραὰμ ὅτι Ἐνευλογηθήσονται ἐν σοὶ πάντα τὰ ἔθνη·
\vs{9}ὥστε οἱ ἐκ πίστεως εὐλογοῦνται σὺν τῷ πιστῷ Ἀβραάμ.

\vs{10}Ὅσοι γὰρ ἐξ ἔργων νόμου εἰσὶν, ὑπὸ κατάραν εἰσίν· γέγραπται γὰρ ὅτι Ἐπικατάρατος πᾶς ὃς οὐκ ἐμμένει πᾶσιν τοῖς γεγραμμένοις ἐν τῷ βιβλίῳ τοῦ νόμου τοῦ ποιῆσαι αὐτά.
\vs{11}ὅτι δὲ ἐν νόμῳ οὐδεὶς δικαιοῦται παρὰ τῷ Θεῷ δῆλον, ὅτι Ὁ δίκαιος ἐκ πίστεως ζήσεται·
\vs{12}ὁ δὲ νόμος οὐκ ἔστιν ἐκ πίστεως, ἀλλ᾽ Ὁ ποιήσας αὐτὰ ζήσεται ἐν αὐτοῖς.
\vs{13}Χριστὸς ἡμᾶς ἐξηγόρασεν ἐκ τῆς κατάρας τοῦ νόμου γενόμενος ὑπὲρ ἡμῶν κατάρα, ὅτι γέγραπται· Ἐπικατάρατος πᾶς ὁ κρεμάμενος ἐπὶ ξύλου,
\vs{14}ἵνα εἰς τὰ ἔθνη ἡ εὐλογία τοῦ Ἀβραὰμ γένηται ἐν Χριστῷ Ἰησοῦ, ἵνα τὴν ἐπαγγελίαν τοῦ Πνεύματος λάβωμεν διὰ τῆς πίστεως.

\vs{15}Ἀδελφοί, κατὰ ἄνθρωπον λέγω· ὅμως ἀνθρώπου κεκυρωμένην διαθήκην οὐδεὶς ἀθετεῖ ἢ ἐπιδιατάσσεται.
\vs{16}τῷ δὲ Ἀβραὰμ ἐρρέθησαν αἱ ἐπαγγελίαι καὶ τῷ σπέρματι αὐτοῦ. οὐ λέγει· Καὶ τοῖς σπέρμασιν, ὡς ἐπὶ πολλῶν ἀλλ᾽ ὡς ἐφ᾽ ἑνός· Καὶ τῷ σπέρματί σου, ὅς ἐστιν Χριστός.
\vs{17}Τοῦτο δὲ λέγω· διαθήκην προκεκυρωμένην ὑπὸ τοῦ Θεοῦ ὁ μετὰ τετρακόσια καὶ τριάκοντα ἔτη γεγονὼς νόμος οὐκ ἀκυροῖ εἰς τὸ καταργῆσαι τὴν ἐπαγγελίαν.
\vs{18}εἰ γὰρ ἐκ νόμου ἡ κληρονομία, οὐκέτι ἐξ ἐπαγγελίας· τῷ δὲ Ἀβραὰμ δι᾽ ἐπαγγελίας κεχάρισται ὁ Θεός.

\vs{19}Τί οὖν ὁ νόμος; τῶν παραβάσεων χάριν προσετέθη, ἄχρις οὗ ἔλθῃ τὸ σπέρμα ᾧ ἐπήγγελται, διαταγεὶς δι᾽ ἀγγέλων ἐν χειρὶ μεσίτου.
\vs{20}ὁ δὲ μεσίτης ἑνὸς οὐκ ἔστιν, ὁ δὲ Θεὸς εἷς ἐστιν.
\vs{21}Ὁ οὖν νόμος κατὰ τῶν ἐπαγγελιῶν τοῦ Θεοῦ; μὴ γένοιτο. εἰ γὰρ ἐδόθη νόμος ὁ δυνάμενος ζωοποιῆσαι, ὄντως ἐκ νόμου ἂν ἦν ἡ δικαιοσύνη·
\vs{22}ἀλλὰ συνέκλεισεν ἡ γραφὴ τὰ πάντα ὑπὸ ἁμαρτίαν, ἵνα ἡ ἐπαγγελία ἐκ πίστεως Ἰησοῦ Χριστοῦ δοθῇ τοῖς πιστεύουσιν.

\vs{23}Πρὸ τοῦ δὲ ἐλθεῖν τὴν πίστιν ὑπὸ νόμον ἐφρουρούμεθα συνκλειόμενοι εἰς τὴν μέλλουσαν πίστιν ἀποκαλυφθῆναι,
\vs{24}ὥστε ὁ νόμος παιδαγωγὸς ἡμῶν γέγονεν εἰς Χριστόν, ἵνα ἐκ πίστεως δικαιωθῶμεν·
\vs{25}ἐλθούσης δὲ τῆς πίστεως οὐκέτι ὑπὸ παιδαγωγόν ἐσμεν.
\vs{26}Πάντες γὰρ υἱοὶ Θεοῦ ἐστε διὰ τῆς πίστεως ἐν Χριστῷ Ἰησοῦ·
\vs{27}ὅσοι γὰρ εἰς Χριστὸν ἐβαπτίσθητε, Χριστὸν ἐνεδύσασθε.
\vs{28}οὐκ ἔνι Ἰουδαῖος οὐδὲ Ἕλλην, οὐκ ἔνι δοῦλος οὐδὲ ἐλεύθερος, οὐκ ἔνι ἄρσεν καὶ θῆλυ· πάντες γὰρ ὑμεῖς εἷς ἐστε ἐν Χριστῷ Ἰησοῦ.
\vs{29}εἰ δὲ ὑμεῖς Χριστοῦ, ἄρα τοῦ Ἀβραὰμ σπέρμα ἐστέ, κατ᾽ ἐπαγγελίαν κληρονόμοι.

\ch{4}
Λέγω δέ, ἐφ᾽ ὅσον χρόνον ὁ κληρονόμος νήπιός ἐστιν, οὐδὲν διαφέρει δούλου κύριος πάντων ὤν,
\vs{2}ἀλλὰ ὑπὸ ἐπιτρόπους ἐστὶν καὶ οἰκονόμους ἄχρι τῆς προθεσμίας τοῦ πατρός.
\vs{3}οὕτως καὶ ἡμεῖς, ὅτε ἦμεν νήπιοι, ὑπὸ τὰ στοιχεῖα τοῦ κόσμου ἤμεθα δεδουλωμένοι·
\vs{4}Ὅτε δὲ ἦλθεν τὸ πλήρωμα τοῦ χρόνου, ἐξαπέστειλεν ὁ Θεὸς τὸν Υἱὸν αὐτοῦ, γενόμενον ἐκ γυναικός, γενόμενον ὑπὸ νόμον,
\vs{5}ἵνα τοὺς ὑπὸ νόμον ἐξαγοράσῃ, ἵνα τὴν υἱοθεσίαν ἀπολάβωμεν.
\vs{6}Ὅτι δέ ἐστε υἱοί, ἐξαπέστειλεν ὁ Θεὸς τὸ Πνεῦμα τοῦ Υἱοῦ αὐτοῦ εἰς τὰς καρδίας ἡμῶν κρᾶζον· Ἀββᾶ ὁ Πατήρ.
\vs{7}ὥστε οὐκέτι εἶ δοῦλος ἀλλὰ υἱός· εἰ δὲ υἱός, καὶ κληρονόμος διὰ Θεοῦ.

\vs{8}Ἀλλὰ τότε μὲν οὐκ εἰδότες Θεὸν ἐδουλεύσατε τοῖς φύσει μὴ οὖσιν θεοῖς·
\vs{9}νῦν δὲ γνόντες Θεόν, μᾶλλον δὲ γνωσθέντες ὑπὸ Θεοῦ, πῶς ἐπιστρέφετε πάλιν ἐπὶ τὰ ἀσθενῆ καὶ πτωχὰ στοιχεῖα οἷς πάλιν ἄνωθεν δουλεύειν θέλετε;
\vs{10}ἡμέρας παρατηρεῖσθε καὶ μῆνας καὶ καιροὺς καὶ ἐνιαυτούς,
\vs{11}φοβοῦμαι ὑμᾶς μή πως εἰκῇ κεκοπίακα εἰς ὑμᾶς.

\vs{12}Γίνεσθε ὡς ἐγώ, ὅτι κἀγὼ ὡς ὑμεῖς, ἀδελφοί, δέομαι ὑμῶν. οὐδέν με ἠδικήσατε·
\vs{13}Οἴδατε δὲ ὅτι δι᾽ ἀσθένειαν τῆς σαρκὸς εὐηγγελισάμην ὑμῖν τὸ πρότερον,
\vs{14}καὶ τὸν πειρασμὸν ὑμῶν ἐν τῇ σαρκί μου οὐκ ἐξουθενήσατε οὐδὲ ἐξεπτύσατε, ἀλλὰ ὡς ἄγγελον Θεοῦ ἐδέξασθέ με, ὡς Χριστὸν Ἰησοῦν.
\vs{15}ποῦ οὖν ὁ μακαρισμὸς ὑμῶν; μαρτυρῶ γὰρ ὑμῖν ὅτι εἰ δυνατὸν τοὺς ὀφθαλμοὺς ὑμῶν ἐξορύξαντες ἐδώκατέ μοι.
\vs{16}ὥστε ἐχθρὸς ὑμῶν γέγονα ἀληθεύων ὑμῖν;
\vs{17}Ζηλοῦσιν ὑμᾶς οὐ καλῶς, ἀλλὰ ἐκκλεῖσαι ὑμᾶς θέλουσιν, ἵνα αὐτοὺς ζηλοῦτε·
\vs{18}καλὸν δὲ ζηλοῦσθαι ἐν καλῷ πάντοτε καὶ μὴ μόνον ἐν τῷ παρεῖναί με πρὸς ὑμᾶς.
\vs{19}Τέκνα μου, οὓς πάλιν ὠδίνω μέχρις οὗ μορφωθῇ Χριστὸς ἐν ὑμῖν·
\vs{20}ἤθελον δὲ παρεῖναι πρὸς ὑμᾶς ἄρτι καὶ ἀλλάξαι τὴν φωνήν μου, ὅτι ἀποροῦμαι ἐν ὑμῖν.

\vs{21}Λέγετέ μοι, οἱ ὑπὸ νόμον θέλοντες εἶναι, τὸν νόμον οὐκ ἀκούετε;
\vs{22}γέγραπται γὰρ ὅτι Ἀβραὰμ δύο υἱοὺς ἔσχεν, ἕνα ἐκ τῆς παιδίσκης καὶ ἕνα ἐκ τῆς ἐλευθέρας.
\vs{23}ἀλλ᾽ ὁ μὲν ἐκ τῆς παιδίσκης κατὰ σάρκα γεγέννηται, ὁ δὲ ἐκ τῆς ἐλευθέρας δι᾽ ἐπαγγελίας.
\vs{24}ἅτινά ἐστιν ἀλληγορούμενα· αὗται γάρ εἰσιν δύο διαθῆκαι, μία μὲν ἀπὸ ὄρους Σινᾶ εἰς δουλείαν γεννῶσα, ἥτις ἐστὶν Ἁγάρ.
\vs{25}τὸ δὲ Ἁγὰρ Σινᾶ ὄρος ἐστὶν ἐν τῇ Ἀραβίᾳ· συστοιχεῖ δὲ τῇ νῦν Ἰερουσαλήμ, δουλεύει γὰρ μετὰ τῶν τέκνων αὐτῆς.
\vs{26}ἡ δὲ ἄνω Ἰερουσαλὴμ ἐλευθέρα ἐστίν, ἥτις ἐστὶν μήτηρ ἡμῶν·
\vs{27}γέγραπται γάρ· 
\begin{poetryblock}

\begin{quote}Εὐφράνθητι, στεῖρα ἡ οὐ τίκτουσα,\end{quote} 

\begin{quote}ῥῆξον καὶ βόησον, ἡ οὐκ ὠδίνουσα·\end{quote} 

\begin{quote}ὅτι πολλὰ τὰ τέκνα τῆς ἐρήμου\end{quote} 

\begin{quote}μᾶλλον ἢ τῆς ἐχούσης τὸν ἄνδρα.\end{quote}
\end{poetryblock}

\vs{28}Ὑμεῖς δέ, ἀδελφοί, κατὰ Ἰσαὰκ ἐπαγγελίας τέκνα ἐστέ.
\vs{29}ἀλλ᾽ ὥσπερ τότε ὁ κατὰ σάρκα γεννηθεὶς ἐδίωκεν τὸν κατὰ Πνεῦμα, οὕτως καὶ νῦν.
\vs{30}Ἀλλὰ τί λέγει ἡ γραφή; Ἔκβαλε τὴν παιδίσκην καὶ τὸν υἱὸν αὐτῆς· οὐ γὰρ μὴ κληρονομήσει ὁ υἱὸς τῆς παιδίσκης μετὰ τοῦ υἱοῦ τῆς ἐλευθέρας.
\vs{31}διό, ἀδελφοί, οὐκ ἐσμὲν παιδίσκης τέκνα ἀλλὰ τῆς ἐλευθέρας.

\ch{5}
Τῇ ἐλευθερίᾳ ἡμᾶς Χριστὸς ἠλευθέρωσεν· στήκετε οὖν καὶ μὴ πάλιν ζυγῷ δουλείας ἐνέχεσθε.
\vs{2}Ἴδε ἐγὼ Παῦλος λέγω ὑμῖν ὅτι ἐὰν περιτέμνησθε, Χριστὸς ὑμᾶς οὐδὲν ὠφελήσει.
\vs{3}μαρτύρομαι δὲ πάλιν παντὶ ἀνθρώπῳ περιτεμνομένῳ ὅτι ὀφειλέτης ἐστὶν ὅλον τὸν νόμον ποιῆσαι.
\vs{4}κατηργήθητε ἀπὸ Χριστοῦ, οἵτινες ἐν νόμῳ δικαιοῦσθε, τῆς χάριτος ἐξεπέσατε.
\vs{5}Ἡμεῖς γὰρ Πνεύματι ἐκ πίστεως ἐλπίδα δικαιοσύνης ἀπεκδεχόμεθα.
\vs{6}ἐν γὰρ Χριστῷ Ἰησοῦ οὔτε περιτομή τι ἰσχύει οὔτε ἀκροβυστία ἀλλὰ πίστις δι᾽ ἀγάπης ἐνεργουμένη.

\vs{7}Ἐτρέχετε καλῶς· τίς ὑμᾶς ἐνέκοψεν τῇ ἀληθείᾳ μὴ πείθεσθαι;
\vs{8}ἡ πεισμονὴ οὐκ ἐκ τοῦ καλοῦντος ὑμᾶς.
\vs{9}μικρὰ ζύμη ὅλον τὸ φύραμα ζυμοῖ.
\vs{10}ἐγὼ πέποιθα εἰς ὑμᾶς ἐν Κυρίῳ ὅτι οὐδὲν ἄλλο φρονήσετε· ὁ δὲ ταράσσων ὑμᾶς βαστάσει τὸ κρίμα, ὅστις ἐὰν ᾖ.
\vs{11}Ἐγὼ δέ, ἀδελφοί, εἰ περιτομὴν ἔτι κηρύσσω, τί ἔτι διώκομαι; ἄρα κατήργηται τὸ σκάνδαλον τοῦ σταυροῦ.
\vs{12}Ὄφελον καὶ ἀποκόψονται οἱ ἀναστατοῦντες ὑμᾶς.

\vs{13}Ὑμεῖς γὰρ ἐπ᾽ ἐλευθερίᾳ ἐκλήθητε, ἀδελφοί· μόνον μὴ τὴν ἐλευθερίαν εἰς ἀφορμὴν τῇ σαρκί, ἀλλὰ διὰ τῆς ἀγάπης δουλεύετε ἀλλήλοις.
\vs{14}ὁ γὰρ πᾶς νόμος ἐν ἑνὶ λόγῳ πεπλήρωται, ἐν τῷ· Ἀγαπήσεις τὸν πλησίον σου ὡς σεαυτόν.
\vs{15}εἰ δὲ ἀλλήλους δάκνετε καὶ κατεσθίετε, βλέπετε μὴ ὑπ᾽ ἀλλήλων ἀναλωθῆτε.
\vs{16}Λέγω δέ, Πνεύματι περιπατεῖτε καὶ ἐπιθυμίαν σαρκὸς οὐ μὴ τελέσητε.
\vs{17}ἡ γὰρ σὰρξ ἐπιθυμεῖ κατὰ τοῦ Πνεύματος, τὸ δὲ Πνεῦμα κατὰ τῆς σαρκός, ταῦτα γὰρ ἀλλήλοις ἀντίκειται, ἵνα μὴ ἃ ἐὰν θέλητε ταῦτα ποιῆτε.
\vs{18}εἰ δὲ Πνεύματι ἄγεσθε, οὐκ ἐστὲ ὑπὸ νόμον.
\vs{19}Φανερὰ δέ ἐστιν τὰ ἔργα τῆς σαρκός, ἅτινά ἐστιν πορνεία, ἀκαθαρσία, ἀσέλγεια,
\vs{20}εἰδωλολατρία, φαρμακεία, ἔχθραι, ἔρις, ζῆλος, θυμοί, ἐριθεῖαι, διχοστασίαι, αἱρέσεις,
\vs{21}φθόνοι, μέθαι, κῶμοι καὶ τὰ ὅμοια τούτοις, ἃ προλέγω ὑμῖν, καθὼς προεῖπον ὅτι οἱ τὰ τοιαῦτα πράσσοντες βασιλείαν Θεοῦ οὐ κληρονομήσουσιν.
\vs{22}Ὁ δὲ καρπὸς τοῦ Πνεύματός ἐστιν ἀγάπη χαρά εἰρήνη, μακροθυμία χρηστότης ἀγαθωσύνη, πίστις
\vs{23}πραΰτης ἐγκράτεια· κατὰ τῶν τοιούτων οὐκ ἔστιν νόμος.
\vs{24}Οἱ δὲ τοῦ Χριστοῦ Ἰησοῦ τὴν σάρκα ἐσταύρωσαν σὺν τοῖς παθήμασιν καὶ ταῖς ἐπιθυμίαις.
\vs{25}Εἰ ζῶμεν Πνεύματι, Πνεύματι καὶ στοιχῶμεν.
\vs{26}μὴ γινώμεθα κενόδοξοι, ἀλλήλους προκαλούμενοι, ἀλλήλοις φθονοῦντες.

\ch{6}
Ἀδελφοί, ἐὰν καὶ προλημφθῇ ἄνθρωπος ἔν τινι παραπτώματι, ὑμεῖς οἱ πνευματικοὶ καταρτίζετε τὸν τοιοῦτον ἐν πνεύματι πραΰτητος, σκοπῶν σεαυτόν μὴ καὶ σὺ πειρασθῇς.
\vs{2}Ἀλλήλων τὰ βάρη βαστάζετε καὶ οὕτως ἀναπληρώσετε τὸν νόμον τοῦ Χριστοῦ.
\vs{3}εἰ γὰρ δοκεῖ τις εἶναί τι μηδὲν ὤν, φρεναπατᾷ ἑαυτόν.
\vs{4}Τὸ δὲ ἔργον ἑαυτοῦ δοκιμαζέτω ἕκαστος, καὶ τότε εἰς ἑαυτὸν μόνον τὸ καύχημα ἕξει καὶ οὐκ εἰς τὸν ἕτερον·
\vs{5}ἕκαστος γὰρ τὸ ἴδιον φορτίον βαστάσει.

\vs{6}Κοινωνείτω δὲ ὁ κατηχούμενος τὸν λόγον τῷ κατηχοῦντι ἐν πᾶσιν ἀγαθοῖς.
\vs{7}Μὴ πλανᾶσθε, Θεὸς οὐ μυκτηρίζεται. ὃ γὰρ ἐὰν σπείρῃ ἄνθρωπος, τοῦτο καὶ θερίσει·
\vs{8}ὅτι ὁ σπείρων εἰς τὴν σάρκα ἑαυτοῦ ἐκ τῆς σαρκὸς θερίσει φθοράν, ὁ δὲ σπείρων εἰς τὸ Πνεῦμα ἐκ τοῦ Πνεύματος θερίσει ζωὴν αἰώνιον.
\vs{9}Τὸ δὲ καλὸν ποιοῦντες μὴ ἐνκακῶμεν, καιρῷ γὰρ ἰδίῳ θερίσομεν μὴ ἐκλυόμενοι.
\vs{10}Ἄρα οὖν ὡς καιρὸν ἔχομεν, ἐργαζώμεθα τὸ ἀγαθὸν πρὸς πάντας, μάλιστα δὲ πρὸς τοὺς οἰκείους τῆς πίστεως.

\vs{11}Ἴδετε πηλίκοις ὑμῖν γράμμασιν ἔγραψα τῇ ἐμῇ χειρί.
\vs{12}Ὅσοι θέλουσιν εὐπροσωπῆσαι ἐν σαρκί, οὗτοι ἀναγκάζουσιν ὑμᾶς περιτέμνεσθαι, μόνον ἵνα τῷ σταυρῷ τοῦ Χριστοῦ μὴ διώκωνται.
\vs{13}οὐδὲ γὰρ οἱ περιτεμνόμενοι αὐτοὶ νόμον φυλάσσουσιν ἀλλὰ θέλουσιν ὑμᾶς περιτέμνεσθαι, ἵνα ἐν τῇ ὑμετέρᾳ σαρκὶ καυχήσωνται.
\vs{14}Ἐμοὶ δὲ μὴ γένοιτο καυχᾶσθαι εἰ μὴ ἐν τῷ σταυρῷ τοῦ Κυρίου ἡμῶν Ἰησοῦ Χριστοῦ, δι᾽ οὗ ἐμοὶ κόσμος ἐσταύρωται κἀγὼ κόσμῳ.
\vs{15}οὔτε γὰρ περιτομή τί ἐστιν οὔτε ἀκροβυστία ἀλλὰ καινὴ κτίσις.
\vs{16}Καὶ ὅσοι τῷ κανόνι τούτῳ στοιχήσουσιν, εἰρήνη ἐπ᾽ αὐτοὺς καὶ ἔλεος καὶ ἐπὶ τὸν Ἰσραὴλ τοῦ Θεοῦ.

\vs{17}Τοῦ λοιποῦ κόπους μοι μηδεὶς παρεχέτω· ἐγὼ γὰρ τὰ στίγματα τοῦ Ἰησοῦ ἐν τῷ σώματί μου βαστάζω.

\vs{18}Ἡ χάρις τοῦ Κυρίου ἡμῶν Ἰησοῦ Χριστοῦ μετὰ τοῦ πνεύματος ὑμῶν, ἀδελφοί· Ἀμήν.


\def\book{ΠΡΟΣ ΕΦΕΣΙΟΥΣ}
\biblebook{ΠΡΟΣ ΕΦΕΣΙΟΥΣ}


\lettrine[lines=2, loversize=0.2, nindent=0em, findent=.25em]{\textcolor{bookheadingcolor}{Π}}{αῦλος} ἀπόστολος Χριστοῦ Ἰησοῦ διὰ θελήματος Θεοῦ Τοῖς ἁγίοις τοῖς οὖσιν ἐν Ἐφέσῳ καὶ πιστοῖς ἐν Χριστῷ Ἰησοῦ,
\vs{2}Χάρις ὑμῖν καὶ εἰρήνη ἀπὸ Θεοῦ Πατρὸς ἡμῶν καὶ Κυρίου Ἰησοῦ Χριστοῦ.

\vs{3}Εὐλογητὸς ὁ Θεὸς καὶ Πατὴρ τοῦ Κυρίου ἡμῶν Ἰησοῦ Χριστοῦ, ὁ εὐλογήσας ἡμᾶς ἐν πάσῃ εὐλογίᾳ πνευματικῇ ἐν τοῖς ἐπουρανίοις ἐν Χριστῷ,
\vs{4}καθὼς ἐξελέξατο ἡμᾶς ἐν αὐτῷ πρὸ καταβολῆς κόσμου εἶναι ἡμᾶς ἁγίους καὶ ἀμώμους κατενώπιον αὐτοῦ ἐν ἀγάπῃ,
\vs{5}προορίσας ἡμᾶς εἰς υἱοθεσίαν διὰ Ἰησοῦ Χριστοῦ εἰς αὐτόν, κατὰ τὴν εὐδοκίαν τοῦ θελήματος αὐτοῦ,
\vs{6}εἰς ἔπαινον δόξης τῆς χάριτος αὐτοῦ ἧς ἐχαρίτωσεν ἡμᾶς ἐν τῷ Ἠγαπημένῳ.
\vs{7}ἐν ᾧ ἔχομεν τὴν ἀπολύτρωσιν διὰ τοῦ αἵματος αὐτοῦ, τὴν ἄφεσιν τῶν παραπτωμάτων, κατὰ τὸ πλοῦτος τῆς χάριτος αὐτοῦ
\vs{8}ἧς ἐπερίσσευσεν εἰς ἡμᾶς, ἐν πάσῃ σοφίᾳ καὶ φρονήσει,
\vs{9}γνωρίσας ἡμῖν τὸ μυστήριον τοῦ θελήματος αὐτοῦ, κατὰ τὴν εὐδοκίαν αὐτοῦ ἣν προέθετο ἐν αὐτῷ
\vs{10}εἰς οἰκονομίαν τοῦ πληρώματος τῶν καιρῶν, ἀνακεφαλαιώσασθαι τὰ πάντα ἐν τῷ Χριστῷ, τὰ ἐπὶ τοῖς οὐρανοῖς καὶ τὰ ἐπὶ τῆς γῆς ἐν αὐτῷ.
\vs{11}ἐν ᾧ καὶ ἐκληρώθημεν προορισθέντες κατὰ πρόθεσιν τοῦ τὰ πάντα ἐνεργοῦντος κατὰ τὴν βουλὴν τοῦ θελήματος αὐτοῦ
\vs{12}εἰς τὸ εἶναι ἡμᾶς εἰς ἔπαινον δόξης αὐτοῦ τοὺς προηλπικότας ἐν τῷ Χριστῷ.
\vs{13}ἐν ᾧ καὶ ὑμεῖς ἀκούσαντες τὸν λόγον τῆς ἀληθείας, τὸ εὐαγγέλιον τῆς σωτηρίας ὑμῶν, ἐν ᾧ καὶ πιστεύσαντες ἐσφραγίσθητε τῷ Πνεύματι τῆς ἐπαγγελίας τῷ Ἁγίῳ,
\vs{14}ὅ ἐστιν ἀρραβὼν τῆς κληρονομίας ἡμῶν, εἰς ἀπολύτρωσιν τῆς περιποιήσεως, εἰς ἔπαινον τῆς δόξης αὐτοῦ.

\vs{15}Διὰ τοῦτο κἀγώ ἀκούσας τὴν καθ᾽ ὑμᾶς πίστιν ἐν τῷ Κυρίῳ Ἰησοῦ καὶ τὴν ἀγάπην τὴν εἰς πάντας τοὺς ἁγίους
\vs{16}οὐ παύομαι εὐχαριστῶν ὑπὲρ ὑμῶν μνείαν ποιούμενος ἐπὶ τῶν προσευχῶν μου,
\vs{17}ἵνα ὁ Θεὸς τοῦ Κυρίου ἡμῶν Ἰησοῦ Χριστοῦ, ὁ Πατὴρ τῆς δόξης, δώῃ ὑμῖν πνεῦμα σοφίας καὶ ἀποκαλύψεως ἐν ἐπιγνώσει αὐτοῦ,
\vs{18}πεφωτισμένους τοὺς ὀφθαλμοὺς τῆς καρδίας ὑμῶν εἰς τὸ εἰδέναι ὑμᾶς τίς ἐστιν ἡ ἐλπὶς τῆς κλήσεως αὐτοῦ, τίς ὁ πλοῦτος τῆς δόξης τῆς κληρονομίας αὐτοῦ ἐν τοῖς ἁγίοις,
\vs{19}καὶ τί τὸ ὑπερβάλλον μέγεθος τῆς δυνάμεως αὐτοῦ εἰς ἡμᾶς τοὺς πιστεύοντας κατὰ τὴν ἐνέργειαν τοῦ κράτους τῆς ἰσχύος αὐτοῦ.
\vs{20}ἣν ἐνήργηκεν ἐν τῷ Χριστῷ ἐγείρας αὐτὸν ἐκ νεκρῶν καὶ καθίσας ἐν δεξιᾷ αὐτοῦ ἐν τοῖς ἐπουρανίοις
\vs{21}ὑπεράνω πάσης ἀρχῆς καὶ ἐξουσίας καὶ δυνάμεως καὶ κυριότητος καὶ παντὸς ὀνόματος ὀνομαζομένου, οὐ μόνον ἐν τῷ αἰῶνι τούτῳ ἀλλὰ καὶ ἐν τῷ μέλλοντι·
\vs{22}καὶ πάντα ὑπέταξεν ὑπὸ τοὺς πόδας αὐτοῦ καὶ αὐτὸν ἔδωκεν κεφαλὴν ὑπὲρ πάντα τῇ ἐκκλησίᾳ,
\vs{23}ἥτις ἐστὶν τὸ σῶμα αὐτοῦ, τὸ πλήρωμα τοῦ τὰ πάντα ἐν πᾶσιν πληρουμένου.

\ch{2}
Καὶ ὑμᾶς ὄντας νεκροὺς τοῖς παραπτώμασιν καὶ ταῖς ἁμαρτίαις ὑμῶν,
\vs{2}ἐν αἷς ποτε περιεπατήσατε κατὰ τὸν αἰῶνα τοῦ κόσμου τούτου, κατὰ τὸν ἄρχοντα τῆς ἐξουσίας τοῦ ἀέρος, τοῦ πνεύματος τοῦ νῦν ἐνεργοῦντος ἐν τοῖς υἱοῖς τῆς ἀπειθείας·
\vs{3}ἐν οἷς καὶ ἡμεῖς πάντες ἀνεστράφημέν ποτε ἐν ταῖς ἐπιθυμίαις τῆς σαρκὸς ἡμῶν ποιοῦντες τὰ θελήματα τῆς σαρκὸς καὶ τῶν διανοιῶν, καὶ ἤμεθα τέκνα φύσει ὀργῆς ὡς καὶ οἱ λοιποί·
\vs{4}Ὁ δὲ Θεὸς πλούσιος ὢν ἐν ἐλέει, διὰ τὴν πολλὴν ἀγάπην αὐτοῦ ἣν ἠγάπησεν ἡμᾶς,
\vs{5}καὶ ὄντας ἡμᾶς νεκροὺς τοῖς παραπτώμασιν συνεζωοποίησεν τῷ Χριστῷ,— χάριτί ἐστε σεσῳσμένοι—
\vs{6}καὶ συνήγειρεν καὶ συνεκάθισεν ἐν τοῖς ἐπουρανίοις ἐν Χριστῷ Ἰησοῦ,
\vs{7}ἵνα ἐνδείξηται ἐν τοῖς αἰῶσιν τοῖς ἐπερχομένοις τὸ ὑπερβάλλον πλοῦτος τῆς χάριτος αὐτοῦ ἐν χρηστότητι ἐφ᾽ ἡμᾶς ἐν Χριστῷ Ἰησοῦ.
\vs{8}Τῇ γὰρ χάριτί ἐστε σεσῳσμένοι διὰ πίστεως· καὶ τοῦτο οὐκ ἐξ ὑμῶν, Θεοῦ τὸ δῶρον·
\vs{9}οὐκ ἐξ ἔργων, ἵνα μή τις καυχήσηται.
\vs{10}αὐτοῦ γάρ ἐσμεν ποίημα, κτισθέντες ἐν Χριστῷ Ἰησοῦ ἐπὶ ἔργοις ἀγαθοῖς οἷς προητοίμασεν ὁ Θεὸς, ἵνα ἐν αὐτοῖς περιπατήσωμεν.
\vs{11}Διὸ μνημονεύετε ὅτι ποτὲ ὑμεῖς τὰ ἔθνη ἐν σαρκί, οἱ λεγόμενοι ἀκροβυστία ὑπὸ τῆς λεγομένης περιτομῆς ἐν σαρκὶ χειροποιήτου,
\vs{12}ὅτι ἦτε τῷ καιρῷ ἐκείνῳ χωρὶς Χριστοῦ, ἀπηλλοτριωμένοι τῆς πολιτείας τοῦ Ἰσραὴλ καὶ ξένοι τῶν διαθηκῶν τῆς ἐπαγγελίας, ἐλπίδα μὴ ἔχοντες καὶ ἄθεοι ἐν τῷ κόσμῳ.
\vs{13}νυνὶ δὲ ἐν Χριστῷ Ἰησοῦ ὑμεῖς οἵ ποτε ὄντες μακρὰν ἐγενήθητε ἐγγὺς ἐν τῷ αἵματι τοῦ Χριστοῦ.
\vs{14}Αὐτὸς γάρ ἐστιν ἡ εἰρήνη ἡμῶν, ὁ ποιήσας τὰ ἀμφότερα ἓν καὶ τὸ μεσότοιχον τοῦ φραγμοῦ λύσας, τὴν ἔχθραν ἐν τῇ σαρκὶ αὐτοῦ,
\vs{15}τὸν νόμον τῶν ἐντολῶν ἐν δόγμασιν καταργήσας, ἵνα τοὺς δύο κτίσῃ ἐν αὑτῷ εἰς ἕνα καινὸν ἄνθρωπον ποιῶν εἰρήνην
\vs{16}καὶ ἀποκαταλλάξῃ τοὺς ἀμφοτέρους ἐν ἑνὶ σώματι τῷ Θεῷ διὰ τοῦ σταυροῦ, ἀποκτείνας τὴν ἔχθραν ἐν αὐτῷ.
\vs{17}Καὶ ἐλθὼν εὐηγγελίσατο εἰρήνην ὑμῖν τοῖς μακρὰν καὶ εἰρήνην τοῖς ἐγγύς·
\vs{18}ὅτι δι᾽ αὐτοῦ ἔχομεν τὴν προσαγωγὴν οἱ ἀμφότεροι ἐν ἑνὶ Πνεύματι πρὸς τὸν Πατέρα.
\vs{19}Ἄρα οὖν οὐκέτι ἐστὲ ξένοι καὶ πάροικοι ἀλλὰ ἐστὲ συμπολῖται τῶν ἁγίων καὶ οἰκεῖοι τοῦ Θεοῦ,
\vs{20}ἐποικοδομηθέντες ἐπὶ τῷ θεμελίῳ τῶν ἀποστόλων καὶ προφητῶν, ὄντος ἀκρογωνιαίου αὐτοῦ Χριστοῦ Ἰησοῦ,
\vs{21}ἐν ᾧ πᾶσα οἰκοδομὴ συναρμολογουμένη αὔξει εἰς ναὸν ἅγιον ἐν Κυρίῳ,
\vs{22}ἐν ᾧ καὶ ὑμεῖς συνοικοδομεῖσθε εἰς κατοικητήριον τοῦ Θεοῦ ἐν Πνεύματι.

\ch{3}
Τούτου χάριν ἐγὼ Παῦλος ὁ δέσμιος τοῦ Χριστοῦ Ἰησοῦ ὑπὲρ ὑμῶν τῶν ἐθνῶν—
\vs{2}Εἴ γε ἠκούσατε τὴν οἰκονομίαν τῆς χάριτος τοῦ Θεοῦ τῆς δοθείσης μοι εἰς ὑμᾶς,
\vs{3}ὅτι κατὰ ἀποκάλυψιν ἐγνωρίσθη μοι τὸ μυστήριον, καθὼς προέγραψα ἐν ὀλίγῳ,
\vs{4}πρὸς ὃ δύνασθε ἀναγινώσκοντες νοῆσαι τὴν σύνεσίν μου ἐν τῷ μυστηρίῳ τοῦ Χριστοῦ,
\vs{5}ὃ ἑτέραις γενεαῖς οὐκ ἐγνωρίσθη τοῖς υἱοῖς τῶν ἀνθρώπων ὡς νῦν ἀπεκαλύφθη τοῖς ἁγίοις ἀποστόλοις αὐτοῦ καὶ προφήταις ἐν Πνεύματι,
\vs{6}εἶναι τὰ ἔθνη συνκληρονόμα καὶ σύσσωμα καὶ συμμέτοχα τῆς ἐπαγγελίας ἐν Χριστῷ Ἰησοῦ διὰ τοῦ εὐαγγελίου,
\vs{7}οὗ ἐγενήθην διάκονος κατὰ τὴν δωρεὰν τῆς χάριτος τοῦ Θεοῦ τῆς δοθείσης μοι κατὰ τὴν ἐνέργειαν τῆς δυνάμεως αὐτοῦ.

\vs{8}Ἐμοὶ τῷ ἐλαχιστοτέρῳ πάντων ἁγίων ἐδόθη ἡ χάρις αὕτη, τοῖς ἔθνεσιν εὐαγγελίσασθαι τὸ ἀνεξιχνίαστον πλοῦτος τοῦ Χριστοῦ
\vs{9}καὶ φωτίσαι πάντας τίς ἡ οἰκονομία τοῦ μυστηρίου τοῦ ἀποκεκρυμμένου ἀπὸ τῶν αἰώνων ἐν τῷ Θεῷ τῷ τὰ πάντα κτίσαντι,
\vs{10}ἵνα γνωρισθῇ νῦν ταῖς ἀρχαῖς καὶ ταῖς ἐξουσίαις ἐν τοῖς ἐπουρανίοις διὰ τῆς ἐκκλησίας ἡ πολυποίκιλος σοφία τοῦ Θεοῦ,
\vs{11}κατὰ πρόθεσιν τῶν αἰώνων ἣν ἐποίησεν ἐν τῷ Χριστῷ Ἰησοῦ τῷ Κυρίῳ ἡμῶν,
\vs{12}ἐν ᾧ ἔχομεν τὴν παρρησίαν καὶ προσαγωγὴν ἐν πεποιθήσει διὰ τῆς πίστεως αὐτοῦ.
\vs{13}Διὸ αἰτοῦμαι μὴ ἐνκακεῖν ἐν ταῖς θλίψεσίν μου ὑπὲρ ὑμῶν, ἥτις ἐστὶν δόξα ὑμῶν.

\vs{14}Τούτου χάριν κάμπτω τὰ γόνατά μου πρὸς τὸν Πατέρα,
\vs{15}ἐξ οὗ πᾶσα πατριὰ ἐν οὐρανοῖς καὶ ἐπὶ γῆς ὀνομάζεται,
\vs{16}ἵνα δῷ ὑμῖν κατὰ τὸ πλοῦτος τῆς δόξης αὐτοῦ δυνάμει κραταιωθῆναι διὰ τοῦ Πνεύματος αὐτοῦ εἰς τὸν ἔσω ἄνθρωπον,
\vs{17}κατοικῆσαι τὸν Χριστὸν διὰ τῆς πίστεως ἐν ταῖς καρδίαις ὑμῶν, ἐν ἀγάπῃ ἐρριζωμένοι καὶ τεθεμελιωμένοι,
\vs{18}ἵνα ἐξισχύσητε καταλαβέσθαι σὺν πᾶσιν τοῖς ἁγίοις τί τὸ πλάτος καὶ μῆκος καὶ ὕψος καὶ βάθος,
\vs{19}γνῶναί τε τὴν ὑπερβάλλουσαν τῆς γνώσεως ἀγάπην τοῦ Χριστοῦ, ἵνα πληρωθῆτε εἰς πᾶν τὸ πλήρωμα τοῦ Θεοῦ.
\vs{20}Τῷ δὲ δυναμένῳ ὑπὲρ πάντα ποιῆσαι ὑπερεκπερισσοῦ ὧν αἰτούμεθα ἢ νοοῦμεν κατὰ τὴν δύναμιν τὴν ἐνεργουμένην ἐν ἡμῖν,
\vs{21}αὐτῷ ἡ δόξα ἐν τῇ ἐκκλησίᾳ καὶ ἐν Χριστῷ Ἰησοῦ εἰς πάσας τὰς γενεὰς τοῦ αἰῶνος τῶν αἰώνων, ἀμήν.

\ch{4}
Παρακαλῶ οὖν ὑμᾶς ἐγὼ ὁ δέσμιος ἐν Κυρίῳ ἀξίως περιπατῆσαι τῆς κλήσεως ἧς ἐκλήθητε,
\vs{2}μετὰ πάσης ταπεινοφροσύνης καὶ πραΰτητος, μετὰ μακροθυμίας, ἀνεχόμενοι ἀλλήλων ἐν ἀγάπῃ,
\vs{3}σπουδάζοντες τηρεῖν τὴν ἑνότητα τοῦ Πνεύματος ἐν τῷ συνδέσμῳ τῆς εἰρήνης·
\vs{4}ἓν σῶμα καὶ ἓν Πνεῦμα, καθὼς καὶ ἐκλήθητε ἐν μιᾷ ἐλπίδι τῆς κλήσεως ὑμῶν·
\begin{poetryblock}

\begin{quote} \vs{5}εἷς Κύριος, μία πίστις, ἓν βάπτισμα,\end{quote}

\begin{quote} \vs{6}εἷς Θεὸς καὶ Πατὴρ πάντων,\end{quote} 

\begin{quote}ὁ ἐπὶ πάντων καὶ διὰ πάντων καὶ ἐν πᾶσιν.\end{quote}
\end{poetryblock}
\vs{7}Ἑνὶ δὲ ἑκάστῳ ἡμῶν ἐδόθη ἡ χάρις κατὰ τὸ μέτρον τῆς δωρεᾶς τοῦ Χριστοῦ.
\vs{8}διὸ λέγει· 
\begin{poetryblock}

\begin{quote}Ἀναβὰς εἰς ὕψος ᾐχμαλώτευσεν αἰχμαλωσίαν,\end{quote} 

\begin{quote}ἔδωκεν δόματα τοῖς ἀνθρώποις.\end{quote}
\end{poetryblock}

\vs{9}Τὸ δὲ Ἀνέβη τί ἐστιν, εἰ μὴ ὅτι καὶ κατέβη εἰς τὰ κατώτερα μέρη τῆς γῆς;
\vs{10}ὁ καταβὰς αὐτός ἐστιν καὶ ὁ ἀναβὰς ὑπεράνω πάντων τῶν οὐρανῶν, ἵνα πληρώσῃ τὰ πάντα.
\vs{11}Καὶ αὐτὸς ἔδωκεν τοὺς μὲν ἀποστόλους, τοὺς δὲ προφήτας, τοὺς δὲ εὐαγγελιστάς, τοὺς δὲ ποιμένας καὶ διδασκάλους,
\vs{12}πρὸς τὸν καταρτισμὸν τῶν ἁγίων εἰς ἔργον διακονίας, εἰς οἰκοδομὴν τοῦ σώματος τοῦ Χριστοῦ,
\vs{13}μέχρι καταντήσωμεν οἱ πάντες εἰς τὴν ἑνότητα τῆς πίστεως καὶ τῆς ἐπιγνώσεως τοῦ Υἱοῦ τοῦ Θεοῦ, εἰς ἄνδρα τέλειον, εἰς μέτρον ἡλικίας τοῦ πληρώματος τοῦ Χριστοῦ,
\vs{14}ἵνα μηκέτι ὦμεν νήπιοι, κλυδωνιζόμενοι καὶ περιφερόμενοι παντὶ ἀνέμῳ τῆς διδασκαλίας ἐν τῇ κυβείᾳ τῶν ἀνθρώπων, ἐν πανουργίᾳ πρὸς τὴν μεθοδείαν τῆς πλάνης,
\vs{15}ἀληθεύοντες δὲ ἐν ἀγάπῃ αὐξήσωμεν εἰς αὐτὸν τὰ πάντα, ὅς ἐστιν ἡ κεφαλή, Χριστός,
\vs{16}ἐξ οὗ πᾶν τὸ σῶμα συναρμολογούμενον καὶ συμβιβαζόμενον διὰ πάσης ἁφῆς τῆς ἐπιχορηγίας κατ᾽ ἐνέργειαν ἐν μέτρῳ ἑνὸς ἑκάστου μέρους τὴν αὔξησιν τοῦ σώματος ποιεῖται εἰς οἰκοδομὴν ἑαυτοῦ ἐν ἀγάπῃ.

\vs{17}Τοῦτο οὖν λέγω καὶ μαρτύρομαι ἐν Κυρίῳ, μηκέτι ὑμᾶς περιπατεῖν, καθὼς καὶ τὰ ἔθνη περιπατεῖ ἐν ματαιότητι τοῦ νοὸς αὐτῶν,
\vs{18}ἐσκοτωμένοι τῇ διανοίᾳ ὄντες, ἀπηλλοτριωμένοι τῆς ζωῆς τοῦ Θεοῦ διὰ τὴν ἄγνοιαν τὴν οὖσαν ἐν αὐτοῖς, διὰ τὴν πώρωσιν τῆς καρδίας αὐτῶν,
\vs{19}οἵτινες ἀπηλγηκότες ἑαυτοὺς παρέδωκαν τῇ ἀσελγείᾳ εἰς ἐργασίαν ἀκαθαρσίας πάσης ἐν πλεονεξίᾳ.

\vs{20}Ὑμεῖς δὲ οὐχ οὕτως ἐμάθετε τὸν Χριστόν,
\vs{21}εἴ γε αὐτὸν ἠκούσατε καὶ ἐν αὐτῷ ἐδιδάχθητε, καθώς ἐστιν ἀλήθεια ἐν τῷ Ἰησοῦ,
\vs{22}ἀποθέσθαι ὑμᾶς κατὰ τὴν προτέραν ἀναστροφὴν τὸν παλαιὸν ἄνθρωπον τὸν φθειρόμενον κατὰ τὰς ἐπιθυμίας τῆς ἀπάτης,
\vs{23}ἀνανεοῦσθαι δὲ τῷ πνεύματι τοῦ νοὸς ὑμῶν
\vs{24}καὶ ἐνδύσασθαι τὸν καινὸν ἄνθρωπον τὸν κατὰ Θεὸν κτισθέντα ἐν δικαιοσύνῃ καὶ ὁσιότητι τῆς ἀληθείας.

\vs{25}Διὸ ἀποθέμενοι τὸ ψεῦδος λαλεῖτε ἀλήθειαν ἕκαστος μετὰ τοῦ πλησίον αὐτοῦ, ὅτι ἐσμὲν ἀλλήλων μέλη.
\vs{26}Ὀργίζεσθε καὶ μὴ ἁμαρτάνετε· ὁ ἥλιος μὴ ἐπιδυέτω ἐπὶ τῷ παροργισμῷ ὑμῶν,
\vs{27}μηδὲ δίδοτε τόπον τῷ διαβόλῳ.
\vs{28}Ὁ κλέπτων μηκέτι κλεπτέτω, μᾶλλον δὲ κοπιάτω ἐργαζόμενος ταῖς ἰδίαις χερσὶν τὸ ἀγαθόν, ἵνα ἔχῃ μεταδιδόναι τῷ χρείαν ἔχοντι.
\vs{29}Πᾶς λόγος σαπρὸς ἐκ τοῦ στόματος ὑμῶν μὴ ἐκπορευέσθω, ἀλλὰ εἴ τις ἀγαθὸς πρὸς οἰκοδομὴν τῆς χρείας, ἵνα δῷ χάριν τοῖς ἀκούουσιν.
\vs{30}Καὶ μὴ λυπεῖτε τὸ Πνεῦμα τὸ Ἅγιον τοῦ Θεοῦ, ἐν ᾧ ἐσφραγίσθητε εἰς ἡμέραν ἀπολυτρώσεως.
\vs{31}Πᾶσα πικρία καὶ θυμὸς καὶ ὀργὴ καὶ κραυγὴ καὶ βλασφημία ἀρθήτω ἀφ᾽ ὑμῶν σὺν πάσῃ κακίᾳ.
\vs{32}γίνεσθε δὲ εἰς ἀλλήλους χρηστοί, εὔσπλαγχνοι, χαριζόμενοι ἑαυτοῖς, καθὼς καὶ ὁ Θεὸς ἐν Χριστῷ ἐχαρίσατο ὑμῖν.

\ch{5}
Γίνεσθε οὖν μιμηταὶ τοῦ Θεοῦ ὡς τέκνα ἀγαπητά
\vs{2}καὶ περιπατεῖτε ἐν ἀγάπῃ, καθὼς καὶ ὁ Χριστὸς ἠγάπησεν ἡμᾶς καὶ παρέδωκεν ἑαυτὸν ὑπὲρ ἡμῶν προσφορὰν καὶ θυσίαν τῷ Θεῷ εἰς ὀσμὴν εὐωδίας.

\vs{3}Πορνεία δὲ καὶ ἀκαθαρσία πᾶσα ἢ πλεονεξία μηδὲ ὀνομαζέσθω ἐν ὑμῖν, καθὼς πρέπει ἁγίοις,
\vs{4}καὶ αἰσχρότης καὶ μωρολογία ἢ εὐτραπελία, ἃ οὐκ ἀνῆκεν, ἀλλὰ μᾶλλον εὐχαριστία.
\vs{5}τοῦτο γὰρ ἴστε γινώσκοντες, ὅτι πᾶς πόρνος ἢ ἀκάθαρτος ἢ πλεονέκτης, ὅ ἐστιν εἰδωλολάτρης, οὐκ ἔχει κληρονομίαν ἐν τῇ βασιλείᾳ τοῦ Χριστοῦ καὶ Θεοῦ.
\vs{6}Μηδεὶς ὑμᾶς ἀπατάτω κενοῖς λόγοις· διὰ ταῦτα γὰρ ἔρχεται ἡ ὀργὴ τοῦ Θεοῦ ἐπὶ τοὺς υἱοὺς τῆς ἀπειθείας.
\vs{7}μὴ οὖν γίνεσθε συμμέτοχοι αὐτῶν·
\vs{8}Ἦτε γάρ ποτε σκότος, νῦν δὲ φῶς ἐν Κυρίῳ· ὡς τέκνα φωτὸς περιπατεῖτε—
\vs{9}ὁ γὰρ καρπὸς τοῦ φωτὸς ἐν πάσῃ ἀγαθωσύνῃ καὶ δικαιοσύνῃ καὶ ἀληθείᾳ—
\vs{10}δοκιμάζοντες τί ἐστιν εὐάρεστον τῷ Κυρίῳ,
\vs{11}Καὶ μὴ συνκοινωνεῖτε τοῖς ἔργοις τοῖς ἀκάρποις τοῦ σκότους, μᾶλλον δὲ καὶ ἐλέγχετε.
\vs{12}τὰ γὰρ κρυφῇ γινόμενα ὑπ᾽ αὐτῶν αἰσχρόν ἐστιν καὶ λέγειν,
\vs{13}τὰ δὲ πάντα ἐλεγχόμενα ὑπὸ τοῦ φωτὸς φανεροῦται,
\vs{14}πᾶν γὰρ τὸ φανερούμενον φῶς ἐστιν. διὸ λέγει· 
\begin{poetryblock}

\begin{quote}Ἔγειρε, ὁ καθεύδων,\end{quote} 

\begin{quote}καὶ ἀνάστα ἐκ τῶν νεκρῶν,\end{quote} 

\begin{quote}καὶ ἐπιφαύσει σοι ὁ Χριστός.\end{quote}
\end{poetryblock}

\vs{15}Βλέπετε οὖν ἀκριβῶς πῶς περιπατεῖτε μὴ ὡς ἄσοφοι ἀλλ᾽ ὡς σοφοί,
\vs{16}ἐξαγοραζόμενοι τὸν καιρόν, ὅτι αἱ ἡμέραι πονηραί εἰσιν.
\vs{17}διὰ τοῦτο μὴ γίνεσθε ἄφρονες, ἀλλὰ συνίετε τί τὸ θέλημα τοῦ Κυρίου.
\vs{18}καὶ μὴ μεθύσκεσθε οἴνῳ, ἐν ᾧ ἐστιν ἀσωτία, ἀλλὰ πληροῦσθε ἐν Πνεύματι,
\vs{19}λαλοῦντες ἑαυτοῖς ἐν ψαλμοῖς καὶ ὕμνοις καὶ ᾠδαῖς πνευματικαῖς, ᾄδοντες καὶ ψάλλοντες τῇ καρδίᾳ ὑμῶν τῷ Κυρίῳ,
\vs{20}εὐχαριστοῦντες πάντοτε ὑπὲρ πάντων ἐν ὀνόματι τοῦ Κυρίου ἡμῶν Ἰησοῦ Χριστοῦ τῷ Θεῷ καὶ Πατρί.
\vs{21}Ὑποτασσόμενοι ἀλλήλοις ἐν φόβῳ Χριστοῦ,
\vs{22}Αἱ γυναῖκες τοῖς ἰδίοις ἀνδράσιν ὡς τῷ Κυρίῳ,
\vs{23}ὅτι ἀνήρ ἐστιν κεφαλὴ τῆς γυναικὸς ὡς καὶ ὁ Χριστὸς κεφαλὴ τῆς ἐκκλησίας, αὐτὸς σωτὴρ τοῦ σώματος·
\vs{24}ἀλλὰ ὡς ἡ ἐκκλησία ὑποτάσσεται τῷ Χριστῷ, οὕτως καὶ αἱ γυναῖκες τοῖς ἀνδράσιν ἐν παντί.

\vs{25}Οἱ ἄνδρες, ἀγαπᾶτε τὰς γυναῖκας, καθὼς καὶ ὁ Χριστὸς ἠγάπησεν τὴν ἐκκλησίαν καὶ ἑαυτὸν παρέδωκεν ὑπὲρ αὐτῆς,
\vs{26}ἵνα αὐτὴν ἁγιάσῃ καθαρίσας τῷ λουτρῷ τοῦ ὕδατος ἐν ῥήματι,
\vs{27}ἵνα παραστήσῃ αὐτὸς ἑαυτῷ ἔνδοξον τὴν ἐκκλησίαν, μὴ ἔχουσαν σπίλον ἢ ῥυτίδα ἤ τι τῶν τοιούτων, ἀλλ᾽ ἵνα ᾖ ἁγία καὶ ἄμωμος.
\vs{28}Οὕτως ὀφείλουσιν καὶ οἱ ἄνδρες ἀγαπᾶν τὰς ἑαυτῶν γυναῖκας ὡς τὰ ἑαυτῶν σώματα. ὁ ἀγαπῶν τὴν ἑαυτοῦ γυναῖκα ἑαυτὸν ἀγαπᾷ.
\vs{29}οὐδεὶς γάρ ποτε τὴν ἑαυτοῦ σάρκα ἐμίσησεν ἀλλὰ ἐκτρέφει καὶ θάλπει αὐτήν, καθὼς καὶ ὁ Χριστὸς τὴν ἐκκλησίαν,
\vs{30}ὅτι μέλη ἐσμὲν τοῦ σώματος αὐτοῦ.
\vs{31}Ἀντὶ τούτου καταλείψει ἄνθρωπος τὸν πατέρα καὶ τὴν μητέρα καὶ προσκολληθήσεται πρὸς τὴν γυναῖκα αὐτοῦ, καὶ ἔσονται οἱ δύο εἰς σάρκα μίαν.
\vs{32}τὸ μυστήριον τοῦτο μέγα ἐστίν· ἐγὼ δὲ λέγω εἰς Χριστὸν καὶ εἰς τὴν ἐκκλησίαν.
\vs{33}πλὴν καὶ ὑμεῖς οἱ καθ᾽ ἕνα, ἕκαστος τὴν ἑαυτοῦ γυναῖκα οὕτως ἀγαπάτω ὡς ἑαυτόν, ἡ δὲ γυνὴ ἵνα φοβῆται τὸν ἄνδρα.

\ch{6}
Τὰ τέκνα, ὑπακούετε τοῖς γονεῦσιν ὑμῶν ἐν Κυρίῳ· τοῦτο γάρ ἐστιν δίκαιον.
\vs{2}Τίμα τὸν πατέρα σου καὶ τὴν μητέρα, ἥτις ἐστὶν ἐντολὴ πρώτη ἐν ἐπαγγελίᾳ,
\vs{3}Ἵνα εὖ σοι γένηται καὶ ἔσῃ μακροχρόνιος ἐπὶ τῆς γῆς.
\vs{4}Καὶ οἱ πατέρες, μὴ παροργίζετε τὰ τέκνα ὑμῶν ἀλλὰ ἐκτρέφετε αὐτὰ ἐν παιδείᾳ καὶ νουθεσίᾳ Κυρίου.

\vs{5}Οἱ δοῦλοι, ὑπακούετε τοῖς κατὰ σάρκα κυρίοις μετὰ φόβου καὶ τρόμου ἐν ἁπλότητι τῆς καρδίας ὑμῶν ὡς τῷ Χριστῷ,
\vs{6}μὴ κατ᾽ ὀφθαλμοδουλίαν ὡς ἀνθρωπάρεσκοι ἀλλ᾽ ὡς δοῦλοι Χριστοῦ ποιοῦντες τὸ θέλημα τοῦ Θεοῦ ἐκ ψυχῆς,
\vs{7}μετ᾽ εὐνοίας δουλεύοντες ὡς τῷ Κυρίῳ καὶ οὐκ ἀνθρώποις,
\vs{8}εἰδότες ὅτι ἕκαστος ἐάν τι ποιήσῃ ἀγαθόν, τοῦτο κομίσεται παρὰ Κυρίου εἴτε δοῦλος εἴτε ἐλεύθερος.
\vs{9}Καὶ οἱ κύριοι, τὰ αὐτὰ ποιεῖτε πρὸς αὐτούς, ἀνιέντες τὴν ἀπειλήν, εἰδότες ὅτι καὶ αὐτῶν καὶ ὑμῶν ὁ Κύριός ἐστιν ἐν οὐρανοῖς καὶ προσωπολημψία οὐκ ἔστιν παρ᾽ αὐτῷ.

\vs{10}Τοῦ λοιποῦ, ἐνδυναμοῦσθε ἐν Κυρίῳ καὶ ἐν τῷ κράτει τῆς ἰσχύος αὐτοῦ.
\vs{11}ἐνδύσασθε τὴν πανοπλίαν τοῦ Θεοῦ πρὸς τὸ δύνασθαι ὑμᾶς στῆναι πρὸς τὰς μεθοδείας τοῦ διαβόλου·
\vs{12}ὅτι οὐκ ἔστιν ἡμῖν ἡ πάλη πρὸς αἷμα καὶ σάρκα ἀλλὰ πρὸς τὰς ἀρχάς, πρὸς τὰς ἐξουσίας, πρὸς τοὺς κοσμοκράτορας τοῦ σκότους τούτου, πρὸς τὰ πνευματικὰ τῆς πονηρίας ἐν τοῖς ἐπουρανίοις.
\vs{13}Διὰ τοῦτο ἀναλάβετε τὴν πανοπλίαν τοῦ Θεοῦ, ἵνα δυνηθῆτε ἀντιστῆναι ἐν τῇ ἡμέρᾳ τῇ πονηρᾷ καὶ ἅπαντα κατεργασάμενοι στῆναι.
\vs{14}στῆτε οὖν περιζωσάμενοι τὴν ὀσφὺν ὑμῶν ἐν ἀληθείᾳ καὶ ἐνδυσάμενοι τὸν θώρακα τῆς δικαιοσύνης
\vs{15}καὶ ὑποδησάμενοι τοὺς πόδας ἐν ἑτοιμασίᾳ τοῦ εὐαγγελίου τῆς εἰρήνης,
\vs{16}ἐν πᾶσιν ἀναλαβόντες τὸν θυρεὸν τῆς πίστεως, ἐν ᾧ δυνήσεσθε πάντα τὰ βέλη τοῦ πονηροῦ τὰ πεπυρωμένα σβέσαι·
\vs{17}καὶ τὴν περικεφαλαίαν τοῦ σωτηρίου δέξασθε καὶ τὴν μάχαιραν τοῦ Πνεύματος, ὅ ἐστιν ῥῆμα Θεοῦ.
\vs{18}διὰ πάσης προσευχῆς καὶ δεήσεως προσευχόμενοι ἐν παντὶ καιρῷ ἐν Πνεύματι, καὶ εἰς αὐτὸ ἀγρυπνοῦντες ἐν πάσῃ προσκαρτερήσει καὶ δεήσει περὶ πάντων τῶν ἁγίων
\vs{19}καὶ ὑπὲρ ἐμοῦ, ἵνα μοι δοθῇ λόγος ἐν ἀνοίξει τοῦ στόματός μου, ἐν παρρησίᾳ γνωρίσαι τὸ μυστήριον τοῦ εὐαγγελίου,
\vs{20}ὑπὲρ οὗ πρεσβεύω ἐν ἁλύσει, ἵνα ἐν αὐτῷ παρρησιάσωμαι ὡς δεῖ με λαλῆσαι.

\vs{21}Ἵνα δὲ εἰδῆτε καὶ ὑμεῖς τὰ κατ᾽ ἐμέ, τί πράσσω, πάντα γνωρίσει ὑμῖν Τυχικὸς ὁ ἀγαπητὸς ἀδελφὸς καὶ πιστὸς διάκονος ἐν Κυρίῳ,
\vs{22}ὃν ἔπεμψα πρὸς ὑμᾶς εἰς αὐτὸ τοῦτο, ἵνα γνῶτε τὰ περὶ ἡμῶν καὶ παρακαλέσῃ τὰς καρδίας ὑμῶν.

\vs{23}Εἰρήνη τοῖς ἀδελφοῖς καὶ ἀγάπη μετὰ πίστεως ἀπὸ Θεοῦ Πατρὸς καὶ Κυρίου Ἰησοῦ Χριστοῦ.
\vs{24}Ἡ χάρις μετὰ πάντων τῶν ἀγαπώντων τὸν Κύριον ἡμῶν Ἰησοῦν Χριστὸν ἐν ἀφθαρσίᾳ.


\def\book{ΠΡΟΣ ΦΙΛΙΠΠΗΣΙΟΥΣ}
\biblebook{ΠΡΟΣ ΦΙΛΙΠΠΗΣΙΟΥΣ}


\lettrine[lines=2, loversize=0.2, nindent=0em, findent=.25em]{\textcolor{bookheadingcolor}{Π}}{αῦλος} καὶ Τιμόθεος δοῦλοι Χριστοῦ Ἰησοῦ Πᾶσιν τοῖς ἁγίοις ἐν Χριστῷ Ἰησοῦ τοῖς οὖσιν ἐν Φιλίπποις σὺν ἐπισκόποις καὶ διακόνοις,
\vs{2}Χάρις ὑμῖν καὶ εἰρήνη ἀπὸ Θεοῦ Πατρὸς ἡμῶν καὶ Κυρίου Ἰησοῦ Χριστοῦ.

\vs{3}Εὐχαριστῶ τῷ Θεῷ μου ἐπὶ πάσῃ τῇ μνείᾳ ὑμῶν
\vs{4}πάντοτε ἐν πάσῃ δεήσει μου ὑπὲρ πάντων ὑμῶν, μετὰ χαρᾶς τὴν δέησιν ποιούμενος,
\vs{5}ἐπὶ τῇ κοινωνίᾳ ὑμῶν εἰς τὸ εὐαγγέλιον ἀπὸ τῆς πρώτης ἡμέρας ἄχρι τοῦ νῦν,
\vs{6}πεποιθὼς αὐτὸ τοῦτο, ὅτι ὁ ἐναρξάμενος ἐν ὑμῖν ἔργον ἀγαθὸν ἐπιτελέσει ἄχρι ἡμέρας Χριστοῦ Ἰησοῦ·
\vs{7}Καθώς ἐστιν δίκαιον ἐμοὶ τοῦτο φρονεῖν ὑπὲρ πάντων ὑμῶν διὰ τὸ ἔχειν με ἐν τῇ καρδίᾳ ὑμᾶς, ἔν τε τοῖς δεσμοῖς μου καὶ ἐν τῇ ἀπολογίᾳ καὶ βεβαιώσει τοῦ εὐαγγελίου συνκοινωνούς μου τῆς χάριτος πάντας ὑμᾶς ὄντας.
\vs{8}μάρτυς γάρ μου ὁ Θεός ὡς ἐπιποθῶ πάντας ὑμᾶς ἐν σπλάγχνοις Χριστοῦ Ἰησοῦ.
\vs{9}Καὶ τοῦτο προσεύχομαι, ἵνα ἡ ἀγάπη ὑμῶν ἔτι μᾶλλον καὶ μᾶλλον περισσεύῃ ἐν ἐπιγνώσει καὶ πάσῃ αἰσθήσει
\vs{10}εἰς τὸ δοκιμάζειν ὑμᾶς τὰ διαφέροντα, ἵνα ἦτε εἰλικρινεῖς καὶ ἀπρόσκοποι εἰς ἡμέραν Χριστοῦ,
\vs{11}πεπληρωμένοι καρπὸν δικαιοσύνης τὸν διὰ Ἰησοῦ Χριστοῦ εἰς δόξαν καὶ ἔπαινον Θεοῦ.

\vs{12}Γινώσκειν δὲ ὑμᾶς βούλομαι, ἀδελφοί, ὅτι τὰ κατ᾽ ἐμὲ μᾶλλον εἰς προκοπὴν τοῦ εὐαγγελίου ἐλήλυθεν,
\vs{13}ὥστε τοὺς δεσμούς μου φανεροὺς ἐν Χριστῷ γενέσθαι ἐν ὅλῳ τῷ πραιτωρίῳ καὶ τοῖς λοιποῖς πᾶσιν,
\vs{14}καὶ τοὺς πλείονας τῶν ἀδελφῶν ἐν Κυρίῳ πεποιθότας τοῖς δεσμοῖς μου περισσοτέρως τολμᾶν ἀφόβως τὸν λόγον λαλεῖν.
\vs{15}Τινὲς μὲν καὶ διὰ φθόνον καὶ ἔριν, τινὲς δὲ καὶ δι᾽ εὐδοκίαν τὸν Χριστὸν κηρύσσουσιν·
\vs{16}οἱ μὲν ἐξ ἀγάπης, εἰδότες ὅτι εἰς ἀπολογίαν τοῦ εὐαγγελίου κεῖμαι,
\vs{17}οἱ δὲ ἐξ ἐριθείας τὸν Χριστὸν καταγγέλλουσιν, οὐχ ἁγνῶς, οἰόμενοι θλῖψιν ἐγείρειν τοῖς δεσμοῖς μου.
\vs{18}Τί γάρ; πλὴν ὅτι παντὶ τρόπῳ, εἴτε προφάσει εἴτε ἀληθείᾳ, Χριστὸς καταγγέλλεται, καὶ ἐν τούτῳ χαίρω.

ἀλλὰ καὶ χαρήσομαι,
\vs{19}οἶδα γὰρ ὅτι τοῦτό μοι ἀποβήσεται εἰς σωτηρίαν διὰ τῆς ὑμῶν δεήσεως καὶ ἐπιχορηγίας τοῦ Πνεύματος Ἰησοῦ Χριστοῦ
\vs{20}κατὰ τὴν ἀποκαραδοκίαν καὶ ἐλπίδα μου, ὅτι ἐν οὐδενὶ αἰσχυνθήσομαι ἀλλ᾽ ἐν πάσῃ παρρησίᾳ ὡς πάντοτε καὶ νῦν μεγαλυνθήσεται Χριστὸς ἐν τῷ σώματί μου, εἴτε διὰ ζωῆς εἴτε διὰ θανάτου.
\vs{21}Ἐμοὶ γὰρ τὸ ζῆν Χριστὸς καὶ τὸ ἀποθανεῖν κέρδος.
\vs{22}εἰ δὲ τὸ ζῆν ἐν σαρκί, τοῦτό μοι καρπὸς ἔργου, καὶ τί αἱρήσομαι οὐ γνωρίζω.
\vs{23}συνέχομαι δὲ ἐκ τῶν δύο, τὴν ἐπιθυμίαν ἔχων εἰς τὸ ἀναλῦσαι καὶ σὺν Χριστῷ εἶναι, πολλῷ γὰρ μᾶλλον κρεῖσσον·
\vs{24}τὸ δὲ ἐπιμένειν ἐν τῇ σαρκὶ ἀναγκαιότερον δι᾽ ὑμᾶς.
\vs{25}Καὶ τοῦτο πεποιθὼς οἶδα ὅτι μενῶ καὶ παραμενῶ πᾶσιν ὑμῖν εἰς τὴν ὑμῶν προκοπὴν καὶ χαρὰν τῆς πίστεως,
\vs{26}ἵνα τὸ καύχημα ὑμῶν περισσεύῃ ἐν Χριστῷ Ἰησοῦ ἐν ἐμοὶ διὰ τῆς ἐμῆς παρουσίας πάλιν πρὸς ὑμᾶς.

\vs{27}Μόνον ἀξίως τοῦ εὐαγγελίου τοῦ Χριστοῦ πολιτεύεσθε, ἵνα εἴτε ἐλθὼν καὶ ἰδὼν ὑμᾶς εἴτε ἀπὼν ἀκούω τὰ περὶ ὑμῶν, ὅτι στήκετε ἐν ἑνὶ πνεύματι, μιᾷ ψυχῇ συναθλοῦντες τῇ πίστει τοῦ εὐαγγελίου
\vs{28}καὶ μὴ πτυρόμενοι ἐν μηδενὶ ὑπὸ τῶν ἀντικειμένων, ἥτις ἐστὶν αὐτοῖς ἔνδειξις ἀπωλείας, ὑμῶν δὲ σωτηρίας, καὶ τοῦτο ἀπὸ Θεοῦ·
\vs{29}ὅτι ὑμῖν ἐχαρίσθη τὸ ὑπὲρ Χριστοῦ, οὐ μόνον τὸ εἰς αὐτὸν πιστεύειν ἀλλὰ καὶ τὸ ὑπὲρ αὐτοῦ πάσχειν,
\vs{30}τὸν αὐτὸν ἀγῶνα ἔχοντες, οἷον εἴδετε ἐν ἐμοὶ καὶ νῦν ἀκούετε ἐν ἐμοί.

\ch{2}
Εἴ τις οὖν παράκλησις ἐν Χριστῷ, εἴ τι παραμύθιον ἀγάπης, εἴ τις κοινωνία Πνεύματος, εἴ τις σπλάγχνα καὶ οἰκτιρμοί,
\vs{2}πληρώσατέ μου τὴν χαρὰν ἵνα τὸ αὐτὸ φρονῆτε, τὴν αὐτὴν ἀγάπην ἔχοντες, σύμψυχοι, τὸ ἓν φρονοῦντες,
\vs{3}μηδὲν κατ᾽ ἐριθείαν μηδὲ κατὰ κενοδοξίαν ἀλλὰ τῇ ταπεινοφροσύνῃ ἀλλήλους ἡγούμενοι ὑπερέχοντας ἑαυτῶν,
\vs{4}μὴ τὰ ἑαυτῶν ἕκαστος σκοποῦντες ἀλλὰ καὶ τὰ ἑτέρων ἕκαστοι.

\vs{5}Τοῦτο φρονεῖτε ἐν ὑμῖν ὃ καὶ ἐν Χριστῷ Ἰησοῦ,
\begin{poetryblock}

\begin{quote} \vs{6}Ὃς ἐν μορφῇ Θεοῦ ὑπάρχων\end{quote} 

\begin{quote}οὐχ ἁρπαγμὸν ἡγήσατο\end{quote} 

\begin{quote}τὸ εἶναι ἴσα Θεῷ,\end{quote}

\begin{quote} \vs{7}ἀλλὰ ἑαυτὸν ἐκένωσεν\end{quote} 

\begin{quote}μορφὴν δούλου λαβών,\end{quote} 

\begin{quote}ἐν ὁμοιώματι ἀνθρώπων γενόμενος·\end{quote} 

\begin{quote}καὶ σχήματι εὑρεθεὶς ὡς ἄνθρωπος\end{quote}

\begin{quote} \vs{8}ἐταπείνωσεν ἑαυτὸν\end{quote} 

\begin{quote}γενόμενος ὑπήκοος μέχρι θανάτου,\end{quote} 

\begin{quote}θανάτου δὲ σταυροῦ.\end{quote}

\begin{quote} \vs{9}Διὸ καὶ ὁ Θεὸς αὐτὸν ὑπερύψωσεν\end{quote} 

\begin{quote}καὶ ἐχαρίσατο αὐτῷ τὸ ὄνομα\end{quote} 

\begin{quote}τὸ ὑπὲρ πᾶν ὄνομα,\end{quote}

\begin{quote} \vs{10}ἵνα ἐν τῷ ὀνόματι Ἰησοῦ\end{quote} 

\begin{quote}πᾶν γόνυ κάμψῃ\end{quote} 

\begin{quote}ἐπουρανίων καὶ ἐπιγείων καὶ καταχθονίων\end{quote}

\begin{quote} \vs{11}καὶ πᾶσα γλῶσσα ἐξομολογήσηται ὅτι\end{quote} 

\begin{quote}ΚΥΡΙΟΣ ΙΗΣΟΥΣ ΧΡΙΣΤΟΣ\end{quote} 

\begin{quote}εἰς δόξαν Θεοῦ Πατρός.\end{quote}
\end{poetryblock}

\vs{12}Ὥστε, ἀγαπητοί μου, καθὼς πάντοτε ὑπηκούσατε, μὴ ὡς ἐν τῇ παρουσίᾳ μου μόνον ἀλλὰ νῦν πολλῷ μᾶλλον ἐν τῇ ἀπουσίᾳ μου, μετὰ φόβου καὶ τρόμου τὴν ἑαυτῶν σωτηρίαν κατεργάζεσθε·
\vs{13}Θεὸς γάρ ἐστιν ὁ ἐνεργῶν ἐν ὑμῖν καὶ τὸ θέλειν καὶ τὸ ἐνεργεῖν ὑπὲρ τῆς εὐδοκίας.
\vs{14}Πάντα ποιεῖτε χωρὶς γογγυσμῶν καὶ διαλογισμῶν,
\vs{15}ἵνα γένησθε ἄμεμπτοι καὶ ἀκέραιοι, τέκνα Θεοῦ ἄμωμα μέσον γενεᾶς σκολιᾶς καὶ διεστραμμένης, ἐν οἷς φαίνεσθε ὡς φωστῆρες ἐν κόσμῳ,
\vs{16}λόγον ζωῆς ἐπέχοντες, εἰς καύχημα ἐμοὶ εἰς ἡμέραν Χριστοῦ, ὅτι οὐκ εἰς κενὸν ἔδραμον οὐδὲ εἰς κενὸν ἐκοπίασα.
\vs{17}Ἀλλὰ εἰ καὶ σπένδομαι ἐπὶ τῇ θυσίᾳ καὶ λειτουργίᾳ τῆς πίστεως ὑμῶν, χαίρω καὶ συνχαίρω πᾶσιν ὑμῖν·
\vs{18}τὸ δὲ αὐτὸ καὶ ὑμεῖς χαίρετε καὶ συνχαίρετέ μοι.

\vs{19}Ἐλπίζω δὲ ἐν Κυρίῳ Ἰησοῦ Τιμόθεον ταχέως πέμψαι ὑμῖν, ἵνα κἀγὼ εὐψυχῶ γνοὺς τὰ περὶ ὑμῶν.
\vs{20}οὐδένα γὰρ ἔχω ἰσόψυχον, ὅστις γνησίως τὰ περὶ ὑμῶν μεριμνήσει·
\vs{21}οἱ πάντες γὰρ τὰ ἑαυτῶν ζητοῦσιν, οὐ τὰ Ἰησοῦ Χριστοῦ.
\vs{22}τὴν δὲ δοκιμὴν αὐτοῦ γινώσκετε, ὅτι ὡς πατρὶ τέκνον σὺν ἐμοὶ ἐδούλευσεν εἰς τὸ εὐαγγέλιον.
\vs{23}Τοῦτον μὲν οὖν ἐλπίζω πέμψαι ὡς ἂν ἀφίδω τὰ περὶ ἐμὲ ἐξαυτῆς·
\vs{24}πέποιθα δὲ ἐν Κυρίῳ ὅτι καὶ αὐτὸς ταχέως ἐλεύσομαι.
\vs{25}Ἀναγκαῖον δὲ ἡγησάμην Ἐπαφρόδιτον τὸν ἀδελφὸν καὶ συνεργὸν καὶ συστρατιώτην μου, ὑμῶν δὲ ἀπόστολον καὶ λειτουργὸν τῆς χρείας μου, πέμψαι πρὸς ὑμᾶς,
\vs{26}ἐπειδὴ ἐπιποθῶν ἦν πάντας ὑμᾶς καὶ ἀδημονῶν, διότι ἠκούσατε ὅτι ἠσθένησεν.
\vs{27}καὶ γὰρ ἠσθένησεν παραπλήσιον θανάτῳ· ἀλλὰ ὁ Θεὸς ἠλέησεν αὐτόν, οὐκ αὐτὸν δὲ μόνον ἀλλὰ καὶ ἐμέ, ἵνα μὴ λύπην ἐπὶ λύπην σχῶ.
\vs{28}Σπουδαιοτέρως οὖν ἔπεμψα αὐτὸν, ἵνα ἰδόντες αὐτὸν πάλιν χαρῆτε κἀγὼ ἀλυπότερος ὦ.
\vs{29}προσδέχεσθε οὖν αὐτὸν ἐν Κυρίῳ μετὰ πάσης χαρᾶς καὶ τοὺς τοιούτους ἐντίμους ἔχετε,
\vs{30}ὅτι διὰ τὸ ἔργον Χριστοῦ μέχρι θανάτου ἤγγισεν παραβολευσάμενος τῇ ψυχῇ, ἵνα ἀναπληρώσῃ τὸ ὑμῶν ὑστέρημα τῆς πρός με λειτουργίας.

\ch{3}
Τὸ λοιπόν, ἀδελφοί μου, χαίρετε ἐν Κυρίῳ. τὰ αὐτὰ γράφειν ὑμῖν ἐμοὶ μὲν οὐκ ὀκνηρόν, ὑμῖν δὲ ἀσφαλές.

\vs{2}Βλέπετε τοὺς κύνας, βλέπετε τοὺς κακοὺς ἐργάτας, βλέπετε τὴν κατατομήν.
\vs{3}ἡμεῖς γάρ ἐσμεν ἡ περιτομή, οἱ Πνεύματι Θεοῦ λατρεύοντες καὶ καυχώμενοι ἐν Χριστῷ Ἰησοῦ καὶ οὐκ ἐν σαρκὶ πεποιθότες,
\vs{4}καίπερ ἐγὼ ἔχων πεποίθησιν καὶ ἐν σαρκί. Εἴ τις δοκεῖ ἄλλος πεποιθέναι ἐν σαρκί, ἐγὼ μᾶλλον·
\vs{5}περιτομῇ ὀκταήμερος, ἐκ γένους Ἰσραήλ, φυλῆς Βενιαμίν, Ἑβραῖος ἐξ Ἑβραίων, κατὰ νόμον Φαρισαῖος,
\vs{6}κατὰ ζῆλος διώκων τὴν ἐκκλησίαν, κατὰ δικαιοσύνην τὴν ἐν νόμῳ γενόμενος ἄμεμπτος.
\vs{7}Ἀλλὰ ἅτινα ἦν μοι κέρδη, ταῦτα ἥγημαι διὰ τὸν Χριστὸν ζημίαν.
\vs{8}ἀλλὰ μενοῦνγε καὶ ἡγοῦμαι πάντα ζημίαν εἶναι διὰ τὸ ὑπερέχον τῆς γνώσεως Χριστοῦ Ἰησοῦ τοῦ Κυρίου μου, δι᾽ ὃν τὰ πάντα ἐζημιώθην, καὶ ἡγοῦμαι σκύβαλα, ἵνα Χριστὸν κερδήσω
\vs{9}καὶ εὑρεθῶ ἐν αὐτῷ, μὴ ἔχων ἐμὴν δικαιοσύνην τὴν ἐκ νόμου ἀλλὰ τὴν διὰ πίστεως Χριστοῦ, τὴν ἐκ Θεοῦ δικαιοσύνην ἐπὶ τῇ πίστει,
\vs{10}τοῦ γνῶναι αὐτὸν καὶ τὴν δύναμιν τῆς ἀναστάσεως αὐτοῦ καὶ τὴν κοινωνίαν τῶν παθημάτων αὐτοῦ, συμμορφιζόμενος τῷ θανάτῳ αὐτοῦ,
\vs{11}εἴ πως καταντήσω εἰς τὴν ἐξανάστασιν τὴν ἐκ νεκρῶν.

\vs{12}Οὐχ ὅτι ἤδη ἔλαβον ἢ ἤδη τετελείωμαι, διώκω δὲ εἰ καὶ καταλάβω, ἐφ᾽ ᾧ καὶ κατελήμφθην ὑπὸ Χριστοῦ Ἰησοῦ.
\vs{13}ἀδελφοί, ἐγὼ ἐμαυτὸν οὐ λογίζομαι κατειληφέναι· ἓν δέ, τὰ μὲν ὀπίσω ἐπιλανθανόμενος τοῖς δὲ ἔμπροσθεν ἐπεκτεινόμενος,
\vs{14}κατὰ σκοπὸν διώκω εἰς τὸ βραβεῖον τῆς ἄνω κλήσεως τοῦ Θεοῦ ἐν Χριστῷ Ἰησοῦ.
\vs{15}Ὅσοι οὖν τέλειοι, τοῦτο φρονῶμεν· καὶ εἴ τι ἑτέρως φρονεῖτε, καὶ τοῦτο ὁ Θεὸς ὑμῖν ἀποκαλύψει·
\vs{16}πλὴν εἰς ὃ ἐφθάσαμεν, τῷ αὐτῷ στοιχεῖν.

\vs{17}Συμμιμηταί μου γίνεσθε, ἀδελφοί, καὶ σκοπεῖτε τοὺς οὕτω περιπατοῦντας καθὼς ἔχετε τύπον ἡμᾶς.
\vs{18}πολλοὶ γὰρ περιπατοῦσιν οὓς πολλάκις ἔλεγον ὑμῖν, νῦν δὲ καὶ κλαίων λέγω, τοὺς ἐχθροὺς τοῦ σταυροῦ τοῦ Χριστοῦ,
\vs{19}ὧν τὸ τέλος ἀπώλεια, ὧν ὁ θεὸς ἡ κοιλία καὶ ἡ δόξα ἐν τῇ αἰσχύνῃ αὐτῶν, οἱ τὰ ἐπίγεια φρονοῦντες.
\vs{20}Ἡμῶν γὰρ τὸ πολίτευμα ἐν οὐρανοῖς ὑπάρχει, ἐξ οὗ καὶ Σωτῆρα ἀπεκδεχόμεθα Κύριον Ἰησοῦν Χριστόν,
\vs{21}ὃς μετασχηματίσει τὸ σῶμα τῆς ταπεινώσεως ἡμῶν σύμμορφον τῷ σώματι τῆς δόξης αὐτοῦ κατὰ τὴν ἐνέργειαν τοῦ δύνασθαι αὐτὸν καὶ ὑποτάξαι αὑτῷ τὰ πάντα.

\ch{4}
Ὥστε, ἀδελφοί μου ἀγαπητοὶ καὶ ἐπιπόθητοι, χαρὰ καὶ στέφανός μου, οὕτως στήκετε ἐν Κυρίῳ, ἀγαπητοί.

\vs{2}Εὐοδίαν παρακαλῶ καὶ Συντύχην παρακαλῶ τὸ αὐτὸ φρονεῖν ἐν Κυρίῳ.
\vs{3}ναὶ ἐρωτῶ καὶ σέ, γνήσιε σύζυγε, συλλαμβάνου αὐταῖς, αἵτινες ἐν τῷ εὐαγγελίῳ συνήθλησάν μοι μετὰ καὶ Κλήμεντος καὶ τῶν λοιπῶν συνεργῶν μου, ὧν τὰ ὀνόματα ἐν βίβλῳ ζωῆς.

\vs{4}Χαίρετε ἐν Κυρίῳ πάντοτε· πάλιν ἐρῶ, χαίρετε.
\vs{5}τὸ ἐπιεικὲς ὑμῶν γνωσθήτω πᾶσιν ἀνθρώποις. ὁ Κύριος ἐγγύς.
\vs{6}Μηδὲν μεριμνᾶτε, ἀλλ᾽ ἐν παντὶ τῇ προσευχῇ καὶ τῇ δεήσει μετὰ εὐχαριστίας τὰ αἰτήματα ὑμῶν γνωριζέσθω πρὸς τὸν Θεόν.
\vs{7}καὶ ἡ εἰρήνη τοῦ Θεοῦ ἡ ὑπερέχουσα πάντα νοῦν φρουρήσει τὰς καρδίας ὑμῶν καὶ τὰ νοήματα ὑμῶν ἐν Χριστῷ Ἰησοῦ.

\vs{8}Τὸ λοιπόν, ἀδελφοί, ὅσα ἐστὶν ἀληθῆ, ὅσα σεμνά, ὅσα δίκαια, ὅσα ἁγνά, ὅσα προσφιλῆ, ὅσα εὔφημα, εἴ τις ἀρετὴ καὶ εἴ τις ἔπαινος, ταῦτα λογίζεσθε·
\vs{9}ἃ καὶ ἐμάθετε καὶ παρελάβετε καὶ ἠκούσατε καὶ εἴδετε ἐν ἐμοί, ταῦτα πράσσετε· καὶ ὁ Θεὸς τῆς εἰρήνης ἔσται μεθ᾽ ὑμῶν.

\vs{10}Ἐχάρην δὲ ἐν Κυρίῳ μεγάλως ὅτι ἤδη ποτὲ ἀνεθάλετε τὸ ὑπὲρ ἐμοῦ φρονεῖν, ἐφ᾽ ᾧ καὶ ἐφρονεῖτε, ἠκαιρεῖσθε δέ.
\vs{11}οὐχ ὅτι καθ᾽ ὑστέρησιν λέγω, ἐγὼ γὰρ ἔμαθον ἐν οἷς εἰμι αὐτάρκης εἶναι.
\vs{12}οἶδα καὶ ταπεινοῦσθαι, οἶδα καὶ περισσεύειν· ἐν παντὶ καὶ ἐν πᾶσιν μεμύημαι, καὶ χορτάζεσθαι καὶ πεινᾶν καὶ περισσεύειν καὶ ὑστερεῖσθαι·
\vs{13}πάντα ἰσχύω ἐν τῷ ἐνδυναμοῦντί με.
\vs{14}Πλὴν καλῶς ἐποιήσατε συνκοινωνήσαντές μου τῇ θλίψει.
\vs{15}οἴδατε δὲ καὶ ὑμεῖς, Φιλιππήσιοι, ὅτι ἐν ἀρχῇ τοῦ εὐαγγελίου, ὅτε ἐξῆλθον ἀπὸ Μακεδονίας, οὐδεμία μοι ἐκκλησία ἐκοινώνησεν εἰς λόγον δόσεως καὶ λήμψεως εἰ μὴ ὑμεῖς μόνοι,
\vs{16}ὅτι καὶ ἐν Θεσσαλονίκῃ καὶ ἅπαξ καὶ δὶς εἰς τὴν χρείαν μοι ἐπέμψατε.
\vs{17}Οὐχ ὅτι ἐπιζητῶ τὸ δόμα, ἀλλὰ ἐπιζητῶ τὸν καρπὸν τὸν πλεονάζοντα εἰς λόγον ὑμῶν.
\vs{18}ἀπέχω δὲ πάντα καὶ περισσεύω· πεπλήρωμαι δεξάμενος παρὰ Ἐπαφροδίτου τὰ παρ᾽ ὑμῶν, ὀσμὴν εὐωδίας, θυσίαν δεκτήν, εὐάρεστον τῷ Θεῷ.
\vs{19}Ὁ δὲ Θεός μου πληρώσει πᾶσαν χρείαν ὑμῶν κατὰ τὸ πλοῦτος αὐτοῦ ἐν δόξῃ ἐν Χριστῷ Ἰησοῦ.
\vs{20}τῷ δὲ Θεῷ καὶ Πατρὶ ἡμῶν ἡ δόξα εἰς τοὺς αἰῶνας τῶν αἰώνων, ἀμήν.

\vs{21}Ἀσπάσασθε πάντα ἅγιον ἐν Χριστῷ Ἰησοῦ. Ἀσπάζονται ὑμᾶς οἱ σὺν ἐμοὶ ἀδελφοί.
\vs{22}Ἀσπάζονται ὑμᾶς πάντες οἱ ἅγιοι, μάλιστα δὲ οἱ ἐκ τῆς Καίσαρος οἰκίας.

\vs{23}Ἡ χάρις τοῦ Κυρίου Ἰησοῦ Χριστοῦ μετὰ τοῦ πνεύματος ὑμῶν.


\def\book{ΠΡΟΣ ΚΟΛΟΣΣΑΕΙΣ}
\biblebook{ΠΡΟΣ ΚΟΛΟΣΣΑΕΙΣ}


\lettrine[lines=2, loversize=0.2, nindent=0em, findent=.25em]{\textcolor{bookheadingcolor}{Π}}{αῦλος} ἀπόστολος Χριστοῦ Ἰησοῦ διὰ θελήματος Θεοῦ καὶ Τιμόθεος ὁ ἀδελφὸς
\vs{2}Τοῖς ἐν Κολοσσαῖς ἁγίοις καὶ πιστοῖς ἀδελφοῖς ἐν Χριστῷ, Χάρις ὑμῖν καὶ εἰρήνη ἀπὸ Θεοῦ Πατρὸς ἡμῶν.
\vs{3}Εὐχαριστοῦμεν τῷ Θεῷ Πατρὶ τοῦ Κυρίου ἡμῶν Ἰησοῦ Χριστοῦ πάντοτε περὶ ὑμῶν προσευχόμενοι,
\vs{4}ἀκούσαντες τὴν πίστιν ὑμῶν ἐν Χριστῷ Ἰησοῦ καὶ τὴν ἀγάπην ἣν ἔχετε εἰς πάντας τοὺς ἁγίους
\vs{5}διὰ τὴν ἐλπίδα τὴν ἀποκειμένην ὑμῖν ἐν τοῖς οὐρανοῖς, ἣν προηκούσατε ἐν τῷ λόγῳ τῆς ἀληθείας τοῦ εὐαγγελίου
\vs{6}τοῦ παρόντος εἰς ὑμᾶς, καθὼς καὶ ἐν παντὶ τῷ κόσμῳ ἐστὶν καρποφορούμενον καὶ αὐξανόμενον καθὼς καὶ ἐν ὑμῖν, ἀφ᾽ ἧς ἡμέρας ἠκούσατε καὶ ἐπέγνωτε τὴν χάριν τοῦ Θεοῦ ἐν ἀληθείᾳ·
\vs{7}καθὼς ἐμάθετε ἀπὸ Ἐπαφρᾶ τοῦ ἀγαπητοῦ συνδούλου ἡμῶν, ὅς ἐστιν πιστὸς ὑπὲρ ἡμῶν διάκονος τοῦ Χριστοῦ,
\vs{8}ὁ καὶ δηλώσας ἡμῖν τὴν ὑμῶν ἀγάπην ἐν Πνεύματι.

\vs{9}Διὰ τοῦτο καὶ ἡμεῖς, ἀφ᾽ ἧς ἡμέρας ἠκούσαμεν, οὐ παυόμεθα ὑπὲρ ὑμῶν προσευχόμενοι καὶ αἰτούμενοι, ἵνα πληρωθῆτε τὴν ἐπίγνωσιν τοῦ θελήματος αὐτοῦ ἐν πάσῃ σοφίᾳ καὶ συνέσει πνευματικῇ,
\vs{10}περιπατῆσαι ἀξίως τοῦ Κυρίου εἰς πᾶσαν ἀρεσκείαν, ἐν παντὶ ἔργῳ ἀγαθῷ καρποφοροῦντες καὶ αὐξανόμενοι τῇ ἐπιγνώσει τοῦ Θεοῦ,
\vs{11}ἐν πάσῃ δυνάμει δυναμούμενοι κατὰ τὸ κράτος τῆς δόξης αὐτοῦ εἰς πᾶσαν ὑπομονὴν καὶ μακροθυμίαν.

μετὰ χαρᾶς
\vs{12}εὐχαριστοῦντες τῷ Πατρὶ τῷ ἱκανώσαντι ὑμᾶς εἰς τὴν μερίδα τοῦ κλήρου τῶν ἁγίων ἐν τῷ φωτί·
\vs{13}ὃς ἐρρύσατο ἡμᾶς ἐκ τῆς ἐξουσίας τοῦ σκότους καὶ μετέστησεν εἰς τὴν βασιλείαν τοῦ Υἱοῦ τῆς ἀγάπης αὐτοῦ,
\vs{14}ἐν ᾧ ἔχομεν τὴν ἀπολύτρωσιν, τὴν ἄφεσιν τῶν ἁμαρτιῶν·
\begin{poetryblock}

\begin{quote} \vs{15}Ὅς ἐστιν εἰκὼν τοῦ Θεοῦ τοῦ ἀοράτου,\end{quote} 

\begin{quote}πρωτότοκος πάσης κτίσεως,\end{quote}

\begin{quote} \vs{16}ὅτι ἐν αὐτῷ ἐκτίσθη τὰ πάντα\end{quote} 

\begin{quote}ἐν τοῖς οὐρανοῖς καὶ ἐπὶ τῆς γῆς,\end{quote} 

\begin{quote}τὰ ὁρατὰ καὶ τὰ ἀόρατα,\end{quote} 

\begin{quote}εἴτε θρόνοι εἴτε κυριότητες\end{quote} 

\begin{quote}εἴτε ἀρχαὶ εἴτε ἐξουσίαι·\end{quote} 

\begin{quote}τὰ πάντα δι᾽ αὐτοῦ καὶ εἰς αὐτὸν ἔκτισται·\end{quote}

\begin{quote} \vs{17}Καὶ αὐτός ἐστιν πρὸ πάντων\end{quote} 

\begin{quote}καὶ τὰ πάντα ἐν αὐτῷ συνέστηκεν,\end{quote}

\begin{quote} \vs{18}καὶ αὐτός ἐστιν ἡ κεφαλὴ τοῦ σώματος τῆς ἐκκλησίας·\end{quote} 

\begin{quote}ὅς ἐστιν ἀρχή,\end{quote} 

\begin{quote}πρωτότοκος ἐκ τῶν νεκρῶν,\end{quote} 

\begin{quote}ἵνα γένηται ἐν πᾶσιν αὐτὸς πρωτεύων,\end{quote}

\begin{quote} \vs{19}ὅτι ἐν αὐτῷ εὐδόκησεν πᾶν τὸ πλήρωμα κατοικῆσαι\end{quote}

\begin{quote} \vs{20}καὶ δι᾽ αὐτοῦ ἀποκαταλλάξαι τὰ πάντα εἰς αὐτόν,\end{quote} 

\begin{quote}εἰρηνοποιήσας διὰ τοῦ αἵματος τοῦ σταυροῦ αὐτοῦ,\end{quote} 

\begin{quote}δι᾽ αὐτοῦ εἴτε τὰ ἐπὶ τῆς γῆς\end{quote} 

\begin{quote}εἴτε τὰ ἐν τοῖς οὐρανοῖς.\end{quote}
\end{poetryblock}

\vs{21}Καὶ ὑμᾶς ποτε ὄντας ἀπηλλοτριωμένους καὶ ἐχθροὺς τῇ διανοίᾳ ἐν τοῖς ἔργοις τοῖς πονηροῖς,
\vs{22}νυνὶ δὲ ἀποκατήλλαξεν ἐν τῷ σώματι τῆς σαρκὸς αὐτοῦ διὰ τοῦ θανάτου παραστῆσαι ὑμᾶς ἁγίους καὶ ἀμώμους καὶ ἀνεγκλήτους κατενώπιον αὐτοῦ,
\vs{23}εἴ γε ἐπιμένετε τῇ πίστει τεθεμελιωμένοι καὶ ἑδραῖοι καὶ μὴ μετακινούμενοι ἀπὸ τῆς ἐλπίδος τοῦ εὐαγγελίου οὗ ἠκούσατε, τοῦ κηρυχθέντος ἐν πάσῃ κτίσει τῇ ὑπὸ τὸν οὐρανόν, οὗ ἐγενόμην ἐγὼ Παῦλος διάκονος.

\vs{24}Νῦν χαίρω ἐν τοῖς παθήμασιν ὑπὲρ ὑμῶν καὶ ἀνταναπληρῶ τὰ ὑστερήματα τῶν θλίψεων τοῦ Χριστοῦ ἐν τῇ σαρκί μου ὑπὲρ τοῦ σώματος αὐτοῦ, ὅ ἐστιν ἡ ἐκκλησία,
\vs{25}ἧς ἐγενόμην ἐγὼ διάκονος κατὰ τὴν οἰκονομίαν τοῦ Θεοῦ τὴν δοθεῖσάν μοι εἰς ὑμᾶς πληρῶσαι τὸν λόγον τοῦ Θεοῦ,
\vs{26}τὸ μυστήριον τὸ ἀποκεκρυμμένον ἀπὸ τῶν αἰώνων καὶ ἀπὸ τῶν γενεῶν— νῦν δὲ ἐφανερώθη τοῖς ἁγίοις αὐτοῦ,
\vs{27}οἷς ἠθέλησεν ὁ Θεὸς γνωρίσαι τί τὸ πλοῦτος τῆς δόξης τοῦ μυστηρίου τούτου ἐν τοῖς ἔθνεσιν, ὅ ἐστιν Χριστὸς ἐν ὑμῖν, ἡ ἐλπὶς τῆς δόξης·
\vs{28}ὃν ἡμεῖς καταγγέλλομεν νουθετοῦντες πάντα ἄνθρωπον καὶ διδάσκοντες πάντα ἄνθρωπον ἐν πάσῃ σοφίᾳ, ἵνα παραστήσωμεν πάντα ἄνθρωπον τέλειον ἐν Χριστῷ·
\vs{29}Εἰς ὃ καὶ κοπιῶ ἀγωνιζόμενος κατὰ τὴν ἐνέργειαν αὐτοῦ τὴν ἐνεργουμένην ἐν ἐμοὶ ἐν δυνάμει.

\ch{2}
Θέλω γὰρ ὑμᾶς εἰδέναι ἡλίκον ἀγῶνα ἔχω ὑπὲρ ὑμῶν καὶ τῶν ἐν Λαοδικείᾳ καὶ ὅσοι οὐχ ἑόρακαν τὸ πρόσωπόν μου ἐν σαρκί,
\vs{2}ἵνα παρακληθῶσιν αἱ καρδίαι αὐτῶν συμβιβασθέντες ἐν ἀγάπῃ καὶ εἰς πᾶν πλοῦτος τῆς πληροφορίας τῆς συνέσεως, εἰς ἐπίγνωσιν τοῦ μυστηρίου τοῦ Θεοῦ, Χριστοῦ,
\vs{3}ἐν ᾧ εἰσιν πάντες οἱ θησαυροὶ τῆς σοφίας καὶ γνώσεως ἀπόκρυφοι.
\vs{4}Τοῦτο λέγω, ἵνα μηδεὶς ὑμᾶς παραλογίζηται ἐν πιθανολογίᾳ.
\vs{5}εἰ γὰρ καὶ τῇ σαρκὶ ἄπειμι, ἀλλὰ τῷ πνεύματι σὺν ὑμῖν εἰμι, χαίρων καὶ βλέπων ὑμῶν τὴν τάξιν καὶ τὸ στερέωμα τῆς εἰς Χριστὸν πίστεως ὑμῶν.

\vs{6}Ὡς οὖν παρελάβετε τὸν Χριστὸν Ἰησοῦν τὸν Κύριον, ἐν αὐτῷ περιπατεῖτε,
\vs{7}ἐρριζωμένοι καὶ ἐποικοδομούμενοι ἐν αὐτῷ καὶ βεβαιούμενοι τῇ πίστει καθὼς ἐδιδάχθητε, περισσεύοντες ἐν εὐχαριστίᾳ.
\vs{8}Βλέπετε μή τις ὑμᾶς ἔσται ὁ συλαγωγῶν διὰ τῆς φιλοσοφίας καὶ κενῆς ἀπάτης κατὰ τὴν παράδοσιν τῶν ἀνθρώπων, κατὰ τὰ στοιχεῖα τοῦ κόσμου καὶ οὐ κατὰ Χριστόν·
\vs{9}ὅτι ἐν αὐτῷ κατοικεῖ πᾶν τὸ πλήρωμα τῆς Θεότητος σωματικῶς,
\vs{10}καὶ ἐστὲ ἐν αὐτῷ πεπληρωμένοι, ὅς ἐστιν ἡ κεφαλὴ πάσης ἀρχῆς καὶ ἐξουσίας.
\vs{11}ἐν ᾧ καὶ περιετμήθητε περιτομῇ ἀχειροποιήτῳ ἐν τῇ ἀπεκδύσει τοῦ σώματος τῆς σαρκός, ἐν τῇ περιτομῇ τοῦ Χριστοῦ,
\vs{12}συνταφέντες αὐτῷ ἐν τῷ βαπτισμῷ, ἐν ᾧ καὶ συνηγέρθητε διὰ τῆς πίστεως τῆς ἐνεργείας τοῦ Θεοῦ τοῦ ἐγείραντος αὐτὸν ἐκ νεκρῶν·
\vs{13}Καὶ ὑμᾶς νεκροὺς ὄντας ἐν τοῖς παραπτώμασιν καὶ τῇ ἀκροβυστίᾳ τῆς σαρκὸς ὑμῶν, συνεζωοποίησεν ὑμᾶς σὺν αὐτῷ, χαρισάμενος ἡμῖν πάντα τὰ παραπτώματα.
\vs{14}ἐξαλείψας τὸ καθ᾽ ἡμῶν χειρόγραφον τοῖς δόγμασιν ὃ ἦν ὑπεναντίον ἡμῖν, καὶ αὐτὸ ἦρκεν ἐκ τοῦ μέσου προσηλώσας αὐτὸ τῷ σταυρῷ·
\vs{15}ἀπεκδυσάμενος τὰς ἀρχὰς καὶ τὰς ἐξουσίας ἐδειγμάτισεν ἐν παρρησίᾳ, θριαμβεύσας αὐτοὺς ἐν αὐτῷ.

\vs{16}Μὴ οὖν τις ὑμᾶς κρινέτω ἐν βρώσει καὶ ἐν πόσει ἢ ἐν μέρει ἑορτῆς ἢ νεομηνίας ἢ σαββάτων·
\vs{17}ἅ ἐστιν σκιὰ τῶν μελλόντων, τὸ δὲ σῶμα τοῦ Χριστοῦ.
\vs{18}μηδεὶς ὑμᾶς καταβραβευέτω θέλων ἐν ταπεινοφροσύνῃ καὶ θρησκείᾳ τῶν ἀγγέλων, ἃ ἑόρακεν ἐμβατεύων, εἰκῇ φυσιούμενος ὑπὸ τοῦ νοὸς τῆς σαρκὸς αὐτοῦ,
\vs{19}καὶ οὐ κρατῶν τὴν Κεφαλήν, ἐξ οὗ πᾶν τὸ σῶμα διὰ τῶν ἁφῶν καὶ συνδέσμων ἐπιχορηγούμενον καὶ συμβιβαζόμενον αὔξει τὴν αὔξησιν τοῦ Θεοῦ.

\vs{20}Εἰ ἀπεθάνετε σὺν Χριστῷ ἀπὸ τῶν στοιχείων τοῦ κόσμου, τί ὡς ζῶντες ἐν κόσμῳ δογματίζεσθε;
\vs{21}Μὴ ἅψῃ μηδὲ γεύσῃ μηδὲ θίγῃς,
\vs{22}ἅ ἐστιν πάντα εἰς φθορὰν τῇ ἀποχρήσει, κατὰ τὰ ἐντάλματα καὶ διδασκαλίας τῶν ἀνθρώπων,
\vs{23}ἅτινά ἐστιν λόγον μὲν ἔχοντα σοφίας ἐν ἐθελοθρησκίᾳ καὶ ταπεινοφροσύνῃ καὶ ἀφειδίᾳ σώματος, οὐκ ἐν τιμῇ τινι πρὸς πλησμονὴν τῆς σαρκός.

\ch{3}
Εἰ οὖν συνηγέρθητε τῷ Χριστῷ, τὰ ἄνω ζητεῖτε, οὗ ὁ Χριστός ἐστιν ἐν δεξιᾷ τοῦ Θεοῦ καθήμενος·
\vs{2}τὰ ἄνω φρονεῖτε, μὴ τὰ ἐπὶ τῆς γῆς.
\vs{3}ἀπεθάνετε γάρ καὶ ἡ ζωὴ ὑμῶν κέκρυπται σὺν τῷ Χριστῷ ἐν τῷ Θεῷ·
\vs{4}ὅταν ὁ Χριστὸς φανερωθῇ, ἡ ζωὴ ὑμῶν, τότε καὶ ὑμεῖς σὺν αὐτῷ φανερωθήσεσθε ἐν δόξῃ.

\vs{5}Νεκρώσατε οὖν τὰ μέλη τὰ ἐπὶ τῆς γῆς, πορνείαν ἀκαθαρσίαν πάθος ἐπιθυμίαν κακήν, καὶ τὴν πλεονεξίαν, ἥτις ἐστὶν εἰδωλολατρία,
\vs{6}δι᾽ ἃ ἔρχεται ἡ ὀργὴ τοῦ Θεοῦ ἐπὶ τοὺς υἱοὺς τῆς ἀπειθείας.
\vs{7}ἐν οἷς καὶ ὑμεῖς περιεπατήσατέ ποτε, ὅτε ἐζῆτε ἐν τούτοις·
\vs{8}νυνὶ δὲ ἀπόθεσθε καὶ ὑμεῖς τὰ πάντα, ὀργήν, θυμόν, κακίαν, βλασφημίαν, αἰσχρολογίαν ἐκ τοῦ στόματος ὑμῶν·
\vs{9}Μὴ ψεύδεσθε εἰς ἀλλήλους, ἀπεκδυσάμενοι τὸν παλαιὸν ἄνθρωπον σὺν ταῖς πράξεσιν αὐτοῦ
\vs{10}καὶ ἐνδυσάμενοι τὸν νέον τὸν ἀνακαινούμενον εἰς ἐπίγνωσιν κατ᾽ εἰκόνα τοῦ κτίσαντος αὐτόν,
\vs{11}ὅπου οὐκ ἔνι Ἕλλην καὶ Ἰουδαῖος, περιτομὴ καὶ ἀκροβυστία, βάρβαρος, Σκύθης, δοῦλος, ἐλεύθερος, ἀλλὰ τὰ πάντα καὶ ἐν πᾶσιν Χριστός.

\vs{12}Ἐνδύσασθε οὖν, ὡς ἐκλεκτοὶ τοῦ Θεοῦ ἅγιοι καὶ ἠγαπημένοι, σπλάγχνα οἰκτιρμοῦ χρηστότητα ταπεινοφροσύνην πραΰτητα μακροθυμίαν,
\vs{13}ἀνεχόμενοι ἀλλήλων καὶ χαριζόμενοι ἑαυτοῖς ἐάν τις πρός τινα ἔχῃ μομφήν· καθὼς καὶ ὁ Κύριος ἐχαρίσατο ὑμῖν, οὕτως καὶ ὑμεῖς·
\vs{14}ἐπὶ πᾶσιν δὲ τούτοις τὴν ἀγάπην, ὅ ἐστιν σύνδεσμος τῆς τελειότητος.
\vs{15}καὶ ἡ εἰρήνη τοῦ Χριστοῦ βραβευέτω ἐν ταῖς καρδίαις ὑμῶν, εἰς ἣν καὶ ἐκλήθητε ἐν ἑνὶ σώματι· καὶ εὐχάριστοι γίνεσθε.
\vs{16}Ὁ λόγος τοῦ Χριστοῦ ἐνοικείτω ἐν ὑμῖν πλουσίως, ἐν πάσῃ σοφίᾳ διδάσκοντες καὶ νουθετοῦντες ἑαυτοὺς, ψαλμοῖς ὕμνοις ᾠδαῖς πνευματικαῖς ἐν τῇ χάριτι ᾄδοντες ἐν ταῖς καρδίαις ὑμῶν τῷ Θεῷ·
\vs{17}καὶ πᾶν ὅ τι ἐὰν ποιῆτε ἐν λόγῳ ἢ ἐν ἔργῳ, πάντα ἐν ὀνόματι Κυρίου Ἰησοῦ, εὐχαριστοῦντες τῷ Θεῷ Πατρὶ δι᾽ αὐτοῦ.

\vs{18}Αἱ γυναῖκες, ὑποτάσσεσθε τοῖς ἀνδράσιν ὡς ἀνῆκεν ἐν Κυρίῳ.
\vs{19}Οἱ ἄνδρες, ἀγαπᾶτε τὰς γυναῖκας καὶ μὴ πικραίνεσθε πρὸς αὐτάς.
\vs{20}Τὰ τέκνα, ὑπακούετε τοῖς γονεῦσιν κατὰ πάντα, τοῦτο γὰρ εὐάρεστόν ἐστιν ἐν Κυρίῳ.
\vs{21}Οἱ πατέρες, μὴ ἐρεθίζετε τὰ τέκνα ὑμῶν, ἵνα μὴ ἀθυμῶσιν.

\vs{22}Οἱ δοῦλοι, ὑπακούετε κατὰ πάντα τοῖς κατὰ σάρκα κυρίοις, μὴ ἐν ὀφθαλμοδουλίαις ὡς ἀνθρωπάρεσκοι, ἀλλ᾽ ἐν ἁπλότητι καρδίας φοβούμενοι τὸν Κύριον.
\vs{23}Ὃ ἐὰν ποιῆτε, ἐκ ψυχῆς ἐργάζεσθε ὡς τῷ Κυρίῳ καὶ οὐκ ἀνθρώποις,
\vs{24}εἰδότες ὅτι ἀπὸ Κυρίου ἀπολήμψεσθε τὴν ἀνταπόδοσιν τῆς κληρονομίας. τῷ Κυρίῳ Χριστῷ δουλεύετε·
\vs{25}ὁ γὰρ ἀδικῶν κομίσεται ὃ ἠδίκησεν, καὶ οὐκ ἔστιν προσωπολημψία.
\inparch{4} \vs{1}Οἱ κύριοι, τὸ δίκαιον καὶ τὴν ἰσότητα τοῖς δούλοις παρέχεσθε, εἰδότες ὅτι καὶ ὑμεῖς ἔχετε Κύριον ἐν οὐρανῷ.

\vs{2}Τῇ προσευχῇ προσκαρτερεῖτε, γρηγοροῦντες ἐν αὐτῇ ἐν εὐχαριστίᾳ,
\vs{3}προσευχόμενοι ἅμα καὶ περὶ ἡμῶν, ἵνα ὁ Θεὸς ἀνοίξῃ ἡμῖν θύραν τοῦ λόγου λαλῆσαι τὸ μυστήριον τοῦ Χριστοῦ, δι᾽ ὃ καὶ δέδεμαι,
\vs{4}ἵνα φανερώσω αὐτὸ ὡς δεῖ με λαλῆσαι.
\vs{5}Ἐν σοφίᾳ περιπατεῖτε πρὸς τοὺς ἔξω τὸν καιρὸν ἐξαγοραζόμενοι.
\vs{6}ὁ λόγος ὑμῶν πάντοτε ἐν χάριτι, ἅλατι ἠρτυμένος, εἰδέναι πῶς δεῖ ὑμᾶς ἑνὶ ἑκάστῳ ἀποκρίνεσθαι.

\vs{7}Τὰ κατ᾽ ἐμὲ πάντα γνωρίσει ὑμῖν Τυχικὸς ὁ ἀγαπητὸς ἀδελφὸς καὶ πιστὸς διάκονος καὶ σύνδουλος ἐν Κυρίῳ,
\vs{8}ὃν ἔπεμψα πρὸς ὑμᾶς εἰς αὐτὸ τοῦτο, ἵνα γνῶτε τὰ περὶ ἡμῶν καὶ παρακαλέσῃ τὰς καρδίας ὑμῶν,
\vs{9}σὺν Ὀνησίμῳ τῷ πιστῷ καὶ ἀγαπητῷ ἀδελφῷ, ὅς ἐστιν ἐξ ὑμῶν· πάντα ὑμῖν γνωρίσουσιν τὰ ὧδε.
\vs{10}Ἀσπάζεται ὑμᾶς Ἀρίσταρχος ὁ συναιχμάλωτός μου καὶ Μᾶρκος ὁ ἀνεψιὸς Βαρνάβα περὶ οὗ ἐλάβετε ἐντολάς, ἐὰν ἔλθῃ πρὸς ὑμᾶς, δέξασθε αὐτόν
\vs{11}καὶ Ἰησοῦς ὁ λεγόμενος Ἰοῦστος, οἱ ὄντες ἐκ περιτομῆς, οὗτοι μόνοι συνεργοὶ εἰς τὴν βασιλείαν τοῦ Θεοῦ, οἵτινες ἐγενήθησάν μοι παρηγορία.
\vs{12}Ἀσπάζεται ὑμᾶς Ἐπαφρᾶς ὁ ἐξ ὑμῶν, δοῦλος Χριστοῦ Ἰησοῦ, πάντοτε ἀγωνιζόμενος ὑπὲρ ὑμῶν ἐν ταῖς προσευχαῖς, ἵνα σταθῆτε τέλειοι καὶ πεπληροφορημένοι ἐν παντὶ θελήματι τοῦ Θεοῦ.
\vs{13}μαρτυρῶ γὰρ αὐτῷ ὅτι ἔχει πολὺν πόνον ὑπὲρ ὑμῶν καὶ τῶν ἐν Λαοδικείᾳ καὶ τῶν ἐν Ἱεραπόλει.
\vs{14}Ἀσπάζεται ὑμᾶς Λουκᾶς ὁ ἰατρὸς ὁ ἀγαπητὸς καὶ Δημᾶς.

\vs{15}Ἀσπάσασθε τοὺς ἐν Λαοδικείᾳ ἀδελφοὺς καὶ Νύμφαν καὶ τὴν κατ᾽ οἶκον αὐτῆς ἐκκλησίαν.
\vs{16}Καὶ ὅταν ἀναγνωσθῇ παρ᾽ ὑμῖν ἡ ἐπιστολή, ποιήσατε ἵνα καὶ ἐν τῇ Λαοδικέων ἐκκλησίᾳ ἀναγνωσθῇ, καὶ τὴν ἐκ Λαοδικείας ἵνα καὶ ὑμεῖς ἀναγνῶτε.
\vs{17}Καὶ εἴπατε Ἀρχίππῳ· Βλέπε τὴν διακονίαν ἣν παρέλαβες ἐν Κυρίῳ, ἵνα αὐτὴν πληροῖς.

\vs{18}Ὁ ἀσπασμὸς τῇ ἐμῇ χειρὶ Παύλου. Μνημονεύετέ μου τῶν δεσμῶν. Ἡ χάρις μεθ᾽ ὑμῶν.


\def\book{ΠΡΟΣ ΘΕΣΣΑΛΟΝΙΚΕΙΣ Α}
\biblebook{ΠΡΟΣ ΘΕΣΣΑΛΟΝΙΚΕΙΣ Α}


\lettrine[lines=2, loversize=0.2, nindent=0em, findent=.25em]{\textcolor{bookheadingcolor}{Π}}{αῦλος} καὶ Σιλουανὸς καὶ Τιμόθεος Τῇ ἐκκλησίᾳ Θεσσαλονικέων ἐν Θεῷ Πατρὶ καὶ Κυρίῳ Ἰησοῦ Χριστῷ, Χάρις ὑμῖν καὶ εἰρήνη.
\vs{2}Εὐχαριστοῦμεν τῷ Θεῷ πάντοτε περὶ πάντων ὑμῶν μνείαν ποιούμενοι ἐπὶ τῶν προσευχῶν ἡμῶν, ἀδιαλείπτως
\vs{3}μνημονεύοντες ὑμῶν τοῦ ἔργου τῆς πίστεως καὶ τοῦ κόπου τῆς ἀγάπης καὶ τῆς ὑπομονῆς τῆς ἐλπίδος τοῦ Κυρίου ἡμῶν Ἰησοῦ Χριστοῦ ἔμπροσθεν τοῦ Θεοῦ καὶ Πατρὸς ἡμῶν,
\vs{4}εἰδότες, ἀδελφοὶ ἠγαπημένοι ὑπὸ τοῦ Θεοῦ, τὴν ἐκλογὴν ὑμῶν,
\vs{5}ὅτι τὸ εὐαγγέλιον ἡμῶν οὐκ ἐγενήθη εἰς ὑμᾶς ἐν λόγῳ μόνον ἀλλὰ καὶ ἐν δυνάμει καὶ ἐν Πνεύματι Ἁγίῳ καὶ ἐν πληροφορίᾳ πολλῇ, καθὼς οἴδατε οἷοι ἐγενήθημεν ἐν ὑμῖν δι᾽ ὑμᾶς.
\vs{6}Καὶ ὑμεῖς μιμηταὶ ἡμῶν ἐγενήθητε καὶ τοῦ Κυρίου, δεξάμενοι τὸν λόγον ἐν θλίψει πολλῇ μετὰ χαρᾶς Πνεύματος Ἁγίου,
\vs{7}ὥστε γενέσθαι ὑμᾶς τύπον πᾶσιν τοῖς πιστεύουσιν ἐν τῇ Μακεδονίᾳ καὶ ἐν τῇ Ἀχαΐᾳ.
\vs{8}ἀφ᾽ ὑμῶν γὰρ ἐξήχηται ὁ λόγος τοῦ Κυρίου οὐ μόνον ἐν τῇ Μακεδονίᾳ καὶ ἐν τῇ Ἀχαΐᾳ, ἀλλ᾽ ἐν παντὶ τόπῳ ἡ πίστις ὑμῶν ἡ πρὸς τὸν Θεὸν ἐξελήλυθεν, ὥστε μὴ χρείαν ἔχειν ἡμᾶς λαλεῖν τι.
\vs{9}αὐτοὶ γὰρ περὶ ἡμῶν ἀπαγγέλλουσιν ὁποίαν εἴσοδον ἔσχομεν πρὸς ὑμᾶς, καὶ πῶς ἐπεστρέψατε πρὸς τὸν Θεὸν ἀπὸ τῶν εἰδώλων δουλεύειν Θεῷ ζῶντι καὶ ἀληθινῷ
\vs{10}καὶ ἀναμένειν τὸν Υἱὸν αὐτοῦ ἐκ τῶν οὐρανῶν, ὃν ἤγειρεν ἐκ τῶν νεκρῶν, Ἰησοῦν τὸν ῥυόμενον ἡμᾶς ἐκ τῆς ὀργῆς τῆς ἐρχομένης.

\ch{2}
Αὐτοὶ γὰρ οἴδατε, ἀδελφοί, τὴν εἴσοδον ἡμῶν τὴν πρὸς ὑμᾶς ὅτι οὐ κενὴ γέγονεν,
\vs{2}ἀλλὰ προπαθόντες καὶ ὑβρισθέντες, καθὼς οἴδατε, ἐν Φιλίπποις ἐπαρρησιασάμεθα ἐν τῷ Θεῷ ἡμῶν λαλῆσαι πρὸς ὑμᾶς τὸ εὐαγγέλιον τοῦ Θεοῦ ἐν πολλῷ ἀγῶνι.
\vs{3}Ἡ γὰρ παράκλησις ἡμῶν οὐκ ἐκ πλάνης οὐδὲ ἐξ ἀκαθαρσίας οὐδὲ ἐν δόλῳ,
\vs{4}ἀλλὰ καθὼς δεδοκιμάσμεθα ὑπὸ τοῦ Θεοῦ πιστευθῆναι τὸ εὐαγγέλιον, οὕτως λαλοῦμεν, οὐχ ὡς ἀνθρώποις ἀρέσκοντες ἀλλὰ Θεῷ τῷ δοκιμάζοντι τὰς καρδίας ἡμῶν.
\vs{5}οὔτε γάρ ποτε ἐν λόγῳ κολακείας ἐγενήθημεν, καθὼς οἴδατε, οὔτε ἐν προφάσει πλεονεξίας, Θεὸς μάρτυς,
\vs{6}οὔτε ζητοῦντες ἐξ ἀνθρώπων δόξαν οὔτε ἀφ᾽ ὑμῶν οὔτε ἀπ᾽ ἄλλων,
\vs{7}δυνάμενοι ἐν βάρει εἶναι ὡς Χριστοῦ ἀπόστολοι. Ἀλλὰ ἐγενήθημεν νήπιοι ἐν μέσῳ ὑμῶν, ὡς ἐὰν τροφὸς θάλπῃ τὰ ἑαυτῆς τέκνα,
\vs{8}οὕτως ὁμειρόμενοι ὑμῶν εὐδοκοῦμεν μεταδοῦναι ὑμῖν οὐ μόνον τὸ εὐαγγέλιον τοῦ Θεοῦ ἀλλὰ καὶ τὰς ἑαυτῶν ψυχάς, διότι ἀγαπητοὶ ἡμῖν ἐγενήθητε.
\vs{9}Μνημονεύετε γάρ, ἀδελφοί, τὸν κόπον ἡμῶν καὶ τὸν μόχθον· νυκτὸς καὶ ἡμέρας ἐργαζόμενοι πρὸς τὸ μὴ ἐπιβαρῆσαί τινα ὑμῶν ἐκηρύξαμεν εἰς ὑμᾶς τὸ εὐαγγέλιον τοῦ Θεοῦ.
\vs{10}ὑμεῖς μάρτυρες καὶ ὁ Θεός, ὡς ὁσίως καὶ δικαίως καὶ ἀμέμπτως ὑμῖν τοῖς πιστεύουσιν ἐγενήθημεν,
\vs{11}καθάπερ οἴδατε, ὡς ἕνα ἕκαστον ὑμῶν ὡς πατὴρ τέκνα ἑαυτοῦ
\vs{12}παρακαλοῦντες ὑμᾶς καὶ παραμυθούμενοι καὶ μαρτυρόμενοι εἰς τὸ περιπατεῖν ὑμᾶς ἀξίως τοῦ Θεοῦ τοῦ καλοῦντος ὑμᾶς εἰς τὴν ἑαυτοῦ βασιλείαν καὶ δόξαν.

\vs{13}Καὶ διὰ τοῦτο καὶ ἡμεῖς εὐχαριστοῦμεν τῷ Θεῷ ἀδιαλείπτως, ὅτι παραλαβόντες λόγον ἀκοῆς παρ᾽ ἡμῶν τοῦ Θεοῦ ἐδέξασθε οὐ λόγον ἀνθρώπων ἀλλὰ καθὼς ἐστὶν ἀληθῶς λόγον Θεοῦ, ὃς καὶ ἐνεργεῖται ἐν ὑμῖν τοῖς πιστεύουσιν.
\vs{14}Ὑμεῖς γὰρ μιμηταὶ ἐγενήθητε, ἀδελφοί, τῶν ἐκκλησιῶν τοῦ Θεοῦ τῶν οὐσῶν ἐν τῇ Ἰουδαίᾳ ἐν Χριστῷ Ἰησοῦ, ὅτι τὰ αὐτὰ ἐπάθετε καὶ ὑμεῖς ὑπὸ τῶν ἰδίων συμφυλετῶν καθὼς καὶ αὐτοὶ ὑπὸ τῶν Ἰουδαίων,
\vs{15}τῶν καὶ τὸν Κύριον ἀποκτεινάντων Ἰησοῦν καὶ τοὺς προφήτας καὶ ἡμᾶς ἐκδιωξάντων καὶ Θεῷ μὴ ἀρεσκόντων καὶ πᾶσιν ἀνθρώποις ἐναντίων,
\vs{16}κωλυόντων ἡμᾶς τοῖς ἔθνεσιν λαλῆσαι ἵνα σωθῶσιν, εἰς τὸ ἀναπληρῶσαι αὐτῶν τὰς ἁμαρτίας πάντοτε. ἔφθασεν δὲ ἐπ᾽ αὐτοὺς ἡ ὀργὴ εἰς τέλος.

\vs{17}Ἡμεῖς δέ, ἀδελφοί, ἀπορφανισθέντες ἀφ᾽ ὑμῶν πρὸς καιρὸν ὥρας, προσώπῳ οὐ καρδίᾳ, περισσοτέρως ἐσπουδάσαμεν τὸ πρόσωπον ὑμῶν ἰδεῖν ἐν πολλῇ ἐπιθυμίᾳ.
\vs{18}διότι ἠθελήσαμεν ἐλθεῖν πρὸς ὑμᾶς, ἐγὼ μὲν Παῦλος καὶ ἅπαξ καὶ δίς, καὶ ἐνέκοψεν ἡμᾶς ὁ Σατανᾶς.
\vs{19}τίς γὰρ ἡμῶν ἐλπὶς ἢ χαρὰ ἢ στέφανος καυχήσεως— ἢ οὐχὶ καὶ ὑμεῖς— ἔμπροσθεν τοῦ Κυρίου ἡμῶν Ἰησοῦ ἐν τῇ αὐτοῦ παρουσίᾳ;
\vs{20}ὑμεῖς γάρ ἐστε ἡ δόξα ἡμῶν καὶ ἡ χαρά.

\ch{3}
Διὸ μηκέτι στέγοντες εὐδοκήσαμεν καταλειφθῆναι ἐν Ἀθήναις μόνοι
\vs{2}καὶ ἐπέμψαμεν Τιμόθεον, τὸν ἀδελφὸν ἡμῶν καὶ συνεργὸν τοῦ Θεοῦ ἐν τῷ εὐαγγελίῳ τοῦ Χριστοῦ, εἰς τὸ στηρίξαι ὑμᾶς καὶ παρακαλέσαι ὑπὲρ τῆς πίστεως ὑμῶν
\vs{3}τὸ μηδένα σαίνεσθαι ἐν ταῖς θλίψεσιν ταύταις. αὐτοὶ γὰρ οἴδατε ὅτι εἰς τοῦτο κείμεθα·
\vs{4}καὶ γὰρ ὅτε πρὸς ὑμᾶς ἦμεν, προελέγομεν ὑμῖν ὅτι μέλλομεν θλίβεσθαι, καθὼς καὶ ἐγένετο καὶ οἴδατε.
\vs{5}διὰ τοῦτο κἀγὼ μηκέτι στέγων ἔπεμψα εἰς τὸ γνῶναι τὴν πίστιν ὑμῶν, μή πως ἐπείρασεν ὑμᾶς ὁ πειράζων καὶ εἰς κενὸν γένηται ὁ κόπος ἡμῶν.

\vs{6}Ἄρτι δὲ ἐλθόντος Τιμοθέου πρὸς ἡμᾶς ἀφ᾽ ὑμῶν καὶ εὐαγγελισαμένου ἡμῖν τὴν πίστιν καὶ τὴν ἀγάπην ὑμῶν καὶ ὅτι ἔχετε μνείαν ἡμῶν ἀγαθὴν πάντοτε, ἐπιποθοῦντες ἡμᾶς ἰδεῖν καθάπερ καὶ ἡμεῖς ὑμᾶς,
\vs{7}διὰ τοῦτο παρεκλήθημεν, ἀδελφοί, ἐφ᾽ ὑμῖν ἐπὶ πάσῃ τῇ ἀνάγκῃ καὶ θλίψει ἡμῶν διὰ τῆς ὑμῶν πίστεως,
\vs{8}ὅτι νῦν ζῶμεν ἐὰν ὑμεῖς στήκετε ἐν Κυρίῳ.
\vs{9}Τίνα γὰρ εὐχαριστίαν δυνάμεθα τῷ Θεῷ ἀνταποδοῦναι περὶ ὑμῶν ἐπὶ πάσῃ τῇ χαρᾷ ᾗ χαίρομεν δι᾽ ὑμᾶς ἔμπροσθεν τοῦ Θεοῦ ἡμῶν,
\vs{10}νυκτὸς καὶ ἡμέρας ὑπερεκπερισσοῦ δεόμενοι εἰς τὸ ἰδεῖν ὑμῶν τὸ πρόσωπον καὶ καταρτίσαι τὰ ὑστερήματα τῆς πίστεως ὑμῶν;

\vs{11}Αὐτὸς δὲ ὁ Θεὸς καὶ Πατὴρ ἡμῶν καὶ ὁ Κύριος ἡμῶν Ἰησοῦς κατευθύναι τὴν ὁδὸν ἡμῶν πρὸς ὑμᾶς·
\vs{12}ὑμᾶς δὲ ὁ Κύριος πλεονάσαι καὶ περισσεύσαι τῇ ἀγάπῃ εἰς ἀλλήλους καὶ εἰς πάντας καθάπερ καὶ ἡμεῖς εἰς ὑμᾶς,
\vs{13}εἰς τὸ στηρίξαι ὑμῶν τὰς καρδίας ἀμέμπτους ἐν ἁγιωσύνῃ ἔμπροσθεν τοῦ Θεοῦ καὶ πατρὸς ἡμῶν ἐν τῇ παρουσίᾳ τοῦ Κυρίου ἡμῶν Ἰησοῦ μετὰ πάντων τῶν ἁγίων αὐτοῦ, ἀμήν.

\ch{4}
Λοιπὸν οὖν, ἀδελφοί, ἐρωτῶμεν ὑμᾶς καὶ παρακαλοῦμεν ἐν Κυρίῳ Ἰησοῦ, ἵνα καθὼς παρελάβετε παρ᾽ ἡμῶν τὸ πῶς δεῖ ὑμᾶς περιπατεῖν καὶ ἀρέσκειν Θεῷ, καθὼς καὶ περιπατεῖτε, ἵνα περισσεύητε μᾶλλον.
\vs{2}οἴδατε γὰρ τίνας παραγγελίας ἐδώκαμεν ὑμῖν διὰ τοῦ Κυρίου Ἰησοῦ.

\vs{3}Τοῦτο γάρ ἐστιν θέλημα τοῦ Θεοῦ, ὁ ἁγιασμὸς ὑμῶν, ἀπέχεσθαι ὑμᾶς ἀπὸ τῆς πορνείας,
\vs{4}εἰδέναι ἕκαστον ὑμῶν τὸ ἑαυτοῦ σκεῦος κτᾶσθαι ἐν ἁγιασμῷ καὶ τιμῇ,
\vs{5}μὴ ἐν πάθει ἐπιθυμίας καθάπερ καὶ τὰ ἔθνη τὰ μὴ εἰδότα τὸν Θεόν,
\vs{6}τὸ μὴ ὑπερβαίνειν καὶ πλεονεκτεῖν ἐν τῷ πράγματι τὸν ἀδελφὸν αὐτοῦ, διότι ἔκδικος Κύριος περὶ πάντων τούτων, καθὼς καὶ προείπαμεν ὑμῖν καὶ διεμαρτυράμεθα.
\vs{7}οὐ γὰρ ἐκάλεσεν ἡμᾶς ὁ Θεὸς ἐπὶ ἀκαθαρσίᾳ ἀλλ᾽ ἐν ἁγιασμῷ.
\vs{8}τοιγαροῦν ὁ ἀθετῶν οὐκ ἄνθρωπον ἀθετεῖ ἀλλὰ τὸν Θεὸν τὸν καὶ διδόντα τὸ Πνεῦμα αὐτοῦ τὸ ἅγιον εἰς ὑμᾶς.

\vs{9}Περὶ δὲ τῆς φιλαδελφίας οὐ χρείαν ἔχετε γράφειν ὑμῖν, αὐτοὶ γὰρ ὑμεῖς θεοδίδακτοί ἐστε εἰς τὸ ἀγαπᾶν ἀλλήλους,
\vs{10}καὶ γὰρ ποιεῖτε αὐτὸ εἰς πάντας τοὺς ἀδελφοὺς τοὺς ἐν ὅλῃ τῇ Μακεδονίᾳ. Παρακαλοῦμεν δὲ ὑμᾶς, ἀδελφοί, περισσεύειν μᾶλλον
\vs{11}καὶ φιλοτιμεῖσθαι ἡσυχάζειν καὶ πράσσειν τὰ ἴδια καὶ ἐργάζεσθαι ταῖς ἰδίαις χερσὶν ὑμῶν, καθὼς ὑμῖν παρηγγείλαμεν,
\vs{12}ἵνα περιπατῆτε εὐσχημόνως πρὸς τοὺς ἔξω καὶ μηδενὸς χρείαν ἔχητε.

\vs{13}Οὐ θέλομεν δὲ ὑμᾶς ἀγνοεῖν, ἀδελφοί, περὶ τῶν κοιμωμένων, ἵνα μὴ λυπῆσθε καθὼς καὶ οἱ λοιποὶ οἱ μὴ ἔχοντες ἐλπίδα.
\vs{14}εἰ γὰρ πιστεύομεν ὅτι Ἰησοῦς ἀπέθανεν καὶ ἀνέστη, οὕτως καὶ ὁ Θεὸς τοὺς κοιμηθέντας διὰ τοῦ Ἰησοῦ ἄξει σὺν αὐτῷ.
\vs{15}Τοῦτο γὰρ ὑμῖν λέγομεν ἐν λόγῳ Κυρίου, ὅτι ἡμεῖς οἱ ζῶντες οἱ περιλειπόμενοι εἰς τὴν παρουσίαν τοῦ Κυρίου οὐ μὴ φθάσωμεν τοὺς κοιμηθέντας·
\vs{16}ὅτι αὐτὸς ὁ Κύριος ἐν κελεύσματι, ἐν φωνῇ ἀρχαγγέλου καὶ ἐν σάλπιγγι Θεοῦ, καταβήσεται ἀπ᾽ οὐρανοῦ καὶ οἱ νεκροὶ ἐν Χριστῷ ἀναστήσονται πρῶτον,
\vs{17}ἔπειτα ἡμεῖς οἱ ζῶντες οἱ περιλειπόμενοι ἅμα σὺν αὐτοῖς ἁρπαγησόμεθα ἐν νεφέλαις εἰς ἀπάντησιν τοῦ Κυρίου εἰς ἀέρα· καὶ οὕτως πάντοτε σὺν Κυρίῳ ἐσόμεθα.
\vs{18}Ὥστε παρακαλεῖτε ἀλλήλους ἐν τοῖς λόγοις τούτοις.

\ch{5}
Περὶ δὲ τῶν χρόνων καὶ τῶν καιρῶν, ἀδελφοί, οὐ χρείαν ἔχετε ὑμῖν γράφεσθαι,
\vs{2}αὐτοὶ γὰρ ἀκριβῶς οἴδατε ὅτι ἡμέρα Κυρίου ὡς κλέπτης ἐν νυκτὶ οὕτως ἔρχεται.
\vs{3}ὅταν λέγωσιν· Εἰρήνη καὶ ἀσφάλεια, τότε αἰφνίδιος αὐτοῖς ἐφίσταται ὄλεθρος ὥσπερ ἡ ὠδὶν τῇ ἐν γαστρὶ ἐχούσῃ, καὶ οὐ μὴ ἐκφύγωσιν.
\vs{4}Ὑμεῖς δέ, ἀδελφοί, οὐκ ἐστὲ ἐν σκότει, ἵνα ἡ ἡμέρα ὑμᾶς ὡς κλέπτης καταλάβῃ·
\vs{5}πάντες γὰρ ὑμεῖς υἱοὶ φωτός ἐστε καὶ υἱοὶ ἡμέρας. Οὐκ ἐσμὲν νυκτὸς οὐδὲ σκότους·
\vs{6}ἄρα οὖν μὴ καθεύδωμεν ὡς οἱ λοιποί ἀλλὰ γρηγορῶμεν καὶ νήφωμεν.
\vs{7}οἱ γὰρ καθεύδοντες νυκτὸς καθεύδουσιν καὶ οἱ μεθυσκόμενοι νυκτὸς μεθύουσιν·
\vs{8}ἡμεῖς δὲ ἡμέρας ὄντες νήφωμεν ἐνδυσάμενοι θώρακα πίστεως καὶ ἀγάπης καὶ περικεφαλαίαν ἐλπίδα σωτηρίας·
\vs{9}ὅτι οὐκ ἔθετο ἡμᾶς ὁ Θεὸς εἰς ὀργὴν ἀλλὰ εἰς περιποίησιν σωτηρίας διὰ τοῦ Κυρίου ἡμῶν Ἰησοῦ Χριστοῦ
\vs{10}τοῦ ἀποθανόντος περὶ ἡμῶν, ἵνα εἴτε γρηγορῶμεν εἴτε καθεύδωμεν ἅμα σὺν αὐτῷ ζήσωμεν.
\vs{11}Διὸ παρακαλεῖτε ἀλλήλους καὶ οἰκοδομεῖτε εἷς τὸν ἕνα, καθὼς καὶ ποιεῖτε.

\vs{12}Ἐρωτῶμεν δὲ ὑμᾶς, ἀδελφοί, εἰδέναι τοὺς κοπιῶντας ἐν ὑμῖν καὶ προϊσταμένους ὑμῶν ἐν Κυρίῳ καὶ νουθετοῦντας ὑμᾶς
\vs{13}καὶ ἡγεῖσθαι αὐτοὺς ὑπερεκπερισσοῦ ἐν ἀγάπῃ διὰ τὸ ἔργον αὐτῶν. εἰρηνεύετε ἐν ἑαυτοῖς.
\vs{14}Παρακαλοῦμεν δὲ ὑμᾶς, ἀδελφοί, νουθετεῖτε τοὺς ἀτάκτους, παραμυθεῖσθε τοὺς ὀλιγοψύχους, ἀντέχεσθε τῶν ἀσθενῶν, μακροθυμεῖτε πρὸς πάντας.
\vs{15}Ὁρᾶτε μή τις κακὸν ἀντὶ κακοῦ τινι ἀποδῷ, ἀλλὰ πάντοτε τὸ ἀγαθὸν διώκετε καὶ εἰς ἀλλήλους καὶ εἰς πάντας.
\begin{poetryblock}

\begin{quote} \vs{16}Πάντοτε χαίρετε,\end{quote}
\end{poetryblock}

\vs{17}ἀδιαλείπτως προσεύχεσθε,

\vs{18}ἐν παντὶ εὐχαριστεῖτε· τοῦτο γὰρ θέλημα Θεοῦ ἐν Χριστῷ Ἰησοῦ εἰς ὑμᾶς.

\vs{19}Τὸ Πνεῦμα μὴ σβέννυτε,

\vs{20}προφητείας μὴ ἐξουθενεῖτε,

\vs{21}πάντα δὲ δοκιμάζετε, τὸ καλὸν κατέχετε,

\vs{22}ἀπὸ παντὸς εἴδους πονηροῦ ἀπέχεσθε.

\vs{23}Αὐτὸς δὲ ὁ Θεὸς τῆς εἰρήνης ἁγιάσαι ὑμᾶς ὁλοτελεῖς, καὶ ὁλόκληρον ὑμῶν τὸ πνεῦμα καὶ ἡ ψυχὴ καὶ τὸ σῶμα ἀμέμπτως ἐν τῇ παρουσίᾳ τοῦ Κυρίου ἡμῶν Ἰησοῦ Χριστοῦ τηρηθείη.
\vs{24}πιστὸς ὁ καλῶν ὑμᾶς, ὃς καὶ ποιήσει.

\vs{25}Ἀδελφοί, προσεύχεσθε καὶ περὶ ἡμῶν.

\vs{26}Ἀσπάσασθε τοὺς ἀδελφοὺς πάντας ἐν φιλήματι ἁγίῳ.
\vs{27}Ἐνορκίζω ὑμᾶς τὸν Κύριον ἀναγνωσθῆναι τὴν ἐπιστολὴν πᾶσιν τοῖς ἀδελφοῖς.

\vs{28}Ἡ χάρις τοῦ Κυρίου ἡμῶν Ἰησοῦ Χριστοῦ μεθ᾽ ὑμῶν.


\def\book{ΠΡΟΣ ΘΕΣΣΑΛΟΝΙΚΕΙΣ Β}
\biblebook{ΠΡΟΣ ΘΕΣΣΑΛΟΝΙΚΕΙΣ Β}


\lettrine[lines=2, loversize=0.2, nindent=0em, findent=.25em]{\textcolor{bookheadingcolor}{Π}}{αῦλος} καὶ Σιλουανὸς καὶ Τιμόθεος Τῇ ἐκκλησίᾳ Θεσσαλονικέων ἐν Θεῷ Πατρὶ ἡμῶν καὶ Κυρίῳ Ἰησοῦ Χριστῷ,
\vs{2}Χάρις ὑμῖν καὶ εἰρήνη ἀπὸ Θεοῦ Πατρὸς ἡμῶν καὶ Κυρίου Ἰησοῦ Χριστοῦ.

\vs{3}Εὐχαριστεῖν ὀφείλομεν τῷ Θεῷ πάντοτε περὶ ὑμῶν, ἀδελφοί, καθὼς ἄξιόν ἐστιν, ὅτι ὑπεραυξάνει ἡ πίστις ὑμῶν καὶ πλεονάζει ἡ ἀγάπη ἑνὸς ἑκάστου πάντων ὑμῶν εἰς ἀλλήλους,
\vs{4}ὥστε αὐτοὺς ἡμᾶς ἐν ὑμῖν ἐνκαυχᾶσθαι ἐν ταῖς ἐκκλησίαις τοῦ Θεοῦ ὑπὲρ τῆς ὑπομονῆς ὑμῶν καὶ πίστεως ἐν πᾶσιν τοῖς διωγμοῖς ὑμῶν καὶ ταῖς θλίψεσιν αἷς ἀνέχεσθε,
\vs{5}Ἔνδειγμα τῆς δικαίας κρίσεως τοῦ Θεοῦ εἰς τὸ καταξιωθῆναι ὑμᾶς τῆς βασιλείας τοῦ Θεοῦ, ὑπὲρ ἧς καὶ πάσχετε,
\vs{6}εἴπερ δίκαιον παρὰ Θεῷ ἀνταποδοῦναι τοῖς θλίβουσιν ὑμᾶς θλῖψιν
\vs{7}καὶ ὑμῖν τοῖς θλιβομένοις ἄνεσιν μεθ᾽ ἡμῶν, ἐν τῇ ἀποκαλύψει τοῦ Κυρίου Ἰησοῦ ἀπ᾽ οὐρανοῦ μετ᾽ ἀγγέλων δυνάμεως αὐτοῦ
\vs{8}ἐν πυρὶ φλογός, διδόντος ἐκδίκησιν τοῖς μὴ εἰδόσιν Θεὸν καὶ τοῖς μὴ ὑπακούουσιν τῷ εὐαγγελίῳ τοῦ Κυρίου ἡμῶν Ἰησοῦ,
\vs{9}οἵτινες δίκην τίσουσιν ὄλεθρον αἰώνιον ἀπὸ προσώπου τοῦ Κυρίου καὶ ἀπὸ τῆς δόξης τῆς ἰσχύος αὐτοῦ,
\vs{10}ὅταν ἔλθῃ ἐνδοξασθῆναι ἐν τοῖς ἁγίοις αὐτοῦ καὶ θαυμασθῆναι ἐν πᾶσιν τοῖς πιστεύσασιν, ὅτι ἐπιστεύθη τὸ μαρτύριον ἡμῶν ἐφ᾽ ὑμᾶς, ἐν τῇ ἡμέρᾳ ἐκείνῃ.
\vs{11}Εἰς ὃ καὶ προσευχόμεθα πάντοτε περὶ ὑμῶν, ἵνα ὑμᾶς ἀξιώσῃ τῆς κλήσεως ὁ Θεὸς ἡμῶν καὶ πληρώσῃ πᾶσαν εὐδοκίαν ἀγαθωσύνης καὶ ἔργον πίστεως ἐν δυνάμει,
\vs{12}ὅπως ἐνδοξασθῇ τὸ ὄνομα τοῦ Κυρίου ἡμῶν Ἰησοῦ ἐν ὑμῖν, καὶ ὑμεῖς ἐν αὐτῷ, κατὰ τὴν χάριν τοῦ Θεοῦ ἡμῶν καὶ Κυρίου Ἰησοῦ Χριστοῦ.

\ch{2}
Ἐρωτῶμεν δὲ ὑμᾶς, ἀδελφοί, ὑπὲρ τῆς παρουσίας τοῦ Κυρίου ἡμῶν Ἰησοῦ Χριστοῦ καὶ ἡμῶν ἐπισυναγωγῆς ἐπ᾽ αὐτόν
\vs{2}εἰς τὸ μὴ ταχέως σαλευθῆναι ὑμᾶς ἀπὸ τοῦ νοὸς μηδὲ θροεῖσθαι, μήτε διὰ πνεύματος μήτε διὰ λόγου μήτε δι᾽ ἐπιστολῆς ὡς δι᾽ ἡμῶν, ὡς ὅτι ἐνέστηκεν ἡ ἡμέρα τοῦ Κυρίου·

\vs{3}μή τις ὑμᾶς ἐξαπατήσῃ κατὰ μηδένα τρόπον. ὅτι ἐὰν μὴ ἔλθῃ ἡ ἀποστασία πρῶτον καὶ ἀποκαλυφθῇ ὁ ἄνθρωπος τῆς ἀνομίας, ὁ υἱὸς τῆς ἀπωλείας,
\vs{4}ὁ ἀντικείμενος καὶ ὑπεραιρόμενος ἐπὶ πάντα λεγόμενον Θεὸν ἢ σέβασμα, ὥστε αὐτὸν εἰς τὸν ναὸν τοῦ Θεοῦ καθίσαι ἀποδεικνύντα ἑαυτὸν ὅτι ἔστιν Θεός.
\vs{5}Οὐ μνημονεύετε ὅτι ἔτι ὢν πρὸς ὑμᾶς ταῦτα ἔλεγον ὑμῖν;
\vs{6}καὶ νῦν τὸ κατέχον οἴδατε εἰς τὸ ἀποκαλυφθῆναι αὐτὸν ἐν τῷ ἑαυτοῦ καιρῷ.
\vs{7}τὸ γὰρ μυστήριον ἤδη ἐνεργεῖται τῆς ἀνομίας· μόνον ὁ κατέχων ἄρτι ἕως ἐκ μέσου γένηται.
\vs{8}καὶ τότε ἀποκαλυφθήσεται ὁ ἄνομος, ὃν ὁ Κύριος Ἰησοῦς ἀνελεῖ τῷ πνεύματι τοῦ στόματος αὐτοῦ καὶ καταργήσει τῇ ἐπιφανείᾳ τῆς παρουσίας αὐτοῦ,
\vs{9}οὗ ἐστιν ἡ παρουσία κατ᾽ ἐνέργειαν τοῦ Σατανᾶ ἐν πάσῃ δυνάμει καὶ σημείοις καὶ τέρασιν ψεύδους
\vs{10}καὶ ἐν πάσῃ ἀπάτῃ ἀδικίας τοῖς ἀπολλυμένοις, ἀνθ᾽ ὧν τὴν ἀγάπην τῆς ἀληθείας οὐκ ἐδέξαντο εἰς τὸ σωθῆναι αὐτούς.
\vs{11}καὶ διὰ τοῦτο πέμπει αὐτοῖς ὁ Θεὸς ἐνέργειαν πλάνης εἰς τὸ πιστεῦσαι αὐτοὺς τῷ ψεύδει,
\vs{12}ἵνα κριθῶσιν πάντες οἱ μὴ πιστεύσαντες τῇ ἀληθείᾳ ἀλλὰ εὐδοκήσαντες τῇ ἀδικίᾳ.

\vs{13}Ἡμεῖς δὲ ὀφείλομεν εὐχαριστεῖν τῷ Θεῷ πάντοτε περὶ ὑμῶν, ἀδελφοὶ ἠγαπημένοι ὑπὸ Κυρίου, ὅτι εἵλατο ὑμᾶς ὁ Θεὸς ἀπαρχὴν εἰς σωτηρίαν ἐν ἁγιασμῷ Πνεύματος καὶ πίστει ἀληθείας,
\vs{14}εἰς ὃ καὶ ἐκάλεσεν ὑμᾶς διὰ τοῦ εὐαγγελίου ἡμῶν εἰς περιποίησιν δόξης τοῦ Κυρίου ἡμῶν Ἰησοῦ Χριστοῦ.

\vs{15}Ἄρα οὖν, ἀδελφοί, στήκετε καὶ κρατεῖτε τὰς παραδόσεις ἃς ἐδιδάχθητε εἴτε διὰ λόγου εἴτε δι᾽ ἐπιστολῆς ἡμῶν.
\vs{16}Αὐτὸς δὲ ὁ Κύριος ἡμῶν Ἰησοῦς Χριστὸς καὶ ὁ Θεὸς ὁ Πατὴρ ἡμῶν ὁ ἀγαπήσας ἡμᾶς καὶ δοὺς παράκλησιν αἰωνίαν καὶ ἐλπίδα ἀγαθὴν ἐν χάριτι,
\vs{17}παρακαλέσαι ὑμῶν τὰς καρδίας καὶ στηρίξαι ἐν παντὶ ἔργῳ καὶ λόγῳ ἀγαθῷ.

\ch{3}
Τὸ λοιπὸν προσεύχεσθε, ἀδελφοί, περὶ ἡμῶν, ἵνα ὁ λόγος τοῦ Κυρίου τρέχῃ καὶ δοξάζηται καθὼς καὶ πρὸς ὑμᾶς,
\vs{2}καὶ ἵνα ῥυσθῶμεν ἀπὸ τῶν ἀτόπων καὶ πονηρῶν ἀνθρώπων· οὐ γὰρ πάντων ἡ πίστις.
\vs{3}Πιστὸς δέ ἐστιν ὁ Κύριος, ὃς στηρίξει ὑμᾶς καὶ φυλάξει ἀπὸ τοῦ πονηροῦ.
\vs{4}πεποίθαμεν δὲ ἐν Κυρίῳ ἐφ᾽ ὑμᾶς, ὅτι ἃ παραγγέλλομεν καὶ ποιεῖτε καὶ ποιήσετε.
\vs{5}Ὁ δὲ Κύριος κατευθύναι ὑμῶν τὰς καρδίας εἰς τὴν ἀγάπην τοῦ Θεοῦ καὶ εἰς τὴν ὑπομονὴν τοῦ Χριστοῦ.

\vs{6}Παραγγέλλομεν δὲ ὑμῖν, ἀδελφοί, ἐν ὀνόματι τοῦ Κυρίου ἡμῶν Ἰησοῦ Χριστοῦ στέλλεσθαι ὑμᾶς ἀπὸ παντὸς ἀδελφοῦ ἀτάκτως περιπατοῦντος καὶ μὴ κατὰ τὴν παράδοσιν ἣν παρελάβοσαν παρ᾽ ἡμῶν.
\vs{7}αὐτοὶ γὰρ οἴδατε πῶς δεῖ μιμεῖσθαι ἡμᾶς, ὅτι οὐκ ἠτακτήσαμεν ἐν ὑμῖν
\vs{8}οὐδὲ δωρεὰν ἄρτον ἐφάγομεν παρά τινος, ἀλλ᾽ ἐν κόπῳ καὶ μόχθῳ νυκτὸς καὶ ἡμέρας ἐργαζόμενοι πρὸς τὸ μὴ ἐπιβαρῆσαί τινα ὑμῶν·
\vs{9}οὐχ ὅτι οὐκ ἔχομεν ἐξουσίαν, ἀλλ᾽ ἵνα ἑαυτοὺς τύπον δῶμεν ὑμῖν εἰς τὸ μιμεῖσθαι ἡμᾶς.
\vs{10}καὶ γὰρ ὅτε ἦμεν πρὸς ὑμᾶς, τοῦτο παρηγγέλλομεν ὑμῖν, ὅτι Εἴ τις οὐ θέλει ἐργάζεσθαι μηδὲ ἐσθιέτω.
\vs{11}Ἀκούομεν γάρ τινας περιπατοῦντας ἐν ὑμῖν ἀτάκτως μηδὲν ἐργαζομένους ἀλλὰ περιεργαζομένους·
\vs{12}τοῖς δὲ τοιούτοις παραγγέλλομεν καὶ παρακαλοῦμεν ἐν Κυρίῳ Ἰησοῦ Χριστῷ, ἵνα μετὰ ἡσυχίας ἐργαζόμενοι τὸν ἑαυτῶν ἄρτον ἐσθίωσιν.
\vs{13}Ὑμεῖς δέ, ἀδελφοί, μὴ ἐγκακήσητε καλοποιοῦντες.

\vs{14}Εἰ δέ τις οὐχ ὑπακούει τῷ λόγῳ ἡμῶν διὰ τῆς ἐπιστολῆς, τοῦτον σημειοῦσθε μὴ συναναμίγνυσθαι αὐτῷ, ἵνα ἐντραπῇ·
\vs{15}καὶ μὴ ὡς ἐχθρὸν ἡγεῖσθε, ἀλλὰ νουθετεῖτε ὡς ἀδελφόν.
\vs{16}Αὐτὸς δὲ ὁ Κύριος τῆς εἰρήνης δῴη ὑμῖν τὴν εἰρήνην διὰ παντὸς ἐν παντὶ τρόπῳ. ὁ Κύριος μετὰ πάντων ὑμῶν.

\vs{17}Ὁ ἀσπασμὸς τῇ ἐμῇ χειρὶ Παύλου, ὅ ἐστιν σημεῖον ἐν πάσῃ ἐπιστολῇ· οὕτως γράφω.
\vs{18}Ἡ χάρις τοῦ Κυρίου ἡμῶν Ἰησοῦ Χριστοῦ μετὰ πάντων ὑμῶν.


\def\book{ΠΡΟΣ ΤΙΜΟΘΕΟΝ Α}
\biblebook{ΠΡΟΣ ΤΙΜΟΘΕΟΝ Α}


\lettrine[lines=2, loversize=0.2, nindent=0em, findent=.25em]{\textcolor{bookheadingcolor}{Π}}{αῦλος} ἀπόστολος Χριστοῦ Ἰησοῦ κατ᾽ ἐπιταγὴν Θεοῦ Σωτῆρος ἡμῶν καὶ Χριστοῦ Ἰησοῦ τῆς ἐλπίδος ἡμῶν
\vs{2}Τιμοθέῳ γνησίῳ τέκνῳ ἐν πίστει, Χάρις ἔλεος εἰρήνη ἀπὸ Θεοῦ πατρὸς καὶ Χριστοῦ Ἰησοῦ τοῦ κυρίου ἡμῶν.

\vs{3}Καθὼς παρεκάλεσά σε προσμεῖναι ἐν Ἐφέσῳ πορευόμενος εἰς Μακεδονίαν, ἵνα παραγγείλῃς τισὶν μὴ ἑτεροδιδασκαλεῖν
\vs{4}μηδὲ προσέχειν μύθοις καὶ γενεαλογίαις ἀπεράντοις, αἵτινες ἐκζητήσεις παρέχουσιν μᾶλλον ἢ οἰκονομίαν Θεοῦ τὴν ἐν πίστει.
\vs{5}Τὸ δὲ τέλος τῆς παραγγελίας ἐστὶν ἀγάπη ἐκ καθαρᾶς καρδίας καὶ συνειδήσεως ἀγαθῆς καὶ πίστεως ἀνυποκρίτου,
\vs{6}ὧν τινες ἀστοχήσαντες ἐξετράπησαν εἰς ματαιολογίαν
\vs{7}θέλοντες εἶναι νομοδιδάσκαλοι, μὴ νοοῦντες μήτε ἃ λέγουσιν μήτε περὶ τίνων διαβεβαιοῦνται.
\vs{8}Οἴδαμεν δὲ ὅτι καλὸς ὁ νόμος, ἐάν τις αὐτῷ νομίμως χρῆται,
\vs{9}εἰδὼς τοῦτο, ὅτι δικαίῳ νόμος οὐ κεῖται, ἀνόμοις δὲ καὶ ἀνυποτάκτοις, ἀσεβέσι καὶ ἁμαρτωλοῖς, ἀνοσίοις καὶ βεβήλοις, πατρολῴαις καὶ μητρολῴαις, ἀνδροφόνοις
\vs{10}πόρνοις ἀρσενοκοίταις ἀνδραποδισταῖς ψεύσταις ἐπιόρκοις, καὶ εἴ τι ἕτερον τῇ ὑγιαινούσῃ διδασκαλίᾳ ἀντίκειται
\vs{11}κατὰ τὸ εὐαγγέλιον τῆς δόξης τοῦ μακαρίου Θεοῦ, ὃ ἐπιστεύθην ἐγώ.

\vs{12}Χάριν ἔχω τῷ ἐνδυναμώσαντί με Χριστῷ Ἰησοῦ τῷ Κυρίῳ ἡμῶν, ὅτι πιστόν με ἡγήσατο θέμενος εἰς διακονίαν
\vs{13}τὸ πρότερον ὄντα βλάσφημον καὶ διώκτην καὶ ὑβριστήν, ἀλλὰ ἠλεήθην, ὅτι ἀγνοῶν ἐποίησα ἐν ἀπιστίᾳ·
\vs{14}ὑπερεπλεόνασεν δὲ ἡ χάρις τοῦ Κυρίου ἡμῶν μετὰ πίστεως καὶ ἀγάπης τῆς ἐν Χριστῷ Ἰησοῦ.
\vs{15}Πιστὸς ὁ λόγος καὶ πάσης ἀποδοχῆς ἄξιος, ὅτι Χριστὸς Ἰησοῦς ἦλθεν εἰς τὸν κόσμον ἁμαρτωλοὺς σῶσαι, ὧν πρῶτός εἰμι ἐγώ.
\vs{16}ἀλλὰ διὰ τοῦτο ἠλεήθην, ἵνα ἐν ἐμοὶ πρώτῳ ἐνδείξηται Χριστὸς Ἰησοῦς τὴν ἅπασαν μακροθυμίαν πρὸς ὑποτύπωσιν τῶν μελλόντων πιστεύειν ἐπ᾽ αὐτῷ εἰς ζωὴν αἰώνιον.
\vs{17}Τῷ δὲ Βασιλεῖ τῶν αἰώνων, ἀφθάρτῳ ἀοράτῳ μόνῳ Θεῷ, τιμὴ καὶ δόξα εἰς τοὺς αἰῶνας τῶν αἰώνων, ἀμήν.

\vs{18}Ταύτην τὴν παραγγελίαν παρατίθεμαί σοι, τέκνον Τιμόθεε, κατὰ τὰς προαγούσας ἐπὶ σὲ προφητείας, ἵνα στρατεύῃ ἐν αὐταῖς τὴν καλὴν στρατείαν
\vs{19}ἔχων πίστιν καὶ ἀγαθὴν συνείδησιν, ἥν τινες ἀπωσάμενοι περὶ τὴν πίστιν ἐναυάγησαν,
\vs{20}ὧν ἐστιν Ὑμέναιος καὶ Ἀλέξανδρος, οὓς παρέδωκα τῷ Σατανᾷ, ἵνα παιδευθῶσιν μὴ βλασφημεῖν.

\ch{2}
Παρακαλῶ οὖν πρῶτον πάντων ποιεῖσθαι δεήσεις προσευχάς ἐντεύξεις εὐχαριστίας ὑπὲρ πάντων ἀνθρώπων,
\vs{2}ὑπὲρ βασιλέων καὶ πάντων τῶν ἐν ὑπεροχῇ ὄντων, ἵνα ἤρεμον καὶ ἡσύχιον βίον διάγωμεν ἐν πάσῃ εὐσεβείᾳ καὶ σεμνότητι.
\vs{3}τοῦτο καλὸν καὶ ἀπόδεκτον ἐνώπιον τοῦ Σωτῆρος ἡμῶν Θεοῦ,
\vs{4}ὃς πάντας ἀνθρώπους θέλει σωθῆναι καὶ εἰς ἐπίγνωσιν ἀληθείας ἐλθεῖν.
\begin{poetryblock}

\begin{quote} \vs{5}Εἷς γὰρ Θεός,\end{quote} 

\begin{quote}εἷς καὶ μεσίτης Θεοῦ καὶ ἀνθρώπων,\end{quote} 

\begin{quote}ἄνθρωπος Χριστὸς Ἰησοῦς,\end{quote}

\begin{quote} \vs{6}ὁ δοὺς ἑαυτὸν ἀντίλυτρον ὑπὲρ πάντων,\end{quote} 

\begin{quote}τὸ μαρτύριον καιροῖς ἰδίοις.\end{quote}
\end{poetryblock}
\vs{7}εἰς ὃ ἐτέθην ἐγὼ κῆρυξ καὶ ἀπόστολος, ἀλήθειαν λέγω οὐ ψεύδομαι, διδάσκαλος ἐθνῶν ἐν πίστει καὶ ἀληθείᾳ.

\vs{8}Βούλομαι οὖν προσεύχεσθαι τοὺς ἄνδρας ἐν παντὶ τόπῳ ἐπαίροντας ὁσίους χεῖρας χωρὶς ὀργῆς καὶ διαλογισμοῦ.
\vs{9}Ὡσαύτως καὶ γυναῖκας ἐν καταστολῇ κοσμίῳ μετὰ αἰδοῦς καὶ σωφροσύνης κοσμεῖν ἑαυτάς, μὴ ἐν πλέγμασιν καὶ χρυσίῳ ἢ μαργαρίταις ἢ ἱματισμῷ πολυτελεῖ,
\vs{10}ἀλλ᾽ ὃ πρέπει γυναιξὶν ἐπαγγελλομέναις θεοσέβειαν, δι᾽ ἔργων ἀγαθῶν.
\vs{11}Γυνὴ ἐν ἡσυχίᾳ μανθανέτω ἐν πάσῃ ὑποταγῇ·
\vs{12}διδάσκειν δὲ γυναικὶ οὐκ ἐπιτρέπω οὐδὲ αὐθεντεῖν ἀνδρός, ἀλλ᾽ εἶναι ἐν ἡσυχίᾳ.
\vs{13}Ἀδὰμ γὰρ πρῶτος ἐπλάσθη, εἶτα Εὕα.
\vs{14}καὶ Ἀδὰμ οὐκ ἠπατήθη, ἡ δὲ γυνὴ ἐξαπατηθεῖσα ἐν παραβάσει γέγονεν·
\vs{15}σωθήσεται δὲ διὰ τῆς τεκνογονίας, ἐὰν μείνωσιν ἐν πίστει καὶ ἀγάπῃ καὶ ἁγιασμῷ μετὰ σωφροσύνης·
\inparch{3} \vs{1}Πιστὸς ὁ λόγος.

εἴ τις ἐπισκοπῆς ὀρέγεται, καλοῦ ἔργου ἐπιθυμεῖ.
\vs{2}δεῖ οὖν τὸν ἐπίσκοπον ἀνεπίλημπτον εἶναι, μιᾶς γυναικὸς ἄνδρα, νηφάλιον σώφρονα κόσμιον φιλόξενον διδακτικόν,
\vs{3}μὴ πάροινον μὴ πλήκτην, ἀλλὰ ἐπιεικῆ ἄμαχον ἀφιλάργυρον,
\vs{4}τοῦ ἰδίου οἴκου καλῶς προϊστάμενον, τέκνα ἔχοντα ἐν ὑποταγῇ, μετὰ πάσης σεμνότητος
\vs{5}εἰ δέ τις τοῦ ἰδίου οἴκου προστῆναι οὐκ οἶδεν, πῶς ἐκκλησίας Θεοῦ ἐπιμελήσεται;,
\vs{6}μὴ νεόφυτον, ἵνα μὴ τυφωθεὶς εἰς κρίμα ἐμπέσῃ τοῦ διαβόλου.
\vs{7}δεῖ δὲ καὶ μαρτυρίαν καλὴν ἔχειν ἀπὸ τῶν ἔξωθεν, ἵνα μὴ εἰς ὀνειδισμὸν ἐμπέσῃ καὶ παγίδα τοῦ διαβόλου.

\vs{8}Διακόνους ὡσαύτως σεμνούς, μὴ διλόγους, μὴ οἴνῳ πολλῷ προσέχοντας, μὴ αἰσχροκερδεῖς,
\vs{9}ἔχοντας τὸ μυστήριον τῆς πίστεως ἐν καθαρᾷ συνειδήσει.
\vs{10}καὶ οὗτοι δὲ δοκιμαζέσθωσαν πρῶτον, εἶτα διακονείτωσαν ἀνέγκλητοι ὄντες.
\vs{11}Γυναῖκας ὡσαύτως σεμνάς, μὴ διαβόλους, νηφαλίους, πιστὰς ἐν πᾶσιν.
\vs{12}Διάκονοι ἔστωσαν μιᾶς γυναικὸς ἄνδρες, τέκνων καλῶς προϊστάμενοι καὶ τῶν ἰδίων οἴκων.
\vs{13}οἱ γὰρ καλῶς διακονήσαντες βαθμὸν ἑαυτοῖς καλὸν περιποιοῦνται καὶ πολλὴν παρρησίαν ἐν πίστει τῇ ἐν Χριστῷ Ἰησοῦ.

\vs{14}Ταῦτά σοι γράφω ἐλπίζων ἐλθεῖν πρὸς σὲ ἐν τάχει·
\vs{15}ἐὰν δὲ βραδύνω, ἵνα εἰδῇς πῶς δεῖ ἐν οἴκῳ Θεοῦ ἀναστρέφεσθαι, ἥτις ἐστὶν ἐκκλησία Θεοῦ ζῶντος, στῦλος καὶ ἑδραίωμα τῆς ἀληθείας.
\vs{16}Καὶ ὁμολογουμένως μέγα ἐστὶν τὸ τῆς εὐσεβείας μυστήριον· 
\begin{poetryblock}

\begin{quote}Ὃς ἐφανερώθη ἐν σαρκί,\end{quote} 

\begin{quote}ἐδικαιώθη ἐν πνεύματι,\end{quote} 

\begin{quote}ὤφθη ἀγγέλοις,\end{quote} 

\begin{quote}ἐκηρύχθη ἐν ἔθνεσιν,\end{quote} 

\begin{quote}ἐπιστεύθη ἐν κόσμῳ,\end{quote} 

\begin{quote}ἀνελήμφθη ἐν δόξῃ.\end{quote}
\end{poetryblock}

\ch{4}
Τὸ δὲ πνεῦμα ῥητῶς λέγει ὅτι ἐν ὑστέροις καιροῖς ἀποστήσονταί τινες τῆς πίστεως προσέχοντες πνεύμασιν πλάνοις καὶ διδασκαλίαις δαιμονίων,
\vs{2}ἐν ὑποκρίσει ψευδολόγων, κεκαυστηριασμένων τὴν ἰδίαν συνείδησιν,
\vs{3}κωλυόντων γαμεῖν, ἀπέχεσθαι βρωμάτων, ἃ ὁ Θεὸς ἔκτισεν εἰς μετάλημψιν μετὰ εὐχαριστίας τοῖς πιστοῖς καὶ ἐπεγνωκόσι τὴν ἀλήθειαν.
\vs{4}ὅτι πᾶν κτίσμα Θεοῦ καλόν καὶ οὐδὲν ἀπόβλητον μετὰ εὐχαριστίας λαμβανόμενον·
\vs{5}ἁγιάζεται γὰρ διὰ λόγου Θεοῦ καὶ ἐντεύξεως.

\vs{6}Ταῦτα ὑποτιθέμενος τοῖς ἀδελφοῖς καλὸς ἔσῃ διάκονος Χριστοῦ Ἰησοῦ, ἐντρεφόμενος τοῖς λόγοις τῆς πίστεως καὶ τῆς καλῆς διδασκαλίας ᾗ παρηκολούθηκας·
\vs{7}Τοὺς δὲ βεβήλους καὶ γραώδεις μύθους παραιτοῦ. γύμναζε δὲ σεαυτὸν πρὸς εὐσέβειαν·
\vs{8}ἡ γὰρ σωματικὴ γυμνασία πρὸς ὀλίγον ἐστὶν ὠφέλιμος, ἡ δὲ εὐσέβεια πρὸς πάντα ὠφέλιμός ἐστιν ἐπαγγελίαν ἔχουσα ζωῆς τῆς νῦν καὶ τῆς μελλούσης.
\vs{9}πιστὸς ὁ λόγος καὶ πάσης ἀποδοχῆς ἄξιος·
\vs{10}Εἰς τοῦτο γὰρ κοπιῶμεν καὶ ἀγωνιζόμεθα, ὅτι ἠλπίκαμεν ἐπὶ Θεῷ ζῶντι, ὅς ἐστιν Σωτὴρ πάντων ἀνθρώπων μάλιστα πιστῶν.

\vs{11}Παράγγελλε ταῦτα καὶ δίδασκε.
\vs{12}Μηδείς σου τῆς νεότητος καταφρονείτω, ἀλλὰ τύπος γίνου τῶν πιστῶν ἐν λόγῳ, ἐν ἀναστροφῇ, ἐν ἀγάπῃ, ἐν πίστει, ἐν ἁγνείᾳ.
\vs{13}ἕως ἔρχομαι πρόσεχε τῇ ἀναγνώσει, τῇ παρακλήσει, τῇ διδασκαλίᾳ.
\vs{14}Μὴ ἀμέλει τοῦ ἐν σοὶ χαρίσματος, ὃ ἐδόθη σοι διὰ προφητείας μετὰ ἐπιθέσεως τῶν χειρῶν τοῦ πρεσβυτερίου.
\vs{15}ταῦτα μελέτα, ἐν τούτοις ἴσθι, ἵνα σου ἡ προκοπὴ φανερὰ ᾖ πᾶσιν.
\vs{16}ἔπεχε σεαυτῷ καὶ τῇ διδασκαλίᾳ, ἐπίμενε αὐτοῖς· τοῦτο γὰρ ποιῶν καὶ σεαυτὸν σώσεις καὶ τοὺς ἀκούοντάς σου.

\ch{5}
Πρεσβυτέρῳ μὴ ἐπιπλήξῃς ἀλλὰ παρακάλει ὡς πατέρα, νεωτέρους ὡς ἀδελφούς,
\vs{2}πρεσβυτέρας ὡς μητέρας, νεωτέρας ὡς ἀδελφὰς ἐν πάσῃ ἁγνείᾳ.

\vs{3}Χήρας τίμα τὰς ὄντως χήρας.
\vs{4}εἰ δέ τις χήρα τέκνα ἢ ἔκγονα ἔχει, μανθανέτωσαν πρῶτον τὸν ἴδιον οἶκον εὐσεβεῖν καὶ ἀμοιβὰς ἀποδιδόναι τοῖς προγόνοις· τοῦτο γάρ ἐστιν ἀπόδεκτον ἐνώπιον τοῦ Θεοῦ.
\vs{5}Ἡ δὲ ὄντως χήρα καὶ μεμονωμένη ἤλπικεν ἐπὶ Θεὸν καὶ προσμένει ταῖς δεήσεσιν καὶ ταῖς προσευχαῖς νυκτὸς καὶ ἡμέρας,
\vs{6}ἡ δὲ σπαταλῶσα ζῶσα τέθνηκεν.
\vs{7}Καὶ ταῦτα παράγγελλε, ἵνα ἀνεπίλημπτοι ὦσιν.
\vs{8}εἰ δέ τις τῶν ἰδίων καὶ μάλιστα οἰκείων οὐ προνοεῖ, τὴν πίστιν ἤρνηται καὶ ἔστιν ἀπίστου χείρων.

\vs{9}Χήρα καταλεγέσθω μὴ ἔλαττον ἐτῶν ἑξήκοντα γεγονυῖα, ἑνὸς ἀνδρὸς γυνή,
\vs{10}ἐν ἔργοις καλοῖς μαρτυρουμένη, εἰ ἐτεκνοτρόφησεν, εἰ ἐξενοδόχησεν, εἰ ἁγίων πόδας ἔνιψεν, εἰ θλιβομένοις ἐπήρκεσεν, εἰ παντὶ ἔργῳ ἀγαθῷ ἐπηκολούθησεν.
\vs{11}Νεωτέρας δὲ χήρας παραιτοῦ· ὅταν γὰρ καταστρηνιάσωσιν τοῦ Χριστοῦ, γαμεῖν θέλουσιν
\vs{12}ἔχουσαι κρίμα ὅτι τὴν πρώτην πίστιν ἠθέτησαν·
\vs{13}ἅμα δὲ καὶ ἀργαὶ μανθάνουσιν περιερχόμεναι τὰς οἰκίας, οὐ μόνον δὲ ἀργαὶ ἀλλὰ καὶ φλύαροι καὶ περίεργοι, λαλοῦσαι τὰ μὴ δέοντα.
\vs{14}Βούλομαι οὖν νεωτέρας γαμεῖν, τεκνογονεῖν, οἰκοδεσποτεῖν, μηδεμίαν ἀφορμὴν διδόναι τῷ ἀντικειμένῳ λοιδορίας χάριν·
\vs{15}ἤδη γάρ τινες ἐξετράπησαν ὀπίσω τοῦ Σατανᾶ.
\vs{16}Εἴ τις πιστὴ ἔχει χήρας, ἐπαρκείτω αὐταῖς καὶ μὴ βαρείσθω ἡ ἐκκλησία, ἵνα ταῖς ὄντως χήραις ἐπαρκέσῃ.

\vs{17}Οἱ καλῶς προεστῶτες πρεσβύτεροι διπλῆς τιμῆς ἀξιούσθωσαν, μάλιστα οἱ κοπιῶντες ἐν λόγῳ καὶ διδασκαλίᾳ.
\vs{18}λέγει γὰρ ἡ γραφή· Βοῦν ἀλοῶντα οὐ φιμώσεις, καί· Ἄξιος ὁ ἐργάτης τοῦ μισθοῦ αὐτοῦ.
\vs{19}Κατὰ πρεσβυτέρου κατηγορίαν μὴ παραδέχου, ἐκτὸς εἰ μὴ ἐπὶ δύο ἢ τριῶν μαρτύρων.
\vs{20}Τοὺς ἁμαρτάνοντας ἐνώπιον πάντων ἔλεγχε, ἵνα καὶ οἱ λοιποὶ φόβον ἔχωσιν.
\vs{21}Διαμαρτύρομαι ἐνώπιον τοῦ Θεοῦ καὶ Χριστοῦ Ἰησοῦ καὶ τῶν ἐκλεκτῶν ἀγγέλων, ἵνα ταῦτα φυλάξῃς χωρὶς προκρίματος, μηδὲν ποιῶν κατὰ πρόσκλισιν.
\vs{22}Χεῖρας ταχέως μηδενὶ ἐπιτίθει μηδὲ κοινώνει ἁμαρτίαις ἀλλοτρίαις· σεαυτὸν ἁγνὸν τήρει.

\vs{23}Μηκέτι ὑδροπότει, ἀλλὰ οἴνῳ ὀλίγῳ χρῶ διὰ τὸν στόμαχον καὶ τὰς πυκνάς σου ἀσθενείας.
\vs{24}Τινῶν ἀνθρώπων αἱ ἁμαρτίαι πρόδηλοί εἰσιν προάγουσαι εἰς κρίσιν, τισὶν δὲ καὶ ἐπακολουθοῦσιν·
\vs{25}ὡσαύτως καὶ τὰ ἔργα τὰ καλὰ πρόδηλα, καὶ τὰ ἄλλως ἔχοντα κρυβῆναι οὐ δύνανται.

\ch{6}
Ὅσοι εἰσὶν ὑπὸ ζυγὸν δοῦλοι, τοὺς ἰδίους δεσπότας πάσης τιμῆς ἀξίους ἡγείσθωσαν, ἵνα μὴ τὸ ὄνομα τοῦ Θεοῦ καὶ ἡ διδασκαλία βλασφημῆται.
\vs{2}οἱ δὲ πιστοὺς ἔχοντες δεσπότας μὴ καταφρονείτωσαν, ὅτι ἀδελφοί εἰσιν, ἀλλὰ μᾶλλον δουλευέτωσαν, ὅτι πιστοί εἰσιν καὶ ἀγαπητοὶ οἱ τῆς εὐεργεσίας ἀντιλαμβανόμενοι.

Ταῦτα δίδασκε καὶ παρακάλει.
\vs{3}Εἴ τις ἑτεροδιδασκαλεῖ καὶ μὴ προσέρχεται ὑγιαίνουσιν λόγοις τοῖς τοῦ Κυρίου ἡμῶν Ἰησοῦ Χριστοῦ καὶ τῇ κατ᾽ εὐσέβειαν διδασκαλίᾳ,
\vs{4}τετύφωται, μηδὲν ἐπιστάμενος, ἀλλὰ νοσῶν περὶ ζητήσεις καὶ λογομαχίας, ἐξ ὧν γίνεται φθόνος ἔρις βλασφημίαι, ὑπόνοιαι πονηραί,
\vs{5}διαπαρατριβαὶ διεφθαρμένων ἀνθρώπων τὸν νοῦν καὶ ἀπεστερημένων τῆς ἀληθείας, νομιζόντων πορισμὸν εἶναι τὴν εὐσέβειαν.

\vs{6}Ἔστιν δὲ πορισμὸς μέγας ἡ εὐσέβεια μετὰ αὐταρκείας·
\begin{poetryblock}

\begin{quote} \vs{7}οὐδὲν γὰρ εἰσηνέγκαμεν εἰς τὸν κόσμον,\end{quote} 

\begin{quote}ὅτι οὐδὲ ἐξενεγκεῖν τι δυνάμεθα·\end{quote}

\begin{quote} \vs{8}ἔχοντες δὲ διατροφὰς καὶ σκεπάσματα,\end{quote} 

\begin{quote}τούτοις ἀρκεσθησόμεθα.\end{quote}
\end{poetryblock}

\vs{9}Οἱ δὲ βουλόμενοι πλουτεῖν ἐμπίπτουσιν εἰς πειρασμὸν καὶ παγίδα καὶ ἐπιθυμίας πολλὰς ἀνοήτους καὶ βλαβεράς, αἵτινες βυθίζουσιν τοὺς ἀνθρώπους εἰς ὄλεθρον καὶ ἀπώλειαν.
\vs{10}ῥίζα γὰρ πάντων τῶν κακῶν ἐστιν ἡ φιλαργυρία, ἧς τινες ὀρεγόμενοι ἀπεπλανήθησαν ἀπὸ τῆς πίστεως καὶ ἑαυτοὺς περιέπειραν ὀδύναις πολλαῖς.

\vs{11}Σὺ δέ, ὦ ἄνθρωπε Θεοῦ, ταῦτα φεῦγε· 
\begin{poetryblock}

\begin{quote}δίωκε δὲ δικαιοσύνην εὐσέβειαν πίστιν,\end{quote} 

\begin{quote}ἀγάπην ὑπομονήν πραϋπαθίαν.\end{quote}

\begin{quote} \vs{12}ἀγωνίζου τὸν καλὸν ἀγῶνα τῆς πίστεως,\end{quote} 

\begin{quote}ἐπιλαβοῦ τῆς αἰωνίου ζωῆς, εἰς ἣν ἐκλήθης\end{quote} 

\begin{quote}καὶ ὡμολόγησας τὴν καλὴν ὁμολογίαν\end{quote} 

\begin{quote}ἐνώπιον πολλῶν μαρτύρων.\end{quote}
\end{poetryblock}

\vs{13}Παραγγέλλω σοι ἐνώπιον τοῦ Θεοῦ τοῦ ζωογονοῦντος τὰ πάντα καὶ Χριστοῦ Ἰησοῦ τοῦ μαρτυρήσαντος ἐπὶ Ποντίου Πιλάτου τὴν καλὴν ὁμολογίαν,
\vs{14}τηρῆσαί σε τὴν ἐντολὴν ἄσπιλον ἀνεπίλημπτον μέχρι τῆς ἐπιφανείας τοῦ Κυρίου ἡμῶν Ἰησοῦ Χριστοῦ,
\vs{15}ἣν καιροῖς ἰδίοις δείξει

ὁ μακάριος καὶ μόνος Δυνάστης, 
\begin{poetryblock}

\begin{quote}ὁ Βασιλεὺς τῶν βασιλευόντων\end{quote} 

\begin{quote}καὶ Κύριος τῶν κυριευόντων,\end{quote}

\begin{quote} \vs{16}ὁ μόνος ἔχων ἀθανασίαν,\end{quote} 

\begin{quote}φῶς οἰκῶν ἀπρόσιτον,\end{quote} 

\begin{quote}ὃν εἶδεν οὐδεὶς ἀνθρώπων οὐδὲ ἰδεῖν δύναται·\end{quote} 

\begin{quote}ᾧ τιμὴ καὶ κράτος αἰώνιον, ἀμήν.\end{quote}
\end{poetryblock}

\vs{17}Τοῖς πλουσίοις ἐν τῷ νῦν αἰῶνι παράγγελλε μὴ ὑψηλοφρονεῖν μηδὲ ἠλπικέναι ἐπὶ πλούτου ἀδηλότητι ἀλλ᾽ ἐπὶ Θεῷ τῷ παρέχοντι ἡμῖν πάντα πλουσίως εἰς ἀπόλαυσιν,
\vs{18}ἀγαθοεργεῖν, πλουτεῖν ἐν ἔργοις καλοῖς, εὐμεταδότους εἶναι, κοινωνικούς,
\vs{19}ἀποθησαυρίζοντας ἑαυτοῖς θεμέλιον καλὸν εἰς τὸ μέλλον, ἵνα ἐπιλάβωνται τῆς ὄντως ζωῆς.

\vs{20}Ὦ Τιμόθεε, τὴν παραθήκην φύλαξον ἐκτρεπόμενος τὰς βεβήλους κενοφωνίας καὶ ἀντιθέσεις τῆς ψευδωνύμου γνώσεως,
\vs{21}ἥν τινες ἐπαγγελλόμενοι περὶ τὴν πίστιν ἠστόχησαν.

Ἡ χάρις μεθ᾽ ὑμῶν.


\def\book{ΠΡΟΣ ΤΙΜΟΘΕΟΝ Β}
\biblebook{ΠΡΟΣ ΤΙΜΟΘΕΟΝ Β}


\lettrine[lines=2, loversize=0.2, nindent=0em, findent=.25em]{\textcolor{bookheadingcolor}{Π}}{αῦλος} ἀπόστολος Χριστοῦ Ἰησοῦ διὰ θελήματος Θεοῦ κατ᾽ ἐπαγγελίαν ζωῆς τῆς ἐν Χριστῷ Ἰησοῦ
\vs{2}Τιμοθέῳ ἀγαπητῷ τέκνῳ, Χάρις ἔλεος εἰρήνη ἀπὸ Θεοῦ Πατρὸς καὶ Χριστοῦ Ἰησοῦ τοῦ Κυρίου ἡμῶν.

\vs{3}Χάριν ἔχω τῷ Θεῷ, ᾧ λατρεύω ἀπὸ προγόνων ἐν καθαρᾷ συνειδήσει, ὡς ἀδιάλειπτον ἔχω τὴν περὶ σοῦ μνείαν ἐν ταῖς δεήσεσίν μου νυκτὸς καὶ ἡμέρας,
\vs{4}ἐπιποθῶν σε ἰδεῖν, μεμνημένος σου τῶν δακρύων, ἵνα χαρᾶς πληρωθῶ,
\vs{5}ὑπόμνησιν λαβὼν τῆς ἐν σοὶ ἀνυποκρίτου πίστεως, ἥτις ἐνῴκησεν πρῶτον ἐν τῇ μάμμῃ σου Λωΐδι καὶ τῇ μητρί σου Εὐνίκῃ, πέπεισμαι δὲ ὅτι καὶ ἐν σοί.

\vs{6}Δι᾽ ἣν αἰτίαν ἀναμιμνῄσκω σε ἀναζωπυρεῖν τὸ χάρισμα τοῦ Θεοῦ, ὅ ἐστιν ἐν σοὶ διὰ τῆς ἐπιθέσεως τῶν χειρῶν μου.
\vs{7}οὐ γὰρ ἔδωκεν ἡμῖν ὁ Θεὸς πνεῦμα δειλίας ἀλλὰ δυνάμεως καὶ ἀγάπης καὶ σωφρονισμοῦ.
\vs{8}Μὴ οὖν ἐπαισχυνθῇς τὸ μαρτύριον τοῦ Κυρίου ἡμῶν μηδὲ ἐμὲ τὸν δέσμιον αὐτοῦ, ἀλλὰ συνκακοπάθησον τῷ εὐαγγελίῳ κατὰ δύναμιν Θεοῦ,

\vs{9}τοῦ σώσαντος ἡμᾶς 
\begin{poetryblock}

\begin{quote}καὶ καλέσαντος κλήσει ἁγίᾳ,\end{quote} 

\begin{quote}οὐ κατὰ τὰ ἔργα ἡμῶν\end{quote} 

\begin{quote}ἀλλὰ κατὰ ἰδίαν πρόθεσιν καὶ χάριν,\end{quote} 

\begin{quote}τὴν δοθεῖσαν ἡμῖν ἐν Χριστῷ Ἰησοῦ\end{quote} 

\begin{quote}πρὸ χρόνων αἰωνίων,\end{quote}

\begin{quote} \vs{10}φανερωθεῖσαν δὲ νῦν\end{quote} 

\begin{quote}διὰ τῆς ἐπιφανείας τοῦ Σωτῆρος ἡμῶν Χριστοῦ Ἰησοῦ,\end{quote} 

\begin{quote}καταργήσαντος μὲν τὸν θάνατον\end{quote} 

\begin{quote}φωτίσαντος δὲ ζωὴν καὶ ἀφθαρσίαν διὰ τοῦ εὐαγγελίου\end{quote}
\end{poetryblock}

\vs{11}εἰς ὃ ἐτέθην ἐγὼ κήρυξ καὶ ἀπόστολος καὶ διδάσκαλος,
\vs{12}Δι᾽ ἣν αἰτίαν καὶ ταῦτα πάσχω· ἀλλ᾽ οὐκ ἐπαισχύνομαι, οἶδα γὰρ ᾧ πεπίστευκα καὶ πέπεισμαι ὅτι δυνατός ἐστιν τὴν παραθήκην μου φυλάξαι εἰς ἐκείνην τὴν ἡμέραν.
\vs{13}Ὑποτύπωσιν ἔχε ὑγιαινόντων λόγων ὧν παρ᾽ ἐμοῦ ἤκουσας ἐν πίστει καὶ ἀγάπῃ τῇ ἐν Χριστῷ Ἰησοῦ·
\vs{14}τὴν καλὴν παραθήκην φύλαξον διὰ Πνεύματος Ἁγίου τοῦ ἐνοικοῦντος ἐν ἡμῖν.

\vs{15}Οἶδας τοῦτο, ὅτι ἀπεστράφησάν με πάντες οἱ ἐν τῇ Ἀσίᾳ, ὧν ἐστιν Φύγελος καὶ Ἑρμογένης.
\vs{16}Δῴη ἔλεος ὁ Κύριος τῷ Ὀνησιφόρου οἴκῳ, ὅτι πολλάκις με ἀνέψυξεν καὶ τὴν ἅλυσίν μου οὐκ ἐπαισχύνθη,
\vs{17}ἀλλὰ γενόμενος ἐν Ῥώμῃ σπουδαίως ἐζήτησέν με καὶ εὗρεν·
\vs{18}Δῴη αὐτῷ ὁ Κύριος εὑρεῖν ἔλεος παρὰ Κυρίου ἐν ἐκείνῃ τῇ ἡμέρᾳ. καὶ ὅσα ἐν Ἐφέσῳ διηκόνησεν, βέλτιον σὺ γινώσκεις.

\ch{2}
Σὺ οὖν, τέκνον μου, ἐνδυναμοῦ ἐν τῇ χάριτι τῇ ἐν Χριστῷ Ἰησοῦ,
\vs{2}καὶ ἃ ἤκουσας παρ᾽ ἐμοῦ διὰ πολλῶν μαρτύρων, ταῦτα παράθου πιστοῖς ἀνθρώποις, οἵτινες ἱκανοὶ ἔσονται καὶ ἑτέρους διδάξαι.
\vs{3}Συνκακοπάθησον ὡς καλὸς στρατιώτης Χριστοῦ Ἰησοῦ.
\vs{4}οὐδεὶς στρατευόμενος ἐμπλέκεται ταῖς τοῦ βίου πραγματείαις, ἵνα τῷ στρατολογήσαντι ἀρέσῃ.
\vs{5}ἐὰν δὲ καὶ ἀθλῇ τις, οὐ στεφανοῦται ἐὰν μὴ νομίμως ἀθλήσῃ.
\vs{6}τὸν κοπιῶντα γεωργὸν δεῖ πρῶτον τῶν καρπῶν μεταλαμβάνειν.
\vs{7}νόει ὃ λέγω· δώσει γάρ σοι ὁ Κύριος σύνεσιν ἐν πᾶσιν.

\vs{8}Μνημόνευε Ἰησοῦν Χριστὸν ἐγηγερμένον ἐκ νεκρῶν, ἐκ σπέρματος Δαυίδ, κατὰ τὸ εὐαγγέλιόν μου,
\vs{9}ἐν ᾧ κακοπαθῶ μέχρι δεσμῶν ὡς κακοῦργος, ἀλλὰ ὁ λόγος τοῦ Θεοῦ οὐ δέδεται·
\vs{10}διὰ τοῦτο πάντα ὑπομένω διὰ τοὺς ἐκλεκτούς, ἵνα καὶ αὐτοὶ σωτηρίας τύχωσιν τῆς ἐν Χριστῷ Ἰησοῦ μετὰ δόξης αἰωνίου.
\vs{11}Πιστὸς ὁ λόγος· 
\begin{poetryblock}

\begin{quote}Εἰ γὰρ συναπεθάνομεν, καὶ συζήσομεν·\end{quote}

\begin{quote} \vs{12}εἰ ὑπομένομεν, καὶ συμβασιλεύσομεν·\end{quote} 

\begin{quote}εἰ ἀρνησόμεθα, κἀκεῖνος ἀρνήσεται ἡμᾶς·\end{quote}

\begin{quote} \vs{13}εἰ ἀπιστοῦμεν, ἐκεῖνος πιστὸς μένει,\end{quote} 

\begin{quote}ἀρνήσασθαι γὰρ ἑαυτὸν οὐ δύναται.\end{quote}
\end{poetryblock}

\vs{14}Ταῦτα ὑπομίμνῃσκε διαμαρτυρόμενος ἐνώπιον τοῦ Θεοῦ μὴ λογομαχεῖν, ἐπ᾽ οὐδὲν χρήσιμον, ἐπὶ καταστροφῇ τῶν ἀκουόντων.
\vs{15}Σπούδασον σεαυτὸν δόκιμον παραστῆσαι τῷ Θεῷ, ἐργάτην ἀνεπαίσχυντον, ὀρθοτομοῦντα τὸν λόγον τῆς ἀληθείας.
\vs{16}Τὰς δὲ βεβήλους κενοφωνίας περιΐστασο· ἐπὶ πλεῖον γὰρ προκόψουσιν ἀσεβείας
\vs{17}καὶ ὁ λόγος αὐτῶν ὡς γάγγραινα νομὴν ἕξει. ὧν ἐστιν Ὑμέναιος καὶ Φιλητός,
\vs{18}οἵτινες περὶ τὴν ἀλήθειαν ἠστόχησαν, λέγοντες τὴν ἀνάστασιν ἤδη γεγονέναι, καὶ ἀνατρέπουσιν τήν τινων πίστιν.
\vs{19}Ὁ μέντοι στερεὸς θεμέλιος τοῦ Θεοῦ ἕστηκεν, ἔχων τὴν σφραγῖδα ταύτην· Ἔγνω Κύριος τοὺς ὄντας αὐτοῦ, καί· Ἀποστήτω ἀπὸ ἀδικίας πᾶς ὁ ὀνομάζων τὸ ὄνομα Κυρίου.
\vs{20}Ἐν μεγάλῃ δὲ οἰκίᾳ οὐκ ἔστιν μόνον σκεύη χρυσᾶ καὶ ἀργυρᾶ ἀλλὰ καὶ ξύλινα καὶ ὀστράκινα, καὶ ἃ μὲν εἰς τιμὴν ἃ δὲ εἰς ἀτιμίαν·
\vs{21}ἐὰν οὖν τις ἐκκαθάρῃ ἑαυτὸν ἀπὸ τούτων, ἔσται σκεῦος εἰς τιμήν, ἡγιασμένον, εὔχρηστον τῷ δεσπότῃ, εἰς πᾶν ἔργον ἀγαθὸν ἡτοιμασμένον.

\vs{22}Τὰς δὲ νεωτερικὰς ἐπιθυμίας φεῦγε, δίωκε δὲ δικαιοσύνην πίστιν ἀγάπην εἰρήνην μετὰ τῶν ἐπικαλουμένων τὸν Κύριον ἐκ καθαρᾶς καρδίας.
\vs{23}Τὰς δὲ μωρὰς καὶ ἀπαιδεύτους ζητήσεις παραιτοῦ, εἰδὼς ὅτι γεννῶσιν μάχας·
\vs{24}δοῦλον δὲ Κυρίου οὐ δεῖ μάχεσθαι ἀλλὰ ἤπιον εἶναι πρὸς πάντας, διδακτικόν, ἀνεξίκακον,
\vs{25}ἐν πραΰτητι παιδεύοντα τοὺς ἀντιδιατιθεμένους, μήποτε δώῃ αὐτοῖς ὁ Θεὸς μετάνοιαν εἰς ἐπίγνωσιν ἀληθείας
\vs{26}καὶ ἀνανήψωσιν ἐκ τῆς τοῦ διαβόλου παγίδος, ἐζωγρημένοι ὑπ᾽ αὐτοῦ εἰς τὸ ἐκείνου θέλημα.

\ch{3}
Τοῦτο δὲ γίνωσκε, ὅτι ἐν ἐσχάταις ἡμέραις ἐνστήσονται καιροὶ χαλεποί·
\vs{2}ἔσονται γὰρ οἱ ἄνθρωποι φίλαυτοι φιλάργυροι ἀλαζόνες ὑπερήφανοι βλάσφημοι, γονεῦσιν ἀπειθεῖς, ἀχάριστοι ἀνόσιοι
\vs{3}ἄστοργοι ἄσπονδοι διάβολοι ἀκρατεῖς ἀνήμεροι ἀφιλάγαθοι
\vs{4}προδόται προπετεῖς τετυφωμένοι, φιλήδονοι μᾶλλον ἢ φιλόθεοι,
\vs{5}ἔχοντες μόρφωσιν εὐσεβείας τὴν δὲ δύναμιν αὐτῆς ἠρνημένοι· καὶ τούτους ἀποτρέπου.
\vs{6}Ἐκ τούτων γάρ εἰσιν οἱ ἐνδύνοντες εἰς τὰς οἰκίας καὶ αἰχμαλωτίζοντες γυναικάρια σεσωρευμένα ἁμαρτίαις, ἀγόμενα ἐπιθυμίαις ποικίλαις,
\vs{7}πάντοτε μανθάνοντα καὶ μηδέποτε εἰς ἐπίγνωσιν ἀληθείας ἐλθεῖν δυνάμενα.
\vs{8}ὃν τρόπον δὲ Ἰάννης καὶ Ἰαμβρῆς ἀντέστησαν Μωϋσεῖ, οὕτως καὶ οὗτοι ἀνθίστανται τῇ ἀληθείᾳ, ἄνθρωποι κατεφθαρμένοι τὸν νοῦν, ἀδόκιμοι περὶ τὴν πίστιν.
\vs{9}ἀλλ᾽ οὐ προκόψουσιν ἐπὶ πλεῖον· ἡ γὰρ ἄνοια αὐτῶν ἔκδηλος ἔσται πᾶσιν, ὡς καὶ ἡ ἐκείνων ἐγένετο.

\vs{10}Σὺ δὲ παρηκολούθησάς μου τῇ διδασκαλίᾳ, τῇ ἀγωγῇ, τῇ προθέσει, τῇ πίστει, τῇ μακροθυμίᾳ, τῇ ἀγάπῃ, τῇ ὑπομονῇ,
\vs{11}τοῖς διωγμοῖς, τοῖς παθήμασιν, οἷά μοι ἐγένετο ἐν Ἀντιοχείᾳ, ἐν Ἰκονίῳ, ἐν Λύστροις, οἵους διωγμοὺς ὑπήνεγκα καὶ ἐκ πάντων με ἐρρύσατο ὁ Κύριος.
\vs{12}καὶ πάντες δὲ οἱ θέλοντες εὐσεβῶς ζῆν ἐν Χριστῷ Ἰησοῦ διωχθήσονται.
\vs{13}πονηροὶ δὲ ἄνθρωποι καὶ γόητες προκόψουσιν ἐπὶ τὸ χεῖρον πλανῶντες καὶ πλανώμενοι.
\vs{14}Σὺ δὲ μένε ἐν οἷς ἔμαθες καὶ ἐπιστώθης, εἰδὼς παρὰ τίνων ἔμαθες,
\vs{15}καὶ ὅτι ἀπὸ βρέφους τὰ ἱερὰ γράμματα οἶδας, τὰ δυνάμενά σε σοφίσαι εἰς σωτηρίαν διὰ πίστεως τῆς ἐν Χριστῷ Ἰησοῦ.
\vs{16}πᾶσα γραφὴ θεόπνευστος καὶ ὠφέλιμος πρὸς διδασκαλίαν, πρὸς ἐλεγμόν, πρὸς ἐπανόρθωσιν, πρὸς παιδείαν τὴν ἐν δικαιοσύνῃ,
\vs{17}ἵνα ἄρτιος ᾖ ὁ τοῦ Θεοῦ ἄνθρωπος, πρὸς πᾶν ἔργον ἀγαθὸν ἐξηρτισμένος.

\ch{4}
Διαμαρτύρομαι ἐνώπιον τοῦ Θεοῦ καὶ Χριστοῦ Ἰησοῦ τοῦ μέλλοντος κρίνειν ζῶντας καὶ νεκρούς, καὶ τὴν ἐπιφάνειαν αὐτοῦ καὶ τὴν βασιλείαν αὐτοῦ·
\vs{2}κήρυξον τὸν λόγον, ἐπίστηθι εὐκαίρως ἀκαίρως, ἔλεγξον, ἐπιτίμησον, παρακάλεσον, ἐν πάσῃ μακροθυμίᾳ καὶ διδαχῇ.
\vs{3}Ἔσται γὰρ καιρὸς ὅτε τῆς ὑγιαινούσης διδασκαλίας οὐκ ἀνέξονται ἀλλὰ κατὰ τὰς ἰδίας ἐπιθυμίας ἑαυτοῖς ἐπισωρεύσουσιν διδασκάλους κνηθόμενοι τὴν ἀκοήν
\vs{4}καὶ ἀπὸ μὲν τῆς ἀληθείας τὴν ἀκοὴν ἀποστρέψουσιν, ἐπὶ δὲ τοὺς μύθους ἐκτραπήσονται.
\vs{5}Σὺ δὲ νῆφε ἐν πᾶσιν, κακοπάθησον, ἔργον ποίησον εὐαγγελιστοῦ, τὴν διακονίαν σου πληροφόρησον.

\vs{6}Ἐγὼ γὰρ ἤδη σπένδομαι, καὶ ὁ καιρὸς τῆς ἀναλύσεώς μου ἐφέστηκεν.
\vs{7}τὸν καλὸν ἀγῶνα ἠγώνισμαι, τὸν δρόμον τετέλεκα, τὴν πίστιν τετήρηκα·
\vs{8}λοιπὸν ἀπόκειταί μοι ὁ τῆς δικαιοσύνης στέφανος, ὃν ἀποδώσει μοι ὁ κύριος ἐν ἐκείνῃ τῇ ἡμέρᾳ, ὁ δίκαιος κριτής, οὐ μόνον δὲ ἐμοὶ ἀλλὰ καὶ πᾶσι τοῖς ἠγαπηκόσι τὴν ἐπιφάνειαν αὐτοῦ.

\vs{9}Σπούδασον ἐλθεῖν πρός με ταχέως·
\vs{10}Δημᾶς γάρ με ἐγκατέλιπεν ἀγαπήσας τὸν νῦν αἰῶνα καὶ ἐπορεύθη εἰς Θεσσαλονίκην, Κρήσκης εἰς Γαλατίαν, Τίτος εἰς Δαλματίαν·
\vs{11}Λουκᾶς ἐστιν μόνος μετ᾽ ἐμοῦ. Μᾶρκον ἀναλαβὼν ἄγε μετὰ σεαυτοῦ, ἔστιν γάρ μοι εὔχρηστος εἰς διακονίαν.
\vs{12}Τυχικὸν δὲ ἀπέστειλα εἰς Ἔφεσον.
\vs{13}τὸν φαιλόνην ὃν ἀπέλιπον ἐν Τρῳάδι παρὰ Κάρπῳ ἐρχόμενος φέρε, καὶ τὰ βιβλία μάλιστα τὰς μεμβράνας.
\vs{14}Ἀλέξανδρος ὁ χαλκεὺς πολλά μοι κακὰ ἐνεδείξατο· ἀποδώσει αὐτῷ ὁ Κύριος κατὰ τὰ ἔργα αὐτοῦ·
\vs{15}ὃν καὶ σὺ φυλάσσου, λίαν γὰρ ἀντέστη τοῖς ἡμετέροις λόγοις.
\vs{16}Ἐν τῇ πρώτῃ μου ἀπολογίᾳ οὐδείς μοι παρεγένετο, ἀλλὰ πάντες με ἐγκατέλιπον· μὴ αὐτοῖς λογισθείη·
\vs{17}ὁ δὲ Κύριός μοι παρέστη καὶ ἐνεδυνάμωσέν με, ἵνα δι᾽ ἐμοῦ τὸ κήρυγμα πληροφορηθῇ καὶ ἀκούσωσιν πάντα τὰ ἔθνη, καὶ ἐρρύσθην ἐκ στόματος λέοντος.
\vs{18}ῥύσεταί με ὁ Κύριος ἀπὸ παντὸς ἔργου πονηροῦ καὶ σώσει εἰς τὴν βασιλείαν αὐτοῦ τὴν ἐπουράνιον· ᾧ ἡ δόξα εἰς τοὺς αἰῶνας τῶν αἰώνων, ἀμήν.

\vs{19}Ἄσπασαι Πρίσκαν καὶ Ἀκύλαν καὶ τὸν Ὀνησιφόρου οἶκον.
\vs{20}Ἔραστος ἔμεινεν ἐν Κορίνθῳ, Τρόφιμον δὲ ἀπέλιπον ἐν Μιλήτῳ ἀσθενοῦντα.
\vs{21}Σπούδασον πρὸ χειμῶνος ἐλθεῖν. Ἀσπάζεταί σε Εὔβουλος καὶ Πούδης καὶ Λίνος καὶ Κλαυδία καὶ οἱ ἀδελφοὶ πάντες.

\vs{22}Ὁ Κύριος μετὰ τοῦ πνεύματός σου. ἡ χάρις μεθ᾽ ὑμῶν.


\def\book{ΠΡΟΣ ΤΙΤΟΝ}
\biblebook{ΠΡΟΣ ΤΙΤΟΝ}


\lettrine[lines=2, loversize=0.2, nindent=0em, findent=.25em]{\textcolor{bookheadingcolor}{Π}}{αῦλος} δοῦλος Θεοῦ, ἀπόστολος δὲ Ἰησοῦ Χριστοῦ κατὰ πίστιν ἐκλεκτῶν Θεοῦ καὶ ἐπίγνωσιν ἀληθείας τῆς κατ᾽ εὐσέβειαν
\vs{2}ἐπ᾽ ἐλπίδι ζωῆς αἰωνίου, ἣν ἐπηγγείλατο ὁ ἀψευδὴς Θεὸς πρὸ χρόνων αἰωνίων,
\vs{3}ἐφανέρωσεν δὲ καιροῖς ἰδίοις τὸν λόγον αὐτοῦ ἐν κηρύγματι, ὃ ἐπιστεύθην ἐγὼ κατ᾽ ἐπιταγὴν τοῦ Σωτῆρος ἡμῶν Θεοῦ,
\vs{4}Τίτῳ γνησίῳ τέκνῳ κατὰ κοινὴν πίστιν, Χάρις καὶ εἰρήνη ἀπὸ Θεοῦ Πατρὸς καὶ Χριστοῦ Ἰησοῦ τοῦ Σωτῆρος ἡμῶν.

\vs{5}Τούτου χάριν ἀπέλιπόν σε ἐν Κρήτῃ, ἵνα τὰ λείποντα ἐπιδιορθώσῃ καὶ καταστήσῃς κατὰ πόλιν πρεσβυτέρους, ὡς ἐγώ σοι διεταξάμην,
\vs{6}εἴ τίς ἐστιν ἀνέγκλητος, μιᾶς γυναικὸς ἀνήρ, τέκνα ἔχων πιστά, μὴ ἐν κατηγορίᾳ ἀσωτίας ἢ ἀνυπότακτα.
\vs{7}Δεῖ γὰρ τὸν ἐπίσκοπον ἀνέγκλητον εἶναι ὡς Θεοῦ οἰκονόμον, μὴ αὐθάδη, μὴ ὀργίλον, μὴ πάροινον, μὴ πλήκτην, μὴ αἰσχροκερδῆ,
\vs{8}ἀλλὰ φιλόξενον φιλάγαθον σώφρονα δίκαιον ὅσιον ἐγκρατῆ,
\vs{9}ἀντεχόμενον τοῦ κατὰ τὴν διδαχὴν πιστοῦ λόγου, ἵνα δυνατὸς ᾖ καὶ παρακαλεῖν ἐν τῇ διδασκαλίᾳ τῇ ὑγιαινούσῃ καὶ τοὺς ἀντιλέγοντας ἐλέγχειν.
\vs{10}Εἰσὶν γὰρ πολλοὶ καὶ ἀνυπότακτοι, ματαιολόγοι καὶ φρεναπάται, μάλιστα οἱ ἐκ τῆς περιτομῆς,
\vs{11}οὓς δεῖ ἐπιστομίζειν, οἵτινες ὅλους οἴκους ἀνατρέπουσιν διδάσκοντες ἃ μὴ δεῖ αἰσχροῦ κέρδους χάριν.
\vs{12}εἶπέν τις ἐξ αὐτῶν ἴδιος αὐτῶν προφήτης· 
\begin{poetryblock}

\begin{quote}Κρῆτες ἀεὶ ψεῦσται, κακὰ θηρία, γαστέρες ἀργαί.\end{quote}
\end{poetryblock}

\vs{13}Ἡ μαρτυρία αὕτη ἐστὶν ἀληθής. δι᾽ ἣν αἰτίαν ἔλεγχε αὐτοὺς ἀποτόμως, ἵνα ὑγιαίνωσιν ἐν τῇ πίστει,
\vs{14}μὴ προσέχοντες Ἰουδαϊκοῖς μύθοις καὶ ἐντολαῖς ἀνθρώπων ἀποστρεφομένων τὴν ἀλήθειαν.
\vs{15}Πάντα καθαρὰ τοῖς καθαροῖς· τοῖς δὲ μεμιαμμένοις καὶ ἀπίστοις οὐδὲν καθαρόν, ἀλλὰ μεμίανται αὐτῶν καὶ ὁ νοῦς καὶ ἡ συνείδησις.
\vs{16}Θεὸν ὁμολογοῦσιν εἰδέναι, τοῖς δὲ ἔργοις ἀρνοῦνται, βδελυκτοὶ ὄντες καὶ ἀπειθεῖς καὶ πρὸς πᾶν ἔργον ἀγαθὸν ἀδόκιμοι.

\ch{2}
Σὺ δὲ λάλει ἃ πρέπει τῇ ὑγιαινούσῃ διδασκαλίᾳ.
\vs{2}Πρεσβύτας νηφαλίους εἶναι, σεμνούς, σώφρονας, ὑγιαίνοντας τῇ πίστει, τῇ ἀγάπῃ, τῇ ὑπομονῇ·
\vs{3}πρεσβύτιδας ὡσαύτως ἐν καταστήματι ἱεροπρεπεῖς, μὴ διαβόλους μηδὲ οἴνῳ πολλῷ δεδουλωμένας, καλοδιδασκάλους,
\vs{4}ἵνα σωφρονίζωσιν τὰς νέας φιλάνδρους εἶναι, φιλοτέκνους
\vs{5}σώφρονας ἁγνάς οἰκουργούς ἀγαθάς, ὑποτασσομένας τοῖς ἰδίοις ἀνδράσιν, ἵνα μὴ ὁ λόγος τοῦ Θεοῦ βλασφημῆται.

\vs{6}Τοὺς νεωτέρους ὡσαύτως παρακάλει σωφρονεῖν
\vs{7}Περὶ πάντα, σεαυτὸν παρεχόμενος τύπον καλῶν ἔργων, ἐν τῇ διδασκαλίᾳ ἀφθορίαν, σεμνότητα,
\vs{8}λόγον ὑγιῆ ἀκατάγνωστον, ἵνα ὁ ἐξ ἐναντίας ἐντραπῇ μηδὲν ἔχων λέγειν περὶ ἡμῶν φαῦλον.

\vs{9}Δούλους ἰδίοις δεσπόταις ὑποτάσσεσθαι ἐν πᾶσιν, εὐαρέστους εἶναι, μὴ ἀντιλέγοντας,
\vs{10}μὴ νοσφιζομένους, ἀλλὰ πᾶσαν πίστιν ἐνδεικνυμένους ἀγαθήν, ἵνα τὴν διδασκαλίαν τὴν τοῦ Σωτῆρος ἡμῶν Θεοῦ κοσμῶσιν ἐν πᾶσιν.

\vs{11}Ἐπεφάνη γὰρ ἡ χάρις τοῦ Θεοῦ σωτήριος πᾶσιν ἀνθρώποις
\vs{12}παιδεύουσα ἡμᾶς, ἵνα ἀρνησάμενοι τὴν ἀσέβειαν καὶ τὰς κοσμικὰς ἐπιθυμίας σωφρόνως καὶ δικαίως καὶ εὐσεβῶς ζήσωμεν ἐν τῷ νῦν αἰῶνι,
\vs{13}προσδεχόμενοι τὴν μακαρίαν ἐλπίδα καὶ ἐπιφάνειαν τῆς δόξης τοῦ μεγάλου Θεοῦ καὶ Σωτῆρος ἡμῶν Ἰησοῦ Χριστοῦ,
\vs{14}ὃς ἔδωκεν ἑαυτὸν ὑπὲρ ἡμῶν, ἵνα λυτρώσηται ἡμᾶς ἀπὸ πάσης ἀνομίας καὶ καθαρίσῃ ἑαυτῷ λαὸν περιούσιον, ζηλωτὴν καλῶν ἔργων.
\vs{15}Ταῦτα λάλει καὶ παρακάλει καὶ ἔλεγχε μετὰ πάσης ἐπιταγῆς· μηδείς σου περιφρονείτω.

\ch{3}
Ὑπομίμνῃσκε αὐτοὺς ἀρχαῖς ἐξουσίαις ὑποτάσσεσθαι, πειθαρχεῖν, πρὸς πᾶν ἔργον ἀγαθὸν ἑτοίμους εἶναι,
\vs{2}μηδένα βλασφημεῖν, ἀμάχους εἶναι, ἐπιεικεῖς, πᾶσαν ἐνδεικνυμένους πραΰτητα πρὸς πάντας ἀνθρώπους.
\vs{3}Ἦμεν γάρ ποτε καὶ ἡμεῖς ἀνόητοι, ἀπειθεῖς, πλανώμενοι, δουλεύοντες ἐπιθυμίαις καὶ ἡδοναῖς ποικίλαις, ἐν κακίᾳ καὶ φθόνῳ διάγοντες, στυγητοί, μισοῦντες ἀλλήλους.

\vs{4}Ὅτε δὲ ἡ χρηστότης καὶ ἡ φιλανθρωπία ἐπεφάνη 
\begin{poetryblock}

\begin{quote}τοῦ Σωτῆρος ἡμῶν Θεοῦ,\end{quote}

\begin{quote} \vs{5}οὐκ ἐξ ἔργων τῶν ἐν δικαιοσύνῃ\end{quote} 

\begin{quote}ἃ ἐποιήσαμεν ἡμεῖς\end{quote} 

\begin{quote}ἀλλὰ κατὰ τὸ αὐτοῦ ἔλεος\end{quote} 

\begin{quote}ἔσωσεν ἡμᾶς διὰ λουτροῦ παλινγενεσίας\end{quote} 

\begin{quote}καὶ ἀνακαινώσεως Πνεύματος Ἁγίου,\end{quote}

\begin{quote} \vs{6}οὗ ἐξέχεεν ἐφ᾽ ἡμᾶς πλουσίως\end{quote} 

\begin{quote}διὰ Ἰησοῦ Χριστοῦ τοῦ Σωτῆρος ἡμῶν,\end{quote}

\begin{quote} \vs{7}ἵνα δικαιωθέντες τῇ ἐκείνου χάριτι\end{quote} 

\begin{quote}κληρονόμοι γενηθῶμεν κατ᾽ ἐλπίδα ζωῆς αἰωνίου.\end{quote}
\end{poetryblock}

\vs{8}Πιστὸς ὁ λόγος· καὶ περὶ τούτων βούλομαί σε διαβεβαιοῦσθαι, ἵνα φροντίζωσιν καλῶν ἔργων προΐστασθαι οἱ πεπιστευκότες Θεῷ· ταῦτά ἐστιν καλὰ καὶ ὠφέλιμα τοῖς ἀνθρώποις.
\vs{9}Μωρὰς δὲ ζητήσεις καὶ γενεαλογίας καὶ ἔρεις καὶ μάχας νομικὰς περιΐστασο· εἰσὶν γὰρ ἀνωφελεῖς καὶ μάταιοι.
\vs{10}αἱρετικὸν ἄνθρωπον μετὰ μίαν καὶ δευτέραν νουθεσίαν παραιτοῦ,
\vs{11}εἰδὼς ὅτι ἐξέστραπται ὁ τοιοῦτος καὶ ἁμαρτάνει ὢν αὐτοκατάκριτος.

\vs{12}Ὅταν πέμψω Ἀρτεμᾶν πρὸς σὲ ἢ Τυχικόν, σπούδασον ἐλθεῖν πρός με εἰς Νικόπολιν, ἐκεῖ γὰρ κέκρικα παραχειμάσαι.
\vs{13}Ζηνᾶν τὸν νομικὸν καὶ Ἀπολλῶν σπουδαίως πρόπεμψον, ἵνα μηδὲν αὐτοῖς λείπῃ.
\vs{14}μανθανέτωσαν δὲ καὶ οἱ ἡμέτεροι καλῶν ἔργων προΐστασθαι εἰς τὰς ἀναγκαίας χρείας, ἵνα μὴ ὦσιν ἄκαρποι.

\vs{15}Ἀσπάζονταί σε οἱ μετ᾽ ἐμοῦ πάντες. Ἄσπασαι τοὺς φιλοῦντας ἡμᾶς ἐν πίστει.

Ἡ χάρις μετὰ πάντων ὑμῶν.


\def\book{ΠΡΟΣ ΦΙΛΗΜΟΝΑ}
\biblebook{ΠΡΟΣ ΦΙΛΗΜΟΝΑ}

 
\lettrine[lines=2, loversize=0.2, nindent=0em, findent=.25em]{\textcolor{bookheadingcolor}{Π}}{αῦλος} δέσμιος Χριστοῦ Ἰησοῦ καὶ Τιμόθεος ὁ ἀδελφὸς Φιλήμονι τῷ ἀγαπητῷ καὶ συνεργῷ ἡμῶν
\vs{2}καὶ Ἀπφίᾳ τῇ ἀδελφῇ καὶ Ἀρχίππῳ τῷ συστρατιώτῃ ἡμῶν καὶ τῇ κατ᾽ οἶκόν σου ἐκκλησίᾳ,
\vs{3}Χάρις ὑμῖν καὶ εἰρήνη ἀπὸ Θεοῦ Πατρὸς ἡμῶν καὶ Κυρίου Ἰησοῦ Χριστοῦ.

\vs{4}Εὐχαριστῶ τῷ Θεῷ μου πάντοτε μνείαν σου ποιούμενος ἐπὶ τῶν προσευχῶν μου,
\vs{5}ἀκούων σου τὴν ἀγάπην καὶ τὴν πίστιν, ἣν ἔχεις πρὸς τὸν Κύριον Ἰησοῦν καὶ εἰς πάντας τοὺς ἁγίους,
\vs{6}ὅπως ἡ κοινωνία τῆς πίστεώς σου ἐνεργὴς γένηται ἐν ἐπιγνώσει παντὸς ἀγαθοῦ τοῦ ἐν ἡμῖν εἰς Χριστόν.
\vs{7}χαρὰν γὰρ πολλὴν ἔσχον καὶ παράκλησιν ἐπὶ τῇ ἀγάπῃ σου, ὅτι τὰ σπλάγχνα τῶν ἁγίων ἀναπέπαυται διὰ σοῦ, ἀδελφέ.

\vs{8}Διό πολλὴν ἐν Χριστῷ παρρησίαν ἔχων ἐπιτάσσειν σοι τὸ ἀνῆκον
\vs{9}διὰ τὴν ἀγάπην μᾶλλον παρακαλῶ, τοιοῦτος ὢν ὡς Παῦλος πρεσβύτης νυνὶ δὲ καὶ δέσμιος Χριστοῦ Ἰησοῦ·
\vs{10}Παρακαλῶ σε περὶ τοῦ ἐμοῦ τέκνου, ὃν ἐγέννησα ἐν τοῖς δεσμοῖς, Ὀνήσιμον,
\vs{11}τόν ποτέ σοι ἄχρηστον νυνὶ δὲ καὶ σοὶ καὶ ἐμοὶ εὔχρηστον,
\vs{12}ὃν ἀνέπεμψά σοι, αὐτόν, τοῦτ᾽ ἔστιν τὰ ἐμὰ σπλάγχνα·
\vs{13}ὃν ἐγὼ ἐβουλόμην πρὸς ἐμαυτὸν κατέχειν, ἵνα ὑπὲρ σοῦ μοι διακονῇ ἐν τοῖς δεσμοῖς τοῦ εὐαγγελίου,
\vs{14}Χωρὶς δὲ τῆς σῆς γνώμης οὐδὲν ἠθέλησα ποιῆσαι, ἵνα μὴ ὡς κατὰ ἀνάγκην τὸ ἀγαθόν σου ᾖ ἀλλὰ κατὰ ἑκούσιον.

\vs{15}τάχα γὰρ διὰ τοῦτο ἐχωρίσθη πρὸς ὥραν, ἵνα αἰώνιον αὐτὸν ἀπέχῃς,
\vs{16}οὐκέτι ὡς δοῦλον ἀλλὰ ὑπὲρ δοῦλον, ἀδελφὸν ἀγαπητόν, μάλιστα ἐμοί, πόσῳ δὲ μᾶλλον σοὶ καὶ ἐν σαρκὶ καὶ ἐν Κυρίῳ.
\vs{17}Εἰ οὖν με ἔχεις κοινωνόν, προσλαβοῦ αὐτὸν ὡς ἐμέ.
\vs{18}εἰ δέ τι ἠδίκησέν σε ἢ ὀφείλει, τοῦτο ἐμοὶ ἐλλόγα.
\vs{19}ἐγὼ Παῦλος ἔγραψα τῇ ἐμῇ χειρί, ἐγὼ ἀποτίσω· ἵνα μὴ λέγω σοι ὅτι καὶ σεαυτόν μοι προσοφείλεις.
\vs{20}Ναί ἀδελφέ, ἐγώ σου ὀναίμην ἐν Κυρίῳ· ἀνάπαυσόν μου τὰ σπλάγχνα ἐν Χριστῷ.

\vs{21}Πεποιθὼς τῇ ὑπακοῇ σου ἔγραψά σοι, εἰδὼς ὅτι καὶ ὑπὲρ ἃ λέγω ποιήσεις.
\vs{22}Ἅμα δὲ καὶ ἑτοίμαζέ μοι ξενίαν· ἐλπίζω γὰρ ὅτι διὰ τῶν προσευχῶν ὑμῶν χαρισθήσομαι ὑμῖν.

\vs{23}Ἀσπάζεταί σε Ἐπαφρᾶς ὁ συναιχμάλωτός μου ἐν Χριστῷ Ἰησοῦ,
\vs{24}Μᾶρκος, Ἀρίσταρχος, Δημᾶς, Λουκᾶς, οἱ συνεργοί μου.

\vs{25}Ἡ χάρις τοῦ Κυρίου Ἰησοῦ Χριστοῦ μετὰ τοῦ πνεύματος ὑμῶν.


\def\book{ΠΡΟΣ ΕΒΡΑΙΟΥΣ}
\biblebook{ΠΡΟΣ ΕΒΡΑΙΟΥΣ}


\lettrine[lines=2, loversize=0.2, nindent=0em, findent=.25em]{\textcolor{bookheadingcolor}{Π}}{ολυμερῶς} καὶ πολυτρόπως πάλαι ὁ Θεὸς λαλήσας τοῖς πατράσιν ἐν τοῖς προφήταις
\vs{2}ἐπ᾽ ἐσχάτου τῶν ἡμερῶν τούτων ἐλάλησεν ἡμῖν ἐν Υἱῷ, ὃν ἔθηκεν κληρονόμον πάντων, δι᾽ οὗ καὶ ἐποίησεν τοὺς αἰῶνας·

\vs{3}ὃς ὢν ἀπαύγασμα τῆς δόξης καὶ χαρακτὴρ τῆς ὑποστάσεως αὐτοῦ, 
\begin{poetryblock}

\begin{quote}φέρων τε τὰ πάντα τῷ ῥήματι τῆς δυνάμεως αὐτοῦ,\end{quote} 

\begin{quote}καθαρισμὸν τῶν ἁμαρτιῶν ποιησάμενος\end{quote} 

\begin{quote}ἐκάθισεν ἐν δεξιᾷ τῆς Μεγαλωσύνης ἐν ὑψηλοῖς,\end{quote}

\begin{quote} \vs{4}τοσούτῳ κρείττων γενόμενος τῶν ἀγγέλων\end{quote} 

\begin{quote}ὅσῳ διαφορώτερον παρ᾽ αὐτοὺς κεκληρονόμηκεν ὄνομα.\end{quote}
\end{poetryblock}

\vs{5}Τίνι γὰρ εἶπέν ποτε τῶν ἀγγέλων· 
\begin{poetryblock}

\begin{quote}Υἱός μου εἶ σύ,\end{quote} 

\begin{quote}ἐγὼ σήμερον γεγέννηκά σε;\end{quote}
\end{poetryblock}

Καὶ πάλιν· 
\begin{poetryblock}

\begin{quote}Ἐγὼ ἔσομαι αὐτῷ εἰς Πατέρα,\end{quote} 

\begin{quote}καὶ αὐτὸς ἔσται μοι εἰς Υἱόν;\end{quote}
\end{poetryblock}

\vs{6}Ὅταν δὲ πάλιν εἰσαγάγῃ τὸν πρωτότοκον εἰς τὴν οἰκουμένην, λέγει· 
\begin{poetryblock}

\begin{quote}Καὶ προσκυνησάτωσαν αὐτῷ πάντες ἄγγελοι Θεοῦ.\end{quote}
\end{poetryblock}

\vs{7}Καὶ πρὸς μὲν τοὺς ἀγγέλους λέγει· 
\begin{poetryblock}

\begin{quote}Ὁ ποιῶν τοὺς ἀγγέλους αὐτοῦ πνεύματα\end{quote} 

\begin{quote}καὶ τοὺς λειτουργοὺς αὐτοῦ πυρὸς φλόγα,\end{quote}
\end{poetryblock}

\vs{8}Πρὸς δὲ τὸν Υἱόν· 
\begin{poetryblock}

\begin{quote}Ὁ θρόνος σου ὁ Θεὸς εἰς τὸν αἰῶνα τοῦ αἰῶνος,\end{quote} 

\begin{quote}καὶ ἡ ῥάβδος τῆς εὐθύτητος ῥάβδος τῆς βασιλείας σου.\end{quote}

\begin{quote} \vs{9}ἠγάπησας δικαιοσύνην καὶ ἐμίσησας ἀνομίαν·\end{quote} 

\begin{quote}διὰ τοῦτο ἔχρισέν σε ὁ Θεός ὁ Θεός σου\end{quote} 

\begin{quote}ἔλαιον ἀγαλλιάσεως παρὰ τοὺς μετόχους σου.\end{quote}
\end{poetryblock}

\vs{10}Καί· 
\begin{poetryblock}

\begin{quote}Σὺ κατ᾽ ἀρχάς, Κύριε, τὴν γῆν ἐθεμελίωσας,\end{quote} 

\begin{quote}καὶ ἔργα τῶν χειρῶν σού εἰσιν οἱ οὐρανοί·\end{quote}

\begin{quote} \vs{11}αὐτοὶ ἀπολοῦνται, σὺ δὲ διαμένεις,\end{quote} 

\begin{quote}καὶ πάντες ὡς ἱμάτιον παλαιωθήσονται,\end{quote}

\begin{quote} \vs{12}καὶ ὡσεὶ περιβόλαιον ἑλίξεις αὐτούς,\end{quote} 

\begin{quote}ὡς ἱμάτιον καὶ ἀλλαγήσονται·\end{quote} 

\begin{quote}σὺ δὲ ὁ αὐτὸς εἶ καὶ τὰ ἔτη σου οὐκ ἐκλείψουσιν.\end{quote}
\end{poetryblock}

\vs{13}Πρὸς τίνα δὲ τῶν ἀγγέλων εἴρηκέν ποτε· 
\begin{poetryblock}

\begin{quote}Κάθου ἐκ δεξιῶν μου,\end{quote} 

\begin{quote}ἕως ἂν θῶ τοὺς ἐχθρούς σου ὑποπόδιον τῶν ποδῶν σου;\end{quote}
\end{poetryblock}

\vs{14}Οὐχὶ πάντες εἰσὶν λειτουργικὰ πνεύματα εἰς διακονίαν ἀποστελλόμενα διὰ τοὺς μέλλοντας κληρονομεῖν σωτηρίαν;

\ch{2}
Διὰ τοῦτο δεῖ περισσοτέρως προσέχειν ἡμᾶς τοῖς ἀκουσθεῖσιν, μήποτε παραρυῶμεν.
\vs{2}εἰ γὰρ ὁ δι᾽ ἀγγέλων λαληθεὶς λόγος ἐγένετο βέβαιος καὶ πᾶσα παράβασις καὶ παρακοὴ ἔλαβεν ἔνδικον μισθαποδοσίαν,
\vs{3}πῶς ἡμεῖς ἐκφευξόμεθα τηλικαύτης ἀμελήσαντες σωτηρίας, ἥτις ἀρχὴν λαβοῦσα λαλεῖσθαι διὰ τοῦ Κυρίου ὑπὸ τῶν ἀκουσάντων εἰς ἡμᾶς ἐβεβαιώθη,
\vs{4}συνεπιμαρτυροῦντος τοῦ Θεοῦ σημείοις τε καὶ τέρασιν καὶ ποικίλαις δυνάμεσιν καὶ Πνεύματος Ἁγίου μερισμοῖς κατὰ τὴν αὐτοῦ θέλησιν;

\vs{5}Οὐ γὰρ ἀγγέλοις ὑπέταξεν τὴν οἰκουμένην τὴν μέλλουσαν, περὶ ἧς λαλοῦμεν.
\vs{6}διεμαρτύρατο δέ πού τις λέγων· 
\begin{poetryblock}

\begin{quote}Τί ἐστιν ἄνθρωπος ὅτι μιμνῄσκῃ αὐτοῦ,\end{quote} 

\begin{quote}ἢ υἱὸς ἀνθρώπου ὅτι ἐπισκέπτῃ αὐτόν;\end{quote}

\begin{quote} \vs{7}ἠλάττωσας αὐτὸν βραχύ τι παρ᾽ ἀγγέλους,\end{quote} 

\begin{quote}δόξῃ καὶ τιμῇ ἐστεφάνωσας αὐτόν,\end{quote}

\begin{quote} \vs{8}πάντα ὑπέταξας ὑποκάτω τῶν ποδῶν αὐτοῦ.\end{quote}
\end{poetryblock}

Ἐν τῷ γὰρ ὑποτάξαι αὐτῷ τὰ πάντα οὐδὲν ἀφῆκεν αὐτῷ ἀνυπότακτον. νῦν δὲ οὔπω ὁρῶμεν αὐτῷ τὰ πάντα ὑποτεταγμένα·
\vs{9}τὸν δὲ βραχύ τι παρ᾽ ἀγγέλους ἠλαττωμένον βλέπομεν Ἰησοῦν διὰ τὸ πάθημα τοῦ θανάτου δόξῃ καὶ τιμῇ ἐστεφανωμένον, ὅπως χάριτι Θεοῦ ὑπὲρ παντὸς γεύσηται θανάτου.

\vs{10}Ἔπρεπεν γὰρ αὐτῷ, δι᾽ ὃν τὰ πάντα καὶ δι᾽ οὗ τὰ πάντα, πολλοὺς υἱοὺς εἰς δόξαν ἀγαγόντα τὸν ἀρχηγὸν τῆς σωτηρίας αὐτῶν διὰ παθημάτων τελειῶσαι.
\vs{11}ὅ τε γὰρ ἁγιάζων καὶ οἱ ἁγιαζόμενοι ἐξ ἑνὸς πάντες· δι᾽ ἣν αἰτίαν οὐκ ἐπαισχύνεται ἀδελφοὺς αὐτοὺς καλεῖν
\vs{12}λέγων· 
\begin{poetryblock}

\begin{quote}Ἀπαγγελῶ τὸ ὄνομά σου τοῖς ἀδελφοῖς μου,\end{quote} 

\begin{quote}ἐν μέσῳ ἐκκλησίας ὑμνήσω σε,\end{quote}
\end{poetryblock}

\vs{13}Καὶ πάλιν· 
\begin{poetryblock}

\begin{quote}Ἐγὼ ἔσομαι πεποιθὼς ἐπ᾽ αὐτῷ,\end{quote}
\end{poetryblock}

Καὶ πάλιν· 
\begin{poetryblock}

\begin{quote}Ἰδοὺ ἐγὼ καὶ τὰ παιδία ἅ μοι ἔδωκεν ὁ Θεός.\end{quote}
\end{poetryblock}

\vs{14}Ἐπεὶ οὖν τὰ παιδία κεκοινώνηκεν αἵματος καὶ σαρκός, καὶ αὐτὸς παραπλησίως μετέσχεν τῶν αὐτῶν, ἵνα διὰ τοῦ θανάτου καταργήσῃ τὸν τὸ κράτος ἔχοντα τοῦ θανάτου, τοῦτ᾽ ἔστιν τὸν διάβολον,
\vs{15}καὶ ἀπαλλάξῃ τούτους, ὅσοι φόβῳ θανάτου διὰ παντὸς τοῦ ζῆν ἔνοχοι ἦσαν δουλείας.
\vs{16}Οὐ γὰρ δήπου ἀγγέλων ἐπιλαμβάνεται ἀλλὰ σπέρματος Ἀβραὰμ ἐπιλαμβάνεται.
\vs{17}ὅθεν ὤφειλεν κατὰ πάντα τοῖς ἀδελφοῖς ὁμοιωθῆναι, ἵνα ἐλεήμων γένηται καὶ πιστὸς ἀρχιερεὺς τὰ πρὸς τὸν Θεόν εἰς τὸ ἱλάσκεσθαι τὰς ἁμαρτίας τοῦ λαοῦ.
\vs{18}ἐν ᾧ γὰρ πέπονθεν αὐτὸς πειρασθείς, δύναται τοῖς πειραζομένοις βοηθῆσαι.

\ch{3}
Ὅθεν, ἀδελφοὶ ἅγιοι, κλήσεως ἐπουρανίου μέτοχοι, κατανοήσατε τὸν Ἀπόστολον καὶ Ἀρχιερέα τῆς ὁμολογίας ἡμῶν Ἰησοῦν,
\vs{2}πιστὸν ὄντα τῷ ποιήσαντι αὐτὸν ὡς καὶ Μωϋσῆς ἐν ὅλῳ τῷ οἴκῳ αὐτοῦ.
\vs{3}Πλείονος γὰρ οὗτος δόξης παρὰ Μωϋσῆν ἠξίωται, καθ᾽ ὅσον πλείονα τιμὴν ἔχει τοῦ οἴκου ὁ κατασκευάσας αὐτόν·
\vs{4}πᾶς γὰρ οἶκος κατασκευάζεται ὑπό τινος, ὁ δὲ πάντα κατασκευάσας Θεός.
\vs{5}Καὶ Μωϋσῆς μὲν πιστὸς ἐν ὅλῳ τῷ οἴκῳ αὐτοῦ ὡς θεράπων εἰς μαρτύριον τῶν λαληθησομένων,
\vs{6}Χριστὸς δὲ ὡς υἱὸς ἐπὶ τὸν οἶκον αὐτοῦ· οὗ οἶκός ἐσμεν ἡμεῖς, ἐὰν τὴν παρρησίαν καὶ τὸ καύχημα τῆς ἐλπίδος κατάσχωμεν.

\vs{7}Διό, καθὼς λέγει τὸ Πνεῦμα τὸ Ἅγιον· 
\begin{poetryblock}

\begin{quote}Σήμερον ἐὰν τῆς φωνῆς αὐτοῦ ἀκούσητε,\end{quote}

\begin{quote} \vs{8}μὴ σκληρύνητε τὰς καρδίας ὑμῶν ὡς ἐν τῷ παραπικρασμῷ\end{quote} 

\begin{quote}κατὰ τὴν ἡμέραν τοῦ πειρασμοῦ ἐν τῇ ἐρήμῳ,\end{quote}

\begin{quote} \vs{9}οὗ ἐπείρασαν οἱ πατέρες ὑμῶν ἐν δοκιμασίᾳ\end{quote} 

\begin{quote}καὶ εἶδον τὰ ἔργα μου\end{quote}
\end{poetryblock}

\vs{10}τεσσεράκοντα ἔτη· 
\begin{poetryblock}

\begin{quote}διὸ προσώχθισα τῇ γενεᾷ ταύτῃ\end{quote} 

\begin{quote}καὶ εἶπον· Ἀεὶ πλανῶνται τῇ καρδίᾳ,\end{quote} 

\begin{quote}αὐτοὶ δὲ οὐκ ἔγνωσαν τὰς ὁδούς μου,\end{quote}

\begin{quote} \vs{11}ὡς ὤμοσα ἐν τῇ ὀργῇ μου·\end{quote} 

\begin{quote}Εἰ εἰσελεύσονται εἰς τὴν κατάπαυσίν μου.\end{quote}
\end{poetryblock}

\vs{12}Βλέπετε, ἀδελφοί, μήποτε ἔσται ἔν τινι ὑμῶν καρδία πονηρὰ ἀπιστίας ἐν τῷ ἀποστῆναι ἀπὸ Θεοῦ ζῶντος,
\vs{13}ἀλλὰ παρακαλεῖτε ἑαυτοὺς καθ᾽ ἑκάστην ἡμέραν, ἄχρις οὗ τὸ Σήμερον καλεῖται, ἵνα μὴ σκληρυνθῇ τις ἐξ ὑμῶν ἀπάτῃ τῆς ἁμαρτίας—
\vs{14}Μέτοχοι γὰρ τοῦ Χριστοῦ γεγόναμεν, ἐάνπερ τὴν ἀρχὴν τῆς ὑποστάσεως μέχρι τέλους βεβαίαν κατάσχωμεν—
\vs{15}ἐν τῷ λέγεσθαι· 
\begin{poetryblock}

\begin{quote}Σήμερον ἐὰν τῆς φωνῆς αὐτοῦ ἀκούσητε,\end{quote} 

\begin{quote}Μὴ σκληρύνητε τὰς καρδίας ὑμῶν ὡς ἐν τῷ παραπικρασμῷ.\end{quote}
\end{poetryblock}

\vs{16}Τίνες γὰρ ἀκούσαντες παρεπίκραναν; ἀλλ᾽ οὐ πάντες οἱ ἐξελθόντες ἐξ Αἰγύπτου διὰ Μωϋσέως;
\vs{17}τίσιν δὲ προσώχθισεν τεσσεράκοντα ἔτη; οὐχὶ τοῖς ἁμαρτήσασιν, ὧν τὰ κῶλα ἔπεσεν ἐν τῇ ἐρήμῳ;
\vs{18}τίσιν δὲ ὤμοσεν μὴ εἰσελεύσεσθαι εἰς τὴν κατάπαυσιν αὐτοῦ εἰ μὴ τοῖς ἀπειθήσασιν;
\vs{19}καὶ βλέπομεν ὅτι οὐκ ἠδυνήθησαν εἰσελθεῖν δι᾽ ἀπιστίαν.

\ch{4}
Φοβηθῶμεν οὖν, μήποτε καταλειπομένης ἐπαγγελίας εἰσελθεῖν εἰς τὴν κατάπαυσιν αὐτοῦ δοκῇ τις ἐξ ὑμῶν ὑστερηκέναι.
\vs{2}καὶ γάρ ἐσμεν εὐηγγελισμένοι καθάπερ κἀκεῖνοι· ἀλλ᾽ οὐκ ὠφέλησεν ὁ λόγος τῆς ἀκοῆς ἐκείνους μὴ συγκεκερασμένους τῇ πίστει τοῖς ἀκούσασιν.
\vs{3}Εἰσερχόμεθα γὰρ εἰς τὴν κατάπαυσιν οἱ πιστεύσαντες, καθὼς εἴρηκεν· 
\begin{poetryblock}

\begin{quote}Ὡς ὤμοσα ἐν τῇ ὀργῇ μου·\end{quote} 

\begin{quote}Εἰ εἰσελεύσονται εἰς τὴν κατάπαυσίν μου,\end{quote}
\end{poetryblock}

Καίτοι τῶν ἔργων ἀπὸ καταβολῆς κόσμου γενηθέντων.
\vs{4}εἴρηκεν γάρ που περὶ τῆς ἑβδόμης οὕτως· Καὶ κατέπαυσεν ὁ Θεὸς ἐν τῇ ἡμέρᾳ τῇ ἑβδόμῃ ἀπὸ πάντων τῶν ἔργων αὐτοῦ,
\vs{5}καὶ ἐν τούτῳ πάλιν· Εἰ εἰσελεύσονται εἰς τὴν κατάπαυσίν μου.
\vs{6}Ἐπεὶ οὖν ἀπολείπεται τινὰς εἰσελθεῖν εἰς αὐτήν, καὶ οἱ πρότερον εὐαγγελισθέντες οὐκ εἰσῆλθον δι᾽ ἀπείθειαν,
\vs{7}πάλιν τινὰ ὁρίζει ἡμέραν, Σήμερον, ἐν Δαυὶδ λέγων μετὰ τοσοῦτον χρόνον, καθὼς προείρηται· 
\begin{poetryblock}

\begin{quote}Σήμερον ἐὰν τῆς φωνῆς αὐτοῦ ἀκούσητε,\end{quote} 

\begin{quote}μὴ σκληρύνητε τὰς καρδίας ὑμῶν.\end{quote}
\end{poetryblock}

\vs{8}Εἰ γὰρ αὐτοὺς Ἰησοῦς κατέπαυσεν, οὐκ ἂν περὶ ἄλλης ἐλάλει μετὰ ταῦτα ἡμέρας.
\vs{9}ἄρα ἀπολείπεται σαββατισμὸς τῷ λαῷ τοῦ Θεοῦ.
\vs{10}ὁ γὰρ εἰσελθὼν εἰς τὴν κατάπαυσιν αὐτοῦ καὶ αὐτὸς κατέπαυσεν ἀπὸ τῶν ἔργων αὐτοῦ ὥσπερ ἀπὸ τῶν ἰδίων ὁ Θεός.
\vs{11}Σπουδάσωμεν οὖν εἰσελθεῖν εἰς ἐκείνην τὴν κατάπαυσιν, ἵνα μὴ ἐν τῷ αὐτῷ τις ὑποδείγματι πέσῃ τῆς ἀπειθείας.
\vs{12}Ζῶν γὰρ ὁ λόγος τοῦ Θεοῦ καὶ ἐνεργὴς καὶ τομώτερος ὑπὲρ πᾶσαν μάχαιραν δίστομον καὶ διϊκνούμενος ἄχρι μερισμοῦ ψυχῆς καὶ πνεύματος, ἁρμῶν τε καὶ μυελῶν, καὶ κριτικὸς ἐνθυμήσεων καὶ ἐννοιῶν καρδίας·
\vs{13}καὶ οὐκ ἔστιν κτίσις ἀφανὴς ἐνώπιον αὐτοῦ, πάντα δὲ γυμνὰ καὶ τετραχηλισμένα τοῖς ὀφθαλμοῖς αὐτοῦ, πρὸς ὃν ἡμῖν ὁ λόγος.

\vs{14}Ἔχοντες οὖν ἀρχιερέα μέγαν διεληλυθότα τοὺς οὐρανούς, Ἰησοῦν τὸν Υἱὸν τοῦ Θεοῦ, κρατῶμεν τῆς ὁμολογίας.
\vs{15}οὐ γὰρ ἔχομεν ἀρχιερέα μὴ δυνάμενον συμπαθῆσαι ταῖς ἀσθενείαις ἡμῶν, πεπειρασμένον δὲ κατὰ πάντα καθ᾽ ὁμοιότητα χωρὶς ἁμαρτίας.
\vs{16}προσερχώμεθα οὖν μετὰ παρρησίας τῷ θρόνῳ τῆς χάριτος, ἵνα λάβωμεν ἔλεος καὶ χάριν εὕρωμεν εἰς εὔκαιρον βοήθειαν.

\ch{5}
Πᾶς γὰρ ἀρχιερεὺς ἐξ ἀνθρώπων λαμβανόμενος ὑπὲρ ἀνθρώπων καθίσταται τὰ πρὸς τὸν Θεόν, ἵνα προσφέρῃ δῶρά τε καὶ θυσίας ὑπὲρ ἁμαρτιῶν,
\vs{2}μετριοπαθεῖν δυνάμενος τοῖς ἀγνοοῦσιν καὶ πλανωμένοις, ἐπεὶ καὶ αὐτὸς περίκειται ἀσθένειαν
\vs{3}καὶ δι᾽ αὐτὴν ὀφείλει, καθὼς περὶ τοῦ λαοῦ, οὕτως καὶ περὶ αὑτοῦ προσφέρειν περὶ ἁμαρτιῶν.
\vs{4}Καὶ οὐχ ἑαυτῷ τις λαμβάνει τὴν τιμήν ἀλλὰ καλούμενος ὑπὸ τοῦ Θεοῦ καθώσπερ καὶ Ἀαρών.

\vs{5}Οὕτως καὶ ὁ Χριστὸς οὐχ ἑαυτὸν ἐδόξασεν γενηθῆναι ἀρχιερέα ἀλλ᾽ ὁ λαλήσας πρὸς αὐτόν· 
\begin{poetryblock}

\begin{quote}Υἱός μου εἶ σύ, ἐγὼ σήμερον γεγέννηκά σε·\end{quote}
\end{poetryblock}

\vs{6}Καθὼς καὶ ἐν ἑτέρῳ λέγει· 
\begin{poetryblock}

\begin{quote}Σὺ ἱερεὺς εἰς τὸν αἰῶνα κατὰ τὴν τάξιν Μελχισέδεκ,\end{quote}
\end{poetryblock}

\vs{7}Ὃς ἐν ταῖς ἡμέραις τῆς σαρκὸς αὐτοῦ δεήσεις τε καὶ ἱκετηρίας πρὸς τὸν δυνάμενον σῴζειν αὐτὸν ἐκ θανάτου μετὰ κραυγῆς ἰσχυρᾶς καὶ δακρύων προσενέγκας καὶ εἰσακουσθεὶς ἀπὸ τῆς εὐλαβείας,
\vs{8}καίπερ ὢν Υἱός, ἔμαθεν ἀφ᾽ ὧν ἔπαθεν τὴν ὑπακοήν,
\vs{9}καὶ τελειωθεὶς ἐγένετο πᾶσιν τοῖς ὑπακούουσιν αὐτῷ αἴτιος σωτηρίας αἰωνίου,
\vs{10}προσαγορευθεὶς ὑπὸ τοῦ Θεοῦ ἀρχιερεὺς κατὰ τὴν τάξιν Μελχισέδεκ.

\vs{11}Περὶ οὗ πολὺς ἡμῖν ὁ λόγος καὶ δυσερμήνευτος λέγειν, ἐπεὶ νωθροὶ γεγόνατε ταῖς ἀκοαῖς.
\vs{12}καὶ γὰρ ὀφείλοντες εἶναι διδάσκαλοι διὰ τὸν χρόνον, πάλιν χρείαν ἔχετε τοῦ διδάσκειν ὑμᾶς τινα τὰ στοιχεῖα τῆς ἀρχῆς τῶν λογίων τοῦ Θεοῦ καὶ γεγόνατε χρείαν ἔχοντες γάλακτος καὶ οὐ στερεᾶς τροφῆς.
\vs{13}πᾶς γὰρ ὁ μετέχων γάλακτος ἄπειρος λόγου δικαιοσύνης, νήπιος γάρ ἐστιν·
\vs{14}τελείων δέ ἐστιν ἡ στερεὰ τροφή, τῶν διὰ τὴν ἕξιν τὰ αἰσθητήρια γεγυμνασμένα ἐχόντων πρὸς διάκρισιν καλοῦ τε καὶ κακοῦ.

\ch{6}
Διὸ ἀφέντες τὸν τῆς ἀρχῆς τοῦ Χριστοῦ λόγον ἐπὶ τὴν τελειότητα φερώμεθα, μὴ πάλιν θεμέλιον καταβαλλόμενοι μετανοίας ἀπὸ νεκρῶν ἔργων καὶ πίστεως ἐπὶ Θεόν,
\vs{2}βαπτισμῶν διδαχὴν ἐπιθέσεώς τε χειρῶν, ἀναστάσεώς τε νεκρῶν καὶ κρίματος αἰωνίου.
\vs{3}καὶ τοῦτο ποιήσομεν, ἐάνπερ ἐπιτρέπῃ ὁ Θεός.

\vs{4}Ἀδύνατον γὰρ τοὺς ἅπαξ φωτισθέντας, γευσαμένους τε τῆς δωρεᾶς τῆς ἐπουρανίου καὶ μετόχους γενηθέντας Πνεύματος Ἁγίου
\vs{5}καὶ καλὸν γευσαμένους Θεοῦ ῥῆμα δυνάμεις τε μέλλοντος αἰῶνος
\vs{6}καὶ παραπεσόντας, πάλιν ἀνακαινίζειν εἰς μετάνοιαν, ἀνασταυροῦντας ἑαυτοῖς τὸν Υἱὸν τοῦ Θεοῦ καὶ παραδειγματίζοντας.
\vs{7}Γῆ γὰρ ἡ πιοῦσα τὸν ἐπ᾽ αὐτῆς ἐρχόμενον πολλάκις ὑετόν καὶ τίκτουσα βοτάνην εὔθετον ἐκείνοις δι᾽ οὓς καὶ γεωργεῖται, μεταλαμβάνει εὐλογίας ἀπὸ τοῦ Θεοῦ·
\vs{8}ἐκφέρουσα δὲ ἀκάνθας καὶ τριβόλους, ἀδόκιμος καὶ κατάρας ἐγγύς, ἧς τὸ τέλος εἰς καῦσιν.

\vs{9}Πεπείσμεθα δὲ περὶ ὑμῶν, ἀγαπητοί, τὰ κρείσσονα καὶ ἐχόμενα σωτηρίας, εἰ καὶ οὕτως λαλοῦμεν.
\vs{10}οὐ γὰρ ἄδικος ὁ Θεὸς ἐπιλαθέσθαι τοῦ ἔργου ὑμῶν καὶ τῆς ἀγάπης ἧς ἐνεδείξασθε εἰς τὸ ὄνομα αὐτοῦ, διακονήσαντες τοῖς ἁγίοις καὶ διακονοῦντες.
\vs{11}Ἐπιθυμοῦμεν δὲ ἕκαστον ὑμῶν τὴν αὐτὴν ἐνδείκνυσθαι σπουδὴν πρὸς τὴν πληροφορίαν τῆς ἐλπίδος ἄχρι τέλους,
\vs{12}ἵνα μὴ νωθροὶ γένησθε, μιμηταὶ δὲ τῶν διὰ πίστεως καὶ μακροθυμίας κληρονομούντων τὰς ἐπαγγελίας.

\vs{13}Τῷ γὰρ Ἀβραὰμ ἐπαγγειλάμενος ὁ Θεός, ἐπεὶ κατ᾽ οὐδενὸς εἶχεν μείζονος ὀμόσαι, ὤμοσεν καθ᾽ ἑαυτοῦ
\vs{14}λέγων· 
\begin{poetryblock}

\begin{quote}Εἰ μὴν εὐλογῶν εὐλογήσω σε καὶ πληθύνων πληθυνῶ σε·\end{quote}
\end{poetryblock}

\vs{15}καὶ οὕτως μακροθυμήσας ἐπέτυχεν τῆς ἐπαγγελίας.
\vs{16}Ἄνθρωποι γὰρ κατὰ τοῦ μείζονος ὀμνύουσιν, καὶ πάσης αὐτοῖς ἀντιλογίας πέρας εἰς βεβαίωσιν ὁ ὅρκος·
\vs{17}ἐν ᾧ περισσότερον βουλόμενος ὁ Θεὸς ἐπιδεῖξαι τοῖς κληρονόμοις τῆς ἐπαγγελίας τὸ ἀμετάθετον τῆς βουλῆς αὐτοῦ ἐμεσίτευσεν ὅρκῳ,
\vs{18}ἵνα διὰ δύο πραγμάτων ἀμεταθέτων, ἐν οἷς ἀδύνατον ψεύσασθαι τὸν Θεόν, ἰσχυρὰν παράκλησιν ἔχωμεν οἱ καταφυγόντες κρατῆσαι τῆς προκειμένης ἐλπίδος·
\vs{19}ἣν ὡς ἄγκυραν ἔχομεν τῆς ψυχῆς ἀσφαλῆ τε καὶ βεβαίαν καὶ εἰσερχομένην εἰς τὸ ἐσώτερον τοῦ καταπετάσματος,
\vs{20}ὅπου πρόδρομος ὑπὲρ ἡμῶν εἰσῆλθεν Ἰησοῦς, κατὰ τὴν τάξιν Μελχισέδεκ ἀρχιερεὺς γενόμενος εἰς τὸν αἰῶνα.

\ch{7}
Οὗτος γὰρ ὁ Μελχισέδεκ, βασιλεὺς Σαλήμ, ἱερεὺς τοῦ Θεοῦ τοῦ Ὑψίστου, ὁ συναντήσας Ἀβραὰμ ὑποστρέφοντι ἀπὸ τῆς κοπῆς τῶν βασιλέων καὶ εὐλογήσας αὐτόν,
\vs{2}ᾧ καὶ δεκάτην ἀπὸ πάντων ἐμέρισεν Ἀβραάμ, πρῶτον μὲν ἑρμηνευόμενος Βασιλεὺς δικαιοσύνης ἔπειτα δὲ καὶ Βασιλεὺς Σαλήμ, ὅ ἐστιν Βασιλεὺς εἰρήνης,
\vs{3}ἀπάτωρ ἀμήτωρ ἀγενεαλόγητος, μήτε ἀρχὴν ἡμερῶν μήτε ζωῆς τέλος ἔχων, ἀφωμοιωμένος δὲ τῷ Υἱῷ τοῦ Θεοῦ, μένει ἱερεὺς εἰς τὸ διηνεκές.

\vs{4}Θεωρεῖτε δὲ πηλίκος οὗτος, ᾧ καὶ δεκάτην Ἀβραὰμ ἔδωκεν ἐκ τῶν ἀκροθινίων ὁ πατριάρχης.
\vs{5}καὶ οἱ μὲν ἐκ τῶν υἱῶν Λευὶ τὴν ἱερατείαν λαμβάνοντες ἐντολὴν ἔχουσιν ἀποδεκατοῦν τὸν λαὸν κατὰ τὸν νόμον, τοῦτ᾽ ἔστιν τοὺς ἀδελφοὺς αὐτῶν, καίπερ ἐξεληλυθότας ἐκ τῆς ὀσφύος Ἀβραάμ·
\vs{6}ὁ δὲ μὴ γενεαλογούμενος ἐξ αὐτῶν δεδεκάτωκεν Ἀβραάμ καὶ τὸν ἔχοντα τὰς ἐπαγγελίας εὐλόγηκεν.
\vs{7}χωρὶς δὲ πάσης ἀντιλογίας τὸ ἔλαττον ὑπὸ τοῦ κρείττονος εὐλογεῖται.
\vs{8}Καὶ ὧδε μὲν δεκάτας ἀποθνῄσκοντες ἄνθρωποι λαμβάνουσιν, ἐκεῖ δὲ μαρτυρούμενος ὅτι ζῇ.
\vs{9}καὶ ὡς ἔπος εἰπεῖν, δι᾽ Ἀβραὰμ καὶ Λευὶ ὁ δεκάτας λαμβάνων δεδεκάτωται·
\vs{10}ἔτι γὰρ ἐν τῇ ὀσφύϊ τοῦ πατρὸς ἦν ὅτε συνήντησεν αὐτῷ Μελχισέδεκ.

\vs{11}Εἰ μὲν οὖν τελείωσις διὰ τῆς Λευιτικῆς ἱερωσύνης ἦν, ὁ λαὸς γὰρ ἐπ᾽ αὐτῆς νενομοθέτηται, τίς ἔτι χρεία κατὰ τὴν τάξιν Μελχισέδεκ ἕτερον ἀνίστασθαι ἱερέα καὶ οὐ κατὰ τὴν τάξιν Ἀαρὼν λέγεσθαι;
\vs{12}μετατιθεμένης γὰρ τῆς ἱερωσύνης ἐξ ἀνάγκης καὶ νόμου μετάθεσις γίνεται.
\vs{13}Ἐφ᾽ ὃν γὰρ λέγεται ταῦτα, φυλῆς ἑτέρας μετέσχηκεν, ἀφ᾽ ἧς οὐδεὶς προσέσχηκεν τῷ θυσιαστηρίῳ·
\vs{14}πρόδηλον γὰρ ὅτι ἐξ Ἰούδα ἀνατέταλκεν ὁ Κύριος ἡμῶν, εἰς ἣν φυλὴν περὶ ἱερέων οὐδὲν Μωϋσῆς ἐλάλησεν.
\vs{15}Καὶ περισσότερον ἔτι κατάδηλόν ἐστιν, εἰ κατὰ τὴν ὁμοιότητα Μελχισέδεκ ἀνίσταται ἱερεὺς ἕτερος,
\vs{16}ὃς οὐ κατὰ νόμον ἐντολῆς σαρκίνης γέγονεν ἀλλὰ κατὰ δύναμιν ζωῆς ἀκαταλύτου.
\vs{17}μαρτυρεῖται γὰρ ὅτι 
\begin{poetryblock}

\begin{quote}Σὺ ἱερεὺς εἰς τὸν αἰῶνα κατὰ τὴν τάξιν Μελχισέδεκ.\end{quote}
\end{poetryblock}

\vs{18}Ἀθέτησις μὲν γὰρ γίνεται προαγούσης ἐντολῆς διὰ τὸ αὐτῆς ἀσθενὲς καὶ ἀνωφελές—
\vs{19}οὐδὲν γὰρ ἐτελείωσεν ὁ νόμος— ἐπεισαγωγὴ δὲ κρείττονος ἐλπίδος δι᾽ ἧς ἐγγίζομεν τῷ Θεῷ.
\vs{20}Καὶ καθ᾽ ὅσον οὐ χωρὶς ὁρκωμοσίας· οἱ μὲν γὰρ χωρὶς ὁρκωμοσίας εἰσὶν ἱερεῖς γεγονότες,
\vs{21}ὁ δὲ μετὰ ὁρκωμοσίας διὰ τοῦ λέγοντος πρὸς αὐτόν· 
\begin{poetryblock}

\begin{quote}Ὤμοσεν Κύριος καὶ οὐ μεταμεληθήσεται·\end{quote} 

\begin{quote}Σὺ ἱερεὺς εἰς τὸν αἰῶνα.\end{quote}
\end{poetryblock}
\vs{22}Κατὰ τοσοῦτο καὶ κρείττονος διαθήκης γέγονεν ἔγγυος Ἰησοῦς.
\vs{23}Καὶ οἱ μὲν πλείονές εἰσιν γεγονότες ἱερεῖς διὰ τὸ θανάτῳ κωλύεσθαι παραμένειν·
\vs{24}ὁ δὲ διὰ τὸ μένειν αὐτὸν εἰς τὸν αἰῶνα ἀπαράβατον ἔχει τὴν ἱερωσύνην·
\vs{25}ὅθεν καὶ σῴζειν εἰς τὸ παντελὲς δύναται τοὺς προσερχομένους δι᾽ αὐτοῦ τῷ Θεῷ, πάντοτε ζῶν εἰς τὸ ἐντυγχάνειν ὑπὲρ αὐτῶν.

\vs{26}Τοιοῦτος γὰρ ἡμῖν καὶ ἔπρεπεν ἀρχιερεύς, ὅσιος ἄκακος ἀμίαντος, κεχωρισμένος ἀπὸ τῶν ἁμαρτωλῶν καὶ ὑψηλότερος τῶν οὐρανῶν γενόμενος,
\vs{27}ὃς οὐκ ἔχει καθ᾽ ἡμέραν ἀνάγκην, ὥσπερ οἱ ἀρχιερεῖς, πρότερον ὑπὲρ τῶν ἰδίων ἁμαρτιῶν θυσίας ἀναφέρειν ἔπειτα τῶν τοῦ λαοῦ· τοῦτο γὰρ ἐποίησεν ἐφάπαξ ἑαυτὸν ἀνενέγκας.
\vs{28}ὁ νόμος γὰρ ἀνθρώπους καθίστησιν ἀρχιερεῖς ἔχοντας ἀσθένειαν, ὁ λόγος δὲ τῆς ὁρκωμοσίας τῆς μετὰ τὸν νόμον Υἱόν εἰς τὸν αἰῶνα τετελειωμένον.

\ch{8}
Κεφάλαιον δὲ ἐπὶ τοῖς λεγομένοις, τοιοῦτον ἔχομεν ἀρχιερέα, ὃς ἐκάθισεν ἐν δεξιᾷ τοῦ θρόνου τῆς Μεγαλωσύνης ἐν τοῖς οὐρανοῖς,
\vs{2}τῶν ἁγίων λειτουργὸς καὶ τῆς σκηνῆς τῆς ἀληθινῆς, ἣν ἔπηξεν ὁ Κύριος, οὐκ ἄνθρωπος.
\vs{3}Πᾶς γὰρ ἀρχιερεὺς εἰς τὸ προσφέρειν δῶρά τε καὶ θυσίας καθίσταται· ὅθεν ἀναγκαῖον ἔχειν τι καὶ τοῦτον ὃ προσενέγκῃ.
\vs{4}εἰ μὲν οὖν ἦν ἐπὶ γῆς, οὐδ᾽ ἂν ἦν ἱερεύς, ὄντων τῶν προσφερόντων κατὰ νόμον τὰ δῶρα·
\vs{5}οἵτινες ὑποδείγματι καὶ σκιᾷ λατρεύουσιν τῶν ἐπουρανίων, καθὼς κεχρημάτισται Μωϋσῆς μέλλων ἐπιτελεῖν τὴν σκηνήν· Ὅρα γάρ φησίν, Ποιήσεις πάντα κατὰ τὸν τύπον τὸν δειχθέντα σοι ἐν τῷ ὄρει·
\vs{6}Νυνὶ δὲ διαφορωτέρας τέτυχεν λειτουργίας, ὅσῳ καὶ κρείττονός ἐστιν διαθήκης μεσίτης, ἥτις ἐπὶ κρείττοσιν ἐπαγγελίαις νενομοθέτηται.

\vs{7}εἰ γὰρ ἡ πρώτη ἐκείνη ἦν ἄμεμπτος, οὐκ ἂν δευτέρας ἐζητεῖτο τόπος.
\vs{8}μεμφόμενος γὰρ αὐτοὺς λέγει· 
\begin{poetryblock}

\begin{quote}Ἰδοὺ ἡμέραι ἔρχονται, λέγει Κύριος,\end{quote} 

\begin{quote}καὶ συντελέσω ἐπὶ τὸν οἶκον Ἰσραὴλ\end{quote} 

\begin{quote}καὶ ἐπὶ τὸν οἶκον Ἰούδα διαθήκην καινήν,\end{quote}

\begin{quote} \vs{9}οὐ κατὰ τὴν διαθήκην, ἣν ἐποίησα τοῖς πατράσιν αὐτῶν\end{quote} 

\begin{quote}ἐν ἡμέρᾳ ἐπιλαβομένου μου τῆς χειρὸς αὐτῶν\end{quote} 

\begin{quote}ἐξαγαγεῖν αὐτοὺς ἐκ γῆς Αἰγύπτου,\end{quote} 

\begin{quote}ὅτι αὐτοὶ οὐκ ἐνέμειναν ἐν τῇ διαθήκῃ μου,\end{quote} 

\begin{quote}κἀγὼ ἠμέλησα αὐτῶν, λέγει Κύριος·\end{quote}

\begin{quote} \vs{10}ὅτι αὕτη ἡ διαθήκη, ἣν διαθήσομαι τῷ οἴκῳ Ἰσραὴλ\end{quote} 

\begin{quote}μετὰ τὰς ἡμέρας ἐκείνας, λέγει Κύριος·\end{quote} 

\begin{quote}διδοὺς νόμους μου εἰς τὴν διάνοιαν αὐτῶν\end{quote} 

\begin{quote}καὶ ἐπὶ καρδίας αὐτῶν ἐπιγράψω αὐτούς,\end{quote} 

\begin{quote}καὶ ἔσομαι αὐτοῖς εἰς Θεόν,\end{quote} 

\begin{quote}καὶ αὐτοὶ ἔσονταί μοι εἰς λαόν·\end{quote}

\begin{quote} \vs{11}καὶ οὐ μὴ διδάξωσιν ἕκαστος τὸν πολίτην αὐτοῦ\end{quote} 

\begin{quote}καὶ ἕκαστος τὸν ἀδελφὸν αὐτοῦ λέγων· Γνῶθι τὸν Κύριον,\end{quote} 

\begin{quote}ὅτι πάντες εἰδήσουσίν με\end{quote} 

\begin{quote}ἀπὸ μικροῦ ἕως μεγάλου αὐτῶν,\end{quote}

\begin{quote} \vs{12}ὅτι ἵλεως ἔσομαι ταῖς ἀδικίαις αὐτῶν\end{quote} 

\begin{quote}καὶ τῶν ἁμαρτιῶν αὐτῶν οὐ μὴ μνησθῶ ἔτι.\end{quote}
\end{poetryblock}

\vs{13}Ἐν τῷ λέγειν Καινὴν πεπαλαίωκεν τὴν πρώτην· τὸ δὲ παλαιούμενον καὶ γηράσκον ἐγγὺς ἀφανισμοῦ.

\ch{9}
Εἶχε μὲν οὖν καὶ ἡ πρώτη δικαιώματα λατρείας τό τε ἅγιον κοσμικόν.
\vs{2}σκηνὴ γὰρ κατεσκευάσθη ἡ πρώτη ἐν ᾗ ἥ τε λυχνία καὶ ἡ τράπεζα καὶ ἡ πρόθεσις τῶν ἄρτων, ἥτις λέγεται Ἅγια·
\vs{3}μετὰ δὲ τὸ δεύτερον καταπέτασμα σκηνὴ ἡ λεγομένη Ἅγια ἁγίων,
\vs{4}χρυσοῦν ἔχουσα θυμιατήριον καὶ τὴν κιβωτὸν τῆς διαθήκης περικεκαλυμμένην πάντοθεν χρυσίῳ, ἐν ᾗ στάμνος χρυσῆ ἔχουσα τὸ μάννα καὶ ἡ ῥάβδος Ἀαρὼν ἡ βλαστήσασα καὶ αἱ πλάκες τῆς διαθήκης,
\vs{5}ὑπεράνω δὲ αὐτῆς Χερουβὶν δόξης κατασκιάζοντα τὸ ἱλαστήριον· περὶ ὧν οὐκ ἔστιν νῦν λέγειν κατὰ μέρος.
\vs{6}Τούτων δὲ οὕτως κατεσκευασμένων εἰς μὲν τὴν πρώτην σκηνὴν διὰ παντὸς εἰσίασιν οἱ ἱερεῖς τὰς λατρείας ἐπιτελοῦντες,
\vs{7}εἰς δὲ τὴν δευτέραν ἅπαξ τοῦ ἐνιαυτοῦ μόνος ὁ ἀρχιερεύς, οὐ χωρὶς αἵματος ὃ προσφέρει ὑπὲρ ἑαυτοῦ καὶ τῶν τοῦ λαοῦ ἀγνοημάτων,
\vs{8}Τοῦτο δηλοῦντος τοῦ Πνεύματος τοῦ Ἁγίου, μήπω πεφανερῶσθαι τὴν τῶν ἁγίων ὁδὸν ἔτι τῆς πρώτης σκηνῆς ἐχούσης στάσιν,
\vs{9}ἥτις παραβολὴ εἰς τὸν καιρὸν τὸν ἐνεστηκότα, καθ᾽ ἣν δῶρά τε καὶ θυσίαι προσφέρονται μὴ δυνάμεναι κατὰ συνείδησιν τελειῶσαι τὸν λατρεύοντα,
\vs{10}μόνον ἐπὶ βρώμασιν καὶ πόμασιν καὶ διαφόροις βαπτισμοῖς, δικαιώματα σαρκὸς μέχρι καιροῦ διορθώσεως ἐπικείμενα.

\vs{11}Χριστὸς δὲ παραγενόμενος ἀρχιερεὺς τῶν γενομένων ἀγαθῶν διὰ τῆς μείζονος καὶ τελειοτέρας σκηνῆς οὐ χειροποιήτου, τοῦτ᾽ ἔστιν οὐ ταύτης τῆς κτίσεως,
\vs{12}οὐδὲ δι᾽ αἵματος τράγων καὶ μόσχων διὰ δὲ τοῦ ἰδίου αἵματος εἰσῆλθεν ἐφάπαξ εἰς τὰ ἅγια αἰωνίαν λύτρωσιν εὑράμενος.
\vs{13}Εἰ γὰρ τὸ αἷμα τράγων καὶ ταύρων καὶ σποδὸς δαμάλεως ῥαντίζουσα τοὺς κεκοινωμένους ἁγιάζει πρὸς τὴν τῆς σαρκὸς καθαρότητα,
\vs{14}πόσῳ μᾶλλον τὸ αἷμα τοῦ Χριστοῦ, ὃς διὰ Πνεύματος αἰωνίου ἑαυτὸν προσήνεγκεν ἄμωμον τῷ Θεῷ, καθαριεῖ τὴν συνείδησιν ἡμῶν ἀπὸ νεκρῶν ἔργων εἰς τὸ λατρεύειν Θεῷ ζῶντι.
\vs{15}Καὶ διὰ τοῦτο διαθήκης καινῆς μεσίτης ἐστίν, ὅπως θανάτου γενομένου εἰς ἀπολύτρωσιν τῶν ἐπὶ τῇ πρώτῃ διαθήκῃ παραβάσεων τὴν ἐπαγγελίαν λάβωσιν οἱ κεκλημένοι τῆς αἰωνίου κληρονομίας.
\vs{16}Ὅπου γὰρ διαθήκη, θάνατον ἀνάγκη φέρεσθαι τοῦ διαθεμένου·
\vs{17}διαθήκη γὰρ ἐπὶ νεκροῖς βεβαία, ἐπεὶ μήποτε ἰσχύει ὅτε ζῇ ὁ διαθέμενος.
\vs{18}Ὅθεν οὐδὲ ἡ πρώτη χωρὶς αἵματος ἐνκεκαίνισται·
\vs{19}λαληθείσης γὰρ πάσης ἐντολῆς κατὰ τὸν νόμον ὑπὸ Μωϋσέως παντὶ τῷ λαῷ, λαβὼν τὸ αἷμα τῶν μόσχων καὶ τῶν τράγων μετὰ ὕδατος καὶ ἐρίου κοκκίνου καὶ ὑσσώπου αὐτό τε τὸ βιβλίον καὶ πάντα τὸν λαὸν ἐράντισεν
\vs{20}λέγων· 
\begin{poetryblock}

\begin{quote}Τοῦτο τὸ αἷμα τῆς διαθήκης ἧς ἐνετείλατο πρὸς ὑμᾶς ὁ Θεός.\end{quote}
\end{poetryblock}

\vs{21}Καὶ τὴν σκηνὴν δὲ καὶ πάντα τὰ σκεύη τῆς λειτουργίας τῷ αἵματι ὁμοίως ἐράντισεν.
\vs{22}καὶ σχεδὸν ἐν αἵματι πάντα καθαρίζεται κατὰ τὸν νόμον καὶ χωρὶς αἱματεκχυσίας οὐ γίνεται ἄφεσις.

\vs{23}Ἀνάγκη οὖν τὰ μὲν ὑποδείγματα τῶν ἐν τοῖς οὐρανοῖς τούτοις καθαρίζεσθαι, αὐτὰ δὲ τὰ ἐπουράνια κρείττοσιν θυσίαις παρὰ ταύτας.
\vs{24}οὐ γὰρ εἰς χειροποίητα εἰσῆλθεν ἅγια Χριστός, ἀντίτυπα τῶν ἀληθινῶν, ἀλλ᾽ εἰς αὐτὸν τὸν οὐρανόν, νῦν ἐμφανισθῆναι τῷ προσώπῳ τοῦ Θεοῦ ὑπὲρ ἡμῶν·
\vs{25}οὐδ᾽ ἵνα πολλάκις προσφέρῃ ἑαυτόν, ὥσπερ ὁ ἀρχιερεὺς εἰσέρχεται εἰς τὰ ἅγια κατ᾽ ἐνιαυτὸν ἐν αἵματι ἀλλοτρίῳ,
\vs{26}Ἐπεὶ ἔδει αὐτὸν πολλάκις παθεῖν ἀπὸ καταβολῆς κόσμου· νυνὶ δὲ ἅπαξ ἐπὶ συντελείᾳ τῶν αἰώνων εἰς ἀθέτησιν τῆς ἁμαρτίας διὰ τῆς θυσίας αὐτοῦ πεφανέρωται.
\vs{27}Καὶ καθ᾽ ὅσον ἀπόκειται τοῖς ἀνθρώποις ἅπαξ ἀποθανεῖν, μετὰ δὲ τοῦτο κρίσις,
\vs{28}οὕτως καὶ ὁ Χριστός ἅπαξ προσενεχθεὶς εἰς τὸ πολλῶν ἀνενεγκεῖν ἁμαρτίας ἐκ δευτέρου χωρὶς ἁμαρτίας ὀφθήσεται τοῖς αὐτὸν ἀπεκδεχομένοις εἰς σωτηρίαν.

\ch{10}
Σκιὰν γὰρ ἔχων ὁ νόμος τῶν μελλόντων ἀγαθῶν, οὐκ αὐτὴν τὴν εἰκόνα τῶν πραγμάτων, κατ᾽ ἐνιαυτὸν ταῖς αὐταῖς θυσίαις ἃς προσφέρουσιν εἰς τὸ διηνεκὲς οὐδέποτε δύναται τοὺς προσερχομένους τελειῶσαι·
\vs{2}ἐπεὶ οὐκ ἂν ἐπαύσαντο προσφερόμεναι διὰ τὸ μηδεμίαν ἔχειν ἔτι συνείδησιν ἁμαρτιῶν τοὺς λατρεύοντας ἅπαξ κεκαθαρισμένους;
\vs{3}Ἀλλ᾽ ἐν αὐταῖς ἀνάμνησις ἁμαρτιῶν κατ᾽ ἐνιαυτόν·
\vs{4}ἀδύνατον γὰρ αἷμα ταύρων καὶ τράγων ἀφαιρεῖν ἁμαρτίας.
\vs{5}Διὸ εἰσερχόμενος εἰς τὸν κόσμον λέγει· 
\begin{poetryblock}

\begin{quote}Θυσίαν καὶ προσφορὰν οὐκ ἠθέλησας,\end{quote} 

\begin{quote}Σῶμα δὲ κατηρτίσω μοι·\end{quote}

\begin{quote} \vs{6}Ὁλοκαυτώματα καὶ περὶ ἁμαρτίας Οὐκ εὐδόκησας.\end{quote}

\begin{quote} \vs{7}Τότε εἶπον· Ἰδοὺ ἥκω,\end{quote} 

\begin{quote}Ἐν κεφαλίδι βιβλίου γέγραπται περὶ ἐμοῦ,\end{quote} 

\begin{quote}Τοῦ ποιῆσαι ὁ Θεός τὸ θέλημά σου.\end{quote}
\end{poetryblock}

\vs{8}Ἀνώτερον λέγων ὅτι 
\begin{poetryblock}

\begin{quote}Θυσίας καὶ προσφορὰς καὶ ὁλοκαυτώματα καὶ περὶ ἁμαρτίας\end{quote} 

\begin{quote}οὐκ ἠθέλησας οὐδὲ εὐδόκησας,\end{quote}
\end{poetryblock}

αἵτινες κατὰ νόμον προσφέρονται,
\vs{9}τότε εἴρηκεν· Ἰδοὺ ἥκω τοῦ ποιῆσαι τὸ θέλημά σου. ἀναιρεῖ τὸ πρῶτον ἵνα τὸ δεύτερον στήσῃ,
\vs{10}ἐν ᾧ θελήματι ἡγιασμένοι ἐσμὲν διὰ τῆς προσφορᾶς τοῦ σώματος Ἰησοῦ Χριστοῦ ἐφάπαξ.

\vs{11}Καὶ πᾶς μὲν ἱερεὺς ἕστηκεν καθ᾽ ἡμέραν λειτουργῶν καὶ τὰς αὐτὰς πολλάκις προσφέρων θυσίας, αἵτινες οὐδέποτε δύνανται περιελεῖν ἁμαρτίας,
\vs{12}οὗτος δὲ μίαν ὑπὲρ ἁμαρτιῶν προσενέγκας θυσίαν εἰς τὸ διηνεκὲς ἐκάθισεν ἐν δεξιᾷ τοῦ Θεοῦ,
\vs{13}τὸ λοιπὸν ἐκδεχόμενος ἕως τεθῶσιν οἱ ἐχθροὶ αὐτοῦ ὑποπόδιον τῶν ποδῶν αὐτοῦ.
\vs{14}μιᾷ γὰρ προσφορᾷ τετελείωκεν εἰς τὸ διηνεκὲς τοὺς ἁγιαζομένους.
\vs{15}Μαρτυρεῖ δὲ ἡμῖν καὶ τὸ Πνεῦμα τὸ Ἅγιον· μετὰ γὰρ τὸ εἰρηκέναι·
\begin{poetryblock}

\begin{quote} \vs{16}Αὕτη ἡ διαθήκη ἣν διαθήσομαι πρὸς αὐτοὺς\end{quote} 

\begin{quote}μετὰ τὰς ἡμέρας ἐκείνας, λέγει Κύριος·\end{quote} 

\begin{quote}διδοὺς νόμους μου ἐπὶ καρδίας αὐτῶν\end{quote} 

\begin{quote}καὶ ἐπὶ τὴν διάνοιαν αὐτῶν ἐπιγράψω αὐτούς,\end{quote}

\begin{quote} \vs{17}Καὶ Τῶν ἁμαρτιῶν αὐτῶν καὶ τῶν ἀνομιῶν αὐτῶν\end{quote} 

\begin{quote}οὐ μὴ μνησθήσομαι ἔτι.\end{quote}
\end{poetryblock}

\vs{18}Ὅπου δὲ ἄφεσις τούτων, οὐκέτι προσφορὰ περὶ ἁμαρτίας.

\vs{19}Ἔχοντες οὖν, ἀδελφοί, παρρησίαν εἰς τὴν εἴσοδον τῶν ἁγίων ἐν τῷ αἵματι Ἰησοῦ,
\vs{20}ἣν ἐνεκαίνισεν ἡμῖν ὁδὸν πρόσφατον καὶ ζῶσαν διὰ τοῦ καταπετάσματος, τοῦτ᾽ ἔστιν τῆς σαρκὸς αὐτοῦ,
\vs{21}καὶ ἱερέα μέγαν ἐπὶ τὸν οἶκον τοῦ Θεοῦ,
\vs{22}προσερχώμεθα μετὰ ἀληθινῆς καρδίας ἐν πληροφορίᾳ πίστεως ῥεραντισμένοι τὰς καρδίας ἀπὸ συνειδήσεως πονηρᾶς καὶ λελουσμένοι τὸ σῶμα ὕδατι καθαρῷ·
\vs{23}Κατέχωμεν τὴν ὁμολογίαν τῆς ἐλπίδος ἀκλινῆ, πιστὸς γὰρ ὁ ἐπαγγειλάμενος,
\vs{24}καὶ κατανοῶμεν ἀλλήλους εἰς παροξυσμὸν ἀγάπης καὶ καλῶν ἔργων,
\vs{25}μὴ ἐγκαταλείποντες τὴν ἐπισυναγωγὴν ἑαυτῶν, καθὼς ἔθος τισίν, ἀλλὰ παρακαλοῦντες, καὶ τοσούτῳ μᾶλλον ὅσῳ βλέπετε ἐγγίζουσαν τὴν ἡμέραν.

\vs{26}Ἑκουσίως γὰρ ἁμαρτανόντων ἡμῶν μετὰ τὸ λαβεῖν τὴν ἐπίγνωσιν τῆς ἀληθείας, οὐκέτι περὶ ἁμαρτιῶν ἀπολείπεται θυσία,
\vs{27}φοβερὰ δέ τις ἐκδοχὴ κρίσεως καὶ πυρὸς ζῆλος ἐσθίειν μέλλοντος τοὺς ὑπεναντίους.
\vs{28}ἀθετήσας τις νόμον Μωϋσέως χωρὶς οἰκτιρμῶν ἐπὶ δυσὶν ἢ τρισὶν μάρτυσιν ἀποθνήσκει·
\vs{29}πόσῳ δοκεῖτε χείρονος ἀξιωθήσεται τιμωρίας ὁ τὸν Υἱὸν τοῦ Θεοῦ καταπατήσας καὶ τὸ αἷμα τῆς διαθήκης κοινὸν ἡγησάμενος, ἐν ᾧ ἡγιάσθη, καὶ τὸ Πνεῦμα τῆς χάριτος ἐνυβρίσας;
\vs{30}Οἴδαμεν γὰρ τὸν εἰπόντα· 
\begin{poetryblock}

\begin{quote}Ἐμοὶ ἐκδίκησις, ἐγὼ ἀνταποδώσω. καὶ πάλιν· Κρινεῖ Κύριος τὸν λαὸν αὐτοῦ.\end{quote}
\end{poetryblock}
\vs{31}φοβερὸν τὸ ἐμπεσεῖν εἰς χεῖρας Θεοῦ ζῶντος.

\vs{32}Ἀναμιμνῄσκεσθε δὲ τὰς πρότερον ἡμέρας, ἐν αἷς φωτισθέντες πολλὴν ἄθλησιν ὑπεμείνατε παθημάτων,
\vs{33}τοῦτο μὲν ὀνειδισμοῖς τε καὶ θλίψεσιν θεατριζόμενοι, τοῦτο δὲ κοινωνοὶ τῶν οὕτως ἀναστρεφομένων γενηθέντες.
\vs{34}καὶ γὰρ τοῖς δεσμίοις συνεπαθήσατε καὶ τὴν ἁρπαγὴν τῶν ὑπαρχόντων ὑμῶν μετὰ χαρᾶς προσεδέξασθε γινώσκοντες ἔχειν ἑαυτοὺς κρείττονα ὕπαρξιν καὶ μένουσαν.
\vs{35}Μὴ ἀποβάλητε οὖν τὴν παρρησίαν ὑμῶν, ἥτις ἔχει μεγάλην μισθαποδοσίαν.
\vs{36}ὑπομονῆς γὰρ ἔχετε χρείαν ἵνα τὸ θέλημα τοῦ Θεοῦ ποιήσαντες κομίσησθε τὴν ἐπαγγελίαν.

\vs{37}ἔτι γὰρ Μικρὸν ὅσον ὅσον, 
\begin{poetryblock}

\begin{quote}ὁ ἐρχόμενος ἥξει καὶ οὐ χρονίσει·\end{quote}

\begin{quote} \vs{38}ὁ δὲ δίκαιός μου ἐκ πίστεως ζήσεται,\end{quote} 

\begin{quote}καὶ ἐὰν ὑποστείληται, οὐκ εὐδοκεῖ ἡ ψυχή μου ἐν αὐτῷ.\end{quote}
\end{poetryblock}

\vs{39}Ἡμεῖς δὲ οὐκ ἐσμὲν ὑποστολῆς εἰς ἀπώλειαν ἀλλὰ πίστεως εἰς περιποίησιν ψυχῆς.

\ch{11}
Ἔστιν δὲ πίστις ἐλπιζομένων ὑπόστασις, πραγμάτων ἔλεγχος οὐ βλεπομένων.
\vs{2}ἐν ταύτῃ γὰρ ἐμαρτυρήθησαν οἱ πρεσβύτεροι.
\vs{3}Πίστει νοοῦμεν κατηρτίσθαι τοὺς αἰῶνας ῥήματι Θεοῦ, εἰς τὸ μὴ ἐκ φαινομένων τὸ βλεπόμενον γεγονέναι.
\vs{4}Πίστει πλείονα θυσίαν Ἅβελ παρὰ Κάϊν προσήνεγκεν τῷ Θεῷ, δι᾽ ἧς ἐμαρτυρήθη εἶναι δίκαιος, μαρτυροῦντος ἐπὶ τοῖς δώροις αὐτοῦ τοῦ Θεοῦ, καὶ δι᾽ αὐτῆς ἀποθανὼν ἔτι λαλεῖ.
\vs{5}Πίστει Ἑνὼχ μετετέθη τοῦ μὴ ἰδεῖν θάνατον, καὶ οὐχ ηὑρίσκετο διότι μετέθηκεν αὐτὸν ὁ Θεός. πρὸ γὰρ τῆς μεταθέσεως μεμαρτύρηται εὐαρεστηκέναι τῷ Θεῷ·
\vs{6}χωρὶς δὲ πίστεως ἀδύνατον εὐαρεστῆσαι· πιστεῦσαι γὰρ δεῖ τὸν προσερχόμενον τῷ Θεῷ ὅτι ἔστιν καὶ τοῖς ἐκζητοῦσιν αὐτὸν μισθαποδότης γίνεται.
\vs{7}Πίστει χρηματισθεὶς Νῶε περὶ τῶν μηδέπω βλεπομένων, εὐλαβηθεὶς κατεσκεύασεν κιβωτὸν εἰς σωτηρίαν τοῦ οἴκου αὐτοῦ δι᾽ ἧς κατέκρινεν τὸν κόσμον, καὶ τῆς κατὰ πίστιν δικαιοσύνης ἐγένετο κληρονόμος.

\vs{8}Πίστει καλούμενος Ἀβραὰμ ὑπήκουσεν ἐξελθεῖν εἰς τόπον ὃν ἤμελλεν λαμβάνειν εἰς κληρονομίαν, καὶ ἐξῆλθεν μὴ ἐπιστάμενος ποῦ ἔρχεται.
\vs{9}Πίστει παρῴκησεν εἰς γῆν τῆς ἐπαγγελίας ὡς ἀλλοτρίαν ἐν σκηναῖς κατοικήσας μετὰ Ἰσαὰκ καὶ Ἰακὼβ τῶν συνκληρονόμων τῆς ἐπαγγελίας τῆς αὐτῆς·
\vs{10}ἐξεδέχετο γὰρ τὴν τοὺς θεμελίους ἔχουσαν πόλιν ἧς τεχνίτης καὶ δημιουργὸς ὁ Θεός.
\vs{11}Πίστει καὶ αὐτῇ Σάρρᾳ στεῖρα δύναμιν εἰς καταβολὴν σπέρματος ἔλαβεν καὶ παρὰ καιρὸν ἡλικίας, ἐπεὶ πιστὸν ἡγήσατο τὸν ἐπαγγειλάμενον.
\vs{12}διὸ καὶ ἀφ᾽ ἑνὸς ἐγεννήθησαν, καὶ ταῦτα νενεκρωμένου, καθὼς τὰ ἄστρα τοῦ οὐρανοῦ τῷ πλήθει καὶ ὡς ἡ ἄμμος ἡ παρὰ τὸ χεῖλος τῆς θαλάσσης ἡ ἀναρίθμητος.

\vs{13}Κατὰ πίστιν ἀπέθανον οὗτοι πάντες, μὴ λαβόντες τὰς ἐπαγγελίας ἀλλὰ πόρρωθεν αὐτὰς ἰδόντες καὶ ἀσπασάμενοι καὶ ὁμολογήσαντες ὅτι ξένοι καὶ παρεπίδημοί εἰσιν ἐπὶ τῆς γῆς.
\vs{14}οἱ γὰρ τοιαῦτα λέγοντες ἐμφανίζουσιν ὅτι πατρίδα ἐπιζητοῦσιν.
\vs{15}καὶ εἰ μὲν ἐκείνης ἐμνημόνευον ἀφ᾽ ἧς ἐξέβησαν, εἶχον ἂν καιρὸν ἀνακάμψαι·
\vs{16}νῦν δὲ κρείττονος ὀρέγονται, τοῦτ᾽ ἔστιν ἐπουρανίου. διὸ οὐκ ἐπαισχύνεται αὐτοὺς ὁ Θεὸς Θεὸς ἐπικαλεῖσθαι αὐτῶν· ἡτοίμασεν γὰρ αὐτοῖς πόλιν.

\vs{17}Πίστει προσενήνοχεν Ἀβραὰμ τὸν Ἰσαὰκ πειραζόμενος καὶ τὸν μονογενῆ προσέφερεν, ὁ τὰς ἐπαγγελίας ἀναδεξάμενος,
\vs{18}πρὸς ὃν ἐλαλήθη ὅτι 
\begin{poetryblock}

\begin{quote}Ἐν Ἰσαὰκ κληθήσεταί σοι σπέρμα,\end{quote}
\end{poetryblock}

\vs{19}λογισάμενος ὅτι καὶ ἐκ νεκρῶν ἐγείρειν δυνατὸς ὁ Θεός, ὅθεν αὐτὸν καὶ ἐν παραβολῇ ἐκομίσατο.
\vs{20}Πίστει καὶ περὶ μελλόντων εὐλόγησεν Ἰσαὰκ τὸν Ἰακὼβ καὶ τὸν Ἠσαῦ.
\vs{21}Πίστει Ἰακὼβ ἀποθνῄσκων ἕκαστον τῶν υἱῶν Ἰωσὴφ εὐλόγησεν καὶ προσεκύνησεν ἐπὶ τὸ ἄκρον τῆς ῥάβδου αὐτοῦ.
\vs{22}Πίστει Ἰωσὴφ τελευτῶν περὶ τῆς ἐξόδου τῶν υἱῶν Ἰσραὴλ ἐμνημόνευσεν καὶ περὶ τῶν ὀστέων αὐτοῦ ἐνετείλατο.

\vs{23}Πίστει Μωϋσῆς γεννηθεὶς ἐκρύβη τρίμηνον ὑπὸ τῶν πατέρων αὐτοῦ, διότι εἶδον ἀστεῖον τὸ παιδίον καὶ οὐκ ἐφοβήθησαν τὸ διάταγμα τοῦ βασιλέως.
\vs{24}Πίστει Μωϋσῆς μέγας γενόμενος ἠρνήσατο λέγεσθαι υἱὸς θυγατρὸς Φαραώ,
\vs{25}μᾶλλον ἑλόμενος συνκακουχεῖσθαι τῷ λαῷ τοῦ Θεοῦ ἢ πρόσκαιρον ἔχειν ἁμαρτίας ἀπόλαυσιν,
\vs{26}μείζονα πλοῦτον ἡγησάμενος τῶν Αἰγύπτου θησαυρῶν τὸν ὀνειδισμὸν τοῦ Χριστοῦ· ἀπέβλεπεν γὰρ εἰς τὴν μισθαποδοσίαν.
\vs{27}Πίστει κατέλιπεν Αἴγυπτον μὴ φοβηθεὶς τὸν θυμὸν τοῦ βασιλέως· τὸν γὰρ ἀόρατον ὡς ὁρῶν ἐκαρτέρησεν.
\vs{28}Πίστει πεποίηκεν τὸ πάσχα καὶ τὴν πρόσχυσιν τοῦ αἵματος, ἵνα μὴ ὁ ὀλοθρεύων τὰ πρωτότοκα θίγῃ αὐτῶν.
\vs{29}Πίστει διέβησαν τὴν Ἐρυθρὰν Θάλασσαν ὡς διὰ ξηρᾶς γῆς, ἧς πεῖραν λαβόντες οἱ Αἰγύπτιοι κατεπόθησαν.
\vs{30}Πίστει τὰ τείχη Ἰεριχὼ ἔπεσαν κυκλωθέντα ἐπὶ ἑπτὰ ἡμέρας.
\vs{31}Πίστει Ῥαὰβ ἡ πόρνη οὐ συναπώλετο τοῖς ἀπειθήσασιν δεξαμένη τοὺς κατασκόπους μετ᾽ εἰρήνης.

\vs{32}Καὶ τί ἔτι λέγω; ἐπιλείψει με γὰρ διηγούμενον ὁ χρόνος περὶ Γεδεών, Βαράκ, Σαμψών, Ἰεφθάε, Δαυίδ τε καὶ Σαμουὴλ καὶ τῶν προφητῶν,
\vs{33}οἳ διὰ πίστεως κατηγωνίσαντο βασιλείας, εἰργάσαντο δικαιοσύνην, ἐπέτυχον ἐπαγγελιῶν, ἔφραξαν στόματα λεόντων,
\vs{34}ἔσβεσαν δύναμιν πυρός, ἔφυγον στόματα μαχαίρης, ἐδυναμώθησαν ἀπὸ ἀσθενείας, ἐγενήθησαν ἰσχυροὶ ἐν πολέμῳ, παρεμβολὰς ἔκλιναν ἀλλοτρίων.
\vs{35}Ἔλαβον γυναῖκες ἐξ ἀναστάσεως τοὺς νεκροὺς αὐτῶν· ἄλλοι δὲ ἐτυμπανίσθησαν οὐ προσδεξάμενοι τὴν ἀπολύτρωσιν, ἵνα κρείττονος ἀναστάσεως τύχωσιν·
\vs{36}ἕτεροι δὲ ἐμπαιγμῶν καὶ μαστίγων πεῖραν ἔλαβον, ἔτι δὲ δεσμῶν καὶ φυλακῆς·
\vs{37}Ἐλιθάσθησαν, ἐπρίσθησαν, ἐν φόνῳ μαχαίρης ἀπέθανον, περιῆλθον ἐν μηλωταῖς, ἐν αἰγείοις δέρμασιν, ὑστερούμενοι, θλιβόμενοι, κακουχούμενοι,
\vs{38}ὧν οὐκ ἦν ἄξιος ὁ κόσμος, ἐπὶ ἐρημίαις πλανώμενοι καὶ ὄρεσιν καὶ σπηλαίοις καὶ ταῖς ὀπαῖς τῆς γῆς.
\vs{39}Καὶ οὗτοι πάντες μαρτυρηθέντες διὰ τῆς πίστεως οὐκ ἐκομίσαντο τὴν ἐπαγγελίαν,
\vs{40}τοῦ Θεοῦ περὶ ἡμῶν κρεῖττόν τι προβλεψαμένου, ἵνα μὴ χωρὶς ἡμῶν τελειωθῶσιν.

\ch{12}
Τοιγαροῦν καὶ ἡμεῖς τοσοῦτον ἔχοντες περικείμενον ἡμῖν νέφος μαρτύρων, ὄγκον ἀποθέμενοι πάντα καὶ τὴν εὐπερίστατον ἁμαρτίαν, δι᾽ ὑπομονῆς τρέχωμεν τὸν προκείμενον ἡμῖν ἀγῶνα
\vs{2}ἀφορῶντες εἰς τὸν τῆς πίστεως ἀρχηγὸν καὶ τελειωτὴν Ἰησοῦν, ὃς ἀντὶ τῆς προκειμένης αὐτῷ χαρᾶς ὑπέμεινεν σταυρὸν αἰσχύνης καταφρονήσας ἐν δεξιᾷ τε τοῦ θρόνου τοῦ Θεοῦ κεκάθικεν.
\vs{3}ἀναλογίσασθε γὰρ τὸν τοιαύτην ὑπομεμενηκότα ὑπὸ τῶν ἁμαρτωλῶν εἰς ἑαυτὸν ἀντιλογίαν, ἵνα μὴ κάμητε ταῖς ψυχαῖς ὑμῶν ἐκλυόμενοι.

\vs{4}Οὔπω μέχρις αἵματος ἀντικατέστητε πρὸς τὴν ἁμαρτίαν ἀνταγωνιζόμενοι.
\vs{5}καὶ ἐκλέλησθε τῆς παρακλήσεως, ἥτις ὑμῖν ὡς υἱοῖς διαλέγεται· 
\begin{poetryblock}

\begin{quote}Υἱέ μου, μὴ ὀλιγώρει παιδείας Κυρίου\end{quote} 

\begin{quote}μηδὲ ἐκλύου ὑπ᾽ αὐτοῦ ἐλεγχόμενος·\end{quote}

\begin{quote} \vs{6}ὃν γὰρ ἀγαπᾷ Κύριος παιδεύει,\end{quote} 

\begin{quote}μαστιγοῖ δὲ πάντα υἱὸν ὃν παραδέχεται.\end{quote}
\end{poetryblock}

\vs{7}Εἰς παιδείαν ὑπομένετε, ὡς υἱοῖς ὑμῖν προσφέρεται ὁ Θεός. τίς γὰρ υἱὸς ὃν οὐ παιδεύει πατήρ;
\vs{8}εἰ δὲ χωρίς ἐστε παιδείας ἧς μέτοχοι γεγόνασιν πάντες, ἄρα νόθοι καὶ οὐχ υἱοί ἐστε.
\vs{9}εἶτα τοὺς μὲν τῆς σαρκὸς ἡμῶν πατέρας εἴχομεν παιδευτὰς καὶ ἐνετρεπόμεθα· οὐ πολὺ δὲ μᾶλλον ὑποταγησόμεθα τῷ Πατρὶ τῶν πνευμάτων καὶ ζήσομεν;
\vs{10}Οἱ μὲν γὰρ πρὸς ὀλίγας ἡμέρας κατὰ τὸ δοκοῦν αὐτοῖς ἐπαίδευον, ὁ δὲ ἐπὶ τὸ συμφέρον εἰς τὸ μεταλαβεῖν τῆς ἁγιότητος αὐτοῦ.
\vs{11}πᾶσα δὲ παιδεία πρὸς μὲν τὸ παρὸν οὐ δοκεῖ χαρᾶς εἶναι ἀλλὰ λύπης, ὕστερον δὲ καρπὸν εἰρηνικὸν τοῖς δι᾽ αὐτῆς γεγυμνασμένοις ἀποδίδωσιν δικαιοσύνης.

\vs{12}Διὸ τὰς παρειμένας χεῖρας καὶ τὰ παραλελυμένα γόνατα ἀνορθώσατε,
\vs{13}καὶ τροχιὰς ὀρθὰς ποιεῖτε τοῖς ποσὶν ὑμῶν, ἵνα μὴ τὸ χωλὸν ἐκτραπῇ, ἰαθῇ δὲ μᾶλλον.
\vs{14}Εἰρήνην διώκετε μετὰ πάντων καὶ τὸν ἁγιασμόν, οὗ χωρὶς οὐδεὶς ὄψεται τὸν Κύριον,
\vs{15}ἐπισκοποῦντες μή τις ὑστερῶν ἀπὸ τῆς χάριτος τοῦ Θεοῦ, μή τις ῥίζα πικρίας ἄνω φύουσα ἐνοχλῇ καὶ δι᾽ αὐτῆς μιανθῶσιν πολλοί,
\vs{16}μή τις πόρνος ἢ βέβηλος ὡς Ἠσαῦ, ὃς ἀντὶ βρώσεως μιᾶς ἀπέδετο τὰ πρωτοτόκια ἑαυτοῦ.
\vs{17}ἴστε γὰρ ὅτι καὶ μετέπειτα θέλων κληρονομῆσαι τὴν εὐλογίαν ἀπεδοκιμάσθη, μετανοίας γὰρ τόπον οὐχ εὗρεν καίπερ μετὰ δακρύων ἐκζητήσας αὐτήν.

\vs{18}Οὐ γὰρ προσεληλύθατε ψηλαφωμένῳ καὶ κεκαυμένῳ πυρὶ καὶ γνόφῳ καὶ ζόφῳ καὶ θυέλλῃ
\vs{19}καὶ σάλπιγγος ἤχῳ καὶ φωνῇ ῥημάτων, ἧς οἱ ἀκούσαντες παρῃτήσαντο μὴ προστεθῆναι αὐτοῖς λόγον,
\vs{20}οὐκ ἔφερον γὰρ τὸ διαστελλόμενον· Κἂν θηρίον θίγῃ τοῦ ὄρους, λιθοβοληθήσεται·
\vs{21}καί, οὕτω φοβερὸν ἦν τὸ φανταζόμενον, Μωϋσῆς εἶπεν· Ἔκφοβός εἰμι καὶ ἔντρομος.
\vs{22}Ἀλλὰ προσεληλύθατε Σιὼν ὄρει καὶ πόλει Θεοῦ ζῶντος, Ἰερουσαλὴμ ἐπουρανίῳ, καὶ μυριάσιν ἀγγέλων, πανηγύρει
\vs{23}καὶ ἐκκλησίᾳ πρωτοτόκων ἀπογεγραμμένων ἐν οὐρανοῖς καὶ Κριτῇ Θεῷ πάντων καὶ πνεύμασι δικαίων τετελειωμένων
\vs{24}καὶ διαθήκης νέας μεσίτῃ Ἰησοῦ καὶ αἵματι ῥαντισμοῦ κρεῖττον λαλοῦντι παρὰ τὸν Ἅβελ.

\vs{25}Βλέπετε μὴ παραιτήσησθε τὸν λαλοῦντα· εἰ γὰρ ἐκεῖνοι οὐκ ἐξέφυγον ἐπὶ γῆς παραιτησάμενοι τὸν χρηματίζοντα, πολὺ μᾶλλον ἡμεῖς οἱ τὸν ἀπ᾽ οὐρανῶν ἀποστρεφόμενοι,
\vs{26}οὗ ἡ φωνὴ τὴν γῆν ἐσάλευσεν τότε, νῦν δὲ ἐπήγγελται λέγων· 
\begin{poetryblock}

\begin{quote}Ἔτι ἅπαξ ἐγὼ σείσω οὐ μόνον τὴν γῆν ἀλλὰ καὶ τὸν οὐρανόν.\end{quote}
\end{poetryblock}

\vs{27}τὸ δὲ Ἔτι ἅπαξ δηλοῖ τὴν τῶν σαλευομένων μετάθεσιν ὡς πεποιημένων, ἵνα μείνῃ τὰ μὴ σαλευόμενα.
\vs{28}Διὸ βασιλείαν ἀσάλευτον παραλαμβάνοντες ἔχωμεν χάριν, δι᾽ ἧς λατρεύωμεν εὐαρέστως τῷ Θεῷ μετὰ εὐλαβείας καὶ δέους·
\vs{29}καὶ γὰρ ὁ Θεὸς ἡμῶν πῦρ καταναλίσκον.

\ch{13}
Ἡ φιλαδελφία μενέτω.
\vs{2}τῆς φιλοξενίας μὴ ἐπιλανθάνεσθε, διὰ ταύτης γὰρ ἔλαθόν τινες ξενίσαντες ἀγγέλους.
\vs{3}μιμνῄσκεσθε τῶν δεσμίων ὡς συνδεδεμένοι, τῶν κακουχουμένων ὡς καὶ αὐτοὶ ὄντες ἐν σώματι.
\vs{4}Τίμιος ὁ γάμος ἐν πᾶσιν καὶ ἡ κοίτη ἀμίαντος, πόρνους γὰρ καὶ μοιχοὺς κρινεῖ ὁ Θεός.
\vs{5}Ἀφιλάργυρος ὁ τρόπος, ἀρκούμενοι τοῖς παροῦσιν. αὐτὸς γὰρ εἴρηκεν· Οὐ μή σε ἀνῶ οὐδ᾽ οὐ μή σε ἐγκαταλίπω,
\vs{6}Ὥστε θαρροῦντας ἡμᾶς λέγειν· 
\begin{poetryblock}

\begin{quote}Κύριος ἐμοὶ βοηθός, καὶ οὐ φοβηθήσομαι,\end{quote} 

\begin{quote}τί ποιήσει μοι ἄνθρωπος;\end{quote}
\end{poetryblock}

\vs{7}Μνημονεύετε τῶν ἡγουμένων ὑμῶν, οἵτινες ἐλάλησαν ὑμῖν τὸν λόγον τοῦ Θεοῦ, ὧν ἀναθεωροῦντες τὴν ἔκβασιν τῆς ἀναστροφῆς μιμεῖσθε τὴν πίστιν.
\vs{8}Ἰησοῦς Χριστὸς ἐχθὲς καὶ σήμερον ὁ αὐτός καὶ εἰς τοὺς αἰῶνας.
\vs{9}Διδαχαῖς ποικίλαις καὶ ξέναις μὴ παραφέρεσθε· καλὸν γὰρ χάριτι βεβαιοῦσθαι τὴν καρδίαν, οὐ βρώμασιν ἐν οἷς οὐκ ὠφελήθησαν οἱ περιπατοῦντες.
\vs{10}Ἔχομεν θυσιαστήριον ἐξ οὗ φαγεῖν οὐκ ἔχουσιν ἐξουσίαν οἱ τῇ σκηνῇ λατρεύοντες.
\vs{11}Ὧν γὰρ εἰσφέρεται ζῴων τὸ αἷμα περὶ ἁμαρτίας εἰς τὰ ἅγια διὰ τοῦ ἀρχιερέως, τούτων τὰ σώματα κατακαίεται ἔξω τῆς παρεμβολῆς.
\vs{12}διὸ καὶ Ἰησοῦς, ἵνα ἁγιάσῃ διὰ τοῦ ἰδίου αἵματος τὸν λαόν, ἔξω τῆς πύλης ἔπαθεν.
\vs{13}τοίνυν ἐξερχώμεθα πρὸς αὐτὸν ἔξω τῆς παρεμβολῆς τὸν ὀνειδισμὸν αὐτοῦ φέροντες·
\vs{14}οὐ γὰρ ἔχομεν ὧδε μένουσαν πόλιν ἀλλὰ τὴν μέλλουσαν ἐπιζητοῦμεν.
\vs{15}Δι᾽ αὐτοῦ οὖν ἀναφέρωμεν θυσίαν αἰνέσεως διὰ παντὸς τῷ Θεῷ, τοῦτ᾽ ἔστιν καρπὸν χειλέων ὁμολογούντων τῷ ὀνόματι αὐτοῦ.
\vs{16}τῆς δὲ εὐποιΐας καὶ κοινωνίας μὴ ἐπιλανθάνεσθε· τοιαύταις γὰρ θυσίαις εὐαρεστεῖται ὁ Θεός.
\vs{17}Πείθεσθε τοῖς ἡγουμένοις ὑμῶν καὶ ὑπείκετε, αὐτοὶ γὰρ ἀγρυπνοῦσιν ὑπὲρ τῶν ψυχῶν ὑμῶν ὡς λόγον ἀποδώσοντες, ἵνα μετὰ χαρᾶς τοῦτο ποιῶσιν καὶ μὴ στενάζοντες· ἀλυσιτελὲς γὰρ ὑμῖν τοῦτο.

\vs{18}Προσεύχεσθε περὶ ἡμῶν· πειθόμεθα γὰρ ὅτι καλὴν συνείδησιν ἔχομεν, ἐν πᾶσιν καλῶς θέλοντες ἀναστρέφεσθαι.
\vs{19}περισσοτέρως δὲ παρακαλῶ τοῦτο ποιῆσαι, ἵνα τάχιον ἀποκατασταθῶ ὑμῖν.

\vs{20}Ὁ δὲ Θεὸς τῆς εἰρήνης, ὁ ἀναγαγὼν ἐκ νεκρῶν τὸν ποιμένα τῶν προβάτων τὸν μέγαν ἐν αἵματι διαθήκης αἰωνίου, τὸν Κύριον ἡμῶν Ἰησοῦν,
\vs{21}καταρτίσαι ὑμᾶς ἐν παντὶ ἀγαθῷ εἰς τὸ ποιῆσαι τὸ θέλημα αὐτοῦ, ποιῶν ἐν ἡμῖν τὸ εὐάρεστον ἐνώπιον αὐτοῦ διὰ Ἰησοῦ Χριστοῦ, ᾧ ἡ δόξα εἰς τοὺς αἰῶνας τῶν αἰώνων, ἀμήν.

\vs{22}Παρακαλῶ δὲ ὑμᾶς, ἀδελφοί, ἀνέχεσθε τοῦ λόγου τῆς παρακλήσεως, καὶ γὰρ διὰ βραχέων ἐπέστειλα ὑμῖν.
\vs{23}Γινώσκετε τὸν ἀδελφὸν ἡμῶν Τιμόθεον ἀπολελυμένον, μεθ᾽ οὗ ἐὰν τάχιον ἔρχηται ὄψομαι ὑμᾶς.

\vs{24}Ἀσπάσασθε πάντας τοὺς ἡγουμένους ὑμῶν καὶ πάντας τοὺς ἁγίους. Ἀσπάζονται ὑμᾶς οἱ ἀπὸ τῆς Ἰταλίας.

\vs{25}Ἡ χάρις μετὰ πάντων ὑμῶν.


\def\book{ΙΑΚΩΒΟΥ}
\biblebook{ΙΑΚΩΒΟΥ}


\lettrine[lines=2, loversize=0.2, nindent=0em, findent=.25em]{\textcolor{bookheadingcolor}{Ἰ}}{άκωβος} Θεοῦ καὶ Κυρίου Ἰησοῦ Χριστοῦ δοῦλος Ταῖς δώδεκα φυλαῖς ταῖς ἐν τῇ Διασπορᾷ Χαίρειν.

\vs{2}Πᾶσαν χαρὰν ἡγήσασθε, ἀδελφοί μου, ὅταν πειρασμοῖς περιπέσητε ποικίλοις,
\vs{3}γινώσκοντες ὅτι τὸ δοκίμιον ὑμῶν τῆς πίστεως κατεργάζεται ὑπομονήν.
\vs{4}ἡ δὲ ὑπομονὴ ἔργον τέλειον ἐχέτω, ἵνα ἦτε τέλειοι καὶ ὁλόκληροι ἐν μηδενὶ λειπόμενοι.

\vs{5}Εἰ δέ τις ὑμῶν λείπεται σοφίας, αἰτείτω παρὰ τοῦ διδόντος Θεοῦ πᾶσιν ἁπλῶς καὶ μὴ ὀνειδίζοντος, καὶ δοθήσεται αὐτῷ.
\vs{6}αἰτείτω δὲ ἐν πίστει μηδὲν διακρινόμενος· ὁ γὰρ διακρινόμενος ἔοικεν κλύδωνι θαλάσσης ἀνεμιζομένῳ καὶ ῥιπιζομένῳ.
\vs{7}μὴ γὰρ οἰέσθω ὁ ἄνθρωπος ἐκεῖνος ὅτι λήμψεταί τι παρὰ τοῦ Κυρίου,
\vs{8}ἀνὴρ δίψυχος, ἀκατάστατος ἐν πάσαις ταῖς ὁδοῖς αὐτοῦ.

\vs{9}Καυχάσθω δὲ ὁ ἀδελφὸς ὁ ταπεινὸς ἐν τῷ ὕψει αὐτοῦ,
\vs{10}ὁ δὲ πλούσιος ἐν τῇ ταπεινώσει αὐτοῦ, ὅτι ὡς ἄνθος χόρτου παρελεύσεται.
\vs{11}ἀνέτειλεν γὰρ ὁ ἥλιος σὺν τῷ καύσωνι καὶ ἐξήρανεν τὸν χόρτον, καὶ τὸ ἄνθος αὐτοῦ ἐξέπεσεν, καὶ ἡ εὐπρέπεια τοῦ προσώπου αὐτοῦ ἀπώλετο· οὕτως καὶ ὁ πλούσιος ἐν ταῖς πορείαις αὐτοῦ μαρανθήσεται.

\vs{12}Μακάριος ἀνὴρ ὃς ὑπομένει πειρασμόν, ὅτι δόκιμος γενόμενος λήμψεται τὸν στέφανον τῆς ζωῆς ὃν ἐπηγγείλατο τοῖς ἀγαπῶσιν αὐτόν.
\vs{13}Μηδεὶς πειραζόμενος λεγέτω ὅτι Ἀπὸ Θεοῦ πειράζομαι· ὁ γὰρ Θεὸς ἀπείραστός ἐστιν κακῶν, πειράζει δὲ αὐτὸς οὐδένα.
\vs{14}ἕκαστος δὲ πειράζεται ὑπὸ τῆς ἰδίας ἐπιθυμίας ἐξελκόμενος καὶ δελεαζόμενος·
\vs{15}εἶτα ἡ ἐπιθυμία συλλαβοῦσα τίκτει ἁμαρτίαν, ἡ δὲ ἁμαρτία ἀποτελεσθεῖσα ἀποκύει θάνατον.

\vs{16}Μὴ πλανᾶσθε, ἀδελφοί μου ἀγαπητοί.
\vs{17}πᾶσα δόσις ἀγαθὴ καὶ πᾶν δώρημα τέλειον ἄνωθέν ἐστιν καταβαῖνον ἀπὸ τοῦ Πατρὸς τῶν φώτων, παρ᾽ ᾧ οὐκ ἔνι παραλλαγὴ ἢ τροπῆς ἀποσκίασμα.
\vs{18}βουληθεὶς ἀπεκύησεν ἡμᾶς λόγῳ ἀληθείας εἰς τὸ εἶναι ἡμᾶς ἀπαρχήν τινα τῶν αὐτοῦ κτισμάτων.

\vs{19}Ἴστε, ἀδελφοί μου ἀγαπητοί· ἔστω δὲ πᾶς ἄνθρωπος ταχὺς εἰς τὸ ἀκοῦσαι, βραδὺς εἰς τὸ λαλῆσαι, βραδὺς εἰς ὀργήν·
\vs{20}ὀργὴ γὰρ ἀνδρὸς δικαιοσύνην Θεοῦ οὐκ ἐργάζεται.
\vs{21}διὸ ἀποθέμενοι πᾶσαν ῥυπαρίαν καὶ περισσείαν κακίας ἐν πραΰτητι δέξασθε τὸν ἔμφυτον λόγον τὸν δυνάμενον σῶσαι τὰς ψυχὰς ὑμῶν.

\vs{22}Γίνεσθε δὲ ποιηταὶ λόγου καὶ μὴ μόνον ἀκροαταὶ παραλογιζόμενοι ἑαυτούς.
\vs{23}ὅτι εἴ τις ἀκροατὴς λόγου ἐστὶν καὶ οὐ ποιητής, οὗτος ἔοικεν ἀνδρὶ κατανοοῦντι τὸ πρόσωπον τῆς γενέσεως αὐτοῦ ἐν ἐσόπτρῳ·
\vs{24}κατενόησεν γὰρ ἑαυτὸν καὶ ἀπελήλυθεν καὶ εὐθέως ἐπελάθετο ὁποῖος ἦν.
\vs{25}ὁ δὲ παρακύψας εἰς νόμον τέλειον τὸν τῆς ἐλευθερίας καὶ παραμείνας οὐκ ἀκροατὴς ἐπιλησμονῆς γενόμενος ἀλλὰ ποιητὴς ἔργου, οὗτος μακάριος ἐν τῇ ποιήσει αὐτοῦ ἔσται.

\vs{26}Εἴ τις δοκεῖ θρησκὸς εἶναι μὴ χαλιναγωγῶν γλῶσσαν αὐτοῦ ἀλλὰ ἀπατῶν καρδίαν αὐτοῦ, τούτου μάταιος ἡ θρησκεία.
\vs{27}θρησκεία καθαρὰ καὶ ἀμίαντος παρὰ τῷ Θεῷ καὶ Πατρὶ αὕτη ἐστίν, ἐπισκέπτεσθαι ὀρφανοὺς καὶ χήρας ἐν τῇ θλίψει αὐτῶν, ἄσπιλον ἑαυτὸν τηρεῖν ἀπὸ τοῦ κόσμου.

\ch{2}
Ἀδελφοί μου, μὴ ἐν προσωπολημψίαις ἔχετε τὴν πίστιν τοῦ Κυρίου ἡμῶν Ἰησοῦ Χριστοῦ τῆς δόξης.
\vs{2}Ἐὰν γὰρ εἰσέλθῃ εἰς συναγωγὴν ὑμῶν ἀνὴρ χρυσοδακτύλιος ἐν ἐσθῆτι λαμπρᾷ, εἰσέλθῃ δὲ καὶ πτωχὸς ἐν ῥυπαρᾷ ἐσθῆτι,
\vs{3}ἐπιβλέψητε δὲ ἐπὶ τὸν φοροῦντα τὴν ἐσθῆτα τὴν λαμπρὰν καὶ εἴπητε· Σὺ κάθου ὧδε καλῶς, καὶ τῷ πτωχῷ εἴπητε· Σὺ στῆθι ἢ Κάθου ἐκεῖ ὑπὸ τὸ ὑποπόδιόν μου,
\vs{4}καὶ οὐ διεκρίθητε ἐν ἑαυτοῖς καὶ ἐγένεσθε κριταὶ διαλογισμῶν πονηρῶν;
\vs{5}Ἀκούσατε, ἀδελφοί μου ἀγαπητοί· οὐχ ὁ Θεὸς ἐξελέξατο τοὺς πτωχοὺς τῷ κόσμῳ πλουσίους ἐν πίστει καὶ κληρονόμους τῆς βασιλείας ἧς ἐπηγγείλατο τοῖς ἀγαπῶσιν αὐτόν;
\vs{6}ὑμεῖς δὲ ἠτιμάσατε τὸν πτωχόν. οὐχ οἱ πλούσιοι καταδυναστεύουσιν ὑμῶν καὶ αὐτοὶ ἕλκουσιν ὑμᾶς εἰς κριτήρια;
\vs{7}οὐκ αὐτοὶ βλασφημοῦσιν τὸ καλὸν ὄνομα τὸ ἐπικληθὲν ἐφ᾽ ὑμᾶς;

\vs{8}Εἰ μέντοι νόμον τελεῖτε βασιλικὸν κατὰ τὴν γραφήν· Ἀγαπήσεις τὸν πλησίον σου ὡς σεαυτόν, καλῶς ποιεῖτε·
\vs{9}εἰ δὲ προσωπολημπτεῖτε, ἁμαρτίαν ἐργάζεσθε ἐλεγχόμενοι ὑπὸ τοῦ νόμου ὡς παραβάται.
\vs{10}Ὅστις γὰρ ὅλον τὸν νόμον τηρήσῃ, πταίσῃ δὲ ἐν ἑνί, γέγονεν πάντων ἔνοχος.
\vs{11}ὁ γὰρ εἰπών· Μὴ μοιχεύσῃς, εἶπεν καί· Μὴ φονεύσῃς· εἰ δὲ οὐ μοιχεύεις, φονεύεις δέ, γέγονας παραβάτης νόμου.

\vs{12}Οὕτως λαλεῖτε καὶ οὕτως ποιεῖτε ὡς διὰ νόμου ἐλευθερίας μέλλοντες κρίνεσθαι.
\vs{13}ἡ γὰρ κρίσις ἀνέλεος τῷ μὴ ποιήσαντι ἔλεος· κατακαυχᾶται ἔλεος κρίσεως.

\vs{14}Τί τὸ ὄφελος, ἀδελφοί μου, ἐὰν πίστιν λέγῃ τις ἔχειν, ἔργα δὲ μὴ ἔχῃ; μὴ δύναται ἡ πίστις σῶσαι αὐτόν;
\vs{15}ἐὰν ἀδελφὸς ἢ ἀδελφὴ γυμνοὶ ὑπάρχωσιν καὶ λειπόμενοι ὦσιν τῆς ἐφημέρου τροφῆς,
\vs{16}εἴπῃ δέ τις αὐτοῖς ἐξ ὑμῶν· Ὑπάγετε ἐν εἰρήνῃ, θερμαίνεσθε καὶ χορτάζεσθε, μὴ δῶτε δὲ αὐτοῖς τὰ ἐπιτήδεια τοῦ σώματος, τί τὸ ὄφελος;
\vs{17}οὕτως καὶ ἡ πίστις, ἐὰν μὴ ἔχῃ ἔργα, νεκρά ἐστιν καθ᾽ ἑαυτήν.

\vs{18}Ἀλλ᾽ ἐρεῖ τις· Σὺ πίστιν ἔχεις, κἀγὼ ἔργα ἔχω. δεῖξόν μοι τὴν πίστιν σου χωρὶς τῶν ἔργων, κἀγώ σοι δείξω ἐκ τῶν ἔργων μου τὴν πίστιν.
\vs{19}σὺ πιστεύεις ὅτι εἷς ἐστιν ὁ Θεός, καλῶς ποιεῖς· καὶ τὰ δαιμόνια πιστεύουσιν καὶ φρίσσουσιν.

\vs{20}Θέλεις δὲ γνῶναι, ὦ ἄνθρωπε κενέ, ὅτι ἡ πίστις χωρὶς τῶν ἔργων ἀργή ἐστιν;
\vs{21}Ἀβραὰμ ὁ πατὴρ ἡμῶν οὐκ ἐξ ἔργων ἐδικαιώθη ἀνενέγκας Ἰσαὰκ τὸν υἱὸν αὐτοῦ ἐπὶ τὸ θυσιαστήριον;
\vs{22}βλέπεις ὅτι ἡ πίστις συνήργει τοῖς ἔργοις αὐτοῦ καὶ ἐκ τῶν ἔργων ἡ πίστις ἐτελειώθη,
\vs{23}καὶ ἐπληρώθη ἡ γραφὴ ἡ λέγουσα· Ἐπίστευσεν δὲ Ἀβραὰμ τῷ Θεῷ, καὶ ἐλογίσθη αὐτῷ εἰς δικαιοσύνην καὶ φίλος Θεοῦ ἐκλήθη.
\vs{24}ὁρᾶτε ὅτι ἐξ ἔργων δικαιοῦται ἄνθρωπος καὶ οὐκ ἐκ πίστεως μόνον.
\vs{25}Ὁμοίως δὲ καὶ Ῥαὰβ ἡ πόρνη οὐκ ἐξ ἔργων ἐδικαιώθη ὑποδεξαμένη τοὺς ἀγγέλους καὶ ἑτέρᾳ ὁδῷ ἐκβαλοῦσα;
\vs{26}ὥσπερ γὰρ τὸ σῶμα χωρὶς πνεύματος νεκρόν ἐστιν, οὕτως καὶ ἡ πίστις χωρὶς ἔργων νεκρά ἐστιν.

\ch{3}
Μὴ πολλοὶ διδάσκαλοι γίνεσθε, ἀδελφοί μου, εἰδότες ὅτι μεῖζον κρίμα λημψόμεθα.
\vs{2}πολλὰ γὰρ πταίομεν ἅπαντες. εἴ τις ἐν λόγῳ οὐ πταίει, οὗτος τέλειος ἀνήρ δυνατὸς χαλιναγωγῆσαι καὶ ὅλον τὸ σῶμα.
\vs{3}Εἰ δὲ τῶν ἵππων τοὺς χαλινοὺς εἰς τὰ στόματα βάλλομεν εἰς τὸ πείθεσθαι αὐτοὺς ἡμῖν, καὶ ὅλον τὸ σῶμα αὐτῶν μετάγομεν.
\vs{4}ἰδοὺ καὶ τὰ πλοῖα τηλικαῦτα ὄντα καὶ ὑπὸ ἀνέμων σκληρῶν ἐλαυνόμενα μετάγεται ὑπὸ ἐλαχίστου πηδαλίου ὅπου ἡ ὁρμὴ τοῦ εὐθύνοντος βούλεται.
\vs{5}Οὕτως καὶ ἡ γλῶσσα μικρὸν μέλος ἐστὶν καὶ μεγάλα αὐχεῖ. ἰδοὺ ἡλίκον πῦρ ἡλίκην ὕλην ἀνάπτει.
\vs{6}καὶ ἡ γλῶσσα πῦρ. ὁ κόσμος τῆς ἀδικίας ἡ γλῶσσα καθίσταται ἐν τοῖς μέλεσιν ἡμῶν ἡ σπιλοῦσα ὅλον τὸ σῶμα καὶ φλογίζουσα τὸν τροχὸν τῆς γενέσεως καὶ φλογιζομένη ὑπὸ τῆς γεέννης.
\vs{7}Πᾶσα γὰρ φύσις θηρίων τε καὶ πετεινῶν, ἑρπετῶν τε καὶ ἐναλίων δαμάζεται καὶ δεδάμασται τῇ φύσει τῇ ἀνθρωπίνῃ,
\vs{8}τὴν δὲ γλῶσσαν οὐδεὶς δαμάσαι δύναται ἀνθρώπων, ἀκατάστατον κακόν, μεστὴ ἰοῦ θανατηφόρου.
\vs{9}Ἐν αὐτῇ εὐλογοῦμεν τὸν Κύριον καὶ Πατέρα καὶ ἐν αὐτῇ καταρώμεθα τοὺς ἀνθρώπους τοὺς καθ᾽ ὁμοίωσιν Θεοῦ γεγονότας·
\vs{10}ἐκ τοῦ αὐτοῦ στόματος ἐξέρχεται εὐλογία καὶ κατάρα. οὐ χρή, ἀδελφοί μου, ταῦτα οὕτως γίνεσθαι.
\vs{11}μήτι ἡ πηγὴ ἐκ τῆς αὐτῆς ὀπῆς βρύει τὸ γλυκὺ καὶ τὸ πικρόν;
\vs{12}μὴ δύναται, ἀδελφοί μου, συκῆ ἐλαίας ποιῆσαι ἢ ἄμπελος σῦκα; οὔτε ἁλυκὸν γλυκὺ ποιῆσαι ὕδωρ.

\vs{13}Τίς σοφὸς καὶ ἐπιστήμων ἐν ὑμῖν; δειξάτω ἐκ τῆς καλῆς ἀναστροφῆς τὰ ἔργα αὐτοῦ ἐν πραΰτητι σοφίας.
\vs{14}εἰ δὲ ζῆλον πικρὸν ἔχετε καὶ ἐριθείαν ἐν τῇ καρδίᾳ ὑμῶν, μὴ κατακαυχᾶσθε καὶ ψεύδεσθε κατὰ τῆς ἀληθείας.
\vs{15}οὐκ ἔστιν αὕτη ἡ σοφία ἄνωθεν κατερχομένη ἀλλὰ ἐπίγειος, ψυχική, δαιμονιώδης.
\vs{16}ὅπου γὰρ ζῆλος καὶ ἐριθεία, ἐκεῖ ἀκαταστασία καὶ πᾶν φαῦλον πρᾶγμα.
\vs{17}Ἡ δὲ ἄνωθεν σοφία πρῶτον μὲν ἁγνή ἐστιν, ἔπειτα εἰρηνική, ἐπιεικής, εὐπειθής, μεστὴ ἐλέους καὶ καρπῶν ἀγαθῶν, ἀδιάκριτος, ἀνυπόκριτος.
\vs{18}καρπὸς δὲ δικαιοσύνης ἐν εἰρήνῃ σπείρεται τοῖς ποιοῦσιν εἰρήνην.

\ch{4}
Πόθεν πόλεμοι καὶ πόθεν μάχαι ἐν ὑμῖν; οὐκ ἐντεῦθεν, ἐκ τῶν ἡδονῶν ὑμῶν τῶν στρατευομένων ἐν τοῖς μέλεσιν ὑμῶν;
\vs{2}ἐπιθυμεῖτε καὶ οὐκ ἔχετε, φονεύετε καὶ ζηλοῦτε καὶ οὐ δύνασθε ἐπιτυχεῖν, μάχεσθε καὶ πολεμεῖτε, οὐκ ἔχετε διὰ τὸ μὴ αἰτεῖσθαι ὑμᾶς,
\vs{3}αἰτεῖτε καὶ οὐ λαμβάνετε, διότι κακῶς αἰτεῖσθε, ἵνα ἐν ταῖς ἡδοναῖς ὑμῶν δαπανήσητε.
\vs{4}Μοιχαλίδες, οὐκ οἴδατε ὅτι ἡ φιλία τοῦ κόσμου ἔχθρα τοῦ Θεοῦ ἐστιν; ὃς ἐὰν οὖν βουληθῇ φίλος εἶναι τοῦ κόσμου, ἐχθρὸς τοῦ Θεοῦ καθίσταται.
\vs{5}ἢ δοκεῖτε ὅτι κενῶς ἡ γραφὴ λέγει· Πρὸς φθόνον ἐπιποθεῖ τὸ πνεῦμα ὃ κατῴκισεν ἐν ἡμῖν,
\vs{6}μείζονα δὲ δίδωσιν χάριν; διὸ λέγει· 
\begin{poetryblock}

\begin{quote}Ὁ Θεὸς ὑπερηφάνοις ἀντιτάσσεται,\end{quote} 

\begin{quote}ταπεινοῖς δὲ δίδωσιν χάριν.\end{quote}
\end{poetryblock}

\vs{7}Ὑποτάγητε οὖν τῷ Θεῷ, ἀντίστητε δὲ τῷ διαβόλῳ, καὶ φεύξεται ἀφ᾽ ὑμῶν·
\vs{8}ἐγγίσατε τῷ Θεῷ καὶ ἐγγιεῖ ὑμῖν. καθαρίσατε χεῖρας, ἁμαρτωλοί, καὶ ἁγνίσατε καρδίας, δίψυχοι.
\vs{9}ταλαιπωρήσατε καὶ πενθήσατε καὶ κλαύσατε. ὁ γέλως ὑμῶν εἰς πένθος μετατραπήτω καὶ ἡ χαρὰ εἰς κατήφειαν.
\vs{10}ταπεινώθητε ἐνώπιον τοῦ Κυρίου καὶ ὑψώσει ὑμᾶς.

\vs{11}Μὴ καταλαλεῖτε ἀλλήλων, ἀδελφοί. ὁ καταλαλῶν ἀδελφοῦ ἢ κρίνων τὸν ἀδελφὸν αὐτοῦ καταλαλεῖ νόμου καὶ κρίνει νόμον· εἰ δὲ νόμον κρίνεις, οὐκ εἶ ποιητὴς νόμου ἀλλὰ κριτής.
\vs{12}εἷς ἐστιν ὁ νομοθέτης καὶ κριτής ὁ δυνάμενος σῶσαι καὶ ἀπολέσαι· σὺ δὲ τίς εἶ ὁ κρίνων τὸν πλησίον;

\vs{13}Ἄγε νῦν οἱ λέγοντες· Σήμερον ἢ αὔριον πορευσόμεθα εἰς τήνδε τὴν πόλιν καὶ ποιήσομεν ἐκεῖ ἐνιαυτὸν καὶ ἐμπορευσόμεθα καὶ κερδήσομεν,
\vs{14}οἵτινες οὐκ ἐπίστασθε τὸ τῆς αὔριον ποία ἡ ζωὴ ὑμῶν— ἀτμὶς γάρ ἐστε ἡ πρὸς ὀλίγον φαινομένη, ἔπειτα καὶ ἀφανιζομένη—
\vs{15}Ἀντὶ τοῦ λέγειν ὑμᾶς· Ἐὰν ὁ Κύριος θελήσῃ καὶ ζήσομεν καὶ ποιήσομεν τοῦτο ἢ ἐκεῖνο.
\vs{16}νῦν δὲ καυχᾶσθε ἐν ταῖς ἀλαζονείαις ὑμῶν· πᾶσα καύχησις τοιαύτη πονηρά ἐστιν.
\vs{17}εἰδότι οὖν καλὸν ποιεῖν καὶ μὴ ποιοῦντι, ἁμαρτία αὐτῷ ἐστιν.

\ch{5}
Ἄγε νῦν οἱ πλούσιοι, κλαύσατε ὀλολύζοντες ἐπὶ ταῖς ταλαιπωρίαις ὑμῶν ταῖς ἐπερχομέναις.
\vs{2}ὁ πλοῦτος ὑμῶν σέσηπεν καὶ τὰ ἱμάτια ὑμῶν σητόβρωτα γέγονεν,
\vs{3}ὁ χρυσὸς ὑμῶν καὶ ὁ ἄργυρος κατίωται καὶ ὁ ἰὸς αὐτῶν εἰς μαρτύριον ὑμῖν ἔσται καὶ φάγεται τὰς σάρκας ὑμῶν ὡς πῦρ. Ἐθησαυρίσατε ἐν ἐσχάταις ἡμέραις.
\vs{4}ἰδοὺ ὁ μισθὸς τῶν ἐργατῶν τῶν ἀμησάντων τὰς χώρας ὑμῶν ὁ ἀφυστερημένος ἀφ᾽ ὑμῶν κράζει, καὶ αἱ βοαὶ τῶν θερισάντων εἰς τὰ ὦτα Κυρίου Σαβαὼθ εἰσεληλύθασιν.
\vs{5}Ἐτρυφήσατε ἐπὶ τῆς γῆς καὶ ἐσπαταλήσατε, ἐθρέψατε τὰς καρδίας ὑμῶν ἐν ἡμέρᾳ σφαγῆς,
\vs{6}κατεδικάσατε, ἐφονεύσατε τὸν δίκαιον· οὐκ ἀντιτάσσεται ὑμῖν.

\vs{7}Μακροθυμήσατε οὖν, ἀδελφοί, ἕως τῆς παρουσίας τοῦ Κυρίου. ἰδοὺ ὁ γεωργὸς ἐκδέχεται τὸν τίμιον καρπὸν τῆς γῆς μακροθυμῶν ἐπ᾽ αὐτῷ, ἕως λάβῃ πρόϊμον καὶ ὄψιμον.
\vs{8}μακροθυμήσατε καὶ ὑμεῖς, στηρίξατε τὰς καρδίας ὑμῶν, ὅτι ἡ παρουσία τοῦ Κυρίου ἤγγικεν.
\vs{9}μὴ στενάζετε, ἀδελφοί, κατ᾽ ἀλλήλων, ἵνα μὴ κριθῆτε· ἰδοὺ ὁ κριτὴς πρὸ τῶν θυρῶν ἕστηκεν.
\vs{10}Ὑπόδειγμα λάβετε, ἀδελφοί, τῆς κακοπαθίας καὶ τῆς μακροθυμίας τοὺς προφήτας οἳ ἐλάλησαν ἐν τῷ ὀνόματι Κυρίου.
\vs{11}ἰδοὺ μακαρίζομεν τοὺς ὑπομείναντας· τὴν ὑπομονὴν Ἰὼβ ἠκούσατε καὶ τὸ τέλος Κυρίου εἴδετε, ὅτι πολύσπλαγχνός ἐστιν ὁ Κύριος καὶ οἰκτίρμων.

\vs{12}Πρὸ πάντων δέ, ἀδελφοί μου, μὴ ὀμνύετε μήτε τὸν οὐρανὸν μήτε τὴν γῆν μήτε ἄλλον τινὰ ὅρκον· ἤτω δὲ ὑμῶν τὸ Ναὶ ναί καὶ τὸ Οὒ οὔ, ἵνα μὴ ὑπὸ κρίσιν πέσητε.

\vs{13}Κακοπαθεῖ τις ἐν ὑμῖν, προσευχέσθω· εὐθυμεῖ τις, ψαλλέτω·
\vs{14}ἀσθενεῖ τις ἐν ὑμῖν, προσκαλεσάσθω τοὺς πρεσβυτέρους τῆς ἐκκλησίας καὶ προσευξάσθωσαν ἐπ᾽ αὐτὸν ἀλείψαντες αὐτὸν ἐλαίῳ ἐν τῷ ὀνόματι τοῦ Κυρίου.
\vs{15}καὶ ἡ εὐχὴ τῆς πίστεως σώσει τὸν κάμνοντα καὶ ἐγερεῖ αὐτὸν ὁ Κύριος· κἂν ἁμαρτίας ᾖ πεποιηκώς, ἀφεθήσεται αὐτῷ.
\vs{16}Ἐξομολογεῖσθε οὖν ἀλλήλοις τὰς ἁμαρτίας καὶ εὔχεσθε ὑπὲρ ἀλλήλων, ὅπως ἰαθῆτε. πολὺ ἰσχύει δέησις δικαίου ἐνεργουμένη.
\vs{17}Ἠλίας ἄνθρωπος ἦν ὁμοιοπαθὴς ἡμῖν καὶ προσευχῇ προσηύξατο τοῦ μὴ βρέξαι, καὶ οὐκ ἔβρεξεν ἐπὶ τῆς γῆς ἐνιαυτοὺς τρεῖς καὶ μῆνας ἕξ·
\vs{18}καὶ πάλιν προσηύξατο, καὶ ὁ οὐρανὸς ὑετὸν ἔδωκεν καὶ ἡ γῆ ἐβλάστησεν τὸν καρπὸν αὐτῆς.

\vs{19}Ἀδελφοί μου, ἐάν τις ἐν ὑμῖν πλανηθῇ ἀπὸ τῆς ἀληθείας καὶ ἐπιστρέψῃ τις αὐτόν,
\vs{20}γινωσκέτω ὅτι ὁ ἐπιστρέψας ἁμαρτωλὸν ἐκ πλάνης ὁδοῦ αὐτοῦ σώσει ψυχὴν αὐτοῦ ἐκ θανάτου καὶ καλύψει πλῆθος ἁμαρτιῶν.


\def\book{ΠΕΤΡΟΥ Α}
\biblebook{ΠΕΤΡΟΥ Α}


\lettrine[lines=2, loversize=0.2, nindent=0em, findent=.25em]{\textcolor{bookheadingcolor}{Π}}{έτρος} ἀπόστολος Ἰησοῦ Χριστοῦ Ἐκλεκτοῖς παρεπιδήμοις Διασπορᾶς Πόντου, Γαλατίας, Καππαδοκίας, Ἀσίας καὶ Βιθυνίας
\vs{2}κατὰ πρόγνωσιν Θεοῦ Πατρός ἐν ἁγιασμῷ Πνεύματος εἰς ὑπακοὴν καὶ ῥαντισμὸν αἵματος Ἰησοῦ Χριστοῦ, Χάρις ὑμῖν καὶ εἰρήνη πληθυνθείη.

\vs{3}Εὐλογητὸς ὁ Θεὸς καὶ Πατὴρ τοῦ Κυρίου ἡμῶν Ἰησοῦ Χριστοῦ ὁ κατὰ τὸ πολὺ αὐτοῦ ἔλεος ἀναγεννήσας ἡμᾶς εἰς ἐλπίδα ζῶσαν δι᾽ ἀναστάσεως Ἰησοῦ Χριστοῦ ἐκ νεκρῶν,
\vs{4}εἰς κληρονομίαν ἄφθαρτον καὶ ἀμίαντον καὶ ἀμάραντον τετηρημένην ἐν οὐρανοῖς εἰς ὑμᾶς
\vs{5}τοὺς ἐν δυνάμει Θεοῦ φρουρουμένους διὰ πίστεως εἰς σωτηρίαν ἑτοίμην ἀποκαλυφθῆναι ἐν καιρῷ ἐσχάτῳ
\vs{6}ἐν ᾧ ἀγαλλιᾶσθε ὀλίγον ἄρτι, εἰ δέον ἐστὶν, λυπηθέντες ἐν ποικίλοις πειρασμοῖς,
\vs{7}ἵνα τὸ δοκίμιον ὑμῶν τῆς πίστεως πολυτιμότερον χρυσίου τοῦ ἀπολλυμένου, διὰ πυρὸς δὲ δοκιμαζομένου εὑρεθῇ εἰς ἔπαινον καὶ δόξαν καὶ τιμὴν ἐν ἀποκαλύψει Ἰησοῦ Χριστοῦ
\vs{8}ὃν οὐκ ἰδόντες ἀγαπᾶτε, εἰς ὃν ἄρτι μὴ ὁρῶντες, πιστεύοντες δὲ ἀγαλλιᾶσθε χαρᾷ ἀνεκλαλήτῳ καὶ δεδοξασμένῃ
\vs{9}κομιζόμενοι τὸ τέλος τῆς πίστεως ὑμῶν σωτηρίαν ψυχῶν.
\vs{10}Περὶ ἧς σωτηρίας ἐξεζήτησαν καὶ ἐξηραύνησαν προφῆται οἱ περὶ τῆς εἰς ὑμᾶς χάριτος προφητεύσαντες
\vs{11}ἐραυνῶντες εἰς τίνα ἢ ποῖον καιρὸν ἐδήλου τὸ ἐν αὐτοῖς Πνεῦμα Χριστοῦ προμαρτυρόμενον τὰ εἰς Χριστὸν παθήματα καὶ τὰς μετὰ ταῦτα δόξας.
\vs{12}οἷς ἀπεκαλύφθη ὅτι οὐχ ἑαυτοῖς, ὑμῖν δὲ διηκόνουν αὐτά ἃ νῦν ἀνηγγέλη ὑμῖν διὰ τῶν εὐαγγελισαμένων ὑμᾶς ἐν Πνεύματι Ἁγίῳ ἀποσταλέντι ἀπ᾽ οὐρανοῦ, εἰς ἃ ἐπιθυμοῦσιν ἄγγελοι παρακύψαι.

\vs{13}Διὸ ἀναζωσάμενοι τὰς ὀσφύας τῆς διανοίας ὑμῶν νήφοντες τελείως ἐλπίσατε ἐπὶ τὴν φερομένην ὑμῖν χάριν ἐν ἀποκαλύψει Ἰησοῦ Χριστοῦ.
\vs{14}ὡς τέκνα ὑπακοῆς μὴ συσχηματιζόμενοι ταῖς πρότερον ἐν τῇ ἀγνοίᾳ ὑμῶν ἐπιθυμίαις,
\vs{15}ἀλλὰ κατὰ τὸν καλέσαντα ὑμᾶς ἅγιον καὶ αὐτοὶ ἅγιοι ἐν πάσῃ ἀναστροφῇ γενήθητε,
\vs{16}διότι γέγραπται· Ἅγιοι ἔσεσθε, ὅτι ἐγὼ ἅγιος.
\vs{17}Καὶ εἰ Πατέρα ἐπικαλεῖσθε τὸν ἀπροσωπολήμπτως κρίνοντα κατὰ τὸ ἑκάστου ἔργον, ἐν φόβῳ τὸν τῆς παροικίας ὑμῶν χρόνον ἀναστράφητε

\vs{18}εἰδότες ὅτι οὐ φθαρτοῖς, ἀργυρίῳ ἢ χρυσίῳ, ἐλυτρώθητε 
\begin{poetryblock}

\begin{quote}ἐκ τῆς ματαίας ὑμῶν ἀναστροφῆς πατροπαραδότου\end{quote}

\begin{quote} \vs{19}ἀλλὰ τιμίῳ αἵματι ὡς ἀμνοῦ ἀμώμου καὶ ἀσπίλου Χριστοῦ\end{quote}

\begin{quote} \vs{20}προεγνωσμένου μὲν πρὸ καταβολῆς κόσμου,\end{quote} 

\begin{quote}φανερωθέντος δὲ ἐπ᾽ ἐσχάτου τῶν χρόνων\end{quote} 

\begin{quote}δι᾽ ὑμᾶς\end{quote}
\end{poetryblock}

\vs{21}τοὺς δι᾽ αὐτοῦ πιστοὺς εἰς Θεὸν 
\begin{poetryblock}

\begin{quote}τὸν ἐγείραντα αὐτὸν ἐκ νεκρῶν καὶ δόξαν αὐτῷ δόντα,\end{quote} 

\begin{quote}ὥστε τὴν πίστιν ὑμῶν καὶ ἐλπίδα εἶναι εἰς Θεόν.\end{quote}
\end{poetryblock}

\vs{22}Τὰς ψυχὰς ὑμῶν ἡγνικότες ἐν τῇ ὑπακοῇ τῆς ἀληθείας εἰς φιλαδελφίαν ἀνυπόκριτον ἐκ καθαρᾶς καρδίας ἀλλήλους ἀγαπήσατε ἐκτενῶς
\vs{23}ἀναγεγεννημένοι οὐκ ἐκ σπορᾶς φθαρτῆς ἀλλὰ ἀφθάρτου διὰ λόγου ζῶντος Θεοῦ καὶ μένοντος.
\vs{24}διότι 
\begin{poetryblock}

\begin{quote}Πᾶσα σὰρξ ὡς χόρτος\end{quote} 

\begin{quote}καὶ πᾶσα δόξα αὐτῆς ὡς ἄνθος χόρτου·\end{quote} 

\begin{quote}ἐξηράνθη ὁ χόρτος καὶ τὸ ἄνθος ἐξέπεσεν·\end{quote}

\begin{quote} \vs{25}τὸ δὲ ῥῆμα Κυρίου μένει εἰς τὸν αἰῶνα.\end{quote}
\end{poetryblock}

Τοῦτο δέ ἐστιν τὸ ῥῆμα τὸ εὐαγγελισθὲν εἰς ὑμᾶς.

\ch{2}
Ἀποθέμενοι οὖν πᾶσαν κακίαν καὶ πάντα δόλον καὶ ὑποκρίσεις καὶ φθόνους καὶ πάσας καταλαλιάς
\vs{2}ὡς ἀρτιγέννητα βρέφη τὸ λογικὸν ἄδολον γάλα ἐπιποθήσατε, ἵνα ἐν αὐτῷ αὐξηθῆτε εἰς σωτηρίαν,
\vs{3}εἰ ἐγεύσασθε ὅτι χρηστὸς ὁ Κύριος.
\vs{4}Πρὸς ὃν προσερχόμενοι λίθον ζῶντα ὑπὸ ἀνθρώπων μὲν ἀποδεδοκιμασμένον, παρὰ δὲ Θεῷ ἐκλεκτὸν ἔντιμον,
\vs{5}καὶ αὐτοὶ ὡς λίθοι ζῶντες οἰκοδομεῖσθε οἶκος πνευματικὸς εἰς ἱεράτευμα ἅγιον ἀνενέγκαι πνευματικὰς θυσίας εὐπροσδέκτους Θεῷ διὰ Ἰησοῦ Χριστοῦ.
\vs{6}διότι περιέχει ἐν γραφῇ· 
\begin{poetryblock}

\begin{quote}Ἰδοὺ τίθημι ἐν Σιὼν λίθον ἀκρογωνιαῖον ἐκλεκτὸν ἔντιμον,\end{quote} 

\begin{quote}καὶ ὁ πιστεύων ἐπ᾽ αὐτῷ οὐ μὴ καταισχυνθῇ.\end{quote}
\end{poetryblock}

\vs{7}Ὑμῖν οὖν ἡ τιμὴ τοῖς πιστεύουσιν, ἀπιστοῦσιν δὲ Λίθος ὃν ἀπεδοκίμασαν οἱ οἰκοδομοῦντες, οὗτος ἐγενήθη εἰς κεφαλὴν γωνίας
\vs{8}Καὶ Λίθος προσκόμματος καὶ πέτρα σκανδάλου· Οἳ προσκόπτουσιν τῷ λόγῳ ἀπειθοῦντες εἰς ὃ καὶ ἐτέθησαν.
\vs{9}Ὑμεῖς δὲ γένος ἐκλεκτόν, βασίλειον ἱεράτευμα, ἔθνος ἅγιον, λαὸς εἰς περιποίησιν, ὅπως τὰς ἀρετὰς ἐξαγγείλητε τοῦ ἐκ σκότους ὑμᾶς καλέσαντος εἰς τὸ θαυμαστὸν αὐτοῦ φῶς·
\vs{10}οἵ ποτε οὐ λαὸς, νῦν δὲ λαὸς Θεοῦ, οἱ οὐκ ἠλεημένοι, νῦν δὲ ἐλεηθέντες.

\vs{11}Ἀγαπητοί, παρακαλῶ ὡς παροίκους καὶ παρεπιδήμους ἀπέχεσθαι τῶν σαρκικῶν ἐπιθυμιῶν αἵτινες στρατεύονται κατὰ τῆς ψυχῆς·
\vs{12}τὴν ἀναστροφὴν ὑμῶν ἐν τοῖς ἔθνεσιν ἔχοντες καλήν, ἵνα ἐν ᾧ καταλαλοῦσιν ὑμῶν ὡς κακοποιῶν ἐκ τῶν καλῶν ἔργων ἐποπτεύοντες δοξάσωσιν τὸν Θεὸν ἐν ἡμέρᾳ ἐπισκοπῆς.

\vs{13}Ὑποτάγητε πάσῃ ἀνθρωπίνῃ κτίσει διὰ τὸν Κύριον, εἴτε βασιλεῖ ὡς ὑπερέχοντι
\vs{14}εἴτε ἡγεμόσιν ὡς δι᾽ αὐτοῦ πεμπομένοις εἰς ἐκδίκησιν κακοποιῶν, ἔπαινον δὲ ἀγαθοποιῶν,
\vs{15}ὅτι οὕτως ἐστὶν τὸ θέλημα τοῦ Θεοῦ ἀγαθοποιοῦντας φιμοῦν τὴν τῶν ἀφρόνων ἀνθρώπων ἀγνωσίαν,
\vs{16}ὡς ἐλεύθεροι καὶ μὴ ὡς ἐπικάλυμμα ἔχοντες τῆς κακίας τὴν ἐλευθερίαν ἀλλ᾽ ὡς Θεοῦ δοῦλοι.
\vs{17}Πάντας τιμήσατε, τὴν ἀδελφότητα ἀγαπᾶτε, τὸν Θεὸν φοβεῖσθε, τὸν βασιλέα τιμᾶτε.

\vs{18}Οἱ οἰκέται ὑποτασσόμενοι ἐν παντὶ φόβῳ τοῖς δεσπόταις, οὐ μόνον τοῖς ἀγαθοῖς καὶ ἐπιεικέσιν ἀλλὰ καὶ τοῖς σκολιοῖς.
\vs{19}τοῦτο γὰρ χάρις, εἰ διὰ συνείδησιν Θεοῦ ὑποφέρει τις λύπας πάσχων ἀδίκως.
\vs{20}ποῖον γὰρ κλέος, εἰ ἁμαρτάνοντες καὶ κολαφιζόμενοι ὑπομενεῖτε; ἀλλ᾽ εἰ ἀγαθοποιοῦντες καὶ πάσχοντες ὑπομενεῖτε, τοῦτο χάρις παρὰ Θεῷ.
\begin{poetryblock}

\begin{quote} \vs{21}Εἰς τοῦτο γὰρ ἐκλήθητε,\end{quote} 

\begin{quote}ὅτι καὶ Χριστὸς ἔπαθεν ὑπὲρ ὑμῶν\end{quote} 

\begin{quote}ὑμῖν ὑπολιμπάνων ὑπογραμμὸν,\end{quote} 

\begin{quote}ἵνα ἐπακολουθήσητε τοῖς ἴχνεσιν αὐτοῦ,\end{quote}

\begin{quote} \vs{22}Ὃς ἁμαρτίαν οὐκ ἐποίησεν\end{quote} 

\begin{quote}οὐδὲ εὑρέθη δόλος ἐν τῷ στόματι αὐτοῦ,\end{quote}

\begin{quote} \vs{23}ὃς λοιδορούμενος οὐκ ἀντελοιδόρει,\end{quote} 

\begin{quote}πάσχων οὐκ ἠπείλει,\end{quote} 

\begin{quote}παρεδίδου δὲ τῷ κρίνοντι δικαίως\end{quote}

\begin{quote} \vs{24}ὃς τὰς ἁμαρτίας ἡμῶν αὐτὸς ἀνήνεγκεν\end{quote} 

\begin{quote}ἐν τῷ σώματι αὐτοῦ ἐπὶ τὸ ξύλον,\end{quote} 

\begin{quote}ἵνα ταῖς ἁμαρτίαις ἀπογενόμενοι\end{quote} 

\begin{quote}τῇ δικαιοσύνῃ ζήσωμεν,\end{quote} 

\begin{quote}Οὗ τῷ μώλωπι ἰάθητε.\end{quote}

\begin{quote} \vs{25}Ἦτε γὰρ ὡς πρόβατα πλανώμενοι,\end{quote} 

\begin{quote}ἀλλὰ ἐπεστράφητε νῦν ἐπὶ τὸν Ποιμένα\end{quote} 

\begin{quote}καὶ Ἐπίσκοπον τῶν ψυχῶν ὑμῶν.\end{quote}
\end{poetryblock}

\ch{3}
Ὁμοίως αἱ γυναῖκες, ὑποτασσόμεναι τοῖς ἰδίοις ἀνδράσιν, ἵνα καὶ εἴ τινες ἀπειθοῦσιν τῷ λόγῳ, διὰ τῆς τῶν γυναικῶν ἀναστροφῆς ἄνευ λόγου κερδηθήσονται
\vs{2}ἐποπτεύσαντες τὴν ἐν φόβῳ ἁγνὴν ἀναστροφὴν ὑμῶν.
\vs{3}ὧν ἔστω οὐχ ὁ ἔξωθεν ἐμπλοκῆς τριχῶν καὶ περιθέσεως χρυσίων ἢ ἐνδύσεως ἱματίων κόσμος,
\vs{4}ἀλλ᾽ ὁ κρυπτὸς τῆς καρδίας ἄνθρωπος ἐν τῷ ἀφθάρτῳ τοῦ πραέως καὶ ἡσυχίου πνεύματος ὅ ἐστιν ἐνώπιον τοῦ Θεοῦ πολυτελές.
\vs{5}Οὕτως γάρ ποτε καὶ αἱ ἅγιαι γυναῖκες αἱ ἐλπίζουσαι εἰς Θεὸν ἐκόσμουν ἑαυτάς ὑποτασσόμεναι τοῖς ἰδίοις ἀνδράσιν,
\vs{6}ὡς Σάρρα ὑπήκουσεν τῷ Ἀβραάμ κύριον αὐτὸν καλοῦσα ἧς ἐγενήθητε τέκνα ἀγαθοποιοῦσαι καὶ μὴ φοβούμεναι μηδεμίαν πτόησιν.

\vs{7}Οἱ ἄνδρες ὁμοίως, συνοικοῦντες κατὰ γνῶσιν ὡς ἀσθενεστέρῳ σκεύει τῷ γυναικείῳ, ἀπονέμοντες τιμήν ὡς καὶ συνκληρονόμοις χάριτος ζωῆς εἰς τὸ μὴ ἐνκόπτεσθαι τὰς προσευχὰς ὑμῶν.

\vs{8}Τὸ δὲ τέλος πάντες ὁμόφρονες, συμπαθεῖς, φιλάδελφοι, εὔσπλαγχνοι, ταπεινόφρονες,
\vs{9}μὴ ἀποδιδόντες κακὸν ἀντὶ κακοῦ ἢ λοιδορίαν ἀντὶ λοιδορίας, τοὐναντίον δὲ εὐλογοῦντες, ὅτι εἰς τοῦτο ἐκλήθητε, ἵνα εὐλογίαν κληρονομήσητε.
\begin{poetryblock}

\begin{quote} \vs{10}Ὁ γὰρ θέλων ζωὴν ἀγαπᾶν\end{quote} 

\begin{quote}καὶ ἰδεῖν ἡμέρας ἀγαθὰς\end{quote} 

\begin{quote}παυσάτω τὴν γλῶσσαν ἀπὸ κακοῦ\end{quote} 

\begin{quote}καὶ χείλη τοῦ μὴ λαλῆσαι δόλον,\end{quote}

\begin{quote} \vs{11}ἐκκλινάτω δὲ ἀπὸ κακοῦ καὶ ποιησάτω ἀγαθόν,\end{quote} 

\begin{quote}ζητησάτω εἰρήνην καὶ διωξάτω αὐτήν·\end{quote}

\begin{quote} \vs{12}ὅτι ὀφθαλμοὶ Κυρίου ἐπὶ δικαίους\end{quote} 

\begin{quote}καὶ ὦτα αὐτοῦ εἰς δέησιν αὐτῶν,\end{quote} 

\begin{quote}πρόσωπον δὲ Κυρίου ἐπὶ ποιοῦντας κακά.\end{quote}
\end{poetryblock}
\vs{13}Καὶ τίς ὁ κακώσων ὑμᾶς, ἐὰν τοῦ ἀγαθοῦ ζηλωταὶ γένησθε;
\vs{14}Ἀλλ᾽ εἰ καὶ πάσχοιτε διὰ δικαιοσύνην, μακάριοι. Τὸν δὲ φόβον αὐτῶν μὴ φοβηθῆτε μηδὲ ταραχθῆτε,
\vs{15}Κύριον δὲ τὸν Χριστὸν ἁγιάσατε ἐν ταῖς καρδίαις ὑμῶν, ἕτοιμοι ἀεὶ πρὸς ἀπολογίαν παντὶ τῷ αἰτοῦντι ὑμᾶς λόγον περὶ τῆς ἐν ὑμῖν ἐλπίδος,
\vs{16}ἀλλὰ μετὰ πραΰτητος καὶ φόβου, συνείδησιν ἔχοντες ἀγαθήν, ἵνα ἐν ᾧ καταλαλεῖσθε καταισχυνθῶσιν οἱ ἐπηρεάζοντες ὑμῶν τὴν ἀγαθὴν ἐν Χριστῷ ἀναστροφήν.
\vs{17}κρεῖττον γὰρ ἀγαθοποιοῦντας, εἰ θέλοι τὸ θέλημα τοῦ Θεοῦ, πάσχειν ἢ κακοποιοῦντας.
\begin{poetryblock}

\begin{quote} \vs{18}ὅτι καὶ Χριστὸς ἅπαξ περὶ ἁμαρτιῶν ἔπαθεν,\end{quote} 

\begin{quote}δίκαιος ὑπὲρ ἀδίκων,\end{quote} 

\begin{quote}ἵνα ὑμᾶς προσαγάγῃ τῷ Θεῷ\end{quote} 

\begin{quote}θανατωθεὶς μὲν σαρκὶ,\end{quote} 

\begin{quote}ζωοποιηθεὶς δὲ πνεύματι·\end{quote}

\begin{quote} \vs{19}ἐν ᾧ καὶ τοῖς ἐν φυλακῇ πνεύμασιν\end{quote} 

\begin{quote}πορευθεὶς ἐκήρυξεν\end{quote}
\end{poetryblock}

\vs{20}ἀπειθήσασίν ποτε, ὅτε ἀπεξεδέχετο ἡ τοῦ Θεοῦ μακροθυμία ἐν ἡμέραις Νῶε κατασκευαζομένης κιβωτοῦ εἰς ἣν ὀλίγοι, τοῦτ᾽ ἔστιν ὀκτὼ ψυχαί, διεσώθησαν δι᾽ ὕδατος
\vs{21}ὃ καὶ ὑμᾶς ἀντίτυπον νῦν σῴζει βάπτισμα, οὐ σαρκὸς ἀπόθεσις ῥύπου ἀλλὰ συνειδήσεως ἀγαθῆς ἐπερώτημα εἰς Θεόν, δι᾽ ἀναστάσεως Ἰησοῦ Χριστοῦ
\vs{22}ὅς ἐστιν ἐν δεξιᾷ τοῦ Θεοῦ πορευθεὶς εἰς οὐρανόν ὑποταγέντων αὐτῷ ἀγγέλων καὶ ἐξουσιῶν καὶ δυνάμεων.

\ch{4}
Χριστοῦ οὖν παθόντος σαρκὶ καὶ ὑμεῖς τὴν αὐτὴν ἔννοιαν ὁπλίσασθε, ὅτι ὁ παθὼν σαρκὶ πέπαυται ἁμαρτίας
\vs{2}εἰς τὸ μηκέτι ἀνθρώπων ἐπιθυμίαις ἀλλὰ θελήματι Θεοῦ τὸν ἐπίλοιπον ἐν σαρκὶ βιῶσαι χρόνον.
\vs{3}ἀρκετὸς γὰρ ὁ παρεληλυθὼς χρόνος τὸ βούλημα τῶν ἐθνῶν κατειργάσθαι πεπορευμένους ἐν ἀσελγείαις, ἐπιθυμίαις, οἰνοφλυγίαις, κώμοις, πότοις καὶ ἀθεμίτοις εἰδωλολατρίαις.
\vs{4}Ἐν ᾧ ξενίζονται μὴ συντρεχόντων ὑμῶν εἰς τὴν αὐτὴν τῆς ἀσωτίας ἀνάχυσιν βλασφημοῦντες,
\vs{5}οἳ ἀποδώσουσιν λόγον τῷ ἑτοίμως ἔχοντι κρῖναι ζῶντας καὶ νεκρούς.
\vs{6}εἰς τοῦτο γὰρ καὶ νεκροῖς εὐηγγελίσθη, ἵνα κριθῶσι μὲν κατὰ ἀνθρώπους σαρκί, ζῶσι δὲ κατὰ Θεὸν πνεύματι.

\vs{7}Πάντων δὲ τὸ τέλος ἤγγικεν. σωφρονήσατε οὖν καὶ νήψατε εἰς προσευχάς
\vs{8}πρὸ πάντων τὴν εἰς ἑαυτοὺς ἀγάπην ἐκτενῆ ἔχοντες, ὅτι ἀγάπη καλύπτει πλῆθος ἁμαρτιῶν,
\vs{9}φιλόξενοι εἰς ἀλλήλους ἄνευ γογγυσμοῦ,
\vs{10}ἕκαστος καθὼς ἔλαβεν χάρισμα εἰς ἑαυτοὺς αὐτὸ διακονοῦντες ὡς καλοὶ οἰκονόμοι ποικίλης χάριτος Θεοῦ.
\vs{11}εἴ τις λαλεῖ, ὡς λόγια Θεοῦ· εἴ τις διακονεῖ, ὡς ἐξ ἰσχύος ἧς χορηγεῖ ὁ Θεός, ἵνα ἐν πᾶσιν δοξάζηται ὁ Θεὸς διὰ Ἰησοῦ Χριστοῦ ᾧ ἐστιν ἡ δόξα καὶ τὸ κράτος εἰς τοὺς αἰῶνας τῶν αἰώνων, ἀμήν.

\vs{12}Ἀγαπητοί, μὴ ξενίζεσθε τῇ ἐν ὑμῖν πυρώσει πρὸς πειρασμὸν ὑμῖν γινομένῃ ὡς ξένου ὑμῖν συμβαίνοντος,
\vs{13}ἀλλὰ καθὸ κοινωνεῖτε τοῖς τοῦ Χριστοῦ παθήμασιν, χαίρετε, ἵνα καὶ ἐν τῇ ἀποκαλύψει τῆς δόξης αὐτοῦ χαρῆτε ἀγαλλιώμενοι.
\vs{14}Εἰ ὀνειδίζεσθε ἐν ὀνόματι Χριστοῦ, μακάριοι, ὅτι τὸ τῆς δόξης καὶ τὸ τοῦ Θεοῦ Πνεῦμα ἐφ᾽ ὑμᾶς ἀναπαύεται.
\vs{15}μὴ γάρ τις ὑμῶν πασχέτω ὡς φονεὺς ἢ κλέπτης ἢ κακοποιὸς ἢ ὡς ἀλλοτριεπίσκοπος·
\vs{16}εἰ δὲ ὡς Χριστιανός, μὴ αἰσχυνέσθω, δοξαζέτω δὲ τὸν Θεὸν ἐν τῷ ὀνόματι τούτῳ.
\vs{17}ὅτι ὁ καιρὸς τοῦ ἄρξασθαι τὸ κρίμα ἀπὸ τοῦ οἴκου τοῦ Θεοῦ· εἰ δὲ πρῶτον ἀφ᾽ ἡμῶν, τί τὸ τέλος τῶν ἀπειθούντων τῷ τοῦ Θεοῦ εὐαγγελίῳ;
\vs{18}καὶ Εἰ ὁ δίκαιος μόλις σώζεται, ὁ ἀσεβὴς καὶ ἁμαρτωλὸς ποῦ φανεῖται;
\vs{19}Ὥστε καὶ οἱ πάσχοντες κατὰ τὸ θέλημα τοῦ Θεοῦ πιστῷ Κτίστῃ παρατιθέσθωσαν τὰς ψυχὰς αὐτῶν ἐν ἀγαθοποιΐᾳ.

\ch{5}
Πρεσβυτέρους τοὺς ἐν ὑμῖν παρακαλῶ ὁ συμπρεσβύτερος καὶ μάρτυς τῶν τοῦ Χριστοῦ παθημάτων, ὁ καὶ τῆς μελλούσης ἀποκαλύπτεσθαι δόξης κοινωνός·
\vs{2}ποιμάνατε τὸ ἐν ὑμῖν ποίμνιον τοῦ Θεοῦ ἐπισκοποῦντες μὴ ἀναγκαστῶς ἀλλὰ ἑκουσίως κατὰ Θεόν, μηδὲ αἰσχροκερδῶς ἀλλὰ προθύμως,
\vs{3}μηδ᾽ ὡς κατακυριεύοντες τῶν κλήρων ἀλλὰ τύποι γινόμενοι τοῦ ποιμνίου·
\vs{4}καὶ φανερωθέντος τοῦ Ἀρχιποίμενος κομιεῖσθε τὸν ἀμαράντινον τῆς δόξης στέφανον.
\vs{5}Ὁμοίως, νεώτεροι, ὑποτάγητε πρεσβυτέροις· πάντες δὲ ἀλλήλοις τὴν ταπεινοφροσύνην ἐγκομβώσασθε, ὅτι Ὁ Θεὸς ὑπερηφάνοις ἀντιτάσσεται, ταπεινοῖς δὲ δίδωσιν χάριν.

\vs{6}Ταπεινώθητε οὖν ὑπὸ τὴν κραταιὰν χεῖρα τοῦ Θεοῦ, ἵνα ὑμᾶς ὑψώσῃ ἐν καιρῷ,
\vs{7}πᾶσαν τὴν μέριμναν ὑμῶν ἐπιρίψαντες ἐπ᾽ αὐτόν, ὅτι αὐτῷ μέλει περὶ ὑμῶν.
\vs{8}Νήψατε, γρηγορήσατε. ὁ ἀντίδικος ὑμῶν διάβολος ὡς λέων ὠρυόμενος περιπατεῖ ζητῶν τινα καταπιεῖν·
\vs{9}ᾧ ἀντίστητε στερεοὶ τῇ πίστει εἰδότες τὰ αὐτὰ τῶν παθημάτων τῇ ἐν κόσμῳ ὑμῶν ἀδελφότητι ἐπιτελεῖσθαι.
\vs{10}Ὁ δὲ Θεὸς πάσης χάριτος, ὁ καλέσας ὑμᾶς εἰς τὴν αἰώνιον αὐτοῦ δόξαν ἐν Χριστῷ ὀλίγον παθόντας αὐτὸς καταρτίσει, στηρίξει, σθενώσει, θεμελιώσει.
\vs{11}αὐτῷ τὸ κράτος εἰς τοὺς αἰῶνας, ἀμήν.

\vs{12}Διὰ Σιλουανοῦ ὑμῖν τοῦ πιστοῦ ἀδελφοῦ, ὡς λογίζομαι, δι᾽ ὀλίγων ἔγραψα παρακαλῶν καὶ ἐπιμαρτυρῶν ταύτην εἶναι ἀληθῆ χάριν τοῦ Θεοῦ εἰς ἣν στῆτε.
\vs{13}Ἀσπάζεται ὑμᾶς ἡ ἐν Βαβυλῶνι συνεκλεκτὴ καὶ Μάρκος ὁ υἱός μου.
\vs{14}Ἀσπάσασθε ἀλλήλους ἐν φιλήματι ἀγάπης.

Εἰρήνη ὑμῖν πᾶσιν τοῖς ἐν Χριστῷ.


\def\book{ΠΕΤΡΟΥ Β}
\biblebook{ΠΕΤΡΟΥ Β}


\lettrine[lines=2, loversize=0.2, nindent=0em, findent=.25em]{\textcolor{bookheadingcolor}{Σ}}{υμεὼν} Πέτρος δοῦλος καὶ ἀπόστολος Ἰησοῦ Χριστοῦ Τοῖς ἰσότιμον ἡμῖν λαχοῦσιν πίστιν ἐν δικαιοσύνῃ τοῦ Θεοῦ ἡμῶν καὶ Σωτῆρος Ἰησοῦ Χριστοῦ,
\vs{2}Χάρις ὑμῖν καὶ εἰρήνη πληθυνθείη ἐν ἐπιγνώσει τοῦ Θεοῦ καὶ Ἰησοῦ τοῦ Κυρίου ἡμῶν.

\vs{3}Ὡς πάντα ἡμῖν τῆς θείας δυνάμεως αὐτοῦ τὰ πρὸς ζωὴν καὶ εὐσέβειαν δεδωρημένης διὰ τῆς ἐπιγνώσεως τοῦ καλέσαντος ἡμᾶς ἰδίᾳ δόξῃ καὶ ἀρετῇ
\vs{4}δι᾽ ὧν τὰ τίμια καὶ μέγιστα ἡμῖν ἐπαγγέλματα δεδώρηται, ἵνα διὰ τούτων γένησθε θείας κοινωνοὶ φύσεως ἀποφυγόντες τῆς ἐν τῷ κόσμῳ ἐν ἐπιθυμίᾳ φθορᾶς.
\vs{5}Καὶ αὐτὸ τοῦτο δὲ σπουδὴν πᾶσαν παρεισενέγκαντες ἐπιχορηγήσατε ἐν τῇ πίστει ὑμῶν τὴν ἀρετήν, ἐν δὲ τῇ ἀρετῇ τὴν γνῶσιν,
\vs{6}ἐν δὲ τῇ γνώσει τὴν ἐγκράτειαν, ἐν δὲ τῇ ἐγκρατείᾳ τὴν ὑπομονήν, ἐν δὲ τῇ ὑπομονῇ τὴν εὐσέβειαν,
\vs{7}ἐν δὲ τῇ εὐσεβείᾳ τὴν φιλαδελφίαν, ἐν δὲ τῇ φιλαδελφίᾳ τὴν ἀγάπην.
\vs{8}ταῦτα γὰρ ὑμῖν ὑπάρχοντα καὶ πλεονάζοντα οὐκ ἀργοὺς οὐδὲ ἀκάρπους καθίστησιν εἰς τὴν τοῦ Κυρίου ἡμῶν Ἰησοῦ Χριστοῦ ἐπίγνωσιν·
\vs{9}ᾧ γὰρ μὴ πάρεστιν ταῦτα, τυφλός ἐστιν μυωπάζων λήθην λαβὼν τοῦ καθαρισμοῦ τῶν πάλαι αὐτοῦ ἁμαρτιῶν.
\vs{10}Διὸ μᾶλλον, ἀδελφοί, σπουδάσατε βεβαίαν ὑμῶν τὴν κλῆσιν καὶ ἐκλογὴν ποιεῖσθαι· ταῦτα γὰρ ποιοῦντες οὐ μὴ πταίσητέ ποτε.
\vs{11}οὕτως γὰρ πλουσίως ἐπιχορηγηθήσεται ὑμῖν ἡ εἴσοδος εἰς τὴν αἰώνιον βασιλείαν τοῦ Κυρίου ἡμῶν καὶ Σωτῆρος Ἰησοῦ Χριστοῦ.

\vs{12}Διὸ μελλήσω ἀεὶ ὑμᾶς ὑπομιμνῄσκειν περὶ τούτων καίπερ εἰδότας καὶ ἐστηριγμένους ἐν τῇ παρούσῃ ἀληθείᾳ.
\vs{13}δίκαιον δὲ ἡγοῦμαι, ἐφ᾽ ὅσον εἰμὶ ἐν τούτῳ τῷ σκηνώματι, διεγείρειν ὑμᾶς ἐν ὑπομνήσει
\vs{14}εἰδὼς ὅτι ταχινή ἐστιν ἡ ἀπόθεσις τοῦ σκηνώματός μου, καθὼς καὶ ὁ Κύριος ἡμῶν Ἰησοῦς Χριστὸς ἐδήλωσέν μοι.
\vs{15}σπουδάσω δὲ καὶ ἑκάστοτε ἔχειν ὑμᾶς μετὰ τὴν ἐμὴν ἔξοδον τὴν τούτων μνήμην ποιεῖσθαι.

\vs{16}Οὐ γὰρ σεσοφισμένοις μύθοις ἐξακολουθήσαντες ἐγνωρίσαμεν ὑμῖν τὴν τοῦ Κυρίου ἡμῶν Ἰησοῦ Χριστοῦ δύναμιν καὶ παρουσίαν ἀλλ᾽ ἐπόπται γενηθέντες τῆς ἐκείνου μεγαλειότητος.
\vs{17}λαβὼν γὰρ παρὰ Θεοῦ Πατρὸς τιμὴν καὶ δόξαν φωνῆς ἐνεχθείσης αὐτῷ τοιᾶσδε ὑπὸ τῆς μεγαλοπρεποῦς δόξης· Ὁ Υἱός μου ὁ ἀγαπητός μου οὗτός ἐστιν εἰς ὃν ἐγὼ εὐδόκησα
\vs{18}καὶ ταύτην τὴν φωνὴν ἡμεῖς ἠκούσαμεν ἐξ οὐρανοῦ ἐνεχθεῖσαν σὺν αὐτῷ ὄντες ἐν τῷ ἁγίῳ ὄρει.
\vs{19}Καὶ ἔχομεν βεβαιότερον τὸν προφητικὸν λόγον ᾧ καλῶς ποιεῖτε προσέχοντες ὡς λύχνῳ φαίνοντι ἐν αὐχμηρῷ τόπῳ, ἕως οὗ ἡμέρα διαυγάσῃ καὶ φωσφόρος ἀνατείλῃ ἐν ταῖς καρδίαις ὑμῶν,
\vs{20}τοῦτο πρῶτον γινώσκοντες ὅτι πᾶσα προφητεία γραφῆς ἰδίας ἐπιλύσεως οὐ γίνεται·
\vs{21}οὐ γὰρ θελήματι ἀνθρώπου ἠνέχθη προφητεία ποτέ, ἀλλὰ ὑπὸ Πνεύματος Ἁγίου φερόμενοι ἐλάλησαν ἀπὸ Θεοῦ ἄνθρωποι.

\ch{2}
Ἐγένοντο δὲ καὶ ψευδοπροφῆται ἐν τῷ λαῷ, ὡς καὶ ἐν ὑμῖν ἔσονται ψευδοδιδάσκαλοι οἵτινες παρεισάξουσιν αἱρέσεις ἀπωλείας καὶ τὸν ἀγοράσαντα αὐτοὺς δεσπότην ἀρνούμενοι ἐπάγοντες ἑαυτοῖς ταχινὴν ἀπώλειαν.
\vs{2}καὶ πολλοὶ ἐξακολουθήσουσιν αὐτῶν ταῖς ἀσελγείαις δι᾽ οὓς ἡ ὁδὸς τῆς ἀληθείας βλασφημηθήσεται,
\vs{3}καὶ ἐν πλεονεξίᾳ πλαστοῖς λόγοις ὑμᾶς ἐμπορεύσονται οἷς τὸ κρίμα ἔκπαλαι οὐκ ἀργεῖ καὶ ἡ ἀπώλεια αὐτῶν οὐ νυστάζει.

\vs{4}Εἰ γὰρ ὁ Θεὸς ἀγγέλων ἁμαρτησάντων οὐκ ἐφείσατο ἀλλὰ σειραῖς ζόφου ταρταρώσας παρέδωκεν εἰς κρίσιν τηρουμένους
\vs{5}καὶ ἀρχαίου κόσμου οὐκ ἐφείσατο ἀλλὰ ὄγδοον Νῶε δικαιοσύνης κήρυκα ἐφύλαξεν κατακλυσμὸν κόσμῳ ἀσεβῶν ἐπάξας
\vs{6}καὶ πόλεις Σοδόμων καὶ Γομόρρας τεφρώσας καταστροφῇ κατέκρινεν ὑπόδειγμα μελλόντων ἀσεβέσιν τεθεικώς
\vs{7}καὶ δίκαιον Λὼτ καταπονούμενον ὑπὸ τῆς τῶν ἀθέσμων ἐν ἀσελγείᾳ ἀναστροφῆς ἐρρύσατο·
\vs{8}βλέμματι γὰρ καὶ ἀκοῇ ὁ δίκαιος ἐνκατοικῶν ἐν αὐτοῖς ἡμέραν ἐξ ἡμέρας ψυχὴν δικαίαν ἀνόμοις ἔργοις ἐβασάνιζεν·
\vs{9}οἶδεν Κύριος εὐσεβεῖς ἐκ πειρασμοῦ ῥύεσθαι, ἀδίκους δὲ εἰς ἡμέραν κρίσεως κολαζομένους τηρεῖν,
\vs{10}μάλιστα δὲ τοὺς ὀπίσω σαρκὸς ἐν ἐπιθυμίᾳ μιασμοῦ πορευομένους καὶ κυριότητος καταφρονοῦντας. τολμηταὶ αὐθάδεις δόξας οὐ τρέμουσιν βλασφημοῦντες,
\vs{11}ὅπου ἄγγελοι ἰσχύϊ καὶ δυνάμει μείζονες ὄντες οὐ φέρουσιν κατ᾽ αὐτῶν παρὰ Κυρίῳ βλάσφημον κρίσιν.

\vs{12}Οὗτοι δέ ὡς ἄλογα ζῷα γεγεννημένα φυσικὰ εἰς ἅλωσιν καὶ φθοράν ἐν οἷς ἀγνοοῦσιν βλασφημοῦντες ἐν τῇ φθορᾷ αὐτῶν καὶ φθαρήσονται
\vs{13}ἀδικούμενοι μισθὸν ἀδικίας ἡδονὴν ἡγούμενοι τὴν ἐν ἡμέρᾳ τρυφήν, σπίλοι καὶ μῶμοι ἐντρυφῶντες ἐν ταῖς ἀπάταις αὐτῶν συνευωχούμενοι ὑμῖν,
\vs{14}ὀφθαλμοὺς ἔχοντες μεστοὺς μοιχαλίδος καὶ ἀκαταπαύστους ἁμαρτίας, δελεάζοντες ψυχὰς ἀστηρίκτους, καρδίαν γεγυμνασμένην πλεονεξίας ἔχοντες, κατάρας τέκνα.
\vs{15}Καταλιπόντες εὐθεῖαν ὁδὸν ἐπλανήθησαν ἐξακολουθήσαντες τῇ ὁδῷ τοῦ Βαλαὰμ τοῦ Βοσὸρ ὃς μισθὸν ἀδικίας ἠγάπησεν,
\vs{16}ἔλεγξιν δὲ ἔσχεν ἰδίας παρανομίας· ὑποζύγιον ἄφωνον ἐν ἀνθρώπου φωνῇ φθεγξάμενον ἐκώλυσεν τὴν τοῦ προφήτου παραφρονίαν.
\vs{17}Οὗτοί εἰσιν πηγαὶ ἄνυδροι καὶ ὁμίχλαι ὑπὸ λαίλαπος ἐλαυνόμεναι οἷς ὁ ζόφος τοῦ σκότους τετήρηται.
\vs{18}ὑπέρογκα γὰρ ματαιότητος φθεγγόμενοι δελεάζουσιν ἐν ἐπιθυμίαις σαρκὸς ἀσελγείαις τοὺς ὀλίγως ἀποφεύγοντας τοὺς ἐν πλάνῃ ἀναστρεφομένους,
\vs{19}ἐλευθερίαν αὐτοῖς ἐπαγγελλόμενοι αὐτοὶ δοῦλοι ὑπάρχοντες τῆς φθορᾶς· ᾧ γάρ τις ἥττηται, τούτῳ δεδούλωται.
\vs{20}Εἰ γὰρ ἀποφυγόντες τὰ μιάσματα τοῦ κόσμου ἐν ἐπιγνώσει τοῦ Κυρίου καὶ Σωτῆρος Ἰησοῦ Χριστοῦ, τούτοις δὲ πάλιν ἐμπλακέντες ἡττῶνται, γέγονεν αὐτοῖς τὰ ἔσχατα χείρονα τῶν πρώτων.
\vs{21}κρεῖττον γὰρ ἦν αὐτοῖς μὴ ἐπεγνωκέναι τὴν ὁδὸν τῆς δικαιοσύνης ἢ ἐπιγνοῦσιν ὑποστρέψαι ἐκ τῆς παραδοθείσης αὐτοῖς ἁγίας ἐντολῆς.
\vs{22}συμβέβηκεν αὐτοῖς τὸ τῆς ἀληθοῦς παροιμίας· Κύων ἐπιστρέψας ἐπὶ τὸ ἴδιον ἐξέραμα, καί· Ὗς λουσαμένη εἰς κυλισμὸν βορβόρου.

\ch{3}
Ταύτην ἤδη, ἀγαπητοί, δευτέραν ὑμῖν γράφω ἐπιστολήν ἐν αἷς διεγείρω ὑμῶν ἐν ὑπομνήσει τὴν εἰλικρινῆ διάνοιαν
\vs{2}μνησθῆναι τῶν προειρημένων ῥημάτων ὑπὸ τῶν ἁγίων προφητῶν καὶ τῆς τῶν ἀποστόλων ὑμῶν ἐντολῆς τοῦ Κυρίου καὶ Σωτῆρος.
\vs{3}τοῦτο πρῶτον γινώσκοντες ὅτι ἐλεύσονται ἐπ᾽ ἐσχάτων τῶν ἡμερῶν ἐν ἐμπαιγμονῇ ἐμπαῖκται κατὰ τὰς ἰδίας ἐπιθυμίας αὐτῶν πορευόμενοι
\vs{4}καὶ λέγοντες· Ποῦ ἐστιν ἡ ἐπαγγελία τῆς παρουσίας αὐτοῦ; ἀφ᾽ ἧς γὰρ οἱ πατέρες ἐκοιμήθησαν, πάντα οὕτως διαμένει ἀπ᾽ ἀρχῆς κτίσεως.
\vs{5}Λανθάνει γὰρ αὐτοὺς τοῦτο θέλοντας ὅτι οὐρανοὶ ἦσαν ἔκπαλαι καὶ γῆ ἐξ ὕδατος καὶ δι᾽ ὕδατος συνεστῶσα τῷ τοῦ Θεοῦ λόγῳ
\vs{6}δι᾽ ὧν ὁ τότε κόσμος ὕδατι κατακλυσθεὶς ἀπώλετο·
\vs{7}οἱ δὲ νῦν οὐρανοὶ καὶ ἡ γῆ τῷ αὐτῷ λόγῳ τεθησαυρισμένοι εἰσὶν πυρὶ τηρούμενοι εἰς ἡμέραν κρίσεως καὶ ἀπωλείας τῶν ἀσεβῶν ἀνθρώπων.

\vs{8}Ἓν δὲ τοῦτο μὴ λανθανέτω ὑμᾶς, ἀγαπητοί, ὅτι μία ἡμέρα παρὰ Κυρίῳ ὡς χίλια ἔτη καὶ χίλια ἔτη ὡς ἡμέρα μία.
\vs{9}οὐ βραδύνει Κύριος τῆς ἐπαγγελίας, ὥς τινες βραδύτητα ἡγοῦνται, ἀλλὰ μακροθυμεῖ εἰς ὑμᾶς μὴ βουλόμενός τινας ἀπολέσθαι ἀλλὰ πάντας εἰς μετάνοιαν χωρῆσαι.

\vs{10}Ἥξει δὲ ἡμέρα Κυρίου ὡς κλέπτης ἐν ᾗ οἱ οὐρανοὶ ῥοιζηδὸν παρελεύσονται, στοιχεῖα δὲ καυσούμενα λυθήσεται, καὶ γῆ καὶ τὰ ἐν αὐτῇ ἔργα οὐχ εὑρεθήσεται.
\vs{11}Τούτων οὕτως πάντων λυομένων ποταποὺς δεῖ ὑπάρχειν ὑμᾶς ἐν ἁγίαις ἀναστροφαῖς καὶ εὐσεβείαις
\vs{12}προσδοκῶντας καὶ σπεύδοντας τὴν παρουσίαν τῆς τοῦ Θεοῦ ἡμέρας δι᾽ ἣν οὐρανοὶ πυρούμενοι λυθήσονται καὶ στοιχεῖα καυσούμενα τήκεται.
\vs{13}καινοὺς δὲ οὐρανοὺς καὶ γῆν καινὴν κατὰ τὸ ἐπάγγελμα αὐτοῦ προσδοκῶμεν ἐν οἷς δικαιοσύνη κατοικεῖ.

\vs{14}Διό, ἀγαπητοί, ταῦτα προσδοκῶντες σπουδάσατε ἄσπιλοι καὶ ἀμώμητοι αὐτῷ εὑρεθῆναι ἐν εἰρήνῃ
\vs{15}Καὶ τὴν τοῦ Κυρίου ἡμῶν μακροθυμίαν σωτηρίαν ἡγεῖσθε, καθὼς καὶ ὁ ἀγαπητὸς ἡμῶν ἀδελφὸς Παῦλος κατὰ τὴν δοθεῖσαν αὐτῷ σοφίαν ἔγραψεν ὑμῖν,
\vs{16}ὡς καὶ ἐν πάσαις ταῖς ἐπιστολαῖς λαλῶν ἐν αὐταῖς περὶ τούτων ἐν αἷς ἐστιν δυσνόητά τινα ἃ οἱ ἀμαθεῖς καὶ ἀστήρικτοι στρεβλοῦσιν ὡς καὶ τὰς λοιπὰς γραφὰς πρὸς τὴν ἰδίαν αὐτῶν ἀπώλειαν.

\vs{17}Ὑμεῖς οὖν, ἀγαπητοί, προγινώσκοντες φυλάσσεσθε, ἵνα μὴ τῇ τῶν ἀθέσμων πλάνῃ συναπαχθέντες ἐκπέσητε τοῦ ἰδίου στηριγμοῦ,
\vs{18}αὐξάνετε δὲ ἐν χάριτι καὶ γνώσει τοῦ Κυρίου ἡμῶν καὶ Σωτῆρος Ἰησοῦ Χριστοῦ. αὐτῷ ἡ δόξα καὶ νῦν καὶ εἰς ἡμέραν αἰῶνος.


\def\book{ΙΩΑΝΝΟΥ Α}
\biblebook{ΙΩΑΝΝΟΥ Α}


\lettrine[lines=2, loversize=0.2, nindent=0em, findent=.25em]{\textcolor{bookheadingcolor}{Ὃ}}{} ἦν ἀπ᾽ ἀρχῆς, ὃ ἀκηκόαμεν, ὃ ἑωράκαμεν τοῖς ὀφθαλμοῖς ἡμῶν, ὃ ἐθεασάμεθα καὶ αἱ χεῖρες ἡμῶν ἐψηλάφησαν περὶ τοῦ Λόγου τῆς ζωῆς—
\vs{2}καὶ ἡ ζωὴ ἐφανερώθη, καὶ ἑωράκαμεν καὶ μαρτυροῦμεν καὶ ἀπαγγέλλομεν ὑμῖν τὴν ζωὴν τὴν αἰώνιον ἥτις ἦν πρὸς τὸν Πατέρα καὶ ἐφανερώθη ἡμῖν—
\vs{3}ὃ ἑωράκαμεν καὶ ἀκηκόαμεν, ἀπαγγέλλομεν καὶ ὑμῖν, ἵνα καὶ ὑμεῖς κοινωνίαν ἔχητε μεθ᾽ ἡμῶν. καὶ ἡ κοινωνία δὲ ἡ ἡμετέρα μετὰ τοῦ Πατρὸς καὶ μετὰ τοῦ Υἱοῦ αὐτοῦ Ἰησοῦ Χριστοῦ.
\vs{4}καὶ ταῦτα γράφομεν ἡμεῖς, ἵνα ἡ χαρὰ ἡμῶν ᾖ πεπληρωμένη.

\vs{5}Καὶ ἔστιν αὕτη ἡ ἀγγελία ἣν ἀκηκόαμεν ἀπ᾽ αὐτοῦ καὶ ἀναγγέλλομεν ὑμῖν, ὅτι ὁ Θεὸς φῶς ἐστιν καὶ σκοτία ἐν αὐτῷ οὐκ ἔστιν οὐδεμία.
\vs{6}Ἐὰν εἴπωμεν ὅτι κοινωνίαν ἔχομεν μετ᾽ αὐτοῦ καὶ ἐν τῷ σκότει περιπατῶμεν, ψευδόμεθα καὶ οὐ ποιοῦμεν τὴν ἀλήθειαν·
\vs{7}ἐὰν ἐν τῷ φωτὶ περιπατῶμεν, ὡς αὐτός ἐστιν ἐν τῷ φωτί, κοινωνίαν ἔχομεν μετ᾽ ἀλλήλων, καὶ τὸ αἷμα Ἰησοῦ τοῦ Υἱοῦ αὐτοῦ καθαρίζει ἡμᾶς ἀπὸ πάσης ἁμαρτίας.
\vs{8}Ἐὰν εἴπωμεν ὅτι ἁμαρτίαν οὐκ ἔχομεν, ἑαυτοὺς πλανῶμεν καὶ ἡ ἀλήθεια οὐκ ἔστιν ἐν ἡμῖν.
\vs{9}ἐὰν ὁμολογῶμεν τὰς ἁμαρτίας ἡμῶν, πιστός ἐστιν καὶ δίκαιος, ἵνα ἀφῇ ἡμῖν τὰς ἁμαρτίας καὶ καθαρίσῃ ἡμᾶς ἀπὸ πάσης ἀδικίας.
\vs{10}ἐὰν εἴπωμεν ὅτι οὐχ ἡμαρτήκαμεν, ψεύστην ποιοῦμεν αὐτὸν, καὶ ὁ λόγος αὐτοῦ οὐκ ἔστιν ἐν ἡμῖν.

\ch{2}
Τεκνία μου, ταῦτα γράφω ὑμῖν ἵνα μὴ ἁμάρτητε. καὶ ἐάν τις ἁμάρτῃ, Παράκλητον ἔχομεν πρὸς τὸν Πατέρα Ἰησοῦν Χριστὸν δίκαιον·
\vs{2}καὶ αὐτὸς ἱλασμός ἐστιν περὶ τῶν ἁμαρτιῶν ἡμῶν, οὐ περὶ τῶν ἡμετέρων δὲ μόνον ἀλλὰ καὶ περὶ ὅλου τοῦ κόσμου.

\vs{3}Καὶ ἐν τούτῳ γινώσκομεν ὅτι ἐγνώκαμεν αὐτόν, ἐὰν τὰς ἐντολὰς αὐτοῦ τηρῶμεν.
\vs{4}ὁ λέγων ὅτι Ἔγνωκα αὐτόν καὶ τὰς ἐντολὰς αὐτοῦ μὴ τηρῶν ψεύστης ἐστίν, καὶ ἐν τούτῳ ἡ ἀλήθεια οὐκ ἔστιν·
\vs{5}ὃς δ᾽ ἂν τηρῇ αὐτοῦ τὸν λόγον, ἀληθῶς ἐν τούτῳ ἡ ἀγάπη τοῦ Θεοῦ τετελείωται· Ἐν τούτῳ γινώσκομεν ὅτι ἐν αὐτῷ ἐσμεν.
\vs{6}ὁ λέγων ἐν αὐτῷ μένειν ὀφείλει, καθὼς ἐκεῖνος περιεπάτησεν, καὶ αὐτὸς οὕτως περιπατεῖν.

\vs{7}Ἀγαπητοί, οὐκ ἐντολὴν καινὴν γράφω ὑμῖν ἀλλ᾽ ἐντολὴν παλαιὰν ἣν εἴχετε ἀπ᾽ ἀρχῆς· ἡ ἐντολὴ ἡ παλαιά ἐστιν ὁ λόγος ὃν ἠκούσατε.
\vs{8}πάλιν ἐντολὴν καινὴν γράφω ὑμῖν ὅ ἐστιν ἀληθὲς ἐν αὐτῷ καὶ ἐν ὑμῖν, ὅτι ἡ σκοτία παράγεται καὶ τὸ φῶς τὸ ἀληθινὸν ἤδη φαίνει.
\vs{9}Ὁ λέγων ἐν τῷ φωτὶ εἶναι καὶ τὸν ἀδελφὸν αὐτοῦ μισῶν ἐν τῇ σκοτίᾳ ἐστὶν ἕως ἄρτι.
\vs{10}ὁ ἀγαπῶν τὸν ἀδελφὸν αὐτοῦ ἐν τῷ φωτὶ μένει, καὶ σκάνδαλον ἐν αὐτῷ οὐκ ἔστιν·
\vs{11}ὁ δὲ μισῶν τὸν ἀδελφὸν αὐτοῦ ἐν τῇ σκοτίᾳ ἐστὶν καὶ ἐν τῇ σκοτίᾳ περιπατεῖ καὶ οὐκ οἶδεν ποῦ ὑπάγει, ὅτι ἡ σκοτία ἐτύφλωσεν τοὺς ὀφθαλμοὺς αὐτοῦ.

\vs{12}Γράφω ὑμῖν, τεκνία, ὅτι ἀφέωνται ὑμῖν αἱ ἁμαρτίαι διὰ τὸ ὄνομα αὐτοῦ.
\begin{poetryblock}

\begin{quote} \vs{13}Γράφω ὑμῖν, πατέρες, ὅτι ἐγνώκατε τὸν ἀπ᾽ ἀρχῆς.\end{quote} 

\begin{quote}Γράφω ὑμῖν, νεανίσκοι, ὅτι νενικήκατε τὸν πονηρόν.\end{quote}

\begin{quote} \vs{14}Ἔγραψα ὑμῖν, παιδία, ὅτι ἐγνώκατε τὸν Πατέρα.\end{quote} 

\begin{quote}Ἔγραψα ὑμῖν, πατέρες, ὅτι ἐγνώκατε τὸν ἀπ᾽ ἀρχῆς.\end{quote} 

\begin{quote}Ἔγραψα ὑμῖν, νεανίσκοι, ὅτι ἰσχυροί ἐστε καὶ ὁ λόγος τοῦ Θεοῦ ἐν ὑμῖν μένει καὶ νενικήκατε τὸν πονηρόν.\end{quote}
\end{poetryblock}

\vs{15}Μὴ ἀγαπᾶτε τὸν κόσμον μηδὲ τὰ ἐν τῷ κόσμῳ. ἐάν τις ἀγαπᾷ τὸν κόσμον, οὐκ ἔστιν ἡ ἀγάπη τοῦ Πατρὸς ἐν αὐτῷ·
\vs{16}ὅτι πᾶν τὸ ἐν τῷ κόσμῳ, ἡ ἐπιθυμία τῆς σαρκὸς καὶ ἡ ἐπιθυμία τῶν ὀφθαλμῶν καὶ ἡ ἀλαζονεία τοῦ βίου, οὐκ ἔστιν ἐκ τοῦ πατρός ἀλλὰ ἐκ τοῦ κόσμου ἐστίν.
\vs{17}καὶ ὁ κόσμος παράγεται καὶ ἡ ἐπιθυμία αὐτοῦ, ὁ δὲ ποιῶν τὸ θέλημα τοῦ Θεοῦ μένει εἰς τὸν αἰῶνα.

\vs{18}Παιδία, ἐσχάτη ὥρα ἐστίν, καὶ καθὼς ἠκούσατε ὅτι ἀντίχριστος ἔρχεται, καὶ νῦν ἀντίχριστοι πολλοὶ γεγόνασιν, ὅθεν γινώσκομεν ὅτι ἐσχάτη ὥρα ἐστίν.
\vs{19}ἐξ ἡμῶν ἐξῆλθαν ἀλλ᾽ οὐκ ἦσαν ἐξ ἡμῶν, εἰ γὰρ ἐξ ἡμῶν ἦσαν, μεμενήκεισαν ἂν μεθ᾽ ἡμῶν— ἀλλ᾽ ἵνα φανερωθῶσιν ὅτι οὐκ εἰσὶν πάντες ἐξ ἡμῶν.
\vs{20}Καὶ ὑμεῖς χρῖσμα ἔχετε ἀπὸ τοῦ Ἁγίου καὶ οἴδατε πάντες.
\vs{21}οὐκ ἔγραψα ὑμῖν ὅτι οὐκ οἴδατε τὴν ἀλήθειαν ἀλλ᾽ ὅτι οἴδατε αὐτήν καὶ ὅτι πᾶν ψεῦδος ἐκ τῆς ἀληθείας οὐκ ἔστιν.

\vs{22}Τίς ἐστιν ὁ ψεύστης εἰ μὴ ὁ ἀρνούμενος ὅτι Ἰησοῦς οὐκ ἔστιν ὁ Χριστός; οὗτός ἐστιν ὁ ἀντίχριστος, ὁ ἀρνούμενος τὸν Πατέρα καὶ τὸν Υἱόν.
\vs{23}πᾶς ὁ ἀρνούμενος τὸν Υἱὸν οὐδὲ τὸν Πατέρα ἔχει, ὁ ὁμολογῶν τὸν Υἱὸν καὶ τὸν Πατέρα ἔχει.
\vs{24}Ὑμεῖς ὃ ἠκούσατε ἀπ᾽ ἀρχῆς, ἐν ὑμῖν μενέτω. ἐὰν ἐν ὑμῖν μείνῃ ὃ ἀπ᾽ ἀρχῆς ἠκούσατε, καὶ ὑμεῖς ἐν τῷ Υἱῷ καὶ ἐν τῷ Πατρὶ μενεῖτε.
\vs{25}καὶ αὕτη ἐστὶν ἡ ἐπαγγελία ἣν αὐτὸς ἐπηγγείλατο ἡμῖν, τὴν ζωὴν τὴν αἰώνιον.

\vs{26}Ταῦτα ἔγραψα ὑμῖν περὶ τῶν πλανώντων ὑμᾶς.
\vs{27}καὶ ὑμεῖς τὸ χρῖσμα ὃ ἐλάβετε ἀπ᾽ αὐτοῦ μένει ἐν ὑμῖν, καὶ οὐ χρείαν ἔχετε ἵνα τις διδάσκῃ ὑμᾶς, ἀλλ᾽ ὡς τὸ αὐτοῦ χρῖσμα διδάσκει ὑμᾶς περὶ πάντων, καὶ ἀληθές ἐστιν καὶ οὐκ ἔστιν ψεῦδος, καὶ καθὼς ἐδίδαξεν ὑμᾶς, μένετε ἐν αὐτῷ.

\vs{28}Καὶ νῦν, τεκνία, μένετε ἐν αὐτῷ, ἵνα ἐὰν φανερωθῇ, σχῶμεν παρρησίαν καὶ μὴ αἰσχυνθῶμεν ἀπ᾽ αὐτοῦ ἐν τῇ παρουσίᾳ αὐτοῦ.
\vs{29}Ἐὰν εἰδῆτε ὅτι δίκαιός ἐστιν, γινώσκετε ὅτι καὶ πᾶς ὁ ποιῶν τὴν δικαιοσύνην ἐξ αὐτοῦ γεγέννηται.

\ch{3}
Ἴδετε ποταπὴν ἀγάπην δέδωκεν ἡμῖν ὁ Πατὴρ, ἵνα τέκνα Θεοῦ κληθῶμεν, καὶ ἐσμέν. διὰ τοῦτο ὁ κόσμος οὐ γινώσκει ἡμᾶς, ὅτι οὐκ ἔγνω αὐτόν.
\vs{2}Ἀγαπητοί νῦν τέκνα Θεοῦ ἐσμεν, καὶ οὔπω ἐφανερώθη τί ἐσόμεθα. οἴδαμεν ὅτι ἐὰν φανερωθῇ, ὅμοιοι αὐτῷ ἐσόμεθα, ὅτι ὀψόμεθα αὐτὸν, καθώς ἐστιν.
\vs{3}καὶ πᾶς ὁ ἔχων τὴν ἐλπίδα ταύτην ἐπ᾽ αὐτῷ ἁγνίζει ἑαυτὸν, καθὼς ἐκεῖνος ἁγνός ἐστιν.

\vs{4}Πᾶς ὁ ποιῶν τὴν ἁμαρτίαν καὶ τὴν ἀνομίαν ποιεῖ, καὶ ἡ ἁμαρτία ἐστὶν ἡ ἀνομία.
\vs{5}καὶ οἴδατε ὅτι ἐκεῖνος ἐφανερώθη, ἵνα τὰς ἁμαρτίας ἄρῃ, καὶ ἁμαρτία ἐν αὐτῷ οὐκ ἔστιν.
\vs{6}πᾶς ὁ ἐν αὐτῷ μένων οὐχ ἁμαρτάνει· πᾶς ὁ ἁμαρτάνων οὐχ ἑώρακεν αὐτὸν οὐδὲ ἔγνωκεν αὐτόν.

\vs{7}Τεκνία, μηδεὶς πλανάτω ὑμᾶς· ὁ ποιῶν τὴν δικαιοσύνην δίκαιός ἐστιν, καθὼς ἐκεῖνος δίκαιός ἐστιν·
\vs{8}ὁ ποιῶν τὴν ἁμαρτίαν ἐκ τοῦ διαβόλου ἐστίν, ὅτι ἀπ᾽ ἀρχῆς ὁ διάβολος ἁμαρτάνει. εἰς τοῦτο ἐφανερώθη ὁ Υἱὸς τοῦ Θεοῦ, ἵνα λύσῃ τὰ ἔργα τοῦ διαβόλου.
\vs{9}Πᾶς ὁ γεγεννημένος ἐκ τοῦ Θεοῦ ἁμαρτίαν οὐ ποιεῖ, ὅτι σπέρμα αὐτοῦ ἐν αὐτῷ μένει, καὶ οὐ δύναται ἁμαρτάνειν, ὅτι ἐκ τοῦ Θεοῦ γεγέννηται.
\vs{10}ἐν τούτῳ φανερά ἐστιν τὰ τέκνα τοῦ Θεοῦ καὶ τὰ τέκνα τοῦ διαβόλου· πᾶς ὁ μὴ ποιῶν δικαιοσύνην οὐκ ἔστιν ἐκ τοῦ Θεοῦ καὶ ὁ μὴ ἀγαπῶν τὸν ἀδελφὸν αὐτοῦ.

\vs{11}Ὅτι αὕτη ἐστὶν ἡ ἀγγελία ἣν ἠκούσατε ἀπ᾽ ἀρχῆς, ἵνα ἀγαπῶμεν ἀλλήλους,
\vs{12}οὐ καθὼς Κάϊν ἐκ τοῦ πονηροῦ ἦν καὶ ἔσφαξεν τὸν ἀδελφὸν αὐτοῦ· καὶ χάριν τίνος ἔσφαξεν αὐτόν; ὅτι τὰ ἔργα αὐτοῦ πονηρὰ ἦν, τὰ δὲ τοῦ ἀδελφοῦ αὐτοῦ δίκαια.

\vs{13}καὶ Μὴ θαυμάζετε, ἀδελφοί, εἰ μισεῖ ὑμᾶς ὁ κόσμος.
\vs{14}Ἡμεῖς οἴδαμεν ὅτι μεταβεβήκαμεν ἐκ τοῦ θανάτου εἰς τὴν ζωήν, ὅτι ἀγαπῶμεν τοὺς ἀδελφούς· ὁ μὴ ἀγαπῶν μένει ἐν τῷ θανάτῳ.
\vs{15}πᾶς ὁ μισῶν τὸν ἀδελφὸν αὐτοῦ ἀνθρωποκτόνος ἐστίν, καὶ οἴδατε ὅτι πᾶς ἀνθρωποκτόνος οὐκ ἔχει ζωὴν αἰώνιον ἐν αὐτῷ μένουσαν.
\vs{16}Ἐν τούτῳ ἐγνώκαμεν τὴν ἀγάπην, ὅτι ἐκεῖνος ὑπὲρ ἡμῶν τὴν ψυχὴν αὐτοῦ ἔθηκεν, καὶ ἡμεῖς ὀφείλομεν ὑπὲρ τῶν ἀδελφῶν τὰς ψυχὰς θεῖναι.
\vs{17}ὃς δ᾽ ἂν ἔχῃ τὸν βίον τοῦ κόσμου καὶ θεωρῇ τὸν ἀδελφὸν αὐτοῦ χρείαν ἔχοντα καὶ κλείσῃ τὰ σπλάγχνα αὐτοῦ ἀπ᾽ αὐτοῦ, πῶς ἡ ἀγάπη τοῦ Θεοῦ μένει ἐν αὐτῷ;
\vs{18}Τεκνία, μὴ ἀγαπῶμεν λόγῳ μηδὲ τῇ γλώσσῃ, ἀλλὰ ἐν ἔργῳ καὶ ἀληθείᾳ,

\vs{19}καὶ ἐν τούτῳ γνωσόμεθα ὅτι ἐκ τῆς ἀληθείας ἐσμέν. καὶ ἔμπροσθεν αὐτοῦ πείσομεν τὴν καρδίαν ἡμῶν,
\vs{20}ὅτι ἐὰν καταγινώσκῃ ἡμῶν ἡ καρδία, ὅτι μείζων ἐστὶν ὁ Θεὸς τῆς καρδίας ἡμῶν καὶ γινώσκει πάντα.
\vs{21}Ἀγαπητοί, ἐὰν ἡ καρδία ἡμῶν μὴ καταγινώσκῃ, παρρησίαν ἔχομεν πρὸς τὸν Θεόν
\vs{22}καὶ ὃ ἐὰν αἰτῶμεν, λαμβάνομεν ἀπ᾽ αὐτοῦ, ὅτι τὰς ἐντολὰς αὐτοῦ τηροῦμεν καὶ τὰ ἀρεστὰ ἐνώπιον αὐτοῦ ποιοῦμεν.

\vs{23}καὶ αὕτη ἐστὶν ἡ ἐντολὴ αὐτοῦ, ἵνα πιστεύσωμεν τῷ ὀνόματι τοῦ Υἱοῦ αὐτοῦ Ἰησοῦ Χριστοῦ καὶ ἀγαπῶμεν ἀλλήλους, καθὼς ἔδωκεν ἐντολὴν ἡμῖν.
\vs{24}καὶ ὁ τηρῶν τὰς ἐντολὰς αὐτοῦ ἐν αὐτῷ μένει καὶ αὐτὸς ἐν αὐτῷ· καὶ ἐν τούτῳ γινώσκομεν ὅτι μένει ἐν ἡμῖν, ἐκ τοῦ Πνεύματος οὗ ἡμῖν ἔδωκεν.

\ch{4}
Ἀγαπητοί, μὴ παντὶ πνεύματι πιστεύετε ἀλλὰ δοκιμάζετε τὰ πνεύματα εἰ ἐκ τοῦ Θεοῦ ἐστιν, ὅτι πολλοὶ ψευδοπροφῆται ἐξεληλύθασιν εἰς τὸν κόσμον.
\vs{2}Ἐν τούτῳ γινώσκετε τὸ Πνεῦμα τοῦ Θεοῦ· πᾶν πνεῦμα ὃ ὁμολογεῖ Ἰησοῦν Χριστὸν ἐν σαρκὶ ἐληλυθότα ἐκ τοῦ Θεοῦ ἐστιν,
\vs{3}καὶ πᾶν πνεῦμα ὃ μὴ ὁμολογεῖ τὸν Ἰησοῦν ἐκ τοῦ Θεοῦ οὐκ ἔστιν· καὶ τοῦτό ἐστιν τὸ τοῦ ἀντιχρίστου ὃ ἀκηκόατε ὅτι ἔρχεται, καὶ νῦν ἐν τῷ κόσμῳ ἐστὶν ἤδη.

\vs{4}Ὑμεῖς ἐκ τοῦ Θεοῦ ἐστε, τεκνία, καὶ νενικήκατε αὐτούς, ὅτι μείζων ἐστὶν ὁ ἐν ὑμῖν ἢ ὁ ἐν τῷ κόσμῳ.
\vs{5}αὐτοὶ ἐκ τοῦ κόσμου εἰσίν, διὰ τοῦτο ἐκ τοῦ κόσμου λαλοῦσιν καὶ ὁ κόσμος αὐτῶν ἀκούει.
\vs{6}ἡμεῖς ἐκ τοῦ Θεοῦ ἐσμεν· ὁ γινώσκων τὸν Θεὸν ἀκούει ἡμῶν· ὃς οὐκ ἔστιν ἐκ τοῦ Θεοῦ οὐκ ἀκούει ἡμῶν. ἐκ τούτου γινώσκομεν τὸ πνεῦμα τῆς ἀληθείας καὶ τὸ πνεῦμα τῆς πλάνης.

\vs{7}Ἀγαπητοί, ἀγαπῶμεν ἀλλήλους, 
\begin{poetryblock}

\begin{quote}ὅτι ἡ ἀγάπη ἐκ τοῦ Θεοῦ ἐστιν,\end{quote} 

\begin{quote}καὶ πᾶς ὁ ἀγαπῶν ἐκ τοῦ Θεοῦ γεγέννηται\end{quote} 

\begin{quote}καὶ γινώσκει τὸν Θεόν.\end{quote}

\begin{quote} \vs{8}ὁ μὴ ἀγαπῶν οὐκ ἔγνω τὸν Θεόν, ὅτι ὁ Θεὸς ἀγάπη ἐστίν.\end{quote}

\begin{quote} \vs{9}Ἐν τούτῳ ἐφανερώθη ἡ ἀγάπη τοῦ Θεοῦ ἐν ἡμῖν,\end{quote} 

\begin{quote}ὅτι τὸν Υἱὸν αὐτοῦ τὸν μονογενῆ ἀπέσταλκεν ὁ Θεὸς\end{quote} 

\begin{quote}εἰς τὸν κόσμον, ἵνα ζήσωμεν δι᾽ αὐτοῦ.\end{quote}

\begin{quote} \vs{10}ἐν τούτῳ ἐστὶν ἡ ἀγάπη, οὐχ ὅτι ἡμεῖς ἠγαπήκαμεν τὸν Θεόν, ἀλλ᾽ ὅτι αὐτὸς ἠγάπησεν ἡμᾶς καὶ ἀπέστειλεν τὸν Υἱὸν αὐτοῦ ἱλασμὸν περὶ τῶν ἁμαρτιῶν ἡμῶν.\end{quote}
\end{poetryblock}
\vs{11}Ἀγαπητοί, εἰ οὕτως ὁ Θεὸς ἠγάπησεν ἡμᾶς, καὶ ἡμεῖς ὀφείλομεν ἀλλήλους ἀγαπᾶν.
\vs{12}Θεὸν οὐδεὶς πώποτε τεθέαται. ἐὰν ἀγαπῶμεν ἀλλήλους, ὁ Θεὸς ἐν ἡμῖν μένει καὶ ἡ ἀγάπη αὐτοῦ ἐν ἡμῖν τετελειωμένη ἐστιν.
\vs{13}ἐν τούτῳ γινώσκομεν ὅτι ἐν αὐτῷ μένομεν καὶ αὐτὸς ἐν ἡμῖν, ὅτι ἐκ τοῦ Πνεύματος αὐτοῦ δέδωκεν ἡμῖν.
\vs{14}καὶ ἡμεῖς τεθεάμεθα καὶ μαρτυροῦμεν ὅτι ὁ Πατὴρ ἀπέσταλκεν τὸν Υἱὸν Σωτῆρα τοῦ κόσμου.
\vs{15}Ὃς ἐὰν ὁμολογήσῃ ὅτι Ἰησοῦς ἐστιν ὁ Υἱὸς τοῦ Θεοῦ, ὁ Θεὸς ἐν αὐτῷ μένει καὶ αὐτὸς ἐν τῷ Θεῷ.
\vs{16}καὶ ἡμεῖς ἐγνώκαμεν καὶ πεπιστεύκαμεν τὴν ἀγάπην ἣν ἔχει ὁ Θεὸς ἐν ἡμῖν.

Ὁ Θεὸς ἀγάπη ἐστίν, καὶ ὁ μένων ἐν τῇ ἀγάπῃ ἐν τῷ Θεῷ μένει, καὶ ὁ Θεὸς ἐν αὐτῷ μένει.
\vs{17}Ἐν τούτῳ τετελείωται ἡ ἀγάπη μεθ᾽ ἡμῶν, ἵνα παρρησίαν ἔχωμεν ἐν τῇ ἡμέρᾳ τῆς κρίσεως, ὅτι καθὼς ἐκεῖνός ἐστιν, καὶ ἡμεῖς ἐσμεν ἐν τῷ κόσμῳ τούτῳ.
\vs{18}Φόβος οὐκ ἔστιν ἐν τῇ ἀγάπῃ, ἀλλ᾽ ἡ τελεία ἀγάπη ἔξω βάλλει τὸν φόβον, ὅτι ὁ φόβος κόλασιν ἔχει, ὁ δὲ φοβούμενος οὐ τετελείωται ἐν τῇ ἀγάπῃ.
\vs{19}Ἡμεῖς ἀγαπῶμεν, ὅτι αὐτὸς πρῶτος ἠγάπησεν ἡμᾶς.
\vs{20}Ἐάν τις εἴπῃ ὅτι Ἀγαπῶ τὸν Θεόν καὶ τὸν ἀδελφὸν αὐτοῦ μισῇ, ψεύστης ἐστίν· ὁ γὰρ μὴ ἀγαπῶν τὸν ἀδελφὸν αὐτοῦ ὃν ἑώρακεν, τὸν Θεὸν ὃν οὐχ ἑώρακεν οὐ δύναται ἀγαπᾶν.
\vs{21}καὶ ταύτην τὴν ἐντολὴν ἔχομεν ἀπ᾽ αὐτοῦ, ἵνα ὁ ἀγαπῶν τὸν Θεὸν ἀγαπᾷ καὶ τὸν ἀδελφὸν αὐτοῦ.

\ch{5}
Πᾶς ὁ πιστεύων ὅτι Ἰησοῦς ἐστιν ὁ Χριστὸς ἐκ τοῦ Θεοῦ γεγέννηται, καὶ πᾶς ὁ ἀγαπῶν τὸν γεννήσαντα ἀγαπᾷ καὶ τὸν γεγεννημένον ἐξ αὐτοῦ.
\vs{2}ἐν τούτῳ γινώσκομεν ὅτι ἀγαπῶμεν τὰ τέκνα τοῦ Θεοῦ, ὅταν τὸν Θεὸν ἀγαπῶμεν καὶ τὰς ἐντολὰς αὐτοῦ ποιῶμεν.
\vs{3}αὕτη γάρ ἐστιν ἡ ἀγάπη τοῦ Θεοῦ, ἵνα τὰς ἐντολὰς αὐτοῦ τηρῶμεν, καὶ αἱ ἐντολαὶ αὐτοῦ βαρεῖαι οὐκ εἰσίν.
\vs{4}ὅτι πᾶν τὸ γεγεννημένον ἐκ τοῦ Θεοῦ νικᾷ τὸν κόσμον· καὶ αὕτη ἐστὶν ἡ νίκη ἡ νικήσασα τὸν κόσμον, ἡ πίστις ἡμῶν.

\vs{5}Τίς δέ ἐστιν ὁ νικῶν τὸν κόσμον εἰ μὴ ὁ πιστεύων ὅτι Ἰησοῦς ἐστιν ὁ Υἱὸς τοῦ Θεοῦ;
\vs{6}οὗτός ἐστιν ὁ ἐλθὼν δι᾽ ὕδατος καὶ αἵματος, Ἰησοῦς Χριστός, οὐκ ἐν τῷ ὕδατι μόνον, ἀλλ᾽ ἐν τῷ ὕδατι καὶ ἐν τῷ αἵματι· καὶ τὸ Πνεῦμά ἐστιν τὸ μαρτυροῦν, ὅτι τὸ Πνεῦμά ἐστιν ἡ ἀλήθεια.
\vs{7}ὅτι τρεῖς εἰσιν οἱ μαρτυροῦντες,
\vs{8}τὸ Πνεῦμα καὶ τὸ ὕδωρ καὶ τὸ αἷμα, καὶ οἱ τρεῖς εἰς τὸ ἕν εἰσιν.
\vs{9}Εἰ τὴν μαρτυρίαν τῶν ἀνθρώπων λαμβάνομεν, ἡ μαρτυρία τοῦ Θεοῦ μείζων ἐστίν· ὅτι αὕτη ἐστὶν ἡ μαρτυρία τοῦ Θεοῦ, ὅτι μεμαρτύρηκεν περὶ τοῦ Υἱοῦ αὐτοῦ.
\vs{10}ὁ πιστεύων εἰς τὸν Υἱὸν τοῦ Θεοῦ ἔχει τὴν μαρτυρίαν ἐν αὑτῷ, ὁ μὴ πιστεύων τῷ Θεῷ ψεύστην πεποίηκεν αὐτόν, ὅτι οὐ πεπίστευκεν εἰς τὴν μαρτυρίαν ἣν μεμαρτύρηκεν ὁ Θεὸς περὶ τοῦ Υἱοῦ αὐτοῦ.
\vs{11}Καὶ αὕτη ἐστὶν ἡ μαρτυρία, ὅτι ζωὴν αἰώνιον ἔδωκεν ἡμῖν ὁ Θεὸς, καὶ αὕτη ἡ ζωὴ ἐν τῷ Υἱῷ αὐτοῦ ἐστιν.
\vs{12}ὁ ἔχων τὸν Υἱὸν ἔχει τὴν ζωήν· ὁ μὴ ἔχων τὸν Υἱὸν τοῦ Θεοῦ τὴν ζωὴν οὐκ ἔχει.

\vs{13}Ταῦτα ἔγραψα ὑμῖν, ἵνα εἰδῆτε ὅτι ζωὴν ἔχετε αἰώνιον, τοῖς πιστεύουσιν εἰς τὸ ὄνομα τοῦ Υἱοῦ τοῦ Θεοῦ.
\vs{14}Καὶ αὕτη ἐστὶν ἡ παρρησία ἣν ἔχομεν πρὸς αὐτόν, ὅτι ἐάν τι αἰτώμεθα κατὰ τὸ θέλημα αὐτοῦ ἀκούει ἡμῶν.
\vs{15}καὶ ἐὰν οἴδαμεν ὅτι ἀκούει ἡμῶν ὃ ἐὰν αἰτώμεθα, οἴδαμεν ὅτι ἔχομεν τὰ αἰτήματα ἃ ᾐτήκαμεν ἀπ᾽ αὐτοῦ.

\vs{16}Ἐάν τις ἴδῃ τὸν ἀδελφὸν αὐτοῦ ἁμαρτάνοντα ἁμαρτίαν μὴ πρὸς θάνατον, αἰτήσει καὶ δώσει αὐτῷ ζωήν, τοῖς ἁμαρτάνουσιν μὴ πρὸς θάνατον. ἔστιν ἁμαρτία πρὸς θάνατον· οὐ περὶ ἐκείνης λέγω ἵνα ἐρωτήσῃ.
\vs{17}πᾶσα ἀδικία ἁμαρτία ἐστίν, καὶ ἔστιν ἁμαρτία οὐ πρὸς θάνατον.

\vs{18}Οἴδαμεν ὅτι πᾶς ὁ γεγεννημένος ἐκ τοῦ Θεοῦ οὐχ ἁμαρτάνει, ἀλλ᾽ ὁ γεννηθεὶς ἐκ τοῦ Θεοῦ τηρεῖ ἑαυτὸν καὶ ὁ πονηρὸς οὐχ ἅπτεται αὐτοῦ.
\vs{19}οἴδαμεν ὅτι ἐκ τοῦ Θεοῦ ἐσμεν καὶ ὁ κόσμος ὅλος ἐν τῷ πονηρῷ κεῖται.
\vs{20}οἴδαμεν δὲ ὅτι ὁ Υἱὸς τοῦ Θεοῦ ἥκει καὶ δέδωκεν ἡμῖν διάνοιαν, ἵνα γινώσκωμεν τὸν ἀληθινόν, καὶ ἐσμὲν ἐν τῷ ἀληθινῷ, ἐν τῷ Υἱῷ αὐτοῦ Ἰησοῦ Χριστῷ. οὗτός ἐστιν ὁ ἀληθινὸς Θεὸς καὶ ζωὴ αἰώνιος.

\vs{21}Τεκνία, φυλάξατε ἑαυτὰ ἀπὸ τῶν εἰδώλων.


\def\book{ΙΩΑΝΝΟΥ Β}
\biblebook{ΙΩΑΝΝΟΥ Β}

 
\lettrine[lines=2, loversize=0.2, nindent=0em, findent=.25em]{\textcolor{bookheadingcolor}{Ὁ}}{} πρεσβύτερος Ἐκλεκτῇ κυρίᾳ καὶ τοῖς τέκνοις αὐτῆς, οὓς ἐγὼ ἀγαπῶ ἐν ἀληθείᾳ, καὶ οὐκ ἐγὼ μόνος ἀλλὰ καὶ πάντες οἱ ἐγνωκότες τὴν ἀλήθειαν,
\vs{2}διὰ τὴν ἀλήθειαν τὴν μένουσαν ἐν ἡμῖν καὶ μεθ᾽ ἡμῶν ἔσται εἰς τὸν αἰῶνα.
\vs{3}Ἔσται μεθ᾽ ἡμῶν χάρις ἔλεος εἰρήνη παρὰ Θεοῦ Πατρός καὶ παρὰ Ἰησοῦ Χριστοῦ τοῦ Υἱοῦ τοῦ Πατρός ἐν ἀληθείᾳ καὶ ἀγάπῃ.

\vs{4}Ἐχάρην λίαν ὅτι εὕρηκα ἐκ τῶν τέκνων σου περιπατοῦντας ἐν ἀληθείᾳ, καθὼς ἐντολὴν ἐλάβομεν παρὰ τοῦ Πατρός.
\vs{5}καὶ νῦν ἐρωτῶ σε, κυρία, οὐχ ὡς ἐντολὴν γράφων σοι καινὴν ἀλλὰ ἣν εἴχομεν ἀπ᾽ ἀρχῆς, ἵνα ἀγαπῶμεν ἀλλήλους.
\vs{6}καὶ αὕτη ἐστὶν ἡ ἀγάπη, ἵνα περιπατῶμεν κατὰ τὰς ἐντολὰς αὐτοῦ· αὕτη ἡ ἐντολή ἐστιν, καθὼς ἠκούσατε ἀπ᾽ ἀρχῆς, ἵνα ἐν αὐτῇ περιπατῆτε.

\vs{7}Ὅτι πολλοὶ πλάνοι ἐξῆλθον εἰς τὸν κόσμον, οἱ μὴ ὁμολογοῦντες Ἰησοῦν Χριστὸν ἐρχόμενον ἐν σαρκί· οὗτός ἐστιν ὁ πλάνος καὶ ὁ ἀντίχριστος.
\vs{8}βλέπετε ἑαυτούς, ἵνα μὴ ἀπολέσητε ἃ εἰργασάμεθα ἀλλὰ μισθὸν πλήρη ἀπολάβητε.

\vs{9}πᾶς ὁ προάγων καὶ μὴ μένων ἐν τῇ διδαχῇ τοῦ Χριστοῦ Θεὸν οὐκ ἔχει· ὁ μένων ἐν τῇ διδαχῇ, οὗτος καὶ τὸν Πατέρα καὶ τὸν Υἱὸν ἔχει.
\vs{10}Εἴ τις ἔρχεται πρὸς ὑμᾶς καὶ ταύτην τὴν διδαχὴν οὐ φέρει, μὴ λαμβάνετε αὐτὸν εἰς οἰκίαν καὶ χαίρειν αὐτῷ μὴ λέγετε·
\vs{11}ὁ λέγων γὰρ αὐτῷ χαίρειν κοινωνεῖ τοῖς ἔργοις αὐτοῦ τοῖς πονηροῖς.

\vs{12}Πολλὰ ἔχων ὑμῖν γράφειν οὐκ ἐβουλήθην διὰ χάρτου καὶ μέλανος, ἀλλὰ ἐλπίζω γενέσθαι πρὸς ὑμᾶς καὶ στόμα πρὸς στόμα λαλῆσαι, ἵνα ἡ χαρὰ ἡμῶν ᾖ πεπληρωμένη.

\vs{13}Ἀσπάζεταί σε τὰ τέκνα τῆς ἀδελφῆς σου τῆς ἐκλεκτῆς.


\def\book{ΙΩΑΝΝΟΥ Γ}
\biblebook{ΙΩΑΝΝΟΥ Γ}

 
\lettrine[lines=2, loversize=0.2, nindent=0em, findent=.25em]{\textcolor{bookheadingcolor}{Ὁ}}{} πρεσβύτερος Γαΐῳ τῷ ἀγαπητῷ, ὃν ἐγὼ ἀγαπῶ ἐν ἀληθείᾳ.

\postdropcapindent\vs{2}Ἀγαπητέ, περὶ πάντων εὔχομαί σε εὐοδοῦσθαι καὶ ὑγιαίνειν, καθὼς εὐοδοῦταί σου ἡ ψυχή.
\vs{3}ἐχάρην γὰρ λίαν ἐρχομένων ἀδελφῶν καὶ μαρτυρούντων σου τῇ ἀληθείᾳ, καθὼς σὺ ἐν ἀληθείᾳ περιπατεῖς
\vs{4}μειζοτέραν τούτων οὐκ ἔχω χαράν, ἵνα ἀκούω τὰ ἐμὰ τέκνα ἐν ἀληθείᾳ περιπατοῦντα.

\vs{5}Ἀγαπητέ, πιστὸν ποιεῖς ὃ ἐὰν ἐργάσῃ εἰς τοὺς ἀδελφοὺς καὶ τοῦτο ξένους,
\vs{6}οἳ ἐμαρτύρησάν σου τῇ ἀγάπῃ ἐνώπιον ἐκκλησίας, οὓς καλῶς ποιήσεις προπέμψας ἀξίως τοῦ Θεοῦ·
\vs{7}ὑπὲρ γὰρ τοῦ Ὀνόματος ἐξῆλθον μηδὲν λαμβάνοντες ἀπὸ τῶν ἐθνικῶν.
\vs{8}ἡμεῖς οὖν ὀφείλομεν ὑπολαμβάνειν τοὺς τοιούτους, ἵνα συνεργοὶ γινώμεθα τῇ ἀληθείᾳ.

\vs{9}Ἔγραψά τι τῇ ἐκκλησίᾳ· ἀλλ᾽ ὁ φιλοπρωτεύων αὐτῶν Διοτρεφὴς οὐκ ἐπιδέχεται ἡμᾶς.
\vs{10}διὰ τοῦτο, ἐὰν ἔλθω, ὑπομνήσω αὐτοῦ τὰ ἔργα ἃ ποιεῖ λόγοις πονηροῖς φλυαρῶν ἡμᾶς, καὶ μὴ ἀρκούμενος ἐπὶ τούτοις οὔτε αὐτὸς ἐπιδέχεται τοὺς ἀδελφοὺς καὶ τοὺς βουλομένους κωλύει καὶ ἐκ τῆς ἐκκλησίας ἐκβάλλει.

\vs{11}Ἀγαπητέ, μὴ μιμοῦ τὸ κακὸν ἀλλὰ τὸ ἀγαθόν. ὁ ἀγαθοποιῶν ἐκ τοῦ Θεοῦ ἐστιν· ὁ κακοποιῶν οὐχ ἑώρακεν τὸν Θεόν.
\vs{12}Δημητρίῳ μεμαρτύρηται ὑπὸ πάντων καὶ ὑπὸ αὐτῆς τῆς ἀληθείας· καὶ ἡμεῖς δὲ μαρτυροῦμεν, καὶ οἶδας ὅτι ἡ μαρτυρία ἡμῶν ἀληθής ἐστιν.

\vs{13}Πολλὰ εἶχον γράψαι σοι ἀλλ᾽ οὐ θέλω διὰ μέλανος καὶ καλάμου σοι γράφειν·
\vs{14}ἐλπίζω δὲ εὐθέως σε ἰδεῖν, καὶ στόμα πρὸς στόμα λαλήσομεν.

\vs{15}Εἰρήνη σοι. Ἀσπάζονταί σε οἱ φίλοι. Ἀσπάζου τοὺς φίλους κατ᾽ ὄνομα.


\def\book{ΙΟΥΔΑ}
\biblebook{ΙΟΥΔΑ}

 
\lettrine[lines=2, loversize=0.2, nindent=0em, findent=.25em]{\textcolor{bookheadingcolor}{Ἰ}}{ούδας} Ἰησοῦ Χριστοῦ δοῦλος, ἀδελφὸς δὲ Ἰακώβου, Τοῖς ἐν Θεῷ Πατρὶ ἠγαπημένοις καὶ Ἰησοῦ Χριστῷ τετηρημένοις κλητοῖς·
\vs{2}Ἔλεος ὑμῖν καὶ εἰρήνη καὶ ἀγάπη πληθυνθείη.

\vs{3}Ἀγαπητοί, πᾶσαν σπουδὴν ποιούμενος γράφειν ὑμῖν περὶ τῆς κοινῆς ἡμῶν σωτηρίας ἀνάγκην ἔσχον γράψαι ὑμῖν παρακαλῶν ἐπαγωνίζεσθαι τῇ ἅπαξ παραδοθείσῃ τοῖς ἁγίοις πίστει.
\vs{4}παρεισέδυσαν γάρ τινες ἄνθρωποι, οἱ πάλαι προγεγραμμένοι εἰς τοῦτο τὸ κρίμα, ἀσεβεῖς, τὴν τοῦ Θεοῦ ἡμῶν χάριτα μετατιθέντες εἰς ἀσέλγειαν καὶ τὸν μόνον Δεσπότην καὶ Κύριον ἡμῶν Ἰησοῦν Χριστὸν ἀρνούμενοι.

\vs{5}Ὑπομνῆσαι δὲ ὑμᾶς βούλομαι, εἰδότας ὑμᾶς ἅπαξ πάντα ὅτι Ἰησοῦς λαὸν ἐκ γῆς Αἰγύπτου σώσας τὸ δεύτερον τοὺς μὴ πιστεύσαντας ἀπώλεσεν,
\vs{6}ἀγγέλους τε τοὺς μὴ τηρήσαντας τὴν ἑαυτῶν ἀρχὴν ἀλλὰ ἀπολιπόντας τὸ ἴδιον οἰκητήριον εἰς κρίσιν μεγάλης ἡμέρας δεσμοῖς ἀϊδίοις ὑπὸ ζόφον τετήρηκεν,
\vs{7}ὡς Σόδομα καὶ Γόμορρα καὶ αἱ περὶ αὐτὰς πόλεις τὸν ὅμοιον τρόπον τούτοις ἐκπορνεύσασαι καὶ ἀπελθοῦσαι ὀπίσω σαρκὸς ἑτέρας, πρόκεινται δεῖγμα πυρὸς αἰωνίου δίκην ὑπέχουσαι.

\vs{8}Ὁμοίως μέντοι καὶ οὗτοι ἐνυπνιαζόμενοι σάρκα μὲν μιαίνουσιν κυριότητα δὲ ἀθετοῦσιν δόξας δὲ βλασφημοῦσιν
\vs{9}ὁ δὲ Μιχαὴλ ὁ ἀρχάγγελος, ὅτε τῷ διαβόλῳ διακρινόμενος διελέγετο περὶ τοῦ Μωϋσέως σώματος, οὐκ ἐτόλμησεν κρίσιν ἐπενεγκεῖν βλασφημίας ἀλλὰ εἶπεν, Ἐπιτιμήσαι σοι Κύριος.
\vs{10}οὗτοι δὲ ὅσα μὲν οὐκ οἴδασιν βλασφημοῦσιν, ὅσα δὲ φυσικῶς ὡς τὰ ἄλογα ζῷα ἐπίστανται, ἐν τούτοις φθείρονται.
\vs{11}οὐαὶ αὐτοῖς, ὅτι τῇ ὁδῷ τοῦ Κάϊν ἐπορεύθησαν καὶ τῇ πλάνῃ τοῦ Βαλαὰμ μισθοῦ ἐξεχύθησαν καὶ τῇ ἀντιλογίᾳ τοῦ Κόρε ἀπώλοντο.
\vs{12}Οὗτοί εἰσιν οἱ ἐν ταῖς ἀγάπαις ὑμῶν σπιλάδες συνευωχούμενοι ἀφόβως, ἑαυτοὺς ποιμαίνοντες, νεφέλαι ἄνυδροι ὑπὸ ἀνέμων παραφερόμεναι, δένδρα φθινοπωρινὰ ἄκαρπα δὶς ἀποθανόντα ἐκριζωθέντα,
\vs{13}κύματα ἄγρια θαλάσσης ἐπαφρίζοντα τὰς ἑαυτῶν αἰσχύνας, ἀστέρες πλανῆται οἷς ὁ ζόφος τοῦ σκότους εἰς αἰῶνα τετήρηται.

\vs{14}Προεφήτευσεν δὲ καὶ τούτοις ἕβδομος ἀπὸ Ἀδὰμ Ἑνὼχ λέγων, Ἰδοὺ ἦλθεν Κύριος ἐν ἁγίαις μυριάσιν αὐτοῦ
\vs{15}ποιῆσαι κρίσιν κατὰ πάντων καὶ ἐλέγξαι πᾶσαν ψυχὴν περὶ πάντων τῶν ἔργων ἀσεβείας αὐτῶν ὧν ἠσέβησαν καὶ περὶ πάντων τῶν σκληρῶν ὧν ἐλάλησαν κατ᾽ αὐτοῦ ἁμαρτωλοὶ ἀσεβεῖς.
\vs{16}Οὗτοί εἰσιν γογγυσταί μεμψίμοιροι κατὰ τὰς ἐπιθυμίας αὐτῶν πορευόμενοι, καὶ τὸ στόμα αὐτῶν λαλεῖ ὑπέρογκα, θαυμάζοντες πρόσωπα ὠφελείας χάριν.

\vs{17}Ὑμεῖς δέ, ἀγαπητοί, μνήσθητε τῶν ῥημάτων τῶν προειρημένων ὑπὸ τῶν ἀποστόλων τοῦ Κυρίου ἡμῶν Ἰησοῦ Χριστοῦ
\vs{18}ὅτι ἔλεγον ὑμῖν Ἐπ᾽ ἐσχάτου χρόνου ἔσονται ἐμπαῖκται κατὰ τὰς ἑαυτῶν ἐπιθυμίας πορευόμενοι τῶν ἀσεβειῶν.
\vs{19}Οὗτοί εἰσιν οἱ ἀποδιορίζοντες, ψυχικοί, Πνεῦμα μὴ ἔχοντες.

\vs{20}Ὑμεῖς δέ, ἀγαπητοί, ἐποικοδομοῦντες ἑαυτοὺς τῇ ἁγιωτάτῃ ὑμῶν πίστει, ἐν Πνεύματι Ἁγίῳ προσευχόμενοι,
\vs{21}ἑαυτοὺς ἐν ἀγάπῃ Θεοῦ τηρήσατε προσδεχόμενοι τὸ ἔλεος τοῦ Κυρίου ἡμῶν Ἰησοῦ Χριστοῦ εἰς ζωὴν αἰώνιον.
\vs{22}Καὶ οὓς μὲν ἐλεᾶτε διακρινομένους,
\vs{23}οὓς δὲ σῴζετε ἐκ πυρὸς ἁρπάζοντες, οὓς δὲ ἐλεᾶτε ἐν φόβῳ μισοῦντες καὶ τὸν ἀπὸ τῆς σαρκὸς ἐσπιλωμένον χιτῶνα.

\vs{24}Τῷ δὲ δυναμένῳ φυλάξαι ὑμᾶς ἀπταίστους καὶ στῆσαι κατενώπιον τῆς δόξης αὐτοῦ ἀμώμους ἐν ἀγαλλιάσει,
\vs{25}μόνῳ Θεῷ Σωτῆρι ἡμῶν διὰ Ἰησοῦ Χριστοῦ τοῦ Κυρίου ἡμῶν δόξα μεγαλωσύνη κράτος καὶ ἐξουσία πρὸ παντὸς τοῦ αἰῶνος καὶ νῦν καὶ εἰς πάντας τοὺς αἰῶνας, Ἀμήν.


\def\book{ΑΠΟΚΑΛΥΨΙΣ ΙΩΑΝΝΟΥ}
\biblebook{ΑΠΟΚΑΛΥΨΙΣ ΙΩΑΝΝΟΥ}


\lettrine[lines=2, loversize=0.2, nindent=0em, findent=.25em]{\textcolor{bookheadingcolor}{Ἀ}}{ποκάλυψις} Ἰησοῦ Χριστοῦ ἣν ἔδωκεν αὐτῷ ὁ Θεός δεῖξαι τοῖς δούλοις αὐτοῦ ἃ δεῖ γενέσθαι ἐν τάχει, καὶ ἐσήμανεν ἀποστείλας διὰ τοῦ ἀγγέλου αὐτοῦ τῷ δούλῳ αὐτοῦ Ἰωάννῃ,
\vs{2}ὃς ἐμαρτύρησεν τὸν λόγον τοῦ Θεοῦ καὶ τὴν μαρτυρίαν Ἰησοῦ Χριστοῦ ὅσα εἶδεν.
\vs{3}Μακάριος ὁ ἀναγινώσκων καὶ οἱ ἀκούοντες τοὺς λόγους τῆς προφητείας καὶ τηροῦντες τὰ ἐν αὐτῇ γεγραμμένα, ὁ γὰρ καιρὸς ἐγγύς.

\vs{4}Ἰωάννης Ταῖς ἑπτὰ ἐκκλησίαις ταῖς ἐν τῇ Ἀσίᾳ· Χάρις ὑμῖν καὶ εἰρήνη ἀπὸ ὁ ὢν καὶ ὁ ἦν καὶ ὁ ἐρχόμενος καὶ ἀπὸ τῶν ἑπτὰ Πνευμάτων ἃ ἐνώπιον τοῦ θρόνου αὐτοῦ
\vs{5}καὶ ἀπὸ Ἰησοῦ Χριστοῦ, ὁ μάρτυς, ὁ πιστός, ὁ πρωτότοκος τῶν νεκρῶν καὶ ὁ ἄρχων τῶν βασιλέων τῆς γῆς. Τῷ ἀγαπῶντι ἡμᾶς καὶ λύσαντι ἡμᾶς ἐκ τῶν ἁμαρτιῶν ἡμῶν ἐν τῷ αἵματι αὐτοῦ,
\vs{6}καὶ ἐποίησεν ἡμᾶς βασιλείαν, ἱερεῖς τῷ Θεῷ καὶ Πατρὶ αὐτοῦ, αὐτῷ ἡ δόξα καὶ τὸ κράτος εἰς τοὺς αἰῶνας τῶν αἰώνων· ἀμήν.

\vs{7}Ἰδοὺ ἔρχεται μετὰ τῶν νεφελῶν, 
\begin{poetryblock}

\begin{quote}καὶ ὄψεται αὐτὸν πᾶς ὀφθαλμὸς\end{quote} 

\begin{quote}καὶ οἵτινες αὐτὸν ἐξεκέντησαν,\end{quote} 

\begin{quote}καὶ κόψονται ἐπ᾽ αὐτὸν πᾶσαι αἱ φυλαὶ τῆς γῆς.\end{quote}
\end{poetryblock}

ναί, ἀμήν.

\vs{8}Ἐγώ εἰμι τὸ Ἄλφα καὶ τὸ Ὦ, λέγει Κύριος ὁ Θεός, ὁ ὢν καὶ ὁ ἦν καὶ ὁ ἐρχόμενος, ὁ Παντοκράτωρ.

\vs{9}Ἐγὼ Ἰωάννης, ὁ ἀδελφὸς ὑμῶν καὶ συνκοινωνὸς ἐν τῇ θλίψει καὶ βασιλείᾳ καὶ ὑπομονῇ ἐν Ἰησοῦ, ἐγενόμην ἐν τῇ νήσῳ τῇ καλουμένῃ Πάτμῳ διὰ τὸν λόγον τοῦ Θεοῦ καὶ τὴν μαρτυρίαν Ἰησοῦ.
\vs{10}ἐγενόμην ἐν Πνεύματι ἐν τῇ κυριακῇ ἡμέρᾳ καὶ ἤκουσα ὀπίσω μου φωνὴν μεγάλην ὡς σάλπιγγος
\vs{11}λεγούσης· Ὃ βλέπεις γράψον εἰς βιβλίον καὶ πέμψον ταῖς ἑπτὰ ἐκκλησίαις, εἰς Ἔφεσον καὶ εἰς Σμύρναν καὶ εἰς Πέργαμον καὶ εἰς Θυάτειρα καὶ εἰς Σάρδεις καὶ εἰς Φιλαδέλφειαν καὶ εἰς Λαοδίκειαν.

\vs{12}Καὶ ἐπέστρεψα βλέπειν τὴν φωνὴν ἥτις ἐλάλει μετ᾽ ἐμοῦ, καὶ ἐπιστρέψας εἶδον ἑπτὰ λυχνίας χρυσᾶς
\vs{13}καὶ ἐν μέσῳ τῶν λυχνιῶν ὅμοιον υἱὸν ἀνθρώπου ἐνδεδυμένον ποδήρη καὶ περιεζωσμένον πρὸς τοῖς μαστοῖς ζώνην χρυσᾶν.
\vs{14}ἡ δὲ κεφαλὴ αὐτοῦ καὶ αἱ τρίχες λευκαὶ ὡς ἔριον λευκόν ὡς χιών καὶ οἱ ὀφθαλμοὶ αὐτοῦ ὡς φλὸξ πυρός
\vs{15}καὶ οἱ πόδες αὐτοῦ ὅμοιοι χαλκολιβάνῳ ὡς ἐν καμίνῳ πεπυρωμένης καὶ ἡ φωνὴ αὐτοῦ ὡς φωνὴ ὑδάτων πολλῶν,
\vs{16}καὶ ἔχων ἐν τῇ δεξιᾷ χειρὶ αὐτοῦ ἀστέρας ἑπτά καὶ ἐκ τοῦ στόματος αὐτοῦ ῥομφαία δίστομος ὀξεῖα ἐκπορευομένη καὶ ἡ ὄψις αὐτοῦ ὡς ὁ ἥλιος φαίνει ἐν τῇ δυνάμει αὐτοῦ.

\vs{17}Καὶ ὅτε εἶδον αὐτόν, ἔπεσα πρὸς τοὺς πόδας αὐτοῦ ὡς νεκρός, καὶ ἔθηκεν τὴν δεξιὰν αὐτοῦ ἐπ᾽ ἐμὲ λέγων·

Μὴ φοβοῦ· ἐγώ εἰμι ὁ πρῶτος καὶ ὁ ἔσχατος
\vs{18}καὶ ὁ Ζῶν, καὶ ἐγενόμην νεκρὸς καὶ ἰδοὺ ζῶν εἰμι εἰς τοὺς αἰῶνας τῶν αἰώνων καὶ ἔχω τὰς κλεῖς τοῦ θανάτου καὶ τοῦ ᾅδου.
\vs{19}Γράψον οὖν ἃ εἶδες καὶ ἃ εἰσὶν καὶ ἃ μέλλει γενέσθαι μετὰ ταῦτα.
\vs{20}τὸ μυστήριον τῶν ἑπτὰ ἀστέρων οὓς εἶδες ἐπὶ τῆς δεξιᾶς μου καὶ τὰς ἑπτὰ λυχνίας τὰς χρυσᾶς· οἱ ἑπτὰ ἀστέρες ἄγγελοι τῶν ἑπτὰ ἐκκλησιῶν εἰσίν καὶ αἱ λυχνίαι αἱ ἑπτὰ ἑπτὰ ἐκκλησίαι εἰσίν.

\ch{2}
Τῷ ἀγγέλῳ τῆς ἐν Ἐφέσῳ ἐκκλησίας γράψον·

Τάδε λέγει ὁ κρατῶν τοὺς ἑπτὰ ἀστέρας ἐν τῇ δεξιᾷ αὐτοῦ, ὁ περιπατῶν ἐν μέσῳ τῶν ἑπτὰ λυχνιῶν τῶν χρυσῶν·
\vs{2}Οἶδα τὰ ἔργα σου καὶ τὸν κόπον καὶ τὴν ὑπομονήν σου καὶ ὅτι οὐ δύνῃ βαστάσαι κακούς, καὶ ἐπείρασας τοὺς λέγοντας ἑαυτοὺς ἀποστόλους καὶ οὐκ εἰσίν καὶ εὗρες αὐτοὺς ψευδεῖς,
\vs{3}καὶ ὑπομονὴν ἔχεις καὶ ἐβάστασας διὰ τὸ ὄνομά μου καὶ οὐ κεκοπίακες.
\vs{4}Ἀλλὰ ἔχω κατὰ σοῦ ὅτι τὴν ἀγάπην σου τὴν πρώτην ἀφῆκες.
\vs{5}μνημόνευε οὖν πόθεν πέπτωκας καὶ μετανόησον καὶ τὰ πρῶτα ἔργα ποίησον· εἰ δὲ μή, ἔρχομαί σοι καὶ κινήσω τὴν λυχνίαν σου ἐκ τοῦ τόπου αὐτῆς, ἐὰν μὴ μετανοήσῃς.
\vs{6}Ἀλλὰ τοῦτο ἔχεις, ὅτι μισεῖς τὰ ἔργα τῶν Νικολαϊτῶν ἃ κἀγὼ μισῶ.

\vs{7}Ὁ ἔχων οὖς ἀκουσάτω τί τὸ Πνεῦμα λέγει ταῖς ἐκκλησίαις. Τῷ νικῶντι δώσω αὐτῷ φαγεῖν ἐκ τοῦ ξύλου τῆς ζωῆς, ὅ ἐστιν ἐν τῷ Παραδείσῳ τοῦ Θεοῦ.
\vs{8}Καὶ τῷ ἀγγέλῳ τῆς ἐν Σμύρνῃ ἐκκλησίας γράψον·

Τάδε λέγει ὁ πρῶτος καὶ ὁ ἔσχατος, ὃς ἐγένετο νεκρὸς καὶ ἔζησεν·
\vs{9}Οἶδά σου τὴν θλῖψιν καὶ τὴν πτωχείαν, ἀλλὰ πλούσιος εἶ, καὶ τὴν βλασφημίαν ἐκ τῶν λεγόντων Ἰουδαίους εἶναι ἑαυτούς καὶ οὐκ εἰσίν ἀλλὰ συναγωγὴ τοῦ Σατανᾶ.
\vs{10}Μηδὲν φοβοῦ ἃ μέλλεις πάσχειν. ἰδοὺ μέλλει βάλλειν ὁ διάβολος ἐξ ὑμῶν εἰς φυλακὴν ἵνα πειρασθῆτε καὶ ἕξετε θλῖψιν ἡμερῶν δέκα. γίνου πιστὸς ἄχρι θανάτου, καὶ δώσω σοι τὸν στέφανον τῆς ζωῆς.

\vs{11}Ὁ ἔχων οὖς ἀκουσάτω τί τὸ Πνεῦμα λέγει ταῖς ἐκκλησίαις. Ὁ νικῶν οὐ μὴ ἀδικηθῇ ἐκ τοῦ θανάτου τοῦ δευτέρου.

\vs{12}Καὶ τῷ ἀγγέλῳ τῆς ἐν Περγάμῳ ἐκκλησίας γράψον·

Τάδε λέγει ὁ ἔχων τὴν ῥομφαίαν τὴν δίστομον τὴν ὀξεῖαν·
\vs{13}Οἶδα ποῦ κατοικεῖς, ὅπου ὁ θρόνος τοῦ Σατανᾶ, καὶ κρατεῖς τὸ ὄνομά μου καὶ οὐκ ἠρνήσω τὴν πίστιν μου καὶ ἐν ταῖς ἡμέραις Ἀντιπᾶς ὁ μάρτυς μου ὁ πιστός μου, ὃς ἀπεκτάνθη παρ᾽ ὑμῖν, ὅπου ὁ Σατανᾶς κατοικεῖ.
\vs{14}Ἀλλ᾽ ἔχω κατὰ σοῦ ὀλίγα ὅτι ἔχεις ἐκεῖ κρατοῦντας τὴν διδαχὴν Βαλαάμ, ὃς ἐδίδασκεν τῷ Βαλὰκ βαλεῖν σκάνδαλον ἐνώπιον τῶν υἱῶν Ἰσραήλ φαγεῖν εἰδωλόθυτα καὶ πορνεῦσαι.
\vs{15}οὕτως ἔχεις καὶ σὺ κρατοῦντας τὴν διδαχὴν τῶν Νικολαϊτῶν ὁμοίως.
\vs{16}μετανόησον οὖν· εἰ δὲ μή, ἔρχομαί σοι ταχύ καὶ πολεμήσω μετ᾽ αὐτῶν ἐν τῇ ῥομφαίᾳ τοῦ στόματός μου.

\vs{17}Ὁ ἔχων οὖς ἀκουσάτω τί τὸ Πνεῦμα λέγει ταῖς ἐκκλησίαις. Τῷ νικῶντι δώσω αὐτῷ τοῦ μάννα τοῦ κεκρυμμένου καὶ δώσω αὐτῷ ψῆφον λευκήν, καὶ ἐπὶ τὴν ψῆφον ὄνομα καινὸν γεγραμμένον ὃ οὐδεὶς οἶδεν εἰ μὴ ὁ λαμβάνων.

\vs{18}Καὶ τῷ ἀγγέλῳ τῆς ἐν Θυατείροις ἐκκλησίας γράψον·

Τάδε λέγει ὁ Υἱὸς τοῦ Θεοῦ, ὁ ἔχων τοὺς ὀφθαλμοὺς αὐτοῦ ὡς φλόγα πυρός καὶ οἱ πόδες αὐτοῦ ὅμοιοι χαλκολιβάνῳ·
\vs{19}Οἶδά σου τὰ ἔργα καὶ τὴν ἀγάπην καὶ τὴν πίστιν καὶ τὴν διακονίαν καὶ τὴν ὑπομονήν σου, καὶ τὰ ἔργα σου τὰ ἔσχατα πλείονα τῶν πρώτων.
\vs{20}Ἀλλὰ ἔχω κατὰ σοῦ ὅτι ἀφεῖς τὴν γυναῖκα Ἰεζάβελ, ἡ λέγουσα ἑαυτὴν προφῆτιν καὶ διδάσκει καὶ πλανᾷ τοὺς ἐμοὺς δούλους πορνεῦσαι καὶ φαγεῖν εἰδωλόθυτα.
\vs{21}καὶ ἔδωκα αὐτῇ χρόνον ἵνα μετανοήσῃ, καὶ οὐ θέλει μετανοῆσαι ἐκ τῆς πορνείας αὐτῆς.
\vs{22}Ἰδοὺ βάλλω αὐτὴν εἰς κλίνην καὶ τοὺς μοιχεύοντας μετ᾽ αὐτῆς εἰς θλῖψιν μεγάλην, ἐὰν μὴ μετανοήσωσιν ἐκ τῶν ἔργων αὐτῆς,
\vs{23}καὶ τὰ τέκνα αὐτῆς ἀποκτενῶ ἐν θανάτῳ. καὶ γνώσονται πᾶσαι αἱ ἐκκλησίαι ὅτι ἐγώ εἰμι ὁ ἐραυνῶν νεφροὺς καὶ καρδίας, καὶ δώσω ὑμῖν ἑκάστῳ κατὰ τὰ ἔργα ὑμῶν.
\vs{24}Ὑμῖν δὲ λέγω τοῖς λοιποῖς τοῖς ἐν Θυατείροις, ὅσοι οὐκ ἔχουσιν τὴν διδαχὴν ταύτην, οἵτινες οὐκ ἔγνωσαν τὰ βαθέα τοῦ Σατανᾶ ὡς λέγουσιν· οὐ βάλλω ἐφ᾽ ὑμᾶς ἄλλο βάρος,
\vs{25}πλὴν ὃ ἔχετε κρατήσατε ἄχρι οὗ ἂν ἥξω.

\vs{26}Καὶ ὁ νικῶν καὶ ὁ τηρῶν ἄχρι τέλους τὰ ἔργα μου, δώσω αὐτῷ ἐξουσίαν ἐπὶ τῶν ἐθνῶν
\vs{27}καὶ ποιμανεῖ αὐτοὺς ἐν ῥάβδῳ σιδηρᾷ ὡς τὰ σκεύη τὰ κεραμικὰ συντρίβεται,
\vs{28}ὡς κἀγὼ εἴληφα παρὰ τοῦ Πατρός μου, καὶ δώσω αὐτῷ τὸν ἀστέρα τὸν πρωϊνόν.
\vs{29}Ὁ ἔχων οὖς ἀκουσάτω τί τὸ Πνεῦμα λέγει ταῖς ἐκκλησίαις.

\ch{3}
Καὶ τῷ ἀγγέλῳ τῆς ἐν Σάρδεσιν ἐκκλησίας γράψον·

Τάδε λέγει ὁ ἔχων τὰ ἑπτὰ Πνεύματα τοῦ Θεοῦ καὶ τοὺς ἑπτὰ ἀστέρας· Οἶδά σου τὰ ἔργα ὅτι ὄνομα ἔχεις ὅτι ζῇς, καὶ νεκρὸς εἶ.
\vs{2}γίνου γρηγορῶν καὶ στήρισον τὰ λοιπὰ ἃ ἔμελλον ἀποθανεῖν, οὐ γὰρ εὕρηκά σου τὰ ἔργα πεπληρωμένα ἐνώπιον τοῦ Θεοῦ μου.
\vs{3}μνημόνευε οὖν πῶς εἴληφας καὶ ἤκουσας καὶ τήρει καὶ μετανόησον. ἐὰν οὖν μὴ γρηγορήσῃς, ἥξω ὡς κλέπτης, καὶ οὐ μὴ γνῷς ποίαν ὥραν ἥξω ἐπὶ σέ.
\vs{4}Ἀλλὰ ἔχεις ὀλίγα ὀνόματα ἐν Σάρδεσιν ἃ οὐκ ἐμόλυναν τὰ ἱμάτια αὐτῶν, καὶ περιπατήσουσιν μετ᾽ ἐμοῦ ἐν λευκοῖς, ὅτι ἄξιοί εἰσιν.

\vs{5}Ὁ νικῶν οὕτως περιβαλεῖται ἐν ἱματίοις λευκοῖς καὶ οὐ μὴ ἐξαλείψω τὸ ὄνομα αὐτοῦ ἐκ τῆς βίβλου τῆς ζωῆς καὶ ὁμολογήσω τὸ ὄνομα αὐτοῦ ἐνώπιον τοῦ Πατρός μου καὶ ἐνώπιον τῶν ἀγγέλων αὐτοῦ.
\vs{6}Ὁ ἔχων οὖς ἀκουσάτω τί τὸ Πνεῦμα λέγει ταῖς ἐκκλησίαις.

\vs{7}Καὶ τῷ ἀγγέλῳ τῆς ἐν Φιλαδελφείᾳ ἐκκλησίας γράψον·

Τάδε λέγει ὁ ἅγιος, ὁ ἀληθινός, ὁ ἔχων τὴν κλεῖν Δαυίδ, ὁ ἀνοίγων καὶ οὐδεὶς κλείσει καὶ κλείων καὶ οὐδεὶς ἀνοίγει·
\vs{8}Οἶδά σου τὰ ἔργα, ἰδοὺ δέδωκα ἐνώπιόν σου θύραν ἠνεῳγμένην, ἣν οὐδεὶς δύναται κλεῖσαι αὐτήν, ὅτι μικρὰν ἔχεις δύναμιν καὶ ἐτήρησάς μου τὸν λόγον καὶ οὐκ ἠρνήσω τὸ ὄνομά μου.
\vs{9}ἰδοὺ διδῶ ἐκ τῆς συναγωγῆς τοῦ Σατανᾶ τῶν λεγόντων ἑαυτοὺς Ἰουδαίους εἶναι, καὶ οὐκ εἰσὶν ἀλλὰ ψεύδονται. ἰδοὺ ποιήσω αὐτοὺς ἵνα ἥξουσιν καὶ προσκυνήσουσιν ἐνώπιον τῶν ποδῶν σου καὶ γνῶσιν ὅτι ἐγὼ ἠγάπησά σε.
\vs{10}Ὅτι ἐτήρησας τὸν λόγον τῆς ὑπομονῆς μου, κἀγώ σε τηρήσω ἐκ τῆς ὥρας τοῦ πειρασμοῦ τῆς μελλούσης ἔρχεσθαι ἐπὶ τῆς οἰκουμένης ὅλης πειράσαι τοὺς κατοικοῦντας ἐπὶ τῆς γῆς.
\vs{11}ἔρχομαι ταχύ· κράτει ὃ ἔχεις, ἵνα μηδεὶς λάβῃ τὸν στέφανόν σου.

\vs{12}Ὁ νικῶν ποιήσω αὐτὸν στῦλον ἐν τῷ ναῷ τοῦ Θεοῦ μου καὶ ἔξω οὐ μὴ ἐξέλθῃ ἔτι καὶ γράψω ἐπ᾽ αὐτὸν τὸ ὄνομα τοῦ Θεοῦ μου καὶ τὸ ὄνομα τῆς πόλεως τοῦ Θεοῦ μου, τῆς καινῆς Ἰερουσαλήμ ἡ καταβαίνουσα ἐκ τοῦ οὐρανοῦ ἀπὸ τοῦ Θεοῦ μου, καὶ τὸ ὄνομά μου τὸ καινόν.
\vs{13}Ὁ ἔχων οὖς ἀκουσάτω τί τὸ Πνεῦμα λέγει ταῖς ἐκκλησίαις.

\vs{14}Καὶ τῷ ἀγγέλῳ τῆς ἐν Λαοδικείᾳ ἐκκλησίας γράψον·

Τάδε λέγει ὁ Ἀμήν, ὁ μάρτυς ὁ πιστὸς καὶ ἀληθινός, ἡ ἀρχὴ τῆς κτίσεως τοῦ Θεοῦ·
\vs{15}Οἶδά σου τὰ ἔργα ὅτι οὔτε ψυχρὸς εἶ οὔτε ζεστός. ὄφελον ψυχρὸς ἦς ἢ ζεστός.
\vs{16}οὕτως ὅτι χλιαρὸς εἶ καὶ οὔτε ζεστὸς οὔτε ψυχρός, μέλλω σε ἐμέσαι ἐκ τοῦ στόματός μου.
\vs{17}Ὅτι λέγεις ὅτι Πλούσιός εἰμι καὶ πεπλούτηκα καὶ οὐδὲν χρείαν ἔχω, καὶ οὐκ οἶδας ὅτι σὺ εἶ ὁ ταλαίπωρος καὶ ἐλεεινὸς καὶ πτωχὸς καὶ τυφλὸς καὶ γυμνός,
\vs{18}συμβουλεύω σοι ἀγοράσαι παρ᾽ ἐμοῦ χρυσίον πεπυρωμένον ἐκ πυρὸς ἵνα πλουτήσῃς, καὶ ἱμάτια λευκὰ ἵνα περιβάλῃ καὶ μὴ φανερωθῇ ἡ αἰσχύνη τῆς γυμνότητός σου, καὶ κολλούριον ἐγχρῖσαι τοὺς ὀφθαλμούς σου ἵνα βλέπῃς.
\vs{19}ἐγὼ ὅσους ἐὰν φιλῶ ἐλέγχω καὶ παιδεύω· ζήλευε οὖν καὶ μετανόησον.
\vs{20}Ἰδοὺ ἕστηκα ἐπὶ τὴν θύραν καὶ κρούω· ἐάν τις ἀκούσῃ τῆς φωνῆς μου καὶ ἀνοίξῃ τὴν θύραν, καὶ εἰσελεύσομαι πρὸς αὐτὸν καὶ δειπνήσω μετ᾽ αὐτοῦ καὶ αὐτὸς μετ᾽ ἐμοῦ.

\vs{21}Ὁ νικῶν δώσω αὐτῷ καθίσαι μετ᾽ ἐμοῦ ἐν τῷ θρόνῳ μου, ὡς κἀγὼ ἐνίκησα καὶ ἐκάθισα μετὰ τοῦ Πατρός μου ἐν τῷ θρόνῳ αὐτοῦ.
\vs{22}Ὁ ἔχων οὖς ἀκουσάτω τί τὸ Πνεῦμα λέγει ταῖς ἐκκλησίαις.

\ch{4}
Μετὰ ταῦτα εἶδον, καὶ ἰδοὺ θύρα ἠνεῳγμένη ἐν τῷ οὐρανῷ, καὶ ἡ φωνὴ ἡ πρώτη ἣν ἤκουσα ὡς σάλπιγγος λαλούσης μετ᾽ ἐμοῦ λέγων· Ἀνάβα ὧδε, καὶ δείξω σοι ἃ δεῖ γενέσθαι μετὰ ταῦτα.

\vs{2}εὐθέως ἐγενόμην ἐν Πνεύματι, καὶ ἰδοὺ θρόνος ἔκειτο ἐν τῷ οὐρανῷ, καὶ ἐπὶ τὸν θρόνον καθήμενος,
\vs{3}καὶ ὁ καθήμενος ὅμοιος ὁράσει λίθῳ ἰάσπιδι καὶ σαρδίῳ, καὶ ἶρις κυκλόθεν τοῦ θρόνου ὅμοιος ὁράσει σμαραγδίνῳ.
\vs{4}καὶ κυκλόθεν τοῦ θρόνου θρόνους εἴκοσι τέσσαρες, καὶ ἐπὶ τοὺς θρόνους εἴκοσι τέσσαρας πρεσβυτέρους καθημένους περιβεβλημένους ἐν ἱματίοις λευκοῖς καὶ ἐπὶ τὰς κεφαλὰς αὐτῶν στεφάνους χρυσοῦς.
\vs{5}Καὶ ἐκ τοῦ θρόνου ἐκπορεύονται ἀστραπαὶ καὶ φωναὶ καὶ βρονταί, καὶ ἑπτὰ λαμπάδες πυρὸς καιόμεναι ἐνώπιον τοῦ θρόνου, ἅ εἰσιν τὰ ἑπτὰ Πνεύματα τοῦ Θεοῦ,
\vs{6}καὶ ἐνώπιον τοῦ θρόνου ὡς θάλασσα ὑαλίνη ὁμοία κρυστάλλῳ. καὶ ἐν μέσῳ τοῦ θρόνου καὶ κύκλῳ τοῦ θρόνου τέσσαρα ζῷα γέμοντα ὀφθαλμῶν ἔμπροσθεν καὶ ὄπισθεν.
\vs{7}καὶ τὸ ζῷον τὸ πρῶτον ὅμοιον λέοντι καὶ τὸ δεύτερον ζῷον ὅμοιον μόσχῳ καὶ τὸ τρίτον ζῷον ἔχων τὸ πρόσωπον ὡς ἀνθρώπου καὶ τὸ τέταρτον ζῷον ὅμοιον ἀετῷ πετομένῳ.
\vs{8}καὶ τὰ τέσσαρα ζῷα, ἓν καθ᾽ ἓν αὐτῶν ἔχων ἀνὰ πτέρυγας ἕξ, κυκλόθεν καὶ ἔσωθεν γέμουσιν ὀφθαλμῶν, καὶ ἀνάπαυσιν οὐκ ἔχουσιν ἡμέρας καὶ νυκτὸς λέγοντες· 
\begin{poetryblock}

\begin{quote}Ἅγιος ἅγιος ἅγιος Κύριος ὁ Θεός ὁ Παντοκράτωρ,\end{quote} 

\begin{quote}ὁ ἦν καὶ ὁ ὢν καὶ ὁ ἐρχόμενος.\end{quote}
\end{poetryblock}

\vs{9}Καὶ ὅταν δώσουσιν τὰ ζῷα δόξαν καὶ τιμὴν καὶ εὐχαριστίαν τῷ καθημένῳ ἐπὶ τῷ θρόνῳ τῷ ζῶντι εἰς τοὺς αἰῶνας τῶν αἰώνων,
\vs{10}πεσοῦνται οἱ εἴκοσι τέσσαρες πρεσβύτεροι ἐνώπιον τοῦ καθημένου ἐπὶ τοῦ θρόνου καὶ προσκυνήσουσιν τῷ ζῶντι εἰς τοὺς αἰῶνας τῶν αἰώνων καὶ βαλοῦσιν τοὺς στεφάνους αὐτῶν ἐνώπιον τοῦ θρόνου λέγοντες·
\begin{poetryblock}

\begin{quote} \vs{11}Ἄξιος εἶ, ὁ Κύριος καὶ ὁ Θεὸς ἡμῶν, λαβεῖν τὴν δόξαν καὶ τὴν τιμὴν καὶ τὴν δύναμιν, ὅτι σὺ ἔκτισας τὰ πάντα καὶ διὰ τὸ θέλημά σου ἦσαν καὶ ἐκτίσθησαν.\end{quote}
\end{poetryblock}

\ch{5}
Καὶ εἶδον ἐπὶ τὴν δεξιὰν τοῦ καθημένου ἐπὶ τοῦ θρόνου βιβλίον γεγραμμένον ἔσωθεν καὶ ὄπισθεν κατεσφραγισμένον σφραγῖσιν ἑπτά.
\vs{2}καὶ εἶδον ἄγγελον ἰσχυρὸν κηρύσσοντα ἐν φωνῇ μεγάλῃ· Τίς ἄξιος ἀνοῖξαι τὸ βιβλίον καὶ λῦσαι τὰς σφραγῖδας αὐτοῦ;
\vs{3}Καὶ οὐδεὶς ἐδύνατο ἐν τῷ οὐρανῷ οὐδὲ ἐπὶ τῆς γῆς οὐδὲ ὑποκάτω τῆς γῆς ἀνοῖξαι τὸ βιβλίον οὔτε βλέπειν αὐτό.
\vs{4}καὶ ἔκλαιον πολὺ, ὅτι οὐδεὶς ἄξιος εὑρέθη ἀνοῖξαι τὸ βιβλίον οὔτε βλέπειν αὐτό.
\vs{5}Καὶ εἷς ἐκ τῶν πρεσβυτέρων λέγει μοι· Μὴ κλαῖε, ἰδοὺ ἐνίκησεν ὁ Λέων ὁ ἐκ τῆς φυλῆς Ἰούδα, ἡ Ῥίζα Δαυίδ, ἀνοῖξαι τὸ βιβλίον καὶ τὰς ἑπτὰ σφραγῖδας αὐτοῦ.

\vs{6}Καὶ εἶδον ἐν μέσῳ τοῦ θρόνου καὶ τῶν τεσσάρων ζῴων καὶ ἐν μέσῳ τῶν πρεσβυτέρων Ἀρνίον ἑστηκὸς ὡς ἐσφαγμένον ἔχων κέρατα ἑπτὰ καὶ ὀφθαλμοὺς ἑπτά οἵ εἰσιν τὰ ἑπτὰ Πνεύματα τοῦ Θεοῦ ἀπεσταλμένοι εἰς πᾶσαν τὴν γῆν.
\vs{7}καὶ ἦλθεν καὶ εἴληφεν ἐκ τῆς δεξιᾶς τοῦ καθημένου ἐπὶ τοῦ θρόνου.

\vs{8}Καὶ ὅτε ἔλαβεν τὸ βιβλίον, τὰ τέσσαρα ζῷα καὶ οἱ εἴκοσι τέσσαρες πρεσβύτεροι ἔπεσαν ἐνώπιον τοῦ Ἀρνίου ἔχοντες ἕκαστος κιθάραν καὶ φιάλας χρυσᾶς γεμούσας θυμιαμάτων, αἵ εἰσιν αἱ προσευχαὶ τῶν ἁγίων,
\vs{9}καὶ ᾄδουσιν ᾠδὴν καινὴν λέγοντες· 
\begin{poetryblock}

\begin{quote}Ἄξιος εἶ λαβεῖν τὸ βιβλίον καὶ ἀνοῖξαι τὰς σφραγῖδας αὐτοῦ,\end{quote} 

\begin{quote}ὅτι ἐσφάγης καὶ ἠγόρασας τῷ Θεῷ ἐν τῷ αἵματί σου\end{quote} 

\begin{quote}ἐκ πάσης φυλῆς καὶ γλώσσης καὶ λαοῦ καὶ ἔθνους\end{quote}

\begin{quote} \vs{10}καὶ ἐποίησας αὐτοὺς τῷ Θεῷ ἡμῶν βασιλείαν καὶ ἱερεῖς,\end{quote} 

\begin{quote}καὶ βασιλεύσουσιν ἐπὶ τῆς γῆς.\end{quote}
\end{poetryblock}

\vs{11}Καὶ εἶδον, καὶ ἤκουσα φωνὴν ἀγγέλων πολλῶν κύκλῳ τοῦ θρόνου καὶ τῶν ζῴων καὶ τῶν πρεσβυτέρων, καὶ ἦν ὁ ἀριθμὸς αὐτῶν μυριάδες μυριάδων καὶ χιλιάδες χιλιάδων
\vs{12}λέγοντες φωνῇ μεγάλῃ· 
\begin{poetryblock}

\begin{quote}Ἄξιόν ἐστιν τὸ Ἀρνίον τὸ ἐσφαγμένον λαβεῖν\end{quote} 

\begin{quote}τὴν δύναμιν καὶ πλοῦτον καὶ σοφίαν\end{quote} 

\begin{quote}καὶ ἰσχὺν καὶ τιμὴν καὶ δόξαν καὶ εὐλογίαν.\end{quote}
\end{poetryblock}

\vs{13}Καὶ πᾶν κτίσμα ὃ ἐν τῷ οὐρανῷ καὶ ἐπὶ τῆς γῆς καὶ ὑποκάτω τῆς γῆς καὶ ἐπὶ τῆς θαλάσσης καὶ τὰ ἐν αὐτοῖς πάντα ἤκουσα λέγοντας· 
\begin{poetryblock}

\begin{quote}Τῷ καθημένῳ ἐπὶ τῷ θρόνῳ καὶ τῷ Ἀρνίῳ\end{quote} 

\begin{quote}ἡ εὐλογία καὶ ἡ τιμὴ καὶ ἡ δόξα καὶ τὸ κράτος\end{quote} 

\begin{quote}εἰς τοὺς αἰῶνας τῶν αἰώνων.\end{quote}
\end{poetryblock}
\vs{14}Καὶ τὰ τέσσαρα ζῷα ἔλεγον· Ἀμήν. καὶ οἱ πρεσβύτεροι ἔπεσαν καὶ προσεκύνησαν.

\ch{6}
Καὶ εἶδον ὅτε ἤνοιξεν τὸ Ἀρνίον μίαν ἐκ τῶν ἑπτὰ σφραγίδων, καὶ ἤκουσα ἑνὸς ἐκ τῶν τεσσάρων ζῴων λέγοντος ὡς φωνῇ βροντῆς· Ἔρχου.
\vs{2}Καὶ εἶδον, καὶ ἰδοὺ ἵππος λευκός, καὶ ὁ καθήμενος ἐπ᾽ αὐτὸν ἔχων τόξον καὶ ἐδόθη αὐτῷ στέφανος καὶ ἐξῆλθεν νικῶν καὶ ἵνα νικήσῃ.

\vs{3}Καὶ ὅτε ἤνοιξεν τὴν σφραγῖδα τὴν δευτέραν, ἤκουσα τοῦ δευτέρου ζῴου λέγοντος· Ἔρχου.
\vs{4}Καὶ ἐξῆλθεν ἄλλος ἵππος πυρρός, καὶ τῷ καθημένῳ ἐπ᾽ αὐτὸν ἐδόθη αὐτῷ λαβεῖν τὴν εἰρήνην ἐκ τῆς γῆς καὶ ἵνα ἀλλήλους σφάξουσιν καὶ ἐδόθη αὐτῷ μάχαιρα μεγάλη.

\vs{5}Καὶ ὅτε ἤνοιξεν τὴν σφραγῖδα τὴν τρίτην, ἤκουσα τοῦ τρίτου ζῴου λέγοντος· Ἔρχου. Καὶ εἶδον, καὶ ἰδοὺ ἵππος μέλας, καὶ ὁ καθήμενος ἐπ᾽ αὐτὸν ἔχων ζυγὸν ἐν τῇ χειρὶ αὐτοῦ.
\vs{6}καὶ ἤκουσα ὡς φωνὴν ἐν μέσῳ τῶν τεσσάρων ζῴων λέγουσαν· Χοῖνιξ σίτου δηναρίου καὶ τρεῖς χοίνικες κριθῶν δηναρίου, καὶ τὸ ἔλαιον καὶ τὸν οἶνον μὴ ἀδικήσῃς.

\vs{7}Καὶ ὅτε ἤνοιξεν τὴν σφραγῖδα τὴν τετάρτην, ἤκουσα φωνὴν τοῦ τετάρτου ζῴου λέγοντος· Ἔρχου.
\vs{8}Καὶ εἶδον, καὶ ἰδοὺ ἵππος χλωρός, καὶ ὁ καθήμενος ἐπάνω αὐτοῦ ὄνομα αὐτῷ Ὁ Θάνατος, καὶ ὁ ᾅδης ἠκολούθει μετ᾽ αὐτοῦ καὶ ἐδόθη αὐτοῖς ἐξουσία ἐπὶ τὸ τέταρτον τῆς γῆς ἀποκτεῖναι ἐν ῥομφαίᾳ καὶ ἐν λιμῷ καὶ ἐν θανάτῳ καὶ ὑπὸ τῶν θηρίων τῆς γῆς.

\vs{9}Καὶ ὅτε ἤνοιξεν τὴν πέμπτην σφραγῖδα, εἶδον ὑποκάτω τοῦ θυσιαστηρίου τὰς ψυχὰς τῶν ἐσφαγμένων διὰ τὸν λόγον τοῦ Θεοῦ καὶ διὰ τὴν μαρτυρίαν ἣν εἶχον.
\vs{10}καὶ ἔκραξαν φωνῇ μεγάλῃ λέγοντες· Ἕως πότε, ὁ Δεσπότης ὁ ἅγιος καὶ ἀληθινός, οὐ κρίνεις καὶ ἐκδικεῖς τὸ αἷμα ἡμῶν ἐκ τῶν κατοικούντων ἐπὶ τῆς γῆς;
\vs{11}Καὶ ἐδόθη αὐτοῖς ἑκάστῳ στολὴ λευκή καὶ ἐρρέθη αὐτοῖς ἵνα ἀναπαύσονται ἔτι χρόνον μικρόν, ἕως πληρωθῶσιν καὶ οἱ σύνδουλοι αὐτῶν καὶ οἱ ἀδελφοὶ αὐτῶν οἱ μέλλοντες ἀποκτέννεσθαι ὡς καὶ αὐτοί.

\vs{12}Καὶ εἶδον ὅτε ἤνοιξεν τὴν σφραγῖδα τὴν ἕκτην, καὶ σεισμὸς μέγας ἐγένετο καὶ ὁ ἥλιος ἐγένετο μέλας ὡς σάκκος τρίχινος καὶ ἡ σελήνη ὅλη ἐγένετο ὡς αἷμα
\vs{13}καὶ οἱ ἀστέρες τοῦ οὐρανοῦ ἔπεσαν εἰς τὴν γῆν, ὡς συκῆ βάλλει τοὺς ὀλύνθους αὐτῆς ὑπὸ ἀνέμου μεγάλου σειομένη,
\vs{14}καὶ ὁ οὐρανὸς ἀπεχωρίσθη ὡς βιβλίον ἑλισσόμενον καὶ πᾶν ὄρος καὶ νῆσος ἐκ τῶν τόπων αὐτῶν ἐκινήθησαν.
\vs{15}Καὶ οἱ βασιλεῖς τῆς γῆς καὶ οἱ μεγιστᾶνες καὶ οἱ χιλίαρχοι καὶ οἱ πλούσιοι καὶ οἱ ἰσχυροὶ καὶ πᾶς δοῦλος καὶ ἐλεύθερος ἔκρυψαν ἑαυτοὺς εἰς τὰ σπήλαια καὶ εἰς τὰς πέτρας τῶν ὀρέων
\vs{16}καὶ λέγουσιν τοῖς ὄρεσιν καὶ ταῖς πέτραις· Πέσετε ἐφ᾽ ἡμᾶς καὶ κρύψατε ἡμᾶς ἀπὸ προσώπου τοῦ καθημένου ἐπὶ τοῦ θρόνου καὶ ἀπὸ τῆς ὀργῆς τοῦ Ἀρνίου,
\vs{17}ὅτι ἦλθεν ἡ ἡμέρα ἡ μεγάλη τῆς ὀργῆς αὐτῶν, καὶ τίς δύναται σταθῆναι;

\ch{7}
Μετὰ τοῦτο εἶδον τέσσαρας ἀγγέλους ἑστῶτας ἐπὶ τὰς τέσσαρας γωνίας τῆς γῆς, κρατοῦντας τοὺς τέσσαρας ἀνέμους τῆς γῆς ἵνα μὴ πνέῃ ἄνεμος ἐπὶ τῆς γῆς μήτε ἐπὶ τῆς θαλάσσης μήτε ἐπὶ πᾶν δένδρον.
\vs{2}καὶ εἶδον ἄλλον ἄγγελον ἀναβαίνοντα ἀπὸ ἀνατολῆς ἡλίου ἔχοντα σφραγῖδα Θεοῦ ζῶντος, καὶ ἔκραξεν φωνῇ μεγάλῃ τοῖς τέσσαρσιν ἀγγέλοις οἷς ἐδόθη αὐτοῖς ἀδικῆσαι τὴν γῆν καὶ τὴν θάλασσαν
\vs{3}λέγων· Μὴ ἀδικήσητε τὴν γῆν μήτε τὴν θάλασσαν μήτε τὰ δένδρα, ἄχρι σφραγίσωμεν τοὺς δούλους τοῦ Θεοῦ ἡμῶν ἐπὶ τῶν μετώπων αὐτῶν.

\vs{4}Καὶ ἤκουσα τὸν ἀριθμὸν τῶν ἐσφραγισμένων, ἑκατὸν τεσσεράκοντα τέσσαρες χιλιάδες, ἐσφραγισμένοι ἐκ πάσης φυλῆς υἱῶν Ἰσραήλ·

\vs{5}Ἐκ φυλῆς Ἰούδα δώδεκα χιλιάδες ἐσφραγισμένοι, 
\begin{poetryblock}

\begin{quote}ἐκ φυλῆς Ῥουβὴν δώδεκα χιλιάδες,\end{quote} 

\begin{quote}ἐκ φυλῆς Γὰδ δώδεκα χιλιάδες,\end{quote}

\begin{quote} \vs{6}Ἐκ φυλῆς Ἀσὴρ δώδεκα χιλιάδες,\end{quote} 

\begin{quote}ἐκ φυλῆς Νεφθαλὶμ δώδεκα χιλιάδες,\end{quote} 

\begin{quote}ἐκ φυλῆς Μανασσῆ δώδεκα χιλιάδες,\end{quote}

\begin{quote} \vs{7}Ἐκ φυλῆς Συμεὼν δώδεκα χιλιάδες,\end{quote} 

\begin{quote}ἐκ φυλῆς Λευὶ δώδεκα χιλιάδες,\end{quote} 

\begin{quote}ἐκ φυλῆς Ἰσσαχὰρ δώδεκα χιλιάδες,\end{quote}

\begin{quote} \vs{8}Ἐκ φυλῆς Ζαβουλὼν δώδεκα χιλιάδες,\end{quote} 

\begin{quote}ἐκ φυλῆς Ἰωσὴφ δώδεκα χιλιάδες,\end{quote} 

\begin{quote}ἐκ φυλῆς Βενιαμὶν δώδεκα χιλιάδες ἐσφραγισμένοι.\end{quote}
\end{poetryblock}

\vs{9}Μετὰ ταῦτα εἶδον, καὶ ἰδοὺ ὄχλος πολύς, ὃν ἀριθμῆσαι αὐτὸν οὐδεὶς ἐδύνατο, ἐκ παντὸς ἔθνους καὶ φυλῶν καὶ λαῶν καὶ γλωσσῶν ἑστῶτες ἐνώπιον τοῦ θρόνου καὶ ἐνώπιον τοῦ Ἀρνίου περιβεβλημένους στολὰς λευκάς καὶ φοίνικες ἐν ταῖς χερσὶν αὐτῶν,
\vs{10}καὶ κράζουσιν φωνῇ μεγάλῃ λέγοντες· 
\begin{poetryblock}

\begin{quote}Ἡ σωτηρία τῷ Θεῷ ἡμῶν τῷ καθημένῳ ἐπὶ τῷ θρόνῳ καὶ τῷ Ἀρνίῳ.\end{quote}
\end{poetryblock}
\vs{11}Καὶ πάντες οἱ ἄγγελοι εἱστήκεισαν κύκλῳ τοῦ θρόνου καὶ τῶν πρεσβυτέρων καὶ τῶν τεσσάρων ζῴων καὶ ἔπεσαν ἐνώπιον τοῦ θρόνου ἐπὶ τὰ πρόσωπα αὐτῶν καὶ προσεκύνησαν τῷ Θεῷ

\vs{12}λέγοντες· 
\begin{poetryblock}

\begin{quote}Ἀμήν, ἡ εὐλογία καὶ ἡ δόξα καὶ ἡ σοφία καὶ ἡ εὐχαριστία καὶ ἡ τιμὴ καὶ ἡ δύναμις καὶ ἡ ἰσχὺς τῷ Θεῷ ἡμῶν εἰς τοὺς αἰῶνας τῶν αἰώνων· ἀμήν.\end{quote}
\end{poetryblock}

\vs{13}Καὶ ἀπεκρίθη εἷς ἐκ τῶν πρεσβυτέρων λέγων μοι· Οὗτοι οἱ περιβεβλημένοι τὰς στολὰς τὰς λευκὰς τίνες εἰσὶν καὶ πόθεν ἦλθον;
\vs{14}Καὶ εἴρηκα αὐτῷ· Κύριέ μου, σὺ οἶδας. Καὶ εἶπέν μοι· 
\begin{poetryblock}

\begin{quote}Οὗτοί εἰσιν οἱ ἐρχόμενοι ἐκ τῆς θλίψεως τῆς μεγάλης\end{quote} 

\begin{quote}καὶ ἔπλυναν τὰς στολὰς αὐτῶν\end{quote} 

\begin{quote}καὶ ἐλεύκαναν αὐτὰς ἐν τῷ αἵματι τοῦ Ἀρνίου.\end{quote}

\begin{quote} \vs{15}διὰ τοῦτό Εἰσιν ἐνώπιον τοῦ θρόνου τοῦ Θεοῦ\end{quote} 

\begin{quote}καὶ λατρεύουσιν αὐτῷ ἡμέρας καὶ νυκτὸς ἐν τῷ ναῷ αὐτοῦ,\end{quote} 

\begin{quote}καὶ ὁ καθήμενος ἐπὶ τοῦ θρόνου σκηνώσει ἐπ᾽ αὐτούς.\end{quote}

\begin{quote} \vs{16}οὐ πεινάσουσιν ἔτι οὐδὲ διψήσουσιν ἔτι\end{quote} 

\begin{quote}οὐδὲ μὴ πέσῃ ἐπ᾽ αὐτοὺς ὁ ἥλιος οὐδὲ πᾶν καῦμα,\end{quote}

\begin{quote} \vs{17}ὅτι τὸ Ἀρνίον τὸ ἀνὰ μέσον τοῦ θρόνου ποιμανεῖ αὐτούς\end{quote} 

\begin{quote}καὶ ὁδηγήσει αὐτοὺς ἐπὶ ζωῆς πηγὰς ὑδάτων,\end{quote} 

\begin{quote}καὶ ἐξαλείψει ὁ Θεὸς πᾶν δάκρυον ἐκ τῶν ὀφθαλμῶν αὐτῶν.\end{quote}
\end{poetryblock}

\ch{8}
Καὶ ὅταν ἤνοιξεν τὴν σφραγῖδα τὴν ἑβδόμην, ἐγένετο σιγὴ ἐν τῷ οὐρανῷ ὡς ἡμιώριον.
\vs{2}καὶ εἶδον τοὺς ἑπτὰ ἀγγέλους οἳ ἐνώπιον τοῦ Θεοῦ ἑστήκασιν, καὶ ἐδόθησαν αὐτοῖς ἑπτὰ σάλπιγγες.
\vs{3}Καὶ ἄλλος ἄγγελος ἦλθεν καὶ ἐστάθη ἐπὶ τοῦ θυσιαστηρίου ἔχων λιβανωτὸν χρυσοῦν, καὶ ἐδόθη αὐτῷ θυμιάματα πολλὰ, ἵνα δώσει ταῖς προσευχαῖς τῶν ἁγίων πάντων ἐπὶ τὸ θυσιαστήριον τὸ χρυσοῦν τὸ ἐνώπιον τοῦ θρόνου.
\vs{4}καὶ ἀνέβη ὁ καπνὸς τῶν θυμιαμάτων ταῖς προσευχαῖς τῶν ἁγίων ἐκ χειρὸς τοῦ ἀγγέλου ἐνώπιον τοῦ Θεοῦ.
\vs{5}Καὶ εἴληφεν ὁ ἄγγελος τὸν λιβανωτόν καὶ ἐγέμισεν αὐτὸν ἐκ τοῦ πυρὸς τοῦ θυσιαστηρίου καὶ ἔβαλεν εἰς τὴν γῆν, καὶ ἐγένοντο βρονταὶ καὶ φωναὶ καὶ ἀστραπαὶ καὶ σεισμός.

\vs{6}Καὶ οἱ ἑπτὰ ἄγγελοι οἱ ἔχοντες τὰς ἑπτὰ σάλπιγγας ἡτοίμασαν αὑτοὺς ἵνα σαλπίσωσιν.

\vs{7}Καὶ ὁ πρῶτος ἐσάλπισεν· καὶ ἐγένετο χάλαζα καὶ πῦρ μεμιγμένα ἐν αἵματι καὶ ἐβλήθη εἰς τὴν γῆν, καὶ τὸ τρίτον τῆς γῆς κατεκάη καὶ τὸ τρίτον τῶν δένδρων κατεκάη καὶ πᾶς χόρτος χλωρὸς κατεκάη.

\vs{8}Καὶ ὁ δεύτερος ἄγγελος ἐσάλπισεν· καὶ ὡς ὄρος μέγα πυρὶ καιόμενον ἐβλήθη εἰς τὴν θάλασσαν, καὶ ἐγένετο τὸ τρίτον τῆς θαλάσσης αἷμα
\vs{9}καὶ ἀπέθανεν τὸ τρίτον τῶν κτισμάτων τῶν ἐν τῇ θαλάσσῃ τὰ ἔχοντα ψυχάς καὶ τὸ τρίτον τῶν πλοίων διεφθάρησαν.

\vs{10}Καὶ ὁ τρίτος ἄγγελος ἐσάλπισεν· καὶ ἔπεσεν ἐκ τοῦ οὐρανοῦ ἀστὴρ μέγας καιόμενος ὡς λαμπάς καὶ ἔπεσεν ἐπὶ τὸ τρίτον τῶν ποταμῶν καὶ ἐπὶ τὰς πηγὰς τῶν ὑδάτων,
\vs{11}καὶ τὸ ὄνομα τοῦ ἀστέρος λέγεται Ὁ Ἄψινθος, καὶ ἐγένετο τὸ τρίτον τῶν ὑδάτων εἰς ἄψινθον καὶ πολλοὶ τῶν ἀνθρώπων ἀπέθανον ἐκ τῶν ὑδάτων ὅτι ἐπικράνθησαν.

\vs{12}Καὶ ὁ τέταρτος ἄγγελος ἐσάλπισεν· καὶ ἐπλήγη τὸ τρίτον τοῦ ἡλίου καὶ τὸ τρίτον τῆς σελήνης καὶ τὸ τρίτον τῶν ἀστέρων, ἵνα σκοτισθῇ τὸ τρίτον αὐτῶν καὶ ἡ ἡμέρα μὴ φάνῃ τὸ τρίτον αὐτῆς καὶ ἡ νὺξ ὁμοίως.

\vs{13}Καὶ εἶδον, καὶ ἤκουσα ἑνὸς ἀετοῦ πετομένου ἐν μεσουρανήματι λέγοντος φωνῇ μεγάλῃ· Οὐαὶ οὐαὶ οὐαὶ τοὺς κατοικοῦντας ἐπὶ τῆς γῆς ἐκ τῶν λοιπῶν φωνῶν τῆς σάλπιγγος τῶν τριῶν ἀγγέλων τῶν μελλόντων σαλπίζειν.

\ch{9}
Καὶ ὁ πέμπτος ἄγγελος ἐσάλπισεν· καὶ εἶδον ἀστέρα ἐκ τοῦ οὐρανοῦ πεπτωκότα εἰς τὴν γῆν, καὶ ἐδόθη αὐτῷ ἡ κλεὶς τοῦ φρέατος τῆς ἀβύσσου
\vs{2}καὶ ἤνοιξεν τὸ φρέαρ τῆς ἀβύσσου, καὶ ἀνέβη καπνὸς ἐκ τοῦ φρέατος ὡς καπνὸς καμίνου μεγάλης, καὶ ἐσκοτώθη ὁ ἥλιος καὶ ὁ ἀὴρ ἐκ τοῦ καπνοῦ τοῦ φρέατος.
\vs{3}Καὶ ἐκ τοῦ καπνοῦ ἐξῆλθον ἀκρίδες εἰς τὴν γῆν, καὶ ἐδόθη αὐταῖς ἐξουσία ὡς ἔχουσιν ἐξουσίαν οἱ σκορπίοι τῆς γῆς.
\vs{4}καὶ ἐρρέθη αὐταῖς ἵνα μὴ ἀδικήσουσιν τὸν χόρτον τῆς γῆς οὐδὲ πᾶν χλωρὸν οὐδὲ πᾶν δένδρον, εἰ μὴ τοὺς ἀνθρώπους οἵτινες οὐκ ἔχουσι τὴν σφραγῖδα τοῦ Θεοῦ ἐπὶ τῶν μετώπων.
\vs{5}καὶ ἐδόθη αὐτοῖς ἵνα μὴ ἀποκτείνωσιν αὐτούς, ἀλλ᾽ ἵνα βασανισθήσονται μῆνας πέντε, καὶ ὁ βασανισμὸς αὐτῶν ὡς βασανισμὸς σκορπίου ὅταν παίσῃ ἄνθρωπον.
\vs{6}καὶ ἐν ταῖς ἡμέραις ἐκείναις ζητήσουσιν οἱ ἄνθρωποι τὸν θάνατον καὶ οὐ μὴ εὑρήσουσιν αὐτόν, καὶ ἐπιθυμήσουσιν ἀποθανεῖν καὶ φεύγει ὁ θάνατος ἀπ᾽ αὐτῶν.

\vs{7}Καὶ τὰ ὁμοιώματα τῶν ἀκρίδων ὅμοια ἵπποις ἡτοιμασμένοις εἰς πόλεμον, καὶ ἐπὶ τὰς κεφαλὰς αὐτῶν ὡς στέφανοι ὅμοιοι χρυσῷ, καὶ τὰ πρόσωπα αὐτῶν ὡς πρόσωπα ἀνθρώπων,
\vs{8}καὶ εἶχον τρίχας ὡς τρίχας γυναικῶν, καὶ οἱ ὀδόντες αὐτῶν ὡς λεόντων ἦσαν,
\vs{9}καὶ εἶχον θώρακας ὡς θώρακας σιδηροῦς, καὶ ἡ φωνὴ τῶν πτερύγων αὐτῶν ὡς φωνὴ ἁρμάτων ἵππων πολλῶν τρεχόντων εἰς πόλεμον,
\vs{10}καὶ ἔχουσιν οὐρὰς ὁμοίας σκορπίοις καὶ κέντρα, καὶ ἐν ταῖς οὐραῖς αὐτῶν ἡ ἐξουσία αὐτῶν ἀδικῆσαι τοὺς ἀνθρώπους μῆνας πέντε,
\vs{11}ἔχουσιν ἐπ᾽ αὐτῶν βασιλέα τὸν ἄγγελον τῆς ἀβύσσου, ὄνομα αὐτῷ Ἑβραϊστί Ἀβαδδών, καὶ ἐν τῇ Ἑλληνικῇ ὄνομα ἔχει Ἀπολλύων.
\vs{12}Ἡ Οὐαὶ ἡ μία ἀπῆλθεν· ἰδοὺ ἔρχεται ἔτι δύο Οὐαὶ μετὰ ταῦτα.

\vs{13}Καὶ ὁ ἕκτος ἄγγελος ἐσάλπισεν· καὶ ἤκουσα φωνὴν μίαν ἐκ τῶν τεσσάρων κεράτων τοῦ θυσιαστηρίου τοῦ χρυσοῦ τοῦ ἐνώπιον τοῦ Θεοῦ,
\vs{14}λέγοντα τῷ ἕκτῳ ἀγγέλῳ, ὁ ἔχων τὴν σάλπιγγα· Λῦσον τοὺς τέσσαρας ἀγγέλους τοὺς δεδεμένους ἐπὶ τῷ ποταμῷ τῷ μεγάλῳ Εὐφράτῃ.
\vs{15}καὶ ἐλύθησαν οἱ τέσσαρες ἄγγελοι οἱ ἡτοιμασμένοι εἰς τὴν ὥραν καὶ ἡμέραν καὶ μῆνα καὶ ἐνιαυτόν, ἵνα ἀποκτείνωσιν τὸ τρίτον τῶν ἀνθρώπων.
\vs{16}καὶ ὁ ἀριθμὸς τῶν στρατευμάτων τοῦ ἱππικοῦ δισμυριάδες μυριάδων, ἤκουσα τὸν ἀριθμὸν αὐτῶν.
\vs{17}Καὶ οὕτως εἶδον τοὺς ἵππους ἐν τῇ ὁράσει καὶ τοὺς καθημένους ἐπ᾽ αὐτῶν, ἔχοντας θώρακας πυρίνους καὶ ὑακινθίνους καὶ θειώδεις, καὶ αἱ κεφαλαὶ τῶν ἵππων ὡς κεφαλαὶ λεόντων, καὶ ἐκ τῶν στομάτων αὐτῶν ἐκπορεύεται πῦρ καὶ καπνὸς καὶ θεῖον.
\vs{18}ἀπὸ τῶν τριῶν πληγῶν τούτων ἀπεκτάνθησαν τὸ τρίτον τῶν ἀνθρώπων, ἐκ τοῦ πυρὸς καὶ τοῦ καπνοῦ καὶ τοῦ θείου τοῦ ἐκπορευομένου ἐκ τῶν στομάτων αὐτῶν.
\vs{19}ἡ γὰρ ἐξουσία τῶν ἵππων ἐν τῷ στόματι αὐτῶν ἐστιν καὶ ἐν ταῖς οὐραῖς αὐτῶν, αἱ γὰρ οὐραὶ αὐτῶν ὅμοιαι ὄφεσιν, ἔχουσαι κεφαλάς καὶ ἐν αὐταῖς ἀδικοῦσιν.

\vs{20}Καὶ οἱ λοιποὶ τῶν ἀνθρώπων, οἳ οὐκ ἀπεκτάνθησαν ἐν ταῖς πληγαῖς ταύταις, οὐδὲ μετενόησαν ἐκ τῶν ἔργων τῶν χειρῶν αὐτῶν, ἵνα μὴ προσκυνήσουσιν τὰ δαιμόνια καὶ τὰ εἴδωλα τὰ χρυσᾶ καὶ τὰ ἀργυρᾶ καὶ τὰ χαλκᾶ καὶ τὰ λίθινα καὶ τὰ ξύλινα, ἃ οὔτε βλέπειν δύνανται οὔτε ἀκούειν οὔτε περιπατεῖν,
\vs{21}καὶ οὐ μετενόησαν ἐκ τῶν φόνων αὐτῶν οὔτε ἐκ τῶν φαρμάκων αὐτῶν οὔτε ἐκ τῆς πορνείας αὐτῶν οὔτε ἐκ τῶν κλεμμάτων αὐτῶν.

\ch{10}
Καὶ εἶδον ἄλλον ἄγγελον ἰσχυρὸν καταβαίνοντα ἐκ τοῦ οὐρανοῦ περιβεβλημένον νεφέλην, καὶ ἡ ἶρις ἐπὶ τῆς κεφαλῆς αὐτοῦ καὶ τὸ πρόσωπον αὐτοῦ ὡς ὁ ἥλιος καὶ οἱ πόδες αὐτοῦ ὡς στῦλοι πυρός,
\vs{2}καὶ ἔχων ἐν τῇ χειρὶ αὐτοῦ βιβλαρίδιον ἠνεῳγμένον. καὶ ἔθηκεν τὸν πόδα αὐτοῦ τὸν δεξιὸν ἐπὶ τῆς θαλάσσης, τὸν δὲ εὐώνυμον ἐπὶ τῆς γῆς,
\vs{3}καὶ ἔκραξεν φωνῇ μεγάλῃ ὥσπερ λέων μυκᾶται. καὶ ὅτε ἔκραξεν, ἐλάλησαν αἱ ἑπτὰ βρονταὶ τὰς ἑαυτῶν φωνάς.
\vs{4}Καὶ ὅτε ἐλάλησαν αἱ ἑπτὰ βρονταί, ἤμελλον γράφειν, καὶ ἤκουσα φωνὴν ἐκ τοῦ οὐρανοῦ λέγουσαν· Σφράγισον ἃ ἐλάλησαν αἱ ἑπτὰ βρονταί, καὶ μὴ αὐτὰ γράψῃς.

\vs{5}Καὶ ὁ ἄγγελος, ὃν εἶδον ἑστῶτα ἐπὶ τῆς θαλάσσης καὶ ἐπὶ τῆς γῆς, ἦρεν τὴν χεῖρα αὐτοῦ τὴν δεξιὰν εἰς τὸν οὐρανόν
\vs{6}καὶ ὤμοσεν ἐν τῷ ζῶντι εἰς τοὺς αἰῶνας τῶν αἰώνων, ὃς ἔκτισεν τὸν οὐρανὸν καὶ τὰ ἐν αὐτῷ καὶ τὴν γῆν καὶ τὰ ἐν αὐτῇ καὶ τὴν θάλασσαν καὶ τὰ ἐν αὐτῇ, ὅτι Χρόνος οὐκέτι ἔσται,
\vs{7}ἀλλ᾽ ἐν ταῖς ἡμέραις τῆς φωνῆς τοῦ ἑβδόμου ἀγγέλου, ὅταν μέλλῃ σαλπίζειν, καὶ ἐτελέσθη τὸ μυστήριον τοῦ Θεοῦ, ὡς εὐηγγέλισεν τοὺς ἑαυτοῦ δούλους τοὺς προφήτας.

\vs{8}Καὶ ἡ φωνὴ ἣν ἤκουσα ἐκ τοῦ οὐρανοῦ πάλιν λαλοῦσαν μετ᾽ ἐμοῦ καὶ λέγουσαν· Ὕπαγε λάβε τὸ βιβλίον τὸ ἠνεῳγμένον ἐν τῇ χειρὶ τοῦ ἀγγέλου τοῦ ἑστῶτος ἐπὶ τῆς θαλάσσης καὶ ἐπὶ τῆς γῆς.
\vs{9}Καὶ ἀπῆλθα πρὸς τὸν ἄγγελον λέγων αὐτῷ Δοῦναί μοι τὸ βιβλαρίδιον. Καὶ λέγει μοι· Λάβε καὶ κατάφαγε αὐτό, καὶ πικρανεῖ σου τὴν κοιλίαν, ἀλλ᾽ ἐν τῷ στόματί σου ἔσται γλυκὺ ὡς μέλι.

\vs{10}Καὶ ἔλαβον τὸ βιβλαρίδιον ἐκ τῆς χειρὸς τοῦ ἀγγέλου καὶ κατέφαγον αὐτό, καὶ ἦν ἐν τῷ στόματί μου ὡς μέλι γλυκύ καὶ ὅτε ἔφαγον αὐτό, ἐπικράνθη ἡ κοιλία μου.
\vs{11}Καὶ λέγουσίν μοι· Δεῖ σε πάλιν προφητεῦσαι ἐπὶ λαοῖς καὶ ἔθνεσιν καὶ γλώσσαις καὶ βασιλεῦσιν πολλοῖς.

\ch{11}
Καὶ ἐδόθη μοι κάλαμος ὅμοιος ῥάβδῳ, λέγων· Ἔγειρε καὶ μέτρησον τὸν ναὸν τοῦ Θεοῦ καὶ τὸ θυσιαστήριον καὶ τοὺς προσκυνοῦντας ἐν αὐτῷ.
\vs{2}καὶ τὴν αὐλὴν τὴν ἔξωθεν τοῦ ναοῦ ἔκβαλε ἔξωθεν καὶ μὴ αὐτὴν μετρήσῃς, ὅτι ἐδόθη τοῖς ἔθνεσιν, καὶ τὴν πόλιν τὴν ἁγίαν πατήσουσιν μῆνας τεσσεράκοντα καὶ δύο.

\vs{3}καὶ δώσω τοῖς δυσὶν μάρτυσίν μου καὶ προφητεύσουσιν ἡμέρας χιλίας διακοσίας ἑξήκοντα περιβεβλημένοι σάκκους.
\vs{4}Οὗτοί εἰσιν αἱ δύο ἐλαῖαι καὶ αἱ δύο λυχνίαι αἱ ἐνώπιον τοῦ Κυρίου τῆς γῆς ἑστῶτες.
\vs{5}καὶ εἴ τις αὐτοὺς θέλει ἀδικῆσαι πῦρ ἐκπορεύεται ἐκ τοῦ στόματος αὐτῶν καὶ κατεσθίει τοὺς ἐχθροὺς αὐτῶν· καὶ εἴ τις θελήσῃ αὐτοὺς ἀδικῆσαι, οὕτως δεῖ αὐτὸν ἀποκτανθῆναι.
\vs{6}οὗτοι ἔχουσιν τὴν ἐξουσίαν κλεῖσαι τὸν οὐρανόν, ἵνα μὴ ὑετὸς βρέχῃ τὰς ἡμέρας τῆς προφητείας αὐτῶν, καὶ ἐξουσίαν ἔχουσιν ἐπὶ τῶν ὑδάτων στρέφειν αὐτὰ εἰς αἷμα καὶ πατάξαι τὴν γῆν ἐν πάσῃ πληγῇ ὁσάκις ἐὰν θελήσωσιν.

\vs{7}Καὶ ὅταν τελέσωσιν τὴν μαρτυρίαν αὐτῶν, τὸ θηρίον τὸ ἀναβαῖνον ἐκ τῆς ἀβύσσου ποιήσει μετ᾽ αὐτῶν πόλεμον καὶ νικήσει αὐτοὺς καὶ ἀποκτενεῖ αὐτούς.
\vs{8}καὶ τὸ πτῶμα αὐτῶν ἐπὶ τῆς πλατείας τῆς πόλεως τῆς μεγάλης, ἥτις καλεῖται πνευματικῶς Σόδομα καὶ Αἴγυπτος, ὅπου καὶ ὁ Κύριος αὐτῶν ἐσταυρώθη.
\vs{9}καὶ βλέπουσιν ἐκ τῶν λαῶν καὶ φυλῶν καὶ γλωσσῶν καὶ ἐθνῶν τὸ πτῶμα αὐτῶν ἡμέρας τρεῖς καὶ ἥμισυ καὶ τὰ πτώματα αὐτῶν οὐκ ἀφίουσιν τεθῆναι εἰς μνῆμα.
\vs{10}καὶ οἱ κατοικοῦντες ἐπὶ τῆς γῆς χαίρουσιν ἐπ᾽ αὐτοῖς καὶ εὐφραίνονται καὶ δῶρα πέμψουσιν ἀλλήλοις, ὅτι οὗτοι οἱ δύο προφῆται ἐβασάνισαν τοὺς κατοικοῦντας ἐπὶ τῆς γῆς.

\vs{11}Καὶ μετὰ τὰς τρεῖς ἡμέρας καὶ ἥμισυ πνεῦμα ζωῆς ἐκ τοῦ Θεοῦ εἰσῆλθεν ἐν αὐτοῖς, καὶ ἔστησαν ἐπὶ τοὺς πόδας αὐτῶν, καὶ φόβος μέγας ἐπέπεσεν ἐπὶ τοὺς θεωροῦντας αὐτούς.
\vs{12}καὶ ἤκουσαν φωνῆς μεγάλης ἐκ τοῦ οὐρανοῦ λεγούσης αὐτοῖς· Ἀνάβατε ὧδε. καὶ ἀνέβησαν εἰς τὸν οὐρανὸν ἐν τῇ νεφέλῃ, καὶ ἐθεώρησαν αὐτοὺς οἱ ἐχθροὶ αὐτῶν.
\vs{13}Καὶ ἐν ἐκείνῃ τῇ ὥρᾳ ἐγένετο σεισμὸς μέγας καὶ τὸ δέκατον τῆς πόλεως ἔπεσεν καὶ ἀπεκτάνθησαν ἐν τῷ σεισμῷ ὀνόματα ἀνθρώπων χιλιάδες ἑπτά καὶ οἱ λοιποὶ ἔμφοβοι ἐγένοντο καὶ ἔδωκαν δόξαν τῷ Θεῷ τοῦ οὐρανοῦ.

\vs{14}Ἡ Οὐαὶ ἡ δευτέρα ἀπῆλθεν· ἰδοὺ ἡ Οὐαὶ ἡ τρίτη ἔρχεται ταχύ.

\vs{15}Καὶ ὁ ἕβδομος ἄγγελος ἐσάλπισεν· καὶ ἐγένοντο φωναὶ μεγάλαι ἐν τῷ οὐρανῷ λέγοντες· 
\begin{poetryblock}

\begin{quote}Ἐγένετο ἡ βασιλεία τοῦ κόσμου τοῦ Κυρίου ἡμῶν\end{quote} 

\begin{quote}καὶ τοῦ Χριστοῦ αὐτοῦ,\end{quote} 

\begin{quote}καὶ βασιλεύσει εἰς τοὺς αἰῶνας τῶν αἰώνων.\end{quote}
\end{poetryblock}

\vs{16}Καὶ οἱ εἴκοσι τέσσαρες πρεσβύτεροι οἱ ἐνώπιον τοῦ Θεοῦ καθήμενοι ἐπὶ τοὺς θρόνους αὐτῶν ἔπεσαν ἐπὶ τὰ πρόσωπα αὐτῶν καὶ προσεκύνησαν τῷ Θεῷ
\vs{17}λέγοντες· 
\begin{poetryblock}

\begin{quote}Εὐχαριστοῦμέν σοι, Κύριε ὁ Θεός ὁ Παντοκράτωρ,\end{quote} 

\begin{quote}ὁ ὢν καὶ ὁ ἦν,\end{quote} 

\begin{quote}ὅτι εἴληφας τὴν δύναμίν σου τὴν μεγάλην\end{quote} 

\begin{quote}καὶ ἐβασίλευσας.\end{quote}

\begin{quote} \vs{18}καὶ τὰ ἔθνη ὠργίσθησαν,\end{quote} 

\begin{quote}καὶ ἦλθεν ἡ ὀργή σου\end{quote} 

\begin{quote}καὶ ὁ καιρὸς τῶν νεκρῶν κριθῆναι\end{quote} 

\begin{quote}καὶ δοῦναι τὸν μισθὸν τοῖς δούλοις σου τοῖς προφήταις\end{quote} 

\begin{quote}καὶ τοῖς ἁγίοις καὶ τοῖς φοβουμένοις τὸ ὄνομά σου,\end{quote} 

\begin{quote}τοὺς μικροὺς καὶ τοὺς μεγάλους,\end{quote} 

\begin{quote}καὶ διαφθεῖραι τοὺς διαφθείροντας τὴν γῆν.\end{quote}
\end{poetryblock}

\vs{19}Καὶ ἠνοίγη ὁ ναὸς τοῦ Θεοῦ ὁ ἐν τῷ οὐρανῷ καὶ ὤφθη ἡ κιβωτὸς τῆς διαθήκης αὐτοῦ ἐν τῷ ναῷ αὐτοῦ, καὶ ἐγένοντο ἀστραπαὶ καὶ φωναὶ καὶ βρονταὶ καὶ σεισμὸς καὶ χάλαζα μεγάλη.

\ch{12}
Καὶ σημεῖον μέγα ὤφθη ἐν τῷ οὐρανῷ, γυνὴ περιβεβλημένη τὸν ἥλιον, καὶ ἡ σελήνη ὑποκάτω τῶν ποδῶν αὐτῆς καὶ ἐπὶ τῆς κεφαλῆς αὐτῆς στέφανος ἀστέρων δώδεκα,
\vs{2}καὶ ἐν γαστρὶ ἔχουσα, καὶ κράζει ὠδίνουσα καὶ βασανιζομένη τεκεῖν.
\vs{3}Καὶ ὤφθη ἄλλο σημεῖον ἐν τῷ οὐρανῷ, καὶ ἰδοὺ δράκων μέγας πυρρός ἔχων κεφαλὰς ἑπτὰ καὶ κέρατα δέκα καὶ ἐπὶ τὰς κεφαλὰς αὐτοῦ ἑπτὰ διαδήματα,
\vs{4}καὶ ἡ οὐρὰ αὐτοῦ σύρει τὸ τρίτον τῶν ἀστέρων τοῦ οὐρανοῦ καὶ ἔβαλεν αὐτοὺς εἰς τὴν γῆν. καὶ ὁ δράκων ἕστηκεν ἐνώπιον τῆς γυναικὸς τῆς μελλούσης τεκεῖν, ἵνα ὅταν τέκῃ τὸ τέκνον αὐτῆς καταφάγῃ.
\vs{5}Καὶ ἔτεκεν υἱόν ἄρσεν, ὃς μέλλει ποιμαίνειν πάντα τὰ ἔθνη ἐν ῥάβδῳ σιδηρᾷ. καὶ ἡρπάσθη τὸ τέκνον αὐτῆς πρὸς τὸν Θεὸν καὶ πρὸς τὸν θρόνον αὐτοῦ.
\vs{6}καὶ ἡ γυνὴ ἔφυγεν εἰς τὴν ἔρημον, ὅπου ἔχει ἐκεῖ τόπον ἡτοιμασμένον ἀπὸ τοῦ Θεοῦ, ἵνα ἐκεῖ τρέφωσιν αὐτὴν ἡμέρας χιλίας διακοσίας ἑξήκοντα.
\vs{7}Καὶ ἐγένετο πόλεμος ἐν τῷ οὐρανῷ, ὁ Μιχαὴλ καὶ οἱ ἄγγελοι αὐτοῦ τοῦ πολεμῆσαι μετὰ τοῦ δράκοντος. καὶ ὁ δράκων ἐπολέμησεν καὶ οἱ ἄγγελοι αὐτοῦ,
\vs{8}καὶ οὐκ ἴσχυσεν οὐδὲ τόπος εὑρέθη αὐτῶν ἔτι ἐν τῷ οὐρανῷ.
\vs{9}καὶ ἐβλήθη ὁ δράκων ὁ μέγας, ὁ ὄφις ὁ ἀρχαῖος, ὁ καλούμενος Διάβολος καὶ Ὁ Σατανᾶς, ὁ πλανῶν τὴν οἰκουμένην ὅλην, ἐβλήθη εἰς τὴν γῆν, καὶ οἱ ἄγγελοι αὐτοῦ μετ᾽ αὐτοῦ ἐβλήθησαν.
\vs{10}Καὶ ἤκουσα φωνὴν μεγάλην ἐν τῷ οὐρανῷ λέγουσαν· 
\begin{poetryblock}

\begin{quote}Ἄρτι ἐγένετο ἡ σωτηρία καὶ ἡ δύναμις\end{quote} 

\begin{quote}καὶ ἡ βασιλεία τοῦ Θεοῦ ἡμῶν\end{quote} 

\begin{quote}καὶ ἡ ἐξουσία τοῦ Χριστοῦ αὐτοῦ,\end{quote} 

\begin{quote}ὅτι ἐβλήθη ὁ κατήγωρ τῶν ἀδελφῶν ἡμῶν,\end{quote} 

\begin{quote}ὁ κατηγορῶν αὐτοὺς ἐνώπιον τοῦ Θεοῦ ἡμῶν ἡμέρας καὶ νυκτός.\end{quote}

\begin{quote} \vs{11}καὶ αὐτοὶ ἐνίκησαν αὐτὸν διὰ τὸ αἷμα τοῦ Ἀρνίου\end{quote} 

\begin{quote}καὶ διὰ τὸν λόγον τῆς μαρτυρίας αὐτῶν\end{quote} 

\begin{quote}καὶ οὐκ ἠγάπησαν τὴν ψυχὴν αὐτῶν ἄχρι θανάτου.\end{quote}

\begin{quote} \vs{12}διὰ τοῦτο εὐφραίνεσθε, οἱ οὐρανοὶ\end{quote} 

\begin{quote}καὶ οἱ ἐν αὐτοῖς σκηνοῦντες.\end{quote} 

\begin{quote}οὐαὶ τὴν γῆν καὶ τὴν θάλασσαν,\end{quote} 

\begin{quote}ὅτι κατέβη ὁ διάβολος πρὸς ὑμᾶς\end{quote} 

\begin{quote}ἔχων θυμὸν μέγαν,\end{quote} 

\begin{quote}εἰδὼς ὅτι ὀλίγον καιρὸν ἔχει.\end{quote}
\end{poetryblock}

\vs{13}Καὶ ὅτε εἶδεν ὁ δράκων ὅτι ἐβλήθη εἰς τὴν γῆν, ἐδίωξεν τὴν γυναῖκα ἥτις ἔτεκεν τὸν ἄρσενα.
\vs{14}καὶ ἐδόθησαν τῇ γυναικὶ αἱ δύο πτέρυγες τοῦ ἀετοῦ τοῦ μεγάλου, ἵνα πέτηται εἰς τὴν ἔρημον εἰς τὸν τόπον αὐτῆς, ὅπου τρέφεται ἐκεῖ καιρὸν καὶ καιροὺς καὶ ἥμισυ καιροῦ ἀπὸ προσώπου τοῦ ὄφεως.
\vs{15}Καὶ ἔβαλεν ὁ ὄφις ἐκ τοῦ στόματος αὐτοῦ ὀπίσω τῆς γυναικὸς ὕδωρ ὡς ποταμόν, ἵνα αὐτὴν ποταμοφόρητον ποιήσῃ.
\vs{16}καὶ ἐβοήθησεν ἡ γῆ τῇ γυναικί καὶ ἤνοιξεν ἡ γῆ τὸ στόμα αὐτῆς καὶ κατέπιεν τὸν ποταμὸν ὃν ἔβαλεν ὁ δράκων ἐκ τοῦ στόματος αὐτοῦ.
\vs{17}καὶ ὠργίσθη ὁ δράκων ἐπὶ τῇ γυναικί καὶ ἀπῆλθεν ποιῆσαι πόλεμον μετὰ τῶν λοιπῶν τοῦ σπέρματος αὐτῆς τῶν τηρούντων τὰς ἐντολὰς τοῦ Θεοῦ καὶ ἐχόντων τὴν μαρτυρίαν Ἰησοῦ.

\vs{18}Καὶ ἐστάθη ἐπὶ τὴν ἄμμον τῆς θαλάσσης.

\ch{13}
Καὶ εἶδον ἐκ τῆς θαλάσσης θηρίον ἀναβαῖνον, ἔχον κέρατα δέκα καὶ κεφαλὰς ἑπτά καὶ ἐπὶ τῶν κεράτων αὐτοῦ δέκα διαδήματα καὶ ἐπὶ τὰς κεφαλὰς αὐτοῦ ὀνόματα βλασφημίας.
\vs{2}καὶ τὸ θηρίον ὃ εἶδον ἦν ὅμοιον παρδάλει καὶ οἱ πόδες αὐτοῦ ὡς ἄρκου καὶ τὸ στόμα αὐτοῦ ὡς στόμα λέοντος. καὶ ἔδωκεν αὐτῷ ὁ δράκων τὴν δύναμιν αὐτοῦ καὶ τὸν θρόνον αὐτοῦ καὶ ἐξουσίαν μεγάλην.
\vs{3}Καὶ μίαν ἐκ τῶν κεφαλῶν αὐτοῦ ὡς ἐσφαγμένην εἰς θάνατον, καὶ ἡ πληγὴ τοῦ θανάτου αὐτοῦ ἐθεραπεύθη.

καὶ ἐθαυμάσθη ὅλη ἡ γῆ ὀπίσω τοῦ θηρίου
\vs{4}καὶ προσεκύνησαν τῷ δράκοντι, ὅτι ἔδωκεν τὴν ἐξουσίαν τῷ θηρίῳ, καὶ προσεκύνησαν τῷ θηρίῳ λέγοντες· Τίς ὅμοιος τῷ θηρίῳ καὶ τίς δύναται πολεμῆσαι μετ᾽ αὐτοῦ;

\vs{5}Καὶ ἐδόθη αὐτῷ στόμα λαλοῦν μεγάλα καὶ βλασφημίας καὶ ἐδόθη αὐτῷ ἐξουσία ποιῆσαι μῆνας τεσσεράκοντα καὶ δύο.
\vs{6}καὶ ἤνοιξεν τὸ στόμα αὐτοῦ εἰς βλασφημίας πρὸς τὸν Θεόν βλασφημῆσαι τὸ ὄνομα αὐτοῦ καὶ τὴν σκηνὴν αὐτοῦ, τοὺς ἐν τῷ οὐρανῷ σκηνοῦντας.
\vs{7}Καὶ ἐδόθη αὐτῷ ποιῆσαι πόλεμον μετὰ τῶν ἁγίων καὶ νικῆσαι αὐτούς, καὶ ἐδόθη αὐτῷ ἐξουσία ἐπὶ πᾶσαν φυλὴν καὶ λαὸν καὶ γλῶσσαν καὶ ἔθνος.
\vs{8}καὶ προσκυνήσουσιν αὐτὸν πάντες οἱ κατοικοῦντες ἐπὶ τῆς γῆς, οὗ οὐ γέγραπται τὸ ὄνομα αὐτοῦ ἐν τῷ βιβλίῳ τῆς ζωῆς τοῦ Ἀρνίου τοῦ ἐσφαγμένου ἀπὸ καταβολῆς κόσμου.

\vs{9}Εἴ τις ἔχει οὖς ἀκουσάτω.

\vs{10}Εἴ τις εἰς αἰχμαλωσίαν, εἰς αἰχμαλωσίαν ὑπάγει· 
\begin{poetryblock}

\begin{quote}εἴ τις ἐν μαχαίρῃ ἀποκτανθῆναι αὐτὸν ἐν μαχαίρῃ ἀποκτανθῆναι.\end{quote}
\end{poetryblock}

Ὧδέ ἐστιν ἡ ὑπομονὴ καὶ ἡ πίστις τῶν ἁγίων.

\vs{11}Καὶ εἶδον ἄλλο θηρίον ἀναβαῖνον ἐκ τῆς γῆς, καὶ εἶχεν κέρατα δύο ὅμοια ἀρνίῳ καὶ ἐλάλει ὡς δράκων.
\vs{12}καὶ τὴν ἐξουσίαν τοῦ πρώτου θηρίου πᾶσαν ποιεῖ ἐνώπιον αὐτοῦ, καὶ ποιεῖ τὴν γῆν καὶ τοὺς ἐν αὐτῇ κατοικοῦντας ἵνα προσκυνήσουσιν τὸ θηρίον τὸ πρῶτον, οὗ ἐθεραπεύθη ἡ πληγὴ τοῦ θανάτου αὐτοῦ.
\vs{13}Καὶ ποιεῖ σημεῖα μεγάλα, ἵνα καὶ πῦρ ποιῇ ἐκ τοῦ οὐρανοῦ καταβαίνειν εἰς τὴν γῆν ἐνώπιον τῶν ἀνθρώπων,
\vs{14}καὶ πλανᾷ τοὺς κατοικοῦντας ἐπὶ τῆς γῆς διὰ τὰ σημεῖα ἃ ἐδόθη αὐτῷ ποιῆσαι ἐνώπιον τοῦ θηρίου, λέγων τοῖς κατοικοῦσιν ἐπὶ τῆς γῆς ποιῆσαι εἰκόνα τῷ θηρίῳ, ὃς ἔχει τὴν πληγὴν τῆς μαχαίρης καὶ ἔζησεν.

\vs{15}καὶ ἐδόθη αὐτῷ δοῦναι πνεῦμα τῇ εἰκόνι τοῦ θηρίου, ἵνα καὶ λαλήσῃ ἡ εἰκὼν τοῦ θηρίου καὶ ποιήσῃ ἵνα ὅσοι ἐὰν μὴ προσκυνήσωσιν τῇ εἰκόνι τοῦ θηρίου ἀποκτανθῶσιν.
\vs{16}Καὶ ποιεῖ πάντας, τοὺς μικροὺς καὶ τοὺς μεγάλους, καὶ τοὺς πλουσίους καὶ τοὺς πτωχούς, καὶ τοὺς ἐλευθέρους καὶ τοὺς δούλους, ἵνα δῶσιν αὐτοῖς χάραγμα ἐπὶ τῆς χειρὸς αὐτῶν τῆς δεξιᾶς ἢ ἐπὶ τὸ μέτωπον αὐτῶν
\vs{17}καὶ ἵνα μή τις δύνηται ἀγοράσαι ἢ πωλῆσαι εἰ μὴ ὁ ἔχων τὸ χάραγμα τὸ ὄνομα τοῦ θηρίου ἢ τὸν ἀριθμὸν τοῦ ὀνόματος αὐτοῦ.

\vs{18}Ὧδε ἡ σοφία ἐστίν. ὁ ἔχων νοῦν ψηφισάτω τὸν ἀριθμὸν τοῦ θηρίου, ἀριθμὸς γὰρ ἀνθρώπου ἐστίν, καὶ ὁ ἀριθμὸς αὐτοῦ ἑξακόσιοι ἑξήκοντα ἕξ.

\ch{14}
Καὶ εἶδον, καὶ ἰδοὺ τὸ Ἀρνίον ἑστὸς ἐπὶ τὸ ὄρος Σιών καὶ μετ᾽ αὐτοῦ ἑκατὸν τεσσεράκοντα τέσσαρες χιλιάδες ἔχουσαι τὸ ὄνομα αὐτοῦ καὶ τὸ ὄνομα τοῦ Πατρὸς αὐτοῦ γεγραμμένον ἐπὶ τῶν μετώπων αὐτῶν.
\vs{2}καὶ ἤκουσα φωνὴν ἐκ τοῦ οὐρανοῦ ὡς φωνὴν ὑδάτων πολλῶν καὶ ὡς φωνὴν βροντῆς μεγάλης, καὶ ἡ φωνὴ ἣν ἤκουσα ὡς κιθαρῳδῶν κιθαριζόντων ἐν ταῖς κιθάραις αὐτῶν.
\vs{3}Καὶ ᾄδουσιν ὡς ᾠδὴν καινὴν ἐνώπιον τοῦ θρόνου καὶ ἐνώπιον τῶν τεσσάρων ζῴων καὶ τῶν πρεσβυτέρων, καὶ οὐδεὶς ἐδύνατο μαθεῖν τὴν ᾠδὴν εἰ μὴ αἱ ἑκατὸν τεσσεράκοντα τέσσαρες χιλιάδες, οἱ ἠγορασμένοι ἀπὸ τῆς γῆς.

\vs{4}οὗτοί εἰσιν οἳ μετὰ γυναικῶν οὐκ ἐμολύνθησαν, παρθένοι γάρ εἰσιν, οὗτοι οἱ ἀκολουθοῦντες τῷ Ἀρνίῳ ὅπου ἂν ὑπάγῃ. οὗτοι ἠγοράσθησαν ἀπὸ τῶν ἀνθρώπων ἀπαρχὴ τῷ Θεῷ καὶ τῷ Ἀρνίῳ,
\vs{5}καὶ ἐν τῷ στόματι αὐτῶν οὐχ εὑρέθη ψεῦδος, ἄμωμοί εἰσιν.

\vs{6}Καὶ εἶδον ἄλλον ἄγγελον πετόμενον ἐν μεσουρανήματι, ἔχοντα εὐαγγέλιον αἰώνιον εὐαγγελίσαι ἐπὶ τοὺς καθημένους ἐπὶ τῆς γῆς καὶ ἐπὶ πᾶν ἔθνος καὶ φυλὴν καὶ γλῶσσαν καὶ λαόν,
\vs{7}λέγων ἐν φωνῇ μεγάλῃ·

Φοβήθητε τὸν Θεὸν καὶ δότε αὐτῷ δόξαν, ὅτι ἦλθεν ἡ ὥρα τῆς κρίσεως αὐτοῦ, καὶ προσκυνήσατε τῷ ποιήσαντι τὸν οὐρανὸν καὶ τὴν γῆν καὶ θάλασσαν καὶ πηγὰς ὑδάτων.

\vs{8}Καὶ ἄλλος ἄγγελος δεύτερος ἠκολούθησεν λέγων· 
\begin{poetryblock}

\begin{quote}Ἔπεσεν ἔπεσεν Βαβυλὼν ἡ μεγάλη ἣ ἐκ τοῦ οἴνου τοῦ θυμοῦ τῆς πορνείας αὐτῆς πεπότικεν πάντα τὰ ἔθνη.\end{quote}
\end{poetryblock}

\vs{9}Καὶ ἄλλος ἄγγελος τρίτος ἠκολούθησεν αὐτοῖς λέγων ἐν φωνῇ μεγάλῃ·

Εἴ τις προσκυνεῖ τὸ θηρίον καὶ τὴν εἰκόνα αὐτοῦ καὶ λαμβάνει χάραγμα ἐπὶ τοῦ μετώπου αὐτοῦ ἢ ἐπὶ τὴν χεῖρα αὐτοῦ,
\vs{10}καὶ αὐτὸς πίεται ἐκ τοῦ οἴνου τοῦ θυμοῦ τοῦ Θεοῦ τοῦ κεκερασμένου ἀκράτου ἐν τῷ ποτηρίῳ τῆς ὀργῆς αὐτοῦ καὶ βασανισθήσεται ἐν πυρὶ καὶ θείῳ ἐνώπιον ἀγγέλων ἁγίων καὶ ἐνώπιον τοῦ Ἀρνίου.
\vs{11}καὶ ὁ καπνὸς τοῦ βασανισμοῦ αὐτῶν εἰς αἰῶνας αἰώνων ἀναβαίνει, καὶ οὐκ ἔχουσιν ἀνάπαυσιν ἡμέρας καὶ νυκτός οἱ προσκυνοῦντες τὸ θηρίον καὶ τὴν εἰκόνα αὐτοῦ καὶ εἴ τις λαμβάνει τὸ χάραγμα τοῦ ὀνόματος αὐτοῦ.
\vs{12}Ὧδε ἡ ὑπομονὴ τῶν ἁγίων ἐστίν, οἱ τηροῦντες τὰς ἐντολὰς τοῦ Θεοῦ καὶ τὴν πίστιν Ἰησοῦ.

\vs{13}Καὶ ἤκουσα φωνῆς ἐκ τοῦ οὐρανοῦ λεγούσης· Γράψον· 
\begin{poetryblock}

\begin{quote}Μακάριοι οἱ νεκροὶ οἱ ἐν Κυρίῳ ἀποθνῄσκοντες ἀπ᾽ ἄρτι. Ναί, λέγει τὸ Πνεῦμα, Ἵνα ἀναπαήσονται ἐκ τῶν κόπων αὐτῶν, τὰ γὰρ ἔργα αὐτῶν ἀκολουθεῖ μετ᾽ αὐτῶν.\end{quote}
\end{poetryblock}

\vs{14}Καὶ εἶδον, καὶ ἰδοὺ νεφέλη λευκή, καὶ ἐπὶ τὴν νεφέλην καθήμενον ὅμοιον υἱὸν ἀνθρώπου, ἔχων ἐπὶ τῆς κεφαλῆς αὐτοῦ στέφανον χρυσοῦν καὶ ἐν τῇ χειρὶ αὐτοῦ δρέπανον ὀξύ.
\vs{15}Καὶ ἄλλος ἄγγελος ἐξῆλθεν ἐκ τοῦ ναοῦ κράζων ἐν φωνῇ μεγάλῃ τῷ καθημένῳ ἐπὶ τῆς νεφέλης·

Πέμψον τὸ δρέπανόν σου καὶ θέρισον, ὅτι ἦλθεν ἡ ὥρα θερίσαι, ὅτι ἐξηράνθη ὁ θερισμὸς τῆς γῆς.
\vs{16}καὶ ἔβαλεν ὁ καθήμενος ἐπὶ τῆς νεφέλης τὸ δρέπανον αὐτοῦ ἐπὶ τὴν γῆν καὶ ἐθερίσθη ἡ γῆ.

\vs{17}Καὶ ἄλλος ἄγγελος ἐξῆλθεν ἐκ τοῦ ναοῦ τοῦ ἐν τῷ οὐρανῷ ἔχων καὶ αὐτὸς δρέπανον ὀξύ.
\vs{18}Καὶ ἄλλος ἄγγελος ἐξῆλθεν ἐκ τοῦ θυσιαστηρίου ὁ ἔχων ἐξουσίαν ἐπὶ τοῦ πυρός, καὶ ἐφώνησεν φωνῇ μεγάλῃ τῷ ἔχοντι τὸ δρέπανον τὸ ὀξὺ λέγων· Πέμψον σου τὸ δρέπανον τὸ ὀξὺ καὶ τρύγησον τοὺς βότρυας τῆς ἀμπέλου τῆς γῆς, ὅτι ἤκμασαν αἱ σταφυλαὶ αὐτῆς.
\vs{19}Καὶ ἔβαλεν ὁ ἄγγελος τὸ δρέπανον αὐτοῦ εἰς τὴν γῆν καὶ ἐτρύγησεν τὴν ἄμπελον τῆς γῆς καὶ ἔβαλεν εἰς τὴν ληνὸν τοῦ θυμοῦ τοῦ Θεοῦ τὸν μέγαν.
\vs{20}καὶ ἐπατήθη ἡ ληνὸς ἔξωθεν τῆς πόλεως καὶ ἐξῆλθεν αἷμα ἐκ τῆς ληνοῦ ἄχρι τῶν χαλινῶν τῶν ἵππων ἀπὸ σταδίων χιλίων ἑξακοσίων.

\ch{15}
Καὶ εἶδον ἄλλο σημεῖον ἐν τῷ οὐρανῷ μέγα καὶ θαυμαστόν, ἀγγέλους ἑπτὰ ἔχοντας πληγὰς ἑπτὰ τὰς ἐσχάτας, ὅτι ἐν αὐταῖς ἐτελέσθη ὁ θυμὸς τοῦ Θεοῦ.

\vs{2}Καὶ εἶδον ὡς θάλασσαν ὑαλίνην μεμιγμένην πυρί καὶ τοὺς νικῶντας ἐκ τοῦ θηρίου καὶ ἐκ τῆς εἰκόνος αὐτοῦ καὶ ἐκ τοῦ ἀριθμοῦ τοῦ ὀνόματος αὐτοῦ ἑστῶτας ἐπὶ τὴν θάλασσαν τὴν ὑαλίνην ἔχοντας κιθάρας τοῦ Θεοῦ.
\vs{3}καὶ ᾄδουσιν τὴν ᾠδὴν Μωϋσέως τοῦ δούλου τοῦ Θεοῦ καὶ τὴν ᾠδὴν τοῦ Ἀρνίου λέγοντες· 
\begin{poetryblock}

\begin{quote}Μεγάλα καὶ θαυμαστὰ τὰ ἔργα σου,\end{quote} 

\begin{quote}Κύριε ὁ Θεός ὁ Παντοκράτωρ·\end{quote} 

\begin{quote}δίκαιαι καὶ ἀληθιναὶ αἱ ὁδοί σου,\end{quote} 

\begin{quote}ὁ Βασιλεὺς τῶν ἐθνῶν·\end{quote}

\begin{quote} \vs{4}τίς οὐ μὴ φοβηθῇ, Κύριε,\end{quote} 

\begin{quote}καὶ δοξάσει τὸ ὄνομά σου;\end{quote} 

\begin{quote}ὅτι μόνος ὅσιος,\end{quote} 

\begin{quote}ὅτι πάντα τὰ ἔθνη ἥξουσιν\end{quote} 

\begin{quote}καὶ προσκυνήσουσιν ἐνώπιόν σου,\end{quote} 

\begin{quote}ὅτι τὰ δικαιώματά σου ἐφανερώθησαν.\end{quote}
\end{poetryblock}

\vs{5}Καὶ μετὰ ταῦτα εἶδον, καὶ ἠνοίγη ὁ ναὸς τῆς σκηνῆς τοῦ μαρτυρίου ἐν τῷ οὐρανῷ,
\vs{6}καὶ ἐξῆλθον οἱ ἑπτὰ ἄγγελοι οἱ ἔχοντες τὰς ἑπτὰ πληγὰς ἐκ τοῦ ναοῦ ἐνδεδυμένοι λίνον καθαρὸν λαμπρὸν καὶ περιεζωσμένοι περὶ τὰ στήθη ζώνας χρυσᾶς.
\vs{7}Καὶ ἓν ἐκ τῶν τεσσάρων ζῴων ἔδωκεν τοῖς ἑπτὰ ἀγγέλοις ἑπτὰ φιάλας χρυσᾶς γεμούσας τοῦ θυμοῦ τοῦ Θεοῦ τοῦ ζῶντος εἰς τοὺς αἰῶνας τῶν αἰώνων.
\vs{8}καὶ ἐγεμίσθη ὁ ναὸς καπνοῦ ἐκ τῆς δόξης τοῦ Θεοῦ καὶ ἐκ τῆς δυνάμεως αὐτοῦ, καὶ οὐδεὶς ἐδύνατο εἰσελθεῖν εἰς τὸν ναὸν ἄχρι τελεσθῶσιν αἱ ἑπτὰ πληγαὶ τῶν ἑπτὰ ἀγγέλων.

\ch{16}
Καὶ ἤκουσα μεγάλης φωνῆς ἐκ τοῦ ναοῦ λεγούσης τοῖς ἑπτὰ ἀγγέλοις· Ὑπάγετε καὶ ἐκχέετε τὰς ἑπτὰ φιάλας τοῦ θυμοῦ τοῦ Θεοῦ εἰς τὴν γῆν.

\vs{2}Καὶ ἀπῆλθεν ὁ πρῶτος καὶ ἐξέχεεν τὴν φιάλην αὐτοῦ εἰς τὴν γῆν, καὶ ἐγένετο ἕλκος κακὸν καὶ πονηρὸν ἐπὶ τοὺς ἀνθρώπους τοὺς ἔχοντας τὸ χάραγμα τοῦ θηρίου καὶ τοὺς προσκυνοῦντας τῇ εἰκόνι αὐτοῦ.

\vs{3}Καὶ ὁ δεύτερος ἐξέχεεν τὴν φιάλην αὐτοῦ εἰς τὴν θάλασσαν, καὶ ἐγένετο αἷμα ὡς νεκροῦ, καὶ πᾶσα ψυχὴ ζωῆς ἀπέθανεν τὰ ἐν τῇ θαλάσσῃ.

\vs{4}Καὶ ὁ τρίτος ἐξέχεεν τὴν φιάλην αὐτοῦ εἰς τοὺς ποταμοὺς καὶ τὰς πηγὰς τῶν ὑδάτων, καὶ ἐγένετο αἷμα.
\vs{5}Καὶ ἤκουσα τοῦ ἀγγέλου τῶν ὑδάτων λέγοντος· 
\begin{poetryblock}

\begin{quote}Δίκαιος εἶ, ὁ ὢν καὶ ὁ ἦν, ὁ Ὅσιος,\end{quote} 

\begin{quote}ὅτι ταῦτα ἔκρινας,\end{quote}

\begin{quote} \vs{6}ὅτι αἷμα ἁγίων καὶ προφητῶν ἐξέχεαν\end{quote} 

\begin{quote}καὶ αἷμα αὐτοῖς δέδωκας πιεῖν,\end{quote} 

\begin{quote}ἄξιοί εἰσιν.\end{quote}
\end{poetryblock}

\vs{7}Καὶ ἤκουσα τοῦ θυσιαστηρίου λέγοντος· 
\begin{poetryblock}

\begin{quote}Ναί Κύριε ὁ Θεός ὁ Παντοκράτωρ,\end{quote} 

\begin{quote}ἀληθιναὶ καὶ δίκαιαι αἱ κρίσεις σου.\end{quote}
\end{poetryblock}

\vs{8}Καὶ ὁ τέταρτος ἐξέχεεν τὴν φιάλην αὐτοῦ ἐπὶ τὸν ἥλιον, καὶ ἐδόθη αὐτῷ καυματίσαι τοὺς ἀνθρώπους ἐν πυρί.
\vs{9}καὶ ἐκαυματίσθησαν οἱ ἄνθρωποι καῦμα μέγα καὶ ἐβλασφήμησαν τὸ ὄνομα τοῦ Θεοῦ τοῦ ἔχοντος τὴν ἐξουσίαν ἐπὶ τὰς πληγὰς ταύτας καὶ οὐ μετενόησαν δοῦναι αὐτῷ δόξαν.

\vs{10}Καὶ ὁ πέμπτος ἐξέχεεν τὴν φιάλην αὐτοῦ ἐπὶ τὸν θρόνον τοῦ θηρίου, καὶ ἐγένετο ἡ βασιλεία αὐτοῦ ἐσκοτωμένη, καὶ ἐμασῶντο τὰς γλώσσας αὐτῶν ἐκ τοῦ πόνου,
\vs{11}καὶ ἐβλασφήμησαν τὸν Θεὸν τοῦ οὐρανοῦ ἐκ τῶν πόνων αὐτῶν καὶ ἐκ τῶν ἑλκῶν αὐτῶν καὶ οὐ μετενόησαν ἐκ τῶν ἔργων αὐτῶν.

\vs{12}Καὶ ὁ ἕκτος ἐξέχεεν τὴν φιάλην αὐτοῦ ἐπὶ τὸν ποταμὸν τὸν μέγαν τὸν Εὐφράτην, καὶ ἐξηράνθη τὸ ὕδωρ αὐτοῦ, ἵνα ἑτοιμασθῇ ἡ ὁδὸς τῶν βασιλέων τῶν ἀπὸ ἀνατολῆς ἡλίου.
\vs{13}Καὶ εἶδον ἐκ τοῦ στόματος τοῦ δράκοντος καὶ ἐκ τοῦ στόματος τοῦ θηρίου καὶ ἐκ τοῦ στόματος τοῦ ψευδοπροφήτου πνεύματα τρία ἀκάθαρτα ὡς βάτραχοι·
\vs{14}εἰσὶν γὰρ πνεύματα δαιμονίων ποιοῦντα σημεῖα, ἃ ἐκπορεύεται ἐπὶ τοὺς βασιλεῖς τῆς οἰκουμένης ὅλης συναγαγεῖν αὐτοὺς εἰς τὸν πόλεμον τῆς ἡμέρας τῆς μεγάλης τοῦ Θεοῦ τοῦ Παντοκράτορος.
\vs{15}Ἰδοὺ ἔρχομαι ὡς κλέπτης. μακάριος ὁ γρηγορῶν καὶ τηρῶν τὰ ἱμάτια αὐτοῦ, ἵνα μὴ γυμνὸς περιπατῇ καὶ βλέπωσιν τὴν ἀσχημοσύνην αὐτοῦ.
\vs{16}Καὶ συνήγαγεν αὐτοὺς εἰς τὸν τόπον τὸν καλούμενον Ἑβραϊστὶ Ἁρμαγεδών.

\vs{17}Καὶ ὁ ἕβδομος ἐξέχεεν τὴν φιάλην αὐτοῦ ἐπὶ τὸν ἀέρα, καὶ ἐξῆλθεν φωνὴ μεγάλη ἐκ τοῦ ναοῦ ἀπὸ τοῦ θρόνου λέγουσα· Γέγονεν.
\vs{18}Καὶ ἐγένοντο ἀστραπαὶ καὶ φωναὶ καὶ βρονταί καὶ σεισμὸς ἐγένετο μέγας, οἷος οὐκ ἐγένετο ἀφ᾽ οὗ ἄνθρωπος ἐγένετο ἐπὶ τῆς γῆς τηλικοῦτος σεισμὸς οὕτω μέγας.
\vs{19}καὶ ἐγένετο ἡ πόλις ἡ μεγάλη εἰς τρία μέρη καὶ αἱ πόλεις τῶν ἐθνῶν ἔπεσαν. καὶ Βαβυλὼν ἡ μεγάλη ἐμνήσθη ἐνώπιον τοῦ Θεοῦ δοῦναι αὐτῇ τὸ ποτήριον τοῦ οἴνου τοῦ θυμοῦ τῆς ὀργῆς αὐτοῦ.
\vs{20}Καὶ πᾶσα νῆσος ἔφυγεν καὶ ὄρη οὐχ εὑρέθησαν.
\vs{21}καὶ χάλαζα μεγάλη ὡς ταλαντιαία καταβαίνει ἐκ τοῦ οὐρανοῦ ἐπὶ τοὺς ἀνθρώπους, καὶ ἐβλασφήμησαν οἱ ἄνθρωποι τὸν Θεὸν ἐκ τῆς πληγῆς τῆς χαλάζης, ὅτι μεγάλη ἐστὶν ἡ πληγὴ αὐτῆς σφόδρα.

\ch{17}
Καὶ ἦλθεν εἷς ἐκ τῶν ἑπτὰ ἀγγέλων τῶν ἐχόντων τὰς ἑπτὰ φιάλας καὶ ἐλάλησεν μετ᾽ ἐμοῦ λέγων· Δεῦρο, δείξω σοι τὸ κρίμα τῆς πόρνης τῆς μεγάλης τῆς καθημένης ἐπὶ ὑδάτων πολλῶν,
\vs{2}μεθ᾽ ἧς ἐπόρνευσαν οἱ βασιλεῖς τῆς γῆς καὶ ἐμεθύσθησαν οἱ κατοικοῦντες τὴν γῆν ἐκ τοῦ οἴνου τῆς πορνείας αὐτῆς.
\vs{3}Καὶ ἀπήνεγκέν με εἰς ἔρημον ἐν Πνεύματι.

καὶ εἶδον γυναῖκα καθημένην ἐπὶ θηρίον κόκκινον, γέμοντα ὀνόματα βλασφημίας, ἔχων κεφαλὰς ἑπτὰ καὶ κέρατα δέκα.
\vs{4}καὶ ἡ γυνὴ ἦν περιβεβλημένη πορφυροῦν καὶ κόκκινον καὶ κεχρυσωμένη χρυσίῳ καὶ λίθῳ τιμίῳ καὶ μαργαρίταις, ἔχουσα ποτήριον χρυσοῦν ἐν τῇ χειρὶ αὐτῆς γέμον βδελυγμάτων καὶ τὰ ἀκάθαρτα τῆς πορνείας αὐτῆς
\vs{5}καὶ ἐπὶ τὸ μέτωπον αὐτῆς ὄνομα γεγραμμένον, μυστήριον, ΒΑΒΥΛΩΝ ἡ ΜΕΓΑΛΗ, ἡ ΜΗΤΗΡ ΤΩΝ ΠΟΡΝΩΝ ΚΑΙ ΤΩΝ ΒΔΕΛΥΓΜΑΤΩΝ ΤΗΣ ΓΗΣ.
\vs{6}Καὶ εἶδον τὴν γυναῖκα μεθύουσαν ἐκ τοῦ αἵματος τῶν ἁγίων καὶ ἐκ τοῦ αἵματος τῶν μαρτύρων Ἰησοῦ. Καὶ ἐθαύμασα ἰδὼν αὐτὴν θαῦμα μέγα.

\vs{7}Καὶ εἶπέν μοι ὁ ἄγγελος· Διὰ τί ἐθαύμασας; ἐγὼ ἐρῶ σοι τὸ μυστήριον τῆς γυναικὸς καὶ τοῦ θηρίου τοῦ βαστάζοντος αὐτήν τοῦ ἔχοντος τὰς ἑπτὰ κεφαλὰς καὶ τὰ δέκα κέρατα.

\vs{8}Τὸ θηρίον ὃ εἶδες ἦν καὶ οὐκ ἔστιν καὶ μέλλει ἀναβαίνειν ἐκ τῆς ἀβύσσου καὶ εἰς ἀπώλειαν ὑπάγει, καὶ θαυμασθήσονται οἱ κατοικοῦντες ἐπὶ τῆς γῆς, ὧν οὐ γέγραπται τὸ ὄνομα ἐπὶ τὸ βιβλίον τῆς ζωῆς ἀπὸ καταβολῆς κόσμου, βλεπόντων τὸ θηρίον ὅτι ἦν καὶ οὐκ ἔστιν καὶ παρέσται.
\vs{9}Ὧδε ὁ νοῦς ὁ ἔχων σοφίαν. αἱ ἑπτὰ κεφαλαὶ ἑπτὰ ὄρη εἰσίν, ὅπου ἡ γυνὴ κάθηται ἐπ᾽ αὐτῶν. καὶ βασιλεῖς ἑπτά εἰσιν·
\vs{10}οἱ πέντε ἔπεσαν, ὁ εἷς ἔστιν, ὁ ἄλλος οὔπω ἦλθεν, καὶ ὅταν ἔλθῃ ὀλίγον αὐτὸν δεῖ μεῖναι.
\vs{11}Καὶ τὸ θηρίον ὃ ἦν καὶ οὐκ ἔστιν καὶ αὐτὸς ὄγδοός ἐστιν καὶ ἐκ τῶν ἑπτά ἐστιν, καὶ εἰς ἀπώλειαν ὑπάγει.
\vs{12}καὶ τὰ δέκα κέρατα ἃ εἶδες δέκα βασιλεῖς εἰσιν, οἵτινες βασιλείαν οὔπω ἔλαβον, ἀλλὰ ἐξουσίαν ὡς βασιλεῖς μίαν ὥραν λαμβάνουσιν μετὰ τοῦ θηρίου.
\vs{13}οὗτοι μίαν γνώμην ἔχουσιν καὶ τὴν δύναμιν καὶ ἐξουσίαν αὐτῶν τῷ θηρίῳ διδόασιν.
\vs{14}Οὗτοι μετὰ τοῦ Ἀρνίου πολεμήσουσιν καὶ τὸ Ἀρνίον νικήσει αὐτούς, ὅτι Κύριος κυρίων ἐστὶν καὶ Βασιλεὺς βασιλέων καὶ οἱ μετ᾽ αὐτοῦ κλητοὶ καὶ ἐκλεκτοὶ καὶ πιστοί.

\vs{15}Καὶ λέγει μοι· Τὰ ὕδατα ἃ εἶδες οὗ ἡ πόρνη κάθηται, λαοὶ καὶ ὄχλοι εἰσὶν καὶ ἔθνη καὶ γλῶσσαι.
\vs{16}καὶ τὰ δέκα κέρατα ἃ εἶδες καὶ τὸ θηρίον οὗτοι μισήσουσιν τὴν πόρνην καὶ ἠρημωμένην ποιήσουσιν αὐτὴν καὶ γυμνήν καὶ τὰς σάρκας αὐτῆς φάγονται καὶ αὐτὴν κατακαύσουσιν ἐν πυρί.
\vs{17}ὁ γὰρ Θεὸς ἔδωκεν εἰς τὰς καρδίας αὐτῶν ποιῆσαι τὴν γνώμην αὐτοῦ καὶ ποιῆσαι μίαν γνώμην καὶ δοῦναι τὴν βασιλείαν αὐτῶν τῷ θηρίῳ ἄχρι τελεσθήσονται οἱ λόγοι τοῦ Θεοῦ.
\vs{18}καὶ ἡ γυνὴ ἣν εἶδες ἔστιν ἡ πόλις ἡ μεγάλη ἡ ἔχουσα βασιλείαν ἐπὶ τῶν βασιλέων τῆς γῆς.

\ch{18}
Μετὰ ταῦτα εἶδον ἄλλον ἄγγελον καταβαίνοντα ἐκ τοῦ οὐρανοῦ ἔχοντα ἐξουσίαν μεγάλην, καὶ ἡ γῆ ἐφωτίσθη ἐκ τῆς δόξης αὐτοῦ.
\vs{2}καὶ ἔκραξεν ἐν ἰσχυρᾷ φωνῇ λέγων·

Ἔπεσεν ἔπεσεν Βαβυλὼν ἡ μεγάλη, καὶ ἐγένετο κατοικητήριον δαιμονίων καὶ φυλακὴ παντὸς πνεύματος ἀκαθάρτου καὶ φυλακὴ παντὸς ὀρνέου ἀκαθάρτου καὶ φυλακὴ παντὸς θηρίου ἀκαθάρτου καὶ μεμισημένου,
\vs{3}ὅτι ἐκ τοῦ οἴνου τοῦ θυμοῦ τῆς πορνείας αὐτῆς πέπωκαν πάντα τὰ ἔθνη καὶ οἱ βασιλεῖς τῆς γῆς μετ᾽ αὐτῆς ἐπόρνευσαν καὶ οἱ ἔμποροι τῆς γῆς ἐκ τῆς δυνάμεως τοῦ στρήνους αὐτῆς ἐπλούτησαν.

\vs{4}Καὶ ἤκουσα ἄλλην φωνὴν ἐκ τοῦ οὐρανοῦ λέγουσαν·

Ἐξέλθατε ὁ λαός μου ἐξ αὐτῆς ἵνα μὴ συνκοινωνήσητε ταῖς ἁμαρτίαις αὐτῆς, καὶ ἐκ τῶν πληγῶν αὐτῆς ἵνα μὴ λάβητε,
\vs{5}ὅτι ἐκολλήθησαν αὐτῆς αἱ ἁμαρτίαι ἄχρι τοῦ οὐρανοῦ καὶ ἐμνημόνευσεν ὁ Θεὸς τὰ ἀδικήματα αὐτῆς.
\vs{6}ἀπόδοτε αὐτῇ ὡς καὶ αὐτὴ ἀπέδωκεν καὶ διπλώσατε τὰ διπλᾶ κατὰ τὰ ἔργα αὐτῆς, ἐν τῷ ποτηρίῳ ᾧ ἐκέρασεν κεράσατε αὐτῇ διπλοῦν,
\vs{7}ὅσα ἐδόξασεν αὑτὴν καὶ ἐστρηνίασεν, τοσοῦτον δότε αὐτῇ βασανισμὸν καὶ πένθος. ὅτι ἐν τῇ καρδίᾳ αὐτῆς λέγει ὅτι Κάθημαι βασίλισσα καὶ χήρα οὐκ εἰμί καὶ πένθος οὐ μὴ ἴδω.
\vs{8}διὰ τοῦτο ἐν μιᾷ ἡμέρᾳ ἥξουσιν αἱ πληγαὶ αὐτῆς, θάνατος καὶ πένθος καὶ λιμός, καὶ ἐν πυρὶ κατακαυθήσεται, ὅτι ἰσχυρὸς Κύριος ὁ Θεὸς ὁ κρίνας αὐτήν.

\vs{9}Καὶ κλαύσουσιν καὶ κόψονται ἐπ᾽ αὐτὴν οἱ βασιλεῖς τῆς γῆς οἱ μετ᾽ αὐτῆς πορνεύσαντες καὶ στρηνιάσαντες, ὅταν βλέπωσιν τὸν καπνὸν τῆς πυρώσεως αὐτῆς,
\vs{10}ἀπὸ μακρόθεν ἑστηκότες διὰ τὸν φόβον τοῦ βασανισμοῦ αὐτῆς λέγοντες· 
\begin{poetryblock}

\begin{quote}Οὐαὶ οὐαί, ἡ πόλις ἡ μεγάλη,\end{quote} 

\begin{quote}Βαβυλὼν ἡ πόλις ἡ ἰσχυρά,\end{quote} 

\begin{quote}ὅτι μιᾷ ὥρᾳ ἦλθεν ἡ κρίσις σου.\end{quote}
\end{poetryblock}

\vs{11}Καὶ οἱ ἔμποροι τῆς γῆς κλαίουσιν καὶ πενθοῦσιν ἐπ᾽ αὐτήν, ὅτι τὸν γόμον αὐτῶν οὐδεὶς ἀγοράζει οὐκέτι
\vs{12}γόμον χρυσοῦ καὶ ἀργύρου καὶ λίθου τιμίου καὶ μαργαριτῶν καὶ βυσσίνου καὶ πορφύρας καὶ σιρικοῦ καὶ κοκκίνου, καὶ πᾶν ξύλον θύϊνον καὶ πᾶν σκεῦος ἐλεφάντινον καὶ πᾶν σκεῦος ἐκ ξύλου τιμιωτάτου καὶ χαλκοῦ καὶ σιδήρου καὶ μαρμάρου,
\vs{13}καὶ κιννάμωμον καὶ ἄμωμον καὶ θυμιάματα καὶ μύρον καὶ λίβανον καὶ οἶνον καὶ ἔλαιον καὶ σεμίδαλιν καὶ σῖτον καὶ κτήνη καὶ πρόβατα, καὶ ἵππων καὶ ῥεδῶν καὶ σωμάτων, καὶ ψυχὰς ἀνθρώπων.
\begin{poetryblock}

\begin{quote} \vs{14}Καὶ ἡ ὀπώρα σου τῆς ἐπιθυμίας τῆς ψυχῆς ἀπῆλθεν ἀπὸ σοῦ,\end{quote} 

\begin{quote}καὶ πάντα τὰ λιπαρὰ καὶ τὰ λαμπρὰ ἀπώλετο ἀπὸ σοῦ\end{quote} 

\begin{quote}καὶ οὐκέτι οὐ μὴ αὐτὰ εὑρήσουσιν.\end{quote}
\end{poetryblock}

\vs{15}Οἱ ἔμποροι τούτων οἱ πλουτήσαντες ἀπ᾽ αὐτῆς ἀπὸ μακρόθεν στήσονται διὰ τὸν φόβον τοῦ βασανισμοῦ αὐτῆς κλαίοντες καὶ πενθοῦντες
\vs{16}λέγοντες· 
\begin{poetryblock}

\begin{quote}Οὐαὶ οὐαί, ἡ πόλις ἡ μεγάλη,\end{quote} 

\begin{quote}ἡ περιβεβλημένη βύσσινον καὶ πορφυροῦν καὶ κόκκινον\end{quote} 

\begin{quote}καὶ κεχρυσωμένη ἐν χρυσίῳ καὶ λίθῳ τιμίῳ καὶ μαργαρίτῃ,\end{quote}
\end{poetryblock}

\vs{17}ὅτι μιᾷ ὥρᾳ ἠρημώθη ὁ τοσοῦτος πλοῦτος.

Καὶ πᾶς κυβερνήτης καὶ πᾶς ὁ ἐπὶ τόπον πλέων καὶ ναῦται καὶ ὅσοι τὴν θάλασσαν ἐργάζονται, ἀπὸ μακρόθεν ἔστησαν
\vs{18}καὶ ἔκραζον βλέποντες τὸν καπνὸν τῆς πυρώσεως αὐτῆς λέγοντες· Τίς ὁμοία τῇ πόλει τῇ μεγάλῃ;
\vs{19}Καὶ ἔβαλον χοῦν ἐπὶ τὰς κεφαλὰς αὐτῶν καὶ ἔκραζον κλαίοντες καὶ πενθοῦντες λέγοντες· 
\begin{poetryblock}

\begin{quote}Οὐαὶ οὐαί, ἡ πόλις ἡ μεγάλη,\end{quote} 

\begin{quote}ἐν ᾗ ἐπλούτησαν πάντες οἱ ἔχοντες τὰ πλοῖα ἐν τῇ θαλάσσῃ ἐκ τῆς τιμιότητος αὐτῆς,\end{quote} 

\begin{quote}ὅτι μιᾷ ὥρᾳ ἠρημώθη.\end{quote}

\begin{quote} \vs{20}Εὐφραίνου ἐπ᾽ αὐτῇ, οὐρανέ\end{quote} 

\begin{quote}καὶ οἱ ἅγιοι καὶ οἱ ἀπόστολοι καὶ οἱ προφῆται,\end{quote} 

\begin{quote}ὅτι ἔκρινεν ὁ Θεὸς τὸ κρίμα ὑμῶν ἐξ αὐτῆς.\end{quote}
\end{poetryblock}

\vs{21}Καὶ ἦρεν εἷς ἄγγελος ἰσχυρὸς λίθον ὡς μύλινον μέγαν καὶ ἔβαλεν εἰς τὴν θάλασσαν λέγων· 
\begin{poetryblock}

\begin{quote}Οὕτως ὁρμήματι βληθήσεται Βαβυλὼν ἡ μεγάλη πόλις\end{quote} 

\begin{quote}καὶ οὐ μὴ εὑρεθῇ ἔτι.\end{quote}

\begin{quote} \vs{22}καὶ φωνὴ κιθαρῳδῶν καὶ μουσικῶν καὶ αὐλητῶν καὶ σαλπιστῶν\end{quote} 

\begin{quote}οὐ μὴ ἀκουσθῇ ἐν σοὶ ἔτι,\end{quote} 

\begin{quote}καὶ πᾶς τεχνίτης πάσης τέχνης\end{quote} 

\begin{quote}οὐ μὴ εὑρεθῇ ἐν σοὶ ἔτι,\end{quote} 

\begin{quote}καὶ φωνὴ μύλου\end{quote} 

\begin{quote}οὐ μὴ ἀκουσθῇ ἐν σοὶ ἔτι,\end{quote}

\begin{quote} \vs{23}καὶ φῶς λύχνου\end{quote} 

\begin{quote}οὐ μὴ φάνῃ ἐν σοὶ ἔτι,\end{quote} 

\begin{quote}καὶ φωνὴ νυμφίου καὶ νύμφης\end{quote} 

\begin{quote}οὐ μὴ ἀκουσθῇ ἐν σοὶ ἔτι·\end{quote} 

\begin{quote}ὅτι οἱ ἔμποροί σου ἦσαν οἱ μεγιστᾶνες τῆς γῆς,\end{quote} 

\begin{quote}ὅτι ἐν τῇ φαρμακείᾳ σου ἐπλανήθησαν πάντα τὰ ἔθνη,\end{quote}

\begin{quote} \vs{24}Καὶ ἐν αὐτῇ αἷμα προφητῶν καὶ ἁγίων εὑρέθη\end{quote} 

\begin{quote}καὶ πάντων τῶν ἐσφαγμένων ἐπὶ τῆς γῆς.\end{quote}
\end{poetryblock}

\ch{19}
Μετὰ ταῦτα ἤκουσα ὡς φωνὴν μεγάλην ὄχλου πολλοῦ ἐν τῷ οὐρανῷ λεγόντων· 
\begin{poetryblock}

\begin{quote}Ἁλληλουϊά·\end{quote} 

\begin{quote}ἡ σωτηρία καὶ ἡ δόξα καὶ ἡ δύναμις τοῦ Θεοῦ ἡμῶν,\end{quote}

\begin{quote} \vs{2}ὅτι ἀληθιναὶ καὶ δίκαιαι αἱ κρίσεις αὐτοῦ·\end{quote} 

\begin{quote}ὅτι ἔκρινεν τὴν πόρνην τὴν μεγάλην\end{quote} 

\begin{quote}ἥτις ἔφθειρεν τὴν γῆν ἐν τῇ πορνείᾳ αὐτῆς,\end{quote} 

\begin{quote}καὶ ἐξεδίκησεν τὸ αἷμα τῶν δούλων αὐτοῦ ἐκ χειρὸς αὐτῆς.\end{quote}
\end{poetryblock}

\vs{3}Καὶ δεύτερον εἴρηκαν· 
\begin{poetryblock}

\begin{quote}Ἁλληλουϊά·\end{quote} 

\begin{quote}καὶ ὁ καπνὸς αὐτῆς ἀναβαίνει εἰς τοὺς αἰῶνας τῶν αἰώνων.\end{quote}
\end{poetryblock}

\vs{4}Καὶ ἔπεσαν οἱ πρεσβύτεροι οἱ εἴκοσι τέσσαρες καὶ τὰ τέσσαρα ζῷα καὶ προσεκύνησαν τῷ Θεῷ τῷ καθημένῳ ἐπὶ τῷ θρόνῳ λέγοντες· 
\begin{poetryblock}

\begin{quote}Ἀμήν Ἁλληλουϊά.\end{quote}
\end{poetryblock}

\vs{5}Καὶ φωνὴ ἀπὸ τοῦ θρόνου ἐξῆλθεν λέγουσα· 
\begin{poetryblock}

\begin{quote}Αἰνεῖτε τῷ Θεῷ ἡμῶν\end{quote} 

\begin{quote}πάντες οἱ δοῦλοι αὐτοῦ\end{quote} 

\begin{quote}καὶ οἱ φοβούμενοι αὐτόν,\end{quote} 

\begin{quote}οἱ μικροὶ καὶ οἱ μεγάλοι.\end{quote}
\end{poetryblock}

\vs{6}Καὶ ἤκουσα ὡς φωνὴν ὄχλου πολλοῦ καὶ ὡς φωνὴν ὑδάτων πολλῶν καὶ ὡς φωνὴν βροντῶν ἰσχυρῶν λεγόντων· 
\begin{poetryblock}

\begin{quote}Ἁλληλουϊά,\end{quote} 

\begin{quote}ὅτι ἐβασίλευσεν Κύριος ὁ Θεός ἡμῶν ὁ Παντοκράτωρ.\end{quote}

\begin{quote} \vs{7}χαίρωμεν καὶ ἀγαλλιῶμεν καὶ δώσομεν τὴν δόξαν αὐτῷ,\end{quote} 

\begin{quote}ὅτι ἦλθεν ὁ γάμος τοῦ Ἀρνίου καὶ ἡ γυνὴ αὐτοῦ ἡτοίμασεν ἑαυτήν\end{quote}
\end{poetryblock}

\vs{8}καὶ ἐδόθη αὐτῇ ἵνα περιβάληται βύσσινον λαμπρὸν καθαρόν·

Τὸ γὰρ βύσσινον τὰ δικαιώματα τῶν ἁγίων ἐστίν.

\vs{9}Καὶ λέγει μοι· Γράψον· Μακάριοι οἱ εἰς τὸ δεῖπνον τοῦ γάμου τοῦ Ἀρνίου κεκλημένοι. καὶ λέγει μοι· Οὗτοι οἱ λόγοι ἀληθινοὶ τοῦ Θεοῦ εἰσιν.
\vs{10}Καὶ ἔπεσα ἔμπροσθεν τῶν ποδῶν αὐτοῦ προσκυνῆσαι αὐτῷ. καὶ λέγει μοι· Ὅρα μή· σύνδουλός σού εἰμι καὶ τῶν ἀδελφῶν σου τῶν ἐχόντων τὴν μαρτυρίαν Ἰησοῦ· τῷ Θεῷ προσκύνησον. ἡ γὰρ μαρτυρία Ἰησοῦ ἐστιν τὸ πνεῦμα τῆς προφητείας.

\vs{11}Καὶ εἶδον τὸν οὐρανὸν ἠνεῳγμένον, καὶ ἰδοὺ ἵππος λευκός καὶ ὁ καθήμενος ἐπ᾽ αὐτὸν καλούμενος Πιστὸς καὶ Ἀληθινός, καὶ ἐν δικαιοσύνῃ κρίνει καὶ πολεμεῖ.
\vs{12}οἱ δὲ ὀφθαλμοὶ αὐτοῦ ὡς φλὸξ πυρός, καὶ ἐπὶ τὴν κεφαλὴν αὐτοῦ διαδήματα πολλά, ἔχων ὄνομα γεγραμμένον ὃ οὐδεὶς οἶδεν εἰ μὴ αὐτός,
\vs{13}καὶ περιβεβλημένος ἱμάτιον βεβαμμένον αἵματι, καὶ κέκληται τὸ ὄνομα αὐτοῦ Ὁ Λόγος τοῦ Θεοῦ.

\vs{14}Καὶ τὰ στρατεύματα τὰ ἐν τῷ οὐρανῷ ἠκολούθει αὐτῷ ἐφ᾽ ἵπποις λευκοῖς, ἐνδεδυμένοι βύσσινον λευκὸν καθαρόν.
\vs{15}καὶ ἐκ τοῦ στόματος αὐτοῦ ἐκπορεύεται ῥομφαία ὀξεῖα, ἵνα ἐν αὐτῇ πατάξῃ τὰ ἔθνη, καὶ αὐτὸς ποιμανεῖ αὐτοὺς ἐν ῥάβδῳ σιδηρᾷ, καὶ αὐτὸς πατεῖ τὴν ληνὸν τοῦ οἴνου τοῦ θυμοῦ τῆς ὀργῆς τοῦ Θεοῦ τοῦ Παντοκράτορος,
\vs{16}καὶ ἔχει ἐπὶ τὸ ἱμάτιον καὶ ἐπὶ τὸν μηρὸν αὐτοῦ ὄνομα γεγραμμένον· ΒΑΣΙΛΕΥΣ ΒΑΣΙΛΕΩΝ ΚΑΙ ΚΥΡΙΟΣ ΚΥΡΙΩΝ.

\vs{17}Καὶ εἶδον ἕνα ἄγγελον ἑστῶτα ἐν τῷ ἡλίῳ καὶ ἔκραξεν ἐν φωνῇ μεγάλῃ λέγων πᾶσιν τοῖς ὀρνέοις τοῖς πετομένοις ἐν μεσουρανήματι·

Δεῦτε συνάχθητε εἰς τὸ δεῖπνον τὸ μέγα τοῦ Θεοῦ
\vs{18}ἵνα φάγητε σάρκας βασιλέων καὶ σάρκας χιλιάρχων καὶ σάρκας ἰσχυρῶν καὶ σάρκας ἵππων καὶ τῶν καθημένων ἐπ᾽ αὐτῶν καὶ σάρκας πάντων ἐλευθέρων τε καὶ δούλων καὶ μικρῶν καὶ μεγάλων.

\vs{19}Καὶ εἶδον τὸ θηρίον καὶ τοὺς βασιλεῖς τῆς γῆς καὶ τὰ στρατεύματα αὐτῶν συνηγμένα ποιῆσαι τὸν πόλεμον μετὰ τοῦ καθημένου ἐπὶ τοῦ ἵππου καὶ μετὰ τοῦ στρατεύματος αὐτοῦ.
\vs{20}καὶ ἐπιάσθη τὸ θηρίον καὶ μετ᾽ αὐτοῦ ὁ ψευδοπροφήτης ὁ ποιήσας τὰ σημεῖα ἐνώπιον αὐτοῦ, ἐν οἷς ἐπλάνησεν τοὺς λαβόντας τὸ χάραγμα τοῦ θηρίου καὶ τοὺς προσκυνοῦντας τῇ εἰκόνι αὐτοῦ· ζῶντες ἐβλήθησαν οἱ δύο εἰς τὴν λίμνην τοῦ πυρὸς τῆς καιομένης ἐν θείῳ.
\vs{21}καὶ οἱ λοιποὶ ἀπεκτάνθησαν ἐν τῇ ῥομφαίᾳ τοῦ καθημένου ἐπὶ τοῦ ἵππου τῇ ἐξελθούσῃ ἐκ τοῦ στόματος αὐτοῦ, Καὶ πάντα τὰ ὄρνεα ἐχορτάσθησαν ἐκ τῶν σαρκῶν αὐτῶν.

\ch{20}
Καὶ εἶδον ἄγγελον καταβαίνοντα ἐκ τοῦ οὐρανοῦ ἔχοντα τὴν κλεῖν τῆς ἀβύσσου καὶ ἅλυσιν μεγάλην ἐπὶ τὴν χεῖρα αὐτοῦ.
\vs{2}καὶ ἐκράτησεν τὸν δράκοντα, ὁ ὄφις ὁ ἀρχαῖος, ὅς ἐστιν Διάβολος καὶ Ὁ Σατανᾶς, καὶ ἔδησεν αὐτὸν χίλια ἔτη
\vs{3}καὶ ἔβαλεν αὐτὸν εἰς τὴν ἄβυσσον καὶ ἔκλεισεν καὶ ἐσφράγισεν ἐπάνω αὐτοῦ, ἵνα μὴ πλανήσῃ ἔτι τὰ ἔθνη ἄχρι τελεσθῇ τὰ χίλια ἔτη. μετὰ ταῦτα δεῖ λυθῆναι αὐτὸν μικρὸν χρόνον.

\vs{4}Καὶ εἶδον θρόνους καὶ ἐκάθισαν ἐπ᾽ αὐτούς καὶ κρίμα ἐδόθη αὐτοῖς, καὶ τὰς ψυχὰς τῶν πεπελεκισμένων διὰ τὴν μαρτυρίαν Ἰησοῦ καὶ διὰ τὸν λόγον τοῦ θεοῦ καὶ οἵτινες οὐ προσεκύνησαν τὸ θηρίον οὐδὲ τὴν εἰκόνα αὐτοῦ καὶ οὐκ ἔλαβον τὸ χάραγμα ἐπὶ τὸ μέτωπον καὶ ἐπὶ τὴν χεῖρα αὐτῶν. καὶ ἔζησαν καὶ ἐβασίλευσαν μετὰ τοῦ χριστοῦ χίλια ἔτη.
\vs{5}Οἱ λοιποὶ τῶν νεκρῶν οὐκ ἔζησαν ἄχρι τελεσθῇ τὰ χίλια ἔτη.

αὕτη ἡ ἀνάστασις ἡ πρώτη.
\vs{6}μακάριος καὶ ἅγιος ὁ ἔχων μέρος ἐν τῇ ἀναστάσει τῇ πρώτῃ· ἐπὶ τούτων ὁ δεύτερος θάνατος οὐκ ἔχει ἐξουσίαν, ἀλλ᾽ ἔσονται ἱερεῖς τοῦ Θεοῦ καὶ τοῦ Χριστοῦ καὶ βασιλεύσουσιν μετ᾽ αὐτοῦ τὰ χίλια ἔτη.

\vs{7}Καὶ ὅταν τελεσθῇ τὰ χίλια ἔτη, λυθήσεται ὁ Σατανᾶς ἐκ τῆς φυλακῆς αὐτοῦ
\vs{8}καὶ ἐξελεύσεται πλανῆσαι τὰ ἔθνη τὰ ἐν ταῖς τέσσαρσιν γωνίαις τῆς γῆς, τὸν Γὼγ καὶ Μαγώγ, συναγαγεῖν αὐτοὺς εἰς τὸν πόλεμον, ὧν ὁ ἀριθμὸς αὐτῶν ὡς ἡ ἄμμος τῆς θαλάσσης.
\vs{9}Καὶ ἀνέβησαν ἐπὶ τὸ πλάτος τῆς γῆς καὶ ἐκύκλευσαν τὴν παρεμβολὴν τῶν ἁγίων καὶ τὴν πόλιν τὴν ἠγαπημένην, καὶ κατέβη πῦρ ἐκ τοῦ οὐρανοῦ καὶ κατέφαγεν αὐτούς.
\vs{10}καὶ ὁ διάβολος ὁ πλανῶν αὐτοὺς ἐβλήθη εἰς τὴν λίμνην τοῦ πυρὸς καὶ θείου ὅπου καὶ τὸ θηρίον καὶ ὁ ψευδοπροφήτης, καὶ βασανισθήσονται ἡμέρας καὶ νυκτὸς εἰς τοὺς αἰῶνας τῶν αἰώνων.

\vs{11}Καὶ εἶδον θρόνον μέγαν λευκὸν καὶ τὸν καθήμενον ἐπ᾽ αὐτόν, οὗ ἀπὸ τοῦ προσώπου ἔφυγεν ἡ γῆ καὶ ὁ οὐρανός καὶ τόπος οὐχ εὑρέθη αὐτοῖς.
\vs{12}καὶ εἶδον τοὺς νεκρούς, τοὺς μεγάλους καὶ τοὺς μικρούς, ἑστῶτας ἐνώπιον τοῦ θρόνου. καὶ βιβλία ἠνοίχθησαν, Καὶ ἄλλο βιβλίον ἠνοίχθη, ὅ ἐστιν τῆς ζωῆς, καὶ ἐκρίθησαν οἱ νεκροὶ ἐκ τῶν γεγραμμένων ἐν τοῖς βιβλίοις κατὰ τὰ ἔργα αὐτῶν.
\vs{13}καὶ ἔδωκεν ἡ θάλασσα τοὺς νεκροὺς τοὺς ἐν αὐτῇ καὶ ὁ θάνατος καὶ ὁ ᾅδης ἔδωκαν τοὺς νεκροὺς τοὺς ἐν αὐτοῖς, καὶ ἐκρίθησαν ἕκαστος κατὰ τὰ ἔργα αὐτῶν.
\vs{14}Καὶ ὁ θάνατος καὶ ὁ ᾅδης ἐβλήθησαν εἰς τὴν λίμνην τοῦ πυρός. οὗτος ὁ θάνατος ὁ δεύτερός ἐστιν, ἡ λίμνη τοῦ πυρός.
\vs{15}καὶ εἴ τις οὐχ εὑρέθη ἐν τῇ βίβλῳ τῆς ζωῆς γεγραμμένος, ἐβλήθη εἰς τὴν λίμνην τοῦ πυρός.

\ch{21}
Καὶ εἶδον οὐρανὸν καινὸν καὶ γῆν καινήν. ὁ γὰρ πρῶτος οὐρανὸς καὶ ἡ πρώτη γῆ ἀπῆλθαν καὶ ἡ θάλασσα οὐκ ἔστιν ἔτι.
\vs{2}καὶ τὴν πόλιν τὴν ἁγίαν Ἰερουσαλὴμ καινὴν εἶδον καταβαίνουσαν ἐκ τοῦ οὐρανοῦ ἀπὸ τοῦ Θεοῦ ἡτοιμασμένην ὡς νύμφην κεκοσμημένην τῷ ἀνδρὶ αὐτῆς.
\vs{3}Καὶ ἤκουσα φωνῆς μεγάλης ἐκ τοῦ θρόνου λεγούσης·

Ἰδοὺ ἡ σκηνὴ τοῦ Θεοῦ μετὰ τῶν ἀνθρώπων, καὶ σκηνώσει μετ᾽ αὐτῶν, καὶ αὐτοὶ λαοὶ αὐτοῦ ἔσονται, καὶ αὐτὸς ὁ Θεὸς μετ᾽ αὐτῶν ἔσται αὐτῶν θεός,
\vs{4}καὶ ἐξαλείψει πᾶν δάκρυον ἐκ τῶν ὀφθαλμῶν αὐτῶν, καὶ ὁ θάνατος οὐκ ἔσται ἔτι οὔτε πένθος οὔτε κραυγὴ οὔτε πόνος οὐκ ἔσται ἔτι, ὅτι τὰ πρῶτα ἀπῆλθαν.

\vs{5}Καὶ εἶπεν ὁ καθήμενος ἐπὶ τῷ θρόνῳ· Ἰδοὺ καινὰ ποιῶ πάντα καὶ λέγει· Γράψον, ὅτι οὗτοι οἱ λόγοι πιστοὶ καὶ ἀληθινοί εἰσιν.
\vs{6}καὶ εἶπέν μοι· Γέγοναν. ἐγὼ εἰμι τὸ Ἄλφα καὶ τὸ Ὦ, ἡ ἀρχὴ καὶ τὸ τέλος. ἐγὼ τῷ διψῶντι δώσω ἐκ τῆς πηγῆς τοῦ ὕδατος τῆς ζωῆς δωρεάν.
\vs{7}ὁ νικῶν κληρονομήσει ταῦτα καὶ ἔσομαι αὐτῷ Θεὸς καὶ αὐτὸς ἔσται μοι υἱός.
\vs{8}Τοῖς δὲ δειλοῖς καὶ ἀπίστοις καὶ ἐβδελυγμένοις καὶ φονεῦσιν καὶ πόρνοις καὶ φαρμάκοις καὶ εἰδωλολάτραις καὶ πᾶσιν τοῖς ψευδέσιν τὸ μέρος αὐτῶν ἐν τῇ λίμνῃ τῇ καιομένῃ πυρὶ καὶ θείῳ, ὅ ἐστιν ὁ θάνατος ὁ δεύτερος.

\vs{9}Καὶ ἦλθεν εἷς ἐκ τῶν ἑπτὰ ἀγγέλων τῶν ἐχόντων τὰς ἑπτὰ φιάλας τῶν γεμόντων τῶν ἑπτὰ πληγῶν τῶν ἐσχάτων καὶ ἐλάλησεν μετ᾽ ἐμοῦ λέγων· Δεῦρο, δείξω σοι τὴν νύμφην τὴν γυναῖκα τοῦ ἀρνίου.
\vs{10}Καὶ ἀπήνεγκέν με ἐν Πνεύματι ἐπὶ ὄρος μέγα καὶ ὑψηλόν, καὶ ἔδειξέν μοι τὴν πόλιν τὴν ἁγίαν Ἰερουσαλὴμ καταβαίνουσαν ἐκ τοῦ οὐρανοῦ ἀπὸ τοῦ Θεοῦ
\vs{11}ἔχουσαν τὴν δόξαν τοῦ Θεοῦ, ὁ φωστὴρ αὐτῆς ὅμοιος λίθῳ τιμιωτάτῳ ὡς λίθῳ ἰάσπιδι κρυσταλλίζοντι.
\vs{12}ἔχουσα τεῖχος μέγα καὶ ὑψηλόν, ἔχουσα πυλῶνας δώδεκα καὶ ἐπὶ τοῖς πυλῶσιν ἀγγέλους δώδεκα καὶ ὀνόματα ἐπιγεγραμμένα, ἅ ἐστιν τὰ ὀνόματα τῶν δώδεκα φυλῶν υἱῶν Ἰσραήλ·
\vs{13}ἀπὸ ἀνατολῆς πυλῶνες τρεῖς καὶ ἀπὸ βορρᾶ πυλῶνες τρεῖς καὶ ἀπὸ νότου πυλῶνες τρεῖς καὶ ἀπὸ δυσμῶν πυλῶνες τρεῖς.
\vs{14}καὶ τὸ τεῖχος τῆς πόλεως ἔχων θεμελίους δώδεκα καὶ ἐπ᾽ αὐτῶν δώδεκα ὀνόματα τῶν δώδεκα ἀποστόλων τοῦ Ἀρνίου.

\vs{15}Καὶ ὁ λαλῶν μετ᾽ ἐμοῦ εἶχεν μέτρον κάλαμον χρυσοῦν, ἵνα μετρήσῃ τὴν πόλιν καὶ τοὺς πυλῶνας αὐτῆς καὶ τὸ τεῖχος αὐτῆς.
\vs{16}καὶ ἡ πόλις τετράγωνος κεῖται καὶ τὸ μῆκος αὐτῆς ὅσον καὶ τὸ πλάτος. καὶ ἐμέτρησεν τὴν πόλιν τῷ καλάμῳ ἐπὶ σταδίων δώδεκα χιλιάδων, τὸ μῆκος καὶ τὸ πλάτος καὶ τὸ ὕψος αὐτῆς ἴσα ἐστίν.
\vs{17}καὶ ἐμέτρησεν τὸ τεῖχος αὐτῆς ἑκατὸν τεσσεράκοντα τεσσάρων πηχῶν μέτρον ἀνθρώπου, ὅ ἐστιν ἀγγέλου.
\vs{18}Καὶ ἡ ἐνδώμησις τοῦ τείχους αὐτῆς ἴασπις καὶ ἡ πόλις χρυσίον καθαρὸν ὅμοιον ὑάλῳ καθαρῷ.
\vs{19}οἱ θεμέλιοι τοῦ τείχους τῆς πόλεως παντὶ λίθῳ τιμίῳ κεκοσμημένοι· ὁ θεμέλιος ὁ πρῶτος ἴασπις, ὁ δεύτερος σάπφιρος, ὁ τρίτος χαλκηδών, ὁ τέταρτος σμάραγδος,
\vs{20}ὁ πέμπτος σαρδόνυξ, ὁ ἕκτος σάρδιον, ὁ ἕβδομος χρυσόλιθος, ὁ ὄγδοος βήρυλλος, ὁ ἔνατος τοπάζιον, ὁ δέκατος χρυσόπρασος, ὁ ἑνδέκατος ὑάκινθος, ὁ δωδέκατος ἀμέθυστος,
\vs{21}Καὶ οἱ δώδεκα πυλῶνες δώδεκα μαργαρῖται, ἀνὰ εἷς ἕκαστος τῶν πυλώνων ἦν ἐξ ἑνὸς μαργαρίτου. καὶ ἡ πλατεῖα τῆς πόλεως χρυσίον καθαρὸν ὡς ὕαλος διαυγής.

\vs{22}Καὶ ναὸν οὐκ εἶδον ἐν αὐτῇ, ὁ γὰρ Κύριος ὁ Θεὸς ὁ Παντοκράτωρ ναὸς αὐτῆς ἐστιν καὶ τὸ Ἀρνίον.
\vs{23}καὶ ἡ πόλις οὐ χρείαν ἔχει τοῦ ἡλίου οὐδὲ τῆς σελήνης ἵνα φαίνωσιν αὐτῇ, ἡ γὰρ δόξα τοῦ Θεοῦ ἐφώτισεν αὐτήν, καὶ ὁ λύχνος αὐτῆς τὸ Ἀρνίον.
\vs{24}καὶ περιπατήσουσιν τὰ ἔθνη διὰ τοῦ φωτὸς αὐτῆς, καὶ οἱ βασιλεῖς τῆς γῆς φέρουσιν τὴν δόξαν αὐτῶν εἰς αὐτήν,
\vs{25}καὶ οἱ πυλῶνες αὐτῆς οὐ μὴ κλεισθῶσιν ἡμέρας, νὺξ γὰρ οὐκ ἔσται ἐκεῖ,
\vs{26}Καὶ οἴσουσιν τὴν δόξαν καὶ τὴν τιμὴν τῶν ἐθνῶν εἰς αὐτήν.
\vs{27}καὶ οὐ μὴ εἰσέλθῃ εἰς αὐτὴν πᾶν κοινὸν καὶ ὁ ποιῶν βδέλυγμα καὶ ψεῦδος εἰ μὴ οἱ γεγραμμένοι ἐν τῷ βιβλίῳ τῆς ζωῆς τοῦ Ἀρνίου.

\ch{22}
Καὶ ἔδειξέν μοι ποταμὸν ὕδατος ζωῆς λαμπρὸν ὡς κρύσταλλον, ἐκπορευόμενον ἐκ τοῦ θρόνου τοῦ Θεοῦ καὶ τοῦ Ἀρνίου.
\vs{2}ἐν μέσῳ τῆς πλατείας αὐτῆς καὶ τοῦ ποταμοῦ ἐντεῦθεν καὶ ἐκεῖθεν ξύλον ζωῆς ποιοῦν καρποὺς δώδεκα, κατὰ μῆνα ἕκαστον ἀποδιδοῦν τὸν καρπὸν αὐτοῦ, καὶ τὰ φύλλα τοῦ ξύλου εἰς θεραπείαν τῶν ἐθνῶν.
\vs{3}Καὶ πᾶν κατάθεμα οὐκ ἔσται ἔτι. καὶ ὁ θρόνος τοῦ Θεοῦ καὶ τοῦ Ἀρνίου ἐν αὐτῇ ἔσται, καὶ οἱ δοῦλοι αὐτοῦ λατρεύσουσιν αὐτῷ
\vs{4}καὶ ὄψονται τὸ πρόσωπον αὐτοῦ, καὶ τὸ ὄνομα αὐτοῦ ἐπὶ τῶν μετώπων αὐτῶν.
\vs{5}καὶ νὺξ οὐκ ἔσται ἔτι καὶ οὐκ ἔχουσιν χρείαν φωτὸς λύχνου καὶ φωτὸς ἡλίου, ὅτι Κύριος ὁ Θεὸς φωτίσει ἐπ᾽ αὐτούς, καὶ βασιλεύσουσιν εἰς τοὺς αἰῶνας τῶν αἰώνων.

\vs{6}Καὶ εἶπέν μοι· Οὗτοι οἱ λόγοι πιστοὶ καὶ ἀληθινοί, καὶ ὁ Κύριος ὁ Θεὸς τῶν πνευμάτων τῶν προφητῶν ἀπέστειλεν τὸν ἄγγελον αὐτοῦ δεῖξαι τοῖς δούλοις αὐτοῦ ἃ δεῖ γενέσθαι ἐν τάχει.
\vs{7}Καὶ Ἰδοὺ ἔρχομαι ταχύ. μακάριος ὁ τηρῶν τοὺς λόγους τῆς προφητείας τοῦ βιβλίου τούτου.

\vs{8}Κἀγὼ Ἰωάννης ὁ ἀκούων καὶ βλέπων ταῦτα. καὶ ὅτε ἤκουσα καὶ ἔβλεψα, ἔπεσα προσκυνῆσαι ἔμπροσθεν τῶν ποδῶν τοῦ ἀγγέλου τοῦ δεικνύοντός μοι ταῦτα.
\vs{9}καὶ λέγει μοι· Ὅρα μή· σύνδουλός σού εἰμι καὶ τῶν ἀδελφῶν σου τῶν προφητῶν καὶ τῶν τηρούντων τοὺς λόγους τοῦ βιβλίου τούτου· τῷ Θεῷ προσκύνησον.

\vs{10}Καὶ λέγει μοι· Μὴ σφραγίσῃς τοὺς λόγους τῆς προφητείας τοῦ βιβλίου τούτου, ὁ καιρὸς γὰρ ἐγγύς ἐστιν.
\vs{11}ὁ ἀδικῶν ἀδικησάτω ἔτι καὶ ὁ ῥυπαρὸς ῥυπανθήτω ἔτι, καὶ ὁ δίκαιος δικαιοσύνην ποιησάτω ἔτι καὶ ὁ ἅγιος ἁγιασθήτω ἔτι.

\vs{12}Ἰδοὺ ἔρχομαι ταχύ, καὶ ὁ μισθός μου μετ᾽ ἐμοῦ ἀποδοῦναι ἑκάστῳ ὡς τὸ ἔργον ἐστὶν αὐτοῦ.
\vs{13}ἐγὼ τὸ Ἄλφα καὶ τὸ Ὦ, ὁ πρῶτος καὶ ὁ ἔσχατος, ἡ ἀρχὴ καὶ τὸ τέλος.

\vs{14}Μακάριοι οἱ πλύνοντες τὰς στολὰς αὐτῶν, ἵνα ἔσται ἡ ἐξουσία αὐτῶν ἐπὶ τὸ ξύλον τῆς ζωῆς καὶ τοῖς πυλῶσιν εἰσέλθωσιν εἰς τὴν πόλιν.
\vs{15}ἔξω οἱ κύνες καὶ οἱ φάρμακοι καὶ οἱ πόρνοι καὶ οἱ φονεῖς καὶ οἱ εἰδωλολάτραι καὶ πᾶς φιλῶν καὶ ποιῶν ψεῦδος.

\vs{16}Ἐγὼ Ἰησοῦς ἔπεμψα τὸν ἄγγελόν μου μαρτυρῆσαι ὑμῖν ταῦτα ἐπὶ ταῖς ἐκκλησίαις. ἐγώ εἰμι ἡ ῥίζα καὶ τὸ γένος Δαυίδ, ὁ ἀστὴρ ὁ λαμπρός ὁ πρωϊνός.

\vs{17}Καὶ τὸ Πνεῦμα καὶ ἡ νύμφη λέγουσιν· Ἔρχου. καὶ ὁ ἀκούων εἰπάτω· Ἔρχου. καὶ ὁ διψῶν ἐρχέσθω, ὁ θέλων λαβέτω ὕδωρ ζωῆς δωρεάν.

\vs{18}Μαρτυρῶ ἐγὼ παντὶ τῷ ἀκούοντι τοὺς λόγους τῆς προφητείας τοῦ βιβλίου τούτου· ἐάν τις ἐπιθῇ ἐπ᾽ αὐτά, ἐπιθήσει ὁ Θεὸς ἐπ᾽ αὐτὸν τὰς πληγὰς τὰς γεγραμμένας ἐν τῷ βιβλίῳ τούτῳ,
\vs{19}καὶ ἐάν τις ἀφέλῃ ἀπὸ τῶν λόγων τοῦ βιβλίου τῆς προφητείας ταύτης, ἀφελεῖ ὁ Θεὸς τὸ μέρος αὐτοῦ ἀπὸ τοῦ ξύλου τῆς ζωῆς καὶ ἐκ τῆς πόλεως τῆς ἁγίας τῶν γεγραμμένων ἐν τῷ βιβλίῳ τούτῳ.

\vs{20}Λέγει ὁ μαρτυρῶν ταῦτα· Ναί, ἔρχομαι ταχύ. Ἀμήν, ἔρχου Κύριε Ἰησοῦ.
\vs{21}Ἡ χάρις τοῦ Κυρίου Ἰησοῦ μετὰ πάντων.


\end{spacing}
\end{document}